% -*- ess-noweb-default-code-mode: f90-mode; noweb-default-code-mode: f90-mode; -*- 
% WHIZARD code as NOWEB source: Header, intro and generic stuff

\documentclass[a4paper]{report} 

\usepackage{amsmath}

\usepackage
  [bookmarks,bookmarksopen=true,bookmarksopenlevel=1,bookmarksnumbered=true]
  {hyperref} 
  
\usepackage{noweb} 
\usepackage{graphics,graphicx}
\usepackage{url}
\setlength{\nwmarginglue}{1em}
\noweboptions{smallcode,noidentxref}
%%% Saving paper:
\def\nwendcode{\endtrivlist\endgroup}
\nwcodepenalty=0
\let\nwdocspar\relax

%\makeindex
 
%%%%%%%%%%%%%%%%%%%%%%%%%%%%%%%%%%%%%%%%%%%%%%%%%%%%%%%%%%%%%%
%%% Macros

\def\tsum{{\textstyle\sum}}

\newcommand{\circeone}{\texttt{CIRCE1}}
\newcommand{\circetwo}{\texttt{CIRCE2}}
\newcommand{\whizard}{\texttt{WHIZARD}}
% Noweb emacs mode: single ' below
\newcommand{\oMega}{\texttt{O'MEGA}}
\newcommand{\vamp}{\texttt{VAMP}}
\newcommand{\pythia}{\texttt{PYTHIA}}
\newcommand{\gosam}{\texttt{GoSam}}

%%%%%%%%%%%%%%%%%%%%%%%%%%%%%%%%%%%%%%%%%%%%%%%%%%%%%%%%%%%%%%
\begin{document}
%%%%%%%%%%%%%%%%%%%%%%%%%%%%%%%%%%%%%%%%%%%%%%%%%%%%%%%%%%%%%%
\def\WhizardVersion{2.2.5}
\def\WhizardDate{Feb 27 2015}
\nwfilename{../../../src/noweb-frame/whizard-prelude.nw}\nwbegincode{1}\moddef{Version}\endmoddef
2.2.5
\nwendcode{}\nwbegincode{2}\moddef{Date}\endmoddef
Feb 27 2015
\nwendcode{}\nwbegindocs{3}\nwdocspar
%%%%%%%%%%%%%%%%%%%%%%%%%%%%%%%%%%%%%%%%%%%%%%%%%%%%%%%%%%%%
\title{%
  \whizard\footnote{The original meaning of the acronym is \emph{$W$,
  Higgs, $Z$, And Respective Decays}.  The current program is much more
  than that, however.}
}
\author{%
  Wolfgang Kilian,%
    \thanks{e-mail: \texttt{kilian@physik.uni-siegen.de}}
  Thorsten Ohl,%
    \thanks{e-mail: \texttt{ohl@physik.uni-wuerzburg.de}}
  J\"urgen Reuter%
    \thanks{e-mail: \texttt{juergen.reuter@desy.de}}}
\date{Version \WhizardVersion, \WhizardDate \\
 \mbox{}
 with contributions from:
 Fabian Bach, Felix Braam, Hans-Werner Boschmann, 
 Bijan Chokouf\'{e} Nejad, David Gordo Gomez, Sebastian Schmidt,  
 Marco Sekulla, Christian Speckner, Christian Weiss, Daniel Wiesler 
 \vspace{1cm}
 \begin{center}
 \includegraphics[width=4cm]{Whizard-Logo}
 \end{center}
 \mbox{} \\
 \vspace{.2cm}} 
\maketitle


\begin{abstract}
\texttt{WHIZARD} is an application of the \texttt{VAMP} algorithm:
Adaptive multi-channel integration and event generation. The bare
\texttt{VAMP} library is augmented by modules for Lorentz algebra,
particles, phase space, etc., such that physical processes with
arbitrary complex final states [well, in principle\ldots] can be
integrated and \emph{unweighted} events be generated.  
\end{abstract}


\newpage

\tableofcontents

\newpage 

%%%%%%%%%%%%%%%%%%%%%%%%%%%%%%%%%%%%%%%%%%%%%%%%%%%%%%%%%%%%
\chapter{Changes}

  {\bf For a comprehensive list of changes confer the
    ChangeLog file or the subversion log.}

\chapter{Preliminaries}
The WHIZARD file header:
\nwenddocs{}\nwbegincode{4}\moddef{File header}\endmoddef
! WHIZARD \LA{}Version\RA{} \LA{}Date\RA{}
! 
! Copyright (C) 1999-2015 by 
!     Wolfgang Kilian <kilian@physik.uni-siegen.de>
!     Thorsten Ohl <ohl@physik.uni-wuerzburg.de>
!     Juergen Reuter <juergen.reuter@desy.de>
!     
!     with contributions from
!     Fabian Bach <fabian.bach@desy.de>
!     Christian Speckner <cnspeckn@googlemail.com> 
!     Christian Weiss <christian.weiss@desy.de>
!     and Hans-Werner Boschmann, Felix Braam, 
!     Sebastian Schmidt, Daniel Wiesler 
!
! WHIZARD is free software; you can redistribute it and/or modify it
! under the terms of the GNU General Public License as published by 
! the Free Software Foundation; either version 2, or (at your option)
! any later version.
!
! WHIZARD is distributed in the hope that it will be useful, but
! WITHOUT ANY WARRANTY; without even the implied warranty of
! MERCHANTABILITY or FITNESS FOR A PARTICULAR PURPOSE.  See the 
! GNU General Public License for more details.
!
! You should have received a copy of the GNU General Public License
! along with this program; if not, write to the Free Software
! Foundation, Inc., 675 Mass Ave, Cambridge, MA 02139, USA.
!
!!!!!!!!!!!!!!!!!!!!!!!!!!!!!!!!!!!!!!!!!!!!!!!!!!!!!!!!!!!!!!!!!!!!!!!!
! This file has been stripped of most comments.  For documentation, refer
! to the source 'whizard.nw'
\nwendcode{}\nwbegindocs{5}\nwdocspar
We are strict with our names:
\nwenddocs{}\nwbegincode{6}\moddef{Standard module head}\endmoddef
  implicit none
  private
\nwendcode{}\nwbegindocs{7}This is the way to envoke the kinds module (not contained in this source)
\nwenddocs{}\nwbegincode{8}\moddef{Use kinds}\endmoddef
  use kinds, only: default
\nwendcode{}\nwbegincode{9}\moddef{Use kinds with double}\endmoddef
  use kinds, only: default, double
\eatline
\nwendcode{}\nwbegindocs{10}And we make heavy use of variable-length strings
\nwenddocs{}\nwbegincode{11}\moddef{Use strings}\endmoddef
  use iso_varying_string, string_t => varying_string
\eatline
\nwendcode{}\nwfilename{../../../src/utilities/utilities.nw}\nwbegindocs{0}% -*- ess-noweb-default-code-mode: f90-mode; noweb-default-code-mode: f90-mode; -*- 
% WHIZARD code as NOWEB source: Utilities

\nwenddocs{}\nwbegindocs{1}\nwdocspar
\chapter{Utilities}

These modules are intended as part of WHIZARD, but in fact they are
generic and could be useful for any purpose.

The modules depend only on modules from the {\tt{}basics} set.
\begin{description}
\item[file\_utils]
  Procedures that deal with external files, if not covered by Fortran
  built-ins.
\item[file\_registries]
  Manage files that are accessed by their name.
\item[string\_utils]
  Some string-handling utilities.  Includes conversion to C string.
\item[format\_utils]
  Utilities for pretty-printing.
\item[format\_defs]
  Predefined format strings.
\end{description}
%%%%%%%%%%%%%%%%%%%%%%%%%%%%%%%%%%%%%%%%%%%%%%%%%%%%%%%%%%%%%%%%%%%%%%%%
\section{File Utilities}

This module provides miscellaneous tools associated with named
external files.  Currently only:
\begin{itemize}
\item
  Delete a named file
\end{itemize}
\nwenddocs{}\nwbegincode{2}\moddef{\code{}file{\char95}utils.f90\edoc{}}\endmoddef
\LA{}File header\RA{}

module file_utils

  use io_units

\LA{}Standard module head\RA{}

\LA{}File utils: public\RA{}

contains

\LA{}File utils: procedures\RA{}

end module file_utils
\eatline
\nwendcode{}\nwbegindocs{3}\nwdocspar
\subsection{Deleting a file}
Fortran does not contain a command for deleting a file.  Here, we
provide a subroutine that deletes a file if it exists.  We do not
handle the subtleties, so we assume that it is writable if it exists.
\nwenddocs{}\nwbegincode{4}\moddef{File utils: public}\endmoddef
  public :: delete_file
\nwendcode{}\nwbegincode{5}\moddef{File utils: procedures}\endmoddef
  subroutine delete_file (name)
    character(*), intent(in) :: name
    logical :: exist
    integer :: u
    inquire (file = name, exist = exist)
    if (exist) then
       u = free_unit ()
       open (unit = u, file = name)
       close (u, status = "delete")
    end if
  end subroutine delete_file
  
\eatline
\nwendcode{}\nwbegindocs{6}\nwdocspar
%%%%%%%%%%%%%%%%%%%%%%%%%%%%%%%%%%%%%%%%%%%%%%%%%%%%%%%%%%%%%%%%%%%%%%%%
\section{File Registries}

This module provides a file-registry facility.  We can open and close
files multiple times without inadvertedly accessing a single file by two
different I/O unit numbers.  Opening a file the first time enters it
into the registry.  Opening again just returns the associated I/O
unit.  The registry maintains a reference count, so closing a file
does not actually complete until the last reference is released.

File access will always be sequential, however.  The file can't be
opened at different positions simultaneously.
\nwenddocs{}\nwbegincode{7}\moddef{\code{}file{\char95}registries.f90\edoc{}}\endmoddef
\LA{}File header\RA{}

module file_registries
  
\LA{}Use strings\RA{}
  use io_units

\LA{}Standard module head\RA{}

\LA{}File registries: public\RA{}

\LA{}File registries: types\RA{}

contains
  
\LA{}File registries: procedures\RA{}

end module file_registries
\eatline
\nwendcode{}\nwbegindocs{8}\nwdocspar
\subsection{File handle}
This object holds a filename (fully qualified), the associated
unit, and a reference count.  The idea is that the object should be
deleted when the reference count drops to zero.
\nwenddocs{}\nwbegincode{9}\moddef{File registries: types}\endmoddef
  type :: file_handle_t
     type(string_t) :: file
     integer :: unit = 0
     integer :: refcount = 0
   contains
   \LA{}File registries: file handle: TBP\RA{}
  end type file_handle_t
  
\eatline
\nwendcode{}\nwbegindocs{10}Debugging output:
\nwenddocs{}\nwbegincode{11}\moddef{File registries: file handle: TBP}\endmoddef
  procedure :: write => file_handle_write
\nwendcode{}\nwbegincode{12}\moddef{File registries: procedures}\endmoddef
  subroutine file_handle_write (handle, u, show_unit)
    class(file_handle_t), intent(in) :: handle
    integer, intent(in) :: u
    logical, intent(in), optional :: show_unit
    logical :: show_u
    show_u = .false.;  if (present (show_unit))  show_u = show_unit
    if (show_u) then
       write (u, "(3x,A,1x,I0,1x,'(',I0,')')")  &
            char (handle%file), handle%unit, handle%refcount
    else
       write (u, "(3x,A,1x,'(',I0,')')")  &
            char (handle%file), handle%refcount
    end if
  end subroutine file_handle_write
  
\eatline
\nwendcode{}\nwbegindocs{13}Initialize with a file name, don't open the file yet:
\nwenddocs{}\nwbegincode{14}\moddef{File registries: file handle: TBP}\plusendmoddef
  procedure :: init => file_handle_init
\nwendcode{}\nwbegincode{15}\moddef{File registries: procedures}\plusendmoddef
  subroutine file_handle_init (handle, file)
    class(file_handle_t), intent(out) :: handle
    type(string_t), intent(in) :: file
    handle%file = file
  end subroutine file_handle_init
  
\eatline
\nwendcode{}\nwbegindocs{16}We check the {\tt{}refcount} before actually opening the file.
\nwenddocs{}\nwbegincode{17}\moddef{File registries: file handle: TBP}\plusendmoddef
  procedure :: open => file_handle_open
\nwendcode{}\nwbegincode{18}\moddef{File registries: procedures}\plusendmoddef
  subroutine file_handle_open (handle)
    class(file_handle_t), intent(inout) :: handle
    if (handle%refcount == 0) then
       handle%unit = free_unit ()
       open (unit = handle%unit, file = char (handle%file), action = "read", &
            status = "old")
    end if
    handle%refcount = handle%refcount + 1
  end subroutine file_handle_open
    
\eatline
\nwendcode{}\nwbegindocs{19}Analogously, close if the refcount drops to zero.  The caller may
then delete the object.
\nwenddocs{}\nwbegincode{20}\moddef{File registries: file handle: TBP}\plusendmoddef
  procedure :: close => file_handle_close
\nwendcode{}\nwbegincode{21}\moddef{File registries: procedures}\plusendmoddef
  subroutine file_handle_close (handle)
    class(file_handle_t), intent(inout) :: handle
    handle%refcount = handle%refcount - 1
    if (handle%refcount == 0) then
       close (handle%unit)
       handle%unit = 0
    end if
  end subroutine file_handle_close
  
\eatline
\nwendcode{}\nwbegindocs{22}The I/O unit will be nonzero when the file is open.
\nwenddocs{}\nwbegincode{23}\moddef{File registries: file handle: TBP}\plusendmoddef
  procedure :: is_open => file_handle_is_open
\nwendcode{}\nwbegincode{24}\moddef{File registries: procedures}\plusendmoddef
  function file_handle_is_open (handle) result (flag)
    class(file_handle_t), intent(in) :: handle
    logical :: flag
    flag = handle%unit /= 0
  end function file_handle_is_open
  
\eatline
\nwendcode{}\nwbegindocs{25}Return the filename, so we can identify the entry.
\nwenddocs{}\nwbegincode{26}\moddef{File registries: file handle: TBP}\plusendmoddef
  procedure :: get_file => file_handle_get_file
\nwendcode{}\nwbegincode{27}\moddef{File registries: procedures}\plusendmoddef
  function file_handle_get_file (handle) result (file)
    class(file_handle_t), intent(in) :: handle
    type(string_t) :: file
    file = handle%file
  end function file_handle_get_file
  
\eatline
\nwendcode{}\nwbegindocs{28}For debugging, return the I/O unit number.
\nwenddocs{}\nwbegincode{29}\moddef{File registries: file handle: TBP}\plusendmoddef
  procedure :: get_unit => file_handle_get_unit
\nwendcode{}\nwbegincode{30}\moddef{File registries: procedures}\plusendmoddef
  function file_handle_get_unit (handle) result (unit)
    class(file_handle_t), intent(in) :: handle
    integer :: unit
    unit = handle%unit
  end function file_handle_get_unit
  
\eatline
\nwendcode{}\nwbegindocs{31}\nwdocspar
\subsection{File handles registry}
This is implemented as a doubly-linked list.  The list exists only
once in the program, as a private module variable.

Extend the handle type to become a list entry:
\nwenddocs{}\nwbegincode{32}\moddef{File registries: types}\plusendmoddef
  type, extends (file_handle_t) :: file_entry_t
     type(file_entry_t), pointer :: prev => null ()
     type(file_entry_t), pointer :: next => null ()
  end type file_entry_t

\eatline
\nwendcode{}\nwbegindocs{33}The actual registry.  We need only the pointer to the first entry.
\nwenddocs{}\nwbegincode{34}\moddef{File registries: public}\endmoddef
  public :: file_registry_t
\nwendcode{}\nwbegincode{35}\moddef{File registries: types}\plusendmoddef
  type :: file_registry_t
     type(file_entry_t), pointer :: first => null ()
   contains
   \LA{}File registries: file registry: TBP\RA{}
  end type file_registry_t
  
\eatline
\nwendcode{}\nwbegindocs{36}Debugging output.
\nwenddocs{}\nwbegincode{37}\moddef{File registries: file registry: TBP}\endmoddef
  procedure :: write => file_registry_write
\nwendcode{}\nwbegincode{38}\moddef{File registries: procedures}\plusendmoddef
  subroutine file_registry_write (registry, unit, show_unit)
    class(file_registry_t), intent(in) :: registry
    integer, intent(in), optional :: unit
    logical, intent(in), optional :: show_unit
    type(file_entry_t), pointer :: entry
    integer :: u
    u = given_output_unit (unit)
    if (associated (registry%first)) then
       write (u, "(1x,A)")  "File registry:"
       entry => registry%first
       do while (associated (entry))
          call entry%write (u, show_unit)
          entry => entry%next
       end do
    else
       write (u, "(1x,A)")  "File registry: [empty]"
    end if
  end subroutine file_registry_write
    
\eatline
\nwendcode{}\nwbegindocs{39}Open a file: find the appropriate entry.  Create a new entry and add
to the list if necessary.  The list is extended at the beginning.
Return the I/O unit number for the records.
\nwenddocs{}\nwbegincode{40}\moddef{File registries: file registry: TBP}\plusendmoddef
  procedure :: open => file_registry_open
\nwendcode{}\nwbegincode{41}\moddef{File registries: procedures}\plusendmoddef
  subroutine file_registry_open (registry, file, unit)
    class(file_registry_t), intent(inout) :: registry
    type(string_t), intent(in) :: file
    integer, intent(out), optional :: unit
    type(file_entry_t), pointer :: entry
    entry => registry%first
    FIND_ENTRY: do while (associated (entry))
       if (entry%get_file () == file)  exit FIND_ENTRY
       entry => entry%next
    end do FIND_ENTRY
    if (.not. associated (entry)) then
       allocate (entry)
       call entry%init (file)
       if (associated (registry%first)) then
          registry%first%prev => entry
          entry%next => registry%first
       end if
       registry%first => entry
    end if
    call entry%open ()
    if (present (unit))  unit = entry%get_unit ()
  end subroutine file_registry_open

\eatline
\nwendcode{}\nwbegindocs{42}Close a file: find the appropriate entry.  Delete the entry if there
is no file connected to it anymore.
\nwenddocs{}\nwbegincode{43}\moddef{File registries: file registry: TBP}\plusendmoddef
  procedure :: close => file_registry_close
\nwendcode{}\nwbegincode{44}\moddef{File registries: procedures}\plusendmoddef
  subroutine file_registry_close (registry, file)
    class(file_registry_t), intent(inout) :: registry
    type(string_t), intent(in) :: file
    type(file_entry_t), pointer :: entry
    entry => registry%first
    FIND_ENTRY: do while (associated (entry))
       if (entry%get_file () == file)  exit FIND_ENTRY
       entry => entry%next
    end do FIND_ENTRY
    if (associated (entry)) then
       call entry%close ()
       if (.not. entry%is_open ()) then
          if (associated (entry%prev)) then
             entry%prev%next => entry%next
          else
             registry%first => entry%next
          end if
          if (associated (entry%next)) then
             entry%next%prev => entry%prev
          end if
          deallocate (entry)
       end if
    end if
  end subroutine file_registry_close
    
\eatline
\nwendcode{}\nwbegindocs{45}\nwdocspar
%%%%%%%%%%%%%%%%%%%%%%%%%%%%%%%%%%%%%%%%%%%%%%%%%%%%%%%%%%%%%%%%%%%%%%%% 
\section{String Utilities}

This module provides tools associated with strings
(built-in and variable).  Currently:
\begin{itemize}
\item
  Upper and lower case for strings
\item
  Convert to null-terminated C string
\end{itemize}
\nwenddocs{}\nwbegincode{46}\moddef{\code{}string{\char95}utils.f90\edoc{}}\endmoddef
\LA{}File header\RA{}

module string_utils
  
  use, intrinsic :: iso_c_binding
\LA{}Use kinds\RA{}  
\LA{}Use strings\RA{}

\LA{}Standard module head\RA{}

\LA{}String utils: public\RA{}

\LA{}String utils: interfaces\RA{}

contains

\LA{}String utils: procedures\RA{}

end module string_utils
\eatline
\nwendcode{}\nwbegindocs{47}\nwdocspar
\subsection{Upper and Lower Case}
These are, unfortunately, not part of Fortran.
\nwenddocs{}\nwbegincode{48}\moddef{String utils: public}\endmoddef
  public :: upper_case
  public :: lower_case
\nwendcode{}\nwbegincode{49}\moddef{String utils: interfaces}\endmoddef
  interface upper_case
     module procedure upper_case_char, upper_case_string
  end interface
  interface lower_case
     module procedure lower_case_char, lower_case_string
  end interface
\nwendcode{}\nwbegincode{50}\moddef{String utils: procedures}\endmoddef
  function upper_case_char (string) result (new_string)
    character(*), intent(in) :: string
    character(len(string)) :: new_string
    integer :: pos, code
    integer, parameter :: offset = ichar('A')-ichar('a')
    do pos = 1, len (string)
       code = ichar (string(pos:pos))
       select case (code)
       case (ichar('a'):ichar('z'))
          new_string(pos:pos) = char (code + offset)
       case default
          new_string(pos:pos) = string(pos:pos)
       end select
    end do
  end function upper_case_char

  function lower_case_char (string) result (new_string)
    character(*), intent(in) :: string
    character(len(string)) :: new_string
    integer :: pos, code
    integer, parameter :: offset = ichar('a')-ichar('A')
    do pos = 1, len (string)
       code = ichar (string(pos:pos))
       select case (code)
       case (ichar('A'):ichar('Z'))
          new_string(pos:pos) = char (code + offset)
       case default
          new_string(pos:pos) = string(pos:pos)
       end select
    end do
  end function lower_case_char

  function upper_case_string (string) result (new_string)
    type(string_t), intent(in) :: string
    type(string_t) :: new_string
    new_string = upper_case_char (char (string))
  end function upper_case_string

  function lower_case_string (string) result (new_string)
    type(string_t), intent(in) :: string
    type(string_t) :: new_string
    new_string = lower_case_char (char (string))
  end function lower_case_string

\eatline
\nwendcode{}\nwbegindocs{51}\nwdocspar
\subsection{C-compatible Output}
Convert a FORTRAN string into a zero terminated C string.
\nwenddocs{}\nwbegincode{52}\moddef{String utils: public}\plusendmoddef
  public :: string_f2c
\nwendcode{}\nwbegincode{53}\moddef{String utils: interfaces}\plusendmoddef
  interface string_f2c
     module procedure string_f2c_char, string_f2c_var_str
  end interface string_f2c
\nwendcode{}\nwbegincode{54}\moddef{String utils: procedures}\plusendmoddef
  pure function string_f2c_char (i) result (o)
    character(*), intent(in) :: i
    character(kind=c_char, len=len (i) + 1) :: o
    o = i // c_null_char
  end function string_f2c_char

  pure function string_f2c_var_str (i) result (o)
    type(string_t), intent(in) :: i
    character(kind=c_char, len=len (i) + 1) :: o
    o = char (i) // c_null_char
  end function string_f2c_var_str

\eatline
\nwendcode{}\nwbegindocs{55}\nwdocspar
\subsection{Number Conversion}
Convert an integer to a string.
\nwenddocs{}\nwbegincode{56}\moddef{String utils: public}\plusendmoddef
  public :: integer_to_string
\nwendcode{}\nwbegincode{57}\moddef{String utils: procedures}\plusendmoddef
  function integer_to_string (x) result (str)
    integer, intent(in) :: x
    type(string_t) :: str
    character(5) :: buf
    write (buf, '(I5)') x
    str = var_str (trim (adjustl (buf)))
  end function integer_to_string

\eatline
\nwendcode{}\nwbegindocs{58}\nwdocspar
Create a string from a number.  We use fixed format for the reals
and variable format for integers.
\nwenddocs{}\nwbegincode{59}\moddef{String utils: public}\plusendmoddef
  public :: str
\nwendcode{}\nwbegincode{60}\moddef{String utils: interfaces}\plusendmoddef
  interface str
     module procedure str_int, str_real
  end interface
\nwendcode{}\nwbegincode{61}\moddef{String utils: procedures}\plusendmoddef
  function str_int (i) result (s)
    integer, intent(in) :: i
    type(string_t) :: s
    character(32) :: buffer
    write (buffer, "(I0)")  i
    s = var_str (trim (buffer))
  end function str_int
  
  function str_real (x) result (s)
    real(default), intent(in) :: x
    type(string_t) :: s
    character(32) :: buffer
    write (buffer, "(ES17.10)")  x
    s = var_str (trim (adjustl (buffer)))
  end function str_real
  
\eatline
\nwendcode{}\nwbegindocs{62}Auxiliary: Read real, integer, string value.
\nwenddocs{}\nwbegincode{63}\moddef{String utils: public}\plusendmoddef
  public :: read_rval
  public :: read_ival
\nwendcode{}\nwbegincode{64}\moddef{String utils: procedures}\plusendmoddef
  function read_rval (s) result (rval)
    type(string_t), intent(in) :: s
    real(default) :: rval
    character(80) :: buffer
    buffer = s
    read (buffer, *)  rval
  end function read_rval
    
  function read_ival (s) result (ival)
    type(string_t), intent(in) :: s
    integer :: ival
    character(80) :: buffer
    buffer = s
    read (buffer, *)  ival
  end function read_ival
    
\eatline
\nwendcode{}\nwbegindocs{65}\nwdocspar
%%%%%%%%%%%%%%%%%%%%%%%%%%%%%%%%%%%%%%%%%%%%%%%%%%%%%%%%%%%%%%%%%%%%%%%% 
\section{Format Utilities}

This module provides miscellaneous tools associated with formatting and
pretty-printing.
\begin{itemize}
\item
  Horizontal separator lines in output
\item
  Indenting an output line
\item
  Formatting a number for \TeX\ output.
\item
  Formatting a number for MetaPost output.
\item
  Alternate numeric formats.
\end{itemize}
\nwenddocs{}\nwbegincode{66}\moddef{\code{}format{\char95}utils.f90\edoc{}}\endmoddef
\LA{}File header\RA{}

module format_utils

\LA{}Use kinds\RA{}
\LA{}Use strings\RA{}
  use string_utils, only: lower_case

\LA{}Standard module head\RA{}

\LA{}Format utils: public\RA{}

contains

\LA{}Format utils: procedures\RA{}

end module format_utils
\eatline
\nwendcode{}\nwbegindocs{67}\nwdocspar
\subsection{Line Output}
Write a separator line.
\nwenddocs{}\nwbegincode{68}\moddef{Format utils: public}\endmoddef
  public :: write_separator
\nwendcode{}\nwbegincode{69}\moddef{Format utils: procedures}\endmoddef
  subroutine write_separator (u, mode)
    integer, intent(in) :: u
    integer, intent(in), optional :: mode
    integer :: m
    m = 1;  if (present (mode))  m = mode
    select case (m)
    case default
       write (u, "(A)")  repeat ("-", 72)
    case (1)
       write (u, "(A)")  repeat ("-", 72)
    case (2)
       write (u, "(A)")  repeat ("=", 72)
    end select
  end subroutine write_separator
  
\eatline
\nwendcode{}\nwbegindocs{70}\nwdocspar
Indent the line with given number of blanks.
\nwenddocs{}\nwbegincode{71}\moddef{Format utils: public}\plusendmoddef
  public :: write_indent
\nwendcode{}\nwbegincode{72}\moddef{Format utils: procedures}\plusendmoddef
  subroutine write_indent (unit, indent)
    integer, intent(in) :: unit
    integer, intent(in), optional :: indent
    if (present (indent)) then
       write (unit, "(1x,A)", advance="no")  repeat ("  ", indent)
    end if
  end subroutine write_indent

\eatline
\nwendcode{}\nwbegindocs{73}\nwdocspar
\subsection{\TeX-compatible Output}
Quote underscore characters for use in \TeX\ output.
\nwenddocs{}\nwbegincode{74}\moddef{Format utils: public}\plusendmoddef
  public :: quote_underscore
\nwendcode{}\nwbegincode{75}\moddef{Format utils: procedures}\plusendmoddef
  function quote_underscore (string) result (quoted)
    type(string_t) :: quoted
    type(string_t), intent(in) :: string
    type(string_t) :: part
    type(string_t) :: buffer
    buffer = string
    quoted = ""
    do
      call split (part, buffer, "_")
      quoted = quoted // part
      if (buffer == "")  exit
      quoted = quoted // "\\_"
    end do
  end function quote_underscore

\eatline
\nwendcode{}\nwbegindocs{76}Format a number with $n$ significant digits for use in \TeX\ documents.
\nwenddocs{}\nwbegincode{77}\moddef{Format utils: public}\plusendmoddef
  public :: tex_format
\nwendcode{}\nwbegincode{78}\moddef{Format utils: procedures}\plusendmoddef
  function tex_format (rval, n_digits) result (string)
    type(string_t) :: string
    real(default), intent(in) :: rval
    integer, intent(in) :: n_digits
    integer :: e, n, w, d
    real(default) :: absval
    real(default) :: mantissa
    character :: sign
    character(20) :: format
    character(80) :: cstr
    n = min (abs (n_digits), 16)
    if (rval == 0) then
       string = "0"
    else
       absval = abs (rval)
       e = log10 (absval)
       if (rval < 0) then
          sign = "-"
       else
          sign = ""
       end if
       select case (e)
       case (:-3)
          d = max (n - 1, 0)
          w = max (d + 2, 2)
          write (format, "('(F',I0,'.',I0,',A,I0,A)')")  w, d
          mantissa = absval * 10._default ** (1 - e)
          write (cstr, fmt=format)  mantissa, "\\times 10^\{", e - 1, "\}"
       case (-2:0)
          d = max (n - e, 1 - e)
          w = max (d + e + 2, d + 2)
          write (format, "('(F',I0,'.',I0,')')")  w, d
          write (cstr, fmt=format)  absval
       case (1:2)
          d = max (n - e - 1, -e, 0)
          w = max (d + e + 2, d + 2, e + 2)
          write (format, "('(F',I0,'.',I0,')')")  w, d
          write (cstr, fmt=format)  absval
       case default
          d = max (n - 1, 0)
          w = max (d + 2, 2)
          write (format, "('(F',I0,'.',I0,',A,I0,A)')")  w, d
          mantissa = absval * 10._default ** (- e)
          write (cstr, fmt=format)  mantissa, "\\times 10^\{", e, "\}"
       end select
       string = sign // trim (cstr)
    end if
  end function tex_format

\eatline
\nwendcode{}\nwbegindocs{79}\nwdocspar
\subsection{Metapost-compatible Output}
Write a number for use in Metapost code:
\nwenddocs{}\nwbegincode{80}\moddef{Format utils: public}\plusendmoddef
  public :: mp_format
\nwendcode{}\nwbegincode{81}\moddef{Format utils: procedures}\plusendmoddef
  function mp_format (rval) result (string)
    type(string_t) :: string
    real(default), intent(in) :: rval
    character(16) :: tmp
    write (tmp, "(G16.8)")  rval
    string = lower_case (trim (adjustl (trim (tmp))))
  end function mp_format

\eatline
\nwendcode{}\nwbegindocs{82}\nwdocspar
\subsection{Conditional Formatting}
Conditional format string, intended for switchable numeric precision.
\nwenddocs{}\nwbegincode{83}\moddef{Format utils: public}\plusendmoddef
  public :: pac_fmt
\nwendcode{}\nwbegincode{84}\moddef{Format utils: procedures}\plusendmoddef
  subroutine pac_fmt (fmt, fmt_orig, fmt_pac, pacify) 
    character(*), intent(in) :: fmt_orig, fmt_pac
    character(*), intent(out) :: fmt
    logical, intent(in), optional :: pacify
    logical :: pacified
    pacified = .false.
    if (present (pacify))  pacified = pacify
    if (pacified) then
       fmt = fmt_pac
    else
       fmt = fmt_orig
    end if
  end subroutine pac_fmt
  
\eatline
\nwendcode{}\nwbegindocs{85}\nwdocspar
%%%%%%%%%%%%%%%%%%%%%%%%%%%%%%%%%%%%%%%%%%%%%%%%%%%%%%%%%%%%%%%%%%%%%%%% 
\section{Format Definitions}

This module provides named integer parameters that specify certain
format strings, used for numerical output.
\nwenddocs{}\nwbegincode{86}\moddef{\code{}format{\char95}defs.f90\edoc{}}\endmoddef
\LA{}File header\RA{}

module format_defs

\LA{}Standard module head\RA{}

\LA{}Format defs: public parameters\RA{}

end module format_defs
\eatline
\nwendcode{}\nwbegindocs{87}We collect format strings for various numerical output formats here.
\nwenddocs{}\nwbegincode{88}\moddef{Format defs: public parameters}\endmoddef
  character(*), parameter, public :: FMT_19 = "ES19.12"
  character(*), parameter, public :: FMT_18 = "ES18.11"
  character(*), parameter, public :: FMT_17 = "ES17.10"  
  character(*), parameter, public :: FMT_16 = "ES16.9"
  character(*), parameter, public :: FMT_15 = "ES15.8"
  character(*), parameter, public :: FMT_14 = "ES14.7"    
  character(*), parameter, public :: FMT_13 = "ES13.6"    
  character(*), parameter, public :: FMT_12 = "ES12.5"
  character(*), parameter, public :: FMT_11 = "ES11.4"
  character(*), parameter, public :: FMT_10 = "ES10.3"

\eatline
\eatline
\nwendcode{}\nwfilename{../../../src/testing/testing.nw}\nwbegindocs{0}% -*- ess-noweb-default-code-mode: f90-mode; noweb-default-code-mode: f90-mode; -*-
% WHIZARD code as NOWEB source

\nwenddocs{}\nwbegindocs{1}\nwdocspar
\chapter{Testing}
This part contains tools for automatic testing.
\begin{description}
\item[unit\_tests]
  A handler that executes test procedures and compares and collects results.
\end{description}

%%%%%%%%%%%%%%%%%%%%%%%%%%%%%%%%%%%%%%%%%%%%%%%%%%%%%%%%%%%%%%%%%%%%%%%%
\section{Unit tests}
We provide functionality for automated unit tests.  Each test is
required to produce output which is compared against a reference
file.  If the two are identical, we signal success.  Otherwise, we
signal failure and write the output to a file.
\nwenddocs{}\nwbegincode{2}\moddef{\code{}unit{\char95}tests.f90\edoc{}}\endmoddef
\LA{}File header\RA{}

module unit_tests

\LA{}Use kinds\RA{}
  use constants, only: zero, tiny_10, tiny_13
\LA{}Use strings\RA{}
  use io_units

\LA{}Standard module head\RA{}

\LA{}Tests: public\RA{}

\LA{}Tests: parameters\RA{}

\LA{}Tests: types\RA{}

\LA{}Tests: interfaces\RA{}

contains

\LA{}Tests: procedures\RA{}

end module unit_tests
\eatline
\nwendcode{}\nwbegindocs{3}\nwdocspar
\subsection{Parameters}
Building blocks of file names.  The directory names and suffixes are
hard-coded here, and they must reflect actual Makefile targets where
applicable.
\nwenddocs{}\nwbegincode{4}\moddef{Tests: parameters}\endmoddef
  character(*), parameter :: ref_prefix = "ref-output/"
  character(*), parameter :: ref = ".ref"

  character(*), parameter :: err_prefix = "err-output/"
  character(*), parameter :: err = ".out"

\eatline
\nwendcode{}\nwbegindocs{5}\nwdocspar
\subsection{Type for storing test results}
We store the results of the individual unit tests in a linked list.
Here is the entry:
\nwenddocs{}\nwbegincode{6}\moddef{Tests: public}\endmoddef
  public :: test_results_t
\nwendcode{}\nwbegincode{7}\moddef{Tests: types}\endmoddef
  type :: test_result_t
     logical :: success = .false.
     type(string_t) :: name
     type(string_t) :: description
     type(test_result_t), pointer :: next => null ()
  end type test_result_t

  type :: test_results_t
     private
     type(test_result_t), pointer :: first => null ()
     type(test_result_t), pointer :: last => null ()
     integer :: n_success = 0
     integer :: n_failure = 0
   contains
   \LA{}Tests: test results: TBP\RA{}
  end type test_results_t

\eatline
\nwendcode{}\nwbegindocs{8}Add a test result.
\nwenddocs{}\nwbegincode{9}\moddef{Tests: test results: TBP}\endmoddef
  procedure, private :: add => test_results_add
\nwendcode{}\nwbegincode{10}\moddef{Tests: procedures}\endmoddef
  subroutine test_results_add (list, name, description, success)
    class(test_results_t), intent(inout) :: list
    character(len=*), intent(in) :: name
    character(len=*), intent(in) :: description
    logical, intent(in) :: success
    type(test_result_t), pointer :: result
    allocate (result)
    result%success = success
    result%name = name
    result%description = description
    if (associated (list%first)) then
       list%last%next => result
    else
       list%first => result
    end if
    list%last => result
    if (success) then
       list%n_success = list%n_success + 1
    else
       list%n_failure = list%n_failure + 1
    end if
  end subroutine test_results_add

\eatline
\nwendcode{}\nwbegindocs{11}Display the current state.
\nwenddocs{}\nwbegincode{12}\moddef{Tests: test results: TBP}\plusendmoddef
  procedure, private :: write => test_results_write
\nwendcode{}\nwbegincode{13}\moddef{Tests: procedures}\plusendmoddef
  subroutine test_results_write (list, u)
    class(test_results_t), intent(in) :: list
    integer, intent(in) :: u
    type(test_result_t), pointer :: result
    write (u, "(A)")  "*** Test Summary ***"
    if (list%n_success > 0) then
       write (u, "(2x,A)")  "Success:"
       result => list%first
       do while (associated (result))
          if (result%success)  write (u, "(4x,A,': ',A)")  &
               char (result%name), char (result%description)
          result => result%next
       end do
    end if
    if (list%n_failure > 0) then
       write (u, "(2x,A)")  "Failure:"
       result => list%first
       do while (associated (result))
          if (.not. result%success)  write (u, "(4x,A,': ',A)")  &
               char (result%name), char (result%description)
          result => result%next
       end do
    end if
    write (u, "(A,I0)")  "Total   = ", list%n_success + list%n_failure
    write (u, "(A,I0)")  "Success = ", list%n_success
    write (u, "(A,I0)")  "Failure = ", list%n_failure
    write (u, "(A)")  "*** End of test Summary ***"
  end subroutine test_results_write

\eatline
\nwendcode{}\nwbegindocs{14}Return true if all tests were successful (or no test).
\nwenddocs{}\nwbegincode{15}\moddef{Tests: test results: TBP}\plusendmoddef
  procedure, private :: report => test_results_report
\nwendcode{}\nwbegincode{16}\moddef{Tests: procedures}\plusendmoddef
  subroutine test_results_report (list, success)
    class(test_results_t), intent(in) :: list
    logical, intent(out) :: success
    success = list%n_failure == 0
  end subroutine test_results_report

\eatline
\nwendcode{}\nwbegindocs{17}Delete the list.
\nwenddocs{}\nwbegincode{18}\moddef{Tests: test results: TBP}\plusendmoddef
  procedure, private :: final => test_results_final
\nwendcode{}\nwbegincode{19}\moddef{Tests: procedures}\plusendmoddef
  subroutine test_results_final (list)
    class(test_results_t), intent(inout) :: list
    type(test_result_t), pointer :: result
    do while (associated (list%first))
       result => list%first
       list%first => result%next
       deallocate (result)
    end do
    list%last => null ()
    list%n_success = 0
    list%n_failure = 0
  end subroutine test_results_final

\eatline
\nwendcode{}\nwbegindocs{20}\nwdocspar
\subsection{Wrapup}
This will write results, report status, and finalize.  This is the
only method which we need to access from outside.
\nwenddocs{}\nwbegincode{21}\moddef{Tests: test results: TBP}\plusendmoddef
  procedure :: wrapup => test_results_wrapup
\nwendcode{}\nwbegincode{22}\moddef{Tests: procedures}\plusendmoddef
  subroutine test_results_wrapup (list, u, success)
    class(test_results_t), intent(inout) :: list
    integer, intent(in) :: u
    logical, intent(out), optional :: success
    call list%write (u)
    if (present (success))  call list%report (success)
    call list%final ()
  end subroutine test_results_wrapup

\eatline
\nwendcode{}\nwbegindocs{23}\nwdocspar
\subsection{Tool for Unit Tests}
This procedure takes a test routine as an argument.  It runs the test,
output directed to a temporary file.  Then, it compares the file
against a reference file.

The test routine must take the output unit as argument.
\nwenddocs{}\nwbegincode{24}\moddef{Tests: interfaces}\endmoddef
  abstract interface
     subroutine unit_test (u)
       integer, intent(in) :: u
     end subroutine unit_test
  end interface

\eatline
\nwendcode{}\nwbegindocs{25}The test routine can print to screen and, optionally, to a logging unit.
\nwenddocs{}\nwbegincode{26}\moddef{Tests: public}\plusendmoddef
  public :: test
\nwendcode{}\nwbegincode{27}\moddef{Tests: procedures}\plusendmoddef
  subroutine test (test_proc, name, description, u_log, results)
    procedure(unit_test) :: test_proc
    character(*), intent(in) :: name
    character(*), intent(in) :: description
    integer, intent(in) :: u_log
    type(test_results_t), intent(inout) :: results
    integer :: u_test, u_ref, u_err
    logical :: exist
    character(256) :: buffer1, buffer2
    integer :: iostat1, iostat2
    logical :: success
    write (*, "(A)", advance="no")  "Running test: " // name
    write (u_log, "(A)")  "Test: " // name
    u_test = free_unit ()
    open (u_test, status="scratch", action="readwrite")
    call test_proc (u_test)
    rewind (u_test)
    inquire (file=ref_prefix//name//ref, exist=exist)
    if (exist) then
       u_ref = free_unit ()
       open (u_ref, file=ref_prefix//name//ref, status="old", action="read")
       COMPARE_FILES: do
          read (u_test, "(A)", iostat=iostat1)  buffer1
          read (u_ref, "(A)", iostat=iostat2)  buffer2
          if (iostat1 /= iostat2) then
             success = .false.
             exit COMPARE_FILES
          else if (iostat1 < 0) then
             success = .true.
             exit COMPARE_FILES
          else if (buffer1 /= buffer2) then
             success = .false.
             exit COMPARE_FILES
          end if
       end do COMPARE_FILES
       close (u_ref)
    else
       write (*, "(A)", advance="no") " ... no reference output available"
       write (u_log, "(A)") "  No reference output available."
       success = .false.
    end if
    if (success) then
       write (*, "(A)") " ... success."
       write (u_log, "(A)")  "  Success."
    else
       write (*, "(A)") " ... failure.  See: " // err_prefix//name//err
       write (u_log, "(A)")  "  Failure."
       rewind (u_test)
       u_err = free_unit ()
       open (u_err, file=err_prefix//name//err, &
            action="write", status="replace")
       WRITE_OUTPUT: do
          read (u_test, "(A)", end=1) buffer1
          write (u_err, "(A)")  trim (buffer1)
       end do WRITE_OUTPUT
1      close (u_err)
    end if
    close (u_test)
    call results%add (name, description, success)
  end subroutine test

\eatline
\nwendcode{}\nwbegindocs{28}\subsection{Convenience functions}
\nwenddocs{}\nwbegindocs{29}Unit tests should include a lot of assertions that only involve IO if
something is wrong. {\tt{}assert} allows to streamline this method.
\nwenddocs{}\nwbegincode{30}\moddef{Tests: public}\plusendmoddef
  public :: assert
\nwendcode{}\nwbegincode{31}\moddef{Tests: procedures}\plusendmoddef
  subroutine assert (unit, ok, description)
    integer, intent(in) :: unit
    logical, intent(in) :: ok
    character(*), intent(in), optional :: description
    if (.not. ok) then
       if (present(description)) then
          write (unit, "(A)") "* FAIL: " // description
       else
          write (unit, "(A)") "* FAIL: Assertion error"
       end if
    end if
  end subroutine assert

\eatline
\nwendcode{}\nwbegindocs{32}Note that this poor man's check will be disabled if someone compiles
with {\tt{}-ffast-math} or similar optimizations.
\nwenddocs{}\nwbegincode{33}\moddef{Tests: procedures}\plusendmoddef
  elemental function ieee_is_nan (x) result (yorn)
    logical :: yorn
    real(default), intent(in) :: x
    yorn = (x /= x)
  end function ieee_is_nan
\nwendcode{}\nwbegindocs{34}This is still not perfect but should work in most cases.  Usually one
wants to compare to a relative epsilon {\tt{}rel{\char95}smallness}, except for
numbers close to zero defined by {\tt{}abs{\char95}smallness}.  Both might need
adaption to specific use cases but have reasonable defaults.
\nwenddocs{}\nwbegincode{35}\moddef{Tests: public}\plusendmoddef
  public :: nearly_equal
\nwendcode{}\nwbegincode{36}\moddef{Tests: procedures}\plusendmoddef
  elemental function nearly_equal (a, b, abs_smallness, rel_smallness) result (r)
    logical :: r
    real(default), intent(in) :: a, b
    real(default), intent(in), optional :: abs_smallness, rel_smallness
    real(default) :: abs_a, abs_b, diff, abs_small, rel_small
    abs_a = abs (a)
    abs_b = abs (b)
    diff = abs (a - b)
    ! shortcut, handles infinities and nans
    if (a == b) then
       r = .true.
       return
    else if (ieee_is_nan (a) .or. ieee_is_nan (b) .or. ieee_is_nan (diff)) then
       r = .false.
       return
    end if
    if (present (abs_smallness)) then
       abs_small = abs_smallness
    else
       abs_small = tiny_13
    end if
    if (present (rel_smallness)) then
       rel_small = rel_smallness
    else
       rel_small = tiny_10
    end if
    if (abs_a < abs_small .and. abs_b < abs_small) then
       r = diff < abs_small
    else
       r = diff / max (abs_a, abs_b) < rel_small
    end if
  end function nearly_equal

\eatline
\nwendcode{}\nwbegindocs{37}Often we will need to check whether floats vanish:
\nwenddocs{}\nwbegincode{38}\moddef{Tests: public}\plusendmoddef
  public:: vanishes
  interface vanishes
     module procedure vanishes_real, vanishes_complex
  end interface
\nwendcode{}\nwbegindocs{39}\nwdocspar
\nwenddocs{}\nwbegincode{40}\moddef{Tests: procedures}\plusendmoddef
  elemental function vanishes_real (x, abs_smallness, rel_smallness) result (r)
    logical :: r
    real(default), intent(in) :: x
    real(default), intent(in), optional :: abs_smallness, rel_smallness
    r = nearly_equal (x, zero, abs_smallness, rel_smallness)
  end function vanishes_real

  elemental function vanishes_complex (x, abs_smallness, rel_smallness) result (r)
    logical :: r
    complex(default), intent(in) :: x
    real(default), intent(in), optional :: abs_smallness, rel_smallness
    r = vanishes_real (abs (x), abs_smallness, rel_smallness)
  end function vanishes_complex

\eatline
\nwendcode{}\nwbegindocs{41}\nwdocspar
\nwenddocs{}\nwfilename{../../../src/system/system.nw}\nwbegindocs{0}% -*- ess-noweb-default-code-mode: f90-mode; noweb-default-code-mode: f90-mode; -*- 
% WHIZARD code as NOWEB source: system interfaces

\nwenddocs{}\nwbegindocs{1}\nwdocspar
\chapter{System: Interfaces and Handlers}
Here, we collect modules that deal with the ``system'':
operating-system interfaces, error handlers and diagnostics.
\begin{description}
\item[system\_defs] 
  Constants relevant for the modules in this set.
\item[diagnostics] 
  Error and diagnostic message handling.  Any
  messages and errors issued by WHIZARD functions are handled by the
  subroutines in this module, if possible.
\item[os\_interface] 
  Execute system calls, build and link external object files and libraries.
\item[cputime]
  Timer data type and methods, for measuring performance.
\end{description}

%%%%%%%%%%%%%%%%%%%%%%%%%%%%%%%%%%%%%%%%%%%%%%%%%%%%%%%%%%%%%%%%%%%%%%%%
\section{Constants}
The parameters here are used in various parts of the program,
starting from the modules in the current chapter.  Some of them may be
modified if the need arises.
\nwenddocs{}\nwbegincode{2}\moddef{\code{}system{\char95}defs.f90\edoc{}}\endmoddef
\LA{}File header\RA{}

module system_defs

  use, intrinsic :: iso_fortran_env, only: iostat_end, iostat_eor !NODEP!

\LA{}Standard module head\RA{}

\LA{}System defs: public parameters\RA{}

end module system_defs
\eatline
\nwendcode{}\nwbegindocs{3}\nwdocspar
\subsection{Version}
The version string is used for checking files.  Note that the string
length MUST NOT be changed, because reading binary files relies on it.
\nwenddocs{}\nwbegincode{4}\moddef{System defs: public parameters}\endmoddef
  integer, parameter, public :: VERSION_STRLEN = 255
  character(len=VERSION_STRLEN), parameter, public :: &
       & VERSION_STRING = "WHIZARD version \LA{}Version\RA{} (\LA{}Date\RA{})"

\eatline
\nwendcode{}\nwbegindocs{5}\nwdocspar
\subsection{Text Buffer}
There is a hard limit on the line length which we should export.  This
buffer size is used both by the message handler, the lexer, and some
further modules.
\nwenddocs{}\nwbegincode{6}\moddef{System defs: public parameters}\plusendmoddef
  integer, parameter, public :: BUFFER_SIZE = 1000

\eatline
\nwendcode{}\nwbegindocs{7}\nwdocspar
\subsection{IOSTAT Codes}
Defined in {\tt{}iso{\char95}fortran{\char95}env}, but we would like to use shorthands.
\nwenddocs{}\nwbegincode{8}\moddef{System defs: public parameters}\plusendmoddef
  integer, parameter, public :: EOF = iostat_end,  EOR = iostat_eor

\eatline
\nwendcode{}\nwbegindocs{9}\nwdocspar
\subsection{Character Codes}
Single-character constants.
\nwenddocs{}\nwbegincode{10}\moddef{System defs: public parameters}\plusendmoddef
  character, parameter, public :: BLANK = ' '
  character, parameter, public :: TAB = achar(9)
  character, parameter, public :: CR = achar(13)
  character, parameter, public :: LF = achar(10)
  character, parameter, public :: BACKSLASH = achar(92)

\eatline
\nwendcode{}\nwbegindocs{11}Character strings that indicate character classes.
\nwenddocs{}\nwbegincode{12}\moddef{System defs: public parameters}\plusendmoddef
  character(*), parameter, public :: WHITESPACE_CHARS = BLANK// TAB // CR // LF
  character(*), parameter, public :: LCLETTERS = "abcdefghijklmnopqrstuvwxyz"
  character(*), parameter, public :: UCLETTERS = "ABCDEFGHIJKLMNOPQRSTUVWXYZ"
  character(*), parameter, public :: DIGITS = "0123456789"

\eatline
\nwendcode{}\nwbegindocs{13}\nwdocspar
%%%%%%%%%%%%%%%%%%%%%%%%%%%%%%%%%%%%%%%%%%%%%%%%%%%%%%%%%%%%%%%%%%%%%%%%
\section{C wrapper for sigaction}
This implements calls to {\tt{}sigaction} and the appropriate signal handlers in
C.  The functionality is needed for the {\tt{}diagnostics} module.
\nwenddocs{}\nwbegincode{14}\moddef{\code{}signal{\char95}interface.c\edoc{}}\endmoddef
/*
\LA{}File header\RA{}
*/
#include <signal.h>
#include <stdlib.h>

extern int wo_sigint;
extern int wo_sigterm;
extern int wo_sigxcpu;
extern int wo_sigxfsz;

static void wo_handler_sigint (int sig) \{
  wo_sigint = sig;
\}

static void wo_handler_sigterm (int sig) \{
  wo_sigterm = sig;
\}

static void wo_handler_sigxcpu (int sig) \{
  wo_sigxcpu = sig;
\}

static void wo_handler_sigxfsz (int sig) \{
  wo_sigxfsz = sig;
\}

int wo_mask_sigint () \{
  struct sigaction sa;
  sigset_t blocks;
  sigfillset (&blocks);
  sa.sa_flags = 0;
  sa.sa_mask = blocks;
  sa.sa_handler = wo_handler_sigint;
  return sigaction(SIGINT, &sa, NULL);
\}

int wo_mask_sigterm () \{
  struct sigaction sa;
  sigset_t blocks;
  sigfillset (&blocks);
  sa.sa_flags = 0;
  sa.sa_mask = blocks;
  sa.sa_handler = wo_handler_sigterm;
  return sigaction(SIGTERM, &sa, NULL);
\}

int wo_mask_sigxcpu () \{
  struct sigaction sa;  
  sigset_t blocks;
  sigfillset (&blocks);
  sa.sa_flags = 0;
  sa.sa_mask = blocks;
  sa.sa_handler = wo_handler_sigxcpu;
  return sigaction(SIGXCPU, &sa, NULL);
\}

int wo_mask_sigxfsz () \{
  struct sigaction sa;
  sigset_t blocks;
  sigfillset (&blocks);
  sa.sa_flags = 0;
  sa.sa_mask = blocks;
  sa.sa_handler = wo_handler_sigxfsz;
  return sigaction(SIGXFSZ, &sa, NULL);
\}

int wo_release_sigint () \{
  struct sigaction sa;  
  sigset_t blocks;
  sigfillset (&blocks);
  sa.sa_flags = 0;
  sa.sa_mask = blocks;
  sa.sa_handler = SIG_DFL;
  return sigaction(SIGINT, &sa, NULL);
\}

int wo_release_sigterm () \{
  struct sigaction sa;
  sigset_t blocks;
  sigfillset (&blocks);
  sa.sa_flags = 0;
  sa.sa_mask = blocks;
  sa.sa_handler = SIG_DFL;
  return sigaction(SIGTERM, &sa, NULL);
\}

int wo_release_sigxcpu () \{
  struct sigaction sa;
  sigset_t blocks;
  sigfillset (&blocks);
  sa.sa_flags = 0;
  sa.sa_mask = blocks;
  sa.sa_handler = SIG_DFL;
  return sigaction(SIGXCPU, &sa, NULL);
\}

int wo_release_sigxfsz () \{
  struct sigaction sa;
  sigset_t blocks;
  sigfillset (&blocks);
  sa.sa_flags = 0;
  sa.sa_mask = blocks;
  sa.sa_handler = SIG_DFL;
  return sigaction(SIGXFSZ, &sa, NULL);
\}
\nwendcode{}\nwbegindocs{15}\nwdocspar
%%%%%%%%%%%%%%%%%%%%%%%%%%%%%%%%%%%%%%%%%%%%%%%%%%%%%%%%%%%%%%%%%%%%%%%%
\section{C wrapper for printf}
The {\tt{}printf} family of functions is implemented in C with an undefined
number of arguments.  This is not supported by the {\tt{}bind(C)} interface.  We
therefore write wrappers for the versions of {\tt{}sprintf} that
we will actually use.

This is used by the {\tt{}formats} module.
\nwenddocs{}\nwbegincode{16}\moddef{\code{}sprintf{\char95}interface.c\edoc{}}\endmoddef
/*
\LA{}File header\RA{}
*/
#include <stdio.h>

int sprintf_none(char* str, const char* format) \{
  return sprintf(str, format);
\}

int sprintf_int(char* str, const char* format, int val) \{
  return sprintf(str, format, val);
\}

int sprintf_double(char* str, const char* format, double val) \{
  return sprintf(str, format, val);
\}

int sprintf_str(char* str, const char* format, const char* val) \{
  return sprintf(str, format, val);
\}

\nwendcode{}\nwbegincode{17}\moddef{sprintf interfaces}\endmoddef
  interface
    function sprintf_none (str, fmt) result (stat) bind(C)
      use iso_c_binding !NODEP!
      integer(c_int) :: stat
      character(c_char), dimension(*), intent(inout) :: str
      character(c_char), dimension(*), intent(in) :: fmt
    end function sprintf_none
  end interface
      
  interface
    function sprintf_int (str, fmt, val) result (stat) bind(C)
      use iso_c_binding !NODEP!
      integer(c_int) :: stat
      character(c_char), dimension(*), intent(inout) :: str
      character(c_char), dimension(*), intent(in) :: fmt
      integer(c_int), value :: val
    end function sprintf_int
  end interface
      
  interface
    function sprintf_double (str, fmt, val) result (stat) bind(C)
      use iso_c_binding !NODEP!
      integer(c_int) :: stat
      character(c_char), dimension(*), intent(inout) :: str
      character(c_char), dimension(*), intent(in) :: fmt
      real(c_double), value :: val
    end function sprintf_double
  end interface
      
  interface
    function sprintf_str(str, fmt, val) result (stat) bind(C)
      use iso_c_binding !NODEP!
      integer(c_int) :: stat
      character(c_char), dimension(*), intent(inout) :: str
      character(c_char), dimension(*), intent(in) :: fmt
      character(c_char), dimension(*), intent(in) :: val
    end function sprintf_str
  end interface

\eatline
\nwendcode{}\nwbegindocs{18}\nwdocspar
%%%%%%%%%%%%%%%%%%%%%%%%%%%%%%%%%%%%%%%%%%%%%%%%%%%%%%%%%%%%%%%%%%%%%%%%
\section{Error, Message and Signal Handling}
We are not so ambitious as to do proper exception handling in
{\tt{}WHIZARD}, but at least it may be useful to have a common interface
for diagnostics: Results, messages, warnings, and such.  As module
variables we keep a buffer where the current message may be written to
and a level indicator which tells which messages should be written on
screen and which ones should be skipped.  Alternatively, a string may
be directly supplied to the message routine: this overrides the
buffer, avoiding the necessety of formatted I/O in trivial cases.
\nwenddocs{}\nwbegincode{19}\moddef{\code{}diagnostics.f90\edoc{}}\endmoddef
\LA{}File header\RA{}

module diagnostics

  use, intrinsic :: iso_c_binding !NODEP! 
  use, intrinsic :: iso_fortran_env, only: output_unit !NODEP!

\LA{}Use kinds\RA{}
\LA{}Use strings\RA{}
  use io_units

  use system_dependencies
  use system_defs, only: BUFFER_SIZE, MAX_ERRORS

\LA{}Standard module head\RA{}

\LA{}Diagnostics: public\RA{}

\LA{}Diagnostics: parameters\RA{}

\LA{}Diagnostics: types\RA{}

\LA{}Diagnostics: variables\RA{}

\LA{}Diagnostics: interfaces\RA{}

contains

\LA{}Diagnostics: procedures\RA{}

end module diagnostics
\eatline
\nwendcode{}\nwbegindocs{20}\nwdocspar
Diagnostics levels:
\nwenddocs{}\nwbegincode{21}\moddef{Diagnostics: parameters}\endmoddef
  integer, parameter :: &
       & TERMINATE=-2, BUG=-1, &
       & FATAL=1, ERROR=2, WARNING=3, MESSAGE=4, RESULT=5, DEBUG=6
\eatline
\nwendcode{}\nwbegincode{22}\moddef{Diagnostics: variables}\endmoddef
  integer, save :: msg_level = RESULT
\eatline
\nwendcode{}\nwbegindocs{23}\nwdocspar
Mask fatal errors so that are treated as normal errors.  Useful for
interactive mode.
\nwenddocs{}\nwbegincode{24}\moddef{Diagnostics: public}\endmoddef
  public :: mask_fatal_errors
\nwendcode{}\nwbegincode{25}\moddef{Diagnostics: variables}\plusendmoddef
  logical, save :: mask_fatal_errors = .false.
\eatline
\nwendcode{}\nwbegindocs{26}\nwdocspar
How to handle bugs and unmasked fatal errors.  Either execute a normal
stop statement, or call the C {\tt{}exit()} function, or try to cause a
program crash by dereferencing a null pointer.
\nwenddocs{}\nwbegincode{27}\moddef{Diagnostics: parameters}\plusendmoddef
  integer, parameter :: TERM_STOP = 0, TERM_EXIT = 1, TERM_CRASH = 2
\eatline
\nwendcode{}\nwbegincode{28}\moddef{Diagnostics: variables}\plusendmoddef
  integer, save :: handle_fatal_errors = TERM_EXIT
\nwendcode{}\nwbegindocs{29}\nwdocspar
Keep track of errors.  This might be used for exception handling,
later.  The counter is incremented only for screen messages, to avoid
double counting.
\nwenddocs{}\nwbegincode{30}\moddef{Diagnostics: public}\plusendmoddef
  public :: msg_count
\nwendcode{}\nwbegincode{31}\moddef{Diagnostics: variables}\plusendmoddef
  integer, dimension(TERMINATE:DEBUG), save :: msg_count = 0
\eatline
\nwendcode{}\nwbegindocs{32}Keep a list of all errors and warnings.  Since we do not know the number of
entries beforehand, we use a linked list.
\nwenddocs{}\nwbegincode{33}\moddef{Diagnostics: types}\endmoddef
  type :: string_list
     character(len=BUFFER_SIZE) :: string
     type(string_list), pointer :: next
  end type string_list
  type :: string_list_pointer
     type(string_list), pointer :: first, last
  end type string_list_pointer
  
\eatline
\nwendcode{}\nwbegincode{34}\moddef{Diagnostics: variables}\plusendmoddef
  type(string_list_pointer), dimension(TERMINATE:WARNING), save :: &
       & msg_list = string_list_pointer (null(), null())
\eatline
\nwendcode{}\nwbegindocs{35}Add the current message buffer contents to the internal list.
\nwenddocs{}\nwbegincode{36}\moddef{Diagnostics: procedures}\endmoddef
  subroutine msg_add (level)
    integer, intent(in) :: level
    type(string_list), pointer :: message
    select case (level)
    case (TERMINATE:WARNING)
       allocate (message)
       message%string = msg_buffer
       nullify (message%next)
       if (.not.associated (msg_list(level)%first)) &
            & msg_list(level)%first => message
       if (associated (msg_list(level)%last)) &
            & msg_list(level)%last%next => message
       msg_list(level)%last => message
       msg_count(level) = msg_count(level) + 1
    end select
  end subroutine msg_add

\eatline
\nwendcode{}\nwbegindocs{37}\nwdocspar
Initialization: 
\nwenddocs{}\nwbegincode{38}\moddef{Diagnostics: public}\plusendmoddef
  public :: msg_list_clear
\nwendcode{}\nwbegincode{39}\moddef{Diagnostics: procedures}\plusendmoddef
  subroutine msg_list_clear
    integer :: level
    type(string_list), pointer :: message
    do level = TERMINATE, WARNING
       do while (associated (msg_list(level)%first))
          message => msg_list(level)%first
          msg_list(level)%first => message%next
          deallocate (message)
       end do
       nullify (msg_list(level)%last)
    end do
    msg_count = 0
  end subroutine msg_list_clear

\eatline
\nwendcode{}\nwbegindocs{40}Display the summary of errors and warnings (no need to count
fatals\ldots)
\nwenddocs{}\nwbegincode{41}\moddef{Diagnostics: public}\plusendmoddef
  public :: msg_summary
\nwendcode{}\nwbegincode{42}\moddef{Diagnostics: procedures}\plusendmoddef
  subroutine msg_summary (unit)
    integer, intent(in), optional :: unit
    call expect_summary (unit)
1   format (A,1x,I2,1x,A,I2,1x,A)
    if (msg_count(ERROR) > 0 .and. msg_count(WARNING) > 0) then
       write (msg_buffer, 1) "There were", &
            & msg_count(ERROR), "error(s) and  ", &
            & msg_count(WARNING), "warning(s)."
       call msg_message (unit=unit)
    else if (msg_count(ERROR) > 0) then
       write (msg_buffer, 1) "There were", &
            & msg_count(ERROR), "error(s) and no warnings."
       call msg_message (unit=unit)
    else if (msg_count(WARNING) > 0) then
       write (msg_buffer, 1) "There were no errors and  ", &
            & msg_count(WARNING), "warning(s)."
       call msg_message (unit=unit)
    end if
  end subroutine msg_summary

\eatline
\nwendcode{}\nwbegindocs{43}Print the list of all messages of a given level.
\nwenddocs{}\nwbegincode{44}\moddef{Diagnostics: public}\plusendmoddef
  public :: msg_listing
\nwendcode{}\nwbegincode{45}\moddef{Diagnostics: procedures}\plusendmoddef
  subroutine msg_listing (level, unit, prefix)
    integer, intent(in) :: level
    integer, intent(in), optional :: unit
    character(len=*), intent(in), optional :: prefix
    type(string_list), pointer :: message
    integer :: u
    u = given_output_unit (unit);  if (u < 0)  return
    if (present (unit))  u = unit
    message => msg_list(level)%first
    do while (associated (message))
       if (present (prefix)) then
          write (u, "(A)") prefix // trim (message%string)
       else
          write (u, "(A)") trim (message%string)
       end if
       message => message%next
    end do
    flush (u)
  end subroutine msg_listing

\eatline
\nwendcode{}\nwbegindocs{46}The message buffer:
\nwenddocs{}\nwbegincode{47}\moddef{Diagnostics: public}\plusendmoddef
  public :: msg_buffer
\nwendcode{}\nwbegincode{48}\moddef{Diagnostics: variables}\plusendmoddef
  character(len=BUFFER_SIZE), save :: msg_buffer = " "
\eatline
\nwendcode{}\nwbegindocs{49}\nwdocspar
After a message is issued, the buffer should be cleared:
\nwenddocs{}\nwbegincode{50}\moddef{Diagnostics: procedures}\plusendmoddef
  subroutine buffer_clear
    msg_buffer = " "
  end subroutine buffer_clear

\eatline
\nwendcode{}\nwbegindocs{51}The generic handler for messages.  If the unit is omitted (or $=6$), the
message is written to standard output if the precedence if
sufficiently high (as determined by the value of
{\tt{}msg{\char95}level}).  If the string is omitted, the buffer is used.
In any case, the buffer is cleared after printing.  In accordance with
FORTRAN custom, the first column in the output is left blank.  For
messages and warnings, an additional exclamation mark and a blank is
prepended.  Furthermore, each message is appended to the internal
message list (without prepending anything).
\nwenddocs{}\nwbegincode{52}\moddef{Diagnostics: procedures}\plusendmoddef
  subroutine message_print (level, string, str_arr, unit, logfile)
    integer, intent(in) :: level
    character(len=*), intent(in), optional :: string
    type(string_t), dimension(:), intent(in), optional :: str_arr
    integer, intent(in), optional :: unit
    logical, intent(in), optional :: logfile   
    type(string_t) :: prep_string, aux_string, head_footer
    integer :: lu, i
    logical :: severe, is_error
    severe = .false.
    head_footer  = "******************************************************************************"
    aux_string = ""
    is_error = .false.
    select case (level)
    case (TERMINATE)
       prep_string = ""
    case (BUG)
       prep_string   = "*** WHIZARD BUG: "
       aux_string    = "***              "
       severe = .true.
       is_error = .true.
    case (FATAL)
       prep_string   = "*** FATAL ERROR: "
       aux_string    = "***              "                     
       severe = .true.
       is_error = .true.
    case (ERROR)
       prep_string = "*** ERROR: "
       aux_string  = "***        "
       is_error = .true.
    case (WARNING)
       prep_string = "Warning: "
    case (MESSAGE, DEBUG)
       prep_string = "| "
    case default
       prep_string = ""
    end select
    if (present(string))  msg_buffer = string
    lu = log_unit
    if (present(unit)) then
       if (unit /= output_unit) then
          if (severe) write (unit, "(A)") char(head_footer) 
          if (is_error) write (unit, "(A)") char(head_footer) 
          write (unit, "(A,A)") char(prep_string), trim(msg_buffer)
          if (present (str_arr)) then
             do i = 1, size(str_arr)
                write (unit, "(A,A)") char(aux_string), char(trim(str_arr(i)))
             end do
          end if
          if (is_error) write (unit, "(A)") char(head_footer) 
          if (severe) write (unit, "(A)") char(head_footer)
          flush (unit)
          lu = -1
       else if (level <= msg_level) then
          if (severe) print "(A)", char(head_footer) 
          if (is_error) print "(A)", char(head_footer) 
          print "(A,A)", char(prep_string), trim(msg_buffer)
          if (present (str_arr)) then
             do i = 1, size(str_arr)
                print "(A,A)", char(aux_string), char(trim(str_arr(i)))
             end do                
          end if
          if (is_error) print "(A)", char(head_footer) 
          if (severe) print "(A)", char(head_footer)
          flush (output_unit)
          if (unit == log_unit)  lu = -1
       end if
    else if (level <= msg_level) then
       if (severe) print "(A)", char(head_footer) 
       if (is_error) print "(A)", char(head_footer) 
       print "(A,A)", char(prep_string), trim(msg_buffer)
          if (present (str_arr)) then
             do i = 1, size(str_arr)
                print "(A,A)", char(aux_string), char(trim(str_arr(i)))
             end do                
          end if
       if (is_error) print "(A)", char(head_footer) 
       if (severe) print "(A)", char(head_footer) 
       flush (output_unit)
    end if
    if (present (logfile)) then
       if (.not. logfile)  lu = -1
    end if
    if (logging .and. lu >= 0) then
       if (severe) write (lu, "(A)") char(head_footer) 
       if (is_error) write (lu, "(A)") char(head_footer) 
       write (lu, "(A,A)")  char(prep_string), trim(msg_buffer)
       if (present (str_arr)) then
          do i = 1, size(str_arr)
             write (lu, "(A,A)") char(aux_string), char(trim(str_arr(i)))
          end do                
       end if
       if (is_error) write (lu, "(A)") char(head_footer) 
       if (severe) write (lu, "(A)") char(head_footer) 
       flush (lu)
    end if
    call msg_add (level)
    call buffer_clear
  end subroutine message_print

\eatline
\nwendcode{}\nwbegindocs{53}\nwdocspar
The number of non-fatal errors that we allow before stopping the
program.  We might trade this later for an adjustable number.
\nwenddocs{}\nwbegincode{54}\moddef{System defs: public parameters}\plusendmoddef
  integer, parameter, public :: MAX_ERRORS = 10
\eatline
\nwendcode{}\nwbegindocs{55}The specific handlers.  In the case of fatal errors, bugs (failed
assertions) and normal termination execution is stopped.  For
non-fatal errors a message is printed to standard output if no unit is
given.  Only if the number of {\tt{}MAX{\char95}ERRORS} errors is reached, we
abort the program.  There are no further actions in the other cases,
but this may change.
\nwenddocs{}\nwbegincode{56}\moddef{Diagnostics: public}\plusendmoddef
  public :: msg_terminate
  public :: msg_bug, msg_fatal, msg_error, msg_warning
  public :: msg_message, msg_result, msg_debug
\nwendcode{}\nwbegincode{57}\moddef{Diagnostics: procedures}\plusendmoddef
  subroutine msg_terminate (string, unit, quit_code)
    integer, intent(in), optional :: unit
    character(len=*), intent(in), optional :: string
    integer, intent(in), optional :: quit_code
    integer(c_int) :: return_code
    call release_term_signals ()
    if (present (quit_code)) then
       return_code = quit_code
    else
       return_code = 0
    end if
    if (present (string)) &
         call message_print (MESSAGE, string, unit=unit)
    call msg_summary (unit)
    if (return_code == 0 .and. expect_failures /= 0) then
       return_code = 5
       call message_print (MESSAGE, &
            "WHIZARD run finished with 'expect' failure(s).", unit=unit)
    else if (return_code == 7) then
       call message_print (MESSAGE, &
            "WHIZARD run finished with failed self-test.", unit=unit)
    else
       call message_print (MESSAGE, "WHIZARD run finished.", unit=unit)
    end if
    call message_print (0, &
         "|=============================================================================|", unit=unit)
    call logfile_final ()
    call msg_list_clear ()
    if (return_code /= 0) then
       call exit (return_code)
    else 
       !!! Should implement WHIZARD exit code (currently only via C)
       ! stop
       call exit (0)
    end if
  end subroutine msg_terminate

  subroutine msg_bug (string, arr, unit)
    integer, intent(in), optional :: unit
    character(len=*), intent(in), optional :: string
    type(string_t), dimension(:), intent(in), optional :: arr
    logical, pointer :: crash_ptr
    call message_print (BUG, string, arr, unit)
    call msg_summary (unit)
    select case (handle_fatal_errors)
    case (TERM_EXIT)
       call message_print (TERMINATE, "WHIZARD run aborted.", unit=unit)
       call exit (-1_c_int)
    case (TERM_CRASH)
       print *, "*** Intentional crash ***"
       crash_ptr => null ()
       print *, crash_ptr
    end select
    stop "WHIZARD run aborted."
  end subroutine msg_bug

  recursive subroutine msg_fatal (string, arr, unit)
    integer, intent(in), optional :: unit
    character(len=*), intent(in), optional :: string
    type(string_t), dimension(:), intent(in), optional :: arr
    logical, pointer :: crash_ptr
    if (mask_fatal_errors) then
       call msg_error (string, arr, unit)
    else
       call message_print (FATAL, string, arr, unit)
       call msg_summary (unit)
       select case (handle_fatal_errors)
       case (TERM_EXIT)
          call message_print (TERMINATE, "WHIZARD run aborted.", unit=unit)
          call exit (1_c_int)
       case (TERM_CRASH)
          print *, "*** Intentional crash ***"
          crash_ptr => null ()
          print *, crash_ptr
       end select
       stop "WHIZARD run aborted."
    end if
  end subroutine msg_fatal

  subroutine msg_error (string, arr, unit)
    integer, intent(in), optional :: unit
    character(len=*), intent(in), optional :: string
    type(string_t), dimension(:), intent(in), optional :: arr
    call message_print (ERROR, string, arr, unit)
    if (msg_count(ERROR) >= MAX_ERRORS) then
       mask_fatal_errors = .false.
       call msg_fatal (" Too many errors encountered.")
    else if (.not.present(unit) .and. .not.mask_fatal_errors)  then
       call message_print (MESSAGE, "            (WHIZARD run continues)")
    end if
  end subroutine msg_error

  subroutine msg_warning (string, arr, unit)
    integer, intent(in), optional :: unit
    character(len=*), intent(in), optional :: string
    type(string_t), dimension(:), intent(in), optional :: arr
    call message_print (WARNING, string, arr, unit)
  end subroutine msg_warning

  subroutine msg_message (string, unit, arr, logfile)
    integer, intent(in), optional :: unit
    character(len=*), intent(in), optional :: string
    type(string_t), dimension(:), intent(in), optional :: arr
    logical, intent(in), optional :: logfile
    call message_print (MESSAGE, string, arr, unit, logfile)
  end subroutine msg_message

  subroutine msg_result (string, arr, unit, logfile)
    integer, intent(in), optional :: unit
    character(len=*), intent(in), optional :: string
    type(string_t), dimension(:), intent(in), optional :: arr
    logical, intent(in), optional :: logfile
    call message_print (RESULT, string, arr, unit, logfile)
  end subroutine msg_result

  subroutine msg_debug (string, arr, unit)
    integer, intent(in), optional :: unit
    character(len=*), intent(in), optional :: string
    type(string_t), dimension(:), intent(in), optional :: arr
    call message_print (DEBUG, string, arr, unit)
  end subroutine msg_debug

\eatline
\nwendcode{}\nwbegindocs{58}Interface to the standard clib exit function
\nwenddocs{}\nwbegincode{59}\moddef{Diagnostics: interfaces}\endmoddef
  interface
     subroutine exit (status) bind (C)
       use iso_c_binding !NODEP!
       integer(c_int), value :: status
     end subroutine exit
  end interface

\eatline
\nwendcode{}\nwbegindocs{60}Print the WHIZARD banner:
\nwenddocs{}\nwbegincode{61}\moddef{Diagnostics: public}\plusendmoddef
  public :: msg_banner
\nwendcode{}\nwbegincode{62}\moddef{Diagnostics: procedures}\plusendmoddef
  subroutine msg_banner (unit)
    integer, intent(in), optional :: unit
    integer :: n_proc
    call message_print (0, "|=============================================================================|", unit=unit)
    call message_print (0, "|                                                                             |", unit=unit)
    call message_print (0, "|    WW             WW  WW   WW  WW  WWWWWW      WW      WWWWW    WWWW        |", unit=unit)
    call message_print (0, "|     WW    WW     WW   WW   WW  WW     WW      WWWW     WW  WW   WW  WW      |", unit=unit)
    call message_print (0, "|      WW  WW WW  WW    WWWWWWW  WW    WW      WW  WW    WWWWW    WW   WW     |", unit=unit)
    call message_print (0, "|       WWWW   WWWW     WW   WW  WW   WW      WWWWWWWW   WW  WW   WW  WW      |", unit=unit)
    call message_print (0, "|        WW     WW      WW   WW  WW  WWWWWW  WW      WW  WW   WW  WWWW        |", unit=unit)
    call message_print (0, "|                                                                             |", unit=unit)
    call message_print (0, "|                                                                             |", unit=unit)
    call message_print (0, "|                                        W                                    |", unit=unit)
    call message_print (0, "|                                       sW                                    |", unit=unit)
    call message_print (0, "|                                       WW                                    |", unit=unit)
    call message_print (0, "|                                      sWW                                    |", unit=unit)
    call message_print (0, "|                                      WWW                                    |", unit=unit)
    call message_print (0, "|                                     wWWW                                    |", unit=unit)
    call message_print (0, "|                                    wWWWW                                    |", unit=unit)
    call message_print (0, "|                                    WW WW                                    |", unit=unit)
    call message_print (0, "|                                    WW WW                                    |", unit=unit)
    call message_print (0, "|                                   wWW WW                                    |", unit=unit)
    call message_print (0, "|                                  wWW  WW                                    |", unit=unit)
    call message_print (0, "|                                  WW   WW                                    |", unit=unit)
    call message_print (0, "|                                  WW   WW                                    |", unit=unit)
    call message_print (0, "|                                 WW    WW                                    |", unit=unit)
    call message_print (0, "|                                 WW    WW                                    |", unit=unit)
    call message_print (0, "|                                WW     WW                                    |", unit=unit)
    call message_print (0, "|                                WW     WW                                    |", unit=unit)
    call message_print (0, "|           wwwwww              WW      WW                                    |", unit=unit)
    call message_print (0, "|              WWWWWww          WW      WW                                    |", unit=unit)
    call message_print (0, "|                 WWWWWwwwww   WW       WW                                    |", unit=unit)
    call message_print (0, "|                     wWWWwwwwwWW       WW                                    |", unit=unit)
    call message_print (0, "|                 wWWWWWWWWWWwWWW       WW                                    |", unit=unit)
    call message_print (0, "|                wWWWWW       wW        WWWWWWW                               |", unit=unit)
    call message_print (0, "|                  WWWW       wW        WW  wWWWWWWWwww                       |", unit=unit)
    call message_print (0, "|                   WWWW                      wWWWWWWWwwww                    |", unit=unit)
    call message_print (0, "|                     WWWW                      WWWW     WWw                  |", unit=unit)
    call message_print (0, "|                       WWWWww                   WWWW                         |", unit=unit)
    call message_print (0, "|                           WWWwwww              WWWW                         |", unit=unit)
    call message_print (0, "|                               wWWWWwww       wWWWWW                         |", unit=unit)
    call message_print (0, "|                                     WwwwwwwwwWWW                            |", unit=unit)
    call message_print (0, "|                                                                             |", unit=unit)
    call message_print (0, "|                                                                             |", unit=unit)
    call message_print (0, "|                                                                             |", unit=unit)
    call message_print (0, "|  by:   Wolfgang Kilian, Thorsten Ohl, Juergen Reuter                        |", unit=unit)
    call message_print (0, "|        with contributions from Christian Speckner                           |", unit=unit)
    call message_print (0, "|        Contact: <whizard@desy.de>                                           |", unit=unit)
    call message_print (0, "|                                                                             |", unit=unit)
    call message_print (0, "|  if you use WHIZARD please cite:                                            |", unit=unit)   
    call message_print (0, "|        W. Kilian, T. Ohl, J. Reuter,  Eur.Phys.J.C71 (2011) 1742            |", unit=unit)
    call message_print (0, "|                                          [arXiv: 0708.4233 [hep-ph]]        |", unit=unit)   
    call message_print (0, "|        M. Moretti, T. Ohl, J. Reuter, arXiv: hep-ph/0102195                 |", unit=unit)
    call message_print (0, "|                                                                             |", unit=unit)
    call message_print (0, "|=============================================================================|", unit=unit)
    call message_print (0, "|                               WHIZARD " // WHIZARD_VERSION, unit=unit)
    call message_print (0, "|=============================================================================|", unit=unit)
  end subroutine msg_banner

\eatline
\nwendcode{}\nwbegindocs{63}\nwdocspar
\subsection{Logfile}
All screen output should be duplicated in the logfile, unless
requested otherwise.
\nwenddocs{}\nwbegincode{64}\moddef{Diagnostics: public}\plusendmoddef
  public :: logging
\nwendcode{}\nwbegincode{65}\moddef{Diagnostics: variables}\plusendmoddef
  integer, save :: log_unit = -1
  logical, target, save :: logging = .false.
\nwendcode{}\nwbegincode{66}\moddef{Diagnostics: public}\plusendmoddef
  public :: logfile_init
\nwendcode{}\nwbegincode{67}\moddef{Diagnostics: procedures}\plusendmoddef
  subroutine logfile_init (filename)
    type(string_t), intent(in) :: filename
    call msg_message ("Writing log to '" // char (filename) // "'")
    if (.not. logging)  call msg_message ("(Logging turned off.)")
    log_unit = free_unit ()
    open (file = char (filename), unit = log_unit, &
          action = "write", status = "replace")
  end subroutine logfile_init

\eatline
\nwendcode{}\nwbegincode{68}\moddef{Diagnostics: public}\plusendmoddef
  public :: logfile_final
\nwendcode{}\nwbegincode{69}\moddef{Diagnostics: procedures}\plusendmoddef
  subroutine logfile_final ()
    if (log_unit >= 0) then
       close (log_unit)
       log_unit = -1
    end if
  end subroutine logfile_final

\eatline
\nwendcode{}\nwbegindocs{70}This returns the valid logfile unit only if the default is write to
screen, and if {\tt{}logfile} is not set false.
\nwenddocs{}\nwbegincode{71}\moddef{Diagnostics: public}\plusendmoddef
  public :: logfile_unit
\nwendcode{}\nwbegincode{72}\moddef{Diagnostics: procedures}\plusendmoddef
  function logfile_unit (unit, logfile)
    integer :: logfile_unit
    integer, intent(in), optional :: unit
    logical, intent(in), optional :: logfile
    if (logging) then
       if (present (unit)) then
          if (unit == output_unit) then
             logfile_unit = log_unit
          else
             logfile_unit = -1
          end if
       else if (present (logfile)) then
          if (logfile) then
             logfile_unit = log_unit
          else
             logfile_unit = -1
          end if
       else
          logfile_unit = log_unit
       end if
    else
       logfile_unit = -1
    end if
  end function logfile_unit

\eatline
\nwendcode{}\nwbegindocs{73}\nwdocspar
\subsection{Checking values}
The {\tt{}expect} function does not just check a value for correctness
(actually, it checks if a logical expression is true); it records its
result here.  If failures are present when the program terminates, the
exit code is nonzero.
\nwenddocs{}\nwbegincode{74}\moddef{Diagnostics: variables}\plusendmoddef
  integer, save :: expect_total = 0
  integer, save :: expect_failures = 0

\eatline
\nwendcode{}\nwbegincode{75}\moddef{Diagnostics: public}\plusendmoddef
  public :: expect_record
\nwendcode{}\nwbegincode{76}\moddef{Diagnostics: procedures}\plusendmoddef
  subroutine expect_record (success)
    logical, intent(in) :: success
    expect_total = expect_total + 1
    if (.not. success)  expect_failures = expect_failures + 1
  end subroutine expect_record

\eatline
\nwendcode{}\nwbegincode{77}\moddef{Diagnostics: public}\plusendmoddef
  public :: expect_clear
\nwendcode{}\nwbegincode{78}\moddef{Diagnostics: procedures}\plusendmoddef
  subroutine expect_clear ()
    expect_total = 0
    expect_failures = 0
  end subroutine expect_clear

\eatline
\nwendcode{}\nwbegincode{79}\moddef{Diagnostics: public}\plusendmoddef
  public :: expect_summary
\nwendcode{}\nwbegincode{80}\moddef{Diagnostics: procedures}\plusendmoddef
  subroutine expect_summary (unit, force)
    integer, intent(in), optional :: unit
    logical, intent(in), optional :: force
    logical :: force_output
    force_output = .false.;  if (present (force))  force_output = force
    if (expect_total /= 0 .or. force_output) then
       call msg_message ("Summary of value checks:", unit)
       write (msg_buffer, "(2x,A,1x,I0,1x,A,1x,A,1x,I0)") &
            "Failures:", expect_failures, "/", "Total:", expect_total
       call msg_message (unit=unit)
    end if
  end subroutine expect_summary

\eatline
\nwendcode{}\nwbegindocs{81}Helpers for converting integers into strings with minimal length.
\nwenddocs{}\nwbegincode{82}\moddef{Diagnostics: public}\plusendmoddef
  public :: int2string
  public :: int2char
  public :: int2fixed
\nwendcode{}\nwbegincode{83}\moddef{Diagnostics: procedures}\plusendmoddef
  pure function int2fixed (i) result (c)
    integer, intent(in) :: i
    character(200) :: c
    c = ""
    write (c, *) i
    c = adjustl (c)
  end function int2fixed

  pure function int2string (i) result (s)
    integer, intent(in) :: i
    type (string_t) :: s
    s = trim (int2fixed (i))
  end function int2string

  pure function int2char (i) result (c)
    integer, intent(in) :: i
    character(len (trim (int2fixed (i)))) :: c
    c = int2fixed (i)
  end function int2char

\eatline
\nwendcode{}\nwbegindocs{84}Dito for reals.
\nwenddocs{}\nwbegincode{85}\moddef{Diagnostics: public}\plusendmoddef
  public :: real2string
  public :: real2char
  public :: real2fixed
\nwendcode{}\nwbegincode{86}\moddef{Diagnostics: interfaces}\plusendmoddef
  interface real2string
     module procedure real2string_list, real2string_fmt
  end interface
  interface real2char
     module procedure real2char_list, real2char_fmt
  end interface
\nwendcode{}\nwbegincode{87}\moddef{Diagnostics: procedures}\plusendmoddef
  pure function real2fixed (x, fmt) result (c)
    real(default), intent(in) :: x
    character(*), intent(in), optional :: fmt
    character(200) :: c
    c = ""
    write (c, *) x
    c = adjustl (c)
  end function real2fixed

  pure function real2fixed_fmt (x, fmt) result (c)
    real(default), intent(in) :: x
    character(*), intent(in) :: fmt
    character(200) :: c
    c = ""
    write (c, fmt)  x
    c = adjustl (c)
  end function real2fixed_fmt

  pure function real2string_list (x) result (s)
    real(default), intent(in) :: x
    type(string_t) :: s
    s = trim (real2fixed (x))
  end function real2string_list

  pure function real2string_fmt (x, fmt) result (s)
    real(default), intent(in) :: x
    character(*), intent(in) :: fmt
    type(string_t) :: s
    s = trim (real2fixed_fmt (x, fmt))
  end function real2string_fmt

  pure function real2char_list (x) result (c)
    real(default), intent(in) :: x
    character(len_trim (real2fixed (x))) :: c
    c = real2fixed (x)
  end function real2char_list

  pure function real2char_fmt (x, fmt) result (c)
    real(default), intent(in) :: x
    character(*), intent(in) :: fmt
    character(len_trim (real2fixed_fmt (x, fmt))) :: c
    c = real2fixed_fmt (x, fmt)
  end function real2char_fmt

\eatline
\nwendcode{}\nwbegindocs{88}Dito for complex values; we do not use the slightly ugly FORTRAN output form
here but instead introduce our own. Ifort and Portland seem to have problems 
with this, therefore temporarily disable it.   
%
\nwenddocs{}\nwbegincode{89}\moddef{CCC Diagnostics: public}\endmoddef
   public :: cmplx2string
   public :: cmplx2char
\nwendcode{}\nwbegincode{90}\moddef{CCC Diagnostics: procedures}\endmoddef
   pure function cmplx2string (x) result (s)
     complex(default), intent(in) :: x
     type(string_t) :: s
     s = real2string (real (x, default))
     if (aimag (x) /= 0) s = s // " + " // real2string (aimag (x)) // " I"
   end function cmplx2string

   pure function cmplx2char (x) result (c)
     complex(default), intent(in) :: x
     character(len (char (cmplx2string (x)))) :: c
     c = char (cmplx2string (x))
   end function cmplx2char
   
\eatline
\nwendcode{}\nwbegindocs{91}\nwdocspar
\subsection{Suppression of numerical noise}
\nwenddocs{}\nwbegincode{92}\moddef{Diagnostics: public}\plusendmoddef
  public :: pacify
\nwendcode{}\nwbegincode{93}\moddef{Diagnostics: interfaces}\plusendmoddef
  interface pacify
     module procedure pacify_real_default
     module procedure pacify_complex_default
  end interface pacify
  
\nwendcode{}\nwbegindocs{94}\nwdocspar
\nwenddocs{}\nwbegincode{95}\moddef{Diagnostics: procedures}\plusendmoddef
  elemental subroutine pacify_real_default (x, tolerance)
    real(default), intent(inout) :: x
    real(default), intent(in) :: tolerance
    if (abs (x) < tolerance)  x = 0
  end subroutine pacify_real_default
  
  elemental subroutine pacify_complex_default (x, tolerance)
    complex(default), intent(inout) :: x
    real(default), intent(in) :: tolerance
    if (abs (real (x)) < tolerance)   &
         x = cmplx (0._default, aimag (x), kind=default)
    if (abs (aimag (x)) < tolerance)  &
         x = cmplx (real (x), 0._default, kind=default)
  end subroutine pacify_complex_default  

\eatline
\nwendcode{}\nwbegindocs{96}\nwdocspar
\subsection{Signal handling}
Killing the program by external signals may leave the files written by it in
an undefined state.  This can be avoided by catching signals and deferring
program termination.  Instead of masking only critical sections, we choose
to mask signals globally (done in the main program) and terminate the program
at predefined checkpoints only.  Checkpoints are after each command, within
the sampling function (so the program can be terminated after each event),
and after each iteration in the phase-space generation algorithm.

Signal handling is done via a C interface to the {\tt{}sigaction} system call.
When a signal is raised that has been masked by the handler, the corresponding
variable is set to the value of the signal.  The variables are visible from
the C signal handler.

The signal SIGINT is for keyboard interrupt (ctrl-C), SIGTERM is for system
interrupt, e.g., at shutdown.  The SIGXCPU and SIGXFSZ signals may be issued
by batch systems.
\nwenddocs{}\nwbegincode{97}\moddef{Diagnostics: public}\plusendmoddef
  public :: wo_sigint
  public :: wo_sigterm
  public :: wo_sigxcpu
  public :: wo_sigxfsz

\nwendcode{}\nwbegincode{98}\moddef{Diagnostics: variables}\plusendmoddef
  integer(c_int), bind(C), volatile :: wo_sigint = 0
  integer(c_int), bind(C), volatile :: wo_sigterm = 0
  integer(c_int), bind(C), volatile :: wo_sigxcpu = 0
  integer(c_int), bind(C), volatile :: wo_sigxfsz = 0

\eatline
\nwendcode{}\nwbegindocs{99}Here are the interfaces to the C functions.  The routine
{\tt{}mask{\char95}term{\char95}signals} forces termination signals to be delayed.
{\tt{}release{\char95}term{\char95}signals} restores normal behavior.  However, the program can be
terminated anytime by calling {\tt{}terminate{\char95}now{\char95}if{\char95}signal} which inspects the
signals and terminates the program if requested..
\nwenddocs{}\nwbegincode{100}\moddef{Diagnostics: public}\plusendmoddef
  public :: mask_term_signals
\nwendcode{}\nwbegincode{101}\moddef{Diagnostics: procedures}\plusendmoddef
  subroutine mask_term_signals ()
    integer(c_int) :: status
    logical :: ok
    wo_sigint = 0
    ok = wo_mask_sigint () == 0
    if (.not. ok)  call msg_error ("Masking SIGINT failed")
    wo_sigterm = 0
    ok = wo_mask_sigterm () == 0
    if (.not. ok)  call msg_error ("Masking SIGTERM failed")
    wo_sigxcpu = 0
    ok = wo_mask_sigxcpu () == 0
    if (.not. ok)  call msg_error ("Masking SIGXCPU failed")
    wo_sigxfsz = 0
    ok = wo_mask_sigxfsz () == 0
    if (.not. ok)  call msg_error ("Masking SIGXFSZ failed")
  end subroutine mask_term_signals

\eatline
\nwendcode{}\nwbegincode{102}\moddef{Diagnostics: interfaces}\plusendmoddef
  interface
     integer(c_int) function wo_mask_sigint () bind(C)
       import
     end function wo_mask_sigint
  end interface
  interface
     integer(c_int) function wo_mask_sigterm () bind(C)
       import
     end function wo_mask_sigterm
  end interface
  interface
     integer(c_int) function wo_mask_sigxcpu () bind(C)
       import
     end function wo_mask_sigxcpu
  end interface
  interface
     integer(c_int) function wo_mask_sigxfsz () bind(C)
       import
     end function wo_mask_sigxfsz
  end interface

\eatline
\nwendcode{}\nwbegincode{103}\moddef{Diagnostics: public}\plusendmoddef
  public :: release_term_signals
\nwendcode{}\nwbegincode{104}\moddef{Diagnostics: procedures}\plusendmoddef
  subroutine release_term_signals ()
    integer(c_int) :: status
    logical :: ok
    ok = wo_release_sigint () == 0
    if (.not. ok)  call msg_error ("Releasing SIGINT failed")
    ok = wo_release_sigterm () == 0
    if (.not. ok)  call msg_error ("Releasing SIGTERM failed")
    ok = wo_release_sigxcpu () == 0
    if (.not. ok)  call msg_error ("Releasing SIGXCPU failed")
    ok = wo_release_sigxfsz () == 0
    if (.not. ok)  call msg_error ("Releasing SIGXFSZ failed")
  end subroutine release_term_signals

\eatline
\nwendcode{}\nwbegincode{105}\moddef{Diagnostics: interfaces}\plusendmoddef
  interface
     integer(c_int) function wo_release_sigint () bind(C)
       import
     end function wo_release_sigint
  end interface
  interface
     integer(c_int) function wo_release_sigterm () bind(C)
       import
     end function wo_release_sigterm
  end interface
  interface
     integer(c_int) function wo_release_sigxcpu () bind(C)
       import
     end function wo_release_sigxcpu
  end interface
  interface
     integer(c_int) function wo_release_sigxfsz () bind(C)
       import
     end function wo_release_sigxfsz
  end interface

\eatline
\eatline
\nwendcode{}\nwbegincode{106}\moddef{Diagnostics: public}\plusendmoddef
  public :: signal_is_pending
\nwendcode{}\nwbegincode{107}\moddef{Diagnostics: procedures}\plusendmoddef
  function signal_is_pending () result (flag)
    logical :: flag
    flag = &
         wo_sigint /= 0 .or. &
         wo_sigterm /= 0 .or. &
         wo_sigxcpu /= 0 .or. &
         wo_sigxfsz /= 0
  end function signal_is_pending
  
\eatline
\nwendcode{}\nwbegincode{108}\moddef{Diagnostics: public}\plusendmoddef
  public :: terminate_now_if_signal
\nwendcode{}\nwbegincode{109}\moddef{Diagnostics: procedures}\plusendmoddef
  subroutine terminate_now_if_signal ()
    if (wo_sigint /= 0) then
       call msg_terminate ("Signal SIGINT (keyboard interrupt) received.", &
          quit_code=int (wo_sigint))
    else if (wo_sigterm /= 0) then
       call msg_terminate ("Signal SIGTERM (termination signal) received.", &
          quit_code=int (wo_sigterm))
    else if (wo_sigxcpu /= 0) then
       call msg_terminate ("Signal SIGXCPU (CPU time limit exceeded) received.", &
          quit_code=int (wo_sigxcpu))
    else if (wo_sigxfsz /= 0) then
       call msg_terminate ("Signal SIGXFSZ (file size limit exceeded) received.", &
          quit_code=int (wo_sigxfsz))
    end if
  end subroutine terminate_now_if_signal    

\eatline
\nwendcode{}\nwbegindocs{110}\nwdocspar
%%%%%%%%%%%%%%%%%%%%%%%%%%%%%%%%%%%%%%%%%%%%%%%%%%%%%%%%%%%%%%%%%%%%%%%%
\section{Operating-system interface}
For specific purposes, we need direct access to the OS (system
calls).  This is, of course, system dependent.  The current version is
valid for GNU/Linux; we expect to use a preprocessor for this module
if different OSs are to be supported.

The current implementation lacks error handling.
\nwenddocs{}\nwbegincode{111}\moddef{\code{}os{\char95}interface.f90\edoc{}}\endmoddef
\LA{}File header\RA{}

module os_interface

  use, intrinsic :: iso_c_binding !NODEP!

\LA{}Use strings\RA{}
  use io_units
  use unit_tests
  use diagnostics
  use system_defs, only: DLERROR_LEN, ENVVAR_LEN
  use system_dependencies
  
\LA{}Standard module head\RA{}

\LA{}OS interface: public\RA{}

\LA{}OS interface: types\RA{}

\LA{}OS interface: interfaces\RA{}

contains

\LA{}OS interface: procedures\RA{}

\LA{}OS interface: tests\RA{}

end module os_interface
\eatline
\nwendcode{}\nwbegindocs{112}\nwdocspar
\subsection{Path variables}
This is a transparent container for storing user-defined path variables.
\nwenddocs{}\nwbegincode{113}\moddef{OS interface: public}\endmoddef
  public :: paths_t
\nwendcode{}\nwbegincode{114}\moddef{OS interface: types}\endmoddef
  type :: paths_t
     type(string_t) :: prefix
     type(string_t) :: exec_prefix
     type(string_t) :: bindir
     type(string_t) :: libdir
     type(string_t) :: includedir
     type(string_t) :: datarootdir
     type(string_t) :: localprefix
     type(string_t) :: libtool
     type(string_t) :: lhapdfdir
  end type paths_t

\eatline
\nwendcode{}\nwbegincode{115}\moddef{OS interface: public}\plusendmoddef
  public :: paths_init
\nwendcode{}\nwbegincode{116}\moddef{OS interface: procedures}\endmoddef
  subroutine paths_init (paths)
    type(paths_t), intent(out) :: paths
    paths%prefix = ""
    paths%exec_prefix = ""
    paths%bindir = ""
    paths%libdir = ""
    paths%includedir = ""
    paths%datarootdir = ""
    paths%localprefix = ""
    paths%libtool = ""
    paths%lhapdfdir = ""
  end subroutine paths_init

\eatline
\nwendcode{}\nwbegindocs{117}\nwdocspar
\subsection{System dependencies}
We store all potentially system- and user/run-dependent data in a
transparent container.  This includes compiler/linker names and flags,
file extensions, etc.
\nwenddocs{}\nwbegincode{118}\moddef{OS interface: public}\plusendmoddef
  public :: os_data_t
\nwendcode{}\nwbegincode{119}\moddef{OS interface: types}\plusendmoddef
  type :: os_data_t
     logical :: use_libtool
     logical :: use_testfiles
     type(string_t) :: fc
     type(string_t) :: fcflags
     type(string_t) :: fcflags_pic
     type(string_t) :: fc_src_ext
     type(string_t) :: cc
     type(string_t) :: cflags
     type(string_t) :: cflags_pic
     type(string_t) :: obj_ext
     type(string_t) :: ld
     type(string_t) :: ldflags
     type(string_t) :: ldflags_so
     type(string_t) :: ldflags_static
     type(string_t) :: ldflags_hepmc
     type(string_t) :: ldflags_hoppet
     type(string_t) :: shlib_ext
     type(string_t) :: makeflags
     type(string_t) :: prefix
     type(string_t) :: exec_prefix
     type(string_t) :: bindir
     type(string_t) :: libdir
     type(string_t) :: includedir
     type(string_t) :: datarootdir
     type(string_t) :: whizard_omega_binpath
     type(string_t) :: whizard_includes
     type(string_t) :: whizard_ldflags
     type(string_t) :: whizard_libtool
     type(string_t) :: whizard_modelpath
     type(string_t) :: whizard_models_libpath
     type(string_t) :: whizard_susypath
     type(string_t) :: whizard_gmlpath
     type(string_t) :: whizard_cutspath
     type(string_t) :: whizard_texpath
     type(string_t) :: whizard_testdatapath
     type(string_t) :: whizard_modelpath_local
     type(string_t) :: whizard_models_libpath_local
     type(string_t) :: whizard_omega_binpath_local
     type(string_t) :: whizard_circe2path
     type(string_t) :: whizard_beamsimpath
     type(string_t) :: whizard_mulipath
     type(string_t) :: pdf_builtin_datapath
     logical :: event_analysis = .false.
     logical :: event_analysis_ps  = .false.
     logical :: event_analysis_pdf = .false.
     type(string_t) :: latex
     type(string_t) :: mpost
     type(string_t) :: gml
     type(string_t) :: dvips
     type(string_t) :: ps2pdf
     type(string_t) :: gosampath
     type(string_t) :: golempath
     type(string_t) :: formpath
     type(string_t) :: qgrafpath
     type(string_t) :: ninjapath
     type(string_t) :: samuraipath
  end type os_data_t

\eatline
\nwendcode{}\nwbegindocs{120}Since all are allocatable strings, explicit initialization is
necessary.
\nwenddocs{}\nwbegincode{121}\moddef{System defs: public parameters}\plusendmoddef
  integer, parameter, public :: ENVVAR_LEN = 1000
\eatline
\nwendcode{}\nwbegincode{122}\moddef{OS interface: public}\plusendmoddef
  public :: os_data_init
\nwendcode{}\nwbegincode{123}\moddef{OS interface: procedures}\plusendmoddef
  subroutine os_data_init (os_data, paths)
    type(os_data_t), intent(out) :: os_data
    type(paths_t), intent(in), optional :: paths
    character(len=ENVVAR_LEN) :: home
    type(string_t) :: localprefix, local_includes
    os_data%use_libtool = .true.
    inquire (file = "TESTFLAG", exist = os_data%use_testfiles)
    call get_environment_variable ("HOME", home)
    if (present(paths)) then
       if (paths%localprefix == "") then
          localprefix = trim (home) // "/.whizard"
       else
          localprefix = paths%localprefix
       end if
    else
       localprefix = trim (home) // "/.whizard"
    end if
    local_includes = localprefix // "/lib/whizard/mod/models"
    os_data%whizard_modelpath_local = localprefix // "/share/whizard/models"
    os_data%whizard_models_libpath_local = localprefix // "/lib/whizard/models"
    os_data%whizard_omega_binpath_local = localprefix // "/bin"
    os_data%fc             = DEFAULT_FC
    os_data%fcflags        = DEFAULT_FCFLAGS
    os_data%fcflags_pic    = DEFAULT_FCFLAGS_PIC
    os_data%fc_src_ext     = DEFAULT_FC_SRC_EXT
    os_data%cc             = DEFAULT_CC
    os_data%cflags         = DEFAULT_CFLAGS
    os_data%cflags_pic     = DEFAULT_CFLAGS_PIC
    os_data%obj_ext        = DEFAULT_OBJ_EXT
    os_data%ld             = DEFAULT_LD
    os_data%ldflags        = DEFAULT_LDFLAGS
    os_data%ldflags_so     = DEFAULT_LDFLAGS_SO
    os_data%ldflags_static = DEFAULT_LDFLAGS_STATIC
    os_data%ldflags_hepmc  = DEFAULT_LDFLAGS_HEPMC
    os_data%ldflags_hoppet = DEFAULT_LDFLAGS_HOPPET
    os_data%shlib_ext      = DEFAULT_SHLIB_EXT
    os_data%makeflags      = DEFAULT_MAKEFLAGS
    os_data%prefix      = PREFIX
    os_data%exec_prefix = EXEC_PREFIX
    os_data%bindir      = BINDIR
    os_data%libdir      = LIBDIR
    os_data%includedir  = INCLUDEDIR
    os_data%datarootdir = DATAROOTDIR
    if (present (paths)) then
       if (paths%prefix      /= "")  os_data%prefix      = paths%prefix
       if (paths%exec_prefix /= "")  os_data%exec_prefix = paths%exec_prefix
       if (paths%bindir      /= "")  os_data%bindir      = paths%bindir
       if (paths%libdir      /= "")  os_data%libdir      = paths%libdir
       if (paths%includedir  /= "")  os_data%includedir  = paths%includedir
       if (paths%datarootdir /= "")  os_data%datarootdir = paths%datarootdir
    end if
    if (os_data%use_testfiles) then
       os_data%whizard_omega_binpath  = WHIZARD_TEST_OMEGA_BINPATH
       os_data%whizard_includes       = WHIZARD_TEST_INCLUDES
       os_data%whizard_ldflags        = WHIZARD_TEST_LDFLAGS
       os_data%whizard_libtool        = WHIZARD_LIBTOOL_TEST
       os_data%whizard_modelpath      = WHIZARD_TEST_MODELPATH
       os_data%whizard_models_libpath = WHIZARD_TEST_MODELS_LIBPATH
       os_data%whizard_susypath       = WHIZARD_TEST_SUSYPATH
       os_data%whizard_gmlpath        = WHIZARD_TEST_GMLPATH
       os_data%whizard_cutspath       = WHIZARD_TEST_CUTSPATH
       os_data%whizard_texpath        = WHIZARD_TEST_TEXPATH
       os_data%whizard_testdatapath   = WHIZARD_TEST_TESTDATAPATH
       os_data%whizard_circe2path     = WHIZARD_TEST_CIRCE2PATH
       os_data%whizard_beamsimpath    = WHIZARD_TEST_BEAMSIMPATH
       os_data%whizard_mulipath       = WHIZARD_TEST_MULIPATH
       os_data%pdf_builtin_datapath   = PDF_BUILTIN_TEST_DATAPATH
    else
       if (os_dir_exist (local_includes)) then
          os_data%whizard_includes = "-I" // local_includes // " "// &
             WHIZARD_INCLUDES 
       else
          os_data%whizard_includes = WHIZARD_INCLUDES
       end if
       os_data%whizard_omega_binpath  = WHIZARD_OMEGA_BINPATH
       os_data%whizard_ldflags        = WHIZARD_LDFLAGS
       os_data%whizard_libtool        = WHIZARD_LIBTOOL
       if(present(paths)) then
          if (paths%libtool /= "")  os_data%whizard_libtool = paths%libtool
       end if
       os_data%whizard_modelpath      = WHIZARD_MODELPATH
       os_data%whizard_models_libpath = WHIZARD_MODELS_LIBPATH
       os_data%whizard_susypath       = WHIZARD_SUSYPATH
       os_data%whizard_gmlpath        = WHIZARD_GMLPATH
       os_data%whizard_cutspath       = WHIZARD_CUTSPATH
       os_data%whizard_texpath        = WHIZARD_TEXPATH
       os_data%whizard_testdatapath   = WHIZARD_TESTDATAPATH
       os_data%whizard_circe2path     = WHIZARD_CIRCE2PATH
       os_data%whizard_beamsimpath    = WHIZARD_BEAMSIMPATH
       os_data%whizard_mulipath       = WHIZARD_MULIPATH
       os_data%pdf_builtin_datapath   = PDF_BUILTIN_DATAPATH
    end if
    os_data%event_analysis     = EVENT_ANALYSIS     == "yes"
    os_data%event_analysis_ps  = EVENT_ANALYSIS_PS  == "yes"
    os_data%event_analysis_pdf = EVENT_ANALYSIS_PDF == "yes"
    os_data%latex  = PRG_LATEX // " " // OPT_LATEX
    os_data%mpost  = PRG_MPOST // " " // OPT_MPOST
    os_data%gml    = os_data%whizard_gmlpath // "/gml" // " " // OPT_MPOST &
         // " " // "--gmldir " // os_data%whizard_gmlpath
    os_data%dvips  = PRG_DVIPS
    os_data%ps2pdf = PRG_PS2PDF
    call os_data_expand_paths (os_data)
    os_data%gosampath = GOSAM_DIR
    os_data%golempath = GOLEM_DIR
    os_data%formpath = FORM_DIR
    os_data%qgrafpath = QGRAF_DIR
    os_data%ninjapath = NINJA_DIR
    os_data%samuraipath = SAMURAI_DIR
  end subroutine os_data_init
    
\eatline
\nwendcode{}\nwbegindocs{124}Replace occurences of GNU path variables (such as {\tt{}{\char36}{\char123}prefix{\char125}}) by their
values.  Do this for all strings that could depend on them, and do the
replacement in reverse order, since the path variables may be defined in terms
of each other.    %% Fooling Noweb Emacs mode: $
\nwenddocs{}\nwbegincode{125}\moddef{OS interface: procedures}\plusendmoddef
  subroutine os_data_expand_paths (os_data)
    type(os_data_t), intent(inout) :: os_data
    integer, parameter :: N_VARIABLES = 6
    type(string_t), dimension(N_VARIABLES) :: variable, value
    variable(1) = "$\{prefix\}";       value(1) = os_data%prefix
    variable(2) = "$\{exec_prefix\}";  value(2) = os_data%exec_prefix
    variable(3) = "$\{bindir\}";       value(3) = os_data%bindir
    variable(4) = "$\{libdir\}";       value(4) = os_data%libdir
    variable(5) = "$\{includedir\}";   value(5) = os_data%includedir
    variable(6) = "$\{datarootdir\}";  value(6) = os_data%datarootdir
    call expand_paths (os_data%whizard_omega_binpath)
    call expand_paths (os_data%whizard_includes)
    call expand_paths (os_data%whizard_ldflags)
    call expand_paths (os_data%whizard_libtool)
    call expand_paths (os_data%whizard_modelpath)
    call expand_paths (os_data%whizard_models_libpath)
    call expand_paths (os_data%whizard_susypath)
    call expand_paths (os_data%whizard_gmlpath)
    call expand_paths (os_data%whizard_cutspath)
    call expand_paths (os_data%whizard_texpath)
    call expand_paths (os_data%whizard_testdatapath)
    call expand_paths (os_data%whizard_circe2path)
    call expand_paths (os_data%whizard_beamsimpath)
    call expand_paths (os_data%whizard_mulipath)
    call expand_paths (os_data%whizard_models_libpath_local)
    call expand_paths (os_data%whizard_modelpath_local)
    call expand_paths (os_data%whizard_omega_binpath_local)
    call expand_paths (os_data%pdf_builtin_datapath)
    call expand_paths (os_data%latex)
    call expand_paths (os_data%mpost)    
    call expand_paths (os_data%gml)
    call expand_paths (os_data%dvips)
    call expand_paths (os_data%ps2pdf)
  contains
    subroutine expand_paths (string)
      type(string_t), intent(inout) :: string
      integer :: i
      do i = N_VARIABLES, 1, -1
         string = replace (string, variable(i), value(i), every=.true.)
      end do
    end subroutine expand_paths
  end subroutine os_data_expand_paths

\eatline
\nwendcode{}\nwbegindocs{126}Write contents
\nwenddocs{}\nwbegincode{127}\moddef{OS interface: public}\plusendmoddef
  public :: os_data_write
\nwendcode{}\nwbegincode{128}\moddef{OS interface: procedures}\plusendmoddef
  subroutine os_data_write (os_data, unit)
    type(os_data_t), intent(in) :: os_data
    integer, intent(in), optional :: unit
    integer :: u
    u = given_output_unit (unit);  if (u < 0)  return
    write (u, "(A)")  "OS data:"
    write (u, *) "use_libtool    = ", os_data%use_libtool
    write (u, *) "use_testfiles  = ", os_data%use_testfiles
    write (u, *) "fc             = ", char (os_data%fc)
    write (u, *) "fcflags        = ", char (os_data%fcflags)
    write (u, *) "fcflags_pic    = ", char (os_data%fcflags_pic)
    write (u, *) "fc_src_ext     = ", char (os_data%fc_src_ext)
    write (u, *) "cc             = ", char (os_data%cc)
    write (u, *) "cflags         = ", char (os_data%cflags)
    write (u, *) "cflags_pic     = ", char (os_data%cflags_pic)
    write (u, *) "obj_ext        = ", char (os_data%obj_ext)
    write (u, *) "ld             = ", char (os_data%ld)
    write (u, *) "ldflags        = ", char (os_data%ldflags)
    write (u, *) "ldflags_so     = ", char (os_data%ldflags_so)
    write (u, *) "ldflags_static = ", char (os_data%ldflags_static)
    write (u, *) "ldflags_hepmc  = ", char (os_data%ldflags_hepmc)
    write (u, *) "ldflags_hoppet = ", char (os_data%ldflags_hoppet)    
    write (u, *) "shlib_ext      = ", char (os_data%shlib_ext)
    write (u, *) "makeflags      = ", char (os_data%makeflags)
    write (u, *) "prefix         = ", char (os_data%prefix)
    write (u, *) "exec_prefix    = ", char (os_data%exec_prefix)
    write (u, *) "bindir         = ", char (os_data%bindir)
    write (u, *) "libdir         = ", char (os_data%libdir)
    write (u, *) "includedir     = ", char (os_data%includedir)
    write (u, *) "datarootdir    = ", char (os_data%datarootdir)
    write (u, *) "whizard_omega_binpath  = ", &
         char (os_data%whizard_omega_binpath)
    write (u, *) "whizard_includes       = ", char (os_data%whizard_includes)
    write (u, *) "whizard_ldflags        = ", char (os_data%whizard_ldflags)
    write (u, *) "whizard_libtool        = ", char (os_data%whizard_libtool)
    write (u, *) "whizard_modelpath      = ", &
         char (os_data%whizard_modelpath)
    write (u, *) "whizard_models_libpath = ", &
         char (os_data%whizard_modelpath)
    write (u, *) "whizard_susypath       = ", char (os_data%whizard_susypath)
    write (u, *) "whizard_gmlpath        = ", char (os_data%whizard_gmlpath)
    write (u, *) "whizard_cutspath       = ", char (os_data%whizard_cutspath)
    write (u, *) "whizard_texpath        = ", char (os_data%whizard_texpath)
    write (u, *) "whizard_circe2path     = ", char (os_data%whizard_circe2path)
    write (u, *) "whizard_beamsimpath    = ", char (os_data%whizard_beamsimpath)
    write (u, *) "whizard_mulipath    = ", char (os_data%whizard_mulipath)
    write (u, *) "whizard_testdatapath  = ", &
         char (os_data%whizard_testdatapath)
    write (u, *) "whizard_modelpath_local      = ", &
         char (os_data%whizard_modelpath_local)
    write (u, *) "whizard_models_libpath_local = ", &
         char (os_data%whizard_models_libpath_local)
    write (u, *) "whizard_omega_binpath_local  = ", &
         char (os_data%whizard_omega_binpath_local)
    write (u, *) "event_analysis     = ", os_data%event_analysis
    write (u, *) "event_analysis_ps  = ", os_data%event_analysis_ps
    write (u, *) "event_analysis_pdf = ", os_data%event_analysis_pdf
    write (u, *) "latex  = ", char (os_data%latex)
    write (u, *) "mpost  = ", char (os_data%mpost)    
    write (u, *) "gml    = ", char (os_data%gml)
    write (u, *) "dvips  = ", char (os_data%dvips)
    write (u, *) "ps2pdf = ", char (os_data%ps2pdf)
    if (os_data%gosampath /= "") then
       write (u, *) "gosam   = ", char (os_data%gosampath)
       write (u, *) "golem   = ", char (os_data%golempath)
       write (u, *) "form    = ", char (os_data%formpath)
       write (u, *) "qgraf   = ", char (os_data%qgrafpath)
       write (u, *) "ninja   = ", char (os_data%ninjapath)
       write (u, *) "samurai = ", char (os_data%samuraipath)
    end if
  end subroutine os_data_write

\eatline
\nwendcode{}\nwbegindocs{129}\nwdocspar
\subsection{Dynamic linking}
We define a type that holds the filehandle for a dynamically linked
library (shared object), together with functions to open and close the
library, and to access functions in this library.
\nwenddocs{}\nwbegincode{130}\moddef{OS interface: public}\plusendmoddef
  public :: dlaccess_t
\nwendcode{}\nwbegincode{131}\moddef{OS interface: types}\plusendmoddef
  type :: dlaccess_t
     private
     type(string_t) :: filename
     type(c_ptr) :: handle = c_null_ptr
     logical :: is_open = .false.
     logical :: has_error = .false.
     type(string_t) :: error
   contains
   \LA{}OS interface: dlaccess: TBP\RA{}
  end type dlaccess_t

\eatline
\nwendcode{}\nwbegindocs{132}Output.  This is called by the output routine for the process
library.
\nwenddocs{}\nwbegincode{133}\moddef{OS interface: dlaccess: TBP}\endmoddef
  procedure :: write => dlaccess_write
\nwendcode{}\nwbegincode{134}\moddef{OS interface: procedures}\plusendmoddef
  subroutine dlaccess_write (object, unit)
    class(dlaccess_t), intent(in) :: object
    integer, intent(in) :: unit
    write (unit, "(1x,A)")  "DL access info:"
    write (unit, "(3x,A,L1)")   "is open   = ", object%is_open
    if (object%has_error) then
       write (unit, "(3x,A,A,A)")  "error     = '", char (object%error), "'"
    else
       write (unit, "(3x,A)")      "error     = [none]"
    end if
  end subroutine dlaccess_write
    
\eatline
\nwendcode{}\nwbegindocs{135}The interface to the library functions:
\nwenddocs{}\nwbegincode{136}\moddef{OS interface: interfaces}\endmoddef
  interface 
     function dlopen (filename, flag) result (handle) bind(C)
       import
       character(c_char), dimension(*) :: filename
       integer(c_int), value :: flag
       type(c_ptr) :: handle
     end function dlopen
  end interface

  interface
     function dlclose (handle) result (status) bind(C)
       import
       type(c_ptr), value :: handle
       integer(c_int) :: status
     end function dlclose
  end interface

  interface 
     function dlerror () result (str) bind(C)
       import
       type(c_ptr) :: str
     end function dlerror
  end interface

  interface 
     function dlsym (handle, symbol) result (fptr) bind(C)
       import
       type(c_ptr), value :: handle
       character(c_char), dimension(*) :: symbol
       type(c_funptr) :: fptr
     end function dlsym
  end interface

\eatline
\nwendcode{}\nwbegindocs{137}This reads an error string and transforms it into a {\tt{}string{\char95}t}
object, if an error has occured.  If not, set the error flag to false
and return an empty string.
\nwenddocs{}\nwbegincode{138}\moddef{System defs: public parameters}\plusendmoddef
  integer, parameter, public :: DLERROR_LEN = 160
\nwendcode{}\nwbegincode{139}\moddef{OS interface: procedures}\plusendmoddef
  subroutine read_dlerror (has_error, error)
    logical, intent(out) :: has_error
    type(string_t), intent(out) :: error
    type(c_ptr) :: err_cptr
    character(len=DLERROR_LEN, kind=c_char), pointer :: err_fptr
    integer :: str_end
    err_cptr = dlerror ()
    if (c_associated (err_cptr)) then
       call c_f_pointer (err_cptr, err_fptr)
       has_error = .true.
       str_end = scan (err_fptr, c_null_char)
       if (str_end > 0) then
          error = err_fptr(1:str_end-1)
       else
          error = err_fptr
       end if
    else
       has_error = .false.
       error = ""
    end if
  end subroutine read_dlerror

\eatline
\nwendcode{}\nwbegindocs{140}This is the Fortran API.  Init/final open and close the file,
i.e., load and unload the library.

Note that a library can be opened more than once, and that for an
ultimate close as many {\tt{}dlclose} calls as {\tt{}dlopen} calls are
necessary.  However, we assume that it is opened and closed only once.
\nwenddocs{}\nwbegincode{141}\moddef{OS interface: public}\plusendmoddef
  public :: dlaccess_init
  public :: dlaccess_final
\nwendcode{}\nwbegincode{142}\moddef{OS interface: dlaccess: TBP}\plusendmoddef
  procedure :: init => dlaccess_init
  procedure :: final => dlaccess_final
\nwendcode{}\nwbegincode{143}\moddef{OS interface: procedures}\plusendmoddef
  subroutine dlaccess_init (dlaccess, prefix, libname, os_data)
    class(dlaccess_t), intent(out) :: dlaccess
    type(string_t), intent(in) :: prefix, libname
    type(os_data_t), intent(in), optional :: os_data
    type(string_t) :: filename
    logical :: exist
    dlaccess%filename = libname
    filename = prefix // "/" // libname
    inquire (file=char(filename), exist=exist)
    if (.not. exist) then
       filename = prefix // "/.libs/" // libname
       inquire (file=char(filename), exist=exist)
       if (.not. exist) then
          dlaccess%has_error = .true.
          dlaccess%error = "Library '" // filename // "' not found"
          return
       end if
    end if
    dlaccess%handle = dlopen (char (filename) // c_null_char, ior ( &
       RTLD_LAZY, RTLD_LOCAL))
    dlaccess%is_open = c_associated (dlaccess%handle)
    call read_dlerror (dlaccess%has_error, dlaccess%error)
  end subroutine dlaccess_init

  subroutine dlaccess_final (dlaccess)
    class(dlaccess_t), intent(inout) :: dlaccess
    integer(c_int) :: status
    if (dlaccess%is_open) then
       status = dlclose (dlaccess%handle)
       dlaccess%is_open = .false.
       call read_dlerror (dlaccess%has_error, dlaccess%error)
    end if
  end subroutine dlaccess_final

\eatline
\nwendcode{}\nwbegindocs{144}Return true if an error has occured.
\nwenddocs{}\nwbegincode{145}\moddef{OS interface: public}\plusendmoddef
  public :: dlaccess_has_error
\nwendcode{}\nwbegincode{146}\moddef{OS interface: procedures}\plusendmoddef
  function dlaccess_has_error (dlaccess) result (flag)
    logical :: flag
    type(dlaccess_t), intent(in) :: dlaccess
    flag = dlaccess%has_error
  end function dlaccess_has_error

\eatline
\nwendcode{}\nwbegindocs{147}Return the error string currently stored in the {\tt{}dlaccess} object.
\nwenddocs{}\nwbegincode{148}\moddef{OS interface: public}\plusendmoddef
  public :: dlaccess_get_error
\nwendcode{}\nwbegincode{149}\moddef{OS interface: procedures}\plusendmoddef
  function dlaccess_get_error (dlaccess) result (error)
    type(string_t) :: error
    type(dlaccess_t), intent(in) :: dlaccess
    error = dlaccess%error
  end function dlaccess_get_error

\eatline
\nwendcode{}\nwbegindocs{150}The symbol handler returns the C address of the function with the
given string name.  (It is a good idea to use {\tt{}bind(C)} for all
functions accessed by this, such that the name string is
well-defined.)  Call {\tt{}c{\char95}f{\char95}procpointer} to cast this into a Fortran
procedure pointer with an appropriate interface.
\nwenddocs{}\nwbegincode{151}\moddef{OS interface: public}\plusendmoddef
  public :: dlaccess_get_c_funptr
\nwendcode{}\nwbegincode{152}\moddef{OS interface: procedures}\plusendmoddef
  function dlaccess_get_c_funptr (dlaccess, fname) result (fptr)
    type(c_funptr) :: fptr
    type(dlaccess_t), intent(inout) :: dlaccess
    type(string_t), intent(in) :: fname
    fptr = dlsym (dlaccess%handle, char (fname) // c_null_char)
    call read_dlerror (dlaccess%has_error, dlaccess%error)
  end function dlaccess_get_c_funptr

\eatline
\nwendcode{}\nwbegindocs{153}\nwdocspar
\subsection{Predicates}
Return true if the library is loaded.  In particular, this is false if
loading was unsuccessful.
\nwenddocs{}\nwbegincode{154}\moddef{OS interface: public}\plusendmoddef
  public :: dlaccess_is_open
\nwendcode{}\nwbegincode{155}\moddef{OS interface: procedures}\plusendmoddef
  function dlaccess_is_open (dlaccess) result (flag)
    logical :: flag
    type(dlaccess_t), intent(in) :: dlaccess
    flag = dlaccess%is_open
  end function dlaccess_is_open

\eatline
\nwendcode{}\nwbegindocs{156}\nwdocspar
\subsection{Shell access}
This is the standard system call for executing a shell command, such
as invoking a compiler.

In F2008 there will be the equivalent built-in command
{\tt{}execute{\char95}command{\char95}line}.
\nwenddocs{}\nwbegincode{157}\moddef{OS interface: public}\plusendmoddef
  public :: os_system_call
\nwendcode{}\nwbegincode{158}\moddef{OS interface: procedures}\plusendmoddef
  subroutine os_system_call (command_string, status, verbose)
    type(string_t), intent(in) :: command_string
    integer, intent(out), optional :: status
    logical, intent(in), optional :: verbose
    logical :: verb
    integer :: stat
    verb = .false.;  if (present (verbose))  verb = verbose
    if (verb) &
         call msg_message ("command: " // char (command_string))
    stat = system (char (command_string) // c_null_char)
    if (present (status)) then
       status = stat
    else if (stat /= 0) then
       if (.not. verb) &
            call msg_message ("command: " // char (command_string))
       write (msg_buffer, "(A,I0)")  "Return code = ", stat
       call msg_message ()
       call msg_fatal ("System command returned with nonzero status code")
    end if
  end subroutine os_system_call

\eatline
\nwendcode{}\nwbegincode{159}\moddef{OS interface: interfaces}\plusendmoddef
  interface
     function system (command) result (status) bind(C)
       import
       integer(c_int) :: status
       character(c_char), dimension(*) :: command
     end function system
  end interface

\eatline
\nwendcode{}\nwbegindocs{160}\nwdocspar
\subsection{Querying for a directory}
This queries for the existence of a directory. There is no standard way to
achieve this in FORTRAN, and if we were to call into {\tt{}libc}, we would need access
to C macros for evaluating the result, so we resort to calling {\tt{}test} as a
system call.
\nwenddocs{}\nwbegincode{161}\moddef{OS interface: public}\plusendmoddef
  public :: os_dir_exist
\nwendcode{}\nwbegincode{162}\moddef{OS interface: procedures}\plusendmoddef
  function os_dir_exist (name) result (res)
    type(string_t), intent(in) :: name
    logical :: res
    integer :: status
    call os_system_call ('test -d "' // name // '"', status=status)
    res = status == 0
  end function os_dir_exist
\eatline
\nwendcode{}\nwbegindocs{163}\nwdocspar
\nwenddocs{}\nwbegincode{164}\moddef{OS interface: public}\plusendmoddef
  public :: os_file_exist
\nwendcode{}\nwbegincode{165}\moddef{OS interface: procedures}\plusendmoddef
  function os_file_exist (name) result (exist)
    type(string_t), intent(in) :: name
!    logical, intent(in), optional :: verb
    logical :: exist
!    integer :: status
!    call  os_system_call ('test -f "' // name // '"', status=status, verbose=verb)
!    res = (status == 0)
    inquire (file = char (name), exist=exist)
  end function os_file_exist

\eatline
\nwendcode{}\nwbegindocs{166}\nwdocspar
\subsection{Fortran compiler and linker}
Compile a single module for use in a shared library, but without
linking.
\nwenddocs{}\nwbegincode{167}\moddef{OS interface: public}\plusendmoddef
  public :: os_compile_shared
\nwendcode{}\nwbegincode{168}\moddef{OS interface: procedures}\plusendmoddef
  subroutine os_compile_shared (src, os_data, status)
    type(string_t), intent(in) :: src
    type(os_data_t), intent(in) :: os_data
    integer, intent(out), optional :: status
    type(string_t) :: command_string
    if (os_data%use_libtool) then
       command_string = &
            os_data%whizard_libtool // " --mode=compile " // &
            os_data%fc // " " // &
            "-c " // &
            os_data%whizard_includes // " " // &
            os_data%fcflags // " " // &
            "'" // src // os_data%fc_src_ext // "'"
    else
       command_string = &
            os_data%fc // " " // &
            "-c  " // &
            os_data%fcflags_pic // " " // &
            os_data%whizard_includes // " " // &
            os_data%fcflags // " " // &
            "'" // src // os_data%fc_src_ext // "'"
    end if
    call os_system_call (command_string, status)
  end subroutine os_compile_shared
   
\eatline
\nwendcode{}\nwbegindocs{169}Link an array of object files to build a shared object library.  In
the libtool case, we have to specify a {\tt{}-rpath}, otherwise only a
static library can be built.  However, since the library is never
installed, this rpath is irrelevant.
\nwenddocs{}\nwbegincode{170}\moddef{OS interface: public}\plusendmoddef
  public :: os_link_shared
\nwendcode{}\nwbegincode{171}\moddef{OS interface: procedures}\plusendmoddef
  subroutine os_link_shared (objlist, lib, os_data, status)
    type(string_t), intent(in) :: objlist, lib
    type(os_data_t), intent(in) :: os_data
    integer, intent(out), optional :: status
    type(string_t) :: command_string
    if (os_data%use_libtool) then
       command_string = &
            os_data%whizard_libtool // " --mode=link " // &
            os_data%fc // " " // &
            "-module " // &
            "-rpath /usr/local/lib" // " " // &
            os_data%fcflags // " " // &
            os_data%whizard_ldflags // " " // &
            os_data%ldflags // " " // &
            "-o '" // lib // ".la' " // &
            objlist
    else
       command_string = &
            os_data%ld // " " // &
            os_data%ldflags_so // " " // &
            os_data%fcflags // " " // &
            os_data%whizard_ldflags // " " // &
            os_data%ldflags // " " // &
            "-o '" // lib // os_data%shlib_ext // "' " // &
            objlist
    end if
    call os_system_call (command_string, status)
  end subroutine os_link_shared

\eatline
\nwendcode{}\nwbegindocs{172}Link an array of object files / libraries to build a static executable.
\nwenddocs{}\nwbegincode{173}\moddef{OS interface: public}\plusendmoddef
  public :: os_link_static
\nwendcode{}\nwbegincode{174}\moddef{OS interface: procedures}\plusendmoddef
  subroutine os_link_static (objlist, exec_name, os_data, status)
    type(string_t), intent(in) :: objlist, exec_name
    type(os_data_t), intent(in) :: os_data
    integer, intent(out), optional :: status
    type(string_t) :: command_string
    if (os_data%use_libtool) then
       command_string = &
            os_data%whizard_libtool // " --mode=link " // &
            os_data%fc // " " // &
            "-static-libtool-libs " // &
            os_data%fcflags // " " // &
            os_data%whizard_ldflags // " " // &
            os_data%ldflags // " " // &
            os_data%ldflags_static // " " // &
            "-o '" // exec_name // "' " // &
            objlist // " " // &
            os_data%ldflags_hepmc // " " // &
            os_data%ldflags_hoppet
    else
       command_string = &
            os_data%ld // " " // &
            os_data%ldflags_so // " " // &
            os_data%fcflags // " " // &
            os_data%whizard_ldflags // " " // &
            os_data%ldflags // " " // &
            os_data%ldflags_static // " " // &
            "-o '" // exec_name // "' " // &
            objlist // " " // &
            os_data%ldflags_hepmc // " " // &
            os_data%ldflags_hoppet
    end if
    call os_system_call (command_string, status)
  end subroutine os_link_static

\eatline
\nwendcode{}\nwbegindocs{175}Determine the name of the shared library to link.  If libtool is
used, this is encoded in the {\tt{}.la} file which resides in place of
the library itself.
\nwenddocs{}\nwbegincode{176}\moddef{OS interface: public}\plusendmoddef
  public :: os_get_dlname
\nwendcode{}\nwbegincode{177}\moddef{OS interface: procedures}\plusendmoddef
  function os_get_dlname (lib, os_data, ignore, silent) result (dlname)
    type(string_t) :: dlname
    type(string_t), intent(in) :: lib
    type(os_data_t), intent(in) :: os_data
    logical, intent(in), optional :: ignore, silent
    type(string_t) :: filename
    type(string_t) :: buffer
    logical :: exist, required, quiet
    integer :: u
    u = free_unit ()
    if (present (ignore)) then
       required = .not. ignore
    else
       required = .true.
    end if
         if (present (silent)) then
       quiet = silent
    else
       quiet = .false.
    end if
    if (os_data%use_libtool) then
       filename = lib // ".la"
       inquire (file=char(filename), exist=exist)
       if (exist) then
          open (unit=u, file=char(filename), action="read", status="old")
          SCAN_LTFILE: do
             call get (u, buffer)
             if (extract (buffer, 1, 7) == "dlname=") then
                dlname = extract (buffer, 9)
                dlname = remove (dlname, len (dlname))
                exit SCAN_LTFILE
             end if
          end do SCAN_LTFILE
          close (u)
       else if (required) then
          if (.not. quiet) call msg_fatal (" Library '" // char (lib) &
               // "': libtool archive not found")
          dlname = ""
       else
          if (.not. quiet) call msg_message ("[No compiled library '" &
               // char (lib) // "']")
          dlname = ""
       end if
    else
       dlname = lib // os_data%shlib_ext
       inquire (file=char(dlname), exist=exist)
       if (.not. exist) then
          if (required) then
             if (.not. quiet) call msg_fatal (" Library '" // char (lib) &
                  // "' not found")
          else
             if (.not. quiet) call msg_message &
                ("[No compiled process library '" // char (lib) // "']")
             dlname = ""
          end if
       end if
    end if
  end function os_get_dlname

\eatline
\nwendcode{}\nwbegindocs{178}\nwdocspar
\subsection{Controlling OpenMP}
OpenMP is handled automatically by the library for the most part.  Here
is a convenience routine for setting the number of threads, with some
diagnostics.
\nwenddocs{}\nwbegincode{179}\moddef{OS interface: public}\plusendmoddef
  public :: openmp_set_num_threads_verbose
\nwendcode{}\nwbegincode{180}\moddef{OS interface: procedures}\plusendmoddef
  subroutine openmp_set_num_threads_verbose (num_threads, openmp_logging)
    integer, intent(in) :: num_threads
    integer :: n_threads
    logical, intent(in), optional :: openmp_logging
    logical :: logging
    if (present (openmp_logging)) then
       logging = openmp_logging
    else
       logging = .true.
    end if
    n_threads = num_threads
    if (openmp_is_active ()) then
       if (num_threads == 1) then
          if (logging) then
             write (msg_buffer, "(A,I0,A)")  "OpenMP: Using ", num_threads, &
                  " thread"
             call msg_message
          end if
          n_threads = num_threads
       else if (num_threads > 1) then
          if (logging) then
             write (msg_buffer, "(A,I0,A)")  "OpenMP: Using ", num_threads, &
                  " threads"
             call msg_message
          end if
          n_threads = num_threads
       else
          if (logging) then
             write (msg_buffer, "(A,I0,A)")  "OpenMP: " &
                  // "Illegal value of openmp_num_threads (", num_threads, &
               ") ignored"
             call msg_error
          end if
          n_threads = openmp_get_default_max_threads ()
          if (logging) then
             write (msg_buffer, "(A,I0,A)")  "OpenMP: Using ", &
                  n_threads, " threads"
             call msg_message          
          end if
       end if
       if (n_threads > openmp_get_default_max_threads ()) then
          if (logging) then
             write (msg_buffer, "(A,I0)")  "OpenMP: " &
                  // "Number of threads is greater than library default of ", &
                  openmp_get_default_max_threads ()
             call msg_warning
          end if
       end if
       call openmp_set_num_threads (n_threads)
    else if (num_threads /= 1) then
       if (logging) then
          write (msg_buffer, "(A,I0,A)")  "openmp_num_threads set to ", &
               num_threads, ", but OpenMP is not active: ignored"
          call msg_warning
       end if
    end if
  end subroutine openmp_set_num_threads_verbose

\eatline
\nwendcode{}\nwbegindocs{181}\nwdocspar
\subsection{Unit tests}
\nwenddocs{}\nwbegincode{182}\moddef{OS interface: public}\plusendmoddef
  public :: os_interface_test
\nwendcode{}\nwbegincode{183}\moddef{OS interface: procedures}\plusendmoddef
  subroutine os_interface_test (u, results)
    integer, intent(in) :: u
    type(test_results_t), intent(inout) :: results
  \LA{}OS interface: execute tests\RA{}
  end subroutine os_interface_test

\eatline
\nwendcode{}\nwbegindocs{184}Write a Fortran source file, compile it to a shared library, load
it, and execute the contained function.
\nwenddocs{}\nwbegincode{185}\moddef{OS interface: execute tests}\endmoddef
  call test (os_interface_1, "os_interface_1", &
       "check OS interface routines", &
       u, results)
\nwendcode{}\nwbegindocs{186}\nwdocspar
\nwenddocs{}\nwbegincode{187}\moddef{OS interface: tests}\endmoddef
  subroutine os_interface_1 (u)
    integer, intent(in) :: u
    type(dlaccess_t) :: dlaccess
    type(string_t) :: fname, libname, ext
    type(os_data_t) :: os_data
    type(string_t) :: filename_src, filename_obj
    abstract interface
       function so_test_proc (i) result (j) bind(C)
         import c_int
         integer(c_int), intent(in) :: i
         integer(c_int) :: j
       end function so_test_proc
    end interface
    procedure(so_test_proc), pointer :: so_test => null ()
    type(c_funptr) :: c_fptr
    integer :: unit
    integer(c_int) :: i
    call os_data_init (os_data)
    fname = "so_test"
    filename_src = fname // os_data%fc_src_ext
    if (os_data%use_libtool) then
       ext = ".lo"
    else
       ext = os_data%obj_ext
    end if
    filename_obj = fname // ext
    libname = fname // os_data%shlib_ext
    
    write (u, "(A)")  "* Test output: OS interface"
    write (u, "(A)")  "*   Purpose: check os_interface routines"
    write (u, "(A)")      
            
    write (u, "(A)")  "* write source file 'so_test.f90'"
    write (u, "(A)")
    unit = free_unit ()
    open (unit=unit, file=char(filename_src), action="write")
    write (unit, "(A)")  "function so_test (i) result (j) bind(C)"
    write (unit, "(A)")  "  use iso_c_binding"
    write (unit, "(A)")  "  integer(c_int), intent(in) :: i"
    write (unit, "(A)")  "  integer(c_int) :: j"
    write (unit, "(A)")  "  j = 2 * i"
    write (unit, "(A)")  "end function so_test"
    close (unit)
    write (u, "(A)")  "* compile and link as 'so_test.so/dylib'"
    write (u, "(A)")
    call os_compile_shared (fname, os_data)
    call os_link_shared (filename_obj, fname, os_data)
    write (u, "(A)")  "* load library 'so_test.so/dylib'"
    write (u, "(A)")
    call dlaccess_init (dlaccess, var_str ("."), libname, os_data)
    if (dlaccess_is_open (dlaccess)) then
       write (u, "(A)") "     success"
    else
       write (u, "(A)") "     failure"
    end if
    write (u, "(A)")  "* load symbol 'so_test'"
    write (u, "(A)")
    c_fptr = dlaccess_get_c_funptr (dlaccess, fname)
    if (c_associated (c_fptr)) then
       write (u, "(A)") "     success"
    else
       write (u, "(A)") "     failure"
    end if
    call c_f_procpointer (c_fptr, so_test)
    write (u, "(A)") "* Execute function from 'so_test.so/dylib'"
    i = 7
    write (u, "(A,1x,I1)")  "     input  = ", i
    write (u, "(A,1x,I1)")  "     result = ", so_test(i)
    if (so_test(i) / i .ne. 2) then
       write (u, "(A)")  "* Compiling and linking ISO C functions failed."
    else
       write (u, "(A)")  "* Successful."
    end if
    write (u, "(A)")
    write (u, "(A)")  "* Cleanup"
    call dlaccess_final (dlaccess)
  end subroutine os_interface_1

\eatline
\nwendcode{}\nwbegindocs{188}\nwdocspar
%%%%%%%%%%%%%%%%%%%%%%%%%%%%%%%%%%%%%%%%%%%%%%%%%%%%%%%%%%%%%%%%%%%%%%%%
\section{Interface for formatted I/O}
For access to formatted printing (possibly input), we interface the C
{\tt{}printf} family of functions.  There are two important issues here:
\begin{enumerate}
\item
  {\tt{}printf} takes an arbitrary number of arguments, relying on the C stack.
  This is not interoperable.  We interface it with C wrappers that output a
  single integer, real or string and restrict the allowed formats accordingly.
\item
  Restricting format strings is essential also for preventing format string
  attacks.  Allowing arbitrary format string would create a real security hole
  in a Fortran program.
\item
  The string returned by {\tt{}sprintf} must be allocated to the right size.
\end{enumerate}
\nwenddocs{}\nwbegincode{189}\moddef{\code{}formats.f90\edoc{}}\endmoddef
\LA{}File header\RA{}

module formats

  use, intrinsic :: iso_c_binding
  
\LA{}Use kinds\RA{}
\LA{}Use strings\RA{}
  use io_units
  use unit_tests
  use diagnostics

\LA{}Standard module head\RA{}

\LA{}Formats: public\RA{}

\LA{}Formats: parameters\RA{}

\LA{}Formats: types\RA{}

\LA{}Formats: interfaces\RA{}

contains

\LA{}Formats: procedures\RA{}

\LA{}Formats: tests\RA{}

end module formats
\eatline
\nwendcode{}\nwbegindocs{190}\nwdocspar
\subsection{Parsing a C format string}
The C format string contains characters and format conversion specifications.
The latter are initiated by a {\tt{}{\char37}} sign.  If the next letter is also a {\tt{}{\char37}},
a percent sign is printed and no conversion is done.  Otherwise, a conversion
is done and applied to the next argument in the argument list.  First comes an
optional flag ({\tt{}{\char35}}, {\tt{}0}, {\tt{}-}, {\tt{}+}, or space), an optional field width
(decimal digits starting not with zero), an optional precision (period, then
another decimal digit string), a length modifier (irrelevant for us, therefore
not supported), and a conversion specifier: {\tt{}d} or {\tt{}i} for integer; {\tt{}e},
{\tt{}f}, {\tt{}g} (also upper case) for double-precision real, {\tt{}s} for a string.

We explicitly exclude all other conversion specifiers, and we check the
specifiers against the actual arguments.

\subsubsection{A type for passing arguments}
This is a polymorphic type that can hold integer, real (double), and string
arguments.
\nwenddocs{}\nwbegincode{191}\moddef{Formats: parameters}\endmoddef
  integer, parameter :: ARGTYPE_NONE = 0
  integer, parameter :: ARGTYPE_LOG = 1
  integer, parameter :: ARGTYPE_INT = 2
  integer, parameter :: ARGTYPE_REAL = 3
  integer, parameter :: ARGTYPE_STR = 4

\eatline
\nwendcode{}\nwbegindocs{192}The integer and real entries are actually scalars, but we avoid relying on
the allocatable-scalar feature and make them one-entry arrays.  The character
entry is a real array which is a copy of the string.

Logical values are mapped to strings (true or false), so this type parameter
value is mostly unused.
\nwenddocs{}\nwbegincode{193}\moddef{Formats: public}\endmoddef
  public :: sprintf_arg_t
\nwendcode{}\nwbegincode{194}\moddef{Formats: types}\endmoddef
  type :: sprintf_arg_t
    private
    integer :: type = ARGTYPE_NONE
    integer(c_int), dimension(:), allocatable :: ival
    real(c_double), dimension(:), allocatable :: rval
    character(c_char), dimension(:), allocatable :: sval
  end type sprintf_arg_t

\eatline
\nwendcode{}\nwbegincode{195}\moddef{Formats: public}\plusendmoddef
  public :: sprintf_arg_init
\nwendcode{}\nwbegincode{196}\moddef{Formats: interfaces}\endmoddef
  interface sprintf_arg_init
     module procedure sprintf_arg_init_log
     module procedure sprintf_arg_init_int
     module procedure sprintf_arg_init_real
     module procedure sprintf_arg_init_str
  end interface

\nwendcode{}\nwbegincode{197}\moddef{Formats: procedures}\endmoddef
  subroutine sprintf_arg_init_log (arg, lval)
    type(sprintf_arg_t), intent(out) :: arg
    logical, intent(in) :: lval
    arg%type = ARGTYPE_STR
    if (lval) then
       allocate (arg%sval (5))
       arg%sval = ['t', 'r', 'u', 'e', c_null_char]
    else
       allocate (arg%sval (6))
       arg%sval = ['f', 'a', 'l', 's', 'e', c_null_char]
    end if
  end subroutine sprintf_arg_init_log

  subroutine sprintf_arg_init_int (arg, ival)
    type(sprintf_arg_t), intent(out) :: arg
    integer, intent(in) :: ival
    arg%type = ARGTYPE_INT
    allocate (arg%ival (1))
    arg%ival = ival
  end subroutine sprintf_arg_init_int

  subroutine sprintf_arg_init_real (arg, rval)
    type(sprintf_arg_t), intent(out) :: arg
    real(default), intent(in) :: rval
    arg%type = ARGTYPE_REAL
    allocate (arg%rval (1))
    arg%rval = rval
  end subroutine sprintf_arg_init_real

  subroutine sprintf_arg_init_str (arg, sval)
    type(sprintf_arg_t), intent(out) :: arg
    type(string_t), intent(in) :: sval
    integer :: i
    arg%type = ARGTYPE_STR
    allocate (arg%sval (len (sval) + 1))
    do i = 1, len (sval)
       arg%sval(i) = extract (sval, i, i)
    end do
    arg%sval(len (sval) + 1) = c_null_char
  end subroutine sprintf_arg_init_str

\eatline
\nwendcode{}\nwbegincode{198}\moddef{Formats: procedures}\plusendmoddef
  subroutine sprintf_arg_write (arg, unit)
    type(sprintf_arg_t), intent(in) :: arg
    integer, intent(in), optional :: unit
    integer :: u
    u = given_output_unit (unit)
    select case (arg%type)
    case (ARGTYPE_NONE)
      write (u, *) "[none]"
    case (ARGTYPE_INT)
      write (u, "(1x,A,1x)", advance = "no")  "[int]"
      write (u, *)  arg%ival
    case (ARGTYPE_REAL)
      write (u, "(1x,A,1x)", advance = "no")  "[real]"
      write (u, *)  arg%rval
    case (ARGTYPE_STR)
      write (u, "(1x,A,1x,A)", advance = "no")  "[string]", '"'
      write (u, *)  arg%rval, '"'
    end select
  end subroutine sprintf_arg_write

\eatline
\nwendcode{}\nwbegindocs{199}Return an upper bound for the length of the printed version; in case of
strings the result is exact.
\nwenddocs{}\nwbegincode{200}\moddef{Formats: procedures}\plusendmoddef
  elemental function sprintf_arg_get_length (arg) result (length)
    integer :: length
    type(sprintf_arg_t), intent(in) :: arg
    select case (arg%type)
    case (ARGTYPE_INT)
       length = log10 (real (huge (arg%ival(1)))) + 2
    case (ARGTYPE_REAL)
       length = log10 (real (radix (arg%rval(1))) ** digits (arg%rval(1))) + 8
    case (ARGTYPE_STR)
       length = size (arg%sval)
    case default
       length = 0
    end select
  end function sprintf_arg_get_length

\eatline
\nwendcode{}\nwbegincode{201}\moddef{Formats: procedures}\plusendmoddef
  subroutine sprintf_arg_apply_sprintf (arg, fmt, result, actual_length)
    type(sprintf_arg_t), intent(in) :: arg
    character(c_char), dimension(:), intent(in) :: fmt
    character(c_char), dimension(:), intent(inout) :: result
    integer, intent(out) :: actual_length
    integer(c_int) :: ival
    real(c_double) :: rval
    select case (arg%type)
    case (ARGTYPE_NONE)
      actual_length = sprintf_none (result, fmt)
    case (ARGTYPE_INT)
      ival = arg%ival(1)
      actual_length = sprintf_int (result, fmt, ival)
    case (ARGTYPE_REAL)
      rval = arg%rval(1)
      actual_length = sprintf_double (result, fmt, rval)
    case (ARGTYPE_STR)
      actual_length = sprintf_str (result, fmt, arg%sval)
    case default
      call msg_bug ("sprintf_arg_apply_sprintf called with illegal type")
    end select
    if (actual_length < 0) then
       write (msg_buffer, *) "Format: '", fmt, "'"
       call msg_message ()
       write (msg_buffer, *) "Output: '", result, "'"
       call msg_message ()
       call msg_error ("I/O error in sprintf call")
       actual_length = 0
    end if
  end subroutine sprintf_arg_apply_sprintf

\eatline
\nwendcode{}\nwbegindocs{202}\nwdocspar
\subsubsection{Container type for the output}
There is a procedure which chops the format string into pieces that contain at
most one conversion specifier.  Pairing this with a {\tt{}sprintf{\char95}arg} object, we
get the actual input to the {\tt{}sprintf} interface.  The type below holds this
input and can allocate the output string.
\nwenddocs{}\nwbegincode{203}\moddef{Formats: types}\plusendmoddef
  type :: sprintf_interface_t
    private
    character(c_char), dimension(:), allocatable :: input_fmt
    type(sprintf_arg_t) :: arg
    character(c_char), dimension(:), allocatable :: output_str
    integer :: output_str_len = 0
  end type sprintf_interface_t

\eatline
\nwendcode{}\nwbegincode{204}\moddef{Formats: procedures}\plusendmoddef
  subroutine sprintf_interface_init (intf, fmt, arg)
    type(sprintf_interface_t), intent(out) :: intf
    type(string_t), intent(in) :: fmt
    type(sprintf_arg_t), intent(in) :: arg
    integer :: fmt_len, i
    fmt_len = len (fmt)
    allocate (intf%input_fmt (fmt_len + 1))
    do i = 1, fmt_len
       intf%input_fmt(i) = extract (fmt, i, i)
    end do
    intf%input_fmt(fmt_len+1) = c_null_char
    intf%arg = arg
    allocate (intf%output_str (len (fmt) + sprintf_arg_get_length (arg) + 1))
  end subroutine sprintf_interface_init 

\eatline
\nwendcode{}\nwbegincode{205}\moddef{Formats: procedures}\plusendmoddef
  subroutine sprintf_interface_write (intf, unit)
    type(sprintf_interface_t), intent(in) :: intf
    integer, intent(in), optional :: unit
    integer :: u
    u = given_output_unit (unit)
    write (u, *) "Format string = ", '"', intf%input_fmt, '"'
    write (u, "(1x,A,1x)", advance = "no") "Argument = "
    call sprintf_arg_write (intf%arg, unit)
    if (intf%output_str_len > 0) then
       write (u, *) "Result string = ", &
            '"', intf%output_str (1:intf%output_str_len), '"'
    end if
  end subroutine sprintf_interface_write

\eatline
\nwendcode{}\nwbegindocs{206}Return the output string:
\nwenddocs{}\nwbegincode{207}\moddef{Formats: procedures}\plusendmoddef
  function sprintf_interface_get_result (intf) result (string)
    type(string_t) :: string
    type(sprintf_interface_t), intent(in) :: intf
    character(kind = c_char, len = max (intf%output_str_len, 0)) :: buffer
    integer :: i
    if (intf%output_str_len > 0) then
       do i = 1, intf%output_str_len
          buffer(i:i) = intf%output_str(i)
       end do
       string = buffer(1:intf%output_str_len)
    else
       string = ""
    end if
  end function sprintf_interface_get_result

\eatline
\nwendcode{}\nwbegincode{208}\moddef{Formats: procedures}\plusendmoddef
  subroutine sprintf_interface_apply_sprintf (intf)
    type(sprintf_interface_t), intent(inout) :: intf
    call sprintf_arg_apply_sprintf &
         (intf%arg, intf%input_fmt, intf%output_str, intf%output_str_len)
  end subroutine sprintf_interface_apply_sprintf

\eatline
\nwendcode{}\nwbegindocs{209}Import the interfaces defined in the previous section:
\nwenddocs{}\nwbegincode{210}\moddef{Formats: interfaces}\plusendmoddef
\LA{}sprintf interfaces\RA{}

\nwendcode{}\nwbegindocs{211}\nwdocspar
\subsubsection{Scan the format string}
Chop it into pieces that contain one conversion
specifier each.  The zero-th piece contains the part before the first
specifier.  Check the specifiers and allow only the subset that we support.
Also check for an exact match between conversion specifiers and input
arguments.  The result is an allocated array of {\tt{}sprintf{\char95}interface} object;
each one contains a piece of the format string and the corresponding
argument.
\nwenddocs{}\nwbegincode{212}\moddef{Formats: procedures}\plusendmoddef
  subroutine chop_and_check_format_string (fmt, arg, intf)
    type(string_t), intent(in) :: fmt
    type(sprintf_arg_t), dimension(:), intent(in) :: arg
    type(sprintf_interface_t), dimension(:), intent(out), allocatable :: intf
    integer :: n_args, i
    type(string_t), dimension(:), allocatable :: split_fmt
    type(string_t) :: word, buffer, separator
    integer :: pos, length, l
    logical :: ok
    type(sprintf_arg_t) :: arg_null
    ok = .true.
    length = 0
    n_args = size (arg)
    allocate (split_fmt (0:n_args))
    split_fmt = ""
    buffer = fmt
    SCAN_ARGS: do i = 1, n_args
       FIND_CONVERSION: do
          call split (buffer, word, "%", separator=separator)
          if (separator == "") then
             call msg_message ('"' // char (fmt) // '"')
             call msg_error ("C-formatting string: " &
                  // "too few conversion specifiers in format string")
             ok = .false.;  exit SCAN_ARGS
          end if
          split_fmt(i-1) = split_fmt(i-1) // word
          if (extract (buffer, 1, 1) /= "%") then
             split_fmt(i) = "%"
             exit FIND_CONVERSION
          else
             split_fmt(i-1) = split_fmt(i-1) // "%"
          end if
       end do FIND_CONVERSION
       pos = verify (buffer, "#0-+ ")   ! Flag characters (zero or more)
       split_fmt(i) = split_fmt(i) // extract (buffer, 1, pos-1)
       buffer = remove (buffer, 1, pos-1)
       pos = verify (buffer, "123456890")  ! Field width
       word = extract (buffer, 1, pos-1)
       if (len (word) /= 0) then
         call read_int_from_string (word, len (word), l)
         length = length + l
       end if
       split_fmt(i) = split_fmt(i) // word
       buffer = remove (buffer, 1, pos-1)
       if (extract (buffer, 1, 1) == ".") then
          buffer = remove (buffer, 1, 1)
          pos = verify (buffer, "1234567890")   ! Precision
          split_fmt(i) = split_fmt(i) // "." // extract (buffer, 1, pos-1)
          buffer = remove (buffer, 1, pos-1)
       end if
       ! Length modifier would come here, but is not allowed
       select case (char (extract (buffer, 1, 1)))  ! conversion specifier
       case ("d", "i")
          if (arg(i)%type /= ARGTYPE_INT) then
             call msg_message ('"' // char (fmt) // '"')
             call msg_error ("C-formatting string: " &
                  // "argument type mismatch: integer value expected")
             ok = .false.;  exit SCAN_ARGS
          end if
       case ("e", "E", "f", "F", "g", "G")
          if (arg(i)%type /= ARGTYPE_REAL) then
             call msg_message ('"' // char (fmt) // '"')
             call msg_error ("C-formatting string: " &
                  // "argument type mismatch: real value expected")
             ok = .false.;  exit SCAN_ARGS
          end if
       case ("s")
          if (arg(i)%type /= ARGTYPE_STR) then
             call msg_message ('"' // char (fmt) // '"')
             call msg_error ("C-formatting string: " &
                  // "argument type mismatch: logical or string value expected")
             ok = .false.;  exit SCAN_ARGS
          end if
       case default
          call msg_message ('"' // char (fmt) // '"')
          call msg_error ("C-formatting string: " &
               // "illegal or incomprehensible conversion specifier")
          ok = .false.;  exit SCAN_ARGS
       end select
       split_fmt(i) = split_fmt(i) // extract (buffer, 1, 1)
       buffer = remove (buffer, 1, 1)
    end do SCAN_ARGS
    if (ok) then
       FIND_EXTRA_CONVERSION: do
          call split (buffer, word, "%", separator=separator)
          split_fmt(n_args) = split_fmt(n_args) // word // separator
          if (separator == "")  exit FIND_EXTRA_CONVERSION
          if (extract (buffer, 1, 1) == "%") then
             split_fmt(n_args) = split_fmt(n_args) // "%"
             buffer = remove (buffer, 1, 1)
          else
             call msg_message ('"' // char (fmt) // '"')
             call msg_error ("C-formatting string: " &
                  // "too many conversion specifiers in format string")
             ok = .false.;  exit FIND_EXTRA_CONVERSION
          end if
       end do FIND_EXTRA_CONVERSION
       split_fmt(n_args) = split_fmt(n_args) // buffer
       allocate (intf (0:n_args))
       call sprintf_interface_init (intf(0), split_fmt(0), arg_null)
       do i = 1, n_args
          call sprintf_interface_init (intf(i), split_fmt(i), arg(i))
       end do
    else
       allocate (intf (0))
    end if
  contains
    subroutine read_int_from_string (word, length, l)
      type(string_t), intent(in) :: word
      integer, intent(in) :: length
      integer, intent(out) :: l
      character(len=length) :: buffer
      buffer = word
      read (buffer, *) l
    end subroutine read_int_from_string
  end subroutine chop_and_check_format_string

\eatline
\nwendcode{}\nwbegindocs{213}\nwdocspar
\subsection{API}
\nwenddocs{}\nwbegincode{214}\moddef{Formats: public}\plusendmoddef
  public :: sprintf
\nwendcode{}\nwbegincode{215}\moddef{Formats: procedures}\plusendmoddef
  function sprintf (fmt, arg) result (string)
    type(string_t) :: string
    type(string_t), intent(in) :: fmt
    type(sprintf_arg_t), dimension(:), intent(in) :: arg
    type(sprintf_interface_t), dimension(:), allocatable :: intf
    integer :: i
    string = ""
    call chop_and_check_format_string (fmt, arg, intf)
    if (size (intf) > 0) then
       do i = 0, ubound (intf, 1)
          call sprintf_interface_apply_sprintf (intf(i))
          string = string // sprintf_interface_get_result (intf(i))
       end do
    end if
  end function sprintf

\eatline
\nwendcode{}\nwbegindocs{216}\nwdocspar
\subsection{Test}
\nwenddocs{}\nwbegincode{217}\moddef{Formats: public}\plusendmoddef
  public :: format_test
\nwendcode{}\nwbegincode{218}\moddef{Formats: procedures}\plusendmoddef
  subroutine format_test (u, results)
    integer, intent(in) :: u
    type(test_results_t), intent(inout) :: results
  \LA{}Formats: execute tests\RA{}
  end subroutine format_test

\eatline
\nwendcode{}\nwbegincode{219}\moddef{Formats: execute tests}\endmoddef
  call test (format_1, "format_1", &
       "check formatting routines", &
       u, results)
\nwendcode{}\nwbegincode{220}\moddef{Formats: tests}\endmoddef
  subroutine format_1 (u)
    integer, intent(in) :: u
    write (u, "(A)")  "*** Test 1: a string ***"
    write (u, "(A)")
    call test_run (var_str("%s"), 1, [4], ['abcdefghij'], u)
    write (u, "(A)")  "*** Test 2: two integers ***"
    write (u, "(A)")
    call test_run (var_str("%d,%d"), 2, [2, 2], ['42', '13'], u)
    write (u, "(A)")  "*** Test 3: floating point number ***"
    write (u, "(A)")
    call test_run (var_str("%8.4f"), 1, [3], ['42567.12345'], u)
    write (u, "(A)")  "*** Test 4: general expression ***"
    call test_run (var_str("%g"), 1, [3], ['3.1415'], u)
    contains
      subroutine test_run (fmt, n_args, type, buffer, unit)
        type(string_t), intent(in) :: fmt        
        integer, intent(in) :: n_args, unit
        logical :: lval
        integer :: ival
        real(default) :: rval
        integer :: i
        type(string_t) :: string
        type(sprintf_arg_t), dimension(:), allocatable :: arg
        integer, dimension(n_args), intent(in) :: type
        character(*), dimension(n_args), intent(in) :: buffer
        write (unit, "(A,A)")   "Format string :", char(fmt)
        write (unit, "(A,I1)")  "Number of args:", n_args
        allocate (arg (n_args))
        do i = 1, n_args
           write (unit, "(A,I1)")  "Argument (type ) = ", type(i)
           select case (type(i))
           case (ARGTYPE_LOG)
              read (buffer(i), *)  lval
              call sprintf_arg_init (arg(i), lval)
           case (ARGTYPE_INT)
              read (buffer(i), *)  ival
              call sprintf_arg_init (arg(i), ival)
           case (ARGTYPE_REAL)
              read (buffer(i), *)  rval
              call sprintf_arg_init (arg(i), rval)
           case (ARGTYPE_STR)
              call sprintf_arg_init (arg(i), var_str (trim (buffer(i))))
           end select
         end do
         string = sprintf (fmt, arg)
         write (unit, "(A,A,A)")  "Result: '", char (string), "'"
         deallocate (arg)
       end subroutine test_run
  end subroutine format_1

\eatline
\nwendcode{}\nwbegindocs{221}\nwdocspar
%%%%%%%%%%%%%%%%%%%%%%%%%%%%%%%%%%%%%%%%%%%%%%%%%%%%%%%%%%%%%%%%%%%%%%%%
\section{CPU timing}

The time is stored in a simple derived type which just holds a
floating-point number.
\nwenddocs{}\nwbegincode{222}\moddef{\code{}cputime.f90\edoc{}}\endmoddef
\LA{}File header\RA{}

module cputime

\LA{}Use kinds\RA{}
  use io_units
\LA{}Use strings\RA{}
  use unit_tests
  use diagnostics

\LA{}Standard module head\RA{}

\LA{}CPU time: public\RA{}

\LA{}CPU time: types\RA{}

\LA{}CPU time: interfaces\RA{}

contains

\LA{}CPU time: procedures\RA{}

\LA{}CPU time: tests\RA{}

end module cputime
\eatline
\nwendcode{}\nwbegindocs{223}The CPU time is a floating-point number with an arbitrary reference time.
It is single precision (default real, not {\tt{}real(default)}).
It is measured in seconds.
\nwenddocs{}\nwbegincode{224}\moddef{CPU time: public}\endmoddef
  public :: time_t
\nwendcode{}\nwbegincode{225}\moddef{CPU time: types}\endmoddef
  type :: time_t
     private
     logical :: known = .false.
     real :: value = 0
   contains
   \LA{}CPU time: time: TBP\RA{}
  end type time_t

\eatline
\nwendcode{}\nwbegincode{226}\moddef{CPU time: time: TBP}\endmoddef
  procedure :: write => time_write
\nwendcode{}\nwbegincode{227}\moddef{CPU time: procedures}\endmoddef
  subroutine time_write (object, unit)
    class(time_t), intent(in) :: object
    integer, intent(in), optional :: unit
    integer :: u
    u = given_output_unit (unit)
    write (u, "(1x,A)", advance="no")  "Time in seconds ="
    if (object%known) then
       write (u, "(1x,ES10.3)")  object%value
    else
       write (u, "(1x,A)")  "[unknown]"
    end if
  end subroutine time_write

\eatline
\nwendcode{}\nwbegindocs{228}Set the current time
\nwenddocs{}\nwbegincode{229}\moddef{CPU time: time: TBP}\plusendmoddef
  procedure :: set_current => time_set_current
\nwendcode{}\nwbegincode{230}\moddef{CPU time: procedures}\plusendmoddef
  subroutine time_set_current (time)
    class(time_t), intent(out) :: time
    integer :: msecs
    call system_clock (msecs)
    time%value = real (msecs) / 1000.
    time%known = time%value > 0
  end subroutine time_set_current

\eatline
\nwendcode{}\nwbegindocs{231}Assign to a {\tt{}real(default} value.  If the time is undefined, return zero.
\nwenddocs{}\nwbegincode{232}\moddef{CPU time: public}\plusendmoddef
  public :: assignment(=)
\nwendcode{}\nwbegincode{233}\moddef{CPU time: interfaces}\endmoddef
  interface assignment(=)
    module procedure real_assign_time
    module procedure real_default_assign_time
  end interface

\nwendcode{}\nwbegincode{234}\moddef{CPU time: procedures}\plusendmoddef
  pure subroutine real_assign_time (r, time)
    real, intent(out) :: r
    class(time_t), intent(in) :: time
    if (time%known) then
       r = time%value
    else
       r = 0
    end if
  end subroutine real_assign_time
    
  pure subroutine real_default_assign_time (r, time)
    real(default), intent(out) :: r
    class(time_t), intent(in) :: time
    if (time%known) then
       r = time%value
    else
       r = 0
    end if
  end subroutine real_default_assign_time
    
\eatline
\nwendcode{}\nwbegindocs{235}Assign an integer or (single precision) real value to the time object.
\nwenddocs{}\nwbegincode{236}\moddef{CPU time: time: TBP}\plusendmoddef
  generic :: assignment(=) => time_assign_from_integer, time_assign_from_real
  procedure, private :: time_assign_from_integer
  procedure, private :: time_assign_from_real
\nwendcode{}\nwbegincode{237}\moddef{CPU time: procedures}\plusendmoddef
  subroutine time_assign_from_integer (time, ival)
    class(time_t), intent(out) :: time
    integer, intent(in) :: ival
    time%value = ival
    time%known = .true.
  end subroutine time_assign_from_integer
  
  subroutine time_assign_from_real (time, rval)
    class(time_t), intent(out) :: time
    real, intent(in) :: rval
    time%value = rval
    time%known = .true.
  end subroutine time_assign_from_real
  
\eatline
\nwendcode{}\nwbegindocs{238}Add times and compute time differences.  If any input value is undefined,
the result is undefined.
\nwenddocs{}\nwbegincode{239}\moddef{CPU time: time: TBP}\plusendmoddef
  generic :: operator(-) => subtract_times
  generic :: operator(+) => add_times
  procedure, private :: subtract_times
  procedure, private :: add_times
\nwendcode{}\nwbegincode{240}\moddef{CPU time: procedures}\plusendmoddef
  pure function subtract_times (t_end, t_begin) result (time)
    type(time_t) :: time
    class(time_t), intent(in) :: t_end, t_begin
    if (t_end%known .and. t_begin%known) then
       time%known = .true.
       time%value = t_end%value - t_begin%value
    end if
  end function subtract_times
    
  pure function add_times (t1, t2) result (time)
    type(time_t) :: time
    class(time_t), intent(in) :: t1, t2
    if (t1%known .and. t2%known) then
       time%known = .true.
       time%value = t1%value + t2%value
    end if
  end function add_times
    
\eatline
\eatline
\nwendcode{}\nwbegindocs{241}Check if a time is known, so we can use it:
\nwenddocs{}\nwbegincode{242}\moddef{CPU time: time: TBP}\plusendmoddef
  procedure :: is_known => time_is_known
\nwendcode{}\nwbegincode{243}\moddef{CPU time: procedures}\plusendmoddef
  function time_is_known (time) result (flag)
    class(time_t), intent(in) :: time
    logical :: flag
    flag = time%known
  end function time_is_known
  
\eatline
\nwendcode{}\nwbegindocs{244}We define functions for converting the time into ss / mm:ss / hh:mm:ss
/ dd:mm:hh:ss.
\nwenddocs{}\nwbegincode{245}\moddef{CPU time: time: TBP}\plusendmoddef
  generic :: expand => time_expand_s, time_expand_ms, &
       time_expand_hms, time_expand_dhms
  procedure, private :: time_expand_s
  procedure, private :: time_expand_ms
  procedure, private :: time_expand_hms
  procedure, private :: time_expand_dhms
\nwendcode{}\nwbegincode{246}\moddef{CPU time: procedures}\plusendmoddef
  subroutine time_expand_s (time, sec)
    class(time_t), intent(in) :: time
    integer, intent(out) :: sec
    if (time%known) then
       sec = time%value
    else
       call msg_bug ("Time: attempt to expand undefined value")
    end if
  end subroutine time_expand_s

  subroutine time_expand_ms (time, min, sec)
    class(time_t), intent(in) :: time
    integer, intent(out) :: min, sec
    if (time%known) then
       if (time%value >= 0) then
          sec = mod (int (time%value), 60)
       else
          sec = - mod (int (- time%value), 60)
       end if
       min = time%value / 60
    else
       call msg_bug ("Time: attempt to expand undefined value")
    end if
  end subroutine time_expand_ms

  subroutine time_expand_hms (time, hour, min, sec)
    class(time_t), intent(in) :: time
    integer, intent(out) :: hour, min, sec
    call time%expand (min, sec)
    hour = min / 60
    if (min >= 0) then
       min = mod (min, 60)
    else
       min = - mod (-min, 60)
    end if
  end subroutine time_expand_hms

  subroutine time_expand_dhms (time, day, hour, min, sec)
    class(time_t), intent(in) :: time
    integer, intent(out) :: day, hour, min, sec
    call time%expand (hour, min, sec)
    day = hour / 24
    if (hour >= 0) then
       hour = mod (hour, 24)
    else
       hour = - mod (- hour, 24)
    end if
  end subroutine time_expand_dhms

\eatline
\nwendcode{}\nwbegindocs{247}Use the above expansions to generate a time string.
\nwenddocs{}\nwbegincode{248}\moddef{CPU time: time: TBP}\plusendmoddef
  procedure :: to_string_s => time_to_string_s
  procedure :: to_string_ms => time_to_string_ms
  procedure :: to_string_hms => time_to_string_hms
  procedure :: to_string_dhms => time_to_string_dhms
\nwendcode{}\nwbegincode{249}\moddef{CPU time: procedures}\plusendmoddef
  function time_to_string_s (time) result (str)
    class(time_t), intent(in) :: time
    type(string_t) :: str
    character(256) :: buffer
    integer :: s
    call time%expand (s)
    write (buffer, "(I0,'s')")  s
    str = trim (buffer)
  end function time_to_string_s

  function time_to_string_ms (time, blank) result (str)
    class(time_t), intent(in) :: time
    logical, intent(in), optional :: blank
    type(string_t) :: str
    character(256) :: buffer    
    integer :: s, m
    logical :: x_out
    x_out = .false.
    if (present (blank))  x_out = blank
    call time%expand (m, s)
    write (buffer, "(I0,'m:',I2.2,'s')")  m, abs (s)
    str = trim (buffer)
    if (x_out) then
       str = replace (str, len(str)-1, "X")
    end if
  end function time_to_string_ms

  function time_to_string_hms (time) result (str)
    class(time_t), intent(in) :: time
    type(string_t) :: str
    character(256) :: buffer
    integer :: s, m, h
    call time%expand (h, m, s)
    write (buffer, "(I0,'h:',I2.2,'m:',I2.2,'s')")  h, abs (m), abs (s)
    str = trim (buffer)
  end function time_to_string_hms

  function time_to_string_dhms (time) result (str)
    class(time_t), intent(in) :: time
    type(string_t) :: str
    character(256) :: buffer
    integer :: s, m, h, d
    call time%expand (d, h, m, s)
    write (buffer, "(I0,'d:',I2.2,'h:',I2.2,'m:',I2.2,'s')")  &
         d, abs (h), abs (m), abs (s)
    str = trim (buffer)
  end function time_to_string_dhms

\eatline
\nwendcode{}\nwbegindocs{250}\nwdocspar
\subsection{Timer}
A timer can measure real (wallclock) time differences.  The base type
corresponds to the result, i.e., time difference.  The object contains
two further times for start and stop time.
\nwenddocs{}\nwbegincode{251}\moddef{CPU time: public}\plusendmoddef
  public :: timer_t
\nwendcode{}\nwbegincode{252}\moddef{CPU time: types}\plusendmoddef
  type, extends (time_t) :: timer_t
     private
     logical :: running = .false.
     type(time_t) :: t1, t2
   contains
   \LA{}CPU time: timer: TBP\RA{}
  end type timer_t
  
\eatline
\nwendcode{}\nwbegindocs{253}Output.  If the timer is running, we indicate this, otherwise write
just the result.
\nwenddocs{}\nwbegincode{254}\moddef{CPU time: timer: TBP}\endmoddef
  procedure :: write => timer_write
\nwendcode{}\nwbegincode{255}\moddef{CPU time: procedures}\plusendmoddef
  subroutine timer_write (object, unit)
    class(timer_t), intent(in) :: object
    integer, intent(in), optional :: unit
    integer :: u
    u = given_output_unit (unit)
    if (object%running) then
       write (u, "(1x,A)")  "Time in seconds = [running]"
    else
       call object%time_t%write (u)
    end if
  end subroutine timer_write

\eatline
\nwendcode{}\nwbegindocs{256}Start the timer: store the current time in the first entry and adapt
the status.  We forget any previous values.
\nwenddocs{}\nwbegincode{257}\moddef{CPU time: timer: TBP}\plusendmoddef
  procedure :: start => timer_start
\nwendcode{}\nwbegincode{258}\moddef{CPU time: procedures}\plusendmoddef
  subroutine timer_start (timer)
    class(timer_t), intent(out) :: timer
    call timer%t1%set_current ()
    timer%running = .true.
  end subroutine timer_start
  
\eatline
\nwendcode{}\nwbegindocs{259}Restart the timer: simply adapt the status, keeping the start time.
\nwenddocs{}\nwbegincode{260}\moddef{CPU time: timer: TBP}\plusendmoddef
  procedure :: restart => timer_restart
\nwendcode{}\nwbegincode{261}\moddef{CPU time: procedures}\plusendmoddef
  subroutine timer_restart (timer)
    class(timer_t), intent(inout) :: timer
    if (timer%t1%known .and. .not. timer%running) then
       timer%running = .true.
    else
       call msg_bug ("Timer: restart attempt from wrong status")
    end if       
  end subroutine timer_restart
  
\eatline
\nwendcode{}\nwbegindocs{262}Stop the timer: store the current time in the second entry, adapt
the status, and compute the elapsed time.
\nwenddocs{}\nwbegincode{263}\moddef{CPU time: timer: TBP}\plusendmoddef
  procedure :: stop => timer_stop
\nwendcode{}\nwbegincode{264}\moddef{CPU time: procedures}\plusendmoddef
  subroutine timer_stop (timer)
    class(timer_t), intent(inout) :: timer
    call timer%t2%set_current ()
    timer%running = .false.
    call timer%evaluate ()
  end subroutine timer_stop
  
\eatline
\nwendcode{}\nwbegindocs{265}This is separate, available for the unit test.
\nwenddocs{}\nwbegincode{266}\moddef{CPU time: timer: TBP}\plusendmoddef
  procedure, private :: evaluate => timer_evaluate
\nwendcode{}\nwbegincode{267}\moddef{CPU time: procedures}\plusendmoddef
  subroutine timer_evaluate (timer)
    class(timer_t), intent(inout) :: timer
    timer%time_t = timer%t2 - timer%t1
  end subroutine timer_evaluate
  
\eatline
\nwendcode{}\nwbegindocs{268}\nwdocspar
\subsection{Unit tests}
\nwenddocs{}\nwbegincode{269}\moddef{CPU time: public}\plusendmoddef
  public :: cputime_test
\nwendcode{}\nwbegincode{270}\moddef{CPU time: tests}\endmoddef
  subroutine cputime_test (u, results)
    integer, intent(in) :: u
    type(test_results_t), intent(inout) :: results
  \LA{}CPU time: execute tests\RA{}
  end subroutine cputime_test
  
\eatline
\nwendcode{}\nwbegindocs{271}\nwdocspar
\subsubsection{Basic tests}
Check basic functions of the time object.  The part which we can't
check is getting the actual time from the system clock, since the
output will not be reproducible.  However, we can check time formats
and operations.
\nwenddocs{}\nwbegincode{272}\moddef{CPU time: execute tests}\endmoddef
  call test (cputime_1, "cputime_1", &
       "time operations", &
       u, results)
\nwendcode{}\nwbegincode{273}\moddef{CPU time: tests}\plusendmoddef
  subroutine cputime_1 (u)
    integer, intent(in) :: u
    type(time_t) :: time, time1, time2
    real :: t
    integer :: d, h, m, s

    write (u, "(A)")  "* Test output: cputime_1"
    write (u, "(A)")  "*   Purpose: check time operations"
    write (u, "(A)")      
    
    write (u, "(A)") "* Undefined time"
    write (u, *)

    call time%write (u)

    write (u, *)
    write (u, "(A)") "* Set time to zero"
    write (u, *)

    time = 0
    call time%write (u)

    write (u, *)
    write (u, "(A)") "* Set time to 1.234 s"
    write (u, *)

    time = 1.234
    call time%write (u)

    t = time
    write (u, "(1x,A,F6.3)")  "Time as real =", t

    write (u, *)
    write (u, "(A)") "* Compute time difference"
    write (u, *)
    
    time1 = 5.33
    time2 = 7.55
    time = time2 - time1
    
    call time1%write (u)
    call time2%write (u)
    call time%write (u)

    write (u, *)
    write (u, "(A)") "* Compute time sum"
    write (u, *)
    
    time = time2 + time1
    
    call time1%write (u)
    call time2%write (u)
    call time%write (u)

    write (u, *)
    write (u, "(A)") "* Expand time"
    write (u, *)
    
    time1 = ((24 + 1) * 60 + 1) * 60 + 1
    time2 = ((3 * 24 + 23) * 60 + 59) * 60 + 59
    
    call time1%expand (s)
    write (u, 1)  "s =", s
    call time1%expand (m,s)
    write (u, 1)  "ms =", m, s
    call time1%expand (h,m,s)
    write (u, 1)  "hms =", h, m, s
    call time1%expand (d,h,m,s)
    write (u, 1)  "dhms =", d, h, m, s
    
    call time2%expand (s)
    write (u, 1)  "s =", s
    call time2%expand (m,s)
    write (u, 1)  "ms =", m, s
    call time2%expand (h,m,s)
    write (u, 1)  "hms =", h, m, s
    call time2%expand (d,h,m,s)
    write (u, 1)  "dhms =", d, h, m, s
    
    write (u, *)
    write (u, "(A)") "* Expand negative time"
    write (u, *)
    
    time1 = - (((24 + 1) * 60 + 1) * 60 + 1)
    time2 = - (((3 * 24 + 23) * 60 + 59) * 60 + 59)
    
    call time1%expand (s)
    write (u, 1)  "s =", s
    call time1%expand (m,s)
    write (u, 1)  "ms =", m, s
    call time1%expand (h,m,s)
    write (u, 1)  "hms =", h, m, s
    call time1%expand (d,h,m,s)
    write (u, 1)  "dhms =", d, h, m, s
    
    call time2%expand (s)
    write (u, 1)  "s =", s
    call time2%expand (m,s)
    write (u, 1)  "ms =", m, s
    call time2%expand (h,m,s)
    write (u, 1)  "hms =", h, m, s
    call time2%expand (d,h,m,s)
    write (u, 1)  "dhms =", d, h, m, s
    
1   format (1x,A,1x,4(I0,:,':'))

    write (u, *)
    write (u, "(A)") "* String from time"
    write (u, *)
    
    time1 = ((24 + 1) * 60 + 1) * 60 + 1
    time2 = ((3 * 24 + 23) * 60 + 59) * 60 + 59
    
    write (u, "(A)")  char (time1%to_string_s ())
    write (u, "(A)")  char (time1%to_string_ms ())
    write (u, "(A)")  char (time1%to_string_hms ())
    write (u, "(A)")  char (time1%to_string_dhms ())

    write (u, "(A)")  char (time2%to_string_s ())
    write (u, "(A)")  char (time2%to_string_ms ())
    write (u, "(A)")  char (time2%to_string_hms ())
    write (u, "(A)")  char (time2%to_string_dhms ())

    write (u, "(A)")
    write (u, "(A)")  "* Blanking out the last second entry"
    write (u, "(A)")
    
    write (u, "(A)")  char (time1%to_string_ms ())
    write (u, "(A)")  char (time1%to_string_ms (.true.))
        
    write (u, *)
    write (u, "(A)") "* String from negative time"
    write (u, *)
    
    time1 = -(((24 + 1) * 60 + 1) * 60 + 1)
    time2 = -(((3 * 24 + 23) * 60 + 59) * 60 + 59)
    
    write (u, "(A)")  char (time1%to_string_s ())
    write (u, "(A)")  char (time1%to_string_ms ())
    write (u, "(A)")  char (time1%to_string_hms ())
    write (u, "(A)")  char (time1%to_string_dhms ())

    write (u, "(A)")  char (time2%to_string_s ())
    write (u, "(A)")  char (time2%to_string_ms ())
    write (u, "(A)")  char (time2%to_string_hms ())
    write (u, "(A)")  char (time2%to_string_dhms ())
    
    write (u, "(A)")
    write (u, "(A)")  "* Test output end: cputime_1"    

  end subroutine cputime_1
  
\eatline
\nwendcode{}\nwbegindocs{274}\nwdocspar
\subsubsection{Timer tests}
Check a timer object.
\nwenddocs{}\nwbegincode{275}\moddef{CPU time: execute tests}\plusendmoddef
  call test (cputime_2, "cputime_2", &
       "timer", &
       u, results)
\nwendcode{}\nwbegincode{276}\moddef{CPU time: tests}\plusendmoddef
  subroutine cputime_2 (u)
    integer, intent(in) :: u
    type(timer_t) :: timer

    write (u, "(A)")  "* Test output: cputime_2"
    write (u, "(A)")  "*   Purpose: check timer"
    write (u, "(A)")      
    
    write (u, "(A)") "* Undefined timer"
    write (u, *)

    call timer%write (u)

    write (u, *)
    write (u, "(A)") "* Start timer"
    write (u, *)

    call timer%start ()
    call timer%write (u)

    write (u, *)
    write (u, "(A)") "* Stop timer (injecting fake timings)"
    write (u, *)

    call timer%stop ()
    timer%t1 = 2
    timer%t2 = 5
    call timer%evaluate ()
    call timer%write (u)

    write (u, *)
    write (u, "(A)") "* Restart timer"
    write (u, *)

    call timer%restart ()
    call timer%write (u)

    write (u, *)
    write (u, "(A)") "* Stop timer again (injecting fake timing)"
    write (u, *)

    call timer%stop ()
    timer%t2 = 10
    call timer%evaluate ()
    call timer%write (u)

    write (u, *)
    write (u, "(A)")  "* Test output end: cputime_2"    

  end subroutine cputime_2
  
\eatline
\nwendcode{}\nwfilename{../../../src/combinatorics/combinatorics.nw}\nwbegindocs{0}% -*- ess-noweb-default-code-mode: f90-mode; noweb-default-code-mode: f90-mode; -*- 
% WHIZARD code as NOWEB source: combinatorics and such


\chapter{Combinatorics}

These modules implement standard algorithms (sorting, hashing, etc.)
that are not available in Fortran.

Fortran doesn't support generic programming, therefore the algorithms
are implemented only for specific data types.
\begin{description}
\item[bytes]
  Derived types for bytes and words.
\item[hashes]
  Types and tools for setting up hashtables.
\item[md5]
  The MD5 algorithm for message digest.
\item[permutations]
  Permuting an array of integers.
\item[sorting]
  Sorting integer and real values.
\end{description}

 
%%%%%%%%%%%%%%%%%%%%%%%%%%%%%%%%%%%%%%%%%%%%%%%%%%%%%%%%%%%%%%%%%%%%%%%%
\section{Bytes and such}
In a few instances we will need the notion of a byte (8-bit) and a word
(32 bit), even a 64-bit word.  A block of 512 bit is also needed (for
MD5).

We rely on integers up to 64 bit being supported by the processor.
The main difference to standard integers is the interpretation as
unsigned integers.
\nwenddocs{}\nwbegincode{1}\moddef{\code{}bytes.f90\edoc{}}\endmoddef
\LA{}File header\RA{}

module bytes

  use kinds, only: i8, i32, i64
  use io_units

\LA{}Standard module head\RA{}

\LA{}Bytes: public\RA{}

\LA{}Bytes: types\RA{}

\LA{}Bytes: parameters\RA{}

\LA{}Bytes: interfaces\RA{}

contains

\LA{}Bytes: procedures\RA{}

end module bytes
\eatline
\nwendcode{}\nwbegindocs{2}\nwdocspar
\subsection{8-bit words: bytes}
This is essentially a wrapper around 8-bit integers.  The wrapper
emphasises their special interpretation as a sequence of bits.
However, we interpret bytes as unsigned integers.
\nwenddocs{}\nwbegincode{3}\moddef{Bytes: public}\endmoddef
  public :: byte_t
\nwendcode{}\nwbegincode{4}\moddef{Bytes: types}\endmoddef
  type :: byte_t
     private
     integer(i8) :: i
  end type byte_t

\eatline
\nwendcode{}\nwbegincode{5}\moddef{Bytes: public}\plusendmoddef
  public :: byte_zero
\nwendcode{}\nwbegincode{6}\moddef{Bytes: parameters}\endmoddef
  type(byte_t), parameter :: byte_zero = byte_t (0_i8)

\eatline
\nwendcode{}\nwbegindocs{7}Set a byte from 8-bit integer:
\nwenddocs{}\nwbegincode{8}\moddef{Bytes: public}\plusendmoddef
  public :: assignment(=)
\nwendcode{}\nwbegincode{9}\moddef{Bytes: interfaces}\endmoddef
  interface assignment(=)
     module procedure set_byte_from_i8
  end interface
\eatline
\nwendcode{}\nwbegincode{10}\moddef{Bytes: procedures}\endmoddef
  subroutine set_byte_from_i8 (b, i)
    type(byte_t), intent(out) :: b
    integer(i8), intent(in) :: i
    b%i = i
  end subroutine set_byte_from_i8

\eatline
\nwendcode{}\nwbegindocs{11}Write a byte in one of two formats: either as a hexadecimal number
(two digits, default) or as a decimal number (one to three digits).
The decimal version is nontrivial because bytes are unsigned integers.
Optionally append a newline.
\nwenddocs{}\nwbegincode{12}\moddef{Bytes: public}\plusendmoddef
  public :: byte_write
\nwendcode{}\nwbegincode{13}\moddef{Bytes: interfaces}\plusendmoddef
  interface byte_write
     module procedure byte_write_unit, byte_write_string
  end interface
\nwendcode{}\nwbegincode{14}\moddef{Bytes: procedures}\plusendmoddef
  subroutine byte_write_unit (b, unit, decimal, newline)
    type(byte_t), intent(in), optional :: b
    integer, intent(in), optional :: unit
    logical, intent(in), optional :: decimal, newline
    logical :: dc, nl
    type(word32_t) :: w
    integer :: u
    u = given_output_unit (unit);  if (u < 0)  return
    dc = .false.;  if (present (decimal))  dc = decimal
    nl = .false.;  if (present (newline))  nl = newline
    if (dc) then
       w = b
       write (u, '(I3)', advance='no')  w%i
    else
       write (u, '(z2.2)', advance='no')  b%i
    end if
    if (nl) write (u, *)
  end subroutine byte_write_unit

\eatline
\nwendcode{}\nwbegindocs{15}The string version is hex-only
\nwenddocs{}\nwbegincode{16}\moddef{Bytes: procedures}\plusendmoddef
  subroutine byte_write_string (b, s)
    type(byte_t), intent(in) :: b
    character(len=2), intent(inout) :: s
    write (s, '(z2.2)')  b%i
  end subroutine byte_write_string

\eatline
\nwendcode{}\nwbegindocs{17}\nwdocspar
\subsection{32-bit words}
This is not exactly a 32-bit integer.  A word is to be filled with
bytes, and it may be partially filled.  The filling is done lowest-byte
first, highest-byte last.  We count the bits, so {\tt{}fill} should be
either 0, 8, 16, 24, or 32.
In printing words, we correspondingly
distinguish between printing zeros and printing blanks.
\nwenddocs{}\nwbegincode{18}\moddef{Bytes: public}\plusendmoddef
  public :: word32_t
\nwendcode{}\nwbegincode{19}\moddef{Bytes: types}\plusendmoddef
  type :: word32_t
     private
     integer(i32) :: i
     integer :: fill = 0
  end type word32_t

\eatline
\nwendcode{}\nwbegindocs{20}Assignment: the word is filled by inserting a 32-bit integer
\nwenddocs{}\nwbegincode{21}\moddef{Bytes: interfaces}\plusendmoddef
  interface assignment(=)
     module procedure word32_set_from_i32
     module procedure word32_set_from_byte
  end interface
\eatline
\nwendcode{}\nwbegincode{22}\moddef{Bytes: procedures}\plusendmoddef
  subroutine word32_set_from_i32 (w, i)
    type(word32_t), intent(out) :: w
    integer(i32), intent(in) :: i
    w%i = i
    w%fill = 32
  end subroutine word32_set_from_i32

\eatline
\nwendcode{}\nwbegindocs{23}Reverse assignment to a 32-bit integer.  We do not check the fill
status.
\nwenddocs{}\nwbegincode{24}\moddef{Bytes: interfaces}\plusendmoddef
  interface assignment(=)
     module procedure i32_from_word32
  end interface
\eatline
\nwendcode{}\nwbegincode{25}\moddef{Bytes: procedures}\plusendmoddef
  subroutine i32_from_word32 (i, w)
    integer(i32), intent(out) :: i
    type(word32_t), intent(in) :: w
    i = w%i
  end subroutine i32_from_word32

\eatline
\nwendcode{}\nwbegindocs{26}Filling with a 8-bit integer is slightly tricky, because in this
interpretation integers are unsigned.
\nwenddocs{}\nwbegincode{27}\moddef{Bytes: procedures}\plusendmoddef
  subroutine word32_set_from_byte (w, b)
    type(word32_t), intent(out) :: w
    type(byte_t), intent(in) :: b
    if (b%i >= 0_i8) then
       w%i = b%i
    else
       w%i = 2_i32*(huge(0_i8)+1_i32) + b%i
    end if
    w%fill = 32
  end subroutine word32_set_from_byte

\eatline
\nwendcode{}\nwbegindocs{28}Check the fill status
\nwenddocs{}\nwbegincode{29}\moddef{Bytes: public}\plusendmoddef
  public :: word32_empty, word32_filled, word32_fill
\nwendcode{}\nwbegincode{30}\moddef{Bytes: procedures}\plusendmoddef
  function word32_empty (w)
    type(word32_t), intent(in) :: w
    logical :: word32_empty
    word32_empty = (w%fill == 0)
  end function word32_empty

  function word32_filled (w)
    type(word32_t), intent(in) :: w
    logical :: word32_filled
    word32_filled = (w%fill == 32)
  end function word32_filled

  function word32_fill (w)
    type(word32_t), intent(in) :: w
    integer :: word32_fill
    word32_fill = w%fill
  end function word32_fill

\eatline
\nwendcode{}\nwbegindocs{31}Partial assignment: append a byte to a partially filled word.
(Note: no assignment if the word is filled, so check this before if
necessary.)  
\nwenddocs{}\nwbegincode{32}\moddef{Bytes: public}\plusendmoddef
  public :: word32_append_byte
\nwendcode{}\nwbegincode{33}\moddef{Bytes: procedures}\plusendmoddef
  subroutine word32_append_byte (w, b)
    type(word32_t), intent(inout) :: w
    type(byte_t), intent(in) :: b
    type(word32_t) :: w1
    if (.not. word32_filled (w)) then
       w1 = b
       call mvbits (w1%i, 0, 8, w%i, w%fill)
       w%fill = w%fill + 8
    end if
  end subroutine word32_append_byte

\eatline
\nwendcode{}\nwbegindocs{34}Extract a byte from a word.  The argument {\tt{}i} is the position,
which may be 0, 1, 2, or 3.

For the final assignment, we set the highest bit separately.
Otherwise, we might trigger an overflow condition for a compiler with
strict checking turned on.
\nwenddocs{}\nwbegincode{35}\moddef{Bytes: public}\plusendmoddef
  public :: byte_from_word32
\nwendcode{}\nwbegincode{36}\moddef{Bytes: procedures}\plusendmoddef
  function byte_from_word32 (w, i) result (b)
    type(word32_t), intent(in) :: w
    integer, intent(in) :: i
    type(byte_t) :: b
    integer(i32) :: j
    j = 0
    if (i >= 0 .and. i*8 < w%fill) then
       call mvbits (w%i, i*8, 8, j, 0)
    end if
    b%i = int (ibclr (j, 7), kind=i8)
    if (btest (j, 7))  b%i = ibset (b%i, 7)
  end function byte_from_word32

\eatline
\nwendcode{}\nwbegindocs{37}Write a word to file or STDOUT.  We understand words as unsigned
integers, therefore we cannot use the built-in routine unchanged.
However, we can make use of the existence of 64-bit integers and their
output routine.

In hexadecimal format, the default version prints eight hex
characters, highest-first.  The {\tt{}bytes} version prints four bytes
(two-hex characters), lowest first, with spaces in-between.  The
decimal bytes version is analogous.  In the {\tt{}bytes} version, missing
bytes are printed as whitespace.
\nwenddocs{}\nwbegincode{38}\moddef{Bytes: public}\plusendmoddef
  public :: word32_write
\nwendcode{}\nwbegincode{39}\moddef{Bytes: interfaces}\plusendmoddef
  interface word32_write
     module procedure word32_write_unit
  end interface
\nwendcode{}\nwbegincode{40}\moddef{Bytes: procedures}\plusendmoddef
  subroutine word32_write_unit (w, unit, bytes, decimal, newline)
    type(word32_t), intent(in) :: w
    integer, intent(in), optional :: unit
    logical, intent(in), optional :: bytes, decimal, newline
    logical :: dc, by, nl
    type(word64_t) :: ww
    integer :: i, u
    u = given_output_unit (unit);  if (u < 0)  return
    by = .false.;  if (present (bytes))   by = bytes
    dc = .false.;  if (present (decimal)) dc = decimal
    nl = .false.;  if (present (newline)) nl = newline
    if (by) then
       do i = 0, 3
          if (i>0)  write (u, '(1x)', advance='no')
          if (8*i < w%fill) then
             call byte_write (byte_from_word32 (w, i), unit, decimal=decimal)
          else if (dc) then
             write (u, '(3x)', advance='no')
          else
             write (u, '(2x)', advance='no')
          end if
       end do
    else if (dc) then
       ww = w
       write (u, '(I10)', advance='no') ww%i
    else
       select case (w%fill)
       case ( 0)
       case ( 8);  write (6, '(1x,z8.2)', advance='no') ibits (w%i, 0, 8)
       case (16);  write (6, '(1x,z8.4)', advance='no') ibits (w%i, 0,16)
       case (24);  write (6, '(1x,z8.6)', advance='no') ibits (w%i, 0,24)
       case (32);  write (6, '(1x,z8.8)', advance='no') ibits (w%i, 0,32)
       end select
    end if
    if (nl) write (u, *)
  end subroutine word32_write_unit

\eatline
\nwendcode{}\nwbegindocs{41}\nwdocspar
\subsection{Operations on 32-bit words}
Define the usual logical operations, as well as addition (mod
$2^{32}$).  We assume that all operands are completely filled.
\nwenddocs{}\nwbegincode{42}\moddef{Bytes: public}\plusendmoddef
  public :: not, ior, ieor, iand, ishft, ishftc
\nwendcode{}\nwbegincode{43}\moddef{Bytes: interfaces}\plusendmoddef
  interface not
     module procedure word_not
  end interface
  interface ior
     module procedure word_or
  end interface
  interface ieor
     module procedure word_eor
  end interface
  interface iand
     module procedure word_and
  end interface
  interface ishft
     module procedure word_shft
  end interface
  interface ishftc
     module procedure word_shftc
  end interface
\eatline
\nwendcode{}\nwbegincode{44}\moddef{Bytes: procedures}\plusendmoddef
  function word_not (w1) result (w2)
    type(word32_t), intent(in) :: w1
    type(word32_t) :: w2
    w2 = not (w1%i)
  end function word_not

  function word_or (w1, w2) result (w3)
    type(word32_t), intent(in) :: w1, w2
    type(word32_t) :: w3
    w3 = ior (w1%i, w2%i)
  end function word_or

  function word_eor (w1, w2) result (w3)
    type(word32_t), intent(in) :: w1, w2
    type(word32_t) :: w3
    w3 = ieor (w1%i, w2%i)
  end function word_eor

  function word_and (w1, w2) result (w3)
    type(word32_t), intent(in) :: w1, w2
    type(word32_t) :: w3
    w3 = iand (w1%i, w2%i)
  end function word_and

  function word_shft (w1, s) result (w2)
    type(word32_t), intent(in) :: w1
    integer, intent(in) :: s
    type(word32_t) :: w2
    w2 = ishft (w1%i, s)
  end function word_shft

  function word_shftc (w1, s) result (w2)
    type(word32_t), intent(in) :: w1
    integer, intent(in) :: s
    type(word32_t) :: w2
    w2 = ishftc (w1%i, s, 32)
  end function word_shftc

\eatline
\nwendcode{}\nwbegindocs{45}Addition is defined mod $2^{32}$, i.e., without overflow checking.
This means that we have to work around a possible overflow check enforced by
the compiler.
\nwenddocs{}\nwbegincode{46}\moddef{Bytes: public}\plusendmoddef
  public :: operator(+)
\nwendcode{}\nwbegincode{47}\moddef{Bytes: interfaces}\plusendmoddef
  interface operator(+)
     module procedure word_add
     module procedure word_add_i8
     module procedure word_add_i32
  end interface
\eatline
\nwendcode{}\nwbegindocs{48}\nwdocspar
\nwenddocs{}\nwbegincode{49}\moddef{Bytes: procedures}\plusendmoddef
  function word_add (w1, w2) result (w3)
    type(word32_t), intent(in) :: w1, w2
    type(word32_t) :: w3
    integer(i64) :: j
    j = int (ibclr (w1%i, 31), i64) + int (ibclr (w2%i, 31), i64)
    w3 = int (ibclr (j, 31), kind=i32)
    if (btest (j, 31)) then
       if (btest (w1%i, 31) .eqv. btest (w2%i, 31))  w3 = ibset (w3%i, 31)
    else
       if (btest (w1%i, 31) .neqv. btest (w2%i, 31))  w3 = ibset (w3%i, 31)
    end if
  end function word_add

  function word_add_i8 (w1, i) result (w3)
    type(word32_t), intent(in) :: w1
    integer(i8), intent(in) :: i
    type(word32_t) :: w3
    integer(i64) :: j
    j = int (ibclr (w1%i, 31), i64) + int (ibclr (i, 7), i64)
    if (btest (i, 7))  j = j + 128
    w3 = int (ibclr (j, 31), kind=i32)
    if (btest (j, 31) .neqv. btest (w1%i, 31))  w3 = ibset (w3%i, 31)
  end function word_add_i8

  function word_add_i32 (w1, i) result (w3)
    type(word32_t), intent(in) :: w1
    integer(i32), intent(in) :: i
    type(word32_t) :: w3
    integer(i64) :: j
    j = int (ibclr (w1%i, 31), i64) + int (ibclr (i, 31), i64)
    w3 = int (ibclr (j, 31), kind=i32)
    if (btest (j, 31)) then
       if (btest (w1%i, 31) .eqv. btest (i, 31))  w3 = ibset (w3%i, 31)
    else
       if (btest (w1%i, 31) .neqv. btest (i, 31))  w3 = ibset (w3%i, 31)
    end if
  end function word_add_i32

\eatline
\nwendcode{}\nwbegindocs{50}\nwdocspar
\subsection{64-bit words}
These objects consist of two 32-bit words.  They thus can hold integer
numbers larger than $2^{32}$ (to be exact, $2^{31}$ since FORTRAN
integers are signed).  The order is low-word, high-word.
\nwenddocs{}\nwbegincode{51}\moddef{Bytes: public}\plusendmoddef
  public :: word64_t
\nwendcode{}\nwbegincode{52}\moddef{Bytes: types}\plusendmoddef
  type :: word64_t
     private
     integer(i64) :: i
  end type word64_t

\eatline
\nwendcode{}\nwbegindocs{53}Set a 64 bit word:
\nwenddocs{}\nwbegincode{54}\moddef{Bytes: interfaces}\plusendmoddef
  interface assignment(=)
     module procedure word64_set_from_i64
     module procedure word64_set_from_word32
  end interface
\eatline
\nwendcode{}\nwbegincode{55}\moddef{Bytes: procedures}\plusendmoddef
  subroutine word64_set_from_i64 (ww, i)
    type(word64_t), intent(out) :: ww
    integer(i64), intent(in) :: i
    ww%i = i
  end subroutine word64_set_from_i64

\eatline
\nwendcode{}\nwbegindocs{56}Filling with a 32-bit word:
\nwenddocs{}\nwbegincode{57}\moddef{Bytes: procedures}\plusendmoddef
  subroutine word64_set_from_word32 (ww, w)
    type(word64_t), intent(out) :: ww
    type(word32_t), intent(in) :: w
    if (w%i >= 0_i32) then
       ww%i = w%i
    else
       ww%i = 2_i64*(huge(0_i32)+1_i64) + w%i
    end if
  end subroutine word64_set_from_word32

\eatline
\nwendcode{}\nwbegindocs{58}Extract a byte from a word.  The argument {\tt{}i} is the position,
which may be between 0 and 7.

For the final assignment, we set the highest bit separately.
Otherwise, we might trigger an overflow condition for a compiler with
strict checking turned on.
\nwenddocs{}\nwbegincode{59}\moddef{Bytes: public}\plusendmoddef
  public :: byte_from_word64, word32_from_word64
\nwendcode{}\nwbegincode{60}\moddef{Bytes: procedures}\plusendmoddef
  function byte_from_word64 (ww, i) result (b)
    type(word64_t), intent(in) :: ww
    integer, intent(in) :: i
    type(byte_t) :: b
    integer(i64) :: j
    j = 0
    if (i >= 0 .and. i*8 < 64) then
       call mvbits (ww%i, i*8, 8, j, 0)
    end if
    b%i = int (ibclr (j, 7), kind=i8)
    if (btest (j, 7))  b%i = ibset (b%i, 7)
  end function byte_from_word64

\eatline
\nwendcode{}\nwbegindocs{61}Extract a 32-bit word from a 64-bit word.  The position is either 0
or 1.
\nwenddocs{}\nwbegincode{62}\moddef{Bytes: procedures}\plusendmoddef
  function word32_from_word64 (ww, i) result (w)
    type(word64_t), intent(in) :: ww
    integer, intent(in) :: i
    type(word32_t) :: w
    integer(i64) :: j
    j = 0
    select case (i)
    case (0);  call mvbits (ww%i,  0, 32, j, 0)
    case (1);  call mvbits (ww%i, 32, 32, j, 0)
    end select
    w = int (ibclr (j, 31), kind=i32)
    if (btest (j, 31))  w = ibset (w%i, 31)
  end function word32_from_word64

\eatline
\nwendcode{}\nwbegindocs{63}Print a 64-bit word.  Decimal version works up to $2^{63}$.
The {\tt{}words} version uses the 'word32' printout, separated by two
spaces.  The low-word is printed first.  The {\tt{}bytes} version also
uses the 'word32' printout.  This implies that the lowest byte is
first.  The default version prints a hexadecimal
number without spaces, highest byte first.
\nwenddocs{}\nwbegincode{64}\moddef{Bytes: public}\plusendmoddef
  public :: word64_write
\nwendcode{}\nwbegincode{65}\moddef{Bytes: interfaces}\plusendmoddef
  interface word64_write
     module procedure word64_write_unit
  end interface
\nwendcode{}\nwbegincode{66}\moddef{Bytes: procedures}\plusendmoddef
  subroutine word64_write_unit (ww, unit, words, bytes, decimal, newline)
    type(word64_t), intent(in) :: ww
    integer, intent(in), optional :: unit
    logical, intent(in), optional :: words, bytes, decimal, newline
    logical :: wo, by, dc, nl
    integer :: u
    u = given_output_unit (unit);  if (u < 0)  return
    wo = .false.;  if (present (words))    wo = words
    by = .false.;  if (present (bytes))    by = bytes
    dc = .false.;  if (present (decimal))  dc = decimal
    nl = .false.;  if (present (newline))  nl = newline
    if (wo .or. by) then
       call word32_write_unit (word32_from_word64 (ww, 0), unit, by, dc)
       write (u, '(2x)', advance='no')
       call word32_write_unit (word32_from_word64 (ww, 1), unit, by, dc)
    else if (dc) then
       write (u, '(I19)', advance='no') ww%i
    else
       write (u, '(Z16)', advance='no') ww%i
    end if
    if (nl) write (u, *)
  end subroutine word64_write_unit

\eatline
\nwendcode{}\nwbegindocs{67}\nwdocspar
%%%%%%%%%%%%%%%%%%%%%%%%%%%%%%%%%%%%%%%%%%%%%%%%%%%%%%%%%%%%%%%%%%%%%%%%
\section{Hashtables}
Hash tables, like lists, are not part of Fortran and must be defined
on a per-case basis.  In this section we define a module that contains
a hash function.  

Furthermore, for reference there is a complete framework of hashtable
type definitions and access functions.  This code is to be replicated
where hash tables are used, mutatis mutandis.
\nwenddocs{}\nwbegincode{68}\moddef{\code{}hashes.f90\edoc{}}\endmoddef
\LA{}File header\RA{}

module hashes

  use kinds, only: i8, i32
  use bytes

\LA{}Standard module head\RA{}

\LA{}Hashes: public\RA{}

contains

\LA{}Hashes: procedures\RA{}

end module hashes
\eatline
\nwendcode{}\nwbegindocs{69}\nwdocspar
\subsection{The hash function} 
This is the one-at-a-time hash function by Bob Jenkins (from
Wikipedia), re-implemented in Fortran.  The function works on an array
of bytes (8-bit integers), as could be produced by, e.g., the
{\tt{}transfer} function, and returns a single 32-bit integer.  For
determining the position in a hashtable, one can pick the lower bits
of the result as appropriate to the hashtable size (which should be a
power of 2).  Note that we are working on signed integers, so the
interpretation of values differs from the C version.  This should not
matter in practice, however.
\nwenddocs{}\nwbegincode{70}\moddef{Hashes: public}\endmoddef
  public :: hash
\nwendcode{}\nwbegincode{71}\moddef{Hashes: procedures}\endmoddef
  function hash (key) result (hashval)
    integer(i32) :: hashval
    integer(i8), dimension(:), intent(in) :: key
    type(word32_t) :: w
    integer :: i
    w = 0_i32
    do i = 1, size (key)
       w = w + key(i)
       w = w + ishft (w, 10)
       w = ieor (w, ishft (w, -6))
    end do
    w = w + ishft (w, 3)
    w = ieor (w, ishft (w, -11))
    w = w + ishft (w, 15)
    hashval = w
  end function hash
    
\eatline
\nwendcode{}\nwbegindocs{72}\nwdocspar
\subsection{The hash table}
We define a generic hashtable type (that depends on the
{\tt{}hash{\char95}data{\char95}t} type) together with associated methods.

This is a template:
\nwenddocs{}\nwbegincode{73}\moddef{Hashtables: types}\endmoddef
  type :: hash_data_t
     integer :: i
  end type hash_data_t

\eatline
\nwendcode{}\nwbegindocs{74}Associated methods:
\nwenddocs{}\nwbegincode{75}\moddef{Hashtables: procedures}\endmoddef
  subroutine hash_data_final (data)
    type(hash_data_t), intent(inout) :: data
  end subroutine hash_data_final

  subroutine hash_data_write (data, unit)
    type(hash_data_t), intent(in) :: data
    integer, intent(in), optional :: unit
    integer :: u
    u = given_output_unit (unit);  if (u < 0)  return
    write (u, *) data%i
  end subroutine hash_data_write

\eatline
\nwendcode{}\nwbegindocs{76}Each hash entry stores the unmasked hash value, the key, and points to
actual data if present.  Note that this could be an allocatable scalar
in principle, but making it a pointer avoids deep copy when expanding
the hashtable.
\nwenddocs{}\nwbegincode{77}\moddef{Hashtables: types}\plusendmoddef
  type :: hash_entry_t
     integer(i32) :: hashval = 0
     integer(i8), dimension(:), allocatable :: key
     type(hash_data_t), pointer :: data => null ()
  end type hash_entry_t

\eatline
\nwendcode{}\nwbegindocs{78}The hashtable object holds the actual table, the number of filled
entries and the number of entries after which the size should be
doubled.  The mask is equal to the table size minus one and thus
coincides with the upper bound of the table index, which starts at zero.
\nwenddocs{}\nwbegincode{79}\moddef{Hashtables: types}\plusendmoddef
  type :: hashtable_t
     integer :: n_entries = 0
     real :: fill_ratio = 0
     integer :: n_entries_max = 0
     integer(i32) :: mask = 0
     type(hash_entry_t), dimension(:), allocatable :: entry
  end type hashtable_t

\eatline
\nwendcode{}\nwbegindocs{80}Initializer: The size has to be a power of two, the fill ratio is a
real (machine-default!) number between 0 and 1.
\nwenddocs{}\nwbegincode{81}\moddef{Hashtables: procedures}\plusendmoddef
  subroutine hashtable_init (hashtable, size, fill_ratio)
    type(hashtable_t), intent(out) :: hashtable
    integer, intent(in) :: size
    real, intent(in) :: fill_ratio
    hashtable%fill_ratio = fill_ratio
    hashtable%n_entries_max = size * fill_ratio
    hashtable%mask = size - 1
    allocate (hashtable%entry (0:hashtable%mask))
  end subroutine hashtable_init

\eatline
\nwendcode{}\nwbegindocs{82}Finalizer: This calls a {\tt{}hash{\char95}data{\char95}final} subroutine which must
exist.
\nwenddocs{}\nwbegincode{83}\moddef{Hashtables: procedures}\plusendmoddef
  subroutine hashtable_final (hashtable)
    type(hashtable_t), intent(inout) :: hashtable
    integer :: i
    do i = 0, hashtable%mask
       if (associated (hashtable%entry(i)%data)) then
          call hash_data_final (hashtable%entry(i)%data)
          deallocate (hashtable%entry(i)%data)
       end if
    end do
    deallocate (hashtable%entry)
  end subroutine hashtable_final

\eatline
\nwendcode{}\nwbegindocs{84}Output.  Here, we refer to a {\tt{}hash{\char95}data{\char95}write} subroutine.
\nwenddocs{}\nwbegincode{85}\moddef{Hashtables: procedures}\plusendmoddef
  subroutine hashtable_write (hashtable, unit)
    type(hashtable_t), intent(in) :: hashtable
    integer, intent(in), optional :: unit
    integer :: u, i
    u = given_output_unit (unit);  if (u < 0)  return
    do i = 0, hashtable%mask
       if (associated (hashtable%entry(i)%data)) then
          write (u, *) i, "(hash =", hashtable%entry(i)%hashval, ")", &
               hashtable%entry(i)%key
          call hash_data_write (hashtable%entry(i)%data, unit)
       end if
    end do
  end subroutine hashtable_write

\eatline
\nwendcode{}\nwbegindocs{86}\nwdocspar
\subsection{Hashtable insertion}
Insert a single entry with the hash value as trial place.  If the
table is filled, first expand it.
\nwenddocs{}\nwbegincode{87}\moddef{Hashtables: procedures}\plusendmoddef
  subroutine hashtable_insert (hashtable, key, data)
    type(hashtable_t), intent(inout) :: hashtable
    integer(i8), dimension(:), intent(in) :: key
    type(hash_data_t), intent(in), target :: data
    integer(i32) :: h
    if (hashtable%n_entries >= hashtable%n_entries_max) &
         call hashtable_expand (hashtable)
    h = hash (key)
    call hashtable_insert_rec (hashtable, h, h, key, data)
  end subroutine hashtable_insert

\eatline
\nwendcode{}\nwbegindocs{88}We need this auxiliary routine for doubling the size of the
hashtable.  We rely on the fact that default assignment copies
the data pointer, not the data themselves.  The temporary array must
not be finalized; it is deallocated automatically together with its
allocatable components.
\nwenddocs{}\nwbegincode{89}\moddef{Hashtables: procedures}\plusendmoddef
  subroutine hashtable_expand (hashtable)
    type(hashtable_t), intent(inout) :: hashtable
    type(hash_entry_t), dimension(:), allocatable :: table_tmp
    integer :: i, s
    allocate (table_tmp (0:hashtable%mask))
    table_tmp = hashtable%entry
    deallocate (hashtable%entry)
    s = 2 * size (table_tmp)
    hashtable%n_entries = 0
    hashtable%n_entries_max = s * hashtable%fill_ratio
    hashtable%mask = s - 1
    allocate (hashtable%entry (0:hashtable%mask))
    do i = 0, ubound (table_tmp, 1)
       if (associated (table_tmp(i)%data)) then
          call hashtable_insert_rec (hashtable, table_tmp(i)%hashval, &
               table_tmp(i)%hashval, table_tmp(i)%key, table_tmp(i)%data)
       end if
    end do
  end subroutine hashtable_expand

\eatline
\nwendcode{}\nwbegindocs{90}Insert a single entry at a trial place {\tt{}h}, reduced to the table
size.  Collision resolution is done simply by choosing the next
element, recursively until the place is empty.  For bookkeeping, we
preserve the original hash value.  For a good hash function, there
should be no clustering.

Note that if the new key exactly matches an existing key, nothing is done.
\nwenddocs{}\nwbegincode{91}\moddef{Hashtables: procedures}\plusendmoddef
  recursive subroutine hashtable_insert_rec (hashtable, h, hashval, key, data)
    type(hashtable_t), intent(inout) :: hashtable
    integer(i32), intent(in) :: h, hashval
    integer(i8), dimension(:), intent(in) :: key
    type(hash_data_t), intent(in), target :: data
    integer(i32) :: i
    i = iand (h, hashtable%mask)
    if (associated (hashtable%entry(i)%data)) then
       if (size (hashtable%entry(i)%key) /= size (key)) then
          call hashtable_insert_rec (hashtable, h + 1, hashval, key, data)
       else if (any (hashtable%entry(i)%key /= key)) then
          call hashtable_insert_rec (hashtable, h + 1, hashval, key, data)
       end if
    else
       hashtable%entry(i)%hashval = hashval
       allocate (hashtable%entry(i)%key (size (key)))
       hashtable%entry(i)%key = key
       hashtable%entry(i)%data => data
       hashtable%n_entries = hashtable%n_entries + 1
    end if
  end subroutine hashtable_insert_rec
    
\eatline
\nwendcode{}\nwbegindocs{92}\nwdocspar
\subsection{Hashtable lookup}
The lookup function has to parallel the insert function.  If the place
is filled, check if the key matches.  Yes: return the pointer; no:
increment the hash value and check again.
\nwenddocs{}\nwbegincode{93}\moddef{Hashtables: procedures}\plusendmoddef
  function hashtable_lookup (hashtable, key) result (ptr)
    type(hash_data_t), pointer :: ptr
    type(hashtable_t), intent(in) :: hashtable
    integer(i8), dimension(:), intent(in) :: key
    ptr => hashtable_lookup_rec (hashtable, hash (key), key)
  end function hashtable_lookup

\eatline
\nwendcode{}\nwbegincode{94}\moddef{Hashtables: procedures}\plusendmoddef
  recursive function hashtable_lookup_rec (hashtable, h, key) result (ptr)
    type(hash_data_t), pointer :: ptr
    type(hashtable_t), intent(in) :: hashtable
    integer(i32), intent(in) :: h
    integer(i8), dimension(:), intent(in) :: key
    integer(i32) :: i
    i = iand (h, hashtable%mask)
    if (associated (hashtable%entry(i)%data)) then
       if (size (hashtable%entry(i)%key) == size (key)) then
          if (all (hashtable%entry(i)%key == key)) then
             ptr => hashtable%entry(i)%data
          else
             ptr => hashtable_lookup_rec (hashtable, h + 1, key)
          end if
       else
          ptr => hashtable_lookup_rec (hashtable, h + 1, key)
       end if
    else
       ptr => null ()
    end if
  end function hashtable_lookup_rec

\eatline
\nwendcode{}\nwbegincode{95}\moddef{Hashtables: public}\endmoddef
  public :: hashtable_test
\nwendcode{}\nwbegincode{96}\moddef{Hashtables: procedures}\plusendmoddef
  subroutine hashtable_test ()
    type(hash_data_t), pointer :: data
    type(hashtable_t) :: hashtable
    integer(i8) :: i
    call hashtable_init (hashtable, 16, 0.25)
    do i = 1, 10
       allocate (data)
       data%i = i*i
       call hashtable_insert (hashtable, [i, i+i], data)
    end do
    call hashtable_insert (hashtable, [2_i8, 4_i8], data)
    call hashtable_write (hashtable)
    data => hashtable_lookup (hashtable, [5_i8, 10_i8])
    if (associated (data)) then
       print *, "lookup:", data%i
    else
       print *, "lookup: --"
    end if
    data => hashtable_lookup (hashtable, [6_i8, 12_i8])
    if (associated (data)) then
       print *, "lookup:", data%i
    else
       print *, "lookup: --"
    end if
    data => hashtable_lookup (hashtable, [4_i8, 9_i8])
    if (associated (data)) then
       print *, "lookup:", data%i
    else
       print *, "lookup: --"
    end if
    call hashtable_final (hashtable)
  end subroutine hashtable_test

\eatline
\nwendcode{}\nwbegindocs{97}\nwdocspar
%%%%%%%%%%%%%%%%%%%%%%%%%%%%%%%%%%%%%%%%%%%%%%%%%%%%%%%%%%%%%%%%%%%%%%%%
\section{MD5 Checksums}
Implementing MD5 checksums allows us
to check input/file integrity on the basis of a well-known standard.
The building blocks have been introduced in the {\tt{}bytes} module.
\nwenddocs{}\nwbegincode{98}\moddef{\code{}md5.f90\edoc{}}\endmoddef
\LA{}File header\RA{}

module md5

  use kinds, only: i8, i32, i64
  use io_units
  use unit_tests
  use system_defs, only: BUFFER_SIZE
  use system_defs, only: LF, EOR, EOF
  use diagnostics
  use bytes
  
\LA{}Standard module head\RA{}

\LA{}MD5: public\RA{}

\LA{}MD5: types\RA{}

\LA{}MD5: variables\RA{}

\LA{}MD5: interfaces\RA{}

contains

\LA{}MD5: procedures\RA{}

\LA{}MD5: tests\RA{}

end module md5
\eatline
\nwendcode{}\nwbegindocs{99}\nwdocspar
\subsection{Blocks}
A block is a sequence of 16 words (64 bytes or 512 bits).  We
anticipate that blocks will be linked, so include a pointer to the
next block.  There is a fill status (word counter), as there is one
for each word.  The fill status is equal to the number of bytes that
are in, so it may be between 0 and 64.
\nwenddocs{}\nwbegincode{100}\moddef{MD5: types}\endmoddef
  type :: block_t
     private
     type(word32_t), dimension(0:15) :: w
     type(block_t), pointer :: next => null ()
     integer :: fill = 0
  end type block_t

\eatline
\nwendcode{}\nwbegindocs{101}Check if a block is completely filled or empty:
\nwenddocs{}\nwbegincode{102}\moddef{MD5: procedures}\endmoddef
  function block_is_empty (b)
    type(block_t), intent(in) :: b
    logical :: block_is_empty
    block_is_empty = (b%fill == 0 .and. word32_empty (b%w(0)))
  end function block_is_empty

  function block_is_filled (b)
    type(block_t), intent(in) :: b
    logical :: block_is_filled
    block_is_filled = (b%fill == 64)
  end function block_is_filled

\eatline
\nwendcode{}\nwbegindocs{103}Append a single byte to a block.  Works only if the block is not yet
filled.
\nwenddocs{}\nwbegincode{104}\moddef{MD5: procedures}\plusendmoddef
  subroutine block_append_byte (bl, by)
    type(block_t), intent(inout) :: bl
    type(byte_t), intent(in) :: by
    if (.not. block_is_filled (bl)) then
       call word32_append_byte (bl%w(bl%fill/4), by)
       bl%fill = bl%fill + 1
    end if
  end subroutine block_append_byte

\eatline
\nwendcode{}\nwbegindocs{105}The printing routine allows for printing as sequences of words or
bytes, decimal or hex.
\nwenddocs{}\nwbegincode{106}\moddef{MD5: interfaces}\endmoddef
  interface block_write
     module procedure block_write_unit
  end interface
\nwendcode{}\nwbegincode{107}\moddef{MD5: procedures}\plusendmoddef
  subroutine block_write_unit (b, unit, bytes, decimal)
    type(block_t), intent(in) :: b
    integer, intent(in), optional :: unit
    logical, intent(in), optional :: bytes, decimal
    logical :: by, dc
    integer :: i, u
    u = given_output_unit (unit);  if (u < 0)  return
    by = .false.;  if (present (bytes))    by = bytes
    dc = .false.;  if (present (decimal))  dc = decimal
    do i = 0, b%fill/4 - 1
       call newline_or_blank (u, i, by, dc)
       call word32_write (b%w(i), unit, bytes, decimal)
    end do
    if (.not. block_is_filled (b)) then
       i = b%fill/4
       if (.not. word32_empty (b%w(i))) then
          call newline_or_blank (u, i, by, dc)
          call word32_write (b%w(i), unit, bytes, decimal)
       end if
    end if
    write (u, *)
  contains
    subroutine newline_or_blank (u, i, bytes, decimal)
      integer, intent(in) :: u, i
      logical, intent(in) :: bytes, decimal
      if (decimal) then
         select case (i)
         case (0)
         case (2,4,6,8,10,12,14);  write (u, *)
         case default
            write (u, '(2x)', advance='no')
         end select
      else if (bytes) then
         select case (i)
         case (0)
         case (4,8,12);  write (u, *)
         case default
            write (u, '(2x)', advance='no')
         end select
      else
         if (i == 8)  write (u, *)
      end if
    end subroutine newline_or_blank
  end subroutine block_write_unit

\eatline
\nwendcode{}\nwbegindocs{108}\nwdocspar
\subsection{Messages}
A message (within this module) is a linked list of blocks.
\nwenddocs{}\nwbegincode{109}\moddef{MD5: types}\plusendmoddef
  type :: message_t
     private
     type(block_t), pointer :: first => null ()
     type(block_t), pointer :: last => null ()
     integer :: n_blocks = 0
  end type message_t

\eatline
\nwendcode{}\nwbegindocs{110}Clear the message list
\nwenddocs{}\nwbegincode{111}\moddef{MD5: procedures}\plusendmoddef
  subroutine message_clear (m)
    type(message_t), intent(inout) :: m
    type(block_t), pointer :: b
    nullify (m%last)
    do
       b => m%first
       if (.not.(associated (b))) exit
       m%first => b%next
       deallocate (b)
    end do
    m%n_blocks = 0
  end subroutine message_clear

\eatline
\nwendcode{}\nwbegindocs{112}Append an empty block to the message list
\nwenddocs{}\nwbegincode{113}\moddef{MD5: procedures}\plusendmoddef
  subroutine message_append_new_block (m)
    type(message_t), intent(inout) :: m
    if (associated (m%last)) then
       allocate (m%last%next)
       m%last => m%last%next
       m%n_blocks = m%n_blocks + 1
    else
       allocate (m%first)
       m%last => m%first
       m%n_blocks = 1
    end if
  end subroutine message_append_new_block

\eatline
\nwendcode{}\nwbegindocs{114}Initialize: clear and allocate the first (empty) block.
\nwenddocs{}\nwbegincode{115}\moddef{MD5: procedures}\plusendmoddef
  subroutine message_init (m)
    type(message_t), intent(inout) :: m
    call message_clear (m)
    call message_append_new_block (m)
  end subroutine message_init

\eatline
\nwendcode{}\nwbegindocs{116}Append a single byte to a message.  If necessary, allocate a new
block.  If the message is empty, initialize it.
\nwenddocs{}\nwbegincode{117}\moddef{MD5: procedures}\plusendmoddef
  subroutine message_append_byte (m, b)
    type(message_t), intent(inout) :: m
    type(byte_t), intent(in) :: b
    if (.not. associated (m%last)) then
       call message_init (m)
    else if (block_is_filled (m%last)) then
       call message_append_new_block (m)
    end if
    call block_append_byte (m%last, b)
  end subroutine message_append_byte

\eatline
\nwendcode{}\nwbegindocs{118}Append zero bytes until the current block is filled up to the required
position.  If we are already beyond that, append a new block and fill
that one.
\nwenddocs{}\nwbegincode{119}\moddef{MD5: procedures}\plusendmoddef
  subroutine message_pad_zero (m, i)
    type(message_t), intent(inout) :: m
    integer, intent(in) :: i
    type(block_t), pointer :: b
    integer :: j
    if (associated (m%last)) then
       b => m%last
       if (b%fill > i) then
          do j = b%fill + 1, 64 + i
             call message_append_byte (m, byte_zero)
          end do
       else
          do j = b%fill + 1, i
             call message_append_byte (m, byte_zero)
          end do
       end if
    end if
  end subroutine message_pad_zero

\eatline
\nwendcode{}\nwbegindocs{120}This returns the number of bits within a message.  We need a 64-bit
word for the result since it may be more than $2^{31}$.  This is also
required by the MD5 standard.
\nwenddocs{}\nwbegincode{121}\moddef{MD5: procedures}\plusendmoddef
  function message_bits (m) result (length)
    type(message_t), intent(in) :: m
    type(word64_t) :: length
    type(block_t), pointer :: b
    integer(i64) :: n_blocks_filled, n_bytes_extra
    if (m%n_blocks > 0) then
       b => m%last
       if (block_is_filled (b)) then
          n_blocks_filled = m%n_blocks
          n_bytes_extra = 0
       else
          n_blocks_filled = m%n_blocks - 1
          n_bytes_extra = b%fill
       end if
       length = n_blocks_filled * 512 + n_bytes_extra * 8
    else
       length = 0_i64
    end if
  end function message_bits

\eatline
\nwendcode{}\nwbegindocs{122}\nwdocspar
\subsection{Message I/O}
Append the contents of a string to a message.  We first cast the
character string into a 8-bit integer array and the append this byte
by byte.
\nwenddocs{}\nwbegincode{123}\moddef{MD5: procedures}\plusendmoddef
  subroutine message_append_string (m, s)
    type(message_t), intent(inout) :: m
    character(len=*), intent(in) :: s
    integer(i64) :: i, n_bytes
    integer(i8), dimension(:), allocatable :: buffer
    integer(i8), dimension(1) :: mold
    type(byte_t) :: b
    n_bytes = size (transfer (s, mold))
    allocate (buffer (n_bytes))
    buffer = transfer (s, mold)
    do i = 1, size (buffer)
       b = buffer(i)
       call message_append_byte (m, b)
    end do
    deallocate (buffer)
  end subroutine message_append_string

\eatline
\nwendcode{}\nwbegindocs{124}Append the contents of a 32-bit integer to a message.  We first cast the
32-bit integer into a 8-bit integer array and the append this byte
by byte.
\nwenddocs{}\nwbegincode{125}\moddef{MD5: procedures}\plusendmoddef
  subroutine message_append_i32 (m, x)
    type(message_t), intent(inout) :: m
    integer(i32), intent(in) :: x
    integer(i8), dimension(4) :: buffer
    type(byte_t) :: b
    integer :: i
    buffer = transfer (x, buffer, size(buffer))
    do i = 1, size (buffer)
       b = buffer(i)
       call message_append_byte (m, b)
    end do
  end subroutine message_append_i32

\eatline
\nwendcode{}\nwbegindocs{126}Append one line from file to a message.  Include the newline character.
\nwenddocs{}\nwbegincode{127}\moddef{MD5: procedures}\plusendmoddef
  subroutine message_append_from_unit (m, u, iostat)
    type(message_t), intent(inout) :: m
    integer, intent(in) :: u
    integer, intent(out) :: iostat
    character(len=BUFFER_SIZE) :: buffer
    read (u, *, iostat=iostat) buffer
    call message_append_string (m, trim (buffer))
    call message_append_string (m, LF)
  end subroutine message_append_from_unit
 
\eatline
\nwendcode{}\nwbegindocs{128}Fill a message from file.  (Each line counts as a string.)
\nwenddocs{}\nwbegincode{129}\moddef{MD5: procedures}\plusendmoddef
  subroutine message_read_from_file (m, f)
    type(message_t), intent(inout) :: m
    character(len=*), intent(in) :: f
    integer :: u, iostat
    u = free_unit ()
    open (file=f, unit=u, action='read')
    do
       call message_append_from_unit (m, u, iostat=iostat)
       if (iostat < 0) exit
    end do
    close (u)
  end subroutine message_read_from_file

\eatline
\nwendcode{}\nwbegindocs{130}Write a message.  After each block, insert an empty line.
\nwenddocs{}\nwbegincode{131}\moddef{MD5: interfaces}\plusendmoddef
  interface message_write
     module procedure message_write_unit
  end interface
\nwendcode{}\nwbegincode{132}\moddef{MD5: procedures}\plusendmoddef
  subroutine message_write_unit (m, unit, bytes, decimal)
    type(message_t), intent(in) :: m
    integer, intent(in), optional :: unit
    logical, intent(in), optional :: bytes, decimal
    type(block_t), pointer :: b
    integer :: u
    u = given_output_unit (unit);  if (u < 0)  return
    b => m%first
    if (associated (b)) then
       do
          call block_write_unit (b, unit, bytes, decimal)
          b => b%next
          if (.not. associated (b))  exit
          write (u, *)
       end do
    end if
  end subroutine message_write_unit

\eatline
\nwendcode{}\nwbegindocs{133}\nwdocspar
\subsection{Auxiliary functions}
These four functions on three words are defined in the MD5 standard:
\nwenddocs{}\nwbegincode{134}\moddef{MD5: procedures}\plusendmoddef
  function ff (x, y, z)
    type(word32_t), intent(in) :: x, y, z
    type(word32_t) :: ff
    ff = ior (iand (x, y), iand (not (x), z))
  end function ff

  function fg (x, y, z)
    type(word32_t), intent(in) :: x, y, z
    type(word32_t) :: fg
    fg = ior (iand (x, z), iand (y, not (z)))
  end function fg

  function fh (x, y, z)
    type(word32_t), intent(in) :: x, y, z
    type(word32_t) :: fh
    fh = ieor (ieor (x, y), z)
  end function fh

  function fi (x, y, z)
    type(word32_t), intent(in) :: x, y, z
    type(word32_t) :: fi
    fi = ieor (y, ior (x, not (z)))
  end function fi

\eatline
\nwendcode{}\nwbegindocs{135}\nwdocspar
\subsection{Auxiliary stuff}
This defines and initializes the table of transformation constants:
\nwenddocs{}\nwbegincode{136}\moddef{MD5: variables}\endmoddef
  type(word32_t), dimension(64), save :: t
  logical, save :: table_initialized = .false.
\eatline
\nwendcode{}\nwbegincode{137}\moddef{MD5: procedures}\plusendmoddef
  subroutine table_init
    type(word64_t) :: ww
    integer :: i
    if (.not.table_initialized) then
       do i = 1, 64
          ww = int (4294967296d0 * abs (sin (i * 1d0)), kind=i64)
          t(i) = word32_from_word64 (ww, 0)
       end do
       table_initialized = .true.
    end if
  end subroutine table_init

\eatline
\nwendcode{}\nwbegindocs{138}This encodes the message digest (4 words) into a 32-character
string.
\nwenddocs{}\nwbegincode{139}\moddef{MD5: procedures}\plusendmoddef
  function digest_string (aa) result (s)
    type(word32_t), dimension (0:3), intent(in) :: aa
    character(len=32) :: s
    integer :: i, j
    do i = 0, 3
       do j = 0, 3
          call byte_write (byte_from_word32 (aa(i), j), s(i*8+j*2+1:i*8+j*2+2))
       end do
    end do
  end function digest_string

\eatline
\nwendcode{}\nwbegindocs{140}\nwdocspar
\subsection{MD5 algorithm}
Pad the message with a byte {\tt{}x80} and then pad zeros up to a full
block minus two words; in these words, insert the message length
(before padding) as a 64-bit word, low-word first.
\nwenddocs{}\nwbegincode{141}\moddef{MD5: procedures}\plusendmoddef
  subroutine message_pad (m)
    type(message_t), intent(inout) :: m
    type(word64_t) :: length
    integer(i8), parameter :: ipad = -128  ! z'80'
    type(byte_t) :: b
    integer :: i
    length = message_bits (m)
    b = ipad
    call message_append_byte (m, b)
    call message_pad_zero (m, 56)
    do i = 0, 7
       call message_append_byte (m, byte_from_word64 (length, i))
    end do
  end subroutine message_pad

\eatline
\nwendcode{}\nwbegindocs{142}Apply a series of transformations onto a state {\tt{}a,b,c,d}, where
the transform function uses each word of the message together with the
predefined words.  Finally, encode the state as a 32-character string.
\nwenddocs{}\nwbegincode{143}\moddef{MD5: procedures}\plusendmoddef
  subroutine message_digest (m, s)
    type(message_t), intent(in) :: m
    character(len=32), intent(out) :: s
    integer(i32), parameter :: ia =  1732584193  ! z'67452301'
    integer(i32), parameter :: ib =  -271733879  ! z'efcdab89'
    integer(i32), parameter :: ic = -1732584194  ! z'98badcfe'
    integer(i32), parameter :: id =   271733878  ! z'10325476'
    type(word32_t) :: a, b, c, d
    type(word32_t) :: aa, bb, cc, dd
    type(word32_t), dimension(0:15) :: x
    type(block_t), pointer :: bl
    call table_init
    a = ia;  b = ib;  c = ic;  d = id
    bl => m%first
    do
       if (.not.associated (bl)) exit
       x = bl%w
       aa = a;  bb = b;  cc = c;  dd = d
       call transform (ff, a, b, c, d,  0,  7,  1)
       call transform (ff, d, a, b, c,  1, 12,  2)
       call transform (ff, c, d, a, b,  2, 17,  3)
       call transform (ff, b, c, d, a,  3, 22,  4)
       call transform (ff, a, b, c, d,  4,  7,  5)
       call transform (ff, d, a, b, c,  5, 12,  6)
       call transform (ff, c, d, a, b,  6, 17,  7)
       call transform (ff, b, c, d, a,  7, 22,  8)
       call transform (ff, a, b, c, d,  8,  7,  9)
       call transform (ff, d, a, b, c,  9, 12, 10)
       call transform (ff, c, d, a, b, 10, 17, 11)
       call transform (ff, b, c, d, a, 11, 22, 12)
       call transform (ff, a, b, c, d, 12,  7, 13)
       call transform (ff, d, a, b, c, 13, 12, 14)
       call transform (ff, c, d, a, b, 14, 17, 15)
       call transform (ff, b, c, d, a, 15, 22, 16)
       call transform (fg, a, b, c, d,  1,  5, 17)
       call transform (fg, d, a, b, c,  6,  9, 18)
       call transform (fg, c, d, a, b, 11, 14, 19)
       call transform (fg, b, c, d, a,  0, 20, 20)
       call transform (fg, a, b, c, d,  5,  5, 21)
       call transform (fg, d, a, b, c, 10,  9, 22)
       call transform (fg, c, d, a, b, 15, 14, 23)
       call transform (fg, b, c, d, a,  4, 20, 24)
       call transform (fg, a, b, c, d,  9,  5, 25)
       call transform (fg, d, a, b, c, 14,  9, 26)
       call transform (fg, c, d, a, b,  3, 14, 27)
       call transform (fg, b, c, d, a,  8, 20, 28)
       call transform (fg, a, b, c, d, 13,  5, 29)
       call transform (fg, d, a, b, c,  2,  9, 30)
       call transform (fg, c, d, a, b,  7, 14, 31)
       call transform (fg, b, c, d, a, 12, 20, 32)
       call transform (fh, a, b, c, d,  5,  4, 33)
       call transform (fh, d, a, b, c,  8, 11, 34)
       call transform (fh, c, d, a, b, 11, 16, 35)
       call transform (fh, b, c, d, a, 14, 23, 36)
       call transform (fh, a, b, c, d,  1,  4, 37)
       call transform (fh, d, a, b, c,  4, 11, 38)
       call transform (fh, c, d, a, b,  7, 16, 39)
       call transform (fh, b, c, d, a, 10, 23, 40)
       call transform (fh, a, b, c, d, 13,  4, 41)
       call transform (fh, d, a, b, c,  0, 11, 42)
       call transform (fh, c, d, a, b,  3, 16, 43)
       call transform (fh, b, c, d, a,  6, 23, 44)
       call transform (fh, a, b, c, d,  9,  4, 45)
       call transform (fh, d, a, b, c, 12, 11, 46)
       call transform (fh, c, d, a, b, 15, 16, 47)
       call transform (fh, b, c, d, a,  2, 23, 48)
       call transform (fi, a, b, c, d,  0,  6, 49)
       call transform (fi, d, a, b, c,  7, 10, 50)
       call transform (fi, c, d, a, b, 14, 15, 51)
       call transform (fi, b, c, d, a,  5, 21, 52)
       call transform (fi, a, b, c, d, 12,  6, 53)
       call transform (fi, d, a, b, c,  3, 10, 54)
       call transform (fi, c, d, a, b, 10, 15, 55)
       call transform (fi, b, c, d, a,  1, 21, 56)
       call transform (fi, a, b, c, d,  8,  6, 57)
       call transform (fi, d, a, b, c, 15, 10, 58)
       call transform (fi, c, d, a, b,  6, 15, 59)
       call transform (fi, b, c, d, a, 13, 21, 60)
       call transform (fi, a, b, c, d,  4,  6, 61)
       call transform (fi, d, a, b, c, 11, 10, 62)
       call transform (fi, c, d, a, b,  2, 15, 63)
       call transform (fi, b, c, d, a,  9, 21, 64)
       a = a + aa
       b = b + bb
       c = c + cc
       d = d + dd
       bl => bl%next
    end do
    s = digest_string ([a, b, c, d])
  contains
  \LA{}MD5: Internal subroutine transform\RA{}
  end subroutine message_digest

\eatline
\nwendcode{}\nwbegindocs{144}And this is the actual transformation that depends on one of the
previous functions, four words, and three integers.  The implicit
arguments are {\tt{}x}, the word from the message to digest, and {\tt{}t},
the entry in the predefined table.
\nwenddocs{}\nwbegincode{145}\moddef{MD5: Internal subroutine transform}\endmoddef
  subroutine transform (f, a, b, c, d, k, s, i)
    interface
       function f (x, y, z)
         import word32_t
         type(word32_t), intent(in) :: x, y, z
         type(word32_t) :: f
       end function f
    end interface
    type(word32_t), intent(inout) :: a
    type(word32_t), intent(in) :: b, c, d
    integer, intent(in) :: k, s, i
    a = b + ishftc (a + f(b, c, d) + x(k) + t(i), s)
  end subroutine transform
  
\eatline
\nwendcode{}\nwbegindocs{146}\nwdocspar
\subsection{User interface}
\nwenddocs{}\nwbegincode{147}\moddef{MD5: public}\endmoddef
  public :: md5sum
\nwendcode{}\nwbegincode{148}\moddef{MD5: interfaces}\plusendmoddef
  interface md5sum
     module procedure md5sum_from_string
     module procedure md5sum_from_unit
  end interface
\eatline
\nwendcode{}\nwbegindocs{149}This function computes the MD5 sum of the input string and returns it
as a 32-character string
\nwenddocs{}\nwbegincode{150}\moddef{MD5: procedures}\plusendmoddef
  function md5sum_from_string (s) result (digest)
    character(len=*), intent(in) :: s
    character(len=32) :: digest
    type(message_t) :: m
    call message_append_string (m, s)
    call message_pad (m)
    call message_digest (m, digest)
    call message_clear (m)
  end function md5sum_from_string

\eatline
\nwendcode{}\nwbegindocs{151}This funct. reads from unit u (an unformmated sequence of 
integers) and computes the MD5 sum.
\nwenddocs{}\nwbegincode{152}\moddef{MD5: procedures}\plusendmoddef
  function md5sum_from_unit (u) result (digest)
    integer, intent(in) :: u
    character(len=32) :: digest
    type(message_t) :: m
    character :: char
    integer :: iostat
    READ_CHARS: do
       read (u, "(A)", advance="no", iostat=iostat)  char
       select case (iostat)
       case (0)
          call message_append_string (m, char)
       case (EOR)
          call message_append_string (m, LF)
       case (EOF)
          exit READ_CHARS
       case default
          call msg_fatal &
               ("Computing MD5 sum: I/O error while reading from scratch file")
       end select
    end do READ_CHARS
    call message_pad (m)
    call message_digest (m, digest)
    call message_clear (m)
  end function md5sum_from_unit

\eatline
\nwendcode{}\nwbegindocs{153}\nwdocspar
\subsection{Unit tests}
\nwenddocs{}\nwbegincode{154}\moddef{MD5: public}\plusendmoddef
  public :: md5_test
\nwendcode{}\nwbegincode{155}\moddef{MD5: procedures}\plusendmoddef
  subroutine md5_test (u, results)
    integer, intent(in) :: u
    type(test_results_t), intent(inout) :: results
  \LA{}MD5: execute tests\RA{}  
  end subroutine md5_test
  
\eatline
\nwendcode{}\nwbegindocs{156}This function checks the implementation by computing the checksum of
certain strings and comparing them with the known values.
\nwenddocs{}\nwbegincode{157}\moddef{MD5: execute tests}\endmoddef
  call test (md5_1, "md5_1", &
       "check MD5 sums", &
       u, results)
\nwendcode{}\nwbegincode{158}\moddef{MD5: tests}\endmoddef
  subroutine md5_1 (u)
    integer, intent(in) :: u
    character(32) :: s
    integer, parameter :: n = 7
    integer :: i
    character(80), dimension(n) :: teststring
    data teststring(1) /""/
    data teststring(2) /"a"/
    data teststring(3) /"abc"/
    data teststring(4) /"message digest"/
    data teststring(5) /"abcdefghijklmnopqrstuvwxyz"/
    data teststring(6) /"ABCDEFGHIJKLMNOPQRSTUVWXYZabcdefghijklmnopqrstuvwxyz0123456789"/
    data teststring(7) /"12345678901234567890123456789012345678901234567890123456789012345678901234567890"/
    character(32), dimension(n) :: result
    data result(1) /"D41D8CD98F00B204E9800998ECF8427E"/
    data result(2) /"0CC175B9C0F1B6A831C399E269772661"/
    data result(3) /"900150983CD24FB0D6963F7D28E17F72"/
    data result(4) /"F96B697D7CB7938D525A2F31AAF161D0"/
    data result(5) /"C3FCD3D76192E4007DFB496CCA67E13B"/
    data result(6) /"D174AB98D277D9F5A5611C2C9F419D9F"/
    data result(7) /"57EDF4A22BE3C955AC49DA2E2107B67A"/
    
    write (u, "(A)")  "* Test output: MD5"
    write (u, "(A)")  "*   Purpose: test MD5 sums"
    write (u, "(A)")        
        
    do i = 1, n
       write (u, "(A)") "MD5 test string = " // '"'// &
            trim (teststring(i)) // '"'
       s = md5sum (trim (teststring(i)))
       write (u, "(A)") "MD5 check sum   = " // trim (s)
       write (u, "(A)") "Ref check sum   = " // result(i)
       if (s == result(i)) then
          call msg_message ("=> ok", u)
       else
          call msg_message ("=> MD5 sum self-test failed", u)
       end if
    end do
    call msg_message ("=============================================================================|", unit=u)
  end subroutine md5_1

\eatline
\nwendcode{}\nwbegindocs{159}\nwdocspar
%%%%%%%%%%%%%%%%%%%%%%%%%%%%%%%%%%%%%%%%%%%%%%%%%%%%%%%%%%%%%%%%%%%%%%%%
\section{Permutations}
Permute arrays of integers (of specific kind).
\nwenddocs{}\nwbegincode{160}\moddef{\code{}permutations.f90\edoc{}}\endmoddef
\LA{}File header\RA{}

module permutations

  use kinds, only: TC

\LA{}Standard module head\RA{}

\LA{}Permutations: public\RA{}

\LA{}Permutations: types\RA{}

\LA{}Permutations: interfaces\RA{}

contains

\LA{}Permutations: procedures\RA{}

end module permutations
\eatline
\nwendcode{}\nwbegindocs{161}\nwdocspar
\subsection{Permutations}
A permutation is an array of integers.  Each integer between one and
{\tt{}size} should occur exactly once.
\nwenddocs{}\nwbegincode{162}\moddef{Permutations: public}\endmoddef
  public :: permutation_t
\nwendcode{}\nwbegincode{163}\moddef{Permutations: types}\endmoddef
  type :: permutation_t
     private
     integer, dimension(:), allocatable :: p
  end type permutation_t

\eatline
\nwendcode{}\nwbegindocs{164}\nwdocspar
Initialize with the identity permutation.
\nwenddocs{}\nwbegincode{165}\moddef{Permutations: public}\plusendmoddef
  public :: permutation_init
  public :: permutation_final
\nwendcode{}\nwbegincode{166}\moddef{Permutations: procedures}\endmoddef
  elemental subroutine permutation_init (p, size)
    type(permutation_t), intent(inout) :: p
    integer, intent(in) :: size
    integer :: i
    allocate (p%p (size))
    forall (i = 1:size)
       p%p(i) = i
    end forall
  end subroutine permutation_init

  elemental subroutine permutation_final (p)
    type(permutation_t), intent(inout) :: p
    deallocate (p%p)
  end subroutine permutation_final

\eatline
\nwendcode{}\nwbegindocs{167}I/O:
\nwenddocs{}\nwbegincode{168}\moddef{Permutations: public}\plusendmoddef
  public :: permutation_write
\nwendcode{}\nwbegincode{169}\moddef{Permutations: procedures}\plusendmoddef
  subroutine permutation_write (p, u)
    type(permutation_t), intent (in) :: p
    integer, intent(in) :: u
    integer :: i
    do i = 1, size (p%p)
       if (size (p%p) < 10) then
          write (u,"(1x,I1)", advance="no") p%p(i)
       else
          write (u,"(1x,I3)", advance="no") p%p(i)
       end if
    end do
    write (u, *)
  end subroutine permutation_write

\eatline
\nwendcode{}\nwbegindocs{170}\nwdocspar
Administration:
\nwenddocs{}\nwbegincode{171}\moddef{Permutations: public}\plusendmoddef
  public :: permutation_size
\nwendcode{}\nwbegincode{172}\moddef{Permutations: procedures}\plusendmoddef
  elemental function permutation_size (perm) result (s)
    type(permutation_t), intent(in) :: perm
    integer :: s
    s = size (perm%p)
  end function permutation_size

\eatline
\nwendcode{}\nwbegindocs{173}Extract an entry in a permutation.
\nwenddocs{}\nwbegincode{174}\moddef{Permutations: public}\plusendmoddef
  public :: permute
\nwendcode{}\nwbegincode{175}\moddef{Permutations: procedures}\plusendmoddef
  elemental function permute (i, p) result (j)
    integer, intent(in) :: i
    type(permutation_t), intent(in) :: p
    integer :: j
    if (i > 0 .and. i <= size (p%p)) then
       j = p%p(i)
    else
       j = 0
    end if
  end function permute

\eatline
\nwendcode{}\nwbegindocs{176}\nwdocspar
Check whether a permutation is valid: Each integer in the range occurs
exactly once.
\nwenddocs{}\nwbegincode{177}\moddef{Permutations: public}\plusendmoddef
  public :: permutation_ok
\nwendcode{}\nwbegincode{178}\moddef{Permutations: procedures}\plusendmoddef
  elemental function permutation_ok (perm) result (ok)
    type(permutation_t), intent(in) :: perm
    logical :: ok
    integer :: i
    logical, dimension(:), allocatable :: set
    ok = .true.
    allocate (set (size (perm%p)))
    set = .false.
    do i = 1, size (perm%p)
       ok = (perm%p(i) > 0 .and. perm%p(i) <= size (perm%p))
       if (.not.ok) return
       set(perm%p(i)) = .true.
    end do
    ok = all (set)
  end function permutation_ok

\eatline
\nwendcode{}\nwbegindocs{179}Find the permutation that transforms the second array into the first
one.  We assume that this is possible and unique and all bounds are
set correctly.

This cannot be elemental.
\nwenddocs{}\nwbegincode{180}\moddef{Permutations: public}\plusendmoddef
  public :: permutation_find
\nwendcode{}\nwbegincode{181}\moddef{Permutations: procedures}\plusendmoddef
  subroutine permutation_find (perm, a1, a2)
    type(permutation_t), intent(inout) :: perm
    integer, dimension(:), intent(in) :: a1, a2
    integer :: i, j
    if (allocated (perm%p))  deallocate (perm%p)
    allocate (perm%p (size (a1)))
    do i = 1, size (a1)
       do j = 1, size (a2)
          if (a1(i) == a2(j)) then
             perm%p(i) = j
             exit
          end if
          perm%p(i) = 0
       end do
    end do
  end subroutine permutation_find

\eatline
\nwendcode{}\nwbegindocs{182}\nwdocspar
Find all permutations that transform an array of integers into
itself.  The resulting permutation list is allocated with the correct
length and filled.

The first step is to count the number of different entries in
{\tt{}code}.  Next, we scan {\tt{}code} again and assign a mask to each
different entry, true for all identical entries.  Finally, we
recursively permute the identity for each possible mask.

The permutation is done as follows:  A list of all permutations of the
initial one with respect to the current mask is generated, then the
permutations are generated in turn for each permutation in this list
with the next mask.  The result is always stored back into the main
list, starting from the end of the current list.
\nwenddocs{}\nwbegincode{183}\moddef{Permutations: public}\plusendmoddef
  public :: permutation_array_make
\nwendcode{}\nwbegincode{184}\moddef{Permutations: procedures}\plusendmoddef
  subroutine permutation_array_make (pa, code)
    type(permutation_t), dimension(:), allocatable, intent(out) :: pa
    integer, dimension(:), intent(in) :: code
    logical, dimension(size(code)) :: mask
    logical, dimension(:,:), allocatable :: imask
    integer, dimension(:), allocatable :: n_i
    type(permutation_t) :: p_init
    type(permutation_t), dimension(:), allocatable :: p_tmp
    integer :: psize, i, j, k, n_different, n, nn_k
    psize = size (code)
    mask = .true.
    n_different = 0
    do i=1, psize
       if (mask(i)) then
          n_different = n_different + 1
          mask = mask .and. (code /= code(i))
       end if
    end do
    allocate (imask(psize, n_different), n_i(n_different))
    mask = .true.
    k = 0
    do i=1, psize
       if (mask(i)) then
          k = k + 1
          imask(:,k) = (code == code(i))
          n_i(k) = factorial (count(imask(:,k)))
          mask = mask .and. (code /= code(i))
       end if
    end do
    n = product (n_i)
    allocate (pa (n))
    call permutation_init (p_init, psize)
    pa(1) = p_init
    nn_k = 1
    do k = 1, n_different
       allocate (p_tmp (n_i(k)))
       do i = nn_k, 1, -1
          call permutation_array_with_mask (p_tmp, imask(:,k), pa(i))
          do j = n_i(k), 1, -1
             pa((i-1)*n_i(k) + j) = p_tmp(j)
          end do
       end do
       deallocate (p_tmp)
       nn_k = nn_k * n_i(k)
    end do
    call permutation_final (p_init)
    deallocate (imask, n_i)
  end subroutine permutation_array_make

\eatline
\nwendcode{}\nwbegindocs{185}Make a list of permutations of the elements marked true in the
{\tt{}mask} array.  The final permutation list must be allocated with the
correct length ($n!$).  The third argument is the initial
permutation to start with, which must have the same length as the
{\tt{}mask} array (this is not checked).
\nwenddocs{}\nwbegincode{186}\moddef{Permutations: procedures}\plusendmoddef
  subroutine permutation_array_with_mask (pa, mask, p_init)
    type(permutation_t), dimension(:), intent(inout) :: pa
    logical, dimension(:), intent(in) :: mask
    type(permutation_t), intent(in) :: p_init
    integer :: plen
    integer :: i, ii, j, fac_i, k, x
    integer, dimension(:), allocatable :: index 
    plen = size (pa)
    allocate (index(count(mask)))
    ii = 0
    do i = 1, size (mask)
       if (mask(i)) then
          ii = ii + 1
          index(ii) = i
       end if
    end do
    pa = p_init
    ii = 0
    fac_i = 1
    do i = 1, size (mask)
       if (mask(i)) then
          ii = ii + 1
          fac_i = fac_i * ii
          x = permute (i, p_init)
          do j = 1, plen
             k = ii - mod (((j-1)*fac_i)/plen, ii)
             call insert (pa(j), x, k, ii, index)
          end do
       end if
    end do
    deallocate (index)
  contains
    subroutine insert (p, x, k, n, index)
      type(permutation_t), intent(inout) :: p
      integer, intent(in) :: x, k, n
      integer, dimension(:), intent(in) :: index
      integer :: i
      do i = n, k+1, -1
         p%p(index(i)) = p%p(index(i-1))
      end do
      p%p(index(k)) = x
    end subroutine insert
  end subroutine permutation_array_with_mask

\eatline
\nwendcode{}\nwbegindocs{187}The factorial function is needed for pre-determining the number of
permutations that will be generated:
\nwenddocs{}\nwbegincode{188}\moddef{Permutations: procedures}\plusendmoddef
  function factorial (n) result (f)
    integer, intent(in) :: n
    integer :: f
    integer :: i
    f = 1
    do i=2, abs(n)
       f = f*i
    end do
  end function factorial

\eatline
\nwendcode{}\nwbegindocs{189}\nwdocspar
\subsection{Operations on binary codes}
Binary codes are needed for phase-space trees.  Since the permutation
function uses permutations, and no other special type is involved, we
put the functions here.

This is needed for phase space trees: permute bits in a tree binary
code.  If no permutation is given, leave as is.  (We may want to
access the permutation directly here if this is efficiency-critical.)
\nwenddocs{}\nwbegincode{190}\moddef{Permutations: public}\plusendmoddef
  public :: tc_permute
\nwendcode{}\nwbegincode{191}\moddef{Permutations: procedures}\plusendmoddef
  function tc_permute (k, perm, mask_in) result (pk)
    integer(TC), intent(in) :: k, mask_in
    type(permutation_t), intent(in) :: perm
    integer(TC) :: pk
    integer :: i
    pk = iand (k, mask_in)
    do i = 1, size (perm%p)
       if (btest(k,i-1))  pk = ibset (pk, perm%p(i)-1)
    end do
  end function tc_permute

\eatline
\nwendcode{}\nwbegindocs{192}\nwdocspar
This routine returns the number of set bits in the tree code value
{\tt{}k}.  Hence, it is the number of externals connected to the current
line.  If {\tt{}mask} is present, the complement of the tree code is also
considered, and the smaller number is returned.  This gives the true
distance from the external states, taking into account the initial
particles.  The complement number is increased by one, since for a
scattering diagram the vertex with the sum of all final-state codes is
still one point apart from the initial particles.
\nwenddocs{}\nwbegincode{193}\moddef{Permutations: public}\plusendmoddef
  public :: tc_decay_level
\nwendcode{}\nwbegincode{194}\moddef{Permutations: interfaces}\endmoddef
  interface tc_decay_level
     module procedure decay_level_simple
     module procedure decay_level_complement
  end interface
\eatline
\nwendcode{}\nwbegincode{195}\moddef{Permutations: procedures}\plusendmoddef
  function decay_level_complement (k, mask) result (l)
    integer(TC), intent(in) :: k, mask
    integer :: l
    l = min (decay_level_simple (k), &
         &   decay_level_simple (ieor (k, mask)) + 1)
  end function decay_level_complement

  function decay_level_simple (k) result(l)
    integer(TC), intent(in) :: k
    integer :: l
    integer :: i
    l = 0
    do i=0, bit_size(k)-1
       if (btest(k,i)) l = l+1
    end do
  end function decay_level_simple

\eatline
\nwendcode{}\nwbegindocs{196}\nwdocspar
%%%%%%%%%%%%%%%%%%%%%%%%%%%%%%%%%%%%%%%%%%%%%%%%%%%%%%%%%%%%%%%%%%%%%%%%
\section{Sorting}
This small module provides functions for sorting integer or real
arrays.
\nwenddocs{}\nwbegincode{197}\moddef{\code{}sorting.f90\edoc{}}\endmoddef
\LA{}File header\RA{}

module sorting

\LA{}Use kinds\RA{}
  use unit_tests
  use diagnostics
  
\LA{}Standard module head\RA{}

\LA{}Sorting: public\RA{}

\LA{}Sorting: interfaces\RA{}

contains

\LA{}Sorting: procedures\RA{}

\LA{}Sorting: tests\RA{}

end module sorting
\eatline
\nwendcode{}\nwbegindocs{198}\nwdocspar
\subsection{Implementation}
The {\tt{}sort} function returns, for a given integer or real array, the
array sorted by increasing value.  The current implementation is
\emph{mergesort}, which has $O(n\ln n)$ behavior in all cases, and is
stable for elements of equal value.

The {\tt{}sort{\char95}abs} variant sorts by increasing absolute value, where for
identical absolute value, the positive number comes first.
\nwenddocs{}\nwbegincode{199}\moddef{Sorting: public}\endmoddef
  public :: sort
  public :: sort_abs
\nwendcode{}\nwbegincode{200}\moddef{Sorting: interfaces}\endmoddef
  interface sort
     module procedure sort_int
     module procedure sort_real
  end interface

  interface sort_abs
     module procedure sort_int_abs
  end interface

\eatline
\nwendcode{}\nwbegindocs{201}This variant of integer sort returns
\nwenddocs{}\nwbegindocs{202}The body is identical, just the interface differs.
\nwenddocs{}\nwbegincode{203}\moddef{Sorting: procedures}\endmoddef
  function sort_int (val_in) result (val)
    integer, dimension(:), intent(in) :: val_in
    integer, dimension(size(val_in)) :: val
  \LA{}Sorting: sort\RA{}
  end function sort_int

  function sort_real (val_in) result (val)
    real(default), dimension(:), intent(in) :: val_in
    real(default), dimension(size(val_in)) :: val
  \LA{}Sorting: sort\RA{}
  end function sort_real

  function sort_int_abs (val_in) result (val)
    integer, dimension(:), intent(in) :: val_in
    integer, dimension(size(val_in)) :: val
  \LA{}Sorting: sort abs\RA{}
  end function sort_int_abs

\eatline
\nwendcode{}\nwbegincode{204}\moddef{Sorting: sort}\endmoddef
  val = val_in( order (val_in) )
\nwendcode{}\nwbegincode{205}\moddef{Sorting: sort abs}\endmoddef
  val = val_in( order_abs (val_in) )
\nwendcode{}\nwbegindocs{206}The {\tt{}order} function returns, for a given integer or real array, the
array of indices of the elements sorted by increasing value.
\nwenddocs{}\nwbegincode{207}\moddef{Sorting: public}\plusendmoddef
  public :: order
  public :: order_abs
\nwendcode{}\nwbegincode{208}\moddef{Sorting: interfaces}\plusendmoddef
  interface order
     module procedure order_int
     module procedure order_real
  end interface

  interface order_abs
     module procedure order_int_abs
  end interface

\eatline
\nwendcode{}\nwbegincode{209}\moddef{Sorting: procedures}\plusendmoddef
  function order_int (val) result (idx)
    integer, dimension(:), intent(in) :: val
    integer, dimension(size(val)) :: idx
  \LA{}Sorting: order\RA{}
  end function order_int

  function order_real (val) result (idx)
    real(default), dimension(:), intent(in) :: val
    integer, dimension(size(val)) :: idx
  \LA{}Sorting: order\RA{}
  end function order_real

  function order_int_abs (val) result (idx)
    integer, dimension(:), intent(in) :: val
    integer, dimension(size(val)) :: idx
  \LA{}Sorting: order abs\RA{}
  end function order_int_abs

\eatline
\nwendcode{}\nwbegindocs{210}We start by individual elements, merge them to pairs, merge those to
four-element subarrays, and so on.  The last subarray can extend only
up to the original array bound, of course, and the second of the
subarrays to merge should contain at least one element.
\nwenddocs{}\nwbegincode{211}\moddef{Sorting: order}\endmoddef
\LA{}Sorting: order1\RA{}
        call merge (idx(b1:e2), idx(b1:e1), idx(b2:e2), val)
\LA{}Sorting: order2\RA{}
\nwendcode{}\nwbegindocs{212}\nwdocspar
\nwenddocs{}\nwbegincode{213}\moddef{Sorting: order abs}\endmoddef
\LA{}Sorting: order1\RA{}
        call merge_abs (idx(b1:e2), idx(b1:e1), idx(b2:e2), val)
\LA{}Sorting: order2\RA{}
\nwendcode{}\nwbegindocs{214}\nwdocspar
\nwenddocs{}\nwbegincode{215}\moddef{Sorting: order1}\endmoddef
  integer :: n, i, s, b1, b2, e1, e2
  n = size (idx)
  forall (i = 1:n)
     idx(i) = i
  end forall
  s = 1
  do while (s < n)
     do b1 = 1, n-s, 2*s
        b2 = b1 + s
        e1 = b2 - 1
        e2 = min (e1 + s, n)
\nwendcode{}\nwbegindocs{216}\nwdocspar
\nwenddocs{}\nwbegincode{217}\moddef{Sorting: order2}\endmoddef
     end do
     s = 2 * s
  end do
\nwendcode{}\nwbegindocs{218}The merging step does the actual sorting.  We take two sorted array
sections and merge them to a sorted result array.  We are working on
the indices, and comparing is done by taking the associated {\tt{}val}
which is real or integer.
\nwenddocs{}\nwbegincode{219}\moddef{Sorting: interfaces}\plusendmoddef
  interface merge
     module procedure merge_int
     module procedure merge_real
  end interface

  interface merge_abs
     module procedure merge_int_abs
  end interface

\eatline
\nwendcode{}\nwbegincode{220}\moddef{Sorting: procedures}\plusendmoddef
  subroutine merge_int (res, src1, src2, val)
    integer, dimension(:), intent(out) :: res
    integer, dimension(:), intent(in) :: src1, src2
    integer, dimension(:), intent(in) :: val
    integer, dimension(size(res)) :: tmp
  \LA{}Sorting: merge\RA{}
  end subroutine merge_int
    
  subroutine merge_real (res, src1, src2, val)
    integer, dimension(:), intent(out) :: res
    integer, dimension(:), intent(in) :: src1, src2
    real(default), dimension(:), intent(in) :: val
    integer, dimension(size(res)) :: tmp
  \LA{}Sorting: merge\RA{}
  end subroutine merge_real

  subroutine merge_int_abs (res, src1, src2, val)
    integer, dimension(:), intent(out) :: res
    integer, dimension(:), intent(in) :: src1, src2
    integer, dimension(:), intent(in) :: val
    integer, dimension(size(res)) :: tmp
  \LA{}Sorting: merge abs\RA{}
  end subroutine merge_int_abs
    
\eatline
\nwendcode{}\nwbegincode{221}\moddef{Sorting: merge}\endmoddef
\LA{}Sorting: merge1\RA{}
     if (val(src1(i1)) <= val(src2(i2))) then
\LA{}Sorting: merge2\RA{}
\nwendcode{}\nwbegindocs{222}We keep the elements if the absolute values are strictly ordered.
If they are equal in magnitude, we keep them if the larger value
comes first, or if they are equal.
\nwenddocs{}\nwbegincode{223}\moddef{Sorting: merge abs}\endmoddef
\LA{}Sorting: merge1\RA{}
     if (abs (val(src1(i1))) < abs (val(src2(i2))) .or. &
        (abs (val(src1(i1))) == abs (val(src2(i2))) .and. &
        val(src1(i1)) >= val(src2(i2)))) then
\LA{}Sorting: merge2\RA{}
\nwendcode{}\nwbegindocs{224}\nwdocspar
\nwenddocs{}\nwbegincode{225}\moddef{Sorting: merge1}\endmoddef
  integer :: i1, i2, i
  i1 = 1
  i2 = 1
  do i = 1, size (tmp)
\nwendcode{}\nwbegindocs{226}\nwdocspar
\nwenddocs{}\nwbegincode{227}\moddef{Sorting: merge2}\endmoddef
        tmp(i) = src1(i1);  i1 = i1 + 1
        if (i1 > size (src1)) then
           tmp(i+1:) = src2(i2:)
           exit
        end if
     else
        tmp(i) = src2(i2);  i2 = i2 + 1
        if (i2 > size (src2)) then
           tmp(i+1:) = src1(i1:)
           exit
        end if
     end if
  end do
  res = tmp
\nwendcode{}\nwbegindocs{228}\nwdocspar
\subsection{Concatenating arrays}
Not precisely a sorting function, but useful: Concatenate two arrays.
\nwenddocs{}\nwbegincode{229}\moddef{Sorting: public}\plusendmoddef
  public :: concat
\nwendcode{}\nwbegincode{230}\moddef{Sorting: interfaces}\plusendmoddef
  interface concat
     module procedure concat_int
     module procedure concat_real
  end interface

\eatline
\nwendcode{}\nwbegincode{231}\moddef{Sorting: procedures}\plusendmoddef
  function concat_int (val1, val2) result (val12)
    integer, dimension(:), intent(in) :: val1, val2
    integer, dimension(size(val1)+size(val2)) :: val12
  \LA{}Sorting: concat\RA{}
  end function concat_int

  function concat_real (val1, val2) result (val12)
    real(default), dimension(:), intent(in) :: val1, val2
    integer, dimension(size(val1)+size(val2)) :: val12
  \LA{}Sorting: concat\RA{}
  end function concat_real

\eatline
\nwendcode{}\nwbegincode{232}\moddef{Sorting: concat}\endmoddef
  val12(:size(val1)) = val1
  val12(size(val1)+1:) = val2
\nwendcode{}\nwbegindocs{233}\nwdocspar
\subsection{Unit tests}
\nwenddocs{}\nwbegincode{234}\moddef{Sorting: public}\plusendmoddef
  public :: sorting_test
\nwendcode{}\nwbegincode{235}\moddef{Sorting: procedures}\plusendmoddef
  subroutine sorting_test (u, results)
    integer, intent(in) :: u
    type(test_results_t), intent(inout) :: results
  \LA{}Sorting: execute tests\RA{}
  end subroutine sorting_test

\eatline
\nwendcode{}\nwbegindocs{236}This checks whether the sorting routine works correctly.  
\nwenddocs{}\nwbegincode{237}\moddef{Sorting: execute tests}\endmoddef
  call test (sorting_1, "sorting_1", &
       "check sorting routines", &
       u, results)
\nwendcode{}\nwbegincode{238}\moddef{Sorting: tests}\endmoddef
  subroutine sorting_1 (u)
    integer, intent(in) :: u
    integer, parameter :: NMAX = 10
    real(default), dimension(NMAX) :: rval
    integer, dimension(NMAX) :: ival
    real, dimension(NMAX,NMAX) :: harvest_r
    integer, dimension(NMAX,NMAX) :: harvest_i
    integer, dimension(NMAX,NMAX) :: harvest_a
    integer :: i, j
    harvest_r(:, 1) = [0.9976, 0., 0., 0., 0., 0., 0., 0., 0., 0.]
    harvest_r(:, 2) = [0.5668, 0.9659, 0., 0., 0., 0., 0., 0., 0., 0.]
    harvest_r(:, 3) = [0.7479, 0.3674, 0.4806, 0., 0., 0., 0., 0., 0., &
         0.]
    harvest_r(:, 4) = [0.0738, 0.0054, 0.3471, 0.3422, 0., 0., 0., 0., &
         0., 0.]
    harvest_r(:, 5) = [0.2180, 0.1332, 0.9005, 0.3868, 0.4455, 0., 0., &
         0., 0., 0.]
    harvest_r(:, 6) = [0.6619, 0.0161, 0.6509, 0.6464, 0.3230, &
         0.8557, 0., 0., 0., 0.]
    harvest_r(:, 7) = [0.4013, 0.2069, 0.9685, 0.5984, 0.6730, &
         0.4569, 0.3300, 0., 0., 0.]
    harvest_r(:, 8) = [0.1004, 0.7555, 0.6057, 0.7190, 0.8973, &
         0.6582, 0.1507, 0.6123, 0., 0.]
    harvest_r(:, 9) = [0.9787, 0.9991, 0.2568, 0.5509, 0.6590, &
         0.5540, 0.9778, 0.9019, 0.6579, 0.]
    harvest_r(:,10) = [0.7289, 0.4025, 0.9286, 0.1478, 0.6745, &
         0.7696, 0.3393, 0.1158, 0.6144, 0.8206]
    
    harvest_i(:, 1) = [18, 0, 0, 0, 0, 0, 0, 0, 0, 0]
    harvest_i(:, 2) = [14, 9, 0, 0, 0, 0, 0, 0, 0, 0]
    harvest_i(:, 3) = [ 7, 8,11, 0, 0, 0, 0, 0, 0, 0]
    harvest_i(:, 4) = [19,19,14,19, 0, 0, 0, 0, 0, 0]
    harvest_i(:, 5) = [ 1,14,15,18,14, 0, 0, 0, 0, 0]
    harvest_i(:, 6) = [16,11, 1, 9,11, 2, 0, 0, 0, 0]
    harvest_i(:, 7) = [11,10,17, 6,13,13,10, 0, 0, 0]
    harvest_i(:, 8) = [ 5, 1, 2,10, 7, 0,15,12, 0, 0]
    harvest_i(:, 9) = [15,19, 2, 6,11, 0, 2, 4, 2, 0]
    harvest_i(:,10) = [ 1, 4, 8, 4,11, 0, 8, 7,19,13]
    
    harvest_a(:, 1) = [-6,  0,  0,  0,  0,  0,  0,  0,  0,  0]
    harvest_a(:, 2) = [-8, -9,  0,  0,  0,  0,  0,  0,  0,  0]
    harvest_a(:, 3) = [ 4, -3,  3,  0,  0,  0,  0,  0,  0,  0]
    harvest_a(:, 4) = [-6,  6,  2, -2,  0,  0,  0,  0,  0,  0]
    harvest_a(:, 5) = [ 1, -2,  0, -6,  8,  0,  0,  0,  0,  0]
    harvest_a(:, 6) = [-2, -1, -8, -5,  8, -5,  0,  0,  0,  0]
    harvest_a(:, 7) = [-9,  0, -6,  2,  5,  3,  2,  0,  0,  0]
    harvest_a(:, 8) = [-5, -7,  6,  7, -3,  0, -7,  4,  0,  0]
    harvest_a(:, 9) = [ 5,  0, -1, -7,  5,  2,  7, -3,  3,  0]
    harvest_a(:,10) = [-9,  2, -6,  3, -9,  5,  5,  7,  5, -9]


    write (u, "(A)")  "* Test output: Sorting"
    write (u, "(A)")  "*   Purpose: test sorting routines"
    write (u, "(A)")      
    
    write (u, "(A)")  "* Sorting real values:"
    
    do i = 1, NMAX
       write (u, "(A)")
       rval(:i) = harvest_r(:i,i)
       write (u, "(10(1x,F7.4))") rval(:i)
       rval(:i) = sort (rval(:i))
       write (u, "(10(1x,F7.4))") rval(:i)
       do j = i, 2, -1
          if (rval(j)-rval(j-1) < 0) &
             write (u, "(A)") "*** Sorting failure. ***"
       end do
    end do
    
    write (u, "(A)")
    write (u, "(A)") "* Sorting integer values:"
    
    do i = 1, NMAX
       write (u, "(A)")
       ival(:i) = harvest_i(:i,i)
       write (u, "(10(1x,I2))") ival(:i)
       ival(:i) = sort (ival(:i))
       write (u, "(10(1x,I2))") ival(:i)
       do j = i, 2, -1
          if (ival(j)-ival(j-1) < 0) &
             write (u, "(A)")  "*** Sorting failure. ***"
       end do
    end do
    
    write (u, "(A)")
    write (u, "(A)") "* Sorting integer values by absolute value:"
    
    do i = 1, NMAX
       write (u, "(A)")
       ival(:i) = harvest_a(:i,i)
       write (u, "(10(1x,I2))") ival(:i)
       ival(:i) = sort_abs (ival(:i))
       write (u, "(10(1x,I2))") ival(:i)
       do j = i, 2, -1
          if (abs(ival(j))-abs(ival(j-1)) < 0 .or. &
               (abs(ival(j))==abs(ival(j-1))) .and. ival(j)>ival(j-1)) &
             write (u, "(A)")  "*** Sorting failure. ***"
       end do
    end do
    
    write (u, "(A)")
    write (u, "(A)")  "* Test output end: sorting_1"    

  end subroutine sorting_1

\eatline
\nwendcode{}\nwbegindocs{239}\nwdocspar
\nwenddocs{}\nwfilename{../../../src/parsing/parsing.nw}\nwbegindocs{0}% -*- ess-noweb-default-code-mode: f90-mode; noweb-default-code-mode: f90-mode; -*- 
% WHIZARD code as NOWEB source: text handling and parsing

\nwenddocs{}\nwbegindocs{1}\nwdocspar
%%%%%%%%%%%%%%%%%%%%%%%%%%%%%%%%%%%%%%%%%%%%%%%%%%%%%%%%%%%%%%%%%%%%%%%%
\chapter{Text handling}

\whizard\ has to handle complex structures in input (and output) data.
Doing this in a generic and transparent way requires a generic lexer
and parser.  The necessary modules are implemented here:
\begin{description}
\item[ifiles]
  Implementation of line-oriented internal files in a more flexible
  way (linked lists of variable-length strings) than the Fortran builtin
  features.
\item[lexers]
  Read text and transform it into a token stream.
\item[syntax\_rules]
  Define the rules for interpreting tokens, to be used by the WHIZARD parser.
\item[parser]
  Categorize tokens (keyword, string, number etc.) and use a set of
  syntax rules to transform the input into a parse tree.
\item[xml]
  Read and parse XML text, separate from the WHIZARD parser.
\end{description}
\nwenddocs{}\nwbegindocs{2}\nwdocspar
\clearpage
%%%%%%%%%%%%%%%%%%%%%%%%%%%%%%%%%%%%%%%%%%%%%%%%%%%%%%%%%%%%%%%%%%%%%%%%
\section{Internal files}
The internal files introduced here ({\tt{}ifile}) are a replacement for
the built-in internal files, which are fixed-size arrays of
fixed-length character strings.  The {\tt{}ifile} type is a doubly-linked
list of variable-length character strings with line numbers.
\nwenddocs{}\nwbegincode{3}\moddef{\code{}ifiles.f90\edoc{}}\endmoddef
\LA{}File header\RA{}

module ifiles

\LA{}Use strings\RA{}
  use io_units
  use system_defs, only: EOF

\LA{}Standard module head\RA{}

\LA{}Ifiles: public\RA{}

\LA{}Ifiles: types\RA{}

\LA{}Ifiles: interfaces\RA{}

contains

\LA{}Ifiles: subroutines\RA{}

end module ifiles
\eatline
\nwendcode{}\nwbegindocs{4}\nwdocspar
\subsection{The line type}
The line entry type is for internal use, it is the list entry to be
collected in an {\tt{}ifile} object.
\nwenddocs{}\nwbegincode{5}\moddef{Ifiles: types}\endmoddef
  type :: line_entry_t
     private
     type(line_entry_t), pointer :: previous => null ()
     type(line_entry_t), pointer :: next => null ()
     type(string_t) :: string
     integer :: index
  end type line_entry_t

\eatline
\nwendcode{}\nwbegindocs{6}Create a new list entry, given a varying string as input.  The line
number and pointers are not set, these make sense only within an
{\tt{}ifile}.
\nwenddocs{}\nwbegincode{7}\moddef{Ifiles: subroutines}\endmoddef
  subroutine line_entry_create (line, string)
    type(line_entry_t), pointer :: line
    type(string_t), intent(in) :: string
    allocate (line)
    line%string = string
  end subroutine line_entry_create

\eatline
\nwendcode{}\nwbegindocs{8}Destroy a single list entry: Since the pointer components should not
be deallocated explicitly, just deallocate the object itself.
\nwenddocs{}\nwbegincode{9}\moddef{Ifiles: subroutines}\plusendmoddef
  subroutine line_entry_destroy (line)
    type(line_entry_t), pointer :: line
    deallocate (line)
  end subroutine line_entry_destroy

\eatline
\nwendcode{}\nwbegindocs{10}\nwdocspar
\subsection{The ifile type}
The internal file is a linked list of line entries.
\nwenddocs{}\nwbegincode{11}\moddef{Ifiles: public}\endmoddef
  public :: ifile_t
\nwendcode{}\nwbegincode{12}\moddef{Ifiles: types}\plusendmoddef
  type :: ifile_t
     private
     type(line_entry_t), pointer :: first => null ()
     type(line_entry_t), pointer :: last => null ()
     integer :: n_lines = 0
  end type ifile_t

\eatline
\nwendcode{}\nwbegindocs{13}We need no explicit initializer, but a routine which recursively
deallocates the contents may be appropriate.  After this, existing
line pointers may become undefined, so they should be nullified before
the file is destroyed.
\nwenddocs{}\nwbegincode{14}\moddef{Ifiles: public}\plusendmoddef
  public :: ifile_clear
\nwendcode{}\nwbegincode{15}\moddef{Ifiles: subroutines}\plusendmoddef
  subroutine ifile_clear (ifile)
    type(ifile_t), intent(inout) :: ifile
    type(line_entry_t), pointer :: current
    do while (associated (ifile%first))
       current => ifile%first
       ifile%first => current%next
       call line_entry_destroy (current)
    end do
    nullify (ifile%last)
    ifile%n_lines = 0
  end subroutine ifile_clear

\eatline
\nwendcode{}\nwbegindocs{16}The finalizer is just an alias for the above.
\nwenddocs{}\nwbegincode{17}\moddef{Ifiles: public}\plusendmoddef
  public :: ifile_final
\nwendcode{}\nwbegincode{18}\moddef{Ifiles: interfaces}\endmoddef
  interface ifile_final
     module procedure ifile_clear
  end interface
\eatline
\nwendcode{}\nwbegindocs{19}\nwdocspar
\subsection{I/O on ifiles}
Fill an ifile from an ordinary external file, i.e., I/O unit.  If the
ifile is not empty, the old contents will be destroyed.  We can read
a fixed-length character string, an ISO varying string, an ordinary
internal file (character-string array), or from an external unit.  In
the latter case, lines are appended until EOF is reached.  Finally,
there is a variant which reads from another ifile, effectively copying it.
\nwenddocs{}\nwbegincode{20}\moddef{Ifiles: public}\plusendmoddef
  public :: ifile_read
\nwendcode{}\nwbegincode{21}\moddef{Ifiles: interfaces}\plusendmoddef
  interface ifile_read
     module procedure ifile_read_from_string
     module procedure ifile_read_from_char
     module procedure ifile_read_from_unit
     module procedure ifile_read_from_char_array
     module procedure ifile_read_from_ifile
  end interface
\nwendcode{}\nwbegincode{22}\moddef{Ifiles: subroutines}\plusendmoddef
  subroutine ifile_read_from_string (ifile, string)
    type(ifile_t), intent(inout) :: ifile
    type(string_t), intent(in) :: string
    call ifile_clear (ifile)
    call ifile_append (ifile, string)
  end subroutine ifile_read_from_string

  subroutine ifile_read_from_char (ifile, char)
    type(ifile_t), intent(inout) :: ifile
    character(*), intent(in) :: char
    call ifile_clear (ifile)
    call ifile_append (ifile, char)
  end subroutine ifile_read_from_char

  subroutine ifile_read_from_char_array (ifile, char)
    type(ifile_t), intent(inout) :: ifile
    character(*), dimension(:), intent(in) :: char
    call ifile_clear (ifile)
    call ifile_append (ifile, char)
  end subroutine ifile_read_from_char_array
    
  subroutine ifile_read_from_unit (ifile, unit, iostat)
    type(ifile_t), intent(inout) :: ifile
    integer, intent(in) :: unit
    integer, intent(out), optional :: iostat
    call ifile_clear (ifile)
    call ifile_append (ifile, unit, iostat)
  end subroutine ifile_read_from_unit
    
  subroutine ifile_read_from_ifile (ifile, ifile_in)
    type(ifile_t), intent(inout) :: ifile
    type(ifile_t), intent(in) :: ifile_in
    call ifile_clear (ifile)
    call ifile_append (ifile, ifile_in)
  end subroutine ifile_read_from_ifile
    
\eatline
\nwendcode{}\nwbegindocs{23}Append to an ifile.  The same as reading, but without resetting the
ifile.  In addition, there is a routine for appending a whole ifile.
\nwenddocs{}\nwbegincode{24}\moddef{Ifiles: public}\plusendmoddef
  public :: ifile_append
\nwendcode{}\nwbegincode{25}\moddef{Ifiles: interfaces}\plusendmoddef
  interface ifile_append
     module procedure ifile_append_from_string
     module procedure ifile_append_from_char
     module procedure ifile_append_from_unit
     module procedure ifile_append_from_char_array
     module procedure ifile_append_from_ifile
  end interface
\nwendcode{}\nwbegincode{26}\moddef{Ifiles: subroutines}\plusendmoddef
  subroutine ifile_append_from_string (ifile, string)
    type(ifile_t), intent(inout) :: ifile
    type(string_t), intent(in) :: string
    type(line_entry_t), pointer :: current
    call line_entry_create (current, string)
    current%index = ifile%n_lines + 1
    if (associated (ifile%last)) then
       current%previous => ifile%last
       ifile%last%next => current
    else
       ifile%first => current
    end if
    ifile%last => current
    ifile%n_lines = current%index
  end subroutine ifile_append_from_string

  subroutine ifile_append_from_char (ifile, char)
    type(ifile_t), intent(inout) :: ifile
    character(*), intent(in) :: char
    call ifile_append_from_string (ifile, var_str (trim (char)))
  end subroutine ifile_append_from_char

  subroutine ifile_append_from_char_array (ifile, char)
    type(ifile_t), intent(inout) :: ifile
    character(*), dimension(:), intent(in) :: char
    integer :: i
    do i = 1, size (char)
       call ifile_append_from_string (ifile, var_str (trim (char(i))))
    end do
  end subroutine ifile_append_from_char_array
    
  subroutine ifile_append_from_unit (ifile, unit, iostat)
    type(ifile_t), intent(inout) :: ifile
    integer, intent(in) :: unit
    integer, intent(out), optional :: iostat
    type(string_t) :: buffer
    integer :: ios
    ios = 0
    READ_LOOP: do
       call get (unit, buffer, iostat = ios)
       if (ios == EOF .or. ios > 0)  exit READ_LOOP
       call ifile_append_from_string (ifile, buffer)
    end do READ_LOOP
    if (present (iostat)) then
       iostat = ios
    else if (ios > 0) then
       call get (unit, buffer)  ! trigger error again
    end if
  end subroutine ifile_append_from_unit
    
  subroutine ifile_append_from_ifile (ifile, ifile_in)
    type(ifile_t), intent(inout) :: ifile
    type(ifile_t), intent(in) :: ifile_in
    type(line_entry_t), pointer :: current
    current => ifile_in%first
    do while (associated (current))
       call ifile_append_from_string (ifile, current%string)
       current => current%next
    end do
  end subroutine ifile_append_from_ifile
    
\eatline
\nwendcode{}\nwbegindocs{27}Write the ifile contents to an external unit
\nwenddocs{}\nwbegincode{28}\moddef{Ifiles: public}\plusendmoddef
  public :: ifile_write
\nwendcode{}\nwbegincode{29}\moddef{Ifiles: subroutines}\plusendmoddef
  subroutine ifile_write (ifile, unit, iostat)
    type(ifile_t), intent(in) :: ifile
    integer, intent(in), optional :: unit
    integer, intent(out), optional :: iostat
    integer :: u
    type(line_entry_t), pointer :: current
    u = given_output_unit (unit);  if (u < 0)  return
    current => ifile%first
    do while (associated (current))
       call put_line (u, current%string, iostat)
       current => current%next
    end do
  end subroutine ifile_write

\eatline
\nwendcode{}\nwbegindocs{30}Convert the ifile to an array of strings, which is allocated by this
function:
\nwenddocs{}\nwbegincode{31}\moddef{Ifiles: public}\plusendmoddef
  public :: ifile_to_string_array
\nwendcode{}\nwbegincode{32}\moddef{Ifiles: subroutines}\plusendmoddef
  subroutine ifile_to_string_array (ifile, string)
    type(ifile_t), intent(in) :: ifile
    type(string_t), dimension(:), intent(inout), allocatable :: string
    type(line_entry_t), pointer :: current
    integer :: i
    allocate (string (ifile_get_length (ifile)))
    current => ifile%first
    do i = 1, ifile_get_length (ifile)
       string(i) = current%string
       current => current%next
    end do
  end subroutine ifile_to_string_array

\eatline
\nwendcode{}\nwbegindocs{33}\nwdocspar
\subsection{Ifile tools}
\nwenddocs{}\nwbegincode{34}\moddef{Ifiles: public}\plusendmoddef
  public :: ifile_get_length
\nwendcode{}\nwbegincode{35}\moddef{Ifiles: subroutines}\plusendmoddef
  function ifile_get_length (ifile) result (length)
    integer :: length
    type(ifile_t), intent(in) :: ifile
    length = ifile%n_lines
  end function ifile_get_length

\eatline
\nwendcode{}\nwbegindocs{36}\nwdocspar
\subsection{Line pointers}
Instead of the implicit pointer used in ordinary file access, we
define explicit pointers, so there can be more than one at a time.
\nwenddocs{}\nwbegincode{37}\moddef{Ifiles: public}\plusendmoddef
  public :: line_p
\nwendcode{}\nwbegincode{38}\moddef{Ifiles: types}\plusendmoddef
  type :: line_p
     private
     type(line_entry_t), pointer :: p => null ()
  end type line_p

\eatline
\nwendcode{}\nwbegindocs{39}Assign a file pointer to the first or last line in an ifile:
\nwenddocs{}\nwbegincode{40}\moddef{Ifiles: public}\plusendmoddef
  public :: line_init
\nwendcode{}\nwbegincode{41}\moddef{Ifiles: subroutines}\plusendmoddef
  subroutine line_init (line, ifile, back)
    type(line_p), intent(inout) :: line
    type(ifile_t), intent(in) :: ifile
    logical, intent(in), optional :: back
    if (present (back)) then
       if (back) then
          line%p => ifile%last
       else
          line%p => ifile%first
       end if
    else
       line%p => ifile%first
    end if
  end subroutine line_init

\eatline
\nwendcode{}\nwbegindocs{42}Remove the pointer association:
\nwenddocs{}\nwbegincode{43}\moddef{Ifiles: public}\plusendmoddef
  public :: line_final
\nwendcode{}\nwbegincode{44}\moddef{Ifiles: subroutines}\plusendmoddef
  subroutine line_final (line)
    type(line_p), intent(inout) :: line
    nullify (line%p)
  end subroutine line_final

\eatline
\nwendcode{}\nwbegindocs{45}Go one step forward
\nwenddocs{}\nwbegincode{46}\moddef{Ifiles: public}\plusendmoddef
  public :: line_advance
\nwendcode{}\nwbegincode{47}\moddef{Ifiles: subroutines}\plusendmoddef
  subroutine line_advance (line)
    type(line_p), intent(inout) :: line
    if (associated (line%p))  line%p => line%p%next
  end subroutine line_advance

\eatline
\nwendcode{}\nwbegindocs{48}Go one step backward
\nwenddocs{}\nwbegincode{49}\moddef{Ifiles: public}\plusendmoddef
  public :: line_backspace
\nwendcode{}\nwbegincode{50}\moddef{Ifiles: subroutines}\plusendmoddef
  subroutine line_backspace (line)
    type(line_p), intent(inout) :: line
    if (associated (line%p))  line%p => line%p%previous
  end subroutine line_backspace

\eatline
\nwendcode{}\nwbegindocs{51}Check whether we are accessing a valid line
\nwenddocs{}\nwbegincode{52}\moddef{Ifiles: public}\plusendmoddef
  public :: line_is_associated
\nwendcode{}\nwbegincode{53}\moddef{Ifiles: subroutines}\plusendmoddef
  function line_is_associated (line) result (ok)
    logical :: ok
    type(line_p), intent(in) :: line
    ok = associated (line%p)
  end function line_is_associated

\eatline
\nwendcode{}\nwbegindocs{54}\nwdocspar
\subsection{Access lines via pointers}
We do not need the ifile as an argument to these functions, because
the {\tt{}line} type will point to an existing ifile.  
\nwenddocs{}\nwbegincode{55}\moddef{Ifiles: public}\plusendmoddef
  public :: line_get_string
\nwendcode{}\nwbegincode{56}\moddef{Ifiles: subroutines}\plusendmoddef
  function line_get_string (line) result (string)
    type(string_t) :: string
    type(line_p), intent(in) :: line
    if (associated (line%p)) then
       string = line%p%string
    else
       string = ""
    end if
  end function line_get_string

\eatline
\nwendcode{}\nwbegindocs{57}Variant where the line pointer is advanced after reading.  
\nwenddocs{}\nwbegincode{58}\moddef{Ifiles: public}\plusendmoddef
  public :: line_get_string_advance
\nwendcode{}\nwbegincode{59}\moddef{Ifiles: subroutines}\plusendmoddef
  function line_get_string_advance (line) result (string)
    type(string_t) :: string
    type(line_p), intent(inout) :: line
    if (associated (line%p)) then
       string = line%p%string
       call line_advance (line)
    else
       string = ""
    end if
  end function line_get_string_advance

\eatline
\nwendcode{}\nwbegincode{60}\moddef{Ifiles: public}\plusendmoddef
  public :: line_get_index
\nwendcode{}\nwbegincode{61}\moddef{Ifiles: subroutines}\plusendmoddef
  function line_get_index (line) result (index)
    integer :: index
    type(line_p), intent(in) :: line
    if (associated (line%p)) then
       index = line%p%index
    else
       index = 0
    end if
  end function line_get_index

\eatline
\nwendcode{}\nwbegincode{62}\moddef{Ifiles: public}\plusendmoddef
  public :: line_get_length
\nwendcode{}\nwbegincode{63}\moddef{Ifiles: subroutines}\plusendmoddef
  function line_get_length (line) result (length)
    integer :: length
    type(line_p), intent(in) :: line
    if (associated (line%p)) then
       length = len (line%p%string)
    else
       length = 0
    end if
  end function line_get_length

\eatline
\nwendcode{}\nwbegindocs{64}\nwdocspar
\clearpage
%%%%%%%%%%%%%%%%%%%%%%%%%%%%%%%%%%%%%%%%%%%%%%%%%%%%%%%%%%%%%%%%%%%%%%%%
\section{Lexer}
The lexer purpose is to read from a line-separated character input
stream (usually a file) and properly chop the stream into lexemes
(tokens).  [The parser will transform lexemes into meaningful tokens,
to be stored in a parse tree, therefore we do not use the term 'token'
here.]  The input is read line-by-line, but interpreted free-form,
except for quotes and the comment syntax.  (Fortran 2003 would allow
us to use a stream type for reading.)

In an object-oriented approach, we can dynamically create and destroy
lexers, including the lexer setup.

The main lexer function is to return a lexeme according to the basic
lexer rules (quotes, comments, whitespace, special classes).  There is
also a routine to write back a lexeme to the input stream (but only
once).

For the rules, we separate the possible characters into classes.
Whitespace usually consists of blank, tab, and line-feed, where any
number of consecutive whitespace is equivalent to one.  Quoted strings
are enclosed by a pair of quote characters, possibly multiline.
Comments are similar to quotes, but interpreted as whitespace.
Numbers are identified (not distinguishing real and integer) but not
interpreted. Other character classes make up identifiers.

\nwenddocs{}\nwbegincode{65}\moddef{\code{}lexers.f90\edoc{}}\endmoddef
\LA{}File header\RA{}

module lexers

\LA{}Use strings\RA{}
  use io_units
  use string_utils
  use unit_tests
  use system_defs, only: EOF, EOR
  use system_defs, only: LF
  use system_defs, only: WHITESPACE_CHARS, LCLETTERS, UCLETTERS, DIGITS
  use diagnostics
  use ifiles, only: ifile_t 
  use ifiles, only: line_p, line_is_associated, line_init, line_final
  use ifiles, only: line_get_string_advance
  
\LA{}Standard module head\RA{}

\LA{}Lexer: public\RA{}

\LA{}Lexer: parameters\RA{}

\LA{}Lexer: types\RA{}

\LA{}Lexer: interfaces\RA{}

contains

\LA{}Lexer: procedures\RA{}

\LA{}Lexer: tests\RA{}

end module lexers
\eatline
\nwendcode{}\nwbegindocs{66}\nwdocspar
\subsection{Input streams}
For flexible input, we define a generic stream type that refers to
either an external file, an external unit which is already open, a
string, an {\tt{}ifile} object (internal file, i.e., string list), or a
line pointer to an {\tt{}ifile} object.  The stream type actually follows
the idea of a formatted external file, which is line-oriented.  Thus,
the stream reader always returns a whole record (input line).

Note that only in the string version, the stream contents are stored
inside the stream object.  In the {\tt{}ifile} version, the stream
contains only the line pointer, while in the external-file case, the
line pointer is implicitly created by the runtime library.
\nwenddocs{}\nwbegincode{67}\moddef{Lexer: public}\endmoddef
  public :: stream_t
\nwendcode{}\nwbegincode{68}\moddef{Lexer: types}\endmoddef
  type :: stream_t
     type(string_t), pointer :: filename => null ()
     integer, pointer :: unit => null ()
     type(string_t), pointer :: string => null ()
     type(ifile_t), pointer :: ifile => null ()
     type(line_p), pointer :: line => null ()
     integer :: record = 0
     logical :: eof = .false.
  end type stream_t

\eatline
\nwendcode{}\nwbegindocs{69}The initializers refer to the specific version.  The stream should
be undefined before calling this.
\nwenddocs{}\nwbegincode{70}\moddef{Lexer: public}\plusendmoddef
  public :: stream_init
\nwendcode{}\nwbegincode{71}\moddef{Lexer: interfaces}\endmoddef
  interface stream_init
     module procedure stream_init_filename
     module procedure stream_init_unit
     module procedure stream_init_string
     module procedure stream_init_ifile
     module procedure stream_init_line
  end interface
\nwendcode{}\nwbegincode{72}\moddef{Lexer: procedures}\endmoddef
  subroutine stream_init_filename (stream, filename)
    type(stream_t), intent(out) :: stream
    character(*), intent(in) :: filename
    integer :: unit
    unit = free_unit ()
    open (unit=unit, file=filename, status="old", action="read")
    call stream_init_unit (stream, unit)
    allocate (stream%filename)
    stream%filename = filename
  end subroutine stream_init_filename

  subroutine stream_init_unit (stream, unit)
    type(stream_t), intent(out) :: stream
    integer, intent(in) :: unit
    allocate (stream%unit)
    stream%unit = unit
    stream%eof = .false.
  end subroutine stream_init_unit
    
  subroutine stream_init_string (stream, string)
    type(stream_t), intent(out) :: stream
    type(string_t), intent(in) :: string
    allocate (stream%string)
    stream%string = string
  end subroutine stream_init_string
    
  subroutine stream_init_ifile (stream, ifile)
    type(stream_t), intent(out) :: stream
    type(ifile_t), intent(in) :: ifile
    type(line_p) :: line
    call line_init (line, ifile)
    call stream_init_line (stream, line)
    allocate (stream%ifile)
    stream%ifile = ifile
  end subroutine stream_init_ifile
    
  subroutine stream_init_line (stream, line)
    type(stream_t), intent(out) :: stream
    type(line_p), intent(in) :: line
    allocate (stream%line)
    stream%line = line
  end subroutine stream_init_line
    
\eatline
\nwendcode{}\nwbegindocs{73}The finalizer restores the initial state.  If an external file was
opened, it is closed.
\nwenddocs{}\nwbegincode{74}\moddef{Lexer: public}\plusendmoddef
  public :: stream_final
\nwendcode{}\nwbegincode{75}\moddef{Lexer: procedures}\plusendmoddef
  subroutine stream_final (stream)
    type(stream_t), intent(inout) :: stream
    if (associated (stream%filename)) then
       close (stream%unit)
       deallocate (stream%unit)
       deallocate (stream%filename)
    else if (associated (stream%unit)) then
       deallocate (stream%unit)
    else if (associated (stream%string)) then
       deallocate (stream%string)
    else if (associated (stream%ifile)) then
       call line_final (stream%line)
       deallocate (stream%line)
       deallocate (stream%ifile)
    else if (associated (stream%line)) then
       call line_final (stream%line)
       deallocate (stream%line)
    end if
  end subroutine stream_final

\eatline
\nwendcode{}\nwbegindocs{76}This returns the next record from the input stream.  Depending on
the stream type, the stream pointers are modified: Reading from
external unit, the external file is advanced (implicitly).  Reading
from string, the string is replaced by an empty string.  Reading from
{\tt{}ifile}, the line pointer is advanced.  Note that the {\tt{}iostat}
argument is mandatory.
\nwenddocs{}\nwbegincode{77}\moddef{Lexer: public}\plusendmoddef
  public :: stream_get_record
\nwendcode{}\nwbegincode{78}\moddef{Lexer: procedures}\plusendmoddef
  subroutine stream_get_record (stream, string, iostat)
    type(stream_t), intent(inout) :: stream
    type(string_t), intent(out) :: string
    integer, intent(out) :: iostat
    if (associated (stream%unit)) then
       if (stream%eof) then
          iostat = EOF
       else
          call get (stream%unit, string, iostat=iostat)
          if (iostat == EOR) then
             iostat = 0
             stream%record = stream%record + 1
          end if
          if (iostat == EOF) then
             iostat = 0
             stream%eof = .true.
             if (len (string) /= 0) stream%record = stream%record + 1
          end if
       end if
    else if (associated (stream%string)) then
       if (len (stream%string) /= 0) then
          string = stream%string
          stream%string = ""
          iostat = 0
          stream%record = stream%record + 1
       else
          string = ""
          iostat = EOF
       end if
    else if (associated (stream%line)) then
       if (line_is_associated (stream%line)) then
          string = line_get_string_advance (stream%line)
          iostat = 0
          stream%record = stream%record + 1
       else
          string = ""
          iostat = EOF
       end if
    else
       call msg_bug (" Attempt to read from uninitialized input stream")
    end if
  end subroutine stream_get_record

\eatline
\nwendcode{}\nwbegindocs{79}Return the current stream source as a message string.
\nwenddocs{}\nwbegincode{80}\moddef{Lexer: public}\plusendmoddef
  public :: stream_get_source_info_string
\nwendcode{}\nwbegincode{81}\moddef{Lexer: procedures}\plusendmoddef
  function stream_get_source_info_string (stream) result (string)
    type(string_t) :: string
    type(stream_t), intent(in) :: stream
    character(20) :: buffer
    if (associated (stream%filename)) then
       string = "File '" // stream%filename // "' (unit = "
       write (buffer, "(I0)")  stream%unit
       string = string // trim (buffer) // ")"
    else if (associated (stream%unit)) then
       write (buffer, "(I0)")  stream%unit
       string = "Unit " // trim (buffer)
    else if (associated (stream%string)) then
       string = "Input string"
    else if (associated (stream%ifile) .or. associated (stream%line)) then
       string = "Internal file"
    else
       string = ""
    end if
  end function stream_get_source_info_string

\eatline
\nwendcode{}\nwbegindocs{82}Return the index of the record just read as a message string.
\nwenddocs{}\nwbegincode{83}\moddef{Lexer: public}\plusendmoddef
  public :: stream_get_record_info_string
\nwendcode{}\nwbegincode{84}\moddef{Lexer: procedures}\plusendmoddef
  function stream_get_record_info_string (stream) result (string)
    type(string_t) :: string
    type(stream_t), intent(in) :: stream
    character(20) :: buffer
    string = stream_get_source_info_string (stream)
    if (string /= "")  string = string // ", "
    write (buffer, "(I0)")  stream%record
    string = string // "line " // trim (buffer)
  end function stream_get_record_info_string

\eatline
\nwendcode{}\nwbegindocs{85}\nwdocspar
\subsection{Keyword list}
The lexer should be capable of identifying a token as a known
keyword.  To this end, we store a list of keywords:
\nwenddocs{}\nwbegincode{86}\moddef{Lexer: public}\plusendmoddef
  public :: keyword_list_t
\nwendcode{}\nwbegincode{87}\moddef{Lexer: types}\plusendmoddef
  type :: keyword_entry_t
     private
     type(string_t) :: string
     type(keyword_entry_t), pointer :: next => null ()
  end type keyword_entry_t

  type :: keyword_list_t
     private
     type(keyword_entry_t), pointer :: first => null ()
     type(keyword_entry_t), pointer :: last => null ()
  end type keyword_list_t

\eatline
\nwendcode{}\nwbegindocs{88}Add a new string to the keyword list, unless it is already there:
\nwenddocs{}\nwbegincode{89}\moddef{Lexer: public}\plusendmoddef
  public :: keyword_list_add
\nwendcode{}\nwbegincode{90}\moddef{Lexer: procedures}\plusendmoddef
  subroutine keyword_list_add (keylist, string)
    type(keyword_list_t), intent(inout) :: keylist
    type(string_t), intent(in) :: string
    type(keyword_entry_t), pointer :: k_entry_new
    if (.not. keyword_list_contains (keylist, string)) then
       allocate (k_entry_new)
       k_entry_new%string = string
       if (associated (keylist%first)) then
          keylist%last%next => k_entry_new
       else
          keylist%first => k_entry_new
       end if
       keylist%last => k_entry_new
    end if
  end subroutine keyword_list_add

\eatline
\nwendcode{}\nwbegindocs{91}Return true if a string is a keyword.
\nwenddocs{}\nwbegincode{92}\moddef{Lexer: public}\plusendmoddef
  public :: keyword_list_contains
\nwendcode{}\nwbegincode{93}\moddef{Lexer: procedures}\plusendmoddef
  function keyword_list_contains (keylist, string) result (found)
    type(keyword_list_t), intent(in) :: keylist
    type(string_t), intent(in) :: string
    logical :: found
    found = .false.
    call check_rec (keylist%first)
  contains
    recursive subroutine check_rec (k_entry)
      type(keyword_entry_t), pointer :: k_entry
      if (associated (k_entry)) then
         if (k_entry%string /= string) then
            call check_rec (k_entry%next)
         else
            found = .true.
         end if
      end if
    end subroutine check_rec
  end function keyword_list_contains

\eatline
\nwendcode{}\nwbegindocs{94}Write the keyword list
\nwenddocs{}\nwbegincode{95}\moddef{Lexer: public}\plusendmoddef
  public :: keyword_list_write
\nwendcode{}\nwbegincode{96}\moddef{Lexer: interfaces}\plusendmoddef
  interface keyword_list_write
     module procedure keyword_list_write_unit
  end interface
\nwendcode{}\nwbegincode{97}\moddef{Lexer: procedures}\plusendmoddef
  subroutine keyword_list_write_unit (keylist, unit)
    type(keyword_list_t), intent(in) :: keylist
    integer, intent(in) :: unit
    write (unit, "(A)") "Keyword list:"
    if (associated (keylist%first)) then
       call keyword_write_rec (keylist%first)
       write (unit, *)
    else
       write (unit, "(1x,A)") "[empty]"
    end if
  contains
    recursive subroutine keyword_write_rec (k_entry)
      type(keyword_entry_t), intent(in), pointer :: k_entry
      if (associated (k_entry)) then
         write (unit, "(1x,A)", advance="no")  char (k_entry%string)
         call keyword_write_rec (k_entry%next)
      end if
    end subroutine keyword_write_rec
  end subroutine keyword_list_write_unit

\eatline
\nwendcode{}\nwbegindocs{98}Clear the keyword list
\nwenddocs{}\nwbegincode{99}\moddef{Lexer: public}\plusendmoddef
  public :: keyword_list_final
\nwendcode{}\nwbegincode{100}\moddef{Lexer: procedures}\plusendmoddef
  subroutine keyword_list_final (keylist)
    type(keyword_list_t), intent(inout) :: keylist
    call keyword_destroy_rec (keylist%first)
    nullify (keylist%last)
  contains
    recursive subroutine keyword_destroy_rec (k_entry)
      type(keyword_entry_t), pointer :: k_entry
      if (associated (k_entry)) then
         call keyword_destroy_rec (k_entry%next)
         deallocate (k_entry)
      end if
    end subroutine keyword_destroy_rec
  end subroutine keyword_list_final

\eatline
\nwendcode{}\nwbegindocs{101}\nwdocspar
\subsection{Lexeme templates}
This type is handled like a rudimentary regular expression.  It
determines the lexer behavior when matching a string.  The actual
objects made from this type and the corresponding matching routines
are listed below.
\nwenddocs{}\nwbegincode{102}\moddef{Lexer: types}\plusendmoddef
  type :: template_t
     private
     integer :: type
     character(256) :: charset1, charset2
     integer :: len1, len2
  end type template_t

\eatline
\nwendcode{}\nwbegindocs{103}These are the types that valid lexemes can have:
\nwenddocs{}\nwbegincode{104}\moddef{Lexer: public}\plusendmoddef
  public :: T_KEYWORD, T_IDENTIFIER, T_QUOTED, T_NUMERIC
\nwendcode{}\nwbegincode{105}\moddef{Lexer: parameters}\endmoddef
  integer, parameter :: T_KEYWORD = 1
  integer, parameter :: T_IDENTIFIER = 2, T_QUOTED = 3, T_NUMERIC = 4
\eatline
\nwendcode{}\nwbegindocs{106}These are special types:
\nwenddocs{}\nwbegincode{107}\moddef{Lexer: parameters}\plusendmoddef
  integer, parameter :: EMPTY = 0, WHITESPACE = 10
  integer, parameter :: NO_MATCH = 11, IO_ERROR = 12, OVERFLOW = 13
  integer, parameter :: UNMATCHED_QUOTE = 14
\eatline
\nwendcode{}\nwbegindocs{108}In addition, we have {\tt{}EOF} which is a negative integer, normally $-1$.
\nwenddocs{}\nwbegindocs{109}Printout for debugging:
\nwenddocs{}\nwbegincode{110}\moddef{Lexer: procedures}\plusendmoddef
  subroutine lexeme_type_write (type, unit)
    integer, intent(in) :: type
    integer, intent(in) :: unit
    select case (type)
    case (EMPTY);       write(unit,"(A)",advance="no") " EMPTY      "
    case (WHITESPACE);  write(unit,"(A)",advance="no") " WHITESPACE "
    case (T_IDENTIFIER);write(unit,"(A)",advance="no") " IDENTIFIER "
    case (T_QUOTED);    write(unit,"(A)",advance="no") " QUOTED     "
    case (T_NUMERIC);   write(unit,"(A)",advance="no") " NUMERIC    "
    case (IO_ERROR);    write(unit,"(A)",advance="no") " IO_ERROR   "
    case (OVERFLOW);    write(unit,"(A)",advance="no") " OVERFLOW   "
    case (UNMATCHED_QUOTE);    write(unit,"(A)",advance="no") " UNMATCHEDQ "
    case (NO_MATCH);    write(unit,"(A)",advance="no") " NO_MATCH   "
    case (EOF);         write(unit,"(A)",advance="no") " EOF        "
    case default;       write(unit,"(A)",advance="no") " [illegal]  "
    end select
  end subroutine lexeme_type_write

  subroutine template_write (tt, unit)
    type(template_t), intent(in) :: tt
    integer, intent(in) :: unit
    call lexeme_type_write (tt%type, unit)
    write (unit, "(A)", advance="no") "'" // tt%charset1(1:tt%len1) // "'"
    write (unit, "(A)", advance="no") " '" // tt%charset2(1:tt%len2) // "'"
  end subroutine template_write
    
\eatline
\nwendcode{}\nwbegindocs{111}The matching functions all return the number of matched characters in
the provided string.  If this number is zero, the match has failed.

The {\tt{}template} functions are declared {\tt{}pure} because they appear
in {\tt{}forall} loops below.

A template for whitespace:
\nwenddocs{}\nwbegincode{112}\moddef{Lexer: procedures}\plusendmoddef
  pure function template_whitespace (chars) result (tt)
    character(*), intent(in) :: chars
    type(template_t) :: tt
    tt = template_t (WHITESPACE, chars, "", len (chars), 0)
  end function template_whitespace

\eatline
\nwendcode{}\nwbegindocs{113}Just match the string against the character set.
\nwenddocs{}\nwbegincode{114}\moddef{Lexer: procedures}\plusendmoddef
  subroutine match_whitespace (tt, s, n)
    type(template_t), intent(in) :: tt
    character(*), intent(in) :: s
    integer, intent(out) :: n
    n = verify (s, tt%charset1(1:tt%len1)) - 1
    if (n < 0)  n = len (s)
  end subroutine match_whitespace

\eatline
\nwendcode{}\nwbegindocs{115}A template for normal identifiers.  To match, a lexeme should have a
first character in class {\tt{}chars1} and an arbitrary number of further
characters in class {\tt{}chars2}.  If the latter is empty, we are
looking for a single-character lexeme.
\nwenddocs{}\nwbegincode{116}\moddef{Lexer: procedures}\plusendmoddef
  pure function template_identifier (chars1, chars2) result (tt)
    character(*), intent(in) :: chars1, chars2
    type(template_t) :: tt
    tt = template_t (T_IDENTIFIER, chars1, chars2, len(chars1), len(chars2))
  end function template_identifier

\eatline
\nwendcode{}\nwbegindocs{117}Here, the first letter must match, the others may or may not.
\nwenddocs{}\nwbegincode{118}\moddef{Lexer: procedures}\plusendmoddef
  subroutine match_identifier (tt, s, n)
    type(template_t), intent(in) :: tt
    character(*), intent(in) :: s
    integer, intent(out) :: n
    if (verify (s(1:1), tt%charset1(1:tt%len1)) == 0) then
       n = verify (s(2:), tt%charset2(1:tt%len2))
       if (n == 0)  n = len (s)
    else
       n = 0
    end if
  end subroutine match_identifier

\eatline
\nwendcode{}\nwbegindocs{119}A template for quoted strings.  The same template applies for
comments.  The first character set indicates the
left quote (could be a sequence of several characters), the second one
the matching right quote.
\nwenddocs{}\nwbegincode{120}\moddef{Lexer: procedures}\plusendmoddef
  pure function template_quoted (chars1, chars2) result (tt)
    character(*), intent(in) :: chars1, chars2
    type(template_t) :: tt
    tt = template_t (T_QUOTED, chars1, chars2, len (chars1), len (chars2))
  end function template_quoted

\eatline
\nwendcode{}\nwbegindocs{121}Here, the beginning of the string must exactly match the first
character set, then we look for the second one.  If found, return.
If there is a first quote but no second one, return a negative number,
indicating this error condition.
\nwenddocs{}\nwbegincode{122}\moddef{Lexer: procedures}\plusendmoddef
  subroutine match_quoted (tt, s, n, range)
    type(template_t), intent(in) :: tt
    character(*), intent(in) :: s
    integer, intent(out) :: n
    integer, dimension(2), intent(out) :: range
    character(tt%len1) :: ch1
    character(tt%len2) :: ch2
    integer :: i
    ch1 = tt%charset1
    if (s(1:tt%len1) == ch1) then
       ch2 = tt%charset2
       do i = tt%len1 + 1, len (s) - tt%len2 + 1
          if (s(i:i+tt%len2-1) == ch2) then
             n = i + tt%len2 - 1
             range(1) = tt%len1 + 1
             range(2) = i - 1
             return
          end if
       end do
       n = -1
       range = 0
    else
       n = 0
       range = 0
    end if
  end subroutine match_quoted

\eatline
\nwendcode{}\nwbegindocs{123}A template for real numbers.  The first character set is the set of
allowed exponent letters.  In accordance with the other functions we
return the lexeme as a string but do not read it.
\nwenddocs{}\nwbegincode{124}\moddef{Lexer: procedures}\plusendmoddef
  pure function template_numeric (chars) result (tt)
    character(*), intent(in) :: chars
    type(template_t) :: tt
    tt = template_t (T_NUMERIC, chars, "", len (chars), 0)
  end function template_numeric

\eatline
\nwendcode{}\nwbegindocs{125}A numeric lexeme may be real or integer.  We purposely do not allow
for a preceding sign.  If the number is followed by an exponent, this
is included, otherwise the rest is ignored.  

There is a possible pitfall with this behavior: while the string
{\tt{}1e3} will be interpreted as a single number, the analogous string
{\tt{}1a3} will be split into the number {\tt{}1} and an identifier {\tt{}a3}.
There is no easy way around such an ambiguity.  We should make sure
that the syntax does not contain identifiers like {\tt{}a3} or {\tt{}e3}.
\nwenddocs{}\nwbegincode{126}\moddef{Lexer: procedures}\plusendmoddef
  subroutine match_numeric (tt, s, n)
    type(template_t), intent(in) :: tt
    character(*), intent(in) :: s
    integer, intent(out) :: n
    integer :: i, n0
    character(10), parameter :: digits = "0123456789"
    character(2), parameter :: signs = "-+"
    n = verify (s, digits) - 1
    if (n < 0) then
       n = 0
       return
    else if (s(n+1:n+1) == ".") then
       i = verify (s(n+2:), digits) - 1
       if (i < 0) then
          n = len (s)
          return
       else if (i > 0 .or. n > 0) then
          n = n + 1 + i
       end if
    end if
    n0 = n
    if (n > 0) then
       if (verify (s(n+1:n+1), tt%charset1(1:tt%len1)) == 0) then
          n = n + 1
          if (verify (s(n+1:n+1), signs) == 0)  n = n + 1
          i = verify (s(n+1:), digits) - 1
          if (i < 0) then
             n = len (s)
          else if (i == 0) then
             n = n0
          else
             n = n + i
          end if
       end if
    end if
  end subroutine match_numeric

\eatline
\nwendcode{}\nwbegindocs{127}The generic matching routine.  With Fortran 2003 we would define
separate types and use a SELECT TYPE instead.
\nwenddocs{}\nwbegincode{128}\moddef{Lexer: procedures}\plusendmoddef
  subroutine match_template (tt, s, n, range)
    type(template_t), intent(in) :: tt
    character(*), intent(in) :: s
    integer, intent(out) :: n
    integer, dimension(2), intent(out) :: range
    select case (tt%type)
    case (WHITESPACE)
       call match_whitespace (tt, s, n)
       range = 0
    case (T_IDENTIFIER)
       call match_identifier (tt, s, n)
       range(1) = 1
       range(2) = len_trim (s)
    case (T_QUOTED)
       call match_quoted (tt, s, n, range)
    case (T_NUMERIC)
       call match_numeric (tt, s, n)
       range(1) = 1
       range(2) = len_trim (s)
    case default
       call msg_bug ("Invalid lexeme template encountered")
    end select
  end subroutine match_template

\eatline
\nwendcode{}\nwbegindocs{129}Match against an array of templates.  Return the index of the first
template that matches together with the number of characters matched
and the range of the relevant substring.  If all fails, these numbers
are zero.
\nwenddocs{}\nwbegincode{130}\moddef{Lexer: procedures}\plusendmoddef
  subroutine match (tt, s, n, range, ii)
    type(template_t), dimension(:), intent(in) :: tt
    character(*), intent(in) :: s
    integer, intent(out) :: n
    integer, dimension(2), intent(out) :: range
    integer, intent(out) :: ii
    integer :: i
    do i = 1, size (tt)
       call match_template (tt(i), s, n, range)
       if (n /= 0) then
          ii = i
          return
       end if
    end do
    n = 0
    ii = 0
  end subroutine match

\eatline
\nwendcode{}\nwbegindocs{131}\nwdocspar
\subsection{The lexer setup}
This object contains information about character classes.  As said
above, one class consists of quoting chars (matching left and right),
another one of comment chars (similar), a class of whitespace, and
several classes of characters that make up identifiers.  When creating
the lexer setup, the character classes are transformed into lexeme
templates which are to be matched in a certain predefined order
against the input stream.

BLANK should always be taken as whitespace, some things may depend on
this.  TAB is also fixed by convention, but may in principle be
modified.  Newline (DOS!) and linefeed are also defined as whitespace.
\nwenddocs{}\nwbegindocs{132}\nwdocspar
The lexer setup, containing the list of lexeme templates.  No
defaults yet.  The type with index zero will be assigned to the
{\tt{}NO{\char95}MATCH} lexeme.

The keyword list is not stored, just a pointer to it.  We anticipate
that the keyword list is part of the syntax table, and the lexer needs
not alter it.  Furthermore, the lexer is typically finished before the
syntax table is.
\nwenddocs{}\nwbegincode{133}\moddef{Lexer: parameters}\plusendmoddef
  integer, parameter :: CASE_KEEP = 0, CASE_UP = 1, CASE_DOWN = 2

\eatline
\nwendcode{}\nwbegincode{134}\moddef{Lexer: types}\plusendmoddef
  type :: lexer_setup_t
     private
     type(template_t), dimension(:), allocatable :: tt
     integer, dimension(:), allocatable :: type
     integer :: keyword_case = CASE_KEEP
     type(keyword_list_t), pointer :: keyword_list => null ()
  end type lexer_setup_t

\eatline
\nwendcode{}\nwbegindocs{135}Fill the lexer setup object.  Some things are hardcoded here
(whitespace, alphanumeric identifiers), some are free: comment chars
(but these must be single, and comments must be terminated by
line-feed), quote chars and matches (must be single), characters to be
read as one-character lexeme, special classes (characters of one class
that should be glued together as identifiers).
\nwenddocs{}\nwbegincode{136}\moddef{Lexer: procedures}\plusendmoddef
  subroutine lexer_setup_init (setup, &
       comment_chars, quote_chars, quote_match, &
       single_chars, special_class, &
       keyword_list, upper_case_keywords)
    type(lexer_setup_t), intent(inout) :: setup
    character(*), intent(in) :: comment_chars
    character(*), intent(in) :: quote_chars, quote_match
    character(*), intent(in) :: single_chars
    character(*), dimension(:), intent(in) :: special_class
    type(keyword_list_t), pointer :: keyword_list
    logical, intent(in), optional :: upper_case_keywords
    integer :: n, i
    if (present (upper_case_keywords)) then
       if (upper_case_keywords) then
          setup%keyword_case = CASE_UP
       else
          setup%keyword_case = CASE_DOWN
       end if
    else
       setup%keyword_case = CASE_KEEP
    end if
    n = 1 + len (comment_chars) + len (quote_chars) + 1 &
         + len (single_chars) + size (special_class) + 1
    allocate (setup%tt(n))
    allocate (setup%type(0:n))
    n = 0
    setup%type(n) = NO_MATCH
    n = n + 1
    setup%tt(n) = template_whitespace (WHITESPACE_CHARS)
    setup%type(n) = EMPTY
    forall (i = 1:len(comment_chars))
       setup%tt(n+i) = template_quoted (comment_chars(i:i), LF)
       setup%type(n+i) = EMPTY
    end forall
    n = n + len (comment_chars)
    forall (i = 1:len(quote_chars))
       setup%tt(n+i) = template_quoted (quote_chars(i:i), quote_match(i:i))
       setup%type(n+i) = T_QUOTED
    end forall
    n = n + len (quote_chars)
    setup%tt(n+1) = template_numeric ("EeDd")
    setup%type(n+1) = T_NUMERIC
    n = n + 1
    forall (i = 1:len (single_chars))
       setup%tt(n+i) = template_identifier (single_chars(i:i), "")
       setup%type(n+i) = T_IDENTIFIER
    end forall
    n = n + len (single_chars)
    forall (i = 1:size (special_class))
       setup%tt(n+i) = template_identifier &
            (trim (special_class(i)), trim (special_class(i)))
       setup%type(n+i) = T_IDENTIFIER
    end forall
    n = n + size (special_class)
    setup%tt(n+1) = template_identifier &
         (LCLETTERS//UCLETTERS, LCLETTERS//DIGITS//"_"//UCLETTERS)
    setup%type(n+1) = T_IDENTIFIER
    n = n + 1
    if (n /= size (setup%tt)) &
         call msg_bug ("Size mismatch in lexer setup")
    setup%keyword_list => keyword_list
  end subroutine lexer_setup_init
  
\eatline
\nwendcode{}\nwbegindocs{137}The destructor is needed only if the object is not itself part of an
allocatable array
\nwenddocs{}\nwbegincode{138}\moddef{Lexer: procedures}\plusendmoddef
  subroutine lexer_setup_final (setup)
    type(lexer_setup_t), intent(inout) :: setup
    deallocate (setup%tt, setup%type)
    setup%keyword_list => null ()
  end subroutine lexer_setup_final

\eatline
\nwendcode{}\nwbegindocs{139}For debugging: Write the lexer setup
\nwenddocs{}\nwbegincode{140}\moddef{Lexer: procedures}\plusendmoddef
  subroutine lexer_setup_write (setup, unit)
    type(lexer_setup_t), intent(in) :: setup
    integer, intent(in) :: unit
    integer :: i
    write (unit, "(A)") "Lexer setup:"
    if (allocated (setup%tt)) then
       do i = 1, size (setup%tt)
          call template_write (setup%tt(i), unit)
          write (unit, '(A)', advance = "no")  " -> "
          call lexeme_type_write (setup%type(i), unit)
          write (unit, *)
       end do
    else
       write (unit, *) "[empty]"
    end if
    if (associated (setup%keyword_list)) then
       call keyword_list_write (setup%keyword_list, unit)
    end if
   end subroutine lexer_setup_write

\eatline
\nwendcode{}\nwbegindocs{141}\nwdocspar
\subsection{The lexeme type}
An object of this type is returned by the lexer.  Apart from the lexeme
string, it gives information about the relevant substring (first and
last character index) and the lexeme type.  Interpreting the string is
up to the parser.
\nwenddocs{}\nwbegincode{142}\moddef{Lexer: public}\plusendmoddef
  public :: lexeme_t
\nwendcode{}\nwbegincode{143}\moddef{Lexer: types}\plusendmoddef
  type :: lexeme_t
     private
     integer :: type = EMPTY
     type(string_t) :: s
     integer :: b = 0, e = 0
  end type lexeme_t

\eatline
\nwendcode{}\nwbegindocs{144}Debugging aid:
\nwenddocs{}\nwbegincode{145}\moddef{Lexer: public}\plusendmoddef
  public :: lexeme_write
\nwendcode{}\nwbegincode{146}\moddef{Lexer: procedures}\plusendmoddef
  subroutine lexeme_write (t, unit)
    type(lexeme_t), intent(in) :: t
    integer, intent(in) :: unit
    integer :: u
    u = given_output_unit (unit);  if (u < 0)  return
    select case (t%type)
    case (T_KEYWORD)
       write (u, *) "KEYWORD:    '" // char (t%s) // "'"
    case (T_IDENTIFIER)
       write (u, *) "IDENTIFIER: '" // char (t%s) // "'"
    case (T_QUOTED)
       write (u, *) "QUOTED:     '" // char (t%s) // "'"
    case (T_NUMERIC)
       write (u, *) "NUMERIC:    '" // char (t%s) // "'"
    case (UNMATCHED_QUOTE)
       write (u, *) "Unmatched quote: "// char (t%s)
    case (OVERFLOW); write (u, *) "Overflow: "// char (t%s)
    case (EMPTY);    write (u, *) "Empty lexeme"
    case (NO_MATCH); write (u, *) "No match"
    case (IO_ERROR); write (u, *) "IO error"
    case (EOF);      write (u, *) "EOF"
    case default
       write (u, *) "Error"
    end select
  end subroutine lexeme_write

\eatline
\nwendcode{}\nwbegindocs{147}Store string and type in a lexeme.  The range determines the
beginning and end of the relevant part of the string.  Check for a
keyword.
\nwenddocs{}\nwbegincode{148}\moddef{Lexer: procedures}\plusendmoddef
  subroutine lexeme_set (t, keyword_list, s, range, type, keyword_case)
    type(lexeme_t), intent(out) :: t
    type(keyword_list_t), pointer :: keyword_list
    type(string_t), intent(in) :: s
    type(string_t) :: keyword
    integer, dimension(2), intent(in) :: range
    integer, intent(in) :: type
    integer, intent(in), optional :: keyword_case
    t%type = type
    if (present (keyword_case)) then
       select case (keyword_case)
       case (CASE_KEEP);   keyword = s
       case (CASE_UP);     keyword = upper_case (s)
       case (CASE_DOWN);   keyword = lower_case (s)
       end select
    else
       keyword = s
    end if
    if (type == T_IDENTIFIER) then
       if (associated (keyword_list)) then
          if (keyword_list_contains (keyword_list, keyword)) &
               t%type = T_KEYWORD
       end if
    end if
    select case (t%type)
    case (T_KEYWORD);  t%s = keyword
    case default;      t%s = s
    end select
    t%b = range(1)
    t%e = range(2)
  end subroutine lexeme_set

  subroutine lexeme_clear (t)
    type(lexeme_t), intent(out) :: t
    t%type = EMPTY
    t%s = ""
  end subroutine lexeme_clear

\eatline
\nwendcode{}\nwbegindocs{149}Retrieve the lexeme string, the relevant part of it, and the type.
The last function returns true if there is a break condition reached
(error or EOF).
\nwenddocs{}\nwbegincode{150}\moddef{Lexer: public}\plusendmoddef
  public :: lexeme_get_string
  public :: lexeme_get_contents
  public :: lexeme_get_delimiters
  public :: lexeme_get_type
\nwendcode{}\nwbegincode{151}\moddef{Lexer: procedures}\plusendmoddef
  function lexeme_get_string (t) result (s)
    type(string_t) :: s
    type(lexeme_t), intent(in) :: t
    s = t%s
  end function lexeme_get_string

  function lexeme_get_contents (t) result (s)
    type(string_t) :: s
    type(lexeme_t), intent(in) :: t
    s = extract (t%s, t%b, t%e)
  end function lexeme_get_contents

  function lexeme_get_delimiters (t) result (del)
    type(string_t), dimension(2) :: del
    type(lexeme_t), intent(in) :: t
    del(1) = extract (t%s, finish = t%b-1)
    del(2) = extract (t%s, start = t%e+1)
  end function lexeme_get_delimiters

  function lexeme_get_type (t) result (type)
    integer :: type
    type(lexeme_t), intent(in) :: t
    type = t%type
  end function lexeme_get_type

\eatline
\nwendcode{}\nwbegindocs{152}Check for a generic break condition (error/eof) and for eof in
particular.
\nwenddocs{}\nwbegincode{153}\moddef{Lexer: public}\plusendmoddef
  public :: lexeme_is_break
  public :: lexeme_is_eof
\nwendcode{}\nwbegincode{154}\moddef{Lexer: procedures}\plusendmoddef
  function lexeme_is_break (t) result (break)
    logical :: break
    type(lexeme_t), intent(in) :: t
    select case (t%type)
    case (EOF, IO_ERROR, OVERFLOW, NO_MATCH)
       break = .true.
    case default
       break = .false.
    end select
  end function lexeme_is_break

  function lexeme_is_eof (t) result (ok)
    logical :: ok
    type(lexeme_t), intent(in) :: t
    ok = t%type == EOF
  end function lexeme_is_eof

\eatline
\nwendcode{}\nwbegindocs{155}\nwdocspar
\subsection{The lexer object}
We store the current lexeme and the current line.  The line buffer is
set each time a new line is read from file.  The working buffer has
one character more, to hold any trailing blank.  Pointers to line and
column are for debugging, they will be used to make up readable error
messages for the parser.
\nwenddocs{}\nwbegincode{156}\moddef{Lexer: public}\plusendmoddef
  public :: lexer_t
\nwendcode{}\nwbegincode{157}\moddef{Lexer: types}\plusendmoddef
  type :: lexer_t
     private
     type(lexer_setup_t) :: setup
     type(stream_t), pointer :: stream => null ()
     type(lexeme_t) :: lexeme
     type(string_t) :: previous_line2
     type(string_t) :: previous_line1
     type(string_t) :: current_line
     integer :: lines_read = 0
     integer :: current_column = 0
     integer :: previous_column = 0
     type(string_t) :: buffer
     type(lexer_t), pointer :: parent => null ()
  end type lexer_t
\eatline
\nwendcode{}\nwbegindocs{158}Create-setup wrapper
\nwenddocs{}\nwbegincode{159}\moddef{Lexer: public}\plusendmoddef
  public :: lexer_init
\nwendcode{}\nwbegincode{160}\moddef{Lexer: procedures}\plusendmoddef
  subroutine lexer_init (lexer, &
       comment_chars, quote_chars, quote_match, &
       single_chars, special_class, &
       keyword_list, upper_case_keywords, &
       parent)
    type(lexer_t), intent(inout) :: lexer
    character(*), intent(in) :: comment_chars
    character(*), intent(in) :: quote_chars, quote_match
    character(*), intent(in) :: single_chars
    character(*), dimension(:), intent(in) :: special_class
    type(keyword_list_t), pointer :: keyword_list
    logical, intent(in), optional :: upper_case_keywords
    type(lexer_t), target, intent(in), optional :: parent
    call lexer_setup_init (lexer%setup, &
         comment_chars = comment_chars, &
         quote_chars = quote_chars, &
         quote_match = quote_match, &
         single_chars = single_chars, &
         special_class = special_class, &
         keyword_list = keyword_list, &
         upper_case_keywords = upper_case_keywords)
    if (present (parent))  lexer%parent => parent
    call lexer_clear (lexer)
  end subroutine lexer_init

\eatline
\nwendcode{}\nwbegindocs{161}Clear the lexer state, but not the setup.  This should be done when
the lexing starts, but it is not known whether the lexer was used
before.
\nwenddocs{}\nwbegincode{162}\moddef{Lexer: public}\plusendmoddef
  public :: lexer_clear
\nwendcode{}\nwbegincode{163}\moddef{Lexer: procedures}\plusendmoddef
  subroutine lexer_clear (lexer)
    type(lexer_t), intent(inout) :: lexer
    call lexeme_clear (lexer%lexeme)
    lexer%previous_line2 = ""
    lexer%previous_line1 = ""
    lexer%current_line = ""
    lexer%lines_read = 0
    lexer%current_column = 0
    lexer%previous_column = 0
    lexer%buffer = ""
  end subroutine lexer_clear

\eatline
\nwendcode{}\nwbegindocs{164}Reset lexer state and delete setup
\nwenddocs{}\nwbegincode{165}\moddef{Lexer: public}\plusendmoddef
  public :: lexer_final
\nwendcode{}\nwbegincode{166}\moddef{Lexer: procedures}\plusendmoddef
  subroutine lexer_final (lexer)
    type(lexer_t), intent(inout) :: lexer
    call lexer_clear (lexer)
    call lexer_setup_final (lexer%setup)
  end subroutine lexer_final

\eatline
\nwendcode{}\nwbegindocs{167}\nwdocspar
\subsection{The lexer routine}
For lexing we need to associate an input stream to the lexer.
\nwenddocs{}\nwbegincode{168}\moddef{Lexer: public}\plusendmoddef
  public :: lexer_assign_stream
\nwendcode{}\nwbegincode{169}\moddef{Lexer: procedures}\plusendmoddef
  subroutine lexer_assign_stream (lexer, stream)
    type(lexer_t), intent(inout) :: lexer
    type(stream_t), intent(in), target :: stream
    lexer%stream => stream
  end subroutine lexer_assign_stream

\eatline
\nwendcode{}\nwbegindocs{170}\nwdocspar
The lexer.  The lexer function takes the lexer and returns the
currently stored lexeme.  If there is none, it is read from buffer,
matching against the lexeme templates in the lexer setup.  Empty
lexemes, i.e., comments and whitespace, are discarded and the buffer
is read again until we have found a nonempty lexeme (which may also be
EOF or an error condition).

The initial state of the lexer contains an empty lexeme, so reading
from buffer is forced.  The empty state is restored after returning
the lexeme.  A nonempty lexeme is present in the lexer only if
{\tt{}lex{\char95}back} has been executed before.

The workspace is the {\tt{}lexer{\char37}buffer}, treated as a sort of input
stream.  We chop off lexemes from the beginning, adjusting the buffer
to the left.  Whenever the buffer is empty, or we are matching against
an open quote which has not terminated, we read a new line and append
it to the right.  This may result in special conditions, which for
simplicity are also returned as lexemes: I/O error, buffer overflow,
end of file.  If the latter happens during reading a quoted string,
we return an unmatched-quote lexeme.  Obviously, the special-condition
lexemes have to be caught by the parser.

Note that reading further lines is only necessary when reading a
quoted string.  Otherwise, the line-feed that ends each line is
interpreted as whitespace which terminates a preceding lexeme, so there
are no other valid multiline lexemes.

To enable meaningful error messages, we also keep track of the line
number of the last line read, and the beginning and the end of the
current lexeme with respect to this line.

The lexer is implemented as a function that returns the next lexeme
(i.e., token).  It uses the {\tt{}lexer} setup and modifies the buffers
and pointers stored within the lexer, a side effect.  The lexer reads
from an input stream object, which also is modified by this reading,
e.g., a line pointer is advanced.
\nwenddocs{}\nwbegincode{171}\moddef{Lexer: public}\plusendmoddef
  public :: lex
\nwendcode{}\nwbegincode{172}\moddef{Lexer: procedures}\plusendmoddef
  subroutine lex (lexeme, lexer)
    type(lexeme_t), intent(out) :: lexeme
    type(lexer_t), intent(inout) :: lexer
    integer :: iostat1, iostat2
    integer :: pos
    integer, dimension(2) :: range
    integer :: template_index, type
    if (.not. associated (lexer%stream)) &
        call msg_bug ("Lexer called without assigned stream")
    GET_LEXEME: do while (lexeme_get_type (lexer%lexeme) == EMPTY)
       if (len (lexer%buffer) /= 0) then
          iostat1 = 0
       else
          call lexer_read_line (lexer, iostat1)
       end if
       select case (iostat1)
       case (0)
          MATCH_BUFFER: do
             call match (lexer%setup%tt, char (lexer%buffer), &
                         pos, range, template_index)
             if (pos >= 0) then
                type = lexer%setup%type(template_index)
                exit MATCH_BUFFER
             else
                pos = 0
                call lexer_read_line (lexer, iostat2)
                select case (iostat2)
                case (EOF); type = UNMATCHED_QUOTE; exit MATCH_BUFFER
                case (1);   type = IO_ERROR;        exit MATCH_BUFFER
                case (2);   type = OVERFLOW;        exit MATCH_BUFFER
                end select
             end if
          end do MATCH_BUFFER
       case (EOF); type = EOF
       case (1);   type = IO_ERROR
       case (2);   type = OVERFLOW
       end select
       call lexeme_set (lexer%lexeme, lexer%setup%keyword_list, &
            extract (lexer%buffer, finish=pos), range, type, &
            lexer%setup%keyword_case)
       lexer%buffer = remove (lexer%buffer, finish=pos)
       lexer%previous_column = lexer%current_column
       lexer%current_column = lexer%current_column + pos
    end do GET_LEXEME
    lexeme = lexer%lexeme
    call lexeme_clear (lexer%lexeme)
  end subroutine lex

\eatline
\nwendcode{}\nwbegindocs{173}Read a line and append it to the input buffer.  If the input buffer
overflows, return {\tt{}iostat=2}.  Otherwise, {\tt{}iostat=1} indicates an
I/O error, and {\tt{}iostat=-1} the EOF.

The input stream may either be an external unit or a {\tt{}ifile}
object.  In the latter case, a line is read and the line pointer is
advanced.

Note that inserting {\tt{}LF} between input lines is the Unix convention.
Since we are doing this explicitly when gluing lines together, we can
pattern-match against {\tt{}LF} without having to worry about the system.
\nwenddocs{}\nwbegincode{174}\moddef{Lexer: procedures}\plusendmoddef
  subroutine lexer_read_line (lexer, iostat)
    type(lexer_t), intent(inout) :: lexer
    integer, intent(out) :: iostat
    type(string_t) :: current_line
    current_line = lexer%current_line
    call stream_get_record (lexer%stream, lexer%current_line, iostat)
    if (iostat == 0) then
       lexer%lines_read = lexer%lines_read + 1
       lexer%previous_line2 = lexer%previous_line1
       lexer%previous_line1 = current_line
       lexer%buffer = lexer%buffer // lexer%current_line // LF
       lexer%previous_column = 0
       lexer%current_column = 0
    end if
  end subroutine lexer_read_line

\eatline
\nwendcode{}\nwbegindocs{175}Once in a while we have read one lexeme to many, which can be
pushed back into the input stream.  Do not do this more than once.
\nwenddocs{}\nwbegincode{176}\moddef{Lexer: public}\plusendmoddef
  public :: lexer_put_back
\nwendcode{}\nwbegincode{177}\moddef{Lexer: procedures}\plusendmoddef
  subroutine lexer_put_back (lexer, lexeme)
    type(lexer_t), intent(inout) :: lexer
    type(lexeme_t), intent(in) :: lexeme
    if (lexeme_get_type (lexer%lexeme) == EMPTY) then
       lexer%lexeme = lexeme
    else
       call msg_bug (" Lexer: lex_back fails; probably called twice")
    end if
  end subroutine lexer_put_back
  
\eatline
\nwendcode{}\nwbegindocs{178}\nwdocspar
\subsection{Diagnostics}
For debugging: print just the setup
\nwenddocs{}\nwbegincode{179}\moddef{Lexer: public}\plusendmoddef
  public :: lexer_write_setup
\nwendcode{}\nwbegincode{180}\moddef{Lexer: procedures}\plusendmoddef
  subroutine lexer_write_setup (lexer, unit)
    type(lexer_t), intent(in) :: lexer
    integer, intent(in), optional :: unit
    integer :: u
    u = given_output_unit (unit);  if (u < 0)  return
    call lexer_setup_write (lexer%setup, u)
  end subroutine lexer_write_setup

\eatline
\nwendcode{}\nwbegindocs{181}\nwdocspar
This is useful for error printing: show the current line with index
and a pointer to the current column within the line.
\nwenddocs{}\nwbegincode{182}\moddef{Lexer: public}\plusendmoddef
  public :: lexer_show_location
\nwendcode{}\nwbegincode{183}\moddef{Lexer: procedures}\plusendmoddef
  subroutine lexer_show_location (lexer)
    type(lexer_t), intent(in) :: lexer
    type(string_t) :: loc_str
    if (associated (lexer%parent)) then
       call lexer_show_source (lexer%parent)
       call msg_message ("[includes]")
    else
       call msg_message ()
    end if
    if (associated (lexer%stream)) then
       call msg_message &
            (char (stream_get_record_info_string (lexer%stream)) // ":")
    end if
    if (lexer%lines_read >= 4)  call msg_result ("[...]")
    if (lexer%lines_read >= 3)  call msg_result (char (lexer%previous_line2))
    if (lexer%lines_read >= 2)  call msg_result (char (lexer%previous_line1))
    if (lexer%lines_read >= 1) then
       call msg_result (char (lexer%current_line))
       loc_str = repeat (" ", lexer%previous_column)
       loc_str = loc_str // "^"
       if (lexer%current_column > lexer%previous_column) then
          loc_str = loc_str &
               // repeat ("-", max (lexer%current_column &
                                    - lexer%previous_column - 1, 0)) &
               // "^"
       end if
       call msg_result (char (loc_str))
    end if
  end subroutine lexer_show_location

\eatline
\nwendcode{}\nwbegindocs{184}This just prints the current stream source.
\nwenddocs{}\nwbegincode{185}\moddef{Lexer: procedures}\plusendmoddef
  recursive subroutine lexer_show_source (lexer)
    type(lexer_t), intent(in) :: lexer
    if (associated (lexer%parent)) then
       call lexer_show_source (lexer%parent)
       call msg_message ("[includes]")
    else
       call msg_message ()
    end if
    if (associated (lexer%stream)) then
       call msg_message &
            (char (stream_get_source_info_string (lexer%stream)) // ":")
    end if
  end subroutine lexer_show_source

\eatline
\nwendcode{}\nwbegindocs{186}\nwdocspar
\subsection{Unit tests}
\nwenddocs{}\nwbegincode{187}\moddef{Lexer: public}\plusendmoddef
  public :: lexer_test
\nwendcode{}\nwbegincode{188}\moddef{Lexer: procedures}\plusendmoddef
  subroutine lexer_test (u, results)
    integer, intent(in) :: u
    type(test_results_t), intent(inout) :: results
  \LA{}Lexer: execute tests\RA{}
  end subroutine lexer_test

\eatline
\nwendcode{}\nwbegindocs{189}Test the lexer by lexing and printing all lexemes from unit {\tt{}u},
one per line, using preset conventions. 
\nwenddocs{}\nwbegincode{190}\moddef{Lexer: execute tests}\endmoddef
  call test (lexer_1, "lexer_1", &
       "check lexer", u, results)
\nwendcode{}\nwbegincode{191}\moddef{Lexer: tests}\endmoddef
  subroutine lexer_1 (u)
    integer, intent(in) :: u
    type(lexer_t), target :: lexer
    type(stream_t), target :: stream
    type(string_t) :: string
    type(lexeme_t) :: lexeme
    string = "abcdefghij"
    call lexer_init (lexer, &
       comment_chars = "", &
       quote_chars = "<'""", &
       quote_match = ">'""", &
       single_chars = "?*+|=,()", &
       special_class = ["."], &
       keyword_list = null ())
    call stream_init (stream, string)
    call lexer_assign_stream (lexer, stream)
    do
       call lex (lexeme, lexer)
       call lexeme_write (lexeme, u)
       if (lexeme_is_break (lexeme))  exit
    end do
    call stream_final (stream)
    call lexer_final (lexer)
  end subroutine lexer_1

\eatline
\nwendcode{}\nwbegindocs{192}\nwdocspar
\clearpage
%%%%%%%%%%%%%%%%%%%%%%%%%%%%%%%%%%%%%%%%%%%%%%%%%%%%%%%%%%%%%%%%%%%%%%%%
\section{Syntax rules}
This module provides tools to handle syntax rules in an abstract way.
\nwenddocs{}\nwbegincode{193}\moddef{\code{}syntax{\char95}rules.f90\edoc{}}\endmoddef
\LA{}File header\RA{}

module syntax_rules

\LA{}Use strings\RA{}
  use io_units
  use diagnostics
  use system_defs, only: LCLETTERS, UCLETTERS, DIGITS
  use ifiles, only: line_p, line_init, line_get_string_advance, line_final
  use ifiles, only: ifile_t, ifile_get_length
  use lexers

\LA{}Standard module head\RA{}

\LA{}Syntax: public\RA{}

\LA{}Syntax: parameters\RA{}

\LA{}Syntax: types\RA{}

\LA{}Syntax: interfaces\RA{}

contains

\LA{}Syntax: subroutines\RA{}

end module syntax_rules
\eatline
\nwendcode{}\nwbegindocs{194}\nwdocspar
\subsection{Syntax rules}
Syntax rules are used by the parser.  They determine how to translate
the stream of lexemes as returned by the lexer into the parse tree
node.  A rule may be terminal, i.e., replace a matching lexeme into a
terminal node.  The node will contain the lexeme interpreted as a
recognized token:
\begin{itemize}
\item a keyword: unquoted fixed character string;
\item a real number, to be determined at runtime;
\item an integer, to be determined at runtime;
\item a boolean value, to be determined at runtime;
\item a quoted token (e.g., string), to be determined at runtime;
\item an identifier (unquoted string that is not a recognized
keyword), to be determined at runtime.
\end{itemize}
It may be nonterminal, i.e., contain a sequence of child rules.
These are matched consecutively (and recursively) against the input
stream; the resulting node will be a branch node.
\begin{itemize}
\item the file, i.e., the input stream as a whole;
\item a sequence of syntax elements, where the last syntax element may be
  optional, or optional repetitive;
\end{itemize}
Sequences carry a flag that tells whether the last child is optional
or may be repeated an arbitrary number of times, correponding to the
regexp modifiers {\tt{}?}, {\tt{}*}, and {\tt{}+}.

We also need an alternative rule; this will be replaced by the node
generated by one of its children that matches; thus, it does not
create a node of its own.
\begin{itemize}
\item an alternative of syntax elements.
\end{itemize}
We also define special types of sequences as convenience macros:
\begin{itemize}
\item a list: a sequence where the elements are separated by a
  separator keyword (e.g., commas), the separators are thrown away
  when parsing the list;
\item a group: a sequence of three tokens, where the first and third
  ones are left and right delimiters, the delimiters are thrown away;
\item an argument list: a delimited list, containing both delimiters
  and separators.
\end{itemize}
It would be great to have a polymorphic type for this purpose, but
until Fortran 2003 is out we have to emulate this.

Here are the syntax element codes:
\nwenddocs{}\nwbegincode{195}\moddef{Syntax: public}\endmoddef
  public :: S_UNKNOWN
  public :: S_LOGICAL, S_INTEGER, S_REAL, S_COMPLEX, S_QUOTED
  public :: S_IDENTIFIER, S_KEYWORD
  public :: S_SEQUENCE, S_LIST, S_GROUP, S_ARGS
  public :: S_ALTERNATIVE
  public :: S_IGNORE
\nwendcode{}\nwbegincode{196}\moddef{Syntax: parameters}\endmoddef
  integer, parameter :: &
       S_UNKNOWN = 0, &
       S_LOGICAL = 1, S_INTEGER = 2, S_REAL = 3, S_COMPLEX = 4, &
       S_QUOTED = 5, S_IDENTIFIER = 6, S_KEYWORD = 7, &
       S_SEQUENCE = 8, S_LIST = 9, S_GROUP = 10, S_ARGS = 11, &
       S_ALTERNATIVE = 12, &
       S_IGNORE = 99
\nwendcode{}\nwbegindocs{197}We need arrays of rule pointers, therefore this construct.
\nwenddocs{}\nwbegincode{198}\moddef{Syntax: types}\endmoddef
  type :: rule_p
     private
     type(syntax_rule_t), pointer :: p => null ()
  end type rule_p

\eatline
\nwendcode{}\nwbegindocs{199}Return the association status of the rule pointer:
\nwenddocs{}\nwbegincode{200}\moddef{Syntax: subroutines}\endmoddef
  elemental function rule_is_associated (rp) result (ok)
    logical :: ok
    type (rule_p), intent(in) :: rp
    ok = associated (rp%p)
  end function rule_is_associated

\eatline
\nwendcode{}\nwbegindocs{201}The rule type is one of the types listed above, represented by an
integer code.  The keyword, for a non-keyword rule, is an identifier
used for the printed syntax table.  The array of children is needed
for nonterminal rules.  In that case, there is a modifier for the last
element (blank, "?", "*", or "+"), mirrored in the flags {\tt{}opt} and
{\tt{}rep}.  Then, we have the character constants used as separators and
delimiters for this rule.  Finally, the {\tt{}used} flag can be set to
indicate that this rule is the child of another rule.
\nwenddocs{}\nwbegincode{202}\moddef{Syntax: types}\plusendmoddef
  public :: syntax_rule_t
\nwendcode{}\nwbegincode{203}\moddef{Syntax: types}\plusendmoddef
  type :: syntax_rule_t
     private
     integer :: type = S_UNKNOWN
     logical :: used = .false.
     type(string_t) :: keyword
     type(string_t) :: separator
     type(string_t), dimension(2) :: delimiter
     type(rule_p), dimension(:), allocatable :: child
     character(1) :: modifier = ""
     logical :: opt = .false., rep = .false.
  end type syntax_rule_t

\eatline
\nwendcode{}\nwbegindocs{204}Initializer: Set type and key for a rule, but do not (yet) allocate
anything.  Finalizer: not needed (no pointer components).
\nwenddocs{}\nwbegincode{205}\moddef{Syntax: subroutines}\plusendmoddef
  subroutine syntax_rule_init (rule, key, type)
    type(syntax_rule_t), intent(inout) :: rule
    type(string_t), intent(in) :: key
    integer, intent(in) :: type
    rule%keyword = key
    rule%type = type
    select case (rule%type)
    case (S_GROUP)
       call syntax_rule_set_delimiter (rule)
    case (S_LIST)
       call syntax_rule_set_separator (rule)
    case (S_ARGS)
       call syntax_rule_set_delimiter (rule)
       call syntax_rule_set_separator (rule)
    end select
  end subroutine syntax_rule_init
       
\eatline
\nwendcode{}\nwbegindocs{206}These characters will not be enclosed in quotes when writing syntax
rules:
\nwenddocs{}\nwbegincode{207}\moddef{Syntax: parameters}\plusendmoddef
  character(*), parameter :: &
           UNQUOTED = "(),|_"//LCLETTERS//UCLETTERS//DIGITS
\nwendcode{}\nwbegindocs{208}\nwdocspar
\subsection{I/O}
Write an account of the rule.  Setting {\tt{}short} true will suppress
the node type.  Setting {\tt{}key{\char95}only} true will suppress the
definition.  Setting {\tt{}advance} false will suppress the trailing
newline.
\nwenddocs{}\nwbegincode{209}\moddef{Syntax: public}\plusendmoddef
  public :: syntax_rule_write
\eatline
\nwendcode{}\nwbegindocs{210}\nwdocspar
\nwenddocs{}\nwbegincode{211}\moddef{Syntax: subroutines}\plusendmoddef
  subroutine syntax_rule_write (rule, unit, short, key_only, advance)
    type(syntax_rule_t), intent(in) :: rule
    integer, intent(in), optional :: unit
    logical, intent(in), optional :: short, key_only, advance
    logical :: typ, def, adv
    integer :: u
    u = given_output_unit (unit);  if (u < 0)  return
    typ = .true.;  if (present (short))    typ = .not. short
    def = .true.;  if (present (key_only)) def = .not. key_only
    adv = .true.;  if (present (advance))  adv = advance
    select case (rule%type)
    case (S_UNKNOWN);    call write_atom ("???", typ)
    case (S_IGNORE);     call write_atom ("IGNORE", typ)
    case (S_LOGICAL);    call write_atom ("LOGICAL", typ)
    case (S_INTEGER);    call write_atom ("INTEGER", typ)
    case (S_REAL);       call write_atom ("REAL", typ)
    case (S_COMPLEX);    call write_atom ("COMPLEX", typ)    
    case (S_IDENTIFIER); call write_atom ("IDENTIFIER", typ)
    case (S_KEYWORD);    call write_atom ("KEYWORD", typ)
    case (S_QUOTED)
       call write_quotes (typ, def, del = rule%delimiter)
    case (S_SEQUENCE)
       call write_sequence ("SEQUENCE", typ, def, size (rule%child))
    case (S_GROUP)
       call write_sequence ("GROUP", typ, def, size (rule%child), &
            del = rule%delimiter)
    case (S_LIST)
       call write_sequence ("LIST", typ, def, size (rule%child), &
            sep = rule%separator)
    case (S_ARGS)
       call write_sequence ("ARGUMENTS", typ, def, size (rule%child), &
            del = rule%delimiter, sep = rule%separator)
    case (S_ALTERNATIVE)
       call write_sequence ("ALTERNATIVE", typ, def, size (rule%child), &
            sep = var_str ("|"))
    end select
    if (adv)  write (u, *)
  contains
    subroutine write_type (type)
      character(*), intent(in) :: type
      character(11) :: str
      str = type
      write (u, "(1x,A)", advance="no")  str
    end subroutine write_type
    subroutine write_key
      write (u, "(1x,A)", advance="no")  char (wkey (rule))
    end subroutine write_key
    subroutine write_atom (type, typ)
      character(*), intent(in) :: type
      logical, intent(in) :: typ
      if (typ)  call write_type (type)
      call write_key
    end subroutine write_atom
    subroutine write_maybe_quoted (string)
      character(*), intent(in) :: string
      character, parameter :: q = "'"
      character, parameter :: qq = '"'
      if (verify (string, UNQUOTED) == 0) then
         write (u, "(1x,A)", advance = "no") trim (string)
      else if (verify (string, q) == 0) then
         write (u, "(1x,A)", advance = "no") qq // trim (string) // qq
      else
         write (u, "(1x,A)", advance = "no") q // trim (string) // q
      end if
    end subroutine write_maybe_quoted
    subroutine write_quotes (typ, def, del)
      logical, intent(in) :: typ, def
      type(string_t), dimension(2), intent(in) :: del
      if (typ)  call write_type ("QUOTED")
      call write_key
      if (def) then
         write (u, "(1x,'=')", advance="no")
         call write_maybe_quoted (char (del(1)))
         write (u, "(1x,A)", advance="no") "..."
         call write_maybe_quoted (char (del(2)))
      end if
    end subroutine write_quotes
    subroutine write_sequence (type, typ, def, n, del, sep)
      character(*), intent(in) :: type
      logical, intent(in) :: typ, def
      integer, intent(in) :: n
      type(string_t), dimension(2), intent(in), optional :: del
      type(string_t), intent(in), optional :: sep
      integer :: i
      if (typ)  call write_type (type)
      call write_key
      if (def) then
         write (u, "(1x,'=')", advance="no")
         if (present (del))  call write_maybe_quoted (char (del(1)))
         do i = 1, n
            if (i > 1 .and. present (sep)) &
                 call write_maybe_quoted (char (sep))
            write (u, "(1x,A)", advance="no")  &
                 char (wkey (syntax_rule_get_sub_ptr(rule, i)))
            if (i == n)  write (u, "(A)", advance="no")  trim (rule%modifier)
         end do
         if (present (del))  call write_maybe_quoted (char (del(2)))
      end if
    end subroutine write_sequence
  end subroutine syntax_rule_write

\eatline
\nwendcode{}\nwbegindocs{212}In the printed representation, the keyword strings are enclosed as
{\tt{}<...>}, unless they are bare keywords.  Bare keywords are enclosed
as {\tt{}'..'} if they contain a character which is not a letter, digit,
or underscore.  If they contain a single-quote character, they are
enclosed as {\tt{}".."}.  (A keyword must not contain both single- and
double-quotes.)
\nwenddocs{}\nwbegincode{213}\moddef{Syntax: subroutines}\plusendmoddef
  function wkey (rule) result (string)
    type(string_t) :: string
    type(syntax_rule_t), intent(in) :: rule
    select case (rule%type)
    case (S_KEYWORD)
       if (verify (rule%keyword, UNQUOTED) == 0) then
          string = rule%keyword
       else if (scan (rule%keyword, "'") == 0) then
          string = "'" // rule%keyword // "'"
       else
          string = '"' // rule%keyword // '"'
       end if
    case default
       string = "<" // rule%keyword // ">"
    end select
  end function wkey

\eatline
\nwendcode{}\nwbegindocs{214}\nwdocspar
\subsection{Completing syntax rules}
Set the separator and delimiter entries, using defaults:
\nwenddocs{}\nwbegincode{215}\moddef{Syntax: subroutines}\plusendmoddef
  subroutine syntax_rule_set_separator (rule, separator)
    type(syntax_rule_t), intent(inout) :: rule
    type(string_t), intent(in), optional :: separator
    if (present (separator)) then
       rule%separator = separator
    else
       rule%separator = ","
    end if
  end subroutine syntax_rule_set_separator

  subroutine syntax_rule_set_delimiter (rule, delimiter)
    type(syntax_rule_t), intent(inout) :: rule
    type(string_t), dimension(2), intent(in), optional :: delimiter
    if (present (delimiter)) then
       rule%delimiter = delimiter
    else
       rule%delimiter = ["(", ")"]
    end if
  end subroutine syntax_rule_set_delimiter

\eatline
\nwendcode{}\nwbegindocs{216}Set the modifier entry and corresponding flags:
\nwenddocs{}\nwbegincode{217}\moddef{Syntax: subroutines}\plusendmoddef
  function is_modifier (string) result (ok)
    logical :: ok
    type(string_t), intent(in) :: string
    select case (char (string))
    case (" ", "?", "*", "+");  ok = .true.
    case default;               ok = .false.
    end select
  end function is_modifier

  subroutine syntax_rule_set_modifier (rule, modifier)
    type(syntax_rule_t), intent(inout) :: rule
    type(string_t), intent(in) :: modifier
    rule%modifier = char (modifier)
    select case (rule%modifier)
    case (" ")
    case ("?");  rule%opt = .true.
    case ("*");  rule%opt = .true.;  rule%rep = .true.
    case ("+");  rule%rep = .true.
    case default
       call msg_bug (" Syntax: sequence modifier '" // rule%modifier &
            // "' is not one of '+' '*' '?'")
    end select
  end subroutine syntax_rule_set_modifier

\eatline
\nwendcode{}\nwbegindocs{218}Check a finalized rule for completeness
\nwenddocs{}\nwbegincode{219}\moddef{Syntax: subroutines}\plusendmoddef
  subroutine syntax_rule_check (rule)
    type(syntax_rule_t), intent(in) :: rule
    if (rule%keyword == "")  call msg_bug ("Rule key not set")
    select case (rule%type)
    case (S_UNKNOWN);  call bug (" Undefined rule")
    case (S_IGNORE, S_LOGICAL, S_INTEGER, S_REAL, S_COMPLEX, &
             S_IDENTIFIER, S_KEYWORD)
    case (S_QUOTED)
       if (any (rule%delimiter == "")) call bug (" Missing quote character(s)")
    case (S_SEQUENCE)
    case (S_GROUP)
       if (any (rule%delimiter == "")) call bug (" Missing delimiter(s)")
    case (S_LIST)
       if (rule%separator == "") call bug (" Missing separator")
    case (S_ARGS)
       if (any (rule%delimiter == "")) call bug (" Missing delimiter(s)")
       if (rule%separator == "") call bug (" Missing separator")
    case (S_ALTERNATIVE)
    case default
       call bug (" Undefined syntax code")
    end select
    select case (rule%type)
    case (S_SEQUENCE, S_GROUP, S_LIST, S_ARGS, S_ALTERNATIVE)
       if (allocated (rule%child)) then
          if (.not.all (rule_is_associated (rule%child))) &
               call bug (" Child rules not all associated")
       else
          call bug (" Parent rule without children")
       end if
    case default
       if (allocated (rule%child))  call bug (" Non-parent rule with children")
    end select
  contains
    subroutine bug (string)
      character(*), intent(in) :: string
      call msg_bug (" Syntax table: Rule " // char (rule%keyword) // ": " &
           // string)
    end subroutine bug
  end subroutine syntax_rule_check

\eatline
\nwendcode{}\nwbegindocs{220}\nwdocspar
\subsection{Accessing rules}
This is the API for syntax rules:
\nwenddocs{}\nwbegincode{221}\moddef{Syntax: public}\plusendmoddef
  public :: syntax_rule_get_type
\nwendcode{}\nwbegincode{222}\moddef{Syntax: subroutines}\plusendmoddef
  function syntax_rule_get_type (rule) result (type)
    integer :: type
    type(syntax_rule_t), intent(in) :: rule
    type = rule%type
  end function syntax_rule_get_type

\eatline
\nwendcode{}\nwbegincode{223}\moddef{Syntax: public}\plusendmoddef
  public :: syntax_rule_get_key
\nwendcode{}\nwbegincode{224}\moddef{Syntax: subroutines}\plusendmoddef
  function syntax_rule_get_key (rule) result (key)
    type(string_t) :: key
    type(syntax_rule_t), intent(in) :: rule
    key = rule%keyword
  end function syntax_rule_get_key

\eatline
\nwendcode{}\nwbegincode{225}\moddef{Syntax: public}\plusendmoddef
  public :: syntax_rule_get_separator
  public :: syntax_rule_get_delimiter
\nwendcode{}\nwbegincode{226}\moddef{Syntax: subroutines}\plusendmoddef
  function syntax_rule_get_separator (rule) result (separator)
    type(string_t) :: separator
    type(syntax_rule_t), intent(in) :: rule
    separator = rule%separator
  end function syntax_rule_get_separator

  function syntax_rule_get_delimiter (rule) result (delimiter)
    type(string_t), dimension(2) :: delimiter
    type(syntax_rule_t), intent(in) :: rule
    delimiter = rule%delimiter
  end function syntax_rule_get_delimiter

\eatline
\nwendcode{}\nwbegindocs{227}Accessing child rules.  If we use {\tt{}syntax{\char95}rule{\char95}get{\char95}n{\char95}sub} for determining
loop bounds, we do not need a check in the second routine.
\nwenddocs{}\nwbegincode{228}\moddef{Syntax: public}\plusendmoddef
  public :: syntax_rule_get_n_sub
  public :: syntax_rule_get_sub_ptr
\nwendcode{}\nwbegincode{229}\moddef{Syntax: subroutines}\plusendmoddef
  function syntax_rule_get_n_sub (rule) result (n)
    integer :: n
    type(syntax_rule_t), intent(in) :: rule
    if (allocated (rule%child)) then
       n = size (rule%child)
    else
       n = 0
    end if
  end function syntax_rule_get_n_sub

  function syntax_rule_get_sub_ptr (rule, i) result (sub)
    type(syntax_rule_t), pointer :: sub
    type(syntax_rule_t), intent(in), target :: rule
    integer, intent(in) :: i
    sub => rule%child(i)%p
  end function syntax_rule_get_sub_ptr

  subroutine syntax_rule_set_sub (rule, i, sub)
    type(syntax_rule_t), intent(inout) :: rule
    integer, intent(in) :: i
    type(syntax_rule_t), intent(in), target :: sub
    rule%child(i)%p => sub
  end subroutine syntax_rule_set_sub

\eatline
\nwendcode{}\nwbegindocs{230}Return the modifier flags:
\nwenddocs{}\nwbegincode{231}\moddef{Syntax: public}\plusendmoddef
  public :: syntax_rule_last_optional
  public :: syntax_rule_last_repetitive
\nwendcode{}\nwbegincode{232}\moddef{Syntax: subroutines}\plusendmoddef
  function syntax_rule_last_optional (rule) result (opt)
    logical :: opt
    type(syntax_rule_t), intent(in) :: rule
    opt = rule%opt
  end function syntax_rule_last_optional
  function syntax_rule_last_repetitive (rule) result (rep)
    logical :: rep
    type(syntax_rule_t), intent(in) :: rule
    rep = rule%rep
  end function syntax_rule_last_repetitive

\eatline
\nwendcode{}\nwbegindocs{233}\nwdocspar
Return true if the rule is atomic, i.e., logical, real, keyword etc.
\nwenddocs{}\nwbegincode{234}\moddef{Syntax: public}\plusendmoddef
  public :: syntax_rule_is_atomic
\nwendcode{}\nwbegincode{235}\moddef{Syntax: subroutines}\plusendmoddef
  function syntax_rule_is_atomic (rule) result (atomic)
    logical :: atomic
    type(syntax_rule_t), intent(in) :: rule
    select case (rule%type)
    case (S_LOGICAL, S_INTEGER, S_REAL, S_COMPLEX, S_IDENTIFIER, &
             S_KEYWORD, S_QUOTED)
       atomic = .true.
    case default
       atomic = .false.
    end select
  end function syntax_rule_is_atomic

\eatline
\nwendcode{}\nwbegindocs{236}\nwdocspar
\subsection{Syntax tables}
A syntax table contains the tree of syntax rules and, for direct
parser access, the list of valid keywords.

\subsubsection{Types}
The syntax contains an array of rules and a list of keywords.  The
array is actually used as a tree, where the top rule is the first
array element, and the other rules are recursively pointed to by this
first rule.  (No rule should be used twice or be unused.)  The keyword
list is derived from the rule tree.  

Objects of this type need the target attribute if they are associated
with a lexer.  The keyword list will be pointed to by this lexer.
\nwenddocs{}\nwbegincode{237}\moddef{Syntax: public}\plusendmoddef
  public :: syntax_t
\nwendcode{}\nwbegincode{238}\moddef{Syntax: types}\plusendmoddef
  type :: syntax_t
     private
     type(syntax_rule_t), dimension(:), allocatable :: rule
     type(keyword_list_t) :: keyword_list
  end type syntax_t

\eatline
\nwendcode{}\nwbegindocs{239}\nwdocspar
\subsubsection{Constructor/destructor}
Initialize and finalize syntax tables
\nwenddocs{}\nwbegincode{240}\moddef{Syntax: public}\plusendmoddef
  public :: syntax_init
  public :: syntax_final
\nwendcode{}\nwbegindocs{241}There are two ways to create a syntax: hard-coded from rules or
dynamically from file.
\nwenddocs{}\nwbegincode{242}\moddef{Syntax: interfaces}\endmoddef
  interface syntax_init
     module procedure syntax_init_from_ifile
  end interface

\eatline
\nwendcode{}\nwbegindocs{243}The syntax definition is read from an {\tt{}ifile} object which
contains the syntax definitions in textual form, one rule per line.
This interface allows for determining the number of rules beforehand.

To parse the rule definitions, we make up a temporary lexer.
Obviously, we cannot use a generic parser yet, so we have to hardcode
the parsing process.
\nwenddocs{}\nwbegincode{244}\moddef{Syntax: subroutines}\plusendmoddef
  subroutine syntax_init_from_ifile (syntax, ifile)
    type(syntax_t), intent(out), target :: syntax
    type(ifile_t), intent(in) :: ifile
    type(lexer_t) :: lexer
    type(line_p) :: line
    type(string_t) :: string
    integer :: n_token
    integer :: i
    call lexer_init (lexer, &
       comment_chars = "", &
       quote_chars = "<'""", &
       quote_match = ">'""", &
       single_chars = "?*+|=,()", &
       special_class = ["."], &
       keyword_list = null ())
    allocate (syntax%rule (ifile_get_length (ifile)))
    call line_init (line, ifile)
    do i = 1, size (syntax%rule)
       string = line_get_string_advance (line)
       call set_rule_type_and_key (syntax%rule(i), string, lexer)
    end do
    call line_init (line, ifile)
    do i = 1, size (syntax%rule)
       string = line_get_string_advance (line)
       select case (syntax%rule(i)%type)
       case (S_QUOTED, S_SEQUENCE, S_GROUP, S_LIST, S_ARGS, S_ALTERNATIVE)
          n_token = get_n_token (string, lexer)
          call set_rule_contents &
               (syntax%rule(i), syntax, n_token, string, lexer)
       end select
    end do
    call line_final (line)
    call lexer_final (lexer)
    call syntax_make_keyword_list (syntax)
    if (.not. all (syntax%rule%used)) then
       do i = 1, size (syntax%rule)
          if (.not. syntax%rule(i)%used) then
             call syntax_rule_write (syntax%rule(i), 6)
          end if
       end do
       call msg_bug (" Syntax table: unused rules")
    end if
  end subroutine syntax_init_from_ifile

\eatline
\nwendcode{}\nwbegindocs{245}For a given rule defined in the input, the first task is to
determine its type and key.  With these, we can initialize the rule in
the table, postponing the association of children.
\nwenddocs{}\nwbegincode{246}\moddef{Syntax: subroutines}\plusendmoddef
  subroutine set_rule_type_and_key (rule, string, lexer)
    type(syntax_rule_t), intent(inout) :: rule
    type(string_t), intent(in) :: string
    type(lexer_t), intent(inout) :: lexer
    type(stream_t), target :: stream
    type(lexeme_t) :: lexeme
    type(string_t) :: key
    character(2) :: type
    call lexer_clear (lexer)
    call stream_init (stream, string)
    call lexer_assign_stream (lexer, stream)
    call lex (lexeme, lexer)
    type = lexeme_get_string (lexeme)
    call lex (lexeme, lexer)
    key = lexeme_get_contents (lexeme)
    call stream_final (stream)
    if (trim (key) /= "") then
       select case (type)
       case ("IG");  call syntax_rule_init (rule, key, S_IGNORE)
       case ("LO");  call syntax_rule_init (rule, key, S_LOGICAL)
       case ("IN");  call syntax_rule_init (rule, key, S_INTEGER)
       case ("RE");  call syntax_rule_init (rule, key, S_REAL)
       case ("CO");  call syntax_rule_init (rule, key, S_COMPLEX)
       case ("ID");  call syntax_rule_init (rule, key, S_IDENTIFIER)
       case ("KE");  call syntax_rule_init (rule, key, S_KEYWORD)
       case ("QU");  call syntax_rule_init (rule, key, S_QUOTED)
       case ("SE");  call syntax_rule_init (rule, key, S_SEQUENCE)
       case ("GR");  call syntax_rule_init (rule, key, S_GROUP)
       case ("LI");  call syntax_rule_init (rule, key, S_LIST)
       case ("AR");  call syntax_rule_init (rule, key, S_ARGS)
       case ("AL");  call syntax_rule_init (rule, key, S_ALTERNATIVE)
       case default
          call lexer_show_location (lexer)
          call msg_bug (" Syntax definition: unknown type '" // type // "'")
       end select
    else
       print *, char (string)
       call msg_bug (" Syntax definition: empty rule key")
    end if
  end subroutine set_rule_type_and_key

\eatline
\nwendcode{}\nwbegindocs{247}This function returns the number of tokens in an input line.
\nwenddocs{}\nwbegincode{248}\moddef{Syntax: subroutines}\plusendmoddef
  function get_n_token (string, lexer) result (n)
    integer :: n
    type(string_t), intent(in) :: string
    type(lexer_t), intent(inout) :: lexer
    type(stream_t), target :: stream
    type(lexeme_t) :: lexeme
    integer :: i
    call lexer_clear (lexer)
    call stream_init (stream, string)
    call lexer_assign_stream (lexer, stream)
    i = 0
    do
       call lex (lexeme, lexer)
       if (lexeme_is_break (lexeme))  exit
       i = i + 1
    end do
    n = i
    call stream_final (stream)
end function get_n_token

\eatline
\nwendcode{}\nwbegindocs{249}This subroutine extracts the rule contents for an input line.  There
are three tasks: (1) determine the number of children, depending on
the rule type; (2) find and set the separator and delimiter strings,
if required; (3) scan the child rules, find them in the syntax
table and associate the parent rule with them.
\nwenddocs{}\nwbegincode{250}\moddef{Syntax: subroutines}\plusendmoddef
  subroutine set_rule_contents (rule, syntax, n_token, string, lexer)
    type(syntax_rule_t), intent(inout) :: rule
    type(syntax_t), intent(in), target :: syntax
    integer, intent(in) :: n_token
    type(string_t), intent(in) :: string
    type(lexer_t), intent(inout) :: lexer
    type(stream_t), target :: stream
    type(lexeme_t), dimension(n_token) :: lexeme
    integer :: i, n_children
    call lexer_clear (lexer)
    call stream_init (stream, string)
    call lexer_assign_stream (lexer, stream)
    do i = 1, n_token
       call lex (lexeme(i), lexer)
    end do
    call stream_final (stream)
    n_children = get_n_children ()
    call set_delimiters
    if (n_children > 1)  call set_separator
    if (n_children > 0)  call set_children
  contains
    function get_n_children () result (n_children)
      integer :: n_children
      select case (rule%type)
      case (S_QUOTED)
         if (n_token /= 6)  call broken_rule (rule)
         n_children = 0
      case (S_GROUP)
         if (n_token /= 6)  call broken_rule (rule)
         n_children = 1
      case (S_SEQUENCE)
         if (is_modifier (lexeme_get_string (lexeme(n_token)))) then
            if (n_token <= 4)  call broken_rule (rule)
            call syntax_rule_set_modifier &
                 (rule, lexeme_get_string (lexeme(n_token)))
            n_children = n_token - 4
         else 
            if (n_token <= 3) call broken_rule (rule)
            n_children = n_token - 3
         end if
      case (S_LIST)
         if (is_modifier (lexeme_get_string (lexeme(n_token)))) then
            if (n_token <= 4 .or. mod (n_token, 2) /= 1) &
                 call broken_rule (rule)
            call syntax_rule_set_modifier &
                 (rule, lexeme_get_string (lexeme(n_token)))
         else if (n_token <= 3 .or. mod (n_token, 2) /= 0) then
            call broken_rule (rule)
         end if
         n_children = (n_token - 2) / 2
      case (S_ARGS)
         if (is_modifier (lexeme_get_string (lexeme(n_token-1)))) then
            if (n_token <= 6 .or. mod (n_token, 2) /= 1) &
                 call broken_rule (rule)
            call syntax_rule_set_modifier &
                 (rule, lexeme_get_string (lexeme(n_token-1)))
         else if (n_token <= 5 .or. mod (n_token, 2) /= 0) then
            call broken_rule (rule)
         end if
         n_children = (n_token - 4) / 2
      case (S_ALTERNATIVE)
         if (n_token <= 3 .or. mod (n_token, 2) /= 0)  call broken_rule (rule)
         n_children = (n_token - 2) / 2
      end select
    end function get_n_children
    subroutine set_delimiters
      type(string_t), dimension(2) :: delimiter
      select case (rule%type)
      case (S_QUOTED, S_GROUP, S_ARGS)
         delimiter(1) = lexeme_get_contents (lexeme(4))
         delimiter(2) = lexeme_get_contents (lexeme(n_token))
         call syntax_rule_set_delimiter (rule, delimiter)
      end select
    end subroutine set_delimiters
    subroutine set_separator
      type(string_t) :: separator
      select case (rule%type)
      case (S_LIST)
         separator = lexeme_get_contents (lexeme(5))
         call syntax_rule_set_separator (rule, separator)
      case (S_ARGS)
         separator = lexeme_get_contents (lexeme(6))
         call syntax_rule_set_separator (rule, separator)
      end select
    end subroutine set_separator
    subroutine set_children
      allocate (rule%child(n_children))
      select case (rule%type)
      case (S_GROUP)
         call syntax_rule_set_sub (rule, 1, syntax_get_rule_ptr (syntax, &
              lexeme_get_contents (lexeme(5))))
      case (S_SEQUENCE)
         do i = 1, n_children
            call syntax_rule_set_sub (rule, i, syntax_get_rule_ptr (syntax, &
                 lexeme_get_contents (lexeme(i+3))))
         end do
      case (S_LIST, S_ALTERNATIVE)
         do i = 1, n_children
            call syntax_rule_set_sub (rule, i, syntax_get_rule_ptr (syntax, &
                 lexeme_get_contents (lexeme(2*i+2))))
         end do
      case (S_ARGS)
         do i = 1, n_children
            call syntax_rule_set_sub (rule, i, syntax_get_rule_ptr (syntax, &
                 lexeme_get_contents (lexeme(2*i+3))))
         end do
      end select
    end subroutine set_children
    subroutine broken_rule (rule)
      type(syntax_rule_t), intent(in) :: rule
      call lexer_show_location (lexer)
      call msg_bug (" Syntax definition: broken rule '" &
           // char (wkey (rule)) // "'")
    end subroutine broken_rule
  end subroutine set_rule_contents

\eatline
\nwendcode{}\nwbegindocs{251}This routine completes the syntax table object. We assume that the
rule array is set up. We associate the top rule with the first entry
in the rule array and build up the keyword list.

The keyword list includes delimiters and separators.  Filling it can
only be done after all rules are set.  We scan the rule tree.  For
each keyword that we find, we try to add it to the keyword list; the
pointer to the last element is carried along with the recursive
scanning.  Before appending a keyword, we check whether it is already
in the list.
\nwenddocs{}\nwbegincode{252}\moddef{Syntax: subroutines}\plusendmoddef
  subroutine syntax_make_keyword_list (syntax)
    type(syntax_t), intent(inout), target :: syntax
    type(syntax_rule_t), pointer :: rule
    rule => syntax%rule(1)
    call rule_scan_rec (rule, syntax%keyword_list)
  contains
    recursive subroutine rule_scan_rec (rule, keyword_list)
      type(syntax_rule_t), pointer :: rule
      type(keyword_list_t), intent(inout) :: keyword_list
      integer :: i
      if (rule%used)  return
      rule%used = .true.
      select case (rule%type)
      case (S_UNKNOWN)
         call msg_bug (" Syntax: rule tree contains undefined rule")
      case (S_KEYWORD)
         call keyword_list_add (keyword_list, rule%keyword)
      end select
      select case (rule%type)
      case (S_LIST, S_ARGS)
         call keyword_list_add (keyword_list, rule%separator)
      end select
      select case (rule%type)
      case (S_GROUP, S_ARGS)
         call keyword_list_add (keyword_list, rule%delimiter(1))
         call keyword_list_add (keyword_list, rule%delimiter(2))
      end select
      select case (rule%type)
      case (S_SEQUENCE, S_GROUP, S_LIST, S_ARGS, S_ALTERNATIVE)
         if (.not. allocated (rule%child)) &
              call msg_bug (" Syntax: Non-terminal rule without children")
      case default
         if (allocated (rule%child)) &
              call msg_bug (" Syntax: Terminal rule with children")
      end select
      if (allocated (rule%child)) then
         do i = 1, size (rule%child)
            call rule_scan_rec (rule%child(i)%p, keyword_list)
         end do
      end if
    end subroutine rule_scan_rec
  end subroutine syntax_make_keyword_list

\eatline
\nwendcode{}\nwbegindocs{253}The finalizer deallocates the rule pointer array and deletes the
keyword list.
\nwenddocs{}\nwbegincode{254}\moddef{Syntax: subroutines}\plusendmoddef
  subroutine syntax_final (syntax)
    type(syntax_t), intent(inout) :: syntax
    if (allocated (syntax%rule))  deallocate (syntax%rule)
    call keyword_list_final (syntax%keyword_list)
  end subroutine syntax_final

\eatline
\nwendcode{}\nwbegindocs{255}\nwdocspar
\subsection{Accessing the syntax table}
Return a pointer to the top rule
\nwenddocs{}\nwbegincode{256}\moddef{Syntax: public}\plusendmoddef
  public :: syntax_get_top_rule_ptr
\nwendcode{}\nwbegincode{257}\moddef{Syntax: subroutines}\plusendmoddef
  function syntax_get_top_rule_ptr (syntax) result (rule)
    type(syntax_rule_t), pointer :: rule
    type(syntax_t), intent(in), target :: syntax
    if (allocated (syntax%rule)) then
       rule => syntax%rule(1)
    else
       rule => null ()
    end if
  end function syntax_get_top_rule_ptr

\eatline
\nwendcode{}\nwbegindocs{258}Assign the pointer to the rule associated with a given key (assumes
that the rule array is allocated)
\nwenddocs{}\nwbegincode{259}\moddef{Syntax: public}\plusendmoddef
  public :: syntax_get_rule_ptr
\nwendcode{}\nwbegincode{260}\moddef{Syntax: subroutines}\plusendmoddef
  function syntax_get_rule_ptr (syntax, key) result (rule)
    type(syntax_rule_t), pointer :: rule
    type(syntax_t), intent(in), target :: syntax
    type(string_t), intent(in) :: key
    integer :: i
    do i = 1, size (syntax%rule)
       if (syntax%rule(i)%keyword == key) then
          rule => syntax%rule(i)
          return
       end if
    end do
    call msg_bug (" Syntax table: Rule " // char (key) // " not found")
  end function syntax_get_rule_ptr

\eatline
\nwendcode{}\nwbegindocs{261}Return a pointer to the keyword list
\nwenddocs{}\nwbegincode{262}\moddef{Syntax: public}\plusendmoddef
  public :: syntax_get_keyword_list_ptr
\nwendcode{}\nwbegincode{263}\moddef{Syntax: subroutines}\plusendmoddef
  function syntax_get_keyword_list_ptr (syntax) result (keyword_list)
    type(keyword_list_t), pointer :: keyword_list
    type(syntax_t), intent(in), target :: syntax
    keyword_list => syntax%keyword_list
  end function syntax_get_keyword_list_ptr

\eatline
\nwendcode{}\nwbegindocs{264}\nwdocspar
\subsection{I/O}
Write a readable representation of the syntax table
\nwenddocs{}\nwbegincode{265}\moddef{Syntax: public}\plusendmoddef
  public :: syntax_write
\nwendcode{}\nwbegincode{266}\moddef{Syntax: subroutines}\plusendmoddef
  subroutine syntax_write (syntax, unit)
    type(syntax_t), intent(in) :: syntax
    integer, intent(in), optional :: unit
    integer :: u
    integer :: i
    u = given_output_unit (unit);  if (u < 0)  return
    write (u, "(A)") "Syntax table:"
    if (allocated (syntax%rule)) then
       do i = 1, size (syntax%rule)
          call syntax_rule_write (syntax%rule(i), u)
       end do
    else
       write (u, "(1x,A)") "[not allocated]"
    end if
    call keyword_list_write (syntax%keyword_list, u)
  end subroutine syntax_write

\eatline
\nwendcode{}\nwbegindocs{267}\nwdocspar
\clearpage
%%%%%%%%%%%%%%%%%%%%%%%%%%%%%%%%%%%%%%%%%%%%%%%%%%%%%%%%%%%%%%%%%%%%%%%%
\section{The parser}
On a small scale, the parser interprets the string tokens returned by
the lexer; they are interpreted as numbers, keywords and such and
stored as a typed object.  On a large scale, a text is read, parsed,
and a syntax rule set is applied such that the tokens are stored as a
parse tree.  Syntax errors are spotted in this process, so the
resulting parse tree is syntactically correct by definition.
\nwenddocs{}\nwbegincode{268}\moddef{\code{}parser.f90\edoc{}}\endmoddef
\LA{}File header\RA{}

module parser

\LA{}Use kinds\RA{}
\LA{}Use strings\RA{}
  use io_units
  use format_defs, only: FMT_19
  use unit_tests
  use system_defs, only: DIGITS
  use diagnostics
  use md5
  use lexers
  use syntax_rules
  
\LA{}Standard module head\RA{}

\LA{}Parser: public\RA{}

\LA{}Parser: types\RA{}

\LA{}Parser: interfaces\RA{}

contains

\LA{}Parser: procedures\RA{}

\LA{}Parser: tests\RA{}

end module parser
\eatline
\nwendcode{}\nwbegindocs{269}\nwdocspar
\subsection{The token type}
Tokens are elements of the parsed input that carry a value: logical,
integer, real, quoted string, (unquoted) identifier, or known
keyword.  Note that non-keyword tokens also have an abstract key
attached to them.

This is an obvious candidate for polymorphism.
\nwenddocs{}\nwbegincode{270}\moddef{Parser: types}\endmoddef
  type :: token_t
     private
     integer :: type = S_UNKNOWN
     logical, pointer :: lval => null ()
     integer, pointer :: ival => null ()
     real(default), pointer :: rval => null ()
     complex(default), pointer :: cval => null ()
     type(string_t), pointer :: sval => null ()
     type(string_t), pointer :: kval => null ()
     type(string_t), dimension(:), pointer :: quote => null ()
  end type token_t

\eatline
\nwendcode{}\nwbegindocs{271}Create a token from the lexeme returned by the lexer: Allocate
storage and try to interpret the lexeme according to the type that is
requested by the parser.  For a keyword token, match the lexeme
against the requested key.  If successful, set the token type, value,
and key.  Otherwise, set the type to {\tt{}S{\char95}UNKNOWN}.
\nwenddocs{}\nwbegincode{272}\moddef{Parser: procedures}\endmoddef
  subroutine token_init (token, lexeme, requested_type, key)
    type(token_t), intent(out) :: token
    type(lexeme_t), intent(in) :: lexeme
    integer, intent(in) :: requested_type
    type(string_t), intent(in) :: key
    integer :: type
    type = lexeme_get_type (lexeme)
    token%type = S_UNKNOWN
    select case (requested_type)
    case (S_LOGICAL)
       if (type == T_IDENTIFIER)  call read_logical &
            (char (lexeme_get_string (lexeme)))
    case (S_INTEGER)
       if (type == T_NUMERIC)  call read_integer &
            (char (lexeme_get_string (lexeme)))
    case (S_REAL)
       if (type == T_NUMERIC)  call read_real &
            (char (lexeme_get_string (lexeme)))
    case (S_COMPLEX)
       if (type == T_NUMERIC)  call read_complex &
            (char (lexeme_get_string (lexeme)))
    case (S_IDENTIFIER)
       if (type == T_IDENTIFIER)  call read_identifier &
            (lexeme_get_string (lexeme))
    case (S_KEYWORD)
       if (type == T_KEYWORD)  call check_keyword &
            (lexeme_get_string (lexeme), key)
    case (S_QUOTED)
       if (type == T_QUOTED)   call read_quoted &
            (lexeme_get_contents (lexeme), lexeme_get_delimiters (lexeme))
    case default
       print *, requested_type
       call msg_bug (" Invalid token type code requested by the parser")
    end select
    if (token%type /= S_UNKNOWN) then
       allocate (token%kval)
       token%kval = key
    end if
  contains
    subroutine read_logical (s)
      character(*), intent(in) :: s
      select case (s)
      case ("t", "T", "true", "TRUE", "y", "Y", "yes", "YES")
         allocate (token%lval)
         token%lval = .true.
         token%type = S_LOGICAL
      case ("f", "F", "false", "FALSE", "n", "N", "no", "NO")
         allocate (token%lval)
         token%lval = .false.
         token%type = S_LOGICAL
      end select
    end subroutine read_logical
    subroutine read_integer (s)
      character(*), intent(in) :: s
      integer :: tmp, iostat
      if (verify (s, DIGITS) == 0) then
         read (s, *, iostat=iostat) tmp
         if (iostat == 0) then
            allocate (token%ival)
            token%ival = tmp
            token%type = S_INTEGER
         end if
      end if
    end subroutine read_integer
    subroutine read_real (s)
      character(*), intent(in) :: s
      real(default) :: tmp
      integer :: iostat
      read (s, *, iostat=iostat) tmp
      if (iostat == 0) then
         allocate (token%rval)
         token%rval = tmp
         token%type = S_REAL
      end if
    end subroutine read_real
    subroutine read_complex (s)
      character(*), intent(in) :: s
      complex(default) :: tmp
      integer :: iostat
      read (s, *, iostat=iostat) tmp
      if (iostat == 0) then
         allocate (token%cval)
         token%cval = tmp
         token%type = S_COMPLEX
      end if
    end subroutine read_complex
    subroutine read_identifier (s)
      type(string_t), intent(in) :: s
      allocate (token%sval)
      token%sval = s
      token%type = S_IDENTIFIER
    end subroutine read_identifier
    subroutine check_keyword (s, key)
      type(string_t), intent(in) :: s
      type(string_t), intent(in) :: key
      if (key == s)  token%type = S_KEYWORD
    end subroutine check_keyword
    subroutine read_quoted (s, del)
      type(string_t), intent(in) :: s
      type(string_t), dimension(2), intent(in) :: del
      allocate (token%sval, token%quote(2))
      token%sval = s
      token%quote(1) = del(1)
      token%quote(2) = del(2)
      token%type = S_QUOTED
    end subroutine read_quoted
  end subroutine token_init

\eatline
\nwendcode{}\nwbegindocs{273}Manually set a token to a keyword.
\nwenddocs{}\nwbegincode{274}\moddef{Parser: procedures}\plusendmoddef
  subroutine token_init_key (token, key)
    type(token_t), intent(out) :: token
    type(string_t), intent(in) :: key
    token%type = S_KEYWORD
    allocate (token%kval)
    token%kval = key
  end subroutine token_init_key
  
\eatline
\nwendcode{}\nwbegindocs{275}Reset a token to an empty state, freeing allocated memory, and
deallocate the token itself.
\nwenddocs{}\nwbegincode{276}\moddef{Parser: procedures}\plusendmoddef
  subroutine token_final (token)
    type(token_t), intent(inout) :: token
    token%type = S_UNKNOWN
    if (associated (token%lval))  deallocate (token%lval)
    if (associated (token%ival))  deallocate (token%ival)
    if (associated (token%rval))  deallocate (token%rval)
    if (associated (token%sval))  deallocate (token%sval)
    if (associated (token%kval))  deallocate (token%kval)
    if (associated (token%quote))  deallocate (token%quote)
  end subroutine token_final

\eatline
\nwendcode{}\nwbegindocs{277}Check for empty=valid token:
\nwenddocs{}\nwbegincode{278}\moddef{Parser: procedures}\plusendmoddef
  function token_is_valid (token) result (valid)
    logical :: valid
    type(token_t), intent(in) :: token
    valid = token%type /= S_UNKNOWN
  end function token_is_valid

\eatline
\nwendcode{}\nwbegindocs{279}Write the contents of a token.
\nwenddocs{}\nwbegincode{280}\moddef{Parser: procedures}\plusendmoddef
  subroutine token_write (token, unit)
    type(token_t), intent(in) :: token
    integer, intent(in), optional :: unit
    integer :: u
    u = given_output_unit (unit);  if (u < 0)  return
    select case (token%type)
    case (S_LOGICAL)
       write (u, "(L1)") token%lval
    case (S_INTEGER)
       write (u, "(I0)") token%ival
    case (S_REAL)
       write (u, "(" // FMT_19 // ")") token%rval
    case (S_COMPLEX)
       write (u, "('('," // FMT_19 // ",','," // FMT_19 // ",')')") token%cval
    case (S_IDENTIFIER)
       write (u, "(A)") char (token%sval)
    case (S_KEYWORD)
       write (u, "(A,A)") '[keyword] ' // char (token%kval)
    case (S_QUOTED)
       write (u, "(A)") &
            char (token%quote(1)) // char (token%sval) // char (token%quote(2))
    case default
       write (u, "(A)") '[empty]'
    end select
  end subroutine token_write

\eatline
\nwendcode{}\nwbegindocs{281}Token assignment via deep copy.  This is useful to avoid confusion
when the token is transferred to some parse-tree node.
\nwenddocs{}\nwbegincode{282}\moddef{Parser: interfaces}\endmoddef
  interface assignment(=)
     module procedure token_assign
     module procedure token_assign_integer
     module procedure token_assign_real
     module procedure token_assign_complex
     module procedure token_assign_logical
     module procedure token_assign_string
  end interface

\eatline
\nwendcode{}\nwbegindocs{283}We need to copy only the contents that are actually assigned, the
other pointers remain disassociated.
\nwenddocs{}\nwbegincode{284}\moddef{Parser: procedures}\plusendmoddef
  subroutine token_assign (token, token_in)
    type(token_t), intent(out) :: token
    type(token_t), intent(in) :: token_in
    token%type = token_in%type
    select case (token%type)
    case (S_LOGICAL);    allocate (token%lval);  token%lval = token_in%lval
    case (S_INTEGER);    allocate (token%ival);  token%ival = token_in%ival
    case (S_REAL);       allocate (token%rval);  token%rval = token_in%rval
    case (S_COMPLEX); allocate (token%cval);  token%cval = token_in%cval    
    case (S_IDENTIFIER); allocate (token%sval);  token%sval = token_in%sval
    case (S_QUOTED);     allocate (token%sval);  token%sval = token_in%sval
       allocate (token%quote(2));  token%quote = token_in%quote
    end select
    if (token%type /= S_UNKNOWN) then
       allocate (token%kval);  token%kval = token_in%kval
    end if
  end subroutine token_assign

\eatline
\nwendcode{}\nwbegindocs{285}We need to copy only the contents that are actually assigned, the
other pointers remain disassociated.
\nwenddocs{}\nwbegincode{286}\moddef{Parser: procedures}\plusendmoddef
  subroutine token_assign_integer (token, ival)
    type(token_t), intent(out) :: token
    integer, intent(in) :: ival
    token%type = S_INTEGER
    allocate (token%ival)
    token%ival = ival
  end subroutine token_assign_integer

  subroutine token_assign_real (token, rval)
    type(token_t), intent(out) :: token
    real(default), intent(in) :: rval
    token%type = S_REAL
    allocate (token%rval)
    token%rval = rval
  end subroutine token_assign_real

  subroutine token_assign_complex (token, cval)
    type(token_t), intent(out) :: token
    complex(default), intent(in) :: cval
    token%type = S_COMPLEX
    allocate (token%cval)
    token%cval = cval
  end subroutine token_assign_complex

  subroutine token_assign_logical (token, lval)
    type(token_t), intent(out) :: token
    logical, intent(in) :: lval
    token%type = S_LOGICAL
    allocate (token%lval)
    token%lval = lval
  end subroutine token_assign_logical

  subroutine token_assign_string (token, sval)
    type(token_t), intent(out) :: token
    type(string_t), intent(in) :: sval
    token%type = S_QUOTED
    allocate (token%sval)
    token%sval = sval
    allocate (token%quote(2));  token%quote = '"'
  end subroutine token_assign_string

\eatline
\nwendcode{}\nwbegindocs{287}\nwdocspar
\subsection{Retrieve token contents}
These functions all do a trivial sanity check that should avoid
crashes.
\nwenddocs{}\nwbegincode{288}\moddef{Parser: procedures}\plusendmoddef
  function token_get_logical (token) result (lval)
    logical :: lval
    type(token_t), intent(in) :: token
    if (associated (token%lval)) then
       lval = token%lval
    else
       call token_mismatch (token, "logical")
    end if
  end function token_get_logical

  function token_get_integer (token) result (ival)
    integer :: ival
    type(token_t), intent(in) :: token
    if (associated (token%ival)) then
       ival = token%ival
    else
       call token_mismatch (token, "integer")
    end if
  end function token_get_integer

  function token_get_real (token) result (rval)
    real(default) :: rval
    type(token_t), intent(in) :: token
    if (associated (token%rval)) then
       rval = token%rval
    else
       call token_mismatch (token, "real")
    end if
  end function token_get_real
  
  function token_get_cmplx (token) result (cval)
    complex(default) :: cval
    type(token_t), intent(in) :: token
    if (associated (token%cval)) then
       cval = token%cval
    else
       call token_mismatch (token, "complex")
    end if
  end function token_get_cmplx

  function token_get_string (token) result (sval)
    type(string_t) :: sval
    type(token_t), intent(in) :: token
    if (associated (token%sval)) then
       sval = token%sval
    else
       call token_mismatch (token, "string")
    end if
  end function token_get_string

  function token_get_key (token) result (kval)
    type(string_t) :: kval
    type(token_t), intent(in) :: token
    if (associated (token%kval)) then
       kval = token%kval
    else
       call token_mismatch (token, "keyword")
    end if
  end function token_get_key

  function token_get_quote (token) result (quote)
    type(string_t), dimension(2) :: quote
    type(token_t), intent(in) :: token
    if (associated (token%quote)) then
       quote = token%quote
    else
       call token_mismatch (token, "quote")
    end if
  end function token_get_quote

\eatline
\eatline
\nwendcode{}\nwbegincode{289}\moddef{Parser: procedures}\plusendmoddef
  subroutine token_mismatch (token, type)
    type(token_t), intent(in) :: token
    character(*), intent(in) :: type
    write (6, "(A)", advance="no")  "Token: "
    call token_write (token)
    call msg_bug (" Token type mismatch; value required as " // type)
  end subroutine token_mismatch

\eatline
\nwendcode{}\nwbegindocs{290}\nwdocspar
\subsection{The parse tree: nodes}
The parser will generate a parse tree from the input stream.  Each
node in this parse tree points to the syntax rule that was applied.
(Since syntax rules are stored in a pointer-type array within the
syntax table, they qualify as targets.)  A leaf node contains a
token.  A branch node has subnodes.  The subnodes are stored as a
list, so each node also has a {\tt{}next} pointer.
\nwenddocs{}\nwbegincode{291}\moddef{Parser: public}\endmoddef
  public :: parse_node_t
\nwendcode{}\nwbegincode{292}\moddef{Parser: types}\plusendmoddef
  type :: parse_node_t
     private
     type(syntax_rule_t), pointer :: rule => null ()
     type(token_t) :: token
     integer :: n_sub = 0
     type(parse_node_t), pointer :: sub_first => null ()
     type(parse_node_t), pointer :: sub_last => null ()
     type(parse_node_t), pointer :: next => null ()
   contains
   \LA{}Parser: parse node: TBP\RA{}
  end type parse_node_t

\eatline
\nwendcode{}\nwbegindocs{293}Container for parse node pointers, useful for creating pointer arrays:
\nwenddocs{}\nwbegincode{294}\moddef{Parser: public}\plusendmoddef
  public :: parse_node_p
\nwendcode{}\nwbegincode{295}\moddef{Parser: types}\plusendmoddef
  type :: parse_node_p
    type(parse_node_t), pointer :: ptr => null ()
  end type parse_node_p

\eatline
\nwendcode{}\nwbegindocs{296}Output.  The first version writes a node together with its
sub-node tree, organized by indentation.
\nwenddocs{}\nwbegincode{297}\moddef{Parser: parse node: TBP}\endmoddef
  procedure :: write => parse_node_write_rec
\nwendcode{}\nwbegincode{298}\moddef{Parser: public}\plusendmoddef
  public :: parse_node_write_rec
\nwendcode{}\nwbegincode{299}\moddef{Parser: procedures}\plusendmoddef
  recursive subroutine parse_node_write_rec (node, unit, short, depth)
    class(parse_node_t), intent(in), target :: node
    integer, intent(in), optional :: unit
    logical, intent(in), optional :: short
    integer, intent(in), optional :: depth
    integer :: u, d
    type(parse_node_t), pointer :: current
    u = given_output_unit (unit);  if (u < 0)  return
    d = 0;  if (present (depth))  d = depth
    call parse_node_write (node, u, short=short)
    current => node%sub_first
    do while (associated (current))
       write (u, "(A)", advance = "no")  repeat ("|  ", d)
       call parse_node_write_rec (current, unit, short, d+1)
       current => current%next
    end do
  end subroutine parse_node_write_rec

\eatline
\nwendcode{}\nwbegindocs{300}This does the actual output for a single node, without recursion.
\nwenddocs{}\nwbegincode{301}\moddef{Parser: public}\plusendmoddef
  public :: parse_node_write
\nwendcode{}\nwbegincode{302}\moddef{Parser: procedures}\plusendmoddef
  subroutine parse_node_write (node, unit, short)
    class(parse_node_t), intent(in) :: node
    integer, intent(in), optional :: unit
    logical, intent(in), optional :: short
    integer :: u
    type(parse_node_t), pointer :: current
    u = given_output_unit (unit);  if (u < 0)  return
    write (u, "('+ ')", advance = "no")
    if (associated (node%rule)) then
       call syntax_rule_write (node%rule, u, &
            short=short, key_only=.true., advance=.false.)
       if (token_is_valid (node%token)) then
          write (u, "('  = ')", advance="no")
          call token_write (node%token, u)
       else if (associated (node%sub_first)) then
          write (u, "('  = ')", advance="no")
          current => node%sub_first
          do while (associated (current))
             call syntax_rule_write (current%rule, u, &
                  short=.true., key_only=.true., advance=.false.)
             current => current%next
          end do
          write (u, *)
       else
          write (u, *)
       end if
    else
       write (u, *) "[empty]"
    end if
  end subroutine parse_node_write

\eatline
\nwendcode{}\nwbegindocs{303}Finalize the token and recursively finalize and deallocate all
sub-nodes.
\nwenddocs{}\nwbegincode{304}\moddef{Parser: public}\plusendmoddef
  public :: parse_node_final
\nwendcode{}\nwbegincode{305}\moddef{Parser: procedures}\plusendmoddef
  recursive subroutine parse_node_final (node, recursive)
    type(parse_node_t), intent(inout) :: node
    type(parse_node_t), pointer :: current
    logical, intent(in), optional :: recursive
    logical :: rec
    rec = .true.;  if (present (recursive))  rec = recursive
    call token_final (node%token)
    if (rec) then
       do while (associated (node%sub_first))
          current => node%sub_first
          node%sub_first => node%sub_first%next
          call parse_node_final (current)
          deallocate (current)
       end do
    end if
  end subroutine parse_node_final

\eatline
\nwendcode{}\nwbegindocs{306}\nwdocspar
\subsection{Filling nodes}
The constructors allocate and initialize the node.  There are two
possible initializers (in a later version, should correspond to
different type extensions).

First, leaf (terminal) nodes.  The token constructor does the actual
work, looking at the requested type and key for the given rule and
matching against the lexeme contents.  If it fails, the token will
keep the type {\tt{}S{\char95}UNKNOWN} and remain empty.  Otherwise, we have a
valid node which contains the new token.

If the lexeme argument is absent, the token is left empty.
\nwenddocs{}\nwbegincode{307}\moddef{Parser: procedures}\plusendmoddef
  subroutine parse_node_create_leaf (node, rule, lexeme)
    type(parse_node_t), pointer :: node
    type(syntax_rule_t), intent(in), target :: rule
    type(lexeme_t), intent(in) :: lexeme
    allocate (node)
    node%rule => rule
    call token_init (node%token, lexeme, &
         syntax_rule_get_type (rule), syntax_rule_get_key (rule))
    if (.not. token_is_valid (node%token))  deallocate (node)
  end subroutine parse_node_create_leaf

\eatline
\nwendcode{}\nwbegindocs{308}This version allows us to manually create a leaf node that holds a keyword.
\nwenddocs{}\nwbegincode{309}\moddef{Parser: public}\plusendmoddef
  public :: parse_node_create_key
\nwendcode{}\nwbegincode{310}\moddef{Parser: procedures}\plusendmoddef
  subroutine parse_node_create_key (node, rule)
    type(parse_node_t), intent(out) :: node
    type(syntax_rule_t), intent(in), target :: rule
    node%rule => rule
    call token_init_key (node%token, syntax_rule_get_key (rule))
  end subroutine parse_node_create_key
    
\eatline
\nwendcode{}\nwbegindocs{311}This version allows us to manually create a leaf node that holds a fixed
value.  Only one of the optional values should be provided.
\nwenddocs{}\nwbegincode{312}\moddef{Parser: public}\plusendmoddef
  public :: parse_node_create_value
\nwendcode{}\nwbegincode{313}\moddef{Parser: procedures}\plusendmoddef
  subroutine parse_node_create_value (node, rule, ival, rval, cval, sval, lval)
    type(parse_node_t), intent(out) :: node
    type(syntax_rule_t), intent(in), target :: rule
    integer, intent(in), optional :: ival
    real(default), intent(in), optional :: rval
    complex(default), intent(in), optional :: cval
    type(string_t), intent(in), optional :: sval
    logical, intent(in), optional :: lval
    node%rule => rule
    call parse_node_set_value (node, ival, rval, cval, sval, lval)
  end subroutine parse_node_create_value
    
\eatline
\nwendcode{}\nwbegindocs{314}Directly set the value without changing anything else.
\nwenddocs{}\nwbegincode{315}\moddef{Parser: public}\plusendmoddef
  public :: parse_node_set_value
\nwendcode{}\nwbegincode{316}\moddef{Parser: procedures}\plusendmoddef
  subroutine parse_node_set_value (node, ival, rval, cval, sval, lval)
    type(parse_node_t), intent(inout) :: node
    integer, intent(in), optional :: ival
    real(default), intent(in), optional :: rval
    complex(default), intent(in), optional :: cval
    type(string_t), intent(in), optional :: sval
    logical, intent(in), optional :: lval
    if (present (ival)) then
       node%token = ival
    else if (present (rval)) then
       node%token = rval
    else if (present (cval)) then
       node%token = cval
    else if (present (lval)) then
       node%token = lval
    else if (present (sval)) then
       node%token = sval
    end if
  end subroutine parse_node_set_value
    
\eatline
\nwendcode{}\nwbegindocs{317}Second, branch nodes.  We first assign the rule:
\nwenddocs{}\nwbegincode{318}\moddef{Parser: public}\plusendmoddef
  public :: parse_node_create_branch
\nwendcode{}\nwbegincode{319}\moddef{Parser: procedures}\plusendmoddef
  subroutine parse_node_create_branch (node, rule)
    type(parse_node_t), pointer :: node
    type(syntax_rule_t), intent(in), target :: rule
    allocate (node)
    node%rule => rule
  end subroutine parse_node_create_branch
    
\eatline
\nwendcode{}\nwbegindocs{320}Append a sub-node.  The sub-node must exist and be a valid target,
otherwise nothing is done.
\nwenddocs{}\nwbegincode{321}\moddef{Parser: public}\plusendmoddef
  public :: parse_node_append_sub
\nwendcode{}\nwbegincode{322}\moddef{Parser: procedures}\plusendmoddef
  subroutine parse_node_append_sub (node, sub)
    type(parse_node_t), intent(inout) :: node
    type(parse_node_t), pointer :: sub
    if (associated (sub)) then
       if (associated (node%sub_last)) then
          node%sub_last%next => sub
       else
          node%sub_first => sub
       end if
       node%sub_last => sub
    end if
  end subroutine parse_node_append_sub
  
\eatline
\nwendcode{}\nwbegindocs{323}For easy access, once the list is complete we count the number of
sub-nodes.  If there are no subnodes, the whole node is deleted.
\nwenddocs{}\nwbegincode{324}\moddef{Parser: public}\plusendmoddef
  public :: parse_node_freeze_branch
\nwendcode{}\nwbegincode{325}\moddef{Parser: procedures}\plusendmoddef
  subroutine parse_node_freeze_branch (node)
    type(parse_node_t), pointer :: node
    type(parse_node_t), pointer :: current
    node%n_sub = 0
    current => node%sub_first
    do while (associated (current))
       node%n_sub = node%n_sub + 1
       current => current%next
    end do
    if (node%n_sub == 0)  deallocate (node)
  end subroutine parse_node_freeze_branch

\eatline
\nwendcode{}\nwbegindocs{326}Replace the syntax rule.  This makes sense only if the parse node
adheres to the syntax of the new rule.
\nwenddocs{}\nwbegincode{327}\moddef{Parser: public}\plusendmoddef
  public :: parse_node_replace_rule
\nwendcode{}\nwbegincode{328}\moddef{Parser: procedures}\plusendmoddef
  subroutine parse_node_replace_rule (node, rule)
    type(parse_node_t), pointer :: node
    type(syntax_rule_t), intent(in), target :: rule
    node%rule => rule
  end subroutine parse_node_replace_rule
  
\eatline
\nwendcode{}\nwbegindocs{329}Replace the last subnode by the target node.  Since the subnodes are
stored as a linked list, we can do this only if we copy the subnodes.
Furthermore, the target node must also be copied, and the {\tt{}next}
pointer of the copy is nullified.  This ensures that we cannot modify
the originals at the subnode level.

All copies are shallow copies.  This implies that further
modifications at the sub-subnode level will affect the original nodes
and must therefore be forbidden.

Use with care, this invites to memory mismanagement.  The copy nodes
can be deallocated, but not finalized, since its subnodes are the same
objects as the subnodes of the target node.
\nwenddocs{}\nwbegincode{330}\moddef{Parser: public}\plusendmoddef
  public :: parse_node_replace_last_sub
\nwendcode{}\nwbegincode{331}\moddef{Parser: procedures}\plusendmoddef
  subroutine parse_node_replace_last_sub (node, pn_target)
    type(parse_node_t), intent(inout), target :: node
    type(parse_node_t), intent(in), target :: pn_target
    type(parse_node_t), pointer :: current, current_copy, previous
    integer :: i
    select case (node%n_sub)
    case (1)
       allocate (current_copy)
       current_copy = pn_target
       node%sub_first => current_copy
    case (2:)
       current => node%sub_first
       allocate (current_copy)
       current_copy = current
       node%sub_first => current_copy
       previous => current_copy
       do i = 1, node%n_sub - 2
          current => current%next
          allocate (current_copy)
          current_copy = current
          previous%next => current_copy
          previous => current_copy
       end do
       allocate (current_copy)
       current_copy = pn_target
       previous%next => current_copy
    case default
       call parse_node_write (node)
       call msg_bug ("'replace_last_sub' called for non-branch parse node")
    end select
    current_copy%next => null ()
    node%sub_last => current_copy
  end subroutine parse_node_replace_last_sub

\eatline
\nwendcode{}\nwbegindocs{332}\nwdocspar
\subsection{Accessing nodes}
Return the node contents.  Check if pointers are associated.  No
check when accessing a sub-node; assume that {\tt{}parse{\char95}node{\char95}n{\char95}sub} is always
used for the upper bound.

The token extractor returns a pointer.
\nwenddocs{}\nwbegincode{333}\moddef{Parser: public}\plusendmoddef
  public :: parse_node_get_rule_ptr
  public :: parse_node_get_n_sub
  public :: parse_node_get_sub_ptr
  public :: parse_node_get_next_ptr
  public :: parse_node_get_last_sub_ptr
\nwendcode{}\nwbegincode{334}\moddef{Parser: procedures}\plusendmoddef
  function parse_node_get_rule_ptr (node) result (rule)
    type(syntax_rule_t), pointer :: rule
    type(parse_node_t), intent(in), target :: node
    if (associated (node%rule)) then
       rule => node%rule
    else
       rule => null ()
       call parse_node_undefined (node, "rule")
    end if
  end function parse_node_get_rule_ptr

  function parse_node_get_n_sub (node) result (n)
    integer :: n
    type(parse_node_t), intent(in) :: node 
    n = node%n_sub
  end function parse_node_get_n_sub

  function parse_node_get_sub_ptr (node, n, tag, required) result (sub)
    type(parse_node_t), pointer :: sub
    type(parse_node_t), intent(in), target :: node
    integer, intent(in), optional :: n
    character(*), intent(in), optional :: tag
    logical, intent(in), optional :: required
    integer :: i
    sub => node%sub_first
    if (present (n)) then
       do i = 2, n
          if (associated (sub)) then
             sub => sub%next
          else
             return
          end if
       end do
    end if
    call parse_node_check (sub, tag, required)
  end function parse_node_get_sub_ptr
  
  function parse_node_get_next_ptr (sub, n, tag, required) result (next)
    type(parse_node_t), pointer :: next
    type(parse_node_t), intent(in), target :: sub
    integer, intent(in), optional :: n
    character(*), intent(in), optional :: tag
    logical, intent(in), optional :: required
    integer :: i
    next => sub%next
    if (present (n)) then
       do i = 2, n
          if (associated (next)) then
             next => next%next
          else
             exit
          end if
       end do
    end if
    call parse_node_check (next, tag, required)
  end function parse_node_get_next_ptr
  
  function parse_node_get_last_sub_ptr (node, tag, required) result (sub)
    type(parse_node_t), pointer :: sub
    type(parse_node_t), intent(in), target :: node
    character(*), intent(in), optional :: tag
    logical, intent(in), optional :: required
    sub => node%sub_last
    call parse_node_check (sub, tag, required)
  end function parse_node_get_last_sub_ptr
  
\eatline
\eatline
\eatline
\eatline
\nwendcode{}\nwbegincode{335}\moddef{Parser: procedures}\plusendmoddef
  subroutine parse_node_undefined (node, obj)
    type(parse_node_t), intent(in) :: node
    character(*), intent(in) :: obj
    call parse_node_write (node, 6)
    call msg_bug (" Parse-tree node: " // obj // " requested, but undefined")
  end subroutine parse_node_undefined

\eatline
\nwendcode{}\nwbegindocs{336}Check if a parse node has a particular tag, and if it is associated:
\nwenddocs{}\nwbegincode{337}\moddef{Parser: public}\plusendmoddef
  public :: parse_node_check
\nwendcode{}\nwbegincode{338}\moddef{Parser: procedures}\plusendmoddef
  subroutine parse_node_check (node, tag, required)
    type(parse_node_t), pointer :: node
    character(*), intent(in), optional :: tag
    logical, intent(in), optional :: required
    if (associated (node)) then
       if (present (tag)) then
          if (parse_node_get_rule_key (node) /= tag) &
               call parse_node_mismatch (tag, node)
       end if
    else
       if (present (required)) then
          if (required) &
               call msg_bug (" Missing node, expected <" // tag // ">")
       end if
    end if
  end subroutine parse_node_check

\eatline
\nwendcode{}\nwbegindocs{339}This is called by a parse-tree scanner if the expected and the
actual nodes do not match
\nwenddocs{}\nwbegincode{340}\moddef{Parser: public}\plusendmoddef
  public :: parse_node_mismatch
\nwendcode{}\nwbegincode{341}\moddef{Parser: procedures}\plusendmoddef
  subroutine parse_node_mismatch (string, parse_node)
    character(*), intent(in) :: string
    type(parse_node_t), intent(in) :: parse_node
    call parse_node_write (parse_node)
    call msg_bug (" Syntax mismatch, expected <" // string // ">.")
  end subroutine parse_node_mismatch

\eatline
\nwendcode{}\nwbegindocs{342}The following functions are wrappers for extracting the token contents.
\nwenddocs{}\nwbegincode{343}\moddef{Parser: public}\plusendmoddef
  public :: parse_node_get_logical
  public :: parse_node_get_integer
  public :: parse_node_get_real
  public :: parse_node_get_cmplx
  public :: parse_node_get_string
  public :: parse_node_get_key
  public :: parse_node_get_rule_key
\nwendcode{}\nwbegincode{344}\moddef{Parser: procedures}\plusendmoddef
  function parse_node_get_logical (node) result (lval)
    logical :: lval
    type(parse_node_t), intent(in), target :: node
    lval = token_get_logical (parse_node_get_token_ptr (node))
  end function parse_node_get_logical

  function parse_node_get_integer (node) result (ival)
    integer :: ival
    type(parse_node_t), intent(in), target :: node
    ival = token_get_integer (parse_node_get_token_ptr (node))
  end function parse_node_get_integer

  function parse_node_get_real (node) result (rval)
    real(default) :: rval
    type(parse_node_t), intent(in), target :: node
    rval = token_get_real (parse_node_get_token_ptr (node))
  end function parse_node_get_real
  
  function parse_node_get_cmplx (node) result (cval)
    complex(default) :: cval
    type(parse_node_t), intent(in), target :: node
    cval = token_get_cmplx (parse_node_get_token_ptr (node))
  end function parse_node_get_cmplx

  function parse_node_get_string (node) result (sval)
    type(string_t) :: sval
    type(parse_node_t), intent(in), target :: node
    sval = token_get_string (parse_node_get_token_ptr (node))
  end function parse_node_get_string

  function parse_node_get_key (node) result (kval)
    type(string_t) :: kval
    type(parse_node_t), intent(in), target :: node
    kval = token_get_key (parse_node_get_token_ptr (node))
  end function parse_node_get_key

  function parse_node_get_rule_key (node) result (kval)
    type(string_t) :: kval
    type(parse_node_t), intent(in), target :: node
    kval = syntax_rule_get_key (parse_node_get_rule_ptr (node))
  end function parse_node_get_rule_key

  function parse_node_get_token_ptr (node) result (token)
    type(token_t), pointer :: token
    type(parse_node_t), intent(in), target :: node
    if (token_is_valid (node%token)) then
       token => node%token
    else
       call parse_node_undefined (node, "token")
    end if
  end function parse_node_get_token_ptr

\eatline
\eatline
\eatline
\nwendcode{}\nwbegindocs{345}Return a MD5 sum for a parse node.  The method is to write the node
to a scratch file and to evaluate the MD5 sum of that file.
\nwenddocs{}\nwbegincode{346}\moddef{Parser: public}\plusendmoddef
  public :: parse_node_get_md5sum
\nwendcode{}\nwbegincode{347}\moddef{Parser: procedures}\plusendmoddef
  function parse_node_get_md5sum (pn) result (md5sum_pn)
    character(32) :: md5sum_pn
    type(parse_node_t), intent(in) :: pn
    integer :: u
    u = free_unit ()
    open (unit = u, status = "scratch", action = "readwrite")
    call parse_node_write_rec (pn, unit=u)
    rewind (u)
    md5sum_pn = md5sum (u)
    close (u)
  end function parse_node_get_md5sum

\eatline
\nwendcode{}\nwbegindocs{348}\nwdocspar
\subsection{The parse tree}
The parse tree is a tree of nodes, where leaf nodes hold a valid
token, while branch nodes point to a list of sub-nodes.
\nwenddocs{}\nwbegincode{349}\moddef{Parser: public}\plusendmoddef
  public :: parse_tree_t
\nwendcode{}\nwbegincode{350}\moddef{Parser: types}\plusendmoddef
  type :: parse_tree_t
     private
     type(parse_node_t), pointer :: root_node => null ()
  end type parse_tree_t

\eatline
\nwendcode{}\nwbegindocs{351}The parser.  Its arguments are the parse tree (which should be empty
initially), the lexer (which should be already set up), the syntax
table (which should be valid), and the input stream.  The input stream
is completely parsed, using the lexer setup and the syntax rules as
given, and the parse tree is built accordingly.

If {\tt{}check{\char95}eof} is absent or true, the parser will complain about
trailing garbage.  Otherwise, it will ignore it.

By default, the input stream is matched against the top rule in the
specified syntax.  If {\tt{}key} is given, it is matched against the rule
with the specified key instead.

Failure at the top level means that no rule could match at all; in
this case the error message will show the top rule.
\nwenddocs{}\nwbegincode{352}\moddef{Parser: public}\plusendmoddef
  public :: parse_tree_init
\nwendcode{}\nwbegincode{353}\moddef{Parser: procedures}\plusendmoddef
  subroutine parse_tree_init &
       (parse_tree, syntax, lexer, key, check_eof)
    type(parse_tree_t), intent(inout) :: parse_tree
    type(lexer_t), intent(inout) :: lexer
    type(syntax_t), intent(in), target :: syntax
    type(string_t), intent(in), optional :: key
    logical, intent(in), optional :: check_eof
    type(syntax_rule_t), pointer :: rule
    type(lexeme_t) :: lexeme
    type(parse_node_t), pointer :: node
    logical :: ok, check
    check = .true.;  if (present (check_eof)) check = check_eof
    call lexer_clear (lexer)
    if (present (key)) then
       rule => syntax_get_rule_ptr (syntax, key)
    else
       rule => syntax_get_top_rule_ptr (syntax)
    end if
    if (associated (rule)) then
       call parse_node_match_rule (node, rule, ok)
       if (ok) then
          parse_tree%root_node => node
       else
          call parse_error (rule, lexeme)
       end if
       if (check) then
          call lex (lexeme, lexer)
          if (.not. lexeme_is_eof (lexeme)) then
             call lexer_show_location (lexer)
             call msg_fatal (" Syntax error " &
                  // "(at or before the location indicated above)")
          end if
       end if
    else
       call msg_bug (" Parser failed because syntax is empty")
    end if
  contains
  \LA{}Parser: internal subroutines of \code{}parse{\char95}tree{\char95}init\edoc{}\RA{}
  end subroutine parse_tree_init

\eatline
\nwendcode{}\nwbegindocs{354}The parser works recursively, following the rule tree, building the
tree of nodes on the fly.  If the given rule matches, the node is
filled on return.  If not, the node remains empty.
\nwenddocs{}\nwbegincode{355}\moddef{Parser: internal subroutines of \code{}parse{\char95}tree{\char95}init\edoc{}}\endmoddef
  recursive subroutine parse_node_match_rule (node, rule, ok)
    type(parse_node_t), pointer :: node
    type(syntax_rule_t), intent(in), target :: rule
    logical, intent(out) :: ok
    logical, parameter :: debug = .false.
    integer :: type
    if (debug)  write (6, "(A)", advance="no") "Parsing rule: "
    if (debug)  call syntax_rule_write (rule, 6)
    node => null ()
    type = syntax_rule_get_type (rule)
    if (syntax_rule_is_atomic (rule)) then
       call lex (lexeme, lexer)
       if (debug)  write (6, "(A)", advance="no") "Token: "
       if (debug)  call lexeme_write (lexeme, 6)
       call parse_node_create_leaf (node, rule, lexeme)
       ok = associated (node)
       if (.not. ok)  call lexer_put_back (lexer, lexeme)
    else
       select case (type)
       case (S_ALTERNATIVE);  call parse_alternative (node, rule, ok)
       case (S_GROUP);        call parse_group (node, rule, ok)
       case (S_SEQUENCE);     call parse_sequence (node, rule, .false., ok)
       case (S_LIST);         call parse_sequence (node, rule, .true., ok)
       case (S_ARGS);         call parse_args (node, rule, ok)
       case (S_IGNORE);       call parse_ignore (node, ok)
       end select
    end if
    if (debug) then
       if (ok) then
          write (6, "(A)", advance="no") "Matched rule: "
       else
          write (6, "(A)", advance="no") "Failed rule: "
       end if
       call syntax_rule_write (rule)
       if (associated (node)) call parse_node_write (node)
    end if
  end subroutine parse_node_match_rule
\eatline
\nwendcode{}\nwbegindocs{356}Parse an alternative: try each case.  If the match succeeds, the
node has been filled, so return.  If nothing works, return failure.
\nwenddocs{}\nwbegincode{357}\moddef{Parser: internal subroutines of \code{}parse{\char95}tree{\char95}init\edoc{}}\plusendmoddef
  recursive subroutine parse_alternative (node, rule, ok)
    type(parse_node_t), pointer :: node
    type(syntax_rule_t), intent(in), target :: rule
    logical, intent(out) :: ok
    integer :: i
    do i = 1, syntax_rule_get_n_sub (rule)
       call parse_node_match_rule (node, syntax_rule_get_sub_ptr (rule, i), ok)
       if (ok)  return
    end do
    ok = .false.
  end subroutine parse_alternative
\eatline
\nwendcode{}\nwbegindocs{358}Parse a group: the first and third lexemes have to be the
delimiters, the second one is parsed as the actual node, using now the
child rule.  If the first match fails, return with failure.  If the
other matches fail, issue an error, since we cannot lex back more than
one item.
\nwenddocs{}\nwbegincode{359}\moddef{Parser: internal subroutines of \code{}parse{\char95}tree{\char95}init\edoc{}}\plusendmoddef
  recursive subroutine parse_group (node, rule, ok)
    type(parse_node_t), pointer :: node
    type(syntax_rule_t), intent(in), target :: rule
    logical, intent(out) :: ok
    type(string_t), dimension(2) :: delimiter
    delimiter = syntax_rule_get_delimiter (rule)
    call lex (lexeme, lexer)
    if (lexeme_get_string (lexeme) == delimiter(1)) then
       call parse_node_match_rule (node, syntax_rule_get_sub_ptr (rule, 1), ok)
       if (ok) then
          call lex (lexeme, lexer)
          if (lexeme_get_string (lexeme) == delimiter(2)) then
             ok = .true.
          else
             call parse_error (rule, lexeme)
          end if
       else
          call parse_error (rule, lexeme)
       end if
    else
       call lexer_put_back (lexer, lexeme)
       ok = .false.
    end if
  end subroutine parse_group
\eatline
\nwendcode{}\nwbegindocs{360}Parsing a sequence.  The last rule element may be special: optional
and/or repetitive.  Each sub-node that matches is appended to the
sub-node list of the parent node.

If {\tt{}sep} is true, we look for a separator after each element.
\nwenddocs{}\nwbegincode{361}\moddef{Parser: internal subroutines of \code{}parse{\char95}tree{\char95}init\edoc{}}\plusendmoddef
  recursive subroutine parse_sequence (node, rule, sep, ok)
    type(parse_node_t), pointer :: node
    type(syntax_rule_t), intent(in), target :: rule
    logical, intent(in) :: sep
    logical, intent(out) :: ok
    type(parse_node_t), pointer :: current
    integer :: i, n
    logical :: opt, rep, cont
    type(string_t) :: separator
    call parse_node_create_branch (node, rule)
    if (sep)  separator = syntax_rule_get_separator (rule)
    n = syntax_rule_get_n_sub (rule)
    opt = syntax_rule_last_optional (rule)
    rep = syntax_rule_last_repetitive (rule)
    ok = .true.
    cont = .true.
    SCAN_RULE: do i = 1, n
       call parse_node_match_rule &
            (current, syntax_rule_get_sub_ptr (rule, i), cont)
       if (cont) then
          call parse_node_append_sub (node, current)
          if (sep .and. (i<n .or. rep)) then
             call lex (lexeme, lexer)
             if (lexeme_get_string (lexeme) /= separator) then
                call lexer_put_back (lexer, lexeme)
                cont = .false.
                exit SCAN_RULE
             end if
          end if
       else
          if (i == n .and. opt) then
             exit SCAN_RULE
          else if (i == 1) then
             ok = .false.
             exit SCAN_RULE
          else             
             call parse_error (rule, lexeme)
          end if
       end if
    end do SCAN_RULE
    if (rep) then
       do while (cont)
          call parse_node_match_rule &
               (current, syntax_rule_get_sub_ptr (rule, n), cont)
          if (cont) then
             call parse_node_append_sub (node, current)
             if (sep) then
                call lex (lexeme, lexer)
                if (lexeme_get_string (lexeme) /= separator) then
                   call lexer_put_back (lexer, lexeme)
                   cont = .false.
                end if
             end if
          else
             if (sep)  call parse_error (rule, lexeme)
          end if
       end do
    end if
    call parse_node_freeze_branch (node)
  end subroutine parse_sequence
\eatline
\nwendcode{}\nwbegindocs{362}Argument list: We use the {\tt{}parse{\char95}group} code, but call
{\tt{}parse{\char95}sequence} inside.
\nwenddocs{}\nwbegincode{363}\moddef{Parser: internal subroutines of \code{}parse{\char95}tree{\char95}init\edoc{}}\plusendmoddef
  recursive subroutine parse_args (node, rule, ok)
    type(parse_node_t), pointer :: node
    type(syntax_rule_t), intent(in), target :: rule
    logical, intent(out) :: ok
    type(string_t), dimension(2) :: delimiter
    delimiter = syntax_rule_get_delimiter (rule)
    call lex (lexeme, lexer)
    if (lexeme_get_string (lexeme) == delimiter(1)) then
       call parse_sequence (node, rule, .true., ok)
       if (ok) then
          call lex (lexeme, lexer)
          if (lexeme_get_string (lexeme) == delimiter(2)) then
             ok = .true.
          else
             call parse_error (rule, lexeme)
          end if
       else
          call parse_error (rule, lexeme)
       end if
    else
       call lexer_put_back (lexer, lexeme)
       ok = .false.
    end if
  end subroutine parse_args
\eatline
\nwendcode{}\nwbegindocs{364}The IGNORE syntax reads one lexeme and discards it if it is numeric,
logical or string/identifier (but not a keyword).  This is a
successful match.  Otherwise, the match fails.  The node
pointer is returned disassociated in any case.
\nwenddocs{}\nwbegincode{365}\moddef{Parser: internal subroutines of \code{}parse{\char95}tree{\char95}init\edoc{}}\plusendmoddef
  subroutine parse_ignore (node, ok)
    type(parse_node_t), pointer :: node
    logical, intent(out) :: ok
    call lex (lexeme, lexer)
    select case (lexeme_get_type (lexeme))
    case (T_NUMERIC, T_IDENTIFIER, T_QUOTED)
       ok = .true.
    case default
       ok = .false.
    end select
    node => null ()
  end subroutine parse_ignore
\eatline
\nwendcode{}\nwbegindocs{366}If the match fails and we cannot step back:
\nwenddocs{}\nwbegincode{367}\moddef{Parser: internal subroutines of \code{}parse{\char95}tree{\char95}init\edoc{}}\plusendmoddef
  subroutine parse_error (rule, lexeme)
    type(syntax_rule_t), intent(in) :: rule
    type(lexeme_t), intent(in) :: lexeme
    character(80) :: buffer
    integer :: u, iostat
    call lexer_show_location (lexer)
    u = free_unit ()
    open (u, status = "scratch")
    write (u, "(A)", advance="no")  "Expected syntax:"
    call syntax_rule_write (rule, u)
    write (u, "(A)", advance="no")  "Found token:"
    call lexeme_write (lexeme, u)
    rewind (u)
    do
       read (u, "(A)", iostat=iostat)  buffer
       if (iostat /= 0)  exit
       call msg_message (trim (buffer))
    end do
    call msg_fatal (" Syntax error " &
         // "(at or before the location indicated above)")
  end subroutine parse_error
\eatline
\nwendcode{}\nwbegindocs{368}The finalizer recursively deallocates all nodes and their contents.
For each node, {\tt{}parse{\char95}node{\char95}final} is called on the sub-nodes, which in
turn deallocates the token or sub-node array contained within each of
them.  At the end, only the top node remains to be deallocated.
\nwenddocs{}\nwbegincode{369}\moddef{Parser: public}\plusendmoddef
  public :: parse_tree_final
\nwendcode{}\nwbegincode{370}\moddef{Parser: procedures}\plusendmoddef
  subroutine parse_tree_final (parse_tree)
    type(parse_tree_t), intent(inout) :: parse_tree
    if (associated (parse_tree%root_node)) then
       call parse_node_final (parse_tree%root_node)
       deallocate (parse_tree%root_node)
    end if
  end subroutine parse_tree_final

\eatline
\nwendcode{}\nwbegindocs{371}Print the parse tree.  Print one token per line, indented according
to the depth of the node.

The {\tt{}verbose} version includes type identifiers for the nodes.
\nwenddocs{}\nwbegincode{372}\moddef{Parser: public}\plusendmoddef
  public :: parse_tree_write
\nwendcode{}\nwbegincode{373}\moddef{Parser: procedures}\plusendmoddef
  subroutine parse_tree_write (parse_tree, unit, verbose)
    type(parse_tree_t), intent(in) :: parse_tree
    integer, intent(in), optional :: unit
    logical, intent(in), optional :: verbose
    integer :: u
    logical :: short
    u = given_output_unit (unit);  if (u < 0)  return
    short = .true.;  if (present (verbose))  short = .not. verbose
    write (u, "(A)") "Parse tree:"
    if (associated (parse_tree%root_node)) then
       call parse_node_write_rec (parse_tree%root_node, unit, short, 1)
    else
       write (u, *) "[empty]"
    end if
  end subroutine parse_tree_write

\eatline
\nwendcode{}\nwbegindocs{374}This is a generic error that can be issued if the parse tree does
not meet the expectaions of the parser.  This most certainly indicates
a bug.
\nwenddocs{}\nwbegincode{375}\moddef{Parser: public}\plusendmoddef
  public :: parse_tree_bug
\nwendcode{}\nwbegincode{376}\moddef{Parser: procedures}\plusendmoddef
  subroutine parse_tree_bug (node, keys)
    type(parse_node_t), intent(in) :: node
    character(*), intent(in) :: keys
    call parse_node_write (node)
    call msg_bug (" Inconsistency in parse tree: expected " // keys)
  end subroutine parse_tree_bug

\eatline
\nwendcode{}\nwbegindocs{377}\nwdocspar
\subsection{Access the parser}
For scanning the parse tree we give access to the top node, as a node
pointer.  Of course, there should be no write access.
\nwenddocs{}\nwbegincode{378}\moddef{Parser: public}\plusendmoddef
  public :: parse_tree_get_root_ptr
\nwendcode{}\nwbegincode{379}\moddef{Parser: procedures}\plusendmoddef
  function parse_tree_get_root_ptr (parse_tree) result (node)
    type(parse_node_t), pointer :: node
    type(parse_tree_t), intent(in), target :: parse_tree
    node => parse_tree%root_node
  end function parse_tree_get_root_ptr

\eatline
\nwendcode{}\nwbegindocs{380}\nwdocspar
\subsection{Tools}
This operation traverses the parse tree and simplifies any occurences of a 
set of syntax rules.  If such a parse node has only one sub-node, it is
replaced by that subnode.  (This makes sense only of the rules to eliminate
have no meaningful token.)
\nwenddocs{}\nwbegincode{381}\moddef{Parser: public}\plusendmoddef
  public :: parse_tree_reduce
\nwendcode{}\nwbegincode{382}\moddef{Parser: procedures}\plusendmoddef
  subroutine parse_tree_reduce (parse_tree, rule_key)
    type(parse_tree_t), intent(inout) :: parse_tree
    type(string_t), dimension(:), intent(in) :: rule_key
    type(parse_node_t), pointer :: pn
    pn => parse_tree%root_node
    if (associated (pn)) then
       call parse_node_reduce (pn, null(), null())
    end if
  contains
    recursive subroutine parse_node_reduce (pn, pn_prev, pn_parent)
      type(parse_node_t), intent(inout), pointer :: pn
      type(parse_node_t), intent(in), pointer :: pn_prev, pn_parent
      type(parse_node_t), pointer :: pn_sub, pn_sub_prev, pn_tmp
      pn_sub_prev => null ()
      pn_sub => pn%sub_first
      do while (associated (pn_sub))
         call parse_node_reduce (pn_sub, pn_sub_prev, pn)
         pn_sub_prev => pn_sub
         pn_sub => pn_sub%next
      end do
      if (parse_node_get_n_sub (pn) == 1) then
         if (matches (parse_node_get_rule_key (pn), rule_key)) then
            pn_tmp => pn
            pn => pn%sub_first
            if (associated (pn_prev)) then
               pn_prev%next => pn
            else if (associated (pn_parent)) then
               pn_parent%sub_first => pn
            else
               parse_tree%root_node => pn
            end if
            if (associated (pn_tmp%next)) then 
               pn%next => pn_tmp%next
            else if (associated (pn_parent)) then
               pn_parent%sub_last => pn
            end if
            call parse_node_final (pn_tmp, recursive=.false.)
            deallocate (pn_tmp)
         end if
      end if
    end subroutine parse_node_reduce
    function matches (key, key_list) result (flag)
      logical :: flag
      type(string_t), intent(in) :: key
      type(string_t), dimension(:), intent(in) :: key_list
      integer :: i
      flag = .true.
      do i = 1, size (key_list)
         if (key == key_list(i))  return
      end do
      flag = .false.
    end function matches
  end subroutine parse_tree_reduce

\eatline
\nwendcode{}\nwbegindocs{383}\nwdocspar
\subsection{Applications}
For a file of the form
\begin{verbatim}
process foo, bar
  <something>
process xyz
  <something>
\end{verbatim}
get the \verb|<something>| entry node for the first matching process
tag.  If no matching entry is found, the node pointer remains
unassociated.
\nwenddocs{}\nwbegincode{384}\moddef{Parser: public}\plusendmoddef
  public :: parse_tree_get_process_ptr
\nwendcode{}\nwbegincode{385}\moddef{Parser: procedures}\plusendmoddef
  function parse_tree_get_process_ptr (parse_tree, process) result (node)
    type(parse_node_t), pointer :: node
    type(parse_tree_t), intent(in), target :: parse_tree
    type(string_t), intent(in) :: process
    type(parse_node_t), pointer :: node_root, node_process_def
    type(parse_node_t), pointer :: node_process_phs, node_process_list
    integer :: j
    node_root => parse_tree_get_root_ptr (parse_tree)
    if (associated (node_root)) then
       node_process_phs => parse_node_get_sub_ptr (node_root)
       SCAN_FILE: do while (associated (node_process_phs))
          node_process_def => parse_node_get_sub_ptr (node_process_phs)
          node_process_list => parse_node_get_sub_ptr (node_process_def, 2)
          do j = 1, parse_node_get_n_sub (node_process_list)
             if (parse_node_get_string &
                  (parse_node_get_sub_ptr (node_process_list, j)) &
                  == process) then
                node => parse_node_get_next_ptr (node_process_def)
                return
             end if
          end do
          node_process_phs => parse_node_get_next_ptr (node_process_phs)
       end do SCAN_FILE
       node => null ()
    else
       node => null ()
    end if
  end function parse_tree_get_process_ptr

\eatline
\nwendcode{}\nwbegindocs{386}\nwdocspar
\subsection{Unit tests}
\nwenddocs{}\nwbegincode{387}\moddef{Parser: public}\plusendmoddef
  public :: parse_test
\nwendcode{}\nwbegincode{388}\moddef{Parser: procedures}\plusendmoddef
  subroutine parse_test (u, results)
    integer, intent(in) :: u
    type(test_results_t), intent(inout) :: results
  \LA{}Parser: execute tests\RA{}  
  end subroutine parse_test
  
\nwendcode{}\nwbegindocs{389} %def parse_test
\nwenddocs{}\nwbegindocs{390}This checks the parser.
\nwenddocs{}\nwbegincode{391}\moddef{Parser: execute tests}\endmoddef
  call test (parse_1, "parse_1", &
       "check the parser", &
       u, results)
\nwendcode{}\nwbegincode{392}\moddef{Parser: tests}\endmoddef
  subroutine parse_1 (u)
    use ifiles
    use lexers
    integer, intent(in) :: u    

    type(ifile_t) :: ifile
    type(syntax_t), target :: syntax
    type(lexer_t) :: lexer
    type(stream_t), target :: stream
    type(parse_tree_t), target :: parse_tree
    
    write (u, "(A)")  "* Test output: Parsing"
    write (u, "(A)")  "*   Purpose: test parse routines"
    write (u, "(A)")      
    
    call ifile_append (ifile, "SEQ expr = term addition*")
    call ifile_append (ifile, "SEQ addition = plus_or_minus term")
    call ifile_append (ifile, "SEQ term = factor multiplication*")
    call ifile_append (ifile, "SEQ multiplication = times_or_over factor")
    call ifile_append (ifile, "SEQ factor = atom exponentiation*")
    call ifile_append (ifile, "SEQ exponentiation = '^' atom")
    call ifile_append (ifile, "ALT atom = real | delimited_expr")
    call ifile_append (ifile, "GRO delimited_expr = ( expr )")
    call ifile_append (ifile, "ALT plus_or_minus = '+' | '-'")
    call ifile_append (ifile, "ALT times_or_over = '*' | '/'")
    call ifile_append (ifile, "KEY '+'")
    call ifile_append (ifile, "KEY '-'")
    call ifile_append (ifile, "KEY '*'")
    call ifile_append (ifile, "KEY '/'")
    call ifile_append (ifile, "KEY '^'")
    call ifile_append (ifile, "REA real")
    
    write (u, "(A)")  "* File contents (syntax definition):"
    call ifile_write (ifile, u)
    write (u, "(A)")  "EOF"
    write (u, "(A)")
    
    call syntax_init (syntax, ifile)
    call ifile_final (ifile)
    call syntax_write (syntax, u)
    write (u, "(A)")
    
    call lexer_init (lexer, &
         comment_chars = "", &
         quote_chars = "'", &
         quote_match = "'", &
         single_chars = "+-*/^()", &
         special_class = [""] , &
         keyword_list = syntax_get_keyword_list_ptr (syntax))
    call lexer_write_setup (lexer, u)
    write (u, "(A)")
    
    call ifile_append (ifile, "(27+8^3-2/3)*(4+7)^2*99")
    write (u, "(A)")  "* File contents (input file):"
    call ifile_write (ifile, u)
    write (u, "(A)")  "EOF"
    print *
    
    call stream_init (stream, ifile)
    call lexer_assign_stream (lexer, stream)
    call parse_tree_init (parse_tree, syntax, lexer)
    call stream_final (stream)
    call parse_tree_write (parse_tree, u, .true.)
    print *
    
    write (u, "(A)")  "* Cleanup, everything should now be empty:"
    write (u, "(A)")
    
    call parse_tree_final (parse_tree)
    call parse_tree_write (parse_tree, u, .true.)
    write (u, "(A)")
    
    call lexer_final (lexer)
    call lexer_write_setup (lexer, u)
    write (u, "(A)")
    
    call ifile_final (ifile)
    write (u, "(A)")  "* File contents:"
    call ifile_write (ifile, u)
    write (u, "(A)")  "EOF"
    write (u, "(A)")
    
    call syntax_final (syntax)
    call syntax_write (syntax, u)

    write (u, "(A)")
    write (u, "(A)")  "* Test output end: parser_1"    
    
  end subroutine parse_1
\eatline
\nwendcode{}\nwbegindocs{393}\nwdocspar
\clearpage
%%%%%%%%%%%%%%%%%%%%%%%%%%%%%%%%%%%%%%%%%%%%%%%%%%%%%%%%%%%%%%%%%%%%%%%%
\section{XML Parser}
The XML parser is actually independent from the previous modules of
lexer and parser.  However, for a generic I/O interface we make use of
the {\tt{}stream{\char95}t} facility.

We need the XML parser for reading and writing LHEF data files.  Only
a subset of XML is actually relevant.  The parser is of the ``pull''
type, i.e., the program steers the reading of XML data in a
context-sensitive manner.
\nwenddocs{}\nwbegincode{394}\moddef{\code{}xml.f90\edoc{}}\endmoddef
\LA{}File header\RA{}

module xml

\LA{}Use strings\RA{}
  use io_units
  use unit_tests
  use system_defs, only: BLANK, TAB
  use diagnostics
  use ifiles
  use lexers

\LA{}Standard module head\RA{}

\LA{}XML: public\RA{}

\LA{}XML: types\RA{}

contains

\LA{}XML: procedures\RA{}

\LA{}XML: tests\RA{}

end module xml
\eatline
\nwendcode{}\nwbegindocs{395}\nwdocspar
\subsection{Cached Stream}
The stream type as defined in the {\tt{}lexer} module is versatile
regarding the choice of input channel, but it does not allow reading a
section more than once.  Here, we define an extension where we can
return a string to the stream, which is stored in a cache variable,
and presented to the caller for the next read.
\nwenddocs{}\nwbegincode{396}\moddef{XML: public}\endmoddef
  public :: cstream_t
\nwendcode{}\nwbegincode{397}\moddef{XML: types}\endmoddef
  type, extends (stream_t) :: cstream_t
     logical :: cache_is_empty = .true.
     type(string_t) :: cache
   contains
   \LA{}XML: cstream: TBP\RA{}
  end type cstream_t
  
\eatline
\nwendcode{}\nwbegindocs{398}The initializers are simply inherited.  We require a generic.
\nwenddocs{}\nwbegincode{399}\moddef{XML: cstream: TBP}\endmoddef
  generic :: init => init_filename, init_unit, init_string, &
       init_ifile, init_line
  procedure :: init_filename => cstream_init_filename
  procedure :: init_unit => cstream_init_unit
  procedure :: init_string => cstream_init_string
  procedure :: init_ifile => cstream_init_ifile
  procedure :: init_line => cstream_init_line
\nwendcode{}\nwbegincode{400}\moddef{XML: procedures}\endmoddef
  subroutine cstream_init_filename (cstream, filename)
    class(cstream_t), intent(out) :: cstream
    character(*), intent(in) :: filename
    call stream_init (cstream%stream_t, filename)
  end subroutine cstream_init_filename

  subroutine cstream_init_unit (cstream, unit)
    class(cstream_t), intent(out) :: cstream
    integer, intent(in) :: unit
    call stream_init (cstream%stream_t, unit)
  end subroutine cstream_init_unit
    
  subroutine cstream_init_string (cstream, string)
    class(cstream_t), intent(out) :: cstream
    type(string_t), intent(in) :: string
    call stream_init (cstream%stream_t, string)
  end subroutine cstream_init_string
    
  subroutine cstream_init_ifile (cstream, ifile)
    class(cstream_t), intent(out) :: cstream
    type(ifile_t), intent(in) :: ifile
    call stream_init (cstream%stream_t, ifile)
  end subroutine cstream_init_ifile
    
  subroutine cstream_init_line (cstream, line)
    class(cstream_t), intent(out) :: cstream
    type(line_p), intent(in) :: line
    call stream_init (cstream%stream_t, line)
  end subroutine cstream_init_line
    
\eatline
\nwendcode{}\nwbegindocs{401}Finalizer: also inherited
\nwenddocs{}\nwbegincode{402}\moddef{XML: cstream: TBP}\plusendmoddef
  procedure :: final => cstream_final
\nwendcode{}\nwbegincode{403}\moddef{XML: procedures}\plusendmoddef
  subroutine cstream_final (cstream)
    class(cstream_t), intent(inout) :: cstream
    cstream%cache_is_empty = .true.
    call stream_final (cstream%stream_t)
  end subroutine cstream_final
  
\eatline
\nwendcode{}\nwbegindocs{404}Get record: now, if there is a cache string, return this instead of
the record from the stream.
\nwenddocs{}\nwbegincode{405}\moddef{XML: cstream: TBP}\plusendmoddef
  procedure :: get_record => cstream_get_record
\nwendcode{}\nwbegincode{406}\moddef{XML: procedures}\plusendmoddef
  subroutine cstream_get_record (cstream, string, iostat)
    class(cstream_t), intent(inout) :: cstream
    type(string_t), intent(out) :: string
    integer, intent(out) :: iostat
    if (cstream%cache_is_empty) then
       call stream_get_record (cstream%stream_t, string, iostat)
    else
       string = cstream%cache
       cstream%cache_is_empty = .true.
       iostat = 0
    end if
  end subroutine cstream_get_record
  
\eatline
\nwendcode{}\nwbegindocs{407}Revert: return the (partially read) record to the stream, putting it
in the cache.
\nwenddocs{}\nwbegincode{408}\moddef{XML: cstream: TBP}\plusendmoddef
  procedure :: revert_record => cstream_revert_record
\nwendcode{}\nwbegincode{409}\moddef{XML: procedures}\plusendmoddef
  subroutine cstream_revert_record (cstream, string)
    class(cstream_t), intent(inout) :: cstream
    type(string_t), intent(in) :: string
    if (cstream%cache_is_empty) then
       cstream%cache = string
       cstream%cache_is_empty = .false.
    else
       call msg_bug ("CStream: attempt to revert twice")
    end if
  end subroutine cstream_revert_record
  
\eatline
\nwendcode{}\nwbegindocs{410}\nwdocspar
\subsection{Attributes}
A tag attribute has a name and a value; both are strings.  When the
attribute is defined, the {\tt{}known} flag indicates this.
\nwenddocs{}\nwbegincode{411}\moddef{XML: types}\plusendmoddef
  type :: attribute_t
     type(string_t) :: name
     type(string_t) :: value
     logical :: known = .false.
   contains
   \LA{}XML: attribute: TBP\RA{}     
  end type attribute_t
  
\eatline
\nwendcode{}\nwbegindocs{412}Output in standard format, non-advancing.  (If the value is unknown,
we indicate this by a question mark, which is non-standard.)
\nwenddocs{}\nwbegincode{413}\moddef{XML: attribute: TBP}\endmoddef
  procedure :: write => attribute_write
\nwendcode{}\nwbegincode{414}\moddef{XML: procedures}\plusendmoddef
  subroutine attribute_write (object, unit)
    class(attribute_t), intent(in) :: object
    integer, intent(in), optional :: unit
    integer :: u
    u = given_output_unit (unit)
    write (u, "(A,'=')", advance = "no")  char (object%name)
    if (object%known) then
       write (u, "(A,A,A)", advance = "no")  '"', char (object%value), '"'
    else
       write (u, "('?')", advance = "no")
    end if
  end subroutine attribute_write
    
\eatline
\nwendcode{}\nwbegindocs{415}This is a genuine constructor.  The value is optional.
\nwenddocs{}\nwbegincode{416}\moddef{XML: public}\plusendmoddef
  public :: xml_attribute
\nwendcode{}\nwbegincode{417}\moddef{XML: procedures}\plusendmoddef
  function xml_attribute (name, value) result (attribute)
    type(string_t), intent(in) :: name
    type(string_t), intent(in), optional :: value
    type(attribute_t) :: attribute
    attribute%name = name
    if (present (value)) then
       attribute%value = value
       attribute%known = .true.
    else
       attribute%known = .false.
    end if
  end function xml_attribute

\eatline
\nwendcode{}\nwbegindocs{418}Set a value explicitly.
\nwenddocs{}\nwbegincode{419}\moddef{XML: attribute: TBP}\plusendmoddef
  procedure :: set_value => attribute_set_value
\nwendcode{}\nwbegincode{420}\moddef{XML: procedures}\plusendmoddef
  subroutine attribute_set_value (attribute, value)
    class(attribute_t), intent(inout) :: attribute
    type(string_t), intent(in) :: value
    attribute%value = value
    attribute%known = .true.
  end subroutine attribute_set_value
  
\eatline
\nwendcode{}\nwbegindocs{421}Extract a value.  If unset, return {\tt{}"?"}
\nwenddocs{}\nwbegincode{422}\moddef{XML: attribute: TBP}\plusendmoddef
  procedure :: get_value => attribute_get_value
\nwendcode{}\nwbegincode{423}\moddef{XML: procedures}\plusendmoddef
  function attribute_get_value (attribute) result (value)
    class(attribute_t), intent(in) :: attribute
    type(string_t) :: value
    if (attribute%known) then
       value = attribute%value
    else
       value = "?"
    end if
  end function attribute_get_value
  
\eatline
\nwendcode{}\nwbegindocs{424}\nwdocspar
\subsection{The Tag Type}
The basic entity is the internal representation of an XML tag.  The
tag has a name, a well-defined set of attributes which may be
mandatory or optional, and a flag that determines whether there is
content or not.  The content itself is not stored in the data structure.
\nwenddocs{}\nwbegincode{425}\moddef{XML: public}\plusendmoddef
  public :: xml_tag_t
\nwendcode{}\nwbegincode{426}\moddef{XML: types}\plusendmoddef
  type :: xml_tag_t
     type(string_t) :: name
     type(attribute_t), dimension(:), allocatable :: attribute
     logical :: has_content = .false.
   contains
   \LA{}XML: tag: TBP\RA{}
  end type xml_tag_t
  
\eatline
\nwendcode{}\nwbegindocs{427}Initialization.  There are different versions, depending on content.
\nwenddocs{}\nwbegincode{428}\moddef{XML: tag: TBP}\endmoddef
  generic :: init => init_no_attributes
  procedure :: init_no_attributes => tag_init_no_attributes
\nwendcode{}\nwbegincode{429}\moddef{XML: procedures}\plusendmoddef
  subroutine tag_init_no_attributes (tag, name, has_content)
    class(xml_tag_t), intent(out) :: tag
    type(string_t), intent(in) :: name
    logical, intent(in), optional :: has_content
    tag%name = name
    allocate (tag%attribute (0))
    if (present (has_content))  tag%has_content = has_content
  end subroutine tag_init_no_attributes
  
\eatline
\nwendcode{}\nwbegindocs{430}This version sets attributes.
\nwenddocs{}\nwbegincode{431}\moddef{XML: tag: TBP}\plusendmoddef
  generic :: init => init_with_attributes
  procedure :: init_with_attributes => tag_init_with_attributes
\nwendcode{}\nwbegincode{432}\moddef{XML: procedures}\plusendmoddef
  subroutine tag_init_with_attributes (tag, name, attribute, has_content)
    class(xml_tag_t), intent(out) :: tag
    type(string_t), intent(in) :: name
    type(attribute_t), dimension(:), intent(in) :: attribute
    logical, intent(in), optional :: has_content
    tag%name = name
    allocate (tag%attribute (size (attribute)))
    tag%attribute = attribute
    if (present (has_content))  tag%has_content = has_content
  end subroutine tag_init_with_attributes
  
\eatline
\nwendcode{}\nwbegindocs{433}Set an attribute value explicitly.
\nwenddocs{}\nwbegincode{434}\moddef{XML: tag: TBP}\plusendmoddef
  procedure :: set_attribute => tag_set_attribute
\nwendcode{}\nwbegincode{435}\moddef{XML: procedures}\plusendmoddef
  subroutine tag_set_attribute (tag, i, value)
    class(xml_tag_t), intent(inout) :: tag
    integer, intent(in) :: i
    type(string_t), intent(in) :: value
    call tag%attribute(i)%set_value (value)
  end subroutine tag_set_attribute
  
\eatline
\nwendcode{}\nwbegindocs{436}Get an attribute value.
\nwenddocs{}\nwbegincode{437}\moddef{XML: tag: TBP}\plusendmoddef
  procedure :: get_attribute => tag_get_attribute
\nwendcode{}\nwbegincode{438}\moddef{XML: procedures}\plusendmoddef
  function tag_get_attribute (tag, i) result (value)
    class(xml_tag_t), intent(in) :: tag
    integer, intent(in) :: i
    type(string_t) :: value
    value = tag%attribute(i)%get_value ()
  end function tag_get_attribute
  
\eatline
\nwendcode{}\nwbegindocs{439}Output to an I/O unit, default STDOUT.  We use non-advancing output.
\nwenddocs{}\nwbegincode{440}\moddef{XML: tag: TBP}\plusendmoddef
  generic :: write => write_without_content
  procedure :: write_without_content => tag_write
\nwendcode{}\nwbegincode{441}\moddef{XML: procedures}\plusendmoddef
  subroutine tag_write (tag, unit)
    class(xml_tag_t), intent(in) :: tag
    integer, intent(in), optional :: unit
    integer :: u, i
    u = given_output_unit (unit)
    write (u, "('<',A)", advance = "no")  char (tag%name)
    do i = 1, size (tag%attribute)
       write (u, "(1x)", advance = "no")
       call tag%attribute(i)%write (u)
    end do
    if (tag%has_content) then
       write (u, "('>')", advance = "no")
    else
       write (u, "(' />')", advance = "no")
    end if
  end subroutine tag_write
  
\eatline
\nwendcode{}\nwbegindocs{442}If there is content, we should write the context next (arbitrary
format), the write the corresponding closing tag.  Again, non-advancing.
\nwenddocs{}\nwbegincode{443}\moddef{XML: tag: TBP}\plusendmoddef
  procedure :: close => tag_close
\nwendcode{}\nwbegincode{444}\moddef{XML: procedures}\plusendmoddef
  subroutine tag_close (tag, unit)
    class(xml_tag_t), intent(in) :: tag
    integer, intent(in), optional :: unit
    integer :: u
    u = given_output_unit (unit)
    write (u, "('</',A,'>')", advance = "no")  char (tag%name)
  end subroutine tag_close
    
\eatline
\nwendcode{}\nwbegindocs{445}Given content as a single string, we can write tag, content, and
closing at once
\nwenddocs{}\nwbegincode{446}\moddef{XML: tag: TBP}\plusendmoddef
  generic :: write => write_with_content
  procedure :: write_with_content => tag_write_with_content
\nwendcode{}\nwbegincode{447}\moddef{XML: procedures}\plusendmoddef
  subroutine tag_write_with_content (tag, content, unit)
    class(xml_tag_t), intent(in) :: tag
    type(string_t), intent(in) :: content
    integer, intent(in), optional :: unit
    integer :: u
    u = given_output_unit (unit)
    call tag%write (u)
    write (u, "(A)", advance = "no")  char (content)
    call tag%close (u)
  end subroutine tag_write_with_content
  
\eatline
\nwendcode{}\nwbegindocs{448}Input from stream.  We know what we are looking for, so we check if
the name matches, then fill attributes.  We report an error if (a) an
I/O error occurs, (b) we reach EOF before encountering the tag, (c)
if the tag is incomplete.

Trailing text after reading a tag is put back to the input stream.

We assume that the tag is not broken across records, and that there is
only one tag within the record.  This is more restricted than standard
XML.
\nwenddocs{}\nwbegincode{449}\moddef{XML: tag: TBP}\plusendmoddef
  procedure :: read => tag_read
\nwendcode{}\nwbegincode{450}\moddef{XML: procedures}\plusendmoddef
  subroutine tag_read (tag, cstream, success)
    class(xml_tag_t), intent(inout) :: tag
    type(cstream_t), intent(inout) :: cstream
    logical, intent(out) :: success
    type(string_t) :: string
    integer :: iostat, p1, p2
    character(2), parameter :: WS = BLANK // TAB
    logical :: done

    ! Skip comments and blank lines
    FIND_NON_COMMENT: do
       FIND_NONEMPTY_RECORD: do
          call cstream%get_record (string, iostat)
          if (iostat /= 0)  call err_io ()
          p1 = verify (string, WS)
          if (p1 > 0)  exit FIND_NONEMPTY_RECORD
       end do FIND_NONEMPTY_RECORD

       ! Look for comment beginning
       p2 = p1 + 3
       if (extract (string, p1, p2) /= "<!--")  exit FIND_NON_COMMENT
       
       ! Look for comment end, then restart
       string = extract (string, p2 + 1)
       FIND_COMMENT_END: do
          do p1 = 1, len (string) - 2
             p2 = p1 + 2
             if (extract (string, p1, p2) == "-->") then

                ! Return trailing text to the stream
                string = extract (string, p2 + 1)
                if (string /= "")  call cstream%revert_record (string)
                exit FIND_COMMENT_END

             end if
          end do
          call cstream%get_record (string, iostat)
          if (iostat /= 0)  call err_io ()
       end do FIND_COMMENT_END
    end do FIND_NON_COMMENT

    ! Look for opening <
    p2 = p1
    if (extract (string, p1, p2) /= "<") then
       call cstream%revert_record (string)
       success = .false.;  return
    else

       ! Look for tag name
       string = extract (string, p2 + 1)
       p1 = verify (string, WS);  if (p1 == 0)  call err_incomplete ()
       p2 = p1 + len (tag%name) - 1
       if (extract (string, p1, p2) /= tag%name) then
          call cstream%revert_record ("<" // string)
          success = .false.;  return
       else

          ! Look for attributes
          string = extract (string, p2 + 1)
          READ_ATTRIBUTES: do
             call tag%read_attribute (string, done)
             if (done)  exit READ_ATTRIBUTES
          end do READ_ATTRIBUTES

          ! Look for closing >
          p1 = verify (string, WS);  if (p1 == 0)  call err_incomplete ()
          p2 = p1
          if (extract (string, p1, p1) == ">") then
             tag%has_content = .true.
          else

             ! Look for closing />
             p2 = p1 + 1
             if (extract (string, p1, p2) /= "/>")  call err_incomplete ()
          end if

          ! Return trailing text to the stream
          string = extract (string, p2 + 1) 
          if (string /= "")  call cstream%revert_record (string)
          success = .true.

       end if
    end if

  contains

    subroutine err_io ()
      select case (iostat)
      case (:-1)
         call msg_fatal ("XML: Error reading tag '" // char (tag%name) &
              // "': end of file")
      case (1:)
         call msg_fatal ("XML: Error reading tag '" // char (tag%name) &
              // "': I/O error")
      end select
      success = .false.
    end subroutine err_io
    
    subroutine err_incomplete ()
      call msg_fatal ("XML: Error reading tag '" // char (tag%name) &
           // "': tag incomplete")
      success = .false.
    end subroutine err_incomplete
    
  end subroutine tag_read

\eatline
\nwendcode{}\nwbegindocs{451}Read a single attribute.  If the attribute is valid, assign the
value.  Setting a value twice should be an error, but is not detected.
If the attribute is unknown, ignore it.  If we reach the closing sign,
report this.
\nwenddocs{}\nwbegincode{452}\moddef{XML: tag: TBP}\plusendmoddef
  procedure :: read_attribute => tag_read_attribute
\nwendcode{}\nwbegincode{453}\moddef{XML: procedures}\plusendmoddef
  subroutine tag_read_attribute (tag, string, done)
    class(xml_tag_t), intent(inout) :: tag
    type(string_t), intent(inout) :: string
    logical, intent(out) :: done
    character(2), parameter :: WS = BLANK // TAB
    type(string_t) :: name, value
    integer :: p1, p2, i
    
    p1 = verify (string, WS);  if (p1 == 0)  call err ()
    p2 = p1

    ! Look for first terminating '>' or '/>'
    if (extract (string, p1, p2) == ">") then
       done = .true.
    else
       p2 = p1 + 1
       if (extract (string, p1, p2) == "/>") then
          done = .true.
       else

          ! Look for '='
          p2 = scan (string, '=')
          if (p2 == 0)  call err ()
          name = trim (extract (string, p1, p2 - 1))
          
          ! Look for '"'
          string = extract (string, p2 + 1)
          p1 = verify (string, WS);  if (p1 == 0)  call err ()
          p2 = p1
          if (extract (string, p1, p2) /= '"')  call err ()

          ! Look for matching '"' and get value
          string = extract (string, p2 + 1)
          p1 = 1
          p2 = scan (string, '"')
          if (p2 == 0)  call err ()
          value = extract (string, p1, p2 - 1)

          SCAN_KNOWN_ATTRIBUTES: do i = 1, size (tag%attribute)
             if (name == tag%attribute(i)%name) then
                call tag%attribute(i)%set_value (value)
                exit SCAN_KNOWN_ATTRIBUTES
             end if
          end do SCAN_KNOWN_ATTRIBUTES
          
          string = extract (string, p2 + 1)
          done = .false.
       end if
    end if

  contains
    
    subroutine err ()
      call msg_fatal ("XML: Error reading attributes of '" // char (tag%name) &
           // "': syntax error")
    end subroutine err
    
  end subroutine tag_read_attribute
    
\eatline
\nwendcode{}\nwbegindocs{454}Read the content string of a tag.  We check for the appropriate
closing tag and report it.  If a closing tag does not match in name,
ignore it.

Note: this assumes that no tag with the same name is embedded in the
current content.  Also, we do not check for XML validity inside the
content.
\nwenddocs{}\nwbegincode{455}\moddef{XML: tag: TBP}\plusendmoddef
  procedure :: read_content => tag_read_content
\nwendcode{}\nwbegincode{456}\moddef{XML: procedures}\plusendmoddef
  subroutine tag_read_content (tag, cstream, content, closing)
    class(xml_tag_t), intent(in) :: tag
    type(cstream_t), intent(inout) :: cstream
    type(string_t), intent(out) :: content
    type(string_t) :: string
    logical, intent(out) :: closing
    integer :: iostat
    integer :: p0, p1, p2
    character(2), parameter :: WS = BLANK // TAB
    call cstream%get_record (content, iostat)
    if (iostat /= 0)  call err_io ()
    closing = .false.
    FIND_CLOSING: do p0 = 1, len (content) - 1

       ! Look for terminating </
       p1 = p0
       p2 = p1 + 1
       if (extract (content, p1, p2) == "</") then

          ! Look for closing tag name
          string = extract (content, p2 + 1)
          p1 = verify (string, WS);  if (p1 == 0)  call err_incomplete ()
          p2 = p1 + len (tag%name) - 1
          if (extract (string, p1, p2) == tag%name) then
             
             ! Tag name matches: look for final >
             string = extract (string, p2 + 1)
             p1 = verify (string, WS);  if (p1 == 0)  call err_incomplete ()
             p2 = p1
             if (extract (string, p1, p2) /= ">")  call err_incomplete ()

             ! Return trailing text to the stream
             string = extract (string, p2 + 1)
             if (string /= "")  call cstream%revert_record (string)
             content = extract (content, 1, p0 -1)
             closing = .true.
             exit FIND_CLOSING

          end if
       end if
    end do FIND_CLOSING
    
  contains
    
    subroutine err_io ()
      select case (iostat)
      case (:-1)
         call msg_fatal ("XML: Error reading content of '" // char (tag%name) &
              // "': end of file")
      case (1:)
         call msg_fatal ("XML: Error reading content of '" // char (tag%name) &
              // "': I/O error")
      end select
      closing = .false.
    end subroutine err_io
    
    subroutine err_incomplete ()
      call msg_fatal ("XML: Error reading content '" // char (tag%name) &
           // "': closing tag incomplete")
      closing = .false.
    end subroutine err_incomplete
    
  end subroutine tag_read_content
          
\eatline
\nwendcode{}\nwbegindocs{457}\nwdocspar
\subsection{Unit tests}
\nwenddocs{}\nwbegincode{458}\moddef{XML: public}\plusendmoddef
  public :: xml_test
\nwendcode{}\nwbegincode{459}\moddef{XML: tests}\endmoddef
  subroutine xml_test (u, results)
    integer, intent(in) :: u
    type(test_results_t), intent(inout) :: results
  \LA{}XML: execute tests\RA{}
end subroutine xml_test
  
\eatline
\nwendcode{}\nwbegindocs{460}\nwdocspar
\subsection{Auxiliary Routines}
Show the contents of a temporary file, i.e., open unit.
\nwenddocs{}\nwbegincode{461}\moddef{XML: tests}\plusendmoddef
  subroutine show (u_tmp, u)
    integer, intent(in) :: u_tmp, u
    character (80) :: buffer
    integer :: iostat
    write (u, "(A)")  "File content:"
    rewind (u_tmp)
    do
       read (u_tmp, "(A)", iostat = iostat)  buffer
       if (iostat /= 0)  exit
       write (u, "(A)")  trim (buffer)
    end do
    rewind (u_tmp)
  end subroutine show

\eatline
\nwendcode{}\nwbegindocs{462}\nwdocspar
\subsection{Basic Tag I/O}
Write a tag and read it again, using a temporary file.
\nwenddocs{}\nwbegincode{463}\moddef{XML: execute tests}\endmoddef
  call test (xml_1, "xml_1", &
       "basic I/O", &
       u, results)
\nwendcode{}\nwbegincode{464}\moddef{XML: tests}\plusendmoddef
  subroutine xml_1 (u)
    integer, intent(in) :: u
    type(xml_tag_t), allocatable :: tag
    integer :: u_tmp
    type(cstream_t) :: cstream
    logical :: success
    
    write (u, "(A)")  "* Test output: xml_1"
    write (u, "(A)")  "*   Purpose: write and read tag"
    write (u, "(A)")
    
    write (u, "(A)")  "* Empty tag"
    write (u, *)

    u_tmp = free_unit ()
    open (u_tmp, status = "scratch", action = "readwrite")

    allocate (tag)
    call tag%init (var_str ("tagname"))
    call tag%write (u_tmp)
    write (u_tmp, *)
    deallocate (tag)

    call show (u_tmp, u)

    write (u, *)
    write (u, "(A)")  "Result from read:"
    call cstream%init (u_tmp)
    allocate (tag)
    call tag%init (var_str ("tagname"))
    call tag%read (cstream, success)
    call tag%write (u)
    write (u, *)
    write (u, "(A,L1)")  "success = ", success
    deallocate (tag)
    call cstream%final ()
    
    write (u, *)
    write (u, "(A)")  "* Tag with preceding blank lines"
    write (u, *)

    u_tmp = free_unit ()
    open (u_tmp, status = "scratch", action = "readwrite")

    allocate (tag)
    call tag%init (var_str ("tagname"))
    write (u_tmp, *)
    write (u_tmp, "(A)")  "    "
    write (u_tmp, "(2x)", advance = "no")
    call tag%write (u_tmp)
    write (u_tmp, *)
    deallocate (tag)

    call show (u_tmp, u)

    write (u, *)
    write (u, "(A)")  "Result from read:"
    call cstream%init (u_tmp)
    allocate (tag)
    call tag%init (var_str ("tagname"))
    call tag%read (cstream, success)
    call tag%write (u)
    write (u, *)
    write (u, "(A,L1)")  "success = ", success
    deallocate (tag)
    call cstream%final ()
    
    write (u, *)
    write (u, "(A)")  "* Tag with preceding comments"
    write (u, *)

    u_tmp = free_unit ()
    open (u_tmp, status = "scratch", action = "readwrite")

    allocate (tag)
    call tag%init (var_str ("tagname"))
    write (u_tmp, "(A)")  "<!-- comment -->"
    write (u_tmp, *)
    write (u_tmp, "(A)")  "<!-- multiline"
    write (u_tmp, "(A)")  "     comment -->"
    call tag%write (u_tmp)
    write (u_tmp, *)
    deallocate (tag)

    call show (u_tmp, u)

    write (u, *)
    write (u, "(A)")  "Result from read:"
    call cstream%init (u_tmp)
    allocate (tag)
    call tag%init (var_str ("tagname"))
    call tag%read (cstream, success)
    call tag%write (u)
    write (u, *)
    write (u, "(A,L1)")  "success = ", success
    deallocate (tag)
    call cstream%final ()
    
    write (u, *)
    write (u, "(A)")  "* Tag with name mismatch"
    write (u, *)

    u_tmp = free_unit ()
    open (u_tmp, status = "scratch", action = "readwrite")

    allocate (tag)
    call tag%init (var_str ("wrongname"))
    call tag%write (u_tmp)
    write (u_tmp, *)
    deallocate (tag)

    call show (u_tmp, u)

    write (u, *)
    write (u, "(A)")  "Result from read:"
    call cstream%init (u_tmp)
    allocate (tag)
    call tag%init (var_str ("tagname"))
    call tag%read (cstream, success)
    call tag%write (u)
    write (u, *)
    write (u, "(A,L1)")  "success = ", success
    deallocate (tag)
    call cstream%final ()
    
    write (u, "(A)")
    write (u, "(A)")  "* Test output end: xml_1"

  end subroutine xml_1

\eatline
\nwendcode{}\nwbegindocs{465}\nwdocspar
\subsection{Optional Tag}
Write and read two tags, one of them optional.
\nwenddocs{}\nwbegincode{466}\moddef{XML: execute tests}\plusendmoddef
  call test (xml_2, "xml_2", &
       "optional tag", &
       u, results)
\nwendcode{}\nwbegincode{467}\moddef{XML: tests}\plusendmoddef
  subroutine xml_2 (u)
    integer, intent(in) :: u
    type(xml_tag_t), allocatable :: tag1, tag2
    integer :: u_tmp
    type(cstream_t) :: cstream
    logical :: success

    write (u, "(A)")  "* Test output: xml_2"
    write (u, "(A)")  "*   Purpose: handle optional tag"
    write (u, "(A)")
    
    write (u, "(A)")  "* Optional tag present"
    write (u, *)
    
    u_tmp = free_unit ()
    open (u_tmp, status = "scratch", action = "readwrite")

    allocate (tag1)
    call tag1%init (var_str ("option"))
    call tag1%write (u_tmp)
    write (u_tmp, *)
    allocate (tag2)
    call tag2%init (var_str ("tagname"))
    call tag2%write (u_tmp)
    write (u_tmp, *)
    deallocate (tag1, tag2)

    call show (u_tmp, u)

    write (u, *)
    write (u, "(A)")  "Result from read:"
    call cstream%init (u_tmp)
    allocate (tag1)
    call tag1%init (var_str ("option"))
    call tag1%read (cstream, success)
    call tag1%write (u)
    write (u, *)
    write (u, "(A,L1)")  "success = ", success
    write (u, *)
    allocate (tag2)
    call tag2%init (var_str ("tagname"))
    call tag2%read (cstream, success)
    call tag2%write (u)
    write (u, *)
    write (u, "(A,L1)")  "success = ", success
    deallocate (tag1, tag2)
    call cstream%final ()
    
    write (u, *)
    write (u, "(A)")  "* Optional tag absent"
    write (u, *)
    
    u_tmp = free_unit ()
    open (u_tmp, status = "scratch", action = "readwrite")

    allocate (tag2)
    call tag2%init (var_str ("tagname"))
    call tag2%write (u_tmp)
    write (u_tmp, *)
    deallocate (tag2)

    call show (u_tmp, u)

    write (u, *)
    write (u, "(A)")  "Result from read:"
    call cstream%init (u_tmp)
    allocate (tag1)
    call tag1%init (var_str ("option"))
    call tag1%read (cstream, success)
    call tag1%write (u)
    write (u, *)
    write (u, "(A,L1)")  "success = ", success
    write (u, *)
    allocate (tag2)
    call tag2%init (var_str ("tagname"))
    call tag2%read (cstream, success)
    call tag2%write (u)
    write (u, *)
    write (u, "(A,L1)")  "success = ", success
    deallocate (tag1, tag2)
    call cstream%final ()
    
    write (u, "(A)")
    write (u, "(A)")  "* Test output end: xml_2"

  end subroutine xml_2

\eatline
\nwendcode{}\nwbegindocs{468}\nwdocspar
\subsection{Optional Tag}
Write and read a tag with single-line content.
\nwenddocs{}\nwbegincode{469}\moddef{XML: execute tests}\plusendmoddef
  call test (xml_3, "xml_3", &
       "content", &
       u, results)
\nwendcode{}\nwbegincode{470}\moddef{XML: tests}\plusendmoddef
  subroutine xml_3 (u)
    integer, intent(in) :: u
    type(xml_tag_t), allocatable :: tag
    integer :: u_tmp
    type(cstream_t) :: cstream
    logical :: success, closing
    type(string_t) :: content

    write (u, "(A)")  "* Test output: xml_3"
    write (u, "(A)")  "*   Purpose: handle tag with content"
    write (u, "(A)")
    
    write (u, "(A)")  "* Tag without content"
    write (u, *)
    
    u_tmp = free_unit ()
    open (u_tmp, status = "scratch", action = "readwrite")

    allocate (tag)
    call tag%init (var_str ("tagname"))
    call tag%write (u_tmp)
    write (u_tmp, *)
    deallocate (tag)

    call show (u_tmp, u)

    write (u, *)
    write (u, "(A)")  "Result from read:"
    call cstream%init (u_tmp)
    allocate (tag)
    call tag%init (var_str ("tagname"))
    call tag%read (cstream, success)
    call tag%write (u)
    write (u, *)
    write (u, "(A,L1)")  "success = ", success
    write (u, "(A,L1)")  "content = ", tag%has_content
    write (u, *)
    deallocate (tag)
    call cstream%final ()
    
    write (u, "(A)")  "* Tag with content"
    write (u, *)
    
    u_tmp = free_unit ()
    open (u_tmp, status = "scratch", action = "readwrite")

    allocate (tag)
    call tag%init (var_str ("tagname"), has_content = .true.)
    call tag%write (var_str ("Content text"), u_tmp)
    write (u_tmp, *)
    deallocate (tag)

    call show (u_tmp, u)

    write (u, *)
    write (u, "(A)")  "Result from read:"
    call cstream%init (u_tmp)
    allocate (tag)
    call tag%init (var_str ("tagname"))
    call tag%read (cstream, success)
    call tag%read_content (cstream, content, closing)
    call tag%write (u)
    write (u, "(A)", advance = "no")  char (content)
    call tag%close (u)
    write (u, *)
    write (u, "(A,L1)")  "success = ", success
    write (u, "(A,L1)")  "content = ", tag%has_content
    write (u, "(A,L1)")  "closing = ", closing
    deallocate (tag)
    call cstream%final ()
    
    write (u, *)
    write (u, "(A)")  "* Tag with multiline content"
    write (u, *)
    
    u_tmp = free_unit ()
    open (u_tmp, status = "scratch", action = "readwrite")

    allocate (tag)
    call tag%init (var_str ("tagname"), has_content = .true.)
    call tag%write (u_tmp)
    write (u_tmp, *)
    write (u_tmp, "(A)")  "Line 1"
    write (u_tmp, "(A)")  "Line 2"
    call tag%close (u_tmp)
    write (u_tmp, *)
    deallocate (tag)

    call show (u_tmp, u)

    write (u, *)
    write (u, "(A)")  "Result from read:"
    call cstream%init (u_tmp)
    allocate (tag)
    call tag%init (var_str ("tagname"))
    call tag%read (cstream, success)
    call tag%write (u)
    write (u, *)
    do
       call tag%read_content (cstream, content, closing)
       if (closing)  exit
       write (u, "(A)")  char (content)
    end do
    call tag%close (u)
    write (u, *)
    write (u, "(A,L1)")  "success = ", success
    write (u, "(A,L1)")  "content = ", tag%has_content
    deallocate (tag)
    call cstream%final ()
    
    write (u, "(A)")
    write (u, "(A)")  "* Test output end: xml_3"

  end subroutine xml_3

\eatline
\nwendcode{}\nwbegindocs{471}\nwdocspar
\subsection{Basic Tag I/O}
Write a tag and read it again, using a temporary file.
\nwenddocs{}\nwbegincode{472}\moddef{XML: execute tests}\plusendmoddef
  call test (xml_4, "xml_4", &
       "attributes", &
       u, results)
\nwendcode{}\nwbegincode{473}\moddef{XML: tests}\plusendmoddef
  subroutine xml_4 (u)
    integer, intent(in) :: u
    type(xml_tag_t), allocatable :: tag
    integer :: u_tmp
    type(cstream_t) :: cstream
    logical :: success
    
    write (u, "(A)")  "* Test output: xml_4"
    write (u, "(A)")  "*   Purpose: handle tag with attributes"
    write (u, "(A)")
    
    write (u, "(A)")  "* Tag with one mandatory and one optional attribute,"
    write (u, "(A)")  "* unknown attribute ignored"
    write (u, *)

    u_tmp = free_unit ()
    open (u_tmp, status = "scratch", action = "readwrite")

    allocate (tag)
    call tag%init (var_str ("tagname"), &
         [xml_attribute (var_str ("a1"), var_str ("foo")), &
          xml_attribute (var_str ("a3"), var_str ("gee"))])
    call tag%write (u_tmp)
    deallocate (tag)

    call show (u_tmp, u)

    write (u, *)
    write (u, "(A)")  "Result from read:"
    call cstream%init (u_tmp)
    allocate (tag)
    call tag%init (var_str ("tagname"), &
         [xml_attribute (var_str ("a1")), &
          xml_attribute (var_str ("a2"), var_str ("bar"))])
    call tag%read (cstream, success)
    call tag%write (u)
    write (u, *)
    deallocate (tag)
    call cstream%final ()
    
    write (u, "(A)")
    write (u, "(A)")  "* Test output end: xml_4"

  end subroutine xml_4

\eatline
\nwendcode{}\nwfilename{../../../src/rng/rng.nw}\nwbegindocs{0}% -*- ess-noweb-default-code-mode: f90-mode; noweb-default-code-mode: f90-mode; -*- 
% WHIZARD code as NOWEB source: combinatorics and such

\chapter{Random-Number Generator}

These modules implement abstract types and tools for random-number
generation. 
\begin{description}
\item[rng\_base]
  Abstract random-number generator and factory
\item[selectors]
  Selection depending on weights and random numbers
\end{description}

Implementation of the RNG abstract types:
\begin{description}
\item[Module {\tt{}rng{\char95}tao}:]
  Interface to the TAO random number generator
  which the VAMP package provides.  Note that VAMP explicitly requests
  this generator.
\end{description}



\clearpage
%%%%%%%%%%%%%%%%%%%%%%%%%%%%%%%%%%%%%%%%%%%%%%%%%%%%%%%%%%%%%%%%%%%%%%%%
\section{Generic Random-Number Generator}
For all generator implementations, we define a {\tt{}rng} type which
represents the state of a random-number generator with the associated
methods that produce a random number.  Furthermore, we define a
{\tt{}rng{\char95}factory} type.  An object of this type is capable of allocating
a sequence of {\tt{}rng} objects.  These generator states should be, if
possible, statistically independent, so they can be used in parallel
in different places of the event-generation chain.
\nwenddocs{}\nwbegincode{1}\moddef{\code{}rng{\char95}base.f90\edoc{}}\endmoddef
\LA{}File header\RA{}

module rng_base

  use kinds
  use io_units
  use format_utils, only: write_indent
  use unit_tests

\LA{}Standard module head\RA{}

\LA{}RNG base: public\RA{}

\LA{}RNG base: types\RA{}

\LA{}RNG base: interfaces\RA{}

\LA{}RNG base: test types\RA{}

contains

\LA{}RNG base: tests\RA{}

end module rng_base
\eatline
\nwendcode{}\nwbegindocs{2}\nwdocspar
\subsection{Generator type}
The rng object is actually the state of the random-number generator.
The methods initialize/reset and call the generator for this state.
\nwenddocs{}\nwbegincode{3}\moddef{RNG base: public}\endmoddef
  public :: rng_t
\nwendcode{}\nwbegincode{4}\moddef{RNG base: types}\endmoddef
  type, abstract :: rng_t
   contains
   \LA{}RNG base: rng: TBP\RA{}
  end type rng_t

\eatline
\nwendcode{}\nwbegindocs{5}The {\tt{}init} method initializes the generator and sets a seed.
We should implement the interface such that a single integer is
sufficient for a seed.

The seed may be omitted.  The behavior without seed is not defined, however.
\nwenddocs{}\nwbegincode{6}\moddef{RNG base: rng: TBP}\endmoddef
  procedure (rng_init), deferred :: init
\nwendcode{}\nwbegincode{7}\moddef{RNG base: interfaces}\endmoddef
  abstract interface
     subroutine rng_init (rng, seed)
       import
       class(rng_t), intent(out) :: rng
       integer, intent(in), optional :: seed
     end subroutine rng_init
  end interface
  
\eatline
\nwendcode{}\nwbegindocs{8}The {\tt{}final} method deallocates memory where necessary and allows
for another call of {\tt{}init} to reset the generator.
\nwenddocs{}\nwbegincode{9}\moddef{RNG base: rng: TBP}\plusendmoddef
  procedure (rng_final), deferred :: final
\nwendcode{}\nwbegincode{10}\moddef{RNG base: interfaces}\plusendmoddef
  abstract interface
     subroutine rng_final (rng)
       import
       class(rng_t), intent(inout) :: rng
     end subroutine rng_final
  end interface
  
\eatline
\nwendcode{}\nwbegindocs{11}Output.  We should, at least, identify the generator.
\nwenddocs{}\nwbegincode{12}\moddef{RNG base: rng: TBP}\plusendmoddef
  procedure (rng_write), deferred :: write
\nwendcode{}\nwbegincode{13}\moddef{RNG base: interfaces}\plusendmoddef
  abstract interface
     subroutine rng_write (rng, unit, indent)
       import
       class(rng_t), intent(in) :: rng
       integer, intent(in), optional :: unit, indent
     end subroutine rng_write
  end interface
  
\eatline
\nwendcode{}\nwbegindocs{14}These routines generate a single and an array of default-precision
random numbers, respectively.
\nwenddocs{}\nwbegincode{15}\moddef{RNG base: rng: TBP}\plusendmoddef
  generic :: generate => generate_single, generate_array
  procedure (rng_generate_single), deferred :: generate_single
  procedure (rng_generate_array), deferred :: generate_array
\nwendcode{}\nwbegincode{16}\moddef{RNG base: interfaces}\plusendmoddef
  abstract interface
     subroutine rng_generate_single (rng, x)
       import
       class(rng_t), intent(inout) :: rng
       real(default), intent(out) :: x
     end subroutine rng_generate_single
  end interface
  
  abstract interface
     subroutine rng_generate_array (rng, x)
       import
       class(rng_t), intent(inout) :: rng
       real(default), dimension(:), intent(out) :: x
     end subroutine rng_generate_array
  end interface
  
\eatline
\nwendcode{}\nwbegindocs{17}\nwdocspar
\subsection{RNG Factory}
A factory object has a {\tt{}make} method that allocates and initializes
a new generator of appropriate type.  It uses a 16-bit integer for
initialization.  For a real-life implementation, the factory should
return a sequence of statistically independent generators, and for
different seeds, the sequences should also be independent.
\nwenddocs{}\nwbegincode{18}\moddef{RNG base: public}\plusendmoddef
  public :: rng_factory_t
\nwendcode{}\nwbegincode{19}\moddef{RNG base: types}\plusendmoddef
  type, abstract :: rng_factory_t
   contains
   \LA{}RNG base: rng factory: TBP\RA{}
  end type rng_factory_t
  
\eatline
\nwendcode{}\nwbegindocs{20}Output.  Should be short, just report the seed and current state of
the factory.
\nwenddocs{}\nwbegincode{21}\moddef{RNG base: rng factory: TBP}\endmoddef
  procedure (rng_factory_write), deferred :: write
\nwendcode{}\nwbegincode{22}\moddef{RNG base: interfaces}\plusendmoddef
  abstract interface
     subroutine rng_factory_write (object, unit)
       import
       class(rng_factory_t), intent(in) :: object
       integer, intent(in), optional :: unit
     end subroutine rng_factory_write
  end interface
  
\eatline
\nwendcode{}\nwbegindocs{23}Initialize.  It should be possible to do this repeatedly, resetting the
state.  The default seed should be 0.
\nwenddocs{}\nwbegincode{24}\moddef{RNG base: rng factory: TBP}\plusendmoddef
  procedure (rng_factory_init), deferred :: init
\nwendcode{}\nwbegincode{25}\moddef{RNG base: interfaces}\plusendmoddef
  abstract interface
     subroutine rng_factory_init (factory, seed)
       import
       class(rng_factory_t), intent(out) :: factory
       integer(i16), intent(in), optional :: seed
     end subroutine rng_factory_init
  end interface
       
\eatline
\nwendcode{}\nwbegindocs{26}Spawn a new generator.
\nwenddocs{}\nwbegincode{27}\moddef{RNG base: rng factory: TBP}\plusendmoddef
  procedure (rng_factory_make), deferred :: make
\nwendcode{}\nwbegincode{28}\moddef{RNG base: interfaces}\plusendmoddef
  abstract interface
     subroutine rng_factory_make (factory, rng)
       import
       class(rng_factory_t), intent(inout) :: factory
       class(rng_t), intent(out), allocatable :: rng
     end subroutine rng_factory_make
  end interface
  
\eatline
\nwendcode{}\nwbegindocs{29}\nwdocspar
\subsection{Unit tests}
\nwenddocs{}\nwbegincode{30}\moddef{RNG base: public}\plusendmoddef
  public :: rng_base_test
\nwendcode{}\nwbegincode{31}\moddef{RNG base: tests}\endmoddef
  subroutine rng_base_test (u, results)
    integer, intent(in) :: u
    type(test_results_t), intent(inout) :: results
  \LA{}RNG base: execute tests\RA{}
  end subroutine rng_base_test
  
\eatline
\nwendcode{}\nwbegindocs{32}\nwdocspar
\subsubsection{Test generator}
The test 'mock' random generator generates a repeating series with the numbers
$0.1, 0.3, 0.5, 0.7, 0.9$.  It has an integer stored as state.  The
integer must be one of $1,3,5,7,9$.
\nwenddocs{}\nwbegincode{33}\moddef{RNG base: public}\plusendmoddef
  public :: rng_test_t
\nwendcode{}\nwbegincode{34}\moddef{RNG base: test types}\endmoddef
  type, extends (rng_t) :: rng_test_t
     integer :: state = 1
   contains
     procedure :: write => rng_test_write
     procedure :: init => rng_test_init
     procedure :: final => rng_test_final
     procedure :: generate_single => rng_test_generate_single
     procedure :: generate_array => rng_test_generate_array
  end type rng_test_t
  
\eatline
\nwendcode{}\nwbegindocs{35}Output.  The state is a single number, so print it.
\nwenddocs{}\nwbegincode{36}\moddef{RNG base: tests}\plusendmoddef
  subroutine rng_test_write (rng, unit, indent)
    class(rng_test_t), intent(in) :: rng
    integer, intent(in), optional :: unit, indent
    integer :: u, ind
    u = given_output_unit (unit)
    ind = 0;  if (present (indent))  ind = indent
    call write_indent (u, ind)
    write (u, "(A,I0,A)")  "Random-number generator: &
         &test (state = ", rng%state, ")"
  end subroutine rng_test_write
  
\eatline
\nwendcode{}\nwbegindocs{37}The default seed is 1.
\nwenddocs{}\nwbegincode{38}\moddef{RNG base: tests}\plusendmoddef
  subroutine rng_test_init (rng, seed)
    class(rng_test_t), intent(out) :: rng
    integer, intent(in), optional :: seed
    if (present (seed))  rng%state = seed
  end subroutine rng_test_init
  
\eatline
\nwendcode{}\nwbegindocs{39}Nothing to finalize:
\nwenddocs{}\nwbegincode{40}\moddef{RNG base: tests}\plusendmoddef
  subroutine rng_test_final (rng)
    class(rng_test_t), intent(inout) :: rng
  end subroutine rng_test_final
  
\eatline
\nwendcode{}\nwbegindocs{41}Generate a single number and advance the state.
\nwenddocs{}\nwbegincode{42}\moddef{RNG base: tests}\plusendmoddef
  subroutine rng_test_generate_single (rng, x)
    class(rng_test_t), intent(inout) :: rng
    real(default), intent(out) :: x
    x = rng%state / 10._default
    rng%state = mod (rng%state + 2, 10)
  end subroutine rng_test_generate_single
  
\eatline
\nwendcode{}\nwbegindocs{43}The array generator calls the single-item generator multiple times.
\nwenddocs{}\nwbegincode{44}\moddef{RNG base: tests}\plusendmoddef
  subroutine rng_test_generate_array (rng, x)
    class(rng_test_t), intent(inout) :: rng
    real(default), dimension(:), intent(out) :: x
    integer :: i
    do i = 1, size (x)
       call rng%generate (x(i))
    end do
  end subroutine rng_test_generate_array
  
\eatline
\nwendcode{}\nwbegindocs{45}\nwdocspar
\subsubsection{Test Factory}
This factory makes {\tt{}rng{\char95}test{\char95}t} generators, initialized with
integers 1, 3, 5, 7, 9 if given the input 0, 1, 2, 3, 4.  The
generators within one sequence are all identical, however.
\nwenddocs{}\nwbegincode{46}\moddef{RNG base: public}\plusendmoddef
  public :: rng_test_factory_t
\nwendcode{}\nwbegincode{47}\moddef{RNG base: test types}\plusendmoddef
  type, extends (rng_factory_t) :: rng_test_factory_t
     integer :: seed = 1
   contains
   \LA{}RNG base: rng test factory: TBP\RA{}
  end type rng_test_factory_t
  
\eatline
\nwendcode{}\nwbegindocs{48}Output.
\nwenddocs{}\nwbegincode{49}\moddef{RNG base: rng test factory: TBP}\endmoddef
  procedure :: write => rng_test_factory_write
\nwendcode{}\nwbegincode{50}\moddef{RNG base: tests}\plusendmoddef
  subroutine rng_test_factory_write (object, unit)
    class(rng_test_factory_t), intent(in) :: object
    integer, intent(in), optional :: unit
    integer :: u
    u = given_output_unit (unit)
    write (u, "(1x,A,I0,A)")  "RNG factory: test (", object%seed, ")"
  end subroutine rng_test_factory_write

\eatline
\nwendcode{}\nwbegindocs{51}Initialize, translating the given seed.
\nwenddocs{}\nwbegincode{52}\moddef{RNG base: rng test factory: TBP}\plusendmoddef
  procedure :: init => rng_test_factory_init
\nwendcode{}\nwbegincode{53}\moddef{RNG base: tests}\plusendmoddef
  subroutine rng_test_factory_init (factory, seed)
    class(rng_test_factory_t), intent(out) :: factory
    integer(i16), intent(in), optional :: seed
    if (present (seed))  factory%seed = mod (seed * 2 + 1, 10)
  end subroutine rng_test_factory_init
    
\eatline
\nwendcode{}\nwbegincode{54}\moddef{RNG base: rng test factory: TBP}\plusendmoddef
  procedure :: make => rng_test_factory_make
\nwendcode{}\nwbegincode{55}\moddef{RNG base: tests}\plusendmoddef
  subroutine rng_test_factory_make (factory, rng)
    class(rng_test_factory_t), intent(inout) :: factory
    class(rng_t), intent(out), allocatable :: rng
    allocate (rng_test_t :: rng)
    select type (rng)
    type is (rng_test_t)
       call rng%init (int (factory%seed))
    end select
  end subroutine rng_test_factory_make

\eatline
\nwendcode{}\nwbegindocs{56}\nwdocspar
\subsubsection{Generator check}
Initialize the generator and draw random numbers.
\nwenddocs{}\nwbegincode{57}\moddef{RNG base: execute tests}\endmoddef
  call test (rng_base_1, "rng_base_1", &
       "rng initialization and call", &
       u, results)
\nwendcode{}\nwbegincode{58}\moddef{RNG base: tests}\plusendmoddef
  subroutine rng_base_1 (u)
    integer, intent(in) :: u
    class(rng_t), allocatable :: rng

    real(default) :: x
    real(default), dimension(2) :: x2
    
    write (u, "(A)")  "* Test output: rng_base_1"
    write (u, "(A)")  "*   Purpose: initialize and call a test random-number &
         &generator"
    write (u, "(A)")
    
    write (u, "(A)")  "* Initialize generator"
    write (u, "(A)")

    allocate (rng_test_t :: rng)
    call rng%init (3)

    call rng%write (u)
    
    write (u, "(A)")
    write (u, "(A)")  "* Get random number"
    write (u, "(A)")
    
    call rng%generate (x)
    write (u, "(A,2(1x,F9.7))")  "x =", x

    write (u, "(A)")
    write (u, "(A)")  "* Get random number pair"
    write (u, "(A)")
    
    call rng%generate (x2)
    write (u, "(A,2(1x,F9.7))")  "x =", x2
    
    write (u, "(A)")
    write (u, "(A)")  "* Cleanup"
        
    call rng%final ()
    
    write (u, "(A)")
    write (u, "(A)")  "* Test output end: rng_base_1"
    
  end subroutine rng_base_1
    
\eatline
\nwendcode{}\nwbegindocs{59}\nwdocspar
\subsubsection{Factory check}
Set up a factory and spawn generators.
\nwenddocs{}\nwbegincode{60}\moddef{RNG base: execute tests}\plusendmoddef
  call test (rng_base_2, "rng_base_2", &
       "rng factory", &
       u, results)
\nwendcode{}\nwbegincode{61}\moddef{RNG base: tests}\plusendmoddef
  subroutine rng_base_2 (u)
    integer, intent(in) :: u
    type(rng_test_factory_t) :: rng_factory
    class(rng_t), allocatable :: rng
    
    write (u, "(A)")  "* Test output: rng_base_2"
    write (u, "(A)")  "*   Purpose: initialize and use a rng factory"
    write (u, "(A)")
    
    write (u, "(A)")  "* Initialize factory"
    write (u, "(A)")

    call rng_factory%init ()
    call rng_factory%write (u)

    write (u, "(A)")
    write (u, "(A)")  "* Make a generator"
    write (u, "(A)")

    call rng_factory%make (rng)
    call rng%write (u)

    write (u, "(A)")
    write (u, "(A)")  "* Cleanup"
        
    call rng%final ()
    
    write (u, "(A)")
    write (u, "(A)")  "* Test output end: rng_base_2"
    
  end subroutine rng_base_2
    
\eatline
\nwendcode{}\nwbegindocs{62}\nwdocspar
\clearpage
%%%%%%%%%%%%%%%%%%%%%%%%%%%%%%%%%%%%%%%%%%%%%%%%%%%%%%%%%%%%%%%%%%%%%%%%
\section{Select from a weighted sample}
\nwenddocs{}\nwbegincode{63}\moddef{\code{}selectors.f90\edoc{}}\endmoddef
\LA{}File header\RA{}

module selectors

  use kinds
  use io_units
  use unit_tests
  use diagnostics
  use rng_base

\LA{}Standard module head\RA{}

\LA{}Selectors: public\RA{}

\LA{}Selectors: types\RA{}

contains
  
\LA{}Selectors: procedures\RA{}

\LA{}Selectors: tests\RA{}

end module selectors
\eatline
\nwendcode{}\nwbegindocs{64}\nwdocspar
\subsection{Selector type}
The rng object is actually the state of the random-number generator.
The methods initialize/reset and call the generator for this state.
\nwenddocs{}\nwbegincode{65}\moddef{Selectors: public}\endmoddef
  public :: selector_t
\nwendcode{}\nwbegincode{66}\moddef{Selectors: types}\endmoddef
  type :: selector_t
     integer, dimension(:), allocatable :: map
     real(default), dimension(:), allocatable :: weight
     real(default), dimension(:), allocatable :: acc
   contains
   \LA{}Selectors: selector: TBP\RA{}
end type selector_t

\eatline
\nwendcode{}\nwbegindocs{67}Display contents.
\nwenddocs{}\nwbegincode{68}\moddef{Selectors: selector: TBP}\endmoddef
  procedure :: write => selector_write
\nwendcode{}\nwbegincode{69}\moddef{Selectors: procedures}\endmoddef
  subroutine selector_write (object, unit)
    class(selector_t), intent(in) :: object
    integer, intent(in), optional :: unit
    integer :: u, i
    u = given_output_unit (unit)
    write (u, "(1x,A)")  "Selector: i, weight, acc. weight"
    if (allocated (object%weight)) then
       do i = 1, size (object%weight)
          write (u, "(3x,I0,1x,ES19.12,1x,ES19.12)") &
               object%map(i), object%weight(i), object%acc(i)
       end do
    else
       write (u, "(3x,A)")  "[undefined]"
    end if
  end subroutine selector_write
 
\eatline
\nwendcode{}\nwbegindocs{70}We pack the input weight array such that zero-weight entries are
removed.  We also normalize it.  This makes a {\tt{}map} array for
mapping the selected weight to the actual entry necessary.
\nwenddocs{}\nwbegincode{71}\moddef{Selectors: selector: TBP}\plusendmoddef
  procedure :: init => selector_init
\nwendcode{}\nwbegincode{72}\moddef{Selectors: procedures}\plusendmoddef
  subroutine selector_init (selector, weight)
    class(selector_t), intent(out) :: selector
    real(default), dimension(:), intent(in) :: weight
    real(default) :: s
    integer :: n, i
    logical, dimension(:), allocatable :: mask
    if (size (weight) == 0) &
         call msg_bug ("Selector init: zero-size weight array")
    if (any (weight < 0)) &
         call msg_bug ("Selector init: negative weight")
    s = sum (weight)
    if (s == 0) &
         call msg_bug ("Selector init: all weights are zero")
    allocate (mask (size (weight)), &
         source = weight /= 0)
    n = count (mask)
    allocate (selector%map (n), &
         source = pack ([(i, i = 1, size (weight))], mask))
    allocate (selector%weight (n), &
         source = pack (weight / s, mask))
    allocate (selector%acc (n))
    selector%acc(1) = selector%weight(1)
    do i = 2, n - 1
       selector%acc(i) = selector%acc(i-1) + selector%weight(i)
    end do
    selector%acc(n) = 1
  end subroutine selector_init
    
\eatline
\nwendcode{}\nwbegindocs{73}Select an entry based upon the number {\tt{}x}, which should be a
uniformly distributed random number between 0 and 1.
\nwenddocs{}\nwbegincode{74}\moddef{Selectors: selector: TBP}\plusendmoddef
  procedure :: select => selector_select
\nwendcode{}\nwbegincode{75}\moddef{Selectors: procedures}\plusendmoddef
  function selector_select (selector, x) result (n)
    class(selector_t), intent(in) :: selector
    real(default), intent(in) :: x
    integer :: n
    integer :: i
    if (x < 0 .or. x > 1) &
         call msg_bug ("Selector: random number out of range")
    do i = 1, size (selector%acc)
       if (x <= selector%acc(i))  exit
    end do
    n = selector%map(i)
  end function selector_select

\eatline
\nwendcode{}\nwbegindocs{76}Use the provided random-number generator to select an entry.  (Unless there
is only one entry.)
\nwenddocs{}\nwbegincode{77}\moddef{Selectors: selector: TBP}\plusendmoddef
  procedure :: generate => selector_generate
\nwendcode{}\nwbegincode{78}\moddef{Selectors: procedures}\plusendmoddef
  subroutine selector_generate (selector, rng, n)
    class(selector_t), intent(in) :: selector
    class(rng_t), intent(inout) :: rng
    integer, intent(out) :: n
    real(default) :: x
    select case (size (selector%acc))
    case (1);  n = 1
    case default
       call rng%generate (x)
       n = selector%select (x)
    end select
  end subroutine selector_generate

\eatline
\nwendcode{}\nwbegindocs{79}Determine the normalized weight for a selected entry.  We use a
linear search for the inverse lookup, assuming that efficiency is not
an issue for this function.
\nwenddocs{}\nwbegincode{80}\moddef{Selectors: selector: TBP}\plusendmoddef
  procedure :: get_weight => selector_get_weight
\nwendcode{}\nwbegincode{81}\moddef{Selectors: procedures}\plusendmoddef
  function selector_get_weight (selector, n) result (weight)
    class(selector_t), intent(in) :: selector
    integer, intent(in) :: n
    real(default) :: weight
    integer :: i
    do i = 1, size (selector%weight)
       if (selector%map(i) == n) then
          weight = selector%weight(i)
          return
       end if
    end do
    weight = 0
  end function selector_get_weight

\eatline
\nwendcode{}\nwbegindocs{82}\nwdocspar
\subsection{Unit tests}
\nwenddocs{}\nwbegincode{83}\moddef{Selectors: public}\plusendmoddef
  public :: selectors_test
\nwendcode{}\nwbegincode{84}\moddef{Selectors: tests}\endmoddef
  subroutine selectors_test (u, results)
    integer, intent(in) :: u
    type(test_results_t), intent(inout) :: results
  \LA{}Selectors: execute tests\RA{}
  end subroutine selectors_test
  
\eatline
\nwendcode{}\nwbegindocs{85}\nwdocspar
\subsubsection{Basic check}
Initialize the selector and draw random numbers.
\nwenddocs{}\nwbegincode{86}\moddef{Selectors: execute tests}\endmoddef
  call test (selectors_1, "selectors_1", &
       "rng initialization and call", &
       u, results)
\nwendcode{}\nwbegincode{87}\moddef{Selectors: tests}\plusendmoddef
  subroutine selectors_1 (u)
    integer, intent(in) :: u
    type(selector_t) :: selector
    class(rng_t), allocatable, target :: rng
    integer :: i, n

    write (u, "(A)")  "* Test output: selectors_1"
    write (u, "(A)")  "*   Purpose: initialize a selector and test it"
    write (u, "(A)")
    
    write (u, "(A)")  "* Initialize selector"
    write (u, "(A)")

    call selector%init &
         ([2._default, 3.5_default, 0._default, &
         2._default, 0.5_default, 2._default])
    call selector%write (u)

    write (u, "(A)")
    write (u, "(A)")  "* Select numbers using predictable test generator"
    write (u, "(A)")

    allocate (rng_test_t :: rng)
    call rng%init (1)

    do i = 1, 5
       call selector%generate (rng, n)
       write (u, "(1x,I0)")  n
    end do

    write (u, "(A)")
    write (u, "(A)")  "* Select numbers using real input number"
    write (u, "(A)")
    
    write (u, "(1x,A,I0)")  "select(0.00) = ", selector%select (0._default)
    write (u, "(1x,A,I0)")  "select(0.77) = ", selector%select (0.77_default)
    write (u, "(1x,A,I0)")  "select(1.00) = ", selector%select (1._default)

    write (u, "(A)")
    write (u, "(A)")  "* Get weight"
    write (u, "(A)")
    
    write (u, "(1x,A,ES19.12)")  "weight(2) =", selector%get_weight(2)
    write (u, "(1x,A,ES19.12)")  "weight(3) =", selector%get_weight(3)
        
    write (u, "(A)")
    write (u, "(A)")  "* Cleanup"
        
    call rng%final ()
    
    write (u, "(A)")
    write (u, "(A)")  "* Test output end: selectors_1"
    
  end subroutine selectors_1
    
\eatline
\nwendcode{}\nwbegindocs{88}\nwdocspar
%%%%%%%%%%%%%%%%%%%%%%%%%%%%%%%%%%%%%%%%%%%%%%%%%%%%%%%%%%%%%%%%%%%%%%%%
\section{TAO Random-Number Generator}
This module provides an implementation for the generic random-number
generator.  Actually, we interface the TAO random-number generator
which is available via the VAMP package.
\nwenddocs{}\nwbegincode{89}\moddef{\code{}rng{\char95}tao.f90\edoc{}}\endmoddef
\LA{}File header\RA{}

module rng_tao

  use kinds
  use io_units
  use format_utils, only: write_indent
  use unit_tests
  use tao_random_numbers !NODEP!

  use rng_base
  
\LA{}Standard module head\RA{}

\LA{}RNG tao: public\RA{}

\LA{}RNG tao: types\RA{}

contains
  
\LA{}RNG tao: procedures\RA{}

\LA{}RNG tao: tests\RA{}

end module rng_tao
\eatline
\nwendcode{}\nwbegindocs{90}\nwdocspar
\subsection{Generator type}
The rng object is actually the state of the random-number generator.
The methods initialize/reset and call the generator for this state.

We keep the seed, in case we want to recover it later, and count the
number of calls since seeding.
\nwenddocs{}\nwbegincode{91}\moddef{RNG tao: public}\endmoddef
  public :: rng_tao_t
\nwendcode{}\nwbegincode{92}\moddef{RNG tao: types}\endmoddef
  type, extends (rng_t) :: rng_tao_t
     integer :: seed = 0
     integer :: n_calls = 0
     type(tao_random_state) :: state
   contains
   \LA{}RNG tao: rng tao: TBP\RA{}
  end type rng_tao_t

\eatline
\nwendcode{}\nwbegindocs{93}Output: Display seed and number of calls.
\nwenddocs{}\nwbegincode{94}\moddef{RNG tao: rng tao: TBP}\endmoddef
  procedure :: write => rng_tao_write
\nwendcode{}\nwbegincode{95}\moddef{RNG tao: procedures}\endmoddef
  subroutine rng_tao_write (rng, unit, indent)
    class(rng_tao_t), intent(in) :: rng
    integer, intent(in), optional :: unit, indent
    integer :: u, ind
    u = given_output_unit (unit)
    ind = 0;  if (present (indent))  ind = indent
    call write_indent (u, ind)
    write (u, "(A)")  "TAO random-number generator:"
    call write_indent (u, ind)
    write (u, "(2x,A,I0)")  "seed  = ", rng%seed
    call write_indent (u, ind)
    write (u, "(2x,A,I0)")  "calls = ", rng%n_calls
  end subroutine rng_tao_write
  
\eatline
\nwendcode{}\nwbegindocs{96}The {\tt{}init} method initializes the generator and sets a seed.
We should implement the interface such that a single integer is
sufficient for a seed.

The seed may be omitted.  The default seed is 0.
\nwenddocs{}\nwbegincode{97}\moddef{RNG tao: rng tao: TBP}\plusendmoddef
  procedure :: init => rng_tao_init
\nwendcode{}\nwbegincode{98}\moddef{RNG tao: procedures}\plusendmoddef
  subroutine rng_tao_init (rng, seed)
    class(rng_tao_t), intent(out) :: rng
    integer, intent(in), optional :: seed
    if (present (seed))  rng%seed = seed
    call tao_random_create (rng%state, rng%seed)
  end subroutine rng_tao_init
  
\eatline
\nwendcode{}\nwbegindocs{99}The {\tt{}final} method deallocates memory where necessary and allows
for another call of {\tt{}init} to reset the generator.
\nwenddocs{}\nwbegincode{100}\moddef{RNG tao: rng tao: TBP}\plusendmoddef
  procedure :: final => rng_tao_final
\nwendcode{}\nwbegincode{101}\moddef{RNG tao: procedures}\plusendmoddef
  subroutine rng_tao_final (rng)
    class(rng_tao_t), intent(inout) :: rng
    call tao_random_destroy (rng%state)
  end subroutine rng_tao_final
  
\eatline
\nwendcode{}\nwbegindocs{102}These routines generate a single and an array of default-precision
random numbers, respectively.

We have to convert from explicit double to abstract default
precision.  Under normal conditions, both are equivalent, however.
Unless, someone decides to do single precision, there is always
an interface for {\tt{}tao{\char95}random{\char95}numbers}.
\nwenddocs{}\nwbegincode{103}\moddef{RNG tao: rng tao: TBP}\plusendmoddef
  procedure :: generate_single => rng_tao_generate_single
  procedure :: generate_array => rng_tao_generate_array
\nwendcode{}\nwbegincode{104}\moddef{RNG tao: procedures}\plusendmoddef
  subroutine rng_tao_generate_single (rng, x)
    class(rng_tao_t), intent(inout) :: rng
    real(default), intent(out) :: x
    real(default) :: r
    call tao_random_number (rng%state, r)
    x = r
    rng%n_calls = rng%n_calls + 1
  end subroutine rng_tao_generate_single
  
  subroutine rng_tao_generate_array (rng, x)
    class(rng_tao_t), intent(inout) :: rng
    real(default), dimension(:), intent(out) :: x
    real(default) :: r
    integer :: i
    do i = 1, size (x)
       call tao_random_number (rng%state, r)
       x(i) = r
    end do
    rng%n_calls = rng%n_calls + size (x)
  end subroutine rng_tao_generate_array
  
\eatline
\nwendcode{}\nwbegindocs{105}\nwdocspar
\subsubsection{Factory}
This factory makes {\tt{}rng{\char95}tao{\char95}t} generators, initialized with the seeds
\begin{equation}
  s_i = s_0 * 2^{16} + i
\end{equation}
where $s_0$ is the seed (a 16-bit integer) given to the factory
object, and $i$ is the index in the generated sequence of generators,
starting with zero.
\nwenddocs{}\nwbegincode{106}\moddef{RNG tao: public}\plusendmoddef
  public :: rng_tao_factory_t
\nwendcode{}\nwbegincode{107}\moddef{RNG tao: types}\plusendmoddef
  type, extends (rng_factory_t) :: rng_tao_factory_t
     integer(i16) :: s = 0
     integer(i16) :: i = 0
   contains
   \LA{}RNG tao: rng tao factory: TBP\RA{}
  end type rng_tao_factory_t
  
\eatline
\nwendcode{}\nwbegindocs{108}Output.
\nwenddocs{}\nwbegincode{109}\moddef{RNG tao: rng tao factory: TBP}\endmoddef
  procedure :: write => rng_tao_factory_write
\nwendcode{}\nwbegincode{110}\moddef{RNG tao: procedures}\plusendmoddef
  subroutine rng_tao_factory_write (object, unit)
    class(rng_tao_factory_t), intent(in) :: object
    integer, intent(in), optional :: unit
    integer :: u
    u = given_output_unit (unit)
    write (u, "(1x,A,2(I0,A))") &
         "RNG factory: tao (", object%s, ",", object%i, ")"
  end subroutine rng_tao_factory_write

\eatline
\nwendcode{}\nwbegindocs{111}Initialize, translating the given seed.
\nwenddocs{}\nwbegincode{112}\moddef{RNG tao: rng tao factory: TBP}\plusendmoddef
  procedure :: init => rng_tao_factory_init
\nwendcode{}\nwbegincode{113}\moddef{RNG tao: procedures}\plusendmoddef
  subroutine rng_tao_factory_init (factory, seed)
    class(rng_tao_factory_t), intent(out) :: factory
    integer(i16), intent(in), optional :: seed
    if (present (seed))  factory%s = seed
  end subroutine rng_tao_factory_init
    
\eatline
\nwendcode{}\nwbegincode{114}\moddef{RNG tao: rng tao factory: TBP}\plusendmoddef
  procedure :: make => rng_tao_factory_make
\nwendcode{}\nwbegincode{115}\moddef{RNG tao: procedures}\plusendmoddef
  subroutine rng_tao_factory_make (factory, rng)
    class(rng_tao_factory_t), intent(inout) :: factory
    class(rng_t), intent(out), allocatable :: rng
    allocate (rng_tao_t :: rng)
    select type (rng)
    type is (rng_tao_t)
       call rng%init (factory%s * 65536 + factory%i)
       factory%i = int (factory%i + 1, kind = i16)
    end select
  end subroutine rng_tao_factory_make

\eatline
\nwendcode{}\nwbegindocs{116}\nwdocspar
\subsection{Unit tests}
\nwenddocs{}\nwbegincode{117}\moddef{RNG tao: public}\plusendmoddef
  public :: rng_tao_test
\nwendcode{}\nwbegincode{118}\moddef{RNG tao: tests}\endmoddef
  subroutine rng_tao_test (u, results)
    integer, intent(in) :: u
    type(test_results_t), intent(inout) :: results
  \LA{}RNG tao: execute tests\RA{}
  end subroutine rng_tao_test
  
\eatline
\nwendcode{}\nwbegindocs{119}\nwdocspar
\subsubsection{Generator check}
Initialize the generator and draw random numbers.
\nwenddocs{}\nwbegincode{120}\moddef{RNG tao: execute tests}\endmoddef
  call test (rng_tao_1, "rng_tao_1", &
       "rng initialization and call", &
       u, results)
\nwendcode{}\nwbegincode{121}\moddef{RNG tao: tests}\plusendmoddef
  subroutine rng_tao_1 (u)
    integer, intent(in) :: u
    class(rng_t), allocatable, target :: rng

    real(default) :: x
    real(default), dimension(2) :: x2
    
    write (u, "(A)")  "* Test output: rng_tao_1"
    write (u, "(A)")  "*   Purpose: initialize and call the TAO random-number &
         &generator"
    write (u, "(A)")
    
    write (u, "(A)")  "* Initialize generator (default seed)"
    write (u, "(A)")

    allocate (rng_tao_t :: rng)
    call rng%init ()
    
    call rng%write (u)
    
    write (u, "(A)")
    write (u, "(A)")  "* Get random number"
    write (u, "(A)")
    
    call rng%generate (x)
    write (u, "(A,2(1x,F9.7))")  "x =", x

    write (u, "(A)")
    write (u, "(A)")  "* Get random number pair"
    write (u, "(A)")
    
    call rng%generate (x2)
    write (u, "(A,2(1x,F9.7))")  "x =", x2
    
    write (u, "(A)")
    call rng%write (u)

    write (u, "(A)")
    write (u, "(A)")  "* Cleanup"
        
    call rng%final ()
    
    write (u, "(A)")
    write (u, "(A)")  "* Test output end: rng_tao_1"
    
  end subroutine rng_tao_1
    
\eatline
\nwendcode{}\nwbegindocs{122}\nwdocspar
\subsubsection{Factory check}
Set up a factory and spawn generators.
\nwenddocs{}\nwbegincode{123}\moddef{RNG tao: execute tests}\plusendmoddef
  call test (rng_tao_2, "rng_tao_2", &
       "rng factory", &
       u, results)
\nwendcode{}\nwbegincode{124}\moddef{RNG tao: tests}\plusendmoddef
  subroutine rng_tao_2 (u)
    integer, intent(in) :: u
    type(rng_tao_factory_t) :: rng_factory
    class(rng_t), allocatable :: rng
    real(default) :: x
    
    write (u, "(A)")  "* Test output: rng_tao_2"
    write (u, "(A)")  "*   Purpose: initialize and use a rng factory"
    write (u, "(A)")
    
    write (u, "(A)")  "* Initialize factory"
    write (u, "(A)")

    call rng_factory%init ()
    call rng_factory%write (u)

    write (u, "(A)")
    write (u, "(A)")  "* Make a generator"
    write (u, "(A)")

    call rng_factory%make (rng)
    call rng%write (u)
    call rng%generate (x)
    write (u, *)
    write (u, "(1x,A,F7.5)")  "x = ", x
    call rng%final ()
    deallocate (rng)

    write (u, "(A)")
    write (u, "(A)")  "* Repeat"
    write (u, "(A)")

    call rng_factory%make (rng)
    call rng%write (u)
    call rng%generate (x)
    write (u, *)
    write (u, "(1x,A,F7.5)")  "x = ", x
    call rng%final ()
    deallocate (rng)

    write (u, *)
    call rng_factory%write (u)

    write (u, "(A)")
    write (u, "(A)")  "* Initialize factory with different seed"
    write (u, "(A)")

    call rng_factory%init (1_i16)
    call rng_factory%write (u)

    write (u, "(A)")
    write (u, "(A)")  "* Make a generator"
    write (u, "(A)")

    call rng_factory%make (rng)
    call rng%write (u)
    call rng%generate (x)
    write (u, *)
    write (u, "(1x,A,F7.5)")  "x = ", x
    call rng%final ()
    deallocate (rng)
    
    write (u, *)
    call rng_factory%write (u)

    write (u, "(A)")
    write (u, "(A)")  "* Test output end: rng_tao_2"
    
  end subroutine rng_tao_2
    
\eatline
\nwendcode{}\nwfilename{../../../src/physics/physics.nw}\nwbegindocs{0}% -*- ess-noweb-default-code-mode: f90-mode; noweb-default-code-mode: f90-mode; -*-
% WHIZARD code as NOWEB source: combinatorics and such

\chapter{Physics}

Here we collect definitions and functions that we need for (particle)
physics in general, to make them available for the more specific needs
of WHIZARD.
\begin{description}
\item[physics\_defs]
  Physical constants.
\item[c\_particles]
  A simple data type for particles which is C compatible.
\item[lorentz]
  Define three-vectors, four-vectors and Lorentz
  transformations and common operations for them.
\item[sm\_physics]
  Here, running functions are stored for special kinematical setup like
  running coupling constants, Catani-Seymour dipoles, or Sudakov factors.
\item[sm\_qcd]
  Definitions and methods for dealing with the running QCD coupling.
\end{description}

\clearpage
%%%%%%%%%%%%%%%%%%%%%%%%%%%%%%%%%%%%%%%%%%%%%%%%%%%%%%%%%%%%%%%%%%%%%%%%
\section{Physics Constants}
There is also the generic {\tt{}constants} module.  The constants listed
here are more specific for particle physics.
\nwenddocs{}\nwbegincode{1}\moddef{\code{}physics{\char95}defs.f90\edoc{}}\endmoddef
\LA{}File header\RA{}

module physics_defs

\LA{}Use kinds\RA{}
  use constants, only: one, two, three

\LA{}Standard module head\RA{}

\LA{}Physics defs: public parameters\RA{}

end module physics_defs
\eatline
\nwendcode{}\nwbegindocs{2}\nwdocspar
\subsection{Units}
Conversion from energy units to cross-section units.
\nwenddocs{}\nwbegincode{3}\moddef{Physics defs: public parameters}\endmoddef
  real(default), parameter, public :: &
       conv = 0.38937966e12_default
\nwendcode{}\nwbegindocs{4}\nwdocspar
Rescaling factor.
\nwenddocs{}\nwbegincode{5}\moddef{Physics defs: public parameters}\plusendmoddef
  real(default), parameter, public :: &
       pb_per_fb = 1.e-3_default
\nwendcode{}\nwbegindocs{6}\nwdocspar
String for the default energy and cross-section units.
\nwenddocs{}\nwbegincode{7}\moddef{Physics defs: public parameters}\plusendmoddef
  character(*), parameter, public :: &
       energy_unit = "GeV"
  character(*), parameter, public :: &
       cross_section_unit = "fb"
\nwendcode{}\nwbegindocs{8}\nwdocspar
\subsection{SM and QCD constants}
\nwenddocs{}\nwbegincode{9}\moddef{Physics defs: public parameters}\plusendmoddef
  real(default), parameter, public :: &
       NC = three, &
       CF = (NC**2 - one) / two / NC, &
       CA = NC, &
       TR = one / two
\nwendcode{}\nwbegindocs{10}\nwdocspar
\subsection{Parameter Reference values}
These are used exclusively in the context of
running QCD parameters.  In other contexts, we rely on the uniform
parameter set as provided by the model definition, modifiable by the
user.
\nwenddocs{}\nwbegincode{11}\moddef{Physics defs: public parameters}\plusendmoddef
  real(default), public, parameter :: MZ_REF = 91.188_default
  real(default), public, parameter :: ALPHA_QCD_MZ_REF = 0.1178_default
  real(default), public, parameter :: LAMBDA_QCD_REF = 200.e-3_default
\eatline
\nwendcode{}\nwbegindocs{12}\nwdocspar
\subsection{Particle codes}
Let us define a few particle codes independent of the model.

We need an UNDEFINED value:
\nwenddocs{}\nwbegincode{13}\moddef{Physics defs: public parameters}\plusendmoddef
  integer, parameter, public :: UNDEFINED = 0

\eatline
\nwendcode{}\nwbegindocs{14}SM fermions:
\nwenddocs{}\nwbegincode{15}\moddef{Physics defs: public parameters}\plusendmoddef
  integer, parameter, public :: ELECTRON = 11

\eatline
\nwendcode{}\nwbegindocs{16}Gauge bosons:
\nwenddocs{}\nwbegincode{17}\moddef{Physics defs: public parameters}\plusendmoddef
  integer, parameter, public :: GLUON = 21
  integer, parameter, public :: PHOTON = 22
  integer, parameter, public :: Z_BOSON = 23
  integer, parameter, public :: W_BOSON = 24

\eatline
\nwendcode{}\nwbegindocs{18}Light mesons:
\nwenddocs{}\nwbegincode{19}\moddef{Physics defs: public parameters}\plusendmoddef
  integer, parameter, public :: PION = 111
  integer, parameter, public :: PIPLUS = 211
  integer, parameter, public :: PIMINUS = - PIPLUS

\eatline
\nwendcode{}\nwbegindocs{20}Di-Quarks:
\nwenddocs{}\nwbegincode{21}\moddef{Physics defs: public parameters}\plusendmoddef
  integer, parameter, public :: UD0 = 2101
  integer, parameter, public :: UD1 = 2103
  integer, parameter, public :: UU1 = 2203

\eatline
\nwendcode{}\nwbegindocs{22}Mesons:
\nwenddocs{}\nwbegincode{23}\moddef{Physics defs: public parameters}\plusendmoddef
  integer, parameter, public :: K0L = 130
  integer, parameter, public :: K0S = 310
  integer, parameter, public :: K0 = 311
  integer, parameter, public :: KPLUS = 321
  integer, parameter, public :: DPLUS = 411
  integer, parameter, public :: D0 = 421
  integer, parameter, public :: B0 = 511
  integer, parameter, public :: BPLUS = 521

\eatline
\nwendcode{}\nwbegindocs{24}Light baryons:
\nwenddocs{}\nwbegincode{25}\moddef{Physics defs: public parameters}\plusendmoddef
  integer, parameter, public :: PROTON = 2212
  integer, parameter, public :: NEUTRON = 2112
  integer, parameter, public :: DELTAPLUSPLUS = 2224
  integer, parameter, public :: DELTAPLUS = 2214
  integer, parameter, public :: DELTA0 = 2114
  integer, parameter, public :: DELTAMINUS = 1114

\eatline
\nwendcode{}\nwbegindocs{26}Strange baryons:
\nwenddocs{}\nwbegincode{27}\moddef{Physics defs: public parameters}\plusendmoddef
  integer, parameter, public :: SIGMAPLUS = 3222
  integer, parameter, public :: SIGMA0 = 3212
  integer, parameter, public :: SIGMAMINUS = 3112

\eatline
\nwendcode{}\nwbegindocs{28}Charmed baryons:
\nwenddocs{}\nwbegincode{29}\moddef{Physics defs: public parameters}\plusendmoddef
  integer, parameter, public :: SIGMACPLUSPLUS = 4222
  integer, parameter, public :: SIGMACPLUS = 4212
  integer, parameter, public :: SIGMAC0 = 4112

\eatline
\nwendcode{}\nwbegindocs{30}Bottom baryons:
\nwenddocs{}\nwbegincode{31}\moddef{Physics defs: public parameters}\plusendmoddef
  integer, parameter, public :: SIGMAB0 = 5212
  integer, parameter, public :: SIGMABPLUS = 5222

\eatline
\nwendcode{}\nwbegindocs{32}81-100 are reserved for internal codes. Hadron and beam remnants:
\nwenddocs{}\nwbegincode{33}\moddef{Physics defs: public parameters}\plusendmoddef
  integer, parameter, public :: BEAM_REMNANT = 9999
  integer, parameter, public :: HADRON_REMNANT = 90
  integer, parameter, public :: HADRON_REMNANT_SINGLET = 91
  integer, parameter, public :: HADRON_REMNANT_TRIPLET = 92
  integer, parameter, public :: HADRON_REMNANT_OCTET = 93

\eatline
\eatline
\nwendcode{}\nwbegindocs{34}\nwdocspar
Further particle codes for internal use:
\nwenddocs{}\nwbegincode{35}\moddef{Physics defs: public parameters}\plusendmoddef
  integer, parameter, public :: INTERNAL = 94
  integer, parameter, public :: INVALID = 97

  integer, parameter, public :: COMPOSITE = 99

\eatline
\nwendcode{}\nwbegindocs{36}\nwdocspar
\subsection{Spin codes}
Somewhat redundant, but for better readability we define named
constants for spin types.  If the mass is nonzero, this is equal to
the number of degrees of freedom.
\nwenddocs{}\nwbegincode{37}\moddef{Physics defs: public parameters}\plusendmoddef
  integer, parameter, public:: UNKNOWN = 0
  integer, parameter, public :: SCALAR = 1, SPINOR = 2, VECTOR = 3, &
                                VECTORSPINOR = 4, TENSOR = 5
\eatline
\nwendcode{}\nwbegindocs{38}Isospin types and charge types are counted in an analogous way,
where charge type 1 is charge 0, 2 is charge 1/3, and so on. Zero
always means unknown. Note that charge and isospin types have an
explicit sign.

Color types are defined as the dimension of the representation.

\clearpage
%%%%%%%%%%%%%%%%%%%%%%%%%%%%%%%%%%%%%%%%%%%%%%%%%%%%%%%%%%%%%%%%%%%%%%%%
\section{C-compatible Particle Type}
For easy communication with C code, we introduce a simple C-compatible
type for particles.   The components are either default C integers or
default C doubles.

The {\tt{}c{\char95}prt} type is transparent, and its contents should be regarded
as part of the interface.
\nwenddocs{}\nwbegincode{39}\moddef{\code{}c{\char95}particles.f90\edoc{}}\endmoddef
\LA{}File header\RA{}

module c_particles

  use, intrinsic :: iso_c_binding !NODEP!

  use io_units
  use format_defs, only: FMT_14, FMT_19

\LA{}Standard module head\RA{}

\LA{}C Particles: public\RA{}

\LA{}C Particles: types\RA{}

contains

\LA{}C Particles: procedures\RA{}
end module c_particles
\eatline
\nwendcode{}\nwbegindocs{40}\nwdocspar
\nwenddocs{}\nwbegincode{41}\moddef{C Particles: public}\endmoddef
  public :: c_prt_t
\nwendcode{}\nwbegincode{42}\moddef{C Particles: types}\endmoddef
  type, bind(C) :: c_prt_t
     integer(c_int) :: type = 0
     integer(c_int) :: pdg = 0
     integer(c_int) :: polarized = 0
     integer(c_int) :: h = 0
     real(c_double) :: pe = 0
     real(c_double) :: px = 0
     real(c_double) :: py = 0
     real(c_double) :: pz = 0
     real(c_double) :: p2 = 0
  end type c_prt_t

\eatline
\nwendcode{}\nwbegindocs{43}This is for debugging only, there is no C binding.  It is a
simplified version of {\tt{}prt{\char95}write}.
\nwenddocs{}\nwbegincode{44}\moddef{C Particles: public}\plusendmoddef
  public :: c_prt_write
\nwendcode{}\nwbegincode{45}\moddef{C Particles: procedures}\endmoddef
  subroutine c_prt_write (prt, unit)
    type(c_prt_t), intent(in) :: prt
    integer, intent(in), optional :: unit
    integer :: u
    u = given_output_unit (unit);  if (u < 0)  return
    write (u, "(1x,A)", advance="no")  "prt("
    write (u, "(I0,':')", advance="no")  prt%type
    if (prt%polarized /= 0) then
       write (u, "(I0,'/',I0,'|')", advance="no")  prt%pdg, prt%h
    else
       write (u, "(I0,'|')", advance="no") prt%pdg
    end if
    write (u, "(" // FMT_14 // ",';'," // FMT_14 // ",','," // &
         FMT_14 // ",','," // FMT_14 // ")", advance="no") &
         prt%pe, prt%px, prt%py, prt%pz
    write (u, "('|'," // FMT_19 // ")", advance="no")  prt%p2
    write (u, "(A)")  ")"
  end subroutine c_prt_write

\eatline
\nwendcode{}\nwbegindocs{46}\nwdocspar
\clearpage
%%%%%%%%%%%%%%%%%%%%%%%%%%%%%%%%%%%%%%%%%%%%%%%%%%%%%%%%%%%%%%%%%%%%%%%%
\section{Lorentz algebra}
Define Lorentz vectors, three-vectors, boosts, and some functions to
manipulate them.

To make maximum use of this, all functions, if possible, are declared
elemental (or pure, if this is not possible).
\nwenddocs{}\nwbegincode{47}\moddef{\code{}lorentz.f90\edoc{}}\endmoddef
\LA{}File header\RA{}

module lorentz

\LA{}Use kinds\RA{}
  use io_units
  use constants, only: pi, twopi, degree, zero, one, eps0
  use format_defs, only: FMT_15, FMT_17, FMT_19
  use format_utils, only: pac_fmt
  use diagnostics
  use c_particles

\LA{}Standard module head\RA{}

\LA{}Lorentz: public\RA{}

\LA{}Lorentz: public operators\RA{}

\LA{}Lorentz: public functions\RA{}

\LA{}Lorentz: types\RA{}

\LA{}Lorentz: parameters\RA{}

\LA{}Lorentz: interfaces\RA{}

contains

\LA{}Lorentz: procedures\RA{}
end module lorentz
\eatline
\nwendcode{}\nwbegindocs{48}\nwdocspar
\subsection{Three-vectors}
First of all, let us introduce three-vectors in a trivial way.  The
functions and overloaded elementary operations clearly are too much
overhead, but we like to keep the interface for three-vectors and
four-vectors exactly parallel.  By the way, we might attach a label to
a vector by extending the type definition later.
\nwenddocs{}\nwbegincode{49}\moddef{Lorentz: public}\endmoddef
  public :: vector3_t
\nwendcode{}\nwbegincode{50}\moddef{Lorentz: types}\endmoddef
  type :: vector3_t
     real(default), dimension(3) :: p
  end type vector3_t

\eatline
\nwendcode{}\nwbegindocs{51}Output a vector
\nwenddocs{}\nwbegincode{52}\moddef{Lorentz: public}\plusendmoddef
  public :: vector3_write
\nwendcode{}\nwbegincode{53}\moddef{Lorentz: procedures}\endmoddef
  subroutine vector3_write (p, unit, testflag)
    type(vector3_t), intent(in) :: p
    integer, intent(in), optional :: unit
    logical, intent(in), optional :: testflag
    character(len=7) :: fmt
    integer :: u
    u = given_output_unit (unit);  if (u < 0)  return
    call pac_fmt (fmt, FMT_19, FMT_15, testflag)
    write(u, "(1x,A,3(1x," // fmt // "))") 'P = ', p%p
  end subroutine vector3_write

\eatline
\nwendcode{}\nwbegindocs{54}This is a three-vector with zero components
\nwenddocs{}\nwbegincode{55}\moddef{Lorentz: public}\plusendmoddef
  public :: vector3_null
\nwendcode{}\nwbegincode{56}\moddef{Lorentz: parameters}\endmoddef
  type(vector3_t), parameter :: vector3_null = &
       vector3_t ([ zero, zero, zero ])

\eatline
\nwendcode{}\nwbegindocs{57}Canonical three-vector:
\nwenddocs{}\nwbegincode{58}\moddef{Lorentz: public}\plusendmoddef
  public :: vector3_canonical
\nwendcode{}\nwbegincode{59}\moddef{Lorentz: procedures}\plusendmoddef
  elemental function vector3_canonical (k) result (p)
    type(vector3_t) :: p
    integer, intent(in) :: k
    p = vector3_null
    p%p(k) = 1
  end function vector3_canonical

\eatline
\nwendcode{}\nwbegindocs{60}A moving particle ($k$-axis, or arbitrary axis).  Note that the
function for the generic momentum cannot be elemental.
\nwenddocs{}\nwbegincode{61}\moddef{Lorentz: public}\plusendmoddef
  public :: vector3_moving
\nwendcode{}\nwbegincode{62}\moddef{Lorentz: interfaces}\endmoddef
  interface vector3_moving
     module procedure vector3_moving_canonical
     module procedure vector3_moving_generic
  end interface
\nwendcode{}\nwbegincode{63}\moddef{Lorentz: procedures}\plusendmoddef
  elemental function vector3_moving_canonical (p, k) result(q)
    type(vector3_t) :: q
    real(default), intent(in) :: p
    integer, intent(in) :: k
    q = vector3_null
    q%p(k) = p
  end function vector3_moving_canonical
  pure function vector3_moving_generic (p) result(q)
    real(default), dimension(3), intent(in) :: p
    type(vector3_t) :: q
    q%p = p
  end function vector3_moving_generic

\eatline
\nwendcode{}\nwbegindocs{64}Equality and inequality
\nwenddocs{}\nwbegincode{65}\moddef{Lorentz: public operators}\endmoddef
  public :: operator(==), operator(/=)
\nwendcode{}\nwbegincode{66}\moddef{Lorentz: interfaces}\plusendmoddef
  interface operator(==)
     module procedure vector3_eq
  end interface
  interface operator(/=)
     module procedure vector3_neq
  end interface
\nwendcode{}\nwbegincode{67}\moddef{Lorentz: procedures}\plusendmoddef
  elemental function vector3_eq (p, q) result (r)
    logical :: r
    type(vector3_t), intent(in) :: p,q
    r = all (abs (p%p - q%p) < eps0)
  end function vector3_eq
  elemental function vector3_neq (p, q) result (r)
    logical :: r
    type(vector3_t), intent(in) :: p,q
    r = any (abs(p%p - q%p) > eps0)
  end function vector3_neq

\eatline
\nwendcode{}\nwbegindocs{68}Define addition and subtraction
\nwenddocs{}\nwbegincode{69}\moddef{Lorentz: public operators}\plusendmoddef
  public :: operator(+), operator(-)
\nwendcode{}\nwbegincode{70}\moddef{Lorentz: interfaces}\plusendmoddef
  interface operator(+)
     module procedure add_vector3
  end interface
  interface operator(-)
     module procedure sub_vector3
  end interface
\nwendcode{}\nwbegincode{71}\moddef{Lorentz: procedures}\plusendmoddef
  elemental function add_vector3 (p, q) result (r)
    type(vector3_t) :: r
    type(vector3_t), intent(in) :: p,q
    r%p = p%p + q%p
  end function add_vector3
  elemental function sub_vector3 (p, q) result (r)
    type(vector3_t) :: r
    type(vector3_t), intent(in) :: p,q
    r%p = p%p - q%p
  end function sub_vector3

\eatline
\nwendcode{}\nwbegindocs{72}The multiplication sign is overloaded with scalar multiplication;
similarly division:
\nwenddocs{}\nwbegincode{73}\moddef{Lorentz: public operators}\plusendmoddef
  public :: operator(*), operator(/)
\nwendcode{}\nwbegincode{74}\moddef{Lorentz: interfaces}\plusendmoddef
  interface operator(*)
     module procedure prod_integer_vector3, prod_vector3_integer
     module procedure prod_real_vector3, prod_vector3_real
  end interface
  interface operator(/)
     module procedure div_vector3_real, div_vector3_integer
  end interface
\nwendcode{}\nwbegincode{75}\moddef{Lorentz: procedures}\plusendmoddef
  elemental function prod_real_vector3 (s, p) result (q)
    type(vector3_t) :: q
    real(default), intent(in) :: s
    type(vector3_t), intent(in) :: p
    q%p = s * p%p
  end function prod_real_vector3
  elemental function prod_vector3_real (p, s) result (q)
    type(vector3_t) :: q
    real(default), intent(in) :: s
    type(vector3_t), intent(in) :: p
    q%p = s * p%p
  end function prod_vector3_real
  elemental function div_vector3_real (p, s) result (q)
    type(vector3_t) :: q
    real(default), intent(in) :: s
    type(vector3_t), intent(in) :: p
    q%p = p%p/s
  end function div_vector3_real
  elemental function prod_integer_vector3 (s, p) result (q)
    type(vector3_t) :: q
    integer, intent(in) :: s
    type(vector3_t), intent(in) :: p
    q%p = s * p%p
  end function prod_integer_vector3
  elemental function prod_vector3_integer (p, s) result (q)
    type(vector3_t) :: q
    integer, intent(in) :: s
    type(vector3_t), intent(in) :: p
    q%p = s * p%p
  end function prod_vector3_integer
  elemental function div_vector3_integer (p, s) result (q)
    type(vector3_t) :: q
    integer, intent(in) :: s
    type(vector3_t), intent(in) :: p
    q%p = p%p/s
  end function div_vector3_integer

\eatline
\nwendcode{}\nwbegindocs{76}The multiplication sign can also indicate scalar products:
\nwenddocs{}\nwbegincode{77}\moddef{Lorentz: interfaces}\plusendmoddef
  interface operator(*)
     module procedure prod_vector3
  end interface
\nwendcode{}\nwbegincode{78}\moddef{Lorentz: procedures}\plusendmoddef
  elemental function prod_vector3 (p, q) result (s)
    real(default) :: s
    type(vector3_t), intent(in) :: p,q
    s = dot_product (p%p, q%p)
  end function prod_vector3

\eatline
\nwendcode{}\nwbegincode{79}\moddef{Lorentz: public functions}\endmoddef
  public :: cross_product
\nwendcode{}\nwbegincode{80}\moddef{Lorentz: interfaces}\plusendmoddef
  interface cross_product
     module procedure vector3_cross_product
  end interface
\nwendcode{}\nwbegincode{81}\moddef{Lorentz: procedures}\plusendmoddef
  elemental function vector3_cross_product (p, q) result (r)
    type(vector3_t) :: r
    type(vector3_t), intent(in) :: p,q
    integer :: i
    do i=1,3
       r%p(i) = dot_product (p%p, matmul(epsilon_three(i,:,:), q%p))
    end do
  end function vector3_cross_product

\eatline
\nwendcode{}\nwbegindocs{82}Exponentiation is defined only for integer powers.  Odd powers mean
take the square root; so {\tt{}p**1} is the length of {\tt{}p}.
\nwenddocs{}\nwbegincode{83}\moddef{Lorentz: public operators}\plusendmoddef
  public :: operator(**)
\nwendcode{}\nwbegincode{84}\moddef{Lorentz: interfaces}\plusendmoddef
  interface operator(**)
     module procedure power_vector3
  end interface
\nwendcode{}\nwbegincode{85}\moddef{Lorentz: procedures}\plusendmoddef
  elemental function power_vector3 (p, e) result (s)
    real(default) :: s
    type(vector3_t), intent(in) :: p
    integer, intent(in) :: e
    s = dot_product (p%p, p%p)
    if (e/=2) then
       if (mod(e,2)==0) then
          s = s**(e/2)
       else
          s = sqrt(s)**e
       end if
    end if
  end function power_vector3

\eatline
\nwendcode{}\nwbegindocs{86}Finally, we need a negation.
\nwenddocs{}\nwbegincode{87}\moddef{Lorentz: interfaces}\plusendmoddef
  interface operator(-)
     module procedure negate_vector3
  end interface
\nwendcode{}\nwbegincode{88}\moddef{Lorentz: procedures}\plusendmoddef
  elemental function negate_vector3 (p) result (q)
    type(vector3_t) :: q
    type(vector3_t), intent(in) :: p
    integer :: i
    do i = 1, 3
       if (abs (p%p(i)) < eps0) then
          q%p(i) = 0
       else
          q%p(i) = -p%p(i)
       end if
    end do
  end function negate_vector3

\eatline
\nwendcode{}\nwbegindocs{89}The sum function can be useful:
\nwenddocs{}\nwbegincode{90}\moddef{Lorentz: public functions}\plusendmoddef
  public :: sum
\nwendcode{}\nwbegincode{91}\moddef{Lorentz: interfaces}\plusendmoddef
  interface sum
     module procedure sum_vector3
  end interface
\eatline
\nwendcode{}\nwbegindocs{92}\nwdocspar
\nwenddocs{}\nwbegincode{93}\moddef{Lorentz: public}\plusendmoddef
  public :: vector3_set_component
\nwendcode{}\nwbegincode{94}\moddef{Lorentz: procedures}\plusendmoddef
  subroutine vector3_set_component (p, i, value)
    type(vector3_t), intent(inout) :: p
    integer, intent(in) :: i
    real(default), intent(in) :: value
    p%p(i) = value
  end subroutine vector3_set_component

\eatline
\nwendcode{}\nwbegindocs{95}\nwdocspar
\nwenddocs{}\nwbegincode{96}\moddef{Lorentz: procedures}\plusendmoddef
  pure function sum_vector3 (p) result (q)
    type(vector3_t) :: q
    type(vector3_t), dimension(:), intent(in) :: p
    integer :: i
    do i=1, 3
       q%p(i) = sum (p%p(i))
    end do
  end function sum_vector3

\eatline
\nwendcode{}\nwbegindocs{97}Any component:
\nwenddocs{}\nwbegincode{98}\moddef{Lorentz: public}\plusendmoddef
  public :: vector3_get_component
\eatline
\nwendcode{}\nwbegincode{99}\moddef{Lorentz: procedures}\plusendmoddef
  elemental function vector3_get_component (p, k) result (c)
    type(vector3_t), intent(in) :: p
    integer, intent(in) :: k
    real(default) :: c
    c = p%p(k)
  end function vector3_get_component

\eatline
\nwendcode{}\nwbegindocs{100}Extract all components.  This is not elemental.
\nwenddocs{}\nwbegincode{101}\moddef{Lorentz: public}\plusendmoddef
  public :: vector3_get_components
\nwendcode{}\nwbegincode{102}\moddef{Lorentz: procedures}\plusendmoddef
  pure function vector3_get_components (p) result (a)
    type(vector3_t), intent(in) :: p
    real(default), dimension(3) :: a
    a = p%p
  end function vector3_get_components

\eatline
\nwendcode{}\nwbegindocs{103}This function returns the direction of a three-vector, i.e., a
normalized three-vector.  If the vector is null, we return a null vector.
\nwenddocs{}\nwbegincode{104}\moddef{Lorentz: public functions}\plusendmoddef
  public :: direction
\nwendcode{}\nwbegincode{105}\moddef{Lorentz: interfaces}\plusendmoddef
  interface direction
     module procedure vector3_get_direction
  end interface
\nwendcode{}\nwbegincode{106}\moddef{Lorentz: procedures}\plusendmoddef
  elemental function vector3_get_direction (p) result (q)
    type(vector3_t) :: q
    type(vector3_t), intent(in) :: p
    real(default) :: pp
    pp = p**1
    if (pp > eps0) then
       q%p = p%p / pp
    else
       q%p = 0
    end if
  end function vector3_get_direction

\eatline
\nwendcode{}\nwbegindocs{107}\nwdocspar
\subsection{Four-vectors}
In four-vectors the zero-component needs special treatment, therefore
we do not use the standard operations.  Sure, we pay for the extra
layer of abstraction by losing efficiency; so we have to assume that
the time-critical applications do not involve four-vector operations.
\nwenddocs{}\nwbegincode{108}\moddef{Lorentz: public}\plusendmoddef
  public :: vector4_t
\nwendcode{}\nwbegincode{109}\moddef{Lorentz: types}\plusendmoddef
  type :: vector4_t
     real(default), dimension(0:3) :: p = &
        [zero, zero, zero, zero]
  end type vector4_t
\eatline
\nwendcode{}\nwbegindocs{110}Output a vector
\nwenddocs{}\nwbegincode{111}\moddef{Lorentz: public}\plusendmoddef
  public :: vector4_write
\nwendcode{}\nwbegincode{112}\moddef{Lorentz: procedures}\plusendmoddef
  subroutine vector4_write (p, unit, show_mass, testflag)
    type(vector4_t), intent(in) :: p
    integer, intent(in), optional :: unit
    logical, intent(in), optional :: show_mass
    logical, intent(in), optional :: testflag
    integer :: u
    character(len=7) :: fmt
    call pac_fmt (fmt, FMT_19, FMT_15, testflag)
    u = given_output_unit (unit);  if (u < 0)  return
    write(u, "(1x,A,1x," // fmt // ")") 'E = ', p%p(0)
    write(u, "(1x,A,3(1x," // fmt // "))") 'P = ', p%p(1:)
    if (present (show_mass)) then
       if (show_mass) &
            write (u, "(1x,A,1x," // fmt // ")") 'M = ', p**1
    end if
  end subroutine vector4_write

\eatline
\nwendcode{}\nwbegindocs{113}Binary I/O
\nwenddocs{}\nwbegincode{114}\moddef{Lorentz: public}\plusendmoddef
  public :: vector4_write_raw
  public :: vector4_read_raw
\nwendcode{}\nwbegincode{115}\moddef{Lorentz: procedures}\plusendmoddef
  subroutine vector4_write_raw (p, u)
    type(vector4_t), intent(in) :: p
    integer, intent(in) :: u
    write (u) p%p
  end subroutine vector4_write_raw

  subroutine vector4_read_raw (p, u, iostat)
    type(vector4_t), intent(out) :: p
    integer, intent(in) :: u
    integer, intent(out), optional :: iostat
    read (u, iostat=iostat) p%p
  end subroutine vector4_read_raw

\eatline
\nwendcode{}\nwbegindocs{116}This is a four-vector with zero components
\nwenddocs{}\nwbegincode{117}\moddef{Lorentz: public}\plusendmoddef
  public :: vector4_null
\nwendcode{}\nwbegincode{118}\moddef{Lorentz: parameters}\plusendmoddef
  type(vector4_t), parameter :: vector4_null = &
       vector4_t ([ zero, zero, zero, zero ])

\eatline
\nwendcode{}\nwbegindocs{119}Canonical four-vector:
\nwenddocs{}\nwbegincode{120}\moddef{Lorentz: public}\plusendmoddef
  public :: vector4_canonical
\nwendcode{}\nwbegincode{121}\moddef{Lorentz: procedures}\plusendmoddef
  elemental function vector4_canonical (k) result (p)
    type(vector4_t) :: p
    integer, intent(in) :: k
    p = vector4_null
    p%p(k) = 1
  end function vector4_canonical

\eatline
\nwendcode{}\nwbegindocs{122}A particle at rest:
\nwenddocs{}\nwbegincode{123}\moddef{Lorentz: public}\plusendmoddef
  public :: vector4_at_rest
\nwendcode{}\nwbegincode{124}\moddef{Lorentz: procedures}\plusendmoddef
  elemental function vector4_at_rest (m) result (p)
    type(vector4_t) :: p
    real(default), intent(in) :: m
    p = vector4_t ([ m, zero, zero, zero ])
  end function vector4_at_rest

\eatline
\nwendcode{}\nwbegindocs{125}A moving particle ($k$-axis, or arbitrary axis)
\nwenddocs{}\nwbegincode{126}\moddef{Lorentz: public}\plusendmoddef
  public :: vector4_moving
\nwendcode{}\nwbegincode{127}\moddef{Lorentz: interfaces}\plusendmoddef
  interface vector4_moving
     module procedure vector4_moving_canonical
     module procedure vector4_moving_generic
  end interface
\nwendcode{}\nwbegincode{128}\moddef{Lorentz: procedures}\plusendmoddef
  elemental function vector4_moving_canonical (E, p, k) result (q)
    type(vector4_t) :: q
    real(default), intent(in) :: E, p
    integer, intent(in) :: k
    q = vector4_at_rest(E)
    q%p(k) = p
  end function vector4_moving_canonical
  elemental function vector4_moving_generic (E, p) result (q)
    type(vector4_t) :: q
    real(default), intent(in) :: E
    type(vector3_t), intent(in) :: p
    q%p(0) = E
    q%p(1:) = p%p
  end function vector4_moving_generic

\eatline
\nwendcode{}\nwbegindocs{129}Equality and inequality
\nwenddocs{}\nwbegincode{130}\moddef{Lorentz: interfaces}\plusendmoddef
  interface operator(==)
     module procedure vector4_eq
  end interface
  interface operator(/=)
     module procedure vector4_neq
  end interface
\nwendcode{}\nwbegincode{131}\moddef{Lorentz: procedures}\plusendmoddef
  elemental function vector4_eq (p, q) result (r)
    logical :: r
    type(vector4_t), intent(in) :: p,q
    r = all (abs (p%p - q%p) < eps0)
  end function vector4_eq
  elemental function vector4_neq (p, q) result (r)
    logical :: r
    type(vector4_t), intent(in) :: p,q
    r = any (abs (p%p - q%p) > eps0)
  end function vector4_neq

\eatline
\nwendcode{}\nwbegindocs{132}Addition and subtraction:
\nwenddocs{}\nwbegincode{133}\moddef{Lorentz: interfaces}\plusendmoddef
  interface operator(+)
     module procedure add_vector4
  end interface
  interface operator(-)
     module procedure sub_vector4
  end interface
\nwendcode{}\nwbegincode{134}\moddef{Lorentz: procedures}\plusendmoddef
  elemental function add_vector4 (p,q) result (r)
    type(vector4_t) :: r
    type(vector4_t), intent(in) :: p,q
    r%p = p%p + q%p
  end function add_vector4
  elemental function sub_vector4 (p,q) result (r)
    type(vector4_t) :: r
    type(vector4_t), intent(in) :: p,q
    r%p = p%p - q%p
  end function sub_vector4

\eatline
\nwendcode{}\nwbegindocs{135}We also need scalar multiplication and division:
\nwenddocs{}\nwbegincode{136}\moddef{Lorentz: interfaces}\plusendmoddef
  interface operator(*)
     module procedure prod_real_vector4, prod_vector4_real
     module procedure prod_integer_vector4, prod_vector4_integer
  end interface
  interface operator(/)
     module procedure div_vector4_real
     module procedure div_vector4_integer
  end interface
\nwendcode{}\nwbegincode{137}\moddef{Lorentz: procedures}\plusendmoddef
  elemental function prod_real_vector4 (s, p) result (q)
    type(vector4_t) :: q
    real(default), intent(in) :: s
    type(vector4_t), intent(in) :: p
    q%p = s * p%p
  end function prod_real_vector4
  elemental function prod_vector4_real (p, s) result (q)
    type(vector4_t) :: q
    real(default), intent(in) :: s
    type(vector4_t), intent(in) :: p
    q%p = s * p%p
  end function prod_vector4_real
  elemental function div_vector4_real (p, s) result (q)
    type(vector4_t) :: q
    real(default), intent(in) :: s
    type(vector4_t), intent(in) :: p
    q%p = p%p/s
  end function div_vector4_real
  elemental function prod_integer_vector4 (s, p) result (q)
    type(vector4_t) :: q
    integer, intent(in) :: s
    type(vector4_t), intent(in) :: p
    q%p = s * p%p
  end function prod_integer_vector4
  elemental function prod_vector4_integer (p, s) result (q)
    type(vector4_t) :: q
    integer, intent(in) :: s
    type(vector4_t), intent(in) :: p
    q%p = s * p%p
  end function prod_vector4_integer
  elemental function div_vector4_integer (p, s) result (q)
    type(vector4_t) :: q
    integer, intent(in) :: s
    type(vector4_t), intent(in) :: p
    q%p = p%p/s
  end function div_vector4_integer

\eatline
\nwendcode{}\nwbegindocs{138}Scalar products and squares in the Minkowski sense:
\nwenddocs{}\nwbegincode{139}\moddef{Lorentz: interfaces}\plusendmoddef
  interface operator(*)
     module procedure prod_vector4
  end interface
  interface operator(**)
     module procedure power_vector4
  end interface
\nwendcode{}\nwbegincode{140}\moddef{Lorentz: procedures}\plusendmoddef
  elemental function prod_vector4 (p, q) result (s)
    real(default) :: s
    type(vector4_t), intent(in) :: p,q
    s = p%p(0)*q%p(0) - dot_product(p%p(1:), q%p(1:))
  end function prod_vector4

\eatline
\nwendcode{}\nwbegindocs{141}The power operation for four-vectors is signed, i.e., {\tt{}p**1} is
positive for timelike and negative for spacelike vectors.  Note that
{\tt{}(p**1)**2} is not necessarily equal to {\tt{}p**2}.
\nwenddocs{}\nwbegincode{142}\moddef{Lorentz: procedures}\plusendmoddef
  elemental function power_vector4 (p, e) result (s)
    real(default) :: s
    type(vector4_t), intent(in) :: p
    integer, intent(in) :: e
    s = p*p
    if (e/=2) then
       if (mod(e,2)==0) then
          s = s**(e/2)
       elseif (s>=0) then
          s = sqrt(s)**e
       else
          s = -(sqrt(abs(s))**e)
       end if
    end if
  end function power_vector4

\eatline
\nwendcode{}\nwbegindocs{143}Finally, we introduce a negation
\nwenddocs{}\nwbegincode{144}\moddef{Lorentz: interfaces}\plusendmoddef
  interface operator(-)
     module procedure negate_vector4
  end interface
\nwendcode{}\nwbegincode{145}\moddef{Lorentz: procedures}\plusendmoddef
  elemental function negate_vector4 (p) result (q)
    type(vector4_t) :: q
    type(vector4_t), intent(in) :: p
    integer :: i
    do i = 0, 3
       if (abs (p%p(i)) < eps0) then
          q%p(i) = 0
       else
          q%p(i) = -p%p(i)
       end if
    end do
  end function negate_vector4

\eatline
\nwendcode{}\nwbegindocs{146}The sum function can be useful:
\nwenddocs{}\nwbegincode{147}\moddef{Lorentz: interfaces}\plusendmoddef
  interface sum
     module procedure sum_vector4
  end interface
\eatline
\nwendcode{}\nwbegindocs{148}\nwdocspar
\nwenddocs{}\nwbegincode{149}\moddef{Lorentz: procedures}\plusendmoddef
  pure function sum_vector4 (p) result (q)
    type(vector4_t) :: q
    type(vector4_t), dimension(:), intent(in) :: p
    integer :: i
    do i=0, 3
       q%p(i) = sum (p%p(i))
    end do
  end function sum_vector4

\eatline
\nwendcode{}\nwbegindocs{150}\nwdocspar
\subsection{Conversions}
Manually set a component of the four-vector:
\nwenddocs{}\nwbegincode{151}\moddef{Lorentz: public}\plusendmoddef
  public :: vector4_set_component
\nwendcode{}\nwbegincode{152}\moddef{Lorentz: procedures}\plusendmoddef
  subroutine vector4_set_component (p, k, c)
    type(vector4_t), intent(inout) :: p
    integer, intent(in) :: k
    real(default), intent(in) :: c
    p%p(k) = c
  end subroutine vector4_set_component

\eatline
\nwendcode{}\nwbegindocs{153}Any component:
\nwenddocs{}\nwbegincode{154}\moddef{Lorentz: public}\plusendmoddef
  public :: vector4_get_component
\nwendcode{}\nwbegincode{155}\moddef{Lorentz: procedures}\plusendmoddef
  elemental function vector4_get_component (p, k) result (c)
    real(default) :: c
    type(vector4_t), intent(in) :: p
    integer, intent(in) :: k
    c = p%p(k)
  end function vector4_get_component

\eatline
\nwendcode{}\nwbegindocs{156}Extract all components.  This is not elemental.
\nwenddocs{}\nwbegincode{157}\moddef{Lorentz: public}\plusendmoddef
  public :: vector4_get_components
\nwendcode{}\nwbegincode{158}\moddef{Lorentz: procedures}\plusendmoddef
  pure function vector4_get_components (p) result (a)
    real(default), dimension(0:3) :: a
    type(vector4_t), intent(in) :: p
    a = p%p
  end function vector4_get_components

\eatline
\nwendcode{}\nwbegindocs{159}This function returns the space part of a four-vector, such that we
can apply three-vector operations on it:
\nwenddocs{}\nwbegincode{160}\moddef{Lorentz: public functions}\plusendmoddef
  public :: space_part
\nwendcode{}\nwbegincode{161}\moddef{Lorentz: interfaces}\plusendmoddef
  interface space_part
     module procedure vector4_get_space_part
  end interface
\nwendcode{}\nwbegincode{162}\moddef{Lorentz: procedures}\plusendmoddef
  elemental function vector4_get_space_part (p) result (q)
    type(vector3_t) :: q
    type(vector4_t), intent(in) :: p
    q%p = p%p(1:)
  end function vector4_get_space_part

\eatline
\nwendcode{}\nwbegindocs{163}This function returns the direction of a four-vector, i.e., a
normalized three-vector.  If the four-vector has zero space part, we
return a null vector.
\nwenddocs{}\nwbegincode{164}\moddef{Lorentz: interfaces}\plusendmoddef
  interface direction
     module procedure vector4_get_direction
  end interface
\nwendcode{}\nwbegincode{165}\moddef{Lorentz: procedures}\plusendmoddef
  elemental function vector4_get_direction (p) result (q)
    type(vector3_t) :: q
    type(vector4_t), intent(in) :: p
    real(default) :: qq
    q%p = p%p(1:)
    qq = q**1
    if (abs(qq) > eps0) then
       q%p = q%p / qq
    else
       q%p = 0
    end if
  end function vector4_get_direction

\eatline
\nwendcode{}\nwbegindocs{166}This function returns the four-vector as an ordinary array.  A
second version for an array of four-vectors.
\nwenddocs{}\nwbegincode{167}\moddef{Lorentz: public}\plusendmoddef
  public :: assignment (=)
\nwendcode{}\nwbegincode{168}\moddef{Lorentz: interfaces}\plusendmoddef
  interface assignment (=)
     module procedure array_from_vector4_1, array_from_vector4_2, &
            array_from_vector3_1, array_from_vector3_2, &
            vector4_from_array, vector3_from_array
  end interface
\nwendcode{}\nwbegincode{169}\moddef{Lorentz: procedures}\plusendmoddef
  pure subroutine array_from_vector4_1 (a, p)
    real(default), dimension(:), intent(out) :: a
    type(vector4_t), intent(in) :: p
    a = p%p
  end subroutine array_from_vector4_1

  pure subroutine array_from_vector4_2 (a, p)
    type(vector4_t), dimension(:), intent(in) :: p
    real(default), dimension(:,:), intent(out) :: a
    integer :: i
    forall (i=1:size(p))
       a(:,i) = p(i)%p
    end forall
  end subroutine array_from_vector4_2

  pure subroutine array_from_vector3_1 (a, p)
    real(default), dimension(:), intent(out) :: a
    type(vector3_t), intent(in) :: p
    a = p%p
  end subroutine array_from_vector3_1

  pure subroutine array_from_vector3_2 (a, p)
    type(vector3_t), dimension(:), intent(in) :: p
    real(default), dimension(:,:), intent(out) :: a
    integer :: i
    forall (i=1:size(p))
       a(:,i) = p(i)%p
    end forall
  end subroutine array_from_vector3_2

  pure subroutine vector4_from_array (p, a)
    type(vector4_t), intent(out) :: p
    real(default), dimension(:), intent(in) :: a
    p%p(0:3) = a
  end subroutine vector4_from_array

  pure subroutine vector3_from_array (p, a)
    type(vector3_t), intent(out) :: p
    real(default), dimension(:), intent(in) :: a
    p%p(1:3) = a
  end subroutine vector3_from_array

\eatline
\nwendcode{}\nwbegindocs{170}Transform the momentum of a {\tt{}c{\char95}prt} object into a four-vector and
vice versa:
\nwenddocs{}\nwbegincode{171}\moddef{Lorentz: interfaces}\plusendmoddef
  interface assignment (=)
     module procedure vector4_from_c_prt, c_prt_from_vector4
  end interface
\nwendcode{}\nwbegincode{172}\moddef{Lorentz: procedures}\plusendmoddef
  pure subroutine vector4_from_c_prt (p, c_prt)
    type(vector4_t), intent(out) :: p
    type(c_prt_t), intent(in) :: c_prt
    p%p(0) = c_prt%pe
    p%p(1) = c_prt%px
    p%p(2) = c_prt%py
    p%p(3) = c_prt%pz
  end subroutine vector4_from_c_prt

  pure subroutine c_prt_from_vector4 (c_prt, p)
    type(c_prt_t), intent(out) :: c_prt
    type(vector4_t), intent(in) :: p
    c_prt%pe = p%p(0)
    c_prt%px = p%p(1)
    c_prt%py = p%p(2)
    c_prt%pz = p%p(3)
    c_prt%p2 = p ** 2
  end subroutine c_prt_from_vector4

\eatline
\nwendcode{}\nwbegindocs{173}Initialize a {\tt{}c{\char95}prt{\char95}t} object with the components of a four-vector
as its kinematical entries.  Compute the invariant mass, or use the
optional mass-squared value instead.
\nwenddocs{}\nwbegincode{174}\moddef{Lorentz: public}\plusendmoddef
  public :: vector4_to_c_prt
\nwendcode{}\nwbegincode{175}\moddef{Lorentz: procedures}\plusendmoddef
  elemental function vector4_to_c_prt (p, p2) result (c_prt)
    type(c_prt_t) :: c_prt
    type(vector4_t), intent(in) :: p
    real(default), intent(in), optional :: p2
    c_prt%pe = p%p(0)
    c_prt%px = p%p(1)
    c_prt%py = p%p(2)
    c_prt%pz = p%p(3)
    if (present (p2)) then
       c_prt%p2 = p2
    else
       c_prt%p2 = p ** 2
    end if
  end function vector4_to_c_prt

\eatline
\nwendcode{}\nwbegindocs{176}\nwdocspar
\subsection{Angles}
Return the angles in a canonical system.  The angle $\phi$ is defined
between $0\leq\phi<2\pi$.  In degenerate cases, return zero.
\nwenddocs{}\nwbegincode{177}\moddef{Lorentz: public functions}\plusendmoddef
  public :: azimuthal_angle
\nwendcode{}\nwbegincode{178}\moddef{Lorentz: interfaces}\plusendmoddef
  interface azimuthal_angle
     module procedure vector3_azimuthal_angle
     module procedure vector4_azimuthal_angle
  end interface
\nwendcode{}\nwbegincode{179}\moddef{Lorentz: procedures}\plusendmoddef
  elemental function vector3_azimuthal_angle (p) result (phi)
    real(default) :: phi
    type(vector3_t), intent(in) :: p
    if (any (abs (p%p(1:2)) > 0)) then
       phi = atan2(p%p(2), p%p(1))
       if (phi < 0) phi = phi + twopi
    else
       phi = 0
    end if
  end function vector3_azimuthal_angle
  elemental function vector4_azimuthal_angle (p) result (phi)
    real(default) :: phi
    type(vector4_t), intent(in) :: p
    phi = vector3_azimuthal_angle (space_part (p))
  end function vector4_azimuthal_angle

\eatline
\nwendcode{}\nwbegindocs{180}Azimuthal angle in degrees
\nwenddocs{}\nwbegincode{181}\moddef{Lorentz: public functions}\plusendmoddef
  public :: azimuthal_angle_deg
\nwendcode{}\nwbegincode{182}\moddef{Lorentz: interfaces}\plusendmoddef
  interface azimuthal_angle_deg
     module procedure vector3_azimuthal_angle_deg
     module procedure vector4_azimuthal_angle_deg
  end interface
\nwendcode{}\nwbegincode{183}\moddef{Lorentz: procedures}\plusendmoddef
  elemental function vector3_azimuthal_angle_deg (p) result (phi)
    real(default) :: phi
    type(vector3_t), intent(in) :: p
    phi = vector3_azimuthal_angle (p) / degree
  end function vector3_azimuthal_angle_deg
  elemental function vector4_azimuthal_angle_deg (p) result (phi)
    real(default) :: phi
    type(vector4_t), intent(in) :: p
    phi = vector4_azimuthal_angle (p) / degree
  end function vector4_azimuthal_angle_deg

\eatline
\nwendcode{}\nwbegindocs{184}The azimuthal distance of two vectors.  This is the difference of
the azimuthal angles, but cannot be larger than $\pi$: The result is
between $-\pi<\Delta\phi\leq\pi$.
\nwenddocs{}\nwbegincode{185}\moddef{Lorentz: public functions}\plusendmoddef
  public :: azimuthal_distance
\nwendcode{}\nwbegincode{186}\moddef{Lorentz: interfaces}\plusendmoddef
  interface azimuthal_distance
     module procedure vector3_azimuthal_distance
     module procedure vector4_azimuthal_distance
  end interface
\nwendcode{}\nwbegincode{187}\moddef{Lorentz: procedures}\plusendmoddef
  elemental function vector3_azimuthal_distance (p, q) result (dphi)
    real(default) :: dphi
    type(vector3_t), intent(in) :: p,q
    dphi = vector3_azimuthal_angle (q) - vector3_azimuthal_angle (p)
    if (dphi <= -pi) then
       dphi = dphi + twopi
    else if (dphi > pi) then
       dphi = dphi - twopi
    end if
  end function vector3_azimuthal_distance
  elemental function vector4_azimuthal_distance (p, q) result (dphi)
    real(default) :: dphi
    type(vector4_t), intent(in) :: p,q
    dphi = vector3_azimuthal_distance &
         (space_part (p), space_part (q))
  end function vector4_azimuthal_distance

\eatline
\nwendcode{}\nwbegindocs{188}The same in degrees:
\nwenddocs{}\nwbegincode{189}\moddef{Lorentz: public functions}\plusendmoddef
  public :: azimuthal_distance_deg
\nwendcode{}\nwbegincode{190}\moddef{Lorentz: interfaces}\plusendmoddef
  interface azimuthal_distance_deg
     module procedure vector3_azimuthal_distance_deg
     module procedure vector4_azimuthal_distance_deg
  end interface
\nwendcode{}\nwbegincode{191}\moddef{Lorentz: procedures}\plusendmoddef
  elemental function vector3_azimuthal_distance_deg (p, q) result (dphi)
    real(default) :: dphi
    type(vector3_t), intent(in) :: p,q
    dphi = vector3_azimuthal_distance (p, q) / degree
  end function vector3_azimuthal_distance_deg
  elemental function vector4_azimuthal_distance_deg (p, q) result (dphi)
    real(default) :: dphi
    type(vector4_t), intent(in) :: p,q
    dphi = vector4_azimuthal_distance (p, q) / degree
  end function vector4_azimuthal_distance_deg

\eatline
\nwendcode{}\nwbegindocs{192}The polar angle is defined $0\leq\theta\leq\pi$.  Note that
{\tt{}ATAN2} has the reversed order of arguments: {\tt{}ATAN2(Y,X)}.  Here,
$x$ is the 3-component while $y$ is the transverse momentum which is
always nonnegative.  Therefore, the result is nonnegative as well.
\nwenddocs{}\nwbegincode{193}\moddef{Lorentz: public functions}\plusendmoddef
  public :: polar_angle
\nwendcode{}\nwbegincode{194}\moddef{Lorentz: interfaces}\plusendmoddef
  interface polar_angle
     module procedure polar_angle_vector3
     module procedure polar_angle_vector4
  end interface
\nwendcode{}\nwbegincode{195}\moddef{Lorentz: procedures}\plusendmoddef
  elemental function polar_angle_vector3 (p) result (theta)
    real(default) :: theta
    type(vector3_t), intent(in) :: p
    if (any (abs (p%p) > 0)) then
       theta = atan2 (sqrt(p%p(1)**2 + p%p(2)**2), p%p(3))
    else
       theta = 0
    end if
  end function polar_angle_vector3
  elemental function polar_angle_vector4 (p) result (theta)
    real(default) :: theta
    type(vector4_t), intent(in) :: p
    theta = polar_angle (space_part (p))
  end function polar_angle_vector4

\eatline
\nwendcode{}\nwbegindocs{196}This is the cosine of the polar angle: $-1\leq\cos\theta\leq 1$.
\nwenddocs{}\nwbegincode{197}\moddef{Lorentz: public functions}\plusendmoddef
  public :: polar_angle_ct
\nwendcode{}\nwbegincode{198}\moddef{Lorentz: interfaces}\plusendmoddef
  interface polar_angle_ct
     module procedure polar_angle_ct_vector3
     module procedure polar_angle_ct_vector4
  end interface
\nwendcode{}\nwbegincode{199}\moddef{Lorentz: procedures}\plusendmoddef
  elemental function polar_angle_ct_vector3 (p) result (ct)
    real(default) :: ct
    type(vector3_t), intent(in) :: p
    if (any (abs (p%p) > 0)) then
       ct = p%p(3) / p**1
    else
       ct = 1
    end if
  end function polar_angle_ct_vector3
  elemental function polar_angle_ct_vector4 (p) result (ct)
    real(default) :: ct
    type(vector4_t), intent(in) :: p
    ct = polar_angle_ct (space_part (p))
  end function polar_angle_ct_vector4

\eatline
\nwendcode{}\nwbegindocs{200}The polar angle in degrees.
\nwenddocs{}\nwbegincode{201}\moddef{Lorentz: public functions}\plusendmoddef
  public :: polar_angle_deg
\nwendcode{}\nwbegincode{202}\moddef{Lorentz: interfaces}\plusendmoddef
  interface polar_angle_deg
     module procedure polar_angle_deg_vector3
     module procedure polar_angle_deg_vector4
  end interface
\nwendcode{}\nwbegincode{203}\moddef{Lorentz: procedures}\plusendmoddef
  elemental function polar_angle_deg_vector3 (p) result (theta)
    real(default) :: theta
    type(vector3_t), intent(in) :: p
    theta = polar_angle (p) / degree
  end function polar_angle_deg_vector3
  elemental function polar_angle_deg_vector4 (p) result (theta)
    real(default) :: theta
    type(vector4_t), intent(in) :: p
    theta = polar_angle (p) / degree
  end function polar_angle_deg_vector4

\eatline
\nwendcode{}\nwbegindocs{204}This is the angle enclosed between two three-momenta.  If one of the
momenta is zero, we return an angle of zero.  The range of the result
is $0\leq\theta\leq\pi$.  If there is only one argument, take the
positive $z$ axis as reference.
\nwenddocs{}\nwbegincode{205}\moddef{Lorentz: public functions}\plusendmoddef
  public :: enclosed_angle
\nwendcode{}\nwbegincode{206}\moddef{Lorentz: interfaces}\plusendmoddef
  interface enclosed_angle
     module procedure enclosed_angle_vector3
     module procedure enclosed_angle_vector4
  end interface
\nwendcode{}\nwbegincode{207}\moddef{Lorentz: procedures}\plusendmoddef
  elemental function enclosed_angle_vector3 (p, q) result (theta)
    real(default) :: theta
    type(vector3_t), intent(in) :: p, q
    theta = acos (enclosed_angle_ct (p, q))
  end function enclosed_angle_vector3
  elemental function enclosed_angle_vector4 (p, q) result (theta)
    real(default) :: theta
    type(vector4_t), intent(in) :: p, q
    theta = enclosed_angle (space_part (p), space_part (q))
  end function enclosed_angle_vector4

\eatline
\nwendcode{}\nwbegindocs{208}The cosine of the enclosed angle.
\nwenddocs{}\nwbegincode{209}\moddef{Lorentz: public functions}\plusendmoddef
  public :: enclosed_angle_ct
\nwendcode{}\nwbegincode{210}\moddef{Lorentz: interfaces}\plusendmoddef
  interface enclosed_angle_ct
     module procedure enclosed_angle_ct_vector3
     module procedure enclosed_angle_ct_vector4
  end interface
\nwendcode{}\nwbegincode{211}\moddef{Lorentz: procedures}\plusendmoddef
  elemental function enclosed_angle_ct_vector3 (p, q) result (ct)
    real(default) :: ct
    type(vector3_t), intent(in) :: p, q
    if (any (abs (p%p) > 0) .and. any (abs (q%p) > 0)) then
       ct = p*q / (p**1 * q**1)
       if (ct>1) then
          ct = 1
       else if (ct<-1) then
          ct = -1
       end if
    else
       ct = 1
    end if
  end function enclosed_angle_ct_vector3
  elemental function enclosed_angle_ct_vector4 (p, q) result (ct)
    real(default) :: ct
    type(vector4_t), intent(in) :: p, q
    ct = enclosed_angle_ct (space_part (p), space_part (q))
  end function enclosed_angle_ct_vector4

\eatline
\nwendcode{}\nwbegindocs{212}The enclosed angle in degrees.
\nwenddocs{}\nwbegincode{213}\moddef{Lorentz: public functions}\plusendmoddef
  public :: enclosed_angle_deg
\nwendcode{}\nwbegincode{214}\moddef{Lorentz: interfaces}\plusendmoddef
  interface enclosed_angle_deg
     module procedure enclosed_angle_deg_vector3
     module procedure enclosed_angle_deg_vector4
  end interface
\nwendcode{}\nwbegincode{215}\moddef{Lorentz: procedures}\plusendmoddef
  elemental function enclosed_angle_deg_vector3 (p, q) result (theta)
    real(default) :: theta
    type(vector3_t), intent(in) :: p, q
    theta = enclosed_angle (p, q) / degree
  end function enclosed_angle_deg_vector3
  elemental function enclosed_angle_deg_vector4 (p, q) result (theta)
    real(default) :: theta
    type(vector4_t), intent(in) :: p, q
    theta = enclosed_angle (p, q) / degree
  end function enclosed_angle_deg_vector4

\eatline
\nwendcode{}\nwbegindocs{216}The polar angle of the first momentum w.r.t.\ the second momentum,
evaluated in the rest frame of the second momentum.  If the second
four-momentum is not timelike, return zero.
\nwenddocs{}\nwbegincode{217}\moddef{Lorentz: public functions}\plusendmoddef
  public :: enclosed_angle_rest_frame
  public :: enclosed_angle_ct_rest_frame
  public :: enclosed_angle_deg_rest_frame
\nwendcode{}\nwbegincode{218}\moddef{Lorentz: interfaces}\plusendmoddef
  interface enclosed_angle_rest_frame
     module procedure enclosed_angle_rest_frame_vector4
  end interface
  interface enclosed_angle_ct_rest_frame
     module procedure enclosed_angle_ct_rest_frame_vector4
  end interface
  interface enclosed_angle_deg_rest_frame
     module procedure enclosed_angle_deg_rest_frame_vector4
  end interface
\nwendcode{}\nwbegincode{219}\moddef{Lorentz: procedures}\plusendmoddef
  elemental function enclosed_angle_rest_frame_vector4 (p, q) result (theta)
    type(vector4_t), intent(in) :: p, q
    real(default) :: theta
    theta = acos (enclosed_angle_ct_rest_frame (p, q))
  end function enclosed_angle_rest_frame_vector4
  elemental function enclosed_angle_ct_rest_frame_vector4 (p, q) result (ct)
    type(vector4_t), intent(in) :: p, q
    real(default) :: ct
    if (invariant_mass(q) > 0) then
       ct = enclosed_angle_ct ( &
            space_part (boost(-q, invariant_mass (q)) * p), &
            space_part (q))
    else
       ct = 1
    end if
  end function enclosed_angle_ct_rest_frame_vector4
  elemental function enclosed_angle_deg_rest_frame_vector4 (p, q) &
       result (theta)
    type(vector4_t), intent(in) :: p, q
    real(default) :: theta
    theta = enclosed_angle_rest_frame (p, q) / degree
  end function enclosed_angle_deg_rest_frame_vector4

\eatline
\eatline
\eatline
\nwendcode{}\nwbegindocs{220}\nwdocspar
\subsection{More kinematical functions (some redundant)}
The scalar transverse momentum (assuming the $z$ axis is longitudinal)
\nwenddocs{}\nwbegincode{221}\moddef{Lorentz: public functions}\plusendmoddef
  public :: transverse_part
\nwendcode{}\nwbegincode{222}\moddef{Lorentz: interfaces}\plusendmoddef
  interface transverse_part
     module procedure transverse_part_vector4
  end interface
\nwendcode{}\nwbegincode{223}\moddef{Lorentz: procedures}\plusendmoddef
  elemental function transverse_part_vector4 (p) result (pT)
    real(default) :: pT
    type(vector4_t), intent(in) :: p
    pT = sqrt(p%p(1)**2 + p%p(2)**2)
  end function transverse_part_vector4

\eatline
\nwendcode{}\nwbegindocs{224}The scalar longitudinal momentum (assuming the $z$ axis is
longitudinal).  Identical to {\tt{}momentum{\char95}z{\char95}component}.
\nwenddocs{}\nwbegincode{225}\moddef{Lorentz: public functions}\plusendmoddef
  public :: longitudinal_part
\nwendcode{}\nwbegincode{226}\moddef{Lorentz: interfaces}\plusendmoddef
  interface longitudinal_part
     module procedure longitudinal_part_vector4
  end interface
\nwendcode{}\nwbegincode{227}\moddef{Lorentz: procedures}\plusendmoddef
  elemental function longitudinal_part_vector4 (p) result (pL)
    real(default) :: pL
    type(vector4_t), intent(in) :: p
    pL = p%p(3)
  end function longitudinal_part_vector4

\eatline
\nwendcode{}\nwbegindocs{228}Absolute value of three-momentum
\nwenddocs{}\nwbegincode{229}\moddef{Lorentz: public functions}\plusendmoddef
  public :: space_part_norm
\nwendcode{}\nwbegincode{230}\moddef{Lorentz: interfaces}\plusendmoddef
  interface space_part_norm
     module procedure space_part_norm_vector4
  end interface
\nwendcode{}\nwbegincode{231}\moddef{Lorentz: procedures}\plusendmoddef
  elemental function space_part_norm_vector4 (p) result (p3)
    real(default) :: p3
    type(vector4_t), intent(in) :: p
    p3 = sqrt (p%p(1)**2 + p%p(2)**2 + p%p(3)**2)
  end function space_part_norm_vector4

\eatline
\nwendcode{}\nwbegindocs{232}The energy (the zeroth component)
\nwenddocs{}\nwbegincode{233}\moddef{Lorentz: public functions}\plusendmoddef
  public :: energy
\nwendcode{}\nwbegincode{234}\moddef{Lorentz: interfaces}\plusendmoddef
  interface energy
     module procedure energy_vector4
     module procedure energy_vector3
     module procedure energy_real
  end interface
\nwendcode{}\nwbegincode{235}\moddef{Lorentz: procedures}\plusendmoddef
  elemental function energy_vector4 (p) result (E)
    real(default) :: E
    type(vector4_t), intent(in) :: p
    E = p%p(0)
  end function energy_vector4

\nwendcode{}\nwbegindocs{236}Alternative: The energy corresponding to a given momentum and mass.
If the mass is omitted, it is zero
\nwenddocs{}\nwbegincode{237}\moddef{Lorentz: procedures}\plusendmoddef
  elemental function energy_vector3 (p, mass) result (E)
    real(default) :: E
    type(vector3_t), intent(in) :: p
    real(default), intent(in), optional :: mass
    if (present (mass)) then
       E = sqrt (p**2 + mass**2)
    else
       E = p**1
    end if
  end function energy_vector3

  elemental function energy_real (p, mass) result (E)
    real(default) :: E
    real(default), intent(in) :: p
    real(default), intent(in), optional :: mass
    if (present (mass)) then
       E = sqrt (p**2 + mass**2)
    else
       E = abs (p)
    end if
  end function energy_real

\eatline
\nwendcode{}\nwbegindocs{238}The invariant mass of four-momenta.  Zero for lightlike, negative for
spacelike momenta.
\nwenddocs{}\nwbegincode{239}\moddef{Lorentz: public functions}\plusendmoddef
  public :: invariant_mass
\nwendcode{}\nwbegincode{240}\moddef{Lorentz: interfaces}\plusendmoddef
  interface invariant_mass
     module procedure invariant_mass_vector4
  end interface
\nwendcode{}\nwbegincode{241}\moddef{Lorentz: procedures}\plusendmoddef
  elemental function invariant_mass_vector4 (p) result (m)
    real(default) :: m
    type(vector4_t), intent(in) :: p
    real(default) :: msq
    msq = p*p
    if (msq >= 0) then
       m = sqrt (msq)
    else
       m = - sqrt (abs (msq))
    end if
  end function invariant_mass_vector4
\eatline
\nwendcode{}\nwbegindocs{242}\nwdocspar
\nwenddocs{}\nwbegindocs{243}The invariant mass squared.  Zero for lightlike, negative for
spacelike momenta.
\nwenddocs{}\nwbegincode{244}\moddef{Lorentz: public functions}\plusendmoddef
  public :: invariant_mass_squared
\nwendcode{}\nwbegincode{245}\moddef{Lorentz: interfaces}\plusendmoddef
  interface invariant_mass_squared
     module procedure invariant_mass_squared_vector4
  end interface
\nwendcode{}\nwbegincode{246}\moddef{Lorentz: procedures}\plusendmoddef
  elemental function invariant_mass_squared_vector4 (p) result (msq)
    real(default) :: msq
    type(vector4_t), intent(in) :: p
    msq = p*p
  end function invariant_mass_squared_vector4

\eatline
\nwendcode{}\nwbegindocs{247}The transverse mass.  If the mass squared is negative, this value
also is negative.
\nwenddocs{}\nwbegincode{248}\moddef{Lorentz: public functions}\plusendmoddef
  public :: transverse_mass
\nwendcode{}\nwbegincode{249}\moddef{Lorentz: interfaces}\plusendmoddef
  interface transverse_mass
     module procedure transverse_mass_vector4
  end interface
\nwendcode{}\nwbegincode{250}\moddef{Lorentz: procedures}\plusendmoddef
  elemental function transverse_mass_vector4 (p) result (m)
    real(default) :: m
    type(vector4_t), intent(in) :: p
    real(default) :: msq
    msq = p%p(0)**2 - p%p(1)**2 - p%p(2)**2
    if (msq >= 0) then
       m = sqrt (msq)
    else
       m = - sqrt (abs (msq))
    end if
  end function transverse_mass_vector4

\eatline
\nwendcode{}\nwbegindocs{251}The rapidity (defined if particle is massive or $p_\perp>0$)
\nwenddocs{}\nwbegincode{252}\moddef{Lorentz: public functions}\plusendmoddef
  public :: rapidity
\nwendcode{}\nwbegincode{253}\moddef{Lorentz: interfaces}\plusendmoddef
  interface rapidity
     module procedure rapidity_vector4
  end interface
\nwendcode{}\nwbegincode{254}\moddef{Lorentz: procedures}\plusendmoddef
  elemental function rapidity_vector4 (p) result (y)
    real(default) :: y
    type(vector4_t), intent(in) :: p
    y = .5 * log( (energy (p) + longitudinal_part (p)) &
         &       /(energy (p) - longitudinal_part (p)))
  end function rapidity_vector4

\eatline
\nwendcode{}\nwbegindocs{255}The pseudorapidity (defined if $p_\perp>0$)
\nwenddocs{}\nwbegincode{256}\moddef{Lorentz: public functions}\plusendmoddef
  public :: pseudorapidity
\nwendcode{}\nwbegincode{257}\moddef{Lorentz: interfaces}\plusendmoddef
  interface pseudorapidity
     module procedure pseudorapidity_vector4
  end interface
\nwendcode{}\nwbegincode{258}\moddef{Lorentz: procedures}\plusendmoddef
  elemental function pseudorapidity_vector4 (p) result (eta)
    real(default) :: eta
    type(vector4_t), intent(in) :: p
    eta = -log( tan (.5 * polar_angle (p)))
  end function pseudorapidity_vector4

\eatline
\nwendcode{}\nwbegindocs{259}The rapidity distance (defined if both $p_\perp>0$)
\nwenddocs{}\nwbegincode{260}\moddef{Lorentz: public functions}\plusendmoddef
  public :: rapidity_distance
\nwendcode{}\nwbegincode{261}\moddef{Lorentz: interfaces}\plusendmoddef
  interface rapidity_distance
     module procedure rapidity_distance_vector4
  end interface
\nwendcode{}\nwbegincode{262}\moddef{Lorentz: procedures}\plusendmoddef
  elemental function rapidity_distance_vector4 (p, q) result (dy)
    type(vector4_t), intent(in) :: p, q
    real(default) :: dy
    dy = rapidity (q) - rapidity (p)
  end function rapidity_distance_vector4

\eatline
\nwendcode{}\nwbegindocs{263}The pseudorapidity distance (defined if both $p_\perp>0$)
\nwenddocs{}\nwbegincode{264}\moddef{Lorentz: public functions}\plusendmoddef
  public :: pseudorapidity_distance
\nwendcode{}\nwbegincode{265}\moddef{Lorentz: interfaces}\plusendmoddef
  interface pseudorapidity_distance
     module procedure pseudorapidity_distance_vector4
  end interface
\nwendcode{}\nwbegincode{266}\moddef{Lorentz: procedures}\plusendmoddef
  elemental function pseudorapidity_distance_vector4 (p, q) result (deta)
    real(default) :: deta
    type(vector4_t), intent(in) :: p, q
    deta = pseudorapidity (q) - pseudorapidity (p)
  end function pseudorapidity_distance_vector4

\eatline
\nwendcode{}\nwbegindocs{267}The distance on the $\eta-\phi$ cylinder:
\nwenddocs{}\nwbegincode{268}\moddef{Lorentz: public functions}\plusendmoddef
  public :: eta_phi_distance
\nwendcode{}\nwbegincode{269}\moddef{Lorentz: interfaces}\plusendmoddef
  interface eta_phi_distance
     module procedure eta_phi_distance_vector4
  end interface
\nwendcode{}\nwbegincode{270}\moddef{Lorentz: procedures}\plusendmoddef
  elemental function eta_phi_distance_vector4 (p, q) result (dr)
    type(vector4_t), intent(in) :: p, q
    real(default) :: dr
    dr = sqrt ( &
         pseudorapidity_distance (p, q)**2 &
         + azimuthal_distance (p, q)**2)
  end function eta_phi_distance_vector4

\eatline
\nwendcode{}\nwbegindocs{271}\nwdocspar
\subsection{Lorentz transformations}
\nwenddocs{}\nwbegincode{272}\moddef{Lorentz: public}\plusendmoddef
  public :: lorentz_transformation_t
\nwendcode{}\nwbegincode{273}\moddef{Lorentz: types}\plusendmoddef
  type :: lorentz_transformation_t
     private
     real(default), dimension(0:3, 0:3) :: L
   contains
   \LA{}Lorentz: lorentz transformation: TBP\RA{}
  end type lorentz_transformation_t

\eatline
\nwendcode{}\nwbegindocs{274}Output:
\nwenddocs{}\nwbegincode{275}\moddef{Lorentz: public}\plusendmoddef
  public :: lorentz_transformation_write
\nwendcode{}\nwbegincode{276}\moddef{Lorentz: lorentz transformation: TBP}\endmoddef
  procedure :: write => lorentz_transformation_write
\nwendcode{}\nwbegincode{277}\moddef{Lorentz: procedures}\plusendmoddef
  subroutine lorentz_transformation_write (L, unit)
    class(lorentz_transformation_t), intent(in) :: L
    integer, intent(in), optional :: unit
    integer :: u
    integer :: i
    u = given_output_unit (unit);  if (u < 0)  return
    write (u, "(1x,A,3(1x," // FMT_19 // "))")  "L00 = ", L%L(0,0)
    write (u, "(1x,A,3(1x," // FMT_19 // "))")  "L0j = ", L%L(0,1:3)
    do i = 1, 3
       write (u, "(1x,A,I0,A,3(1x," // FMT_19 // "))")  &
            "L", i, "0 = ", L%L(i,0)
       write (u, "(1x,A,I0,A,3(1x," // FMT_19 // "))")  &
            "L", i, "j = ", L%L(i,1:3)
    end do
  end subroutine lorentz_transformation_write

\eatline
\nwendcode{}\nwbegindocs{278}Extract all components:
\nwenddocs{}\nwbegincode{279}\moddef{Lorentz: public}\plusendmoddef
  public :: lorentz_transformation_get_components
\nwendcode{}\nwbegincode{280}\moddef{Lorentz: procedures}\plusendmoddef
  pure function lorentz_transformation_get_components (L) result (a)
    type(lorentz_transformation_t), intent(in) :: L
    real(default), dimension(0:3,0:3) :: a
    a = L%L
  end function lorentz_transformation_get_components

\eatline
\nwendcode{}\nwbegindocs{281}\nwdocspar
\subsection{Functions of Lorentz transformations}
For the inverse, we make use of the fact that
$\Lambda^{\mu\nu}\Lambda_{\mu\rho}=\delta^\nu_\rho$.  So, lowering the
indices and transposing is sufficient.
\nwenddocs{}\nwbegincode{282}\moddef{Lorentz: public functions}\plusendmoddef
  public :: inverse
\nwendcode{}\nwbegincode{283}\moddef{Lorentz: interfaces}\plusendmoddef
  interface inverse
     module procedure lorentz_transformation_inverse
  end interface
\nwendcode{}\nwbegincode{284}\moddef{Lorentz: procedures}\plusendmoddef
  elemental function lorentz_transformation_inverse (L) result (IL)
    type(lorentz_transformation_t) :: IL
    type(lorentz_transformation_t), intent(in) :: L
    IL%L(0,0) = L%L(0,0)
    IL%L(0,1:) = -L%L(1:,0)
    IL%L(1:,0) = -L%L(0,1:)
    IL%L(1:,1:) = transpose(L%L(1:,1:))
  end function lorentz_transformation_inverse

\eatline
\eatline
\nwendcode{}\nwbegindocs{285}\nwdocspar
\subsection{Invariants}
These are used below.  The first array index is varying fastest in
{\tt{}FORTRAN}; therefore the extra minus in the odd-rank tensor
epsilon.
\nwenddocs{}\nwbegincode{286}\moddef{Lorentz: parameters}\plusendmoddef
  integer, dimension(3,3), parameter :: delta_three = &
       & reshape( source = [ 1,0,0, 0,1,0, 0,0,1 ], &
       &          shape  = [3,3] )
  integer, dimension(3,3,3), parameter :: epsilon_three = &
       & reshape( source = [ 0, 0,0,  0,0,-1,   0,1,0, &
       &                     0, 0,1,  0,0, 0,  -1,0,0, &
       &                     0,-1,0,  1,0, 0,   0,0,0 ],&
       &          shape = [3,3,3] )
\eatline
\nwendcode{}\nwbegindocs{287}This could be of some use:
\nwenddocs{}\nwbegincode{288}\moddef{Lorentz: public}\plusendmoddef
  public :: identity
\nwendcode{}\nwbegincode{289}\moddef{Lorentz: parameters}\plusendmoddef
  type(lorentz_transformation_t), parameter :: &
       & identity = &
       & lorentz_transformation_t ( &
       & reshape( source = [ one, zero, zero, zero, &
       &                     zero, one, zero, zero, &
       &                     zero, zero, one, zero, &
       &                     zero, zero, zero, one ],&
       &          shape = [4,4] ) )
\eatline
\nwendcode{}\nwbegincode{290}\moddef{Lorentz: public}\plusendmoddef
  public :: space_reflection
\nwendcode{}\nwbegincode{291}\moddef{Lorentz: parameters}\plusendmoddef
  type(lorentz_transformation_t), parameter :: &
       & space_reflection = &
       & lorentz_transformation_t ( &
       & reshape( source = [ one, zero, zero, zero, &
       &                     zero,-one, zero, zero, &
       &                     zero, zero,-one, zero, &
       &                     zero, zero, zero,-one ],&
       &          shape = [4,4] ) )
\eatline
\nwendcode{}\nwbegindocs{292}Builds a unit vector orthogal to the input vector in the xy-plane.
\nwenddocs{}\nwbegincode{293}\moddef{Lorentz: public functions}\plusendmoddef
  public :: create_orthogonal
\nwendcode{}\nwbegincode{294}\moddef{Lorentz: procedures}\plusendmoddef
  function create_orthogonal (p_in) result (p_out)
    type(vector3_t), intent(in) :: p_in
    type(vector3_t) :: p_out
    real(default) :: ab
    ab = sqrt (p_in%p(1)**2 + p_in%p(2)**2)
    if (abs (ab) < eps0) then
      p_out%p(1) = 1
      p_out%p(2) = 0
      p_out%p(3) = 0
    else
      p_out%p(1) = p_in%p(2)
      p_out%p(2) = -p_in%p(1)
      p_out%p(3) = 0
      p_out = p_out / ab
    end if
  end function create_orthogonal

\eatline
\nwendcode{}\nwbegindocs{295}\nwdocspar
\nwenddocs{}\nwbegincode{296}\moddef{Lorentz: public functions}\plusendmoddef
  public :: create_unit_vector
\nwendcode{}\nwbegincode{297}\moddef{Lorentz: procedures}\plusendmoddef
  function create_unit_vector (p_in) result (p_out)
    type(vector4_t), intent(in) :: p_in
    type(vector3_t) :: p_out
    real(default) :: abs
    abs = space_part_norm (p_in)
    p_out%p(1) = p_in%p(1)/abs
    p_out%p(2) = p_in%p(2)/abs
    p_out%p(3) = p_in%p(3)/abs
  end function create_unit_vector

\eatline
\nwendcode{}\nwbegindocs{298}\nwdocspar
\nwenddocs{}\nwbegincode{299}\moddef{Lorentz: public functions}\plusendmoddef
  public :: normalize
\nwendcode{}\nwbegincode{300}\moddef{Lorentz: procedures}\plusendmoddef
  subroutine normalize(p)
    type(vector3_t), intent(inout) :: p
    real(default) :: abs
    abs = sqrt (p%p(1)**2 + p%p(2)**2 + p%p(3)**2)
    p = p/abs
  end subroutine normalize

\eatline
\nwendcode{}\nwbegindocs{301}\nwdocspar
\subsection{Boosts}
We build Lorentz transformations from boosts and rotations.  In both
cases we can supply a three-vector which defines the axis and (hyperbolic)
angle.  For a boost, this is the vector $\vec\beta=\vec p/E$,
such that a particle at rest with mass $m$ is boosted to a particle
with three-vector $\vec p$.  Here, we have
\begin{equation}
  \beta = \tanh\chi = p/E, \qquad
  \gamma = \cosh\chi = E/m, \qquad
  \beta\gamma = \sinh\chi = p/m
\end{equation}
\nwenddocs{}\nwbegincode{302}\moddef{Lorentz: public functions}\plusendmoddef
  public :: boost
\nwendcode{}\nwbegincode{303}\moddef{Lorentz: interfaces}\plusendmoddef
  interface boost
     module procedure boost_from_rest_frame
     module procedure boost_from_rest_frame_vector3
     module procedure boost_generic
     module procedure boost_canonical
  end interface
\eatline
\nwendcode{}\nwbegindocs{304}In the first form, the argument is some four-momentum, the space
part of which determines a direction, and the associated mass (which
is not checked against the four-momentum).  The boost vector
$\gamma\vec\beta$ is then given by $\vec p/m$.  This boosts from the
rest frame of a particle to the current frame.  To be explicit, if
$\vec p$ is the momentum of a particle and $m$ its mass, $L(\vec p/m)$
is the transformation that turns $(m;\vec 0)$ into $(E;\vec p)$.
Conversely, the inverse transformation boosts a vector \emph{into} the
rest frame of a particle, in particular $(E;\vec p)$ into $(m;\vec
0)$.
\nwenddocs{}\nwbegincode{305}\moddef{Lorentz: procedures}\plusendmoddef
  elemental function boost_from_rest_frame (p, m) result (L)
    type(lorentz_transformation_t) :: L
    type(vector4_t), intent(in) :: p
    real(default), intent(in) :: m
    L = boost_from_rest_frame_vector3 (space_part (p), m)
  end function boost_from_rest_frame
  elemental function boost_from_rest_frame_vector3 (p, m) result (L)
    type(lorentz_transformation_t) :: L
    type(vector3_t), intent(in) :: p
    real(default), intent(in) :: m
    type(vector3_t) :: beta_gamma
    real(default) :: bg2, g, c
    integer :: i,j
    if (m > eps0) then
       beta_gamma = p / m
       bg2 = beta_gamma**2
    else
       bg2 = 0
    end if
    if (bg2 > eps0) then
       g = sqrt(1 + bg2);  c = (g-1)/bg2
       L%L(0,0)  = g
       L%L(0,1:) = beta_gamma%p
       L%L(1:,0) = L%L(0,1:)
       do i=1,3
          do j=1,3
             L%L(i,j) = delta_three(i,j) + c*beta_gamma%p(i)*beta_gamma%p(j)
          end do
       end do
    else
       L = identity
    end if
  end function boost_from_rest_frame_vector3
\eatline
\nwendcode{}\nwbegindocs{306}A canonical boost is a boost along one of the coordinate axes, which
we may supply as an integer argument.  Here, $\gamma\beta$ is scalar.
\nwenddocs{}\nwbegincode{307}\moddef{Lorentz: procedures}\plusendmoddef
  elemental function boost_canonical (beta_gamma, k) result (L)
    type(lorentz_transformation_t) :: L
    real(default), intent(in) :: beta_gamma
    integer, intent(in) :: k
    real(default) :: g
    g = sqrt(1 + beta_gamma**2)
    L = identity
    L%L(0,0) = g
    L%L(0,k) = beta_gamma
    L%L(k,0) = L%L(0,k)
    L%L(k,k) = L%L(0,0)
  end function boost_canonical
\eatline
\nwendcode{}\nwbegindocs{308}Instead of a canonical axis, we can supply an arbitrary axis which
need not be normalized.  If it is zero, return the unit matrix.
\nwenddocs{}\nwbegincode{309}\moddef{Lorentz: procedures}\plusendmoddef
  elemental function boost_generic (beta_gamma, axis) result (L)
    type(lorentz_transformation_t) :: L
    real(default), intent(in) :: beta_gamma
    type(vector3_t), intent(in) :: axis
    if (any (abs (axis%p) > 0)) then
       L = boost_from_rest_frame_vector3 (beta_gamma * axis, axis**1)
    else
       L = identity
    end if
  end function boost_generic

\eatline
\nwendcode{}\nwbegindocs{310}\nwdocspar
\subsection{Rotations}
For a rotation, the vector defines the rotation axis, and its length
the rotation angle.
\nwenddocs{}\nwbegincode{311}\moddef{Lorentz: public functions}\plusendmoddef
  public :: rotation
\nwendcode{}\nwbegincode{312}\moddef{Lorentz: interfaces}\plusendmoddef
  interface rotation
     module procedure rotation_generic
     module procedure rotation_canonical
     module procedure rotation_generic_cs
     module procedure rotation_canonical_cs
  end interface
\eatline
\nwendcode{}\nwbegindocs{313}If $\cos\phi$ and $\sin\phi$ is already known, we do not have to
calculate them.  Of course, the user has to ensure that
$\cos^2\phi+\sin^2\phi=1$, and that the given axis {\tt{}n} is normalized to
one.  In the second form, the length of {\tt{}axis} is the rotation
angle.
\nwenddocs{}\nwbegincode{314}\moddef{Lorentz: procedures}\plusendmoddef
  elemental function rotation_generic_cs (cp, sp, axis) result (R)
    type(lorentz_transformation_t) :: R
    real(default), intent(in) :: cp, sp
    type(vector3_t), intent(in) :: axis
    integer :: i,j
    R = identity
    do i=1,3
       do j=1,3
          R%L(i,j) = cp*delta_three(i,j) + (1-cp)*axis%p(i)*axis%p(j)  &
               &   - sp*dot_product(epsilon_three(i,j,:), axis%p)
       end do
    end do
  end function rotation_generic_cs
  elemental function rotation_generic (axis) result (R)
    type(lorentz_transformation_t) :: R
    type(vector3_t), intent(in) :: axis
    real(default) :: phi
    if (any (abs(axis%p) > 0)) then
       phi = abs(axis**1)
       R = rotation_generic_cs (cos(phi), sin(phi), axis/phi)
    else
       R = identity
    end if
  end function rotation_generic
\eatline
\nwendcode{}\nwbegindocs{315}Alternatively, give just the angle and label the coordinate axis by
an integer.
\nwenddocs{}\nwbegincode{316}\moddef{Lorentz: procedures}\plusendmoddef
  elemental function rotation_canonical_cs (cp, sp, k) result (R)
    type(lorentz_transformation_t) :: R
    real(default), intent(in) :: cp, sp
    integer, intent(in) :: k
    integer :: i,j
    R = identity
    do i=1,3
       do j=1,3
          R%L(i,j) = -sp*epsilon_three(i,j,k)
       end do
       R%L(i,i) = cp
    end do
    R%L(k,k) = 1
  end function rotation_canonical_cs
  elemental function rotation_canonical (phi, k) result (R)
    type(lorentz_transformation_t) :: R
    real(default), intent(in) :: phi
    integer, intent(in) :: k
    R = rotation_canonical_cs(cos(phi), sin(phi), k)
  end function rotation_canonical
\eatline
\nwendcode{}\nwbegindocs{317}\nwdocspar
This is viewed as a method for the first argument (three-vector):
Reconstruct the rotation that rotates it into the second three-vector.
\nwenddocs{}\nwbegincode{318}\moddef{Lorentz: public functions}\plusendmoddef
  public :: rotation_to_2nd
\nwendcode{}\nwbegincode{319}\moddef{Lorentz: interfaces}\plusendmoddef
  interface rotation_to_2nd
     module procedure rotation_to_2nd_generic
     module procedure rotation_to_2nd_canonical
  end interface
\nwendcode{}\nwbegincode{320}\moddef{Lorentz: procedures}\plusendmoddef
  elemental function rotation_to_2nd_generic (p, q) result (R)
    type(lorentz_transformation_t) :: R
    type(vector3_t), intent(in) :: p, q
    type(vector3_t) :: a, b, ab
    real(default) :: ct, st
    if (any (abs (p%p) > 0) .and. any (abs (q%p) > 0)) then
       a = direction (p)
       b = direction (q)
       ab = cross_product(a,b)
       ct = a*b;  st = ab**1
       if (abs(st) > eps0) then
          R = rotation_generic_cs (ct, st, ab/st)
       else if (ct < 0) then
          R = space_reflection
       else
          R = identity
       end if
    else
       R = identity
    end if
  end function rotation_to_2nd_generic
\eatline
\nwendcode{}\nwbegindocs{321}\nwdocspar
The same for a canonical axis: The function returns the transformation that
rotates the $k$-axis into the direction of $p$.
\nwenddocs{}\nwbegincode{322}\moddef{Lorentz: procedures}\plusendmoddef
  elemental function rotation_to_2nd_canonical (k, p) result (R)
    type(lorentz_transformation_t) :: R
    integer, intent(in) :: k
    type(vector3_t), intent(in) :: p
    type(vector3_t) :: b, ab
    real(default) :: ct, st
    integer :: i, j
    if (any (abs (p%p) > 0)) then
       b = direction (p)
       ab%p = 0
       do i = 1, 3
          do j = 1, 3
             ab%p(j) = ab%p(j) + b%p(i) * epsilon_three(i,j,k)
          end do
       end do
       ct = b%p(k);  st = ab**1
       if (abs(st) > eps0) then
          R = rotation_generic_cs (ct, st, ab/st)
       else if (ct < 0) then
          R = space_reflection
       else
          R = identity
       end if
    else
       R = identity
    end if
  end function rotation_to_2nd_canonical

\eatline
\nwendcode{}\nwbegindocs{323}\nwdocspar
\subsection{Composite Lorentz transformations}
This function returns the transformation that, given a pair of vectors
$p_{1,2}$, (a) boosts from the rest frame of the c.m. system (with
invariant mass $m$) into the lab frame where $p_i$ are defined, and
(b) turns the given axis (or the canonical vectors $\pm
e_k$) in the rest frame into the directions of $p_{1,2}$ in the lab frame.
Note that the energy components are not used; for a
consistent result one should have $(p_1+p_2)^2 = m^2$.
\nwenddocs{}\nwbegincode{324}\moddef{Lorentz: public functions}\plusendmoddef
  public :: transformation
\nwendcode{}\nwbegincode{325}\moddef{Lorentz: interfaces}\plusendmoddef
  interface transformation
     module procedure transformation_rec_generic
     module procedure transformation_rec_canonical
  end interface
\eatline
\nwendcode{}\nwbegincode{326}\moddef{Lorentz: procedures}\plusendmoddef
  elemental function transformation_rec_generic (axis, p1, p2, m) result (L)
    type(vector3_t), intent(in) :: axis
    type(vector4_t), intent(in) :: p1, p2
    real(default), intent(in) :: m
    type(lorentz_transformation_t) :: L
    L = boost (p1 + p2, m)
    L = L * rotation_to_2nd (axis, space_part (inverse (L) * p1))
  end function transformation_rec_generic
  elemental function transformation_rec_canonical (k, p1, p2, m) result (L)
    integer, intent(in) :: k
    type(vector4_t), intent(in) :: p1, p2
    real(default), intent(in) :: m
    type(lorentz_transformation_t) :: L
    L = boost (p1 + p2, m)
    L = L * rotation_to_2nd (k, space_part (inverse (L) * p1))
  end function transformation_rec_canonical
\eatline
\nwendcode{}\nwbegindocs{327}\nwdocspar
\subsection{Applying Lorentz transformations}
Multiplying vectors and Lorentz transformations is straightforward.
\nwenddocs{}\nwbegincode{328}\moddef{Lorentz: interfaces}\plusendmoddef
  interface operator(*)
     module procedure prod_LT_vector4
     module procedure prod_LT_LT
     module procedure prod_vector4_LT
  end interface
\nwendcode{}\nwbegincode{329}\moddef{Lorentz: procedures}\plusendmoddef
  elemental function prod_LT_vector4 (L, p) result (np)
    type(vector4_t) :: np
    type(lorentz_transformation_t), intent(in) :: L
    type(vector4_t), intent(in) :: p
    np%p = matmul (L%L, p%p)
  end function prod_LT_vector4
  elemental function prod_LT_LT (L1, L2) result (NL)
    type(lorentz_transformation_t) :: NL
    type(lorentz_transformation_t), intent(in) :: L1,L2
    NL%L = matmul (L1%L, L2%L)
  end function prod_LT_LT
  elemental function prod_vector4_LT (p, L) result (np)
    type(vector4_t) :: np
    type(vector4_t), intent(in) :: p
    type(lorentz_transformation_t), intent(in) :: L
    np%p = matmul (p%p, L%L)
  end function prod_vector4_LT

\eatline
\nwendcode{}\nwbegindocs{330}\nwdocspar
\subsection{Special Lorentz transformations}
These routines have their application in the generation and extraction
of angles in the phase-space sampling routine.  Since this part of the
program is time-critical, we calculate the composition of
transformations directly instead of multiplying rotations and boosts.

This Lorentz transformation is the composition of a rotation by $\phi$
around the $3$ axis, a rotation by $\theta$ around the $2$ axis, and a
boost along the $3$ axis:
\begin{equation}
  L = B_3(\beta\gamma)\,R_2(\theta)\,R_3(\phi)
\end{equation}
Instead of the angles we provide sine and cosine.
\nwenddocs{}\nwbegincode{331}\moddef{Lorentz: public functions}\plusendmoddef
  public :: LT_compose_r3_r2_b3
\nwendcode{}\nwbegincode{332}\moddef{Lorentz: procedures}\plusendmoddef
  elemental function LT_compose_r3_r2_b3 &
       (cp, sp, ct, st, beta_gamma) result (L)
    type(lorentz_transformation_t) :: L
    real(default), intent(in) :: cp, sp, ct, st, beta_gamma
    real(default) :: gamma
    if (abs(beta_gamma) < eps0) then
       L%L(0,0)  = 1
       L%L(1:,0) = 0
       L%L(0,1:) = 0
       L%L(1,1:) = [  ct*cp, -ct*sp, st ]
       L%L(2,1:) = [     sp,     cp,  zero ]
       L%L(3,1:) = [ -st*cp,  st*sp, ct ]
    else
       gamma = sqrt(1 + beta_gamma**2)
       L%L(0,0)  = gamma
       L%L(1,0)  = 0
       L%L(2,0)  = 0
       L%L(3,0)  = beta_gamma
       L%L(0,1:) = beta_gamma * [ -st*cp,  st*sp, ct ]
       L%L(1,1:) =              [  ct*cp, -ct*sp, st ]
       L%L(2,1:) =              [     sp,     cp, zero ]
       L%L(3,1:) = gamma      * [ -st*cp,  st*sp, ct ]
    end if
  end function LT_compose_r3_r2_b3

\eatline
\nwendcode{}\nwbegindocs{333}Different ordering:
\begin{equation}
  L = B_3(\beta\gamma)\,R_3(\phi)\,R_2(\theta)
\end{equation}
\nwenddocs{}\nwbegincode{334}\moddef{Lorentz: public functions}\plusendmoddef
  public :: LT_compose_r2_r3_b3
\nwendcode{}\nwbegincode{335}\moddef{Lorentz: procedures}\plusendmoddef
  elemental function LT_compose_r2_r3_b3 &
       (ct, st, cp, sp, beta_gamma) result (L)
    type(lorentz_transformation_t) :: L
    real(default), intent(in) :: ct, st, cp, sp, beta_gamma
    real(default) :: gamma
    if (abs(beta_gamma) < eps0) then
       L%L(0,0)  = 1
       L%L(1:,0) = 0
       L%L(0,1:) = 0
       L%L(1,1:) = [  ct*cp,    -sp,     st*cp ]
       L%L(2,1:) = [  ct*sp,     cp,     st*sp ]
       L%L(3,1:) = [ -st   ,   zero,     ct    ]
    else
       gamma = sqrt(1 + beta_gamma**2)
       L%L(0,0)  = gamma
       L%L(1,0)  = 0
       L%L(2,0)  = 0
       L%L(3,0)  = beta_gamma
       L%L(0,1:) = beta_gamma * [ -st   ,   zero,     ct    ]
       L%L(1,1:) =              [  ct*cp,    -sp,     st*cp ]
       L%L(2,1:) =              [  ct*sp,     cp,     st*sp ]
       L%L(3,1:) = gamma      * [ -st   ,   zero,     ct    ]
    end if
  end function LT_compose_r2_r3_b3

\eatline
\nwendcode{}\nwbegindocs{336}This function returns the previous Lorentz transformation applied to
an arbitrary four-momentum and extracts the space part of the result:
\begin{equation}
  \vec n = [B_3(\beta\gamma)\,R_2(\theta)\,R_3(\phi)\,p]_{\rm space\ part}
\end{equation}
The second variant applies if there is no rotation
\nwenddocs{}\nwbegincode{337}\moddef{Lorentz: public functions}\plusendmoddef
  public :: axis_from_p_r3_r2_b3, axis_from_p_b3
\nwendcode{}\nwbegincode{338}\moddef{Lorentz: procedures}\plusendmoddef
  elemental function axis_from_p_r3_r2_b3 &
       (p, cp, sp, ct, st, beta_gamma) result (n)
    type(vector3_t) :: n
    type(vector4_t), intent(in) :: p
    real(default), intent(in) :: cp, sp, ct, st, beta_gamma
    real(default) :: gamma, px, py
    px = cp * p%p(1) - sp * p%p(2)
    py = sp * p%p(1) + cp * p%p(2)
    n%p(1) =  ct * px + st * p%p(3)
    n%p(2) = py
    n%p(3) = -st * px + ct * p%p(3)
    if (abs(beta_gamma) > eps0) then
       gamma = sqrt(1 + beta_gamma**2)
       n%p(3) = n%p(3) * gamma + p%p(0) * beta_gamma
    end if
  end function axis_from_p_r3_r2_b3

  elemental function axis_from_p_b3 (p, beta_gamma) result (n)
    type(vector3_t) :: n
    type(vector4_t), intent(in) :: p
    real(default), intent(in) :: beta_gamma
    real(default) :: gamma
    n%p = p%p(1:3)
    if (abs(beta_gamma) > eps0) then
       gamma = sqrt(1 + beta_gamma**2)
       n%p(3) = n%p(3) * gamma + p%p(0) * beta_gamma
    end if
  end function axis_from_p_b3

\eatline
\nwendcode{}\nwbegindocs{339}\nwdocspar
\subsection{Special functions}
The standard phase space function:
\nwenddocs{}\nwbegincode{340}\moddef{Lorentz: public functions}\plusendmoddef
  public :: lambda
\nwendcode{}\nwbegincode{341}\moddef{Lorentz: procedures}\plusendmoddef
  elemental function lambda (m1sq, m2sq, m3sq)
    real(default) :: lambda
    real(default), intent(in) :: m1sq, m2sq, m3sq
    lambda = (m1sq - m2sq - m3sq)**2 - 4*m2sq*m3sq
  end function lambda

\eatline
\nwendcode{}\nwbegindocs{342}Return a pair of head-to-head colliding momenta, given the collider
energy, particle masses, and optionally the momentum of the
c.m. system.
\nwenddocs{}\nwbegincode{343}\moddef{Lorentz: public functions}\plusendmoddef
  public :: colliding_momenta
\nwendcode{}\nwbegincode{344}\moddef{Lorentz: procedures}\plusendmoddef
  function colliding_momenta (sqrts, m, p_cm) result (p)
    type(vector4_t), dimension(2) :: p
    real(default), intent(in) :: sqrts
    real(default), dimension(2), intent(in), optional :: m
    real(default), intent(in), optional :: p_cm
    real(default), dimension(2) :: dmsq
    real(default) :: ch, sh
    real(default), dimension(2) :: E0, p0
    integer, dimension(2), parameter :: sgn = [1, -1]
    if (abs(sqrts) < eps0) then
       call msg_fatal (" Colliding beams: sqrts is zero (please set sqrts)")
       p = vector4_null;  return
    else if (sqrts <= 0) then
       call msg_fatal (" Colliding beams: sqrts is negative")
       p = vector4_null;  return
    end if
    if (present (m)) then
       dmsq = sgn * (m(1)**2-m(2)**2)
       E0 = (sqrts + dmsq/sqrts) / 2
       if (any (E0 < m)) then
          call msg_fatal &
               (" Colliding beams: beam energy is less than particle mass")
          p = vector4_null;  return
       end if
       p0 = sgn * sqrt (E0**2 - m**2)
    else
       E0 = sqrts / 2
       p0 = sgn * E0
    end if
    if (present (p_cm)) then
       sh = p_cm / sqrts
       ch = sqrt (1 + sh**2)
       p = vector4_moving (E0 * ch + p0 * sh, E0 * sh + p0 * ch, 3)
    else
       p = vector4_moving (E0, p0, 3)
    end if
  end function colliding_momenta

\eatline
\nwendcode{}\nwbegindocs{345}This subroutine is for the purpose of numerical checks and
comparisons.  The idea is to set a number to zero if it is numerically
equivalent with zero.  The equivalence is established by comparing
with a {\tt{}tolerance} argument.  We implement this for vectors and
transformations.
\nwenddocs{}\nwbegincode{346}\moddef{Lorentz: public functions}\plusendmoddef
  public :: pacify
\nwendcode{}\nwbegincode{347}\moddef{Lorentz: interfaces}\plusendmoddef
  interface pacify
     module procedure pacify_vector3
     module procedure pacify_vector4
     module procedure pacify_LT
  end interface pacify

\nwendcode{}\nwbegincode{348}\moddef{Lorentz: procedures}\plusendmoddef
  elemental subroutine pacify_vector3 (p, tolerance)
    type(vector3_t), intent(inout) :: p
    real(default), intent(in) :: tolerance
    where (abs (p%p) < tolerance)  p%p = 0
  end subroutine pacify_vector3

  elemental subroutine pacify_vector4 (p, tolerance)
    type(vector4_t), intent(inout) :: p
    real(default), intent(in) :: tolerance
    where (abs (p%p) < tolerance)  p%p = 0
  end subroutine pacify_vector4

  elemental subroutine pacify_LT (LT, tolerance)
    type(lorentz_transformation_t), intent(inout) :: LT
    real(default), intent(in) :: tolerance
    where (abs (LT%L) < tolerance)  LT%L = 0
  end subroutine pacify_LT

\eatline
\nwendcode{}\nwbegindocs{349}\nwdocspar
\nwenddocs{}\nwbegincode{350}\moddef{Lorentz: public}\plusendmoddef
  public :: vector_set_reshuffle
\nwendcode{}\nwbegincode{351}\moddef{Lorentz: procedures}\plusendmoddef
  subroutine vector_set_reshuffle (p1, list, p2)
    type(vector4_t), intent(in), dimension(:), allocatable :: p1
    integer, intent(in), dimension(:), allocatable :: list
    type(vector4_t), intent(out), dimension(:), allocatable :: p2
    integer :: n, n_p
    n_p = size (p1)
    if (size (list) /= n_p) return
    allocate (p2 (n_p))
    do n = 1, n_p
      p2(n) = p1(list(n))
    end do
  end subroutine vector_set_reshuffle

\eatline
\nwendcode{}\nwbegindocs{352}\nwdocspar
\nwenddocs{}\nwbegincode{353}\moddef{Lorentz: public}\plusendmoddef
  public :: vector4_write_set
\nwendcode{}\nwbegincode{354}\moddef{Lorentz: procedures}\plusendmoddef
  subroutine vector4_write_set (p, unit, show_mass)
    type(vector4_t), intent(in), dimension(:) :: p
    integer, intent(in), optional :: unit
    logical, intent(in), optional :: show_mass
    integer :: i, j
    real(default), dimension(4) :: p_tot
    character(len=7) :: fmt
    integer :: u
    u = given_output_unit (unit);  if (u < 0)  return
    p_tot = 0
    do i = 1, size (p)
      forall (j=1:4) p_tot(j) = p_tot(j) + vector4_get_component(p(i),j-1)
      call vector4_write (p(i), u, show_mass)
    end do
    call pac_fmt (fmt, FMT_19, FMT_15)
    write(u, "(A5)") 'Total: '
    write(u, "(1x,A,1x," // fmt // ")") 'E = ', p_tot(1)
    write(u, "(1x,A,3(1x," // fmt // "))") 'P = ', p_tot(2:)
  end subroutine vector4_write_set

\eatline
\nwendcode{}\nwbegindocs{355}This computes the quantities
\begin{align*}
  \langle ij \rangle &= \sqrt{|S_{ij}|} e^{i\phi_{ij}},
  [ij] &= \sqrt{|S_{ij}|} e^{\i\tilde{\phi}_{ij}},
\end{align*}
with $S_{ij} = \left(p_i + p_j\right)^2$. The phase space factor
$\phi_{ij}$ is determined by
\begin{align*}
  \cos\phi_{ij} &= \frac{p_i^1p_j^+ - p_j^1p_i^+}{\sqrt{p_i^+p_j^+S_{ij}}},
  \sin\phi_{ij} &= \frac{p_i^2p_j^+ - p_j^2p_i^+}{\sqrt{p_i^+p_j^+S_{ij}}}.
\end{align*}
After $\langle ij \rangle$ has been computed according to these
formulae, $[ij]$ can be obtained by using the relation $S_{ij} =
\langle ij \rangle [ji]$ and taking into account that $[ij] =
-[ji]$. Thus, a minus-sign has to be applied.
\nwenddocs{}\nwbegincode{356}\moddef{Lorentz: public}\plusendmoddef
  public :: spinor_product
\nwendcode{}\nwbegincode{357}\moddef{Lorentz: procedures}\plusendmoddef
  subroutine spinor_product (p1, p2, prod1, prod2)
    type(vector4_t), intent(in) :: p1, p2
    complex(default), intent(out) :: prod1, prod2
    real(default) :: sij
    complex(default) :: phase
    real(default) :: pp_1, pp_2
    pp_1 = p1%p(0) + p1%p(3)
    pp_2 = p2%p(0) + p2%p(3)
    sij = (p1+p2)**2
    phase = cmplx ((p1%p(1)*pp_2 - p2%p(1)*pp_1)/sqrt (sij*pp_1*pp_2), &
                   (p1%p(2)*pp_2 - p2%p(2)*pp_1)/sqrt (sij*pp_1*pp_2), &
                    default)
    !!! <ij>
    prod1 = sqrt (sij) * phase
    !!! [ij]
    if (abs(prod1) > 0) then
       prod2 = - sij / prod1
    else
       prod2 = 0
    end if
  end subroutine spinor_product

\eatline
\nwendcode{}\nwbegindocs{358}\nwdocspar
\clearpage
%%%%%%%%%%%%%%%%%%%%%%%%%%%%%%%%%%%%%%%%%%%%%%%%%%%%%%%%%%%%%%%%%%%%%%%%
\section{Special Physics functions}
Here, we declare functions that are specific for the Standard Model,
including QCD: fixed and running $\alpha_s$, Catani-Seymour
dipole terms, loop functions, etc.

To make maximum use of this, all functions, if possible, are declared
elemental (or pure, if this is not possible).
\nwenddocs{}\nwbegincode{359}\moddef{\code{}sm{\char95}physics.f90\edoc{}}\endmoddef
\LA{}File header\RA{}

module sm_physics

\LA{}Use kinds\RA{}
  use io_units
  use constants
  use unit_tests
  use diagnostics
  use physics_defs
  use lorentz

\LA{}Standard module head\RA{}

\LA{}SM physics: public\RA{}

\LA{}SM physics: parameters\RA{}

contains

\LA{}SM physics: procedures\RA{}

end module sm_physics
\eatline
\nwendcode{}\nwbegindocs{360}\nwdocspar
\subsection{Running $\alpha_s$}
\nwenddocs{}\nwbegindocs{361}Then we define the coefficients of the beta function of QCD (as a
reference cf. the Particle Data Group), where $n_f$ is the number of
active flavors in two different schemes:
\begin{align}
  \beta_0 &=\; 11 - \frac23 n_f \\
  \beta_1 &=\; 51 - \frac{19}{3} n_f \\
  \beta_2 &=\; 2857 - \frac{5033}{9} n_f + \frac{325}{27} n_f^2
\end{align}
\begin{align}
  b_0 &=\; \frac{1}{12 \pi} \left( 11 C_A  - 2 n_f \right) \\
  b_1 &=\; \frac{1}{24 \pi^2} \left( 17 C_A^2 - 5 C_A n_f - 3 C_F n_f \right) \\
  b_2 &=\; \frac{1}{(4\pi)^3} \biggl( \frac{2857}{54} C_A^3 -
  \frac{1415}{54} * C_A^2 n_f - \frac{205}{18} C_A C_F n_f + C_F^2 n_f
  + \frac{79}{54} C_A n_f**2 + \frac{11}{9} C_F n_f**2 \biggr)
\end{align}
\nwenddocs{}\nwbegincode{362}\moddef{SM physics: public}\endmoddef
  public :: beta0, beta1, beta2, coeff_b0, coeff_b1, coeff_b2
\nwendcode{}\nwbegincode{363}\moddef{SM physics: procedures}\endmoddef
  pure function beta0 (nf)
    real(default), intent(in) :: nf
    real(default) :: beta0
    beta0 = 11.0_default - two/three * nf
  end function beta0

  pure function beta1 (nf)
    real(default), intent(in) :: nf
    real(default) :: beta1
    beta1 = 51.0_default - 19.0_default/three * nf
  end function beta1

  pure function beta2 (nf)
    real(default), intent(in) :: nf
    real(default) :: beta2
    beta2 = 2857.0_default - 5033.0_default / 9.0_default * &
                    nf + 325.0_default/27.0_default * nf**2
  end function beta2

  pure function coeff_b0 (nf)
    real(default), intent(in) :: nf
    real(default) :: coeff_b0
    coeff_b0 = (11.0_default * CA - two * nf) / (12.0_default * pi)
  end function coeff_b0

  pure function coeff_b1 (nf)
    real(default), intent(in) :: nf
    real(default) :: coeff_b1
    coeff_b1 = (17.0_default * CA**2 - five * CA * nf - three * CF * nf) / &
               (24.0_default * pi**2)
  end function coeff_b1

  pure function coeff_b2 (nf)
    real(default), intent(in) :: nf
    real(default) :: coeff_b2
    coeff_b2 = (2857.0_default/54.0_default * CA**3 - &
                    1415.0_default/54.0_default * &
                    CA**2 * nf - 205.0_default/18.0_default * CA*CF*nf &
                    + 79.0_default/54.0_default * CA*nf**2 + &
                    11.0_default/9.0_default * CF * nf**2) / (four*pi)**3
  end function coeff_b2

\eatline
\eatline
\nwendcode{}\nwbegindocs{364}There should be two versions of running $\alpha_s$, one which takes
the scale and $\Lambda_{\text{QCD}}$ as input, and one which takes the
scale and e.g. $\alpha_s(m_Z)$ as input. Here, we take the one which
takes the QCD scale and scale as inputs from the PDG book.
\nwenddocs{}\nwbegincode{365}\moddef{SM physics: public}\plusendmoddef
  public :: running_as, running_as_lam
\nwendcode{}\nwbegincode{366}\moddef{SM physics: procedures}\plusendmoddef
  pure function running_as (scale, al_mz, mz, order, nf) result (ascale)
    real(default), intent(in) :: scale
    real(default), intent(in), optional :: al_mz, nf, mz
    integer, intent(in), optional :: order
    integer :: ord
    real(default) :: az, m_z, as_log, n_f, b0, b1, b2, ascale
    real(default) :: as0, as1
    if (present (mz)) then
       m_z = mz
    else
       m_z = MZ_REF
    end if
    if (present (order)) then
       ord = order
    else
       ord = 0
    end if
    if (present (al_mz)) then
       az = al_mz
    else
       az = ALPHA_QCD_MZ_REF
    end if
    if (present (nf)) then
       n_f = nf
    else
       n_f = 5
    end if
    b0 = coeff_b0 (n_f)
    b1 = coeff_b1 (n_f)
    b2 = coeff_b2 (n_f)
    as_log = one + b0 * az * log(scale**2/m_z**2)
    as0 = az / as_log
    as1 = as0 - as0**2 * b1/b0 * log(as_log)
    select case (ord)
    case (0)
       ascale = as0
    case (1)
       ascale = as1
    case (2)
       ascale = as1 + as0**3 * (b1**2/b0**2 * ((log(as_log))**2 - &
            log(as_log) + as_log - one) - b2/b0 * (as_log - one))
    case default
       ascale = as0
    end select
  end function running_as

  pure function running_as_lam (nf, scale, lambda, order) result (ascale)
    real(default), intent(in) :: nf, scale
    real(default), intent(in), optional :: lambda
    integer, intent(in), optional :: order
    real(default) :: lambda_qcd
    real(default) :: as0, as1, logmul, b0, b1, b2, ascale
    integer :: ord
    if (present (lambda)) then
       lambda_qcd = lambda
    else
       lambda_qcd = LAMBDA_QCD_REF
    end if
    if (present (order)) then
       ord = order
    else
       ord = 0
    end if
    b0 = beta0(nf)
    logmul = log(scale**2/lambda_qcd**2)
    as0 = four*pi / b0 / logmul
    if (ord > 0) then
       b1 = beta1(nf)
       as1 = as0 * (one - two* b1 / b0**2 * log(logmul) / logmul)
    end if
    select case (ord)
    case (0)
       ascale = as0
    case (1)
       ascale = as1
    case (2)
       b2 = beta2(nf)
       ascale = as1 + as0 * four * b1**2/b0**4/logmul**2 * &
            ((log(logmul) - 0.5_default)**2 + &
             b2*b0/8.0_default/b1**2 - five/four)
    case default
       ascale = as0
    end select
  end function running_as_lam

\eatline
\eatline
\nwendcode{}\nwbegindocs{367}\nwdocspar
\subsection{Catani-Seymour Parameters}
These are fundamental constants of the Catani-Seymour dipole formalism.
Since the corresponding parameters for the gluon case depend on the
number of flavors which is treated as an argument, there we do have
functions and not parameters.
\begin{equation}
  \gamma_q = \gamma_{\bar q} = \frac{3}{2} C_F \qquad \gamma_g =
  \frac{11}{6} C_A - \frac{2}{3} T_R N_f
\end{equation}
\begin{equation}
  K_q = K_{\bar q} = \left(  \frac{7}{2} - \frac{\pi^2}{6} \right) C_F \qquad
  K_g = \left( \frac{67}{18} - \frac{\pi^2}{6} \right) C_A -
  \frac{10}{9} T_R N_f
\end{equation}
\nwenddocs{}\nwbegincode{368}\moddef{SM physics: parameters}\endmoddef
  real(kind=default), parameter, public ::  gamma_q = three/two * CF, &
     k_q = (7.0_default/two - pi**2/6.0_default) * CF

\eatline
\nwendcode{}\nwbegindocs{369}\nwdocspar
\nwenddocs{}\nwbegincode{370}\moddef{SM physics: public}\plusendmoddef
  public :: gamma_g, k_g
\nwendcode{}\nwbegincode{371}\moddef{SM physics: procedures}\plusendmoddef
  elemental function gamma_g (nf) result (gg)
     real(kind=default), intent(in) :: nf
     real(kind=default) :: gg
     gg = 11.0_default/6.0_default * CA - two/three * TR * nf
  end function gamma_g

  elemental function k_g (nf) result (kg)
     real(kind=default), intent(in) :: nf
     real(kind=default) :: kg
     kg = (67.0_default/18.0_default - pi**2/6.0_default) * CA - &
                10.0_default/9.0_default * TR * nf
  end function k_g

\eatline
\eatline
\nwendcode{}\nwbegindocs{372}\nwdocspar
\subsection{Mathematical Functions}
The dilogarithm.  This simplified version is bound to double
precision, and restricted to argument values less or equal to unity,
so we do not need complex algebra.  The wrapper converts it to default
precision (which is, of course, a no-op if double=default).

The routine calculates the dilogarithm through mapping on the area
where there is a quickly convergent series (adapted from an F77
routine by Hans Kuijf, 1988): Map $x$ such that $x$ is not in the
neighbourhood of $1$.  Note that $|z|=-\ln(1-x)$ is always smaller
than $1.10$, but $\frac{1.10^{19}}{19!}{\rm Bernoulli}_{19}=2.7\times
10^{-15}$.
\nwenddocs{}\nwbegincode{373}\moddef{SM physics: public}\plusendmoddef
  public :: Li2
\nwendcode{}\nwbegincode{374}\moddef{SM physics: procedures}\plusendmoddef
  function Li2 (x)
      use kinds, only: double
      real(default), intent(in) :: x
      real(default) :: Li2
      Li2 = real( Li2_double (real(x, kind=double)), kind=default)
  end function Li2

\nwendcode{}\nwbegindocs{375}%def: Li2
\nwenddocs{}\nwbegindocs{376}\nwdocspar
\nwenddocs{}\nwbegincode{377}\moddef{SM physics: procedures}\plusendmoddef
  function Li2_double (x)  result (Li2)
    use kinds, only: double
    real(kind=double), intent(in) :: x
    real(kind=double) :: Li2
    real(kind=double), parameter :: pi2_6 = pi**2/6
    if (abs(1-x) < 1.E-13_double) then
       Li2 = pi2_6
    else if (abs(1-x) <  0.5_double) then
       Li2 = pi2_6 - log(1-x) * log(x) - Li2_restricted (1-x)
    else if (abs(x) > 1.d0) then
       ! Li2 = 0
       ! call msg_bug (" Dilogarithm called outside of defined range.")
       !!! Reactivate Dilogarithm identity
        Li2 = -pi2_6 - 0.5_default * log(-x) * log(-x) - Li2_restricted (1/x)
    else
       Li2 = Li2_restricted (x)
    end if
  contains
    function Li2_restricted (x) result (Li2)
      real(kind=double), intent(in) :: x
      real(kind=double) :: Li2
      real(kind=double) :: tmp, z, z2
      z = - log (1-x)
      z2 = z**2
! Horner's rule for the powers z^3 through z^19
      tmp = 43867._double/798._double
      tmp = tmp * z2 /342._double - 3617._double/510._double
      tmp = tmp * z2 /272._double + 7._double/6._double
      tmp = tmp * z2 /210._double - 691._double/2730._double
      tmp = tmp * z2 /156._double + 5._double/66._double
      tmp = tmp * z2 /110._double - 1._double/30._double
      tmp = tmp * z2 / 72._double + 1._double/42._double
      tmp = tmp * z2 / 42._double - 1._double/30._double
      tmp = tmp * z2 / 20._double + 1._double/6._double
! The first three terms of the power series
      Li2 = z2 * z * tmp / 6._double - 0.25_double * z2 + z
    end function Li2_restricted
  end function Li2_double

\eatline
\nwendcode{}\nwbegindocs{378}\nwdocspar
\subsection{Loop Integrals}
These functions appear in the calculation of the effective one-loop coupling of
a (pseudo)scalar to a vector boson pair.
\nwenddocs{}\nwbegincode{379}\moddef{SM physics: public}\plusendmoddef
  public :: faux
\nwendcode{}\nwbegincode{380}\moddef{SM physics: procedures}\plusendmoddef
  elemental function faux (x) result (y)
    real(default), intent(in) :: x
    complex(default) :: y
    if (1 <= x) then
       y = asin(sqrt(1/x))**2
    else
       y = - 1/4.0_default * (log((1 + sqrt(1 - x))/ &
            (1 - sqrt(1 - x))) - cmplx (0.0_default, pi, kind=default))**2
    end if
  end function faux

\eatline
\nwendcode{}\nwbegindocs{381}\nwdocspar
\nwenddocs{}\nwbegincode{382}\moddef{SM physics: public}\plusendmoddef
  public :: fonehalf
\nwendcode{}\nwbegincode{383}\moddef{SM physics: procedures}\plusendmoddef
  elemental function fonehalf (x) result (y)
    real(default), intent(in) :: x
    complex(default) :: y
    if (abs(x) < eps0) then
       y = 0
    else
       y = - 2.0_default * x * (1 + (1 - x) * faux(x))
    end if
  end function fonehalf

\eatline
\nwendcode{}\nwbegindocs{384}\nwdocspar
\nwenddocs{}\nwbegincode{385}\moddef{SM physics: public}\plusendmoddef
  public :: fonehalf_pseudo
\nwendcode{}\nwbegincode{386}\moddef{SM physics: procedures}\plusendmoddef
  function fonehalf_pseudo (x) result (y)
    real(default), intent(in) :: x
    complex(default) :: y
    if (abs(x) < eps0) then
       y = 0
    else
       y = - 2.0_default * x * faux(x)
    end if
  end function fonehalf_pseudo

\eatline
\nwendcode{}\nwbegindocs{387}\nwdocspar
\nwenddocs{}\nwbegincode{388}\moddef{SM physics: public}\plusendmoddef
  public :: fone
\nwendcode{}\nwbegincode{389}\moddef{SM physics: procedures}\plusendmoddef
  elemental function fone (x) result  (y)
    real(default), intent(in) :: x
    complex(default) :: y
    if (abs(x) < eps0) then
       y = 2.0_default
    else
       y = 2.0_default + 3.0_default * x + &
            3.0_default * x * (2.0_default - x) * &
            faux(x)
    end if
  end function fone

\eatline
\nwendcode{}\nwbegindocs{390}\nwdocspar
\nwenddocs{}\nwbegincode{391}\moddef{SM physics: public}\plusendmoddef
  public :: gaux
\nwendcode{}\nwbegincode{392}\moddef{SM physics: procedures}\plusendmoddef
  elemental function gaux (x) result (y)
    real(default), intent(in) :: x
    complex(default) :: y
    if (1 <= x) then
       y = sqrt(x - 1) * asin(sqrt(1/x))
    else
       y = sqrt(1 - x) * (log((1 + sqrt(1 - x)) / &
            (1 - sqrt(1 - x))) - &
            cmplx (0.0_default, pi, kind=default)) / 2.0_default
    end if
  end function gaux

\eatline
\nwendcode{}\nwbegindocs{393}\nwdocspar
\nwenddocs{}\nwbegincode{394}\moddef{SM physics: public}\plusendmoddef
  public :: tri_i1
\nwendcode{}\nwbegincode{395}\moddef{SM physics: procedures}\plusendmoddef
  elemental function tri_i1 (a,b) result (y)
    real(default), intent(in) :: a,b
    complex(default) :: y
    if (a < eps0 .or. b < eps0) then
       y = 0
    else
       y = a*b/2.0_default/(a-b) + a**2 * b**2/2.0_default/(a-b)**2 * &
            (faux(a) - faux(b)) + &
            a**2 * b/(a-b)**2 * (gaux(a) - gaux(b))
    end if
  end function tri_i1

\eatline
\nwendcode{}\nwbegindocs{396}\nwdocspar
\nwenddocs{}\nwbegincode{397}\moddef{SM physics: public}\plusendmoddef
  public :: tri_i2
\nwendcode{}\nwbegincode{398}\moddef{SM physics: procedures}\plusendmoddef
  elemental function tri_i2 (a,b) result (y)
    real(default), intent(in) :: a,b
    complex(default) :: y
    if (a < eps0 .or. b < eps0) then
       y = 0
    else
       y = - a * b / 2.0_default / (a-b) * (faux(a) - faux(b))
    end if
  end function tri_i2

\eatline
\nwendcode{}\nwbegindocs{399}\nwdocspar
\subsection{More on $\alpha_s$}
These functions are for the running of the strong coupling constants,
$\alpha_s$.
\nwenddocs{}\nwbegincode{400}\moddef{SM physics: public}\plusendmoddef
  public :: run_b0
\nwendcode{}\nwbegincode{401}\moddef{SM physics: procedures}\plusendmoddef
  elemental function run_b0 (nf) result (bnull)
    integer, intent(in) :: nf
    real(default) :: bnull
    bnull = 33.0_default - 2.0_default * nf
  end function run_b0

\eatline
\nwendcode{}\nwbegindocs{402}\nwdocspar
\nwenddocs{}\nwbegincode{403}\moddef{SM physics: public}\plusendmoddef
  public :: run_b1
\nwendcode{}\nwbegincode{404}\moddef{SM physics: procedures}\plusendmoddef
  elemental function run_b1 (nf) result (bone)
      integer, intent(in) :: nf
      real(default) :: bone
      bone = 6.0_default * (153.0_default - 19.0_default * nf)/run_b0(nf)**2
  end function run_b1

\eatline
\nwendcode{}\nwbegindocs{405}\nwdocspar
\nwenddocs{}\nwbegincode{406}\moddef{SM physics: public}\plusendmoddef
  public :: run_aa
\nwendcode{}\nwbegincode{407}\moddef{SM physics: procedures}\plusendmoddef
  elemental function run_aa (nf) result (aaa)
      integer, intent(in) :: nf
      real(default) :: aaa
      aaa = 12.0_default * PI / run_b0(nf)
  end function run_aa

\eatline
\nwendcode{}\nwbegindocs{408}\nwdocspar
\nwenddocs{}\nwbegincode{409}\moddef{SM physics: pubic functions}\endmoddef
  public :: run_bb
\nwendcode{}\nwbegincode{410}\moddef{SM physics: procedures}\plusendmoddef
  elemental function run_bb (nf) result (bbb)
    integer, intent(in) :: nf
    real(default) :: bbb
    bbb = run_b1(nf) / run_aa(nf)
  end function run_bb

\eatline
\nwendcode{}\nwbegindocs{411}\nwdocspar
\subsection{Functions for Catani-Seymour dipoles}

For the automated Catani-Seymour dipole subtraction, we need the
following functions.

\nwenddocs{}\nwbegincode{412}\moddef{SM physics: public}\plusendmoddef
  public :: ff_dipole
\nwendcode{}\nwbegincode{413}\moddef{SM physics: procedures}\plusendmoddef
  pure subroutine ff_dipole (v_ijk,y_ijk,p_ij,pp_k,p_i,p_j,p_k)
      type(vector4_t), intent(in) :: p_i, p_j, p_k
      type(vector4_t), intent(out) :: p_ij, pp_k
      real(kind=default), intent(out) :: y_ijk
      real(kind=default) :: z_i
      real(kind=default), intent(out) :: v_ijk
      z_i   = (p_i*p_k) / ((p_k*p_j) + (p_k*p_i))
      y_ijk = (p_i*p_j) / ((p_i*p_j) + (p_i*p_k) + (p_j*p_k))
      p_ij  = p_i + p_j - y_ijk/(1.0_default - y_ijk) * p_k
      pp_k  = (1.0/(1.0_default - y_ijk)) * p_k
      !!! We don't multiply by alpha_s right here:
      v_ijk = 8.0_default * PI * CF * &
           (2.0 / (1.0 - z_i*(1.0 - y_ijk)) - (1.0 + z_i))
  end subroutine ff_dipole

\eatline
\nwendcode{}\nwbegindocs{414}\nwdocspar
\nwenddocs{}\nwbegincode{415}\moddef{SM physics: public}\plusendmoddef
  public :: fi_dipole
\nwendcode{}\nwbegincode{416}\moddef{SM physics: procedures}\plusendmoddef
  pure subroutine fi_dipole (v_ija,x_ija,p_ij,pp_a,p_i,p_j,p_a)
     type(vector4_t), intent(in) :: p_i, p_j, p_a
     type(vector4_t), intent(out) :: p_ij, pp_a
     real(kind=default), intent(out) :: x_ija
     real(kind=default) :: z_i
     real(kind=default), intent(out) :: v_ija
     z_i   = (p_i*p_a) / ((p_a*p_j) + (p_a*p_i))
     x_ija = ((p_i*p_a) + (p_j*p_a) - (p_i*p_j)) &
          / ((p_i*p_a) + (p_j*p_a))
     p_ij  = p_i + p_j - (1.0_default - x_ija) * p_a
     pp_a  = x_ija * p_a
     !!! We don't not multiply by alpha_s right here:
     v_ija = 8.0_default * PI * CF * &
          (2.0 / (1.0 - z_i + (1.0 - x_ija)) - (1.0 + z_i)) / x_ija
  end subroutine fi_dipole

\eatline
\nwendcode{}\nwbegindocs{417}\nwdocspar
\nwenddocs{}\nwbegincode{418}\moddef{SM physics: public}\plusendmoddef
  public :: if_dipole
\nwendcode{}\nwbegincode{419}\moddef{SM physics: procedures}\plusendmoddef
  pure subroutine if_dipole (v_kja,u_j,p_aj,pp_k,p_k,p_j,p_a)
     type(vector4_t), intent(in) :: p_k, p_j, p_a
     type(vector4_t), intent(out) :: p_aj, pp_k
     real(kind=default), intent(out) :: u_j
     real(kind=default) :: x_kja
     real(kind=default), intent(out) :: v_kja
     u_j   = (p_a*p_j) / ((p_a*p_j) + (p_a*p_k))
     x_kja = ((p_a*p_k) + (p_a*p_j) - (p_j*p_k)) &
          / ((p_a*p_j) + (p_a*p_k))
     p_aj  = x_kja * p_a
     pp_k  = p_k + p_j - (1.0_default - x_kja) * p_a
     v_kja = 8.0_default * PI * CF * &
          (2.0 / (1.0 - x_kja + u_j) - (1.0 + x_kja)) / x_kja
  end subroutine if_dipole

\eatline
\nwendcode{}\nwbegindocs{420}This function depends on a variable number of final state particles
whose kinematics all get changed by the initial-initial dipole insertion.
\nwenddocs{}\nwbegincode{421}\moddef{SM physics: public}\plusendmoddef
  public :: ii_dipole
\nwendcode{}\nwbegincode{422}\moddef{SM physics: procedures}\plusendmoddef
  pure subroutine ii_dipole (v_jab,v_j,p_in,p_out,flag_1or2)
      type(vector4_t), dimension(:), intent(in) :: p_in
      type(vector4_t), dimension(size(p_in)-1), intent(out) :: p_out
      logical, intent(in) :: flag_1or2
      real(kind=default), intent(out) :: v_j
      real(kind=default), intent(out) :: v_jab
      type(vector4_t) :: p_a, p_b, p_j
      type(vector4_t) :: k, kk
      type(vector4_t) :: p_aj
      real(kind=default) :: x_jab
      integer :: i
      !!! flag_1or2 decides whether this a 12 or 21 dipole
      if (flag_1or2) then
         p_a = p_in(1)
         p_b = p_in(2)
      else
         p_b = p_in(1)
         p_a = p_in(2)
      end if
      !!! We assume that the unresolved particle has always the last
      !!! momentum
      p_j = p_in(size(p_in))
      x_jab = ((p_a*p_b) - (p_a*p_j) - (p_b*p_j)) / (p_a*p_b)
      v_j = (p_a*p_j) / (p_a * p_b)
      p_aj  = x_jab * p_a
      k     = p_a + p_b - p_j
      kk    = p_aj + p_b
      do i = 3, size(p_in)-1
         p_out(i) = p_in(i) - 2.0*((k+kk)*p_in(i))/((k+kk)*(k+kk)) * (k+kk) + &
              (2.0 * (k*p_in(i)) / (k*k)) * kk
      end do
      if (flag_1or2) then
         p_out(1) = p_aj
         p_out(2) = p_b
      else
         p_out(1) = p_b
         p_out(2) = p_aj
      end if
      v_jab = 8.0_default * PI * CF * &
           (2.0 / (1.0 - x_jab) - (1.0 + x_jab)) / x_jab
  end subroutine ii_dipole
\eatline
\nwendcode{}\nwbegindocs{423}\nwdocspar
\subsection{Distributions for integrated dipoles and such}
Note that the following formulae are only meaningful for
$0 \leq x \leq 1$.

The Dirac delta distribution, modified for Monte-Carlo sampling,
centered at $x=1-\frac{\epsilon}{2}$:

\nwenddocs{}\nwbegincode{424}\moddef{SM physics: public}\plusendmoddef
  public :: delta
\nwendcode{}\nwbegincode{425}\moddef{SM physics: procedures}\plusendmoddef
  elemental function delta (x,eps) result (z)
     real(kind=default), intent(in) :: x, eps
     real(kind=default) :: z
     if (x > one - eps) then
        z = one / eps
     else
        z = 0
     end if
  end function delta

\eatline
\nwendcode{}\nwbegindocs{426}The $+$-distribution, $P_+(x) = \left( \frac{1}{1-x}\right)_+$, for
the regularization of soft-collinear singularities. The constant part
for the Monte-Carlo sampling is the integral over the splitting
function divided by the weight for the WHIZARD numerical integration
over the interval.
\nwenddocs{}\nwbegincode{427}\moddef{SM physics: public}\plusendmoddef
  public :: plus_distr
\nwendcode{}\nwbegincode{428}\moddef{SM physics: procedures}\plusendmoddef
  elemental function plus_distr (x,eps) result (plusd)
     real(kind=default), intent(in) :: x, eps
     real(kind=default) :: plusd
     if (x > one - eps) then
        plusd = log(eps) / eps
     else
        plusd = one / (one - x)
     end if
  end function plus_distr

\eatline
\nwendcode{}\nwbegindocs{429}The splitting function in $D=4$ dimensions, regularized as
$+$-distributions if necessary:
\begin{align}
  P^{qq} (x) = P^{\bar q\bar q} (x) &= \; C_F \cdot \left( \frac{1 +
      x^2}{1-x} \right)_+ \\
  P^{qg} (x) = P^{\bar q g}     (x) &= \; C_F  \cdot \frac{1 + (1-x)^2}{x}\\
  P^{gq} (x) = P^{g \bar q}     (x) &= \; T_R \cdot \left[ x^2 + (1-x)^2
  \right] \\
  P^{gg} (x) &= \; 2 C_A \biggl[ \left( \frac{1}{1-x} \right)_+ +
                 \frac{1-x}{x} - 1 + x(1-x) \biggl] \notag{}\\
             &\quad + \delta(1-x) \left( \frac{11}{6} C_A -
                                          \frac{2}{3} N_f T_R \right)
\end{align}
Since the number of flavors summed over in the gluon splitting
function might depend on the physics case under consideration, it is
implemented as an input variable.
\nwenddocs{}\nwbegincode{430}\moddef{SM physics: public}\plusendmoddef
  public :: pqq
\nwendcode{}\nwbegincode{431}\moddef{SM physics: procedures}\plusendmoddef
  elemental function pqq (x,eps) result (pqqx)
     real(kind=default), intent(in) :: x, eps
     real(kind=default) :: pqqx
     if (x > (1.0_default - eps)) then
        pqqx = (eps - one) / two + two * log(eps) / eps - &
             three * (eps - one) / eps / two
     else
        pqqx = (one + x**2) / (one - x)
     end if
     pqqx = CF * pqqx
  end function pqq

\eatline
\nwendcode{}\nwbegindocs{432}\nwdocspar
\nwenddocs{}\nwbegincode{433}\moddef{SM physics: public}\plusendmoddef
  public :: pgq
\nwendcode{}\nwbegincode{434}\moddef{SM physics: procedures}\plusendmoddef
  elemental function pgq (x) result (pgqx)
     real(kind=default), intent(in) :: x
     real(kind=default) :: pgqx
     pgqx = TR * (x**2 + (one - x)**2)
  end function pgq

\eatline
\nwendcode{}\nwbegindocs{435}\nwdocspar
\nwenddocs{}\nwbegincode{436}\moddef{SM physics: public}\plusendmoddef
  public :: pqg
\nwendcode{}\nwbegincode{437}\moddef{SM physics: procedures}\plusendmoddef
  elemental function pqg (x) result (pqgx)
     real(kind=default), intent(in) :: x
     real(kind=default) :: pqgx
     pqgx = CF * (one + (one - x)**2) / x
  end function pqg

\eatline
\nwendcode{}\nwbegindocs{438}\nwdocspar
\nwenddocs{}\nwbegincode{439}\moddef{SM physics: public}\plusendmoddef
  public :: pgg
\nwendcode{}\nwbegincode{440}\moddef{SM physics: procedures}\plusendmoddef
  elemental function pgg (x, nf, eps) result (pggx)
    real(kind=default), intent(in) :: x, nf, eps
    real(kind=default) :: pggx
    pggx = two * CA * ( plus_distr (x, eps) + (one-x)/x - one + &
                   x*(one-x)) + delta (x, eps)  * gamma_g(nf)
  end function pgg

\eatline
\nwendcode{}\nwbegindocs{441}For the $qq$ and $gg$ cases, there exist ``regularized'' versions of
the splitting functions:
\begin{align}
  P^{qq}_{\text{reg}} (x) &= - C_F \cdot (1 + x) \\
  P^{gg}_{\text{reg}} (x) &= 2 C_A \left[ \frac{1-x}{x} - 1 + x(1-x) \right]
\end{align}
\nwenddocs{}\nwbegincode{442}\moddef{SM physics: public}\plusendmoddef
  public :: pqq_reg
\nwendcode{}\nwbegincode{443}\moddef{SM physics: procedures}\plusendmoddef
  elemental function pqq_reg (x) result (pqqregx)
     real(kind=default), intent(in) :: x
     real(kind=default) :: pqqregx
     pqqregx = - CF * (one + x)
  end function pqq_reg

\eatline
\nwendcode{}\nwbegindocs{444}\nwdocspar
\nwenddocs{}\nwbegincode{445}\moddef{SM physics: public}\plusendmoddef
  public :: pgg_reg
\nwendcode{}\nwbegincode{446}\moddef{SM physics: procedures}\plusendmoddef
  elemental function pgg_reg (x) result (pggregx)
     real(kind=default), intent(in) :: x
     real(kind=default) :: pggregx
     pggregx = two * CA * ((one - x)/x - one + x*(one - x))
  end function pgg_reg

\eatline
\nwendcode{}\nwbegindocs{447}Here, we collect the expressions needed for integrated
Catani-Seymour dipoles, and the so-called flavor kernels. We always
distinguish between the ``ordinary'' Catani-Seymour version, and the
one including a phase-space slicing parameter, $\alpha$.

The standard flavor kernels $\overline{K}^{ab}$ are:
\begin{align}
  \overline{K}^{qg} (x) = \overline{K}^{\bar q g} (x) & = \;
  P^{qg} (x) \log ((1-x)/x) + CF \times x \\
  %%%
  \overline{K}^{gq} (x) = \overline{K}^{g \bar q} (x) & = \;
  P^{gq} (x) \log ((1-x)/x) + TR \times 2x(1-x) \\
  %%%
  \overline{K}^{qq} &=\; C_F \biggl[ \left( \frac{2}{1-x} \log
  \frac{1-x}{x} \right)_+ - (1+x) \log ((1-x)/x)  +
  (1-x) \biggr] \notag{}\\
                    &\quad - (5 - \pi^2) \cdot C_F \cdot \delta(1-x) \\
  %%%
  \overline{K}^{gg} &=\; 2 C_A \biggl[ \left( \frac{1}{1-x} \log
    \frac{1-x}{x} \right)_+ + \left( \frac{1-x}{x} - 1 + x(1-x)
  \right) \log((1-x)/x) \biggr] \notag{}\\
                    &\quad - \delta(1-x) \biggl[ \left(
    \frac{50}{9} - \pi^2 \right) C_A - \frac{16}{9} T_R N_f \biggr]
\end{align}
\nwenddocs{}\nwbegincode{448}\moddef{SM physics: public}\plusendmoddef
  public :: kbarqg
\nwendcode{}\nwbegincode{449}\moddef{SM physics: procedures}\plusendmoddef
  function kbarqg (x) result (kbarqgx)
    real(kind=default), intent(in) :: x
    real(kind=default) :: kbarqgx
    kbarqgx = pqg(x) * log((one-x)/x) + CF * x
  end function kbarqg

\eatline
\nwendcode{}\nwbegindocs{450}\nwdocspar
\nwenddocs{}\nwbegincode{451}\moddef{SM physics: public}\plusendmoddef
  public :: kbargq
\nwendcode{}\nwbegincode{452}\moddef{SM physics: procedures}\plusendmoddef
  function kbargq (x) result (kbargqx)
    real(kind=default), intent(in) :: x
    real(kind=default) :: kbargqx
    kbargqx = pgq(x) * log((one-x)/x) + two * TR * x * (one - x)
  end function kbargq

\eatline
\nwendcode{}\nwbegindocs{453}\nwdocspar
\nwenddocs{}\nwbegincode{454}\moddef{SM physics: public}\plusendmoddef
  public :: kbarqq
\nwendcode{}\nwbegincode{455}\moddef{SM physics: procedures}\plusendmoddef
  function kbarqq (x,eps) result (kbarqqx)
    real(kind=default), intent(in) :: x, eps
    real(kind=default) :: kbarqqx
    kbarqqx = CF*(log_plus_distr(x,eps) - (one+x) * log((one-x)/x) + (one - &
         x) - (five - pi**2) * delta(x,eps))
  end function kbarqq

\eatline
\nwendcode{}\nwbegindocs{456}\nwdocspar
\nwenddocs{}\nwbegincode{457}\moddef{SM physics: public}\plusendmoddef
  public :: kbargg
\nwendcode{}\nwbegincode{458}\moddef{SM physics: procedures}\plusendmoddef
  function kbargg (x,eps,nf) result (kbarggx)
    real(kind=default), intent(in) :: x, eps, nf
    real(kind=default) :: kbarggx
    kbarggx = CA * (log_plus_distr(x,eps) + two * ((one-x)/x - one + &
                         x*(one-x) * log((1-x)/x))) - delta(x,eps) * &
                         ((50.0_default/9.0_default - pi**2) * CA - &
                         16.0_default/9.0_default * TR * nf)
  end function kbargg

\eatline
\nwendcode{}\nwbegindocs{459}The $\tilde{K}$ are used when two identified hadrons participate:
\begin{equation}
    \tilde{K}^{ab} (x) = P^{ab}_{\text{reg}} (x) \cdot \log (1-x) +
    \delta^{ab} \mathbf{T}_a^2 \biggl[ \left( \frac{2}{1-x} \log (1-x)
    \right)_+ - \frac{\pi^2}{3} \delta(1-x) \biggr]
\end{equation}

\nwenddocs{}\nwbegincode{460}\moddef{SM physics: public}\plusendmoddef
  public :: ktildeqq
\nwendcode{}\nwbegincode{461}\moddef{SM physics: procedures}\plusendmoddef
  function ktildeqq (x,eps) result (ktildeqqx)
    real(kind=default), intent(in) :: x, eps
    real(kind=default) :: ktildeqqx
    ktildeqqx = pqq_reg (x) * log(one-x) + CF * ( - log2_plus_distr (x,eps) &
                          - pi**2/three * delta(x,eps))
  end function ktildeqq

\eatline
\nwendcode{}\nwbegindocs{462}\nwdocspar
\nwenddocs{}\nwbegincode{463}\moddef{SM physics: public}\plusendmoddef
  public :: ktildeqg
\nwendcode{}\nwbegincode{464}\moddef{SM physics: procedures}\plusendmoddef
  function ktildeqg (x,eps) result (ktildeqgx)
    real(kind=default), intent(in) :: x, eps
    real(kind=default) :: ktildeqgx
    ktildeqgx = pqg (x) * log(one-x)
  end function ktildeqg

\eatline
\nwendcode{}\nwbegindocs{465}\nwdocspar
\nwenddocs{}\nwbegincode{466}\moddef{SM physics: public}\plusendmoddef
  public :: ktildegq
\nwendcode{}\nwbegincode{467}\moddef{SM physics: procedures}\plusendmoddef
  function ktildegq (x,eps) result (ktildegqx)
    real(kind=default), intent(in) :: x, eps
    real(kind=default) :: ktildegqx
    ktildegqx = pgq (x) * log(one-x)
  end function ktildegq

\eatline
\nwendcode{}\nwbegindocs{468}\nwdocspar
\nwenddocs{}\nwbegincode{469}\moddef{SM physics: public}\plusendmoddef
  public :: ktildegg
\nwendcode{}\nwbegincode{470}\moddef{SM physics: procedures}\plusendmoddef
  function ktildegg (x,eps) result (ktildeggx)
    real(kind=default), intent(in) :: x, eps
    real(kind=default) :: ktildeggx
    ktildeggx = pgg_reg (x) * log(one-x) + CA * ( - &
       log2_plus_distr (x,eps) - pi**2/three * delta(x,eps))
  end function ktildegg

\eatline
\nwendcode{}\nwbegindocs{471}The insertion operator might not be necessary for a GOLEM interface
but is demanded by the Les Houches NLO accord. It is a
three-dimensional array, where the index always gives the inverse
power of the DREG expansion parameter, $\epsilon$.
\nwenddocs{}\nwbegincode{472}\moddef{SM physics: public}\plusendmoddef
  public :: insert_q
\nwendcode{}\nwbegincode{473}\moddef{SM physics: procedures}\plusendmoddef
  pure function insert_q ()
    real(kind=default), dimension(0:2) :: insert_q
    insert_q(0) = gamma_q + k_q - pi**2/three * CF
    insert_q(1) = gamma_q
    insert_q(2) = CF
  end function insert_q

\eatline
\nwendcode{}\nwbegindocs{474}\nwdocspar
\nwenddocs{}\nwbegincode{475}\moddef{SM physics: public}\plusendmoddef
  public :: insert_g
\nwendcode{}\nwbegincode{476}\moddef{SM physics: procedures}\plusendmoddef
  pure function insert_g (nf)
    real(kind=default), intent(in) :: nf
    real(kind=default), dimension(0:2) :: insert_g
    insert_g(0) = gamma_g (nf) + k_g (nf) - pi**2/three * CA
    insert_g(1) = gamma_g (nf)
    insert_g(2) = CA
  end function insert_g

\eatline
\nwendcode{}\nwbegindocs{477}For better convergence, one can exclude regions of phase space with
a slicing parameter from the dipole subtraction procedure. First of
all, the $K$ functions get modified:
\begin{equation}
  K_i (\alpha) = K_i - \mathbf{T}_i^2 \log^2 \alpha + \gamma_i (
  \alpha - 1 - \log\alpha)
\end{equation}
\nwenddocs{}\nwbegincode{478}\moddef{SM physics: public}\plusendmoddef
  public :: k_q_al, k_g_al
\nwendcode{}\nwbegincode{479}\moddef{SM physics: procedures}\plusendmoddef
  pure function k_q_al (alpha)
    real(kind=default), intent(in) :: alpha
    real(kind=default) :: k_q_al
    k_q_al = k_q - CF * (log(alpha))**2 + gamma_q * &
                      (alpha - one - log(alpha))
  end function k_q_al

  pure function k_g_al (alpha, nf)
    real(kind=default), intent(in) :: alpha, nf
    real(kind=default) :: k_g_al
    k_g_al = k_g (nf) - CA * (log(alpha))**2 + gamma_g (nf) * &
                     (alpha - one - log(alpha))
  end function k_g_al

\eatline
\eatline
\nwendcode{}\nwbegindocs{480}The $+$-distribution, but with a phase-space slicing parameter,
$\alpha$,  $P_{1-\alpha}(x) = \left( \frac{1}{1-x}
\right)_{1-x}$. Since we need the fatal error message here, this
function cannot be elemental.
\nwenddocs{}\nwbegincode{481}\moddef{SM physics: public}\plusendmoddef
  public :: plus_distr_al
\nwendcode{}\nwbegincode{482}\moddef{SM physics: procedures}\plusendmoddef
  function plus_distr_al (x,alpha,eps) result (plusd_al)
     real(kind=default), intent(in) :: x,  eps, alpha
     real(kind=default) :: plusd_al
     if ((one - alpha) >= (one - eps)) then
        plusd_al = zero
        call msg_fatal ('sm_physics, plus_distr_al: alpha and epsilon chosen wrongly')
     elseif (x < (1.0_default - alpha)) then
        plusd_al = 0
     else if (x > (1.0_default - eps)) then
        plusd_al = log(eps/alpha)/eps
     else
        plusd_al = one/(one-x)
     end if
   end function plus_distr_al

\eatline
\nwendcode{}\nwbegindocs{483}Introducing phase-space slicing parameters, these standard flavor
kernels $\overline{K}^{ab}$ become:
\begin{align}
  \overline{K}^{qg}_\alpha (x) = \overline{K}^{\bar q g}_\alpha (x) & = \;
  P^{qg} (x) \log (\alpha (1-x)/x) + C_F \times x \\
  %%%
  \overline{K}^{gq}_\alpha (x) = \overline{K}^{g \bar q}_\alpha (x) & = \;
  P^{gq} (x) \log (\alpha (1-x)/x) + T_R \times 2x(1-x) \\
  %%%
  \overline{K}^{qq}_\alpha &=
      C_F  (1 - x) + P^{qq}_{\text{reg}} (x) \log \frac{\alpha(1-x)}{x}
      \notag{}\\           &\quad
      + C_F \delta (1 - x) \log^2 \alpha
      + C_F \left( \frac{2}{1-x} \log \frac{1-x}{x} \right)_+ \notag{}\\
                           &\quad
      - \left( \gamma_q + K_q(\alpha) - \frac56 \pi^2 C_F \right) \cdot
      \delta(1-x) \; C_F \Bigl[ + \frac{2}{1-x}  \log \left(
      \frac{\alpha (2-x)}{1+\alpha-x} \right)
      - \theta(1 - \alpha - x) \cdot \left( \frac{2}{1-x} \log
      \frac{2-x}{1-x} \right) \Bigr] \\
  %%%
  \overline{K}^{gg}_\alpha &=\;
       P^{gg}_{\text{reg}} (x) \log \frac{\alpha(1-x)}{x}
      + C_A \delta (1 - x) \log^2 \alpha  \notag{}\\
                           &\quad
      + C_A \left( \frac{2}{1-x} \log \frac{1-x}{x} \right)_+
       - \left( \gamma_g + K_g(\alpha) - \frac56 \pi^2 C_A \right) \cdot
       \delta(1-x) \; C_A \Bigl[ + \frac{2}{1-x}  \log \left(
       \frac{\alpha (2-x)}{1+\alpha-x} \right)
      - \theta(1 - \alpha - x) \cdot \left( \frac{2}{1-x} \log
      \frac{2-x}{1-x} \right) \Bigr]
\end{align}
\nwenddocs{}\nwbegincode{484}\moddef{SM physics: public}\plusendmoddef
  public :: kbarqg_al
\nwendcode{}\nwbegincode{485}\moddef{SM physics: procedures}\plusendmoddef
  function kbarqg_al (x,alpha,eps) result (kbarqgx)
    real(kind=default), intent(in) :: x, alpha, eps
    real(kind=default) :: kbarqgx
    kbarqgx = pqg (x) * log(alpha*(one-x)/x) + CF * x
  end function kbarqg_al
\eatline
\nwendcode{}\nwbegindocs{486}\nwdocspar

\nwenddocs{}\nwbegincode{487}\moddef{SM physics: public}\plusendmoddef
  public :: kbargq_al
\nwendcode{}\nwbegincode{488}\moddef{SM physics: procedures}\plusendmoddef
  function kbargq_al (x,alpha,eps) result (kbargqx)
    real(kind=default), intent(in) :: x, alpha, eps
    real(kind=default) :: kbargqx
    kbargqx = pgq (x) * log(alpha*(one-x)/x) + two * TR * x * (one-x)
  end function kbargq_al
\eatline
\nwendcode{}\nwbegindocs{489}\nwdocspar

\nwenddocs{}\nwbegincode{490}\moddef{SM physics: public}\plusendmoddef
  public :: kbarqq_al
\nwendcode{}\nwbegincode{491}\moddef{SM physics: procedures}\plusendmoddef
  function kbarqq_al (x,alpha,eps) result (kbarqqx)
     real(kind=default), intent(in) :: x, alpha, eps
     real(kind=default) :: kbarqqx
     kbarqqx = CF * (one - x) + pqq_reg(x) * log(alpha*(one-x)/x) &
              + CF * log_plus_distr(x,eps) &
             - (gamma_q + k_q_al(alpha) - CF * &
              five/6.0_default  * pi**2 - CF * (log(alpha))**2) * &
              delta(x,eps) + &
              CF * two/(one -x)*log(alpha*(two-x)/(one+alpha-x))
     if (x < (one-alpha)) then
        kbarqqx = kbarqqx - CF * two/(one-x) * log((two-x)/(one-x))
     end if
  end function kbarqq_al

\eatline
\nwendcode{}\nwbegincode{492}\moddef{SM physics: public}\plusendmoddef
  public :: kbargg_al
\nwendcode{}\nwbegincode{493}\moddef{SM physics: procedures}\plusendmoddef
  function kbargg_al (x,alpha,eps,nf) result (kbarggx)
     real(kind=default), intent(in) :: x, alpha, eps, nf
     real(kind=default) :: kbarggx
     kbarggx = pgg_reg(x) * log(alpha*(one-x)/x) &
              + CA * log_plus_distr(x,eps) &
             - (gamma_g(nf) + k_g_al(alpha,nf) - CA * &
              five/6.0_default  * pi**2 - CA * (log(alpha))**2) * &
              delta(x,eps) + &
              CA * two/(one -x)*log(alpha*(two-x)/(one+alpha-x))
     if (x < (one-alpha)) then
        kbarggx = kbarggx - CA * two/(one-x) * log((two-x)/(one-x))
     end if
  end function kbargg_al

\eatline
\nwendcode{}\nwbegindocs{494}The $\tilde{K}$ flavor kernels in the presence of a phase-space slicing
parameter, are:
\begin{equation}
  \tilde{K}^{ab} (x,\alpha) = P^{qq, \text{reg}} (x)
  \log\frac{1-x}{\alpha} + ..........
\end{equation}
\nwenddocs{}\nwbegincode{495}\moddef{SM physics: public}\plusendmoddef
  public :: ktildeqq_al
\nwendcode{}\nwbegincode{496}\moddef{SM physics: procedures}\plusendmoddef
  function ktildeqq_al (x,alpha,eps) result (ktildeqqx)
    real(kind=default), intent(in) :: x, eps, alpha
    real(kind=default) :: ktildeqqx
    ktildeqqx = pqq_reg(x) * log((one-x)/alpha) + CF*( &
         - log2_plus_distr_al(x,alpha,eps) - Pi**2/three * delta(x,eps) &
         + (one+x**2)/(one-x) * log(min(one,(alpha/(one-x)))) &
         + two/(one-x) * log((one+alpha-x)/alpha))
    if (x > (one-alpha)) then
       ktildeqqx = ktildeqqx - CF*two/(one-x)*log(two-x)
    end if
  end function ktildeqq_al

\eatline
\nwendcode{}\nwbegindocs{497}This is a logarithmic $+$-distribution, $\left(
  \frac{\log((1-x)/x)}{1-x} \right)_+$. For the sampling, we need the
integral over this function over the incomplete sampling interval
$[0,1-\epsilon]$, which is $\log^2(x) + 2 Li_2(x) -
\frac{\pi^2}{3}$. As this function is negative definite for $\epsilon
> 0.1816$, we take a hard upper limit for that sampling parameter,
irrespective of the fact what the user chooses.
\nwenddocs{}\nwbegincode{498}\moddef{SM physics: public}\plusendmoddef
  public :: log_plus_distr
\nwendcode{}\nwbegincode{499}\moddef{SM physics: procedures}\plusendmoddef
  function log_plus_distr (x,eps) result (lpd)
     real(kind=default), intent(in) :: x, eps
     real(kind=default) :: lpd, eps2
     eps2 = min (eps, 0.1816_default)
     if (x > (1.0_default - eps2)) then
        lpd = ((log(eps2))**2 + two*Li2(eps2) - pi**2/three)/eps2
     else
        lpd = two*log((one-x)/x)/(one-x)
     end if
  end function log_plus_distr

\eatline
\nwendcode{}\nwbegindocs{500}Logarithmic $+$-distribution, $2 \left( \frac{\log(1/(1-x))}{1-x} \right)_+$.
\nwenddocs{}\nwbegincode{501}\moddef{SM physics: public}\plusendmoddef
  public :: log2_plus_distr
\nwendcode{}\nwbegincode{502}\moddef{SM physics: procedures}\plusendmoddef
  function log2_plus_distr (x,eps) result (lpd)
      real(kind=default), intent(in) :: x, eps
      real(kind=default) :: lpd
      if (x > (1.0_default - eps)) then
         lpd = - (log(eps))**2/eps
      else
         lpd = two*log(one/(one-x))/(one-x)
      end if
  end function log2_plus_distr

\eatline
\nwendcode{}\nwbegindocs{503}Logarithmic $+$-distribution with phase-space slicing parameter, $2
\left( \frac{\log(1/(1-x))}{1-x} \right)_{1-\alpha}$.
\nwenddocs{}\nwbegincode{504}\moddef{SM physics: public}\plusendmoddef
  public :: log2_plus_distr_al
\nwendcode{}\nwbegincode{505}\moddef{SM physics: procedures}\plusendmoddef
  function log2_plus_distr_al (x,alpha,eps) result (lpd_al)
    real(kind=default), intent(in) :: x, eps, alpha
    real(kind=default) :: lpd_al
    if ((one - alpha) >= (one - eps)) then
       lpd_al = zero
       call msg_fatal ('alpha and epsilon chosen wrongly')
    elseif (x < (one - alpha)) then
       lpd_al = 0
    elseif (x > (1.0_default - eps)) then
       lpd_al = - ((log(eps))**2 - (log(alpha))**2)/eps
    else
       lpd_al = two*log(one/(one-x))/(one-x)
    end if
  end function log2_plus_distr_al

\eatline
\nwendcode{}\nwbegindocs{506}\nwdocspar
\subsection{Splitting Functions}
\nwenddocs{}\nwbegindocs{507}Analogue to the regularized distributions of the last subsection, we
give here the unregularized splitting functions, relevant for the parton
shower algorithm. We can use this unregularized version since there will
be a cut-off $\epsilon$ that ensures that $\{z,1-z\}>\epsilon(t)$. This
cut-off seperates resolvable from unresolvable emissions.

{\tt{}p{\char95}xxx} are the kernels that are summed over helicity:
\nwenddocs{}\nwbegincode{508}\moddef{SM physics: public}\plusendmoddef
  public :: p_qqg
  public :: p_gqq
  public :: p_ggg
\nwendcode{}\nwbegindocs{509}$q\to q g$
\nwenddocs{}\nwbegincode{510}\moddef{SM physics: procedures}\plusendmoddef
  elemental function p_qqg (z) result (P)
    real(default), intent(in) :: z
    real(default) :: P
    P = CF * (one + z**2) / (one - z)
  end function p_qqg
\nwendcode{}\nwbegindocs{511}$g\to q \bar{q}$
\nwenddocs{}\nwbegincode{512}\moddef{SM physics: procedures}\plusendmoddef
  elemental function p_gqq (z) result (P)
    real(default), intent(in) :: z
    real(default) :: P
    P = TR * (z**2 + (one - z)**2)
  end function p_gqq
\nwendcode{}\nwbegindocs{513}$g\to g g$
\nwenddocs{}\nwbegincode{514}\moddef{SM physics: procedures}\plusendmoddef
  elemental function p_ggg (z) result (P)
    real(default), intent(in) :: z
    real(default) :: P
    P = NC * ((one - z) / z + z / (one - z) + z * (one - z))
  end function p_ggg

\eatline
\nwendcode{}\nwbegindocs{515}Analytically integrated splitting kernels:
\nwenddocs{}\nwbegincode{516}\moddef{SM physics: public}\plusendmoddef
  public :: integral_over_p_qqg
  public :: integral_over_p_gqq
  public :: integral_over_p_ggg
\nwendcode{}\nwbegincode{517}\moddef{SM physics: procedures}\plusendmoddef
  pure function integral_over_p_qqg (zmin, zmax) result (integral)
    real(default), intent(in) :: zmin, zmax
    real(default) :: integral
    integral = (two / three) * (- zmax**2 + zmin**2 - &
         two * (zmax - zmin) + four * log((one - zmin) / (one - zmax)))
  end function integral_over_p_qqg

  pure function integral_over_p_gqq (zmin, zmax) result (integral)
    real(default), intent(in) :: zmin, zmax
    real(default) :: integral
    integral = 0.5_default * ((two / three) * &
         (zmax**3 - zmin**3) - (zmax**2 - zmin**2) + (zmax - zmin))
  end function integral_over_p_gqq

  pure function integral_over_p_ggg (zmin, zmax) result (integral)
    real(default), intent(in) :: zmin, zmax
    real(default) :: integral
    integral = three * ((log(zmax) - two * zmax - &
         log(one - zmax) + zmax**2 / two - zmax**3 / three) - &
         (log(zmin) - zmin - zmin - log(one - zmin) + zmin**2 &
         / two - zmin**3 / three) )
  end function integral_over_p_ggg

\eatline
\nwendcode{}\nwbegindocs{518}We can also use (massless) helicity dependent splitting functions:
\nwenddocs{}\nwbegincode{519}\moddef{SM physics: public}\plusendmoddef
  public :: p_qqg_pol
\nwendcode{}\nwbegindocs{520}$q_a\to q_b g_c$, the helicity of the quark is not changed by gluon
emission and the gluon is preferably polarized in the branching plane
($l_c=1$):
\nwenddocs{}\nwbegincode{521}\moddef{SM physics: procedures}\plusendmoddef
  elemental function p_qqg_pol (z, l_a, l_b, l_c) result (P)
    real(default), intent(in) :: z
    integer, intent(in) :: l_a, l_b, l_c
    real(default) :: P
    if (l_a /= l_b) then
       P = zero
       return
    end if
    if (l_c == -1) then
       P = one - z
    else
       P = (one + z)**2 / (one - z)
    end if
    P = P * CF
  end function p_qqg_pol

\nwendcode{}\nwbegindocs{522}\nwdocspar
\subsection{Unit tests}
\nwenddocs{}\nwbegincode{523}\moddef{SM qcd: public}\endmoddef
  public :: sm_physics_test
\nwendcode{}\nwbegincode{524}\moddef{SM qcd: tests}\endmoddef
  subroutine sm_physics_test (u, results)
    integer, intent(in) :: u
    type(test_results_t), intent(inout) :: results
  \LA{}SM qcd: execute tests\RA{}
  end subroutine sm_physics_test

\eatline
\nwendcode{}\nwbegindocs{525}\nwdocspar
\subsubsection{Splitting functions}
\nwenddocs{}\nwbegincode{526}\moddef{SM physics: execute tests}\endmoddef
  call test (sm_physics_1, "sm_physics_1", &
       "running alpha_s", &
       u, results)
\nwendcode{}\nwbegincode{527}\moddef{SM physics: tests}\endmoddef
  subroutine sm_physics_1 (u)
    integer, intent(in) :: u
    real(default) :: z = 0.75_default

    write (u, "(A)")  "* Test output: sm_physics_1"
    write (u, "(A)")  "*   Purpose: check analytic properties"
    write (u, "(A)")

    write (u, "(A)")  "* Splitting functions:"
    write (u, "(A)")

    call assert (u, vanishes (p_qqg_pol (z, +1, -1, +1)))
    call assert (u, vanishes (p_qqg_pol (z, +1, -1, -1)))
    call assert (u, vanishes (p_qqg_pol (z, -1, +1, +1)))
    call assert (u, vanishes (p_qqg_pol (z, -1, +1, -1)))

    call assert (u, nearly_equal ( &
         p_qqg_pol (z, +1, +1, -1) + p_qqg_pol (z, +1, +1, +1),
         p_qqg (z))

    write (u, "(A)")
    write (u, "(A)")  "* Test output end: sm_physics_1"

  end subroutine sm_physics_1

\eatline
\nwendcode{}\nwbegindocs{528}\nwdocspar
\clearpage
%%%%%%%%%%%%%%%%%%%%%%%%%%%%%%%%%%%%%%%%%%%%%%%%%%%%%%%%%%%%%%%%%%%%%%%
\section{QCD Coupling}
We provide various distinct implementations of the QCD coupling.  In
this module, we define an abstract data type and three
implementations: fixed, running with $\alpha_s(M_Z)$ as input, and
running with $\Lambda_{\text{QCD}}$ as input.  We use the functions
defined above in the module {\tt{}sm{\char95}physics} but provide a common
interface.  Later modules may define additional implementations.
\nwenddocs{}\nwbegincode{529}\moddef{\code{}sm{\char95}qcd.f90\edoc{}}\endmoddef
\LA{}File header\RA{}

module sm_qcd

\LA{}Use kinds\RA{}
  use io_units
  use format_defs, only: FMT_12
  use unit_tests
  use diagnostics
  use md5
  use physics_defs
  use sm_physics

\LA{}Standard module head\RA{}

\LA{}SM qcd: public\RA{}

\LA{}SM qcd: types\RA{}

\LA{}SM qcd: interfaces\RA{}

contains

\LA{}SM qcd: procedures\RA{}

\LA{}SM qcd: tests\RA{}

end module sm_qcd
\eatline
\nwendcode{}\nwbegindocs{530}\nwdocspar
\subsection{Coupling: Abstract Data Type}
This is the abstract version of the QCD coupling implementation.
\nwenddocs{}\nwbegincode{531}\moddef{SM qcd: public}\plusendmoddef
  public :: alpha_qcd_t
\nwendcode{}\nwbegincode{532}\moddef{SM qcd: types}\endmoddef
  type, abstract :: alpha_qcd_t
   contains
   \LA{}SM qcd: alpha qcd: TBP\RA{}
  end type alpha_qcd_t

\eatline
\nwendcode{}\nwbegindocs{533}There must be an output routine.
\nwenddocs{}\nwbegincode{534}\moddef{SM qcd: alpha qcd: TBP}\endmoddef
  procedure (alpha_qcd_write), deferred :: write
\nwendcode{}\nwbegincode{535}\moddef{SM qcd: interfaces}\endmoddef
  abstract interface
     subroutine alpha_qcd_write (object, unit)
       import
       class(alpha_qcd_t), intent(in) :: object
       integer, intent(in), optional :: unit
     end subroutine alpha_qcd_write
  end interface

\eatline
\nwendcode{}\nwbegindocs{536}This method computes the running coupling, given a certain scale.  All
parameters (reference value, order of the approximation, etc.) must be
set before calling this.
\nwenddocs{}\nwbegincode{537}\moddef{SM qcd: alpha qcd: TBP}\plusendmoddef
  procedure (alpha_qcd_get), deferred :: get
\nwendcode{}\nwbegincode{538}\moddef{SM qcd: interfaces}\plusendmoddef
  abstract interface
     function alpha_qcd_get (alpha_qcd, scale) result (alpha)
       import
       class(alpha_qcd_t), intent(in) :: alpha_qcd
       real(default), intent(in) :: scale
       real(default) :: alpha
     end function alpha_qcd_get
  end interface

\eatline
\nwendcode{}\nwbegindocs{539}\nwdocspar
\subsection{Fixed Coupling}
In this version, the $\alpha_s$ value is fixed, the {\tt{}scale} argument
of the {\tt{}get} method is ignored.  There is only one parameter, the
value.  By default, this is the value at $M_Z$.
\nwenddocs{}\nwbegincode{540}\moddef{SM qcd: public}\plusendmoddef
  public :: alpha_qcd_fixed_t
\nwendcode{}\nwbegincode{541}\moddef{SM qcd: types}\plusendmoddef
  type, extends (alpha_qcd_t) :: alpha_qcd_fixed_t
     real(default) :: val = ALPHA_QCD_MZ_REF
   contains
   \LA{}SM qcd: alpha qcd fixed: TBP\RA{}
  end type alpha_qcd_fixed_t

\eatline
\nwendcode{}\nwbegindocs{542}Output.
\nwenddocs{}\nwbegincode{543}\moddef{SM qcd: alpha qcd fixed: TBP}\endmoddef
  procedure :: write => alpha_qcd_fixed_write
\nwendcode{}\nwbegincode{544}\moddef{SM qcd: procedures}\endmoddef
  subroutine alpha_qcd_fixed_write (object, unit)
    class(alpha_qcd_fixed_t), intent(in) :: object
    integer, intent(in), optional :: unit
    integer :: u
    u = given_output_unit (unit)
    write (u, "(3x,A)")  "QCD parameters (fixed coupling):"
    write (u, "(5x,A," // FMT_12 // ")")  "alpha = ", object%val
  end subroutine alpha_qcd_fixed_write

\eatline
\nwendcode{}\nwbegindocs{545}Calculation: the scale is ignored in this case.
\nwenddocs{}\nwbegincode{546}\moddef{SM qcd: alpha qcd fixed: TBP}\plusendmoddef
  procedure :: get => alpha_qcd_fixed_get
\nwendcode{}\nwbegincode{547}\moddef{SM qcd: procedures}\plusendmoddef
  function alpha_qcd_fixed_get (alpha_qcd, scale) result (alpha)
    class(alpha_qcd_fixed_t), intent(in) :: alpha_qcd
    real(default), intent(in) :: scale
    real(default) :: alpha
    alpha = alpha_qcd%val
  end function alpha_qcd_fixed_get

\eatline
\nwendcode{}\nwbegindocs{548}\nwdocspar
\subsection{Running Coupling}
In this version, the $\alpha_s$ value runs relative to the value at a
given reference scale.  There are two parameters: the value of this
scale (default: $M_Z$), the value of $\alpha_s$ at this scale, and the
number of effective flavors.  Furthermore, we have the order of the
approximation.
\nwenddocs{}\nwbegincode{549}\moddef{SM qcd: public}\plusendmoddef
  public :: alpha_qcd_from_scale_t
\nwendcode{}\nwbegincode{550}\moddef{SM qcd: types}\plusendmoddef
  type, extends (alpha_qcd_t) :: alpha_qcd_from_scale_t
     real(default) :: mu_ref = MZ_REF
     real(default) :: ref = ALPHA_QCD_MZ_REF
     integer :: order = 0
     integer :: nf = 5
   contains
   \LA{}SM qcd: alpha qcd from scale: TBP\RA{}
  end type alpha_qcd_from_scale_t

\eatline
\nwendcode{}\nwbegindocs{551}Output.
\nwenddocs{}\nwbegincode{552}\moddef{SM qcd: alpha qcd from scale: TBP}\endmoddef
  procedure :: write => alpha_qcd_from_scale_write
\nwendcode{}\nwbegincode{553}\moddef{SM qcd: procedures}\plusendmoddef
  subroutine alpha_qcd_from_scale_write (object, unit)
    class(alpha_qcd_from_scale_t), intent(in) :: object
    integer, intent(in), optional :: unit
    integer :: u
    u = given_output_unit (unit)
    write (u, "(3x,A)")  "QCD parameters (running coupling):"
    write (u, "(5x,A," // FMT_12 // ")")  "Scale mu  = ", object%mu_ref
    write (u, "(5x,A," // FMT_12 // ")")  "alpha(mu) = ", object%ref
    write (u, "(5x,A,I0)")      "LL order  = ", object%order
    write (u, "(5x,A,I0)")      "N(flv)    = ", object%nf
  end subroutine alpha_qcd_from_scale_write

\eatline
\nwendcode{}\nwbegindocs{554}Calculation: here, we call the function for running $\alpha_s$ that
was defined in {\tt{}sm{\char95}physics} above.  The function does not take into
account thresholds, so the number of flavors should be the correct one
for the chosen scale.  Normally, this should be the $Z$ boson mass.
\nwenddocs{}\nwbegincode{555}\moddef{SM qcd: alpha qcd from scale: TBP}\plusendmoddef
  procedure :: get => alpha_qcd_from_scale_get
\nwendcode{}\nwbegincode{556}\moddef{SM qcd: procedures}\plusendmoddef
  function alpha_qcd_from_scale_get (alpha_qcd, scale) result (alpha)
    class(alpha_qcd_from_scale_t), intent(in) :: alpha_qcd
    real(default), intent(in) :: scale
    real(default) :: alpha
    alpha = running_as (scale, &
         alpha_qcd%ref, alpha_qcd%mu_ref, alpha_qcd%order, &
         real (alpha_qcd%nf, kind=default))
  end function alpha_qcd_from_scale_get

\eatline
\nwendcode{}\nwbegindocs{557}\nwdocspar
\subsection{Running Coupling, determined by $\Lambda_{\text{QCD}}$}
In this version, the input are the value $\Lambda_{\text{QCD}}$ and
the order of the approximation.
\nwenddocs{}\nwbegincode{558}\moddef{SM qcd: public}\plusendmoddef
  public :: alpha_qcd_from_lambda_t
\nwendcode{}\nwbegincode{559}\moddef{SM qcd: types}\plusendmoddef
  type, extends (alpha_qcd_t) :: alpha_qcd_from_lambda_t
     real(default) :: lambda = LAMBDA_QCD_REF
     integer :: order = 0
     integer :: nf = 5
   contains
   \LA{}SM qcd: alpha qcd from lambda: TBP\RA{}
  end type alpha_qcd_from_lambda_t

\eatline
\nwendcode{}\nwbegindocs{560}Output.
\nwenddocs{}\nwbegincode{561}\moddef{SM qcd: alpha qcd from lambda: TBP}\endmoddef
  procedure :: write => alpha_qcd_from_lambda_write
\nwendcode{}\nwbegincode{562}\moddef{SM qcd: procedures}\plusendmoddef
  subroutine alpha_qcd_from_lambda_write (object, unit)
    class(alpha_qcd_from_lambda_t), intent(in) :: object
    integer, intent(in), optional :: unit
    integer :: u
    u = given_output_unit (unit)
    write (u, "(3x,A)")  "QCD parameters (Lambda_QCD as input):"
    write (u, "(5x,A," // FMT_12 // ")")  "Lambda_QCD = ", object%lambda
    write (u, "(5x,A,I0)")      "LL order   = ", object%order
    write (u, "(5x,A,I0)")      "N(flv)     = ", object%nf
  end subroutine alpha_qcd_from_lambda_write

\eatline
\nwendcode{}\nwbegindocs{563}Calculation: here, we call the second function for running $\alpha_s$ that
was defined in {\tt{}sm{\char95}physics} above.  The $\Lambda$ value should be
the one that is appropriate for the chosen number of effective
flavors.  Again, thresholds are not incorporated.
\nwenddocs{}\nwbegincode{564}\moddef{SM qcd: alpha qcd from lambda: TBP}\plusendmoddef
  procedure :: get => alpha_qcd_from_lambda_get
\nwendcode{}\nwbegincode{565}\moddef{SM qcd: procedures}\plusendmoddef
  function alpha_qcd_from_lambda_get (alpha_qcd, scale) result (alpha)
    class(alpha_qcd_from_lambda_t), intent(in) :: alpha_qcd
    real(default), intent(in) :: scale
    real(default) :: alpha
    alpha = running_as_lam (real (alpha_qcd%nf, kind=default), scale, &
         alpha_qcd%lambda, alpha_qcd%order)
  end function alpha_qcd_from_lambda_get

\eatline
\nwendcode{}\nwbegindocs{566}\nwdocspar
\subsection{Wrapper type}
We could get along with a polymorphic QCD type, but a monomorphic wrapper type
with a polymorphic component is easier to handle and probably safer
(w.r.t.\ compiler bugs).  However, we keep the object transparent, so we can
set the type-specific parameters directly (by a {\tt{}dispatch} routine).
\nwenddocs{}\nwbegincode{567}\moddef{SM qcd: public}\plusendmoddef
  public :: qcd_t
\nwendcode{}\nwbegincode{568}\moddef{SM qcd: types}\plusendmoddef
  type :: qcd_t
     class(alpha_qcd_t), allocatable :: alpha
     character(32) :: md5sum = ""
   contains
   \LA{}SM qcd: qcd: TBP\RA{}
  end type qcd_t

\eatline
\nwendcode{}\nwbegindocs{569}Output.  We first print the polymorphic {\tt{}alpha} which contains a headline,
then any extra components.
\nwenddocs{}\nwbegincode{570}\moddef{SM qcd: qcd: TBP}\endmoddef
  procedure :: write => qcd_write
\nwendcode{}\nwbegincode{571}\moddef{SM qcd: procedures}\plusendmoddef
  subroutine qcd_write (qcd, unit, show_md5sum)
    class(qcd_t), intent(in) :: qcd
    integer, intent(in), optional :: unit
    logical, intent(in), optional :: show_md5sum
    logical :: show_md5
    integer :: u
    u = given_output_unit (unit)
    show_md5 = .true.;  if (present (show_md5sum))  show_md5 = show_md5sum
    if (allocated (qcd%alpha)) then
       call qcd%alpha%write (u)
    else
       write (u, "(3x,A)")  "QCD parameters (coupling undefined)"
    end if
    if (show_md5 .and. qcd%md5sum /= "") &
         write (u, "(5x,A,A,A)") "md5sum = '", qcd%md5sum, "'"
  end subroutine qcd_write

\eatline
\nwendcode{}\nwbegindocs{572}Compute a MD5 sum for the {\tt{}alpha{\char95}s} setup.  This is
done by writing them to a temporary file, using a standard format.
\nwenddocs{}\nwbegincode{573}\moddef{SM qcd: qcd: TBP}\plusendmoddef
  procedure :: compute_alphas_md5sum => qcd_compute_alphas_md5sum
\nwendcode{}\nwbegincode{574}\moddef{SM qcd: procedures}\plusendmoddef
  subroutine qcd_compute_alphas_md5sum (qcd)
    class(qcd_t), intent(inout) :: qcd
    integer :: unit
    if (allocated (qcd%alpha)) then
       unit = free_unit ()
       open (unit, status="scratch", action="readwrite")
       call qcd%alpha%write (unit)
       rewind (unit)
       qcd%md5sum = md5sum (unit)
       close (unit)
    end if
  end subroutine qcd_compute_alphas_md5sum

\eatline
\nwendcode{}\nwbegindocs{575}\nwdocspar
\nwenddocs{}\nwbegindocs{576}Retrieve the MD5 sum of the qcd setup.
\nwenddocs{}\nwbegincode{577}\moddef{SM qcd: qcd: TBP}\plusendmoddef
  procedure :: get_md5sum => qcd_get_md5sum
\nwendcode{}\nwbegincode{578}\moddef{SM qcd: procedures}\plusendmoddef
  function qcd_get_md5sum (qcd) result (md5sum)
    character(32) :: md5sum
    class(qcd_t), intent(inout) :: qcd
    md5sum = qcd%md5sum
  end function qcd_get_md5sum

\eatline
\nwendcode{}\nwbegindocs{579}\nwdocspar
\subsection{Unit tests}
\nwenddocs{}\nwbegincode{580}\moddef{SM qcd: public}\plusendmoddef
  public :: sm_qcd_test
\nwendcode{}\nwbegincode{581}\moddef{SM qcd: tests}\plusendmoddef
  subroutine sm_qcd_test (u, results)
    integer, intent(in) :: u
    type(test_results_t), intent(inout) :: results
  \LA{}SM qcd: execute tests\RA{}
  end subroutine sm_qcd_test

\eatline
\nwendcode{}\nwbegindocs{582}\nwdocspar
\subsubsection{QCD Coupling}
We check two different implementations of the abstract QCD coupling.
\nwenddocs{}\nwbegincode{583}\moddef{SM qcd: execute tests}\endmoddef
  call test (sm_qcd_1, "sm_qcd_1", &
       "running alpha_s", &
       u, results)
\nwendcode{}\nwbegincode{584}\moddef{SM qcd: tests}\plusendmoddef
  subroutine sm_qcd_1 (u)
    integer, intent(in) :: u
    type(qcd_t) :: qcd

    write (u, "(A)")  "* Test output: sm_qcd_1"
    write (u, "(A)")  "*   Purpose: compute running alpha_s"
    write (u, "(A)")

    write (u, "(A)")  "* Fixed:"
    write (u, "(A)")

    allocate (alpha_qcd_fixed_t :: qcd%alpha)
    call qcd%compute_alphas_md5sum ()

    call qcd%write (u)
    write (u, *)
    write (u, "(1x,A,F10.7)")  "alpha_s (mz)    =", &
         qcd%alpha%get (MZ_REF)
    write (u, "(1x,A,F10.7)")  "alpha_s (1 TeV) =", &
         qcd%alpha%get (1000._default)
    write (u, *)
    deallocate (qcd%alpha)

    write (u, "(A)")  "* Running from MZ (LO):"
    write (u, "(A)")

    allocate (alpha_qcd_from_scale_t :: qcd%alpha)
    call qcd%compute_alphas_md5sum ()

    call qcd%write (u)
    write (u, *)
    write (u, "(1x,A,F10.7)")  "alpha_s (mz)    =", &
         qcd%alpha%get (MZ_REF)
    write (u, "(1x,A,F10.7)")  "alpha_s (1 TeV) =", &
         qcd%alpha%get (1000._default)
    write (u, *)

    write (u, "(A)")  "* Running from MZ (NLO):"
    write (u, "(A)")

    select type (alpha => qcd%alpha)
    type is (alpha_qcd_from_scale_t)
       alpha%order = 1
    end select
    call qcd%compute_alphas_md5sum ()

    call qcd%write (u)
    write (u, *)
    write (u, "(1x,A,F10.7)")  "alpha_s (mz)    =", &
         qcd%alpha%get (MZ_REF)
    write (u, "(1x,A,F10.7)")  "alpha_s (1 TeV) =", &
         qcd%alpha%get (1000._default)
    write (u, *)

    write (u, "(A)")  "* Running from MZ (NNLO):"
    write (u, "(A)")

    select type (alpha => qcd%alpha)
    type is (alpha_qcd_from_scale_t)
       alpha%order = 2
    end select
    call qcd%compute_alphas_md5sum ()

    call qcd%write (u)
    write (u, *)
    write (u, "(1x,A,F10.7)")  "alpha_s (mz)    =", &
         qcd%alpha%get (MZ_REF)
    write (u, "(1x,A,F10.7)")  "alpha_s (1 TeV) =", &
         qcd%alpha%get (1000._default)
    write (u, *)

    deallocate (qcd%alpha)
    write (u, "(A)")  "* Running from Lambda_QCD (LO):"
    write (u, "(A)")

    allocate (alpha_qcd_from_lambda_t :: qcd%alpha)
    call qcd%compute_alphas_md5sum ()

    call qcd%write (u)
    write (u, *)
    write (u, "(1x,A,F10.7)")  "alpha_s (mz)    =", &
         qcd%alpha%get (MZ_REF)
    write (u, "(1x,A,F10.7)")  "alpha_s (1 TeV) =", &
         qcd%alpha%get (1000._default)
    write (u, *)

    write (u, "(A)")  "* Running from Lambda_QCD (NLO):"
    write (u, "(A)")

    select type (alpha => qcd%alpha)
    type is (alpha_qcd_from_lambda_t)
       alpha%order = 1
    end select
    call qcd%compute_alphas_md5sum ()

    call qcd%write (u)
    write (u, *)
    write (u, "(1x,A,F10.7)")  "alpha_s (mz)    =", &
         qcd%alpha%get (MZ_REF)
    write (u, "(1x,A,F10.7)")  "alpha_s (1 TeV) =", &
         qcd%alpha%get (1000._default)
    write (u, *)

    write (u, "(A)")  "* Running from Lambda_QCD (NNLO):"
    write (u, "(A)")

    select type (alpha => qcd%alpha)
    type is (alpha_qcd_from_lambda_t)
       alpha%order = 2
    end select
    call qcd%compute_alphas_md5sum ()

    call qcd%write (u)
    write (u, *)
    write (u, "(1x,A,F10.7)")  "alpha_s (mz)    =", &
         qcd%alpha%get (MZ_REF)
    write (u, "(1x,A,F10.7)")  "alpha_s (1 TeV) =", &
         qcd%alpha%get (1000._default)

    write (u, "(A)")
    write (u, "(A)")  "* Test output end: sm_qcd_1"

  end subroutine sm_qcd_1

\eatline
\nwendcode{}\nwfilename{../../../src/qft/qft.nw}\nwbegindocs{0}%% -*- ess-noweb-default-code-mode: f90-mode; noweb-default-code-mode: f90-mode; -*- 
% WHIZARD code as NOWEB source: Quantum Field Theory concepts
%%%%%%%%%%%%%%%%%%%%%%%%%%%%%%%%%%%%%%%%%%%%%%%%%%%%%%%%%%%%%%%%%%%%%%% 
\chapter{Quantum Field Theory Concepts}

The objects and methods defined here implement concepts and data for
the underlying quantum field theory that we use for computing matrix
elements and processes.
\begin{description}
\item[model\_data]
  Fields and coupling parameters, operators as vertex structures, for
  a specific model.
\item[model\_testbed]
  Provide hooks to deal with a {\tt{}model{\char95}data} extension without
  referencing it explicitly.
\item[helicities]
  Types and methods for spin density matrices.
\item[colors]
  Dealing with colored particles, using the color-flow representation.
\item[flavors]
  PDG codes and particle properties, depends on the model.
\item[quantum\_numbers]
  Quantum numbers and density matrices for entangled particle systems.
\end{description}

\clearpage
%%%%%%%%%%%%%%%%%%%%%%%%%%%%%%%%%%%%%%%%%%%%%%%%%%%%%%%%%%%%%%%%%%%%%%%%
\section{Model Data}

These data represent a specific Lagrangian in numeric terms.  That is,
we have the fields with their quantum numbers, the masses, widths and
couplings as numerical values, and the vertices as arrays of fields.

We do not store the relations between coupling parameters.  They
should be represented by expressions for evaluation, implemented as
Sindarin objects in a distinct data structure.  Neither do we need the
algebraic structure of vertices.  The field content of vertices is
required for the sole purpose of setting up phase space.
\nwenddocs{}\nwbegincode{1}\moddef{\code{}model{\char95}data.f90\edoc{}}\endmoddef
\LA{}File header\RA{}

module model_data

  use, intrinsic :: iso_c_binding !NODEP!
  
\LA{}Use kinds\RA{}
  use kinds, only: i8, i32
  use kinds, only: c_default_float
\LA{}Use strings\RA{}
  use format_defs, only: FMT_19
  use io_units
  use unit_tests
  use diagnostics
  use md5
  use hashes, only: hash
  use physics_defs, only: UNDEFINED, SCALAR

\LA{}Standard module head\RA{}

\LA{}Model data: public\RA{}

\LA{}Model data: parameters\RA{}

\LA{}Model data: types\RA{}

\LA{}Model data: interfaces\RA{}

contains

\LA{}Model data: procedures\RA{}

\LA{}Model data: tests\RA{}

end module model_data
\eatline
\nwendcode{}\nwbegindocs{2}\nwdocspar
\subsection{Physics Parameters}
Couplings, masses, and widths are physics parameters.  Each parameter
has a unique name (used, essentially, for diagnostics output and
debugging) and a value.  The value may be a real or a complex number,
so we provide to implementations of an abstract type.
\nwenddocs{}\nwbegincode{3}\moddef{Model data: public}\endmoddef
  public :: modelpar_data_t
\nwendcode{}\nwbegincode{4}\moddef{Model data: types}\endmoddef
  type, abstract :: modelpar_data_t
     private
     type(string_t) :: name
   contains
   \LA{}Model data: par data: TBP\RA{}
  end type modelpar_data_t
  
  type, extends (modelpar_data_t) :: modelpar_real_t
     private
     real(default) :: value
  end type modelpar_real_t
  
  type, extends (modelpar_data_t) :: modelpar_complex_t
     private
     complex(default) :: value
  end type modelpar_complex_t
  
\eatline
\nwendcode{}\nwbegindocs{5}\nwdocspar
Output for diagnostics.  Non-advancing.
\nwenddocs{}\nwbegincode{6}\moddef{Model data: par data: TBP}\endmoddef
  procedure :: write => par_write
\nwendcode{}\nwbegincode{7}\moddef{Model data: procedures}\endmoddef
  subroutine par_write (par, unit)
    class(modelpar_data_t), intent(in) :: par
    integer, intent(in), optional :: unit
    integer :: u
    u = given_output_unit (unit)
    write (u, "(1x,A,1x,A)", advance="no")  char (par%name), "= "
    select type (par)
    class is (modelpar_real_t)
       write (u, "(" // FMT_19 // ")", advance="no")  par%value
    class is (modelpar_complex_t)
       write (u, "(" // FMT_19 // ",1x,'+',1x," // FMT_19 // ",1x,'I')", &
            advance="no")  par%value
    end select
  end subroutine par_write
    
\eatline
\nwendcode{}\nwbegindocs{8}\nwdocspar
Pretty-printed on separate line, with fixed line length
\nwenddocs{}\nwbegincode{9}\moddef{Model data: par data: TBP}\plusendmoddef
  procedure :: show => par_show
\nwendcode{}\nwbegincode{10}\moddef{Model data: procedures}\plusendmoddef
  subroutine par_show (par, l, u)
    class(modelpar_data_t), intent(in) :: par
    integer, intent(in) :: l, u
    character(len=l) :: buffer
    buffer = par%name
    select type (par)
    class is (modelpar_real_t)
       write (u, "(4x,A,1x,'=',1x," // FMT_19 // ")")  buffer, par%value
    class is (modelpar_complex_t)
       write (u, "(4x,A,1x,'=',1x," // FMT_19 // ",1x,'+',1x," &
            // FMT_19 // ",1x,'I')")  buffer, par%value
    end select
  end subroutine par_show
    
\eatline
\nwendcode{}\nwbegindocs{11}\nwdocspar
Initialize with name and value.  The type depends on the argument
type.  If the type does not match, the value is converted following
Fortran rules.
\nwenddocs{}\nwbegincode{12}\moddef{Model data: par data: TBP}\plusendmoddef
  generic :: init => modelpar_data_init_real, modelpar_data_init_complex
  procedure, private :: modelpar_data_init_real
  procedure, private :: modelpar_data_init_complex
\nwendcode{}\nwbegincode{13}\moddef{Model data: procedures}\plusendmoddef
  subroutine modelpar_data_init_real (par, name, value)
    class(modelpar_data_t), intent(out) :: par
    type(string_t), intent(in) :: name
    real(default), intent(in) :: value
    par%name = name
    par = value
  end subroutine modelpar_data_init_real

  subroutine modelpar_data_init_complex (par, name, value)
    class(modelpar_data_t), intent(out) :: par
    type(string_t), intent(in) :: name
    complex(default), intent(in) :: value
    par%name = name
    par = value
  end subroutine modelpar_data_init_complex

\eatline
\nwendcode{}\nwbegindocs{14}\nwdocspar
Modify the value.  We assume that the parameter has been
initialized.  The type (real or complex) must not be changed, and the
name is also fixed.  
\nwenddocs{}\nwbegincode{15}\moddef{Model data: par data: TBP}\plusendmoddef
  generic :: assignment(=) => modelpar_data_set_real, modelpar_data_set_complex
  procedure, private :: modelpar_data_set_real
  procedure, private :: modelpar_data_set_complex
\nwendcode{}\nwbegincode{16}\moddef{Model data: procedures}\plusendmoddef
  elemental subroutine modelpar_data_set_real (par, value)
    class(modelpar_data_t), intent(inout) :: par
    real(default), intent(in) :: value
    select type (par)
    class is (modelpar_real_t)
       par%value = value
    class is (modelpar_complex_t)
       par%value = value
    end select
  end subroutine modelpar_data_set_real

  elemental subroutine modelpar_data_set_complex (par, value)
    class(modelpar_data_t), intent(inout) :: par
    complex(default), intent(in) :: value
    select type (par)
    class is (modelpar_real_t)
       par%value = value
    class is (modelpar_complex_t)
       par%value = value
    end select
  end subroutine modelpar_data_set_complex

\eatline
\nwendcode{}\nwbegindocs{17}\nwdocspar
Return the parameter name.
\nwenddocs{}\nwbegincode{18}\moddef{Model data: par data: TBP}\plusendmoddef
  procedure :: get_name => modelpar_data_get_name
\nwendcode{}\nwbegincode{19}\moddef{Model data: procedures}\plusendmoddef
  function modelpar_data_get_name (par) result (name)
    class(modelpar_data_t), intent(in) :: par
    type(string_t) :: name
    name = par%name
  end function modelpar_data_get_name
  
\eatline
\nwendcode{}\nwbegindocs{20}\nwdocspar
Return the value.  In case of a type mismatch, follow Fortran conventions.
\nwenddocs{}\nwbegincode{21}\moddef{Model data: par data: TBP}\plusendmoddef
  procedure, pass :: get_real => modelpar_data_get_real
  procedure, pass :: get_complex => modelpar_data_get_complex
\nwendcode{}\nwbegincode{22}\moddef{Model data: procedures}\plusendmoddef
  elemental function modelpar_data_get_real (par) result (value)
    class(modelpar_data_t), intent(in), target :: par
    real(default) :: value
    select type (par)
    class is (modelpar_real_t)
       value = par%value
    class is (modelpar_complex_t)
       value = par%value
    end select
  end function modelpar_data_get_real
  
  elemental function modelpar_data_get_complex (par) result (value)
    class(modelpar_data_t), intent(in), target :: par
    complex(default) :: value
    select type (par)
    class is (modelpar_real_t)
       value = par%value
    class is (modelpar_complex_t)
       value = par%value
    end select
  end function modelpar_data_get_complex
  
\eatline
\eatline
\nwendcode{}\nwbegindocs{23}\nwdocspar
Return a pointer to the value.  This makes sense only for matching types.
\nwenddocs{}\nwbegincode{24}\moddef{Model data: par data: TBP}\plusendmoddef
  procedure :: get_real_ptr => modelpar_data_get_real_ptr
  procedure :: get_complex_ptr => modelpar_data_get_complex_ptr
\nwendcode{}\nwbegincode{25}\moddef{Model data: procedures}\plusendmoddef
  function modelpar_data_get_real_ptr (par) result (ptr)
    class(modelpar_data_t), intent(in), target :: par
    real(default), pointer :: ptr
    select type (par)
    class is (modelpar_real_t)
       ptr => par%value
    class default
       ptr => null ()
    end select
  end function modelpar_data_get_real_ptr
  
  function modelpar_data_get_complex_ptr (par) result (ptr)
    class(modelpar_data_t), intent(in), target :: par
    complex(default), pointer :: ptr
    select type (par)
    class is (modelpar_complex_t)
       ptr => par%value
    class default
       ptr => null ()
    end select
  end function modelpar_data_get_complex_ptr
  
\eatline
\eatline
\nwendcode{}\nwbegindocs{26}\nwdocspar
\subsection{Field Data}
The field-data type holds all information that pertains to a particular field
(or particle) within a particular model.  Information such as spin type,
particle code etc.\ is stored within the object itself, while mass and width
are associated to parameters, otherwise assumed zero.
\nwenddocs{}\nwbegincode{27}\moddef{Model data: public}\plusendmoddef
  public :: field_data_t
\nwendcode{}\nwbegincode{28}\moddef{Model data: types}\plusendmoddef
  type :: field_data_t
     private
     type(string_t) :: longname
     integer :: pdg = UNDEFINED
     logical :: visible = .true.
     logical :: parton = .false.
     logical :: gauge = .false.
     logical :: left_handed = .false.
     logical :: right_handed = .false.
     logical :: has_anti = .false.
     logical :: p_is_stable = .true.
     logical :: p_decays_isotropically = .false.
     logical :: p_decays_diagonal = .false.
     logical :: p_has_decay_helicity = .false.
     integer :: p_decay_helicity = 0
     logical :: a_is_stable = .true.
     logical :: a_decays_isotropically = .false.
     logical :: a_decays_diagonal = .false.
     logical :: a_has_decay_helicity = .false.
     integer :: a_decay_helicity = 0
     logical :: p_polarized = .false.
     logical :: a_polarized = .false.
     type(string_t), dimension(:), allocatable :: name, anti
     type(string_t) :: tex_name, tex_anti
     integer :: spin_type = UNDEFINED
     integer :: isospin_type = 1
     integer :: charge_type = 1
     integer :: color_type = 1
     real(default), pointer :: mass_val => null ()
     class(modelpar_data_t), pointer :: mass_data => null ()
     real(default), pointer :: width_val => null ()
     class(modelpar_data_t), pointer :: width_data => null ()
     integer :: multiplicity = 1
     type(string_t), dimension(:), allocatable :: p_decay
     type(string_t), dimension(:), allocatable :: a_decay
   contains
   \LA{}Model data: field data: TBP\RA{}
  end type field_data_t

\eatline
\nwendcode{}\nwbegindocs{29}Initialize field data with PDG long name and PDG code.  \TeX\
names should be initialized to avoid issues with accessing
unallocated string contents.
\nwenddocs{}\nwbegincode{30}\moddef{Model data: field data: TBP}\endmoddef
  procedure :: init => field_data_init
\nwendcode{}\nwbegincode{31}\moddef{Model data: procedures}\plusendmoddef
  subroutine field_data_init (prt, longname, pdg)
    class(field_data_t), intent(out) :: prt
    type(string_t), intent(in) :: longname
    integer, intent(in) :: pdg
    prt%longname = longname
    prt%pdg = pdg
    prt%tex_name = ""
    prt%tex_anti = ""
  end subroutine field_data_init
    
\eatline
\nwendcode{}\nwbegindocs{32}Copy quantum numbers from another particle.  Do not compute the multiplicity
yet, because this depends on the association of the {\tt{}mass{\char95}data} pointer.
\nwenddocs{}\nwbegincode{33}\moddef{Model data: field data: TBP}\plusendmoddef
  procedure :: copy_from => field_data_copy_from
\nwendcode{}\nwbegincode{34}\moddef{Model data: procedures}\plusendmoddef
  subroutine field_data_copy_from (prt, prt_src)
    class(field_data_t), intent(inout) :: prt
    class(field_data_t), intent(in) :: prt_src
    prt%visible = prt_src%visible
    prt%parton = prt_src%parton
    prt%gauge = prt_src%gauge
    prt%left_handed = prt_src%left_handed
    prt%right_handed = prt_src%right_handed
    prt%p_is_stable =             prt_src%p_is_stable
    prt%p_decays_isotropically =  prt_src%p_decays_isotropically
    prt%p_decays_diagonal =       prt_src%p_decays_diagonal   
    prt%p_has_decay_helicity =    prt_src%p_has_decay_helicity
    prt%p_decay_helicity =        prt_src%p_decay_helicity
    prt%p_decays_diagonal =       prt_src%p_decays_diagonal   
    prt%a_is_stable =             prt_src%a_is_stable        
    prt%a_decays_isotropically =  prt_src%a_decays_isotropically
    prt%a_decays_diagonal =       prt_src%a_decays_diagonal   
    prt%a_has_decay_helicity =    prt_src%a_has_decay_helicity
    prt%a_decay_helicity =        prt_src%a_decay_helicity
    prt%p_polarized =             prt_src%p_polarized        
    prt%a_polarized =             prt_src%a_polarized
    prt%spin_type = prt_src%spin_type
    prt%isospin_type = prt_src%isospin_type
    prt%charge_type = prt_src%charge_type
    prt%color_type = prt_src%color_type
    prt%has_anti = prt_src%has_anti
    if (allocated (prt_src%name)) then
       if (allocated (prt%name))  deallocate (prt%name)
       allocate (prt%name (size (prt_src%name)), source = prt_src%name)
    end if
    if (allocated (prt_src%anti)) then
       if (allocated (prt%anti))  deallocate (prt%anti)
       allocate (prt%anti (size (prt_src%anti)), source = prt_src%anti)
    end if
    prt%tex_name = prt_src%tex_name
    prt%tex_anti = prt_src%tex_anti
    if (allocated (prt_src%p_decay)) then
       if (allocated (prt%p_decay))  deallocate (prt%p_decay)
       allocate (prt%p_decay (size (prt_src%p_decay)), source = prt_src%p_decay)
    end if
    if (allocated (prt_src%a_decay)) then
       if (allocated (prt%a_decay))  deallocate (prt%a_decay)
       allocate (prt%a_decay (size (prt_src%a_decay)), source = prt_src%a_decay)
    end if
  end subroutine field_data_copy_from

\eatline
\nwendcode{}\nwbegindocs{35}Set particle quantum numbers.
\nwenddocs{}\nwbegincode{36}\moddef{Model data: field data: TBP}\plusendmoddef
  procedure :: set => field_data_set
\nwendcode{}\nwbegincode{37}\moddef{Model data: procedures}\plusendmoddef
  subroutine field_data_set (prt, &
       is_visible, is_parton, is_gauge, is_left_handed, is_right_handed, &
       p_is_stable, p_decays_isotropically, p_decays_diagonal, &
       p_decay_helicity, &
       a_is_stable, a_decays_isotropically, a_decays_diagonal, &
       a_decay_helicity, &
       p_polarized, a_polarized, &
       name, anti, tex_name, tex_anti, &
       spin_type, isospin_type, charge_type, color_type, &
       mass_data, width_data, &
       p_decay, a_decay)
    class(field_data_t), intent(inout) :: prt
    logical, intent(in), optional :: is_visible, is_parton, is_gauge
    logical, intent(in), optional :: is_left_handed, is_right_handed
    logical, intent(in), optional :: p_is_stable
    logical, intent(in), optional :: p_decays_isotropically, p_decays_diagonal
    integer, intent(in), optional :: p_decay_helicity
    logical, intent(in), optional :: a_is_stable
    logical, intent(in), optional :: a_decays_isotropically, a_decays_diagonal
    integer, intent(in), optional :: a_decay_helicity
    logical, intent(in), optional :: p_polarized, a_polarized
    type(string_t), dimension(:), intent(in), optional :: name, anti
    type(string_t), intent(in), optional :: tex_name, tex_anti
    integer, intent(in), optional :: spin_type, isospin_type
    integer, intent(in), optional :: charge_type, color_type
    class(modelpar_data_t), intent(in), pointer, optional :: mass_data, width_data
    type(string_t), dimension(:), intent(in), optional :: p_decay, a_decay
    if (present (is_visible))  prt%visible = is_visible
    if (present (is_parton))  prt%parton = is_parton
    if (present (is_gauge))  prt%gauge = is_gauge
    if (present (is_left_handed))  prt%left_handed = is_left_handed
    if (present (is_right_handed))  prt%right_handed = is_right_handed
    if (present (p_is_stable))  prt%p_is_stable = p_is_stable
    if (present (p_decays_isotropically)) &
          prt%p_decays_isotropically = p_decays_isotropically
    if (present (p_decays_diagonal)) &
          prt%p_decays_diagonal = p_decays_diagonal
    if (present (p_decay_helicity)) then
       prt%p_has_decay_helicity = .true.
       prt%p_decay_helicity = p_decay_helicity
    end if
    if (present (a_is_stable))  prt%a_is_stable = a_is_stable
    if (present (a_decays_isotropically)) &
          prt%a_decays_isotropically = a_decays_isotropically
    if (present (a_decays_diagonal)) &
          prt%a_decays_diagonal = a_decays_diagonal
    if (present (a_decay_helicity)) then
       prt%a_has_decay_helicity = .true.
       prt%a_decay_helicity = a_decay_helicity
    end if
    if (present (p_polarized)) prt%p_polarized = p_polarized
    if (present (a_polarized)) prt%a_polarized = a_polarized
    if (present (name)) then
       if (allocated (prt%name))  deallocate (prt%name)
       allocate (prt%name (size (name)), source = name)
    end if
    if (present (anti)) then
       if (allocated (prt%anti))  deallocate (prt%anti)
       allocate (prt%anti (size (anti)), source = anti)
       prt%has_anti = .true.
    end if
    if (present (tex_name))  prt%tex_name = tex_name
    if (present (tex_anti))  prt%tex_anti = tex_anti
    if (present (spin_type))  prt%spin_type = spin_type
    if (present (isospin_type))  prt%isospin_type = isospin_type
    if (present (charge_type))  prt%charge_type = charge_type
    if (present (color_type))  prt%color_type = color_type
    if (present (mass_data)) then
       prt%mass_data => mass_data
       if (associated (mass_data)) then
          prt%mass_val => mass_data%get_real_ptr ()
       else
          prt%mass_val => null ()
       end if
    end if
    if (present (width_data)) then
       prt%width_data => width_data
       if (associated (width_data)) then
          prt%width_val => width_data%get_real_ptr ()
       else
          prt%width_val => null ()
       end if
    end if
    if (present (spin_type) .or. present (mass_data)) then
       call prt%set_multiplicity ()
    end if
    if (present (p_decay)) then
       if (allocated (prt%p_decay))  deallocate (prt%p_decay)
       if (size (p_decay) > 0) &
            allocate (prt%p_decay (size (p_decay)), source = p_decay)
    end if
    if (present (a_decay)) then
       if (allocated (prt%a_decay))  deallocate (prt%a_decay)
       if (size (a_decay) > 0) &
            allocate (prt%a_decay (size (a_decay)), source = a_decay)
    end if
  end subroutine field_data_set

\eatline
\nwendcode{}\nwbegindocs{38}Calculate the multiplicity given spin type and mass.
\nwenddocs{}\nwbegincode{39}\moddef{Model data: field data: TBP}\plusendmoddef
  procedure, private :: &
       set_multiplicity => field_data_set_multiplicity
\nwendcode{}\nwbegincode{40}\moddef{Model data: procedures}\plusendmoddef
  subroutine field_data_set_multiplicity (prt)
    class(field_data_t), intent(inout) :: prt
    if (prt%spin_type /= SCALAR) then
       if (associated (prt%mass_data)) then
          prt%multiplicity = prt%spin_type
       else if (prt%left_handed .or. prt%right_handed) then
          prt%multiplicity = 1
       else
          prt%multiplicity = 2
       end if
    end if
  end subroutine field_data_set_multiplicity

\eatline
\nwendcode{}\nwbegindocs{41}Set the mass/width value (not the pointer).  The mass/width pointer
must be allocated.
\nwenddocs{}\nwbegincode{42}\moddef{Model data: field data: TBP}\plusendmoddef
  procedure, private :: set_mass => field_data_set_mass
  procedure, private :: set_width => field_data_set_width
\nwendcode{}\nwbegincode{43}\moddef{Model data: procedures}\plusendmoddef
  subroutine field_data_set_mass (prt, mass)
    class(field_data_t), intent(inout) :: prt
    real(default), intent(in) :: mass
    if (associated (prt%mass_val))  prt%mass_val = mass
  end subroutine field_data_set_mass
    
  subroutine field_data_set_width (prt, width)
    class(field_data_t), intent(inout) :: prt
    real(default), intent(in) :: width
    if (associated (prt%width_val))  prt%width_val = width
  end subroutine field_data_set_width
    
\eatline
\nwendcode{}\nwbegindocs{44}Loose ends: name arrays should be allocated.
\nwenddocs{}\nwbegincode{45}\moddef{Model data: field data: TBP}\plusendmoddef
  procedure :: freeze => field_data_freeze
\nwendcode{}\nwbegincode{46}\moddef{Model data: procedures}\plusendmoddef
  elemental subroutine field_data_freeze (prt)
    class(field_data_t), intent(inout) :: prt
    if (.not. allocated (prt%name))  allocate (prt%name (0))
    if (.not. allocated (prt%anti))  allocate (prt%anti (0))
  end subroutine field_data_freeze

\eatline
\nwendcode{}\nwbegindocs{47}Output
\nwenddocs{}\nwbegincode{48}\moddef{Model data: field data: TBP}\plusendmoddef
  procedure :: write => field_data_write
\nwendcode{}\nwbegincode{49}\moddef{Model data: procedures}\plusendmoddef
  subroutine field_data_write (prt, unit)
    class(field_data_t), intent(in) :: prt
    integer, intent(in), optional :: unit
    integer :: u, i
    u = given_output_unit (unit);  if (u < 0)  return
    write (u, "(3x,A,1x,A)", advance="no") "particle", char (prt%longname)
    write (u, "(1x,I0)", advance="no") prt%pdg
    if (.not. prt%visible) write (u, "(2x,A)", advance="no") "invisible"
    if (prt%parton)  write (u, "(2x,A)", advance="no") "parton"
    if (prt%gauge)  write (u, "(2x,A)", advance="no") "gauge"
    if (prt%left_handed)  write (u, "(2x,A)", advance="no") "left"
    if (prt%right_handed)  write (u, "(2x,A)", advance="no") "right"
    write (u, *)
    write (u, "(5x,A)", advance="no") "name"
    if (allocated (prt%name)) then
       do i = 1, size (prt%name)
          write (u, "(1x,A)", advance="no")  '"' // char (prt%name(i)) // '"'
       end do
       write (u, *)
       if (prt%has_anti) then
          write (u, "(5x,A)", advance="no") "anti"
          do i = 1, size (prt%anti)
             write (u, "(1x,A)", advance="no")  '"' // char (prt%anti(i)) // '"'
          end do
          write (u, *)
       end if
       if (prt%tex_name /= "") then
          write (u, "(5x,A)")  &
               "tex_name " // '"' // char (prt%tex_name) // '"'
       end if
       if (prt%has_anti .and. prt%tex_anti /= "") then
          write (u, "(5x,A)")  &
               "tex_anti " // '"' // char (prt%tex_anti) // '"'
       end if
    else
       write (u, "(A)")  "???"
    end if
    write (u, "(5x,A)", advance="no") "spin "
    select case (mod (prt%spin_type - 1, 2))
    case (0);  write (u, "(I0)", advance="no") (prt%spin_type-1) / 2
    case default;  write (u, "(I0,A)", advance="no") prt%spin_type-1, "/2"
    end select
    ! write (u, "(2x,A,I1,A)") "! [multiplicity = ", prt%multiplicity, "]"
    if (abs (prt%isospin_type) /= 1) then
       write (u, "(2x,A)", advance="no") "isospin "
       select case (mod (abs (prt%isospin_type) - 1, 2))
       case (0);  write (u, "(I0)", advance="no") &
            sign (abs (prt%isospin_type) - 1, prt%isospin_type) / 2
       case default;  write (u, "(I0,A)", advance="no") &
            sign (abs (prt%isospin_type) - 1, prt%isospin_type), "/2"
       end select
    end if
    if (abs (prt%charge_type) /= 1) then
       write (u, "(2x,A)", advance="no") "charge "
       select case (mod (abs (prt%charge_type) - 1, 3))
       case (0);  write (u, "(I0)", advance="no") &
            sign (abs (prt%charge_type) - 1, prt%charge_type) / 3
       case default;  write (u, "(I0,A)", advance="no") &
            sign (abs (prt%charge_type) - 1, prt%charge_type), "/3"
       end select
    end if
    if (prt%color_type /= 1) then
       write (u, "(2x,A,I0)", advance="no") "color ", prt%color_type
    end if
    write (u, *)
    if (associated (prt%mass_data)) then
       write (u, "(5x,A)", advance="no") &
            "mass " // char (prt%mass_data%get_name ())
       if (associated (prt%width_data)) then
          write (u, "(2x,A)") &
               "width " // char (prt%width_data%get_name ())
       else
          write (u, *)
       end if
    end if
    call prt%write_decays (u)
  end subroutine field_data_write

\eatline
\nwendcode{}\nwbegindocs{50}Write decay and polarization data.
\nwenddocs{}\nwbegincode{51}\moddef{Model data: field data: TBP}\plusendmoddef
  procedure :: write_decays => field_data_write_decays
\nwendcode{}\nwbegincode{52}\moddef{Model data: procedures}\plusendmoddef
  subroutine field_data_write_decays (prt, unit)
    class(field_data_t), intent(in) :: prt
    integer, intent(in), optional :: unit
    integer :: u, i
    u = given_output_unit (unit)
    if (.not. prt%p_is_stable) then
       if (allocated (prt%p_decay)) then
          write (u, "(5x,A)", advance="no") "p_decay"
          do i = 1, size (prt%p_decay)
             write (u, "(1x,A)", advance="no")  char (prt%p_decay(i))
          end do
          if (prt%p_decays_isotropically) then
             write (u, "(1x,A)", advance="no")  "isotropic"
          else if (prt%p_decays_diagonal) then
             write (u, "(1x,A)", advance="no")  "diagonal"
          else if (prt%p_has_decay_helicity) then
             write (u, "(1x,A,I0)", advance="no")  "helicity = ", &
                  prt%p_decay_helicity
          end if
          write (u, *)
       end if
    else if (prt%p_polarized) then
       write (u, "(5x,A)")  "p_polarized"
    end if
    if (.not. prt%a_is_stable) then
       if (allocated (prt%a_decay)) then
          write (u, "(5x,A)", advance="no") "a_decay"
          do i = 1, size (prt%a_decay)
             write (u, "(1x,A)", advance="no")  char (prt%a_decay(i))
          end do
          if (prt%a_decays_isotropically) then
             write (u, "(1x,A)", advance="no")  "isotropic"
          else if (prt%a_decays_diagonal) then
             write (u, "(1x,A)", advance="no")  "diagonal"
          else if (prt%a_has_decay_helicity) then
             write (u, "(1x,A,I0)", advance="no")  "helicity = ", &
                  prt%a_decay_helicity
          end if
          write (u, *)
       end if
    else if (prt%a_polarized) then
       write (u, "(5x,A)")  "a_polarized"
    end if
  end subroutine field_data_write_decays
  
\eatline
\nwendcode{}\nwbegindocs{53}Screen version of output.
\nwenddocs{}\nwbegincode{54}\moddef{Model data: field data: TBP}\plusendmoddef
  procedure :: show => field_data_show
\nwendcode{}\nwbegincode{55}\moddef{Model data: procedures}\plusendmoddef
  subroutine field_data_show (prt, l, u)
    class(field_data_t), intent(in) :: prt
    integer, intent(in) :: l, u
    character(len=l) :: buffer
    integer :: i
    type(string_t), dimension(:), allocatable :: decay
    buffer = prt%get_name (.false.)
    write (u, "(4x,A,1x,I8)", advance="no")  buffer, &
         prt%get_pdg ()
    if (prt%is_polarized ()) then
       write (u, "(3x,A)")  "polarized"
    else if (.not. prt%is_stable ()) then
       write (u, "(3x,A)", advance="no")  "decays:"
       call prt%get_decays (decay)
       do i = 1, size (decay)
          write (u, "(1x,A)", advance="no")  char (decay(i))
       end do
       write (u, *)
    else
       write (u, *)
    end if
    if (prt%has_antiparticle ()) then
       buffer = prt%get_name (.true.)
       write (u, "(4x,A,1x,I8)", advance="no")  buffer, &
            prt%get_pdg_anti ()
       if (prt%is_polarized (.true.)) then
          write (u, "(3x,A)")  "polarized"
       else if (.not. prt%is_stable (.true.)) then
          write (u, "(3x,A)", advance="no")  "decays:"
          call prt%get_decays (decay, .true.)
          do i = 1, size (decay)
             write (u, "(1x,A)", advance="no")  char (decay(i))
          end do
          write (u, *)
       else
          write (u, *)
       end if
    end if
  end subroutine field_data_show
  
\eatline
\nwendcode{}\nwbegindocs{56}Retrieve data:
\nwenddocs{}\nwbegincode{57}\moddef{Model data: field data: TBP}\plusendmoddef
  procedure :: get_pdg => field_data_get_pdg
  procedure :: get_pdg_anti => field_data_get_pdg_anti
\nwendcode{}\nwbegincode{58}\moddef{Model data: procedures}\plusendmoddef
  elemental function field_data_get_pdg (prt) result (pdg)
    integer :: pdg
    class(field_data_t), intent(in) :: prt
    pdg = prt%pdg
  end function field_data_get_pdg

  elemental function field_data_get_pdg_anti (prt) result (pdg)
    integer :: pdg
    class(field_data_t), intent(in) :: prt
    if (prt%has_anti) then
       pdg = - prt%pdg
    else
       pdg = prt%pdg
    end if
  end function field_data_get_pdg_anti

\eatline
\nwendcode{}\nwbegindocs{59}Predicates:
\nwenddocs{}\nwbegincode{60}\moddef{Model data: field data: TBP}\plusendmoddef
  procedure :: is_visible => field_data_is_visible
  procedure :: is_parton => field_data_is_parton
  procedure :: is_gauge => field_data_is_gauge
  procedure :: is_left_handed => field_data_is_left_handed
  procedure :: is_right_handed => field_data_is_right_handed
  procedure :: has_antiparticle => field_data_has_antiparticle
  procedure :: is_stable => field_data_is_stable
  procedure :: get_decays => field_data_get_decays
  procedure :: decays_isotropically => field_data_decays_isotropically
  procedure :: decays_diagonal => field_data_decays_diagonal
  procedure :: has_decay_helicity => field_data_has_decay_helicity
  procedure :: decay_helicity => field_data_decay_helicity
  procedure :: is_polarized => field_data_is_polarized
\nwendcode{}\nwbegincode{61}\moddef{Model data: procedures}\plusendmoddef
  elemental function field_data_is_visible (prt) result (flag)
    logical :: flag
    class(field_data_t), intent(in) :: prt
    flag = prt%visible
  end function field_data_is_visible

  elemental function field_data_is_parton (prt) result (flag)
    logical :: flag
    class(field_data_t), intent(in) :: prt
    flag = prt%parton
  end function field_data_is_parton

  elemental function field_data_is_gauge (prt) result (flag)
    logical :: flag
    class(field_data_t), intent(in) :: prt
    flag = prt%gauge
  end function field_data_is_gauge

  elemental function field_data_is_left_handed (prt) result (flag)
    logical :: flag
    class(field_data_t), intent(in) :: prt
    flag = prt%left_handed
  end function field_data_is_left_handed

  elemental function field_data_is_right_handed (prt) result (flag)
    logical :: flag
    class(field_data_t), intent(in) :: prt
    flag = prt%right_handed
  end function field_data_is_right_handed

  elemental function field_data_has_antiparticle (prt) result (flag)
    logical :: flag
    class(field_data_t), intent(in) :: prt
    flag = prt%has_anti
  end function field_data_has_antiparticle

  elemental function field_data_is_stable (prt, anti) result (flag)
    logical :: flag
    class(field_data_t), intent(in) :: prt
    logical, intent(in), optional :: anti
    if (present (anti)) then
       if (anti) then
          flag = prt%a_is_stable
       else
          flag = prt%p_is_stable
       end if
    else
       flag = prt%p_is_stable
    end if
  end function field_data_is_stable

  subroutine field_data_get_decays (prt, decay, anti)
    class(field_data_t), intent(in) :: prt
    type(string_t), dimension(:), intent(out), allocatable :: decay
    logical, intent(in), optional :: anti
    if (present (anti)) then
       if (anti) then
          allocate (decay (size (prt%a_decay)), source = prt%a_decay)
       else
          allocate (decay (size (prt%p_decay)), source = prt%p_decay)
       end if
    else
       allocate (decay (size (prt%p_decay)), source = prt%p_decay)
    end if
  end subroutine field_data_get_decays

  elemental function field_data_decays_isotropically &
       (prt, anti) result (flag)
    logical :: flag
    class(field_data_t), intent(in) :: prt
    logical, intent(in), optional :: anti
    if (present (anti)) then
       if (anti) then
          flag = prt%a_decays_isotropically
       else
          flag = prt%p_decays_isotropically
       end if
    else
       flag = prt%p_decays_isotropically
    end if
  end function field_data_decays_isotropically

  elemental function field_data_decays_diagonal &
       (prt, anti) result (flag)
    logical :: flag
    class(field_data_t), intent(in) :: prt
    logical, intent(in), optional :: anti
    if (present (anti)) then
       if (anti) then
          flag = prt%a_decays_diagonal
       else
          flag = prt%p_decays_diagonal
       end if
    else
       flag = prt%p_decays_diagonal
    end if
  end function field_data_decays_diagonal

  elemental function field_data_has_decay_helicity &
       (prt, anti) result (flag)
    logical :: flag
    class(field_data_t), intent(in) :: prt
    logical, intent(in), optional :: anti
    if (present (anti)) then
       if (anti) then
          flag = prt%a_has_decay_helicity
       else
          flag = prt%p_has_decay_helicity
       end if
    else
       flag = prt%p_has_decay_helicity
    end if
  end function field_data_has_decay_helicity

  elemental function field_data_decay_helicity &
       (prt, anti) result (hel)
    integer :: hel
    class(field_data_t), intent(in) :: prt
    logical, intent(in), optional :: anti
    if (present (anti)) then
       if (anti) then
          hel = prt%a_decay_helicity
       else
          hel = prt%p_decay_helicity
       end if
    else
       hel = prt%p_decay_helicity
    end if
  end function field_data_decay_helicity

  elemental function field_data_is_polarized (prt, anti) result (flag)
    logical :: flag
    class(field_data_t), intent(in) :: prt
    logical, intent(in), optional :: anti
    logical :: a
    if (present (anti)) then
       a = anti
    else
       a = .false.
    end if
    if (a) then
       flag = prt%a_polarized
    else
       flag = prt%p_polarized
    end if
  end function field_data_is_polarized
       
\eatline
\eatline
\eatline
\eatline
\eatline
\eatline
\eatline
\eatline
\eatline
\eatline
\nwendcode{}\nwbegindocs{62}Names.  Return the first name in the list (or the first antiparticle name)
\nwenddocs{}\nwbegincode{63}\moddef{Model data: field data: TBP}\plusendmoddef
  procedure :: get_longname => field_data_get_longname
  procedure :: get_name => field_data_get_name
  procedure :: get_name_array => field_data_get_name_array
\nwendcode{}\nwbegincode{64}\moddef{Model data: procedures}\plusendmoddef
  pure function field_data_get_longname (prt) result (name)
    type(string_t) :: name
    class(field_data_t), intent(in) :: prt
    name = prt%longname
  end function field_data_get_longname

  pure function field_data_get_name (prt, is_antiparticle) result (name)
    type(string_t) :: name
    class(field_data_t), intent(in) :: prt
    logical, intent(in) :: is_antiparticle
    name = prt%longname
    if (is_antiparticle) then
       if (prt%has_anti) then
          if (allocated (prt%anti)) then
             if (size(prt%anti) > 0) name = prt%anti(1)
          end if
       else
          if (allocated (prt%name)) then
             if (size (prt%name) > 0) name = prt%name(1)
          end if
       end if
    else
       if (allocated (prt%name)) then
          if (size (prt%name) > 0) name = prt%name(1)
       end if
    end if
  end function field_data_get_name

  subroutine field_data_get_name_array (prt, is_antiparticle, name)
    class(field_data_t), intent(in) :: prt
    logical, intent(in) :: is_antiparticle
    type(string_t), dimension(:), allocatable, intent(inout) :: name
    if (allocated (name))  deallocate (name)
    if (is_antiparticle) then
       if (prt%has_anti) then
          allocate (name (size (prt%anti)))
          name = prt%anti
       else
          allocate (name (0))
       end if
    else
       allocate (name (size (prt%name)))
       name = prt%name
    end if
  end subroutine field_data_get_name_array

\eatline
\nwendcode{}\nwbegindocs{65}Same for the \TeX\ name.
\nwenddocs{}\nwbegincode{66}\moddef{Model data: field data: TBP}\plusendmoddef
  procedure :: get_tex_name => field_data_get_tex_name
\nwendcode{}\nwbegincode{67}\moddef{Model data: procedures}\plusendmoddef
  elemental function field_data_get_tex_name &
       (prt, is_antiparticle) result (name)
    type(string_t) :: name
    class(field_data_t), intent(in) :: prt
    logical, intent(in) :: is_antiparticle
    if (is_antiparticle) then
       if (prt%has_anti) then
          name = prt%tex_anti
       else
          name = prt%tex_name
       end if
    else
       name = prt%tex_name
    end if
    if (name == "")  name = prt%get_name (is_antiparticle)
  end function field_data_get_tex_name

\eatline
\nwendcode{}\nwbegindocs{68}Check if any of the field names matches the given string.
\nwenddocs{}\nwbegincode{69}\moddef{Model data: field data: TBP}\plusendmoddef
  procedure, private :: matches_name => field_data_matches_name
\nwendcode{}\nwbegincode{70}\moddef{Model data: procedures}\plusendmoddef
  function field_data_matches_name (field, name, is_antiparticle) result (flag)
    class(field_data_t), intent(in) :: field
    type(string_t), intent(in) :: name
    logical, intent(in) :: is_antiparticle
    logical :: flag
    if (is_antiparticle) then
       if (field%has_anti) then
          flag = any (name == field%anti)
       else
          flag = .false.
       end if
    else
       flag = name == field%longname .or. any (name == field%name)
    end if
  end function field_data_matches_name
  
\eatline
\nwendcode{}\nwbegindocs{71}Quantum numbers
\nwenddocs{}\nwbegincode{72}\moddef{Model data: field data: TBP}\plusendmoddef
  procedure :: get_spin_type => field_data_get_spin_type
  procedure :: get_multiplicity => field_data_get_multiplicity
  procedure :: get_isospin_type => field_data_get_isospin_type
  procedure :: get_charge_type => field_data_get_charge_type
  procedure :: get_color_type => field_data_get_color_type
\nwendcode{}\nwbegincode{73}\moddef{Model data: procedures}\plusendmoddef
  elemental function field_data_get_spin_type (prt) result (type)
    integer :: type
    class(field_data_t), intent(in) :: prt
    type = prt%spin_type
  end function field_data_get_spin_type

  elemental function field_data_get_multiplicity (prt) result (type)
    integer :: type
    class(field_data_t), intent(in) :: prt
    type = prt%multiplicity
  end function field_data_get_multiplicity

  elemental function field_data_get_isospin_type (prt) result (type)
    integer :: type
    class(field_data_t), intent(in) :: prt
    type = prt%isospin_type
  end function field_data_get_isospin_type

  elemental function field_data_get_charge_type (prt) result (type)
    integer :: type
    class(field_data_t), intent(in) :: prt
    type = prt%charge_type
  end function field_data_get_charge_type

  elemental function field_data_get_color_type (prt) result (type)
    integer :: type
    class(field_data_t), intent(in) :: prt
    type = prt%color_type
  end function field_data_get_color_type

\eatline
\eatline
\eatline
\eatline
\eatline
\nwendcode{}\nwbegindocs{74}In the MSSM, neutralinos can have a negative mass.   This is
relevant for computing matrix elements.  However, within the
\whizard\ main program we are interested only in kinematics, therefore
we return the absolute value of the particle mass.  If desired, we can
extract the sign separately.
\nwenddocs{}\nwbegincode{75}\moddef{Model data: field data: TBP}\plusendmoddef
  procedure :: get_charge => field_data_get_charge
  procedure :: get_isospin => field_data_get_isospin
  procedure :: get_mass => field_data_get_mass
  procedure :: get_mass_sign => field_data_get_mass_sign
  procedure :: get_width => field_data_get_width
\nwendcode{}\nwbegincode{76}\moddef{Model data: procedures}\plusendmoddef
  elemental function field_data_get_charge (prt) result (charge)
    real(default) :: charge
    class(field_data_t), intent(in) :: prt
    if (prt%charge_type /= 0) then
       charge = real (sign ((abs(prt%charge_type) - 1), &
             prt%charge_type), default) / 3
    else
       charge = 0
    end if
  end function field_data_get_charge

  elemental function field_data_get_isospin (prt) result (isospin)
    real(default) :: isospin
    class(field_data_t), intent(in) :: prt
    if (prt%isospin_type /= 0) then
       isospin = real (sign (abs(prt%isospin_type) - 1, &
            prt%isospin_type), default) / 2
    else
       isospin = 0
    end if
  end function field_data_get_isospin

  elemental function field_data_get_mass (prt) result (mass)
    real(default) :: mass
    class(field_data_t), intent(in) :: prt
    if (associated (prt%mass_val)) then
       mass = abs (prt%mass_val)
    else
       mass = 0
    end if
  end function field_data_get_mass

  elemental function field_data_get_mass_sign (prt) result (sgn)
    integer :: sgn
    class(field_data_t), intent(in) :: prt
    if (associated (prt%mass_val)) then
       sgn = sign (1._default, prt%mass_val)
    else
       sgn = 0
    end if
  end function field_data_get_mass_sign

  elemental function field_data_get_width (prt) result (width)
    real(default) :: width
    class(field_data_t), intent(in) :: prt
    if (associated (prt%width_val)) then
       width = prt%width_val
    else
       width = 0
    end if
  end function field_data_get_width
  
\eatline
\eatline
\eatline
\nwendcode{}\nwbegindocs{77}\nwdocspar
\subsection{Vertex data}
The vertex object contains an array of particle-data pointers, for
which we need a separate type.  (We could use the flavor type defined
in another module.)

The program does not (yet?) make use of vertex definitions, so they
are not stored here.
\nwenddocs{}\nwbegincode{78}\moddef{Model data: types}\plusendmoddef
  type :: field_data_p
     type(field_data_t), pointer :: p => null ()
  end type field_data_p

\eatline
\nwendcode{}\nwbegincode{79}\moddef{Model data: types}\plusendmoddef
  type :: vertex_t
     private
     logical :: trilinear
     integer, dimension(:), allocatable :: pdg
     type(field_data_p), dimension(:), allocatable :: prt
   contains
   \LA{}Model data: vertex: TBP\RA{}
  end type vertex_t

\eatline
\nwendcode{}\nwbegincode{80}\moddef{Model data: vertex: TBP}\endmoddef
  procedure :: write => vertex_write
\nwendcode{}\nwbegincode{81}\moddef{Model data: procedures}\plusendmoddef
  subroutine vertex_write (vtx, unit)
    class(vertex_t), intent(in) :: vtx
    integer, intent(in), optional :: unit
    integer :: u, i
    u = given_output_unit (unit)
    write (u, "(3x,A)", advance="no")  "vertex"
    do i = 1, size (vtx%prt)
       if (associated (vtx%prt(i)%p)) then
          write (u, "(1x,A)", advance="no") &
               '"' // char (vtx%prt(i)%p%get_name (vtx%pdg(i) < 0)) &
                   // '"'
       else
          write (u, "(1x,I7)", advance="no") vtx%pdg(i)
       end if
    end do
    write (u, *)
  end subroutine vertex_write

\eatline
\nwendcode{}\nwbegindocs{82}Initialize using PDG codes.  The model is used for finding particle
data pointers associated with the pdg codes.
\nwenddocs{}\nwbegincode{83}\moddef{Model data: vertex: TBP}\plusendmoddef
  procedure :: init => vertex_init
\nwendcode{}\nwbegincode{84}\moddef{Model data: procedures}\plusendmoddef
  subroutine vertex_init (vtx, pdg, model)
    class(vertex_t), intent(out) :: vtx
    integer, dimension(:), intent(in) :: pdg
    type(model_data_t), intent(in), target, optional :: model
    integer :: i
    allocate (vtx%pdg (size (pdg)))
    allocate (vtx%prt (size (pdg)))
    vtx%trilinear = size (pdg) == 3
    vtx%pdg = pdg
    if (present (model)) then
       do i = 1, size (pdg)
          vtx%prt(i)%p => model%get_field_ptr (pdg(i))
       end do
    end if
  end subroutine vertex_init

\eatline
\nwendcode{}\nwbegindocs{85}Copy vertex: we must reassign the field-data pointer to a new model.
\nwenddocs{}\nwbegincode{86}\moddef{Model data: vertex: TBP}\plusendmoddef
  procedure :: copy_from => vertex_copy_from
\nwendcode{}\nwbegincode{87}\moddef{Model data: procedures}\plusendmoddef
  subroutine vertex_copy_from (vtx, old_vtx, new_model)
    class(vertex_t), intent(out) :: vtx
    class(vertex_t), intent(in) :: old_vtx
    type(model_data_t), intent(in), target, optional :: new_model
    call vtx%init (old_vtx%pdg, new_model)
  end subroutine vertex_copy_from
    
\eatline
\nwendcode{}\nwbegindocs{88}Single-particle lookup: Given a particle code, we return matching
codes if present, otherwise zero.  Actually, we return the
antiparticles of the matching codes, as appropriate for computing
splittings.
\nwenddocs{}\nwbegincode{89}\moddef{Model data: vertex: TBP}\plusendmoddef
  procedure :: get_match => vertex_get_match
\nwendcode{}\nwbegincode{90}\moddef{Model data: procedures}\plusendmoddef
  subroutine vertex_get_match (vtx, pdg1, pdg2)
    class(vertex_t), intent(in) :: vtx
    integer, intent(in) :: pdg1
    integer, dimension(:), allocatable, intent(out) :: pdg2
    integer :: i, j
    do i = 1, size (vtx%pdg)
       if (vtx%pdg(i) == pdg1) then
          allocate (pdg2 (size (vtx%pdg) - 1))
          do j = 1, i-1
             pdg2(j) = anti (j)
          end do
          do j = i, size (pdg2)
             pdg2(j) = anti (j+1)
          end do
          exit
       end if
    end do
  contains
    function anti (i) result (pdg)
      integer, intent(in) :: i
      integer :: pdg
      if (vtx%prt(i)%p%has_antiparticle ()) then
         pdg = - vtx%pdg(i)
      else
         pdg = vtx%pdg(i)
      end if
    end function anti
  end subroutine vertex_get_match

\eatline
\nwendcode{}\nwbegindocs{91}To access this from the outside, we create an iterator.  The iterator has
the sole purpose of returning the matching particles for a given array of PDG
codes.
\nwenddocs{}\nwbegincode{92}\moddef{Model data: public}\plusendmoddef
  public :: vertex_iterator_t
\nwendcode{}\nwbegincode{93}\moddef{Model data: types}\plusendmoddef
  type :: vertex_iterator_t
     private
     class(model_data_t), pointer :: model => null ()
     integer, dimension(:), allocatable :: pdg
     integer :: vertex_index = 0
   contains
     procedure :: init => vertex_iterator_init
     procedure :: get_next_match => vertex_iterator_get_next_match
  end type vertex_iterator_t
  
\eatline
\nwendcode{}\nwbegindocs{94}We initialize the iterator for a particular model with the {\tt{}pdg} index of
the particle we are looking at.
\nwenddocs{}\nwbegincode{95}\moddef{Model data: procedures}\plusendmoddef
  subroutine vertex_iterator_init (it, model, pdg)
    class(vertex_iterator_t), intent(out) :: it
    class(model_data_t), intent(in), target :: model
    integer, dimension(:), intent(in) :: pdg
    it%model => model
    allocate (it%pdg (size (pdg)), source = pdg)
  end subroutine vertex_iterator_init
  
  subroutine vertex_iterator_get_next_match (it, pdg_match)
    class(vertex_iterator_t), intent(inout) :: it
    integer, dimension(:), allocatable, intent(out) :: pdg_match
    integer :: i, j
    do i = it%vertex_index + 1, size (it%model%vtx)
       do j = 1, size (it%pdg)
          call vertex_get_match (it%model%vtx(i), it%pdg(j), pdg_match)
          if (allocated (pdg_match)) then
             it%vertex_index = i
             return
          end if
       end do
    end do
    it%vertex_index = 0
  end subroutine vertex_iterator_get_next_match
  
\eatline
\nwendcode{}\nwbegindocs{96}\nwdocspar
\subsection{Vertex lookup table}
The vertex lookup table is a hash table: given two particle codes, we
check which codes are allowed for the third one.

The size of the hash table should be large enough that collisions are
rare.  We first select a size based on the number of vertices
(multiplied by six because all permutations count), with some margin,
and then choose the smallest integer power of two larger than this.
\nwenddocs{}\nwbegincode{97}\moddef{Model data: parameters}\endmoddef
  integer, parameter :: VERTEX_TABLE_SCALE_FACTOR = 60
\eatline
\nwendcode{}\nwbegincode{98}\moddef{Model data: procedures}\plusendmoddef
  function vertex_table_size (n_vtx) result (n)
    integer(i32) :: n
    integer, intent(in) :: n_vtx
    integer :: i, s
    s = VERTEX_TABLE_SCALE_FACTOR * n_vtx
    n = 1
    do i = 1, 31
       n = ishft (n, 1)
       s = ishft (s,-1)
       if (s == 0)  exit
    end do
  end function vertex_table_size

\eatline
\nwendcode{}\nwbegindocs{99}The specific hash function takes two particle codes (arbitrary
integers) and returns a 32-bit integer.  It makes use of the universal
function {\tt{}hash} which operates on a byte array.
\nwenddocs{}\nwbegincode{100}\moddef{Model data: procedures}\plusendmoddef
  function hash2 (pdg1, pdg2)
    integer(i32) :: hash2
    integer, intent(in) :: pdg1, pdg2
    integer(i8), dimension(1) :: mold
    hash2 = hash (transfer ([pdg1, pdg2], mold))
  end function hash2

\eatline
\nwendcode{}\nwbegindocs{101}Each entry in the vertex table stores the two particle codes and an
array of possibilities for the third code.
\nwenddocs{}\nwbegincode{102}\moddef{Model data: types}\plusendmoddef
  type :: vertex_table_entry_t
     private
     integer :: pdg1 = 0, pdg2 = 0
     integer :: n = 0
     integer, dimension(:), allocatable :: pdg3
  end type vertex_table_entry_t

\eatline
\nwendcode{}\nwbegindocs{103}The vertex table:
\nwenddocs{}\nwbegincode{104}\moddef{Model data: types}\plusendmoddef
  type :: vertex_table_t
     type(vertex_table_entry_t), dimension(:), allocatable :: entry
     integer :: n_collisions = 0
     integer(i32) :: mask
   contains
   \LA{}Model data: vertex table: TBP\RA{}
  end type vertex_table_t

\eatline
\nwendcode{}\nwbegindocs{105}Output.
\nwenddocs{}\nwbegincode{106}\moddef{Model data: vertex table: TBP}\endmoddef
  procedure :: write => vertex_table_write
\nwendcode{}\nwbegincode{107}\moddef{Model data: procedures}\plusendmoddef
  subroutine vertex_table_write (vt, unit)
    class(vertex_table_t), intent(in) :: vt
    integer, intent(in), optional :: unit
    integer :: u, i
    character(9) :: size_pdg3
    u = given_output_unit (unit)
    write (u, "(A)") "vertex hash table:"
    write (u, "(A,I7)") "  size = ", size (vt%entry)
    write (u, "(A,I7)") "  used = ", count (vt%entry%n /= 0)
    write (u, "(A,I7)") "  coll = ", vt%n_collisions
    do i = lbound (vt%entry, 1), ubound (vt%entry, 1)
       if (vt%entry(i)%n /= 0) then
          write (size_pdg3, "(I7)") size (vt%entry(i)%pdg3)
          write (u, "(A,1x,I7,1x,A,2(1x,I7),A," // &
               size_pdg3 // "(1x,I7))")  &
               "  ", i, ":", vt%entry(i)%pdg1, &
               vt%entry(i)%pdg2, "->", vt%entry(i)%pdg3
       end if
    end do
  end subroutine vertex_table_write

\eatline
\nwendcode{}\nwbegindocs{108}Initializing the vertex table: This is done in two passes.  First,
we scan all permutations for all vertices and count the number of
entries in each bucket of the hashtable.  Then, the buckets are
allocated accordingly and filled.

Collision resolution is done by just incrementing the hash value until
an empty bucket is found.  The vertex table size is fixed, since we
know from the beginning the number of entries.
\nwenddocs{}\nwbegincode{109}\moddef{Model data: vertex table: TBP}\plusendmoddef
  procedure :: init => vertex_table_init
\nwendcode{}\nwbegincode{110}\moddef{Model data: procedures}\plusendmoddef
  subroutine vertex_table_init (vt, prt, vtx)
    class(vertex_table_t), intent(out) :: vt
    type(field_data_t), dimension(:), intent(in) :: prt
    type(vertex_t), dimension(:), intent(in) :: vtx
    integer :: n_vtx, vt_size, i, p1, p2, p3
    integer, dimension(3) :: p
    n_vtx = size (vtx)
    vt_size = vertex_table_size (count (vtx%trilinear))
    vt%mask = vt_size - 1
    allocate (vt%entry (0:vt_size-1))
    do i = 1, n_vtx
       if (vtx(i)%trilinear) then
          p = vtx(i)%pdg
          p1 = p(1);  p2 = p(2)
          call create (hash2 (p1, p2))
          if (p(2) /= p(3)) then
             p2 = p(3)
             call create (hash2 (p1, p2))
          end if
          if (p(1) /= p(2)) then
             p1 = p(2);  p2 = p(1)
             call create (hash2 (p1, p2))
             if (p(1) /= p(3)) then
                p2 = p(3)
                call create (hash2 (p1, p2))
             end if
          end if
          if (p(1) /= p(3)) then
             p1 = p(3);  p2 = p(1)
             call create (hash2 (p1, p2))
             if (p(1) /= p(2)) then
                p2 = p(2)
                call create (hash2 (p1, p2))
             end if
          end if
       end if
    end do
    do i = 0, vt_size - 1
       allocate (vt%entry(i)%pdg3 (vt%entry(i)%n))
    end do
    vt%entry%n = 0
    do i = 1, n_vtx
       if (vtx(i)%trilinear) then
          p = vtx(i)%pdg
          p1 = p(1);  p2 = p(2);  p3 = p(3)
          call register (hash2 (p1, p2))
          if (p(2) /= p(3)) then
             p2 = p(3);  p3 = p(2)
             call register (hash2 (p1, p2))
          end if
          if (p(1) /= p(2)) then
             p1 = p(2);  p2 = p(1);  p3 = p(3)
             call register (hash2 (p1, p2))
             if (p(1) /= p(3)) then
                p2 = p(3);  p3 = p(1)
                call register (hash2 (p1, p2))
             end if
          end if
          if (p(1) /= p(3)) then
             p1 = p(3);  p2 = p(1);  p3 = p(2)
             call register (hash2 (p1, p2))
             if (p(1) /= p(2)) then
                p2 = p(2);  p3 = p(1)
                call register (hash2 (p1, p2))
             end if
          end if
       end if
    end do
  contains
    recursive subroutine create (hashval)
      integer(i32), intent(in) :: hashval
      integer :: h
      h = iand (hashval, vt%mask)
      if (vt%entry(h)%n == 0) then
         vt%entry(h)%pdg1 = p1
         vt%entry(h)%pdg2 = p2
         vt%entry(h)%n = 1
      else if (vt%entry(h)%pdg1 == p1 .and. vt%entry(h)%pdg2 == p2) then
         vt%entry(h)%n = vt%entry(h)%n + 1
      else
         vt%n_collisions = vt%n_collisions + 1
         call create (hashval + 1)
      end if
    end subroutine create
    recursive subroutine register (hashval)
      integer(i32), intent(in) :: hashval
      integer :: h
      h = iand (hashval, vt%mask)
      if (vt%entry(h)%pdg1 == p1 .and. vt%entry(h)%pdg2 == p2) then
         vt%entry(h)%n = vt%entry(h)%n + 1
         vt%entry(h)%pdg3(vt%entry(h)%n) = p3
      else
         call register (hashval + 1)
      end if
    end subroutine register
  end subroutine vertex_table_init

\eatline
\nwendcode{}\nwbegindocs{111}Return the array of particle codes that match the given pair.
\nwenddocs{}\nwbegincode{112}\moddef{Model data: vertex table: TBP}\plusendmoddef
  procedure :: match => vertex_table_match
\nwendcode{}\nwbegincode{113}\moddef{Model data: procedures}\plusendmoddef
  subroutine vertex_table_match (vt, pdg1, pdg2, pdg3)
    class(vertex_table_t), intent(in) :: vt
    integer, intent(in) :: pdg1, pdg2
    integer, dimension(:), allocatable, intent(out) :: pdg3
    call match (hash2 (pdg1, pdg2))
  contains
    recursive subroutine match (hashval)
      integer(i32), intent(in) :: hashval
      integer :: h
      h = iand (hashval, vt%mask)
      if (vt%entry(h)%n == 0) then
         allocate (pdg3 (0))
      else if (vt%entry(h)%pdg1 == pdg1 .and. vt%entry(h)%pdg2 == pdg2) then
         allocate (pdg3 (size (vt%entry(h)%pdg3)))
         pdg3 = vt%entry(h)%pdg3
      else
         call match (hashval + 1)
      end if
    end subroutine match
  end subroutine vertex_table_match

\eatline
\nwendcode{}\nwbegindocs{114}Return true if the triplet is represented as a vertex.
\nwenddocs{}\nwbegincode{115}\moddef{Model data: vertex table: TBP}\plusendmoddef
  procedure :: check => vertex_table_check
\nwendcode{}\nwbegincode{116}\moddef{Model data: procedures}\plusendmoddef
  function vertex_table_check (vt, pdg1, pdg2, pdg3) result (flag)
    class(vertex_table_t), intent(in) :: vt
    integer, intent(in) :: pdg1, pdg2, pdg3
    logical :: flag
    flag = check (hash2 (pdg1, pdg2))
  contains
    recursive function check (hashval) result (flag)
      integer(i32), intent(in) :: hashval
      integer :: h
      logical :: flag
      h = iand (hashval, vt%mask)
      if (vt%entry(h)%n == 0) then
         flag = .false.
      else if (vt%entry(h)%pdg1 == pdg1 .and. vt%entry(h)%pdg2 == pdg2) then
         flag = any (vt%entry(h)%pdg3 == pdg3)
      else
         flag = check (hashval + 1)
      end if
    end function check
  end function vertex_table_check

\eatline
\nwendcode{}\nwbegindocs{117}\nwdocspar
\subsection{Model Data Record}
This type collects the model data as defined above.

We deliberately implement the parameter arrays as pointer arrays.  We
thus avoid keeping track of TARGET attributes.
\nwenddocs{}\nwbegincode{118}\moddef{Model data: public}\plusendmoddef
  public :: model_data_t
\nwendcode{}\nwbegincode{119}\moddef{Model data: types}\plusendmoddef
  type :: model_data_t
     private
     type(string_t) :: name
     type(modelpar_real_t), dimension(:), pointer :: par_real => null ()
     type(modelpar_complex_t), dimension(:), pointer :: par_complex => null ()
     type(field_data_t), dimension(:), allocatable :: field
     type(vertex_t), dimension(:), allocatable :: vtx
     type(vertex_table_t) :: vt
   contains
   \LA{}Model data: model data: TBP\RA{}
  end type model_data_t
     
\eatline
\nwendcode{}\nwbegindocs{120}Finalizer, deallocate pointer arrays.
\nwenddocs{}\nwbegincode{121}\moddef{Model data: model data: TBP}\endmoddef
  procedure :: final => model_data_final
\nwendcode{}\nwbegincode{122}\moddef{Model data: procedures}\plusendmoddef
  subroutine model_data_final (model)
    class(model_data_t), intent(inout) :: model
    deallocate (model%par_real)
    deallocate (model%par_complex)
  end subroutine model_data_final
  
\eatline
\nwendcode{}\nwbegindocs{123}Output.  The signature matches the signature of the high-level
{\tt{}model{\char95}write} procedure, so some of the options don't actually apply.
\nwenddocs{}\nwbegincode{124}\moddef{Model data: model data: TBP}\plusendmoddef
  procedure :: write => model_data_write
\nwendcode{}\nwbegincode{125}\moddef{Model data: procedures}\plusendmoddef
  subroutine model_data_write (model, unit, verbose, &
       show_md5sum, show_variables, show_parameters, &
       show_particles, show_vertices)
    class(model_data_t), intent(in) :: model
    integer, intent(in), optional :: unit
    logical, intent(in), optional :: verbose
    logical, intent(in), optional :: show_md5sum
    logical, intent(in), optional :: show_variables
    logical, intent(in), optional :: show_parameters
    logical, intent(in), optional :: show_particles
    logical, intent(in), optional :: show_vertices
    logical :: show_par, show_prt, show_vtx
    integer :: u, i
    u = given_output_unit (unit)
    show_par = .true.;  if (present (show_parameters)) &
         show_par = show_parameters
    show_prt = .true.;  if (present (show_particles)) &
         show_prt = show_particles
    show_vtx = .true.;  if (present (show_vertices)) &
         show_vtx = show_vertices
    if (show_par) then
       do i = 1, size (model%par_real)
          call model%par_real(i)%write (u)
          write (u, *)
       end do
       do i = 1, size (model%par_complex)
          call model%par_complex(i)%write (u)
          write (u, *)
       end do
    end if
    if (show_prt) then
       write (u, *)
       call model%write_fields (u)
    end if
    if (show_vtx) then
       write (u, *)
       call model%write_vertices (u, verbose)
    end if
  end subroutine model_data_write
  
\eatline
\nwendcode{}\nwbegindocs{126}Initialize, allocating pointer arrays.  The second version makes a
deep copy.
\nwenddocs{}\nwbegincode{127}\moddef{Model data: model data: TBP}\plusendmoddef
  generic :: init => model_data_init
  procedure, private :: model_data_init
\nwendcode{}\nwbegincode{128}\moddef{Model data: procedures}\plusendmoddef
  subroutine model_data_init (model, name, &
       n_par_real, n_par_complex, n_field, n_vtx)
    class(model_data_t), intent(out) :: model
    type(string_t), intent(in) :: name
    integer, intent(in) :: n_par_real, n_par_complex
    integer, intent(in) :: n_field
    integer, intent(in) :: n_vtx
    model%name = name
    allocate (model%par_real (n_par_real))
    allocate (model%par_complex (n_par_complex))
    allocate (model%field (n_field))
    allocate (model%vtx (n_vtx))
  end subroutine model_data_init
  
\eatline
\nwendcode{}\nwbegindocs{129}Complete model data initialization.
\nwenddocs{}\nwbegincode{130}\moddef{Model data: model data: TBP}\plusendmoddef
  procedure :: freeze_fields => model_data_freeze_fields
\nwendcode{}\nwbegincode{131}\moddef{Model data: procedures}\plusendmoddef
  subroutine model_data_freeze_fields (model)
    class(model_data_t), intent(inout) :: model
    call model%field%freeze ()
  end subroutine model_data_freeze_fields

\eatline
\nwendcode{}\nwbegindocs{132}Deep copy.  The new model should already be initialized, so we do
not allocate memory.
\nwenddocs{}\nwbegincode{133}\moddef{Model data: model data: TBP}\plusendmoddef
  procedure :: copy_from => model_data_copy
\nwendcode{}\nwbegincode{134}\moddef{Model data: procedures}\plusendmoddef
  subroutine model_data_copy (model, src)
    class(model_data_t), intent(inout), target :: model
    class(model_data_t), intent(in), target :: src
    class(modelpar_data_t), pointer :: data, src_data
    integer :: i
    model%par_real = src%par_real
    model%par_complex = src%par_complex
    do i = 1, size (src%field)
       associate (field => model%field(i), src_field => src%field(i))
         call field%init (src_field%get_longname (), src_field%get_pdg ())
         call field%copy_from (src_field)
         src_data => src_field%mass_data
         if (associated (src_data)) then
            data => model%get_par_data_ptr (src_data%get_name ())
            call field%set (mass_data = data)
         end if
         src_data => src_field%width_data
         if (associated (src_data)) then
            data => model%get_par_data_ptr (src_data%get_name ())
            call field%set (width_data = data)
         end if
         call field%set_multiplicity ()
       end associate
    end do
    do i = 1, size (src%vtx)
       call model%vtx(i)%copy_from (src%vtx(i), model)
    end do
    call model%freeze_vertices ()
  end subroutine model_data_copy

\eatline
\nwendcode{}\nwbegindocs{135}Return the model name.
\nwenddocs{}\nwbegincode{136}\moddef{Model data: model data: TBP}\plusendmoddef
  procedure :: get_name => model_data_get_name
\nwendcode{}\nwbegincode{137}\moddef{Model data: procedures}\plusendmoddef
  function model_data_get_name (model) result (name)
    class(model_data_t), intent(in) :: model
    type(string_t) :: name
    name = model%name
  end function model_data_get_name
  
\eatline
\nwendcode{}\nwbegindocs{138}Retrieve a MD5 sum for the current model parameter values and
decay/polarization settings.  This is
done by writing them to a temporary file, using a standard format.
\nwenddocs{}\nwbegincode{139}\moddef{Model data: model data: TBP}\plusendmoddef
  procedure :: get_parameters_md5sum => model_data_get_parameters_md5sum
\nwendcode{}\nwbegincode{140}\moddef{Model data: procedures}\plusendmoddef
  function model_data_get_parameters_md5sum (model) result (par_md5sum)
    character(32) :: par_md5sum
    class(model_data_t), intent(in) :: model
    real(default), dimension(:), allocatable :: par
    type(field_data_t), pointer :: field
    integer :: unit, i
    allocate (par (model%get_n_real ()))
    call model%real_parameters_to_array (par)
    unit = free_unit ()
    open (unit, status="scratch", action="readwrite")
    write (unit, "(" // FMT_19 // ")")  par
    do i = 1, model%get_n_field ()
       field => model%get_field_ptr_by_index (i)
       if (.not. field%is_stable (.false.) .or. .not. field%is_stable (.true.) &
            .or. field%is_polarized (.false.) .or. field%is_polarized (.true.))&
            then
          write (unit, "(3x,A)") char (field%get_longname ())
          call field%write_decays (unit)
       end if
    end do
    rewind (unit)
    par_md5sum = md5sum (unit)
    close (unit)
  end function model_data_get_parameters_md5sum

\eatline
\nwendcode{}\nwbegindocs{141}Return the MD5 sum.  This is a placeholder, to be overwritten
for the complete model definition.
\nwenddocs{}\nwbegincode{142}\moddef{Model data: model data: TBP}\plusendmoddef
  procedure :: get_md5sum => model_data_get_md5sum
\nwendcode{}\nwbegincode{143}\moddef{Model data: procedures}\plusendmoddef
  function model_data_get_md5sum (model) result (md5sum)
    class(model_data_t), intent(in) :: model
    character(32) :: md5sum
    md5sum = model%get_parameters_md5sum ()
  end function model_data_get_md5sum
  
\eatline
\nwendcode{}\nwbegindocs{144}Initialize a real or complex parameter.
\nwenddocs{}\nwbegincode{145}\moddef{Model data: model data: TBP}\plusendmoddef
  generic :: init_par => model_data_init_par_real, model_data_init_par_complex
  procedure, private :: model_data_init_par_real
  procedure, private :: model_data_init_par_complex
\nwendcode{}\nwbegincode{146}\moddef{Model data: procedures}\plusendmoddef
  subroutine model_data_init_par_real (model, i, name, value)
    class(model_data_t), intent(inout) :: model
    integer, intent(in) :: i
    type(string_t), intent(in) :: name
    real(default), intent(in) :: value
    call model%par_real(i)%init (name, value)
  end subroutine model_data_init_par_real
  
  subroutine model_data_init_par_complex (model, i, name, value)
    class(model_data_t), intent(inout) :: model
    integer, intent(in) :: i
    type(string_t), intent(in) :: name
    complex(default), intent(in) :: value
    call model%par_complex(i)%init (name, value)
  end subroutine model_data_init_par_complex
  
\eatline
\nwendcode{}\nwbegindocs{147}After initialization, return size of parameter array.
\nwenddocs{}\nwbegincode{148}\moddef{Model data: model data: TBP}\plusendmoddef
  procedure :: get_n_real => model_data_get_n_real
  procedure :: get_n_complex => model_data_get_n_complex
\nwendcode{}\nwbegincode{149}\moddef{Model data: procedures}\plusendmoddef
  function model_data_get_n_real (model) result (n)
    class(model_data_t), intent(in) :: model
    integer :: n
    n = size (model%par_real)
  end function model_data_get_n_real
  
  function model_data_get_n_complex (model) result (n)
    class(model_data_t), intent(in) :: model
    integer :: n
    n = size (model%par_complex)
  end function model_data_get_n_complex
  
\eatline
\eatline
\nwendcode{}\nwbegindocs{150}After initialization, extract the whole parameter array.
\nwenddocs{}\nwbegincode{151}\moddef{Model data: model data: TBP}\plusendmoddef
  procedure :: real_parameters_to_array &
       => model_data_real_par_to_array
  procedure :: complex_parameters_to_array &
       => model_data_complex_par_to_array
\nwendcode{}\nwbegincode{152}\moddef{Model data: procedures}\plusendmoddef
  subroutine model_data_real_par_to_array (model, array)
    class(model_data_t), intent(in) :: model
    real(default), dimension(:), intent(inout) :: array
    array = model%par_real%get_real ()
  end subroutine model_data_real_par_to_array
  
  subroutine model_data_complex_par_to_array (model, array)
    class(model_data_t), intent(in) :: model
    complex(default), dimension(:), intent(inout) :: array
    array = model%par_complex%get_complex ()
  end subroutine model_data_complex_par_to_array
  
\eatline
\eatline
\nwendcode{}\nwbegindocs{153}After initialization, set the whole parameter array.
\nwenddocs{}\nwbegincode{154}\moddef{Model data: model data: TBP}\plusendmoddef
  procedure :: real_parameters_from_array &
       => model_data_real_par_from_array
  procedure :: complex_parameters_from_array &
       => model_data_complex_par_from_array
\nwendcode{}\nwbegincode{155}\moddef{Model data: procedures}\plusendmoddef
  subroutine model_data_real_par_from_array (model, array)
    class(model_data_t), intent(inout) :: model
    real(default), dimension(:), intent(in) :: array
    model%par_real = array
  end subroutine model_data_real_par_from_array
  
  subroutine model_data_complex_par_from_array (model, array)
    class(model_data_t), intent(inout) :: model
    complex(default), dimension(:), intent(in) :: array
    model%par_complex = array
  end subroutine model_data_complex_par_from_array
  
\eatline
\eatline
\nwendcode{}\nwbegindocs{156}Analogous, for a C parameter array.
\nwenddocs{}\nwbegincode{157}\moddef{Model data: model data: TBP}\plusendmoddef
  procedure :: real_parameters_to_c_array &
       => model_data_real_par_to_c_array
\nwendcode{}\nwbegincode{158}\moddef{Model data: procedures}\plusendmoddef
  subroutine model_data_real_par_to_c_array (model, array)
    class(model_data_t), intent(in) :: model
    real(c_default_float), dimension(:), intent(inout) :: array
    array = model%par_real%get_real ()
  end subroutine model_data_real_par_to_c_array
  
\eatline
\nwendcode{}\nwbegindocs{159}After initialization, set the whole parameter array.
\nwenddocs{}\nwbegincode{160}\moddef{Model data: model data: TBP}\plusendmoddef
  procedure :: real_parameters_from_c_array &
       => model_data_real_par_from_c_array
\nwendcode{}\nwbegincode{161}\moddef{Model data: procedures}\plusendmoddef
  subroutine model_data_real_par_from_c_array (model, array)
    class(model_data_t), intent(inout) :: model
    real(c_default_float), dimension(:), intent(in) :: array
    model%par_real = real (array, default)
  end subroutine model_data_real_par_from_c_array
  
\eatline
\nwendcode{}\nwbegindocs{162}After initialization, get pointer to a real or complex parameter,
directly by index.
\nwenddocs{}\nwbegincode{163}\moddef{Model data: model data: TBP}\plusendmoddef
  procedure :: get_par_real_ptr => model_data_get_par_real_ptr_index
  procedure :: get_par_complex_ptr => model_data_get_par_complex_ptr_index
\nwendcode{}\nwbegincode{164}\moddef{Model data: procedures}\plusendmoddef
  function model_data_get_par_real_ptr_index (model, i) result (ptr)
    class(model_data_t), intent(inout) :: model
    integer, intent(in) :: i
    class(modelpar_data_t), pointer :: ptr
    ptr => model%par_real(i)
  end function model_data_get_par_real_ptr_index
  
  function model_data_get_par_complex_ptr_index (model, i) result (ptr)
    class(model_data_t), intent(inout) :: model
    integer, intent(in) :: i
    class(modelpar_data_t), pointer :: ptr
    ptr => model%par_complex(i)
  end function model_data_get_par_complex_ptr_index
  
\eatline
\nwendcode{}\nwbegindocs{165}After initialization, get pointer to a parameter by name.
\nwenddocs{}\nwbegincode{166}\moddef{Model data: model data: TBP}\plusendmoddef
  procedure :: get_par_data_ptr => model_data_get_par_data_ptr_name
\nwendcode{}\nwbegincode{167}\moddef{Model data: procedures}\plusendmoddef
  function model_data_get_par_data_ptr_name (model, name) result (ptr)
    class(model_data_t), intent(in) :: model
    type(string_t), intent(in) :: name
    class(modelpar_data_t), pointer :: ptr
    integer :: i
    do i = 1, size (model%par_real)
       if (model%par_real(i)%name == name) then
          ptr => model%par_real(i)
          return
       end if
    end do
    do i = 1, size (model%par_complex)
       if (model%par_complex(i)%name == name) then
          ptr => model%par_complex(i)
          return
       end if
    end do
    ptr => null ()
  end function model_data_get_par_data_ptr_name
  
\eatline
\nwendcode{}\nwbegindocs{168}Return the value by name.  Again, type conversion is allowed.
\nwenddocs{}\nwbegincode{169}\moddef{Model data: model data: TBP}\plusendmoddef
  procedure :: get_real => model_data_get_par_real_value
  procedure :: get_complex => model_data_get_par_complex_value
\nwendcode{}\nwbegincode{170}\moddef{Model data: procedures}\plusendmoddef
  function model_data_get_par_real_value (model, name) result (value)
    class(model_data_t), intent(in) :: model
    type(string_t), intent(in) :: name
    class(modelpar_data_t), pointer :: par
    real(default) :: value
    par => model%get_par_data_ptr (name)
    value = par%get_real ()
  end function model_data_get_par_real_value
  
  function model_data_get_par_complex_value (model, name) result (value)
    class(model_data_t), intent(in) :: model
    type(string_t), intent(in) :: name
    class(modelpar_data_t), pointer :: par
    complex(default) :: value
    par => model%get_par_data_ptr (name)
    value = par%get_complex ()
  end function model_data_get_par_complex_value
  
\eatline
\eatline
\nwendcode{}\nwbegindocs{171}Modify a real or complex parameter.
\nwenddocs{}\nwbegincode{172}\moddef{Model data: model data: TBP}\plusendmoddef
  generic :: set_par => model_data_set_par_real, model_data_set_par_complex
  procedure, private :: model_data_set_par_real
  procedure, private :: model_data_set_par_complex
\nwendcode{}\nwbegincode{173}\moddef{Model data: procedures}\plusendmoddef
  subroutine model_data_set_par_real (model, name, value)
    class(model_data_t), intent(inout) :: model
    type(string_t), intent(in) :: name
    real(default), intent(in) :: value
    class(modelpar_data_t), pointer :: par
    par => model%get_par_data_ptr (name)
    par = value
  end subroutine model_data_set_par_real
  
  subroutine model_data_set_par_complex (model, name, value)
    class(model_data_t), intent(inout) :: model
    type(string_t), intent(in) :: name
    complex(default), intent(in) :: value
    class(modelpar_data_t), pointer :: par
    par => model%get_par_data_ptr (name)
    par = value
  end subroutine model_data_set_par_complex
  
\eatline
\nwendcode{}\nwbegindocs{174}List all fields in the model.
\nwenddocs{}\nwbegincode{175}\moddef{Model data: model data: TBP}\plusendmoddef
  procedure :: write_fields => model_data_write_fields
\nwendcode{}\nwbegincode{176}\moddef{Model data: procedures}\plusendmoddef
  subroutine model_data_write_fields (model, unit)
    class(model_data_t), intent(in) :: model
    integer, intent(in), optional :: unit
    integer :: i
    do i = 1, size (model%field)
       call model%field(i)%write (unit)
    end do
  end subroutine model_data_write_fields
    
\eatline
\nwendcode{}\nwbegindocs{177}After initialization, return number of fields (particles):
\nwenddocs{}\nwbegincode{178}\moddef{Model data: model data: TBP}\plusendmoddef
  procedure :: get_n_field => model_data_get_n_field
\nwendcode{}\nwbegincode{179}\moddef{Model data: procedures}\plusendmoddef
  function model_data_get_n_field (model) result (n)
    class(model_data_t), intent(in) :: model
    integer :: n
    n = size (model%field)
  end function model_data_get_n_field
  
\eatline
\nwendcode{}\nwbegindocs{180}Return the PDG code of a field.  The field is identified by name or
by index.  If the field is not found, return zero.
\nwenddocs{}\nwbegincode{181}\moddef{Model data: model data: TBP}\plusendmoddef
  generic :: get_pdg => &
       model_data_get_field_pdg_index, &
       model_data_get_field_pdg_name
  procedure, private :: model_data_get_field_pdg_index
  procedure, private :: model_data_get_field_pdg_name
\nwendcode{}\nwbegincode{182}\moddef{Model data: procedures}\plusendmoddef
  function model_data_get_field_pdg_index (model, i) result (pdg)
    class(model_data_t), intent(in) :: model
    integer, intent(in) :: i
    integer :: pdg
    pdg = model%field(i)%get_pdg ()
  end function model_data_get_field_pdg_index
  
  function model_data_get_field_pdg_name (model, name, check) result (pdg)
    class(model_data_t), intent(in) :: model
    type(string_t), intent(in) :: name
    logical, intent(in), optional :: check
    integer :: pdg
    integer :: i
    do i = 1, size (model%field)
       associate (field => model%field(i))
         if (field%matches_name (name, .false.)) then
            pdg = field%get_pdg ()
            return
         else if (field%matches_name (name, .true.)) then
            pdg = - field%get_pdg ()
            return
         end if
       end associate
    end do
    pdg = 0
    call model%field_error (check, name)
  end function model_data_get_field_pdg_name
    
\eatline
\nwendcode{}\nwbegindocs{183}Return an array of all PDG codes, including antiparticles.  The antiparticle
are sorted after all particles.
\nwenddocs{}\nwbegincode{184}\moddef{Model data: model data: TBP}\plusendmoddef
  procedure :: get_all_pdg => model_data_get_all_pdg
\nwendcode{}\nwbegincode{185}\moddef{Model data: procedures}\plusendmoddef
  subroutine model_data_get_all_pdg (model, pdg)
    class(model_data_t), intent(in) :: model
    integer, dimension(:), allocatable, intent(inout) :: pdg
    integer :: n0, n1, i, k
    n0 = size (model%field)
    n1 = n0 + count (model%field%has_antiparticle ())
    allocate (pdg (n1))
    pdg(1:n0) = model%field%get_pdg ()
    k = n0
    do i = 1, size (model%field)
       associate (field => model%field(i))
         if (field%has_antiparticle ()) then
            k = k + 1
            pdg(k) = - field%get_pdg ()
         end if
       end associate
    end do
  end subroutine model_data_get_all_pdg
    
\eatline
\nwendcode{}\nwbegindocs{186}Return pointer to the field array.
\nwenddocs{}\nwbegincode{187}\moddef{Model data: model data: TBP}\plusendmoddef
  procedure :: get_field_array_ptr => model_data_get_field_array_ptr
\nwendcode{}\nwbegincode{188}\moddef{Model data: procedures}\plusendmoddef
  function model_data_get_field_array_ptr (model) result (ptr)
    class(model_data_t), intent(in), target :: model
    type(field_data_t), dimension(:), pointer :: ptr
    ptr => model%field
  end function model_data_get_field_array_ptr
  
\eatline
\nwendcode{}\nwbegindocs{189}Return pointer to a field.  The identifier should be the unique long
name, the PDG code, or the index.

We can issue an error message, if the {\tt{}check} flag is set.  We never return
an error if the PDG code is zero, this yields just a null pointer.
\nwenddocs{}\nwbegincode{190}\moddef{Model data: model data: TBP}\plusendmoddef
  generic :: get_field_ptr => &
       model_data_get_field_ptr_name, &
       model_data_get_field_ptr_pdg
  procedure, private :: model_data_get_field_ptr_name
  procedure, private :: model_data_get_field_ptr_pdg
  procedure :: get_field_ptr_by_index => model_data_get_field_ptr_index
\nwendcode{}\nwbegincode{191}\moddef{Model data: procedures}\plusendmoddef
  function model_data_get_field_ptr_name (model, name, check) result (ptr)
    class(model_data_t), intent(in), target :: model
    type(string_t), intent(in) :: name
    logical, intent(in), optional :: check
    type(field_data_t), pointer :: ptr
    integer :: i
    do i = 1, size (model%field)
       if (model%field(i)%matches_name (name, .false.)) then
          ptr => model%field(i)
          return
       else if (model%field(i)%matches_name (name, .true.)) then
          ptr => model%field(i)
          return
       end if
    end do
    ptr => null ()
    call model%field_error (check, name)
  end function model_data_get_field_ptr_name

  function model_data_get_field_ptr_pdg (model, pdg, check) result (ptr)
    class(model_data_t), intent(in), target :: model
    integer, intent(in) :: pdg
    logical, intent(in), optional :: check
    type(field_data_t), pointer :: ptr
    integer :: i, pdg_abs
    if (pdg == 0) then
       ptr => null ()
       return
    end if
    pdg_abs = abs (pdg)
    do i = 1, size (model%field)
       if (model%field(i)%get_pdg () == pdg_abs) then
          ptr => model%field(i)
          return
       end if
    end do
    ptr => null ()
    call model%field_error (check, pdg=pdg)
  end function model_data_get_field_ptr_pdg

  function model_data_get_field_ptr_index (model, i) result (ptr)
    class(model_data_t), intent(in), target :: model
    integer, intent(in) :: i
    type(field_data_t), pointer :: ptr
    ptr => model%field(i)
  end function model_data_get_field_ptr_index

\eatline
\nwendcode{}\nwbegindocs{192}Don't assign a pointer, just check.
\nwenddocs{}\nwbegincode{193}\moddef{Model data: model data: TBP}\plusendmoddef
  procedure :: test_field => model_data_test_field_pdg
\nwendcode{}\nwbegincode{194}\moddef{Model data: procedures}\plusendmoddef
  function model_data_test_field_pdg (model, pdg, check) result (exist)
    class(model_data_t), intent(in), target :: model
    integer, intent(in) :: pdg
    logical, intent(in), optional :: check
    logical :: exist
    exist = associated (model%get_field_ptr (pdg, check))
  end function model_data_test_field_pdg
  
\eatline
\nwendcode{}\nwbegindocs{195}Error message, if {\tt{}check} is set.
\nwenddocs{}\nwbegincode{196}\moddef{Model data: model data: TBP}\plusendmoddef
  procedure :: field_error => model_data_field_error
\nwendcode{}\nwbegincode{197}\moddef{Model data: procedures}\plusendmoddef
  subroutine model_data_field_error (model, check, name, pdg)
    class(model_data_t), intent(in) :: model
    logical, intent(in), optional :: check
    type(string_t), intent(in), optional :: name
    integer, intent(in), optional :: pdg
    if (present (check)) then
       if (check) then
          if (present (name)) then
             write (msg_buffer, "(A,1x,A,1x,A,1x,A)") &
                  "No particle with name", char (name), &
                  "is contained in model", char (model%name)
          else if (present (pdg)) then
             write (msg_buffer, "(A,1x,I0,1x,A,1x,A)") &
                  "No particle with PDG code", pdg, &
                  "is contained in model", char (model%name)
          else
             write (msg_buffer, "(A,1x,A,1x,A)") &
                  "Particle missing", &
                  "in model", char (model%name)
          end if
          call msg_fatal ()
       end if
    end if
  end subroutine model_data_field_error
          
\eatline
\nwendcode{}\nwbegindocs{198}Assign mass and width value, which are associated via pointer.
Identify the particle via pdg.
\nwenddocs{}\nwbegincode{199}\moddef{Model data: model data: TBP}\plusendmoddef
  procedure :: set_field_mass => model_data_set_field_mass_pdg
  procedure :: set_field_width => model_data_set_field_width_pdg
\nwendcode{}\nwbegincode{200}\moddef{Model data: procedures}\plusendmoddef
  subroutine model_data_set_field_mass_pdg (model, pdg, value)
    class(model_data_t), intent(inout) :: model
    integer, intent(in) :: pdg
    real(default), intent(in) :: value
    type(field_data_t), pointer :: field
    field => model%get_field_ptr (pdg, check = .true.)
    call field%set_mass (value)
  end subroutine model_data_set_field_mass_pdg
    
  subroutine model_data_set_field_width_pdg (model, pdg, value)
    class(model_data_t), intent(inout) :: model
    integer, intent(in) :: pdg
    real(default), intent(in) :: value
    type(field_data_t), pointer :: field
    field => model%get_field_ptr (pdg, check = .true.)
    call field%set_width (value)
  end subroutine model_data_set_field_width_pdg
    
\eatline
\eatline
\nwendcode{}\nwbegindocs{201}Mark a particle as unstable and provide a list of names for its
decay processes.  In contrast with the previous subroutine which is
for internal use, we address the particle by its PDG code.  If the
index is negative, we address the antiparticle.
\nwenddocs{}\nwbegincode{202}\moddef{Model data: model data: TBP}\plusendmoddef
  procedure :: set_unstable => model_data_set_unstable
  procedure :: set_stable => model_data_set_stable
\nwendcode{}\nwbegincode{203}\moddef{Model data: procedures}\plusendmoddef
  subroutine model_data_set_unstable &
       (model, pdg, decay, isotropic, diagonal, decay_helicity)
    class(model_data_t), intent(inout), target :: model
    integer, intent(in) :: pdg
    type(string_t), dimension(:), intent(in) :: decay
    logical, intent(in), optional :: isotropic, diagonal
    integer, intent(in), optional :: decay_helicity
    type(field_data_t), pointer :: field
    field => model%get_field_ptr (pdg)
    if (pdg > 0) then
       call field%set ( &
            p_is_stable = .false., p_decay = decay, &
            p_decays_isotropically = isotropic, &
            p_decays_diagonal = diagonal, &
            p_decay_helicity = decay_helicity)
    else
       call field%set ( &
            a_is_stable = .false., a_decay = decay, &
            a_decays_isotropically = isotropic, &
            a_decays_diagonal = diagonal, &
            a_decay_helicity = decay_helicity)
    end if
  end subroutine model_data_set_unstable
       
  subroutine model_data_set_stable (model, pdg)
    class(model_data_t), intent(inout), target :: model
    integer, intent(in) :: pdg
    type(field_data_t), pointer :: field
    field => model%get_field_ptr (pdg)
    if (pdg > 0) then
       call field%set (p_is_stable = .true.)
    else
       call field%set (a_is_stable = .true.)
    end if
  end subroutine model_data_set_stable
       
\eatline
\eatline
\nwendcode{}\nwbegindocs{204}Mark a particle as polarized.
\nwenddocs{}\nwbegincode{205}\moddef{Model data: model data: TBP}\plusendmoddef
  procedure :: set_polarized => model_data_set_polarized
  procedure :: set_unpolarized => model_data_set_unpolarized
\nwendcode{}\nwbegincode{206}\moddef{Model data: procedures}\plusendmoddef
  subroutine model_data_set_polarized (model, pdg)
    class(model_data_t), intent(inout), target :: model
    integer, intent(in) :: pdg
    type(field_data_t), pointer :: field
    field => model%get_field_ptr (pdg)
    if (pdg > 0) then
       call field%set (p_polarized = .true.)
    else
       call field%set (a_polarized = .true.)
    end if
  end subroutine model_data_set_polarized
    
  subroutine model_data_set_unpolarized (model, pdg)
    class(model_data_t), intent(inout), target :: model
    integer, intent(in) :: pdg
    type(field_data_t), pointer :: field
    field => model%get_field_ptr (pdg)
    if (pdg > 0) then
       call field%set (p_polarized = .false.)
    else
       call field%set (a_polarized = .false.)
    end if
  end subroutine model_data_set_unpolarized
    
\eatline
\eatline
\nwendcode{}\nwbegindocs{207}Revert all polarized (unstable) particles to unpolarized (stable)
status, respectively.
\nwenddocs{}\nwbegincode{208}\moddef{Model data: model data: TBP}\plusendmoddef
  procedure :: clear_unstable => model_clear_unstable
  procedure :: clear_polarized => model_clear_polarized
\nwendcode{}\nwbegincode{209}\moddef{Model data: procedures}\plusendmoddef
  subroutine model_clear_unstable (model)
    class(model_data_t), intent(inout), target :: model
    integer :: i
    type(field_data_t), pointer :: field
    do i = 1, model%get_n_field ()
       field => model%get_field_ptr_by_index (i)
       call field%set (p_is_stable = .true.)
       if (field%has_antiparticle ()) then
          call field%set (a_is_stable = .true.)
       end if
    end do
  end subroutine model_clear_unstable
  
  subroutine model_clear_polarized (model)
    class(model_data_t), intent(inout), target :: model
    integer :: i
    type(field_data_t), pointer :: field
    do i = 1, model%get_n_field ()
       field => model%get_field_ptr_by_index (i)
       call field%set (p_polarized = .false.)
       if (field%has_antiparticle ()) then
          call field%set (a_polarized = .false.)
       end if
    end do
  end subroutine model_clear_polarized
  
\eatline
\eatline
\nwendcode{}\nwbegindocs{210}List all vertices, optionally also the hash table.
\nwenddocs{}\nwbegincode{211}\moddef{Model data: model data: TBP}\plusendmoddef
  procedure :: write_vertices => model_data_write_vertices
\nwendcode{}\nwbegincode{212}\moddef{Model data: procedures}\plusendmoddef
  subroutine model_data_write_vertices (model, unit, verbose)
    class(model_data_t), intent(in) :: model
    integer, intent(in), optional :: unit
    logical, intent(in), optional :: verbose
    integer :: i, u
    u = given_output_unit (unit)
    do i = 1, size (model%vtx)
       call vertex_write (model%vtx(i), unit)
    end do
    if (present (verbose)) then
       if (verbose) then
          write (u, *)
          call vertex_table_write (model%vt, unit)
       end if
    end if
  end subroutine model_data_write_vertices

\eatline
\nwendcode{}\nwbegindocs{213}Vertex definition.
\nwenddocs{}\nwbegincode{214}\moddef{Model data: model data: TBP}\plusendmoddef
  generic :: set_vertex => &
       model_data_set_vertex_pdg, model_data_set_vertex_names
  procedure, private :: model_data_set_vertex_pdg
  procedure, private :: model_data_set_vertex_names
\nwendcode{}\nwbegincode{215}\moddef{Model data: procedures}\plusendmoddef
  subroutine model_data_set_vertex_pdg (model, i, pdg)
    class(model_data_t), intent(inout), target :: model
    integer, intent(in) :: i
    integer, dimension(:), intent(in) :: pdg
    call vertex_init (model%vtx(i), pdg, model)
  end subroutine model_data_set_vertex_pdg

  subroutine model_data_set_vertex_names (model, i, name)
    class(model_data_t), intent(inout), target :: model
    integer, intent(in) :: i
    type(string_t), dimension(:), intent(in) :: name
    integer, dimension(size(name)) :: pdg
    integer :: j
    do j = 1, size (name)
       pdg(j) = model%get_pdg (name(j))
    end do
    call model%set_vertex (i, pdg)
  end subroutine model_data_set_vertex_names

\eatline
\nwendcode{}\nwbegindocs{216}Finalize vertex definition: set up the hash table.
\nwenddocs{}\nwbegincode{217}\moddef{Model data: model data: TBP}\plusendmoddef
  procedure :: freeze_vertices => model_data_freeze_vertices
\nwendcode{}\nwbegincode{218}\moddef{Model data: procedures}\plusendmoddef
  subroutine model_data_freeze_vertices (model)
    class(model_data_t), intent(inout) :: model
    call model%vt%init (model%field, model%vtx)
  end subroutine model_data_freeze_vertices
  
\eatline
\nwendcode{}\nwbegindocs{219}Number of vertices in model
\nwenddocs{}\nwbegincode{220}\moddef{Model data: model data: TBP}\plusendmoddef
  procedure :: get_n_vtx => model_data_get_n_vtx
\nwendcode{}\nwbegincode{221}\moddef{Model data: procedures}\plusendmoddef
  function model_data_get_n_vtx (model) result (n)
    class(model_data_t), intent(in) :: model
    integer :: n
    n = size (model%vtx)
  end function model_data_get_n_vtx
  
\eatline
\nwendcode{}\nwbegindocs{222}Lookup functions
\nwenddocs{}\nwbegincode{223}\moddef{Model data: model data: TBP}\plusendmoddef
  procedure :: match_vertex => model_data_match_vertex
\nwendcode{}\nwbegincode{224}\moddef{Model data: procedures}\plusendmoddef
  subroutine model_data_match_vertex (model, pdg1, pdg2, pdg3)
    class(model_data_t), intent(in) :: model
    integer, intent(in) :: pdg1, pdg2
    integer, dimension(:), allocatable, intent(out) :: pdg3
    call model%vt%match (pdg1, pdg2, pdg3)
  end subroutine model_data_match_vertex

\eatline
\nwendcode{}\nwbegincode{225}\moddef{Model data: model data: TBP}\plusendmoddef
  procedure :: check_vertex => model_data_check_vertex
\nwendcode{}\nwbegincode{226}\moddef{Model data: procedures}\plusendmoddef
  function model_data_check_vertex (model, pdg1, pdg2, pdg3) result (flag)
    logical :: flag
    class(model_data_t), intent(in) :: model
    integer, intent(in) :: pdg1, pdg2, pdg3
    flag = model%vt%check (pdg1, pdg2, pdg3)
  end function model_data_check_vertex

\eatline
\nwendcode{}\nwbegindocs{227}\nwdocspar
\subsection{Toy Models}
This is a stripped-down version of the (already trivial) model 'Test'.
\nwenddocs{}\nwbegincode{228}\moddef{Model data: model data: TBP}\plusendmoddef
  procedure :: init_test => model_data_init_test
\nwendcode{}\nwbegincode{229}\moddef{Model data: procedures}\plusendmoddef
  subroutine model_data_init_test (model)
    class(model_data_t), intent(out) :: model
    type(field_data_t), pointer :: field
    integer, parameter :: n_real = 4
    integer, parameter :: n_field = 2
    integer, parameter :: n_vertex = 2
    integer :: i
    call model%init (var_str ("Test"), &
         n_real, 0, n_field, n_vertex)
    i = 0
    i = i + 1
    call model%init_par (i, var_str ("gy"), 1._default)
    i = i + 1
    call model%init_par (i, var_str ("ms"), 125._default)
    i = i + 1
    call model%init_par (i, var_str ("ff"), 1.5_default)
    i = i + 1
    call model%init_par (i, var_str ("mf"), 1.5_default * 125._default)
    i = 0
    i = i + 1
    field => model%get_field_ptr_by_index (i)
    call field%init (var_str ("SCALAR"), 25)
    call field%set (spin_type=1) 
    call field%set (mass_data=model%get_par_real_ptr (2))
    call field%set (name = [var_str ("s")])
    i = i + 1
    field => model%get_field_ptr_by_index (i)
    call field%init (var_str ("FERMION"), 6)
    call field%set (spin_type=2) 
    call field%set (mass_data=model%get_par_real_ptr (4))
    call field%set (name = [var_str ("f")], anti = [var_str ("fbar")])
    call model%freeze_fields ()
    i = 0
    i = i + 1
    call model%set_vertex (i, [var_str ("fbar"), var_str ("f"), var_str ("s")])
    i = i + 1
    call model%set_vertex (i, [var_str ("s"), var_str ("s"), var_str ("s")])
    call model%freeze_vertices ()
  end subroutine model_data_init_test
  
\eatline
\nwendcode{}\nwbegindocs{230}\nwdocspar
This procedure prepares a subset of QED for testing purposes.
\nwenddocs{}\nwbegincode{231}\moddef{Model data: model data: TBP}\plusendmoddef
  procedure :: init_qed_test => model_data_init_qed_test
\nwendcode{}\nwbegincode{232}\moddef{Model data: procedures}\plusendmoddef
  subroutine model_data_init_qed_test (model)
    class(model_data_t), intent(out) :: model
    type(field_data_t), pointer :: field
    integer, parameter :: n_real = 1
    integer, parameter :: n_field = 2
    integer :: i
    call model%init (var_str ("QED_test"), &
         n_real, 0, n_field, 0)
    i = 0
    i = i + 1
    call model%init_par (i, var_str ("me"), 0.000510997_default)
    i = 0
    i = i + 1
    field => model%get_field_ptr_by_index (i)
    call field%init (var_str ("E_LEPTON"), 11)
    call field%set (spin_type=2, charge_type=-4)
    call field%set (mass_data=model%get_par_real_ptr (1))
    call field%set (name = [var_str ("e-")], anti = [var_str ("e+")])
    i = i + 1
    field => model%get_field_ptr_by_index (i)
    call field%init (var_str ("PHOTON"), 22)
    call field%set (spin_type=3)
    call field%set (name = [var_str ("A")])
    call model%freeze_fields ()
    call model%freeze_vertices ()
  end subroutine model_data_init_qed_test
  
\eatline
\nwendcode{}\nwbegindocs{233}\nwdocspar
This procedure prepares a subset of the Standard Model for testing purposes.
We can thus avoid dependencies on model I/O, which is not defined here.
\nwenddocs{}\nwbegincode{234}\moddef{Model data: model data: TBP}\plusendmoddef
  procedure :: init_sm_test => model_data_init_sm_test
\nwendcode{}\nwbegincode{235}\moddef{Model data: procedures}\plusendmoddef
  subroutine model_data_init_sm_test (model)
    class(model_data_t), intent(out) :: model
    type(field_data_t), pointer :: field
    integer, parameter :: n_real = 11
    integer, parameter :: n_field = 19
    integer, parameter :: n_vtx = 9
    integer :: i
    call model%init (var_str ("SM_test"), &
         n_real, 0, n_field, n_vtx)
    i = 0
    i = i + 1
    call model%init_par (i, var_str ("mZ"), 91.1882_default)
    i = i + 1
    call model%init_par (i, var_str ("mW"), 80.419_default)
    i = i + 1
    call model%init_par (i, var_str ("me"), 0.000510997_default)
    i = i + 1
    call model%init_par (i, var_str ("mmu"), 0.105658389_default)
    i = i + 1
    call model%init_par (i, var_str ("mb"), 4.2_default)
    i = i + 1
    call model%init_par (i, var_str ("mtop"), 173.1_default)
    i = i + 1
    call model%init_par (i, var_str ("wZ"), 2.443_default)
    i = i + 1
    call model%init_par (i, var_str ("wW"), 2.049_default)
    i = i + 1
    call model%init_par (i, var_str ("ee"), 0.3079561542961_default)
    i = i + 1
    call model%init_par (i, var_str ("cw"), 8.819013863636E-01_default)
    i = i + 1
    call model%init_par (i, var_str ("sw"), 4.714339240339E-01_default)
    i = 0
    i = i + 1
    field => model%get_field_ptr_by_index (i)
    call field%init (var_str ("D_QUARK"), 1)
    call field%set (spin_type=2, color_type=3, charge_type=-2, isospin_type=-2)
    call field%set (name = [var_str ("d")], anti = [var_str ("dbar")])
    i = i + 1
    field => model%get_field_ptr_by_index (i)
    call field%init (var_str ("U_QUARK"), 2)
    call field%set (spin_type=2, color_type=3, charge_type=3, isospin_type=2)
    call field%set (name = [var_str ("u")], anti = [var_str ("ubar")])
    i = i + 1
    field => model%get_field_ptr_by_index (i)
    call field%init (var_str ("S_QUARK"), 3)
    call field%set (spin_type=2, color_type=3, charge_type=-2, isospin_type=-2)
    call field%set (name = [var_str ("s")], anti = [var_str ("sbar")])
    i = i + 1
    field => model%get_field_ptr_by_index (i)
    call field%init (var_str ("C_QUARK"), 4)
    call field%set (spin_type=2, color_type=3, charge_type=3, isospin_type=2)
    call field%set (name = [var_str ("c")], anti = [var_str ("cbar")])
    i = i + 1
    field => model%get_field_ptr_by_index (i)
    call field%init (var_str ("B_QUARK"), 5)
    call field%set (spin_type=2, color_type=3, charge_type=-2, isospin_type=-2)
    call field%set (mass_data=model%get_par_real_ptr (5))
    call field%set (name = [var_str ("b")], anti = [var_str ("bbar")])
    i = i + 1
    field => model%get_field_ptr_by_index (i)
    call field%init (var_str ("T_QUARK"), 6)
    call field%set (spin_type=2, color_type=3, charge_type=3, isospin_type=2)
    call field%set (mass_data=model%get_par_real_ptr (6))
    call field%set (name = [var_str ("t")], anti = [var_str ("tbar")])
    i = i + 1
    field => model%get_field_ptr_by_index (i)
    call field%init (var_str ("E_LEPTON"), 11)
    call field%set (spin_type=2)
    call field%set (mass_data=model%get_par_real_ptr (3))
    call field%set (name = [var_str ("e-")], anti = [var_str ("e+")])
    i = i + 1
    field => model%get_field_ptr_by_index (i)
    call field%init (var_str ("E_NEUTRINO"), 12)
    call field%set (spin_type=2, is_left_handed=.true.)
    call field%set (name = [var_str ("nue")], anti = [var_str ("nuebar")])
    i = i + 1
    field => model%get_field_ptr_by_index (i)
    call field%init (var_str ("MU_LEPTON"), 13)
    call field%set (spin_type=2)
    call field%set (mass_data=model%get_par_real_ptr (4))
    call field%set (name = [var_str ("mu-")], anti = [var_str ("mu+")])
    i = i + 1
    field => model%get_field_ptr_by_index (i)
    call field%init (var_str ("MU_NEUTRINO"), 14)
    call field%set (spin_type=2, is_left_handed=.true.)
    call field%set (name = [var_str ("numu")], anti = [var_str ("numubar")])
    i = i + 1
    field => model%get_field_ptr_by_index (i)
    call field%init (var_str ("GLUON"), 21)
    call field%set (spin_type=3, color_type=8)
    call field%set (name = [var_str ("gl")])
    i = i + 1
    field => model%get_field_ptr_by_index (i)
    call field%init (var_str ("PHOTON"), 22)
    call field%set (spin_type=3)
    call field%set (name = [var_str ("A")])
    i = i + 1
    field => model%get_field_ptr_by_index (i)
    call field%init (var_str ("Z_BOSON"), 23)
    call field%set (spin_type=3)
    call field%set (mass_data=model%get_par_real_ptr (1))
    call field%set (width_data=model%get_par_real_ptr (7))
    call field%set (name = [var_str ("Z")])
    i = i + 1
    field => model%get_field_ptr_by_index (i)
    call field%init (var_str ("W_BOSON"), 24)
    call field%set (spin_type=3)
    call field%set (mass_data=model%get_par_real_ptr (2))
    call field%set (width_data=model%get_par_real_ptr (8))
    call field%set (name = [var_str ("W+")], anti = [var_str ("W-")])
    i = i + 1
    field => model%get_field_ptr_by_index (i)
    call field%init (var_str ("HIGGS"), 25)
    call field%set (spin_type=1)
!    call field%set (mass_data=model%get_par_real_ptr (2))
!    call field%set (width_data=model%get_par_real_ptr (8))
    call field%set (name = [var_str ("H")])
    i = i + 1
    field => model%get_field_ptr_by_index (i)
    call field%init (var_str ("PROTON"), 2212)
    call field%set (spin_type=2)
    call field%set (name = [var_str ("p")], anti = [var_str ("pbar")])
    i = i + 1
    field => model%get_field_ptr_by_index (i)
    call field%init (var_str ("HADRON_REMNANT_SINGLET"), 91)
    call field%set (color_type=1)
    call field%set (name = [var_str ("hr1")])
    i = i + 1
    field => model%get_field_ptr_by_index (i)
    call field%init (var_str ("HADRON_REMNANT_TRIPLET"), 92)
    call field%set (color_type=3)
    call field%set (name = [var_str ("hr3")], anti = [var_str ("hr3bar")])
    i = i + 1
    field => model%get_field_ptr_by_index (i)
    call field%init (var_str ("HADRON_REMNANT_OCTET"), 93)
    call field%set (color_type=8)
    call field%set (name = [var_str ("hr8")])
    call model%freeze_fields ()
    i = 0
    i = i + 1
    call model%set_vertex (i, [var_str ("dbar"), var_str ("d"), var_str ("A")])
    i = i + 1
    call model%set_vertex (i, [var_str ("ubar"), var_str ("u"), var_str ("A")])
    i = i + 1
    call model%set_vertex (i, [var_str ("gl"), var_str ("gl"), var_str ("gl")])
    i = i + 1
    call model%set_vertex (i, [var_str ("dbar"), var_str ("d"), var_str ("gl")])
    i = i + 1
    call model%set_vertex (i, [var_str ("ubar"), var_str ("u"), var_str ("gl")])
    i = i + 1
    call model%set_vertex (i, [var_str ("dbar"), var_str ("d"), var_str ("Z")])
    i = i + 1
    call model%set_vertex (i, [var_str ("ubar"), var_str ("u"), var_str ("Z")])
    i = i + 1
    call model%set_vertex (i, [var_str ("ubar"), var_str ("d"), var_str ("W+")])
    i = i + 1
    call model%set_vertex (i, [var_str ("dbar"), var_str ("u"), var_str ("W-")])
    call model%freeze_vertices ()
  end subroutine model_data_init_sm_test
  
\eatline
\nwendcode{}\nwbegindocs{236}\nwdocspar
\clearpage
%%%%%%%%%%%%%%%%%%%%%%%%%%%%%%%%%%%%%%%%%%%%%%%%%%%%%%%%%%%%%%%%%%%%%%%%
\section{Model Testbed}
The standard way of defining a model uses concrete variables and expressions to
interpret the model file.  Some of this is not available at the point of use.  This
is no problem for the \whizard\ program as a whole, but unit tests are
kept local to their respective module and don't access all definitions.

Instead, we introduce a separate module that provides hooks, one for
initializing a model and one for finalizing a model.  The main program can
assign real routines to the hooks (procedure pointers of abstract type) before
unit tests are called.  The unit tests can call the abstract routines without
knowing about their implementation.
\nwenddocs{}\nwbegincode{237}\moddef{\code{}model{\char95}testbed.f90\edoc{}}\endmoddef
  module model_testbed

\LA{}File header\RA{}

\LA{}Use strings\RA{}
  use model_data
  use var_base
 
\LA{}Standard module head\RA{}

\LA{}Model testbed: public\RA{}

\LA{}Model testbed: variables\RA{}

\LA{}Model testbed: interfaces\RA{}

end module model_testbed
\eatline
\nwendcode{}\nwbegindocs{238}\nwdocspar
\subsection{Abstract Model Handlers}
Both routines take a polymorphic model (data) pointer, which
is allocated/deallocated inside the subroutine.  The model constructor
{\tt{}prepare{\char95}model} requires the model name as input.  It can, optionally,
return a link to the variable list of the model.
\nwenddocs{}\nwbegincode{239}\moddef{Model testbed: public}\endmoddef
  public :: prepare_model
  public :: cleanup_model
\nwendcode{}\nwbegincode{240}\moddef{Model testbed: variables}\endmoddef
  procedure (prepare_model_proc), pointer :: prepare_model => null ()
  procedure (cleanup_model_proc), pointer :: cleanup_model => null ()
\nwendcode{}\nwbegincode{241}\moddef{Model testbed: interfaces}\endmoddef
  abstract interface
     subroutine prepare_model_proc (model, name, vars)
       import
       class(model_data_t), pointer, intent(inout) :: model
       type(string_t), intent(in) :: name
       class(vars_t), pointer, intent(out), optional :: vars
     end subroutine prepare_model_proc
  end interface
       
  abstract interface
     subroutine cleanup_model_proc (model)
       import
       class(model_data_t), pointer, intent(inout) :: model
     end subroutine cleanup_model_proc
  end interface
       
\eatline
\eatline
\nwendcode{}\nwbegindocs{242}\nwdocspar
\clearpage
%%%%%%%%%%%%%%%%%%%%%%%%%%%%%%%%%%%%%%%%%%%%%%%%%%%%%%%%%%%%%%%%%%%%%%%%
\section{Helicities}
This module defines types and tools for dealing with helicity
information.
\nwenddocs{}\nwbegincode{243}\moddef{\code{}helicities.f90\edoc{}}\endmoddef
\LA{}File header\RA{}

module helicities

  use io_units

\LA{}Standard module head\RA{}

\LA{}Helicities: public\RA{}

\LA{}Helicities: types\RA{}

\LA{}Helicities: interfaces\RA{}

contains

\LA{}Helicities: procedures\RA{}

end module helicities
\eatline
\nwendcode{}\nwbegindocs{244}\nwdocspar
\subsection{Helicity types}
Helicities may be defined or undefined, corresponding to a polarized
or unpolarized state.  Each helicity is actually a pair of helicities,
corresponding to an entry in the spin density matrix.  Obviously,
diagonal entries are distinguished.
\nwenddocs{}\nwbegincode{245}\moddef{Helicities: public}\endmoddef
  public :: helicity_t
\nwendcode{}\nwbegincode{246}\moddef{Helicities: types}\endmoddef
  type :: helicity_t
     private
     logical :: defined = .false.
     integer :: h1, h2
   contains
   \LA{}Helicities: helicity: TBP\RA{}
  end type helicity_t

\eatline
\nwendcode{}\nwbegindocs{247}Constructor functions, for convenience:
\nwenddocs{}\nwbegincode{248}\moddef{Helicities: public}\plusendmoddef
  public :: helicity
\nwendcode{}\nwbegincode{249}\moddef{Helicities: interfaces}\endmoddef
  interface helicity
     module procedure helicity0, helicity1, helicity2
  end interface helicity
  
\nwendcode{}\nwbegincode{250}\moddef{Helicities: procedures}\endmoddef
  pure function helicity0 () result (hel)
    type(helicity_t) :: hel
  end function helicity0
  
  elemental function helicity1 (h) result (hel)
    type(helicity_t) :: hel
    integer, intent(in) :: h
    call hel%init (h)
  end function helicity1
  
  elemental function helicity2 (h2, h1) result (hel)
    type(helicity_t) :: hel
    integer, intent(in) :: h1, h2
    call hel%init (h2, h1)
  end function helicity2
  
\eatline
\nwendcode{}\nwbegindocs{251}Initializers.

Note: conceptually, the argument to initializers should be INTENT(OUT).
However, Interp.\ F08/0033 prohibited this.  The reason is that, in principle,
the call could result in the execution of an impure finalizer for a type
extension of {\tt{}hel} (ugh).
\nwenddocs{}\nwbegincode{252}\moddef{Helicities: helicity: TBP}\endmoddef
  generic :: init => helicity_init0, helicity_init1, helicity_init2
  procedure, private :: helicity_init0
  procedure, private :: helicity_init1
  procedure, private :: helicity_init2
\nwendcode{}\nwbegincode{253}\moddef{Helicities: procedures}\plusendmoddef
  elemental subroutine helicity_init0 (hel)
    class(helicity_t), intent(inout) :: hel
    hel%defined = .false.
  end subroutine helicity_init0

  elemental subroutine helicity_init1 (hel, h)
    class(helicity_t), intent(inout) :: hel
    integer, intent(in) :: h
    hel%defined = .true.
    hel%h1 = h
    hel%h2 = h
  end subroutine helicity_init1

  elemental subroutine helicity_init2 (hel, h2, h1)
    class(helicity_t), intent(inout) :: hel
    integer, intent(in) :: h1, h2
    hel%defined = .true.
    hel%h2 = h2
    hel%h1 = h1
  end subroutine helicity_init2

\eatline
\nwendcode{}\nwbegindocs{254}Undefine:
\nwenddocs{}\nwbegincode{255}\moddef{Helicities: helicity: TBP}\plusendmoddef
  procedure :: undefine => helicity_undefine
\nwendcode{}\nwbegincode{256}\moddef{Helicities: procedures}\plusendmoddef
  elemental subroutine helicity_undefine (hel)
    class(helicity_t), intent(inout) :: hel
    hel%defined = .false.
  end subroutine helicity_undefine

\eatline
\nwendcode{}\nwbegindocs{257}Diagonalize by removing the second entry (use with care!)
\nwenddocs{}\nwbegincode{258}\moddef{Helicities: helicity: TBP}\plusendmoddef
  procedure :: diagonalize => helicity_diagonalize
\nwendcode{}\nwbegincode{259}\moddef{Helicities: procedures}\plusendmoddef
  elemental subroutine helicity_diagonalize (hel)
    class(helicity_t), intent(inout) :: hel
    hel%h2 = hel%h1
  end subroutine helicity_diagonalize

\eatline
\nwendcode{}\nwbegindocs{260}Output (no linebreak).  No output if undefined.
\nwenddocs{}\nwbegincode{261}\moddef{Helicities: helicity: TBP}\plusendmoddef
  procedure :: write => helicity_write
\nwendcode{}\nwbegincode{262}\moddef{Helicities: procedures}\plusendmoddef
  subroutine helicity_write (hel, unit)
    class(helicity_t), intent(in) :: hel
    integer, intent(in), optional :: unit
    integer :: u
    u = given_output_unit (unit);  if (u < 0)  return
    if (hel%defined) then
       write (u, "(A)", advance="no")  "h("
       write (u, "(I0)", advance="no")  hel%h1
       if (hel%h1 /= hel%h2) then
          write (u, "(A)", advance="no") "|"
          write (u, "(I0)", advance="no")  hel%h2
       end if
       write (u, "(A)", advance="no")  ")"
    end if
  end subroutine helicity_write

\eatline
\nwendcode{}\nwbegindocs{263}Binary I/O.  Write contents only if defined.
\nwenddocs{}\nwbegincode{264}\moddef{Helicities: helicity: TBP}\plusendmoddef
  procedure :: write_raw => helicity_write_raw
  procedure :: read_raw => helicity_read_raw
\nwendcode{}\nwbegincode{265}\moddef{Helicities: procedures}\plusendmoddef
  subroutine helicity_write_raw (hel, u)
    class(helicity_t), intent(in) :: hel
    integer, intent(in) :: u
    write (u) hel%defined
    if (hel%defined) then
       write (u) hel%h1, hel%h2
    end if
  end subroutine helicity_write_raw

  subroutine helicity_read_raw (hel, u, iostat)
    class(helicity_t), intent(out) :: hel
    integer, intent(in) :: u
    integer, intent(out), optional :: iostat
    read (u, iostat=iostat) hel%defined
    if (hel%defined) then
       read (u, iostat=iostat) hel%h1, hel%h2
    end if
  end subroutine helicity_read_raw

\eatline
\nwendcode{}\nwbegindocs{266}\nwdocspar
\subsection{Predicates}
Check if the helicity is defined:
\nwenddocs{}\nwbegincode{267}\moddef{Helicities: helicity: TBP}\plusendmoddef
  procedure :: is_defined => helicity_is_defined
\nwendcode{}\nwbegincode{268}\moddef{Helicities: procedures}\plusendmoddef
  elemental function helicity_is_defined (hel) result (defined)
    logical :: defined
    class(helicity_t), intent(in) :: hel
    defined = hel%defined
  end function helicity_is_defined

\eatline
\nwendcode{}\nwbegindocs{269}Return true if the two helicities are equal or the particle is unpolarized:
\nwenddocs{}\nwbegincode{270}\moddef{Helicities: helicity: TBP}\plusendmoddef
  procedure :: is_diagonal => helicity_is_diagonal
\nwendcode{}\nwbegincode{271}\moddef{Helicities: procedures}\plusendmoddef
  elemental function helicity_is_diagonal (hel) result (diagonal)
    logical :: diagonal
    class(helicity_t), intent(in) :: hel
    if (hel%defined) then
       diagonal = hel%h1 == hel%h2
    else
       diagonal = .true.
    end if
  end function helicity_is_diagonal

\eatline
\nwendcode{}\nwbegindocs{272}\nwdocspar
\subsection{Accessing contents}
This returns a two-element array and thus cannot be elemental.  The
result is unpredictable if the helicity is undefined.
\nwenddocs{}\nwbegincode{273}\moddef{Helicities: helicity: TBP}\plusendmoddef
  procedure :: to_pair => helicity_to_pair
\nwendcode{}\nwbegincode{274}\moddef{Helicities: procedures}\plusendmoddef
  pure function helicity_to_pair (hel) result (h)
    integer, dimension(2) :: h
    class(helicity_t), intent(in) :: hel
    h(1) = hel%h2
    h(2) = hel%h1
  end function helicity_to_pair

\eatline
\nwendcode{}\nwbegindocs{275}\nwdocspar
\subsection{Comparisons}
When comparing helicities, if either one is undefined, they are
considered to match.  In other words, an unpolarized particle matches
any polarization.  In the {\tt{}dmatch} variant, it matches only diagonal
helicity.
\nwenddocs{}\nwbegincode{276}\moddef{Helicities: helicity: TBP}\plusendmoddef
  generic :: operator(.match.) => helicity_match
  generic :: operator(.dmatch.) => helicity_match_diagonal
  generic :: operator(==) => helicity_eq
  generic :: operator(/=) => helicity_neq
  procedure, private ::  helicity_match
  procedure, private ::  helicity_match_diagonal
  procedure, private ::  helicity_eq
  procedure, private ::  helicity_neq
\eatline
\nwendcode{}\nwbegincode{277}\moddef{Helicities: procedures}\plusendmoddef
  elemental function helicity_match (hel1, hel2) result (eq)
    logical :: eq
    class(helicity_t), intent(in) :: hel1, hel2
    if (hel1%defined .and. hel2%defined) then
       eq = (hel1%h1 == hel2%h1) .and. (hel1%h2 == hel2%h2)
    else
       eq = .true.
    end if
  end function helicity_match

  elemental function helicity_match_diagonal (hel1, hel2) result (eq)
    logical :: eq
    class(helicity_t), intent(in) :: hel1, hel2
    if (hel1%defined .and. hel2%defined) then
       eq = (hel1%h1 == hel2%h1) .and. (hel1%h2 == hel2%h2)
    else if (hel1%defined) then
       eq = hel1%h1 == hel1%h2
    else if (hel2%defined) then
       eq = hel2%h1 == hel2%h2
    else
       eq = .true.
    end if
  end function helicity_match_diagonal

\eatline
\nwendcode{}\nwbegincode{278}\moddef{Helicities: procedures}\plusendmoddef
  elemental function helicity_eq (hel1, hel2) result (eq)
    logical :: eq
    class(helicity_t), intent(in) :: hel1, hel2
    if (hel1%defined .and. hel2%defined) then
       eq = (hel1%h1 == hel2%h1) .and. (hel1%h2 == hel2%h2)
    else if (.not. hel1%defined .and. .not. hel2%defined) then
       eq = .true.
    else
       eq = .false.
    end if
  end function helicity_eq

\eatline
\nwendcode{}\nwbegincode{279}\moddef{Helicities: procedures}\plusendmoddef
  elemental function helicity_neq (hel1, hel2) result (neq)
    logical :: neq
    class(helicity_t), intent(in) :: hel1, hel2
    if (hel1%defined .and. hel2%defined) then
       neq = (hel1%h1 /= hel2%h1) .or. (hel1%h2 /= hel2%h2)
    else if (.not. hel1%defined .and. .not. hel2%defined) then
       neq = .false.
    else
       neq = .true.
    end if
  end function helicity_neq

\eatline
\nwendcode{}\nwbegindocs{280}\nwdocspar
\subsection{Tools}
Merge two helicity objects by taking the first entry from the first and
the second entry from the second argument.  Makes sense only if the
input helicities were defined and diagonal.  The handling of ghost
flags is not well-defined; one should verify beforehand that they
match.
\nwenddocs{}\nwbegincode{281}\moddef{Helicities: helicity: TBP}\plusendmoddef
  generic :: operator(.merge.) => merge_helicities
  procedure, private ::  merge_helicities
\eatline
\nwendcode{}\nwbegincode{282}\moddef{Helicities: procedures}\plusendmoddef
  elemental function merge_helicities (hel1, hel2) result (hel)
    type(helicity_t) :: hel
    class(helicity_t), intent(in) :: hel1, hel2
    if (hel1%defined .and. hel2%defined) then
       call helicity_init2 (hel, hel2%h1, hel1%h1)       
    else if (hel1%defined) then
       call helicity_init2 (hel, hel1%h2, hel1%h1)
    else if (hel2%defined) then
       call helicity_init2 (hel, hel2%h2, hel2%h1)
    end if
  end function merge_helicities

\eatline
\nwendcode{}\nwbegindocs{283}\nwdocspar
\clearpage
%%%%%%%%%%%%%%%%%%%%%%%%%%%%%%%%%%%%%%%%%%%%%%%%%%%%%%%%%%%%%%%%%%%%%%%%
\section{Colors}
This module defines a type and tools for dealing with color information.

Each particle can have zero or more (in practice, usually not more
than two) color indices.  Color indices are positive; flow direction
can be determined from the particle nature.

While parton shower matrix elements are diagonal in color, some
special applications (e.g., subtractions for NLO matrix elements)
require non-diagonal color matrices.
\nwenddocs{}\nwbegincode{284}\moddef{\code{}colors.f90\edoc{}}\endmoddef
\LA{}File header\RA{}

module colors

\LA{}Use kinds\RA{}
  use io_units
  use unit_tests
  use diagnostics
  
\LA{}Standard module head\RA{}

\LA{}Colors: public\RA{}

\LA{}Colors: types\RA{}

\LA{}Colors: interfaces\RA{}

contains

\LA{}Colors: procedures\RA{}

\LA{}Colors: tests\RA{}

end module colors
\eatline
\nwendcode{}\nwbegindocs{285}\nwdocspar
\subsection{The color type}
A particle may have an arbitrary number of color indices (in practice,
from zero to two, but more are possible).  This object acts as a
container.  (The current implementation has a fixed array of length two.)

The fact that color comes as an array prohibits elemental procedures
in some places.  (May add interfaces and multi versions where
necessary.)

The color may be undefined.

NOTE: Due to a compiler bug in nagfor 5.2, we do not use allocatable
but fixed-size arrays with dimension 2.  Only nonzero entries count.
This may be more efficient anyway, but gives up some flexibility.
However, the squaring algorithm currently works only for singlets,
(anti)triplets and octets anyway, so two components are enough.

This type has to be generalized (abstract type and specific
implementations) when trying to pursue generalized color flows or
Monte Carlo over continuous color.
\nwenddocs{}\nwbegincode{286}\moddef{Colors: public}\endmoddef
  public :: color_t
\nwendcode{}\nwbegincode{287}\moddef{Colors: types}\endmoddef
  type :: color_t
     private
     logical :: defined = .false.
     integer, dimension(2) :: c1 = 0, c2 = 0
     logical :: ghost = .false.
   contains
   \LA{}Colors: color: TBP\RA{}
  end type color_t

\eatline
\nwendcode{}\nwbegindocs{288}Initializers:
\nwenddocs{}\nwbegincode{289}\moddef{Colors: color: TBP}\endmoddef
  generic :: init => &
       color_init_trivial, color_init_trivial_ghost, &
       color_init_array, color_init_array_ghost, &
       color_init_arrays, color_init_arrays_ghost
  procedure, private :: color_init_trivial
  procedure, private :: color_init_trivial_ghost
  procedure, private :: color_init_array
  procedure, private :: color_init_array_ghost
  procedure, private :: color_init_arrays
  procedure, private :: color_init_arrays_ghost
\nwendcode{}\nwbegindocs{290}Undefined color: array remains unallocated
\nwenddocs{}\nwbegincode{291}\moddef{Colors: procedures}\endmoddef
  pure subroutine color_init_trivial (col)
    class(color_t), intent(inout) :: col
    col%defined = .true.
    col%c1 = 0
    col%c2 = 0
    col%ghost = .false.
  end subroutine color_init_trivial

  pure subroutine color_init_trivial_ghost (col, ghost)
    class(color_t), intent(inout) :: col
    logical, intent(in) :: ghost
    col%defined = .true.
    col%c1 = 0
    col%c2 = 0
    col%ghost = ghost
  end subroutine color_init_trivial_ghost

\nwendcode{}\nwbegindocs{292}This defines color from an arbitrary length color array, suitable
for any representation.  We may have two color arrays (non-diagonal
matrix elements).  This cannot be elemental.  The third version
assigns an array of colors, using a two-dimensional array as input.  
\nwenddocs{}\nwbegincode{293}\moddef{Colors: procedures}\plusendmoddef
  pure subroutine color_init_array (col, c1)
    class(color_t), intent(inout) :: col
    integer, dimension(:), intent(in) :: c1
    col%defined = .true.
    col%c1 = pack (c1, c1 /= 0, [0,0])
    col%c2 = col%c1
    col%ghost = .false.
  end subroutine color_init_array

  pure subroutine color_init_array_ghost (col, c1, ghost)
    class(color_t), intent(inout) :: col
    integer, dimension(:), intent(in) :: c1
    logical, intent(in) :: ghost
    call color_init_array (col, c1)
    col%ghost = ghost
  end subroutine color_init_array_ghost

  pure subroutine color_init_arrays (col, c1, c2)
    class(color_t), intent(inout) :: col
    integer, dimension(:), intent(in) :: c1, c2
    col%defined = .true.
    if (size (c1) == size (c2)) then
       col%c1 = pack (c1, c1 /= 0, [0,0])
       col%c2 = pack (c2, c2 /= 0, [0,0])
    else if (size (c1) /= 0) then
       col%c1 = pack (c1, c1 /= 0, [0,0])
       col%c2 = col%c1
    else if (size (c2) /= 0) then
       col%c1 = pack (c2, c2 /= 0, [0,0])
       col%c2 = col%c1
    end if
    col%ghost = .false.
  end subroutine color_init_arrays

  pure subroutine color_init_arrays_ghost (col, c1, c2, ghost)
    class(color_t), intent(inout) :: col
    integer, dimension(:), intent(in) :: c1, c2
    logical, intent(in) :: ghost
    call color_init_arrays (col, c1, c2)
    col%ghost = ghost
  end subroutine color_init_arrays_ghost

\eatline
\nwendcode{}\nwbegindocs{294}This version is restricted to singlets, triplets, antitriplets, and
octets: The input contains the color and anticolor index, each of the
may be zero.
\nwenddocs{}\nwbegincode{295}\moddef{Colors: color: TBP}\plusendmoddef
  procedure :: init_col_acl => color_init_col_acl
\nwendcode{}\nwbegincode{296}\moddef{Colors: procedures}\plusendmoddef
  elemental subroutine color_init_col_acl (col, col_in, acl_in)
    class(color_t), intent(inout) :: col
    integer, intent(in) :: col_in, acl_in
    integer, dimension(0) :: null_array
    select case (col_in)
    case (0)
       select case (acl_in)
       case (0)
          call color_init_array (col, null_array)
       case default
          call color_init_array (col, [-acl_in])
       end select
    case default
       select case (acl_in)
       case (0)
          call color_init_array (col, [col_in])
       case default
          call color_init_array (col, [col_in, -acl_in])
       end select
    end select
  end subroutine color_init_col_acl

\eatline
\nwendcode{}\nwbegindocs{297}This version is used for the external interface.  We convert a
fixed-size array of colors (for each particle) to the internal form by
packing only the nonzero entries.

Some of these procedures produce an arry, so they can't be all
type-bound.  We implement them as ordinary procedures.
\nwenddocs{}\nwbegincode{298}\moddef{Colors: public}\plusendmoddef
  public :: color_init_from_array
\nwendcode{}\nwbegincode{299}\moddef{Colors: interfaces}\endmoddef
  interface color_init_from_array
     module procedure color_init_from_array1
     module procedure color_init_from_array1g
     module procedure color_init_from_array2
     module procedure color_init_from_array2g
  end interface color_init_from_array
  
\eatline
\nwendcode{}\nwbegincode{300}\moddef{Colors: procedures}\plusendmoddef
  pure subroutine color_init_from_array1 (col, c1)
    type(color_t), intent(inout) :: col
    integer, dimension(:), intent(in) :: c1
    logical, dimension(size(c1)) :: mask
    mask = c1 /= 0
    col%defined = .true.
    col%c1 = pack (c1, mask, col%c1)
    col%c2 = col%c1
    col%ghost = .false.
  end subroutine color_init_from_array1

  pure subroutine color_init_from_array1g (col, c1, ghost)
    type(color_t), intent(inout) :: col
    integer, dimension(:), intent(in) :: c1
    logical, intent(in) :: ghost
    call color_init_from_array1 (col, c1)
    col%ghost = ghost
  end subroutine color_init_from_array1g

  pure subroutine color_init_from_array2 (col, c1)
    integer, dimension(:,:), intent(in) :: c1
    type(color_t), dimension(:), intent(inout) :: col
    integer :: i
    do i = 1, size (c1,2)
       call color_init_from_array1 (col(i), c1(:,i))
    end do
  end subroutine color_init_from_array2

  pure subroutine color_init_from_array2g (col, c1, ghost)
    integer, dimension(:,:), intent(inout) :: c1
    type(color_t), dimension(:), intent(out) :: col
    logical, intent(in), dimension(:) :: ghost
    call color_init_from_array2 (col, c1)
    col%ghost = ghost
  end subroutine color_init_from_array2g

\eatline
\nwendcode{}\nwbegindocs{301}Set the ghost property
\nwenddocs{}\nwbegincode{302}\moddef{Colors: color: TBP}\plusendmoddef
  procedure :: set_ghost => color_set_ghost
\nwendcode{}\nwbegincode{303}\moddef{Colors: procedures}\plusendmoddef
  elemental subroutine color_set_ghost (col, ghost)
    class(color_t), intent(inout) :: col
    logical, intent(in) :: ghost
    col%ghost = ghost
  end subroutine color_set_ghost

\eatline
\nwendcode{}\nwbegindocs{304}Undefine the color state:
\nwenddocs{}\nwbegincode{305}\moddef{Colors: color: TBP}\plusendmoddef
  procedure :: undefine => color_undefine
\nwendcode{}\nwbegincode{306}\moddef{Colors: procedures}\plusendmoddef
  elemental subroutine color_undefine (col, undefine_ghost)
    class(color_t), intent(inout) :: col
    logical, intent(in), optional :: undefine_ghost
    col%defined = .false.
    if (present (undefine_ghost)) then
       if (undefine_ghost)  col%ghost = .false.
    else
       col%ghost = .false.
    end if
  end subroutine color_undefine

\eatline
\nwendcode{}\nwbegindocs{307}Output.  As dense as possible, no linebreak.  If color is undefined,
no output.

The separate version for a color array suggest two distinct interfaces.
\nwenddocs{}\nwbegincode{308}\moddef{Colors: public}\plusendmoddef
  public :: color_write
\nwendcode{}\nwbegincode{309}\moddef{Colors: interfaces}\plusendmoddef
  interface color_write
     module procedure color_write_single
     module procedure color_write_array
  end interface color_write
  
\nwendcode{}\nwbegincode{310}\moddef{Colors: color: TBP}\plusendmoddef
  procedure :: write => color_write_single
\nwendcode{}\nwbegincode{311}\moddef{Colors: procedures}\plusendmoddef
  subroutine color_write_single (col, unit)
    class(color_t), intent(in) :: col
    integer, intent(in), optional :: unit
    integer :: u
    u = given_output_unit (unit);  if (u < 0)  return
    if (col%ghost) then
       write (u, "(A)", advance="no")  "c*"
    else if (col%defined) then
       write (u, "(A)", advance="no")  "c("
       if (col%c1(1) /= 0)  write (u, "(I0)", advance="no")  col%c1(1)
       if (any (col%c1 /= 0))  write (u, "(1x)", advance="no")
       if (col%c1(2) /= 0)  write (u, "(I0)", advance="no")  col%c1(2)
       if (.not. col%is_diagonal ()) then
          write (u, "(A)", advance="no")  "|"
          if (col%c2(1) /= 0)  write (u, "(I0)", advance="no")  col%c2(1)
          if (any (col%c2 /= 0))  write (u, "(1x)", advance="no")
          if (col%c2(2) /= 0)  write (u, "(I0)", advance="no")  col%c2(2)
       end if
       write (u, "(A)", advance="no") ")"
    end if
  end subroutine color_write_single

  subroutine color_write_array (col, unit)
    type(color_t), dimension(:), intent(in) :: col
    integer, intent(in), optional :: unit
    integer :: u
    integer :: i
    u = given_output_unit (unit);  if (u < 0)  return
    write (u, "(A)", advance="no") "["
    do i = 1, size (col)
       if (i > 1)  write (u, "(1x)", advance="no")
       call color_write_single (col(i), u)
    end do
    write (u, "(A)", advance="no") "]"
  end subroutine color_write_array

\eatline
\nwendcode{}\nwbegindocs{312}Binary I/O.  For allocatable colors, this would have to be modified.
\nwenddocs{}\nwbegincode{313}\moddef{Colors: color: TBP}\plusendmoddef
  procedure :: write_raw => color_write_raw
  procedure :: read_raw => color_read_raw
\nwendcode{}\nwbegincode{314}\moddef{Colors: procedures}\plusendmoddef
  subroutine color_write_raw (col, u)
    class(color_t), intent(in) :: col
    integer, intent(in) :: u
    logical :: defined
    defined = col%is_defined () .or. col%is_ghost ()
    write (u) defined
    if (defined) then
       write (u) col%c1, col%c2
       write (u) col%ghost
    end if
  end subroutine color_write_raw
    
  subroutine color_read_raw (col, u, iostat)
    class(color_t), intent(inout) :: col
    integer, intent(in) :: u
    integer, intent(out), optional :: iostat
    logical :: defined
    read (u, iostat=iostat) col%defined
    if (col%defined) then
       read (u, iostat=iostat) col%c1, col%c2
       read (u, iostat=iostat) col%ghost
    end if
  end subroutine color_read_raw

\eatline
\nwendcode{}\nwbegindocs{315}\nwdocspar
\subsection{Predicates}
Return the definition status.  A color state may be defined but trivial.
\nwenddocs{}\nwbegincode{316}\moddef{Colors: color: TBP}\plusendmoddef
  procedure :: is_defined => color_is_defined
  procedure :: is_nonzero => color_is_nonzero
\nwendcode{}\nwbegincode{317}\moddef{Colors: procedures}\plusendmoddef
  elemental function color_is_defined (col) result (defined)
    logical :: defined
    class(color_t), intent(in) :: col
    defined = col%defined
  end function color_is_defined

  elemental function color_is_nonzero (col) result (flag)
    logical :: flag
    class(color_t), intent(in) :: col
    flag = col%defined &
         .and. .not. col%ghost &
         .and. any (col%c1 /= 0 .or. col%c2 /= 0)
  end function color_is_nonzero

\eatline
\eatline
\nwendcode{}\nwbegindocs{318}Diagonal color objects have only one array allocated:
\nwenddocs{}\nwbegincode{319}\moddef{Colors: color: TBP}\plusendmoddef
  procedure :: is_diagonal => color_is_diagonal
\nwendcode{}\nwbegincode{320}\moddef{Colors: procedures}\plusendmoddef
  elemental function color_is_diagonal (col) result (diagonal)
    logical :: diagonal
    class(color_t), intent(in) :: col
    if (col%defined) then
       diagonal = all (col%c1 == col%c2)
    else
       diagonal = .true.
    end if
  end function color_is_diagonal

\eatline
\nwendcode{}\nwbegindocs{321}Return the ghost flag
\nwenddocs{}\nwbegincode{322}\moddef{Colors: color: TBP}\plusendmoddef
  procedure :: is_ghost => color_is_ghost
\nwendcode{}\nwbegincode{323}\moddef{Colors: procedures}\plusendmoddef
  elemental function color_is_ghost (col) result (ghost)
    logical :: ghost
    class(color_t), intent(in) :: col
    ghost = col%ghost
  end function color_is_ghost

\eatline
\nwendcode{}\nwbegindocs{324}The ghost parity: true if the color-ghost flag is set.  Again, no
TBP since this is an array.
\nwenddocs{}\nwbegincode{325}\moddef{Colors: procedures}\plusendmoddef
  pure function color_ghost_parity (col) result (parity)
    type(color_t), dimension(:), intent(in) :: col
    logical :: parity
    parity = mod (count (col%ghost), 2) == 1
  end function color_ghost_parity

\eatline
\nwendcode{}\nwbegindocs{326}\nwdocspar
\subsection{Accessing contents}
Return the number of color indices.  We assume that it is identical
for both arrays.
\nwenddocs{}\nwbegincode{327}\moddef{Colors: color: TBP}\plusendmoddef
  procedure, private :: get_number_of_indices => color_get_number_of_indices
\nwendcode{}\nwbegincode{328}\moddef{Colors: procedures}\plusendmoddef
  elemental function color_get_number_of_indices (col) result (n)
    integer :: n
    class(color_t), intent(in) :: col
    if (col%defined .and. .not. col%ghost) then
       n = count (col%c1 /= 0)
    else
       n = 0
    end if
  end function color_get_number_of_indices

\eatline
\nwendcode{}\nwbegindocs{329}Return the (first) color/anticolor entry (assuming that color is
diagonal).  The result is a positive color index.
\nwenddocs{}\nwbegincode{330}\moddef{Colors: color: TBP}\plusendmoddef
  procedure :: get_col => color_get_col
  procedure :: get_acl => color_get_acl
\nwendcode{}\nwbegincode{331}\moddef{Colors: procedures}\plusendmoddef
  elemental function color_get_col (col) result (c)
    integer :: c
    class(color_t), intent(in) :: col
    integer :: i
    if (col%defined .and. .not. col%ghost) then
       do i = 1, size (col%c1)
          if (col%c1(i) > 0) then
             c = col%c1(i)
             return
          end if
       end do
    end if
    c = 0
  end function color_get_col

  elemental function color_get_acl (col) result (c)
    integer :: c
    class(color_t), intent(in) :: col
    integer :: i
    if (col%defined .and. .not. col%ghost) then
       do i = 1, size (col%c1)
          if (col%c1(i) < 0) then
             c = - col%c1(i)
             return
          end if
       end do
    end if
    c = 0
  end function color_get_acl

\eatline
\nwendcode{}\nwbegindocs{332}Return the color index with highest absolute value
\nwenddocs{}\nwbegincode{333}\moddef{Colors: public}\plusendmoddef
  public :: color_get_max_value
\nwendcode{}\nwbegincode{334}\moddef{Colors: interfaces}\plusendmoddef
  interface color_get_max_value
     module procedure color_get_max_value0
     module procedure color_get_max_value1
     module procedure color_get_max_value2
  end interface color_get_max_value
  
\nwendcode{}\nwbegincode{335}\moddef{Colors: procedures}\plusendmoddef
  elemental function color_get_max_value0 (col) result (cmax)
    integer :: cmax
    type(color_t), intent(in) :: col
    if (col%defined .and. .not. col%ghost) then
       cmax = maxval (abs (col%c1))
    else
       cmax = 0
    end if
  end function color_get_max_value0

  pure function color_get_max_value1 (col) result (cmax)
    integer :: cmax
    type(color_t), dimension(:), intent(in) :: col
    cmax = maxval (color_get_max_value0 (col))
  end function color_get_max_value1

  pure function color_get_max_value2 (col) result (cmax)
    integer :: cmax
    type(color_t), dimension(:,:), intent(in) :: col
    integer, dimension(size(col, 2)) :: cm
    integer :: i
    forall (i = 1:size(col, 2))
       cm(i) = color_get_max_value1 (col(:,i))
    end forall
    cmax = maxval (cm)
  end function color_get_max_value2

\eatline
\nwendcode{}\nwbegindocs{336}\nwdocspar
\subsection{Comparisons}
Similar to helicities, colors match if they are equal, or if either
one is undefined.
\nwenddocs{}\nwbegincode{337}\moddef{Colors: color: TBP}\plusendmoddef
  generic :: operator(.match.) => color_match
  generic :: operator(==) => color_eq
  generic :: operator(/=) => color_neq
  procedure, private ::  color_match
  procedure, private ::  color_eq
  procedure, private ::  color_neq
\eatline
\nwendcode{}\nwbegincode{338}\moddef{Colors: procedures}\plusendmoddef
  elemental function color_match (col1, col2) result (eq)
    logical :: eq
    class(color_t), intent(in) :: col1, col2
    if (col1%defined .and. col2%defined) then
       if (col1%ghost .and. col2%ghost) then
          eq = .true.
       else if (.not. col1%ghost .and. .not. col2%ghost) then
          eq = all (col1%c1 == col2%c1) .and. all (col1%c2 == col2%c2)
       else
          eq = .false.
       end if
    else
       eq = .true.
    end if
  end function color_match

  elemental function color_eq (col1, col2) result (eq)
    logical :: eq
    class(color_t), intent(in) :: col1, col2
    if (col1%defined .and. col2%defined) then
       if (col1%ghost .and. col2%ghost) then
          eq = .true.
       else if (.not. col1%ghost .and. .not. col2%ghost) then
          eq = all (col1%c1 == col2%c1) .and. all (col1%c2 == col2%c2)
       else
          eq = .false.
       end if
    else if (.not. col1%defined &
       .and. .not. col2%defined) then
       eq = col1%ghost .eqv. col2%ghost
    else
       eq = .false.
    end if
  end function color_eq

\eatline
\nwendcode{}\nwbegincode{339}\moddef{Colors: procedures}\plusendmoddef
  elemental function color_neq (col1, col2) result (neq)
    logical :: neq
    class(color_t), intent(in) :: col1, col2
    if (col1%defined .and. col2%defined) then
       if (col1%ghost .and. col2%ghost) then
          neq = .false.
       else if (.not. col1%ghost .and. .not. col2%ghost) then
          neq = any (col1%c1 /= col2%c1) .or. any (col1%c2 /= col2%c2)
       else
          neq = .true.
       end if
    else if (.not. col1%defined &
         .and. .not. col2%defined) then
       neq = col1%ghost .neqv. col2%ghost
    else
       neq = .true.
    end if
  end function color_neq

\eatline
\nwendcode{}\nwbegindocs{340}\nwdocspar
\subsection{Tools}
Shift color indices by a common offset.
\nwenddocs{}\nwbegincode{341}\moddef{Colors: color: TBP}\plusendmoddef
  procedure :: add_offset => color_add_offset
\nwendcode{}\nwbegincode{342}\moddef{Colors: procedures}\plusendmoddef
  elemental subroutine color_add_offset (col, offset)
    class(color_t), intent(inout) :: col
    integer, intent(in) :: offset
    if (col%defined .and. .not. col%ghost) then
       where (col%c1 /= 0)  col%c1 = col%c1 + sign (offset, col%c1)
       where (col%c2 /= 0)  col%c2 = col%c2 + sign (offset, col%c2)
    end if
  end subroutine color_add_offset

\eatline
\nwendcode{}\nwbegindocs{343}Reassign color indices for an array of colored particle in canonical
order.  The allocated size of the color map is such that two colors
per particle can be accomodated.

The algorithm works directly on the contents of the color objects, it
\nwenddocs{}\nwbegincode{344}\moddef{Colors: public}\plusendmoddef
  public :: color_canonicalize
\nwendcode{}\nwbegincode{345}\moddef{Colors: procedures}\plusendmoddef
  subroutine color_canonicalize (col)
    type(color_t), dimension(:), intent(inout) :: col
    integer, dimension(2*size(col)) :: map
    integer :: n_col, i, j, k
    n_col = 0
    do i = 1, size (col)
       if (col(i)%defined .and. .not. col(i)%ghost) then
          do j = 1, size (col(i)%c1)
             if (col(i)%c1(j) /= 0) then
                k = find (abs (col(i)%c1(j)), map(:n_col))
                if (k == 0) then
                   n_col = n_col + 1
                   map(n_col) = abs (col(i)%c1(j))
                   k = n_col
                end if
                col(i)%c1(j) = sign (k, col(i)%c1(j))
             end if
             if (col(i)%c2(j) /= 0) then
                k = find (abs (col(i)%c2(j)), map(:n_col))
                if (k == 0) then
                   n_col = n_col + 1
                   map(n_col) = abs (col(i)%c2(j))
                   k = n_col
                end if
                col(i)%c2(j) = sign (k, col(i)%c2(j))
             end if
          end do
       end if
    end do
  contains
    function find (c, array) result (k)
      integer :: k
      integer, intent(in) :: c
      integer, dimension(:), intent(in) :: array
      integer :: i
      k = 0
      do i = 1, size (array)
         if (c == array (i)) then
            k = i
            return
         end if
      end do
    end function find
  end subroutine color_canonicalize

\eatline
\nwendcode{}\nwbegindocs{346}Return an array of different color indices from an array of colors.
The last argument is a pseudo-color array, where the color entries
correspond to the position of the corresponding index entry in the
index array.  The colors are assumed to be diagonal.

The algorithm works directly on the contents of the color objects.
\nwenddocs{}\nwbegincode{347}\moddef{Colors: procedures}\plusendmoddef
  subroutine extract_color_line_indices (col, c_index, col_pos)
    type(color_t), dimension(:), intent(in) :: col
    integer, dimension(:), intent(out), allocatable :: c_index
    type(color_t), dimension(size(col)), intent(out) :: col_pos
    integer, dimension(:), allocatable :: c_tmp
    integer :: i, j, k, n, c
    allocate (c_tmp (sum (col%get_number_of_indices ())), source=0)
    n = 0
    SCAN1: do i = 1, size (col)
       if (col(i)%defined .and. .not. col(i)%ghost) then
          SCAN2: do j = 1, 2
             c = abs (col(i)%c1(j))
             if (c /= 0) then
                do k = 1, n
                   if (c_tmp(k) == c) then
                      col_pos(i)%c1(j) = k
                      cycle SCAN2
                   end if
                end do
                n = n + 1
                c_tmp(n) = c
                col_pos(i)%c1(j) = n
             end if
          end do SCAN2
       end if
    end do SCAN1
    allocate (c_index (n))
    c_index = c_tmp(1:n)
  end subroutine extract_color_line_indices

\eatline
\nwendcode{}\nwbegindocs{348}Given a color array, pairwise contract the color lines in all
possible ways and return the resulting array of arrays.  The input
color array must be diagonal, and each color should occur exactly
twice, once as color and once as anticolor.

Gluon entries with equal color and anticolor are explicitly excluded.

This algorithm is generic, but for long arrays it is neither
efficient, nor does it avoid duplicates.  It is intended for small
arrays, in particular for the state matrix of a structure-function
pair.

The algorithm works directly on the contents of the color objects, it
thus depends on the implementation.
\nwenddocs{}\nwbegincode{349}\moddef{Colors: public}\plusendmoddef
  public :: color_array_make_contractions
\nwendcode{}\nwbegincode{350}\moddef{Colors: procedures}\plusendmoddef
  subroutine color_array_make_contractions (col_in, col_out)
    type(color_t), dimension(:), intent(in) :: col_in
    type(color_t), dimension(:,:), intent(out), allocatable :: col_out
    type :: entry_t
       integer, dimension(:), allocatable :: map
       type(color_t), dimension(:), allocatable :: col
       type(entry_t), pointer :: next => null ()
       logical :: nlo_event = .false.
    end type entry_t
    type :: list_t
       integer :: n = 0
       type(entry_t), pointer :: first => null ()
       type(entry_t), pointer :: last => null ()
    end type list_t
    type(list_t) :: list
    type(entry_t), pointer :: entry
    integer, dimension(:), allocatable :: c_index
    type(color_t), dimension(size(col_in)) :: col_pos
    integer :: n_prt, n_c_index
    integer, dimension(:), allocatable :: map
    integer :: i, j, c
    n_prt = size (col_in)
    call extract_color_line_indices (col_in, c_index, col_pos)
    ! print *, c_index     !!! Debugging
    n_c_index = size (c_index)
    allocate (map (n_c_index))
    map = 0
    call list_append_if_valid (list, map)
    entry => list%first
    do while (associated (entry))
       do i = 1, n_c_index
          if (entry%map(i) == 0) then
             c = c_index(i)
             do j = i + 1, n_c_index
                if (entry%map(j) == 0) then
                   map = entry%map
                   map(i) = c
                   map(j) = c
                   call list_append_if_valid (list, map)
                end if
             end do
          end if
       end do
       entry => entry%next
    end do
    call list_to_array (list, col_out)
  contains
    subroutine list_append_if_valid (list, map)
      type(list_t), intent(inout) :: list
      integer, dimension(:), intent(in) :: map
      type(entry_t), pointer :: entry
      integer :: i, j, c, p
      entry => list%first
      do while (associated (entry))
         if (all (map == entry%map))  return
         entry => entry%next
      end do
      allocate (entry)
      allocate (entry%map (n_c_index))
      entry%map = map
      allocate (entry%col (n_prt))
      do i = 1, n_prt
         do j = 1, 2
            c = col_in(i)%c1(j)
            if (c /= 0) then
               p = col_pos(i)%c1(j)
               entry%col(i)%defined = .true.
               if (map(p) /= 0) then
                  entry%col(i)%c1(j) = sign (map(p), c)
               else
                  entry%col(i)%c1(j) = c
               endif
               entry%col(i)%c2(j) = entry%col(i)%c1(j)
            end if
         end do
         if (any (entry%col(i)%c1 /= 0) .and. &
              entry%col(i)%c1(1) == - entry%col(i)%c1(2))  return
      end do
      ! call color_write (entry%col); print *, map     !!! Debugging
      if (associated (list%last)) then
         list%last%next => entry
      else
         list%first => entry
      end if
      list%last => entry
      list%n = list%n + 1
    end subroutine list_append_if_valid
    subroutine list_to_array (list, col)
      type(list_t), intent(inout) :: list
      type(color_t), dimension(:,:), intent(out), allocatable :: col
      type(entry_t), pointer :: entry
      integer :: i
      allocate (col (n_prt, list%n - 1))
      do i = 0, list%n - 1
         entry => list%first
         list%first => list%first%next
         if (i /= 0)  col(:,i) = entry%col
         deallocate (entry)
      end do
      list%last => null ()
    end subroutine list_to_array
  end subroutine color_array_make_contractions

\eatline
\nwendcode{}\nwbegindocs{351}Invert the color index, switching from particle to antiparticle.
For gluons, we have to swap the order of color entries.
\nwenddocs{}\nwbegincode{352}\moddef{Colors: color: TBP}\plusendmoddef
  procedure :: invert => color_invert
\nwendcode{}\nwbegincode{353}\moddef{Colors: procedures}\plusendmoddef
  elemental subroutine color_invert (col)
    class(color_t), intent(inout) :: col
    if (col%defined .and. .not. col%ghost) then
       col%c1 = - col%c1
       col%c2 = - col%c2
       if (col%c1(1) < 0 .and. col%c1(2) > 0) then
          col%c1 = col%c1(2:1:-1)
          col%c2 = col%c2(2:1:-1)
       end if
    end if
  end subroutine color_invert

\eatline
\nwendcode{}\nwbegindocs{354}Make a color map for two matching color arrays.  The result is an
array of integer pairs.
\nwenddocs{}\nwbegincode{355}\moddef{Colors: public}\plusendmoddef
  public :: make_color_map
\nwendcode{}\nwbegincode{356}\moddef{Colors: interfaces}\plusendmoddef
  interface make_color_map
     module procedure color_make_color_map
  end interface make_color_map
  
\nwendcode{}\nwbegincode{357}\moddef{Colors: procedures}\plusendmoddef
  subroutine color_make_color_map (map, col1, col2)
    integer, dimension(:,:), intent(out), allocatable :: map
    type(color_t), dimension(:), intent(in) :: col1, col2
    integer, dimension(:,:), allocatable :: map1
    integer :: i, j, k
    allocate (map1 (2, 2 * sum (col1%get_number_of_indices ())))
    k = 0
    do i = 1, size (col1)
       if (col1(i)%defined .and. .not. col1(i)%ghost) then
          do j = 1, size (col1(i)%c1)
             if (col1(i)%c1(j) /= 0 &
                  .and. all (map1(1,:k) /= abs (col1(i)%c1(j)))) then
                k = k + 1
                map1(1,k) = abs (col1(i)%c1(j))
                map1(2,k) = abs (col2(i)%c1(j))
             end if
             if (col1(i)%c2(j) /= 0 &
                  .and. all (map1(1,:k) /= abs (col1(i)%c2(j)))) then
                k = k + 1
                map1(1,k) = abs (col1(i)%c2(j))
                map1(2,k) = abs (col2(i)%c2(j))
             end if
          end do
       end if
    end do
    allocate (map (2, k))
    map(:,:) = map1(:,:k)
  end subroutine color_make_color_map

\eatline
\nwendcode{}\nwbegindocs{358}Translate colors which have a match in the translation table (an
array of integer pairs).  Color that do not match an entry are simply
transferred; this is done by first transferring all components, then
modifiying entries where appropriate.
\nwenddocs{}\nwbegincode{359}\moddef{Colors: public}\plusendmoddef
  public :: color_translate
\nwendcode{}\nwbegincode{360}\moddef{Colors: interfaces}\plusendmoddef
  interface color_translate
     module procedure color_translate0
     module procedure color_translate0_offset
     module procedure color_translate1
  end interface color_translate
  
\nwendcode{}\nwbegincode{361}\moddef{Colors: procedures}\plusendmoddef
  subroutine color_translate0 (col, map)
    type(color_t), intent(inout) :: col
    integer, dimension(:,:), intent(in) :: map
    type(color_t) :: col_tmp
    integer :: i
    if (col%defined .and. .not. col%ghost) then
       col_tmp = col
       do i = 1, size (map,2)
          where (abs (col%c1) == map(1,i))  
             col_tmp%c1 = sign (map(2,i), col%c1)
          end where
          where (abs (col%c2) == map(1,i))  
             col_tmp%c2 = sign (map(2,i), col%c2)
          end where
       end do
       col = col_tmp
    end if
  end subroutine color_translate0

  subroutine color_translate0_offset (col, map, offset)
    type(color_t), intent(inout) :: col
    integer, dimension(:,:), intent(in) :: map
    integer, intent(in) :: offset
    logical, dimension(size(col%c1)) :: mask1, mask2
    type(color_t) :: col_tmp
    integer :: i
    if (col%defined .and. .not. col%ghost) then
       col_tmp = col
       mask1 = col%c1 /= 0
       mask2 = col%c2 /= 0
       do i = 1, size (map,2)
          where (abs (col%c1) == map(1,i))  
             col_tmp%c1 = sign (map(2,i), col%c1)
             mask1 = .false.
          end where
          where (abs (col%c2) == map(1,i))  
             col_tmp%c2 = sign (map(2,i), col%c2)
             mask2 = .false.
          end where
       end do
       col = col_tmp
       where (mask1)  col%c1 = sign (abs (col%c1) + offset, col%c1)
       where (mask2)  col%c2 = sign (abs (col%c2) + offset, col%c2)
    end if
  end subroutine color_translate0_offset

  subroutine color_translate1 (col, map, offset)
    type(color_t), dimension(:), intent(inout) :: col
    integer, dimension(:,:), intent(in) :: map
    integer, intent(in), optional :: offset
    integer :: i
    if (present (offset)) then
       do i = 1, size (col)
          call color_translate0_offset (col(i), map, offset)
       end do
    else
       do i = 1, size (col)
          call color_translate0 (col(i), map)
       end do
    end if
  end subroutine color_translate1

\eatline
\nwendcode{}\nwbegindocs{362}Merge two color objects by taking the first entry from the first and
the first entry from the second argument.  Makes sense only if the
input colors are defined (and diagonal).  If either one is undefined,
transfer the defined one.
\nwenddocs{}\nwbegincode{363}\moddef{Colors: color: TBP}\plusendmoddef
  generic :: operator(.merge.) => merge_colors
  procedure, private ::  merge_colors
\eatline
\nwendcode{}\nwbegincode{364}\moddef{Colors: procedures}\plusendmoddef
  elemental function merge_colors (col1, col2) result (col)
    type(color_t) :: col
    class(color_t), intent(in) :: col1, col2
    if (color_is_defined (col1) .and. color_is_defined (col2)) then
       if (color_is_ghost (col1) .and. color_is_ghost (col2)) then
          call color_init_trivial_ghost (col, .true.)
       else
          call color_init_arrays (col, col1%c1, col2%c1)
       end if
    else if (color_is_defined (col1)) then
       call color_init_array (col, col1%c1)
    else if (color_is_defined (col2)) then
       call color_init_array (col, col2%c1)
    end if
  end function merge_colors

\eatline
\nwendcode{}\nwbegindocs{365}Compute the color factor, given two interfering color arrays.
\nwenddocs{}\nwbegincode{366}\moddef{Colors: public}\plusendmoddef
  public :: compute_color_factor
\nwendcode{}\nwbegincode{367}\moddef{Colors: procedures}\plusendmoddef
  function compute_color_factor (col1, col2, nc) result (factor)
    real(default) :: factor
    type(color_t), dimension(:), intent(in) :: col1, col2
    integer, intent(in), optional :: nc
    type(color_t), dimension(size(col1)) :: col
    integer :: ncol, nloops, nghost
    ncol = 3;  if (present (nc))  ncol = nc
    col = col1 .merge. col2
    nloops = count_color_loops (col)
    nghost = count (col%is_ghost ())
    factor = real (ncol, default) ** (nloops - nghost)
    if (color_ghost_parity (col))  factor = - factor
  end function compute_color_factor

\eatline
\nwendcode{}\nwbegindocs{368}\nwdocspar
We have a pair of color index arrays which corresponds to a squared
matrix element.  We want to determine the number of color loops in
this square matrix element.  So we first copy the colors (stored in a
single color array with a pair of color lists in each entry) to a
temporary where the color indices are shifted by some offset.  We then
recursively follow each loop, starting at the first color that has the
offset, resetting the first color index to the loop index and each
further index to zero as we go.  We check that (a) each color index
occurs twice within the left (right) color array, (b) the loops are
closed, so we always come back to a line which has the loop index.

In order for the algorithm to work we have to conjugate the colors of
initial state particles (one for decays, two for scatterings) into
their corresponding anticolors of outgoing particles.
\nwenddocs{}\nwbegincode{369}\moddef{Colors: public}\plusendmoddef
  public :: count_color_loops
\nwendcode{}\nwbegincode{370}\moddef{Colors: procedures}\plusendmoddef
  function count_color_loops (col) result (count)
    integer :: count
    type(color_t), dimension(:), intent(in) :: col
    type(color_t), dimension(size(col)) :: cc
    integer :: i, n, offset
    ! print *, "Count color loops:"      !!! Debugging
    ! call color_write (col); print *    !!! Debugging
    cc = col
    n = size (cc)
    offset = n
    call color_add_offset (cc, offset)
    ! print *, offset                    !!! Debugging
    ! call color_write (cc); print *     !!! Debugging
    count = 0
    SCAN_LOOPS: do
       do i = 1, n
          ! print *, i, ':', cc(i)%c1    !!! Debugging 
          if (color_is_nonzero (cc(i))) then
             if (any (cc(i)%c1 > offset)) then
                ! print *, 'start', i    !!! Debugging
                count = count + 1
                call follow_line1 (pick_new_line (cc(i)%c1, count, 1))
                cycle SCAN_LOOPS
             end if
          end if
       end do
       exit SCAN_LOOPS
    end do SCAN_LOOPS
  contains
    function pick_new_line (c, reset_val, sgn) result (line)
      integer :: line
      integer, dimension(:), intent(inout) :: c
      integer, intent(in) :: reset_val
      integer, intent(in) :: sgn
      integer :: i
      if (any (c == count)) then
         line = count
      else
         do i = 1, size (c)
            if (sign (1, c(i)) == sgn .and. abs (c(i)) > offset) then
               line = c(i)
               c(i) = reset_val
               return
            end if
         end do
         call color_mismatch
      end if
    end function pick_new_line
    subroutine reset_line (c, line)
      integer, dimension(:), intent(inout) :: c
      integer, intent(in) :: line
      integer :: i
      do i = 1, size (c)
         if (c(i) == line) then
            c(i) = 0
            return
         end if
      end do
    end subroutine reset_line
    recursive subroutine follow_line1 (line)
      integer, intent(in) :: line
      integer :: i
      ! print *, 'follow line 1:', line     !!! Debugging
      if (line == count) then
         ! print *, 'loop closed'           !!! Debugging
         return
      end if
      do i = 1, n
         if (any (cc(i)%c1 == -line)) then
            call reset_line (cc(i)%c1, -line)
            ! print *, 'found', -line, ' resetting c1:'   !!! Debugging
            ! call color_write (cc); print *              !!! Debugging
            call follow_line2 (pick_new_line (cc(i)%c2, 0, sign (1, -line)))
            return
         end if
      end do
      call color_mismatch ()
    end subroutine follow_line1
    recursive subroutine follow_line2 (line)
      integer, intent(in) :: line
      integer :: i
      ! print *, 'follow line 2:', line     !!! Debugging   
      do i = 1, n
         if (any (cc(i)%c2 == -line)) then
            call reset_line (cc(i)%c2, -line)
            ! print *, 'found', -line, ' resetting c2:'   !!! Debugging
            ! call color_write (cc); print *              !!! Debugging
            call follow_line1 (pick_new_line (cc(i)%c1, 0, sign (1, -line)))
            return
         end if
      end do
      call color_mismatch ()
    end subroutine follow_line2
    subroutine color_mismatch ()
      call color_write (col)
      print *
      call msg_bug (" Color flow mismatch (color loops should be closed)")
    end subroutine color_mismatch
  end function count_color_loops

\eatline
\nwendcode{}\nwbegindocs{371}\nwdocspar
\subsection{Unit tests}
\nwenddocs{}\nwbegincode{372}\moddef{Colors: public}\plusendmoddef
  public :: color_test
\nwendcode{}\nwbegincode{373}\moddef{Colors: procedures}\plusendmoddef
  subroutine color_test (u, results)
    integer, intent(in) :: u
    type(test_results_t), intent(inout) :: results
  \LA{}Colors: execute tests\RA{}  
  end subroutine color_test
  
\eatline
\nwendcode{}\nwbegindocs{374}This is a color counting test.
\nwenddocs{}\nwbegincode{375}\moddef{Colors: execute tests}\endmoddef
  call test (color_1, "color_1", &
       "check color counting", &
       u, results)
\nwendcode{}\nwbegincode{376}\moddef{Colors: tests}\endmoddef
  subroutine color_1 (u)
    integer, intent(in) :: u
    type(color_t), dimension(4) :: col1, col2, col
    type(color_t), dimension(:), allocatable :: col3
    type(color_t), dimension(:,:), allocatable :: col_array
    integer :: count, i
    call col1%init_col_acl ([1, 0, 2, 3], [0, 1, 3, 2])
    col2 = col1
    call color_write (col1, u)
    write (u, "(A)")
    call color_write (col2, u)
    write (u, "(A)")
    col = col1 .merge. col2
    call color_write (col, u)
    write (u, "(A)")
    count = count_color_loops (col)
    write (u, "(A,I1)") "Number of color loops (3): ", count
    call col2%init_col_acl ([1, 0, 2, 3], [0, 2, 3, 1])
    call color_write (col1, u)
    write (u, "(A)")
    call color_write (col2, u)
    write (u, "(A)")
    col = col1 .merge. col2
    call color_write (col, u)
    write (u, "(A)")
    count = count_color_loops (col)
    write (u, "(A,I1)")  "Number of color loops (2): ", count
    write (u, "(A)")
    allocate (col3 (4))
    call color_init_from_array (col3, &
         reshape ([1, 0,   0, -1,  2, -3,  3, -2], & 
                  [2, 4]))
    call color_write (col3, u)
    write (u, "(A)")
    call color_array_make_contractions (col3, col_array)
    write (u, "(A)")  "Contractions:"
    do i = 1, size (col_array, 2)
       call color_write (col_array(:,i), u)
       write (u, "(A)")
    end do
    deallocate (col3)
    write (u, "(A)")
    allocate (col3 (6))
    call color_init_from_array (col3, &
         reshape ([1, -2,   3, 0,  0, -1,  2, -4,  -3, 0,  4, 0], & 
                  [2, 6]))
    call color_write (col3, u)
    write (u, "(A)")
    call color_array_make_contractions (col3, col_array)
    write (u, "(A)")  "Contractions:"
    do i = 1, size (col_array, 2)
       call color_write (col_array(:,i), u)
       write (u, "(A)")
    end do
  end subroutine color_1

\eatline
\nwendcode{}\nwbegindocs{377}\nwdocspar
%%%%%%%%%%%%%%%%%%%%%%%%%%%%%%%%%%%%%%%%%%%%%%%%%%%%%%%%%%%%%%%%%%%%%%%%
\subsection{The Madgraph color model}
This section describes the method for matrix element and color
flow calculation within Madgraph.

For each Feynman diagram, the colorless amplitude for a specified
helicity and momentum configuration (in- and out- combined) is
computed:
\begin{equation}
  A_d(p,h)
\end{equation}
Inserting color, the squared matrix element for definite helicity and
momentum is
\begin{equation}
  M^2(p,h) = \sum_{dd'} A_{d}(p,h)\,C_{dd'} A_{d'}^*(p,h)
\end{equation}
where $C_{dd'}$ describes the color interference of the two diagrams
$A_d$ and $A_d'$, which is independent of momentum and helicity and
can be calculated for each Feynman diagram pair by reducing it to the
corresponding color graph.  Obviously, one could combine all diagrams
with identical color structure, such that the index $d$ runs only over
different color graphs.  For colorless diagrams all elements of
$C_{dd'}$ are equal to unity.

The hermitian matrix $C_{dd'}$ is diagonalized once and for all, such
that it can be written in the form
\begin{equation}
  C_{dd'} = \sum_\lambda c_d^\lambda \lambda\, c_d^\lambda{}^*,
\end{equation}
where the eigenvectors $c_d$ are normalized,
\begin{equation}
  \sum_d |c_d^\lambda|^2 = 1,
\end{equation}
and the $\lambda$ values are the corresponding eigenvalues.  In the
colorless case, this means $c_d = 1/\sqrt{N_d}$ for all diagrams
($N_d=$ number of diagrams), and $\lambda=N_d$ is the only nonzero
eigenvalue.

Consequently, the squared matrix element for definite helicity and
momentum can also be written as
\begin{equation}
  M^2(p,h) = \sum_\lambda A_\lambda(p,h)\, \lambda\, A_\lambda(p,h)^*
\end{equation}
with
\begin{equation}
  A_\lambda(p,h) = \sum_d c_d^\lambda A_d(p,h).
\end{equation}
For generic spin density matrices, this is easily generalized to
\begin{equation}
  M^2(p,h,h') = \sum_\lambda A_\lambda(p,h)\, \lambda\, A_\lambda(p,h')^*
\end{equation}

To determine the color flow probabilities of a given momentum-helicity
configuration, the color flow amplitudes are calculated as
\begin{equation}
  a_f(p,h) = \sum_d \beta^f_d A_d(p,h),
\end{equation}
where the coefficients $\beta^f_d$ describe the amplitude for a given
Feynman diagram (or color graph) $d$ to correspond to a definite color
flow~$f$.  They are computed from $C_{dd'}$ by transforming this
matrix into the color flow basis and neglecting all off-diagonal
elements.  Again, these coefficients do not depend on momentum or
helicity and can therefore be calculated in advance.  This gives the
color flow transition matrix
\begin{equation}
  F^f(p,h,h') = a_f(p,h)\, a^*_f(p,h')
\end{equation}
which is assumed diagonal in color flow space and is separate from the
color-summed transition matrix $M^2$.  They are, however, equivalent
(up to a factor) to leading order in $1/N_c$, and using the color flow
transition matrix is appropriate for matching to hadronization.

Note that the color flow transition matrix is not normalized at this
stage.  To make use of it, we have to fold it with the in-state
density matrix to get a pseudo density matrix
\begin{equation}
  \hat\rho_{\rm out}^f(p,h_{\rm out},h'_{\rm out})
  = \sum_{h_{\rm in} h'_{\rm in}} F^f(p,h,h')\,
    \rho_{\rm in}(p,h_{\rm in},h'_{\rm in})
\end{equation}
which gets a meaning only after contracted with projections on the
outgoing helicity states $k_{\rm out}$, given as linear combinations
of helicity states with the unitary coefficient matrix $c(k_{\rm out},
h_{\rm out})$.  Then the probability of finding color flow $f$ when
the helicity state $k_{\rm out}$ is measured is given by
\begin{equation}
  P^f(p, k_{\rm out}) = Q^f(p, k_{\rm out}) / \sum_f Q^f(p, k_{\rm out})
\end{equation}
where
\begin{equation}
  Q^f(p, k_{\rm out}) = \sum_{h_{\rm out} h'_{\rm out}}
    c(k_{\rm out}, h_{\rm out})\,
    \hat\rho_{\rm out}^f(p,h_{\rm out},h'_{\rm out})\,
    c^*(k_{\rm out}, h'_{\rm out})
\end{equation}
However, if we can assume that the out-state helicity basis is the
canonical one, we can throw away the off diagonal elements in the
color flow density matrix and normalize the ones on the diagonal to obtain
\begin{equation}
  P^f(p, h_{\rm out}) = 
    \hat\rho_{\rm out}^f(p,h_{\rm out},h_{\rm out}) /
    \sum_f \hat\rho_{\rm out}^f(p,h_{\rm out},h_{\rm out})
\end{equation}

Finally, the color-summed out-state density matrix is computed by the
scattering formula
\begin{align}
  {\rho_{\rm out}(p,h_{\rm out},h'_{\rm out})}
  &=
    \sum_{h_{\rm in} h'_{\rm in}} M^2(p,h,h')\,
    \rho_{\rm in}(p,h_{\rm in},h'_{\rm in}) \\
  &= \sum_{h_{\rm in} h'_{\rm in} \lambda} 
     A_\lambda(p,h)\, \lambda\, A_\lambda(p,h')^*
     \rho_{\rm in}(p,h_{\rm in},h'_{\rm in}),
\end{align}
The trace of $\rho_{\rm out}$ is the squared matrix element, summed
over all internal degrees of freedom.  To get the squared matrix
element for a definite helicity $k_{\rm out}$ and color flow $f$, one
has to project the density matrix onto the given helicity state and
multiply with $P^f(p, k_{\rm out})$.

For diagonal helicities the out-state density reduces to
\begin{equation}
  \rho_{\rm out}(p,h_{\rm out})
  = \sum_{h_{\rm in}\lambda} 
     \lambda|A_\lambda(p,h)|^2 \rho_{\rm in}(p,h_{\rm in}).
\end{equation}
Since no basis transformation is involved, we can use the normalized
color flow probability $P^f(p, h_{\rm out})$ and express the result as
\begin{align}
  \rho_{\rm out}^f(p,h_{\rm out})
  &= \rho_{\rm out}(p,h_{\rm out})\,P^f(p, h_{\rm out}) \\
  &= \sum_{h_{\rm in}\lambda} 
     \frac{|a^f(p,h)|^2}{\sum_f|a^f(p,h)|^2}
     \lambda|A_\lambda(p,h)|^2 \rho_{\rm in}(p,h_{\rm in}).
\end{align}

From these considerations, the following calculation strategy can be
derived:
\begin{itemize}
\item
  Before the first event is generated, the color interference matrix
  $C_{dd'}$ is computed and diagonalized, so the eigenvectors
  $c^\lambda_d$, eigenvalues $\lambda$ and color flow coefficients
  $\beta^f_d$ are obtained.  In practice, these calculations are
  done when the matrix element code is generated, and the results are
  hardcoded in the matrix element subroutine as {\tt{}DATA} statements.
\item
  For each event, one loops over helicities once and stores the
  matrices $A_\lambda(p,h)$ and $a^f(p,h)$.  The allowed color flows,
  helicity combinations and eigenvalues are each labeled by integer
  indices, so one has to store complex matrices of dimension
  $N_\lambda\times N_h$ and $N_f\times N_h$, respectively.
\item
  The further strategy depends on the requested information.
  \begin{enumerate}
  \item
    If colorless diagonal helicity amplitudes are required, the
    eigenvalues $A_\lambda(p,h)$ are squared, summed with weight
    $\lambda$, and the result contracted with the in-state probability
    vector $\rho_{\rm in}(p, h_{\rm in})$.  The result is a
    probability vector $\rho_{\rm out}(p, h_{\rm out})$.
  \item
    For colored diagonal helicity amplitudes, the color coefficients
    $a^f(p,h)$ are also squared and used as weights to obtain the color-flow
    probability vector $\rho_{\rm out}^f(p, h_{\rm out})$.
  \item
    For colorless non-diagonal helicity amplitudes, we contract the
    tensor product of $A_\lambda(p,h)$ with $A_\lambda(p,h')$,
    weighted with $\lambda$, with the correlated in-state density
    matrix, to obtain a correlated out-state density matrix.
  \item
    In the general (colored, non-diagonal) case, we do the same as
    in the colorless case, but return the un-normalized color flow
    density matrix $\hat\rho_{\rm out}^f(p,h_{\rm out},h'_{\rm out})$
    in addition.  When the relevant helicity basis is known, the
    latter can be used by the caller program to determine flow
    probabilities.  (In reality, we assume the canonical basis and
    reduce the correlated out-state density to its diagonal immediately.)
  \end{enumerate}
\end{itemize}
\nwenddocs{}\nwbegindocs{378}\nwdocspar
\clearpage
%%%%%%%%%%%%%%%%%%%%%%%%%%%%%%%%%%%%%%%%%%%%%%%%%%%%%%%%%%%%%%%%%%%%%%%%
\section{Flavors: Particle properties}
This module contains a type for holding the flavor code, and all
functions that depend on the model, i.e., that determine particle
properties.

The PDG code is packed in a special {\tt{}flavor} type.  (This prohibits
meaningless operations, and it allows for a different implementation,
e.g., some non-PDG scheme internally, if appropiate at some point.)

There are lots of further particle properties that depend on the
model.  Implementing a flyweight pattern, the associated field data
object is to be stored in a central area, the {\tt{}flavor} object just
receives a pointer to this, so all queries can be delegated.
\nwenddocs{}\nwbegincode{379}\moddef{\code{}flavors.f90\edoc{}}\endmoddef
\LA{}File header\RA{}

module flavors

\LA{}Use kinds\RA{}
\LA{}Use strings\RA{}
  use io_units
  use diagnostics
  use physics_defs, only: UNDEFINED
  use physics_defs, only: INVALID
  use physics_defs, only: HADRON_REMNANT
  use physics_defs, only: HADRON_REMNANT_SINGLET
  use physics_defs, only: HADRON_REMNANT_TRIPLET
  use physics_defs, only: HADRON_REMNANT_OCTET
  use model_data
  use colors, only: color_t

\LA{}Standard module head\RA{}

\LA{}Flavors: public\RA{}

\LA{}Flavors: types\RA{}

\LA{}Flavors: interfaces\RA{}
       
contains

\LA{}Flavors: procedures\RA{}

end module flavors
\eatline
\nwendcode{}\nwbegindocs{380}\nwdocspar
\subsection{The flavor type}
The flavor type is an integer representing the PDG code, or
undefined (zero).  Negative codes represent antiflavors.  They should
be used only for particles which do have a distinct antiparticle.

The {\tt{}radiated} flag can be set for particles that are the result of
a beam-structure interaction (hadron beam remnant, ISR photon, etc.),
not of the hard interaction itself.

Further properties of the given flavor can be retrieved via the
particle-data pointer, if it is associated.
\nwenddocs{}\nwbegincode{381}\moddef{Flavors: public}\endmoddef
  public :: flavor_t
\nwendcode{}\nwbegincode{382}\moddef{Flavors: types}\endmoddef
  type :: flavor_t
     private
     integer :: f = UNDEFINED
     logical :: radiated = .false.
     type(field_data_t), pointer :: field_data => null ()
   contains
   \LA{}Flavors: flavor: TBP\RA{}
  end type flavor_t

\eatline
\nwendcode{}\nwbegindocs{383}Initializer form.  If the model is assigned, the procedure is
impure, therefore we have to define a separate array version.

Note: The pure elemental subroutines can't have an intent(out) CLASS
argument (because of the potential for an impure finalizer in a type
extension), so we stick to intent(inout) and (re)set all components
explicitly.
\nwenddocs{}\nwbegincode{384}\moddef{Flavors: flavor: TBP}\endmoddef
  generic :: init => &
       flavor_init0_empty, &
       flavor_init0, &
       flavor_init0_field_data, &
       flavor_init0_model, &
       flavor_init0_model_alt, &
       flavor_init0_name_model
  procedure, private :: flavor_init0_empty
  procedure, private :: flavor_init0
  procedure, private :: flavor_init0_field_data
  procedure, private :: flavor_init0_model
  procedure, private :: flavor_init0_model_alt
  procedure, private :: flavor_init0_name_model
\nwendcode{}\nwbegincode{385}\moddef{Flavors: procedures}\endmoddef
  elemental subroutine flavor_init0_empty (flv)
    class(flavor_t), intent(inout) :: flv
    flv%f = UNDEFINED
    flv%radiated = .false.
    flv%field_data => null ()
  end subroutine flavor_init0_empty

  elemental subroutine flavor_init0 (flv, f)
    class(flavor_t), intent(inout) :: flv
    integer, intent(in) :: f
    flv%f = f
    flv%radiated = .false.
    flv%field_data => null ()
  end subroutine flavor_init0

  impure elemental subroutine flavor_init0_field_data (flv, field_data)
    class(flavor_t), intent(inout) :: flv
    type(field_data_t), intent(in), target :: field_data
    flv%f = field_data%get_pdg ()
    flv%radiated = .false.
    flv%field_data => field_data
  end subroutine flavor_init0_field_data

  impure elemental subroutine flavor_init0_model (flv, f, model)
    class(flavor_t), intent(inout) :: flv
    integer, intent(in) :: f
    class(model_data_t), intent(in), target :: model
    flv%f = f
    flv%radiated = .false.
    flv%field_data => model%get_field_ptr (f, check=.true.)
  end subroutine flavor_init0_model

  impure elemental subroutine flavor_init0_model_alt (flv, f, model, alt_model)
    class(flavor_t), intent(inout) :: flv
    integer, intent(in) :: f
    class(model_data_t), intent(in), target :: model, alt_model
    flv%f = f
    flv%radiated = .false.
    flv%field_data => model%get_field_ptr (f, check=.false.)
    if (.not. associated (flv%field_data)) then
       flv%field_data => alt_model%get_field_ptr (f, check=.false.)
       if (.not. associated (flv%field_data)) then
          write (msg_buffer, "(A,1x,I0,1x,A,1x,A,1x,A,1x,A)") &
               "Particle with code", f, &
               "found neither in model", char (model%get_name ()), &
               "nor in model", char (alt_model%get_name ())
          call msg_fatal ()
       end if
    end if
  end subroutine flavor_init0_model_alt

!   subroutine flavor_init1_model (flv, f, model)
!     type(flavor_t), dimension(:), intent(out) :: flv
!     integer, dimension(:), intent(in) :: f
!     class(model_data_t), intent(in), target :: model
!     integer :: i
!     do i = 1, size (f)
!        call flavor_init0_model (flv(i), f(i), model)
!     end do
!   end subroutine flavor_init1_model
! 
!   subroutine flavor_init1_model_alt (flv, f, model, alt_model)
!     type(flavor_t), dimension(:), intent(out) :: flv
!     integer, dimension(:), intent(in) :: f
!     class(model_data_t), intent(in), target :: model, alt_model
!     integer :: i
!     do i = 1, size (f)
!        call flavor_init0_model_alt (flv(i), f(i), model, alt_model)
!     end do
!   end subroutine flavor_init1_model_alt
! 
!   subroutine flavor_init2_model (flv, f, model)
!     type(flavor_t), dimension(:,:), intent(out) :: flv
!     integer, dimension(:,:), intent(in) :: f
!     class(model_data_t), intent(in), target :: model
!     integer :: i
!     do i = 1, size (f, 2)
!        call flavor_init1_model (flv(:,i), f(:,i), model)
!     end do
!   end subroutine flavor_init2_model
! 
!   subroutine flavor_init2_model_alt (flv, f, model, alt_model)
!     type(flavor_t), dimension(:,:), intent(out) :: flv
!     integer, dimension(:,:), intent(in) :: f
!     class(model_data_t), intent(in), target :: model, alt_model
!     integer :: i
!     do i = 1, size (f, 2)
!        call flavor_init1_model_alt (flv(:,i), f(:,i), model, alt_model)
!     end do
!   end subroutine flavor_init2_model_alt
! 
  impure elemental subroutine flavor_init0_name_model (flv, name, model)
    class(flavor_t), intent(inout) :: flv
    type(string_t), intent(in) :: name
    class(model_data_t), intent(in), target :: model
    flv%f = model%get_pdg (name)
    flv%radiated = .false.
    flv%field_data => model%get_field_ptr (name, check=.true.)
  end subroutine flavor_init0_name_model
! 
!   subroutine flavor_init1_name_model (flv, name, model)
!     type(flavor_t), dimension(:), intent(out) :: flv
!     type(string_t), dimension(:), intent(in) :: name
!     class(model_data_t), intent(in), target :: model
!     integer :: i
!     do i = 1, size (name)
!        call flavor_init0_name_model (flv(i), name(i), model)
!     end do
!   end subroutine flavor_init1_name_model

\eatline
\nwendcode{}\nwbegindocs{386}Set the {\tt{}radiated} flag.
\nwenddocs{}\nwbegincode{387}\moddef{Flavors: flavor: TBP}\plusendmoddef
  procedure :: tag_radiated => flavor_tag_radiated
\nwendcode{}\nwbegincode{388}\moddef{Flavors: procedures}\plusendmoddef
  elemental subroutine flavor_tag_radiated (flv)
    class(flavor_t), intent(inout) :: flv
    flv%radiated = .true.
  end subroutine flavor_tag_radiated
  
\eatline
\nwendcode{}\nwbegindocs{389}Undefine the flavor state:
\nwenddocs{}\nwbegincode{390}\moddef{Flavors: flavor: TBP}\plusendmoddef
  procedure :: undefine => flavor_undefine
\nwendcode{}\nwbegincode{391}\moddef{Flavors: procedures}\plusendmoddef
  elemental subroutine flavor_undefine (flv)
    class(flavor_t), intent(inout) :: flv
    flv%f = UNDEFINED
    flv%field_data => null ()
  end subroutine flavor_undefine

\eatline
\nwendcode{}\nwbegindocs{392}Output: dense, no linebreak
\nwenddocs{}\nwbegincode{393}\moddef{Flavors: flavor: TBP}\plusendmoddef
  procedure :: write => flavor_write
\nwendcode{}\nwbegincode{394}\moddef{Flavors: procedures}\plusendmoddef
  subroutine flavor_write (flv, unit)
    class(flavor_t), intent(in) :: flv
    integer, intent(in), optional :: unit
    integer :: u
    u = given_output_unit (unit);  if (u < 0)  return
    if (associated (flv%field_data)) then
       write (u, "(A)", advance="no")  "f("
    else
       write (u, "(A)", advance="no")  "p("
    end if
    write (u, "(I0)", advance="no")  flv%f
    if (flv%radiated) then
       write (u, "('*')", advance="no")
    end if
    write (u, "(A)", advance="no")  ")"
  end subroutine flavor_write

\eatline
\nwendcode{}\nwbegindocs{395}Binary I/O.  Currently, the model information is not written/read,
so after reading the particle-data pointer is empty.
\nwenddocs{}\nwbegincode{396}\moddef{Flavors: flavor: TBP}\plusendmoddef
  procedure :: write_raw => flavor_write_raw
  procedure :: read_raw => flavor_read_raw
\nwendcode{}\nwbegincode{397}\moddef{Flavors: procedures}\plusendmoddef
  subroutine flavor_write_raw (flv, u)
    class(flavor_t), intent(in) :: flv
    integer, intent(in) :: u
    write (u) flv%f
    write (u) flv%radiated
  end subroutine flavor_write_raw

  subroutine flavor_read_raw (flv, u, iostat)
    class(flavor_t), intent(out) :: flv
    integer, intent(in) :: u
    integer, intent(out), optional :: iostat
    read (u, iostat=iostat) flv%f
    if (present (iostat)) then
       if (iostat /= 0)  return
    end if
    read (u, iostat=iostat)  flv%radiated
  end subroutine flavor_read_raw

\eatline
\nwendcode{}\nwbegindocs{398}\nwdocspar
\subsubsection{Assignment}
Default assignment of flavor objects is possible, but cannot be used
in pure procedures, because a pointer assignment is involved.

Assign the particle pointer separately.  This cannot be elemental, so
we define a scalar and an array version explicitly.  We refer to an
array of flavors, not an array of models.
\nwenddocs{}\nwbegincode{399}\moddef{Flavors: flavor: TBP}\plusendmoddef
  procedure :: set_model => flavor_set_model_single
\nwendcode{}\nwbegincode{400}\moddef{Flavors: procedures}\plusendmoddef
  impure elemental subroutine flavor_set_model_single (flv, model)
    class(flavor_t), intent(inout) :: flv
    class(model_data_t), intent(in), target :: model
    if (flv%f /= UNDEFINED) &
         flv%field_data => model%get_field_ptr (flv%f)
  end subroutine flavor_set_model_single

\eatline
\nwendcode{}\nwbegindocs{401}\nwdocspar
\subsubsection{Predicates}
Return the definition status.  By definition, the flavor object is
defined if the flavor PDG code is nonzero.
\nwenddocs{}\nwbegincode{402}\moddef{Flavors: flavor: TBP}\plusendmoddef
  procedure :: is_defined => flavor_is_defined
\nwendcode{}\nwbegincode{403}\moddef{Flavors: procedures}\plusendmoddef
  elemental function flavor_is_defined (flv) result (defined)
    class(flavor_t), intent(in) :: flv
    logical :: defined
    defined = flv%f /= UNDEFINED
  end function flavor_is_defined

\eatline
\nwendcode{}\nwbegindocs{404}Check for valid flavor (including undefined).  This is distinct from
the {\tt{}is{\char95}defined} status.  Invalid flavor is actually a specific PDG
code.
\nwenddocs{}\nwbegincode{405}\moddef{Flavors: flavor: TBP}\plusendmoddef
  procedure :: is_valid => flavor_is_valid
\nwendcode{}\nwbegincode{406}\moddef{Flavors: procedures}\plusendmoddef
  elemental function flavor_is_valid (flv) result (valid)
    class(flavor_t), intent(in) :: flv
    logical :: valid
    valid = flv%f /= INVALID
  end function flavor_is_valid

\eatline
\nwendcode{}\nwbegindocs{407}Return true if the particle-data pointer is associated.  (Debugging aid)
\nwenddocs{}\nwbegincode{408}\moddef{Flavors: flavor: TBP}\plusendmoddef
  procedure :: is_associated => flavor_is_associated
\nwendcode{}\nwbegincode{409}\moddef{Flavors: procedures}\plusendmoddef
  elemental function flavor_is_associated (flv) result (flag)
    class(flavor_t), intent(in) :: flv
    logical :: flag
    flag = associated (flv%field_data)
  end function flavor_is_associated

\eatline
\nwendcode{}\nwbegindocs{410}Check the {\tt{}radiated} flag.  A radiated particle has a definite PDG
flavor status, but it is actually a pseudoparticle (a beam remnant)
which may be subject to fragmentation.
\nwenddocs{}\nwbegincode{411}\moddef{Flavors: flavor: TBP}\plusendmoddef
  procedure :: is_radiated => flavor_is_radiated
\nwendcode{}\nwbegincode{412}\moddef{Flavors: procedures}\plusendmoddef
  elemental function flavor_is_radiated (flv) result (flag)
    class(flavor_t), intent(in) :: flv
    logical :: flag
    flag = flv%radiated
  end function flavor_is_radiated

\eatline
\nwendcode{}\nwbegindocs{413}\nwdocspar
\subsubsection{Accessing contents}
With the exception of the PDG code, all particle property enquiries are
delegated to the {\tt{}field{\char95}data} pointer.  If this is unassigned, some
access function will crash.

Return the flavor as an integer
\nwenddocs{}\nwbegincode{414}\moddef{Flavors: flavor: TBP}\plusendmoddef
  procedure :: get_pdg => flavor_get_pdg
\nwendcode{}\nwbegincode{415}\moddef{Flavors: procedures}\plusendmoddef
  elemental function flavor_get_pdg (flv) result (f)
    integer :: f
    class(flavor_t), intent(in) :: flv
    f = flv%f
  end function flavor_get_pdg

\eatline
\nwendcode{}\nwbegindocs{416}Return the flavor of the antiparticle
\nwenddocs{}\nwbegincode{417}\moddef{Flavors: flavor: TBP}\plusendmoddef
  procedure :: get_pdg_anti => flavor_get_pdg_anti
\nwendcode{}\nwbegincode{418}\moddef{Flavors: procedures}\plusendmoddef
  elemental function flavor_get_pdg_anti (flv) result (f)
    integer :: f
    class(flavor_t), intent(in) :: flv
    if (associated (flv%field_data)) then
       if (flv%field_data%has_antiparticle ()) then
          f = -flv%f
       else
          f = flv%f
       end if
    else
       f = 0
    end if
  end function flavor_get_pdg_anti

\eatline
\nwendcode{}\nwbegindocs{419}\nwdocspar
Absolute value:
\nwenddocs{}\nwbegincode{420}\moddef{Flavors: flavor: TBP}\plusendmoddef
  procedure :: get_pdg_abs => flavor_get_pdg_abs
\nwendcode{}\nwbegincode{421}\moddef{Flavors: procedures}\plusendmoddef
  elemental function flavor_get_pdg_abs (flv) result (f)
    integer :: f
    class(flavor_t), intent(in) :: flv
    f = abs (flv%f)
  end function flavor_get_pdg_abs

\eatline
\nwendcode{}\nwbegindocs{422}\nwdocspar
Generic properties
\nwenddocs{}\nwbegincode{423}\moddef{Flavors: flavor: TBP}\plusendmoddef
  procedure :: is_visible => flavor_is_visible
  procedure :: is_parton => flavor_is_parton
  procedure :: is_beam_remnant => flavor_is_beam_remnant
  procedure :: is_gauge => flavor_is_gauge
  procedure :: is_left_handed => flavor_is_left_handed
  procedure :: is_right_handed => flavor_is_right_handed
  procedure :: is_antiparticle => flavor_is_antiparticle
  procedure :: has_antiparticle => flavor_has_antiparticle
  procedure :: is_stable => flavor_is_stable
  procedure :: get_decays => flavor_get_decays
  procedure :: decays_isotropically => flavor_decays_isotropically
  procedure :: decays_diagonal => flavor_decays_diagonal
  procedure :: has_decay_helicity => flavor_has_decay_helicity
  procedure :: get_decay_helicity => flavor_get_decay_helicity
  procedure :: is_polarized => flavor_is_polarized
\nwendcode{}\nwbegincode{424}\moddef{Flavors: procedures}\plusendmoddef
  elemental function flavor_is_visible (flv) result (flag)
    logical :: flag
    class(flavor_t), intent(in) :: flv
    if (associated (flv%field_data)) then
       flag = flv%field_data%is_visible ()
    else
       flag = .false.
    end if
  end function flavor_is_visible

  elemental function flavor_is_parton (flv) result (flag)
    logical :: flag
    class(flavor_t), intent(in) :: flv
    if (associated (flv%field_data)) then
       flag = flv%field_data%is_parton ()
    else
       flag = .false.
    end if
  end function flavor_is_parton

  elemental function flavor_is_beam_remnant (flv) result (flag)
    logical :: flag
    class(flavor_t), intent(in) :: flv
    select case (abs (flv%f))
    case (HADRON_REMNANT, &
         HADRON_REMNANT_SINGLET, HADRON_REMNANT_TRIPLET, HADRON_REMNANT_OCTET)
       flag = .true.
    case default
       flag = .false.
    end select
  end function flavor_is_beam_remnant

  elemental function flavor_is_gauge (flv) result (flag)
    logical :: flag
    class(flavor_t), intent(in) :: flv
    if (associated (flv%field_data)) then
       flag = flv%field_data%is_gauge ()
    else
       flag = .false.
    end if
  end function flavor_is_gauge

  elemental function flavor_is_left_handed (flv) result (flag)
    logical :: flag
    class(flavor_t), intent(in) :: flv
    if (associated (flv%field_data)) then
       if (flv%f > 0) then
          flag = flv%field_data%is_left_handed ()
       else
          flag = flv%field_data%is_right_handed ()
       end if
    else
       flag = .false.
    end if
  end function flavor_is_left_handed

  elemental function flavor_is_right_handed (flv) result (flag)
    logical :: flag
    class(flavor_t), intent(in) :: flv
    if (associated (flv%field_data)) then
       if (flv%f > 0) then
          flag = flv%field_data%is_right_handed ()
       else
          flag = flv%field_data%is_left_handed ()
       end if
    else
       flag = .false.
    end if
  end function flavor_is_right_handed

  elemental function flavor_is_antiparticle (flv) result (flag)
    logical :: flag
    class(flavor_t), intent(in) :: flv
    flag = flv%f < 0
  end function flavor_is_antiparticle

  elemental function flavor_has_antiparticle (flv) result (flag)
    logical :: flag
    class(flavor_t), intent(in) :: flv
    if (associated (flv%field_data)) then
       flag = flv%field_data%has_antiparticle ()
    else
       flag = .false.
    end if
  end function flavor_has_antiparticle

  elemental function flavor_is_stable (flv) result (flag)
    logical :: flag
    class(flavor_t), intent(in) :: flv
    if (associated (flv%field_data)) then
       flag = flv%field_data%is_stable (anti = flv%f < 0)
    else
       flag = .true.
    end if
  end function flavor_is_stable

  subroutine flavor_get_decays (flv, decay)
    class(flavor_t), intent(in) :: flv
    type(string_t), dimension(:), intent(out), allocatable :: decay
    logical :: anti
    anti = flv%f < 0
    if (.not. flv%field_data%is_stable (anti)) then
       call flv%field_data%get_decays (decay, anti)
    end if
  end subroutine flavor_get_decays

  elemental function flavor_decays_isotropically (flv) result (flag)
    logical :: flag
    class(flavor_t), intent(in) :: flv
    if (associated (flv%field_data)) then
       flag = flv%field_data%decays_isotropically (anti = flv%f < 0)
    else
       flag = .true.
    end if
  end function flavor_decays_isotropically

  elemental function flavor_decays_diagonal (flv) result (flag)
    logical :: flag
    class(flavor_t), intent(in) :: flv
    if (associated (flv%field_data)) then
       flag = flv%field_data%decays_diagonal (anti = flv%f < 0)
    else
       flag = .true.
    end if
  end function flavor_decays_diagonal

  elemental function flavor_has_decay_helicity (flv) result (flag)
    logical :: flag
    class(flavor_t), intent(in) :: flv
    if (associated (flv%field_data)) then
       flag = flv%field_data%has_decay_helicity (anti = flv%f < 0)
    else
       flag = .false.
    end if
  end function flavor_has_decay_helicity

  elemental function flavor_get_decay_helicity (flv) result (hel)
    integer :: hel
    class(flavor_t), intent(in) :: flv
    if (associated (flv%field_data)) then
       hel = flv%field_data%decay_helicity (anti = flv%f < 0)
    else
       hel = 0
    end if
  end function flavor_get_decay_helicity

  elemental function flavor_is_polarized (flv) result (flag)
    logical :: flag
    class(flavor_t), intent(in) :: flv
    if (associated (flv%field_data)) then
       flag = flv%field_data%is_polarized (anti = flv%f < 0)
    else
       flag = .false.
    end if
  end function flavor_is_polarized

\eatline
\eatline
\eatline
\eatline
\eatline
\eatline
\eatline
\eatline
\eatline
\eatline
\eatline
\eatline
\eatline
\eatline
\eatline
\nwendcode{}\nwbegindocs{425}Names:
\nwenddocs{}\nwbegincode{426}\moddef{Flavors: flavor: TBP}\plusendmoddef
  procedure :: get_name => flavor_get_name
  procedure :: get_tex_name => flavor_get_tex_name
\nwendcode{}\nwbegincode{427}\moddef{Flavors: procedures}\plusendmoddef
  elemental function flavor_get_name (flv) result (name)
    type(string_t) :: name
    class(flavor_t), intent(in) :: flv
    if (associated (flv%field_data)) then
       name = flv%field_data%get_name (flv%f < 0)
    else
       name = "?"
    end if
  end function flavor_get_name
       
  elemental function flavor_get_tex_name (flv) result (name)
    type(string_t) :: name
    class(flavor_t), intent(in) :: flv
    if (associated (flv%field_data)) then
       name = flv%field_data%get_tex_name (flv%f < 0)
    else
       name = "?"
    end if
  end function flavor_get_tex_name
       
\eatline
\nwendcode{}\nwbegincode{428}\moddef{Flavors: flavor: TBP}\plusendmoddef
  procedure :: get_spin_type => flavor_get_spin_type
  procedure :: get_multiplicity => flavor_get_multiplicity
  procedure :: get_isospin_type => flavor_get_isospin_type
  procedure :: get_charge_type => flavor_get_charge_type
  procedure :: get_color_type => flavor_get_color_type
\nwendcode{}\nwbegincode{429}\moddef{Flavors: procedures}\plusendmoddef
  elemental function flavor_get_spin_type (flv) result (type)
    integer :: type
    class(flavor_t), intent(in) :: flv
    if (associated (flv%field_data)) then
       type = flv%field_data%get_spin_type ()
    else
       type = 1
    end if
  end function flavor_get_spin_type

  elemental function flavor_get_multiplicity (flv) result (type)
    integer :: type
    class(flavor_t), intent(in) :: flv
    if (associated (flv%field_data)) then
       type = flv%field_data%get_multiplicity ()
    else
       type = 1
    end if
  end function flavor_get_multiplicity

  elemental function flavor_get_isospin_type (flv) result (type)
    integer :: type
    class(flavor_t), intent(in) :: flv
    if (associated (flv%field_data)) then
       type = flv%field_data%get_isospin_type ()
    else
       type = 1
    end if
  end function flavor_get_isospin_type

  elemental function flavor_get_charge_type (flv) result (type)
    integer :: type
    class(flavor_t), intent(in) :: flv
    if (associated (flv%field_data)) then
       type = flv%field_data%get_charge_type ()
    else
       type = 1
    end if
  end function flavor_get_charge_type

  elemental function flavor_get_color_type (flv) result (type)
    integer :: type
    class(flavor_t), intent(in) :: flv
    if (associated (flv%field_data)) then
       if (flavor_is_antiparticle (flv)) then
          type = - flv%field_data%get_color_type ()
       else
          type = flv%field_data%get_color_type ()
       end if
       select case (type)
       case (-1,-8);  type = abs (type)
       end select
    else
       type = 1
    end if
  end function flavor_get_color_type

\eatline
\eatline
\eatline
\eatline
\eatline
\nwendcode{}\nwbegindocs{430}These functions return real values:
\nwenddocs{}\nwbegincode{431}\moddef{Flavors: flavor: TBP}\plusendmoddef
  procedure :: get_charge => flavor_get_charge
  procedure :: get_mass => flavor_get_mass
  procedure :: get_width => flavor_get_width
  procedure :: get_isospin => flavor_get_isospin
\nwendcode{}\nwbegincode{432}\moddef{Flavors: procedures}\plusendmoddef
  elemental function flavor_get_charge (flv) result (charge)
    real(default) :: charge
    class(flavor_t), intent(in) :: flv
    if (associated (flv%field_data)) then
       if (flavor_is_antiparticle (flv)) then
          charge = flv%field_data%get_charge ()
       else
          charge = - flv%field_data%get_charge ()
       end if
    else
       charge = 0
    end if
  end function flavor_get_charge

  elemental function flavor_get_mass (flv) result (mass)
    real(default) :: mass
    class(flavor_t), intent(in) :: flv
    if (associated (flv%field_data)) then
       mass = flv%field_data%get_mass ()
    else
       mass = 0
    end if
  end function flavor_get_mass

  elemental function flavor_get_width (flv) result (width)
    real(default) :: width
    class(flavor_t), intent(in) :: flv
    if (associated (flv%field_data)) then
       width = flv%field_data%get_width ()
    else
       width = 0
    end if
  end function flavor_get_width

  elemental function flavor_get_isospin (flv) result (isospin)
    real(default) :: isospin
    class(flavor_t), intent(in) :: flv
    if (associated (flv%field_data)) then
       if (flavor_is_antiparticle (flv)) then
          isospin = flv%field_data%get_isospin ()
       else
          isospin = - flv%field_data%get_isospin ()
       end if
    else
       isospin = 0
    end if
  end function flavor_get_isospin

\eatline
\eatline
\nwendcode{}\nwbegindocs{433}\nwdocspar
\subsubsection{Comparisons}
If one of the flavors is undefined, the other defined, they match.
\nwenddocs{}\nwbegincode{434}\moddef{Flavors: flavor: TBP}\plusendmoddef
  generic :: operator(.match.) => flavor_match
  generic :: operator(==) => flavor_eq
  generic :: operator(/=) => flavor_neq
  procedure, private :: flavor_match
  procedure, private :: flavor_eq
  procedure, private :: flavor_neq
\eatline
\nwendcode{}\nwbegincode{435}\moddef{Flavors: procedures}\plusendmoddef
  elemental function flavor_match (flv1, flv2) result (eq)
    logical :: eq
    class(flavor_t), intent(in) :: flv1, flv2
    if (flv1%f /= UNDEFINED .and. flv2%f /= UNDEFINED) then
       eq = flv1%f == flv2%f
    else
       eq = .true.
    end if
  end function flavor_match

  elemental function flavor_eq (flv1, flv2) result (eq)
    logical :: eq
    class(flavor_t), intent(in) :: flv1, flv2
    if (flv1%f /= UNDEFINED .and. flv2%f /= UNDEFINED) then
       eq = flv1%f == flv2%f
    else if (flv1%f == UNDEFINED .and. flv2%f == UNDEFINED) then
       eq = .true.
    else
       eq = .false.
    end if
  end function flavor_eq

\eatline
\nwendcode{}\nwbegincode{436}\moddef{Flavors: procedures}\plusendmoddef
  elemental function flavor_neq (flv1, flv2) result (neq)
    logical :: neq
    class(flavor_t), intent(in) :: flv1, flv2
    if (flv1%f /= UNDEFINED .and. flv2%f /= UNDEFINED) then
       neq = flv1%f /= flv2%f
    else if (flv1%f == UNDEFINED .and. flv2%f == UNDEFINED) then
       neq = .false.
    else
       neq = .true.
    end if
  end function flavor_neq

\eatline
\nwendcode{}\nwbegindocs{437}\nwdocspar
\subsubsection{Tools}
Merge two flavor indices.  This works only if both are equal or either
one is undefined, because we have no off-diagonal flavor entries.
Otherwise, generate an invalid flavor.

We cannot use elemental procedures because of the pointer component.
\nwenddocs{}\nwbegincode{438}\moddef{Flavors: public}\plusendmoddef
  public :: operator(.merge.)
\nwendcode{}\nwbegincode{439}\moddef{Flavors: interfaces}\endmoddef
  interface operator(.merge.)
     module procedure merge_flavors0
     module procedure merge_flavors1
  end interface

\eatline
\nwendcode{}\nwbegincode{440}\moddef{Flavors: procedures}\plusendmoddef
  function merge_flavors0 (flv1, flv2) result (flv)
    type(flavor_t) :: flv
    type(flavor_t), intent(in) :: flv1, flv2
    if (flavor_is_defined (flv1) .and. flavor_is_defined (flv2)) then
       if (flv1 == flv2) then
          flv = flv1
       else
          flv%f = INVALID
       end if
    else if (flavor_is_defined (flv1)) then
       flv = flv1
    else if (flavor_is_defined (flv2)) then
       flv = flv2
    end if
  end function merge_flavors0

  function merge_flavors1 (flv1, flv2) result (flv)
    type(flavor_t), dimension(:), intent(in) :: flv1, flv2
    type(flavor_t), dimension(size(flv1)) :: flv
    integer :: i
    do i = 1, size (flv1)
       flv(i) = flv1(i) .merge. flv2(i)
    end do
  end function merge_flavors1

\eatline
\nwendcode{}\nwbegindocs{441}Generate consecutive color indices for a given flavor.  The indices
are counted starting with the stored value of c, so new indices are
created each time this (impure) function is called.  The counter can
be reset by the optional argument {\tt{}c{\char95}seed} if desired.  The optional
flag {\tt{}reverse} is used only for octets.  If set, the color and
anticolor entries of the octet particle are exchanged.
\nwenddocs{}\nwbegincode{442}\moddef{Flavors: public}\plusendmoddef
  public :: color_from_flavor
\nwendcode{}\nwbegincode{443}\moddef{Flavors: interfaces}\plusendmoddef
  interface color_from_flavor
     module procedure color_from_flavor0
     module procedure color_from_flavor1
  end interface
\nwendcode{}\nwbegincode{444}\moddef{Flavors: procedures}\plusendmoddef
  function color_from_flavor0 (flv, c_seed, reverse) result (col)
    type(color_t) :: col
    type(flavor_t), intent(in) :: flv
    integer, intent(in), optional :: c_seed
    logical, intent(in), optional :: reverse
    integer, save :: c = 1
    logical :: rev
    if (present (c_seed))  c = c_seed
    rev = .false.;  if (present (reverse)) rev = reverse
    select case (flavor_get_color_type (flv))
    case (1)
       call col%init ()
    case (3)
       call col%init ([c]);  c = c + 1
    case (-3)
       call col%init ([-c]);  c = c + 1
    case (8)
       if (rev) then
          call col%init ([c+1, -c]);  c = c + 2
       else
          call col%init ([c, -(c+1)]);  c = c + 2
       end if
    end select
  end function color_from_flavor0

  function color_from_flavor1 (flv, c_seed, reverse) result (col)
    type(flavor_t), dimension(:), intent(in) :: flv
    integer, intent(in), optional :: c_seed
    logical, intent(in), optional :: reverse
    type(color_t), dimension(size(flv)) :: col
    integer :: i
    col(1) = color_from_flavor0 (flv(1), c_seed, reverse)
    do i = 2, size (flv)
       col(i) = color_from_flavor0 (flv(i), reverse=reverse)
    end do
  end function color_from_flavor1

\eatline
\nwendcode{}\nwbegindocs{445}This procedure returns the flavor object for the antiparticle.  The
antiparticle code may either be the same code or its negative.
\nwenddocs{}\nwbegincode{446}\moddef{Flavors: flavor: TBP}\plusendmoddef
  procedure :: anti => flavor_anti
\nwendcode{}\nwbegincode{447}\moddef{Flavors: procedures}\plusendmoddef
  function flavor_anti (flv) result (aflv)
    type(flavor_t) :: aflv
    class(flavor_t), intent(in) :: flv
    if (flavor_has_antiparticle (flv)) then
       aflv%f = - flv%f
    else
       aflv%f = flv%f
    end if
    aflv%field_data => flv%field_data
  end function flavor_anti

\eatline
\nwendcode{}\nwbegindocs{448}\nwdocspar
\clearpage
%%%%%%%%%%%%%%%%%%%%%%%%%%%%%%%%%%%%%%%%%%%%%%%%%%%%%%%%%%%%%%%%%%%%%%%%
\section{Quantum numbers}
This module collects helicity, color, and flavor in a single type and
defines procedures
\nwenddocs{}\nwbegincode{449}\moddef{\code{}quantum{\char95}numbers.f90\edoc{}}\endmoddef
\LA{}File header\RA{}

module quantum_numbers

  use io_units
  use model_data
  use helicities
  use colors
  use flavors

\LA{}Standard module head\RA{}

\LA{}Quantum numbers: public\RA{}

\LA{}Quantum numbers: types\RA{}

\LA{}Quantum numbers: interfaces\RA{}

contains

\LA{}Quantum numbers: procedures\RA{}

end module quantum_numbers
\eatline
\nwendcode{}\nwbegindocs{450}\nwdocspar
\subsection{The quantum number type}
\nwenddocs{}\nwbegincode{451}\moddef{Quantum numbers: public}\endmoddef
  public :: quantum_numbers_t
\nwendcode{}\nwbegincode{452}\moddef{Quantum numbers: types}\endmoddef
  type :: quantum_numbers_t
     private
     type(flavor_t) :: f
     type(color_t) :: c
     type(helicity_t) :: h
   contains
   \LA{}Quantum numbers: quantum numbers: TBP\RA{}
  end type quantum_numbers_t

\eatline
\nwendcode{}\nwbegindocs{453}Define quantum numbers: Initializer form.  All arguments may be
present or absent.

Some elemental initializers are impure because they set the {\tt{}flv}
component.  This implies transfer of a pointer behind the scenes.
\nwenddocs{}\nwbegincode{454}\moddef{Quantum numbers: quantum numbers: TBP}\endmoddef
  generic :: init => &
     quantum_numbers_init0_f, &
     quantum_numbers_init0_c, &
     quantum_numbers_init0_h, &
     quantum_numbers_init0_fc, &
     quantum_numbers_init0_fh, &
     quantum_numbers_init0_ch, &
     quantum_numbers_init0_fch
  procedure, private :: quantum_numbers_init0_f
  procedure, private :: quantum_numbers_init0_c
  procedure, private :: quantum_numbers_init0_h
  procedure, private :: quantum_numbers_init0_fc
  procedure, private :: quantum_numbers_init0_fh
  procedure, private :: quantum_numbers_init0_ch
  procedure, private :: quantum_numbers_init0_fch
\nwendcode{}\nwbegincode{455}\moddef{Quantum numbers: procedures}\endmoddef
  impure elemental subroutine quantum_numbers_init0_f (qn, flv)
    class(quantum_numbers_t), intent(inout) :: qn
    type(flavor_t), intent(in) :: flv
    qn%f = flv
    call qn%c%undefine ()
    call qn%h%undefine ()
  end subroutine quantum_numbers_init0_f

  elemental subroutine quantum_numbers_init0_c (qn, col)
    class(quantum_numbers_t), intent(inout) :: qn
    type(color_t), intent(in) :: col
    call qn%f%undefine ()
    qn%c = col
    call qn%h%undefine ()
  end subroutine quantum_numbers_init0_c

  elemental subroutine quantum_numbers_init0_h (qn, hel)
    class(quantum_numbers_t), intent(inout) :: qn
    type(helicity_t), intent(in) :: hel
    call qn%f%undefine ()
    call qn%c%undefine ()
    qn%h = hel
  end subroutine quantum_numbers_init0_h

  impure elemental subroutine quantum_numbers_init0_fc (qn, flv, col)
    class(quantum_numbers_t), intent(inout) :: qn
    type(flavor_t), intent(in) :: flv
    type(color_t), intent(in) :: col
    qn%f = flv
    qn%c = col
    call qn%h%undefine ()
  end subroutine quantum_numbers_init0_fc

  impure elemental subroutine quantum_numbers_init0_fh (qn, flv, hel)
    class(quantum_numbers_t), intent(inout) :: qn
    type(flavor_t), intent(in) :: flv
    type(helicity_t), intent(in) :: hel
    qn%f = flv
    call qn%c%undefine ()
    qn%h = hel
  end subroutine quantum_numbers_init0_fh

  elemental subroutine quantum_numbers_init0_ch (qn, col, hel)
    class(quantum_numbers_t), intent(inout) :: qn
    type(color_t), intent(in) :: col
    type(helicity_t), intent(in) :: hel
    call qn%f%undefine ()
    qn%c = col
    qn%h = hel
  end subroutine quantum_numbers_init0_ch

  impure elemental subroutine quantum_numbers_init0_fch (qn, flv, col, hel)
    class(quantum_numbers_t), intent(inout) :: qn
    type(flavor_t), intent(in) :: flv
    type(color_t), intent(in) :: col
    type(helicity_t), intent(in) :: hel
    qn%f = flv
    qn%c = col
    qn%h = hel
  end subroutine quantum_numbers_init0_fch

\eatline
\nwendcode{}\nwbegindocs{456}\nwdocspar
\subsection{I/O}
Write the quantum numbers in condensed form, enclosed by square
brackets.  Color is written only if nontrivial.  For convenience,
introduce also an array version.

If the {\tt{}col{\char95}verbose} option is set, show the quantum number color also
if it is zero, but defined.  Otherwise, suppress zero color.
\nwenddocs{}\nwbegincode{457}\moddef{Quantum numbers: public}\plusendmoddef
  public :: quantum_numbers_write
\nwendcode{}\nwbegincode{458}\moddef{Quantum numbers: quantum numbers: TBP}\plusendmoddef
  procedure :: write => quantum_numbers_write_single
\nwendcode{}\nwbegincode{459}\moddef{Quantum numbers: interfaces}\endmoddef
  interface quantum_numbers_write
     module procedure quantum_numbers_write_single
     module procedure quantum_numbers_write_array
  end interface
\nwendcode{}\nwbegincode{460}\moddef{Quantum numbers: procedures}\plusendmoddef
  subroutine quantum_numbers_write_single (qn, unit, col_verbose)
    class(quantum_numbers_t), intent(in) :: qn
    integer, intent(in), optional :: unit
    logical, intent(in), optional :: col_verbose
    integer :: u
    logical :: col_verb
    u = given_output_unit (unit);  if (u < 0)  return
    col_verb = .false.;  if (present (col_verbose)) col_verb = col_verbose
    write (u, "(A)", advance="no")  "["
    if (qn%f%is_defined ()) then
       call qn%f%write (u)
       if (qn%c%is_nonzero () .or. qn%h%is_defined ()) &
            write (u, "(1x)", advance="no")
    end if
    if (col_verb) then
       if (qn%c%is_defined () .or. qn%c%is_ghost ()) then
          call color_write (qn%c, u)
          if (qn%h%is_defined ())  write (u, "(1x)", advance="no")
       end if
    else
       if (qn%c%is_nonzero () .or. qn%c%is_ghost ()) then
          call color_write (qn%c, u)
          if (qn%h%is_defined ())  write (u, "(1x)", advance="no")
       end if
    end if
    if (qn%h%is_defined ()) then
       call qn%h%write (u)
    end if
    write (u, "(A)", advance="no")  "]"
  end subroutine quantum_numbers_write_single

  subroutine quantum_numbers_write_array (qn, unit, col_verbose)
    type(quantum_numbers_t), dimension(:), intent(in) :: qn
    integer, intent(in), optional :: unit
    logical, intent(in), optional :: col_verbose
    integer :: i
    integer :: u
    logical :: col_verb
    u = given_output_unit (unit);  if (u < 0)  return
    col_verb = .false.;  if (present (col_verbose)) col_verb = col_verbose
    write (u, "(A)", advance="no")  "["
    do i = 1, size (qn)
       if (i > 1)  write (u, "(A)", advance="no")  " / "
       if (qn(i)%f%is_defined ()) then
          call qn(i)%f%write (u)
          if (qn(i)%c%is_nonzero () .or. qn(i)%h%is_defined ()) &
               write (u, "(1x)", advance="no")
       end if
       if (col_verb) then
          if (qn(i)%c%is_defined () .or. qn(i)%c%is_ghost ()) then
             call color_write (qn(i)%c, u)
             if (qn(i)%h%is_defined ())  write (u, "(1x)", advance="no")
          end if
       else
          if (qn(i)%c%is_nonzero () .or. qn(i)%c%is_ghost ()) then
             call color_write (qn(i)%c, u)
             if (qn(i)%h%is_defined ())  write (u, "(1x)", advance="no")
          end if
       end if
       if (qn(i)%h%is_defined ()) then
          call qn(i)%h%write (u)
       end if
    end do
    write (u, "(A)", advance="no")  "]"
  end subroutine quantum_numbers_write_array

\eatline
\nwendcode{}\nwbegindocs{461}Binary I/O.
\nwenddocs{}\nwbegincode{462}\moddef{Quantum numbers: quantum numbers: TBP}\plusendmoddef
  procedure :: write_raw => quantum_numbers_write_raw
  procedure :: read_raw => quantum_numbers_read_raw
\nwendcode{}\nwbegincode{463}\moddef{Quantum numbers: procedures}\plusendmoddef
  subroutine quantum_numbers_write_raw (qn, u)
    class(quantum_numbers_t), intent(in) :: qn
    integer, intent(in) :: u
    call qn%f%write_raw (u)
    call qn%c%write_raw (u)
    call qn%h%write_raw (u)
  end subroutine quantum_numbers_write_raw

  subroutine quantum_numbers_read_raw (qn, u, iostat)
    class(quantum_numbers_t), intent(out) :: qn
    integer, intent(in) :: u
    integer, intent(out), optional :: iostat
    call qn%f%read_raw (u, iostat=iostat)
    call qn%c%read_raw (u, iostat=iostat)
    call qn%h%read_raw (u, iostat=iostat)
  end subroutine quantum_numbers_read_raw

\eatline
\nwendcode{}\nwbegindocs{464}\nwdocspar
\subsection{Accessing contents}
Color and helicity can be done by elemental functions.  Flavor needs
impure elemental.  We export also the functions directly, this allows
us to avoid temporaries in some places.
\nwenddocs{}\nwbegincode{465}\moddef{Quantum numbers: public}\plusendmoddef
  public :: quantum_numbers_get_flavor
  public :: quantum_numbers_get_color
  public :: quantum_numbers_get_helicity
\nwendcode{}\nwbegincode{466}\moddef{Quantum numbers: quantum numbers: TBP}\plusendmoddef
  procedure :: get_flavor => quantum_numbers_get_flavor
  procedure :: get_color => quantum_numbers_get_color
  procedure :: get_helicity => quantum_numbers_get_helicity
\nwendcode{}\nwbegincode{467}\moddef{Quantum numbers: procedures}\plusendmoddef
  impure elemental function quantum_numbers_get_flavor (qn) result (flv)
    type(flavor_t) :: flv
    class(quantum_numbers_t), intent(in) :: qn
    flv = qn%f
  end function quantum_numbers_get_flavor

  elemental function quantum_numbers_get_color (qn) result (col)
    type(color_t) :: col
    class(quantum_numbers_t), intent(in) :: qn
    col = qn%c
  end function quantum_numbers_get_color

  elemental function quantum_numbers_get_helicity (qn) result (hel)
    type(helicity_t) :: hel
    class(quantum_numbers_t), intent(in) :: qn
    hel = qn%h
  end function quantum_numbers_get_helicity

\eatline
\eatline
\eatline
\nwendcode{}\nwbegindocs{468}This just resets the ghost property of the color part:
\nwenddocs{}\nwbegincode{469}\moddef{Quantum numbers: quantum numbers: TBP}\plusendmoddef
  procedure :: set_color_ghost => quantum_numbers_set_color_ghost
\nwendcode{}\nwbegincode{470}\moddef{Quantum numbers: procedures}\plusendmoddef
  elemental subroutine quantum_numbers_set_color_ghost (qn, ghost)
    class(quantum_numbers_t), intent(inout) :: qn
    logical, intent(in) :: ghost
    call qn%c%set_ghost (ghost)
  end subroutine quantum_numbers_set_color_ghost

\eatline
\nwendcode{}\nwbegindocs{471}Assign a model to the flavor part of quantum numbers.
\nwenddocs{}\nwbegincode{472}\moddef{Quantum numbers: quantum numbers: TBP}\plusendmoddef
  procedure :: set_model => quantum_numbers_set_model
\nwendcode{}\nwbegincode{473}\moddef{Quantum numbers: procedures}\plusendmoddef
  impure elemental subroutine quantum_numbers_set_model (qn, model)
    class(quantum_numbers_t), intent(inout) :: qn
    class(model_data_t), intent(in), target :: model
    call qn%f%set_model (model)
  end subroutine quantum_numbers_set_model

\eatline
\nwendcode{}\nwbegindocs{474}Set the {\tt{}radiated} flag for the flavor component.
\nwenddocs{}\nwbegincode{475}\moddef{Quantum numbers: quantum numbers: TBP}\plusendmoddef
  procedure :: tag_radiated => quantum_numbers_tag_radiated
\nwendcode{}\nwbegincode{476}\moddef{Quantum numbers: procedures}\plusendmoddef
  elemental subroutine quantum_numbers_tag_radiated (qn)
    class(quantum_numbers_t), intent(inout) :: qn
    call qn%f%tag_radiated ()
  end subroutine quantum_numbers_tag_radiated
  
\eatline
\nwendcode{}\nwbegindocs{477}This is a convenience function: return the color type for the flavor
(array).

Note: keep the public version temporarily, this will be used in a
complicated expression which triggers a compiler bug (nagfor 5.3) in
the TBP version.
\nwenddocs{}\nwbegincode{478}\moddef{Quantum numbers: public}\plusendmoddef
  public :: quantum_numbers_get_color_type
\nwendcode{}\nwbegincode{479}\moddef{Quantum numbers: quantum numbers: TBP}\plusendmoddef
  procedure :: get_color_type => quantum_numbers_get_color_type
\nwendcode{}\nwbegincode{480}\moddef{Quantum numbers: procedures}\plusendmoddef
  elemental function quantum_numbers_get_color_type (qn) result (color_type)
    integer :: color_type
    class(quantum_numbers_t), intent(in) :: qn
    color_type = qn%f%get_color_type ()
  end function quantum_numbers_get_color_type

\eatline
\nwendcode{}\nwbegindocs{481}\nwdocspar
\subsection{Predicates}
Check if the flavor index is valid (including UNDEFINED).
\nwenddocs{}\nwbegincode{482}\moddef{Quantum numbers: quantum numbers: TBP}\plusendmoddef
  procedure :: are_valid => quantum_numbers_are_valid
\nwendcode{}\nwbegincode{483}\moddef{Quantum numbers: procedures}\plusendmoddef
  elemental function quantum_numbers_are_valid (qn) result (valid)
    logical :: valid
    class(quantum_numbers_t), intent(in) :: qn
    valid = qn%f%is_valid ()
  end function quantum_numbers_are_valid

\eatline
\nwendcode{}\nwbegindocs{484}Check if the flavor part has its particle-data pointer associated
(debugging aid).
\nwenddocs{}\nwbegincode{485}\moddef{Quantum numbers: quantum numbers: TBP}\plusendmoddef
  procedure :: are_associated => quantum_numbers_are_associated
\nwendcode{}\nwbegincode{486}\moddef{Quantum numbers: procedures}\plusendmoddef
  elemental function quantum_numbers_are_associated (qn) result (flag)
    logical :: flag
    class(quantum_numbers_t), intent(in) :: qn
    flag = qn%f%is_associated ()
  end function quantum_numbers_are_associated

\eatline
\nwendcode{}\nwbegindocs{487}Check if the helicity and color quantum numbers are
diagonal. (Unpolarized/colorless also counts as diagonal.)  Flavor is
diagonal by definition.
\nwenddocs{}\nwbegincode{488}\moddef{Quantum numbers: quantum numbers: TBP}\plusendmoddef
  procedure :: are_diagonal => quantum_numbers_are_diagonal
\nwendcode{}\nwbegincode{489}\moddef{Quantum numbers: procedures}\plusendmoddef
  elemental function quantum_numbers_are_diagonal (qn) result (diagonal)
    logical :: diagonal
    class(quantum_numbers_t), intent(in) :: qn
    diagonal = qn%h%is_diagonal () .and. qn%c%is_diagonal ()
  end function quantum_numbers_are_diagonal

\eatline
\nwendcode{}\nwbegindocs{490}Check if the color part has the ghost property.
\nwenddocs{}\nwbegincode{491}\moddef{Quantum numbers: quantum numbers: TBP}\plusendmoddef
  procedure :: is_color_ghost => quantum_numbers_is_color_ghost
\nwendcode{}\nwbegincode{492}\moddef{Quantum numbers: procedures}\plusendmoddef
  elemental function quantum_numbers_is_color_ghost (qn) result (ghost)
    logical :: ghost
    class(quantum_numbers_t), intent(in) :: qn
    ghost = qn%c%is_ghost ()
  end function quantum_numbers_is_color_ghost

\eatline
\nwendcode{}\nwbegindocs{493}\nwdocspar
\subsection{Comparisons}
Matching and equality is derived from the individual quantum numbers.
The variant {\tt{}fhmatch} matches only flavor and helicity.  The variant
{\tt{}dhmatch} matches only diagonal helicity, if the matching helicity is
undefined.
\nwenddocs{}\nwbegincode{494}\moddef{Quantum numbers: quantum numbers: TBP}\plusendmoddef
  generic :: operator(.match.) => quantum_numbers_match
  generic :: operator(.fmatch.) => quantum_numbers_match_f
  generic :: operator(.fhmatch.) => quantum_numbers_match_fh
  generic :: operator(.dhmatch.) => quantum_numbers_match_hel_diag
  generic :: operator(==) => quantum_numbers_eq
  generic :: operator(/=) => quantum_numbers_neq
  procedure, private :: quantum_numbers_match
  procedure, private :: quantum_numbers_match_f
  procedure, private :: quantum_numbers_match_fh
  procedure, private :: quantum_numbers_match_hel_diag
  procedure, private :: quantum_numbers_eq
  procedure, private :: quantum_numbers_neq
\eatline
\nwendcode{}\nwbegincode{495}\moddef{Quantum numbers: procedures}\plusendmoddef
  elemental function quantum_numbers_match (qn1, qn2) result (match)
    logical :: match
    class(quantum_numbers_t), intent(in) :: qn1, qn2
    match = (qn1%f .match. qn2%f) .and. &
         (qn1%c .match. qn2%c) .and. &
         (qn1%h .match. qn2%h)
  end function quantum_numbers_match

  elemental function quantum_numbers_match_f (qn1, qn2) result (match)
    logical :: match
    class(quantum_numbers_t), intent(in) :: qn1, qn2
    match = (qn1%f .match. qn2%f)
  end function quantum_numbers_match_f

  elemental function quantum_numbers_match_fh (qn1, qn2) result (match)
    logical :: match
    class(quantum_numbers_t), intent(in) :: qn1, qn2
    match = (qn1%f .match. qn2%f) .and. &
         (qn1%h .match. qn2%h)
  end function quantum_numbers_match_fh

  elemental function quantum_numbers_match_hel_diag (qn1, qn2) result (match)
    logical :: match
    class(quantum_numbers_t), intent(in) :: qn1, qn2
    match = (qn1%f .match. qn2%f) .and. &
         (qn1%c .match. qn2%c) .and. &
         (qn1%h .dmatch. qn2%h)
  end function quantum_numbers_match_hel_diag

  elemental function quantum_numbers_eq (qn1, qn2) result (eq)
    logical :: eq
    class(quantum_numbers_t), intent(in) :: qn1, qn2
    eq = (qn1%f == qn2%f) .and. &
         (qn1%c == qn2%c) .and. &
         (qn1%h == qn2%h)
  end function quantum_numbers_eq

  elemental function quantum_numbers_neq (qn1, qn2) result (neq)
    logical :: neq
    class(quantum_numbers_t), intent(in) :: qn1, qn2
    neq = (qn1%f /= qn2%f) .or. &
         (qn1%c /= qn2%c) .or. &
         (qn1%h /= qn2%h)
  end function quantum_numbers_neq

\eatline
\eatline
\eatline
\nwendcode{}\nwbegindocs{496}Two sets of quantum numbers are compatible if the individual quantum numbers
are compatible, depending on the mask.  Flavor has to match, regardless of the
flavor mask.  

If the color flag is set, color is compatible if the ghost property is
identical.  If the color flag is unset, color has to be identical.  I.e., if
the flag is set, the color amplitudes can interfere.  If it is not set, they
must be identical, and there must be no ghost.  The latter property is used
for expanding physical color flows.

Helicity is compatible if the mask is unset, otherwise it has to match.  This
determines if two amplitudes can be multiplied (no mask) or traced (mask).
\nwenddocs{}\nwbegincode{497}\moddef{Quantum numbers: public}\plusendmoddef
  public :: quantum_numbers_are_compatible
\nwendcode{}\nwbegincode{498}\moddef{Quantum numbers: procedures}\plusendmoddef
  elemental function quantum_numbers_are_compatible (qn1, qn2, mask) &
      result (flag)
    logical :: flag
    type(quantum_numbers_t), intent(in) :: qn1, qn2
    type(quantum_numbers_mask_t), intent(in) :: mask
    if (mask%h .or. mask%hd) then
       flag = (qn1%f .match. qn2%f) .and. (qn1%h .match. qn2%h)
    else
       flag = (qn1%f .match. qn2%f)
    end if
    if (mask%c) then
       flag = flag .and. (qn1%c%is_ghost () .eqv. qn2%c%is_ghost ())
    else
       flag = flag .and. &
            .not. (qn1%c%is_ghost () .or. qn2%c%is_ghost ()) .and. &
            (qn1%c == qn2%c)
    end if
  end function quantum_numbers_are_compatible

\eatline
\nwendcode{}\nwbegindocs{499}This is the analog for a single quantum-number set.  We just check for color
ghosts; they are excluded if the color mask is unset (color-flow expansion).
\nwenddocs{}\nwbegincode{500}\moddef{Quantum numbers: public}\plusendmoddef
  public :: quantum_numbers_are_physical
\nwendcode{}\nwbegincode{501}\moddef{Quantum numbers: procedures}\plusendmoddef
  elemental function quantum_numbers_are_physical (qn, mask) result (flag)
    logical :: flag
    type(quantum_numbers_t), intent(in) :: qn
    type(quantum_numbers_mask_t), intent(in) :: mask
    if (mask%c) then
       flag = .true.
    else
       flag = .not. qn%c%is_ghost ()
    end if
  end function quantum_numbers_are_physical

\eatline
\nwendcode{}\nwbegindocs{502}\nwdocspar
\subsection{Operations}
Inherited from the color component: reassign color indices in
canonical order.
\nwenddocs{}\nwbegincode{503}\moddef{Quantum numbers: public}\plusendmoddef
  public :: quantum_numbers_canonicalize_color
\nwendcode{}\nwbegincode{504}\moddef{Quantum numbers: procedures}\plusendmoddef
  subroutine quantum_numbers_canonicalize_color (qn)
    type(quantum_numbers_t), dimension(:), intent(inout) :: qn
    call color_canonicalize (qn%c)
  end subroutine quantum_numbers_canonicalize_color

\eatline
\nwendcode{}\nwbegindocs{505}Inherited from the color component: make a color map for two matching
quantum-number arrays.
\nwenddocs{}\nwbegincode{506}\moddef{Quantum numbers: public}\plusendmoddef
  public :: make_color_map
\nwendcode{}\nwbegincode{507}\moddef{Quantum numbers: interfaces}\plusendmoddef
  interface make_color_map
     module procedure quantum_numbers_make_color_map
  end interface make_color_map
  
\nwendcode{}\nwbegincode{508}\moddef{Quantum numbers: procedures}\plusendmoddef
  subroutine quantum_numbers_make_color_map (map, qn1, qn2)
    integer, dimension(:,:), intent(out), allocatable :: map
    type(quantum_numbers_t), dimension(:), intent(in) :: qn1, qn2
    call make_color_map (map, qn1%c, qn2%c)
  end subroutine quantum_numbers_make_color_map

\eatline
\nwendcode{}\nwbegindocs{509}Inherited from the color component: translate the color part using a
color-map array
\nwenddocs{}\nwbegincode{510}\moddef{Quantum numbers: public}\plusendmoddef
  public :: quantum_numbers_translate_color
\nwendcode{}\nwbegincode{511}\moddef{Quantum numbers: interfaces}\plusendmoddef
  interface quantum_numbers_translate_color
     module procedure quantum_numbers_translate_color0
     module procedure quantum_numbers_translate_color1
  end interface

\nwendcode{}\nwbegincode{512}\moddef{Quantum numbers: procedures}\plusendmoddef
  subroutine quantum_numbers_translate_color0 (qn, map, offset)
    type(quantum_numbers_t), intent(inout) :: qn
    integer, dimension(:,:), intent(in) :: map
    integer, intent(in), optional :: offset
    call color_translate (qn%c, map, offset)
  end subroutine quantum_numbers_translate_color0
  
  subroutine quantum_numbers_translate_color1 (qn, map, offset)
    type(quantum_numbers_t), dimension(:), intent(inout) :: qn
    integer, dimension(:,:), intent(in) :: map
    integer, intent(in), optional :: offset
    call color_translate (qn%c, map, offset)
  end subroutine quantum_numbers_translate_color1

\eatline
\nwendcode{}\nwbegindocs{513}Inherited from the color component: return the color index with
highest absolute value.

Since the algorithm is not elemental, we keep the separate
procedures for different array rank.
\nwenddocs{}\nwbegincode{514}\moddef{Quantum numbers: public}\plusendmoddef
  public :: quantum_numbers_get_max_color_value
\nwendcode{}\nwbegincode{515}\moddef{Quantum numbers: interfaces}\plusendmoddef
  interface quantum_numbers_get_max_color_value
     module procedure quantum_numbers_get_max_color_value0
     module procedure quantum_numbers_get_max_color_value1
     module procedure quantum_numbers_get_max_color_value2
  end interface

\nwendcode{}\nwbegincode{516}\moddef{Quantum numbers: procedures}\plusendmoddef
  pure function quantum_numbers_get_max_color_value0 (qn) result (cmax)
    integer :: cmax
    type(quantum_numbers_t), intent(in) :: qn
    cmax = color_get_max_value (qn%c)
  end function quantum_numbers_get_max_color_value0
    
  pure function quantum_numbers_get_max_color_value1 (qn) result (cmax)
    integer :: cmax
    type(quantum_numbers_t), dimension(:), intent(in) :: qn
    cmax = color_get_max_value (qn%c)
  end function quantum_numbers_get_max_color_value1
    
  pure function quantum_numbers_get_max_color_value2 (qn) result (cmax)
    integer :: cmax
    type(quantum_numbers_t), dimension(:,:), intent(in) :: qn
    cmax = color_get_max_value (qn%c)
  end function quantum_numbers_get_max_color_value2
    
\nwendcode{}\nwbegindocs{517}Inherited from the color component: add an offset to the indices of
the color part
\nwenddocs{}\nwbegincode{518}\moddef{Quantum numbers: quantum numbers: TBP}\plusendmoddef
  procedure :: add_color_offset => quantum_numbers_add_color_offset
\nwendcode{}\nwbegincode{519}\moddef{Quantum numbers: procedures}\plusendmoddef
  elemental subroutine quantum_numbers_add_color_offset (qn, offset)
    class(quantum_numbers_t), intent(inout) :: qn
    integer, intent(in) :: offset
    call qn%c%add_offset (offset)
  end subroutine quantum_numbers_add_color_offset

\eatline
\nwendcode{}\nwbegindocs{520}Given a quantum number array, return all possible color
contractions, leaving the other quantum numbers intact.
\nwenddocs{}\nwbegincode{521}\moddef{Quantum numbers: public}\plusendmoddef
  public :: quantum_number_array_make_color_contractions
\nwendcode{}\nwbegincode{522}\moddef{Quantum numbers: procedures}\plusendmoddef
  subroutine quantum_number_array_make_color_contractions (qn_in, qn_out)
    type(quantum_numbers_t), dimension(:), intent(in) :: qn_in
    type(quantum_numbers_t), dimension(:,:), intent(out), allocatable :: qn_out
    type(color_t), dimension(:,:), allocatable :: col
    integer :: i
    call color_array_make_contractions (qn_in%c, col)
    allocate (qn_out (size (col, 1), size (col, 2)))
    do i = 1, size (qn_out, 2)
       qn_out(:,i)%f = qn_in%f
       qn_out(:,i)%c = col(:,i)
       qn_out(:,i)%h = qn_in%h
    end do
  end subroutine quantum_number_array_make_color_contractions

\eatline
\nwendcode{}\nwbegindocs{523}Inherited from the color component: invert the color, switching
particle/antiparticle.
\nwenddocs{}\nwbegincode{524}\moddef{Quantum numbers: quantum numbers: TBP}\plusendmoddef
  procedure :: invert_color => quantum_numbers_invert_color
\nwendcode{}\nwbegincode{525}\moddef{Quantum numbers: procedures}\plusendmoddef
  elemental subroutine quantum_numbers_invert_color (qn)
    class(quantum_numbers_t), intent(inout) :: qn
    call qn%c%invert ()
  end subroutine quantum_numbers_invert_color

\eatline
\nwendcode{}\nwbegindocs{526}\nwdocspar
Merge two quantum number sets: for each entry, if both are defined,
combine them to an off-diagonal entry (meaningful only if the input
was diagonal).  If either entry is undefined, take the defined
one.

For flavor, off-diagonal entries are invalid, so both
flavors must be equal, otherwise an invalid flavor is inserted.
\nwenddocs{}\nwbegincode{527}\moddef{Quantum numbers: public}\plusendmoddef
  public :: operator(.merge.)
\nwendcode{}\nwbegincode{528}\moddef{Quantum numbers: interfaces}\plusendmoddef
  interface operator(.merge.)
     module procedure merge_quantum_numbers0
     module procedure merge_quantum_numbers1
  end interface

\nwendcode{}\nwbegincode{529}\moddef{Quantum numbers: procedures}\plusendmoddef
  function merge_quantum_numbers0 (qn1, qn2) result (qn3)
    type(quantum_numbers_t) :: qn3
    type(quantum_numbers_t), intent(in) :: qn1, qn2
    qn3%f = qn1%f .merge. qn2%f
    qn3%c = qn1%c .merge. qn2%c
    qn3%h = qn1%h .merge. qn2%h
  end function merge_quantum_numbers0

  function merge_quantum_numbers1 (qn1, qn2) result (qn3)
    type(quantum_numbers_t), dimension(:), intent(in) :: qn1, qn2
    type(quantum_numbers_t), dimension(size(qn1)) :: qn3
    qn3%f = qn1%f .merge. qn2%f
    qn3%c = qn1%c .merge. qn2%c
    qn3%h = qn1%h .merge. qn2%h
  end function merge_quantum_numbers1

\eatline
\nwendcode{}\nwbegindocs{530}\nwdocspar
\subsection{The quantum number mask}
The quantum numbers mask is true for quantum numbers that should be
ignored or summed over.  The three mandatory entries correspond to
flavor, color, and helicity, respectively.  

There is an additional entry {\tt{}cg}: If false, the color-ghosts
property should be kept even if color is ignored.  This is relevant
only if {\tt{}c} is set, otherwise it is always false.

The flag {\tt{}hd} tells that only diagonal entries in helicity should be
kept.  If {\tt{}h} is set, {\tt{}hd} is irrelevant and will be kept
{\tt{}.false.}
\nwenddocs{}\nwbegincode{531}\moddef{Quantum numbers: public}\plusendmoddef
  public :: quantum_numbers_mask_t
\nwendcode{}\nwbegincode{532}\moddef{Quantum numbers: types}\plusendmoddef
  type :: quantum_numbers_mask_t
     private
     logical :: f = .false.
     logical :: c = .false.
     logical :: cg = .false.
     logical :: h = .false.
     logical :: hd = .false.
   contains
   \LA{}Quantum numbers: quantum numbers mask: TBP\RA{}
  end type quantum_numbers_mask_t

\eatline
\nwendcode{}\nwbegindocs{533}Define a quantum number mask: Constructor form
\nwenddocs{}\nwbegincode{534}\moddef{Quantum numbers: public}\plusendmoddef
  public :: quantum_numbers_mask
\nwendcode{}\nwbegincode{535}\moddef{Quantum numbers: procedures}\plusendmoddef
  elemental function quantum_numbers_mask &
       (mask_f, mask_c, mask_h, mask_cg, mask_hd) result (mask)
    type(quantum_numbers_mask_t) :: mask
    logical, intent(in) :: mask_f, mask_c, mask_h
    logical, intent(in), optional :: mask_cg
    logical, intent(in), optional :: mask_hd
    call quantum_numbers_mask_init &
         (mask, mask_f, mask_c, mask_h, mask_cg, mask_hd)
  end function quantum_numbers_mask

\eatline
\nwendcode{}\nwbegindocs{536}Define quantum numbers: Initializer form
\nwenddocs{}\nwbegincode{537}\moddef{Quantum numbers: quantum numbers mask: TBP}\endmoddef
  procedure :: init => quantum_numbers_mask_init
\nwendcode{}\nwbegincode{538}\moddef{Quantum numbers: procedures}\plusendmoddef
  elemental subroutine quantum_numbers_mask_init &
       (mask, mask_f, mask_c, mask_h, mask_cg, mask_hd)
    class(quantum_numbers_mask_t), intent(inout) :: mask
    logical, intent(in) :: mask_f, mask_c, mask_h
    logical, intent(in), optional :: mask_cg, mask_hd
    mask%f = mask_f
    mask%c = mask_c
    mask%h = mask_h
    mask%cg = .false.
    if (present (mask_cg)) then
       if (mask%c)  mask%cg = mask_cg
    else
       mask%cg = mask_c
    end if
    mask%hd = .false.
    if (present (mask_hd)) then
       if (.not. mask%h)  mask%hd = mask_hd
    end if
  end subroutine quantum_numbers_mask_init

\eatline
\nwendcode{}\nwbegindocs{539}Write a quantum numbers mask.  We need the stand-alone subroutine for the
array case.
\nwenddocs{}\nwbegincode{540}\moddef{Quantum numbers: public}\plusendmoddef
  public :: quantum_numbers_mask_write
\nwendcode{}\nwbegincode{541}\moddef{Quantum numbers: interfaces}\plusendmoddef
  interface quantum_numbers_mask_write
     module procedure quantum_numbers_mask_write_single
     module procedure quantum_numbers_mask_write_array
  end interface
\nwendcode{}\nwbegincode{542}\moddef{Quantum numbers: quantum numbers mask: TBP}\plusendmoddef
  procedure :: write => quantum_numbers_mask_write_single
\nwendcode{}\nwbegincode{543}\moddef{Quantum numbers: procedures}\plusendmoddef
  subroutine quantum_numbers_mask_write_single (mask, unit)
    class(quantum_numbers_mask_t), intent(in) :: mask
    integer, intent(in), optional :: unit
    integer :: u
    u = given_output_unit (unit);  if (u < 0)  return
    write (u, "(A)", advance="no") "["
    write (u, "(L1)", advance="no")  mask%f
    write (u, "(L1)", advance="no")  mask%c
    if (.not.mask%cg)  write (u, "('g')", advance="no")
    write (u, "(L1)", advance="no")  mask%h
    if (mask%hd)  write (u, "('d')", advance="no")
    write (u, "(A)", advance="no") "]"
  end subroutine quantum_numbers_mask_write_single

  subroutine quantum_numbers_mask_write_array (mask, unit)
    type(quantum_numbers_mask_t), dimension(:), intent(in) :: mask
    integer, intent(in), optional :: unit
    integer :: u, i
    u = given_output_unit (unit);  if (u < 0)  return
    write (u, "(A)", advance="no") "["
    do i = 1, size (mask)
       if (i > 1)  write (u, "(A)", advance="no")  "/"
       write (u, "(L1)", advance="no")  mask(i)%f
       write (u, "(L1)", advance="no")  mask(i)%c
       if (.not.mask(i)%cg)  write (u, "('g')", advance="no")
       write (u, "(L1)", advance="no")  mask(i)%h
       if (mask(i)%hd)  write (u, "('d')", advance="no")
    end do
    write (u, "(A)", advance="no") "]"
  end subroutine quantum_numbers_mask_write_array

\eatline
\nwendcode{}\nwbegindocs{544}\nwdocspar
\subsection{Setting mask components}
\nwenddocs{}\nwbegincode{545}\moddef{Quantum numbers: quantum numbers mask: TBP}\plusendmoddef
  procedure :: set_flavor => quantum_numbers_mask_set_flavor
  procedure :: set_color => quantum_numbers_mask_set_color
  procedure :: set_helicity => quantum_numbers_mask_set_helicity
\nwendcode{}\nwbegincode{546}\moddef{Quantum numbers: procedures}\plusendmoddef
  elemental subroutine quantum_numbers_mask_set_flavor (mask, mask_f)
    class(quantum_numbers_mask_t), intent(inout) :: mask
    logical, intent(in) :: mask_f
    mask%f = mask_f
  end subroutine quantum_numbers_mask_set_flavor

  elemental subroutine quantum_numbers_mask_set_color (mask, mask_c, mask_cg)
    class(quantum_numbers_mask_t), intent(inout) :: mask
    logical, intent(in) :: mask_c
    logical, intent(in), optional :: mask_cg
    mask%c = mask_c
    if (present (mask_cg)) then
       if (mask%c)  mask%cg = mask_cg
    else
       mask%cg = mask_c
    end if
  end subroutine quantum_numbers_mask_set_color

  elemental subroutine quantum_numbers_mask_set_helicity (mask, mask_h, mask_hd)
    class(quantum_numbers_mask_t), intent(inout) :: mask
    logical, intent(in) :: mask_h
    logical, intent(in), optional :: mask_hd
    mask%h = mask_h
    if (present (mask_hd)) then
       if (.not. mask%h)  mask%hd = mask_hd
    end if
  end subroutine quantum_numbers_mask_set_helicity

\eatline
\eatline
\eatline
\nwendcode{}\nwbegindocs{547}The following routines assign part of a mask, depending on the flags given.
\nwenddocs{}\nwbegincode{548}\moddef{Quantum numbers: quantum numbers mask: TBP}\plusendmoddef
  procedure :: assign => quantum_numbers_mask_assign
\nwendcode{}\nwbegincode{549}\moddef{Quantum numbers: procedures}\plusendmoddef
  elemental subroutine quantum_numbers_mask_assign &
       (mask, mask_in, flavor, color, helicity)
    class(quantum_numbers_mask_t), intent(inout) :: mask
    class(quantum_numbers_mask_t), intent(in) :: mask_in
    logical, intent(in), optional :: flavor, color, helicity
    if (present (flavor)) then
       if (flavor) then
          mask%f = mask_in%f
       end if
    end if
    if (present (color)) then
       if (color) then
          mask%c = mask_in%c
          mask%cg = mask_in%cg
       end if
    end if
    if (present (helicity)) then
       if (helicity) then
          mask%h = mask_in%h
          mask%hd = mask_in%hd
       end if
    end if
  end subroutine quantum_numbers_mask_assign

\eatline
\nwendcode{}\nwbegindocs{550}\nwdocspar
\subsection{Mask predicates}
Return true if either one of the entries is set:
\nwenddocs{}\nwbegincode{551}\moddef{Quantum numbers: public}\plusendmoddef
  public :: any
\nwendcode{}\nwbegincode{552}\moddef{Quantum numbers: interfaces}\plusendmoddef
  interface any
     module procedure quantum_numbers_mask_any
  end interface
\nwendcode{}\nwbegincode{553}\moddef{Quantum numbers: procedures}\plusendmoddef
  function quantum_numbers_mask_any (mask) result (match)
    logical :: match
    type(quantum_numbers_mask_t), intent(in) :: mask
    match = mask%f .or. mask%c .or. mask%h .or. mask%hd
  end function quantum_numbers_mask_any

\eatline
\nwendcode{}\nwbegindocs{554}\nwdocspar
\subsection{Operators}
The OR operation is applied to all components.
\nwenddocs{}\nwbegincode{555}\moddef{Quantum numbers: quantum numbers mask: TBP}\plusendmoddef
  generic :: operator(.or.) => quantum_numbers_mask_or
  procedure, private :: quantum_numbers_mask_or
\eatline
\nwendcode{}\nwbegincode{556}\moddef{Quantum numbers: procedures}\plusendmoddef
  elemental function quantum_numbers_mask_or (mask1, mask2) result (mask)
    type(quantum_numbers_mask_t) :: mask
    class(quantum_numbers_mask_t), intent(in) :: mask1, mask2
    mask%f = mask1%f .or. mask2%f
    mask%c = mask1%c .or. mask2%c
    if (mask%c)  mask%cg = mask1%cg .or. mask2%cg
    mask%h = mask1%h .or. mask2%h
    if (.not. mask%h)  mask%hd = mask1%hd .or. mask2%hd
  end function quantum_numbers_mask_or

\eatline
\nwendcode{}\nwbegindocs{557}\nwdocspar
\subsection{Mask comparisons}
Return true if the two masks are equivalent / differ:
\nwenddocs{}\nwbegincode{558}\moddef{Quantum numbers: quantum numbers mask: TBP}\plusendmoddef
  generic :: operator(.eqv.) => quantum_numbers_mask_eqv
  generic :: operator(.neqv.) => quantum_numbers_mask_neqv
  procedure, private :: quantum_numbers_mask_eqv
  procedure, private :: quantum_numbers_mask_neqv
\nwendcode{}\nwbegincode{559}\moddef{Quantum numbers: procedures}\plusendmoddef
  elemental function quantum_numbers_mask_eqv (mask1, mask2) result (eqv)
    logical :: eqv
    class(quantum_numbers_mask_t), intent(in) :: mask1, mask2
    eqv = (mask1%f .eqv. mask2%f) .and. &
         (mask1%c .eqv. mask2%c) .and. &
         (mask1%cg .eqv. mask2%cg) .and. &
         (mask1%h .eqv. mask2%h) .and. &
         (mask1%hd .eqv. mask2%hd)
  end function quantum_numbers_mask_eqv

  elemental function quantum_numbers_mask_neqv (mask1, mask2) result (neqv)
    logical :: neqv
    class(quantum_numbers_mask_t), intent(in) :: mask1, mask2
    neqv = (mask1%f .neqv. mask2%f) .or. &
         (mask1%c .neqv. mask2%c) .or. &
         (mask1%cg .neqv. mask2%cg) .or. &
         (mask1%h .neqv. mask2%h) .or. &
         (mask1%hd .neqv. mask2%hd)
  end function quantum_numbers_mask_neqv

\eatline
\nwendcode{}\nwbegindocs{560}\nwdocspar
\subsection{Apply a mask}
Applying a mask to the quantum number object means undefining those
entries where the mask is set.  The others remain unaffected.

The {\tt{}hd} mask has the special property that it ``diagonalizes''
helicity, i.e., the second helicity entry is dropped and the result is
a diagonal helicity quantum number.
\nwenddocs{}\nwbegincode{561}\moddef{Quantum numbers: quantum numbers: TBP}\plusendmoddef
  procedure :: undefine => quantum_numbers_undefine
  procedure :: undefined => quantum_numbers_undefined0
\nwendcode{}\nwbegincode{562}\moddef{Quantum numbers: public}\plusendmoddef
  public :: quantum_numbers_undefined
\nwendcode{}\nwbegincode{563}\moddef{Quantum numbers: interfaces}\plusendmoddef
  interface quantum_numbers_undefined
     module procedure quantum_numbers_undefined0
     module procedure quantum_numbers_undefined1
     module procedure quantum_numbers_undefined11
  end interface

\nwendcode{}\nwbegincode{564}\moddef{Quantum numbers: procedures}\plusendmoddef
  elemental subroutine quantum_numbers_undefine (qn, mask)
    class(quantum_numbers_t), intent(inout) :: qn
    type(quantum_numbers_mask_t), intent(in) :: mask
    if (mask%f)  call qn%f%undefine ()
    if (mask%c)  call qn%c%undefine (undefine_ghost=mask%cg)
    if (mask%h) then
       call qn%h%undefine ()
    else if (mask%hd) then
       if (.not. qn%h%is_diagonal ()) then
          call qn%h%diagonalize ()
       end if
    end if
  end subroutine quantum_numbers_undefine

  function quantum_numbers_undefined0 (qn, mask) result (qn_new)
    class(quantum_numbers_t), intent(in) :: qn
    type(quantum_numbers_mask_t), intent(in) :: mask
    type(quantum_numbers_t) :: qn_new
    select type (qn)
    type is (quantum_numbers_t);  qn_new = qn
    end select
    call quantum_numbers_undefine (qn_new, mask)
  end function quantum_numbers_undefined0

  function quantum_numbers_undefined1 (qn, mask) result (qn_new)
    type(quantum_numbers_t), dimension(:), intent(in) :: qn
    type(quantum_numbers_mask_t), intent(in) :: mask
    type(quantum_numbers_t), dimension(size(qn)) :: qn_new
    qn_new = qn
    call quantum_numbers_undefine (qn_new, mask)
  end function quantum_numbers_undefined1

  function quantum_numbers_undefined11 (qn, mask) result (qn_new)
    type(quantum_numbers_t), dimension(:), intent(in) :: qn
    type(quantum_numbers_mask_t), dimension(:), intent(in) :: mask
    type(quantum_numbers_t), dimension(size(qn)) :: qn_new
    qn_new = qn
    call quantum_numbers_undefine (qn_new, mask)
  end function quantum_numbers_undefined11

\eatline
\eatline
\nwendcode{}\nwbegindocs{565}Return true if the input quantum number set has entries that would
be removed by the applied mask, e.g., if polarization is defined but
{\tt{}mask{\char37}h} is set:
\nwenddocs{}\nwbegincode{566}\moddef{Quantum numbers: quantum numbers: TBP}\plusendmoddef
  procedure :: are_redundant => quantum_numbers_are_redundant
\nwendcode{}\nwbegincode{567}\moddef{Quantum numbers: procedures}\plusendmoddef
  elemental function quantum_numbers_are_redundant (qn, mask) &
       result (redundant)
    logical :: redundant
    class(quantum_numbers_t), intent(in) :: qn
    type(quantum_numbers_mask_t), intent(in) :: mask
    redundant = .false.
    if (mask%f) then
       redundant = qn%f%is_defined ()
    end if
    if (mask%c) then
       redundant = qn%c%is_defined ()
    end if
    if (mask%h) then
       redundant = qn%h%is_defined ()
    else if (mask%hd) then
       redundant = .not. qn%h%is_diagonal ()
    end if
  end function quantum_numbers_are_redundant

\eatline
\nwendcode{}\nwbegindocs{568}Return true if the helicity flag is set or the diagonal-helicity flag is
set.
\nwenddocs{}\nwbegincode{569}\moddef{Quantum numbers: quantum numbers mask: TBP}\plusendmoddef
  procedure :: diagonal_helicity => quantum_numbers_mask_diagonal_helicity
\nwendcode{}\nwbegincode{570}\moddef{Quantum numbers: procedures}\plusendmoddef
  elemental function quantum_numbers_mask_diagonal_helicity (mask) &
       result (flag)
    logical :: flag
    class(quantum_numbers_mask_t), intent(in) :: mask
    flag = mask%h .or. mask%hd
  end function quantum_numbers_mask_diagonal_helicity

\eatline
\nwendcode{}\nwbegindocs{571}\nwdocspar
%%%%%%%%%%%%%%%%%%%%%%%%%%%%%%%%%%%%%%%%%%%%%%%%%%%%%%%%%%%%%%%%%%%%%%%%
\chapter{Transition Matrices and Evaluation}

The modules in this chapter implement transition matrices and calculations.
The functionality is broken down in three modules
\begin{description}
\item[state\_matrices]
  represent state and transition density matrices built from particle quantum
  numbers (helicity, color, flavor)
\item[interactions]
  extend state matrices with the record of particle momenta.  They also
  distinguish in- and out-particles and store parent-child relations.
\item[evaluators]
  These objects extend interaction objects by the information how to calculate
  matrix elements from products and squares of other interactions.  They
  implement the methods to actually compute those matrix elements.
\end{description}

\clearpage
%%%%%%%%%%%%%%%%%%%%%%%%%%%%%%%%%%%%%%%%%%%%%%%%%%%%%%%%%%%%%%%%%%%%%%%%
\section{State matrices}
This module deals with the internal state of a particle system, i.e.,
with its density matrix in flavor, color, and helicity space.
\nwenddocs{}\nwbegincode{572}\moddef{\code{}state{\char95}matrices.f90\edoc{}}\endmoddef
\LA{}File header\RA{}

module state_matrices

\LA{}Use kinds\RA{}
  use io_units
  use format_utils, only: pac_fmt
  use format_defs, only: FMT_17, FMT_19
  use unit_tests
  use diagnostics
  use sorting
  use model_data
  use flavors
  use colors
  use helicities
  use quantum_numbers

\LA{}Standard module head\RA{}

\LA{}State matrices: public\RA{}

\LA{}State matrices: parameters\RA{}
       
\LA{}State matrices: types\RA{}

\LA{}State matrices: interfaces\RA{}

contains

\LA{}State matrices: procedures\RA{}

\LA{}State matrices: tests\RA{}

end module state_matrices
\eatline
\nwendcode{}\nwbegindocs{573}\nwdocspar
\subsection{Nodes of the quantum state trie}
A quantum state object represents an unnormalized density matrix,
i.e., an array of possibilities for flavor, color, and helicity
indices with associated complex values.  Physically, the trace of this
matrix is the summed squared matrix element for an interaction, and
the matrix elements divided by this value correspond to the
flavor-color-helicity density matrix.  (Flavor and color are
diagonal.)

We store density matrices as tries, that is, as trees where each
branching represents the possible quantum numbers of a particle.  The
first branching is the first particle in the system.  A leaf (the node
corresponding to the last particle) contains the value of the matrix
element.

Each node contains a flavor, color, and helicity entry.  Note that
each of those entries may be actually undefined, so we can also represent,
e.g., unpolarized particles.

The value is meaningful only for leaves, which have no child nodes.
There is a pointer to the parent node which allows for following the
trie downwards from a leaf, it is null for a root node.  The child
nodes are implemented as a list, so there is a pointer to the first
and last child, and each node also has a {\tt{}next} pointer to the next
sibling.

The root node does not correspond to a particle, only its children do.
The quantum numbers of the root node are irrelevant and will not be
set.  However, we use a common type for the three classes (root,
branch, leaf); they may easily be distinguished by the association
status of parent and child.

\subsubsection{Node type}
The node is linked in all directions: the parent, the first and last
in the list of children, and the previous and next sibling.  This allows
us for adding and removing nodes and whole branches anywhere in the
trie.  (Circular links are not allowed, however.).  The node holds its
associated set of quantum numbers.  The integer index, which is set
only for leaf nodes, is the index of the corresponding matrix element
value within the state matrix.

Temporarily, matrix-element values may be stored within a leaf node.
This is used during state-matrix factorization.  When the state matrix
is {\tt{}freeze}d, these values are transferred to the matrix-element
array within the host state matrix.
\nwenddocs{}\nwbegincode{574}\moddef{State matrices: types}\endmoddef
  type :: node_t
     private
     type(quantum_numbers_t) :: qn
     type(node_t), pointer :: parent => null ()
     type(node_t), pointer :: child_first => null ()
     type(node_t), pointer :: child_last => null ()
     type(node_t), pointer :: next => null ()
     type(node_t), pointer :: previous => null ()
     integer :: me_index = 0
     integer, dimension(:), allocatable :: me_count
     complex(default) :: me = 0
  end type node_t

\eatline
\nwendcode{}\nwbegindocs{575}\nwdocspar
\subsubsection{Operations on nodes}
Recursively deallocate all children of the current
node.  This includes any values associated with the children.
\nwenddocs{}\nwbegincode{576}\moddef{State matrices: procedures}\endmoddef
  pure recursive subroutine node_delete_offspring (node)
    type(node_t), pointer :: node
    type(node_t), pointer :: child
    child => node%child_first
    do while (associated (child))
       node%child_first => node%child_first%next
       call node_delete_offspring (child)
       deallocate (child)
       child => node%child_first
    end do
    node%child_last => null ()
  end subroutine node_delete_offspring

\eatline
\nwendcode{}\nwbegindocs{577}Remove a node including its offspring.  Adjust the pointers of
parent and siblings, if necessary.
\nwenddocs{}\nwbegincode{578}\moddef{State matrices: procedures}\plusendmoddef
  pure subroutine node_delete (node)
    type(node_t), pointer :: node
    call node_delete_offspring (node)
    if (associated (node%previous)) then
       node%previous%next => node%next
    else if (associated (node%parent)) then
       node%parent%child_first => node%next
    end if
    if (associated (node%next)) then
       node%next%previous => node%previous
    else if (associated (node%parent)) then
       node%parent%child_last => node%previous
    end if
    deallocate (node)
  end subroutine node_delete

\eatline
\nwendcode{}\nwbegindocs{579}Append a child node
\nwenddocs{}\nwbegincode{580}\moddef{State matrices: procedures}\plusendmoddef
  subroutine node_append_child (node, child)
    type(node_t), target, intent(inout) :: node
    type(node_t), pointer :: child
    allocate (child)
    if (associated (node%child_last)) then
       node%child_last%next => child
       child%previous => node%child_last
    else
       node%child_first => child
    end if
    node%child_last => child
    child%parent => node
  end subroutine node_append_child

\eatline
\nwendcode{}\nwbegindocs{581}\nwdocspar
\subsubsection{I/O}
Output of a single node, no recursion.  We print the quantum numbers
in square brackets, then the value (if any).
\nwenddocs{}\nwbegincode{582}\moddef{State matrices: procedures}\plusendmoddef
  subroutine node_write (node, me_array, verbose, unit, col_verbose, testflag)
    type(node_t), intent(in) :: node
    complex(default), dimension(:), intent(in), optional :: me_array
    logical, intent(in), optional :: verbose, col_verbose, testflag
    integer, intent(in), optional :: unit
    logical :: verb
    integer :: u
    character(len=7) :: fmt
    call pac_fmt (fmt, FMT_19, FMT_17, testflag)
    verb = .false.;  if (present (verbose)) verb = verbose
    u = given_output_unit (unit);  if (u < 0)  return
    call node%qn%write (u, col_verbose)
    if (node%me_index /= 0) then
       write (u, "(A,I0,A)", advance="no")  " => ME(", node%me_index, ")"
       if (present (me_array)) then
          write (u, "(A)", advance="no")  " = "
          write (u, "('('," // fmt // ",','," // fmt // ",')')", &
               advance="no") pacify_complex (me_array(node%me_index))
       end if
    end if
    write (u, *)
    if (verb) then
       call ptr_write ("parent     ", node%parent)
       call ptr_write ("child_first", node%child_first)
       call ptr_write ("child_last ", node%child_last)
       call ptr_write ("next       ", node%next)
       call ptr_write ("previous   ", node%previous)
    end if
  contains
    subroutine ptr_write (label, node)
      character(*), intent(in) :: label
      type(node_t), pointer :: node
      if (associated (node)) then
         write (u, "(10x,A,1x,'->',1x)", advance="no") label
         call node%qn%write (u, col_verbose)
         write (u, *)
      end if
    end subroutine ptr_write
  end subroutine node_write

\eatline
\nwendcode{}\nwbegindocs{583}Recursive output of a node:
\nwenddocs{}\nwbegincode{584}\moddef{State matrices: procedures}\plusendmoddef
  recursive subroutine node_write_rec (node, me_array, verbose, &
        indent, unit, col_verbose, testflag)
    type(node_t), intent(in), target :: node
    complex(default), dimension(:), intent(in), optional :: me_array
    logical, intent(in), optional :: verbose, col_verbose, testflag
    integer, intent(in), optional :: indent
    integer, intent(in), optional :: unit
    type(node_t), pointer :: current
    logical :: verb
    integer :: i, u
    verb = .false.;  if (present (verbose))  verb = verbose
    i = 0;  if (present (indent)) i = indent
    u = given_output_unit (unit);  if (u < 0)  return
    current => node%child_first
    do while (associated (current))
       write (u, "(A)", advance="no")  repeat (" ", i)
       call node_write (current, me_array, verbose=verb, &
          unit=u, col_verbose=col_verbose, testflag=testflag)
       call node_write_rec (current, me_array, verbose=verb, &
          indent=i+2, unit=u, col_verbose=col_verbose, testflag=testflag)
       current => current%next
    end do
  end subroutine node_write_rec
  
\eatline
\nwendcode{}\nwbegindocs{585}Binary I/O.  Matrix elements are written only for leaf nodes.
\nwenddocs{}\nwbegincode{586}\moddef{State matrices: procedures}\plusendmoddef
  recursive subroutine node_write_raw_rec (node, u)
    type(node_t), intent(in), target :: node
    integer, intent(in) :: u
    logical :: associated_child_first, associated_next
    call node%qn%write_raw (u)
    associated_child_first = associated (node%child_first)
    write (u) associated_child_first
    associated_next = associated (node%next)
    write (u) associated_next
    if (associated_child_first) then
       call node_write_raw_rec (node%child_first, u)
    else
       write (u)  node%me_index
       write (u)  node%me
    end if
    if (associated_next) then
       call node_write_raw_rec (node%next, u)
    end if
  end subroutine node_write_raw_rec

  recursive subroutine node_read_raw_rec (node, u, parent, iostat)
    type(node_t), intent(out), target :: node
    integer, intent(in) :: u
    type(node_t), intent(in), optional, target :: parent
    integer, intent(out), optional :: iostat
    logical :: associated_child_first, associated_next
    type(node_t), pointer :: child
    call node%qn%read_raw (u, iostat=iostat)
    read (u, iostat=iostat) associated_child_first
    read (u, iostat=iostat) associated_next
    if (present (parent))  node%parent => parent
    if (associated_child_first) then
       allocate (child)
       node%child_first => child
       node%child_last => null ()
       call node_read_raw_rec (child, u, node, iostat=iostat)
       do while (associated (child))
          child%previous => node%child_last
          node%child_last => child
          child => child%next
       end do
    else
       read (u, iostat=iostat)  node%me_index
       read (u, iostat=iostat)  node%me
    end if
    if (associated_next) then
       allocate (node%next)
       call node_read_raw_rec (node%next, u, parent, iostat=iostat)
    end if
  end subroutine node_read_raw_rec

\nwendcode{}\nwbegindocs{587}\nwdocspar
\subsection{State matrix}
\subsubsection{Definition}
The quantum state object is a container that keeps and hides the root
node.  For direct accessibility of values, they are stored
in a separate array.  The leaf nodes of the quantum-number tree point to those
values, once the state matrix is finalized.

The {\tt{}norm} component is redefined if a common factor is extracted from all
nodes.
\nwenddocs{}\nwbegincode{588}\moddef{State matrices: public}\endmoddef
  public :: state_matrix_t
\nwendcode{}\nwbegincode{589}\moddef{State matrices: types}\plusendmoddef
  type :: state_matrix_t
     private
     type(node_t), pointer :: root => null ()
     integer :: depth = 0
     integer :: n_matrix_elements = 0
     logical :: leaf_nodes_store_values = .false.
     integer :: n_counters = 0
     complex(default), dimension(:), allocatable :: me
     real(default) :: norm = 1
   contains
   \LA{}State matrices: state matrix: TBP\RA{}
  end type state_matrix_t

\eatline
\nwendcode{}\nwbegindocs{590}This initializer allocates the root node but does not fill
anything.  We declare whether values are stored within the nodes
during state-matrix construction, and how many counters should be
maintained (default: none).
\nwenddocs{}\nwbegincode{591}\moddef{State matrices: state matrix: TBP}\endmoddef
  procedure :: init => state_matrix_init
\nwendcode{}\nwbegincode{592}\moddef{State matrices: procedures}\plusendmoddef
  subroutine state_matrix_init (state, store_values, n_counters)
    class(state_matrix_t), intent(out) :: state
    logical, intent(in), optional :: store_values
    integer, intent(in), optional :: n_counters
    allocate (state%root)
    if (present (store_values)) &
       state%leaf_nodes_store_values = store_values
    if (present (n_counters)) state%n_counters = n_counters
  end subroutine state_matrix_init

\eatline
\nwendcode{}\nwbegindocs{593}This recursively deletes all children of the root node, restoring
the initial state.  The matrix element array is not finalized, since
it does not contain physical entries, just pointers.
\nwenddocs{}\nwbegincode{594}\moddef{State matrices: state matrix: TBP}\plusendmoddef
  procedure :: final => state_matrix_final
\nwendcode{}\nwbegincode{595}\moddef{State matrices: procedures}\plusendmoddef
  subroutine state_matrix_final (state)
    class(state_matrix_t), intent(inout) :: state
    if (allocated (state%me))  deallocate (state%me)
    if (associated (state%root))  call node_delete (state%root)
    state%depth = 0
    state%n_matrix_elements = 0
  end subroutine state_matrix_final

\eatline
\nwendcode{}\nwbegindocs{596}Output: Present the tree as a nested list with appropriate
indentation.
\nwenddocs{}\nwbegincode{597}\moddef{State matrices: state matrix: TBP}\plusendmoddef
  procedure :: write => state_matrix_write
\nwendcode{}\nwbegincode{598}\moddef{State matrices: procedures}\plusendmoddef
  subroutine state_matrix_write (state, unit, write_value_list, &
        verbose, col_verbose, testflag)
    class(state_matrix_t), intent(in) :: state
    logical, intent(in), optional :: write_value_list, verbose, col_verbose
    logical, intent(in), optional :: testflag
    integer, intent(in), optional :: unit
    complex(default) :: me_dum
    character(len=7) :: fmt
    integer :: u
    integer :: i
    call pac_fmt (fmt, FMT_19, FMT_17, testflag)        
    u = given_output_unit (unit);  if (u < 0)  return
    write (u, "(1x,A," // fmt // ")") "State matrix:  norm = ", state%norm
    if (associated (state%root)) then
       if (allocated (state%me)) then
          call node_write_rec (state%root, state%me, verbose=verbose, &
             indent=1, unit=u, col_verbose=col_verbose, testflag=testflag)
       else
          call node_write_rec (state%root, verbose=verbose, indent=1, &
             unit=u, col_verbose=col_verbose, testflag=testflag)
       end if
    end if
    if (present (write_value_list)) then
       if (write_value_list .and. allocated (state%me)) then
          do i = 1, size (state%me)
             write (u, "(1x,I0,A)", advance="no")  i, ":"
             me_dum = state%me(i)             
             if (real(state%me(i)) == -real(state%me(i))) then
                me_dum = &
                     cmplx (0._default, aimag(me_dum), kind=default)
             end if
             if (aimag(me_dum) == -aimag(me_dum)) then
                me_dum = &
                     cmplx (real(me_dum), 0._default, kind=default)
             end if
             write (u, "('('," // fmt // ",','," // fmt // &
                  ",')')")  me_dum
          end do
       end if
    end if
  end subroutine state_matrix_write

\eatline
\nwendcode{}\nwbegindocs{599}Binary I/O.  The auxiliary matrix-element array is not written, but
reconstructed after reading the tree.

Note: To be checked.  Might be broken, don't use (unless trivial).
\nwenddocs{}\nwbegincode{600}\moddef{State matrices: state matrix: TBP}\plusendmoddef
  procedure :: write_raw => state_matrix_write_raw
  procedure :: read_raw => state_matrix_read_raw
\nwendcode{}\nwbegincode{601}\moddef{State matrices: procedures}\plusendmoddef
  subroutine state_matrix_write_raw (state, u)
    class(state_matrix_t), intent(in), target :: state
    integer, intent(in) :: u
    logical :: is_defined
    integer :: depth, j
    type(state_iterator_t) :: it
    type(quantum_numbers_t), dimension(:), allocatable :: qn
    is_defined = state%is_defined ()
    write (u)  is_defined
    if (is_defined) then
       write (u)  state%get_norm ()
       write (u)  state%get_n_leaves ()
       depth = state%get_depth ()
       write (u)  depth
       allocate (qn (depth))
       call it%init (state)
       do while (it%is_valid ())
          qn = it%get_quantum_numbers ()
          do j = 1, depth
             call qn(j)%write_raw (u)
          end do
          write (u)  it%get_me_index ()
          write (u)  it%get_matrix_element ()
          call it%advance ()
       end do
    end if
  end subroutine state_matrix_write_raw

  subroutine state_matrix_read_raw (state, u, iostat)
    class(state_matrix_t), intent(out) :: state
    integer, intent(in) :: u
    integer, intent(out) :: iostat
    logical :: is_defined
    real(default) :: norm
    integer :: n_leaves, depth, i, j
    type(quantum_numbers_t), dimension(:), allocatable :: qn
    integer :: me_index
    complex(default) :: me
    read (u, iostat=iostat)  is_defined
    if (iostat /= 0)  goto 1
    if (is_defined) then
       call state%init (store_values = .true.)
       read (u, iostat=iostat)  norm
       if (iostat /= 0)  goto 1
       call state_matrix_set_norm (state, norm)
       read (u)  n_leaves
       if (iostat /= 0)  goto 1
       read (u)  depth
       if (iostat /= 0)  goto 1
       allocate (qn (depth))
       do i = 1, n_leaves
          do j = 1, depth
             call qn(j)%read_raw (u, iostat=iostat)
             if (iostat /= 0)  goto 1
          end do
          read (u, iostat=iostat)  me_index
          if (iostat /= 0)  goto 1
          read (u, iostat=iostat)  me
          if (iostat /= 0)  goto 1
          call state%add_state (qn, index = me_index, value = me)
       end do
       call state_matrix_freeze (state)
    end if
    return

    ! Clean up on error
1   continue
    call state%final ()
  end subroutine state_matrix_read_raw

\eatline
\nwendcode{}\nwbegindocs{602}Assign a model pointer to all flavor entries.  This will become
necessary when we have read a state matrix from file.
\nwenddocs{}\nwbegincode{603}\moddef{State matrices: state matrix: TBP}\plusendmoddef
  procedure :: set_model => state_matrix_set_model
\nwendcode{}\nwbegincode{604}\moddef{State matrices: procedures}\plusendmoddef
  subroutine state_matrix_set_model (state, model)
    class(state_matrix_t), intent(inout), target :: state
    class(model_data_t), intent(in), target :: model
    type(state_iterator_t) :: it
    call it%init (state)
    do while (it%is_valid ())
       call it%set_model (model)
       call it%advance ()
    end do
  end subroutine state_matrix_set_model
    
\eatline
\nwendcode{}\nwbegindocs{605}\nwdocspar
\subsubsection{Properties of the quantum state}
A state is defined if its root is allocated:
\nwenddocs{}\nwbegincode{606}\moddef{State matrices: state matrix: TBP}\plusendmoddef
  procedure :: is_defined => state_matrix_is_defined
\nwendcode{}\nwbegincode{607}\moddef{State matrices: procedures}\plusendmoddef
  elemental function state_matrix_is_defined (state) result (defined)
    logical :: defined
    class(state_matrix_t), intent(in) :: state
    defined = associated (state%root)
  end function state_matrix_is_defined

\eatline
\nwendcode{}\nwbegindocs{608}A state is empty if its depth is zero:
\nwenddocs{}\nwbegincode{609}\moddef{State matrices: state matrix: TBP}\plusendmoddef
  procedure :: is_empty => state_matrix_is_empty
\nwendcode{}\nwbegincode{610}\moddef{State matrices: procedures}\plusendmoddef
  elemental function state_matrix_is_empty (state) result (flag)
    logical :: flag
    class(state_matrix_t), intent(in) :: state
    flag = state%depth == 0
  end function state_matrix_is_empty

\eatline
\nwendcode{}\nwbegindocs{611}Return the number of matrix-element values.
\nwenddocs{}\nwbegincode{612}\moddef{State matrices: state matrix: TBP}\plusendmoddef
  procedure :: get_n_matrix_elements => state_matrix_get_n_matrix_elements
\nwendcode{}\nwbegincode{613}\moddef{State matrices: procedures}\plusendmoddef
  function state_matrix_get_n_matrix_elements (state) result (n)
    integer :: n
    class(state_matrix_t), intent(in) :: state
    n = state%n_matrix_elements
  end function state_matrix_get_n_matrix_elements

\eatline
\nwendcode{}\nwbegindocs{614}Return the number of leaves.  This can be larger than the number of
independent matrix elements.
\nwenddocs{}\nwbegincode{615}\moddef{State matrices: state matrix: TBP}\plusendmoddef
  procedure :: get_n_leaves => state_matrix_get_n_leaves
\nwendcode{}\nwbegincode{616}\moddef{State matrices: procedures}\plusendmoddef
  function state_matrix_get_n_leaves (state) result (n)
    integer :: n
    class(state_matrix_t), intent(in) :: state
    type(state_iterator_t) :: it
    n = 0
    call it%init (state)
    do while (it%is_valid ())
       n = n + 1
       call it%advance ()
    end do
  end function state_matrix_get_n_leaves

\eatline
\nwendcode{}\nwbegindocs{617}Return the depth:
\nwenddocs{}\nwbegincode{618}\moddef{State matrices: state matrix: TBP}\plusendmoddef
  procedure :: get_depth => state_matrix_get_depth
\nwendcode{}\nwbegincode{619}\moddef{State matrices: procedures}\plusendmoddef
  function state_matrix_get_depth (state) result (depth)
    integer :: depth
    class(state_matrix_t), intent(in) :: state
    depth = state%depth
  end function state_matrix_get_depth

\eatline
\nwendcode{}\nwbegindocs{620}Return the norm:
\nwenddocs{}\nwbegincode{621}\moddef{State matrices: state matrix: TBP}\plusendmoddef
  procedure :: get_norm => state_matrix_get_norm
\nwendcode{}\nwbegincode{622}\moddef{State matrices: procedures}\plusendmoddef
  function state_matrix_get_norm (state) result (norm)
    real(default) :: norm
    class(state_matrix_t), intent(in) :: state
    norm = state%norm
  end function state_matrix_get_norm

\eatline
\nwendcode{}\nwbegindocs{623}\nwdocspar
\subsubsection{Retrieving contents}
Return the quantum number array, using an index.  We have to scan the
state matrix since there is no shortcut.
\nwenddocs{}\nwbegincode{624}\moddef{State matrices: state matrix: TBP}\plusendmoddef
  procedure :: get_quantum_numbers => state_matrix_get_quantum_numbers
\nwendcode{}\nwbegincode{625}\moddef{State matrices: procedures}\plusendmoddef
  function state_matrix_get_quantum_numbers (state, i) result (qn)
    class(state_matrix_t), intent(in), target :: state
    integer, intent(in) :: i
    type(quantum_numbers_t), dimension(state%depth) :: qn
    type(state_iterator_t) :: it
    integer :: k
    k = 0
    call it%init (state)
    do while (it%is_valid ())
       k = k + 1
       if (k == i) then
          qn = it%get_quantum_numbers ()
          return
       end if
       call it%advance ()
    end do
  end function state_matrix_get_quantum_numbers

\eatline
\nwendcode{}\nwbegindocs{626}Return a single matrix element using its index.  Works only if the
shortcut array is allocated.
\nwenddocs{}\nwbegincode{627}\moddef{State matrices: state matrix: TBP}\plusendmoddef
  procedure :: get_matrix_element => state_matrix_get_matrix_element
\nwendcode{}\nwbegincode{628}\moddef{State matrices: procedures}\plusendmoddef
  function state_matrix_get_matrix_element (state, i) result (me)
    complex(default) :: me
    class(state_matrix_t), intent(in) :: state
    integer, intent(in) :: i
    if (allocated (state%me)) then
       me = state%me(i)
    else
       me = 0
    end if
  end function state_matrix_get_matrix_element

\eatline
\nwendcode{}\nwbegindocs{629}Return the color index with maximum absolute value that is present within
the state matrix.
\nwenddocs{}\nwbegincode{630}\moddef{State matrices: state matrix: TBP}\plusendmoddef
  procedure :: get_max_color_value => state_matrix_get_max_color_value
\nwendcode{}\nwbegincode{631}\moddef{State matrices: procedures}\plusendmoddef
  function state_matrix_get_max_color_value (state) result (cmax)
    integer :: cmax
    class(state_matrix_t), intent(in) :: state
    if (associated (state%root)) then
       cmax = node_get_max_color_value (state%root)
    else
       cmax = 0
    end if
  contains
    recursive function node_get_max_color_value (node) result (cmax)
      integer :: cmax
      type(node_t), intent(in), target :: node
      type(node_t), pointer :: current
      cmax = quantum_numbers_get_max_color_value (node%qn)
      current => node%child_first
      do while (associated (current))
         cmax = max (cmax, node_get_max_color_value (current))
         current => current%next
      end do
    end function node_get_max_color_value
  end function state_matrix_get_max_color_value

\eatline
\nwendcode{}\nwbegindocs{632}\nwdocspar
\subsubsection{Building the quantum state}
The procedure generates a branch associated to the input array of
quantum numbers.  If the branch exists already, it is used.

Optionally, we set the matrix-element index, a value (which may be
added to the previous one), and increment one of the possible
counters.  We may also return the matrix element index of the current
node.
\nwenddocs{}\nwbegincode{633}\moddef{State matrices: state matrix: TBP}\plusendmoddef
  procedure :: add_state => state_matrix_add_state
\nwendcode{}\nwbegincode{634}\moddef{State matrices: procedures}\plusendmoddef
  subroutine state_matrix_add_state &
       (state, qn, index, value, sum_values, counter_index, me_index)
    class(state_matrix_t), intent(inout) :: state
    type(quantum_numbers_t), dimension(:), intent(in) :: qn
    integer, intent(in), optional :: index
    complex(default), intent(in), optional :: value
    logical, intent(in), optional :: sum_values
    integer, intent(in), optional :: counter_index
    integer, intent(out), optional :: me_index
    logical :: set_index, get_index, add
    set_index = present (index)
    get_index = present (me_index)
    add = .false.;  if (present (sum_values))  add = sum_values
    if (state%depth == 0) then
       state%depth = size (qn)
    else if (state%depth /= size (qn)) then
       call state%write ()
       call msg_bug ("State matrix: depth mismatch")
    end if
    if (size (qn) > 0)  call node_make_branch (state%root, qn)
  contains
     recursive subroutine node_make_branch (parent, qn)
       type(node_t), pointer :: parent
       type(quantum_numbers_t), dimension(:), intent(in) :: qn
       type(node_t), pointer :: child
       logical :: match
       match = .false.
       child => parent%child_first
       SCAN_CHILDREN: do while (associated (child))
          match = child%qn == qn(1)
          if (match)  exit SCAN_CHILDREN
          child => child%next
       end do SCAN_CHILDREN
       if (.not. match) then
          call node_append_child (parent, child)
          child%qn = qn(1)
       end if
       select case (size (qn))
       case (1)
          if (.not. match) then
             state%n_matrix_elements = state%n_matrix_elements + 1
             child%me_index = state%n_matrix_elements
          end if
          if (set_index) then
             child%me_index = index
          end if
          if (get_index) then
             me_index = child%me_index
          end if
          if (present (counter_index)) then
             if (.not. allocated (child%me_count)) then
                allocate (child%me_count (state%n_counters))
                child%me_count = 0
             end if
             child%me_count(counter_index) = child%me_count(counter_index) + 1
          end if
          if (present (value)) then
             if (add) then
                child%me = child%me + value
             else
                child%me = value
             end if
          end if
       case (2:)
          call node_make_branch (child, qn(2:))
       end select
     end subroutine node_make_branch
   end subroutine state_matrix_add_state

\eatline
\nwendcode{}\nwbegindocs{635}Remove irrelevant flavor/color/helicity labels and the corresponding
branchings.  The masks indicate which particles are affected; the
masks length should coincide with the depth of the trie (without the
root node).  Recursively scan the whole tree, starting from the leaf
nodes and working up to the root node.  If a mask entry is set for the
current tree level, scan the children there.  For each child within
that level make a new empty branch where the masked quantum number is
undefined.  Then recursively combine all following children with
matching quantum number into this new node and move on.
\nwenddocs{}\nwbegincode{636}\moddef{State matrices: state matrix: TBP}\plusendmoddef
  procedure :: collapse => state_matrix_collapse
\nwendcode{}\nwbegincode{637}\moddef{State matrices: procedures}\plusendmoddef
  subroutine state_matrix_collapse (state, mask)
    class(state_matrix_t), intent(inout) :: state
    type(quantum_numbers_mask_t), dimension(:), intent(in) :: mask
    type(state_matrix_t) :: red_state
    if (state%is_defined ()) then
       call state%reduce (mask, red_state)
       call state%final ()
       state = red_state
    end if
   end subroutine state_matrix_collapse

\eatline
\nwendcode{}\nwbegindocs{638}Transform the given state matrix into a reduced state matrix where
some quantum numbers are removed, as indicated by the mask.  The
procedure creates a new state matrix, so the old one can be deleted
after this if it is no longer used.
\nwenddocs{}\nwbegincode{639}\moddef{State matrices: state matrix: TBP}\plusendmoddef
  procedure :: reduce => state_matrix_reduce
\nwendcode{}\nwbegincode{640}\moddef{State matrices: procedures}\plusendmoddef
  subroutine state_matrix_reduce (state, mask, red_state)
    class(state_matrix_t), intent(in), target :: state
    type(quantum_numbers_mask_t), dimension(:), intent(in) :: mask
    type(state_matrix_t), intent(out) :: red_state
    type(state_iterator_t) :: it
    type(quantum_numbers_t), dimension(size(mask)) :: qn
    call red_state%init ()
    call it%init (state)
    do while (it%is_valid ())
       qn = it%get_quantum_numbers ()
       call qn%undefine (mask)
       call red_state%add_state (qn)
       call it%advance ()
    end do
  end subroutine state_matrix_reduce

\eatline
\nwendcode{}\nwbegindocs{641}This subroutine sets up the matrix-element array.  The leaf nodes
aquire the index values that point to the appropriate matrix-element
entry.

We recursively scan the trie.  Once we arrive at a leaf node, the
index is increased and associated to that node.  Finally, we allocate
the matrix-element array with the appropriate size.

If matrix element values are temporarily stored within the leaf nodes,
we scan the state again and transfer them to the matrix-element array.
\nwenddocs{}\nwbegincode{642}\moddef{State matrices: state matrix: TBP}\plusendmoddef
  procedure :: freeze => state_matrix_freeze
\nwendcode{}\nwbegincode{643}\moddef{State matrices: procedures}\plusendmoddef
  subroutine state_matrix_freeze (state)
    class(state_matrix_t), intent(inout), target :: state
    type(state_iterator_t) :: it
    if (associated (state%root)) then
       if (allocated (state%me))  deallocate (state%me)
       allocate (state%me (state%n_matrix_elements))
       state%me = 0
    end if
    if (state%leaf_nodes_store_values) then
       call it%init (state)
       do while (it%is_valid ())
          state%me(it%get_me_index ()) = it%get_matrix_element ()
          call it%advance ()
       end do
       state%leaf_nodes_store_values = .false.
    end if
  end subroutine state_matrix_freeze

\eatline
\nwendcode{}\nwbegindocs{644}\nwdocspar
\subsubsection{Direct access to the value array}
Several methods for setting a value directly are summarized in this
generic:
\nwenddocs{}\nwbegincode{645}\moddef{State matrices: state matrix: TBP}\plusendmoddef
  generic :: set_matrix_element => set_matrix_element_qn
  generic :: set_matrix_element => set_matrix_element_all
  generic :: set_matrix_element => set_matrix_element_array
  generic :: set_matrix_element => set_matrix_element_single
  generic :: set_matrix_element => set_matrix_element_clone
  procedure :: set_matrix_element_qn => state_matrix_set_matrix_element_qn
  procedure :: set_matrix_element_all => state_matrix_set_matrix_element_all
  procedure :: set_matrix_element_array => &
       state_matrix_set_matrix_element_array 
  procedure :: set_matrix_element_single => &
       state_matrix_set_matrix_element_single
  procedure :: set_matrix_element_clone => &
     state_matrix_set_matrix_element_clone
\eatline
\nwendcode{}\nwbegindocs{646}Set a value that corresponds to a quantum number array:
\nwenddocs{}\nwbegincode{647}\moddef{State matrices: procedures}\plusendmoddef
  subroutine state_matrix_set_matrix_element_qn (state, qn, value)
    class(state_matrix_t), intent(inout) :: state
    type(quantum_numbers_t), dimension(:), intent(in) :: qn
    complex(default), intent(in) :: value
    type(state_iterator_t) :: it
    if (.not. allocated (it%state%me)) then
       allocate (it%state%me (size(qn)))
    end if
    call it%init (state)
    do while (it%is_valid ())
       if (all (qn == it%get_quantum_numbers ())) then
          call it%set_matrix_element (value)
          return
       end if
       call it%advance ()
    end do
  end subroutine state_matrix_set_matrix_element_qn

\eatline
\nwendcode{}\nwbegindocs{648}Set all matrix elements to a single value
\nwenddocs{}\nwbegincode{649}\moddef{State matrices: procedures}\plusendmoddef
  subroutine state_matrix_set_matrix_element_all (state, value)
    class(state_matrix_t), intent(inout) :: state
    complex(default), intent(in) :: value
    if (.not. allocated (state%me)) then
       allocate (state%me (state%n_matrix_elements))    
    end if
    state%me = value
  end subroutine state_matrix_set_matrix_element_all

\eatline
\nwendcode{}\nwbegindocs{650}Set the matrix-element array directly.
\nwenddocs{}\nwbegincode{651}\moddef{State matrices: procedures}\plusendmoddef
  subroutine state_matrix_set_matrix_element_array (state, value)
    class(state_matrix_t), intent(inout) :: state
    complex(default), dimension(:), intent(in) :: value
    if (.not. allocated (state%me)) then
       allocate (state%me (size (value)))
    end if
    state%me = value
  end subroutine state_matrix_set_matrix_element_array

  pure subroutine state_matrix_set_matrix_element_single (state, i, value)
    class(state_matrix_t), intent(inout) :: state
    integer, intent(in) :: i
    complex(default), intent(in) :: value
    if (.not. allocated (state%me)) then
       allocate (state%me (state%n_matrix_elements))    
    end if
    state%me(i) = value
  end subroutine state_matrix_set_matrix_element_single

\eatline
\eatline
\nwendcode{}\nwbegindocs{652}Clone the matrix elements from another (matching) state matrix.
\nwenddocs{}\nwbegincode{653}\moddef{State matrices: procedures}\plusendmoddef
  subroutine state_matrix_set_matrix_element_clone (state, state1)
    class(state_matrix_t), intent(inout) :: state
    type(state_matrix_t), intent(in) :: state1
    if (.not. allocated (state1%me)) return
    if (.not. allocated (state%me)) allocate (state%me (size (state1%me)))
    state%me = state1%me
  end subroutine state_matrix_set_matrix_element_clone

\eatline
\nwendcode{}\nwbegindocs{654}Add a value to a matrix element
\nwenddocs{}\nwbegincode{655}\moddef{State matrices: state matrix: TBP}\plusendmoddef
  procedure :: add_to_matrix_element => state_matrix_add_to_matrix_element
\nwendcode{}\nwbegincode{656}\moddef{State matrices: procedures}\plusendmoddef
  subroutine state_matrix_add_to_matrix_element (state, i, value)
    class(state_matrix_t), intent(inout) :: state
    integer, intent(in) :: i
    complex(default), intent(in) :: value
    state%me(i) = state%me(i) + value
  end subroutine state_matrix_add_to_matrix_element

\eatline
\nwendcode{}\nwbegindocs{657}\nwdocspar
\subsection{State iterators}
Accessing the quantum state from outside is best done using a
specialized iterator, i.e., a pointer to a particular branch of the
quantum state trie.  Technically, the iterator contains a pointer to a
leaf node, but via parent pointers it allows to access the whole
branch where the leaf is attached.  For quick access, we also keep the
branch depth (which is assumed to be universal for a quantum state).
\nwenddocs{}\nwbegincode{658}\moddef{State matrices: public}\plusendmoddef
  public :: state_iterator_t
\nwendcode{}\nwbegincode{659}\moddef{State matrices: types}\plusendmoddef
  type :: state_iterator_t
     private
     integer :: depth = 0
     type(state_matrix_t), pointer :: state => null ()
     type(node_t), pointer :: node => null ()
   contains
   \LA{}State matrices: state iterator: TBP\RA{}  
  end type state_iterator_t

\eatline
\nwendcode{}\nwbegindocs{660}The initializer: Point at the first branch.  Note that this cannot
be pure, thus not be elemental, because the iterator can be used to
manipulate data in the state matrix.
\nwenddocs{}\nwbegincode{661}\moddef{State matrices: state iterator: TBP}\endmoddef
  procedure :: init => state_iterator_init
\nwendcode{}\nwbegincode{662}\moddef{State matrices: procedures}\plusendmoddef
  subroutine state_iterator_init (it, state)
    class(state_iterator_t), intent(out) :: it
    type(state_matrix_t), intent(in), target :: state
    it%state => state
    it%depth = state%depth
    if (state%is_defined ()) then
       it%node => state%root
       do while (associated (it%node%child_first))
          it%node => it%node%child_first
       end do
    else
       it%node => null ()
    end if
  end subroutine state_iterator_init
    
\eatline
\nwendcode{}\nwbegindocs{663}Go forward.  Recursively programmed: if the next node does not
exist, go back to the parent node and look at its successor (if
present), etc.

There is a possible pitfall in the implementation: If the dummy
pointer argument to the {\tt{}find{\char95}next} routine is used directly, we
still get the correct result for the iterator, but calling the
recursion on {\tt{}node{\char37}parent} means that we manipulate a parent pointer
in the original state in addition to the iterator.  Making a local
copy of the pointer avoids this.  Using pointer intent would be
helpful, but we do not yet rely on this F2003 feature.
\nwenddocs{}\nwbegincode{664}\moddef{State matrices: state iterator: TBP}\plusendmoddef
  procedure :: advance => state_iterator_advance
\nwendcode{}\nwbegincode{665}\moddef{State matrices: procedures}\plusendmoddef
  subroutine state_iterator_advance (it)
    class(state_iterator_t), intent(inout) :: it
    call find_next (it%node)
  contains
    recursive subroutine find_next (node_in)
      type(node_t), intent(in), target :: node_in
      type(node_t), pointer :: node
      node => node_in
      if (associated (node%next)) then
         node => node%next
         do while (associated (node%child_first))
            node => node%child_first
         end do
         it%node => node
      else if (associated (node%parent)) then
         call find_next (node%parent)
      else
         it%node => null ()
      end if
    end subroutine find_next
  end subroutine state_iterator_advance

\eatline
\nwendcode{}\nwbegindocs{666}If all has been scanned, the iterator is at an undefined state.
Check for this:
\nwenddocs{}\nwbegincode{667}\moddef{State matrices: state iterator: TBP}\plusendmoddef
  procedure :: is_valid => state_iterator_is_valid
\nwendcode{}\nwbegincode{668}\moddef{State matrices: procedures}\plusendmoddef
  function state_iterator_is_valid (it) result (defined)
    logical :: defined
    class(state_iterator_t), intent(in) :: it
    defined = associated (it%node)
  end function state_iterator_is_valid

\eatline
\nwendcode{}\nwbegindocs{669}Return the matrix-element index that corresponds to the current node
\nwenddocs{}\nwbegincode{670}\moddef{State matrices: state iterator: TBP}\plusendmoddef
  procedure :: get_me_index => state_iterator_get_me_index
\nwendcode{}\nwbegincode{671}\moddef{State matrices: procedures}\plusendmoddef
  function state_iterator_get_me_index (it) result (n)
    integer :: n
    class(state_iterator_t), intent(in) :: it
    n = it%node%me_index
  end function state_iterator_get_me_index

\eatline
\nwendcode{}\nwbegindocs{672}Return the number of times this quantum-number state has been added
(noting that it is physically inserted only the first time).  Note
that for each state, there is an array of counters.
\nwenddocs{}\nwbegincode{673}\moddef{State matrices: state iterator: TBP}\plusendmoddef
  procedure :: get_me_count => state_iterator_get_me_count
\nwendcode{}\nwbegincode{674}\moddef{State matrices: procedures}\plusendmoddef
  function state_iterator_get_me_count (it) result (n)
    integer, dimension(:), allocatable :: n
    class(state_iterator_t), intent(in) :: it
    if (allocated (it%node%me_count)) then
       allocate (n (size (it%node%me_count)))
       n = it%node%me_count
    else
       allocate (n (0))
    end if
  end function state_iterator_get_me_count

\eatline
\nwendcode{}\nwbegindocs{675}Use the iterator to retrieve quantum-number information:
\nwenddocs{}\nwbegincode{676}\moddef{State matrices: state iterator: TBP}\plusendmoddef
  generic :: get_quantum_numbers => get_qn_multi, get_qn_slice, &
       get_qn_range, get_qn_single
  generic :: get_flavor => get_flv_multi, get_flv_slice, &
       get_flv_range, get_flv_single
  generic :: get_color => get_col_multi, get_col_slice, &
       get_col_range, get_col_single
  generic :: get_helicity => get_hel_multi, get_hel_slice, &
       get_hel_range, get_hel_single
\nwendcode{}\nwbegincode{677}\moddef{State matrices: state iterator: TBP}\plusendmoddef
  procedure :: get_qn_multi => state_iterator_get_qn_multi
  procedure :: get_qn_slice => state_iterator_get_qn_slice
  procedure :: get_qn_range => state_iterator_get_qn_range
  procedure :: get_qn_single => state_iterator_get_qn_single
  procedure :: get_flv_multi => state_iterator_get_flv_multi
  procedure :: get_flv_slice => state_iterator_get_flv_slice
  procedure :: get_flv_range => state_iterator_get_flv_range
  procedure :: get_flv_single => state_iterator_get_flv_single
  procedure :: get_col_multi => state_iterator_get_col_multi
  procedure :: get_col_slice => state_iterator_get_col_slice
  procedure :: get_col_range => state_iterator_get_col_range
  procedure :: get_col_single => state_iterator_get_col_single
  procedure :: get_hel_multi => state_iterator_get_hel_multi
  procedure :: get_hel_slice => state_iterator_get_hel_slice
  procedure :: get_hel_range => state_iterator_get_hel_range
  procedure :: get_hel_single => state_iterator_get_hel_single
\nwendcode{}\nwbegindocs{678}These versions return the whole quantum number array
\nwenddocs{}\nwbegincode{679}\moddef{State matrices: procedures}\plusendmoddef
  function state_iterator_get_qn_multi (it) result (qn)
    class(state_iterator_t), intent(in) :: it
    type(quantum_numbers_t), dimension(it%depth) :: qn
    type(node_t), pointer :: node
    integer :: i
    node => it%node
    do i = it%depth, 1, -1
       qn(i) = node%qn
       node => node%parent
    end do
  end function state_iterator_get_qn_multi

  function state_iterator_get_flv_multi (it) result (flv)
    class(state_iterator_t), intent(in) :: it
    type(flavor_t), dimension(it%depth) :: flv
    flv = quantum_numbers_get_flavor &
         (it%get_quantum_numbers ())
  end function state_iterator_get_flv_multi

  function state_iterator_get_col_multi (it) result (col)
    class(state_iterator_t), intent(in) :: it
    type(color_t), dimension(it%depth) :: col
    col = quantum_numbers_get_color &
         (it%get_quantum_numbers ())
  end function state_iterator_get_col_multi

  function state_iterator_get_hel_multi (it) result (hel)
    class(state_iterator_t), intent(in) :: it
    type(helicity_t), dimension(it%depth) :: hel
    hel = quantum_numbers_get_helicity &
         (it%get_quantum_numbers ())
  end function state_iterator_get_hel_multi

\nwendcode{}\nwbegindocs{680}An array slice (derived from the above).
\nwenddocs{}\nwbegincode{681}\moddef{State matrices: procedures}\plusendmoddef
  function state_iterator_get_qn_slice (it, index) result (qn)
    class(state_iterator_t), intent(in) :: it
    integer, dimension(:), intent(in) :: index
    type(quantum_numbers_t), dimension(size(index)) :: qn
    type(quantum_numbers_t), dimension(it%depth) :: qn_tmp
    qn_tmp = state_iterator_get_qn_multi (it)
    qn = qn_tmp(index)
  end function state_iterator_get_qn_slice

  function state_iterator_get_flv_slice (it, index) result (flv)
    class(state_iterator_t), intent(in) :: it
    integer, dimension(:), intent(in) :: index
    type(flavor_t), dimension(size(index)) :: flv
    flv = quantum_numbers_get_flavor &
         (it%get_quantum_numbers (index))
  end function state_iterator_get_flv_slice

  function state_iterator_get_col_slice (it, index) result (col)
    class(state_iterator_t), intent(in) :: it
    integer, dimension(:), intent(in) :: index
    type(color_t), dimension(size(index)) :: col
    col = quantum_numbers_get_color &
         (it%get_quantum_numbers (index))
  end function state_iterator_get_col_slice

  function state_iterator_get_hel_slice (it, index) result (hel)
    class(state_iterator_t), intent(in) :: it
    integer, dimension(:), intent(in) :: index
    type(helicity_t), dimension(size(index)) :: hel
    hel = quantum_numbers_get_helicity &
         (it%get_quantum_numbers (index))
  end function state_iterator_get_hel_slice

\nwendcode{}\nwbegindocs{682}An array range (implemented directly).
\nwenddocs{}\nwbegincode{683}\moddef{State matrices: procedures}\plusendmoddef
  function state_iterator_get_qn_range (it, k1, k2) result (qn)
    class(state_iterator_t), intent(in) :: it
    integer, intent(in) :: k1, k2
    type(quantum_numbers_t), dimension(k2-k1+1) :: qn
    type(node_t), pointer :: node
    integer :: i
    node => it%node
    SCAN: do i = it%depth, 1, -1
       if (k1 <= i .and. i <= k2) then
          qn(i-k1+1) = node%qn
       else
          node => node%parent
       end if
    end do SCAN
  end function state_iterator_get_qn_range

  function state_iterator_get_flv_range (it, k1, k2) result (flv)
    class(state_iterator_t), intent(in) :: it
    integer, intent(in) :: k1, k2
    type(flavor_t), dimension(k2-k1+1) :: flv
    flv = quantum_numbers_get_flavor &
         (it%get_quantum_numbers (k1, k2))
  end function state_iterator_get_flv_range

  function state_iterator_get_col_range (it, k1, k2) result (col)
    class(state_iterator_t), intent(in) :: it
    integer, intent(in) :: k1, k2
    type(color_t), dimension(k2-k1+1) :: col
    col = quantum_numbers_get_color &
         (it%get_quantum_numbers (k1, k2))
  end function state_iterator_get_col_range

  function state_iterator_get_hel_range (it, k1, k2) result (hel)
    class(state_iterator_t), intent(in) :: it
    integer, intent(in) :: k1, k2
    type(helicity_t), dimension(k2-k1+1) :: hel
    hel = quantum_numbers_get_helicity &
         (it%get_quantum_numbers (k1, k2))
  end function state_iterator_get_hel_range

\nwendcode{}\nwbegindocs{684}Just a specific single element
\nwenddocs{}\nwbegincode{685}\moddef{State matrices: procedures}\plusendmoddef
  function state_iterator_get_qn_single (it, k) result (qn)
    class(state_iterator_t), intent(in) :: it
    integer, intent(in) :: k
    type(quantum_numbers_t) :: qn
    type(node_t), pointer :: node
    integer :: i
    node => it%node
    SCAN: do i = it%depth, 1, -1
       if (i == k) then
          qn = node%qn
          exit SCAN
       else
          node => node%parent
       end if
    end do SCAN
  end function state_iterator_get_qn_single

  function state_iterator_get_flv_single (it, k) result (flv)
    class(state_iterator_t), intent(in) :: it
    integer, intent(in) :: k
    type(flavor_t) :: flv
    flv = quantum_numbers_get_flavor &
         (it%get_quantum_numbers (k))
  end function state_iterator_get_flv_single

  function state_iterator_get_col_single (it, k) result (col)
    class(state_iterator_t), intent(in) :: it
    integer, intent(in) :: k
    type(color_t) :: col
    col = quantum_numbers_get_color &
         (it%get_quantum_numbers (k))
  end function state_iterator_get_col_single

  function state_iterator_get_hel_single (it, k) result (hel)
    class(state_iterator_t), intent(in) :: it
    integer, intent(in) :: k
    type(helicity_t) :: hel
    hel = quantum_numbers_get_helicity &
         (it%get_quantum_numbers (k))
  end function state_iterator_get_hel_single

\eatline
\eatline
\eatline
\eatline
\nwendcode{}\nwbegindocs{686}Assign a model pointer to the current flavor entries.
\nwenddocs{}\nwbegincode{687}\moddef{State matrices: state iterator: TBP}\plusendmoddef
  procedure :: set_model => state_iterator_set_model
\nwendcode{}\nwbegincode{688}\moddef{State matrices: procedures}\plusendmoddef
  subroutine state_iterator_set_model (it, model)
    class(state_iterator_t), intent(inout) :: it
    class(model_data_t), intent(in), target :: model
    type(node_t), pointer :: node
    integer :: i
    node => it%node
    do i = it%depth, 1, -1
       call node%qn%set_model (model)
       node => node%parent
    end do
  end subroutine state_iterator_set_model
  
\eatline
\nwendcode{}\nwbegindocs{689}Retrieve the matrix element value associated with the current node.
\nwenddocs{}\nwbegincode{690}\moddef{State matrices: state iterator: TBP}\plusendmoddef
  procedure :: get_matrix_element => state_iterator_get_matrix_element
\nwendcode{}\nwbegincode{691}\moddef{State matrices: procedures}\plusendmoddef
  function state_iterator_get_matrix_element (it) result (me)
    complex(default) :: me
    class(state_iterator_t), intent(in) :: it
    if (it%state%leaf_nodes_store_values) then
       me = it%node%me
    else if (it%node%me_index /= 0) then
       me = it%state%me(it%node%me_index)
    else
       me = 0
    end if
  end function state_iterator_get_matrix_element

\eatline
\nwendcode{}\nwbegindocs{692}Set the matrix element value using the state iterator.
\nwenddocs{}\nwbegincode{693}\moddef{State matrices: state iterator: TBP}\plusendmoddef
  procedure :: set_matrix_element => state_iterator_set_matrix_element
\nwendcode{}\nwbegincode{694}\moddef{State matrices: procedures}\plusendmoddef
  subroutine state_iterator_set_matrix_element (it, value)
    class(state_iterator_t), intent(inout) :: it
    complex(default), intent(in) :: value
    if (it%node%me_index /= 0) then
       it%state%me(it%node%me_index) = value
    end if
  end subroutine state_iterator_set_matrix_element

\eatline
\nwendcode{}\nwbegindocs{695}\nwdocspar
\subsection{Operations on quantum states}
Return a deep copy of a state matrix.
\nwenddocs{}\nwbegincode{696}\moddef{State matrices: public}\plusendmoddef
  public :: assignment(=)
\nwendcode{}\nwbegincode{697}\moddef{State matrices: interfaces}\endmoddef
  interface assignment(=)
     module procedure state_matrix_assign
  end interface

\nwendcode{}\nwbegincode{698}\moddef{State matrices: procedures}\plusendmoddef
  subroutine state_matrix_assign (state_out, state_in)
    type(state_matrix_t), intent(out) :: state_out
    type(state_matrix_t), intent(in), target :: state_in
    type(state_iterator_t) :: it
    if (.not. state_in%is_defined ())  return
    call state_out%init ()
    call it%init (state_in)
    do while (it%is_valid ())
       call state_out%add_state (it%get_quantum_numbers (), &
            it%get_me_index ())
       call it%advance ()
    end do
    if (allocated (state_in%me)) then
       allocate (state_out%me (size (state_in%me)))
       state_out%me = state_in%me
    end if
  end subroutine state_matrix_assign

\eatline
\nwendcode{}\nwbegindocs{699}Determine the indices of all diagonal matrix elements.
\nwenddocs{}\nwbegincode{700}\moddef{State matrices: state matrix: TBP}\plusendmoddef
  procedure :: get_diagonal_entries => state_matrix_get_diagonal_entries
\nwendcode{}\nwbegincode{701}\moddef{State matrices: procedures}\plusendmoddef
  subroutine state_matrix_get_diagonal_entries (state, i)
    class(state_matrix_t), intent(in) :: state
    integer, dimension(:), allocatable, intent(out) :: i
    integer, dimension(state%n_matrix_elements) :: tmp
    integer :: n
    type(state_iterator_t) :: it
    type(quantum_numbers_t), dimension(:), allocatable :: qn
    n = 0
    call it%init (state)
    allocate (qn (it%depth))
    do while (it%is_valid ())
       qn = it%get_quantum_numbers ()
       if (all (qn%are_diagonal ())) then
          n = n + 1
          tmp(n) = it%get_me_index ()
       end if
       call it%advance ()
    end do
    allocate (i(n))
    if (n > 0) i = tmp(:n)
  end subroutine state_matrix_get_diagonal_entries

\eatline
\nwendcode{}\nwbegindocs{702}Normalize all matrix elements, i.e., multiply by a common factor.
Assuming that the factor is nonzero, of course.
\nwenddocs{}\nwbegincode{703}\moddef{State matrices: state matrix: TBP}\plusendmoddef
  procedure :: renormalize => state_matrix_renormalize
\nwendcode{}\nwbegincode{704}\moddef{State matrices: procedures}\plusendmoddef
  subroutine state_matrix_renormalize (state, factor)
    class(state_matrix_t), intent(inout) :: state
    complex(default), intent(in) :: factor
    state%me = state%me * factor
  end subroutine state_matrix_renormalize

\eatline
\nwendcode{}\nwbegindocs{705}Renormalize the state matrix by its trace, if nonzero.  The renormalization
is reflected in the state-matrix norm.
\nwenddocs{}\nwbegincode{706}\moddef{State matrices: state matrix: TBP}\plusendmoddef
  procedure :: normalize_by_trace => state_matrix_normalize_by_trace
\nwendcode{}\nwbegincode{707}\moddef{State matrices: procedures}\plusendmoddef
  subroutine state_matrix_normalize_by_trace (state)
    class(state_matrix_t), intent(inout) :: state
    real(default) :: trace
    trace = state%trace ()
    if (trace /= 0) then
       state%me = state%me / trace
       state%norm = state%norm * trace
    end if
  end subroutine state_matrix_normalize_by_trace

\eatline
\nwendcode{}\nwbegindocs{708}Analogous, but renormalize by maximal (absolute) value.
\nwenddocs{}\nwbegincode{709}\moddef{State matrices: state matrix: TBP}\plusendmoddef
  procedure :: normalize_by_max => state_matrix_normalize_by_max
\nwendcode{}\nwbegincode{710}\moddef{State matrices: procedures}\plusendmoddef
  subroutine state_matrix_normalize_by_max (state)
    class(state_matrix_t), intent(inout) :: state
    real(default) :: m
    m = maxval (abs (state%me))
    if (m /= 0) then
       state%me = state%me / m
       state%norm = state%norm * m
    end if
  end subroutine state_matrix_normalize_by_max

\eatline
\nwendcode{}\nwbegindocs{711}Explicitly set the norm of a state matrix.
\nwenddocs{}\nwbegincode{712}\moddef{State matrices: state matrix: TBP}\plusendmoddef
  procedure :: set_norm => state_matrix_set_norm
\nwendcode{}\nwbegincode{713}\moddef{State matrices: procedures}\plusendmoddef
  subroutine state_matrix_set_norm (state, norm)
    class(state_matrix_t), intent(inout) :: state
    real(default), intent(in) :: norm
    state%norm = norm
  end subroutine state_matrix_set_norm
  
\eatline
\nwendcode{}\nwbegindocs{714}Return the sum of all matrix element values.
\nwenddocs{}\nwbegincode{715}\moddef{State matrices: state matrix: TBP}\plusendmoddef
  procedure :: sum => state_matrix_sum
\nwendcode{}\nwbegincode{716}\moddef{State matrices: procedures}\plusendmoddef
  function state_matrix_sum (state) result (value)
    complex(default) :: value
    class(state_matrix_t), intent(in) :: state
    value = sum (state%me)
  end function state_matrix_sum

\eatline
\nwendcode{}\nwbegindocs{717}Return the trace of a state matrix, i.e., the sum over all diagonal
values.  If {\tt{}qn{\char95}in} is provided, only branches that match this
quantum-numbers array are considered.
\nwenddocs{}\nwbegincode{718}\moddef{State matrices: state matrix: TBP}\plusendmoddef
  procedure :: trace => state_matrix_trace
\nwendcode{}\nwbegincode{719}\moddef{State matrices: procedures}\plusendmoddef
  function state_matrix_trace (state, qn_in) result (trace)
    complex(default) :: trace
    class(state_matrix_t), intent(in), target :: state
    type(quantum_numbers_t), dimension(:), intent(in), optional :: qn_in
    type(quantum_numbers_t), dimension(:), allocatable :: qn
    type(state_iterator_t) :: it
    allocate (qn (state%get_depth ()))
    trace = 0
    call it%init (state)
    do while (it%is_valid ())
       qn = it%get_quantum_numbers ()
       if (present (qn_in)) then
          if (.not. all (qn .match. qn_in)) then
             call it%advance ();  cycle
          end if
       end if
       if (all (qn%are_diagonal ())) then
          trace = trace + it%get_matrix_element ()
       end if
       call it%advance ()
    end do
  end function state_matrix_trace

\eatline
\nwendcode{}\nwbegindocs{720}Append new states which are color-contracted versions of the
existing states.  The matrix element index of each color contraction
coincides with the index of its origin, so no new matrix elements are
generated.  After this operation, no {\tt{}freeze} must be performed
anymore.
\nwenddocs{}\nwbegincode{721}\moddef{State matrices: state matrix: TBP}\plusendmoddef
  procedure :: add_color_contractions => state_matrix_add_color_contractions
\nwendcode{}\nwbegincode{722}\moddef{State matrices: procedures}\plusendmoddef
  subroutine state_matrix_add_color_contractions (state)
    class(state_matrix_t), intent(inout), target :: state
    type(state_iterator_t) :: it
    type(quantum_numbers_t), dimension(:,:), allocatable :: qn
    type(quantum_numbers_t), dimension(:,:), allocatable :: qn_con
    integer, dimension(:), allocatable :: me_index
    integer :: depth, n_me, i, j
    depth = state%get_depth ()
    n_me = state%get_n_matrix_elements ()
    allocate (qn (depth, n_me))
    allocate (me_index (n_me))
    i = 0
    call it%init (state)
    do while (it%is_valid ())
       i = i + 1
       qn(:,i) = it%get_quantum_numbers ()
       me_index(i) = it%get_me_index ()
       call it%advance ()
    end do
    do i = 1, n_me
       call quantum_number_array_make_color_contractions (qn(:,i), qn_con)
       do j = 1, size (qn_con, 2)
          call state%add_state (qn_con(:,j), index = me_index(i))
       end do
    end do
  end subroutine state_matrix_add_color_contractions

\eatline
\nwendcode{}\nwbegindocs{723}This procedure merges two state matrices of equal depth.  For each
quantum number (flavor, color, helicity), we take the entry from the
first argument where defined, otherwise the second one.  (If both are
defined, we get an off-diagonal matrix.) The resulting
trie combines the information of the input tries in all possible ways.
Note that values are ignored, all values in the result are zero.
\nwenddocs{}\nwbegincode{724}\moddef{State matrices: public}\plusendmoddef
  public :: merge_state_matrices
\nwendcode{}\nwbegincode{725}\moddef{State matrices: procedures}\plusendmoddef
  subroutine merge_state_matrices (state1, state2, state3)
    type(state_matrix_t), intent(in), target :: state1, state2
    type(state_matrix_t), intent(out) :: state3
    type(state_iterator_t) :: it1, it2
    type(quantum_numbers_t), dimension(state1%depth) :: qn1, qn2
    if (state1%depth /= state2%depth) then
       call state1%write ()
       call state2%write ()
       call msg_bug ("State matrices merge impossible: incompatible depths")
    end if
    call state3%init ()
    call it1%init (state1)
    do while (it1%is_valid ())
       qn1 = it1%get_quantum_numbers ()
       call it2%init (state2)
       do while (it2%is_valid ())
          qn2 = it2%get_quantum_numbers ()
          call state3%add_state (qn1 .merge. qn2)
          call it2%advance ()
       end do
       call it1%advance ()
    end do
    call state3%freeze ()
  end subroutine merge_state_matrices

\eatline
\nwendcode{}\nwbegindocs{726}Multiply matrix elements from two state matrices.  Choose the elements
as given by the integer index arrays, multiply them and store the sum
of products in the indicated matrix element.  The suffixes mean:
c=conjugate first factor; f=include weighting factor.

Note that the {\tt{}dot{\char95}product} intrinsic function conjugates its first
complex argument.  This is intended for the {\tt{}c} suffix case, but
must be reverted for the plain-product case.

We provide analogous subroutines for just summing over state matrix
entries.  The {\tt{}evaluate{\char95}sum} variant includes the state-matrix norm
in the evaluation, the {\tt{}evaluate{\char95}me{\char95}sum} takes into account just the
matrix elements proper.
\nwenddocs{}\nwbegincode{727}\moddef{State matrices: state matrix: TBP}\plusendmoddef
  procedure :: evaluate_product => state_matrix_evaluate_product
  procedure :: evaluate_product_cf => state_matrix_evaluate_product_cf
  procedure :: evaluate_square_c => state_matrix_evaluate_square_c
  procedure :: evaluate_sum => state_matrix_evaluate_sum
  procedure :: evaluate_me_sum => state_matrix_evaluate_me_sum
\nwendcode{}\nwbegincode{728}\moddef{State matrices: procedures}\plusendmoddef
  pure subroutine state_matrix_evaluate_product &
       (state, i, state1, state2, index1, index2)
    class(state_matrix_t), intent(inout) :: state
    integer, intent(in) :: i
    type(state_matrix_t), intent(in) :: state1, state2
    integer, dimension(:), intent(in) :: index1, index2
    state%me(i) = &
         dot_product (conjg (state1%me(index1)), state2%me(index2))
    state%norm = state1%norm * state2%norm
  end subroutine state_matrix_evaluate_product

  pure subroutine state_matrix_evaluate_product_cf &
       (state, i, state1, state2, index1, index2, factor)
    class(state_matrix_t), intent(inout) :: state
    integer, intent(in) :: i
    type(state_matrix_t), intent(in) :: state1, state2
    integer, dimension(:), intent(in) :: index1, index2
    complex(default), dimension(:), intent(in) :: factor
    state%me(i) = &
         dot_product (state1%me(index1), factor * state2%me(index2))
    state%norm = state1%norm * state2%norm
  end subroutine state_matrix_evaluate_product_cf

  pure subroutine state_matrix_evaluate_square_c (state, i, state1, index1)
    class(state_matrix_t), intent(inout) :: state
    integer, intent(in) :: i
    type(state_matrix_t), intent(in) :: state1
    integer, dimension(:), intent(in) :: index1
    state%me(i) = &
         dot_product (state1%me(index1), state1%me(index1))
    state%norm = abs (state1%norm) ** 2
  end subroutine state_matrix_evaluate_square_c

  pure subroutine state_matrix_evaluate_sum (state, i, state1, index1)
    class(state_matrix_t), intent(inout) :: state
    integer, intent(in) :: i
    type(state_matrix_t), intent(in) :: state1
    integer, dimension(:), intent(in) :: index1
    state%me(i) = &
         sum (state1%me(index1)) * state1%norm
  end subroutine state_matrix_evaluate_sum

  pure subroutine state_matrix_evaluate_me_sum (state, i, state1, index1)
    class(state_matrix_t), intent(inout) :: state
    integer, intent(in) :: i
    type(state_matrix_t), intent(in) :: state1
    integer, dimension(:), intent(in) :: index1
    state%me(i) = sum (state1%me(index1))
  end subroutine state_matrix_evaluate_me_sum

\eatline
\eatline
\eatline
\eatline
\eatline
\nwendcode{}\nwbegindocs{729}Outer product (of states and matrix elements):
\nwenddocs{}\nwbegincode{730}\moddef{State matrices: public}\plusendmoddef
  public :: outer_multiply
\nwendcode{}\nwbegincode{731}\moddef{State matrices: interfaces}\plusendmoddef
  interface outer_multiply
     module procedure outer_multiply_pair
     module procedure outer_multiply_array
  end interface

\eatline
\nwendcode{}\nwbegindocs{732}This procedure constructs the outer product of two state matrices.
\nwenddocs{}\nwbegincode{733}\moddef{State matrices: procedures}\plusendmoddef
  subroutine outer_multiply_pair (state1, state2, state3)
    type(state_matrix_t), intent(in), target :: state1, state2
    type(state_matrix_t), intent(out) :: state3
    type(state_iterator_t) :: it1, it2
    type(quantum_numbers_t), dimension(state1%depth) :: qn1
    type(quantum_numbers_t), dimension(state2%depth) :: qn2
    type(quantum_numbers_t), dimension(state1%depth+state2%depth) :: qn3
    complex(default) :: val1, val2
    call state3%init (store_values=.true.)
    call it1%init (state1)
    do while (it1%is_valid ())
       qn1 = it1%get_quantum_numbers ()
       val1 = it1%get_matrix_element ()
       call it2%init (state2)
       do while (it2%is_valid ())
          qn2 = it2%get_quantum_numbers ()
          val2 = it2%get_matrix_element ()
          qn3(:state1%depth) = qn1
          qn3(state1%depth+1:) = qn2
          call state3%add_state (qn3, value=val1 * val2)
          call it2%advance ()
       end do
       call it1%advance ()
    end do
    call state3%freeze ()
  end subroutine outer_multiply_pair

\eatline
\nwendcode{}\nwbegindocs{734}This executes the above routine iteratively for an arbitrary number
of state matrices.
\nwenddocs{}\nwbegincode{735}\moddef{State matrices: procedures}\plusendmoddef
  subroutine outer_multiply_array (state_in, state_out)
    type(state_matrix_t), dimension(:), intent(in), target :: state_in
    type(state_matrix_t), intent(out) :: state_out
    type(state_matrix_t), dimension(:), allocatable, target :: state_tmp
    integer :: i, n
    n = size (state_in)
    select case (n)
    case (0)
       call state_out%init ()
    case (1)
       state_out = state_in(1)
    case (2)
       call outer_multiply_pair (state_in(1), state_in(2), state_out)
    case default
       allocate (state_tmp (n-2))
       call outer_multiply_pair (state_in(1), state_in(2), state_tmp(1))
       do i = 2, n - 2
          call outer_multiply_pair (state_tmp(i-1), state_in(i+1), state_tmp(i))
       end do
       call outer_multiply_pair (state_tmp(n-2), state_in(n), state_out)
       do i = 1, size(state_tmp)
          call state_tmp(i)%final ()
       end do
    end select
  end subroutine outer_multiply_array

\eatline
\eatline
\nwendcode{}\nwbegindocs{736}\nwdocspar
\subsection{Factorization}
In physical events, the state matrix is factorized into
single-particle state matrices.  This is essentially a measurement.

In a simulation, we select one particular branch of the state matrix
with a probability that is determined by the matrix elements at the
leaves.  (This makes sense only if the state matrix represents a
squared amplitude.)  The selection is based on a (random) value {\tt{}x}
between 0 and one that is provided as the third argument.

For flavor and color, we select a unique value for each particle.  For
polarization, we have three options (modes).  Option 1 is to drop
helicity information altogether and sum over all diagonal helicities.
Option 2 is to select a unique diagonal helicity in the same way as
flavor and color.  Option 3 is, for each particle, to trace over all
remaining helicities in order to obtain an array of independent
single-particle helicity matrices.

Only branches that match the given quantum-number array {\tt{}qn{\char95}in}, if
present, are considered.  For this array, color is ignored.

If the optional {\tt{}correlated{\char95}state} is provided, it is assigned the
correlated density matrix for the selected flavor-color branch, so
multi-particle spin correlations remain available even if they are
dropped in the single-particle density matrices.

The algorithm is as follows: First, we determine the normalization by
summing over all diagonal matrix elements.  In a second scan, we
select one of the diagonal matrix elements by a cumulative comparison
with the normalized random number.  In the corresponding quantum
number array, we undefine the helicity entries.  Then, we scan the
third time.  For each branch that matches the selected quantum number
array (i.e., definite flavor and color, arbitrary helicity), we
determine its contribution to any of the single-particle state
matrices.  The matrix-element value is added if all other quantum
numbers are diagonal, while the helicity of the chosen particle may be
arbitrary; this helicity determines the branch in the single-particle
state.  

As a result, flavor and color quantum numbers are selected with the
correct probability.  Within this subset of states, each
single-particle state matrix results from tracing over all other
particles.  Note that the single-particle state matrices are not
normalized.

The flag {\tt{}ok} is set to false if the matrix element sum is zero, so
factorization is not possible.  This can happen if an event did not pass
cuts.
\nwenddocs{}\nwbegincode{737}\moddef{State matrices: parameters}\endmoddef
  integer, parameter, public :: FM_IGNORE_HELICITY = 1
  integer, parameter, public :: FM_SELECT_HELICITY = 2
  integer, parameter, public :: FM_FACTOR_HELICITY = 3

\eatline
\nwendcode{}\nwbegincode{738}\moddef{State matrices: state matrix: TBP}\plusendmoddef
  procedure :: factorize => state_matrix_factorize
\nwendcode{}\nwbegincode{739}\moddef{State matrices: procedures}\plusendmoddef
  subroutine state_matrix_factorize &
       (state, mode, x, ok, single_state, correlated_state, qn_in)
    class(state_matrix_t), intent(in), target :: state
    integer, intent(in) :: mode
    real(default), intent(in) :: x
    logical, intent(out) :: ok
    type(state_matrix_t), &
         dimension(:), allocatable, intent(out) :: single_state
    type(state_matrix_t), intent(out), optional :: correlated_state
    type(quantum_numbers_t), dimension(:), intent(in), optional :: qn_in
    type(state_iterator_t) :: it
    real(default) :: s, xt
    complex(default) :: value
    integer :: i, depth
    type(quantum_numbers_t), dimension(:), allocatable :: qn, qn1
    type(quantum_numbers_mask_t), dimension(:), allocatable :: qn_mask
    logical, dimension(:), allocatable :: diagonal
    logical, dimension(:,:), allocatable :: mask
    ok = .true.
    if (x /= 0) then
       xt = x * state%trace (qn_in)
    else
       xt = 0
    end if
    s = 0
    depth = state%get_depth ()
    allocate (qn (depth), qn1 (depth), diagonal (depth))
    call it%init (state)
    do while (it%is_valid ())
       qn = it%get_quantum_numbers ()
       if (present (qn_in)) then
          if (.not. all (qn .fhmatch. qn_in)) then
             call it%advance (); cycle
          end if
       end if
       if (all (qn%are_diagonal ())) then
          value = it%get_matrix_element ()
          if (real (value, default) < 0) then
             call state%write ()
             print *, value
             call msg_bug ("Event generation: " &
                  // "Negative real part of squared matrix element value")
             value = 0
          end if
          s = s + value
          if (s > xt)  exit
       end if
       call it%advance ()
    end do
    if (.not. it%is_valid ()) then
       if (s == 0)  ok = .false.
       call it%init (state)
    end if
    allocate (single_state (depth))
    do i = 1, depth
       call single_state(i)%init (store_values=.true.)
    end do
    if (present (correlated_state)) &
         call correlated_state%init (store_values=.true.)
    qn = it%get_quantum_numbers ()
    select case (mode)
    case (FM_SELECT_HELICITY)  ! single branch selected; shortcut
       do i = 1, depth
          call single_state(i)%add_state ([qn(i)], value=value)
       end do
       if (.not. present (correlated_state)) then       
          do i = 1, size(single_state)
             call single_state(i)%freeze ()
          end do
          return
       end if
    end select
    allocate (qn_mask (depth))
    call qn_mask%init (.false., .false., .false., .true.)
    call qn%undefine (qn_mask)
    select case (mode)
    case (FM_FACTOR_HELICITY)
       allocate (mask (depth, depth))
       mask = .false.
       forall (i = 1:depth)  mask(i,i) = .true.
    end select
    call it%init (state)
    do while (it%is_valid ())
       qn1 = it%get_quantum_numbers ()
       if (all (qn .match. qn1)) then
          diagonal = qn1%are_diagonal ()
          value = it%get_matrix_element ()
          select case (mode)
          case (FM_IGNORE_HELICITY)  ! trace over diagonal states that match qn
             if (all (diagonal)) then
                do i = 1, depth
                   call single_state(i)%add_state &
                        ([qn(i)], value=value, sum_values=.true.)
                end do
             end if
          case (FM_FACTOR_HELICITY)  ! trace over all other particles
             do i = 1, depth
                if (all (diagonal .or. mask(:,i))) then
                   call single_state(i)%add_state &
                        ([qn1(i)], value=value, sum_values=.true.)
                end if
             end do
          end select
          if (present (correlated_state)) &
               call correlated_state%add_state (qn1, value=value)
       end if
       call it%advance ()
    end do
    do i = 1, depth
       call single_state(i)%freeze ()
    end do
    if (present (correlated_state)) &
         call correlated_state%freeze ()
  end subroutine state_matrix_factorize

\eatline
\nwendcode{}\nwbegindocs{740}\nwdocspar
\subsubsection{Quantum-number matching}
This feature allows us to check whether a given string of PDG values
matches, in any ordering, any of the flavor combinations that the
state matrix provides.  We will also request the permutation of the
successful match.

This type provides an account of the state's flavor content.  We store
all flavor combinations, as {\tt{}pdg} values, in an array, assuming that
the length is uniform.  

We check only the entries selected by {\tt{}mask{\char95}match}.  Among those,
only the entries selected by {\tt{}mask{\char95}sort} are sorted and thus matched
without respecting array element order. The entries that correspond to
a true value in the associated {\tt{}mask} are sorted.  The mapping from
the original state to the sorted state is given by the index array
{\tt{}map}.
\nwenddocs{}\nwbegincode{741}\moddef{State matrices: public}\plusendmoddef
  public :: state_flv_content_t
\nwendcode{}\nwbegincode{742}\moddef{State matrices: types}\plusendmoddef
  type :: state_flv_content_t
     private
     integer, dimension(:,:), allocatable :: pdg
     integer, dimension(:,:), allocatable :: map
     logical, dimension(:), allocatable :: mask
   contains
   \LA{}State matrices: state flv content: TBP\RA{}
  end type state_flv_content_t
  
\eatline
\nwendcode{}\nwbegindocs{743}Output (debugging aid).
\nwenddocs{}\nwbegincode{744}\moddef{State matrices: state flv content: TBP}\endmoddef
  procedure :: write => state_flv_content_write
\nwendcode{}\nwbegincode{745}\moddef{State matrices: procedures}\plusendmoddef
  subroutine state_flv_content_write (state_flv, unit)
    class(state_flv_content_t), intent(in), target :: state_flv
    integer, intent(in), optional :: unit
    type(state_iterator_t) :: it
    integer :: u, n, d, i, j
    u = given_output_unit (unit)
    d = size (state_flv%pdg, 1)
    n = size (state_flv%pdg, 2)
    do i = 1, n
       write (u, "(2x,'PDG =')", advance="no")
       do j = 1, d
          write (u, "(1x,I0)", advance="no")  state_flv%pdg(j,i)
       end do
       write (u, "(' :: map = (')", advance="no")
       do j = 1, d
          write (u, "(1x,I0)", advance="no")  state_flv%map(j,i)
       end do
       write (u, "(' )')")
    end do
  end subroutine state_flv_content_write
    
\eatline
\nwendcode{}\nwbegindocs{746}Initialize with table length and mask.  Each row of the {\tt{}map}
array, of length $d$, is initialized with $(0,1,\ldots,d)$.
\nwenddocs{}\nwbegincode{747}\moddef{State matrices: state flv content: TBP}\plusendmoddef
  procedure :: init => state_flv_content_init
\nwendcode{}\nwbegincode{748}\moddef{State matrices: procedures}\plusendmoddef
  subroutine state_flv_content_init (state_flv, n, mask)
    class(state_flv_content_t), intent(out) :: state_flv
    integer, intent(in) :: n
    logical, dimension(:), intent(in) :: mask
    integer :: d, i
    d = size (mask)
    allocate (state_flv%pdg (d, n), source = 0)
    allocate (state_flv%map (d, n), source = spread ([(i, i = 1, d)], 2, n))
    allocate (state_flv%mask (d), source = mask)
  end subroutine state_flv_content_init
    
\eatline
\nwendcode{}\nwbegindocs{749}Manually fill the entries, one flavor set and mapping at a time.
\nwenddocs{}\nwbegincode{750}\moddef{State matrices: state flv content: TBP}\plusendmoddef
  procedure :: set_entry => state_flv_content_set_entry
\nwendcode{}\nwbegincode{751}\moddef{State matrices: procedures}\plusendmoddef
  subroutine state_flv_content_set_entry (state_flv, i, pdg, map)
    class(state_flv_content_t), intent(inout) :: state_flv
    integer, intent(in) :: i
    integer, dimension(:), intent(in) :: pdg, map
    state_flv%pdg(:,i) = pdg
    where (map /= 0)
       state_flv%map(:,i) = map
    end where
  end subroutine state_flv_content_set_entry
  
\eatline
\nwendcode{}\nwbegindocs{752}Given a state matrix, determine the flavor content.  That is, scan
the state matrix and extract flavor only, build a new state matrix
from that.
\nwenddocs{}\nwbegincode{753}\moddef{State matrices: state flv content: TBP}\plusendmoddef
  procedure :: fill => state_flv_content_fill
\nwendcode{}\nwbegincode{754}\moddef{State matrices: procedures}\plusendmoddef
  subroutine state_flv_content_fill &
       (state_flv, state_full, mask)
    class(state_flv_content_t), intent(out) :: state_flv
    type(state_matrix_t), intent(in), target :: state_full
    logical, dimension(:), intent(in) :: mask
    type(state_matrix_t), target :: state_tmp
    type(state_iterator_t) :: it
    type(flavor_t), dimension(:), allocatable :: flv
    integer, dimension(:), allocatable :: pdg, pdg_subset
    integer, dimension(:), allocatable :: idx, map_subset, idx_subset, map
    type(quantum_numbers_t), dimension(:), allocatable :: qn
    integer :: n, d, c, i
    call state_tmp%init ()
    d = state_full%get_depth ()
    allocate (flv (d), qn (d), pdg (d), idx (d), map (d))
    idx = [(i, i = 1, d)]
    c = count (mask)
    allocate (pdg_subset (c), map_subset (c), idx_subset (c))
    call it%init (state_full)
    do while (it%is_valid ())
       flv = it%get_flavor ()
       call qn%init (flv)
       call state_tmp%add_state (qn)
       call it%advance ()
    end do
    n = state_tmp%get_n_leaves ()
    call state_flv%init (n, mask)
    i = 0
    call it%init (state_tmp)
    do while (it%is_valid ())
       i = i + 1
       flv = it%get_flavor ()
       pdg = flv%get_pdg ()
       idx_subset = pack (idx, mask)
       pdg_subset = pack (pdg, mask)
       map_subset = order_abs (pdg_subset)
       map = unpack (idx_subset (map_subset), mask, idx)
       call state_flv%set_entry (i, &
            unpack (pdg_subset(map_subset), mask, pdg), &
            order (map))
       call it%advance ()
    end do
    call state_tmp%final ()
  end subroutine state_flv_content_fill

\eatline
\nwendcode{}\nwbegindocs{755}Match a given flavor string against the flavor content.  We sort the
input string and check whether it matches any of the stored strings.
If yes, return the mapping.

Only PDG entries under the preset mask are sorted before matching.  The
other entries must match exactly (i.e., without reordering).  A zero
entry matches anything.  In any case, the length of the PDG string
must be equal to the length $d$ of the individual flavor-state entries.
\nwenddocs{}\nwbegincode{756}\moddef{State matrices: state flv content: TBP}\plusendmoddef
  procedure :: match => state_flv_content_match
\nwendcode{}\nwbegincode{757}\moddef{State matrices: procedures}\plusendmoddef
  subroutine state_flv_content_match (state_flv, pdg, success, map)
    class(state_flv_content_t), intent(in) :: state_flv
    integer, dimension(:), intent(in) :: pdg
    logical, intent(out) :: success
    integer, dimension(:), intent(out) :: map
    integer, dimension(:), allocatable :: pdg_subset, pdg_sorted, map1, map2
    integer, dimension(:), allocatable :: idx, map_subset, idx_subset
    integer :: i, n, c, d
    c = count (state_flv%mask)
    d = size (state_flv%pdg, 1)
    n = size (state_flv%pdg, 2)
    allocate (idx (d), source = [(i, i = 1, d)])
    allocate (idx_subset (c), pdg_subset (c), map_subset (c))
    allocate (pdg_sorted (d), map1 (d), map2 (d))
    idx_subset = pack (idx, state_flv%mask)
    pdg_subset = pack (pdg, state_flv%mask)
    map_subset = order_abs (pdg_subset)
    pdg_sorted = unpack (pdg_subset(map_subset), state_flv%mask, pdg)
    success = .false.
    do i = 1, n
       if (all (pdg_sorted == state_flv%pdg(:,i) &
            .or. pdg_sorted == 0)) then
          success = .true.
          exit
       end if
    end do
    if (success) then
       map1 = state_flv%map(:,i)
       map2 = unpack (idx_subset(map_subset), state_flv%mask, idx)
       map = map2(map1)
       where (pdg == 0)  map = 0
    end if
  end subroutine state_flv_content_match
    
\eatline
\nwendcode{}\nwbegindocs{758}\nwdocspar
\nwenddocs{}\nwbegincode{759}\moddef{State matrices: procedures}\plusendmoddef
  elemental function pacify_complex (c_in) result (c_pac)
    complex(default), intent(in) :: c_in
    complex(default) :: c_pac
    c_pac = c_in
    if (real(c_pac) == -real(c_pac)) then
       c_pac = &
            cmplx (0._default, aimag(c_pac), kind=default)
    end if
    if (aimag(c_pac) == -aimag(c_pac)) then
       c_pac = &
            cmplx (real(c_pac), 0._default, kind=default)
    end if
  end function pacify_complex
  
\eatline
\nwendcode{}\nwbegindocs{760}\nwdocspar
\subsection{Unit tests}
\nwenddocs{}\nwbegincode{761}\moddef{State matrices: public}\plusendmoddef
  public :: state_matrix_test
\nwendcode{}\nwbegincode{762}\moddef{State matrices: procedures}\plusendmoddef
  subroutine state_matrix_test (u, results)
    integer, intent(in) :: u
    type(test_results_t), intent(inout) :: results
  \LA{}State matrices: execute tests\RA{}  
  end subroutine state_matrix_test 
  
\eatline
\nwendcode{}\nwbegindocs{763}Create two quantum states of equal depth and merge them.
\nwenddocs{}\nwbegincode{764}\moddef{State matrices: execute tests}\endmoddef
  call test (state_matrix_1, "state_matrix_1", &
       "check merge of quantum states of equal depth", &
       u, results)
\nwendcode{}\nwbegincode{765}\moddef{State matrices: tests}\endmoddef
  subroutine state_matrix_1 (u)
    integer, intent(in) :: u
    type(state_matrix_t) :: state1, state2, state3
    type(flavor_t), dimension(3) :: flv
    type(color_t), dimension(3) :: col
    type(quantum_numbers_t), dimension(3) :: qn
    
    write (u, "(A)")  "* Test output: state_matrix_1"
    write (u, "(A)")  "*   Purpose: create and merge two quantum states"
    write (u, "(A)")

    write (u, "(A)")  "*  Initialization"
    write (u, "(A)")    
        
    write (u, "(A)")  "*  State matrix 1"
    write (u, "(A)")        
    
    call state1%init ()
    call flv%init ([1, 2, 11])
    call qn%init (flv, helicity ([ 1, 1, 1]))
    call state1%add_state (qn)
    call qn%init (flv, helicity ([ 1, 1, 1], [-1, 1, -1]))
    call state1%add_state (qn)
    call state1%freeze ()
    call state1%write (u)

    write (u, "(A)")
    write (u, "(A)")  "*  State matrix 2"
    write (u, "(A)")

    call state2%init ()
    call col(1)%init ([501])
    call col(2)%init ([-501])
    call col(3)%init ([0])
    call qn%init (col, helicity ([-1, -1, 0]))
    call state2%add_state (qn)
    call col(3)%init ([99])
    call qn%init (col, helicity ([-1, -1, 0]))
    call state2%add_state (qn)
    call state2%freeze ()
    call state2%write (u)

    write (u, "(A)")
    write (u, "(A)")  "* Merge the state matrices"   
    write (u, "(A)")
        
    call merge_state_matrices (state1, state2, state3)
    call state3%write (u)

    write (u, "(A)")
    write (u, "(A)")  "* Collapse the state matrix"    
    write (u, "(A)")
    
    call state3%collapse (quantum_numbers_mask (.false., .false., &
         [.true.,.false.,.false.]))
    call state3%write (u)
    
    write (u, "(A)")
    write (u, "(A)")  "* Cleanup"    
    write (u, "(A)")    
    
    call state1%final ()
    call state2%final ()
    call state3%final ()
    
    write (u, "(A)")
    write (u, "(A)")  "* Test output end: state_matrix_1"    
    write (u, "(A)")    
    
  end subroutine state_matrix_1

\eatline
\nwendcode{}\nwbegindocs{766}Create a correlated three-particle state matrix and factorize it.
\nwenddocs{}\nwbegincode{767}\moddef{State matrices: execute tests}\plusendmoddef
  call test (state_matrix_2, "state_matrix_2", &
       "check factorizing 3-particle state matrix", &
       u, results)
\nwendcode{}\nwbegincode{768}\moddef{State matrices: tests}\plusendmoddef
  subroutine state_matrix_2 (u)
    integer, intent(in) :: u
    type(state_matrix_t) :: state
    type(state_matrix_t), dimension(:), allocatable :: single_state
    type(state_matrix_t) :: correlated_state
    complex(default) :: z, val
    complex(default), dimension(-1:1) :: v
    integer :: f, h11, h12, h21, h22, i, mode
    type(flavor_t), dimension(2) :: flv
    type(color_t), dimension(2) :: col
    type(helicity_t), dimension(2) :: hel
    type(quantum_numbers_t), dimension(2) :: qn
    logical :: ok
    
    write (u, "(A)")
    write (u, "(A)")  "* Test output: state_matrix_2"
    write (u, "(A)")  "*   Purpose: factorize correlated 3-particle state"
    write (u, "(A)")        
    
    write (u, "(A)")  "*  Initialization"
    write (u, "(A)")    
        
    z = 1 / 2._default
    v(-1) = (0.6_default, 0._default)
    v( 1) = (0._default, 0.8_default)
    call state%init ()
    do f = 1, 2
       do h11 = -1, 1, 2
          do h12 = -1, 1, 2
             do h21 = -1, 1, 2
                do h22 = -1, 1, 2
                   call flv%init ([f, -f])
                   call col(1)%init ([1])
                   call col(2)%init ([-1])
                   call hel%init ([h11,h12], [h21, h22])
                   call qn%init (flv, col, hel)
                   val = z * v(h11) * v(h12) * conjg (v(h21) * v(h22))
                   call state%add_state (qn)
                end do
             end do
          end do
       end do
    end do
    call state%freeze ()
    call state%write (u)

    write (u, "(A)")
    write (u, "(A,'('," // FMT_19 // ",','," // FMT_19 // ",')')") &
         "* Trace = ", state%trace ()
    write (u, "(A)")
    
    do mode = 1, 3
       write (u, "(A)")
       write (u, "(A,I1)")  "* Mode = ", mode
       call state%factorize &
            (mode, 0.15_default, ok, single_state, correlated_state)
       do i = 1, size (single_state)
          write (u, "(A)")
          call single_state(i)%write (u)
          write (u, "(A,'('," // FMT_19 // ",','," // FMT_19 // ",')')") &
               "Trace = ", single_state(i)%trace ()
       end do
       write (u, "(A)")
       call correlated_state%write (u)
       write (u, "(A,'('," // FMT_19 // ",','," // FMT_19 // ",')')")  &
            "Trace = ", correlated_state%trace ()
       do i = 1, size(single_state)
          call single_state(i)%final ()
       end do
       call correlated_state%final ()
    end do
    
    write (u, "(A)")
    write (u, "(A)")  "* Cleanup"
    
    call state%final ()
    
    write (u, "(A)")
    write (u, "(A)")  "* Test output end: state_matrix_2"
    
  end subroutine state_matrix_2

\eatline
\nwendcode{}\nwbegindocs{769}Create a colored state matrix and add color contractions.
\nwenddocs{}\nwbegincode{770}\moddef{State matrices: execute tests}\plusendmoddef
  call test (state_matrix_3, "state_matrix_3", &
       "check factorizing 3-particle state matrix", &
       u, results)
\nwendcode{}\nwbegincode{771}\moddef{State matrices: tests}\plusendmoddef
  subroutine state_matrix_3 (u)
    use physics_defs, only: HADRON_REMNANT_TRIPLET, HADRON_REMNANT_OCTET
    integer, intent(in) :: u
    type(state_matrix_t) :: state
    type(flavor_t), dimension(4) :: flv
    type(color_t), dimension(4) :: col
    type(quantum_numbers_t), dimension(4) :: qn
    
    write (u, "(A)")  "* Test output: state_matrix_3"
    write (u, "(A)")  "*   Purpose: add color connections to colored state"
    write (u, "(A)")    
       
    write (u, "(A)")  "*  Initialization"
    write (u, "(A)")    
    
    call state%init ()
    call flv%init ([ 1, -HADRON_REMNANT_TRIPLET, -1, HADRON_REMNANT_TRIPLET ])
    call col(1)%init ([17])
    call col(2)%init ([-17])
    call col(3)%init ([-19])
    call col(4)%init ([19])
    call qn%init (flv, col)
    call state%add_state (qn)
    call flv%init ([ 1, -HADRON_REMNANT_TRIPLET, 21, HADRON_REMNANT_OCTET ])
    call col(1)%init ([17])
    call col(2)%init ([-17])
    call col(3)%init ([3, -5])
    call col(4)%init ([5, -3])
    call qn%init (flv, col)
    call state%add_state (qn)
    call state%freeze ()

    write (u, "(A)") "* State:"
    write (u, "(A)") 
    
    call state%write (u)
    call state%add_color_contractions ()

    write (u, "(A)") "* State with contractions:"
    write (u, "(A)")
    
    call state%write (u)
    
    write (u, "(A)")
    write (u, "(A)")  "* Cleanup"
        
    call state%final ()
    
    write (u, "(A)")
    write (u, "(A)")  "* Test output end: state_matrx_3"    
    
  end subroutine state_matrix_3

\eatline
\nwendcode{}\nwbegindocs{772}Create a correlated three-particle state matrix, write it to file
and read again.
\nwenddocs{}\nwbegincode{773}\moddef{State matrices: execute tests}\plusendmoddef
  call test (state_matrix_4, "state_matrix_4", &
       "check raw I/O", &
       u, results)
\nwendcode{}\nwbegincode{774}\moddef{State matrices: tests}\plusendmoddef
  subroutine state_matrix_4 (u)
    integer, intent(in) :: u
    type(state_matrix_t), allocatable :: state
    complex(default) :: z, val
    complex(default), dimension(-1:1) :: v
    integer :: f, h11, h12, h21, h22, i, mode
    type(flavor_t), dimension(2) :: flv
    type(color_t), dimension(2) :: col
    type(helicity_t), dimension(2) :: hel
    type(quantum_numbers_t), dimension(2) :: qn
    integer :: unit, iostat
    
    write (u, "(A)")
    write (u, "(A)")  "* Test output: state_matrix_4"
    write (u, "(A)")  "*   Purpose: raw I/O for correlated 3-particle state"
    write (u, "(A)")        
    
    write (u, "(A)")  "*  Initialization"
    write (u, "(A)")    
        
    allocate (state)

    z = 1 / 2._default
    v(-1) = (0.6_default, 0._default)
    v( 1) = (0._default, 0.8_default)
    call state%init ()
    do f = 1, 2
       do h11 = -1, 1, 2
          do h12 = -1, 1, 2
             do h21 = -1, 1, 2
                do h22 = -1, 1, 2
                   call flv%init ([f, -f])
                   call col(1)%init ([1])
                   call col(2)%init ([-1])
                   call hel%init ([h11,h12], [h21, h22])
                   call qn%init (flv, col, hel)
                   val = z * v(h11) * v(h12) * conjg (v(h21) * v(h22))
                   call state%add_state (qn)
                end do
             end do
          end do
       end do
    end do
    call state%freeze ()

    call state%set_norm (3._default)
    do i = 1, state%get_n_leaves ()
       call state%set_matrix_element (i, cmplx (2 * i, 2 * i + 1, default))
    end do
    
    call state%write (u)
    
    write (u, "(A)")
    write (u, "(A)")  "* Write to file and read again "
    write (u, "(A)")
    
    unit = free_unit ()
    open (unit, action="readwrite", form="unformatted", status="scratch")
    call state%write_raw (unit)
    call state%final ()
    deallocate (state)
    
    allocate(state)
    rewind (unit)
    call state%read_raw (unit, iostat=iostat)
    close (unit)
    
    call state%write (u)

    write (u, "(A)")
    write (u, "(A)")  "* Cleanup"
    
    call state%final ()
    deallocate (state)
    
    write (u, "(A)")
    write (u, "(A)")  "* Test output end: state_matrix_4"
    
  end subroutine state_matrix_4

\eatline
\nwendcode{}\nwbegindocs{775}\nwdocspar
Create a flavor-content object for a given state matrix and match it
against trial flavor (i.e., PDG) strings.
\nwenddocs{}\nwbegincode{776}\moddef{State matrices: execute tests}\plusendmoddef
  call test (state_matrix_5, "state_matrix_5", &
       "check flavor content", &
       u, results)
\nwendcode{}\nwbegincode{777}\moddef{State matrices: tests}\plusendmoddef
  subroutine state_matrix_5 (u)
    integer, intent(in) :: u
    type(state_matrix_t), allocatable, target :: state
    type(state_iterator_t) :: it
    type(state_flv_content_t), allocatable :: state_flv
    type(flavor_t), dimension(4) :: flv1, flv2, flv3, flv4
    type(color_t), dimension(4) :: col1, col2
    type(helicity_t), dimension(4) :: hel1, hel2, hel3
    type(quantum_numbers_t), dimension(4) :: qn
    logical, dimension(4) :: mask
    
    write (u, "(A)")  "* Test output: state_matrix_5"
    write (u, "(A)")  "*   Purpose: check flavor-content state"
    write (u, "(A)")        
    
    write (u, "(A)")  "* Set up arbitrary state matrix"
    write (u, "(A)")    
    
    call flv1%init ([1, 4, 2, 7])
    call flv2%init ([1, 3,-3, 8])
    call flv3%init ([5, 6, 3, 7])
    call flv4%init ([6, 3, 5, 8])
    call hel1%init ([0, 1, -1, 0])
    call hel2%init ([0, 1, 1, 1])
    call hel3%init ([1, 0, 0, 0])
    call col1(1)%init ([0])
    call col1(2)%init ([0])
    call col1(3)%init ([0])
    call col1(4)%init ([0])
    call col2(1)%init ([5, -6])
    call col2(2)%init ([0])
    call col2(3)%init ([6, -5])
    call col2(4)%init ([0])

    allocate (state)
    call state%init ()
    call qn%init (flv1, col1, hel1)
    call state%add_state (qn)
    call qn%init (flv1, col1, hel2)
    call state%add_state (qn)
    call qn%init (flv3, col1, hel3)
    call state%add_state (qn)
    call qn%init (flv4, col1, hel3)
    call state%add_state (qn)
    call qn%init (flv1, col2, hel3)
    call state%add_state (qn)
    call qn%init (flv2, col2, hel2)
    call state%add_state (qn)
    call qn%init (flv2, col2, hel1)
    call state%add_state (qn)
    call qn%init (flv2, col1, hel1)
    call state%add_state (qn)
    call qn%init (flv3, col1, hel1)
    call state%add_state (qn)
    call qn%init (flv3, col2, hel3)
    call state%add_state (qn)
    call qn%init (flv1, col1, hel1)
    call state%add_state (qn)
    
    write (u, "(A)")  "* Quantum number content"
    write (u, "(A)")
    
    call it%init (state)
    do while (it%is_valid ())
       call quantum_numbers_write (it%get_quantum_numbers (), u)
       write (u, *)
       call it%advance ()
    end do
    
    write (u, "(A)")    
    write (u, "(A)")  "* Extract the flavor content"
    write (u, "(A)")
    
    mask = [.true., .true., .true., .false.]

    allocate (state_flv)
    call state_flv%fill (state, mask)
    call state_flv%write (u)

    write (u, "(A)")    
    write (u, "(A)")  "* Match trial sets"
    write (u, "(A)")
   
    call check ([1, 2, 3, 0])
    call check ([1, 4, 2, 0])
    call check ([4, 2, 1, 0])
    call check ([1, 3, -3, 0])
    call check ([1, -3, 3, 0])
    call check ([6, 3, 5, 0])

    write (u, "(A)")    
    write (u, "(A)")  "* Determine the flavor content with mask"
    write (u, "(A)")
    
    mask = [.false., .true., .true., .false.]

    call state_flv%fill (state, mask)
    call state_flv%write (u)
    
    write (u, "(A)")    
    write (u, "(A)")  "* Match trial sets"
    write (u, "(A)")
   
    call check ([1, 2, 3, 0])
    call check ([1, 4, 2, 0])
    call check ([4, 2, 1, 0])
    call check ([1, 3, -3, 0])
    call check ([1, -3, 3, 0])
    call check ([6, 3, 5, 0])

    write (u, "(A)")
    write (u, "(A)")  "* Cleanup"
    
    deallocate (state_flv)
    
    call state%final ()
    deallocate (state)
    
    write (u, "(A)")
    write (u, "(A)")  "* Test output end: state_matrix_5"
    
  contains

    subroutine check (pdg)
      integer, dimension(4), intent(in) :: pdg
      integer, dimension(4) :: map
      logical :: success
      call state_flv%match (pdg, success, map)
      write (u, "(2x,4(1x,I0),':',1x,L1)", advance="no")  pdg, success
      if (success) then
         write (u, "(2x,'map = (',4(1x,I0),' )')")  map
      else
         write (u, *)
      end if
    end subroutine check

  end subroutine state_matrix_5

\eatline
\nwendcode{}\nwbegindocs{778}\nwdocspar
%%%%%%%%%%%%%%%%%%%%%%%%%%%%%%%%%%%%%%%%%%%%%%%%%%%%%%%%%%%%%%%%%%%%%%%%
\section{Interactions}
This module defines the {\tt{}interaction{\char95}t} type.  It is an extension of
the {\tt{}state{\char95}matrix{\char95}t} type.

The state matrix is a representation of a multi-particle density
matrix.  It implements all possible flavor, color, and quantum-number
assignments of the entries in a generic density matrix, and it can
hold a complex matrix element for each entry.  (Note that this matrix
can hold non-diagonal entries in color and helicity space.)  The
{\tt{}interaction{\char95}t} object associates this with a list of momenta, such
that the whole object represents a multi-particle state.

The {\tt{}interaction{\char95}t} holds information about which particles are
incoming, virtual (i.e., kept for the records), or outgoing.  Each
particle can be associated to a source within another interaction.
This allows us to automatically fill those interaction momenta which
have been computed or defined elsewhere.  It also contains internal
parent-child relations and flags for (virtual) particles which are to
be treated as resonances.

A quantum-number mask array summarizes, for each particle within the
interaction, the treatment of flavor, color, or helicity (expose or
ignore).  A list of locks states which particles are bound to have an
identical quantum-number mask.  This is useful when the mask is
changed at one place.
\nwenddocs{}\nwbegincode{779}\moddef{\code{}interactions.f90\edoc{}}\endmoddef
\LA{}File header\RA{}

module interactions

\LA{}Use kinds\RA{}
  use io_units
  use unit_tests
  use diagnostics
  use sorting
  use lorentz
  use flavors
  use colors
  use helicities
  use quantum_numbers
  use state_matrices

\LA{}Standard module head\RA{}

\LA{}Interactions: public\RA{}

\LA{}Interactions: types\RA{}

\LA{}Interactions: interfaces\RA{}

contains

\LA{}Interactions: procedures\RA{}

\LA{}Interactions: tests\RA{}

end module interactions
\eatline
\nwendcode{}\nwbegindocs{780}\nwdocspar
\subsection{External interaction links}
Each particle in an interaction can have a link to a corresponding
particle in another interaction.  This allows to fetch the momenta of
incoming or virtual particles from the interaction where they are
defined.  The link object consists of a pointer to the interaction and
an index.  
\nwenddocs{}\nwbegincode{781}\moddef{Interactions: types}\endmoddef
  type :: external_link_t
     private
     type(interaction_t), pointer :: int => null ()
     integer :: i
  end type external_link_t

\eatline
\nwendcode{}\nwbegindocs{782}Set an external link. 
\nwenddocs{}\nwbegincode{783}\moddef{Interactions: procedures}\endmoddef
  subroutine external_link_set (link, int, i)
    type(external_link_t), intent(out) :: link
    type(interaction_t), target, intent(in) :: int
    integer, intent(in) :: i
    if (i /= 0) then
       link%int => int
       link%i = i
    end if
  end subroutine external_link_set

\eatline
\nwendcode{}\nwbegindocs{784}Reassign an external link to a new interaction (which should be an
image of the original target).
\nwenddocs{}\nwbegincode{785}\moddef{Interactions: procedures}\plusendmoddef
  subroutine external_link_reassign (link, int_src, int_target)
    type(external_link_t), intent(inout) :: link
    type(interaction_t), intent(in) :: int_src
    type(interaction_t), intent(in), target :: int_target
    if (associated (link%int)) then
       if (link%int%tag == int_src%tag)  link%int => int_target
    end if
  end subroutine external_link_reassign

\eatline
\nwendcode{}\nwbegindocs{786}Return true if the link is set
\nwenddocs{}\nwbegincode{787}\moddef{Interactions: procedures}\plusendmoddef
  function external_link_is_set (link) result (flag)
    logical :: flag
    type(external_link_t), intent(in) :: link
    flag = associated (link%int)
  end function external_link_is_set

\eatline
\nwendcode{}\nwbegindocs{788}Return the interaction pointer.
\nwenddocs{}\nwbegincode{789}\moddef{Interactions: public}\endmoddef
  public :: external_link_get_ptr
\nwendcode{}\nwbegincode{790}\moddef{Interactions: procedures}\plusendmoddef
  function external_link_get_ptr (link) result (int)
    type(interaction_t), pointer :: int
    type(external_link_t), intent(in) :: link
    int => link%int
  end function external_link_get_ptr

\eatline
\nwendcode{}\nwbegindocs{791}Return the index within that interaction
\nwenddocs{}\nwbegincode{792}\moddef{Interactions: public}\plusendmoddef
  public :: external_link_get_index
\nwendcode{}\nwbegincode{793}\moddef{Interactions: procedures}\plusendmoddef
  function external_link_get_index (link) result (i)
    integer :: i
    type(external_link_t), intent(in) :: link
    i = link%i
  end function external_link_get_index

\eatline
\nwendcode{}\nwbegindocs{794}Return a pointer to the momentum of the corresponding particle.  If
there is no association, return a null pointer.
\nwenddocs{}\nwbegincode{795}\moddef{Interactions: procedures}\plusendmoddef
  function external_link_get_momentum_ptr (link) result (p)
    type(vector4_t), pointer :: p
    type(external_link_t), intent(in) :: link
    if (associated (link%int)) then
       p => link%int%p(link%i)
    else
       p => null ()
    end if
  end function external_link_get_momentum_ptr

\eatline
\nwendcode{}\nwbegindocs{796}\nwdocspar
\subsection{Internal relations}
In addition to the external links, particles within the interaction
have parent-child relations.  Here, more than one link is possible,
and we set up an array.
\nwenddocs{}\nwbegincode{797}\moddef{Interactions: types}\plusendmoddef
  type :: internal_link_list_t
     private
     integer :: length = 0
     integer, dimension(:), allocatable :: link
   contains
   \LA{}Interactions: internal link list: TBP\RA{}
  end type internal_link_list_t

\eatline
\nwendcode{}\nwbegindocs{798}Output, non-advancing.
\nwenddocs{}\nwbegincode{799}\moddef{Interactions: internal link list: TBP}\endmoddef
  procedure :: write => internal_link_list_write
\nwendcode{}\nwbegincode{800}\moddef{Interactions: procedures}\plusendmoddef
  subroutine internal_link_list_write (object, unit)
    class(internal_link_list_t), intent(in) :: object
    integer, intent(in), optional :: unit
    integer :: u, i
    u = given_output_unit (unit)
    do i = 1, object%length
       write (u, "(1x,I0)", advance="no")  object%link(i)
    end do
  end subroutine internal_link_list_write
  
\eatline
\nwendcode{}\nwbegindocs{801}Append an item.  Start with an array size of 2 and double the size
if necessary.

Make sure that the indices are stored in ascending order.  To this
end, shift the existing entries right, starting from the end, as long
as they are larger than the new entry.
\nwenddocs{}\nwbegincode{802}\moddef{Interactions: internal link list: TBP}\plusendmoddef
  procedure :: append => internal_link_list_append
\nwendcode{}\nwbegincode{803}\moddef{Interactions: procedures}\plusendmoddef
  subroutine internal_link_list_append (link_list, link)
    class(internal_link_list_t), intent(inout) :: link_list
    integer, intent(in) :: link
    integer :: l, j
    integer, dimension(:), allocatable :: tmp
    l = link_list%length
    if (allocated (link_list%link)) then
       if (l == size (link_list%link)) then
          allocate (tmp (2 * l))
          tmp(:l) = link_list%link
          call move_alloc (from = tmp, to = link_list%link)
       end if
    else
       allocate (link_list%link (2))
    end if
    link_list%link(l+1) = link
    SHIFT_LINK_IN_PLACE: do j = l, 1, -1
       if (link >= link_list%link(j)) then
          exit SHIFT_LINK_IN_PLACE
       else
          link_list%link(j+1) = link_list%link(j)
          link_list%link(j) = link
       end if
    end do SHIFT_LINK_IN_PLACE
    link_list%length = l + 1
  end subroutine internal_link_list_append

\eatline
\nwendcode{}\nwbegindocs{804}Return true if the link list is nonempty:
\nwenddocs{}\nwbegincode{805}\moddef{Interactions: internal link list: TBP}\plusendmoddef
  procedure :: has_entries => internal_link_list_has_entries
\nwendcode{}\nwbegincode{806}\moddef{Interactions: procedures}\plusendmoddef
  function internal_link_list_has_entries (link_list) result (flag)
    class(internal_link_list_t), intent(in) :: link_list
    logical :: flag
    flag = link_list%length > 0
  end function internal_link_list_has_entries

\eatline
\nwendcode{}\nwbegindocs{807}Return the list length
\nwenddocs{}\nwbegincode{808}\moddef{Interactions: internal link list: TBP}\plusendmoddef
  procedure :: get_length => internal_link_list_get_length
\nwendcode{}\nwbegincode{809}\moddef{Interactions: procedures}\plusendmoddef
  function internal_link_list_get_length (link_list) result (length)
    class(internal_link_list_t), intent(in) :: link_list
    integer :: length
    length = link_list%length
  end function internal_link_list_get_length

\eatline
\nwendcode{}\nwbegindocs{810}Return an entry.
\nwenddocs{}\nwbegincode{811}\moddef{Interactions: internal link list: TBP}\plusendmoddef
  procedure :: get_link => internal_link_list_get_link
\nwendcode{}\nwbegincode{812}\moddef{Interactions: procedures}\plusendmoddef
  function internal_link_list_get_link (link_list, i) result (link)
    class(internal_link_list_t), intent(in) :: link_list
    integer, intent(in) :: i
    integer :: link
    if (i <= link_list%length) then
       link = link_list%link(i)
    else
       call msg_bug ("Internal link list: out of bounds")
    end if
  end function internal_link_list_get_link

\eatline
\nwendcode{}\nwbegindocs{813}\nwdocspar
\subsection{The interaction type}
An interaction is an entangled system of particles.  Thus, the
interaction object consists of two parts: the subevent, and the
quantum state which technically is a trie.  The subnode levels beyond
the trie root node are in correspondence to the subevent, so
both should be traversed in parallel.

The subevent is implemented as an allocatable array of
four-momenta.  The first {\tt{}n{\char95}in} particles are incoming, {\tt{}n{\char95}vir}
particles in-between can be kept for bookkeeping, and the last
{\tt{}n{\char95}out} particles are outgoing.

Distinct interactions are linked by their particles: for each
particle, we have the possibility of links to corresponding particles
in other interactions.  Furthermore, for bookkeeping purposes we have
a self-link array {\tt{}relations} where the parent-child relations are
kept, and a flag array {\tt{}resonant} which is set for an intermediate
resonance.

Each momentum is associated with masks for flavor, color, and
helicity.  If a mask entry is set, the associated quantum number is to
be ignored for that particle.  If any mask has changed, the flag
{\tt{}update} is set.

We can have particle pairs locked together.  If this is the case, the
corresponding mask entries are bound to be equal.  This is useful for
particles that go through the interaction.

The interaction tag serves bookkeeping purposes.  In particular, it
identifies links in printout.
\nwenddocs{}\nwbegincode{814}\moddef{Interactions: public}\plusendmoddef
  public :: interaction_t
\nwendcode{}\nwbegincode{815}\moddef{Interactions: types}\plusendmoddef
  type :: interaction_t
     private
     integer :: tag = 0
     type(state_matrix_t) :: state_matrix
     integer :: n_in = 0
     integer :: n_vir = 0
     integer :: n_out = 0
     integer :: n_tot = 0
     logical, dimension(:), allocatable :: p_is_known
     type(vector4_t), dimension(:), allocatable :: p
     type(external_link_t), dimension(:), allocatable :: source
     type(internal_link_list_t), dimension(:), allocatable :: parents
     type(internal_link_list_t), dimension(:), allocatable :: children
     logical, dimension(:), allocatable :: resonant
     type(quantum_numbers_mask_t), dimension(:), allocatable :: mask
     integer, dimension(:), allocatable :: hel_lock
     logical :: update_state_matrix = .false.
     logical :: update_values = .false.
   contains
   \LA{}Interactions: interaction: TBP\RA{}
  end type interaction_t

\eatline
\nwendcode{}\nwbegindocs{816}Initialize the particle array with a fixed size.  The first {\tt{}n{\char95}in}
particles are incoming, the rest outgoing.  Masks are optional.  There
is also an optional tag.  The interaction still needs fixing the
values, but that is to be done after all branches have been added.

Interaction tags are assigned consecutively, using a {\tt{}save}d
variable local to this procedure.  If desired, we can provide a seed
for the interaction tags.  Such a seed should be positive.  The
default seed is one. {\tt{}tag=0} indicates an empty interaction.

If {\tt{}set{\char95}relations} is set and true, we establish parent-child
relations for all incoming and outgoing particles.  Virtual particles
are skipped; this option is normally used only for interations without
virtual particles.
\nwenddocs{}\nwbegincode{817}\moddef{Interactions: public}\plusendmoddef
  public :: interaction_init
\nwendcode{}\nwbegincode{818}\moddef{Interactions: procedures}\plusendmoddef
  subroutine interaction_init &
       (int, n_in, n_vir, n_out, &
        tag, resonant, mask, hel_lock, set_relations, store_values)
    type(interaction_t), intent(out) :: int
    integer, intent(in) :: n_in, n_vir, n_out
    integer, intent(in), optional :: tag
    logical, dimension(:), intent(in), optional :: resonant
    type(quantum_numbers_mask_t), dimension(:), intent(in), optional :: mask
    integer, dimension(:), intent(in), optional :: hel_lock
    logical, intent(in), optional :: set_relations, store_values
    logical :: set_rel
    integer :: i, j
    set_rel = .false.;  if (present (set_relations))  set_rel = set_relations
    call interaction_set_tag (int, tag)
    call int%state_matrix%init (store_values)
    int%n_in = n_in
    int%n_vir = n_vir
    int%n_out = n_out
    int%n_tot = n_in + n_vir + n_out
    allocate (int%p_is_known (int%n_tot))
    int%p_is_known = .false.
    allocate (int%p (int%n_tot))
    allocate (int%source (int%n_tot))
    allocate (int%parents (int%n_tot))
    allocate (int%children (int%n_tot))
    allocate (int%resonant (int%n_tot))
    if (present (resonant)) then
       int%resonant = resonant
    else
       int%resonant = .false.
    end if
    allocate (int%mask (int%n_tot))
    allocate (int%hel_lock (int%n_tot))
    if (present (mask)) then
       int%mask = mask
    end if
    if (present (hel_lock)) then
       int%hel_lock = hel_lock
    else
       int%hel_lock = 0
    end if
    int%update_state_matrix = .false.
    int%update_values = .true.
    if (set_rel) then
       do i = 1, n_in
          do j = 1, n_out
             call interaction_relate (int, i, n_in + j)
          end do
       end do
    end if
  end subroutine interaction_init

\eatline
\nwendcode{}\nwbegindocs{819}Set or create a unique tag for the interaction.  Without
interaction, reset the tag counter.
\nwenddocs{}\nwbegincode{820}\moddef{Interactions: procedures}\plusendmoddef
  subroutine interaction_set_tag (int, tag)
    type(interaction_t), intent(inout), optional :: int
    integer, intent(in), optional :: tag
    integer, save :: stored_tag = 1
    if (present (int)) then
       if (present (tag)) then
          int%tag = tag
       else
          int%tag = stored_tag
          stored_tag = stored_tag + 1
       end if
    else if (present (tag)) then
       stored_tag = tag
    else
       stored_tag = 1
    end if
  end subroutine interaction_set_tag

\eatline
\nwendcode{}\nwbegindocs{821}The public interface for the previous procedure only covers the
reset functionality.
\nwenddocs{}\nwbegincode{822}\moddef{Interactions: public}\plusendmoddef
  public :: reset_interaction_counter
\nwendcode{}\nwbegincode{823}\moddef{Interactions: procedures}\plusendmoddef
  subroutine reset_interaction_counter (tag)
    integer, intent(in), optional :: tag
    call interaction_set_tag (tag=tag)
  end subroutine reset_interaction_counter

\eatline
\nwendcode{}\nwbegindocs{824}Finalizer: The state-matrix object contains pointers.
\nwenddocs{}\nwbegincode{825}\moddef{Interactions: public}\plusendmoddef
  public :: interaction_final
\nwendcode{}\nwbegincode{826}\moddef{Interactions: procedures}\plusendmoddef
  subroutine interaction_final (int)
    type(interaction_t), intent(inout) :: int
    call int%state_matrix%final ()
  end subroutine interaction_final

\eatline
\nwendcode{}\nwbegindocs{827}Output.  The {\tt{}verbose} option refers to the state matrix output.
\nwenddocs{}\nwbegincode{828}\moddef{Interactions: public}\plusendmoddef
  public :: interaction_write
\nwendcode{}\nwbegincode{829}\moddef{Interactions: procedures}\plusendmoddef
  subroutine interaction_write &
       (int, unit, verbose, show_momentum_sum, show_mass, show_state, &
       col_verbose, testflag)
    type(interaction_t), intent(in) :: int
    integer, intent(in), optional :: unit
    logical, intent(in), optional :: verbose, show_momentum_sum, show_mass
    logical, intent(in), optional :: show_state, col_verbose, testflag
    integer :: u
    integer :: i, index_link
    type(interaction_t), pointer :: int_link
    logical :: show_st
    u = given_output_unit (unit);  if (u < 0)  return
    show_st = .true.;  if (present (show_state))  show_st = show_state
    if (int%tag /= 0) then
       write (u, "(1x,A,I0)")  "Interaction: ", int%tag
       do i = 1, int%n_tot
          if (i == 1 .and. int%n_in > 0) then
             write (u, "(1x,A)") "Incoming:"
          else if (i == int%n_in + 1 .and. int%n_vir > 0) then
             write (u, "(1x,A)") "Virtual:"
          else if (i == int%n_in + int%n_vir + 1 .and. int%n_out > 0) then
             write (u, "(1x,A)") "Outgoing:"
          end if
          write (u, "(1x,A,1x,I0)", advance="no") "Particle", i
          if (allocated (int%resonant)) then
             if (int%resonant(i)) then
                write (u, "(A)") "[r]"
             else
                write (u, *)
             end if
          else
             write (u, *)
          end if
          if (allocated (int%p)) then
             if (int%p_is_known(i)) then
                call vector4_write (int%p(i), u, show_mass, testflag)
             else
                write (u, "(A)")  "  [momentum undefined]"
             end if
          else
             write (u, "(A)") "  [momentum not allocated]"
          end if
          if (allocated (int%mask)) then
             write (u, "(1x,A)", advance="no")  "mask [fch] = "
             call int%mask(i)%write (u)
             write (u, *)
          end if
          if (int%parents(i)%has_entries () &
               .or. int%children(i)%has_entries ()) then
             write (u, "(1x,A)", advance="no") "internal links:"
             call int%parents(i)%write (u)
             if (int%parents(i)%has_entries ()) &
                  write (u, "(1x,A)", advance="no") "=>"
             write (u, "(1x,A)", advance="no") "X"
             if (int%children(i)%has_entries ()) &
                  write (u, "(1x,A)", advance="no") "=>" 
             call int%children(i)%write (u)
             write (u, *)
          end if
          if (allocated (int%hel_lock)) then
             if (int%hel_lock(i) /= 0) then
                write (u, "(1x,A,1x,I0)")  "helicity lock:", int%hel_lock(i)
             end if
          end if
          if (external_link_is_set (int%source(i))) then
             write (u, "(1x,A)", advance="no") "source:"
             int_link => external_link_get_ptr (int%source(i))
             index_link = external_link_get_index (int%source(i))
             write (u, "(1x,'(',I0,')',I0)", advance="no") &
                  int_link%tag, index_link
             write (u, *)
          end if
       end do
       if (present (show_momentum_sum)) then
          if (allocated (int%p) .and. show_momentum_sum) then
             write (u, "(1x,A)") "Incoming particles (sum):"
             call vector4_write &
                  (sum (int%p(1:int%n_in)), u, show_mass = show_mass)
             write (u, "(1x,A)") "Outgoing particles (sum):"
             call vector4_write &
                  (sum (int%p(int%n_in+int%n_vir+1:)), u, show_mass = show_mass)
             write (u, *)
          end if
       end if
       if (show_st) then
          call int%state_matrix%write (write_value_list=verbose, &
               verbose=verbose, unit=unit, col_verbose=col_verbose, &
               testflag = testflag)
       end if
    else
       write (u, "(1x,A)") "Interaction: [empty]"
    end if
  end subroutine interaction_write

\eatline
\nwendcode{}\nwbegindocs{830}Assignment: We implement this as a deep copy.  This applies, in
particular, to the state-matrix and internal-link components.
Furthermore, the new interaction acquires a new tag.
\nwenddocs{}\nwbegincode{831}\moddef{Interactions: public}\plusendmoddef
  public :: assignment(=)
\nwendcode{}\nwbegincode{832}\moddef{Interactions: interfaces}\endmoddef
  interface assignment(=)
     module procedure interaction_assign
  end interface

\nwendcode{}\nwbegincode{833}\moddef{Interactions: procedures}\plusendmoddef
  subroutine interaction_assign (int_out, int_in)
    type(interaction_t), intent(out) :: int_out
    type(interaction_t), intent(in), target :: int_in
    call interaction_set_tag (int_out)
    int_out%state_matrix = int_in%state_matrix
    int_out%n_in  = int_in%n_in
    int_out%n_out = int_in%n_out
    int_out%n_vir = int_in%n_vir
    int_out%n_tot = int_in%n_tot
    if (allocated (int_in%p_is_known)) then
       allocate (int_out%p_is_known (size (int_in%p_is_known)))
       int_out%p_is_known = int_in%p_is_known
    end if
    if (allocated (int_in%p)) then
       allocate (int_out%p (size (int_in%p)))
       int_out%p = int_in%p
    end if
    if (allocated (int_in%source)) then
       allocate (int_out%source (size (int_in%source)))
       int_out%source = int_in%source
    end if
    if (allocated (int_in%parents)) then
       allocate (int_out%parents (size (int_in%parents)))
       int_out%parents = int_in%parents
    end if
    if (allocated (int_in%children)) then
       allocate (int_out%children (size (int_in%children)))
       int_out%children = int_in%children
    end if
    if (allocated (int_in%resonant)) then
       allocate (int_out%resonant (size (int_in%resonant)))
       int_out%resonant = int_in%resonant
    end if
    if (allocated (int_in%mask)) then
       allocate (int_out%mask (size (int_in%mask)))
       int_out%mask = int_in%mask
    end if
    if (allocated (int_in%hel_lock)) then
       allocate (int_out%hel_lock (size (int_in%hel_lock)))
       int_out%hel_lock = int_in%hel_lock
    end if
    int_out%update_state_matrix = int_in%update_state_matrix
    int_out%update_values = int_in%update_values
  end subroutine interaction_assign

\eatline
\nwendcode{}\nwbegindocs{834}\nwdocspar
\subsection{Methods inherited from the state matrix member}
Until F2003 is standard, we cannot implement inheritance directly.
Therefore, we need wrappers for ``inherited'' methods.

Make a new branch in the state matrix if it does not yet exist.  This
is not just a wrapper but it introduces the interaction mask: where a
quantum number is masked, it is not transferred but set undefined.
After this, the value array has to be updated.
\nwenddocs{}\nwbegincode{835}\moddef{Interactions: public}\plusendmoddef
  public :: interaction_add_state
\nwendcode{}\nwbegincode{836}\moddef{Interactions: procedures}\plusendmoddef
  subroutine interaction_add_state &
       (int, qn, index, value, sum_values, counter_index, me_index)
    type(interaction_t), intent(inout) :: int
    type(quantum_numbers_t), dimension(:), intent(in) :: qn
    integer, intent(in), optional :: index
    complex(default), intent(in), optional :: value
    logical, intent(in), optional :: sum_values
    integer, intent(in), optional :: counter_index
    integer, intent(out), optional :: me_index
    type(quantum_numbers_t), dimension(size(qn)) :: qn_tmp
    qn_tmp = qn
    call qn_tmp%undefine (int%mask)
    call int%state_matrix%add_state (qn_tmp, index, value, sum_values, &
         counter_index, me_index)
    int%update_values = .true.
  end subroutine interaction_add_state

\eatline
\nwendcode{}\nwbegindocs{837}Freeze the quantum state: First collapse the quantum state, i.e.,
remove quantum numbers if any mask has changed, then fix the array of
value pointers.
\nwenddocs{}\nwbegincode{838}\moddef{Interactions: public}\plusendmoddef
  public :: interaction_freeze
\nwendcode{}\nwbegincode{839}\moddef{Interactions: procedures}\plusendmoddef
  subroutine interaction_freeze (int)
    type(interaction_t), intent(inout) :: int
    if (int%update_state_matrix) then
       call int%state_matrix%collapse (int%mask)
       int%update_state_matrix = .false.
       int%update_values = .true.
    end if
    if (int%update_values) then
       call int%state_matrix%freeze ()
       int%update_values = .false.
    end if
  end subroutine interaction_freeze

\eatline
\nwendcode{}\nwbegindocs{840}Return true if the state matrix is empty.
\nwenddocs{}\nwbegincode{841}\moddef{Interactions: public}\plusendmoddef
  public :: interaction_is_empty
\nwendcode{}\nwbegincode{842}\moddef{Interactions: procedures}\plusendmoddef
  function interaction_is_empty (int) result (flag)
    logical :: flag
    type(interaction_t), intent(in) :: int
    flag = int%state_matrix%is_empty ()
  end function interaction_is_empty

\eatline
\nwendcode{}\nwbegindocs{843}Get the number of values stored in the state matrix:
\nwenddocs{}\nwbegincode{844}\moddef{Interactions: public}\plusendmoddef
  public :: interaction_get_n_matrix_elements
\nwendcode{}\nwbegincode{845}\moddef{Interactions: procedures}\plusendmoddef
  function interaction_get_n_matrix_elements (int) result (n)
    integer :: n
    type(interaction_t), intent(in) :: int
    n = int%state_matrix%get_n_matrix_elements ()
  end function interaction_get_n_matrix_elements

\eatline
\nwendcode{}\nwbegindocs{846}Get the norm of the state matrix (if the norm has been taken out, otherwise
this would be unity).
\nwenddocs{}\nwbegincode{847}\moddef{Interactions: public}\plusendmoddef
  public :: interaction_get_norm
\nwendcode{}\nwbegincode{848}\moddef{Interactions: procedures}\plusendmoddef
  function interaction_get_norm (int) result (norm)
    real(default) :: norm
    type(interaction_t), intent(in) :: int
    norm = int%state_matrix%get_norm ()
  end function interaction_get_norm

\eatline
\nwendcode{}\nwbegindocs{849}Get the quantum number array that corresponds to a given index.
\nwenddocs{}\nwbegincode{850}\moddef{Interactions: public}\plusendmoddef
  public :: interaction_get_quantum_numbers
\nwendcode{}\nwbegincode{851}\moddef{Interactions: procedures}\plusendmoddef
  function interaction_get_quantum_numbers (int, i) result (qn)
    type(quantum_numbers_t), dimension(:), allocatable :: qn
    type(interaction_t), intent(in), target :: int
    integer, intent(in) :: i
    allocate (qn (int%state_matrix%get_depth ()))
    qn = int%state_matrix%get_quantum_numbers (i)
  end function interaction_get_quantum_numbers

\eatline
\nwendcode{}\nwbegindocs{852}Get the matrix element that corresponds to a set of quantum
numbers, a given index, or return the whole array.
\nwenddocs{}\nwbegincode{853}\moddef{Interactions: public}\plusendmoddef
  public :: interaction_get_matrix_element
\nwendcode{}\nwbegincode{854}\moddef{Interactions: procedures}\plusendmoddef
  function interaction_get_matrix_element (int, i) result (me)
    complex(default) :: me
    type(interaction_t), intent(in) :: int
    integer, intent(in) :: i
    me = int%state_matrix%get_matrix_element (i)
  end function interaction_get_matrix_element

\eatline
\nwendcode{}\nwbegindocs{855}Set the complex value(s) stored in the quantum state.
\nwenddocs{}\nwbegincode{856}\moddef{Interactions: public}\plusendmoddef
  public :: interaction_set_matrix_element
\nwendcode{}\nwbegincode{857}\moddef{Interactions: interfaces}\plusendmoddef
  interface interaction_set_matrix_element
     module procedure interaction_set_matrix_element_qn
     module procedure interaction_set_matrix_element_all
     module procedure interaction_set_matrix_element_array 
     module procedure interaction_set_matrix_element_single
     module procedure interaction_set_matrix_element_clone
  end interface
\eatline
\nwendcode{}\nwbegindocs{858}Indirect access via the quantum number array:
\nwenddocs{}\nwbegincode{859}\moddef{Interactions: procedures}\plusendmoddef
  subroutine interaction_set_matrix_element_qn (int, qn, val)
    type(interaction_t), intent(inout) :: int
    type(quantum_numbers_t), dimension(:), intent(in) :: qn
    complex(default), intent(in) :: val
    call int%state_matrix%set_matrix_element (qn, val)
  end subroutine interaction_set_matrix_element_qn

\eatline
\nwendcode{}\nwbegindocs{860}Set all entries of the matrix-element array to a given value.
\nwenddocs{}\nwbegincode{861}\moddef{Interactions: procedures}\plusendmoddef
  subroutine interaction_set_matrix_element_all (int, value)
    type(interaction_t), intent(inout) :: int
    complex(default), intent(in) :: value
    call int%state_matrix%set_matrix_element (value)
  end subroutine interaction_set_matrix_element_all

\eatline
\nwendcode{}\nwbegindocs{862}Set the matrix-element array directly.
\nwenddocs{}\nwbegincode{863}\moddef{Interactions: procedures}\plusendmoddef
  subroutine interaction_set_matrix_element_array (int, value)
    type(interaction_t), intent(inout) :: int
    complex(default), dimension(:), intent(in) :: value
    call int%state_matrix%set_matrix_element (value)
  end subroutine interaction_set_matrix_element_array

  pure subroutine interaction_set_matrix_element_single (int, i, value)
    type(interaction_t), intent(inout) :: int
    integer, intent(in) :: i
    complex(default), intent(in) :: value
    call int%state_matrix%set_matrix_element (i, value)
  end subroutine interaction_set_matrix_element_single

\eatline
\eatline
\nwendcode{}\nwbegindocs{864}Clone from another (matching) interaction.
\nwenddocs{}\nwbegincode{865}\moddef{Interactions: procedures}\plusendmoddef
  subroutine interaction_set_matrix_element_clone (int, int1)
    type(interaction_t), intent(inout) :: int
    type(interaction_t), intent(in) :: int1
    call int%state_matrix%set_matrix_element (int1%state_matrix)
  end subroutine interaction_set_matrix_element_clone

\eatline
\nwendcode{}\nwbegindocs{866}Get the indices of any diagonal matrix elements.
\nwenddocs{}\nwbegincode{867}\moddef{Interactions: public}\plusendmoddef
  public :: interaction_get_diagonal_entries
\nwendcode{}\nwbegincode{868}\moddef{Interactions: procedures}\plusendmoddef
  subroutine interaction_get_diagonal_entries (int, i)
    type(interaction_t), intent(in) :: int
    integer, dimension(:), allocatable, intent(out) :: i
    call int%state_matrix%get_diagonal_entries (i)
  end subroutine interaction_get_diagonal_entries

\eatline
\nwendcode{}\nwbegindocs{869}Renormalize the state matrix by its trace, if nonzero.  The renormalization
is reflected in the state-matrix norm.
\nwenddocs{}\nwbegincode{870}\moddef{Interactions: public}\plusendmoddef
  public :: interaction_normalize_by_trace
\nwendcode{}\nwbegincode{871}\moddef{Interactions: procedures}\plusendmoddef
  subroutine interaction_normalize_by_trace (int)
    type(interaction_t), intent(inout) :: int
    call int%state_matrix%normalize_by_trace ()
  end subroutine interaction_normalize_by_trace

\eatline
\nwendcode{}\nwbegindocs{872}Analogous, but renormalize by maximal (absolute) value.
\nwenddocs{}\nwbegincode{873}\moddef{Interactions: public}\plusendmoddef
  public :: interaction_normalize_by_max
\nwendcode{}\nwbegincode{874}\moddef{Interactions: procedures}\plusendmoddef
  subroutine interaction_normalize_by_max (int)
    type(interaction_t), intent(inout) :: int
    call int%state_matrix%normalize_by_max ()
  end subroutine interaction_normalize_by_max

\eatline
\nwendcode{}\nwbegindocs{875}Explicitly set the norm value (of the state matrix).
\nwenddocs{}\nwbegincode{876}\moddef{Interactions: public}\plusendmoddef
  public :: interaction_set_norm
\nwendcode{}\nwbegincode{877}\moddef{Interactions: procedures}\plusendmoddef
  subroutine interaction_set_norm (int, norm)
    type(interaction_t), intent(inout) :: int
    real(default), intent(in) :: norm
    call int%state_matrix%set_norm (norm)
  end subroutine interaction_set_norm

\eatline
\nwendcode{}\nwbegindocs{878}Return the maximum absolute value of color indices.
\nwenddocs{}\nwbegincode{879}\moddef{Interactions: public}\plusendmoddef
  public :: interaction_get_max_color_value
\nwendcode{}\nwbegincode{880}\moddef{Interactions: procedures}\plusendmoddef
  function interaction_get_max_color_value (int) result (cmax)
    integer :: cmax
    type(interaction_t), intent(in) :: int
    cmax = int%state_matrix%get_max_color_value ()
  end function interaction_get_max_color_value

\eatline
\nwendcode{}\nwbegindocs{881}Factorize the state matrix into single-particle state matrices, the
branch selection depending on a (random) value between 0 and 1;
optionally also return a correlated state matrix.
\nwenddocs{}\nwbegincode{882}\moddef{Interactions: public}\plusendmoddef
  public :: interaction_factorize
\nwendcode{}\nwbegincode{883}\moddef{Interactions: procedures}\plusendmoddef
  subroutine interaction_factorize &
       (int, mode, x, ok, single_state, correlated_state, qn_in)
    type(interaction_t), intent(in), target :: int
    integer, intent(in) :: mode
    real(default), intent(in) :: x
    logical, intent(out) :: ok
    type(state_matrix_t), &
         dimension(:), allocatable, intent(out) :: single_state
    type(state_matrix_t), intent(out), optional :: correlated_state
    type(quantum_numbers_t), dimension(:), intent(in), optional :: qn_in
    call int%state_matrix%factorize &
         (mode, x, ok, single_state, correlated_state, qn_in)
  end subroutine interaction_factorize

\eatline
\nwendcode{}\nwbegindocs{884}Sum all matrix element values
\nwenddocs{}\nwbegincode{885}\moddef{Interactions: public}\plusendmoddef
  public :: interaction_sum
\nwendcode{}\nwbegincode{886}\moddef{Interactions: procedures}\plusendmoddef
  function interaction_sum (int) result (value)
    complex(default) :: value
    type(interaction_t), intent(in) :: int
    value = int%state_matrix%sum ()
  end function interaction_sum

\eatline
\nwendcode{}\nwbegindocs{887}Append new states which are color-contracted versions of the
existing states.  The matrix element index of each color contraction
coincides with the index of its origin, so no new matrix elements are
generated.  After this operation, no {\tt{}freeze} must be performed
anymore.
\nwenddocs{}\nwbegincode{888}\moddef{Interactions: public}\plusendmoddef
  public :: interaction_add_color_contractions
\nwendcode{}\nwbegincode{889}\moddef{Interactions: procedures}\plusendmoddef
  subroutine interaction_add_color_contractions (int)
    type(interaction_t), intent(inout) :: int
    call int%state_matrix%add_color_contractions ()
  end subroutine interaction_add_color_contractions

\eatline
\nwendcode{}\nwbegindocs{890}Multiply matrix elements from two interactions.  Choose the elements
as given by the integer index arrays, multiply them and store the sum
of products in the indicated matrix element.  The suffixes mean:
c=conjugate first factor; f=include weighting factor.
\nwenddocs{}\nwbegincode{891}\moddef{Interactions: public}\plusendmoddef
  public :: interaction_evaluate_product
  public :: interaction_evaluate_product_cf
  public :: interaction_evaluate_square_c
  public :: interaction_evaluate_sum
  public :: interaction_evaluate_me_sum
\nwendcode{}\nwbegincode{892}\moddef{Interactions: procedures}\plusendmoddef
  pure subroutine interaction_evaluate_product &
       (int, i, int1, int2, index1, index2)
    type(interaction_t), intent(inout) :: int
    integer, intent(in) :: i
    type(interaction_t), intent(in) :: int1, int2
    integer, dimension(:), intent(in) :: index1, index2
    call int%state_matrix%evaluate_product &
         (i, int1%state_matrix, int2%state_matrix, &
          index1, index2)
  end subroutine interaction_evaluate_product

  pure subroutine interaction_evaluate_product_cf &
       (int, i, int1, int2, index1, index2, factor)
    type(interaction_t), intent(inout) :: int
    integer, intent(in) :: i
    type(interaction_t), intent(in) :: int1, int2
    integer, dimension(:), intent(in) :: index1, index2
    complex(default), dimension(:), intent(in) :: factor
    call int%state_matrix%evaluate_product_cf &
         (i, int1%state_matrix, int2%state_matrix, &
          index1, index2, factor)
  end subroutine interaction_evaluate_product_cf

  pure subroutine interaction_evaluate_square_c (int, i, int1, index1)
    type(interaction_t), intent(inout) :: int
    integer, intent(in) :: i
    type(interaction_t), intent(in) :: int1
    integer, dimension(:), intent(in) :: index1
    call int%state_matrix%evaluate_square_c (i, int1%state_matrix, index1)
  end subroutine interaction_evaluate_square_c

  pure subroutine interaction_evaluate_sum (int, i, int1, index1)
    type(interaction_t), intent(inout) :: int
    integer, intent(in) :: i
    type(interaction_t), intent(in) :: int1
    integer, dimension(:), intent(in) :: index1
    call int%state_matrix%evaluate_sum (i, int1%state_matrix, index1)
  end subroutine interaction_evaluate_sum

  pure subroutine interaction_evaluate_me_sum (int, i, int1, index1)
    type(interaction_t), intent(inout) :: int
    integer, intent(in) :: i
    type(interaction_t), intent(in) :: int1
    integer, dimension(:), intent(in) :: index1
    call int%state_matrix%evaluate_me_sum (i, int1%state_matrix, index1)
  end subroutine interaction_evaluate_me_sum

\eatline
\eatline
\eatline
\eatline
\eatline
\nwendcode{}\nwbegindocs{893}\nwdocspar
\subsection{Accessing contents}
Return the integer tag.
\nwenddocs{}\nwbegincode{894}\moddef{Interactions: public}\plusendmoddef
  public :: interaction_get_tag
\nwendcode{}\nwbegincode{895}\moddef{Interactions: procedures}\plusendmoddef
  function interaction_get_tag (int) result (tag)
    integer :: tag
    type(interaction_t), intent(in) :: int
    tag = int%tag
  end function interaction_get_tag

\eatline
\nwendcode{}\nwbegindocs{896}Return the number of particles.
\nwenddocs{}\nwbegincode{897}\moddef{Interactions: public}\plusendmoddef
  public :: interaction_get_n_tot
  public :: interaction_get_n_in
  public :: interaction_get_n_vir
  public :: interaction_get_n_out
\nwendcode{}\nwbegincode{898}\moddef{Interactions: procedures}\plusendmoddef
  function interaction_get_n_tot (int) result (n_tot)
    integer :: n_tot
    type(interaction_t), intent(in) :: int
    n_tot = int%n_tot
  end function interaction_get_n_tot

  function interaction_get_n_in (int) result (n_in)
    integer :: n_in
    type(interaction_t), intent(in) :: int
    n_in = int%n_in
  end function interaction_get_n_in

  function interaction_get_n_vir (int) result (n_vir)
    integer :: n_vir
    type(interaction_t), intent(in) :: int
    n_vir = int%n_vir
  end function interaction_get_n_vir

  function interaction_get_n_out (int) result (n_out)
    integer :: n_out
    type(interaction_t), intent(in) :: int
    n_out = int%n_out
  end function interaction_get_n_out

\eatline
\eatline
\nwendcode{}\nwbegindocs{899}Return a momentum index.  The flags specify whether to keep/drop
incoming, virtual, or outgoing momenta.  Check for illegal values.
\nwenddocs{}\nwbegincode{900}\moddef{Interactions: procedures}\plusendmoddef
  function idx (int, i, outgoing)
    integer :: idx
    type(interaction_t), intent(in) :: int
    integer, intent(in) :: i
    logical, intent(in), optional :: outgoing
    logical :: in, vir, out
    if (present (outgoing)) then
       in  = .not. outgoing
       vir = .false.
       out = outgoing
    else
       in = .true.
       vir = .true.
       out = .true.
    end if
    idx = 0
    if (in) then
       if (vir) then
          if (out) then
             if (i <= int%n_tot)  idx = i
          else
             if (i <= int%n_in + int%n_vir)  idx = i
          end if
       else if (out) then
          if (i <= int%n_in) then
             idx = i
          else if (i <= int%n_in + int%n_out) then
             idx = int%n_vir + i
          end if
       else
          if (i <= int%n_in)  idx = i
       end if
    else if (vir) then
       if (out) then
          if (i <= int%n_vir + int%n_out)  idx = int%n_in + i
       else
          if (i <= int%n_vir)  idx = int%n_in + i
       end if
    else if (out) then
       if (i <= int%n_out)  idx = int%n_in + int%n_vir + i
    end if
    if (idx == 0) then
       call interaction_write (int)
       print *, i, in, vir, out
       call msg_bug (" Momentum index is out of range for this interaction")
    end if
  end function idx

\eatline
\nwendcode{}\nwbegindocs{901}Return all or just a specific four-momentum.
\nwenddocs{}\nwbegincode{902}\moddef{Interactions: public}\plusendmoddef
  public :: interaction_get_momenta
  public :: interaction_get_momentum
\nwendcode{}\nwbegincode{903}\moddef{Interactions: interfaces}\plusendmoddef
  interface interaction_get_momenta
     module procedure interaction_get_momenta_all
     module procedure interaction_get_momenta_idx
  end interface
\nwendcode{}\nwbegincode{904}\moddef{Interactions: procedures}\plusendmoddef
  function interaction_get_momenta_all (int, outgoing) result (p)
    type(vector4_t), dimension(:), allocatable :: p
    type(interaction_t), intent(in) :: int
    logical, intent(in), optional :: outgoing
    integer :: i
    if (present (outgoing)) then
       if (outgoing) then
          allocate (p (int%n_out))
       else
          allocate (p (int%n_in))
       end if
    else
       allocate (p (int%n_tot))
    end if
    do i = 1, size (p)
       p(i) = int%p(idx (int, i, outgoing))
    end do
  end function interaction_get_momenta_all

  function interaction_get_momenta_idx (int, jj) result (p)
    type(vector4_t), dimension(:), allocatable :: p
    type(interaction_t), intent(in) :: int
    integer, dimension(:), intent(in) :: jj
    allocate (p (size (jj)))
    p = int%p(jj)
  end function interaction_get_momenta_idx

  function interaction_get_momentum (int, i, outgoing) result (p)
    type(vector4_t) :: p
    type(interaction_t), intent(in) :: int
    integer, intent(in) :: i
    logical, intent(in), optional :: outgoing
    p = int%p(idx (int, i, outgoing))
  end function interaction_get_momentum

\eatline
\nwendcode{}\nwbegindocs{905}This is a variant as a subroutine.  Redundant, but the function
above fails at times for gfortran 4.5.0 (double allocation, compiler
bug). 
\nwenddocs{}\nwbegincode{906}\moddef{Interactions: public}\plusendmoddef
  public :: interaction_get_momenta_sub
\nwendcode{}\nwbegincode{907}\moddef{Interactions: procedures}\plusendmoddef
  subroutine interaction_get_momenta_sub (int, p, outgoing)
    type(vector4_t), dimension(:), intent(out) :: p
    type(interaction_t), intent(in) :: int
    logical, intent(in), optional :: outgoing
    integer :: i
    do i = 1, size (p)
       p(i) = int%p(idx (int, i, outgoing))
    end do
  end subroutine interaction_get_momenta_sub

\eatline
\nwendcode{}\nwbegindocs{908}Return a shallow copy of the state matrix:
\nwenddocs{}\nwbegincode{909}\moddef{Interactions: public}\plusendmoddef
  public :: interaction_get_state_matrix_ptr
\nwendcode{}\nwbegincode{910}\moddef{Interactions: procedures}\plusendmoddef
  function interaction_get_state_matrix_ptr (int) result (state)
    type(state_matrix_t), pointer :: state
    type(interaction_t), intent(in), target :: int
    state => int%state_matrix
  end function interaction_get_state_matrix_ptr

\eatline
\nwendcode{}\nwbegindocs{911}Return the array of resonance flags
\nwenddocs{}\nwbegincode{912}\moddef{Interactions: public}\plusendmoddef
  public :: interaction_get_resonance_flags
\nwendcode{}\nwbegincode{913}\moddef{Interactions: procedures}\plusendmoddef
  function interaction_get_resonance_flags (int) result (resonant)
    type(interaction_t), intent(in) :: int
    logical, dimension(size(int%resonant)) :: resonant
    resonant = int%resonant
  end function interaction_get_resonance_flags

\eatline
\nwendcode{}\nwbegindocs{914}Return the quantum-numbers mask (or part of it)
\nwenddocs{}\nwbegincode{915}\moddef{Interactions: public}\plusendmoddef
  public :: interaction_get_mask
\nwendcode{}\nwbegincode{916}\moddef{Interactions: interfaces}\plusendmoddef
  interface interaction_get_mask
     module procedure interaction_get_mask_all
     module procedure interaction_get_mask_slice
  end interface

\nwendcode{}\nwbegincode{917}\moddef{Interactions: procedures}\plusendmoddef
  function interaction_get_mask_all (int) result (mask)
    type(interaction_t), intent(in) :: int
    type(quantum_numbers_mask_t), dimension(size(int%mask)) :: mask
    mask = int%mask
  end function interaction_get_mask_all

  function interaction_get_mask_slice (int, index) result (mask)
    type(interaction_t), intent(in) :: int
    integer, dimension(:), intent(in) :: index
    type(quantum_numbers_mask_t), dimension(size(index)) :: mask
    mask = int%mask(index)
  end function interaction_get_mask_slice

\eatline
\nwendcode{}\nwbegindocs{918}Compute the invariant mass squared of the incoming particles (if any,
otherwise outgoing).
\nwenddocs{}\nwbegincode{919}\moddef{Interactions: public}\plusendmoddef
  public :: interaction_get_s
\nwendcode{}\nwbegincode{920}\moddef{Interactions: procedures}\plusendmoddef
  function interaction_get_s (int) result (s)
    real(default) :: s
    type(interaction_t), intent(in) :: int
    if (int%n_in /= 0) then
       s = sum (int%p(:int%n_in)) ** 2
    else
       s = sum (int%p(int%n_vir+1:)) ** 2
    end if
  end function interaction_get_s

\eatline
\nwendcode{}\nwbegindocs{921}Compute the Lorentz transformation that transforms the incoming
particles from the center-of-mass frame to the lab frame where they
are given.  If the c.m. mass squared is negative or zero, return the
identity.
\nwenddocs{}\nwbegincode{922}\moddef{Interactions: public}\plusendmoddef
  public :: interaction_get_cm_transformation
\nwendcode{}\nwbegincode{923}\moddef{Interactions: procedures}\plusendmoddef
  function interaction_get_cm_transformation (int) result (lt)
    type(lorentz_transformation_t) :: lt
    type(interaction_t), intent(in) :: int
    type(vector4_t) :: p_cm
    real(default) :: s
    if (int%n_in /= 0) then
       p_cm = sum (int%p(:int%n_in))
    else
       p_cm = sum (int%p(int%n_vir+1:))
    end if
    s = p_cm ** 2
    if (s > 0) then
       lt = boost (p_cm, sqrt (s))
    else
       lt = identity
    end if
  end function interaction_get_cm_transformation

\eatline
\nwendcode{}\nwbegindocs{924}Return flavor, momentum, and position of the first outgoing
unstable particle present in the interaction.  Note that we need not
iterate through the state matrix; if there is an unstable particle, it
will be present in all state-matrix entries.
\nwenddocs{}\nwbegincode{925}\moddef{Interactions: public}\plusendmoddef
  public :: interaction_get_unstable_particle
\nwendcode{}\nwbegincode{926}\moddef{Interactions: procedures}\plusendmoddef
  subroutine interaction_get_unstable_particle (int, flv, p, i)
    type(interaction_t), intent(in), target :: int
    type(flavor_t), intent(out) :: flv
    type(vector4_t), intent(out) :: p
    integer, intent(out) :: i
    type(state_iterator_t) :: it
    type(flavor_t), dimension(int%n_tot) :: flv_array
    call it%init (int%state_matrix)
    flv_array = it%get_flavor ()
    do i = int%n_in + int%n_vir + 1, int%n_tot
       if (.not. flv_array(i)%is_stable ()) then
          flv = flv_array(i)
          p = int%p(i)
          return
       end if
    end do
  end subroutine interaction_get_unstable_particle

\eatline
\nwendcode{}\nwbegindocs{927}Return the complete set of \emph{outgoing} flavors, assuming that
the flavor quantum number is not suppressed.
\nwenddocs{}\nwbegincode{928}\moddef{Interactions: public}\plusendmoddef
  public :: interaction_get_flv_out
\nwendcode{}\nwbegincode{929}\moddef{Interactions: procedures}\plusendmoddef
  subroutine interaction_get_flv_out (int, flv)
    type(interaction_t), intent(in), target :: int
    type(flavor_t), dimension(:,:), allocatable, intent(out) :: flv
    type(state_iterator_t) :: it
    type(flavor_t), dimension(:), allocatable :: flv_state
    integer :: n_in, n_vir, n_out, n_tot, n_state, i
    n_in = interaction_get_n_in (int)
    n_vir = interaction_get_n_vir (int)
    n_out = interaction_get_n_out (int)
    n_tot = interaction_get_n_tot (int)
    n_state = interaction_get_n_matrix_elements (int)
    allocate (flv (n_out, n_state))
    allocate (flv_state (n_tot))
    i = 1
    call it%init (interaction_get_state_matrix_ptr (int))
    do while (it%is_valid ())
       flv_state = it%get_flavor ()
       flv(:,i) = flv_state(n_in+n_vir+1:)
       i = i + 1
       call it%advance ()
    end do
  end subroutine interaction_get_flv_out
  
\eatline
\nwendcode{}\nwbegindocs{930}Determine the flavor content of the interaction.  We analyze the
state matrix for this, and we select the outgoing particles of the
hard process only for the required mask, which indicates the particles
that can appear in any order in a matching event record.

We have to assume that any radiated particles (beam remnants) appear
at the beginning of the particles marked as outgoing.
\nwenddocs{}\nwbegincode{931}\moddef{Interactions: public}\plusendmoddef
  public :: interaction_get_flv_content
\nwendcode{}\nwbegincode{932}\moddef{Interactions: procedures}\plusendmoddef
  subroutine interaction_get_flv_content (int, state_flv, n_out_hard)
    type(interaction_t), intent(in), target :: int
    type(state_flv_content_t), intent(out) :: state_flv
    integer, intent(in) :: n_out_hard
    logical, dimension(:), allocatable :: mask
    integer :: n_tot
    n_tot = interaction_get_n_tot (int)
    allocate (mask (n_tot), source = .false.)
    mask(n_tot-n_out_hard+1:) = .true.
    call state_flv%fill (interaction_get_state_matrix_ptr (int), mask)
  end subroutine interaction_get_flv_content
  
\eatline
\nwendcode{}\nwbegindocs{933}\nwdocspar
\subsection{Modifying contents}
Set the quantum numbers mask.
\nwenddocs{}\nwbegincode{934}\moddef{Interactions: public}\plusendmoddef
  public :: interaction_set_mask
\nwendcode{}\nwbegincode{935}\moddef{Interactions: procedures}\plusendmoddef
  subroutine interaction_set_mask (int, mask)
    type(interaction_t), intent(inout) :: int
    type(quantum_numbers_mask_t), dimension(:), intent(in) :: mask
    int%mask = mask
    int%update_state_matrix = .true.
  end subroutine interaction_set_mask

\eatline
\nwendcode{}\nwbegindocs{936}Merge a particular mask entry, respecting a possible helicity lock for this
entry.  We apply an OR relation, which means that quantum numbers are
summed over if either of the two masks requires it.
\nwenddocs{}\nwbegincode{937}\moddef{Interactions: procedures}\plusendmoddef
  subroutine interaction_merge_mask_entry (int, i, mask)
    type(interaction_t), intent(inout) :: int
    integer, intent(in) :: i
    type(quantum_numbers_mask_t), intent(in) :: mask
    type(quantum_numbers_mask_t) :: mask_tmp
    integer :: ii
    ii = idx (int, i)
    if (int%mask(ii) .neqv. mask) then
       int%mask(ii) = int%mask(ii) .or. mask
       if (int%hel_lock(ii) /= 0) then
          call mask_tmp%assign (mask, helicity=.true.)
          int%mask(int%hel_lock(ii)) = int%mask(int%hel_lock(ii)) .or. mask_tmp
       end if
    end if
    int%update_state_matrix = .true.
  end subroutine interaction_merge_mask_entry

\eatline
\nwendcode{}\nwbegindocs{938}Fill the momenta array, do not care about the quantum numbers of
particles.
\nwenddocs{}\nwbegincode{939}\moddef{Interactions: public}\plusendmoddef
  public :: interaction_reset_momenta
  public :: interaction_set_momenta
  public :: interaction_set_momentum
\nwendcode{}\nwbegincode{940}\moddef{Interactions: procedures}\plusendmoddef
  subroutine interaction_reset_momenta (int)
    type(interaction_t), intent(inout) :: int
    int%p = vector4_null
    int%p_is_known = .true.
  end subroutine interaction_reset_momenta

  subroutine interaction_set_momenta (int, p, outgoing)
    type(interaction_t), intent(inout) :: int
    type(vector4_t), dimension(:), intent(in) :: p
    logical, intent(in), optional :: outgoing
    integer :: i, index
    do i = 1, size (p)
       index = idx (int, i, outgoing)
       int%p(index) = p(i)
       int%p_is_known(index) = .true.
    end do
  end subroutine interaction_set_momenta

  subroutine interaction_set_momentum (int, p, i, outgoing)
    type(interaction_t), intent(inout) :: int
    type(vector4_t), intent(in) :: p
    integer, intent(in) :: i
    logical, intent(in), optional :: outgoing
    integer :: index
    index = idx (int, i, outgoing)
    int%p(index) = p
    int%p_is_known(index) = .true.
  end subroutine interaction_set_momentum

\eatline
\eatline
\nwendcode{}\nwbegindocs{941}This more sophisticated version of setting values is used for
structure functions, in particular if nontrivial flavor, color, and
helicity may be present: set values selectively for the given flavors.
If there is more than one flavor, scan the interaction and check for a
matching flavor at the specified particle location.  If it matches,
insert the value that corresponds to this flavor.
\nwenddocs{}\nwbegincode{942}\moddef{Interactions: public}\plusendmoddef
  public :: interaction_set_flavored_values
\nwendcode{}\nwbegincode{943}\moddef{Interactions: procedures}\plusendmoddef
  subroutine interaction_set_flavored_values (int, value, flv_in, pos)
    type(interaction_t), intent(inout) :: int
    complex(default), dimension(:), intent(in) :: value
    type(flavor_t), dimension(:), intent(in) :: flv_in
    integer, intent(in) :: pos
    type(state_iterator_t) :: it
    type(flavor_t) :: flv
    integer :: i
    if (size (value) == 1) then
       call interaction_set_matrix_element (int, value(1))
    else
       call it%init (int%state_matrix)
       do while (it%is_valid ())
          flv = it%get_flavor (pos)
          SCAN_FLV: do i = 1, size (value)
             if (flv == flv_in(i)) then
                call it%set_matrix_element (value(i))
                exit SCAN_FLV
             end if
          end do SCAN_FLV
          call it%advance ()
       end do
    end if
  end subroutine interaction_set_flavored_values

\eatline
\nwendcode{}\nwbegindocs{944}\nwdocspar
\subsection{Handling Linked interactions}
Store relations between corresponding particles within one
interaction.  The first particle is the parent, the second one the
child.  Links are established in both directions.

These relations have no effect on the propagation of momenta etc.,
they are rather used for mother-daughter relations in event output.
\nwenddocs{}\nwbegincode{945}\moddef{Interactions: public}\plusendmoddef
  public :: interaction_relate
\nwendcode{}\nwbegincode{946}\moddef{Interactions: procedures}\plusendmoddef
  subroutine interaction_relate (int, i1, i2)
    type(interaction_t), intent(inout), target :: int
    integer, intent(in) :: i1, i2
    if (i1 /= 0 .and. i2 /= 0) then
       call int%children(i1)%append (i2)
       call int%parents(i2)%append (i1)
    end if
  end subroutine interaction_relate

\eatline
\nwendcode{}\nwbegindocs{947}Transfer internal parent-child relations defined within interaction
{\tt{}int1} to a new interaction {\tt{}int} where the particle indices are
mapped to.  Some particles in {\tt{}int1} may have no image in {\tt{}int}.
In that case, a child entry maps to zero, and we skip this relation.

Also transfer resonance flags.
\nwenddocs{}\nwbegincode{948}\moddef{Interactions: public}\plusendmoddef
  public :: interaction_transfer_relations
\nwendcode{}\nwbegincode{949}\moddef{Interactions: procedures}\plusendmoddef
  subroutine interaction_transfer_relations (int1, int2, map)
    type(interaction_t), intent(in) :: int1
    type(interaction_t), intent(inout), target :: int2
    integer, dimension(:), intent(in) :: map
    integer :: i, j, k
    do i = 1, size (map)
       do j = 1, int1%parents(i)%get_length ()
          k = int1%parents(i)%get_link (j)
          call interaction_relate (int2, map(k), map(i))
       end do
       if (map(i) /= 0) then
          int2%resonant(map(i)) = int1%resonant(i)
       end if
    end do
  end subroutine interaction_transfer_relations

\eatline
\nwendcode{}\nwbegindocs{950}Make up internal parent-child relations for the particle(s) that are
connected to a new interaction {\tt{}int}.

If {\tt{}resonant} is defined and true, the connections are marked as
resonant in the result interaction
\nwenddocs{}\nwbegincode{951}\moddef{Interactions: public}\plusendmoddef
  public :: interaction_relate_connections
\nwendcode{}\nwbegincode{952}\moddef{Interactions: procedures}\plusendmoddef
  subroutine interaction_relate_connections &
       (int, int_in, connection_index, &
        map, map_connections, resonant)
    type(interaction_t), intent(inout), target :: int
    type(interaction_t), intent(in) :: int_in
    integer, dimension(:), intent(in) :: connection_index
    integer, dimension(:), intent(in) :: map, map_connections
    logical, intent(in), optional :: resonant
    logical :: reson
    integer :: i, j, i2, k2
    reson = .false.;  if (present (resonant))  reson = resonant 
    do i = 1, size (map_connections)
       k2 = connection_index(i)
       do j = 1, int_in%children(k2)%get_length ()
          i2 = int_in%children(k2)%get_link (j)
          call interaction_relate (int, map_connections(i), map(i2))
       end do
       int%resonant(map_connections(i)) = reson
    end do
  end subroutine interaction_relate_connections

\eatline
\nwendcode{}\nwbegindocs{953}Return the number of source/target links of the internal connections of
particle {\tt{}i}.
\nwenddocs{}\nwbegincode{954}\moddef{Interactions: public}\plusendmoddef
  public :: interaction_get_n_children
  public :: interaction_get_n_parents
\nwendcode{}\nwbegincode{955}\moddef{Interactions: procedures}\plusendmoddef
  function interaction_get_n_children (int, i) result (n)
    integer :: n
    type(interaction_t), intent(in) :: int
    integer, intent(in) :: i
    n = int%children(i)%get_length ()
  end function interaction_get_n_children

  function interaction_get_n_parents (int, i) result (n)
    integer :: n
    type(interaction_t), intent(in) :: int
    integer, intent(in) :: i
    n = int%parents(i)%get_length ()
  end function interaction_get_n_parents

\eatline
\nwendcode{}\nwbegindocs{956}Return the source/target links of the internal connections of
particle {\tt{}i} as an array.
\nwenddocs{}\nwbegincode{957}\moddef{Interactions: public}\plusendmoddef
  public :: interaction_get_children
  public :: interaction_get_parents
\nwendcode{}\nwbegincode{958}\moddef{Interactions: procedures}\plusendmoddef
  function interaction_get_children (int, i) result (idx)
    integer, dimension(:), allocatable :: idx
    type(interaction_t), intent(in) :: int
    integer, intent(in) :: i
    integer :: k, l
    l = int%children(i)%get_length ()
    allocate (idx (l))
    do k = 1, l
       idx(k) = int%children(i)%get_link (k)
    end do
  end function interaction_get_children

  function interaction_get_parents (int, i) result (idx)
    integer, dimension(:), allocatable :: idx
    type(interaction_t), intent(in) :: int
    integer, intent(in) :: i
    integer :: k, l
    l = int%parents(i)%get_length ()
    allocate (idx (l))
    do k = 1, l
       idx(k) = int%parents(i)%get_link (k)
    end do
  end function interaction_get_parents

\eatline
\nwendcode{}\nwbegindocs{959}Add a source link from an interaction to a corresponding particle
within another interaction.  These links affect the propagation of
particles: the two linked particles are considered as the same
particle, outgoing and incoming.
\nwenddocs{}\nwbegincode{960}\moddef{Interactions: public}\plusendmoddef
  public :: interaction_set_source_link
\nwendcode{}\nwbegincode{961}\moddef{Interactions: interfaces}\plusendmoddef
  interface interaction_set_source_link
     module procedure interaction_set_source_link_int
  end interface
\eatline
\nwendcode{}\nwbegincode{962}\moddef{Interactions: procedures}\plusendmoddef
  subroutine interaction_set_source_link_int (int, i, int1, i1)
    type(interaction_t), intent(inout) :: int
    integer, intent(in) :: i
    type(interaction_t), intent(in), target :: int1
    integer, intent(in) :: i1
    if (i /= 0)  call external_link_set (int%source(i), int1, i1)
  end subroutine interaction_set_source_link_int

\eatline
\nwendcode{}\nwbegindocs{963}Reassign links to a new interaction (which is an image of the
current interaction).
\nwenddocs{}\nwbegincode{964}\moddef{Interactions: public}\plusendmoddef
  public :: interaction_reassign_links
\nwendcode{}\nwbegincode{965}\moddef{Interactions: procedures}\plusendmoddef
  subroutine interaction_reassign_links (int, int_src, int_target)
    type(interaction_t), intent(inout) :: int
    type(interaction_t), intent(in) :: int_src
    type(interaction_t), intent(in), target :: int_target
    integer :: i
    if (allocated (int%source)) then
       do i = 1, size (int%source)
          call external_link_reassign (int%source(i), int_src, int_target)
       end do
    end if
  end subroutine interaction_reassign_links

\eatline
\nwendcode{}\nwbegindocs{966}Since links are one-directional, if we want to follow them backwards
we have to scan all possibilities.  This procedure returns the index
of the particle within {\tt{}int} which points to the particle {\tt{}i1}
within interaction {\tt{}int1}.  If unsuccessful, return zero.
\nwenddocs{}\nwbegincode{967}\moddef{Interactions: public}\plusendmoddef
  public :: interaction_find_link
\nwendcode{}\nwbegincode{968}\moddef{Interactions: procedures}\plusendmoddef
  function interaction_find_link (int, int1, i1) result (i)
    integer :: i
    type(interaction_t), intent(in) :: int, int1
    integer, intent(in) :: i1
    type(interaction_t), pointer :: int_tmp
    do i = 1, int%n_tot
       int_tmp => external_link_get_ptr (int%source(i))
       if (int_tmp%tag == int1%tag) then
          if (external_link_get_index (int%source(i)) == i1)  return
       end if
    end do
    i = 0
  end function interaction_find_link

\eatline
\nwendcode{}\nwbegindocs{969}The inverse: return interaction pointer and index for the ultimate source of
{\tt{}i} within {\tt{}int}.
\nwenddocs{}\nwbegincode{970}\moddef{Interactions: public}\plusendmoddef
  public :: interaction_find_source
\nwendcode{}\nwbegincode{971}\moddef{Interactions: procedures}\plusendmoddef
  subroutine interaction_find_source (int, i, int1, i1)
    type(interaction_t), intent(in) :: int
    integer, intent(in) :: i
    type(interaction_t), intent(out), pointer :: int1
    integer, intent(out) :: i1
    type(external_link_t) :: link
    link = interaction_get_ultimate_source (int, i)
    int1 => external_link_get_ptr (link)
    i1 = external_link_get_index (link)
  end subroutine interaction_find_source
    
\eatline
\nwendcode{}\nwbegindocs{972}Follow source links recursively to return the ultimate source of a particle.
\nwenddocs{}\nwbegincode{973}\moddef{Interactions: procedures}\plusendmoddef
  function interaction_get_ultimate_source (int, i) result (link)
    type(external_link_t) :: link
    type(interaction_t), intent(in) :: int
    integer, intent(in) :: i
    type(interaction_t), pointer :: int_src
    integer :: i_src
    link = int%source(i)
    if (external_link_is_set (link)) then
       do
          int_src => external_link_get_ptr (link)
          i_src = external_link_get_index (link)
          if (external_link_is_set (int_src%source(i_src))) then
             link = int_src%source(i_src)
          else
             exit
          end if
       end do
    end if
  end function interaction_get_ultimate_source

\eatline
\nwendcode{}\nwbegindocs{974}Update mask entries by merging them with corresponding masks in
interactions linked to the current one.  The mask determines quantum
numbers which are summed over.

Note that both the mask of the current interaction and the mask of the
linked interaction are updated (side effect!).  This ensures that both
agree for the linked particle.
\nwenddocs{}\nwbegincode{975}\moddef{Interactions: public}\plusendmoddef
  public :: interaction_exchange_mask
\nwendcode{}\nwbegincode{976}\moddef{Interactions: procedures}\plusendmoddef
  subroutine interaction_exchange_mask (int)
    type(interaction_t), intent(inout) :: int
    integer :: i, index_link
    type(interaction_t), pointer :: int_link
    do i = 1, int%n_tot
       if (external_link_is_set (int%source(i))) then
          int_link => external_link_get_ptr (int%source(i))
          index_link = external_link_get_index (int%source(i))
          call interaction_merge_mask_entry &
               (int, i, int_link%mask(index_link))
          call interaction_merge_mask_entry &
               (int_link, index_link, int%mask(i))
       end if
    end do
    call interaction_freeze (int)
  end subroutine interaction_exchange_mask

\eatline
\nwendcode{}\nwbegindocs{977}Copy momenta from interactions linked to the current one.
\nwenddocs{}\nwbegincode{978}\moddef{Interactions: public}\plusendmoddef
  public :: interaction_receive_momenta
\nwendcode{}\nwbegincode{979}\moddef{Interactions: procedures}\plusendmoddef
  subroutine interaction_receive_momenta (int)
    type(interaction_t), intent(inout) :: int
    integer :: i, index_link
    type(interaction_t), pointer :: int_link
    do i = 1, int%n_tot
       if (external_link_is_set (int%source(i))) then
          int_link => external_link_get_ptr (int%source(i))
          index_link = external_link_get_index (int%source(i))
          call interaction_set_momentum (int, int_link%p(index_link), i)
       end if
    end do
  end subroutine interaction_receive_momenta

\eatline
\nwendcode{}\nwbegindocs{980}The inverse operation: Copy momenta back to the interactions linked
to the current one.
\nwenddocs{}\nwbegincode{981}\moddef{Interactions: public}\plusendmoddef
  public :: interaction_send_momenta
\nwendcode{}\nwbegincode{982}\moddef{Interactions: procedures}\plusendmoddef
  subroutine interaction_send_momenta (int)
    type(interaction_t), intent(in) :: int
    integer :: i, index_link
    type(interaction_t), pointer :: int_link
    do i = 1, int%n_tot
       if (external_link_is_set (int%source(i))) then
          int_link => external_link_get_ptr (int%source(i))
          index_link = external_link_get_index (int%source(i))
          call interaction_set_momentum (int_link, int%p(i), index_link)
       end if
    end do
  end subroutine interaction_send_momenta

\eatline
\nwendcode{}\nwbegindocs{983}For numerical comparisons: pacify all momenta in an interaction.
\nwenddocs{}\nwbegincode{984}\moddef{Interactions: public}\plusendmoddef
  public :: interaction_pacify_momenta
\nwendcode{}\nwbegincode{985}\moddef{Interactions: procedures}\plusendmoddef
  subroutine interaction_pacify_momenta (int, acc)
    type(interaction_t), intent(inout) :: int
    real(default), intent(in) :: acc
    integer :: i
    do i = 1, int%n_tot
       call pacify (int%p(i), acc)
    end do
  end subroutine interaction_pacify_momenta

\nwendcode{}\nwbegindocs{986}\nwdocspar
%
\subsection{Recovering connections}
When creating an evaluator for two interactions, we have to know by
which particles they are connected.  The connection indices can be
determined if we have two linked interactions.  We assume that
{\tt{}int1} is the source and {\tt{}int2} the target, so the connections of
interest are stored within {\tt{}int2}.  A connection is found if either the
source is {\tt{}int1}, or the (ultimate)
source of a particle within {\tt{}int2} coincides with the (ultimate) source of a
aparticle within {\tt{}int1}.  The result is an array of
index pairs.

To make things simple, we scan the interaction twice,
once for counting hits, then allocate the array, then scan again and
store the connections.

The connections are scanned for {\tt{}int2}, which has sources in {\tt{}int1}.  It
may happen that the order of connections is interchanged (crossed).  We
require the indices in {\tt{}int1} to be sorted, so we reorder both index arrays
correspondingly before returning them.  (After this, the indices in {\tt{}int2}
may be out of order.)
\nwenddocs{}\nwbegincode{987}\moddef{Interactions: public}\plusendmoddef
  public :: find_connections
\nwendcode{}\nwbegincode{988}\moddef{Interactions: procedures}\plusendmoddef
  subroutine find_connections (int1, int2, n, connection_index)
    type(interaction_t), intent(in) :: int1, int2
    integer, intent(out) :: n
    integer, dimension(:,:), intent(out), allocatable :: connection_index
    integer, dimension(:,:), allocatable :: conn_index_tmp
    integer, dimension(:), allocatable :: ordering
    integer :: i, j, k
    type(external_link_t) :: link2, link1
    type(interaction_t), pointer :: int_link, int_link1
    n = 0
    do i = 1, size (int2%source)
       link2 = interaction_get_ultimate_source (int2, i)
       if (external_link_is_set (link2)) then
          int_link => external_link_get_ptr (link2)
          if (int_link%tag == int1%tag) then
             n = n + 1
          else
             k = external_link_get_index (link2)
             do j = 1, size (int1%source)
                link1 = interaction_get_ultimate_source (int1, j)
                if (external_link_is_set (link1)) then
                   int_link1 => external_link_get_ptr (link1)
                   if (int_link1%tag == int_link%tag) then
                      if (external_link_get_index (link1) == k) then
                         n = n + 1
                      end if
                   end if
                end if
             end do
          end if
       end if
    end do
    allocate (conn_index_tmp (n, 2))
    n = 0
    do i = 1, size (int2%source)
       link2 = interaction_get_ultimate_source (int2, i)
       if (external_link_is_set (link2)) then
          int_link => external_link_get_ptr (link2)
          if (int_link%tag == int1%tag) then
             n = n + 1
             conn_index_tmp(n,1) = external_link_get_index (int2%source(i))
             conn_index_tmp(n,2) = i
          else
             k = external_link_get_index (link2)
             do j = 1, size (int1%source)
                link1 = interaction_get_ultimate_source (int1, j)
                if (external_link_is_set (link1)) then
                   int_link1 => external_link_get_ptr (link1)
                   if (int_link1%tag == int_link%tag) then
                      if (external_link_get_index (link1) == k) then
                         n = n + 1
                         conn_index_tmp(n,1) = j
                         conn_index_tmp(n,2) = i
                      end if
                   end if
                end if
             end do
          end if
       end if
    end do
    allocate (connection_index (n, 2))
    if (n > 1) then
       allocate (ordering (n))
       ordering = order (conn_index_tmp(:,1))
       connection_index = conn_index_tmp(ordering,:)
    else               
       connection_index = conn_index_tmp
    end if
  end subroutine find_connections

\eatline
\nwendcode{}\nwbegindocs{989}\nwdocspar
\subsection{Unit tests}
\nwenddocs{}\nwbegincode{990}\moddef{Interactions: public}\plusendmoddef
  public :: interaction_test
\nwendcode{}\nwbegincode{991}\moddef{Interactions: procedures}\plusendmoddef
  subroutine interaction_test (u, results)
    integer, intent(in) :: u
    type(test_results_t), intent(inout) :: results
  \LA{}Interactions: execute tests\RA{}
  end subroutine interaction_test

\eatline
\nwendcode{}\nwbegindocs{992}Generate an interaction of a polarized virtual photon and a colored
quark which may be either up or down.  Remove the quark polarization.
Generate another interaction for the quark radiating a photon and link
this to the first interation.  The radiation ignores polarization;
transfer this information to the first interaction to simplify it.
Then, transfer the momentum to the radiating quark and perform a
splitting.
\nwenddocs{}\nwbegincode{993}\moddef{Interactions: execute tests}\endmoddef
  call test (interaction_1, "interaction_1", &
       "check interaction setup", &
       u, results)
\nwendcode{}\nwbegincode{994}\moddef{Interactions: tests}\endmoddef
  subroutine interaction_1 (u)
    integer, intent(in) :: u
    type(interaction_t), target :: int, rad
    type(vector4_t), dimension(3) :: p
    type(quantum_numbers_mask_t), dimension(3) :: mask
    p(2) = vector4_moving (500._default, 500._default, 1)
    p(3) = vector4_moving (500._default,-500._default, 1)
    p(1) = p(2) + p(3)

    write (u, "(A)")  "* Test output: interaction"
    write (u, "(A)")  "*   Purpose: check routines for interactions"
    write (u, "(A)")      
           
    call interaction_init (int, 1, 0, 2, set_relations=.true., &
         store_values = .true. )
    call int_set (int, 1, -1, 1, 1, &
         cmplx (0.3_default, 0.1_default, kind=default))
    call int_set (int, 1, -1,-1, 1, &
         cmplx (0.5_default,-0.7_default, kind=default))
    call int_set (int, 1, 1, 1, 1, &
         cmplx (0.1_default, 0._default, kind=default))
    call int_set (int, -1, 1, -1, 2, &
         cmplx (0.4_default, -0.1_default, kind=default))
    call int_set (int, 1, 1, 1, 2, &
         cmplx (0.2_default, 0._default, kind=default))
    call interaction_freeze (int)
    call interaction_set_momenta (int, p)
    mask = quantum_numbers_mask (.false.,.false., [.true.,.true.,.true.])
    call interaction_init (rad, 1, 0, 2, &
         mask=mask, set_relations=.true., store_values = .true.)
    call rad_set (1)
    call rad_set (2)
    call interaction_set_source_link (rad, 1, int, 2)
    call interaction_exchange_mask (rad)
    call interaction_receive_momenta (rad)
    p(1) = interaction_get_momentum (rad, 1)
    p(2) = 0.4_default * p(1)
    p(3) = p(1) - p(2)
    call interaction_set_momenta (rad, p(2:3), outgoing=.true.)
    call interaction_freeze (int)
    call interaction_freeze (rad)
    call interaction_set_matrix_element &
         (rad, cmplx (0._default, 0._default, kind=default))
    call interaction_write (int, u)
    write (u, "(A)")
    call interaction_write (rad, u)
    write (u, "(A)")
    write (u, "(A)")  "* Cleanup"
    call interaction_final (int)
    call interaction_final (rad) 
    write (u, "(A)")
    write (u, "(A)")  "* Test interaction_1: successful."   
  contains
    subroutine int_set (int, h1, h2, hq, q, val)
      type(interaction_t), target, intent(inout) :: int
      integer, intent(in) :: h1, h2, hq, q
      type(flavor_t), dimension(3) :: flv
      type(color_t), dimension(3) :: col
      type(helicity_t), dimension(3) :: hel
      type(quantum_numbers_t), dimension(3) :: qn
      complex(default), intent(in) :: val
      call flv%init ([21, q, -q])
      call col(2)%init_col_acl (5, 0)
      call col(3)%init_col_acl (0, 5)
      call hel%init ([h1, hq, -hq], [h2, hq, -hq])
      call qn%init (flv, col, hel)
      call interaction_add_state (int, qn)
      call interaction_set_matrix_element (int, val)
    end subroutine int_set
    subroutine rad_set (q)
      integer, intent(in) :: q
      type(flavor_t), dimension(3) :: flv
      type(quantum_numbers_t), dimension(3) :: qn
      call flv%init ([ q, q, 21 ])
      call qn%init (flv)
      call interaction_add_state (rad, qn)
    end subroutine rad_set
  end subroutine interaction_1
  
\eatline
\nwendcode{}\nwbegindocs{995}\nwdocspar
%%%%%%%%%%%%%%%%%%%%%%%%%%%%%%%%%%%%%%%%%%%%%%%%%%%%%%%%%%%%%%%%%%%%%%%%
\section{Matrix element evaluation}

The {\tt{}evaluator{\char95}t} type is an extension of the {\tt{}interaction{\char95}t}
type.  It represents either a density matrix as the square of a
transition matrix element, or the product of two density matrices.
Usually, some quantum numbers are summed over in the result.

The {\tt{}interaction{\char95}t} subobject represents a multi-particle
interaction with incoming, virtual, and outgoing particles and the
associated (not necessarily diagonal) density matrix of quantum
state.  When the evaluator is initialized, this interaction is
constructed from the input interaction(s).

In addition, the initialization process sets up a multiplication
table.  For each matrix element of the result, it states which matrix
elements are to be taken from the input interaction(s), multiplied
(optionally, with an additional weight factor) and summed over.

Eventually, to a processes we associate a chain of evaluators which
are to be evaluated sequentially.  The physical event and its matrix
element value(s) can be extracted from the last evaluator in such a
chain.

Evaluators are constructed only once (as long as this is possible)
during an initialization step.  Then, for each event, momenta
are computed and transferred among evaluators using the links within
the interaction subobject.  The multiplication tables enable fast
evaluation of the result without looking at quantum numbers anymore.
\nwenddocs{}\nwbegincode{996}\moddef{\code{}evaluators.f90\edoc{}}\endmoddef
\LA{}File header\RA{}

module evaluators

\LA{}Use kinds\RA{}
\LA{}Use strings\RA{}
  use io_units
  use format_defs, only: FMT_19
  use unit_tests
  use diagnostics
  use lorentz
  use model_data
  use flavors
  use colors
  use helicities
  use quantum_numbers
  use state_matrices
  use interactions

\LA{}Standard module head\RA{}

\LA{}Evaluators: public\RA{}

\LA{}Evaluators: parameters\RA{}

\LA{}Evaluators: types\RA{}

\LA{}Evaluators: interfaces\RA{}

contains

\LA{}Evaluators: procedures\RA{}

\LA{}Evaluators: tests\RA{}

end module evaluators
\eatline
\nwendcode{}\nwbegindocs{997}\nwdocspar
\subsection{Array of pairings}
The evaluator contains an array of {\tt{}pairing{\char95}array} objects.  This
makes up the multiplication table.

Each pairing array contains two list of matrix element indices and a
list of numerical factors.  The matrix element indices correspond to
the input interactions.  The corresponding matrix elements are to be
multiplied and optionally multiplied by a factor.  The results are
summed over to yield one specific matrix element of the result
evaluator.
\nwenddocs{}\nwbegincode{998}\moddef{Evaluators: types}\endmoddef
  type :: pairing_array_t
     integer, dimension(:), allocatable :: i1, i2
     complex(default), dimension(:), allocatable :: factor
  end type pairing_array_t
     
\eatline
\nwendcode{}\nwbegincode{999}\moddef{Evaluators: procedures}\endmoddef
  elemental subroutine pairing_array_init (pa, n, has_i2, has_factor)
    type(pairing_array_t), intent(out) :: pa
    integer, intent(in) :: n
    logical, intent(in) :: has_i2, has_factor
    allocate (pa%i1 (n))
    if (has_i2)  allocate (pa%i2 (n))
    if (has_factor)  allocate (pa%factor (n))
  end subroutine pairing_array_init

\eatline
\nwendcode{}\nwbegindocs{1000}\nwdocspar
\subsection{The evaluator type}
Possible variants of evaluators:
\nwenddocs{}\nwbegincode{1001}\moddef{Evaluators: parameters}\endmoddef
  integer, parameter :: &
       EVAL_UNDEFINED = 0, &
       EVAL_PRODUCT = 1, &
       EVAL_SQUARED_FLOWS = 2, &
       EVAL_SQUARE_WITH_COLOR_FACTORS = 3, &
       EVAL_COLOR_CONTRACTION = 4, &
       EVAL_IDENTITY = 5, &
       EVAL_QN_SUM = 6
\eatline
\eatline
\nwendcode{}\nwbegindocs{1002}The evaluator type contains the result interaction and an array of
pairing lists, one for each matrix element in the result interaction.
\nwenddocs{}\nwbegincode{1003}\moddef{Evaluators: public}\endmoddef
  public :: evaluator_t
\nwendcode{}\nwbegincode{1004}\moddef{Evaluators: types}\plusendmoddef
  type :: evaluator_t
     private
     integer :: type = EVAL_UNDEFINED
     type(interaction_t), pointer :: int_in1 => null ()
     type(interaction_t), pointer :: int_in2 => null ()
     type(interaction_t) :: int
     type(pairing_array_t), dimension(:), allocatable :: pairing_array
   contains
   \LA{}Evaluators: evaluator: TBP\RA{}
  end type evaluator_t

\eatline
\nwendcode{}\nwbegindocs{1005}Output.
\nwenddocs{}\nwbegincode{1006}\moddef{Evaluators: evaluator: TBP}\endmoddef
  procedure :: write => evaluator_write
\nwendcode{}\nwbegincode{1007}\moddef{Evaluators: procedures}\plusendmoddef
  subroutine evaluator_write (eval, unit, &
       verbose, show_momentum_sum, show_mass, show_state, show_table, &
       col_verbose, testflag)
    class(evaluator_t), intent(in) :: eval
    integer, intent(in), optional :: unit
    logical, intent(in), optional :: verbose, show_momentum_sum, show_mass
    logical, intent(in), optional :: show_state, show_table, col_verbose
    logical, intent(in), optional :: testflag
    logical :: conjugate, square, show_tab
    integer :: u, i, j
    u = given_output_unit (unit);  if (u < 0)  return
    show_tab = .true.;  if (present (show_table))  show_tab = .false.
    ! write (u, "(1x,A)")  "Evaluator:"     !!! Debugging
    call interaction_write &
         (eval%int, unit, verbose, show_momentum_sum, show_mass, & 
            show_state, col_verbose, testflag)
    if (show_tab) then
       write (u, "(1x,A)")  "Matrix-element multiplication"
       write (u, "(2x,A)", advance="no")  "Input interaction 1:"
       if (associated (eval%int_in1)) then
          write (u, "(1x,I0)")  interaction_get_tag (eval%int_in1)
       else
          write (u, "(A)")  " [undefined]"
       end if
       write (u, "(2x,A)", advance="no")  "Input interaction 2:"
       if (associated (eval%int_in2)) then
          write (u, "(1x,I0)")  interaction_get_tag (eval%int_in2)
       else
          write (u, "(A)")  " [undefined]"
       end if
       select case (eval%type)
       case (EVAL_SQUARED_FLOWS, EVAL_SQUARE_WITH_COLOR_FACTORS)
          conjugate = .true.
          square = .true.
       case default
          conjugate = .false.
          square = .false.
       end select
       if (eval%type == EVAL_IDENTITY) then
          write (u, "(1X,A)") "Identity evaluator, pairing array unused"
          return
       end if
       if (allocated (eval%pairing_array)) then
          do i = 1, size (eval%pairing_array)
             write (u, "(2x,A,I0,A)")  "ME(", i, ") = "
             do j = 1, size (eval%pairing_array(i)%i1)
                write (u, "(4x,A)", advance="no")  "+"
                if (allocated (eval%pairing_array(i)%i2)) then
                   write (u, "(1x,A,I0,A)", advance="no")  &
                        "ME1(", eval%pairing_array(i)%i1(j), ")"
                   if (conjugate) then
                      write (u, "(A)", advance="no")  "* x"
                   else
                      write (u, "(A)", advance="no")  " x"
                   end if
                   write (u, "(1x,A,I0,A)", advance="no")  &
                        "ME2(", eval%pairing_array(i)%i2(j), ")"
                else if (square) then
                   write (u, "(1x,A)", advance="no")  "|"
                   write (u, "(A,I0,A)", advance="no")  &
                        "ME1(", eval%pairing_array(i)%i1(j), ")"
                   write (u, "(A)", advance="no")  "|^2"
                else
                   write (u, "(1x,A,I0,A)", advance="no")  &
                        "ME1(", eval%pairing_array(i)%i1(j), ")"
                end if
                if (allocated (eval%pairing_array(i)%factor)) then
                   write (u, "(1x,A)", advance="no")  "x"
                   write (u, "(1x,'('," // FMT_19 // ",','," // FMT_19 // &
                        ",')')") eval%pairing_array(i)%factor(j)
                else
                   write (u, *)
                end if
             end do
          end do
       end if
       ! print *, size (eval%pairing_array)     !!! Debugging
    end if
  end subroutine evaluator_write

\eatline
\nwendcode{}\nwbegindocs{1008}Assignment: Deep copy of the interaction component.
\nwenddocs{}\nwbegincode{1009}\moddef{Evaluators: public}\plusendmoddef
  public :: assignment(=)
\nwendcode{}\nwbegincode{1010}\moddef{Evaluators: interfaces}\endmoddef
  interface assignment(=)
     module procedure evaluator_assign
  end interface

\nwendcode{}\nwbegincode{1011}\moddef{Evaluators: procedures}\plusendmoddef
  subroutine evaluator_assign (eval_out, eval_in)
    type(evaluator_t), intent(out) :: eval_out
    type(evaluator_t), intent(in) :: eval_in
    eval_out%type = eval_in%type
    eval_out%int_in1 => eval_in%int_in1
    eval_out%int_in2 => eval_in%int_in2
    eval_out%int = eval_in%int
    if (allocated (eval_in%pairing_array)) then
       allocate (eval_out%pairing_array (size (eval_in%pairing_array)))
       eval_out%pairing_array = eval_in%pairing_array
    end if
  end subroutine evaluator_assign

\eatline
\nwendcode{}\nwbegindocs{1012}\nwdocspar
\subsection{Auxiliary structures for evaluator creation}
Creating an evaluator that properly handles all quantum numbers requires some
bookkeeping.  In this section, we define several auxiliary types and methods
that organize and simplify this task.  More structures are defined within the
specific initializers (as local types and internal subroutines).

These types are currently implemented in a partial object-oriented way: We
define some basic methods for initialization etc.\ here, but the evaluator
routines below do access their internals as well.  This simplifies some things
such as index addressing using array slices, at the expense of losing some
clarity.

\subsubsection{Index mapping}
Index mapping are abundant when constructing an evaluator.  To have arrays of
index mappings, we define this:
\nwenddocs{}\nwbegincode{1013}\moddef{Evaluators: types}\plusendmoddef
  type :: index_map_t
     integer, dimension(:), allocatable :: entry
  end type index_map_t

\eatline
\nwendcode{}\nwbegincode{1014}\moddef{Evaluators: procedures}\plusendmoddef
  elemental subroutine index_map_init (map, n)
    type(index_map_t), intent(out) :: map
    integer, intent(in) :: n
    allocate (map%entry (n))
  end subroutine index_map_init
    
\eatline
\nwendcode{}\nwbegincode{1015}\moddef{Evaluators: procedures}\plusendmoddef
  function index_map_exists (map) result (flag)
    logical :: flag
    type(index_map_t), intent(in) :: map
    flag = allocated (map%entry)
  end function index_map_exists

\eatline
\nwendcode{}\nwbegincode{1016}\moddef{Evaluators: interfaces}\plusendmoddef
  interface size
     module procedure index_map_size
  end interface

\eatline
\nwendcode{}\nwbegincode{1017}\moddef{Evaluators: procedures}\plusendmoddef
  function index_map_size (map) result (s)
    integer :: s
    type(index_map_t), intent(in) :: map
    if (allocated (map%entry)) then
       s = size (map%entry)
    else
       s = 0
    end if
  end function index_map_size

\eatline
\nwendcode{}\nwbegincode{1018}\moddef{Evaluators: interfaces}\plusendmoddef
  interface assignment(=)
     module procedure index_map_assign_int
     module procedure index_map_assign_array
  end interface

\eatline
\nwendcode{}\nwbegincode{1019}\moddef{Evaluators: procedures}\plusendmoddef
  elemental subroutine index_map_assign_int (map, ival)
    type(index_map_t), intent(inout) :: map
    integer, intent(in) :: ival
    map%entry = ival
  end subroutine index_map_assign_int

  subroutine index_map_assign_array (map, array)
    type(index_map_t), intent(inout) :: map
    integer, dimension(:), intent(in) :: array
    map%entry = array
  end subroutine index_map_assign_array

\eatline
\nwendcode{}\nwbegincode{1020}\moddef{Evaluators: procedures}\plusendmoddef
  elemental subroutine index_map_set_entry (map, i, ival)
    type(index_map_t), intent(inout) :: map
    integer, intent(in) :: i
    integer, intent(in) :: ival
    map%entry(i) = ival
  end subroutine index_map_set_entry
    
\eatline
\nwendcode{}\nwbegincode{1021}\moddef{Evaluators: procedures}\plusendmoddef
  elemental function index_map_get_entry (map, i) result (ival)
    integer :: ival
    type(index_map_t), intent(in) :: map
    integer, intent(in) :: i
    ival = map%entry(i)
  end function index_map_get_entry
    
\eatline
\nwendcode{}\nwbegindocs{1022}\nwdocspar
\subsubsection{Index mapping (two-dimensional)}
This is a variant with a square matrix instead of an array.
\nwenddocs{}\nwbegincode{1023}\moddef{Evaluators: types}\plusendmoddef
  type :: index_map2_t
     integer :: s = 0
     integer, dimension(:,:), allocatable :: entry
  end type index_map2_t

\eatline
\nwendcode{}\nwbegincode{1024}\moddef{Evaluators: procedures}\plusendmoddef
  elemental subroutine index_map2_init (map, n)
    type(index_map2_t), intent(out) :: map
    integer, intent(in) :: n
    map%s = n
    allocate (map%entry (n, n))
  end subroutine index_map2_init
    
\eatline
\nwendcode{}\nwbegincode{1025}\moddef{Evaluators: procedures}\plusendmoddef
  function index_map2_exists (map) result (flag)
    logical :: flag
    type(index_map2_t), intent(in) :: map
    flag = allocated (map%entry)
  end function index_map2_exists

\eatline
\nwendcode{}\nwbegincode{1026}\moddef{Evaluators: interfaces}\plusendmoddef
  interface size
     module procedure index_map2_size
  end interface

\eatline
\nwendcode{}\nwbegincode{1027}\moddef{Evaluators: procedures}\plusendmoddef
  function index_map2_size (map) result (s)
    integer :: s
    type(index_map2_t), intent(in) :: map
    s = map%s
  end function index_map2_size

\eatline
\nwendcode{}\nwbegincode{1028}\moddef{Evaluators: interfaces}\plusendmoddef
  interface assignment(=)
     module procedure index_map2_assign_int
  end interface

\eatline
\nwendcode{}\nwbegincode{1029}\moddef{Evaluators: procedures}\plusendmoddef
  elemental subroutine index_map2_assign_int (map, ival)
    type(index_map2_t), intent(inout) :: map
    integer, intent(in) :: ival
    map%entry = ival
  end subroutine index_map2_assign_int

\eatline
\nwendcode{}\nwbegincode{1030}\moddef{Evaluators: procedures}\plusendmoddef
  elemental subroutine index_map2_set_entry (map, i, j, ival)
    type(index_map2_t), intent(inout) :: map
    integer, intent(in) :: i, j
    integer, intent(in) :: ival
    map%entry(i,j) = ival
  end subroutine index_map2_set_entry
    
\eatline
\nwendcode{}\nwbegincode{1031}\moddef{Evaluators: procedures}\plusendmoddef
  elemental function index_map2_get_entry (map, i, j) result (ival)
    integer :: ival
    type(index_map2_t), intent(in) :: map
    integer, intent(in) :: i, j
    ival = map%entry(i,j)
  end function index_map2_get_entry
    
\eatline
\nwendcode{}\nwbegindocs{1032}\nwdocspar
\subsubsection{Auxiliary structures: particle mask}
This is a simple container of a logical array.
\nwenddocs{}\nwbegincode{1033}\moddef{Evaluators: types}\plusendmoddef
  type :: prt_mask_t
     logical, dimension(:), allocatable :: entry
  end type prt_mask_t

\eatline
\nwendcode{}\nwbegincode{1034}\moddef{Evaluators: procedures}\plusendmoddef
  subroutine prt_mask_init (mask, n)
    type(prt_mask_t), intent(out) :: mask
    integer, intent(in) :: n
    allocate (mask%entry (n))
  end subroutine prt_mask_init

\eatline
\nwendcode{}\nwbegincode{1035}\moddef{Evaluators: interfaces}\plusendmoddef
  interface size
     module procedure prt_mask_size
  end interface

\eatline
\nwendcode{}\nwbegincode{1036}\moddef{Evaluators: procedures}\plusendmoddef
  function prt_mask_size (mask) result (s)
    integer :: s
    type(prt_mask_t), intent(in) :: mask
    s = size (mask%entry)
  end function prt_mask_size

\eatline
\nwendcode{}\nwbegindocs{1037}\nwdocspar
\subsubsection{Quantum number containers}
Trivial transparent containers:
\nwenddocs{}\nwbegincode{1038}\moddef{Evaluators: types}\plusendmoddef
  type :: qn_list_t
     type(quantum_numbers_t), dimension(:,:), allocatable :: qn
  end type qn_list_t

  type :: qn_mask_array_t
     type(quantum_numbers_mask_t), dimension(:), allocatable :: mask
  end type qn_mask_array_t

\eatline
\nwendcode{}\nwbegindocs{1039}\nwdocspar
\subsubsection{Auxiliary structures: connection entries}
This type is used as intermediate storage when computing the product of two
evaluators or the square of an evaluator.  The quantum-number array {\tt{}qn}
corresponds to the particles common to both interactions, but irrelevant
quantum numbers (color) masked out.  The index arrays {\tt{}index{\char95}in} determine
the entries in the input interactions that contribute to this connection.
{\tt{}n{\char95}index} is the size of these arrays, and {\tt{}count} is used while filling
the entries.  Finally, the quantum-number arrays {\tt{}qn{\char95}in{\char95}list} are the actual
entries in the input interaction that contribute.  In the product case, they
exclude the connected quantum numbers.

Each evaluator has its own {\tt{}connection{\char95}table} which contains an array of
{\tt{}connection{\char95}entry} objects, but also has stuff that specifically applies to
the evaluator type.  Hence, we do not generalize the {\tt{}connection{\char95}table{\char95}t}
type. 

The filling procedure {\tt{}connection{\char95}entry{\char95}add{\char95}state} is specific to the
various evaluator types.
\nwenddocs{}\nwbegincode{1040}\moddef{Evaluators: types}\plusendmoddef
  type :: connection_entry_t
     type(quantum_numbers_t), dimension(:), allocatable :: qn_conn
     integer, dimension(:), allocatable :: n_index
     integer, dimension(:), allocatable :: count
     type(index_map_t), dimension(:), allocatable :: index_in
     type(qn_list_t), dimension(:), allocatable :: qn_in_list
  end type connection_entry_t

\eatline
\nwendcode{}\nwbegincode{1041}\moddef{Evaluators: procedures}\plusendmoddef
  subroutine connection_entry_init &
      (entry, n_count, n_map, qn_conn, count, n_rest)
    type(connection_entry_t), intent(out) :: entry
    integer, intent(in) :: n_count, n_map
    type(quantum_numbers_t), dimension(:), intent(in) :: qn_conn
    integer, dimension(n_count), intent(in) :: count
    integer, dimension(n_count), intent(in) :: n_rest
    integer :: i
    allocate (entry%qn_conn (size (qn_conn)))
    allocate (entry%n_index (n_count))
    allocate (entry%count (n_count))
    allocate (entry%index_in (n_map))
    allocate (entry%qn_in_list (n_count))
    entry%qn_conn = qn_conn
    entry%n_index = count
    entry%count = 0
    if (size (entry%index_in) == size (count)) then
       call index_map_init (entry%index_in, count)
    else
       call index_map_init (entry%index_in, count(1))
    end if
    do i = 1, n_count
      allocate (entry%qn_in_list(i)%qn (n_rest(i), count(i)))
    end do
  end subroutine connection_entry_init

\eatline
\nwendcode{}\nwbegincode{1042}\moddef{Evaluators: procedures}\plusendmoddef
  subroutine connection_entry_write (entry, unit)
    type(connection_entry_t), intent(in) :: entry
    integer, intent(in), optional :: unit
    integer :: i, j
    integer :: u
    u = given_output_unit (unit)
    call quantum_numbers_write (entry%qn_conn, unit)
    write (u, *)
    do i = 1, size (entry%n_index)
       write (u, *)  "Input interaction", i
       do j = 1, entry%n_index(i)
          if (size (entry%n_index) == size (entry%index_in)) then
             write (u, "(2x,I0,4x,I0,2x)", advance = "no") &
                  j, index_map_get_entry (entry%index_in(i), j)
          else
             write (u, "(2x,I0,4x,I0,2x,I0,2x)", advance = "no") &
                  j, index_map_get_entry (entry%index_in(1), j), &
                     index_map_get_entry (entry%index_in(2), j)
          end if
          call quantum_numbers_write (entry%qn_in_list(i)%qn(:,j), unit)
          write (u, *)
       end do
    end do
  end subroutine connection_entry_write

\eatline
\nwendcode{}\nwbegindocs{1043}\nwdocspar
\subsubsection{Color handling}
For managing color-factor computation, we introduce this local type.   The
{\tt{}index} is the index in the color table that corresponds to a given matrix
element index in the input interaction.  The {\tt{}col} array stores the color
assignments in rows.  The {\tt{}factor} array associates a complex number with
each pair of arrays in the color table.  The {\tt{}factor{\char95}is{\char95}known} array reveals
whether a given factor is known already or still has to be computed.
\nwenddocs{}\nwbegincode{1044}\moddef{Evaluators: types}\plusendmoddef
  type :: color_table_t
     integer, dimension(:), allocatable :: index
     type(color_t), dimension(:,:), allocatable :: col
     logical, dimension(:,:), allocatable :: factor_is_known
     complex(default), dimension(:,:), allocatable :: factor
  end type color_table_t

\eatline
\nwendcode{}\nwbegindocs{1045}This is the initializer.  We extract the color states from the given state
matrices, establish index mappings between the two states (implemented by the
array {\tt{}me{\char95}index}), make an array of color states, and initialize the
color-factor table.  The latter is two-dimensional (includes interference) and
not yet filled.
\nwenddocs{}\nwbegincode{1046}\moddef{Evaluators: procedures}\plusendmoddef
  subroutine color_table_init (color_table, state, n_tot)
    type(color_table_t), intent(out) :: color_table
    type(state_matrix_t), intent(in) :: state
    integer, intent(in) :: n_tot
    type(state_iterator_t) :: it
    type(quantum_numbers_t), dimension(:), allocatable :: qn
    type(state_matrix_t) :: state_col
    integer :: index, n_col_state
    allocate (color_table%index (state%get_n_matrix_elements ()))
    color_table%index = 0
    allocate (qn (n_tot))
    call state_col%init ()
    call it%init (state)
    do while (it%is_valid ())
       index = it%get_me_index ()
       call qn%init (col = it%get_color ())
       call state_col%add_state (qn, me_index = color_table%index(index))
       call it%advance ()
    end do
    n_col_state = state_col%get_n_matrix_elements ()
    allocate (color_table%col (n_tot, n_col_state))
    call it%init (state_col)
    do while (it%is_valid ())
       index = it%get_me_index ()
       color_table%col(:,index) = it%get_color ()
       call it%advance ()
    end do
    call state_col%final ()
    allocate (color_table%factor_is_known (n_col_state, n_col_state))
    allocate (color_table%factor (n_col_state, n_col_state))
    color_table%factor_is_known = .false.
  end subroutine color_table_init

\eatline
\nwendcode{}\nwbegindocs{1047}Output (debugging use):
\nwenddocs{}\nwbegincode{1048}\moddef{Evaluators: procedures}\plusendmoddef
  subroutine color_table_write (color_table, unit)
    type(color_table_t), intent(in) :: color_table
    integer, intent(in), optional :: unit
    integer :: i, j
    integer :: u
    u = given_output_unit (unit)
    write (u, *) "Color table:"
    if (allocated (color_table%index)) then
       write (u, *) "  Index mapping state => color table:"
       do i = 1, size (color_table%index)
          write (u, "(3x,I0,2x,I0,2x)")  i, color_table%index(i)
       end do
       write (u, *) "  Color table:"
       do i = 1, size (color_table%col, 2)
          write (u, "(3x,I0,2x)", advance = "no")  i
          call color_write (color_table%col(:,i), unit)
          write (u, *)
       end do
       write (u, *) "  Defined color factors:"
       do i = 1, size (color_table%factor, 1)
          do j = 1, size (color_table%factor, 2)
             if (color_table%factor_is_known(i,j)) then
                write (u, *)  i, j, color_table%factor(i,j)
             end if
          end do
       end do
    end if
  end subroutine color_table_write

\eatline
\nwendcode{}\nwbegindocs{1049}This subroutine sets color factors, based on information from the hard
matrix element: the list of pairs of color-flow indices (in the basis of the
matrix element code), the list of corresponding factors, and the list of
mappings from the matrix element index in the input interaction to the
color-flow index in the hard matrix element object.

We first determine the mapping of color-flow indices from the hard matrix
element code to the current color table.  The mapping could be nontrivial
because the latter is derived from iterating over a state matrix, which may
return states in non-canonical order.  The translation table can be determined
because we have, for the complete state matrix, both the mapping to the hard
interaction (the input {\tt{}col{\char95}index{\char95}hi}) and the mapping to the current
color table (the component {\tt{}color{\char95}table{\char37}index}).

Once this mapping is known, we scan the list of index pairs
{\tt{}color{\char95}flow{\char95}index} and translate them to valid color-table index pairs.  For
this pair, the color factor is set using the corresponding entry in the list
{\tt{}col{\char95}factor}.
\nwenddocs{}\nwbegincode{1050}\moddef{Evaluators: procedures}\plusendmoddef
  subroutine color_table_set_color_factors (color_table, &
       col_flow_index, col_factor, col_index_hi)
    type(color_table_t), intent(inout) :: color_table
    integer, dimension(:,:), intent(in) :: col_flow_index
    complex(default), dimension(:), intent(in) :: col_factor
    integer, dimension(:), intent(in) :: col_index_hi
    integer, dimension(:), allocatable :: hi_to_ct
    integer :: n_cflow
    integer :: hi_index, me_index, ct_index, cf_index
    integer, dimension(2) :: hi_index_pair, ct_index_pair
    n_cflow = size (col_index_hi)
    if (size (color_table%index) /= n_cflow) &
         call msg_bug ("Mismatch between hard matrix element and color table")
    allocate (hi_to_ct (n_cflow))
    do me_index = 1, size (color_table%index)
       ct_index = color_table%index(me_index)
       hi_index = col_index_hi(me_index)
       hi_to_ct(hi_index) = ct_index
    end do
    do cf_index = 1, size (col_flow_index, 2)
       hi_index_pair = col_flow_index(:,cf_index)
       ct_index_pair = hi_to_ct(hi_index_pair)
       color_table%factor(ct_index_pair(1), ct_index_pair(2)) = &
            col_factor(cf_index)
       color_table%factor_is_known(ct_index_pair(1), ct_index_pair(2)) = .true.
    end do
  end subroutine color_table_set_color_factors

\eatline
\nwendcode{}\nwbegindocs{1051}This function returns a color factor, given two indices which point to the
matrix elements of the initial state matrix.  Internally, we can map them to
the corresponding indices in the color table.  As a side effect, we store the
color factor in the color table for later lookup.  (I.e., this function is
impure.)
\nwenddocs{}\nwbegincode{1052}\moddef{Evaluators: procedures}\plusendmoddef
  function color_table_get_color_factor (color_table, index1, index2, nc) &
      result (factor)
    real(default) :: factor
    type(color_table_t), intent(inout) :: color_table
    integer, intent(in) :: index1, index2
    integer, intent(in), optional :: nc
    integer :: i1, i2
    ! print *, "compute color factor ", index1, index2   !!! Debugging
    i1 = color_table%index(index1)
    i2 = color_table%index(index2)
    ! print *, "  indices = ", i1, i2                    !!! Debugging
    if (color_table%factor_is_known(i1,i2)) then
       factor = color_table%factor(i1,i2)
       ! print *, "  is known : ", factor                !!! Debugging
    else
       factor = compute_color_factor &
            (color_table%col(:,i1), color_table%col(:,i2), nc)
       color_table%factor(i1,i2) = factor
       color_table%factor_is_known(i1,i2) = .true.
       ! print *, "  computed : ", factor                !!! Debugging 
    end if
  end function color_table_get_color_factor

\eatline
\nwendcode{}\nwbegindocs{1053}\nwdocspar
\subsection{Creating an evaluator: Matrix multiplication}
The evaluator for matrix multiplication is the most complicated
variant.

The initializer takes two input interactions and constructs the result
evaluator, which consists of the interaction and the multiplication
table for the product (or convolution) of the two.  Normally, the
input interactions are connected by one or more common particles
(e.g., decay, structure function convolution).  

In the result interaction, quantum numbers of the connections can be
summed over.  This is determined by the {\tt{}qn{\char95}mask{\char95}conn} argument.
The {\tt{}qn{\char95}mask{\char95}rest} argument is its analog for the other particles
within the result interaction.  (E.g., for the trace of the state
matrix, all quantum numbers are summed over.)  Finally, the
{\tt{}connections{\char95}are{\char95}resonant} argument tells whether the connecting
particles should be marked as resonant in the final event record.
This is useful for decays.

The algorithm consists of the following steps:
\begin{enumerate}
\item
  {\tt{}find{\char95}connections}: Find the particles which are connected, i.e.,
  common to both input interactions.  Either they are directly linked,
  or both are linked to a common source.
\item
  {\tt{}compute{\char95}index{\char95}bounds{\char95}and{\char95}mappings}: Compute the mappings of
  particle indices from the input interactions to the result
  interaction.  There is a separate mapping for the connected
  particles.
\item
  {\tt{}accumulate{\char95}connected{\char95}states}: Create an auxiliary state matrix
  which lists the possible quantum numbers for the connected
  particles.  When building this matrix, count the number of times
  each assignment is contained in any of the input states and, for
  each of the input states, record the index of the matrix element
  within the new state matrix.  For the connected particles, reassign
  color indices such that no color state is present twice in different
  color-index assignment.  Note that helicity assignments of the
  connected state can be (and will be) off-diagonal, so no spin
  correlations are lost in decays.

  Do this for both input interactions.
\item
  {\tt{}allocate{\char95}connection{\char95}entries}: Allocate a table of connections.
  Each table row corresponds to one state in the auxiliary matrix, and
  to multiple states of the input interactions.  It collects all
  states of the unconnected particles in the two input interactions
  that are associated with the particular state (quantum-number
  assignment) of the connected particles.
\item
  {\tt{}fill{\char95}connection{\char95}table}: Fill the table of connections by scanning
  both input interactions.  When copying states, reassign color
  indices for the unconnected particles such that they match between
  all involved particle sets (interaction 1, interaction 2, and
  connected particles).
\item
  {\tt{}make{\char95}product{\char95}interaction}: Scan the table of connections we have
  just built.  For each entry, construct all possible pairs of states
  of the unconnected particles and combine them with the specific
  connected-particle state.  This is a possible quantum-number
  assignment of the result interaction.  Now mask all quantum numbers
  that should be summed over, and append this to the result state
  matrix.  Record the matrix element index of the result.  We now have
  the result interaction.
\item
  {\tt{}make{\char95}pairing{\char95}array}: First allocate the pairing array with the
  number of entries of the result interaction.  Then scan the table of
  connections again.  For each entry, record the indices of the matrix
  elements which have to be multiplied and summed over in order to
  compute this particular matrix element.  This makes up the
  multiplication table.
\item
  {\tt{}record{\char95}links}: Transfer all source pointers from the input
  interactions to the result interaction.  Do the same for the
  internal parent-child relations and resonance assignments.  For the
  connected particles, make up appropriate additional parent-child
  relations.  This allows for fetching momenta from other interactions
  when a new event is filled, and to reconstruct the event history
  when the event is analyzed.
\end{enumerate}

After all this is done, for each event, we just have to evaluate the
pairing arrays (multiplication tables) in order to compute the result
matrix elements in their proper positions.  The quantum-number
assignments remain fixed from now on.
\nwenddocs{}\nwbegincode{1054}\moddef{Evaluators: public}\plusendmoddef
  public :: evaluator_init_product
\nwendcode{}\nwbegincode{1055}\moddef{Evaluators: interfaces}\plusendmoddef
  interface evaluator_init_product
     module procedure evaluator_init_product_ii
     module procedure evaluator_init_product_ie
     module procedure evaluator_init_product_ei
     module procedure evaluator_init_product_ee
  end interface

\nwendcode{}\nwbegincode{1056}\moddef{Evaluators: procedures}\plusendmoddef
  subroutine evaluator_init_product_ii &
       (eval, int_in1, int_in2, qn_mask_conn, qn_filter_conn, qn_mask_rest, &
        connections_are_resonant)

    type(evaluator_t), intent(out), target :: eval
    type(interaction_t), intent(in), target :: int_in1, int_in2
    type(quantum_numbers_mask_t), intent(in) :: qn_mask_conn
    type(quantum_numbers_t), intent(in), optional :: qn_filter_conn
    type(quantum_numbers_mask_t), intent(in), optional :: qn_mask_rest
    logical, intent(in), optional :: connections_are_resonant

    type(qn_mask_array_t), dimension(2) :: qn_mask_in
    type(state_matrix_t), pointer :: state_in1, state_in2

    type :: connection_table_t
       integer :: n_conn = 0
       integer, dimension(2) :: n_rest = 0
       integer :: n_tot = 0
       integer :: n_me_conn = 0
       type(state_matrix_t) :: state
       type(index_map_t), dimension(:), allocatable :: index_conn
       type(connection_entry_t), dimension(:), allocatable :: entry
       type(index_map_t) :: index_result
    end type connection_table_t
    type(connection_table_t) :: connection_table

    integer :: n_in, n_vir, n_out, n_tot
    integer, dimension(2) :: n_rest
    integer :: n_conn

    integer, dimension(:,:), allocatable :: connection_index
    type(index_map_t), dimension(2) :: prt_map_in
    type(index_map_t) :: prt_map_conn
    type(prt_mask_t), dimension(2) :: prt_is_connected
    type(quantum_numbers_mask_t), dimension(:), allocatable :: &
         qn_mask_conn_initial

    integer :: i

    eval%type = EVAL_PRODUCT
    eval%int_in1 => int_in1
    eval%int_in2 => int_in2
    ! print *, "Evaluator product"          !!! Debugging      
    ! print *, "First interaction"          !!! Debugging      
    ! call interaction_write (int_in1)      !!! Debugging      
    ! print *                               !!! Debugging      
    ! print *, "Second interaction"         !!! Debugging      
    ! call interaction_write (int_in2)      !!! Debugging      
    ! print *                               !!! Debugging      

    state_in1 => interaction_get_state_matrix_ptr (int_in1)
    state_in2 => interaction_get_state_matrix_ptr (int_in2)

    call find_connections (int_in1, int_in2, n_conn, connection_index)
    if (n_conn == 0) then
       call msg_message ("First interaction:")
       call interaction_write (int_in1, col_verbose=.true.)
       call msg_message ("Second interaction:")
       call interaction_write (int_in2, col_verbose=.true.)
       call msg_fatal ("Evaluator product: no connections found between factors")
    end if
    call compute_index_bounds_and_mappings &
         (int_in1, int_in2, n_conn, &
          n_in, n_vir, n_out, n_tot, &
          n_rest, prt_map_in, prt_map_conn)

    call prt_mask_init (prt_is_connected(1), interaction_get_n_tot (int_in1))
    call prt_mask_init (prt_is_connected(2), interaction_get_n_tot (int_in2))
    do i = 1, 2
       prt_is_connected(i)%entry = .true.
       prt_is_connected(i)%entry(connection_index(:,i)) = .false.
    end do
    allocate (qn_mask_conn_initial (n_conn))
    qn_mask_conn_initial = &
         interaction_get_mask (int_in1, connection_index(:,1)) .or. &
         interaction_get_mask (int_in2, connection_index(:,2))
    allocate (qn_mask_in(1)%mask (interaction_get_n_tot (int_in1)))
    allocate (qn_mask_in(2)%mask (interaction_get_n_tot (int_in2)))
    qn_mask_in(1)%mask = interaction_get_mask (int_in1)
    qn_mask_in(2)%mask = interaction_get_mask (int_in2)

    call connection_table_init (connection_table, &
         state_in1, state_in2, &
         qn_mask_conn_initial,  &
         n_conn, connection_index, n_rest, &
         qn_filter_conn)
    call connection_table_fill (connection_table, &
         state_in1, state_in2, &
         connection_index, prt_is_connected)
    call make_product_interaction (eval%int, &
         n_in, n_vir, n_out, &
         connection_table, &
         prt_map_in, prt_is_connected, &
         qn_mask_in, qn_mask_conn_initial, &
         qn_mask_conn, qn_filter_conn, qn_mask_rest)
    ! call connection_table_write (connection_table)    !!! Debugging
    call make_pairing_array (eval%pairing_array, &
         interaction_get_n_matrix_elements (eval%int), &
         connection_table)
    call record_links (eval%int, &
         int_in1, int_in2, connection_index, prt_map_in, prt_map_conn, &
         prt_is_connected, connections_are_resonant)
    call connection_table_final (connection_table)

    ! print *, "Result evaluator"                !!! Debugging
    ! call evaluator_write (eval)                !!! Debugging

    if (interaction_get_n_matrix_elements (eval%int) == 0) then
       print *, "Evaluator product"
       print *, "First interaction"
       call interaction_write (int_in1, col_verbose=.true.)
       print *
       print *, "Second interaction"
       call interaction_write (int_in2, col_verbose=.true.)
       print *
       call msg_fatal ("Product of density matrices is empty", &
           [var_str ("   --------------------------------------------"), &
            var_str ("This happens when two density matrices are convoluted "), &
            var_str ("but the processes they belong to (e.g., production "), &
            var_str ("and decay) do not match. This could happen if the "), &
            var_str ("beam specification does not match the hard "), &
            var_str ("process. Or it may indicate a WHIZARD bug.")])
    end if

  contains

    subroutine compute_index_bounds_and_mappings &
         (int1, int2, n_conn, &
          n_in, n_vir, n_out, n_tot, &
          n_rest, prt_map_in, prt_map_conn)
      type(interaction_t), intent(in) :: int1, int2
      integer, intent(in) :: n_conn
      integer, intent(out) :: n_in, n_vir, n_out, n_tot
      integer, dimension(2), intent(out) :: n_rest
      type(index_map_t), dimension(2), intent(out) :: prt_map_in
      type(index_map_t), intent(out) :: prt_map_conn
      integer, dimension(:), allocatable :: index
      integer :: n_in1, n_vir1, n_out1
      integer :: n_in2, n_vir2, n_out2
      integer :: k
      n_in1  = interaction_get_n_in  (int1)
      n_vir1 = interaction_get_n_vir (int1)
      n_out1 = interaction_get_n_out (int1) - n_conn
      n_rest(1) = n_in1 + n_vir1 + n_out1
      n_in2  = interaction_get_n_in  (int2) - n_conn
      n_vir2 = interaction_get_n_vir (int2)
      n_out2 = interaction_get_n_out (int2)
      n_rest(2) = n_in2 + n_vir2 + n_out2
      n_in  = n_in1  + n_in2
      n_vir = n_vir1 + n_vir2 + n_conn
      n_out = n_out1 + n_out2
      n_tot = n_in + n_vir + n_out
      call index_map_init (prt_map_in, n_rest)
      call index_map_init (prt_map_conn, n_conn)
      allocate (index (n_tot))
      index = [ (i, i = 1, n_tot) ]
      prt_map_in(1)%entry(1 : n_in1) = index(  1 :   n_in1)
      k =     n_in1
      prt_map_in(2)%entry(1 : n_in2) = index(k+1 : k+n_in2)
      k = k + n_in2
      prt_map_in(1)%entry(n_in1+1 : n_in1+n_vir1) = index(k+1 : k+n_vir1)
      k = k + n_vir1
      prt_map_in(2)%entry(n_in2+1 : n_in2+n_vir2) = index(k+1 : k+n_vir2)
      k = k + n_vir2
      prt_map_conn%entry = index(k+1 : k+n_conn)
      k = k + n_conn
      prt_map_in(1)%entry(n_in1+n_vir1+1 : n_rest(1)) = index(k+1 : k+n_out1)
      k = k + n_out1
      prt_map_in(2)%entry(n_in2+n_vir2+1 : n_rest(2)) = index(k+1 : k+n_out2)
    end subroutine compute_index_bounds_and_mappings

    subroutine connection_table_init &
        (connection_table, state_in1, state_in2, qn_mask_conn, &
         n_conn, connection_index, n_rest, &
         qn_filter_conn)
      type(connection_table_t), intent(out) :: connection_table
      type(state_matrix_t), intent(in), target :: state_in1, state_in2
      type(quantum_numbers_mask_t), dimension(:), intent(in) :: qn_mask_conn
      integer, intent(in) :: n_conn
      integer, dimension(:,:), intent(in) :: connection_index
      integer, dimension(2), intent(in) :: n_rest
      type(quantum_numbers_t), intent(in), optional :: qn_filter_conn
      integer, dimension(2) :: n_me_in
      type(state_iterator_t) :: it
      type(quantum_numbers_t), dimension(n_conn) :: qn
      integer :: i, me_index_in, me_index_conn, n_me_conn
      integer, dimension(2) :: me_count
      connection_table%n_conn = n_conn
      connection_table%n_rest = n_rest
      n_me_in(1) = state_in1%get_n_matrix_elements ()
      n_me_in(2) = state_in2%get_n_matrix_elements ()
      allocate (connection_table%index_conn (2))
      call index_map_init (connection_table%index_conn, n_me_in)
      connection_table%index_conn = 0
      call connection_table%state%init (n_counters=2)
      do i = 1, 2
         select case (i)
         case (1);  call it%init (state_in1)
         case (2);  call it%init (state_in2)
         end select
         do while (it%is_valid ())
            qn = it%get_quantum_numbers (connection_index(:,i))
            call qn%undefine (qn_mask_conn)
            if (present (qn_filter_conn)) then
               if (.not. all (qn .match. qn_filter_conn)) then
                  call it%advance ();  cycle
               end if
            end if
            call quantum_numbers_canonicalize_color (qn)
            me_index_in = it%get_me_index ()
            call connection_table%state%add_state (qn, &
                counter_index = i, me_index = me_index_conn)
            call index_map_set_entry (connection_table%index_conn(i), &
                 me_index_in, me_index_conn)
            call it%advance ()
         end do
      end do
      n_me_conn = connection_table%state%get_n_matrix_elements ()
      connection_table%n_me_conn = n_me_conn
      allocate (connection_table%entry (n_me_conn))
      call it%init (connection_table%state)
      do while (it%is_valid ())
         i = it%get_me_index ()
         me_count = it%get_me_count ()
         call connection_entry_init (connection_table%entry(i), 2, 2, &
              it%get_quantum_numbers (), me_count, n_rest)
         call it%advance ()
      end do
    end subroutine connection_table_init

    subroutine connection_table_final (connection_table)
      type(connection_table_t), intent(inout) :: connection_table
      call connection_table%state%final ()
    end subroutine connection_table_final
   
    subroutine connection_table_write (connection_table, unit)
      type(connection_table_t), intent(in) :: connection_table
      integer, intent(in), optional :: unit
      integer :: i, j
      integer :: u
      u = given_output_unit (unit)
      write (u, *) "Connection table:"
      call connection_table%state%write (unit)
      if (allocated (connection_table%index_conn)) then
         write (u, *) "  Index mapping input => connection table:"
         do i = 1, size (connection_table%index_conn)
            write (u, *) "  Input state", i
            do j = 1, size (connection_table%index_conn(i))
               write (u, *)  j, &
                    index_map_get_entry (connection_table%index_conn(i), j)
            end do
         end do
      end if
      if (allocated (connection_table%entry)) then
         write (u, *) "  Connection table contents:"
         do i = 1, size (connection_table%entry)
            call connection_entry_write (connection_table%entry(i), unit)
         end do
      end if
      if (index_map_exists (connection_table%index_result)) then
         write (u, *) "  Index mapping connection table => output:"
         do i = 1, size (connection_table%index_result)
            write (u, *)  i, &
                 index_map_get_entry (connection_table%index_result, i)
         end do
      end if
    end subroutine connection_table_write

    subroutine connection_table_fill &
        (connection_table, state_in1, state_in2, &
         connection_index, prt_is_connected)
      type(connection_table_t), intent(inout) :: connection_table
      type(state_matrix_t), intent(in), target :: state_in1, state_in2
      integer, dimension(:,:), intent(in) :: connection_index
      type(prt_mask_t), dimension(2), intent(in) :: prt_is_connected
      type(state_iterator_t) :: it
      integer :: index_in, index_conn
      integer :: color_offset
      integer :: n_result_entries
      integer :: i, k
      color_offset = connection_table%state%get_max_color_value ()
      do i = 1, 2
         select case (i)
         case (1);  call it%init (state_in1)
         case (2);  call it%init (state_in2)
         end select
         do while (it%is_valid ())
            index_in = it%get_me_index ()
            index_conn = index_map_get_entry &
                              (connection_table%index_conn(i), index_in)
            if (index_conn /= 0) then
               call connection_entry_add_state &
                    (connection_table%entry(index_conn), i, &
                     index_in, it%get_quantum_numbers (), &
                     connection_index(:,i), prt_is_connected(i), &
                     color_offset)
            end if
            call it%advance ()
         end do
         color_offset = color_offset + state_in1%get_max_color_value ()
      end do
      n_result_entries = 0
      do k = 1, size (connection_table%entry)
         n_result_entries = &
              n_result_entries + product (connection_table%entry(k)%n_index)
      end do
      call index_map_init (connection_table%index_result, n_result_entries)
    end subroutine connection_table_fill
      
    subroutine connection_entry_add_state &
        (entry, i, index_in, qn_in, connection_index, prt_is_connected, &
         color_offset)
      type(connection_entry_t), intent(inout) :: entry
      integer, intent(in) :: i
      integer, intent(in) :: index_in
      type(quantum_numbers_t), dimension(:), intent(in) :: qn_in
      integer, dimension(:), intent(in) :: connection_index
      type(prt_mask_t), intent(in) :: prt_is_connected
      integer, intent(in) :: color_offset
      integer :: c
      integer, dimension(:,:), allocatable :: color_map
      entry%count(i) = entry%count(i) + 1
      c = entry%count(i)
      call make_color_map (color_map, &
           qn_in(connection_index), entry%qn_conn)
      call index_map_set_entry (entry%index_in(i), c, index_in)
      entry%qn_in_list(i)%qn(:,c) = pack (qn_in, prt_is_connected%entry)
      call quantum_numbers_translate_color &
           (entry%qn_in_list(i)%qn(:,c), color_map, color_offset)
    end subroutine connection_entry_add_state

    subroutine make_product_interaction (int, &
         n_in, n_vir, n_out, &
         connection_table, &
         prt_map_in, prt_is_connected, &
         qn_mask_in, qn_mask_conn_initial, &
         qn_mask_conn, qn_filter_conn, qn_mask_rest)
      type(interaction_t), intent(out), target :: int
      integer, intent(in) :: n_in, n_vir, n_out
      type(connection_table_t), intent(inout), target :: connection_table
      type(index_map_t), dimension(2), intent(in) :: prt_map_in
      type(prt_mask_t), dimension(2), intent(in) :: prt_is_connected
      type(qn_mask_array_t), dimension(2), intent(in) :: qn_mask_in
      type(quantum_numbers_mask_t), dimension(:), intent(in) :: &
           qn_mask_conn_initial
      type(quantum_numbers_mask_t), intent(in) :: qn_mask_conn
      type(quantum_numbers_t), intent(in), optional :: qn_filter_conn
      type(quantum_numbers_mask_t), intent(in), optional :: qn_mask_rest
      type(index_map_t), dimension(2) :: prt_index_in
      type(index_map_t) :: prt_index_conn
      integer :: n_tot, n_conn
      integer, dimension(2) :: n_rest
      integer :: i, j, k, m
      type(quantum_numbers_t), dimension(:), allocatable :: qn
      type(quantum_numbers_mask_t), dimension(:), allocatable :: qn_mask
      type(connection_entry_t), pointer :: entry
      integer :: result_index
      n_conn = connection_table%n_conn
      n_rest = connection_table%n_rest
      n_tot = sum (n_rest) + n_conn
      allocate (qn (n_tot), qn_mask (n_tot))
      do i = 1, 2
         call index_map_init (prt_index_in(i), n_rest(i))
         prt_index_in(i) = &
              prt_map_in(i)%entry ([ (j, j = 1, n_rest(i)) ])
      end do
      call index_map_init (prt_index_conn, n_conn)
      prt_index_conn = prt_map_conn%entry ([ (j, j = 1, n_conn) ])
      do i = 1, 2
         if (present (qn_mask_rest)) then
            qn_mask(prt_index_in(i)%entry) = &
                 pack (qn_mask_in(i)%mask, prt_is_connected(i)%entry) &
                 .or. qn_mask_rest
         else
            qn_mask(prt_index_in(i)%entry) = &
                 pack (qn_mask_in(i)%mask, prt_is_connected(i)%entry)
         end if
      end do
      qn_mask(prt_index_conn%entry) = qn_mask_conn_initial .or. qn_mask_conn
      call interaction_init (eval%int, n_in, n_vir, n_out, mask=qn_mask)
      m = 1
      do i = 1, connection_table%n_me_conn
         entry => connection_table%entry(i)
         qn(prt_index_conn%entry) = &
              quantum_numbers_undefined (entry%qn_conn, qn_mask_conn)
         if (present (qn_filter_conn)) then
            if (.not. all (qn(prt_index_conn%entry) .match. qn_filter_conn)) &
                 cycle
         end if
         do j = 1, entry%n_index(1)
            qn(prt_index_in(1)%entry) = entry%qn_in_list(1)%qn(:,j)
            do k = 1, entry%n_index(2)
               qn(prt_index_in(2)%entry) = entry%qn_in_list(2)%qn(:,k)
               call interaction_add_state (int, qn, me_index = result_index)
               call index_map_set_entry &
                    (connection_table%index_result, m, result_index)
               m = m + 1
            end do
         end do
      end do
      call interaction_freeze (int)
    end subroutine make_product_interaction

    subroutine make_pairing_array (pa, n_matrix_elements, connection_table)
      type(pairing_array_t), dimension(:), intent(out), allocatable :: pa
      integer, intent(in) :: n_matrix_elements
      type(connection_table_t), intent(in), target :: connection_table
      type(connection_entry_t), pointer :: entry
      integer, dimension(:), allocatable :: n_entries
      integer :: i, j, k, m, r
      allocate (pa (n_matrix_elements))
      allocate (n_entries (n_matrix_elements))
      n_entries = 0
      do m = 1, size (connection_table%index_result)
         r = index_map_get_entry (connection_table%index_result, m)
         n_entries(r) = n_entries(r) + 1
      end do
      call pairing_array_init &
           (pa, n_entries, has_i2=.true., has_factor=.false.)
      m = 1
      n_entries = 0
      do i = 1, connection_table%n_me_conn
         entry => connection_table%entry(i)
         do j = 1, entry%n_index(1)
            do k = 1, entry%n_index(2)
               r = index_map_get_entry (connection_table%index_result, m)
               n_entries(r) = n_entries(r) + 1
               pa(r)%i1(n_entries(r)) = &
                    index_map_get_entry (entry%index_in(1), j)
               pa(r)%i2(n_entries(r)) = &
                    index_map_get_entry (entry%index_in(2), k)
               m = m + 1
            end do
         end do
      end do
    end subroutine make_pairing_array

    subroutine record_links (int, &
         int_in1, int_in2, connection_index, prt_map_in, prt_map_conn, &
         prt_is_connected, connections_are_resonant)
      type(interaction_t), intent(inout) :: int
      type(interaction_t), intent(in), target :: int_in1, int_in2
      integer, dimension(:,:), intent(in) :: connection_index
      type(index_map_t), dimension(2), intent(in) :: prt_map_in
      type(index_map_t), intent(in) :: prt_map_conn
      type(prt_mask_t), dimension(2), intent(in) :: prt_is_connected
      logical, intent(in), optional :: connections_are_resonant
      type(index_map_t), dimension(2) :: prt_map_all
      integer :: i, j, k, ival
      call index_map_init (prt_map_all(1), size (prt_is_connected(1)))
      k = 0
      j = 0
      do i = 1, size (prt_is_connected(1))
         if (prt_is_connected(1)%entry(i)) then
            j = j + 1
            ival = index_map_get_entry (prt_map_in(1), j)
            call index_map_set_entry (prt_map_all(1), i, ival)
         else
            k = k + 1
            ival = index_map_get_entry (prt_map_conn, k)
            call index_map_set_entry (prt_map_all(1), i, ival)
         end if
         call interaction_set_source_link (int, ival, int_in1, i)
      end do
      call interaction_transfer_relations (int_in1, int, prt_map_all(1)%entry)
      call index_map_init (prt_map_all(2), size (prt_is_connected(2)))
      j = 0
      do i = 1, size (prt_is_connected(2))
         if (prt_is_connected(2)%entry(i)) then
            j = j + 1
            ival = index_map_get_entry (prt_map_in(2), j)
            call index_map_set_entry (prt_map_all(2), i, ival)
            call interaction_set_source_link (int, ival, int_in2, i)
         else
            call index_map_set_entry (prt_map_all(2), i, 0)
         end if
      end do
      call interaction_transfer_relations (int_in2, int, prt_map_all(2)%entry)
      call interaction_relate_connections (int, &
           int_in2, connection_index(:,2), prt_map_all(2)%entry, &
           prt_map_conn%entry, connections_are_resonant)
    end subroutine record_links

  end subroutine evaluator_init_product_ii

  subroutine evaluator_init_product_ie &
       (eval, int_in1, eval_in2, qn_mask_conn, qn_filter_conn, qn_mask_rest, &
        connections_are_resonant)
    type(evaluator_t), intent(out), target :: eval
    type(interaction_t), intent(in), target :: int_in1
    type(evaluator_t), intent(in), target :: eval_in2
    type(quantum_numbers_mask_t), intent(in) :: qn_mask_conn
    type(quantum_numbers_t), intent(in), optional :: qn_filter_conn
    type(quantum_numbers_mask_t), intent(in), optional :: qn_mask_rest
    logical, intent(in), optional :: connections_are_resonant
    call evaluator_init_product_ii &
         (eval, int_in1, eval_in2%int, &
          qn_mask_conn, qn_filter_conn, qn_mask_rest, &
          connections_are_resonant)
  end subroutine evaluator_init_product_ie

  subroutine evaluator_init_product_ei &
       (eval, eval_in1, int_in2, qn_mask_conn, qn_filter_conn, qn_mask_rest, &
        connections_are_resonant)
    type(evaluator_t), intent(out), target :: eval
    type(evaluator_t), intent(in), target :: eval_in1
    type(interaction_t), intent(in), target :: int_in2
    type(quantum_numbers_mask_t), intent(in) :: qn_mask_conn
    type(quantum_numbers_t), intent(in), optional :: qn_filter_conn
    type(quantum_numbers_mask_t), intent(in), optional :: qn_mask_rest
    logical, intent(in), optional :: connections_are_resonant
    call evaluator_init_product_ii &
         (eval, eval_in1%int, int_in2, &
          qn_mask_conn, qn_filter_conn, qn_mask_rest, &
         connections_are_resonant)
  end subroutine evaluator_init_product_ei

  subroutine evaluator_init_product_ee &
       (eval, eval_in1, eval_in2, qn_mask_conn, qn_filter_conn, qn_mask_rest, &
        connections_are_resonant)
    type(evaluator_t), intent(out), target :: eval
    type(evaluator_t), intent(in), target :: eval_in1, eval_in2
    type(quantum_numbers_mask_t), intent(in) :: qn_mask_conn
    type(quantum_numbers_t), intent(in), optional :: qn_filter_conn
    type(quantum_numbers_mask_t), intent(in), optional :: qn_mask_rest
    logical, intent(in), optional :: connections_are_resonant
    call evaluator_init_product_ii &
         (eval, eval_in1%int, eval_in2%int, &
          qn_mask_conn, qn_filter_conn, qn_mask_rest, &
          connections_are_resonant)
  end subroutine evaluator_init_product_ee

\eatline
\nwendcode{}\nwbegindocs{1057}\nwdocspar
\subsection{Creating an evaluator: square}
The generic initializer for an evaluator that squares a matrix element.
Depending on the provided mask, we select the appropriate specific initializer
for either diagonal or non-diagonal helicity density matrices.
\nwenddocs{}\nwbegincode{1058}\moddef{Evaluators: public}\plusendmoddef
  public :: evaluator_init_square
\nwendcode{}\nwbegincode{1059}\moddef{Evaluators: procedures}\plusendmoddef
  subroutine evaluator_init_square (eval, int_in, qn_mask, &
       col_flow_index, col_factor, col_index_hi, expand_color_flows, nc)
    type(evaluator_t), intent(out), target :: eval
    type(interaction_t), intent(in), target :: int_in
    type(quantum_numbers_mask_t), dimension(:), intent(in) :: qn_mask
    integer, dimension(:,:), intent(in), optional :: col_flow_index
    complex(default), dimension(:), intent(in), optional :: col_factor
    integer, dimension(:), intent(in), optional :: col_index_hi
    logical, intent(in), optional :: expand_color_flows
    integer, intent(in), optional :: nc
    if (all (qn_mask%diagonal_helicity ())) then
       call evaluator_init_square_diag (eval, int_in, qn_mask, &
            col_flow_index, col_factor, col_index_hi, expand_color_flows, nc)
    else
       call evaluator_init_square_nondiag (eval, int_in, qn_mask, &
            col_flow_index, col_factor, col_index_hi, expand_color_flows, nc)
    end if
  end subroutine evaluator_init_square

\eatline
\nwendcode{}\nwbegindocs{1060}\nwdocspar
\subsubsection{Color-summed squared matrix (diagonal helicities)}
The initializer for an evaluator that squares a matrix element,
including color factors.  The mask must be such that off-diagonal matrix
elements are excluded.

If {\tt{}color{\char95}flows} is set, the evaluator keeps color-flow entries
separate and drops all interfering color structures.  The color factors are
set to unity in this case.

There is only one input interaction.  The quantum-number mask is an
array, one entry for each particle, so they can be treated
individually.  For academic purposes, we allow for the number of
colors being different from three (but 3 is the default).

The algorithm is analogous to multiplication, with a few notable
differences:
\begin{enumerate}
\item 
  The connected particles are known, the correspondence is
  one-to-one.  All particles are connected, and the mapping of indices
  is trivial, which simplifies the following steps.
\item
  {\tt{}accumulate{\char95}connected{\char95}states}: The matrix of connected states
  encompasses all particles, but color indices are removed.  However,
  ghost states are still kept separate from physical color states.  No
  color-index reassignment is necessary.
\item
  The table of connections contains single index and quantum-number
  arrays instead of pairs of them.  They are paired with themselves
  in all possible ways.
\item
  {\tt{}make{\char95}squared{\char95}interaction}: Now apply the predefined
  quantum-numbers mask, which usually collects all color states
  (physical and ghosts), and possibly a helicity sum.
\item
  {\tt{}make{\char95}pairing{\char95}array}: For each pair of input states, compute the
  color factor (including a potential ghost-parity sign) and store
  this in the pairing array together with the matrix-element indices
  for multiplication.
\item
  {\tt{}record{\char95}links}: This is again trivial due to the one-to-one
  correspondence.
\end{enumerate}

\nwenddocs{}\nwbegincode{1061}\moddef{Evaluators: public}\plusendmoddef
  public :: evaluator_init_square_diag
\nwendcode{}\nwbegincode{1062}\moddef{Evaluators: procedures}\plusendmoddef
  subroutine evaluator_init_square_diag (eval, int_in, qn_mask, & 
       col_flow_index, col_factor, col_index_hi, expand_color_flows, nc)

    type(evaluator_t), intent(out), target :: eval
    type(interaction_t), intent(in), target :: int_in
    type(quantum_numbers_mask_t), dimension(:), intent(in) :: qn_mask
    integer, dimension(:,:), intent(in), optional :: col_flow_index
    complex(default), dimension(:), intent(in), optional :: col_factor
    integer, dimension(:), intent(in), optional :: col_index_hi
    logical, intent(in), optional :: expand_color_flows
    integer, intent(in), optional :: nc

    integer :: n_in, n_vir, n_out, n_tot
    type(quantum_numbers_mask_t), dimension(:), allocatable :: qn_mask_initial
    type(state_matrix_t), pointer :: state_in

    type :: connection_table_t
      integer :: n_tot = 0
      integer :: n_me_conn = 0
      type(state_matrix_t) :: state
      type(index_map_t) :: index_conn
      type(connection_entry_t), dimension(:), allocatable :: entry
      type(index_map_t) :: index_result
    end type connection_table_t
    type(connection_table_t) :: connection_table

    logical :: sum_colors
    type(color_table_t) :: color_table

    if (present (expand_color_flows)) then
       sum_colors = .not. expand_color_flows
    else
       sum_colors = .true.
    end if

    if (sum_colors) then
       eval%type = EVAL_SQUARE_WITH_COLOR_FACTORS
    else
       eval%type = EVAL_SQUARED_FLOWS
    end if
    eval%int_in1 => int_in

    ! print *, "Interaction square with color factors (diag)"  !!! Debugging
    ! print *, "Input interaction"                             !!! Debugging
    ! call interaction_write (int_in)                          !!! Debugging
    
    n_in  = interaction_get_n_in  (int_in)
    n_vir = interaction_get_n_vir (int_in)
    n_out = interaction_get_n_out (int_in)
    n_tot = interaction_get_n_tot (int_in)

    state_in => interaction_get_state_matrix_ptr (int_in)

    allocate (qn_mask_initial (n_tot))
    qn_mask_initial = interaction_get_mask (int_in)
    call qn_mask_initial%set_color (sum_colors, mask_cg=.false.)
    if (sum_colors) then
       call color_table_init (color_table, state_in, n_tot)
       if (present (col_flow_index) .and. present (col_factor) &
           .and. present (col_index_hi)) then
          call color_table_set_color_factors &
               (color_table, col_flow_index, col_factor, col_index_hi)
       end if
       ! call color_table_write (color_table)     !!! Debugging
    end if

    call connection_table_init (connection_table, state_in, &
         qn_mask_initial, qn_mask, n_tot)
    call connection_table_fill (connection_table, state_in)
    call make_squared_interaction (eval%int, &
         n_in, n_vir, n_out, n_tot, &
         connection_table, sum_colors, qn_mask_initial .or. qn_mask)
    call make_pairing_array (eval%pairing_array, &
         interaction_get_n_matrix_elements (eval%int), &
         connection_table, sum_colors, color_table, n_in, n_tot, nc)
    call record_links (eval%int, int_in, n_tot)
    call connection_table_final (connection_table)
    ! print *, "Result evaluator:"     !!! Debugging
    ! call evaluator_write (eval)      !!! Debugging

  contains
    
    subroutine connection_table_init &
         (connection_table, state_in, qn_mask_in, qn_mask, n_tot)
      type(connection_table_t), intent(out) :: connection_table
      type(state_matrix_t), intent(in), target :: state_in
      type(quantum_numbers_mask_t), dimension(:), intent(in) :: qn_mask_in
      type(quantum_numbers_mask_t), dimension(:), intent(in) :: qn_mask
      integer, intent(in) :: n_tot
      type(quantum_numbers_t), dimension(n_tot) :: qn
      type(state_iterator_t) :: it
      integer :: i, n_me_in, me_index_in
      integer :: me_index_conn, n_me_conn
      integer, dimension(1) :: me_count
      logical :: qn_passed
      connection_table%n_tot = n_tot
      n_me_in = state_in%get_n_matrix_elements ()
      call index_map_init (connection_table%index_conn, n_me_in)
      connection_table%index_conn = 0
      call connection_table%state%init (n_counters=1)
      call it%init (state_in)
      do while (it%is_valid ())
         qn = it%get_quantum_numbers ()
         if (all (quantum_numbers_are_physical (qn, qn_mask))) then
            call qn%undefine (qn_mask_in)
            qn_passed = .true.
            if (qn_passed) then
               me_index_in = it%get_me_index ()
               call connection_table%state%add_state (qn, &
                    counter_index = 1, me_index = me_index_conn)
               call index_map_set_entry (connection_table%index_conn, &
                    me_index_in, me_index_conn)
            end if
         end if
         call it%advance ()
      end do
      n_me_conn = connection_table%state%get_n_matrix_elements ()
      connection_table%n_me_conn = n_me_conn
      allocate (connection_table%entry (n_me_conn))
      call it%init (connection_table%state)
      do while (it%is_valid ())
         i = it%get_me_index ()
         me_count = it%get_me_count ()
         call connection_entry_init (connection_table%entry(i), 1, 2, &
              it%get_quantum_numbers (), me_count, [n_tot])
         call it%advance ()
      end do
    end subroutine connection_table_init
         
    subroutine connection_table_final (connection_table)
      type(connection_table_t), intent(inout) :: connection_table
      call connection_table%state%final ()
    end subroutine connection_table_final

    subroutine connection_table_write (connection_table, unit)
      type(connection_table_t), intent(in) :: connection_table
      integer, intent(in), optional :: unit
      integer :: i
      integer :: u
      u = given_output_unit (unit)
      write (u, *) "Connection table:"
      call connection_table%state%write (unit)
      if (index_map_exists (connection_table%index_conn)) then
         write (u, *) "  Index mapping input => connection table:"
         do i = 1, size (connection_table%index_conn)
            write (u, *)  i, &
                   index_map_get_entry (connection_table%index_conn, i)
         end do
      end if
      if (allocated (connection_table%entry)) then
         write (u, *) "  Connection table contents"
         do i = 1, size (connection_table%entry)
            call connection_entry_write (connection_table%entry(i), unit)
         end do
      end if
      if (index_map_exists (connection_table%index_result)) then
         write (u, *) "  Index mapping connection table => output"
         do i = 1, size (connection_table%index_result)
            write (u, *)  i, &
                  index_map_get_entry (connection_table%index_result, i)
         end do
      end if
    end subroutine connection_table_write

    subroutine connection_table_fill (connection_table, state)
      type(connection_table_t), intent(inout) :: connection_table
      type(state_matrix_t), intent(in), target :: state
      integer :: index_in, index_conn, n_result_entries
      type(state_iterator_t) :: it
      integer :: k
      call it%init (state)
      do while (it%is_valid ())
         index_in = it%get_me_index ()
         index_conn = &
              index_map_get_entry (connection_table%index_conn, index_in)
         if (index_conn /= 0) then
            call connection_entry_add_state &
                 (connection_table%entry(index_conn), &
                  index_in, it%get_quantum_numbers ())
         end if
         call it%advance ()
      end do
      n_result_entries = 0
      do k = 1, size (connection_table%entry)
         n_result_entries = &
              n_result_entries + connection_table%entry(k)%n_index(1) ** 2
      end do
      call index_map_init (connection_table%index_result, n_result_entries)
      connection_table%index_result = 0
    end subroutine connection_table_fill

    subroutine connection_entry_add_state (entry, index_in, qn_in)
      type(connection_entry_t), intent(inout) :: entry
      integer, intent(in) :: index_in
      type(quantum_numbers_t), dimension(:), intent(in) :: qn_in
      integer :: c
      entry%count = entry%count + 1
      c = entry%count(1)
      call index_map_set_entry (entry%index_in(1), c, index_in)
      entry%qn_in_list(1)%qn(:,c) = qn_in
    end subroutine connection_entry_add_state

    subroutine make_squared_interaction (int, &
         n_in, n_vir, n_out, n_tot, &
         connection_table, sum_colors, qn_mask)
      type(interaction_t), intent(out), target :: int
      integer, intent(in) :: n_in, n_vir, n_out, n_tot
      type(connection_table_t), intent(inout), target :: connection_table
      logical, intent(in) :: sum_colors
      type(quantum_numbers_mask_t), dimension(:), intent(in) :: qn_mask
      type(connection_entry_t), pointer :: entry
      integer :: result_index, n_contrib
      integer :: i, m
      type(quantum_numbers_t), dimension(n_tot) :: qn
      call interaction_init (eval%int, n_in, n_vir, n_out, mask=qn_mask)
      m = 0
      do i = 1, connection_table%n_me_conn
         entry => connection_table%entry(i)
         qn = quantum_numbers_undefined (entry%qn_conn, qn_mask)
         if (.not. sum_colors)   call qn(1:n_in)%invert_color ()
         call interaction_add_state (int, qn, me_index = result_index)
         n_contrib = entry%n_index(1) ** 2
         connection_table%index_result%entry(m+1:m+n_contrib) = result_index
         m = m + n_contrib
      end do
      call interaction_freeze (int)
    end subroutine make_squared_interaction

    subroutine make_pairing_array (pa, &
         n_matrix_elements, connection_table, sum_colors, color_table, &
         n_in, n_tot, nc)
      type(pairing_array_t), dimension(:), intent(out), allocatable :: pa
      integer, intent(in) :: n_matrix_elements
      type(connection_table_t), intent(in), target :: connection_table
      logical, intent(in) :: sum_colors
      type(color_table_t), intent(inout) :: color_table
      type(connection_entry_t), pointer :: entry
      integer, intent(in) :: n_in, n_tot
      integer, intent(in), optional :: nc
      integer, dimension(:), allocatable :: n_entries
      integer :: i, k, l, ks, ls, m, r
      integer :: color_multiplicity_in
      allocate (pa (n_matrix_elements))
      allocate (n_entries (n_matrix_elements))
      n_entries = 0
      do m = 1, size (connection_table%index_result)
         r = index_map_get_entry (connection_table%index_result, m)
         n_entries(r) = n_entries(r) + 1
      end do
      call pairing_array_init &
           (pa, n_entries, has_i2 = sum_colors, has_factor = sum_colors)
      m = 1
      n_entries = 0
      do i = 1, connection_table%n_me_conn
         entry => connection_table%entry(i)
         do k = 1, entry%n_index(1)
            if (sum_colors) then
               color_multiplicity_in = &
                     product (abs (quantum_numbers_get_color_type &
                                     (entry%qn_in_list(1)%qn(:n_in, k))))
               do l = 1, entry%n_index(1)
                  r = index_map_get_entry (connection_table%index_result, m)
                  n_entries(r) = n_entries(r) + 1
                  ks = index_map_get_entry (entry%index_in(1), k)
                  ls = index_map_get_entry (entry%index_in(1), l)
                  pa(r)%i1(n_entries(r)) = ks
                  pa(r)%i2(n_entries(r)) = ls
                  pa(r)%factor(n_entries(r)) = &
                       color_table_get_color_factor (color_table, ks, ls, nc) &
                       / color_multiplicity_in
                  m = m + 1
               end do
            else
               r = index_map_get_entry (connection_table%index_result, m)
               n_entries(r) = n_entries(r) + 1
               ks = index_map_get_entry (entry%index_in(1), k)
               pa(r)%i1(n_entries(r)) = ks
               m = m + 1  
            end if               
         end do
      end do
    end subroutine make_pairing_array

    subroutine record_links (int, int_in, n_tot)
      type(interaction_t), intent(inout) :: int
      type(interaction_t), intent(in), target :: int_in
      integer, intent(in) :: n_tot
      integer, dimension(n_tot) :: map
      integer :: i
      do i = 1, n_tot
         call interaction_set_source_link (int, i, int_in, i)
      end do
      map = [ (i, i = 1, n_tot) ]
      call interaction_transfer_relations (int_in, int, map)
    end subroutine record_links

  end subroutine evaluator_init_square_diag

\eatline
\nwendcode{}\nwbegindocs{1063}\nwdocspar
\subsubsection{Color-summed squared matrix (support nodiagonal helicities)}
The initializer for an evaluator that squares a matrix element,
including color factors.  Unless requested otherwise by the
quantum-number mask, the result contains off-diagonal matrix elements.
(The input interaction must be diagonal since it represents an
amplitude, not a density matrix.)

There is only one input interaction.  The quantum-number mask is an
array, one entry for each particle, so they can be treated
individually.  For academic purposes, we allow for the number of
colors being different from three (but 3 is the default).

The algorithm is analogous to the previous one, with some additional
complications due to the necessity to loop over two helicity indices.
\nwenddocs{}\nwbegincode{1064}\moddef{Evaluators: public}\plusendmoddef
  public :: evaluator_init_square_nondiag
\nwendcode{}\nwbegincode{1065}\moddef{Evaluators: procedures}\plusendmoddef
  subroutine evaluator_init_square_nondiag (eval, int_in, qn_mask, & 
       col_flow_index, col_factor, col_index_hi, expand_color_flows, nc)

    type(evaluator_t), intent(out), target :: eval
    type(interaction_t), intent(in), target :: int_in
    type(quantum_numbers_mask_t), dimension(:), intent(in) :: qn_mask
    integer, dimension(:,:), intent(in), optional :: col_flow_index
    complex(default), dimension(:), intent(in), optional :: col_factor
    integer, dimension(:), intent(in), optional :: col_index_hi
    logical, intent(in), optional :: expand_color_flows
    integer, intent(in), optional :: nc

    integer :: n_in, n_vir, n_out, n_tot
    type(quantum_numbers_mask_t), dimension(:), allocatable :: qn_mask_initial
    type(state_matrix_t), pointer :: state_in

    type :: connection_table_t
      integer :: n_tot = 0
      integer :: n_me_conn = 0
      type(state_matrix_t) :: state
      type(index_map2_t) :: index_conn
      type(connection_entry_t), dimension(:), allocatable :: entry
      type(index_map_t) :: index_result
    end type connection_table_t
    type(connection_table_t) :: connection_table
       
    logical :: sum_colors
    type(color_table_t) :: color_table

    if (present (expand_color_flows)) then
       sum_colors = .not. expand_color_flows
    else
       sum_colors = .true.
    end if

    if (sum_colors) then
       eval%type = EVAL_SQUARE_WITH_COLOR_FACTORS
    else
       eval%type = EVAL_SQUARED_FLOWS
    end if
    eval%int_in1 => int_in

    ! print *, "Interaction square with color factors (nondiag)"  !!! Debugging
    ! print *, "Input interaction"                                !!! Debugging
    ! call interaction_write (int_in)                             !!! Debugging
    n_in  = interaction_get_n_in  (int_in)
    n_vir = interaction_get_n_vir (int_in)
    n_out = interaction_get_n_out (int_in)
    n_tot = interaction_get_n_tot (int_in)

    state_in => interaction_get_state_matrix_ptr (int_in)

    allocate (qn_mask_initial (n_tot))
    qn_mask_initial = interaction_get_mask (int_in)
    call qn_mask_initial%set_color (sum_colors, mask_cg=.false.)
    if (sum_colors) then
       call color_table_init (color_table, state_in, n_tot)
       if (present (col_flow_index) .and. present (col_factor) &
           .and. present (col_index_hi)) then
          call color_table_set_color_factors &
               (color_table, col_flow_index, col_factor, col_index_hi)
       end if
       ! call color_table_write (color_table)    !!! Debugging
    end if

    call connection_table_init (connection_table, state_in, &
         qn_mask_initial, qn_mask, n_tot)
    call connection_table_fill (connection_table, state_in)
    call make_squared_interaction (eval%int, &
         n_in, n_vir, n_out, n_tot, &
         connection_table, sum_colors, qn_mask_initial .or. qn_mask)
    ! call connection_table_write (connection_table)     !!! Debugging
    call make_pairing_array (eval%pairing_array, &
         interaction_get_n_matrix_elements (eval%int), &
         connection_table, sum_colors, color_table, n_in, n_tot, nc)
    call record_links (eval%int, int_in, n_tot)
    call connection_table_final (connection_table)

    ! print *, "Result evaluator:"     !!! Debugging
    ! call evaluator_write (eval)      !!! Debugging

  contains
    
    subroutine connection_table_init &
         (connection_table, state_in, qn_mask_in, qn_mask, n_tot)
      type(connection_table_t), intent(out) :: connection_table
      type(state_matrix_t), intent(in), target :: state_in
      type(quantum_numbers_mask_t), dimension(:), intent(in) :: qn_mask_in
      type(quantum_numbers_mask_t), dimension(:), intent(in) :: qn_mask
      integer, intent(in) :: n_tot
      type(quantum_numbers_t), dimension(n_tot) :: qn1, qn2, qn
      type(state_iterator_t) :: it1, it2, it
      integer :: i, n_me_in, me_index_in1, me_index_in2
      integer :: me_index_conn, n_me_conn
      integer, dimension(1) :: me_count
      logical :: qn_passed
      connection_table%n_tot = n_tot
      n_me_in = state_in%get_n_matrix_elements ()
      call index_map2_init (connection_table%index_conn, n_me_in)
      connection_table%index_conn = 0
      call connection_table%state%init (n_counters=1)
      call it1%init (state_in)
      do while (it1%is_valid ())
         qn1 = it1%get_quantum_numbers ()
         me_index_in1 = it1%get_me_index ()
         call it2%init (state_in)
         do while (it2%is_valid ())
            qn2 = it2%get_quantum_numbers ()
            if (all (quantum_numbers_are_compatible (qn1, qn2, qn_mask))) then
               qn = qn1 .merge. qn2
               call qn%undefine (qn_mask_in)
               qn_passed = .true.
               if (qn_passed) then
                  me_index_in2 = it2%get_me_index ()
                  call connection_table%state%add_state (qn, &
                       counter_index = 1, me_index = me_index_conn)
                  call index_map2_set_entry (connection_table%index_conn, &
                       me_index_in1, me_index_in2, me_index_conn)
               end if
            end if
            call it2%advance ()
         end do
         call it1%advance ()
      end do
      n_me_conn = connection_table%state%get_n_matrix_elements ()
      connection_table%n_me_conn = n_me_conn
      allocate (connection_table%entry (n_me_conn))
      call it%init (connection_table%state)
      do while (it%is_valid ())
         i = it%get_me_index ()
         me_count = it%get_me_count ()
         call connection_entry_init (connection_table%entry(i), 1, 2, &
              it%get_quantum_numbers (), me_count, [n_tot])
         call it%advance ()
      end do
    end subroutine connection_table_init
         
    subroutine connection_table_final (connection_table)
      type(connection_table_t), intent(inout) :: connection_table
      call connection_table%state%final ()
    end subroutine connection_table_final

    subroutine connection_table_write (connection_table, unit)
      type(connection_table_t), intent(in) :: connection_table
      integer, intent(in), optional :: unit
      integer :: i, j
      integer :: u
      u = given_output_unit (unit)
      write (u, *) "Connection table:"
      call connection_table%state%write (unit)
      if (index_map2_exists (connection_table%index_conn)) then
         write (u, *) "  Index mapping input => connection table:"
         do i = 1, size (connection_table%index_conn)
            do j = 1, size (connection_table%index_conn)
               write (u, *)  i, j, &
                    index_map2_get_entry (connection_table%index_conn, i, j)
            end do
         end do
      end if
      if (allocated (connection_table%entry)) then
         write (u, *) "  Connection table contents"
         do i = 1, size (connection_table%entry)
            call connection_entry_write (connection_table%entry(i), unit)
         end do
      end if
      if (index_map_exists (connection_table%index_result)) then
         write (u, *) "  Index mapping connection table => output"
         do i = 1, size (connection_table%index_result)
            write (u, *)  i, &
                 index_map_get_entry (connection_table%index_result, i)
         end do
      end if
    end subroutine connection_table_write

    subroutine connection_table_fill (connection_table, state)
      type(connection_table_t), intent(inout), target :: connection_table
      type(state_matrix_t), intent(in), target :: state
      integer :: index1_in, index2_in, index_conn, n_result_entries
      type(state_iterator_t) :: it1, it2
      integer :: k
      call it1%init (state)
      do while (it1%is_valid ())
         index1_in = it1%get_me_index ()
         call it2%init (state)
         do while (it2%is_valid ())
            index2_in = it2%get_me_index ()
            index_conn = index_map2_get_entry &
                            (connection_table%index_conn, index1_in, index2_in)
            if (index_conn /= 0) then
               call connection_entry_add_state &
                    (connection_table%entry(index_conn), &
                     index1_in, index2_in, &
                     it1%get_quantum_numbers () &
                     .merge. &
                     it2%get_quantum_numbers ())
            end if
            call it2%advance ()
         end do
         call it1%advance ()
      end do
      n_result_entries = 0
      do k = 1, size (connection_table%entry)
         n_result_entries = &
              n_result_entries + connection_table%entry(k)%n_index(1)
      end do
      call index_map_init (connection_table%index_result, n_result_entries)
      connection_table%index_result = 0
    end subroutine connection_table_fill

    subroutine connection_entry_add_state (entry, index1_in, index2_in, qn_in)
      type(connection_entry_t), intent(inout) :: entry
      integer, intent(in) :: index1_in, index2_in
      type(quantum_numbers_t), dimension(:), intent(in) :: qn_in
      integer :: c
      entry%count = entry%count + 1
      c = entry%count(1)
      call index_map_set_entry (entry%index_in(1), c, index1_in)
      call index_map_set_entry (entry%index_in(2), c, index2_in)
      entry%qn_in_list(1)%qn(:,c) = qn_in
    end subroutine connection_entry_add_state

    subroutine make_squared_interaction (int, &
         n_in, n_vir, n_out, n_tot, &
         connection_table, sum_colors, qn_mask)
      type(interaction_t), intent(out), target :: int
      integer, intent(in) :: n_in, n_vir, n_out, n_tot
      type(connection_table_t), intent(inout), target :: connection_table
      logical, intent(in) :: sum_colors
      type(quantum_numbers_mask_t), dimension(:), intent(in) :: qn_mask
      type(connection_entry_t), pointer :: entry
      integer :: result_index
      integer :: i, k, m
      type(quantum_numbers_t), dimension(n_tot) :: qn
      call interaction_init (eval%int, n_in, n_vir, n_out, mask=qn_mask)
      m = 0
      do i = 1, connection_table%n_me_conn
         entry => connection_table%entry(i)
         do k = 1, size (entry%qn_in_list(1)%qn, 2)
            qn = quantum_numbers_undefined &
                    (entry%qn_in_list(1)%qn(:,k), qn_mask)
            if (.not. sum_colors)  call qn(1:n_in)%invert_color ()
            call interaction_add_state (int, qn, me_index = result_index)
            call index_map_set_entry (connection_table%index_result, m + 1, &
                 result_index)
            m = m + 1
         end do
      end do
      call interaction_freeze (int)
    end subroutine make_squared_interaction

    subroutine make_pairing_array (pa, &
         n_matrix_elements, connection_table, sum_colors, color_table, &
         n_in, n_tot, nc)
      type(pairing_array_t), dimension(:), intent(out), allocatable :: pa
      integer, intent(in) :: n_matrix_elements
      type(connection_table_t), intent(in), target :: connection_table
      logical, intent(in) :: sum_colors
      type(color_table_t), intent(inout) :: color_table
      type(connection_entry_t), pointer :: entry
      integer, intent(in) :: n_in, n_tot
      integer, intent(in), optional :: nc
      integer, dimension(:), allocatable :: n_entries
      integer :: i, k, k1s, k2s, m, r
      integer :: color_multiplicity_in
      allocate (pa (n_matrix_elements))
      allocate (n_entries (n_matrix_elements))
      n_entries = 0
      do m = 1, size (connection_table%index_result)
         r = index_map_get_entry (connection_table%index_result, m)
         n_entries(r) = n_entries(r) + 1
      end do
      call pairing_array_init &
           (pa, n_entries, has_i2 = sum_colors, has_factor = sum_colors)
      m = 1
      n_entries = 0
      do i = 1, connection_table%n_me_conn
         entry => connection_table%entry(i)
         do k = 1, entry%n_index(1)
            r = index_map_get_entry (connection_table%index_result, m)
            n_entries(r) = n_entries(r) + 1
            if (sum_colors) then
               k1s = index_map_get_entry (entry%index_in(1), k)
               k2s = index_map_get_entry (entry%index_in(2), k)
               pa(r)%i1(n_entries(r)) = k1s
               pa(r)%i2(n_entries(r)) = k2s
               color_multiplicity_in = &
                     product (abs (quantum_numbers_get_color_type &
                                     (entry%qn_in_list(1)%qn(:n_in, k))))
               pa(r)%factor(n_entries(r)) = &
                    color_table_get_color_factor (color_table, k1s, k2s, nc) &
                    / color_multiplicity_in
            else
               k1s = index_map_get_entry (entry%index_in(1), k)
               pa(r)%i1(n_entries(r)) = k1s
            end if               
            m = m + 1
         end do
      end do
    end subroutine make_pairing_array

    subroutine record_links (int, int_in, n_tot)
      type(interaction_t), intent(inout) :: int
      type(interaction_t), intent(in), target :: int_in
      integer, intent(in) :: n_tot
      integer, dimension(n_tot) :: map
      integer :: i
      do i = 1, n_tot
         call interaction_set_source_link (int, i, int_in, i)
      end do
      map = [ (i, i = 1, n_tot) ]
      call interaction_transfer_relations (int_in, int, map)
    end subroutine record_links

  end subroutine evaluator_init_square_nondiag

\eatline
\nwendcode{}\nwbegindocs{1066}\nwdocspar
\subsubsection{Copy with additional contracted color states}
This evaluator involves no square or multiplication, its matrix
elements are just copies of the (single) input interaction.  However,
the state matrix of the interaction contains additional states that
have color indices contracted.  This is used for copies of the beam or
structure-function interactions that need to match the hard
interaction also in the case where its color indices coincide.
\nwenddocs{}\nwbegincode{1067}\moddef{Evaluators: public}\plusendmoddef
  public :: evaluator_init_color_contractions
\nwendcode{}\nwbegincode{1068}\moddef{Evaluators: procedures}\plusendmoddef
  subroutine evaluator_init_color_contractions (eval, int_in)
    type(evaluator_t), intent(out), target :: eval
    type(interaction_t), intent(in), target :: int_in
    integer :: n_in, n_vir, n_out, n_tot
    type(state_matrix_t) :: state_with_contractions
    integer, dimension(:), allocatable :: me_index
    integer, dimension(:), allocatable :: result_index
    eval%type = EVAL_COLOR_CONTRACTION
    eval%int_in1 => int_in
    ! print *, "Interaction with additional color contractions"  !!! Debugging
    ! print *, "Input interaction"                               !!! Debugging
    ! call interaction_write (int_in)                            !!! Debugging
    n_in  = interaction_get_n_in  (int_in)
    n_vir = interaction_get_n_vir (int_in)
    n_out = interaction_get_n_out (int_in)
    n_tot = interaction_get_n_tot (int_in)
    state_with_contractions = interaction_get_state_matrix_ptr (int_in)
    call state_with_contractions%add_color_contractions ()
    call make_contracted_interaction (eval%int, &
         me_index, result_index, &
         n_in, n_vir, n_out, n_tot, &
         state_with_contractions, interaction_get_mask (int_in))
    call make_pairing_array (eval%pairing_array, me_index, result_index)
    call record_links (eval%int, int_in, n_tot)
    call state_with_contractions%final ()
    ! print *, "Result evaluator:"     !!! Debugging
    ! call evaluator_write (eval)      !!! Debugging

  contains

    subroutine make_contracted_interaction (int, &
         me_index, result_index, &
         n_in, n_vir, n_out, n_tot, state, qn_mask)
      type(interaction_t), intent(out), target :: int
      integer, dimension(:), intent(out), allocatable :: me_index
      integer, dimension(:), intent(out), allocatable :: result_index
      integer, intent(in) :: n_in, n_vir, n_out, n_tot
      type(state_matrix_t), intent(in) :: state
      type(quantum_numbers_mask_t), dimension(:), intent(in) :: qn_mask
      type(state_iterator_t) :: it
      integer :: n_me, i
      type(quantum_numbers_t), dimension(n_tot) :: qn
      call interaction_init (int, n_in, n_vir, n_out, mask=qn_mask)
      n_me = state%get_n_leaves ()
      allocate (me_index (n_me))
      allocate (result_index (n_me))
      call it%init (state)
      i = 0
      do while (it%is_valid ())
         i = i + 1
         me_index(i) = it%get_me_index ()
         qn = it%get_quantum_numbers ()
         call interaction_add_state (int, qn, me_index = result_index(i))
         call it%advance ()
      end do
      call interaction_freeze (int)
    end subroutine make_contracted_interaction

    subroutine make_pairing_array (pa, me_index, result_index)
      type(pairing_array_t), dimension(:), intent(out), allocatable :: pa
      integer, dimension(:), intent(in) :: me_index, result_index
      integer, dimension(:), allocatable :: n_entries
      integer :: n_matrix_elements, r, i
      n_matrix_elements = size (me_index)
      allocate (pa (n_matrix_elements))
      allocate (n_entries (n_matrix_elements))
      n_entries = 1
      call pairing_array_init &
           (pa, n_entries, has_i2=.false., has_factor=.false.)
      do i = 1, n_matrix_elements
         r = result_index(i)
         pa(r)%i1(1) = me_index(i)
      end do
    end subroutine make_pairing_array

    subroutine record_links (int, int_in, n_tot)
      type(interaction_t), intent(inout) :: int
      type(interaction_t), intent(in), target :: int_in
      integer, intent(in) :: n_tot
      integer, dimension(n_tot) :: map
      integer :: i
      do i = 1, n_tot
         call interaction_set_source_link (int, i, int_in, i)
      end do
      map = [ (i, i = 1, n_tot) ]
      call interaction_transfer_relations (int_in, int, map)
    end subroutine record_links

  end subroutine evaluator_init_color_contractions

\eatline
\nwendcode{}\nwbegindocs{1069}\nwdocspar
\subsubsection{Auxiliary procedure for initialization}
This will become a standard procedure in F2008.  The result is true if
the number of true values in the mask is odd.  We use the function for
determining the ghost parity of a quantum-number array.
\nwenddocs{}\nwbegincode{1070}\moddef{Evaluators: procedures}\plusendmoddef
  function parity (mask)
    logical :: parity
    logical, dimension(:) :: mask
    integer :: i
    parity = .false.
    do i = 1, size (mask)
       if (mask(i))  parity = .not. parity
    end do
  end function parity

\eatline
\nwendcode{}\nwbegindocs{1071}\nwdocspar
\subsection{Accessing contents}
Return the interaction component, as a pointer to avoid any copying.
\nwenddocs{}\nwbegincode{1072}\moddef{Evaluators: public}\plusendmoddef
  public :: evaluator_get_int_ptr
\nwendcode{}\nwbegincode{1073}\moddef{Evaluators: procedures}\plusendmoddef
  function evaluator_get_int_ptr (eval) result (int)
    type(interaction_t), pointer :: int
    type(evaluator_t), intent(in), target :: eval
    int => eval%int
  end function evaluator_get_int_ptr

\eatline
\nwendcode{}\nwbegindocs{1074}\nwdocspar
\subsection{Inherited procedures}
Return true if the state matrix within the interaction is empty.
\nwenddocs{}\nwbegincode{1075}\moddef{Evaluators: public}\plusendmoddef
  public :: evaluator_is_empty
\nwendcode{}\nwbegincode{1076}\moddef{Evaluators: procedures}\plusendmoddef
  function evaluator_is_empty (eval) result (flag)
    logical :: flag
    type(evaluator_t), intent(in) :: eval
    flag = interaction_is_empty (eval%int)
  end function evaluator_is_empty

\eatline
\nwendcode{}\nwbegindocs{1077}Return the integer tag.
\nwenddocs{}\nwbegincode{1078}\moddef{Evaluators: public}\plusendmoddef
  public :: evaluator_get_tag
\nwendcode{}\nwbegincode{1079}\moddef{Evaluators: procedures}\plusendmoddef
  function evaluator_get_tag (eval) result (tag)
    integer :: tag
    type(evaluator_t), intent(in) :: eval
    tag = interaction_get_tag (eval%int)
  end function evaluator_get_tag

\eatline
\nwendcode{}\nwbegindocs{1080}Get the norm of the state matrix (if the norm has been taken out, otherwise
this would be unity).
\nwenddocs{}\nwbegincode{1081}\moddef{Evaluators: public}\plusendmoddef
  public :: evaluator_get_norm
\nwendcode{}\nwbegincode{1082}\moddef{Evaluators: procedures}\plusendmoddef
  function evaluator_get_norm (eval) result (norm)
    real(default) :: norm
    type(evaluator_t), intent(in) :: eval
    norm = interaction_get_norm (eval%int)
  end function evaluator_get_norm

\eatline
\nwendcode{}\nwbegindocs{1083}Return the number of particles.
\nwenddocs{}\nwbegincode{1084}\moddef{Evaluators: public}\plusendmoddef
  public :: evaluator_get_n_tot
  public :: evaluator_get_n_in
  public :: evaluator_get_n_vir
  public :: evaluator_get_n_out
\nwendcode{}\nwbegincode{1085}\moddef{Evaluators: procedures}\plusendmoddef
  function evaluator_get_n_tot (eval) result (n_tot)
    integer :: n_tot
    type(evaluator_t), intent(in) :: eval
    n_tot = interaction_get_n_tot (eval%int)
  end function evaluator_get_n_tot

  function evaluator_get_n_in (eval) result (n_in)
    integer :: n_in
    type(evaluator_t), intent(in) :: eval
    n_in = interaction_get_n_in (eval%int)
  end function evaluator_get_n_in

  function evaluator_get_n_vir (eval) result (n_vir)
    integer :: n_vir
    type(evaluator_t), intent(in) :: eval
    n_vir = interaction_get_n_vir (eval%int)
  end function evaluator_get_n_vir

  function evaluator_get_n_out (eval) result (n_out)
    integer :: n_out
    type(evaluator_t), intent(in) :: eval
    n_out = interaction_get_n_out (eval%int)
  end function evaluator_get_n_out

\eatline
\eatline
\nwendcode{}\nwbegindocs{1086}Return a specific matrix element.
\nwenddocs{}\nwbegincode{1087}\moddef{Evaluators: public}\plusendmoddef
  public :: evaluator_get_matrix_element
\nwendcode{}\nwbegincode{1088}\moddef{Evaluators: procedures}\plusendmoddef
  function evaluator_get_matrix_element (eval, i) result (value)
    complex(default) :: value
    type(evaluator_t), intent(in) :: eval
    integer :: i
    value = interaction_get_matrix_element (eval%int, i)
  end function evaluator_get_matrix_element

\eatline
\nwendcode{}\nwbegindocs{1089}Sum all matrix element values.
\nwenddocs{}\nwbegincode{1090}\moddef{Evaluators: public}\plusendmoddef
  public :: evaluator_sum
\nwendcode{}\nwbegincode{1091}\moddef{Evaluators: procedures}\plusendmoddef
  function evaluator_sum (eval) result (value)
    complex(default) :: value
    type(evaluator_t), intent(in) :: eval
    value = interaction_sum (eval%int)
  end function evaluator_sum

\eatline
\nwendcode{}\nwbegindocs{1092}Renormalize the state matrix by its trace, if nonzero.  The renormalization
is reflected in the state-matrix norm.
\nwenddocs{}\nwbegincode{1093}\moddef{Evaluators: public}\plusendmoddef
  public :: evaluator_normalize_by_trace
\nwendcode{}\nwbegincode{1094}\moddef{Evaluators: procedures}\plusendmoddef
  subroutine evaluator_normalize_by_trace (eval)
    type(evaluator_t), intent(inout) :: eval
    call interaction_normalize_by_trace (eval%int)
  end subroutine evaluator_normalize_by_trace

\eatline
\nwendcode{}\nwbegindocs{1095}Analogous, but renormalize by maximal (absolute) value.
\nwenddocs{}\nwbegincode{1096}\moddef{Evaluators: public}\plusendmoddef
  public :: evaluator_normalize_by_max
\nwendcode{}\nwbegincode{1097}\moddef{Evaluators: procedures}\plusendmoddef
  subroutine evaluator_normalize_by_max (eval)
    type(evaluator_t), intent(inout) :: eval
    call interaction_normalize_by_max (eval%int)
  end subroutine evaluator_normalize_by_max

\eatline
\nwendcode{}\nwbegindocs{1098}Append color-contracted states.  Matrix element array and
multiplication table are unaffected by this.
\nwenddocs{}\nwbegincode{1099}\moddef{Evaluators: public}\plusendmoddef
  public :: evaluator_add_color_contractions
\nwendcode{}\nwbegincode{1100}\moddef{Evaluators: procedures}\plusendmoddef
  subroutine evaluator_add_color_contractions (eval)
    type(evaluator_t), intent(inout) :: eval
    call interaction_add_color_contractions (eval%int)
  end subroutine evaluator_add_color_contractions

\eatline
\nwendcode{}\nwbegindocs{1101}Return the quantum-numbers mask of the enclosed interaction.
\nwenddocs{}\nwbegincode{1102}\moddef{Evaluators: public}\plusendmoddef
  public :: evaluator_get_mask
\nwendcode{}\nwbegincode{1103}\moddef{Evaluators: procedures}\plusendmoddef
  function evaluator_get_mask (eval) result (mask)
    type(quantum_numbers_mask_t), dimension(:), allocatable :: mask
    type(evaluator_t), intent(in), target :: eval
    allocate (mask (interaction_get_n_tot (eval%int)))
    mask = interaction_get_mask (eval%int)
  end function evaluator_get_mask

\eatline
\nwendcode{}\nwbegindocs{1104}Extend the linking of interactions to evaluators.
\nwenddocs{}\nwbegincode{1105}\moddef{Evaluators: public}\plusendmoddef
  public :: interaction_set_source_link
  public :: evaluator_set_source_link
\nwendcode{}\nwbegincode{1106}\moddef{Evaluators: interfaces}\plusendmoddef
  interface interaction_set_source_link
     module procedure interaction_set_source_link_eval
  end interface
  interface evaluator_set_source_link
     module procedure evaluator_set_source_link_int
     module procedure evaluator_set_source_link_eval
  end interface
\eatline
\nwendcode{}\nwbegincode{1107}\moddef{Evaluators: procedures}\plusendmoddef
  subroutine interaction_set_source_link_eval (int, i, eval1, i1)
    type(interaction_t), intent(inout) :: int
    type(evaluator_t), intent(in), target :: eval1
    integer, intent(in) :: i, i1
    call interaction_set_source_link (int, i, eval1%int, i1)
  end subroutine interaction_set_source_link_eval

  subroutine evaluator_set_source_link_int (eval, i, int1, i1)
    type(evaluator_t), intent(inout) :: eval
    type(interaction_t), intent(in), target :: int1
    integer, intent(in) :: i, i1
    call interaction_set_source_link (eval%int, i, int1, i1)
  end subroutine evaluator_set_source_link_int

  subroutine evaluator_set_source_link_eval (eval, i, eval1, i1)
    type(evaluator_t), intent(inout) :: eval
    type(evaluator_t), intent(in), target :: eval1
    integer, intent(in) :: i, i1
    call interaction_set_source_link (eval%int, i, eval1%int, i1)
  end subroutine evaluator_set_source_link_eval

\eatline
\eatline
\eatline
\nwendcode{}\nwbegindocs{1108}Receive from / send momenta to the linked interactions.
\nwenddocs{}\nwbegincode{1109}\moddef{Evaluators: public}\plusendmoddef
  public :: evaluator_receive_momenta
  public :: evaluator_send_momenta
\nwendcode{}\nwbegincode{1110}\moddef{Evaluators: procedures}\plusendmoddef
  subroutine evaluator_receive_momenta (eval)
    type(evaluator_t), intent(inout) :: eval
    call interaction_receive_momenta (eval%int)
  end subroutine evaluator_receive_momenta

  subroutine evaluator_send_momenta (eval)
    type(evaluator_t), intent(in) :: eval
    call interaction_send_momenta (eval%int)
  end subroutine evaluator_send_momenta

\eatline
\eatline
\nwendcode{}\nwbegindocs{1111}Reassign external source links from one to another.
\nwenddocs{}\nwbegincode{1112}\moddef{Evaluators: public}\plusendmoddef
  public :: evaluator_reassign_links
\nwendcode{}\nwbegincode{1113}\moddef{Evaluators: interfaces}\plusendmoddef
  interface evaluator_reassign_links
     module procedure evaluator_reassign_links_eval
     module procedure evaluator_reassign_links_int
  end interface

\nwendcode{}\nwbegincode{1114}\moddef{Evaluators: procedures}\plusendmoddef
  subroutine evaluator_reassign_links_eval (eval, eval_src, eval_target)
    type(evaluator_t), intent(inout) :: eval
    type(evaluator_t), intent(in) :: eval_src
    type(evaluator_t), intent(in), target :: eval_target
    if (associated (eval%int_in1)) then
       if (interaction_get_tag (eval%int_in1) &
            == interaction_get_tag (eval_src%int)) then
          eval%int_in1 => eval_target%int
       end if
    end if
    if (associated (eval%int_in2)) then
       if (interaction_get_tag (eval%int_in2) &
            == interaction_get_tag (eval_src%int)) then
          eval%int_in2 => eval_target%int
       end if
    end if
    call interaction_reassign_links (eval%int, eval_src%int, eval_target%int)
  end subroutine evaluator_reassign_links_eval

  subroutine evaluator_reassign_links_int (eval, int_src, int_target)
    type(evaluator_t), intent(inout) :: eval
    type(interaction_t), intent(in) :: int_src
    type(interaction_t), intent(in), target :: int_target
    if (associated (eval%int_in1)) then
       if (interaction_get_tag (eval%int_in1) &
            == interaction_get_tag (int_src)) then
          eval%int_in1 => int_target
       end if
    end if
    if (associated (eval%int_in2)) then
       if (interaction_get_tag (eval%int_in2) &
            == interaction_get_tag (int_src)) then
          eval%int_in2 => int_target
       end if
    end if
    call interaction_reassign_links (eval%int, int_src, int_target)
  end subroutine evaluator_reassign_links_int

\eatline
\nwendcode{}\nwbegindocs{1115}Return flavor, momentum, and position of the first unstable particle
present in the interaction.
\nwenddocs{}\nwbegincode{1116}\moddef{Evaluators: public}\plusendmoddef
  public :: evaluator_get_unstable_particle
\nwendcode{}\nwbegincode{1117}\moddef{Evaluators: procedures}\plusendmoddef
  subroutine evaluator_get_unstable_particle (eval, flv, p, i)
    type(evaluator_t), intent(in) :: eval
    type(flavor_t), intent(out) :: flv
    type(vector4_t), intent(out) :: p
    integer, intent(out) :: i
    call interaction_get_unstable_particle (eval%int, flv, p, i)
  end subroutine evaluator_get_unstable_particle

\eatline
\nwendcode{}\nwbegindocs{1118}\nwdocspar
\subsection{Deleting the evaluator}
Only the interaction component needs finalization.
\nwenddocs{}\nwbegincode{1119}\moddef{Evaluators: public}\plusendmoddef
  public :: evaluator_final
\nwendcode{}\nwbegincode{1120}\moddef{Evaluators: procedures}\plusendmoddef
  subroutine evaluator_final (eval)
    type(evaluator_t), intent(inout) :: eval
    call interaction_final (eval%int)
  end subroutine evaluator_final

\eatline
\nwendcode{}\nwbegindocs{1121}\nwdocspar
\subsection{Creating an evaluator: identity}
The identity evaluator creates a copy of the first input evaluator; the second
input is not used.

All particles link back to the input evaluatorand  the internal
relations are copied. As evaluation does take a shortcut by cloning the matrix
elements, the pairing array is not used and does not have to be set up.

\nwenddocs{}\nwbegincode{1122}\moddef{Evaluators: public}\plusendmoddef
  public :: evaluator_init_identity
\nwendcode{}\nwbegincode{1123}\moddef{Evaluators: interfaces}\plusendmoddef
  interface evaluator_init_identity
     module procedure evaluator_init_identity_i
     module procedure evaluator_init_identity_e
  end interface

\nwendcode{}\nwbegincode{1124}\moddef{Evaluators: procedures}\plusendmoddef
  subroutine evaluator_init_identity_i (eval, int)
    type(evaluator_t), intent(out), target :: eval
    type(interaction_t), intent(in), target :: int
    integer :: n_in, n_out, n_vir, n_tot
    integer :: i
    integer, dimension(:), allocatable :: map
    type(state_matrix_t), pointer :: state
    type(state_iterator_t) :: it

    eval%type = EVAL_IDENTITY
    eval%int_in1 => int
    nullify (eval%int_in2)
    n_in = interaction_get_n_in (int)
    n_out = interaction_get_n_out (int)
    n_vir = interaction_get_n_vir (int)
    n_tot = interaction_get_n_tot (int)
    call interaction_init (eval%int, n_in, n_vir, n_out, &
       mask=interaction_get_mask (int), &
       resonant=interaction_get_resonance_flags (int))
    do i = 1, n_tot
       call interaction_set_source_link (eval%int, i, int, i)
    end do
    allocate (map(n_tot))
    map = [(i, i = 1, n_tot)]
    call interaction_transfer_relations (int, eval%int, map)
    state => interaction_get_state_matrix_ptr (int)
    call it%init (state)
    do while (it%is_valid ())
       call interaction_add_state (eval%int, &
          it%get_quantum_numbers (), it%get_me_index ())
       call it%advance ()
    end do
    call interaction_freeze (eval%int)

  end subroutine evaluator_init_identity_i

  subroutine evaluator_init_identity_e (eval, eval1)
    type(evaluator_t), intent(out), target :: eval
    type(evaluator_t), intent(in), target :: eval1
    call evaluator_init_identity_i (eval, eval1%int)
  end subroutine evaluator_init_identity_e

\eatline
\eatline
\eatline
\nwendcode{}\nwbegindocs{1125}\nwdocspar
\subsection {Creating an evaluator: quantum number sum}
This evaluator operates on the diagonal of a density matrix and sums over the
quantum numbers specified by the mask. The optional argument {\tt{}drop} allows to
drop a particle from the resulting density matrix. The handling of virtuals is
not completely sane, especially in connection with dropping particles.

When summing over matrix element entries, we keep the separation into
entries and normalization (in the corresponding evaluation routine below).
\nwenddocs{}\nwbegincode{1126}\moddef{Evaluators: public}\plusendmoddef
  public :: evaluator_init_qn_sum
\nwendcode{}\nwbegincode{1127}\moddef{Evaluators: interfaces}\plusendmoddef
  interface evaluator_init_qn_sum
     module procedure evaluator_init_qn_sum_i
     module procedure evaluator_init_qn_sum_e
  end interface
\nwendcode{}\nwbegincode{1128}\moddef{Evaluators: procedures}\plusendmoddef
  subroutine evaluator_init_qn_sum_i (eval, int, qn_mask, drop)
    type(evaluator_t), intent(out), target :: eval
    type(interaction_t), target, intent(in) :: int
    type(quantum_numbers_mask_t), dimension(:), intent(in) :: qn_mask
    logical, intent(in), optional, dimension(:) :: drop
    type(state_iterator_t) :: it_old, it_new
    integer, dimension(:,:), allocatable :: pairing_proto
    integer, dimension(:), allocatable :: pairing_sizes
    integer, dimension(:), allocatable :: map
    integer :: n_in, n_out, n_vir, n_tot, n_me_old, n_me_new
    integer :: i, j
    type(state_matrix_t), pointer :: state_new, state_old
    type(quantum_numbers_t), dimension(:), allocatable :: qn
    logical :: matched
    logical, dimension(size (qn_mask)) :: dropped
    integer :: ndropped
    integer, dimension(:), allocatable :: inotdropped
    type(quantum_numbers_mask_t), dimension(:), allocatable :: mask
    logical, dimension(:), allocatable :: resonant
 
    eval%type = EVAL_QN_SUM
    eval%int_in1 => int
    nullify (eval%int_in2)
    if (present (drop)) then
       dropped = drop
    else
       dropped = .false.
    end if
    ndropped = count (dropped)

    n_in = interaction_get_n_in (int)
    n_out = interaction_get_n_out (int) - ndropped
    n_vir = interaction_get_n_vir (int)
    n_tot = interaction_get_n_tot (int) - ndropped

    allocate (inotdropped (n_tot))
    i = 1
    do j = 1, n_tot + ndropped
       if (dropped (j)) cycle
       inotdropped(i) = j
       i = i + 1
    end do

    allocate (mask(n_tot + ndropped))
    mask = interaction_get_mask (int)
    allocate (resonant(n_tot + ndropped))
    resonant = interaction_get_resonance_flags (int)
    call interaction_init (eval%int, n_in, n_vir, n_out, &
       mask = mask(inotdropped) .or. qn_mask(inotdropped), &
       resonant = resonant(inotdropped))
    i = 1
    do j = 1, n_tot + ndropped
       if (dropped(j)) cycle
       call interaction_set_source_link (eval%int, i, int, j)
       i = i + 1
    end do    
    allocate (map(n_tot + ndropped))
    i = 1
    do j = 1, n_tot + ndropped
       if (dropped (j)) then
          map(j) = 0
       else
          map(j) = i
          i = i + 1
       end if
    end do
    call interaction_transfer_relations (int, eval%int, map)

    n_me_old = interaction_get_n_matrix_elements (int)
    allocate (pairing_proto(n_me_old, n_me_old))
    allocate (pairing_sizes(n_me_old))
    pairing_sizes = 0
    state_old => interaction_get_state_matrix_ptr (int)
    state_new => interaction_get_state_matrix_ptr (eval%int)
    call it_old%init (state_old)
    allocate (qn(n_tot + ndropped))
    do while (it_old%is_valid ())
       qn = it_old%get_quantum_numbers ()
       if (.not. all (qn%are_diagonal ())) then
          call it_old%advance ()
          cycle
       end if
       matched = .false.
       call it_new%init (state_new)
       if (interaction_get_n_matrix_elements (eval%int) > 0) then
          do while (it_new%is_valid ())
             if (all (qn(inotdropped) .match. &
                it_new%get_quantum_numbers ())) &
             then
                matched = .true.
                i = it_new%get_me_index ()
                exit
             end if
             call it_new%advance ()
          end do
       end if
       if (.not. matched) then
          call interaction_add_state (eval%int, qn(inotdropped))
          i = interaction_get_n_matrix_elements (eval%int)
       end if
       pairing_sizes(i) = pairing_sizes(i) + 1
       pairing_proto(pairing_sizes(i), i) = it_old%get_me_index ()
       ! print *, &                         !!! Debugging
       !     i, pairing_sizes(i), it_old%get_me_index ()
       call it_old%advance ()
    end do
    call interaction_freeze (eval%int)

    n_me_new = interaction_get_n_matrix_elements (eval%int)
    allocate (eval%pairing_array(n_me_new))
    do i = 1, n_me_new
       ! print *, i, pairing_sizes(i)     !!! Debugging
       call pairing_array_init (eval%pairing_array(i), &
          pairing_sizes(i), .false., .false.)
       eval%pairing_array(i)%i1 = pairing_proto(:pairing_sizes(i), i) 
    end do

  end subroutine evaluator_init_qn_sum_i

  subroutine evaluator_init_qn_sum_e (eval, eval1, qn_mask, drop)
    type(evaluator_t), intent(out) :: eval
    type(evaluator_t), intent(in), target :: eval1
    type(quantum_numbers_mask_t), dimension(:), intent(in) :: qn_mask
    logical, dimension(:), optional, intent(in) :: drop
    call evaluator_init_qn_sum_i (eval, eval1%int, qn_mask, drop)
  end subroutine

\eatline
\nwendcode{}\nwbegindocs{1129}\nwdocspar

\subsection{Evaluation}
When the input interactions (which are pointed to in the pairings
stored within the evaluator) are filled with values, we can activate
the evaluator, i.e., calculate the result values which are stored in
the interaction.

The evaluation of matrix elements can be done in parallel.  A
{\tt{}forall} construct is not appropriate, however.  We would need 
{\tt{}do\ concurrent} here.  Nevertheless, the evaluation functions are
marked as {\tt{}pure}.
\nwenddocs{}\nwbegincode{1130}\moddef{Evaluators: evaluator: TBP}\plusendmoddef
  procedure :: evaluate => evaluator_evaluate
\nwendcode{}\nwbegincode{1131}\moddef{Evaluators: procedures}\plusendmoddef
  subroutine evaluator_evaluate (eval)
    class(evaluator_t), intent(inout), target :: eval
    integer :: i
    select case (eval%type)
    case (EVAL_PRODUCT)
       do i = 1, size(eval%pairing_array)
          call interaction_evaluate_product (eval%int, i, &
               eval%int_in1, eval%int_in2, &
               eval%pairing_array(i)%i1, eval%pairing_array(i)%i2)
       end do
    case (EVAL_SQUARE_WITH_COLOR_FACTORS)
       do i = 1, size(eval%pairing_array)
          call interaction_evaluate_product_cf (eval%int, i, &
               eval%int_in1, eval%int_in1, &
               eval%pairing_array(i)%i1, eval%pairing_array(i)%i2, &
               eval%pairing_array(i)%factor)
       end do
    case (EVAL_SQUARED_FLOWS)
       do i = 1, size(eval%pairing_array)
          call interaction_evaluate_square_c (eval%int, i, &
               eval%int_in1, &
               eval%pairing_array(i)%i1)
       end do
    case (EVAL_COLOR_CONTRACTION)
       do i = 1, size(eval%pairing_array)
          call interaction_evaluate_sum (eval%int, i, &
               eval%int_in1, &
               eval%pairing_array(i)%i1)
       end do
    case (EVAL_IDENTITY)
       call interaction_set_matrix_element (eval%int, eval%int_in1)
    case (EVAL_QN_SUM)
       do i = 1, size (eval%pairing_array)
          call interaction_evaluate_me_sum (eval%int, i, &
             eval%int_in1, eval%pairing_array(i)%i1)
          call interaction_set_norm &
               (eval%int, interaction_get_norm (eval%int_in1))
       end do
    end select
  end subroutine evaluator_evaluate

\eatline
\nwendcode{}\nwbegindocs{1132}\nwdocspar
\subsection{Unit tests}
Test: Create two interactions.  The interactions are twofold
connected.  The first connection has a helicity index that is kept,
the second connection has a helicity index that is summed over.
Concatenate the interactions in an evaluator, which thus contains a
result interaction.  Fill the input interactions with values, activate
the evaluator and print the result.
\nwenddocs{}\nwbegincode{1133}\moddef{Evaluators: public}\plusendmoddef
  public :: evaluator_test
\nwendcode{}\nwbegincode{1134}\moddef{Evaluators: procedures}\plusendmoddef
  subroutine evaluator_test (u, results)
    integer, intent(in) :: u
    type(test_results_t), intent(inout) :: results
  \LA{}Evaluators: execute tests\RA{}
  end subroutine evaluator_test

\eatline
\nwendcode{}\nwbegincode{1135}\moddef{Evaluators: execute tests}\endmoddef
  call test (evaluator_1, "evaluator_1", &
       "check evaluators (1)", &
       u, results)
\nwendcode{}\nwbegincode{1136}\moddef{Evaluators: tests}\endmoddef
  subroutine evaluator_1 (u) 
    integer, intent(in) :: u
    type(model_data_t), target :: model       
    type(interaction_t), target :: int_qqtt, int_tbw, int1, int2
    type(flavor_t), dimension(:), allocatable :: flv
    type(color_t), dimension(:), allocatable :: col
    type(helicity_t), dimension(:), allocatable :: hel
    type(quantum_numbers_t), dimension(:), allocatable :: qn
    integer :: f, c, h1, h2, h3
    type(vector4_t), dimension(4) :: p
    type(vector4_t), dimension(2) :: q
    type(quantum_numbers_mask_t) :: qn_mask_conn
    type(quantum_numbers_mask_t), dimension(:), allocatable :: qn_mask2
    type(evaluator_t), target :: eval, eval2, eval3

    call model%init_sm_test ()
    
    write (u, "(A)")   "*** Evaluator for matrix product"
    write (u, "(A)")   "***   Construct interaction for qq -> tt"
    write (u, "(A)")
    call interaction_init (int_qqtt, 2, 0, 2, set_relations=.true.)
    allocate (flv (4), col (4), hel (4), qn (4))
    allocate (qn_mask2 (4))
    do c = 1, 2
       select case (c)
       case (1)
          call col%init_col_acl ([1, 0, 1, 0], [0, 2, 0, 2])
       case (2)
          call col%init_col_acl ([1, 0, 2, 0], [0, 1, 0, 2])
       end select
       do f = 1, 2
          call flv%init ([f, -f, 6, -6], model)
          do h1 = -1, 1, 2
             call hel(3)%init (h1)
             do h2 = -1, 1, 2
                call hel(4)%init (h2)
                call qn%init (flv, col, hel)
                call interaction_add_state (int_qqtt, qn)
             end do
          end do
       end do
    end do
    call interaction_freeze (int_qqtt)
    deallocate (flv, col, hel, qn)
    write (u, "(A)")  "***   Construct interaction for t -> bW"
    call interaction_init (int_tbw, 1, 0, 2, set_relations=.true.)
    allocate (flv (3), col (3), hel (3), qn (3))
    call flv%init ([6, 5, 24], model)
    call col%init_col_acl ([1, 1, 0], [0, 0, 0])
    do h1 = -1, 1, 2
       call hel(1)%init (h1)
       do h2 = -1, 1, 2
          call hel(2)%init (h2)
          do h3 = -1, 1
             call hel(3)%init (h3)
             call qn%init (flv, col, hel)
             call interaction_add_state (int_tbw, qn)
          end do
       end do
    end do
    call interaction_freeze (int_tbw)
    deallocate (flv, col, hel, qn)
    write (u, "(A)")  "***   Link interactions"
    call interaction_set_source_link (int_tbw, 1, int_qqtt, 3)
    qn_mask_conn = quantum_numbers_mask (.false.,.false.,.true.)
    write (u, "(A)")
    write (u, "(A)")  "***   Show input"
    call interaction_write (int_qqtt, unit = u)
    write (u, "(A)")
    call interaction_write (int_tbw, unit = u)
    write (u, "(A)")
    write (u, "(A)")  "***   Evaluate product"
    call evaluator_init_product &
         (eval, int_qqtt, int_tbw, qn_mask_conn)
    call eval%write (unit = u)

     call interaction_init (int1, 2, 0, 2, set_relations=.true.)
     call interaction_init (int2, 1, 0, 2, set_relations=.true.)
     p(1) = vector4_moving (1000._default, 1000._default, 3)
     p(2) = vector4_moving (200._default, 200._default, 2)
     p(3) = vector4_moving (100._default, 200._default, 1)
     p(4) = p(1) - p(2) - p(3)
     call interaction_set_momenta (int1, p)
     q(1) = vector4_moving (50._default,-50._default, 3)
     q(2) = p(2) + p(4) - q(1)
     call interaction_set_momenta (int2, q, outgoing=.true.)
     call interaction_set_matrix_element &
          (int1, [(2._default,0._default), &
          (4._default,1._default), (-3._default,0._default)])
     call interaction_set_matrix_element &
          (int2, [(-3._default,0._default), &
          (0._default,1._default), (1._default,2._default)])
     call evaluator_receive_momenta (eval)
     call eval%evaluate ()
     call interaction_write (int1, unit = u)
     write (u, "(A)")
     call interaction_write (int2, unit = u)
     write (u, "(A)")
     call eval%write (unit = u)
     write (u, "(A)")
     call interaction_final (int1)
     call interaction_final (int2)
     call evaluator_final (eval)

     write (u, "(A)")
     write (u, "(A)")   "*** Evaluator for matrix square"
     allocate (flv(4), col(4), qn(4))
     call interaction_init (int1, 2, 0, 2, set_relations=.true.)
     call flv%init ([1, -1, 21, 21], model)
     call col(1)%init ([1])
     call col(2)%init ([-2])
     call col(3)%init ([2, -3])
     call col(4)%init ([3, -1])
     call qn%init (flv, col)
     call interaction_add_state (int1, qn)
     call col(3)%init ([3, -1])
     call col(4)%init ([2, -3])
     call qn%init (flv, col)
     call interaction_add_state (int1, qn)
     call col(3)%init ([2, -1])
     call col(4)%init (.true.)
     call qn%init (flv, col)
     call interaction_add_state (int1, qn)
     call interaction_freeze (int1)
     ! [qn_mask2 not set since default is false]
     call evaluator_init_square (eval, int1, qn_mask2, nc=3)
     call evaluator_init_square_nondiag (eval2, int1, qn_mask2)
     qn_mask2 = quantum_numbers_mask (.false., .true., .true.)
     call evaluator_init_square_diag (eval3, eval%int, qn_mask2)
     call interaction_set_matrix_element &
          (int1, [(2._default,0._default), &
          (4._default,1._default), (-3._default,0._default)])
     call interaction_set_momenta (int1, p)
     call interaction_write (int1, unit = u)
     write (u, "(A)")
     call evaluator_receive_momenta (eval)
     call eval%evaluate ()
     call eval%write (unit = u)
     write (u, "(A)")
     call evaluator_receive_momenta (eval2)
     call eval2%evaluate ()
     call eval2%write (unit = u)
     write (u, "(A)")
     call evaluator_receive_momenta (eval3)
     call eval3%evaluate ()
     call eval3%write (unit = u)
     call interaction_final (int1)
     call evaluator_final (eval)
     call evaluator_final (eval2)
     call evaluator_final (eval3)
     
     call model%final ()
  end subroutine evaluator_1

\eatline
\nwendcode{}\nwbegindocs{1137}\nwdocspar
\nwenddocs{}\nwbegincode{1138}\moddef{Evaluators: execute tests}\plusendmoddef
  call test (evaluator_2, "evaluator_2", &
       "check evaluators (2)", &
       u, results)
\nwendcode{}\nwbegincode{1139}\moddef{Evaluators: tests}\plusendmoddef
  subroutine evaluator_2 (u)
    integer, intent(in) :: u
    type(model_data_t), target :: model       
    type(interaction_t), target :: int
    integer :: h1, h2, h3, h4
    type(helicity_t), dimension(4) :: hel
    type(color_t), dimension(4) :: col
    type(flavor_t), dimension(4) :: flv
    type(quantum_numbers_t), dimension(4) :: qn
    type(vector4_t), dimension(4) :: p
    type(evaluator_t) :: eval
    integer :: i
    
    call model%init_sm_test ()
        
    write (u, "(A)") "*** Creating interaction for e+ e- -> W+ W-"
    write (u, "(A)") 
    
    call flv%init ([11, -11, 24, -24], model)
    do i = 1, 4
       call col(i)%init ()
    end do
    call interaction_init (int, 2, 0, 2, set_relations=.true.)
    do h1 = -1, 1, 2
       call hel(1)%init (h1)
       do h2 = -1, 1, 2
          call hel(2)%init (h2)
          do h3 = -1, 1
             call hel(3)%init (h3)
             do h4 = -1, 1
                call hel(4)%init (h4)
                call qn%init (flv, col, hel)
                call interaction_add_state (int, qn)
             end do
          end do
       end do
    end do
    call interaction_freeze (int)
    call interaction_set_matrix_element (int, &
       [(cmplx (i, kind=default), i = 1, 36)])
    p(1) = vector4_moving (1000._default, 1000._default, 3)    
    p(2) = vector4_moving (1000._default, -1000._default, 3)
    p(3) = vector4_moving (1000._default, &
       sqrt (1E6_default - 80._default**2), 3)
    p(4) = p(1) + p(2) - p(3)
    call interaction_set_momenta (int, p)
    write (u, "(A)") "*** Setting up evaluator"
    write (u, "(A)")
    
    call evaluator_init_identity (eval, int)
    write (u, "(A)") "*** Transferring momenta and evaluating"
    write (u, "(A)")
    
    call evaluator_receive_momenta (eval)
    call eval%evaluate ()
    write (u, "(A)")  "*******************************************************"
    write (u, "(A)")  "   Interaction dump"
    write (u, "(A)")  "*******************************************************"
    call interaction_write (int, unit = u)
    write (u, "(A)")  
    write (u, "(A)")  "*******************************************************"
    write (u, "(A)")  "   Evaluator dump"
    write (u, "(A)")  "*******************************************************"
    call eval%write (unit = u)
    write (u, "(A)")  
    write (u, "(A)")   "*** cleaning up"
    call interaction_final (int)
    call evaluator_final (eval)
    
    call model%final ()
  end subroutine evaluator_2

\nwendcode{}\nwbegindocs{1140} %def evaluator_2
\nwenddocs{}\nwbegindocs{1141}\nwdocspar
\nwenddocs{}\nwbegincode{1142}\moddef{Evaluators: execute tests}\plusendmoddef
  call test (evaluator_3, "evaluator_3", &
       "check evaluators (3)", &
       u, results)
\nwendcode{}\nwbegincode{1143}\moddef{Evaluators: tests}\plusendmoddef
  subroutine evaluator_3 (u)   
    integer, intent(in) :: u
    type(model_data_t), target :: model
    type(interaction_t), target :: int
    integer :: h1, h2, h3, h4
    type(helicity_t), dimension(4) :: hel
    type(color_t), dimension(4) :: col
    type(flavor_t), dimension(4) :: flv1, flv2
    type(quantum_numbers_t), dimension(4) :: qn
    type(vector4_t), dimension(4) :: p
    type(evaluator_t) :: eval1, eval2, eval3
    type(quantum_numbers_mask_t), dimension(4) :: qn_mask
    integer :: i
    
    call model%init_sm_test ()
            
    write (u, "(A)")  "*** Creating interaction for e+/mu+ e-/mu- -> W+ W-"
    call flv1%init ([11, -11, 24, -24], model)
    call flv2%init ([13, -13, 24, -24], model)
    do i = 1, 4
       call col (i)%init ()
    end do
    call interaction_init (int, 2, 0, 2, set_relations=.true.)
    do h1 = -1, 1, 2
       call hel(1)%init (h1)
       do h2 = -1, 1, 2
          call hel(2)%init (h2)
          do h3 = -1, 1
             call hel(3)%init (h3)
             do h4 = -1, 1
                call hel(4)%init (h4)
                call qn%init (flv1, col, hel)
                call interaction_add_state (int, qn)
                call qn%init (flv2, col, hel)
                call interaction_add_state (int, qn)
             end do
          end do
       end do
    end do
    call interaction_freeze (int)
    call interaction_set_matrix_element (int, &
       [(cmplx (1, kind=default), i = 1, 72)])
    p(1) = vector4_moving (1000._default, 1000._default, 3)    
    p(2) = vector4_moving (1000._default, -1000._default, 3)
    p(3) = vector4_moving (1000._default, &
       sqrt (1E6_default - 80._default**2), 3)
    p(4) = p(1) + p(2) - p(3)
    call interaction_set_momenta (int, p)
    write (u, "(A)")  "*** Setting up evaluators"
    call qn_mask%init (.false., .true., .true.)
    call evaluator_init_qn_sum (eval1, int, qn_mask)
    call qn_mask%init (.true., .true., .true.)
    call evaluator_init_qn_sum (eval2, int, qn_mask)
    call qn_mask%init (.false., .true., .false.)
    call evaluator_init_qn_sum (eval3, int, qn_mask, &
      [.false., .false., .false., .true.])
    write (u, "(A)")  "*** Transferring momenta and evaluating"
    call evaluator_receive_momenta (eval1)
    call eval1%evaluate ()
    call evaluator_receive_momenta (eval2)
    call eval2%evaluate ()
    call evaluator_receive_momenta (eval3)
    call eval3%evaluate ()
    write (u, "(A)")  "*******************************************************"
    write (u, "(A)")  "   Interaction dump"
    write (u, "(A)")  "*******************************************************"
    call interaction_write (int, unit = u)
    write (u, "(A)")  
    write (u, "(A)")  "*******************************************************"
    write (u, "(A)")  "   Evaluator dump --- spin sum"
    write (u, "(A)")  "*******************************************************"
    call eval1%write (unit = u)
    call interaction_write (evaluator_get_int_ptr (eval1), unit = u)
    write (u, "(A)")  "*******************************************************"
    write (u, "(A)")  "   Evaluator dump --- spin / flavor sum"
    write (u, "(A)")  "*******************************************************"
    call eval2%write (unit = u)
    call interaction_write (evaluator_get_int_ptr (eval2), unit = u)
    write (u, "(A)")  "*******************************************************"
    write (u, "(A)")  "   Evaluator dump --- flavor sum, drop last W"
    write (u, "(A)")  "*******************************************************"
    call eval3%write (unit = u)
    call interaction_write (evaluator_get_int_ptr (eval3), unit = u)
    write (u, "(A)")  
    write (u, "(A)")  "*** cleaning up"
    call interaction_final (int)
    call evaluator_final (eval1)
    call evaluator_final (eval2)
    call evaluator_final (eval3)
    
    call model%final ()
  end subroutine evaluator_3

\eatline
\nwendcode{}\nwbegindocs{1144}This test evaluates a product with different quantum-number masks and
filters for the linked entry.
\nwenddocs{}\nwbegincode{1145}\moddef{Evaluators: execute tests}\plusendmoddef
  call test (evaluator_4, "evaluator_4", &
       "check evaluator product with filter", &
       u, results)
\nwendcode{}\nwbegincode{1146}\moddef{Evaluators: tests}\plusendmoddef
  subroutine evaluator_4 (u)   
    integer, intent(in) :: u
    type(model_data_t), target :: model
    type(interaction_t), target :: int1, int2
    integer :: h1, h2, h3
    type(helicity_t), dimension(3) :: hel
    type(color_t), dimension(3) :: col
    type(flavor_t), dimension(2) :: flv1, flv2
    type(flavor_t), dimension(3) :: flv3, flv4
    type(quantum_numbers_t), dimension(3) :: qn
    type(evaluator_t) :: eval1, eval2, eval3, eval4
    type(quantum_numbers_mask_t) :: qn_mask
    type(flavor_t) :: flv_filter
    type(helicity_t) :: hel_filter
    type(color_t) :: col_filter
    type(quantum_numbers_t) :: qn_filter
    integer :: i
    
    write (u, "(A)")  "* Test output: evaluator_4"
    write (u, "(A)")  "*   Purpose: test evaluator products &
         &with mask and filter"
    write (u, "(A)")

    call model%init_sm_test ()
            
    write (u, "(A)")  "* Creating interaction for e- -> W+/Z"
    write (u, "(A)")
    
    call flv1%init ([11, 24], model)
    call flv2%init ([11, 23], model)
    do i = 1, 3
       call col(i)%init ()
    end do
    call interaction_init (int1, 1, 0, 1, set_relations=.true.)
    do h1 = -1, 1, 2
       call hel(1)%init (h1)
       do h2 = -1, 1
          call hel(2)%init (h2)
          call qn(:2)%init (flv1, col(:2), hel(:2))
          call interaction_add_state (int1, qn(:2))
          call qn(:2)%init (flv2, col(:2), hel(:2))
          call interaction_add_state (int1, qn(:2))
       end do
    end do
    call interaction_freeze (int1)
    call interaction_write (int1, u)

    write (u, "(A)")
    write (u, "(A)")  "* Creating interaction for W+/Z -> u ubar/dbar"
    write (u, "(A)")

    call flv3%init ([24, 2, -1], model)
    call flv4%init ([23, 2, -2], model)

    call interaction_init (int2, 1, 0, 2, set_relations=.true.)
    do h1 = -1, 1
       call hel(1)%init (h1)
       do h2 = -1, 1, 2
          call hel(2)%init (h2)
          do h3 = -1, 1, 2
             call hel(3)%init (h3)
             call qn(:3)%init (flv3, col(:3), hel(:3))
             call interaction_add_state (int2, qn(:3))
             call qn(:3)%init (flv4, col(:3), hel(:3))
             call interaction_add_state (int2, qn(:3))
          end do
       end do
    end do
    call interaction_freeze (int2)

    call interaction_set_source_link (int2, 1, int1, 2)
    call interaction_write (int2, u)

    write (u, "(A)")
    write (u, "(A)")  "* Product evaluator"
    write (u, "(A)")

    call qn_mask%init (.false., .false., .false.)
    call evaluator_init_product (eval1, int1, int2, &
         qn_mask_conn = qn_mask)
    call evaluator_write (eval1, u)

    write (u, "(A)")
    write (u, "(A)")  "* Product evaluator with helicity mask"
    write (u, "(A)")

    call qn_mask%init (.false., .false., .true.)
    call evaluator_init_product (eval2, int1, int2, &
         qn_mask_conn = qn_mask)
    call evaluator_write (eval2, u)

    write (u, "(A)")
    write (u, "(A)")  "* Product with flavor filter and helicity mask"
    write (u, "(A)")

    call qn_mask%init (.false., .false., .true.)
    call flv_filter%init (24, model)
    call hel_filter%init ()
    call col_filter%init ()
    call qn_filter%init (flv_filter, col_filter, hel_filter)
    call evaluator_init_product (eval3, int1, int2, &
         qn_mask_conn = qn_mask, qn_filter_conn = qn_filter)
    call evaluator_write (eval3, u)

    write (u, "(A)")
    write (u, "(A)")  "* Product with helicity filter and mask"
    write (u, "(A)")

    call qn_mask%init (.false., .false., .true.)
    call flv_filter%init ()
    call hel_filter%init (0)
    call col_filter%init ()
    call qn_filter%init (flv_filter, col_filter, hel_filter)
    call evaluator_init_product (eval4, int1, int2, &
         qn_mask_conn = qn_mask, qn_filter_conn = qn_filter)
    call evaluator_write (eval4, u)

    write (u, "(A)")
    write (u, "(A)")  "* Cleanup"
    
    call evaluator_final (eval1)
    call evaluator_final (eval2)
    call evaluator_final (eval3)
    call evaluator_final (eval4)
    
    call interaction_final (int1)
    call interaction_final (int2)
    
    call model%final ()

    write (u, "(A)")
    write (u, "(A)")  "* Test output end: evaluator_4"

  end subroutine evaluator_4

\eatline
\nwendcode{}\nwbegindocs{1147}\nwdocspar
  
\nwenddocs{}\nwfilename{../../../src/types/types.nw}\nwbegindocs{0}%% -*- ess-noweb-default-code-mode: f90-mode; noweb-default-code-mode: f90-mode; -*- 
% WHIZARD code as NOWEB source: common types and objects
%%%%%%%%%%%%%%%%%%%%%%%%%%%%%%%%%%%%%%%%%%%%%%%%%%%%%%%%%%%%%%%%%%%%%%% 
\chapter{Sindarin Built-In Types}

Here, we define a couple of types and objects which are useful both
internally for \whizard, and visible to the user, so they correspond
to Sindarin types.
\begin{description}
\item[particle\_specifiers]
  Expressions for particles and particle alternatives, involving
  particle names.
\item[pdg\_arrays]
  Integer (PDG) codes for particles.  Useful for particle aliases
  (e.g., 'quark' for $u,d,s$ etc.).
\item[jets]
  Define (pseudo)jets as objects.  Functional only if the {\tt{}fastjet} library
  is linked.  (This may change in the future.)
\item[subevents]
  Particle collections built from event records, for use in analysis and other
  Sindarin expressions
\item[analysis]
  Observables, histograms, and plots.
\end{description}

%%%%%%%%%%%%%%%%%%%%%%%%%%%%%%%%%%%%%%%%%%%%%%%%%%%%%%%%%%%%%%%%%%%%%%%%
\section{Particle Specifiers}
In this module we introduce a type for specifying a particle or particle
alternative.  In addition to the particle specifiers (strings separated by
colons), the type contains an optional flag {\tt{}polarized} and a string
{\tt{}decay}.  If the {\tt{}polarized} flag is set, particle polarization
information should be kept when generating events for this process.  If the
{\tt{}decay} string is set, it is the ID of a decay process which should be
applied to this particle when generating events.

In input/output form, the {\tt{}polarized} flag is indicated by an asterisk
{\tt{}(*)} in brackets, and the {\tt{}decay} is indicated by its ID in brackets.

The {\tt{}read} and {\tt{}write} procedures in this module are not type-bound but
generic procedures which handle scalar and array arguments.
\nwenddocs{}\nwbegincode{1}\moddef{\code{}particle{\char95}specifiers.f90\edoc{}}\endmoddef
\LA{}File header\RA{}

module particle_specifiers

\LA{}Use strings\RA{}
  use io_units
  use unit_tests
  use diagnostics

\LA{}Standard module head\RA{}

\LA{}Particle specifiers: public\RA{}

\LA{}Particle specifiers: types\RA{}

\LA{}Particle specifiers: interfaces\RA{}

contains

\LA{}Particle specifiers: procedures\RA{}

\LA{}Particle specifiers: tests\RA{}

end module particle_specifiers
\eatline
\nwendcode{}\nwbegindocs{2}\nwdocspar
\subsection{Base type}
This is an abstract type which can hold a single particle or an expression.
\nwenddocs{}\nwbegincode{3}\moddef{Particle specifiers: types}\endmoddef
  type, abstract :: prt_spec_expr_t
   contains
   \LA{}Particle specifiers: prt spec expr: TBP\RA{}
  end type prt_spec_expr_t
  
\eatline
\nwendcode{}\nwbegindocs{4}Output, as a string.
\nwenddocs{}\nwbegincode{5}\moddef{Particle specifiers: prt spec expr: TBP}\endmoddef
  procedure (prt_spec_expr_to_string), deferred :: to_string
\nwendcode{}\nwbegincode{6}\moddef{Particle specifiers: interfaces}\endmoddef
  abstract interface
     function prt_spec_expr_to_string (object) result (string)
       import
       class(prt_spec_expr_t), intent(in) :: object
       type(string_t) :: string
     end function prt_spec_expr_to_string
  end interface
  
\eatline
\nwendcode{}\nwbegindocs{7}Call an {\tt{}expand} method for all enclosed subexpressions (before handling
the current expression).
\nwenddocs{}\nwbegincode{8}\moddef{Particle specifiers: prt spec expr: TBP}\plusendmoddef
  procedure (prt_spec_expr_expand_sub), deferred :: expand_sub
\nwendcode{}\nwbegincode{9}\moddef{Particle specifiers: interfaces}\plusendmoddef
  abstract interface
     subroutine prt_spec_expr_expand_sub (object)
       import
       class(prt_spec_expr_t), intent(inout) :: object
     end subroutine prt_spec_expr_expand_sub
  end interface

\eatline
\nwendcode{}\nwbegindocs{10}\nwdocspar
\subsection{Wrapper type}
This wrapper can hold a particle expression of any kind.  We need it so we can
make variadic arrays.
\nwenddocs{}\nwbegincode{11}\moddef{Particle specifiers: public}\endmoddef
  public :: prt_expr_t
\nwendcode{}\nwbegincode{12}\moddef{Particle specifiers: types}\plusendmoddef
  type :: prt_expr_t
     class(prt_spec_expr_t), allocatable :: x
   contains
   \LA{}Particle specifiers: prt expr: TBP\RA{}
  end type prt_expr_t
  
\eatline
\nwendcode{}\nwbegindocs{13}Output as a string: delegate.
\nwenddocs{}\nwbegincode{14}\moddef{Particle specifiers: prt expr: TBP}\endmoddef
  procedure :: to_string => prt_expr_to_string
\nwendcode{}\nwbegincode{15}\moddef{Particle specifiers: procedures}\endmoddef
  recursive function prt_expr_to_string (object) result (string)
    class(prt_expr_t), intent(in) :: object
    type(string_t) :: string
    if (allocated (object%x)) then
       string = object%x%to_string ()
    else
       string = ""
    end if
  end function prt_expr_to_string
  
\eatline
\nwendcode{}\nwbegindocs{16}Allocate the expression as a particle specifier and copy the value.
\nwenddocs{}\nwbegincode{17}\moddef{Particle specifiers: prt expr: TBP}\plusendmoddef
  procedure :: init_spec => prt_expr_init_spec
\nwendcode{}\nwbegincode{18}\moddef{Particle specifiers: procedures}\plusendmoddef
  subroutine prt_expr_init_spec (object, spec)
    class(prt_expr_t), intent(out) :: object
    type(prt_spec_t), intent(in) :: spec
    allocate (prt_spec_t :: object%x)
    select type (x => object%x)
    type is (prt_spec_t)
       x = spec
    end select
  end subroutine prt_expr_init_spec
  
\eatline
\nwendcode{}\nwbegindocs{19}Allocate as a list/sum and allocate for a given length
\nwenddocs{}\nwbegincode{20}\moddef{Particle specifiers: prt expr: TBP}\plusendmoddef
  procedure :: init_list => prt_expr_init_list
  procedure :: init_sum => prt_expr_init_sum
\nwendcode{}\nwbegincode{21}\moddef{Particle specifiers: procedures}\plusendmoddef
  subroutine prt_expr_init_list (object, n)
    class(prt_expr_t), intent(out) :: object
    integer, intent(in) :: n
    allocate (prt_spec_list_t :: object%x)
    select type (x => object%x)
    type is (prt_spec_list_t)
       allocate (x%expr (n))
    end select
  end subroutine prt_expr_init_list
  
  subroutine prt_expr_init_sum (object, n)
    class(prt_expr_t), intent(out) :: object
    integer, intent(in) :: n
    allocate (prt_spec_sum_t :: object%x)
    select type (x => object%x)
    type is (prt_spec_sum_t)
       allocate (x%expr (n))
    end select
  end subroutine prt_expr_init_sum
  
\eatline
\eatline
\nwendcode{}\nwbegindocs{22}Return the number of terms.  This is unity, except if the expression is a
sum.
\nwenddocs{}\nwbegincode{23}\moddef{Particle specifiers: prt expr: TBP}\plusendmoddef
  procedure :: get_n_terms => prt_expr_get_n_terms
\nwendcode{}\nwbegincode{24}\moddef{Particle specifiers: procedures}\plusendmoddef
  function prt_expr_get_n_terms (object) result (n)
    class(prt_expr_t), intent(in) :: object
    integer :: n
    if (allocated (object%x)) then
       select type (x => object%x)
       type is (prt_spec_sum_t)
          n = size (x%expr)
       class default
          n = 1
       end select
    else
       n = 0
    end if
  end function prt_expr_get_n_terms
  
\eatline
\nwendcode{}\nwbegindocs{25}Transform one of the terms, as returned by the previous method, to an array
of particle specifiers.  The array has more than one entry if the selected
term is a list.  This makes sense only if the expression has been completely
expanded, so the list contains only atoms.
\nwenddocs{}\nwbegincode{26}\moddef{Particle specifiers: prt expr: TBP}\plusendmoddef
  procedure :: term_to_array => prt_expr_term_to_array
\nwendcode{}\nwbegincode{27}\moddef{Particle specifiers: procedures}\plusendmoddef
  recursive subroutine prt_expr_term_to_array (object, array, i)
    class(prt_expr_t), intent(in) :: object
    type(prt_spec_t), dimension(:), intent(inout), allocatable :: array
    integer, intent(in) :: i
    integer :: j
    if (allocated (array))  deallocate (array)
    select type (x => object%x)
    type is (prt_spec_t)
       allocate (array (1))
       array(1) = x
    type is (prt_spec_list_t)
       allocate (array (size (x%expr)))
       do j = 1, size (array)
          select type (y => x%expr(j)%x)
          type is (prt_spec_t)
             array(j) = y
          end select
       end do
    type is (prt_spec_sum_t)
       call x%expr(i)%term_to_array (array, 1)
    end select
  end subroutine prt_expr_term_to_array
  
\eatline
\nwendcode{}\nwbegindocs{28}\nwdocspar
\subsection{The atomic type}
The trivial case is a single particle, including optional decay and
polarization attributes.

\subsubsection{Definition}
The particle is unstable if the {\tt{}decay} array is allocated.  The
{\tt{}polarized} flag and decays may not be set simultaneously.
\nwenddocs{}\nwbegincode{29}\moddef{Particle specifiers: public}\plusendmoddef
  public :: prt_spec_t
\nwendcode{}\nwbegincode{30}\moddef{Particle specifiers: types}\plusendmoddef
  type, extends (prt_spec_expr_t) :: prt_spec_t
     private
     type(string_t) :: name
     logical :: polarized = .false.
     type(string_t), dimension(:), allocatable :: decay
   contains
   \LA{}Particle specifiers: prt spec: TBP\RA{}
  end type prt_spec_t
  
\eatline
\nwendcode{}\nwbegindocs{31}\nwdocspar
\subsubsection{I/O}
Output.  Old-style subroutines.
\nwenddocs{}\nwbegincode{32}\moddef{Particle specifiers: public}\plusendmoddef
  public :: prt_spec_write
\nwendcode{}\nwbegincode{33}\moddef{Particle specifiers: interfaces}\plusendmoddef
  interface prt_spec_write
     module procedure prt_spec_write1
     module procedure prt_spec_write2
  end interface prt_spec_write
\nwendcode{}\nwbegincode{34}\moddef{Particle specifiers: procedures}\plusendmoddef
  subroutine prt_spec_write1 (object, unit, advance)
    type(prt_spec_t), intent(in) :: object
    integer, intent(in), optional :: unit
    character(len=*), intent(in), optional :: advance
    character(3) :: adv
    integer :: u
    u = given_output_unit (unit)
    adv = "yes";  if (present (advance))  adv = advance
    write (u, "(A)", advance = adv)  char (object%to_string ())
  end subroutine prt_spec_write1
       
\eatline
\nwendcode{}\nwbegindocs{35}Write an array as a list of particle specifiers.
\nwenddocs{}\nwbegincode{36}\moddef{Particle specifiers: procedures}\plusendmoddef
  subroutine prt_spec_write2 (prt_spec, unit, advance)
    type(prt_spec_t), dimension(:), intent(in) :: prt_spec
    integer, intent(in), optional :: unit
    character(len=*), intent(in), optional :: advance
    character(3) :: adv
    integer :: u, i
    u = given_output_unit (unit)
    adv = "yes";  if (present (advance))  adv = advance
    do i = 1, size (prt_spec)
       if (i > 1)  write (u, "(A)", advance="no")  ", "
       call prt_spec_write (prt_spec(i), u, advance="no")
    end do
    write (u, "(A)", advance = adv)
  end subroutine prt_spec_write2
       
\eatline
\nwendcode{}\nwbegindocs{37}Read.  Input may be string or array of strings.
\nwenddocs{}\nwbegincode{38}\moddef{Particle specifiers: public}\plusendmoddef
  public :: prt_spec_read
\nwendcode{}\nwbegincode{39}\moddef{Particle specifiers: interfaces}\plusendmoddef
  interface prt_spec_read
     module procedure prt_spec_read1
     module procedure prt_spec_read2
  end interface prt_spec_read
\nwendcode{}\nwbegindocs{40}Read a single particle specifier
\nwenddocs{}\nwbegincode{41}\moddef{Particle specifiers: procedures}\plusendmoddef
  pure subroutine prt_spec_read1 (prt_spec, string)
    type(prt_spec_t), intent(out) :: prt_spec
    type(string_t), intent(in) :: string
    type(string_t) :: arg, buffer
    integer :: b1, b2, c, n, i
    b1 = scan (string, "(")
    b2 = scan (string, ")")
    if (b1 == 0) then
       prt_spec%name = trim (adjustl (string))
    else
       prt_spec%name = trim (adjustl (extract (string, 1, b1-1)))
       arg = trim (adjustl (extract (string, b1+1, b2-1)))
       if (arg == "*") then
          prt_spec%polarized = .true.
       else
          n = 0
          buffer = arg
          do
             if (verify (buffer, " ") == 0)  exit
             n = n + 1
             c = scan (buffer, "+")
             if (c == 0)  exit
             buffer = extract (buffer, c+1)
          end do
          allocate (prt_spec%decay (n))
          buffer = arg
          do i = 1, n
             c = scan (buffer, "+")
             if (c == 0)  c = len (buffer) + 1
             prt_spec%decay(i) = trim (adjustl (extract (buffer, 1, c-1)))
             buffer = extract (buffer, c+1)
          end do
       end if
    end if
  end subroutine prt_spec_read1

\eatline
\nwendcode{}\nwbegindocs{42}Read a particle specifier array, given as a single string.  The
array is allocated to the correct size.
\nwenddocs{}\nwbegincode{43}\moddef{Particle specifiers: procedures}\plusendmoddef
  pure subroutine prt_spec_read2 (prt_spec, string)
    type(prt_spec_t), dimension(:), intent(out), allocatable :: prt_spec
    type(string_t), intent(in) :: string
    type(string_t) :: buffer
    integer :: c, i, n
    n = 0
    buffer = string
    do
       n = n + 1
       c = scan (buffer, ",")
       if (c == 0)  exit
       buffer = extract (buffer, c+1)
    end do
    allocate (prt_spec (n))
    buffer = string
    do i = 1, size (prt_spec)
       c = scan (buffer, ",")
       if (c == 0)  c = len (buffer) + 1
       call prt_spec_read (prt_spec(i), &
            trim (adjustl (extract (buffer, 1, c-1))))
       buffer = extract (buffer, c+1)
    end do
  end subroutine prt_spec_read2
  
\eatline
\nwendcode{}\nwbegindocs{44}\nwdocspar
\subsubsection{Constructor}
Initialize a particle specifier.
\nwenddocs{}\nwbegincode{45}\moddef{Particle specifiers: public}\plusendmoddef
  public :: new_prt_spec
\nwendcode{}\nwbegincode{46}\moddef{Particle specifiers: interfaces}\plusendmoddef
  interface new_prt_spec
     module procedure new_prt_spec
     module procedure new_prt_spec_polarized
     module procedure new_prt_spec_unstable
  end interface new_prt_spec
\nwendcode{}\nwbegincode{47}\moddef{Particle specifiers: procedures}\plusendmoddef
  elemental function new_prt_spec (name) result (prt_spec)
    type(string_t), intent(in) :: name
    type(prt_spec_t) :: prt_spec
    prt_spec%name = name
  end function new_prt_spec

  elemental function new_prt_spec_polarized (name, polarized) result (prt_spec)
    type(string_t), intent(in) :: name
    logical, intent(in) :: polarized
    type(prt_spec_t) :: prt_spec
    prt_spec%name = name
    prt_spec%polarized = polarized
  end function new_prt_spec_polarized

  pure function new_prt_spec_unstable (name, decay) result (prt_spec)
    type(string_t), intent(in) :: name
    type(string_t), dimension(:), intent(in) :: decay
    type(prt_spec_t) :: prt_spec
    prt_spec%name = name
    allocate (prt_spec%decay (size (decay)))
    prt_spec%decay = decay
  end function new_prt_spec_unstable

\eatline
\nwendcode{}\nwbegindocs{48}\nwdocspar
\subsubsection{Access Methods}
Return the particle name without qualifiers
\nwenddocs{}\nwbegincode{49}\moddef{Particle specifiers: prt spec: TBP}\endmoddef
  procedure :: get_name => prt_spec_get_name
\nwendcode{}\nwbegincode{50}\moddef{Particle specifiers: procedures}\plusendmoddef
  elemental function prt_spec_get_name (prt_spec) result (name)
    class(prt_spec_t), intent(in) :: prt_spec
    type(string_t) :: name
    name = prt_spec%name
  end function prt_spec_get_name
  
\eatline
\nwendcode{}\nwbegindocs{51}Return the name with qualifiers
\nwenddocs{}\nwbegincode{52}\moddef{Particle specifiers: prt spec: TBP}\plusendmoddef
  procedure :: to_string => prt_spec_to_string
\nwendcode{}\nwbegincode{53}\moddef{Particle specifiers: procedures}\plusendmoddef
  function prt_spec_to_string (object) result (string)
    class(prt_spec_t), intent(in) :: object
    type(string_t) :: string
    integer :: i
    string = object%name
    if (allocated (object%decay)) then
       string = string // "("
       do i = 1, size (object%decay)
          if (i > 1)  string = string // " + "
          string = string // object%decay(i)
       end do
       string = string // ")"
    else if (object%polarized) then
       string = string // "(*)"
    end if
  end function prt_spec_to_string
  
\eatline
\nwendcode{}\nwbegindocs{54}Return the polarization flag
\nwenddocs{}\nwbegincode{55}\moddef{Particle specifiers: prt spec: TBP}\plusendmoddef
  procedure :: is_polarized => prt_spec_is_polarized
\nwendcode{}\nwbegincode{56}\moddef{Particle specifiers: procedures}\plusendmoddef
  elemental function prt_spec_is_polarized (prt_spec) result (flag)
    class(prt_spec_t), intent(in) :: prt_spec
    logical :: flag
    flag = prt_spec%polarized
  end function prt_spec_is_polarized
  
\eatline
\nwendcode{}\nwbegindocs{57}The particle is unstable if there is a decay array.
\nwenddocs{}\nwbegincode{58}\moddef{Particle specifiers: prt spec: TBP}\plusendmoddef
  procedure :: is_unstable => prt_spec_is_unstable
\nwendcode{}\nwbegincode{59}\moddef{Particle specifiers: procedures}\plusendmoddef
  elemental function prt_spec_is_unstable (prt_spec) result (flag)
    class(prt_spec_t), intent(in) :: prt_spec
    logical :: flag
    flag = allocated (prt_spec%decay)
  end function prt_spec_is_unstable
  
\eatline
\nwendcode{}\nwbegindocs{60}Return the number of decay channels
\nwenddocs{}\nwbegincode{61}\moddef{Particle specifiers: prt spec: TBP}\plusendmoddef
  procedure :: get_n_decays => prt_spec_get_n_decays
\nwendcode{}\nwbegincode{62}\moddef{Particle specifiers: procedures}\plusendmoddef
  elemental function prt_spec_get_n_decays (prt_spec) result (n)
    class(prt_spec_t), intent(in) :: prt_spec
    integer :: n
    if (allocated (prt_spec%decay)) then
       n = size (prt_spec%decay)
    else
       n = 0
    end if
  end function prt_spec_get_n_decays
  
\eatline
\nwendcode{}\nwbegindocs{63}Return the decay channels
\nwenddocs{}\nwbegincode{64}\moddef{Particle specifiers: prt spec: TBP}\plusendmoddef
  procedure :: get_decays => prt_spec_get_decays
\nwendcode{}\nwbegincode{65}\moddef{Particle specifiers: procedures}\plusendmoddef
  subroutine prt_spec_get_decays (prt_spec, decay)
    class(prt_spec_t), intent(in) :: prt_spec
    type(string_t), dimension(:), allocatable, intent(out) :: decay
    if (allocated (prt_spec%decay)) then
       allocate (decay (size (prt_spec%decay)))
       decay = prt_spec%decay
    else
       allocate (decay (0))
    end if
  end subroutine prt_spec_get_decays
  
\eatline
\nwendcode{}\nwbegindocs{66}\nwdocspar
\subsubsection{Miscellaneous}
There is nothing to expand here:
\nwenddocs{}\nwbegincode{67}\moddef{Particle specifiers: prt spec: TBP}\plusendmoddef
  procedure :: expand_sub => prt_spec_expand_sub
\nwendcode{}\nwbegincode{68}\moddef{Particle specifiers: procedures}\plusendmoddef
  subroutine prt_spec_expand_sub (object)
    class(prt_spec_t), intent(inout) :: object
  end subroutine prt_spec_expand_sub

\eatline
\nwendcode{}\nwbegindocs{69}\nwdocspar
\subsection{List}
A list of particle specifiers, indicating, e.g., the final state of a
process.
\nwenddocs{}\nwbegincode{70}\moddef{Particle specifiers: public}\plusendmoddef
  public :: prt_spec_list_t
\nwendcode{}\nwbegincode{71}\moddef{Particle specifiers: types}\plusendmoddef
  type, extends (prt_spec_expr_t) :: prt_spec_list_t
     type(prt_expr_t), dimension(:), allocatable :: expr
   contains
   \LA{}Particle specifiers: prt spec list: TBP\RA{}
  end type prt_spec_list_t
  
\eatline
\nwendcode{}\nwbegindocs{72}Output: Concatenate the components.  Insert brackets if the component is
also a list.  The components of the {\tt{}expr} array, if any, should all be
filled.
\nwenddocs{}\nwbegincode{73}\moddef{Particle specifiers: prt spec list: TBP}\endmoddef
  procedure :: to_string => prt_spec_list_to_string
\nwendcode{}\nwbegincode{74}\moddef{Particle specifiers: procedures}\plusendmoddef
  recursive function prt_spec_list_to_string (object) result (string)
    class(prt_spec_list_t), intent(in) :: object
    type(string_t) :: string
    integer :: i
    string = ""
    if (allocated (object%expr)) then
       do i = 1, size (object%expr)
          if (i > 1)  string = string // ", "
          select type (x => object%expr(i)%x)
          type is (prt_spec_list_t)
             string = string // "(" // x%to_string () // ")"
          class default
             string = string // x%to_string ()
          end select
       end do
    end if
  end function prt_spec_list_to_string

\eatline
\nwendcode{}\nwbegindocs{75}Flatten: if there is a subexpression which is also a list, include the
components as direct members of the current list.
\nwenddocs{}\nwbegincode{76}\moddef{Particle specifiers: prt spec list: TBP}\plusendmoddef
  procedure :: flatten => prt_spec_list_flatten
\nwendcode{}\nwbegincode{77}\moddef{Particle specifiers: procedures}\plusendmoddef
  subroutine prt_spec_list_flatten (object)
    class(prt_spec_list_t), intent(inout) :: object
    type(prt_expr_t), dimension(:), allocatable :: tmp_expr
    integer :: i, n_flat, i_flat
    n_flat = 0
    do i = 1, size (object%expr)
       select type (y => object%expr(i)%x)
       type is (prt_spec_list_t)
          n_flat = n_flat + size (y%expr)
       class default
          n_flat = n_flat + 1
       end select
    end do
    if (n_flat > size (object%expr)) then
       allocate (tmp_expr (n_flat))
       i_flat = 0
       do i = 1, size (object%expr)
          select type (y => object%expr(i)%x)
          type is (prt_spec_list_t)
             tmp_expr (i_flat + 1 : i_flat + size (y%expr)) = y%expr
             i_flat = i_flat + size (y%expr)
          class default
             tmp_expr (i_flat + 1) = object%expr(i)
             i_flat = i_flat + 1
          end select
       end do
    end if
    if (allocated (tmp_expr)) &
         call move_alloc (from = tmp_expr, to = object%expr)
  end subroutine prt_spec_list_flatten
    
\eatline
\nwendcode{}\nwbegindocs{78}Convert a list of sums into a sum of lists.  (Subexpressions which are not
sums are left untouched.)
\nwenddocs{}\nwbegincode{79}\moddef{Particle specifiers: procedures}\plusendmoddef
  subroutine distribute_prt_spec_list (object)
    class(prt_spec_expr_t), intent(inout), allocatable :: object
    class(prt_spec_expr_t), allocatable :: new_object
    integer, dimension(:), allocatable :: n, ii
    integer :: k, n_expr, n_terms, i_term
    select type (object)
    type is (prt_spec_list_t)
       n_expr = size (object%expr)
       allocate (n (n_expr), source = 1)
       allocate (ii (n_expr), source = 1)
       do k = 1, size (object%expr)
          select type (y => object%expr(k)%x)
          type is (prt_spec_sum_t)
             n(k) = size (y%expr)
          end select
       end do
       n_terms = product (n)
       if (n_terms > 1) then
          allocate (prt_spec_sum_t :: new_object)
          select type (new_object)
          type is (prt_spec_sum_t)
             allocate (new_object%expr (n_terms))
             do i_term = 1, n_terms
                allocate (prt_spec_list_t :: new_object%expr(i_term)%x)
                select type (x => new_object%expr(i_term)%x)
                type is (prt_spec_list_t)
                   allocate (x%expr (n_expr))
                   do k = 1, n_expr
                      select type (y => object%expr(k)%x)
                      type is (prt_spec_sum_t)
                         x%expr(k) = y%expr(ii(k))
                      class default
                         x%expr(k) = object%expr(k)
                      end select
                   end do
                end select
                INCR_INDEX: do k = n_expr, 1, -1
                   if (ii(k) < n(k)) then
                      ii(k) = ii(k) + 1
                      exit INCR_INDEX
                   else
                      ii(k) = 1
                   end if
                end do INCR_INDEX
             end do
          end select
       end if
    end select
    if (allocated (new_object)) call move_alloc (from = new_object, to = object)
  end subroutine distribute_prt_spec_list
    
\eatline
\nwendcode{}\nwbegindocs{80}Apply {\tt{}expand} to all components of the list.
\nwenddocs{}\nwbegincode{81}\moddef{Particle specifiers: prt spec list: TBP}\plusendmoddef
  procedure :: expand_sub => prt_spec_list_expand_sub
\nwendcode{}\nwbegincode{82}\moddef{Particle specifiers: procedures}\plusendmoddef
  recursive subroutine prt_spec_list_expand_sub (object)
    class(prt_spec_list_t), intent(inout) :: object
    integer :: i
    if (allocated (object%expr)) then
       do i = 1, size (object%expr)
          call object%expr(i)%expand ()
       end do
    end if
  end subroutine prt_spec_list_expand_sub

\eatline
\nwendcode{}\nwbegindocs{83}\nwdocspar
\subsection{Sum}
A sum of particle specifiers, indicating, e.g., a sum of final states.
\nwenddocs{}\nwbegincode{84}\moddef{Particle specifiers: public}\plusendmoddef
  public :: prt_spec_sum_t
\nwendcode{}\nwbegincode{85}\moddef{Particle specifiers: types}\plusendmoddef
  type, extends (prt_spec_expr_t) :: prt_spec_sum_t
     type(prt_expr_t), dimension(:), allocatable :: expr
   contains
   \LA{}Particle specifiers: prt spec sum: TBP\RA{}
  end type prt_spec_sum_t
  
\eatline
\nwendcode{}\nwbegindocs{86}Output: Concatenate the components.  Insert brackets if the component is
a list or also a sum.  The components of the {\tt{}expr} array, if any, should
all be filled.
\nwenddocs{}\nwbegincode{87}\moddef{Particle specifiers: prt spec sum: TBP}\endmoddef
  procedure :: to_string => prt_spec_sum_to_string
\nwendcode{}\nwbegincode{88}\moddef{Particle specifiers: procedures}\plusendmoddef
  recursive function prt_spec_sum_to_string (object) result (string)
    class(prt_spec_sum_t), intent(in) :: object
    type(string_t) :: string
    integer :: i
    string = ""
    if (allocated (object%expr)) then
       do i = 1, size (object%expr)
          if (i > 1)  string = string // " + "
          select type (x => object%expr(i)%x)
          type is (prt_spec_list_t)
             string = string // "(" // x%to_string () // ")"
          type is (prt_spec_sum_t)
             string = string // "(" // x%to_string () // ")"
          class default
             string = string // x%to_string ()
          end select
       end do
    end if
  end function prt_spec_sum_to_string

\eatline
\nwendcode{}\nwbegindocs{89}Flatten: if there is a subexpression which is also a sum, include the
components as direct members of the current sum.

This is identical to {\tt{}prt{\char95}spec{\char95}list{\char95}flatten} above, except for the type.
\nwenddocs{}\nwbegincode{90}\moddef{Particle specifiers: prt spec sum: TBP}\plusendmoddef
  procedure :: flatten => prt_spec_sum_flatten
\nwendcode{}\nwbegincode{91}\moddef{Particle specifiers: procedures}\plusendmoddef
  subroutine prt_spec_sum_flatten (object)
    class(prt_spec_sum_t), intent(inout) :: object
    type(prt_expr_t), dimension(:), allocatable :: tmp_expr
    integer :: i, n_flat, i_flat
    n_flat = 0
    do i = 1, size (object%expr)
       select type (y => object%expr(i)%x)
       type is (prt_spec_sum_t)
          n_flat = n_flat + size (y%expr)
          class default
          n_flat = n_flat + 1
       end select
    end do
    if (n_flat > size (object%expr)) then
       allocate (tmp_expr (n_flat))
       i_flat = 0
       do i = 1, size (object%expr)
          select type (y => object%expr(i)%x)
          type is (prt_spec_sum_t)
             tmp_expr (i_flat + 1 : i_flat + size (y%expr)) = y%expr
             i_flat = i_flat + size (y%expr)
          class default
             tmp_expr (i_flat + 1) = object%expr(i)
             i_flat = i_flat + 1
          end select
       end do
    end if
    if (allocated (tmp_expr)) &
         call move_alloc (from = tmp_expr, to = object%expr)
  end subroutine prt_spec_sum_flatten
    
\eatline
\nwendcode{}\nwbegindocs{92}Apply {\tt{}expand} to all terms in the sum.
\nwenddocs{}\nwbegincode{93}\moddef{Particle specifiers: prt spec sum: TBP}\plusendmoddef
  procedure :: expand_sub => prt_spec_sum_expand_sub
\nwendcode{}\nwbegincode{94}\moddef{Particle specifiers: procedures}\plusendmoddef
  recursive subroutine prt_spec_sum_expand_sub (object)
    class(prt_spec_sum_t), intent(inout) :: object
    integer :: i
    if (allocated (object%expr)) then
       do i = 1, size (object%expr)
          call object%expr(i)%expand ()
       end do
    end if
  end subroutine prt_spec_sum_expand_sub

\eatline
\nwendcode{}\nwbegindocs{95}\nwdocspar
\subsection{Expression Expansion}
The {\tt{}expand} method transforms each particle specifier expression into a sum
of lists, according to the rules
\begin{align}
  a, (b, c) &\to a, b, c
\\
  a + (b + c) &\to a + b + c
\\
  a, b + c &\to (a, b) + (a, c)
\end{align}
Note that the precedence of comma and plus are opposite to this expansion, so
the parentheses in the final expression are necessary.

We assume that subexpressions are filled, i.e., arrays are allocated.
\nwenddocs{}\nwbegincode{96}\moddef{Particle specifiers: prt expr: TBP}\plusendmoddef
  procedure :: expand => prt_expr_expand
\nwendcode{}\nwbegincode{97}\moddef{Particle specifiers: procedures}\plusendmoddef
  recursive subroutine prt_expr_expand (expr)
    class(prt_expr_t), intent(inout) :: expr
    if (allocated (expr%x)) then
       call distribute_prt_spec_list (expr%x)
       call expr%x%expand_sub ()
       select type (x => expr%x)
       type is (prt_spec_list_t)
          call x%flatten ()
       type is (prt_spec_sum_t)
          call x%flatten ()
       end select
    end if
  end subroutine prt_expr_expand

\eatline
\nwendcode{}\nwbegindocs{98}\nwdocspar
\subsection{Test}
This is the master for calling self-test procedures.
\nwenddocs{}\nwbegincode{99}\moddef{Particle specifiers: public}\plusendmoddef
  public :: particle_specifiers_test
\nwendcode{}\nwbegincode{100}\moddef{Particle specifiers: tests}\endmoddef
  subroutine particle_specifiers_test (u, results)
    integer, intent(in) :: u
    type(test_results_t), intent(inout) :: results
  \LA{}Particle specifiers: execute tests\RA{}
  end subroutine particle_specifiers_test

\eatline
\nwendcode{}\nwbegindocs{101}\nwdocspar
\subsubsection{Particle specifier array}
Define, read and write an array of particle specifiers.
\nwenddocs{}\nwbegincode{102}\moddef{Particle specifiers: execute tests}\endmoddef
  call test (particle_specifiers_1, "particle_specifiers_1", &
       "Handle particle specifiers", &
       u, results)
\nwendcode{}\nwbegincode{103}\moddef{Particle specifiers: tests}\plusendmoddef
  subroutine particle_specifiers_1 (u)
    integer, intent(in) :: u
    type(prt_spec_t), dimension(:), allocatable :: prt_spec
    type(string_t), dimension(:), allocatable :: decay
    type(string_t), dimension(0) :: no_decay
    integer :: i, j

    write (u, "(A)")  "* Test output: particle_specifiers_1"
    write (u, "(A)")  "*   Purpose: Read and write a particle specifier array"
    write (u, "(A)")

    allocate (prt_spec (5))
    prt_spec = [ &
         new_prt_spec (var_str ("a")), &
         new_prt_spec (var_str ("b"), .true.), &
         new_prt_spec (var_str ("c"), [var_str ("dec1")]), &
         new_prt_spec (var_str ("d"), [var_str ("dec1"), var_str ("dec2")]), &
         new_prt_spec (var_str ("e"), no_decay) &
         ]
    do i = 1, size (prt_spec)
       write (u, "(A)")  char (prt_spec(i)%to_string ())
    end do
    write (u, "(A)")

    call prt_spec_read (prt_spec, &
         var_str (" a, b( *), c( dec1), d (dec1 + dec2 ), e()"))
    call prt_spec_write (prt_spec, u)
    
    do i = 1, size (prt_spec)
       write (u, "(A)")
       write (u, "(A,A)")  char (prt_spec(i)%get_name ()), ":"
       write (u, "(A,L1)")  "polarized = ", prt_spec(i)%is_polarized ()
       write (u, "(A,L1)")  "unstable  = ", prt_spec(i)%is_unstable ()
       write (u, "(A,I0)")  "n_decays  = ", prt_spec(i)%get_n_decays ()
       call prt_spec(i)%get_decays (decay)
       write (u, "(A)", advance="no") "decays    ="
       do j = 1, size (decay)
          write (u, "(1x,A)", advance="no") char (decay(j))
       end do
       write (u, "(A)")
    end do

    write (u, "(A)")
    write (u, "(A)")  "* Test output end: particle_specifiers_1"
  end subroutine particle_specifiers_1
  
\eatline
\nwendcode{}\nwbegindocs{104}\nwdocspar
\subsubsection{Particle specifier expressions}
Nested expressions (only basic particles, no decay specs).
\nwenddocs{}\nwbegincode{105}\moddef{Particle specifiers: execute tests}\plusendmoddef
  call test (particle_specifiers_2, "particle_specifiers_2", &
       "Particle specifier expressions", &
       u, results)
\nwendcode{}\nwbegincode{106}\moddef{Particle specifiers: tests}\plusendmoddef
  subroutine particle_specifiers_2 (u)
    integer, intent(in) :: u
    type(prt_spec_t) :: a, b, c, d, e, f
    type(prt_expr_t) :: pe1, pe2, pe3
    type(prt_expr_t) :: pe4, pe5, pe6, pe7, pe8, pe9
    integer :: i
    type(prt_spec_t), dimension(:), allocatable :: pa

    write (u, "(A)")  "* Test output: particle_specifiers_2"
    write (u, "(A)")  "*   Purpose: Create and display particle expressions"
    write (u, "(A)")

    write (u, "(A)")  "* Basic expressions"
    write (u, *)

    a = new_prt_spec (var_str ("a"))
    b = new_prt_spec (var_str ("b"))
    c = new_prt_spec (var_str ("c"))
    d = new_prt_spec (var_str ("d"))
    e = new_prt_spec (var_str ("e"))
    f = new_prt_spec (var_str ("f"))
    
    call pe1%init_spec (a)
    write (u, "(A)")  char (pe1%to_string ())

    call pe2%init_sum (2)
    select type (x => pe2%x)
    type is (prt_spec_sum_t)
       call x%expr(1)%init_spec (a)
       call x%expr(2)%init_spec (b)
    end select
    write (u, "(A)")  char (pe2%to_string ())

    call pe3%init_list (2)
    select type (x => pe3%x)
    type is (prt_spec_list_t)
       call x%expr(1)%init_spec (a)
       call x%expr(2)%init_spec (b)
    end select
    write (u, "(A)")  char (pe3%to_string ())

    write (u, *)
    write (u, "(A)")  "* Nested expressions"
    write (u, *)
    
    call pe4%init_list (2)
    select type (x => pe4%x)
    type is (prt_spec_list_t)
       call x%expr(1)%init_sum (2)
       select type (y => x%expr(1)%x)
       type is (prt_spec_sum_t)
          call y%expr(1)%init_spec (a)
          call y%expr(2)%init_spec (b)
       end select
       call x%expr(2)%init_spec (c)
    end select
    write (u, "(A)")  char (pe4%to_string ())
          
    call pe5%init_list (2)
    select type (x => pe5%x)
    type is (prt_spec_list_t)
       call x%expr(1)%init_list (2)
       select type (y => x%expr(1)%x)
       type is (prt_spec_list_t)
          call y%expr(1)%init_spec (a)
          call y%expr(2)%init_spec (b)
       end select
       call x%expr(2)%init_spec (c)
    end select
    write (u, "(A)")  char (pe5%to_string ())
          
    call pe6%init_sum (2)
    select type (x => pe6%x)
    type is (prt_spec_sum_t)
       call x%expr(1)%init_spec (a)
       call x%expr(2)%init_sum (2)
       select type (y => x%expr(2)%x)
       type is (prt_spec_sum_t)
          call y%expr(1)%init_spec (b)
          call y%expr(2)%init_spec (c)
       end select
    end select
    write (u, "(A)")  char (pe6%to_string ())
          
    call pe7%init_list (2)
    select type (x => pe7%x)
    type is (prt_spec_list_t)
       call x%expr(1)%init_sum (2)
       select type (y => x%expr(1)%x)
       type is (prt_spec_sum_t)
          call y%expr(1)%init_spec (a)
          call y%expr(2)%init_list (2)
          select type (z => y%expr(2)%x)
          type is (prt_spec_list_t)
             call z%expr(1)%init_spec (b)
             call z%expr(2)%init_spec (c)
          end select
       end select
       call x%expr(2)%init_spec (d)
    end select
    write (u, "(A)")  char (pe7%to_string ())
          
    call pe8%init_sum (2)
    select type (x => pe8%x)
    type is (prt_spec_sum_t)
       call x%expr(1)%init_list (2)
       select type (y => x%expr(1)%x)
       type is (prt_spec_list_t)
          call y%expr(1)%init_spec (a)
          call y%expr(2)%init_spec (b)
       end select
       call x%expr(2)%init_list (2)
       select type (y => x%expr(2)%x)
       type is (prt_spec_list_t)
          call y%expr(1)%init_spec (c)
          call y%expr(2)%init_spec (d)
       end select
    end select
    write (u, "(A)")  char (pe8%to_string ())

    call pe9%init_list (3)
    select type (x => pe9%x)
    type is (prt_spec_list_t)
       call x%expr(1)%init_sum (2)
       select type (y => x%expr(1)%x)
       type is (prt_spec_sum_t)
          call y%expr(1)%init_spec (a)
          call y%expr(2)%init_spec (b)
       end select
       call x%expr(2)%init_spec (c)
       call x%expr(3)%init_sum (3)
       select type (y => x%expr(3)%x)
       type is (prt_spec_sum_t)
          call y%expr(1)%init_spec (d)
          call y%expr(2)%init_spec (e)
          call y%expr(3)%init_spec (f)
       end select
    end select
    write (u, "(A)")  char (pe9%to_string ())
          
    write (u, *)
    write (u, "(A)")  "* Expand as sum"
    write (u, *)
    
    call pe1%expand ()
    write (u, "(A)")  char (pe1%to_string ())

    call pe4%expand ()
    write (u, "(A)")  char (pe4%to_string ())

    call pe5%expand ()
    write (u, "(A)")  char (pe5%to_string ())

    call pe6%expand ()
    write (u, "(A)")  char (pe6%to_string ())

    call pe7%expand ()
    write (u, "(A)")  char (pe7%to_string ())
    
    call pe8%expand ()
    write (u, "(A)")  char (pe8%to_string ())
    
    call pe9%expand ()
    write (u, "(A)")  char (pe9%to_string ())
    
    write (u, *)
    write (u, "(A)")  "* Transform to arrays:"

    write (u, "(A)")  "* Atomic specifier"
    do i = 1, pe1%get_n_terms ()
       call pe1%term_to_array (pa, i)
       call prt_spec_write (pa, u)
    end do
    
    write (u, *)
    write (u, "(A)")  "* List"
    do i = 1, pe5%get_n_terms ()
       call pe5%term_to_array (pa, i)
       call prt_spec_write (pa, u)
    end do

    write (u, *)
    write (u, "(A)")  "* Sum of atoms"
    do i = 1, pe6%get_n_terms ()
       call pe6%term_to_array (pa, i)
       call prt_spec_write (pa, u)
    end do

    write (u, *)
    write (u, "(A)")  "* Sum of lists"
    do i = 1, pe9%get_n_terms ()
       call pe9%term_to_array (pa, i)
       call prt_spec_write (pa, u)
    end do

    write (u, "(A)")
    write (u, "(A)")  "* Test output end: particle_specifiers_2"
  end subroutine particle_specifiers_2
  
\eatline
\nwendcode{}\nwbegindocs{107}\nwdocspar
%%%%%%%%%%%%%%%%%%%%%%%%%%%%%%%%%%%%%%%%%%%%%%%%%%%%%%%%%%%%%%%%%%%%%%%%
\section{PDG arrays}

For defining aliases, we introduce a special type which holds a set of
(integer) PDG codes. 
\nwenddocs{}\nwbegincode{108}\moddef{\code{}pdg{\char95}arrays.f90\edoc{}}\endmoddef
\LA{}File header\RA{}

module pdg_arrays

  use io_units
  use unit_tests
  use sorting
  use physics_defs, only: UNDEFINED

\LA{}Standard module head\RA{}

\LA{}PDG arrays: public\RA{}

\LA{}PDG arrays: types\RA{}

\LA{}PDG arrays: interfaces\RA{}

contains

\LA{}PDG arrays: procedures\RA{}

\LA{}PDG arrays: tests\RA{}

end module pdg_arrays
\eatline
\nwendcode{}\nwbegindocs{109}\nwdocspar
\subsection{Type definition}
Using an allocatable array eliminates the need for initializer and/or
finalizer.
\nwenddocs{}\nwbegincode{110}\moddef{PDG arrays: public}\endmoddef
  public :: pdg_array_t
\nwendcode{}\nwbegincode{111}\moddef{PDG arrays: types}\endmoddef
  type :: pdg_array_t
     private
     integer, dimension(:), allocatable :: pdg
   contains
   \LA{}PDG arrays: pdg array: TBP\RA{}
  end type pdg_array_t

\eatline
\nwendcode{}\nwbegindocs{112}Output
\nwenddocs{}\nwbegincode{113}\moddef{PDG arrays: public}\plusendmoddef
  public :: pdg_array_write
\nwendcode{}\nwbegincode{114}\moddef{PDG arrays: pdg array: TBP}\endmoddef
  procedure :: write => pdg_array_write
\nwendcode{}\nwbegincode{115}\moddef{PDG arrays: procedures}\endmoddef
  subroutine pdg_array_write (aval, unit)
    class(pdg_array_t), intent(in) :: aval
    integer, intent(in), optional :: unit
    integer :: u, i
    u = given_output_unit (unit);  if (u < 0)  return
    write (u, "(A)", advance="no")  "PDG("
    if (allocated (aval%pdg)) then
       do i = 1, size (aval%pdg)
          if (i > 1)  write (u, "(A)", advance="no")  ", "
          write (u, "(I0)", advance="no")  aval%pdg(i)
       end do
    end if
    write (u, "(A)", advance="no")  ")"
  end subroutine pdg_array_write

\eatline
\nwendcode{}\nwbegindocs{116}\nwdocspar
\subsection{Basic operations}
Assignment.  We define assignment from and to an integer array.
Note that the integer array, if it is the l.h.s., must be declared
allocatable by the caller.
\nwenddocs{}\nwbegincode{117}\moddef{PDG arrays: public}\plusendmoddef
  public :: assignment(=)
\nwendcode{}\nwbegincode{118}\moddef{PDG arrays: interfaces}\endmoddef
  interface assignment(=)
     module procedure pdg_array_from_int_array
     module procedure pdg_array_from_int
     module procedure int_array_from_pdg_array
  end interface

\nwendcode{}\nwbegincode{119}\moddef{PDG arrays: procedures}\plusendmoddef
  subroutine pdg_array_from_int_array (aval, iarray)
    type(pdg_array_t), intent(out) :: aval
    integer, dimension(:), intent(in) :: iarray
    allocate (aval%pdg (size (iarray)))
    aval%pdg = iarray
  end subroutine pdg_array_from_int_array

  elemental subroutine pdg_array_from_int (aval, int)
    type(pdg_array_t), intent(out) :: aval
    integer, intent(in) :: int
    allocate (aval%pdg (1))
    aval%pdg = int
  end subroutine pdg_array_from_int

  subroutine int_array_from_pdg_array (iarray, aval)
    integer, dimension(:), allocatable, intent(out) :: iarray
    type(pdg_array_t), intent(in) :: aval
    if (allocated (aval%pdg)) then
       allocate (iarray (size (aval%pdg)))
       iarray = aval%pdg
    else
       allocate (iarray (0))
    end if
  end subroutine int_array_from_pdg_array

\eatline
\nwendcode{}\nwbegindocs{120}Allocate space for a PDG array
\nwenddocs{}\nwbegincode{121}\moddef{PDG arrays: public}\plusendmoddef
  public :: pdg_array_init
\nwendcode{}\nwbegincode{122}\moddef{PDG arrays: procedures}\plusendmoddef
   subroutine pdg_array_init (aval, n_elements)
     type(pdg_array_t), intent(inout) :: aval
     integer, intent(in) :: n_elements
     allocate(aval%pdg(n_elements))
   end subroutine pdg_array_init

\eatline
\nwendcode{}\nwbegindocs{123}Deallocate a previously allocated pdg array
\nwenddocs{}\nwbegincode{124}\moddef{PDG arrays: public}\plusendmoddef
  public :: pdg_array_delete
\nwendcode{}\nwbegincode{125}\moddef{PDG arrays: procedures}\plusendmoddef
  subroutine pdg_array_delete (aval)
    type(pdg_array_t), intent(inout) :: aval
    if (allocated (aval%pdg)) deallocate (aval%pdg)
  end subroutine pdg_array_delete

\eatline
\nwendcode{}\nwbegindocs{126}Merge two pdg arrays, i.e. append a particle string to another leaving out doublettes
\nwenddocs{}\nwbegincode{127}\moddef{PDG arrays: public}\plusendmoddef
  public :: pdg_array_merge
\nwendcode{}\nwbegincode{128}\moddef{PDG arrays: procedures}\plusendmoddef
  subroutine pdg_array_merge (aval1, aval2)
    type(pdg_array_t), intent(inout) :: aval1
    type(pdg_array_t), intent(in) :: aval2
    type(pdg_array_t) :: aval
    if (allocated (aval1%pdg) .and. allocated (aval2%pdg)) then
      if (.not. any (aval1%pdg == aval2%pdg)) aval = aval1 // aval2
    else if (allocated (aval1%pdg)) then
      aval = aval1
    else if (allocated (aval2%pdg)) then
      aval = aval2
    end if
    call pdg_array_delete (aval1)
    aval1 = aval%pdg
  end subroutine pdg_array_merge

\eatline
\nwendcode{}\nwbegindocs{129}Length of the array.
\nwenddocs{}\nwbegincode{130}\moddef{PDG arrays: public}\plusendmoddef
  public :: pdg_array_get_length
\nwendcode{}\nwbegincode{131}\moddef{PDG arrays: pdg array: TBP}\plusendmoddef
  procedure :: get_length => pdg_array_get_length
\nwendcode{}\nwbegincode{132}\moddef{PDG arrays: procedures}\plusendmoddef
  elemental function pdg_array_get_length (aval) result (n)
    class(pdg_array_t), intent(in) :: aval
    integer :: n
    if (allocated (aval%pdg)) then
       n = size (aval%pdg)
    else
       n = 0
    end if
  end function pdg_array_get_length

\eatline
\nwendcode{}\nwbegindocs{133}Return the element with index i.
\nwenddocs{}\nwbegincode{134}\moddef{PDG arrays: public}\plusendmoddef
  public :: pdg_array_get
\nwendcode{}\nwbegincode{135}\moddef{PDG arrays: pdg array: TBP}\plusendmoddef
  procedure :: get => pdg_array_get
\nwendcode{}\nwbegincode{136}\moddef{PDG arrays: procedures}\plusendmoddef
  elemental function pdg_array_get (aval, i) result (pdg)
    class(pdg_array_t), intent(in) :: aval
    integer, intent(in), optional :: i
    integer :: pdg
    if (present (i)) then
       pdg = aval%pdg(i)
    else 
       pdg = aval%pdg(1)
    end if
  end function pdg_array_get

\eatline
\nwendcode{}\nwbegindocs{137}Explicitly set the element with index i.
\nwenddocs{}\nwbegincode{138}\moddef{PDG arrays: pdg array: TBP}\plusendmoddef
  procedure :: set => pdg_array_set
\nwendcode{}\nwbegincode{139}\moddef{PDG arrays: procedures}\plusendmoddef
  subroutine pdg_array_set (aval, i, pdg)
    class(pdg_array_t), intent(inout) :: aval
    integer, intent(in) :: i
    integer, intent(in) :: pdg
    aval%pdg(i) = pdg
  end subroutine pdg_array_set

\eatline
\nwendcode{}\nwbegindocs{140}Replace element with index {\tt{}i} by a new array of elements.
\nwenddocs{}\nwbegincode{141}\moddef{PDG arrays: public}\plusendmoddef
  public :: pdg_array_replace
\nwendcode{}\nwbegincode{142}\moddef{PDG arrays: pdg array: TBP}\plusendmoddef
  procedure :: replace => pdg_array_replace
\nwendcode{}\nwbegincode{143}\moddef{PDG arrays: procedures}\plusendmoddef
  function pdg_array_replace (aval, i, pdg_new) result (aval_new)
    class(pdg_array_t), intent(in) :: aval
    integer, intent(in) :: i
    integer, dimension(:), intent(in) :: pdg_new
    type(pdg_array_t) :: aval_new
    integer :: n, l
    n = size (aval%pdg)
    l = size (pdg_new)
    allocate (aval_new%pdg (n + l - 1))
    aval_new%pdg(:i-1) = aval%pdg(:i-1)
    aval_new%pdg(i:i+l-1) = pdg_new
    aval_new%pdg(i+l:) = aval%pdg(i+1:)
  end function pdg_array_replace
    
\eatline
\nwendcode{}\nwbegindocs{144}Concatenate two PDG arrays
\nwenddocs{}\nwbegincode{145}\moddef{PDG arrays: public}\plusendmoddef
  public :: operator(//)
\nwendcode{}\nwbegincode{146}\moddef{PDG arrays: interfaces}\plusendmoddef
  interface operator(//)
     module procedure concat_pdg_arrays
  end interface

\nwendcode{}\nwbegincode{147}\moddef{PDG arrays: procedures}\plusendmoddef
  function concat_pdg_arrays (aval1, aval2) result (aval)
    type(pdg_array_t) :: aval
    type(pdg_array_t), intent(in) :: aval1, aval2
    integer :: n1, n2
    if (allocated (aval1%pdg) .and. allocated (aval2%pdg)) then
       n1 = size (aval1%pdg)
       n2 = size (aval2%pdg)
       allocate (aval%pdg (n1 + n2))
       aval%pdg(:n1) = aval1%pdg
       aval%pdg(n1+1:) = aval2%pdg
    else if (allocated (aval1%pdg)) then
       aval = aval1
    else if (allocated (aval2%pdg)) then
       aval = aval2
    end if
  end function concat_pdg_arrays

\eatline
\nwendcode{}\nwbegindocs{148}\nwdocspar
\subsection{Matching}
A PDG array matches a given PDG code if the code is present within the
array.  If either one is zero (UNDEFINED), the match also succeeds.
\nwenddocs{}\nwbegincode{149}\moddef{PDG arrays: public}\plusendmoddef
  public :: operator(.match.)
\nwendcode{}\nwbegincode{150}\moddef{PDG arrays: interfaces}\plusendmoddef
  interface operator(.match.)
     module procedure pdg_array_match_integer
     module procedure pdg_array_match_pdg_array
  end interface

\eatline
\nwendcode{}\nwbegindocs{151}Match a single code against the array.
\nwenddocs{}\nwbegincode{152}\moddef{PDG arrays: procedures}\plusendmoddef
  elemental function pdg_array_match_integer (aval, pdg) result (flag)
    logical :: flag
    type(pdg_array_t), intent(in) :: aval
    integer, intent(in) :: pdg
    if (allocated (aval%pdg)) then
       flag = pdg == UNDEFINED &
            .or. any (aval%pdg == UNDEFINED) &
            .or. any (aval%pdg == pdg)
    else
       flag = .false.
    end if
  end function pdg_array_match_integer

\eatline
\nwendcode{}\nwbegindocs{153}Check if the pdg-number corresponds to a quark
\nwenddocs{}\nwbegincode{154}\moddef{PDG arrays: public}\plusendmoddef
  public :: is_quark
\nwendcode{}\nwbegincode{155}\moddef{PDG arrays: procedures}\plusendmoddef
  function is_quark (pdg_nr) result(res)
      integer, intent(in) :: pdg_nr
      logical :: res
      if (pdg_nr >= 1 .and. pdg_nr <= 6) then 
        res = .true.
      else
        res = .false.
      end if
   end function is_quark

\eatline
\nwendcode{}\nwbegindocs{156}Check if pdg-number corresponds to a gluon
\nwenddocs{}\nwbegincode{157}\moddef{PDG arrays: public}\plusendmoddef
  public :: is_gluon
\nwendcode{}\nwbegincode{158}\moddef{PDG arrays: procedures}\plusendmoddef
   function is_gluon (pdg_nr) result(res)
     integer, intent(in) :: pdg_nr
     logical :: res
     if (pdg_nr == 21) then
       res = .true.
     else
       res = .false.
     end if
   end function is_gluon

\eatline
\nwendcode{}\nwbegindocs{159}% Check if the pdg-number corresponds to a lepton
\nwenddocs{}\nwbegincode{160}\moddef{PDG arrays: public}\plusendmoddef
  public :: is_lepton
\nwendcode{}\nwbegincode{161}\moddef{PDG arrays: procedures}\plusendmoddef
  function is_lepton (pdg_nr) result(res)
    integer, intent(in) :: pdg_nr
    logical :: res
    if (pdg_nr >= 11 .and. pdg_nr <= 16) then
      res = .true.
    else
      res = .false.
    end if
  end function is_lepton

\eatline
\nwendcode{}\nwbegindocs{162}Check if the pdg-number corresponds to a massless vector boson
\nwenddocs{}\nwbegincode{163}\moddef{PDG arrays: public}\plusendmoddef
  public :: is_massless_vector
\nwendcode{}\nwbegincode{164}\moddef{PDG arrays: procedures}\plusendmoddef
  function is_massless_vector (pdg_nr) result (res)
    integer, intent(in) :: pdg_nr
    logical :: res
    if (pdg_nr == 21 .or. pdg_nr == 22) then
      res = .true.
    else
      res = .false.
    end if
  end function is_massless_vector

\eatline
\nwendcode{}\nwbegindocs{165}Check if pdg-number corresponds to a massive vector boson
\nwenddocs{}\nwbegincode{166}\moddef{PDG arrays: public}\plusendmoddef
  public :: is_massive_vector
\nwendcode{}\nwbegincode{167}\moddef{PDG arrays: procedures}\plusendmoddef
  function is_massive_vector (pdg_nr) result (res)
    integer, intent(in) :: pdg_nr
    logical :: res
    if (pdg_nr == 23 .or. pdg_nr == 24) then
      res = .true.
    else
      res = .false.
    end if
  end function is_massive_vector

\eatline
\nwendcode{}\nwbegindocs{168}Check if particle is strongly interacting
\nwenddocs{}\nwbegincode{169}\moddef{PDG arrays: public}\plusendmoddef
  public :: is_qcd_particle
\nwendcode{}\nwbegincode{170}\moddef{PDG arrays: procedures}\plusendmoddef
  function is_qcd_particle (pdg_nr) result (res)
    integer, intent(in) :: pdg_nr
    logical :: res
    res = .false.
    if (is_quark (abs (pdg_nr)) .or. is_gluon (pdg_nr)) &
      res = .true.
  end function is_qcd_particle

\eatline
\nwendcode{}\nwbegindocs{171}Match two arrays.  Succeeds if any pair of entries matches.
\nwenddocs{}\nwbegincode{172}\moddef{PDG arrays: procedures}\plusendmoddef
  function pdg_array_match_pdg_array (aval1, aval2) result (flag)
    logical :: flag
    type(pdg_array_t), intent(in) :: aval1, aval2
    if (allocated (aval1%pdg) .and. allocated (aval2%pdg)) then
       flag = any (aval1 .match. aval2%pdg)
    else
       flag = .false.
    end if
  end function pdg_array_match_pdg_array

\eatline
\nwendcode{}\nwbegindocs{173}Comparison.  Here, we take the PDG arrays as-is, assuming that they
are sorted.

The ordering is a bit odd: first, we look only at the absolute values
of the PDG codes.  If they all match, the particle comes before the
antiparticle, scanning from left to right.
\nwenddocs{}\nwbegincode{174}\moddef{PDG arrays: public}\plusendmoddef
  public :: operator(<)
  public :: operator(>)
  public :: operator(<=)
  public :: operator(>=)
  public :: operator(==)
  public :: operator(/=)
\nwendcode{}\nwbegincode{175}\moddef{PDG arrays: interfaces}\plusendmoddef
  interface operator(<)
     module procedure pdg_array_lt
  end interface
  interface operator(>)
     module procedure pdg_array_gt
  end interface
  interface operator(<=)
     module procedure pdg_array_le
  end interface
  interface operator(>=)
     module procedure pdg_array_ge
  end interface
  interface operator(==)
     module procedure pdg_array_eq
  end interface
  interface operator(/=)
     module procedure pdg_array_ne
  end interface

\nwendcode{}\nwbegincode{176}\moddef{PDG arrays: procedures}\plusendmoddef
  elemental function pdg_array_lt (aval1, aval2) result (flag)
    type(pdg_array_t), intent(in) :: aval1, aval2
    logical :: flag
    integer :: i
    if (size (aval1%pdg) /= size (aval2%pdg)) then
       flag = size (aval1%pdg) < size (aval2%pdg)
    else
       do i = 1, size (aval1%pdg)
          if (abs (aval1%pdg(i)) /= abs (aval2%pdg(i))) then
             flag = abs (aval1%pdg(i)) < abs (aval2%pdg(i))
             return
          end if
       end do
       do i = 1, size (aval1%pdg)
          if (aval1%pdg(i) /= aval2%pdg(i)) then
             flag = aval1%pdg(i) > aval2%pdg(i)
             return
          end if
       end do
       flag = .false.
    end if
  end function pdg_array_lt

  elemental function pdg_array_gt (aval1, aval2) result (flag)
    type(pdg_array_t), intent(in) :: aval1, aval2
    logical :: flag
    flag = .not. (aval1 < aval2 .or. aval1 == aval2)
  end function pdg_array_gt

  elemental function pdg_array_le (aval1, aval2) result (flag)
    type(pdg_array_t), intent(in) :: aval1, aval2
    logical :: flag
    flag = aval1 < aval2 .or. aval1 == aval2
  end function pdg_array_le

  elemental function pdg_array_ge (aval1, aval2) result (flag)
    type(pdg_array_t), intent(in) :: aval1, aval2
    logical :: flag
    flag = .not. (aval1 < aval2)
  end function pdg_array_ge

  elemental function pdg_array_eq (aval1, aval2) result (flag)
    type(pdg_array_t), intent(in) :: aval1, aval2
    logical :: flag
    if (size (aval1%pdg) /= size (aval2%pdg)) then
       flag = .false.
    else
       flag = all (aval1%pdg == aval2%pdg)
    end if
  end function pdg_array_eq

  elemental function pdg_array_ne (aval1, aval2) result (flag)
    type(pdg_array_t), intent(in) :: aval1, aval2
    logical :: flag
    flag = .not. (aval1 == aval2)
  end function pdg_array_ne

\nwendcode{}\nwbegindocs{177}Equivalence.  Two PDG arrays are equivalent if either one contains
{\tt{}UNDEFINED} or if each element of array 1 is present in array 2, and
vice versa.
\nwenddocs{}\nwbegincode{178}\moddef{PDG arrays: public}\plusendmoddef
  public :: operator(.eqv.)
  public :: operator(.neqv.)
\nwendcode{}\nwbegincode{179}\moddef{PDG arrays: interfaces}\plusendmoddef
  interface operator(.eqv.)
     module procedure pdg_array_equivalent
  end interface
  interface operator(.neqv.)
     module procedure pdg_array_inequivalent
  end interface

\nwendcode{}\nwbegincode{180}\moddef{PDG arrays: procedures}\plusendmoddef
  elemental function pdg_array_equivalent (aval1, aval2) result (eq)
    logical :: eq
    type(pdg_array_t), intent(in) :: aval1, aval2
    logical, dimension(:), allocatable :: match1, match2
    integer :: i
    if (allocated (aval1%pdg) .and. allocated (aval2%pdg)) then
       eq = any (aval1%pdg == UNDEFINED) &
            .or. any (aval2%pdg == UNDEFINED)
       if (.not. eq) then
          allocate (match1 (size (aval1%pdg)))
          allocate (match2 (size (aval2%pdg)))
          match1 = .false.
          match2 = .false.
          do i = 1, size (aval1%pdg)
             match2 = match2 .or. aval1%pdg(i) == aval2%pdg
          end do
          do i = 1, size (aval2%pdg)
             match1 = match1 .or. aval2%pdg(i) == aval1%pdg
          end do
          eq = all (match1) .and. all (match2)
       end if
    else
       eq = .false.
    end if
  end function pdg_array_equivalent

  elemental function pdg_array_inequivalent (aval1, aval2) result (neq)
    logical :: neq
    type(pdg_array_t), intent(in) :: aval1, aval2
    neq = .not. pdg_array_equivalent (aval1, aval2)
  end function pdg_array_inequivalent

\eatline
\nwendcode{}\nwbegindocs{181}\nwdocspar
\subsection{Sorting}
Sort a PDG array by absolute value, particle before antiparticle.  After
sorting, we eliminate double entries.
\nwenddocs{}\nwbegincode{182}\moddef{PDG arrays: public}\plusendmoddef
  public :: sort_abs
\nwendcode{}\nwbegincode{183}\moddef{PDG arrays: interfaces}\plusendmoddef
  interface sort_abs
     module procedure pdg_array_sort_abs
  end interface
  
\nwendcode{}\nwbegincode{184}\moddef{PDG arrays: pdg array: TBP}\plusendmoddef
  procedure :: sort_abs => pdg_array_sort_abs
\nwendcode{}\nwbegincode{185}\moddef{PDG arrays: procedures}\plusendmoddef
  function pdg_array_sort_abs (aval1, unique) result (aval2)
    class(pdg_array_t), intent(in) :: aval1
    logical, intent(in), optional :: unique
    type(pdg_array_t) :: aval2
    integer, dimension(:), allocatable :: tmp
    logical, dimension(:), allocatable :: mask
    integer :: i, n
    logical :: uni
    uni = .false.;  if (present (unique))  uni = unique
    n = size (aval1%pdg)
    if (uni) then
       allocate (tmp (n), mask(n))
       tmp = sort_abs (aval1%pdg)
       mask(1) = .true.
       do i = 2, n
          mask(i) = tmp(i) /= tmp(i-1)
       end do
       allocate (aval2%pdg (count (mask)))
       aval2%pdg = pack (tmp, mask)
    else
       allocate (aval2%pdg (n))
       aval2%pdg = sort_abs (aval1%pdg)
    end if
  end function pdg_array_sort_abs

\eatline
\nwendcode{}\nwbegindocs{186}\nwdocspar
\nwenddocs{}\nwbegincode{187}\moddef{PDG arrays: pdg array: TBP}\plusendmoddef
  procedure :: intersect => pdg_array_intersect
\nwendcode{}\nwbegincode{188}\moddef{PDG arrays: procedures}\plusendmoddef
  function pdg_array_intersect (aval1, match) result (aval2)
   class(pdg_array_t), intent(in) :: aval1
   integer, dimension(:) :: match
   type(pdg_array_t) :: aval2
   integer, dimension(:), allocatable :: isec
   integer :: i
   isec = pack (aval1%pdg, [(any(aval1%pdg(i) == match), i=1,size(aval1%pdg))])
   aval2 = isec
  end function pdg_array_intersect 

\eatline
\nwendcode{}\nwbegindocs{189}\nwdocspar
\nwenddocs{}\nwbegincode{190}\moddef{PDG arrays: pdg array: TBP}\plusendmoddef
  procedure :: search_for_particle => pdg_array_search_for_particle
\nwendcode{}\nwbegincode{191}\moddef{PDG arrays: procedures}\plusendmoddef
  function pdg_array_search_for_particle (pdg, i_part) result (found)
    class(pdg_array_t), intent(in) :: pdg
    integer, intent(in) :: i_part
    logical :: found
    found = any (pdg%pdg == i_part)
  end function pdg_array_search_for_particle

\eatline
\nwendcode{}\nwbegindocs{192}\nwdocspar
\subsection{PDG array list}
A PDG array list, or PDG list, is an array of PDG-array objects with
some convenience methods.
\nwenddocs{}\nwbegincode{193}\moddef{PDG arrays: public}\plusendmoddef
  public :: pdg_list_t
\nwendcode{}\nwbegincode{194}\moddef{PDG arrays: types}\plusendmoddef
  type :: pdg_list_t
     type(pdg_array_t), dimension(:), allocatable :: a
   contains
   \LA{}PDG arrays: pdg list: TBP\RA{}
  end type pdg_list_t
  
\eatline
\nwendcode{}\nwbegindocs{195}Output, as a comma-separated list without advancing I/O.
\nwenddocs{}\nwbegincode{196}\moddef{PDG arrays: pdg list: TBP}\endmoddef
  procedure :: write => pdg_list_write
\nwendcode{}\nwbegincode{197}\moddef{PDG arrays: procedures}\plusendmoddef
  subroutine pdg_list_write (object, unit)
    class(pdg_list_t), intent(in) :: object
    integer, intent(in), optional :: unit
    integer :: u, i
    u = given_output_unit (unit)
    if (allocated (object%a)) then
       do i = 1, size (object%a)
          if (i > 1)  write (u, "(A)", advance="no")  ", "
          call object%a(i)%write (u)
       end do
    end if
  end subroutine pdg_list_write
    
\eatline
\nwendcode{}\nwbegindocs{198}Initialize for a certain size.  The entries are initially empty PDG arrays.
\nwenddocs{}\nwbegincode{199}\moddef{PDG arrays: pdg list: TBP}\plusendmoddef
  generic :: init => pdg_list_init_size
  procedure, private :: pdg_list_init_size
\nwendcode{}\nwbegincode{200}\moddef{PDG arrays: procedures}\plusendmoddef
  subroutine pdg_list_init_size (pl, n)
    class(pdg_list_t), intent(out) :: pl
    integer, intent(in) :: n
    allocate (pl%a (n))
  end subroutine pdg_list_init_size
  
\eatline
\nwendcode{}\nwbegindocs{201}Initialize with a definite array of PDG codes.  That is, each entry
in the list becomes a single-particle PDG array.
\nwenddocs{}\nwbegincode{202}\moddef{PDG arrays: pdg list: TBP}\plusendmoddef
  generic :: init => pdg_list_init_int_array
  procedure, private :: pdg_list_init_int_array
\nwendcode{}\nwbegincode{203}\moddef{PDG arrays: procedures}\plusendmoddef
  subroutine pdg_list_init_int_array (pl, pdg)
    class(pdg_list_t), intent(out) :: pl
    integer, dimension(:), intent(in) :: pdg
    integer :: i
    allocate (pl%a (size (pdg)))
    do i = 1, size (pdg)
       pl%a(i) = pdg(i)
    end do
  end subroutine pdg_list_init_int_array
  
\eatline
\nwendcode{}\nwbegindocs{204}Set one of the entries.  No bounds-check.
\nwenddocs{}\nwbegincode{205}\moddef{PDG arrays: pdg list: TBP}\plusendmoddef
  generic :: set => pdg_list_set_int
  generic :: set => pdg_list_set_int_array
  generic :: set => pdg_list_set_pdg_array
  procedure, private :: pdg_list_set_int
  procedure, private :: pdg_list_set_int_array
  procedure, private :: pdg_list_set_pdg_array
\nwendcode{}\nwbegincode{206}\moddef{PDG arrays: procedures}\plusendmoddef
  subroutine pdg_list_set_int (pl, i, pdg)
    class(pdg_list_t), intent(inout) :: pl
    integer, intent(in) :: i
    integer, intent(in) :: pdg
    pl%a(i) = pdg
  end subroutine pdg_list_set_int
  
  subroutine pdg_list_set_int_array (pl, i, pdg)
    class(pdg_list_t), intent(inout) :: pl
    integer, intent(in) :: i
    integer, dimension(:), intent(in) :: pdg
    pl%a(i) = pdg
  end subroutine pdg_list_set_int_array
  
  subroutine pdg_list_set_pdg_array (pl, i, pa)
    class(pdg_list_t), intent(inout) :: pl
    integer, intent(in) :: i
    type(pdg_array_t), intent(in) :: pa
    pl%a(i) = pa
  end subroutine pdg_list_set_pdg_array
  
\eatline
\nwendcode{}\nwbegindocs{207}Array size, not the length of individual entries
\nwenddocs{}\nwbegincode{208}\moddef{PDG arrays: pdg list: TBP}\plusendmoddef
  procedure :: get_size => pdg_list_get_size
\nwendcode{}\nwbegincode{209}\moddef{PDG arrays: procedures}\plusendmoddef
  function pdg_list_get_size (pl) result (n)
    class(pdg_list_t), intent(in) :: pl
    integer :: n
    if (allocated (pl%a)) then
       n = size (pl%a)
    else
       n = 0
    end if
  end function pdg_list_get_size
  
\eatline
\nwendcode{}\nwbegindocs{210}Return an entry, as a PDG array.
\nwenddocs{}\nwbegincode{211}\moddef{PDG arrays: pdg list: TBP}\plusendmoddef
  procedure :: get => pdg_list_get
\nwendcode{}\nwbegincode{212}\moddef{PDG arrays: procedures}\plusendmoddef
  function pdg_list_get (pl, i) result (pa)
    class(pdg_list_t), intent(in) :: pl
    integer, intent(in) :: i
    type(pdg_array_t) :: pa
    pa = pl%a(i)
  end function pdg_list_get
  
\eatline
\nwendcode{}\nwbegindocs{213}Check if the list entries are all either mutually disjoint or identical.
The individual entries (PDG arrays) should already be sorted, so we can test
for equality.
\nwenddocs{}\nwbegincode{214}\moddef{PDG arrays: pdg list: TBP}\plusendmoddef
  procedure :: is_regular => pdg_list_is_regular
\nwendcode{}\nwbegincode{215}\moddef{PDG arrays: procedures}\plusendmoddef
  function pdg_list_is_regular (pl) result (flag)
    class(pdg_list_t), intent(in) :: pl
    logical :: flag
    integer :: i, j, s
    s = pl%get_size ()
    flag = .true.
    do i = 1, s
       do j = i + 1, s
          if (pl%a(i) .match. pl%a(j)) then
             if (pl%a(i) /= pl%a(j)) then
                flag = .false.
                return
             end if
          end if
       end do
    end do
  end function pdg_list_is_regular
  
\eatline
\nwendcode{}\nwbegindocs{216}Sort the list.  First, each entry gets sorted, including elimination
of doublers.  Then, we sort the list, using the first member of each
PDG array as the marker.  No removal of doublers at this stage.

If {\tt{}n{\char95}in} is supplied, we do not reorder the first {\tt{}n{\char95}in} particle
entries.
\nwenddocs{}\nwbegincode{217}\moddef{PDG arrays: pdg list: TBP}\plusendmoddef
  procedure :: sort_abs => pdg_list_sort_abs
\nwendcode{}\nwbegincode{218}\moddef{PDG arrays: procedures}\plusendmoddef
  function pdg_list_sort_abs (pl, n_in) result (pl_sorted)
    class(pdg_list_t), intent(in) :: pl
    integer, intent(in), optional :: n_in
    type(pdg_list_t) :: pl_sorted
    type(pdg_array_t), dimension(:), allocatable :: pa
    integer, dimension(:), allocatable :: pdg, map
    integer :: i, n0
    call pl_sorted%init (pl%get_size ())
    if (allocated (pl%a)) then
       allocate (pa (size (pl%a)))
       do i = 1, size (pl%a)
          pa(i) = pl%a(i)%sort_abs (unique = .true.)
       end do
       allocate (pdg (size (pa)), source = 0)
       do i = 1, size (pa)
          if (allocated (pa(i)%pdg)) then
             if (size (pa(i)%pdg) > 0) then
                pdg(i) = pa(i)%pdg(1)
             end if
          end if
       end do
       if (present (n_in)) then
          n0 = n_in
       else
          n0 = 0
       end if
       allocate (map (size (pdg)))
       map(:n0) = [(i, i = 1, n0)]
       map(n0+1:) = n0 + order_abs (pdg(n0+1:))
       do i = 1, size (pa)
          call pl_sorted%set (i, pa(map(i)))
       end do
    end if
  end function pdg_list_sort_abs
    
\eatline
\nwendcode{}\nwbegindocs{219}Compare sorted lists: equality.  The result is undefined if some entries
are not allocated.
\nwenddocs{}\nwbegincode{220}\moddef{PDG arrays: pdg list: TBP}\plusendmoddef
  generic :: operator (==) => pdg_list_eq
  procedure, private :: pdg_list_eq
\nwendcode{}\nwbegincode{221}\moddef{PDG arrays: procedures}\plusendmoddef
  function pdg_list_eq (pl1, pl2) result (flag)
    class(pdg_list_t), intent(in) :: pl1, pl2
    logical :: flag
    integer :: i
    flag = .false.
    if (allocated (pl1%a) .and. allocated (pl2%a)) then
       if (size (pl1%a) == size (pl2%a)) then
          do i = 1, size (pl1%a)
             associate (a1 => pl1%a(i), a2 => pl2%a(i))
               if (allocated (a1%pdg) .and. allocated (a2%pdg)) then
                  if (size (a1%pdg) == size (a2%pdg)) then
                     if (size (a1%pdg) > 0) then
                        if (a1%pdg(1) /= a2%pdg(1)) return
                     end if
                  else
                     return
                  end if
               else
                  return
               end if
             end associate
          end do
          flag = .true.
       end if
    end if
  end function pdg_list_eq
  
\eatline
\nwendcode{}\nwbegindocs{222}Compare sorted lists.  The result is undefined if some entries
are not allocated.

The ordering is quite complicated.  First, a shorter list comes before
a longer list.  Comparing entry by entry, a shorter entry comes
first.  Next, we check the first PDG code within corresponding
entries.  This is compared by absolute value.  If equal, particle
comes before antiparticle.  Finally, if all is equal, the result is
false.
\nwenddocs{}\nwbegincode{223}\moddef{PDG arrays: pdg list: TBP}\plusendmoddef
  generic :: operator (<) => pdg_list_lt
  procedure, private :: pdg_list_lt
\nwendcode{}\nwbegincode{224}\moddef{PDG arrays: procedures}\plusendmoddef
  function pdg_list_lt (pl1, pl2) result (flag)
    class(pdg_list_t), intent(in) :: pl1, pl2
    logical :: flag
    integer :: i
    flag = .false.
    if (allocated (pl1%a) .and. allocated (pl2%a)) then
       if (size (pl1%a) < size (pl2%a)) then
          flag = .true.;  return
       else if (size (pl1%a) > size (pl2%a)) then
          return
       else
          do i = 1, size (pl1%a)
             associate (a1 => pl1%a(i), a2 => pl2%a(i))
               if (allocated (a1%pdg) .and. allocated (a2%pdg)) then
                  if (size (a1%pdg) < size (a2%pdg)) then
                     flag = .true.;  return
                  else if (size (a1%pdg) > size (a2%pdg)) then
                     return
                  else
                     if (size (a1%pdg) > 0) then
                        if (abs (a1%pdg(1)) < abs (a2%pdg(1))) then
                           flag = .true.;  return
                        else if (abs (a1%pdg(1)) > abs (a2%pdg(1))) then
                           return
                        else if (a1%pdg(1) > 0 .and. a2%pdg(1) < 0) then
                           flag = .true.;  return
                        else if (a1%pdg(1) < 0 .and. a2%pdg(1) > 0) then
                           return
                        end if
                     end if
                  end if
               else
                  return
               end if
             end associate
          end do
          flag = .false.
       end if
    end if
  end function pdg_list_lt
  
\eatline
\nwendcode{}\nwbegindocs{225}Replace an entry.  In the result, the entry {\tt{}{\char35}i} is replaced by
the contents of the second argument.  The result is not sorted.

If {\tt{}n{\char95}in} is also set and {\tt{}i} is less or equal to {\tt{}n{\char95}in},
replace {\tt{}{\char35}i} only by the first entry of {\tt{}pl{\char95}insert}, and insert
the remainder after entry {\tt{}n{\char95}in}.
\nwenddocs{}\nwbegincode{226}\moddef{PDG arrays: pdg list: TBP}\plusendmoddef
  procedure :: replace => pdg_list_replace
\nwendcode{}\nwbegincode{227}\moddef{PDG arrays: procedures}\plusendmoddef
  function pdg_list_replace (pl, i, pl_insert, n_in) result (pl_out)
    class(pdg_list_t), intent(in) :: pl
    integer, intent(in) :: i
    class(pdg_list_t), intent(in) :: pl_insert
    integer, intent(in), optional :: n_in
    type(pdg_list_t) :: pl_out
    integer :: n, n_insert, n_out, k
    n = pl%get_size ()
    n_insert = pl_insert%get_size ()
    n_out = n + n_insert - 1
    call pl_out%init (n_out)
!    if (allocated (pl%a)) then
       do k = 1, i - 1
          pl_out%a(k) = pl%a(k)
       end do
!    end if
    if (present (n_in)) then
       pl_out%a(i) = pl_insert%a(1)
       do k = i + 1, n_in
          pl_out%a(k) = pl%a(k)
       end do
       do k = 1, n_insert - 1
          pl_out%a(n_in+k) = pl_insert%a(1+k)
       end do
       do k = 1, n - n_in
          pl_out%a(n_in+k+n_insert-1) = pl%a(n_in+k)
       end do
    else
!       if (allocated (pl_insert%a)) then
          do k = 1, n_insert
             pl_out%a(i-1+k) = pl_insert%a(k)
          end do
!       end if
!       if (allocated (pl%a)) then
          do k = 1, n - i
             pl_out%a(i+n_insert-1+k) = pl%a(i+k)
          end do
       end if
!    end if
  end function pdg_list_replace
    
\eatline
\nwendcode{}\nwbegindocs{228}Replace the entries of {\tt{}pl} by the matching entries of {\tt{}pl{\char95}match}, one by
one.  This is done in-place.  If there is no match, return failure.
\nwenddocs{}\nwbegincode{229}\moddef{PDG arrays: pdg list: TBP}\plusendmoddef
  procedure :: match_replace => pdg_list_match_replace
\nwendcode{}\nwbegincode{230}\moddef{PDG arrays: procedures}\plusendmoddef
  subroutine pdg_list_match_replace (pl, pl_match, success)
    class(pdg_list_t), intent(inout) :: pl
    class(pdg_list_t), intent(in) :: pl_match
    logical, intent(out) :: success
    integer :: i, j
    success = .true.
    SCAN_ENTRIES: do i = 1, size (pl%a)
       do j = 1, size (pl_match%a)
          if (pl%a(i) .match. pl_match%a(j)) then
             pl%a(i) = pl_match%a(j)
             cycle SCAN_ENTRIES
          end if
       end do
       success = .false.
       return
    end do SCAN_ENTRIES
  end subroutine pdg_list_match_replace
 
\eatline
\nwendcode{}\nwbegindocs{231}Just check if a PDG array matches any entry in the PDG list.  The second
version returns the position of the match within the list.  An optional mask
indicates the list elements that should be checked.
\nwenddocs{}\nwbegincode{232}\moddef{PDG arrays: pdg list: TBP}\plusendmoddef
  generic :: operator (.match.) => pdg_list_match_pdg_array
  procedure, private :: pdg_list_match_pdg_array
  procedure :: find_match => pdg_list_find_match_pdg_array
\nwendcode{}\nwbegincode{233}\moddef{PDG arrays: procedures}\plusendmoddef
  function pdg_list_match_pdg_array (pl, pa) result (flag)
    class(pdg_list_t), intent(in) :: pl
    type(pdg_array_t), intent(in) :: pa
    logical :: flag
    flag = pl%find_match (pa) /= 0
  end function pdg_list_match_pdg_array
  
  function pdg_list_find_match_pdg_array (pl, pa, mask) result (i)
    class(pdg_list_t), intent(in) :: pl
    type(pdg_array_t), intent(in) :: pa
    logical, dimension(:), intent(in), optional :: mask
    integer :: i
    do i = 1, size (pl%a)
       if (present (mask)) then
          if (.not. mask(i))  cycle
       end if
       if (pl%a(i) .match. pa)  return
    end do
    i = 0
  end function pdg_list_find_match_pdg_array
  
\eatline
\eatline
\nwendcode{}\nwbegindocs{234}Some old compilers have problems with allocatable arrays as
intent(out) or as function result, so be conservative here:
\nwenddocs{}\nwbegincode{235}\moddef{PDG arrays: pdg list: TBP}\plusendmoddef
  procedure :: create_pdg_array => pdg_list_create_pdg_array
\nwendcode{}\nwbegincode{236}\moddef{PDG arrays: procedures}\plusendmoddef
  subroutine pdg_list_create_pdg_array (pl, pdg)
    class(pdg_list_t), intent(in) :: pl
    type(pdg_array_t), dimension(:), intent(inout), allocatable :: pdg
    integer :: n_elements
    integer :: i
    associate (a => pl%a)
      n_elements = size (a)
      if (allocated (pdg))  deallocate (pdg)
      allocate (pdg (n_elements))
      do i = 1, n_elements
         pdg(i) = a(i)
      end do
    end associate
  end subroutine pdg_list_create_pdg_array

\eatline
\nwendcode{}\nwbegindocs{237}\nwdocspar
\nwenddocs{}\nwbegincode{238}\moddef{PDG arrays: pdg list: TBP}\plusendmoddef
  procedure :: search_for_particle => pdg_list_search_for_particle
\nwendcode{}\nwbegincode{239}\moddef{PDG arrays: procedures}\plusendmoddef
  function pdg_list_search_for_particle (pl, i_part) result (found)
    class(pdg_list_t), intent(in) :: pl
    integer, intent(in) :: i_part
    logical :: found
    integer :: i_pl
    do i_pl = 1, size (pl%a)
       found = pl%a(i_pl)%search_for_particle (i_part)
       if (found) return
    end do
  end function pdg_list_search_for_particle

\eatline
\nwendcode{}\nwbegindocs{240}\nwdocspar
\subsection{Unit tests}
\nwenddocs{}\nwbegincode{241}\moddef{PDG arrays: public}\plusendmoddef
  public :: pdg_arrays_test 
\nwendcode{}\nwbegincode{242}\moddef{PDG arrays: procedures}\plusendmoddef
  subroutine pdg_arrays_test (u, results)
    integer, intent(in) :: u
    type (test_results_t), intent(inout) :: results
  \LA{}PDG arrays: execute tests\RA{}  
  end subroutine pdg_arrays_test

\eatline
\nwendcode{}\nwbegindocs{243}Basic functionality.
\nwenddocs{}\nwbegincode{244}\moddef{PDG arrays: execute tests}\endmoddef
  call test (pdg_arrays_1, "pdg_arrays_1", &
       "create and sort PDG array", &
       u, results) 
\nwendcode{}\nwbegincode{245}\moddef{PDG arrays: tests}\endmoddef
  subroutine pdg_arrays_1 (u)
    integer, intent(in) :: u

    type(pdg_array_t) :: pa, pa1, pa2, pa3, pa4, pa5, pa6
    integer, dimension(:), allocatable :: pdg

    write (u, "(A)")  "* Test output: pdg_arrays_1"
    write (u, "(A)")  "*   Purpose: create and sort PDG arrays"
    write (u, "(A)")
    
    write (u, "(A)")  "* Assignment"
    write (u, "(A)")
    
    call pa%write (u)
    write (u, *)
    write (u, "(A,I0)")  "length = ", pa%get_length ()
    pdg = pa
    write (u, "(A,3(1x,I0))")  "contents = ", pdg
    
    write (u, *)
    pa = 1
    call pa%write (u)
    write (u, *)
    write (u, "(A,I0)")  "length = ", pa%get_length ()
    pdg = pa
    write (u, "(A,3(1x,I0))")  "contents = ", pdg
    
    write (u, *)
    pa = [1, 2, 3]
    call pa%write (u)
    write (u, *)
    write (u, "(A,I0)")  "length = ", pa%get_length ()
    pdg = pa
    write (u, "(A,3(1x,I0))")  "contents = ", pdg
    write (u, "(A,I0)")  "element #2 = ", pa%get (2)
    
    write (u, *)
    write (u, "(A)")  "* Replace"
    write (u, *)

    pa = pa%replace (2, [-5, 5, -7])
    call pa%write (u)
    write (u, *)
    
    write (u, *)
    write (u, "(A)")  "* Sort"
    write (u, *)

    pa = [1, -7, 3, -5, 5, 3]
    call pa%write (u)
    write (u, *)
    pa1 = pa%sort_abs ()
    pa2 = pa%sort_abs (unique = .true.)
    call pa1%write (u)
    write (u, *)
    call pa2%write (u)
    write (u, *)
    
    write (u, *)
    write (u, "(A)")  "* Compare"
    write (u, *)

    pa1 = [1, 3]
    pa2 = [1, 2, -2]
    pa3 = [1, 2, 4]
    pa4 = [1, 2, 4]
    pa5 = [1, 2, -4]
    pa6 = [1, 2, -3]
    
    write (u, "(A,6(1x,L1))")  "< ", &
         pa1 < pa2, pa2 < pa3, pa3 < pa4, pa4 < pa5, pa5 < pa6, pa6 < pa1
    write (u, "(A,6(1x,L1))")  "> ", &
         pa1 > pa2, pa2 > pa3, pa3 > pa4, pa4 > pa5, pa5 > pa6, pa6 > pa1
    write (u, "(A,6(1x,L1))")  "<=", &
         pa1 <= pa2, pa2 <= pa3, pa3 <= pa4, pa4 <= pa5, pa5 <= pa6, pa6 <= pa1
    write (u, "(A,6(1x,L1))")  ">=", &
         pa1 >= pa2, pa2 >= pa3, pa3 >= pa4, pa4 >= pa5, pa5 >= pa6, pa6 >= pa1
    write (u, "(A,6(1x,L1))")  "==", &
         pa1 == pa2, pa2 == pa3, pa3 == pa4, pa4 == pa5, pa5 == pa6, pa6 == pa1
    write (u, "(A,6(1x,L1))")  "/=", &
         pa1 /= pa2, pa2 /= pa3, pa3 /= pa4, pa4 /= pa5, pa5 /= pa6, pa6 /= pa1
   
    write (u, *)
    pa1 = [0]
    pa2 = [1, 2]
    pa3 = [1, -2]
    
    write (u, "(A,6(1x,L1))")  "eqv ", &
         pa1 .eqv. pa1, pa1 .eqv. pa2, &
         pa2 .eqv. pa2, pa2 .eqv. pa3
    
    write (u, "(A,6(1x,L1))")  "neqv", &
         pa1 .neqv. pa1, pa1 .neqv. pa2, &
         pa2 .neqv. pa2, pa2 .neqv. pa3
    

    write (u, *)
    write (u, "(A,6(1x,L1))")  "match", &
         pa1 .match. 0, pa1 .match. 1, &
         pa2 .match. 0, pa2 .match. 1, pa2 .match. 3

    write (u, "(A)")
    write (u, "(A)")  "* Test output end: pdg_arrays_1"        
    
  end subroutine pdg_arrays_1
  
\eatline
\nwendcode{}\nwbegindocs{246}PDG array list, i.e., arrays of arrays.
\nwenddocs{}\nwbegincode{247}\moddef{PDG arrays: execute tests}\plusendmoddef
  call test (pdg_arrays_2, "pdg_arrays_2", &
       "create and sort PDG lists", &
       u, results) 
\nwendcode{}\nwbegincode{248}\moddef{PDG arrays: tests}\plusendmoddef
  subroutine pdg_arrays_2 (u)
    integer, intent(in) :: u

    type(pdg_array_t) :: pa
    type(pdg_list_t) :: pl, pl1

    write (u, "(A)")  "* Test output: pdg_arrays_2"
    write (u, "(A)")  "*   Purpose: create and sort PDG lists"
    write (u, "(A)")
    
    write (u, "(A)")  "* Assignment"
    write (u, "(A)")
    
    call pl%init (3)
    call pl%set (1, 42)
    call pl%set (2, [3, 2])
    pa = [5, -5]
    call pl%set (3, pa)
    call pl%write (u)
    write (u, *)
    write (u, "(A,I0)")  "size = ", pl%get_size ()

    write (u, "(A)")
    write (u, "(A)")  "* Sort"
    write (u, "(A)")
    
    pl = pl%sort_abs ()
    call pl%write (u)
    write (u, *)

    write (u, "(A)")
    write (u, "(A)")  "* Extract item #3"
    write (u, "(A)")
    
    pa = pl%get (3)
    call pa%write (u)
    write (u, *)

    write (u, "(A)")
    write (u, "(A)")  "* Replace item #3"
    write (u, "(A)")
    
    call pl1%init (2)
    call pl1%set (1, [2, 4])
    call pl1%set (2, -7)
    
    pl = pl%replace (3, pl1)
    call pl%write (u)
    write (u, *)

    write (u, "(A)")
    write (u, "(A)")  "* Test output end: pdg_arrays_2"        
    
  end subroutine pdg_arrays_2
  
\eatline
\nwendcode{}\nwbegindocs{249}Check if a (sorted) PDG array lists is regular.  The entries (PDG arrays)
must not overlap, unless they are identical.
\nwenddocs{}\nwbegincode{250}\moddef{PDG arrays: execute tests}\plusendmoddef
  call test (pdg_arrays_3, "pdg_arrays_3", &
       "check PDG lists", &
       u, results) 
\nwendcode{}\nwbegincode{251}\moddef{PDG arrays: tests}\plusendmoddef
  subroutine pdg_arrays_3 (u)
    integer, intent(in) :: u

    type(pdg_list_t) :: pl

    write (u, "(A)")  "* Test output: pdg_arrays_3"
    write (u, "(A)")  "*   Purpose: check for regular PDG lists"
    write (u, "(A)")
    
    write (u, "(A)")  "* Regular list"
    write (u, "(A)")
    
    call pl%init (4)
    call pl%set (1, [1, 2])
    call pl%set (2, [1, 2])
    call pl%set (3, [5, -5])
    call pl%set (4, 42)
    call pl%write (u)
    write (u, *)
    write (u, "(L1)") pl%is_regular ()

    write (u, "(A)")
    write (u, "(A)")  "* Irregular list"
    write (u, "(A)")
    
    call pl%init (4)
    call pl%set (1, [1, 2])
    call pl%set (2, [1, 2])
    call pl%set (3, [2, 5, -5])
    call pl%set (4, 42)
    call pl%write (u)
    write (u, *)
    write (u, "(L1)") pl%is_regular ()

    write (u, "(A)")
    write (u, "(A)")  "* Test output end: pdg_arrays_3"        
    
  end subroutine pdg_arrays_3
  
\eatline
\nwendcode{}\nwbegindocs{252}Compare PDG array lists.  The lists must be regular, i.e., sorted and with
non-overlapping (or identical) entries.
\nwenddocs{}\nwbegincode{253}\moddef{PDG arrays: execute tests}\plusendmoddef
  call test (pdg_arrays_4, "pdg_arrays_4", &
       "compare PDG lists", &
       u, results) 
\nwendcode{}\nwbegincode{254}\moddef{PDG arrays: tests}\plusendmoddef
  subroutine pdg_arrays_4 (u)
    integer, intent(in) :: u

    type(pdg_list_t) :: pl1, pl2, pl3

    write (u, "(A)")  "* Test output: pdg_arrays_4"
    write (u, "(A)")  "*   Purpose: check for regular PDG lists"
    write (u, "(A)")
    
    write (u, "(A)")  "* Create lists"
    write (u, "(A)")
    
    call pl1%init (4)
    call pl1%set (1, [1, 2])
    call pl1%set (2, [1, 2])
    call pl1%set (3, [5, -5])
    call pl1%set (4, 42)
    write (u, "(I1,1x)", advance = "no")  1
    call pl1%write (u)
    write (u, *)

    call pl2%init (2)
    call pl2%set (1, 3)
    call pl2%set (2, [5, -5])
    write (u, "(I1,1x)", advance = "no")  2
    call pl2%write (u)
    write (u, *)

    call pl3%init (2)
    call pl3%set (1, 4)
    call pl3%set (2, [5, -5])
    write (u, "(I1,1x)", advance = "no")  3
    call pl3%write (u)
    write (u, *)

    write (u, "(A)")
    write (u, "(A)")  "* a == b"
    write (u, "(A)")
    
    write (u, "(2x,A)")  "123"
    write (u, *)
    write (u, "(I1,1x,4L1)")  1, pl1 == pl1, pl1 == pl2, pl1 == pl3
    write (u, "(I1,1x,4L1)")  2, pl2 == pl1, pl2 == pl2, pl2 == pl3
    write (u, "(I1,1x,4L1)")  3, pl3 == pl1, pl3 == pl2, pl3 == pl3

    write (u, "(A)")
    write (u, "(A)")  "* a < b"
    write (u, "(A)")
    
    write (u, "(2x,A)")  "123"
    write (u, *)
    write (u, "(I1,1x,4L1)")  1, pl1 < pl1, pl1 < pl2, pl1 < pl3
    write (u, "(I1,1x,4L1)")  2, pl2 < pl1, pl2 < pl2, pl2 < pl3
    write (u, "(I1,1x,4L1)")  3, pl3 < pl1, pl3 < pl2, pl3 < pl3

    write (u, "(A)")
    write (u, "(A)")  "* Test output end: pdg_arrays_4"        
    
  end subroutine pdg_arrays_4
  
\eatline
\nwendcode{}\nwbegindocs{255}Match-replace: translate all entries in the first list into the
matching entries of the second list, if there is a match.
\nwenddocs{}\nwbegincode{256}\moddef{PDG arrays: execute tests}\plusendmoddef
  call test (pdg_arrays_5, "pdg_arrays_5", &
       "match PDG lists", &
       u, results) 
\nwendcode{}\nwbegincode{257}\moddef{PDG arrays: tests}\plusendmoddef
  subroutine pdg_arrays_5 (u)
    integer, intent(in) :: u

    type(pdg_list_t) :: pl1, pl2, pl3
    logical :: success

    write (u, "(A)")  "* Test output: pdg_arrays_5"
    write (u, "(A)")  "*   Purpose: match-replace"
    write (u, "(A)")
    
    write (u, "(A)")  "* Create lists"
    write (u, "(A)")
    
    call pl1%init (2)
    call pl1%set (1, [1, 2])
    call pl1%set (2, 42)
    call pl1%write (u)
    write (u, *)
    call pl3%init (2)
    call pl3%set (1, [42, -42])
    call pl3%set (2, [1, 2, 3, 4])
    call pl1%match_replace (pl3, success)
    call pl3%write (u)
    write (u, "(1x,A,1x,L1,':',1x)", advance="no")  "=>", success
    call pl1%write (u)
    write (u, *)

    write (u, *)

    call pl2%init (2)
    call pl2%set (1, 9)
    call pl2%set (2, 42)
    call pl2%write (u)
    write (u, *)
    call pl2%match_replace (pl3, success)
    call pl3%write (u)
    write (u, "(1x,A,1x,L1,':',1x)", advance="no")  "=>", success
    call pl2%write (u)
    write (u, *)

    write (u, "(A)")
    write (u, "(A)")  "* Test output end: pdg_arrays_5"        
    
  end subroutine pdg_arrays_5
  
\eatline
\nwendcode{}\nwbegindocs{258}\nwdocspar
\clearpage
%%%%%%%%%%%%%%%%%%%%%%%%%%%%%%%%%%%%%%%%%%%%%%%%%%%%%%%%%%%%%%%%%%%%%%%%
\section{Jets}

The FastJet library is linked externally, if available.  The wrapper code is
also in a separate directory.  Here, we define \whizard-specific procedures
and tests.
\nwenddocs{}\nwbegincode{259}\moddef{\code{}jets.f90\edoc{}}\endmoddef
\LA{}File header\RA{}

module jets
    
\LA{}Use kinds\RA{}
  use unit_tests
  use fastjet !NODEP!
  
\LA{}Standard module head\RA{}

\LA{}Jets: public\RA{}

contains
  
\LA{}Jets: procedures\RA{}

\LA{}Jets: tests\RA{}

end module jets
\eatline
\nwendcode{}\nwbegindocs{260}\nwdocspar
\subsection{Re-exported symbols}
We use this module as a proxy for the FastJet interface, therefore we
re-export some symbols.
\nwenddocs{}\nwbegincode{261}\moddef{Jets: public}\endmoddef
  public :: fastjet_available
  public :: fastjet_init
  public :: jet_definition_t
  public :: pseudojet_t
  public :: pseudojet_vector_t
  public :: cluster_sequence_t
  public :: assignment (=)

\eatline
\nwendcode{}\nwbegindocs{262}The initialization routine prints the banner.
\nwenddocs{}\nwbegincode{263}\moddef{Jets: procedures}\endmoddef
  subroutine fastjet_init ()
    call print_banner ()
  end subroutine fastjet_init
  
\eatline
\nwendcode{}\nwbegindocs{264}The jet algorithm codes (actually, integers)
\nwenddocs{}\nwbegincode{265}\moddef{Jets: public}\plusendmoddef
  public :: kt_algorithm
  public :: cambridge_algorithm
  public :: antikt_algorithm
  public :: genkt_algorithm
  public :: cambridge_for_passive_algorithm
  public :: genkt_for_passive_algorithm
  public :: ee_kt_algorithm
  public :: ee_genkt_algorithm
  public :: plugin_algorithm
  public :: undefined_jet_algorithm
  
\nwendcode{}\nwbegindocs{266}\nwdocspar
\subsection{Unit tests}
\nwenddocs{}\nwbegincode{267}\moddef{Jets: public}\plusendmoddef
  public :: jets_test 
\nwendcode{}\nwbegincode{268}\moddef{Jets: procedures}\plusendmoddef
  subroutine jets_test (u, results)
    integer, intent(in) :: u
    type (test_results_t), intent(inout) :: results
  \LA{}Jets: execute tests\RA{}  
  end subroutine jets_test

\eatline
\nwendcode{}\nwbegindocs{269}This test is actually the minimal example from the FastJet manual,
translated to Fortran.

Note that FastJet creates pseudojet vectors, which we mirror in the
{\tt{}pseudojet{\char95}vector{\char95}t}, but immediately assign to pseudojet arrays.  Without
automatic finalization available in the compilers, we should avoid this in
actual code and rather introduce intermediate variables for those objects,
which we can finalize explicitly.
\nwenddocs{}\nwbegincode{270}\moddef{Jets: execute tests}\endmoddef
  call test (jets_1, "jets_1", &
       "basic FastJet functionality", &
       u, results) 
\nwendcode{}\nwbegincode{271}\moddef{Jets: tests}\endmoddef
  subroutine jets_1 (u)
    integer, intent(in) :: u

    type(pseudojet_t), dimension(:), allocatable :: prt, jets, constituents
    type(jet_definition_t) :: jet_def
    type(cluster_sequence_t) :: cs

    integer, parameter :: dp = default
    integer :: i, j

    write (u, "(A)")  "* Test output: jets_1"
    write (u, "(A)")  "*   Purpose: test basic FastJet functionality"
    write (u, "(A)")    

    write (u, "(A)")  "* Print banner"
    call print_banner ()
    
    write (u, *)
    write (u, "(A)")  "* Prepare input particles"
    allocate (prt (3))
    call prt(1)%init ( 99._dp, 0.1_dp, 0._dp, 100._dp)
    call prt(2)%init (  4._dp,-0.1_dp, 0._dp,   5._dp)
    call prt(3)%init (-99._dp, 0._dp,  0._dp,  99._dp)

    write (u, *)
    write (u, "(A)")  "* Define jet algorithm"
    call jet_def%init (antikt_algorithm, 0.7_dp)

    write (u, *)
    write (u, "(A)")  "* Cluster particles according to jet algorithm"

    write (u, *)
    write (u, "(A,A)")  "Clustering with ", jet_def%description ()
    call cs%init (pseudojet_vector (prt), jet_def)

    write (u, *)
    write (u, "(A)")  "* Sort output jets"
    jets = sorted_by_pt (cs%inclusive_jets ())

    write (u, *)
    write (u, "(A)")  "* Print jet observables and constituents"
    write (u, *)
    write (u, "(4x,3(7x,A3))") "pt", "y", "phi"
    do i = 1, size (jets)
       write (u, "(A,1x,I0,A,3(1x,F9.5))") &
            "jet", i, ":", jets(i)%perp (), jets(i)%rap (), jets(i)%phi ()
       constituents = jets(i)%constituents ()
       do j = 1, size (constituents)
          write (u, "(4x,A,1x,I0,A,F9.5)") &
               "constituent", j, "'s pt:", constituents(j)%perp ()
       end do
       do j = 1, size (constituents)
          call constituents(j)%final ()
       end do
    end do
    
    write (u, *)
    write (u, "(A)")  "* Cleanup"
    
    do i = 1, size (prt)
       call prt(i)%final ()
    end do
    do i = 1, size (jets)
       call jets(i)%final ()
    end do
    call jet_def%final ()
    call cs%final ()

    write (u, "(A)")
    write (u, "(A)")  "* Test output end: jets_1"        
    
  end subroutine jets_1
  
\eatline
\nwendcode{}\nwbegindocs{272}\nwdocspar
\clearpage
%%%%%%%%%%%%%%%%%%%%%%%%%%%%%%%%%%%%%%%%%%%%%%%%%%%%%%%%%%%%%%%%%%%%%%%%
\section{Subevents}

The purpose of subevents is to store the relevant part of the physical
event (either partonic or hadronic), and to hold particle selections
and combinations which are constructed in cut or analysis expressions.
\nwenddocs{}\nwbegincode{273}\moddef{\code{}subevents.f90\edoc{}}\endmoddef
\LA{}File header\RA{}

module subevents

  use, intrinsic :: iso_c_binding !NODEP!
  
\LA{}Use kinds\RA{}
  use io_units
  use format_defs, only: FMT_14, FMT_19
  use format_utils, only: pac_fmt
  use sorting
  use c_particles
  use lorentz
  use pdg_arrays
  use jets

\LA{}Standard module head\RA{}

\LA{}Subevents: public\RA{}

\LA{}Subevents: parameters\RA{}

\LA{}Subevents: types\RA{}

\LA{}Subevents: interfaces\RA{}

contains

\LA{}Subevents: procedures\RA{}

end module subevents
\eatline
\nwendcode{}\nwbegindocs{274}\nwdocspar
\subsection{Particles}
For the purpose of this module, a particle has a type which can
indicate a beam, incoming, outgoing, or composite particle, flavor and
helicity codes (integer, undefined for composite), four-momentum and
invariant mass squared.  (Other particles types are used in extended
event types, but also defined here.)  Furthermore, each particle has
an allocatable array of ancestors -- particle indices which indicate
the building blocks of a composite particle.  For an incoming/outgoing
particle, the array contains only the index of the particle itself.

For incoming particles, the momentum is inverted before storing it in
the particle object.
\nwenddocs{}\nwbegincode{275}\moddef{Subevents: parameters}\endmoddef
  integer, parameter, public :: PRT_UNDEFINED = 0 
  integer, parameter, public :: PRT_BEAM = -9
  integer, parameter, public :: PRT_INCOMING = 1
  integer, parameter, public :: PRT_OUTGOING = 2
  integer, parameter, public :: PRT_COMPOSITE = 3
  integer, parameter, public :: PRT_VIRTUAL = 4
  integer, parameter, public :: PRT_RESONANT = 5
  integer, parameter, public :: PRT_BEAM_REMNANT = 9

\eatline
\eatline
\eatline
\eatline
\nwendcode{}\nwbegindocs{276}\nwdocspar
\subsubsection{The type}
We initialize only the type here and mark as unpolarized.  The
initializers below do the rest.
\nwenddocs{}\nwbegincode{277}\moddef{Subevents: public}\endmoddef
  public :: prt_t
\nwendcode{}\nwbegincode{278}\moddef{Subevents: types}\endmoddef
  type :: prt_t
     private
     integer :: type = PRT_UNDEFINED
     integer :: pdg
     logical :: polarized = .false.
     integer :: h
     type(vector4_t) :: p
     real(default) :: p2
     integer, dimension(:), allocatable :: src
  end type prt_t

\eatline
\nwendcode{}\nwbegindocs{279}Initializers.  Polarization is set separately.  Finalizers are not
needed.
\nwenddocs{}\nwbegincode{280}\moddef{Subevents: procedures}\endmoddef
  subroutine prt_init_beam (prt, pdg, p, p2, src)
    type(prt_t), intent(out) :: prt
    integer, intent(in) :: pdg
    type(vector4_t), intent(in) :: p
    real(default), intent(in) :: p2
    integer, dimension(:), intent(in) :: src
    prt%type = PRT_BEAM
    call prt_set (prt, pdg, - p, p2, src)
  end subroutine prt_init_beam
    
  subroutine prt_init_incoming (prt, pdg, p, p2, src)
    type(prt_t), intent(out) :: prt
    integer, intent(in) :: pdg
    type(vector4_t), intent(in) :: p
    real(default), intent(in) :: p2
    integer, dimension(:), intent(in) :: src
    prt%type = PRT_INCOMING
    call prt_set (prt, pdg, - p, p2, src)
  end subroutine prt_init_incoming
    
  subroutine prt_init_outgoing (prt, pdg, p, p2, src)
    type(prt_t), intent(out) :: prt
    integer, intent(in) :: pdg
    type(vector4_t), intent(in) :: p
    real(default), intent(in) :: p2
    integer, dimension(:), intent(in) :: src
    prt%type = PRT_OUTGOING
    call prt_set (prt, pdg, p, p2, src)
  end subroutine prt_init_outgoing
    
  subroutine prt_init_composite (prt, p, src)
    type(prt_t), intent(out) :: prt
    type(vector4_t), intent(in) :: p
    integer, dimension(:), intent(in) :: src
    prt%type = PRT_COMPOSITE
    call prt_set (prt, 0, p, p**2, src)
  end subroutine prt_init_composite

\eatline
\nwendcode{}\nwbegindocs{281}\nwdocspar
This version is for temporary particle objects, so the {\tt{}src} array
is not set.
\nwenddocs{}\nwbegincode{282}\moddef{Subevents: public}\plusendmoddef
  public :: prt_init_combine
\nwendcode{}\nwbegincode{283}\moddef{Subevents: procedures}\plusendmoddef
  subroutine prt_init_combine (prt, prt1, prt2)
    type(prt_t), intent(out) :: prt
    type(prt_t), intent(in) :: prt1, prt2
    type(vector4_t) :: p
    integer, dimension(0) :: src
    prt%type = PRT_COMPOSITE
    p = prt1%p + prt2%p
    call prt_set (prt, 0, p, p**2, src)
  end subroutine prt_init_combine

\eatline
\nwendcode{}\nwbegindocs{284}Init from a pseudojet object.
\nwenddocs{}\nwbegincode{285}\moddef{Subevents: procedures}\plusendmoddef
  subroutine prt_init_pseudojet (prt, jet, src)
    type(prt_t), intent(out) :: prt
    type(pseudojet_t), intent(in) :: jet
    integer, dimension(:), intent(in) :: src
    type(vector4_t) :: p
    prt%type = PRT_COMPOSITE
    p = vector4_moving (jet%e(), &
         vector3_moving ([jet%px(), jet%py(), jet%pz()]))
    call prt_set (prt, 0, p, p**2, src)
  end subroutine prt_init_pseudojet
  
\eatline
\nwendcode{}\nwbegindocs{286}\nwdocspar
\subsubsection{Accessing contents}
\nwenddocs{}\nwbegincode{287}\moddef{Subevents: public}\plusendmoddef
  public :: prt_get_pdg
\nwendcode{}\nwbegincode{288}\moddef{Subevents: procedures}\plusendmoddef
  elemental function prt_get_pdg (prt) result (pdg)
    integer :: pdg
    type(prt_t), intent(in) :: prt
    pdg = prt%pdg
  end function prt_get_pdg

\eatline
\nwendcode{}\nwbegincode{289}\moddef{Subevents: public}\plusendmoddef
  public :: prt_get_momentum
\nwendcode{}\nwbegincode{290}\moddef{Subevents: procedures}\plusendmoddef
  elemental function prt_get_momentum (prt) result (p)
    type(vector4_t) :: p
    type(prt_t), intent(in) :: prt
    p = prt%p
  end function prt_get_momentum

\eatline
\nwendcode{}\nwbegincode{291}\moddef{Subevents: public}\plusendmoddef
  public :: prt_get_msq
\nwendcode{}\nwbegincode{292}\moddef{Subevents: procedures}\plusendmoddef
  elemental function prt_get_msq (prt) result (msq)
    real(default) :: msq
    type(prt_t), intent(in) :: prt
    msq = prt%p2
  end function prt_get_msq

\eatline
\nwendcode{}\nwbegincode{293}\moddef{Subevents: public}\plusendmoddef
  public :: prt_is_polarized
\nwendcode{}\nwbegincode{294}\moddef{Subevents: procedures}\plusendmoddef
  elemental function prt_is_polarized (prt) result (flag)
    logical :: flag
    type(prt_t), intent(in) :: prt
    flag = prt%polarized
  end function prt_is_polarized

\eatline
\nwendcode{}\nwbegincode{295}\moddef{Subevents: public}\plusendmoddef
  public :: prt_get_helicity
\nwendcode{}\nwbegincode{296}\moddef{Subevents: procedures}\plusendmoddef
  elemental function prt_get_helicity (prt) result (h)
    integer :: h
    type(prt_t), intent(in) :: prt
    h = prt%h
  end function prt_get_helicity

\eatline
\nwendcode{}\nwbegindocs{297}\nwdocspar
\subsubsection{Setting data}
Set the PDG, momentum and momentum squared, and ancestors.  If
allocate-on-assignment is available, this can be simplified.
\nwenddocs{}\nwbegincode{298}\moddef{Subevents: procedures}\plusendmoddef
  subroutine prt_set (prt, pdg, p, p2, src)
    type(prt_t), intent(inout) :: prt
    integer, intent(in) :: pdg
    type(vector4_t), intent(in) :: p
    real(default), intent(in) :: p2
    integer, dimension(:), intent(in) :: src
    prt%pdg = pdg
    prt%p = p
    prt%p2 = p2
    if (allocated (prt%src)) then
       if (size (src) /= size (prt%src)) then
          deallocate (prt%src)
          allocate (prt%src (size (src)))
       end if
    else
       allocate (prt%src (size (src)))
    end if
    prt%src = src
  end subroutine prt_set

\eatline
\nwendcode{}\nwbegindocs{299}Set the particle PDG code separately.
\nwenddocs{}\nwbegincode{300}\moddef{Subevents: procedures}\plusendmoddef
  elemental subroutine prt_set_pdg (prt, pdg)
    type(prt_t), intent(inout) :: prt
    integer, intent(in) :: pdg
    prt%pdg = pdg
  end subroutine prt_set_pdg

\eatline
\nwendcode{}\nwbegindocs{301}Set the momentum separately.
\nwenddocs{}\nwbegincode{302}\moddef{Subevents: procedures}\plusendmoddef
  elemental subroutine prt_set_p (prt, p)
    type(prt_t), intent(inout) :: prt
    type(vector4_t), intent(in) :: p
    prt%p = p
  end subroutine prt_set_p

\eatline
\nwendcode{}\nwbegindocs{303}Set the squared invariant mass separately.
\nwenddocs{}\nwbegincode{304}\moddef{Subevents: procedures}\plusendmoddef
  elemental subroutine prt_set_p2 (prt, p2)
    type(prt_t), intent(inout) :: prt
    real(default), intent(in) :: p2
    prt%p2 = p2
  end subroutine prt_set_p2

\eatline
\nwendcode{}\nwbegindocs{305}Set helicity (optional).
\nwenddocs{}\nwbegincode{306}\moddef{Subevents: procedures}\plusendmoddef
  subroutine prt_polarize (prt, h)
    type(prt_t), intent(inout) :: prt
    integer, intent(in) :: h
    prt%polarized = .true.
    prt%h = h
  end subroutine prt_polarize

\eatline
\nwendcode{}\nwbegindocs{307}\nwdocspar
\subsubsection{Conversion}
Transform a {\tt{}prt{\char95}t} object into a {\tt{}c{\char95}prt{\char95}t} object.
\nwenddocs{}\nwbegincode{308}\moddef{Subevents: public}\plusendmoddef
  public :: c_prt
\nwendcode{}\nwbegincode{309}\moddef{Subevents: interfaces}\endmoddef
  interface c_prt
     module procedure c_prt_from_prt
  end interface

\eatline
\nwendcode{}\nwbegincode{310}\moddef{Subevents: procedures}\plusendmoddef
  elemental function c_prt_from_prt (prt) result (c_prt)
    type(c_prt_t) :: c_prt
    type(prt_t), intent(in) :: prt
    c_prt%type = prt%type
    c_prt%pdg = prt%pdg
    if (prt%polarized) then
       c_prt%polarized = 1
    else
       c_prt%polarized = 0
    end if
    c_prt%h = prt%h
    c_prt%pe = energy (prt%p)
    c_prt%px = vector4_get_component (prt%p, 1)
    c_prt%py = vector4_get_component (prt%p, 2)
    c_prt%pz = vector4_get_component (prt%p, 3)
    c_prt%p2 = prt%p2
  end function c_prt_from_prt

\eatline
\nwendcode{}\nwbegindocs{311}\nwdocspar
\subsubsection{Output}
\nwenddocs{}\nwbegincode{312}\moddef{Subevents: public}\plusendmoddef
  public :: prt_write
\nwendcode{}\nwbegincode{313}\moddef{Subevents: procedures}\plusendmoddef
  subroutine prt_write (prt, unit, pacified)
    type(prt_t), intent(in) :: prt
    integer, intent(in), optional :: unit
    logical, intent(in), optional :: pacified
    real(default), dimension(0:3) :: arr
    logical :: num_pac
    type(prt_t) :: tmp
    character(len=7) :: fmt   
    integer :: u, i
    call pac_fmt (fmt, FMT_19, FMT_14, pacified)
    u = given_output_unit (unit);  if (u < 0)  return
    num_pac = .false. ; if (present (pacified))  num_pac = pacified
    write (u, "(1x,A)", advance="no")  "prt("
    select case (prt%type)
    case (PRT_UNDEFINED);    write (u, "('?')", advance="no")
    case (PRT_BEAM);         write (u, "('b:')", advance="no")
    case (PRT_INCOMING);     write (u, "('i:')", advance="no")
    case (PRT_OUTGOING);     write (u, "('o:')", advance="no")
    case (PRT_COMPOSITE);    write (u, "('c:')", advance="no")
    end select
    select case (prt%type)
    case (PRT_BEAM, PRT_INCOMING, PRT_OUTGOING)
       if (prt%polarized) then
          write (u, "(I0,'/',I0,'|')", advance="no")  prt%pdg, prt%h
       else
          write (u, "(I0,'|')", advance="no") prt%pdg
       end if
    end select
    select case (prt%type)
    case (PRT_BEAM, PRT_INCOMING, PRT_OUTGOING, PRT_COMPOSITE)
       tmp = prt
       if (num_pac) then
          call pacify (tmp)
       end if
       arr = tmp%p
       write (u, "(" // FMT_14 // ",';'," // FMT_14 // ",','," // &
            FMT_14 // ",','," // FMT_14 // ")", advance="no") arr
       write (u, "('|'," // fmt // ")", advance="no") &
            tmp%p2
    end select
    if (allocated (prt%src)) then
       write (u, "('|')", advance="no")
       do i = 1, size (prt%src)
          write (u, "(1x,I0)", advance="no")  prt%src(i)
       end do
    end if
    write (u, "(A)")  ")"
  end subroutine prt_write

\eatline
\nwendcode{}\nwbegindocs{314}\nwdocspar
\subsubsection{Tools}
Two particles match if their {\tt{}src} arrays are the same.
\nwenddocs{}\nwbegincode{315}\moddef{Subevents: interfaces}\plusendmoddef
  interface operator(.match.)
     module procedure prt_match
  end interface

\eatline
\nwendcode{}\nwbegincode{316}\moddef{Subevents: procedures}\plusendmoddef
  elemental function prt_match (prt1, prt2) result (match)
    logical :: match
    type(prt_t), intent(in) :: prt1, prt2
    if (size (prt1%src) == size (prt2%src)) then
       match = all (prt1%src == prt2%src)
    else
       match = .false.
    end if
  end function prt_match

\eatline
\nwendcode{}\nwbegindocs{317}The combine operation makes a pseudoparticle whose momentum is the
result of adding (the momenta of) the pair of input particles.  We
trace the particles from which a particle is built by storing a
{\tt{}src} array.  Each particle entry in the {\tt{}src} list contains a
list of indices which indicates its building blocks.  The indices
refer to an original list of particles.  Index lists are sorted, and
they contain no element more than once.

We thus require that in a given pseudoparticle, each original particle
occurs at most once.

The result is intent(inout), so it will not be initialized when the
subroutine is entered.
\nwenddocs{}\nwbegincode{318}\moddef{Subevents: procedures}\plusendmoddef
  subroutine prt_combine (prt, prt_in1, prt_in2, ok)
    type(prt_t), intent(inout) :: prt
    type(prt_t), intent(in) :: prt_in1, prt_in2
    logical :: ok
    integer, dimension(:), allocatable :: src
    call combine_index_lists (src, prt_in1%src, prt_in2%src)
    ok = allocated (src)
    if (ok)  call prt_init_composite (prt, prt_in1%p + prt_in2%p, src)
  end subroutine prt_combine

\eatline
\nwendcode{}\nwbegindocs{319}This variant does not produce the combined particle, it just checks
whether the combination is valid (no common {\tt{}src} entry).
\nwenddocs{}\nwbegincode{320}\moddef{Subevents: public}\plusendmoddef
  public :: are_disjoint
\nwendcode{}\nwbegincode{321}\moddef{Subevents: procedures}\plusendmoddef
  function are_disjoint (prt_in1, prt_in2) result (flag)
    logical :: flag
    type(prt_t), intent(in) :: prt_in1, prt_in2
    flag = index_lists_are_disjoint (prt_in1%src, prt_in2%src)
  end function are_disjoint

\eatline
\nwendcode{}\nwbegindocs{322}{\tt{}src} Lists with length $>1$ are built by a {\tt{}combine} operation
which merges the lists in a sorted manner.  If the result would have a
duplicate entry, it is discarded, and the result is unallocated.
\nwenddocs{}\nwbegincode{323}\moddef{Subevents: procedures}\plusendmoddef
  subroutine combine_index_lists (res, src1, src2)
    integer, dimension(:), intent(in) :: src1, src2
    integer, dimension(:), allocatable :: res
    integer :: i1, i2, i
    allocate (res (size (src1) + size (src2)))
    if (size (src1) == 0) then
       res = src2
       return
    else if (size (src2) == 0) then
       res = src1
       return
    end if
    i1 = 1
    i2 = 1
    LOOP: do i = 1, size (res)
       if (src1(i1) < src2(i2)) then
          res(i) = src1(i1);  i1 = i1 + 1
          if (i1 > size (src1)) then
             res(i+1:) = src2(i2:)
             exit LOOP
          end if
       else if (src1(i1) > src2(i2)) then
          res(i) = src2(i2);  i2 = i2 + 1
          if (i2 > size (src2)) then
             res(i+1:) = src1(i1:)
             exit LOOP
          end if
       else
          deallocate (res)
          exit LOOP
       end if
    end do LOOP
  end subroutine combine_index_lists

\eatline
\nwendcode{}\nwbegindocs{324}This function is similar, but it does not actually merge the list,
it just checks whether they are disjoint (no common {\tt{}src} entry).
\nwenddocs{}\nwbegincode{325}\moddef{Subevents: procedures}\plusendmoddef
  function index_lists_are_disjoint (src1, src2) result (flag)
    logical :: flag
    integer, dimension(:), intent(in) :: src1, src2
    integer :: i1, i2, i
    flag = .true.
    i1 = 1
    i2 = 1
    LOOP: do i = 1, size (src1) + size (src2)
       if (src1(i1) < src2(i2)) then
          i1 = i1 + 1
          if (i1 > size (src1)) then
             exit LOOP
          end if
       else if (src1(i1) > src2(i2)) then
          i2 = i2 + 1
          if (i2 > size (src2)) then
             exit LOOP
          end if
       else
          flag = .false.
          exit LOOP
       end if
    end do LOOP
  end function index_lists_are_disjoint

\eatline
\nwendcode{}\nwbegindocs{326}\nwdocspar
\subsection{subevents}
Particles are collected in subevents.  This type is implemented as a
dynamically allocated array, which need not be completely filled.  The
value {\tt{}n{\char95}active} determines the number of meaningful entries.

\subsubsection{Type definition}
\nwenddocs{}\nwbegincode{327}\moddef{Subevents: public}\plusendmoddef
  public :: subevt_t
\nwendcode{}\nwbegincode{328}\moddef{Subevents: types}\plusendmoddef
  type :: subevt_t
     private
     integer :: n_active = 0
     type(prt_t), dimension(:), allocatable :: prt
   contains
   \LA{}Subevents: subevt: TBP\RA{}
  end type subevt_t

\eatline
\nwendcode{}\nwbegindocs{329}Initialize, allocating with size zero (default) or given size.  The
number of contained particles is set equal to the size.
\nwenddocs{}\nwbegincode{330}\moddef{Subevents: public}\plusendmoddef
  public :: subevt_init
\nwendcode{}\nwbegincode{331}\moddef{Subevents: procedures}\plusendmoddef
  subroutine subevt_init (subevt, n_active)
    type(subevt_t), intent(out) :: subevt
    integer, intent(in), optional :: n_active
    if (present (n_active))  subevt%n_active = n_active
    allocate (subevt%prt (subevt%n_active))
  end subroutine subevt_init

\eatline
\nwendcode{}\nwbegindocs{332}(Re-)allocate the subevent with some given size.  If the size
is greater than the previous one, do a real reallocation.  Otherwise,
just reset the recorded size. Contents are untouched, but become
invalid.
\nwenddocs{}\nwbegincode{333}\moddef{Subevents: public}\plusendmoddef
  public :: subevt_reset
\nwendcode{}\nwbegincode{334}\moddef{Subevents: procedures}\plusendmoddef
  subroutine subevt_reset (subevt, n_active)
    type(subevt_t), intent(inout) :: subevt
    integer, intent(in) :: n_active
    subevt%n_active = n_active
    if (subevt%n_active > size (subevt%prt)) then
       deallocate (subevt%prt)
       allocate (subevt%prt (subevt%n_active))
    end if
  end subroutine subevt_reset

\eatline
\nwendcode{}\nwbegindocs{335}Output.  No prefix for the headline 'subevt', because this will usually be
printed appending to a previous line.
\nwenddocs{}\nwbegincode{336}\moddef{Subevents: public}\plusendmoddef
  public :: subevt_write
\nwendcode{}\nwbegincode{337}\moddef{Subevents: subevt: TBP}\endmoddef
  procedure :: write => subevt_write
\nwendcode{}\nwbegincode{338}\moddef{Subevents: procedures}\plusendmoddef
  subroutine subevt_write (object, unit, prefix, pacified)
    class(subevt_t), intent(in) :: object
    integer, intent(in), optional :: unit
    character(*), intent(in), optional :: prefix
    logical, intent(in), optional :: pacified
    integer :: u, i
    u = given_output_unit (unit);  if (u < 0)  return
    write (u, "(1x,A)") "subevent:"
    do i = 1, object%n_active
       if (present (prefix))  write (u, "(A)", advance="no") prefix
       write (u, "(1x,I0)", advance="no")  i
       call prt_write (object%prt(i), unit = unit, pacified = pacified)
    end do
  end subroutine subevt_write

\eatline
\nwendcode{}\nwbegindocs{339}Defined assignment: transfer only meaningful entries.  This is a
deep copy (as would be default assignment).
\nwenddocs{}\nwbegincode{340}\moddef{Subevents: interfaces}\plusendmoddef
  interface assignment(=)
     module procedure subevt_assign
  end interface

\eatline
\nwendcode{}\nwbegincode{341}\moddef{Subevents: procedures}\plusendmoddef
  subroutine subevt_assign (subevt, subevt_in)
    type(subevt_t), intent(inout) :: subevt
    type(subevt_t), intent(in) :: subevt_in
    if (.not. allocated (subevt%prt)) then
       call subevt_init (subevt, subevt_in%n_active)
    else
       call subevt_reset (subevt, subevt_in%n_active)
    end if
    subevt%prt(:subevt%n_active) = subevt_in%prt(:subevt%n_active)
  end subroutine subevt_assign

\eatline
\nwendcode{}\nwbegindocs{342}\nwdocspar
\subsubsection{Fill contents}
Store incoming/outgoing particles which are completely defined.
\nwenddocs{}\nwbegincode{343}\moddef{Subevents: public}\plusendmoddef
  public :: subevt_set_beam
  public :: subevt_set_incoming
  public :: subevt_set_outgoing
  public :: subevt_set_composite
\nwendcode{}\nwbegincode{344}\moddef{Subevents: procedures}\plusendmoddef
  subroutine subevt_set_beam (subevt, i, pdg, p, p2, src)
    type(subevt_t), intent(inout) :: subevt
    integer, intent(in) :: i
    integer, intent(in) :: pdg
    type(vector4_t), intent(in) :: p
    real(default), intent(in) :: p2
    integer, dimension(:), intent(in), optional :: src
    if (present (src)) then
       call prt_init_beam (subevt%prt(i), pdg, p, p2, src)
    else
       call prt_init_beam (subevt%prt(i), pdg, p, p2, [i])
    end if
  end subroutine subevt_set_beam

  subroutine subevt_set_incoming (subevt, i, pdg, p, p2, src)
    type(subevt_t), intent(inout) :: subevt
    integer, intent(in) :: i
    integer, intent(in) :: pdg
    type(vector4_t), intent(in) :: p
    real(default), intent(in) :: p2
    integer, dimension(:), intent(in), optional :: src
    if (present (src)) then
       call prt_init_incoming (subevt%prt(i), pdg, p, p2, src)
    else
       call prt_init_incoming (subevt%prt(i), pdg, p, p2, [i])
    end if
  end subroutine subevt_set_incoming

  subroutine subevt_set_outgoing (subevt, i, pdg, p, p2, src)
    type(subevt_t), intent(inout) :: subevt
    integer, intent(in) :: i
    integer, intent(in) :: pdg
    type(vector4_t), intent(in) :: p
    real(default), intent(in) :: p2
    integer, dimension(:), intent(in), optional :: src
    if (present (src)) then
       call prt_init_outgoing (subevt%prt(i), pdg, p, p2, src)
    else
       call prt_init_outgoing (subevt%prt(i), pdg, p, p2, [i])
    end if
  end subroutine subevt_set_outgoing

  subroutine subevt_set_composite (subevt, i, p, src)
    type(subevt_t), intent(inout) :: subevt
    integer, intent(in) :: i
    type(vector4_t), intent(in) :: p
    integer, dimension(:), intent(in) :: src
    call prt_init_composite (subevt%prt(i), p, src)
  end subroutine subevt_set_composite

\eatline
\nwendcode{}\nwbegindocs{345}Separately assign flavors, simultaneously for all incoming/outgoing
particles.
\nwenddocs{}\nwbegincode{346}\moddef{Subevents: public}\plusendmoddef
  public :: subevt_set_pdg_beam
  public :: subevt_set_pdg_incoming
  public :: subevt_set_pdg_outgoing
\nwendcode{}\nwbegincode{347}\moddef{Subevents: procedures}\plusendmoddef
  subroutine subevt_set_pdg_beam (subevt, pdg)
    type(subevt_t), intent(inout) :: subevt
    integer, dimension(:), intent(in) :: pdg
    integer :: i, j
    j = 1
    do i = 1, subevt%n_active
       if (subevt%prt(i)%type == PRT_BEAM) then
          call prt_set_pdg (subevt%prt(i), pdg(j))
          j = j + 1
          if (j > size (pdg))  exit
       end if
    end do
  end subroutine subevt_set_pdg_beam

  subroutine subevt_set_pdg_incoming (subevt, pdg)
    type(subevt_t), intent(inout) :: subevt
    integer, dimension(:), intent(in) :: pdg
    integer :: i, j
    j = 1
    do i = 1, subevt%n_active
       if (subevt%prt(i)%type == PRT_INCOMING) then
          call prt_set_pdg (subevt%prt(i), pdg(j))
          j = j + 1
          if (j > size (pdg))  exit
       end if
    end do
  end subroutine subevt_set_pdg_incoming

  subroutine subevt_set_pdg_outgoing (subevt, pdg)
    type(subevt_t), intent(inout) :: subevt
    integer, dimension(:), intent(in) :: pdg
    integer :: i, j
    j = 1
    do i = 1, subevt%n_active
       if (subevt%prt(i)%type == PRT_OUTGOING) then
          call prt_set_pdg (subevt%prt(i), pdg(j))
          j = j + 1
          if (j > size (pdg))  exit
       end if
    end do
  end subroutine subevt_set_pdg_outgoing

\eatline
\eatline
\eatline
\nwendcode{}\nwbegindocs{348}Separately assign momenta, simultaneously for all incoming/outgoing
particles.
\nwenddocs{}\nwbegincode{349}\moddef{Subevents: public}\plusendmoddef
  public :: subevt_set_p_beam
  public :: subevt_set_p_incoming
  public :: subevt_set_p_outgoing
\nwendcode{}\nwbegincode{350}\moddef{Subevents: procedures}\plusendmoddef
  subroutine subevt_set_p_beam (subevt, p)
    type(subevt_t), intent(inout) :: subevt
    type(vector4_t), dimension(:), intent(in) :: p
    integer :: i, j
    j = 1
    do i = 1, subevt%n_active
       if (subevt%prt(i)%type == PRT_BEAM) then
          call prt_set_p (subevt%prt(i), p(j))
          j = j + 1
          if (j > size (p))  exit
       end if
    end do
  end subroutine subevt_set_p_beam

  subroutine subevt_set_p_incoming (subevt, p)
    type(subevt_t), intent(inout) :: subevt
    type(vector4_t), dimension(:), intent(in) :: p
    integer :: i, j
    j = 1
    do i = 1, subevt%n_active
       if (subevt%prt(i)%type == PRT_INCOMING) then
          call prt_set_p (subevt%prt(i), p(j))
          j = j + 1
          if (j > size (p))  exit
       end if
    end do
  end subroutine subevt_set_p_incoming

  subroutine subevt_set_p_outgoing (subevt, p)
    type(subevt_t), intent(inout) :: subevt
    type(vector4_t), dimension(:), intent(in) :: p
    integer :: i, j
    j = 1
    do i = 1, subevt%n_active
       if (subevt%prt(i)%type == PRT_OUTGOING) then
          call prt_set_p (subevt%prt(i), p(j))
          j = j + 1
          if (j > size (p))  exit
       end if
    end do
  end subroutine subevt_set_p_outgoing

\eatline
\eatline
\eatline
\nwendcode{}\nwbegindocs{351}Separately assign the squared invariant mass, simultaneously for all
incoming/outgoing particles.
\nwenddocs{}\nwbegincode{352}\moddef{Subevents: public}\plusendmoddef
  public :: subevt_set_p2_beam
  public :: subevt_set_p2_incoming
  public :: subevt_set_p2_outgoing
\nwendcode{}\nwbegincode{353}\moddef{Subevents: procedures}\plusendmoddef
  subroutine subevt_set_p2_beam (subevt, p2)
    type(subevt_t), intent(inout) :: subevt
    real(default), dimension(:), intent(in) :: p2
    integer :: i, j
    j = 1
    do i = 1, subevt%n_active
       if (subevt%prt(i)%type == PRT_BEAM) then
          call prt_set_p2 (subevt%prt(i), p2(j))
          j = j + 1
          if (j > size (p2))  exit
       end if
    end do
  end subroutine subevt_set_p2_beam

  subroutine subevt_set_p2_incoming (subevt, p2)
    type(subevt_t), intent(inout) :: subevt
    real(default), dimension(:), intent(in) :: p2
    integer :: i, j
    j = 1
    do i = 1, subevt%n_active
       if (subevt%prt(i)%type == PRT_INCOMING) then
          call prt_set_p2 (subevt%prt(i), p2(j))
          j = j + 1
          if (j > size (p2))  exit
       end if
    end do
  end subroutine subevt_set_p2_incoming

  subroutine subevt_set_p2_outgoing (subevt, p2)
    type(subevt_t), intent(inout) :: subevt
    real(default), dimension(:), intent(in) :: p2
    integer :: i, j
    j = 1
    do i = 1, subevt%n_active
       if (subevt%prt(i)%type == PRT_OUTGOING) then
          call prt_set_p2 (subevt%prt(i), p2(j))
          j = j + 1
          if (j > size (p2))  exit
       end if
    end do
  end subroutine subevt_set_p2_outgoing

\eatline
\eatline
\eatline
\nwendcode{}\nwbegindocs{354}Set polarization for an entry
\nwenddocs{}\nwbegincode{355}\moddef{Subevents: public}\plusendmoddef
  public :: subevt_polarize
\nwendcode{}\nwbegincode{356}\moddef{Subevents: procedures}\plusendmoddef
  subroutine subevt_polarize (subevt, i, h)
    type(subevt_t), intent(inout) :: subevt
    integer, intent(in) :: i, h
    call prt_polarize (subevt%prt(i), h)
  end subroutine subevt_polarize

\eatline
\nwendcode{}\nwbegindocs{357}\nwdocspar
\subsubsection{Accessing contents}
Return true if the subevent has entries.
\nwenddocs{}\nwbegincode{358}\moddef{Subevents: public}\plusendmoddef
  public :: subevt_is_nonempty
\nwendcode{}\nwbegincode{359}\moddef{Subevents: procedures}\plusendmoddef
  function subevt_is_nonempty (subevt) result (flag)
    logical :: flag
    type(subevt_t), intent(in) :: subevt
    flag = subevt%n_active /= 0
  end function subevt_is_nonempty

\eatline
\nwendcode{}\nwbegindocs{360}Return the number of entries
\nwenddocs{}\nwbegincode{361}\moddef{Subevents: public}\plusendmoddef
  public :: subevt_get_length
\nwendcode{}\nwbegincode{362}\moddef{Subevents: procedures}\plusendmoddef
  function subevt_get_length (subevt) result (length)
    integer :: length
    type(subevt_t), intent(in) :: subevt
    length = subevt%n_active
  end function subevt_get_length

\eatline
\nwendcode{}\nwbegindocs{363}Return a specific particle.  The index is not checked for validity.
\nwenddocs{}\nwbegincode{364}\moddef{Subevents: public}\plusendmoddef
  public :: subevt_get_prt
\nwendcode{}\nwbegincode{365}\moddef{Subevents: procedures}\plusendmoddef
  function subevt_get_prt (subevt, i) result (prt)
    type(prt_t) :: prt
    type(subevt_t), intent(in) :: subevt
    integer, intent(in) :: i
    prt = subevt%prt(i)
  end function subevt_get_prt

\eatline
\nwendcode{}\nwbegindocs{366}Return the partonic energy squared.  We take the particles with flag
{\tt{}PRT{\char95}INCOMING} and compute their total invariant mass.
\nwenddocs{}\nwbegincode{367}\moddef{Subevents: public}\plusendmoddef
  public :: subevt_get_sqrts_hat
\nwendcode{}\nwbegincode{368}\moddef{Subevents: procedures}\plusendmoddef
  function subevt_get_sqrts_hat (subevt) result (sqrts_hat)
    type(subevt_t), intent(in) :: subevt
    real(default) :: sqrts_hat
    type(vector4_t) :: p
    integer :: i
    do i = 1, subevt%n_active
       if (subevt%prt(i)%type == PRT_INCOMING) then
          p = p + prt_get_momentum (subevt%prt(i))
       end if
    end do
    sqrts_hat = p ** 1
  end function subevt_get_sqrts_hat
    
\eatline
\nwendcode{}\nwbegindocs{369}Return the number of incoming (outgoing) particles, respectively.
Beam particles or composites are not counted.
\nwenddocs{}\nwbegincode{370}\moddef{Subevents: public}\plusendmoddef
  public :: subevt_get_n_in
  public :: subevt_get_n_out
\nwendcode{}\nwbegincode{371}\moddef{Subevents: procedures}\plusendmoddef
  function subevt_get_n_in (subevt) result (n_in)
    type(subevt_t), intent(in) :: subevt
    integer :: n_in
    n_in = count (subevt%prt(:subevt%n_active)%type == PRT_INCOMING)
  end function subevt_get_n_in
    
  function subevt_get_n_out (subevt) result (n_out)
    type(subevt_t), intent(in) :: subevt
    integer :: n_out
    n_out = count (subevt%prt(:subevt%n_active)%type == PRT_OUTGOING)
  end function subevt_get_n_out
    
\eatline
\eatline
\nwendcode{}\nwbegindocs{372}\nwdocspar
\nwenddocs{}\nwbegincode{373}\moddef{Subevents: interfaces}\plusendmoddef
  interface c_prt
     module procedure c_prt_from_subevt
     module procedure c_prt_array_from_subevt
  end interface

\eatline
\nwendcode{}\nwbegincode{374}\moddef{Subevents: procedures}\plusendmoddef
  function c_prt_from_subevt (subevt, i) result (c_prt)
    type(c_prt_t) :: c_prt
    type(subevt_t), intent(in) :: subevt
    integer, intent(in) :: i
    c_prt = c_prt_from_prt (subevt%prt(i))
  end function c_prt_from_subevt

  function c_prt_array_from_subevt (subevt) result (c_prt_array)
    type(subevt_t), intent(in) :: subevt
    type(c_prt_t), dimension(subevt%n_active) :: c_prt_array
    c_prt_array = c_prt_from_prt (subevt%prt(1:subevt%n_active))
  end function c_prt_array_from_subevt

\eatline
\eatline
\nwendcode{}\nwbegindocs{375}\nwdocspar
\subsubsection{Operations with subevents}
The join operation joins two subevents.  When appending the
elements of the second list, we check for each particle whether it is
already in the first list.  If yes, it is discarded.  The result list
should be initialized already.

If a mask is present, it refers to the second subevent.
Particles where the mask is not set are discarded.
\nwenddocs{}\nwbegincode{376}\moddef{Subevents: public}\plusendmoddef
  public :: subevt_join
\nwendcode{}\nwbegincode{377}\moddef{Subevents: procedures}\plusendmoddef
  subroutine subevt_join (subevt, pl1, pl2, mask2)
    type(subevt_t), intent(inout) :: subevt
    type(subevt_t), intent(in) :: pl1, pl2
    logical, dimension(:), intent(in), optional :: mask2
    integer :: n1, n2, i, n
    n1 = pl1%n_active
    n2 = pl2%n_active
    call subevt_reset (subevt, n1 + n2)
    subevt%prt(:n1)   = pl1%prt(:n1)
    n = n1
    if (present (mask2)) then
       do i = 1, pl2%n_active
          if (mask2(i)) then
             if (disjoint (i)) then
                n = n + 1
                subevt%prt(n) = pl2%prt(i)
             end if
          end if
       end do
    else
       do i = 1, pl2%n_active
          if (disjoint (i)) then
             n = n + 1
             subevt%prt(n) = pl2%prt(i)
          end if
       end do
    end if
    subevt%n_active = n
  contains
    function disjoint (i) result (flag)
      integer, intent(in) :: i
      logical :: flag
      integer :: j
      do j = 1, pl1%n_active
         if (.not. are_disjoint (pl1%prt(j), pl2%prt(i))) then
            flag = .false.
            return
         end if
      end do
      flag = .true.
    end function disjoint
  end subroutine subevt_join

\eatline
\nwendcode{}\nwbegindocs{378}The combine operation makes a subevent whose entries are the
result of adding (the momenta of) each pair of particles in the input
lists.  We trace the particles from which a particles is built by
storing a {\tt{}src} array.  Each particle entry in the {\tt{}src} list
contains a list of indices which indicates its building blocks.  The
indices refer to an original list of particles.  Index lists are sorted,
and they contain no element more than once.

We thus require that in a given pseudoparticle, each original particle
occurs at most once.
\nwenddocs{}\nwbegincode{379}\moddef{Subevents: public}\plusendmoddef
  public :: subevt_combine
\nwendcode{}\nwbegincode{380}\moddef{Subevents: procedures}\plusendmoddef
  subroutine subevt_combine (subevt, pl1, pl2, mask12)
    type(subevt_t), intent(inout) :: subevt
    type(subevt_t), intent(in) :: pl1, pl2
    logical, dimension(:,:), intent(in), optional :: mask12
    integer :: n1, n2, i1, i2, n, j
    logical :: ok
    n1 = pl1%n_active
    n2 = pl2%n_active
    call subevt_reset (subevt, n1 * n2)
    n = 1
    do i1 = 1, n1
       do i2 = 1, n2
          if (present (mask12)) then
             ok = mask12(i1,i2)
          else
             ok = .true.
          end if
          if (ok)  call prt_combine &
               (subevt%prt(n), pl1%prt(i1), pl2%prt(i2), ok)
          if (ok) then
             CHECK_DOUBLES: do j = 1, n - 1
                if (subevt%prt(n) .match. subevt%prt(j)) then
                   ok = .false.;  exit CHECK_DOUBLES
                end if
             end do CHECK_DOUBLES
             if (ok)  n = n + 1
          end if
       end do
    end do
    subevt%n_active = n - 1
  end subroutine subevt_combine

\eatline
\nwendcode{}\nwbegindocs{381}The collect operation makes a single-entry subevent which
results from combining (the momenta of) all particles in the input
list.  As above, the result does not contain an original particle more
than once; this is checked for each particle when it is collected.
Furthermore, each entry has a mask; where the mask is false, the entry
is dropped.

(Thus, if the input particles are already composite, there is some
chance that the result depends on the order of the input list and is
not as expected.  This situation should be avoided.)
\nwenddocs{}\nwbegincode{382}\moddef{Subevents: public}\plusendmoddef
  public :: subevt_collect
\nwendcode{}\nwbegincode{383}\moddef{Subevents: procedures}\plusendmoddef
  subroutine subevt_collect (subevt, pl1, mask1)
    type(subevt_t), intent(inout) :: subevt
    type(subevt_t), intent(in) :: pl1
    logical, dimension(:), intent(in) :: mask1
    type(prt_t) :: prt
    integer :: i
    logical :: ok
    call subevt_reset (subevt, 1)
    subevt%n_active = 0
    do i = 1, pl1%n_active
       if (mask1(i)) then
          if (subevt%n_active == 0) then
             subevt%n_active = 1
             subevt%prt(1) = pl1%prt(i)
          else
             call prt_combine (prt, subevt%prt(1), pl1%prt(i), ok)
             if (ok)  subevt%prt(1) = prt
          end if
       end if
    end do
  end subroutine subevt_collect

\eatline
\nwendcode{}\nwbegindocs{384}The cluster operation is similar to {\tt{}collect}, but applies a jet
algorithm.  The result is a subevent consisting of jets and, possibly,
unclustered extra particles.  As above, the result does not contain an
original particle more than once; this is checked for each particle when it is
collected.  Furthermore, each entry has a mask; where the mask is false, the
entry is dropped.

The algorithm: first determine the (pseudo)particles that participate in the
clustering.  They should not overlap, and the mask entry must be set.  We then
cluster the particles, using the given jet definition.  The result particles are
retrieved from the cluster sequence.  We still have to determine the source
indices for each jet: for each input particle, we get the jet index.
Accumulating the source entries for all particles that are part of a given
jet, we derive the jet source entries.  Finally, we delete the C structures
that have been constructed by FastJet and its interface.
\nwenddocs{}\nwbegincode{385}\moddef{Subevents: public}\plusendmoddef
  public :: subevt_cluster
\nwendcode{}\nwbegincode{386}\moddef{Subevents: procedures}\plusendmoddef
  subroutine subevt_cluster (subevt, pl1, mask1, jet_def)
    type(subevt_t), intent(inout) :: subevt
    type(subevt_t), intent(in) :: pl1
    logical, dimension(:), intent(in) :: mask1
    type(jet_definition_t), intent(in) :: jet_def
    integer, dimension(:), allocatable :: src, src_tmp
    integer, dimension(:), allocatable :: map, jet_idx
    type(pseudojet_t), dimension(:), allocatable :: jet_in, jet_out
    type(pseudojet_vector_t) :: jv_in, jv_out
    type(cluster_sequence_t) :: cs
    integer :: i, j, k, n_src, n_active
    n_active = 0
    allocate (map (pl1%n_active), source = 0)
    allocate (src (0))
    do i = 1, pl1%n_active
       if (mask1(i)) then
          call combine_index_lists (src_tmp, src, pl1%prt(i)%src)
          if (allocated (src_tmp)) then
             call move_alloc (from=src_tmp, to=src)
             n_active = n_active + 1
             map(n_active) = i
          end if
       end if
    end do
    allocate (jet_in (count (map /= 0)))
    do i = 1, size (jet_in)
       call jet_in(i)%init (prt_get_momentum (pl1%prt(map(i))))
    end do
    call jv_in%init (jet_in)
    call cs%init (jv_in, jet_def)
    jv_out = cs%inclusive_jets ()
    allocate (jet_idx (size (jet_in)))
    call cs%assign_jet_indices (jv_out, jet_idx)
    allocate (jet_out (jv_out%size ()))
    jet_out = jv_out
    call subevt_reset (subevt, size (jet_out))
    do i = 1, size (jet_out)
       src = 0
       n_src = 0
       do j = 1, size (jet_idx)
          if (jet_idx(j) == i) then
             associate (prt => pl1%prt(map(j)))
               do k = 1, size (prt%src)
                  src(n_src + k) = prt%src(k)
               end do
               n_src = n_src + size (prt%src)
             end associate
          end if
       end do
       call prt_init_pseudojet (subevt%prt(i), jet_out(i), src(:n_src))
    end do
    do i = 1, size (jet_out)
       call jet_out(i)%final ()
    end do
    call jv_out%final ()
    call cs%final ()
    call jv_in%final ()
    do i = 1, size (jet_in)
       call jet_in(i)%final ()
    end do
  end subroutine subevt_cluster

\eatline
\nwendcode{}\nwbegindocs{387}Return a list of all particles for which the mask is true.
\nwenddocs{}\nwbegincode{388}\moddef{Subevents: public}\plusendmoddef
  public :: subevt_select
\nwendcode{}\nwbegincode{389}\moddef{Subevents: procedures}\plusendmoddef
  subroutine subevt_select (subevt, pl, mask1)
    type(subevt_t), intent(inout) :: subevt
    type(subevt_t), intent(in) :: pl
    logical, dimension(:), intent(in) :: mask1
    integer :: i, n
    call subevt_reset (subevt, pl%n_active)
    n = 0
    do i = 1, pl%n_active
       if (mask1(i)) then
          n = n + 1
          subevt%prt(n) = pl%prt(i)
       end if
    end do
    subevt%n_active = n
  end subroutine subevt_select

\eatline
\nwendcode{}\nwbegindocs{390}Return a subevent which consists of the single particle with
specified {\tt{}index}.  If {\tt{}index} is negative, count from the end.
If it is out of bounds, return an empty list.
\nwenddocs{}\nwbegincode{391}\moddef{Subevents: public}\plusendmoddef
  public :: subevt_extract
\nwendcode{}\nwbegincode{392}\moddef{Subevents: procedures}\plusendmoddef
  subroutine subevt_extract (subevt, pl, index)
    type(subevt_t), intent(inout) :: subevt
    type(subevt_t), intent(in) :: pl
    integer, intent(in) :: index
    if (index > 0) then
       if (index <= pl%n_active) then
          call subevt_reset (subevt, 1)
          subevt%prt(1) = pl%prt(index)
       else
          call subevt_reset (subevt, 0)
       end if
    else if (index < 0) then
       if (abs (index) <= pl%n_active) then
          call subevt_reset (subevt, 1)
          subevt%prt(1) = pl%prt(pl%n_active + 1 + index)
       else
          call subevt_reset (subevt, 0)
       end if
    else
       call subevt_reset (subevt, 0)
    end if
  end subroutine subevt_extract

\eatline
\nwendcode{}\nwbegindocs{393}Return the list of particles sorted according to increasing values
of the provided integer or real array.  If no array is given, sort by
PDG value.
\nwenddocs{}\nwbegincode{394}\moddef{Subevents: public}\plusendmoddef
  public :: subevt_sort
\nwendcode{}\nwbegincode{395}\moddef{Subevents: interfaces}\plusendmoddef
  interface subevt_sort
     module procedure subevt_sort_pdg
     module procedure subevt_sort_int
     module procedure subevt_sort_real
  end interface

\nwendcode{}\nwbegincode{396}\moddef{Subevents: procedures}\plusendmoddef
  subroutine subevt_sort_pdg (subevt, pl)
    type(subevt_t), intent(inout) :: subevt
    type(subevt_t), intent(in) :: pl
    integer :: n
    n = subevt%n_active
    call subevt_sort_int (subevt, pl, abs (3 * subevt%prt(:n)%pdg - 1))
  end subroutine subevt_sort_pdg

  subroutine subevt_sort_int (subevt, pl, ival)
    type(subevt_t), intent(inout) :: subevt
    type(subevt_t), intent(in) :: pl
    integer, dimension(:), intent(in) :: ival
    call subevt_reset (subevt, pl%n_active)
    subevt%n_active = pl%n_active
    subevt%prt = pl%prt( order (ival) )
  end subroutine subevt_sort_int

  subroutine subevt_sort_real (subevt, pl, rval)
    type(subevt_t), intent(inout) :: subevt
    type(subevt_t), intent(in) :: pl
    real(default), dimension(:), intent(in) :: rval
    call subevt_reset (subevt, pl%n_active)
    subevt%n_active = pl%n_active
    subevt%prt = pl%prt( order (rval) )
  end subroutine subevt_sort_real

\eatline
\nwendcode{}\nwbegindocs{397}Return the list of particles which have any of the specified PDG
codes (and optionally particle type: beam, incoming, outgoing).

The {\tt{}pack} command was buggy in some gfortran versions, therefore it
is unrolled.  The unrolled version may be more efficient, actually.
\nwenddocs{}\nwbegincode{398}\moddef{Subevents: public}\plusendmoddef
  public :: subevt_select_pdg_code
\nwendcode{}\nwbegincode{399}\moddef{Subevents: procedures}\plusendmoddef
  subroutine subevt_select_pdg_code (subevt, aval, subevt_in, prt_type)
    type(subevt_t), intent(inout) :: subevt
    type(pdg_array_t), intent(in) :: aval
    type(subevt_t), intent(in) :: subevt_in
    integer, intent(in), optional :: prt_type
    integer :: n_active, n_match
    logical, dimension(:), allocatable :: mask
    integer :: i, j
    n_active = subevt_in%n_active
    allocate (mask (n_active))
    forall (i = 1:n_active) &
         mask(i) = aval .match. subevt_in%prt(i)%pdg
    if (present (prt_type)) &
         mask = mask .and. subevt_in%prt(:n_active)%type == prt_type
    n_match = count (mask)
    call subevt_reset (subevt, n_match)
    !!! !!! !!! Workaround for gfortran compiler bug
    ! subevt%prt(:n_match) = pack (subevt_in%prt(:n_active), mask)
    j = 0
    do i = 1, n_active
       if (mask(i)) then
          j = j + 1
          subevt%prt(j) = subevt_in%prt(i)
       end if
    end do
  end subroutine subevt_select_pdg_code

\eatline
\nwendcode{}\nwbegindocs{400}\nwdocspar
\subsection{Eliminate numerical noise}
This is useful for testing purposes: set entries to zero that are smaller in
absolute values than a given tolerance parameter.

Note: instead of setting the tolerance in terms of EPSILON
(kind-dependent), we fix it to $10^{-16}$, which is the typical value
for double precision.  The reason is that there are situations where
intermediate representations (external libraries, files) are limited
to double precision, even if the main program uses higher precision.
\nwenddocs{}\nwbegincode{401}\moddef{Subevents: public}\plusendmoddef
  public :: pacify
\nwendcode{}\nwbegincode{402}\moddef{Subevents: interfaces}\plusendmoddef
  interface pacify
     module procedure pacify_prt
     module procedure pacify_subevt
  end interface pacify
  
\eatline
\nwendcode{}\nwbegincode{403}\moddef{Subevents: procedures}\plusendmoddef
  subroutine pacify_prt (prt)
    class(prt_t), intent(inout) :: prt
    real(default) :: e
    e = 1e-13_default * energy (prt%p)
    call pacify (prt%p, e)
    call pacify (prt%p2, 1e3 * e)
  end subroutine pacify_prt
  
  subroutine pacify_subevt (subevt)
    class(subevt_t), intent(inout) :: subevt
    integer :: i
    do i = 1, subevt%n_active
       call pacify (subevt%prt(i))
    end do
  end subroutine pacify_subevt

\eatline
\eatline
\nwendcode{}\nwbegindocs{404}\nwdocspar
\clearpage
%%%%%%%%%%%%%%%%%%%%%%%%%%%%%%%%%%%%%%%%%%%%%%%%%%%%%%%%%%%%%%%%%%%%%%% 
\section{Analysis tools}

This module defines structures useful for data analysis.  These
include observables, histograms, and plots.

Observables are quantities that are calculated and summed up event by
event.  At the end, one can compute the average and error.

Histograms have their bins in addition to the observable properties.
Histograms are usually written out in tables and displayed
graphically.

In plots, each record creates its own entry in a table.  This can be
used for scatter plots if called event by event, or for plotting
dependencies on parameters if called once per integration run.

Graphs are container for histograms and plots, which carry their own graphics
options.

The type layout is still somewhat obfuscated.  This would become much simpler
if type extension could be used.
\nwenddocs{}\nwbegincode{405}\moddef{\code{}analysis.f90\edoc{}}\endmoddef
\LA{}File header\RA{}

module analysis

\LA{}Use kinds\RA{}
\LA{}Use strings\RA{}
  use io_units
  use format_utils, only: quote_underscore, tex_format
  use format_defs, only: FMT_19
  use unit_tests
  use diagnostics
  use os_interface
  use ifiles

\LA{}Standard module head\RA{}

\LA{}Analysis: public\RA{}

\LA{}Analysis: parameters\RA{}

\LA{}Analysis: types\RA{}

\LA{}Analysis: interfaces\RA{}

\LA{}Analysis: variables\RA{}

contains

\LA{}Analysis: procedures\RA{}

\LA{}Analysis: tests\RA{}

end module analysis
\eatline
\nwendcode{}\nwbegindocs{406}\nwdocspar
\subsection{Output formats}
These formats share a common field width (alignment).
\nwenddocs{}\nwbegincode{407}\moddef{Analysis: parameters}\endmoddef
  character(*), parameter, public :: HISTOGRAM_HEAD_FORMAT = "1x,A15,3x"
  character(*), parameter, public :: HISTOGRAM_INTG_FORMAT = "3x,I9,3x"
  character(*), parameter, public :: HISTOGRAM_DATA_FORMAT = "ES19.12"

\eatline
\nwendcode{}\nwbegindocs{408}\nwdocspar
\subsection{Graph options}
These parameters are used for displaying data.  They apply to a whole graph,
which may contain more than one plot element.

The GAMELAN code chunks are part of both {\tt{}graph{\char95}options} and
{\tt{}drawing{\char95}options}.  The {\tt{}drawing{\char95}options} copy is used in histograms and
plots, also as graph elements.  The {\tt{}graph{\char95}options} copy is used for
{\tt{}graph} objects as a whole.  Both copies are usually identical.
\nwenddocs{}\nwbegincode{409}\moddef{Analysis: public}\endmoddef
  public :: graph_options_t
\nwendcode{}\nwbegincode{410}\moddef{Analysis: types}\endmoddef
  type :: graph_options_t
     private
     type(string_t) :: id
     type(string_t) :: title
     type(string_t) :: description
     type(string_t) :: x_label
     type(string_t) :: y_label
     integer :: width_mm = 130
     integer :: height_mm = 90
     logical :: x_log = .false.
     logical :: y_log = .false.
     real(default) :: x_min = 0
     real(default) :: x_max = 1
     real(default) :: y_min = 0
     real(default) :: y_max = 1
     logical :: x_min_set = .false.
     logical :: x_max_set = .false.
     logical :: y_min_set = .false.
     logical :: y_max_set = .false.
     type(string_t) :: gmlcode_bg
     type(string_t) :: gmlcode_fg
  end type graph_options_t

\eatline
\nwendcode{}\nwbegindocs{411}Initialize the record, all strings are empty.  The limits are undefined.
\nwenddocs{}\nwbegincode{412}\moddef{Analysis: public}\plusendmoddef
  public :: graph_options_init
\nwendcode{}\nwbegincode{413}\moddef{Analysis: procedures}\endmoddef
  subroutine graph_options_init (graph_options)
    type(graph_options_t), intent(out) :: graph_options
    graph_options%id = ""
    graph_options%title = ""
    graph_options%description = ""
    graph_options%x_label = ""
    graph_options%y_label = ""
    graph_options%gmlcode_bg = ""
    graph_options%gmlcode_fg = ""
  end subroutine graph_options_init

\eatline
\nwendcode{}\nwbegindocs{414}Set individual options.
\nwenddocs{}\nwbegincode{415}\moddef{Analysis: public}\plusendmoddef
  public :: graph_options_set
\nwendcode{}\nwbegincode{416}\moddef{Analysis: procedures}\plusendmoddef
  subroutine graph_options_set (graph_options, id, &
       title, description, x_label, y_label, width_mm, height_mm, &
       x_log, y_log, x_min, x_max, y_min, y_max, &
       gmlcode_bg, gmlcode_fg)
    type(graph_options_t), intent(inout) :: graph_options
    type(string_t), intent(in), optional :: id
    type(string_t), intent(in), optional :: title
    type(string_t), intent(in), optional :: description
    type(string_t), intent(in), optional :: x_label, y_label
    integer, intent(in), optional :: width_mm, height_mm
    logical, intent(in), optional :: x_log, y_log
    real(default), intent(in), optional :: x_min, x_max, y_min, y_max
    type(string_t), intent(in), optional :: gmlcode_bg, gmlcode_fg
    if (present (id))  graph_options%id = id
    if (present (title))  graph_options%title = title
    if (present (description))  graph_options%description = description
    if (present (x_label))  graph_options%x_label = x_label
    if (present (y_label))  graph_options%y_label = y_label
    if (present (width_mm))   graph_options%width_mm  = width_mm
    if (present (height_mm))  graph_options%height_mm = height_mm
    if (present (x_log))  graph_options%x_log = x_log
    if (present (y_log))  graph_options%y_log = y_log
    if (present (x_min))  graph_options%x_min = x_min
    if (present (x_max))  graph_options%x_max = x_max
    if (present (y_min))  graph_options%y_min = y_min
    if (present (y_max))  graph_options%y_max = y_max
    if (present (x_min))  graph_options%x_min_set = .true.
    if (present (x_max))  graph_options%x_max_set = .true.
    if (present (y_min))  graph_options%y_min_set = .true.
    if (present (y_max))  graph_options%y_max_set = .true.
    if (present (gmlcode_bg))  graph_options%gmlcode_bg = gmlcode_bg
    if (present (gmlcode_fg))  graph_options%gmlcode_fg = gmlcode_fg
  end subroutine graph_options_set

\eatline
\nwendcode{}\nwbegindocs{417}Write a simple account of all options.
\nwenddocs{}\nwbegincode{418}\moddef{Analysis: public}\plusendmoddef
  public :: graph_options_write
\nwendcode{}\nwbegincode{419}\moddef{Analysis: procedures}\plusendmoddef
  subroutine graph_options_write (gro, unit)
    type(graph_options_t), intent(in) :: gro
    integer, intent(in), optional :: unit
    integer :: u
    u = given_output_unit (unit)
1   format (A,1x,'"',A,'"')
2   format (A,1x,L1)
3   format (A,1x,ES19.12)
4   format (A,1x,I0)
5   format (A,1x,'[undefined]')
    write (u, 1)  "title       =", char (gro%title)
    write (u, 1)  "description =", char (gro%description)
    write (u, 1)  "x_label     =", char (gro%x_label)
    write (u, 1)  "y_label     =", char (gro%y_label)
    write (u, 2)  "x_log       =", gro%x_log
    write (u, 2)  "y_log       =", gro%y_log
    if (gro%x_min_set) then
       write (u, 3)  "x_min       =", gro%x_min
    else
       write (u, 5)  "x_min       ="
    end if
    if (gro%x_max_set) then
       write (u, 3)  "x_max       =", gro%x_max
    else
       write (u, 5)  "x_max       ="
    end if
    if (gro%y_min_set) then
       write (u, 3)  "y_min       =", gro%y_min
    else
       write (u, 5)  "y_min       ="
    end if
    if (gro%y_max_set) then
       write (u, 3)  "y_max       =", gro%y_max
    else
       write (u, 5)  "y_max       ="
    end if
    write (u, 4)  "width_mm    =", gro%width_mm
    write (u, 4)  "height_mm   =", gro%height_mm
    write (u, 1)  "gmlcode_bg  =", char (gro%gmlcode_bg)
    write (u, 1)  "gmlcode_fg  =", char (gro%gmlcode_fg)
  end subroutine graph_options_write
  
\eatline
\nwendcode{}\nwbegindocs{420}Write a \LaTeX\ header/footer for the analysis file.
\nwenddocs{}\nwbegincode{421}\moddef{Analysis: procedures}\plusendmoddef
  subroutine graph_options_write_tex_header (gro, unit)
    type(graph_options_t), intent(in) :: gro
    integer, intent(in), optional :: unit
    integer :: u
    u = given_output_unit (unit)
    if (gro%title /= "") then
       write (u, "(A)")
       write (u, "(A)")  "\\section\{" // char (gro%title) // "\}"
    else
       write (u, "(A)")  "\\section\{" // char (quote_underscore (gro%id)) // "\}"
    end if
    if (gro%description /= "") then
       write (u, "(A)")  char (gro%description)
       write (u, *)
       write (u, "(A)")  "\\vspace*\{\\baselineskip\}"
    end if
    write (u, "(A)")  "\\vspace*\{\\baselineskip\}"
    write (u, "(A)")  "\\unitlength 1mm"
    write (u, "(A,I0,',',I0,A)")  &
         "\\begin\{gmlgraph*\}(", &
         gro%width_mm, gro%height_mm, &
         ")[dat]"
  end subroutine graph_options_write_tex_header

  subroutine graph_options_write_tex_footer (gro, unit)
    type(graph_options_t), intent(in) :: gro
    integer, intent(in), optional :: unit
    integer :: u, width, height
    width = gro%width_mm - 10
    height = gro%height_mm - 10
    u = given_output_unit (unit)
    write (u, "(A)")  "  begingmleps ""Whizard-Logo.eps"";"
    write (u, "(A,I0,A,I0,A)")  &
         "    base := (", width, "*unitlength,", height, "*unitlength);"
    write (u, "(A)")  "    height := 9.6*unitlength;"
    write (u, "(A)")  "    width := 11.2*unitlength;"
    write (u, "(A)")  "  endgmleps;"
    write (u, "(A)")  "\\end\{gmlgraph*\}"
  end subroutine graph_options_write_tex_footer

\eatline
\eatline
\nwendcode{}\nwbegindocs{422}Return the analysis object ID.
\nwenddocs{}\nwbegincode{423}\moddef{Analysis: procedures}\plusendmoddef
  function graph_options_get_id (gro) result (id)
    type(string_t) :: id
    type(graph_options_t), intent(in) :: gro
    id = gro%id
  end function graph_options_get_id

\eatline
\nwendcode{}\nwbegindocs{424}Create an appropriate {\tt{}setup} command (linear/log).
\nwenddocs{}\nwbegincode{425}\moddef{Analysis: procedures}\plusendmoddef
  function graph_options_get_gml_setup (gro) result (cmd)
    type(string_t) :: cmd
    type(graph_options_t), intent(in) :: gro
    type(string_t) :: x_str, y_str
    if (gro%x_log) then
       x_str = "log"
    else
       x_str = "linear"
    end if
    if (gro%y_log) then
       y_str = "log"
    else
       y_str = "linear"
    end if
    cmd = "setup (" // x_str // ", " // y_str // ");"
  end function graph_options_get_gml_setup

\eatline
\nwendcode{}\nwbegindocs{426}Return the labels in GAMELAN form.
\nwenddocs{}\nwbegincode{427}\moddef{Analysis: procedures}\plusendmoddef
  function graph_options_get_gml_x_label (gro) result (cmd)
    type(string_t) :: cmd
    type(graph_options_t), intent(in) :: gro
    cmd = 'label.bot (<' // '<' // gro%x_label // '>' // '>, out);'
  end function graph_options_get_gml_x_label
    
  function graph_options_get_gml_y_label (gro) result (cmd)
    type(string_t) :: cmd
    type(graph_options_t), intent(in) :: gro
    cmd = 'label.ulft (<' // '<' // gro%y_label // '>' // '>, out);'
  end function graph_options_get_gml_y_label
    
\eatline
\eatline
\nwendcode{}\nwbegindocs{428}Create an appropriate {\tt{}graphrange} statement for the given graph options.
Where the graph options are not set, use the supplied arguments, if any,
otherwise set the undefined value.
\nwenddocs{}\nwbegincode{429}\moddef{Analysis: procedures}\plusendmoddef
  function graph_options_get_gml_graphrange &
       (gro, x_min, x_max, y_min, y_max) result (cmd)
    type(string_t) :: cmd
    type(graph_options_t), intent(in) :: gro
    real(default), intent(in), optional :: x_min, x_max, y_min, y_max
    type(string_t) :: x_min_str, x_max_str, y_min_str, y_max_str
    character(*), parameter :: fmt = "(ES15.8)"
    if (gro%x_min_set) then
       x_min_str = "#" // trim (adjustl (real2string (gro%x_min, fmt)))
    else if (present (x_min)) then
       x_min_str = "#" // trim (adjustl (real2string (x_min, fmt)))
    else
       x_min_str = "??"
    end if
    if (gro%x_max_set) then
       x_max_str = "#" // trim (adjustl (real2string (gro%x_max, fmt)))
    else if (present (x_max)) then
       x_max_str = "#" // trim (adjustl (real2string (x_max, fmt)))
    else
       x_max_str = "??"
    end if
    if (gro%y_min_set) then
       y_min_str = "#" // trim (adjustl (real2string (gro%y_min, fmt)))
    else if (present (y_min)) then
       y_min_str = "#" // trim (adjustl (real2string (y_min, fmt)))
    else
       y_min_str = "??"
    end if
    if (gro%y_max_set) then
       y_max_str = "#" // trim (adjustl (real2string (gro%y_max, fmt)))
    else if (present (y_max)) then
       y_max_str = "#" // trim (adjustl (real2string (y_max, fmt)))
    else
       y_max_str = "??"
    end if
    cmd = "graphrange (" // x_min_str // ", " // y_min_str // "), " &
         // "(" // x_max_str // ", " // y_max_str // ");"
  end function graph_options_get_gml_graphrange

\eatline
\nwendcode{}\nwbegindocs{430}Get extra GAMELAN code to be executed before and after the usual drawing
commands.
\nwenddocs{}\nwbegincode{431}\moddef{Analysis: procedures}\plusendmoddef
  function graph_options_get_gml_bg_command (gro) result (cmd)
    type(string_t) :: cmd
    type(graph_options_t), intent(in) :: gro
    cmd = gro%gmlcode_bg
  end function graph_options_get_gml_bg_command

  function graph_options_get_gml_fg_command (gro) result (cmd)
    type(string_t) :: cmd
    type(graph_options_t), intent(in) :: gro
    cmd = gro%gmlcode_fg
  end function graph_options_get_gml_fg_command

\eatline
\eatline
\nwendcode{}\nwbegindocs{432}Append the header for generic data output in ifile format.  We print only
labels, not graphics parameters.
\nwenddocs{}\nwbegincode{433}\moddef{Analysis: procedures}\plusendmoddef
  subroutine graph_options_get_header (pl, header, comment)
    type(graph_options_t), intent(in) :: pl
    type(ifile_t), intent(inout) :: header
    type(string_t), intent(in), optional :: comment
    type(string_t) :: c
    if (present (comment)) then
       c = comment
    else
       c = ""
    end if
    call ifile_append (header, &
         c // "ID: " // pl%id)
    call ifile_append (header, &
         c // "title: " // pl%title)
    call ifile_append (header, &
         c // "description: " // pl%description)
    call ifile_append (header, &
         c // "x axis label: " // pl%x_label)
    call ifile_append (header, &
         c // "y axis label: " // pl%y_label)
  end subroutine graph_options_get_header

\eatline
\nwendcode{}\nwbegindocs{434}\nwdocspar
\subsection{Drawing options}
These options apply to an individual graph element (histogram or plot).
\nwenddocs{}\nwbegincode{435}\moddef{Analysis: public}\plusendmoddef
  public :: drawing_options_t
\nwendcode{}\nwbegincode{436}\moddef{Analysis: types}\plusendmoddef
  type :: drawing_options_t
     type(string_t) :: dataset
     logical :: with_hbars = .false.
     logical :: with_base = .false.
     logical :: piecewise = .false.
     logical :: fill = .false.
     logical :: draw = .false.
     logical :: err = .false.
     logical :: symbols = .false.
     type(string_t) :: fill_options
     type(string_t) :: draw_options
     type(string_t) :: err_options
     type(string_t) :: symbol
     type(string_t) :: gmlcode_bg
     type(string_t) :: gmlcode_fg
  end type drawing_options_t

\eatline
\nwendcode{}\nwbegindocs{437}Write a simple account of all options.
\nwenddocs{}\nwbegincode{438}\moddef{Analysis: public}\plusendmoddef
  public :: drawing_options_write
\nwendcode{}\nwbegincode{439}\moddef{Analysis: procedures}\plusendmoddef
  subroutine drawing_options_write (dro, unit)
     type(drawing_options_t), intent(in) :: dro
    integer, intent(in), optional :: unit
    integer :: u
    u = given_output_unit (unit)
1   format (A,1x,'"',A,'"')
2   format (A,1x,L1)
    write (u, 2)  "with_hbars  =", dro%with_hbars
    write (u, 2)  "with_base   =", dro%with_base
    write (u, 2)  "piecewise   =", dro%piecewise
    write (u, 2)  "fill        =", dro%fill
    write (u, 2)  "draw        =", dro%draw
    write (u, 2)  "err         =", dro%err
    write (u, 2)  "symbols     =", dro%symbols
    write (u, 1)  "fill_options=", char (dro%fill_options)
    write (u, 1)  "draw_options=", char (dro%draw_options)
    write (u, 1)  "err_options =", char (dro%err_options)
    write (u, 1)  "symbol      =", char (dro%symbol)
    write (u, 1)  "gmlcode_bg  =", char (dro%gmlcode_bg)
    write (u, 1)  "gmlcode_fg  =", char (dro%gmlcode_fg)
  end subroutine drawing_options_write
  
\eatline
\nwendcode{}\nwbegindocs{440}Init with empty strings and default options, appropriate for either
histogram or plot.
\nwenddocs{}\nwbegincode{441}\moddef{Analysis: public}\plusendmoddef
  public :: drawing_options_init_histogram
  public :: drawing_options_init_plot
\nwendcode{}\nwbegincode{442}\moddef{Analysis: procedures}\plusendmoddef
  subroutine drawing_options_init_histogram (dro)
    type(drawing_options_t), intent(out) :: dro
    dro%dataset = "dat"
    dro%with_hbars = .true.
    dro%with_base = .true.
    dro%piecewise = .true.
    dro%fill = .true.
    dro%draw = .true.
    dro%fill_options = "withcolor col.default"
    dro%draw_options = ""
    dro%err_options = ""
    dro%symbol = "fshape(circle scaled 1mm)()"
    dro%gmlcode_bg = ""
    dro%gmlcode_fg = ""
  end subroutine drawing_options_init_histogram

  subroutine drawing_options_init_plot (dro)
    type(drawing_options_t), intent(out) :: dro
    dro%dataset = "dat"
    dro%draw = .true.
    dro%fill_options = "withcolor col.default"
    dro%draw_options = ""
    dro%err_options = ""
    dro%symbol = "fshape(circle scaled 1mm)()"
    dro%gmlcode_bg = ""
    dro%gmlcode_fg = ""
  end subroutine drawing_options_init_plot

\eatline
\eatline
\nwendcode{}\nwbegindocs{443}Set individual options.
\nwenddocs{}\nwbegincode{444}\moddef{Analysis: public}\plusendmoddef
  public :: drawing_options_set
\nwendcode{}\nwbegincode{445}\moddef{Analysis: procedures}\plusendmoddef
  subroutine drawing_options_set (dro, dataset, &
       with_hbars, with_base, piecewise, fill, draw, err, symbols, &
       fill_options, draw_options, err_options, symbol, &
       gmlcode_bg, gmlcode_fg)
    type(drawing_options_t), intent(inout) :: dro
    type(string_t), intent(in), optional :: dataset
    logical, intent(in), optional :: with_hbars, with_base, piecewise
    logical, intent(in), optional :: fill, draw, err, symbols
    type(string_t), intent(in), optional :: fill_options, draw_options
    type(string_t), intent(in), optional :: err_options, symbol
    type(string_t), intent(in), optional :: gmlcode_bg, gmlcode_fg
    if (present (dataset))  dro%dataset = dataset
    if (present (with_hbars))  dro%with_hbars = with_hbars
    if (present (with_base))  dro%with_base = with_base
    if (present (piecewise))  dro%piecewise = piecewise
    if (present (fill))  dro%fill = fill
    if (present (draw))  dro%draw = draw
    if (present (err))  dro%err = err
    if (present (symbols))  dro%symbols = symbols
    if (present (fill_options))  dro%fill_options = fill_options
    if (present (draw_options))  dro%draw_options = draw_options
    if (present (err_options))  dro%err_options = err_options
    if (present (symbol))  dro%symbol = symbol
    if (present (gmlcode_bg))  dro%gmlcode_bg = gmlcode_bg
    if (present (gmlcode_fg))  dro%gmlcode_fg = gmlcode_fg
  end subroutine drawing_options_set

\eatline
\nwendcode{}\nwbegindocs{446}There are sepate commands for drawing the
curve and for drawing errors.  The symbols are applied to the latter.  First
of all, we may have to compute a baseline:
\nwenddocs{}\nwbegincode{447}\moddef{Analysis: procedures}\plusendmoddef
  function drawing_options_get_calc_command (dro) result (cmd)
    type(string_t) :: cmd
    type(drawing_options_t), intent(in) :: dro
    if (dro%with_base) then
       cmd = "calculate " // dro%dataset // ".base (" // dro%dataset // ") " &
            // "(x, #0);"
    else
       cmd = ""
    end if
  end function drawing_options_get_calc_command

\eatline
\nwendcode{}\nwbegindocs{448}Return the drawing command.  
\nwenddocs{}\nwbegincode{449}\moddef{Analysis: procedures}\plusendmoddef
  function drawing_options_get_draw_command (dro) result (cmd)
    type(string_t) :: cmd
    type(drawing_options_t), intent(in) :: dro
    if (dro%fill) then
       cmd = "fill"
    else if (dro%draw) then
       cmd = "draw"
    else
       cmd = ""
    end if
    if (dro%fill .or. dro%draw) then
       if (dro%piecewise)  cmd = cmd // " piecewise"
       if (dro%draw .and. dro%with_base)  cmd = cmd // " cyclic"
       cmd = cmd // " from (" // dro%dataset
       if (dro%with_base) then
          if (dro%piecewise) then
             cmd = cmd // ", " // dro%dataset // ".base/\\"  ! "
          else
             cmd = cmd // " ~ " // dro%dataset // ".base\\"  ! "
          end if
       end if
       cmd = cmd // ")"
       if (dro%fill) then
          cmd = cmd // " " // dro%fill_options
          if (dro%draw)  cmd = cmd // " outlined"
       end if
       if (dro%draw)  cmd = cmd // " " // dro%draw_options
       cmd = cmd // ";"
    end if
  end function drawing_options_get_draw_command

\eatline
\nwendcode{}\nwbegindocs{450}The error command draws error bars, if any.
\nwenddocs{}\nwbegincode{451}\moddef{Analysis: procedures}\plusendmoddef
  function drawing_options_get_err_command (dro) result (cmd)
    type(string_t) :: cmd
    type(drawing_options_t), intent(in) :: dro
    if (dro%err) then
       cmd = "draw piecewise " &
            // "from (" // dro%dataset // ".err)" &
            // " " // dro%err_options // ";"
    else
       cmd = ""
    end if
  end function drawing_options_get_err_command

\eatline
\nwendcode{}\nwbegindocs{452}The symbol command draws symbols, if any.
\nwenddocs{}\nwbegincode{453}\moddef{Analysis: procedures}\plusendmoddef
  function drawing_options_get_symb_command (dro) result (cmd)
    type(string_t) :: cmd
    type(drawing_options_t), intent(in) :: dro
    if (dro%symbols) then
       cmd = "phantom" &
            // " from (" // dro%dataset // ")" &
            // " withsymbol (" // dro%symbol // ");"
    else
       cmd = ""
    end if
  end function drawing_options_get_symb_command

\eatline
\nwendcode{}\nwbegindocs{454}Get extra GAMELAN code to be executed before and after the usual drawing
commands.
\nwenddocs{}\nwbegincode{455}\moddef{Analysis: procedures}\plusendmoddef
  function drawing_options_get_gml_bg_command (dro) result (cmd)
    type(string_t) :: cmd
    type(drawing_options_t), intent(in) :: dro
    cmd = dro%gmlcode_bg
  end function drawing_options_get_gml_bg_command

  function drawing_options_get_gml_fg_command (dro) result (cmd)
    type(string_t) :: cmd
    type(drawing_options_t), intent(in) :: dro
    cmd = dro%gmlcode_fg
  end function drawing_options_get_gml_fg_command

\eatline
\eatline
\nwendcode{}\nwbegindocs{456}\nwdocspar
\subsection{Observables}
The observable type holds the accumulated observable values and weight
sums which are necessary for proper averaging.
\nwenddocs{}\nwbegincode{457}\moddef{Analysis: types}\plusendmoddef
  type :: observable_t
     private
     real(default) :: sum_values = 0
     real(default) :: sum_squared_values = 0
     real(default) :: sum_weights = 0
     real(default) :: sum_squared_weights = 0
     integer :: count = 0
     type(string_t) :: obs_label
     type(string_t) :: obs_unit
     type(graph_options_t) :: graph_options
  end type observable_t

\eatline
\nwendcode{}\nwbegindocs{458}Initialize with defined properties
\nwenddocs{}\nwbegincode{459}\moddef{Analysis: procedures}\plusendmoddef
  subroutine observable_init (obs, obs_label, obs_unit, graph_options)
    type(observable_t), intent(out) :: obs
    type(string_t), intent(in), optional :: obs_label, obs_unit
    type(graph_options_t), intent(in), optional :: graph_options
    if (present (obs_label)) then
       obs%obs_label = obs_label
    else
       obs%obs_label = ""
    end if
    if (present (obs_unit)) then
       obs%obs_unit = obs_unit
    else
       obs%obs_unit = ""
    end if
    if (present (graph_options)) then
       obs%graph_options = graph_options
    else
       call graph_options_init (obs%graph_options)
    end if
  end subroutine observable_init

\eatline
\nwendcode{}\nwbegindocs{460}Reset all numeric entries.
\nwenddocs{}\nwbegincode{461}\moddef{Analysis: procedures}\plusendmoddef
  subroutine observable_clear (obs)
    type(observable_t), intent(inout) :: obs
    obs%sum_values = 0
    obs%sum_squared_values = 0
    obs%sum_weights = 0
    obs%sum_squared_weights = 0
    obs%count = 0
  end subroutine observable_clear

\eatline
\nwendcode{}\nwbegindocs{462}Record a a value.  Always successful for observables.
\nwenddocs{}\nwbegincode{463}\moddef{Analysis: interfaces}\endmoddef
  interface observable_record_value
     module procedure observable_record_value_unweighted
     module procedure observable_record_value_weighted
  end interface

\nwendcode{}\nwbegincode{464}\moddef{Analysis: procedures}\plusendmoddef
  subroutine observable_record_value_unweighted (obs, value, success)
    type(observable_t), intent(inout) :: obs
    real(default), intent(in) :: value
    logical, intent(out), optional :: success
    obs%sum_values = obs%sum_values + value
    obs%sum_squared_values = obs%sum_squared_values + value**2
    obs%sum_weights = obs%sum_weights + 1
    obs%sum_squared_weights = obs%sum_squared_weights + 1
    obs%count = obs%count + 1
    if (present (success))  success = .true.
  end subroutine observable_record_value_unweighted

  subroutine observable_record_value_weighted (obs, value, weight, success)
    type(observable_t), intent(inout) :: obs
    real(default), intent(in) :: value, weight
    logical, intent(out), optional :: success
    obs%sum_values = obs%sum_values + value * weight
    obs%sum_squared_values = obs%sum_squared_values + value**2 * weight
    obs%sum_weights = obs%sum_weights + abs (weight)
    obs%sum_squared_weights = obs%sum_squared_weights + weight**2
    obs%count = obs%count + 1
    if (present (success))  success = .true.
  end subroutine observable_record_value_weighted

\eatline
\nwendcode{}\nwbegindocs{465}Here are the statistics formulas:
\begin{enumerate}
\item Unweighted case:
  Given a sample of $n$ values $x_i$, the average is
  \begin{equation}
    \langle x \rangle = \frac{\sum x_i}{n}
  \end{equation}
  and the error estimate
  \begin{align}
    \Delta x &= \sqrt{\frac{1}{n-1}\langle{\sum(x_i - \langle x\rangle)^2}}
    \\
    &= \sqrt{\frac{1}{n-1}
             \left(\frac{\sum x_i^2}{n} - \frac{(\sum x_i)^2}{n^2}\right)}
  \end{align}
\item Weighted case:
  Instead of weight 1, each event comes with weight $w_i$.
  \begin{equation}
    \langle x \rangle = \frac{\sum x_i w_i}{\sum w_i}
  \end{equation}
  and
  \begin{equation}
    \Delta x 
    = \sqrt{\frac{1}{n-1}
            \left(\frac{\sum x_i^2 w_i}{\sum w_i} 
                  - \frac{(\sum x_i w_i)^2}{(\sum w_i)^2}\right)}
  \end{equation}
  For $w_i=1$, this specializes to the previous formula.
\end{enumerate}
\nwenddocs{}\nwbegincode{466}\moddef{Analysis: procedures}\plusendmoddef
  function observable_get_n_entries (obs) result (n)
    integer :: n
    type(observable_t), intent(in) :: obs
    n = obs%count
  end function observable_get_n_entries

  function observable_get_average (obs) result (avg)
    real(default) :: avg
    type(observable_t), intent(in) :: obs
    if (obs%sum_weights /= 0) then
       avg = obs%sum_values / obs%sum_weights
    else
       avg = 0
    end if
  end function observable_get_average

  function observable_get_error (obs) result (err)
    real(default) :: err
    type(observable_t), intent(in) :: obs
    real(default) :: var, n
    if (obs%sum_weights /= 0) then
       select case (obs%count)
       case (0:1)
          err = 0
       case default
          n = obs%count
          var = obs%sum_squared_values / obs%sum_weights &
                - (obs%sum_values / obs%sum_weights) ** 2
          err = sqrt (max (var, 0._default) / (n - 1))
       end select
    else
       err = 0
    end if
  end function observable_get_error

\eatline
\eatline
\eatline
\eatline
\nwendcode{}\nwbegindocs{467}Write label and/or physical unit to a string.
\nwenddocs{}\nwbegincode{468}\moddef{Analysis: procedures}\plusendmoddef
  function observable_get_label (obs, wl, wu) result (string)
    type(string_t) :: string
    type(observable_t), intent(in) :: obs 
    logical, intent(in) :: wl, wu
    type(string_t) :: obs_label, obs_unit
    if (wl) then
       if (obs%obs_label /= "") then
          obs_label = obs%obs_label
       else
          obs_label = "\\textrm\{Observable\}"
       end if
    else
       obs_label = ""
    end if
    if (wu) then
       if (obs%obs_unit /= "") then
          if (wl) then
             obs_unit = "\\;[" // obs%obs_unit // "]"
          else
             obs_unit = obs%obs_unit
          end if
       else
          obs_unit = ""
       end if
    else
       obs_unit = ""
    end if
    string = obs_label // obs_unit
  end function observable_get_label

\eatline
\nwendcode{}\nwbegindocs{469}\nwdocspar
\subsection{Output}
\nwenddocs{}\nwbegincode{470}\moddef{Analysis: procedures}\plusendmoddef
  subroutine observable_write (obs, unit)
    type(observable_t), intent(in) :: obs
    integer, intent(in), optional :: unit
    real(default) :: avg, err, relerr
    integer :: n
    integer :: u
    u = given_output_unit (unit);  if (u < 0)  return
    avg = observable_get_average (obs)
    err = observable_get_error (obs)
    if (avg /= 0) then
       relerr = err / avg
    else
       relerr = 0
    end if
    n = observable_get_n_entries (obs)
    if (obs%graph_options%title /= "") then
       write (u, "(A,1x,3A)") &
            "title       =", '"', char (obs%graph_options%title), '"'
    end if
    if (obs%graph_options%title /= "") then
       write (u, "(A,1x,3A)") &
            "description =", '"', char (obs%graph_options%description), '"'
    end if
    write (u, "(A,1x," // HISTOGRAM_DATA_FORMAT // ")", advance = "no") &
         "average     =", avg
    call write_unit ()
    write (u, "(A,1x," // HISTOGRAM_DATA_FORMAT // ")", advance = "no") &
         "error[abs]  =", err
    call write_unit ()
    write (u, "(A,1x," // HISTOGRAM_DATA_FORMAT // ")") &
         "error[rel]  =", relerr
    write (u, "(A,1x,I0)") &
         "n_entries   =", n
  contains
    subroutine write_unit ()
      if (obs%obs_unit /= "") then
         write (u, "(1x,A)")  char (obs%obs_unit)
      else
         write (u, *)
      end if
    end subroutine write_unit
  end subroutine observable_write

\eatline
\nwendcode{}\nwbegindocs{471}\LaTeX\ output.
\nwenddocs{}\nwbegincode{472}\moddef{Analysis: procedures}\plusendmoddef
  subroutine observable_write_driver (obs, unit, write_heading)
    type(observable_t), intent(in) :: obs
    integer, intent(in), optional :: unit
    logical, intent(in), optional :: write_heading
    real(default) :: avg, err
    integer :: n_digits
    logical :: heading
    integer :: u
    u = given_output_unit (unit);  if (u < 0)  return
    heading = .true.;  if (present (write_heading))  heading = write_heading
    avg = observable_get_average (obs)
    err = observable_get_error (obs)
    if (avg /= 0 .and. err /= 0) then
       n_digits = max (2, 2 - int (log10 (abs (err / real (avg, default)))))
    else if (avg /= 0) then
       n_digits = 100
    else
       n_digits = 1
    end if
    if (heading) then
       write (u, "(A)") 
       if (obs%graph_options%title /= "") then
          write (u, "(A)")  "\\section\{" // char (obs%graph_options%title) &
               // "\}"
       else
          write (u, "(A)")  "\\section\{Observable\}"
       end if
       if (obs%graph_options%description /= "") then
          write (u, "(A)")  char (obs%graph_options%description)
          write (u, *)
       end if
       write (u, "(A)")  "\\begin\{flushleft\}"
    end if
    write (u, "(A)", advance="no")  "  $\\langle\{" ! $ sign
    write (u, "(A)", advance="no")  char (observable_get_label (obs, wl=.true., wu=.false.))
    write (u, "(A)", advance="no")  "\}\\rangle = "
    write (u, "(A)", advance="no")  char (tex_format (avg, n_digits))
    write (u, "(A)", advance="no")  "\\pm"
    write (u, "(A)", advance="no")  char (tex_format (err, 2))
    write (u, "(A)", advance="no")  "\\;\{"
    write (u, "(A)", advance="no")  char (observable_get_label (obs, wl=.false., wu=.true.))
    write (u, "(A)")  "\}"
    write (u, "(A)", advance="no")  "     \\quad[n_\{\\text\{entries\}\} = "
    write (u, "(I0)",advance="no")  observable_get_n_entries (obs)
    write (u, "(A)")  "]$"          ! $ fool Emacs' noweb mode 
    if (heading) then
       write (u, "(A)") "\\end\{flushleft\}"
    end if
  end subroutine observable_write_driver

\eatline
\nwendcode{}\nwbegindocs{473}\nwdocspar
\subsection{Histograms}
\subsubsection{Bins}
\nwenddocs{}\nwbegincode{474}\moddef{Analysis: types}\plusendmoddef
  type :: bin_t
     private
     real(default) :: midpoint = 0
     real(default) :: width = 0
     real(default) :: sum_weights = 0
     real(default) :: sum_squared_weights = 0
     real(default) :: sum_excess_weights = 0
     integer :: count = 0
  end type bin_t

\eatline
\nwendcode{}\nwbegincode{475}\moddef{Analysis: procedures}\plusendmoddef
  subroutine bin_init (bin, midpoint, width)
    type(bin_t), intent(out) :: bin
    real(default), intent(in) :: midpoint, width
    bin%midpoint = midpoint
    bin%width = width
  end subroutine bin_init

\eatline
\nwendcode{}\nwbegincode{476}\moddef{Analysis: procedures}\plusendmoddef
  elemental subroutine bin_clear (bin)
    type(bin_t), intent(inout) :: bin
    bin%sum_weights = 0
    bin%sum_squared_weights = 0
    bin%sum_excess_weights = 0
    bin%count = 0
  end subroutine bin_clear

\eatline
\nwendcode{}\nwbegincode{477}\moddef{Analysis: procedures}\plusendmoddef
  subroutine bin_record_value (bin, normalize, weight, excess)
    type(bin_t), intent(inout) :: bin
    logical, intent(in) :: normalize
    real(default), intent(in) :: weight
    real(default), intent(in), optional :: excess
    real(default) :: w, e
    if (normalize) then
       if (bin%width /= 0) then
          w = weight / bin%width
          if (present (excess))  e = excess / bin%width
       else
          w = 0
          if (present (excess))  e = 0
       end if
    else
       w = weight
       if (present (excess))  e = excess
    end if
    bin%sum_weights = bin%sum_weights + abs (w)
    bin%sum_squared_weights = bin%sum_squared_weights + w ** 2
    if (present (excess)) &
         bin%sum_excess_weights = bin%sum_excess_weights + abs (e)
    bin%count = bin%count + 1
  end subroutine bin_record_value

\eatline
\nwendcode{}\nwbegincode{478}\moddef{Analysis: procedures}\plusendmoddef
  function bin_get_midpoint (bin) result (x)
    real(default) :: x
    type(bin_t), intent(in) :: bin
    x = bin%midpoint
  end function bin_get_midpoint

  function bin_get_width (bin) result (w)
    real(default) :: w
    type(bin_t), intent(in) :: bin
    w = bin%width
  end function bin_get_width

  function bin_get_n_entries (bin) result (n)
    integer :: n
    type(bin_t), intent(in) :: bin
    n = bin%count
  end function bin_get_n_entries
  
  function bin_get_sum (bin) result (s)
    real(default) :: s
    type(bin_t), intent(in) :: bin
    s = bin%sum_weights
  end function bin_get_sum

  function bin_get_error (bin) result (err)
    real(default) :: err
    type(bin_t), intent(in) :: bin
    err = sqrt (bin%sum_squared_weights)
  end function bin_get_error

  function bin_get_excess (bin) result (excess)
    real(default) :: excess
    type(bin_t), intent(in) :: bin
    excess = bin%sum_excess_weights
  end function bin_get_excess

\eatline
\eatline
\eatline
\eatline
\eatline
\eatline
\nwendcode{}\nwbegincode{479}\moddef{Analysis: procedures}\plusendmoddef
  subroutine bin_write_header (unit)
    integer, intent(in), optional :: unit
    character(120) :: buffer
    integer :: u
    u = given_output_unit (unit);  if (u < 0)  return
    write (buffer, "(A,4(1x," //HISTOGRAM_HEAD_FORMAT // "),2x,A)") &
         "#", "bin midpoint", "value    ", "error    ", &
         "excess     ", "n"
    write (u, "(A)")  trim (buffer)
  end subroutine bin_write_header

  subroutine bin_write (bin, unit)
    type(bin_t), intent(in) :: bin
    integer, intent(in), optional :: unit
    integer :: u
    u = given_output_unit (unit);  if (u < 0)  return
    write (u, "(1x,4(1x," // HISTOGRAM_DATA_FORMAT // "),2x,I0)") &
         bin_get_midpoint (bin), &
         bin_get_sum (bin), &
         bin_get_error (bin), &
         bin_get_excess (bin), &
         bin_get_n_entries (bin)
  end subroutine bin_write

\eatline
\eatline
\nwendcode{}\nwbegindocs{480}\nwdocspar
\subsubsection{Histograms}
\nwenddocs{}\nwbegincode{481}\moddef{Analysis: types}\plusendmoddef
  type :: histogram_t
     private
     real(default) :: lower_bound = 0
     real(default) :: upper_bound = 0
     real(default) :: width = 0
     integer :: n_bins = 0
     logical :: normalize_bins = .false.
     type(observable_t) :: obs
     type(observable_t) :: obs_within_bounds
     type(bin_t) :: underflow
     type(bin_t), dimension(:), allocatable :: bin
     type(bin_t) :: overflow
     type(graph_options_t) :: graph_options
     type(drawing_options_t) :: drawing_options
  end type histogram_t

\eatline
\nwendcode{}\nwbegindocs{482}\nwdocspar
\subsubsection{Initializer/finalizer}
Initialize a histogram.  We may provide either the bin width or the
number of bins.  A finalizer is not needed, since the histogram contains no
pointer (sub)components.
\nwenddocs{}\nwbegincode{483}\moddef{Analysis: interfaces}\plusendmoddef
  interface histogram_init
     module procedure histogram_init_n_bins
     module procedure histogram_init_bin_width
  end interface

\nwendcode{}\nwbegincode{484}\moddef{Analysis: procedures}\plusendmoddef
  subroutine histogram_init_n_bins (h, id, &
       lower_bound, upper_bound, n_bins, normalize_bins, &
       obs_label, obs_unit, graph_options, drawing_options)
    type(histogram_t), intent(out) :: h
    type(string_t), intent(in) :: id
    real(default), intent(in) :: lower_bound, upper_bound
    integer, intent(in) :: n_bins
    logical, intent(in) :: normalize_bins
    type(string_t), intent(in), optional :: obs_label, obs_unit
    type(graph_options_t), intent(in), optional :: graph_options
    type(drawing_options_t), intent(in), optional :: drawing_options
    real(default) :: bin_width
    integer :: i
    call observable_init (h%obs_within_bounds, obs_label, obs_unit) 
    call observable_init (h%obs, obs_label, obs_unit) 
    h%lower_bound = lower_bound
    h%upper_bound = upper_bound
    h%n_bins = max (n_bins, 1)
    h%width = h%upper_bound - h%lower_bound    
    h%normalize_bins = normalize_bins
    bin_width = h%width / h%n_bins
    allocate (h%bin (h%n_bins))
    call bin_init (h%underflow, h%lower_bound, 0._default)
    do i = 1, h%n_bins
       call bin_init (h%bin(i), &
            h%lower_bound - bin_width/2 + i * bin_width, bin_width)
    end do
    call bin_init (h%overflow, h%upper_bound, 0._default)
    if (present (graph_options)) then
       h%graph_options = graph_options
    else
       call graph_options_init (h%graph_options)
    end if
    call graph_options_set (h%graph_options, id = id)
    if (present (drawing_options)) then
       h%drawing_options = drawing_options
    else
       call drawing_options_init_histogram (h%drawing_options)
    end if
  end subroutine histogram_init_n_bins

  subroutine histogram_init_bin_width (h, id, &
       lower_bound, upper_bound, bin_width, normalize_bins, &
       obs_label, obs_unit, graph_options, drawing_options)
    type(histogram_t), intent(out) :: h
    type(string_t), intent(in) :: id
    real(default), intent(in) :: lower_bound, upper_bound, bin_width
    logical, intent(in) :: normalize_bins
    type(string_t), intent(in), optional :: obs_label, obs_unit
    type(graph_options_t), intent(in), optional :: graph_options
    type(drawing_options_t), intent(in), optional :: drawing_options
    integer :: n_bins
    if (bin_width /= 0) then
       n_bins = nint ((upper_bound - lower_bound) / bin_width)
    else
       n_bins = 1
    end if
    call histogram_init_n_bins (h, id, &
         lower_bound, upper_bound, n_bins, normalize_bins, &
         obs_label, obs_unit, graph_options, drawing_options)
  end subroutine histogram_init_bin_width

\eatline
\nwendcode{}\nwbegindocs{485}Initialize a histogram by copying another one.

Since {\tt{}h} has no pointer (sub)components, intrinsic assignment is
sufficient.  Optionally, we replace the drawing options.
\nwenddocs{}\nwbegincode{486}\moddef{Analysis: procedures}\plusendmoddef
  subroutine histogram_init_histogram (h, h_in, drawing_options)
    type(histogram_t), intent(out) :: h
    type(histogram_t), intent(in) :: h_in
    type(drawing_options_t), intent(in), optional :: drawing_options
    h = h_in
    if (present (drawing_options)) then
       h%drawing_options = drawing_options
    end if
  end subroutine histogram_init_histogram

\eatline
\nwendcode{}\nwbegindocs{487}\nwdocspar
\subsubsection{Fill histograms}
Clear the histogram contents, but do not modify the structure.
\nwenddocs{}\nwbegincode{488}\moddef{Analysis: procedures}\plusendmoddef
  subroutine histogram_clear (h)
    type(histogram_t), intent(inout) :: h
    call observable_clear (h%obs)
    call observable_clear (h%obs_within_bounds)
    call bin_clear (h%underflow)
    if (allocated (h%bin))  call bin_clear (h%bin)
    call bin_clear (h%overflow)
  end subroutine histogram_clear

\eatline
\nwendcode{}\nwbegindocs{489}Record a value.  Successful if the value is within bounds, otherwise
it is recorded as under-/overflow.  Optionally, we may provide an
excess weight that could be returned by the unweighting procedure.
\nwenddocs{}\nwbegincode{490}\moddef{Analysis: procedures}\plusendmoddef
  subroutine histogram_record_value_unweighted (h, value, excess, success)
    type(histogram_t), intent(inout) :: h
    real(default), intent(in) :: value
    real(default), intent(in), optional :: excess
    logical, intent(out), optional :: success
    integer :: i_bin
    call observable_record_value (h%obs, value)
    if (h%width /= 0) then
       i_bin = floor (((value - h%lower_bound) / h%width) * h%n_bins) + 1
    else
       i_bin = 0
    end if
    if (i_bin <= 0) then
       call bin_record_value (h%underflow, .false., 1._default, excess)
       if (present (success))  success = .false.
    else if (i_bin <= h%n_bins) then
       call observable_record_value (h%obs_within_bounds, value)
       call bin_record_value &
            (h%bin(i_bin), h%normalize_bins, 1._default, excess)
       if (present (success))  success = .true.
    else
       call bin_record_value (h%overflow, .false., 1._default, excess)
       if (present (success))  success = .false.
    end if
  end subroutine histogram_record_value_unweighted

\eatline
\nwendcode{}\nwbegindocs{491}Weighted events: analogous, but no excess weight.
\nwenddocs{}\nwbegincode{492}\moddef{Analysis: procedures}\plusendmoddef
  subroutine histogram_record_value_weighted (h, value, weight, success)
    type(histogram_t), intent(inout) :: h
    real(default), intent(in) :: value, weight
    logical, intent(out), optional :: success
    integer :: i_bin
    call observable_record_value (h%obs, value, weight)
    if (h%width /= 0) then
       i_bin = floor (((value - h%lower_bound) / h%width) * h%n_bins) + 1
    else
       i_bin = 0
    end if
    if (i_bin <= 0) then
       call bin_record_value (h%underflow, .false., weight)
       if (present (success))  success = .false.
    else if (i_bin <= h%n_bins) then
       call observable_record_value (h%obs_within_bounds, value, weight)
       call bin_record_value (h%bin(i_bin), h%normalize_bins, weight)
       if (present (success))  success = .true.
    else
       call bin_record_value (h%overflow, .false., weight)
       if (present (success))  success = .false.
    end if
  end subroutine histogram_record_value_weighted

\eatline
\nwendcode{}\nwbegindocs{493}\nwdocspar
\subsubsection{Access contents}
Inherited from the observable component (all-over average etc.)
\nwenddocs{}\nwbegincode{494}\moddef{Analysis: procedures}\plusendmoddef
  function histogram_get_n_entries (h) result (n)
    integer :: n
    type(histogram_t), intent(in) :: h
    n = observable_get_n_entries (h%obs)
  end function histogram_get_n_entries

  function histogram_get_average (h) result (avg)
    real(default) :: avg
    type(histogram_t), intent(in) :: h
    avg = observable_get_average (h%obs)
  end function histogram_get_average

  function histogram_get_error (h) result (err)
    real(default) :: err
    type(histogram_t), intent(in) :: h
    err = observable_get_error (h%obs)
  end function histogram_get_error

\eatline
\eatline
\eatline
\nwendcode{}\nwbegindocs{495}Analogous, but applied only to events within bounds.
\nwenddocs{}\nwbegincode{496}\moddef{Analysis: procedures}\plusendmoddef
  function histogram_get_n_entries_within_bounds (h) result (n)
    integer :: n
    type(histogram_t), intent(in) :: h
    n = observable_get_n_entries (h%obs_within_bounds)
  end function histogram_get_n_entries_within_bounds

  function histogram_get_average_within_bounds (h) result (avg)
    real(default) :: avg
    type(histogram_t), intent(in) :: h
    avg = observable_get_average (h%obs_within_bounds)
  end function histogram_get_average_within_bounds

  function histogram_get_error_within_bounds (h) result (err)
    real(default) :: err
    type(histogram_t), intent(in) :: h
    err = observable_get_error (h%obs_within_bounds)
  end function histogram_get_error_within_bounds

\eatline
\eatline
\eatline
\nwendcode{}\nwbegindocs{497}Get the number of bins
\nwenddocs{}\nwbegincode{498}\moddef{Analysis: procedures}\plusendmoddef
  function histogram_get_n_bins (h) result (n)
    type(histogram_t), intent(in) :: h
    integer :: n
    n = h%n_bins
  end function histogram_get_n_bins

\eatline
\nwendcode{}\nwbegindocs{499}Check bins.  If the index is zero or above the limit, return the
results for underflow or overflow, respectively.
\nwenddocs{}\nwbegincode{500}\moddef{Analysis: procedures}\plusendmoddef
  function histogram_get_n_entries_for_bin (h, i) result (n)
    integer :: n
    type(histogram_t), intent(in) :: h
    integer, intent(in) :: i
    if (i <= 0) then
       n = bin_get_n_entries (h%underflow)
    else if (i <= h%n_bins) then
       n = bin_get_n_entries (h%bin(i))
    else
       n = bin_get_n_entries (h%overflow)
    end if
  end function histogram_get_n_entries_for_bin

  function histogram_get_sum_for_bin (h, i) result (avg)
    real(default) :: avg
    type(histogram_t), intent(in) :: h
    integer, intent(in) :: i
    if (i <= 0) then
       avg = bin_get_sum (h%underflow)
    else if (i <= h%n_bins) then
       avg = bin_get_sum (h%bin(i))
    else
       avg = bin_get_sum (h%overflow)
    end if
  end function histogram_get_sum_for_bin

  function histogram_get_error_for_bin (h, i) result (err)
    real(default) :: err
    type(histogram_t), intent(in) :: h
    integer, intent(in) :: i
    if (i <= 0) then
       err = bin_get_error (h%underflow)
    else if (i <= h%n_bins) then
       err = bin_get_error (h%bin(i))
    else
       err = bin_get_error (h%overflow)
    end if
  end function histogram_get_error_for_bin

  function histogram_get_excess_for_bin (h, i) result (err)
    real(default) :: err
    type(histogram_t), intent(in) :: h
    integer, intent(in) :: i
    if (i <= 0) then
       err = bin_get_excess (h%underflow)
    else if (i <= h%n_bins) then
       err = bin_get_excess (h%bin(i))
    else
       err = bin_get_excess (h%overflow)
    end if
  end function histogram_get_excess_for_bin

\eatline
\eatline
\eatline
\eatline
\nwendcode{}\nwbegindocs{501}Return a pointer to the graph options.
\nwenddocs{}\nwbegincode{502}\moddef{Analysis: procedures}\plusendmoddef
  function histogram_get_graph_options_ptr (h) result (ptr)
    type(graph_options_t), pointer :: ptr
    type(histogram_t), intent(in), target :: h
    ptr => h%graph_options
  end function histogram_get_graph_options_ptr

\eatline
\nwendcode{}\nwbegindocs{503}Return a pointer to the drawing options.
\nwenddocs{}\nwbegincode{504}\moddef{Analysis: procedures}\plusendmoddef
  function histogram_get_drawing_options_ptr (h) result (ptr)
    type(drawing_options_t), pointer :: ptr
    type(histogram_t), intent(in), target :: h
    ptr => h%drawing_options
  end function histogram_get_drawing_options_ptr

\eatline
\nwendcode{}\nwbegindocs{505}\nwdocspar
\subsubsection{Output}
\nwenddocs{}\nwbegincode{506}\moddef{Analysis: procedures}\plusendmoddef
  subroutine histogram_write (h, unit)
    type(histogram_t), intent(in) :: h
    integer, intent(in), optional :: unit
    integer :: u, i
    u = given_output_unit (unit);  if (u < 0)  return
    call bin_write_header (u)
    if (allocated (h%bin)) then
       do i = 1, h%n_bins
          call bin_write (h%bin(i), u)
       end do
    end if
    write (u, *)
    write (u, "(A,1x,A)")  "#", "Underflow:"
    call bin_write (h%underflow, u)
    write (u, *)
    write (u, "(A,1x,A)")  "#", "Overflow:"
    call bin_write (h%overflow, u)
    write (u, *)
    write (u, "(A,1x,A)")  "#", "Summary: data within bounds"
    call observable_write (h%obs_within_bounds, u)
    write (u, *)
    write (u, "(A,1x,A)")  "#", "Summary: all data"
    call observable_write (h%obs, u)
    write (u, *)
  end subroutine histogram_write

\eatline
\nwendcode{}\nwbegindocs{507}Write the GAMELAN reader for histogram contents.
\nwenddocs{}\nwbegincode{508}\moddef{Analysis: procedures}\plusendmoddef
  subroutine histogram_write_gml_reader (h, filename, unit)
    type(histogram_t), intent(in) :: h
    type(string_t), intent(in) :: filename
    integer, intent(in), optional :: unit
    character(*), parameter :: fmt = "(ES15.8)"    
    integer :: u
    u = given_output_unit (unit);  if (u < 0)  return
    write (u, "(2x,A)")  'fromfile "' // char (filename) // '":'
    write (u, "(4x,A)")  'key "# Histogram:";'
    write (u, "(4x,A)")  'dx := #' &
         // real2char (h%width / h%n_bins / 2, fmt) // ';'
    write (u, "(4x,A)")  'for i withinblock:' 
    write (u, "(6x,A)")  'get x, y, y.d, y.n, y.e;'
    if (h%drawing_options%with_hbars) then
       write (u, "(6x,A)")  'plot (' // char (h%drawing_options%dataset) &
            // ') (x,y) hbar dx;'
    else
       write (u, "(6x,A)")  'plot (' // char (h%drawing_options%dataset) &
            // ') (x,y);'
    end if
    if (h%drawing_options%err) then
       write (u, "(6x,A)")  'plot (' // char (h%drawing_options%dataset) &
         // '.err) ' &
            // '(x,y) vbar y.d;'
    end if
    !!! Future excess options for plots
    ! write (u, "(6x,A)")  'if show_excess: ' // &
    !            & 'plot(dat.e)(x, y plus y.e) hbar dx; fi'
    write (u, "(4x,A)")  'endfor'
    write (u, "(2x,A)")  'endfrom'
  end subroutine histogram_write_gml_reader

\eatline
\nwendcode{}\nwbegindocs{509}\LaTeX\ and GAMELAN output.
\nwenddocs{}\nwbegincode{510}\moddef{Analysis: procedures}\plusendmoddef
  subroutine histogram_write_gml_driver (h, filename, unit)
    type(histogram_t), intent(in) :: h
    type(string_t), intent(in) :: filename
    integer, intent(in), optional :: unit
    type(string_t) :: calc_cmd, bg_cmd, draw_cmd, err_cmd, symb_cmd, fg_cmd
    integer :: u
    u = given_output_unit (unit);  if (u < 0)  return
    call graph_options_write_tex_header (h%graph_options, unit)
    write (u, "(2x,A)")  char (graph_options_get_gml_setup (h%graph_options))
    write (u, "(2x,A)")  char (graph_options_get_gml_graphrange &
         (h%graph_options, x_min=h%lower_bound, x_max=h%upper_bound))
    call histogram_write_gml_reader (h, filename, unit)
    calc_cmd = drawing_options_get_calc_command (h%drawing_options)
    if (calc_cmd /= "")  write (u, "(2x,A)")  char (calc_cmd)
    bg_cmd = drawing_options_get_gml_bg_command (h%drawing_options)
    if (bg_cmd /= "")  write (u, "(2x,A)")  char (bg_cmd)
    draw_cmd = drawing_options_get_draw_command (h%drawing_options)
    if (draw_cmd /= "")  write (u, "(2x,A)")  char (draw_cmd)
    err_cmd = drawing_options_get_err_command (h%drawing_options)
    if (err_cmd /= "")  write (u, "(2x,A)")  char (err_cmd)
    symb_cmd = drawing_options_get_symb_command (h%drawing_options)
    if (symb_cmd /= "")  write (u, "(2x,A)")  char (symb_cmd)
    fg_cmd = drawing_options_get_gml_fg_command (h%drawing_options)
    if (fg_cmd /= "")  write (u, "(2x,A)")  char (fg_cmd)
    write (u, "(2x,A)")  char (graph_options_get_gml_x_label (h%graph_options))
    write (u, "(2x,A)")  char (graph_options_get_gml_y_label (h%graph_options))
    call graph_options_write_tex_footer (h%graph_options, unit)
    write (u, "(A)") "\\vspace*\{2\\baselineskip\}"
    write (u, "(A)") "\\begin\{flushleft\}"
    write (u, "(A)") "\\textbf\{Data within bounds:\} \\\\"
    call observable_write_driver (h%obs_within_bounds, unit, &
                                  write_heading=.false.)
    write (u, "(A)") "\\\\[0.5\\baselineskip]"
    write (u, "(A)") "\\textbf\{All data:\} \\\\"
    call observable_write_driver (h%obs, unit, write_heading=.false.)
    write (u, "(A)") "\\end\{flushleft\}"
  end subroutine histogram_write_gml_driver

\eatline
\nwendcode{}\nwbegindocs{511}Return the header for generic data output as an ifile.
\nwenddocs{}\nwbegincode{512}\moddef{Analysis: procedures}\plusendmoddef
  subroutine histogram_get_header (h, header, comment)
    type(histogram_t), intent(in) :: h
    type(ifile_t), intent(inout) :: header
    type(string_t), intent(in), optional :: comment
    type(string_t) :: c
    if (present (comment)) then
       c = comment
    else
       c = ""
    end if
    call ifile_append (header, c // "WHIZARD histogram data")
    call graph_options_get_header (h%graph_options, header, comment)
    call ifile_append (header, &
         c // "range: " // real2string (h%lower_bound) &
         // " - " // real2string (h%upper_bound))
    call ifile_append (header, &
         c // "counts total: " &
         // int2char (histogram_get_n_entries_within_bounds (h)))
    call ifile_append (header, &
         c // "total average: " &
         // real2string (histogram_get_average_within_bounds (h)) // " +- " &
         // real2string (histogram_get_error_within_bounds (h)))
  end subroutine histogram_get_header

\eatline
\nwendcode{}\nwbegindocs{513}\nwdocspar
\subsection{Plots}
\subsubsection{Points}
\nwenddocs{}\nwbegincode{514}\moddef{Analysis: types}\plusendmoddef
  type :: point_t
     private
     real(default) :: x = 0
     real(default) :: y = 0
     real(default) :: yerr = 0
     real(default) :: xerr = 0
     type(point_t), pointer :: next => null ()
  end type point_t

\eatline
\nwendcode{}\nwbegincode{515}\moddef{Analysis: interfaces}\plusendmoddef
  interface point_init
    module procedure point_init_contents
    module procedure point_init_point
  end interface
\nwendcode{}\nwbegincode{516}\moddef{Analysis: procedures}\plusendmoddef
  subroutine point_init_contents (point, x, y, yerr, xerr)
    type(point_t), intent(out) :: point
    real(default), intent(in) :: x, y
    real(default), intent(in), optional :: yerr, xerr
    point%x = x
    point%y = y
    if (present (yerr))  point%yerr = yerr
    if (present (xerr))  point%xerr = xerr
  end subroutine point_init_contents

  subroutine point_init_point (point, point_in)
    type(point_t), intent(out) :: point
    type(point_t), intent(in) :: point_in
    point%x = point_in%x
    point%y = point_in%y
    point%yerr = point_in%yerr
    point%xerr = point_in%xerr
  end subroutine point_init_point

\eatline
\nwendcode{}\nwbegincode{517}\moddef{Analysis: procedures}\plusendmoddef
  function point_get_x (point) result (x)
    real(default) :: x
    type(point_t), intent(in) :: point
    x = point%x
  end function point_get_x

  function point_get_y (point) result (y)
    real(default) :: y
    type(point_t), intent(in) :: point
    y = point%y
  end function point_get_y

  function point_get_xerr (point) result (xerr)
    real(default) :: xerr
    type(point_t), intent(in) :: point
    xerr = point%xerr
  end function point_get_xerr

  function point_get_yerr (point) result (yerr)
    real(default) :: yerr
    type(point_t), intent(in) :: point
    yerr = point%yerr
  end function point_get_yerr

\eatline
\eatline
\eatline
\eatline
\nwendcode{}\nwbegincode{518}\moddef{Analysis: procedures}\plusendmoddef
  subroutine point_write_header (unit)
    integer, intent(in) :: unit
    character(120) :: buffer
    integer :: u
    u = given_output_unit (unit);  if (u < 0)  return
    write (buffer, "(A,4(1x," // HISTOGRAM_HEAD_FORMAT // "))") &
         "#", "x       ", "y       ", "yerr     ", "xerr     "
    write (u, "(A)")  trim (buffer)
  end subroutine point_write_header

  subroutine point_write (point, unit)
    type(point_t), intent(in) :: point
    integer, intent(in), optional :: unit
    integer :: u
    u = given_output_unit (unit);  if (u < 0)  return
    write (u, "(1x,4(1x," // HISTOGRAM_DATA_FORMAT // "))") &
         point_get_x (point), &
         point_get_y (point), &
         point_get_yerr (point), &
         point_get_xerr (point)
  end subroutine point_write

\eatline
\nwendcode{}\nwbegindocs{519}\nwdocspar
\subsubsection{Plots}
\nwenddocs{}\nwbegincode{520}\moddef{Analysis: types}\plusendmoddef
  type :: plot_t
     private
     type(point_t), pointer :: first => null ()
     type(point_t), pointer :: last => null ()
     integer :: count = 0
     type(graph_options_t) :: graph_options
     type(drawing_options_t) :: drawing_options
  end type plot_t

\eatline
\nwendcode{}\nwbegindocs{521}\nwdocspar
\subsubsection{Initializer/finalizer}
Initialize a plot.  We provide the lower and upper bound in the $x$
direction.
\nwenddocs{}\nwbegincode{522}\moddef{Analysis: interfaces}\plusendmoddef
  interface plot_init
     module procedure plot_init_empty
     module procedure plot_init_plot
  end interface
\nwendcode{}\nwbegincode{523}\moddef{Analysis: procedures}\plusendmoddef
  subroutine plot_init_empty (p, id, graph_options, drawing_options)
    type(plot_t), intent(out) :: p
    type(string_t), intent(in) :: id
    type(graph_options_t), intent(in), optional :: graph_options
    type(drawing_options_t), intent(in), optional :: drawing_options
    if (present (graph_options)) then
       p%graph_options = graph_options
    else
       call graph_options_init (p%graph_options)
    end if
    call graph_options_set (p%graph_options, id = id)
    if (present (drawing_options)) then
       p%drawing_options = drawing_options
    else
       call drawing_options_init_plot (p%drawing_options)
    end if
  end subroutine plot_init_empty

\eatline
\nwendcode{}\nwbegindocs{524}Initialize a plot by copying another one, optionally merging in a new
set of drawing options.

Since {\tt{}p} has pointer (sub)components, we have to explicitly deep-copy the
original.
\nwenddocs{}\nwbegincode{525}\moddef{Analysis: procedures}\plusendmoddef
  subroutine plot_init_plot (p, p_in, drawing_options)
    type(plot_t), intent(out) :: p
    type(plot_t), intent(in) :: p_in
    type(drawing_options_t), intent(in), optional :: drawing_options
    type(point_t), pointer :: current, new
    current => p_in%first
    do while (associated (current))
       allocate (new)
       call point_init (new, current)
       if (associated (p%last)) then
          p%last%next => new
       else
          p%first => new
       end if
       p%last => new
       current => current%next
    end do
    p%count = p_in%count
    p%graph_options = p_in%graph_options
    if (present (drawing_options)) then
       p%drawing_options = drawing_options
    else
       p%drawing_options = p_in%drawing_options
    end if
  end subroutine plot_init_plot

\eatline
\nwendcode{}\nwbegindocs{526}Finalize the plot by deallocating the list of points.
\nwenddocs{}\nwbegincode{527}\moddef{Analysis: procedures}\plusendmoddef
  subroutine plot_final (plot)
    type(plot_t), intent(inout) :: plot
    type(point_t), pointer :: current
    do while (associated (plot%first))
       current => plot%first
       plot%first => current%next
       deallocate (current)
    end do
    plot%last => null ()
  end subroutine plot_final

\eatline
\nwendcode{}\nwbegindocs{528}\nwdocspar
\subsubsection{Fill plots}
Clear the plot contents, but do not modify the structure.
\nwenddocs{}\nwbegincode{529}\moddef{Analysis: procedures}\plusendmoddef
  subroutine plot_clear (plot)
    type(plot_t), intent(inout) :: plot
    plot%count = 0
    call plot_final (plot)
  end subroutine plot_clear

\eatline
\nwendcode{}\nwbegindocs{530}Record a value.  Successful if the value is within bounds, otherwise
it is recorded as under-/overflow.
\nwenddocs{}\nwbegincode{531}\moddef{Analysis: procedures}\plusendmoddef
  subroutine plot_record_value (plot, x, y, yerr, xerr, success)
    type(plot_t), intent(inout) :: plot
    real(default), intent(in) :: x, y
    real(default), intent(in), optional :: yerr, xerr
    logical, intent(out), optional :: success
    type(point_t), pointer :: point
    plot%count = plot%count + 1
    allocate (point)
    call point_init (point, x, y, yerr, xerr)
    if (associated (plot%first)) then
       plot%last%next => point
    else
       plot%first => point
    end if
    plot%last => point
    if (present (success))  success = .true.
  end subroutine plot_record_value

\eatline
\nwendcode{}\nwbegindocs{532}\nwdocspar
\subsubsection{Access contents}
The number of points.
\nwenddocs{}\nwbegincode{533}\moddef{Analysis: procedures}\plusendmoddef
  function plot_get_n_entries (plot) result (n)
    integer :: n
    type(plot_t), intent(in) :: plot
    n = plot%count
  end function plot_get_n_entries

\eatline
\nwendcode{}\nwbegindocs{534}Return a pointer to the graph options.
\nwenddocs{}\nwbegincode{535}\moddef{Analysis: procedures}\plusendmoddef
  function plot_get_graph_options_ptr (p) result (ptr)
    type(graph_options_t), pointer :: ptr
    type(plot_t), intent(in), target :: p
    ptr => p%graph_options
  end function plot_get_graph_options_ptr

\eatline
\nwendcode{}\nwbegindocs{536}Return a pointer to the drawing options.
\nwenddocs{}\nwbegincode{537}\moddef{Analysis: procedures}\plusendmoddef
  function plot_get_drawing_options_ptr (p) result (ptr)
    type(drawing_options_t), pointer :: ptr
    type(plot_t), intent(in), target :: p
    ptr => p%drawing_options
  end function plot_get_drawing_options_ptr

\eatline
\nwendcode{}\nwbegindocs{538}\nwdocspar
\subsubsection{Output}
This output format is used by the GAMELAN driver below.
\nwenddocs{}\nwbegincode{539}\moddef{Analysis: procedures}\plusendmoddef
  subroutine plot_write (plot, unit)
    type(plot_t), intent(in) :: plot
    integer, intent(in), optional :: unit
    type(point_t), pointer :: point
    integer :: u
    u = given_output_unit (unit);  if (u < 0)  return
    call point_write_header (u)
    point => plot%first
    do while (associated (point))
       call point_write (point, unit)
       point => point%next
    end do
    write (u, *)
    write (u, "(A,1x,A)")  "#", "Summary:"
    write (u, "(A,1x,I0)") &
         "n_entries =", plot_get_n_entries (plot)
    write (u, *)
  end subroutine plot_write

\eatline
\nwendcode{}\nwbegindocs{540}Write the GAMELAN reader for plot contents.
\nwenddocs{}\nwbegincode{541}\moddef{Analysis: procedures}\plusendmoddef
  subroutine plot_write_gml_reader (p, filename, unit)
    type(plot_t), intent(in) :: p
    type(string_t), intent(in) :: filename
    integer, intent(in), optional :: unit
    integer :: u
    u = given_output_unit (unit);  if (u < 0)  return
    write (u, "(2x,A)")  'fromfile "' // char (filename) // '":'
    write (u, "(4x,A)")  'key "# Plot:";'
    write (u, "(4x,A)")  'for i withinblock:' 
    write (u, "(6x,A)")  'get x, y, y.err, x.err;'
    write (u, "(6x,A)")  'plot (' // char (p%drawing_options%dataset) &
         // ') (x,y);'
    if (p%drawing_options%err) then
       write (u, "(6x,A)")  'plot (' // char (p%drawing_options%dataset) &
            // '.err) (x,y) vbar y.err hbar x.err;'
    end if
    write (u, "(4x,A)")  'endfor'
    write (u, "(2x,A)")  'endfrom'
  end subroutine plot_write_gml_reader

\eatline
\nwendcode{}\nwbegindocs{542}\LaTeX\ and GAMELAN output.  Analogous to histogram output.
\nwenddocs{}\nwbegincode{543}\moddef{Analysis: procedures}\plusendmoddef
  subroutine plot_write_gml_driver (p, filename, unit)
    type(plot_t), intent(in) :: p
    type(string_t), intent(in) :: filename
    integer, intent(in), optional :: unit
    type(string_t) :: calc_cmd, bg_cmd, draw_cmd, err_cmd, symb_cmd, fg_cmd
    integer :: u
    u = given_output_unit (unit);  if (u < 0)  return
    call graph_options_write_tex_header (p%graph_options, unit)
    write (u, "(2x,A)") &
         char (graph_options_get_gml_setup (p%graph_options))
    write (u, "(2x,A)") &
         char (graph_options_get_gml_graphrange (p%graph_options))
    call plot_write_gml_reader (p, filename, unit)
    calc_cmd = drawing_options_get_calc_command (p%drawing_options)
    if (calc_cmd /= "")  write (u, "(2x,A)")  char (calc_cmd)
    bg_cmd = drawing_options_get_gml_bg_command (p%drawing_options)
    if (bg_cmd /= "")  write (u, "(2x,A)")  char (bg_cmd)
    draw_cmd = drawing_options_get_draw_command (p%drawing_options)
    if (draw_cmd /= "")  write (u, "(2x,A)")  char (draw_cmd)
    err_cmd = drawing_options_get_err_command (p%drawing_options)
    if (err_cmd /= "")  write (u, "(2x,A)")  char (err_cmd)
    symb_cmd = drawing_options_get_symb_command (p%drawing_options)
    if (symb_cmd /= "")  write (u, "(2x,A)")  char (symb_cmd)
    fg_cmd = drawing_options_get_gml_fg_command (p%drawing_options)
    if (fg_cmd /= "")  write (u, "(2x,A)")  char (fg_cmd)
    write (u, "(2x,A)")  char (graph_options_get_gml_x_label (p%graph_options))
    write (u, "(2x,A)")  char (graph_options_get_gml_y_label (p%graph_options))
    call graph_options_write_tex_footer (p%graph_options, unit)
  end subroutine plot_write_gml_driver

\eatline
\nwendcode{}\nwbegindocs{544}Append header for generic data output in ifile format.
\nwenddocs{}\nwbegincode{545}\moddef{Analysis: procedures}\plusendmoddef
  subroutine plot_get_header (plot, header, comment)
    type(plot_t), intent(in) :: plot
    type(ifile_t), intent(inout) :: header
    type(string_t), intent(in), optional :: comment
    type(string_t) :: c
    if (present (comment)) then
       c = comment
    else
       c = ""
    end if
    call ifile_append (header, c // "WHIZARD plot data")
    call graph_options_get_header (plot%graph_options, header, comment)
    call ifile_append (header, &
         c // "number of points: " &
         // int2char (plot_get_n_entries (plot)))
  end subroutine plot_get_header

\eatline
\nwendcode{}\nwbegindocs{546}\nwdocspar
\subsection{Graphs}
A graph is a container for several graph elements.  Each graph element is
either a plot or a histogram.  There is an appropriate base type below
(the {\tt{}analysis{\char95}object{\char95}t}), but to avoid recursion, we define a separate base
type here.  Note that there is no actual recursion: a graph is an analysis
object, but a graph cannot contain graphs.

(If we could use type extension, the implementation would be much more
transparent.)

\subsubsection{Graph elements}
Graph elements cannot be filled by the {\tt{}record} command directly.  The
contents are always copied from elementary histograms or plots.
\nwenddocs{}\nwbegincode{547}\moddef{Analysis: types}\plusendmoddef
  type :: graph_element_t
     private
     integer :: type = AN_UNDEFINED
     type(histogram_t), pointer :: h => null ()
     type(plot_t), pointer :: p => null ()
  end type graph_element_t

\eatline
\nwendcode{}\nwbegincode{548}\moddef{Analysis: procedures}\plusendmoddef
  subroutine graph_element_final (el)
    type(graph_element_t), intent(inout) :: el
    select case (el%type)
    case (AN_HISTOGRAM)
       deallocate (el%h)
    case (AN_PLOT)
       call plot_final (el%p)
       deallocate (el%p)
    end select
    el%type = AN_UNDEFINED
  end subroutine graph_element_final

\eatline
\nwendcode{}\nwbegindocs{549}Return the number of entries in the graph element:
\nwenddocs{}\nwbegincode{550}\moddef{Analysis: procedures}\plusendmoddef
  function graph_element_get_n_entries (el) result (n)
    integer :: n
    type(graph_element_t), intent(in) :: el
    select case (el%type)
    case (AN_HISTOGRAM);  n = histogram_get_n_entries (el%h)
    case (AN_PLOT);       n = plot_get_n_entries (el%p)
    case default;         n = 0
    end select
  end function graph_element_get_n_entries

\eatline
\nwendcode{}\nwbegindocs{551}Return a pointer to the graph / drawing options.
\nwenddocs{}\nwbegincode{552}\moddef{Analysis: procedures}\plusendmoddef
  function graph_element_get_graph_options_ptr (el) result (ptr)
    type(graph_options_t), pointer :: ptr
    type(graph_element_t), intent(in) :: el
    select case (el%type)
    case (AN_HISTOGRAM);  ptr => histogram_get_graph_options_ptr (el%h)
    case (AN_PLOT);       ptr => plot_get_graph_options_ptr (el%p)
    case default;         ptr => null ()
    end select
  end function graph_element_get_graph_options_ptr

  function graph_element_get_drawing_options_ptr (el) result (ptr)
    type(drawing_options_t), pointer :: ptr
    type(graph_element_t), intent(in) :: el
    select case (el%type)
    case (AN_HISTOGRAM);  ptr => histogram_get_drawing_options_ptr (el%h)
    case (AN_PLOT);       ptr => plot_get_drawing_options_ptr (el%p)
    case default;         ptr => null ()
    end select
  end function graph_element_get_drawing_options_ptr

\eatline
\eatline
\nwendcode{}\nwbegindocs{553}Output, simple wrapper for the plot/histogram writer.
\nwenddocs{}\nwbegincode{554}\moddef{Analysis: procedures}\plusendmoddef
  subroutine graph_element_write (el, unit)
    type(graph_element_t), intent(in) :: el
    integer, intent(in), optional :: unit
    type(graph_options_t), pointer :: gro
    type(string_t) :: id
    integer :: u
    u = given_output_unit (unit);  if (u < 0)  return
    gro => graph_element_get_graph_options_ptr (el)
    id = graph_options_get_id (gro)
    write (u, "(A,A)")  '#', repeat ("-", 78)
    select case (el%type)
    case (AN_HISTOGRAM)
       write (u, "(A)", advance="no")  "# Histogram: "
       write (u, "(1x,A)")  char (id)
       call histogram_write (el%h, unit)
    case (AN_PLOT)
       write (u, "(A)", advance="no")  "# Plot: "
       write (u, "(1x,A)")  char (id)
       call plot_write (el%p, unit)
    end select
  end subroutine graph_element_write

\eatline
\nwendcode{}\nwbegincode{555}\moddef{Analysis: procedures}\plusendmoddef
  subroutine graph_element_write_gml_reader (el, filename, unit)
    type(graph_element_t), intent(in) :: el
    type(string_t), intent(in) :: filename
    integer, intent(in), optional :: unit
    select case (el%type)
    case (AN_HISTOGRAM);  call histogram_write_gml_reader (el%h, filename, unit)
    case (AN_PLOT);       call plot_write_gml_reader (el%p, filename, unit)
    end select
  end subroutine graph_element_write_gml_reader

\eatline
\nwendcode{}\nwbegindocs{556}\nwdocspar
\subsubsection{The graph type}
The actual graph type contains its own {\tt{}graph{\char95}options}, which override the
individual settings.  The {\tt{}drawing{\char95}options} are set in the graph elements.
This distinction motivates the separation of the two types.
\nwenddocs{}\nwbegincode{557}\moddef{Analysis: types}\plusendmoddef
  type :: graph_t
     private
     type(graph_element_t), dimension(:), allocatable :: el
     type(graph_options_t) :: graph_options
  end type graph_t

\eatline
\nwendcode{}\nwbegindocs{558}\nwdocspar
\subsubsection{Initializer/finalizer}
The graph is created with a definite number of elements.  The elements are
filled one by one, optionally with modified drawing options.
\nwenddocs{}\nwbegincode{559}\moddef{Analysis: procedures}\plusendmoddef
  subroutine graph_init (g, id, n_elements, graph_options)
    type(graph_t), intent(out) :: g
    type(string_t), intent(in) :: id
    integer, intent(in) :: n_elements
    type(graph_options_t), intent(in), optional :: graph_options
    allocate (g%el (n_elements))
    if (present (graph_options)) then
       g%graph_options = graph_options
    else
       call graph_options_init (g%graph_options)
    end if
    call graph_options_set (g%graph_options, id = id)
  end subroutine graph_init

\eatline
\nwendcode{}\nwbegincode{560}\moddef{Analysis: procedures}\plusendmoddef
  subroutine graph_insert_histogram (g, i, h, drawing_options)
    type(graph_t), intent(inout), target :: g
    integer, intent(in) :: i
    type(histogram_t), intent(in) :: h
    type(drawing_options_t), intent(in), optional :: drawing_options
    type(graph_options_t), pointer :: gro
    type(drawing_options_t), pointer :: dro
    type(string_t) :: id
    g%el(i)%type = AN_HISTOGRAM
    allocate (g%el(i)%h)
    call histogram_init_histogram (g%el(i)%h, h, drawing_options)
    gro => histogram_get_graph_options_ptr (g%el(i)%h)
    dro => histogram_get_drawing_options_ptr (g%el(i)%h)
    id = graph_options_get_id (gro)
    call drawing_options_set (dro, dataset = "dat." // id)
  end subroutine graph_insert_histogram

\eatline
\nwendcode{}\nwbegincode{561}\moddef{Analysis: procedures}\plusendmoddef
  subroutine graph_insert_plot (g, i, p, drawing_options)
    type(graph_t), intent(inout) :: g
    integer, intent(in) :: i
    type(plot_t), intent(in) :: p
    type(drawing_options_t), intent(in), optional :: drawing_options
    type(graph_options_t), pointer :: gro
    type(drawing_options_t), pointer :: dro
    type(string_t) :: id
    g%el(i)%type = AN_PLOT
    allocate (g%el(i)%p)
    call plot_init_plot (g%el(i)%p, p, drawing_options)
    gro => plot_get_graph_options_ptr (g%el(i)%p)
    dro => plot_get_drawing_options_ptr (g%el(i)%p)
    id = graph_options_get_id (gro)
    call drawing_options_set (dro, dataset = "dat." // id)
  end subroutine graph_insert_plot

\eatline
\nwendcode{}\nwbegindocs{562}Finalizer.
\nwenddocs{}\nwbegincode{563}\moddef{Analysis: procedures}\plusendmoddef
  subroutine graph_final (g)
    type(graph_t), intent(inout) :: g
    integer :: i
    do i = 1, size (g%el)
       call graph_element_final (g%el(i))
    end do
    deallocate (g%el)
  end subroutine graph_final

\eatline
\nwendcode{}\nwbegindocs{564}\nwdocspar
\subsubsection{Access contents}
The number of elements.
\nwenddocs{}\nwbegincode{565}\moddef{Analysis: procedures}\plusendmoddef
  function graph_get_n_elements (graph) result (n)
    integer :: n
    type(graph_t), intent(in) :: graph
    n = size (graph%el)
  end function graph_get_n_elements

\eatline
\nwendcode{}\nwbegindocs{566}Retrieve a pointer to the drawing options of an element, so they can be
modified.  (The {\tt{}target} attribute is not actually needed because the
components are pointers.)
\nwenddocs{}\nwbegincode{567}\moddef{Analysis: procedures}\plusendmoddef
  function graph_get_drawing_options_ptr (g, i) result (ptr)
    type(drawing_options_t), pointer :: ptr
    type(graph_t), intent(in), target :: g
    integer, intent(in) :: i
    ptr => graph_element_get_drawing_options_ptr (g%el(i))
  end function graph_get_drawing_options_ptr

\eatline
\nwendcode{}\nwbegindocs{568}\nwdocspar
\subsubsection{Output}
The default output format just writes histogram and plot data.
\nwenddocs{}\nwbegincode{569}\moddef{Analysis: procedures}\plusendmoddef
  subroutine graph_write (graph, unit)
    type(graph_t), intent(in) :: graph
    integer, intent(in), optional :: unit
    integer :: i
    do i = 1, size (graph%el)
       call graph_element_write (graph%el(i), unit)
    end do
  end subroutine graph_write

\eatline
\nwendcode{}\nwbegindocs{570}The GAMELAN driver is not a simple wrapper, but it writes the plot/histogram
contents embedded the complete graph.  First, data are read in, global
background commands next, then individual elements, then global foreground
commands.
\nwenddocs{}\nwbegincode{571}\moddef{Analysis: procedures}\plusendmoddef
  subroutine graph_write_gml_driver (g, filename, unit)
    type(graph_t), intent(in) :: g
    type(string_t), intent(in) :: filename
    type(string_t) :: calc_cmd, bg_cmd, draw_cmd, err_cmd, symb_cmd, fg_cmd
    integer, intent(in), optional :: unit
    type(drawing_options_t), pointer :: dro
    integer :: u, i
    u = given_output_unit (unit);  if (u < 0)  return
    call graph_options_write_tex_header (g%graph_options, unit)
    write (u, "(2x,A)") &
         char (graph_options_get_gml_setup (g%graph_options))
    write (u, "(2x,A)") &
         char (graph_options_get_gml_graphrange (g%graph_options))
    do i = 1, size (g%el)
       call graph_element_write_gml_reader (g%el(i), filename, unit)
       calc_cmd = drawing_options_get_calc_command &
                       (graph_element_get_drawing_options_ptr (g%el(i)))
       if (calc_cmd /= "")  write (u, "(2x,A)")  char (calc_cmd)
    end do
    bg_cmd = graph_options_get_gml_bg_command (g%graph_options)
    if (bg_cmd /= "")  write (u, "(2x,A)")  char (bg_cmd)
    do i = 1, size (g%el)
       dro => graph_element_get_drawing_options_ptr (g%el(i))
       bg_cmd = drawing_options_get_gml_bg_command (dro)
       if (bg_cmd /= "")  write (u, "(2x,A)")  char (bg_cmd)
       draw_cmd = drawing_options_get_draw_command (dro)
       if (draw_cmd /= "")  write (u, "(2x,A)")  char (draw_cmd)
       err_cmd = drawing_options_get_err_command (dro)
       if (err_cmd /= "")  write (u, "(2x,A)")  char (err_cmd)
       symb_cmd = drawing_options_get_symb_command (dro)
       if (symb_cmd /= "")  write (u, "(2x,A)")  char (symb_cmd)
       fg_cmd = drawing_options_get_gml_fg_command (dro)
       if (fg_cmd /= "")  write (u, "(2x,A)")  char (fg_cmd)
    end do
    fg_cmd = graph_options_get_gml_fg_command (g%graph_options)
    if (fg_cmd /= "")  write (u, "(2x,A)")  char (fg_cmd)
    write (u, "(2x,A)")  char (graph_options_get_gml_x_label (g%graph_options))
    write (u, "(2x,A)")  char (graph_options_get_gml_y_label (g%graph_options))
    call graph_options_write_tex_footer (g%graph_options, unit)
  end subroutine graph_write_gml_driver

\eatline
\nwendcode{}\nwbegindocs{572}Append header for generic data output in ifile format.
\nwenddocs{}\nwbegincode{573}\moddef{Analysis: procedures}\plusendmoddef
  subroutine graph_get_header (graph, header, comment)
    type(graph_t), intent(in) :: graph
    type(ifile_t), intent(inout) :: header
    type(string_t), intent(in), optional :: comment
    type(string_t) :: c
    if (present (comment)) then
       c = comment
    else
       c = ""
    end if
    call ifile_append (header, c // "WHIZARD graph data")
    call graph_options_get_header (graph%graph_options, header, comment)
    call ifile_append (header, &
         c // "number of graph elements: " &
         // int2char (graph_get_n_elements (graph)))
  end subroutine graph_get_header

\eatline
\nwendcode{}\nwbegindocs{574}\nwdocspar
\subsection{Analysis objects}
This data structure holds all observables, histograms and such that
are currently active.  We have one global store; individual items are
identified by their ID strings.

(This should rather be coded by type extension.)
\nwenddocs{}\nwbegincode{575}\moddef{Analysis: parameters}\plusendmoddef
  integer, parameter :: AN_UNDEFINED = 0
  integer, parameter :: AN_OBSERVABLE = 1
  integer, parameter :: AN_HISTOGRAM = 2
  integer, parameter :: AN_PLOT = 3
  integer, parameter :: AN_GRAPH = 4
\nwendcode{}\nwbegincode{576}\moddef{Analysis: public}\plusendmoddef
  public :: AN_UNDEFINED, AN_HISTOGRAM, AN_OBSERVABLE, AN_PLOT, AN_GRAPH

\eatline
\eatline
\nwendcode{}\nwbegincode{577}\moddef{Analysis: types}\plusendmoddef
  type :: analysis_object_t
     private
     type(string_t) :: id
     integer :: type = AN_UNDEFINED
     type(observable_t), pointer :: obs => null ()
     type(histogram_t), pointer :: h => null ()
     type(plot_t), pointer :: p => null ()
     type(graph_t), pointer :: g => null ()
     type(analysis_object_t), pointer :: next => null ()
  end type analysis_object_t

\eatline
\nwendcode{}\nwbegindocs{578}\nwdocspar
\subsubsection{Initializer/finalizer}
Allocate with the correct type but do not fill initial values.
\nwenddocs{}\nwbegincode{579}\moddef{Analysis: procedures}\plusendmoddef
  subroutine analysis_object_init (obj, id, type)
    type(analysis_object_t), intent(out) :: obj
    type(string_t), intent(in) :: id
    integer, intent(in) :: type
    obj%id = id
    obj%type = type
    select case (obj%type)
    case (AN_OBSERVABLE);  allocate (obj%obs)
    case (AN_HISTOGRAM);   allocate (obj%h)
    case (AN_PLOT);        allocate (obj%p)
    case (AN_GRAPH);       allocate (obj%g)
    end select
  end subroutine analysis_object_init

\eatline
\nwendcode{}\nwbegincode{580}\moddef{Analysis: procedures}\plusendmoddef
  subroutine analysis_object_final (obj)
    type(analysis_object_t), intent(inout) :: obj
    select case (obj%type)
    case (AN_OBSERVABLE)
       deallocate (obj%obs)
    case (AN_HISTOGRAM)
       deallocate (obj%h)
    case (AN_PLOT)
       call plot_final (obj%p)
       deallocate (obj%p)
    case (AN_GRAPH)
       call graph_final (obj%g)
       deallocate (obj%g)
    end select
    obj%type = AN_UNDEFINED
  end subroutine analysis_object_final

\eatline
\nwendcode{}\nwbegindocs{581}Clear the analysis object, i.e., reset it to its initial state.  Not
applicable to graphs, which are always combinations of other existing
objects.
\nwenddocs{}\nwbegincode{582}\moddef{Analysis: procedures}\plusendmoddef
  subroutine analysis_object_clear (obj)
    type(analysis_object_t), intent(inout) :: obj
    select case (obj%type)
    case (AN_OBSERVABLE)
       call observable_clear (obj%obs)
    case (AN_HISTOGRAM)
       call histogram_clear (obj%h)
    case (AN_PLOT)
       call plot_clear (obj%p)
    end select
  end subroutine analysis_object_clear

\eatline
\nwendcode{}\nwbegindocs{583}\nwdocspar
\subsubsection{Fill with data}
Record data.  The effect depends on the type of analysis object.
\nwenddocs{}\nwbegincode{584}\moddef{Analysis: procedures}\plusendmoddef
  subroutine analysis_object_record_data (obj, &
       x, y, yerr, xerr, weight, excess, success)
    type(analysis_object_t), intent(inout) :: obj
    real(default), intent(in) :: x
    real(default), intent(in), optional :: y, yerr, xerr, weight, excess
    logical, intent(out), optional :: success
    select case (obj%type)
    case (AN_OBSERVABLE)
       if (present (weight)) then
          call observable_record_value_weighted (obj%obs, x, weight, success)
       else
          call observable_record_value_unweighted (obj%obs, x, success)
       end if
    case (AN_HISTOGRAM)
       if (present (weight)) then
          call histogram_record_value_weighted (obj%h, x, weight, success)
       else
          call histogram_record_value_unweighted (obj%h, x, excess, success)
       end if
    case (AN_PLOT)
       if (present (y)) then
          call plot_record_value (obj%p, x, y, yerr, xerr, success)
       else
          if (present (success))  success = .false.
       end if
    case default
       if (present (success))  success = .false.
    end select
  end subroutine analysis_object_record_data

\eatline
\nwendcode{}\nwbegindocs{585}Explicitly set the pointer to the next object in the list.
\nwenddocs{}\nwbegincode{586}\moddef{Analysis: procedures}\plusendmoddef
  subroutine analysis_object_set_next_ptr (obj, next)
    type(analysis_object_t), intent(inout) :: obj
    type(analysis_object_t), pointer :: next
    obj%next => next
  end subroutine analysis_object_set_next_ptr

\eatline
\nwendcode{}\nwbegindocs{587}\nwdocspar
\subsubsection{Access contents}
Return a pointer to the next object in the list.
\nwenddocs{}\nwbegincode{588}\moddef{Analysis: procedures}\plusendmoddef
  function analysis_object_get_next_ptr (obj) result (next)
    type(analysis_object_t), pointer :: next
    type(analysis_object_t), intent(in) :: obj
    next => obj%next
  end function analysis_object_get_next_ptr

\eatline
\nwendcode{}\nwbegindocs{589}Return data as appropriate for the object type.
\nwenddocs{}\nwbegincode{590}\moddef{Analysis: procedures}\plusendmoddef
  function analysis_object_get_n_elements (obj) result (n)
    integer :: n
    type(analysis_object_t), intent(in) :: obj
    select case (obj%type)
    case (AN_HISTOGRAM)
       n = 1
    case (AN_PLOT)
       n = 1
    case (AN_GRAPH)
       n = graph_get_n_elements (obj%g)
    case default
       n = 0
    end select
  end function analysis_object_get_n_elements

  function analysis_object_get_n_entries (obj, within_bounds) result (n)
    integer :: n
    type(analysis_object_t), intent(in) :: obj
    logical, intent(in), optional :: within_bounds
    logical :: wb
    select case (obj%type)
    case (AN_OBSERVABLE)
       n = observable_get_n_entries (obj%obs)
    case (AN_HISTOGRAM)
       wb = .false.;  if (present (within_bounds)) wb = within_bounds
       if (wb) then
          n = histogram_get_n_entries_within_bounds (obj%h)
       else
          n = histogram_get_n_entries (obj%h)
       end if
    case (AN_PLOT)
       n = plot_get_n_entries (obj%p)
    case default
       n = 0
    end select
  end function analysis_object_get_n_entries

  function analysis_object_get_average (obj, within_bounds) result (avg)
    real(default) :: avg
    type(analysis_object_t), intent(in) :: obj
    logical, intent(in), optional :: within_bounds
    logical :: wb
    select case (obj%type)
    case (AN_OBSERVABLE)
       avg = observable_get_average (obj%obs)
    case (AN_HISTOGRAM)
       wb = .false.;  if (present (within_bounds)) wb = within_bounds
       if (wb) then
          avg = histogram_get_average_within_bounds (obj%h)
       else
          avg = histogram_get_average (obj%h)
       end if
    case default
       avg = 0
    end select
  end function analysis_object_get_average

  function analysis_object_get_error (obj, within_bounds) result (err)
    real(default) :: err
    type(analysis_object_t), intent(in) :: obj
    logical, intent(in), optional :: within_bounds
    logical :: wb
    select case (obj%type)
    case (AN_OBSERVABLE)
       err = observable_get_error (obj%obs)
    case (AN_HISTOGRAM)
       wb = .false.;  if (present (within_bounds)) wb = within_bounds
       if (wb) then
          err = histogram_get_error_within_bounds (obj%h)
       else
          err = histogram_get_error (obj%h)
       end if
    case default
       err = 0
    end select
  end function analysis_object_get_error

\eatline
\eatline
\eatline
\eatline
\nwendcode{}\nwbegindocs{591}Return pointers to the actual contents:
\nwenddocs{}\nwbegincode{592}\moddef{Analysis: procedures}\plusendmoddef
  function analysis_object_get_observable_ptr (obj) result (obs)
    type(observable_t), pointer :: obs
    type(analysis_object_t), intent(in) :: obj
    select case (obj%type)
    case (AN_OBSERVABLE);  obs => obj%obs
    case default;          obs => null ()
    end select
  end function analysis_object_get_observable_ptr

  function analysis_object_get_histogram_ptr (obj) result (h)
    type(histogram_t), pointer :: h
    type(analysis_object_t), intent(in) :: obj
    select case (obj%type)
    case (AN_HISTOGRAM);  h => obj%h
    case default;         h => null ()
    end select
  end function analysis_object_get_histogram_ptr

  function analysis_object_get_plot_ptr (obj) result (plot)
    type(plot_t), pointer :: plot
    type(analysis_object_t), intent(in) :: obj
    select case (obj%type)
    case (AN_PLOT);  plot => obj%p
    case default;    plot => null ()
    end select
  end function analysis_object_get_plot_ptr

  function analysis_object_get_graph_ptr (obj) result (g)
    type(graph_t), pointer :: g
    type(analysis_object_t), intent(in) :: obj
    select case (obj%type)
    case (AN_GRAPH);  g => obj%g
    case default;     g => null ()
    end select
  end function analysis_object_get_graph_ptr

\eatline
\eatline
\eatline
\eatline
\nwendcode{}\nwbegindocs{593}Return true if the object has a graphical representation:
\nwenddocs{}\nwbegincode{594}\moddef{Analysis: procedures}\plusendmoddef
  function analysis_object_has_plot (obj) result (flag)
    logical :: flag
    type(analysis_object_t), intent(in) :: obj
    select case (obj%type)
    case (AN_HISTOGRAM);  flag = .true.
    case (AN_PLOT);       flag = .true.
    case (AN_GRAPH);      flag = .true.
    case default;         flag = .false.
    end select
  end function analysis_object_has_plot

\eatline
\nwendcode{}\nwbegindocs{595}\nwdocspar
\subsubsection{Output}
\nwenddocs{}\nwbegincode{596}\moddef{Analysis: procedures}\plusendmoddef
  subroutine analysis_object_write (obj, unit, verbose)
    type(analysis_object_t), intent(in) :: obj
    integer, intent(in), optional :: unit
    logical, intent(in), optional :: verbose
    logical :: verb
    integer :: u
    u = given_output_unit (unit);  if (u < 0)  return
    verb = .false.;  if (present (verbose))  verb = verbose
    write (u, "(A)")  repeat ("#", 79)
    select case (obj%type)
    case (AN_OBSERVABLE)
       write (u, "(A)", advance="no")  "# Observable:"
    case (AN_HISTOGRAM)
       write (u, "(A)", advance="no")  "# Histogram: "
    case (AN_PLOT)
       write (u, "(A)", advance="no")  "# Plot: "
    case (AN_GRAPH)
       write (u, "(A)", advance="no")  "# Graph: "
    case default
       write (u, "(A)") "# [undefined analysis object]"
       return
    end select
    write (u, "(1x,A)")  char (obj%id)
    select case (obj%type)
    case (AN_OBSERVABLE)
       call observable_write (obj%obs, unit)
    case (AN_HISTOGRAM)
       if (verb) then
          call graph_options_write (obj%h%graph_options, unit)
          write (u, *)
          call drawing_options_write (obj%h%drawing_options, unit)
          write (u, *)
       end if
       call histogram_write (obj%h, unit)
    case (AN_PLOT)
       if (verb) then
          call graph_options_write (obj%p%graph_options, unit)
          write (u, *)
          call drawing_options_write (obj%p%drawing_options, unit)
          write (u, *)
       end if
       call plot_write (obj%p, unit)
    case (AN_GRAPH)
       call graph_write (obj%g, unit)
    end select
  end subroutine analysis_object_write

\eatline
\nwendcode{}\nwbegindocs{597}Write the object part of the \LaTeX\ driver file.
\nwenddocs{}\nwbegincode{598}\moddef{Analysis: procedures}\plusendmoddef
  subroutine analysis_object_write_driver (obj, filename, unit)
    type(analysis_object_t), intent(in) :: obj
    type(string_t), intent(in) :: filename
    integer, intent(in), optional :: unit
    select case (obj%type)
    case (AN_OBSERVABLE)
       call observable_write_driver (obj%obs, unit)
    case (AN_HISTOGRAM)
       call histogram_write_gml_driver (obj%h, filename, unit)
    case (AN_PLOT)
       call plot_write_gml_driver (obj%p, filename, unit)
    case (AN_GRAPH)
       call graph_write_gml_driver (obj%g, filename, unit)
    end select
  end subroutine analysis_object_write_driver

\eatline
\nwendcode{}\nwbegindocs{599}Return a data header for external formats, in ifile form.
\nwenddocs{}\nwbegincode{600}\moddef{Analysis: procedures}\plusendmoddef
  subroutine analysis_object_get_header (obj, header, comment)
    type(analysis_object_t), intent(in) :: obj
    type(ifile_t), intent(inout) :: header
    type(string_t), intent(in), optional :: comment
    select case (obj%type)
    case (AN_HISTOGRAM)
       call histogram_get_header (obj%h, header, comment)
    case (AN_PLOT)
       call plot_get_header (obj%p, header, comment)
    end select
  end subroutine analysis_object_get_header

\eatline
\nwendcode{}\nwbegindocs{601}\nwdocspar
\subsection{Analysis object iterator}
Analysis objects are containers which have iterable data structures:
histograms/bins and plots/points.  If they are to be treated on a common
basis, it is useful to have an iterator which hides the implementation
details.

The iterator is used only for elementary analysis objects that contain plot
data: histograms or plots.  It is invalid for meta-objects (graphs) and
non-graphical objects (observables).
\nwenddocs{}\nwbegincode{602}\moddef{Analysis: public}\plusendmoddef
  public :: analysis_iterator_t
\nwendcode{}\nwbegincode{603}\moddef{Analysis: types}\plusendmoddef
  type :: analysis_iterator_t
    private
    integer :: type = AN_UNDEFINED
    type(analysis_object_t), pointer :: object => null ()
    integer :: index = 1
    type(point_t), pointer :: point => null ()
  end type

\eatline
\nwendcode{}\nwbegindocs{604}The initializer places the iterator at the beginning of the analysis object.
\nwenddocs{}\nwbegincode{605}\moddef{Analysis: procedures}\plusendmoddef
  subroutine analysis_iterator_init (iterator, object)
    type(analysis_iterator_t), intent(out) :: iterator
    type(analysis_object_t), intent(in), target :: object
    iterator%object => object
    if (associated (iterator%object)) then
       iterator%type = iterator%object%type
       select case (iterator%type)
       case (AN_PLOT)
          iterator%point => iterator%object%p%first
       end select
    end if
  end subroutine analysis_iterator_init

\eatline
\nwendcode{}\nwbegindocs{606}The iterator is valid as long as it points to an existing entry.  An
iterator for a data object without array data (observable) is always invalid.
\nwenddocs{}\nwbegincode{607}\moddef{Analysis: public}\plusendmoddef
  public :: analysis_iterator_is_valid
\nwendcode{}\nwbegincode{608}\moddef{Analysis: procedures}\plusendmoddef
  function analysis_iterator_is_valid (iterator) result (valid)
    logical :: valid
    type(analysis_iterator_t), intent(in) :: iterator
    if (associated (iterator%object)) then
       select case (iterator%type)
       case (AN_HISTOGRAM)
          valid = iterator%index <= histogram_get_n_bins (iterator%object%h)
       case (AN_PLOT)
          valid = associated (iterator%point)
       case default
          valid = .false.
       end select
    else
       valid = .false.
    end if
  end function analysis_iterator_is_valid

\eatline
\nwendcode{}\nwbegindocs{609}Advance the iterator.
\nwenddocs{}\nwbegincode{610}\moddef{Analysis: public}\plusendmoddef
  public :: analysis_iterator_advance
\nwendcode{}\nwbegincode{611}\moddef{Analysis: procedures}\plusendmoddef
  subroutine analysis_iterator_advance (iterator)
    type(analysis_iterator_t), intent(inout) :: iterator
    if (associated (iterator%object)) then
       select case (iterator%type)
       case (AN_PLOT)
          iterator%point => iterator%point%next
       end select
       iterator%index = iterator%index + 1
    end if
  end subroutine analysis_iterator_advance

\eatline
\nwendcode{}\nwbegindocs{612}Retrieve the object type:
\nwenddocs{}\nwbegincode{613}\moddef{Analysis: public}\plusendmoddef
  public :: analysis_iterator_get_type
\nwendcode{}\nwbegincode{614}\moddef{Analysis: procedures}\plusendmoddef
  function analysis_iterator_get_type (iterator) result (type)
    integer :: type
    type(analysis_iterator_t), intent(in) :: iterator
    type = iterator%type
  end function analysis_iterator_get_type

\eatline
\nwendcode{}\nwbegindocs{615}Use the iterator to retrieve data.  We implement a common routine which
takes the data descriptors as optional arguments.  Data which do not occur in
the selected type trigger to an error condition.

The iterator must point to a valid entry.
\nwenddocs{}\nwbegincode{616}\moddef{Analysis: public}\plusendmoddef
  public :: analysis_iterator_get_data
\nwendcode{}\nwbegincode{617}\moddef{Analysis: procedures}\plusendmoddef
  subroutine analysis_iterator_get_data (iterator, &
       x, y, yerr, xerr, width, excess, index, n_total)
    type(analysis_iterator_t), intent(in) :: iterator
    real(default), intent(out), optional :: x, y, yerr, xerr, width, excess
    integer, intent(out), optional :: index, n_total
    select case (iterator%type)
    case (AN_HISTOGRAM)
       if (present (x)) &
            x = bin_get_midpoint (iterator%object%h%bin(iterator%index))
       if (present (y)) &
            y = bin_get_sum (iterator%object%h%bin(iterator%index))
       if (present (yerr)) &
            yerr = bin_get_error (iterator%object%h%bin(iterator%index))
       if (present (xerr)) &
            call invalid ("histogram", "xerr")
       if (present (width)) &
            width = bin_get_width (iterator%object%h%bin(iterator%index))
       if (present (excess)) &
            excess = bin_get_excess (iterator%object%h%bin(iterator%index))
       if (present (index)) &
            index = iterator%index
       if (present (n_total)) &
            n_total = histogram_get_n_bins (iterator%object%h)
    case (AN_PLOT)
       if (present (x)) &
            x = point_get_x (iterator%point)
       if (present (y)) &
            y = point_get_y (iterator%point)
       if (present (yerr)) &
            yerr = point_get_yerr (iterator%point)
       if (present (xerr)) &
            xerr = point_get_xerr (iterator%point)
       if (present (width)) &
            call invalid ("plot", "width")
       if (present (excess)) &
            call invalid ("plot", "excess")
       if (present (index)) &
            index = iterator%index
       if (present (n_total)) &
            n_total = plot_get_n_entries (iterator%object%p)
    case default
       call msg_bug ("analysis_iterator_get_data: called " &
            // "for unsupported analysis object type")
    end select
  contains
    subroutine invalid (typestr, objstr)
      character(*), intent(in) :: typestr, objstr
      call msg_bug ("analysis_iterator_get_data: attempt to get '" &
           // objstr // "' for type '" // typestr // "'")
    end subroutine invalid
  end subroutine analysis_iterator_get_data

\eatline
\nwendcode{}\nwbegindocs{618}\nwdocspar
\subsection{Analysis store}
This data structure holds all observables, histograms and such that
are currently active.  We have one global store; individual items are
identified by their ID strings and types.
\nwenddocs{}\nwbegincode{619}\moddef{Analysis: variables}\endmoddef
  type(analysis_store_t), save :: analysis_store

\eatline
\nwendcode{}\nwbegincode{620}\moddef{Analysis: types}\plusendmoddef
  type :: analysis_store_t
     private
     type(analysis_object_t), pointer :: first => null ()
     type(analysis_object_t), pointer :: last => null ()
  end type analysis_store_t

\eatline
\nwendcode{}\nwbegindocs{621}Delete the analysis store
\nwenddocs{}\nwbegincode{622}\moddef{Analysis: public}\plusendmoddef
  public :: analysis_final
\nwendcode{}\nwbegincode{623}\moddef{Analysis: procedures}\plusendmoddef
  subroutine analysis_final ()
    type(analysis_object_t), pointer :: current
    do while (associated (analysis_store%first))
       current => analysis_store%first
       analysis_store%first => current%next
       call analysis_object_final (current)
    end do
    analysis_store%last => null ()
  end subroutine analysis_final

\eatline
\nwendcode{}\nwbegindocs{624}Append a new analysis object
\nwenddocs{}\nwbegincode{625}\moddef{Analysis: procedures}\plusendmoddef
  subroutine analysis_store_append_object (id, type)
    type(string_t), intent(in) :: id
    integer, intent(in) :: type
    type(analysis_object_t), pointer :: obj
    allocate (obj)
    call analysis_object_init (obj, id, type)
    if (associated (analysis_store%last)) then
       analysis_store%last%next => obj
    else
       analysis_store%first => obj
    end if
    analysis_store%last => obj
  end subroutine analysis_store_append_object

\eatline
\nwendcode{}\nwbegindocs{626}Return a pointer to the analysis object with given ID.
\nwenddocs{}\nwbegincode{627}\moddef{Analysis: procedures}\plusendmoddef
  function analysis_store_get_object_ptr (id) result (obj) 
    type(string_t), intent(in) :: id
    type(analysis_object_t), pointer :: obj
    obj => analysis_store%first
    do while (associated (obj))
       if (obj%id == id)  return
       obj => obj%next
    end do
  end function analysis_store_get_object_ptr

\eatline
\nwendcode{}\nwbegindocs{628}Initialize an analysis object: either reset it if present, or append
a new entry.
\nwenddocs{}\nwbegincode{629}\moddef{Analysis: procedures}\plusendmoddef
  subroutine analysis_store_init_object (id, type, obj)
    type(string_t), intent(in) :: id
    integer, intent(in) :: type
    type(analysis_object_t), pointer :: obj, next
    obj => analysis_store_get_object_ptr (id)
    if (associated (obj)) then
       next => analysis_object_get_next_ptr (obj)
       call analysis_object_final (obj)
       call analysis_object_init (obj, id, type)
       call analysis_object_set_next_ptr (obj, next)
    else
       call analysis_store_append_object (id, type)
       obj => analysis_store%last
    end if
  end subroutine analysis_store_init_object

\eatline
\nwendcode{}\nwbegindocs{630}Get the type of a analysis object
\nwenddocs{}\nwbegincode{631}\moddef{Analysis: public}\plusendmoddef
  public :: analysis_store_get_object_type
\nwendcode{}\nwbegincode{632}\moddef{Analysis: procedures}\plusendmoddef
  function analysis_store_get_object_type (id) result (type)
    type(string_t), intent(in) :: id
    integer :: type
    type(analysis_object_t), pointer :: object
    object => analysis_store_get_object_ptr (id)
    if (associated (object)) then
       type = object%type
    else
       type = AN_UNDEFINED
    end if
  end function analysis_store_get_object_type

\eatline
\nwendcode{}\nwbegindocs{633}Return the number of objects in the store.
\nwenddocs{}\nwbegincode{634}\moddef{Analysis: procedures}\plusendmoddef
  function analysis_store_get_n_objects () result (n)
    integer :: n
    type(analysis_object_t), pointer :: current
    n = 0
    current => analysis_store%first
    do while (associated (current))
       n = n + 1
       current => current%next
    end do
  end function analysis_store_get_n_objects

\eatline
\nwendcode{}\nwbegindocs{635}Allocate an array and fill it with all existing IDs.
\nwenddocs{}\nwbegincode{636}\moddef{Analysis: public}\plusendmoddef
  public :: analysis_store_get_ids
\nwendcode{}\nwbegincode{637}\moddef{Analysis: procedures}\plusendmoddef
  subroutine analysis_store_get_ids (id)
    type(string_t), dimension(:), allocatable, intent(out) :: id
    type(analysis_object_t), pointer :: current
    integer :: i
    allocate (id (analysis_store_get_n_objects()))
    i = 0
    current => analysis_store%first
    do while (associated (current))
       i = i + 1
       id(i) = current%id
       current => current%next
    end do
  end subroutine analysis_store_get_ids

\eatline
\nwendcode{}\nwbegindocs{638}\nwdocspar
\subsection{\LaTeX\ driver file}
Write a driver file for all objects in the store.
\nwenddocs{}\nwbegincode{639}\moddef{Analysis: procedures}\plusendmoddef
  subroutine analysis_store_write_driver_all (filename_data, unit)
    type(string_t), intent(in) :: filename_data
    integer, intent(in), optional :: unit
    type(analysis_object_t), pointer :: obj
    call analysis_store_write_driver_header (unit)
    obj => analysis_store%first
    do while (associated (obj))
       call analysis_object_write_driver (obj, filename_data, unit)
       obj => obj%next
    end do
    call analysis_store_write_driver_footer (unit)
  end subroutine analysis_store_write_driver_all

\eatline
\nwendcode{}\nwbegindocs{640}\nwdocspar
Write a driver file for an array of objects.
\nwenddocs{}\nwbegincode{641}\moddef{Analysis: procedures}\plusendmoddef
  subroutine analysis_store_write_driver_obj (filename_data, id, unit)
    type(string_t), intent(in) :: filename_data
    type(string_t), dimension(:), intent(in) :: id
    integer, intent(in), optional :: unit
    type(analysis_object_t), pointer :: obj
    integer :: i
    call analysis_store_write_driver_header (unit)
    do i = 1, size (id)
       obj => analysis_store_get_object_ptr (id(i))
       if (associated (obj))  &
            call analysis_object_write_driver (obj, filename_data, unit)
    end do
    call analysis_store_write_driver_footer (unit)
  end subroutine analysis_store_write_driver_obj

\eatline
\nwendcode{}\nwbegindocs{642}The beginning of the driver file.
\nwenddocs{}\nwbegincode{643}\moddef{Analysis: procedures}\plusendmoddef
  subroutine analysis_store_write_driver_header (unit)
    integer, intent(in), optional :: unit
    integer :: u
    u = given_output_unit (unit);  if (u < 0)  return
    write (u, '(A)') "\\documentclass[12pt]\{article\}"
    write (u, *)
    write (u, '(A)') "\\usepackage\{gamelan\}"
    write (u, '(A)') "\\usepackage\{amsmath\}"
    write (u, '(A)') "\\usepackage\{ifpdf\}"
    write (u, '(A)') "\\ifpdf"
    write (u, '(A)') "   \\DeclareGraphicsRule\{*\}\{mps\}\{*\}\{\}"
    write (u, '(A)') "\\else"
    write (u, '(A)') "   \\DeclareGraphicsRule\{*\}\{eps\}\{*\}\{\}"
    write (u, '(A)') "\\fi"    
    write (u, *)
    write (u, '(A)') "\\begin\{document\}"
    write (u, '(A)') "\\begin\{gmlfile\}"
    write (u, *)
    write (u, '(A)') "\\begin\{gmlcode\}"
    write (u, '(A)') "  color col.default, col.excess;"
    write (u, '(A)') "  col.default = 0.9white;"
    write (u, '(A)') "  col.excess  = red;"
    write (u, '(A)') "  boolean show_excess;"
    !!! Future excess options for plots
    ! if (mcs(1)%plot_excess .and. mcs(1)%unweighted) then
    !    write (u, '(A)') "  show_excess = true;"
    ! else
    write (u, '(A)') "  show_excess = false;"
    ! end if
    write (u, '(A)') "\\end\{gmlcode\}"
    write (u, *)
  end subroutine analysis_store_write_driver_header

\eatline
\nwendcode{}\nwbegindocs{644}The end of the driver file.
\nwenddocs{}\nwbegincode{645}\moddef{Analysis: procedures}\plusendmoddef
  subroutine analysis_store_write_driver_footer (unit)
    integer, intent(in), optional :: unit
    integer :: u
    u = given_output_unit (unit);  if (u < 0)  return
    write(u, *)
    write(u, '(A)') "\\end\{gmlfile\}"
    write(u, '(A)') "\\end\{document\}"
  end subroutine analysis_store_write_driver_footer

\eatline
\nwendcode{}\nwbegindocs{646}\nwdocspar
\subsection{API}
\subsubsection{Creating new objects}
The specific versions below:
\nwenddocs{}\nwbegincode{647}\moddef{Analysis: public}\plusendmoddef
  public :: analysis_init_observable
\nwendcode{}\nwbegincode{648}\moddef{Analysis: procedures}\plusendmoddef
  subroutine analysis_init_observable (id, obs_label, obs_unit, graph_options)
    type(string_t), intent(in) :: id
    type(string_t), intent(in), optional :: obs_label, obs_unit
    type(graph_options_t), intent(in), optional :: graph_options
    type(analysis_object_t), pointer :: obj
    type(observable_t), pointer :: obs
    call analysis_store_init_object (id, AN_OBSERVABLE, obj)
    obs => analysis_object_get_observable_ptr (obj)
    call observable_init (obs, obs_label, obs_unit, graph_options)
  end subroutine analysis_init_observable

\eatline
\nwendcode{}\nwbegincode{649}\moddef{Analysis: public}\plusendmoddef
  public :: analysis_init_histogram
\nwendcode{}\nwbegincode{650}\moddef{Analysis: interfaces}\plusendmoddef
  interface analysis_init_histogram
     module procedure analysis_init_histogram_n_bins
     module procedure analysis_init_histogram_bin_width
  end interface

\nwendcode{}\nwbegincode{651}\moddef{Analysis: procedures}\plusendmoddef
  subroutine analysis_init_histogram_n_bins &
       (id, lower_bound, upper_bound, n_bins, normalize_bins, &
        obs_label, obs_unit, graph_options, drawing_options)
    type(string_t), intent(in) :: id
    real(default), intent(in) :: lower_bound, upper_bound
    integer, intent(in) :: n_bins
    logical, intent(in) :: normalize_bins
    type(string_t), intent(in), optional :: obs_label, obs_unit
    type(graph_options_t), intent(in), optional :: graph_options
    type(drawing_options_t), intent(in), optional :: drawing_options
    type(analysis_object_t), pointer :: obj   
    type(histogram_t), pointer :: h
    call analysis_store_init_object (id, AN_HISTOGRAM, obj)
    h => analysis_object_get_histogram_ptr (obj)
    call histogram_init (h, id, &
         lower_bound, upper_bound, n_bins, normalize_bins, &
         obs_label, obs_unit, graph_options, drawing_options)
  end subroutine analysis_init_histogram_n_bins

  subroutine analysis_init_histogram_bin_width &
       (id, lower_bound, upper_bound, bin_width, normalize_bins, &
        obs_label, obs_unit, graph_options, drawing_options)
    type(string_t), intent(in) :: id
    real(default), intent(in) :: lower_bound, upper_bound, bin_width
    logical, intent(in) :: normalize_bins
    type(string_t), intent(in), optional :: obs_label, obs_unit
    type(graph_options_t), intent(in), optional :: graph_options
    type(drawing_options_t), intent(in), optional :: drawing_options
    type(analysis_object_t), pointer :: obj
    type(histogram_t), pointer :: h
    call analysis_store_init_object (id, AN_HISTOGRAM, obj)
    h => analysis_object_get_histogram_ptr (obj)
    call histogram_init (h, id, &
         lower_bound, upper_bound, bin_width, normalize_bins, &
         obs_label, obs_unit, graph_options, drawing_options)
  end subroutine analysis_init_histogram_bin_width

\eatline
\eatline
\nwendcode{}\nwbegincode{652}\moddef{Analysis: public}\plusendmoddef
  public :: analysis_init_plot
\nwendcode{}\nwbegincode{653}\moddef{Analysis: procedures}\plusendmoddef
  subroutine analysis_init_plot (id, graph_options, drawing_options)
    type(string_t), intent(in) :: id
    type(graph_options_t), intent(in), optional :: graph_options    
    type(drawing_options_t), intent(in), optional :: drawing_options
    type(analysis_object_t), pointer :: obj
    type(plot_t), pointer :: plot
    call analysis_store_init_object (id, AN_PLOT, obj)
    plot => analysis_object_get_plot_ptr (obj)
    call plot_init (plot, id, graph_options, drawing_options)
  end subroutine analysis_init_plot

\eatline
\nwendcode{}\nwbegincode{654}\moddef{Analysis: public}\plusendmoddef
  public :: analysis_init_graph
\nwendcode{}\nwbegincode{655}\moddef{Analysis: procedures}\plusendmoddef
  subroutine analysis_init_graph (id, n_elements, graph_options)
    type(string_t), intent(in) :: id
    integer, intent(in) :: n_elements
    type(graph_options_t), intent(in), optional :: graph_options    
    type(analysis_object_t), pointer :: obj
    type(graph_t), pointer :: graph
    call analysis_store_init_object (id, AN_GRAPH, obj)
    graph => analysis_object_get_graph_ptr (obj)
    call graph_init (graph, id, n_elements, graph_options)
  end subroutine analysis_init_graph

\eatline
\nwendcode{}\nwbegindocs{656}\nwdocspar
\subsubsection{Recording data}
This procedure resets an object or the whole store to its initial
state.
\nwenddocs{}\nwbegincode{657}\moddef{Analysis: public}\plusendmoddef
  public :: analysis_clear
\nwendcode{}\nwbegincode{658}\moddef{Analysis: interfaces}\plusendmoddef
  interface analysis_clear
     module procedure analysis_store_clear_obj
     module procedure analysis_store_clear_all
  end interface

\nwendcode{}\nwbegincode{659}\moddef{Analysis: procedures}\plusendmoddef
  subroutine analysis_store_clear_obj (id)
    type(string_t), intent(in) :: id
    type(analysis_object_t), pointer :: obj
    obj => analysis_store_get_object_ptr (id)
    if (associated (obj)) then
       call analysis_object_clear (obj)
    end if
  end subroutine analysis_store_clear_obj

  subroutine analysis_store_clear_all ()
    type(analysis_object_t), pointer :: obj
    obj => analysis_store%first
    do while (associated (obj))
       call analysis_object_clear (obj)
       obj => obj%next
    end do
  end subroutine analysis_store_clear_all

\eatline
\nwendcode{}\nwbegindocs{660}\nwdocspar
There is one generic recording function whose behavior depends on the
type of analysis object.
\nwenddocs{}\nwbegincode{661}\moddef{Analysis: public}\plusendmoddef
  public :: analysis_record_data
\nwendcode{}\nwbegincode{662}\moddef{Analysis: procedures}\plusendmoddef
  subroutine analysis_record_data (id, x, y, yerr, xerr, &
       weight, excess, success, exist)
    type(string_t), intent(in) :: id
    real(default), intent(in) :: x
    real(default), intent(in), optional :: y, yerr, xerr, weight, excess
    logical, intent(out), optional :: success, exist
    type(analysis_object_t), pointer :: obj
    obj => analysis_store_get_object_ptr (id)
    if (associated (obj)) then
       call analysis_object_record_data (obj, x, y, yerr, xerr, &
            weight, excess, success)
       if (present (exist))  exist = .true.
    else
       if (present (success))  success = .false.
       if (present (exist))  exist = .false.
    end if
  end subroutine analysis_record_data

\eatline
\nwendcode{}\nwbegindocs{663}\nwdocspar
\subsubsection{Build a graph}
This routine sets up the array of graph elements by copying the graph elements
given as input.  The object must exist and already be initialized as a graph.
\nwenddocs{}\nwbegincode{664}\moddef{Analysis: public}\plusendmoddef
  public :: analysis_fill_graph
\nwendcode{}\nwbegincode{665}\moddef{Analysis: procedures}\plusendmoddef
  subroutine analysis_fill_graph (id, i, id_in, drawing_options)
    type(string_t), intent(in) :: id
    integer, intent(in) :: i
    type(string_t), intent(in) :: id_in
    type(drawing_options_t), intent(in), optional :: drawing_options
    type(analysis_object_t), pointer :: obj
    type(graph_t), pointer :: g
    type(histogram_t), pointer :: h
    type(plot_t), pointer :: p
    obj => analysis_store_get_object_ptr (id)
    g => analysis_object_get_graph_ptr (obj)
    obj => analysis_store_get_object_ptr (id_in)
    if (associated (obj)) then
       select case (obj%type)
       case (AN_HISTOGRAM)
          h => analysis_object_get_histogram_ptr (obj)
          call graph_insert_histogram (g, i, h, drawing_options)
       case (AN_PLOT)
          p => analysis_object_get_plot_ptr (obj)
          call graph_insert_plot (g, i, p, drawing_options)
       case default
          call msg_error ("Graph '" // char (id) // "': Element '" &
               // char (id_in) // "' is neither histogram nor plot.")
       end select
    else
       call msg_error ("Graph '" // char (id) // "': Element '" &
               // char (id_in) // "' is undefined.")
    end if
  end subroutine analysis_fill_graph

\eatline
\nwendcode{}\nwbegindocs{666}\nwdocspar
\subsubsection{Retrieve generic results}
Check if a named object exists.
\nwenddocs{}\nwbegincode{667}\moddef{Analysis: public}\plusendmoddef
  public :: analysis_exists
\nwendcode{}\nwbegincode{668}\moddef{Analysis: procedures}\plusendmoddef
  function analysis_exists (id) result (flag) 
    type(string_t), intent(in) :: id
    logical :: flag
    type(analysis_object_t), pointer :: obj
    flag = .true.
    obj => analysis_store%first
    do while (associated (obj))
       if (obj%id == id)  return
       obj => obj%next
    end do
    flag = .false.
  end function analysis_exists

\eatline
\nwendcode{}\nwbegindocs{669}The following functions should work for all kinds of analysis object:
\nwenddocs{}\nwbegincode{670}\moddef{Analysis: public}\plusendmoddef
  public :: analysis_get_n_elements
  public :: analysis_get_n_entries
  public :: analysis_get_average
  public :: analysis_get_error
\nwendcode{}\nwbegincode{671}\moddef{Analysis: procedures}\plusendmoddef
  function analysis_get_n_elements (id) result (n)
    integer :: n
    type(string_t), intent(in) :: id
    type(analysis_object_t), pointer :: obj
    obj => analysis_store_get_object_ptr (id)
    if (associated (obj)) then
       n = analysis_object_get_n_elements (obj)
    else
       n = 0
    end if
  end function analysis_get_n_elements

  function analysis_get_n_entries (id, within_bounds) result (n)
    integer :: n
    type(string_t), intent(in) :: id
    logical, intent(in), optional :: within_bounds
    type(analysis_object_t), pointer :: obj
    obj => analysis_store_get_object_ptr (id)
    if (associated (obj)) then
       n = analysis_object_get_n_entries (obj, within_bounds)
    else
       n = 0
    end if
  end function analysis_get_n_entries

  function analysis_get_average (id, within_bounds) result (avg)
    real(default) :: avg
    type(string_t), intent(in) :: id
    type(analysis_object_t), pointer :: obj
    logical, intent(in), optional :: within_bounds
    obj => analysis_store_get_object_ptr (id)
    if (associated (obj)) then
       avg = analysis_object_get_average (obj, within_bounds)
    else
       avg = 0
    end if
  end function analysis_get_average

  function analysis_get_error (id, within_bounds) result (err)
    real(default) :: err
    type(string_t), intent(in) :: id
    type(analysis_object_t), pointer :: obj
    logical, intent(in), optional :: within_bounds
    obj => analysis_store_get_object_ptr (id)
    if (associated (obj)) then
       err = analysis_object_get_error (obj, within_bounds)
    else
       err = 0
    end if
  end function analysis_get_error

\eatline
\eatline
\eatline
\eatline
\nwendcode{}\nwbegindocs{672}Return true if any analysis object is graphical
\nwenddocs{}\nwbegincode{673}\moddef{Analysis: public}\plusendmoddef
  public :: analysis_has_plots
\nwendcode{}\nwbegincode{674}\moddef{Analysis: interfaces}\plusendmoddef
  interface analysis_has_plots
     module procedure analysis_has_plots_any
     module procedure analysis_has_plots_obj
  end interface

\nwendcode{}\nwbegincode{675}\moddef{Analysis: procedures}\plusendmoddef
  function analysis_has_plots_any () result (flag)
    logical :: flag
    type(analysis_object_t), pointer :: obj
    flag = .false.
    obj => analysis_store%first
    do while (associated (obj))
       flag = analysis_object_has_plot (obj)
       if (flag)  return
    end do
  end function analysis_has_plots_any

  function analysis_has_plots_obj (id) result (flag)
    logical :: flag
    type(string_t), dimension(:), intent(in) :: id
    type(analysis_object_t), pointer :: obj
    integer :: i
    flag = .false.
    do i = 1, size (id)
       obj => analysis_store_get_object_ptr (id(i))
       if (associated (obj)) then
          flag = analysis_object_has_plot (obj)
          if (flag)  return
       end if
    end do
  end function analysis_has_plots_obj

\eatline
\nwendcode{}\nwbegindocs{676}\nwdocspar
\subsubsection{Iterators}
Initialize an iterator for the given object.  If the object does not exist or
has wrong type, the iterator will be invalid.
\nwenddocs{}\nwbegincode{677}\moddef{Analysis: public}\plusendmoddef
  public :: analysis_init_iterator
\nwendcode{}\nwbegincode{678}\moddef{Analysis: procedures}\plusendmoddef
  subroutine analysis_init_iterator (id, iterator)
    type(string_t), intent(in) :: id
    type(analysis_iterator_t), intent(out) :: iterator
    type(analysis_object_t), pointer :: obj
    obj => analysis_store_get_object_ptr (id)
    if (associated (obj))  call analysis_iterator_init (iterator, obj)
  end subroutine analysis_init_iterator

\eatline
\nwendcode{}\nwbegindocs{679}\nwdocspar
\subsubsection{Output}
\nwenddocs{}\nwbegincode{680}\moddef{Analysis: public}\plusendmoddef
  public :: analysis_write
\nwendcode{}\nwbegincode{681}\moddef{Analysis: interfaces}\plusendmoddef
  interface analysis_write
     module procedure analysis_write_object
     module procedure analysis_write_all
  end interface

\eatline
\nwendcode{}\nwbegincode{682}\moddef{Analysis: procedures}\plusendmoddef
  subroutine analysis_write_object (id, unit, verbose)
    type(string_t), intent(in) :: id
    integer, intent(in), optional :: unit
    logical, intent(in), optional :: verbose
    type(analysis_object_t), pointer :: obj
    obj => analysis_store_get_object_ptr (id)
    if (associated (obj)) then
       call analysis_object_write (obj, unit, verbose)
    else
       call msg_error ("Analysis object '" // char (id) // "' not found")
    end if
  end subroutine analysis_write_object

  subroutine analysis_write_all (unit, verbose)
    integer, intent(in), optional :: unit
    logical, intent(in), optional :: verbose
    type(analysis_object_t), pointer :: obj
    integer :: u
    u = given_output_unit (unit);  if (u < 0)  return
    obj => analysis_store%first
    do while (associated (obj))
       call analysis_object_write (obj, unit, verbose)
       obj => obj%next
    end do
  end subroutine analysis_write_all

\eatline
\eatline
\nwendcode{}\nwbegincode{683}\moddef{Analysis: public}\plusendmoddef
  public :: analysis_write_driver
\nwendcode{}\nwbegincode{684}\moddef{Analysis: procedures}\plusendmoddef
  subroutine analysis_write_driver (filename_data, id, unit)
    type(string_t), intent(in) :: filename_data
    type(string_t), dimension(:), intent(in), optional :: id
    integer, intent(in), optional :: unit
    if (present (id)) then
       call analysis_store_write_driver_obj (filename_data, id, unit)
    else
       call analysis_store_write_driver_all (filename_data, unit)
    end if
  end subroutine analysis_write_driver

\eatline
\nwendcode{}\nwbegincode{685}\moddef{Analysis: public}\plusendmoddef
  public :: analysis_compile_tex
\nwendcode{}\nwbegincode{686}\moddef{Analysis: procedures}\plusendmoddef
  subroutine analysis_compile_tex (file, has_gmlcode, os_data)
    type(string_t), intent(in) :: file
    logical, intent(in) :: has_gmlcode
    type(os_data_t), intent(in) :: os_data
    type(string_t) :: setenv
    integer :: status
    if (os_data%event_analysis_ps) then
       BLOCK: do
          if (os_data%whizard_texpath /= "") then
             setenv = "TEXINPUTS=" // os_data%whizard_texpath // ":$TEXINPUTS "
          else
             setenv = ""
          end if
          call os_system_call (setenv // os_data%latex // " " // file, status)
          if (status /= 0)  exit BLOCK
          if (has_gmlcode) then
             call os_system_call (os_data%gml // " " // file, status)
             if (status /= 0)  exit BLOCK
             call os_system_call (setenv // os_data%latex // " " // file, &
                  status)
             if (status /= 0)  exit BLOCK
          end if
          call os_system_call (os_data%dvips // " -o " // file // ".ps " &
               // file, status)
          if (status /= 0)  exit BLOCK
          if (os_data%event_analysis_pdf) then
             call os_system_call (os_data%ps2pdf // " " // file // ".ps", &
                                  status)
             if (status /= 0)  exit BLOCK
          end if
          exit BLOCK
       end do BLOCK
       if (status /= 0) then
          call msg_error ("Unable to compile analysis output file")
       end if
    else
       call msg_warning ("Skipping results display because " &
            // "latex/mpost/dvips is not available")
    end if
  end subroutine analysis_compile_tex

\eatline
\nwendcode{}\nwbegindocs{687}Write header for generic data output to an ifile.
\nwenddocs{}\nwbegincode{688}\moddef{Analysis: public}\plusendmoddef
  public :: analysis_get_header
\nwendcode{}\nwbegincode{689}\moddef{Analysis: procedures}\plusendmoddef
  subroutine analysis_get_header (id, header, comment)
    type(string_t), intent(in) :: id
    type(ifile_t), intent(inout) :: header
    type(string_t), intent(in), optional :: comment
    type(analysis_object_t), pointer :: object
    object => analysis_store_get_object_ptr (id)
    if (associated (object)) then
       call analysis_object_get_header (object, header, comment)
    end if
  end subroutine analysis_get_header

\eatline
\nwendcode{}\nwbegindocs{690}\nwdocspar
\subsection{Unit tests}
\nwenddocs{}\nwbegincode{691}\moddef{Analysis: public}\plusendmoddef
  public :: analysis_test
\nwendcode{}\nwbegincode{692}\moddef{Analysis: procedures}\plusendmoddef
  subroutine analysis_test (u, results)
    integer, intent(in) :: u
    type(test_results_t), intent(inout) :: results
  \LA{}Analysis: execute tests\RA{}  
  end subroutine analysis_test

\eatline
\nwendcode{}\nwbegincode{693}\moddef{Analysis: execute tests}\endmoddef
  call test (analysis_1, "analysis_1", &
       "check elementary analysis building blocks", &
       u, results)
\nwendcode{}\nwbegincode{694}\moddef{Analysis: tests}\endmoddef
  subroutine analysis_1 (u)
    integer, intent(in) :: u
    type(string_t) :: id1, id2, id3, id4
    integer :: i
    id1 = "foo"
    id2 = "bar"
    id3 = "hist"
    id4 = "plot"
    
    write (u, "(A)")  "* Test output: Analysis"
    write (u, "(A)")  "*   Purpose: test the analysis routines"
    write (u, "(A)")        
    
    call analysis_init_observable (id1)
    call analysis_init_observable (id2)
    call analysis_init_histogram_bin_width &
         (id3, 0.5_default, 5.5_default, 1._default, normalize_bins=.false.)
    call analysis_init_plot (id4)
    do i = 1, 3
       write (u, "(A,1x," // FMT_19 // ")")  "data = ", real(i,default)
       call analysis_record_data (id1, real(i,default))
       call analysis_record_data (id2, real(i,default), &
                                        weight=real(i,default))
       call analysis_record_data (id3, real(i,default))
       call analysis_record_data (id4, real(i,default), real(i,default)**2)
    end do
    write (u, "(A,10(1x,I5))") "n_entries = ", &
         analysis_get_n_entries (id1), &
         analysis_get_n_entries (id2), &
         analysis_get_n_entries (id3), &
         analysis_get_n_entries (id3, within_bounds = .true.), &
         analysis_get_n_entries (id4), &
         analysis_get_n_entries (id4, within_bounds = .true.)
    write (u, "(A,10(1x," // FMT_19 // "))")  "average   = ", &
         analysis_get_average (id1), &
         analysis_get_average (id2), &
         analysis_get_average (id3), &
         analysis_get_average (id3, within_bounds = .true.)
    write (u, "(A,10(1x," // FMT_19 // "))")  "error     = ", &
         analysis_get_error (id1), &
         analysis_get_error (id2), &
         analysis_get_error (id3), &
         analysis_get_error (id3, within_bounds = .true.)

    write (u, "(A)")
    write (u, "(A)") "* Clear analysis #2"
    write (u, "(A)")
    
    call analysis_clear (id2)
    do i = 4, 6
       print *, "data = ", real(i,default)
       call analysis_record_data (id1, real(i,default))
       call analysis_record_data (id2, real(i,default), &
                                        weight=real(i,default))
       call analysis_record_data (id3, real(i,default))
       call analysis_record_data (id4, real(i,default), real(i,default)**2)
    end do
    write (u, "(A,10(1x,I5))")  "n_entries = ", &
         analysis_get_n_entries (id1), &
         analysis_get_n_entries (id2), &
         analysis_get_n_entries (id3), &
         analysis_get_n_entries (id3, within_bounds = .true.), &
         analysis_get_n_entries (id4), &
         analysis_get_n_entries (id4, within_bounds = .true.)
    write (u, "(A,10(1x," // FMT_19 // "))")  "average   = ", &
         analysis_get_average (id1), &
         analysis_get_average (id2), &
         analysis_get_average (id3), &
         analysis_get_average (id3, within_bounds = .true.)
    write (u, "(A,10(1x," // FMT_19 // "))")  "error     = ", &
         analysis_get_error (id1), &
         analysis_get_error (id2), &
         analysis_get_error (id3), &
         analysis_get_error (id3, within_bounds = .true.)
    write (u, "(A)")
    call analysis_write (u)

    write (u, "(A)")
    write (u, "(A)")  "* Cleanup"    
       
    call analysis_clear ()
    call analysis_final ()
    
    write (u, "(A)")
    write (u, "(A)")  "* Test output end: analysis_1"            
  end subroutine analysis_1

\eatline
\nwendcode{}\nwfilename{../../../src/matrix_elements/matrix_elements.nw}\nwbegindocs{0}% -*- ess-noweb-default-code-mode: f90-mode; noweb-default-code-mode: f90-mode; -*- 
% WHIZARD code as NOWEB source: matrix elements and process libraries

\chapter{Matrix Element Handling}
In this chapter, we support internal and external matrix elements:
initialization, automatic generation where necessary, and numerical
evaluation.  We provide the interface for code generation and linking.
Matrix-element code is organized in processes and process libraries.
\begin{description}
\item[process\_constants]
  A record of static process properties, for easy transfer between
  various \whizard\ modules.
\item[prclib\_interfaces]
  This module deals with matrix-element code which is accessible via
  external libraries (Fortran libraries or generic C-compatible
  libraries) and must either be generated by the program or provided
  by the user explicitly.

  The module defines and uses an abstract type {\tt{}prc{\char95}writer{\char95}t} and two
  abstract extensions, one for a Fortran module and one for a C-compatible
  library.   The implementation provides the specific methods for writing the  
  appropriate parts in external matrix element code.
\item[prc\_core\_def]
  This module defines the abstract types {\tt{}prc{\char95}core{\char95}def{\char95}t} and
  {\tt{}prc{\char95}driver{\char95}t}.  The implementation of the former provides the
  configuration for processes of a certain class, while the latter accesses
  the corresponding matrix element, in particular those generated by the
  appropriate {\tt{}prc{\char95}writer{\char95}t} object.
\item[process\_libraries] 
  This module combines the functionality of
  the previous module with the means for holding processes definitions
  (the internal counterpart of appropriate declarations in the user
  interface), for handling matrix elements which do not need external
  code, and for accessing the matrix elements by the procedures for
  matrix-element evaluation, integration and event generation.
\item[prclib\_stacks]
  Collect process libraries.
\item[test\_me] This module provides a test implementation for the abstract
  types in the {\tt{}prc{\char95}core{\char95}def} module.  The implementation is intended for
  self-tests of several later modules.  The implementation is internal, i.e.,
  no external code has is generated.
\end{description}
All data structures which are specific for a particular way of
generating code or evaluating matrix element are kept abstract and
thus generic.  Later modules such as {\tt{}prc{\char95}omega} provide
implementations, in the form of type extensions for the various
abstract types.

\clearpage
%%%%%%%%%%%%%%%%%%%%%%%%%%%%%%%%%%%%%%%%%%%%%%%%%%%%%%%%%%%%%%%%%%%%%%%%
\section{Process data block}

We define a simple transparent type that contains universal constant
process data.  We will reference objects of this type for the
phase-space setup, for interfacing with process libraries, for
implementing matrix-element generation, and in the master
process-handling module.
\nwenddocs{}\nwbegincode{1}\moddef{\code{}process{\char95}constants.f90\edoc{}}\endmoddef
\LA{}File header\RA{}

module process_constants
  
\LA{}Use kinds\RA{}
\LA{}Use strings\RA{}

\LA{}Standard module head\RA{}

\LA{}Process constants: public\RA{}

\LA{}Process constants: types\RA{}

contains

\LA{}Process constants: procedures\RA{}

end module process_constants
\eatline
\nwendcode{}\nwbegindocs{2}\nwdocspar
The data type is just a block of public objects, only elementary
types, no type-bound procedures.
\nwenddocs{}\nwbegincode{3}\moddef{Process constants: public}\endmoddef
  public :: process_constants_t
\nwendcode{}\nwbegincode{4}\moddef{Process constants: types}\endmoddef
  type :: process_constants_t
     type(string_t) :: id
     type(string_t) :: model_name
     character(32) :: md5sum = ""
     logical :: openmp_supported = .false.
     integer :: n_in  = 0
     integer :: n_out = 0
     integer :: n_flv = 0
     integer :: n_hel = 0
     integer :: n_col = 0
     integer :: n_cin = 0
     integer :: n_cf  = 0
     integer, dimension(:,:), allocatable :: flv_state
     integer, dimension(:,:), allocatable :: hel_state
     integer, dimension(:,:,:), allocatable :: col_state
     logical, dimension(:,:), allocatable :: ghost_flag
     complex(default), dimension(:), allocatable :: color_factors
     integer, dimension(:,:), allocatable :: cf_index
  contains
  \LA{}Process constants: prc const: TBP\RA{}
  end type process_constants_t

\eatline
\nwendcode{}\nwbegindocs{5}\nwdocspar
\nwenddocs{}\nwbegincode{6}\moddef{Process constants: prc const: TBP}\endmoddef
  procedure :: get_flv_state => process_constants_get_flv_state
\nwendcode{}\nwbegincode{7}\moddef{Process constants: procedures}\endmoddef
  function process_constants_get_flv_state (prc_const) result (flv_state)
    class(process_constants_t), intent(inout) :: prc_const
    integer, dimension(:,:), allocatable :: flv_state
    flv_state = prc_const%flv_state
  end function process_constants_get_flv_state

\eatline
\nwendcode{}\nwbegindocs{8}\nwdocspar
\nwenddocs{}\nwbegincode{9}\moddef{Process constants: prc const: TBP}\plusendmoddef
  procedure :: get_hel_state => process_constants_get_hel_state
\nwendcode{}\nwbegincode{10}\moddef{Process constants: procedures}\plusendmoddef
  subroutine process_constants_get_hel_state (prc_const, hel_state)
    class(process_constants_t), intent(inout) :: prc_const
    integer, dimension(:,:), allocatable :: hel_state
    hel_state = prc_const%hel_state
  end subroutine process_constants_get_hel_state

\eatline
\nwendcode{}\nwbegindocs{11}\nwdocspar
\nwenddocs{}\nwbegincode{12}\moddef{Process constants: prc const: TBP}\plusendmoddef
  procedure :: get_col_state => process_constants_get_col_state
\nwendcode{}\nwbegincode{13}\moddef{Process constants: procedures}\plusendmoddef
  subroutine process_constants_get_col_state (prc_const, col_state)
    class(process_constants_t), intent(inout) :: prc_const
    integer, dimension(:,:,:), allocatable :: col_state
    col_state = prc_const%col_state
  end subroutine process_constants_get_col_state

\eatline
\nwendcode{}\nwbegindocs{14}\nwdocspar
\nwenddocs{}\nwbegincode{15}\moddef{Process constants: prc const: TBP}\plusendmoddef
  procedure :: get_ghost_flag => process_constants_get_ghost_flag
\nwendcode{}\nwbegincode{16}\moddef{Process constants: procedures}\plusendmoddef
  function process_constants_get_ghost_flag (prc_const) result(ghost_flag)
   class(process_constants_t), intent(inout) :: prc_const
   logical, dimension(:,:), allocatable :: ghost_flag
   ghost_flag = prc_const%ghost_flag
  end function process_constants_get_ghost_flag

\eatline
\nwendcode{}\nwbegindocs{17}\nwdocspar
\nwenddocs{}\nwbegincode{18}\moddef{Process constants: prc const: TBP}\plusendmoddef
  procedure :: get_color_factors => process_constants_get_color_factors
\nwendcode{}\nwbegincode{19}\moddef{Process constants: procedures}\plusendmoddef
  subroutine process_constants_get_color_factors (prc_const, col_facts)
   class(process_constants_t), intent(inout) :: prc_const
   complex(default), intent(inout), dimension(:), allocatable :: col_facts
   col_facts = prc_const%color_factors
  end subroutine process_constants_get_color_factors

\eatline
\nwendcode{}\nwbegindocs{20}\nwdocspar
\nwenddocs{}\nwbegincode{21}\moddef{Process constants: prc const: TBP}\plusendmoddef
  procedure :: get_cf_index => process_constants_get_cf_index
\nwendcode{}\nwbegincode{22}\moddef{Process constants: procedures}\plusendmoddef
  subroutine process_constants_get_cf_index (prc_const, cf_index)
    class(process_constants_t), intent(inout) :: prc_const
    integer, intent(inout), dimension(:,:), allocatable :: cf_index
    cf_index = prc_const%cf_index
  end subroutine process_constants_get_cf_index

\eatline
\nwendcode{}\nwbegindocs{23}\nwdocspar
\nwenddocs{}\nwbegincode{24}\moddef{Process constants: prc const: TBP}\plusendmoddef
  procedure :: set_flv_state => process_constants_set_flv_state
\nwendcode{}\nwbegincode{25}\moddef{Process constants: procedures}\plusendmoddef
  subroutine process_constants_set_flv_state (prc_const, flv_state)
    class(process_constants_t), intent(inout) :: prc_const
    integer, intent(in), dimension(:,:), allocatable :: flv_state
    prc_const%flv_state = flv_state
  end subroutine process_constants_set_flv_state

\eatline
\nwendcode{}\nwbegindocs{26}\nwdocspar
\nwenddocs{}\nwbegincode{27}\moddef{Process constants: prc const: TBP}\plusendmoddef
  procedure :: set_col_state => process_constants_set_col_state
\nwendcode{}\nwbegincode{28}\moddef{Process constants: procedures}\plusendmoddef
  subroutine process_constants_set_col_state (prc_const, col_state)
    class(process_constants_t), intent(inout) :: prc_const
    integer, intent(in), dimension(:,:,:), allocatable :: col_state
    prc_const%col_state = col_state
  end subroutine process_constants_set_col_state

\eatline
\nwendcode{}\nwbegindocs{29}\nwdocspar
\nwenddocs{}\nwbegincode{30}\moddef{Process constants: prc const: TBP}\plusendmoddef
  procedure :: set_cf_index => process_constants_set_cf_index
\nwendcode{}\nwbegincode{31}\moddef{Process constants: procedures}\plusendmoddef
  subroutine process_constants_set_cf_index (prc_const, cf_index)
    class(process_constants_t), intent(inout) :: prc_const
    integer, dimension(:,:), intent(in), allocatable :: cf_index
    prc_const%cf_index = cf_index
  end subroutine process_constants_set_cf_index

\eatline
\nwendcode{}\nwbegindocs{32}\nwdocspar
\nwenddocs{}\nwbegincode{33}\moddef{Process constants: prc const: TBP}\plusendmoddef
  procedure :: set_color_factors => process_constants_set_color_factors
\nwendcode{}\nwbegincode{34}\moddef{Process constants: procedures}\plusendmoddef
  subroutine process_constants_set_color_factors (prc_const, color_factors)
    class(process_constants_t), intent(inout) :: prc_const
    complex(default), dimension(:), intent(in), allocatable :: color_factors
    prc_const%color_factors = color_factors
  end subroutine process_constants_set_color_factors

\eatline
\nwendcode{}\nwbegindocs{35}\nwdocspar
\nwenddocs{}\nwbegincode{36}\moddef{Process constants: prc const: TBP}\plusendmoddef
  procedure :: set_ghost_flag => process_constants_set_ghost_flag
\nwendcode{}\nwbegincode{37}\moddef{Process constants: procedures}\plusendmoddef
  subroutine process_constants_set_ghost_flag (prc_const, ghost_flag)
    class(process_constants_t), intent(inout) :: prc_const
    logical, intent(in), dimension(:,:), allocatable :: ghost_flag
    prc_const%ghost_flag = ghost_flag
  end subroutine process_constants_set_ghost_flag
\eatline
\nwendcode{}\nwbegindocs{38}\nwdocspar
%%%%%%%%%%%%%%%%%%%%%%%%%%%%%%%%%%%%%%%%%%%%%%%%%%%%%%%%%%%%%%%%%%%%%%%% 
\section{Process library interface}

The module {\tt{}prclib{\char95}interfaces} handles external matrix-element code.


\subsection{Overview}
The top-level data structure is the {\tt{}prclib{\char95}driver{\char95}t} data type.
The associated type-bound procedures deal with the generation of
external code, compilation and linking, and accessing the active
external library.

An object of type {\tt{}prclib{\char95}driver{\char95}t} consists of the following parts:
\begin{enumerate}
\item\ Metadata that identify name and status of the library driver,
  etc.
\item\ An array of process records ({\tt{}prclib{\char95}driver{\char95}record{\char95}t}), one
  for each external matrix element.
\item\ A record of type {\tt{}dlaccess{\char95}t} which handles the
  operating-system part of linking a dynamically loadable library.
\item\ A collection of procedure pointers which have a counterpart in
  the external library interface.  Given the unique identifier of a
  matrix element, the procedures retrieve generic matrix-element
  information such as the particle content and helicity combination
  tables.  There is also a procedure which returns pointers to the
  more specific procedures that a matrix element provides, called
  \emph{features}.
\end{enumerate}

The process records of type {\tt{}prclib{\char95}driver{\char95}record{\char95}t} handle the
individual matrix elements.  Each record identifies a process by name
({\tt{}id}), names the physics model to be loaded for this process, lists
the features that the associated matrix-element code provides, and
holds a {\tt{}writer} object which handles all operations that depend on
the process type.  The numbering of process records is identical to
the numbering of matrix-element codes in the external library.

The writer object is of abstract type {\tt{}prc{\char95}writer{\char95}t}.  The module
defines two basic, also abstract, extensions:
{\tt{}prc{\char95}writer{\char95}f{\char95}module{\char95}t} and {\tt{}prc{\char95}writer{\char95}c{\char95}lib{\char95}t}.  The first
version is for matrix-element code that is available in form of
Fortran modules.  The writer contains type-bound procedures which
create appropriate {\tt{}use} directives and {\tt{}C}-compatible wrapper
functions for the given set of Fortran modules and their features.
The second version is for matrix-element code that is available in
form of a C-compatible library (this includes Fortran libraries with
proper C bindings).  The writer needs not write wrapper function, but
explicit interface blocks for the matrix-element features.

Each matrix-element variant is encoded in an appropriate extension of
{\tt{}prc{\char95}writer{\char95}t}.  For instance, \oMega\ matrix elements provide an
implementation {\tt{}omega{\char95}writer{\char95}t} which extends
{\tt{}prc{\char95}writer{\char95}f{\char95}module{\char95}t}.

\subsection{Workflow}
We expect that the functionality provided by this module is called in
the following order:
\begin{enumerate}
\item
  The caller initializes the {\tt{}prclib{\char95}driver{\char95}t} object and fills the
  array of {\tt{}prclib{\char95}record{\char95}t} entries with the appropriate process
  data and process-specific writer objects.
\item
  It calls the {\tt{}generate{\char95}makefile} method to set up an appropriate
  makefile in the current directory.  The makefile will handle source
  generation, compilation and linking both for the individual matrix
  elements (unless this has to be done manually) and for the common
  external driver code which interfaces those matrix element.
\item
  The {\tt{}generate{\char95}driver{\char95}code} writes the common driver as source code
  to file.
\item
  The methods {\tt{}make{\char95}source}, {\tt{}make{\char95}compile}, and {\tt{}make{\char95}link}
  individually perform the corresponding steps in building the
  library.   Wherever possible, they simply use the generated makefile.
  By calling {\tt{}make}, we make sure that we can avoid
  unnecessary recompilation.  For the
  compilation and linking steps, the makefile will employ {\tt{}libtool}.
\item
  The {\tt{}load} method loads the library procedures into the
  corresponding procedure pointers, using the {\tt{}dlopen} mechanism via
  the {\tt{}dlaccess} subobject.
\end{enumerate}

\subsection{The module}
\nwenddocs{}\nwbegincode{39}\moddef{\code{}prclib{\char95}interfaces.f90\edoc{}}\endmoddef
\LA{}File header\RA{}

module prclib_interfaces

  use, intrinsic :: iso_c_binding !NODEP!
  
  use kinds
\LA{}Use strings\RA{}
  use io_units
  use unit_tests
  use system_defs, only: TAB
  use diagnostics
  use os_interface

\LA{}Standard module head\RA{}

\LA{}Prclib interfaces: public\RA{}

\LA{}Prclib interfaces: types\RA{}

\LA{}Prclib interfaces: interfaces\RA{}

\LA{}Prclib interfaces: test types\RA{}

contains
  
\LA{}Prclib interfaces: procedures\RA{}

\LA{}Prclib interfaces: tests\RA{}

end module prclib_interfaces
\eatline
\nwendcode{}\nwbegindocs{40}\nwdocspar
\subsection{Writers}
External matrix element code provides externally visible procedures,
which we denote as \emph{features}.  The features consist of
informational subroutines and functions which are mandatory (universal
features) and matrix-element specific subroutines and functions
(specific features).  The driver interfaces the
generic features directly, while it returns the specific features in
form of bind(C) procedure pointers to the caller.  For instance,
function {\tt{}n{\char95}in} is generic, while the matrix matrix-element value
itself is specific.

To implement these tasks, the driver needs {\tt{}use} directives for
Fortran module procedures, interface blocks for other external stuff,
wrapper code, and Makefile snippets.

\subsubsection{Generic writer}
In the {\tt{}prc{\char95}writer{\char95}t} data type, we collect the procedures which
implement the writing tasks.  The type is abstract.  The
concrete implementations are defined by an extension which is specific
for the process type.

The MD5 sum stored here should be the MD5 checksum of the current process
component, which can be calculated once the process is configured completely.
It can be used by implementations which work with external files, such as
\oMega.
\nwenddocs{}\nwbegincode{41}\moddef{Prclib interfaces: public}\endmoddef
  public :: prc_writer_t
\nwendcode{}\nwbegincode{42}\moddef{Prclib interfaces: types}\endmoddef
  type, abstract :: prc_writer_t
     character(32) :: md5sum = ""
   contains
   \LA{}Prclib interfaces: prc writer: TBP\RA{}
  end type prc_writer_t

\eatline
\nwendcode{}\nwbegindocs{43}In any case, it is useful to have a string representation of the
writer type.  This must be implemented by all extensions.
\nwenddocs{}\nwbegincode{44}\moddef{Prclib interfaces: prc writer: TBP}\endmoddef
  procedure(get_const_string), nopass, deferred :: type_name
\nwendcode{}\nwbegincode{45}\moddef{Prclib interfaces: interfaces}\endmoddef
  abstract interface
     function get_const_string () result (string)
       import
       type(string_t) :: string
     end function get_const_string
  end interface

\eatline
\nwendcode{}\nwbegindocs{46}Return the name of a procedure that implements a given feature, as
it is provided by the external matrix-element code.  For a reasonable
default, we take the feature name unchanged.
\nwenddocs{}\nwbegincode{47}\moddef{Prclib interfaces: prc writer: TBP}\plusendmoddef
  procedure, nopass :: get_procname => prc_writer_get_procname
\nwendcode{}\nwbegincode{48}\moddef{Prclib interfaces: procedures}\endmoddef
  function prc_writer_get_procname (feature) result (name)
    type(string_t) :: name
    type(string_t), intent(in) :: feature
    name = feature
  end function prc_writer_get_procname
  
\eatline
\nwendcode{}\nwbegindocs{49}Return the name of a procedure that implements a given feature with
the bind(C) property, so it can be accessed via a C procedure pointer and
handled by dlopen.  We need this for all special features of a matrix
element, since the interface has to return a C function pointer for it.
For a default implementation, we prefix the external procedure name by
the process ID.
\nwenddocs{}\nwbegincode{50}\moddef{Prclib interfaces: prc writer: TBP}\plusendmoddef
  procedure :: get_c_procname => prc_writer_get_c_procname
\nwendcode{}\nwbegincode{51}\moddef{Prclib interfaces: procedures}\plusendmoddef
  function prc_writer_get_c_procname (writer, id, feature) result (name)
    class(prc_writer_t), intent(in) :: writer
    type(string_t), intent(in) :: id, feature
    type(string_t) :: name
    name = id // "_" // feature
  end function prc_writer_get_c_procname
  
\eatline
\nwendcode{}\nwbegindocs{52}Common signature of code-writing procedures.  The procedure may
use the process ID, and the feature name.
(Not necessarily all of them.)
\nwenddocs{}\nwbegincode{53}\moddef{Prclib interfaces: interfaces}\plusendmoddef
  abstract interface
     subroutine write_code_file (writer, id)
       import
       class(prc_writer_t), intent(in) :: writer
       type(string_t), intent(in) :: id
     end subroutine write_code_file
  end interface
  
  abstract interface
     subroutine write_code (writer, unit, id)
       import
       class(prc_writer_t), intent(in) :: writer
       integer, intent(in) :: unit
       type(string_t), intent(in) :: id
     end subroutine write_code
  end interface
  
  abstract interface
     subroutine write_code_os (writer, unit, id, os_data, testflag)
       import
       class(prc_writer_t), intent(in) :: writer
       integer, intent(in) :: unit
       type(string_t), intent(in) :: id
       type(os_data_t), intent(in) :: os_data
       logical, intent(in), optional :: testflag
     end subroutine write_code_os
  end interface
  
  abstract interface
     subroutine write_feature_code (writer, unit, id, feature)
       import
       class(prc_writer_t), intent(in) :: writer
       integer, intent(in) :: unit
       type(string_t), intent(in) :: id, feature
     end subroutine write_feature_code
  end interface
  
\eatline
\nwendcode{}\nwbegindocs{54}There must be a procedure which writes an interface block for a
given feature.  If the external matrix element is implemented as a
Fortran module, this is required only for the specific features which
are returned as procedure pointers.
\nwenddocs{}\nwbegincode{55}\moddef{Prclib interfaces: prc writer: TBP}\plusendmoddef
  procedure(write_feature_code), deferred :: write_interface
\eatline
\nwendcode{}\nwbegindocs{56}There must also be a procedure which writes Makefile code which is
specific for the current process, but not the feature.
\nwenddocs{}\nwbegincode{57}\moddef{Prclib interfaces: prc writer: TBP}\plusendmoddef
  procedure(write_code_os), deferred :: write_makefile_code
\eatline
\nwendcode{}\nwbegindocs{58}This procedure writes code process-specific source-code file
(which need not be Fortran).  It may be a no-op, if the source code is
generated by Make instead.
\nwenddocs{}\nwbegincode{59}\moddef{Prclib interfaces: prc writer: TBP}\plusendmoddef
  procedure(write_code_file), deferred :: write_source_code
\eatline
\nwendcode{}\nwbegindocs{60}\nwdocspar
\subsubsection{Writer for Fortran-module matrix elements}
If the matrix element is available as a Fortran module, we have
specific requirements: (i) the features are imported via {\tt{}use}
directives, (ii) the specific features require bind(C) wrappers.

The type is still abstract, all methods must be implemented explicitly
for a specific matrix-element variant.
\nwenddocs{}\nwbegincode{61}\moddef{Prclib interfaces: public}\plusendmoddef
  public :: prc_writer_f_module_t
\nwendcode{}\nwbegincode{62}\moddef{Prclib interfaces: types}\plusendmoddef
  type, extends (prc_writer_t), abstract :: prc_writer_f_module_t
   contains
   \LA{}Prclib interfaces: prc writer f module: TBP\RA{}
  end type prc_writer_f_module_t
  
\eatline
\nwendcode{}\nwbegindocs{63}Return the name of the Fortran module.  As a default
implementation, we take the process ID unchanged.
\nwenddocs{}\nwbegincode{64}\moddef{Prclib interfaces: prc writer f module: TBP}\endmoddef
  procedure, nopass :: get_module_name => prc_writer_get_module_name
\nwendcode{}\nwbegincode{65}\moddef{Prclib interfaces: procedures}\plusendmoddef
  function prc_writer_get_module_name (id) result (name)
    type(string_t) :: name
    type(string_t), intent(in) :: id
    name = id
  end function prc_writer_get_module_name
  
\eatline
\nwendcode{}\nwbegindocs{66}Write a {\tt{}use} directive that associates the driver reference with
the procedure in the matrix element code.  By default, we use the C
name for this.
\nwenddocs{}\nwbegincode{67}\moddef{Prclib interfaces: prc writer f module: TBP}\plusendmoddef
  procedure :: write_use_line => prc_writer_write_use_line
\nwendcode{}\nwbegincode{68}\moddef{Prclib interfaces: procedures}\plusendmoddef
  subroutine prc_writer_write_use_line (writer, unit, id, feature)
    class(prc_writer_f_module_t), intent(in) :: writer
    integer, intent(in) :: unit
    type(string_t) :: id, feature
    write (unit, "(2x,9A)")  "use ", char (writer%get_module_name (id)), &
         ", only: ", char (writer%get_c_procname (id, feature)), &
         " => ", char (writer%get_procname (feature))
  end subroutine prc_writer_write_use_line

\eatline
\nwendcode{}\nwbegindocs{69}Write a wrapper routine for a feature.  This also associates a C
name the module procedure.  The details depend on the writer variant.
\nwenddocs{}\nwbegincode{70}\moddef{Prclib interfaces: prc writer f module: TBP}\plusendmoddef
  procedure(prc_write_wrapper), deferred :: write_wrapper
\nwendcode{}\nwbegincode{71}\moddef{Prclib interfaces: interfaces}\plusendmoddef
  abstract interface
     subroutine prc_write_wrapper (writer, unit, id, feature)
       import
       class(prc_writer_f_module_t), intent(in) :: writer
       integer, intent(in) :: unit
       type(string_t), intent(in) :: id, feature
     end subroutine prc_write_wrapper
  end interface

\eatline
\nwendcode{}\nwbegindocs{72}This is used for testing only: initialize the writer with a specific MD5 sum
string.
\nwenddocs{}\nwbegincode{73}\moddef{Prclib interfaces: prc writer: TBP}\plusendmoddef
  procedure :: init_test => prc_writer_init_test
\nwendcode{}\nwbegincode{74}\moddef{Prclib interfaces: procedures}\plusendmoddef
  subroutine prc_writer_init_test (writer)
    class(prc_writer_t), intent(out) :: writer
    writer%md5sum = "1234567890abcdef1234567890abcdef"
  end subroutine prc_writer_init_test
  
\eatline
\nwendcode{}\nwbegindocs{75}\nwdocspar
\subsubsection{Writer for C-library matrix elements}
This applies if the matrix element is available as a C library or a Fortran
library with bind(C) compatible interface.  We can use the basic
version.

The type is still abstract, all methods must be implemented explicitly
for a specific matrix-element variant.
\nwenddocs{}\nwbegincode{76}\moddef{Prclib interfaces: types}\plusendmoddef
  type, extends (prc_writer_t), abstract :: prc_writer_c_lib_t
   contains
   \LA{}Prclib interfaces: prc writer c lib: TBP\RA{}
  end type prc_writer_c_lib_t
  
\eatline
\nwendcode{}\nwbegindocs{77}\nwdocspar
\subsection{Process records in the library driver}
A process record holds the process (component) {\tt{}ID}, the physics
{\tt{}model{\char95}name}, and the array of {\tt{}feature}s that are
implemented by the corresponding matrix element code.

The {\tt{}writer} component holds procedures.  The procedures write
source code for the current record, either for the driver or for the
Makefile.
\nwenddocs{}\nwbegincode{78}\moddef{Prclib interfaces: types}\plusendmoddef
  type :: prclib_driver_record_t
     type(string_t) :: id
     type(string_t) :: model_name
     type(string_t), dimension(:), allocatable :: feature
     class(prc_writer_t), pointer :: writer => null ()
   contains
   \LA{}Prclib interfaces: prclib driver record: TBP\RA{}
  end type prclib_driver_record_t

\eatline
\nwendcode{}\nwbegindocs{79}Output routine.  We indent the output, so it smoothly integrates
into the output routine for the whole driver.

Note: the pointer {\tt{}writer} is introduced as a workaround for a NAG compiler
bug.
\nwenddocs{}\nwbegincode{80}\moddef{Prclib interfaces: prclib driver record: TBP}\endmoddef
  procedure :: write => prclib_driver_record_write
\nwendcode{}\nwbegincode{81}\moddef{Prclib interfaces: procedures}\plusendmoddef
  subroutine prclib_driver_record_write (object, unit)
    class(prclib_driver_record_t), intent(in) :: object
    integer, intent(in) :: unit
    integer :: j
    class(prc_writer_t), pointer :: writer
    write (unit, "(3x,A,2x,'[',A,']')")  &
         char (object%id), char (object%model_name)
    if (allocated (object%feature)) then 
       writer => object%writer
       write (unit, "(5x,A,A)", advance="no") &
            char (writer%type_name ()), ":"
       do j = 1, size (object%feature)
          write (unit, "(1x,A)", advance="no") &
               char (object%feature(j))
       end do
       write (unit, *)
    end if
  end subroutine prclib_driver_record_write

\eatline
\nwendcode{}\nwbegindocs{82}Get the C procedure name for a feature.
\nwenddocs{}\nwbegincode{83}\moddef{Prclib interfaces: prclib driver record: TBP}\plusendmoddef
  procedure :: get_c_procname => prclib_driver_record_get_c_procname
\nwendcode{}\nwbegincode{84}\moddef{Prclib interfaces: procedures}\plusendmoddef
  function prclib_driver_record_get_c_procname (record, feature) result (name)
    type(string_t) :: name
    class(prclib_driver_record_t), intent(in) :: record
    type(string_t), intent(in) :: feature
    name = record%writer%get_c_procname (record%id, feature)
  end function prclib_driver_record_get_c_procname
  
\eatline
\nwendcode{}\nwbegindocs{85}Write a USE directive for a given feature.  Applies only if the
record corresponds to a Fortran module.
\nwenddocs{}\nwbegincode{86}\moddef{Prclib interfaces: prclib driver record: TBP}\plusendmoddef
  procedure :: write_use_line => prclib_driver_record_write_use_line
\nwendcode{}\nwbegincode{87}\moddef{Prclib interfaces: procedures}\plusendmoddef
  subroutine prclib_driver_record_write_use_line (record, unit, feature)
    class(prclib_driver_record_t), intent(in) :: record
    integer, intent(in) :: unit
    type(string_t), intent(in) :: feature
    select type (writer => record%writer)
    class is (prc_writer_f_module_t)
       call writer%write_use_line (unit, record%id, feature)
    end select
  end subroutine prclib_driver_record_write_use_line
  
\eatline
\nwendcode{}\nwbegindocs{88}The alternative: write an interface block for a given feature,
unless the record corresponds to a Fortran module.
\nwenddocs{}\nwbegincode{89}\moddef{Prclib interfaces: prclib driver record: TBP}\plusendmoddef
  procedure :: write_interface => prclib_driver_record_write_interface
\nwendcode{}\nwbegincode{90}\moddef{Prclib interfaces: procedures}\plusendmoddef
  subroutine prclib_driver_record_write_interface (record, unit, feature)
    class(prclib_driver_record_t), intent(in) :: record
    integer, intent(in) :: unit
    type(string_t), intent(in) :: feature
    select type (writer => record%writer)
    class is (prc_writer_f_module_t)
    class default
       call writer%write_interface (unit, record%id, feature)
    end select
  end subroutine prclib_driver_record_write_interface
  
\eatline
\nwendcode{}\nwbegindocs{91}Write all special feature interfaces for the current record.  Do
this for all process variants.
\nwenddocs{}\nwbegincode{92}\moddef{Prclib interfaces: prclib driver record: TBP}\plusendmoddef
  procedure :: write_interfaces => prclib_driver_record_write_interfaces
\nwendcode{}\nwbegincode{93}\moddef{Prclib interfaces: procedures}\plusendmoddef
  subroutine prclib_driver_record_write_interfaces (record, unit)
    class(prclib_driver_record_t), intent(in) :: record
    integer, intent(in) :: unit
    integer :: i
    do i = 1, size (record%feature)
       call record%writer%write_interface (unit, record%id, record%feature(i))
    end do
  end subroutine prclib_driver_record_write_interfaces
  
\eatline
\nwendcode{}\nwbegindocs{94}Write the wrapper routines for this record, if it corresponds to a
Fortran module.
\nwenddocs{}\nwbegincode{95}\moddef{Prclib interfaces: prclib driver record: TBP}\plusendmoddef
  procedure :: write_wrappers => prclib_driver_record_write_wrappers
\nwendcode{}\nwbegincode{96}\moddef{Prclib interfaces: procedures}\plusendmoddef
  subroutine prclib_driver_record_write_wrappers (record, unit)
    class(prclib_driver_record_t), intent(in) :: record
    integer, intent(in) :: unit
    integer :: i
    select type (writer => record%writer)
    class is (prc_writer_f_module_t)
       do i = 1, size (record%feature)
          call writer%write_wrapper (unit, record%id, record%feature(i))
       end do
    end select
  end subroutine prclib_driver_record_write_wrappers

\eatline
\nwendcode{}\nwbegindocs{97}Write the Makefile code for this record.
\nwenddocs{}\nwbegincode{98}\moddef{Prclib interfaces: prclib driver record: TBP}\plusendmoddef
  procedure :: write_makefile_code => prclib_driver_record_write_makefile_code
\nwendcode{}\nwbegincode{99}\moddef{Prclib interfaces: procedures}\plusendmoddef
  subroutine prclib_driver_record_write_makefile_code &
       (record, unit, os_data, testflag)
    class(prclib_driver_record_t), intent(in) :: record
    integer, intent(in) :: unit
    type(os_data_t), intent(in) :: os_data
    logical, intent(in), optional :: testflag
    call record%writer%write_makefile_code (unit, record%id, os_data, testflag)
  end subroutine prclib_driver_record_write_makefile_code

\eatline
\nwendcode{}\nwbegindocs{100}Write source-code files for this record.  (Does nothing if the
source code is handled by Makefile rules.)
\nwenddocs{}\nwbegincode{101}\moddef{Prclib interfaces: prclib driver record: TBP}\plusendmoddef
  procedure :: write_source_code => prclib_driver_record_write_source_code
\nwendcode{}\nwbegincode{102}\moddef{Prclib interfaces: procedures}\plusendmoddef
  subroutine prclib_driver_record_write_source_code (record)
    class(prclib_driver_record_t), intent(in) :: record
    call record%writer%write_source_code (record%id)
  end subroutine prclib_driver_record_write_source_code

\eatline
\nwendcode{}\nwbegindocs{103}\nwdocspar
\subsection{The process library driver object}
A {\tt{}prclib{\char95}driver{\char95}t} object provides the interface to external matrix element
code.  The code is provided by an external library which is either
statically or dynamically linked.

The dynamic and static versions of the library are two different
implementations of the abstract base type.

The {\tt{}basename} identifies the library, both by file names and by Fortran
variable names.

The {\tt{}loaded} flag becomes true once all procedure pointers to the
matrix element have been assigned.

For a dynamical external library, the communication proceeds via a
{\tt{}dlaccess} object.

{\tt{}n{\char95}processes} is the number of external process code components that
are referenced by this library.  The code is addressed by index ({\tt{}i{\char95}lib}
in the process library entry above).  This number should be equal to
the number returned by {\tt{}get{\char95}n{\char95}prc}.

For each external process, there is a separate {\tt{}record} which holds
the data that are needed for the driver parts which are specific
for a given process component.  The actual pointers for the loaded
library will be assigned elsewhere.

The remainder is a collection of procedure pointers, which can be
assigned once all external code has been compiled and linked.
The procedure pointers all take a process component code
index as an argument.  Most return information about the process
component that should match the process definition.  The {\tt{}get{\char95}fptr}
procedures return a function pointer, which is the actual means to
compute matrix elements or retrieve associated data.

Finally, the {\tt{}unload{\char95}hook} and {\tt{}reload{\char95}hook} pointers allow for
the insertion of additional code when a library is loaded.
\nwenddocs{}\nwbegincode{104}\moddef{Prclib interfaces: public}\plusendmoddef
  public :: prclib_driver_t
\nwendcode{}\nwbegincode{105}\moddef{Prclib interfaces: types}\plusendmoddef
  type, abstract :: prclib_driver_t
     type(string_t) :: basename
     character(32) :: md5sum = ""
     logical :: loaded = .false.
     type(string_t) :: libname
     type(string_t) :: modellibs_ldflags
     integer :: n_processes = 0
     type(prclib_driver_record_t), dimension(:), allocatable :: record
     procedure(prc_get_n_processes), nopass, pointer :: &
          get_n_processes => null ()
     procedure(prc_get_stringptr), nopass, pointer :: &
          get_process_id_ptr => null ()
     procedure(prc_get_stringptr), nopass, pointer :: &
          get_model_name_ptr => null ()
     procedure(prc_get_stringptr), nopass, pointer :: &
          get_md5sum_ptr => null ()
     procedure(prc_get_log), nopass, pointer :: &
          get_openmp_status => null ()
     procedure(prc_get_int), nopass, pointer :: get_n_in  => null ()
     procedure(prc_get_int), nopass, pointer :: get_n_out => null ()
     procedure(prc_get_int), nopass, pointer :: get_n_flv => null ()
     procedure(prc_get_int), nopass, pointer :: get_n_hel => null ()
     procedure(prc_get_int), nopass, pointer :: get_n_col => null ()
     procedure(prc_get_int), nopass, pointer :: get_n_cin => null ()
     procedure(prc_get_int), nopass, pointer :: get_n_cf  => null ()
     procedure(prc_set_int_tab1), nopass, pointer :: &
          set_flv_state_ptr => null ()
     procedure(prc_set_int_tab1), nopass, pointer :: &
          set_hel_state_ptr => null ()
     procedure(prc_set_col_state), nopass, pointer :: &
          set_col_state_ptr => null ()
     procedure(prc_set_color_factors), nopass, pointer :: &
          set_color_factors_ptr => null ()
     procedure(prc_get_fptr), nopass, pointer :: get_fptr => null ()
   contains
   \LA{}Prclib interfaces: prclib driver: TBP\RA{}
  end type prclib_driver_t
  
\eatline
\nwendcode{}\nwbegindocs{106}This is the dynamic version.  It contains a {\tt{}dlaccess} object for
communicating with the OS.
\nwenddocs{}\nwbegincode{107}\moddef{Prclib interfaces: types}\plusendmoddef
  type, extends (prclib_driver_t) :: prclib_driver_dynamic_t
     type(dlaccess_t) :: dlaccess
   contains
   \LA{}Prclib interfaces: prclib driver dynamic: TBP\RA{}
  end type prclib_driver_dynamic_t
  
\eatline
\nwendcode{}\nwbegindocs{108}Print just the metadata.  Procedure pointers cannot be printed.
\nwenddocs{}\nwbegincode{109}\moddef{Prclib interfaces: prclib driver: TBP}\endmoddef
  procedure :: write => prclib_driver_write
\nwendcode{}\nwbegincode{110}\moddef{Prclib interfaces: procedures}\plusendmoddef
  subroutine prclib_driver_write (object, unit, libpath)
    class(prclib_driver_t), intent(in) :: object
    integer, intent(in) :: unit
    logical, intent(in), optional :: libpath
    logical :: write_lib
    integer :: i    
    write_lib = .true.    
    if (present (libpath))  write_lib = libpath
    write (unit, "(1x,A,A)")  &
         "External matrix-element code library: ", char (object%basename)
    select type (object)
    type is (prclib_driver_dynamic_t)
       write (unit, "(3x,A,L1)")  "static    = F"
    class default
       write (unit, "(3x,A,L1)")  "static    = T"
    end select
    write (unit, "(3x,A,L1)")  "loaded    = ", object%loaded
    write (unit, "(3x,A,A,A)") "MD5 sum   = '", object%md5sum, "'"
    if (write_lib) then
       write (unit, "(3x,A,A,A)") "Mdl flags = '", &
            char (object%modellibs_ldflags), "'"
    end if
    select type (object)
    type is (prclib_driver_dynamic_t)
       write (unit, *)
       call object%dlaccess%write (unit)
    end select
    write (unit, *)
    if (allocated (object%record)) then
       write (unit, "(1x,A)")  "Matrix-element code entries:"
       do i = 1, object%n_processes
          call object%record(i)%write (unit)
       end do
    else
       write (unit, "(1x,A)")  "Matrix-element code entries: [undefined]"
    end if
  end subroutine prclib_driver_write
  
\eatline
\nwendcode{}\nwbegindocs{111}Allocate a library as either static or dynamic.  For static
libraries, the procedure defers control to an external procedure which
knows about the available static libraries.  By default, this
procedure is empty, but when we build a stand-alone executable, we
replace the dummy by an actual dispatcher for the available
static libraries.  If the static dispatcher was not successful, we
allocate a dynamic library.

The default version of {\tt{}dispatch{\char95}prclib{\char95}static} resides in the
{\tt{}prebuilt} section of the \whizard\ tree, in a separate
library.  It does nothing, but can be replaced by a different
procedure that allocates a static library driver if requested by name.

Note: {\tt{}intent(out)} for the {\tt{}driver} argument segfaults with
gfortran 4.7.
\nwenddocs{}\nwbegincode{112}\moddef{Prclib interfaces: public}\plusendmoddef
  public :: dispatch_prclib_driver
\nwendcode{}\nwbegincode{113}\moddef{Prclib interfaces: procedures}\plusendmoddef
  subroutine dispatch_prclib_driver &
       (driver, basename, modellibs_ldflags)
    class(prclib_driver_t), intent(inout), allocatable :: driver
    type(string_t), intent(in) :: basename
    type(string_t), intent(in), optional :: modellibs_ldflags
    procedure(dispatch_prclib_driver) :: dispatch_prclib_static
    if (allocated (driver))  deallocate (driver)
    call dispatch_prclib_static (driver, basename)
    if (.not. allocated (driver)) then
       allocate (prclib_driver_dynamic_t :: driver)
    end if
    driver%basename = basename
    driver%modellibs_ldflags = modellibs_ldflags
  end subroutine dispatch_prclib_driver
  
\eatline
\nwendcode{}\nwbegindocs{114}Initialize the ID array and set {\tt{}n{\char95}processes} accordingly.
\nwenddocs{}\nwbegincode{115}\moddef{Prclib interfaces: prclib driver: TBP}\plusendmoddef
  procedure :: init => prclib_driver_init
\nwendcode{}\nwbegincode{116}\moddef{Prclib interfaces: procedures}\plusendmoddef
  subroutine prclib_driver_init (driver, n_processes)
    class(prclib_driver_t), intent(inout) :: driver
    integer, intent(in) :: n_processes
    driver%n_processes = n_processes
    allocate (driver%record (n_processes))
  end subroutine prclib_driver_init
  
\eatline
\nwendcode{}\nwbegindocs{117}Set the MD5 sum.  This is separate because the MD5 sum may be known only
after initialization.
\nwenddocs{}\nwbegincode{118}\moddef{Prclib interfaces: prclib driver: TBP}\plusendmoddef
  procedure :: set_md5sum => prclib_driver_set_md5sum
\nwendcode{}\nwbegincode{119}\moddef{Prclib interfaces: procedures}\plusendmoddef
  subroutine prclib_driver_set_md5sum (driver, md5sum)
    class(prclib_driver_t), intent(inout) :: driver
    character(32), intent(in) :: md5sum
    driver%md5sum = md5sum
  end subroutine prclib_driver_set_md5sum
  
\eatline
\nwendcode{}\nwbegindocs{120}Set the process record for a specific library entry.  If the index
is zero, we do nothing.

\nwenddocs{}\nwbegincode{121}\moddef{Prclib interfaces: prclib driver: TBP}\plusendmoddef
  procedure :: set_record => prclib_driver_set_record
\nwendcode{}\nwbegincode{122}\moddef{Prclib interfaces: procedures}\plusendmoddef
  subroutine prclib_driver_set_record (driver, i, &
       id, model_name, features, writer)
    class(prclib_driver_t), intent(inout) :: driver
    integer, intent(in) :: i
    type(string_t), intent(in) :: id
    type(string_t), intent(in) :: model_name
    type(string_t), dimension(:), intent(in) :: features
    class(prc_writer_t), intent(in), pointer :: writer
    if (i > 0) then
       associate (record => driver%record(i))
         record%id = id
         record%model_name = model_name
         allocate (record%feature (size (features)))
         record%feature = features
         record%writer => writer
       end associate
    end if
  end subroutine prclib_driver_set_record
  
\eatline
\nwendcode{}\nwbegindocs{123}Write all USE directives for a given feature, scanning the array of
processes.  Only Fortran-module processes count.  Then, write
interface blocks for the remaining processes.

The {\tt{}implicit\ none} statement must go in-between.
\nwenddocs{}\nwbegincode{124}\moddef{Prclib interfaces: prclib driver: TBP}\plusendmoddef
  procedure :: write_interfaces => prclib_driver_write_interfaces
\nwendcode{}\nwbegincode{125}\moddef{Prclib interfaces: procedures}\plusendmoddef
  subroutine prclib_driver_write_interfaces (driver, unit, feature)
    class(prclib_driver_t), intent(in) :: driver
    integer, intent(in) :: unit
    type(string_t), intent(in) :: feature
    integer :: i
    do i = 1, driver%n_processes
       call driver%record(i)%write_use_line (unit, feature)
    end do
    write (unit, "(2x,9A)")  "implicit none"
    do i = 1, driver%n_processes
       call driver%record(i)%write_interface (unit, feature)
    end do
  end subroutine prclib_driver_write_interfaces

\eatline
\nwendcode{}\nwbegindocs{126}\nwdocspar
\subsection{Write makefile}
The makefile contains constant parts, parts that depend on the library
name, and parts that depend on the specific processes and their types.
\nwenddocs{}\nwbegincode{127}\moddef{Prclib interfaces: prclib driver: TBP}\plusendmoddef
  procedure :: generate_makefile => prclib_driver_generate_makefile
\nwendcode{}\nwbegincode{128}\moddef{Prclib interfaces: procedures}\plusendmoddef
  subroutine prclib_driver_generate_makefile (driver, unit, os_data, testflag)
    class(prclib_driver_t), intent(in) :: driver
    integer, intent(in) :: unit
    type(os_data_t), intent(in) :: os_data
    logical, intent(in), optional :: testflag
    integer :: i
    write (unit, "(A)")  "# WHIZARD: Makefile for process library '" &
         // char (driver%basename) // "'"
    write (unit, "(A)")  "# Automatically generated file, do not edit"
    write (unit, "(A)")  ""
    write (unit, "(A)")  "# Integrity check (don't modify the following line!)"
    write (unit, "(A)")  "MD5SUM = '" // driver%md5sum // "'"
    write (unit, "(A)")  ""
    write (unit, "(A)")  "# Library name"
    write (unit, "(A)")  "BASE = " // char (driver%basename)
    write (unit, "(A)")  ""
    write (unit, "(A)")  "# Compiler"
    write (unit, "(A)")  "FC = " // char (os_data%fc)
    write (unit, "(A)")  "CC = " // char (os_data%cc)
    write (unit, "(A)")  ""    
    write (unit, "(A)")  "# Included libraries"
    write (unit, "(A)")  "FCINCL = " // char (os_data%whizard_includes)
    write (unit, "(A)")  ""
    write (unit, "(A)")  "# Compiler flags"
    write (unit, "(A)")  "FCFLAGS = " // char (os_data%fcflags)
    write (unit, "(A)")  "FCFLAGS_PIC = " // char (os_data%fcflags_pic)
    write (unit, "(A)")  "CFLAGS = " // char (os_data%cflags)
    write (unit, "(A)")  "CFLAGS_PIC = " // char (os_data%cflags_pic)
    write (unit, "(A)")  "LDFLAGS = " // char (os_data%whizard_ldflags) &
         // " " // char (os_data%ldflags) // " " // &
         char (driver%modellibs_ldflags)
    write (unit, "(A)")  ""
    write (unit, "(A)")  "# LaTeX setup"
    write (unit, "(A)")  "LATEX = " // char (os_data%latex)
    write (unit, "(A)")  "MPOST = " // char (os_data%mpost)
    write (unit, "(A)")  "DVIPS = " // char (os_data%dvips)
    write (unit, "(A)")  "PS2PDF = " // char (os_data%ps2pdf)    
    write (unit, "(A)")  'TEX_FLAGS = "$$TEXINPUTS:' // &
         char(os_data%whizard_texpath) // '"'
    write (unit, "(A)")  'MP_FLAGS  = "$$MPINPUTS:' // &
         char(os_data%whizard_texpath) // '"'
    write (unit, "(A)")  ""    
    write (unit, "(A)")  "# Libtool"
    write (unit, "(A)")  "LIBTOOL = " // char (os_data%whizard_libtool)
    write (unit, "(A)")  "FCOMPILE = $(LIBTOOL) --tag=FC --mode=compile"
    write (unit, "(A)")  "CCOMPILE = $(LIBTOOL) --tag=CC --mode=compile"
    write (unit, "(A)")  "LINK = $(LIBTOOL) --tag=FC --mode=link"
    write (unit, "(A)")  ""
    write (unit, "(A)")  "# Compile commands (default)"
    write (unit, "(A)")  "LTFCOMPILE = $(FCOMPILE) $(FC) -c &
         &$(FCINCL) $(FCFLAGS) $(FCFLAGS_PIC)"
    write (unit, "(A)")  "LTCCOMPILE = $(CCOMPILE) $(CC) -c &
         &$(CFLAGS) $(CFLAGS_PIC)"
    write (unit, "(A)")  ""
    write (unit, "(A)")  "# Default target"
    write (unit, "(A)")  "all: link diags"
    write (unit, "(A)")  ""
    write (unit, "(A)")  "# Matrix-element code files"
    do i = 1, size (driver%record)
       call driver%record(i)%write_makefile_code (unit, os_data, testflag)
    end do
    write (unit, "(A)")  ""
    write (unit, "(A)")  "# Library driver"
    write (unit, "(A)")  "$(BASE).lo: $(BASE).f90 $(OBJECTS)"
    write (unit, "(A)")  TAB // "$(LTFCOMPILE) $<"
    write (unit, "(A)")  ""
    write (unit, "(A)")  "# Library"
    write (unit, "(A)")  "$(BASE).la: $(BASE).lo $(OBJECTS)"
    write (unit, "(A)")  TAB // "$(LINK) $(FC) -module -rpath /dev/null &
         &$(FCFLAGS) $(LDFLAGS) -o $(BASE).la $^"
    write (unit, "(A)")  ""
    write (unit, "(A)")  "# Main targets"
    write (unit, "(A)")  "link: compile $(BASE).la"
    write (unit, "(A)")  "compile: source $(OBJECTS) $(TEX_OBJECTS) $(BASE).lo"
    write (unit, "(A)")  "compile_tex: $(TEX_OBJECTS)"    
    write (unit, "(A)")  "source: $(SOURCES) $(BASE).f90 $(TEX_SOURCES)"
    write (unit, "(A)")  ".PHONY: link diags compile compile_tex source"
    write (unit, "(A)")  ""
    write (unit, "(A)")  "# Specific cleanup targets"
    do i = 1, size (driver%record)
       write (unit, "(A)")  "clean-" // char (driver%record(i)%id) // ":"
       write (unit, "(A)")  ".PHONY: clean-" // char (driver%record(i)%id)
    end do
    write (unit, "(A)")  ""
    write (unit, "(A)")  "# Generic cleanup targets"
    write (unit, "(A)")  "clean-library:"
    write (unit, "(A)")  TAB // "rm -f $(BASE).la"
    write (unit, "(A)")  "clean-objects:"
    write (unit, "(A)")  TAB // "rm -f $(BASE).lo $(BASE)_driver.mod &
         &$(CLEAN_OBJECTS)"
    write (unit, "(A)")  "clean-source:"
    write (unit, "(A)")  TAB // "rm -f $(CLEAN_SOURCES)"
    write (unit, "(A)")  "clean-driver:"
    write (unit, "(A)")  TAB // "rm -f $(BASE).f90"
    write (unit, "(A)")  "clean-makefile:"
    write (unit, "(A)")  TAB // "rm -f $(BASE).makefile"
    write (unit, "(A)")  ".PHONY: clean-library clean-objects &
         &clean-source clean-driver clean-makefile"
    write (unit, "(A)")  ""
    write (unit, "(A)")  "clean: clean-library clean-objects clean-source"
    write (unit, "(A)")  "distclean: clean clean-driver clean-makefile"
    write (unit, "(A)")  ".PHONY: clean distclean"
  end subroutine prclib_driver_generate_makefile

\eatline
\nwendcode{}\nwbegindocs{129}\nwdocspar
\subsection{Write driver file}
This procedure writes the process library driver source code to the
specified output unit.  The individual routines for writing
source-code procedures are given below.
\nwenddocs{}\nwbegincode{130}\moddef{Prclib interfaces: prclib driver: TBP}\plusendmoddef
  procedure :: generate_driver_code => prclib_driver_generate_code
\nwendcode{}\nwbegincode{131}\moddef{Prclib interfaces: procedures}\plusendmoddef
  subroutine prclib_driver_generate_code (driver, unit)
    class(prclib_driver_t), intent(in) :: driver
    integer, intent(in) :: unit
    type(string_t) :: prefix
    integer :: i

    prefix = driver%basename // "_"

    write (unit, "(A)")  "! WHIZARD matrix-element code interface"
    write (unit, "(A)")  "!"
    write (unit, "(A)")  "! Automatically generated file, do not edit"
    call driver%write_module (unit, prefix)
    call driver%write_lib_md5sum_fun (unit, prefix)
    call driver%write_get_n_processes_fun (unit, prefix)
    call driver%write_get_process_id_fun (unit, prefix)
    call driver%write_get_model_name_fun (unit, prefix)
    call driver%write_get_md5sum_fun (unit, prefix)
    call driver%write_string_to_array_fun (unit, prefix)
    call driver%write_get_openmp_status_fun (unit, prefix)
    call driver%write_get_int_fun (unit, prefix, var_str ("n_in"))
    call driver%write_get_int_fun (unit, prefix, var_str ("n_out"))
    call driver%write_get_int_fun (unit, prefix, var_str ("n_flv"))
    call driver%write_get_int_fun (unit, prefix, var_str ("n_hel"))
    call driver%write_get_int_fun (unit, prefix, var_str ("n_col"))
    call driver%write_get_int_fun (unit, prefix, var_str ("n_cin"))
    call driver%write_get_int_fun (unit, prefix, var_str ("n_cf"))
    call driver%write_set_int_sub (unit, prefix, var_str ("flv_state"))
    call driver%write_set_int_sub (unit, prefix, var_str ("hel_state"))
    call driver%write_set_col_state_sub (unit, prefix)
    call driver%write_set_color_factors_sub (unit, prefix)
    call driver%write_get_fptr_sub (unit, prefix)
    do i = 1, driver%n_processes
       call driver%record(i)%write_wrappers (unit)
    end do
  end subroutine prclib_driver_generate_code
  
\eatline
\nwendcode{}\nwbegindocs{132}The driver module is used and required \emph{only} if we intend to
link the library statically.  Then, it provides the (static) driver
type as a concrete implementation of the abstract library driver.
This type contains the internal dispatcher for assigning the library
procedures to their appropriate procedure pointers.  In the dynamical
case, the assignment is done via the base-type dispatcher which invokes
the DL mechanism.  

However, compiling this together with the rest in any case should not
do any harm.
\nwenddocs{}\nwbegincode{133}\moddef{Prclib interfaces: prclib driver: TBP}\plusendmoddef
  procedure, nopass :: write_module => prclib_driver_write_module
\nwendcode{}\nwbegincode{134}\moddef{Prclib interfaces: procedures}\plusendmoddef
  subroutine prclib_driver_write_module (unit, prefix)
    integer, intent(in) :: unit
    type(string_t), intent(in) :: prefix
    write (unit, "(A)")  ""
    write (unit, "(A)")  "! Module: define library driver as an extension &
         &of the abstract driver type."
    write (unit, "(A)")  "! This is used _only_ by the library dispatcher &
         &of a static executable."
    write (unit, "(A)")  "! For a dynamical library, the stand-alone proce&
         &dures are linked via libdl."
    write (unit, "(A)")  ""
    write (unit, "(A)")  "module " &
         // char (prefix) // "driver"
    write (unit, "(A)")  "  use iso_c_binding"
    write (unit, "(A)")  "  use iso_varying_string, string_t => varying_string"
    write (unit, "(A)")  "  use diagnostics"
    write (unit, "(A)")  "  use prclib_interfaces"
    write (unit, "(A)")  ""
    write (unit, "(A)")  "  implicit none"
    write (unit, "(A)")  ""
    write (unit, "(A)")  "  type, extends (prclib_driver_t) :: " &
         // char (prefix) // "driver_t"
    write (unit, "(A)")  "   contains"
    write (unit, "(A)")  "     procedure :: get_c_funptr => " &
         // char (prefix) // "driver_get_c_funptr"
    write (unit, "(A)")  "  end type " &
         // char (prefix) // "driver_t"
    write (unit, "(A)")  ""
    write (unit, "(A)")  "contains"
    write (unit, "(A)")  ""
    write (unit, "(A)")  "  function " &
         // char (prefix) // "driver_get_c_funptr (driver, feature) result &
         &(c_fptr)"
    write (unit, "(A)")  "    class(" &
         // char (prefix) // "driver_t), intent(inout) :: driver"
    write (unit, "(A)")  "    type(string_t), intent(in) :: feature"
    write (unit, "(A)")  "    type(c_funptr) :: c_fptr"
    call write_decl ("get_n_processes", "get_n_processes")
    call write_decl ("get_stringptr", "get_process_id_ptr")
    call write_decl ("get_stringptr", "get_model_name_ptr")
    call write_decl ("get_stringptr", "get_md5sum_ptr")
    call write_decl ("get_log", "get_openmp_status")
    call write_decl ("get_int", "get_n_in")
    call write_decl ("get_int", "get_n_out")
    call write_decl ("get_int", "get_n_flv")
    call write_decl ("get_int", "get_n_hel")
    call write_decl ("get_int", "get_n_col")
    call write_decl ("get_int", "get_n_cin")
    call write_decl ("get_int", "get_n_cf")
    call write_decl ("set_int_tab1", "set_flv_state_ptr")
    call write_decl ("set_int_tab1", "set_hel_state_ptr")
    call write_decl ("set_col_state", "set_col_state_ptr")
    call write_decl ("set_color_factors", "set_color_factors_ptr")
    call write_decl ("get_fptr", "get_fptr")
    write (unit, "(A)")  "    select case (char (feature))"
    call write_case ("get_n_processes")
    call write_case ("get_process_id_ptr")
    call write_case ("get_model_name_ptr")
    call write_case ("get_md5sum_ptr")
    call write_case ("get_openmp_status")
    call write_case ("get_n_in")
    call write_case ("get_n_out")
    call write_case ("get_n_flv")
    call write_case ("get_n_hel")
    call write_case ("get_n_col")
    call write_case ("get_n_cin")
    call write_case ("get_n_cf")
    call write_case ("set_flv_state_ptr")
    call write_case ("set_hel_state_ptr")
    call write_case ("set_col_state_ptr")
    call write_case ("set_color_factors_ptr")
    call write_case ("get_fptr")
    write (unit, "(A)")  "    case default"
    write (unit, "(A)")  "       call msg_bug ('prclib2 driver setup: unknown &
         &function name')"
    write (unit, "(A)")  "    end select"
    write (unit, "(A)")  "  end function " &
         // char (prefix) // "driver_get_c_funptr"
    write (unit, "(A)")  ""
    write (unit, "(A)")  "end module " &
         // char (prefix) // "driver"
    write (unit, "(A)")  ""
    write (unit, "(A)")  "! Stand-alone external procedures: used for both &
         &static and dynamic linkage"
  contains
    subroutine write_decl (template, feature)
      character(*), intent(in) :: template, feature
      write (unit, "(A)")  "    procedure(prc_" // template // ") &"
      write (unit, "(A)")  "         :: " &
           // char (prefix) // feature
    end subroutine write_decl
    subroutine write_case (feature)
      character(*), intent(in) :: feature
      write (unit, "(A)")  "    case ('" // feature // "')"
      write (unit, "(A)")  "       c_fptr = c_funloc (" &
           // char (prefix) // feature // ")"
    end subroutine write_case
  end subroutine prclib_driver_write_module

\eatline
\nwendcode{}\nwbegindocs{135}This function provides the overall library MD5sum.  The function is for
internal use (therefore not bind(C)), the external interface is via the
{\tt{}get{\char95}md5sum{\char95}ptr} procedure with index 0.
\nwenddocs{}\nwbegincode{136}\moddef{Prclib interfaces: prclib driver: TBP}\plusendmoddef
  procedure :: write_lib_md5sum_fun => prclib_driver_write_lib_md5sum_fun
\nwendcode{}\nwbegincode{137}\moddef{Prclib interfaces: procedures}\plusendmoddef
  subroutine prclib_driver_write_lib_md5sum_fun (driver, unit, prefix)
    class(prclib_driver_t), intent(in) :: driver
    integer, intent(in) :: unit
    type(string_t), intent(in) :: prefix
    write (unit, "(A)")  ""
    write (unit, "(A)")  "! The MD5 sum of the library"
    write (unit, "(A)")  "function " // char (prefix) &
         // "md5sum () result (md5sum)"
    write (unit, "(A)")  "  implicit none"
    write (unit, "(A)")  "  character(32) :: md5sum"
    write (unit, "(A)")  "  md5sum = '" // driver%md5sum // "'"
    write (unit, "(A)")  "end function " // char (prefix) // "md5sum"
  end subroutine prclib_driver_write_lib_md5sum_fun
    
\eatline
\nwendcode{}\nwbegindocs{138}\nwdocspar
\subsection{Interface bodies for informational functions}
These interfaces implement the communication between WHIZARD (the main
program) and the process-library driver.  The procedures are all
BIND(C), so they can safely be exposed by the library and handled by
the {\tt{}dlopen} mechanism, which apparently understands only C calling
conventions.

In the sections below, for each procedure, we provide both the
interface itself and a procedure that writes the correponding
procedure as source code to the process library driver.

\subsubsection{Process count}
Return the number of processes contained in the library.
\nwenddocs{}\nwbegincode{139}\moddef{Prclib interfaces: public}\plusendmoddef
  public :: prc_get_n_processes
\nwendcode{}\nwbegincode{140}\moddef{Prclib interfaces: interfaces}\plusendmoddef
  abstract interface
     function prc_get_n_processes () result (n) bind(C)
       import
       integer(c_int) :: n
     end function prc_get_n_processes
  end interface
\eatline
\nwendcode{}\nwbegindocs{141}Here is the code.
\nwenddocs{}\nwbegincode{142}\moddef{Prclib interfaces: prclib driver: TBP}\plusendmoddef
  procedure :: write_get_n_processes_fun
\nwendcode{}\nwbegincode{143}\moddef{Prclib interfaces: procedures}\plusendmoddef
  subroutine write_get_n_processes_fun (driver, unit, prefix)
    class(prclib_driver_t), intent(in) :: driver
    integer, intent(in) :: unit
    type(string_t), intent(in) :: prefix
    write (unit, "(A)")  ""
    write (unit, "(A)")  "! Return the number of processes in this library"
    write (unit, "(A)")  "function " // char (prefix) &
         // "get_n_processes () result (n) bind(C)"
    write (unit, "(A)")  "  use iso_c_binding"
    write (unit, "(A)")  "  implicit none"
    write (unit, "(A)")  "  integer(c_int) :: n"
    write (unit, "(A,I0)")  "  n = ", driver%n_processes
    write (unit, "(A)")  "end function " // char (prefix) &
         // "get_n_processes"
  end subroutine write_get_n_processes_fun

\eatline
\nwendcode{}\nwbegindocs{144}\nwdocspar
\subsubsection{Informational string functions}
These functions return constant information about the matrix-element
code.

The following procedures have to return strings.  With the BIND(C)
constraint, we choose to return the C pointer to a string, and its
length, so the procedures implement this interface.  They are actually
subroutines.
\nwenddocs{}\nwbegincode{145}\moddef{Prclib interfaces: public}\plusendmoddef
  public :: prc_get_stringptr
\nwendcode{}\nwbegincode{146}\moddef{Prclib interfaces: interfaces}\plusendmoddef
  abstract interface
     subroutine prc_get_stringptr (i, cptr, len) bind(C)
       import
       integer(c_int), intent(in) :: i
       type(c_ptr), intent(out) :: cptr
       integer(c_int), intent(out) :: len
     end subroutine prc_get_stringptr
  end interface
\eatline
\nwendcode{}\nwbegindocs{147}To hide this complication, we introduce a subroutine that converts the
returned C pointer to a {\tt{}string{\char95}t} object.  As a side effect, we
deallocate the original after conversion -- otherwise, we might have a
memory leak.

For the conversion, we first pointer-convert the C pointer to a
Fortran character array pointer, length 1 and size {\tt{}len}.  Using
argument association and an internal subroutine, we convert this to a
character array with length {\tt{}len} and size 1.  Using ordinary
assignment, we finally convert this to {\tt{}string{\char95}t}.

The function takes the pointer-returning function as an argument.  The
index {\tt{}i} identifies the process in the library.
\nwenddocs{}\nwbegincode{148}\moddef{Prclib interfaces: procedures}\plusendmoddef
  subroutine get_string_via_cptr (string, i, get_stringptr)
    type(string_t), intent(out) :: string
    integer, intent(in) :: i
    procedure(prc_get_stringptr) :: get_stringptr
    type(c_ptr) :: cptr
    integer(c_int) :: pid, len
    character(kind=c_char), dimension(:), pointer :: c_array
    pid = i
    call get_stringptr (pid, cptr, len)
    if (c_associated (cptr)) then
       call c_f_pointer (cptr, c_array, shape = [len])
       call set_string (c_array)
       call get_stringptr (0_c_int, cptr, len)
    else
       string = ""
    end if
  contains
    subroutine set_string (buffer)
      character(len, kind=c_char), dimension(1), intent(in) :: buffer
      string = buffer(1)
    end subroutine set_string
  end subroutine get_string_via_cptr
    
\eatline
\nwendcode{}\nwbegindocs{149}Since the module procedures return Fortran strings, we have to
convert them.  This is the necessary auxiliary routine.  The routine
is not BIND(C), it is not accessed from outside.
\nwenddocs{}\nwbegincode{150}\moddef{Prclib interfaces: prclib driver: TBP}\plusendmoddef
  procedure, nopass :: write_string_to_array_fun
\nwendcode{}\nwbegincode{151}\moddef{Prclib interfaces: procedures}\plusendmoddef
  subroutine write_string_to_array_fun (unit, prefix)
    integer, intent(in) :: unit
    type(string_t), intent(in) :: prefix
    write (unit, "(A)")  ""
    write (unit, "(A)")  "! Auxiliary: convert character string &
         &to array pointer"
    write (unit, "(A)")  "subroutine " // char (prefix) &
         // "string_to_array (string, a)"
    write (unit, "(A)")  "  use iso_c_binding"
    write (unit, "(A)")  "  implicit none"
    write (unit, "(A)")  "  character(*), intent(in) :: string"
    write (unit, "(A)")  "  character(kind=c_char), dimension(:), &
         &allocatable, intent(out) :: a"
    write (unit, "(A)")  "  integer :: i"
    write (unit, "(A)")  "  allocate (a (len (string)))"
    write (unit, "(A)")  "  do i = 1, size (a)"
    write (unit, "(A)")  "     a(i) = string(i:i)"
    write (unit, "(A)")  "  end do"
    write (unit, "(A)")  "end subroutine " // char (prefix) &
         // "string_to_array"
  end subroutine write_string_to_array_fun

\eatline
\nwendcode{}\nwbegindocs{152}The above routine is called by other functions.  It is not in a
module, so they need its interface explicitly.
\nwenddocs{}\nwbegincode{153}\moddef{Prclib interfaces: procedures}\plusendmoddef
  subroutine write_string_to_array_interface (unit, prefix)
    integer, intent(in) :: unit
    type(string_t), intent(in) :: prefix
    write (unit, "(2x,A)")  "interface"
    write (unit, "(2x,A)")  "   subroutine " // char (prefix) &
         // "string_to_array (string, a)"
    write (unit, "(2x,A)")  "     use iso_c_binding"
    write (unit, "(2x,A)")  "     implicit none"
    write (unit, "(2x,A)")  "     character(*), intent(in) :: string"
    write (unit, "(2x,A)")  "     character(kind=c_char), dimension(:), &
         &allocatable, intent(out) :: a"
    write (unit, "(2x,A)")  "   end subroutine " // char (prefix) &
         // "string_to_array"
    write (unit, "(2x,A)")  "end interface"
  end subroutine write_string_to_array_interface

\eatline
\nwendcode{}\nwbegindocs{154}\nwdocspar
Here are the info functions which return strings, implementing the interface
{\tt{}prc{\char95}get{\char95}stringptr}.

Return the process ID for each process.
\nwenddocs{}\nwbegincode{155}\moddef{Prclib interfaces: prclib driver: TBP}\plusendmoddef
  procedure :: write_get_process_id_fun
\nwendcode{}\nwbegincode{156}\moddef{Prclib interfaces: procedures}\plusendmoddef
  subroutine write_get_process_id_fun (driver, unit, prefix)
    class(prclib_driver_t), intent(in) :: driver
    integer, intent(in) :: unit
    type(string_t), intent(in) :: prefix
    integer :: i
    write (unit, "(A)")  ""
    write (unit, "(A)")  "! Return the process ID of process #i &
         &(as a C pointer to a character array)"
    write (unit, "(A)")  "subroutine " // char (prefix) &
         // "get_process_id_ptr (i, cptr, len) bind(C)"
    write (unit, "(A)")  "  use iso_c_binding"
    write (unit, "(A)")  "  implicit none"
    write (unit, "(A)")  "  integer(c_int), intent(in) :: i"
    write (unit, "(A)")  "  type(c_ptr), intent(out) :: cptr"
    write (unit, "(A)")  "  integer(c_int), intent(out) :: len"
    write (unit, "(A)")  "  character(kind=c_char), dimension(:), &
         &allocatable, target, save :: a"
    call write_string_to_array_interface (unit, prefix)
    write (unit, "(A)")  "  select case (i)"
    write (unit, "(A)")  "  case (0);  if (allocated (a))  deallocate (a)"
    do i = 1, driver%n_processes
       write (unit, "(A,I0,9A)")  "  case (", i, ");  ", &
            "call ", char (prefix), "string_to_array ('", &
            char (driver%record(i)%id), "', a)"
    end do
    write (unit, "(A)")  "  end select"
    write (unit, "(A)")  "  if (allocated (a)) then"
    write (unit, "(A)")  "     cptr = c_loc (a)"
    write (unit, "(A)")  "     len = size (a)"
    write (unit, "(A)")  "  else"
    write (unit, "(A)")  "     cptr = c_null_ptr"
    write (unit, "(A)")  "     len = 0"
    write (unit, "(A)")  "  end if"
    write (unit, "(A)")  "end subroutine " // char (prefix) &
         // "get_process_id_ptr"
  end subroutine write_get_process_id_fun
  
\eatline
\nwendcode{}\nwbegindocs{157}Return the model name, given explicitly.
\nwenddocs{}\nwbegincode{158}\moddef{Prclib interfaces: prclib driver: TBP}\plusendmoddef
  procedure :: write_get_model_name_fun
\nwendcode{}\nwbegincode{159}\moddef{Prclib interfaces: procedures}\plusendmoddef
  subroutine write_get_model_name_fun (driver, unit, prefix)
    class(prclib_driver_t), intent(in) :: driver
    integer, intent(in) :: unit
    type(string_t), intent(in) :: prefix
    integer :: i
    write (unit, "(A)")  ""
    write (unit, "(A)")  "! Return the model name for process #i &
         &(as a C pointer to a character array)"
    write (unit, "(A)")  "subroutine " // char (prefix) &
         // "get_model_name_ptr (i, cptr, len) bind(C)"
    write (unit, "(A)")  "  use iso_c_binding"
    write (unit, "(A)")  "  implicit none"
    write (unit, "(A)")  "  integer(c_int), intent(in) :: i"
    write (unit, "(A)")  "  type(c_ptr), intent(out) :: cptr"
    write (unit, "(A)")  "  integer(c_int), intent(out) :: len"
    write (unit, "(A)")  "  character(kind=c_char), dimension(:), &
         &allocatable, target, save :: a"
    call write_string_to_array_interface (unit, prefix)
    write (unit, "(A)")  "  select case (i)"
    write (unit, "(A)")  "  case (0);  if (allocated (a))  deallocate (a)"
    do i = 1, driver%n_processes
       write (unit, "(A,I0,9A)")  "  case (", i, ");  ", &
            "call ", char (prefix), "string_to_array ('" , &
            char (driver%record(i)%model_name), &
            "', a)"
    end do
    write (unit, "(A)")  "  end select"
    write (unit, "(A)")  "  if (allocated (a)) then"
    write (unit, "(A)")  "     cptr = c_loc (a)"
    write (unit, "(A)")  "     len = size (a)"
    write (unit, "(A)")  "  else"
    write (unit, "(A)")  "     cptr = c_null_ptr"
    write (unit, "(A)")  "     len = 0"
    write (unit, "(A)")  "  end if"
    write (unit, "(A)")  "end subroutine " // char (prefix) &
         // "get_model_name_ptr"
  end subroutine write_get_model_name_fun

\eatline
\nwendcode{}\nwbegindocs{160}Call the MD5 sum function for the process.  The function calls the
corresponding function of the matrix-element code, and it returns the
C address of a character array with length 32.
\nwenddocs{}\nwbegincode{161}\moddef{Prclib interfaces: prclib driver: TBP}\plusendmoddef
  procedure :: write_get_md5sum_fun
\nwendcode{}\nwbegincode{162}\moddef{Prclib interfaces: procedures}\plusendmoddef
  subroutine write_get_md5sum_fun (driver, unit, prefix)
    class(prclib_driver_t), intent(in) :: driver
    integer, intent(in) :: unit
    type(string_t), intent(in) :: prefix
    integer :: i
    write (unit, "(A)")  ""
    write (unit, "(A)")  "! Return the MD5 sum for the process configuration &
         &(as a C pointer to a character array)"
    write (unit, "(A)")  "subroutine " // char (prefix) &
         // "get_md5sum_ptr (i, cptr, len) bind(C)"
    write (unit, "(A)")  "  use iso_c_binding"
    call driver%write_interfaces (unit, var_str ("md5sum"))
    write (unit, "(A)")  "  interface"
    write (unit, "(A)")  "     function " // char (prefix) &
         // "md5sum () result (md5sum)"
    write (unit, "(A)")  "       character(32) :: md5sum"
    write (unit, "(A)")  "     end function " // char (prefix) // "md5sum"
    write (unit, "(A)")  "  end interface"
    write (unit, "(A)")  "  integer(c_int), intent(in) :: i"
    write (unit, "(A)")  "  type(c_ptr), intent(out) :: cptr"
    write (unit, "(A)")  "  integer(c_int), intent(out) :: len"
    write (unit, "(A)")  "  character(kind=c_char), dimension(32), &
         &target, save :: md5sum"
    write (unit, "(A)")  "  select case (i)"
    write (unit, "(A)")  "  case (0)"
    write (unit, "(A)")  "     call copy (" // char (prefix) // "md5sum ())"
    write (unit, "(A)")  "     cptr = c_loc (md5sum)"
    do i = 1, driver%n_processes
       write (unit, "(A,I0,A)")  "  case (", i, ")"
       call driver%record(i)%write_md5sum_call (unit)
    end do
    write (unit, "(A)")  "  case default"
    write (unit, "(A)")  "     cptr = c_null_ptr"
    write (unit, "(A)")  "  end select"
    write (unit, "(A)")  "  len = 32"
    write (unit, "(A)")  "contains"
    write (unit, "(A)")  "  subroutine copy (md5sum_tmp)"
    write (unit, "(A)")  "    character, dimension(32), intent(in) :: &
         &md5sum_tmp"
    write (unit, "(A)")  "    md5sum = md5sum_tmp"
    write (unit, "(A)")  "  end subroutine copy"
    write (unit, "(A)")  "end subroutine " // char (prefix) &
           // "get_md5sum_ptr"
  end subroutine write_get_md5sum_fun

\eatline
\nwendcode{}\nwbegindocs{163}The actual call depends on the type of matrix element.
\nwenddocs{}\nwbegincode{164}\moddef{Prclib interfaces: prclib driver record: TBP}\plusendmoddef
  procedure :: write_md5sum_call => prclib_driver_record_write_md5sum_call
\nwendcode{}\nwbegincode{165}\moddef{Prclib interfaces: procedures}\plusendmoddef
  subroutine prclib_driver_record_write_md5sum_call (record, unit)
    class(prclib_driver_record_t), intent(in) :: record
    integer, intent(in) :: unit
    call record%writer%write_md5sum_call (unit, record%id)
  end subroutine prclib_driver_record_write_md5sum_call
  
\eatline
\nwendcode{}\nwbegindocs{166}The interface goes into the writer base type:
\nwenddocs{}\nwbegincode{167}\moddef{Prclib interfaces: prc writer: TBP}\plusendmoddef
  procedure(write_code), deferred :: write_md5sum_call
\eatline
\nwendcode{}\nwbegindocs{168}In the Fortran module case, we take a detour.  The string returned
by the Fortran function is copied into a fixed-size array.  The copy
routine is an internal subroutine of {\tt{}get{\char95}md5sum{\char95}ptr}.  We
return the C address of the target array.
\nwenddocs{}\nwbegincode{169}\moddef{Prclib interfaces: prc writer f module: TBP}\plusendmoddef
  procedure :: write_md5sum_call => prc_writer_f_module_write_md5sum_call
\nwendcode{}\nwbegincode{170}\moddef{Prclib interfaces: procedures}\plusendmoddef
  subroutine prc_writer_f_module_write_md5sum_call (writer, unit, id)
    class(prc_writer_f_module_t), intent(in) :: writer
    integer, intent(in) :: unit
    type(string_t), intent(in) :: id
    write (unit, "(5x,9A)")  "call copy (", &
         char (writer%get_c_procname (id, var_str ("md5sum"))), " ())"
    write (unit, "(5x,9A)")  "cptr = c_loc (md5sum)"
  end subroutine prc_writer_f_module_write_md5sum_call
  
\eatline
\nwendcode{}\nwbegindocs{171}In the C library case, the library function returns a C pointer,
which we can just copy.
\nwenddocs{}\nwbegincode{172}\moddef{Prclib interfaces: prc writer c lib: TBP}\endmoddef
  procedure :: write_md5sum_call => prc_writer_c_lib_write_md5sum_call
\nwendcode{}\nwbegincode{173}\moddef{Prclib interfaces: procedures}\plusendmoddef
  subroutine prc_writer_c_lib_write_md5sum_call (writer, unit, id)
    class(prc_writer_c_lib_t), intent(in) :: writer
    integer, intent(in) :: unit
    type(string_t), intent(in) :: id
    write (unit, "(5x,9A)") &
         "cptr =  ", &
         char (writer%get_c_procname (id, var_str ("get_md5sum"))), " ()"
  end subroutine prc_writer_c_lib_write_md5sum_call
  
\eatline
\nwendcode{}\nwbegindocs{174}\nwdocspar
\subsubsection{Actual references to the info functions}
The string-valued info functions return C character arrays.  For the
API of the library driver, we provide convenience functions which
(re)convert those arrays into {\tt{}string{\char95}t} objects.
\nwenddocs{}\nwbegincode{175}\moddef{Prclib interfaces: prclib driver: TBP}\plusendmoddef
  procedure :: get_process_id => prclib_driver_get_process_id
  procedure :: get_model_name => prclib_driver_get_model_name
  procedure :: get_md5sum => prclib_driver_get_md5sum
\nwendcode{}\nwbegincode{176}\moddef{Prclib interfaces: procedures}\plusendmoddef
  function prclib_driver_get_process_id (driver, i) result (string)
    type(string_t) :: string
    class(prclib_driver_t), intent(in) :: driver
    integer, intent(in) :: i
    call get_string_via_cptr (string, i, driver%get_process_id_ptr)
  end function prclib_driver_get_process_id
  
  function prclib_driver_get_model_name (driver, i) result (string)
    type(string_t) :: string
    class(prclib_driver_t), intent(in) :: driver
    integer, intent(in) :: i
    call get_string_via_cptr (string, i, driver%get_model_name_ptr)
  end function prclib_driver_get_model_name
  
  function prclib_driver_get_md5sum (driver, i) result (string)
    type(string_t) :: string
    class(prclib_driver_t), intent(in) :: driver
    integer, intent(in) :: i
    call get_string_via_cptr (string, i, driver%get_md5sum_ptr)
  end function prclib_driver_get_md5sum
  
\eatline
\eatline
\eatline
\nwendcode{}\nwbegindocs{177}\nwdocspar
\subsubsection{Informational logical functions}
When returning a logical value, we use the C boolean type, which
may differ from Fortran.
\nwenddocs{}\nwbegincode{178}\moddef{Prclib interfaces: public}\plusendmoddef
  public :: prc_get_log
\nwendcode{}\nwbegincode{179}\moddef{Prclib interfaces: interfaces}\plusendmoddef
  abstract interface
     function prc_get_log (pid) result (l) bind(C)
       import
       integer(c_int), intent(in) :: pid
       logical(c_bool) :: l
     end function prc_get_log
  end interface
\eatline
\nwendcode{}\nwbegindocs{180}Return a logical flag which tells whether OpenMP is supported for a
specific process code.
\nwenddocs{}\nwbegincode{181}\moddef{Prclib interfaces: prclib driver: TBP}\plusendmoddef
  procedure :: write_get_openmp_status_fun
\nwendcode{}\nwbegincode{182}\moddef{Prclib interfaces: procedures}\plusendmoddef
  subroutine write_get_openmp_status_fun (driver, unit, prefix)
    class(prclib_driver_t), intent(in) :: driver
    integer, intent(in) :: unit
    type(string_t), intent(in) :: prefix
    integer :: i
    write (unit, "(A)")  ""
    write (unit, "(A)")  "! Return the OpenMP support status"
    write (unit, "(A)")  "function " // char (prefix) &
         // "get_openmp_status (i) result (openmp_status) bind(C)"
    write (unit, "(A)")  "  use iso_c_binding"
    call driver%write_interfaces (unit, var_str ("openmp_supported"))
    write (unit, "(A)")  "  integer(c_int), intent(in) :: i"
    write (unit, "(A)")  "  logical(c_bool) :: openmp_status"
    write (unit, "(A)")  "  select case (i)"
    do i = 1, driver%n_processes
       write (unit, "(A,I0,9A)")  "  case (", i, ");  ", &
            "openmp_status = ", &
            char (driver%record(i)%get_c_procname &
            (var_str ("openmp_supported"))), " ()"
    end do
    write (unit, "(A)")  "  end select"
    write (unit, "(A)")  "end function " // char (prefix) &
         // "get_openmp_status"
  end subroutine write_get_openmp_status_fun

\eatline
\nwendcode{}\nwbegindocs{183}\nwdocspar
\subsubsection{Informational integer functions}
Various process metadata are integer values.  We can use a single
interface for all of them.
\nwenddocs{}\nwbegincode{184}\moddef{Prclib interfaces: public}\plusendmoddef
  public :: prc_get_int
\nwendcode{}\nwbegincode{185}\moddef{Prclib interfaces: interfaces}\plusendmoddef
  abstract interface
     function prc_get_int (pid) result (n) bind(C)
       import
       integer(c_int), intent(in) :: pid
       integer(c_int) :: n
     end function prc_get_int
  end interface
\eatline
\nwendcode{}\nwbegindocs{186}This function returns any data of type integer, for each process.
\nwenddocs{}\nwbegincode{187}\moddef{Prclib interfaces: prclib driver: TBP}\plusendmoddef
  procedure :: write_get_int_fun
\nwendcode{}\nwbegincode{188}\moddef{Prclib interfaces: procedures}\plusendmoddef
  subroutine write_get_int_fun (driver, unit, prefix, feature)
    class(prclib_driver_t), intent(in) :: driver
    integer, intent(in) :: unit
    type(string_t), intent(in) :: prefix
    type(string_t), intent(in) :: feature
    integer :: i
    write (unit, "(A)")  ""
    write (unit, "(9A)")  "! Return the value of ", char (feature)
    write (unit, "(9A)")  "function ", char (prefix), &
         "get_", char (feature), " (pid)", &
         " result (", char (feature), ") bind(C)"
    write (unit, "(9A)")  "  use iso_c_binding"
    call driver%write_interfaces (unit, feature)
    write (unit, "(9A)")  "  integer(c_int), intent(in) :: pid"
    write (unit, "(9A)")  "  integer(c_int) :: ", char (feature)
    write (unit, "(9A)")  "  select case (pid)"
    do i = 1, driver%n_processes
       write (unit, "(2x,A,I0,9A)")  "case (", i, ");  ", &
            char (feature), " = ", &
            char (driver%record(i)%get_c_procname (feature)), &
            " ()"
    end do
    write (unit, "(9A)")  "  end select"
    write (unit, "(9A)")  "end function ", char (prefix), &
         "get_", char (feature)
  end subroutine write_get_int_fun
  
\eatline
\nwendcode{}\nwbegindocs{189}Write a {\tt{}case} line that assigns the value of the external function
to the current return value.
\nwenddocs{}\nwbegincode{190}\moddef{Prclib interfaces: procedures}\plusendmoddef
  subroutine write_case_int_fun (record, unit, i, feature)
    class(prclib_driver_record_t), intent(in) :: record
    integer, intent(in) :: unit
    integer, intent(in) :: i
    type(string_t), intent(in) :: feature
    write (unit, "(5x,A,I0,9A)")  "case (", i, ");  ", &
         char (feature), " = ", char (record%get_c_procname (feature))
  end subroutine write_case_int_fun
  
\eatline
\nwendcode{}\nwbegindocs{191}\nwdocspar
\subsubsection{Flavor and helicity tables}
Transferring tables is more complicated.  First, a two-dimensional array.
\nwenddocs{}\nwbegincode{192}\moddef{Prclib interfaces: public}\plusendmoddef
  public :: prc_set_int_tab1
\nwendcode{}\nwbegincode{193}\moddef{Prclib interfaces: interfaces}\plusendmoddef
  abstract interface
     subroutine prc_set_int_tab1 (pid, tab, shape) bind(C)
       import
       integer(c_int), intent(in) :: pid
       integer(c_int), dimension(*), intent(out) :: tab
       integer(c_int), dimension(2), intent(in) :: shape
     end subroutine prc_set_int_tab1
  end interface
\eatline
\nwendcode{}\nwbegindocs{194}This subroutine returns a table of integers.
\nwenddocs{}\nwbegincode{195}\moddef{Prclib interfaces: prclib driver: TBP}\plusendmoddef
  procedure :: write_set_int_sub
\nwendcode{}\nwbegincode{196}\moddef{Prclib interfaces: procedures}\plusendmoddef
  subroutine write_set_int_sub (driver, unit, prefix, feature)
    class(prclib_driver_t), intent(in) :: driver
    integer, intent(in) :: unit
    type(string_t), intent(in) :: prefix
    type(string_t), intent(in) :: feature
    integer :: i
    write (unit, "(A)")  ""
    write (unit, "(9A)")  "! Set table: ", char (feature)
    write (unit, "(9A)")  "subroutine ", char (prefix), &
         "set_", char (feature), "_ptr (pid, ", char (feature), &
         ", shape) bind(C)"
    write (unit, "(9A)")  "  use iso_c_binding"
    call driver%write_interfaces (unit, feature)
    write (unit, "(9A)")  "  integer(c_int), intent(in) :: pid"
    write (unit, "(9A)")  "  integer(c_int), dimension(*), intent(out) :: ", &
         char (feature)
    write (unit, "(9A)")  "  integer(c_int), dimension(2), intent(in) :: shape"
    write (unit, "(9A)")  "  integer, dimension(:,:), allocatable :: ", &
         char (feature), "_tmp"
    write (unit, "(9A)")  "  integer :: i, j"
    write (unit, "(9A)")  "  select case (pid)"
    do i = 1, driver%n_processes
       write (unit, "(2x,A,I0,A)")  "case (", i, ")"
       call driver%record(i)%write_int_sub_call (unit, feature)
    end do
    write (unit, "(9A)")  "  end select"
    write (unit, "(9A)")  "end subroutine ", char (prefix), &
         "set_", char (feature), "_ptr"
  end subroutine write_set_int_sub
  
\eatline
\nwendcode{}\nwbegindocs{197}The actual call depends on the type of matrix element.
\nwenddocs{}\nwbegincode{198}\moddef{Prclib interfaces: prclib driver record: TBP}\plusendmoddef
  procedure :: write_int_sub_call => prclib_driver_record_write_int_sub_call
\nwendcode{}\nwbegincode{199}\moddef{Prclib interfaces: procedures}\plusendmoddef
  subroutine prclib_driver_record_write_int_sub_call (record, unit, feature)
    class(prclib_driver_record_t), intent(in) :: record
    integer, intent(in) :: unit
    type(string_t), intent(in) :: feature
    call record%writer%write_int_sub_call (unit, record%id, feature)
  end subroutine prclib_driver_record_write_int_sub_call
  
\eatline
\nwendcode{}\nwbegindocs{200}The interface goes into the writer base type:
\nwenddocs{}\nwbegincode{201}\moddef{Prclib interfaces: prc writer: TBP}\plusendmoddef
  procedure(write_feature_code), deferred :: write_int_sub_call
\eatline
\nwendcode{}\nwbegindocs{202}In the Fortran module case, we need an extra copy in the
(academical) situation where default integer and {\tt{}c{\char95}int} differ.
Otherwise, we just associate a Fortran array with the C pointer and
let the matrix-element subroutine fill the array.
\nwenddocs{}\nwbegincode{203}\moddef{Prclib interfaces: prc writer f module: TBP}\plusendmoddef
  procedure :: write_int_sub_call => prc_writer_f_module_write_int_sub_call
\nwendcode{}\nwbegincode{204}\moddef{Prclib interfaces: procedures}\plusendmoddef
  subroutine prc_writer_f_module_write_int_sub_call (writer, unit, id, feature)
    class(prc_writer_f_module_t), intent(in) :: writer
    integer, intent(in) :: unit
    type(string_t), intent(in) :: id, feature
    write (unit, "(5x,9A)")  "allocate (",  char (feature), "_tmp ", &
         "(shape(1), shape(2)))"
    write (unit, "(5x,9A)")  "call ", &
         char (writer%get_c_procname (id, feature)), &
         " (", char (feature), "_tmp)"
    write (unit, "(5x,9A)")  "forall (i=1:shape(1), j=1:shape(2)) "
    write (unit, "(8x,9A)")  char (feature), "(i + shape(1)*(j-1)) = ", &
         char (feature), "_tmp", "(i,j)"
    write (unit, "(5x,9A)")  "end forall"
  end subroutine prc_writer_f_module_write_int_sub_call
  
\eatline
\nwendcode{}\nwbegindocs{205}In the C library case, we just transfer the C pointer to the library
function.
\nwenddocs{}\nwbegincode{206}\moddef{Prclib interfaces: prc writer c lib: TBP}\plusendmoddef
  procedure :: write_int_sub_call => prc_writer_c_lib_write_int_sub_call
\nwendcode{}\nwbegincode{207}\moddef{Prclib interfaces: procedures}\plusendmoddef
  subroutine prc_writer_c_lib_write_int_sub_call (writer, unit, id, feature)
    class(prc_writer_c_lib_t), intent(in) :: writer
    integer, intent(in) :: unit
    type(string_t), intent(in) :: id, feature
    write (unit, "(5x,9A)")  "call ", &
         char (writer%get_c_procname (id, feature)), " (", char (feature), ")"
  end subroutine prc_writer_c_lib_write_int_sub_call
  
\eatline
\nwendcode{}\nwbegindocs{208}\nwdocspar
\subsubsection{Color state table}
The color-state specification needs a table of integers (one array per
color flow) and a corresponding array of color-ghost flags.
\nwenddocs{}\nwbegincode{209}\moddef{Prclib interfaces: public}\plusendmoddef
  public :: prc_set_col_state
\nwendcode{}\nwbegincode{210}\moddef{Prclib interfaces: interfaces}\plusendmoddef
  abstract interface
     subroutine prc_set_col_state (pid, col_state, ghost_flag, shape) bind(C)
       import
       integer(c_int), intent(in) :: pid
       integer(c_int), dimension(*), intent(out) :: col_state
       logical(c_bool), dimension(*), intent(out) :: ghost_flag
       integer(c_int), dimension(3), intent(in) :: shape
     end subroutine prc_set_col_state
  end interface
\eatline
\nwendcode{}\nwbegindocs{211}\nwdocspar
\nwenddocs{}\nwbegincode{212}\moddef{Prclib interfaces: prclib driver: TBP}\plusendmoddef
  procedure :: write_set_col_state_sub
\nwendcode{}\nwbegincode{213}\moddef{Prclib interfaces: procedures}\plusendmoddef
  subroutine write_set_col_state_sub (driver, unit, prefix)
    class(prclib_driver_t), intent(in) :: driver
    integer, intent(in) :: unit
    type(string_t), intent(in) :: prefix
    integer :: i
    type(string_t) :: feature
    feature = "col_state"
    write (unit, "(A)")  ""
    write (unit, "(9A)")  "! Set tables: col_state, ghost_flag"
    write (unit, "(9A)")  "subroutine ", char (prefix), &
         "set_col_state_ptr (pid, col_state, ghost_flag, shape) bind(C)"
    write (unit, "(9A)")  "  use iso_c_binding"
    call driver%write_interfaces (unit, feature)
    write (unit, "(9A)")  "  integer(c_int), intent(in) :: pid"
    write (unit, "(9A)") &
         "  integer(c_int), dimension(*), intent(out) :: col_state"
    write (unit, "(9A)") &
         "  logical(c_bool), dimension(*), intent(out) :: ghost_flag"
    write (unit, "(9A)") &
         "  integer(c_int), dimension(3), intent(in) :: shape"
    write (unit, "(9A)") &
         "  integer, dimension(:,:,:), allocatable :: col_state_tmp"
    write (unit, "(9A)") &
         "  logical, dimension(:,:), allocatable :: ghost_flag_tmp"
    write (unit, "(9A)") "  integer :: i, j, k"
    write (unit, "(A)")  "  select case (pid)"
    do i = 1, driver%n_processes
       write (unit, "(A,I0,A)")  "  case (", i, ")"
       call driver%record(i)%write_col_state_call (unit)
    end do
    write (unit, "(A)")  "  end select"
    write (unit, "(9A)")  "end subroutine ", char (prefix), &
         "set_col_state_ptr"
  end subroutine write_set_col_state_sub

\eatline
\nwendcode{}\nwbegindocs{214}The actual call depends on the type of matrix element.
\nwenddocs{}\nwbegincode{215}\moddef{Prclib interfaces: prclib driver record: TBP}\plusendmoddef
  procedure :: write_col_state_call => prclib_driver_record_write_col_state_call
\nwendcode{}\nwbegincode{216}\moddef{Prclib interfaces: procedures}\plusendmoddef
  subroutine prclib_driver_record_write_col_state_call (record, unit)
    class(prclib_driver_record_t), intent(in) :: record
    integer, intent(in) :: unit
    call record%writer%write_col_state_call (unit, record%id)
  end subroutine prclib_driver_record_write_col_state_call
  
\eatline
\nwendcode{}\nwbegindocs{217}The interface goes into the writer base type:
\nwenddocs{}\nwbegincode{218}\moddef{Prclib interfaces: prc writer: TBP}\plusendmoddef
  procedure(write_code), deferred :: write_col_state_call
\eatline
\nwendcode{}\nwbegindocs{219}In the Fortran module case, we need an extra copy in the
(academical) situation where default integer and {\tt{}c{\char95}int} differ.
Otherwise, we just associate a Fortran array with the C pointer and
let the matrix-element subroutine fill the array.
\nwenddocs{}\nwbegincode{220}\moddef{Prclib interfaces: prc writer f module: TBP}\plusendmoddef
  procedure :: write_col_state_call => prc_writer_f_module_write_col_state_call
\nwendcode{}\nwbegincode{221}\moddef{Prclib interfaces: procedures}\plusendmoddef
  subroutine prc_writer_f_module_write_col_state_call (writer, unit, id)
    class(prc_writer_f_module_t), intent(in) :: writer
    integer, intent(in) :: unit
    type(string_t), intent(in) :: id
    write (unit, "(9A)")  "  allocate (col_state_tmp ", &
            "(shape(1), shape(2), shape(3)))"
    write (unit, "(5x,9A)")  "allocate (ghost_flag_tmp ", &
               "(shape(2), shape(3)))"
    write (unit, "(5x,9A)")  "call ", &
         char (writer%get_c_procname (id, var_str ("col_state"))), &
         " (col_state_tmp, ghost_flag_tmp)"
    write (unit, "(5x,9A)")  "forall (i = 1:shape(2), j = 1:shape(3))"
    write (unit, "(8x,9A)")  "forall (k = 1:shape(1))"
    write (unit, "(11x,9A)")  &
         "col_state(k + shape(1) * (i + shape(2)*(j-1) - 1)) ", &
         "= col_state_tmp(k,i,j)"
    write (unit, "(8x,9A)")  "end forall"
    write (unit, "(8x,9A)")  &
         "ghost_flag(i + shape(2)*(j-1)) = ghost_flag_tmp(i,j)"
    write (unit, "(5x,9A)")  "end forall"
  end subroutine prc_writer_f_module_write_col_state_call
  
\eatline
\nwendcode{}\nwbegindocs{222}In the C library case, we just transfer the C pointer to the library
function.
\nwenddocs{}\nwbegincode{223}\moddef{Prclib interfaces: prc writer c lib: TBP}\plusendmoddef
  procedure :: write_col_state_call => prc_writer_c_lib_write_col_state_call
\nwendcode{}\nwbegincode{224}\moddef{Prclib interfaces: procedures}\plusendmoddef
  subroutine prc_writer_c_lib_write_col_state_call (writer, unit, id)
    class(prc_writer_c_lib_t), intent(in) :: writer
    integer, intent(in) :: unit
    type(string_t), intent(in) :: id
    write (unit, "(5x,9A)")  "call ", &
         char (writer%get_c_procname (id, var_str ("col_state"))), &
         " (col_state, ghost_flag)"
  end subroutine prc_writer_c_lib_write_col_state_call
  
\eatline
\nwendcode{}\nwbegindocs{225}\nwdocspar
\subsubsection{Color factors}
For the color-factor information, we return two integer arrays and a
complex array. 
\nwenddocs{}\nwbegincode{226}\moddef{Prclib interfaces: public}\plusendmoddef
  public :: prc_set_color_factors
\nwendcode{}\nwbegincode{227}\moddef{Prclib interfaces: interfaces}\plusendmoddef
  abstract interface
     subroutine prc_set_color_factors &
          (pid, cf_index1, cf_index2, color_factors, shape) bind(C)
       import
       integer(c_int), intent(in) :: pid
       integer(c_int), dimension(*), intent(out) :: cf_index1, cf_index2
       complex(c_default_complex), dimension(*), intent(out) :: color_factors
       integer(c_int), dimension(1), intent(in) :: shape
     end subroutine prc_set_color_factors
  end interface

\eatline
\nwendcode{}\nwbegindocs{228}This subroutine returns the color-flavor factor table.
\nwenddocs{}\nwbegincode{229}\moddef{Prclib interfaces: prclib driver: TBP}\plusendmoddef
  procedure :: write_set_color_factors_sub
\nwendcode{}\nwbegincode{230}\moddef{Prclib interfaces: procedures}\plusendmoddef
  subroutine write_set_color_factors_sub (driver, unit, prefix)
    class(prclib_driver_t), intent(in) :: driver
    integer, intent(in) :: unit
    type(string_t), intent(in) :: prefix
    integer :: i
    type(string_t) :: feature
    feature = "color_factors"
    write (unit, "(A)")  ""
    write (unit, "(A)")  "! Set tables: color factors"
    write (unit, "(9A)")  "subroutine ", char (prefix), &
         "set_color_factors_ptr (pid, cf_index1, cf_index2, color_factors, ", &
         "shape) bind(C)"
    write (unit, "(A)")  "  use iso_c_binding"
    write (unit, "(A)")  "  use kinds"
    write (unit, "(A)")  "  use omega_color"
    call driver%write_interfaces (unit, feature)
    write (unit, "(A)")  "  integer(c_int), intent(in) :: pid"
    write (unit, "(A)")  "  integer(c_int), dimension(1), intent(in) :: shape"
    write (unit, "(A)")  "  integer(c_int), dimension(*), intent(out) :: &
         &cf_index1, cf_index2"
    write (unit, "(A)")  "  complex(c_default_complex), dimension(*), &
         &intent(out) :: color_factors"
    write (unit, "(A)")  "  type(omega_color_factor), dimension(:), &
         &allocatable :: cf"
    write (unit, "(A)")  "  select case (pid)"
    do i = 1, driver%n_processes
       write (unit, "(2x,A,I0,A)")  "case (", i, ")"
       call driver%record(i)%write_color_factors_call (unit)
    end do
    write (unit, "(A)")  "  end select"
    write (unit, "(A)")  "end subroutine " // char (prefix) &
         // "set_color_factors_ptr"
  end subroutine write_set_color_factors_sub
  
\eatline
\nwendcode{}\nwbegindocs{231}The actual call depends on the type of matrix element.
\nwenddocs{}\nwbegincode{232}\moddef{Prclib interfaces: prclib driver record: TBP}\plusendmoddef
  procedure :: write_color_factors_call => prclib_driver_record_write_color_factors_call
\nwendcode{}\nwbegincode{233}\moddef{Prclib interfaces: procedures}\plusendmoddef
  subroutine prclib_driver_record_write_color_factors_call (record, unit)
    class(prclib_driver_record_t), intent(in) :: record
    integer, intent(in) :: unit
    call record%writer%write_color_factors_call (unit, record%id)
  end subroutine prclib_driver_record_write_color_factors_call
  
\eatline
\nwendcode{}\nwbegindocs{234}The interface goes into the writer base type:
\nwenddocs{}\nwbegincode{235}\moddef{Prclib interfaces: prc writer: TBP}\plusendmoddef
  procedure(write_code), deferred :: write_color_factors_call
\eatline
\nwendcode{}\nwbegindocs{236}In the Fortran module case, the matrix-element procedure fills an
array of {\tt{}omega{\char95}color{\char95}factor} elements.  We distribute this array
among two integer arrays and one complex-valued array, for which we
have the C pointers.
\nwenddocs{}\nwbegincode{237}\moddef{Prclib interfaces: prc writer f module: TBP}\plusendmoddef
  procedure :: write_color_factors_call => prc_writer_f_module_write_color_factors_call
\nwendcode{}\nwbegincode{238}\moddef{Prclib interfaces: procedures}\plusendmoddef
  subroutine prc_writer_f_module_write_color_factors_call (writer, unit, id)
    class(prc_writer_f_module_t), intent(in) :: writer
    integer, intent(in) :: unit
    type(string_t), intent(in) :: id
    write (unit, "(5x,A)")  "allocate (cf (shape(1)))"
    write (unit, "(5x,9A)")  "call ", &
         char (writer%get_c_procname (id, var_str ("color_factors"))), " (cf)"
    write (unit, "(5x,9A)")  "cf_index1(1:shape(1)) = cf%i1"
    write (unit, "(5x,9A)")  "cf_index2(1:shape(1)) = cf%i2"
    write (unit, "(5x,9A)")  "color_factors(1:shape(1)) = cf%factor"
  end subroutine prc_writer_f_module_write_color_factors_call
  
\eatline
\nwendcode{}\nwbegindocs{239}In the C library case, we just transfer the C pointers to the library
function.  There are three arrays.
\nwenddocs{}\nwbegincode{240}\moddef{Prclib interfaces: prc writer c lib: TBP}\plusendmoddef
  procedure :: write_color_factors_call => &
       prc_writer_c_lib_write_color_factors_call
\nwendcode{}\nwbegincode{241}\moddef{Prclib interfaces: procedures}\plusendmoddef
  subroutine prc_writer_c_lib_write_color_factors_call (writer, unit, id)
    class(prc_writer_c_lib_t), intent(in) :: writer
    integer, intent(in) :: unit
    type(string_t), intent(in) :: id
    write (unit, "(5x,9A)")  "call ", &
         char (writer%get_c_procname (id, var_str ("color_factors"))), &
         " (cf_index1, cf_index2, color_factors)"
  end subroutine prc_writer_c_lib_write_color_factors_call
  
\eatline
\nwendcode{}\nwbegindocs{242}\nwdocspar
\subsection{Interfaces for C-library matrix element}
If the matrix element code is not provided as a Fortran module but as
a C or bind(C) Fortran library, we need explicit interfaces for the
library functions.  They are not identical to the Fortran module
versions.  They transfer pointers directly.

The implementation is part of the {\tt{}prc{\char95}writer{\char95}c{\char95}lib} type, which
serves as base type for all C-library writers.  It writes specific
interfaces depending on the feature.

We bind this as the method {\tt{}write{\char95}standard{\char95}interface} instead of
{\tt{}write{\char95}interface}, because we have to override the latter.
Otherwise we could not call the method because the writer type is
abstract.
\nwenddocs{}\nwbegincode{243}\moddef{Prclib interfaces: prc writer c lib: TBP}\plusendmoddef
  procedure :: write_standard_interface => prc_writer_c_lib_write_interface
\nwendcode{}\nwbegincode{244}\moddef{Prclib interfaces: procedures}\plusendmoddef
  subroutine prc_writer_c_lib_write_interface (writer, unit, id, feature)
    class(prc_writer_c_lib_t), intent(in) :: writer
    integer, intent(in) :: unit
    type(string_t), intent(in) :: id, feature
    select case (char (feature))
    case ("md5sum")
       write (unit, "(2x,9A)")  "interface"
       write (unit, "(5x,9A)")  "function ", &
            char (writer%get_c_procname (id, var_str ("get_md5sum"))), &
            " () result (cptr) bind(C)"
       write (unit, "(7x,9A)")  "import"
       write (unit, "(7x,9A)")  "implicit none"
       write (unit, "(7x,9A)")  "type(c_ptr) :: cptr"    
       write (unit, "(5x,9A)")  "end function ", &
            char (writer%get_c_procname (id, var_str ("get_md5sum")))
       write (unit, "(2x,9A)")  "end interface"
    case ("openmp_supported")
       write (unit, "(2x,9A)")  "interface"
       write (unit, "(5x,9A)")  "function ", &
            char (writer%get_c_procname (id, feature)), &
            " () result (status) bind(C)"
       write (unit, "(7x,9A)")  "import"
       write (unit, "(7x,9A)")  "implicit none"
       write (unit, "(7x,9A)")  "logical(c_bool) :: status"    
       write (unit, "(5x,9A)")  "end function ", &
            char (writer%get_c_procname (id, feature))
       write (unit, "(2x,9A)")  "end interface"
    case ("n_in", "n_out", "n_flv", "n_hel", "n_col", "n_cin", "n_cf")
       write (unit, "(2x,9A)")  "interface"
       write (unit, "(5x,9A)")  "function ", &
            char (writer%get_c_procname (id, feature)), &
            " () result (n) bind(C)"
       write (unit, "(7x,9A)")  "import"
       write (unit, "(7x,9A)")  "implicit none"
       write (unit, "(7x,9A)")  "integer(c_int) :: n"    
       write (unit, "(5x,9A)")  "end function ", &
            char (writer%get_c_procname (id, feature))
       write (unit, "(2x,9A)")  "end interface"
    case ("flv_state", "hel_state")
       write (unit, "(2x,9A)")  "interface"
       write (unit, "(5x,9A)")  "subroutine ", &
            char (writer%get_c_procname (id, feature)), &
            " (", char (feature), ") bind(C)"
       write (unit, "(7x,9A)")  "import"
       write (unit, "(7x,9A)")  "implicit none"
       write (unit, "(7x,9A)")  "integer(c_int), dimension(*), intent(out) ", &
            ":: ", char (feature)
       write (unit, "(5x,9A)")  "end subroutine ", &
            char (writer%get_c_procname (id, feature))
       write (unit, "(2x,9A)")  "end interface"
    case ("col_state")
       write (unit, "(2x,9A)")  "interface"
       write (unit, "(5x,9A)")  "subroutine ", &
            char (writer%get_c_procname (id, feature)), &
            " (col_state, ghost_flag) bind(C)"
       write (unit, "(7x,9A)")  "import"
       write (unit, "(7x,9A)")  "implicit none"
       write (unit, "(7x,9A)")  "integer(c_int), dimension(*), intent(out) ", &
            ":: col_state"
       write (unit, "(7x,9A)")  "logical(c_bool), dimension(*), intent(out) ", &
            ":: ghost_flag"
       write (unit, "(5x,9A)")  "end subroutine ", &
            char (writer%get_c_procname (id, feature))
       write (unit, "(2x,9A)")  "end interface"
    case ("color_factors")
       write (unit, "(2x,9A)")  "interface"
       write (unit, "(5x,9A)")  "subroutine ", &
            char (writer%get_c_procname (id, feature)), &
            " (cf_index1, cf_index2, color_factors) bind(C)"
       write (unit, "(7x,9A)")  "import"
       write (unit, "(7x,9A)")  "implicit none"
       write (unit, "(7x,9A)")  "integer(c_int), dimension(*), &
            &intent(out) :: cf_index1"
       write (unit, "(7x,9A)")  "integer(c_int), dimension(*), &
            &intent(out) :: cf_index2"
       write (unit, "(7x,9A)")  "complex(c_default_complex), dimension(*), &
            &intent(out) :: color_factors"
       write (unit, "(5x,9A)")  "end subroutine ", &
            char (writer%get_c_procname (id, feature))
       write (unit, "(2x,9A)")  "end interface"
    end select
  end subroutine prc_writer_c_lib_write_interface

\eatline
\nwendcode{}\nwbegindocs{245}\nwdocspar
\subsection{Retrieving the tables}
In the previous section we had the writer routines for procedures that
return tables, actually C pointers to tables.  Here, we write
convenience routines that unpack them and move the contents to
suitable Fortran arrays.

The flavor and helicity tables are two-dimensional integer arrays.  We
use intermediate storage for correctly transforming C to Fortran data
types.
\nwenddocs{}\nwbegincode{246}\moddef{Prclib interfaces: prclib driver: TBP}\plusendmoddef
  procedure :: set_flv_state => prclib_driver_set_flv_state
  procedure :: set_hel_state => prclib_driver_set_hel_state
\nwendcode{}\nwbegincode{247}\moddef{Prclib interfaces: procedures}\plusendmoddef
  subroutine prclib_driver_set_flv_state (driver, i, flv_state)
    class(prclib_driver_t), intent(in) :: driver
    integer, intent(in) :: i
    integer, dimension(:,:), allocatable, intent(out) :: flv_state
    integer :: n_tot, n_flv
    integer(c_int) :: pid
    integer(c_int), dimension(:,:), allocatable :: c_flv_state
    pid = i
    n_tot = driver%get_n_in (pid) + driver%get_n_out (pid)
    n_flv = driver%get_n_flv (pid)
    allocate (flv_state (n_tot, n_flv))
    allocate (c_flv_state (n_tot, n_flv))
    call driver%set_flv_state_ptr &
         (pid, c_flv_state, int ([n_tot, n_flv], kind=c_int))
    flv_state = c_flv_state
  end subroutine prclib_driver_set_flv_state
    
  subroutine prclib_driver_set_hel_state (driver, i, hel_state)
    class(prclib_driver_t), intent(in) :: driver
    integer, intent(in) :: i
    integer, dimension(:,:), allocatable, intent(out) :: hel_state
    integer :: n_tot, n_hel
    integer(c_int) :: pid
    integer(c_int), dimension(:,:), allocatable, target :: c_hel_state
    pid = i
    n_tot = driver%get_n_in (pid) + driver%get_n_out (pid)
    n_hel = driver%get_n_hel (pid)
    allocate (hel_state (n_tot, n_hel))
    allocate (c_hel_state (n_tot, n_hel))
    call driver%set_hel_state_ptr &
         (pid, c_hel_state, int ([n_tot, n_hel], kind=c_int))
    hel_state = c_hel_state
  end subroutine prclib_driver_set_hel_state
    
\eatline
\eatline
\nwendcode{}\nwbegindocs{248}The color-flow table is three-dimensional, otherwise similar.  We
simultaneously set the ghost-flag table, which consists of logical
entries.
\nwenddocs{}\nwbegincode{249}\moddef{Prclib interfaces: prclib driver: TBP}\plusendmoddef
  procedure :: set_col_state => prclib_driver_set_col_state
\nwendcode{}\nwbegincode{250}\moddef{Prclib interfaces: procedures}\plusendmoddef
  subroutine prclib_driver_set_col_state (driver, i, col_state, ghost_flag)
    class(prclib_driver_t), intent(in) :: driver
    integer, intent(in) :: i
    integer, dimension(:,:,:), allocatable, intent(out) :: col_state
    logical, dimension(:,:), allocatable, intent(out) :: ghost_flag
    integer :: n_cin, n_tot, n_col
    integer(c_int) :: pid
    integer(c_int), dimension(:,:,:), allocatable :: c_col_state
    logical(c_bool), dimension(:,:), allocatable :: c_ghost_flag
    pid = i
    n_cin = driver%get_n_cin (pid)
    n_tot = driver%get_n_in (pid) + driver%get_n_out (pid)
    n_col = driver%get_n_col (pid)
    allocate (col_state (n_cin, n_tot, n_col))
    allocate (c_col_state (n_cin, n_tot, n_col))
    allocate (ghost_flag (n_tot, n_col))
    allocate (c_ghost_flag (n_tot, n_col))
    call driver%set_col_state_ptr (pid, &
         c_col_state, c_ghost_flag, int ([n_cin, n_tot, n_col], kind=c_int))
    col_state = c_col_state
    ghost_flag = c_ghost_flag
  end subroutine prclib_driver_set_col_state
    
\eatline
\nwendcode{}\nwbegindocs{251}The color-factor table is a sparse matrix: a two-column array of indices and
one array which contains the corresponding factors.
\nwenddocs{}\nwbegincode{252}\moddef{Prclib interfaces: prclib driver: TBP}\plusendmoddef
  procedure :: set_color_factors => prclib_driver_set_color_factors
\nwendcode{}\nwbegincode{253}\moddef{Prclib interfaces: procedures}\plusendmoddef
  subroutine prclib_driver_set_color_factors (driver, i, color_factors, cf_index)
    class(prclib_driver_t), intent(in) :: driver
    integer, intent(in) :: i
    complex(default), dimension(:), allocatable, intent(out) :: color_factors
    integer, dimension(:,:), allocatable, intent(out) :: cf_index
    integer :: n_cf
    integer(c_int) :: pid
    complex(c_default_complex), dimension(:), allocatable, target :: c_color_factors
    integer(c_int), dimension(:), allocatable, target :: c_cf_index1
    integer(c_int), dimension(:), allocatable, target :: c_cf_index2
    pid = i
    n_cf = driver%get_n_cf (pid)
    allocate (color_factors (n_cf))
    allocate (c_color_factors (n_cf))
    allocate (c_cf_index1 (n_cf))
    allocate (c_cf_index2 (n_cf))
    call driver%set_color_factors_ptr (pid, &
         c_cf_index1, c_cf_index2, &
         c_color_factors, int ([n_cf], kind=c_int))
    color_factors = c_color_factors
    allocate (cf_index (2, n_cf))
    cf_index(1,:) = c_cf_index1
    cf_index(2,:) = c_cf_index2
  end subroutine prclib_driver_set_color_factors
    
\eatline
\nwendcode{}\nwbegindocs{254}\nwdocspar
\subsection{Returning a procedure pointer}
The functions that directly access the matrix element, event by event,
are assigned to a process-specific driver object as procedure
pointers.  For the {\tt{}dlopen} interface, we use C function pointers.
This subroutine returns such a pointer:
\nwenddocs{}\nwbegincode{255}\moddef{Prclib interfaces: public}\plusendmoddef
  public :: prc_get_fptr
\nwendcode{}\nwbegincode{256}\moddef{Prclib interfaces: interfaces}\plusendmoddef
  abstract interface
     subroutine prc_get_fptr (pid, fid, fptr) bind(C)
       import
       integer(c_int), intent(in) :: pid
       integer(c_int), intent(in) :: fid
       type(c_funptr), intent(out) :: fptr
     end subroutine prc_get_fptr
  end interface
  
\eatline
\nwendcode{}\nwbegindocs{257}This procedure writes the source code for the procedure pointer
returning subroutine.

All C functions that are provided by the matrix element code of a
specific process are handled here.  The selection consists of a double
layered {\tt{}select} {\tt{}case} construct. 
\nwenddocs{}\nwbegincode{258}\moddef{Prclib interfaces: prclib driver: TBP}\plusendmoddef
  procedure :: write_get_fptr_sub
\nwendcode{}\nwbegincode{259}\moddef{Prclib interfaces: procedures}\plusendmoddef
  subroutine write_get_fptr_sub (driver, unit, prefix)
    class(prclib_driver_t), intent(in) :: driver
    integer, intent(in) :: unit
    type(string_t), intent(in) :: prefix
    integer :: i, j
    write (unit, "(A)")  ""
    write (unit, "(A)")  "! Return C pointer to a procedure:"
    write (unit, "(A)")  "! pid = process index;  fid = function index"
    write (unit, "(4A)")  "subroutine ", char (prefix), "get_fptr ", &
         "(pid, fid, fptr) bind(C)"
    write (unit, "(A)")  "  use iso_c_binding"
    write (unit, "(A)")  "  use kinds"
    write (unit, "(A)")  "  implicit none"
    write (unit, "(A)")  "  integer(c_int), intent(in) :: pid"
    write (unit, "(A)")  "  integer(c_int), intent(in) :: fid"
    write (unit, "(A)")  "  type(c_funptr), intent(out) :: fptr"
    do i = 1, driver%n_processes
       call driver%record(i)%write_interfaces (unit)
    end do
    write (unit, "(A)")  "  select case (pid)"
    do i = 1, driver%n_processes
       write (unit, "(2x,A,I0,A)")  "case (", i, ")"
       write (unit, "(5x,A)")  "select case (fid)"
       associate (record => driver%record(i))
         do j = 1, size (record%feature)
            write (unit, "(5x,A,I0,9A)")  "case (", j, ");  ", &
                 "fptr = c_funloc (", &
                 char (record%get_c_procname (record%feature(j))), &
                 ")"
         end do
       end associate
       write (unit, "(5x,A)")  "end select"
    end do
    write (unit, "(A)")  "  end select"
    write (unit, "(3A)")  "end subroutine ", char (prefix), "get_fptr"
  end subroutine write_get_fptr_sub

\eatline
\nwendcode{}\nwbegindocs{260}The procedures for which we want to return a pointer (the 'features'
of the matrix element code) are actually Fortran module procedures.
If we want to have a C signature, we must write wrapper functions for
all of them.  The procedures, their signatures, and the appropriate
writer routines are specific for the process type.

To keep this generic, we do not provide the writer routines here, but
just the interface for a writer routine.  The actual routines are
stored in the process record.

The {\tt{}prefix} indicates the library, the {\tt{}id} indicates the
process, and {\tt{}procname} is the bare name of the procedure to be
written.
\nwenddocs{}\nwbegincode{261}\moddef{Prclib interfaces: public}\plusendmoddef
  public :: write_driver_code
\nwendcode{}\nwbegincode{262}\moddef{Prclib interfaces: interfaces}\plusendmoddef
  abstract interface
     subroutine write_driver_code (unit, prefix, id, procname)
       import
       integer, intent(in) :: unit
       type(string_t), intent(in) :: prefix
       type(string_t), intent(in) :: id
       type(string_t), intent(in) :: procname
     end subroutine write_driver_code
  end interface

\eatline
\nwendcode{}\nwbegindocs{263}\nwdocspar
\subsection{Hooks}
Interface for additional library unload / reload hooks (currently unused!)
\nwenddocs{}\nwbegincode{264}\moddef{Prclib interfaces: public}\plusendmoddef
  public :: prclib_unload_hook
  public :: prclib_reload_hook
\nwendcode{}\nwbegincode{265}\moddef{Prclib interfaces: interfaces}\plusendmoddef
  abstract interface
     subroutine prclib_unload_hook (libname)
       import
       type(string_t), intent(in) :: libname
     end subroutine prclib_unload_hook

     subroutine prclib_reload_hook (libname)
       import
       type(string_t), intent(in) :: libname
     end subroutine prclib_reload_hook
  end interface
\eatline
\eatline
\nwendcode{}\nwbegindocs{266}\nwdocspar
\subsection{Make source, compile, link}
Since we should have written a Makefile, these tasks amount to simple
{\tt{}make} calls.  Note that the Makefile targets depend on each other,
so calling {\tt{}link} executes also the {\tt{}source} and {\tt{}compile}
steps, when necessary.

The first routine writes source-code files for the individual
processes.  First it calls the writer routines directly for each
process, then it calls {\tt{}make\ source}.  The make command may either
post-process the files, or it may do the complete work, e.g., calling
an external program the generates the files.
\nwenddocs{}\nwbegincode{267}\moddef{Prclib interfaces: prclib driver: TBP}\plusendmoddef
  procedure :: make_source => prclib_driver_make_source
\nwendcode{}\nwbegincode{268}\moddef{Prclib interfaces: procedures}\plusendmoddef
  subroutine prclib_driver_make_source (driver, os_data)
    class(prclib_driver_t), intent(in) :: driver
    type(os_data_t), intent(in) :: os_data
    integer :: i
    do i = 1, driver%n_processes
       call driver%record(i)%write_source_code ()
    end do
    call os_system_call ("make source " // os_data%makeflags &
         // " -f " // driver%basename // ".makefile")
  end subroutine prclib_driver_make_source
  
\eatline
\nwendcode{}\nwbegindocs{269}Compile matrix element source code and the driver source code.
\nwenddocs{}\nwbegincode{270}\moddef{Prclib interfaces: prclib driver: TBP}\plusendmoddef
  procedure :: make_compile => prclib_driver_make_compile
\nwendcode{}\nwbegincode{271}\moddef{Prclib interfaces: procedures}\plusendmoddef
  subroutine prclib_driver_make_compile (driver, os_data)
    class(prclib_driver_t), intent(in) :: driver
    type(os_data_t), intent(in) :: os_data
    call os_system_call ("make compile " // os_data%makeflags &
         // " -f " // driver%basename // ".makefile")
  end subroutine prclib_driver_make_compile
  
\eatline
\nwendcode{}\nwbegindocs{272}Combine all matrix-element code together with the driver in a
process library on disk.
\nwenddocs{}\nwbegincode{273}\moddef{Prclib interfaces: prclib driver: TBP}\plusendmoddef
  procedure :: make_link => prclib_driver_make_link
\nwendcode{}\nwbegincode{274}\moddef{Prclib interfaces: procedures}\plusendmoddef
  subroutine prclib_driver_make_link (driver, os_data)
    class(prclib_driver_t), intent(in) :: driver
    type(os_data_t), intent(in) :: os_data
    call os_system_call ("make link " // os_data%makeflags &
         // " -f " // driver%basename // ".makefile")
  end subroutine prclib_driver_make_link
  
\eatline
\nwendcode{}\nwbegindocs{275}\nwdocspar
\subsection{Clean up generated files}
The task of cleaning any generated files should also be deferred to
Makefile targets.  Apart from removing everything, removing specific
files may be useful for partial rebuilds.  (Note that removing the
makefile itself can only be done once, for obvious reasons.)

If there is no makefile, do nothing.
\nwenddocs{}\nwbegincode{276}\moddef{Prclib interfaces: prclib driver: TBP}\plusendmoddef
  procedure :: clean_library => prclib_driver_clean_library
  procedure :: clean_objects => prclib_driver_clean_objects
  procedure :: clean_source => prclib_driver_clean_source
  procedure :: clean_driver => prclib_driver_clean_driver
  procedure :: clean_makefile => prclib_driver_clean_makefile
  procedure :: clean => prclib_driver_clean
  procedure :: distclean => prclib_driver_distclean
\nwendcode{}\nwbegincode{277}\moddef{Prclib interfaces: procedures}\plusendmoddef
  subroutine prclib_driver_clean_library (driver, os_data)
    class(prclib_driver_t), intent(in) :: driver
    type(os_data_t), intent(in) :: os_data
    if (driver%makefile_exists ()) then
       call os_system_call ("make clean-library " // os_data%makeflags &
            // " -f " // driver%basename // ".makefile")
    end if
  end subroutine prclib_driver_clean_library
  
  subroutine prclib_driver_clean_objects (driver, os_data)
    class(prclib_driver_t), intent(in) :: driver
    type(os_data_t), intent(in) :: os_data
    if (driver%makefile_exists ()) then
       call os_system_call ("make clean-objects " // os_data%makeflags &
            // " -f " // driver%basename // ".makefile")
    end if
  end subroutine prclib_driver_clean_objects
  
  subroutine prclib_driver_clean_source (driver, os_data)
    class(prclib_driver_t), intent(in) :: driver
    type(os_data_t), intent(in) :: os_data
    if (driver%makefile_exists ()) then
       call os_system_call ("make clean-source " // os_data%makeflags &
            // " -f " // driver%basename // ".makefile")
    end if
  end subroutine prclib_driver_clean_source
  
  subroutine prclib_driver_clean_driver (driver, os_data)
    class(prclib_driver_t), intent(in) :: driver
    type(os_data_t), intent(in) :: os_data
    if (driver%makefile_exists ()) then
       call os_system_call ("make clean-driver " // os_data%makeflags &
            // " -f " // driver%basename // ".makefile")
    end if
  end subroutine prclib_driver_clean_driver
  
  subroutine prclib_driver_clean_makefile (driver, os_data)
    class(prclib_driver_t), intent(in) :: driver
    type(os_data_t), intent(in) :: os_data
    if (driver%makefile_exists ()) then
       call os_system_call ("make clean-makefile " // os_data%makeflags &
            // " -f " // driver%basename // ".makefile")
    end if
  end subroutine prclib_driver_clean_makefile
  
  subroutine prclib_driver_clean (driver, os_data)
    class(prclib_driver_t), intent(in) :: driver
    type(os_data_t), intent(in) :: os_data
    if (driver%makefile_exists ()) then
       call os_system_call ("make clean " // os_data%makeflags &
            // " -f " // driver%basename // ".makefile")
    end if
  end subroutine prclib_driver_clean
  
  subroutine prclib_driver_distclean (driver, os_data)
    class(prclib_driver_t), intent(in) :: driver
    type(os_data_t), intent(in) :: os_data
    if (driver%makefile_exists ()) then
       call os_system_call ("make distclean " // os_data%makeflags &
            // " -f " // driver%basename // ".makefile")
    end if
  end subroutine prclib_driver_distclean
  
\eatline
\eatline
\eatline
\eatline
\eatline
\eatline
\eatline
\nwendcode{}\nwbegindocs{278}This Make target should remove all files that apply to a specific process.
We execute this when we want to force remaking source code.  Note that source
targets need not have prerequisites, so just calling {\tt{}make{\char95}source} would not
do anything if the files exist.
\nwenddocs{}\nwbegincode{279}\moddef{Prclib interfaces: prclib driver: TBP}\plusendmoddef
  procedure :: clean_proc => prclib_driver_clean_proc
\nwendcode{}\nwbegincode{280}\moddef{Prclib interfaces: procedures}\plusendmoddef
  subroutine prclib_driver_clean_proc (driver, i, os_data)
    class(prclib_driver_t), intent(in) :: driver
    integer, intent(in) :: i
    type(os_data_t), intent(in) :: os_data
    type(string_t) :: id
    if (driver%makefile_exists ()) then
       id = driver%record(i)%id
      call os_system_call ("make clean-" // driver%record(i)%id // " " &
           // os_data%makeflags &
           // " -f " // driver%basename // ".makefile")
    end if
  end subroutine prclib_driver_clean_proc
  
\eatline
\nwendcode{}\nwbegindocs{281}\nwdocspar
\subsection{Further Tools}
Check for the appropriate makefile.
\nwenddocs{}\nwbegincode{282}\moddef{Prclib interfaces: prclib driver: TBP}\plusendmoddef
  procedure :: makefile_exists => prclib_driver_makefile_exists
\nwendcode{}\nwbegincode{283}\moddef{Prclib interfaces: procedures}\plusendmoddef
  function prclib_driver_makefile_exists (driver) result (flag)
    class(prclib_driver_t), intent(in) :: driver
    logical :: flag
    inquire (file = char (driver%basename) // ".makefile", exist = flag)
  end function prclib_driver_makefile_exists
  
\eatline
\nwendcode{}\nwbegindocs{284}\nwdocspar
\subsection{Load the library}
Once the library has been linked, we can dlopen it and assign all
procedure pointers to their proper places in the library driver
object.  The {\tt{}loaded} flag is set only if all required pointers
have become assigned.
\nwenddocs{}\nwbegincode{285}\moddef{Prclib interfaces: prclib driver: TBP}\plusendmoddef
  procedure :: load => prclib_driver_load
\nwendcode{}\nwbegincode{286}\moddef{Prclib interfaces: procedures}\plusendmoddef
  subroutine prclib_driver_load (driver, os_data, noerror)
    class(prclib_driver_t), intent(inout) :: driver
    type(os_data_t), intent(in) :: os_data
    logical, intent(in), optional :: noerror
    type(c_funptr) :: c_fptr
    logical :: ignore
    
    ignore = .false.;  if (present (noerror))  ignore = noerror
    
    driver%libname = os_get_dlname (driver%basename, os_data, noerror, noerror)
    if (driver%libname == "")  return
    select type (driver)
    type is (prclib_driver_dynamic_t)
       if (.not. dlaccess_is_open (driver%dlaccess)) then
          call dlaccess_init &
               (driver%dlaccess, var_str ("."), driver%libname, os_data)
          if (.not. ignore)  call driver%check_dlerror ()
       end if
       driver%loaded = dlaccess_is_open (driver%dlaccess)
    class default
       driver%loaded = .true.
    end select
    if (.not. driver%loaded)  return

    c_fptr = driver%get_c_funptr (var_str ("get_n_processes"))
    call c_f_procpointer (c_fptr, driver%get_n_processes)
    driver%loaded = driver%loaded .and. associated (driver%get_n_processes)

    c_fptr = driver%get_c_funptr (var_str ("get_process_id_ptr"))
    call c_f_procpointer (c_fptr, driver%get_process_id_ptr)
    driver%loaded = driver%loaded .and. associated (driver%get_process_id_ptr)

    c_fptr = driver%get_c_funptr (var_str ("get_model_name_ptr"))
    call c_f_procpointer (c_fptr, driver%get_model_name_ptr)
    driver%loaded = driver%loaded .and. associated (driver%get_model_name_ptr)

    c_fptr = driver%get_c_funptr (var_str ("get_md5sum_ptr"))
    call c_f_procpointer (c_fptr, driver%get_md5sum_ptr)
    driver%loaded = driver%loaded .and. associated (driver%get_md5sum_ptr)

    c_fptr = driver%get_c_funptr (var_str ("get_openmp_status"))
    call c_f_procpointer (c_fptr, driver%get_openmp_status)
    driver%loaded = driver%loaded .and. associated (driver%get_openmp_status)

    c_fptr = driver%get_c_funptr (var_str ("get_n_in"))
    call c_f_procpointer (c_fptr, driver%get_n_in)
    driver%loaded = driver%loaded .and. associated (driver%get_n_in)

    c_fptr = driver%get_c_funptr (var_str ("get_n_out"))
    call c_f_procpointer (c_fptr, driver%get_n_out)
    driver%loaded = driver%loaded .and. associated (driver%get_n_out)

    c_fptr = driver%get_c_funptr (var_str ("get_n_flv"))
    call c_f_procpointer (c_fptr, driver%get_n_flv)
    driver%loaded = driver%loaded .and. associated (driver%get_n_flv)

    c_fptr = driver%get_c_funptr (var_str ("get_n_hel"))
    call c_f_procpointer (c_fptr, driver%get_n_hel)
    driver%loaded = driver%loaded .and. associated (driver%get_n_hel)

    c_fptr = driver%get_c_funptr (var_str ("get_n_col"))
    call c_f_procpointer (c_fptr, driver%get_n_col)
    driver%loaded = driver%loaded .and. associated (driver%get_n_col)

    c_fptr = driver%get_c_funptr (var_str ("get_n_cin"))
    call c_f_procpointer (c_fptr, driver%get_n_cin)
    driver%loaded = driver%loaded .and. associated (driver%get_n_cin)

    c_fptr = driver%get_c_funptr (var_str ("get_n_cf"))
    call c_f_procpointer (c_fptr, driver%get_n_cf)
    driver%loaded = driver%loaded .and. associated (driver%get_n_cf)

    c_fptr = driver%get_c_funptr (var_str ("set_flv_state_ptr"))
    call c_f_procpointer (c_fptr, driver%set_flv_state_ptr)
    driver%loaded = driver%loaded .and. associated (driver%set_flv_state_ptr)

    c_fptr = driver%get_c_funptr (var_str ("set_hel_state_ptr"))
    call c_f_procpointer (c_fptr, driver%set_hel_state_ptr)
    driver%loaded = driver%loaded .and. associated (driver%set_hel_state_ptr)

    c_fptr = driver%get_c_funptr (var_str ("set_col_state_ptr"))
    call c_f_procpointer (c_fptr, driver%set_col_state_ptr)
    driver%loaded = driver%loaded .and. associated (driver%set_col_state_ptr)

    c_fptr = driver%get_c_funptr (var_str ("set_color_factors_ptr"))
    call c_f_procpointer (c_fptr, driver%set_color_factors_ptr)
    driver%loaded = driver%loaded .and. associated (driver%set_color_factors_ptr)

    c_fptr = driver%get_c_funptr (var_str ("get_fptr"))
    call c_f_procpointer (c_fptr, driver%get_fptr)
    driver%loaded = driver%loaded .and. associated (driver%get_fptr)

  end subroutine prclib_driver_load
  
\eatline
\nwendcode{}\nwbegindocs{287}Unload.  To be sure, nullify the procedure pointers.
\nwenddocs{}\nwbegincode{288}\moddef{Prclib interfaces: prclib driver: TBP}\plusendmoddef
  procedure :: unload => prclib_driver_unload
\nwendcode{}\nwbegincode{289}\moddef{Prclib interfaces: procedures}\plusendmoddef
  subroutine prclib_driver_unload (driver)
    class(prclib_driver_t), intent(inout) :: driver
    select type (driver)
    type is (prclib_driver_dynamic_t)
       if (dlaccess_is_open (driver%dlaccess)) then
          call dlaccess_final (driver%dlaccess)
          call driver%check_dlerror ()
       end if
    end select
    driver%loaded = .false.
    nullify (driver%get_n_processes)
    nullify (driver%get_process_id_ptr)
    nullify (driver%get_model_name_ptr)
    nullify (driver%get_md5sum_ptr)
    nullify (driver%get_openmp_status)
    nullify (driver%get_n_in)
    nullify (driver%get_n_out)
    nullify (driver%get_n_flv)
    nullify (driver%get_n_hel)
    nullify (driver%get_n_col)
    nullify (driver%get_n_cin)
    nullify (driver%get_n_cf)
    nullify (driver%set_flv_state_ptr)
    nullify (driver%set_hel_state_ptr)
    nullify (driver%set_col_state_ptr)
    nullify (driver%set_color_factors_ptr)
    nullify (driver%get_fptr)
  end subroutine prclib_driver_unload
    
\eatline
\nwendcode{}\nwbegindocs{290}This subroutine checks the {\tt{}dlerror} content and issues a fatal
error if it finds an error there.
\nwenddocs{}\nwbegincode{291}\moddef{Prclib interfaces: prclib driver dynamic: TBP}\endmoddef
  procedure :: check_dlerror => prclib_driver_check_dlerror
\nwendcode{}\nwbegincode{292}\moddef{Prclib interfaces: procedures}\plusendmoddef
  subroutine prclib_driver_check_dlerror (driver)
    class(prclib_driver_dynamic_t), intent(in) :: driver
    if (dlaccess_has_error (driver%dlaccess)) then
       call msg_fatal (char (dlaccess_get_error (driver%dlaccess)))
    end if
  end subroutine prclib_driver_check_dlerror
  
\eatline
\nwendcode{}\nwbegindocs{293}Get the handle (C function pointer) for a given ``feature'' of the
matrix element code, so it can be assigned to the appropriate
procedure pointer slot.  In the static case, this is a
trivial pointer assignment, hard-coded into the driver type
implementation.
\nwenddocs{}\nwbegincode{294}\moddef{Prclib interfaces: prclib driver: TBP}\plusendmoddef
  procedure (prclib_driver_get_c_funptr), deferred :: get_c_funptr
\nwendcode{}\nwbegincode{295}\moddef{Prclib interfaces: interfaces}\plusendmoddef
  abstract interface
     function prclib_driver_get_c_funptr (driver, feature) result (c_fptr)
       import
       class(prclib_driver_t), intent(inout) :: driver
       type(string_t), intent(in) :: feature
       type(c_funptr) :: c_fptr
     end function prclib_driver_get_c_funptr
  end interface

\eatline
\nwendcode{}\nwbegindocs{296}In the dynamic-library case, we call the DL interface to retrieve the C
pointer to a named procedure.  
\nwenddocs{}\nwbegincode{297}\moddef{Prclib interfaces: prclib driver dynamic: TBP}\plusendmoddef
  procedure :: get_c_funptr => prclib_driver_dynamic_get_c_funptr
\nwendcode{}\nwbegincode{298}\moddef{Prclib interfaces: procedures}\plusendmoddef
  function prclib_driver_dynamic_get_c_funptr (driver, feature) result (c_fptr)
    class(prclib_driver_dynamic_t), intent(inout) :: driver
    type(string_t), intent(in) :: feature
    type(c_funptr) :: c_fptr
    type(string_t) :: prefix, full_name
    prefix = driver%basename // "_"
    full_name = prefix // feature
    c_fptr = dlaccess_get_c_funptr (driver%dlaccess, full_name)
    call driver%check_dlerror ()
  end function prclib_driver_dynamic_get_c_funptr

\eatline
\nwendcode{}\nwbegindocs{299}\nwdocspar
\subsection{MD5 sums}
Recall the MD5 sum written in the Makefile
\nwenddocs{}\nwbegincode{300}\moddef{Prclib interfaces: prclib driver: TBP}\plusendmoddef
  procedure :: get_md5sum_makefile => prclib_driver_get_md5sum_makefile
\nwendcode{}\nwbegincode{301}\moddef{Prclib interfaces: procedures}\plusendmoddef
  function prclib_driver_get_md5sum_makefile (driver) result (md5sum)
    class(prclib_driver_t), intent(in) :: driver
    character(32) :: md5sum
    type(string_t) :: filename
    character(80) :: buffer
    logical :: exist
    integer :: u, iostat
    md5sum = ""
    filename = driver%basename // ".makefile"
    inquire (file = char (filename), exist = exist)
    if (exist) then
       u = free_unit ()
       open (u, file = char (filename), action = "read", status = "old")
       iostat = 0
       do
          read (u, "(A)", iostat = iostat)  buffer
          if (iostat /= 0)  exit
          buffer = adjustl (buffer)
          select case (buffer(1:9))
          case ("MD5SUM = ")
             read (buffer(11:), "(A32)")  md5sum
             exit
          end select
       end do
       close (u)
    end if
  end function prclib_driver_get_md5sum_makefile
    
\eatline
\nwendcode{}\nwbegindocs{302}Recall the MD5 sum written in the driver source code.
\nwenddocs{}\nwbegincode{303}\moddef{Prclib interfaces: prclib driver: TBP}\plusendmoddef
  procedure :: get_md5sum_driver => prclib_driver_get_md5sum_driver
\nwendcode{}\nwbegincode{304}\moddef{Prclib interfaces: procedures}\plusendmoddef
  function prclib_driver_get_md5sum_driver (driver) result (md5sum)
    class(prclib_driver_t), intent(in) :: driver
    character(32) :: md5sum
    type(string_t) :: filename
    character(80) :: buffer
    logical :: exist
    integer :: u, iostat
    md5sum = ""
    filename = driver%basename // ".f90"
    inquire (file = char (filename), exist = exist)
    if (exist) then
       u = free_unit ()
       open (u, file = char (filename), action = "read", status = "old")
       iostat = 0
       do
          read (u, "(A)", iostat = iostat)  buffer
          if (iostat /= 0)  exit
          buffer = adjustl (buffer)
          select case (buffer(1:9))
          case ("md5sum = ")
             read (buffer(11:), "(A32)")  md5sum
             exit
          end select
       end do
       close (u)
    end if
  end function prclib_driver_get_md5sum_driver
    
\eatline
\nwendcode{}\nwbegindocs{305}Recall the MD5 sum written in the matrix element source code.
\nwenddocs{}\nwbegincode{306}\moddef{Prclib interfaces: prclib driver: TBP}\plusendmoddef
  procedure :: get_md5sum_source => prclib_driver_get_md5sum_source
\nwendcode{}\nwbegincode{307}\moddef{Prclib interfaces: procedures}\plusendmoddef
  function prclib_driver_get_md5sum_source (driver, i) result (md5sum)
    class(prclib_driver_t), intent(in) :: driver
    integer, intent(in) :: i
    character(32) :: md5sum
    type(string_t) :: filename
    character(80) :: buffer
    logical :: exist
    integer :: u, iostat
    md5sum = ""
    
    filename = driver%record(i)%id // ".f90"
    inquire (file = char (filename), exist = exist)
    if (exist) then
       u = free_unit ()
       open (u, file = char (filename), action = "read", status = "old")
       iostat = 0
       do
          read (u, "(A)", iostat = iostat)  buffer
          if (iostat /= 0)  exit
          buffer = adjustl (buffer)
          select case (buffer(1:9))
          case ("md5sum = ")
             read (buffer(11:), "(A32)")  md5sum
             exit
          end select
       end do
       close (u)
    end if
  end function prclib_driver_get_md5sum_source
    
\eatline
\nwendcode{}\nwbegindocs{308}\nwdocspar
\subsection{Test}
This is the master for calling self-test procedures.
\nwenddocs{}\nwbegincode{309}\moddef{Prclib interfaces: public}\plusendmoddef
  public :: prclib_interfaces_test
\nwendcode{}\nwbegincode{310}\moddef{Prclib interfaces: tests}\endmoddef
  subroutine prclib_interfaces_test (u, results)
    integer, intent(in) :: u
    type(test_results_t), intent(inout) :: results
  \LA{}Prclib interfaces: execute tests\RA{}
  end subroutine prclib_interfaces_test

\eatline
\nwendcode{}\nwbegindocs{311}\nwdocspar
\subsubsection{Empty process list}
Test 1: Create a driver object and display its contents.  One of the
feature lists references a writer procedure; this is just a dummy that
does nothing useful.
\nwenddocs{}\nwbegincode{312}\moddef{Prclib interfaces: execute tests}\endmoddef
  call test (prclib_interfaces_1, "prclib_interfaces_1", &
       "create driver object", &
       u, results)
\nwendcode{}\nwbegincode{313}\moddef{Prclib interfaces: tests}\plusendmoddef
  subroutine prclib_interfaces_1 (u)
    integer, intent(in) :: u
    class(prclib_driver_t), allocatable :: driver
    character(32), parameter :: md5sum = "prclib_interfaces_1_md5sum      "
    class(prc_writer_t), pointer :: test_writer_1
  

    write (u, "(A)")  "* Test output: prclib_interfaces_1"
    write (u, "(A)")  "*   Purpose: display the driver object contents"
    write (u, *)
    write (u, "(A)")  "* Create a prclib driver object"
    write (u, "(A)")

    call dispatch_prclib_driver (driver, var_str ("prclib"), var_str (""))
    call driver%init (3)
    call driver%set_md5sum (md5sum)

    allocate (test_writer_1_t :: test_writer_1)

    call driver%set_record (1, var_str ("test1"), var_str ("test_model"), &
         [var_str ("init")], test_writer_1)

    call driver%set_record (2, var_str ("test2"), var_str ("foo_model"), &
         [var_str ("another_proc")], test_writer_1)

    call driver%set_record (3, var_str ("test3"), var_str ("test_model"), &
         [var_str ("init"), var_str ("some_proc")], test_writer_1)

    call driver%write (u)
    
    deallocate (test_writer_1)

    write (u, "(A)")
    write (u, "(A)")  "* Test output end: prclib_interfaces_1"
  end subroutine prclib_interfaces_1
  
\eatline
\nwendcode{}\nwbegindocs{314}The writer: the procedures write just comment lines.  We can fix an
instance of this as a parameter (since it has no mutable content) and
just reference the fixed parameter.

NOTE: temporarily made public.
\nwenddocs{}\nwbegincode{315}\moddef{Prclib interfaces: test types}\endmoddef
  type, extends (prc_writer_t) :: test_writer_1_t
   contains
     procedure, nopass :: type_name => test_writer_1_type_name
     procedure :: write_makefile_code => test_writer_1_mk
     procedure :: write_source_code => test_writer_1_src
     procedure :: write_interface => test_writer_1_if
     procedure :: write_md5sum_call => test_writer_1_md5sum
     procedure :: write_int_sub_call => test_writer_1_int_sub
     procedure :: write_col_state_call => test_writer_1_col_state
     procedure :: write_color_factors_call => test_writer_1_col_factors
  end type test_writer_1_t
  
\eatline
\nwendcode{}\nwbegindocs{316}\nwdocspar
\nwenddocs{}\nwbegincode{317}\moddef{Prclib interfaces: tests}\plusendmoddef
  function test_writer_1_type_name () result (string)
    type(string_t) :: string
    string = "test_1"
  end function test_writer_1_type_name

  subroutine test_writer_1_mk (writer, unit, id, os_data, testflag)
    class(test_writer_1_t), intent(in) :: writer
    integer, intent(in) :: unit
    type(string_t), intent(in) :: id
    type(os_data_t), intent(in) :: os_data
    logical, intent(in), optional :: testflag
    write (unit, "(5A)")  "# Makefile code for process ", char (id), &
         " goes here."
  end subroutine test_writer_1_mk
  
  subroutine test_writer_1_src (writer, id)
    class(test_writer_1_t), intent(in) :: writer
    type(string_t), intent(in) :: id
  end subroutine test_writer_1_src
  
  subroutine test_writer_1_if (writer, unit, id, feature)
    class(test_writer_1_t), intent(in) :: writer
    integer, intent(in) :: unit
    type(string_t), intent(in) :: id, feature
    write (unit, "(2x,9A)")  "! Interface code for ", &
       char (id), "_", char (writer%get_procname (feature)), &
       " goes here."
  end subroutine test_writer_1_if

  subroutine test_writer_1_md5sum (writer, unit, id)
    class(test_writer_1_t), intent(in) :: writer
    integer, intent(in) :: unit
    type(string_t), intent(in) :: id
    write (unit, "(5x,9A)")  "! MD5sum call for ", char (id), " goes here."
  end subroutine test_writer_1_md5sum

  subroutine test_writer_1_int_sub (writer, unit, id, feature)
    class(test_writer_1_t), intent(in) :: writer
    integer, intent(in) :: unit
    type(string_t), intent(in) :: id, feature
    write (unit, "(5x,9A)")  "! ", char (feature), " call for ", &
         char (id), " goes here."
  end subroutine test_writer_1_int_sub

  subroutine test_writer_1_col_state (writer, unit, id)
    class(test_writer_1_t), intent(in) :: writer
    integer, intent(in) :: unit
    type(string_t), intent(in) :: id
    write (unit, "(5x,9A)")  "! col_state call for ", &
         char (id), " goes here."
  end subroutine test_writer_1_col_state

  subroutine test_writer_1_col_factors (writer, unit, id)
    class(test_writer_1_t), intent(in) :: writer
    integer, intent(in) :: unit
    type(string_t), intent(in) :: id
    write (unit, "(5x,9A)")  "! color_factors call for ", &
         char (id), " goes here."
  end subroutine test_writer_1_col_factors

\eatline
\eatline
\eatline
\eatline
\nwendcode{}\nwbegindocs{318}\nwdocspar
\subsubsection{Process library driver file}
Test 2: Write the driver file for a test case with two processes.  The
first process needs no wrapper (C library), the second one needs
wrappers (Fortran module library).
\nwenddocs{}\nwbegincode{319}\moddef{Prclib interfaces: execute tests}\plusendmoddef
  call test (prclib_interfaces_2, "prclib_interfaces_2", &
       "write driver file", &
       u, results)
\nwendcode{}\nwbegincode{320}\moddef{Prclib interfaces: tests}\plusendmoddef
  subroutine prclib_interfaces_2 (u)
    integer, intent(in) :: u
    class(prclib_driver_t), allocatable :: driver
    character(32), parameter :: md5sum = "prclib_interfaces_2_md5sum      "
    class(prc_writer_t), pointer :: test_writer_1, test_writer_2

    write (u, "(A)")  "* Test output: prclib_interfaces_2"
    write (u, "(A)")  "*   Purpose: check the generated driver source code"
    write (u, "(A)")
    write (u, "(A)")  "* Create a prclib driver object (2 processes)"
    write (u, "(A)")

    call dispatch_prclib_driver (driver, var_str ("prclib2"), var_str (""))
    call driver%init (2)
    call driver%set_md5sum (md5sum)

    allocate (test_writer_1_t :: test_writer_1)
    allocate (test_writer_2_t :: test_writer_2)
    
    call driver%set_record (1, var_str ("test1"), var_str ("Test_model"), &
         [var_str ("proc1")], test_writer_1)

    call driver%set_record (2, var_str ("test2"), var_str ("Test_model"), &
         [var_str ("proc1"), var_str ("proc2")], test_writer_2)
    
    call driver%write (u)

    write (u, "(A)")
    write (u, "(A)")  "* Write the driver file"
    write (u, "(A)")  "* File contents:"
    write (u, "(A)")

    call driver%generate_driver_code (u)
    
    deallocate (test_writer_1)
    deallocate (test_writer_2)

    write (u, "(A)")
    write (u, "(A)")  "* Test output end: prclib_interfaces_2"
  end subroutine prclib_interfaces_2
  
\eatline
\nwendcode{}\nwbegindocs{321}A writer with wrapper code: the procedures again write just comment
lines.  Since all procedures are NOPASS, we can reuse two of the TBP.
\nwenddocs{}\nwbegincode{322}\moddef{Prclib interfaces: test types}\plusendmoddef
  type, extends (prc_writer_f_module_t) :: test_writer_2_t
   contains
     procedure, nopass :: type_name => test_writer_2_type_name
     procedure :: write_makefile_code => test_writer_2_mk
     procedure :: write_source_code => test_writer_2_src
     procedure :: write_interface => test_writer_2_if
     procedure :: write_wrapper => test_writer_2_wr
  end type test_writer_2_t
  
\eatline
\nwendcode{}\nwbegindocs{323}\nwdocspar
\nwenddocs{}\nwbegincode{324}\moddef{Prclib interfaces: tests}\plusendmoddef
  function test_writer_2_type_name () result (string)
    type(string_t) :: string
    string = "test_2"
  end function test_writer_2_type_name

  subroutine test_writer_2_mk (writer, unit, id, os_data, testflag)
    class(test_writer_2_t), intent(in) :: writer
    integer, intent(in) :: unit
    type(string_t), intent(in) :: id
    type(os_data_t), intent(in) :: os_data
    logical, intent(in), optional :: testflag
    write (unit, "(5A)")  "# Makefile code for process ", char (id), &
         " goes here."
  end subroutine test_writer_2_mk
  
  subroutine test_writer_2_src (writer, id)
    class(test_writer_2_t), intent(in) :: writer
    type(string_t), intent(in) :: id
  end subroutine test_writer_2_src
  
  subroutine test_writer_2_if (writer, unit, id, feature)
    class(test_writer_2_t), intent(in) :: writer
    integer, intent(in) :: unit
    type(string_t), intent(in) :: id, feature
    write (unit, "(2x,9A)")  "! Interface code for ", &
       char (writer%get_module_name (id)), "_", &
       char (writer%get_procname (feature)), " goes here."
  end subroutine test_writer_2_if

  subroutine test_writer_2_wr (writer, unit, id, feature)
    class(test_writer_2_t), intent(in) :: writer
    integer, intent(in) :: unit
    type(string_t), intent(in) :: id, feature
    write (unit, *)
    write (unit, "(9A)")  "! Wrapper code for ", &
       char (writer%get_c_procname (id, feature)), " goes here."
  end subroutine test_writer_2_wr

\eatline
\nwendcode{}\nwbegindocs{325}\nwdocspar
\subsubsection{Process library makefile}
Test 3: Write the makefile for compiling and linking the process
library (processes and driver code).  There are two processes, one
with one method, one with two methods.

To have predictable output, we reset the system-dependent initial
components of {\tt{}os{\char95}data} to known values.
\nwenddocs{}\nwbegincode{326}\moddef{Prclib interfaces: execute tests}\plusendmoddef
  call test (prclib_interfaces_3, "prclib_interfaces_3", &
       "write makefile", &
       u, results)
\nwendcode{}\nwbegincode{327}\moddef{Prclib interfaces: tests}\plusendmoddef
  subroutine prclib_interfaces_3 (u)
    integer, intent(in) :: u
    class(prclib_driver_t), allocatable :: driver
    type(os_data_t) :: os_data
    character(32), parameter :: md5sum = "prclib_interfaces_3_md5sum      "
    class(prc_writer_t), pointer :: test_writer_1, test_writer_2

    call os_data_init (os_data)
    os_data%fc = "fortran-compiler"
    os_data%whizard_includes = "-I module-dir"
    os_data%fcflags = "-C=all"
    os_data%fcflags_pic = "-PIC"
    os_data%cc = "c-compiler"
    os_data%cflags = "-I include-dir"
    os_data%cflags_pic = "-PIC"
    os_data%whizard_ldflags = ""
    os_data%ldflags = ""
    os_data%whizard_libtool = "my-libtool"
    os_data%latex = "latex -halt-on-error"
    os_data%mpost = "mpost --math=scaled -halt-on-error"
    os_data%dvips = "dvips"
    os_data%ps2pdf = "ps2pdf14"
    os_data%whizard_texpath = ""

    write (u, "(A)")  "* Test output: prclib_interfaces_3"
    write (u, "(A)")  "*   Purpose: check the generated Makefile"
    write (u, *)
    write (u, "(A)")  "* Create a prclib driver object (2 processes)"
    write (u, "(A)")

    call dispatch_prclib_driver (driver, var_str ("prclib3"), var_str (""))
    call driver%init (2)
    call driver%set_md5sum (md5sum)

    allocate (test_writer_1_t :: test_writer_1)
    allocate (test_writer_2_t :: test_writer_2)

    call driver%set_record (1, var_str ("test1"), var_str ("Test_model"), &
         [var_str ("proc1")], test_writer_1)

    call driver%set_record (2, var_str ("test2"), var_str ("Test_model"), &
         [var_str ("proc1"), var_str ("proc2")], test_writer_2)
    
    call driver%write (u)

    write (u, "(A)")
    write (u, "(A)")  "* Write Makefile"
    write (u, "(A)")  "* File contents:"
    write (u, "(A)")
    
    call driver%generate_makefile (u, os_data)
    
    deallocate (test_writer_1)
    deallocate (test_writer_2)

    write (u, "(A)")
    write (u, "(A)")  "* Test output end: prclib_interfaces_3"
  end subroutine prclib_interfaces_3
  
\eatline
\nwendcode{}\nwbegindocs{328}\nwdocspar
\subsubsection{Compile test with Fortran module}
Test 4: Write driver and makefile and try to compile and link the
library driver.

There is a single test process with a single feature.  The process
code is provided as a Fortran module, therefore we need a wrapper for
the featured procedure.
\nwenddocs{}\nwbegincode{329}\moddef{Prclib interfaces: execute tests}\plusendmoddef
  call test (prclib_interfaces_4, "prclib_interfaces_4", &
       "compile and link (Fortran module)", &
       u, results)
\nwendcode{}\nwbegincode{330}\moddef{Prclib interfaces: tests}\plusendmoddef
  subroutine prclib_interfaces_4 (u)
    integer, intent(in) :: u
    class(prclib_driver_t), allocatable :: driver
    class(prc_writer_t), pointer :: test_writer_4
    type(os_data_t) :: os_data
    integer :: u_file

    integer, dimension(:,:), allocatable :: flv_state
    integer, dimension(:,:), allocatable :: hel_state
    integer, dimension(:,:,:), allocatable :: col_state
    logical, dimension(:,:), allocatable :: ghost_flag
    integer, dimension(:,:), allocatable :: cf_index
    complex(default), dimension(:), allocatable :: color_factors
    character(32), parameter :: md5sum = "prclib_interfaces_4_md5sum      "
    character(32) :: md5sum_file

    type(c_funptr) :: proc1_ptr
    interface
       subroutine proc1_t (n) bind(C)
         import
         integer(c_int), intent(out) :: n
       end subroutine proc1_t
    end interface
    procedure(proc1_t), pointer :: proc1
    integer(c_int) :: n
    
    write (u, "(A)")  "* Test output: prclib_interfaces_4"
    write (u, "(A)")  "*   Purpose: compile, link, and load process library"
    write (u, "(A)")  "*            with (fake) matrix-element code &
         &as a Fortran module"
    write (u, *)
    write (u, "(A)")  "* Create a prclib driver object (1 process)"
    write (u, "(A)")

    call os_data_init (os_data)

    allocate (test_writer_4_t :: test_writer_4)
    call test_writer_4%init_test ()

    call dispatch_prclib_driver (driver, var_str ("prclib4"), var_str (""))
    call driver%init (1)
    call driver%set_md5sum (md5sum)

    call driver%set_record (1, var_str ("test4"), var_str ("Test_model"), &
         [var_str ("proc1")], test_writer_4)

    call driver%write (u)

    write (u, *)
    write (u, "(A)")  "* Write Makefile"
    u_file = free_unit ()
    open (u_file, file="prclib4.makefile", status="replace", action="write")
    call driver%generate_makefile (u_file, os_data)
    close (u_file)
    
    write (u, "(A)")
    write (u, "(A)")  "* Recall MD5 sum from Makefile"
    write (u, "(A)")
    
    md5sum_file = driver%get_md5sum_makefile ()
    write (u, "(1x,A,A,A)")  "MD5 sum = '", md5sum_file, "'"

    write (u, "(A)")
    write (u, "(A)")  "* Write driver source code"

    u_file = free_unit ()
    open (u_file, file="prclib4.f90", status="replace", action="write")
    call driver%generate_driver_code (u_file)
    close (u_file)
    
    write (u, "(A)")
    write (u, "(A)")  "* Recall MD5 sum from driver source"
    write (u, "(A)")
    
    md5sum_file = driver%get_md5sum_driver ()
    write (u, "(1x,A,A,A)")  "MD5 sum = '", md5sum_file, "'"

    write (u, "(A)")
    write (u, "(A)")  "* Write matrix-element source code"
    call driver%make_source (os_data)

    write (u, "(A)")
    write (u, "(A)")  "* Recall MD5 sum from matrix-element source"
    write (u, "(A)")
    
    md5sum_file = driver%get_md5sum_source (1)
    write (u, "(1x,A,A,A)")  "MD5 sum = '", md5sum_file, "'"

    write (u, "(A)")
    write (u, "(A)")  "* Compile source code"
    call driver%make_compile (os_data)
    
    write (u, "(A)")  "* Link library"
    call driver%make_link (os_data)
    
    write (u, "(A)")  "* Load library"
    call driver%load (os_data)

    write (u, *)
    call driver%write (u)
    write (u, *)
    
    if (driver%loaded) then
       write (u, "(A)")  "* Call library functions:"
       write (u, *)
       write (u, "(1x,A,I0)")  "n_processes   = ", driver%get_n_processes ()
       write (u, "(1x,A,A,A)")  "process_id    = '", &
            char (driver%get_process_id (1)), "'"
       write (u, "(1x,A,A,A)")  "model_name    = '", &
            char (driver%get_model_name (1)), "'"
       write (u, "(1x,A,A,A)")  "md5sum (lib)  = '", &
            char (driver%get_md5sum (0)), "'"
       write (u, "(1x,A,A,A)")  "md5sum (proc) = '", &
            char (driver%get_md5sum (1)), "'"
       write (u, "(1x,A,L1)")  "openmp_status = ", driver%get_openmp_status (1)
       write (u, "(1x,A,I0)")  "n_in  = ", driver%get_n_in (1)
       write (u, "(1x,A,I0)")  "n_out = ", driver%get_n_out (1)
       write (u, "(1x,A,I0)")  "n_flv = ", driver%get_n_flv (1)
       write (u, "(1x,A,I0)")  "n_hel = ", driver%get_n_hel (1)
       write (u, "(1x,A,I0)")  "n_col = ", driver%get_n_col (1)
       write (u, "(1x,A,I0)")  "n_cin = ", driver%get_n_cin (1)
       write (u, "(1x,A,I0)")  "n_cf  = ", driver%get_n_cf (1)

       call driver%set_flv_state (1, flv_state)
       write (u, "(1x,A,10(1x,I0))")  "flv_state =", flv_state

       call driver%set_hel_state (1, hel_state)
       write (u, "(1x,A,10(1x,I0))")  "hel_state =", hel_state

       call driver%set_col_state (1, col_state, ghost_flag)
       write (u, "(1x,A,10(1x,I0))")  "col_state =", col_state
       write (u, "(1x,A,10(1x,L1))")  "ghost_flag =", ghost_flag

       call driver%set_color_factors (1, color_factors, cf_index)
       write (u, "(1x,A,10(1x,F5.3))")  "color_factors =", color_factors
       write (u, "(1x,A,10(1x,I0))")  "cf_index =", cf_index

       call driver%get_fptr (1, 1, proc1_ptr)
       call c_f_procpointer (proc1_ptr, proc1)
       if (associated (proc1)) then
          write (u, *)
          call proc1 (n)
          write (u, "(1x,A,I0)")  "proc1(1) = ", n
       end if
       
    end if
    
    deallocate (test_writer_4)

    write (u, "(A)")
    write (u, "(A)")  "* Test output end: prclib_interfaces_4"
  end subroutine prclib_interfaces_4
  
\eatline
\nwendcode{}\nwbegindocs{331}This version of test-code writer actually writes an interface and
wrapper code.  The wrapped function is a no-parameter function with integer
result.

The stored MD5 sum may be modified.

We will reuse this later, therefore public.
\nwenddocs{}\nwbegincode{332}\moddef{Prclib interfaces: public}\plusendmoddef
  public :: test_writer_4_t
\nwendcode{}\nwbegindocs{333}\nwdocspar
\nwenddocs{}\nwbegincode{334}\moddef{Prclib interfaces: test types}\plusendmoddef
  type, extends (prc_writer_f_module_t) :: test_writer_4_t
   contains
     procedure, nopass :: type_name => test_writer_4_type_name
     procedure, nopass :: get_module_name => &
          test_writer_4_get_module_name
     procedure :: write_makefile_code => test_writer_4_mk
     procedure :: write_source_code => test_writer_4_src
     procedure :: write_interface => test_writer_4_if
     procedure :: write_wrapper => test_writer_4_wr
  end type test_writer_4_t
  
\eatline
\nwendcode{}\nwbegindocs{335}\nwdocspar
\nwenddocs{}\nwbegincode{336}\moddef{Prclib interfaces: tests}\plusendmoddef
  function test_writer_4_type_name () result (string)
    type(string_t) :: string
    string = "test_4"
  end function test_writer_4_type_name
  
  function test_writer_4_get_module_name (id) result (name)
    type(string_t), intent(in) :: id
    type(string_t) :: name
    name = "tpr_" // id
  end function test_writer_4_get_module_name

  subroutine test_writer_4_mk (writer, unit, id, os_data, testflag)
    class(test_writer_4_t), intent(in) :: writer
    integer, intent(in) :: unit
    type(string_t), intent(in) :: id
    type(os_data_t), intent(in) :: os_data
    logical, intent(in), optional :: testflag
    write (unit, "(5A)")  "SOURCES += ", char (id), ".f90"
    write (unit, "(5A)")  "OBJECTS += ", char (id), ".lo"
    write (unit, "(5A)")  "CLEAN_SOURCES += ", char (id), ".f90"
    write (unit, "(5A)")  "CLEAN_OBJECTS += tpr_", char (id), ".mod"
    write (unit, "(5A)")  "CLEAN_OBJECTS += ", char (id), ".lo"
    write (unit, "(5A)")  char (id), ".lo: ", char (id), ".f90"
    write (unit, "(5A)")  TAB, "$(LTFCOMPILE) $<"    
  end subroutine test_writer_4_mk
  
  subroutine test_writer_4_src (writer, id)
    class(test_writer_4_t), intent(in) :: writer
    type(string_t), intent(in) :: id
    call write_test_module_file (id, var_str ("proc1"), writer%md5sum)
  end subroutine test_writer_4_src
  
  subroutine test_writer_4_if (writer, unit, id, feature)
    class(test_writer_4_t), intent(in) :: writer
    integer, intent(in) :: unit
    type(string_t), intent(in) :: id, feature
    write (unit, "(2x,9A)")  "interface"
    write (unit, "(5x,9A)")  "subroutine ", &
       char (writer%get_c_procname (id, feature)), &
       " (n) bind(C)"
    write (unit, "(7x,9A)")  "import"
    write (unit, "(7x,9A)")  "implicit none"
    write (unit, "(7x,9A)")  "integer(c_int), intent(out) :: n"    
    write (unit, "(5x,9A)")  "end subroutine ", &
       char (writer%get_c_procname (id, feature))
    write (unit, "(2x,9A)")  "end interface"
  end subroutine test_writer_4_if

  subroutine test_writer_4_wr (writer, unit, id, feature)
    class(test_writer_4_t), intent(in) :: writer
    integer, intent(in) :: unit
    type(string_t), intent(in) :: id, feature
    write (unit, *)
    write (unit, "(9A)")  "subroutine ", &
         char (writer%get_c_procname (id, feature)), &
       " (n) bind(C)"
    write (unit, "(2x,9A)")  "use iso_c_binding"
    write (unit, "(2x,9A)")  "use tpr_", char (id), ", only: ", &
         char (writer%get_procname (feature))
    write (unit, "(2x,9A)")  "implicit none"
    write (unit, "(2x,9A)")  "integer(c_int), intent(out) :: n"    
    write (unit, "(2x,9A)")  "call ", char (feature), " (n)"
    write (unit, "(9A)")  "end subroutine ", &
       char (writer%get_c_procname (id, feature))
  end subroutine test_writer_4_wr

\eatline
\nwendcode{}\nwbegindocs{337}\nwdocspar
We need a test module file (actually, one for each process in the test
above) that allows us to check compilation and linking.  The test
module implements a colorless $1\to 2$ process, and it implements one
additional function (feature), the name given as an argument.
\nwenddocs{}\nwbegincode{338}\moddef{Prclib interfaces: tests}\plusendmoddef
  subroutine write_test_module_file (basename, feature, md5sum)
    type(string_t), intent(in) :: basename
    type(string_t), intent(in) :: feature
    character(32), intent(in) :: md5sum
    integer :: u
    u = free_unit ()
    open (u, file = char (basename) // ".f90", &
         status = "replace", action = "write")
    write (u, "(A)")  "! (Pseudo) matrix element code file &
         &for WHIZARD self-test"
    write (u, *)
    write (u, "(A)")  "module tpr_" // char (basename)
    write (u, *)
    write (u, "(2x,A)")  "use kinds"
    write (u, "(2x,A)")  "use omega_color, OCF => omega_color_factor"
    write (u, *)
    write (u, "(2x,A)")  "implicit none"
    write (u, "(2x,A)")  "private"
    write (u, *)
    call write_test_me_code_1 (u)
    write (u, *)
    write (u, "(2x,A)")  "public :: " // char (feature)
    write (u, *)
    write (u, "(A)")  "contains"
    write (u, *)
    call write_test_me_code_2 (u, md5sum)
    write (u, *)
    write (u, "(2x,A)")  "subroutine " // char (feature) // " (n)"
    write (u, "(2x,A)")  "  integer, intent(out) :: n"
    write (u, "(2x,A)")  "  n = 42"
    write (u, "(2x,A)")  "end subroutine " // char (feature)
    write (u, *)
    write (u, "(A)")  "end module tpr_" // char (basename)
    close (u)
  end subroutine write_test_module_file
    
\eatline
\nwendcode{}\nwbegindocs{339}\nwdocspar
The following two subroutines provide building blocks for a
matrix-element source code file, useful only for testing the
workflow.  The first routine writes the header part, the other routine
the implementation of the procedures listed in the header.
\nwenddocs{}\nwbegincode{340}\moddef{Prclib interfaces: tests}\plusendmoddef
  subroutine write_test_me_code_1 (u)
    integer, intent(in) :: u
    write (u, "(2x,A)")  "public :: md5sum"
    write (u, "(2x,A)")  "public :: openmp_supported"
    write (u, *)
    write (u, "(2x,A)")  "public :: n_in"
    write (u, "(2x,A)")  "public :: n_out"
    write (u, "(2x,A)")  "public :: n_flv"
    write (u, "(2x,A)")  "public :: n_hel"
    write (u, "(2x,A)")  "public :: n_cin"
    write (u, "(2x,A)")  "public :: n_col"
    write (u, "(2x,A)")  "public :: n_cf"
    write (u, *)
    write (u, "(2x,A)")  "public :: flv_state"
    write (u, "(2x,A)")  "public :: hel_state"
    write (u, "(2x,A)")  "public :: col_state"
    write (u, "(2x,A)")  "public :: color_factors"
  end subroutine write_test_me_code_1
    
  subroutine write_test_me_code_2 (u, md5sum)
    integer, intent(in) :: u
    character(32), intent(in) :: md5sum
    write (u, "(2x,A)")  "pure function md5sum ()"
    write (u, "(2x,A)")  "  character(len=32) :: md5sum"
    write (u, "(2x,A)")  "  md5sum = '" // md5sum // "'"
    write (u, "(2x,A)")  "end function md5sum"
    write (u, *)
    write (u, "(2x,A)")  "pure function openmp_supported () result (status)"
    write (u, "(2x,A)")  "  logical :: status"
    write (u, "(2x,A)")  "  status = .false."
    write (u, "(2x,A)")  "end function openmp_supported"
    write (u, *)
    write (u, "(2x,A)")  "pure function n_in () result (n)"
    write (u, "(2x,A)")  "  integer :: n"
    write (u, "(2x,A)")  "  n = 1"
    write (u, "(2x,A)")  "end function n_in"
    write (u, *)
    write (u, "(2x,A)")  "pure function n_out () result (n)"
    write (u, "(2x,A)")  "  integer :: n"
    write (u, "(2x,A)")  "  n = 2"
    write (u, "(2x,A)")  "end function n_out"
    write (u, *)
    write (u, "(2x,A)")  "pure function n_flv () result (n)"
    write (u, "(2x,A)")  "  integer :: n"
    write (u, "(2x,A)")  "  n = 1"
    write (u, "(2x,A)")  "end function n_flv"
    write (u, *)
    write (u, "(2x,A)")  "pure function n_hel () result (n)"
    write (u, "(2x,A)")  "  integer :: n"
    write (u, "(2x,A)")  "  n = 1"
    write (u, "(2x,A)")  "end function n_hel"
    write (u, *)
    write (u, "(2x,A)")  "pure function n_cin () result (n)"
    write (u, "(2x,A)")  "  integer :: n"
    write (u, "(2x,A)")  "  n = 2"
    write (u, "(2x,A)")  "end function n_cin"
    write (u, *)
    write (u, "(2x,A)")  "pure function n_col () result (n)"
    write (u, "(2x,A)")  "  integer :: n"
    write (u, "(2x,A)")  "  n = 1"
    write (u, "(2x,A)")  "end function n_col"
    write (u, *)
    write (u, "(2x,A)")  "pure function n_cf () result (n)"
    write (u, "(2x,A)")  "  integer :: n"
    write (u, "(2x,A)")  "  n = 1"
    write (u, "(2x,A)")  "end function n_cf"
    write (u, *)
    write (u, "(2x,A)")  "pure subroutine flv_state (a)"
    write (u, "(2x,A)")  "  integer, dimension(:,:), intent(out) :: a"
    write (u, "(2x,A)")  "  a = reshape ([1,2,3], [3,1])"
    write (u, "(2x,A)")  "end subroutine flv_state"
    write (u, *)
    write (u, "(2x,A)")  "pure subroutine hel_state (a)"
    write (u, "(2x,A)")  "  integer, dimension(:,:), intent(out) :: a"
    write (u, "(2x,A)")  "  a = reshape ([0,0,0], [3,1])"
    write (u, "(2x,A)")  "end subroutine hel_state"
    write (u, *)
    write (u, "(2x,A)")  "pure subroutine col_state (a, g)"
    write (u, "(2x,A)")  "  integer, dimension(:,:,:), intent(out) :: a"
    write (u, "(2x,A)")  "  logical, dimension(:,:), intent(out) :: g"
    write (u, "(2x,A)")  "  a = reshape ([0,0, 0,0, 0,0], [2,3,1])"
    write (u, "(2x,A)")  "  g = reshape ([.false., .false., .false.], [3,1])"
    write (u, "(2x,A)")  "end subroutine col_state"
    write (u, *)
    write (u, "(2x,A)")  "pure subroutine color_factors (cf)"
    write (u, "(2x,A)")  "  type(OCF), dimension(:), intent(out) :: cf"
    write (u, "(2x,A)")  "  cf = [ OCF(1,1,+1._default) ]"
    write (u, "(2x,A)")  "end subroutine color_factors"
  end subroutine write_test_me_code_2
    
\eatline
\nwendcode{}\nwbegindocs{341}\nwdocspar
\subsubsection{Compile test with Fortran bind(C) library}
Test 5: Write driver and makefile and try to compile and link the
library driver.

There is a single test process with a single feature.  The process
code is provided as a Fortran library of independent procedures.
These procedures are bind(C).
\nwenddocs{}\nwbegincode{342}\moddef{Prclib interfaces: execute tests}\plusendmoddef
  call test (prclib_interfaces_5, "prclib_interfaces_5", &
       "compile and link (Fortran library)", &
       u, results)
\nwendcode{}\nwbegincode{343}\moddef{Prclib interfaces: tests}\plusendmoddef
  subroutine prclib_interfaces_5 (u)
    integer, intent(in) :: u
    class(prclib_driver_t), allocatable :: driver
    class(prc_writer_t), pointer :: test_writer_5
    type(os_data_t) :: os_data
    integer :: u_file

    integer, dimension(:,:), allocatable :: flv_state
    integer, dimension(:,:), allocatable :: hel_state
    integer, dimension(:,:,:), allocatable :: col_state
    logical, dimension(:,:), allocatable :: ghost_flag
    integer, dimension(:,:), allocatable :: cf_index
    complex(default), dimension(:), allocatable :: color_factors
    character(32), parameter :: md5sum = "prclib_interfaces_5_md5sum      "

    type(c_funptr) :: proc1_ptr
    interface
       subroutine proc1_t (n) bind(C)
         import
         integer(c_int), intent(out) :: n
       end subroutine proc1_t
    end interface
    procedure(proc1_t), pointer :: proc1
    integer(c_int) :: n
    
    write (u, "(A)")  "* Test output: prclib_interfaces_5"
    write (u, "(A)")  "*   Purpose: compile, link, and load process library"
    write (u, "(A)")  "*            with (fake) matrix-element code &
         &as a Fortran bind(C) library"
    write (u, *)
    write (u, "(A)")  "* Create a prclib driver object (1 process)"
    write (u, "(A)")

    call os_data_init (os_data)
    allocate (test_writer_5_t :: test_writer_5)

    call dispatch_prclib_driver (driver, var_str ("prclib5"), var_str (""))
    call driver%init (1)
    call driver%set_md5sum (md5sum)

    call driver%set_record (1, var_str ("test5"), var_str ("Test_model"), &
         [var_str ("proc1")], test_writer_5)

    call driver%write (u)

    write (u, *)
    write (u, "(A)")  "* Write makefile"
    u_file = free_unit ()
    open (u_file, file="prclib5.makefile", status="replace", action="write")
    call driver%generate_makefile (u_file, os_data)
    close (u_file)
    
    write (u, "(A)")  "* Write driver source code"
    u_file = free_unit ()
    open (u_file, file="prclib5.f90", status="replace", action="write")
    call driver%generate_driver_code (u_file)
    close (u_file)
    
    write (u, "(A)")  "* Write matrix-element source code"
    call driver%make_source (os_data)

    write (u, "(A)")  "* Compile source code"
    call driver%make_compile (os_data)
    
    write (u, "(A)")  "* Link library"
    call driver%make_link (os_data)
    
    write (u, "(A)")  "* Load library"
    call driver%load (os_data)

    write (u, *)
    call driver%write (u)
    write (u, *)
    
    if (driver%loaded) then
       write (u, "(A)")  "* Call library functions:"
       write (u, *)
       write (u, "(1x,A,I0)")  "n_processes   = ", driver%get_n_processes ()
       write (u, "(1x,A,A)")  "process_id    = ", &
            char (driver%get_process_id (1))
       write (u, "(1x,A,A)")  "model_name    = ", &
            char (driver%get_model_name (1))
       write (u, "(1x,A,A)")  "md5sum        = ", &
            char (driver%get_md5sum (1))
       write (u, "(1x,A,L1)")  "openmp_status = ", driver%get_openmp_status (1)
       write (u, "(1x,A,I0)")  "n_in  = ", driver%get_n_in (1)
       write (u, "(1x,A,I0)")  "n_out = ", driver%get_n_out (1)
       write (u, "(1x,A,I0)")  "n_flv = ", driver%get_n_flv (1)
       write (u, "(1x,A,I0)")  "n_hel = ", driver%get_n_hel (1)
       write (u, "(1x,A,I0)")  "n_col = ", driver%get_n_col (1)
       write (u, "(1x,A,I0)")  "n_cin = ", driver%get_n_cin (1)
       write (u, "(1x,A,I0)")  "n_cf  = ", driver%get_n_cf (1)

       call driver%set_flv_state (1, flv_state)
       write (u, "(1x,A,10(1x,I0))")  "flv_state =", flv_state

       call driver%set_hel_state (1, hel_state)
       write (u, "(1x,A,10(1x,I0))")  "hel_state =", hel_state

       call driver%set_col_state (1, col_state, ghost_flag)
       write (u, "(1x,A,10(1x,I0))")  "col_state =", col_state
       write (u, "(1x,A,10(1x,L1))")  "ghost_flag =", ghost_flag

       call driver%set_color_factors (1, color_factors, cf_index)
       write (u, "(1x,A,10(1x,F5.3))")  "color_factors =", color_factors
       write (u, "(1x,A,10(1x,I0))")  "cf_index =", cf_index

       call driver%get_fptr (1, 1, proc1_ptr)
       call c_f_procpointer (proc1_ptr, proc1)
       if (associated (proc1)) then
          write (u, *)
          call proc1 (n)
          write (u, "(1x,A,I0)")  "proc1(1) = ", n
       end if
       
    end if
    
    deallocate (test_writer_5)

    write (u, "(A)")
    write (u, "(A)")  "* Test output end: prclib_interfaces_5"
  end subroutine prclib_interfaces_5
  
\eatline
\nwendcode{}\nwbegindocs{344}This version of test-code writer writes interfaces for all standard
features plus one specific feature.  The interfaces are all bind(C),
so no wrapper is needed.
\nwenddocs{}\nwbegincode{345}\moddef{Prclib interfaces: test types}\plusendmoddef
  type, extends (prc_writer_c_lib_t) :: test_writer_5_t
   contains
     procedure, nopass :: type_name => test_writer_5_type_name
     procedure :: write_makefile_code => test_writer_5_mk
     procedure :: write_source_code => test_writer_5_src
     procedure :: write_interface => test_writer_5_if
  end type test_writer_5_t
  
\eatline
\nwendcode{}\nwbegindocs{346}The 
\nwenddocs{}\nwbegincode{347}\moddef{Prclib interfaces: tests}\plusendmoddef
  function test_writer_5_type_name () result (string)
    type(string_t) :: string
    string = "test_5"
  end function test_writer_5_type_name

  subroutine test_writer_5_mk (writer, unit, id, os_data, testflag)
    class(test_writer_5_t), intent(in) :: writer
    integer, intent(in) :: unit
    type(string_t), intent(in) :: id
    type(os_data_t), intent(in) :: os_data
    logical, intent(in), optional :: testflag 
    write (unit, "(5A)")  "SOURCES += ", char (id), ".f90"
    write (unit, "(5A)")  "OBJECTS += ", char (id), ".lo"
    write (unit, "(5A)")  char (id), ".lo: ", char (id), ".f90"
    write (unit, "(5A)")  TAB, "$(LTFCOMPILE) $<"
  end subroutine test_writer_5_mk
  
  subroutine test_writer_5_src (writer, id)
    class(test_writer_5_t), intent(in) :: writer
    type(string_t), intent(in) :: id
    call write_test_f_lib_file (id, var_str ("proc1"))
  end subroutine test_writer_5_src
  
  subroutine test_writer_5_if (writer, unit, id, feature)
    class(test_writer_5_t), intent(in) :: writer
    integer, intent(in) :: unit
    type(string_t), intent(in) :: id, feature
    select case (char (feature))
    case ("proc1")
       write (unit, "(2x,9A)")  "interface"
       write (unit, "(5x,9A)")  "subroutine ", &
            char (writer%get_c_procname (id, feature)), &
            " (n) bind(C)"
       write (unit, "(7x,9A)")  "import"
       write (unit, "(7x,9A)")  "implicit none"
       write (unit, "(7x,9A)")  "integer(c_int), intent(out) :: n"    
       write (unit, "(5x,9A)")  "end subroutine ", &
            char (writer%get_c_procname (id, feature))
       write (unit, "(2x,9A)")  "end interface"
    case default
       call writer%write_standard_interface (unit, id, feature)
    end select
  end subroutine test_writer_5_if

\eatline
\eatline
\nwendcode{}\nwbegindocs{348}\nwdocspar
We need a test module file (actually, one for each process in the test
above) that allows us to check compilation and linking.  The test
module implements a colorless $1\to 2$ process, and it implements one
additional function (feature), the name given as an argument.
\nwenddocs{}\nwbegincode{349}\moddef{Prclib interfaces: tests}\plusendmoddef
  subroutine write_test_f_lib_file (basename, feature)
    type(string_t), intent(in) :: basename
    type(string_t), intent(in) :: feature
    integer :: u
    u = free_unit ()
    open (u, file = char (basename) // ".f90", &
         status = "replace", action = "write")
    write (u, "(A)")  "! (Pseudo) matrix element code file &
         &for WHIZARD self-test"
    call write_test_me_code_3 (u, char (basename))
    write (u, *)
    write (u, "(A)")  "subroutine " // char (basename) // "_" &
         // char (feature) // " (n) bind(C)"
    write (u, "(A)")  "  use iso_c_binding"
    write (u, "(A)")  "  implicit none"
    write (u, "(A)")  "  integer(c_int), intent(out) :: n"
    write (u, "(A)")  "  n = 42"
    write (u, "(A)")  "end subroutine " // char (basename) // "_" &
         // char (feature)
    close (u)
  end subroutine write_test_f_lib_file
    
\eatline
\nwendcode{}\nwbegindocs{350}\nwdocspar
The following matrix-element source code is identical to the previous
one, but modified such as to provide independent procedures without a
module envelope.
\nwenddocs{}\nwbegincode{351}\moddef{Prclib interfaces: tests}\plusendmoddef
  subroutine write_test_me_code_3 (u, id)
    integer, intent(in) :: u
    character(*), intent(in) :: id
    write (u, "(A)")  "function " // id // "_get_md5sum () &
         &result (cptr) bind(C)"
    write (u, "(A)")  "  use iso_c_binding"
    write (u, "(A)")  "  implicit none"
    write (u, "(A)")  "  type(c_ptr) :: cptr"
    write (u, "(A)")  "  character(c_char), dimension(32), &
         &target, save :: md5sum"
    write (u, "(A)")  "  md5sum = copy (c_char_&
         &'1234567890abcdef1234567890abcdef')"
    write (u, "(A)")  "  cptr = c_loc (md5sum)"
    write (u, "(A)")  "contains"
    write (u, "(A)")  "  function copy (md5sum)"
    write (u, "(A)")  "    character(c_char), dimension(32) :: copy"
    write (u, "(A)")  "    character(c_char), dimension(32), intent(in) :: &
         &md5sum"
    write (u, "(A)")  "    copy = md5sum"
    write (u, "(A)")  "  end function copy"
    write (u, "(A)")  "end function " // id // "_get_md5sum"
    write (u, *)
    write (u, "(A)")  "function " // id // "_openmp_supported () &
         &result (status) bind(C)"
    write (u, "(A)")  "  use iso_c_binding"
    write (u, "(A)")  "  implicit none"
    write (u, "(A)")  "  logical(c_bool) :: status"
    write (u, "(A)")  "  status = .false."
    write (u, "(A)")  "end function " // id // "_openmp_supported"
    write (u, *)
    write (u, "(A)")  "function " // id // "_n_in () result (n) bind(C)"
    write (u, "(A)")  "  use iso_c_binding"
    write (u, "(A)")  "  implicit none"
    write (u, "(A)")  "  integer(c_int) :: n"
    write (u, "(A)")  "  n = 1"
    write (u, "(A)")  "end function " // id // "_n_in"
    write (u, *)
    write (u, "(A)")  "function " // id // "_n_out () result (n) bind(C)"
    write (u, "(A)")  "  use iso_c_binding"
    write (u, "(A)")  "  implicit none"
    write (u, "(A)")  "  integer(c_int) :: n"
    write (u, "(A)")  "  n = 2"
    write (u, "(A)")  "end function " // id // "_n_out"
    write (u, *)
    write (u, "(A)")  "function " // id // "_n_flv () result (n) bind(C)"
    write (u, "(A)")  "  use iso_c_binding"
    write (u, "(A)")  "  implicit none"
    write (u, "(A)")  "  integer(c_int) :: n"
    write (u, "(A)")  "  n = 1"
    write (u, "(A)")  "end function " // id // "_n_flv"
    write (u, *)
    write (u, "(A)")  "function " // id // "_n_hel () result (n) bind(C)"
    write (u, "(A)")  "  use iso_c_binding"
    write (u, "(A)")  "  implicit none"
    write (u, "(A)")  "  integer(c_int) :: n"
    write (u, "(A)")  "  n = 1"
    write (u, "(A)")  "end function " // id // "_n_hel"
    write (u, *)
    write (u, "(A)")  "function " // id // "_n_cin () result (n) bind(C)"
    write (u, "(A)")  "  use iso_c_binding"
    write (u, "(A)")  "  implicit none"
    write (u, "(A)")  "  integer(c_int) :: n"
    write (u, "(A)")  "  n = 2"
    write (u, "(A)")  "end function " // id // "_n_cin"
    write (u, *)
    write (u, "(A)")  "function " // id // "_n_col () result (n) bind(C)"
    write (u, "(A)")  "  use iso_c_binding"
    write (u, "(A)")  "  implicit none"
    write (u, "(A)")  "  integer(c_int) :: n"
    write (u, "(A)")  "  n = 1"
    write (u, "(A)")  "end function " // id // "_n_col"
    write (u, *)
    write (u, "(A)")  "function " // id // "_n_cf () result (n) bind(C)"
    write (u, "(A)")  "  use iso_c_binding"
    write (u, "(A)")  "  implicit none"
    write (u, "(A)")  "  integer(c_int) :: n"
    write (u, "(A)")  "  n = 1"
    write (u, "(A)")  "end function " // id // "_n_cf"
    write (u, *)
    write (u, "(A)")  "subroutine " // id // "_flv_state (flv_state) bind(C)"
    write (u, "(A)")  "  use iso_c_binding"
    write (u, "(A)")  "  implicit none"
    write (u, "(A)")  "  integer(c_int), dimension(*), intent(out) :: flv_state"
    write (u, "(A)")  "  flv_state(1:3) = [1,2,3]"
    write (u, "(A)")  "end subroutine " // id // "_flv_state"
    write (u, *)
    write (u, "(A)")  "subroutine " // id // "_hel_state (hel_state) bind(C)"
    write (u, "(A)")  "  use iso_c_binding"
    write (u, "(A)")  "  implicit none"
    write (u, "(A)")  "  integer(c_int), dimension(*), intent(out) :: hel_state"
    write (u, "(A)")  "  hel_state(1:3) = [0,0,0]"
    write (u, "(A)")  "end subroutine " // id // "_hel_state"
    write (u, *)
    write (u, "(A)")  "subroutine " // id // "_col_state &
         &(col_state, ghost_flag) bind(C)"
    write (u, "(A)")  "  use iso_c_binding"
    write (u, "(A)")  "  implicit none"
    write (u, "(A)")  "  integer(c_int), dimension(*), intent(out) &
         &:: col_state"
    write (u, "(A)")  "  logical(c_bool), dimension(*), intent(out) &
         &:: ghost_flag"
    write (u, "(A)")  "  col_state(1:6) = [0,0, 0,0, 0,0]"
    write (u, "(A)")  "  ghost_flag(1:3) = [.false., .false., .false.]"
    write (u, "(A)")  "end subroutine " // id // "_col_state"
    write (u, *)
    write (u, "(A)")  "subroutine " // id // "_color_factors &
         &(cf_index1, cf_index2, color_factors) bind(C)"
    write (u, "(A)")  "  use iso_c_binding"
    write (u, "(A)")  "  use kinds"
    write (u, "(A)")  "  implicit none"
    write (u, "(A)")  "  integer(c_int), dimension(*), intent(out) :: cf_index1"
    write (u, "(A)")  "  integer(c_int), dimension(*), intent(out) :: cf_index2"
    write (u, "(A)")  "  complex(c_default_complex), dimension(*), &
         &intent(out) :: color_factors"
    write (u, "(A)")  "  cf_index1(1:1) = [1]"
    write (u, "(A)")  "  cf_index2(1:1) = [1]"
    write (u, "(A)")  "  color_factors(1:1) = [1]"
    write (u, "(A)")  "end subroutine " // id // "_color_factors"
  end subroutine write_test_me_code_3
    
\eatline
\nwendcode{}\nwbegindocs{352}\nwdocspar
\subsubsection{Compile test with genuine C library}
Test 6: Write driver and makefile and try to compile and link the
library driver.

There is a single test process with a single feature.  The process
code is provided as a C library of independent procedures.
These procedures should match the Fortran bind(C) interface.
\nwenddocs{}\nwbegincode{353}\moddef{Prclib interfaces: execute tests}\plusendmoddef
  call test (prclib_interfaces_6, "prclib_interfaces_6", &
       "compile and link (C library)", &
       u, results)
\nwendcode{}\nwbegincode{354}\moddef{Prclib interfaces: tests}\plusendmoddef
  subroutine prclib_interfaces_6 (u)
    integer, intent(in) :: u
    class(prclib_driver_t), allocatable :: driver
    class(prc_writer_t), pointer :: test_writer_6
    type(os_data_t) :: os_data
    integer :: u_file

    integer, dimension(:,:), allocatable :: flv_state
    integer, dimension(:,:), allocatable :: hel_state
    integer, dimension(:,:,:), allocatable :: col_state
    logical, dimension(:,:), allocatable :: ghost_flag
    integer, dimension(:,:), allocatable :: cf_index
    complex(default), dimension(:), allocatable :: color_factors
    character(32), parameter :: md5sum = "prclib_interfaces_6_md5sum      "

    type(c_funptr) :: proc1_ptr
    interface
       subroutine proc1_t (n) bind(C)
         import
         integer(c_int), intent(out) :: n
       end subroutine proc1_t
    end interface
    procedure(proc1_t), pointer :: proc1
    integer(c_int) :: n
    
    write (u, "(A)")  "* Test output: prclib_interfaces_6"
    write (u, "(A)")  "*   Purpose: compile, link, and load process library"
    write (u, "(A)")  "*            with (fake) matrix-element code &
         &as a C library"
    write (u, *)
    write (u, "(A)")  "* Create a prclib driver object (1 process)"
    write (u, "(A)")

    call os_data_init (os_data)
    allocate (test_writer_6_t :: test_writer_6)

    call dispatch_prclib_driver (driver, var_str ("prclib6"), var_str (""))
    call driver%init (1)
    call driver%set_md5sum (md5sum)

    call driver%set_record (1, var_str ("test6"), var_str ("Test_model"), &
         [var_str ("proc1")], test_writer_6)

    call driver%write (u)

    write (u, *)
    write (u, "(A)")  "* Write makefile"
    u_file = free_unit ()
    open (u_file, file="prclib6.makefile", status="replace", action="write")
    call driver%generate_makefile (u_file, os_data)
    close (u_file)
    
    write (u, "(A)")  "* Write driver source code"
    u_file = free_unit ()
    open (u_file, file="prclib6.f90", status="replace", action="write")
    call driver%generate_driver_code (u_file)
    close (u_file)
    
    write (u, "(A)")  "* Write matrix-element source code"
    call driver%make_source (os_data)

    write (u, "(A)")  "* Compile source code"
    call driver%make_compile (os_data)
    
    write (u, "(A)")  "* Link library"
    call driver%make_link (os_data)
    
    write (u, "(A)")  "* Load library"
    call driver%load (os_data)

    write (u, *)
    call driver%write (u)
    write (u, *)
    
    if (driver%loaded) then
       write (u, "(A)")  "* Call library functions:"
       write (u, *)
       write (u, "(1x,A,I0)")  "n_processes   = ", driver%get_n_processes ()
       write (u, "(1x,A,A)")  "process_id    = ", &
            char (driver%get_process_id (1))
       write (u, "(1x,A,A)")  "model_name    = ", &
            char (driver%get_model_name (1))
       write (u, "(1x,A,A)")  "md5sum        = ", &
            char (driver%get_md5sum (1))
       write (u, "(1x,A,L1)")  "openmp_status = ", driver%get_openmp_status (1)
       write (u, "(1x,A,I0)")  "n_in  = ", driver%get_n_in (1)
       write (u, "(1x,A,I0)")  "n_out = ", driver%get_n_out (1)
       write (u, "(1x,A,I0)")  "n_flv = ", driver%get_n_flv (1)
       write (u, "(1x,A,I0)")  "n_hel = ", driver%get_n_hel (1)
       write (u, "(1x,A,I0)")  "n_col = ", driver%get_n_col (1)
       write (u, "(1x,A,I0)")  "n_cin = ", driver%get_n_cin (1)
       write (u, "(1x,A,I0)")  "n_cf  = ", driver%get_n_cf (1)

       call driver%set_flv_state (1, flv_state)
       write (u, "(1x,A,10(1x,I0))")  "flv_state =", flv_state

       call driver%set_hel_state (1, hel_state)
       write (u, "(1x,A,10(1x,I0))")  "hel_state =", hel_state

       call driver%set_col_state (1, col_state, ghost_flag)
       write (u, "(1x,A,10(1x,I0))")  "col_state =", col_state
       write (u, "(1x,A,10(1x,L1))")  "ghost_flag =", ghost_flag

       call driver%set_color_factors (1, color_factors, cf_index)
       write (u, "(1x,A,10(1x,F5.3))")  "color_factors =", color_factors
       write (u, "(1x,A,10(1x,I0))")  "cf_index =", cf_index

       call driver%get_fptr (1, 1, proc1_ptr)
       call c_f_procpointer (proc1_ptr, proc1)
       if (associated (proc1)) then
          write (u, *)
          call proc1 (n)
          write (u, "(1x,A,I0)")  "proc1(1) = ", n
       end if
       
    end if
    
    deallocate (test_writer_6)

    write (u, "(A)")
    write (u, "(A)")  "* Test output end: prclib_interfaces_6"
  end subroutine prclib_interfaces_6
  
\eatline
\nwendcode{}\nwbegindocs{355}This version of test-code writer writes interfaces for all standard
features plus one specific feature.  The interfaces are all bind(C),
so no wrapper is needed.

The driver part is identical to the Fortran case, so we simply extend
the previous {\tt{}test{\char95}writer{\char95}5} type.  We only have to override the
Makefile writer.
\nwenddocs{}\nwbegincode{356}\moddef{Prclib interfaces: test types}\plusendmoddef
  type, extends (test_writer_5_t) :: test_writer_6_t
   contains
     procedure, nopass :: type_name => test_writer_6_type_name
     procedure :: write_makefile_code => test_writer_6_mk
     procedure :: write_source_code => test_writer_6_src
  end type test_writer_6_t
  
\eatline
\nwendcode{}\nwbegindocs{357}\nwdocspar
\nwenddocs{}\nwbegincode{358}\moddef{Prclib interfaces: tests}\plusendmoddef
  function test_writer_6_type_name () result (string)
    type(string_t) :: string
    string = "test_6"
  end function test_writer_6_type_name

  subroutine test_writer_6_mk (writer, unit, id, os_data, testflag)
    class(test_writer_6_t), intent(in) :: writer
    integer, intent(in) :: unit
    type(string_t), intent(in) :: id
    type(os_data_t), intent(in) :: os_data
    logical, intent(in), optional :: testflag
    write (unit, "(5A)")  "SOURCES += ", char (id), ".c"
    write (unit, "(5A)")  "OBJECTS += ", char (id), ".lo"
    write (unit, "(5A)")  char (id), ".lo: ", char (id), ".c"
    write (unit, "(5A)")  TAB, "$(LTCCOMPILE) $<"
  end subroutine test_writer_6_mk

  subroutine test_writer_6_src (writer, id)
    class(test_writer_6_t), intent(in) :: writer
    type(string_t), intent(in) :: id
    call write_test_c_lib_file (id, var_str ("proc1"))
  end subroutine test_writer_6_src
  
\eatline
\nwendcode{}\nwbegindocs{359}\nwdocspar
We need a test module file (actually, one for each process in the test
above) that allows us to check compilation and linking.  The test
module implements a colorless $1\to 2$ process, and it implements one
additional function (feature), the name given as an argument.
\nwenddocs{}\nwbegincode{360}\moddef{Prclib interfaces: tests}\plusendmoddef
  subroutine write_test_c_lib_file (basename, feature)
    type(string_t), intent(in) :: basename
    type(string_t), intent(in) :: feature
    integer :: u
    u = free_unit ()
    open (u, file = char (basename) // ".c", &
         status = "replace", action = "write")
    write (u, "(A)")  "/* (Pseudo) matrix element code file &
         &for WHIZARD self-test */"
    write (u, "(A)")  "#include <stdbool.h>"
    write (u, *)
    call write_test_me_code_4 (u, char (basename))
    write (u, *)
    write (u, "(A)")  "void " // char (basename) // "_" &
         // char (feature) // "(int* n) \{"
    write (u, "(A)")  "  *n = 42;"
    write (u, "(A)")  "\}"
    close (u)
  end subroutine write_test_c_lib_file
    
\eatline
\nwendcode{}\nwbegindocs{361}\nwdocspar
The following matrix-element source code is equivalent to the code in
the previous example, but coded in C.
\nwenddocs{}\nwbegincode{362}\moddef{Prclib interfaces: tests}\plusendmoddef
  subroutine write_test_me_code_4 (u, id)
    integer, intent(in) :: u
    character(*), intent(in) :: id
    write (u, "(A)")  "char* " // id // "_get_md5sum() \{"
    write (u, "(A)")  "  return ""1234567890abcdef1234567890abcdef"";"
    write (u, "(A)")  "\}"
    write (u, *)
    write (u, "(A)")  "bool " // id // "_openmp_supported() \{"
    write (u, "(A)")  "  return false;"
    write (u, "(A)")  "\}"
    write (u, *)
    write (u, "(A)")  "int " // id // "_n_in() \{"
    write (u, "(A)")  "  return 1;"
    write (u, "(A)")  "\}"
    write (u, *)
    write (u, "(A)")  "int " // id // "_n_out() \{"
    write (u, "(A)")  "  return 2;"
    write (u, "(A)")  "\}"
    write (u, *)
    write (u, "(A)")  "int " // id // "_n_flv() \{"
    write (u, "(A)")  "  return 1;"
    write (u, "(A)")  "\}"
    write (u, *)
    write (u, "(A)")  "int " // id // "_n_hel() \{"
    write (u, "(A)")  "  return 1;"
    write (u, "(A)")  "\}"
    write (u, *)
    write (u, "(A)")  "int " // id // "_n_cin() \{"
    write (u, "(A)")  "  return 2;"
    write (u, "(A)")  "\}"
    write (u, *)
    write (u, "(A)")  "int " // id // "_n_col() \{"
    write (u, "(A)")  "  return 1;"
    write (u, "(A)")  "\}"
    write (u, *)
    write (u, "(A)")  "int " // id // "_n_cf() \{"
    write (u, "(A)")  "  return 1;"
    write (u, "(A)")  "\}"
    write (u, *)
    write (u, "(A)")  "void " // id // "_flv_state( int (*a)[] ) \{"
    write (u, "(A)")  "  static int flv_state[1][3] =  \{ \{ 1, 2, 3 \} \};"
    write (u, "(A)")  "  int j;"
    write (u, "(A)")  "  for (j = 0; j < 3; j++) \{ (*a)[j] &
         &= flv_state[0][j]; \}"
    write (u, "(A)")  "\}"
    write (u, *)
    write (u, "(A)")  "void " // id // "_hel_state( int (*a)[] ) \{"
    write (u, "(A)")  "  static int hel_state[1][3] =  \{ \{ 0, 0, 0 \} \};"
    write (u, "(A)")  "  int j;"
    write (u, "(A)")  "  for (j = 0; j < 3; j++) \{ (*a)[j] &
         &= hel_state[0][j]; \}"
    write (u, "(A)")  "\}"
    write (u, *)
    write (u, "(A)")  "void " // id // "_col_state&
         &( int (*a)[], bool (*g)[] ) \{"
    write (u, "(A)")  "  static int col_state[1][3][2] = &
         &\{ \{ \{0, 0\}, \{0, 0\}, \{0, 0\} \} \};"
    write (u, "(A)")  "  static bool ghost_flag[1][3] =  &
         &\{ \{ false, false, false \} \};"
    write (u, "(A)")  "  int j,k;"
    write (u, "(A)")  "  for (j = 0; j < 3; j++) \{"
    write (u, "(A)")  "    for (k = 0; k < 2; k++) \{"
    write (u, "(A)")  "       (*a)[j*2+k] = col_state[0][j][k];"
    write (u, "(A)")  "    \}"
    write (u, "(A)")  "    (*g)[j] = ghost_flag[0][j];"
    write (u, "(A)")  "  \}"
    write (u, "(A)")  "\}"
    write (u, *)
    if (c_default_complex == c_long_double_complex) then
       write (u, "(A)")  "void " // id // "_color_factors&
            &( int (*cf_index1)[], int (*cf_index2)[], &
            &long double _Complex (*color_factors)[] ) \{"
    else
       write (u, "(A)")  "void " // id // "_color_factors&
            &( int (*cf_index1)[], int (*cf_index2)[], &
            &double _Complex (*color_factors)[] ) \{"
    end if
    write (u, "(A)")  "  (*color_factors)[0] = 1;"
    write (u, "(A)")  "  (*cf_index1)[0] = 1;"
    write (u, "(A)")  "  (*cf_index2)[0] = 1;"
    write (u, "(A)")  "\}"
  end subroutine write_test_me_code_4
    
\eatline
\nwendcode{}\nwbegindocs{363}\nwdocspar
\subsubsection{Test cleanup targets}
Test 7: Repeat test 4 (create, compile, link Fortran module and
driver) and properly clean up all generated files.
\nwenddocs{}\nwbegincode{364}\moddef{Prclib interfaces: execute tests}\plusendmoddef
  call test (prclib_interfaces_7, "prclib_interfaces_7", &
       "cleanup", &
       u, results)
\nwendcode{}\nwbegincode{365}\moddef{Prclib interfaces: tests}\plusendmoddef
  subroutine prclib_interfaces_7 (u)
    integer, intent(in) :: u
    class(prclib_driver_t), allocatable :: driver
    class(prc_writer_t), pointer :: test_writer_4
    type(os_data_t) :: os_data
    integer :: u_file
    character(32), parameter :: md5sum = "1234567890abcdef1234567890abcdef"

    write (u, "(A)")  "* Test output: prclib_interfaces_7"
    write (u, "(A)")  "*   Purpose: compile and link process library"
    write (u, "(A)")  "*            with (fake) matrix-element code &
         &as a Fortran module"
    write (u, "(A)")  "*            then clean up generated files"
    write (u, *)
    write (u, "(A)")  "* Create a prclib driver object (1 process)"

    allocate (test_writer_4_t :: test_writer_4)

    call os_data_init (os_data)
    call dispatch_prclib_driver (driver, var_str ("prclib7"), var_str (""))
    call driver%init (1)
    call driver%set_md5sum (md5sum)
    call driver%set_record (1, var_str ("test7"), var_str ("Test_model"), &
         [var_str ("proc1")], test_writer_4)

    write (u, "(A)")  "* Write makefile"
    u_file = free_unit ()
    open (u_file, file="prclib7.makefile", status="replace", action="write")
    call driver%generate_makefile (u_file, os_data)
    close (u_file)
    
    write (u, "(A)")  "* Write driver source code"
    u_file = free_unit ()
    open (u_file, file="prclib7.f90", status="replace", action="write")
    call driver%generate_driver_code (u_file)
    close (u_file)
    
    write (u, "(A)")  "* Write matrix-element source code"
    call driver%make_source (os_data)

    write (u, "(A)")  "* Compile source code"
    call driver%make_compile (os_data)
    
    write (u, "(A)")  "* Link library"
    call driver%make_link (os_data)

    
    write (u, "(A)")  "* File check"
    write (u, *)
    call check_file (u, "test7.f90")
    call check_file (u, "tpr_test7.mod")
    call check_file (u, "test7.lo")
    call check_file (u, "prclib7.makefile")
    call check_file (u, "prclib7.f90")
    call check_file (u, "prclib7.lo")
    call check_file (u, "prclib7.la")

    write (u, *)
    write (u, "(A)")  "* Delete library"
    write (u, *)
    call driver%clean_library (os_data)
    call check_file (u, "prclib7.la")

    write (u, *)
    write (u, "(A)")  "* Delete object code"
    write (u, *)
    call driver%clean_objects (os_data)
    call check_file (u, "test7.lo")
    call check_file (u, "tpr_test7.mod")
    call check_file (u, "prclib7.lo")

    write (u, *)
    write (u, "(A)")  "* Delete source code"
    write (u, *)
    call driver%clean_source (os_data)
    call check_file (u, "test7.f90")

    write (u, *)
    write (u, "(A)")  "* Delete driver source code"
    write (u, *)
    call driver%clean_driver (os_data)
    call check_file (u, "prclib7.f90")

    write (u, *)
    write (u, "(A)")  "* Delete makefile"
    write (u, *)
    call driver%clean_makefile (os_data)
    call check_file (u, "prclib7.makefile")

    deallocate (test_writer_4)

    write (u, *)
    write (u, "(A)")  "* Test output end: prclib_interfaces_7"
  end subroutine prclib_interfaces_7
  
\eatline
\nwendcode{}\nwbegindocs{366}Auxiliary routine: check and report existence of a file
\nwenddocs{}\nwbegincode{367}\moddef{Prclib interfaces: tests}\plusendmoddef
  subroutine check_file (u, file)
    integer, intent(in) :: u
    character(*), intent(in) :: file
    logical :: exist
    inquire (file=file, exist=exist)
    write (u, "(2x,A,A,L1)")  file, " = ", exist
  end subroutine check_file
  
\eatline
\nwendcode{}\nwbegindocs{368}\nwdocspar
%%%%%%%%%%%%%%%%%%%%%%%%%%%%%%%%%%%%%%%%%%%%%%%%%%%%%%%%%%%%%%%%%%%%%%%%
\section{Abstract process core configuration}
In this module, we define abstract data types that handle the method-specific
part of defining a process (including all of its options) and accessing an
external matrix element.

There are no unit tests, these are deferred to the {\tt{}process{\char95}libraries}
module below. 
\nwenddocs{}\nwbegincode{369}\moddef{\code{}prc{\char95}core{\char95}def.f90\edoc{}}\endmoddef
\LA{}File header\RA{}

module prc_core_def

\LA{}Use strings\RA{}
  use io_units
  use diagnostics

  use process_constants
  use prclib_interfaces
  
\LA{}Standard module head\RA{}

\LA{}Prc core def: public\RA{}

\LA{}Prc core def: types\RA{}

\LA{}Prc core def: interfaces\RA{}

contains

\LA{}Prc core def: procedures\RA{}

end module prc_core_def
\eatline
\nwendcode{}\nwbegindocs{370}\nwdocspar
\subsection{Process core definition type}
For storing configuration data that depend on the specific process
variant, we introduce a polymorphic type.  At this point, we just
declare an abstract base type.  This allows us to defer the
implementation to later modules.

There should be no components that need explicit finalization,
otherwise we would have to call a finalizer from the
{\tt{}process{\char95}component{\char95}def{\char95}t} wrapper.
\nwenddocs{}\nwbegincode{371}\moddef{Prc core def: public}\endmoddef
  public :: prc_core_def_t
\nwendcode{}\nwbegincode{372}\moddef{Prc core def: types}\endmoddef
  type, abstract :: prc_core_def_t
     class(prc_writer_t), allocatable :: writer
   contains
   \LA{}Prc core def: process core def: TBP\RA{}
  end type prc_core_def_t
  
\eatline
\nwendcode{}\nwbegindocs{373}Interfaces for the deferred methods.  

This returns a string.  No passed argument; the string is constant and
depends just on the type.
\nwenddocs{}\nwbegincode{374}\moddef{Prc core def: process core def: TBP}\endmoddef
  procedure (prc_core_def_get_string), nopass, deferred :: type_string
\nwendcode{}\nwbegincode{375}\moddef{Prc core def: interfaces}\endmoddef
  abstract interface
     function prc_core_def_get_string () result (string)
       import
       type(string_t) :: string
     end function prc_core_def_get_string
  end interface
  
\eatline
\nwendcode{}\nwbegindocs{376}\nwdocspar
The {\tt{}write} method should
display the content completely.
\nwenddocs{}\nwbegincode{377}\moddef{Prc core def: process core def: TBP}\plusendmoddef
  procedure (prc_core_def_write), deferred :: write
\nwendcode{}\nwbegincode{378}\moddef{Prc core def: interfaces}\plusendmoddef
  abstract interface
     subroutine prc_core_def_write (object, unit)
       import
       class(prc_core_def_t), intent(in) :: object
       integer, intent(in) :: unit
     end subroutine prc_core_def_write
  end interface

\eatline
\nwendcode{}\nwbegindocs{379}\nwdocspar
The {\tt{}read} method should
fill the content completely.
\nwenddocs{}\nwbegincode{380}\moddef{Prc core def: process core def: TBP}\plusendmoddef
  procedure (prc_core_def_read), deferred :: read
\nwendcode{}\nwbegincode{381}\moddef{Prc core def: interfaces}\plusendmoddef
  abstract interface
     subroutine prc_core_def_read (object, unit)
       import
       class(prc_core_def_t), intent(out) :: object
       integer, intent(in) :: unit
     end subroutine prc_core_def_read
  end interface

\eatline
\nwendcode{}\nwbegindocs{382}This communicates a MD5 checksum to the writer inside the {\tt{}core{\char95}def}
object, if there is any.  Usually, this checksum is not yet known at the time
when the writer is initialized.
\nwenddocs{}\nwbegincode{383}\moddef{Prc core def: process core def: TBP}\plusendmoddef
  procedure :: set_md5sum => prc_core_def_set_md5sum
\nwendcode{}\nwbegincode{384}\moddef{Prc core def: procedures}\endmoddef
  subroutine prc_core_def_set_md5sum (core_def, md5sum)
    class(prc_core_def_t), intent(inout) :: core_def
    character(32) :: md5sum
    if (allocated (core_def%writer))  core_def%writer%md5sum = md5sum
  end subroutine prc_core_def_set_md5sum
  
\eatline
\nwendcode{}\nwbegindocs{385}Allocate an appropriate driver object which corresponds to the
chosen process core definition.

For internal matrix element (i.e., those which do not need external
code), the driver should have access to all matrix element information
from the beginning.  In short, it is the matrix-element code.  

For external matrix elements, the driver will get access to the
external matrix element code.
\nwenddocs{}\nwbegincode{386}\moddef{Prc core def: process core def: TBP}\plusendmoddef
  procedure(prc_core_def_allocate_driver), deferred :: allocate_driver
\nwendcode{}\nwbegincode{387}\moddef{Prc core def: interfaces}\plusendmoddef
  abstract interface
     subroutine prc_core_def_allocate_driver (object, driver, basename)
       import
       class(prc_core_def_t), intent(in) :: object
       class(prc_core_driver_t), intent(out), allocatable :: driver
       type(string_t), intent(in) :: basename
     end subroutine prc_core_def_allocate_driver
  end interface
  
\eatline
\nwendcode{}\nwbegindocs{388}This flag tells whether the particular variant needs external code.
We implement a default function which returns false.  The flag
depends only on the type, therefore we implement it as {\tt{}nopass}.
\nwenddocs{}\nwbegincode{389}\moddef{Prc core def: process core def: TBP}\plusendmoddef
  procedure, nopass :: needs_code => prc_core_def_needs_code
\nwendcode{}\nwbegincode{390}\moddef{Prc core def: procedures}\plusendmoddef
  function prc_core_def_needs_code () result (flag)
    logical :: flag
    flag = .false.
  end function prc_core_def_needs_code
  
\eatline
\nwendcode{}\nwbegindocs{391}This subroutine allocates an array which holds the name of all
features that this process core implements.  This feature
applies to matrix element code that is not coded as a Fortran module
but communicates via independent library functions, which follow the C
calling conventions.  The addresses of those functions are returned as
C function pointers, which can be converted into Fortran procedure
pointers.  The conversion is done in code specific for the process
variant; here we just retrieve the C function pointer.

The array returned here serves the purpose of writing specific
driver code.  The driver interfaces only those C functions which are
supported for the given process core.

If the process core does not require external code, this array is
meaningless.
\nwenddocs{}\nwbegincode{392}\moddef{Prc core def: process core def: TBP}\plusendmoddef
  procedure(prc_core_def_get_features), nopass, deferred &
       :: get_features
\nwendcode{}\nwbegincode{393}\moddef{Prc core def: interfaces}\plusendmoddef
  abstract interface
     subroutine prc_core_def_get_features (features)
       import
       type(string_t), dimension(:), allocatable, intent(out) :: features
     end subroutine prc_core_def_get_features
  end interface

\eatline
\nwendcode{}\nwbegindocs{394}Assign pointers to the process-specific procedures to the driver, if
the process is external.
\nwenddocs{}\nwbegincode{395}\moddef{Prc core def: process core def: TBP}\plusendmoddef
  procedure(prc_core_def_connect), deferred :: connect
\nwendcode{}\nwbegincode{396}\moddef{Prc core def: interfaces}\plusendmoddef
  abstract interface
     subroutine prc_core_def_connect (def, lib_driver, i, proc_driver)
       import
       class(prc_core_def_t), intent(in) :: def
       class(prclib_driver_t), intent(in) :: lib_driver
       integer, intent(in) :: i
       class(prc_core_driver_t), intent(inout) :: proc_driver
     end subroutine prc_core_def_connect
  end interface

\eatline
\nwendcode{}\nwbegindocs{397}\nwdocspar
\subsection{Process core template}
We must be able to automatically allocate a process core definition object
with the appropriate type, given only the type name.

To this end, we introduce a {\tt{}prc{\char95}template{\char95}t} type which is simply a wrapper
for an empty {\tt{}prc{\char95}core{\char95}def{\char95}t} object.  Choosing one of the templates from an
array, we can allocate the target object.
\nwenddocs{}\nwbegincode{398}\moddef{Prc core def: public}\plusendmoddef
  public :: prc_template_t
\nwendcode{}\nwbegincode{399}\moddef{Prc core def: types}\plusendmoddef
  type :: prc_template_t
     class(prc_core_def_t), allocatable :: core_def
  end type prc_template_t
  
\eatline
\nwendcode{}\nwbegindocs{400}The allocation routine.  We use the {\tt{}source} option of the {\tt{}allocate}
statement.   The {\tt{}mold} option would probably more appropriate, but is a
F2008 feature.
\nwenddocs{}\nwbegincode{401}\moddef{Prc core def: public}\plusendmoddef
  public :: allocate_core_def
\nwendcode{}\nwbegincode{402}\moddef{Prc core def: procedures}\plusendmoddef
  subroutine allocate_core_def (template, name, core_def)
    type(prc_template_t), dimension(:), intent(in) :: template
    type(string_t), intent(in) :: name
    class(prc_core_def_t), allocatable :: core_def
    integer :: i
    do i = 1, size (template)
       if (template(i)%core_def%type_string () == name) then
          allocate (core_def, source = template(i)%core_def)
          return
       end if
    end do
  end subroutine allocate_core_def
  
\eatline
\nwendcode{}\nwbegindocs{403}\nwdocspar
\subsection{Process driver}
For each process component, we implement a driver object which holds
the calls to the matrix element and various auxiliary routines as
procedure pointers.  Any actual calculation will use this object to
communicate with the process.

Depending on the type of process (as described by a corresponding
{\tt{}prc{\char95}core{\char95}def} object), the procedure pointers may refer to
external or internal code, and there may be additional procedures for
certain types.  The base type defined here is abstract.
\nwenddocs{}\nwbegincode{404}\moddef{Prc core def: public}\plusendmoddef
  public :: prc_core_driver_t
\nwendcode{}\nwbegincode{405}\moddef{Prc core def: types}\plusendmoddef
  type, abstract :: prc_core_driver_t
   contains
   \LA{}Prc core def: process driver: TBP\RA{}
  end type prc_core_driver_t
  
\eatline
\nwendcode{}\nwbegindocs{406}This returns the process type.  No reference to contents.
\nwenddocs{}\nwbegincode{407}\moddef{Prc core def: process driver: TBP}\endmoddef
  procedure(prc_core_driver_type_name), nopass, deferred :: type_name
\nwendcode{}\nwbegincode{408}\moddef{Prc core def: interfaces}\plusendmoddef
  abstract interface
     function prc_core_driver_type_name () result (type)
       import
       type(string_t) :: type
     end function prc_core_driver_type_name
  end interface
  
\eatline
\nwendcode{}\nwbegindocs{409}\nwdocspar
\subsection{Process driver for intrinsic process}
This is an abstract extension for the driver type.  It has one
additional method, namely a subroutine that fills the record of
constant process data.  For an external process, this task is
performed by the external library driver instead.
\nwenddocs{}\nwbegincode{410}\moddef{Prc core def: public}\plusendmoddef
  public :: process_driver_internal_t
\nwendcode{}\nwbegincode{411}\moddef{Prc core def: types}\plusendmoddef
  type, extends (prc_core_driver_t), abstract :: process_driver_internal_t
   contains
   \LA{}Prc core def: process driver internal: TBP\RA{}
  end type process_driver_internal_t
  
\eatline
\nwendcode{}\nwbegincode{412}\moddef{Prc core def: process driver internal: TBP}\endmoddef
  procedure(process_driver_fill_constants), deferred :: fill_constants
\nwendcode{}\nwbegincode{413}\moddef{Prc core def: interfaces}\plusendmoddef
  abstract interface
     subroutine process_driver_fill_constants (driver, data)
       import
       class(process_driver_internal_t), intent(in) :: driver
       type(process_constants_t), intent(out) :: data
     end subroutine process_driver_fill_constants
  end interface
  
\eatline
\nwendcode{}\nwbegindocs{414}\nwdocspar
%%%%%%%%%%%%%%%%%%%%%%%%%%%%%%%%%%%%%%%%%%%%%%%%%%%%%%%%%%%%%%%%%%%%%%%%
\section{Process library access}
\label{sec:process_libraries}
Processes (the code and data that are necessary for evaluating matrix
elements of a particular process or process component) are organized
in process libraries.  In full form, process libraries contain
generated and dynamically compiled and linked code, so they are actual
libraries on the OS level.  Alternatively, there may be simple
processes that can be generated without referring to external
libraries, and external libraries that are just linked in.

This module interfaces the OS to create, build, and use process
libraries.

We work with two related data structures.  There is the list of
process configurations that stores the user input and data derived
from it.  A given process configuration list is scanned for creating a
process library, which consists of both data and code.  The creation
step involves calling external programs and incorporating external
code.

For the subsequent integration and event generation steps, we read the
process library.  We also support partial (re)creation of the process
library.  To this end, we should be able to reconstruct the
configuration data records from the process library.
\nwenddocs{}\nwbegincode{415}\moddef{\code{}process{\char95}libraries.f90\edoc{}}\endmoddef
\LA{}File header\RA{}

module process_libraries

  use iso_c_binding !NODEP!
  use kinds
\LA{}Use strings\RA{}
  use io_units
  use unit_tests
  use diagnostics
  use md5
  use os_interface
  use lexers
  use model_data
  use particle_specifiers
  use process_constants
  use prclib_interfaces
  use prc_core_def

\LA{}Standard module head\RA{}

\LA{}Process libraries: public\RA{}

\LA{}Process libraries: parameters\RA{}

\LA{}Process libraries: types\RA{}

\LA{}Process libraries: interfaces\RA{}

\LA{}Process libraries: test declarations\RA{}

contains

\LA{}Process libraries: procedures\RA{}

\LA{}Process libraries: tests\RA{}

end module process_libraries
\eatline
\nwendcode{}\nwbegindocs{416}\nwdocspar
\subsection{Auxiliary stuff}
Here is a small subroutine that strips the left-hand side and the
equals sign off an equation.
\nwenddocs{}\nwbegincode{417}\moddef{Process libraries: procedures}\endmoddef
  subroutine strip_prefix (buffer)
    character(*), intent(inout) :: buffer
    type(string_t) :: string, prefix
    string = buffer
    call split (string, prefix, "=")
    buffer = string
  end subroutine strip_prefix
    
\eatline
\nwendcode{}\nwbegindocs{418}\nwdocspar
\subsection{Process definition objects}
We collect process configuration data in a derived type,
{\tt{}process{\char95}def{\char95}t}.  A process can be a collection of several
components which are treated as a single entity for the purpose of
observables and event generation.  Multiple process components may
initially be defined by the user.  The system may add additional
components, e.g., subtraction terms.  The common data type is
{\tt{}process{\char95}component{\char95}def{\char95}t}.  Within each component, there are several
universal data items, and a part which depend on the particular
process variant.  The latter is covered by an abstract type
{\tt{}prc{\char95}core{\char95}def{\char95}t} and its extensions.

\nwenddocs{}\nwbegindocs{419}\nwdocspar
\subsubsection{Wrapper for components}
We define a wrapper type for the configuration of individual
components.

The string {\tt{}basename} is used for building file, module, and
function names for the current process component.  Initially, it will
be built from the corresponding process basename by appending an
alphanumeric suffix.

The logical {\tt{}initial} tells whether this is a user-defined (true) or
system-generated (false) configuration.

The numbers {\tt{}n{\char95}in}, {\tt{}n{\char95}out}, and {\tt{}n{\char95}tot} denote the incoming,
outgoing and total number of particles (partons) participating in the
process component, respectively.  These are the nominal particles, as
input by the user (recombination may change the particle content, for
the output events).

The string arrays {\tt{}prt{\char95}in} and {\tt{}prt{\char95}out} hold the particle
specifications as provided by the user.  For a system-generated
process component, they remain deallocated.

The {\tt{}method} string is used to determine the type of process matrix
element and how it is obtained.

The {\tt{}description} string collects the information about particle
content and method in a single human-readable string.

The pointer object {\tt{}core{\char95}def} is allocated according to the
actual process variant, which depends on the method.  The subobject
holds any additional configuration data that is relevant for the
process component.

We assume that no finalizer is needed.
\nwenddocs{}\nwbegincode{420}\moddef{Process libraries: public}\endmoddef
  public :: process_component_def_t
\nwendcode{}\nwbegincode{421}\moddef{Process libraries: types}\endmoddef
  type :: process_component_def_t
     private
     type(string_t) :: basename
     logical :: initial = .false.
     integer :: n_in = 0
     integer :: n_out = 0
     integer :: n_tot = 0
     type(prt_spec_t), dimension(:), allocatable :: prt_in
     type(prt_spec_t), dimension(:), allocatable :: prt_out
     type(string_t) :: method
     type(string_t) :: description
     class(prc_core_def_t), allocatable :: core_def
     character(32) :: md5sum = ""
     type(string_t) :: nlo_type
     integer :: associated_born = 0
     integer :: associated_real = 0
     integer :: associated_virt = 0
     integer :: associated_sub = 0
     logical :: active_nlo_component
   contains
   \LA{}Process libraries: process component def: TBP\RA{}
  end type process_component_def_t
  
\eatline
\nwendcode{}\nwbegindocs{422}Display the complete content.
\nwenddocs{}\nwbegincode{423}\moddef{Process libraries: process component def: TBP}\endmoddef
  procedure :: write => process_component_def_write
\nwendcode{}\nwbegincode{424}\moddef{Process libraries: procedures}\plusendmoddef
  subroutine process_component_def_write (object, unit)
    class(process_component_def_t), intent(in) :: object
    integer, intent(in), optional :: unit
    integer :: u
    u = given_output_unit (unit)
    write (u, "(3x,A,A)")  "Component ID        = ", char (object%basename)
    write (u, "(3x,A,L1)") "Initial component   = ", object%initial
    write (u, "(3x,A,I0,1x,I0,1x,I0)") "N (in, out, tot)    = ", &
         object%n_in, object%n_out, object%n_tot
    write (u, "(3x,A)", advance="no") "Particle content    = "
    if (allocated (object%prt_in)) then
       call prt_spec_write (object%prt_in, u, advance="no")
    else
       write (u, "(A)", advance="no")  "[undefined]"
    end if
    write (u, "(A)", advance="no") " => "
    if (allocated (object%prt_out)) then
       call prt_spec_write (object%prt_out, u, advance="no")
    else
       write (u, "(A)", advance="no")  "[undefined]"
    end if
    write (u, "(A)")
    if (object%method /= "") then
       write (u, "(3x,A,A)")  "Method              = ", &
            char (object%method)
    else
       write (u, "(3x,A)")  "Method              = [undefined]"
    end if
    if (allocated (object%core_def)) then
       write (u, "(3x,A,A)")  "Process variant     = ", &
            char (object%core_def%type_string ())
       call object%core_def%write (u)
    else
       write (u, "(3x,A)")  "Process variant     = [undefined]"
    end if
    write (u, "(3x,A,A,A)") "MD5 sum (def)       = '", object%md5sum, "'"
  end subroutine process_component_def_write
    
\eatline
\nwendcode{}\nwbegindocs{425}Read the process component definition.  Allocate the process variant
definition with appropriate type, matching the type name on file with
the provided templates.
\nwenddocs{}\nwbegincode{426}\moddef{Process libraries: process component def: TBP}\plusendmoddef
  procedure :: read => process_component_def_read
\nwendcode{}\nwbegincode{427}\moddef{Process libraries: procedures}\plusendmoddef
  subroutine process_component_def_read (component, unit, core_def_templates)
    class(process_component_def_t), intent(out) :: component
    integer, intent(in) :: unit
    type(prc_template_t), dimension(:), intent(in) :: core_def_templates
    character(80) :: buffer
    type(string_t) :: var_buffer, prefix, in_state, out_state
    type(string_t) :: variant_type
    
    read (unit, "(A)")  buffer
    call strip_prefix (buffer)
    component%basename = trim (adjustl (buffer))
    
    read (unit, "(A)")  buffer
    call strip_prefix (buffer)
    read (buffer, *)  component%initial
    
    read (unit, "(A)")  buffer
    call strip_prefix (buffer)
    read (buffer, *)  component%n_in, component%n_out, component%n_tot
    
    call get (unit, var_buffer)
    call split (var_buffer, prefix, "=")   ! keeps 'in => out'
    call split (var_buffer, prefix, "=")   ! actually: separator is '=>'

    in_state = prefix
    if (component%n_in > 0) then
       call prt_spec_read (component%prt_in, in_state)
    end if

    out_state = extract (var_buffer, 2)
    if (component%n_out > 0) then
       call prt_spec_read (component%prt_out, out_state)
    end if
    
    read (unit, "(A)")  buffer
    call strip_prefix (buffer)
    component%method = trim (adjustl (buffer))
    if (component%method == "[undefined]") &
         component%method = ""
    
    read (unit, "(A)")  buffer
    call strip_prefix (buffer)
    variant_type = trim (adjustl (buffer))
    call allocate_core_def &
         (core_def_templates, variant_type, component%core_def)
    if (allocated (component%core_def)) then
       call component%core_def%read (unit)
    end if

    read (unit, "(A)")  buffer
    call strip_prefix (buffer)
    read (buffer(3:34), "(A32)")  component%md5sum
    
  end subroutine process_component_def_read
    
\eatline
\nwendcode{}\nwbegindocs{428}Short account.
\nwenddocs{}\nwbegincode{429}\moddef{Process libraries: process component def: TBP}\plusendmoddef
  procedure :: show => process_component_def_show
\nwendcode{}\nwbegincode{430}\moddef{Process libraries: procedures}\plusendmoddef
  subroutine process_component_def_show (object, unit)
    class(process_component_def_t), intent(in) :: object
    integer, intent(in), optional :: unit
    integer :: u
    u = given_output_unit (unit)
    write (u, "(6x,A)", advance="no")  char (object%basename)
    if (.not. object%initial) &
         write (u, "('*')", advance="no")
    write (u, "(':',1x)", advance="no")
    if (allocated (object%prt_in)) then
       call prt_spec_write (object%prt_in, u, advance="no")
    else
       write (u, "(A)", advance="no")  "[undefined]"
    end if
    write (u, "(A)", advance="no") " => "
    if (allocated (object%prt_out)) then
       call prt_spec_write (object%prt_out, u, advance="no")
    else
       write (u, "(A)", advance="no")  "[undefined]"
    end if
    if (object%method /= "") then
       write (u, "(2x,'[',A,']')")  char (object%method)
    else
       write (u, *)
    end if
  end subroutine process_component_def_show
    
\eatline
\nwendcode{}\nwbegindocs{431}Compute the MD5 sum of a process component.  We reset the stored MD5
sum to the empty string (so a previous value is not included in the
calculation), the write a temporary file and calculate the MD5 sum of
that file.

This implies that all data that are displayed by the {\tt{}write} method
become part of the MD5 sum calculation.

The {\tt{}model} is not part of the object, but must be included in the MD5 sum.
Otherwise, modifying the model and nothing else would not trigger remaking the
process-component source.  Note that the model parameters may change later and
therefore are not incorporated.

After the MD5 sum of the component has been computed, we communicate it to the
{\tt{}writer} subobject of the specific {\tt{}core{\char95}def} component.  Although these
types are abstract, the MD5-related features are valid for the abstract
types.
\nwenddocs{}\nwbegincode{432}\moddef{Process libraries: process component def: TBP}\plusendmoddef
  procedure :: compute_md5sum => process_component_def_compute_md5sum
\nwendcode{}\nwbegincode{433}\moddef{Process libraries: procedures}\plusendmoddef
  subroutine process_component_def_compute_md5sum (component, model)
    class(process_component_def_t), intent(inout) :: component
    class(model_data_t), intent(in), optional, target :: model
    integer :: u
    component%md5sum = ""
    u = free_unit ()
    open (u, status = "scratch", action = "readwrite")
    if (present (model))  write (u, "(A32)")  model%get_md5sum ()
    call component%write (u)
    rewind (u)
    component%md5sum = md5sum (u)
    close (u)
    if (allocated (component%core_def)) then
       call component%core_def%set_md5sum (component%md5sum)
    end if
  end subroutine process_component_def_compute_md5sum

\eatline
\nwendcode{}\nwbegindocs{434}Allocate the process driver (with a suitable type) for a process
component.  For internal processes, we may set all data already at
this stage.
\nwenddocs{}\nwbegincode{435}\moddef{Process libraries: process component def: TBP}\plusendmoddef
  procedure :: allocate_driver => process_component_def_allocate_driver
\nwendcode{}\nwbegincode{436}\moddef{Process libraries: procedures}\plusendmoddef
  subroutine process_component_def_allocate_driver (component, driver)
    class(process_component_def_t), intent(in) :: component
    class(prc_core_driver_t), intent(out), allocatable :: driver
    if (allocated (component%core_def)) then
       call component%core_def%allocate_driver (driver, component%basename)
    end if
  end subroutine process_component_def_allocate_driver
    
\eatline
\nwendcode{}\nwbegindocs{437}Tell whether the process core needs external code.
\nwenddocs{}\nwbegincode{438}\moddef{Process libraries: process component def: TBP}\plusendmoddef
  procedure :: needs_code => process_component_def_needs_code
\nwendcode{}\nwbegincode{439}\moddef{Process libraries: procedures}\plusendmoddef
  function process_component_def_needs_code (component) result (flag)
    class(process_component_def_t), intent(in) :: component
    logical :: flag
    flag = component%core_def%needs_code ()
  end function process_component_def_needs_code
  
\eatline
\nwendcode{}\nwbegindocs{440}If there is external code, the {\tt{}core{\char95}def} subobject should
provide a writer object.  This method returns a pointer to the writer.
\nwenddocs{}\nwbegincode{441}\moddef{Process libraries: process component def: TBP}\plusendmoddef
  procedure :: get_writer_ptr => process_component_def_get_writer_ptr
\nwendcode{}\nwbegincode{442}\moddef{Process libraries: procedures}\plusendmoddef
  function process_component_def_get_writer_ptr (component) result (writer)
    class(process_component_def_t), intent(in), target :: component
    class(prc_writer_t), pointer :: writer
    writer => component%core_def%writer
  end function process_component_def_get_writer_ptr
  
\eatline
\nwendcode{}\nwbegindocs{443}Return an array which holds the names of all C functions that this
process component implements.
\nwenddocs{}\nwbegincode{444}\moddef{Process libraries: process component def: TBP}\plusendmoddef
  procedure :: get_features => process_component_def_get_features
\nwendcode{}\nwbegincode{445}\moddef{Process libraries: procedures}\plusendmoddef
  function process_component_def_get_features (component) result (features)
    class(process_component_def_t), intent(in) :: component
    type(string_t), dimension(:), allocatable :: features
    call component%core_def%get_features (features)
  end function process_component_def_get_features

\eatline
\nwendcode{}\nwbegindocs{446}Assign procedure pointers in the {\tt{}driver} component (external
processes).  For internal processes, this is meaningless.
\nwenddocs{}\nwbegincode{447}\moddef{Process libraries: process component def: TBP}\plusendmoddef
  procedure :: connect => process_component_def_connect
\nwendcode{}\nwbegincode{448}\moddef{Process libraries: procedures}\plusendmoddef
  subroutine process_component_def_connect &
       (component, lib_driver, i, proc_driver)
    class(process_component_def_t), intent(in) :: component
    class(prclib_driver_t), intent(in) :: lib_driver
    integer, intent(in) :: i
    class(prc_core_driver_t), intent(inout) :: proc_driver
    select type (proc_driver)
    class is (process_driver_internal_t)
       ! nothing to do
    class default
       call component%core_def%connect (lib_driver, i, proc_driver)
    end select
  end subroutine process_component_def_connect

\eatline
\nwendcode{}\nwbegindocs{449}Return a pointer to the process core definition, which is of
abstract type.
\nwenddocs{}\nwbegincode{450}\moddef{Process libraries: process component def: TBP}\plusendmoddef
  procedure :: get_core_def_ptr => process_component_get_core_def_ptr
\nwendcode{}\nwbegincode{451}\moddef{Process libraries: procedures}\plusendmoddef
  function process_component_get_core_def_ptr (component) result (ptr)
    class(process_component_def_t), intent(in), target :: component
    class(prc_core_def_t), pointer :: ptr
    ptr => component%core_def
  end function process_component_get_core_def_ptr
  
\eatline
\nwendcode{}\nwbegindocs{452}Return nominal particle counts, as input by the user.
\nwenddocs{}\nwbegincode{453}\moddef{Process libraries: process component def: TBP}\plusendmoddef
  procedure :: get_n_in  => process_component_def_get_n_in
  procedure :: get_n_out => process_component_def_get_n_out
  procedure :: get_n_tot => process_component_def_get_n_tot
\nwendcode{}\nwbegincode{454}\moddef{Process libraries: procedures}\plusendmoddef
  function process_component_def_get_n_in (component) result (n_in)
    class(process_component_def_t), intent(in) :: component
    integer :: n_in
    n_in = component%n_in
  end function process_component_def_get_n_in
  
  function process_component_def_get_n_out (component) result (n_out)
    class(process_component_def_t), intent(in) :: component
    integer :: n_out
    n_out = component%n_out
  end function process_component_def_get_n_out
  
  function process_component_def_get_n_tot (component) result (n_tot)
    class(process_component_def_t), intent(in) :: component
    integer :: n_tot
    n_tot = component%n_tot
  end function process_component_def_get_n_tot
  
\eatline
\eatline
\eatline
\nwendcode{}\nwbegindocs{455}Allocate and return string arrays for the incoming and outgoing particles.
\nwenddocs{}\nwbegincode{456}\moddef{Process libraries: process component def: TBP}\plusendmoddef
  procedure :: get_prt_in => process_component_def_get_prt_in
  procedure :: get_prt_out => process_component_def_get_prt_out
\nwendcode{}\nwbegincode{457}\moddef{Process libraries: procedures}\plusendmoddef
  subroutine process_component_def_get_prt_in (component, prt)
    class(process_component_def_t), intent(in) :: component
    type(string_t), dimension(:), intent(out), allocatable :: prt
    integer :: i
    allocate (prt (component%n_in))
    do i = 1, component%n_in
       prt(i) = component%prt_in(i)%to_string ()
    end do
  end subroutine process_component_def_get_prt_in
  
  subroutine process_component_def_get_prt_out (component, prt)
    class(process_component_def_t), intent(in) :: component
    type(string_t), dimension(:), intent(out), allocatable :: prt
    integer :: i
    allocate (prt (component%n_out))
    do i = 1, component%n_out
       prt(i) = component%prt_out(i)%to_string ()
    end do
  end subroutine process_component_def_get_prt_out
  
\eatline
\eatline
\nwendcode{}\nwbegindocs{458}Return the MD5 sum.
\nwenddocs{}\nwbegincode{459}\moddef{Process libraries: process component def: TBP}\plusendmoddef
  procedure :: get_md5sum => process_component_def_get_md5sum
\nwendcode{}\nwbegincode{460}\moddef{Process libraries: procedures}\plusendmoddef
  function process_component_def_get_md5sum (component) result (md5sum)
    class(process_component_def_t), intent(in) :: component
    character(32) :: md5sum
    md5sum = component%md5sum
  end function process_component_def_get_md5sum
  
\eatline
\nwendcode{}\nwbegindocs{461}Get NLO data
\nwenddocs{}\nwbegincode{462}\moddef{Process libraries: process component def: TBP}\plusendmoddef
  procedure :: get_nlo_type => process_component_def_get_nlo_type
  procedure :: get_associated_born &
               => process_component_def_get_associated_born
  procedure :: get_association_list &
               => process_component_def_get_association_list
  procedure :: is_active_nlo_component &
               => process_component_def_is_active_nlo_component
\nwendcode{}\nwbegincode{463}\moddef{Process libraries: procedures}\plusendmoddef
  function process_component_def_get_nlo_type (component) result (nlo_type)
    class(process_component_def_t), intent(in) :: component
    type(string_t) :: nlo_type
    if (component%nlo_type == "") then
      nlo_type = 'Born'
    else
      nlo_type = component%nlo_type
    end if
  end function process_component_def_get_nlo_type

  function process_component_def_get_associated_born (component) result (i_born)
    class(process_component_def_t), intent(in) :: component
    integer :: i_born
    i_born = component%associated_born
  end function process_component_def_get_associated_born

  function process_component_def_is_active_nlo_component (component) result (active)
    class(process_component_def_t), intent(in) :: component
    logical :: active
    active = component%active_nlo_component
  end function process_component_def_is_active_nlo_component

  function process_component_def_get_association_list (component) result (list)
    class(process_component_def_t), intent(in) :: component
    integer, dimension(4) :: list
    list(1) = component%associated_born
    list(2) = component%associated_real
    list(3) = component%associated_virt
    list(4) = component%associated_sub
  end function process_component_def_get_association_list

\eatline
\eatline
\eatline
\nwendcode{}\nwbegindocs{464}\nwdocspar
\subsubsection{Process definition}
The process component definitions are collected in a common process
definition object.

The {\tt{}id} is the ID string that the user has provided for identifying
this process.  It must be a string that is allowed as part of a
Fortran variable name, since it may be used for generating code.

The number {\tt{}n{\char95}in} is 1 or 2 for a decay or scattering process,
respectively.  This must be identical to {\tt{}n{\char95}in} for all components.

The initial and extra component definitions (see above) are allocated as the
{\tt{}initial} and {\tt{}extra} arrays, respectively.  The latter
are determined from the former.

The {\tt{}md5sum} is used to verify the integrity of the configuration.
\nwenddocs{}\nwbegincode{465}\moddef{Process libraries: public}\plusendmoddef
  public :: process_def_t
\nwendcode{}\nwbegincode{466}\moddef{Process libraries: types}\plusendmoddef
  type :: process_def_t
     private
     type(string_t) :: id
     integer :: num_id = 0
     class(model_data_t), pointer :: model => null ()
     type(string_t) :: model_name
     integer :: n_in  = 0
     integer :: n_initial = 0
     integer :: n_extra = 0
     type(process_component_def_t), dimension(:), allocatable :: initial
     type(process_component_def_t), dimension(:), allocatable :: extra
     character(32) :: md5sum = ""
     logical :: nlo_process = .false.
     logical :: combined_nlo_integration = .false.
   contains
   \LA{}Process libraries: process def: TBP\RA{}
  end type process_def_t

\eatline
\nwendcode{}\nwbegindocs{467}Write the process definition including components:
\nwenddocs{}\nwbegincode{468}\moddef{Process libraries: process def: TBP}\endmoddef
  procedure :: write => process_def_write
\nwendcode{}\nwbegincode{469}\moddef{Process libraries: procedures}\plusendmoddef
  subroutine process_def_write (object, unit)
    class(process_def_t), intent(in) :: object
    integer, intent(in) :: unit
    integer :: i
    write (unit, "(1x,A,A,A)") "ID = '", char (object%id), "'"
    if (object%num_id /= 0) &
         write (unit, "(1x,A,I0)")  "ID(num) = ", object%num_id
    select case (object%n_in)
    case (1);  write (unit, "(1x,A)")  "Decay"
    case (2);  write (unit, "(1x,A)")  "Scattering"
    case default
       write (unit, "(1x,A)")  "[Undefined process]"
       return
    end select
    if (object%model_name /= "") then
       write (unit, "(1x,A,A)")  "Model = ", char (object%model_name)
    else
       write (unit, "(1x,A)")  "Model = [undefined]"
    end if
    write (unit, "(1x,A,I0)")  "Initially defined component(s) = ", &
         object%n_initial
    write (unit, "(1x,A,I0)")  "Extra generated component(s)   = ", &
         object%n_extra
    write (unit, "(1x,A,A,A)") "MD5 sum   = '", object%md5sum, "'"
    if (allocated (object%initial)) then
       do i = 1, size (object%initial)
          write (unit, "(1x,A,I0)")  "Component #", i
          call object%initial(i)%write (unit)
       end do
    end if
    if (allocated (object%extra)) then
       do i = 1, size (object%extra)
          write (unit, "(1x,A,I0)")  "Component #", object%n_initial + i
          call object%extra(i)%write (unit)
       end do
    end if
  end subroutine process_def_write
    
\eatline
\nwendcode{}\nwbegindocs{470}Read the process definition including components.
\nwenddocs{}\nwbegincode{471}\moddef{Process libraries: process def: TBP}\plusendmoddef
  procedure :: read => process_def_read
\nwendcode{}\nwbegincode{472}\moddef{Process libraries: procedures}\plusendmoddef
  subroutine process_def_read (object, unit, core_def_templates)
    class(process_def_t), intent(out) :: object
    integer, intent(in) :: unit
    type(prc_template_t), dimension(:), intent(in) :: core_def_templates
    integer :: i, i1, i2
    character(80) :: buffer, ref
    read (unit, "(A)")  buffer
    call strip_prefix (buffer)
    i1 = scan (buffer, "'")
    i2 = scan (buffer, "'", back=.true.)
    if (i2 > i1) then
       object%id = buffer(i1+1:i2-1)
    else
       object%id = ""
    end if

    read (unit, "(A)")  buffer
    select case (buffer(2:11))
    case ("Decay     "); object%n_in = 1
    case ("Scattering"); object%n_in = 2
    case default
       return
    end select

    read (unit, "(A)")  buffer
    call strip_prefix (buffer)
    object%model_name = trim (adjustl (buffer))
    if (object%model_name == "[undefined]")  object%model_name = ""
          
    read (unit, "(A)")  buffer
    call strip_prefix (buffer)
    read (buffer, *)  object%n_initial

    read (unit, "(A)")  buffer
    call strip_prefix (buffer)
    read (buffer, *)  object%n_extra

    read (unit, "(A)")  buffer
    call strip_prefix (buffer)
    read (buffer(3:34), "(A32)")  object%md5sum
    
    if (object%n_initial > 0) then
       allocate (object%initial (object%n_initial))
       do i = 1, object%n_initial
          read (unit, "(A)")  buffer
          write (ref, "(1x,A,I0)")  "Component #", i
          if (buffer /= ref)  return                ! Wrong component header
          call object%initial(i)%read (unit, core_def_templates)
       end do
    end if
    
  end subroutine process_def_read
    
\eatline
\nwendcode{}\nwbegindocs{473}Short account.
\nwenddocs{}\nwbegincode{474}\moddef{Process libraries: process def: TBP}\plusendmoddef
  procedure :: show => process_def_show
\nwendcode{}\nwbegincode{475}\moddef{Process libraries: procedures}\plusendmoddef
  subroutine process_def_show (object, unit)
    class(process_def_t), intent(in) :: object
    integer, intent(in) :: unit
    integer :: i
    write (unit, "(4x,A)", advance="no") char (object%id)
    if (object%num_id /= 0) &
         write (unit, "(1x,'(',I0,')')", advance="no")  object%num_id
    if (object%model_name /= "") &
         write (unit, "(1x,'[',A,']')")  char (object%model_name)
    if (allocated (object%initial)) then
       do i = 1, size (object%initial)
          call object%initial(i)%show (unit)
       end do
    end if
    if (allocated (object%extra)) then
       do i = 1, size (object%extra)
          call object%extra(i)%show (unit)
       end do
    end if
  end subroutine process_def_show
    
\eatline
\nwendcode{}\nwbegindocs{476}Initialize an entry (initialize the process definition inside).  We
allocate the 'initial' set of components.  Extra components remain
unallocated.

The model should be present as a pointer.  This allows us to retrieve the
model's MD5 sum.  However, for various tests it is sufficient to have the
name.

We create the basenames for the process components by appending a
suffix which we increment for each component.
\nwenddocs{}\nwbegincode{477}\moddef{Process libraries: process def: TBP}\plusendmoddef
  procedure :: init => process_def_init
\nwendcode{}\nwbegincode{478}\moddef{Process libraries: procedures}\plusendmoddef
  subroutine process_def_init (def, id, &
       model, model_name, n_in, n_components, num_id, nlo_process)
    class(process_def_t), intent(out) :: def
    type(string_t), intent(in), optional :: id
    class(model_data_t), intent(in), optional, target :: model
    type(string_t), intent(in), optional :: model_name
    integer, intent(in), optional :: n_in
    integer, intent(in), optional :: n_components
    integer, intent(in), optional :: num_id
    logical, intent(in), optional :: nlo_process
    character(16) :: suffix
    integer :: i
    if (present (id)) then
       def%id = id
    else
       def%id = ""
    end if
    if (present (num_id)) then
       def%num_id = num_id
    end if
    if (present (model)) then
       def%model => model
       def%model_name = model%get_name ()
    else
       def%model => null ()
       if (present (model_name)) then
          def%model_name = model_name
       else
          def%model_name = ""
       end if
    end if
    if (present (n_in))  def%n_in = n_in
    if (present (n_components)) then
       def%n_initial = n_components
       allocate (def%initial (n_components))
    end if
    if (present (nlo_process)) def%nlo_process = nlo_process
    def%initial%initial = .true.
    def%initial%method     = ""
    do i = 1, def%n_initial
       write (suffix, "(A,I0)")  "_i", i
       def%initial(i)%basename = def%id // trim (suffix)
    end do
    def%initial%description = ""
  end subroutine process_def_init
  
\eatline
\nwendcode{}\nwbegindocs{479}Initialize an initial component.  The particle content
must be specified.  The process core block is not (yet) allocated.

We assume that the particle arrays match the {\tt{}n{\char95}in} and
{\tt{}n{\char95}out} values in size.  The model is referred to by name; it is
identified as an existing model later.  The index {\tt{}i} must refer to
an existing element of the component array.

Data specific for the process core of a component are imported as
the {\tt{}core{\char95}def} argument.  We should allocate an object of class
{\tt{}prc{\char95}core{\char95}def{\char95}t} with the appropriate specific type, fill it,
and transfer it to the process component definition here.  The
allocation is moved, so the original allocated object is returned empty.
\nwenddocs{}\nwbegincode{480}\moddef{Process libraries: process def: TBP}\plusendmoddef
  procedure :: import_component => process_def_import_component
\nwendcode{}\nwbegincode{481}\moddef{Process libraries: procedures}\plusendmoddef
  subroutine process_def_import_component (def, &
       i, n_out, prt_in, prt_out, method, variant, &
       nlo_type, active)
    class(process_def_t), intent(inout) :: def
    integer, intent(in) :: i
    integer, intent(in), optional :: n_out
    type(prt_spec_t), dimension(:), intent(in), optional :: prt_in
    type(prt_spec_t), dimension(:), intent(in), optional :: prt_out
    type(string_t), intent(in), optional :: method
    type(string_t), intent(in), optional :: nlo_type
    logical, intent(in), optional :: active
    class(prc_core_def_t), &
         intent(inout), allocatable, optional :: variant
    integer :: p
    associate (comp => def%initial(i))
      if (present (n_out)) then
         comp%n_in  = def%n_in
         comp%n_out = n_out
         comp%n_tot = def%n_in + n_out
      end if
      if (present (prt_in)) then
         allocate (comp%prt_in (size (prt_in)))
         comp%prt_in = prt_in
      end if
      if (present (prt_out)) then
         allocate (comp%prt_out (size (prt_out)))
         comp%prt_out = prt_out
      end if
      if (present (method))  comp%method = method
      if (present (variant)) then
         call move_alloc (variant, comp%core_def)
      end if
      if (present (nlo_type)) then
        comp%nlo_type = nlo_type
      end if
      if (present (active)) then
         comp%active_nlo_component = active
      else
         comp%active_nlo_component = .true.
      end if
      if (allocated (comp%prt_in) .and. allocated (comp%prt_out)) then
         associate (d => comp%description)
           d = ""
           do p = 1, size (prt_in)
              if (p > 1)  d = d // ", "
              d = d // comp%prt_in(p)%to_string ()
           end do
           d = d // " => "
           do p = 1, size (prt_out)
              if (p > 1)  d = d // ", "
              d = d // comp%prt_out(p)%to_string ()
           end do
           if (comp%method /= "") then
              d = d // " [" // comp%method // "]"
           end if
           if (comp%nlo_type /= "") then
             d = d // ", [" // comp%nlo_type // "]"
           end if
         end associate
      end if
    end associate
  end subroutine process_def_import_component

\eatline
\nwendcode{}\nwbegindocs{482}\nwdocspar
\nwenddocs{}\nwbegincode{483}\moddef{Process libraries: process def: TBP}\plusendmoddef
  procedure :: get_n_components => process_def_get_n_components
\nwendcode{}\nwbegincode{484}\moddef{Process libraries: procedures}\plusendmoddef
  function process_def_get_n_components (def) result (n)
    class(process_def_t), intent(in) :: def
    integer :: n
    n = size (def%initial)
  end function process_def_get_n_components

\eatline
\nwendcode{}\nwbegindocs{485}\nwdocspar
\nwenddocs{}\nwbegincode{486}\moddef{Process libraries: process def: TBP}\plusendmoddef
  procedure :: set_associated_components => &
                   process_def_set_associated_components
\nwendcode{}\nwbegincode{487}\moddef{Process libraries: procedures}\plusendmoddef
  subroutine process_def_set_associated_components (def, i, &
                     i_born, i_real, i_virt, i_sub)
    class(process_def_t), intent(inout) :: def
    integer, intent(in) :: i
    integer, intent(in) :: i_born, i_real
    integer, intent(in) :: i_virt, i_sub
    associate (comp => def%initial(i))
       comp%associated_born = i_born
       comp%associated_real = i_real
       comp%associated_virt = i_virt
       comp%associated_sub = i_sub
    end associate
  end subroutine process_def_set_associated_components

\eatline
\nwendcode{}\nwbegindocs{488}\nwdocspar
Compute the MD5 sum for this process definition.  We compute the MD5
sums for all components individually, than concatenate a string of
those and compute the MD5 sum of this string.  We also include the
model name.  All other data part of the component definitions.
\nwenddocs{}\nwbegincode{489}\moddef{Process libraries: process def: TBP}\plusendmoddef
  procedure :: compute_md5sum => process_def_compute_md5sum
\nwendcode{}\nwbegincode{490}\moddef{Process libraries: procedures}\plusendmoddef
  subroutine process_def_compute_md5sum (def, model)
    class(process_def_t), intent(inout) :: def
    class(model_data_t), intent(in), optional, target :: model
    integer :: i
    type(string_t) :: buffer
    buffer = def%model_name
    do i = 1, def%n_initial
       call def%initial(i)%compute_md5sum (model)
       buffer = buffer // def%initial(i)%md5sum
    end do
    do i = 1, def%n_extra
       call def%extra(i)%compute_md5sum (model)
       buffer = buffer // def%initial(i)%md5sum
    end do
    def%md5sum = md5sum (char (buffer))
  end subroutine process_def_compute_md5sum
    
\eatline
\nwendcode{}\nwbegindocs{491}Return the MD5 sum of the process or of a process component.
\nwenddocs{}\nwbegincode{492}\moddef{Process libraries: process def: TBP}\plusendmoddef
  procedure :: get_md5sum => process_def_get_md5sum
\nwendcode{}\nwbegincode{493}\moddef{Process libraries: procedures}\plusendmoddef
  function process_def_get_md5sum (def, i_component) result (md5sum)
    class(process_def_t), intent(in) :: def
    integer, intent(in), optional :: i_component
    character(32) :: md5sum
    if (present (i_component)) then
       md5sum = def%initial(i_component)%md5sum
    else
       md5sum = def%md5sum
    end if
  end function process_def_get_md5sum
  
\eatline
\nwendcode{}\nwbegindocs{494}\nwdocspar
This query tells whether a specific process component relies on
external code.  This includes all traditional WHIZARD matrix elements
which rely on \oMega\ for code generation.  Other process components
(trivial decays, subtraction terms) do not require external code.

NOTE: Implemented only for initial component.

The query is passed to the process component.
\nwenddocs{}\nwbegincode{495}\moddef{Process libraries: process def: TBP}\plusendmoddef
  procedure :: needs_code => process_def_needs_code
\nwendcode{}\nwbegincode{496}\moddef{Process libraries: procedures}\plusendmoddef
  function process_def_needs_code (def, i_component) result (flag)
    class(process_def_t), intent(in) :: def
    integer, intent(in) :: i_component
    logical :: flag
    flag = def%initial(i_component)%needs_code ()
  end function process_def_needs_code
  
\eatline
\nwendcode{}\nwbegindocs{497}\nwdocspar
\subsubsection{Process definition list}
A list of process definitions is the starting point for creating a
process library.  The list is built when reading the user input.  When
reading an existing process library, the list is used for
cross-checking and updating the configuration.

We need a type for the list entry.  The simplest way is to extend the
process definition type, so all methods apply to the process
definition directly.
\nwenddocs{}\nwbegincode{498}\moddef{Process libraries: public}\plusendmoddef
  public :: process_def_entry_t
\nwendcode{}\nwbegincode{499}\moddef{Process libraries: types}\plusendmoddef
  type, extends (process_def_t) :: process_def_entry_t
     private
     type(process_def_entry_t), pointer :: next => null ()
  end type process_def_entry_t
  
\eatline
\nwendcode{}\nwbegindocs{500}This is the type for the list itself.
\nwenddocs{}\nwbegincode{501}\moddef{Process libraries: types}\plusendmoddef
  type :: process_def_list_t
     private
     type(process_def_entry_t), pointer :: first => null ()
     type(process_def_entry_t), pointer :: last => null ()
   contains
   \LA{}Process libraries: process def list: TBP\RA{}
  end type process_def_list_t

\eatline
\nwendcode{}\nwbegindocs{502}The deallocates the list iteratively.  We assume that the list
entries do not need finalization themselves.
\nwenddocs{}\nwbegincode{503}\moddef{Process libraries: process def list: TBP}\endmoddef
  procedure :: final => process_def_list_final
\nwendcode{}\nwbegincode{504}\moddef{Process libraries: procedures}\plusendmoddef
  subroutine process_def_list_final (list)
    class(process_def_list_t), intent(inout) :: list
    type(process_def_entry_t), pointer :: current
    nullify (list%last)
    do while (associated (list%first))
       current => list%first
       list%first => current%next
       deallocate (current)
    end do
  end subroutine process_def_list_final
  
\eatline
\nwendcode{}\nwbegindocs{505}Write the complete list.
\nwenddocs{}\nwbegincode{506}\moddef{Process libraries: process def list: TBP}\plusendmoddef
  procedure :: write => process_def_list_write
\nwendcode{}\nwbegincode{507}\moddef{Process libraries: procedures}\plusendmoddef
  subroutine process_def_list_write (object, unit, libpath)
    class(process_def_list_t), intent(in) :: object
    integer, intent(in), optional :: unit
    logical, intent(in), optional :: libpath
    type(process_def_entry_t), pointer :: entry
    integer :: i, u
    u = given_output_unit (unit)
    if (associated (object%first)) then
       i = 1
       entry => object%first
       do while (associated (entry))
          write (u, "(1x,A,I0,A)")  "Process #", i, ":"
          call entry%write (u)
          i = i + 1
          entry => entry%next
          if (associated (entry))  write (u, *)
       end do
    else
       write (u, "(1x,A)")  "Process definition list: [empty]"
    end if
  end subroutine process_def_list_write

\eatline
\nwendcode{}\nwbegindocs{508}Short account.
\nwenddocs{}\nwbegincode{509}\moddef{Process libraries: process def list: TBP}\plusendmoddef
  procedure :: show => process_def_list_show
\nwendcode{}\nwbegincode{510}\moddef{Process libraries: procedures}\plusendmoddef
  subroutine process_def_list_show (object, unit)
    class(process_def_list_t), intent(in) :: object
    integer, intent(in), optional :: unit
    type(process_def_entry_t), pointer :: entry
    integer :: u
    u = given_output_unit (unit)
    if (associated (object%first)) then
       write (u, "(2x,A)")  "Processes:"
       entry => object%first
       do while (associated (entry))
          call entry%show (u)
          entry => entry%next
       end do
    else
       write (u, "(2x,A)")  "Processes: [empty]"
    end if
  end subroutine process_def_list_show

\eatline
\nwendcode{}\nwbegindocs{511}Read the complete list.  We need an array of templates for the
component subobjects of abstract {\tt{}prc{\char95}core{\char95}t} type, to
allocate them with the correct specific type.

NOTE: Error handling is missing.  Reading will just be aborted on
error, or an I/O error occurs.
\nwenddocs{}\nwbegincode{512}\moddef{Process libraries: process def list: TBP}\plusendmoddef
  procedure :: read => process_def_list_read
\nwendcode{}\nwbegincode{513}\moddef{Process libraries: procedures}\plusendmoddef
  subroutine process_def_list_read (object, unit, core_def_templates)
    class(process_def_list_t), intent(out) :: object
    integer, intent(in) :: unit
    type(prc_template_t), dimension(:), intent(in) :: core_def_templates
    type(process_def_entry_t), pointer :: entry
    character(80) :: buffer, ref
    integer :: i
    read (unit, "(A)")  buffer
    write (ref, "(1x,A)")  "Process definition list: [empty]"  
    if (buffer == ref)  return         ! OK: empty library
    backspace (unit)
    READ_ENTRIES: do i = 1, huge (0)
       if (i > 1) read (unit, *, end=1)
       read (unit, "(A)")  buffer

       write (ref, "(1x,A,I0,A)")  "Process #", i, ":"
       if (buffer /= ref)  return      ! Wrong process header: done.
       allocate (entry)
       call entry%read (unit, core_def_templates)
       call object%append (entry)
    end do READ_ENTRIES
1   continue                           ! EOF: done
  end subroutine process_def_list_read

\eatline
\nwendcode{}\nwbegindocs{514}Append an entry to the list.  The entry should be allocated as a
pointer, and the pointer allocation is transferred.  The original
pointer is returned null.
\nwenddocs{}\nwbegincode{515}\moddef{Process libraries: process def list: TBP}\plusendmoddef
  procedure :: append => process_def_list_append
\nwendcode{}\nwbegincode{516}\moddef{Process libraries: procedures}\plusendmoddef
  subroutine process_def_list_append (list, entry)
    class(process_def_list_t), intent(inout) :: list
    type(process_def_entry_t), intent(inout), pointer :: entry
    if (list%contains (entry%id)) then
       call msg_fatal ("Recording process: '" // char (entry%id) &
            // "' has already been defined")
    end if
    if (associated (list%first)) then
       list%last%next => entry
    else
       list%first => entry
    end if
    list%last => entry
    entry => null ()
  end subroutine process_def_list_append
  
\eatline
\nwendcode{}\nwbegindocs{517}\nwdocspar
\subsubsection{Probe the process definition list}
Return the number of processes supported by the library.
\nwenddocs{}\nwbegincode{518}\moddef{Process libraries: process def list: TBP}\plusendmoddef
  procedure :: get_n_processes => process_def_list_get_n_processes
\nwendcode{}\nwbegincode{519}\moddef{Process libraries: procedures}\plusendmoddef
  function process_def_list_get_n_processes (list) result (n)
    integer :: n
    class(process_def_list_t), intent(in) :: list
    type(process_def_entry_t), pointer :: current
    n = 0
    current => list%first
    do while (associated (current))
       n = n + 1
       current => current%next
    end do
  end function process_def_list_get_n_processes
  
\eatline
\nwendcode{}\nwbegindocs{520}Allocate an array with the process IDs supported by the library.
\nwenddocs{}\nwbegincode{521}\moddef{Process libraries: process def list: TBP}\plusendmoddef
  procedure :: get_process_id_list => process_def_list_get_process_id_list
\nwendcode{}\nwbegincode{522}\moddef{Process libraries: procedures}\plusendmoddef
  subroutine process_def_list_get_process_id_list (list, id)
    class(process_def_list_t), intent(in) :: list
    type(string_t), dimension(:), allocatable, intent(out) :: id
    type(process_def_entry_t), pointer :: current
    integer :: i
    allocate (id (list%get_n_processes ()))
    i = 0
    current => list%first
    do while (associated (current))
       i = i + 1
       id(i) = current%id
       current => current%next
    end do
  end subroutine process_def_list_get_process_id_list
  
\eatline
\nwendcode{}\nwbegindocs{523}Return true if a given process is in the library.
\nwenddocs{}\nwbegincode{524}\moddef{Process libraries: process def list: TBP}\plusendmoddef
  procedure :: contains => process_def_list_contains
\nwendcode{}\nwbegincode{525}\moddef{Process libraries: procedures}\plusendmoddef
  function process_def_list_contains (list, id) result (flag)
    logical :: flag
    class(process_def_list_t), intent(in) :: list
    type(string_t), intent(in) :: id
    type(process_def_entry_t), pointer :: current
    current => list%first
    do while (associated (current))
       if (id == current%id) then
          flag = .true.;  return
       end if
       current => current%next
    end do
    flag = .false.
  end function process_def_list_contains

\eatline
\nwendcode{}\nwbegindocs{526}Return the index of the entry that corresponds to a given process.
\nwenddocs{}\nwbegincode{527}\moddef{Process libraries: process def list: TBP}\plusendmoddef
  procedure :: get_entry_index => process_def_list_get_entry_index
\nwendcode{}\nwbegincode{528}\moddef{Process libraries: procedures}\plusendmoddef
  function process_def_list_get_entry_index (list, id) result (n)
    integer :: n
    class(process_def_list_t), intent(in) :: list
    type(string_t), intent(in) :: id
    type(process_def_entry_t), pointer :: current
    n = 0
    current => list%first
    do while (associated (current))
       n = n + 1
       if (id == current%id) then
          return
       end if
       current => current%next
    end do
    n = 0
  end function process_def_list_get_entry_index
    
\eatline
\nwendcode{}\nwbegindocs{529}Return the numerical ID for a process.
\nwenddocs{}\nwbegincode{530}\moddef{Process libraries: process def list: TBP}\plusendmoddef
  procedure :: get_num_id => process_def_list_get_num_id
\nwendcode{}\nwbegincode{531}\moddef{Process libraries: procedures}\plusendmoddef
  function process_def_list_get_num_id (list, id) result (num_id)
    integer :: num_id
    class(process_def_list_t), intent(in) :: list
    type(string_t), intent(in) :: id
    type(process_def_entry_t), pointer :: current
    current => list%first
    do while (associated (current))
       if (id == current%id) then
          num_id = current%num_id
          return
       end if
       current => current%next
    end do
    num_id = 0
  end function process_def_list_get_num_id
    
\eatline
\nwendcode{}\nwbegindocs{532}Return the model name for a given process in the library.
\nwenddocs{}\nwbegincode{533}\moddef{Process libraries: process def list: TBP}\plusendmoddef
  procedure :: get_model_name => process_def_list_get_model_name
\nwendcode{}\nwbegincode{534}\moddef{Process libraries: procedures}\plusendmoddef
  function process_def_list_get_model_name (list, id) result (model_name)
    type(string_t) :: model_name
    class(process_def_list_t), intent(in) :: list
    type(string_t), intent(in) :: id
    type(process_def_entry_t), pointer :: current
    current => list%first
    do while (associated (current))
       if (id == current%id) then
          model_name = current%model_name
          return
       end if
       current => current%next
    end do
    model_name = ""
  end function process_def_list_get_model_name
  
\eatline
\nwendcode{}\nwbegindocs{535}Return the number of incoming particles of a given process in the library.
This tells us whether the process is a decay or a scattering.
\nwenddocs{}\nwbegincode{536}\moddef{Process libraries: process def list: TBP}\plusendmoddef
  procedure :: get_n_in => process_def_list_get_n_in
\nwendcode{}\nwbegincode{537}\moddef{Process libraries: procedures}\plusendmoddef
  function process_def_list_get_n_in (list, id) result (n)
    integer :: n
    class(process_def_list_t), intent(in) :: list
    type(string_t), intent(in) :: id
    type(process_def_entry_t), pointer :: current
    current => list%first
    do while (associated (current))
       if (id == current%id) then
          n = current%n_in
          return
       end if
       current => current%next
    end do
  end function process_def_list_get_n_in
  
\eatline
\nwendcode{}\nwbegindocs{538}Return the number of components of a given process in the library.
\nwenddocs{}\nwbegincode{539}\moddef{Process libraries: process def list: TBP}\plusendmoddef
  procedure :: get_n_components => process_def_list_get_n_components
\nwendcode{}\nwbegincode{540}\moddef{Process libraries: procedures}\plusendmoddef
  function process_def_list_get_n_components (list, id) result (n)
    integer :: n
    class(process_def_list_t), intent(in) :: list
    type(string_t), intent(in) :: id
    type(process_def_entry_t), pointer :: current
    current => list%first
    do while (associated (current))
       if (id == current%id) then
          n = current%n_initial + current%n_extra
          return
       end if
       current => current%next
    end do
  end function process_def_list_get_n_components
  
\eatline
\nwendcode{}\nwbegindocs{541}Return a pointer to a specific process component definition.
\nwenddocs{}\nwbegincode{542}\moddef{Process libraries: process def list: TBP}\plusendmoddef
  procedure :: get_component_def_ptr => process_def_list_get_component_def_ptr
\nwendcode{}\nwbegincode{543}\moddef{Process libraries: procedures}\plusendmoddef
  function process_def_list_get_component_def_ptr (list, id, i) result (ptr)
    class(process_def_list_t), intent(in) :: list
    type(string_t), intent(in) :: id
    integer, intent(in) :: i
    type(process_component_def_t), pointer :: ptr
    type(process_def_entry_t), pointer :: current
    ptr => null ()
    current => list%first
    do while (associated (current))
       if (id == current%id) then
          if (i <= current%n_initial) then
             ptr => current%initial(i)
          else if (i <= current%n_initial + current%n_extra) then
             ptr => current%extra(i-current%n_initial)
          end if
          return
       end if
       current => current%next
    end do
  end function process_def_list_get_component_def_ptr
  
\eatline
\nwendcode{}\nwbegindocs{544}Return the list of component IDs of a given process in the library.
\nwenddocs{}\nwbegincode{545}\moddef{Process libraries: process def list: TBP}\plusendmoddef
  procedure :: get_component_list => process_def_list_get_component_list
\nwendcode{}\nwbegincode{546}\moddef{Process libraries: procedures}\plusendmoddef
  subroutine process_def_list_get_component_list (list, id, cid)
    class(process_def_list_t), intent(in) :: list
    type(string_t), intent(in) :: id
    type(string_t), dimension(:), allocatable, intent(out) :: cid
    type(process_def_entry_t), pointer :: current
    integer :: i, n
    current => list%first
    do while (associated (current))
       if (id == current%id) then
          allocate (cid (current%n_initial + current%n_extra))
          do i = 1, current%n_initial
             cid(i) = current%initial(i)%basename
          end do
          n = current%n_initial
          do i = 1, current%n_extra
             cid(n + i) = current%extra(i)%basename
          end do
          return
       end if
       current => current%next
    end do
  end subroutine process_def_list_get_component_list
  
\eatline
\nwendcode{}\nwbegindocs{547}Return the list of component description strings for a given process
in the library.
\nwenddocs{}\nwbegincode{548}\moddef{Process libraries: process def list: TBP}\plusendmoddef
  procedure :: get_component_description_list => &
       process_def_list_get_component_description_list
\nwendcode{}\nwbegincode{549}\moddef{Process libraries: procedures}\plusendmoddef
  subroutine process_def_list_get_component_description_list &
       (list, id, description)
    class(process_def_list_t), intent(in) :: list
    type(string_t), intent(in) :: id
    type(string_t), dimension(:), allocatable, intent(out) :: description
    type(process_def_entry_t), pointer :: current
    integer :: i, n
    current => list%first
    do while (associated (current))
       if (id == current%id) then
          allocate (description (current%n_initial + current%n_extra))
          do i = 1, current%n_initial
             description(i) = current%initial(i)%description
          end do
          n = current%n_initial
          do i = 1, current%n_extra
             description(n + i) = current%extra(i)%description
          end do
          return
       end if
       current => current%next
    end do
  end subroutine process_def_list_get_component_description_list
  
\eatline
\nwendcode{}\nwbegindocs{550}Return whether the entry corresponds to an NLO-process
\nwenddocs{}\nwbegincode{551}\moddef{Process libraries: process def list: TBP}\plusendmoddef
  procedure :: get_nlo_process => process_def_list_get_nlo_process
\nwendcode{}\nwbegincode{552}\moddef{Process libraries: procedures}\plusendmoddef
  function process_def_list_get_nlo_process (list, id) result (nlo)
    class(process_def_list_t), intent(in) :: list
    type(string_t), intent(in) :: id
    logical :: nlo
    type(process_def_entry_t), pointer :: current
    current => list%first
    do while (associated (current))
      if (id == current%id) then
        nlo = current%nlo_process
        return
      end if
      current => current%next
    end do
  end function process_def_list_get_nlo_process

\eatline
\nwendcode{}\nwbegindocs{553}\nwdocspar
\subsection{Process library}
The process library object is the interface between the process
definition data, as provided by the user, generated or linked process
code on file, and the process run data that reference the process
code.

\subsubsection{Process library entry}
For each process component that is part of the library, there is a
separate library entry ({\tt{}process{\char95}library{\char95}entry{\char95}t}.  The library
entry connects a process definition with the specific code (if any) in
the compiled driver library.

The {\tt{}status} indicates how far the process has been
processed by the system (definition, code generation, compilation,
linking).  A process with status {\tt{}STAT{\char95}LOADED} is accessible for
computing matrix elements.

The {\tt{}def} pointer identifies the corresponding process definition.
The process component within that definition is identified by the
{\tt{}i{\char95}component} index.

The {\tt{}i{\char95}external} index refers to the compiled library driver.  If it is zero,
there is no associated matrix-element code.

The {\tt{}driver} component holds the pointers to the matrix-element
specific functions, in particular the matrix element function itself.
\nwenddocs{}\nwbegincode{554}\moddef{Process libraries: types}\plusendmoddef
  type :: process_library_entry_t
     private
     integer :: status = STAT_UNKNOWN
     type(process_def_t), pointer :: def => null ()
     integer :: i_component = 0
     integer :: i_external = 0
     class(prc_core_driver_t), allocatable :: driver
   contains
   \LA{}Process libraries: process library entry: TBP\RA{}
  end type process_library_entry_t

\eatline
\nwendcode{}\nwbegindocs{555}Here are the available status codes.  An entry starts with
{\tt{}UNKNOWN} status.  Once the association with a valid process
definition is established, the status becomes {\tt{}CONFIGURED}.
If matrix element source code is to be generated by the system or
provided from elsewhere, {\tt{}CODE{\char95}GENERATED} indicates that this is
done.  The {\tt{}COMPILED} status is next, it also applies to
processes which are accessed as precompiled binaries.  Finally, the
library is linked and process pointers are set; this is marked as
{\tt{}LOADED}.

For a process library, the initial status is {\tt{}OPEN}, since process
definitions may be added.  After configuration, the process content is fixed
and the status becomes {\tt{}CONFIGURED}.  The further states are as above,
always referring to the lowest status among the process entries.
\nwenddocs{}\nwbegincode{556}\moddef{Process libraries: parameters}\endmoddef
  integer, parameter :: STAT_UNKNOWN = 0
  integer, parameter :: STAT_OPEN = 1
  integer, parameter :: STAT_CONFIGURED = 2
  integer, parameter :: STAT_SOURCE = 3
  integer, parameter :: STAT_COMPILED = 4
  integer, parameter :: STAT_LINKED = 5
  integer, parameter :: STAT_ACTIVE = 6

\eatline
\eatline
\nwendcode{}\nwbegindocs{557}These are the associated code letters, for output:
\nwenddocs{}\nwbegincode{558}\moddef{Process libraries: parameters}\plusendmoddef
  character, dimension(0:6), parameter :: STATUS_LETTER = &
       ["?", "o", "f", "s", "c", "l", "a"]

\eatline
\nwendcode{}\nwbegindocs{559}This produces a condensed account of the library entry.  The status
is indicated by a letter in brackets, then the ID and component index
of the associated process definition, finally the library index, if available.
\nwenddocs{}\nwbegincode{560}\moddef{Process libraries: process library entry: TBP}\endmoddef
  procedure :: to_string => process_library_entry_to_string
\nwendcode{}\nwbegincode{561}\moddef{Process libraries: procedures}\plusendmoddef
  function process_library_entry_to_string (object) result (string)
    type(string_t) :: string
    class(process_library_entry_t), intent(in) :: object
    character(32) :: buffer
    string = "[" // STATUS_LETTER(object%status) // "]"
    select case (object%status)
    case (STAT_UNKNOWN)
    case default
       if (associated (object%def)) then
          write (buffer, "(I0)")  object%i_component
          string = string // " " // object%def%id // "." // trim (buffer)
       end if
       if (object%i_external /= 0) then
          write (buffer, "(I0)")  object%i_external
          string = string // " = ext:" // trim (buffer)
       else
          string = string // " = int"
       end if
       if (allocated (object%driver)) then
          string = string // " (" // object%driver%type_name () // ")"
       end if
    end select
  end function process_library_entry_to_string
  
\eatline
\nwendcode{}\nwbegindocs{562}Initialize with data.  Used for the unit tests.
\nwenddocs{}\nwbegincode{563}\moddef{Process libraries: process library entry: TBP}\plusendmoddef
  procedure :: init => process_library_entry_init
\nwendcode{}\nwbegincode{564}\moddef{Process libraries: procedures}\plusendmoddef
  subroutine process_library_entry_init (object, &
       status, def, i_component, i_external)
    class(process_library_entry_t), intent(out) :: object
    integer, intent(in) :: status
    type(process_def_t), target, intent(in) :: def
    integer, intent(in) :: i_component
    integer, intent(in) :: i_external
    object%status = status
    object%def => def
    object%i_component = i_component
    object%i_external = i_external
  end subroutine process_library_entry_init
  
\eatline
\nwendcode{}\nwbegindocs{565}Assign pointers for all process-specific features.  We have to
combine the method from the {\tt{}core{\char95}def} specification, the
assigned pointers within the library driver, the index within that
driver, and the process driver which should receive the links.
\nwenddocs{}\nwbegincode{566}\moddef{Process libraries: process library entry: TBP}\plusendmoddef
  procedure :: connect => process_library_entry_connect
\nwendcode{}\nwbegincode{567}\moddef{Process libraries: procedures}\plusendmoddef
  subroutine process_library_entry_connect (entry, lib_driver, i)
    class(process_library_entry_t), intent(inout) :: entry
    class(prclib_driver_t), intent(in) :: lib_driver
    integer, intent(in) :: i
    call entry%def%initial(entry%i_component)%connect &
         (lib_driver, i, entry%driver)
  end subroutine process_library_entry_connect

\eatline
\nwendcode{}\nwbegindocs{568}\nwdocspar
\subsubsection{The process library object}
The {\tt{}process{\char95}library{\char95}t} type is an extension of the
{\tt{}process{\char95}def{\char95}list{\char95}t} type.  Thus, it automatically contains the
process definition list.  

The {\tt{}basename} identifies the library generically.

The {\tt{}external} flag is true if any process within the library needs external
code, so the library must correspond to an actual code library (statically or
dynamically linked).

The {\tt{}entry} array contains all process components that can be handled by this
library. Each entry refers to the process (component) definition and to the
associated external matrix element code, if there is any.

The {\tt{}driver} object is needed only if {\tt{}external} is true.  This object
handles all interactions with external matrix-element code.

The {\tt{}md5sum} summarizes the complete {\tt{}process{\char95}def{\char95}list{\char95}t} base
object.  It can be used to check if the library configuration has changed.
\nwenddocs{}\nwbegincode{569}\moddef{Process libraries: public}\plusendmoddef
  public :: process_library_t
\nwendcode{}\nwbegincode{570}\moddef{Process libraries: types}\plusendmoddef
  type, extends (process_def_list_t) :: process_library_t
     private
     type(string_t) :: basename
     integer :: n_entries = 0
     logical :: external = .false.
     integer :: status = STAT_UNKNOWN
     logical :: static = .false.
     logical :: driver_exists = .false.
     logical :: makefile_exists = .false.
     type(process_library_entry_t), dimension(:), allocatable :: entry
     class(prclib_driver_t), allocatable :: driver
     character(32) :: md5sum = ""
   contains
   \LA{}Process libraries: process library: TBP\RA{}
  end type process_library_t

\eatline
\nwendcode{}\nwbegindocs{571}For the output, we write first the metadata and the DL access
record, then the library entries in short form, and finally the
process definition list which is the base object.

Don't write the MD5 sum since this is used to generate it.
\nwenddocs{}\nwbegincode{572}\moddef{Process libraries: process library: TBP}\endmoddef
  procedure :: write => process_library_write
\nwendcode{}\nwbegincode{573}\moddef{Process libraries: procedures}\plusendmoddef
  subroutine process_library_write (object, unit, libpath)
    class(process_library_t), intent(in) :: object
    integer, intent(in), optional :: unit
    logical, intent(in), optional :: libpath
    integer :: i, u
    u = given_output_unit (unit)
    write (u, "(1x,A,A)")  "Process library: ", char (object%basename)
    write (u, "(3x,A,L1)")   "external        = ", object%external
    write (u, "(3x,A,L1)")   "makefile exists = ", object%makefile_exists
    write (u, "(3x,A,L1)")   "driver exists   = ", object%driver_exists
    write (u, "(3x,A,A1)")   "code status     = ", &
         STATUS_LETTER (object%status)
    write (u, *)
    if (allocated (object%entry)) then
       write (u, "(1x,A)", advance="no")  "Process library entries:"
       write (u, "(1x,I0)")  object%n_entries
       do i = 1, size (object%entry)
          write (u, "(1x,A,I0,A,A)")  "Entry #", i, ": ", &
               char (object%entry(i)%to_string ())
       end do
       write (u, *)
    end if
    if (object%external) then
       call object%driver%write (u, libpath)
       write (u, *)
    end if
    call object%process_def_list_t%write (u)
  end subroutine process_library_write
         
\eatline
\nwendcode{}\nwbegindocs{574}Condensed version for screen output.
\nwenddocs{}\nwbegincode{575}\moddef{Process libraries: process library: TBP}\plusendmoddef
  procedure :: show => process_library_show
\nwendcode{}\nwbegincode{576}\moddef{Process libraries: procedures}\plusendmoddef
  subroutine process_library_show (object, unit)
    class(process_library_t), intent(in) :: object
    integer, intent(in), optional :: unit
    integer :: u
    u = given_output_unit (unit)
    write (u, "(A,A)")  "Process library: ", char (object%basename)
    write (u, "(2x,A,L1)")   "external        = ", object%external
    if (object%static) then
       write (u, "(2x,A,L1)")   "static          = ", .true.
    else
       write (u, "(2x,A,L1)")   "makefile exists = ", object%makefile_exists
       write (u, "(2x,A,L1)")   "driver exists   = ", object%driver_exists
    end if
    write (u, "(2x,A,A1)", advance="no")   "code status     = "
    select case (object%status)
    case (STAT_UNKNOWN);    write (u, "(A)")  "[unknown]"
    case (STAT_OPEN);       write (u, "(A)")  "open"
    case (STAT_CONFIGURED); write (u, "(A)")  "configured"
    case (STAT_SOURCE);     write (u, "(A)")  "source code exists"
    case (STAT_COMPILED);   write (u, "(A)")  "compiled"
    case (STAT_LINKED);     write (u, "(A)")  "linked"
    case (STAT_ACTIVE);     write (u, "(A)")  "active"
    end select
    call object%process_def_list_t%show (u)
  end subroutine process_library_show
         
\eatline
\nwendcode{}\nwbegindocs{577}\nwdocspar
The initializer defines just the basename.  We may now add process definitions
to the library.
\nwenddocs{}\nwbegincode{578}\moddef{Process libraries: process library: TBP}\plusendmoddef
  procedure :: init => process_library_init
\nwendcode{}\nwbegincode{579}\moddef{Process libraries: procedures}\plusendmoddef
  subroutine process_library_init (lib, basename)
    class(process_library_t), intent(out) :: lib
    type(string_t), intent(in) :: basename
    lib%basename = basename
    lib%status = STAT_OPEN
    call msg_message ("Process library '" // char (basename) &
         // "': initialized")
  end subroutine process_library_init
  
\eatline
\nwendcode{}\nwbegindocs{580}\nwdocspar
This alternative initializer declares the library as static.  We
should now add process definitions to the library, but all external
process code exists already.  We need the driver object, and we should
check the defined processes against the stored ones.
\nwenddocs{}\nwbegincode{581}\moddef{Process libraries: process library: TBP}\plusendmoddef
  procedure :: init_static => process_library_init_static
\nwendcode{}\nwbegincode{582}\moddef{Process libraries: procedures}\plusendmoddef
  subroutine process_library_init_static (lib, basename)
    class(process_library_t), intent(out) :: lib
    type(string_t), intent(in) :: basename
    lib%basename = basename
    lib%status = STAT_OPEN
    lib%static = .true.
    call msg_message ("Static process library '" // char (basename) &
         // "': initialized")
  end subroutine process_library_init_static
  
\eatline
\nwendcode{}\nwbegindocs{583}The {\tt{}configure} procedure scans the allocated entries in the process
definition list.  The configuration proceeds in three passes.

In the first pass, we scan the process definition list and count the
number of process components and the number of components which need
external code.  This is used to allocate the {\tt{}entry} array.

In the second pass, we initialize the {\tt{}entry} elements which connect
process definitions, process driver objects, and external code.

In the third pass, we initialize the library driver object, allocating
an entry for each external matrix element.

NOTE: Currently we handle only {\tt{}initial} process components; {\tt{}extra}
components are ignored.
\nwenddocs{}\nwbegincode{584}\moddef{Process libraries: process library: TBP}\plusendmoddef
  procedure :: configure => process_library_configure
\nwendcode{}\nwbegincode{585}\moddef{Process libraries: procedures}\plusendmoddef
  subroutine process_library_configure (lib, os_data)
    class(process_library_t), intent(inout) :: lib
    type(os_data_t), intent(in) :: os_data
    type(process_def_entry_t), pointer :: def_entry
    integer :: n_entries, n_external, i_entry, i_external
    type(string_t) :: model_name
    integer :: i_component

    n_entries = 0
    n_external = 0
    if (allocated (lib%entry))  deallocate (lib%entry)
    
    def_entry => lib%first
    do while (associated (def_entry))
       do i_component = 1, def_entry%n_initial
          n_entries = n_entries + 1
          if (def_entry%initial(i_component)%needs_code ()) then
             n_external = n_external + 1
             lib%external = .true.
          end if
       end do
       def_entry => def_entry%next
    end do
    lib%n_entries = n_entries

    allocate (lib%entry (n_entries))
    i_entry = 0
    i_external = 0
    def_entry => lib%first
    do while (associated (def_entry))
       do i_component = 1, def_entry%n_initial
          i_entry = i_entry + 1
          associate (lib_entry => lib%entry(i_entry))
            lib_entry%status = STAT_CONFIGURED
            lib_entry%def => def_entry%process_def_t
            lib_entry%i_component = i_component
            if (def_entry%initial(i_component)%needs_code ()) then
               i_external = i_external + 1
               lib_entry%i_external = i_external
            end if
            call def_entry%initial(i_component)%allocate_driver &
                 (lib_entry%driver)
          end associate
       end do
       def_entry => def_entry%next
    end do

    call dispatch_prclib_driver (lib%driver, &
         lib%basename, lib%get_modellibs_ldflags (os_data))
    call lib%driver%init (n_external)
    do i_entry = 1, n_entries
       associate (lib_entry => lib%entry(i_entry))
         i_component = lib_entry%i_component
         model_name = lib_entry%def%model_name
         associate (def => lib_entry%def%initial(i_component))
           if (def%needs_code ()) then
              call lib%driver%set_record (lib_entry%i_external, &
                   def%basename, &
                   model_name, &
                   def%get_features (), def%get_writer_ptr ())
           end if
         end associate
       end associate
    end do
    
    if (lib%static) then
       if (lib%n_entries /= 0)  lib%entry%status = STAT_LINKED
       lib%status = STAT_LINKED
    else if (lib%external) then
       where (lib%entry%i_external == 0)  lib%entry%status = STAT_LINKED
       lib%status = STAT_CONFIGURED
       lib%makefile_exists = .false.
       lib%driver_exists = .false.
    else
       if (lib%n_entries /= 0)  lib%entry%status = STAT_LINKED
       lib%status = STAT_LINKED
    end if
  end subroutine process_library_configure
  
\eatline
\nwendcode{}\nwbegindocs{586}Compute the MD5 sum.  We concatenate the individual MD5 sums of all
processes (which, in turn, are derived from the MD5 sums of their
components) and compute the MD5 sum of that.

This should be executed \emph{after} configuration, where the driver was
initialized, since otherwise the MD5 sum stored in the driver would be
overwritten.
\nwenddocs{}\nwbegincode{587}\moddef{Process libraries: process library: TBP}\plusendmoddef
  procedure :: compute_md5sum => process_library_compute_md5sum
\nwendcode{}\nwbegincode{588}\moddef{Process libraries: procedures}\plusendmoddef
  subroutine process_library_compute_md5sum (lib, model)
    class(process_library_t), intent(inout) :: lib
    class(model_data_t), intent(in), optional, target :: model
    type(process_def_entry_t), pointer :: def_entry
    type(string_t) :: buffer
    buffer = lib%basename
    def_entry => lib%first
    do while (associated (def_entry))
       call def_entry%compute_md5sum (model)
       buffer = buffer // def_entry%md5sum
       def_entry => def_entry%next
    end do
    lib%md5sum = md5sum (char (buffer))
    call lib%driver%set_md5sum (lib%md5sum)
  end subroutine process_library_compute_md5sum
  
\eatline
\nwendcode{}\nwbegindocs{589}Write an appropriate makefile, if there are external processes.  Unless
{\tt{}force} is in effect, first check if there is already a makefile with the
correct MD5 sum.  If yes, do nothing.
\nwenddocs{}\nwbegincode{590}\moddef{Process libraries: process library: TBP}\plusendmoddef
  procedure :: write_makefile => process_library_write_makefile
\nwendcode{}\nwbegincode{591}\moddef{Process libraries: procedures}\plusendmoddef
  subroutine process_library_write_makefile (lib, os_data, force, testflag)
    class(process_library_t), intent(inout) :: lib
    type(os_data_t), intent(in) :: os_data
    logical, intent(in) :: force
    logical, intent(in), optional :: testflag
    character(32) :: md5sum_file
    logical :: generate
    integer :: unit
    if (lib%external .and. .not. lib%static) then
       generate = .true.
       if (.not. force) then
          md5sum_file = lib%driver%get_md5sum_makefile ()
          if (lib%md5sum == md5sum_file) then
             call msg_message ("Process library '" // char (lib%basename) &
                  // "': keeping makefile")
             generate = .false.
          end if
       end if
       if (generate) then
          call msg_message ("Process library '" // char (lib%basename) &
               // "': writing makefile")
          unit = free_unit ()
          open (unit, file = char (lib%driver%basename // ".makefile"), &
               status="replace", action="write")
          call lib%driver%generate_makefile (unit, os_data, testflag)
          close (unit)
       end if
       lib%makefile_exists = .true.
    end if
  end subroutine process_library_write_makefile
  
\eatline
\nwendcode{}\nwbegindocs{592}Write the driver source code for the library to file, if there are
external processes.
\nwenddocs{}\nwbegincode{593}\moddef{Process libraries: process library: TBP}\plusendmoddef
  procedure :: write_driver => process_library_write_driver
\nwendcode{}\nwbegincode{594}\moddef{Process libraries: procedures}\plusendmoddef
  subroutine process_library_write_driver (lib, force)
    class(process_library_t), intent(inout) :: lib
    logical, intent(in) :: force
    character(32) :: md5sum_file
    logical :: generate
    integer :: unit
    if (lib%external .and. .not. lib%static) then
       generate = .true.
       if (.not. force) then
          md5sum_file = lib%driver%get_md5sum_driver ()
          if (lib%md5sum == md5sum_file) then
             call msg_message ("Process library '" // char (lib%basename) &
                  // "': keeping driver")
             generate = .false.
          end if
       end if
       if (generate) then
          call msg_message ("Process library '" // char (lib%basename) &
               // "': writing driver")
          unit = free_unit ()
          open (unit, file = char (lib%driver%basename // ".f90"), &
               status="replace", action="write")
          call lib%driver%generate_driver_code (unit)
          close (unit)
       end if
       lib%driver_exists = .true.
    end if
  end subroutine process_library_write_driver

\eatline
\nwendcode{}\nwbegindocs{595}Update the compilation status of an external library.

Strictly speaking, this is not necessary for a one-time run, since the
individual library methods will update the status themselves.
However, it allows us to identify compilation steps that we can skip
because the file exists or is already loaded, for the whole library or
for particular entries.

Independently, the building process is controlled by a makefile.
Thus, previous files are reused if they are not modified by the
current compilation.
\begin{enumerate}
\item
  If it is not already loaded, attempt to load the library.  If successful,
  check the overall MD5 sum.  If it matches, just keep it loaded and mark as
  ACTIVE.  If not, check the MD5 sum for all linked process components.
  Where it matches, mark the entry as COMPILED.  Then, unload the library and
  mark as CONFIGURED.
  
  Thus, we can identify compiled files for all matrix elements which are
  accessible via the previous compiled library, even if it is no longer up to
  date.
\item
  If the library is now in CONFIGURED state, look for valid source files.
  Each entry that is just in CONFIGURED state will advance to SOURCE if the
  MD5 sum matches.  Finally, advance the whole library to SOURCE if all
  entries are at least in this condition.
\end{enumerate}
\nwenddocs{}\nwbegincode{596}\moddef{Process libraries: process library: TBP}\plusendmoddef
  procedure :: update_status => process_library_update_status
\nwendcode{}\nwbegincode{597}\moddef{Process libraries: procedures}\plusendmoddef
  subroutine process_library_update_status (lib, os_data)
    class(process_library_t), intent(inout) :: lib
    type(os_data_t), intent(in) :: os_data
    character(32) :: md5sum_file
    integer :: i, i_external, i_component
    if (lib%external) then
       select case (lib%status)
       case (STAT_CONFIGURED:STAT_LINKED)
          call lib%driver%load (os_data, noerror=.true.)
       end select
       if (lib%driver%loaded) then
          md5sum_file = lib%driver%get_md5sum (0)
          if (lib%md5sum == md5sum_file) then
             call lib%load_entries ()
             lib%entry%status = STAT_ACTIVE
             lib%status = STAT_ACTIVE
             call msg_message ("Process library '" // char (lib%basename) &
                  // "': active")
          else
             do i = 1, lib%n_entries
                associate (entry => lib%entry(i))
                  i_external = entry%i_external
                  i_component = entry%i_component
                  if (i_external /= 0) then
                     md5sum_file = lib%driver%get_md5sum (i_external)
                     if (entry%def%get_md5sum (i_component) == md5sum_file) then
                        entry%status = STAT_COMPILED
                     else
                        entry%status = STAT_CONFIGURED
                     end if
                  end if
                end associate
             end do
             call lib%driver%unload ()
             lib%status = STAT_CONFIGURED
          end if
       end if
       select case (lib%status)
       case (STAT_CONFIGURED)
          do i = 1, lib%n_entries
             associate (entry => lib%entry(i))
               i_external = entry%i_external
               i_component = entry%i_component
               if (i_external /= 0) then
                  select case (entry%status)
                  case (STAT_CONFIGURED)
                     md5sum_file = lib%driver%get_md5sum_source (i_external)
                     if (entry%def%get_md5sum (i_component) == md5sum_file) then
                        entry%status = STAT_SOURCE
                     end if
                  end select
               end if
             end associate
          end do
          if (all (lib%entry%status >= STAT_SOURCE)) then
             md5sum_file = lib%driver%get_md5sum_driver ()
             if (lib%md5sum == md5sum_file) then
                lib%status = STAT_SOURCE
             end if
          end if
       end select
    end if
  end subroutine process_library_update_status

\eatline
\nwendcode{}\nwbegindocs{598}\nwdocspar
This procedure triggers code generation for all processes where this
is possible.

We generate code only for external processes of status
{\tt{}STAT{\char95}CONFIGURED}, which then advance to {\tt{}STAT{\char95}SOURCE}.  If, for a
particular process, the status is already advanced, we do not remove previous
files, so {\tt{}make} will consider them as up to date if they exist.  Otherwise,
we remove those files to force a fresh {\tt{}make}.

Finally, if any source code has been generated, we need a driver file.
\nwenddocs{}\nwbegincode{599}\moddef{Process libraries: process library: TBP}\plusendmoddef
  procedure :: make_source => process_library_make_source
\nwendcode{}\nwbegincode{600}\moddef{Process libraries: procedures}\plusendmoddef
  subroutine process_library_make_source (lib, os_data, keep_old_source)
    class(process_library_t), intent(inout) :: lib
    type(os_data_t), intent(in) :: os_data
    logical, intent(in), optional :: keep_old_source
    logical :: keep_old
    integer :: i, i_external
    keep_old = .false.
    if (present (keep_old_source))  keep_old = keep_old_source
    if (lib%external .and. .not. lib%static) then
       select case (lib%status)
       case (STAT_CONFIGURED)
          if (keep_old) then
             call msg_message ("Process library '" // char (lib%basename) &
                  // "': keeping source code")
          else
             call msg_message ("Process library '" // char (lib%basename) &
                  // "': creating source code")
             do i = 1, size (lib%entry)
                associate (entry => lib%entry(i))
                  i_external = entry%i_external
                  if (i_external /= 0 &
                       .and. lib%entry(i)%status == STAT_CONFIGURED) then
                     call lib%driver%clean_proc (i_external, os_data)
                  end if
                end associate
                if (signal_is_pending ())  return
             end do
             call lib%driver%make_source (os_data)
          end if
          lib%status = STAT_SOURCE
          where (lib%entry%i_external /= 0 &
               .and. lib%entry%status == STAT_CONFIGURED)
             lib%entry%status = STAT_SOURCE
          end where
          lib%status = STAT_SOURCE
       end select
    end if
  end subroutine process_library_make_source
  
\eatline
\nwendcode{}\nwbegindocs{601}Compile the generated code and update the status codes.  Try to make
the sources first, just in case. This includes compiling possible \LaTeX
Feynman diagram files. 
\nwenddocs{}\nwbegincode{602}\moddef{Process libraries: process library: TBP}\plusendmoddef
  procedure :: make_compile => process_library_make_compile
\nwendcode{}\nwbegincode{603}\moddef{Process libraries: procedures}\plusendmoddef
  subroutine process_library_make_compile (lib, os_data, keep_old_source)
    class(process_library_t), intent(inout) :: lib
    type(os_data_t), intent(in) :: os_data
    logical, intent(in), optional :: keep_old_source
    if (lib%external .and. .not. lib%static) then
       select case (lib%status)
       case (STAT_CONFIGURED)
          call lib%make_source (os_data, keep_old_source)
       end select
       if (signal_is_pending ())  return
       select case (lib%status)
       case (STAT_SOURCE)
          call msg_message ("Process library '" // char (lib%basename) &
               // "': compiling sources")
          call lib%driver%make_compile (os_data)
          where (lib%entry%i_external /= 0 &
               .and. lib%entry%status == STAT_SOURCE)
             lib%entry%status = STAT_COMPILED
          end where
          lib%status = STAT_COMPILED
       end select
    end if
  end subroutine process_library_make_compile
  
\eatline
\nwendcode{}\nwbegindocs{604}Link the process library.  Try to compile first, just in case.
\nwenddocs{}\nwbegincode{605}\moddef{Process libraries: process library: TBP}\plusendmoddef
  procedure :: make_link => process_library_make_link
\nwendcode{}\nwbegincode{606}\moddef{Process libraries: procedures}\plusendmoddef
  subroutine process_library_make_link (lib, os_data, keep_old_source)
    class(process_library_t), intent(inout) :: lib
    type(os_data_t), intent(in) :: os_data
    logical, intent(in), optional :: keep_old_source
    if (lib%external .and. .not. lib%static) then
       select case (lib%status)
       case (STAT_CONFIGURED:STAT_SOURCE)
          call lib%make_compile (os_data, keep_old_source)
       end select
       if (signal_is_pending ())  return
       select case (lib%status)
       case (STAT_COMPILED)
          call msg_message ("Process library '" // char (lib%basename) &
               // "': linking")
          call lib%driver%make_link (os_data)
          lib%entry%status = STAT_LINKED
          lib%status = STAT_LINKED
       end select
    end if
  end subroutine process_library_make_link
  
\eatline
\nwendcode{}\nwbegindocs{607}Load the process library, i.e., assign pointers to the library
functions.
\nwenddocs{}\nwbegincode{608}\moddef{Process libraries: process library: TBP}\plusendmoddef
  procedure :: load => process_library_load
\nwendcode{}\nwbegincode{609}\moddef{Process libraries: procedures}\plusendmoddef
  subroutine process_library_load (lib, os_data, keep_old_source)
    class(process_library_t), intent(inout) :: lib
    type(os_data_t), intent(in) :: os_data
    logical, intent(in), optional :: keep_old_source
    select case (lib%status)
    case (STAT_CONFIGURED:STAT_COMPILED)
       call lib%make_link (os_data, keep_old_source)
    end select
    if (signal_is_pending ())  return
    select case (lib%status)
    case (STAT_LINKED)
       if (lib%external) then
          call msg_message ("Process library '" // char (lib%basename) &
               // "': loading")
          call lib%driver%load (os_data)
          call lib%load_entries ()
       end if
       lib%entry%status = STAT_ACTIVE
       lib%status = STAT_ACTIVE
    end select
  end subroutine process_library_load
  
\eatline
\nwendcode{}\nwbegindocs{610}This is the actual loading part for the process methods.
\nwenddocs{}\nwbegincode{611}\moddef{Process libraries: process library: TBP}\plusendmoddef
  procedure :: load_entries => process_library_load_entries
\nwendcode{}\nwbegincode{612}\moddef{Process libraries: procedures}\plusendmoddef
  subroutine process_library_load_entries (lib)
    class(process_library_t), intent(inout) :: lib
    integer :: i
    do i = 1, size (lib%entry)
       associate (entry => lib%entry(i))
         if (entry%i_external /= 0) then
            call entry%connect (lib%driver, entry%i_external)
         end if
       end associate
    end do
  end subroutine process_library_load_entries

\eatline
\nwendcode{}\nwbegindocs{613}Unload the library, if possible.  This reverts the status to ``linked''.
\nwenddocs{}\nwbegincode{614}\moddef{Process libraries: process library: TBP}\plusendmoddef
  procedure :: unload => process_library_unload
\nwendcode{}\nwbegincode{615}\moddef{Process libraries: procedures}\plusendmoddef
  subroutine process_library_unload (lib)
    class(process_library_t), intent(inout) :: lib
    select case (lib%status)
    case (STAT_ACTIVE)
       if (lib%external) then
          call msg_message ("Process library '" // char (lib%basename) &
               // "': unloading")
          call lib%driver%unload ()
       end if
       lib%entry%status = STAT_LINKED
       lib%status = STAT_LINKED
    end select
  end subroutine process_library_unload

\eatline
\nwendcode{}\nwbegindocs{616}Unload, clean all generated files and revert the library status.  If
{\tt{}distclean} is set, also remove the makefile and the driver source.
\nwenddocs{}\nwbegincode{617}\moddef{Process libraries: process library: TBP}\plusendmoddef
  procedure :: clean => process_library_clean
\nwendcode{}\nwbegincode{618}\moddef{Process libraries: procedures}\plusendmoddef
  subroutine process_library_clean (lib, os_data, distclean)
    class(process_library_t), intent(inout) :: lib
    type(os_data_t), intent(in) :: os_data
    logical, intent(in) :: distclean
    call lib%unload ()
    if (lib%external .and. .not. lib%static) then
       call msg_message ("Process library '" // char (lib%basename) &
            // "': removing old files")
       if (distclean) then
          call lib%driver%distclean (os_data)
       else
          call lib%driver%clean (os_data)
       end if
    end if
    where (lib%entry%i_external /= 0)
       lib%entry%status = STAT_CONFIGURED
    elsewhere
       lib%entry%status = STAT_LINKED
    end where
    if (lib%external) then
       lib%status = STAT_CONFIGURED
    else
       lib%status = STAT_LINKED
    end if
  end subroutine process_library_clean
  
\eatline
\nwendcode{}\nwbegindocs{619}Unload and clean the compiled library and all generated object files, so
force a recompilation.
\nwenddocs{}\nwbegincode{620}\moddef{XXX Process libraries: process library: TBP}\endmoddef
  procedure :: clean_objects => process_library_clean_objects
\nwendcode{}\nwbegincode{621}\moddef{XXX Process libraries: procedures}\endmoddef
  subroutine process_library_clean_objects (lib, os_data)
    class(process_library_t), intent(inout) :: lib
    type(os_data_t), intent(in) :: os_data
    call lib%unload ()
    if (lib%external .and. .not. lib%static) then
       call msg_message ("Process library '" // char (lib%basename) &
            // "': removing compiled objects")
       call lib%driver%clean_library (os_data)
       call lib%driver%clean_objects (os_data)
    end if
    where (lib%entry%i_external /= 0)
       lib%entry%status = STAT_CONFIGURED
    elsewhere
       lib%entry%status = STAT_LINKED
    end where
    if (lib%external) then
       lib%status = STAT_CONFIGURED
    else
       lib%status = STAT_LINKED
    end if
  end subroutine process_library_clean_objects
  
\eatline
\nwendcode{}\nwbegindocs{622}Unload and revert the library status to INITIAL.  This allows for
appending new processes.  No files are deleted.
\nwenddocs{}\nwbegincode{623}\moddef{Process libraries: process library: TBP}\plusendmoddef
  procedure :: open => process_library_open
\nwendcode{}\nwbegincode{624}\moddef{Process libraries: procedures}\plusendmoddef
  subroutine process_library_open (lib)
    class(process_library_t), intent(inout) :: lib
    select case (lib%status)
    case (STAT_OPEN)
    case default
       call lib%unload ()
       if (.not. lib%static) then
          lib%entry%status = STAT_OPEN
          lib%status = STAT_OPEN
          call msg_message ("Process library '" // char (lib%basename) &
               // "': open")
       else
          call msg_error ("Static process library '" // char (lib%basename) &
               // "': processes can't be appended")
       end if
    end select
  end subroutine process_library_open
  
\eatline
\nwendcode{}\nwbegindocs{625}\nwdocspar
\subsection{Use the library}
Return the base name of the library
\nwenddocs{}\nwbegincode{626}\moddef{Process libraries: process library: TBP}\plusendmoddef
  procedure :: get_name => process_library_get_name
\nwendcode{}\nwbegincode{627}\moddef{Process libraries: procedures}\plusendmoddef
  function process_library_get_name (lib) result (name)
    class(process_library_t), intent(in) :: lib
    type(string_t) :: name
    name = lib%basename
  end function process_library_get_name
  
\eatline
\nwendcode{}\nwbegindocs{628}\nwdocspar
Once activated, we view the process library object as an interface for
accessing the matrix elements.
\nwenddocs{}\nwbegincode{629}\moddef{Process libraries: process library: TBP}\plusendmoddef
  procedure :: is_active => process_library_is_active
\nwendcode{}\nwbegincode{630}\moddef{Process libraries: procedures}\plusendmoddef
  function process_library_is_active (lib) result (flag)
    logical :: flag
    class(process_library_t), intent(in) :: lib
    flag = lib%status == STAT_ACTIVE
  end function process_library_is_active
  
\eatline
\nwendcode{}\nwbegindocs{631}Retrieve constants using the process library driver.  We assume that
the process code has been loaded, if external.
\nwenddocs{}\nwbegincode{632}\moddef{Process libraries: process library entry: TBP}\plusendmoddef
  procedure :: fill_constants => process_library_entry_fill_constants
\nwendcode{}\nwbegincode{633}\moddef{Process libraries: procedures}\plusendmoddef
  subroutine process_library_entry_fill_constants (entry, driver, data)
    class(process_library_entry_t), intent(in) :: entry
    class(prclib_driver_t), intent(in) :: driver
    type(process_constants_t), intent(out) :: data
    integer :: i
    if (entry%i_external /= 0) then
       i = entry%i_external
       data%id         = driver%get_process_id (i)
       data%model_name = driver%get_model_name (i)
       data%md5sum     = driver%get_md5sum (i)
       data%openmp_supported = driver%get_openmp_status (i)
       data%n_in  = driver%get_n_in  (i)
       data%n_out = driver%get_n_out (i)
       data%n_flv = driver%get_n_flv (i)
       data%n_hel = driver%get_n_hel (i)
       data%n_col = driver%get_n_col (i)
       data%n_cin = driver%get_n_cin (i)
       data%n_cf  = driver%get_n_cf  (i)
       call driver%set_flv_state (i, data%flv_state)
       call driver%set_hel_state (i, data%hel_state)
       call driver%set_col_state (i, data%col_state, data%ghost_flag)
       call driver%set_color_factors (i, data%color_factors, data%cf_index)
    else
       select type (proc_driver => entry%driver)
       class is (process_driver_internal_t)
          call proc_driver%fill_constants (data)
       end select
    end if
  end subroutine process_library_entry_fill_constants
  
\eatline
\nwendcode{}\nwbegindocs{634}Retrieve the constants and a connected driver for a process,
identified by a process ID and a subprocess index.  We
scan the process entries until we have found a match.
\nwenddocs{}\nwbegincode{635}\moddef{Process libraries: process library: TBP}\plusendmoddef
  procedure :: connect_process => process_library_connect_process
\nwendcode{}\nwbegincode{636}\moddef{Process libraries: procedures}\plusendmoddef
  subroutine process_library_connect_process &
       (lib, id, i_component, data, proc_driver)
    class(process_library_t), intent(in) :: Lib
    type(string_t), intent(in) :: id
    integer, intent(in) :: i_component
    type(process_constants_t), intent(out) :: data
    class(prc_core_driver_t), allocatable, intent(out) :: proc_driver
    integer :: i
    do i = 1, size (lib%entry)
       associate (entry => lib%entry(i))
         if (entry%def%id == id .and. entry%i_component == i_component) then
            call entry%fill_constants (lib%driver, data)
            allocate (proc_driver, source=entry%driver)
            return
         end if
       end associate
    end do
    call msg_fatal ("Process library '" // char (lib%basename) &
               // "': process '" // char (id) // "' not found")
  end subroutine process_library_connect_process

\eatline
\nwendcode{}\nwbegindocs{637}\nwdocspar
\nwenddocs{}\nwbegindocs{638}\nwdocspar
\subsection{Collect model-specific libraries}
This returns appropriate linker flags for the model parameter libraries that
are used by the generated matrix element.  At the end, the main libwhizard is
appended (again), because functions from that may be reqired.

Extra models in the local user space need to be treated individually.
\nwenddocs{}\nwbegincode{639}\moddef{Process libraries: process library: TBP}\plusendmoddef
  procedure :: get_modellibs_ldflags => process_library_get_modellibs_ldflags
\nwendcode{}\nwbegincode{640}\moddef{Process libraries: procedures}\plusendmoddef
  function process_library_get_modellibs_ldflags (prc_lib, os_data) result (flags)
    class(process_library_t), intent(in) :: prc_lib
    type(os_data_t), intent(in) :: os_data
    type(string_t) :: flags
    type(string_t), dimension(:), allocatable :: models
    type(string_t) :: modelname, modellib, modellib_full
    logical :: exist
    integer :: i, j, mi
    flags = " -lomega"
    if ((.not. os_data%use_testfiles) .and. &
               os_dir_exist (os_data%whizard_models_libpath_local)) &
                 flags = flags // " -L" // os_data%whizard_models_libpath_local
    flags = flags // " -L" // os_data%whizard_models_libpath
    allocate (models(prc_lib%n_entries + 1))
    models = ""
    mi = 1
    if (allocated (prc_lib%entry)) then
       SCAN: do i = 1, prc_lib%n_entries
          if (associated (prc_lib%entry(i)%def)) then
             if (prc_lib%entry(i)%def%model_name /= "") then
                modelname = prc_lib%entry(i)%def%model_name
             else
                cycle SCAN
             end if
          else
             cycle SCAN
          end if
          do j = 1, mi
             if (models(mi) == modelname) cycle SCAN
          end do
          models(mi) = modelname
          mi = mi + 1
          if (os_data%use_libtool) then
             modellib = "libparameters_" // modelname // ".la"
          else
             modellib = "libparameters_" // modelname // ".a"
          end if          
          exist = .false.
          if (.not. os_data%use_testfiles) then
             modellib_full = os_data%whizard_models_libpath_local &
                  // "/" // modellib
             inquire (file=char (modellib_full), exist=exist)
          end if
          if (.not. exist) then
             modellib_full = os_data%whizard_models_libpath &
                  // "/" // modellib
             inquire (file=char (modellib_full), exist=exist)
          end if
          if (exist) flags = flags // " -lparameters_" // modelname
       end do SCAN
    end if
    deallocate (models)
    flags = flags // " -lwhizard"
  end function process_library_get_modellibs_ldflags

\eatline
\nwendcode{}\nwbegindocs{641}\nwdocspar
\nwenddocs{}\nwbegincode{642}\moddef{Process libraries: process library: TBP}\plusendmoddef
  procedure :: get_static_modelname => process_library_get_static_modelname
\nwendcode{}\nwbegincode{643}\moddef{Process libraries: procedures}\plusendmoddef
  function process_library_get_static_modelname (prc_lib, os_data) result (name)
    class(process_library_t), intent(in) :: prc_lib
    type(os_data_t), intent(in) :: os_data
    type(string_t) :: name
    type(string_t), dimension(:), allocatable :: models
    type(string_t) :: modelname, modellib, modellib_full
    logical :: exist
    integer :: i, j, mi
    name = ""
    allocate (models(prc_lib%n_entries + 1))
    models = ""
    mi = 1
    if (allocated (prc_lib%entry)) then
       SCAN: do i = 1, prc_lib%n_entries
          if (associated (prc_lib%entry(i)%def)) then
             if (prc_lib%entry(i)%def%model_name /= "") then
                modelname = prc_lib%entry(i)%def%model_name
             else
                cycle SCAN
             end if
          else
             cycle SCAN
          end if
          do j = 1, mi
             if (models(mi) == modelname) cycle SCAN
          end do
          models(mi) = modelname
          mi = mi + 1
          modellib = "libparameters_" // modelname // ".a"
          exist = .false.
          if (.not. os_data%use_testfiles) then
             modellib_full = os_data%whizard_models_libpath_local &
                  // "/" // modellib
             inquire (file=char (modellib_full), exist=exist)
          end if
          if (.not. exist) then
             modellib_full = os_data%whizard_models_libpath &
                  // "/" // modellib
             inquire (file=char (modellib_full), exist=exist)
          end if
          if (exist) name = name // " " // modellib_full
       end do SCAN
    end if
    deallocate (models)
  end function process_library_get_static_modelname

\eatline
\nwendcode{}\nwbegindocs{644}\nwdocspar
\subsection{Test}
This is the master for calling self-test procedures.
\nwenddocs{}\nwbegincode{645}\moddef{Process libraries: public}\plusendmoddef
  public :: process_libraries_test
\nwendcode{}\nwbegincode{646}\moddef{Process libraries: tests}\endmoddef
  subroutine process_libraries_test (u, results)
    integer, intent(in) :: u
    type(test_results_t), intent(inout) :: results
  \LA{}Process libraries: execute tests\RA{}
  end subroutine process_libraries_test

\eatline
\nwendcode{}\nwbegindocs{647}\nwdocspar
\subsubsection{Empty process list}
Test 1: Write an empty process list.
\nwenddocs{}\nwbegincode{648}\moddef{Process libraries: execute tests}\endmoddef
  call test (process_libraries_1, "process_libraries_1", &
       "empty process list", &
       u, results)
\nwendcode{}\nwbegincode{649}\moddef{Process libraries: tests}\plusendmoddef
  subroutine process_libraries_1 (u)
    integer, intent(in) :: u
    type(process_def_list_t) :: process_def_list

    write (u, "(A)")  "* Test output: process_libraries_1"
    write (u, "(A)")  "*   Purpose: Display an empty process definition list"
    write (u, "(A)")

    call process_def_list%write (u)

    write (u, "(A)")
    write (u, "(A)")  "* Test output end: process_libraries_1"
  end subroutine process_libraries_1
  
\eatline
\nwendcode{}\nwbegindocs{650}\nwdocspar
\subsubsection{Process definition list}
Test 2: Process definition list with processes and components.
Construct the list, write to file, read it in again, and display.
Finalize and delete the list after use.

We define a trivial 'test' type for the process variant.  The test
type contains just one (meaningless) data item, which is an integer.
\nwenddocs{}\nwbegincode{651}\moddef{Process libraries: test declarations}\endmoddef
  type, extends (prc_core_def_t) :: prcdef_2_t
     integer :: data = 0
     logical :: file = .false.
   contains
   \LA{}Process libraries: prcdef 2: TBP\RA{}
  end type prcdef_2_t
  
\eatline
\nwendcode{}\nwbegindocs{652}The process variant is named 'test'.
\nwenddocs{}\nwbegincode{653}\moddef{Process libraries: prcdef 2: TBP}\endmoddef
  procedure, nopass :: type_string => prcdef_2_type_string
\nwendcode{}\nwbegincode{654}\moddef{Process libraries: tests}\plusendmoddef
  function prcdef_2_type_string () result (string)
    type(string_t) :: string
    string = "test"
  end function prcdef_2_type_string

\eatline
\nwendcode{}\nwbegindocs{655}Write the contents (the integer value).
\nwenddocs{}\nwbegincode{656}\moddef{Process libraries: prcdef 2: TBP}\plusendmoddef
  procedure :: write => prcdef_2_write
\nwendcode{}\nwbegincode{657}\moddef{Process libraries: tests}\plusendmoddef
  subroutine prcdef_2_write (object, unit)
    class(prcdef_2_t), intent(in) :: object
    integer, intent(in) :: unit
    write (unit, "(3x,A,I0)")  "Test data         = ", object%data
  end subroutine prcdef_2_write
  
\eatline
\nwendcode{}\nwbegindocs{658}Recover the integer value.
\nwenddocs{}\nwbegincode{659}\moddef{Process libraries: prcdef 2: TBP}\plusendmoddef
  procedure :: read => prcdef_2_read
\nwendcode{}\nwbegincode{660}\moddef{Process libraries: tests}\plusendmoddef
  subroutine prcdef_2_read (object, unit)
    class(prcdef_2_t), intent(out) :: object
    integer, intent(in) :: unit
    character(80) :: buffer
    read (unit, "(A)")  buffer
    call strip_prefix (buffer)
    read (buffer, *)  object%data
  end subroutine prcdef_2_read
  
\eatline
\nwendcode{}\nwbegindocs{661}No external procedures.
\nwenddocs{}\nwbegincode{662}\moddef{Process libraries: prcdef 2: TBP}\plusendmoddef
  procedure, nopass :: get_features => prcdef_2_get_features
\nwendcode{}\nwbegincode{663}\moddef{Process libraries: tests}\plusendmoddef
  subroutine prcdef_2_get_features (features)
    type(string_t), dimension(:), allocatable, intent(out) :: features
    allocate (features (0))
  end subroutine prcdef_2_get_features

\eatline
\nwendcode{}\nwbegindocs{664}No code generated.
\nwenddocs{}\nwbegincode{665}\moddef{Process libraries: prcdef 2: TBP}\plusendmoddef
  procedure :: generate_code => prcdef_2_generate_code
\nwendcode{}\nwbegincode{666}\moddef{Process libraries: tests}\plusendmoddef
  subroutine prcdef_2_generate_code (object, &
       basename, model_name, prt_in, prt_out)
    class(prcdef_2_t), intent(in) :: object
    type(string_t), intent(in) :: basename
    type(string_t), intent(in) :: model_name
    type(string_t), dimension(:), intent(in) :: prt_in
    type(string_t), dimension(:), intent(in) :: prt_out
  end subroutine prcdef_2_generate_code
  
\eatline
\nwendcode{}\nwbegindocs{667}Allocate the driver with the appropriate type.
\nwenddocs{}\nwbegincode{668}\moddef{Process libraries: prcdef 2: TBP}\plusendmoddef
  procedure :: allocate_driver => prcdef_2_allocate_driver
\nwendcode{}\nwbegincode{669}\moddef{Process libraries: tests}\plusendmoddef
  subroutine prcdef_2_allocate_driver (object, driver, basename)
    class(prcdef_2_t), intent(in) :: object
    class(prc_core_driver_t), intent(out), allocatable :: driver
    type(string_t), intent(in) :: basename
    allocate (prctest_2_t :: driver)
  end subroutine prcdef_2_allocate_driver
  
\eatline
\nwendcode{}\nwbegindocs{670}Nothing to connect.
\nwenddocs{}\nwbegincode{671}\moddef{Process libraries: prcdef 2: TBP}\plusendmoddef
  procedure :: connect => prcdef_2_connect
\nwendcode{}\nwbegincode{672}\moddef{Process libraries: tests}\plusendmoddef
  subroutine prcdef_2_connect (def, lib_driver, i, proc_driver)
    class(prcdef_2_t), intent(in) :: def
    class(prclib_driver_t), intent(in) :: lib_driver
    integer, intent(in) :: i
    class(prc_core_driver_t), intent(inout) :: proc_driver
  end subroutine prcdef_2_connect

\eatline
\nwendcode{}\nwbegindocs{673}The associated driver type.
\nwenddocs{}\nwbegincode{674}\moddef{Process libraries: types}\plusendmoddef
  type, extends (process_driver_internal_t) :: prctest_2_t
   contains
   \LA{}Process libraries: prctest 2: TBP\RA{}
  end type prctest_2_t
  
\eatline
\nwendcode{}\nwbegindocs{675}Return the type name.
\nwenddocs{}\nwbegincode{676}\moddef{Process libraries: prctest 2: TBP}\endmoddef
  procedure, nopass :: type_name => prctest_2_type_name
\nwendcode{}\nwbegincode{677}\moddef{Process libraries: procedures}\plusendmoddef
  function prctest_2_type_name () result (type)
    type(string_t) :: type
    type = "test"
  end function prctest_2_type_name
  
\eatline
\nwendcode{}\nwbegindocs{678}This should fill constant process data.  We do not check those here,
however, therefore nothing done.
\nwenddocs{}\nwbegincode{679}\moddef{Process libraries: prctest 2: TBP}\plusendmoddef
  procedure :: fill_constants => prctest_2_fill_constants
\nwendcode{}\nwbegincode{680}\moddef{Process libraries: procedures}\plusendmoddef
  subroutine prctest_2_fill_constants (driver, data)
    class(prctest_2_t), intent(in) :: driver
    type(process_constants_t), intent(out) :: data
  end subroutine prctest_2_fill_constants

\eatline
\nwendcode{}\nwbegindocs{681}\nwdocspar
Here is the actual test.

For reading, we need a list of templates, i.e., an array containing
allocated objects for all available process variants.  This is the
purpose of {\tt{}process{\char95}core{\char95}templates}.  Here, we have only a single
template for the 'test' variant.
\nwenddocs{}\nwbegincode{682}\moddef{Process libraries: execute tests}\plusendmoddef
  call test (process_libraries_2, "process_libraries_2", &
       "process definition list", &
       u, results)
\nwendcode{}\nwbegincode{683}\moddef{Process libraries: tests}\plusendmoddef
  subroutine process_libraries_2 (u)
    integer, intent(in) :: u
    type(prc_template_t), dimension(:), allocatable :: process_core_templates
    type(process_def_list_t) :: process_def_list
    type(process_def_entry_t), pointer :: entry => null ()
    class(prc_core_def_t), allocatable :: test_def
    integer :: scratch_unit

    write (u, "(A)")  "* Test output: process_libraries_2"
    write (u, "(A)")  "* Purpose: Construct a process definition list,"
    write (u, "(A)")  "*          write it to file and reread it"
    write (u, "(A)")  ""
    write (u, "(A)")  "* Construct a process definition list"
    write (u, "(A)")  "*   First process definition: empty"
    write (u, "(A)")  "*   Second process definition: two components"
    write (u, "(A)")  "*     First component: empty"
    write (u, "(A)")  "*     Second component: test data"
    write (u, "(A)")  "*   Third process definition:"
    write (u, "(A)")  "*     Embedded decays and polarization"
    write (u, "(A)")

    allocate (process_core_templates (1))
    allocate (prcdef_2_t :: process_core_templates(1)%core_def)
    
    allocate (entry)
    call entry%init (var_str ("first"), n_in = 0, n_components = 0)
    call entry%compute_md5sum ()
    call process_def_list%append (entry)
    
    allocate (entry)
    call entry%init (var_str ("second"), model_name = var_str ("Test"), &
         n_in = 1, n_components = 2)
    allocate (prcdef_2_t :: test_def)
    select type (test_def)
    type is (prcdef_2_t);  test_def%data = 42
    end select
    call entry%import_component (2, n_out = 2, &
         prt_in  = new_prt_spec ([var_str ("a")]), &
         prt_out = new_prt_spec ([var_str ("b"), var_str ("c")]), &
         method  = var_str ("test"), &
         variant = test_def)
    call entry%compute_md5sum ()
    call process_def_list%append (entry)

    allocate (entry)
    call entry%init (var_str ("third"), model_name = var_str ("Test"), &
         n_in = 2, n_components = 1)
    allocate (prcdef_2_t :: test_def)
    call entry%import_component (1, n_out = 3, &
         prt_in  = &
           new_prt_spec ([var_str ("a"), var_str ("b")]), &
         prt_out = &
           [new_prt_spec (var_str ("c")), &
            new_prt_spec (var_str ("d"), .true.), &
            new_prt_spec (var_str ("e"), [var_str ("e_decay")])], &
         method  = var_str ("test"), &
         variant = test_def)
    call entry%compute_md5sum ()
    call process_def_list%append (entry)
    call process_def_list%write (u)

    write (u, "(A)")  ""
    write (u, "(A)")  "* Write the process definition list to (scratch) file"

    scratch_unit = free_unit ()
    open (unit = scratch_unit, status="scratch", action = "readwrite")
    call process_def_list%write (scratch_unit)
    call process_def_list%final ()
    
    write (u, "(A)")  "* Reread it"
    write (u, "(A)")  ""

    rewind (scratch_unit)
    call process_def_list%read (scratch_unit, process_core_templates)
    close (scratch_unit)
    
    call process_def_list%write (u)
    call process_def_list%final ()

    write (u, "(A)")
    write (u, "(A)")  "* Test output end: process_libraries_2"
  end subroutine process_libraries_2
  
\eatline
\nwendcode{}\nwbegindocs{684}\nwdocspar
\subsubsection{Process library object}
Test 3: Process library object with several process definitions and
library entries.  Just construct the object, modify some initial
content, and write the result.  The modifications are mostly applied
directly, so we do not test anything but the contents and the output
routine.
\nwenddocs{}\nwbegincode{685}\moddef{Process libraries: execute tests}\plusendmoddef
  call test (process_libraries_3, "process_libraries_3", &
       "recover process definition list from file", &
       u, results)
\nwendcode{}\nwbegincode{686}\moddef{Process libraries: tests}\plusendmoddef
  subroutine process_libraries_3 (u)
    integer, intent(in) :: u
    type(process_library_t) :: lib
    type(process_def_entry_t), pointer :: entry

    write (u, "(A)")  "* Test output: process_libraries_3"
    write (u, "(A)")  "* Purpose: Construct a process library object &
         &with entries"
    write (u, "(A)")  ""
    write (u, "(A)")  "* Construct and display a process library object"
    write (u, "(A)")  "*   with 5 entries"
    write (u, "(A)")  "*   associated with 3 matrix element codes"
    write (u, "(A)")  "*   corresponding to 3 process definitions"
    write (u, "(A)")  "*   with 2, 1, 1 components, respectively"
    write (u, "(A)")

    call lib%init (var_str ("testlib"))

    lib%status = STAT_ACTIVE
    lib%n_entries = 5
    allocate (lib%entry (lib%n_entries))
   
    allocate (entry)
    call entry%init (var_str ("test_a"), n_in = 2, n_components = 2)
    call lib%entry(3)%init (STAT_SOURCE, entry%process_def_t, 2, 2)
    allocate (prctest_2_t :: lib%entry(3)%driver)
    call lib%entry(4)%init (STAT_COMPILED, entry%process_def_t, 1, 0)
    call lib%append (entry)

    allocate (entry)
    call entry%init (var_str ("test_b"), n_in = 2, n_components = 1)
    call lib%entry(2)%init (STAT_CONFIGURED, entry%process_def_t, 1, 1)
    call lib%append (entry)

    allocate (entry)
    call entry%init (var_str ("test_c"), n_in = 2, n_components = 1)
    call lib%entry(5)%init (STAT_LINKED, entry%process_def_t, 1, 3)
    allocate (prctest_2_t :: lib%entry(5)%driver)
    call lib%append (entry)
    
    call lib%write (u)
    call lib%final ()
    
    write (u, "(A)")
    write (u, "(A)")  "* Test output end: process_libraries_3"
  end subroutine process_libraries_3
  
\eatline
\nwendcode{}\nwbegindocs{687}\nwdocspar
\subsubsection{Process library for test matrix element (no file)}
Test 4: We proceed through the library generation and loading phases
with a test matrix element type that needs no code written on file.
\nwenddocs{}\nwbegincode{688}\moddef{Process libraries: execute tests}\plusendmoddef
  call test (process_libraries_4, "process_libraries_4", &
       "build and load internal process library", &
       u, results)
\nwendcode{}\nwbegincode{689}\moddef{Process libraries: tests}\plusendmoddef
  subroutine process_libraries_4 (u)
    integer, intent(in) :: u
    type(process_library_t) :: lib
    type(process_def_entry_t), pointer :: entry
    class(prc_core_def_t), allocatable :: core_def
    type(os_data_t) :: os_data

    write (u, "(A)")  "* Test output: process_libraries_4"
    write (u, "(A)")  "* Purpose: build a process library with an &
         &internal (pseudo) matrix element"
    write (u, "(A)")  "*          No Makefile or code should be generated"
    write (u, "(A)")

    write (u, "(A)")  "* Initialize a process library with one entry &
         &(no external code)"
    write (u, "(A)")
    call os_data_init (os_data)
    call lib%init (var_str ("proclibs4"))

    allocate (prcdef_2_t :: core_def)

    allocate (entry)
    call entry%init (var_str ("proclibs4_a"), n_in = 1, n_components = 1)
    call entry%import_component (1, n_out = 2, variant = core_def)
    call lib%append (entry)

    write (u, "(A)")  "* Configure library"
    write (u, "(A)")
    call lib%configure (os_data)

    write (u, "(A)")  "* Compute MD5 sum"
    write (u, "(A)")
    call lib%compute_md5sum ()

    write (u, "(A)")  "* Write makefile (no-op)"
    write (u, "(A)")
    call lib%write_makefile (os_data, force = .true.)

    write (u, "(A)")  "* Write driver source code (no-op)"
    write (u, "(A)")
    call lib%write_driver (force = .true.)

    write (u, "(A)")  "* Write process source code (no-op)"
    write (u, "(A)")
    call lib%make_source (os_data)

    write (u, "(A)")  "* Compile (no-op)"
    write (u, "(A)")
    call lib%make_compile (os_data)

    write (u, "(A)")  "* Link (no-op)"
    write (u, "(A)")
    call lib%make_link (os_data)

    write (u, "(A)")  "* Load (no-op)"
    write (u, "(A)")
    call lib%load (os_data)

    call lib%write (u)
    call lib%final ()
    
    write (u, "(A)")
    write (u, "(A)")  "* Test output end: process_libraries_4"
  end subroutine process_libraries_4
  
\eatline
\nwendcode{}\nwbegindocs{690}\nwdocspar
\subsubsection{Build workflow for test matrix element}
Test 5: We write source code for a dummy process.

We define another trivial type for the process variant.  The test
type contains just no variable data, but produces code on file.
\nwenddocs{}\nwbegincode{691}\moddef{Process libraries: test declarations}\plusendmoddef
  type, extends (prc_core_def_t) :: prcdef_5_t
   contains
   \LA{}Process libraries: prcdef 5: TBP\RA{}
  end type prcdef_5_t
  
\eatline
\nwendcode{}\nwbegindocs{692}The process variant is named {\tt{}test{\char95}file}.
\nwenddocs{}\nwbegincode{693}\moddef{Process libraries: prcdef 5: TBP}\endmoddef
  procedure, nopass :: type_string => prcdef_5_type_string
\nwendcode{}\nwbegincode{694}\moddef{Process libraries: tests}\plusendmoddef
  function prcdef_5_type_string () result (string)
    type(string_t) :: string
    string = "test_file"
  end function prcdef_5_type_string

\eatline
\nwendcode{}\nwbegindocs{695}We reuse the writer {\tt{}test{\char95}writer{\char95}4} from the previous module.
\nwenddocs{}\nwbegincode{696}\moddef{Process libraries: prcdef 5: TBP}\plusendmoddef
  procedure :: init => prcdef_5_init
\nwendcode{}\nwbegincode{697}\moddef{Process libraries: tests}\plusendmoddef
  subroutine prcdef_5_init (object)
    class(prcdef_5_t), intent(out) :: object
    allocate (test_writer_4_t :: object%writer)
  end subroutine prcdef_5_init
  
\eatline
\nwendcode{}\nwbegindocs{698}Nothing to write.
\nwenddocs{}\nwbegincode{699}\moddef{Process libraries: prcdef 5: TBP}\plusendmoddef
  procedure :: write => prcdef_5_write
\nwendcode{}\nwbegincode{700}\moddef{Process libraries: tests}\plusendmoddef
  subroutine prcdef_5_write (object, unit)
    class(prcdef_5_t), intent(in) :: object
    integer, intent(in) :: unit
  end subroutine prcdef_5_write
  
\eatline
\nwendcode{}\nwbegindocs{701}Nothing to read.
\nwenddocs{}\nwbegincode{702}\moddef{Process libraries: prcdef 5: TBP}\plusendmoddef
  procedure :: read => prcdef_5_read
\nwendcode{}\nwbegincode{703}\moddef{Process libraries: tests}\plusendmoddef
  subroutine prcdef_5_read (object, unit)
    class(prcdef_5_t), intent(out) :: object
    integer, intent(in) :: unit
  end subroutine prcdef_5_read
  
\eatline
\nwendcode{}\nwbegindocs{704}Allocate the driver with the appropriate type.
\nwenddocs{}\nwbegincode{705}\moddef{Process libraries: prcdef 5: TBP}\plusendmoddef
  procedure :: allocate_driver => prcdef_5_allocate_driver
\nwendcode{}\nwbegincode{706}\moddef{Process libraries: tests}\plusendmoddef
  subroutine prcdef_5_allocate_driver (object, driver, basename)
    class(prcdef_5_t), intent(in) :: object
    class(prc_core_driver_t), intent(out), allocatable :: driver
    type(string_t), intent(in) :: basename
    allocate (prctest_5_t :: driver)
  end subroutine prcdef_5_allocate_driver
  
\eatline
\nwendcode{}\nwbegindocs{707}This time we need code:
\nwenddocs{}\nwbegincode{708}\moddef{Process libraries: prcdef 5: TBP}\plusendmoddef
  procedure, nopass :: needs_code => prcdef_5_needs_code
\nwendcode{}\nwbegincode{709}\moddef{Process libraries: tests}\plusendmoddef
  function prcdef_5_needs_code () result (flag)
    logical :: flag
    flag = .true.
  end function prcdef_5_needs_code
  
\eatline
\nwendcode{}\nwbegindocs{710}For the test case, we implement a single feature {\tt{}proc1}.
\nwenddocs{}\nwbegincode{711}\moddef{Process libraries: prcdef 5: TBP}\plusendmoddef
  procedure, nopass :: get_features => prcdef_5_get_features
\nwendcode{}\nwbegincode{712}\moddef{Process libraries: tests}\plusendmoddef
  subroutine prcdef_5_get_features (features)
    type(string_t), dimension(:), allocatable, intent(out) :: features
    allocate (features (1))
    features = [ var_str ("proc1") ]
  end subroutine prcdef_5_get_features

\eatline
\nwendcode{}\nwbegindocs{713}Nothing to connect.
\nwenddocs{}\nwbegincode{714}\moddef{Process libraries: prcdef 5: TBP}\plusendmoddef
  procedure :: connect => prcdef_5_connect
\nwendcode{}\nwbegincode{715}\moddef{Process libraries: tests}\plusendmoddef
  subroutine prcdef_5_connect (def, lib_driver, i, proc_driver)
    class(prcdef_5_t), intent(in) :: def
    class(prclib_driver_t), intent(in) :: lib_driver
    integer, intent(in) :: i
    class(prc_core_driver_t), intent(inout) :: proc_driver
  end subroutine prcdef_5_connect

\eatline
\nwendcode{}\nwbegindocs{716}The driver type.
\nwenddocs{}\nwbegincode{717}\moddef{Process libraries: types}\plusendmoddef
  type, extends (prc_core_driver_t) :: prctest_5_t
   contains
   \LA{}Process libraries: prctest 5: TBP\RA{}
  end type prctest_5_t
  
\eatline
\nwendcode{}\nwbegindocs{718}Return the type name.
\nwenddocs{}\nwbegincode{719}\moddef{Process libraries: prctest 5: TBP}\endmoddef
  procedure, nopass :: type_name => prctest_5_type_name
\nwendcode{}\nwbegincode{720}\moddef{Process libraries: tests}\plusendmoddef
  function prctest_5_type_name () result (type)
    type(string_t) :: type
    type = "test_file"
  end function prctest_5_type_name
  
\eatline
\nwendcode{}\nwbegindocs{721}\nwdocspar
Here is the actual test:
\nwenddocs{}\nwbegincode{722}\moddef{Process libraries: execute tests}\plusendmoddef
  call test (process_libraries_5, "process_libraries_5", &
       "build external process library", &
       u, results)
\nwendcode{}\nwbegincode{723}\moddef{Process libraries: tests}\plusendmoddef
  subroutine process_libraries_5 (u)
    integer, intent(in) :: u
    type(process_library_t) :: lib
    type(process_def_entry_t), pointer :: entry
    class(prc_core_def_t), allocatable :: core_def
    type(os_data_t) :: os_data

    write (u, "(A)")  "* Test output: process_libraries_5"
    write (u, "(A)")  "* Purpose: build a process library with an &
         &external (pseudo) matrix element"
    write (u, "(A)")

    write (u, "(A)")  "* Initialize a process library with one entry"
    write (u, "(A)")
    call lib%init (var_str ("proclibs5"))
    call os_data_init (os_data)

    allocate (prcdef_5_t :: core_def)
    select type (core_def)
    type is (prcdef_5_t)
       call core_def%init ()
    end select

    allocate (entry)
    call entry%init (var_str ("proclibs5_a"), &
         model_name = var_str ("Test_Model"), &
         n_in = 1, n_components = 1)
    call entry%import_component (1, n_out = 2, &
         prt_in  = new_prt_spec ([var_str ("a")]), &
         prt_out = new_prt_spec ([var_str ("b"), var_str ("c")]), &
         method  = var_str ("test"), &
         variant = core_def)
    call lib%append (entry)
    
    write (u, "(A)")  "* Configure library"
    write (u, "(A)")
    call lib%configure (os_data)

    write (u, "(A)")  "* Compute MD5 sum"
    write (u, "(A)")
    call lib%compute_md5sum ()

    write (u, "(A)")  "* Write makefile"
    write (u, "(A)")
    call lib%write_makefile (os_data, force = .true.)
    
    write (u, "(A)")  "* Write driver source code"
    write (u, "(A)")
    call lib%write_driver (force = .true.)

    write (u, "(A)")  "* Write process source code"
    write (u, "(A)")
    call lib%make_source (os_data)

    write (u, "(A)")  "* Compile"
    write (u, "(A)")
    call lib%make_compile (os_data)

    write (u, "(A)")  "* Link"
    write (u, "(A)")
    call lib%make_link (os_data)

    call lib%write (u, libpath = .false.)
    
    call lib%final ()

    write (u, "(A)")
    write (u, "(A)")  "* Test output end: process_libraries_5"
  end subroutine process_libraries_5
  
\eatline
\nwendcode{}\nwbegindocs{724}\nwdocspar
\subsubsection{Build and load library with test matrix element}
Test 6: We write source code for a dummy process.

This process variant is identical to the previous case, but it
supports a driver for the test procedure 'proc1'.
\nwenddocs{}\nwbegincode{725}\moddef{Process libraries: test declarations}\plusendmoddef
  type, extends (prc_core_def_t) :: prcdef_6_t
   contains
   \LA{}Process libraries: prcdef 6: TBP\RA{}
  end type prcdef_6_t
  
\eatline
\nwendcode{}\nwbegindocs{726}The process variant is named {\tt{}test{\char95}file}.
\nwenddocs{}\nwbegincode{727}\moddef{Process libraries: prcdef 6: TBP}\endmoddef
  procedure, nopass :: type_string => prcdef_6_type_string
\nwendcode{}\nwbegincode{728}\moddef{Process libraries: tests}\plusendmoddef
  function prcdef_6_type_string () result (string)
    type(string_t) :: string
    string = "test_file"
  end function prcdef_6_type_string

\eatline
\nwendcode{}\nwbegindocs{729}We reuse the writer {\tt{}test{\char95}writer{\char95}4} from the previous module.
\nwenddocs{}\nwbegincode{730}\moddef{Process libraries: prcdef 6: TBP}\plusendmoddef
  procedure :: init => prcdef_6_init
\nwendcode{}\nwbegincode{731}\moddef{Process libraries: tests}\plusendmoddef
  subroutine prcdef_6_init (object)
    class(prcdef_6_t), intent(out) :: object
    allocate (test_writer_4_t :: object%writer)
    call object%writer%init_test ()
  end subroutine prcdef_6_init
  
\eatline
\nwendcode{}\nwbegindocs{732}Nothing to write.
\nwenddocs{}\nwbegincode{733}\moddef{Process libraries: prcdef 6: TBP}\plusendmoddef
  procedure :: write => prcdef_6_write
\nwendcode{}\nwbegincode{734}\moddef{Process libraries: tests}\plusendmoddef
  subroutine prcdef_6_write (object, unit)
    class(prcdef_6_t), intent(in) :: object
    integer, intent(in) :: unit
  end subroutine prcdef_6_write
  
\eatline
\nwendcode{}\nwbegindocs{735}Nothing to read.
\nwenddocs{}\nwbegincode{736}\moddef{Process libraries: prcdef 6: TBP}\plusendmoddef
  procedure :: read => prcdef_6_read
\nwendcode{}\nwbegincode{737}\moddef{Process libraries: tests}\plusendmoddef
  subroutine prcdef_6_read (object, unit)
    class(prcdef_6_t), intent(out) :: object
    integer, intent(in) :: unit
  end subroutine prcdef_6_read
  
\eatline
\nwendcode{}\nwbegindocs{738}Allocate the driver with the appropriate type.
\nwenddocs{}\nwbegincode{739}\moddef{Process libraries: prcdef 6: TBP}\plusendmoddef
  procedure :: allocate_driver => prcdef_6_allocate_driver
\nwendcode{}\nwbegincode{740}\moddef{Process libraries: tests}\plusendmoddef
  subroutine prcdef_6_allocate_driver (object, driver, basename)
    class(prcdef_6_t), intent(in) :: object
    class(prc_core_driver_t), intent(out), allocatable :: driver
    type(string_t), intent(in) :: basename
    allocate (prctest_6_t :: driver)
  end subroutine prcdef_6_allocate_driver
  
\eatline
\nwendcode{}\nwbegindocs{741}This time we need code:
\nwenddocs{}\nwbegincode{742}\moddef{Process libraries: prcdef 6: TBP}\plusendmoddef
  procedure, nopass :: needs_code => prcdef_6_needs_code
\nwendcode{}\nwbegincode{743}\moddef{Process libraries: tests}\plusendmoddef
  function prcdef_6_needs_code () result (flag)
    logical :: flag
    flag = .true.
  end function prcdef_6_needs_code
  
\eatline
\nwendcode{}\nwbegindocs{744}For the test case, we implement a single feature {\tt{}proc1}.
\nwenddocs{}\nwbegincode{745}\moddef{Process libraries: prcdef 6: TBP}\plusendmoddef
  procedure, nopass :: get_features => prcdef_6_get_features
\nwendcode{}\nwbegincode{746}\moddef{Process libraries: tests}\plusendmoddef
  subroutine prcdef_6_get_features (features)
    type(string_t), dimension(:), allocatable, intent(out) :: features
    allocate (features (1))
    features = [ var_str ("proc1") ]
  end subroutine prcdef_6_get_features

\eatline
\nwendcode{}\nwbegindocs{747}The interface of the only specific feature.
\nwenddocs{}\nwbegincode{748}\moddef{Process libraries: interfaces}\endmoddef
  abstract interface
     subroutine proc1_t (n) bind(C)
       import
       integer(c_int), intent(out) :: n
     end subroutine proc1_t
  end interface

\eatline
\nwendcode{}\nwbegindocs{749}Connect the feature {\tt{}proc1} with the process driver.
\nwenddocs{}\nwbegincode{750}\moddef{Process libraries: prcdef 6: TBP}\plusendmoddef
  procedure :: connect => prcdef_6_connect
\nwendcode{}\nwbegincode{751}\moddef{Process libraries: tests}\plusendmoddef
  subroutine prcdef_6_connect (def, lib_driver, i, proc_driver)
    class(prcdef_6_t), intent(in) :: def
    class(prclib_driver_t), intent(in) :: lib_driver
    integer, intent(in) :: i
    class(prc_core_driver_t), intent(inout) :: proc_driver
    integer(c_int) :: pid, fid
    type(c_funptr) :: fptr
    select type (proc_driver)
    type is  (prctest_6_t)
       pid = i
       fid = 1
       call lib_driver%get_fptr (pid, fid, fptr)
       call c_f_procpointer (fptr, proc_driver%proc1)
    end select
  end subroutine prcdef_6_connect

\eatline
\nwendcode{}\nwbegindocs{752}\nwdocspar
The driver type.
\nwenddocs{}\nwbegincode{753}\moddef{Process libraries: types}\plusendmoddef
  type, extends (prc_core_driver_t) :: prctest_6_t
     procedure(proc1_t), nopass, pointer :: proc1 => null ()
   contains
   \LA{}Process libraries: prctest 6: TBP\RA{}
  end type prctest_6_t
  
\eatline
\nwendcode{}\nwbegindocs{754}Return the type name.
\nwenddocs{}\nwbegincode{755}\moddef{Process libraries: prctest 6: TBP}\endmoddef
  procedure, nopass :: type_name => prctest_6_type_name
\nwendcode{}\nwbegincode{756}\moddef{Process libraries: tests}\plusendmoddef
  function prctest_6_type_name () result (type)
    type(string_t) :: type
    type = "test_file"
  end function prctest_6_type_name
  
\eatline
\nwendcode{}\nwbegindocs{757}\nwdocspar
Here is the actual test:
\nwenddocs{}\nwbegincode{758}\moddef{Process libraries: execute tests}\plusendmoddef
  call test (process_libraries_6, "process_libraries_6", &
       "build and load external process library", &
       u, results)
\nwendcode{}\nwbegincode{759}\moddef{Process libraries: tests}\plusendmoddef
  subroutine process_libraries_6 (u)
    integer, intent(in) :: u
    type(process_library_t) :: lib
    type(process_def_entry_t), pointer :: entry
    class(prc_core_def_t), allocatable :: core_def
    type(os_data_t) :: os_data
    type(string_t), dimension(:), allocatable :: name_list
    type(process_constants_t) :: data
    class(prc_core_driver_t), allocatable :: proc_driver
    integer :: i
    integer(c_int) :: n

    write (u, "(A)")  "* Test output: process_libraries_6"
    write (u, "(A)")  "* Purpose: build and load a process library"
    write (u, "(A)")  "*          with an external (pseudo) matrix element"
    write (u, "(A)")  "*          Check single-call linking"
    write (u, "(A)")

    write (u, "(A)")  "* Initialize a process library with one entry"
    write (u, "(A)")
    call lib%init (var_str ("proclibs6"))
    call os_data_init (os_data)

    allocate (prcdef_6_t :: core_def)
    select type (core_def)
    type is (prcdef_6_t)
       call core_def%init ()
    end select

    allocate (entry)
    call entry%init (var_str ("proclibs6_a"), &
         model_name = var_str ("Test_model"), &
         n_in = 1, n_components = 1)
    call entry%import_component (1, n_out = 2, &
         prt_in  = new_prt_spec ([var_str ("a")]), &
         prt_out = new_prt_spec ([var_str ("b"), var_str ("c")]), &
         method  = var_str ("test"), &
         variant = core_def)
    call lib%append (entry)

    write (u, "(A)")  "* Configure library"
    write (u, "(A)")
    call lib%configure (os_data)

    write (u, "(A)")  "* Write makefile"
    write (u, "(A)")
    call lib%write_makefile (os_data, force = .true.)
    
    write (u, "(A)")  "* Write driver source code"
    write (u, "(A)")
    call lib%write_driver (force = .true.)

    write (u, "(A)")  "* Write process source code, compile, link, load"
    write (u, "(A)")
    call lib%load (os_data)

    call lib%write (u, libpath = .false.)
    
    write (u, "(A)")
    write (u, "(A)")  "* Probe library API:"
    write (u, "(A)")
       
    write (u, "(1x,A,A,A)")  "name                      = '", &
         char (lib%get_name ()), "'"
    write (u, "(1x,A,L1)")  "is active                 = ", &
         lib%is_active ()
    write (u, "(1x,A,I0)")  "n_processes               = ", &
         lib%get_n_processes ()
    write (u, "(1x,A)", advance="no")  "processes                 ="
    call lib%get_process_id_list (name_list)
    do i = 1, size (name_list)
       write (u, "(1x,A)", advance="no")  char (name_list(i))
    end do
    write (u, *)
    write (u, "(1x,A,L1)")  "proclibs6_a is process    = ", &
         lib%contains (var_str ("proclibs6_a"))
    write (u, "(1x,A,I0)")  "proclibs6_a has index     = ", &
         lib%get_entry_index (var_str ("proclibs6_a"))
    write (u, "(1x,A,L1)")  "foobar is process         = ", &
         lib%contains (var_str ("foobar"))
    write (u, "(1x,A,I0)")  "foobar has index          = ", &
         lib%get_entry_index (var_str ("foobar"))
    write (u, "(1x,A,I0)")  "n_in(proclibs6_a)         = ", &
         lib%get_n_in (var_str ("proclibs6_a"))
    write (u, "(1x,A,A)")   "model_name(proclibs6_a)   = ", &
         char (lib%get_model_name (var_str ("proclibs6_a")))
    write (u, "(1x,A,I0)")  "n_components(proclibs6_a) = ", &
         lib%get_n_components (var_str ("proclibs6_a"))
    write (u, "(1x,A)", advance="no")  "components(proclibs6_a)   ="
    call lib%get_component_list (var_str ("proclibs6_a"), name_list)
    do i = 1, size (name_list)
       write (u, "(1x,A)", advance="no")  char (name_list(i))
    end do
    write (u, *)
    
    write (u, "(A)")
    write (u, "(A)")  "* Constants of proclibs6_a_i1:"
    write (u, "(A)")

    call lib%connect_process (var_str ("proclibs6_a"), 1, data, proc_driver)

    write (u, "(1x,A,A)")  "component ID     = ", char (data%id)
    write (u, "(1x,A,A)")  "model name       = ", char (data%model_name)
    write (u, "(1x,A,A,A)")  "md5sum           = '", data%md5sum, "'"
    write (u, "(1x,A,L1)") "openmp supported = ", data%openmp_supported
    write (u, "(1x,A,I0)") "n_in  = ", data%n_in
    write (u, "(1x,A,I0)") "n_out = ", data%n_out
    write (u, "(1x,A,I0)") "n_flv = ", data%n_flv
    write (u, "(1x,A,I0)") "n_hel = ", data%n_hel
    write (u, "(1x,A,I0)") "n_col = ", data%n_col
    write (u, "(1x,A,I0)") "n_cin = ", data%n_cin
    write (u, "(1x,A,I0)") "n_cf  = ", data%n_cf
    write (u, "(1x,A,10(1x,I0))") "flv state =", data%flv_state
    write (u, "(1x,A,10(1x,I0))") "hel state =", data%hel_state
    write (u, "(1x,A,10(1x,I0))") "col state =", data%col_state
    write (u, "(1x,A,10(1x,L1))") "ghost flag =", data%ghost_flag
    write (u, "(1x,A,10(1x,F5.3))") "color factors =", data%color_factors
    write (u, "(1x,A,10(1x,I0))") "cf index =", data%cf_index
       
    write (u, "(A)")
    write (u, "(A)")  "* Call feature of proclibs6_a:"
    write (u, "(A)")

    select type (proc_driver)
    type is (prctest_6_t)
       call proc_driver%proc1 (n)
       write (u, "(1x,A,I0)") "proc1 = ", n
    end select

    call lib%final ()

    write (u, "(A)")
    write (u, "(A)")  "* Test output end: process_libraries_6"
  end subroutine process_libraries_6
  
\eatline
\nwendcode{}\nwbegindocs{760}\nwdocspar
\subsubsection{MD5 sums}
Check MD5 sum calculation.
\nwenddocs{}\nwbegincode{761}\moddef{Process libraries: execute tests}\plusendmoddef
  call test (process_libraries_7, "process_libraries_7", &
       "process definition list", &
       u, results)
\nwendcode{}\nwbegincode{762}\moddef{Process libraries: tests}\plusendmoddef
  subroutine process_libraries_7 (u)
    integer, intent(in) :: u
    type(prc_template_t), dimension(:), allocatable :: process_core_templates
    type(process_def_entry_t) :: entry
    class(prc_core_def_t), allocatable :: test_def

    write (u, "(A)")  "* Test output: process_libraries_7"
    write (u, "(A)")  "* Purpose: Construct a process definition list &
         &and check MD5 sums"
    write (u, "(A)")
    write (u, "(A)")  "* Construct a process definition list"
    write (u, "(A)")  "*   Process: two components"
    write (u, "(A)")

    allocate (process_core_templates (1))
    allocate (prcdef_2_t :: process_core_templates(1)%core_def)
    
    call entry%init (var_str ("first"), model_name = var_str ("Test"), &
         n_in = 1, n_components = 2)
    allocate (prcdef_2_t :: test_def)
    select type (test_def)
    type is (prcdef_2_t);  test_def%data = 31
    end select
    call entry%import_component (1, n_out = 3, &
         prt_in  = new_prt_spec ([var_str ("a")]), &
         prt_out = new_prt_spec ([var_str ("b"), var_str ("c"), &
                                  var_str ("e")]), &
         method  = var_str ("test"), &
         variant = test_def)
    allocate (prcdef_2_t :: test_def)
    select type (test_def)
    type is (prcdef_2_t);  test_def%data = 42
    end select
    call entry%import_component (2, n_out = 2, &
         prt_in  = new_prt_spec ([var_str ("a")]), &
         prt_out = new_prt_spec ([var_str ("b"), var_str ("c")]), &
         method  = var_str ("test"), &
         variant = test_def)
    call entry%write (u)
    
    write (u, "(A)")
    write (u, "(A)")  "* Compute MD5 sums"
    write (u, "(A)")

    call entry%compute_md5sum ()
    call entry%write (u)

    write (u, "(A)")
    write (u, "(A)")  "* Recalculate MD5 sums (should be identical)"
    write (u, "(A)")

    call entry%compute_md5sum ()
    call entry%write (u)

    write (u, "(A)")
    write (u, "(A)")  "* Modify a component and recalculate MD5 sums"
    write (u, "(A)")

    select type (test_def => entry%initial(2)%core_def)
    type is (prcdef_2_t)
       test_def%data = 54
    end select
    call entry%compute_md5sum ()
    call entry%write (u)

    write (u, "(A)")
    write (u, "(A)")  "* Modify the model and recalculate MD5 sums"
    write (u, "(A)")

    entry%model_name = "foo"
    call entry%compute_md5sum ()
    call entry%write (u)

    write (u, "(A)")
    write (u, "(A)")  "* Test output end: process_libraries_7"
  end subroutine process_libraries_7
  
\eatline
\nwendcode{}\nwbegindocs{763}\nwdocspar
Here is the actual test:
\nwenddocs{}\nwbegincode{764}\moddef{Process libraries: execute tests}\plusendmoddef
  call test (process_libraries_8, "process_libraries_8", &
       "library status checks", &
       u, results)
\nwendcode{}\nwbegincode{765}\moddef{Process libraries: tests}\plusendmoddef
  subroutine process_libraries_8 (u)
    integer, intent(in) :: u
    type(process_library_t) :: lib
    type(process_def_entry_t), pointer :: entry
    class(prc_core_def_t), allocatable :: core_def
    type(os_data_t) :: os_data

    write (u, "(A)")  "* Test output: process_libraries_8"
    write (u, "(A)")  "* Purpose: build and load a process library"
    write (u, "(A)")  "*          with an external (pseudo) matrix element"
    write (u, "(A)")  "*          Check status updates"
    write (u, "(A)")

    write (u, "(A)")  "* Initialize a process library with one entry"
    write (u, "(A)")
    call lib%init (var_str ("proclibs8"))
    call os_data_init (os_data)

    allocate (prcdef_6_t :: core_def)
    select type (core_def)
    type is (prcdef_6_t)
       call core_def%init ()
    end select

    allocate (entry)
    call entry%init (var_str ("proclibs8_a"), &
         model_name = var_str ("Test_model"), &
         n_in = 1, n_components = 1)
    call entry%import_component (1, n_out = 2, &
         prt_in  = new_prt_spec ([var_str ("a")]), &
         prt_out = new_prt_spec ([var_str ("b"), var_str ("c")]), &
         method  = var_str ("test"), &
         variant = core_def)
    call lib%append (entry)

    write (u, "(A)")  "* Configure library"
    write (u, "(A)")

    call lib%configure (os_data)
    call lib%compute_md5sum ()
    
    associate (def => lib%entry(1)%def%initial(1))
      select type (writer => lib%driver%record(1)%writer)
      type is (test_writer_4_t)
         writer%md5sum = def%md5sum
      end select
    end associate

    write (u, "(1x,A,L1)")  "library loaded = ", lib%driver%loaded
    write (u, "(1x,A,I0)")  "lib status   = ", lib%status
    write (u, "(1x,A,I0)")  "proc1 status = ", lib%entry(1)%status

    write (u, "(A)")
    write (u, "(A)")  "* Write makefile"
    write (u, "(A)")
    call lib%write_makefile (os_data, force = .true.)
    
    write (u, "(A)")  "* Update status"
    write (u, "(A)")
    
    call lib%update_status (os_data)
    write (u, "(1x,A,L1)")  "library loaded = ", lib%driver%loaded
    write (u, "(1x,A,I0)")  "lib status   = ", lib%status
    write (u, "(1x,A,I0)")  "proc1 status = ", lib%entry(1)%status
    
    write (u, "(A)")
    write (u, "(A)")  "* Write driver source code"
    write (u, "(A)")
    call lib%write_driver (force = .false.)

    write (u, "(A)")  "* Write process source code"
    write (u, "(A)")
    call lib%make_source (os_data)

    write (u, "(1x,A,L1)")  "library loaded = ", lib%driver%loaded
    write (u, "(1x,A,I0)")  "lib status   = ", lib%status
    write (u, "(1x,A,I0)")  "proc1 status = ", lib%entry(1)%status

    write (u, "(A)")
    write (u, "(A)")  "* Compile and load"
    write (u, "(A)")

    call lib%load (os_data)
    write (u, "(1x,A,L1)")  "library loaded = ", lib%driver%loaded
    write (u, "(1x,A,I0)")  "lib status   = ", lib%status
    write (u, "(1x,A,I0)")  "proc1 status = ", lib%entry(1)%status

    write (u, "(A)")
    write (u, "(A)")  "* Append process and reconfigure"
    write (u, "(A)")
    
    allocate (prcdef_6_t :: core_def)
    select type (core_def)
    type is (prcdef_6_t)
       call core_def%init ()
    end select

    allocate (entry)
    call entry%init (var_str ("proclibs8_b"), &
         model_name = var_str ("Test_model"), &
         n_in = 1, n_components = 1)
    call entry%import_component (1, n_out = 2, &
         prt_in  = new_prt_spec ([var_str ("a")]), &
         prt_out = new_prt_spec ([var_str ("b"), var_str ("d")]), &
         method  = var_str ("test"), &
         variant = core_def)
    call lib%append (entry)

    call lib%configure (os_data)
    call lib%compute_md5sum ()
    associate (def => lib%entry(2)%def%initial(1))
      select type (writer => lib%driver%record(2)%writer)
      type is (test_writer_4_t)
         writer%md5sum = def%md5sum
      end select
    end associate
    call lib%write_makefile (os_data, force = .false.)
    call lib%write_driver (force = .false.)

    write (u, "(1x,A,L1)")  "library loaded = ", lib%driver%loaded
    write (u, "(1x,A,I0)")  "lib status   = ", lib%status
    write (u, "(1x,A,I0)")  "proc1 status = ", lib%entry(1)%status
    write (u, "(1x,A,I0)")  "proc2 status = ", lib%entry(2)%status

    write (u, "(A)")
    write (u, "(A)")  "* Update status"
    write (u, "(A)")
    
    call lib%update_status (os_data)
    write (u, "(1x,A,L1)")  "library loaded = ", lib%driver%loaded
    write (u, "(1x,A,I0)")  "lib status   = ", lib%status
    write (u, "(1x,A,I0)")  "proc1 status = ", lib%entry(1)%status
    write (u, "(1x,A,I0)")  "proc2 status = ", lib%entry(2)%status
    
    write (u, "(A)")
    write (u, "(A)")  "* Write source code"
    write (u, "(A)")
    
    call lib%make_source (os_data)
    write (u, "(1x,A,L1)")  "library loaded = ", lib%driver%loaded
    write (u, "(1x,A,I0)")  "lib status   = ", lib%status
    write (u, "(1x,A,I0)")  "proc1 status = ", lib%entry(1)%status
    write (u, "(1x,A,I0)")  "proc2 status = ", lib%entry(2)%status
    
    write (u, "(A)")
    write (u, "(A)")  "* Reset status"
    write (u, "(A)")
    
    lib%status = STAT_CONFIGURED
    lib%entry%status = STAT_CONFIGURED
    write (u, "(1x,A,L1)")  "library loaded = ", lib%driver%loaded
    write (u, "(1x,A,I0)")  "lib status   = ", lib%status
    write (u, "(1x,A,I0)")  "proc1 status = ", lib%entry(1)%status
    write (u, "(1x,A,I0)")  "proc2 status = ", lib%entry(2)%status
    
    write (u, "(A)")
    write (u, "(A)")  "* Update status"
    write (u, "(A)")
    
    call lib%update_status (os_data)
    write (u, "(1x,A,L1)")  "library loaded = ", lib%driver%loaded
    write (u, "(1x,A,I0)")  "lib status   = ", lib%status
    write (u, "(1x,A,I0)")  "proc1 status = ", lib%entry(1)%status
    write (u, "(1x,A,I0)")  "proc2 status = ", lib%entry(2)%status
    
    write (u, "(A)")
    write (u, "(A)")  "* Partial cleanup"
    write (u, "(A)")

    call lib%clean (os_data, distclean = .false.)
    write (u, "(1x,A,L1)")  "library loaded = ", lib%driver%loaded
    write (u, "(1x,A,I0)")  "lib status   = ", lib%status
    write (u, "(1x,A,I0)")  "proc1 status = ", lib%entry(1)%status
    write (u, "(1x,A,I0)")  "proc2 status = ", lib%entry(2)%status
    

    write (u, "(A)")
    write (u, "(A)")  "* Update status"
    write (u, "(A)")
    
    call lib%update_status (os_data)
    write (u, "(1x,A,L1)")  "library loaded = ", lib%driver%loaded
    write (u, "(1x,A,I0)")  "lib status   = ", lib%status
    write (u, "(1x,A,I0)")  "proc1 status = ", lib%entry(1)%status
    write (u, "(1x,A,I0)")  "proc2 status = ", lib%entry(2)%status
    
    write (u, "(A)")
    write (u, "(A)")  "* Complete cleanup"

    call lib%clean (os_data, distclean = .true.)
    call lib%final ()

    write (u, "(A)")
    write (u, "(A)")  "* Test output end: process_libraries_8"
  end subroutine process_libraries_8
  
\eatline
\nwendcode{}\nwbegindocs{766}\nwdocspar
%%%%%%%%%%%%%%%%%%%%%%%%%%%%%%%%%%%%%%%%%%%%%%%%%%%%%%%%%%%%%%%%%%%%%%%%
\section{Process Library Stacks}

For storing and handling multiple libraries, we define process library stacks.
These are ordinary stacks where new entries are pushed onto the top.
\nwenddocs{}\nwbegincode{767}\moddef{\code{}prclib{\char95}stacks.f90\edoc{}}\endmoddef
\LA{}File header\RA{}

module prclib_stacks
  
\LA{}Use strings\RA{}
  use io_units
  use format_utils, only: write_separator
  use unit_tests
  use os_interface
  use process_libraries

\LA{}Standard module head\RA{}

\LA{}Prclib stacks: public\RA{}

\LA{}Prclib stacks: types\RA{}

contains
  
\LA{}Prclib stacks: procedures\RA{}

\LA{}Prclib stacks: tests\RA{}

end module prclib_stacks
\eatline
\nwendcode{}\nwbegindocs{768}\nwdocspar
\subsection{The stack entry type}
A stack entry is a process library object, augmented by a pointer to the
next entry.  We do not need specific methods, all relevant methods are
inherited.

On higher level, process libraries should be prepared as process entry objects.
\nwenddocs{}\nwbegincode{769}\moddef{Prclib stacks: public}\endmoddef
  public :: prclib_entry_t
\nwendcode{}\nwbegincode{770}\moddef{Prclib stacks: types}\endmoddef
  type, extends (process_library_t) :: prclib_entry_t
     type(prclib_entry_t), pointer :: next => null ()
  end type prclib_entry_t

\eatline
\nwendcode{}\nwbegindocs{771}\nwdocspar
\subsection{The prclib stack type}
For easy conversion and lookup it is useful to store the filling
number in the object.  The content is stored as a linked list.
\nwenddocs{}\nwbegincode{772}\moddef{Prclib stacks: public}\plusendmoddef
  public :: prclib_stack_t
\nwendcode{}\nwbegincode{773}\moddef{Prclib stacks: types}\plusendmoddef
  type :: prclib_stack_t
     integer :: n = 0
     type(prclib_entry_t), pointer :: first => null ()
   contains
   \LA{}Prclib stacks: prclib stack: TBP\RA{}
  end type prclib_stack_t
  
\eatline
\nwendcode{}\nwbegindocs{774}Finalizer.  Iteratively deallocate the stack entries.  The resulting
empty stack can be immediately recycled, if necessary.
\nwenddocs{}\nwbegincode{775}\moddef{Prclib stacks: prclib stack: TBP}\endmoddef
  procedure :: final => prclib_stack_final
\nwendcode{}\nwbegincode{776}\moddef{Prclib stacks: procedures}\endmoddef
  subroutine prclib_stack_final (object)
    class(prclib_stack_t), intent(inout) :: object
    type(prclib_entry_t), pointer :: lib
    do while (associated (object%first))
       lib => object%first
       object%first => lib%next
       call lib%final ()
       deallocate (lib)
    end do
    object%n = 0
  end subroutine prclib_stack_final
  
\eatline
\nwendcode{}\nwbegindocs{777}Output.  The entries on the stack will be ordered LIFO, i.e., backwards.
\nwenddocs{}\nwbegincode{778}\moddef{Prclib stacks: prclib stack: TBP}\plusendmoddef
  procedure :: write => prclib_stack_write
\nwendcode{}\nwbegincode{779}\moddef{Prclib stacks: procedures}\plusendmoddef
  subroutine prclib_stack_write (object, unit, libpath)
    class(prclib_stack_t), intent(in) :: object
    integer, intent(in), optional :: unit
    logical, intent(in), optional :: libpath
    type(prclib_entry_t), pointer :: lib
    integer :: u
    u = given_output_unit (unit)
    call write_separator (u, 2)
    select case (object%n)
    case (0)
       write (u, "(1x,A)")  "Process library stack: [empty]"
    case default
       write (u, "(1x,A)")  "Process library stack:"
       lib => object%first
       do while (associated (lib))
          call write_separator (u)
          call lib%write (u, libpath)
          lib => lib%next
       end do
    end select
    call write_separator (u, 2)
  end subroutine prclib_stack_write

\eatline
\nwendcode{}\nwbegindocs{780}\nwdocspar
\subsection{Operating on Stacks}
We take a library entry pointer and push it onto the stack.  The previous
pointer is nullified.  Subsequently, the library entry is `owned' by the
stack and will be finalized when the stack is deleted.
\nwenddocs{}\nwbegincode{781}\moddef{Prclib stacks: prclib stack: TBP}\plusendmoddef
  procedure :: push => prclib_stack_push
\nwendcode{}\nwbegincode{782}\moddef{Prclib stacks: procedures}\plusendmoddef
  subroutine prclib_stack_push (stack, lib)
    class(prclib_stack_t), intent(inout) :: stack
    type(prclib_entry_t), intent(inout), pointer :: lib
    lib%next => stack%first
    stack%first => lib
    lib => null ()
    stack%n = stack%n + 1
  end subroutine prclib_stack_push
  
\eatline
\nwendcode{}\nwbegindocs{783}\nwdocspar
\subsection{Accessing Contents}
Return a pointer to the topmost stack element.  The result type is
just the bare {\tt{}process{\char95}library{\char95}t}.  There is no {\tt{}target} attribute
required since the stack elements are allocated via pointers.
\nwenddocs{}\nwbegincode{784}\moddef{Prclib stacks: prclib stack: TBP}\plusendmoddef
  procedure :: get_first_ptr => prclib_stack_get_first_ptr
\nwendcode{}\nwbegincode{785}\moddef{Prclib stacks: procedures}\plusendmoddef
  function prclib_stack_get_first_ptr (stack) result (ptr)
    class(prclib_stack_t), intent(in) :: stack
    type(process_library_t), pointer :: ptr
    if (associated (stack%first)) then
       ptr => stack%first%process_library_t
    else
       ptr => null ()
    end if
  end function prclib_stack_get_first_ptr

\eatline
\nwendcode{}\nwbegindocs{786}Return a complete list of the libraries (names) in the stack.  The list is
in the order in which the elements got pushed onto the stack, so the 'first'
entry is listed last.
\nwenddocs{}\nwbegincode{787}\moddef{Prclib stacks: prclib stack: TBP}\plusendmoddef
  procedure :: get_names => prclib_stack_get_names
\nwendcode{}\nwbegincode{788}\moddef{Prclib stacks: procedures}\plusendmoddef
  subroutine prclib_stack_get_names (stack, libname)
    class(prclib_stack_t), intent(in) :: stack
    type(string_t), dimension(:), allocatable, intent(out) :: libname
    type(prclib_entry_t), pointer :: lib
    integer :: i
    allocate (libname (stack%n))
    i = stack%n
    lib => stack%first
    do while (associated (lib))
       libname(i) = lib%get_name ()
       i = i - 1
       lib => lib%next
    end do
  end subroutine prclib_stack_get_names

\eatline
\nwendcode{}\nwbegindocs{789}Return a pointer to the library with given name.
\nwenddocs{}\nwbegincode{790}\moddef{Prclib stacks: prclib stack: TBP}\plusendmoddef
  procedure :: get_library_ptr => prclib_stack_get_library_ptr
\nwendcode{}\nwbegincode{791}\moddef{Prclib stacks: procedures}\plusendmoddef
  function prclib_stack_get_library_ptr (stack, libname) result (ptr)
    class(prclib_stack_t), intent(in) :: stack
    type(string_t), intent(in) :: libname
    type(process_library_t), pointer :: ptr
    type(prclib_entry_t), pointer :: current
    current => stack%first
    do while (associated (current))
       if (current%get_name () == libname) then
          ptr => current%process_library_t
          return
       end if
       current => current%next
    end do
    ptr => null ()
  end function prclib_stack_get_library_ptr

\eatline
\nwendcode{}\nwbegindocs{792}\nwdocspar
\subsection{Unit tests}
\nwenddocs{}\nwbegincode{793}\moddef{Prclib stacks: public}\plusendmoddef
  public :: prclib_stacks_test
\nwendcode{}\nwbegincode{794}\moddef{Prclib stacks: tests}\endmoddef
  subroutine prclib_stacks_test (u, results)
    integer, intent(in) :: u
    type(test_results_t), intent(inout) :: results
  \LA{}Prclib stacks: execute tests\RA{}
  end subroutine prclib_stacks_test
  
\eatline
\nwendcode{}\nwbegindocs{795}\nwdocspar
\subsubsection{Write an empty process library stack}
The most trivial test is to write an uninitialized process library stack.
\nwenddocs{}\nwbegincode{796}\moddef{Prclib stacks: execute tests}\endmoddef
  call test (prclib_stacks_1, "prclib_stacks_1", &
       "write an empty process library stack", &
       u, results)
\nwendcode{}\nwbegincode{797}\moddef{Prclib stacks: tests}\plusendmoddef
  subroutine prclib_stacks_1 (u)
    integer, intent(in) :: u
    type(prclib_stack_t) :: stack

    write (u, "(A)")  "* Test output: prclib_stacks_1"
    write (u, "(A)")  "*   Purpose: display an empty process library stack"
    write (u, "(A)")

    call stack%write (u)
    
    write (u, "(A)")
    write (u, "(A)")  "* Test output end: prclib_stacks_1"
    
  end subroutine prclib_stacks_1
  
\eatline
\nwendcode{}\nwbegindocs{798}\nwdocspar
\subsubsection{Fill a process library stack}
Fill a process library stack with two (identical) processes.
\nwenddocs{}\nwbegincode{799}\moddef{Prclib stacks: execute tests}\plusendmoddef
  call test (prclib_stacks_2, "prclib_stacks_2", &
       "fill a process library stack", &
       u, results)
\nwendcode{}\nwbegincode{800}\moddef{Prclib stacks: tests}\plusendmoddef
  subroutine prclib_stacks_2 (u)
    integer, intent(in) :: u
    type(prclib_stack_t) :: stack
    type(prclib_entry_t), pointer :: lib

    write (u, "(A)")  "* Test output: prclib_stacks_2"
    write (u, "(A)")  "*   Purpose: fill a process library stack"
    write (u, "(A)")

    write (u, "(A)")  "* Initialize two (empty) libraries &
         &and push them on the stack"
    write (u, "(A)")

    allocate (lib)
    call lib%init (var_str ("lib1"))
    call stack%push (lib)

    allocate (lib)
    call lib%init (var_str ("lib2"))
    call stack%push (lib)

    call stack%write (u)

    write (u, "(A)")
    write (u, "(A)")  "* Cleanup"

    call stack%final ()
    
    write (u, "(A)")
    write (u, "(A)")  "* Test output end: prclib_stacks_2"
    
  end subroutine prclib_stacks_2
  
\eatline
\nwendcode{}\nwbegindocs{801}\nwdocspar
%%%%%%%%%%%%%%%%%%%%%%%%%%%%%%%%%%%%%%%%%%%%%%%%%%%%%%%%%%%%%%%%%%%%%%%%
\section{Trivial matrix element for tests}

For the purpose of testing the workflow, we implement here two matrix
elements with the simplest possible structure.

This matrix element generator can only generate a single scattering
process and a single decay process.  The scattering process is a
quartic interaction of a massless, neutral and colorless scalar {\tt{}s}
with unit coupling results in a trivial $2\to 2$ scattering process.
The matrix element is implemented internally, so we do not need the
machinery of external process libraries.  The decay process is a decay
of {\tt{}s} into a pair of colored fermions {\tt{}f}.

\nwenddocs{}\nwbegincode{802}\moddef{\code{}prc{\char95}test.f90\edoc{}}\endmoddef
\LA{}File header\RA{}

module prc_test
  
  use iso_c_binding !NODEP!
  use kinds
\LA{}Use strings\RA{}
  use unit_tests
  use os_interface

  use process_constants
  use prclib_interfaces
  use prc_core_def
  use particle_specifiers
  use process_libraries

\LA{}Standard module head\RA{}

\LA{}Test ME: public\RA{}

\LA{}Test ME: types\RA{}

contains
  
\LA{}Test ME: procedures\RA{}

\LA{}Test ME: tests\RA{}

end module prc_test
\eatline
\nwendcode{}\nwbegindocs{803}\nwdocspar
\subsection{Process definition}
For the process definition we implement an extension of the
{\tt{}prc{\char95}core{\char95}def{\char95}t} abstract type.
\nwenddocs{}\nwbegincode{804}\moddef{Test ME: public}\endmoddef
  public :: prc_test_def_t
\nwendcode{}\nwbegincode{805}\moddef{Test ME: types}\endmoddef
  type, extends (prc_core_def_t) :: prc_test_def_t
     type(string_t) :: model_name
     type(string_t), dimension(:), allocatable :: prt_in
     type(string_t), dimension(:), allocatable :: prt_out
   contains
   \LA{}Test ME: test me def: TBP\RA{}
  end type prc_test_def_t
  
\eatline
\nwendcode{}\nwbegincode{806}\moddef{Test ME: test me def: TBP}\endmoddef
  procedure, nopass :: type_string => prc_test_def_type_string
\nwendcode{}\nwbegincode{807}\moddef{Test ME: procedures}\endmoddef
  function prc_test_def_type_string () result (string)
    type(string_t) :: string
    string = "test_me"
  end function prc_test_def_type_string

\eatline
\nwendcode{}\nwbegindocs{808}There is no 'feature' here since there is no external code.
\nwenddocs{}\nwbegincode{809}\moddef{Test ME: test me def: TBP}\plusendmoddef
  procedure, nopass :: get_features => prc_test_def_get_features
\nwendcode{}\nwbegincode{810}\moddef{Test ME: procedures}\plusendmoddef
  subroutine prc_test_def_get_features (features)
    type(string_t), dimension(:), allocatable, intent(out) :: features
    allocate (features (0))
  end subroutine prc_test_def_get_features

\eatline
\nwendcode{}\nwbegindocs{811}Initialization: set some data (not really useful).
\nwenddocs{}\nwbegincode{812}\moddef{Test ME: test me def: TBP}\plusendmoddef
  procedure :: init => prc_test_def_init
\nwendcode{}\nwbegincode{813}\moddef{Test ME: procedures}\plusendmoddef
  subroutine prc_test_def_init (object, model_name, prt_in, prt_out)
    class(prc_test_def_t), intent(out) :: object
    type(string_t), intent(in) :: model_name
    type(string_t), dimension(:), intent(in) :: prt_in
    type(string_t), dimension(:), intent(in) :: prt_out
    object%model_name = model_name
    allocate (object%prt_in (size (prt_in)))
    object%prt_in = prt_in
    allocate (object%prt_out (size (prt_out)))
    object%prt_out = prt_out
  end subroutine prc_test_def_init

\eatline
\nwendcode{}\nwbegindocs{814}Write/read process- and method-specific data.  (No-op)
\nwenddocs{}\nwbegincode{815}\moddef{Test ME: test me def: TBP}\plusendmoddef
  procedure :: write => prc_test_def_write
\nwendcode{}\nwbegincode{816}\moddef{Test ME: procedures}\plusendmoddef
  subroutine prc_test_def_write (object, unit)
    class(prc_test_def_t), intent(in) :: object
    integer, intent(in) :: unit
  end subroutine prc_test_def_write
  
\eatline
\nwendcode{}\nwbegindocs{817}\nwdocspar
\nwenddocs{}\nwbegincode{818}\moddef{Test ME: test me def: TBP}\plusendmoddef
  procedure :: read => prc_test_def_read
\nwendcode{}\nwbegincode{819}\moddef{Test ME: procedures}\plusendmoddef
  subroutine prc_test_def_read (object, unit)
    class(prc_test_def_t), intent(out) :: object
    integer, intent(in) :: unit
  end subroutine prc_test_def_read
  
\eatline
\nwendcode{}\nwbegindocs{820}Allocate the driver for test ME matrix elements.  We get the
actual component ID (basename), and we can transfer all
process-specific data from the process definition.
\nwenddocs{}\nwbegincode{821}\moddef{Test ME: test me def: TBP}\plusendmoddef
  procedure :: allocate_driver => prc_test_def_allocate_driver
\nwendcode{}\nwbegincode{822}\moddef{Test ME: procedures}\plusendmoddef
  subroutine prc_test_def_allocate_driver (object, driver, basename)
    class(prc_test_def_t), intent(in) :: object
    class(prc_core_driver_t), intent(out), allocatable :: driver
    type(string_t), intent(in) :: basename
    allocate (prc_test_t :: driver)
    select type (driver)
    type is (prc_test_t)
       driver%id = basename
       driver%model_name = object%model_name
       select case (size (object%prt_in))
       case (1);  driver%scattering = .false.
       case (2);  driver%scattering = .true.
       end select
    end select
  end subroutine prc_test_def_allocate_driver
  
\eatline
\nwendcode{}\nwbegindocs{823}Nothing to connect.  This subroutine will not be called.
\nwenddocs{}\nwbegincode{824}\moddef{Test ME: test me def: TBP}\plusendmoddef
  procedure :: connect => prc_test_def_connect
\nwendcode{}\nwbegincode{825}\moddef{Test ME: procedures}\plusendmoddef
  subroutine prc_test_def_connect (def, lib_driver, i, proc_driver)
    class(prc_test_def_t), intent(in) :: def
    class(prclib_driver_t), intent(in) :: lib_driver
    integer, intent(in) :: i
    class(prc_core_driver_t), intent(inout) :: proc_driver
  end subroutine prc_test_def_connect

\eatline
\nwendcode{}\nwbegindocs{826}\nwdocspar
\subsection{Driver}
\nwenddocs{}\nwbegincode{827}\moddef{Test ME: public}\plusendmoddef
  public :: prc_test_t
\nwendcode{}\nwbegincode{828}\moddef{Test ME: types}\plusendmoddef
  type, extends (process_driver_internal_t) :: prc_test_t
     type(string_t) :: id
     type(string_t) :: model_name
     logical :: scattering = .true.
   contains
   \LA{}Test ME: test me driver: TBP\RA{}
  end type prc_test_t

\eatline
\nwendcode{}\nwbegindocs{829}In contrast to generic matrix-element implementations, we can
hard-wire the amplitude method as a type-bound procedure.
\nwenddocs{}\nwbegincode{830}\moddef{Test ME: test me driver: TBP}\endmoddef
  procedure, nopass :: get_amplitude => prc_test_get_amplitude
\nwendcode{}\nwbegincode{831}\moddef{Test ME: procedures}\plusendmoddef
  function prc_test_get_amplitude (p) result (amp)
    complex(default) :: amp
    real(default), dimension(:,:), intent(in) :: p
    amp = 1
  end function prc_test_get_amplitude

\eatline
\nwendcode{}\nwbegindocs{832}The reported type is the same as for the {\tt{}prc{\char95}test{\char95}def{\char95}t} type.
\nwenddocs{}\nwbegincode{833}\moddef{Test ME: test me driver: TBP}\plusendmoddef
  procedure, nopass :: type_name => prc_test_type_name
\nwendcode{}\nwbegincode{834}\moddef{Test ME: procedures}\plusendmoddef
  function prc_test_type_name () result (string)
    type(string_t) :: string
    string = "test_me"
  end function prc_test_type_name

\eatline
\nwendcode{}\nwbegindocs{835}Fill process constants.
\nwenddocs{}\nwbegincode{836}\moddef{Test ME: test me driver: TBP}\plusendmoddef
  procedure :: fill_constants => prc_test_fill_constants
\nwendcode{}\nwbegincode{837}\moddef{Test ME: procedures}\plusendmoddef
  subroutine prc_test_fill_constants (driver, data)
    class(prc_test_t), intent(in) :: driver
    type(process_constants_t), intent(out) :: data
    data%id = driver%id
    data%model_name = driver%model_name
    if (driver%scattering) then
       data%n_in  = 2
       data%n_out = 2
       data%n_flv = 1
       data%n_hel = 1
       data%n_col = 1
       data%n_cin = 2
       data%n_cf  = 1
       allocate (data%flv_state (4, 1))
       data%flv_state = 25
       allocate (data%hel_state (4, 1))
       data%hel_state = 0
       allocate (data%col_state (2, 4, 1))
       data%col_state = 0
       allocate (data%ghost_flag (4, 1))
       data%ghost_flag = .false.
       allocate (data%color_factors (1))
       data%color_factors = 1
       allocate (data%cf_index (2, 1))
       data%cf_index = 1
    else
       data%n_in  = 1
       data%n_out = 2
       data%n_flv = 1
       data%n_hel = 2
       data%n_col = 1
       data%n_cin = 2
       data%n_cf  = 1
       allocate (data%flv_state (3, 1))
       data%flv_state(:,1) = [25, 6, -6]
       allocate (data%hel_state (3, 2))
       data%hel_state(:,1) = [0, 1,-1]
       data%hel_state(:,2) = [0,-1, 1]
       allocate (data%col_state (2, 3, 1))
       data%col_state = reshape ([0,0, 1,0, 0,-1], [2,3,1])
       allocate (data%ghost_flag (3, 1))
       data%ghost_flag = .false.
       allocate (data%color_factors (1))
       data%color_factors = 3
       allocate (data%cf_index (2, 1))
       data%cf_index = 1
    end if
  end subroutine prc_test_fill_constants
  
\eatline
\nwendcode{}\nwbegindocs{838}\nwdocspar
\subsection{Shortcut}
Since this module is there for testing purposes, we set up a
subroutine that does all the work at once: create a library with the
two processes (scattering and decay), configure and load, and set up
the driver.
\nwenddocs{}\nwbegincode{839}\moddef{Test ME: public}\plusendmoddef
  public :: prc_test_create_library
\nwendcode{}\nwbegincode{840}\moddef{Test ME: procedures}\plusendmoddef
  subroutine prc_test_create_library &
       (libname, lib, scattering, decay, procname1, procname2)
    type(string_t), intent(in) :: libname
    type(process_library_t), intent(out) :: lib
    logical, intent(in), optional :: scattering, decay
    type(string_t), intent(in), optional :: procname1, procname2
    type(string_t) :: model_name, procname
    type(string_t), dimension(:), allocatable :: prt_in, prt_out
    class(prc_core_def_t), allocatable :: def
    type(process_def_entry_t), pointer :: entry
    type(os_data_t) :: os_data
    logical :: sca, dec
    sca = .true.;   if (present (scattering))  sca = scattering
    dec = .false.;  if (present (decay))       dec = decay

    call os_data_init (os_data)
    call lib%init (libname)
    model_name = "Test"

    if (sca) then
       if (present (procname1)) then
          procname = procname1
       else
          procname = libname
       end if
       allocate (prt_in (2), prt_out (2))
       prt_in  = [var_str ("s"), var_str ("s")]
       prt_out = [var_str ("s"), var_str ("s")]
       allocate (prc_test_def_t :: def)
       select type (def)
       type is (prc_test_def_t)
          call def%init (model_name, prt_in, prt_out)
       end select
       allocate (entry)
       call entry%init (procname, model_name = model_name, &
            n_in = 2, n_components = 1)
       call entry%import_component (1, n_out = size (prt_out), &
            prt_in  = new_prt_spec (prt_in), &
            prt_out = new_prt_spec (prt_out), &
            method  = var_str ("test_me"), &
            variant = def)
       call lib%append (entry)
    end if

    if (dec) then
       if (present (procname2)) then
          procname = procname2
       else
          procname = libname
       end if
       if (allocated (prt_in))  deallocate (prt_in, prt_out)
       allocate (prt_in (1), prt_out (2))
       prt_in  = [var_str ("s")]
       prt_out = [var_str ("f"), var_str ("fbar")]
       allocate (prc_test_def_t :: def)
       select type (def)
       type is (prc_test_def_t)
          call def%init (model_name, prt_in, prt_out)
       end select
       allocate (entry)
       call entry%init (procname, model_name = model_name, &
            n_in = 1, n_components = 1)
       call entry%import_component (1, n_out = size (prt_out), &
            prt_in  = new_prt_spec (prt_in), &
            prt_out = new_prt_spec (prt_out), &
            method  = var_str ("test_decay"), &
            variant = def)
       call lib%append (entry)
    end if
    
    call lib%configure (os_data)
    call lib%load (os_data)
  end subroutine prc_test_create_library
  
\eatline
\nwendcode{}\nwbegindocs{841}\nwdocspar
\subsection{Test}
This is the master for calling self-test procedures.
\nwenddocs{}\nwbegincode{842}\moddef{Test ME: public}\plusendmoddef
  public :: prc_test_test
\nwendcode{}\nwbegincode{843}\moddef{Test ME: tests}\endmoddef
  subroutine prc_test_test (u, results)
    integer, intent(in) :: u
    type(test_results_t), intent(inout) :: results
  \LA{}Test ME: execute tests\RA{}
end subroutine prc_test_test

\eatline
\nwendcode{}\nwbegindocs{844}\nwdocspar
\subsubsection{Generate and load the scattering process}
The process is $s s \to s s$, where $s$ is a trivial scalar particle,
for vanishing mass and unit coupling.  We initialize the process,
build the library, and compute the particular matrix element for
momenta of unit energy and right-angle scattering.  (The scattering is
independent of angle.)  The matrix element is equal to unity.
\nwenddocs{}\nwbegincode{845}\moddef{Test ME: execute tests}\endmoddef
  call test (prc_test_1, "prc_test_1", &
       "build and load trivial process", &
       u, results)
\nwendcode{}\nwbegincode{846}\moddef{Test ME: tests}\plusendmoddef
  subroutine prc_test_1 (u)
    integer, intent(in) :: u
    type(os_data_t) :: os_data
    type(process_library_t) :: lib
    class(prc_core_def_t), allocatable :: def
    type(process_def_entry_t), pointer :: entry
    type(string_t) :: model_name
    type(string_t), dimension(:), allocatable :: prt_in, prt_out
    type(process_constants_t) :: data
    class(prc_core_driver_t), allocatable :: driver
    real(default), dimension(0:3,4) :: p
    integer :: i
    
    write (u, "(A)")  "* Test output: prc_test_1"
    write (u, "(A)")  "*   Purpose: create a trivial process"
    write (u, "(A)")  "*            build a library and &
         &access the matrix element"
    write (u, "(A)")

    write (u, "(A)")  "* Initialize a process library with one entry"
    write (u, "(A)")
    call os_data_init (os_data)
    call lib%init (var_str ("prc_test1"))

    model_name = "Test"
    allocate (prt_in (2), prt_out (2))
    prt_in  = [var_str ("s"), var_str ("s")]
    prt_out = [var_str ("s"), var_str ("s")]
    
    allocate (prc_test_def_t :: def)
    select type (def)
    type is (prc_test_def_t)
       call def%init (model_name, prt_in, prt_out)
    end select
    allocate (entry)
    call entry%init (var_str ("prc_test1_a"), model_name = model_name, &
         n_in = 2, n_components = 1)
    call entry%import_component (1, n_out = size (prt_out), &
         prt_in  = new_prt_spec (prt_in), &
         prt_out = new_prt_spec (prt_out), &
         method  = var_str ("test_me"), &
         variant = def)
    call lib%append (entry)
    
    write (u, "(A)")  "* Configure library"
    write (u, "(A)")
    call lib%configure (os_data)
    
    write (u, "(A)")  "* Load library"
    write (u, "(A)")
    call lib%load (os_data)

    call lib%write (u)
    
    write (u, "(A)")
    write (u, "(A)")  "* Probe library API:"
    write (u, "(A)")
       
    write (u, "(1x,A,L1)")  "is active                 = ", &
         lib%is_active ()
    write (u, "(1x,A,I0)")  "n_processes               = ", &
         lib%get_n_processes ()

    write (u, "(A)")
    write (u, "(A)")  "* Constants of prc_test1_a_i1:"
    write (u, "(A)")

    call lib%connect_process (var_str ("prc_test1_a"), 1, data, driver)

    write (u, "(1x,A,A)")  "component ID     = ", char (data%id)
    write (u, "(1x,A,A)")  "model name       = ", char (data%model_name)
    write (u, "(1x,A,A,A)")  "md5sum           = '", data%md5sum, "'"
    write (u, "(1x,A,L1)") "openmp supported = ", data%openmp_supported
    write (u, "(1x,A,I0)") "n_in  = ", data%n_in
    write (u, "(1x,A,I0)") "n_out = ", data%n_out
    write (u, "(1x,A,I0)") "n_flv = ", data%n_flv
    write (u, "(1x,A,I0)") "n_hel = ", data%n_hel
    write (u, "(1x,A,I0)") "n_col = ", data%n_col
    write (u, "(1x,A,I0)") "n_cin = ", data%n_cin
    write (u, "(1x,A,I0)") "n_cf  = ", data%n_cf
    write (u, "(1x,A,10(1x,I0))") "flv state =", data%flv_state
    write (u, "(1x,A,10(1x,I2))") "hel state =", data%hel_state(:,1)
    write (u, "(1x,A,10(1x,I0))") "col state =", data%col_state
    write (u, "(1x,A,10(1x,L1))") "ghost flag =", data%ghost_flag
    write (u, "(1x,A,10(1x,F5.3))") "color factors =", data%color_factors
    write (u, "(1x,A,10(1x,I0))") "cf index =", data%cf_index

    write (u, "(A)")
    write (u, "(A)")  "* Set kinematics:"
    write (u, "(A)")
    
    p = reshape ([ &
         1.0_default, 0.0_default, 0.0_default, 1.0_default, &
         1.0_default, 0.0_default, 0.0_default,-1.0_default, &
         1.0_default, 1.0_default, 0.0_default, 0.0_default, &
         1.0_default,-1.0_default, 0.0_default, 0.0_default &
         ], [4,4])
    do i = 1, 4
       write (u, "(2x,A,I0,A,4(1x,F7.4))")  "p", i, " =", p(:,i)
    end do

    write (u, "(A)")
    write (u, "(A)")  "* Compute matrix element:"
    write (u, "(A)")

    select type (driver)
    type is (prc_test_t)
       write (u, "(1x,A,1x,E11.4)") "|amp| =", abs (driver%get_amplitude (p))
    end select

    call lib%final ()
    
    write (u, "(A)")
    write (u, "(A)")  "* Test output end: prc_test_1"
    
  end subroutine prc_test_1
  
\eatline
\nwendcode{}\nwbegindocs{847}\nwdocspar
\subsubsection{Shortcut}
This is identical to the previous test, but we create the library be a single
command.  This is handy for other modules which use the test process.
\nwenddocs{}\nwbegincode{848}\moddef{Test ME: execute tests}\plusendmoddef
  call test (prc_test_2, "prc_test_2", &
       "build and load trivial process using shortcut", &
       u, results)
\nwendcode{}\nwbegincode{849}\moddef{Test ME: tests}\plusendmoddef
  subroutine prc_test_2 (u)
    integer, intent(in) :: u
    type(process_library_t) :: lib
    class(prc_core_driver_t), allocatable :: driver
    type(process_constants_t) :: data
    real(default), dimension(0:3,4) :: p
    
    write (u, "(A)")  "* Test output: prc_test_2"
    write (u, "(A)")  "*   Purpose: create a trivial process"
    write (u, "(A)")  "*            build a library and &
         &access the matrix element"
    write (u, "(A)")

    write (u, "(A)")  "* Build and load a process library with one entry"

    call prc_test_create_library (var_str ("prc_test2"), lib)
    call lib%connect_process (var_str ("prc_test2"), 1, data, driver)

    p = reshape ([ &
         1.0_default, 0.0_default, 0.0_default, 1.0_default, &
         1.0_default, 0.0_default, 0.0_default,-1.0_default, &
         1.0_default, 1.0_default, 0.0_default, 0.0_default, &
         1.0_default,-1.0_default, 0.0_default, 0.0_default &
         ], [4,4])

    write (u, "(A)")
    write (u, "(A)")  "* Compute matrix element:"
    write (u, "(A)")

    select type (driver)
    type is (prc_test_t)
       write (u, "(1x,A,1x,E11.4)") "|amp| =", abs (driver%get_amplitude (p))
    end select

    call lib%final ()
    
    write (u, "(A)")
    write (u, "(A)")  "* Test output end: prc_test_2"
    
  end subroutine prc_test_2
  
\eatline
\nwendcode{}\nwbegindocs{850}\nwdocspar
\subsubsection{Generate and load the decay process}
The process is $s \to f\bar f$, where $s$ is a trivial scalar particle
and $f$ is a colored fermion.  We initialize the process,
build the library, and compute the particular matrix element for a
fixed momentum configuration.  (The decay is
independent of angle.)  The matrix element is equal to unity.
\nwenddocs{}\nwbegincode{851}\moddef{Test ME: execute tests}\plusendmoddef
  call test (prc_test_3, "prc_test_3", &
       "build and load trivial decay", &
       u, results)
\nwendcode{}\nwbegincode{852}\moddef{Test ME: tests}\plusendmoddef
  subroutine prc_test_3 (u)
    integer, intent(in) :: u
    type(os_data_t) :: os_data
    type(process_library_t) :: lib
    class(prc_core_def_t), allocatable :: def
    type(process_def_entry_t), pointer :: entry
    type(string_t) :: model_name
    type(string_t), dimension(:), allocatable :: prt_in, prt_out
    type(process_constants_t) :: data
    class(prc_core_driver_t), allocatable :: driver
    real(default), dimension(0:3,3) :: p
    integer :: i
    
    write (u, "(A)")  "* Test output: prc_test_3"
    write (u, "(A)")  "*   Purpose: create a trivial decay process"
    write (u, "(A)")  "*            build a library and &
         &access the matrix element"
    write (u, "(A)")

    write (u, "(A)")  "* Initialize a process library with one entry"
    write (u, "(A)")
    call os_data_init (os_data)
    call lib%init (var_str ("prc_test3"))

    model_name = "Test"
    allocate (prt_in (2), prt_out (2))
    prt_in  = [var_str ("s")]
    prt_out = [var_str ("f"), var_str ("F")]
    
    allocate (prc_test_def_t :: def)
    select type (def)
    type is (prc_test_def_t)
       call def%init (model_name, prt_in, prt_out)
    end select
    allocate (entry)
    call entry%init (var_str ("prc_test3_a"), model_name = model_name, &
         n_in = 1, n_components = 1)
    call entry%import_component (1, n_out = size (prt_out), &
         prt_in  = new_prt_spec (prt_in), &
         prt_out = new_prt_spec (prt_out), &
         method  = var_str ("test_me"), &
         variant = def)
    call lib%append (entry)
    
    write (u, "(A)")  "* Configure library"
    write (u, "(A)")
    call lib%configure (os_data)
    
    write (u, "(A)")  "* Load library"
    write (u, "(A)")
    call lib%load (os_data)

    call lib%write (u)
    
    write (u, "(A)")
    write (u, "(A)")  "* Probe library API:"
    write (u, "(A)")
       
    write (u, "(1x,A,L1)")  "is active                 = ", &
         lib%is_active ()
    write (u, "(1x,A,I0)")  "n_processes               = ", &
         lib%get_n_processes ()

    write (u, "(A)")
    write (u, "(A)")  "* Constants of prc_test3_a_i1:"
    write (u, "(A)")

    call lib%connect_process (var_str ("prc_test3_a"), 1, data, driver)

    write (u, "(1x,A,A)")  "component ID     = ", char (data%id)
    write (u, "(1x,A,A)")  "model name       = ", char (data%model_name)
    write (u, "(1x,A,A,A)")  "md5sum           = '", data%md5sum, "'"
    write (u, "(1x,A,L1)") "openmp supported = ", data%openmp_supported
    write (u, "(1x,A,I0)") "n_in  = ", data%n_in
    write (u, "(1x,A,I0)") "n_out = ", data%n_out
    write (u, "(1x,A,I0)") "n_flv = ", data%n_flv
    write (u, "(1x,A,I0)") "n_hel = ", data%n_hel
    write (u, "(1x,A,I0)") "n_col = ", data%n_col
    write (u, "(1x,A,I0)") "n_cin = ", data%n_cin
    write (u, "(1x,A,I0)") "n_cf  = ", data%n_cf
    write (u, "(1x,A,10(1x,I0))") "flv state =", data%flv_state
    write (u, "(1x,A,10(1x,I2))") "hel state =", data%hel_state(:,1)
    write (u, "(1x,A,10(1x,I2))") "hel state =", data%hel_state(:,2)
    write (u, "(1x,A,10(1x,I0))") "col state =", data%col_state
    write (u, "(1x,A,10(1x,L1))") "ghost flag =", data%ghost_flag
    write (u, "(1x,A,10(1x,F5.3))") "color factors =", data%color_factors
    write (u, "(1x,A,10(1x,I0))") "cf index =", data%cf_index

    write (u, "(A)")
    write (u, "(A)")  "* Set kinematics:"
    write (u, "(A)")
    
    p = reshape ([ &
         125._default, 0.0_default, 0.0_default, 0.0_default, &
         62.5_default, 0.0_default, 0.0_default, 62.5_default, &
         62.5_default, 0.0_default, 0.0_default,-62.5_default &
         ], [4,3])
    do i = 1, 3
       write (u, "(2x,A,I0,A,4(1x,F8.4))")  "p", i, " =", p(:,i)
    end do

    write (u, "(A)")
    write (u, "(A)")  "* Compute matrix element:"
    write (u, "(A)")

    select type (driver)
    type is (prc_test_t)
       write (u, "(1x,A,1x,E11.4)") "|amp| =", abs (driver%get_amplitude (p))
    end select

    call lib%final ()
    
    write (u, "(A)")
    write (u, "(A)")  "* Test output end: prc_test_3"
    
  end subroutine prc_test_3
  
\eatline
\nwendcode{}\nwbegindocs{853}\nwdocspar
\subsubsection{Shortcut}
This is identical to the previous test, but we create the library be a single
command.  This is handy for other modules which use the test process.
\nwenddocs{}\nwbegincode{854}\moddef{Test ME: execute tests}\plusendmoddef
  call test (prc_test_4, "prc_test_4", &
       "build and load trivial decay using shortcut", &
       u, results)
\nwendcode{}\nwbegincode{855}\moddef{Test ME: tests}\plusendmoddef
  subroutine prc_test_4 (u)
    integer, intent(in) :: u
    type(process_library_t) :: lib
    class(prc_core_driver_t), allocatable :: driver
    type(process_constants_t) :: data
    real(default), dimension(0:3,3) :: p
    
    write (u, "(A)")  "* Test output: prc_test_4"
    write (u, "(A)")  "*   Purpose: create a trivial decay process"
    write (u, "(A)")  "*            build a library and &
         &access the matrix element"
    write (u, "(A)")

    write (u, "(A)")  "* Build and load a process library with one entry"

    call prc_test_create_library (var_str ("prc_test4"), lib, &
         scattering=.false., decay=.true.)
    call lib%connect_process (var_str ("prc_test4"), 1, data, driver)

    p = reshape ([ &
         125._default, 0.0_default, 0.0_default, 0.0_default, &
         62.5_default, 0.0_default, 0.0_default, 62.5_default, &
         62.5_default, 0.0_default, 0.0_default,-62.5_default &
         ], [4,3])

    write (u, "(A)")
    write (u, "(A)")  "* Compute matrix element:"
    write (u, "(A)")

    select type (driver)
    type is (prc_test_t)
       write (u, "(1x,A,1x,E11.4)") "|amp| =", abs (driver%get_amplitude (p))
    end select

    call lib%final ()
    
    write (u, "(A)")
    write (u, "(A)")  "* Test output end: prc_test_4"
    
  end subroutine prc_test_4
  
\eatline
\nwendcode{}\nwfilename{../../../src/particles/particles.nw}\nwbegindocs{0}%% -*- ess-noweb-default-code-mode: f90-mode; noweb-default-code-mode: f90-mode; -*- 
% WHIZARD code as NOWEB source: particle objects
%%%%%%%%%%%%%%%%%%%%%%%%%%%%%%%%%%%%%%%%%%%%%%%%%%%%%%%%%%%%%%%%%%%%%%% 
\chapter{Particles}

This chapter collects modules that implement particle objects, for use in
event records.

While within interactions, all correlations are
manifest, a particle array is derived by selecting a particular
quantum number set.  This involves tracing over all other particles,
as far as polarization is concerned.  Thus, a particle has definite
flavor, color, and a single-particle density matrix for polarization.
\begin{description}
\item[polarizations]
  The polarization object is based on a {\tt{}state{\char95}matrix{\char95}t} object.  It
  describes the spin density matrix of a particle.
\item[particles]
  Particle objects and particle lists, as the base of event records.
\end{description}

\clearpage
%%%%%%%%%%%%%%%%%%%%%%%%%%%%%%%%%%%%%%%%%%%%%%%%%%%%%%%%%%%%%%%%%%%%%%%%
\section{Polarization}
Particle polarization is determined by a particular quantum state
which has just helicity information.  For defining polarizations, we
adopt the phase convention for a spin-1/2 particle that
\begin{equation}
  \rho = \tfrac12(1 + \vec\alpha\cdot\vec\sigma)
\end{equation}
with the polarization axis $\vec\alpha$.  Using this, we define
\begin{enumerate}
\item Trivial polarization: $\vec\alpha=0$. [This is unpolarized, but
  distinct from the particular undefined polarization matrix which has
  the same meaning.]
\item Circular polarization: $\vec\alpha$ points in $\pm z$ direction.
\item Transversal polarization: $\vec\alpha$ points orthogonal to the
  $z$ direction, with a phase $\phi$ that is $0$ for the $x$ axis, and
  $\pi/2=90^\circ$ for the $y$ axis.  For antiparticles, the phase
  switches sign, corresponding to complex conjugation.
\item Axis polarization, where we explicitly give $\vec\alpha$.
\end{enumerate}
For higher spin, we retain this definition, but apply it to the two
components with maximum and minimum weight.  For massless particles,
this is sufficient.  For massive particles, we add the possibilities:
\begin{enumerate}\setcounter{enumi}{4}
\item Longitudinal polarization: Only the 0-component is set.  This is
  possible only for bosons.
\item Diagonal polarization: Explicitly specify all components in the
  helicity basis.
\end{enumerate}
Obviously, this does not exhaust the possible density matrices for
higher spin, but it should cover all practical applications.
\nwenddocs{}\nwbegincode{1}\moddef{\code{}polarizations.f90\edoc{}}\endmoddef
\LA{}File header\RA{}

module polarizations

\LA{}Use kinds\RA{}
\LA{}Use strings\RA{}
  use constants, only: imago
  use io_units
  use format_defs, only: FMT_12, FMT_19
  use unit_tests
  use diagnostics
  use physics_defs, only: SCALAR
  use lorentz
  use model_data
  use flavors
  use colors
  use helicities
  use quantum_numbers
  use state_matrices
!  use models

\LA{}Standard module head\RA{}

\LA{}Polarizations: public\RA{}

\LA{}Polarizations: types\RA{}

\LA{}Polarizations: interfaces\RA{}

contains

\LA{}Polarizations: procedures\RA{}

\LA{}Polarizations: tests\RA{}

end module polarizations
\eatline
\nwendcode{}\nwbegindocs{2}\nwdocspar
\subsection{The polarization type}
This is not an extension, but rather a restriction of the quantum
state.  Flavor and color are ignored, there is just a one-particle
helicity density matrix.
\nwenddocs{}\nwbegincode{3}\moddef{Polarizations: public}\endmoddef
  public :: polarization_t
\nwendcode{}\nwbegincode{4}\moddef{Polarizations: types}\endmoddef
  type :: polarization_t
     logical :: polarized = .false.
     integer :: spin_type = 0
     integer :: multiplicity = 0
     type(state_matrix_t) :: state
  end type polarization_t

\eatline
\nwendcode{}\nwbegindocs{5}\nwdocspar
\subsection{Basic initializer and finalizer}
We need the particle flavor for determining the allowed helicity
values.  The density matrix is not set, but prepared to be filled
later.  This is private.
\nwenddocs{}\nwbegincode{6}\moddef{Polarizations: procedures}\endmoddef
  subroutine polarization_init (pol, flv)
    type(polarization_t), intent(out) :: pol
    type(flavor_t), intent(in) :: flv
    pol%spin_type = flv%get_spin_type ()
    pol%multiplicity = flv%get_multiplicity ()
    call pol%state%init (store_values = .true.)
  end subroutine polarization_init
    
\eatline
\nwendcode{}\nwbegindocs{7}The finalizer has to be public.  The quantum state contains memory
allocated to pointers.
\nwenddocs{}\nwbegincode{8}\moddef{Polarizations: public}\plusendmoddef
  public :: polarization_final
\nwendcode{}\nwbegincode{9}\moddef{Polarizations: procedures}\plusendmoddef
  subroutine polarization_final (pol)
    type(polarization_t), intent(inout) :: pol
    call pol%state%final ()
  end subroutine polarization_final

\eatline
\nwendcode{}\nwbegindocs{10}\nwdocspar
\subsection{I/O}
\nwenddocs{}\nwbegincode{11}\moddef{Polarizations: public}\plusendmoddef
  public :: polarization_write
\nwendcode{}\nwbegincode{12}\moddef{Polarizations: procedures}\plusendmoddef
  subroutine polarization_write (pol, unit)
    type(polarization_t), intent(in) :: pol
    integer, intent(in), optional :: unit
    integer :: u
    u = given_output_unit (unit);  if (u < 0)  return
    write (u, "(1x,A,I1,A,I1,A)")  &
         "Polarization: [spin_type = ", pol%spin_type, &
         ", mult = ", pol%multiplicity, "]"
    call pol%state%write (unit=unit)
  end subroutine polarization_write

\eatline
\nwendcode{}\nwbegindocs{13}Defined assignment: deep copy
\nwenddocs{}\nwbegincode{14}\moddef{Polarizations: public}\plusendmoddef
  public :: assignment(=)
\nwendcode{}\nwbegincode{15}\moddef{Polarizations: interfaces}\endmoddef
  interface assignment(=)
     module procedure polarization_assign
  end interface

\nwendcode{}\nwbegincode{16}\moddef{Polarizations: procedures}\plusendmoddef
  subroutine polarization_assign (pol_out, pol_in)
    type(polarization_t), intent(out) :: pol_out
    type(polarization_t), intent(in) :: pol_in
    pol_out%polarized = pol_in%polarized
    pol_out%spin_type = pol_in%spin_type
    pol_out%multiplicity = pol_in%multiplicity
    pol_out%state = pol_in%state
  end subroutine polarization_assign

\eatline
\nwendcode{}\nwbegindocs{17}Binary I/O.
\nwenddocs{}\nwbegincode{18}\moddef{Polarizations: public}\plusendmoddef
  public :: polarization_write_raw
  public :: polarization_read_raw
\nwendcode{}\nwbegincode{19}\moddef{Polarizations: procedures}\plusendmoddef
  subroutine polarization_write_raw (pol, u)
    type(polarization_t), intent(in) :: pol
    integer, intent(in) :: u
    write (u) pol%polarized
    write (u) pol%spin_type
    write (u) pol%multiplicity
    call pol%state%write_raw (u)
  end subroutine polarization_write_raw

  subroutine polarization_read_raw (pol, u, iostat)
    type(polarization_t), intent(out) :: pol
    integer, intent(in) :: u
    integer, intent(out), optional :: iostat
    read (u, iostat=iostat) pol%polarized
    read (u, iostat=iostat) pol%spin_type
    read (u, iostat=iostat) pol%multiplicity
    call pol%state%read_raw (u, iostat=iostat)
  end subroutine polarization_read_raw

\eatline
\nwendcode{}\nwbegindocs{20}\nwdocspar
\subsection{Accessing contents}
Return true if the particle is polarized.  This is the case if the
first (and only) entry in the quantum state has undefined helicity.
\nwenddocs{}\nwbegincode{21}\moddef{Polarizations: public}\plusendmoddef
  public :: polarization_is_polarized
\nwendcode{}\nwbegincode{22}\moddef{Polarizations: procedures}\plusendmoddef
  elemental function polarization_is_polarized (pol) result (polarized)
    logical :: polarized
    type(polarization_t), intent(in) :: pol
    polarized = pol%polarized
  end function polarization_is_polarized

\eatline
\nwendcode{}\nwbegindocs{23}Return true if the polarization is diagonal, i.e., all entries in
the density matrix are diagonal.
\nwenddocs{}\nwbegincode{24}\moddef{Polarizations: public}\plusendmoddef
  public :: polarization_is_diagonal
\nwendcode{}\nwbegincode{25}\moddef{Polarizations: interfaces}\plusendmoddef
  interface polarization_is_diagonal
     module procedure polarization_is_diagonal0
     module procedure polarization_is_diagonal1
  end interface

\nwendcode{}\nwbegincode{26}\moddef{Polarizations: procedures}\plusendmoddef
  function polarization_is_diagonal0 (pol) result (diagonal)
    logical :: diagonal
    type(polarization_t), intent(in) :: pol
    type(state_iterator_t) :: it
    type(quantum_numbers_t), dimension(:), allocatable :: qn
    diagonal = .true.
    allocate (qn (pol%state%get_depth ()))
    call it%init (pol%state)
    do while (it%is_valid ())
       qn = it%get_quantum_numbers ()
       diagonal = all (qn%are_diagonal ())
       if (.not. diagonal) exit
       call it%advance ()
    end do
  end function polarization_is_diagonal0

  function polarization_is_diagonal1 (pol) result (diagonal)
    type(polarization_t), dimension(:), intent(in) :: pol
    logical, dimension(size(pol)) :: diagonal
    integer :: i
    do i = 1, size (pol)
       diagonal(i) = polarization_is_diagonal0 (pol(i))
    end do
  end function polarization_is_diagonal1

\eatline
\nwendcode{}\nwbegindocs{27}\nwdocspar
\subsection{Initialization from state matrix}
Here, the state matrix is already known (but not necessarily
normalized).  The result will be either unpolarized, or a normalized
spin density matrix.
\nwenddocs{}\nwbegincode{28}\moddef{Polarizations: public}\plusendmoddef
  public :: polarization_init_state_matrix
\nwendcode{}\nwbegincode{29}\moddef{Polarizations: procedures}\plusendmoddef
  subroutine polarization_init_state_matrix (pol, state)
    type(polarization_t), intent(out) :: pol
    type(state_matrix_t), intent(in), target :: state
    type(state_iterator_t) :: it
    type(flavor_t) :: flv
    type(helicity_t) :: hel
    type(quantum_numbers_t), dimension(1) :: qn
    complex(default) :: value, t
    call it%init (state)
    flv = it%get_flavor (1)
    hel = it%get_helicity (1)
    if (hel%is_defined ()) then
       call polarization_init (pol, flv)
       pol%polarized = .true.
       t = 0
       do while (it%is_valid ())
          hel = it%get_helicity (1)
          call qn(1)%init (hel)
          value = it%get_matrix_element ()
          call pol%state%add_state (qn, value=value)
          if (hel%is_diagonal ())  t = t + value
          call it%advance ()
       end do
       call pol%state%freeze ()
       if (t /= 0)  call pol%state%renormalize (1._default / t)
    else
       call polarization_init_unpolarized (pol, flv)
    end if
  end subroutine polarization_init_state_matrix

\eatline
\nwendcode{}\nwbegindocs{30}\nwdocspar
\subsection{Specific initializers}
Unpolarized particle, no helicity labels in the density matrix.  The
value is specified as $1/N$, where $N$ is the multiplicity.

Exception: for left-handed or right-handed particles (neutrinos),
polarization is always circular with fraction unity.
\nwenddocs{}\nwbegincode{31}\moddef{Polarizations: public}\plusendmoddef
  public :: polarization_init_unpolarized
\nwendcode{}\nwbegincode{32}\moddef{Polarizations: procedures}\plusendmoddef
  subroutine polarization_init_unpolarized (pol, flv)
    type(polarization_t), intent(inout) :: pol
    type(flavor_t), intent(in) :: flv
    type(quantum_numbers_t), dimension(1) :: qn
    complex(default) :: value
    if (flv%is_left_handed ()) then
       call polarization_init_circular (pol, flv, -1._default)
    else if (flv%is_right_handed ()) then
       call polarization_init_circular (pol, flv, 1._default)
    else
       call polarization_init (pol, flv)
       value = 1._default / flv%get_multiplicity ()
       call pol%state%add_state (qn)
       call pol%state%freeze ()
       call pol%state%set_matrix_element (value)
    end if
  end subroutine polarization_init_unpolarized
    
\eatline
\nwendcode{}\nwbegindocs{33}Unpolarized particle, but explicit density matrix with helicity
states allocated according to given flavor.  Note that fermions have
even spin type, bosons odd.  The spin density matrix entries are
scaled by {\tt{}fraction}.  This is used for initializing other
polarizations:
\begin{equation*}
  \rho(f) = 
  \frac{|f|}{N}\mathbf{1}.
\end{equation*}
\nwenddocs{}\nwbegincode{34}\moddef{Polarizations: public}\plusendmoddef
  public :: polarization_init_trivial
\nwendcode{}\nwbegincode{35}\moddef{Polarizations: procedures}\plusendmoddef
  subroutine polarization_init_trivial (pol, flv, fraction)
    type(polarization_t), intent(out) :: pol
    type(flavor_t), intent(in) :: flv
    real(default), intent(in), optional :: fraction
    type(helicity_t) :: hel
    type(quantum_numbers_t), dimension(1) :: qn
    integer :: h, hmax
    logical :: fermion
    complex(default) :: value
    call polarization_init (pol, flv)
    pol%polarized = .true.
    if (present (fraction)) then
       value = fraction / pol%multiplicity
    else
       value = 1._default / pol%multiplicity
    end if
    fermion = mod (pol%spin_type, 2) == 0
    hmax = pol%spin_type / 2
    select case (pol%multiplicity)
    case (1)
       if (flv%is_left_handed ()) then
          call hel%init (-hmax)
       else if (flv%is_right_handed ()) then
          call hel%init (hmax)
       else
          call hel%init (0)
       end if
       call qn(1)%init (hel)
       call pol%state%add_state (qn)
    case (2)
       do h = -hmax, hmax, 2*hmax
          call hel%init (h)
          call qn(1)%init (hel)
          call pol%state%add_state (qn)
       end do
    case default
       do h = -hmax, hmax
          if (fermion .and. h == 0)  cycle
          call hel%init (h)
          call qn(1)%init (hel)
          call pol%state%add_state (qn)
       end do
    end select
    call pol%state%freeze ()
    call pol%state%set_matrix_element (value)
  end subroutine polarization_init_trivial

\eatline
\nwendcode{}\nwbegindocs{36}The following three modes are useful mainly for spin-1/2 particle
and massless particles of any nonzero spin.  Only the highest-weight
components are filled.

Circular polarization: The density matrix of the two highest-weight
states is
\begin{equation*}
  \rho(f) = 
  \frac{1-|f|}{2}\mathbf{1} +
  |f| \times
  \begin{cases}
    \begin{pmatrix} 1 & 0 \\ 0 & 0 \end{pmatrix}, & f > 0; \\[6pt]
    \begin{pmatrix} 0 & 0 \\ 0 & 1 \end{pmatrix}, & f < 0,
  \end{cases}
\end{equation*}
If the polarization fraction $|f|$ is unity, we need only one entry in the
density matrix.
\nwenddocs{}\nwbegincode{37}\moddef{Polarizations: public}\plusendmoddef
  public :: polarization_init_circular
\nwendcode{}\nwbegincode{38}\moddef{Polarizations: procedures}\plusendmoddef
  subroutine polarization_init_circular (pol, flv, fraction)
    type(polarization_t), intent(out) :: pol
    type(flavor_t), intent(in) :: flv
    real(default), intent(in) :: fraction
    type(helicity_t), dimension(2) :: hel
    type(quantum_numbers_t), dimension(1) :: qn
    complex(default) :: value
    integer :: hmax
    call polarization_init (pol, flv)
    pol%polarized = .true.
    hmax = pol%spin_type / 2
    call hel(1)%init ( hmax)
    call hel(2)%init (-hmax)
    if (abs (fraction) /= 1) then
       value = (1 + fraction) / 2
       call qn(1)%init (hel(1))
       call pol%state%add_state (qn, value=value)
       value = (1 - fraction) / 2
       call qn(1)%init (hel(2))
       call pol%state%add_state (qn, value=value)
    else
       value = abs (fraction)
       if (fraction > 0) then
          call qn(1)%init (hel(1))
       else
          call qn(1)%init (hel(2))
       end if
       call pol%state%add_state (qn, value=value)
    end if
    call pol%state%freeze ()
  end subroutine polarization_init_circular

\eatline
\nwendcode{}\nwbegindocs{39}Transversal polarization is analogous to circular, but we get a
density matrix
\begin{equation*}
  \rho(f,\phi) = 
  \frac{1-|f|}{2}\mathbf{1}
  + \frac{|f|}{2}  \begin{pmatrix} 1 & e^{-i\phi} \\ e^{i\phi} & 1
  \end{pmatrix}.
\end{equation*}
The phase is $\phi=0$ for the $x$-axis, $\phi=90^\circ$ for the $y$
axis as polarization vector.  For an antiparticle, the phase switches
sign, and for $f<0$, the off-diagonal elements switch sign.  
\nwenddocs{}\nwbegincode{40}\moddef{Polarizations: public}\plusendmoddef
  public :: polarization_init_transversal
\nwendcode{}\nwbegincode{41}\moddef{Polarizations: procedures}\plusendmoddef
  subroutine polarization_init_transversal (pol, flv, phi, fraction)
    type(polarization_t), intent(inout) :: pol
    type(flavor_t), intent(in) :: flv
    real(default), intent(in) :: phi, fraction
    call polarization_init_axis &
         (pol, flv, fraction * [ cos (phi), sin (phi), 0._default])
  end subroutine polarization_init_transversal

\eatline
\nwendcode{}\nwbegindocs{42}For axis polarization, we again set only the entries with maximum weight.
\begin{equation*}
  \rho(f,\phi) = 
  \frac{1}{2} \begin{pmatrix} 
                 1 + \alpha_3 & \alpha_1 - i\alpha_2 \\ 
                 \alpha_1 + i\alpha_2 & 1 - \alpha_3 
              \end{pmatrix}.
\end{equation*}
For an antiparticle, $\alpha_2$ switches sign (complex conjugate).
\nwenddocs{}\nwbegincode{43}\moddef{Polarizations: public}\plusendmoddef
  public :: polarization_init_axis
\nwendcode{}\nwbegincode{44}\moddef{Polarizations: procedures}\plusendmoddef
  subroutine polarization_init_axis (pol, flv, alpha)
    type(polarization_t), intent(out) :: pol
    type(flavor_t), intent(in) :: flv
    real(default), dimension(3), intent(in) :: alpha
    type(quantum_numbers_t), dimension(1) :: qn
    type(helicity_t), dimension(2,2) :: hel
    complex(default), dimension(2,2) :: value
    integer :: hmax
    call polarization_init (pol, flv)
    pol%polarized = .true.
    hmax = pol%spin_type / 2
    call hel(1,1)%init ( hmax, hmax)
    call hel(1,2)%init ( hmax,-hmax)
    call hel(2,1)%init (-hmax, hmax)
    call hel(2,2)%init (-hmax,-hmax)
    value(1,1) = (1 + alpha(3)) / 2
    value(2,2) = (1 - alpha(3)) / 2
    if (flv%is_antiparticle ()) then
       value(1,2) = (alpha(1) + imago * alpha(2)) / 2
    else
       value(1,2) = (alpha(1) - imago * alpha(2)) / 2
    end if
    value(2,1) = conjg (value(1,2))
    if (value(1,1) /= 0) then
       call qn(1)%init (hel(1,1))
       call pol%state%add_state (qn, value=value(1,1))
    end if
    if (value(2,2) /= 0) then
       call qn(1)%init (hel(2,2))
       call pol%state%add_state (qn, value=value(2,2))
    end if
    if (value(1,2) /= 0) then
       call qn(1)%init (hel(1,2))
       call pol%state%add_state (qn, value=value(1,2))
       call qn(1)%init (hel(2,1))
       call pol%state%add_state (qn, value=value(2,1))
    end if
    call pol%state%freeze ()
  end subroutine polarization_init_axis

\eatline
\nwendcode{}\nwbegindocs{45}This version specifies the polarization axis in terms of $r$
(polarization degree) and $\theta,\phi$ (polar and azimuthal angles).

If one of the angles is a nonzero multiple of $\pi$, roundoff errors
typically will result in tiny contributions to unwanted components.
Therefore, include a catch for small numbers.
\nwenddocs{}\nwbegincode{46}\moddef{Polarizations: public}\plusendmoddef
  public :: polarization_init_angles
\nwendcode{}\nwbegincode{47}\moddef{Polarizations: procedures}\plusendmoddef
  subroutine polarization_init_angles (pol, flv, r, theta, phi)
    type(polarization_t), intent(out) :: pol
    type(flavor_t), intent(in) :: flv
    real(default), intent(in) :: r, theta, phi
    real(default), dimension(3) :: alpha
    real(default), parameter :: eps = 10 * epsilon (1._default)
    alpha(1) = r * sin (theta) * cos (phi)
    alpha(2) = r * sin (theta) * sin (phi)
    alpha(3) = r * cos (theta)
    where (abs (alpha) < eps)  alpha = 0
    call polarization_init_axis (pol, flv, alpha)
  end subroutine polarization_init_angles

\eatline
\nwendcode{}\nwbegindocs{48}Longitudinal polarization is defined only for massive bosons.  Only
the zero component is filled.  Otherwise, unpolarized.
\nwenddocs{}\nwbegincode{49}\moddef{Polarizations: public}\plusendmoddef
  public :: polarization_init_longitudinal
\nwendcode{}\nwbegincode{50}\moddef{Polarizations: procedures}\plusendmoddef
  subroutine polarization_init_longitudinal (pol, flv, fraction)
    type(polarization_t), intent(out) :: pol
    type(flavor_t), intent(in) :: flv
    real(default), intent(in) :: fraction
    integer :: spin_type, multiplicity
    type(helicity_t) :: hel
    type(quantum_numbers_t), dimension(1) :: qn
    complex(default) :: value
    integer :: n_values
    value = abs (fraction)
    spin_type = flv%get_spin_type ()
    multiplicity = flv%get_multiplicity ()
    if (mod (spin_type, 2) == 1 .and. multiplicity > 2) then
       if (fraction /= 1) then
          call polarization_init_trivial (pol, flv, 1 - fraction)
          n_values = pol%state%get_n_matrix_elements ()
          call pol%state%add_to_matrix_element (n_values/2 + 1, value)
       else
          call polarization_init (pol, flv)
          pol%polarized = .true.
          call hel%init (0)
          call qn(1)%init (hel)
          call pol%state%add_state (qn)
          call pol%state%freeze ()
          call pol%state%set_matrix_element (value)
       end if
    else
       call polarization_init_unpolarized (pol, flv)
    end if
  end subroutine polarization_init_longitudinal

\eatline
\nwendcode{}\nwbegindocs{51}This is diagonal polarization: we specify all components explicitly.
We use only the positive components.  The sum is normalized to unity.
We assume that the length of {\tt{}alpha} is equal to the particle
multiplicity.
\nwenddocs{}\nwbegincode{52}\moddef{Polarizations: public}\plusendmoddef
  public :: polarization_init_diagonal
\nwendcode{}\nwbegincode{53}\moddef{Polarizations: procedures}\plusendmoddef
  subroutine polarization_init_diagonal (pol, flv, alpha)
    type(polarization_t), intent(inout) :: pol
    type(flavor_t), intent(in) :: flv
    real(default), dimension(:), intent(in) :: alpha
    type(helicity_t) :: hel
    type(quantum_numbers_t), dimension(1) :: qn
    logical, dimension(size(alpha)) :: mask
    real(default) :: norm
    complex(default), dimension(:), allocatable :: value
    logical :: fermion
    integer :: h, hmax, i
    mask = alpha > 0
    norm = sum (alpha, mask);  if (norm == 0)  norm = 1
    allocate (value (count (mask)))
    value = pack (alpha / norm, mask)
    call polarization_init (pol, flv)
    pol%polarized = .true.
    fermion = mod (pol%spin_type, 2) == 0
    hmax = pol%spin_type / 2
    i = 0
    select case (pol%multiplicity)
    case (1)
       if (flv%is_left_handed ()) then
          call hel%init (-hmax)
       else if (flv%is_right_handed ()) then
          call hel%init ( hmax)
       else
          call hel%init (0)
       end if
       call qn(1)%init (hel)
       call pol%state%add_state (qn)
    case (2)
       do h = -hmax, hmax, 2*hmax
          i = i + 1
          if (mask(i)) then
             call hel%init (h)
             call qn(1)%init (hel)
             call pol%state%add_state (qn)
          end if
       end do
    case default
       do h = -hmax, hmax
          if (fermion .and. h == 0)  cycle
          i = i + 1
          if (mask(i)) then
             call hel%init (h)
             call qn(1)%init (hel)
             call pol%state%add_state (qn)
          end if
       end do
    end select
    call pol%state%freeze ()
    call pol%state%set_matrix_element (value)
  end subroutine polarization_init_diagonal

\eatline
\nwendcode{}\nwbegindocs{54}Generic polarization: we generate all possible density matrix
entries, but the values are left zero.
\nwenddocs{}\nwbegincode{55}\moddef{Polarizations: public}\plusendmoddef
  public :: polarization_init_generic
\nwendcode{}\nwbegincode{56}\moddef{Polarizations: procedures}\plusendmoddef
  subroutine polarization_init_generic (pol, flv)
    type(polarization_t), intent(out) :: pol
    type(flavor_t), intent(in) :: flv
    type(helicity_t) :: hel
    type(quantum_numbers_t), dimension(1) :: qn
    logical :: fermion
    integer :: hmax, h1, h2
    call polarization_init (pol, flv)
    pol%polarized = .true.
    fermion = mod (pol%spin_type, 2) == 0
    hmax = pol%spin_type / 2
    select case (pol%multiplicity)
    case (1)
       if (flv%is_left_handed ()) then
          call hel%init (-hmax)
       else if (flv%is_right_handed ()) then
          call hel%init ( hmax)
       else
          call hel%init (0)
       end if
       call qn(1)%init (hel)
       call pol%state%add_state (qn)
    case (2)
       do h1 = -hmax, hmax, 2*hmax
          do h2 = -hmax, hmax, 2*hmax
             call hel%init (h1, h2)
             call qn(1)%init (hel)
             call pol%state%add_state (qn)
          end do
       end do
    case default
       do h1 = -hmax, hmax
          if (fermion .and. h1 == 0)  cycle
          do h2 = -hmax, hmax
             if (fermion .and. h2 == 0)  cycle
             call hel%init (h1, h2)
             call qn(1)%init (hel)
             call pol%state%add_state (qn)
          end do
       end do
    end select
    call pol%state%freeze ()
  end subroutine polarization_init_generic

\eatline
\nwendcode{}\nwbegindocs{57}\nwdocspar
\subsection{Operations}
Combine polarization states by computing the outer product of the
state matrices.
\nwenddocs{}\nwbegincode{58}\moddef{Polarizations: public}\plusendmoddef
  public :: combine_polarization_states
\nwendcode{}\nwbegincode{59}\moddef{Polarizations: procedures}\plusendmoddef
  subroutine combine_polarization_states (pol, state)
    type(polarization_t), dimension(:), intent(in), target :: pol
    type(state_matrix_t), intent(out) :: state
    call outer_multiply (pol%state, state)
  end subroutine combine_polarization_states

\eatline
\nwendcode{}\nwbegindocs{60}\nwdocspar
Transform a polarization density matrix into a polarization vector.
This is possible without information loss only for spin-1/2 and for
massless particles.  To get a unique answer in all cases, we consider
only the components with highest weight.  Obviously, this loses the
longitudinal component of a massive vector, for instance.

This is the inverse operation of {\tt{}polarization{\char95}init{\char95}axis} above,
where the polarization fraction is set to unity.
\nwenddocs{}\nwbegincode{61}\moddef{Polarizations: public}\plusendmoddef
  public :: polarization_get_axis
\nwendcode{}\nwbegincode{62}\moddef{Polarizations: procedures}\plusendmoddef
  function polarization_get_axis (pol) result (alpha)
    real(default), dimension(3) :: alpha
    type(polarization_t), intent(in), target :: pol
    type(state_iterator_t) :: it
    complex(default), dimension(2,2) :: value
    type(helicity_t), dimension(2,2) :: hel
    type(helicity_t), dimension(1) :: hel1
    integer :: hmax, i, j
    if (pol%polarized) then
       hmax = pol%spin_type / 2
       call hel(1,1)%init ( hmax, hmax)
       call hel(1,2)%init ( hmax,-hmax)
       call hel(2,1)%init (-hmax, hmax)
       call hel(2,2)%init (-hmax,-hmax)
       value = 0
       call it%init (pol%state)
       do while (it%is_valid ())
          hel1 = it%get_helicity ()
          SCAN_HEL: do i = 1, 2
             do j = 1, 2
                if (hel1(1) == hel(i,j)) then
                   value(i,j) = it%get_matrix_element ()
                   exit SCAN_HEL
                end if
             end do
          end do SCAN_HEL
          call it%advance ()
       end do
       alpha(1) = real(value(1,2) + value(2,1))
       alpha(2) = - aimag(value(1,2) - value(2,1))
       alpha(3) = real(value(1,1) - value(2,2))
    else
       alpha = 0
    end if
  end function polarization_get_axis

\eatline
\nwendcode{}\nwbegindocs{63}This function returns polarization degree and polar and azimuthal
angles ($\theta,\phi$) of the polarization axis.
\nwenddocs{}\nwbegincode{64}\moddef{Polarizations: public}\plusendmoddef
  public :: polarization_to_angles
\nwendcode{}\nwbegincode{65}\moddef{Polarizations: procedures}\plusendmoddef
  subroutine polarization_to_angles (pol, r, theta, phi)
    type(polarization_t), intent(in) :: pol
    real(default), intent(out) :: r, theta, phi
    real(default), dimension(3) :: alpha
    real(default) :: r12
    if (pol%polarized) then
       alpha = polarization_get_axis (pol)
       r = sqrt (sum (alpha**2))
       if (any (alpha /= 0)) then
          r12 = sqrt (alpha(1)**2 + alpha(2)**2)
          theta = atan2 (r12, alpha(3))
          if (any (alpha(1:2) /= 0)) then
             phi = atan2 (alpha(2), alpha(1))
          else
             phi = 0
          end if
       else
          theta = 0
       end if
    else
       r = 0
       theta = 0
       phi = 0
    end if
  end subroutine polarization_to_angles

\eatline
\nwendcode{}\nwbegindocs{66}\nwdocspar
\subsection{Sparse Matrix}
We introduce a simpler implementation of a sparse matrix that should represent
polarization.  It consists of an integer array that represents the index
values, and a complex array that represents the nonvanishing entries.  The
number of nonvanishing entries must be known for initialization, but the
entries are filled one at a time.

This should eventually replace the more baroque polarization definition above,
but in the first stage we introduce it just as an intermediage storage
container.

Here is a base type without the special properties of a spin-density matrix.
\nwenddocs{}\nwbegincode{67}\moddef{Polarizations: public}\plusendmoddef
  public :: smatrix_t
\nwendcode{}\nwbegincode{68}\moddef{Polarizations: types}\plusendmoddef
  type :: smatrix_t
     private
     integer :: dim = 0
     integer :: n_entry = 0
     integer, dimension(:,:), allocatable :: index
     complex(default), dimension(:), allocatable :: value
   contains
   \LA{}Polarizations: smatrix: TBP\RA{}
  end type smatrix_t
  
\eatline
\nwendcode{}\nwbegindocs{69}Output.
\nwenddocs{}\nwbegincode{70}\moddef{Polarizations: smatrix: TBP}\endmoddef
  procedure :: write => smatrix_write
\nwendcode{}\nwbegincode{71}\moddef{Polarizations: procedures}\plusendmoddef
  subroutine smatrix_write (object, unit, indent)
    class(smatrix_t), intent(in) :: object
    integer, intent(in), optional :: unit, indent
    integer :: u, i, ind
    u = given_output_unit (unit)
    ind = 0;  if (present (indent))  ind = indent
    if (allocated (object%value)) then
       if (size (object%value) > 0) then
          do i = 1, object%n_entry
             write (u, "(1x,A,'@(')", advance="no")  repeat ("  ", ind)
             write (u, "(SP,9999(I2.1,':',1x))", advance="no") &
                  object%index(:,i)
             write (u, "('('," // FMT_19 // ",','," // FMT_19 // &
                  ",'))')")  object%value(i)
          end do
       else
          write (u, "(1x,A)", advance="no")  repeat ("  ", ind)
          write (u, "(A)")  "[empty matrix]"
       end if
    else
       write (u, "(1x,A)", advance="no")  repeat ("  ", ind)
       write (u, "(A)")  "[undefined matrix]"
    end if
  end subroutine smatrix_write
  
\eatline
\nwendcode{}\nwbegindocs{72}Initialization: allocate arrays to the correct size.  We specify both the
dimension of the matrix (if different from two, this is rather a generic
tensor) and the number of nonvanishing entries.
\nwenddocs{}\nwbegincode{73}\moddef{Polarizations: smatrix: TBP}\plusendmoddef
  procedure :: init => smatrix_init
\nwendcode{}\nwbegincode{74}\moddef{Polarizations: procedures}\plusendmoddef
  subroutine smatrix_init (smatrix, dim, n_entry)
    class(smatrix_t), intent(out) :: smatrix
    integer, intent(in) :: dim
    integer, intent(in) :: n_entry
    smatrix%dim = dim
    smatrix%n_entry = n_entry
    allocate (smatrix%index (dim, n_entry))
    allocate (smatrix%value (n_entry))
  end subroutine smatrix_init
  
\eatline
\nwendcode{}\nwbegindocs{75}Fill: one entry at a time.
\nwenddocs{}\nwbegincode{76}\moddef{Polarizations: smatrix: TBP}\plusendmoddef
  procedure :: set_entry => smatrix_set_entry
\nwendcode{}\nwbegincode{77}\moddef{Polarizations: procedures}\plusendmoddef
  subroutine smatrix_set_entry (smatrix, i, index, value)
    class(smatrix_t), intent(inout) :: smatrix
    integer, intent(in) :: i
    integer, dimension(:), intent(in) :: index
    complex(default), intent(in) :: value
    smatrix%index(:,i) = index
    smatrix%value(i) = value
  end subroutine smatrix_set_entry
  
\eatline
\nwendcode{}\nwbegindocs{78}\nwdocspar
\subsection{Polarization Matrix}
As an extension of the more generic {\tt{}smatrix} type, we implement a proper
spin-density matrix.  After the matrix has been filled, we can fix spin type
and multiplicity for a particle, check the matrix for consistency, and
normalize it if necessary.
\nwenddocs{}\nwbegincode{79}\moddef{Polarizations: public}\plusendmoddef
  public :: pmatrix_t
\nwendcode{}\nwbegincode{80}\moddef{Polarizations: types}\plusendmoddef
  type, extends (smatrix_t) :: pmatrix_t
     private
     integer :: spin_type = 0
     integer :: multiplicity = 0
     logical :: massive = .true.
     integer :: chirality = 0
     real(default) :: degree = 1
     logical :: pure = .false.
   contains
   \LA{}Polarizations: pmatrix: TBP\RA{}
  end type pmatrix_t
  
\eatline
\nwendcode{}\nwbegindocs{81}Output, including extra data.  (The {\tt{}indent} argument is ignored.)
\nwenddocs{}\nwbegincode{82}\moddef{Polarizations: pmatrix: TBP}\endmoddef
  procedure :: write => pmatrix_write
\nwendcode{}\nwbegincode{83}\moddef{Polarizations: procedures}\plusendmoddef
  subroutine pmatrix_write (object, unit, indent)
    class(pmatrix_t), intent(in) :: object
    integer, intent(in), optional :: unit, indent
    integer :: u
    u = given_output_unit (unit)
    write (u, "(1x,A)")  "Polarization: spin density matrix"
    write (u, "(3x,A,I0)")  "spin type     = ", object%spin_type
    write (u, "(3x,A,I0)")  "multiplicity  = ", object%multiplicity
    write (u, "(3x,A,L1)")  "massive       = ", object%massive
    write (u, "(3x,A,I0)")  "chirality     = ", object%chirality
    write (u, "(3x,A,F10.7)")  "pol.degree    =", object%degree
    write (u, "(3x,A,L1)")  "pure state    = ", object%pure
    call object%smatrix_t%write (u, 1)
  end subroutine pmatrix_write
  
\eatline
\nwendcode{}\nwbegindocs{84}This assignment is trivial, but must be coded explicitly.
\nwenddocs{}\nwbegincode{85}\moddef{Polarizations: pmatrix: TBP}\plusendmoddef
  generic :: assignment(=) => pmatrix_assign_from_smatrix
  procedure, private :: pmatrix_assign_from_smatrix
\nwendcode{}\nwbegincode{86}\moddef{Polarizations: procedures}\plusendmoddef
  subroutine pmatrix_assign_from_smatrix (pmatrix, smatrix)
    class(pmatrix_t), intent(out) :: pmatrix
    type(smatrix_t), intent(in) :: smatrix
    pmatrix%smatrix_t = smatrix
  end subroutine pmatrix_assign_from_smatrix
  
\eatline
\nwendcode{}\nwbegindocs{87}Declare spin, multiplicity, and polarization degree.  Check whether all
entries fit, and whether this is a valid matrix.

The required properties are:
\begin{enumerate}
\item all entries apply to the given spin and mass type
\item the diagonal is real
\item only the upper of corresponding off-diagonal elements is specified,
  i.e., the row index is less than the column index
\item the trace is nonnegative and equal to the polarization degree (the
  remainder, proportional to the unit matrix, is understood to be present)
\item the trace of the matrix square is positive and less or equal
  to the trace of the matrix itself, which is the polarization degree.  (If it
  is equal, we are dealing with a pure state).
\end{enumerate}
\nwenddocs{}\nwbegincode{88}\moddef{Polarizations: pmatrix: TBP}\plusendmoddef
  procedure :: normalize => pmatrix_normalize
\nwendcode{}\nwbegincode{89}\moddef{Polarizations: procedures}\plusendmoddef
  subroutine pmatrix_normalize (pmatrix, flv, degree, tolerance)
    class(pmatrix_t), intent(inout) :: pmatrix
    type(flavor_t), intent(in) :: flv
    real(default), intent(in), optional :: degree
    real(default), intent(in), optional :: tolerance
    integer :: i, hmax
    logical :: fermion, ok
    real(default) :: trace, trace_sq
    real(default) :: tol
    tol = 0;  if (present (tolerance))  tol = tolerance
    pmatrix%spin_type = flv%get_spin_type ()
    pmatrix%massive = flv%get_mass () /= 0
    if (.not. pmatrix%massive) then
       if (flv%is_left_handed ()) then
          pmatrix%chirality = -1
       else if (flv%is_right_handed ()) then
          pmatrix%chirality = +1
       end if
    end if
    if (pmatrix%spin_type == SCALAR) then
       pmatrix%multiplicity = 1
    else if (pmatrix%massive) then
       pmatrix%multiplicity = pmatrix%spin_type
    else if (pmatrix%chirality == 0) then
       pmatrix%multiplicity = 2
    else
       pmatrix%multiplicity = 1
    end if
    if (present (degree)) then
       if (degree < 0 .or. degree > 1) &
            call msg_error ("polarization degree must be between 0 and 1")
       pmatrix%degree = degree
    end if
    if (size (pmatrix%index, 1) /= 2)  call error ("wrong array rank")
    fermion = mod (pmatrix%spin_type, 2) == 0
    hmax = pmatrix%spin_type / 2
    if (pmatrix%n_entry > 0) then
       if (fermion) then
          if (pmatrix%massive) then
             ok = all (pmatrix%index /= 0) &
                  .and. all (abs (pmatrix%index) <= hmax)
          else if (pmatrix%chirality == -1) then
             ok = all (pmatrix%index == -hmax)
          else if (pmatrix%chirality == +1) then
             ok = all (pmatrix%index == +hmax)
          else
             ok = all (abs (pmatrix%index) == hmax)
          end if
       else
          if (pmatrix%massive) then
             ok = all (abs (pmatrix%index) <= hmax)
          else
             ok = all (abs (pmatrix%index) == hmax)
          end if
       end if
       if (.not. ok)  call error ("illegal index value")
    else
       pmatrix%degree = 0
       pmatrix%pure = pmatrix%multiplicity == 1
       return
    end if
    trace = 0
    do i = 1, pmatrix%n_entry
       associate (index => pmatrix%index(:,i), value => pmatrix%value(i))
         if (index(1) == index(2)) then
            if (abs (aimag (value)) > tol)  call error ("diagonal must be real")
            value = real (value, kind=default)
            trace = trace + value
            
         else if (any (pmatrix%index(1,:) == index(2) &
              .and.    pmatrix%index(2,:) == index(1))) then
            call error ("redundant off-diagonal entry")
         else if (index(2) < index (1)) then
            index = index([2,1])
            value = conjg (value)
         end if
       end associate
    end do
    if (abs (trace) <= tol)  call error ("trace must not vanish")
    trace = real (trace, kind=default)
    pmatrix%value = pmatrix%value / trace * pmatrix%degree
    trace_sq = (1 - pmatrix%degree ** 2) / pmatrix%multiplicity
    do i = 1, pmatrix%n_entry
       associate (index => pmatrix%index(:,i), value => pmatrix%value(i))
         if (index(1) == index(2)) then
            trace_sq = trace_sq + abs (value) ** 2
         else
            trace_sq = trace_sq + 2 * abs (value) ** 2
         end if
       end associate
    end do
    if (pmatrix%multiplicity == 1) then
       pmatrix%pure = .true.
    else if (abs (trace_sq - 1) <= tol) then
       pmatrix%pure = .true.
    else if (trace_sq - 1 > tol .or. trace_sq < -tol) then
       print *, "Trace of matrix square = ", trace_sq
       call error ("not permissible as density matrix")
    end if
  contains
    subroutine error (msg)
      character(*), intent(in) :: msg
      call pmatrix%write ()
      call msg_fatal ("Spin density matrix: " // msg)
    end subroutine error
  end subroutine pmatrix_normalize
  
\eatline
\nwendcode{}\nwbegindocs{90}\nwdocspar
A polarized matrix is defined as one with a positive polarization degree, even
if the actual matrix is trivial.
\nwenddocs{}\nwbegincode{91}\moddef{Polarizations: pmatrix: TBP}\plusendmoddef
  procedure :: is_polarized => pmatrix_is_polarized
\nwendcode{}\nwbegincode{92}\moddef{Polarizations: procedures}\plusendmoddef
  elemental function pmatrix_is_polarized (pmatrix) result (flag)
    class(pmatrix_t), intent(in) :: pmatrix
    logical :: flag
    flag = pmatrix%degree > 0
  end function pmatrix_is_polarized
  
\eatline
\nwendcode{}\nwbegindocs{93}\nwdocspar
Check if there are only diagonal entries.
\nwenddocs{}\nwbegincode{94}\moddef{Polarizations: pmatrix: TBP}\plusendmoddef
  procedure :: is_diagonal => pmatrix_is_diagonal
\nwendcode{}\nwbegincode{95}\moddef{Polarizations: procedures}\plusendmoddef
  elemental function pmatrix_is_diagonal (pmatrix) result (flag)
    class(pmatrix_t), intent(in) :: pmatrix
    logical :: flag
    flag = all (pmatrix%index(1,:) == pmatrix%index(2,:))
  end function pmatrix_is_diagonal
  
\eatline
\nwendcode{}\nwbegindocs{96}\nwdocspar
\subsection{Data Translation}
Create a {\tt{}polarization{\char95}t} object from the contents of a normalized
{\tt{}pmatrix{\char95}t} object.

Note for off diagonal entries: the convention for the arguments of
{\tt{}helicity{\char95}init} is {\tt{}(h2,h1)}, so the helicity of the conjugate state is
inserted first.
\nwenddocs{}\nwbegincode{97}\moddef{Polarizations: public}\plusendmoddef
  public :: polarization_init_pmatrix
\nwendcode{}\nwbegincode{98}\moddef{Polarizations: procedures}\plusendmoddef
  subroutine polarization_init_pmatrix (pol, pmatrix)
    type(polarization_t), intent(out) :: pol
    type(pmatrix_t), intent(in) :: pmatrix
    type(quantum_numbers_t), dimension(1) :: qn
    type(helicity_t) :: hel
    integer :: i, h, h1, h2, hmin, hmax, dh
    logical :: fermion
    complex(default) :: value
    pol%polarized = .true.
    pol%spin_type = pmatrix%spin_type
    pol%multiplicity = pmatrix%multiplicity
    call pol%state%init (store_values = .true.)
    fermion = mod (pol%spin_type, 2) == 0
    h = pol%spin_type / 2
    select case (pmatrix%chirality)
    case (-1)
       hmin = -h
       hmax = -h
    case (0)
       hmin = -h
       hmax = h
    case (1)
       hmin = h
       hmax = h
    end select
    if (pol%multiplicity == 1) then
       dh = 1
    else if (pol%multiplicity == 2) then
       dh = hmax - hmin
    else
       dh = 1
    end if
    if (pmatrix%degree < 1) then
       value = (1 - pmatrix%degree) / pol%multiplicity
       do h = hmin, hmax, dh
          if (h == 0 .and. fermion)  cycle
          call hel%init (h)
          call qn(1)%init (hel)
          call pol%state%add_state (qn, value = value)
       end do
    end if
    do i = 1, pmatrix%n_entry
       associate (index => pmatrix%index(:,i), value => pmatrix%value(i))
         h1 = index(1)
         h2 = index(2)
         if (h1 == h2) then
            call hel%init (h1)
            call qn(1)%init (hel)
            call pol%state%add_state (qn, value = value, &
                 sum_values = .true.)
         else
            call hel%init (h2, h1)
            call qn(1)%init (hel)
            call pol%state%add_state (qn, value = value)
            call hel%init (h1, h2)
            call qn(1)%init (hel)
            call pol%state%add_state (qn, value = conjg (value))
         end if
       end associate
    end do
    call pol%state%freeze ()
  end subroutine polarization_init_pmatrix

\eatline
\nwendcode{}\nwbegindocs{99}\subsection{Unit tests}
\nwenddocs{}\nwbegincode{100}\moddef{Polarizations: public}\plusendmoddef
  public :: polarization_test
\nwendcode{}\nwbegincode{101}\moddef{Polarizations: procedures}\plusendmoddef
  subroutine polarization_test (u, results)
    integer, intent(in) :: u
    type(test_results_t), intent(inout) :: results
  \LA{}Polarizations: execute tests\RA{}
  end subroutine polarization_test

\nwendcode{}\nwbegindocs{102} %def polarization_test
\nwenddocs{}\nwbegindocs{103}\nwdocspar
\subsubsection{Polarization type}
Checking the setup for polarization. 
\nwenddocs{}\nwbegincode{104}\moddef{Polarizations: execute tests}\endmoddef
  call test (polarization_1, "polarization_1", &
       "check polarization setup", &
       u, results)
\nwendcode{}\nwbegincode{105}\moddef{Polarizations: tests}\endmoddef
  subroutine polarization_1 (u)
    use os_interface
    integer, intent(in) :: u
    type(model_data_t), target :: model
    type(polarization_t) :: pol
    type(flavor_t) :: flv
    real(default), dimension(3) :: alpha
    real(default) :: r, theta, phi

    write (u, "(A)")  "* Test output: polarization_1"
    write (u, "(A)")  "*   Purpose: test polarization setup"
    write (u, "(A)")      
    
    write (u, "(A)")  "* Reading model file"
    write (u, "(A)")      
    
    call model%init_sm_test ()

    write (u, "(A)") "* Unpolarized fermion"
    write (u, "(A)")
    
    call flv%init (1, model)
    call polarization_init_unpolarized (pol, flv)
    call polarization_write (pol, u)
    write (u, "(A,L1)")  "   diagonal =", polarization_is_diagonal (pol)
    call polarization_final (pol)
    
    write (u, "(A)") 
    write (u, "(A)")  "* Unpolarized fermion"
    write (u, "(A)") 
    
    call polarization_init_circular (pol, flv, 0._default)
    call polarization_write (pol, u)
    call polarization_final (pol)
    
    write (u, "(A)")
    write (u, "(A)")  "* Transversally polarized fermion, phi=0"
    write (u, "(A)")
    
    call polarization_init_transversal (pol, flv, 0._default, 1._default)
    call polarization_write (pol, u)
    write (u, "(A,L1)")  "   diagonal =", polarization_is_diagonal (pol)
    call polarization_final (pol)
    
    write (u, "(A)") 
    write (u, "(A)")  "* Transversally polarized fermion, phi=0.9, frac=0.8"
    write (u, "(A)")
    
    call polarization_init_transversal (pol, flv, 0.9_default, 0.8_default)
    call polarization_write (pol, u)
    write (u, "(A,L1)")  "   diagonal =", polarization_is_diagonal (pol)
    call polarization_final (pol)
    
    write (u, "(A)")
    write (u, "(A)") "* All polarization directions of a fermion"
    write (u, "(A)")
    
    call polarization_init_generic (pol, flv)
    call polarization_write (pol, u)
    call polarization_final (pol)
    call flv%init (21, model)
    
    write (u, "(A)") 
    write (u, "(A)")  "* Circularly polarized gluon, frac=0.3"
    write (u, "(A)") 
    
    call polarization_init_circular (pol, flv, 0.3_default)
    call polarization_write (pol, u)
    call polarization_final (pol)   
    call flv%init (23, model)
    
    write (u, "(A)") 
    write (u, "(A)") "* Circularly polarized massive vector, frac=-0.7"
    write (u, "(A)") 
    
    call polarization_init_circular (pol, flv,  -0.7_default)
    call polarization_write (pol, u)
    call polarization_final (pol)
    
    write (u, "(A)") 
    write (u, "(A)")  "* Circularly polarized massive vector"
    write (u, "(A)")
    
    call polarization_init_circular (pol, flv, 1._default)
    call polarization_write (pol, u)
    call polarization_final (pol)
    
    write (u, "(A)") 
    write (u, "(A)")  "* Longitudinally polarized massive vector, frac=0.4"
    write (u, "(A)")
    
    call polarization_init_longitudinal (pol, flv, 0.4_default)
    call polarization_write (pol, u)
    call polarization_final (pol)
    
    write (u, "(A)") 
    write (u, "(A)")  "* Longitudinally polarized massive vector"
    write (u, "(A)") 
    
    call polarization_init_longitudinal (pol, flv, 1._default)
    call polarization_write (pol, u)
    call polarization_final (pol)
    
    write (u, "(A)") 
    write (u, "(A)")  "* Diagonally polarized massive vector"
    write (u, "(A)")
    
    call polarization_init_diagonal &
         (pol, flv, [0._default, 1._default, 2._default])
    call polarization_write (pol, u)
    call polarization_final (pol)
    
    write (u, "(A)") 
    write (u, "(A)")  "* All polarization directions of a massive vector"
    write (u, "(A)") 

    call polarization_init_generic (pol, flv)
    call polarization_write (pol, u)
    call polarization_final (pol)
    call flv%init (21, model)
    
    write (u, "(A)") 
    write (u, "(A)")  "* Axis polarization (0.2, 0.4, 0.6)"
    write (u, "(A)") 
    
    alpha = [0.2_default, 0.4_default, 0.6_default]
    call polarization_init_axis (pol, flv, alpha)
    call polarization_write (pol, u)
    
    write (u, "(A)")  "   Recovered axis:"
    alpha = polarization_get_axis (pol)
    write (u, "(A)")  "   Angle polarization (0.5, 0.6, -1)"
    r = 0.5_default
    theta = 0.6_default
    phi = -1._default
    call polarization_init_angles (pol, flv, r, theta, phi)
    call polarization_write (pol, u)
    write (u, "(A)")  "   Recovered parameters (r, theta, phi):"
    call polarization_to_angles (pol, r, theta, phi)
    write (u, "(A,3(1x," // FMT_12 // "))")  "     ", r, theta, phi
    call polarization_final (pol)
    
    call model%final ()

    write (u, "(A)")
    write (u, "(A)")  "* Test output end: polarization_1"    
      
  end subroutine polarization_1

\eatline
\nwendcode{}\nwbegindocs{106}\nwdocspar
\subsubsection{Sparse-Matrix type}
Use a sparse density matrix universally as the input for setting up
polarization.
\nwenddocs{}\nwbegincode{107}\moddef{Polarizations: execute tests}\plusendmoddef
  call test (polarization_2, "polarization_2", &
       "matrix polarization setup", &
       u, results)
\nwendcode{}\nwbegincode{108}\moddef{Polarizations: tests}\plusendmoddef
  subroutine polarization_2 (u)
    use os_interface
    integer, intent(in) :: u
    type(model_data_t), target :: model
    type(flavor_t) :: flv
    type(polarization_t) :: pol
    real(default), dimension(3) :: alpha
    type(pmatrix_t) :: pmatrix
    real(default), parameter :: tolerance = 1e-8_default

    write (u, "(A)")  "* Test output: polarization_2"
    write (u, "(A)")  "*   Purpose: matrix polarization setup"
    write (u, "(A)")      
    
    write (u, "(A)")  "* Reading model file"
    write (u, "(A)") 
    
    call model%init_sm_test ()

    write (u, "(A)") "* Unpolarized fermion"
    write (u, "(A)")
    
    call flv%init (1, model)
    call pmatrix%init (2, 0)
    call pmatrix%normalize (flv, 0._default, tolerance)
    call pmatrix%write (u)
    write (u, *) 
    write (u, "(1x,A,L1)")  "polarized = ", pmatrix%is_polarized ()
    write (u, "(1x,A,L1)")  "diagonal = ", pmatrix%is_diagonal ()
    write (u, *) 

    call polarization_init_pmatrix (pol, pmatrix)
    call polarization_write (pol, u)
    call polarization_final (pol)
    
    write (u, "(A)")
    write (u, "(A)")  "* Transversally polarized fermion, phi=0"
    write (u, "(A)")
    
    call pmatrix%init (2, 3)
    call pmatrix%set_entry (1, [-1,-1], (1._default, 0._default))
    call pmatrix%set_entry (2, [+1,+1], (1._default, 0._default))
    call pmatrix%set_entry (3, [-1,+1], (1._default, 0._default))
    call pmatrix%normalize (flv, 1._default, tolerance)
    call pmatrix%write (u)
    write (u, *) 
    write (u, "(1x,A,L1)")  "polarized = ", pmatrix%is_polarized ()
    write (u, "(1x,A,L1)")  "diagonal = ", pmatrix%is_diagonal ()
    write (u, *) 

    call polarization_init_pmatrix (pol, pmatrix)
    call polarization_write (pol, u)
    call polarization_final (pol)
    
    write (u, "(A)") 
    write (u, "(A)")  "* Transversally polarized fermion, phi=0.9, frac=0.8"
    write (u, "(A)")
    
    call pmatrix%init (2, 3)
    call pmatrix%set_entry (1, [-1,-1], (1._default, 0._default))
    call pmatrix%set_entry (2, [+1,+1], (1._default, 0._default))
    call pmatrix%set_entry (3, [-1,+1], exp ((0._default, -0.9_default)))
    call pmatrix%normalize (flv, 0.8_default, tolerance)
    call pmatrix%write (u)
    write (u, *) 

    call polarization_init_pmatrix (pol, pmatrix)
    call polarization_write (pol, u)
    call polarization_final (pol)
    
    write (u, "(A)") 
    write (u, "(A)")  "* Left-handed massive fermion, frac=1"
    write (u, "(A)")
    
    call flv%init (11, model)
    call pmatrix%init (2, 1)
    call pmatrix%set_entry (1, [-1,-1], (1._default, 0._default))
    call pmatrix%normalize (flv, 1._default, tolerance)
    call pmatrix%write (u)
    write (u, *) 
    write (u, "(1x,A,L1)")  "polarized = ", pmatrix%is_polarized ()
    write (u, "(1x,A,L1)")  "diagonal = ", pmatrix%is_diagonal ()
    write (u, *) 

    call polarization_init_pmatrix (pol, pmatrix)
    call polarization_write (pol, u)
    call polarization_final (pol)
    
    write (u, "(A)") 
    write (u, "(A)")  "* Left-handed massive fermion, frac=0.8"
    write (u, "(A)")
    
    call flv%init (11, model)
    call pmatrix%init (2, 1)
    call pmatrix%set_entry (1, [-1,-1], (1._default, 0._default))
    call pmatrix%normalize (flv, 0.8_default, tolerance)
    call pmatrix%write (u)
    write (u, *) 

    call polarization_init_pmatrix (pol, pmatrix)
    call polarization_write (pol, u)
    call polarization_final (pol)
    
    write (u, "(A)") 
    write (u, "(A)")  "* Left-handed massless fermion"
    write (u, "(A)")
    
    call flv%init (12, model)
    call pmatrix%init (2, 0)
    call pmatrix%normalize (flv, 1._default, tolerance)
    call pmatrix%write (u)
    write (u, *) 

    call polarization_init_pmatrix (pol, pmatrix)
    call polarization_write (pol, u)
    call polarization_final (pol)
    
    write (u, "(A)") 
    write (u, "(A)")  "* Right-handed massless fermion, frac=0.5"
    write (u, "(A)")
    
    call flv%init (-12, model)
    call pmatrix%init (2, 1)
    call pmatrix%set_entry (1, [1,1], (1._default, 0._default))
    call pmatrix%normalize (flv, 0.5_default, tolerance)
    call pmatrix%write (u)
    write (u, *) 

    call polarization_init_pmatrix (pol, pmatrix)
    call polarization_write (pol, u)
    call polarization_final (pol)
    
    write (u, "(A)") 
    write (u, "(A)")  "* Circularly polarized gluon, frac=0.3"
    write (u, "(A)") 
    
    call flv%init (21, model)
    call pmatrix%init (2, 1)
    call pmatrix%set_entry (1, [1,1], (1._default, 0._default))
    call pmatrix%normalize (flv, 0.3_default, tolerance)
    call pmatrix%write (u)
    write (u, *) 

    call polarization_init_pmatrix (pol, pmatrix)
    call polarization_write (pol, u)
    call polarization_final (pol)   
    
    write (u, "(A)") 
    write (u, "(A)") "* Circularly polarized massive vector, frac=0.7"
    write (u, "(A)") 
    
    call flv%init (23, model)
    call pmatrix%init (2, 1)
    call pmatrix%set_entry (1, [1,1], (1._default, 0._default))
    call pmatrix%normalize (flv, 0.7_default, tolerance)
    call pmatrix%write (u)
    write (u, *) 

    call polarization_init_pmatrix (pol, pmatrix)
    call polarization_write (pol, u)
    call polarization_final (pol)
    
    write (u, "(A)") 
    write (u, "(A)")  "* Circularly polarized massive vector"
    write (u, "(A)")
    
    call flv%init (23, model)
    call pmatrix%init (2, 1)
    call pmatrix%set_entry (1, [1,1], (1._default, 0._default))
    call pmatrix%normalize (flv, 1._default, tolerance)
    call pmatrix%write (u)
    write (u, *) 

    call polarization_init_pmatrix (pol, pmatrix)
    call polarization_write (pol, u)
    call polarization_final (pol)
    
    write (u, "(A)") 
    write (u, "(A)")  "* Longitudinally polarized massive vector, frac=0.4"
    write (u, "(A)")
    
    call flv%init (23, model)
    call pmatrix%init (2, 1)
    call pmatrix%set_entry (1, [0,0], (1._default, 0._default))
    call pmatrix%normalize (flv, 0.4_default, tolerance)
    call pmatrix%write (u)
    write (u, *) 
    write (u, "(1x,A,L1)")  "polarized = ", pmatrix%is_polarized ()
    write (u, "(1x,A,L1)")  "diagonal = ", pmatrix%is_diagonal ()
    write (u, *) 

    call polarization_init_pmatrix (pol, pmatrix)
    call polarization_write (pol, u)
    call polarization_final (pol)
    
    write (u, "(A)") 
    write (u, "(A)")  "* Longitudinally polarized massive vector"
    write (u, "(A)") 
    
    call flv%init (23, model)
    call pmatrix%init (2, 1)
    call pmatrix%set_entry (1, [0,0], (1._default, 0._default))
    call pmatrix%normalize (flv, 1._default, tolerance)
    call pmatrix%write (u)
    write (u, *) 
    write (u, "(1x,A,L1)")  "polarized = ", pmatrix%is_polarized ()
    write (u, "(1x,A,L1)")  "diagonal = ", pmatrix%is_diagonal ()
    write (u, *) 

    call polarization_init_pmatrix (pol, pmatrix)
    call polarization_write (pol, u)
    call polarization_final (pol)
    
    write (u, "(A)") 
    write (u, "(A)")  "* Axis polarization (0.2, 0.4, 0.6)"
    write (u, "(A)") 
    
    call flv%init (11, model)
    alpha = [0.2_default, 0.4_default, 0.6_default]
    alpha = alpha / sqrt (sum (alpha**2))
    call pmatrix%init (2, 3)
    call pmatrix%set_entry (1, [-1,-1], cmplx (1 - alpha(3), kind=default))
    call pmatrix%set_entry (2, [1,-1], &
         cmplx (alpha(1), -alpha(2), kind=default))
    call pmatrix%set_entry (3, [1,1], cmplx (1 + alpha(3), kind=default))
    call pmatrix%normalize (flv, 1._default, tolerance)
    call pmatrix%write (u)
    write (u, *) 

    call polarization_init_pmatrix (pol, pmatrix)
    call polarization_write (pol, u)
    call polarization_final (pol)
    
    call model%final ()

    write (u, "(A)")
    write (u, "(A)")  "* Test output end: polarization_2"    
      
  end subroutine polarization_2

\eatline
\nwendcode{}\nwbegindocs{109}\nwdocspar
%%%%%%%%%%%%%%%%%%%%%%%%%%%%%%%%%%%%%%%%%%%%%%%%%%%%%%%%%%%%%%%%%%%%%%%% 
\section{Particles}
This module defines the {\tt{}particle{\char95}t} object type, and the methods
and operations that deal with it.
\nwenddocs{}\nwbegincode{110}\moddef{\code{}particles.f90\edoc{}}\endmoddef
\LA{}File header\RA{}

module particles

\LA{}Use kinds\RA{}
\LA{}Use strings\RA{}
  use io_units
  use format_utils, only: write_separator
  use format_utils, only: pac_fmt
  use format_defs, only: FMT_16, FMT_17, FMT_19
  use unit_tests
  use diagnostics
  use lorentz
  use model_data
  use flavors
  use colors
  use helicities
  use quantum_numbers
  use state_matrices
  use interactions
  use evaluators
  use subevents
  use polarizations

\LA{}Standard module head\RA{}

\LA{}Particles: public\RA{}

\LA{}Particles: parameters\RA{}
       
\LA{}Particles: types\RA{}

\LA{}Particles: interfaces\RA{}

contains

\LA{}Particles: procedures\RA{}

\LA{}Particles: tests\RA{}

end module particles
\eatline
\nwendcode{}\nwbegindocs{111}\nwdocspar
\subsection{The particle type}
\subsubsection{Particle status codes}
The overall status codes (incoming/outgoing etc.) are inherited from
the module {\tt{}subevents}.

Polarization status:
\nwenddocs{}\nwbegincode{112}\moddef{Particles: parameters}\endmoddef
  integer, parameter, public :: PRT_UNPOLARIZED = 0
  integer, parameter, public :: PRT_DEFINITE_HELICITY = 1
  integer, parameter, public :: PRT_GENERIC_POLARIZATION = 2

\eatline
\nwendcode{}\nwbegindocs{113}\nwdocspar
\subsubsection{Definition}
The quantum numbers are flavor (from which invariant particle
properties can be derived), color, and polarization.  The particle may
be unpolarized.  In this case, {\tt{}hel} and {\tt{}pol} are unspecified.
If it has a definite helicity, the {\tt{}hel} component is defined.  If
it has a generic polarization, the {\tt{}pol} component is defined.  For
each particle we store the four-momentum and the invariant mass
squared, i.e., the squared norm of the four-momentum.  There is also
an optional list of parent and child particles, for bookkeeping in
physical events.
\nwenddocs{}\nwbegincode{114}\moddef{Particles: public}\endmoddef
  public :: particle_t
\nwendcode{}\nwbegincode{115}\moddef{Particles: types}\endmoddef
  type :: particle_t
     private
     integer :: status = PRT_UNDEFINED
     integer :: polarization = PRT_UNPOLARIZED
     type(flavor_t) :: flv
     type(color_t) :: col
     type(helicity_t) :: hel
     type(polarization_t) :: pol
     type(vector4_t) :: p = vector4_null
     real(default) :: p2 = 0
     integer, dimension(:), allocatable :: parent
     integer, dimension(:), allocatable :: child
   contains
  \LA{}Particles: particle: TBP\RA{}  
  end type particle_t

\eatline
\nwendcode{}\nwbegindocs{116}Copy a particle.  (Deep copy)  This excludes the parent-child
relations.
\nwenddocs{}\nwbegincode{117}\moddef{Particles: partic}\endmoddef
  generic :: init => init_particle
  procedure :: init_particle => particle_init_particle
\nwendcode{}\nwbegincode{118}\moddef{Particles: procedures}\endmoddef
  subroutine particle_init_particle (prt_out, prt_in)
    type(particle_t), intent(in) :: prt_in
    class(particle_t), intent(out) :: prt_out
    prt_out%status = prt_in%status
    prt_out%polarization = prt_in%polarization
    prt_out%flv = prt_in%flv
    prt_out%col = prt_in%col
    prt_out%hel = prt_in%hel
    prt_out%pol = prt_in%pol
    prt_out%p = prt_in%p
    prt_out%p2 = prt_in%p2
  end subroutine particle_init_particle

\eatline
\nwendcode{}\nwbegindocs{119}Initialize a particle using a single-particle state matrix which
determines flavor, color, and polarization.  The state matrix must
have unique flavor and color.  The factorization mode determines
whether the particle is unpolarized, has definite helicity, or generic
polarization.  This mode is translated into the polarization status.
\nwenddocs{}\nwbegincode{120}\moddef{Particles: particle: TBP}\endmoddef
  generic :: init => init_state
  procedure :: init_state => particle_init_state
\nwendcode{}\nwbegincode{121}\moddef{Particles: procedures}\plusendmoddef
  subroutine particle_init_state (prt, state, status, mode)
    class(particle_t), intent(out) :: prt
    type(state_matrix_t), intent(in) :: state
    integer, intent(in) :: status, mode
    type(state_iterator_t) :: it
    prt%status = status
    call it%init (state)
    prt%flv = it%get_flavor (1)
    if (prt%flv%is_radiated ())  prt%status = PRT_BEAM_REMNANT
    prt%col = it%get_color (1)
    select case (mode)
    case (FM_SELECT_HELICITY)
       prt%hel = it%get_helicity (1)
       if (prt%hel%is_defined ()) then
          prt%polarization = PRT_DEFINITE_HELICITY
       end if
    case (FM_FACTOR_HELICITY)
       call polarization_init_state_matrix (prt%pol, state)
       prt%polarization = PRT_GENERIC_POLARIZATION
    end select
  end subroutine particle_init_state

\eatline
\nwendcode{}\nwbegindocs{122}Finalizer.  The polarization component has pointers allocated.
\nwenddocs{}\nwbegincode{123}\moddef{Particles: particle: TBP}\plusendmoddef
  procedure :: final => particle_final
\nwendcode{}\nwbegincode{124}\moddef{Particles: procedures}\plusendmoddef
  subroutine particle_final (prt)
    class(particle_t), intent(inout) :: prt
    call polarization_final (prt%pol)
  end subroutine particle_final

\eatline
\nwendcode{}\nwbegindocs{125}\nwdocspar
\subsubsection{I/O}
\nwenddocs{}\nwbegincode{126}\moddef{Particles: particle: TBP}\plusendmoddef
  procedure :: write => particle_write
\nwendcode{}\nwbegincode{127}\moddef{Particles: procedures}\plusendmoddef
  subroutine particle_write (prt, unit, testflag)
    class(particle_t), intent(in) :: prt
    integer, intent(in), optional :: unit
    logical, intent(in), optional :: testflag
    logical :: pacified
    integer :: u
    real(default) :: pp2
    character(len=7) :: fmt
    pacified = .false.
    if (present (testflag))  pacified = testflag
    call pac_fmt (fmt, FMT_19, FMT_16, testflag)
    u = given_output_unit (unit);  if (u < 0)  return
    select case (prt%status)
    case (PRT_UNDEFINED);    write (u, "(1x, A)", advance="no")  "[-]"
    case (PRT_BEAM);         write (u, "(1x, A)", advance="no")  "[b]"
    case (PRT_INCOMING);     write (u, "(1x, A)", advance="no")  "[i]"
    case (PRT_OUTGOING);     write (u, "(1x, A)", advance="no")  "[o]"
    case (PRT_VIRTUAL);      write (u, "(1x, A)", advance="no")  "[v]"
    case (PRT_RESONANT);     write (u, "(1x, A)", advance="no")  "[r]"
    case (PRT_BEAM_REMNANT); write (u, "(1x, A)", advance="no")  "[x]"
    end select
    write (u, "(1x)", advance="no")
    call prt%flv%write (unit)
    if (prt%col%is_nonzero ()) then
       call color_write (prt%col, unit)
    end if
    select case (prt%polarization)
    case (PRT_DEFINITE_HELICITY)
       call prt%hel%write (unit)
       write (u, *)
    case (PRT_GENERIC_POLARIZATION)
       write (u, *)
       call polarization_write (prt%pol, unit)
    case default
       write (u, *)
    end select    
    call vector4_write (prt%p, unit, testflag = testflag)
    pp2 = prt%p2        
    if (pacified)  call pacify (pp2, tolerance = 1E-10_default) 
    write (u, "(1x,A,1x," // fmt // ")")  "T = ", pp2
    if (allocated (prt%parent)) then
       if (size (prt%parent) /= 0) then
          write (u, "(1x,A,40(1x,I0))")  "Parents: ", prt%parent
       end if
    end if
    if (allocated (prt%child)) then
       if (size (prt%child) /= 0) then
          write (u, "(1x,A,40(1x,I0))")  "Children:", prt%child
       end if
    end if
  end subroutine particle_write

\eatline
\nwendcode{}\nwbegindocs{128}Binary I/O:
\nwenddocs{}\nwbegincode{129}\moddef{Particles: particle: TBP}\plusendmoddef
  procedure :: write_raw => particle_write_raw
  procedure :: read_raw => particle_read_raw  
\nwendcode{}\nwbegincode{130}\moddef{Particles: procedures}\plusendmoddef
  subroutine particle_write_raw (prt, u)
    class(particle_t), intent(in) :: prt
    integer, intent(in) :: u
    write (u) prt%status, prt%polarization
    call prt%flv%write_raw (u)
    call prt%col%write_raw (u)
    select case (prt%polarization)
    case (PRT_DEFINITE_HELICITY)
       call prt%hel%write_raw (u)
    case (PRT_GENERIC_POLARIZATION)
       call polarization_write_raw (prt%pol, u)
    end select
    call vector4_write_raw (prt%p, u)
    write (u) prt%p2
    write (u) allocated (prt%parent)
    if (allocated (prt%parent)) then
       write (u) size (prt%parent)
       write (u) prt%parent
    end if
    write (u) allocated (prt%child)
    if (allocated (prt%child)) then
       write (u) size (prt%child)
       write (u) prt%child
    end if
  end subroutine particle_write_raw

  subroutine particle_read_raw (prt, u, iostat)
    class(particle_t), intent(out) :: prt
    integer, intent(in) :: u
    integer, intent(out) :: iostat
    logical :: allocated_parent, allocated_child
    integer :: size_parent, size_child
    read (u, iostat=iostat) prt%status, prt%polarization
    call prt%flv%read_raw (u, iostat=iostat)
    call prt%col%read_raw (u, iostat=iostat)
    select case (prt%polarization)
    case (PRT_DEFINITE_HELICITY)
       call prt%hel%read_raw (u, iostat=iostat)
    case (PRT_GENERIC_POLARIZATION)
       call polarization_read_raw (prt%pol, u, iostat=iostat)
    end select
    call vector4_read_raw (prt%p, u, iostat=iostat)
    read (u, iostat=iostat) prt%p2
    read (u, iostat=iostat) allocated_parent
    if (allocated_parent) then
       read (u, iostat=iostat) size_parent
       allocate (prt%parent (size_parent))
       read (u, iostat=iostat) prt%parent
    end if
    read (u, iostat=iostat) allocated_child
    if (allocated_child) then
       read (u, iostat=iostat) size_child
       allocate (prt%child (size_child))
       read (u, iostat=iostat) prt%child
    end if
  end subroutine particle_read_raw

\eatline
\nwendcode{}\nwbegindocs{131}\nwdocspar
\subsubsection{Setting contents}
Reset the status code.  Where applicable, set $p^2$ assuming that the
particle is on-shell.
\nwenddocs{}\nwbegincode{132}\moddef{Particles: particle: TBP}\plusendmoddef
  procedure :: reset_status => particle_reset_status
\nwendcode{}\nwbegincode{133}\moddef{Particles: procedures}\plusendmoddef
  elemental subroutine particle_reset_status (prt, status)
    class(particle_t), intent(inout) :: prt
    integer, intent(in) :: status
    prt%status = status
    select case (status)
    case (PRT_BEAM, PRT_INCOMING, PRT_OUTGOING)
       prt%p2 = prt%flv%get_mass () ** 2
    end select
  end subroutine particle_reset_status

\eatline
\nwendcode{}\nwbegindocs{134}The color can be given explicitly.
\nwenddocs{}\nwbegincode{135}\moddef{Particles: particle: TBP}\plusendmoddef
  procedure :: set_color => particle_set_color
\nwendcode{}\nwbegincode{136}\moddef{Particles: procedures}\plusendmoddef
  elemental subroutine particle_set_color (prt, col)
    class(particle_t), intent(inout) :: prt
    type(color_t), intent(in) :: col
    prt%col = col
  end subroutine particle_set_color

\eatline
\nwendcode{}\nwbegindocs{137}The flavor can be given explicitly.
\nwenddocs{}\nwbegincode{138}\moddef{Particles: particle: TBP}\plusendmoddef
  procedure :: set_flavor => particle_set_flavor
\nwendcode{}\nwbegincode{139}\moddef{Particles: procedures}\plusendmoddef
  subroutine particle_set_flavor (prt, flv)
    class(particle_t), intent(inout) :: prt
    type(flavor_t), intent(in) :: flv
    prt%flv = flv
  end subroutine particle_set_flavor

\eatline
\nwendcode{}\nwbegindocs{140}As can the helicity.
\nwenddocs{}\nwbegincode{141}\moddef{Particles: particle: TBP}\plusendmoddef
  procedure :: set_helicity => particle_set_helicity
\nwendcode{}\nwbegincode{142}\moddef{Particles: procedures}\plusendmoddef
  subroutine particle_set_helicity (prt, hel)
    class(particle_t), intent(inout) :: prt
    type(helicity_t), intent(in) :: hel
    prt%hel = hel
  end subroutine particle_set_helicity

\eatline
\nwendcode{}\nwbegindocs{143}And the polarization.
\nwenddocs{}\nwbegincode{144}\moddef{Particles: particle: TBP}\plusendmoddef
  procedure :: set_pol => particle_set_pol
\nwendcode{}\nwbegincode{145}\moddef{Particles: procedures}\plusendmoddef
  subroutine particle_set_pol (prt, pol)
    class(particle_t), intent(inout) :: prt
    type(polarization_t), intent(in) :: pol
    prt%pol = pol
  end subroutine particle_set_pol

\eatline
\nwendcode{}\nwbegindocs{146}Manually set the model for the particle flavor.  This is required, e.g., if
the particle has been read from file.
\nwenddocs{}\nwbegincode{147}\moddef{Particles: particle: TBP}\plusendmoddef
  procedure :: set_model => particle_set_model
\nwendcode{}\nwbegincode{148}\moddef{Particles: procedures}\plusendmoddef
  subroutine particle_set_model (prt, model)
    class(particle_t), intent(inout) :: prt
    class(model_data_t), intent(in), target :: model
    call prt%flv%set_model (model)
  end subroutine particle_set_model
  
\eatline
\nwendcode{}\nwbegindocs{149}The momentum is set independent of the quantum numbers.
\nwenddocs{}\nwbegincode{150}\moddef{Particles: particle: TBP}\plusendmoddef
  procedure :: set_momentum => particle_set_momentum
\nwendcode{}\nwbegincode{151}\moddef{Particles: procedures}\plusendmoddef
  elemental subroutine particle_set_momentum (prt, p, p2, on_shell)
    class(particle_t), intent(inout) :: prt
    type(vector4_t), intent(in) :: p
    real(default), intent(in), optional :: p2
    logical, intent(in), optional :: on_shell
    prt%p = p
    if (present (on_shell)) then
       if (on_shell) then
          if (prt%flv%is_associated ()) then
             prt%p2 = prt%flv%get_mass () ** 2
             return
          end if
       end if
    end if
    if (present (p2)) then
       prt%p2 = p2
    else
       prt%p2 = p ** 2
    end if
  end subroutine particle_set_momentum

\eatline
\nwendcode{}\nwbegindocs{152}Set resonance information.  This should be done after momentum
assignment, because we need to know wheter the particle is spacelike
or timelike.  The resonance flag is defined only for virtual
particles. 
\nwenddocs{}\nwbegincode{153}\moddef{Particle: particle: TBP}\endmoddef
  procedure :: set_resonance_flag => particle_set_resonance_flag
\nwendcode{}\nwbegincode{154}\moddef{Particles: procedures}\plusendmoddef
  elemental subroutine particle_set_resonance_flag (prt, resonant)
    class(particle_t), intent(inout) :: prt
    logical, intent(in) :: resonant
    select case (prt%status)
    case (PRT_VIRTUAL)
       if (resonant)  prt%status = PRT_RESONANT
    end select
  end subroutine particle_set_resonance_flag

\eatline
\nwendcode{}\nwbegindocs{155}Set children and parents information.
\nwenddocs{}\nwbegincode{156}\moddef{Particles: particle: TBP}\plusendmoddef
  procedure :: set_children => particle_set_children
  procedure :: set_parents => particle_set_parents
\nwendcode{}\nwbegincode{157}\moddef{Particles: procedures}\plusendmoddef
  subroutine particle_set_children (prt, idx)
    class(particle_t), intent(inout) :: prt
    integer, dimension(:), intent(in) :: idx
    if (allocated (prt%child))  deallocate (prt%child)
    allocate (prt%child (count (idx /= 0)))
    prt%child = pack (idx, idx /= 0)
  end subroutine particle_set_children

  subroutine particle_set_parents (prt, idx)
    class(particle_t), intent(inout) :: prt
    integer, dimension(:), intent(in) :: idx
    if (allocated (prt%parent))  deallocate (prt%parent) 
    allocate (prt%parent (count (idx /= 0)))
    prt%parent = pack (idx, idx /= 0)
  end subroutine particle_set_parents

\eatline
\nwendcode{}\nwbegindocs{158}\nwdocspar
\nwenddocs{}\nwbegincode{159}\moddef{Particles: particle: TBP}\plusendmoddef
  procedure :: add_child => particle_add_child
\nwendcode{}\nwbegincode{160}\moddef{Particles: procedures}\plusendmoddef
  subroutine particle_add_child (prt, new_child)
    class(particle_t), intent(inout) :: prt
    integer, intent(in) :: new_child
    integer, dimension(:), allocatable :: idx
    integer :: n
    n = prt%get_n_children()
    if (n == 0) then
       call prt%set_children ([new_child])
    else
       allocate (idx (1:n+1))
       idx(1:n) = prt%get_children ()
       idx(n+1) = new_child
       call prt%set_children (idx)
    end if
  end subroutine particle_add_child
\nwendcode{}\nwbegindocs{161}\nwdocspar
\nwenddocs{}\nwbegincode{162}\moddef{Particles: particle: TBP}\plusendmoddef
  procedure :: set_status => particle_set_status
\nwendcode{}\nwbegincode{163}\moddef{Particles: procedures}\plusendmoddef
  elemental subroutine particle_set_status (prt, status)
    class(particle_t), intent(inout) :: prt
    integer, intent(in) :: status
    prt%status = status
  end subroutine particle_set_status

\eatline
\nwendcode{}\nwbegindocs{164}\nwdocspar
\nwenddocs{}\nwbegincode{165}\moddef{Particles: particle: TBP}\plusendmoddef
  procedure :: set_polarization => particle_set_polarization
\nwendcode{}\nwbegincode{166}\moddef{Particles: procedures}\plusendmoddef
  subroutine particle_set_polarization (prt, polarization)
    class(particle_t), intent(inout) :: prt
    integer, intent(in) :: polarization
    prt%polarization = polarization
  end subroutine particle_set_polarization

\eatline
\nwendcode{}\nwbegindocs{167}\nwdocspar
\subsubsection{Accessing contents}
The status code.
\nwenddocs{}\nwbegincode{168}\moddef{Particles: particle: TBP}\plusendmoddef
  procedure :: get_status => particle_get_status
\nwendcode{}\nwbegincode{169}\moddef{Particles: procedures}\plusendmoddef
  elemental function particle_get_status (prt) result (status)
    integer :: status
    class(particle_t), intent(in) :: prt
    status = prt%status
  end function particle_get_status

\eatline
\nwendcode{}\nwbegindocs{170}Return true if the status is either {\tt{}INCOMING},
{\tt{}OUTGOING} or {\tt{}RESONANT}. {\tt{}BEAM} is kept, if 
{\tt{}keep{\char95}beams} is set true.
\nwenddocs{}\nwbegincode{171}\moddef{Particles: particle: TBP}\plusendmoddef
  procedure :: is_real => particle_is_real
\nwendcode{}\nwbegincode{172}\moddef{Particles: procedures}\plusendmoddef
  elemental function particle_is_real (prt, keep_beams) result (flag)
    logical :: flag, kb
    class(particle_t), intent(in) :: prt
    logical, intent(in), optional :: keep_beams
    kb = .false.
    if (present (keep_beams)) kb = keep_beams
    select case (prt%status)
    case (PRT_INCOMING, PRT_OUTGOING, PRT_RESONANT)
       flag = .true.
    case (PRT_BEAM)
       flag = kb 
    case default
       flag = .false.
    end select
  end function particle_is_real

\eatline
\nwendcode{}\nwbegindocs{173}Polarization status.
\nwenddocs{}\nwbegincode{174}\moddef{Particles: particle: TBP}\plusendmoddef
  procedure :: get_polarization_status => particle_get_polarization_status
\nwendcode{}\nwbegincode{175}\moddef{Particles: procedures}\plusendmoddef
  elemental function particle_get_polarization_status (prt) result (status)
    integer :: status
    class(particle_t), intent(in) :: prt
    status = prt%polarization
  end function particle_get_polarization_status

\eatline
\nwendcode{}\nwbegindocs{176}Return the PDG code from the flavor component directly.
\nwenddocs{}\nwbegincode{177}\moddef{Particles: particle: TBP}\plusendmoddef
  procedure :: get_pdg => particle_get_pdg
\nwendcode{}\nwbegincode{178}\moddef{Particles: procedures}\plusendmoddef
  elemental function particle_get_pdg (prt) result (pdg)
    integer :: pdg
    class(particle_t), intent(in) :: prt
    pdg = prt%flv%get_pdg ()
  end function particle_get_pdg

\eatline
\nwendcode{}\nwbegindocs{179}Return the color and anticolor quantum numbers.
\nwenddocs{}\nwbegincode{180}\moddef{Particles: particle: TBP}\plusendmoddef
  procedure :: get_color => particle_get_color
\nwendcode{}\nwbegincode{181}\moddef{Particles: procedures}\plusendmoddef
  pure function particle_get_color (prt) result (col)
    integer, dimension(2) :: col
    class(particle_t), intent(in) :: prt
    col(1) = prt%col%get_col ()
    col(2) = prt%col%get_acl ()
  end function particle_get_color
  
\eatline
\nwendcode{}\nwbegindocs{182}Return the polarization density matrix (as a shallow copy).
\nwenddocs{}\nwbegincode{183}\moddef{Particles: particle: TBP}\plusendmoddef
  procedure :: get_polarization => particle_get_polarization
\nwendcode{}\nwbegincode{184}\moddef{Particles: procedures}\plusendmoddef
  function particle_get_polarization (prt) result (pol)
    class(particle_t), intent(in) :: prt
    class(polarization_t), allocatable :: pol
    pol = prt%pol
  end function particle_get_polarization

\eatline
\nwendcode{}\nwbegindocs{185}Return the flavor, color and helicity.
\nwenddocs{}\nwbegincode{186}\moddef{Particles: particle: TBP}\plusendmoddef
  procedure :: get_flv => particle_get_flv
  procedure :: get_col => particle_get_col 
  procedure :: get_hel => particle_get_hel
\nwendcode{}\nwbegincode{187}\moddef{Particles: procedures}\plusendmoddef
  function particle_get_flv (prt) result (flv)
    class(particle_t), intent(in) :: prt
    type(flavor_t) :: flv
    flv = prt%flv
  end function particle_get_flv 
  
  function particle_get_col (prt) result (col)
    class(particle_t), intent(in) :: prt
    type(color_t) :: col
    col = prt%col
  end function particle_get_col
  
  function particle_get_hel (prt) result (hel)
    class(particle_t), intent(in) :: prt
    type(helicity_t) :: hel
    hel = prt%hel
  end function particle_get_hel
  
\eatline
\nwendcode{}\nwbegindocs{188}Return the helicity (if defined and diagonal).
\nwenddocs{}\nwbegincode{189}\moddef{Particles: particle: TBP}\plusendmoddef
  procedure :: get_helicity => particle_get_helicity
\nwendcode{}\nwbegincode{190}\moddef{Particles: procedures}\plusendmoddef
  elemental function particle_get_helicity (prt) result (hel)
    integer :: hel
    integer, dimension(2) :: hel_arr
    class(particle_t), intent(in) :: prt    
    hel = 0
    if (prt%hel%is_defined () .and. prt%hel%is_diagonal ()) then
       hel_arr = prt%hel%to_pair ()
       hel = hel_arr (1)
    end if
  end function particle_get_helicity  
  
\eatline
\nwendcode{}\nwbegindocs{191}Return the number of children/parents
\nwenddocs{}\nwbegincode{192}\moddef{Particles: particle: TBP}\plusendmoddef
  procedure :: get_n_parents => particle_get_n_parents
  procedure :: get_n_children => particle_get_n_children
\nwendcode{}\nwbegincode{193}\moddef{Particles: procedures}\plusendmoddef
  elemental function particle_get_n_parents (prt) result (n)
    integer :: n
    class(particle_t), intent(in) :: prt
    if (allocated (prt%parent)) then
       n = size (prt%parent)
    else
       n = 0
    end if
  end function particle_get_n_parents
    
  elemental function particle_get_n_children (prt) result (n)
    integer :: n
    class(particle_t), intent(in) :: prt
    if (allocated (prt%child)) then
       n = size (prt%child)
    else
       n = 0
    end if
  end function particle_get_n_children
    
\eatline
\nwendcode{}\nwbegindocs{194}Return the array of parents/children.
\nwenddocs{}\nwbegincode{195}\moddef{Particles: particle: TBP}\plusendmoddef
  procedure :: get_parents => particle_get_parents
  procedure :: get_children => particle_get_children
\nwendcode{}\nwbegincode{196}\moddef{Particles: procedures}\plusendmoddef
  function particle_get_parents (prt) result (parent)
    class(particle_t), intent(in) :: prt
    integer, dimension(:), allocatable :: parent
    if (allocated (prt%parent)) then
       allocate (parent (size (prt%parent)))
       parent = prt%parent
    else
       allocate (parent (0))
    end if
  end function particle_get_parents

  function particle_get_children (prt) result (child)
    class(particle_t), intent(in) :: prt
    integer, dimension(:), allocatable :: child
    if (allocated (prt%child)) then
       allocate (child (size (prt%child)))
       child = prt%child
    else
       allocate (child (0))
    end if
  end function particle_get_children

\eatline
\nwendcode{}\nwbegindocs{197}Return momentum and momentum squared.
\nwenddocs{}\nwbegincode{198}\moddef{Particles: particle: TBP}\plusendmoddef
  procedure :: get_momentum => particle_get_momentum
  procedure :: get_p2 => particle_get_p2
\nwendcode{}\nwbegincode{199}\moddef{Particles: procedures}\plusendmoddef
  elemental function particle_get_momentum (prt) result (p)
    type(vector4_t) :: p
    class(particle_t), intent(in) :: prt
    p = prt%p
  end function particle_get_momentum

  elemental function particle_get_p2 (prt) result (p2)
    real(default) :: p2
    class(particle_t), intent(in) :: prt
    p2 = prt%p2
  end function particle_get_p2

\eatline
\nwendcode{}\nwbegindocs{200}\nwdocspar
\subsection{Particle sets}
A particle set is what is usually called an event: an array of
particles.  The individual particle entries carry momentum, quantum
numbers, polarization, and optionally connections.  There is (also
optionally) a correlated state-density matrix that maintains spin
correlations that are lost in the individual particle entries.
\nwenddocs{}\nwbegincode{201}\moddef{Particles: public}\plusendmoddef
  public :: particle_set_t
\nwendcode{}\nwbegincode{202}\moddef{Particles: types}\plusendmoddef
  type :: particle_set_t
     integer :: n_beam = 0
     integer :: n_in  = 0
     integer :: n_vir = 0
     integer :: n_out = 0
     integer :: n_tot = 0
     type(particle_t), dimension(:), allocatable :: prt
     type(state_matrix_t) :: correlated_state
   contains
   \LA{}Particles: particle set: TBP\RA{}
  end type particle_set_t

\eatline
\nwendcode{}\nwbegindocs{203}A particle set can be initialized from an interaction or from a
HepMC event record.
\nwenddocs{}\nwbegincode{204}\moddef{Particles: particle set: TBP}\endmoddef
  generic :: init => init_interaction
  procedure :: init_interaction => particle_set_init_interaction
\eatline
\nwendcode{}\nwbegindocs{205}When a particle set is initialized from a given interaction, we have
to determine the branch within the original state matrix that fixes
the particle quantum numbers.  This is done with the appropriate
probabilities, based on a random number {\tt{}x}.  The {\tt{}mode}
determines whether the individual particles become unpolarized, or
take a definite (diagonal) helicity, or acquire single-particle
polarization matrices.  The flag {\tt{}keep{\char95}correlations} tells whether
the spin-correlation matrix is to be calculated and stored in addition
to the particles.  The flag {\tt{}keep{\char95}virtual} tells whether virtual
particles should be dropped.  Note that if virtual particles are
dropped, the spin-correlation matrix makes no sense, and parent-child
relations are not set.

For a correct disentangling of color and flavor (in the presence of
helicity), we consider two interactions.  {\tt{}int} has no color
information, and is used to select a flavor state.  Consequently, we
trace over helicities here.  {\tt{}int{\char95}flows} contains color-flow and
potentially helicity information, but is useful only after the flavor
combination has been chosen.  So this interaction is used to select
helicity and color, but restricted to the selected flavor combination.

{\tt{}int} and {\tt{}int{\char95}flows} may be identical if there is only a single
(or no) color flow.  If there is just a single flavor combination,
{\tt{}x(1)} can be set to zero.

The current algorithm of evaluator convolution requires that the beam
particles are assumed outgoing (in the beam interaction) and become
virtual in all derived interactions.  In the particle set they should
be re-identified as incoming.  The optional integer {\tt{}n{\char95}incoming}
can be used to perform this correction.
    
The flag {\tt{}is{\char95}valid} is false if factorization of the state is not
possible, in particular if the squared matrix element is zero.
\nwenddocs{}\nwbegincode{206}\moddef{Particles: procedures}\plusendmoddef
  subroutine particle_set_init_interaction &
       (particle_set, is_valid, int, int_flows, mode, x, &
        keep_correlations, keep_virtual, n_incoming)
    class(particle_set_t), intent(out) :: particle_set
    logical, intent(out) :: is_valid
    type(interaction_t), intent(in), target :: int, int_flows
    integer, intent(in) :: mode
    real(default), dimension(2), intent(in) :: x
    logical, intent(in) :: keep_correlations, keep_virtual
    integer, intent(in), optional :: n_incoming
    type(state_matrix_t), dimension(:), allocatable, target :: flavor_state
    type(state_matrix_t), dimension(:), allocatable, target :: single_state
    integer :: n_in, n_vir, n_out, n_tot
    type(quantum_numbers_t), dimension(:,:), allocatable :: qn
    logical :: ok
    integer :: i, j
    if (present (n_incoming)) then
       n_in  = n_incoming
       n_vir = interaction_get_n_vir (int) - n_incoming
    else
       n_in  = interaction_get_n_in  (int)
       n_vir = interaction_get_n_vir (int)
    end if
    n_out = interaction_get_n_out (int)
    n_tot = interaction_get_n_tot (int)
    particle_set%n_in  = n_in
    particle_set%n_out = n_out
    if (keep_virtual) then
       particle_set%n_vir = n_vir
       particle_set%n_tot = n_tot
    else
       particle_set%n_vir = 0
       particle_set%n_tot = n_in + n_out
    end if
    call interaction_factorize &
         (int, FM_IGNORE_HELICITY, x(1), is_valid, flavor_state)
    allocate (qn (n_tot,1))
    do i = 1, n_tot
       qn(i,:) = flavor_state(i)%get_quantum_numbers (1)
    end do
    if (keep_correlations .and. keep_virtual) then
       call interaction_factorize (int_flows, mode, x(2), ok, &
            single_state, particle_set%correlated_state, qn(:,1))
    else
       call interaction_factorize (int_flows, mode, x(2), ok, &
            single_state, qn_in=qn(:,1))
    end if
    is_valid = is_valid .and. ok
    allocate (particle_set%prt (particle_set%n_tot))
    j = 1
    do i = 1, n_tot
       if (i <= n_in) then
          call particle_set%prt(j)%init (single_state(i), PRT_INCOMING, mode)
          call particle_set%prt(j)%set_momentum &
               (interaction_get_momentum (int, i))
       else if (i <= n_in + n_vir) then
          if (.not. keep_virtual)  cycle
          call particle_set%prt(j)%init &
               (single_state(i), PRT_VIRTUAL, mode) 
          call particle_set%prt(j)%set_momentum &
               (interaction_get_momentum (int, i))
       else
          call particle_set%prt(j)%init (single_state(i), PRT_OUTGOING, mode)
          call particle_set%prt(j)%set_momentum &
               (interaction_get_momentum (int, i), on_shell = .true.)
       end if
       if (keep_virtual) then
          call particle_set%prt(j)%set_children &
               (interaction_get_children (int, i))
          call particle_set%prt(j)%set_parents &
               (interaction_get_parents (int, i))
       end if
       j = j + 1
    end do
    if (keep_virtual) then
       call particle_set_resonance_flag &
            (particle_set%prt, interaction_get_resonance_flags (int))
    end if
    do i = i, size(flavor_state)
       call flavor_state(i)%final ()
    end do
    do i = i, size(single_state)
       call single_state(i)%final ()
    end do
  end subroutine particle_set_init_interaction

\eatline
\nwendcode{}\nwbegindocs{207}Manually set the model for the stored particles.
\nwenddocs{}\nwbegincode{208}\moddef{Particles: particle set: TBP}\plusendmoddef
  procedure :: set_model => particle_set_set_model
\nwendcode{}\nwbegincode{209}\moddef{Particles: procedures}\plusendmoddef
  subroutine particle_set_set_model (particle_set, model)
    class(particle_set_t), intent(inout) :: particle_set
    class(model_data_t), intent(in), target :: model
    integer :: i
    do i = 1, particle_set%n_tot
       call particle_set%prt(i)%set_model (model)
    end do
    call particle_set%correlated_state%set_model (model)
  end subroutine particle_set_set_model
    
\eatline
\nwendcode{}\nwbegindocs{210}Pointer components are hidden inside the particle polarization, and
in the correlated state matrix.
\nwenddocs{}\nwbegincode{211}\moddef{Particles: particle set: TBP}\plusendmoddef
  procedure :: final => particle_set_final
\nwendcode{}\nwbegincode{212}\moddef{Particles: procedures}\plusendmoddef
  subroutine particle_set_final (particle_set)
    class(particle_set_t), intent(inout) :: particle_set
    integer :: i
    if (allocated (particle_set%prt)) then
       do i = 1, size(particle_set%prt)
          call particle_set%prt(i)%final ()
       end do
       deallocate (particle_set%prt)
    end if
    call particle_set%correlated_state%final ()
  end subroutine particle_set_final

\eatline
\nwendcode{}\nwbegindocs{213}Output (default format)
\nwenddocs{}\nwbegincode{214}\moddef{Particles: particle set: TBP}\plusendmoddef
  procedure :: write => particle_set_write
\nwendcode{}\nwbegincode{215}\moddef{Particles: procedures}\plusendmoddef
  subroutine particle_set_write (particle_set, unit, testflag)
    class(particle_set_t), intent(in) :: particle_set
    integer, intent(in), optional :: unit
    logical, intent(in), optional :: testflag
    integer :: u, i
    u = given_output_unit (unit);  if (u < 0)  return
    write (u, "(1x,A)") "Particle set:"
    call write_separator (u)
    if (particle_set%n_tot /= 0) then
       do i = 1, particle_set%n_tot
          write (u, "(1x,A,1x,I0)", advance="no") "Particle", i
          call particle_set%prt(i)%write (u, testflag)
       end do
       if (particle_set%correlated_state%is_defined ()) then
          call write_separator (u)
          write (u, *) "Correlated state density matrix:"
          call particle_set%correlated_state%write (u)
       end if
    else
       write (u, "(3x,A)") "[empty]"
    end if
  end subroutine particle_set_write

\eatline
\nwendcode{}\nwbegindocs{216}\nwdocspar
\subsection{I/O formats}
Here, we define input/output of particle sets in various formats.
This is the right place since particle sets contain most of the event
information.

All write/read routines take as first argument the object, as second
argument the I/O unit which in this case is a mandatory argument.
Then follow further event data.

\subsubsection{Internal binary format}
This format is supposed to contain the complete information, so
the particle data set can be fully reconstructed.  The exception is
the model part of the particle flavors; this is unassigned for the
flavor values read from file.
\nwenddocs{}\nwbegincode{217}\moddef{Particles: particle set: TBP}\plusendmoddef
  procedure :: write_raw => particle_set_write_raw 
  procedure :: read_raw => particle_set_read_raw 
\nwendcode{}\nwbegincode{218}\moddef{Particles: procedures}\plusendmoddef
  subroutine particle_set_write_raw (particle_set, u)
    class(particle_set_t), intent(in) :: particle_set
    integer, intent(in) :: u
    integer :: i
    write (u) &
         particle_set%n_beam, particle_set%n_in, &
         particle_set%n_vir, particle_set%n_out
    write (u) particle_set%n_tot
    do i = 1, particle_set%n_tot
       call particle_set%prt(i)%write_raw (u)
    end do
    call particle_set%correlated_state%write_raw (u)
  end subroutine particle_set_write_raw

  subroutine particle_set_read_raw (particle_set, u, iostat)
    class(particle_set_t), intent(out) :: particle_set
    integer, intent(in) :: u
    integer, intent(out) :: iostat
    integer :: i
    read (u, iostat=iostat) &
         particle_set%n_beam, particle_set%n_in, &
         particle_set%n_vir, particle_set%n_out
    read (u, iostat=iostat) particle_set%n_tot
    allocate (particle_set%prt (particle_set%n_tot))
    do i = 1, size (particle_set%prt)
       call particle_set%prt(i)%read_raw (u, iostat=iostat)
    end do
    call particle_set%correlated_state%read_raw (u, iostat=iostat)
  end subroutine particle_set_read_raw

\eatline
\nwendcode{}\nwbegindocs{219}\nwdocspar
\subsubsection{Get contents}
Find parents/children of a particular particle recursively; the
search terminates if a parent/child has status {\tt{}BEAM}, {\tt{}INCOMING},
{\tt{}OUTGOING} or {\tt{}RESONANT}.
\nwenddocs{}\nwbegincode{220}\moddef{Particles: particle set: TBP}\plusendmoddef
  procedure :: get_real_parents => particle_set_get_real_parents
  procedure :: get_real_children => particle_set_get_real_children
\nwendcode{}\nwbegincode{221}\moddef{Particles: procedures}\plusendmoddef
  function particle_set_get_real_parents (pset, i, keep_beams) result (parent)
    integer, dimension(:), allocatable :: parent
    class(particle_set_t), intent(in) :: pset
    integer, intent(in) :: i
    logical, intent(in), optional :: keep_beams
    logical, dimension(:), allocatable :: is_real
    logical, dimension(:), allocatable :: is_parent, is_real_parent
    logical :: kb
    integer :: j, k
    kb = .false.
    if (present (keep_beams)) kb = keep_beams
    allocate (is_real (pset%n_tot))
    is_real = pset%prt%is_real (kb)
    allocate (is_parent (pset%n_tot), is_real_parent (pset%n_tot))
    is_real_parent = .false.
    is_parent = .false.
    is_parent(pset%prt(i)%get_parents()) = .true.
    do while (any (is_parent))
       where (is_real .and. is_parent)
          is_real_parent = .true.
          is_parent = .false.
       end where
       mark_next_parent: do j = size (is_parent), 1, -1
          if (is_parent(j)) then
             is_parent(pset%prt(j)%get_parents()) = .true.
             is_parent(j) = .false.
             exit mark_next_parent
          end if
       end do mark_next_parent
    end do
    allocate (parent (count (is_real_parent)))
    j = 0
    do k = 1, size (is_parent)
       if (is_real_parent(k)) then
          j = j + 1
          parent(j) = k
       end if
    end do
  end function particle_set_get_real_parents

  function particle_set_get_real_children (pset, i, keep_beams) result (child)
    integer, dimension(:), allocatable :: child
    class(particle_set_t), intent(in) :: pset
    integer, intent(in) :: i
    logical, dimension(:), allocatable :: is_real
    logical, dimension(:), allocatable :: is_child, is_real_child
    logical, intent(in), optional :: keep_beams
    integer :: j, k
    logical :: kb
    kb = .false.
    if (present (keep_beams)) kb = keep_beams
    allocate (is_real (pset%n_tot))
    is_real = pset%prt%is_real (kb)
    allocate (is_child (pset%n_tot), is_real_child (pset%n_tot))
    is_real_child = .false.
    is_child = .false.
    is_child(pset%prt(i)%get_children()) = .true.
    do while (any (is_child))
       where (is_real .and. is_child)
          is_real_child = .true.
          is_child = .false.
       end where
       mark_next_child: do j = 1, size (is_child)
          if (is_child(j)) then
             is_child(pset%prt(j)%get_children()) = .true.
             is_child(j) = .false.
             exit mark_next_child
          end if
       end do mark_next_child
    end do
    allocate (child (count (is_real_child)))
    j = 0
    do k = 1, size (is_child)
       if (is_real_child(k)) then
          j = j + 1
          child(j) = k
       end if
    end do
  end function particle_set_get_real_children

\eatline
\eatline
\nwendcode{}\nwbegindocs{222}Get the {\tt{}n{\char95}tot}, {\tt{}n{\char95}in}, and {\tt{}n{\char95}out} values out of the
particle set.
\nwenddocs{}\nwbegincode{223}\moddef{Particles: particle set: TBP}\plusendmoddef
  procedure :: get_n_beam => particle_set_get_n_beam
  procedure :: get_n_in => particle_set_get_n_in
  procedure :: get_n_vir => particle_set_get_n_vir
  procedure :: get_n_out => particle_set_get_n_out
  procedure :: get_n_tot => particle_set_get_n_tot
\nwendcode{}\nwbegincode{224}\moddef{Particles: procedures}\plusendmoddef
  function particle_set_get_n_beam (pset) result (n_beam)
     class(particle_set_t), intent(in) :: pset
     integer :: n_beam
     n_beam = pset%n_beam
  end function particle_set_get_n_beam

  function particle_set_get_n_in (pset) result (n_in)
     class(particle_set_t), intent(in) :: pset
     integer :: n_in
     n_in = pset%n_in
  end function particle_set_get_n_in

  function particle_set_get_n_vir (pset) result (n_vir)
     class(particle_set_t), intent(in) :: pset
     integer :: n_vir
     n_vir = pset%n_vir
   end function particle_set_get_n_vir

  function particle_set_get_n_out (pset) result (n_out)
     class(particle_set_t), intent(in) :: pset
     integer :: n_out
     n_out = pset%n_out
  end function particle_set_get_n_out
  
  function particle_set_get_n_tot (pset) result (n_tot)
     class(particle_set_t), intent(in) :: pset
     integer :: n_tot
     n_tot = pset%n_tot
  end function particle_set_get_n_tot

  function particle_set_get_n_rad (pset) result (n_rad)
    class(particle_set_t), intent(in) :: pset
    integer :: n_rad
    n_rad = count (pset%prt%get_status () == PRT_BEAM_REMNANT)
  end function particle_set_get_n_rad

\eatline
\eatline
\eatline
\eatline
\eatline
\eatline
\nwendcode{}\nwbegindocs{225}Return a pointer to the particle corresponding to the number
\nwenddocs{}\nwbegincode{226}\moddef{Particles: particle set: TBP}\plusendmoddef
  procedure :: get_particle => particle_set_get_particle
\nwendcode{}\nwbegincode{227}\moddef{Particles: procedures}\plusendmoddef
  function particle_set_get_particle (pset, index) result (particle)
    class(particle_set_t), intent(in) :: pset
    integer, intent(in) :: index
    type(particle_t) :: particle
    particle = pset%prt(index)
  end function particle_set_get_particle

\eatline
\nwendcode{}\nwbegindocs{228}\nwdocspar
\subsubsection{Tools}
Given a subevent, reset status codes.  If the new status is beam,
incoming, or outgoing, we also make sure that the stored $p^2$ value
is equal to the on-shell mass squared.
\nwenddocs{}\nwbegincode{229}\moddef{Particles: particle set: TBP}\plusendmoddef
  procedure :: reset_status => particle_set_reset_status
\nwendcode{}\nwbegincode{230}\moddef{Particles: procedures}\plusendmoddef
  subroutine particle_set_reset_status (particle_set, index, status)
    class(particle_set_t), intent(inout) :: particle_set
    integer, dimension(:), intent(in) :: index
    integer, intent(in) :: status
    integer :: i
    if (allocated (particle_set%prt)) then
       do i = 1, size (index)
          call particle_set%prt(index(i))%reset_status (status)
       end do
    end if
    particle_set%n_beam  = &
         count (particle_set%prt%get_status () == PRT_BEAM)
    particle_set%n_in  = &
         count (particle_set%prt%get_status () == PRT_INCOMING)
    particle_set%n_out = &
         count (particle_set%prt%get_status () == PRT_OUTGOING)
    particle_set%n_vir = particle_set%n_tot &
         - particle_set%n_beam - particle_set%n_in - particle_set%n_out
  end subroutine particle_set_reset_status

\eatline
\nwendcode{}\nwbegindocs{231}Reduce a particle set to the essential entries.  The entries kept
are those with status {\tt{}INCOMING}, {\tt{}OUTGOING} or
{\tt{}RESONANT}. {\tt{}BEAM} is kept if {\tt{}keep{\char95}beams} is true. Other
entries are skipped. The correlated state matrix, if any, is also
ignored. 
\nwenddocs{}\nwbegincode{232}\moddef{Particles: particle set: TBP}\plusendmoddef
  procedure :: reduce => particle_set_reduce
\nwendcode{}\nwbegincode{233}\moddef{Particles: procedures}\plusendmoddef
  subroutine particle_set_reduce (pset_in, pset_out, keep_beams)
    class(particle_set_t), intent(in) :: pset_in
    type(particle_set_t), intent(out) :: pset_out
    logical, intent(in), optional :: keep_beams
    integer, dimension(:), allocatable :: status, map
    integer :: i, j
    logical :: kb
    kb = .false.;  if (present (keep_beams))  kb = keep_beams
    allocate (status (pset_in%n_tot))    
    status = pset_in%prt%get_status ()
    if (kb)  pset_out%n_beam  = count (status == PRT_BEAM)
    pset_out%n_in  = count (status == PRT_INCOMING)
    pset_out%n_vir = count (status == PRT_RESONANT)
    pset_out%n_out = count (status == PRT_OUTGOING)
    pset_out%n_tot = &
         pset_out%n_beam + pset_out%n_in + pset_out%n_vir + pset_out%n_out
    allocate (pset_out%prt (pset_out%n_tot))
    allocate (map (pset_in%n_tot))
    map = 0
    j = 0
    if (kb) call copy_particles (PRT_BEAM)
    call copy_particles (PRT_INCOMING)
    call copy_particles (PRT_RESONANT)
    call copy_particles (PRT_OUTGOING)
    do i = 1, pset_in%n_tot
       if (map(i) == 0)  cycle
       !!! !!! triggers nagfor bug!
       !!!  call particle_set_parents (pset_out%prt(map(i)), &
       !!!       map (particle_set_get_real_parents (pset_in, i)))
       !!!  call particle_set_children (pset_out%prt(map(i)), &
       !!!       map (particle_set_get_real_children (pset_in, i)))
       !!! !!! workaround:
       call pset_out%prt(map(i))%set_parents &
            (pset_in%get_real_parents (i, kb))
       call pset_out%prt(map(i))%set_parents &
            (map (pset_out%prt(map(i))%parent))
       call pset_out%prt(map(i))%set_children &
            (pset_in%get_real_children (i, kb))
       call pset_out%prt(map(i))%set_children &
            (map (pset_out%prt(map(i))%child))
    end do
  contains
    subroutine copy_particles (stat)
      integer, intent(in) :: stat
      integer :: i
      do i = 1, pset_in%n_tot
         if (status(i) == stat) then
            j = j + 1
            map(i) = j
            call particle_init_particle (pset_out%prt(j), pset_in%prt(i))
         end if
      end do
    end subroutine copy_particles
  end subroutine particle_set_reduce

\eatline
\nwendcode{}\nwbegindocs{234}Remove the beam particles and beam remnants from the particle set if the
keep beams flag is false. If keep beams is not given, the beam particles
and the beam remnants are removed.
The correlated state matrix, if any, is also ignored.
\nwenddocs{}\nwbegincode{235}\moddef{Particles: particle set: TBP}\plusendmoddef
  procedure :: apply_keep_beams => particle_set_apply_keep_beams
\nwendcode{}\nwbegincode{236}\moddef{Particles: procedures}\plusendmoddef
  subroutine particle_set_apply_keep_beams (pset_in, pset_out, keep_beams)
    class(particle_set_t), intent(in) :: pset_in
    type(particle_set_t), intent(out) :: pset_out
    logical, intent(in), optional :: keep_beams
    integer, dimension(:), allocatable :: status, map
    integer :: i, j
    logical :: kb
    kb = .false.;  if (present (keep_beams))  kb = keep_beams
    allocate (status (pset_in%n_tot))    
    status = particle_get_status (pset_in%prt)
    if (kb)  pset_out%n_beam  = count (status == PRT_BEAM)
    pset_out%n_in  = count (status == PRT_INCOMING)
    if (kb) then
       pset_out%n_vir = count (status == PRT_VIRTUAL) + &
            count (status == PRT_RESONANT) + &
            count (status == PRT_BEAM_REMNANT)
    else 
       pset_out%n_vir = count (status == PRT_VIRTUAL) + &
            count (status == PRT_RESONANT)
    end if
    pset_out%n_out = count (status == PRT_OUTGOING)
    pset_out%n_tot = &
         pset_out%n_beam + pset_out%n_in + pset_out%n_vir + pset_out%n_out
    allocate (pset_out%prt (pset_out%n_tot))
    allocate (map (pset_in%n_tot))
    map = 0
    j = 0
    if (kb) call copy_particles (PRT_BEAM)
    call copy_particles (PRT_INCOMING)
    if (kb) call copy_particles (PRT_BEAM_REMNANT)
    call copy_particles (PRT_RESONANT)
    call copy_particles (PRT_VIRTUAL)
    call copy_particles (PRT_OUTGOING)
    do i = 1, pset_in%n_tot
       if (map(i) == 0)  cycle
       !!! !!! triggers nagfor bug!
       !!!  call particle_set_parents (pset_out%prt(map(i)), &
       !!!       map (particle_set_get_real_parents (pset_in, i)))
       !!!  call particle_set_children (pset_out%prt(map(i)), &
       !!!       map (particle_set_get_real_children (pset_in, i)))
       !!! !!! workaround:
       call pset_out%prt(map(i))%set_parents &
            (pset_in%get_real_parents (i, kb))
       call pset_out%prt(map(i))%set_parents &
            (map (pset_out%prt(map(i))%parent))
       call pset_out%prt(map(i))%set_children &
            (pset_in%get_real_children (i, kb))
       call pset_out%prt(map(i))%set_children &
            (map (pset_out%prt(map(i))%child))
    end do
  contains
    subroutine copy_particles (stat)
      integer, intent(in) :: stat
      integer :: i
      do i = 1, pset_in%n_tot
         if (status(i) == stat) then
            j = j + 1
            map(i) = j
            call particle_init_particle (pset_out%prt(j), pset_in%prt(i))
         end if
      end do
    end subroutine copy_particles
  end subroutine particle_set_apply_keep_beams

\eatline
\nwendcode{}\nwbegindocs{237}Transform a particle set into HEPEVT-compatible form.  In this form, for each
particle, the parents and the children are contiguous in the particle array.
Usually, this requires to clone some particles.

We do not know in advance how many particles the canonical form will have.
To be on the safe side, allocate four times the original size.
\nwenddocs{}\nwbegincode{238}\moddef{Particles: particle set: TBP}\plusendmoddef
  procedure :: to_hepevt_form => particle_set_to_hepevt_form
\nwendcode{}\nwbegincode{239}\moddef{Particles: procedures}\plusendmoddef
  subroutine particle_set_to_hepevt_form (pset_in, pset_out, keep_beams)
    class(particle_set_t), intent(in) :: pset_in
    type(particle_set_t), intent(out) :: pset_out
    logical, intent(in), optional :: keep_beams
    type(particle_set_t) :: pset
    type :: particle_entry_t
       integer :: src = 0
       integer :: status = 0
       integer :: orig = 0
       integer :: copy = 0
    end type particle_entry_t
    type(particle_entry_t), dimension(:), allocatable :: prt
    integer, dimension(:), allocatable :: map1, map2
    integer, dimension(:), allocatable :: parent, child
    integer :: n_tot, n_parents, n_children, i, j, c, n

    call pset_in%apply_keep_beams (pset, keep_beams)
    n_tot = pset%n_tot
    allocate (prt (4 * n_tot))
    allocate (map1(4 * n_tot))
    allocate (map2(4 * n_tot))
    map1 = 0
    map2 = 0
    allocate (child (n_tot))
    allocate (parent (n_tot))
    n = 0
    do i = 1, n_tot
       if (pset%prt(i)%get_n_parents () == 0) then
          call append (i)
       end if
    end do
    do i = 1, n_tot
       n_children = pset%prt(i)%get_n_children ()
       if (n_children > 0) then
          child(1:n_children) = pset%prt(i)%get_children ()
          c = child(1)
          if (map1(c) == 0) then
             n_parents = pset%prt(c)%get_n_parents ()
             if (n_parents > 1) then
                parent(1:n_parents) = pset%prt(c)%get_parents ()
                if (i == parent(1) .and. &
                    any( [(map1(i)+j-1, j=1,n_parents)] /= &
                           map1(parent(1:n_parents)))) &
                    then
                   do j = 1, n_parents
                      call append (parent(j))
                   end do
                end if
             else if (map1(i) == 0) then
                call append (i)
             end if
             do j = 1, n_children
                call append (child(j))
             end do
          end if
       else if (map1(i) == 0) then
          call append (i)
       end if
    end do
    do i = n, 1, -1
       if (prt(i)%status /= PRT_OUTGOING) then
          do j = 1, i-1
             if (prt(j)%status == PRT_OUTGOING) then
                call append(prt(j)%src)
             end if
          end do
          exit
       end if
    end do
    pset_out%n_beam = count (prt(1:n)%status == PRT_BEAM)
    pset_out%n_in   = count (prt(1:n)%status == PRT_INCOMING)
    pset_out%n_vir  = count (prt(1:n)%status == PRT_RESONANT)
    pset_out%n_out  = count (prt(1:n)%status == PRT_OUTGOING)
    pset_out%n_tot = n
    allocate (pset_out%prt (n))
    do i = 1, n
       call particle_init_particle (pset_out%prt(i), pset%prt(prt(i)%src))
       call particle_reset_status (pset_out%prt(i), prt(i)%status)
       if (prt(i)%orig == 0) then
          !!! !!! This causes nagfor 5.2 (770) Panic
          !!!  call particle_set_parents &
          !!!       (pset_out%prt(i), &
          !!!        map2 (particle_get_parents (pset%prt(prt(i)%src))))
          !!! !!! Workaround
          n_parents = particle_get_n_parents (pset%prt(prt(i)%src))
          parent(1:n_parents) = particle_get_parents (pset%prt(prt(i)%src))
          call particle_set_parents (pset_out%prt(i), &
                                     map2(parent(1:n_parents)))
       else
          call particle_set_parents (pset_out%prt(i), [ prt(i)%orig ])
       end if
       if (prt(i)%copy == 0) then
          !!! !!! This causes nagfor 5.2 (770) Panic
          !!!  call particle_set_children &
          !!!       (pset_out%prt(i), &
          !!!        map1 (particle_get_children (pset%prt(prt(i)%src))))
          !!! !!! Workaround
          n_children = particle_get_n_children (pset%prt(prt(i)%src))
          child(1:n_children) = particle_get_children (pset%prt(prt(i)%src))
          call particle_set_children (pset_out%prt(i), &
                                      map1(child(1:n_children)))
       else
          call particle_set_children (pset_out%prt(i), [ prt(i)%copy ])
       end if
    end do
  contains
    subroutine append (i)
      integer, intent(in) :: i
      n = n + 1
      if (n > size (prt)) &
           call msg_bug ("Particle set transform to HEPEVT: insufficient space")
      prt(n)%src = i
      prt(n)%status = pset%prt(i)%get_status ()
      if (map1(i) == 0) then
         map1(i) = n
      else
         prt(map2(i))%status = PRT_VIRTUAL
         prt(map2(i))%copy = n
         prt(n)%orig = map2(i)
      end if
      map2(i) = n
    end subroutine append
  end subroutine particle_set_to_hepevt_form

\eatline
\nwendcode{}\nwbegindocs{240}This procedure aims at reconstructing the momenta of an interaction,
given a particle set.  Since the particle orderings
\nwenddocs{}\nwbegincode{241}\moddef{Particles: particle set: TBP}\plusendmoddef
  procedure :: fill_interaction => particle_set_fill_interaction
\nwendcode{}\nwbegincode{242}\moddef{Particles: procedures}\plusendmoddef
  subroutine particle_set_fill_interaction &
       (pset, int, n_in, recover_beams, check_match, state_flv)
    class(particle_set_t), intent(in) :: pset
    type(interaction_t), intent(inout) :: int
    integer, intent(in) :: n_in
    logical, intent(in), optional :: recover_beams, check_match
    type(state_flv_content_t), intent(in), optional :: state_flv
    integer, dimension(:), allocatable :: map, pdg
    integer, dimension(:), allocatable :: i_in, i_out, p_in, p_out
    logical, dimension(:), allocatable :: i_set
    integer :: n_out, i, p
    logical :: r_beams, check
    r_beams = .false.;  if (present (recover_beams))  r_beams = recover_beams
    check = .true.;  if (present (check_match))  check = check_match
    if (check) then
       call find_hard_process_in_int  (i_in, i_out)
       call find_hard_process_in_pset (p_in, p_out)
       n_out = size (i_out)
       if (size (i_in) /= n_in)  call err_int_n_in
       if (size (p_in) /= n_in)  call err_pset_n_in
       if (size (p_out) /= n_out)  call err_pset_n_out
       call extract_hard_process_from_pset (pdg)
       call determine_map_for_hard_process (map, state_flv)
       if (.not. r_beams) then
          select case (n_in)
          case (1)
             call recover_parents (p_in(1), map)
          case (2)
             do i = 1, 2
                call recover_parents (p_in(i), map)
             end do
             do p = 1, 2
                call recover_radiation (p, map)
             end do
          end select
       end if
    else
       allocate (map (interaction_get_n_tot (int)))
       map = [(i, i = 1, size (map))]
       r_beams = .false.
    end if
    allocate (i_set (interaction_get_n_tot (int)), source = .false.)
    do p = 1, size (map)
       if (map(p) /= 0) then
          i_set(map(p)) = .true.
          call interaction_set_momentum (int, &
               pset%prt(p)%get_momentum (), map(p))
       end if
    end do
    if (r_beams) then
       do i = 1, n_in
          call reconstruct_beam_and_radiation (i, i_set)
       end do
    end if
    if (any (.not. i_set))  call err_map
  contains
    subroutine find_hard_process_in_pset (p_in, p_out)
      integer, dimension(:), allocatable, intent(out) :: p_in, p_out
      integer, dimension(:), allocatable :: p_status, p_idx
      integer :: n_in_p, n_out_p
      integer :: i
      allocate (p_status (pset%n_tot), p_idx (pset%n_tot))
      p_status = particle_get_status (pset%prt)
      p_idx = [(i, i = 1, pset%n_tot)]
      n_in_p = count (p_status == PRT_INCOMING)
      allocate (p_in (n_in))
      p_in = pack (p_idx, p_status == PRT_INCOMING)
      if (size (p_in) == 0)  call err_pset_hard
      i = p_in(1)
      n_out_p = particle_get_n_children (pset%prt(i))
      allocate (p_out (n_out_p))
      p_out = particle_get_children (pset%prt(i))
    end subroutine find_hard_process_in_pset
    subroutine find_hard_process_in_int (i_in, i_out)
      integer, dimension(:), allocatable, intent(out) :: i_in, i_out
      integer, dimension(:), allocatable :: p_status, p_idx
      integer :: n_in_i, n_out_i
      integer :: i
      i = interaction_get_n_tot (int)
      n_in_i = interaction_get_n_parents (int, i)
      if (n_in_i /= n_in)  call err_int_n_in
      allocate (i_in (n_in))
      i_in = interaction_get_parents (int, i)
      i = i_in(1)
      n_out = interaction_get_n_children (int, i)
      allocate (i_out (n_out))
      i_out = interaction_get_children (int, i)
    end subroutine find_hard_process_in_int
    subroutine extract_hard_process_from_pset (pdg)
      integer, dimension(:), allocatable, intent(out) :: pdg
      integer, dimension(:), allocatable :: pdg_p
      logical, dimension(:), allocatable :: mask_p
      allocate (pdg_p (pset%n_tot))
      pdg_p = particle_get_pdg (pset%prt)
      allocate (mask_p (pset%n_tot), source = .false.)
      mask_p (p_in) = .true.
      mask_p (p_out) = .true.
      allocate (pdg (n_in + n_out))
      pdg = pack (pdg_p, mask_p)
    end subroutine extract_hard_process_from_pset
    subroutine determine_map_for_hard_process (map, state_flv)
      integer, dimension(:), allocatable, intent(out) :: map
      type(state_flv_content_t), intent(in), optional :: state_flv
      integer, dimension(:), allocatable :: pdg_i, map_i
      integer :: n_tot
      logical, dimension(:), allocatable :: mask_i, mask_p
      logical :: success
      n_tot = interaction_get_n_tot (int)
      allocate (map (n_tot), source = 0)
      if (present (state_flv)) then
         allocate (mask_i (n_tot), source = .false.)
         mask_i (i_in) = .true.
         mask_i (i_out) = .true.
         allocate (pdg_i (n_tot), map_i (n_tot))
         pdg_i = unpack (pdg, mask_i, 0)
         call state_flv%match (pdg_i, success, map_i)
         allocate (mask_p (pset%n_tot), source = .false.)
         mask_p (p_in) = .true.
         mask_p (p_out) = .true.
         map = unpack (pack (map_i, mask_i), mask_p, 0)
         if (.not. success)  call err_mismatch
      else
         map(p_in) = i_in
         map(p_out) = i_out
      end if
    end subroutine determine_map_for_hard_process
    recursive subroutine recover_parents (p, map)
      integer, intent(in) :: p
      integer, dimension(:), intent(inout) :: map
      integer :: i, n, n_p, q, k
      integer, dimension(:), allocatable :: i_parents, p_parents
      integer, dimension(1) :: pp
      i = map(p)
      n = interaction_get_n_parents (int, i)
      q = p
      n_p = particle_get_n_parents (pset%prt(q))
      do while (n_p == 1)
         pp = particle_get_parents (pset%prt(q))
         if (pset%prt(pp(1))%get_n_children () > 1)  exit
         q = pp(1)
         n_p = pset%prt(q)%get_n_parents ()
      end do
      if (n_p /= n)  call err_map
      allocate (i_parents (n), p_parents (n))
      i_parents = interaction_get_parents (int, i)
      p_parents = pset%prt(q)%get_parents ()
      do k = 1, n
         q = p_parents(k)
         if (map(q) == 0) then
            map(q) = i_parents(k)
            call recover_parents (q, map)
         end if
      end do
    end subroutine recover_parents
    recursive subroutine recover_radiation (p, map)
      integer, intent(in) :: p
      integer, dimension(:), intent(inout) :: map
      integer :: i, n, n_p, q, k
      integer, dimension(:), allocatable :: i_children, p_children
      if (particle_get_status (pset%prt(p)) == PRT_INCOMING)  return
      i = map(p)
      n = interaction_get_n_children (int, i)
      n_p = pset%prt(p)%get_n_children ()
      if (n_p /= n)  call err_map
      allocate (i_children (n), p_children (n))
      i_children = interaction_get_children (int, i)
      p_children = pset%prt(p)%get_children ()
      do k = 1, n
         q = p_children(k)
         if (map(q) == 0) then
            i = i_children(k)
            if (interaction_get_n_children (int, i) == 0) then
               map(q) = i
            else
               select case (n)
               case (2)
                  select case (k)
                  case (1);  map(q) = i_children(2)
                  case (2);  map(q) = i_children(1)
                  end select
               case (4)
                  select case (k)
                  case (1);  map(q) = i_children(3)
                  case (2);  map(q) = i_children(4)
                  case (3);  map(q) = i_children(1)
                  case (4);  map(q) = i_children(2)
                  end select
               case default
                  call err_radiation
               end select
            end if
         else
            call recover_radiation (q, map)
         end if
      end do
    end subroutine recover_radiation
    subroutine reconstruct_beam_and_radiation (k, i_set)
      integer, intent(in) :: k
      logical, dimension(:), intent(inout) :: i_set
      integer :: k_src, k_in, k_rad
      type(interaction_t), pointer :: int_src
      integer, dimension(2) :: i_child
      call interaction_find_source (int, k, int_src, k_src)
      call interaction_set_momentum (int, &
           interaction_get_momentum (int_src, k_src), k)
      i_set(k) = .true.
      if (n_in == 2) then
         i_child = interaction_get_children (int, k)
         if (interaction_get_n_children (int, i_child(1)) > 0) then
            k_in = i_child(1);  k_rad = i_child(2)
         else
            k_in = i_child(2);  k_rad = i_child(1)
         end if
         if (.not. i_set(k_in))  call err_beams
         call interaction_set_momentum (int, &
              interaction_get_momentum (int, k) &
              - interaction_get_momentum (int, k_in), k_rad)
         i_set(k_rad) = .true.
      end if
    end subroutine reconstruct_beam_and_radiation
    subroutine err_pset_hard
      call msg_fatal ("Reading particle set: no particles marked as incoming")
    end subroutine err_pset_hard
    subroutine err_int_n_in
      integer :: n
      if (allocated (i_in)) then
         n = size (i_in)
      else
         n = 0
      end if
      write (msg_buffer, "(A,I0,A,I0)") &
           "Filling hard process from particle set: expect ", n_in, &
           " incoming particle(s), found ", n
      call msg_bug
    end subroutine err_int_n_in
    subroutine err_pset_n_in
      write (msg_buffer, "(A,I0,A,I0)") &
           "Reading hard-process particle set: should contain ", n_in, &
           " incoming particle(s), found ", size (p_in)
      call msg_fatal
    end subroutine err_pset_n_in
    subroutine err_pset_n_out
      write (msg_buffer, "(A,I0,A,I0)") &
           "Reading hard-process particle set: should contain ", n_out, &
           " outgoing particle(s), found ", size (p_out)
      call msg_fatal
    end subroutine err_pset_n_out
    subroutine err_mismatch
      call pset%write ()
      call state_flv%write ()
      call msg_fatal ("Reading particle set: Flavor combination " &
           // "does not match requested process")
    end subroutine err_mismatch
    subroutine err_map
      call pset%write ()
      call interaction_write (int)
      call msg_fatal ("Reading hard-process particle set: " &
           // "Incomplete mapping from particle set to interaction")
    end subroutine err_map
    subroutine err_beams
      call pset%write ()
      call interaction_write (int)
      call msg_fatal ("Reading particle set: Beam structure " &
           // "does not match requested process")
    end subroutine err_beams
    subroutine err_radiation
      call interaction_write (int)
      call msg_bug ("Reading particle set: Interaction " &
           // "contains inconsistent radiation pattern.")
    end subroutine err_radiation
  end subroutine particle_set_fill_interaction
    
\eatline
\nwendcode{}\nwbegindocs{243}\nwdocspar
This procedure reconstructs an array of vertex indices from the
parent-child information in the particle entries, according to the
HepMC scheme.  For each particle, we determine which vertex it comes
from and which vertex it goes to.  We return the two arrays and the
maximum vertex index.

For each particle in the list, we first check its parents.  If for any
parent the vertex where it goes to is already known, this vertex index
is assigned as the current 'from' vertex.  Otherwise, a new index is
created, assigned as the current 'from' vertex, and as the 'to' vertex
for all parents.

Then, the analogous procedure is done for the children.
\nwenddocs{}\nwbegincode{244}\moddef{Particles: particle set: TBP}\plusendmoddef
  procedure :: assign_vertices => particle_set_assign_vertices
\nwendcode{}\nwbegincode{245}\moddef{Particles: procedures}\plusendmoddef
  subroutine particle_set_assign_vertices &
       (particle_set, v_from, v_to, n_vertices)
    class(particle_set_t), intent(in) :: particle_set
    integer, dimension(:), intent(out) :: v_from, v_to
    integer, intent(out) :: n_vertices
    integer, dimension(:), allocatable :: parent, child
    integer :: n_parents, n_children, vf, vt
    integer :: i, j, v
    v_from = 0
    v_to = 0
    vf = 0
    vt = 0
    do i = 1, particle_set%n_tot
       n_parents = particle_set%prt(i)%get_n_parents ()
       if (n_parents /= 0) then
          allocate (parent (n_parents))
          parent = particle_set%prt(i)%get_parents ()
          SCAN_PARENTS: do j = 1, size (parent)
             v = v_to(parent(j))
             if (v /= 0) then
                v_from(i) = v;  exit SCAN_PARENTS
             end if
          end do SCAN_PARENTS
          if (v_from(i) == 0) then
             vf = vf + 1;  v_from(i) = vf
             v_to(parent) = vf
          end if
          deallocate (parent)
       end if
       n_children = particle_set%prt(i)%get_n_children ()
       if (n_children /= 0) then
          allocate (child (n_children))
          child = particle_set%prt(i)%get_children ()
          SCAN_CHILDREN: do j = 1, size (child)
             v = v_from(child(j))
             if (v /= 0) then
                v_to(i) = v;  exit SCAN_CHILDREN
             end if
          end do SCAN_CHILDREN
          if (v_to(i) == 0) then
             vt = vt + 1;  v_to(i) = vt
             v_from(child) = vt
          end if
          deallocate (child)
       end if
    end do
    n_vertices = max (vf, vt)
  end subroutine particle_set_assign_vertices

\eatline
\nwendcode{}\nwbegindocs{246}\nwdocspar
\subsection{Expression interface}
This converts a {\tt{}particle{\char95}set} object as defined here to a more
concise {\tt{}subevt} object that can be used as the event root of an
expression.  In particular, the latter lacks virtual particles, spin
correlations and parent-child relations.

We keep beam particles, incoming partons, and outgoing partons.
Furthermore, we keep radiated particles (a.k.a.\ beam remnants) if
they have no children in the current particle set, and mark them as
outgoing particles.
\nwenddocs{}\nwbegincode{247}\moddef{Particles: particle set: TBP}\plusendmoddef
  procedure :: to_subevt => particle_set_to_subevt
\nwendcode{}\nwbegincode{248}\moddef{Particles: procedures}\plusendmoddef
  subroutine particle_set_to_subevt (particle_set, subevt)
    class(particle_set_t), intent(in) :: particle_set
    type(subevt_t), intent(out) :: subevt
    integer :: n_tot, n_beam, n_in, n_out, n_rad
    integer :: i, k, n_active
    integer, dimension(2) :: hel
    logical :: keep
    n_tot  = particle_set_get_n_tot  (particle_set)
    n_beam = particle_set_get_n_beam (particle_set)
    n_in   = particle_set_get_n_in   (particle_set)
    n_out  = particle_set_get_n_out  (particle_set)
    n_rad  = particle_set_get_n_rad  (particle_set)
    call subevt_init (subevt, n_beam + n_rad + n_in + n_out)
    k = 0
    do i = 1, n_tot
       associate (prt => particle_set%prt(i))
         keep = .false.
         select case (particle_get_status (prt))
         case (PRT_BEAM)
            k = k + 1
            call subevt_set_beam (subevt, k, &
                 particle_get_pdg (prt), &
                 particle_get_momentum (prt), &
                 particle_get_p2 (prt))
            keep = .true.
         case (PRT_INCOMING)
            k = k + 1
            call subevt_set_incoming (subevt, k, &
                 particle_get_pdg (prt), &
                 particle_get_momentum (prt), &
                 particle_get_p2 (prt))
            keep = .true.
         case (PRT_OUTGOING)
            k = k + 1
            call subevt_set_outgoing (subevt, k, &
                 particle_get_pdg (prt), &
                 particle_get_momentum (prt), &
                 particle_get_p2 (prt))
            keep = .true.
         case (PRT_BEAM_REMNANT)
            if (particle_get_n_children (prt) == 0) then
               k = k + 1
               call subevt_set_outgoing (subevt, k, &
                    particle_get_pdg (prt), &
                    particle_get_momentum (prt), &
                    particle_get_p2 (prt))
               keep = .true.
            end if
         end select
         if (keep) then
            if (prt%polarization == PRT_DEFINITE_HELICITY) then
               if (prt%hel%is_diagonal ()) then
                  hel = prt%hel%to_pair ()
                  call subevt_polarize (subevt, k, hel(1))
               end if
            end if
         end if
       end associate
       n_active = k
    end do
    call subevt_reset (subevt, n_active)
  end subroutine particle_set_to_subevt

\eatline
\nwendcode{}\nwbegindocs{249}\nwdocspar
This replaces the {\tt{}particle{\char92}{\char95}set{\char92}{\char37}prt\ array} with a given array of particles
\nwenddocs{}\nwbegincode{250}\moddef{Particles: particle set: TBP}\plusendmoddef
  procedure :: replace => particle_set_replace
\nwendcode{}\nwbegincode{251}\moddef{Particles: procedures}\plusendmoddef
  subroutine particle_set_replace (particle_set, newprt)
    class(particle_set_t), intent(inout) :: particle_set
    type(particle_t), intent(in), dimension(:), allocatable :: newprt
    if (allocated (particle_set%prt))  deallocate (particle_set%prt)
    allocate (particle_set%prt(size (newprt)))
    particle_set%prt = newprt
    particle_set%n_tot = size (newprt)
    particle_set%n_beam = count (particle_get_status (newprt) == PRT_BEAM)
    particle_set%n_in = count (particle_get_status (newprt) == PRT_INCOMING)
    particle_set%n_out = count (particle_get_status (newprt) == PRT_OUTGOING)
    particle_set%n_vir = particle_set%n_tot &
         - particle_set%n_beam - particle_set%n_in - particle_set%n_out
  end subroutine particle_set_replace

\eatline
\nwendcode{}\nwbegindocs{252}\nwdocspar
Eliminate numerical noise
\nwenddocs{}\nwbegincode{253}\moddef{Particles: public}\plusendmoddef
  public :: pacify
\nwendcode{}\nwbegincode{254}\moddef{Particles: interfaces}\endmoddef
  interface pacify
     module procedure pacify_particle
     module procedure pacify_particle_set
  end interface pacify

\nwendcode{}\nwbegincode{255}\moddef{Particles: procedures}\plusendmoddef
  subroutine pacify_particle (prt)
    class(particle_t), intent(inout) :: prt
    real(default) :: e
    e = epsilon (1._default) * energy (prt%p)
    call pacify (prt%p, 10 * e)
    call pacify (prt%p2, 1e4 * e)
  end subroutine pacify_particle
  
  subroutine pacify_particle_set (pset)
    class(particle_set_t), intent(inout) :: pset
    integer :: i
    do i = 1, pset%n_tot
       call pacify (pset%prt(i))
    end do
  end subroutine pacify_particle_set

\eatline
\nwendcode{}\nwbegindocs{256}\nwdocspar
\subsection{Unit tests}
Set up a chain of production and decay and factorize the result into
particles.  The process is $d\bar d \to Z \to q\bar q$.
\nwenddocs{}\nwbegincode{257}\moddef{Particles: public}\plusendmoddef
  public :: particles_test
\nwendcode{}\nwbegincode{258}\moddef{Particles: procedures}\plusendmoddef
  subroutine particles_test (u, results)
    integer, intent(in) :: u
    type(test_results_t), intent(inout) :: results
  \LA{}Particles: execute tests\RA{}  
  end subroutine particles_test

\nwendcode{}\nwbegindocs{259} %def particles_test  
\nwenddocs{}\nwbegindocs{260}Check the basic setup of the {\tt{}particle{\char95}set{\char95}t} type.
\nwenddocs{}\nwbegincode{261}\moddef{Particles: execute tests}\endmoddef
  call test (particles_1, "particles_1", &
       "check particle_set routines", &
       u, results)
\nwendcode{}\nwbegincode{262}\moddef{Particles: tests}\endmoddef
  subroutine particles_1 (u)
    use os_interface
    integer, intent(in) :: u
    type(model_data_t), target :: model
    type(flavor_t), dimension(3) :: flv
    type(color_t), dimension(3) :: col
    type(helicity_t), dimension(3) :: hel
    type(quantum_numbers_t), dimension(3) :: qn
    type(vector4_t), dimension(3) :: p
    type(interaction_t), target :: int1, int2
    type(quantum_numbers_mask_t) :: qn_mask_conn, qn_rest
    type(evaluator_t), target :: eval
    type(interaction_t), pointer :: int
    type(particle_set_t) :: particle_set1, particle_set2
    type(particle_set_t) :: particle_set3, particle_set4
    type(subevt_t) :: subevt
    logical :: ok
    integer :: unit, iostat      

    write (u, "(A)")  "* Test output: Particles"
    write (u, "(A)")  "*   Purpose: test particle_set routines"
    write (u, "(A)")      

    write (u, "(A)")  "* Reading model file"

    call model%init_sm_test ()

    write (u, "(A)")
    write (u, "(A)")  "* Initializing production process"

    call interaction_init (int1, 2, 0, 1, set_relations=.true.)
    call flv%init ([1, -1, 23], model)
    call col%init_col_acl ([0, 0, 0], [0, 0, 0])
    call hel(3)%init (1, 1)
    call qn%init (flv, col, hel)
    call interaction_add_state (int1, qn, value=(0.25_default, 0._default))
    call hel(3)%init (1,-1)
    call qn%init (flv, col, hel)
    call interaction_add_state (int1, qn, value=(0._default, 0.25_default))
    call hel(3)%init (-1, 1)
    call qn%init (flv, col, hel)
    call interaction_add_state (int1, qn, value=(0._default,-0.25_default))
    call hel(3)%init (-1,-1)
    call qn%init (flv, col, hel)
    call interaction_add_state (int1, qn, value=(0.25_default, 0._default))
    call hel(3)%init (0, 0)
    call qn%init (flv, col, hel)
    call interaction_add_state (int1, qn, value=(0.5_default, 0._default))
    call interaction_freeze (int1)
    p(1) = vector4_moving (45._default, 45._default, 3)
    p(2) = vector4_moving (45._default,-45._default, 3)
    p(3) = p(1) + p(2)
    call interaction_set_momenta (int1, p)

    write (u, "(A)")
    write (u, "(A)")  "* Setup decay process"

    call interaction_init (int2, 1, 0, 2, set_relations=.true.)
    call flv%init ([23, 1, -1], model)
    call col%init_col_acl ([0, 501, 0], [0, 0, 501])
    call hel%init ([1, 1, 1], [1, 1, 1])
    call qn%init (flv, col, hel)
    call interaction_add_state (int2, qn, value=(1._default, 0._default))
    call hel%init ([1, 1, 1], [-1,-1,-1])
    call qn%init (flv, col, hel)
    call interaction_add_state (int2, qn, value=(0._default, 0.1_default))
    call hel%init ([-1,-1,-1], [1, 1, 1])
    call qn%init (flv, col, hel)
    call interaction_add_state (int2, qn, value=(0._default,-0.1_default))
    call hel%init ([-1,-1,-1], [-1,-1,-1])
    call qn%init (flv, col, hel)
    call interaction_add_state (int2, qn, value=(1._default, 0._default))
    call hel%init ([0, 1,-1], [0, 1,-1])
    call qn%init (flv, col, hel)
    call interaction_add_state (int2, qn, value=(4._default, 0._default))
    call hel%init ([0,-1, 1], [0, 1,-1])
    call qn%init (flv, col, hel)
    call interaction_add_state (int2, qn, value=(2._default, 0._default))
    call hel%init ([0, 1,-1], [0,-1, 1])
    call qn%init (flv, col, hel)
    call interaction_add_state (int2, qn, value=(2._default, 0._default))
    call hel%init ([0,-1, 1], [0,-1, 1])
    call qn%init (flv, col, hel)
    call interaction_add_state (int2, qn, value=(4._default, 0._default))
    call flv%init ([23, 2, -2], model)
    call hel%init ([0, 1,-1], [0, 1,-1])
    call qn%init (flv, col, hel)
    call interaction_add_state (int2, qn, value=(0.5_default, 0._default))
    call hel%init ([0,-1, 1], [0,-1, 1])
    call qn%init (flv, col, hel)
    call interaction_add_state (int2, qn, value=(0.5_default, 0._default))
    call interaction_freeze (int2)
    p(2) = vector4_moving (45._default, 45._default, 2)
    p(3) = vector4_moving (45._default,-45._default, 2)
    call interaction_set_momenta (int2, p)
    call interaction_set_source_link (int2, 1, int1, 3)
    call interaction_write (int1, u)
    call interaction_write (int2, u)

    write (u, "(A)")
    write (u, "(A)")  "* Concatenate production and decay"

    call evaluator_init_product (eval, int1, int2, qn_mask_conn, &
         connections_are_resonant=.true.)
    call evaluator_receive_momenta (eval)
    call eval%evaluate ()
    call eval%write (u)
    
    write (u, "(A)")
    write (u, "(A)")  "* Factorize as subevent (complete, polarized)"
    write (u, "(A)")
    
    int => evaluator_get_int_ptr (eval)
    call particle_set1%init &
         (ok, int, int, FM_FACTOR_HELICITY, &
          [0.2_default, 0.2_default], .false., .true.)
    call particle_set_write (particle_set1, u)

    write (u, "(A)")
    write (u, "(A)")  "* Factorize as subevent (in/out only, selected helicity)"
    write (u, "(A)")
    
    int => evaluator_get_int_ptr (eval)
    call particle_set2%init &
         (ok, int, int, FM_SELECT_HELICITY, &
          [0.9_default, 0.9_default], .false., .false.)
    call particle_set2%write (u)
    call particle_set2%final ()

    write (u, "(A)")
    write (u, "(A)")  "* Factorize as subevent (complete, selected helicity)"
    write (u, "(A)") 
    
    int => evaluator_get_int_ptr (eval)
    call particle_set2%init &
         (ok, int, int, FM_SELECT_HELICITY, &
          [0.7_default, 0.7_default], .false., .true.)
    call particle_set2%write (u)      
        
    write (u, "(A)")
    write (u, "(A)")  &
         "* Factorize (complete, polarized, correlated); write and read again"
    write (u, "(A)")
    
    int => evaluator_get_int_ptr (eval)
    call particle_set3%init &
         (ok, int, int, FM_FACTOR_HELICITY, &
          [0.7_default, 0.7_default], .true., .true.)
    call particle_set3%write (u)

    unit = free_unit ()
    open (unit, action="readwrite", form="unformatted", status="scratch")
    call particle_set3%write_raw (unit)
    rewind (unit)
    call particle_set4%read_raw (unit, iostat=iostat)
    call particle_set4%set_model (model)
    close (unit)
    
    write (u, "(A)")
    write (u, "(A)")  "* Result from reading"
    write (u, "(A)")

    call particle_set4%write (u)

    !!! write (u, "(A)")
    !!! write (u, "(A)")  "* Transform to a subevt object"
    !!! write (u, "(A)")
    !!! 
    !!! call particle_set4%to_subevt (subevt)
    !!! call subevt_write (subevt, u)

    write (u, "(A)")
    write (u, "(A)")  "* Cleanup"
    
    call particle_set1%final ()
    call particle_set2%final ()
    call particle_set3%final ()
    call particle_set4%final ()
    call evaluator_final (eval)
    call interaction_final (int1)
    call interaction_final (int2)

    call model%final ()
       
    write (u, "(A)")
    write (u, "(A)")  "* Test output end: particles_1"        
    
  end subroutine particles_1

\eatline
\nwendcode{}\nwbegindocs{263}\nwdocspar
Reconstruct a hard interaction from a particle set.
\nwenddocs{}\nwbegincode{264}\moddef{Particles: execute tests}\plusendmoddef
  call test (particles_2, "particles_2", &
       "reconstruct hard interaction", &
       u, results)
\nwendcode{}\nwbegincode{265}\moddef{Particles: tests}\plusendmoddef
  subroutine particles_2 (u)

    integer, intent(in) :: u
    type(interaction_t) :: int
    type(state_flv_content_t) :: state_flv
    type(particle_set_t) :: pset
    type(flavor_t), dimension(:), allocatable :: flv
    type(quantum_numbers_t), dimension(:), allocatable :: qn
    integer :: i, j
    
    write (u, "(A)")  "* Test output: Particles"
    write (u, "(A)")  "*   Purpose: reconstruct simple interaction"
    write (u, "(A)")      

    write (u, "(A)")  "* Set up a 2 -> 3 interaction"
    write (u, "(A)")  "    + incoming partons marked as virtual"
    write (u, "(A)")  "    + no quantum numbers"
    write (u, "(A)")      

    call reset_interaction_counter ()
    call interaction_init (int, 0, 2, 3)
    do i = 1, 2
       do j = 3, 5
          call interaction_relate (int, i, j)
       end do
    end do

    allocate (qn (5))
    call interaction_add_state (int, qn)
    call interaction_freeze (int)

    call interaction_write (int, u)

    write (u, "(A)")      
    write (u, "(A)")  "* Manually set up a flavor-content record"
    write (u, "(A)")      

    call state_flv%init (1, &
         mask = [.false., .false., .true., .true., .true.])
    call state_flv%set_entry (1, &
         pdg = [11, 12, 3, 4, 5], &
         map = [1, 2, 3, 4, 5])
    
    call state_flv%write (u)

    write (u, "(A)")      
    write (u, "(A)")  "* Manually create a matching particle set"
    write (u, "(A)")      

    pset%n_beam = 0
    pset%n_in   = 2
    pset%n_vir  = 0
    pset%n_out  = 3
    pset%n_tot  = 5
    allocate (pset%prt (pset%n_tot))
    do i = 1, 2
       call pset%prt(i)%reset_status (PRT_INCOMING)
       call pset%prt(i)%set_children ([3,4,5])
    end do
    do i = 3, 5
       call particle_reset_status (pset%prt(i), PRT_OUTGOING)
       call particle_set_parents (pset%prt(i), [1,2])
    end do
    call pset%prt(1)%set_momentum (vector4_at_rest (1._default))
    call pset%prt(2)%set_momentum (vector4_at_rest (2._default))
    call pset%prt(3)%set_momentum (vector4_at_rest (5._default))
    call pset%prt(4)%set_momentum (vector4_at_rest (4._default))
    call pset%prt(5)%set_momentum (vector4_at_rest (3._default))

    allocate (flv (5))
    call flv%init ([11,12,5,4,3])
    do i = 1, 5
       call pset%prt(i)%set_flavor (flv(i))
    end do

    call pset%write (u)

    write (u, "(A)")      
    write (u, "(A)")  "*   Fill interaction from particle set"
    write (u, "(A)")

    call pset%fill_interaction (int, 2, state_flv=state_flv)
    call interaction_write (int, u)

    write (u, "(A)")
    write (u, "(A)")  "* Cleanup"

    call interaction_final (int)
    call pset%final ()
    
    write (u, "(A)")
    write (u, "(A)")  "* Test output end: particles_2"

  end subroutine particles_2
  
\eatline
\nwendcode{}\nwbegindocs{266}\nwdocspar
Reconstruct an interaction with beam structure, e.g., a hadronic
interaction, from a particle set.
\nwenddocs{}\nwbegincode{267}\moddef{Particles: execute tests}\plusendmoddef
  call test (particles_3, "particles_3", &
       "reconstruct interaction with beam structure", &
       u, results)
\nwendcode{}\nwbegincode{268}\moddef{Particles: tests}\plusendmoddef
  subroutine particles_3 (u)

    integer, intent(in) :: u
    type(interaction_t) :: int
    type(state_flv_content_t) :: state_flv
    type(particle_set_t) :: pset
    type(flavor_t), dimension(:), allocatable :: flv
    type(quantum_numbers_t), dimension(:), allocatable :: qn
    integer :: i, j
    
    write (u, "(A)")  "* Test output: Particles"
    write (u, "(A)")  "*   Purpose: reconstruct simple interaction"
    write (u, "(A)")      

    write (u, "(A)")  "* Set up a 2 -> 2 -> 3 interaction with radiation"
    write (u, "(A)")  "    + no quantum numbers"
    write (u, "(A)")      

    call reset_interaction_counter ()
    call interaction_init (int, 0, 6, 3)
    call interaction_relate (int, 1, 3)
    call interaction_relate (int, 1, 4)
    call interaction_relate (int, 2, 5)
    call interaction_relate (int, 2, 6)
    do i = 4, 6, 2
       do j = 7, 9
          call interaction_relate (int, i, j)
       end do
    end do

    allocate (qn (9))
    call interaction_add_state (int, qn)
    call interaction_freeze (int)

    call interaction_write (int, u)

    write (u, "(A)")      
    write (u, "(A)")  "* Manually set up a flavor-content record"
    write (u, "(A)")      

    call state_flv%init (1, &
         mask = [.false., .false., .false., .false., .false., .false., &
         .true., .true., .true.])
    call state_flv%set_entry (1, &
         pdg = [2011, 2012, 91, 11, 92, 12, 3, 4, 5], &
         map = [1, 2, 3, 4, 5, 6, 7, 8, 9])
    
    call state_flv%write (u)

    write (u, "(A)")      
    write (u, "(A)")  "* Manually create a matching particle set"
    write (u, "(A)")      

    pset%n_beam = 2
    pset%n_in   = 2
    pset%n_vir  = 2
    pset%n_out  = 3
    pset%n_tot  = 9

    allocate (pset%prt (pset%n_tot))
    call pset%prt(1)%reset_status (PRT_BEAM)
    call pset%prt(2)%reset_status (PRT_BEAM)
    call pset%prt(3)%reset_status (PRT_INCOMING)
    call pset%prt(4)%reset_status (PRT_INCOMING)
    call pset%prt(5)%reset_status (PRT_OUTGOING)
    call pset%prt(6)%reset_status (PRT_OUTGOING)
    call pset%prt(7)%reset_status (PRT_OUTGOING)
    call pset%prt(8)%reset_status (PRT_OUTGOING)
    call pset%prt(9)%reset_status (PRT_OUTGOING)

    call pset%prt(1)%set_children ([3,5])
    call pset%prt(2)%set_children ([4,6])
    call pset%prt(3)%set_children ([7,8,9])
    call pset%prt(4)%set_children ([7,8,9])

    call pset%prt(3)%set_parents ([1])
    call pset%prt(4)%set_parents ([2])
    call pset%prt(5)%set_parents ([1])
    call pset%prt(6)%set_parents ([2])
    call pset%prt(7)%set_parents ([5,6])
    call pset%prt(8)%set_parents ([5,6])
    call pset%prt(9)%set_parents ([5,6])

    call pset%prt(1)%set_momentum (vector4_at_rest (1._default))
    call pset%prt(2)%set_momentum (vector4_at_rest (2._default))
    call pset%prt(3)%set_momentum (vector4_at_rest (4._default))
    call pset%prt(4)%set_momentum (vector4_at_rest (6._default))
    call pset%prt(5)%set_momentum (vector4_at_rest (3._default))
    call pset%prt(6)%set_momentum (vector4_at_rest (5._default))
    call pset%prt(7)%set_momentum (vector4_at_rest (7._default))
    call pset%prt(8)%set_momentum (vector4_at_rest (8._default))
    call pset%prt(9)%set_momentum (vector4_at_rest (9._default))

    allocate (flv (9))
    call flv%init ([2011, 2012, 11, 12, 91, 92, 3, 4, 5])
    do i = 1, 9
       call pset%prt(i)%set_flavor (flv(i))
    end do

    call pset%write (u)
 
    write (u, "(A)")      
    write (u, "(A)")  "*   Fill interaction from particle set"
    write (u, "(A)")

    call particle_set_fill_interaction (pset, int, 2, state_flv=state_flv)
    call interaction_write (int, u)

    write (u, "(A)")
    write (u, "(A)")  "* Cleanup"

    call interaction_final (int)
    call pset%final ()
    
    write (u, "(A)")
    write (u, "(A)")  "* Test output end: particles_3"

  end subroutine particles_3
  
\eatline
\nwendcode{}\nwbegindocs{269}\nwdocspar
Reconstruct an interaction with beam structure, e.g., a hadronic
interaction, from a particle set that is missing the beam information.
\nwenddocs{}\nwbegincode{270}\moddef{Particles: execute tests}\plusendmoddef
  call test (particles_4, "particles_4", &
       "reconstruct interaction with missing beams", &
       u, results)
\nwendcode{}\nwbegincode{271}\moddef{Particles: tests}\plusendmoddef
  subroutine particles_4 (u)

    integer, intent(in) :: u
    type(interaction_t) :: int
    type(interaction_t), target :: int_beams
    type(state_flv_content_t) :: state_flv
    type(particle_set_t) :: pset
    type(flavor_t), dimension(:), allocatable :: flv
    type(quantum_numbers_t), dimension(:), allocatable :: qn
    integer :: i, j
    
    write (u, "(A)")  "* Test output: Particles"
    write (u, "(A)")  "*   Purpose: reconstruct beams"
    write (u, "(A)")      

    call reset_interaction_counter ()

    write (u, "(A)")  "* Set up an interaction that contains beams only"
    write (u, "(A)")      

    call interaction_init (int_beams, 0, 0, 2)
    call interaction_set_momentum (int_beams, &
         vector4_at_rest (1._default), 1)
    call interaction_set_momentum (int_beams, &
         vector4_at_rest (2._default), 2)
    allocate (qn (2))
    call interaction_add_state (int_beams, qn)
    call interaction_freeze (int_beams)
    
    call interaction_write (int_beams, u)

    write (u, "(A)")      
    write (u, "(A)")  "* Set up a 2 -> 2 -> 3 interaction with radiation"
    write (u, "(A)")  "    + no quantum numbers"
    write (u, "(A)")      

    call interaction_init (int, 0, 6, 3)
    call interaction_relate (int, 1, 3)
    call interaction_relate (int, 1, 4)
    call interaction_relate (int, 2, 5)
    call interaction_relate (int, 2, 6)
    do i = 4, 6, 2
       do j = 7, 9
          call interaction_relate (int, i, j)
       end do
    end do
    do i = 1, 2
       call interaction_set_source_link (int, i, int_beams, i)
    end do

    deallocate (qn)
    allocate (qn (9))
    call interaction_add_state (int, qn)
    call interaction_freeze (int)

    call interaction_write (int, u)

    write (u, "(A)")      
    write (u, "(A)")  "* Manually set up a flavor-content record"
    write (u, "(A)")      

    call state_flv%init (1, &
         mask = [.false., .false., .false., .false., .false., .false., &
         .true., .true., .true.])
    call state_flv%set_entry (1, &
         pdg = [2011, 2012, 91, 11, 92, 12, 3, 4, 5], &
         map = [1, 2, 3, 4, 5, 6, 7, 8, 9])
    
    call state_flv%write (u)

    write (u, "(A)")      
    write (u, "(A)")  "* Manually create a matching particle set"
    write (u, "(A)")      

    pset%n_beam = 0
    pset%n_in   = 2
    pset%n_vir  = 0
    pset%n_out  = 3
    pset%n_tot  = 5

    allocate (pset%prt (pset%n_tot))
    call pset%prt(1)%reset_status (PRT_INCOMING)
    call pset%prt(2)%reset_status (PRT_INCOMING)
    call pset%prt(3)%reset_status (PRT_OUTGOING)
    call pset%prt(4)%reset_status (PRT_OUTGOING)
    call pset%prt(5)%reset_status (PRT_OUTGOING)

    call pset%prt(1)%set_children ([3,4,5])
    call pset%prt(2)%set_children ([3,4,5])

    call pset%prt(3)%set_parents ([1,2])
    call pset%prt(4)%set_parents ([1,2])
    call pset%prt(5)%set_parents ([1,2])

    call pset%prt(1)%set_momentum (vector4_at_rest (6._default))
    call pset%prt(2)%set_momentum (vector4_at_rest (6._default))
    call pset%prt(3)%set_momentum (vector4_at_rest (3._default))
    call pset%prt(4)%set_momentum (vector4_at_rest (4._default))
    call pset%prt(5)%set_momentum (vector4_at_rest (5._default))

    allocate (flv (5))
    call flv%init ([11, 12, 3, 4, 5])
    do i = 1, 5
       call pset%prt(i)%set_flavor (flv(i))
    end do

    call pset%write (u)
 
    write (u, "(A)")      
    write (u, "(A)")  "*   Fill interaction from particle set"
    write (u, "(A)")

    call pset%fill_interaction (int, 2, state_flv=state_flv, &
         recover_beams = .true.)
    call interaction_write (int, u)

    write (u, "(A)")
    write (u, "(A)")  "* Cleanup"

    call interaction_final (int)
    call pset%final ()
    
    write (u, "(A)")
    write (u, "(A)")  "* Test output end: particles_4"

  end subroutine particles_4
  
\eatline
\nwendcode{}\nwbegindocs{272}\nwdocspar
Reconstruct an interaction with beam structure and cloned particles
(radiated particles repeated in the event record, to maintain some
canonical ordering).
\nwenddocs{}\nwbegincode{273}\moddef{Particles: execute tests}\plusendmoddef
  call test (particles_5, "particles_5", &
       "reconstruct interaction with beams and duplicate entries", &
       u, results)
\nwendcode{}\nwbegincode{274}\moddef{Particles: tests}\plusendmoddef
  subroutine particles_5 (u)

    integer, intent(in) :: u
    type(interaction_t) :: int
    type(state_flv_content_t) :: state_flv
    type(particle_set_t) :: pset
    type(flavor_t), dimension(:), allocatable :: flv
    type(quantum_numbers_t), dimension(:), allocatable :: qn
    integer :: i, j
    
    write (u, "(A)")  "* Test output: Particles"
    write (u, "(A)")  "*   Purpose: reconstruct event with duplicate entries"
    write (u, "(A)")      

    write (u, "(A)")  "* Set up a 2 -> 2 -> 3 interaction with radiation"
    write (u, "(A)")  "    + no quantum numbers"
    write (u, "(A)")      

    call reset_interaction_counter ()
    call interaction_init (int, 0, 6, 3)
    call interaction_relate (int, 1, 3)
    call interaction_relate (int, 1, 4)
    call interaction_relate (int, 2, 5)
    call interaction_relate (int, 2, 6)
    do i = 4, 6, 2
       do j = 7, 9
          call interaction_relate (int, i, j)
       end do
    end do

    allocate (qn (9))
    call interaction_add_state (int, qn)
    call interaction_freeze (int)

    call interaction_write (int, u)

    write (u, "(A)")      
    write (u, "(A)")  "* Manually set up a flavor-content record"
    write (u, "(A)")      

    call state_flv%init (1, &
         mask = [.false., .false., .false., .false., .false., .false., &
         .true., .true., .true.])
    call state_flv%set_entry (1, &
         pdg = [2011, 2012, 91, 11, 92, 12, 3, 4, 5], &
         map = [1, 2, 3, 4, 5, 6, 7, 8, 9])
    
    call state_flv%write (u)

    write (u, "(A)")      
    write (u, "(A)")  "* Manually create a matching particle set"
    write (u, "(A)")      

    pset%n_beam = 2
    pset%n_in   = 2
    pset%n_vir  = 4
    pset%n_out  = 5
    pset%n_tot  = 13

    allocate (pset%prt (pset%n_tot))
    call pset%prt(1)%reset_status (PRT_BEAM)
    call pset%prt(2)%reset_status (PRT_BEAM)
    call pset%prt(3)%reset_status (PRT_VIRTUAL)
    call pset%prt(4)%reset_status (PRT_VIRTUAL)
    call pset%prt(5)%reset_status (PRT_VIRTUAL)
    call pset%prt(6)%reset_status (PRT_VIRTUAL)
    call pset%prt(7)%reset_status (PRT_INCOMING)
    call pset%prt(8)%reset_status (PRT_INCOMING)
    call pset%prt( 9)%reset_status (PRT_OUTGOING)
    call pset%prt(10)%reset_status (PRT_OUTGOING)
    call pset%prt(11)%reset_status (PRT_OUTGOING)
    call pset%prt(12)%reset_status (PRT_OUTGOING)
    call pset%prt(13)%reset_status (PRT_OUTGOING)

    call pset%prt(1)%set_children ([3,4])
    call pset%prt(2)%set_children ([5,6])
    call pset%prt(3)%set_children ([ 7])
    call pset%prt(4)%set_children ([ 9])
    call pset%prt(5)%set_children ([ 8])
    call pset%prt(6)%set_children ([10])
    call pset%prt(7)%set_children ([11,12,13])
    call pset%prt(8)%set_children ([11,12,13])

    call pset%prt(3)%set_parents ([1])
    call pset%prt(4)%set_parents ([1])
    call pset%prt(5)%set_parents ([2])
    call pset%prt(6)%set_parents ([2])
    call pset%prt( 7)%set_parents ([3])
    call pset%prt( 8)%set_parents ([5])
    call pset%prt( 9)%set_parents ([4])
    call pset%prt(10)%set_parents ([6])
    call pset%prt(11)%set_parents ([7,8])
    call pset%prt(12)%set_parents ([7,8])
    call pset%prt(13)%set_parents ([7,8])

    call pset%prt(1)%set_momentum (vector4_at_rest (1._default))
    call pset%prt(2)%set_momentum (vector4_at_rest (2._default))
    call pset%prt(3)%set_momentum (vector4_at_rest (4._default))
    call pset%prt(4)%set_momentum (vector4_at_rest (3._default))
    call pset%prt(5)%set_momentum (vector4_at_rest (6._default))
    call pset%prt(6)%set_momentum (vector4_at_rest (5._default))
    call pset%prt(7)%set_momentum (vector4_at_rest (4._default))
    call pset%prt(8)%set_momentum (vector4_at_rest (6._default))
    call pset%prt( 9)%set_momentum (vector4_at_rest (3._default))
    call pset%prt(10)%set_momentum (vector4_at_rest (5._default))
    call pset%prt(11)%set_momentum (vector4_at_rest (7._default))
    call pset%prt(12)%set_momentum (vector4_at_rest (8._default))
    call pset%prt(13)%set_momentum (vector4_at_rest (9._default))

    allocate (flv (13))
    call flv%init ([2011, 2012, 11, 91, 12, 92, 11, 12, 91, 92, 3, 4, 5])
    do i = 1, 13
       call pset%prt(i)%set_flavor (flv(i))
    end do

    call pset%write (u)
 
    write (u, "(A)")      
    write (u, "(A)")  "*   Fill interaction from particle set"
    write (u, "(A)")

    call pset%fill_interaction (int, 2, state_flv=state_flv)
    call interaction_write (int, u)

    write (u, "(A)")
    write (u, "(A)")  "* Cleanup"

    call interaction_final (int)
    call pset%final ()
    
    write (u, "(A)")
    write (u, "(A)")  "* Test output end: particles_5"

  end subroutine particles_5
  
\eatline
\nwendcode{}\nwbegindocs{275}\nwdocspar
Reconstruct an interaction with pair spectrum, e.g., beamstrahlung from a
particle set. 
\nwenddocs{}\nwbegincode{276}\moddef{Particles: execute tests}\plusendmoddef
  call test (particles_6, "particles_6", &
       "reconstruct interaction with pair spectrum", &
       u, results)
\nwendcode{}\nwbegincode{277}\moddef{Particles: tests}\plusendmoddef
  subroutine particles_6 (u)

    integer, intent(in) :: u
    type(interaction_t) :: int
    type(state_flv_content_t) :: state_flv
    type(particle_set_t) :: pset
    type(flavor_t), dimension(:), allocatable :: flv
    type(quantum_numbers_t), dimension(:), allocatable :: qn
    integer :: i, j
    
    write (u, "(A)")  "* Test output: Particles"
    write (u, "(A)")  "*   Purpose: reconstruct interaction with pair spectrum"
    write (u, "(A)")      

    write (u, "(A)")  "* Set up a 2 -> 2 -> 3 interaction with radiation"
    write (u, "(A)")  "    + no quantum numbers"
    write (u, "(A)")      

    call reset_interaction_counter ()
    call interaction_init (int, 0, 6, 3)
    do i = 1, 2
       do j = 3, 6
          call interaction_relate (int, i, j)
       end do
    end do
    do i = 5, 6
       do j = 7, 9
          call interaction_relate (int, i, j)
       end do
    end do

    allocate (qn (9))
    call interaction_add_state (int, qn)
    call interaction_freeze (int)

    call interaction_write (int, u)

    write (u, "(A)")      
    write (u, "(A)")  "* Manually set up a flavor-content record"
    write (u, "(A)")      

    call state_flv%init (1, &
         mask = [.false., .false., .false., .false., .false., .false., &
         .true., .true., .true.])
    call state_flv%set_entry (1, &
         pdg = [1011, 1012, 21, 22, 11, 12, 3, 4, 5], &
         map = [1, 2, 3, 4, 5, 6, 7, 8, 9])
    
    call state_flv%write (u)

    write (u, "(A)")      
    write (u, "(A)")  "* Manually create a matching particle set"
    write (u, "(A)")      

    pset%n_beam = 2
    pset%n_in   = 2
    pset%n_vir  = 2
    pset%n_out  = 3
    pset%n_tot  = 9

    allocate (pset%prt (pset%n_tot))
    call pset%prt(1)%reset_status (PRT_BEAM)
    call pset%prt(2)%reset_status (PRT_BEAM)
    call pset%prt(3)%reset_status (PRT_INCOMING)
    call pset%prt(4)%reset_status (PRT_INCOMING)
    call pset%prt(5)%reset_status (PRT_OUTGOING)
    call pset%prt(6)%reset_status (PRT_OUTGOING)
    call pset%prt(7)%reset_status (PRT_OUTGOING)
    call pset%prt(8)%reset_status (PRT_OUTGOING)
    call pset%prt(9)%reset_status (PRT_OUTGOING)

    call pset%prt(1)%set_children ([3,4,5,6])
    call pset%prt(2)%set_children ([3,4,5,6])
    call pset%prt(3)%set_children ([7,8,9])
    call pset%prt(4)%set_children ([7,8,9])

    call pset%prt(3)%set_parents ([1,2])
    call pset%prt(4)%set_parents ([1,2])
    call pset%prt(5)%set_parents ([1,2])
    call pset%prt(6)%set_parents ([1,2])
    call pset%prt(7)%set_parents ([3,4])
    call pset%prt(8)%set_parents ([3,4])
    call pset%prt(9)%set_parents ([3,4])

    call pset%prt(1)%set_momentum (vector4_at_rest (1._default))
    call pset%prt(2)%set_momentum (vector4_at_rest (2._default))
    call pset%prt(3)%set_momentum (vector4_at_rest (5._default))
    call pset%prt(4)%set_momentum (vector4_at_rest (6._default))
    call pset%prt(5)%set_momentum (vector4_at_rest (3._default))
    call pset%prt(6)%set_momentum (vector4_at_rest (4._default))
    call pset%prt(7)%set_momentum (vector4_at_rest (7._default))
    call pset%prt(8)%set_momentum (vector4_at_rest (8._default))
    call pset%prt(9)%set_momentum (vector4_at_rest (9._default))

    allocate (flv (9))
    call flv%init ([1011, 1012, 11, 12, 21, 22, 3, 4, 5])
    do i = 1, 9
       call pset%prt(i)%set_flavor (flv(i))
    end do

    call pset%write (u)
 
    write (u, "(A)")      
    write (u, "(A)")  "*   Fill interaction from particle set"
    write (u, "(A)")

    call particle_set_fill_interaction (pset, int, 2, state_flv=state_flv)
    call interaction_write (int, u)

    write (u, "(A)")
    write (u, "(A)")  "* Cleanup"

    call interaction_final (int)
    call pset%final ()
    
    write (u, "(A)")
    write (u, "(A)")  "* Test output end: particles_6"

  end subroutine particles_6
  
\eatline
\nwendcode{}\nwbegindocs{278}\nwdocspar
Reconstruct a hard decay interaction from a shuffled particle set.
\nwenddocs{}\nwbegincode{279}\moddef{Particles: execute tests}\plusendmoddef
  call test (particles_7, "particles_7", &
       "reconstruct decay interaction with reordering", &
       u, results)
\nwendcode{}\nwbegincode{280}\moddef{Particles: tests}\plusendmoddef
  subroutine particles_7 (u)

    integer, intent(in) :: u
    type(interaction_t) :: int
    type(state_flv_content_t) :: state_flv
    type(particle_set_t) :: pset
    type(flavor_t), dimension(:), allocatable :: flv
    type(quantum_numbers_t), dimension(:), allocatable :: qn
    integer :: i, j
    
    write (u, "(A)")  "* Test output: Particles"
    write (u, "(A)")  "*   Purpose: reconstruct decay interaction with reordering"
    write (u, "(A)")      

    write (u, "(A)")  "* Set up a 1 -> 3 interaction"
    write (u, "(A)")  "    + no quantum numbers"
    write (u, "(A)")      

    call reset_interaction_counter ()
    call interaction_init (int, 0, 1, 3)
    do j = 2, 4
       call interaction_relate (int, 1, j)
    end do

    allocate (qn (4))
    call interaction_add_state (int, qn)
    call interaction_freeze (int)

    call interaction_write (int, u)

    write (u, "(A)")      
    write (u, "(A)")  "* Manually set up a flavor-content record"
    write (u, "(A)")  "*   assumed interaction: 6 12 5 -11"
    write (u, "(A)")      

    call state_flv%init (1, &
         mask = [.false., .true., .true., .true.])
    call state_flv%set_entry (1, &
         pdg = [6, 5, -11, 12], &
         map = [1, 4, 2, 3])
    
    call state_flv%write (u)

    write (u, "(A)")      
    write (u, "(A)")  "* Manually create a matching particle set"
    write (u, "(A)")      

    pset%n_beam = 0
    pset%n_in   = 1
    pset%n_vir  = 0
    pset%n_out  = 3
    pset%n_tot  = 4
    allocate (pset%prt (pset%n_tot))
    do i = 1, 1
       call pset%prt(i)%reset_status (PRT_INCOMING)
       call pset%prt(i)%set_children ([2,3,4])
    end do
    do i = 2, 4
       call pset%prt(i)%reset_status (PRT_OUTGOING)
       call pset%prt(i)%set_parents ([1])
    end do
    call pset%prt(1)%set_momentum (vector4_at_rest (1._default))
    call pset%prt(2)%set_momentum (vector4_at_rest (3._default))
    call pset%prt(3)%set_momentum (vector4_at_rest (2._default))
    call pset%prt(4)%set_momentum (vector4_at_rest (4._default))

    allocate (flv (5))
    call flv%init ([6,5,12,-11])
    do i = 1, 4
       call pset%prt(i)%set_flavor (flv(i))
    end do

    call pset%write (u)

    write (u, "(A)")      
    write (u, "(A)")  "*   Fill interaction from particle set"
    write (u, "(A)")

    call particle_set_fill_interaction (pset, int, 1, state_flv=state_flv)
    call interaction_write (int, u)

    write (u, "(A)")
    write (u, "(A)")  "* Cleanup"

    call interaction_final (int)
    call pset%final ()
    
    write (u, "(A)")
    write (u, "(A)")  "* Test output end: particles_7"

  end subroutine particles_7
  
\eatline
\nwendcode{}\nwbegindocs{281}\nwdocspar
  

\nwenddocs{}\nwfilename{../../../src/beams/beams.nw}\nwbegindocs{0}%% -*- ess-noweb-default-code-mode: f90-mode; noweb-default-code-mode: f90-mode; -*- 
% WHIZARD code as NOWEB source: beams and beam structure
%%%%%%%%%%%%%%%%%%%%%%%%%%%%%%%%%%%%%%%%%%%%%%%%%%%%%%%%%%%%%%%%%%%%%%%%
\chapter{Beams}

These modules implement beam configuration and beam structure, the
latter in abstract terms.
\begin{description}
\item[beam\_structures]
  The {\tt{}beam{\char95}structure{\char95}t} type is a messenger type that communicates
  the user settings to the \whizard\ core.
\item[beams]
  Beam configuration.
\item[sf\_aux]
  Tools for handling structure functions and splitting
\item[sf\_mappings]
  Mapping functions, useful for structure function implementation
\item[sf\_base]
  The abstract structure-function interaction and structure-function
  chain types.
\end{description}

These are the implementation modules, the concrete counterparts of
{\tt{}sf{\char95}base}:
\begin{description}
\item[sf\_isr]
  ISR structure function (photon radiation inclusive and resummed in
  collinear and IR regions).
\item[sf\_epa]
  Effective Photon Approximation.
\item[sf\_ewa]
  Effective $W$ (and $Z$) approximation.
\item[sf\_escan]
  Energy spectrum that emulates a uniform energy scan.
\item[sf\_beam\_events]
  Beam-event generator that reads its input from an external file.
\item[sf\_circe1]
  CIRCE1 beam spectra for electrons and photons.
\item[sf\_circe2]
  CIRCE2 beam spectra for electrons and photons.
\item[hoppet\_interface]
  Support for $b$-quark matching, addon to PDF modules.
\item[sf\_pdf\_builtin]
  Direct support for selected hadron PDFs.
\item[sf\_lhapdf]
  LHAPDF library support.
\end{description}

\clearpage
%%%%%%%%%%%%%%%%%%%%%%%%%%%%%%%%%%%%%%%%%%%%%%%%%%%%%%%%%%%%%%%%%%%%%%%%
\section{Beam structure}

This module stores the beam structure definition as it is declared in
the SINDARIN script.  The structure definition is not analyzed, just
recorded for later use.

We do not capture any numerical parameters, just names of particles and
structure functions.
\nwenddocs{}\nwbegincode{1}\moddef{\code{}beam{\char95}structures.f90\edoc{}}\endmoddef
\LA{}File header\RA{}

module beam_structures
  
\LA{}Use kinds\RA{}
\LA{}Use strings\RA{}
  use io_units
  use format_defs, only: FMT_19
  use unit_tests
  use diagnostics
  use lorentz
  use polarizations

\LA{}Standard module head\RA{}

\LA{}Beam structures: public\RA{}

\LA{}Beam structures: types\RA{}

\LA{}Beam structures: interfaces\RA{}

contains
  
\LA{}Beam structures: procedures\RA{}

\LA{}Beam structures: tests\RA{}

end module beam_structures
\eatline
\nwendcode{}\nwbegindocs{2}\nwdocspar
\subsection{Beam structure elements}
An entry in a beam-structure record consists of a string
that denotes a type of structure function.
\nwenddocs{}\nwbegincode{3}\moddef{Beam structures: types}\endmoddef
  type :: beam_structure_entry_t
     logical :: is_valid = .false.
     type(string_t) :: name
   contains
   \LA{}Beam structures: beam structure entry: TBP\RA{}
  end type beam_structure_entry_t
  
\eatline
\nwendcode{}\nwbegindocs{4}Output.
\nwenddocs{}\nwbegincode{5}\moddef{Beam structures: beam structure entry: TBP}\endmoddef
  procedure :: to_string => beam_structure_entry_to_string
\nwendcode{}\nwbegincode{6}\moddef{Beam structures: procedures}\endmoddef
  function beam_structure_entry_to_string (object) result (string)
    class(beam_structure_entry_t), intent(in) :: object
    type(string_t) :: string
    if (object%is_valid) then
       string = object%name
    else
       string = "none"
    end if
  end function beam_structure_entry_to_string

\eatline
\nwendcode{}\nwbegindocs{7}\nwdocspar
A record in the beam-structure sequence denotes either a
structure-function entry, a pair of such entries, or a pair spectrum.
\nwenddocs{}\nwbegincode{8}\moddef{Beam structures: types}\plusendmoddef
  type :: beam_structure_record_t
     type(beam_structure_entry_t), dimension(:), allocatable :: entry
  end type beam_structure_record_t

\eatline
\nwendcode{}\nwbegindocs{9}\nwdocspar
\subsection{Beam structure type}
The beam-structure object contains the beam particle(s) as simple strings.
The sequence of records indicates the structure functions by name.  No
numerical parameters are stored.
\nwenddocs{}\nwbegincode{10}\moddef{Beam structures: public}\endmoddef
  public :: beam_structure_t
\nwendcode{}\nwbegincode{11}\moddef{Beam structures: types}\plusendmoddef
  type :: beam_structure_t
     private
     integer :: n_beam = 0
     type(string_t), dimension(:), allocatable :: prt
     type(beam_structure_record_t), dimension(:), allocatable :: record
     type(smatrix_t), dimension(:), allocatable :: smatrix
     real(default), dimension(:), allocatable :: pol_f
     real(default), dimension(:), allocatable :: p
     real(default), dimension(:), allocatable :: theta
     real(default), dimension(:), allocatable :: phi
   contains
   \LA{}Beam structures: beam structure: TBP\RA{}
  end type beam_structure_t
  
\eatline
\nwendcode{}\nwbegindocs{12}The finalizer deletes all contents explicitly, so we can continue
with an empty beam record.  (It is not needed for deallocation.)  We
have distinct finalizers for the independent parts of the beam structure.
\nwenddocs{}\nwbegincode{13}\moddef{Beam structures: beam structure: TBP}\endmoddef
  procedure :: final_sf => beam_structure_final_sf
\nwendcode{}\nwbegincode{14}\moddef{Beam structures: procedures}\plusendmoddef
  subroutine beam_structure_final_sf (object)
    class(beam_structure_t), intent(inout) :: object
    if (allocated (object%prt))  deallocate (object%prt)
    if (allocated (object%record))  deallocate (object%record)
    object%n_beam = 0
  end subroutine beam_structure_final_sf
  
\eatline
\nwendcode{}\nwbegindocs{15}Output.  The actual information fits in a single line, therefore we can
provide a {\tt{}to{\char95}string} method.  The {\tt{}show} method also lists the
current values of relevant global variables.
\nwenddocs{}\nwbegincode{16}\moddef{Beam structures: beam structure: TBP}\plusendmoddef
  procedure :: write => beam_structure_write
  procedure :: to_string => beam_structure_to_string
\nwendcode{}\nwbegincode{17}\moddef{Beam structures: procedures}\plusendmoddef
  subroutine beam_structure_write (object, unit)
    class(beam_structure_t), intent(in) :: object
    integer, intent(in), optional :: unit
    integer :: u, i
    u = given_output_unit (unit)
    write (u, "(1x,A,A)")  "Beam structure: ", char (object%to_string ())
    if (allocated (object%smatrix)) then
       do i = 1, size (object%smatrix)
          write (u, "(3x,A,I0,A)") "polarization (beam ", i, "):"
          call object%smatrix(i)%write (u, indent=2)
       end do
    end if
    if (allocated (object%pol_f)) then
       write (u, "(3x,A,F10.7,:,',',F10.7)")  "polarization degree =", &
            object%pol_f
    end if
    if (allocated (object%p)) then
       write (u, "(3x,A," // FMT_19 // ",:,','," // FMT_19 // &
            ")")  "momentum =", object%p
    end if
    if (allocated (object%theta)) then
       write (u, "(3x,A," // FMT_19 // ",:,','," // FMT_19 // &
            ")")  "angle th =", object%theta
    end if
    if (allocated (object%phi)) then
       write (u, "(3x,A," // FMT_19 // ",:,','," // FMT_19 // &
            ")")  "angle ph =", object%phi
    end if
  end subroutine beam_structure_write
  
  function beam_structure_to_string (object, sf_only) result (string)
    class(beam_structure_t), intent(in) :: object
    logical, intent(in), optional :: sf_only
    type(string_t) :: string
    integer :: i, j
    logical :: with_beams
    with_beams = .true.;  if (present (sf_only))  with_beams = .not. sf_only
    select case (object%n_beam)
    case (1)
       if (with_beams) then
          string = object%prt(1)
       else
          string = ""
       end if
    case (2)
       if (with_beams) then
          string = object%prt(1) // ", " // object%prt(2)
       else
          string = ""
       end if
       if (allocated (object%record)) then
          if (size (object%record) > 0) then
             if (with_beams)  string = string // " => "
             do i = 1, size (object%record)
                if (i > 1)  string = string // " => "
                do j = 1, size (object%record(i)%entry)
                   if (j > 1)  string = string // ", "
                   string = string // object%record(i)%entry(j)%to_string ()
                end do
             end do
          end if
       end if
    case default
       string = "[any particles]"
    end select
  end function beam_structure_to_string
    
\eatline
\nwendcode{}\nwbegindocs{18}Initializer: dimension the beam structure record.  Each array
element denotes the number of entries for a record within the
beam-structure sequence.  The number of entries is either one or two,
while the number of records is unlimited.
\nwenddocs{}\nwbegincode{19}\moddef{Beam structures: beam structure: TBP}\plusendmoddef
  procedure :: init_sf => beam_structure_init_sf
\nwendcode{}\nwbegincode{20}\moddef{Beam structures: procedures}\plusendmoddef
  subroutine beam_structure_init_sf (beam_structure, prt, dim_array)
    class(beam_structure_t), intent(inout) :: beam_structure
    type(string_t), dimension(:), intent(in) :: prt
    integer, dimension(:), intent(in), optional :: dim_array
    integer :: i
    call beam_structure%final_sf ()
    beam_structure%n_beam = size (prt)
    allocate (beam_structure%prt (size (prt)))
    beam_structure%prt = prt
    if (present (dim_array)) then
       allocate (beam_structure%record (size (dim_array)))
       do i = 1, size (dim_array)
          allocate (beam_structure%record(i)%entry (dim_array(i)))
       end do
    else
       allocate (beam_structure%record (0))
    end if
  end subroutine beam_structure_init_sf
    
\eatline
\nwendcode{}\nwbegindocs{21}Set an entry, specified by record number and entry number.
\nwenddocs{}\nwbegincode{22}\moddef{Beam structures: beam structure: TBP}\plusendmoddef
  procedure :: set_sf => beam_structure_set_sf
\nwendcode{}\nwbegincode{23}\moddef{Beam structures: procedures}\plusendmoddef
  subroutine beam_structure_set_sf (beam_structure, i, j, name)
    class(beam_structure_t), intent(inout) :: beam_structure
    integer, intent(in) :: i, j
    type(string_t), intent(in) :: name
    associate (entry => beam_structure%record(i)%entry(j))
      entry%name = name
      entry%is_valid = .true.
    end associate
  end subroutine beam_structure_set_sf

\eatline
\nwendcode{}\nwbegindocs{24}Expand the beam-structure object.  (i) For a pair spectrum, keep the
entry.  (ii) For a single-particle structure function written as a
single entry, replace this by a record with two entries.
(ii) For a record with two nontrivial entries, separate this into two
records with one trivial entry each.

To achieve this, we need a function that tells us whether an entry is
a spectrum or a structure function.  It returns 0 for a trivial entry,
1 for a single-particle structure function, and 2 for a two-particle
spectrum.
\nwenddocs{}\nwbegincode{25}\moddef{Beam structures: interfaces}\endmoddef
  abstract interface
     function strfun_mode_fun (name) result (n)
       import
       type(string_t), intent(in) :: name
       integer :: n
     end function strfun_mode_fun
  end interface
  
\eatline
\nwendcode{}\nwbegindocs{26}Algorithm: (1) Mark entries as invalid where necessary.  (2) Count
the number of entries that we will need.  (3) Expand and copy
entries to a new record array.  (4) Replace the old array by the new one.
\nwenddocs{}\nwbegincode{27}\moddef{Beam structures: beam structure: TBP}\plusendmoddef
  procedure :: expand => beam_structure_expand
\nwendcode{}\nwbegincode{28}\moddef{Beam structures: procedures}\plusendmoddef
  subroutine beam_structure_expand (beam_structure, strfun_mode)
    class(beam_structure_t), intent(inout) :: beam_structure
    procedure(strfun_mode_fun) :: strfun_mode
    type(beam_structure_record_t), dimension(:), allocatable :: new
    integer :: n_record, i, j
    if (.not. allocated (beam_structure%record))  return
    do i = 1, size (beam_structure%record)
       associate (entry => beam_structure%record(i)%entry)
         do j = 1, size (entry)
            select case (strfun_mode (entry(j)%name))
            case (0);  entry(j)%is_valid = .false.
            end select
         end do
       end associate
    end do
    n_record = 0
    do i = 1, size (beam_structure%record)
       associate (entry => beam_structure%record(i)%entry)
         select case (size (entry))
         case (1)
            if (entry(1)%is_valid) then
               select case (strfun_mode (entry(1)%name))
               case (1);  n_record = n_record + 2
               case (2);  n_record = n_record + 1
               end select
            end if
         case (2)
            do j = 1, 2
               if (entry(j)%is_valid) then
                  select case (strfun_mode (entry(j)%name))
                  case (1);  n_record = n_record + 1
                  case (2)
                     call beam_structure%write ()
                     call msg_fatal ("Pair spectrum used as &
                          &single-particle structure function")
                  end select
               end if
            end do
         end select
       end associate
    end do
    allocate (new (n_record))
    n_record = 0
    do i = 1, size (beam_structure%record)
       associate (entry => beam_structure%record(i)%entry)
         select case (size (entry))
         case (1)
            if (entry(1)%is_valid) then
               select case (strfun_mode (entry(1)%name))
               case (1)
                  n_record = n_record + 1
                  allocate (new(n_record)%entry (2))
                  new(n_record)%entry(1) = entry(1)
                  n_record = n_record + 1
                  allocate (new(n_record)%entry (2))
                  new(n_record)%entry(2) = entry(1)
               case (2)
                  n_record = n_record + 1
                  allocate (new(n_record)%entry (1))
                  new(n_record)%entry(1) = entry(1)
               end select
            end if
         case (2)
            do j = 1, 2
               if (entry(j)%is_valid) then
                  n_record = n_record + 1
                  allocate (new(n_record)%entry (2))
                  new(n_record)%entry(j) = entry(j)
               end if
            end do
         end select
       end associate
    end do
    call move_alloc (from = new, to = beam_structure%record)
  end subroutine beam_structure_expand
    
\eatline
\nwendcode{}\nwbegindocs{29}\nwdocspar
\subsection{Polarization}
To record polarization, we provide an allocatable array of {\tt{}smatrix}
objects, sparse matrices.  The polarization structure is independent of the
structure-function setup, they are combined only when an actual beam object is
constructed.
\nwenddocs{}\nwbegincode{30}\moddef{Beam structures: beam structure: TBP}\plusendmoddef
  procedure :: final_pol => beam_structure_final_pol
  procedure :: init_pol => beam_structure_init_pol
\nwendcode{}\nwbegincode{31}\moddef{Beam structures: procedures}\plusendmoddef
  subroutine beam_structure_final_pol (beam_structure)
    class(beam_structure_t), intent(inout) :: beam_structure
    if (allocated (beam_structure%smatrix))  deallocate (beam_structure%smatrix)
    if (allocated (beam_structure%pol_f))  deallocate (beam_structure%pol_f)
  end subroutine beam_structure_final_pol
    
  subroutine beam_structure_init_pol (beam_structure, n)
    class(beam_structure_t), intent(inout) :: beam_structure
    integer, intent(in) :: n
    if (allocated (beam_structure%smatrix))  deallocate (beam_structure%smatrix)
    allocate (beam_structure%smatrix (n))
    if (.not. allocated (beam_structure%pol_f)) &
         allocate (beam_structure%pol_f (n), source = 1._default)
  end subroutine beam_structure_init_pol
    
\eatline
\eatline
\nwendcode{}\nwbegindocs{32}Directly copy the spin density matrices.
\nwenddocs{}\nwbegincode{33}\moddef{Beam structures: beam structure: TBP}\plusendmoddef
  procedure :: set_smatrix => beam_structure_set_smatrix
\nwendcode{}\nwbegincode{34}\moddef{Beam structures: procedures}\plusendmoddef
  subroutine beam_structure_set_smatrix (beam_structure, i, smatrix)
    class(beam_structure_t), intent(inout) :: beam_structure
    integer, intent(in) :: i
    type(smatrix_t), intent(in) :: smatrix
    beam_structure%smatrix(i) = smatrix
  end subroutine beam_structure_set_smatrix
  
\eatline
\nwendcode{}\nwbegindocs{35}Initialize one of the spin density matrices manually.
\nwenddocs{}\nwbegincode{36}\moddef{Beam structures: beam structure: TBP}\plusendmoddef
  procedure :: init_smatrix => beam_structure_init_smatrix
\nwendcode{}\nwbegincode{37}\moddef{Beam structures: procedures}\plusendmoddef
  subroutine beam_structure_init_smatrix (beam_structure, i, n_entry)
    class(beam_structure_t), intent(inout) :: beam_structure
    integer, intent(in) :: i
    integer, intent(in) :: n_entry
    call beam_structure%smatrix(i)%init (2, n_entry)
  end subroutine beam_structure_init_smatrix
  
\eatline
\nwendcode{}\nwbegindocs{38}Set a polarization entry.
\nwenddocs{}\nwbegincode{39}\moddef{Beam structures: beam structure: TBP}\plusendmoddef
  procedure :: set_sentry => beam_structure_set_sentry
\nwendcode{}\nwbegincode{40}\moddef{Beam structures: procedures}\plusendmoddef
  subroutine beam_structure_set_sentry &
       (beam_structure, i, i_entry, index, value)
    class(beam_structure_t), intent(inout) :: beam_structure
    integer, intent(in) :: i
    integer, intent(in) :: i_entry
    integer, dimension(:), intent(in) :: index
    complex(default), intent(in) :: value
    call beam_structure%smatrix(i)%set_entry (i_entry, index, value)
  end subroutine beam_structure_set_sentry
  
\eatline
\nwendcode{}\nwbegindocs{41}Set the array of polarization fractions.
\nwenddocs{}\nwbegincode{42}\moddef{Beam structures: beam structure: TBP}\plusendmoddef
  procedure :: set_pol_f => beam_structure_set_pol_f
\nwendcode{}\nwbegincode{43}\moddef{Beam structures: procedures}\plusendmoddef
  subroutine beam_structure_set_pol_f (beam_structure, f)
    class(beam_structure_t), intent(inout) :: beam_structure
    real(default), dimension(:), intent(in) :: f
    if (allocated (beam_structure%pol_f))  deallocate (beam_structure%pol_f)
    allocate (beam_structure%pol_f (size (f)), source = f)
  end subroutine beam_structure_set_pol_f
    
\eatline
\nwendcode{}\nwbegindocs{44}\nwdocspar
\subsection{Beam momenta}
By default, beam momenta are deduced from the {\tt{}sqrts} value or from
the mass of the decaying particle, assuming a c.m.\ setup.  Here we
set them explicitly.
\nwenddocs{}\nwbegincode{45}\moddef{Beam structures: beam structure: TBP}\plusendmoddef
  procedure :: final_mom => beam_structure_final_mom
\nwendcode{}\nwbegincode{46}\moddef{Beam structures: procedures}\plusendmoddef
  subroutine beam_structure_final_mom (beam_structure)
    class(beam_structure_t), intent(inout) :: beam_structure
    if (allocated (beam_structure%p))  deallocate (beam_structure%p)
    if (allocated (beam_structure%theta))  deallocate (beam_structure%theta)
    if (allocated (beam_structure%phi))  deallocate (beam_structure%phi)
  end subroutine beam_structure_final_mom

\eatline
\nwendcode{}\nwbegincode{47}\moddef{Beam structures: beam structure: TBP}\plusendmoddef
  procedure :: set_momentum => beam_structure_set_momentum
  procedure :: set_theta => beam_structure_set_theta
  procedure :: set_phi => beam_structure_set_phi
\nwendcode{}\nwbegincode{48}\moddef{Beam structures: procedures}\plusendmoddef
  subroutine beam_structure_set_momentum (beam_structure, p)
    class(beam_structure_t), intent(inout) :: beam_structure
    real(default), dimension(:), intent(in) :: p
    if (allocated (beam_structure%p))  deallocate (beam_structure%p)
    allocate (beam_structure%p (size (p)), source = p)
  end subroutine beam_structure_set_momentum
    
  subroutine beam_structure_set_theta (beam_structure, theta)
    class(beam_structure_t), intent(inout) :: beam_structure
    real(default), dimension(:), intent(in) :: theta
    if (allocated (beam_structure%theta))  deallocate (beam_structure%theta)
    allocate (beam_structure%theta (size (theta)), source = theta)
  end subroutine beam_structure_set_theta
    
  subroutine beam_structure_set_phi (beam_structure, phi)
    class(beam_structure_t), intent(inout) :: beam_structure
    real(default), dimension(:), intent(in) :: phi
    if (allocated (beam_structure%phi))  deallocate (beam_structure%phi)
    allocate (beam_structure%phi (size (phi)), source = phi)
  end subroutine beam_structure_set_phi
    
\eatline
\eatline
\eatline
\nwendcode{}\nwbegindocs{49}\nwdocspar
\subsection{Get contents}
Look at the incoming particles.  We may also have the case that beam
particles are not specified, but polarization.
\nwenddocs{}\nwbegincode{50}\moddef{Beam structures: beam structure: TBP}\plusendmoddef
  procedure :: is_set => beam_structure_is_set
  procedure :: get_n_beam => beam_structure_get_n_beam
  procedure :: get_prt => beam_structure_get_prt
\nwendcode{}\nwbegincode{51}\moddef{Beam structures: procedures}\plusendmoddef
  function beam_structure_is_set (beam_structure) result (flag)
    class(beam_structure_t), intent(in) :: beam_structure
    logical :: flag
    flag = beam_structure%n_beam > 0 .or. beam_structure%asymmetric ()
  end function beam_structure_is_set

  function beam_structure_get_n_beam (beam_structure) result (n)
    class(beam_structure_t), intent(in) :: beam_structure
    integer :: n
    n = beam_structure%n_beam
  end function beam_structure_get_n_beam

  function beam_structure_get_prt (beam_structure) result (prt)
    class(beam_structure_t), intent(in) :: beam_structure
    type(string_t), dimension(:), allocatable :: prt
    allocate (prt (size (beam_structure%prt)))
    prt = beam_structure%prt
  end function beam_structure_get_prt

\eatline
\eatline
\eatline
\nwendcode{}\nwbegindocs{52}\nwdocspar
Return the number of records.
\nwenddocs{}\nwbegincode{53}\moddef{Beam structures: beam structure: TBP}\plusendmoddef
  procedure :: get_n_record => beam_structure_get_n_record
\nwendcode{}\nwbegincode{54}\moddef{Beam structures: procedures}\plusendmoddef
  function beam_structure_get_n_record (beam_structure) result (n)
    class(beam_structure_t), intent(in) :: beam_structure
    integer :: n
    if (allocated (beam_structure%record)) then
       n = size (beam_structure%record)
    else
       n = 0
    end if
  end function beam_structure_get_n_record
  
\eatline
\nwendcode{}\nwbegindocs{55}Return an array consisting of the beam indices affected by the valid
entries within a record.  After expansion, there should be exactly one
valid entry per record.
\nwenddocs{}\nwbegincode{56}\moddef{Beam structures: beam structure: TBP}\plusendmoddef
  procedure :: get_i_entry => beam_structure_get_i_entry
\nwendcode{}\nwbegincode{57}\moddef{Beam structures: procedures}\plusendmoddef
  function beam_structure_get_i_entry (beam_structure, i) result (i_entry)
    class(beam_structure_t), intent(in) :: beam_structure
    integer, intent(in) :: i
    integer, dimension(:), allocatable :: i_entry
    associate (record => beam_structure%record(i))
      select case (size (record%entry))
      case (1)
         if (record%entry(1)%is_valid) then
            allocate (i_entry (2), source = [1, 2])
         else
            allocate (i_entry (0))
         end if
      case (2)
         if (all (record%entry%is_valid)) then
            allocate (i_entry (2), source = [1, 2])
         else if (record%entry(1)%is_valid) then
            allocate (i_entry (1), source = [1])
         else if (record%entry(2)%is_valid) then
            allocate (i_entry (1), source = [2])
         else
            allocate (i_entry (0))
         end if
      end select
    end associate
  end function beam_structure_get_i_entry
  
\eatline
\nwendcode{}\nwbegindocs{58}Return the name of the first valid entry within a record.  After
expansion, there should be exactly one valid entry per record.
\nwenddocs{}\nwbegincode{59}\moddef{Beam structures: beam structure: TBP}\plusendmoddef
  procedure :: get_name => beam_structure_get_name
\nwendcode{}\nwbegincode{60}\moddef{Beam structures: procedures}\plusendmoddef
  function beam_structure_get_name (beam_structure, i) result (name)
    class(beam_structure_t), intent(in) :: beam_structure
    integer, intent(in) :: i
    type(string_t) :: name
    associate (record => beam_structure%record(i))
      if (record%entry(1)%is_valid) then
         name = record%entry(1)%name
      else if (size (record%entry) == 2) then
         name = record%entry(2)%name
      end if
    end associate
  end function beam_structure_get_name
  
\eatline
\nwendcode{}\nwbegindocs{61}Return true if the beam structure contains a particular structure
function identifier (such as {\tt{}lhapdf}, {\tt{}isr}, etc.)
\nwenddocs{}\nwbegincode{62}\moddef{Beam structures: beam structure: TBP}\plusendmoddef
  procedure :: contains => beam_structure_contains
\nwendcode{}\nwbegincode{63}\moddef{Beam structures: procedures}\plusendmoddef
  function beam_structure_contains (beam_structure, name) result (flag)
    class(beam_structure_t), intent(in) :: beam_structure
    character(*), intent(in) :: name
    logical :: flag
    integer :: i, j
    flag = .false.
    if (allocated (beam_structure%record)) then
       do i = 1, size (beam_structure%record)
          do j = 1, size (beam_structure%record(i)%entry)
             flag = beam_structure%record(i)%entry(j)%name == name
             if (flag)  return
          end do
       end do
    end if
  end function beam_structure_contains

\eatline
\nwendcode{}\nwbegindocs{64}Return polarization data.
\nwenddocs{}\nwbegincode{65}\moddef{Beam structures: beam structure: TBP}\plusendmoddef
  procedure :: polarized => beam_structure_polarized
  procedure :: get_smatrix => beam_structure_get_smatrix
  procedure :: get_pol_f => beam_structure_get_pol_f
  procedure :: asymmetric => beam_structure_asymmetric
\nwendcode{}\nwbegincode{66}\moddef{Beam structures: procedures}\plusendmoddef
  function beam_structure_polarized (beam_structure) result (flag)
    class(beam_structure_t), intent(in) :: beam_structure
    logical :: flag
    flag = allocated (beam_structure%smatrix)
  end function beam_structure_polarized
  
  function beam_structure_get_smatrix (beam_structure) result (smatrix)
    class(beam_structure_t), intent(in) :: beam_structure
    type(smatrix_t), dimension(:), allocatable :: smatrix
    allocate (smatrix (size (beam_structure%smatrix)), &
         source = beam_structure%smatrix)
  end function beam_structure_get_smatrix
  
  function beam_structure_get_pol_f (beam_structure) result (pol_f)
    class(beam_structure_t), intent(in) :: beam_structure
    real(default), dimension(:), allocatable :: pol_f
    allocate (pol_f (size (beam_structure%pol_f)), &
         source = beam_structure%pol_f)
  end function beam_structure_get_pol_f
  
  function beam_structure_asymmetric (beam_structure) result (flag)
    class(beam_structure_t), intent(in) :: beam_structure
    logical :: flag
    flag = allocated (beam_structure%p) &
         .or. allocated (beam_structure%theta) &
         .or. allocated (beam_structure%phi)
  end function beam_structure_asymmetric
  
\eatline
\eatline
\eatline
\eatline
\nwendcode{}\nwbegindocs{67}Return the beam momenta (the space part, i.e., three-momenta).  This
is meaningful only if momenta and, optionally, angles have been set.
\nwenddocs{}\nwbegincode{68}\moddef{Beam structures: beam structure: TBP}\plusendmoddef
  procedure :: get_momenta => beam_structure_get_momenta
\nwendcode{}\nwbegincode{69}\moddef{Beam structures: procedures}\plusendmoddef
  function beam_structure_get_momenta (beam_structure) result (p)
    class(beam_structure_t), intent(in) :: beam_structure
    type(vector3_t), dimension(:), allocatable :: p
    real(default), dimension(:), allocatable :: theta, phi
    integer :: n, i
    if (allocated (beam_structure%p)) then
       n = size (beam_structure%p)
       if (allocated (beam_structure%theta)) then
          if (size (beam_structure%theta) == n) then
             allocate (theta (n), source = beam_structure%theta)
          else
             call msg_fatal ("Beam structure: mismatch in momentum vs. &
                  &angle theta specification")
          end if
       else
          allocate (theta (n), source = 0._default)
       end if
       if (allocated (beam_structure%phi)) then
          if (size (beam_structure%phi) == n) then
             allocate (phi (n), source = beam_structure%phi)
          else
             call msg_fatal ("Beam structure: mismatch in momentum vs. &
                  &angle phi specification")
          end if
       else
          allocate (phi (n), source = 0._default)
       end if
       allocate (p (n))
       do i = 1, n
          p(i) = beam_structure%p(i) * vector3_moving ([ &
               sin (theta(i)) * cos (phi(i)), &
               sin (theta(i)) * sin (phi(i)), &
               cos (theta(i))])
       end do
       if (n == 2)  p(2) = - p(2)
    else
       call msg_fatal ("Beam structure: angle theta/phi specified but &
            &momentum/a p undefined")
    end if
  end function beam_structure_get_momenta
    
\eatline
\nwendcode{}\nwbegindocs{70}\nwdocspar
\subsection{Unit Tests}
\nwenddocs{}\nwbegincode{71}\moddef{Beam structures: public}\plusendmoddef
  public :: beam_structures_test
\nwendcode{}\nwbegincode{72}\moddef{Beam structures: tests}\endmoddef
  subroutine beam_structures_test (u, results)
    integer, intent(in) :: u
    type(test_results_t), intent(inout) :: results
  \LA{}Beam structures: execute tests\RA{}
end subroutine beam_structures_test
  
\eatline
\nwendcode{}\nwbegindocs{73}\nwdocspar
\subsubsection{Empty structure}
\nwenddocs{}\nwbegincode{74}\moddef{Beam structures: execute tests}\endmoddef
  call test (beam_structures_1, "beam_structures_1", &
       "empty beam structure record", &
       u, results)
\nwendcode{}\nwbegincode{75}\moddef{Beam structures: tests}\plusendmoddef
  subroutine beam_structures_1 (u)
    integer, intent(in) :: u
    type(beam_structure_t) :: beam_structure
    
    write (u, "(A)")  "* Test output: beam_structures_1"
    write (u, "(A)")  "*   Purpose: display empty beam structure record"
    write (u, "(A)")

    call beam_structure%write (u)
    
    write (u, "(A)")
    write (u, "(A)")  "* Test output end: beam_structures_1"
    
  end subroutine beam_structures_1
  
\eatline
\nwendcode{}\nwbegindocs{76}\nwdocspar
\subsubsection{Nontrivial configurations}
\nwenddocs{}\nwbegincode{77}\moddef{Beam structures: execute tests}\plusendmoddef
  call test (beam_structures_2, "beam_structures_2", &
       "beam structure records", &
       u, results)
\nwendcode{}\nwbegincode{78}\moddef{Beam structures: tests}\plusendmoddef
  subroutine beam_structures_2 (u)
    integer, intent(in) :: u
    type(beam_structure_t) :: beam_structure
    integer, dimension(0) :: empty_array
    type(string_t) :: s
    
    write (u, "(A)")  "* Test output: beam_structures_2"
    write (u, "(A)")  "*   Purpose: setup beam structure records"
    write (u, "(A)")

    s = "s"

    call beam_structure%init_sf ([s], empty_array)
    call beam_structure%write (u)
    
    write (u, "(A)")

    call beam_structure%init_sf ([s, s], [1])
    call beam_structure%set_sf (1, 1, var_str ("a"))
    call beam_structure%write (u)
    
    write (u, "(A)")

    call beam_structure%init_sf ([s, s], [2])
    call beam_structure%set_sf (1, 1, var_str ("a"))
    call beam_structure%set_sf (1, 2, var_str ("b"))
    call beam_structure%write (u)
    
    write (u, "(A)")

    call beam_structure%init_sf ([s, s], [2, 1])
    call beam_structure%set_sf (1, 1, var_str ("a"))
    call beam_structure%set_sf (1, 2, var_str ("b"))
    call beam_structure%set_sf (2, 1, var_str ("c"))
    call beam_structure%write (u)

    write (u, "(A)")
    write (u, "(A)")  "* Test output end: beam_structures_2"
    
  end subroutine beam_structures_2
  
\eatline
\nwendcode{}\nwbegindocs{79}\nwdocspar
\subsubsection{Expansion}
Provide a function that tells, for the dummy structure function names
used here, whether they are considered a two-particle spectrum or a
single-particle structure function:
\nwenddocs{}\nwbegincode{80}\moddef{Beam structures: tests}\plusendmoddef
  function test_strfun_mode (name) result (n)
    type(string_t), intent(in) :: name
    integer :: n
    select case (char (name))
    case ("a");  n = 2
    case ("b");  n = 1
    case default;  n = 0
    end select
  end function test_strfun_mode

\eatline
\nwendcode{}\nwbegindocs{81}\nwdocspar
\nwenddocs{}\nwbegincode{82}\moddef{Beam structures: execute tests}\plusendmoddef
  call test (beam_structures_3, "beam_structures_3", &
       "beam structure expansion", &
       u, results)
\nwendcode{}\nwbegincode{83}\moddef{Beam structures: tests}\plusendmoddef
  subroutine beam_structures_3 (u)
    integer, intent(in) :: u
    type(beam_structure_t) :: beam_structure
    type(string_t) :: s
    
    write (u, "(A)")  "* Test output: beam_structures_3"
    write (u, "(A)")  "*   Purpose: expand beam structure records"
    write (u, "(A)")

    s = "s"

    write (u, "(A)")  "* Pair spectrum (keep as-is)"
    write (u, "(A)")

    call beam_structure%init_sf ([s, s], [1])
    call beam_structure%set_sf (1, 1, var_str ("a"))
    call beam_structure%write (u)

    write (u, "(A)")

    call beam_structure%expand (test_strfun_mode)
    call beam_structure%write (u)
    
    write (u, "(A)")
    write (u, "(A)")  "* Structure function pair (expand)"
    write (u, "(A)")

    call beam_structure%init_sf ([s, s], [2])
    call beam_structure%set_sf (1, 1, var_str ("b"))
    call beam_structure%set_sf (1, 2, var_str ("b"))
    call beam_structure%write (u)

    write (u, "(A)")

    call beam_structure%expand (test_strfun_mode)
    call beam_structure%write (u)
    
    write (u, "(A)")
    write (u, "(A)")  "* Structure function (separate and expand)"
    write (u, "(A)")

    call beam_structure%init_sf ([s, s], [1])
    call beam_structure%set_sf (1, 1, var_str ("b"))
    call beam_structure%write (u)

    write (u, "(A)")

    call beam_structure%expand (test_strfun_mode)
    call beam_structure%write (u)
    
    write (u, "(A)")
    write (u, "(A)")  "* Combination"
    write (u, "(A)")

    call beam_structure%init_sf ([s, s], [1, 1])
    call beam_structure%set_sf (1, 1, var_str ("a"))
    call beam_structure%set_sf (2, 1, var_str ("b"))
    call beam_structure%write (u)

    write (u, "(A)")

    call beam_structure%expand (test_strfun_mode)
    call beam_structure%write (u)
    
    write (u, "(A)")
    write (u, "(A)")  "* Test output end: beam_structures_3"
    
  end subroutine beam_structures_3
  
\eatline
\nwendcode{}\nwbegindocs{84}\nwdocspar
\subsubsection{Public methods}
Check the methods that can be called to get the beam-structure
contents.
\nwenddocs{}\nwbegincode{85}\moddef{Beam structures: execute tests}\plusendmoddef
  call test (beam_structures_4, "beam_structures_4", &
       "beam structure contents", &
       u, results)
\nwendcode{}\nwbegincode{86}\moddef{Beam structures: tests}\plusendmoddef
  subroutine beam_structures_4 (u)
    integer, intent(in) :: u
    type(beam_structure_t) :: beam_structure
    type(string_t) :: s
    type(string_t), dimension(2) :: prt
    integer :: i
    
    write (u, "(A)")  "* Test output: beam_structures_4"
    write (u, "(A)")  "*   Purpose: check the API"
    write (u, "(A)")

    s = "s"

    write (u, "(A)")  "* Structure-function combination"
    write (u, "(A)")

    call beam_structure%init_sf ([s, s], [1, 2, 2])
    call beam_structure%set_sf (1, 1, var_str ("a"))
    call beam_structure%set_sf (2, 1, var_str ("b"))
    call beam_structure%set_sf (3, 2, var_str ("c"))
    call beam_structure%write (u)

    write (u, *)
    write (u, "(1x,A,I0)")  "n_beam = ", beam_structure%get_n_beam ()
    prt = beam_structure%get_prt ()
    write (u, "(1x,A,2(1x,A))")  "prt =", char (prt(1)), char (prt(2))
    
    write (u, *)
    write (u, "(1x,A,I0)")  "n_record = ", beam_structure%get_n_record ()

    do i = 1, 3
       write (u, "(A)")
       write (u, "(1x,A,I0,A,A)")  "name(", i, ") = ", &
            char (beam_structure%get_name (i))
       write (u, "(1x,A,I0,A,2(1x,I0))")  "i_entry(", i, ") =", &
            beam_structure%get_i_entry (i)
    end do
    
    write (u, "(A)")
    write (u, "(A)")  "* Test output end: beam_structures_4"
    
  end subroutine beam_structures_4
  
\eatline
\nwendcode{}\nwbegindocs{87}\nwdocspar
\subsubsection{Polarization}
The polarization properties are independent from the structure-function setup.
\nwenddocs{}\nwbegincode{88}\moddef{Beam structures: execute tests}\plusendmoddef
  call test (beam_structures_5, "beam_structures_5", &
       "polarization", &
       u, results)
\nwendcode{}\nwbegincode{89}\moddef{Beam structures: tests}\plusendmoddef
  subroutine beam_structures_5 (u)
    integer, intent(in) :: u
    type(beam_structure_t) :: beam_structure
    integer, dimension(0) :: empty_array
    type(string_t) :: s
    
    write (u, "(A)")  "* Test output: beam_structures_5"
    write (u, "(A)")  "*   Purpose: setup polarization in beam structure records"
    write (u, "(A)")

    s = "s"

    call beam_structure%init_sf ([s], empty_array)
    call beam_structure%init_pol (1)
    call beam_structure%init_smatrix (1, 1)
    call beam_structure%set_sentry (1, 1, [0,0], (1._default, 0._default))
    call beam_structure%set_pol_f ([0.5_default])
    call beam_structure%write (u)
    
    
    write (u, "(A)")
    call beam_structure%final_sf ()
    call beam_structure%final_pol ()

    call beam_structure%init_sf ([s, s], [1])
    call beam_structure%set_sf (1, 1, var_str ("a"))
    call beam_structure%init_pol (2)
    call beam_structure%init_smatrix (1, 2)
    call beam_structure%set_sentry (1, 1, [-1,1], (0.5_default,-0.5_default))
    call beam_structure%set_sentry (1, 2, [ 1,1], (1._default, 0._default))
    call beam_structure%init_smatrix (2, 0)
    call beam_structure%write (u)
    
    write (u, "(A)")
    write (u, "(A)")  "* Test output end: beam_structures_5"
    
  end subroutine beam_structures_5
  
\eatline
\nwendcode{}\nwbegindocs{90}\nwdocspar
\subsubsection{Momenta}
The momenta are independent from the structure-function setup.
\nwenddocs{}\nwbegincode{91}\moddef{Beam structures: execute tests}\plusendmoddef
  call test (beam_structures_6, "beam_structures_6", &
       "momenta", &
       u, results)
\nwendcode{}\nwbegincode{92}\moddef{Beam structures: tests}\plusendmoddef
  subroutine beam_structures_6 (u)
    integer, intent(in) :: u
    type(beam_structure_t) :: beam_structure
    integer, dimension(0) :: empty_array
    type(string_t) :: s
    
    write (u, "(A)")  "* Test output: beam_structures_6"
    write (u, "(A)")  "*   Purpose: setup momenta in beam structure records"
    write (u, "(A)")

    s = "s"

    call beam_structure%init_sf ([s], empty_array)
    call beam_structure%set_momentum ([500._default])
    call beam_structure%write (u)
    
    
    write (u, "(A)")
    call beam_structure%final_sf ()
    call beam_structure%final_mom ()

    call beam_structure%init_sf ([s, s], [1])
    call beam_structure%set_momentum ([500._default, 700._default])
    call beam_structure%set_theta ([0._default, 0.1_default])
    call beam_structure%set_phi ([0._default, 1.51_default])
    call beam_structure%write (u)
    
    write (u, "(A)")
    write (u, "(A)")  "* Test output end: beam_structures_6"
    
  end subroutine beam_structures_6
  
\eatline
\nwendcode{}\nwbegindocs{93}\nwdocspar
%%%%%%%%%%%%%%%%%%%%%%%%%%%%%%%%%%%%%%%%%%%%%%%%%%%%%%%%%%%%%%%%%%%%%%%%
\section{Beams for collisions and decays}

\nwenddocs{}\nwbegincode{94}\moddef{\code{}beams.f90\edoc{}}\endmoddef
\LA{}File header\RA{}

module beams

\LA{}Use kinds\RA{}
\LA{}Use strings\RA{}
  use io_units
  use constants
  use format_defs, only: FMT_19
  use unit_tests
  use diagnostics
  use md5
  use lorentz
  use model_data
  use flavors
  use colors
  use quantum_numbers
  use state_matrices
  use interactions
  use polarizations
  use beam_structures

\LA{}Standard module head\RA{}

\LA{}Beams: public\RA{}

\LA{}Beams: types\RA{}

\LA{}Beams: interfaces\RA{}

contains

\LA{}Beams: procedures\RA{}

\LA{}Beams: tests\RA{}

end module beams
\eatline
\nwendcode{}\nwbegindocs{95}\nwdocspar
\subsection{Beam data}
The beam data type contains beam data for one or two beams, depending
on whether we are dealing with beam collisions or particle decay.  In
addition, it holds the c.m.\ energy {\tt{}sqrts}, the Lorentz
transformation {\tt{}L} that transforms the c.m.\ system into the lab
system, and the pair of c.m.\ momenta.
\nwenddocs{}\nwbegincode{96}\moddef{Beams: public}\endmoddef
  public :: beam_data_t
\nwendcode{}\nwbegincode{97}\moddef{Beams: types}\endmoddef
  type :: beam_data_t
     logical :: initialized = .false.
     integer :: n = 0
     type(flavor_t), dimension(:), allocatable :: flv
     real(default), dimension(:), allocatable :: mass
     type(pmatrix_t), dimension(:), allocatable :: pmatrix
     logical :: lab_is_cm_frame = .true.
     type(vector4_t), dimension(:), allocatable :: p_cm
     type(vector4_t), dimension(:), allocatable :: p
     type(lorentz_transformation_t), allocatable  :: L_cm_to_lab
     real(default) :: sqrts = 0
     character(32) :: md5sum = ""
   contains
   \LA{}Beams: beam data: TBP\RA{}
  end type beam_data_t

\eatline
\nwendcode{}\nwbegindocs{98}Generic initializer.  This is called by the specific initializers
below.  Initialize either for decay or for collision.
\nwenddocs{}\nwbegincode{99}\moddef{Beams: procedures}\endmoddef
  subroutine beam_data_init (beam_data, n)
    type(beam_data_t), intent(out) :: beam_data
    integer, intent(in) :: n
    beam_data%n = n
    allocate (beam_data%flv (n))
    allocate (beam_data%mass (n))
    allocate (beam_data%pmatrix (n))
    allocate (beam_data%p_cm (n))
    allocate (beam_data%p (n))
    beam_data%initialized = .true.
  end subroutine beam_data_init

\eatline
\nwendcode{}\nwbegindocs{100}Finalizer: needed for the polarization components of the beams.
\nwenddocs{}\nwbegincode{101}\moddef{Beams: public}\plusendmoddef
  public :: beam_data_final
\nwendcode{}\nwbegincode{102}\moddef{Beams: procedures}\plusendmoddef
  subroutine beam_data_final (beam_data)
    type(beam_data_t), intent(inout) :: beam_data
    beam_data%initialized = .false.
  end subroutine beam_data_final

\eatline
\nwendcode{}\nwbegindocs{103}The verbose (default) version is for debugging.  The short version
is for screen output in the UI.
\nwenddocs{}\nwbegincode{104}\moddef{Beams: public}\plusendmoddef
  public :: beam_data_write
\nwendcode{}\nwbegincode{105}\moddef{Beams: beam data: TBP}\endmoddef
  procedure :: write => beam_data_write
\nwendcode{}\nwbegincode{106}\moddef{Beams: procedures}\plusendmoddef
  subroutine beam_data_write (beam_data, unit, verbose, write_md5sum)
    class(beam_data_t), intent(in) :: beam_data
    integer, intent(in), optional :: unit
    logical, intent(in), optional :: verbose, write_md5sum
    integer :: prt_name_len
    logical :: verb, write_md5
    integer :: u
    u = given_output_unit (unit);  if (u < 0)  return
    verb = .false.;  if (present (verbose))  verb = verbose
    write_md5 = verb;  if (present (write_md5sum)) write_md5 = write_md5sum
    if (.not. beam_data%initialized) then
       write (u, "(1x,A)") "Beam data: [undefined]"
       return
    end if
    prt_name_len = maxval (len (beam_data%flv%get_name ()))
    select case (beam_data%n)
    case (1)
       write (u, "(1x,A)") "Beam data (decay):"
       if (verb) then
          call write_prt (1)
          call beam_data%pmatrix(1)%write (u)          
          write (u, *) "R.f. momentum:"
          call vector4_write (beam_data%p_cm(1), u)
          write (u, *) "Lab momentum:"
          call vector4_write (beam_data%p(1), u)
       else
          call write_prt (1)
       end if
    case (2)
       write (u, "(1x,A)") "Beam data (collision):"
       if (verb) then
          call write_prt (1)
          call beam_data%pmatrix(1)%write (u)          
          call write_prt (2)
          call beam_data%pmatrix(2)%write (u)
          call write_sqrts
          write (u, *) "C.m. momenta:"
          call vector4_write (beam_data%p_cm(1), u)
          call vector4_write (beam_data%p_cm(2), u)
          write (u, *) "Lab momenta:"
          call vector4_write (beam_data%p(1), u)
          call vector4_write (beam_data%p(2), u)
       else
          call write_prt (1)
          call write_prt (2)
          call write_sqrts
       end if
    end select
    if (allocated (beam_data%L_cm_to_lab)) then
       if (verb) then
          call lorentz_transformation_write (beam_data%L_cm_to_lab, u)
       else
          write (u, "(1x,A)")  "Beam structure: lab and c.m. frame differ"
       end if
    end if
    if (write_md5) then
       write (u, *) "MD5 sum: ", beam_data%md5sum
    end if
  contains
    subroutine write_sqrts
      character(80) :: sqrts_str
      write (sqrts_str, "(" // FMT_19 // ")")  beam_data%sqrts
      write (u, "(3x,A)")  "sqrts = " // trim (adjustl (sqrts_str)) // " GeV"
    end subroutine write_sqrts
    subroutine write_prt (i)
      integer, intent(in) :: i
      character(80) :: name_str, mass_str
      write (name_str, "(A)")  char (beam_data%flv(i)%get_name ())
      write (mass_str, "(ES13.7)")  beam_data%mass(i)
      write (u, "(3x,A)", advance="no") &
           name_str(:prt_name_len) // "  (mass = " &
           // trim (adjustl (mass_str)) // " GeV)"
      if (beam_data%pmatrix(i)%is_polarized ()) then
         write (u, "(2x,A)")  "polarized"
      else
         write (u, *)
      end if
    end subroutine write_prt
  end subroutine beam_data_write

\eatline
\nwendcode{}\nwbegindocs{107}Return initialization status:
\nwenddocs{}\nwbegincode{108}\moddef{Beams: public}\plusendmoddef
  public :: beam_data_are_valid
\nwendcode{}\nwbegincode{109}\moddef{Beams: procedures}\plusendmoddef
  function beam_data_are_valid (beam_data) result (flag)
    logical :: flag
    type(beam_data_t), intent(in) :: beam_data
    flag = beam_data%initialized
  end function beam_data_are_valid

\eatline
\nwendcode{}\nwbegindocs{110}Check whether beam data agree with the current values of relevant
parameters.
\nwenddocs{}\nwbegincode{111}\moddef{Beams: public}\plusendmoddef
  public :: beam_data_check_scattering
\nwendcode{}\nwbegincode{112}\moddef{Beams: procedures}\plusendmoddef
  subroutine beam_data_check_scattering (beam_data, sqrts)
    type(beam_data_t), intent(in) :: beam_data
    real(default), intent(in), optional :: sqrts
    if (beam_data_are_valid (beam_data)) then
       if (present (sqrts)) then
          if (sqrts /= beam_data%sqrts) then
             call msg_error ("Current setting of sqrts is inconsistent " &
                  // "with beam setup (ignored).")
          end if
       end if
    else
       call msg_bug ("Beam setup: invalid beam data")
    end if
  end subroutine beam_data_check_scattering

\eatline
\nwendcode{}\nwbegindocs{113}Return the number of beams (1 for decays, 2 for collisions).
\nwenddocs{}\nwbegincode{114}\moddef{Beams: public}\plusendmoddef
  public :: beam_data_get_n_in
\nwendcode{}\nwbegincode{115}\moddef{Beams: procedures}\plusendmoddef
  function beam_data_get_n_in (beam_data) result (n_in)
    integer :: n_in
    type(beam_data_t), intent(in) :: beam_data
    n_in = beam_data%n
  end function beam_data_get_n_in

\eatline
\nwendcode{}\nwbegindocs{116}Return the beam flavor
\nwenddocs{}\nwbegincode{117}\moddef{Beams: public}\plusendmoddef
  public :: beam_data_get_flavor
\nwendcode{}\nwbegincode{118}\moddef{Beams: procedures}\plusendmoddef
  function beam_data_get_flavor (beam_data) result (flv)
    type(flavor_t), dimension(:), allocatable :: flv
    type(beam_data_t), intent(in) :: beam_data
    allocate (flv (beam_data%n))
    flv = beam_data%flv
  end function beam_data_get_flavor

\eatline
\nwendcode{}\nwbegindocs{119}Return the beam energies
\nwenddocs{}\nwbegincode{120}\moddef{Beams: public}\plusendmoddef
  public :: beam_data_get_energy
\nwendcode{}\nwbegincode{121}\moddef{Beams: procedures}\plusendmoddef
  function beam_data_get_energy (beam_data) result (e)
    real(default), dimension(:), allocatable :: e
    type(beam_data_t), intent(in) :: beam_data
    allocate (e (beam_data%n))
    if (beam_data%initialized) then
       e = energy (beam_data%p)
    else
       e = 0
    end if
  end function beam_data_get_energy

\eatline
\nwendcode{}\nwbegindocs{122}Return the c.m.\ energy.
\nwenddocs{}\nwbegincode{123}\moddef{Beams: public}\plusendmoddef
  public :: beam_data_get_sqrts
\nwendcode{}\nwbegincode{124}\moddef{Beams: procedures}\plusendmoddef
  function beam_data_get_sqrts (beam_data) result (sqrts)
    real(default) :: sqrts
    type(beam_data_t), intent(in) :: beam_data
    sqrts = beam_data%sqrts
  end function beam_data_get_sqrts
  
\eatline
\nwendcode{}\nwbegindocs{125}Return true if the lab and c.m.\ frame are specified as identical.
\nwenddocs{}\nwbegincode{126}\moddef{Beams: public}\plusendmoddef
  public :: beam_data_cm_frame
\nwendcode{}\nwbegincode{127}\moddef{Beams: procedures}\plusendmoddef
  function beam_data_cm_frame (beam_data) result (flag)
    type(beam_data_t), intent(in) :: beam_data
    logical :: flag
    flag = beam_data%lab_is_cm_frame
  end function beam_data_cm_frame
  
\eatline
\nwendcode{}\nwbegindocs{128}Return a MD5 checksum for beam data.  If no checksum is present
(because beams have not been initialized), compute the checksum of the
sqrts value.
\nwenddocs{}\nwbegincode{129}\moddef{Beams: public}\plusendmoddef
  public :: beam_data_get_md5sum
\nwendcode{}\nwbegincode{130}\moddef{Beams: procedures}\plusendmoddef
  function beam_data_get_md5sum (beam_data, sqrts) result (md5sum_beams)
    type(beam_data_t), intent(in) :: beam_data
    real(default), intent(in) :: sqrts
    character(32) :: md5sum_beams
    character(80) :: buffer
    if (beam_data%md5sum /= "") then
       md5sum_beams = beam_data%md5sum
    else
       write (buffer, *)  sqrts
       md5sum_beams = md5sum (buffer)
    end if
  end function beam_data_get_md5sum

\eatline
\nwendcode{}\nwbegindocs{131}\nwdocspar
\subsection{Initializers: beam structure}
Initialize the beam data object from a beam structure object, given energy and
model.
\nwenddocs{}\nwbegincode{132}\moddef{Beams: public}\plusendmoddef
  public :: beam_data_init_structure
\nwendcode{}\nwbegincode{133}\moddef{Beams: procedures}\plusendmoddef
  subroutine beam_data_init_structure &
       (beam_data, structure, sqrts, model, decay_rest_frame)
    type(beam_data_t), intent(out) :: beam_data
    type(beam_structure_t), intent(in) :: structure
    integer :: n_beam
    real(default), intent(in) :: sqrts
    class(model_data_t), intent(in), target :: model
    logical, intent(in), optional :: decay_rest_frame
    type(flavor_t), dimension(:), allocatable :: flv
    n_beam = structure%get_n_beam ()
    allocate (flv (n_beam))
    call flv%init (structure%get_prt (), model)
    if (structure%asymmetric ()) then
       if (structure%polarized ()) then
          call beam_data_init_momenta (beam_data, &
               structure%get_momenta (), flv, &
               structure%get_smatrix (), structure%get_pol_f ())
       else
          call beam_data_init_momenta (beam_data, &
               structure%get_momenta (), flv)
       end if
    else
       select case (n_beam)
       case (1)
          if (structure%polarized ()) then
             call beam_data_init_decay (beam_data, flv, &
                  structure%get_smatrix (), structure%get_pol_f (), &
                  rest_frame = decay_rest_frame)
          else
             call beam_data_init_decay (beam_data, flv, &
                  rest_frame = decay_rest_frame)
          end if
       case (2)
          if (structure%polarized ()) then
             call beam_data_init_sqrts (beam_data, sqrts, flv, &
                  structure%get_smatrix (), structure%get_pol_f ())
          else
             call beam_data_init_sqrts (beam_data, sqrts, flv)
          end if
       case default
          call msg_bug ("Beam data: invalid beam structure object")
       end select
    end if
  end subroutine beam_data_init_structure
    
\eatline
\nwendcode{}\nwbegindocs{134}\nwdocspar
\subsection{Initializers: collisions}
This is the simplest one: just the two flavors, c.m.\ energy,
polarization.  Color is inferred from flavor.  Beam momenta and c.m.\
momenta coincide.
\nwenddocs{}\nwbegincode{135}\moddef{Beams: public}\plusendmoddef
  public :: beam_data_init_sqrts
\nwendcode{}\nwbegincode{136}\moddef{Beams: procedures}\plusendmoddef
  subroutine beam_data_init_sqrts (beam_data, sqrts, flv, smatrix, pol_f)
    type(beam_data_t), intent(out) :: beam_data
    real(default), intent(in) :: sqrts
    type(flavor_t), dimension(:), intent(in) :: flv
    type(smatrix_t), dimension(:), intent(in), optional :: smatrix
    real(default), dimension(:), intent(in), optional :: pol_f
    real(default), dimension(size(flv)) :: E, p
    call beam_data_init (beam_data, size (flv))
    beam_data%sqrts = sqrts
    beam_data%lab_is_cm_frame = .true.
    select case (beam_data%n)
    case (1)
       E = sqrts;  p = 0
       beam_data%p_cm = vector4_moving (E, p, 3)
       beam_data%p = beam_data%p_cm
    case (2)
       beam_data%p_cm = colliding_momenta (sqrts, flv%get_mass ())
       beam_data%p = colliding_momenta (sqrts, flv%get_mass ())
    end select
    call beam_data_finish_initialization (beam_data, flv, smatrix, pol_f)
  end subroutine beam_data_init_sqrts

\eatline
\nwendcode{}\nwbegindocs{137}This version sets beam momenta directly, assuming that they are
asymmetric, i.e., lab frame and c.m.\ frame do not coincide.
Polarization info is deferred to a common initializer.

The Lorentz transformation that we compute here is not actually used
in the calculation; instead, it will be recomputed for each event in
the subroutine {\tt{}phs{\char95}set{\char95}incoming{\char95}momenta}.  We compute it here for
the nominal beam setup nevertheless, so we can print it and, in
particular, include it in the MD5 sum.
\nwenddocs{}\nwbegincode{138}\moddef{Beams: public}\plusendmoddef
  public :: beam_data_init_momenta
\nwendcode{}\nwbegincode{139}\moddef{Beams: procedures}\plusendmoddef
  subroutine beam_data_init_momenta (beam_data, p3, flv, smatrix, pol_f)
    type(beam_data_t), intent(out) :: beam_data
    type(vector3_t), dimension(:), intent(in) :: p3
    type(flavor_t), dimension(:), intent(in) :: flv
    type(smatrix_t), dimension(:), intent(in), optional :: smatrix
    real(default), dimension(:), intent(in), optional :: pol_f
    type(vector4_t) :: p0
    type(vector4_t), dimension(:), allocatable :: p, p_cm_rot
    real(default), dimension(size(p3)) :: e
    real(default), dimension(size(flv)) :: m
    type(lorentz_transformation_t) :: L_boost, L_rot
    call beam_data_init (beam_data, size (flv))
    m = flv%get_mass ()
    e = sqrt (p3 ** 2 + m ** 2)
    allocate (p (beam_data%n))
    p = vector4_moving (e, p3)
    p0 = sum (p)
    beam_data%p = p
    beam_data%lab_is_cm_frame = .false.
    beam_data%sqrts = p0 ** 1
    L_boost = boost (p0, beam_data%sqrts)
    allocate (p_cm_rot (beam_data%n))
    p_cm_rot = inverse (L_boost) * p
    allocate (beam_data%L_cm_to_lab)
    select case (beam_data%n)
    case (1)
       beam_data%L_cm_to_lab = L_boost
       beam_data%p_cm = vector4_at_rest (beam_data%sqrts)
    case (2)
       L_rot = rotation_to_2nd (3, space_part (p_cm_rot(1)))
       beam_data%L_cm_to_lab = L_boost * L_rot
       beam_data%p_cm = &
            colliding_momenta (beam_data%sqrts, flv%get_mass ())
    end select
    call beam_data_finish_initialization (beam_data, flv, smatrix, pol_f)
  end subroutine beam_data_init_momenta
    
\eatline
\nwendcode{}\nwbegindocs{140}\nwdocspar
Final steps:
If requested, rotate the beams in the lab frame, and set
the beam-data components.
\nwenddocs{}\nwbegincode{141}\moddef{Beams: procedures}\plusendmoddef
  subroutine beam_data_finish_initialization (beam_data, flv, smatrix, pol_f)
    type(beam_data_t), intent(inout) :: beam_data
    type(flavor_t), dimension(:), intent(in) :: flv
    type(smatrix_t), dimension(:), intent(in), optional :: smatrix
    real(default), dimension(:), intent(in), optional :: pol_f
    integer :: i
    do i = 1, beam_data%n
       beam_data%flv(i) = flv(i)
       beam_data%mass(i) = flv(i)%get_mass ()
       if (present (smatrix)) then
          if (size (smatrix) /= beam_data%n) &
               call msg_fatal ("Beam data: &
               &polarization density array has wrong dimension")
          beam_data%pmatrix(i) = smatrix(i)
          if (present (pol_f)) then
             if (size (pol_f) /= size (smatrix)) &
                  call msg_fatal ("Beam data: &
                  &polarization fraction array has wrong dimension")
             call beam_data%pmatrix(i)%normalize (flv(i), pol_f(i))
          else
             call beam_data%pmatrix(i)%normalize (flv(i), 1._default)
          end if
       else
          call beam_data%pmatrix(i)%init (2, 0)
          call beam_data%pmatrix(i)%normalize (flv(i), 0._default)
       end if
    end do
    call beam_data_compute_md5sum (beam_data)
  end subroutine beam_data_finish_initialization

\eatline
\nwendcode{}\nwbegindocs{142}\nwdocspar
The MD5 sum is stored within the beam-data record, so it can be
checked for integrity in subsequent runs. 
\nwenddocs{}\nwbegincode{143}\moddef{Beams: procedures}\plusendmoddef
  subroutine beam_data_compute_md5sum (beam_data)
    type(beam_data_t), intent(inout) :: beam_data
    integer :: unit
    unit = free_unit ()
    open (unit = unit, status = "scratch", action = "readwrite")
    call beam_data_write (beam_data, unit, write_md5sum = .false., &
       verbose = .true.)
    rewind (unit)
    beam_data%md5sum = md5sum (unit)
    close (unit)
  end subroutine beam_data_compute_md5sum

\eatline
\nwendcode{}\nwbegindocs{144}\nwdocspar
\subsection{Initializers: decays}
This is the simplest one: decay in rest frame.  We need just flavor
and polarization.  Color is inferred from flavor.  Beam momentum and
c.m.\ momentum coincide.
\nwenddocs{}\nwbegincode{145}\moddef{Beams: public}\plusendmoddef
  public :: beam_data_init_decay
\nwendcode{}\nwbegincode{146}\moddef{Beams: procedures}\plusendmoddef
  subroutine beam_data_init_decay (beam_data, flv, smatrix, pol_f, rest_frame)
    type(beam_data_t), intent(out) :: beam_data
    type(flavor_t), dimension(1), intent(in) :: flv
    type(smatrix_t), dimension(1), intent(in), optional :: smatrix
    real(default), dimension(:), intent(in), optional :: pol_f
    logical, intent(in), optional :: rest_frame
    real(default), dimension(1) :: m
    m = flv%get_mass ()
    if (present (smatrix)) then
       call beam_data_init_sqrts (beam_data, m(1), flv, smatrix, pol_f)
    else
       call beam_data_init_sqrts (beam_data, m(1), flv, smatrix, pol_f)
    end if
    if (present (rest_frame))  beam_data%lab_is_cm_frame = rest_frame
  end subroutine beam_data_init_decay

\eatline
\nwendcode{}\nwbegindocs{147}\nwdocspar
\subsection{Sanity check}
After the beams have been set, the initial-particle masses may have
been modified.  This can be checked here.
\nwenddocs{}\nwbegincode{148}\moddef{Beams: public}\plusendmoddef
  public :: beam_data_masses_are_consistent
\nwendcode{}\nwbegincode{149}\moddef{Beams: procedures}\plusendmoddef
  function beam_data_masses_are_consistent (beam_data) result (flag)
    logical :: flag
    type(beam_data_t), intent(in) :: beam_data
    flag = all (beam_data%mass == beam_data%flv%get_mass ())
  end function beam_data_masses_are_consistent

\eatline
\nwendcode{}\nwbegindocs{150}\nwdocspar
\subsection{The beams type}
Beam objects are interaction objects that contain the actual beam
data including polarization and density matrix.  For collisions, the
beam object actually contains two beams.
\nwenddocs{}\nwbegincode{151}\moddef{Beams: public}\plusendmoddef
  public :: beam_t
\nwendcode{}\nwbegincode{152}\moddef{Beams: types}\plusendmoddef
  type :: beam_t
     private
     type(interaction_t) :: int
  end type beam_t

\eatline
\nwendcode{}\nwbegindocs{153}The constructor contains code that converts beam data into the
(entangled) particle-pair quantum state.  First, we set the number of
particles and polarization mask.  (The polarization mask is handed
over to all later interactions, so if helicity is diagonal or absent, this fact
is used when constructing the hard-interaction events.)  Then, we
construct the entangled state that combines helicity, flavor and color
of the two particles (where flavor and color are unique, while several
helicity states are possible).  Then, we transfer this state together
with the associated values from the spin density matrix into the
{\tt{}interaction{\char95}t} object.
\nwenddocs{}\nwbegincode{154}\moddef{Beams: public}\plusendmoddef
  public :: beam_init
\nwendcode{}\nwbegincode{155}\moddef{Beams: procedures}\plusendmoddef
  subroutine beam_init (beam, beam_data)
    type(beam_t), intent(out) :: beam
    type(beam_data_t), intent(in), target :: beam_data
    type(quantum_numbers_mask_t), dimension(beam_data%n) :: mask
    type(state_matrix_t), target :: state_hel, state_fc, state_tmp
    type(state_iterator_t) :: it_hel, it_tmp
    type(quantum_numbers_t), dimension(:), allocatable :: qn
    type(polarization_t), dimension(:), allocatable :: pol
    integer :: i
    mask = quantum_numbers_mask (.false., .false., &
         .not. beam_data%pmatrix%is_polarized (), &
         mask_hd = beam_data%pmatrix%is_diagonal ())
    call interaction_init &
         (beam%int, 0, 0, beam_data%n, mask=mask, store_values=.true.)
    allocate (pol (beam_data%n))
    do i = 1, size (pol)
       call polarization_init_pmatrix (pol(i), beam_data%pmatrix(i))
    end do
    call combine_polarization_states (pol, state_hel)
    do i = 1, size (pol)
       call polarization_final (pol(i))
    end do
    allocate (qn (beam_data%n))
    call qn%init (beam_data%flv, color_from_flavor (beam_data%flv, 1))
    call state_fc%init ()
    call state_fc%add_state (qn)
    call merge_state_matrices (state_hel, state_fc, state_tmp)
    call it_hel%init (state_hel)
    call it_tmp%init (state_tmp)
    do while (it_hel%is_valid ())
       call interaction_add_state (beam%int, &
            it_tmp%get_quantum_numbers (), &
            value=it_hel%get_matrix_element ())
       call it_hel%advance ()
       call it_tmp%advance ()
    end do
    call interaction_freeze (beam%int)
    call interaction_set_momenta &
         (beam%int, beam_data%p, outgoing = .true.)
    call state_hel%final ()
    call state_fc%final ()
    call state_tmp%final ()
  end subroutine beam_init

\eatline
\nwendcode{}\nwbegindocs{156}Finalizer:
\nwenddocs{}\nwbegincode{157}\moddef{Beams: public}\plusendmoddef
  public :: beam_final
\nwendcode{}\nwbegincode{158}\moddef{Beams: procedures}\plusendmoddef
  subroutine beam_final (beam)
    type(beam_t), intent(inout) :: beam
    call interaction_final (beam%int)
  end subroutine beam_final

\eatline
\nwendcode{}\nwbegindocs{159}I/O:
\nwenddocs{}\nwbegincode{160}\moddef{Beams: public}\plusendmoddef
  public :: beam_write
\nwendcode{}\nwbegincode{161}\moddef{Beams: procedures}\plusendmoddef
  subroutine beam_write (beam, unit, verbose, show_momentum_sum, show_mass)
    type(beam_t), intent(in) :: beam
    integer, intent(in), optional :: unit
    logical, intent(in), optional :: verbose, show_momentum_sum, show_mass
    integer :: u
    u = given_output_unit (unit);  if (u < 0)  return
    select case (interaction_get_n_out (beam%int))
    case (1);  write (u, *) "Decaying particle:"
    case (2);  write (u, *) "Colliding beams:"
    end select
    call interaction_write &
         (beam%int, unit, verbose = verbose, show_momentum_sum = &
            show_momentum_sum, show_mass = show_mass)
  end subroutine beam_write

\eatline
\nwendcode{}\nwbegindocs{162}Defined assignment: deep copy
\nwenddocs{}\nwbegincode{163}\moddef{Beams: public}\plusendmoddef
  public :: assignment(=)
\nwendcode{}\nwbegincode{164}\moddef{Beams: interfaces}\endmoddef
  interface assignment(=)
     module procedure beam_assign
  end interface

\nwendcode{}\nwbegincode{165}\moddef{Beams: procedures}\plusendmoddef
  subroutine beam_assign (beam_out, beam_in)
    type(beam_t), intent(out) :: beam_out
    type(beam_t), intent(in) :: beam_in
    beam_out%int = beam_in%int
  end subroutine beam_assign

\eatline
\nwendcode{}\nwbegindocs{166}\nwdocspar
\subsection{Inherited procedures}
\nwenddocs{}\nwbegincode{167}\moddef{Beams: public}\plusendmoddef
  public :: interaction_set_source_link
\nwendcode{}\nwbegincode{168}\moddef{Beams: interfaces}\plusendmoddef
  interface interaction_set_source_link
     module procedure interaction_set_source_link_beam
  end interface
\nwendcode{}\nwbegincode{169}\moddef{Beams: procedures}\plusendmoddef
  subroutine interaction_set_source_link_beam (int, i, beam1, i1)
    type(interaction_t), intent(inout) :: int
    type(beam_t), intent(in), target :: beam1
    integer, intent(in) :: i, i1
    call interaction_set_source_link (int, i, beam1%int, i1)
  end subroutine interaction_set_source_link_beam

\eatline
\nwendcode{}\nwbegindocs{170}\nwdocspar
\subsection{Accessing contents}
Return the interaction component -- as a pointer, to avoid any copying.
\nwenddocs{}\nwbegincode{171}\moddef{Beams: public}\plusendmoddef
  public :: beam_get_int_ptr
\nwendcode{}\nwbegincode{172}\moddef{Beams: procedures}\plusendmoddef
  function beam_get_int_ptr (beam) result (int)
    type(interaction_t), pointer :: int
    type(beam_t), intent(in), target :: beam
    int => beam%int
  end function beam_get_int_ptr

\eatline
\nwendcode{}\nwbegindocs{173}Set beam momenta directly.  (Used for cascade decays.)
\nwenddocs{}\nwbegincode{174}\moddef{Beams: public}\plusendmoddef
  public :: beam_set_momenta
\nwendcode{}\nwbegincode{175}\moddef{Beams: procedures}\plusendmoddef
  subroutine beam_set_momenta (beam, p)
    type(beam_t), intent(inout) :: beam
    type(vector4_t), dimension(:), intent(in) :: p
    call interaction_set_momenta (beam%int, p)
  end subroutine beam_set_momenta

\eatline
\nwendcode{}\nwbegindocs{176}\subsection{Unit tests}
\nwenddocs{}\nwbegincode{177}\moddef{Beams: public}\plusendmoddef
  public :: beam_test
\nwendcode{}\nwbegincode{178}\moddef{Beams: procedures}\plusendmoddef
  subroutine beam_test (u, results)
    integer, intent(in) :: u
    type(test_results_t), intent(inout) :: results
  \LA{}Beams: execute tests\RA{}
  end subroutine beam_test

\nwendcode{}\nwbegindocs{179} %def beam_test
\nwenddocs{}\nwbegindocs{180}Test the basic beam setup.
\nwenddocs{}\nwbegincode{181}\moddef{Beams: execute tests}\endmoddef
  call test (beam_1, "beam_1", &
       "check basic beam setup", &
       u, results)
\nwendcode{}\nwbegincode{182}\moddef{Beams: tests}\endmoddef
  subroutine beam_1 (u)
    integer, intent(in) :: u
    type(beam_data_t), target :: beam_data
    type(beam_t) :: beam
    real(default) :: sqrts
    type(flavor_t), dimension(2) :: flv
    type(smatrix_t), dimension(2) :: smatrix
    real(default), dimension(2) :: pol_f
    type(model_data_t), target :: model
    
    write (u, "(A)")  "* Test output: beam_1"
    write (u, "(A)")  "*   Purpose: test basic beam setup"
    write (u, "(A)")      
        
    write (u, "(A)")  "* Reading model file"
    write (u, "(A)") 

    call reset_interaction_counter ()

    call model%init_sm_test ()

    write (u, "(A)")  "* 1: Scattering process"
    write (u, "(A)")
    
    sqrts = 500
    call flv%init ([1,-1], model)

    call smatrix(1)%init (2, 1)
    call smatrix(1)%set_entry (1, [1,1], (1._default, 0._default))
    pol_f(1) = 0.5_default

    !!! 2.1 version:
    ! call polarization_init_circular (pol(1), flv(1), 0.5_default)

    call smatrix(2)%init (2, 3)
    call smatrix(2)%set_entry (1, [1,1], (1._default, 0._default))
    call smatrix(2)%set_entry (2, [-1,-1], (1._default, 0._default))
    call smatrix(2)%set_entry (3, [-1,1], (1._default, 0._default))
    pol_f(2) = 1._default

    !!! 2.1 version:
    ! call polarization_init_transversal (pol(2), flv(2), 0._default, 1._default)
    call beam_data_init_sqrts (beam_data, sqrts, flv, smatrix, pol_f)
    call beam_data_write (beam_data, u)
    write (u, "(A)")
    call beam_init (beam, beam_data)
    call beam_write (beam, u)
    call beam_final (beam)
    call beam_data_final (beam_data)
    
    write (u, "(A)")
    write (u, "(A)")  "* 2: Decay"
    write (u, "(A)")
    call flv(1)%init (23, model)
    call smatrix(1)%init (2, 1)
    call smatrix(1)%set_entry (1, [0,0], (1._default, 0._default))
    pol_f(1) = 0.4_default

    !!! 2.1 version:
    ! call polarization_init_longitudinal (pol(1), flv(1), 0.4_default)
    call beam_data_init_decay (beam_data, flv(1:1), smatrix(1:1), pol_f(1:1))
    call beam_data_write (beam_data, u)
    write (u, "(A)")
    call beam_init (beam, beam_data)
    call beam_write (beam, u)
    
    write (u, "(A)")
    write (u, "(A)")  "* Cleanup"            
       
    call beam_final (beam)
    call beam_data_final (beam_data)

    call model%final ()

    write (u, "(A)")
    write (u, "(A)")  "* Test output end: beam_1"        
    
  end subroutine beam_1

\eatline
\nwendcode{}\nwbegindocs{183}Test advanced beam setup.
\nwenddocs{}\nwbegincode{184}\moddef{Beams: execute tests}\plusendmoddef
  call test (beam_2, "beam_2", &
       "beam initialization", &
       u, results)
\nwendcode{}\nwbegincode{185}\moddef{Beams: tests}\plusendmoddef
  subroutine beam_2 (u)
    integer, intent(in) :: u
    type(beam_data_t), target :: beam_data
    type(beam_t) :: beam
    real(default) :: sqrts
    type(flavor_t), dimension(2) :: flv
    integer, dimension(0) :: no_records
    type(beam_structure_t) :: beam_structure
    type(model_data_t), target :: model
    
    write (u, "(A)")  "* Test output: beam_2"
    write (u, "(A)")  "*   Purpose: transfer beam polarization using &
         &beam structure"
    write (u, "(A)")      
        
    write (u, "(A)")  "* Reading model file"
    write (u, "(A)") 

    call reset_interaction_counter ()

    call model%init_sm_test ()

    write (u, "(A)")  "* 1: Scattering process"
    write (u, "(A)")
    
    sqrts = 500
    call flv%init ([1,-1], model)
    call beam_structure%init_sf (flv%get_name (), no_records)

    call beam_structure%init_pol (2)

    call beam_structure%init_smatrix (1, 1)
    call beam_structure%set_sentry (1, 1, [1,1], (1._default, 0._default))

    call beam_structure%init_smatrix (2, 3)
    call beam_structure%set_sentry (2, 1, [1,1], (1._default, 0._default))
    call beam_structure%set_sentry (2, 2, [-1,-1], (1._default, 0._default))
    call beam_structure%set_sentry (2, 3, [-1,1], (1._default, 0._default))

    call beam_structure%set_pol_f ([0.5_default, 1._default])
    call beam_structure%write (u)
    write (u, *)
    
    call beam_data_init_structure (beam_data, beam_structure, sqrts, model)
    call beam_data_write (beam_data, u)
    write (u, *)
 
    call beam_init (beam, beam_data)
    call beam_write (beam, u)

    call beam_final (beam)
    call beam_data_final (beam_data)
    call beam_structure%final_pol ()
    call beam_structure%final_sf ()
    
    write (u, "(A)")
    write (u, "(A)")  "* 2: Decay"
    write (u, "(A)")

    call flv(1)%init (23, model)
    call beam_structure%init_sf ([flv(1)%get_name ()], no_records)

    call beam_structure%init_pol (1)

    call beam_structure%init_smatrix (1, 1)
    call beam_structure%set_sentry (1, 1, [0,0], (1._default, 0._default))
    call beam_structure%set_pol_f ([0.4_default])
    call beam_structure%write (u)
    write (u, *)
    
    call beam_data_init_structure (beam_data, beam_structure, sqrts, model)
    call beam_data_write (beam_data, u)
    write (u, "(A)")
    call beam_init (beam, beam_data)
    call beam_write (beam, u)
    
    write (u, "(A)")
    write (u, "(A)")  "* Cleanup"            
       
    call beam_final (beam)
    call beam_data_final (beam_data)

    call model%final ()

    write (u, "(A)")
    write (u, "(A)")  "* Test output end: beam_2"        
    
  end subroutine beam_2

\eatline
\nwendcode{}\nwbegindocs{186}Test advanced beam setup, completely arbitrary momenta.
\nwenddocs{}\nwbegincode{187}\moddef{Beams: execute tests}\plusendmoddef
  call test (beam_3, "beam_3", &
       "generic beam momenta", &
       u, results)
\nwendcode{}\nwbegincode{188}\moddef{Beams: tests}\plusendmoddef
  subroutine beam_3 (u)
    integer, intent(in) :: u
    type(beam_data_t), target :: beam_data
    type(beam_t) :: beam
    type(flavor_t), dimension(2) :: flv
    integer, dimension(0) :: no_records
    type(model_data_t), target :: model
    type(beam_structure_t) :: beam_structure
    type(vector3_t), dimension(2) :: p3
    type(vector4_t), dimension(2) :: p
    
    write (u, "(A)")  "* Test output: beam_3"
    write (u, "(A)")  "*   Purpose: set up beams with generic momenta"
    write (u, "(A)")      
        
    write (u, "(A)")  "* Reading model file"
    write (u, "(A)") 

    call reset_interaction_counter ()

    call model%init_sm_test ()

    write (u, "(A)")  "* 1: Scattering process"
    write (u, "(A)")
    
    call flv%init ([2212,2212], model)

    p3(1) = vector3_moving ([5._default, 0._default, 10._default])
    p3(2) = -vector3_moving ([1._default, 1._default, -10._default])

    call beam_structure%init_sf (flv%get_name (), no_records)
    call beam_structure%set_momentum (p3 ** 1)
    call beam_structure%set_theta (polar_angle (p3))
    call beam_structure%set_phi (azimuthal_angle (p3))
    call beam_structure%write (u)
    write (u, *)

    call beam_data_init_structure (beam_data, beam_structure, 0._default, model)
    call beam_data_write (beam_data, u, verbose = .true.)
    write (u, *)
 
    write (u, "(1x,A)")  "Beam momenta reconstructed from LT:"
    p = beam_data%L_cm_to_lab * beam_data%p_cm
    call pacify (p, 1e-12_default)
    call vector4_write (p(1), u)
    call vector4_write (p(2), u)
    write (u, "(A)")

    call beam_init (beam, beam_data)
    call beam_write (beam, u)

    call beam_final (beam)
    call beam_data_final (beam_data)
    call beam_structure%final_sf ()
    call beam_structure%final_mom ()
    
    write (u, "(A)")
    write (u, "(A)")  "* 2: Decay"
    write (u, "(A)")

    call flv(1)%init (23, model)
    p3(1) = vector3_moving ([10._default, 5._default, 50._default])
    
    call beam_structure%init_sf ([flv(1)%get_name ()], no_records)
    call beam_structure%set_momentum ([p3(1) ** 1])
    call beam_structure%set_theta ([polar_angle (p3(1))])
    call beam_structure%set_phi ([azimuthal_angle (p3(1))])
    call beam_structure%write (u)
    write (u, *)

    call beam_data_init_structure (beam_data, beam_structure, 0._default, model)
    call beam_data_write (beam_data, u, verbose = .true.)
    write (u, "(A)")

    write (u, "(1x,A)")  "Beam momentum reconstructed from LT:"
    p(1) = beam_data%L_cm_to_lab * beam_data%p_cm(1)
    call pacify (p(1), 1e-12_default)
    call vector4_write (p(1), u)
    write (u, "(A)")

    call beam_init (beam, beam_data)
    call beam_write (beam, u)
    
    write (u, "(A)")
    write (u, "(A)")  "* Cleanup"            
       
    call beam_final (beam)
    call beam_data_final (beam_data)
    call beam_structure%final_sf ()
    call beam_structure%final_mom ()

    call model%final ()

    write (u, "(A)")
    write (u, "(A)")  "* Test output end: beam_3"        
    
  end subroutine beam_3

\eatline
\nwendcode{}\nwbegindocs{189}\nwdocspar
%%%%%%%%%%%%%%%%%%%%%%%%%%%%%%%%%%%%%%%%%%%%%%%%%%%%%%%%%%%%%%%%%%%%%%%%
\section{Tools}

This module contains auxiliary procedures that can be accessed by the
structure function code.

\nwenddocs{}\nwbegincode{190}\moddef{\code{}sf{\char95}aux.f90\edoc{}}\endmoddef
\LA{}File header\RA{}

module sf_aux

\LA{}Use kinds\RA{}
  use io_units
  use constants, only: twopi
  use unit_tests

  use lorentz

\LA{}Standard module head\RA{}

\LA{}SF aux: public\RA{}

\LA{}SF aux: parameters\RA{}

\LA{}SF aux: types\RA{}

contains

\LA{}SF aux: procedures\RA{}

\LA{}SF aux: tests\RA{}

end module sf_aux

\eatline
\nwendcode{}\nwbegindocs{191}\nwdocspar
\subsection{Momentum splitting}
Let us consider first an incoming parton with momentum $k$ and
invariant mass squared $s=k^2$ that splits into two partons with
momenta $q,p$ and invariant masses $t=q^2$ and $u=p^2$.  (This is an
abuse of the Mandelstam notation.  $t$ is actually the momentum
transfer, assuming that $p$ is radiated and $q$ initiates the hard
process.) The energy is split among the partons such that if $E=k^0$,
we have $q^0 = xE$ and $p^0=\bar x E$, where $\bar x\equiv 1-x$.

We define the angle $\theta$ as the polar angle of $p$ w.r.t.\ the
momentum axis of the incoming momentum $k$.  Ignoring azimuthal angle,
we can write the four-momenta in the basis $(E,p_T,p_L)$ as
\begin{equation}
  k = 
  \begin{pmatrix}
    E \\ 0 \\ p
  \end{pmatrix},
  \qquad
  p =
  \begin{pmatrix}
    \bar x E \\ \bar x\bar p\sin\theta \\ \bar x\bar p\cos\theta
  \end{pmatrix},
  \qquad
  q =
  \begin{pmatrix}
    x E \\ -\bar x\bar p\sin\theta \\ p - \bar x\bar p\cos\theta
  \end{pmatrix},
\end{equation}
where the first two mass-shell conditions are
\begin{equation}
  p^2 = E^2 - s,
  \qquad
  \bar p^2 = E^2 - \frac{u}{\bar x^2}.
\end{equation}
The second condition implies that, for positive $u$, $\bar x^2 >
u/E^2$, or equivalently
\begin{equation}
  x < 1 - \sqrt{u} / E.
\end{equation}

We are interested in the third mass-shell conditions: $s$ and $u$ are
fixed, so we need $t$ as a function of $\cos\theta$:
\begin{equation}
  t = -2\bar x \left(E^2 - p\bar p\cos\theta\right) + s + u.
\end{equation}
Solving for $\cos\theta$, we get
\begin{equation}
  \cos\theta = \frac{2\bar x E^2 + t - s - u}{2\bar x p\bar p}.
\end{equation}
We can compute $\sin\theta$ numerically as
$\sin^2\theta=1-\cos^2\theta$, but it is important to reexpress this
in view of numerical stability.  To this end, we first determine the
bounds for $t$.  The cosine must be between $-1$ and $1$, so the
bounds are
\begin{align}
  t_0 &= -2\bar x\left(E^2 + p\bar p\right) + s + u,
\\
  t_1 &= -2\bar x\left(E^2 - p\bar p\right) + s + u.
\end{align}
Computing $\sin^2\theta$ from $\cos\theta$ above, we observe that the
numerator is a quadratic polynomial in $t$ which has the zeros $t_0$
and $t_1$, while the common denominator is given by $(2\bar x p\bar
p)^2$.  Hence, we can write
\begin{equation}
  \sin^2\theta = -\frac{(t - t_0)(t - t_1)}{(2\bar x p\bar p)^2}
  \qquad\text{and}\qquad
  \cos\theta = \frac{(t-t_0) + (t-t_1)}{4\bar x p\bar p},
\end{equation}
which is free of large cancellations near $t=t_0$ or $t=t_1$.

If all is massless, i.e., $s=u=0$, this simplifies to
\begin{align}
  t_0 &= -4\bar x E^2, 
&
  t_1 &= 0,
\\
  \sin^2\theta &= -\frac{t}{\bar x E^2}
                   \left(1 + \frac{t}{4\bar x E^2}\right),
&
  \cos\theta &= 1 + \frac{t}{2\bar x E^2}.
\end{align}

Here is the implementation.  First, we define a container for the
kinematical integration limits and some further data.
\nwenddocs{}\nwbegincode{192}\moddef{SF aux: public}\endmoddef
  public :: splitting_data_t
\nwendcode{}\nwbegincode{193}\moddef{SF aux: types}\endmoddef
  type :: splitting_data_t
     private
     logical :: collinear = .false.
     real(default) :: x0 = 0
     real(default) :: x1
     real(default) :: t0
     real(default) :: t1
     real(default) :: phi0 = 0
     real(default) :: phi1 = twopi
     real(default) :: E, p, s, u, m2
     real(default) :: x, xb, pb
     real(default) :: t = 0
     real(default) :: phi = 0
   contains
   \LA{}SF aux: splitting data: TBP\RA{}
  end type splitting_data_t

\eatline
\nwendcode{}\nwbegindocs{194}I/O for debugging:
\nwenddocs{}\nwbegincode{195}\moddef{SF aux: splitting data: TBP}\endmoddef
  procedure :: write => splitting_data_write
\nwendcode{}\nwbegincode{196}\moddef{SF aux: procedures}\endmoddef
  subroutine splitting_data_write (d, unit)
    class(splitting_data_t), intent(in) :: d
    integer, intent(in), optional :: unit
    integer :: u
    u = given_output_unit (unit);  if (u < 0)  return
    write (u, "(A)") "Splitting data:"
    write (u, "(2x,A,L1)")  "collinear = ", d%collinear
1   format (2x,A,1x,ES15.8)
    write (u, 1) "x0   =", d%x0
    write (u, 1) "x    =", d%x
    write (u, 1) "xb   =", d%xb
    write (u, 1) "x1   =", d%x1
    write (u, 1) "t0   =", d%t0
    write (u, 1) "t    =", d%t
    write (u, 1) "t1   =", d%t1
    write (u, 1) "phi0 =", d%phi0
    write (u, 1) "phi  =", d%phi
    write (u, 1) "phi1 =", d%phi1
    write (u, 1) "E    =", d%E
    write (u, 1) "p    =", d%p
    write (u, 1) "pb   =", d%pb
    write (u, 1) "s    =", d%s
    write (u, 1) "u    =", d%u
    write (u, 1) "m2   =", d%m2
  end subroutine splitting_data_write

\eatline
\nwendcode{}\nwbegindocs{197}\nwdocspar
\subsection{Constant data}
This is the initializer for the data.  The input consists of the
incoming momentum, its invariant mass squared, and the invariant mass
squared of the radiated particle.  $m2$ is the \emph{physical} mass
squared of the outgoing particle.  The $t$ bounds depend on the chosen $x$
value and cannot be determined yet.
\nwenddocs{}\nwbegincode{198}\moddef{SF aux: splitting data: TBP}\plusendmoddef
  procedure :: init => splitting_data_init
\nwendcode{}\nwbegincode{199}\moddef{SF aux: procedures}\plusendmoddef
  subroutine splitting_data_init (d, k, mk2, mr2, mo2, collinear)
    class(splitting_data_t), intent(out) :: d
    type(vector4_t), intent(in) :: k
    real(default), intent(in) :: mk2, mr2, mo2
    logical, intent(in), optional :: collinear
    if (present (collinear))  d%collinear = collinear
    d%E = energy (k)
    d%x1 = 1 - sqrt (max (mr2, 0._default)) / d%E
    d%p = sqrt (d%E**2 - mk2)
    d%s = mk2
    d%u = mr2
    d%m2 = mo2
  end subroutine splitting_data_init

\eatline
\nwendcode{}\nwbegindocs{200}Retrieve the $x$ bounds, if needed for $x$ sampling.  Generating an
$x$ value is done by the caller, since this is the part that depends
on the nature of the structure function.
\nwenddocs{}\nwbegincode{201}\moddef{SF aux: splitting data: TBP}\plusendmoddef
  procedure :: get_x_bounds => splitting_get_x_bounds
\nwendcode{}\nwbegincode{202}\moddef{SF aux: procedures}\plusendmoddef
  function splitting_get_x_bounds (d) result (x)
    class(splitting_data_t), intent(in) :: d
    real(default), dimension(2) :: x
    x = [ d%x0, d%x1 ]
  end function splitting_get_x_bounds

\eatline
\nwendcode{}\nwbegindocs{203}Now set the momentum fraction and compute $t_0$ and $t_1$.

[The calculation of $t_1$ is subject to numerical problems.  The exact
formula is ($s=m_i^2$, $u=m_r^2$)
\begin{equation}
  t_1 = -2\bar x E^2 + m_i^2 + m_r^2 
  + 2\bar x \sqrt{E^2-m_i^2}\,\sqrt{E^2 - m_r^2/\bar x^2}.
\end{equation}
The structure-function paradigm is useful only if $E\gg m_i,m_r$.  In
a Taylor expansion for large $E$, the leading term cancels.  The
expansion of the square roots (to subleading order) yields
\begin{equation}
  t_1 = xm_i^2 - \frac{x}{\bar x}m_r^2.
\end{equation}
There are two cases of interest: $m_i=m_o$ and $m_r=0$,
\begin{equation}
  t_1 = xm_o^2
\end{equation}
and $m_i=m_r$ and $m_o=0$,
\begin{equation}
  t_1 = -\frac{x^2}{\bar x}m_i^2.
\end{equation}
In both cases, $t_1\leq m_o^2$.]

That said, it turns out that taking the $t_1$ evaluation at face value
leads to less problems than the approximation.  We express the angles
in terms of $t-t_0$ and $t-t_1$.  Numerical noise in $t_1$ can then be
tolerated.
\nwenddocs{}\nwbegincode{204}\moddef{SF aux: splitting data: TBP}\plusendmoddef
  procedure :: set_t_bounds => splitting_set_t_bounds
\nwendcode{}\nwbegincode{205}\moddef{SF aux: procedures}\plusendmoddef
  elemental subroutine splitting_set_t_bounds (d, x, xb)
    class(splitting_data_t), intent(inout) :: d
    real(default), intent(in), optional :: x, xb
    real(default) :: tp, tm
    if (present (x))  d%x = x
    if (present (xb)) d%xb = xb
    if (d%xb /= 0) then
       d%pb = sqrt (max (d%E**2 - d%u / d%xb**2, 0._default))
    else
       d%pb = 0
    end if
    tp = -2 * d%xb * d%E**2 + d%s + d%u
    tm = -2 * d%xb * d%p * d%pb
    d%t0 = tp + tm
    d%t1 = tp - tm
    d%t = d%t1
  end subroutine splitting_set_t_bounds

\eatline
\nwendcode{}\nwbegindocs{206}\nwdocspar
\subsection{Sampling recoil}
Compute a value for the momentum transfer $t$, using a random number
$r$.  We assume a logarithmic distribution for $t-m^2$, corresponding
to the propagator $1/(t-m^2)$ with the physical mass $m$ for the
outgoing particle.  Optionally, we can narrow the kinematical bounds.

If all three masses in the splitting vanish, the upper limit for $t$
is zero.  In that case, the $t$ value is set to zero and the splitting
will be collinear.
\nwenddocs{}\nwbegincode{207}\moddef{SF aux: splitting data: TBP}\plusendmoddef
  procedure :: sample_t => splitting_sample_t
\nwendcode{}\nwbegincode{208}\moddef{SF aux: procedures}\plusendmoddef
  subroutine splitting_sample_t (d, r, t0, t1)
    class(splitting_data_t), intent(inout) :: d
    real(default), intent(in) :: r
    real(default), intent(in), optional :: t0, t1
    real(default) :: tt0, tt1, tt0m, tt1m
    if (d%collinear) then
       d%t = d%t1
    else
       tt0 = d%t0;  if (present (t0))  tt0 = max (t0, tt0)
       tt1 = d%t1;  if (present (t1))  tt1 = min (t1, tt1)
       tt0m = tt0 - d%m2
       tt1m = tt1 - d%m2
       if (tt0m < 0 .and. tt1m < 0 .and. abs(tt0m) > &
            epsilon(tt0m) .and. abs(tt1m) > epsilon(tt0m)) then
          d%t = d%m2 + tt0m * exp (r * log (tt1m / tt0m))
       else
          d%t = tt1
       end if
    end if
  end subroutine splitting_sample_t

\eatline
\nwendcode{}\nwbegindocs{209}The inverse operation: Given $t$, we recover the value of $r$ that
would have produced this value.
\nwenddocs{}\nwbegincode{210}\moddef{SF aux: splitting data: TBP}\plusendmoddef
  procedure :: inverse_t => splitting_inverse_t
\nwendcode{}\nwbegincode{211}\moddef{SF aux: procedures}\plusendmoddef
  subroutine splitting_inverse_t (d, r, t0, t1)
    class(splitting_data_t), intent(in) :: d
    real(default), intent(out) :: r
    real(default), intent(in), optional :: t0, t1
    real(default) :: tt0, tt1, tt0m, tt1m
    if (d%collinear) then
       r = 0
    else
       tt0 = d%t0;  if (present (t0))  tt0 = max (t0, tt0)
       tt1 = d%t1;  if (present (t1))  tt1 = min (t1, tt1)
       tt0m = tt0 - d%m2
       tt1m = tt1 - d%m2
       if (tt0m < 0 .and. tt1m < 0) then
          r = log ((d%t - d%m2) / tt0m) / log (tt1m / tt0m)
       else
          r = 0
       end if
    end if
  end subroutine splitting_inverse_t
    
\eatline
\nwendcode{}\nwbegindocs{212}This is trivial, but provided for convenience:
\nwenddocs{}\nwbegincode{213}\moddef{SF aux: splitting data: TBP}\plusendmoddef
  procedure :: sample_phi => splitting_sample_phi
\nwendcode{}\nwbegincode{214}\moddef{SF aux: procedures}\plusendmoddef
  subroutine splitting_sample_phi (d, r)
    class(splitting_data_t), intent(inout) :: d
    real(default), intent(in) :: r
    if (d%collinear) then
       d%phi = 0
    else
       d%phi = (1-r) * d%phi0 + r * d%phi1
    end if
  end subroutine splitting_sample_phi

\eatline
\nwendcode{}\nwbegindocs{215}Inverse:
\nwenddocs{}\nwbegincode{216}\moddef{SF aux: splitting data: TBP}\plusendmoddef
  procedure :: inverse_phi => splitting_inverse_phi
\nwendcode{}\nwbegincode{217}\moddef{SF aux: procedures}\plusendmoddef
  subroutine splitting_inverse_phi (d, r)
    class(splitting_data_t), intent(in) :: d
    real(default), intent(out) :: r
    if (d%collinear) then
       r = 0
    else
       r = (d%phi - d%phi0) / (d%phi1 - d%phi0)
    end if
  end subroutine splitting_inverse_phi

\eatline
\nwendcode{}\nwbegindocs{218}\nwdocspar
\subsection{Splitting}
In this function, we actually perform the splitting.  The incoming momentum
$k$ is split into (if no recoil) $q_1=(1-x)k$ and $q_2=xk$.

Apart from the splitting data, we need the incoming momentum $k$, the momentum
transfer $t$, and the azimuthal angle $\phi$.  The momentum fraction $x$ is
already known here.

Alternatively, we can split without recoil.  The azimuthal angle is
irrelevant, and the momentum transfer is always equal to the upper
limit $t_1$, so the polar angle is zero.  Obviously, if there are
nonzero masses it is not possible to keep both energy-momentum
conservation and at the same time all particles on shell.  We choose
for dropping the on-shell condition here.
\nwenddocs{}\nwbegincode{219}\moddef{SF aux: splitting data: TBP}\plusendmoddef
  procedure :: split_momentum => splitting_split_momentum
\nwendcode{}\nwbegincode{220}\moddef{SF aux: procedures}\plusendmoddef
  function splitting_split_momentum (d, k) result (q)
    class(splitting_data_t), intent(in) :: d
    type(vector4_t), dimension(2) :: q
    type(vector4_t), intent(in) :: k
    real(default) :: st2, ct2, st, ct, cp, sp
    type(lorentz_transformation_t) :: rot
    real(default) :: tt0, tt1, den
    type(vector3_t) :: kk, q1, q2
    if (d%collinear) then
       if (d%s == 0 .and. d%u == 0) then
          q(1) = d%xb * k
          q(2) = d%x * k
       else
          kk = space_part (k)
          q1 = d%xb * (d%pb / d%p) * kk
          q2 = kk - q1
          q(1) = vector4_moving (d%xb * d%E, q1)
          q(2) = vector4_moving (d%x * d%E, q2)
       end if
    else       
       den = 2 * d%xb * d%p * d%pb       
       tt0 = max (d%t - d%t0, 0._default)
       tt1 = min (d%t - d%t1, 0._default)
       if (den**2 <= epsilon(den)) then
          st2 = 1
       else
          st2 = - (tt0 * tt1) / den ** 2
       end if
       if (st2 > 1) then 
          st2 = 1
       end if
       ct2 = 1 - st2
       st = sqrt (max (st2, 0._default))
       ct = sqrt (max (ct2, 0._default))
       sp = sin (d%phi)
       cp = cos (d%phi)
       rot = rotation_to_2nd (3, space_part (k))
       q1 = vector3_moving (d%xb * d%pb * [st * cp, st * sp, ct])
       q2 = vector3_moving (d%p, 3) - q1
       q(1) = rot * vector4_moving (d%xb * d%E, q1)
       q(2) = rot * vector4_moving (d%x * d%E, q2)
    end if
  end function splitting_split_momentum
    
\eatline
\nwendcode{}\nwbegindocs{221}\nwdocspar
Momenta generated by splitting will in general be off-shell.  They are
on-shell only if they are collinear and massless.  This subroutine
puts them on shell by brute force, violating either momentum or energy
conservation.  The direction of three-momentum is always retained.

If the energy is below mass shell, we return a zero momentum.
\nwenddocs{}\nwbegincode{222}\moddef{SF aux: parameters}\endmoddef
  integer, parameter, public :: KEEP_ENERGY = 0, KEEP_MOMENTUM = 1
\eatline
\nwendcode{}\nwbegincode{223}\moddef{SF aux: public}\plusendmoddef
  public :: on_shell
\nwendcode{}\nwbegincode{224}\moddef{SF aux: procedures}\plusendmoddef
  elemental subroutine on_shell (p, m2, keep)
    type(vector4_t), intent(inout) :: p
    real(default), intent(in) :: m2
    integer, intent(in) :: keep
    real(default) :: E, E2, pn
    select case (keep)
    case (KEEP_ENERGY)
       E = energy (p)
       E2 = E ** 2
       if (E2 >= m2) then
          pn = sqrt (E2 - m2)
          p = vector4_moving (E, pn * direction (space_part (p)))
       else
          p = vector4_null
       end if
    case (KEEP_MOMENTUM)
       E = sqrt (space_part (p) ** 2 + m2)
       p = vector4_moving (E, space_part (p))
    end select
  end subroutine on_shell

\eatline
\nwendcode{}\nwbegindocs{225}\nwdocspar
\subsection{Recovering the splitting}
This is the inverse problem.  We have on-shell momenta and want to
deduce the splitting parameters $x$, $t$, and $\phi$.
\nwenddocs{}\nwbegincode{226}\moddef{SF aux: splitting data: TBP}\plusendmoddef
  procedure :: recover => splitting_recover
\nwendcode{}\nwbegincode{227}\moddef{SF aux: procedures}\plusendmoddef
  subroutine splitting_recover (d, k, q, keep)
    class(splitting_data_t), intent(inout) :: d
    type(vector4_t), intent(in) :: k, q
    integer, intent(in) :: keep
    type(lorentz_transformation_t) :: rot
    type(vector4_t) :: q0, k0
    real(default) :: p1, p2, p3, pt2, pp2, pl
    real(default) :: aux, den, norm
    real(default) :: st2, ct2, ct
    rot = inverse (rotation_to_2nd (3, space_part (k)))
    q0 = rot * q
    p1 = vector4_get_component (q0, 1)
    p2 = vector4_get_component (q0, 2)
    p3 = vector4_get_component (q0, 3)
    pt2 = p1 ** 2 + p2 ** 2
    pp2 = p1 ** 2 + p2 ** 2 + p3 ** 2
    pl = abs (p3)
    k0 = vector4_moving (d%E, d%p, 3)
    select case (keep)
    case (KEEP_ENERGY)
       d%x = energy (q0) / d%E
       d%xb = 1 - d%x
       call d%set_t_bounds ()
       if (.not. d%collinear) then
          aux = (d%xb * d%pb) ** 2 * pp2 - d%p ** 2 * pt2
          den = d%p ** 2 - (d%xb * d%pb) ** 2
          if (aux >= 0 .and. den > 0) then
             norm = (d%p * pl + sqrt (aux)) / den
          else
             norm = 1
          end if
       end if
    case (KEEP_MOMENTUM)
       d%xb = sqrt (space_part (k0 - q0) ** 2 + d%u) / d%E
       d%x = 1 - d%xb
       call d%set_t_bounds ()
       norm = 1
    end select
    if (d%collinear) then
       d%t = d%t1
       d%phi = 0
    else
       if ((d%xb * d%pb * norm)**2 < epsilon(d%xb)) then
          st2 = 1
       else
          st2 = pt2 / (d%xb * d%pb * norm ) ** 2
       end if
       if (st2 > 1) then
          st2 = 1
       end if
       ct2 = 1 - st2
       ct = sqrt (max (ct2, 0._default))
       if (ct /= -1) then
          d%t = d%t1 - 2 * d%xb * d%p * d%pb * st2 / (1 + ct)
       else
          d%t = d%t0
       end if
       if (p1 /= 0 .or. p2 /= 0) then
          d%phi = atan2 (-p2, -p1)
       else
          d%phi = 0
       end if
    end if
  end subroutine splitting_recover
       
\eatline
\nwendcode{}\nwbegindocs{228}\nwdocspar
\subsection{Extract data}
\nwenddocs{}\nwbegincode{229}\moddef{SF aux: splitting data: TBP}\plusendmoddef
  procedure :: get_x => splitting_get_x
\nwendcode{}\nwbegincode{230}\moddef{SF aux: procedures}\plusendmoddef
  function splitting_get_x (sd) result (x)
    class(splitting_data_t), intent(in) :: sd
    real(default) :: x
    x = sd%x
  end function splitting_get_x

\eatline
\nwendcode{}\nwbegindocs{231}\nwdocspar
\subsection{Unit tests}
\nwenddocs{}\nwbegincode{232}\moddef{SF aux: public}\plusendmoddef
  public :: sf_aux_test
\nwendcode{}\nwbegincode{233}\moddef{SF aux: tests}\endmoddef
  subroutine sf_aux_test (u, results)
    integer, intent(in) :: u
    type(test_results_t), intent(inout) :: results
  \LA{}SF aux: execute tests\RA{}
  end subroutine sf_aux_test
  
\eatline
\nwendcode{}\nwbegindocs{234}\nwdocspar
\subsubsection{Momentum splitting: massless radiation}
Compute momentum splitting for generic kinematics.  It turns out that
for $x=0.5$, where $t-m^2$ is the geometric mean between its upper and
lower bounds (this can be directly seen from the logarithmic
distribution in the function {\tt{}sample{\char95}t} for $r \equiv x = 1 - x =
0.5$), we arrive at an exact number $t=-0.15$ for the given 
input values.
\nwenddocs{}\nwbegincode{235}\moddef{SF aux: execute tests}\endmoddef
  call test (sf_aux_1, "sf_aux_1", &
       "massless radiation", &
       u, results)
\nwendcode{}\nwbegincode{236}\moddef{SF aux: tests}\plusendmoddef
  subroutine sf_aux_1 (u)
    integer, intent(in) :: u
    type(splitting_data_t) :: sd
    type(vector4_t) :: k
    type(vector4_t), dimension(2) :: q, q0
    real(default) :: E, mk, mp, mq
    real(default) :: x, r1, r2, r1o, r2o
    real(default) :: k2, q0_2, q1_2, q2_2
    
    write (u, "(A)")  "* Test output: sf_aux_1"
    write (u, "(A)")  "*   Purpose: compute momentum splitting"
    write (u, "(A)")  "             (massless radiated particle)"
    write (u, "(A)")

    E = 1
    mk = 0.3_default
    mp = 0
    mq = mk

    k = vector4_moving (E, sqrt (E**2 - mk**2), 3)
    k2 = k ** 2;  call pacify (k2, 1e-10_default)

    x = 0.6_default
    r1 = 0.5_default
    r2 = 0.125_default
    
    write (u, "(A)")  "* (1) Non-collinear setup"
    write (u, "(A)")

    call sd%init (k, mk**2, mp**2, mq**2)
    call sd%set_t_bounds (x, 1 - x)
    call sd%sample_t (r1)
    call sd%sample_phi (r2)

    call sd%write (u)
    
    q = sd%split_momentum (k)
    q1_2 = q(1) ** 2;  call pacify (q1_2, 1e-10_default)
    q2_2 = q(2) ** 2;  call pacify (q2_2, 1e-10_default)

    write (u, "(A)")
    write (u, "(A)")  "Incoming momentum k ="
    call vector4_write (k, u)
    write (u, "(A)")
    write (u, "(A)")  "Outgoing momentum sum p + q ="
    call vector4_write (sum (q), u)
    write (u, "(A)")
    write (u, "(A)")  "Radiated momentum p ="
    call vector4_write (q(1), u)
    write (u, "(A)")
    write (u, "(A)")  "Outgoing momentum q ="
    call vector4_write (q(2), u)
    write (u, "(A)")

    write (u, "(A)")  "Compare: s"
    write (u, "(2(1x,F11.8))")  sd%s, k2
    
    write (u, "(A)")  "Compare: t"
    write (u, "(2(1x,F11.8))")  sd%t, q2_2
    
    write (u, "(A)")  "Compare: u"
    write (u, "(2(1x,F11.8))")  sd%u, q1_2
    
    write (u, "(A)")  "Compare: x"
    write (u, "(2(1x,F11.8))")  sd%x, energy (q(2)) / energy (k)
    
    write (u, "(A)")  "Compare: 1-x"
    write (u, "(2(1x,F11.8))")  sd%xb, energy (q(1)) / energy (k)
    
    write (u, "(A)")
    write (u, "(A)")  "* Project on-shell (keep energy)"

    q0 = q
    call on_shell (q0, [mp**2, mq**2], KEEP_ENERGY)

    write (u, "(A)")
    write (u, "(A)")  "Incoming momentum k ="
    call vector4_write (k, u)
    write (u, "(A)")
    write (u, "(A)")  "Outgoing momentum sum p + q ="
    call vector4_write (sum (q0), u)
    write (u, "(A)")
    write (u, "(A)")  "Radiated momentum p ="
    call vector4_write (q0(1), u)
    write (u, "(A)")
    write (u, "(A)")  "Outgoing momentum q ="
    call vector4_write (q0(2), u)
    write (u, "(A)")

    write (u, "(A)")  "Compare: mo^2"
    q0_2 = q0(2) ** 2;  call pacify (q0_2, 1e-10_default)
    write (u, "(2(1x,F11.8))")  sd%m2, q0_2
    write (u, "(A)")
    
    write (u, "(A)")  "* Recover parameters from outgoing momentum"
    write (u, "(A)")

    call sd%init (k, mk**2, mp**2, mq**2)
    call sd%recover (k, q0(2), KEEP_ENERGY)

    write (u, "(A)")  "Compare: x"
    write (u, "(2(1x,F11.8))")  x, sd%x
    write (u, "(A)")  "Compare: t"
    write (u, "(2(1x,F11.8))")  q2_2, sd%t
    

    call sd%inverse_t (r1o)
    
    write (u, "(A)")  "Compare: r1"
    write (u, "(2(1x,F11.8))")  r1, r1o

    call sd%inverse_phi (r2o)
    
    write (u, "(A)")  "Compare: r2"
    write (u, "(2(1x,F11.8))")  r2, r2o

    write (u, "(A)")
    call sd%write (u)

    write (u, "(A)")
    write (u, "(A)")  "* Project on-shell (keep momentum)"

    q0 = q
    call on_shell (q0, [mp**2, mq**2], KEEP_MOMENTUM)

    write (u, "(A)")
    write (u, "(A)")  "Incoming momentum k ="
    call vector4_write (k, u)
    write (u, "(A)")
    write (u, "(A)")  "Outgoing momentum sum p + q ="
    call vector4_write (sum (q0), u)
    write (u, "(A)")
    write (u, "(A)")  "Radiated momentum p ="
    call vector4_write (q0(1), u)
    write (u, "(A)")
    write (u, "(A)")  "Outgoing momentum q ="
    call vector4_write (q0(2), u)
    write (u, "(A)")

    write (u, "(A)")  "Compare: mo^2"
    q0_2 = q0(2) ** 2;  call pacify (q0_2, 1e-10_default)
    write (u, "(2(1x,F11.8))")  sd%m2, q0_2
    write (u, "(A)")
    
    write (u, "(A)")  "* Recover parameters from outgoing momentum"
    write (u, "(A)")

    call sd%init (k, mk**2, mp**2, mq**2)
    call sd%recover (k, q0(2), KEEP_MOMENTUM)

    write (u, "(A)")  "Compare: x"
    write (u, "(2(1x,F11.8))")  x, sd%x
    write (u, "(A)")  "Compare: t"
    write (u, "(2(1x,F11.8))")  q2_2, sd%t

    call sd%inverse_t (r1o)

    write (u, "(A)")  "Compare: r1"
    write (u, "(2(1x,F11.8))")  r1, r1o

    call sd%inverse_phi (r2o)

    write (u, "(A)")  "Compare: r2"
    write (u, "(2(1x,F11.8))")  r2, r2o

    write (u, "(A)")
    call sd%write (u)

    write (u, "(A)")
    write (u, "(A)")  "* (2) Collinear setup"
    write (u, "(A)")

    call sd%init (k, mk**2, mp**2, mq**2, collinear = .true.)
    call sd%set_t_bounds (x, 1 - x)

    call sd%write (u)
    
    q = sd%split_momentum (k)
    q1_2 = q(1) ** 2;  call pacify (q1_2, 1e-10_default)
    q2_2 = q(2) ** 2;  call pacify (q2_2, 1e-10_default)

    write (u, "(A)")
    write (u, "(A)")  "Incoming momentum k ="
    call vector4_write (k, u)
    write (u, "(A)")
    write (u, "(A)")  "Outgoing momentum sum p + q ="
    call vector4_write (sum (q), u)
    write (u, "(A)")
    write (u, "(A)")  "Radiated momentum p ="
    call vector4_write (q(1), u)
    write (u, "(A)")
    write (u, "(A)")  "Outgoing momentum q ="
    call vector4_write (q(2), u)
    write (u, "(A)")

    write (u, "(A)")  "Compare: s"
    write (u, "(2(1x,F11.8))")  sd%s, k2
    
    write (u, "(A)")  "Compare: t"
    write (u, "(2(1x,F11.8))")  sd%t, q2_2
    
    write (u, "(A)")  "Compare: u"
    write (u, "(2(1x,F11.8))")  sd%u, q1_2
    
    write (u, "(A)")  "Compare: x"
    write (u, "(2(1x,F11.8))")  sd%x, energy (q(2)) / energy (k)
    
    write (u, "(A)")  "Compare: 1-x"
    write (u, "(2(1x,F11.8))")  sd%xb, energy (q(1)) / energy (k)
    
    write (u, "(A)")
    write (u, "(A)")  "* Project on-shell (keep energy)"

    q0 = q
    call on_shell (q0, [mp**2, mq**2], KEEP_ENERGY)

    write (u, "(A)")
    write (u, "(A)")  "Incoming momentum k ="
    call vector4_write (k, u)
    write (u, "(A)")
    write (u, "(A)")  "Outgoing momentum sum p + q ="
    call vector4_write (sum (q0), u)
    write (u, "(A)")
    write (u, "(A)")  "Radiated momentum p ="
    call vector4_write (q0(1), u)
    write (u, "(A)")
    write (u, "(A)")  "Outgoing momentum q ="
    call vector4_write (q0(2), u)
    write (u, "(A)")

    write (u, "(A)")  "Compare: mo^2"
    q0_2 = q0(2) ** 2;  call pacify (q0_2, 1e-10_default)
    write (u, "(2(1x,F11.8))")  sd%m2, q0_2
    write (u, "(A)")
    
    write (u, "(A)")  "* Recover parameters from outgoing momentum"
    write (u, "(A)")

    call sd%init (k, mk**2, mp**2, mq**2)
    call sd%recover (k, q0(2), KEEP_ENERGY)

    write (u, "(A)")  "Compare: x"
    write (u, "(2(1x,F11.8))")  x, sd%x
    write (u, "(A)")  "Compare: t"
    write (u, "(2(1x,F11.8))")  q2_2, sd%t

    write (u, "(A)")
    call sd%write (u)

    write (u, "(A)")
    write (u, "(A)")  "* Project on-shell (keep momentum)"

    q0 = q
    call on_shell (q0, [mp**2, mq**2], KEEP_MOMENTUM)

    write (u, "(A)")
    write (u, "(A)")  "Incoming momentum k ="
    call vector4_write (k, u)
    write (u, "(A)")
    write (u, "(A)")  "Outgoing momentum sum p + q ="
    call vector4_write (sum (q0), u)
    write (u, "(A)")
    write (u, "(A)")  "Radiated momentum p ="
    call vector4_write (q0(1), u)
    write (u, "(A)")
    write (u, "(A)")  "Outgoing momentum q ="
    call vector4_write (q0(2), u)
    write (u, "(A)")

    write (u, "(A)")  "Compare: mo^2"
    q0_2 = q0(2) ** 2;  call pacify (q0_2, 1e-10_default)
    write (u, "(2(1x,F11.8))")  sd%m2, q0_2
    write (u, "(A)")
    
    write (u, "(A)")  "* Recover parameters from outgoing momentum"
    write (u, "(A)")

    call sd%init (k, mk**2, mp**2, mq**2)
    call sd%recover (k, q0(2), KEEP_MOMENTUM)

    write (u, "(A)")  "Compare: x"
    write (u, "(2(1x,F11.8))")  x, sd%x
    write (u, "(A)")  "Compare: t"
    write (u, "(2(1x,F11.8))")  q2_2, sd%t

    write (u, "(A)")
    call sd%write (u)

    write (u, "(A)")
    write (u, "(A)")  "* Test output end: sf_aux_1"

  end subroutine sf_aux_1
  
\eatline
\nwendcode{}\nwbegindocs{237}\nwdocspar
\subsubsection{Momentum splitting: massless parton}
Compute momentum splitting for generic kinematics.  It turns out that
for $x=0.5$, where $t-m^2$ is the geometric mean between its upper and
lower bounds, we arrive at an exact number $t=-0.36$ for the given
input values.
\nwenddocs{}\nwbegincode{238}\moddef{SF aux: execute tests}\plusendmoddef
  call test (sf_aux_2, "sf_aux_2", &
       "massless parton", &
       u, results)
\nwendcode{}\nwbegincode{239}\moddef{SF aux: tests}\plusendmoddef
  subroutine sf_aux_2 (u)
    integer, intent(in) :: u
    type(splitting_data_t) :: sd
    type(vector4_t) :: k
    type(vector4_t), dimension(2) :: q, q0
    real(default) :: E, mk, mp, mq
    real(default) :: x, r1, r2, r1o, r2o
    real(default) :: k2, q02_2, q1_2, q2_2
    
    write (u, "(A)")  "* Test output: sf_aux_2"
    write (u, "(A)")  "*   Purpose: compute momentum splitting"
    write (u, "(A)")  "             (massless outgoing particle)"
    write (u, "(A)")

    E = 1
    mk = 0.3_default
    mp = mk
    mq = 0

    k = vector4_moving (E, sqrt (E**2 - mk**2), 3)
    k2 = k ** 2;  call pacify (k2, 1e-10_default)

    x = 0.6_default
    r1 = 0.5_default
    r2 = 0.125_default
    
    write (u, "(A)")  "* (1) Non-collinear setup"
    write (u, "(A)")

    call sd%init (k, mk**2, mp**2, mq**2)
    call sd%set_t_bounds (x, 1 - x)
    call sd%sample_t (r1)
    call sd%sample_phi (r2)

    call sd%write (u)
    
    q = sd%split_momentum (k)
    q1_2 = q(1) ** 2;  call pacify (q1_2, 1e-10_default)
    q2_2 = q(2) ** 2;  call pacify (q2_2, 1e-10_default)

    write (u, "(A)")
    write (u, "(A)")  "Incoming momentum k ="
    call vector4_write (k, u)
    write (u, "(A)")
    write (u, "(A)")  "Outgoing momentum sum p + q ="
    call vector4_write (sum (q), u)
    write (u, "(A)")
    write (u, "(A)")  "Radiated momentum p ="
    call vector4_write (q(1), u)
    write (u, "(A)")
    write (u, "(A)")  "Outgoing momentum q ="
    call vector4_write (q(2), u)
    write (u, "(A)")

    write (u, "(A)")  "Compare: s"
    write (u, "(2(1x,F11.8))")  sd%s, k2
    
    write (u, "(A)")  "Compare: t"
    write (u, "(2(1x,F11.8))")  sd%t, q2_2
    
    write (u, "(A)")  "Compare: u"
    write (u, "(2(1x,F11.8))")  sd%u, q1_2
    
    write (u, "(A)")  "Compare: x"
    write (u, "(2(1x,F11.8))")  sd%x, energy (q(2)) / energy (k)
    
    write (u, "(A)")  "Compare: 1-x"
    write (u, "(2(1x,F11.8))")  sd%xb, energy (q(1)) / energy (k)
    
    write (u, "(A)")
    write (u, "(A)")  "* Project on-shell (keep energy)"

    q0 = q
    call on_shell (q0, [mp**2, mq**2], KEEP_ENERGY)

    write (u, "(A)")
    write (u, "(A)")  "Incoming momentum k ="
    call vector4_write (k, u)
    write (u, "(A)")
    write (u, "(A)")  "Outgoing momentum sum p + q ="
    call vector4_write (sum (q0), u)
    write (u, "(A)")
    write (u, "(A)")  "Radiated momentum p ="
    call vector4_write (q0(1), u)
    write (u, "(A)")
    write (u, "(A)")  "Outgoing momentum q ="
    call vector4_write (q0(2), u)
    write (u, "(A)")

    write (u, "(A)")  "Compare: mo^2"
    q02_2 = q0(2) ** 2;  call pacify (q02_2, 1e-10_default)
    write (u, "(2(1x,F11.8))")  sd%m2, q02_2
    write (u, "(A)")
    
    write (u, "(A)")  "* Recover parameters from outgoing momentum"
    write (u, "(A)")

    call sd%init (k, mk**2, mp**2, mq**2)
    call sd%set_t_bounds (x, 1 - x)
    call sd%recover (k, q0(2), KEEP_ENERGY)

    write (u, "(A)")  "Compare: x"
    write (u, "(2(1x,F11.8))")  x, sd%x
    write (u, "(A)")  "Compare: t"
    write (u, "(2(1x,F11.8))")  q2_2, sd%t
    

    call sd%inverse_t (r1o)
    
    write (u, "(A)")  "Compare: r1"
    write (u, "(2(1x,F11.8))")  r1, r1o

    call sd%inverse_phi (r2o)
    
    write (u, "(A)")  "Compare: r2"
    write (u, "(2(1x,F11.8))")  r2, r2o

    write (u, "(A)")
    call sd%write (u)

    write (u, "(A)")
    write (u, "(A)")  "* Project on-shell (keep momentum)"

    q0 = q
    call on_shell (q0, [mp**2, mq**2], KEEP_MOMENTUM)

    write (u, "(A)")
    write (u, "(A)")  "Incoming momentum k ="
    call vector4_write (k, u)
    write (u, "(A)")
    write (u, "(A)")  "Outgoing momentum sum p + q ="
    call vector4_write (sum (q0), u)
    write (u, "(A)")
    write (u, "(A)")  "Radiated momentum p ="
    call vector4_write (q0(1), u)
    write (u, "(A)")
    write (u, "(A)")  "Outgoing momentum q ="
    call vector4_write (q0(2), u)
    write (u, "(A)")

    write (u, "(A)")  "Compare: mo^2"
    q02_2 = q0(2) ** 2;  call pacify (q02_2, 1e-10_default)
    write (u, "(2(1x,F11.8))")  sd%m2, q02_2
    write (u, "(A)")
    
    write (u, "(A)")  "* Recover parameters from outgoing momentum"
    write (u, "(A)")

    call sd%init (k, mk**2, mp**2, mq**2)
    call sd%set_t_bounds (x, 1 - x)
    call sd%recover (k, q0(2), KEEP_MOMENTUM)

    write (u, "(A)")  "Compare: x"
    write (u, "(2(1x,F11.8))")  x, sd%x
    write (u, "(A)")  "Compare: t"
    write (u, "(2(1x,F11.8))")  q2_2, sd%t

    call sd%inverse_t (r1o)

    write (u, "(A)")  "Compare: r1"
    write (u, "(2(1x,F11.8))")  r1, r1o

    call sd%inverse_phi (r2o)

    write (u, "(A)")  "Compare: r2"
    write (u, "(2(1x,F11.8))")  r2, r2o

    write (u, "(A)")
    call sd%write (u)

    write (u, "(A)")
    write (u, "(A)")  "* (2) Collinear setup"
    write (u, "(A)")

    call sd%init (k, mk**2, mp**2, mq**2, collinear = .true.)
    call sd%set_t_bounds (x, 1 - x)

    call sd%write (u)
    
    q = sd%split_momentum (k)
    q1_2 = q(1) ** 2;  call pacify (q1_2, 1e-10_default)
    q2_2 = q(2) ** 2;  call pacify (q2_2, 1e-10_default)

    write (u, "(A)")
    write (u, "(A)")  "Incoming momentum k ="
    call vector4_write (k, u)
    write (u, "(A)")
    write (u, "(A)")  "Outgoing momentum sum p + q ="
    call vector4_write (sum (q), u)
    write (u, "(A)")
    write (u, "(A)")  "Radiated momentum p ="
    call vector4_write (q(1), u)
    write (u, "(A)")
    write (u, "(A)")  "Outgoing momentum q ="
    call vector4_write (q(2), u)
    write (u, "(A)")

    write (u, "(A)")  "Compare: s"
    write (u, "(2(1x,F11.8))")  sd%s, k2
    
    write (u, "(A)")  "Compare: t"
    write (u, "(2(1x,F11.8))")  sd%t, q2_2
    
    write (u, "(A)")  "Compare: u"
    write (u, "(2(1x,F11.8))")  sd%u, q1_2
    
    write (u, "(A)")  "Compare: x"
    write (u, "(2(1x,F11.8))")  sd%x, energy (q(2)) / energy (k)
    
    write (u, "(A)")  "Compare: 1-x"
    write (u, "(2(1x,F11.8))")  sd%xb, energy (q(1)) / energy (k)
    
    write (u, "(A)")
    write (u, "(A)")  "* Project on-shell (keep energy)"

    q0 = q
    call on_shell (q0, [mp**2, mq**2], KEEP_ENERGY)

    write (u, "(A)")
    write (u, "(A)")  "Incoming momentum k ="
    call vector4_write (k, u)
    write (u, "(A)")
    write (u, "(A)")  "Outgoing momentum sum p + q ="
    call vector4_write (sum (q0), u)
    write (u, "(A)")
    write (u, "(A)")  "Radiated momentum p ="
    call vector4_write (q0(1), u)
    write (u, "(A)")
    write (u, "(A)")  "Outgoing momentum q ="
    call vector4_write (q0(2), u)
    write (u, "(A)")

    write (u, "(A)")  "Compare: mo^2"
    q02_2 = q0(2) ** 2;  call pacify (q02_2, 1e-10_default)
    write (u, "(2(1x,F11.8))")  sd%m2, q02_2
    write (u, "(A)")
    
    write (u, "(A)")  "* Recover parameters from outgoing momentum"
    write (u, "(A)")

    call sd%init (k, mk**2, mp**2, mq**2)
    call sd%set_t_bounds (x, 1 - x)
    call sd%recover (k, q0(2), KEEP_ENERGY)

    write (u, "(A)")  "Compare: x"
    write (u, "(2(1x,F11.8))")  x, sd%x
    write (u, "(A)")  "Compare: t"
    write (u, "(2(1x,F11.8))")  q2_2, sd%t

    write (u, "(A)")
    call sd%write (u)

    write (u, "(A)")
    write (u, "(A)")  "* Project on-shell (keep momentum)"

    q0 = q
    call on_shell (q0, [mp**2, mq**2], KEEP_MOMENTUM)

    write (u, "(A)")
    write (u, "(A)")  "Incoming momentum k ="
    call vector4_write (k, u)
    write (u, "(A)")
    write (u, "(A)")  "Outgoing momentum sum p + q ="
    call vector4_write (sum (q0), u)
    write (u, "(A)")
    write (u, "(A)")  "Radiated momentum p ="
    call vector4_write (q0(1), u)
    write (u, "(A)")
    write (u, "(A)")  "Outgoing momentum q ="
    call vector4_write (q0(2), u)
    write (u, "(A)")

    write (u, "(A)")  "Compare: mo^2"
    q02_2 = q0(2) ** 2;  call pacify (q02_2, 1e-10_default)
    write (u, "(2(1x,F11.8))")  sd%m2, q02_2
    write (u, "(A)")
    
    write (u, "(A)")  "* Recover parameters from outgoing momentum"
    write (u, "(A)")

    call sd%init (k, mk**2, mp**2, mq**2)
    call sd%set_t_bounds (x, 1 - x)
    call sd%recover (k, q0(2), KEEP_MOMENTUM)

    write (u, "(A)")  "Compare: x"
    write (u, "(2(1x,F11.8))")  x, sd%x
    write (u, "(A)")  "Compare: t"
    write (u, "(2(1x,F11.8))")  q2_2, sd%t

    write (u, "(A)")
    call sd%write (u)

    write (u, "(A)")
    write (u, "(A)")  "* Test output end: sf_aux_2"

  end subroutine sf_aux_2
  
\eatline
\nwendcode{}\nwbegindocs{240}\nwdocspar
\subsubsection{Momentum splitting: all massless}
Compute momentum splitting for massless kinematics.  In the non-collinear
case, we need a lower cutoff for $|t|$, otherwise a logarithmic distribution
is not possible.
\nwenddocs{}\nwbegincode{241}\moddef{SF aux: execute tests}\plusendmoddef
  call test (sf_aux_3, "sf_aux_3", &
       "massless parton", &
       u, results)
\nwendcode{}\nwbegincode{242}\moddef{SF aux: tests}\plusendmoddef
  subroutine sf_aux_3 (u)
    integer, intent(in) :: u
    type(splitting_data_t) :: sd
    type(vector4_t) :: k
    type(vector4_t), dimension(2) :: q, q0
    real(default) :: E, mk, mp, mq, qmin, qmax
    real(default) :: x, r1, r2, r1o, r2o
    real(default) :: k2, q02_2, q1_2, q2_2
    
    write (u, "(A)")  "* Test output: sf_aux_3"
    write (u, "(A)")  "*   Purpose: compute momentum splitting"
    write (u, "(A)")  "             (all massless, q cuts)"
    write (u, "(A)")

    E = 1
    mk = 0
    mp = 0
    mq = 0
    qmin = 1e-2_default
    qmax = 1e0_default

    k = vector4_moving (E, sqrt (E**2 - mk**2), 3)
    k2 = k ** 2;  call pacify (k2, 1e-10_default)

    x = 0.6_default
    r1 = 0.5_default
    r2 = 0.125_default
    
    write (u, "(A)")  "* (1) Non-collinear setup"
    write (u, "(A)")

    call sd%init (k, mk**2, mp**2, mq**2)
    call sd%set_t_bounds (x, 1 - x)
    call sd%sample_t (r1, t1 = - qmin ** 2, t0 = - qmax **2)
    call sd%sample_phi (r2)

    call sd%write (u)
    
    q = sd%split_momentum (k)
    q1_2 = q(1) ** 2;  call pacify (q1_2, 1e-10_default)
    q2_2 = q(2) ** 2;  call pacify (q2_2, 1e-10_default)

    write (u, "(A)")
    write (u, "(A)")  "Incoming momentum k ="
    call vector4_write (k, u)
    write (u, "(A)")
    write (u, "(A)")  "Outgoing momentum sum p + q ="
    call vector4_write (sum (q), u)
    write (u, "(A)")
    write (u, "(A)")  "Radiated momentum p ="
    call vector4_write (q(1), u)
    write (u, "(A)")
    write (u, "(A)")  "Outgoing momentum q ="
    call vector4_write (q(2), u)
    write (u, "(A)")

    write (u, "(A)")  "Compare: s"
    write (u, "(2(1x,F11.8))")  sd%s, k2
    
    write (u, "(A)")  "Compare: t"
    write (u, "(2(1x,F11.8))")  sd%t, q2_2
    
    write (u, "(A)")  "Compare: u"
    write (u, "(2(1x,F11.8))")  sd%u, q1_2
    
    write (u, "(A)")  "Compare: x"
    write (u, "(2(1x,F11.8))")  sd%x, energy (q(2)) / energy (k)
    
    write (u, "(A)")  "Compare: 1-x"
    write (u, "(2(1x,F11.8))")  sd%xb, energy (q(1)) / energy (k)
    
    write (u, "(A)")
    write (u, "(A)")  "* Project on-shell (keep energy)"

    q0 = q
    call on_shell (q0, [mp**2, mq**2], KEEP_ENERGY)

    write (u, "(A)")
    write (u, "(A)")  "Incoming momentum k ="
    call vector4_write (k, u)
    write (u, "(A)")
    write (u, "(A)")  "Outgoing momentum sum p + q ="
    call vector4_write (sum (q0), u)
    write (u, "(A)")
    write (u, "(A)")  "Radiated momentum p ="
    call vector4_write (q0(1), u)
    write (u, "(A)")
    write (u, "(A)")  "Outgoing momentum q ="
    call vector4_write (q0(2), u)
    write (u, "(A)")

    write (u, "(A)")  "Compare: mo^2"
    q02_2 = q0(2) ** 2;  call pacify (q02_2, 1e-10_default)
    write (u, "(2(1x,F11.8))")  sd%m2, q02_2
    write (u, "(A)")
    
    write (u, "(A)")  "* Recover parameters from outgoing momentum"
    write (u, "(A)")

    call sd%init (k, mk**2, mp**2, mq**2)
    call sd%set_t_bounds (x, 1 - x)
    call sd%recover (k, q0(2), KEEP_ENERGY)

    write (u, "(A)")  "Compare: x"
    write (u, "(2(1x,F11.8))")  x, sd%x
    write (u, "(A)")  "Compare: t"
    write (u, "(2(1x,F11.8))")  q2_2, sd%t
    

    call sd%inverse_t (r1o, t1 = - qmin ** 2, t0 = - qmax **2)
    
    write (u, "(A)")  "Compare: r1"
    write (u, "(2(1x,F11.8))")  r1, r1o

    call sd%inverse_phi (r2o)
    
    write (u, "(A)")  "Compare: r2"
    write (u, "(2(1x,F11.8))")  r2, r2o

    write (u, "(A)")
    call sd%write (u)

    write (u, "(A)")
    write (u, "(A)")  "* Project on-shell (keep momentum)"

    q0 = q
    call on_shell (q0, [mp**2, mq**2], KEEP_MOMENTUM)

    write (u, "(A)")
    write (u, "(A)")  "Incoming momentum k ="
    call vector4_write (k, u)
    write (u, "(A)")
    write (u, "(A)")  "Outgoing momentum sum p + q ="
    call vector4_write (sum (q0), u)
    write (u, "(A)")
    write (u, "(A)")  "Radiated momentum p ="
    call vector4_write (q0(1), u)
    write (u, "(A)")
    write (u, "(A)")  "Outgoing momentum q ="
    call vector4_write (q0(2), u)
    write (u, "(A)")

    write (u, "(A)")  "Compare: mo^2"
    q02_2 = q0(2) ** 2;  call pacify (q02_2, 1e-10_default)
    write (u, "(2(1x,F11.8))")  sd%m2, q02_2
    write (u, "(A)")
    
    write (u, "(A)")  "* Recover parameters from outgoing momentum"
    write (u, "(A)")

    call sd%init (k, mk**2, mp**2, mq**2)
    call sd%set_t_bounds (x, 1 - x)
    call sd%recover (k, q0(2), KEEP_MOMENTUM)

    write (u, "(A)")  "Compare: x"
    write (u, "(2(1x,F11.8))")  x, sd%x
    write (u, "(A)")  "Compare: t"
    write (u, "(2(1x,F11.8))")  q2_2, sd%t

    call sd%inverse_t (r1o, t1 = - qmin ** 2, t0 = - qmax **2)

    write (u, "(A)")  "Compare: r1"
    write (u, "(2(1x,F11.8))")  r1, r1o

    call sd%inverse_phi (r2o)

    write (u, "(A)")  "Compare: r2"
    write (u, "(2(1x,F11.8))")  r2, r2o

    write (u, "(A)")
    call sd%write (u)

    write (u, "(A)")
    write (u, "(A)")  "* (2) Collinear setup"
    write (u, "(A)")

    call sd%init (k, mk**2, mp**2, mq**2, collinear = .true.)
    call sd%set_t_bounds (x, 1 - x)

    call sd%write (u)
    
    q = sd%split_momentum (k)
    q1_2 = q(1) ** 2;  call pacify (q1_2, 1e-10_default)
    q2_2 = q(2) ** 2;  call pacify (q2_2, 1e-10_default)

    write (u, "(A)")
    write (u, "(A)")  "Incoming momentum k ="
    call vector4_write (k, u)
    write (u, "(A)")
    write (u, "(A)")  "Outgoing momentum sum p + q ="
    call vector4_write (sum (q), u)
    write (u, "(A)")
    write (u, "(A)")  "Radiated momentum p ="
    call vector4_write (q(1), u)
    write (u, "(A)")
    write (u, "(A)")  "Outgoing momentum q ="
    call vector4_write (q(2), u)
    write (u, "(A)")

    write (u, "(A)")  "Compare: s"
    write (u, "(2(1x,F11.8))")  sd%s, k2
    
    write (u, "(A)")  "Compare: t"
    write (u, "(2(1x,F11.8))")  sd%t, q2_2
    
    write (u, "(A)")  "Compare: u"
    write (u, "(2(1x,F11.8))")  sd%u, q1_2
    
    write (u, "(A)")  "Compare: x"
    write (u, "(2(1x,F11.8))")  sd%x, energy (q(2)) / energy (k)
    
    write (u, "(A)")  "Compare: 1-x"
    write (u, "(2(1x,F11.8))")  sd%xb, energy (q(1)) / energy (k)
    
    write (u, "(A)")
    write (u, "(A)")  "* Project on-shell (keep energy)"

    q0 = q
    call on_shell (q0, [mp**2, mq**2], KEEP_ENERGY)

    write (u, "(A)")
    write (u, "(A)")  "Incoming momentum k ="
    call vector4_write (k, u)
    write (u, "(A)")
    write (u, "(A)")  "Outgoing momentum sum p + q ="
    call vector4_write (sum (q0), u)
    write (u, "(A)")
    write (u, "(A)")  "Radiated momentum p ="
    call vector4_write (q0(1), u)
    write (u, "(A)")
    write (u, "(A)")  "Outgoing momentum q ="
    call vector4_write (q0(2), u)
    write (u, "(A)")

    write (u, "(A)")  "Compare: mo^2"
    q02_2 = q0(2) ** 2;  call pacify (q02_2, 1e-10_default)
    write (u, "(2(1x,F11.8))")  sd%m2, q02_2
    write (u, "(A)")
    
    write (u, "(A)")  "* Recover parameters from outgoing momentum"
    write (u, "(A)")

    call sd%init (k, mk**2, mp**2, mq**2)
    call sd%set_t_bounds (x, 1 - x)
    call sd%recover (k, q0(2), KEEP_ENERGY)

    write (u, "(A)")  "Compare: x"
    write (u, "(2(1x,F11.8))")  x, sd%x
    write (u, "(A)")  "Compare: t"
    write (u, "(2(1x,F11.8))")  q2_2, sd%t

    write (u, "(A)")
    call sd%write (u)

    write (u, "(A)")
    write (u, "(A)")  "* Project on-shell (keep momentum)"

    q0 = q
    call on_shell (q0, [mp**2, mq**2], KEEP_MOMENTUM)

    write (u, "(A)")
    write (u, "(A)")  "Incoming momentum k ="
    call vector4_write (k, u)
    write (u, "(A)")
    write (u, "(A)")  "Outgoing momentum sum p + q ="
    call vector4_write (sum (q0), u)
    write (u, "(A)")
    write (u, "(A)")  "Radiated momentum p ="
    call vector4_write (q0(1), u)
    write (u, "(A)")
    write (u, "(A)")  "Outgoing momentum q ="
    call vector4_write (q0(2), u)
    write (u, "(A)")

    write (u, "(A)")  "Compare: mo^2"
    q02_2 = q0(2) ** 2;  call pacify (q02_2, 1e-10_default)
    write (u, "(2(1x,F11.8))")  sd%m2, q02_2
    write (u, "(A)")
    
    write (u, "(A)")  "* Recover parameters from outgoing momentum"
    write (u, "(A)")

    call sd%init (k, mk**2, mp**2, mq**2)
    call sd%set_t_bounds (x, 1 - x)
    call sd%recover (k, q0(2), KEEP_MOMENTUM)

    write (u, "(A)")  "Compare: x"
    write (u, "(2(1x,F11.8))")  x, sd%x
    write (u, "(A)")  "Compare: t"
    write (u, "(2(1x,F11.8))")  q2_2, sd%t

    write (u, "(A)")
    call sd%write (u)

    write (u, "(A)")
    write (u, "(A)")  "* Test output end: sf_aux_3"

  end subroutine sf_aux_3
  
\eatline
\nwendcode{}\nwbegindocs{243}\nwdocspar
%%%%%%%%%%%%%%%%%%%%%%%%%%%%%%%%%%%%%%%%%%%%%%%%%%%%%%%%%%%%%%%%%%%%%%%%
\section{Mappings for structure functions}
In this module, we provide a wrapper for useful mappings of the unit
(hyper-)square that we can apply to a set of structure functions.

In some cases it is useful, or even mandatory, to map the MC input
parameters nontrivially onto a set of structure functions for the two
beams.  In all cases considered here, instead of $x_1,x_2,\ldots$ as
parameters for the beams, we generate one parameter that is equal, or
related to, the product $x_1x_2\cdots$ (so it directly corresponds to
$\sqrt{s}$).  The other parameters describe the distribution of energy
(loss) between beams and radiations.
\nwenddocs{}\nwbegincode{244}\moddef{\code{}sf{\char95}mappings.f90\edoc{}}\endmoddef
\LA{}File header\RA{}

module sf_mappings

\LA{}Use kinds\RA{}
  use kinds, only: double
  use io_units
  use constants, only: pi
  use format_defs, only: FMT_12, FMT_13, FMT_14, FMT_15, FMT_16
  use unit_tests
  use diagnostics

\LA{}Standard module head\RA{}

\LA{}SF mappings: public\RA{}

\LA{}SF mappings: parameters\RA{}

\LA{}SF mappings: types\RA{}

\LA{}SF mappings: interfaces\RA{}

contains

\LA{}SF mappings: procedures\RA{}

\LA{}SF mappings: tests\RA{}

end module sf_mappings
\eatline
\nwendcode{}\nwbegindocs{245}\nwdocspar
\subsection{Base type}
First, we define an abstract base type for the mapping.  In all cases
we need to store the indices of the parameters on which the mapping
applies.   Additional parameters can be stored in the extensions of
this type.
\nwenddocs{}\nwbegincode{246}\moddef{SF mappings: public}\endmoddef
  public :: sf_mapping_t
\nwendcode{}\nwbegincode{247}\moddef{SF mappings: types}\endmoddef
  type, abstract :: sf_mapping_t
     integer, dimension(:), allocatable :: i
   contains
   \LA{}SF mappings: sf mapping: TBP\RA{}
  end type sf_mapping_t
     
\eatline
\nwendcode{}\nwbegindocs{248}The output routine is deferred:
\nwenddocs{}\nwbegincode{249}\moddef{SF mappings: sf mapping: TBP}\endmoddef
  procedure (sf_mapping_write), deferred :: write
\nwendcode{}\nwbegincode{250}\moddef{SF mappings: interfaces}\endmoddef
  abstract interface
     subroutine sf_mapping_write (object, unit)
       import
       class(sf_mapping_t), intent(in) :: object
       integer, intent(in), optional :: unit
     end subroutine sf_mapping_write
  end interface

\eatline
\nwendcode{}\nwbegindocs{251}Initializer for the base type.  The array of parameter indices is
allocated but initialized to zero.
\nwenddocs{}\nwbegincode{252}\moddef{SF mappings: sf mapping: TBP}\plusendmoddef
  procedure :: base_init => sf_mapping_base_init
\nwendcode{}\nwbegincode{253}\moddef{SF mappings: procedures}\endmoddef
  subroutine sf_mapping_base_init (mapping, n_par)
    class(sf_mapping_t), intent(out) :: mapping
    integer, intent(in) :: n_par
    allocate (mapping%i (n_par))
    mapping%i = 0
  end subroutine sf_mapping_base_init
    
\eatline
\nwendcode{}\nwbegindocs{254}Set an index value.
\nwenddocs{}\nwbegincode{255}\moddef{SF mappings: sf mapping: TBP}\plusendmoddef
  procedure :: set_index => sf_mapping_set_index
\nwendcode{}\nwbegincode{256}\moddef{SF mappings: procedures}\plusendmoddef
  subroutine sf_mapping_set_index (mapping, j, i)
    class(sf_mapping_t), intent(inout) :: mapping
    integer, intent(in) :: j, i
    mapping%i(j) = i
  end subroutine sf_mapping_set_index
    
\eatline
\nwendcode{}\nwbegindocs{257}Return the dimensionality, i.e., the number of parameters.
\nwenddocs{}\nwbegincode{258}\moddef{SF mappings: sf mapping: TBP}\plusendmoddef
  procedure :: get_n_dim => sf_mapping_get_n_dim
\nwendcode{}\nwbegincode{259}\moddef{SF mappings: procedures}\plusendmoddef
  function sf_mapping_get_n_dim (mapping) result (n)
    class(sf_mapping_t), intent(in) :: mapping
    integer :: n
    n = size (mapping%i)
  end function sf_mapping_get_n_dim
  
\eatline
\nwendcode{}\nwbegindocs{260}Computation: the values {\tt{}p} are the input parameters, the values
{\tt{}r} are the output parameters.  The values {\tt{}rb} are defined as
$\bar r = 1 - r$, but provided explicitly.  They allow us to avoid
numerical problems near $r=1$.  

The extra parameter {\tt{}x{\char95}free}
indicates that the total energy has already been renormalized by this
factor.  We have to take such a factor into account in a resonance or
on-shell mapping.

The Jacobian is {\tt{}f}.  We modify only
the two parameters indicated by the indices {\tt{}i}.
\nwenddocs{}\nwbegincode{261}\moddef{SF mappings: sf mapping: TBP}\plusendmoddef
  procedure (sf_mapping_compute), deferred :: compute
\nwendcode{}\nwbegincode{262}\moddef{SF mappings: interfaces}\plusendmoddef
  abstract interface
     subroutine sf_mapping_compute (mapping, r, rb, f, p, pb, x_free)
       import
       class(sf_mapping_t), intent(inout) :: mapping
       real(default), dimension(:), intent(out) :: r, rb
       real(default), intent(out) :: f
       real(default), dimension(:), intent(in) :: p, pb
       real(default), intent(inout), optional :: x_free
     end subroutine sf_mapping_compute
  end interface
  
\eatline
\nwendcode{}\nwbegindocs{263}The inverse mapping.  Use {\tt{}r} and/or {\tt{}rb} to reconstruct {\tt{}p}
and also compute {\tt{}f}.
\nwenddocs{}\nwbegincode{264}\moddef{SF mappings: sf mapping: TBP}\plusendmoddef
  procedure (sf_mapping_inverse), deferred :: inverse
\nwendcode{}\nwbegincode{265}\moddef{SF mappings: interfaces}\plusendmoddef
  abstract interface
     subroutine sf_mapping_inverse (mapping, r, rb, f, p, pb, x_free)
       import
       class(sf_mapping_t), intent(inout) :: mapping
       real(default), dimension(:), intent(in) :: r, rb
       real(default), intent(out) :: f
       real(default), dimension(:), intent(out) :: p, pb
       real(default), intent(inout), optional :: x_free
     end subroutine sf_mapping_inverse
  end interface
  
\eatline
\nwendcode{}\nwbegindocs{266}\nwdocspar
\subsection{Methods for self-tests}
This is a shorthand for: inject parameters, compute the mapping,
display results, compute the inverse, display again.  We provide an
output format for the parameters and, optionally, a different output
format for the Jacobians.
\nwenddocs{}\nwbegincode{267}\moddef{SF mappings: sf mapping: TBP}\plusendmoddef
  procedure :: check => sf_mapping_check
\nwendcode{}\nwbegincode{268}\moddef{SF mappings: procedures}\plusendmoddef
  subroutine sf_mapping_check (mapping, u, p_in, pb_in, fmt_p, fmt_f)
    class(sf_mapping_t), intent(inout) :: mapping
    integer, intent(in) :: u
    real(default), dimension(:), intent(in) :: p_in, pb_in
    character(*), intent(in) :: fmt_p
    character(*), intent(in), optional :: fmt_f
    real(default), dimension(size(p_in)) :: p, pb, r, rb
    real(default) :: f, tolerance
    tolerance = 1.5E-17
    p = p_in
    pb= pb_in
    call mapping%compute (r, rb, f, p, pb)
    call pacify (p, tolerance)  
    call pacify (pb, tolerance) 
    call pacify (r, tolerance)  
    call pacify (rb, tolerance) 
    write (u, "(3x,A,9(1x," // fmt_p // "))")  "p =", p
    write (u, "(3x,A,9(1x," // fmt_p // "))")  "pb=", pb
    write (u, "(3x,A,9(1x," // fmt_p // "))")  "r =", r
    write (u, "(3x,A,9(1x," // fmt_p // "))")  "rb=", rb
    if (present (fmt_f)) then
       write (u, "(3x,A,9(1x," // fmt_f // "))")  "f =", f
    else
       write (u, "(3x,A,9(1x," // fmt_p // "))")  "f =", f
    end if
    write (u, *)
    call mapping%inverse (r, rb, f, p, pb)
    call pacify (p, tolerance)  
    call pacify (pb, tolerance) 
    call pacify (r, tolerance)  
    call pacify (rb, tolerance) 
    write (u, "(3x,A,9(1x," // fmt_p // "))")  "p =", p
    write (u, "(3x,A,9(1x," // fmt_p // "))")  "pb=", pb
    write (u, "(3x,A,9(1x," // fmt_p // "))")  "r =", r
    write (u, "(3x,A,9(1x," // fmt_p // "))")  "rb=", rb
    if (present (fmt_f)) then
       write (u, "(3x,A,9(1x," // fmt_f // "))")  "f =", f
    else
       write (u, "(3x,A,9(1x," // fmt_p // "))")  "f =", f
    end if
    write (u, *)
    write (u, "(3x,A,9(1x," // fmt_p // "))")  "*r=", product (r)
  end subroutine sf_mapping_check
    
\eatline
\nwendcode{}\nwbegindocs{269}This is a consistency check for the self-tests: the integral over the unit
square should be unity.  We estimate this by a simple binning and adding up
the values; this should be sufficient for a self-test.

The argument is the requested number of sampling points.  We take the square
root for binning in both dimensions, so the precise number might be
different.
\nwenddocs{}\nwbegincode{270}\moddef{SF mappings: sf mapping: TBP}\plusendmoddef
  procedure :: integral => sf_mapping_integral
\nwendcode{}\nwbegincode{271}\moddef{SF mappings: procedures}\plusendmoddef
  function sf_mapping_integral (mapping, n_calls) result (integral)
    class(sf_mapping_t), intent(inout) :: mapping
    integer, intent(in) :: n_calls
    real(default) :: integral
    integer :: n_dim, n_bin, k
    real(default), dimension(:), allocatable :: p, pb, r, rb
    integer, dimension(:), allocatable :: ii
    real(default) :: dx, f, s

    n_dim = mapping%get_n_dim ()
    allocate (p (n_dim))
    allocate (pb(n_dim))
    allocate (r (n_dim))
    allocate (rb(n_dim))
    allocate (ii(n_dim))
    n_bin = nint (real (n_calls, default) ** (1._default / n_dim))
    dx = 1._default / n_bin
    s = 0
    ii = 1

    SAMPLE: do
       do k = 1, n_dim
          p(k)  = ii(k) * dx - dx/2
          pb(k) = (n_bin - ii(k)) * dx + dx/2
       end do
       call mapping%compute (r, rb, f, p, pb)
       s = s + f
       INCR: do k = 1, n_dim
          ii(k) = ii(k) + 1
          if (ii(k) <= n_bin) then
             exit INCR
          else if (k < n_dim) then
             ii(k) = 1
          else
             exit SAMPLE
          end if
       end do INCR
    end do SAMPLE

    integral = s / real (n_bin, default) ** n_dim

  end function sf_mapping_integral

\eatline
\nwendcode{}\nwbegindocs{272}\nwdocspar
\subsection{Implementation: standard mapping}
This maps the unit square ($r_1,r_2$) such that $p_1$ is the product $r_1r_2$,
while $p_2$ is related to the ratio.
\nwenddocs{}\nwbegincode{273}\moddef{SF mappings: public}\plusendmoddef
  public :: sf_s_mapping_t
\nwendcode{}\nwbegincode{274}\moddef{SF mappings: types}\plusendmoddef
  type, extends (sf_mapping_t) :: sf_s_mapping_t
     logical :: power_set = .false.
     real(default) :: power = 1
   contains
   \LA{}SF mappings: sf standard mapping: TBP\RA{}
  end type sf_s_mapping_t
  
\eatline
\nwendcode{}\nwbegindocs{275}Output.
\nwenddocs{}\nwbegincode{276}\moddef{SF mappings: sf standard mapping: TBP}\endmoddef
  procedure :: write => sf_s_mapping_write
\nwendcode{}\nwbegincode{277}\moddef{SF mappings: procedures}\plusendmoddef
  subroutine sf_s_mapping_write (object, unit)
    class(sf_s_mapping_t), intent(in) :: object
    integer, intent(in), optional :: unit
    integer :: u
    u = given_output_unit (unit)
    write (u, "(1x,A)", advance="no")  "map"
    if (any (object%i /= 0)) then
       write (u, "('(',I0,',',I0,')')", advance="no")  object%i
    end if
    write (u, "(A,F7.5,A)")  ": standard (", object%power, ")"
  end subroutine sf_s_mapping_write
  
\eatline
\nwendcode{}\nwbegindocs{278}Initialize: index pair and power parameter.
\nwenddocs{}\nwbegincode{279}\moddef{SF mappings: sf standard mapping: TBP}\plusendmoddef
  procedure :: init => sf_s_mapping_init
\nwendcode{}\nwbegincode{280}\moddef{SF mappings: procedures}\plusendmoddef
  subroutine sf_s_mapping_init (mapping, power)
    class(sf_s_mapping_t), intent(out) :: mapping
    real(default), intent(in), optional :: power
    call mapping%base_init (2)
    if (present (power)) then
       mapping%power_set = .true.
       mapping%power = power
    end if
  end subroutine sf_s_mapping_init
    
\eatline
\nwendcode{}\nwbegindocs{281}Apply mapping.
\nwenddocs{}\nwbegincode{282}\moddef{SF mappings: sf standard mapping: TBP}\plusendmoddef
  procedure :: compute => sf_s_mapping_compute
\nwendcode{}\nwbegincode{283}\moddef{SF mappings: procedures}\plusendmoddef
  subroutine sf_s_mapping_compute (mapping, r, rb, f, p, pb, x_free)
    class(sf_s_mapping_t), intent(inout) :: mapping
    real(default), dimension(:), intent(out) :: r, rb
    real(default), intent(out) :: f
    real(default), dimension(:), intent(in) :: p, pb
    real(default), intent(inout), optional :: x_free
    real(default), dimension(2) :: r2
    integer :: j
    if (mapping%power_set) then
       call map_unit_square (r2, f, p(mapping%i), mapping%power)
    else
       call map_unit_square (r2, f, p(mapping%i))
    end if
    r = p
    rb= pb
    do j = 1, 2
       r (mapping%i(j)) = r2(j)
       rb(mapping%i(j)) = 1 - r2(j)
    end do
  end subroutine sf_s_mapping_compute

\eatline
\nwendcode{}\nwbegindocs{284}Apply inverse.
\nwenddocs{}\nwbegincode{285}\moddef{SF mappings: sf standard mapping: TBP}\plusendmoddef
  procedure :: inverse => sf_s_mapping_inverse
\nwendcode{}\nwbegincode{286}\moddef{SF mappings: procedures}\plusendmoddef
  subroutine sf_s_mapping_inverse (mapping, r, rb, f, p, pb, x_free)
    class(sf_s_mapping_t), intent(inout) :: mapping
    real(default), dimension(:), intent(in) :: r, rb
    real(default), intent(out) :: f
    real(default), dimension(:), intent(out) :: p, pb
    real(default), intent(inout), optional :: x_free
    real(default), dimension(2) :: p2
    integer :: j
    if (mapping%power_set) then
       call map_unit_square_inverse (r(mapping%i), f, p2, mapping%power)
    else
       call map_unit_square_inverse (r(mapping%i), f, p2)
    end if
    p = r
    pb= rb
    do j = 1, 2
       p (mapping%i(j)) = p2(j)
       pb(mapping%i(j)) = 1 - p2(j)
    end do
  end subroutine sf_s_mapping_inverse

\eatline
\nwendcode{}\nwbegindocs{287}\nwdocspar
\subsection{Implementation: resonance pair mapping}
This maps the unit square ($r_1,r_2$) such that $p_1$ is the product $r_1r_2$,
while $p_2$ is related to the ratio, then it maps $p_1$ to itself
according to a Breit-Wigner shape, i.e., a flat prior distribution in $p_1$
results in a Breit-Wigner distribution.  Mass and width of the BW are
rescaled by the energy, thus dimensionless fractions.
\nwenddocs{}\nwbegincode{288}\moddef{SF mappings: public}\plusendmoddef
  public :: sf_res_mapping_t
\nwendcode{}\nwbegincode{289}\moddef{SF mappings: types}\plusendmoddef
  type, extends (sf_mapping_t) :: sf_res_mapping_t
     real(default) :: m = 0
     real(default) :: w = 0
   contains
   \LA{}SF mappings: sf resonance mapping: TBP\RA{}
  end type sf_res_mapping_t
  
\eatline
\nwendcode{}\nwbegindocs{290}Output.
\nwenddocs{}\nwbegincode{291}\moddef{SF mappings: sf resonance mapping: TBP}\endmoddef
  procedure :: write => sf_res_mapping_write
\nwendcode{}\nwbegincode{292}\moddef{SF mappings: procedures}\plusendmoddef
  subroutine sf_res_mapping_write (object, unit)
    class(sf_res_mapping_t), intent(in) :: object
    integer, intent(in), optional :: unit
    integer :: u
    u = given_output_unit (unit)
    write (u, "(1x,A)", advance="no")  "map"
    if (any (object%i /= 0)) then
       write (u, "('(',I0,',',I0,')')", advance="no")  object%i
    end if
    write (u, "(A,F7.5,', ',F7.5,A)")  ": resonance (", object%m, object%w, ")"
  end subroutine sf_res_mapping_write
  
\eatline
\nwendcode{}\nwbegindocs{293}Initialize: index pair and dimensionless mass and width parameters.
\nwenddocs{}\nwbegincode{294}\moddef{SF mappings: sf resonance mapping: TBP}\plusendmoddef
  procedure :: init => sf_res_mapping_init
\nwendcode{}\nwbegincode{295}\moddef{SF mappings: procedures}\plusendmoddef
  subroutine sf_res_mapping_init (mapping, m, w)
    class(sf_res_mapping_t), intent(out) :: mapping
    real(default), intent(in) :: m, w
    call mapping%base_init (2)
    mapping%m = m
    mapping%w = w
  end subroutine sf_res_mapping_init
    
\eatline
\nwendcode{}\nwbegindocs{296}Apply mapping.
\nwenddocs{}\nwbegincode{297}\moddef{SF mappings: sf resonance mapping: TBP}\plusendmoddef
  procedure :: compute => sf_res_mapping_compute
\nwendcode{}\nwbegincode{298}\moddef{SF mappings: procedures}\plusendmoddef
  subroutine sf_res_mapping_compute (mapping, r, rb, f, p, pb, x_free)
    class(sf_res_mapping_t), intent(inout) :: mapping
    real(default), dimension(:), intent(out) :: r, rb
    real(default), intent(out) :: f
    real(default), dimension(:), intent(in) :: p, pb
    real(default), intent(inout), optional :: x_free
    real(default), dimension(2) :: r2, p2
    real(default) :: fbw, f2, p1m
    integer :: j
    p2 = p(mapping%i)
    call map_breit_wigner &
         (p1m, fbw, p2(1), mapping%m, mapping%w, x_free)
    call map_unit_square (r2, f2, [p1m, p2(2)])
    f = fbw * f2
    r = p
    rb= pb
    do j = 1, 2
       r (mapping%i(j)) = r2(j)
       rb(mapping%i(j)) = 1 - r2(j)
    end do
  end subroutine sf_res_mapping_compute

\eatline
\nwendcode{}\nwbegindocs{299}Apply inverse.
\nwenddocs{}\nwbegincode{300}\moddef{SF mappings: sf resonance mapping: TBP}\plusendmoddef
  procedure :: inverse => sf_res_mapping_inverse
\nwendcode{}\nwbegincode{301}\moddef{SF mappings: procedures}\plusendmoddef
  subroutine sf_res_mapping_inverse (mapping, r, rb, f, p, pb, x_free)
    class(sf_res_mapping_t), intent(inout) :: mapping
    real(default), dimension(:), intent(in) :: r, rb
    real(default), intent(out) :: f
    real(default), dimension(:), intent(out) :: p, pb
    real(default), intent(inout), optional :: x_free
    real(default), dimension(2) :: p2
    real(default) :: fbw, f2, p1m
    call map_unit_square_inverse (r(mapping%i), f2, p2)
    call map_breit_wigner_inverse &
         (p2(1), fbw, p1m, mapping%m, mapping%w, x_free)
    p = r
    pb= rb
    p (mapping%i(1)) = p1m
    pb(mapping%i(1)) = 1 - p1m
    p (mapping%i(2)) = p2(2)
    pb(mapping%i(2)) = 1 - p2(2)
    f = fbw * f2
  end subroutine sf_res_mapping_inverse

\eatline
\nwendcode{}\nwbegindocs{302}\nwdocspar
\subsection{Implementation: on-shell mapping}
This is a degenerate version of the unit-square mapping where the
product $r_1r_2$ is constant.  This product is given by the rescaled
squared mass.  We introduce an artificial first parameter $p_1$ to
keep the counting, but nothing depends on it.  The second parameter is
the same $p_2$ as for the standard unit-square mapping for $\alpha=1$,
it parameterizes the ratio of $r_1$ and $r_2$.
\nwenddocs{}\nwbegincode{303}\moddef{SF mappings: public}\plusendmoddef
  public :: sf_os_mapping_t
\nwendcode{}\nwbegincode{304}\moddef{SF mappings: types}\plusendmoddef
  type, extends (sf_mapping_t) :: sf_os_mapping_t
     real(default) :: m = 0
     real(default) :: lm2 = 0
   contains
   \LA{}SF mappings: sf on-shell mapping: TBP\RA{}
  end type sf_os_mapping_t
  
\eatline
\nwendcode{}\nwbegindocs{305}Output.
\nwenddocs{}\nwbegincode{306}\moddef{SF mappings: sf on-shell mapping: TBP}\endmoddef
  procedure :: write => sf_os_mapping_write
\nwendcode{}\nwbegincode{307}\moddef{SF mappings: procedures}\plusendmoddef
  subroutine sf_os_mapping_write (object, unit)
    class(sf_os_mapping_t), intent(in) :: object
    integer, intent(in), optional :: unit
    integer :: u
    u = given_output_unit (unit)
    write (u, "(1x,A)", advance="no")  "map"
    if (any (object%i /= 0)) then
       write (u, "('(',I0,',',I0,')')", advance="no")  object%i
    end if
    write (u, "(A,F7.5,A)")  ": on-shell (", object%m, ")"
  end subroutine sf_os_mapping_write
  
\eatline
\nwendcode{}\nwbegindocs{308}Initialize: index pair and dimensionless mass parameter.
\nwenddocs{}\nwbegincode{309}\moddef{SF mappings: sf on-shell mapping: TBP}\plusendmoddef
  procedure :: init => sf_os_mapping_init
\nwendcode{}\nwbegincode{310}\moddef{SF mappings: procedures}\plusendmoddef
  subroutine sf_os_mapping_init (mapping, m)
    class(sf_os_mapping_t), intent(out) :: mapping
    real(default), intent(in) :: m
    call mapping%base_init (2)
    mapping%m = m
    mapping%lm2 = abs (2 * log (mapping%m))
  end subroutine sf_os_mapping_init
    
\eatline
\nwendcode{}\nwbegindocs{311}Apply mapping.  The {\tt{}x{\char95}free} parameter rescales the total energy,
which must be accounted for in the enclosed mapping.
\nwenddocs{}\nwbegincode{312}\moddef{SF mappings: sf on-shell mapping: TBP}\plusendmoddef
  procedure :: compute => sf_os_mapping_compute
\nwendcode{}\nwbegincode{313}\moddef{SF mappings: procedures}\plusendmoddef
  subroutine sf_os_mapping_compute (mapping, r, rb, f, p, pb, x_free)
    class(sf_os_mapping_t), intent(inout) :: mapping
    real(default), dimension(:), intent(out) :: r, rb
    real(default), intent(out) :: f
    real(default), dimension(:), intent(in) :: p, pb
    real(default), intent(inout), optional :: x_free
    real(default), dimension(2) :: r2, p2
    integer :: j
    p2 = p(mapping%i)
    call map_on_shell (r2, f, p2, mapping%lm2, x_free)
    r = p
    rb= pb
    do j = 1, 2
       r (mapping%i(j)) = r2(j)
       rb(mapping%i(j)) = 1 - r2(j)
    end do
  end subroutine sf_os_mapping_compute

\eatline
\nwendcode{}\nwbegindocs{314}Apply inverse.  The irrelevant parameter $p_1$ is always set zero.
\nwenddocs{}\nwbegincode{315}\moddef{SF mappings: sf on-shell mapping: TBP}\plusendmoddef
  procedure :: inverse => sf_os_mapping_inverse
\nwendcode{}\nwbegincode{316}\moddef{SF mappings: procedures}\plusendmoddef
  subroutine sf_os_mapping_inverse (mapping, r, rb, f, p, pb, x_free)
    class(sf_os_mapping_t), intent(inout) :: mapping
    real(default), dimension(:), intent(in) :: r, rb
    real(default), intent(out) :: f
    real(default), dimension(:), intent(out) :: p, pb
    real(default), intent(inout), optional :: x_free
    real(default), dimension(2) :: p2, r2
    r2 = r(mapping%i)
    call map_on_shell_inverse (r2, f, p2, mapping%lm2, x_free)
    p = r
    pb= rb
    p (mapping%i(1)) = p2(1)
    pb(mapping%i(1)) = 1 - p2(1)
    p (mapping%i(2)) = p2(2)
    pb(mapping%i(2)) = 1 - p2(2)
  end subroutine sf_os_mapping_inverse

\eatline
\nwendcode{}\nwbegindocs{317}\nwdocspar
\subsection{Implementation: endpoint mapping}
This maps the unit square ($r_1,r_2$) such that $p_1$ is the product $r_1r_2$,
while $p_2$ is related to the ratio.  Furthermore, we enhance the
region at $r_1=1$ and $r_2=1$, which translates into $p_1=1$ and
$p_2=0,1$.  The enhancement is such that any power-like singularity is
caught.  This is useful for beamstrahlung spectra.

In addition, we allow for a delta-function singularity in $r_1$ and/or
$r_2$.  The singularity is smeared to an interval of width
$\epsilon$.  If nonzero, we distinguish the kinematical momentum
fractions $r_i$ from effective values $x_i$, which should go into the
structure-function evaluation.  A bin of width $\epsilon$ in $r$ is
mapped to $x=1$ exactly, while the interval $(0,1-\epsilon)$ is mapped
to $(0,1)$ in $x$.  The Jacobian reflects this distinction, and the
logical {\tt{}in{\char95}peak} allows for an unambiguous distinction.

The delta-peak fraction is used only for the integration self-test.
\nwenddocs{}\nwbegincode{318}\moddef{SF mappings: public}\plusendmoddef
  public :: sf_ep_mapping_t
\nwendcode{}\nwbegincode{319}\moddef{SF mappings: types}\plusendmoddef
  type, extends (sf_mapping_t) :: sf_ep_mapping_t
     real(default) :: a = 1
   contains
   \LA{}SF mappings: sf endpoint mapping: TBP\RA{}
  end type sf_ep_mapping_t
  
\eatline
\nwendcode{}\nwbegindocs{320}Output.
\nwenddocs{}\nwbegincode{321}\moddef{SF mappings: sf endpoint mapping: TBP}\endmoddef
  procedure :: write => sf_ep_mapping_write
\nwendcode{}\nwbegincode{322}\moddef{SF mappings: procedures}\plusendmoddef
  subroutine sf_ep_mapping_write (object, unit)
    class(sf_ep_mapping_t), intent(in) :: object
    integer, intent(in), optional :: unit
    integer :: u
    u = given_output_unit (unit)
    write (u, "(1x,A)", advance="no")  "map"
    if (any (object%i /= 0)) then
       write (u, "('(',I0,',',I0,')')", advance="no")  object%i
    end if
    write (u, "(A,ES12.5,A)")  ": endpoint (a =", object%a, ")"
  end subroutine sf_ep_mapping_write
  
\eatline
\nwendcode{}\nwbegindocs{323}Initialize: no extra parameters.
\nwenddocs{}\nwbegincode{324}\moddef{SF mappings: sf endpoint mapping: TBP}\plusendmoddef
  procedure :: init => sf_ep_mapping_init
\nwendcode{}\nwbegincode{325}\moddef{SF mappings: procedures}\plusendmoddef
  subroutine sf_ep_mapping_init (mapping, a)
    class(sf_ep_mapping_t), intent(out) :: mapping
    real(default), intent(in), optional :: a
    call mapping%base_init (2)
    if (present (a))  mapping%a = a
  end subroutine sf_ep_mapping_init
    
\eatline
\nwendcode{}\nwbegindocs{326}Apply mapping.
\nwenddocs{}\nwbegincode{327}\moddef{SF mappings: sf endpoint mapping: TBP}\plusendmoddef
  procedure :: compute => sf_ep_mapping_compute
\nwendcode{}\nwbegincode{328}\moddef{SF mappings: procedures}\plusendmoddef
  subroutine sf_ep_mapping_compute (mapping, r, rb, f, p, pb, x_free)
    class(sf_ep_mapping_t), intent(inout) :: mapping
    real(default), dimension(:), intent(out) :: r, rb
    real(default), intent(out) :: f
    real(default), dimension(:), intent(in) :: p, pb
    real(default), intent(inout), optional :: x_free
    real(default), dimension(2) :: px, r2
    real(default) :: f1, f2
    integer :: j
    call map_endpoint_1 (px(1), f1, p(mapping%i(1)), mapping%a)
    call map_endpoint_01 (px(2), f2, p(mapping%i(2)), mapping%a)
    call map_unit_square (r2, f, px)
    f = f * f1 * f2
    r = p
    rb= pb
    do j = 1, 2
       r (mapping%i(j)) = r2(j)
       rb(mapping%i(j)) = 1 - r2(j)
    end do
  end subroutine sf_ep_mapping_compute

\eatline
\nwendcode{}\nwbegindocs{329}Apply inverse.
\nwenddocs{}\nwbegincode{330}\moddef{SF mappings: sf endpoint mapping: TBP}\plusendmoddef
  procedure :: inverse => sf_ep_mapping_inverse
\nwendcode{}\nwbegincode{331}\moddef{SF mappings: procedures}\plusendmoddef
  subroutine sf_ep_mapping_inverse (mapping, r, rb, f, p, pb, x_free)
    class(sf_ep_mapping_t), intent(inout) :: mapping
    real(default), dimension(:), intent(in) :: r, rb
    real(default), intent(out) :: f
    real(default), dimension(:), intent(out) :: p, pb
    real(default), intent(inout), optional :: x_free
    real(default), dimension(2) :: r2, px, p2
    real(default) :: f1, f2
    integer :: j
    do j = 1, 2
       r2(j) = r(mapping%i(j))
    end do
    call map_unit_square_inverse (r2, f, px)
    call map_endpoint_inverse_1 (px(1), f1, p2(1), mapping%a)
    call map_endpoint_inverse_01 (px(2), f2, p2(2), mapping%a)
    f = f * f1 * f2
    p = r
    pb= rb
    do j = 1, 2
       p (mapping%i(j)) = p2(j)
       pb(mapping%i(j)) = 1 - p2(j)
    end do
  end subroutine sf_ep_mapping_inverse

\eatline
\nwendcode{}\nwbegindocs{332}\nwdocspar
\subsection{Implementation: endpoint mapping with resonance}
Like the endpoint mapping for $p_2$, but replace the endpoint mapping
by a Breit-Wigner mapping for $p_1$.  This covers resonance production
in the presence of beamstrahlung.

If the flag {\tt{}resonance} is unset, we skip the resonance mapping, so
the parameter $p_1$ remains equal to $r_1r_2$, as in the standard
s-channel mapping.
\nwenddocs{}\nwbegincode{333}\moddef{SF mappings: public}\plusendmoddef
  public :: sf_epr_mapping_t
\nwendcode{}\nwbegincode{334}\moddef{SF mappings: types}\plusendmoddef
  type, extends (sf_mapping_t) :: sf_epr_mapping_t
     real(default) :: a = 1
     real(default) :: m = 0
     real(default) :: w = 0
     logical :: resonance = .true.
   contains
   \LA{}SF mappings: sf endpoint/res mapping: TBP\RA{}
  end type sf_epr_mapping_t
  
\eatline
\nwendcode{}\nwbegindocs{335}Output.
\nwenddocs{}\nwbegincode{336}\moddef{SF mappings: sf endpoint/res mapping: TBP}\endmoddef
  procedure :: write => sf_epr_mapping_write
\nwendcode{}\nwbegincode{337}\moddef{SF mappings: procedures}\plusendmoddef
  subroutine sf_epr_mapping_write (object, unit)
    class(sf_epr_mapping_t), intent(in) :: object
    integer, intent(in), optional :: unit
    integer :: u
    u = given_output_unit (unit)
    write (u, "(1x,A)", advance="no")  "map"
    if (any (object%i /= 0)) then
       write (u, "('(',I0,',',I0,')')", advance="no")  object%i
    end if
    if (object%resonance) then
       write (u, "(A,F7.5,A,F7.5,', ',F7.5,A)")  ": ep/res (a = ", object%a, &
            " | ", object%m, object%w, ")"
    else
       write (u, "(A,F7.5,A)")  ": ep/nores (a = ", object%a, ")"
    end if
  end subroutine sf_epr_mapping_write
  
\eatline
\nwendcode{}\nwbegindocs{338}Initialize: if mass and width are not given, we initialize a
non-resonant version of the mapping.
\nwenddocs{}\nwbegincode{339}\moddef{SF mappings: sf endpoint/res mapping: TBP}\plusendmoddef
  procedure :: init => sf_epr_mapping_init
\nwendcode{}\nwbegincode{340}\moddef{SF mappings: procedures}\plusendmoddef
  subroutine sf_epr_mapping_init (mapping, a, m, w)
    class(sf_epr_mapping_t), intent(out) :: mapping
    real(default), intent(in) :: a
    real(default), intent(in), optional :: m, w
    call mapping%base_init (2)
    mapping%a = a
    if (present (m) .and. present (w)) then
       mapping%m = m
       mapping%w = w
    else
       mapping%resonance = .false.
    end if
  end subroutine sf_epr_mapping_init
    
\eatline
\nwendcode{}\nwbegindocs{341}Apply mapping.
\nwenddocs{}\nwbegincode{342}\moddef{SF mappings: sf endpoint/res mapping: TBP}\plusendmoddef
  procedure :: compute => sf_epr_mapping_compute
\nwendcode{}\nwbegincode{343}\moddef{SF mappings: procedures}\plusendmoddef
  subroutine sf_epr_mapping_compute (mapping, r, rb, f, p, pb, x_free)
    class(sf_epr_mapping_t), intent(inout) :: mapping
    real(default), dimension(:), intent(out) :: r, rb
    real(default), intent(out) :: f
    real(default), dimension(:), intent(in) :: p, pb
    real(default), intent(inout), optional :: x_free
    real(default), dimension(2) :: px, r2
    real(default) :: f1, f2
    integer :: j
    if (mapping%resonance) then
       call map_breit_wigner &
            (px(1), f1, p(mapping%i(1)), mapping%m, mapping%w, x_free)
    else
       px(1) = p(mapping%i(1))
       f1 = 1
    end if
    call map_endpoint_01 (px(2), f2, p(mapping%i(2)), mapping%a)
    call map_unit_square (r2, f, px)
    f = f * f1 * f2
    r = p
    rb= pb
    do j = 1, 2
       r (mapping%i(j)) = r2(j)
       rb(mapping%i(j)) = 1 - r2(j)
    end do
  end subroutine sf_epr_mapping_compute

\eatline
\nwendcode{}\nwbegindocs{344}Apply inverse.
\nwenddocs{}\nwbegincode{345}\moddef{SF mappings: sf endpoint/res mapping: TBP}\plusendmoddef
  procedure :: inverse => sf_epr_mapping_inverse
\nwendcode{}\nwbegincode{346}\moddef{SF mappings: procedures}\plusendmoddef
  subroutine sf_epr_mapping_inverse (mapping, r, rb, f, p, pb, x_free)
    class(sf_epr_mapping_t), intent(inout) :: mapping
    real(default), dimension(:), intent(in) :: r, rb
    real(default), intent(out) :: f
    real(default), dimension(:), intent(out) :: p, pb
    real(default), intent(inout), optional :: x_free
    real(default), dimension(2) :: px, p2
    real(default) :: f1, f2
    integer :: j
    call map_unit_square_inverse (r(mapping%i), f, px)
    if (mapping%resonance) then
       call map_breit_wigner_inverse &
            (px(1), f1, p2(1), mapping%m, mapping%w, x_free)
    else
       p2(1) = px(1)
       f1 = 1
    end if
    call map_endpoint_inverse_01 (px(2), f2, p2(2), mapping%a)
    f = f * f1 * f2
    p = r
    pb= rb
    do j = 1, 2
       p (mapping%i(j)) = p2(j)
       pb(mapping%i(j)) = 1 - p2(j)
    end do
  end subroutine sf_epr_mapping_inverse

\eatline
\nwendcode{}\nwbegindocs{347}\nwdocspar
\subsection{Implementation: endpoint mapping for on-shell particle}
Analogous to the resonance mapping, but the $p_1$ input is ignored
altogether.  This covers on-shell particle production
in the presence of beamstrahlung.
\nwenddocs{}\nwbegincode{348}\moddef{SF mappings: public}\plusendmoddef
  public :: sf_epo_mapping_t
\nwendcode{}\nwbegincode{349}\moddef{SF mappings: types}\plusendmoddef
  type, extends (sf_mapping_t) :: sf_epo_mapping_t
     real(default) :: a = 1
     real(default) :: m = 0
     real(default) :: lm2 = 0
   contains
   \LA{}SF mappings: sf endpoint/os mapping: TBP\RA{}
  end type sf_epo_mapping_t
  
\eatline
\nwendcode{}\nwbegindocs{350}Output.
\nwenddocs{}\nwbegincode{351}\moddef{SF mappings: sf endpoint/os mapping: TBP}\endmoddef
  procedure :: write => sf_epo_mapping_write
\nwendcode{}\nwbegincode{352}\moddef{SF mappings: procedures}\plusendmoddef
  subroutine sf_epo_mapping_write (object, unit)
    class(sf_epo_mapping_t), intent(in) :: object
    integer, intent(in), optional :: unit
    integer :: u
    u = given_output_unit (unit)
    write (u, "(1x,A)", advance="no")  "map"
    if (any (object%i /= 0)) then
       write (u, "('(',I0,',',I0,')')", advance="no")  object%i
    end if
    write (u, "(A,F7.5,A,F7.5,A)")  ": ep/on-shell (a = ", object%a, &
         " | ", object%m, ")"
  end subroutine sf_epo_mapping_write
  
\eatline
\nwendcode{}\nwbegindocs{353}Initialize: no extra parameters.
\nwenddocs{}\nwbegincode{354}\moddef{SF mappings: sf endpoint/os mapping: TBP}\plusendmoddef
  procedure :: init => sf_epo_mapping_init
\nwendcode{}\nwbegincode{355}\moddef{SF mappings: procedures}\plusendmoddef
  subroutine sf_epo_mapping_init (mapping, a, m)
    class(sf_epo_mapping_t), intent(out) :: mapping
    real(default), intent(in) :: a, m
    call mapping%base_init (2)
    mapping%a = a
    mapping%m = m
    mapping%lm2 = abs (2 * log (mapping%m))
  end subroutine sf_epo_mapping_init
    
\eatline
\nwendcode{}\nwbegindocs{356}Apply mapping.
\nwenddocs{}\nwbegincode{357}\moddef{SF mappings: sf endpoint/os mapping: TBP}\plusendmoddef
  procedure :: compute => sf_epo_mapping_compute
\nwendcode{}\nwbegincode{358}\moddef{SF mappings: procedures}\plusendmoddef
  subroutine sf_epo_mapping_compute (mapping, r, rb, f, p, pb, x_free)
    class(sf_epo_mapping_t), intent(inout) :: mapping
    real(default), dimension(:), intent(out) :: r, rb
    real(default), intent(out) :: f
    real(default), dimension(:), intent(in) :: p, pb
    real(default), intent(inout), optional :: x_free
    real(default), dimension(2) :: px, r2
    real(default) :: f2
    integer :: j
    px(1) = 0
    call map_endpoint_01 (px(2), f2, p(mapping%i(2)), mapping%a)
    call map_on_shell (r2, f, px, mapping%lm2)
    f = f * f2
    r = p
    rb= pb
    do j = 1, 2
       r (mapping%i(j)) = r2(j)
       rb(mapping%i(j)) = 1 - r2(j)
    end do
  end subroutine sf_epo_mapping_compute

\eatline
\nwendcode{}\nwbegindocs{359}Apply inverse.
\nwenddocs{}\nwbegincode{360}\moddef{SF mappings: sf endpoint/os mapping: TBP}\plusendmoddef
  procedure :: inverse => sf_epo_mapping_inverse
\nwendcode{}\nwbegincode{361}\moddef{SF mappings: procedures}\plusendmoddef
  subroutine sf_epo_mapping_inverse (mapping, r, rb, f, p, pb, x_free)
    class(sf_epo_mapping_t), intent(inout) :: mapping
    real(default), dimension(:), intent(in) :: r, rb
    real(default), intent(out) :: f
    real(default), dimension(:), intent(out) :: p, pb
    real(default), intent(inout), optional :: x_free
    real(default), dimension(2) :: px, p2
    real(default) :: f2
    integer :: j
    call map_on_shell_inverse (r(mapping%i), f, px, mapping%lm2)
    p2(1) = 0
    call map_endpoint_inverse_01 (px(2), f2, p2(2), mapping%a)
    f = f * f2
    p = r
    pb= rb
    do j = 1, 2
       p (mapping%i(j)) = p2(j)
       pb(mapping%i(j)) = 1 - p2(j)
    end do
  end subroutine sf_epo_mapping_inverse

\eatline
\nwendcode{}\nwbegindocs{362}\nwdocspar
\subsection{Implementation: ISR endpoint mapping}
Similar to the endpoint mapping above: This maps the unit square
($r_1,r_2$) such that $p_1$ is the product $r_1r_2$, while $p_2$ is
related to the ratio.  Furthermore, we enhance the region at $r_1=1$
and $r_2=1$, which translates into $p_1=1$ and $p_2=0,1$.  

The enhancement is such that ISR singularity $(1-x)^{-1+\epsilon}$ is
flattened.  This would be easy in one dimension, but becomes
nontrivial in two dimensions.
\nwenddocs{}\nwbegincode{363}\moddef{SF mappings: public}\plusendmoddef
  public :: sf_ip_mapping_t
\nwendcode{}\nwbegincode{364}\moddef{SF mappings: types}\plusendmoddef
  type, extends (sf_mapping_t) :: sf_ip_mapping_t
     real(default) :: eps = 0
   contains
   \LA{}SF mappings: sf power mapping: TBP\RA{}
  end type sf_ip_mapping_t
  
\eatline
\nwendcode{}\nwbegindocs{365}Output.
\nwenddocs{}\nwbegincode{366}\moddef{SF mappings: sf power mapping: TBP}\endmoddef
  procedure :: write => sf_ip_mapping_write
\nwendcode{}\nwbegincode{367}\moddef{SF mappings: procedures}\plusendmoddef
  subroutine sf_ip_mapping_write (object, unit)
    class(sf_ip_mapping_t), intent(in) :: object
    integer, intent(in), optional :: unit
    integer :: u
    u = given_output_unit (unit)
    write (u, "(1x,A)", advance="no")  "map"
    if (any (object%i /= 0)) then
       write (u, "('(',I0,',',I0,')')", advance="no")  object%i
    end if
    write (u, "(A,ES12.5,A)")  ": isr (eps =", object%eps, ")"
  end subroutine sf_ip_mapping_write
  
\eatline
\nwendcode{}\nwbegindocs{368}Initialize: no extra parameters.
\nwenddocs{}\nwbegincode{369}\moddef{SF mappings: sf power mapping: TBP}\plusendmoddef
  procedure :: init => sf_ip_mapping_init
\nwendcode{}\nwbegincode{370}\moddef{SF mappings: procedures}\plusendmoddef
  subroutine sf_ip_mapping_init (mapping, eps)
    class(sf_ip_mapping_t), intent(out) :: mapping
    real(default), intent(in), optional :: eps
    call mapping%base_init (2)
    if (present (eps))  mapping%eps = eps
    if (mapping%eps <= 0) &
         call msg_fatal ("ISR mapping: regulator epsilon must not be zero")
  end subroutine sf_ip_mapping_init
    
\eatline
\nwendcode{}\nwbegindocs{371}Apply mapping.
\nwenddocs{}\nwbegincode{372}\moddef{SF mappings: sf power mapping: TBP}\plusendmoddef
  procedure :: compute => sf_ip_mapping_compute
\nwendcode{}\nwbegincode{373}\moddef{SF mappings: procedures}\plusendmoddef
  subroutine sf_ip_mapping_compute (mapping, r, rb, f, p, pb, x_free)
    class(sf_ip_mapping_t), intent(inout) :: mapping
    real(default), dimension(:), intent(out) :: r, rb
    real(default), intent(out) :: f
    real(default), dimension(:), intent(in) :: p, pb
    real(default), intent(inout), optional :: x_free
    real(default), dimension(2) :: px, pxb, r2, r2b
    real(default) :: f1, f2, xb, y, yb
    integer :: j
    call map_power_1 (xb, f1, pb(mapping%i(1)), 2 * mapping%eps)
    call map_power_01 (y, yb, f2, pb(mapping%i(2)), mapping%eps)
    px(1)  = 1 - xb
    pxb(1) = xb
    px(2)  = y
    pxb(2) = yb
    call map_unit_square_prec (r2, r2b, f, px, pxb)
    f = f * f1 * f2
    r = p
    rb= pb
    do j = 1, 2
       r (mapping%i(j)) = r2 (j)
       rb(mapping%i(j)) = r2b(j)
    end do
  end subroutine sf_ip_mapping_compute

\eatline
\nwendcode{}\nwbegindocs{374}Apply inverse.
\nwenddocs{}\nwbegincode{375}\moddef{SF mappings: sf power mapping: TBP}\plusendmoddef
  procedure :: inverse => sf_ip_mapping_inverse
\nwendcode{}\nwbegincode{376}\moddef{SF mappings: procedures}\plusendmoddef
  subroutine sf_ip_mapping_inverse (mapping, r, rb, f, p, pb, x_free)
    class(sf_ip_mapping_t), intent(inout) :: mapping
    real(default), dimension(:), intent(in) :: r, rb
    real(default), intent(out) :: f
    real(default), dimension(:), intent(out) :: p, pb
    real(default), intent(inout), optional :: x_free
    real(default), dimension(2) :: r2, r2b, px, pxb, p2, p2b
    real(default) :: f1, f2, xb, y, yb
    integer :: j
    do j = 1, 2
       r2 (j) = r (mapping%i(j))
       r2b(j) = rb(mapping%i(j))
    end do
    call map_unit_square_inverse_prec (r2, r2b, f, px, pxb)
    xb = pxb(1)
    if (px(1) > 0) then
       y  = px(2)
       yb = pxb(2)
    else
       y  = 0.5_default
       yb = 0.5_default
    end if
    call map_power_inverse_1 (xb, f1, p2b(1), 2 * mapping%eps)
    call map_power_inverse_01 (y, yb, f2, p2b(2), mapping%eps)
    p2 = 1 - p2b
    f = f * f1 * f2
    p = r
    pb= rb
    do j = 1, 2
       p (mapping%i(j)) = p2(j)
       pb(mapping%i(j)) = p2b(j)
    end do
  end subroutine sf_ip_mapping_inverse

\eatline
\nwendcode{}\nwbegindocs{377}\nwdocspar
\subsection{Implementation: ISR endpoint mapping, resonant}
Similar to the endpoint mapping above: This maps the unit square
($r_1,r_2$) such that $p_1$ is the product $r_1r_2$, while $p_2$ is
related to the ratio.  Furthermore, we enhance the region at $r_1=1$
and $r_2=1$, which translates into $p_1=1$ and $p_2=0,1$.  

The enhancement is such that ISR singularity $(1-x)^{-1+\epsilon}$ is
flattened.  This would be easy in one dimension, but becomes
nontrivial in two dimensions.

The resonance can be turned off by the flag {\tt{}resonance}.
\nwenddocs{}\nwbegincode{378}\moddef{SF mappings: public}\plusendmoddef
  public :: sf_ipr_mapping_t
\nwendcode{}\nwbegincode{379}\moddef{SF mappings: types}\plusendmoddef
  type, extends (sf_mapping_t) :: sf_ipr_mapping_t
     real(default) :: eps = 0
     real(default) :: m = 0
     real(default) :: w = 0
     logical :: resonance = .true.
   contains
   \LA{}SF mappings: sf power/res mapping: TBP\RA{}
  end type sf_ipr_mapping_t
  
\eatline
\nwendcode{}\nwbegindocs{380}Output.
\nwenddocs{}\nwbegincode{381}\moddef{SF mappings: sf power/res mapping: TBP}\endmoddef
  procedure :: write => sf_ipr_mapping_write
\nwendcode{}\nwbegincode{382}\moddef{SF mappings: procedures}\plusendmoddef
  subroutine sf_ipr_mapping_write (object, unit)
    class(sf_ipr_mapping_t), intent(in) :: object
    integer, intent(in), optional :: unit
    integer :: u
    u = given_output_unit (unit)
    write (u, "(1x,A)", advance="no")  "map"
    if (any (object%i /= 0)) then
       write (u, "('(',I0,',',I0,')')", advance="no")  object%i
    end if
    if (object%resonance) then
       write (u, "(A,F7.5,A,F7.5,', ',F7.5,A)")  ": isr/res (eps = ", &
            object%eps, " | ", object%m, object%w, ")"
    else
       write (u, "(A,F7.5,A)")  ": isr/res (eps = ", object%eps, ")"
    end if
  end subroutine sf_ipr_mapping_write
  
\eatline
\nwendcode{}\nwbegindocs{383}Initialize:
\nwenddocs{}\nwbegincode{384}\moddef{SF mappings: sf power/res mapping: TBP}\plusendmoddef
  procedure :: init => sf_ipr_mapping_init
\nwendcode{}\nwbegincode{385}\moddef{SF mappings: procedures}\plusendmoddef
  subroutine sf_ipr_mapping_init (mapping, eps, m, w)
    class(sf_ipr_mapping_t), intent(out) :: mapping
    real(default), intent(in), optional :: eps, m, w
    call mapping%base_init (2)
    if (present (eps))  mapping%eps = eps
    if (mapping%eps <= 0) &
         call msg_fatal ("ISR mapping: regulator epsilon must not be zero")
    if (present (m) .and. present (w)) then
       mapping%m = m
       mapping%w = w
    else
       mapping%resonance = .false.
    end if
  end subroutine sf_ipr_mapping_init
    
\eatline
\nwendcode{}\nwbegindocs{386}Apply mapping.
\nwenddocs{}\nwbegincode{387}\moddef{SF mappings: sf power/res mapping: TBP}\plusendmoddef
  procedure :: compute => sf_ipr_mapping_compute
\nwendcode{}\nwbegincode{388}\moddef{SF mappings: procedures}\plusendmoddef
  subroutine sf_ipr_mapping_compute (mapping, r, rb, f, p, pb, x_free)
    class(sf_ipr_mapping_t), intent(inout) :: mapping
    real(default), dimension(:), intent(out) :: r, rb
    real(default), intent(out) :: f
    real(default), dimension(:), intent(in) :: p, pb
    real(default), intent(inout), optional :: x_free
    real(default), dimension(2) :: px, pxb, r2, r2b
    real(default) :: f1, f2, y, yb
    integer :: j
    if (mapping%resonance) then
       call map_breit_wigner &
            (px(1), f1, p(mapping%i(1)), mapping%m, mapping%w, x_free)
    else
       px(1) = p(mapping%i(1))
       f1 = 1
    end if
    call map_power_01 (y, yb, f2, pb(mapping%i(2)), mapping%eps)
    pxb(1) = 1 - px(1)
    px(2)  = y
    pxb(2) = yb
    call map_unit_square_prec (r2, r2b, f, px, pxb)
    f = f * f1 * f2
    r = p
    rb= pb
    do j = 1, 2
       r (mapping%i(j)) = r2 (j)
       rb(mapping%i(j)) = r2b(j)
    end do
  end subroutine sf_ipr_mapping_compute

\eatline
\nwendcode{}\nwbegindocs{389}Apply inverse.
\nwenddocs{}\nwbegincode{390}\moddef{SF mappings: sf power/res mapping: TBP}\plusendmoddef
  procedure :: inverse => sf_ipr_mapping_inverse
\nwendcode{}\nwbegincode{391}\moddef{SF mappings: procedures}\plusendmoddef
  subroutine sf_ipr_mapping_inverse (mapping, r, rb, f, p, pb, x_free)
    class(sf_ipr_mapping_t), intent(inout) :: mapping
    real(default), dimension(:), intent(in) :: r, rb
    real(default), intent(out) :: f
    real(default), dimension(:), intent(out) :: p, pb
    real(default), intent(inout), optional :: x_free
    real(default), dimension(2) :: r2, r2b, px, pxb, p2, p2b
    real(default) :: f1, f2, y, yb
    integer :: j
    do j = 1, 2
       r2 (j) = r (mapping%i(j))
       r2b(j) = rb(mapping%i(j))
    end do
    call map_unit_square_inverse_prec (r2, r2b, f, px, pxb)
    if (px(1) > 0) then
       y  = px(2)
       yb = pxb(2)
    else
       y  = 0.5_default
       yb = 0.5_default
    end if
    if (mapping%resonance) then
       call map_breit_wigner_inverse &
            (px(1), f1, p2(1), mapping%m, mapping%w, x_free)
    else
       p2(1) = px(1)
       f1 = 1
    end if
    call map_power_inverse_01 (y, yb, f2, p2b(2), mapping%eps)
    p2b(1) = 1 - p2(1)
    p2 (2) = 1 - p2b(2)
    f = f * f1 * f2
    p = r
    pb= rb
    do j = 1, 2
       p (mapping%i(j)) = p2(j)
       pb(mapping%i(j)) = p2b(j)
    end do
  end subroutine sf_ipr_mapping_inverse

\eatline
\nwendcode{}\nwbegindocs{392}\nwdocspar
\subsection{Implementation: ISR on-shell mapping}
Similar to the endpoint mapping above: This maps the unit square
($r_1,r_2$) such that $p_1$ is ignored while the product $r_1r_2$ is
constant.  $p_2$ is related to the ratio.  Furthermore, we enhance the
region at $r_1=1$ and $r_2=1$, which translates into $p_1=1$ and
$p_2=0,1$.

The enhancement is such that ISR singularity $(1-x)^{-1+\epsilon}$ is
flattened.  This would be easy in one dimension, but becomes
nontrivial in two dimensions.
\nwenddocs{}\nwbegincode{393}\moddef{SF mappings: public}\plusendmoddef
  public :: sf_ipo_mapping_t
\nwendcode{}\nwbegincode{394}\moddef{SF mappings: types}\plusendmoddef
  type, extends (sf_mapping_t) :: sf_ipo_mapping_t
     real(default) :: eps = 0
     real(default) :: m = 0
   contains
   \LA{}SF mappings: sf power/os mapping: TBP\RA{}
  end type sf_ipo_mapping_t
  
\eatline
\nwendcode{}\nwbegindocs{395}Output.
\nwenddocs{}\nwbegincode{396}\moddef{SF mappings: sf power/os mapping: TBP}\endmoddef
  procedure :: write => sf_ipo_mapping_write
\nwendcode{}\nwbegincode{397}\moddef{SF mappings: procedures}\plusendmoddef
  subroutine sf_ipo_mapping_write (object, unit)
    class(sf_ipo_mapping_t), intent(in) :: object
    integer, intent(in), optional :: unit
    integer :: u
    u = given_output_unit (unit)
    write (u, "(1x,A)", advance="no")  "map"
    if (any (object%i /= 0)) then
       write (u, "('(',I0,',',I0,')')", advance="no")  object%i
    end if
    write (u, "(A,F7.5,A,F7.5,A)")  ": isr/os (eps = ", object%eps, &
         " | ", object%m, ")"
  end subroutine sf_ipo_mapping_write
  
\eatline
\nwendcode{}\nwbegindocs{398}Initialize: no extra parameters.
\nwenddocs{}\nwbegincode{399}\moddef{SF mappings: sf power/os mapping: TBP}\plusendmoddef
  procedure :: init => sf_ipo_mapping_init
\nwendcode{}\nwbegincode{400}\moddef{SF mappings: procedures}\plusendmoddef
  subroutine sf_ipo_mapping_init (mapping, eps, m)
    class(sf_ipo_mapping_t), intent(out) :: mapping
    real(default), intent(in), optional :: eps, m
    call mapping%base_init (2)
    if (present (eps))  mapping%eps = eps
    if (mapping%eps <= 0) &
         call msg_fatal ("ISR mapping: regulator epsilon must not be zero")
    mapping%m = m
  end subroutine sf_ipo_mapping_init
    
\eatline
\nwendcode{}\nwbegindocs{401}Apply mapping.
\nwenddocs{}\nwbegincode{402}\moddef{SF mappings: sf power/os mapping: TBP}\plusendmoddef
  procedure :: compute => sf_ipo_mapping_compute
\nwendcode{}\nwbegincode{403}\moddef{SF mappings: procedures}\plusendmoddef
  subroutine sf_ipo_mapping_compute (mapping, r, rb, f, p, pb, x_free)
    class(sf_ipo_mapping_t), intent(inout) :: mapping
    real(default), dimension(:), intent(out) :: r, rb
    real(default), intent(out) :: f
    real(default), dimension(:), intent(in) :: p, pb
    real(default), intent(inout), optional :: x_free
    real(default), dimension(2) :: px, pxb, r2, r2b
    real(default) :: f1, f2, y, yb
    integer :: j
    call map_power_01 (y, yb, f2, pb(mapping%i(2)), mapping%eps)
    px(1)  = mapping%m ** 2
    if (present (x_free))  px(1) = px(1) / x_free
    pxb(1) = 1 - px(1)
    px(2)  = y
    pxb(2) = yb
    call map_unit_square_prec (r2, r2b, f1, px, pxb)
    f = f1 * f2
    r = p
    rb= pb
    do j = 1, 2
       r (mapping%i(j)) = r2 (j)
       rb(mapping%i(j)) = r2b(j)
    end do
  end subroutine sf_ipo_mapping_compute

\eatline
\nwendcode{}\nwbegindocs{404}Apply inverse.
\nwenddocs{}\nwbegincode{405}\moddef{SF mappings: sf power/os mapping: TBP}\plusendmoddef
  procedure :: inverse => sf_ipo_mapping_inverse
\nwendcode{}\nwbegincode{406}\moddef{SF mappings: procedures}\plusendmoddef
  subroutine sf_ipo_mapping_inverse (mapping, r, rb, f, p, pb, x_free)
    class(sf_ipo_mapping_t), intent(inout) :: mapping
    real(default), dimension(:), intent(in) :: r, rb
    real(default), intent(out) :: f
    real(default), dimension(:), intent(out) :: p, pb
    real(default), intent(inout), optional :: x_free
    real(default), dimension(2) :: r2, r2b, px, pxb, p2, p2b
    real(default) :: f1, f2, y, yb
    integer :: j
    do j = 1, 2
       r2 (j) = r (mapping%i(j))
       r2b(j) = rb(mapping%i(j))
    end do
    call map_unit_square_inverse_prec (r2, r2b, f1, px, pxb)
    y  = px(2)
    yb = pxb(2)
    call map_power_inverse_01 (y, yb, f2, p2b(2), mapping%eps)
    p2(1) = 0
    p2b(1)= 1
    p2(2) = 1 - p2b(2)
    f = f1 * f2
    p = r
    pb= rb
    do j = 1, 2
       p (mapping%i(j)) = p2(j)
       pb(mapping%i(j)) = p2b(j)
    end do
  end subroutine sf_ipo_mapping_inverse

\eatline
\nwendcode{}\nwbegindocs{407}\nwdocspar
\subsection{Implementation: Endpoint + ISR power mapping}
This is a combination of endpoint (i.e., beamstrahlung) and ISR power
mapping.  The first two parameters apply to the beamstrahlung
spectrum, the last two to the ISR function for the first and second
beam, respectively.
\nwenddocs{}\nwbegincode{408}\moddef{SF mappings: public}\plusendmoddef
  public :: sf_ei_mapping_t
\nwendcode{}\nwbegincode{409}\moddef{SF mappings: types}\plusendmoddef
  type, extends (sf_mapping_t) :: sf_ei_mapping_t
     type(sf_ep_mapping_t) :: ep
     type(sf_ip_mapping_t) :: ip
   contains
   \LA{}SF mappings: sf ep-ip mapping: TBP\RA{}
  end type sf_ei_mapping_t
  
\eatline
\nwendcode{}\nwbegindocs{410}Output.
\nwenddocs{}\nwbegincode{411}\moddef{SF mappings: sf ep-ip mapping: TBP}\endmoddef
  procedure :: write => sf_ei_mapping_write
\nwendcode{}\nwbegincode{412}\moddef{SF mappings: procedures}\plusendmoddef
  subroutine sf_ei_mapping_write (object, unit)
    class(sf_ei_mapping_t), intent(in) :: object
    integer, intent(in), optional :: unit
    integer :: u
    u = given_output_unit (unit)
    write (u, "(1x,A)", advance="no")  "map"
    if (any (object%i /= 0)) then
       write (u, "('(',I0,3(',',I0),')')", advance="no")  object%i
    end if
    write (u, "(A,ES12.5,A,ES12.5,A)")  ": ep/isr (a =", object%ep%a, &
         ", eps =", object%ip%eps, ")"
  end subroutine sf_ei_mapping_write
  
\eatline
\nwendcode{}\nwbegindocs{413}Initialize: no extra parameters.
\nwenddocs{}\nwbegincode{414}\moddef{SF mappings: sf ep-ip mapping: TBP}\plusendmoddef
  procedure :: init => sf_ei_mapping_init
\nwendcode{}\nwbegincode{415}\moddef{SF mappings: procedures}\plusendmoddef
  subroutine sf_ei_mapping_init (mapping, a, eps)
    class(sf_ei_mapping_t), intent(out) :: mapping
    real(default), intent(in), optional :: a, eps
    call mapping%base_init (4)
    call mapping%ep%init (a)
    call mapping%ip%init (eps)
  end subroutine sf_ei_mapping_init
    
\eatline
\nwendcode{}\nwbegindocs{416}Set an index value.  We should communicate the appropriate indices to the
enclosed sub-mappings, therefore override the method.
\nwenddocs{}\nwbegincode{417}\moddef{SF mappings: sf ep-ip mapping: TBP}\plusendmoddef
  procedure :: set_index => sf_ei_mapping_set_index
\nwendcode{}\nwbegincode{418}\moddef{SF mappings: procedures}\plusendmoddef
  subroutine sf_ei_mapping_set_index (mapping, j, i)
    class(sf_ei_mapping_t), intent(inout) :: mapping
    integer, intent(in) :: j, i
    mapping%i(j) = i
    select case (j)
    case (1:2);  call mapping%ep%set_index (j, i)
    case (3:4);  call mapping%ip%set_index (j-2, i)
    end select
  end subroutine sf_ei_mapping_set_index
    
\eatline
\nwendcode{}\nwbegindocs{419}Apply mapping.  Now, the beamstrahlung and ISR mappings are
independent of each other.  The parameter subsets that are actually
used should not overlap.  The Jacobians are multiplied.
\nwenddocs{}\nwbegincode{420}\moddef{SF mappings: sf ep-ip mapping: TBP}\plusendmoddef
  procedure :: compute => sf_ei_mapping_compute
\nwendcode{}\nwbegincode{421}\moddef{SF mappings: procedures}\plusendmoddef
  subroutine sf_ei_mapping_compute (mapping, r, rb, f, p, pb, x_free)
    class(sf_ei_mapping_t), intent(inout) :: mapping
    real(default), dimension(:), intent(out) :: r, rb
    real(default), intent(out) :: f
    real(default), dimension(:), intent(in) :: p, pb
    real(default), intent(inout), optional :: x_free
    real(default), dimension(size(p)) :: q, qb
    real(default) :: f1, f2
    call mapping%ep%compute (q, qb, f1, p, pb, x_free)
    call mapping%ip%compute (r, rb, f2, q, qb, x_free)
    f = f1 * f2
  end subroutine sf_ei_mapping_compute

\eatline
\nwendcode{}\nwbegindocs{422}Apply inverse.
\nwenddocs{}\nwbegincode{423}\moddef{SF mappings: sf ep-ip mapping: TBP}\plusendmoddef
  procedure :: inverse => sf_ei_mapping_inverse
\nwendcode{}\nwbegincode{424}\moddef{SF mappings: procedures}\plusendmoddef
  subroutine sf_ei_mapping_inverse (mapping, r, rb, f, p, pb, x_free)
    class(sf_ei_mapping_t), intent(inout) :: mapping
    real(default), dimension(:), intent(in) :: r, rb
    real(default), intent(out) :: f
    real(default), dimension(:), intent(out) :: p, pb
    real(default), intent(inout), optional :: x_free
    real(default), dimension(size(p)) :: q, qb
    real(default) :: f1, f2
    call mapping%ip%inverse (r, rb, f2, q, qb, x_free)
    call mapping%ep%inverse (q, qb, f1, p, pb, x_free)
    f = f1 * f2
  end subroutine sf_ei_mapping_inverse

\eatline
\nwendcode{}\nwbegindocs{425}\nwdocspar
\subsection{Implementation: Endpoint + ISR + resonance}
This is a combination of endpoint (i.e., beamstrahlung) and ISR power
mapping, adapted for an s-channel resonance.  The first two internal
parameters apply to the beamstrahlung spectrum, the last two to the
ISR function for the first and second beam, respectively.  The first
and third parameters are the result of an overall resonance mapping,
so on the outside, the first parameter is the total momentum fraction,
the third one describes the distribution between beamstrahlung and ISR.
\nwenddocs{}\nwbegincode{426}\moddef{SF mappings: public}\plusendmoddef
  public :: sf_eir_mapping_t
\nwendcode{}\nwbegincode{427}\moddef{SF mappings: types}\plusendmoddef
  type, extends (sf_mapping_t) :: sf_eir_mapping_t
     type(sf_res_mapping_t) :: res
     type(sf_epr_mapping_t) :: ep
     type(sf_ipr_mapping_t) :: ip
   contains
   \LA{}SF mappings: sf ep-ip-res mapping: TBP\RA{}
  end type sf_eir_mapping_t
  
\eatline
\nwendcode{}\nwbegindocs{428}Output.
\nwenddocs{}\nwbegincode{429}\moddef{SF mappings: sf ep-ip-res mapping: TBP}\endmoddef
  procedure :: write => sf_eir_mapping_write
\nwendcode{}\nwbegincode{430}\moddef{SF mappings: procedures}\plusendmoddef
  subroutine sf_eir_mapping_write (object, unit)
    class(sf_eir_mapping_t), intent(in) :: object
    integer, intent(in), optional :: unit
    integer :: u
    u = given_output_unit (unit)
    write (u, "(1x,A)", advance="no")  "map"
    if (any (object%i /= 0)) then
       write (u, "('(',I0,3(',',I0),')')", advance="no")  object%i
    end if
    write (u, "(A,F7.5,A,F7.5,A,F7.5,', ',F7.5,A)")  &
         ": ep/isr/res (a =", object%ep%a, &
         ", eps =", object%ip%eps, " | ", object%res%m, object%res%w, ")"
  end subroutine sf_eir_mapping_write
  
\eatline
\nwendcode{}\nwbegindocs{431}Initialize: no extra parameters.
\nwenddocs{}\nwbegincode{432}\moddef{SF mappings: sf ep-ip-res mapping: TBP}\plusendmoddef
  procedure :: init => sf_eir_mapping_init
\nwendcode{}\nwbegincode{433}\moddef{SF mappings: procedures}\plusendmoddef
  subroutine sf_eir_mapping_init (mapping, a, eps, m, w)
    class(sf_eir_mapping_t), intent(out) :: mapping
    real(default), intent(in) :: a, eps, m, w
    call mapping%base_init (4)
    call mapping%res%init (m, w)
    call mapping%ep%init (a)
    call mapping%ip%init (eps)
  end subroutine sf_eir_mapping_init
    
\eatline
\nwendcode{}\nwbegindocs{434}Set an index value.  We should communicate the appropriate indices to the
enclosed sub-mappings, therefore override the method.
\nwenddocs{}\nwbegincode{435}\moddef{SF mappings: sf ep-ip-res mapping: TBP}\plusendmoddef
  procedure :: set_index => sf_eir_mapping_set_index
\nwendcode{}\nwbegincode{436}\moddef{SF mappings: procedures}\plusendmoddef
  subroutine sf_eir_mapping_set_index (mapping, j, i)
    class(sf_eir_mapping_t), intent(inout) :: mapping
    integer, intent(in) :: j, i
    mapping%i(j) = i
    select case (j)
    case (1);  call mapping%res%set_index (1, i)
    case (3);  call mapping%res%set_index (2, i)
    end select
    select case (j)
    case (1:2);  call mapping%ep%set_index (j, i)
    case (3:4);  call mapping%ip%set_index (j-2, i)
    end select
  end subroutine sf_eir_mapping_set_index
    
\eatline
\nwendcode{}\nwbegindocs{437}Apply mapping.  Now, the beamstrahlung and ISR mappings are
independent of each other.  The parameter subsets that are actually
used should not overlap.  The Jacobians are multiplied.
\nwenddocs{}\nwbegincode{438}\moddef{SF mappings: sf ep-ip-res mapping: TBP}\plusendmoddef
  procedure :: compute => sf_eir_mapping_compute
\nwendcode{}\nwbegincode{439}\moddef{SF mappings: procedures}\plusendmoddef
  subroutine sf_eir_mapping_compute (mapping, r, rb, f, p, pb, x_free)
    class(sf_eir_mapping_t), intent(inout) :: mapping
    real(default), dimension(:), intent(out) :: r, rb
    real(default), intent(out) :: f
    real(default), dimension(:), intent(in) :: p, pb
    real(default), intent(inout), optional :: x_free
    real(default), dimension(size(p)) :: px, pxb, q, qb
    real(default) :: f0, f1, f2
    call mapping%res%compute (px, pxb, f0, p, pb, x_free)
    call mapping%ep%compute (q, qb, f1, px, pxb, x_free)
    call mapping%ip%compute (r, rb, f2, q, qb, x_free)
    f = f0 * f1 * f2
  end subroutine sf_eir_mapping_compute

\eatline
\nwendcode{}\nwbegindocs{440}Apply inverse.
\nwenddocs{}\nwbegincode{441}\moddef{SF mappings: sf ep-ip-res mapping: TBP}\plusendmoddef
  procedure :: inverse => sf_eir_mapping_inverse
\nwendcode{}\nwbegincode{442}\moddef{SF mappings: procedures}\plusendmoddef
  subroutine sf_eir_mapping_inverse (mapping, r, rb, f, p, pb, x_free)
    class(sf_eir_mapping_t), intent(inout) :: mapping
    real(default), dimension(:), intent(in) :: r, rb
    real(default), intent(out) :: f
    real(default), dimension(:), intent(out) :: p, pb
    real(default), intent(inout), optional :: x_free
    real(default), dimension(size(p)) :: px, pxb, q, qb
    real(default) :: f0, f1, f2
    call mapping%ip%inverse (r, rb, f2, q, qb, x_free)
    call mapping%ep%inverse (q, qb, f1, px, pxb, x_free)
    call mapping%res%inverse (px, pxb, f0, p, pb, x_free)
    f = f0 * f1 * f2
  end subroutine sf_eir_mapping_inverse

\eatline
\nwendcode{}\nwbegindocs{443}\nwdocspar
\subsection{Implementation: Endpoint + ISR power mapping, on-shell}
This is a combination of endpoint (i.e., beamstrahlung) and ISR power
mapping.  The first two parameters apply to the beamstrahlung
spectrum, the last two to the ISR function for the first and second
beam, respectively.  On top of that, we map the first and third parameter
such that the product is constant.  From the outside, the first
parameter is irrelevant while the third parameter describes the
distribution of energy (loss) among beamstrahlung and ISR.
\nwenddocs{}\nwbegincode{444}\moddef{SF mappings: public}\plusendmoddef
  public :: sf_eio_mapping_t
\nwendcode{}\nwbegincode{445}\moddef{SF mappings: types}\plusendmoddef
  type, extends (sf_mapping_t) :: sf_eio_mapping_t
     type(sf_os_mapping_t) :: os
     type(sf_epr_mapping_t) :: ep
     type(sf_ipr_mapping_t) :: ip
   contains
   \LA{}SF mappings: sf ep-ip-os mapping: TBP\RA{}
  end type sf_eio_mapping_t
  
\eatline
\nwendcode{}\nwbegindocs{446}Output.
\nwenddocs{}\nwbegincode{447}\moddef{SF mappings: sf ep-ip-os mapping: TBP}\endmoddef
  procedure :: write => sf_eio_mapping_write
\nwendcode{}\nwbegincode{448}\moddef{SF mappings: procedures}\plusendmoddef
  subroutine sf_eio_mapping_write (object, unit)
    class(sf_eio_mapping_t), intent(in) :: object
    integer, intent(in), optional :: unit
    integer :: u
    u = given_output_unit (unit)
    write (u, "(1x,A)", advance="no")  "map"
    if (any (object%i /= 0)) then
       write (u, "('(',I0,3(',',I0),')')", advance="no")  object%i
    end if
    write (u, "(A,F7.5,A,F7.5,A,F7.5,A)")  ": ep/isr/os (a =", object%ep%a, &
         ", eps =", object%ip%eps, " | ", object%os%m, ")"
  end subroutine sf_eio_mapping_write
  
\eatline
\nwendcode{}\nwbegindocs{449}Initialize: no extra parameters.
\nwenddocs{}\nwbegincode{450}\moddef{SF mappings: sf ep-ip-os mapping: TBP}\plusendmoddef
  procedure :: init => sf_eio_mapping_init
\nwendcode{}\nwbegincode{451}\moddef{SF mappings: procedures}\plusendmoddef
  subroutine sf_eio_mapping_init (mapping, a, eps, m)
    class(sf_eio_mapping_t), intent(out) :: mapping
    real(default), intent(in), optional :: a, eps, m
    call mapping%base_init (4)
    call mapping%os%init (m)
    call mapping%ep%init (a)
    call mapping%ip%init (eps)
  end subroutine sf_eio_mapping_init
    
\eatline
\nwendcode{}\nwbegindocs{452}Set an index value.  We should communicate the appropriate indices to the
enclosed sub-mappings, therefore override the method.
\nwenddocs{}\nwbegincode{453}\moddef{SF mappings: sf ep-ip-os mapping: TBP}\plusendmoddef
  procedure :: set_index => sf_eio_mapping_set_index
\nwendcode{}\nwbegincode{454}\moddef{SF mappings: procedures}\plusendmoddef
  subroutine sf_eio_mapping_set_index (mapping, j, i)
    class(sf_eio_mapping_t), intent(inout) :: mapping
    integer, intent(in) :: j, i
    mapping%i(j) = i
    select case (j)
    case (1);  call mapping%os%set_index (1, i)
    case (3);  call mapping%os%set_index (2, i)
    end select
    select case (j)
    case (1:2);  call mapping%ep%set_index (j, i)
    case (3:4);  call mapping%ip%set_index (j-2, i)
    end select
  end subroutine sf_eio_mapping_set_index
    
\eatline
\nwendcode{}\nwbegindocs{455}Apply mapping.  Now, the beamstrahlung and ISR mappings are
independent of each other.  The parameter subsets that are actually
used should not overlap.  The Jacobians are multiplied.
\nwenddocs{}\nwbegincode{456}\moddef{SF mappings: sf ep-ip-os mapping: TBP}\plusendmoddef
  procedure :: compute => sf_eio_mapping_compute
\nwendcode{}\nwbegincode{457}\moddef{SF mappings: procedures}\plusendmoddef
  subroutine sf_eio_mapping_compute (mapping, r, rb, f, p, pb, x_free)
    class(sf_eio_mapping_t), intent(inout) :: mapping
    real(default), dimension(:), intent(out) :: r, rb
    real(default), intent(out) :: f
    real(default), dimension(:), intent(in) :: p, pb
    real(default), intent(inout), optional :: x_free
    real(default), dimension(size(p)) :: px, pxb, q, qb
    real(default) :: f0, f1, f2
    call mapping%os%compute (px, pxb, f0, p, pb, x_free)
    call mapping%ep%compute (q, qb, f1, px, pxb, x_free)
    call mapping%ip%compute (r, rb, f2, q, qb, x_free)
    f = f0 * f1 * f2
  end subroutine sf_eio_mapping_compute

\eatline
\nwendcode{}\nwbegindocs{458}Apply inverse.
\nwenddocs{}\nwbegincode{459}\moddef{SF mappings: sf ep-ip-os mapping: TBP}\plusendmoddef
  procedure :: inverse => sf_eio_mapping_inverse
\nwendcode{}\nwbegincode{460}\moddef{SF mappings: procedures}\plusendmoddef
  subroutine sf_eio_mapping_inverse (mapping, r, rb, f, p, pb, x_free)
    class(sf_eio_mapping_t), intent(inout) :: mapping
    real(default), dimension(:), intent(in) :: r, rb
    real(default), intent(out) :: f
    real(default), dimension(:), intent(out) :: p, pb
    real(default), intent(inout), optional :: x_free
    real(default), dimension(size(p)) :: px, pxb, q, qb
    real(default) :: f0, f1, f2
    call mapping%ip%inverse (r, rb, f2, q, qb, x_free)
    call mapping%ep%inverse (q, qb, f1, px, pxb, x_free)
    call mapping%os%inverse (px, pxb, f0, p, pb, x_free)
    f = f0 * f1 * f2
  end subroutine sf_eio_mapping_inverse

\eatline
\nwendcode{}\nwbegindocs{461}\nwdocspar
\subsection{Basic formulas}
\subsubsection{Standard mapping of the unit square}
This mapping of the unit square is appropriate in particular for
structure functions which are concentrated at the lower end.  Instead
of a rectangular grid, one set of grid lines corresponds to constant
parton c.m. energy.  The other set is chosen such that the jacobian is
only mildly singular ($\ln x$ which is zero at $x=1$), corresponding
to an initial concentration of sampling points at the maximum energy.
If {\tt{}power} is greater than one (the default), points are also
concentrated at the lower end.

The formula is ({\tt{}power}=$\alpha$):
\begin{align}
  r_1 &= (p_1 ^ {p_2})^\alpha \\
  r_2 &= (p_1 ^ {1 - p_2})^\alpha\\
  f   &= \alpha^2 p_1 ^ {\alpha - 1} |\log p_1|
\end{align}
and for the default case $\alpha=1$:
\begin{align}
  r_1 &= p_1 ^ {p_2} \\
  r_2 &= p_1 ^ {1 - p_2} \\
  f   &= |\log p_1|
\end{align}
\nwenddocs{}\nwbegincode{462}\moddef{SF mappings: procedures}\plusendmoddef
  subroutine map_unit_square (r, factor, p, power)
    real(default), dimension(2), intent(out) :: r
    real(default), intent(out) :: factor
    real(default), dimension(2), intent(in) :: p
    real(default), intent(in), optional :: power
    real(default) :: xx, yy
    factor = 1
    xx = p(1)
    yy = p(2)
    if (present(power)) then
       if (p(1) > 0 .and. power > 1) then
          xx = p(1)**power
          factor = factor * power * xx / p(1)
       end if
    end if
    if (xx /= 0) then
       r(1) = xx ** yy
       r(2) = xx / r(1)
       factor = factor * abs (log (xx))
    else
       r = 0
    end if
  end subroutine map_unit_square

\eatline
\nwendcode{}\nwbegindocs{463}This is the inverse mapping.
\nwenddocs{}\nwbegincode{464}\moddef{SF mappings: procedures}\plusendmoddef
  subroutine map_unit_square_inverse (r, factor, p, power)
    real(kind=default), dimension(2), intent(in) :: r
    real(kind=default), intent(out) :: factor
    real(kind=default), dimension(2), intent(out) :: p
    real(kind=default), intent(in), optional :: power
    real(kind=default) :: lg, xx, yy
    factor = 1
    xx = r(1) * r(2)
    if (xx /= 0) then
       lg = log (xx)
       if (lg /= 0) then
          yy = log (r(1)) / lg
       else
          yy = 0
       end if
       p(2) = yy
       factor = factor * abs (lg)
       if (present(power)) then
          p(1) = xx**(1._default/power)
          factor = factor * power * xx / p(1)
       else
          p(1) = xx
       end if
    else
       p = 0
    end if
  end subroutine map_unit_square_inverse

\eatline
\nwendcode{}\nwbegindocs{465}\nwdocspar
\subsubsection{Precise mapping of the unit square}
A more precise version (with unit power parameter).  This version
should be numerically stable near $x=1$ and $y=0,1$.  The formulas are again
\begin{equation}
  r_1 = p_1^{p_2}, \qquad
  r_2 = p_1^{\bar p_2}, \qquad
  f = - \log p_1
\end{equation}
but we compute both $r_i$ and $\bar r_i$ simultaneously and make
direct use of both $p_i$ and $\bar p_i$ as appropriate.
\nwenddocs{}\nwbegincode{466}\moddef{SF mappings: procedures}\plusendmoddef
  subroutine map_unit_square_prec (r, rb, factor, p, pb)
    real(default), dimension(2), intent(out) :: r
    real(default), dimension(2), intent(out) :: rb
    real(default), intent(out) :: factor
    real(default), dimension(2), intent(in) :: p
    real(default), dimension(2), intent(in) :: pb
    if (p(1) > 0.5_default) then
       call compute_prec_xy_1 (r(1), rb(1), p(1), pb(1), p (2))
       call compute_prec_xy_1 (r(2), rb(2), p(1), pb(1), pb(2))
       factor = - log_prec (p(1), pb(1))
    else if (p(1) > 0) then
       call compute_prec_xy_0 (r(1), rb(1), p(1), pb(1), p (2))
       call compute_prec_xy_0 (r(2), rb(2), p(1), pb(1), pb(2))
       factor = - log_prec (p(1), pb(1))
    else
       r  = 0
       rb = 1
       factor = 0
    end if
  end subroutine map_unit_square_prec

\eatline
\nwendcode{}\nwbegindocs{467}This is the inverse mapping.
\nwenddocs{}\nwbegincode{468}\moddef{SF mappings: procedures}\plusendmoddef
  subroutine map_unit_square_inverse_prec (r, rb, factor, p, pb)
    real(default), dimension(2), intent(in) :: r
    real(default), dimension(2), intent(in) :: rb
    real(default), intent(out) :: factor
    real(default), dimension(2), intent(out) :: p
    real(default), dimension(2), intent(out) :: pb
    call inverse_prec_x (r, rb, p(1), pb(1))
    if (all (r > 0)) then
       if (rb(1) < rb(2)) then
          call inverse_prec_y (r, rb, p(2), pb(2))
       else
          call inverse_prec_y ([r(2),r(1)], [rb(2),rb(1)], pb(2), p(2))
       end if
       factor = - log_prec (p(1), pb(1))
    else
       p(1)  = 0
       pb(1) = 1
       p(2)  = 0.5_default
       pb(2) = 0.5_default
       factor = 0
    end if
  end subroutine map_unit_square_inverse_prec

\eatline
\nwendcode{}\nwbegindocs{469}This is an auxiliary function: evaluate the expression $\bar z = 1 -
x^y$ in a numerically stable way.  Instabilities occur for $y=0$ and
$x=1$.  The idea is to replace the bracket by the first terms of its
Taylor expansion around $x=1$ (read $\bar x\equiv 1 -x$)
\begin{equation}
  1 - x^y = y\bar x\left(1 + \frac12(1-y)\bar x +
    \frac16(2-y)(1-y)\bar x^2\right)
\end{equation}
whenever this is the better approximation.  Actually, the relative
numerical error of the exact formula is about $\eta/(y\bar x)$ where
$\eta$ is given by {\tt{}epsilon(KIND)} in Fortran.  The relative error
of the approximation is better than the last included term divided by
$(y\bar x)$.

The first subroutine computes $z$ and $\bar z$ near $x=1$ where $\log
x$ should be expanded, the second one near $x=0$ where $\log x$ can be
kept.
\nwenddocs{}\nwbegincode{470}\moddef{SF mappings: procedures}\plusendmoddef
  subroutine compute_prec_xy_1 (z, zb, x, xb, y)
    real(default), intent(out) :: z, zb
    real(default), intent(in) :: x, xb, y
    real(default) :: a1, a2, a3
    a1 = y * xb
    a2 = a1 * (1 - y) * xb / 2
    a3 = a2 * (2 - y) * xb / 3
    if (abs (a3) < epsilon (a3)) then
       zb = a1 + a2 + a3
       z = 1 - zb
    else
       z = x ** y
       zb = 1 - z
    end if
  end subroutine compute_prec_xy_1
   
  subroutine compute_prec_xy_0 (z, zb, x, xb, y)
    real(default), intent(out) :: z, zb
    real(default), intent(in) :: x, xb, y
    real(default) :: a1, a2, a3, lx
    lx = -log (x)
    a1 = y * lx
    a2 = a1 * y * lx / 2
    a3 = a2 * y * lx / 3
    if (abs (a3) < epsilon (a3)) then
       zb = a1 + a2 + a3
       z = 1 - zb
    else
       z = x ** y
       zb = 1 - z
    end if
  end subroutine compute_prec_xy_0
   
\eatline
\eatline
\nwendcode{}\nwbegindocs{471}For the inverse calculation, we evaluate $x=r_1r_2$ in a stable way.
Since it is just a polynomial, the expansion near $x=1$ is
analytically exact, and we don't need to choose based on precision.
\nwenddocs{}\nwbegincode{472}\moddef{SF mappings: procedures}\plusendmoddef
  subroutine inverse_prec_x (r, rb, x, xb)
    real(default), dimension(2), intent(in) :: r, rb
    real(default), intent(out) :: x, xb
    real(default) :: a0, a1
    a0 = rb(1) + rb(2)
    a1 = rb(1) * rb(2)
    if (a0 > 0.5_default) then
       xb = a0 - a1
       x = 1 - xb
    else
       x = r(1) * r(2)
       xb = 1 - x
    end if
  end subroutine inverse_prec_x
    
\eatline
\nwendcode{}\nwbegindocs{473}The inverse calculation for the relative momentum fraction
\begin{equation}
  y = \frac{1}{1 + \frac{\log{r_2}}{\log{r_1}}}
\end{equation}
is slightly more complicated.  We should take the precise form of the
logarithm, so we are safe near $r_i=1$.  A series expansion is
required if $r_1\ll r_2$, since then $y$ becomes small.  (We assume
$r_1<r_2$ here; for the opposite case, the arguments can be
exchanged.)
\nwenddocs{}\nwbegincode{474}\moddef{SF mappings: procedures}\plusendmoddef
  subroutine inverse_prec_y (r, rb, y, yb)
    real(default), dimension(2), intent(in) :: r, rb
    real(default), intent(out) :: y, yb
    real(default) :: log1, log2, a1, a2, a3
    log1 = log_prec (r(1), rb(1))
    log2 = log_prec (r(2), rb(2))
    a1 = - rb(1) / log2
    a2 = - rb(1) ** 2 * (1 / log2**2 + 1 / (2 * log2))
    a3 = - rb(1) ** 3 * (1 / log2**3 + 1 / log2**2 + 1 / (3 * log2))
    if (abs (a3) < epsilon (a3)) then
       y  = a1 + a2 + a3
       yb = 1 - y
    else
       y  = 1 / (1 + log2/log1)
       yb = 1 / (1 + log1/log2)
    end if
  end subroutine inverse_prec_y
  
\eatline
\nwendcode{}\nwbegindocs{475}We also need an evaluation of $\log x$ which is stable near $x=1$.
\nwenddocs{}\nwbegincode{476}\moddef{SF mappings: public}\plusendmoddef
  public :: log_prec
\nwendcode{}\nwbegincode{477}\moddef{SF mappings: procedures}\plusendmoddef
  function log_prec (x, xb) result (lx)
    real(default), intent(in) :: x, xb
    real(default) :: a1, a2, a3, lx
    a1 = xb
    a2 = a1 * xb / 2
    a3 = a2 * xb * 2 / 3
    if (abs (a3) < epsilon (a3)) then
       lx = - a1 - a2 - a3
    else
       lx = log (x)
    end if
  end function log_prec
  
\eatline
\nwendcode{}\nwbegindocs{478}\nwdocspar
\subsubsection{Mapping for on-shell s-channel}
The limiting case, if the product $r_1r_2$ is fixed for on-shell
production.  The parameter $p_1$ is ignored.  In the inverse mapping,
it is returned zero.

The parameter {\tt{}x{\char95}free}, if present, rescales the total energy.  If
it is less than one, the rescaled mass parameter $m^2$ should be increased
accordingly.
\nwenddocs{}\nwbegincode{479}\moddef{SF mappings: procedures}\plusendmoddef
  subroutine map_on_shell (r, factor, p, lm2, x_free)
    real(default), dimension(2), intent(out) :: r
    real(default), intent(out) :: factor
    real(default), dimension(2), intent(in) :: p
    real(default), intent(in) :: lm2
    real(default), intent(in), optional :: x_free
    real(default) :: lx
    lx = lm2;  if (present (x_free))  lx = lx + log (x_free)
    r(1) = exp (- p(2) * lx)
    r(2) = exp (- (1 - p(2)) * lx)
    factor = lx
  end subroutine map_on_shell

  subroutine map_on_shell_inverse (r, factor, p, lm2, x_free)
    real(default), dimension(2), intent(in) :: r
    real(default), intent(out) :: factor
    real(default), dimension(2), intent(out) :: p
    real(default), intent(in) :: lm2
    real(default), intent(in), optional :: x_free
    real(default) :: lx
    lx = lm2;  if (present (x_free))  lx = lx + log (x_free)
    p(1) = 0
    p(2) = abs (log (r(1))) / lx
    factor = lx
  end subroutine map_on_shell_inverse
    
\eatline
\eatline
\nwendcode{}\nwbegindocs{480}This is the standard Breit-Wigner mapping.  We apply it to a single
variable, independently of or in addition to a unit-square mapping.  We
assume here that the limits for the variable are 0 and 1, and that the
mass $m$ and width $w$ are rescaled appropriately, so they are
dimensionless and usually between 0 and 1.

If {\tt{}x{\char95}free} is set, it rescales the total energy and thus mass and
width, since these are defined with respect to the total energy.
\nwenddocs{}\nwbegincode{481}\moddef{SF mappings: procedures}\plusendmoddef
  subroutine map_breit_wigner (r, factor, p, m, w, x_free)
    real(default), intent(out) :: r
    real(default), intent(out) :: factor
    real(default), intent(in) :: p
    real(default), intent(in) :: m
    real(default), intent(in) :: w
    real(default), intent(in), optional :: x_free
    real(default) :: m2, mw, a1, a2, a3, z, tmp
    m2 = m ** 2
    mw = m * w
    if (present (x_free)) then
       m2 = m2 / x_free
       mw = mw / x_free
    end if
    a1 = atan (- m2 / mw)
    a2 = atan ((1 - m2) / mw)
    a3 = (a2 - a1) * mw
    z = (1-p) * a1 + p * a2
    if (-pi/2 < z .and. z < pi/2) then
       tmp = tan (z)
       r = max (m2 + mw * tmp, 0._default)
       factor = a3 * (1 + tmp ** 2) 
    else
       r = 0
       factor = 0
    end if
  end subroutine map_breit_wigner
    
  subroutine map_breit_wigner_inverse (r, factor, p, m, w, x_free)
    real(default), intent(in) :: r
    real(default), intent(out) :: factor
    real(default), intent(out) :: p
    real(default), intent(in) :: m
    real(default), intent(in) :: w
    real(default) :: m2, mw, a1, a2, a3, tmp
    real(default), intent(in), optional :: x_free
    m2 = m ** 2
    mw = m * w
    if (present (x_free)) then
       m2 = m2 / x_free
       mw = mw / x_free
    end if
    a1 = atan (- m2 / mw)
    a2 = atan ((1 - m2) / mw)
    a3 = (a2 - a1) * mw
    tmp = (r - m2) / mw
    p = (atan (tmp) - a1) / (a2 - a1)
    factor = a3 * (1 + tmp ** 2)
  end subroutine map_breit_wigner_inverse

\eatline
\eatline
\nwendcode{}\nwbegindocs{482}\nwdocspar
\subsubsection{Mapping with endpoint enhancement}
This is a mapping which is close to the unit mapping, except that at
the endpoint(s), the output values are exponentially enhanced.
\begin{equation}
  y = \tanh (a \tan (\frac{\pi}{2}x))
\end{equation}
We have two variants: one covers endpoints at $0$ and $1$
symmetrically, while the other one (which essentially maps one-half of
the range), covers only the endpoint at $1$.
\nwenddocs{}\nwbegincode{483}\moddef{SF mappings: procedures}\plusendmoddef
  subroutine map_endpoint_1 (x3, factor, x1, a)
    real(default), intent(out) :: x3, factor
    real(default), intent(in) :: x1
    real(default), intent(in) :: a
    real(default) :: x2
    if (abs (x1) < 1) then
       x2 = tan (x1 * pi / 2)
       x3 = tanh (a * x2)
       factor = a * pi/2 * (1 + x2 ** 2) * (1 - x3 ** 2)
    else
       x3 = x1
       factor = 0
    end if
  end subroutine map_endpoint_1

  subroutine map_endpoint_inverse_1 (x3, factor, x1, a)
    real(default), intent(in) :: x3
    real(default), intent(out) :: x1, factor
    real(default), intent(in) :: a
    real(default) :: x2
    if (abs (x3) < 1) then
       x2 = atanh (x3) / a
       x1 = 2 / pi * atan (x2)
       factor = a * pi/2 * (1 + x2 ** 2) * (1 - x3 ** 2)
    else
       x1 = x3
       factor = 0
    end if
  end subroutine map_endpoint_inverse_1

  subroutine map_endpoint_01 (x4, factor, x0, a)
    real(default), intent(out) :: x4, factor
    real(default), intent(in) :: x0
    real(default), intent(in) :: a
    real(default) :: x1, x3
    x1 = 2 * x0 - 1
    call map_endpoint_1 (x3, factor, x1, a)
    x4 = (x3 + 1) / 2
  end subroutine map_endpoint_01

  subroutine map_endpoint_inverse_01 (x4, factor, x0, a)
    real(default), intent(in) :: x4
    real(default), intent(out) :: x0, factor
    real(default), intent(in) :: a
    real(default) :: x1, x3
    x3 = 2 * x4 - 1
    call map_endpoint_inverse_1 (x3, factor, x1, a)
    x0 = (x1 + 1) / 2
  end subroutine map_endpoint_inverse_01

\eatline
\eatline
\eatline
\eatline
\nwendcode{}\nwbegindocs{484}\nwdocspar
\subsubsection{Mapping with endpoint enhancement (ISR)}
This is another endpoint mapping.  It is designed to flatten the ISR
singularity which is of power type at $x=1$, i.e., if
\begin{equation}
  \sigma = \int_0^1 dx\,f(x)\,G(x)
  = \int_0^1 dx\,\epsilon(1-x)^{-1+\epsilon} G(x),
\end{equation}
we replace this by
\begin{equation}
  r = x^\epsilon \quad\Longrightarrow\quad
  \sigma = \int_0^1 dr\,G(1- (1-r)^{1/\epsilon}).
\end{equation}
We expect that $\epsilon$ is small.

The actual mapping is $r\to x$ (so $x$ emerges closer to $1$).  The
Jacobian that we return is thus $1/f(x)$.  We compute the mapping in
terms of $\bar x\equiv 1 - x$, so we can achieve the required precision.
Because some compilers show quite wild numeric fluctuations, we
internally convert numeric types to explicit {\tt{}double} precision.
\nwenddocs{}\nwbegincode{485}\moddef{SF mappings: public}\plusendmoddef
  public :: map_power_1
  public :: map_power_inverse_1
\nwendcode{}\nwbegincode{486}\moddef{SF mappings: procedures}\plusendmoddef
  subroutine map_power_1 (xb, factor, rb, eps)
    real(default), intent(out) :: xb, factor
    real(default), intent(in) :: rb
    real(double) :: rb_db, factor_db, eps_db, xb_db
    real(default), intent(in) :: eps
    rb_db = real (rb, kind=double)
    eps_db = real (eps, kind=double)
    xb_db = rb_db ** (1 / eps_db)
    if (rb_db > 0) then
       factor_db = xb_db / rb_db / eps_db
       factor = real (factor_db, kind=default)
    else
       factor = 0
    end if
    xb = real (xb_db, kind=default)
  end subroutine map_power_1
    
  subroutine map_power_inverse_1 (xb, factor, rb, eps)
    real(default), intent(in) :: xb
    real(default), intent(out) :: rb, factor
    real(double) :: xb_db, factor_db, eps_db, rb_db
    real(default), intent(in) :: eps
    xb_db = real (xb, kind=double)
    eps_db = real (eps, kind=double)
    rb_db = xb_db ** eps_db
    if (xb_db > 0) then
       factor_db = xb_db / rb_db / eps_db
       factor = real (factor_db, kind=default)
    else
       factor = 0
    end if
    rb = real (rb_db, kind=default)
  end subroutine map_power_inverse_1
    
\eatline
\eatline
\nwendcode{}\nwbegindocs{487}Here we apply a power mapping to both endpoints.  We divide the
interval in two equal halves and apply the power mapping for the
nearest endpoint, either $0$ or $1$.
\nwenddocs{}\nwbegincode{488}\moddef{SF mappings: procedures}\plusendmoddef
  subroutine map_power_01 (y, yb, factor, r, eps)
    real(default), intent(out) :: y, yb, factor
    real(default), intent(in) :: r
    real(default), intent(in) :: eps
    real(default) :: u, ub, zp, zm
    u = 2 * r - 1
    if (u > 0) then
       ub = 2 * (1 - r)
       call map_power_1 (zm, factor, ub, eps)
       zp = 2 - zm
    else if (u < 0) then
       ub = 2 * r
       call map_power_1 (zp, factor, ub, eps)
       zm = 2 - zp
    else
       factor = 1 / eps
       zp = 1
       zm = 1
    end if
    y  = zp / 2
    yb = zm / 2
  end subroutine map_power_01
    
  subroutine map_power_inverse_01 (y, yb, factor, r, eps)
    real(default), intent(in) :: y, yb
    real(default), intent(out) :: r, factor
    real(default), intent(in) :: eps
    real(default) :: ub, zp, zm
    zp = 2 * y
    zm = 2 * yb
    if (zm < zp) then
       call map_power_inverse_1 (zm, factor, ub, eps)
       r = 1 - ub / 2
    else if (zp < zm) then
       call map_power_inverse_1 (zp, factor, ub, eps)
       r = ub / 2
    else
       factor = 1 / eps
       ub = 1
       r = ub / 2
    end if
  end subroutine map_power_inverse_01
    
\eatline
\eatline
\nwendcode{}\nwbegindocs{489}\nwdocspar
\subsubsection{Structure-function channels}
A structure-function chain parameterization (channel) may contain a
mapping that applies to multiple structure functions.  This is
described by an extension of the {\tt{}sf{\char95}mapping{\char95}t} type.  In addition,
it may contain mappings that apply to (other) individual structure
functions.  The details of these mappings are implementation-specific.

The {\tt{}sf{\char95}channel{\char95}t} type combines this information.  It contains an
array of map codes, one for each structure-function entry.  The code
values are:
\begin{description}
\item[none] MC input parameters $r$ directly become energy fractions $x$
\item[single] default mapping for a single structure-function entry
\item[multi/s] map $r\to x$ such that one MC input parameter is $\hat s/s$
\item[multi/resonance] as before, adapted to s-channel resonance
\item[multi/on-shell] as before, adapted to an on-shell particle in
  the s channel
\item[multi/endpoint] like multi/s, but enhance the region near $r_i=1$
\item[multi/endpoint/res] endpoint mapping with resonance
\item[multi/endpoint/os] endpoint mapping for on-shell
\item[multi/power/os] like multi/endpoint, regulating a power singularity
\end{description}
\nwenddocs{}\nwbegincode{490}\moddef{SF mappings: parameters}\endmoddef
  integer, parameter :: SFMAP_NONE = 0
  integer, parameter :: SFMAP_SINGLE = 1
  integer, parameter :: SFMAP_MULTI_S = 2
  integer, parameter :: SFMAP_MULTI_RES = 3
  integer, parameter :: SFMAP_MULTI_ONS = 4
  integer, parameter :: SFMAP_MULTI_EP = 5
  integer, parameter :: SFMAP_MULTI_EPR = 6
  integer, parameter :: SFMAP_MULTI_EPO = 7
  integer, parameter :: SFMAP_MULTI_IP = 8
  integer, parameter :: SFMAP_MULTI_IPR = 9
  integer, parameter :: SFMAP_MULTI_IPO = 10
  integer, parameter :: SFMAP_MULTI_EI = 11

\eatline
\eatline
\eatline
\eatline
\eatline
\nwendcode{}\nwbegindocs{491}Then, it contains an allocatable entry for the multi mapping.  This
entry holds the MC-parameter indices on which the mapping applies
(there may be more than one MC parameter per structure-function entry)
and any parameters associated with the mapping.

There can be only one multi-mapping per channel.
\nwenddocs{}\nwbegincode{492}\moddef{SF mappings: public}\plusendmoddef
  public :: sf_channel_t
\nwendcode{}\nwbegincode{493}\moddef{SF mappings: types}\plusendmoddef
  type :: sf_channel_t
     integer, dimension(:), allocatable :: map_code
     class(sf_mapping_t), allocatable :: multi_mapping
   contains
   \LA{}SF mappings: sf channel: TBP\RA{}
  end type sf_channel_t
  
\eatline
\nwendcode{}\nwbegindocs{494}The output format prints a single character for each
structure-function entry and, if applicable, an account of the mapping
parameters.
\nwenddocs{}\nwbegincode{495}\moddef{SF mappings: sf channel: TBP}\endmoddef
  procedure :: write => sf_channel_write
\nwendcode{}\nwbegincode{496}\moddef{SF mappings: procedures}\plusendmoddef
  subroutine sf_channel_write (object, unit)
    class(sf_channel_t), intent(in) :: object
    integer, intent(in), optional :: unit
    integer :: u, i
    u = given_output_unit (unit)
    if (allocated (object%map_code)) then
       do i = 1, size (object%map_code)
          select case (object%map_code (i))
          case (SFMAP_NONE)
             write (u, "(1x,A)", advance="no") "-"
          case (SFMAP_SINGLE)
             write (u, "(1x,A)", advance="no") "+"
          case (SFMAP_MULTI_S)
             write (u, "(1x,A)", advance="no") "s"
          case (SFMAP_MULTI_RES)
             write (u, "(1x,A)", advance="no") "r"
          case (SFMAP_MULTI_ONS)
             write (u, "(1x,A)", advance="no") "o"
          case (SFMAP_MULTI_EP)
             write (u, "(1x,A)", advance="no") "e"
          case (SFMAP_MULTI_EPR)
             write (u, "(1x,A)", advance="no") "p"
          case (SFMAP_MULTI_EPO)
             write (u, "(1x,A)", advance="no") "q"
          case (SFMAP_MULTI_IP)
             write (u, "(1x,A)", advance="no") "i"
          case (SFMAP_MULTI_IPR)
             write (u, "(1x,A)", advance="no") "i"
          case (SFMAP_MULTI_IPO)
             write (u, "(1x,A)", advance="no") "i"
          case (SFMAP_MULTI_EI)
             write (u, "(1x,A)", advance="no") "i"
          end select
       end do
    else
       write (u, "(1x,A)", advance="no") "-"
    end if
    if (allocated (object%multi_mapping)) then
       write (u, "(1x,'/')", advance="no")
       call object%multi_mapping%write (u)
    else
       write (u, *)
    end if
  end subroutine sf_channel_write
       
\eatline
\nwendcode{}\nwbegindocs{497}Initializer for a single {\tt{}sf{\char95}channel} object.
\nwenddocs{}\nwbegincode{498}\moddef{SF mappings: sf channel: TBP}\plusendmoddef
  procedure :: init => sf_channel_init
\nwendcode{}\nwbegincode{499}\moddef{SF mappings: procedures}\plusendmoddef
  subroutine sf_channel_init (channel, n_strfun)
    class(sf_channel_t), intent(out) :: channel
    integer, intent(in) :: n_strfun
    allocate (channel%map_code (n_strfun))
    channel%map_code = SFMAP_NONE
  end subroutine sf_channel_init
    
\eatline
\nwendcode{}\nwbegindocs{500}Assignment.  This merely copies intrinsic assignment, but apparently
the latter is bugged in gfortran 4.6.3, causing memory corruption.
\nwenddocs{}\nwbegincode{501}\moddef{SF mappings: sf channel: TBP}\plusendmoddef
  generic :: assignment (=) => sf_channel_assign
  procedure :: sf_channel_assign
\nwendcode{}\nwbegincode{502}\moddef{SF mappings: procedures}\plusendmoddef
  subroutine sf_channel_assign (copy, original)
    class(sf_channel_t), intent(out) :: copy
    type(sf_channel_t), intent(in) :: original
    allocate (copy%map_code (size (original%map_code)))
    copy%map_code = original%map_code
    if (allocated (original%multi_mapping)) then
       allocate (copy%multi_mapping, source = original%multi_mapping)
    end if
  end subroutine sf_channel_assign
  
\eatline
\nwendcode{}\nwbegindocs{503}This initializer allocates an array of channels with common number of
structure-function entries, therefore it is not a type-bound procedure.
\nwenddocs{}\nwbegincode{504}\moddef{SF mappings: public}\plusendmoddef
  public :: allocate_sf_channels
\nwendcode{}\nwbegincode{505}\moddef{SF mappings: procedures}\plusendmoddef
  subroutine allocate_sf_channels (channel, n_channel, n_strfun)
    type(sf_channel_t), dimension(:), intent(out), allocatable :: channel
    integer, intent(in) :: n_channel
    integer, intent(in) :: n_strfun
    integer :: c
    allocate (channel (n_channel))
    do c = 1, n_channel
       call channel(c)%init (n_strfun)
    end do
  end subroutine allocate_sf_channels

\eatline
\nwendcode{}\nwbegindocs{506}This marks a given subset of indices as single-mapping.
\nwenddocs{}\nwbegincode{507}\moddef{SF mappings: sf channel: TBP}\plusendmoddef
  procedure :: activate_mapping => sf_channel_activate_mapping
\nwendcode{}\nwbegincode{508}\moddef{SF mappings: procedures}\plusendmoddef
  subroutine sf_channel_activate_mapping (channel, i_sf)
    class(sf_channel_t), intent(inout) :: channel
    integer, dimension(:), intent(in) :: i_sf
    channel%map_code(i_sf) = SFMAP_SINGLE
  end subroutine sf_channel_activate_mapping
  
\eatline
\nwendcode{}\nwbegindocs{509}This sets an s-channel multichannel mapping.  The length of the
{\tt{}i{\char95}sf} array must be 2.  The parameter indices are not yet set.
\nwenddocs{}\nwbegincode{510}\moddef{SF mappings: sf channel: TBP}\plusendmoddef
  procedure :: set_s_mapping => sf_channel_set_s_mapping
\nwendcode{}\nwbegincode{511}\moddef{SF mappings: procedures}\plusendmoddef
  subroutine sf_channel_set_s_mapping (channel, i_sf, power)
    class(sf_channel_t), intent(inout) :: channel
    integer, dimension(:), intent(in) :: i_sf
    real(default), intent(in), optional :: power
    channel%map_code(i_sf) = SFMAP_MULTI_S
    allocate (sf_s_mapping_t :: channel%multi_mapping)
    select type (mapping => channel%multi_mapping)
    type is (sf_s_mapping_t)
       call mapping%init (power)
    end select
  end subroutine sf_channel_set_s_mapping

\eatline
\nwendcode{}\nwbegindocs{512}This sets an s-channel resonance multichannel mapping.
\nwenddocs{}\nwbegincode{513}\moddef{SF mappings: sf channel: TBP}\plusendmoddef
  procedure :: set_res_mapping => sf_channel_set_res_mapping
\nwendcode{}\nwbegincode{514}\moddef{SF mappings: procedures}\plusendmoddef
  subroutine sf_channel_set_res_mapping (channel, i_sf, m, w)
    class(sf_channel_t), intent(inout) :: channel
    integer, dimension(:), intent(in) :: i_sf
    real(default), intent(in) :: m, w
    channel%map_code(i_sf) = SFMAP_MULTI_RES
    allocate (sf_res_mapping_t :: channel%multi_mapping)
    select type (mapping => channel%multi_mapping)
    type is (sf_res_mapping_t)
       call mapping%init (m, w)
    end select
  end subroutine sf_channel_set_res_mapping

\eatline
\nwendcode{}\nwbegindocs{515}This sets an s-channel on-shell multichannel mapping.  The length of the
{\tt{}i{\char95}sf} array must be 2.  (The first parameter actually becomes an
irrelevant dummy.)
\nwenddocs{}\nwbegincode{516}\moddef{SF mappings: sf channel: TBP}\plusendmoddef
  procedure :: set_os_mapping => sf_channel_set_os_mapping
\nwendcode{}\nwbegincode{517}\moddef{SF mappings: procedures}\plusendmoddef
  subroutine sf_channel_set_os_mapping (channel, i_sf, m)
    class(sf_channel_t), intent(inout) :: channel
    integer, dimension(:), intent(in) :: i_sf
    real(default), intent(in) :: m
    channel%map_code(i_sf) = SFMAP_MULTI_ONS
    allocate (sf_os_mapping_t :: channel%multi_mapping)
    select type (mapping => channel%multi_mapping)
    type is (sf_os_mapping_t)
       call mapping%init (m)
    end select
  end subroutine sf_channel_set_os_mapping

\eatline
\nwendcode{}\nwbegindocs{518}This sets an s-channel endpoint mapping.  The parameter $a$ is the
slope parameter (default 1); increasing it moves the endpoint region
(at $x=1$ to lower values in the input parameter.
region even more.
\nwenddocs{}\nwbegincode{519}\moddef{SF mappings: sf channel: TBP}\plusendmoddef
  procedure :: set_ep_mapping => sf_channel_set_ep_mapping
\nwendcode{}\nwbegincode{520}\moddef{SF mappings: procedures}\plusendmoddef
  subroutine sf_channel_set_ep_mapping (channel, i_sf, a)
    class(sf_channel_t), intent(inout) :: channel
    integer, dimension(:), intent(in) :: i_sf
    real(default), intent(in), optional :: a
    channel%map_code(i_sf) = SFMAP_MULTI_EP
    allocate (sf_ep_mapping_t :: channel%multi_mapping)
    select type (mapping => channel%multi_mapping)
    type is (sf_ep_mapping_t)
       call mapping%init (a = a)
    end select
  end subroutine sf_channel_set_ep_mapping

\eatline
\nwendcode{}\nwbegindocs{521}This sets a resonant endpoint mapping.
\nwenddocs{}\nwbegincode{522}\moddef{SF mappings: sf channel: TBP}\plusendmoddef
  procedure :: set_epr_mapping => sf_channel_set_epr_mapping
\nwendcode{}\nwbegincode{523}\moddef{SF mappings: procedures}\plusendmoddef
  subroutine sf_channel_set_epr_mapping (channel, i_sf, a, m, w)
    class(sf_channel_t), intent(inout) :: channel
    integer, dimension(:), intent(in) :: i_sf
    real(default), intent(in) :: a, m, w
    channel%map_code(i_sf) = SFMAP_MULTI_EPR
    allocate (sf_epr_mapping_t :: channel%multi_mapping)
    select type (mapping => channel%multi_mapping)
    type is (sf_epr_mapping_t)
       call mapping%init (a, m, w)
    end select
  end subroutine sf_channel_set_epr_mapping

\eatline
\nwendcode{}\nwbegindocs{524}This sets an on-shell endpoint mapping.
\nwenddocs{}\nwbegincode{525}\moddef{SF mappings: sf channel: TBP}\plusendmoddef
  procedure :: set_epo_mapping => sf_channel_set_epo_mapping
\nwendcode{}\nwbegincode{526}\moddef{SF mappings: procedures}\plusendmoddef
  subroutine sf_channel_set_epo_mapping (channel, i_sf, a, m)
    class(sf_channel_t), intent(inout) :: channel
    integer, dimension(:), intent(in) :: i_sf
    real(default), intent(in) :: a, m
    channel%map_code(i_sf) = SFMAP_MULTI_EPO
    allocate (sf_epo_mapping_t :: channel%multi_mapping)
    select type (mapping => channel%multi_mapping)
    type is (sf_epo_mapping_t)
       call mapping%init (a, m)
    end select
  end subroutine sf_channel_set_epo_mapping

\eatline
\nwendcode{}\nwbegindocs{527}This sets an s-channel power mapping, regulating a singularity of
type $(1-x)^{-1+\epsilon}$.  The parameter $\epsilon$ depends on the
structure function.
\nwenddocs{}\nwbegincode{528}\moddef{SF mappings: sf channel: TBP}\plusendmoddef
  procedure :: set_ip_mapping => sf_channel_set_ip_mapping
\nwendcode{}\nwbegincode{529}\moddef{SF mappings: procedures}\plusendmoddef
  subroutine sf_channel_set_ip_mapping (channel, i_sf, eps)
    class(sf_channel_t), intent(inout) :: channel
    integer, dimension(:), intent(in) :: i_sf
    real(default), intent(in), optional :: eps
    channel%map_code(i_sf) = SFMAP_MULTI_IP
    allocate (sf_ip_mapping_t :: channel%multi_mapping)
    select type (mapping => channel%multi_mapping)
    type is (sf_ip_mapping_t)
       call mapping%init (eps)
    end select
  end subroutine sf_channel_set_ip_mapping

\eatline
\nwendcode{}\nwbegindocs{530}This sets an s-channel resonant power mapping, regulating a
singularity of type $(1-x)^{-1+\epsilon}$ in the presence of an
s-channel resonance.  The parameter $\epsilon$ depends on the
structure function.
\nwenddocs{}\nwbegincode{531}\moddef{SF mappings: sf channel: TBP}\plusendmoddef
  procedure :: set_ipr_mapping => sf_channel_set_ipr_mapping
\nwendcode{}\nwbegincode{532}\moddef{SF mappings: procedures}\plusendmoddef
  subroutine sf_channel_set_ipr_mapping (channel, i_sf, eps, m, w)
    class(sf_channel_t), intent(inout) :: channel
    integer, dimension(:), intent(in) :: i_sf
    real(default), intent(in), optional :: eps, m, w
    channel%map_code(i_sf) = SFMAP_MULTI_IPR
    allocate (sf_ipr_mapping_t :: channel%multi_mapping)
    select type (mapping => channel%multi_mapping)
    type is (sf_ipr_mapping_t)
       call mapping%init (eps, m, w)
    end select
  end subroutine sf_channel_set_ipr_mapping

\eatline
\nwendcode{}\nwbegindocs{533}This sets an on-shell power mapping, regulating a
singularity of type $(1-x)^{-1+\epsilon}$ for the production of a
single on-shell particle..  The parameter $\epsilon$ depends on the
structure function.
\nwenddocs{}\nwbegincode{534}\moddef{SF mappings: sf channel: TBP}\plusendmoddef
  procedure :: set_ipo_mapping => sf_channel_set_ipo_mapping
\nwendcode{}\nwbegincode{535}\moddef{SF mappings: procedures}\plusendmoddef
  subroutine sf_channel_set_ipo_mapping (channel, i_sf, eps, m)
    class(sf_channel_t), intent(inout) :: channel
    integer, dimension(:), intent(in) :: i_sf
    real(default), intent(in), optional :: eps, m
    channel%map_code(i_sf) = SFMAP_MULTI_IPO
    allocate (sf_ipo_mapping_t :: channel%multi_mapping)
    select type (mapping => channel%multi_mapping)
    type is (sf_ipo_mapping_t)
       call mapping%init (eps, m)
    end select
  end subroutine sf_channel_set_ipo_mapping

\eatline
\nwendcode{}\nwbegindocs{536}This sets a combined endpoint/ISR mapping.
\nwenddocs{}\nwbegincode{537}\moddef{SF mappings: sf channel: TBP}\plusendmoddef
  procedure :: set_ei_mapping => sf_channel_set_ei_mapping
\nwendcode{}\nwbegincode{538}\moddef{SF mappings: procedures}\plusendmoddef
  subroutine sf_channel_set_ei_mapping (channel, i_sf, a, eps)
    class(sf_channel_t), intent(inout) :: channel
    integer, dimension(:), intent(in) :: i_sf
    real(default), intent(in), optional :: a, eps
    channel%map_code(i_sf) = SFMAP_MULTI_EI
    allocate (sf_ei_mapping_t :: channel%multi_mapping)
    select type (mapping => channel%multi_mapping)
    type is (sf_ei_mapping_t)
       call mapping%init (a, eps)
    end select
  end subroutine sf_channel_set_ei_mapping

\eatline
\nwendcode{}\nwbegindocs{539}This sets a combined endpoint/ISR mapping with resonance.
\nwenddocs{}\nwbegincode{540}\moddef{SF mappings: sf channel: TBP}\plusendmoddef
  procedure :: set_eir_mapping => sf_channel_set_eir_mapping
\nwendcode{}\nwbegincode{541}\moddef{SF mappings: procedures}\plusendmoddef
  subroutine sf_channel_set_eir_mapping (channel, i_sf, a, eps, m, w)
    class(sf_channel_t), intent(inout) :: channel
    integer, dimension(:), intent(in) :: i_sf
    real(default), intent(in), optional :: a, eps, m, w
    channel%map_code(i_sf) = SFMAP_MULTI_EI
    allocate (sf_eir_mapping_t :: channel%multi_mapping)
    select type (mapping => channel%multi_mapping)
    type is (sf_eir_mapping_t)
       call mapping%init (a, eps, m, w)
    end select
  end subroutine sf_channel_set_eir_mapping

\eatline
\nwendcode{}\nwbegindocs{542}This sets a combined endpoint/ISR mapping, on-shell.
\nwenddocs{}\nwbegincode{543}\moddef{SF mappings: sf channel: TBP}\plusendmoddef
  procedure :: set_eio_mapping => sf_channel_set_eio_mapping
\nwendcode{}\nwbegincode{544}\moddef{SF mappings: procedures}\plusendmoddef
  subroutine sf_channel_set_eio_mapping (channel, i_sf, a, eps, m)
    class(sf_channel_t), intent(inout) :: channel
    integer, dimension(:), intent(in) :: i_sf
    real(default), intent(in), optional :: a, eps, m
    channel%map_code(i_sf) = SFMAP_MULTI_EI
    allocate (sf_eio_mapping_t :: channel%multi_mapping)
    select type (mapping => channel%multi_mapping)
    type is (sf_eio_mapping_t)
       call mapping%init (a, eps, m)
    end select
  end subroutine sf_channel_set_eio_mapping

\eatline
\nwendcode{}\nwbegindocs{545}Return true if the mapping code at position {\tt{}i{\char95}sf} is {\tt{}SFMAP{\char95}SINGLE}.
\nwenddocs{}\nwbegincode{546}\moddef{SF mappings: sf channel: TBP}\plusendmoddef
  procedure :: is_single_mapping => sf_channel_is_single_mapping
\nwendcode{}\nwbegincode{547}\moddef{SF mappings: procedures}\plusendmoddef
  function sf_channel_is_single_mapping (channel, i_sf) result (flag)
    class(sf_channel_t), intent(in) :: channel
    integer, intent(in) :: i_sf
    logical :: flag
    flag = channel%map_code(i_sf) == SFMAP_SINGLE
  end function sf_channel_is_single_mapping 

\eatline
\nwendcode{}\nwbegindocs{548}Return true if the mapping code at position {\tt{}i{\char95}sf} is any of the
{\tt{}SFMAP{\char95}MULTI} mappings.
\nwenddocs{}\nwbegincode{549}\moddef{SF mappings: sf channel: TBP}\plusendmoddef
  procedure :: is_multi_mapping => sf_channel_is_multi_mapping
\nwendcode{}\nwbegincode{550}\moddef{SF mappings: procedures}\plusendmoddef
  function sf_channel_is_multi_mapping (channel, i_sf) result (flag)
    class(sf_channel_t), intent(in) :: channel
    integer, intent(in) :: i_sf
    logical :: flag
    select case (channel%map_code(i_sf))
    case (SFMAP_NONE, SFMAP_SINGLE)
       flag = .false.
    case default
       flag = .true.
    end select
  end function sf_channel_is_multi_mapping 

\eatline
\nwendcode{}\nwbegindocs{551}Return true if there is any nontrivial mapping in any of the channels.

Note: we provide an explicit public function.  gfortran 4.6.3 has
problems with the alternative implementation as a type-bound
procedure for an array base object.
\nwenddocs{}\nwbegincode{552}\moddef{SF mappings: public}\plusendmoddef
  public :: any_sf_channel_has_mapping
\nwendcode{}\nwbegincode{553}\moddef{SF mappings: procedures}\plusendmoddef
  function any_sf_channel_has_mapping (channel) result (flag)
    type(sf_channel_t), dimension(:), intent(in) :: channel
    logical :: flag
    integer :: c
    flag = .false.
    do c = 1, size (channel)
       flag = flag .or. any (channel(c)%map_code /= SFMAP_NONE)
    end do
  end function any_sf_channel_has_mapping
  
\eatline
\nwendcode{}\nwbegindocs{554}Set a parameter index for an active multi mapping.  We assume that
the index array is allocated properly.
\nwenddocs{}\nwbegincode{555}\moddef{SF mappings: sf channel: TBP}\plusendmoddef
  procedure :: set_par_index => sf_channel_set_par_index
\nwendcode{}\nwbegincode{556}\moddef{SF mappings: procedures}\plusendmoddef
  subroutine sf_channel_set_par_index (channel, j, i_par)
    class(sf_channel_t), intent(inout) :: channel
    integer, intent(in) :: j
    integer, intent(in) :: i_par
    call channel%multi_mapping%set_index (j, i_par)
  end subroutine sf_channel_set_par_index
    
\eatline
\nwendcode{}\nwbegindocs{557}\nwdocspar
\subsection{Unit tests}
\nwenddocs{}\nwbegincode{558}\moddef{SF mappings: public}\plusendmoddef
  public :: sf_mappings_test
\nwendcode{}\nwbegincode{559}\moddef{SF mappings: tests}\endmoddef
  subroutine sf_mappings_test (u, results)
    integer, intent(in) :: u
    type(test_results_t), intent(inout) :: results
  \LA{}SF mappings: execute tests\RA{}
  end subroutine sf_mappings_test
  
\eatline
\nwendcode{}\nwbegindocs{560}\nwdocspar
\subsubsection{Check standard mapping}
Probe the standard mapping of the unit square for different parameter
values.  Also calculates integrals.  For a finite number of bins, they differ
slightly from $1$, but the result is well-defined because we are not using
random points.
\nwenddocs{}\nwbegincode{561}\moddef{SF mappings: execute tests}\endmoddef
  call test (sf_mappings_1, "sf_mappings_1", &
       "standard pair mapping", &
       u, results)
\nwendcode{}\nwbegincode{562}\moddef{SF mappings: tests}\plusendmoddef
  subroutine sf_mappings_1 (u)
    integer, intent(in) :: u
    class(sf_mapping_t), allocatable :: mapping
    real(default), dimension(2) :: p
    
    write (u, "(A)")  "* Test output: sf_mappings_1"
    write (u, "(A)")  "*   Purpose: probe standard mapping"
    write (u, "(A)")
    
    allocate (sf_s_mapping_t :: mapping)
    select type (mapping)
    type is (sf_s_mapping_t)
       call mapping%init ()
       call mapping%set_index (1, 1)
       call mapping%set_index (2, 2)
    end select
       
    call mapping%write (u)
    
    write (u, *)
    write (u, "(A)")  "Probe at (0,0):"
    p = [0._default, 0._default]
    call mapping%check (u, p, 1-p, "F7.5")

    write (u, *)
    write (u, "(A)")  "Probe at (0.5,0.5):"
    p = [0.5_default, 0.5_default]
    call mapping%check (u, p, 1-p, "F7.5")

    write (u, *)
    write (u, "(A)")  "Probe at (0.1,0.5):"
    p = [0.1_default, 0.5_default]
    call mapping%check (u, p, 1-p, "F7.5")

    write (u, *)
    write (u, "(A)")  "Probe at (0.1,0.1):"
    p = [0.1_default, 0.1_default]
    call mapping%check (u, p, 1-p, "F7.5")

    write (u, *)
    write (u, "(A)")  "Compute integral:"
    write (u, "(3x,A,1x,F7.5)") "I =", mapping%integral (100000)

    deallocate (mapping)
    allocate (sf_s_mapping_t :: mapping)
    select type (mapping)
    type is (sf_s_mapping_t)
       call mapping%init (power=2._default)
       call mapping%set_index (1, 1)
       call mapping%set_index (2, 2)
    end select
       
    write (u, *)
    call mapping%write (u)
    
    write (u, *)
    write (u, "(A)")  "Probe at (0,0):"
    p = [0._default, 0._default]
    call mapping%check (u, p, 1-p, "F7.5")

    write (u, *)
    write (u, "(A)")  "Probe at (0.5,0.5):"
    p = [0.5_default, 0.5_default]
    call mapping%check (u, p, 1-p, "F7.5")

    write (u, *)
    write (u, "(A)")  "Probe at (0.1,0.5):"
    p = [0.1_default, 0.5_default]
    call mapping%check (u, p, 1-p, "F7.5")

    write (u, *)
    write (u, "(A)")  "Probe at (0.1,0.1):"
    p = [0.1_default, 0.1_default]
    call mapping%check (u, p, 1-p, "F7.5")

    write (u, *)
    write (u, "(A)")  "Compute integral:"
    write (u, "(3x,A,1x,F7.5)")  "I =", mapping%integral (100000)

    write (u, "(A)")
    write (u, "(A)")  "* Test output end: sf_mappings_1"

  end subroutine sf_mappings_1
  
\eatline
\nwendcode{}\nwbegindocs{563}\nwdocspar
\subsubsection{Channel entries}
Construct channel entries and print them.
\nwenddocs{}\nwbegincode{564}\moddef{SF mappings: execute tests}\plusendmoddef
  call test (sf_mappings_2, "sf_mappings_2", &
       "structure-function mapping channels", &
       u, results)
\nwendcode{}\nwbegincode{565}\moddef{SF mappings: tests}\plusendmoddef
  subroutine sf_mappings_2 (u)
    integer, intent(in) :: u
    type(sf_channel_t), dimension(:), allocatable :: channel
    integer :: c
    
    write (u, "(A)")  "* Test output: sf_mappings_2"
    write (u, "(A)")  "*   Purpose: construct and display &
         &mapping-channel objects"
    write (u, "(A)")
    
    call allocate_sf_channels (channel, n_channel = 6, n_strfun = 2)
    call channel(2)%activate_mapping ([1])
    call channel(3)%set_s_mapping ([1,2])
    call channel(4)%set_s_mapping ([1,2], power=2._default)
    call channel(5)%set_res_mapping ([1,2], m = 0.5_default, w = 0.1_default)
    call channel(6)%set_os_mapping ([1,2], m = 0.5_default)
    
    call channel(3)%set_par_index (1, 1)
    call channel(3)%set_par_index (2, 4)

    call channel(4)%set_par_index (1, 1)
    call channel(4)%set_par_index (2, 4)

    call channel(5)%set_par_index (1, 1)
    call channel(5)%set_par_index (2, 3)

    call channel(6)%set_par_index (1, 1)
    call channel(6)%set_par_index (2, 2)

    do c = 1, size (channel)
       write (u, "(I0,':')", advance="no")  c
       call channel(c)%write (u)
    end do

    write (u, "(A)")
    write (u, "(A)")  "* Test output end: sf_mappings_2"

  end subroutine sf_mappings_2
  
\eatline
\nwendcode{}\nwbegindocs{566}\nwdocspar
\subsubsection{Check resonance mapping}
Probe the resonance mapping of the unit square for different parameter
values.  Also calculates integrals.  For a finite number of bins, they differ
slightly from $1$, but the result is well-defined because we are not using
random points.

The resonance mass is at $1/2$ the energy, the width is $1/10$.
\nwenddocs{}\nwbegincode{567}\moddef{SF mappings: execute tests}\plusendmoddef
  call test (sf_mappings_3, "sf_mappings_3", &
       "resonant pair mapping", &
       u, results)
\nwendcode{}\nwbegincode{568}\moddef{SF mappings: tests}\plusendmoddef
  subroutine sf_mappings_3 (u)
    integer, intent(in) :: u
    class(sf_mapping_t), allocatable :: mapping
    real(default), dimension(2) :: p
    
    write (u, "(A)")  "* Test output: sf_mappings_3"
    write (u, "(A)")  "*   Purpose: probe resonance pair mapping"
    write (u, "(A)")
    
    allocate (sf_res_mapping_t :: mapping)
    select type (mapping)
    type is (sf_res_mapping_t)
       call mapping%init (0.5_default, 0.1_default)
       call mapping%set_index (1, 1)
       call mapping%set_index (2, 2)
    end select
       
    call mapping%write (u)
    
    write (u, *)
    write (u, "(A)")  "Probe at (0,0):"
    p = [0._default, 0._default]
    call mapping%check (u, p, 1-p, "F7.5")

    write (u, *)
    write (u, "(A)")  "Probe at (0.5,0.5):"
    p = [0.5_default, 0.5_default]
    call mapping%check (u, p, 1-p, "F7.5")

    write (u, *)
    write (u, "(A)")  "Probe at (0.1,0.5):"
    p = [0.1_default, 0.5_default]
    call mapping%check (u, p, 1-p, "F7.5")

    write (u, *)
    write (u, "(A)")  "Probe at (0.1,0.1):"
    p = [0.1_default, 0.1_default]
    call mapping%check (u, p, 1-p, "F7.5")

    write (u, *)
    write (u, "(A)")  "Compute integral:"
    write (u, "(3x,A,1x,F7.5)") "I =", mapping%integral (100000)

    deallocate (mapping)

    write (u, "(A)")
    write (u, "(A)")  "* Test output end: sf_mappings_3"

  end subroutine sf_mappings_3
  
\eatline
\nwendcode{}\nwbegindocs{569}\nwdocspar
\subsubsection{Check on-shell mapping}
Probe the on-shell mapping of the unit square for different parameter
values.  Also calculates integrals.  In this case, the Jacobian is
constant and given by $|\log m^2|$, so this is also the value of the
integral.  The factor results from the variable change in the $\delta$
function $\delta (m^2 - x_1x_2)$ which multiplies the cross section
for the case at hand.

For the test, the (rescaled) resonance mass is set at $1/2$ the
energy.
\nwenddocs{}\nwbegincode{570}\moddef{SF mappings: execute tests}\plusendmoddef
  call test (sf_mappings_4, "sf_mappings_4", &
       "on-shell pair mapping", &
       u, results)
\nwendcode{}\nwbegincode{571}\moddef{SF mappings: tests}\plusendmoddef
  subroutine sf_mappings_4 (u)
    integer, intent(in) :: u
    class(sf_mapping_t), allocatable :: mapping
    real(default), dimension(2) :: p
    
    write (u, "(A)")  "* Test output: sf_mappings_4"
    write (u, "(A)")  "*   Purpose: probe on-shell pair mapping"
    write (u, "(A)")
    
    allocate (sf_os_mapping_t :: mapping)
    select type (mapping)
    type is (sf_os_mapping_t)
       call mapping%init (0.5_default)
       call mapping%set_index (1, 1)
       call mapping%set_index (2, 2)
    end select
       
    call mapping%write (u)
    
    write (u, *)
    write (u, "(A)")  "Probe at (0,0):"
    p = [0._default, 0._default]
    call mapping%check (u, p, 1-p, "F7.5")

    write (u, *)
    write (u, "(A)")  "Probe at (0.5,0.5):"
    p = [0.5_default, 0.5_default]
    call mapping%check (u, p, 1-p, "F7.5")

    write (u, *)
    write (u, "(A)")  "Probe at (0,0.1):"
    p = [0._default, 0.1_default]
    call mapping%check (u, p, 1-p, "F7.5")

    write (u, *)
    write (u, "(A)")  "Probe at (0,1.0):"
    p = [0._default, 1.0_default]
    call mapping%check (u, p, 1-p, "F7.5")

    write (u, *)
    write (u, "(A)")  "Compute integral:"
    write (u, "(3x,A,1x,F7.5)") "I =", mapping%integral (100000)

    deallocate (mapping)

    write (u, "(A)")
    write (u, "(A)")  "* Test output end: sf_mappings_4"

  end subroutine sf_mappings_4
  
\eatline
\nwendcode{}\nwbegindocs{572}\nwdocspar
\subsubsection{Check endpoint mapping}
Probe the endpoint mapping of the unit square for different parameter
values.  Also calculates integrals.  For a finite number of bins, they differ
slightly from $1$, but the result is well-defined because we are not using
random points.
\nwenddocs{}\nwbegincode{573}\moddef{SF mappings: execute tests}\plusendmoddef
  call test (sf_mappings_5, "sf_mappings_5", &
       "endpoint pair mapping", &
       u, results)
\nwendcode{}\nwbegincode{574}\moddef{SF mappings: tests}\plusendmoddef
  subroutine sf_mappings_5 (u)
    integer, intent(in) :: u
    class(sf_mapping_t), allocatable :: mapping
    real(default), dimension(2) :: p
    
    write (u, "(A)")  "* Test output: sf_mappings_5"
    write (u, "(A)")  "*   Purpose: probe endpoint pair mapping"
    write (u, "(A)")
    
    allocate (sf_ep_mapping_t :: mapping)
    select type (mapping)
    type is (sf_ep_mapping_t)
       call mapping%init ()
       call mapping%set_index (1, 1)
       call mapping%set_index (2, 2)
    end select
       
    call mapping%write (u)
    
    write (u, *)
    write (u, "(A)")  "Probe at (0,0):"
    p = [0._default, 0._default]
    call mapping%check (u, p, 1-p, "F7.5")

    write (u, *)
    write (u, "(A)")  "Probe at (0.5,0.5):"
    p = [0.5_default, 0.5_default]
    call mapping%check (u, p, 1-p, "F7.5")

    write (u, *)
    write (u, "(A)")  "Probe at (0.1,0.5):"
    p = [0.1_default, 0.5_default]
    call mapping%check (u, p, 1-p, "F7.5")

    write (u, *)
    write (u, "(A)")  "Probe at (0.7,0.2):"
    p = [0.7_default, 0.2_default]
    call mapping%check (u, p, 1-p, "F7.5")

    write (u, *)
    write (u, "(A)")  "Compute integral:"
    write (u, "(3x,A,1x,F7.5)") "I =", mapping%integral (100000)

    deallocate (mapping)

    write (u, "(A)")
    write (u, "(A)")  "* Test output end: sf_mappings_5"

  end subroutine sf_mappings_5
  
\eatline
\nwendcode{}\nwbegindocs{575}\nwdocspar
\subsubsection{Check endpoint resonant mapping}
Probe the endpoint mapping with resonance.  Also calculates integrals.
\nwenddocs{}\nwbegincode{576}\moddef{SF mappings: execute tests}\plusendmoddef
  call test (sf_mappings_6, "sf_mappings_6", &
       "endpoint resonant mapping", &
       u, results)
\nwendcode{}\nwbegincode{577}\moddef{SF mappings: tests}\plusendmoddef
  subroutine sf_mappings_6 (u)
    integer, intent(in) :: u
    class(sf_mapping_t), allocatable :: mapping
    real(default), dimension(2) :: p
    
    write (u, "(A)")  "* Test output: sf_mappings_6"
    write (u, "(A)")  "*   Purpose: probe endpoint resonant mapping"
    write (u, "(A)")
    
    allocate (sf_epr_mapping_t :: mapping)
    select type (mapping)
    type is (sf_epr_mapping_t)
       call mapping%init (a = 1._default, m = 0.5_default, w = 0.1_default)
       call mapping%set_index (1, 1)
       call mapping%set_index (2, 2)
    end select
       
    call mapping%write (u)
    
    write (u, *)
    write (u, "(A)")  "Probe at (0,0):"
    p = [0._default, 0._default]
    call mapping%check (u, p, 1-p, "F7.5")

    write (u, *)
    write (u, "(A)")  "Probe at (0.5,0.5):"
    p = [0.5_default, 0.5_default]
    call mapping%check (u, p, 1-p, "F7.5")

    write (u, *)
    write (u, "(A)")  "Probe at (0.1,0.5):"
    p = [0.1_default, 0.5_default]
    call mapping%check (u, p, 1-p, "F7.5")

    write (u, *)
    write (u, "(A)")  "Probe at (0.7,0.2):"
    p = [0.7_default, 0.2_default]
    call mapping%check (u, p, 1-p, "F7.5")

    write (u, *)
    write (u, "(A)")  "Compute integral:"
    write (u, "(3x,A,1x,F7.5)") "I =", mapping%integral (100000)

    deallocate (mapping)

    write (u, "(A)")
    write (u, "(A)")  "* Same mapping without resonance:"
    write (u, "(A)")

    allocate (sf_epr_mapping_t :: mapping)
    select type (mapping)
    type is (sf_epr_mapping_t)
       call mapping%init (a = 1._default)
       call mapping%set_index (1, 1)
       call mapping%set_index (2, 2)
    end select
       
    call mapping%write (u)
    
    write (u, *)
    write (u, "(A)")  "Probe at (0,0):"
    p = [0._default, 0._default]
    call mapping%check (u, p, 1-p, "F7.5")

    write (u, *)
    write (u, "(A)")  "Probe at (0.5,0.5):"
    p = [0.5_default, 0.5_default]
    call mapping%check (u, p, 1-p, "F7.5")

    write (u, *)
    write (u, "(A)")  "Probe at (0.1,0.5):"
    p = [0.1_default, 0.5_default]
    call mapping%check (u, p, 1-p, "F7.5")

    write (u, *)
    write (u, "(A)")  "Probe at (0.7,0.2):"
    p = [0.7_default, 0.2_default]
    call mapping%check (u, p, 1-p, "F7.5")

    write (u, *)
    write (u, "(A)")  "Compute integral:"
    write (u, "(3x,A,1x,F7.5)") "I =", mapping%integral (100000)

    deallocate (mapping)

    write (u, "(A)")
    write (u, "(A)")  "* Test output end: sf_mappings_6"

  end subroutine sf_mappings_6
  
\eatline
\nwendcode{}\nwbegindocs{578}\nwdocspar
\subsubsection{Check endpoint on-shell mapping}
Probe the endpoint mapping with an on-shell particle.  Also calculates
integrals.
\nwenddocs{}\nwbegincode{579}\moddef{SF mappings: execute tests}\plusendmoddef
  call test (sf_mappings_7, "sf_mappings_7", &
       "endpoint on-shell mapping", &
       u, results)
\nwendcode{}\nwbegincode{580}\moddef{SF mappings: tests}\plusendmoddef
  subroutine sf_mappings_7 (u)
    integer, intent(in) :: u
    class(sf_mapping_t), allocatable :: mapping
    real(default), dimension(2) :: p
    
    write (u, "(A)")  "* Test output: sf_mappings_7"
    write (u, "(A)")  "*   Purpose: probe endpoint on-shell mapping"
    write (u, "(A)")
    
    allocate (sf_epo_mapping_t :: mapping)
    select type (mapping)
    type is (sf_epo_mapping_t)
       call mapping%init (a = 1._default, m = 0.5_default)
       call mapping%set_index (1, 1)
       call mapping%set_index (2, 2)
    end select
       
    call mapping%write (u)
    
    write (u, *)
    write (u, "(A)")  "Probe at (0,0):"
    p = [0._default, 0._default]
    call mapping%check (u, p, 1-p, "F7.5")

    write (u, *)
    write (u, "(A)")  "Probe at (0.5,0.5):"
    p = [0.5_default, 0.5_default]
    call mapping%check (u, p, 1-p, "F7.5")

    write (u, *)
    write (u, "(A)")  "Probe at (0.1,0.5):"
    p = [0.1_default, 0.5_default]
    call mapping%check (u, p, 1-p, "F7.5")

    write (u, *)
    write (u, "(A)")  "Probe at (0.7,0.2):"
    p = [0.7_default, 0.2_default]
    call mapping%check (u, p, 1-p, "F7.5")

    write (u, *)
    write (u, "(A)")  "Compute integral:"
    write (u, "(3x,A,1x,F7.5)") "I =", mapping%integral (100000)

    deallocate (mapping)

    write (u, "(A)")
    write (u, "(A)")  "* Test output end: sf_mappings_7"

  end subroutine sf_mappings_7
  
\eatline
\nwendcode{}\nwbegindocs{581}\nwdocspar
\subsubsection{Check power mapping}
Probe the power mapping of the unit square for different parameter
values.  Also calculates integrals.  For a finite number of bins, they differ
slightly from $1$, but the result is well-defined because we are not using
random points.
\nwenddocs{}\nwbegincode{582}\moddef{SF mappings: execute tests}\plusendmoddef
  call test (sf_mappings_8, "sf_mappings_8", &
       "power pair mapping", &
       u, results)
\nwendcode{}\nwbegincode{583}\moddef{SF mappings: tests}\plusendmoddef
  subroutine sf_mappings_8 (u)
    integer, intent(in) :: u
    class(sf_mapping_t), allocatable :: mapping
    real(default), dimension(2) :: p, pb
    
    write (u, "(A)")  "* Test output: sf_mappings_8"
    write (u, "(A)")  "*   Purpose: probe power pair mapping"
    write (u, "(A)")
    
    allocate (sf_ip_mapping_t :: mapping)
    select type (mapping)
    type is (sf_ip_mapping_t)
       call mapping%init (eps = 0.1_default)
       call mapping%set_index (1, 1)
       call mapping%set_index (2, 2)
    end select
       
    call mapping%write (u)
    
    write (u, *)
    write (u, "(A)")  "Probe at (0,0.5):"
    p = [0._default, 0.5_default]
    pb= [1._default, 0.5_default]
    call mapping%check (u, p, pb, FMT_16)

    write (u, *)
    write (u, "(A)")  "Probe at (0.5,0.5):"
    p = [0.5_default, 0.5_default]
    pb= [0.5_default, 0.5_default]
    call mapping%check (u, p, pb, FMT_16)

    write (u, *)
    write (u, "(A)")  "Probe at (0.9,0.5):"
    p = [0.9_default, 0.5_default]
    pb= [0.1_default, 0.5_default]
    call mapping%check (u, p, pb, FMT_16)

    write (u, *)
    write (u, "(A)")  "Probe at (0.7,0.2):"
    p = [0.7_default, 0.2_default]
    pb= [0.3_default, 0.8_default]
    call mapping%check (u, p, pb, FMT_16)

    write (u, *)
    write (u, "(A)")  "Probe at (0.7,0.8):"
    p = [0.7_default, 0.8_default]
    pb= [0.3_default, 0.2_default]
    call mapping%check (u, p, pb, FMT_16)

    write (u, *)
    write (u, "(A)")  "Probe at (0.99,0.02):"
    p = [0.99_default, 0.02_default]
    pb= [0.01_default, 0.98_default]
    call mapping%check (u, p, pb, FMT_14, FMT_12)

    write (u, *)
    write (u, "(A)")  "Probe at (0.99,0.98):"
    p = [0.99_default, 0.98_default]
    pb= [0.01_default, 0.02_default]
    call mapping%check (u, p, pb, FMT_14, FMT_12)

    write (u, *)
    write (u, "(A)")  "Compute integral:"
    write (u, "(3x,A,1x,F7.5)") "I =", mapping%integral (100000)

    deallocate (mapping)

    write (u, "(A)")
    write (u, "(A)")  "* Test output end: sf_mappings_8"

  end subroutine sf_mappings_8
  
\eatline
\nwendcode{}\nwbegindocs{584}\nwdocspar
\subsubsection{Check resonant power mapping}
Probe the power mapping of the unit square, adapted for an s-channel
resonance, for different parameter values.  Also calculates integrals.
For a finite number of bins, they differ slightly from $1$, but the
result is well-defined because we are not using random points.
\nwenddocs{}\nwbegincode{585}\moddef{SF mappings: execute tests}\plusendmoddef
  call test (sf_mappings_9, "sf_mappings_9", &
       "power resonance mapping", &
       u, results)
\nwendcode{}\nwbegincode{586}\moddef{SF mappings: tests}\plusendmoddef
  subroutine sf_mappings_9 (u)
    integer, intent(in) :: u
    class(sf_mapping_t), allocatable :: mapping
    real(default), dimension(2) :: p, pb
    
    write (u, "(A)")  "* Test output: sf_mappings_9"
    write (u, "(A)")  "*   Purpose: probe power resonant pair mapping"
    write (u, "(A)")
    
    allocate (sf_ipr_mapping_t :: mapping)
    select type (mapping)
    type is (sf_ipr_mapping_t)
       call mapping%init (eps = 0.1_default, m = 0.5_default, w = 0.1_default)
       call mapping%set_index (1, 1)
       call mapping%set_index (2, 2)
    end select
       
    call mapping%write (u)
    
    write (u, *)
    write (u, "(A)")  "Probe at (0,0.5):"
    p = [0._default, 0.5_default]
    pb= [1._default, 0.5_default]
    call mapping%check (u, p, pb, FMT_16)

    write (u, *)
    write (u, "(A)")  "Probe at (0.5,0.5):"
    p = [0.5_default, 0.5_default]
    pb= [0.5_default, 0.5_default]
    call mapping%check (u, p, pb, FMT_16)

    write (u, *)
    write (u, "(A)")  "Probe at (0.9,0.5):"
    p = [0.9_default, 0.5_default]
    pb= [0.1_default, 0.5_default]
    call mapping%check (u, p, pb, FMT_16)

    write (u, *)
    write (u, "(A)")  "Probe at (0.7,0.2):"
    p = [0.7_default, 0.2_default]
    pb= [0.3_default, 0.8_default]
    call mapping%check (u, p, pb, FMT_16)

    write (u, *)
    write (u, "(A)")  "Probe at (0.7,0.8):"
    p = [0.7_default, 0.8_default]
    pb= [0.3_default, 0.2_default]
    call mapping%check (u, p, pb, FMT_16)

    write (u, *)
    write (u, "(A)")  "Probe at (0.9999,0.02):"
    p = [0.9999_default, 0.02_default]
    pb= [0.0001_default, 0.98_default]
    call mapping%check (u, p, pb, FMT_15, FMT_12)

    write (u, *)
    write (u, "(A)")  "Probe at (0.9999,0.98):"
    p = [0.9999_default, 0.98_default]
    pb= [0.0001_default, 0.02_default]
    call mapping%check (u, p, pb, FMT_15, FMT_12)

    write (u, *)
    write (u, "(A)")  "Compute integral:"
    write (u, "(3x,A,1x,F7.5)") "I =", mapping%integral (100000)

    deallocate (mapping)

    write (u, "(A)")
    write (u, "(A)")  "* Same mapping without resonance:"
    write (u, "(A)")
    
    allocate (sf_ipr_mapping_t :: mapping)
    select type (mapping)
    type is (sf_ipr_mapping_t)
       call mapping%init (eps = 0.1_default)
       call mapping%set_index (1, 1)
       call mapping%set_index (2, 2)
    end select
       
    call mapping%write (u)
    
    write (u, *)
    write (u, "(A)")  "Probe at (0,0.5):"
    p = [0._default, 0.5_default]
    pb= [1._default, 0.5_default]
    call mapping%check (u, p, pb, FMT_16)

    write (u, *)
    write (u, "(A)")  "Probe at (0.5,0.5):"
    p = [0.5_default, 0.5_default]
    pb= [0.5_default, 0.5_default]
    call mapping%check (u, p, pb, FMT_16)

    write (u, *)
    write (u, "(A)")  "Probe at (0.9,0.5):"
    p = [0.9_default, 0.5_default]
    pb= [0.1_default, 0.5_default]
    call mapping%check (u, p, pb, FMT_16)

    write (u, *)
    write (u, "(A)")  "Probe at (0.7,0.2):"
    p = [0.7_default, 0.2_default]
    pb= [0.3_default, 0.8_default]
    call mapping%check (u, p, pb, FMT_16)

    write (u, *)
    write (u, "(A)")  "Probe at (0.7,0.8):"
    p = [0.7_default, 0.8_default]
    pb= [0.3_default, 0.2_default]
    call mapping%check (u, p, pb, FMT_16)

    write (u, *)
    write (u, "(A)")  "Compute integral:"
    write (u, "(3x,A,1x,F7.5)") "I =", mapping%integral (100000)

    deallocate (mapping)

    write (u, "(A)")
    write (u, "(A)")  "* Test output end: sf_mappings_9"

  end subroutine sf_mappings_9
  
\eatline
\nwendcode{}\nwbegindocs{587}\nwdocspar
\subsubsection{Check on-shell power mapping}
Probe the power mapping of the unit square, adapted for
single-particle production, for different parameter values.  Also
calculates integrals.  For a finite number of bins, they differ
slightly from $1$, but the result is well-defined because we are not
using random points.
\nwenddocs{}\nwbegincode{588}\moddef{SF mappings: execute tests}\plusendmoddef
  call test (sf_mappings_10, "sf_mappings_10", &
       "power on-shell mapping", &
       u, results)
\nwendcode{}\nwbegincode{589}\moddef{SF mappings: tests}\plusendmoddef
  subroutine sf_mappings_10 (u)
    integer, intent(in) :: u
    class(sf_mapping_t), allocatable :: mapping
    real(default), dimension(2) :: p, pb
    
    write (u, "(A)")  "* Test output: sf_mappings_10"
    write (u, "(A)")  "*   Purpose: probe power on-shell mapping"
    write (u, "(A)")
    
    allocate (sf_ipo_mapping_t :: mapping)
    select type (mapping)
    type is (sf_ipo_mapping_t)
       call mapping%init (eps = 0.1_default, m = 0.5_default)
       call mapping%set_index (1, 1)
       call mapping%set_index (2, 2)
    end select
       
    call mapping%write (u)
    
    write (u, *)
    write (u, "(A)")  "Probe at (0,0.5):"
    p = [0._default, 0.5_default]
    pb= [1._default, 0.5_default]
    call mapping%check (u, p, pb, FMT_16)

    write (u, *)
    write (u, "(A)")  "Probe at (0,0.02):"
    p = [0._default, 0.02_default]
    pb= [1._default, 0.98_default]
    call mapping%check (u, p, pb, FMT_15, FMT_12)

    write (u, *)
    write (u, "(A)")  "Probe at (0,0.98):"
    p = [0._default, 0.98_default]
    pb= [1._default, 0.02_default]
    call mapping%check (u, p, pb, FMT_15, FMT_12)

    write (u, *)
    write (u, "(A)")  "Compute integral:"
    write (u, "(3x,A,1x,F7.5)") "I =", mapping%integral (100000)

    deallocate (mapping)

    write (u, "(A)")
    write (u, "(A)")  "* Test output end: sf_mappings_10"

  end subroutine sf_mappings_10
  
\eatline
\nwendcode{}\nwbegindocs{590}\nwdocspar
\subsubsection{Check combined endpoint-power mapping}
Probe the mapping for the beamstrahlung/ISR combination.
\nwenddocs{}\nwbegincode{591}\moddef{SF mappings: execute tests}\plusendmoddef
  call test (sf_mappings_11, "sf_mappings_11", &
       "endpoint/power combined mapping", &
       u, results)
\nwendcode{}\nwbegincode{592}\moddef{SF mappings: tests}\plusendmoddef
  subroutine sf_mappings_11 (u)
    integer, intent(in) :: u
    class(sf_mapping_t), allocatable :: mapping
    real(default), dimension(4) :: p, pb
    
    write (u, "(A)")  "* Test output: sf_mappings_11"
    write (u, "(A)")  "*   Purpose: probe power pair mapping"
    write (u, "(A)")
    
    allocate (sf_ei_mapping_t :: mapping)
    select type (mapping)
    type is (sf_ei_mapping_t)
       call mapping%init (eps = 0.1_default)
       call mapping%set_index (1, 1)
       call mapping%set_index (2, 2)
       call mapping%set_index (3, 3)
       call mapping%set_index (4, 4)
    end select
       
    call mapping%write (u)
    
    write (u, *)
    write (u, "(A)")  "Probe at (0.5, 0.5, 0.5, 0.5):"
    p = [0.5_default, 0.5_default, 0.5_default, 0.5_default]
    pb= [0.5_default, 0.5_default, 0.5_default, 0.5_default]
    call mapping%check (u, p, pb, FMT_16)

    write (u, *)
    write (u, "(A)")  "Probe at (0.7, 0.2, 0.4, 0.8):"
    p = [0.7_default, 0.2_default, 0.4_default, 0.8_default]
    pb= [0.3_default, 0.8_default, 0.6_default, 0.2_default]
    call mapping%check (u, p, pb, FMT_16)

    write (u, *)
    write (u, "(A)")  "Probe at (0.9, 0.06, 0.95, 0.1):"
    p = [0.9_default, 0.06_default, 0.95_default, 0.1_default]
    pb= [0.1_default, 0.94_default, 0.05_default, 0.9_default]
    call mapping%check (u, p, pb, FMT_13, FMT_12)

    write (u, *)
    write (u, "(A)")  "Compute integral:"
    write (u, "(3x,A,1x,F7.5)") "I =", mapping%integral (100000)

    deallocate (mapping)

    write (u, "(A)")
    write (u, "(A)")  "* Test output end: sf_mappings_11"

  end subroutine sf_mappings_11
  
\eatline
\nwendcode{}\nwbegindocs{593}\nwdocspar
\subsubsection{Check resonant endpoint-power mapping}
Probe the mapping for the beamstrahlung/ISR combination.
\nwenddocs{}\nwbegincode{594}\moddef{SF mappings: execute tests}\plusendmoddef
  call test (sf_mappings_12, "sf_mappings_12", &
       "endpoint/power resonant combined mapping", &
       u, results)
\nwendcode{}\nwbegincode{595}\moddef{SF mappings: tests}\plusendmoddef
  subroutine sf_mappings_12 (u)
    integer, intent(in) :: u
    class(sf_mapping_t), allocatable :: mapping
    real(default), dimension(4) :: p, pb
    
    write (u, "(A)")  "* Test output: sf_mappings_12"
    write (u, "(A)")  "*   Purpose: probe resonant combined mapping"
    write (u, "(A)")
    
    allocate (sf_eir_mapping_t :: mapping)
    select type (mapping)
    type is (sf_eir_mapping_t)
       call mapping%init (a = 1._default, &
            eps = 0.1_default, m = 0.5_default, w = 0.1_default)
       call mapping%set_index (1, 1)
       call mapping%set_index (2, 2)
       call mapping%set_index (3, 3)
       call mapping%set_index (4, 4)
    end select
       
    call mapping%write (u)
    
    write (u, *)
    write (u, "(A)")  "Probe at (0.5, 0.5, 0.5, 0.5):"
    p = [0.5_default, 0.5_default, 0.5_default, 0.5_default]
    pb= [0.5_default, 0.5_default, 0.5_default, 0.5_default]
    call mapping%check (u, p, pb, FMT_16)

    write (u, *)
    write (u, "(A)")  "Probe at (0.7, 0.2, 0.4, 0.8):"
    p = [0.7_default, 0.2_default, 0.4_default, 0.8_default]
    pb= [0.3_default, 0.8_default, 0.6_default, 0.2_default]
    call mapping%check (u, p, pb, FMT_16)

    write (u, *)
    write (u, "(A)")  "Probe at (0.9, 0.06, 0.95, 0.1):"
    p = [0.9_default, 0.06_default, 0.95_default, 0.1_default]
    pb= [0.1_default, 0.94_default, 0.05_default, 0.9_default]
    call mapping%check (u, p, pb, FMT_15, FMT_12)

    write (u, *)
    write (u, "(A)")  "Compute integral:"
    write (u, "(3x,A,1x,F7.5)") "I =", mapping%integral (100000)

    deallocate (mapping)

    write (u, "(A)")
    write (u, "(A)")  "* Test output end: sf_mappings_12"

  end subroutine sf_mappings_12
  
\eatline
\nwendcode{}\nwbegindocs{596}\nwdocspar
\subsubsection{Check on-shell endpoint-power mapping}
Probe the mapping for the beamstrahlung/ISR combination.
\nwenddocs{}\nwbegincode{597}\moddef{SF mappings: execute tests}\plusendmoddef
  call test (sf_mappings_13, "sf_mappings_13", &
       "endpoint/power on-shell combined mapping", &
       u, results)
\nwendcode{}\nwbegincode{598}\moddef{SF mappings: tests}\plusendmoddef
  subroutine sf_mappings_13 (u)
    integer, intent(in) :: u
    class(sf_mapping_t), allocatable :: mapping
    real(default), dimension(4) :: p, pb
    
    write (u, "(A)")  "* Test output: sf_mappings_13"
    write (u, "(A)")  "*   Purpose: probe on-shell combined mapping"
    write (u, "(A)")
    
    allocate (sf_eio_mapping_t :: mapping)
    select type (mapping)
    type is (sf_eio_mapping_t)
       call mapping%init (a = 1._default, eps = 0.1_default, m = 0.5_default)
       call mapping%set_index (1, 1)
       call mapping%set_index (2, 2)
       call mapping%set_index (3, 3)
       call mapping%set_index (4, 4)
    end select
       
    call mapping%write (u)
    
    write (u, *)
    write (u, "(A)")  "Probe at (0.5, 0.5, 0.5, 0.5):"
    p = [0.5_default, 0.5_default, 0.5_default, 0.5_default]
    pb= [0.5_default, 0.5_default, 0.5_default, 0.5_default]
    call mapping%check (u, p, pb, FMT_16)

    write (u, *)
    write (u, "(A)")  "Probe at (0.7, 0.2, 0.4, 0.8):"
    p = [0.7_default, 0.2_default, 0.4_default, 0.8_default]
    pb= [0.3_default, 0.8_default, 0.6_default, 0.2_default]
    call mapping%check (u, p, pb, FMT_16)

    write (u, *)
    write (u, "(A)")  "Probe at (0.9, 0.06, 0.95, 0.1):"
    p = [0.9_default, 0.06_default, 0.95_default, 0.1_default]
    pb= [0.1_default, 0.94_default, 0.05_default, 0.9_default]
    call mapping%check (u, p, pb, FMT_14, FMT_12)

    write (u, *)
    write (u, "(A)")  "Compute integral:"
    write (u, "(3x,A,1x,F7.5)") "I =", mapping%integral (100000)

    deallocate (mapping)

    write (u, "(A)")
    write (u, "(A)")  "* Test output end: sf_mappings_13"

  end subroutine sf_mappings_13
  
\eatline
\nwendcode{}\nwbegindocs{599}\nwdocspar
\subsubsection{Check rescaling}
Check the rescaling factor in on-shell basic mapping.
\nwenddocs{}\nwbegincode{600}\moddef{SF mappings: execute tests}\plusendmoddef
  call test (sf_mappings_14, "sf_mappings_14", &
       "rescaled on-shell mapping", &
       u, results)
\nwendcode{}\nwbegincode{601}\moddef{SF mappings: tests}\plusendmoddef
  subroutine sf_mappings_14 (u)
    integer, intent(in) :: u
    real(default), dimension(2) :: p, r
    real(default) :: f, x_free, m2
    
    write (u, "(A)")  "* Test output: sf_mappings_14"
    write (u, "(A)")  "*   Purpose: probe rescaling in os mapping"
    write (u, "(A)")
    
    p = [0.1_default, 0.2_default]
    x_free = 0.9_default
    m2 = 0.5_default
    
    call map_on_shell (r, f, p, -log (m2), x_free)
    
    write (u, "(A,9(1x," // FMT_14 // "))")  "p =", p
    write (u, "(A,9(1x," // FMT_14 // "))")  "r =", r
    write (u, "(A,9(1x," // FMT_14 // "))")  "f =", f
    write (u, "(A,9(1x," // FMT_14 // "))")  "*r=", x_free * product (r)

    write (u, *)

    call map_on_shell_inverse (r, f, p, -log (m2), x_free)
    
    write (u, "(A,9(1x," // FMT_14 // "))")  "p =", p
    write (u, "(A,9(1x," // FMT_14 // "))")  "r =", r
    write (u, "(A,9(1x," // FMT_14 // "))")  "f =", f
    write (u, "(A,9(1x," // FMT_14 // "))")  "*r=", x_free * product (r)

    write (u, "(A)")
    write (u, "(A)")  "* Test output end: sf_mappings_14"

  end subroutine sf_mappings_14
  
\eatline
\nwendcode{}\nwbegindocs{602}\nwdocspar
%%%%%%%%%%%%%%%%%%%%%%%%%%%%%%%%%%%%%%%%%%%%%%%%%%%%%%%%%%%%%%%%%%%%%%%%
\section{Structure function base}

\nwenddocs{}\nwbegincode{603}\moddef{\code{}sf{\char95}base.f90\edoc{}}\endmoddef
\LA{}File header\RA{}

module sf_base

\LA{}Use kinds\RA{}
\LA{}Use strings\RA{}
  use io_units
  use format_utils, only: write_separator
  use format_defs, only: FMT_17, FMT_19
  use unit_tests
  use diagnostics
  use lorentz
  use model_data
  use flavors
  use colors
  use helicities
  use quantum_numbers
  use state_matrices
  use interactions
  use evaluators
  use pdg_arrays
  use particles
  use beams
  use sf_aux
  use sf_mappings

\LA{}Standard module head\RA{}

\LA{}SF base: public\RA{}

\LA{}SF base: parameters\RA{}

\LA{}SF base: types\RA{}

\LA{}SF base: interfaces\RA{}

\LA{}SF base: test types\RA{}

contains

\LA{}SF base: procedures\RA{}

\LA{}SF base: tests\RA{}

end module sf_base
\eatline
\nwendcode{}\nwbegindocs{604}\nwdocspar
\subsection{Abstract structure-function data type}
This type should hold all configuration data for a specific type of
structure function.  The base object is empty; the implementations
will fill it.
\nwenddocs{}\nwbegincode{605}\moddef{SF base: public}\endmoddef
  public :: sf_data_t
\nwendcode{}\nwbegincode{606}\moddef{SF base: types}\endmoddef
  type, abstract :: sf_data_t
   contains
   \LA{}SF base: sf data: TBP\RA{}
  end type sf_data_t

\eatline
\nwendcode{}\nwbegindocs{607}Output.
\nwenddocs{}\nwbegincode{608}\moddef{SF base: sf data: TBP}\endmoddef
  procedure (sf_data_write), deferred :: write
\nwendcode{}\nwbegincode{609}\moddef{SF base: interfaces}\endmoddef
  abstract interface
     subroutine sf_data_write (data, unit, verbose)
       import
       class(sf_data_t), intent(in) :: data
       integer, intent(in), optional :: unit
       logical, intent(in), optional :: verbose
     end subroutine sf_data_write
  end interface
  
\eatline
\nwendcode{}\nwbegindocs{610}Return true if this structure function is in generator mode.  In
that case, all parameters are free, otherwise bound.  (We do not
support mixed cases.)  Default is: no generator.
\nwenddocs{}\nwbegincode{611}\moddef{SF base: sf data: TBP}\plusendmoddef
  procedure :: is_generator => sf_data_is_generator
\nwendcode{}\nwbegincode{612}\moddef{SF base: procedures}\endmoddef
  function sf_data_is_generator (data) result (flag)
    class(sf_data_t), intent(in) :: data
    logical :: flag
    flag = .false.
  end function sf_data_is_generator
  
\eatline
\nwendcode{}\nwbegindocs{613}Return the number of input parameters that determine the
structure function.
\nwenddocs{}\nwbegincode{614}\moddef{SF base: sf data: TBP}\plusendmoddef
  procedure (sf_data_get_int), deferred :: get_n_par
\nwendcode{}\nwbegincode{615}\moddef{SF base: interfaces}\plusendmoddef
  abstract interface
     function sf_data_get_int (data) result (n)
       import
       class(sf_data_t), intent(in) :: data
       integer :: n
     end function sf_data_get_int
  end interface

\eatline
\nwendcode{}\nwbegindocs{616}Return the outgoing particle PDG codes for the current setup.  The codes can
be an array of particles, for each beam.
\nwenddocs{}\nwbegincode{617}\moddef{SF base: sf data: TBP}\plusendmoddef
  procedure (sf_data_get_pdg_out), deferred :: get_pdg_out
\nwendcode{}\nwbegincode{618}\moddef{SF base: interfaces}\plusendmoddef
  abstract interface
     subroutine sf_data_get_pdg_out (data, pdg_out)
       import
       class(sf_data_t), intent(in) :: data
       type(pdg_array_t), dimension(:), intent(inout) :: pdg_out
     end subroutine sf_data_get_pdg_out
  end interface
  
\eatline
\nwendcode{}\nwbegindocs{619}Allocate a matching structure function interaction object and
properly initialize it.

\nwenddocs{}\nwbegincode{620}\moddef{SF base: sf data: TBP}\plusendmoddef
  procedure (sf_data_allocate_sf_int), deferred :: allocate_sf_int
\nwendcode{}\nwbegincode{621}\moddef{SF base: interfaces}\plusendmoddef
  abstract interface
     subroutine sf_data_allocate_sf_int (data, sf_int)
       import
       class(sf_data_t), intent(in) :: data
       class(sf_int_t), intent(inout), allocatable :: sf_int
     end subroutine sf_data_allocate_sf_int
  end interface

\eatline
\nwendcode{}\nwbegindocs{622}Return the PDF set index, if applicable.  We implement a default
method which returns zero.  The PDF (builtin and LHA) implementations
will override this.
\nwenddocs{}\nwbegincode{623}\moddef{SF base: sf data: TBP}\plusendmoddef
  procedure :: get_pdf_set => sf_data_get_pdf_set
\nwendcode{}\nwbegincode{624}\moddef{SF base: procedures}\plusendmoddef
  function sf_data_get_pdf_set (data) result (pdf_set)
    class(sf_data_t), intent(in) :: data
    integer :: pdf_set
    pdf_set = 0
  end function sf_data_get_pdf_set
  
\eatline
\nwendcode{}\nwbegindocs{625}\nwdocspar
\subsection{Structure-function chain configuration}
This is the data type that the {\tt{}processes} module uses for setting
up its structure-function chain.  For each structure function described
by the beam data, there is an entry.  The {\tt{}i} array indicates the
beam(s) to which this structure function applies, and the {\tt{}data}
object contains the actual configuration data.
\nwenddocs{}\nwbegincode{626}\moddef{SF base: public}\plusendmoddef
  public :: sf_config_t
\nwendcode{}\nwbegincode{627}\moddef{SF base: types}\plusendmoddef
  type :: sf_config_t
     integer, dimension(:), allocatable :: i
     class(sf_data_t), allocatable :: data
   contains
   \LA{}SF base: sf config: TBP\RA{}
  end type sf_config_t
  
\eatline
\nwendcode{}\nwbegindocs{628}Output:
\nwenddocs{}\nwbegincode{629}\moddef{SF base: sf config: TBP}\endmoddef
  procedure :: write => sf_config_write
\nwendcode{}\nwbegincode{630}\moddef{SF base: procedures}\plusendmoddef
  subroutine sf_config_write (object, unit)
    class(sf_config_t), intent(in) :: object
    integer, intent(in), optional :: unit
    integer :: u
    u = given_output_unit (unit)
    if (allocated (object%i)) then
       write (u, "(1x,A,2(1x,I0))")  "Structure-function configuration: &
            &beam(s)", object%i
       if (allocated (object%data))  call object%data%write (u)
    else
       write (u, "(1x,A)")  "Structure-function configuration: [undefined]"
    end if
  end subroutine sf_config_write
       
\eatline
\nwendcode{}\nwbegindocs{631}Initialize.
\nwenddocs{}\nwbegincode{632}\moddef{SF base: sf config: TBP}\plusendmoddef
  procedure :: init => sf_config_init
\nwendcode{}\nwbegincode{633}\moddef{SF base: procedures}\plusendmoddef
  subroutine sf_config_init (sf_config, i_beam, sf_data)
    class(sf_config_t), intent(out) :: sf_config
    integer, dimension(:), intent(in) :: i_beam
    class(sf_data_t), intent(in) :: sf_data
    allocate (sf_config%i (size (i_beam)), source = i_beam)
    allocate (sf_config%data, source = sf_data)
  end subroutine sf_config_init
  
\eatline
\nwendcode{}\nwbegindocs{634}Return the PDF set, if any.
\nwenddocs{}\nwbegincode{635}\moddef{SF base: sf config: TBP}\plusendmoddef
  procedure :: get_pdf_set => sf_config_get_pdf_set
\nwendcode{}\nwbegincode{636}\moddef{SF base: procedures}\plusendmoddef
  function sf_config_get_pdf_set (sf_config) result (pdf_set)
    class(sf_config_t), intent(in) :: sf_config
    integer :: pdf_set
    pdf_set = sf_config%data%get_pdf_set ()
  end function sf_config_get_pdf_set
  
\eatline
\nwendcode{}\nwbegindocs{637}\nwdocspar
\subsection{Structure-function instance}
The {\tt{}sf{\char95}int{\char95}t} data type contains an {\tt{}interaction{\char95}t} object (it is
an extension of this type) and a pointer to the {\tt{}sf{\char95}data{\char95}t}
configuration data.  This interaction, or copies of it, is used to
implement structure-function kinematics and dynamics in the context of
process evaluation.

The status code {\tt{}status} tells whether the interaction is undefined,
has defined kinematics (but matrix elements invalid), or is completely
defined.  There is also a status code for failure.  The implementation
is responsible for updating the status.

The entries {\tt{}mi2}, {\tt{}mr2}, and {\tt{}mo2} hold the squared
invariant masses of the incoming, radiated, and outgoing particle,
respectively.  They are supposed to be set upon initialization, but
could also be varied event by event.

If the radiated or outgoing mass is nonzero, we may need to apply an
on-shell projection.  The projection mode is stored as
{\tt{}on{\char95}shell{\char95}mode}.

The array {\tt{}beam{\char95}index} is the list of beams on which this structure
function applies ($1$, $2$, or both).  The arrays {\tt{}incoming},
{\tt{}radiated}, and {\tt{}outgoing} contain the indices of the respective
particle sets within the interaction, for convenient lookup.  The
array {\tt{}par{\char95}index} indicates the MC input parameters that this entry
will use up in the structure-function chain.  The first parameter (or
the first two, for a spectrum) in this array determines the momentum
fraction and is thus subject to global mappings.

In the abstract base type, we do not implement the data pointer.  This
allows us to restrict its type in the implementations.
\nwenddocs{}\nwbegincode{638}\moddef{SF base: public}\plusendmoddef
  public :: sf_int_t
\nwendcode{}\nwbegincode{639}\moddef{SF base: types}\plusendmoddef
  type, abstract, extends (interaction_t) :: sf_int_t
     integer :: status = SF_UNDEFINED
     real(default), dimension(:), allocatable :: mi2
     real(default), dimension(:), allocatable :: mr2
     real(default), dimension(:), allocatable :: mo2
     integer :: on_shell_mode = KEEP_ENERGY
     logical :: qmin_defined = .false.
     logical :: qmax_defined = .false.
     real(default), dimension(:), allocatable :: qmin
     real(default), dimension(:), allocatable :: qmax
     integer, dimension(:), allocatable :: beam_index
     integer, dimension(:), allocatable :: incoming
     integer, dimension(:), allocatable :: radiated
     integer, dimension(:), allocatable :: outgoing
     integer, dimension(:), allocatable :: par_index
     integer, dimension(:), allocatable :: par_primary
   contains
   \LA{}SF base: sf int: TBP\RA{}
  end type sf_int_t
  
\eatline
\nwendcode{}\nwbegindocs{640}Status codes.  The codes that refer to links, masks, and
connections, apply to structure-function chains only.

The status codes are public.
\nwenddocs{}\nwbegincode{641}\moddef{SF base: parameters}\endmoddef
  integer, parameter, public :: SF_UNDEFINED = 0
  integer, parameter, public :: SF_INITIAL = 1
  integer, parameter, public :: SF_DONE_LINKS = 2
  integer, parameter, public :: SF_FAILED_MASK = 3
  integer, parameter, public :: SF_DONE_MASK = 4
  integer, parameter, public :: SF_FAILED_CONNECTIONS = 5
  integer, parameter, public :: SF_DONE_CONNECTIONS = 6
  integer, parameter, public :: SF_SEED_KINEMATICS = 10
  integer, parameter, public :: SF_FAILED_KINEMATICS = 11
  integer, parameter, public :: SF_DONE_KINEMATICS = 12
  integer, parameter, public :: SF_FAILED_EVALUATION = 13
  integer, parameter, public :: SF_EVALUATED = 20

\eatline
\eatline
\eatline
\eatline
\eatline
\nwendcode{}\nwbegindocs{642}Write a string version of the status code:
\nwenddocs{}\nwbegincode{643}\moddef{SF base: procedures}\plusendmoddef
  subroutine write_sf_status (status, u)
    integer, intent(in) :: status
    integer, intent(in) :: u
    select case (status)
    case (SF_UNDEFINED)
       write (u, "(1x,'[',A,']')")  "undefined"
    case (SF_INITIAL)
       write (u, "(1x,'[',A,']')")  "initialized"
    case (SF_DONE_LINKS)
       write (u, "(1x,'[',A,']')")  "links set"
    case (SF_FAILED_MASK)
       write (u, "(1x,'[',A,']')")  "mask mismatch"
    case (SF_DONE_MASK)
       write (u, "(1x,'[',A,']')")  "mask set"
    case (SF_FAILED_CONNECTIONS)
       write (u, "(1x,'[',A,']')")  "connections failed"
    case (SF_DONE_CONNECTIONS)
       write (u, "(1x,'[',A,']')")  "connections set"
    case (SF_SEED_KINEMATICS)
       write (u, "(1x,'[',A,']')")  "incoming momenta set"
    case (SF_FAILED_KINEMATICS)
       write (u, "(1x,'[',A,']')")  "kinematics failed"
    case (SF_DONE_KINEMATICS)
       write (u, "(1x,'[',A,']')")  "kinematics set"
    case (SF_FAILED_EVALUATION)
       write (u, "(1x,'[',A,']')")  "evaluation failed"
    case (SF_EVALUATED)
       write (u, "(1x,'[',A,']')")  "evaluated"
    end select
  end subroutine write_sf_status

\eatline
\nwendcode{}\nwbegindocs{644}Finalizer.  Required because the base interaction needs a
finalizer.  We provide an implementation, which need not be
overwritten if the implementation does not add finalizable
subobjects.
\nwenddocs{}\nwbegincode{645}\moddef{SF base: sf int: TBP}\endmoddef
  procedure :: final => sf_int_final
\nwendcode{}\nwbegincode{646}\moddef{SF base: procedures}\plusendmoddef
  subroutine sf_int_final (object)
    class(sf_int_t), intent(inout) :: object
    call interaction_final (object%interaction_t)
  end subroutine sf_int_final

\eatline
\nwendcode{}\nwbegindocs{647}This is the basic output routine.  Display status and interaction.
\nwenddocs{}\nwbegincode{648}\moddef{SF base: sf int: TBP}\plusendmoddef
  procedure :: base_write => sf_int_base_write
\nwendcode{}\nwbegincode{649}\moddef{SF base: procedures}\plusendmoddef
  subroutine sf_int_base_write (object, unit, testflag)
    class(sf_int_t), intent(in) :: object
    integer, intent(in), optional :: unit
    logical, intent(in), optional :: testflag
    integer :: u
    u = given_output_unit (unit)
    write (u, "(1x,A)", advance="no")  "SF instance:"
    call write_sf_status (object%status, u)
    if (allocated (object%beam_index)) &
         write (u, "(3x,A,2(1x,I0))")  "beam      =", object%beam_index
    if (allocated (object%incoming)) &
         write (u, "(3x,A,2(1x,I0))")  "incoming  =", object%incoming
    if (allocated (object%radiated)) &
         write (u, "(3x,A,2(1x,I0))")  "radiated  =", object%radiated
    if (allocated (object%outgoing)) &
         write (u, "(3x,A,2(1x,I0))")  "outgoing  =", object%outgoing
    if (allocated (object%par_index)) &
         write (u, "(3x,A,2(1x,I0))")  "parameter =", object%par_index
    if (object%qmin_defined) &
         write (u, "(3x,A,1x," // FMT_19 // ")")  "q_min     =", object%qmin
    if (object%qmax_defined) &
         write (u, "(3x,A,1x," // FMT_19 // ")")  "q_max     =", object%qmax
    call interaction_write (object%interaction_t, u, testflag = testflag)
  end subroutine sf_int_base_write
  
\eatline
\nwendcode{}\nwbegindocs{650}The type string identifies the structure function class, and possibly more
details about the structure function.
\nwenddocs{}\nwbegincode{651}\moddef{SF base: sf int: TBP}\plusendmoddef
  procedure (sf_int_type_string), deferred :: type_string
\nwendcode{}\nwbegincode{652}\moddef{SF base: interfaces}\plusendmoddef
  abstract interface
     function sf_int_type_string (object) result (string)
       import
       class(sf_int_t), intent(in) :: object
       type(string_t) :: string
     end function sf_int_type_string
  end interface
  
\eatline
\nwendcode{}\nwbegindocs{653}Output of the concrete object.  We should not forget to call the
output routine for the base type.
\nwenddocs{}\nwbegincode{654}\moddef{SF base: sf int: TBP}\plusendmoddef
  procedure (sf_int_write), deferred :: write
\nwendcode{}\nwbegincode{655}\moddef{SF base: interfaces}\plusendmoddef
  abstract interface
     subroutine sf_int_write (object, unit, testflag)
       import
       class(sf_int_t), intent(in) :: object
       integer, intent(in), optional :: unit
       logical, intent(in), optional :: testflag
     end subroutine sf_int_write
  end interface
  
\eatline
\nwendcode{}\nwbegindocs{656}Basic initialization: set the invariant masses for the particles and
initialize the interaction.  The caller should then add states to the
interaction and freeze it.

The dimension of the mask should be equal to the sum of the dimensions
of the mass-squared arrays, which determine incoming, radiated, and
outgoing particles, respectively.

Optionally, we can define minimum and maximum values for the momentum
transfer to the outgoing particle(s).  If all masses are zero, this is
actually required for non-collinear splitting.
\nwenddocs{}\nwbegincode{657}\moddef{SF base: sf int: TBP}\plusendmoddef
  procedure :: base_init => sf_int_base_init
\nwendcode{}\nwbegincode{658}\moddef{SF base: procedures}\plusendmoddef
  subroutine sf_int_base_init &
       (sf_int, mask, mi2, mr2, mo2, qmin, qmax, hel_lock)
    class(sf_int_t), intent(out) :: sf_int
    type (quantum_numbers_mask_t), dimension(:), intent(in) :: mask
    real(default), dimension(:), intent(in) :: mi2, mr2, mo2
    real(default), dimension(:), intent(in), optional :: qmin, qmax
    integer, dimension(:), intent(in), optional :: hel_lock
    allocate (sf_int%mi2 (size (mi2)))
    sf_int%mi2 = mi2
    allocate (sf_int%mr2 (size (mr2)))
    sf_int%mr2 = mr2
    allocate (sf_int%mo2 (size (mo2)))
    sf_int%mo2 = mo2
    if (present (qmin)) then
       sf_int%qmin_defined = .true.
       allocate (sf_int%qmin (size (qmin)))
       sf_int%qmin = qmin
    end if
    if (present (qmax)) then
       sf_int%qmax_defined = .true.
       allocate (sf_int%qmax (size (qmax)))
       sf_int%qmax = qmax
    end if
    call interaction_init (sf_int%interaction_t, &
         size (mi2), 0, size (mr2) + size (mo2), &
         mask = mask, hel_lock = hel_lock, set_relations = .true.)
  end subroutine sf_int_base_init
    
\eatline
\nwendcode{}\nwbegindocs{659}Set the indices of the incoming, radiated, and outgoing particles,
respectively.
\nwenddocs{}\nwbegincode{660}\moddef{SF base: sf int: TBP}\plusendmoddef
  procedure :: set_incoming => sf_int_set_incoming
  procedure :: set_radiated => sf_int_set_radiated
  procedure :: set_outgoing => sf_int_set_outgoing
\nwendcode{}\nwbegincode{661}\moddef{SF base: procedures}\plusendmoddef
  subroutine sf_int_set_incoming (sf_int, incoming)
    class(sf_int_t), intent(inout) :: sf_int
    integer, dimension(:), intent(in) :: incoming
    allocate (sf_int%incoming (size (incoming)))
    sf_int%incoming = incoming
  end subroutine sf_int_set_incoming

  subroutine sf_int_set_radiated (sf_int, radiated)
    class(sf_int_t), intent(inout) :: sf_int
    integer, dimension(:), intent(in) :: radiated
    allocate (sf_int%radiated (size (radiated)))
    sf_int%radiated = radiated
  end subroutine sf_int_set_radiated

  subroutine sf_int_set_outgoing (sf_int, outgoing)
    class(sf_int_t), intent(inout) :: sf_int
    integer, dimension(:), intent(in) :: outgoing
    allocate (sf_int%outgoing (size (outgoing)))
    sf_int%outgoing = outgoing
  end subroutine sf_int_set_outgoing

\eatline
\eatline
\eatline
\nwendcode{}\nwbegindocs{662}Initialization.  This proceeds via an abstract data object, which
for the actual implementation should have the matching concrete type.
Since all implementations have the same signature, we can prepare a
deferred procedure.  The data object will become the target of a
corresponding pointer within the {\tt{}sf{\char95}int{\char95}t} implementation.

This should call the previous procedure.
\nwenddocs{}\nwbegincode{663}\moddef{SF base: sf int: TBP}\plusendmoddef
  procedure (sf_int_init), deferred :: init
\nwendcode{}\nwbegincode{664}\moddef{SF base: interfaces}\plusendmoddef
  abstract interface
     subroutine sf_int_init (sf_int, data)
       import
       class(sf_int_t), intent(out) :: sf_int
       class(sf_data_t), intent(in), target :: data
     end subroutine sf_int_init
  end interface
     
\eatline
\nwendcode{}\nwbegindocs{665}Complete initialization.  This routine contains initializations that can
only be performed after the interaction object got its final shape, i.e.,
redundant helicities have been eliminated by matching with beams and process.

The default implementation does nothing.
\nwenddocs{}\nwbegincode{666}\moddef{SF base: sf int: TBP}\plusendmoddef
  procedure :: setup_constants => sf_int_setup_constants
\nwendcode{}\nwbegincode{667}\moddef{SF base: procedures}\plusendmoddef
  subroutine sf_int_setup_constants (sf_int)
    class(sf_int_t), intent(inout) :: sf_int
  end subroutine sf_int_setup_constants
  
\eatline
\nwendcode{}\nwbegindocs{668}Set beam indices, i.e., the beam(s) on which
this structure function applies.
\nwenddocs{}\nwbegincode{669}\moddef{SF base: sf int: TBP}\plusendmoddef
  procedure :: set_beam_index => sf_int_set_beam_index
\nwendcode{}\nwbegincode{670}\moddef{SF base: procedures}\plusendmoddef
  subroutine sf_int_set_beam_index (sf_int, beam_index)
    class(sf_int_t), intent(inout) :: sf_int
    integer, dimension(:), intent(in) :: beam_index
    allocate (sf_int%beam_index (size (beam_index)))
    sf_int%beam_index = beam_index
  end subroutine sf_int_set_beam_index

\eatline
\nwendcode{}\nwbegindocs{671}Set parameter indices, indicating which MC input parameters are to
be used for evaluating this structure function.
\nwenddocs{}\nwbegincode{672}\moddef{SF base: sf int: TBP}\plusendmoddef
  procedure :: set_par_index => sf_int_set_par_index
\nwendcode{}\nwbegincode{673}\moddef{SF base: procedures}\plusendmoddef
  subroutine sf_int_set_par_index (sf_int, par_index)
    class(sf_int_t), intent(inout) :: sf_int
    integer, dimension(:), intent(in) :: par_index
    allocate (sf_int%par_index (size (par_index)))
    sf_int%par_index = par_index
  end subroutine sf_int_set_par_index

\eatline
\nwendcode{}\nwbegindocs{674}Initialize the structure-function kinematics, setting incoming
momenta.  We assume that array shapes match.

Three versions.  The first version relies on the momenta being linked
to another interaction.  The second version sets the momenta
explicitly.  In the third version, we first compute momenta for the
specified energies and store those.
\nwenddocs{}\nwbegincode{675}\moddef{SF base: sf int: TBP}\plusendmoddef
  generic :: seed_kinematics => sf_int_receive_momenta
  generic :: seed_kinematics => sf_int_seed_momenta
  generic :: seed_kinematics => sf_int_seed_energies
  procedure :: sf_int_receive_momenta
  procedure :: sf_int_seed_momenta
  procedure :: sf_int_seed_energies
\nwendcode{}\nwbegincode{676}\moddef{SF base: procedures}\plusendmoddef
  subroutine sf_int_receive_momenta (sf_int)
    class(sf_int_t), intent(inout) :: sf_int
    if (sf_int%status >= SF_INITIAL) then
       call interaction_receive_momenta (sf_int%interaction_t)
       sf_int%status = SF_SEED_KINEMATICS
    end if
  end subroutine sf_int_receive_momenta

  subroutine sf_int_seed_momenta (sf_int, k)
    class(sf_int_t), intent(inout) :: sf_int
    type(vector4_t), dimension(:), intent(in) :: k
    if (sf_int%status >= SF_INITIAL) then
       call interaction_set_momenta (sf_int%interaction_t, k, &
            outgoing=.false.)
       sf_int%status = SF_SEED_KINEMATICS
    end if
  end subroutine sf_int_seed_momenta
  
  subroutine sf_int_seed_energies (sf_int, E)
    class(sf_int_t), intent(inout) :: sf_int
    real(default), dimension(:), intent(in) :: E
    type(vector4_t), dimension(:), allocatable :: k
    integer :: j
    if (sf_int%status >= SF_INITIAL) then
       allocate (k (size (E)))
       if (all (E**2 >= sf_int%mi2)) then
          do j = 1, size (E)
             k(j) = vector4_moving (E(j), &
                  (3-2*j) * sqrt (E(j)**2 - sf_int%mi2(j)), 3)
          end do
          call sf_int%seed_kinematics (k)
       end if
    end if
  end subroutine sf_int_seed_energies
  
\eatline
\eatline
\nwendcode{}\nwbegindocs{677}Tell if in generator mode.  By default, this is false.  To be
overridden where appropriate; we may refer to the {\tt{}is{\char95}generator}
method of the {\tt{}data} component in the concrete type.
\nwenddocs{}\nwbegincode{678}\moddef{SF base: sf int: TBP}\plusendmoddef
  procedure :: is_generator => sf_int_is_generator
\nwendcode{}\nwbegincode{679}\moddef{SF base: procedures}\plusendmoddef
  function sf_int_is_generator (sf_int) result (flag)
    class(sf_int_t), intent(in) :: sf_int
    logical :: flag
    flag = .false.
  end function sf_int_is_generator

\eatline
\nwendcode{}\nwbegindocs{680}Generate free parameters {\tt{}r}.  Parameters are free if they do not
correspond to integration parameters (i.e., are bound), but are
generated by the structure function object itself.  By default, all
parameters are bound, and the output values of this procedure will be
discarded.  With free parameters, we have to override this procedure.  

The value {\tt{}x{\char95}free} is the renormalization factor of the total energy
that corresponds to the free parameters.  If there are no free
parameters, the procedure will not change its value, which starts as
unity.  Otherwise, the fraction is typically decreased, but may also
be increased in some cases.
\nwenddocs{}\nwbegincode{681}\moddef{SF base: sf int: TBP}\plusendmoddef
  procedure :: generate_free => sf_int_generate_free
\nwendcode{}\nwbegincode{682}\moddef{SF base: procedures}\plusendmoddef
  subroutine sf_int_generate_free (sf_int, r, rb,  x_free)
    class(sf_int_t), intent(inout) :: sf_int
    real(default), dimension(:), intent(out) :: r, rb
    real(default), intent(inout) :: x_free
    r = 0
    rb= 1
  end subroutine sf_int_generate_free
    
\eatline
\nwendcode{}\nwbegindocs{683}Complete the structure-function kinematics, derived from an input
parameter (array) $r$ between 0 and 1.  The interaction momenta are
calculated, and we return $x$ (the momentum fraction), and $f$ (the
Jacobian factor for the map $r\to x$), if {\tt{}map} is set.

If the {\tt{}map} flag is unset, $r$ and $x$ values will coincide, and $f$ will
become unity.  If it is set, the structure-function implementation chooses a
convenient mapping from $r$ to $x$ with Jacobian $f$.

In the {\tt{}inverse{\char95}kinematics} variant, we exchange the intent of {\tt{}x}
and {\tt{}r}.  The momenta are calculated only if the optional flag
{\tt{}set{\char95}momenta} is present and set.
\nwenddocs{}\nwbegincode{684}\moddef{SF base: sf int: TBP}\plusendmoddef
  procedure (sf_int_complete_kinematics), deferred :: complete_kinematics
  procedure (sf_int_inverse_kinematics), deferred :: inverse_kinematics
\nwendcode{}\nwbegincode{685}\moddef{SF base: interfaces}\plusendmoddef
  abstract interface
     subroutine sf_int_complete_kinematics (sf_int, x, f, r, rb, map)
       import
       class(sf_int_t), intent(inout) :: sf_int
       real(default), dimension(:), intent(out) :: x
       real(default), intent(out) :: f
       real(default), dimension(:), intent(in) :: r
       real(default), dimension(:), intent(in) :: rb
       logical, intent(in) :: map
     end subroutine sf_int_complete_kinematics
  end interface
  
  abstract interface
     subroutine sf_int_inverse_kinematics (sf_int, x, f, r, rb, map, &
          set_momenta)
       import
       class(sf_int_t), intent(inout) :: sf_int
       real(default), dimension(:), intent(in) :: x
       real(default), intent(out) :: f
       real(default), dimension(:), intent(out) :: r
       real(default), dimension(:), intent(out) :: rb
       logical, intent(in) :: map
       logical, intent(in), optional :: set_momenta
     end subroutine sf_int_inverse_kinematics
  end interface
  
\eatline
\eatline
\nwendcode{}\nwbegindocs{686}Single splitting: compute momenta, given $x$ input parameters.  We
assume that the incoming momentum is set.  The status code is set to
{\tt{}SF{\char95}FAILED{\char95}KINEMATICS} if
the $x$ array does not correspond to a valid momentum configuration.
Otherwise, it is updated to {\tt{}SF{\char95}DONE{\char95}KINEMATICS}.

We force the outgoing particle on-shell.  The on-shell projection is
determined by the {\tt{}on{\char95}shell{\char95}mode}.  The radiated particle should already be
on shell.
\nwenddocs{}\nwbegincode{687}\moddef{SF base: sf int: TBP}\plusendmoddef
  procedure :: split_momentum => sf_int_split_momentum
\nwendcode{}\nwbegincode{688}\moddef{SF base: procedures}\plusendmoddef
  subroutine sf_int_split_momentum (sf_int, x, xb1)
    class(sf_int_t), intent(inout) :: sf_int
    real(default), dimension(:), intent(in) :: x
    real(default), intent(in) :: xb1
    type(vector4_t) :: k
    type(vector4_t), dimension(2) :: q
    type(splitting_data_t) :: sd
    real(default) :: E1, E2
    logical :: fail
    if (sf_int%status >= SF_SEED_KINEMATICS) then
       k = interaction_get_momentum (sf_int%interaction_t, 1)
       call sd%init (k, &
            sf_int%mi2(1), sf_int%mr2(1), sf_int%mo2(1), &
            collinear = size (x) == 1)
       call sd%set_t_bounds (x(1), xb1)
       select case (size (x))
       case (1)
       case (3)
          if (sf_int%qmax_defined) then
             if (sf_int%qmin_defined) then
                call sd%sample_t (x(2), &
                     t0 = - sf_int%qmax(1) ** 2, t1 = - sf_int%qmin(1) ** 2)
             else
                call sd%sample_t (x(2), &
                     t0 = - sf_int%qmax(1) ** 2)
             end if
          else
             if (sf_int%qmin_defined) then
                call sd%sample_t (x(2), t1 = - sf_int%qmin(1) ** 2)
             else
                call sd%sample_t (x(2))
             end if
          end if
          call sd%sample_phi (x(3))
       case default
          call msg_bug ("Structure function: impossible number of parameters")
       end select
       q = sd%split_momentum (k)
       call on_shell (q, [sf_int%mr2, sf_int%mo2], &
            sf_int%on_shell_mode)
       call interaction_set_momenta (sf_int%interaction_t, &
            q, outgoing=.true.)
       E1 = energy (q(1))
       E2 = energy (q(2))
       fail = E1 < 0 .or. E2 < 0 &
            .or. E1 ** 2 < sf_int%mr2(1) &
            .or. E2 ** 2 < sf_int%mo2(1)
       if (fail) then
          sf_int%status = SF_FAILED_KINEMATICS
       else
          sf_int%status = SF_DONE_KINEMATICS
       end if
    end if
  end subroutine sf_int_split_momentum
    
\eatline
\nwendcode{}\nwbegindocs{689}Pair splitting: two incoming momenta, two radiated, two outgoing.
This is simple because we insist on all momenta being collinear.
\nwenddocs{}\nwbegincode{690}\moddef{SF base: sf int: TBP}\plusendmoddef
  procedure :: split_momenta => sf_int_split_momenta
\nwendcode{}\nwbegincode{691}\moddef{SF base: procedures}\plusendmoddef
  subroutine sf_int_split_momenta (sf_int, x, xb1)
    class(sf_int_t), intent(inout) :: sf_int
    real(default), dimension(:), intent(in) :: x
    real(default), dimension(:), intent(in) :: xb1
    type(vector4_t), dimension(2) :: k
    type(vector4_t), dimension(4) :: q
    real(default), dimension(4) :: E
    logical :: fail
    if (sf_int%status >= SF_SEED_KINEMATICS) then
       select case (size (x))
       case (2)
       case default
          call msg_bug ("Pair structure function: recoil requested &
               &but not implemented yet")
       end select
       k(1) = interaction_get_momentum (sf_int%interaction_t, 1)
       k(2) = interaction_get_momentum (sf_int%interaction_t, 2)
       q(1:2) = xb1 * k
       q(3:4) = x * k
       select case (size (sf_int%mr2))
       case (2)
          call on_shell (q, &
               [sf_int%mr2(1), sf_int%mr2(2), &
               sf_int%mo2(1), sf_int%mo2(2)], &
               sf_int%on_shell_mode)
          call interaction_set_momenta (sf_int%interaction_t, &
               q, outgoing=.true.)
          E = energy (q)
          fail = any (E < 0) &
               .or. any (E(1:2) ** 2 < sf_int%mr2) &
               .or. any (E(3:4) ** 2 < sf_int%mo2)
       case default;  call msg_bug ("split momenta: incorrect use")
       end select
       if (fail) then
          sf_int%status = SF_FAILED_KINEMATICS
       else
          sf_int%status = SF_DONE_KINEMATICS
       end if
    end if
  end subroutine sf_int_split_momenta
    
\eatline
\nwendcode{}\nwbegindocs{692}Pair spectrum: the reduced version of the previous splitting,
without radiated momenta.
\nwenddocs{}\nwbegincode{693}\moddef{SF base: sf int: TBP}\plusendmoddef
  procedure :: reduce_momenta => sf_int_reduce_momenta
\nwendcode{}\nwbegincode{694}\moddef{SF base: procedures}\plusendmoddef
  subroutine sf_int_reduce_momenta (sf_int, x)
    class(sf_int_t), intent(inout) :: sf_int
    real(default), dimension(:), intent(in) :: x
    type(vector4_t), dimension(2) :: k
    type(vector4_t), dimension(2) :: q
    real(default), dimension(2) :: E
    logical :: fail
    if (sf_int%status >= SF_SEED_KINEMATICS) then
       select case (size (x))
       case (2)
       case default
          call msg_bug ("Pair spectrum: recoil requested &
               &but not implemented yet")
       end select
       k(1) = interaction_get_momentum (sf_int%interaction_t, 1)
       k(2) = interaction_get_momentum (sf_int%interaction_t, 2)
       q = x * k
       call on_shell (q, &
            [sf_int%mo2(1), sf_int%mo2(2)], &
            sf_int%on_shell_mode)
       call interaction_set_momenta (sf_int%interaction_t, &
            q, outgoing=.true.)
       E = energy (q)
       fail = any (E < 0) &
            .or. any (E ** 2 < sf_int%mo2)
       if (fail) then
          sf_int%status = SF_FAILED_KINEMATICS
       else
          sf_int%status = SF_DONE_KINEMATICS
       end if
    end if
  end subroutine sf_int_reduce_momenta
    
\eatline
\nwendcode{}\nwbegindocs{695}The inverse procedure: we compute the {\tt{}x} array from the momentum
configuration.  In an overriding TBP, we may also set internal data
that depend on this, for convenience.

NOTE: Here and above, the single-particle case is treated in detail,
allowing for non-collinearity and non-vanishing masses and nontrivial
momentum-transfer bounds.  For the pair case, we currently implement
only collinear splitting and assume massless particles.  This should
be improved.

NOTE: We should also improve this in the sense of extracting both $x$
and $\bar x$, where for the latter we may consult the beam remnant.
The current implementation is not sensitive to small $\bar x$ values.
\nwenddocs{}\nwbegincode{696}\moddef{SF base: sf int: TBP}\plusendmoddef
  procedure :: recover_x => sf_int_recover_x
  procedure :: base_recover_x => sf_int_recover_x
\nwendcode{}\nwbegincode{697}\moddef{SF base: procedures}\plusendmoddef
  subroutine sf_int_recover_x (sf_int, x, x_free)
    class(sf_int_t), intent(inout) :: sf_int
    real(default), dimension(:), intent(out) :: x
    real(default), intent(inout), optional :: x_free
    type(vector4_t), dimension(:), allocatable :: k
    type(vector4_t), dimension(:), allocatable :: q
    type(splitting_data_t) :: sd
    if (sf_int%status >= SF_SEED_KINEMATICS) then
       allocate (k (interaction_get_n_in (sf_int%interaction_t)))
       allocate (q (interaction_get_n_out (sf_int%interaction_t)))
       k = interaction_get_momenta (sf_int%interaction_t, outgoing=.false.)
       q = interaction_get_momenta (sf_int%interaction_t, outgoing=.true.)
       select case (size (k))
       case (1)
          call sd%init (k(1), &
               sf_int%mi2(1), sf_int%mr2(1), sf_int%mo2(1), &
               collinear = size (x) == 1) 
          call sd%recover (k(1), q(2), sf_int%on_shell_mode)
          x(1) = sd%get_x ()
          select case (size (x))
          case (1)
          case (3)
             if (sf_int%qmax_defined) then
                if (sf_int%qmin_defined) then
                   call sd%inverse_t (x(2), &
                        t0 = - sf_int%qmax(1) ** 2, t1 = - sf_int%qmin(1) ** 2)
                else
                   call sd%inverse_t (x(2), &
                        t0 = - sf_int%qmax(1) ** 2)
                end if
             else
                if (sf_int%qmin_defined) then
                   call sd%inverse_t (x(2), t1 = - sf_int%qmin(1) ** 2)
                else
                   call sd%inverse_t (x(2)) 
                end if
             end if
             call sd%inverse_phi (x(3))
          case default
             call msg_bug ("Structure function: impossible number &
                  &of parameters")
          end select
       case (2)
          select case (size (x))
          case (2)
          case default
             call msg_bug ("Pair structure function: recoil requested &
                  &but not implemented yet")
          end select
          select case (sf_int%on_shell_mode)
          case (KEEP_ENERGY)
             select case (size (q))
             case (4)
                x = energy (q(3:4)) / energy (k)
             case (2)
                x = energy (q) / energy (k)
             end select
          case (KEEP_MOMENTUM)
             select case (size (q))
             case (4)
                x = longitudinal_part (q(3:4)) / longitudinal_part (k)
             case (2)
                x = longitudinal_part (q) / longitudinal_part (k)
             end select
          end select
       end select
    end if
  end subroutine sf_int_recover_x
  
\eatline
\nwendcode{}\nwbegindocs{698}Apply the structure function, i.e., evaluate the interaction.  For
the calculation, we may use the stored momenta, any further
information stored inside the {\tt{}sf{\char95}int} implementation during
kinematics setup, and the given energy scale.  It may happen that for
the given kinematics the value is not defined.  This should be
indicated by the status code.
\nwenddocs{}\nwbegincode{699}\moddef{SF base: sf int: TBP}\plusendmoddef
  procedure (sf_int_apply), deferred :: apply
\nwendcode{}\nwbegincode{700}\moddef{SF base: interfaces}\plusendmoddef
  abstract interface
     subroutine sf_int_apply (sf_int, scale)
       import
       class(sf_int_t), intent(inout) :: sf_int
       real(default), intent(in) :: scale
     end subroutine sf_int_apply
  end interface

\eatline
\nwendcode{}\nwbegindocs{701}\nwdocspar
\subsection{Accessing the structure function}
Return metadata.  Once {\tt{}interaction{\char95}t} is rewritten in OO, some of this will
be inherited.

The number of outgoing is equal to the number of incoming particles.  The
radiated particles are the difference.
\nwenddocs{}\nwbegincode{702}\moddef{SF base: sf int: TBP}\plusendmoddef
  procedure :: get_n_in => sf_int_get_n_in
  procedure :: get_n_rad => sf_int_get_n_rad
  procedure :: get_n_out => sf_int_get_n_out
\nwendcode{}\nwbegincode{703}\moddef{SF base: procedures}\plusendmoddef
  function sf_int_get_n_in (sf_int) result (n_in)
    class(sf_int_t), intent(in) :: sf_int
    integer :: n_in
    n_in = interaction_get_n_in (sf_int%interaction_t)
  end function sf_int_get_n_in
  
  function sf_int_get_n_rad (sf_int) result (n_rad)
    class(sf_int_t), intent(in) :: sf_int
    integer :: n_rad
    n_rad = interaction_get_n_out (sf_int%interaction_t) &
         - interaction_get_n_in (sf_int%interaction_t)
  end function sf_int_get_n_rad
  
  function sf_int_get_n_out (sf_int) result (n_out)
    class(sf_int_t), intent(in) :: sf_int
    integer :: n_out
    n_out = interaction_get_n_in (sf_int%interaction_t)
  end function sf_int_get_n_out
  
\eatline
\eatline
\eatline
\nwendcode{}\nwbegindocs{704}Number of matrix element entries in the interaction:
\nwenddocs{}\nwbegincode{705}\moddef{SF base: sf int: TBP}\plusendmoddef
  procedure :: get_n_states => sf_int_get_n_states
\nwendcode{}\nwbegincode{706}\moddef{SF base: procedures}\plusendmoddef
  function sf_int_get_n_states (sf_int) result (n_states)
    class(sf_int_t), intent(in) :: sf_int
    integer :: n_states
    n_states = interaction_get_n_matrix_elements (sf_int%interaction_t)
  end function sf_int_get_n_states
  
\eatline
\nwendcode{}\nwbegindocs{707}Return a specific state as a quantum-number array.
\nwenddocs{}\nwbegincode{708}\moddef{SF base: sf int: TBP}\plusendmoddef
  procedure :: get_state => sf_int_get_state
\nwendcode{}\nwbegincode{709}\moddef{SF base: procedures}\plusendmoddef
  function sf_int_get_state (sf_int, i) result (qn)
    class(sf_int_t), intent(in) :: sf_int
    type(quantum_numbers_t), dimension(:), allocatable :: qn
    integer, intent(in) :: i
    allocate (qn (interaction_get_n_tot (sf_int%interaction_t)))
    qn = interaction_get_quantum_numbers (sf_int%interaction_t, i)
  end function sf_int_get_state

\eatline
\nwendcode{}\nwbegindocs{710}Return the matrix-element values for all states.  We can assume that
the matrix elements are real, so we take the real part.
\nwenddocs{}\nwbegincode{711}\moddef{SF base: sf int: TBP}\plusendmoddef
  procedure :: get_values => sf_int_get_values
\nwendcode{}\nwbegincode{712}\moddef{SF base: procedures}\plusendmoddef
  subroutine sf_int_get_values (sf_int, value)
    class(sf_int_t), intent(in) :: sf_int
    real(default), dimension(:), intent(out) :: value
    integer :: i
    if (sf_int%status >= SF_EVALUATED) then
       do i = 1, size (value)
          value(i) = interaction_get_matrix_element &
               (sf_int%interaction_t, i)
       end do
    else
       value = 0
    end if
  end subroutine sf_int_get_values

\eatline
\nwendcode{}\nwbegindocs{713}\nwdocspar
\subsection{Direct calculations}
Compute a structure function value (array) directly, given an array of $x$
values and a scale.  If the energy is also given, we initialize the
kinematics for that energy, otherwise take it from a previous run.

We assume that the {\tt{}E} array has dimension {\tt{}n{\char95}in}, and the {\tt{}x}
array has {\tt{}n{\char95}par}. 
\nwenddocs{}\nwbegincode{714}\moddef{SF base: sf int: TBP}\plusendmoddef
  procedure :: compute_values => sf_int_compute_values
\nwendcode{}\nwbegincode{715}\moddef{SF base: procedures}\plusendmoddef
  subroutine sf_int_compute_values (sf_int, value, x, xb, scale, E)
    class(sf_int_t), intent(inout) :: sf_int
    real(default), dimension(:), intent(out) :: value
    real(default), dimension(:), intent(in) :: x
    real(default), dimension(:), intent(in) :: xb
    real(default), intent(in) :: scale
    real(default), dimension(:), intent(in), optional :: E
    real(default), dimension(size (x)) :: xx
    real(default) :: f
    if (present (E))  call sf_int%seed_kinematics (E)
    if (sf_int%status >= SF_SEED_KINEMATICS) then
       call sf_int%complete_kinematics (xx, f, x, xb, map=.false.)
       call sf_int%apply (scale)
       call sf_int%get_values (value)
       value = value * f
    else
       value = 0
    end if
  end subroutine sf_int_compute_values

\eatline
\nwendcode{}\nwbegindocs{716}Compute just a single value for one of the states, i.e., throw the
others away.
\nwenddocs{}\nwbegincode{717}\moddef{SF base: sf int: TBP}\plusendmoddef
  procedure :: compute_value => sf_int_compute_value
\nwendcode{}\nwbegincode{718}\moddef{SF base: procedures}\plusendmoddef
  subroutine sf_int_compute_value &
       (sf_int, i_state, value, x, xb, scale, E)
    class(sf_int_t), intent(inout) :: sf_int
    integer, intent(in) :: i_state
    real(default), intent(out) :: value
    real(default), dimension(:), intent(in) :: x
    real(default), dimension(:), intent(in) :: xb
    real(default), intent(in) :: scale
    real(default), dimension(:), intent(in), optional :: E
    real(default), dimension(:), allocatable :: value_array
    if (sf_int%status >= SF_INITIAL) then
       allocate (value_array (sf_int%get_n_states ()))
       call sf_int%compute_values (value_array, x, xb, scale, E)
       value = value_array(i_state)
    else
       value = 0
    end if
  end subroutine sf_int_compute_value

\eatline
\nwendcode{}\nwbegindocs{719}\nwdocspar
\subsection{Structure-function instance}
This is a wrapper for {\tt{}sf{\char95}int{\char95}t} objects, such that we can
build an array with different structure-function types.  The
structure-function contains an array (a sequence) of {\tt{}sf{\char95}int{\char95}t}
objects.

The object, it holds the evaluator that connects the preceding part of the
structure-function chain to the current interaction.

It also stores the input and output parameter values for the
contained structure function.  The {\tt{}r} array has a second dimension,
corresponding to the mapping channels in a multi-channel
configuration.  There is a Jacobian entry {\tt{}f} for each channel.  The
corresponding logical array {\tt{}mapping} tells whether we apply the
mapping appropriate for the current structure function in this channel.
The {\tt{}x} parameter values (energy fractions) are common to all
channels. 
\nwenddocs{}\nwbegincode{720}\moddef{SF base: types}\plusendmoddef
  type :: sf_instance_t
     class(sf_int_t), allocatable :: int
     type(evaluator_t) :: eval
     real(default), dimension(:,:), allocatable :: r
     real(default), dimension(:,:), allocatable :: rb
     real(default), dimension(:), allocatable :: f
     logical, dimension(:), allocatable :: m
     real(default), dimension(:), allocatable :: x
  end type sf_instance_t
  
\eatline
\nwendcode{}\nwbegindocs{721}\nwdocspar
\subsection{Structure-function chain}
A chain is an array of structure functions {\tt{}sf}, initiated by a beam setup.
We do not use this directly for evaluation, but create instances with
copies of the contained interactions.

{\tt{}n{\char95}par} is the total number of parameters that is necessary for
completely determining the structure-function chain.  {\tt{}n{\char95}bound} is
the number of MC input parameters that are requested from the
integrator.  The difference of {\tt{}n{\char95}par} and {\tt{}n{\char95}bound} is the number
of free parameters, which are generated by a structure-function
object in generator mode.
\nwenddocs{}\nwbegincode{722}\moddef{SF base: public}\plusendmoddef
  public :: sf_chain_t
\nwendcode{}\nwbegincode{723}\moddef{SF base: types}\plusendmoddef
  type, extends (beam_t) :: sf_chain_t
     type(beam_data_t), pointer :: beam_data => null ()
     integer :: n_in = 0
     integer :: n_strfun = 0
     integer :: n_par = 0
     integer :: n_bound = 0
     type(sf_instance_t), dimension(:), allocatable :: sf
     logical :: trace_enable = .false.
     integer :: trace_unit = 0
   contains
   \LA{}SF base: sf chain: TBP\RA{}
  end type sf_chain_t
     
\eatline
\nwendcode{}\nwbegindocs{724}Finalizer.
\nwenddocs{}\nwbegincode{725}\moddef{SF base: sf chain: TBP}\endmoddef
  procedure :: final => sf_chain_final
\nwendcode{}\nwbegincode{726}\moddef{SF base: procedures}\plusendmoddef
  subroutine sf_chain_final (object)
    class(sf_chain_t), intent(inout) :: object
    integer :: i
    call object%final_tracing ()
    if (allocated (object%sf)) then 
       do i = 1, size (object%sf, 1)
          associate (sf => object%sf(i))
            if (allocated (sf%int)) then
               call sf%int%final ()
            end if
          end associate
       end do
    end if
    call beam_final (object%beam_t)
  end subroutine sf_chain_final

\eatline
\nwendcode{}\nwbegindocs{727}Output.
\nwenddocs{}\nwbegincode{728}\moddef{SF base: sf chain: TBP}\plusendmoddef
  procedure :: write => sf_chain_write
\nwendcode{}\nwbegincode{729}\moddef{SF base: procedures}\plusendmoddef
  subroutine sf_chain_write (object, unit)
    class(sf_chain_t), intent(in) :: object
    integer, intent(in), optional :: unit
    integer :: u, i
    u = given_output_unit (unit)
    write (u, "(1x,A)")  "Incoming particles / structure-function chain:"
    if (associated (object%beam_data)) then
       write (u, "(3x,A,I0)")  "n_in      = ", object%n_in
       write (u, "(3x,A,I0)")  "n_strfun  = ", object%n_strfun
       write (u, "(3x,A,I0)")  "n_par     = ", object%n_par
       if (object%n_par /= object%n_bound) then
          write (u, "(3x,A,I0)")  "n_bound   = ", object%n_bound
       end if
       call beam_data_write (object%beam_data, u)
       call write_separator (u)
       call beam_write (object%beam_t, u)
       if (allocated (object%sf)) then
          do i = 1, object%n_strfun
             associate (sf => object%sf(i))
               call write_separator (u)
               if (allocated (sf%int)) then
                  call sf%int%write (u)
               else
                  write (u, "(1x,A)")  "SF instance: [undefined]"
               end if
             end associate
          end do
       end if
    else
       write (u, "(3x,A)")  "[undefined]"
    end if
  end subroutine sf_chain_write
  
\eatline
\nwendcode{}\nwbegindocs{730}Initialize: setup beams.  The {\tt{}beam{\char95}data} target must remain valid
for the lifetime of the chain, since we just establish a pointer.  The
structure-function configuration array is used to initialize the
individual structure-function entries.  The target attribute is needed
because the {\tt{}sf{\char95}int} entries establish pointers to the configuration data.
\nwenddocs{}\nwbegincode{731}\moddef{SF base: sf chain: TBP}\plusendmoddef
  procedure :: init => sf_chain_init
\nwendcode{}\nwbegincode{732}\moddef{SF base: procedures}\plusendmoddef
  subroutine sf_chain_init (sf_chain, beam_data, sf_config)
    class(sf_chain_t), intent(out) :: sf_chain
    type(beam_data_t), intent(in), target :: beam_data
    type(sf_config_t), dimension(:), intent(in), optional, target :: sf_config
    integer :: i
    sf_chain%beam_data => beam_data
    sf_chain%n_in = beam_data_get_n_in (beam_data)
    call beam_init (sf_chain%beam_t, beam_data)
    if (present (sf_config)) then
       sf_chain%n_strfun = size (sf_config)
       allocate (sf_chain%sf (sf_chain%n_strfun))
       do i = 1, sf_chain%n_strfun
          call sf_chain%set_strfun (i, sf_config(i)%i, sf_config(i)%data)
       end do
    end if
  end subroutine sf_chain_init
  
\eatline
\nwendcode{}\nwbegindocs{733}Receive the beam momenta from a source to which the beam interaction
is linked.
\nwenddocs{}\nwbegincode{734}\moddef{SF base: sf chain: TBP}\plusendmoddef
  procedure :: receive_beam_momenta => sf_chain_receive_beam_momenta
\nwendcode{}\nwbegincode{735}\moddef{SF base: procedures}\plusendmoddef
  subroutine sf_chain_receive_beam_momenta (sf_chain)
    class(sf_chain_t), intent(inout), target :: sf_chain
    type(interaction_t), pointer :: beam_int
    beam_int => sf_chain%get_beam_int_ptr ()
    call interaction_receive_momenta (beam_int)
  end subroutine sf_chain_receive_beam_momenta
  
\eatline
\nwendcode{}\nwbegindocs{736}Explicitly set the beam momenta.
\nwenddocs{}\nwbegincode{737}\moddef{SF base: sf chain: TBP}\plusendmoddef
  procedure :: set_beam_momenta => sf_chain_set_beam_momenta
\nwendcode{}\nwbegincode{738}\moddef{SF base: procedures}\plusendmoddef
  subroutine sf_chain_set_beam_momenta (sf_chain, p)
    class(sf_chain_t), intent(inout) :: sf_chain
    type(vector4_t), dimension(:), intent(in) :: p
    call beam_set_momenta (sf_chain%beam_t, p)
  end subroutine sf_chain_set_beam_momenta

\eatline
\nwendcode{}\nwbegindocs{739}Set a structure-function entry.  We use the {\tt{}data} input to
allocate the {\tt{}int} structure-function instance with appropriate
type, then initialize the entry.  The entry establishes a pointer to
{\tt{}data}.

The index {\tt{}i} is the structure-function index in the chain.
\nwenddocs{}\nwbegincode{740}\moddef{SF base: sf chain: TBP}\plusendmoddef
  procedure :: set_strfun => sf_chain_set_strfun
\nwendcode{}\nwbegincode{741}\moddef{SF base: procedures}\plusendmoddef
  subroutine sf_chain_set_strfun (sf_chain, i, beam_index, data)
    class(sf_chain_t), intent(inout) :: sf_chain
    integer, intent(in) :: i
    integer, dimension(:), intent(in) :: beam_index
    class(sf_data_t), intent(in), target :: data
    integer :: n_par, j
    n_par = data%get_n_par ()
    call data%allocate_sf_int (sf_chain%sf(i)%int)
    associate (sf_int => sf_chain%sf(i)%int)
      call sf_int%init (data)
      call sf_int%set_beam_index (beam_index)
      call sf_int%set_par_index &
           ([(j, j = sf_chain%n_par + 1, sf_chain%n_par + n_par)])
      sf_chain%n_par = sf_chain%n_par + n_par
      if (.not. data%is_generator ()) then
         sf_chain%n_bound = sf_chain%n_bound + n_par
      end if
    end associate
  end subroutine sf_chain_set_strfun
    
\eatline
\nwendcode{}\nwbegindocs{742}Return the number of structure-function parameters.
\nwenddocs{}\nwbegincode{743}\moddef{SF base: sf chain: TBP}\plusendmoddef
  procedure :: get_n_par => sf_chain_get_n_par
  procedure :: get_n_bound => sf_chain_get_n_bound
\nwendcode{}\nwbegincode{744}\moddef{SF base: procedures}\plusendmoddef
  function sf_chain_get_n_par (sf_chain) result (n)
    class(sf_chain_t), intent(in) :: sf_chain
    integer :: n
    n = sf_chain%n_par
  end function sf_chain_get_n_par
  
  function sf_chain_get_n_bound (sf_chain) result (n)
    class(sf_chain_t), intent(in) :: sf_chain
    integer :: n
    n = sf_chain%n_bound
  end function sf_chain_get_n_bound
  
\eatline
\eatline
\nwendcode{}\nwbegindocs{745}Return a pointer to the beam interaction.
\nwenddocs{}\nwbegincode{746}\moddef{SF base: sf chain: TBP}\plusendmoddef
  procedure :: get_beam_int_ptr => sf_chain_get_beam_int_ptr
\nwendcode{}\nwbegincode{747}\moddef{SF base: procedures}\plusendmoddef
  function sf_chain_get_beam_int_ptr (sf_chain) result (int)
    class(sf_chain_t), intent(in), target :: sf_chain
    type(interaction_t), pointer :: int
    int => beam_get_int_ptr (sf_chain%beam_t)
  end function sf_chain_get_beam_int_ptr
  
\eatline
\nwendcode{}\nwbegindocs{748}Enable the trace feature: record structure function data (input
parameters, $x$ values, evaluation result) to an external file.
\nwenddocs{}\nwbegincode{749}\moddef{SF base: sf chain: TBP}\plusendmoddef
  procedure :: setup_tracing => sf_chain_setup_tracing
  procedure :: final_tracing => sf_chain_final_tracing
\nwendcode{}\nwbegincode{750}\moddef{SF base: procedures}\plusendmoddef
  subroutine sf_chain_setup_tracing (sf_chain, file)
    class(sf_chain_t), intent(inout) :: sf_chain
    type(string_t), intent(in) :: file
    if (sf_chain%n_strfun > 0) then
       sf_chain%trace_enable = .true.
       sf_chain%trace_unit = free_unit ()
       open (sf_chain%trace_unit, file = char (file), action = "write", &
            status = "replace")
       call sf_chain%write_trace_header ()
    else
       call msg_error ("Beam structure: no structure functions, tracing &
            &disabled")
    end if
  end subroutine sf_chain_setup_tracing

  subroutine sf_chain_final_tracing (sf_chain)
    class(sf_chain_t), intent(inout) :: sf_chain
    if (sf_chain%trace_enable) then
       close (sf_chain%trace_unit)
       sf_chain%trace_enable = .false.
    end if
  end subroutine sf_chain_final_tracing

\eatline
\eatline
\nwendcode{}\nwbegindocs{751}Write the header for the tracing file.
\nwenddocs{}\nwbegincode{752}\moddef{SF base: sf chain: TBP}\plusendmoddef
  procedure :: write_trace_header => sf_chain_write_trace_header
\nwendcode{}\nwbegincode{753}\moddef{SF base: procedures}\plusendmoddef
  subroutine sf_chain_write_trace_header (sf_chain)
    class(sf_chain_t), intent(in) :: sf_chain
    integer :: u
    if (sf_chain%trace_enable) then
       u = sf_chain%trace_unit
       write (u, "('# ',A)")  "WHIZARD output: &
            &structure-function sampling data"
       write (u, "('# ',A,1x,I0)")  "Number of sf records:", sf_chain%n_strfun
       write (u, "('# ',A,1x,I0)")  "Number of parameters:", sf_chain%n_par
       write (u, "('# ',A)")  "Columns: channel, p(n_par), x(n_par), f, Jac * f"
    end if
  end subroutine sf_chain_write_trace_header
    
\eatline
\nwendcode{}\nwbegindocs{754}Write a record which collects the structure function data for the
current data point.  For the selected channel, we print first the
input integration parameters, then the $x$ values, then the
structure-function value summed over all quantum numbers, then the
structure function value times the mapping Jacobian.
\nwenddocs{}\nwbegincode{755}\moddef{SF base: sf chain: TBP}\plusendmoddef
  procedure :: trace => sf_chain_trace
\nwendcode{}\nwbegincode{756}\moddef{SF base: procedures}\plusendmoddef
  subroutine sf_chain_trace (sf_chain, c_sel, p, x, f, sf_sum)
    class(sf_chain_t), intent(in) :: sf_chain
    integer, intent(in) :: c_sel
    real(default), dimension(:,:), intent(in) :: p
    real(default), dimension(:), intent(in) :: x
    real(default), dimension(:), intent(in) :: f
    real(default), intent(in) :: sf_sum
    integer :: u, i
    if (sf_chain%trace_enable) then
       u = sf_chain%trace_unit
       write (u, "(1x,I0)", advance="no")  c_sel
       write (u, "(2x)", advance="no")
       do i = 1, sf_chain%n_par
          write (u, "(1x," // FMT_17 // ")", advance="no")  p(i,c_sel)
       end do
       write (u, "(2x)", advance="no")
       do i = 1, sf_chain%n_par
          write (u, "(1x," // FMT_17 // ")", advance="no")  x(i)
       end do
       write (u, "(2x)", advance="no")
       write (u, "(2(1x," // FMT_17 // "))")  sf_sum, f(c_sel) * sf_sum
    end if
  end subroutine sf_chain_trace
  
\eatline
\nwendcode{}\nwbegindocs{757}\nwdocspar
\subsection{Chain instances}
A structure-function chain instance contains copies of the
interactions in the configuration chain, suitably linked to each other
and connected by evaluators.

After initialization, {\tt{}out{\char95}sf} should point, for each beam, to the
last structure function that affects this beam.  {\tt{}out{\char95}sf{\char95}i} should
indicate the index of the corresponding outgoing particle within that
structure-function interaction.

Analogously, {\tt{}out{\char95}eval} is the last evaluator in the
structure-function chain, which contains the complete set of outgoing
particles. {\tt{}out{\char95}eval{\char95}i} should indicate the index of the outgoing
particles, within that evaluator, which will initiate the collision.

When calculating actual kinematics, we fill the {\tt{}p}, {\tt{}r}, and
{\tt{}x} arrays and the {\tt{}f} factor.  The {\tt{}p} array denotes the MC
input parameters as they come from the random-number generator.  The
{\tt{}r} array results from applying global mappings.  The {\tt{}x} array
results from applying structure-function local mappings.  The $x$
values can be interpreted directly as momentum fractions (or angle
fractions, where recoil is involved).  The {\tt{}f} factor is the
Jacobian that results from applying all mappings.

The {\tt{}mapping} entry may store a global mapping that is applied to a
combination of $x$ values and structure functions, as opposed to mappings that
affect only a single structure function.  It is applied before the latter
mappings, in the transformation from the {\tt{}p} array to the {\tt{}r} array.  For
parameters affected by this mapping, we should ensure that they are not
involved in a local mapping.
\nwenddocs{}\nwbegincode{758}\moddef{SF base: public}\plusendmoddef
  public :: sf_chain_instance_t
\nwendcode{}\nwbegincode{759}\moddef{SF base: types}\plusendmoddef
  type, extends (beam_t) :: sf_chain_instance_t
     type(sf_chain_t), pointer :: config => null ()
     integer :: status = SF_UNDEFINED
     type(sf_instance_t), dimension(:), allocatable :: sf
     integer, dimension(:), allocatable :: out_sf
     integer, dimension(:), allocatable :: out_sf_i
     integer :: out_eval = 0
     integer, dimension(:), allocatable :: out_eval_i
     integer :: selected_channel = 0
     real(default), dimension(:,:), allocatable :: p, pb
     real(default), dimension(:,:), allocatable :: r, rb
     real(default), dimension(:), allocatable :: f
     real(default), dimension(:), allocatable :: x
     logical, dimension(:), allocatable :: bound
     real(default) :: x_free = 1
     type(sf_channel_t), dimension(:), allocatable :: channel
   contains
   \LA{}SF base: sf chain instance: TBP\RA{}
  end type sf_chain_instance_t
     
\eatline
\nwendcode{}\nwbegindocs{760}Finalizer.
\nwenddocs{}\nwbegincode{761}\moddef{SF base: sf chain instance: TBP}\endmoddef
  procedure :: final => sf_chain_instance_final
\nwendcode{}\nwbegincode{762}\moddef{SF base: procedures}\plusendmoddef
  subroutine sf_chain_instance_final (object)
    class(sf_chain_instance_t), intent(inout) :: object
    integer :: i
    if (allocated (object%sf)) then
       do i = 1, size (object%sf, 1)
          associate (sf => object%sf(i))
            if (allocated (sf%int)) then
               call evaluator_final (sf%eval)
               call sf%int%final ()
            end if
          end associate
       end do
    end if
    call beam_final (object%beam_t)
  end subroutine sf_chain_instance_final

\eatline
\nwendcode{}\nwbegindocs{763}Output.

Note: nagfor 5.3.1 appears to be slightly confused with the allocation
status.  We check both for allocation and nonzero size.
\nwenddocs{}\nwbegincode{764}\moddef{SF base: sf chain instance: TBP}\plusendmoddef
  procedure :: write => sf_chain_instance_write
\nwendcode{}\nwbegincode{765}\moddef{SF base: procedures}\plusendmoddef
  subroutine sf_chain_instance_write (object, unit)
    class(sf_chain_instance_t), intent(in) :: object
    integer, intent(in), optional :: unit
    integer :: u, i, c
    u = given_output_unit (unit)
    write (u, "(1x,A)", advance="no")  "Structure-function chain instance:"
    call write_sf_status (object%status, u)
    if (allocated (object%out_sf)) then
       write (u, "(3x,A)", advance="no")  "outgoing (interactions) ="
       do i = 1, size (object%out_sf)
          write (u, "(1x,I0,':',I0)", advance="no") &
               object%out_sf(i), object%out_sf_i(i)
       end do
       write (u, *)
    end if
    if (object%out_eval /= 0) then
       write (u, "(3x,A)", advance="no")  "outgoing (evaluators)   ="
       do i = 1, size (object%out_sf)
          write (u, "(1x,I0,':',I0)", advance="no") &
               object%out_eval, object%out_eval_i(i)
       end do
       write (u, *)
    end if
    if (allocated (object%sf)) then
       if (size (object%sf) /= 0) then
          write (u, "(1x,A)")  "Structure-function parameters:"
          do c = 1, size (object%f)
             write (u, "(1x,A,I0,A)", advance="no")  "Channel #", c, ":"
             if (c == object%selected_channel) then
                write (u, "(1x,A)")  "[selected]"
             else
                write (u, *)
             end if
             write (u, "(3x,A,9(1x,F9.7))")  "p =", object%p(:,c)
             write (u, "(3x,A,9(1x,F9.7))")  "r =", object%r(:,c)
             write (u, "(3x,A,9(1x,ES13.7))")  "f =", object%f(c)
             write (u, "(3x,A)", advance="no") "m ="
             call object%channel(c)%write (u)
          end do
          write (u, "(3x,A,9(1x,F9.7))")  "x =", object%x
          if (.not. all (object%bound)) then
             write (u, "(3x,A,9(1x,L1))")  "bound =", object%bound
          end if
       end if
    end if
    call write_separator (u)
    call beam_write (object%beam_t, u)
    if (allocated (object%sf)) then
       do i = 1, size (object%sf)
          associate (sf => object%sf(i))
            call write_separator (u)
            if (allocated (sf%int)) then
               if (allocated (sf%r)) then
                  write (u, "(1x,A)")  "Structure-function parameters:"
                  do c = 1, size (sf%f)
                     write (u, "(1x,A,I0,A)", advance="no")  "Channel #", c, ":"
                     if (c == object%selected_channel) then
                        write (u, "(1x,A)")  "[selected]"
                     else
                        write (u, *)
                     end if
                     write (u, "(3x,A,9(1x,F9.7))")  "r =", sf%r(:,c)
                     write (u, "(3x,A,9(1x,ES13.7))")  "f =", sf%f(c)
                     write (u, "(3x,A,9(1x,L1,7x))") "m =", sf%m(c)
                  end do
                  write (u, "(3x,A,9(1x,F9.7))")  "x =", sf%x
               end if
               call sf%int%write (u)
               if (.not. evaluator_is_empty (sf%eval)) then
                     call sf%eval%write (u)
               end if
            end if
          end associate
       end do
    end if
  end subroutine sf_chain_instance_write
  
\eatline
\nwendcode{}\nwbegindocs{766}Initialize.  This creates a copy of the interactions in the
configuration chain, assumed to be properly initialized.  In the copy,
we allocate the {\tt{}p} etc.\ arrays.

The brute-force assignment of the {\tt{}sf} component would be
straightforward, but at least gfortran 4.6.3 would like a more
fine-grained copy.  In any case, the copy is deep
as far as allocatables are concerned, but for the contained
{\tt{}interaction{\char95}t} objects the copy is shallow, as long as we do not
bind defined assignment to the type.  Therefore, we have to re-assign
the {\tt{}interaction{\char95}t} components explicitly, this time calling the
proper defined assignment.  Furthermore, we allocate the parameter
arrays for each structure function.
\nwenddocs{}\nwbegincode{767}\moddef{SF base: sf chain instance: TBP}\plusendmoddef
  procedure :: init => sf_chain_instance_init
\nwendcode{}\nwbegincode{768}\moddef{SF base: procedures}\plusendmoddef
  subroutine sf_chain_instance_init (chain, config, n_channel)
    class(sf_chain_instance_t), intent(out), target :: chain
    type(sf_chain_t), intent(in), target :: config
    integer, intent(in) :: n_channel
    integer :: i, j
    integer :: n_par_tot, n_par, n_strfun
    chain%config => config
    n_strfun = config%n_strfun
    chain%beam_t = config%beam_t
    allocate (chain%out_sf (config%n_in), chain%out_sf_i (config%n_in))
    allocate (chain%out_eval_i (config%n_in))
    chain%out_sf = 0
    chain%out_sf_i = [(i, i = 1, config%n_in)]
    chain%out_eval_i = chain%out_sf_i
    n_par_tot = 0
    if (n_strfun /= 0) then
       allocate (chain%sf (n_strfun))
       do i = 1, n_strfun
          associate (sf => chain%sf(i))
            allocate (sf%int, source=config%sf(i)%int)
            sf%int%interaction_t = config%sf(i)%int%interaction_t
            n_par = size (sf%int%par_index)
            allocate (sf%r (n_par, n_channel));  sf%r = 0
            allocate (sf%rb(n_par, n_channel));  sf%rb= 0
            allocate (sf%f (n_channel));         sf%f = 0
            allocate (sf%m (n_channel));         sf%m = .false.
            allocate (sf%x (n_par));             sf%x = 0
            n_par_tot = n_par_tot + n_par
          end associate
       end do
       allocate (chain%p (n_par_tot, n_channel));  chain%p = 0
       allocate (chain%pb(n_par_tot, n_channel));  chain%pb= 0
       allocate (chain%r (n_par_tot, n_channel));  chain%r = 0
       allocate (chain%rb(n_par_tot, n_channel));  chain%rb= 0
       allocate (chain%f (n_channel));             chain%f = 0
       allocate (chain%x (n_par_tot));             chain%x = 0
       call allocate_sf_channels &
            (chain%channel, n_channel=n_channel, n_strfun=n_strfun)
    end if
    allocate (chain%bound (n_par_tot), source = .true.)
    do i = 1, n_strfun
       associate (sf => chain%sf(i))
         if (sf%int%is_generator ()) then
            do j = 1, size (sf%int%par_index)
               chain%bound(sf%int%par_index(j)) = .false.
            end do
         end if
       end associate
    end do
    chain%status = SF_INITIAL
  end subroutine sf_chain_instance_init
  
\eatline
\nwendcode{}\nwbegindocs{769}Manually select a channel.
\nwenddocs{}\nwbegincode{770}\moddef{SF base: sf chain instance: TBP}\plusendmoddef
  procedure :: select_channel => sf_chain_instance_select_channel
\nwendcode{}\nwbegincode{771}\moddef{SF base: procedures}\plusendmoddef
  subroutine sf_chain_instance_select_channel (chain, channel)
    class(sf_chain_instance_t), intent(inout) :: chain
    integer, intent(in), optional :: channel
    if (present (channel)) then
       chain%selected_channel = channel
    else
       chain%selected_channel = 0
    end if
  end subroutine sf_chain_instance_select_channel
  
\eatline
\nwendcode{}\nwbegindocs{772}Copy a channel-mapping object to the structure-function
chain instance.  We assume that assignment is sufficient, i.e., any
non-static components of the {\tt{}channel} object are allocatable und
thus recursively copied.

After the copy, we extract the single-entry mappings and activate them
for the individual structure functions.  If there is a multi-entry
mapping, we obtain the corresponding MC parameter indices and set them
in the copy of the channel object.
\nwenddocs{}\nwbegincode{773}\moddef{SF base: sf chain instance: TBP}\plusendmoddef
  procedure :: set_channel => sf_chain_instance_set_channel
\nwendcode{}\nwbegincode{774}\moddef{SF base: procedures}\plusendmoddef
  subroutine sf_chain_instance_set_channel (chain, c, channel)
    class(sf_chain_instance_t), intent(inout) :: chain
    integer, intent(in) :: c
    type(sf_channel_t), intent(in) :: channel
    integer :: i, j, k
    if (chain%status >= SF_INITIAL) then
       chain%channel(c) = channel
       j = 0
       do i = 1, chain%config%n_strfun
          associate (sf => chain%sf(i))
            sf%m(c) = channel%is_single_mapping (i)
            if (channel%is_multi_mapping (i)) then
               do k = 1, size (sf%int%beam_index)
                  j = j + 1
                  call chain%channel(c)%set_par_index (j, sf%int%par_index(k))
               end do
            end if
          end associate
       end do
       chain%status = SF_INITIAL
    end if
  end subroutine sf_chain_instance_set_channel
  
\eatline
\nwendcode{}\nwbegindocs{775}Link the interactions in the chain.  First, link the beam instance
to its template in the configuration chain, which should have the
appropriate momenta fixed.

Then, we follow the chain via the
arrays {\tt{}out{\char95}sf} and {\tt{}out{\char95}sf{\char95}i}.  The arrays are (up to)
two-dimensional, the entries correspond to the beam particle(s).
For each beam, the entry {\tt{}out{\char95}sf} points to the last interaction
that affected this beam, and {\tt{}out{\char95}sf{\char95}i} is the
out-particle index within that interaction.  For the initial beam,
{\tt{}out{\char95}sf} is zero by definition.

For each entry in the chain, we scan the affected beams (one or two).
We look for {\tt{}out{\char95}sf} and link the out-particle there to the
corresponding in-particle in the current interaction.  Then, we update
the entry in {\tt{}out{\char95}sf} and {\tt{}out{\char95}sf{\char95}i} to point to the current
interaction.
\nwenddocs{}\nwbegincode{776}\moddef{SF base: sf chain instance: TBP}\plusendmoddef
  procedure :: link_interactions => sf_chain_instance_link_interactions
\nwendcode{}\nwbegincode{777}\moddef{SF base: procedures}\plusendmoddef
  subroutine sf_chain_instance_link_interactions (chain)
    class(sf_chain_instance_t), intent(inout), target :: chain
    type(interaction_t), pointer :: int
    integer :: i, j, b
    if (chain%status >= SF_INITIAL) then
       do b = 1, chain%config%n_in
          int => beam_get_int_ptr (chain%beam_t)
          call interaction_set_source_link (int, b, &
               chain%config%beam_t, b)
       end do
       if (allocated (chain%sf)) then
          do i = 1, size (chain%sf)
             associate (sf_int => chain%sf(i)%int)
               do j = 1, size (sf_int%beam_index)
                  b = sf_int%beam_index(j)
                  call link (sf_int%interaction_t, b, sf_int%incoming(j))
                  chain%out_sf(b) = i
                  chain%out_sf_i(b) = sf_int%outgoing(j)
               end do
             end associate
          end do
       end if
       chain%status = SF_DONE_LINKS
    end if
  contains
    subroutine link (int, b, in_index)
      type(interaction_t), intent(inout) :: int
      integer, intent(in) :: b, in_index
      integer :: i
      i = chain%out_sf(b)
      select case (i)
      case (0)
         call interaction_set_source_link (int, in_index, &
              chain%beam_t, chain%out_sf_i(b))
      case default
         call interaction_set_source_link (int, in_index, &
              chain%sf(i)%int%interaction_t, chain%out_sf_i(b))
      end select
    end subroutine link
  end subroutine sf_chain_instance_link_interactions
  
\eatline
\nwendcode{}\nwbegindocs{778}Exchange the quantum-number masks between the interactions in the
chain, so we can combine redundant entries and detect any obvious mismatch.

We proceed first in the forward direction and then backwards again.

After this is finished, we finalize initialization by calling the
{\tt{}setup{\char95}constants} method, which prepares constant data that depend on the
matrix element structure.
\nwenddocs{}\nwbegincode{779}\moddef{SF base: sf chain instance: TBP}\plusendmoddef
  procedure :: exchange_mask => sf_chain_exchange_mask
\nwendcode{}\nwbegincode{780}\moddef{SF base: procedures}\plusendmoddef
  subroutine sf_chain_exchange_mask (chain)
    class(sf_chain_instance_t), intent(inout), target :: chain
    type(interaction_t), pointer :: int
    type(quantum_numbers_mask_t), dimension(:), allocatable :: mask
    integer :: i
    if (chain%status >= SF_DONE_LINKS) then
       if (allocated (chain%sf)) then
          int => beam_get_int_ptr (chain%beam_t)
          allocate (mask (interaction_get_n_out (int)))
          mask = interaction_get_mask (int)
          if (size (chain%sf) /= 0) then
             do i = 1, size (chain%sf) - 1
                call interaction_exchange_mask (chain%sf(i)%int%interaction_t)
             end do
             do i = size (chain%sf), 1, -1
                call interaction_exchange_mask (chain%sf(i)%int%interaction_t)
             end do
             if (any (mask .neqv. interaction_get_mask (int))) then
                chain%status = SF_FAILED_MASK
                return
             end if
             do i = 1, size (chain%sf)
                call chain%sf(i)%int%setup_constants ()
             end do
          end if
       end if
       chain%status = SF_DONE_MASK
    end if
  end subroutine sf_chain_exchange_mask
  
\eatline
\nwendcode{}\nwbegindocs{781}Initialize the evaluators that connect the interactions in the
chain.
\nwenddocs{}\nwbegincode{782}\moddef{SF base: sf chain instance: TBP}\plusendmoddef
  procedure :: init_evaluators => sf_chain_instance_init_evaluators
\nwendcode{}\nwbegincode{783}\moddef{SF base: procedures}\plusendmoddef
  subroutine sf_chain_instance_init_evaluators (chain)
    class(sf_chain_instance_t), intent(inout), target :: chain
    type(interaction_t), pointer :: int
    type(quantum_numbers_mask_t) :: mask
    integer :: i
    if (chain%status >= SF_DONE_MASK) then
       if (allocated (chain%sf)) then
          if (size (chain%sf) /= 0) then
             mask = quantum_numbers_mask (.false., .false., .true.)
             int => beam_get_int_ptr (chain%beam_t)
             do i = 1, size (chain%sf)
                associate (sf => chain%sf(i))
                  call evaluator_init_product (sf%eval, &
                       int, sf%int%interaction_t, &
                       mask)
                  if (evaluator_is_empty (sf%eval)) then
                     chain%status = SF_FAILED_CONNECTIONS
                     return
                  end if
                  int => evaluator_get_int_ptr (sf%eval)
                end associate
             end do
             call find_outgoing_particles ()
          end if
       end if
       chain%status = SF_DONE_CONNECTIONS
    end if
  contains
  \LA{}SF base: init evaluators: find outgoing particles\RA{}
  end subroutine sf_chain_instance_init_evaluators
  
\eatline
\nwendcode{}\nwbegindocs{784}This is an internal subroutine of the previous one: After evaluators
are set, trace the outgoing particles to the last evaluator.  We only
need the first channel, all channels are equivalent for this purpose.

For each beam, the outgoing particle is located by {\tt{}out{\char95}sf} (the
structure-function object where it originates) and {\tt{}out{\char95}sf{\char95}i} (the
index within that object).  This particle is referenced by the
corresponding evaluator, which in turn is referenced by the next
evaluator, until we are at the end of the chain.  We can trace back
references by {\tt{}interaction{\char95}find{\char95}link}.  Knowing that {\tt{}out{\char95}eval} is
the index of the last evaluator, we thus determine {\tt{}out{\char95}eval{\char95}i}, the
index of the outgoing particle within that evaluator.
\nwenddocs{}\nwbegincode{785}\moddef{SF base: init evaluators: find outgoing particles}\endmoddef
  subroutine find_outgoing_particles ()
    type(interaction_t), pointer :: int, int_next
    integer :: i, j, out_sf, out_i
    chain%out_eval = size (chain%sf)
    do j = 1, size (chain%out_eval_i)
       out_sf = chain%out_sf(j)
       out_i = chain%out_sf_i(j)
       if (out_sf == 0) then
          int => beam_get_int_ptr (chain%beam_t)
          out_sf = 1
       else
          int => chain%sf(out_sf)%int%interaction_t
       end if
       do i = out_sf, chain%out_eval
          int_next => evaluator_get_int_ptr (chain%sf(i)%eval)
          out_i = interaction_find_link (int_next, int, out_i)
          int => int_next
       end do
       chain%out_eval_i(j) = out_i
    end do
  end subroutine find_outgoing_particles
\eatline
\nwendcode{}\nwbegindocs{786}Compute the kinematics in the chain instance.  We can assume that
the seed momenta are set in the configuration beams.  Scanning the
chain, we first transfer the incoming momenta.  Then, the use up the MC input
parameter array {\tt{}p} to compute the radiated and outgoing momenta.

In the multi-channel case, {\tt{}c{\char95}sel} is the channel which we use for
computing the kinematics and the {\tt{}x} values.  In the other channels,
we invert the kinematics in order to recover the corresponding rows in
the {\tt{}r} array, and the Jacobian {\tt{}f}.

We first apply any global mapping to transform the input {\tt{}p} into
the array {\tt{}r}.  This is then given to the structure functions which
compute the final array {\tt{}x} and Jacobian factors {\tt{}f}, which we
multiply to obtain the overall Jacobian.
\nwenddocs{}\nwbegincode{787}\moddef{SF base: sf chain instance: TBP}\plusendmoddef
  procedure :: compute_kinematics => sf_chain_instance_compute_kinematics
\nwendcode{}\nwbegincode{788}\moddef{SF base: procedures}\plusendmoddef
  subroutine sf_chain_instance_compute_kinematics (chain, c_sel, p_in)
    class(sf_chain_instance_t), intent(inout), target :: chain
    integer, intent(in) :: c_sel
    real(default), dimension(:), intent(in) :: p_in
    type(interaction_t), pointer :: int
    real(default) :: f_mapping
    logical, dimension(size (chain%bound)) :: bound
    integer :: i, j, c
    if (chain%status >= SF_DONE_CONNECTIONS) then
       call chain%select_channel (c_sel)
       int => beam_get_int_ptr (chain%beam_t)
       call interaction_receive_momenta (int)
       if (allocated (chain%sf)) then
          if (size (chain%sf) /= 0) then
             forall (i = 1:size (chain%sf))  chain%sf(i)%int%status = SF_INITIAL
             !!! Bug in nagfor 5.3.1(907), fixed in 5.3.1(982)
             ! chain%p (:,c_sel) = unpack (p_in, chain%bound, 0._default)
             !!! Workaround:
             bound = chain%bound
             chain%p (:,c_sel) = unpack (p_in, bound, 0._default)
             chain%pb(:,c_sel) = 1 - chain%p(:,c_sel)
             chain%f = 1
             chain%x_free = 1
             do i = 1, size (chain%sf)
                associate (sf => chain%sf(i))
                  call sf%int%generate_free (sf%r(:,c_sel), sf%rb(:,c_sel), &
                       chain%x_free)
                  do j = 1, size (sf%x)
                     if (.not. chain%bound(sf%int%par_index(j))) then
                        chain%p (sf%int%par_index(j),c_sel) = sf%r (j,c_sel)
                        chain%pb(sf%int%par_index(j),c_sel) = sf%rb(j,c_sel)
                     end if
                  end do
                end associate
             end do
             if (allocated (chain%channel(c_sel)%multi_mapping)) then
                call chain%channel(c_sel)%multi_mapping%compute &
                     (chain%r(:,c_sel), chain%rb(:,c_sel), &
                      f_mapping, &
                      chain%p(:,c_sel), chain%pb(:,c_sel), &
                      chain%x_free)
                chain%f(c_sel) = f_mapping
             else
                chain%r (:,c_sel) = chain%p (:,c_sel)
                chain%rb(:,c_sel) = chain%pb(:,c_sel)
                chain%f(c_sel) = 1
             end if
             do i = 1, size (chain%sf)
                associate (sf => chain%sf(i))
                  call sf%int%seed_kinematics ()
                  do j = 1, size (sf%x)
                     sf%r (j,c_sel) = chain%r (sf%int%par_index(j),c_sel)
                     sf%rb(j,c_sel) = chain%rb(sf%int%par_index(j),c_sel)
                  end do
                  call sf%int%complete_kinematics &
                       (sf%x, sf%f(c_sel), sf%r(:,c_sel), sf%rb(:,c_sel), &
                        sf%m(c_sel))
                  do j = 1, size (sf%x)
                     chain%x(sf%int%par_index(j)) = sf%x(j)
                  end do
                  if (sf%int%status <= SF_FAILED_KINEMATICS) then
                     chain%status = SF_FAILED_KINEMATICS
                     return
                  end if
                  do c = 1, size (sf%f)
                     if (c /= c_sel) then
                        call sf%int%inverse_kinematics &
                             (sf%x, sf%f(c), sf%r(:,c), sf%rb(:,c), sf%m(c))
                        do j = 1, size (sf%x)
                           chain%r (sf%int%par_index(j),c) = sf%r (j,c)
                           chain%rb(sf%int%par_index(j),c) = sf%rb(j,c)
                        end do
                     end if
                     chain%f(c) = chain%f(c) * sf%f(c)
                  end do
                  if (.not. evaluator_is_empty (sf%eval)) then
                     call evaluator_receive_momenta (sf%eval)
                  end if
                end associate
             end do
             do c = 1, size (chain%f)
                if (c /= c_sel) then
                   if (allocated (chain%channel(c)%multi_mapping)) then
                      call chain%channel(c)%multi_mapping%inverse &
                           (chain%r(:,c), chain%rb(:,c), &
                            f_mapping, &
                            chain%p(:,c), chain%pb(:,c), &
                            chain%x_free)
                      chain%f(c) = chain%f(c) * f_mapping
                   else
                      chain%p (:,c) = chain%r (:,c)
                      chain%pb(:,c) = chain%rb(:,c)
                   end if
                end if
             end do
          end if
       end if
       chain%status = SF_DONE_KINEMATICS
    end if
  end subroutine sf_chain_instance_compute_kinematics
  
\eatline
\nwendcode{}\nwbegindocs{789}This is a variant of the previous procedure.  We know the $x$ parameters and
reconstruct the momenta and the MC input parameters {\tt{}p}.  We do not need to
select a channel.

Note: this is probably redundant, since the method we actually want
starts from the momenta, recovers all $x$ parameters, and then
inverts mappings.  See below {\tt{}recover{\char95}kinematics}.
\nwenddocs{}\nwbegincode{790}\moddef{SF base: sf chain instance: TBP}\plusendmoddef
  procedure :: inverse_kinematics => sf_chain_instance_inverse_kinematics
\nwendcode{}\nwbegincode{791}\moddef{SF base: procedures}\plusendmoddef
  subroutine sf_chain_instance_inverse_kinematics (chain, x)
    class(sf_chain_instance_t), intent(inout), target :: chain
    real(default), dimension(:), intent(in) :: x
    type(interaction_t), pointer :: int
    real(default) :: f_mapping
    integer :: i, j, c
    if (chain%status >= SF_DONE_CONNECTIONS) then
       call chain%select_channel ()
       int => beam_get_int_ptr (chain%beam_t)
       call interaction_receive_momenta (int)
       if (allocated (chain%sf)) then
          chain%f = 1
          if (size (chain%sf) /= 0) then
             forall (i = 1:size (chain%sf))  chain%sf(i)%int%status = SF_INITIAL
             chain%x = x
             do i = 1, size (chain%sf)
                associate (sf => chain%sf(i))
                  call sf%int%seed_kinematics ()
                  do j = 1, size (sf%x)
                     sf%x(j) = chain%x(sf%int%par_index(j))
                  end do
                  do c = 1, size (sf%f)
                     call sf%int%inverse_kinematics &
                          (sf%x, sf%f(c), sf%r(:,c), sf%rb(:,c), sf%m(c), c==1)
                     chain%f(c) = chain%f(c) * sf%f(c)
                     do j = 1, size (sf%x)
                        chain%r (sf%int%par_index(j),c) = sf%r (j,c)
                        chain%rb(sf%int%par_index(j),c) = sf%rb(j,c)
                     end do
                  end do
                  if (.not. evaluator_is_empty (sf%eval)) then
                     call evaluator_receive_momenta (sf%eval)
                  end if
                end associate
             end do
             do c = 1, size (chain%f)
                if (allocated (chain%channel(c)%multi_mapping)) then
                   call chain%channel(c)%multi_mapping%inverse &
                        (chain%r(:,c), chain%rb(:,c), &
                        f_mapping, &
                        chain%p(:,c), chain%pb(:,c), &
                        chain%x_free)
                   chain%f(c) = chain%f(c) * f_mapping
                else
                   chain%p (:,c) = chain%r (:,c)
                   chain%pb(:,c) = chain%rb(:,c)
                end if
             end do
          end if
       end if
       chain%status = SF_DONE_KINEMATICS
    end if
  end subroutine sf_chain_instance_inverse_kinematics
  
\eatline
\nwendcode{}\nwbegindocs{792}Recover the kinematics: assuming that the last evaluator has
been filled with a valid set of momenta, we travel the momentum links
backwards and fill the preceding evaluators and, as a side effect,
interactions.  We stop at the beam interaction.

After all momenta are set, apply the {\tt{}inverse{\char95}kinematics} procedure
above, suitably modified, to recover the $x$ and $p$ parameters and
the Jacobian factors.

The {\tt{}c{\char95}sel} (channel) argument is just used to mark a selected
channel for the records, otherwise the recovery procedure is
independent of this.
\nwenddocs{}\nwbegincode{793}\moddef{SF base: sf chain instance: TBP}\plusendmoddef
  procedure :: recover_kinematics => sf_chain_instance_recover_kinematics
\nwendcode{}\nwbegincode{794}\moddef{SF base: procedures}\plusendmoddef
  subroutine sf_chain_instance_recover_kinematics (chain, c_sel)
    class(sf_chain_instance_t), intent(inout), target :: chain
    integer, intent(in) :: c_sel
    real(default) :: f_mapping
    integer :: i, j, c
    if (chain%status >= SF_DONE_CONNECTIONS) then
       call chain%select_channel (c_sel)
       if (allocated (chain%sf)) then
          do i = size (chain%sf), 1, -1
             associate (sf => chain%sf(i))
               if (.not. evaluator_is_empty (sf%eval)) then
                  call evaluator_send_momenta (sf%eval)
               end if
             end associate
          end do
          chain%f = 1
          if (size (chain%sf) /= 0) then
             forall (i = 1:size (chain%sf))  chain%sf(i)%int%status = SF_INITIAL
             chain%x_free = 1
             do i = 1, size (chain%sf)
                associate (sf => chain%sf(i))
                  call sf%int%seed_kinematics ()
                  call sf%int%recover_x (sf%x, chain%x_free)
                  do j = 1, size (sf%x)
                     chain%x(sf%int%par_index(j)) = sf%x(j)
                  end do
                  do c = 1, size (sf%f)
                     call sf%int%inverse_kinematics &
                          (sf%x, sf%f(c), sf%r(:,c), sf%rb(:,c), sf%m(c), c==1)
                     chain%f(c) = chain%f(c) * sf%f(c)
                     do j = 1, size (sf%x)
                        chain%r (sf%int%par_index(j),c) = sf%r (j,c)
                        chain%rb(sf%int%par_index(j),c) = sf%rb(j,c)
                     end do
                  end do
                end associate
             end do
             do c = 1, size (chain%f)
                if (allocated (chain%channel(c)%multi_mapping)) then
                   call chain%channel(c)%multi_mapping%inverse &
                        (chain%r(:,c), chain%rb(:,c), &
                        f_mapping, &
                        chain%p(:,c), chain%pb(:,c), &
                        chain%x_free)
                   chain%f(c) = chain%f(c) * f_mapping
                else
                   chain%p (:,c) = chain%r (:,c)
                   chain%pb(:,c) = chain%rb(:,c)
                end if
             end do
          end if
       end if
       chain%status = SF_DONE_KINEMATICS
    end if
  end subroutine sf_chain_instance_recover_kinematics

\eatline
\nwendcode{}\nwbegindocs{795}Return the initial beam momenta to their source, thus completing
kinematics recovery.  Obviously, this works as a side effect.
\nwenddocs{}\nwbegincode{796}\moddef{SF base: sf chain instance: TBP}\plusendmoddef
  procedure :: return_beam_momenta => sf_chain_instance_return_beam_momenta
\nwendcode{}\nwbegincode{797}\moddef{SF base: procedures}\plusendmoddef
  subroutine sf_chain_instance_return_beam_momenta (chain)
    class(sf_chain_instance_t), intent(in), target :: chain
    type(interaction_t), pointer :: int
    if (chain%status >= SF_DONE_KINEMATICS) then
       int => beam_get_int_ptr (chain%beam_t)
       call interaction_send_momenta (int)
    end if
  end subroutine sf_chain_instance_return_beam_momenta

\eatline
\nwendcode{}\nwbegindocs{798}Evaluate all interactions in the chain and the product evaluators.
We provide a {\tt{}scale} argument that is given to all structure
functions in the chain.
\nwenddocs{}\nwbegincode{799}\moddef{SF base: sf chain instance: TBP}\plusendmoddef
  procedure :: evaluate => sf_chain_instance_evaluate
\nwendcode{}\nwbegincode{800}\moddef{SF base: procedures}\plusendmoddef
  subroutine sf_chain_instance_evaluate (chain, scale)
    class(sf_chain_instance_t), intent(inout), target :: chain
    real(default), intent(in) :: scale
    type(interaction_t), pointer :: out_int
    real(default) :: sf_sum
    integer :: i
    if (chain%status >= SF_DONE_KINEMATICS) then
       if (allocated (chain%sf)) then
          if (size (chain%sf) /= 0) then
             do i = 1, size (chain%sf)
                associate (sf => chain%sf(i))
                  call sf%int%apply (scale)
                  if (sf%int%status <= SF_FAILED_EVALUATION) then
                     chain%status = SF_FAILED_EVALUATION
                     return
                  end if
                  if (.not. evaluator_is_empty (sf%eval)) then
                     call sf%eval%evaluate ()
                  end if
                end associate
             end do
             out_int => chain%get_out_int_ptr ()
             sf_sum = interaction_sum (out_int)
             call chain%config%trace &
                  (chain%selected_channel, chain%p, chain%x, chain%f, sf_sum)
          end if
       end if
       chain%status = SF_EVALUATED
    end if
  end subroutine sf_chain_instance_evaluate
  
\eatline
\nwendcode{}\nwbegindocs{801}\nwdocspar
\subsection{Access to the chain instance}
Transfer the outgoing momenta to the array {\tt{}p}.  We assume that
array sizes match.
\nwenddocs{}\nwbegincode{802}\moddef{SF base: sf chain instance: TBP}\plusendmoddef
  procedure :: get_out_momenta => sf_chain_instance_get_out_momenta
\nwendcode{}\nwbegincode{803}\moddef{SF base: procedures}\plusendmoddef
  subroutine sf_chain_instance_get_out_momenta (chain, p)
    class(sf_chain_instance_t), intent(in), target :: chain
    type(vector4_t), dimension(:), intent(out) :: p
    type(interaction_t), pointer :: int
    integer :: i, j
    if (chain%status >= SF_DONE_KINEMATICS) then
       do j = 1, size (chain%out_sf)
          i = chain%out_sf(j)
          select case (i)
          case (0)
             int => beam_get_int_ptr (chain%beam_t)
          case default
             int => chain%sf(i)%int%interaction_t
          end select
          p(j) = interaction_get_momentum (int, chain%out_sf_i(j))
       end do
    end if
  end subroutine sf_chain_instance_get_out_momenta
       
\eatline
\nwendcode{}\nwbegindocs{804}Return a pointer to the last evaluator in the chain (to the interaction).
\nwenddocs{}\nwbegincode{805}\moddef{SF base: sf chain instance: TBP}\plusendmoddef
  procedure :: get_out_int_ptr => sf_chain_instance_get_out_int_ptr
\nwendcode{}\nwbegincode{806}\moddef{SF base: procedures}\plusendmoddef
  function sf_chain_instance_get_out_int_ptr (chain) result (int)
    class(sf_chain_instance_t), intent(in), target :: chain
    type(interaction_t), pointer :: int
    if (chain%out_eval == 0) then
       int => beam_get_int_ptr (chain%beam_t)
    else
       int => evaluator_get_int_ptr (chain%sf(chain%out_eval)%eval)
    end if
  end function sf_chain_instance_get_out_int_ptr

\eatline
\nwendcode{}\nwbegindocs{807}Return the index of the {\tt{}j}-th outgoing particle, within the last
evaluator.
\nwenddocs{}\nwbegincode{808}\moddef{SF base: sf chain instance: TBP}\plusendmoddef
  procedure :: get_out_i => sf_chain_instance_get_out_i
\nwendcode{}\nwbegincode{809}\moddef{SF base: procedures}\plusendmoddef
  function sf_chain_instance_get_out_i (chain, j) result (i)
    class(sf_chain_instance_t), intent(in) :: chain
    integer, intent(in) :: j
    integer :: i
    i = chain%out_eval_i(j)
  end function sf_chain_instance_get_out_i
    
\eatline
\nwendcode{}\nwbegindocs{810}Return the mask for the outgoing particle(s), within the last evaluator.
\nwenddocs{}\nwbegincode{811}\moddef{SF base: sf chain instance: TBP}\plusendmoddef
  procedure :: get_out_mask => sf_chain_instance_get_out_mask
\nwendcode{}\nwbegincode{812}\moddef{SF base: procedures}\plusendmoddef
  function sf_chain_instance_get_out_mask (chain) result (mask)
    class(sf_chain_instance_t), intent(in), target :: chain
    type(quantum_numbers_mask_t), dimension(:), allocatable :: mask
    type(interaction_t), pointer :: int
    allocate (mask (chain%config%n_in))
    int => chain%get_out_int_ptr ()
    mask = interaction_get_mask (int, chain%out_eval_i)
  end function sf_chain_instance_get_out_mask
    
\eatline
\nwendcode{}\nwbegindocs{813}Return the array of MC input parameters that corresponds to channel {\tt{}c}.
This is the {\tt{}p} array, the parameters before all mappings.

The {\tt{}p} array may be deallocated.  This should correspond to a
zero-size {\tt{}r} argument, so nothing to do then.
\nwenddocs{}\nwbegincode{814}\moddef{SF base: sf chain instance: TBP}\plusendmoddef
  procedure :: get_mcpar => sf_chain_instance_get_mcpar
\nwendcode{}\nwbegincode{815}\moddef{SF base: procedures}\plusendmoddef
  subroutine sf_chain_instance_get_mcpar (chain, c, r)
    class(sf_chain_instance_t), intent(in) :: chain
    integer, intent(in) :: c
    real(default), dimension(:), intent(out) :: r
    if (allocated (chain%p))  r = pack (chain%p(:,c), chain%bound)
  end subroutine sf_chain_instance_get_mcpar
  
\eatline
\nwendcode{}\nwbegindocs{816}Return the Jacobian factor that corresponds to channel {\tt{}c}.
\nwenddocs{}\nwbegincode{817}\moddef{SF base: sf chain instance: TBP}\plusendmoddef
  procedure :: get_f => sf_chain_instance_get_f
\nwendcode{}\nwbegincode{818}\moddef{SF base: procedures}\plusendmoddef
  function sf_chain_instance_get_f (chain, c) result (f)
    class(sf_chain_instance_t), intent(in) :: chain
    integer, intent(in) :: c
    real(default) :: f
    if (allocated (chain%f)) then
       f = chain%f(c)
    else
       f = 1
    end if
  end function sf_chain_instance_get_f
  
\eatline
\nwendcode{}\nwbegindocs{819}Return the evaluation status.
\nwenddocs{}\nwbegincode{820}\moddef{SF base: sf chain instance: TBP}\plusendmoddef
  procedure :: get_status => sf_chain_instance_get_status
\nwendcode{}\nwbegincode{821}\moddef{SF base: procedures}\plusendmoddef
  function sf_chain_instance_get_status (chain) result (status)
    class(sf_chain_instance_t), intent(in) :: chain
    integer :: status
    status = chain%status
  end function sf_chain_instance_get_status
  
\eatline
\nwendcode{}\nwbegindocs{822}\nwdocspar
\subsection{Unit tests}
\nwenddocs{}\nwbegincode{823}\moddef{SF base: public}\plusendmoddef
  public :: sf_base_test
\nwendcode{}\nwbegincode{824}\moddef{SF base: tests}\endmoddef
  subroutine sf_base_test (u, results)
    integer, intent(in) :: u
    type(test_results_t), intent(inout) :: results
  \LA{}SF base: execute tests\RA{}
  end subroutine sf_base_test
  
\eatline
\nwendcode{}\nwbegindocs{825}\nwdocspar
\subsection{Test implementation: structure function}
This is a template for the actual structure-function implementation
which will be defined in separate modules.

\subsubsection{Configuration data}
The test structure function uses the {\tt{}Test} model.  It describes a
scalar within an arbitrary initial particle, which is given in the
initialization.  The radiated particle is also a scalar, the same one,
but we set its mass artificially to zero.
\nwenddocs{}\nwbegincode{826}\moddef{SF base: public}\plusendmoddef
  public :: sf_test_data_t
\nwendcode{}\nwbegincode{827}\moddef{SF base: test types}\endmoddef
  type, extends (sf_data_t) :: sf_test_data_t
     class(model_data_t), pointer :: model => null ()
     integer :: mode = 0
     type(flavor_t) :: flv_in
     type(flavor_t) :: flv_out
     type(flavor_t) :: flv_rad
     real(default) :: m = 0
     logical :: collinear = .true.
     real(default), dimension(:), allocatable :: qbounds
   contains
   \LA{}SF base: sf test data: TBP\RA{}
  end type sf_test_data_t
  
\eatline
\nwendcode{}\nwbegindocs{828}Output.
\nwenddocs{}\nwbegincode{829}\moddef{SF base: sf test data: TBP}\endmoddef
  procedure :: write => sf_test_data_write
\nwendcode{}\nwbegincode{830}\moddef{SF base: procedures}\plusendmoddef
  subroutine sf_test_data_write (data, unit, verbose)
    class(sf_test_data_t), intent(in) :: data
    integer, intent(in), optional :: unit
    logical, intent(in), optional :: verbose
    integer :: u
    u = given_output_unit (unit)
    write (u, "(1x,A)")  "SF test data:"
    write (u, "(3x,A,A)") "model     = ", char (data%model%get_name ())
    write (u, "(3x,A)", advance="no") "incoming  = "
    call data%flv_in%write (u);  write (u, *)
    write (u, "(3x,A)", advance="no") "outgoing  = "
    call data%flv_out%write (u);  write (u, *)
    write (u, "(3x,A)", advance="no") "radiated  = "
    call data%flv_rad%write (u);  write (u, *)
    write (u, "(3x,A," // FMT_19 // ")")  "mass      = ", data%m
    write (u, "(3x,A,L1)")  "collinear = ", data%collinear
    if (.not. data%collinear .and. allocated (data%qbounds)) then
       write (u, "(3x,A," // FMT_19 // ")")  "qmin      = ", data%qbounds(1)
       write (u, "(3x,A," // FMT_19 // ")")  "qmax      = ", data%qbounds(2)
    end if
  end subroutine sf_test_data_write
    
\eatline
\nwendcode{}\nwbegindocs{831}Initialization.
\nwenddocs{}\nwbegincode{832}\moddef{SF base: sf test data: TBP}\plusendmoddef
  procedure :: init => sf_test_data_init
\nwendcode{}\nwbegincode{833}\moddef{SF base: procedures}\plusendmoddef
  subroutine sf_test_data_init (data, model, pdg_in, collinear, qbounds, mode)
    class(sf_test_data_t), intent(out) :: data
    class(model_data_t), intent(in), target :: model
    type(pdg_array_t), intent(in) :: pdg_in
    logical, intent(in), optional :: collinear
    real(default), dimension(2), intent(in), optional :: qbounds
    integer, intent(in), optional :: mode
    data%model => model
    if (present (mode))  data%mode = mode
    if (pdg_array_get (pdg_in, 1) /= 25) then
       call msg_fatal ("Test spectrum function: input flavor must be 's'")
    end if
    call data%flv_in%init (25, model)
    data%m = data%flv_in%get_mass ()
    if (present (collinear))  data%collinear = collinear
    call data%flv_out%init (25, model)
    call data%flv_rad%init (25, model)
    if (present (qbounds)) then
       allocate (data%qbounds (2))
       data%qbounds = qbounds
    end if
  end subroutine sf_test_data_init
  
\eatline
\nwendcode{}\nwbegindocs{834}Return the number of parameters: 1 if only consider collinear
splitting, 3 otherwise.
\nwenddocs{}\nwbegincode{835}\moddef{SF base: sf test data: TBP}\plusendmoddef
  procedure :: get_n_par => sf_test_data_get_n_par
\nwendcode{}\nwbegincode{836}\moddef{SF base: procedures}\plusendmoddef
  function sf_test_data_get_n_par (data) result (n)
    class(sf_test_data_t), intent(in) :: data
    integer :: n
    if (data%collinear) then
       n = 1
    else
       n = 3
    end if
  end function sf_test_data_get_n_par
  
\eatline
\nwendcode{}\nwbegindocs{837}Return the outgoing particle PDG code: 25
\nwenddocs{}\nwbegincode{838}\moddef{SF base: sf test data: TBP}\plusendmoddef
  procedure :: get_pdg_out => sf_test_data_get_pdg_out
\nwendcode{}\nwbegincode{839}\moddef{SF base: procedures}\plusendmoddef
  subroutine sf_test_data_get_pdg_out (data, pdg_out)
    class(sf_test_data_t), intent(in) :: data
    type(pdg_array_t), dimension(:), intent(inout) :: pdg_out
    pdg_out(1) = 25
  end subroutine sf_test_data_get_pdg_out
  
\eatline
\nwendcode{}\nwbegindocs{840}Allocate the matching interaction.
\nwenddocs{}\nwbegincode{841}\moddef{SF base: sf test data: TBP}\plusendmoddef
  procedure :: allocate_sf_int => sf_test_data_allocate_sf_int
\nwendcode{}\nwbegincode{842}\moddef{SF base: procedures}\plusendmoddef
  subroutine sf_test_data_allocate_sf_int (data, sf_int)
    class(sf_test_data_t), intent(in) :: data
    class(sf_int_t), intent(inout), allocatable :: sf_int
    allocate (sf_test_t :: sf_int)
  end subroutine sf_test_data_allocate_sf_int
    
\eatline
\nwendcode{}\nwbegindocs{843}\nwdocspar
\subsubsection{Interaction}
\nwenddocs{}\nwbegincode{844}\moddef{SF base: test types}\plusendmoddef
  type, extends (sf_int_t) :: sf_test_t
     type(sf_test_data_t), pointer :: data => null ()
     real(default) :: x = 0
   contains
   \LA{}SF base: sf test int: TBP\RA{}
  end type sf_test_t
  
\eatline
\nwendcode{}\nwbegindocs{845}Type string: constant
\nwenddocs{}\nwbegincode{846}\moddef{SF base: sf test int: TBP}\endmoddef
  procedure :: type_string => sf_test_type_string
\nwendcode{}\nwbegincode{847}\moddef{SF base: procedures}\plusendmoddef
  function sf_test_type_string (object) result (string)
    class(sf_test_t), intent(in) :: object
    type(string_t) :: string
    string = "Test"
  end function sf_test_type_string
  
\eatline
\nwendcode{}\nwbegindocs{848}Output.  Call the interaction routine after displaying the configuration.
\nwenddocs{}\nwbegincode{849}\moddef{SF base: sf test int: TBP}\plusendmoddef
  procedure :: write => sf_test_write
\nwendcode{}\nwbegincode{850}\moddef{SF base: procedures}\plusendmoddef
  subroutine sf_test_write (object, unit, testflag)
    class(sf_test_t), intent(in) :: object
    integer, intent(in), optional :: unit
    logical, intent(in), optional :: testflag
    integer :: u
    u = given_output_unit (unit)
    if (associated (object%data)) then
       call object%data%write (u)
       call object%base_write (u, testflag)
    else
       write (u, "(1x,A)")  "SF test data: [undefined]"
    end if
  end subroutine sf_test_write
    
\eatline
\nwendcode{}\nwbegindocs{851}Initialize.  We know that {\tt{}data} will be of concrete type
{\tt{}sf{\char95}test{\char95}data{\char95}t}, but we have to cast this explicitly.

For this implementation, we set the incoming and outgoing masses equal
to the physical particle mass, but keep the radiated mass zero.

Optionally, we can provide minimum and maximum values for the momentum
transfer.
\nwenddocs{}\nwbegincode{852}\moddef{SF base: sf test int: TBP}\plusendmoddef
  procedure :: init => sf_test_init
\nwendcode{}\nwbegincode{853}\moddef{SF base: procedures}\plusendmoddef
  subroutine sf_test_init (sf_int, data)
    class(sf_test_t), intent(out) :: sf_int
    class(sf_data_t), intent(in), target :: data
    type(quantum_numbers_mask_t), dimension(3) :: mask
    type(helicity_t) :: hel0
    type(color_t) :: col0
    type(quantum_numbers_t), dimension(3) :: qn
    mask = quantum_numbers_mask (.false., .false., .false.)
    select type (data)
    type is (sf_test_data_t)
       if (allocated (data%qbounds)) then
          call sf_int%base_init (mask, &
               [data%m**2], [0._default], [data%m**2], &
               [data%qbounds(1)], [data%qbounds(2)])
       else
          call sf_int%base_init (mask, &
               [data%m**2], [0._default], [data%m**2])
       end if
       sf_int%data => data
       call hel0%init (0)
       call col0%init ()
       call qn(1)%init (data%flv_in,  col0, hel0)
       call qn(2)%init (data%flv_rad, col0, hel0)
       call qn(3)%init (data%flv_out, col0, hel0)
       call interaction_add_state (sf_int%interaction_t, qn)
       call interaction_freeze (sf_int%interaction_t)
       call sf_int%set_incoming ([1])
       call sf_int%set_radiated ([2])
       call sf_int%set_outgoing ([3])
    end select
    sf_int%status = SF_INITIAL
  end subroutine sf_test_init

\eatline
\nwendcode{}\nwbegindocs{854}Set kinematics.  If {\tt{}map} is unset, the $r$ and $x$ values
coincide, and the Jacobian $f(r)$ is trivial.

If {\tt{}map} is set, we are asked to provide an efficient mapping.
For the test case, we set $x=r^2$ and consequently $f(r)=2r$.
\nwenddocs{}\nwbegincode{855}\moddef{SF base: sf test int: TBP}\plusendmoddef
  procedure :: complete_kinematics => sf_test_complete_kinematics
\nwendcode{}\nwbegincode{856}\moddef{SF base: procedures}\plusendmoddef
  subroutine sf_test_complete_kinematics (sf_int, x, f, r, rb, map)
    class(sf_test_t), intent(inout) :: sf_int
    real(default), dimension(:), intent(out) :: x
    real(default), intent(out) :: f
    real(default), dimension(:), intent(in) :: r
    real(default), dimension(:), intent(in) :: rb
    logical, intent(in) :: map
    real(default) :: xb1
    if (map) then
       x(1) = r(1)**2
       f = 2 * r(1)
    else
       x(1) = r(1)
       f = 1
    end if
    xb1 = 1 - x(1)
    if (size (x) == 3)  x(2:3) = r(2:3)
    call sf_int%split_momentum (x, xb1)
    sf_int%x = x(1)
    select case (sf_int%status)
    case (SF_FAILED_KINEMATICS);  f = 0
    end select
  end subroutine sf_test_complete_kinematics

\eatline
\nwendcode{}\nwbegindocs{857}Compute inverse kinematics.  Here, we start with the $x$ array and
compute the ``input'' $r$ values and the Jacobian $f$.  After this, we
can set momenta by the same formula as for normal kinematics.
\nwenddocs{}\nwbegincode{858}\moddef{SF base: sf test int: TBP}\plusendmoddef
  procedure :: inverse_kinematics => sf_test_inverse_kinematics
\nwendcode{}\nwbegincode{859}\moddef{SF base: procedures}\plusendmoddef
  subroutine sf_test_inverse_kinematics (sf_int, x, f, r, rb, map, set_momenta)
    class(sf_test_t), intent(inout) :: sf_int
    real(default), dimension(:), intent(in) :: x
    real(default), intent(out) :: f
    real(default), dimension(:), intent(out) :: r
    real(default), dimension(:), intent(out) :: rb
    logical, intent(in) :: map
    logical, intent(in), optional :: set_momenta
    real(default) :: xb1
    logical :: set_mom
    set_mom = .false.;  if (present (set_momenta))  set_mom = set_momenta
    if (map) then
       r(1) = sqrt (x(1))
       f = 2 * r(1)
    else
       r(1) = x(1)
       f = 1
    end if
    xb1 = 1 - x(1)
    if (size (x) == 3)  r(2:3) = x(2:3)
    rb = 1 - r
    sf_int%x = x(1)
    if (set_mom) then
       call sf_int%split_momentum (x, xb1)
       select case (sf_int%status)
       case (SF_FAILED_KINEMATICS);  f = 0
       end select
    end if
  end subroutine sf_test_inverse_kinematics

\eatline
\nwendcode{}\nwbegindocs{860}Apply the structure function.  The matrix element becomes unity and
the application always succeeds.

If the {\tt{}mode} indicator is one, the matrix element is equal to the
parameter~$x$.
\nwenddocs{}\nwbegincode{861}\moddef{SF base: sf test int: TBP}\plusendmoddef
  procedure :: apply => sf_test_apply
\nwendcode{}\nwbegincode{862}\moddef{SF base: procedures}\plusendmoddef
  subroutine sf_test_apply (sf_int, scale)
    class(sf_test_t), intent(inout) :: sf_int
    real(default), intent(in) :: scale
    select case (sf_int%data%mode)
    case (0)
       call interaction_set_matrix_element (sf_int%interaction_t, &
            cmplx (1._default, kind=default))
    case (1)
       call interaction_set_matrix_element (sf_int%interaction_t, &
            cmplx (sf_int%x, kind=default))
    end select
    sf_int%status = SF_EVALUATED
  end subroutine sf_test_apply

\eatline
\nwendcode{}\nwbegindocs{863}\nwdocspar
\subsection{Test implementation: pair spectrum}
Another template, this time for a incoming particle pair, splitting
into two radiated and two outgoing particles.

\subsubsection{Configuration data}
For simplicity, the spectrum contains two mirror images of the
previous structure-function configuration: the incoming and all
outgoing particles are test scalars.

We have two versions, one with radiated particles, one without.
\nwenddocs{}\nwbegincode{864}\moddef{SF base: test types}\plusendmoddef
  type, extends (sf_data_t) :: sf_test_spectrum_data_t
     class(model_data_t), pointer :: model => null ()
     type(flavor_t) :: flv_in
     type(flavor_t) :: flv_out
     type(flavor_t) :: flv_rad
     logical :: with_radiation = .true.
     real(default) :: m = 0
   contains
   \LA{}SF base: sf test spectrum data: TBP\RA{}
  end type sf_test_spectrum_data_t
  
\eatline
\nwendcode{}\nwbegindocs{865}Output.
\nwenddocs{}\nwbegincode{866}\moddef{SF base: sf test spectrum data: TBP}\endmoddef
  procedure :: write => sf_test_spectrum_data_write
\nwendcode{}\nwbegincode{867}\moddef{SF base: procedures}\plusendmoddef
  subroutine sf_test_spectrum_data_write (data, unit, verbose)
    class(sf_test_spectrum_data_t), intent(in) :: data
    integer, intent(in), optional :: unit
    logical, intent(in), optional :: verbose
    integer :: u
    u = given_output_unit (unit)
    write (u, "(1x,A)")  "SF test spectrum data:"
    write (u, "(3x,A,A)") "model     = ", char (data%model%get_name ())
    write (u, "(3x,A)", advance="no") "incoming  = "
    call data%flv_in%write (u);  write (u, *)
    write (u, "(3x,A)", advance="no") "outgoing  = "
    call data%flv_out%write (u);  write (u, *)
    write (u, "(3x,A)", advance="no") "radiated  = "
    call data%flv_rad%write (u);  write (u, *)
    write (u, "(3x,A," // FMT_19 // ")")  "mass      = ", data%m
  end subroutine sf_test_spectrum_data_write
    
\eatline
\nwendcode{}\nwbegindocs{868}Initialization.
\nwenddocs{}\nwbegincode{869}\moddef{SF base: sf test spectrum data: TBP}\plusendmoddef
  procedure :: init => sf_test_spectrum_data_init
\nwendcode{}\nwbegincode{870}\moddef{SF base: procedures}\plusendmoddef
  subroutine sf_test_spectrum_data_init (data, model, pdg_in, with_radiation)
    class(sf_test_spectrum_data_t), intent(out) :: data
    class(model_data_t), intent(in), target :: model
    type(pdg_array_t), intent(in) :: pdg_in
    logical, intent(in) :: with_radiation
    data%model => model
    data%with_radiation = with_radiation
    if (pdg_array_get (pdg_in, 1) /= 25) then
       call msg_fatal ("Test structure function: input flavor must be 's'")
    end if
    call data%flv_in%init (25, model)
    data%m = data%flv_in%get_mass ()
    call data%flv_out%init (25, model)
    if (with_radiation) then
       call data%flv_rad%init (25, model)
    end if
  end subroutine sf_test_spectrum_data_init
  
\eatline
\nwendcode{}\nwbegindocs{871}Return the number of parameters: 2, since we have only collinear
splitting here.
\nwenddocs{}\nwbegincode{872}\moddef{SF base: sf test spectrum data: TBP}\plusendmoddef
  procedure :: get_n_par => sf_test_spectrum_data_get_n_par
\nwendcode{}\nwbegincode{873}\moddef{SF base: procedures}\plusendmoddef
  function sf_test_spectrum_data_get_n_par (data) result (n)
    class(sf_test_spectrum_data_t), intent(in) :: data
    integer :: n
    n = 2
  end function sf_test_spectrum_data_get_n_par
  
\eatline
\nwendcode{}\nwbegindocs{874}Return the outgoing particle PDG codes: 25
\nwenddocs{}\nwbegincode{875}\moddef{SF base: sf test spectrum data: TBP}\plusendmoddef
  procedure :: get_pdg_out => sf_test_spectrum_data_get_pdg_out
\nwendcode{}\nwbegincode{876}\moddef{SF base: procedures}\plusendmoddef
  subroutine sf_test_spectrum_data_get_pdg_out (data, pdg_out)
    class(sf_test_spectrum_data_t), intent(in) :: data
    type(pdg_array_t), dimension(:), intent(inout) :: pdg_out
    pdg_out(1) = 25
    pdg_out(2) = 25
  end subroutine sf_test_spectrum_data_get_pdg_out
  
\eatline
\nwendcode{}\nwbegindocs{877}Allocate the matching interaction.
\nwenddocs{}\nwbegincode{878}\moddef{SF base: sf test spectrum data: TBP}\plusendmoddef
  procedure :: allocate_sf_int => &
       sf_test_spectrum_data_allocate_sf_int
\nwendcode{}\nwbegincode{879}\moddef{SF base: procedures}\plusendmoddef
  subroutine sf_test_spectrum_data_allocate_sf_int (data, sf_int)
    class(sf_test_spectrum_data_t), intent(in) :: data
    class(sf_int_t), intent(inout), allocatable :: sf_int
    allocate (sf_test_spectrum_t :: sf_int)
  end subroutine sf_test_spectrum_data_allocate_sf_int
    
\eatline
\nwendcode{}\nwbegindocs{880}\nwdocspar
\subsubsection{Interaction}
\nwenddocs{}\nwbegincode{881}\moddef{SF base: test types}\plusendmoddef
  type, extends (sf_int_t) :: sf_test_spectrum_t
     type(sf_test_spectrum_data_t), pointer :: data => null ()
   contains
   \LA{}SF base: sf test spectrum: TBP\RA{}
  end type sf_test_spectrum_t
  
\eatline
\nwendcode{}\nwbegincode{882}\moddef{SF base: sf test spectrum: TBP}\endmoddef
  procedure :: type_string => sf_test_spectrum_type_string
\nwendcode{}\nwbegincode{883}\moddef{SF base: procedures}\plusendmoddef
  function sf_test_spectrum_type_string (object) result (string)
    class(sf_test_spectrum_t), intent(in) :: object
    type(string_t) :: string
    string = "Test Spectrum"
  end function sf_test_spectrum_type_string
  
\eatline
\nwendcode{}\nwbegindocs{884}Output.  Call the interaction routine after displaying the configuration.
\nwenddocs{}\nwbegincode{885}\moddef{SF base: sf test spectrum: TBP}\plusendmoddef
  procedure :: write => sf_test_spectrum_write
\nwendcode{}\nwbegincode{886}\moddef{SF base: procedures}\plusendmoddef
  subroutine sf_test_spectrum_write (object, unit, testflag)
    class(sf_test_spectrum_t), intent(in) :: object
    integer, intent(in), optional :: unit
    logical, intent(in), optional :: testflag
    integer :: u
    u = given_output_unit (unit)
    if (associated (object%data)) then
       call object%data%write (u)
       call object%base_write (u, testflag)
    else
       write (u, "(1x,A)")  "SF test spectrum data: [undefined]"
    end if
  end subroutine sf_test_spectrum_write
    
\eatline
\nwendcode{}\nwbegindocs{887}Initialize.  We know that {\tt{}data} will be of concrete type
{\tt{}sf{\char95}test{\char95}spectrum{\char95}data{\char95}t}, but we have to cast this explicitly.

For this implementation, we set the incoming and outgoing masses equal
to the physical particle mass, but keep the radiated mass zero.

Optionally, we can provide minimum and maximum values for the momentum
transfer.
\nwenddocs{}\nwbegincode{888}\moddef{SF base: sf test spectrum: TBP}\plusendmoddef
  procedure :: init => sf_test_spectrum_init
\nwendcode{}\nwbegincode{889}\moddef{SF base: procedures}\plusendmoddef
  subroutine sf_test_spectrum_init (sf_int, data)
    class(sf_test_spectrum_t), intent(out) :: sf_int
    class(sf_data_t), intent(in), target :: data
    type(quantum_numbers_mask_t), dimension(6) :: mask
    type(helicity_t) :: hel0
    type(color_t) :: col0
    type(quantum_numbers_t), dimension(6) :: qn
    mask = quantum_numbers_mask (.false., .false., .false.)
    select type (data)
    type is (sf_test_spectrum_data_t)
       if (data%with_radiation) then
          call sf_int%base_init (mask(1:6), &
               [data%m**2, data%m**2], &
               [0._default, 0._default], &
               [data%m**2, data%m**2])
          sf_int%data => data
          call hel0%init (0)
          call col0%init ()
          call qn(1)%init (data%flv_in,  col0, hel0)
          call qn(2)%init (data%flv_in,  col0, hel0)
          call qn(3)%init (data%flv_rad, col0, hel0)
          call qn(4)%init (data%flv_rad, col0, hel0)
          call qn(5)%init (data%flv_out, col0, hel0)
          call qn(6)%init (data%flv_out, col0, hel0)
          call interaction_add_state (sf_int%interaction_t, qn(1:6))
          call sf_int%set_incoming ([1,2])
          call sf_int%set_radiated ([3,4])
          call sf_int%set_outgoing ([5,6])
       else
          call sf_int%base_init (mask(1:4), &
               [data%m**2, data%m**2], &
               [real(default) :: ], &
               [data%m**2, data%m**2])
          sf_int%data => data
          call hel0%init (0)
          call col0%init ()
          call qn(1)%init (data%flv_in,  col0, hel0)
          call qn(2)%init (data%flv_in,  col0, hel0)
          call qn(3)%init (data%flv_out, col0, hel0)
          call qn(4)%init (data%flv_out, col0, hel0)
          call interaction_add_state (sf_int%interaction_t, qn(1:4))
          call sf_int%set_incoming ([1,2])
          call sf_int%set_outgoing ([3,4])
       end if
       call interaction_freeze (sf_int%interaction_t)
    end select
    sf_int%status = SF_INITIAL
  end subroutine sf_test_spectrum_init

\eatline
\nwendcode{}\nwbegindocs{890}Set kinematics.  If {\tt{}map} is unset, the $r$ and $x$ values
coincide, and the Jacobian $f(r)$ is trivial.

If {\tt{}map} is set, we are asked to provide an efficient mapping.
For the test case, we set $x=r^2$ (as above) for both $x$ parameters
and consequently $f(r)=4r_1r_2$.

\nwenddocs{}\nwbegincode{891}\moddef{SF base: sf test spectrum: TBP}\plusendmoddef
  procedure :: complete_kinematics => sf_test_spectrum_complete_kinematics
\nwendcode{}\nwbegincode{892}\moddef{SF base: procedures}\plusendmoddef
  subroutine sf_test_spectrum_complete_kinematics (sf_int, x, f, r, rb, map)
    class(sf_test_spectrum_t), intent(inout) :: sf_int
    real(default), dimension(:), intent(out) :: x
    real(default), intent(out) :: f
    real(default), dimension(:), intent(in) :: r
    real(default), dimension(:), intent(in) :: rb
    logical, intent(in) :: map
    real(default), dimension(2) :: xb1
    if (map) then
       x = r**2
       f = 4 * r(1) * r(2)
    else
       x = r
       f = 1
    end if
    if (sf_int%data%with_radiation) then
       xb1 = 1 - x
       call sf_int%split_momenta (x, xb1)
    else
       call sf_int%reduce_momenta (x)
    end if
    select case (sf_int%status)
    case (SF_FAILED_KINEMATICS);  f = 0
    end select
  end subroutine sf_test_spectrum_complete_kinematics

\eatline
\nwendcode{}\nwbegindocs{893}Compute inverse kinematics.  Here, we start with the $x$ array and
compute the ``input'' $r$ values and the Jacobian $f$.  After this, we
can set momenta by the same formula as for normal kinematics.
\nwenddocs{}\nwbegincode{894}\moddef{SF base: sf test spectrum: TBP}\plusendmoddef
  procedure :: inverse_kinematics => sf_test_spectrum_inverse_kinematics
\nwendcode{}\nwbegincode{895}\moddef{SF base: procedures}\plusendmoddef
  subroutine sf_test_spectrum_inverse_kinematics &
       (sf_int, x, f, r, rb, map, set_momenta)
    class(sf_test_spectrum_t), intent(inout) :: sf_int
    real(default), dimension(:), intent(in) :: x
    real(default), intent(out) :: f
    real(default), dimension(:), intent(out) :: r
    real(default), dimension(:), intent(out) :: rb
    logical, intent(in) :: map
    logical, intent(in), optional :: set_momenta
    real(default), dimension(2) :: xb1
    logical :: set_mom
    set_mom = .false.;  if (present (set_momenta))  set_mom = set_momenta
    if (map) then
       r = sqrt (x)
       f = 4 * r(1) * r(2)
    else
       r = x
       f = 1
    end if
    rb = 1 - r
    if (set_mom)  then
       if (sf_int%data%with_radiation) then
          xb1 = 1 - x
          call sf_int%split_momenta (x, xb1)
       else
          call sf_int%reduce_momenta (x)
       end if
       select case (sf_int%status)
       case (SF_FAILED_KINEMATICS);  f = 0
       end select
    end if
  end subroutine sf_test_spectrum_inverse_kinematics

\eatline
\nwendcode{}\nwbegindocs{896}Apply the structure function.  The matrix element becomes unity and
the application always succeeds.
\nwenddocs{}\nwbegincode{897}\moddef{SF base: sf test spectrum: TBP}\plusendmoddef
  procedure :: apply => sf_test_spectrum_apply
\nwendcode{}\nwbegincode{898}\moddef{SF base: procedures}\plusendmoddef
  subroutine sf_test_spectrum_apply (sf_int, scale)
    class(sf_test_spectrum_t), intent(inout) :: sf_int
    real(default), intent(in) :: scale
    call interaction_set_matrix_element (sf_int%interaction_t, &
         cmplx (1._default, kind=default))
    sf_int%status = SF_EVALUATED
  end subroutine sf_test_spectrum_apply

\eatline
\nwendcode{}\nwbegindocs{899}\nwdocspar
\subsection{Test implementation: generator spectrum}
A generator for two beams, no radiation (for simplicity).

\subsubsection{Configuration data}
For simplicity, the spectrum contains two mirror images of the
previous structure-function configuration: the incoming and all
outgoing particles are test scalars.

We have two versions, one with radiated particles, one without.
\nwenddocs{}\nwbegincode{900}\moddef{SF base: test types}\plusendmoddef
  type, extends (sf_data_t) :: sf_test_generator_data_t
     class(model_data_t), pointer :: model => null ()
     type(flavor_t) :: flv_in
     type(flavor_t) :: flv_out
     type(flavor_t) :: flv_rad
     real(default) :: m = 0
   contains
   \LA{}SF base: sf test generator data: TBP\RA{}
  end type sf_test_generator_data_t
  
\eatline
\nwendcode{}\nwbegindocs{901}Output.
\nwenddocs{}\nwbegincode{902}\moddef{SF base: sf test generator data: TBP}\endmoddef
  procedure :: write => sf_test_generator_data_write
\nwendcode{}\nwbegincode{903}\moddef{SF base: procedures}\plusendmoddef
  subroutine sf_test_generator_data_write (data, unit, verbose)
    class(sf_test_generator_data_t), intent(in) :: data
    integer, intent(in), optional :: unit
    logical, intent(in), optional :: verbose
    integer :: u
    u = given_output_unit (unit)
    write (u, "(1x,A)")  "SF test generator data:"
    write (u, "(3x,A,A)") "model     = ", char (data%model%get_name ())
    write (u, "(3x,A)", advance="no") "incoming  = "
    call data%flv_in%write (u);  write (u, *)
    write (u, "(3x,A)", advance="no") "outgoing  = "
    call data%flv_out%write (u);  write (u, *)
    write (u, "(3x,A," // FMT_19 // ")")  "mass      = ", data%m
  end subroutine sf_test_generator_data_write
    
\eatline
\nwendcode{}\nwbegindocs{904}Initialization.
\nwenddocs{}\nwbegincode{905}\moddef{SF base: sf test generator data: TBP}\plusendmoddef
  procedure :: init => sf_test_generator_data_init
\nwendcode{}\nwbegincode{906}\moddef{SF base: procedures}\plusendmoddef
  subroutine sf_test_generator_data_init (data, model, pdg_in)
    class(sf_test_generator_data_t), intent(out) :: data
    class(model_data_t), intent(in), target :: model
    type(pdg_array_t), intent(in) :: pdg_in
    data%model => model
    if (pdg_array_get (pdg_in, 1) /= 25) then
       call msg_fatal ("Test generator: input flavor must be 's'")
    end if
    call data%flv_in%init (25, model)
    data%m = data%flv_in%get_mass ()
    call data%flv_out%init (25, model)
  end subroutine sf_test_generator_data_init
  
\eatline
\nwendcode{}\nwbegindocs{907}This structure function is a generator.
\nwenddocs{}\nwbegincode{908}\moddef{SF base: sf test generator data: TBP}\plusendmoddef
  procedure :: is_generator => sf_test_generator_data_is_generator
\nwendcode{}\nwbegincode{909}\moddef{SF base: procedures}\plusendmoddef
  function sf_test_generator_data_is_generator (data) result (flag)
    class(sf_test_generator_data_t), intent(in) :: data
    logical :: flag
    flag = .true.
  end function sf_test_generator_data_is_generator
  
\eatline
\nwendcode{}\nwbegindocs{910}Return the number of parameters: 2, since we have only collinear
splitting here.
\nwenddocs{}\nwbegincode{911}\moddef{SF base: sf test generator data: TBP}\plusendmoddef
  procedure :: get_n_par => sf_test_generator_data_get_n_par
\nwendcode{}\nwbegincode{912}\moddef{SF base: procedures}\plusendmoddef
  function sf_test_generator_data_get_n_par (data) result (n)
    class(sf_test_generator_data_t), intent(in) :: data
    integer :: n
    n = 2
  end function sf_test_generator_data_get_n_par
  
\eatline
\nwendcode{}\nwbegindocs{913}Return the outgoing particle PDG codes: 25
\nwenddocs{}\nwbegincode{914}\moddef{SF base: sf test generator data: TBP}\plusendmoddef
  procedure :: get_pdg_out => sf_test_generator_data_get_pdg_out
\nwendcode{}\nwbegincode{915}\moddef{SF base: procedures}\plusendmoddef
  subroutine sf_test_generator_data_get_pdg_out (data, pdg_out)
    class(sf_test_generator_data_t), intent(in) :: data
    type(pdg_array_t), dimension(:), intent(inout) :: pdg_out
    pdg_out(1) = 25
    pdg_out(2) = 25
  end subroutine sf_test_generator_data_get_pdg_out
  
\eatline
\nwendcode{}\nwbegindocs{916}Allocate the matching interaction.
\nwenddocs{}\nwbegincode{917}\moddef{SF base: sf test generator data: TBP}\plusendmoddef
  procedure :: allocate_sf_int => &
       sf_test_generator_data_allocate_sf_int
\nwendcode{}\nwbegincode{918}\moddef{SF base: procedures}\plusendmoddef
  subroutine sf_test_generator_data_allocate_sf_int (data, sf_int)
    class(sf_test_generator_data_t), intent(in) :: data
    class(sf_int_t), intent(inout), allocatable :: sf_int
    allocate (sf_test_generator_t :: sf_int)
  end subroutine sf_test_generator_data_allocate_sf_int
    
\eatline
\nwendcode{}\nwbegindocs{919}\nwdocspar
\subsubsection{Interaction}
\nwenddocs{}\nwbegincode{920}\moddef{SF base: test types}\plusendmoddef
  type, extends (sf_int_t) :: sf_test_generator_t
     type(sf_test_generator_data_t), pointer :: data => null ()
   contains
   \LA{}SF base: sf test generator: TBP\RA{}
  end type sf_test_generator_t
  
\eatline
\nwendcode{}\nwbegincode{921}\moddef{SF base: sf test generator: TBP}\endmoddef
  procedure :: type_string => sf_test_generator_type_string
\nwendcode{}\nwbegincode{922}\moddef{SF base: procedures}\plusendmoddef
  function sf_test_generator_type_string (object) result (string)
    class(sf_test_generator_t), intent(in) :: object
    type(string_t) :: string
    string = "Test Generator"
  end function sf_test_generator_type_string
  
\eatline
\nwendcode{}\nwbegindocs{923}Output.  Call the interaction routine after displaying the configuration.
\nwenddocs{}\nwbegincode{924}\moddef{SF base: sf test generator: TBP}\plusendmoddef
  procedure :: write => sf_test_generator_write
\nwendcode{}\nwbegincode{925}\moddef{SF base: procedures}\plusendmoddef
  subroutine sf_test_generator_write (object, unit, testflag)
    class(sf_test_generator_t), intent(in) :: object
    integer, intent(in), optional :: unit
    logical, intent(in), optional :: testflag
    integer :: u
    u = given_output_unit (unit)
    if (associated (object%data)) then
       call object%data%write (u)
       call object%base_write (u, testflag)
    else
       write (u, "(1x,A)")  "SF test generator data: [undefined]"
    end if
  end subroutine sf_test_generator_write
    
\eatline
\nwendcode{}\nwbegindocs{926}Initialize.  We know that {\tt{}data} will be of concrete type
{\tt{}sf{\char95}test{\char95}generator{\char95}data{\char95}t}, but we have to cast this explicitly.

For this implementation, we set the incoming and outgoing masses equal
to the physical particle mass.  No radiation.
\nwenddocs{}\nwbegincode{927}\moddef{SF base: sf test generator: TBP}\plusendmoddef
  procedure :: init => sf_test_generator_init
\nwendcode{}\nwbegincode{928}\moddef{SF base: procedures}\plusendmoddef
  subroutine sf_test_generator_init (sf_int, data)
    class(sf_test_generator_t), intent(out) :: sf_int
    class(sf_data_t), intent(in), target :: data
    type(quantum_numbers_mask_t), dimension(4) :: mask
    type(helicity_t) :: hel0
    type(color_t) :: col0
    type(quantum_numbers_t), dimension(4) :: qn
    mask = quantum_numbers_mask (.false., .false., .false.)
    select type (data)
    type is (sf_test_generator_data_t)
       call sf_int%base_init (mask(1:4), &
            [data%m**2, data%m**2], &
            [real(default) :: ], &
            [data%m**2, data%m**2])
       sf_int%data => data
       call hel0%init (0)
       call col0%init ()
       call qn(1)%init (data%flv_in,  col0, hel0)
       call qn(2)%init (data%flv_in,  col0, hel0)
       call qn(3)%init (data%flv_out, col0, hel0)
       call qn(4)%init (data%flv_out, col0, hel0)
       call interaction_add_state (sf_int%interaction_t, qn(1:4))
       call sf_int%set_incoming ([1,2])
       call sf_int%set_outgoing ([3,4])
       call interaction_freeze (sf_int%interaction_t)
    end select
    sf_int%status = SF_INITIAL
  end subroutine sf_test_generator_init

\eatline
\nwendcode{}\nwbegindocs{929}This structure function is a generator.
\nwenddocs{}\nwbegincode{930}\moddef{SF base: sf test generator: TBP}\plusendmoddef
  procedure :: is_generator => sf_test_generator_is_generator
\nwendcode{}\nwbegincode{931}\moddef{SF base: procedures}\plusendmoddef
  function sf_test_generator_is_generator (sf_int) result (flag)
    class(sf_test_generator_t), intent(in) :: sf_int
    logical :: flag
    flag = sf_int%data%is_generator ()
  end function sf_test_generator_is_generator
  
\eatline
\nwendcode{}\nwbegindocs{932}Generate free parameters.  This mock generator always produces the
nubmers 0.8 and 0.5.
\nwenddocs{}\nwbegincode{933}\moddef{SF base: sf test generator: TBP}\plusendmoddef
  procedure :: generate_free => sf_test_generator_generate_free
\nwendcode{}\nwbegincode{934}\moddef{SF base: procedures}\plusendmoddef
  subroutine sf_test_generator_generate_free (sf_int, r, rb,  x_free)
    class(sf_test_generator_t), intent(inout) :: sf_int
    real(default), dimension(:), intent(out) :: r, rb
    real(default), intent(inout) :: x_free
    r = [0.8, 0.5]
    rb= 1 - r
    x_free = x_free * product (r)
  end subroutine sf_test_generator_generate_free
    
\eatline
\nwendcode{}\nwbegindocs{935}Recover momentum fractions.  Since the x values are free, we also set the {\tt{}x{\char95}free} parameter.
\nwenddocs{}\nwbegincode{936}\moddef{SF base: sf test generator: TBP}\plusendmoddef
  procedure :: recover_x => sf_test_generator_recover_x
\nwendcode{}\nwbegincode{937}\moddef{SF base: procedures}\plusendmoddef
  subroutine sf_test_generator_recover_x (sf_int, x, x_free)
    class(sf_test_generator_t), intent(inout) :: sf_int
    real(default), dimension(:), intent(out) :: x
    real(default), intent(inout), optional :: x_free
    call sf_int%base_recover_x (x)
    if (present (x_free))  x_free = x_free * product (x)
  end subroutine sf_test_generator_recover_x
  
\eatline
\nwendcode{}\nwbegindocs{938}Set kinematics.  Since this is a generator, just transfer input to output.
\nwenddocs{}\nwbegincode{939}\moddef{SF base: sf test generator: TBP}\plusendmoddef
  procedure :: complete_kinematics => sf_test_generator_complete_kinematics
\nwendcode{}\nwbegincode{940}\moddef{SF base: procedures}\plusendmoddef
  subroutine sf_test_generator_complete_kinematics (sf_int, x, f, r, rb, map)
    class(sf_test_generator_t), intent(inout) :: sf_int
    real(default), dimension(:), intent(out) :: x
    real(default), intent(out) :: f
    real(default), dimension(:), intent(in) :: r
    real(default), dimension(:), intent(in) :: rb
    logical, intent(in) :: map
    x = r
    f = 1
    call sf_int%reduce_momenta (x)
  end subroutine sf_test_generator_complete_kinematics

\eatline
\nwendcode{}\nwbegindocs{941}Compute inverse kinematics.  Here, we start with the $x$ array and
compute the ``input'' $r$ values and the Jacobian $f$.  After this, we
can set momenta by the same formula as for normal kinematics.
\nwenddocs{}\nwbegincode{942}\moddef{SF base: sf test generator: TBP}\plusendmoddef
  procedure :: inverse_kinematics => sf_test_generator_inverse_kinematics
\nwendcode{}\nwbegincode{943}\moddef{SF base: procedures}\plusendmoddef
  subroutine sf_test_generator_inverse_kinematics &
       (sf_int, x, f, r, rb, map, set_momenta)
    class(sf_test_generator_t), intent(inout) :: sf_int
    real(default), dimension(:), intent(in) :: x
    real(default), intent(out) :: f
    real(default), dimension(:), intent(out) :: r
    real(default), dimension(:), intent(out) :: rb
    logical, intent(in) :: map
    logical, intent(in), optional :: set_momenta
    logical :: set_mom
    set_mom = .false.;  if (present (set_momenta))  set_mom = set_momenta
    r = x
    rb= 1 - x
    f = 1
    if (set_mom)  call sf_int%reduce_momenta (x)
  end subroutine sf_test_generator_inverse_kinematics

\eatline
\nwendcode{}\nwbegindocs{944}Apply the structure function.  The matrix element becomes unity and
the application always succeeds.
\nwenddocs{}\nwbegincode{945}\moddef{SF base: sf test generator: TBP}\plusendmoddef
  procedure :: apply => sf_test_generator_apply
\nwendcode{}\nwbegincode{946}\moddef{SF base: procedures}\plusendmoddef
  subroutine sf_test_generator_apply (sf_int, scale)
    class(sf_test_generator_t), intent(inout) :: sf_int
    real(default), intent(in) :: scale
    call interaction_set_matrix_element (sf_int%interaction_t, &
         cmplx (1._default, kind=default))
    sf_int%status = SF_EVALUATED
  end subroutine sf_test_generator_apply

\eatline
\nwendcode{}\nwbegindocs{947}\nwdocspar
\subsubsection{Test structure function data}
Construct and display a test structure function data object.
\nwenddocs{}\nwbegincode{948}\moddef{SF base: execute tests}\endmoddef
  call test (sf_base_1, "sf_base_1", &
       "structure function configuration", &
       u, results)
\nwendcode{}\nwbegincode{949}\moddef{SF base: tests}\plusendmoddef
  subroutine sf_base_1 (u)
    integer, intent(in) :: u
    type(model_data_t), target :: model
    type(pdg_array_t) :: pdg_in
    type(pdg_array_t), dimension(1) :: pdg_out
    integer, dimension(:), allocatable :: pdg1
    class(sf_data_t), allocatable :: data
    
    write (u, "(A)")  "* Test output: sf_base_1"
    write (u, "(A)")  "*   Purpose: initialize and display &
         &test structure function data"
    write (u, "(A)")
    
    call model%init_test ()
    pdg_in = 25

    allocate (sf_test_data_t :: data)
    select type (data)
    type is (sf_test_data_t)
       call data%init (model, pdg_in)
    end select
       
    call data%write (u)

    write (u, "(A)") 

    write (u, "(1x,A)")  "Outgoing particle code:"
    call data%get_pdg_out (pdg_out)
    pdg1 = pdg_out(1)
    write (u, "(2x,99(1x,I0))")  pdg1
    
    call model%final ()
    
    write (u, "(A)")
    write (u, "(A)")  "* Test output end: sf_base_1"

  end subroutine sf_base_1

\eatline
\nwendcode{}\nwbegindocs{950}\nwdocspar
\subsubsection{Test and probe structure function}
Construct and display a structure function object based on the test
structure function.
\nwenddocs{}\nwbegincode{951}\moddef{SF base: execute tests}\plusendmoddef
  call test (sf_base_2, "sf_base_2", &
       "structure function instance", &
       u, results)
\nwendcode{}\nwbegincode{952}\moddef{SF base: tests}\plusendmoddef
  subroutine sf_base_2 (u)
    integer, intent(in) :: u
    type(model_data_t), target :: model
    type(flavor_t) :: flv
    type(pdg_array_t) :: pdg_in
    class(sf_data_t), allocatable, target :: data
    class(sf_int_t), allocatable :: sf_int
    type(vector4_t) :: k
    type(vector4_t), dimension(2) :: q
    real(default) :: E
    real(default), dimension(:), allocatable :: r, rb, x
    real(default) :: f
    
    write (u, "(A)")  "* Test output: sf_base_2"
    write (u, "(A)")  "*   Purpose: initialize and fill &
         &test structure function object"
    write (u, "(A)")
    
    write (u, "(A)")  "* Initialize configuration data"
    write (u, "(A)")

    call model%init_test ()
    pdg_in = 25
    call flv%init (25, model)

    call reset_interaction_counter ()
    
    allocate (sf_test_data_t :: data)
    select type (data)
    type is (sf_test_data_t)
       call data%init (model, pdg_in)
    end select
       
    write (u, "(A)")  "* Initialize structure-function object"
    write (u, "(A)")
    
    call data%allocate_sf_int (sf_int)
    call sf_int%init (data)
    call sf_int%set_beam_index ([1])
    
    call sf_int%write (u)

    write (u, "(A)")
    write (u, "(A)")  "* Initialize incoming momentum with E=500"
    write (u, "(A)")
    E = 500
    k = vector4_moving (E, sqrt (E**2 - flv%get_mass ()**2), 3)
    call vector4_write (k, u)
    call sf_int%seed_kinematics ([k])

    write (u, "(A)")
    write (u, "(A)")  "* Set kinematics for x=0"
    write (u, "(A)")

    allocate (r (data%get_n_par ()))
    allocate (rb(size (r)))
    allocate (x (size (r)))

    r = 0
    rb = 1 - r
    call sf_int%complete_kinematics (x, f, r, rb, map=.false.)
    call sf_int%write (u)

    write (u, "(A)")
    write (u, "(A,9(1x,F10.7))")  "x =", x
    write (u, "(A,9(1x,F10.7))")  "f =", f

    write (u, "(A)")
    write (u, "(A)")  "* Set kinematics for x=1"
    write (u, "(A)")

    r = 1
    rb = 1 - r
    call sf_int%complete_kinematics (x, f, r, rb, map=.false.)
    call sf_int%write (u)

    write (u, "(A)")
    write (u, "(A,9(1x,F10.7))")  "x =", x
    write (u, "(A,9(1x,F10.7))")  "f =", f

    write (u, "(A)")
    write (u, "(A)")  "* Set kinematics for x=0.5"
    write (u, "(A)")

    r = 0.5_default
    rb = 1 - r
    call sf_int%complete_kinematics (x, f, r, rb, map=.false.)
    call sf_int%write (u)

    write (u, "(A)")
    write (u, "(A,9(1x,F10.7))")  "x =", x
    write (u, "(A,9(1x,F10.7))")  "f =", f

    write (u, "(A)")
    write (u, "(A)")  "* Set kinematics with mapping for r=0.8"
    write (u, "(A)")

    r = 0.8_default
    rb = 1 - r
    call sf_int%complete_kinematics (x, f, r, rb, map=.true.)
    call sf_int%write (u)

    write (u, "(A)")
    write (u, "(A,9(1x,F10.7))")  "x =", x
    write (u, "(A,9(1x,F10.7))")  "f =", f

    write (u, "(A)")
    write (u, "(A)")  "* Recover x from momenta"
    write (u, "(A)")

    q = interaction_get_momenta (sf_int%interaction_t, outgoing=.true.)
    call sf_int%final ()
    deallocate (sf_int)

    call data%allocate_sf_int (sf_int)
    call sf_int%init (data)
    call sf_int%set_beam_index ([1])

    call sf_int%seed_kinematics ([k])
    call interaction_set_momenta (sf_int%interaction_t, q, outgoing=.true.)
    call sf_int%recover_x (x)

    write (u, "(A,9(1x,F10.7))")  "x =", x

    write (u, "(A)")
    write (u, "(A)")  "* Compute inverse kinematics for x=0.64 and evaluate"
    write (u, "(A)")

    x = 0.64_default
    call sf_int%inverse_kinematics (x, f, r, rb, map=.true.)
    call sf_int%apply (scale=0._default)
    call sf_int%write (u)

    write (u, "(A)")
    write (u, "(A,9(1x,F10.7))")  "r =", r
    write (u, "(A,9(1x,F10.7))")  "f =", f

    write (u, "(A)")
    write (u, "(A)")  "* Cleanup"

    call sf_int%final ()
    call model%final ()
    
    write (u, "(A)")
    write (u, "(A)")  "* Test output end: sf_base_2"

  end subroutine sf_base_2

\eatline
\nwendcode{}\nwbegindocs{953}\nwdocspar
\subsubsection{Collinear kinematics}
Scan over the possibilities for mass assignment and on-shell
projections, collinear case.
\nwenddocs{}\nwbegincode{954}\moddef{SF base: execute tests}\plusendmoddef
  call test (sf_base_3, "sf_base_3", &
       "alternatives for collinear kinematics", &
       u, results)
\nwendcode{}\nwbegincode{955}\moddef{SF base: tests}\plusendmoddef
  subroutine sf_base_3 (u)
    integer, intent(in) :: u
    type(model_data_t), target :: model
    type(pdg_array_t) :: pdg_in
    type(flavor_t) :: flv
    class(sf_data_t), allocatable, target :: data
    class(sf_int_t), allocatable :: sf_int
    type(vector4_t) :: k
    real(default) :: E
    real(default), dimension(:), allocatable :: r, rb, x
    real(default) :: f
    
    write (u, "(A)")  "* Test output: sf_base_3"
    write (u, "(A)")  "*   Purpose: check various kinematical setups"
    write (u, "(A)")  "*            for collinear structure-function splitting."
    write (u, "(A)")  "             (two masses equal, one zero)"
    write (u, "(A)")
    
    write (u, "(A)")  "* Initialize configuration data"
    write (u, "(A)")

    call model%init_test ()
    pdg_in = 25
    call flv%init (25, model)

    call reset_interaction_counter ()
    
    allocate (sf_test_data_t :: data)
    select type (data)
    type is (sf_test_data_t)
       call data%init (model, pdg_in)
    end select
       
    write (u, "(A)")  "* Initialize structure-function object"
    write (u, "(A)")
    
    call data%allocate_sf_int (sf_int)
    call sf_int%init (data)
    
    call sf_int%write (u)

    allocate (r (data%get_n_par ()))
    allocate (rb(size (r)))
    allocate (x (size (r)))

    write (u, "(A)")
    write (u, "(A)")  "* Initialize incoming momentum with E=500"

    E = 500
    k = vector4_moving (E, sqrt (E**2 - flv%get_mass ()**2), 3)
    call sf_int%seed_kinematics ([k])

    write (u, "(A)")
    write (u, "(A)")  "* Set radiated mass to zero"

    sf_int%mr2 = 0
    sf_int%mo2 = sf_int%mi2
    
    write (u, "(A)")
    write (u, "(A)")  "* Set kinematics for x=0.5, keeping energy"
    write (u, "(A)")

    r = 0.5_default
    rb = 1 - r
    sf_int%on_shell_mode = KEEP_ENERGY
    call sf_int%complete_kinematics (x, f, r, rb, map=.false.)
    call sf_int%write (u)

    write (u, "(A)")
    write (u, "(A)")  "* Recover x and r"
    write (u, "(A)")

    call sf_int%recover_x (x)
    call sf_int%inverse_kinematics (x, f, r, rb, map=.false.)
    write (u, "(A,9(1x,F10.7))")  "x =", x
    write (u, "(A,9(1x,F10.7))")  "r =", r
    
    write (u, "(A)")
    write (u, "(A)")  "* Set kinematics for x=0.5, keeping momentum"
    write (u, "(A)")

    r = 0.5_default
    rb = 1 - r
    sf_int%on_shell_mode = KEEP_MOMENTUM
    call sf_int%complete_kinematics (x, f, r, rb, map=.false.)
    call sf_int%write (u)

    write (u, "(A)")
    write (u, "(A)")  "* Recover x and r"
    write (u, "(A)")

    call sf_int%recover_x (x)
    call sf_int%inverse_kinematics (x, f, r, rb, map=.false.)
    write (u, "(A,9(1x,F10.7))")  "x =", x
    write (u, "(A,9(1x,F10.7))")  "r =", r
    
    write (u, "(A)")
    write (u, "(A)")  "* Set outgoing mass to zero"

    sf_int%mr2 = sf_int%mi2
    sf_int%mo2 = 0
    
    write (u, "(A)")
    write (u, "(A)")  "* Set kinematics for x=0.5, keeping energy"
    write (u, "(A)")

    r = 0.5_default
    rb = 1 - r
    sf_int%on_shell_mode = KEEP_ENERGY
    call sf_int%complete_kinematics (x, f, r, rb, map=.false.)
    call sf_int%write (u)

    write (u, "(A)")
    write (u, "(A)")  "* Recover x and r"
    write (u, "(A)")

    call sf_int%recover_x (x)
    call sf_int%inverse_kinematics (x, f, r, rb, map=.false.)
    write (u, "(A,9(1x,F10.7))")  "x =", x
    write (u, "(A,9(1x,F10.7))")  "r =", r
    
    write (u, "(A)")
    write (u, "(A)")  "* Set kinematics for x=0.5, keeping momentum"
    write (u, "(A)")

    r = 0.5_default
    rb = 1 - r
    sf_int%on_shell_mode = KEEP_MOMENTUM
    call sf_int%complete_kinematics (x, f, r, rb, map=.false.)
    call sf_int%write (u)

    write (u, "(A)")
    write (u, "(A)")  "* Recover x and r"
    write (u, "(A)")

    call sf_int%recover_x (x)
    call sf_int%inverse_kinematics (x, f, r, rb, map=.false.)
    write (u, "(A,9(1x,F10.7))")  "x =", x
    write (u, "(A,9(1x,F10.7))")  "r =", r
    
    write (u, "(A)")
    write (u, "(A)")  "* Set incoming mass to zero"

    k = vector4_moving (E, E, 3)
    call sf_int%seed_kinematics ([k])

    sf_int%mr2 = sf_int%mi2
    sf_int%mo2 = sf_int%mi2
    sf_int%mi2 = 0
    
    write (u, "(A)")
    write (u, "(A)")  "* Set kinematics for x=0.5, keeping energy"
    write (u, "(A)")

    r = 0.5_default
    rb = 1 - r
    sf_int%on_shell_mode = KEEP_ENERGY
    call sf_int%complete_kinematics (x, f, r, rb, map=.false.)
    call sf_int%write (u)

    write (u, "(A)")
    write (u, "(A)")  "* Recover x and r"
    write (u, "(A)")

    call sf_int%recover_x (x)
    call sf_int%inverse_kinematics (x, f, r, rb, map=.false.)
    write (u, "(A,9(1x,F10.7))")  "x =", x
    write (u, "(A,9(1x,F10.7))")  "r =", r
    
    write (u, "(A)")
    write (u, "(A)")  "* Set kinematics for x=0.5, keeping momentum"
    write (u, "(A)")

    r = 0.5_default
    rb = 1 - r
    sf_int%on_shell_mode = KEEP_MOMENTUM
    call sf_int%complete_kinematics (x, f, r, rb, map=.false.)
    call sf_int%write (u)

    write (u, "(A)")
    write (u, "(A)")  "* Recover x and r"
    write (u, "(A)")

    call sf_int%recover_x (x)
    call sf_int%inverse_kinematics (x, f, r, rb, map=.false.)
    write (u, "(A,9(1x,F10.7))")  "x =", x
    write (u, "(A,9(1x,F10.7))")  "r =", r
    
    write (u, "(A)")
    write (u, "(A)")  "* Set all masses to zero"

    sf_int%mr2 = 0
    sf_int%mo2 = 0
    sf_int%mi2 = 0
    
    write (u, "(A)")
    write (u, "(A)")  "* Set kinematics for x=0.5, keeping energy"
    write (u, "(A)")

    r = 0.5_default
    rb = 1 - r
    sf_int%on_shell_mode = KEEP_ENERGY
    call sf_int%complete_kinematics (x, f, r, rb, map=.false.)
    call sf_int%write (u)

    write (u, "(A)")
    write (u, "(A)")  "* Recover x and r"
    write (u, "(A)")

    call sf_int%recover_x (x)
    call sf_int%inverse_kinematics (x, f, r, rb, map=.false.)
    write (u, "(A,9(1x,F10.7))")  "x =", x
    write (u, "(A,9(1x,F10.7))")  "r =", r
    
    write (u, "(A)")
    write (u, "(A)")  "* Set kinematics for x=0.5, keeping momentum"
    write (u, "(A)")

    r = 0.5_default
    rb = 1 - r
    sf_int%on_shell_mode = KEEP_MOMENTUM
    call sf_int%complete_kinematics (x, f, r, rb, map=.false.)
    call sf_int%write (u)

    write (u, "(A)")
    write (u, "(A)")  "* Recover x and r"
    write (u, "(A)")

    call sf_int%recover_x (x)
    call sf_int%inverse_kinematics (x, f, r, rb, map=.false.)
    write (u, "(A,9(1x,F10.7))")  "x =", x
    write (u, "(A,9(1x,F10.7))")  "r =", r
    
    write (u, "(A)")
    write (u, "(A)")  "* Cleanup"

    call sf_int%final ()
    call model%final ()
    
    write (u, "(A)")
    write (u, "(A)")  "* Test output end: sf_base_3"

  end subroutine sf_base_3

\eatline
\nwendcode{}\nwbegindocs{956}\nwdocspar
\subsubsection{Non-collinear kinematics}
Scan over the possibilities for mass assignment and on-shell
projections, non-collinear case.
\nwenddocs{}\nwbegincode{957}\moddef{SF base: execute tests}\plusendmoddef
  call test (sf_base_4, "sf_base_4", &
       "alternatives for non-collinear kinematics", &
       u, results)
\nwendcode{}\nwbegincode{958}\moddef{SF base: tests}\plusendmoddef
  subroutine sf_base_4 (u)
    integer, intent(in) :: u
    type(model_data_t), target :: model
    type(pdg_array_t) :: pdg_in
    type(flavor_t) :: flv
    class(sf_data_t), allocatable, target :: data
    class(sf_int_t), allocatable :: sf_int
    type(vector4_t) :: k
    real(default) :: E
    real(default), dimension(:), allocatable :: r, rb, x
    real(default) :: f
    
    write (u, "(A)")  "* Test output: sf_base_4"
    write (u, "(A)")  "*   Purpose: check various kinematical setups"
    write (u, "(A)")  "*            for free structure-function splitting."
    write (u, "(A)")  "             (two masses equal, one zero)"
    write (u, "(A)")
    
    write (u, "(A)")  "* Initialize configuration data"
    write (u, "(A)")

    call model%init_test ()
    pdg_in = 25
    call flv%init (25, model)

    call reset_interaction_counter ()
    
    allocate (sf_test_data_t :: data)
    select type (data)
    type is (sf_test_data_t)
       call data%init (model, pdg_in, collinear=.false.)
    end select
       
    write (u, "(A)")  "* Initialize structure-function object"
    write (u, "(A)")
    
    call data%allocate_sf_int (sf_int)
    call sf_int%init (data)
    
    call sf_int%write (u)

    allocate (r (data%get_n_par ()))
    allocate (rb(size (r)))
    allocate (x (size (r)))

    write (u, "(A)")
    write (u, "(A)")  "* Initialize incoming momentum with E=500"

    E = 500
    k = vector4_moving (E, sqrt (E**2 - flv%get_mass ()**2), 3)
    call sf_int%seed_kinematics ([k])

    write (u, "(A)")
    write (u, "(A)")  "* Set radiated mass to zero"

    sf_int%mr2 = 0
    sf_int%mo2 = sf_int%mi2
    
    write (u, "(A)")
    write (u, "(A)")  "* Set kinematics for x=0.5/0.5/0.125, keeping energy"
    write (u, "(A)")

    r = [0.5_default, 0.5_default, 0.125_default]
    rb = 1 - r
    sf_int%on_shell_mode = KEEP_ENERGY
    call sf_int%complete_kinematics (x, f, r, rb, map=.false.)
    call sf_int%write (u)

    write (u, "(A)")
    write (u, "(A)")  "* Recover x and r"
    write (u, "(A)")

    call sf_int%recover_x (x)
    call sf_int%inverse_kinematics (x, f, r, rb, map=.false.)
    write (u, "(A,9(1x,F10.7))")  "x =", x
    write (u, "(A,9(1x,F10.7))")  "r =", r
    
    write (u, "(A)")
    write (u, "(A)")  "* Set kinematics for x=0.5/0.5/0.125, keeping momentum"
    write (u, "(A)")

    r = [0.5_default, 0.5_default, 0.125_default]
    rb = 1 - r
    sf_int%on_shell_mode = KEEP_MOMENTUM
    call sf_int%complete_kinematics (x, f, r, rb, map=.false.)
    call sf_int%write (u)

    write (u, "(A)")
    write (u, "(A)")  "* Recover x and r"
    write (u, "(A)")

    call sf_int%recover_x (x)
    call sf_int%inverse_kinematics (x, f, r, rb, map=.false.)
    write (u, "(A,9(1x,F10.7))")  "x =", x
    write (u, "(A,9(1x,F10.7))")  "r =", r
    
    write (u, "(A)")
    write (u, "(A)")  "* Set outgoing mass to zero"

    sf_int%mr2 = sf_int%mi2
    sf_int%mo2 = 0
    
    write (u, "(A)")
    write (u, "(A)")  "* Set kinematics for x=0.5/0.5/0.125, keeping energy"
    write (u, "(A)")

    r = [0.5_default, 0.5_default, 0.125_default]
    rb = 1 - r
    sf_int%on_shell_mode = KEEP_ENERGY
    call sf_int%complete_kinematics (x, f, r, rb, map=.false.)
    call sf_int%write (u)

    write (u, "(A)")
    write (u, "(A)")  "* Recover x and r"
    write (u, "(A)")

    call sf_int%recover_x (x)
    call sf_int%inverse_kinematics (x, f, r, rb, map=.false.)
    write (u, "(A,9(1x,F10.7))")  "x =", x
    write (u, "(A,9(1x,F10.7))")  "r =", r
    
    write (u, "(A)")
    write (u, "(A)")  "* Set kinematics for x=0.5/0.5/0.125, keeping momentum"
    write (u, "(A)")

    r = [0.5_default, 0.5_default, 0.125_default]
    rb = 1 - r
    sf_int%on_shell_mode = KEEP_MOMENTUM
    call sf_int%complete_kinematics (x, f, r, rb, map=.false.)
    call sf_int%write (u)

    write (u, "(A)")
    write (u, "(A)")  "* Recover x and r"
    write (u, "(A)")

    call sf_int%recover_x (x)
    call sf_int%inverse_kinematics (x, f, r, rb, map=.false.)
    write (u, "(A,9(1x,F10.7))")  "x =", x
    write (u, "(A,9(1x,F10.7))")  "r =", r
    
    write (u, "(A)")
    write (u, "(A)")  "* Set incoming mass to zero"

    k = vector4_moving (E, E, 3)
    call sf_int%seed_kinematics ([k])

    sf_int%mr2 = sf_int%mi2
    sf_int%mo2 = sf_int%mi2
    sf_int%mi2 = 0
    
    write (u, "(A)")
    write (u, "(A)")  "* Set kinematics for x=0.5/0.5/0.125, keeping energy"
    write (u, "(A)")

    r = [0.5_default, 0.5_default, 0.125_default]
    rb = 1 - r
    sf_int%on_shell_mode = KEEP_ENERGY
    call sf_int%complete_kinematics (x, f, r, rb, map=.false.)
    call sf_int%write (u)

    write (u, "(A)")
    write (u, "(A)")  "* Recover x and r"
    write (u, "(A)")

    call sf_int%recover_x (x)
    call sf_int%inverse_kinematics (x, f, r, rb, map=.false.)
    write (u, "(A,9(1x,F10.7))")  "x =", x
    write (u, "(A,9(1x,F10.7))")  "r =", r
    
    write (u, "(A)")
    write (u, "(A)")  "* Set kinematics for x=0.5/0.5/0.125, keeping momentum"
    write (u, "(A)")

    r = [0.5_default, 0.5_default, 0.125_default]
    rb = 1 - r
    sf_int%on_shell_mode = KEEP_MOMENTUM
    call sf_int%complete_kinematics (x, f, r, rb, map=.false.)
    call sf_int%write (u)

    write (u, "(A)")
    write (u, "(A)")  "* Recover x and r"
    write (u, "(A)")

    call sf_int%recover_x (x)
    call sf_int%inverse_kinematics (x, f, r, rb, map=.false.)
    write (u, "(A,9(1x,F10.7))")  "x =", x
    write (u, "(A,9(1x,F10.7))")  "r =", r
    
    write (u, "(A)")
    write (u, "(A)")  "* Set all masses to zero"

    sf_int%mr2 = 0
    sf_int%mo2 = 0
    sf_int%mi2 = 0
    
    write (u, "(A)")
    write (u, "(A)")  "* Re-Initialize structure-function object with Q bounds"
    
    call reset_interaction_counter ()
    
    select type (data)
    type is (sf_test_data_t)
       call data%init (model, pdg_in, collinear=.false., &
            qbounds = [1._default, 100._default])
    end select
       
    call sf_int%init (data)
    call sf_int%seed_kinematics ([k])

    write (u, "(A)")
    write (u, "(A)")  "* Set kinematics for x=0.5/0.5/0.125, keeping energy"
    write (u, "(A)")

    r = [0.5_default, 0.5_default, 0.125_default]
    rb = 1 - r
    sf_int%on_shell_mode = KEEP_ENERGY
    call sf_int%complete_kinematics (x, f, r, rb, map=.false.)
    call sf_int%write (u)

    write (u, "(A)")
    write (u, "(A)")  "* Recover x and r"
    write (u, "(A)")

    call sf_int%recover_x (x)
    call sf_int%inverse_kinematics (x, f, r, rb, map=.false.)
    write (u, "(A,9(1x,F10.7))")  "x =", x
    write (u, "(A,9(1x,F10.7))")  "r =", r
    
    write (u, "(A)")
    write (u, "(A)")  "* Set kinematics for x=0.5/0.5/0.125, keeping momentum"
    write (u, "(A)")

    r = [0.5_default, 0.5_default, 0.125_default]
    rb = 1 - r
    sf_int%on_shell_mode = KEEP_MOMENTUM
    call sf_int%complete_kinematics (x, f, r, rb, map=.false.)
    call sf_int%write (u)

    write (u, "(A)")
    write (u, "(A)")  "* Recover x and r"
    write (u, "(A)")

    call sf_int%recover_x (x)
    call sf_int%inverse_kinematics (x, f, r, rb, map=.false.)
    write (u, "(A,9(1x,F10.7))")  "x =", x
    write (u, "(A,9(1x,F10.7))")  "r =", r
    
    write (u, "(A)")
    write (u, "(A)")  "* Cleanup"

    call sf_int%final ()
    call model%final ()
    
    write (u, "(A)")
    write (u, "(A)")  "* Test output end: sf_base_4"

  end subroutine sf_base_4

\eatline
\nwendcode{}\nwbegindocs{959}\nwdocspar
\subsubsection{Pair spectrum}
Construct and display a structure function object for a pair spectrum
(a structure function involving two particles simultaneously).
\nwenddocs{}\nwbegincode{960}\moddef{SF base: execute tests}\plusendmoddef
  call test (sf_base_5, "sf_base_5", &
       "pair spectrum with radiation", &
       u, results)
\nwendcode{}\nwbegincode{961}\moddef{SF base: tests}\plusendmoddef
  subroutine sf_base_5 (u)
    integer, intent(in) :: u
    type(model_data_t), target :: model
    type(pdg_array_t) :: pdg_in
    type(pdg_array_t), dimension(2) :: pdg_out
    integer, dimension(:), allocatable :: pdg1, pdg2
    type(flavor_t) :: flv
    class(sf_data_t), allocatable, target :: data
    class(sf_int_t), allocatable :: sf_int
    type(vector4_t), dimension(2) :: k
    type(vector4_t), dimension(4) :: q
    real(default) :: E
    real(default), dimension(:), allocatable :: r, rb, x
    real(default) :: f
    
    write (u, "(A)")  "* Test output: sf_base_5"
    write (u, "(A)")  "*   Purpose: initialize and fill &
         &a pair spectrum object"
    write (u, "(A)")
    
    write (u, "(A)")  "* Initialize configuration data"
    write (u, "(A)")

    call model%init_test ()
    call flv%init (25, model)
    pdg_in = 25

    call reset_interaction_counter ()
    
    allocate (sf_test_spectrum_data_t :: data)
    select type (data)
    type is (sf_test_spectrum_data_t)
       call data%init (model, pdg_in, with_radiation=.true.)
    end select
       
    write (u, "(1x,A)")  "Outgoing particle codes:"
    call data%get_pdg_out (pdg_out)
    pdg1 = pdg_out(1)
    pdg2 = pdg_out(2)
    write (u, "(2x,99(1x,I0))")  pdg1, pdg2
    
    write (u, "(A)") 
    write (u, "(A)")  "* Initialize spectrum object"
    write (u, "(A)")
    
    call data%allocate_sf_int (sf_int)
    call sf_int%init (data)
    
    call sf_int%write (u)

    write (u, "(A)")
    write (u, "(A)")  "* Initialize incoming momenta with sqrts=1000"

    E = 500
    k(1) = vector4_moving (E, sqrt (E**2 - flv%get_mass ()**2), 3)
    k(2) = vector4_moving (E, sqrt (E**2 - flv%get_mass ()**2), 3)
    call sf_int%seed_kinematics (k)

    write (u, "(A)")
    write (u, "(A)")  "* Set kinematics for x=0.4,0.8"
    write (u, "(A)")

    allocate (r (data%get_n_par ()))
    allocate (rb(size (r)))
    allocate (x (size (r)))

    r = [0.4_default, 0.8_default]
    rb = 1 - r
    call sf_int%complete_kinematics (x, f, r, rb, map=.false.)
    call sf_int%write (u)

    write (u, "(A)")
    write (u, "(A,9(1x,F10.7))")  "x =", x
    write (u, "(A,9(1x,F10.7))")  "f =", f

    write (u, "(A)")
    write (u, "(A)")  "* Set kinematics with mapping for r=0.6,0.8"
    write (u, "(A)")

    r = [0.6_default, 0.8_default]
    rb = 1 - r
    call sf_int%complete_kinematics (x, f, r, rb, map=.true.)
    call sf_int%write (u)

    write (u, "(A)")
    write (u, "(A,9(1x,F10.7))")  "x =", x
    write (u, "(A,9(1x,F10.7))")  "f =", f

    write (u, "(A)")
    write (u, "(A)")  "* Recover x from momenta"
    write (u, "(A)")

    q = interaction_get_momenta (sf_int%interaction_t, outgoing=.true.)
    call sf_int%final ()
    deallocate (sf_int)

    call reset_interaction_counter ()
    call data%allocate_sf_int (sf_int)
    call sf_int%init (data)

    call sf_int%seed_kinematics (k)
    call interaction_set_momenta (sf_int%interaction_t, q, outgoing=.true.)
    call sf_int%recover_x (x)
    write (u, "(A,9(1x,F10.7))")  "x =", x

    write (u, "(A)")
    write (u, "(A)")  "* Compute inverse kinematics for x=0.36,0.64 &
         &and evaluate"
    write (u, "(A)")

    x = [0.36_default, 0.64_default]
    call sf_int%inverse_kinematics (x, f, r, rb, map=.true.)
    call sf_int%apply (scale=0._default)
    call sf_int%write (u)

    write (u, "(A)")
    write (u, "(A,9(1x,F10.7))")  "r =", r
    write (u, "(A,9(1x,F10.7))")  "f =", f

    write (u, "(A)")
    write (u, "(A)")  "* Cleanup"

    call sf_int%final ()
    call model%final ()
    
    write (u, "(A)")
    write (u, "(A)")  "* Test output end: sf_base_5"

  end subroutine sf_base_5

\eatline
\nwendcode{}\nwbegindocs{962}\nwdocspar
\subsubsection{Pair spectrum without radiation}
Construct and display a structure function object for a pair spectrum
(a structure function involving two particles simultaneously).
\nwenddocs{}\nwbegincode{963}\moddef{SF base: execute tests}\plusendmoddef
  call test (sf_base_6, "sf_base_6", &
       "pair spectrum without radiation", &
       u, results)
\nwendcode{}\nwbegincode{964}\moddef{SF base: tests}\plusendmoddef
  subroutine sf_base_6 (u)
    integer, intent(in) :: u
    type(model_data_t), target :: model
    type(pdg_array_t) :: pdg_in
    type(flavor_t) :: flv
    class(sf_data_t), allocatable, target :: data
    class(sf_int_t), allocatable :: sf_int
    type(vector4_t), dimension(2) :: k
    type(vector4_t), dimension(2) :: q
    real(default) :: E
    real(default), dimension(:), allocatable :: r, rb, x
    real(default) :: f
    
    write (u, "(A)")  "* Test output: sf_base_6"
    write (u, "(A)")  "*   Purpose: initialize and fill &
         &a pair spectrum object"
    write (u, "(A)")
    
    write (u, "(A)")  "* Initialize configuration data"
    write (u, "(A)")

    call model%init_test ()
    call flv%init (25, model)
    pdg_in = 25

    call reset_interaction_counter ()
    
    allocate (sf_test_spectrum_data_t :: data)
    select type (data)
    type is (sf_test_spectrum_data_t)
       call data%init (model, pdg_in, with_radiation=.false.)
    end select
       
    write (u, "(A)")  "* Initialize spectrum object"
    write (u, "(A)")
    
    call data%allocate_sf_int (sf_int)
    call sf_int%init (data)
    
    write (u, "(A)")  "* Initialize incoming momenta with sqrts=1000"

    E = 500
    k(1) = vector4_moving (E, sqrt (E**2 - flv%get_mass ()**2), 3)
    k(2) = vector4_moving (E, sqrt (E**2 - flv%get_mass ()**2), 3)
    call sf_int%seed_kinematics (k)

    write (u, "(A)")
    write (u, "(A)")  "* Set kinematics for x=0.4,0.8"
    write (u, "(A)")

    allocate (r (data%get_n_par ()))
    allocate (rb(size (r)))
    allocate (x (size (r)))

    r = [0.4_default, 0.8_default]
    rb = 1 - r
    call sf_int%complete_kinematics (x, f, r, rb, map=.false.)
    call sf_int%write (u)

    write (u, "(A)")
    write (u, "(A,9(1x,F10.7))")  "x =", x
    write (u, "(A,9(1x,F10.7))")  "f =", f

    write (u, "(A)")
    write (u, "(A)")  "* Recover x from momenta"
    write (u, "(A)")

    q = interaction_get_momenta (sf_int%interaction_t, outgoing=.true.)
    call sf_int%final ()
    deallocate (sf_int)

    call reset_interaction_counter ()
    call data%allocate_sf_int (sf_int)
    call sf_int%init (data)

    call sf_int%seed_kinematics (k)
    call interaction_set_momenta (sf_int%interaction_t, q, outgoing=.true.)
    call sf_int%recover_x (x)
    write (u, "(A,9(1x,F10.7))")  "x =", x

    write (u, "(A)")
    write (u, "(A)")  "* Compute inverse kinematics for x=0.4,0.8 &
         &and evaluate"
    write (u, "(A)")

    x = [0.4_default, 0.8_default]
    call sf_int%inverse_kinematics (x, f, r, rb, map=.false.)
    call sf_int%apply (scale=0._default)
    call sf_int%write (u)

    write (u, "(A)")
    write (u, "(A,9(1x,F10.7))")  "r =", r
    write (u, "(A,9(1x,F10.7))")  "f =", f

    write (u, "(A)")
    write (u, "(A)")  "* Cleanup"

    call sf_int%final ()
    call model%final ()
    
    write (u, "(A)")
    write (u, "(A)")  "* Test output end: sf_base_6"

  end subroutine sf_base_6

\eatline
\nwendcode{}\nwbegindocs{965}\nwdocspar
\subsubsection{Direct access to structure function}
Probe a structure function directly.
\nwenddocs{}\nwbegincode{966}\moddef{SF base: execute tests}\plusendmoddef
  call test (sf_base_7, "sf_base_7", &
       "direct access", &
       u, results)
\nwendcode{}\nwbegincode{967}\moddef{SF base: tests}\plusendmoddef
  subroutine sf_base_7 (u)
    integer, intent(in) :: u
    type(model_data_t), target :: model
    type(pdg_array_t) :: pdg_in
    type(flavor_t) :: flv
    class(sf_data_t), allocatable, target :: data
    class(sf_int_t), allocatable :: sf_int
    real(default), dimension(:), allocatable :: value
    
    write (u, "(A)")  "* Test output: sf_base_7"
    write (u, "(A)")  "*   Purpose: check direct access method"
    write (u, "(A)")
    
    write (u, "(A)")  "* Initialize configuration data"
    write (u, "(A)")

    call model%init_test ()
    call flv%init (25, model)
    pdg_in = 25

    call reset_interaction_counter ()
    
    write (u, "(A)")  "* Initialize structure-function object"
    write (u, "(A)")
    
    allocate (sf_test_data_t :: data)
    select type (data)
    type is (sf_test_data_t)
       call data%init (model, pdg_in)
    end select
       
    call data%allocate_sf_int (sf_int)
    call sf_int%init (data)

    write (u, "(A)")  "* Probe structure function: states"
    write (u, "(A)")
    
    write (u, "(A,I0)")  "n_states = ", sf_int%get_n_states ()
    write (u, "(A,I0)")  "n_in     = ", sf_int%get_n_in ()
    write (u, "(A,I0)")  "n_rad    = ", sf_int%get_n_rad ()
    write (u, "(A,I0)")  "n_out    = ", sf_int%get_n_out ()
    write (u, "(A)")
    write (u, "(A)", advance="no")  "state(1)  = "
    call quantum_numbers_write (sf_int%get_state (1), u)
    write (u, *)
    
    allocate (value (sf_int%get_n_states ()))
    call sf_int%compute_values (value, &
         E=[500._default], x=[0.5_default], xb=[0.5_default], scale=0._default)

    write (u, "(A)")
    write (u, "(A)", advance="no")  "value (E=500, x=0.5) ="
    write (u, "(9(1x," // FMT_19 // "))")  value

    call sf_int%compute_values (value, &
         x=[0.1_default], xb=[0.9_default], scale=0._default)

    write (u, "(A)")
    write (u, "(A)", advance="no")  "value (E=500, x=0.1) ="
    write (u, "(9(1x," // FMT_19 // "))")  value


    write (u, "(A)")
    write (u, "(A)")  "* Initialize spectrum object"
    write (u, "(A)")
    
    deallocate (value)
    call sf_int%final ()
    deallocate (sf_int)
    deallocate (data)
    
    allocate (sf_test_spectrum_data_t :: data)
    select type (data)
    type is (sf_test_spectrum_data_t)
       call data%init (model, pdg_in, with_radiation=.false.)
    end select
       
    call data%allocate_sf_int (sf_int)
    call sf_int%init (data)

    write (u, "(A)")  "* Probe spectrum: states"
    write (u, "(A)")
    
    write (u, "(A,I0)")  "n_states = ", sf_int%get_n_states ()
    write (u, "(A,I0)")  "n_in     = ", sf_int%get_n_in ()
    write (u, "(A,I0)")  "n_rad    = ", sf_int%get_n_rad ()
    write (u, "(A,I0)")  "n_out    = ", sf_int%get_n_out ()
    write (u, "(A)")
    write (u, "(A)", advance="no")  "state(1)  = "
    call quantum_numbers_write (sf_int%get_state (1), u)
    write (u, *)
    
    allocate (value (sf_int%get_n_states ()))
    call sf_int%compute_value (1, value(1), &
         E = [500._default, 500._default], &
         x = [0.5_default, 0.6_default], &
         xb= [0.5_default, 0.4_default], &
         scale = 0._default)

    write (u, "(A)")
    write (u, "(A)", advance="no")  "value (E=500,500, x=0.5,0.6) ="
    write (u, "(9(1x," // FMT_19 // "))")  value

    write (u, "(A)")
    write (u, "(A)")  "* Cleanup"

    call sf_int%final ()
    call model%final ()
    
    write (u, "(A)")
    write (u, "(A)")  "* Test output end: sf_base_7"

  end subroutine sf_base_7

\eatline
\nwendcode{}\nwbegindocs{968}\nwdocspar
\subsubsection{Structure function chain configuration}
\nwenddocs{}\nwbegincode{969}\moddef{SF base: execute tests}\plusendmoddef
  call test (sf_base_8, "sf_base_8", &
       "structure function chain configuration", &
       u, results)
\nwendcode{}\nwbegincode{970}\moddef{SF base: tests}\plusendmoddef
  subroutine sf_base_8 (u)
    integer, intent(in) :: u
    type(model_data_t), target :: model
    type(flavor_t) :: flv
    type(pdg_array_t) :: pdg_in
    type(beam_data_t), target :: beam_data
    class(sf_data_t), allocatable, target :: data_strfun
    class(sf_data_t), allocatable, target :: data_spectrum
    type(sf_config_t), dimension(:), allocatable :: sf_config
    type(sf_chain_t) :: sf_chain
    
    write (u, "(A)")  "* Test output: sf_base_8"
    write (u, "(A)")  "*   Purpose: set up a structure-function chain"
    write (u, "(A)")
    
    write (u, "(A)")  "* Initialize configuration data"
    write (u, "(A)")

    call model%init_test ()
    call flv%init (25, model)
    pdg_in = 25

    call reset_interaction_counter ()
    
    call beam_data_init_sqrts (beam_data, &
         1000._default, [flv, flv])

    allocate (sf_test_data_t :: data_strfun)
    select type (data_strfun)
    type is (sf_test_data_t)
       call data_strfun%init (model, pdg_in)
    end select
       
    allocate (sf_test_spectrum_data_t :: data_spectrum)
    select type (data_spectrum)
    type is (sf_test_spectrum_data_t)
       call data_spectrum%init (model, pdg_in, with_radiation=.true.)
    end select
       
    write (u, "(A)")  "* Set up chain with beams only"
    write (u, "(A)")
    
    call sf_chain%init (beam_data)
    call write_separator (u, 2)
    call sf_chain%write (u)
    call write_separator (u, 2)
    call sf_chain%final ()
    
    write (u, "(A)")
    write (u, "(A)")  "* Set up chain with structure function"
    write (u, "(A)")
    
    allocate (sf_config (1))
    call sf_config(1)%init ([1], data_strfun)
    call sf_chain%init (beam_data, sf_config)

    call write_separator (u, 2)
    call sf_chain%write (u)
    call write_separator (u, 2)
    call sf_chain%final ()

    write (u, "(A)")
    write (u, "(A)")  "* Set up chain with spectrum and structure function"
    write (u, "(A)")
    
    deallocate (sf_config)
    allocate (sf_config (2))
    call sf_config(1)%init ([1,2], data_spectrum)
    call sf_config(2)%init ([2], data_strfun)
    call sf_chain%init (beam_data, sf_config)

    call write_separator (u, 2)
    call sf_chain%write (u)
    call write_separator (u, 2)
    call sf_chain%final ()

    write (u, "(A)")
    write (u, "(A)")  "* Cleanup"

    call model%final ()
    
    write (u, "(A)")
    write (u, "(A)")  "* Test output end: sf_base_8"

  end subroutine sf_base_8

\eatline
\nwendcode{}\nwbegindocs{971}\nwdocspar
\subsubsection{Structure function instance configuration}
We create a structure-function chain instance which implements a
configured structure-function chain.  We link the momentum entries in
the interactions and compute kinematics.

We do not actually connect the interactions and create evaluators.  We
skip this step and manually advance the status of the chain instead.
\nwenddocs{}\nwbegincode{972}\moddef{SF base: execute tests}\plusendmoddef
  call test (sf_base_9, "sf_base_9", &
       "structure function chain instance", &
       u, results)
\nwendcode{}\nwbegincode{973}\moddef{SF base: tests}\plusendmoddef
  subroutine sf_base_9 (u)
    integer, intent(in) :: u
    type(model_data_t), target :: model
    type(flavor_t) :: flv
    type(pdg_array_t) :: pdg_in
    type(beam_data_t), target :: beam_data
    class(sf_data_t), allocatable, target :: data_strfun
    class(sf_data_t), allocatable, target :: data_spectrum
    type(sf_config_t), dimension(:), allocatable, target :: sf_config
    type(sf_chain_t), target :: sf_chain
    type(sf_chain_instance_t), target :: sf_chain_instance
    type(sf_channel_t), dimension(2) :: sf_channel
    type(vector4_t), dimension(2) :: p
    integer :: j
    
    write (u, "(A)")  "* Test output: sf_base_9"
    write (u, "(A)")  "*   Purpose: set up a structure-function chain &
         &and create an instance"
    write (u, "(A)")  "*            compute kinematics"
    write (u, "(A)")
   
    write (u, "(A)")  "* Initialize configuration data"
    write (u, "(A)")

    call model%init_test ()
    call flv%init (25, model)
    pdg_in = 25

    call reset_interaction_counter ()
    
    call beam_data_init_sqrts (beam_data, &
         1000._default, [flv, flv])

    allocate (sf_test_data_t :: data_strfun)
    select type (data_strfun)
    type is (sf_test_data_t)
       call data_strfun%init (model, pdg_in)
    end select
       
    allocate (sf_test_spectrum_data_t :: data_spectrum)
    select type (data_spectrum)
    type is (sf_test_spectrum_data_t)
       call data_spectrum%init (model, pdg_in, with_radiation=.true.)
    end select
       
    write (u, "(A)")  "* Set up chain with beams only"
    write (u, "(A)")
    
    call sf_chain%init (beam_data)

    call sf_chain_instance%init (sf_chain, n_channel = 1)

    call sf_chain_instance%link_interactions ()
    sf_chain_instance%status = SF_DONE_CONNECTIONS
    call sf_chain_instance%compute_kinematics (1, [real(default) ::])

    call write_separator (u, 2)
    call sf_chain%write (u)
    call write_separator (u, 2)
    call sf_chain_instance%write (u)
    call write_separator (u, 2)

    call sf_chain_instance%get_out_momenta (p)
    
    write (u, "(A)")
    write (u, "(A)")  "* Outgoing momenta:"
    
    do j = 1, 2
       write (u, "(A)")
       call vector4_write (p(j), u)
    end do
  
    call sf_chain_instance%final ()
    call sf_chain%final ()

    write (u, "(A)")
    write (u, "(A)")  "* Set up chain with structure function"
    write (u, "(A)")
    
    allocate (sf_config (1))
    call sf_config(1)%init ([1], data_strfun)
    call sf_chain%init (beam_data, sf_config)

    call sf_chain_instance%init (sf_chain, n_channel = 1)
    
    call sf_channel(1)%init (1)
    call sf_channel(1)%activate_mapping ([1])
    call sf_chain_instance%set_channel (1, sf_channel(1))

    call sf_chain_instance%link_interactions ()
    sf_chain_instance%status = SF_DONE_CONNECTIONS
    call sf_chain_instance%compute_kinematics (1, [0.8_default])
    
    call write_separator (u, 2)
    call sf_chain%write (u)
    call write_separator (u, 2)
    call sf_chain_instance%write (u)
    call write_separator (u, 2)

    call sf_chain_instance%get_out_momenta (p)
    
    write (u, "(A)")
    write (u, "(A)")  "* Outgoing momenta:"
    
    do j = 1, 2
       write (u, "(A)")
       call vector4_write (p(j), u)
    end do
    
    call sf_chain_instance%final ()
    call sf_chain%final ()
    
    write (u, "(A)")
    write (u, "(A)")  "* Set up chain with spectrum and structure function"
    write (u, "(A)")
    
    deallocate (sf_config)
    allocate (sf_config (2))
    call sf_config(1)%init ([1,2], data_spectrum)
    call sf_config(2)%init ([2], data_strfun)
    call sf_chain%init (beam_data, sf_config)

    call sf_chain_instance%init (sf_chain, n_channel = 1)
    
    call sf_channel(2)%init (2)
    call sf_channel(2)%activate_mapping ([2])
    call sf_chain_instance%set_channel (1, sf_channel(2))

    call sf_chain_instance%link_interactions ()
    sf_chain_instance%status = SF_DONE_CONNECTIONS
    call sf_chain_instance%compute_kinematics &
         (1, [0.5_default, 0.6_default, 0.8_default])
    
    call write_separator (u, 2)
    call sf_chain%write (u)
    call write_separator (u, 2)
    call sf_chain_instance%write (u)
    call write_separator (u, 2)

    call sf_chain_instance%get_out_momenta (p)
    
    write (u, "(A)")
    write (u, "(A)")  "* Outgoing momenta:"
    
    do j = 1, 2
       write (u, "(A)")
       call vector4_write (p(j), u)
    end do
    
    call sf_chain_instance%final ()
    call sf_chain%final ()

    write (u, "(A)")
    write (u, "(A)")  "* Cleanup"

    call model%final ()
    
    write (u, "(A)")
    write (u, "(A)")  "* Test output end: sf_base_9"

  end subroutine sf_base_9

\eatline
\nwendcode{}\nwbegindocs{974}\nwdocspar
\subsubsection{Structure function chain mappings}
Set up a structure function chain instance with a pair of
single-particle structure functions.  We test different global
mappings for this setup.

Again, we skip evaluators.
\nwenddocs{}\nwbegincode{975}\moddef{SF base: execute tests}\plusendmoddef
 call test (sf_base_10, "sf_base_10", &
      "structure function chain mapping", &
      u, results)
\nwendcode{}\nwbegincode{976}\moddef{SF base: tests}\plusendmoddef
  subroutine sf_base_10 (u)
    integer, intent(in) :: u
    type(model_data_t), target :: model
    type(flavor_t) :: flv
    type(pdg_array_t) :: pdg_in
    type(beam_data_t), target :: beam_data
    class(sf_data_t), allocatable, target :: data_strfun
    type(sf_config_t), dimension(:), allocatable, target :: sf_config
    type(sf_chain_t), target :: sf_chain
    type(sf_chain_instance_t), target :: sf_chain_instance
    type(sf_channel_t), dimension(2) :: sf_channel
    real(default), dimension(2) :: x_saved
    
    write (u, "(A)")  "* Test output: sf_base_10"
    write (u, "(A)")  "*   Purpose: set up a structure-function chain"
    write (u, "(A)")  "*            and check mappings"
    write (u, "(A)")
    
    write (u, "(A)")  "* Initialize configuration data"
    write (u, "(A)")

    call model%init_test ()
    call flv%init (25, model)
    pdg_in = 25

    call reset_interaction_counter ()
    
    call beam_data_init_sqrts (beam_data, &
         1000._default, [flv, flv])

    allocate (sf_test_data_t :: data_strfun)
    select type (data_strfun)
    type is (sf_test_data_t)
       call data_strfun%init (model, pdg_in)
    end select
       
    write (u, "(A)")  "* Set up chain with structure function pair &
         &and standard mapping"
    write (u, "(A)")
    
    allocate (sf_config (2))
    call sf_config(1)%init ([1], data_strfun)
    call sf_config(2)%init ([2], data_strfun)
    call sf_chain%init (beam_data, sf_config)

    call sf_chain_instance%init (sf_chain, n_channel = 1)

    call sf_channel(1)%init (2)
    call sf_channel(1)%set_s_mapping ([1,2])
    call sf_chain_instance%set_channel (1, sf_channel(1))

    call sf_chain_instance%link_interactions ()
    sf_chain_instance%status = SF_DONE_CONNECTIONS
    call sf_chain_instance%compute_kinematics (1, [0.8_default, 0.6_default])
    
    call write_separator (u, 2)
    call sf_chain_instance%write (u)
    call write_separator (u, 2)

    write (u, "(A)")
    write (u, "(A)")  "* Invert the kinematics calculation"
    write (u, "(A)")

    x_saved = sf_chain_instance%x

    call sf_chain_instance%init (sf_chain, n_channel = 1)

    call sf_channel(2)%init (2)
    call sf_channel(2)%set_s_mapping ([1, 2])
    call sf_chain_instance%set_channel (1, sf_channel(2))

    call sf_chain_instance%link_interactions ()
    sf_chain_instance%status = SF_DONE_CONNECTIONS
    call sf_chain_instance%inverse_kinematics (x_saved)
    
    call write_separator (u, 2)
    call sf_chain_instance%write (u)
    call write_separator (u, 2)

    
    call sf_chain_instance%final ()
    call sf_chain%final ()
    
    write (u, "(A)")
    write (u, "(A)")  "* Cleanup"

    call model%final ()
    
    write (u, "(A)")
    write (u, "(A)")  "* Test output end: sf_base_10"

  end subroutine sf_base_10

\eatline
\nwendcode{}\nwbegindocs{977}\nwdocspar
\subsubsection{Structure function chain evaluation}
Here, we test the complete workflow for structure-function chains.
First, we create the template chain, then initialize an instance.  We
set up links, mask, and evaluators.  Finally, we set kinematics and
evaluate the matrix elements and their products.
\nwenddocs{}\nwbegincode{978}\moddef{SF base: execute tests}\plusendmoddef
  call test (sf_base_11, "sf_base_11", &
       "structure function chain evaluation", &
       u, results)
\nwendcode{}\nwbegincode{979}\moddef{SF base: tests}\plusendmoddef
  subroutine sf_base_11 (u)
    integer, intent(in) :: u
    type(model_data_t), target :: model
    type(flavor_t) :: flv
    type(pdg_array_t) :: pdg_in
    type(beam_data_t), target :: beam_data
    class(sf_data_t), allocatable, target :: data_strfun
    class(sf_data_t), allocatable, target :: data_spectrum
    type(sf_config_t), dimension(:), allocatable, target :: sf_config
    type(sf_chain_t), target :: sf_chain
    type(sf_chain_instance_t), target :: sf_chain_instance
    type(sf_channel_t), dimension(2) :: sf_channel
    type(particle_set_t) :: pset
    type(interaction_t), pointer :: int
    logical :: ok
    
    write (u, "(A)")  "* Test output: sf_base_11"
    write (u, "(A)")  "*   Purpose: set up a structure-function chain"
    write (u, "(A)")  "*            create an instance and evaluate"
    write (u, "(A)")
    
    write (u, "(A)")  "* Initialize configuration data"
    write (u, "(A)")

    call model%init_test ()
    call flv%init (25, model)
    pdg_in = 25

    call reset_interaction_counter ()
    
    call beam_data_init_sqrts (beam_data, &
         1000._default, [flv, flv])

    allocate (sf_test_data_t :: data_strfun)
    select type (data_strfun)
    type is (sf_test_data_t)
       call data_strfun%init (model, pdg_in)
    end select
       
    allocate (sf_test_spectrum_data_t :: data_spectrum)
    select type (data_spectrum)
    type is (sf_test_spectrum_data_t)
       call data_spectrum%init (model, pdg_in, with_radiation=.true.)
    end select
       
    write (u, "(A)")  "* Set up chain with beams only"
    write (u, "(A)")
    
    call sf_chain%init (beam_data)

    call sf_chain_instance%init (sf_chain, n_channel = 1)
    call sf_chain_instance%link_interactions ()
    call sf_chain_instance%exchange_mask ()
    call sf_chain_instance%init_evaluators ()
    
    call sf_chain_instance%compute_kinematics (1, [real(default) ::])
    call sf_chain_instance%evaluate (scale=0._default)

    call write_separator (u, 2)
    call sf_chain_instance%write (u)
    call write_separator (u, 2)

    int => sf_chain_instance%get_out_int_ptr ()
    call pset%init (ok, int, int, FM_IGNORE_HELICITY, &
         [0._default, 0._default], .false., .true.)
    call sf_chain_instance%final ()

    write (u, "(A)")
    write (u, "(A)")  "* Particle content:"
    write (u, "(A)")

    call write_separator (u)
    call pset%write (u)
    call write_separator (u)

    write (u, "(A)")
    write (u, "(A)")  "* Recover chain:"
    write (u, "(A)")
    
    call sf_chain_instance%init (sf_chain, n_channel = 1)
    call sf_chain_instance%link_interactions ()
    call sf_chain_instance%exchange_mask ()
    call sf_chain_instance%init_evaluators ()

    int => sf_chain_instance%get_out_int_ptr ()
    call pset%fill_interaction (int, 2, check_match=.false.)

    call sf_chain_instance%recover_kinematics (1)
    call sf_chain_instance%evaluate (scale=0._default)

    call write_separator (u, 2)
    call sf_chain_instance%write (u)
    call write_separator (u, 2)

    call pset%final ()
    call sf_chain_instance%final ()
    call sf_chain%final ()

    write (u, "(A)")
    write (u, "(A)")
    write (u, "(A)")
    write (u, "(A)")  "* Set up chain with structure function"
    write (u, "(A)")
    
    allocate (sf_config (1))
    call sf_config(1)%init ([1], data_strfun)
    call sf_chain%init (beam_data, sf_config)

    call sf_chain_instance%init (sf_chain, n_channel = 1)
    call sf_channel(1)%init (1)
    call sf_channel(1)%activate_mapping ([1])
    call sf_chain_instance%set_channel (1, sf_channel(1))
    call sf_chain_instance%link_interactions ()
    call sf_chain_instance%exchange_mask ()
    call sf_chain_instance%init_evaluators ()

    call sf_chain_instance%compute_kinematics (1, [0.8_default])
    call sf_chain_instance%evaluate (scale=0._default)
    
    call write_separator (u, 2)
    call sf_chain_instance%write (u)
    call write_separator (u, 2)

    int => sf_chain_instance%get_out_int_ptr ()
    call pset%init (ok, int, int, FM_IGNORE_HELICITY, &
         [0._default, 0._default], .false., .true.)
    call sf_chain_instance%final ()

    write (u, "(A)")
    write (u, "(A)")  "* Particle content:"
    write (u, "(A)")

    call write_separator (u)
    call pset%write (u)
    call write_separator (u)

    write (u, "(A)")
    write (u, "(A)")  "* Recover chain:"
    write (u, "(A)")
    
    call sf_chain_instance%init (sf_chain, n_channel = 1)
    call sf_channel(1)%init (1)
    call sf_channel(1)%activate_mapping ([1])
    call sf_chain_instance%set_channel (1, sf_channel(1))
    call sf_chain_instance%link_interactions ()
    call sf_chain_instance%exchange_mask ()
    call sf_chain_instance%init_evaluators ()

    int => sf_chain_instance%get_out_int_ptr ()
    call pset%fill_interaction (int, 2, check_match=.false.)

    call sf_chain_instance%recover_kinematics (1)
    call sf_chain_instance%evaluate (scale=0._default)

    call write_separator (u, 2)
    call sf_chain_instance%write (u)
    call write_separator (u, 2)

    call pset%final ()
    call sf_chain_instance%final ()
    call sf_chain%final ()

    write (u, "(A)")
    write (u, "(A)")
    write (u, "(A)")
    write (u, "(A)")  "* Set up chain with spectrum and structure function"
    write (u, "(A)")
    
    deallocate (sf_config)
    allocate (sf_config (2))
    call sf_config(1)%init ([1,2], data_spectrum)
    call sf_config(2)%init ([2], data_strfun)
    call sf_chain%init (beam_data, sf_config)
    
    call sf_chain_instance%init (sf_chain, n_channel = 1)
    call sf_channel(2)%init (2)
    call sf_channel(2)%activate_mapping ([2])
    call sf_chain_instance%set_channel (1, sf_channel(2))
    call sf_chain_instance%link_interactions ()
    call sf_chain_instance%exchange_mask ()
    call sf_chain_instance%init_evaluators ()

    call sf_chain_instance%compute_kinematics &
         (1, [0.5_default, 0.6_default, 0.8_default])
    call sf_chain_instance%evaluate (scale=0._default)
    
    call write_separator (u, 2)
    call sf_chain_instance%write (u)
    call write_separator (u, 2)

    int => sf_chain_instance%get_out_int_ptr ()
    call pset%init (ok, int, int, FM_IGNORE_HELICITY, &
         [0._default, 0._default], .false., .true.)
    call sf_chain_instance%final ()

    write (u, "(A)")
    write (u, "(A)")  "* Particle content:"
    write (u, "(A)")

    call write_separator (u)
    call pset%write (u)
    call write_separator (u)

    write (u, "(A)")
    write (u, "(A)")  "* Recover chain:"
    write (u, "(A)")
    
    call sf_chain_instance%init (sf_chain, n_channel = 1)
    call sf_channel(2)%init (2)
    call sf_channel(2)%activate_mapping ([2])
    call sf_chain_instance%set_channel (1, sf_channel(2))
    call sf_chain_instance%link_interactions ()
    call sf_chain_instance%exchange_mask ()
    call sf_chain_instance%init_evaluators ()

    int => sf_chain_instance%get_out_int_ptr ()
    call pset%fill_interaction (int, 2, check_match=.false.)

    call sf_chain_instance%recover_kinematics (1)
    call sf_chain_instance%evaluate (scale=0._default)

    call write_separator (u, 2)
    call sf_chain_instance%write (u)
    call write_separator (u, 2)

    call pset%final ()
    call sf_chain_instance%final ()
    call sf_chain%final ()

    write (u, "(A)")
    write (u, "(A)")  "* Cleanup"

    call model%final ()
    
    write (u, "(A)")
    write (u, "(A)")  "* Test output end: sf_base_11"

  end subroutine sf_base_11

\eatline
\nwendcode{}\nwbegindocs{980}\nwdocspar
\subsubsection{Multichannel case}
We set up a structure-function chain as before, but with three
different parameterizations.  The first instance is without mappings,
the second one with single-particle mappings, and the third one with
two-particle mappings.
\nwenddocs{}\nwbegincode{981}\moddef{SF base: execute tests}\plusendmoddef
  call test (sf_base_12, "sf_base_12", &
       "multi-channel structure function chain", &
       u, results)
\nwendcode{}\nwbegincode{982}\moddef{SF base: tests}\plusendmoddef
  subroutine sf_base_12 (u)
    integer, intent(in) :: u
    type(model_data_t), target :: model
    type(flavor_t) :: flv
    type(pdg_array_t) :: pdg_in
    type(beam_data_t), target :: beam_data
    class(sf_data_t), allocatable, target :: data
    type(sf_config_t), dimension(:), allocatable, target :: sf_config
    type(sf_chain_t), target :: sf_chain
    type(sf_chain_instance_t), target :: sf_chain_instance
    real(default), dimension(2) :: x_saved
    real(default), dimension(2,3) :: p_saved
    type(sf_channel_t), dimension(:), allocatable :: sf_channel
    
    write (u, "(A)")  "* Test output: sf_base_12"
    write (u, "(A)")  "*   Purpose: set up and evaluate a multi-channel &
         &structure-function chain"
    write (u, "(A)")
    
    write (u, "(A)")  "* Initialize configuration data"
    write (u, "(A)")

    call model%init_test ()
    call flv%init (25, model)
    pdg_in = 25

    call reset_interaction_counter ()
    
    call beam_data_init_sqrts (beam_data, &
         1000._default, [flv, flv])

    allocate (sf_test_data_t :: data)
    select type (data)
    type is (sf_test_data_t)
       call data%init (model, pdg_in)
    end select
       
    write (u, "(A)")  "* Set up chain with structure function pair &
         &and three different mappings"
    write (u, "(A)")
    
    allocate (sf_config (2))
    call sf_config(1)%init ([1], data)
    call sf_config(2)%init ([2], data) 
    call sf_chain%init (beam_data, sf_config)

    call sf_chain_instance%init (sf_chain, n_channel = 3)

    call allocate_sf_channels (sf_channel, n_channel = 3, n_strfun = 2)

    ! channel 1: no mapping
    call sf_chain_instance%set_channel (1, sf_channel(1))

    ! channel 2: single-particle mappings
    call sf_channel(2)%activate_mapping ([1,2])
    ! call sf_chain_instance%activate_mapping (2, [1,2])
    call sf_chain_instance%set_channel (2, sf_channel(2))
   
    ! channel 3: two-particle mapping
    call sf_channel(3)%set_s_mapping ([1,2])
    ! call sf_chain_instance%set_s_mapping (3, [1, 2])
    call sf_chain_instance%set_channel (3, sf_channel(3))

    call sf_chain_instance%link_interactions ()
    call sf_chain_instance%exchange_mask ()
    call sf_chain_instance%init_evaluators ()

    write (u, "(A)")  "* Compute kinematics in channel 1 and evaluate"
    write (u, "(A)")

    call sf_chain_instance%compute_kinematics (1, [0.8_default, 0.6_default])
    call sf_chain_instance%evaluate (scale=0._default)
    
    call write_separator (u, 2)
    call sf_chain_instance%write (u)
    call write_separator (u, 2)

    write (u, "(A)")
    write (u, "(A)")  "* Invert the kinematics calculation"
    write (u, "(A)")

    x_saved = sf_chain_instance%x

    call sf_chain_instance%inverse_kinematics (x_saved)
    call sf_chain_instance%evaluate (scale=0._default)
    
    call write_separator (u, 2)
    call sf_chain_instance%write (u)
    call write_separator (u, 2)

    write (u, "(A)")
    write (u, "(A)")  "* Compute kinematics in channel 2 and evaluate"
    write (u, "(A)")

    p_saved = sf_chain_instance%p

    call sf_chain_instance%compute_kinematics (2, p_saved(:,2))
    call sf_chain_instance%evaluate (scale=0._default)
    
    call write_separator (u, 2)
    call sf_chain_instance%write (u)
    call write_separator (u, 2)

    write (u, "(A)")
    write (u, "(A)")  "* Compute kinematics in channel 3 and evaluate"
    write (u, "(A)")

    call sf_chain_instance%compute_kinematics (3, p_saved(:,3))
    call sf_chain_instance%evaluate (scale=0._default)
    
    call write_separator (u, 2)
    call sf_chain_instance%write (u)
    call write_separator (u, 2)

    write (u, "(A)")
    write (u, "(A)")  "* Cleanup"

    call sf_chain_instance%final ()
    call sf_chain%final ()

    call model%final ()
    
    write (u, "(A)")
    write (u, "(A)")  "* Test output end: sf_base_12"

  end subroutine sf_base_12

\eatline
\nwendcode{}\nwbegindocs{983}\nwdocspar
\subsubsection{Generated spectrum}
Construct and evaluate a structure function object for a pair spectrum
which is evaluated as a beam-event generator.
\nwenddocs{}\nwbegincode{984}\moddef{SF base: execute tests}\plusendmoddef
  call test (sf_base_13, "sf_base_13", &
       "pair spectrum generator", &
       u, results)
\nwendcode{}\nwbegincode{985}\moddef{SF base: tests}\plusendmoddef
  subroutine sf_base_13 (u)
    integer, intent(in) :: u
    type(model_data_t), target :: model
    type(flavor_t) :: flv
    type(pdg_array_t) :: pdg_in
    class(sf_data_t), allocatable, target :: data
    class(sf_int_t), allocatable :: sf_int
    type(vector4_t), dimension(2) :: k
    type(vector4_t), dimension(2) :: q
    real(default) :: E
    real(default), dimension(:), allocatable :: r, rb, x
    real(default) :: f, x_free
    
    write (u, "(A)")  "* Test output: sf_base_13"
    write (u, "(A)")  "*   Purpose: initialize and fill &
         &a pair generator object"
    write (u, "(A)")
    
    write (u, "(A)")  "* Initialize configuration data"
    write (u, "(A)")

    call model%init_test ()
    call flv%init (25, model)
    pdg_in = 25

    call reset_interaction_counter ()
    
    allocate (sf_test_generator_data_t :: data)
    select type (data)
    type is (sf_test_generator_data_t)
       call data%init (model, pdg_in)
    end select
       
    write (u, "(A)")  "* Initialize generator object"
    write (u, "(A)")
    
    call data%allocate_sf_int (sf_int)
    call sf_int%init (data)

    allocate (r (data%get_n_par ()))
    allocate (rb(size (r)))
    allocate (x (size (r)))
    
    write (u, "(A)")  "* Generate free r values"
    write (u, "(A)")

    x_free = 1
    call sf_int%generate_free (r, rb, x_free)

    write (u, "(A)")  "* Initialize incoming momenta with sqrts=1000"

    E = 500
    k(1) = vector4_moving (E, sqrt (E**2 - flv%get_mass ()**2), 3)
    k(2) = vector4_moving (E, sqrt (E**2 - flv%get_mass ()**2), 3)
    call sf_int%seed_kinematics (k)

    write (u, "(A)")
    write (u, "(A)")  "* Complete kinematics"
    write (u, "(A)")

    call sf_int%complete_kinematics (x, f, r, rb, map=.false.)
    call sf_int%write (u)

    write (u, "(A)")
    write (u, "(A,9(1x,F10.7))")  "x =", x
    write (u, "(A,9(1x,F10.7))")  "f =", f
    write (u, "(A,9(1x,F10.7))")  "xf=", x_free

    write (u, "(A)")
    write (u, "(A)")  "* Recover x from momenta"
    write (u, "(A)")

    q = interaction_get_momenta (sf_int%interaction_t, outgoing=.true.)
    call sf_int%final ()
    deallocate (sf_int)

    call reset_interaction_counter ()
    call data%allocate_sf_int (sf_int)
    call sf_int%init (data)

    call sf_int%seed_kinematics (k)
    call interaction_set_momenta (sf_int%interaction_t, q, outgoing=.true.)
    x_free = 1
    call sf_int%recover_x (x, x_free)
    write (u, "(A,9(1x,F10.7))")  "x =", x
    write (u, "(A,9(1x,F10.7))")  "xf=", x_free

    write (u, "(A)")
    write (u, "(A)")  "* Compute inverse kinematics &
         &and evaluate"
    write (u, "(A)")

    call sf_int%inverse_kinematics (x, f, r, rb, map=.false.)
    call sf_int%apply (scale=0._default)
    call sf_int%write (u)

    write (u, "(A)")
    write (u, "(A,9(1x,F10.7))")  "r =", r
    write (u, "(A,9(1x,F10.7))")  "f =", f

    write (u, "(A)")
    write (u, "(A)")  "* Cleanup"

    call sf_int%final ()
    call model%final ()
    
    write (u, "(A)")
    write (u, "(A)")  "* Test output end: sf_base_13"

  end subroutine sf_base_13

\eatline
\nwendcode{}\nwbegindocs{986}\nwdocspar
\subsubsection{Structure function chain evaluation}
Here, we test the complete workflow for a structure-function chain
with generator.  First, we create the template chain, then initialize
an instance.  We set up links, mask, and evaluators.  Finally, we set
kinematics and evaluate the matrix elements and their products.
\nwenddocs{}\nwbegincode{987}\moddef{SF base: execute tests}\plusendmoddef
  call test (sf_base_14, "sf_base_14", &
       "structure function generator evaluation", &
       u, results)
\nwendcode{}\nwbegincode{988}\moddef{SF base: tests}\plusendmoddef
  subroutine sf_base_14 (u)
    integer, intent(in) :: u
    type(model_data_t), target :: model
    type(flavor_t) :: flv
    type(pdg_array_t) :: pdg_in
    type(beam_data_t), target :: beam_data
    class(sf_data_t), allocatable, target :: data_strfun
    class(sf_data_t), allocatable, target :: data_generator
    type(sf_config_t), dimension(:), allocatable, target :: sf_config
    real(default), dimension(:), allocatable :: p_in
    type(sf_chain_t), target :: sf_chain
    type(sf_chain_instance_t), target :: sf_chain_instance
    
    write (u, "(A)")  "* Test output: sf_base_14"
    write (u, "(A)")  "*   Purpose: set up a structure-function chain"
    write (u, "(A)")  "*            create an instance and evaluate"
    write (u, "(A)")
    
    write (u, "(A)")  "* Initialize configuration data"
    write (u, "(A)")

    call model%init_test ()
    call flv%init (25, model)
    pdg_in = 25

    call reset_interaction_counter ()
    
    call beam_data_init_sqrts (beam_data, &
         1000._default, [flv, flv])

    allocate (sf_test_data_t :: data_strfun)
    select type (data_strfun)
    type is (sf_test_data_t)
       call data_strfun%init (model, pdg_in)
    end select
       
    allocate (sf_test_generator_data_t :: data_generator)
    select type (data_generator)
    type is (sf_test_generator_data_t)
       call data_generator%init (model, pdg_in)
    end select
       
    write (u, "(A)")  "* Set up chain with generator and structure function"
    write (u, "(A)")
    
    allocate (sf_config (2))
    call sf_config(1)%init ([1,2], data_generator)
    call sf_config(2)%init ([2], data_strfun)
    call sf_chain%init (beam_data, sf_config)
    
    call sf_chain_instance%init (sf_chain, n_channel = 1)
    call sf_chain_instance%link_interactions ()
    call sf_chain_instance%exchange_mask ()
    call sf_chain_instance%init_evaluators ()

    write (u, "(A)")  "* Inject integration parameter"
    write (u, "(A)")

    allocate (p_in (sf_chain%get_n_bound ()), source = 0.9_default)
    write (u, "(A,9(1x,F10.7))")  "p_in =", p_in
    
    write (u, "(A)")
    write (u, "(A)")  "* Evaluate"
    write (u, "(A)")

    call sf_chain_instance%compute_kinematics (1, p_in)
    call sf_chain_instance%evaluate (scale=0._default)
    
    call sf_chain_instance%write (u)

    write (u, "(A)")
    write (u, "(A)")  "* Extract integration parameter"
    write (u, "(A)")

    call sf_chain_instance%get_mcpar (1, p_in)
    write (u, "(A,9(1x,F10.7))")  "p_in =", p_in
    
    call sf_chain_instance%final ()
    call sf_chain%final ()

    write (u, "(A)")
    write (u, "(A)")  "* Cleanup"

    call model%final ()
    
    write (u, "(A)")
    write (u, "(A)")  "* Test output end: sf_base_14"

  end subroutine sf_base_14

\eatline
\nwendcode{}\nwbegindocs{989}\nwdocspar
%%%%%%%%%%%%%%%%%%%%%%%%%%%%%%%%%%%%%%%%%%%%%%%%%%%%%%%%%%%%%%%%%%%%%%%%
\section{Photon radiation: ISR}

\nwenddocs{}\nwbegincode{990}\moddef{\code{}sf{\char95}isr.f90\edoc{}}\endmoddef
\LA{}File header\RA{}

module sf_isr
 
\LA{}Use kinds\RA{}
\LA{}Use strings\RA{}
  use io_units
  use constants, only: pi
  use format_defs, only: FMT_12, FMT_17, FMT_19
  use unit_tests
  use diagnostics
  use physics_defs, only: ELECTRON, PHOTON
  use lorentz
  use sm_physics, only: Li2
  use pdg_arrays
  use model_data
  use flavors
  use colors
  use quantum_numbers
  use state_matrices
  use interactions
  use polarizations
  use sf_mappings
  use sf_aux
  use sf_base

\LA{}Standard module head\RA{}

\LA{}SF isr: public\RA{}

\LA{}SF isr: parameters\RA{}

\LA{}SF isr: types\RA{}

contains

\LA{}SF isr: procedures\RA{}

\LA{}SF isr: tests\RA{}

end module sf_isr
\eatline
\nwendcode{}\nwbegindocs{991}\nwdocspar
\subsection{Physics}
The ISR structure function is in the most crude approximation (LLA
without $\alpha$ corrections, i.e. $\epsilon^0$)
\begin{equation}
  f_0(x) = \epsilon (1-x)^{-1+\epsilon} \qquad\text{with}\qquad
  \epsilon = \frac{\alpha}{\pi}q_e^2\ln\frac{s}{m^2},
\end{equation}
where $m$ is the mass of the incoming (and outgoing) particle, which
is initially assumed on-shell.

Here, the form of $\epsilon$ results from the kinematical bounds for
the momentum squared of the outgoing particle, which in the limit
$m^2\ll s$ are given by
\begin{align}
  t_0 &= -2\bar xE(E+p) + m^2 \approx -\bar x s,
\\
  t_1 &= -2\bar xE(E-p) + m^2 \approx x m^2,
\end{align}
so the integration over the propagator $1/(t-m^2)$ yields
\begin{equation}
  \ln\frac{t_0-m^2}{t_1-m^2} = \ln\frac{s}{m^2}.
\end{equation}

In $f_0(x)$, there is an integrable singularity at $x=1$ which does
not spoil the integration, but would lead to an unbounded $f_{\rm
max}$.  Therefore, we map this singularity like
\begin{equation}\label{ISR-mapping}
  x = 1 - (1-x')^{1/\epsilon}
\end{equation}
such that
\begin{equation}
  \int dx\,f_0(x) = \int dx'
\end{equation}
The structure function has three parameters: $\alpha$, $m_{\rm in}$ of
the incoming particle and $s$, the hard scale.  Internally, we store
the exponent $\epsilon$ which is the relevant parameter.  (In
conventional notation, $\epsilon=\beta/2$.)  As defaults, we take the
actual values of $\alpha$ (which is probably $\alpha(s)$), the actual
mass $m_{\rm in}$ and the squared total c.m. energy $s$.

Including $\epsilon$, $\epsilon^2$, and $\epsilon^3$ corrections, the
successive approximation of the ISR structure function read
\begin{align}
  f_0(x) &= \epsilon(1-x)^{-1+\epsilon} \\
  f_1(x) &= g_1(\epsilon)\,f_0(x) - \frac{\epsilon}{2}(1+x) \\
\begin{split}
  f_2(x) &= g_2(\epsilon)\,f_0(x) - \frac{\epsilon}{2}(1+x) \\
  &\quad - \frac{\epsilon^2}{8}\left(
    \frac{1+3x^2}{1-x}\ln x + 4(1+x) \ln(1-x) + 5 + x \right) 
\end{split} \\
\begin{split}
  f_3(x) &= g_3(\epsilon)\,f_0(x) - \frac{\epsilon}{2}(1+x) \\
  &\quad - \frac{\epsilon^2}{8}\left(
    \frac{1+3x^2}{1-x}\ln x + 4(1+x) \ln(1-x) + 5 + x \right) \\
  &\quad - \frac{\epsilon^3}{48}\left( \vphantom{\frac{1}{1-x}}
    (1+x)\left[6\mathop{\rm Li_2}(x) + 12\ln^2(1-x) - 3\pi^2\right]\right.
     + 6(x+5)\ln(1-x) \\
  &\qquad\qquad + \frac{1}{1-x}\left[\frac32(1+8x+3x^2)\ln x
                               + 12(1+x^2)\ln x\ln(1-x) \right. \\
  &\qquad\qquad\qquad\qquad 
       \left.\left. - \frac12(1+7x^2)\ln^2x + \frac14(39-24x-15x^2)\right]
       \vphantom{\frac{1}{1-x}} \right)
\end{split}
\end{align}
where the successive approximations to the prefactor of the leading
singularity
\begin{equation}
  g(\epsilon) = \frac{\exp\left(\epsilon(-\gamma_E + \tfrac34)\right)}
              {\Gamma(1 + \epsilon)},
\end{equation}
are given by
\begin{align}
  g_0(\epsilon) &= 1 \\
  g_1(\epsilon) &= 1 + \frac34\epsilon \\
  g_2(\epsilon) &= 1 + \frac34\epsilon 
    + \frac{27 - 8\pi^2}{96}\epsilon^2 \\
  g_3(\epsilon) &= 1 + \frac34\epsilon 
    + \frac{27 - 8\pi^2}{96}\epsilon^2
    + \frac{27 - 24\pi^2 + 128 \zeta(3)}{384}\epsilon^3,
\end{align}
where, numerically
\begin{equation}
  \zeta(3) = 1.20205690315959428539973816151\ldots
\end{equation}
Although one could calculate the function $g(\epsilon)$ exactly,
truncating its Taylor expansion ensures the exact normalization of the
truncated structure function at each given order:
\begin{equation}
  \int_0^1 dx\,f_i(x) = 1 \qquad\text{for all $i$.}
\end{equation}

Effectively, the $O(\epsilon)$ correction reduces the low-$x$ tail of
the structure function by $50\%$ while increasing the coefficient of
the singularity by $O(\epsilon)$.  Relative to this, the
$O(\epsilon^2)$ correction slightly enhances $x>\frac12$ compared to
$x<\frac12$.  At $x=0$, $f_2(x)$ introduces a logarithmic singularity
which should be cut off at $x_0=O(e^{-1/\epsilon})$: for lower $x$ the
perturbative series breaks down.  The $f_3$ correction is slightly
positive for low $x$ values and negative near $x=1$, where the
$\mathop{\rm Li_2}$ piece slightly softens the singularity at $x=1$.

Instead of the definition for $\epsilon$ given above, it is customary
to include a universal nonlogarithmic piece:
\begin{equation}
  \epsilon = \frac{\alpha}{\pi}q_e^2\left(\ln\tfrac{s}{m^2} - 1\right)
\end{equation}


\subsection{Implementation}

In the concrete implementation, the zeroth order mapping
(\ref{ISR-mapping}) is implemented, and the Jacobian is equal to
$f_i(x)/f_0(x)$.  This can be written as
\begin{align}
  \frac{f_0(x)}{f_0(x)} &= 1 \\
  \frac{f_1(x)}{f_0(x)} &= 1 + \frac34\epsilon - \frac{1-x^2}{2(1-x')} \\
\begin{split}\label{ISR-f2}
  \frac{f_2(x)}{f_0(x)} &= 1 + \frac34\epsilon
    + \frac{27 - 8\pi^2}{96}\epsilon^2 
    - \frac{1-x^2}{2(1-x')} \\
  &\quad - \frac{(1+3x^2)\ln x 
    + (1-x)\left(4(1+x)\ln(1-x) + 5 + x\right)}{8(1-x')}\epsilon
\end{split}
\end{align}
%'
For $x=1$ (i.e., numerically indistinguishable from $1$), this reduces to
\begin{align}
  \frac{f_0(x)}{f_0(x)} &= 1 \\
  \frac{f_1(x)}{f_0(x)} &= 1 + \frac34\epsilon \\
  \frac{f_2(x)}{f_0(x)} &= 1 + \frac34\epsilon
    + \frac{27 - 8\pi^2}{96}\epsilon^2 
\end{align}
The last line in (\ref{ISR-f2}) is zero for
\begin{equation}
  x_{\rm min} = 0.00714053329734592839549879772019
\end{equation}
(Mathematica result), independent of $\epsilon$.  For $x$ values less
than this we ignore this correction because of the logarithmic
singularity which should in principle be resummed.


\subsection{The ISR data block}
\nwenddocs{}\nwbegincode{992}\moddef{SF isr: public}\endmoddef
  public :: isr_data_t
\nwendcode{}\nwbegincode{993}\moddef{SF isr: types}\endmoddef
  type, extends (sf_data_t) :: isr_data_t
     private
     class(model_data_t), pointer :: model => null ()
     type(flavor_t), dimension(:), allocatable :: flv_in
     real(default) :: alpha = 0
     real(default) :: q_max = 0
     real(default) :: real_mass = 0
     real(default) :: mass = 0
     real(default) :: eps = 0
     real(default) :: log = 0
     logical :: recoil = .false.
     integer :: order = 3
     integer :: error = NONE
   contains
   \LA{}SF isr: isr data: TBP\RA{}  
  end type isr_data_t

\eatline
\nwendcode{}\nwbegindocs{994}Error codes
\nwenddocs{}\nwbegincode{995}\moddef{SF isr: parameters}\endmoddef
  integer, parameter :: NONE = 0
  integer, parameter :: ZERO_MASS = 1
  integer, parameter :: Q_MAX_TOO_SMALL = 2
  integer, parameter :: EPS_TOO_LARGE = 3
  integer, parameter :: INVALID_ORDER = 4
  integer, parameter :: CHARGE_MIX = 5
  integer, parameter :: CHARGE_ZERO = 6
  integer, parameter :: MASS_MIX = 7
\nwendcode{}\nwbegindocs{996}Generate flavor-dependent ISR data:
\nwenddocs{}\nwbegincode{997}\moddef{SF isr: isr data: TBP}\endmoddef
  procedure :: init => isr_data_init
\nwendcode{}\nwbegincode{998}\moddef{SF isr: procedures}\endmoddef
  subroutine isr_data_init &
       (data, model, pdg_in, alpha, q_max, mass, order, recoil)
    class(isr_data_t), intent(out) :: data
    class(model_data_t), intent(in), target :: model
    type(pdg_array_t), intent(in) :: pdg_in
    real(default), intent(in) :: alpha
    real(default), intent(in) :: q_max
    real(default), intent(in), optional :: mass
    integer, intent(in), optional :: order
    logical, intent(in), optional :: recoil
    integer :: i, n_flv
    real(default) :: charge
    data%model => model
    n_flv = pdg_array_get_length (pdg_in)
    allocate (data%flv_in (n_flv))
    do i = 1, n_flv
       call data%flv_in(i)%init (pdg_array_get (pdg_in, i), model)
    end do
    data%alpha = alpha
    data%q_max = q_max
    if (present (order)) then
       call data%set_order (order)
    end if
    if (present (recoil)) then
       data%recoil = recoil
    end if
    data%real_mass = data%flv_in(1)%get_mass ()
    if (present (mass)) then
       if (mass > 0) then
          data%mass = mass
       else
          data%mass = data%real_mass
          if (any (data%flv_in%get_mass () /= data%mass)) then
             data%error = MASS_MIX;  return
          end if
       end if
    else
       data%mass = data%real_mass
       if (any (data%flv_in%get_mass () /= data%mass)) then
          data%error = MASS_MIX;  return
       end if
    end if
    if (data%mass == 0) then
       data%error = ZERO_MASS;  return
    else if (data%mass >= data%q_max) then
       data%error = Q_MAX_TOO_SMALL;  return
    end if
    data%log = log (1 + (data%q_max / data%mass)**2)
    charge = data%flv_in(1)%get_charge ()
    if (any (abs (data%flv_in%get_charge ()) /= abs (charge))) then
       data%error = CHARGE_MIX;  return
    else if (charge == 0) then
       data%error = CHARGE_ZERO;  return
    end if
    data%eps = data%alpha / pi * charge ** 2 &
         * (2 * log (data%q_max / data%mass) - 1)
    if (data%eps > 1) then
       data%error = EPS_TOO_LARGE;  return
    end if
  end subroutine isr_data_init

\eatline
\nwendcode{}\nwbegindocs{999}Explicitly set ISR order
\nwenddocs{}\nwbegincode{1000}\moddef{SF isr: isr data: TBP}\plusendmoddef
  procedure :: set_order => isr_data_set_order
\nwendcode{}\nwbegincode{1001}\moddef{SF isr: procedures}\plusendmoddef
  elemental subroutine isr_data_set_order (data, order)
    class(isr_data_t), intent(inout) :: data
    integer, intent(in) :: order
    if (order < 0 .or. order > 3) then
       data%error = INVALID_ORDER
    else
       data%order = order
    end if
  end subroutine isr_data_set_order

\eatline
\nwendcode{}\nwbegindocs{1002}Handle error conditions.  Should always be done after
initialization, unless we are sure everything is ok.
\nwenddocs{}\nwbegincode{1003}\moddef{SF isr: isr data: TBP}\plusendmoddef
  procedure :: check => isr_data_check
\nwendcode{}\nwbegincode{1004}\moddef{SF isr: procedures}\plusendmoddef
  subroutine isr_data_check (data)
    class(isr_data_t), intent(in) :: data
    select case (data%error)
    case (ZERO_MASS)
       call msg_fatal ("ISR: Particle mass is zero")
    case (Q_MAX_TOO_SMALL)
       call msg_fatal ("ISR: Particle mass exceeds Qmax")
    case (EPS_TOO_LARGE)
       call msg_fatal ("ISR: Expansion parameter too large, " // &
            "perturbative expansion breaks down")
    case (INVALID_ORDER)
       call msg_error ("ISR: LLA order invalid (valid values are 0,1,2,3)")
    case (MASS_MIX)
       call msg_fatal ("ISR: Incoming particle masses must be uniform")
    case (CHARGE_MIX)
       call msg_fatal ("ISR: Incoming particle charges must be uniform")
    case (CHARGE_ZERO)
       call msg_fatal ("ISR: Incoming particle must be charged")
    end select
  end subroutine isr_data_check

\eatline
\nwendcode{}\nwbegindocs{1005}Output
\nwenddocs{}\nwbegincode{1006}\moddef{SF isr: isr data: TBP}\plusendmoddef
  procedure :: write => isr_data_write
\nwendcode{}\nwbegincode{1007}\moddef{SF isr: procedures}\plusendmoddef
  subroutine isr_data_write (data, unit, verbose) 
    class(isr_data_t), intent(in) :: data
    integer, intent(in), optional :: unit
    logical, intent(in), optional :: verbose
    integer :: u, i
    u = given_output_unit (unit);  if (u < 0)  return
    write (u, "(1x,A)") "ISR data:"
    if (allocated (data%flv_in)) then
       write (u, "(3x,A)", advance="no") "  flavor =  "
       do i = 1, size (data%flv_in)
          if (i > 1)  write (u, "(',',1x)", advance="no")
          call data%flv_in(i)%write (u)
       end do
       write (u, *)    
       write (u, "(3x,A," // FMT_19 // ")") "  alpha  = ", data%alpha
       write (u, "(3x,A," // FMT_19 // ")") "  q_max  = ", data%q_max
       write (u, "(3x,A," // FMT_19 // ")") "  mass   = ", data%mass
       write (u, "(3x,A," // FMT_19 // ")") "  eps    = ", data%eps
       write (u, "(3x,A," // FMT_19 // ")") "  log    = ", data%log
       write (u, "(3x,A,I2)")      "  order  = ", data%order
       write (u, "(3x,A,L2)")      "  recoil = ", data%recoil
    else
       write (u, "(3x,A)") "[undefined]"       
    end if
  end subroutine isr_data_write

\eatline
\nwendcode{}\nwbegindocs{1008}For ISR, there is the option to generate transverse momentum is
generated. Hence, there can be up to three parameters, $x$, and two
angles. 
\nwenddocs{}\nwbegincode{1009}\moddef{SF isr: isr data: TBP}\plusendmoddef
  procedure :: get_n_par => isr_data_get_n_par
\nwendcode{}\nwbegincode{1010}\moddef{SF isr: procedures}\plusendmoddef
  function isr_data_get_n_par (data) result (n)
    class(isr_data_t), intent(in) :: data
    integer :: n
    if (data%recoil) then
       n = 3
    else
       n = 1
    end if
  end function isr_data_get_n_par
  
\eatline
\nwendcode{}\nwbegindocs{1011}Return the outgoing particles PDG codes.  For ISR, these are
identical to the incoming particles.
\nwenddocs{}\nwbegincode{1012}\moddef{SF isr: isr data: TBP}\plusendmoddef
  procedure :: get_pdg_out => isr_data_get_pdg_out
\nwendcode{}\nwbegincode{1013}\moddef{SF isr: procedures}\plusendmoddef
  subroutine isr_data_get_pdg_out (data, pdg_out)
    class(isr_data_t), intent(in) :: data
    type(pdg_array_t), dimension(:), intent(inout) :: pdg_out
    pdg_out(1) = data%flv_in%get_pdg ()
  end subroutine isr_data_get_pdg_out
  
\eatline
\nwendcode{}\nwbegindocs{1014}Return the {\tt{}eps} value.  We need it for an appropriate mapping of
structure-function parameters.
\nwenddocs{}\nwbegincode{1015}\moddef{SF isr: isr data: TBP}\plusendmoddef
  procedure :: get_eps => isr_data_get_eps
\nwendcode{}\nwbegincode{1016}\moddef{SF isr: procedures}\plusendmoddef
  function isr_data_get_eps (data) result (eps)
    class(isr_data_t), intent(in) :: data
    real(default) :: eps
    eps = data%eps
  end function isr_data_get_eps

\eatline
\nwendcode{}\nwbegindocs{1017}Allocate the interaction record.
\nwenddocs{}\nwbegincode{1018}\moddef{SF isr: isr data: TBP}\plusendmoddef
  procedure :: allocate_sf_int => isr_data_allocate_sf_int
\nwendcode{}\nwbegincode{1019}\moddef{SF isr: procedures}\plusendmoddef
  subroutine isr_data_allocate_sf_int (data, sf_int)
    class(isr_data_t), intent(in) :: data
    class(sf_int_t), intent(inout), allocatable :: sf_int
    allocate (isr_t :: sf_int)
  end subroutine isr_data_allocate_sf_int
  
\eatline
\nwendcode{}\nwbegindocs{1020}\nwdocspar
\subsection{The ISR object}
The {\tt{}isr{\char95}t} data type is a $1\to 2$ interaction, i.e., we allow for
single-photon emission only (but use the multi-photon resummed
radiator function).  The particles are ordered as (incoming, photon,
outgoing).

There is no need to handle several flavors (and data blocks) in
parallel, since ISR is always applied immediately after beam
collision.  (ISR for partons is accounted for by the PDFs themselves.)
Polarization is carried through, i.e., we retain the polarization of
the incoming particle and treat the emitted photon as unpolarized.
Color is trivially carried through.  This implies that particles 1 and
3 should be locked together. For ISR we don't need the q variable.
\nwenddocs{}\nwbegincode{1021}\moddef{SF isr: types}\plusendmoddef
  type, extends (sf_int_t) :: isr_t
     type(isr_data_t), pointer :: data => null ()
     real(default) :: x = 0
     real(default) :: xb= 0
   contains
   \LA{}SF isr: isr: TBP\RA{}
  end type isr_t 
  
\eatline
\nwendcode{}\nwbegindocs{1022}Type string: has to be here, but there is no string variable on which ISR
depends. Hence, a dummy routine.
\nwenddocs{}\nwbegincode{1023}\moddef{SF isr: isr: TBP}\endmoddef
  procedure :: type_string => isr_type_string
\nwendcode{}\nwbegincode{1024}\moddef{SF isr: procedures}\plusendmoddef
  function isr_type_string (object) result (string)
    class(isr_t), intent(in) :: object
    type(string_t) :: string
    if (associated (object%data)) then
       string = "ISR: e+ e- ISR spectrum" 
    else
       string = "ISR: [undefined]"
    end if
  end function isr_type_string
  
\eatline
\nwendcode{}\nwbegindocs{1025}Output.  Call the interaction routine after displaying the configuration.
\nwenddocs{}\nwbegincode{1026}\moddef{SF isr: isr: TBP}\plusendmoddef
  procedure :: write => isr_write
\nwendcode{}\nwbegincode{1027}\moddef{SF isr: procedures}\plusendmoddef
  subroutine isr_write (object, unit, testflag)
    class(isr_t), intent(in) :: object
    integer, intent(in), optional :: unit
    logical, intent(in), optional :: testflag
    integer :: u
    u = given_output_unit (unit)
    if (associated (object%data)) then
       call object%data%write (u)
       if (object%status >= SF_DONE_KINEMATICS) then
          write (u, "(1x,A)")  "SF parameters:"
          write (u, "(3x,A," // FMT_17 // ")")  "x =", object%x
          write (u, "(3x,A," // FMT_17 // ")")  "xb=", object%xb
       end if
       call object%base_write (u, testflag)
    else
       write (u, "(1x,A)")  "ISR data: [undefined]"
    end if
  end subroutine isr_write
    
\eatline
\nwendcode{}\nwbegindocs{1028}\nwdocspar
\subsection{Kinematics}
Set kinematics.  If {\tt{}map} is unset, the $r$ and $x$ values
coincide, and the Jacobian $f(r)$ were trivial. The ISR structure
function allows for a straightforward mapping of the unit interval.
So, to leading order, the structure function value is unity, but the
$x$ value is transformed.  Higher orders affect the function value.

The structure function implementation applies the above mapping to the
input (random) number {\tt{}r} to generate the momentum fraction {\tt{}x}
and the function value {\tt{}f}.  For numerical stability reasons, we
also output {\tt{}xb}, which is $\bar x=1-x$.

For the ISR structure function, the mapping Jacobian cancels the
structure function (to order zero).  We apply the cancellation
explicitly, therefore both the Jacobian {\tt{}f} and the zeroth-order value
(see the {\tt{}apply} method) are unity if mapping is turned on.  If
mapping is turned off, the Jacobian {\tt{}f} includes the value of the
(zeroth-order) structure function, and strongly peaked.
\nwenddocs{}\nwbegincode{1029}\moddef{SF isr: isr: TBP}\plusendmoddef
  procedure :: complete_kinematics => isr_complete_kinematics
\nwendcode{}\nwbegincode{1030}\moddef{SF isr: procedures}\plusendmoddef
  subroutine isr_complete_kinematics (sf_int, x, f, r, rb, map)
    class(isr_t), intent(inout) :: sf_int
    real(default), dimension(:), intent(out) :: x
    real(default), intent(out) :: f
    real(default), dimension(:), intent(in) :: r
    real(default), dimension(:), intent(in) :: rb
    logical, intent(in) :: map
    real(default) :: eps
    eps = sf_int%data%eps
    if (map) then
       call map_power_1 (sf_int%xb, f, rb(1), eps)
    else
       sf_int%xb = rb(1)
       if (rb(1) > 0) then
          f = 1
       else
          f = 0
       end if
    end if
    sf_int%x = 1 - sf_int%xb
    x(1) = sf_int%x
    if (size (x) == 3)  x(2:3) = r(2:3)
    call sf_int%split_momentum (x, sf_int%xb)
    select case (sf_int%status)
    case (SF_FAILED_KINEMATICS)
       sf_int%x = 0
       sf_int%xb= 0
       f = 0
    end select
  end subroutine isr_complete_kinematics

\eatline
\nwendcode{}\nwbegindocs{1031}Overriding the default method: we compute the {\tt{}x} array from the
momentum configuration.  In the specific case of ISR, we also set the
internally stored $x$ and $\bar x$ values, so they can be used in the
following routine.

Note: the extraction of $\bar x$ is not numerically safe, but it can't
be as long as the base {\tt{}recover{\char95}x} isn't.
\nwenddocs{}\nwbegincode{1032}\moddef{SF isr: isr: TBP}\plusendmoddef
  procedure :: recover_x => sf_isr_recover_x
\nwendcode{}\nwbegincode{1033}\moddef{SF isr: procedures}\plusendmoddef
  subroutine sf_isr_recover_x (sf_int, x, x_free)
    class(isr_t), intent(inout) :: sf_int
    real(default), dimension(:), intent(out) :: x
    real(default), intent(inout), optional :: x_free
    call sf_int%base_recover_x (x, x_free)
    sf_int%x  = x(1)
    sf_int%xb = 1 - x(1)
  end subroutine sf_isr_recover_x
  
\eatline
\nwendcode{}\nwbegindocs{1034}Compute inverse kinematics.  Here, we start with the $x$ array and
compute the ``input'' $r$ values and the Jacobian $f$.  After this, we
can set momenta by the same formula as for normal kinematics.

For extracting $x$, we rely on the stored $\bar x$ value, since the
$x$ value in the argument is likely imprecise.  This means that either
{\tt{}complete{\char95}kinematics} or {\tt{}recover{\char95}x} must be called first, for the
current sampling point (but maybe another channel).
\nwenddocs{}\nwbegincode{1035}\moddef{SF isr: isr: TBP}\plusendmoddef
  procedure :: inverse_kinematics => isr_inverse_kinematics
\nwendcode{}\nwbegincode{1036}\moddef{SF isr: procedures}\plusendmoddef
  subroutine isr_inverse_kinematics (sf_int, x, f, r, rb, map, set_momenta)
    class(isr_t), intent(inout) :: sf_int
    real(default), dimension(:), intent(in) :: x
    real(default), intent(out) :: f
    real(default), dimension(:), intent(out) :: r
    real(default), dimension(:), intent(out) :: rb
    logical, intent(in) :: map
    logical, intent(in), optional :: set_momenta
    real(default) :: eps
    logical :: set_mom
    set_mom = .false.;  if (present (set_momenta))  set_mom = set_momenta
    eps = sf_int%data%eps
    if (map) then
       call map_power_inverse_1 (sf_int%xb, f, rb(1), eps)
    else
       rb(1) = sf_int%xb
       if (rb(1) > 0) then
          f = 1
       else
          f = 0
       end if
    end if
    r(1) = 1 - rb(1)
    if (size(r) == 3) then
       r(2:3) = x(2:3)
       rb(2:3)= 1 - r(2:3)
    end if
    if (set_mom) then
       call sf_int%split_momentum (x, sf_int%xb)
       select case (sf_int%status)
       case (SF_FAILED_KINEMATICS)
          r = 0
          rb= 0
          f = 0
       end select
    end if
  end subroutine isr_inverse_kinematics

\eatline
\nwendcode{}\nwbegindocs{1037}\nwdocspar
\nwenddocs{}\nwbegincode{1038}\moddef{SF isr: isr: TBP}\plusendmoddef
  procedure :: init => isr_init
\nwendcode{}\nwbegincode{1039}\moddef{SF isr: procedures}\plusendmoddef
  subroutine isr_init (sf_int, data)
    class(isr_t), intent(out) :: sf_int
    class(sf_data_t), intent(in), target :: data
    type(quantum_numbers_mask_t), dimension(3) :: mask
    integer, dimension(3) :: hel_lock
    type(polarization_t) :: pol
    type(quantum_numbers_t), dimension(1) :: qn_fc, qn_hel
    type(flavor_t) :: flv_photon
    type(color_t) :: col_photon
    type(quantum_numbers_t) :: qn_photon, qn
    type(state_iterator_t) :: it_hel
    real(default) :: m2
    integer :: i
    mask = quantum_numbers_mask (.false., .false., &
         mask_h = [.false., .true., .false.])
    hel_lock = [3, 0, 1]
    select type (data)
    type is (isr_data_t)   
       m2 = data%mass**2
       call sf_int%base_init (mask, [m2], [0._default], [m2], &
            hel_lock = hel_lock)
       sf_int%data => data              
       call flv_photon%init (PHOTON, data%model)
       call col_photon%init ()
       call qn_photon%init (flv_photon, col_photon)
       call qn_photon%tag_radiated ()
       do i = 1, size (data%flv_in)
          call polarization_init_generic (pol, data%flv_in(i))
          call qn_fc(1)%init (&
               flv = data%flv_in(i), &
               col = color_from_flavor (data%flv_in(i), 1))
          call it_hel%init (pol%state)
          do while (it_hel%is_valid ())
             qn_hel = it_hel%get_quantum_numbers ()
             qn = qn_hel(1) .merge. qn_fc(1)
             call interaction_add_state &
                  (sf_int%interaction_t, [qn, qn_photon, qn])
             call it_hel%advance ()
          end do
          call polarization_final (pol)
       end do
       call interaction_freeze (sf_int%interaction_t)
       call sf_int%set_incoming ([1])
       call sf_int%set_radiated ([2])
       call sf_int%set_outgoing ([3])
       sf_int%status = SF_INITIAL
    end select
  end subroutine isr_init
    
\eatline
\nwendcode{}\nwbegindocs{1040}\nwdocspar
\subsection{ISR application}
For ISR, we could in principle compute kinematics and function value
in a single step. In order to be able to reweight matrix elements
including structure functions we split kinematics and structure
function calculation. The structure function works on a single beam,
assuming that the input momentum has been set. 

For the structure-function evaluation, we rely on the fact that the
power mapping, which we apply in the kinematics method (if the {\tt{}map}
flag is set), has a Jacobian which is just the inverse lowest-order
structure function.  With mapping active, the two should cancel
exactly.

After splitting momenta, we set the outgoing momenta on-shell.  We
choose to conserve momentum, so energy conservation may be violated.
\nwenddocs{}\nwbegincode{1041}\moddef{SF isr: isr: TBP}\plusendmoddef
  procedure :: apply => isr_apply
\nwendcode{}\nwbegincode{1042}\moddef{SF isr: procedures}\plusendmoddef
  subroutine isr_apply (sf_int, scale)
    class(isr_t), intent(inout) :: sf_int
    real(default), intent(in) :: scale
    real(default) :: f, finv, x, xb, eps, rb
    real(default) :: log_x, log_xb, x_2
    real(default), parameter :: &
         & xmin = 0.00714053329734592839549879772019_default
    real(default), parameter :: &
         & zeta3 = 1.20205690315959428539973816151_default
    real(default), parameter :: &
         & g1 = 3._default / 4._default, &
         & g2 = (27 - 8*pi**2) / 96._default, &
         & g3 = (27 - 24*pi**2 + 128*zeta3) / 384._default    
    associate (data => sf_int%data)
      eps = sf_int%data%eps
      x = sf_int%x
      xb = sf_int%xb
      call map_power_inverse_1 (xb, finv, rb, eps)
      if (finv > 0) then
         f = 1 / finv
      else
         f = 0
      end if   
      if (f > 0 .and. data%order > 0) then
         f = f * (1 + g1 * eps)
         x_2 = x*x
         if (rb>0)  f = f * (1 - (1-x_2) / (2 * rb))
         if (data%order > 1) then
            f = f * (1 + g2 * eps**2)
            if (rb>0 .and. xb>0 .and. x>xmin) then
               log_x  = log_prec (x, xb)
               log_xb = log_prec (xb, x)
               f = f * (1 - ((1+3*x_2)*log_x + xb * (4*(1+x)*log_xb + 5 + x)) &
                    / ( 8 * rb) * eps)
            end if
            if (data%order > 2) then
               f = f * (1 + g3 * eps**3)
               if (rb > 0 .and. xb > 0 .and. x > xmin) then
                  f = f * (1 - ((1+x) * xb &
                       * (6 * Li2(x) + 12 * log_xb**2 - 3 * pi**2) &
                       + 1.5_default * (1 + 8*x + 3*x_2) * log_x &
                       + 6 * (x+5) * xb * log_xb &
                       + 12 * (1+x_2) * log_x * log_xb &
                       - (1 + 7*x_2) * log_x**2 / 2 &
                       + (39 - 24*x - 15*x_2) / 4) &
                       / ( 48 * rb) * eps**2)
               end if
            end if
         end if
      end if
    end associate
    call interaction_set_matrix_element &
           (sf_int%interaction_t, cmplx (f, kind=default))    
    sf_int%status = SF_EVALUATED
  end subroutine isr_apply

\eatline
\nwendcode{}\nwbegindocs{1043}\nwdocspar
\subsection{Unit tests}
\nwenddocs{}\nwbegincode{1044}\moddef{SF isr: public}\plusendmoddef
  public :: sf_isr_test
\nwendcode{}\nwbegincode{1045}\moddef{SF isr: tests}\endmoddef
  subroutine sf_isr_test (u, results)
    integer, intent(in) :: u
    type(test_results_t), intent(inout) :: results
  \LA{}SF isr: execute tests\RA{}
  end subroutine sf_isr_test
  
\eatline
\nwendcode{}\nwbegindocs{1046}\nwdocspar
\subsubsection{Test structure function data}
Construct and display a test structure function data object.
\nwenddocs{}\nwbegincode{1047}\moddef{SF isr: execute tests}\endmoddef
  call test (sf_isr_1, "sf_isr_1", &
       "structure function configuration", &
       u, results)
\nwendcode{}\nwbegincode{1048}\moddef{SF isr: tests}\plusendmoddef
  subroutine sf_isr_1 (u)
    integer, intent(in) :: u
    type(model_data_t), target :: model
    type(pdg_array_t) :: pdg_in
    type(pdg_array_t), dimension(1) :: pdg_out
    integer, dimension(:), allocatable :: pdg1
    class(sf_data_t), allocatable :: data
    
    write (u, "(A)")  "* Test output: sf_isr_1"
    write (u, "(A)")  "*   Purpose: initialize and display &
         &test structure function data"
    write (u, "(A)")
    
    write (u, "(A)")  "* Create empty data object"
    write (u, "(A)")

    call model%init_qed_test ()
    pdg_in = ELECTRON

    allocate (isr_data_t :: data)
    call data%write (u)

    write (u, "(A)")
    write (u, "(A)")  "* Initialize"
    write (u, "(A)")

    select type (data)
    type is (isr_data_t)
       call data%init (model, pdg_in, 1./137._default, 10._default, &
            0.000511_default, order = 3, recoil = .false.)
    end select

    call data%write (u)

    write (u, "(A)")

    write (u, "(1x,A)")  "Outgoing particle codes:"
    call data%get_pdg_out (pdg_out)
    pdg1 = pdg_out(1)
    write (u, "(2x,99(1x,I0))")  pdg1
        
    call model%final ()
    
    write (u, "(A)")
    write (u, "(A)")  "* Test output end: sf_isr_1"

  end subroutine sf_isr_1

\eatline
\nwendcode{}\nwbegindocs{1049}\nwdocspar
\subsubsection{Structure function without mapping}
Direct ISR evaluation.  This is the use case for a double-beam
structure function.  The parameter pair is mapped in the calling program.
\nwenddocs{}\nwbegincode{1050}\moddef{SF isr: execute tests}\plusendmoddef
  call test (sf_isr_2, "sf_isr_2", &
       "no ISR mapping", &
       u, results)
\nwendcode{}\nwbegincode{1051}\moddef{SF isr: tests}\plusendmoddef
  subroutine sf_isr_2 (u)
    integer, intent(in) :: u
    type(model_data_t), target :: model
    type(pdg_array_t) :: pdg_in
    type(flavor_t) :: flv
    class(sf_data_t), allocatable, target :: data
    class(sf_int_t), allocatable :: sf_int
    type(vector4_t) :: k
    real(default) :: E
    real(default), dimension(:), allocatable :: r, rb, x
    real(default) :: f, f_isr
    
    write (u, "(A)")  "* Test output: sf_isr_2"
    write (u, "(A)")  "*   Purpose: initialize and fill &
         &test structure function object"
    write (u, "(A)")
    
    write (u, "(A)")  "* Initialize configuration data"
    write (u, "(A)")

    call model%init_qed_test ()
    pdg_in = ELECTRON
    call flv%init (ELECTRON, model)

    call reset_interaction_counter ()
    
    allocate (isr_data_t :: data)
    select type (data)
    type is (isr_data_t)
       call data%init (model, pdg_in, 1./137._default, 500._default, &
            0.000511_default, order = 3, recoil = .false.)
    end select
       
    write (u, "(A)")  "* Initialize structure-function object"
    write (u, "(A)")
    
    call data%allocate_sf_int (sf_int)
    call sf_int%init (data)
    call sf_int%set_beam_index ([1])

    write (u, "(A)")  "* Initialize incoming momentum with E=500"
    write (u, "(A)")
    E = 500
    k = vector4_moving (E, sqrt (E**2 - flv%get_mass ()**2), 3)
    call pacify (k, 1e-10_default)
    call vector4_write (k, u)
    call sf_int%seed_kinematics ([k])
    
    write (u, "(A)")
    write (u, "(A)")  "* Set kinematics for r=0.9, no ISR mapping, &
         &collinear"
    write (u, "(A)")
    
    allocate (r (data%get_n_par ()))
    allocate (rb(size (r)))
    allocate (x (size (r)))
    
    r = 0.9_default
    rb = 1 - r
    write (u, "(A,9(1x," // FMT_12 // "))")  "r =", r
    write (u, "(A,9(1x," // FMT_12 // "))")  "rb=", rb

    call sf_int%complete_kinematics (x, f, r, rb, map=.false.)
    
    write (u, "(A)")
    write (u, "(A,9(1x," // FMT_12 // "))")  "x =", x
    write (u, "(A,9(1x," // FMT_12 // "))")  "f =", f
    
    write (u, "(A)")
    write (u, "(A)")  "* Invert kinematics"
    write (u, "(A)")
    
    call sf_int%inverse_kinematics (x, f, r, rb, map=.false.) 
    write (u, "(A,9(1x," // FMT_12 // "))")  "r =", r
    write (u, "(A,9(1x," // FMT_12 // "))")  "rb=", rb
    write (u, "(A,9(1x," // FMT_12 // "))")  "f =", f

    write (u, "(A)")
    write (u, "(A)")  "* Evaluate ISR structure function"
    write (u, "(A)")
    
    call sf_int%apply (scale = 100._default)
    call sf_int%write (u)
    
    write (u, "(A)")
    write (u, "(A)")  "* Structure-function value, default order"
    write (u, "(A)")

    f_isr = interaction_get_matrix_element (sf_int%interaction_t, 1)
    
    write (u, "(A,9(1x," // FMT_12 // "))")  "f_isr         =", f_isr
    write (u, "(A,9(1x," // FMT_12 // "))")  "f_isr * f_map =", f_isr * f

    write (u, "(A)")
    write (u, "(A)")  "* Re-evaluate structure function, leading order"
    write (u, "(A)")
    
    select type (sf_int)
    type is (isr_t)
       sf_int%data%order = 0
    end select
    call sf_int%apply (scale = 100._default)
    f_isr = interaction_get_matrix_element (sf_int%interaction_t, 1)
    
    write (u, "(A,9(1x," // FMT_12 // "))")  "f_isr         =", f_isr
    write (u, "(A,9(1x," // FMT_12 // "))")  "f_isr * f_map =", f_isr * f

    write (u, "(A)")
    write (u, "(A)")  "* Cleanup"
    
    call sf_int%final ()
    call model%final ()
    
    write (u, "(A)")
    write (u, "(A)")  "* Test output end: sf_isr_2"

  end subroutine sf_isr_2

\eatline
\nwendcode{}\nwbegindocs{1052}\nwdocspar
\subsubsection{Structure function with mapping}
Apply the optimal ISR mapping.  This is the use case for a single-beam
structure function.
\nwenddocs{}\nwbegincode{1053}\moddef{SF isr: execute tests}\plusendmoddef
  call test (sf_isr_3, "sf_isr_3", &
       "ISR mapping", &
       u, results)
\nwendcode{}\nwbegincode{1054}\moddef{SF isr: tests}\plusendmoddef
  subroutine sf_isr_3 (u)
    integer, intent(in) :: u
    type(model_data_t), target :: model
    type(flavor_t) :: flv
    type(pdg_array_t) :: pdg_in
    class(sf_data_t), allocatable, target :: data
    class(sf_int_t), allocatable :: sf_int
    type(vector4_t) :: k
    real(default) :: E
    real(default), dimension(:), allocatable :: r, rb, x
    real(default) :: f, f_isr
    
    write (u, "(A)")  "* Test output: sf_isr_3"
    write (u, "(A)")  "*   Purpose: initialize and fill &
         &test structure function object"
    write (u, "(A)")
    
    write (u, "(A)")  "* Initialize configuration data"
    write (u, "(A)")

    call model%init_qed_test ()
    call flv%init (ELECTRON, model)
    pdg_in = ELECTRON

    call reset_interaction_counter ()
    
    allocate (isr_data_t :: data)
    select type (data)
    type is (isr_data_t)
       call data%init (model, pdg_in, 1./137._default, 500._default, &
            0.000511_default, order = 3, recoil = .false.)
    end select
       
    write (u, "(A)")  "* Initialize structure-function object"
    write (u, "(A)")
    
    call data%allocate_sf_int (sf_int)
    call sf_int%init (data)
    call sf_int%set_beam_index ([1])

    write (u, "(A)")  "* Initialize incoming momentum with E=500"
    write (u, "(A)")
    E = 500
    k = vector4_moving (E, sqrt (E**2 - flv%get_mass ()**2), 3)
    call pacify (k, 1e-10_default)
    call vector4_write (k, u)
    call sf_int%seed_kinematics ([k])
    
    write (u, "(A)")
    write (u, "(A)")  "* Set kinematics for r=0.7, with ISR mapping, &
         &collinear"
    write (u, "(A)")
    
    allocate (r (data%get_n_par ()))
    allocate (rb(size (r)))
    allocate (x (size (r)))
    
    r = 0.7_default
    rb = 1 - r
    write (u, "(A,9(1x," // FMT_12 // "))")  "r =", r
    write (u, "(A,9(1x," // FMT_12 // "))")  "rb=", rb

    call sf_int%complete_kinematics (x, f, r, rb, map=.true.)
    
    write (u, "(A)")
    write (u, "(A,9(1x," // FMT_12 // "))")  "x =", x
    write (u, "(A,9(1x," // FMT_12 // "))")  "f =", f
    
    write (u, "(A)")
    write (u, "(A)")  "* Invert kinematics"
    write (u, "(A)")
    
    call sf_int%inverse_kinematics (x, f, r, rb, map=.true.) 
    write (u, "(A,9(1x," // FMT_12 // "))")  "r =", r
    write (u, "(A,9(1x," // FMT_12 // "))")  "rb=", rb
    write (u, "(A,9(1x," // FMT_12 // "))")  "f =", f

    write (u, "(A)")
    write (u, "(A)")  "* Evaluate ISR structure function"
    write (u, "(A)")
    
    call sf_int%apply (scale = 100._default)
    call sf_int%write (u)
    
    write (u, "(A)")
    write (u, "(A)")  "* Structure-function value, default order"
    write (u, "(A)")

    f_isr = interaction_get_matrix_element (sf_int%interaction_t, 1)
    
    write (u, "(A,9(1x," // FMT_12 // "))")  "f_isr         =", f_isr
    write (u, "(A,9(1x," // FMT_12 // "))")  "f_isr * f_map =", f_isr * f

    write (u, "(A)")
    write (u, "(A)")  "* Re-evaluate structure function, leading order"
    write (u, "(A)")
    
    select type (sf_int)
    type is (isr_t)
       sf_int%data%order = 0
    end select
    call sf_int%apply (scale = 100._default)
    f_isr = interaction_get_matrix_element (sf_int%interaction_t, 1)
    
    write (u, "(A,9(1x," // FMT_12 // "))")  "f_isr         =", f_isr
    write (u, "(A,9(1x," // FMT_12 // "))")  "f_isr * f_map =", f_isr * f

    write (u, "(A)")
    write (u, "(A)")  "* Cleanup"
    
    call sf_int%final ()
    call model%final ()
    
    write (u, "(A)")
    write (u, "(A)")  "* Test output end: sf_isr_3"

  end subroutine sf_isr_3

\eatline
\nwendcode{}\nwbegindocs{1055}\nwdocspar
\subsubsection{Non-collinear ISR splitting}
Construct and display a structure function object based on the ISR
structure function. We blank out numerical fluctuations for 32bit.
\nwenddocs{}\nwbegincode{1056}\moddef{SF isr: execute tests}\plusendmoddef
  call test (sf_isr_4, "sf_isr_4", &
       "ISR non-collinear", &
       u, results)
\nwendcode{}\nwbegincode{1057}\moddef{SF isr: tests}\plusendmoddef
  subroutine sf_isr_4 (u)
    integer, intent(in) :: u
    type(model_data_t), target :: model
    type(flavor_t) :: flv
    type(pdg_array_t) :: pdg_in
    class(sf_data_t), allocatable, target :: data
    class(sf_int_t), allocatable :: sf_int
    type(vector4_t) :: k
    type(vector4_t), dimension(2) :: q
    real(default) :: E
    real(default), dimension(:), allocatable :: r, rb, x
    real(default) :: f, f_isr
    character(len=80) :: buffer
    integer :: u_scratch, iostat
    
    write (u, "(A)")  "* Test output: sf_isr_4"
    write (u, "(A)")  "*   Purpose: initialize and fill &
         &test structure function object"
    write (u, "(A)")
    
    write (u, "(A)")  "* Initialize configuration data"
    write (u, "(A)")

    call model%init_qed_test ()
    call flv%init (ELECTRON, model)
    pdg_in = ELECTRON

    call reset_interaction_counter ()
    
    write (u, "(A)")
    write (u, "(A)")  "* Initialize structure-function object"
    write (u, "(A)")

    allocate (isr_data_t :: data)
    select type (data)
    type is (isr_data_t)
       call data%init (model, pdg_in, 1./137._default, 500._default, &
            0.000511_default, order = 3, recoil = .true.)
    end select    
    
    call data%allocate_sf_int (sf_int)
    call sf_int%init (data)
    call sf_int%set_beam_index ([1])

    write (u, "(A)")
    write (u, "(A)")  "* Initialize incoming momentum with E=500"
    write (u, "(A)")
    E = 500
    k = vector4_moving (E, sqrt (E**2 - flv%get_mass ()**2), 3)
    call pacify (k, 1e-10_default)
    call vector4_write (k, u)
    call sf_int%seed_kinematics ([k])
        
    write (u, "(A)")
    write (u, "(A)")  "* Set kinematics for x=0.5/0.5/0.25, with ISR mapping, "
    write (u, "(A)")  "          non-coll., keeping energy"
    write (u, "(A)")
    
    allocate (r (data%get_n_par ()))
    allocate (rb(size (r)))
    allocate (x (size (r)))
    
    r = [0.5_default, 0.5_default, 0.25_default]
    rb = 1 - r
    sf_int%on_shell_mode = KEEP_ENERGY
    call sf_int%complete_kinematics (x, f, r, rb, map=.true.)
    call interaction_pacify_momenta (sf_int%interaction_t, 1e-10_default)
    
    write (u, "(A,9(1x,F10.7))")  "x =", x
    write (u, "(A,9(1x,F10.7))")  "f =", f
    
    write (u, "(A)")
    write (u, "(A)")  "* Recover x and r from momenta"
    write (u, "(A)")
    
    q = interaction_get_momenta (sf_int%interaction_t, outgoing=.true.)
    call sf_int%final ()
    deallocate (sf_int)
    
    call data%allocate_sf_int (sf_int)
    call sf_int%init (data)
    call sf_int%set_beam_index ([1])
    
    call sf_int%seed_kinematics ([k])
    call interaction_set_momenta (sf_int%interaction_t, q, outgoing=.true.)
    call sf_int%recover_x (x)
    call sf_int%inverse_kinematics (x, f, r, rb, map=.true.)    
    
    write (u, "(A,9(1x,F10.7))")  "x =", x
    write (u, "(A,9(1x,F10.7))")  "r =", r
    
    write (u, "(A)")
    write (u, "(A)")  "* Evaluate ISR structure function"
    write (u, "(A)")
    
    call sf_int%complete_kinematics (x, f, r, rb, map=.true.) 
    call interaction_pacify_momenta (sf_int%interaction_t, 1e-10_default)    
    call sf_int%apply (scale = 10._default)
    u_scratch = free_unit ()
    open (u_scratch, status="scratch", action = "readwrite")
    call sf_int%write (u_scratch, testflag = .true.)
    rewind (u_scratch)
    do 
       read (u_scratch, "(A)", iostat=iostat) buffer    
       if (iostat /= 0) exit
       if (buffer(1:27) == " P =   0.00000000E+00  9.57") then
          buffer = replace (buffer, 28, "XXXXXX")
       end if
       if (buffer(1:27) == " P =   0.00000000E+00 -9.57") then
          buffer = replace (buffer, 28, "XXXXXX")
       end if       
       write (u, "(A)") buffer
    end do
    close (u_scratch)
       
    write (u, "(A)")
    write (u, "(A)")  "* Structure-function value"
    write (u, "(A)")

    f_isr = interaction_get_matrix_element (sf_int%interaction_t, 1)
    
    write (u, "(A,9(1x," // FMT_12 // "))")  "f_isr         =", f_isr
    write (u, "(A,9(1x," // FMT_12 // "))")  "f_isr * f_map =", f_isr * f

    write (u, "(A)")
    write (u, "(A)")  "* Cleanup"
    
    call sf_int%final ()
    call model%final ()
    
    write (u, "(A)")
    write (u, "(A)")  "* Test output end: sf_isr_4"

  end subroutine sf_isr_4

\eatline
\nwendcode{}\nwbegindocs{1058}\nwdocspar
\subsubsection{Structure function pair with mapping}
Apply the ISR mapping for a ISR pair.
structure function.
\nwenddocs{}\nwbegincode{1059}\moddef{SF isr: execute tests}\plusendmoddef
  call test (sf_isr_5, "sf_isr_5", &
       "ISR pair mapping", &
       u, results)
\nwendcode{}\nwbegincode{1060}\moddef{SF isr: tests}\plusendmoddef
  subroutine sf_isr_5 (u)
    integer, intent(in) :: u
    type(model_data_t), target :: model
    type(flavor_t) :: flv
    type(pdg_array_t) :: pdg_in
    class(sf_data_t), allocatable, target :: data
    class(sf_mapping_t), allocatable :: mapping
    class(sf_int_t), dimension(:), allocatable :: sf_int
    type(vector4_t), dimension(2) :: k
    real(default) :: E, f_map
    real(default), dimension(:), allocatable :: p, pb, r, rb, x
    real(default), dimension(2) :: f, f_isr
    integer :: i
    
    write (u, "(A)")  "* Test output: sf_isr_5"
    write (u, "(A)")  "*   Purpose: initialize and fill &
         &test structure function object"
    write (u, "(A)")
    
    write (u, "(A)")  "* Initialize configuration data"
    write (u, "(A)")

    call model%init_qed_test ()
    call flv%init (ELECTRON, model)
    pdg_in = ELECTRON

    call reset_interaction_counter ()
    
    allocate (isr_data_t :: data)
    select type (data)
    type is (isr_data_t)
       call data%init (model, pdg_in, 1./137._default, 500._default, &
            0.000511_default, order = 3, recoil = .false.)
    end select
       
    allocate (sf_ip_mapping_t :: mapping)
    select type (mapping)
    type is (sf_ip_mapping_t)
       select type (data)
       type is (isr_data_t)
          call mapping%init (eps = data%eps)
       end select
       call mapping%set_index (1, 1)
       call mapping%set_index (2, 2)
    end select

    call mapping%write (u)
    
    write (u, "(A)")
    write (u, "(A)")  "* Initialize structure-function object"
    write (u, "(A)")

    allocate (isr_t :: sf_int (2))

    do i = 1, 2
       call sf_int(i)%init (data)
       call sf_int(i)%set_beam_index ([i])
    end do

    write (u, "(A)")  "* Initialize incoming momenta with E=500"
    write (u, "(A)")
    E = 500
    k(1) = vector4_moving (E,   sqrt (E**2 - flv%get_mass ()**2), 3)
    k(2) = vector4_moving (E, - sqrt (E**2 - flv%get_mass ()**2), 3)
    call pacify (k, 1e-10_default)
    do i = 1, 2
       call vector4_write (k(i), u)
       call sf_int(i)%seed_kinematics (k(i:i))
    end do
    
    write (u, "(A)")
    write (u, "(A)")  "* Set kinematics for p=[0.7,0.4], collinear"
    write (u, "(A)")
    
    allocate (p (2 * data%get_n_par ()))
    allocate (pb(size (p)))
    allocate (r (size (p)))
    allocate (rb(size (p)))
    allocate (x (size (p)))
    
    p = [0.7_default, 0.4_default]
    pb= 1 - p
    call mapping%compute (r, rb, f_map, p, pb)

    write (u, "(A,9(1x," // FMT_12 // "))")  "p =", p
    write (u, "(A,9(1x," // FMT_12 // "))")  "pb=", pb
    write (u, "(A,9(1x," // FMT_12 // "))")  "r =", r
    write (u, "(A,9(1x," // FMT_12 // "))")  "rb=", rb
    write (u, "(A,9(1x," // FMT_12 // "))")  "fm=", f_map

    do i = 1, 2
       call sf_int(i)%complete_kinematics (x(i:i), f(i), r(i:i), rb(i:i), &
            map=.false.)
    end do
    
    write (u, "(A)")
    write (u, "(A,9(1x," // FMT_12 // "))")  "x =", x
    write (u, "(A,9(1x," // FMT_12 // "))")  "f =", f
    
    write (u, "(A)")
    write (u, "(A)")  "* Invert kinematics"
    write (u, "(A)")
    
    do i = 1, 2
       call sf_int(i)%inverse_kinematics (x(i:i), f(i), r(i:i), rb(i:i), &
            map=.false.) 
    end do
    call mapping%inverse (r, rb, f_map, p, pb)

    write (u, "(A,9(1x," // FMT_12 // "))")  "p =", p
    write (u, "(A,9(1x," // FMT_12 // "))")  "pb=", pb
    write (u, "(A,9(1x," // FMT_12 // "))")  "r =", r
    write (u, "(A,9(1x," // FMT_12 // "))")  "rb=", rb
    write (u, "(A,9(1x," // FMT_12 // "))")  "fm=", f_map

    write (u, "(A)")
    write (u, "(A)")  "* Evaluate ISR structure function"
    
    call sf_int(1)%apply (scale = 100._default)
    call sf_int(2)%apply (scale = 100._default)

    write (u, "(A)")
    write (u, "(A)")  "* Structure function #1"
    write (u, "(A)")
    call sf_int(1)%write (u)
    
    write (u, "(A)")
    write (u, "(A)")  "* Structure function #2"
    write (u, "(A)")
    call sf_int(2)%write (u)
    
    write (u, "(A)")
    write (u, "(A)")  "* Structure-function value, default order"
    write (u, "(A)")

    do i = 1, 2
       f_isr(i) = interaction_get_matrix_element (sf_int(i)%interaction_t, 1)
    end do
    
    write (u, "(A,9(1x," // FMT_12 // "))")  "f_isr         =", &
         product (f_isr)
    write (u, "(A,9(1x," // FMT_12 // "))")  "f_isr * f_map =", &
         product (f_isr * f) * f_map

    write (u, "(A)")
    write (u, "(A)")  "* Cleanup"
    
    do i = 1, 2
       call sf_int(i)%final ()
    end do
    call model%final ()
    
    write (u, "(A)")
    write (u, "(A)")  "* Test output end: sf_isr_5"

  end subroutine sf_isr_5

\eatline
\nwendcode{}\nwbegindocs{1061}\nwdocspar
\clearpage
%------------------------------------------------------------------------
\section{EPA}

\nwenddocs{}\nwbegincode{1062}\moddef{\code{}sf{\char95}epa.f90\edoc{}}\endmoddef
\LA{}File header\RA{}

module sf_epa

\LA{}Use kinds\RA{}
\LA{}Use strings\RA{}
  use io_units
  use constants, only: pi
  use format_defs, only: FMT_17, FMT_19
  use unit_tests
  use diagnostics
  use physics_defs, only: ELECTRON, PHOTON
  use lorentz
  use pdg_arrays
  use model_data
  use flavors
  use colors
  use quantum_numbers
  use state_matrices
  use polarizations
  use interactions
  use sf_aux
  use sf_base
  
\LA{}Standard module head\RA{}

\LA{}SF epa: public\RA{}

\LA{}SF epa: parameters\RA{}

\LA{}SF epa: types\RA{}

contains

\LA{}SF epa: procedures\RA{}

\LA{}SF epa: tests\RA{}

end module sf_epa
\eatline
\nwendcode{}\nwbegindocs{1063}\nwdocspar
\subsection{Physics}
The EPA structure function for a photon inside an (elementary)
particle $p$ with energy $E$, mass $m$ and charge $q_p$ (e.g.,
electron) is given by ($\bar x \equiv 1-x$)
%% %\cite{Budnev:1974de}
%% \bibitem{Budnev:1974de}
%%   V.~M.~Budnev, I.~F.~Ginzburg, G.~V.~Meledin and V.~G.~Serbo,
%%   %``The Two photon particle production mechanism. Physical problems.
%%   %Applications. Equivalent photon approximation,''
%%   Phys.\ Rept.\  {\bf 15} (1974) 181.
%%   %%CITATION = PRPLC,15,181;%%
\begin{multline}
\label{EPA}
  f(x) = 
  \frac{\alpha}{\pi}\,q_p^2\,
  \frac{1}{x}\, 
  \biggl[\left(\bar x + \frac{x^2}{2}\right) 
        \ln\frac{Q^2_{\rm max}}{Q^2_{\rm min}} 
        \\
        - \left(1 - \frac{x}{2}\right)^2 
          \ln\frac{x^2+\frac{Q^2_{\rm max}}{E^2}}
                  {x^2+\frac{Q^2_{\rm min}}{E^2}} 
        - x^2\frac{m^2}{Q^2_{\rm min}}
          \left(1 - \frac{Q^2_{\rm min}}{Q^2_{\rm max}}\right) 
  \biggr].
\end{multline}
If no explicit $Q$ bounds are provided, the kinematical bounds are
\begin{align}
  -Q^2_{\rm max} &= t_0 = -2\bar x(E^2+p\bar p) + 2m^2 \approx -4\bar x E^2,
\\
  -Q^2_{\rm min} &= t_1 = -2\bar x(E^2-p\bar p) + 2m^2 
  \approx
  -\frac{x^2}{\bar x}m^2.
\end{align}

The second and third terms in (\ref{EPA}) are negative definite (and
subleading).  Noting that $\bar x + x^2/2$ is bounded between
$1/2$ and $1$, we derive that $f(x)$ is always smaller than
\begin{equation}
  \bar f(x) = \frac{\alpha}{\pi}\,q_p^2\,\frac{L - 2\ln x}{x}
  \qquad\text{where}\qquad 
  L = \ln\frac{\min(4E_{\rm max}^2,Q^2_{\rm max})}{\max(m^2,Q_{\rm min}^2)},
\end{equation}
where we allow for explicit $Q$ bounds that narrow the kinematical range.
Therefore, we generate this distribution:
\begin{equation}\label{EPA-subst}
  \int_{x_0}^{x_1} dx\,\bar f(x) = C(x_0,x_1)\int_0^1 dx'
\end{equation}
We set
\begin{equation}\label{EPA-x(x')}
  \ln x = \frac12\left\{ L - \sqrt{L^2 - 4\left[ x'\ln x_1(L-\ln x_1)
        + \bar x'\ln x_0(L-\ln x_0) \right]} \right\}
\end{equation}
such that $x(0)=x_0$ and $x(1)=x_1$ and
\begin{equation}
  \frac{dx}{dx'} = \left(\frac{\alpha}{\pi} q_p^2 \right)^{-1}
        x\frac{C(x_0,x_1)}{L - 2\ln x}
\end{equation}
with
\begin{equation}
  C(x_0,x_1) = \frac{\alpha}{\pi} q_p^2\,\left[\ln x_1(L-\ln x_1) - \ln
        x_0(L-\ln x_0)\right]
\end{equation}
such that (\ref{EPA-subst}) is satisfied.  Finally, we have
\begin{equation}
  \int_{x_0}^{x_1} dx\,f(x) = C(x_0,x_1)\int_0^1 dx'\,
        \frac{f(x(x'))}{\bar f(x(x'))}
\end{equation}
where $x'$ is calculated from $x$ via (\ref{EPA-x(x')}). 

The structure of the mapping is most obvious from:
\begin{equation}
  x'(x) = \frac{\log x ( L - \log x) - \log x_0 (L - \log x_0)}
  {\log x_1 ( L - \log x_1) - \log x_0 (L - \log x_0)} \; .
\end{equation}



\subsection{The EPA data block}
The EPA parameters are: $\alpha$, $E_{\rm max}$, $m$, $Q_{\rm min}$, and
$x_{\rm min}$.  Instead of $m$ we can use the incoming particle PDG
code as input; from this we can deduce the mass and charge.

Internally we store in addition $C_{0/1} = \frac{\alpha}{\pi}q_e^2\ln
x_{0/1} (L - \ln x_{0/1})$, the c.m. energy squared and the incoming
particle mass.
\nwenddocs{}\nwbegincode{1064}\moddef{SF epa: public}\endmoddef
  public :: epa_data_t
\nwendcode{}\nwbegincode{1065}\moddef{SF epa: types}\endmoddef
  type, extends(sf_data_t) :: epa_data_t
     private
     class(model_data_t), pointer :: model => null ()
     type(flavor_t), dimension(:), allocatable :: flv_in
     real(default) :: alpha
     real(default) :: x_min
     real(default) :: x_max
     real(default) :: q_min
     real(default) :: q_max
     real(default) :: E_max
     real(default) :: mass
     real(default) :: log
     real(default) :: a
     real(default) :: c0
     real(default) :: c1
     real(default) :: dc
     integer :: error = NONE
     logical :: recoil = .false.
   contains
   \LA{}SF epa: epa data: TBP\RA{}  
  end type epa_data_t

\eatline
\nwendcode{}\nwbegindocs{1066}Error codes
\nwenddocs{}\nwbegincode{1067}\moddef{SF epa: parameters}\endmoddef
  integer, parameter :: NONE = 0
  integer, parameter :: ZERO_QMIN = 1
  integer, parameter :: Q_MAX_TOO_SMALL = 2
  integer, parameter :: ZERO_XMIN = 3
  integer, parameter :: MASS_MIX = 4
  integer, parameter :: NO_EPA = 5
\nwendcode{}\nwbegincode{1068}\moddef{SF epa: epa data: TBP}\endmoddef
  procedure :: init => epa_data_init
\nwendcode{}\nwbegincode{1069}\moddef{SF epa: procedures}\endmoddef
  subroutine epa_data_init &
       (data, model, pdg_in, alpha, x_min, q_min, E_max, mass, recoil)
    class(epa_data_t), intent(inout) :: data
    class(model_data_t), intent(in), target :: model
    type(pdg_array_t), intent(in) :: pdg_in
    real(default), intent(in) :: alpha, x_min, q_min, E_max
    real(default), intent(in), optional :: mass
    logical, intent(in), optional :: recoil
    integer :: n_flv, i
    data%model => model
    n_flv = pdg_array_get_length (pdg_in)
    allocate (data%flv_in (n_flv))
    do i = 1, n_flv
       call data%flv_in(i)%init (pdg_array_get (pdg_in, i), model)
    end do
    data%alpha = alpha
    data%E_max = E_max
    data%x_min = x_min
    data%x_max = 1
    if (data%x_min == 0) then
       data%error = ZERO_XMIN;  return
    end if
    data%q_min = q_min
    data%q_max = 2 * data%E_max
    select case (char (data%model%get_name ()))
    case ("QCD","Test")
       data%error = NO_EPA;  return
    end select     
    if (present (recoil)) then
       data%recoil = recoil
    end if
    if (present (mass)) then
       data%mass = mass
    else
       data%mass = data%flv_in(1)%get_mass ()
       if (any (data%flv_in%get_mass () /= data%mass)) then
          data%error = MASS_MIX;  return
       end if 
    end if
    if (max (data%mass, data%q_min) == 0) then
       data%error = ZERO_QMIN;  return
    else if (max (data%mass, data%q_min) >= data%E_max) then
       data%error = Q_MAX_TOO_SMALL;  return
    end if
    data%log = log (4 * (data%E_max / max (data%mass, data%q_min)) ** 2 )
    data%a  = data%alpha / pi
    data%c0 = log (data%x_min) * (data%log - log (data%x_min))
    data%c1 = log (data%x_max) * (data%log - log (data%x_max))
    data%dc = data%c1 - data%c0
  end subroutine epa_data_init

\eatline
\nwendcode{}\nwbegindocs{1070}Handle error conditions.  Should always be done after
initialization, unless we are sure everything is ok.
\nwenddocs{}\nwbegincode{1071}\moddef{SF epa: epa data: TBP}\plusendmoddef
  procedure :: check => epa_data_check
\nwendcode{}\nwbegincode{1072}\moddef{SF epa: procedures}\plusendmoddef
  subroutine epa_data_check (data)
    class(epa_data_t), intent(in) :: data
    select case (data%error)
    case (NO_EPA)
       call msg_fatal ("EPA structure function not available for model " &
            // char (data%model%get_name ()) // ".")
    case (ZERO_QMIN)
       call msg_fatal ("EPA: Particle mass is zero")
    case (Q_MAX_TOO_SMALL)
       call msg_fatal ("EPA: Particle mass exceeds Qmax")
    case (ZERO_XMIN)
       call msg_fatal ("EPA: x_min must be larger than zero")
    case (MASS_MIX)
       call msg_fatal ("EPA: incoming particle masses must be uniform")
    end select
  end subroutine epa_data_check

\eatline
\nwendcode{}\nwbegindocs{1073}Output
\nwenddocs{}\nwbegincode{1074}\moddef{SF epa: epa data: TBP}\plusendmoddef
  procedure :: write => epa_data_write
\nwendcode{}\nwbegincode{1075}\moddef{SF epa: procedures}\plusendmoddef
  subroutine epa_data_write (data, unit, verbose) 
    class(epa_data_t), intent(in) :: data
    integer, intent(in), optional :: unit
    logical, intent(in), optional :: verbose
    integer :: u, i
    u = given_output_unit (unit);  if (u < 0)  return
    write (u, "(1x,A)") "EPA data:"
    if (allocated (data%flv_in)) then
       write (u, "(3x,A)", advance="no") "  flavor =  "
       do i = 1, size (data%flv_in)
          if (i > 1)  write (u, "(',',1x)", advance="no")
          call data%flv_in(i)%write (u)
       end do
       write (u, *)
       write (u, "(3x,A," // FMT_19 // ")") "  alpha  = ", data%alpha
       write (u, "(3x,A," // FMT_19 // ")") "  x_min  = ", data%x_min
       write (u, "(3x,A," // FMT_19 // ")") "  x_max  = ", data%x_max
       write (u, "(3x,A," // FMT_19 // ")") "  q_min  = ", data%q_min
       write (u, "(3x,A," // FMT_19 // ")") "  q_max  = ", data%q_max
       write (u, "(3x,A," // FMT_19 // ")") "  E_max  = ", data%e_max
       write (u, "(3x,A," // FMT_19 // ")") "  mass   = ", data%mass
       write (u, "(3x,A," // FMT_19 // ")") "  a      = ", data%a
       write (u, "(3x,A," // FMT_19 // ")") "  c0     = ", data%c0
       write (u, "(3x,A," // FMT_19 // ")") "  c1     = ", data%c1
       write (u, "(3x,A," // FMT_19 // ")") "  log    = ", data%log
       write (u, "(3x,A,L2)")      "  recoil = ", data%recoil
    else
       write (u, "(3x,A)") "[undefined]"
    end if
  end subroutine epa_data_write

\eatline
\nwendcode{}\nwbegindocs{1076}The number of kinematic parameters.
\nwenddocs{}\nwbegincode{1077}\moddef{SF epa: epa data: TBP}\plusendmoddef
  procedure :: get_n_par => epa_data_get_n_par
\nwendcode{}\nwbegincode{1078}\moddef{SF epa: procedures}\plusendmoddef
  function epa_data_get_n_par (data) result (n)
    class(epa_data_t), intent(in) :: data
    integer :: n
    if (data%recoil) then
       n = 3
    else
       n = 1
    end if
  end function epa_data_get_n_par
  
\eatline
\nwendcode{}\nwbegindocs{1079}Return the outgoing particles PDG codes. The outgoing particle is always
the photon while the radiated particle is identical to the incoming one.
\nwenddocs{}\nwbegincode{1080}\moddef{SF epa: epa data: TBP}\plusendmoddef
  procedure :: get_pdg_out => epa_data_get_pdg_out
\nwendcode{}\nwbegincode{1081}\moddef{SF epa: procedures}\plusendmoddef
  subroutine epa_data_get_pdg_out (data, pdg_out)
    class(epa_data_t), intent(in) :: data
    type(pdg_array_t), dimension(:), intent(inout) :: pdg_out
    pdg_out(1) = PHOTON
  end subroutine epa_data_get_pdg_out
  
\eatline
\nwendcode{}\nwbegindocs{1082}Allocate the interaction record.
\nwenddocs{}\nwbegincode{1083}\moddef{SF epa: epa data: TBP}\plusendmoddef
  procedure :: allocate_sf_int => epa_data_allocate_sf_int
\nwendcode{}\nwbegincode{1084}\moddef{SF epa: procedures}\plusendmoddef
  subroutine epa_data_allocate_sf_int (data, sf_int)
    class(epa_data_t), intent(in) :: data
    class(sf_int_t), intent(inout), allocatable :: sf_int
    allocate (epa_t :: sf_int)
  end subroutine epa_data_allocate_sf_int
  
\eatline
\nwendcode{}\nwbegindocs{1085}\nwdocspar
\subsection{The EPA object}
The {\tt{}epa{\char95}t} data type is a $1\to 2$ interaction.  We should be able
to handle several flavors in parallel, since EPA is not necessarily
applied immediately after beam collision:  Photons may be radiated
from quarks.  In that case, the partons are massless and $q_{\rm min}$
applies instead, so we do not need to generate several kinematical
configurations in parallel.

The squared charge values multiply the matrix elements, depending on the
flavour.  We scan the interaction after building it, so we have the correct
assignments.

The particles are ordered as (incoming, radiated, photon), where the
photon initiates the hard interaction.

We generate an unpolarized photon and transfer initial polarization to
the radiated parton.  Color is transferred in the same way.
\nwenddocs{}\nwbegincode{1086}\moddef{SF epa: types}\plusendmoddef
  type, extends (sf_int_t) :: epa_t
     type(epa_data_t), pointer :: data => null ()
     real(default) :: x  = 0
     real(default) :: xb = 0
     real(default) :: E  = 0     
     real(default), dimension(:), allocatable :: charge2
   contains
   \LA{}SF epa: epa: TBP\RA{}
  end type epa_t 
  
\eatline
\nwendcode{}\nwbegindocs{1087}Type string: has to be here, but there is no string variable on which EPA
depends. Hence, a dummy routine.
\nwenddocs{}\nwbegincode{1088}\moddef{SF epa: epa: TBP}\endmoddef
  procedure :: type_string => epa_type_string
\nwendcode{}\nwbegincode{1089}\moddef{SF epa: procedures}\plusendmoddef
  function epa_type_string (object) result (string)
    class(epa_t), intent(in) :: object
    type(string_t) :: string
    if (associated (object%data)) then
       string = "EPA: equivalent photon approx." 
    else
       string = "EPA: [undefined]"
    end if
  end function epa_type_string
  
\eatline
\nwendcode{}\nwbegindocs{1090}Output.  Call the interaction routine after displaying the configuration.
\nwenddocs{}\nwbegincode{1091}\moddef{SF epa: epa: TBP}\plusendmoddef
  procedure :: write => epa_write
\nwendcode{}\nwbegincode{1092}\moddef{SF epa: procedures}\plusendmoddef
  subroutine epa_write (object, unit, testflag)
    class(epa_t), intent(in) :: object
    integer, intent(in), optional :: unit
    logical, intent(in), optional :: testflag
    integer :: u
    u = given_output_unit (unit)
    if (associated (object%data)) then
       call object%data%write (u)
       if (object%status >= SF_DONE_KINEMATICS) then
          write (u, "(1x,A)")  "SF parameters:"
          write (u, "(3x,A," // FMT_17 // ")")  "x =", object%x
          if (object%status >= SF_FAILED_EVALUATION) then
             write (u, "(3x,A," // FMT_17 // ")")  "E =", object%E
          end if          
       end if
       call object%base_write (u, testflag)
    else
       write (u, "(1x,A)")  "EPA data: [undefined]"
    end if
  end subroutine epa_write
    
\eatline
\nwendcode{}\nwbegindocs{1093}Prepare the interaction object.  We have to construct transition matrix
elements for all flavor and helicity combinations.
\nwenddocs{}\nwbegincode{1094}\moddef{SF epa: epa: TBP}\plusendmoddef
  procedure :: init => epa_init
\nwendcode{}\nwbegincode{1095}\moddef{SF epa: procedures}\plusendmoddef
  subroutine epa_init (sf_int, data)
    class(epa_t), intent(out) :: sf_int
    class(sf_data_t), intent(in), target :: data
    type(quantum_numbers_mask_t), dimension(3) :: mask
    integer, dimension(3) :: hel_lock
    type(polarization_t) :: pol
    type(quantum_numbers_t), dimension(1) :: qn_fc, qn_hel
    type(flavor_t) :: flv_photon
    type(color_t) :: col_photon
    type(quantum_numbers_t) :: qn_photon, qn, qn_rad
    type(state_iterator_t) :: it_hel
    integer :: i
    mask = quantum_numbers_mask (.false., .false., &
         mask_h = [.false., .false., .true.])
    hel_lock = [2, 1, 0]
    select type (data)
    type is (epa_data_t)
       call sf_int%base_init (mask, [data%mass**2], &
            [data%mass**2], [0._default], hel_lock = hel_lock)       
       sf_int%data => data
       call flv_photon%init (PHOTON, data%model)
       call col_photon%init ()
       call qn_photon%init (flv_photon, col_photon)
       do i = 1, size (data%flv_in)
          call polarization_init_generic (pol, data%flv_in(i))
          call qn_fc(1)%init ( &
               flv = data%flv_in(i), &
               col = color_from_flavor (data%flv_in(i), 1))
          call it_hel%init (pol%state)
          do while (it_hel%is_valid ())
             qn_hel = it_hel%get_quantum_numbers ()
             qn = qn_hel(1) .merge. qn_fc(1)
             qn_rad = qn
             call qn_rad%tag_radiated ()
             call interaction_add_state (sf_int%interaction_t, &
                  [qn, qn_rad, qn_photon])
             call it_hel%advance ()
          end do
          call polarization_final (pol)
       end do
       call interaction_freeze (sf_int%interaction_t)
       call sf_int%set_incoming ([1])
       call sf_int%set_radiated ([2])
       call sf_int%set_outgoing ([3])
    end select
  end subroutine epa_init
    
\eatline
\nwendcode{}\nwbegindocs{1096}Prepare the charge array.  This is separate from the previous routine since
the state matrix may be helicity-contracted.
\nwenddocs{}\nwbegincode{1097}\moddef{SF epa: epa: TBP}\plusendmoddef
  procedure :: setup_constants => epa_setup_constants
\nwendcode{}\nwbegincode{1098}\moddef{SF epa: procedures}\plusendmoddef
  subroutine epa_setup_constants (sf_int)
    class(epa_t), intent(inout) :: sf_int
    type(state_iterator_t) :: it
    type(flavor_t) :: flv
    integer :: i, n_me
    n_me = interaction_get_n_matrix_elements (sf_int%interaction_t)
    allocate (sf_int%charge2 (n_me))
    call it%init (interaction_get_state_matrix_ptr (sf_int%interaction_t))
    do while (it%is_valid ())
       i = it%get_me_index ()
       flv = it%get_flavor (1)
       sf_int%charge2(i) = flv%get_charge () ** 2
       call it%advance ()
    end do
    sf_int%status = SF_INITIAL
  end subroutine epa_setup_constants
  
\eatline
\nwendcode{}\nwbegindocs{1099}\nwdocspar
\subsection{Kinematics}
Set kinematics.  If {\tt{}map} is unset, the $r$ and $x$ values
coincide, and the Jacobian $f(r)$ is trivial.

The EPA structure function allows for a straightforward mapping of the
unit interval. The $x$ value is transformed, and the mapped structure
function becomes unity at its upper boundary. 

The structure function implementation applies the above mapping to the
input (random) number {\tt{}r} to generate the momentum fraction {\tt{}x}
and the function value {\tt{}f}.  For numerical stability reasons, we
also output {\tt{}xb}, which is $\bar x=1-x$. 
\nwenddocs{}\nwbegincode{1100}\moddef{SF epa: epa: TBP}\plusendmoddef
  procedure :: complete_kinematics => epa_complete_kinematics
\nwendcode{}\nwbegincode{1101}\moddef{SF epa: procedures}\plusendmoddef
  subroutine epa_complete_kinematics (sf_int, x, f, r, rb, map)
    class(epa_t), intent(inout) :: sf_int
    real(default), dimension(:), intent(out) :: x
    real(default), intent(out) :: f
    real(default), dimension(:), intent(in) :: r
    real(default), dimension(:), intent(in) :: rb
    logical, intent(in) :: map
    real(default) :: xb1
    real(default) :: delta, sqrt_delta, lx
    if (map) then                                   
       associate (data => sf_int%data)
         delta = data%log ** 2 -  4 * (r(1) * data%c1 + rb(1) * data%c0)
         if (delta > 0) then
            sqrt_delta = sqrt (delta)
            lx = (data%log - sqrt_delta) / 2
         else
            sf_int%status = SF_FAILED_KINEMATICS          
            f = 0
            return
         end if
         x(1) = exp (lx)
         f = x(1) * data%dc / sqrt_delta
       end associate
    else
       x(1) = r(1)
       if (sf_int%data%x_min < x(1) .and. x(1) < sf_int%data%x_max) then
          f = 1
       else
          sf_int%status = SF_FAILED_KINEMATICS
          f = 0
          return
       end if       
    end if       
    xb1 = 1 - x(1)
    if (size(x) == 3)  x(2:3) = r(2:3)
    call sf_int%split_momentum (x, xb1) 
    select case (sf_int%status)
    case (SF_DONE_KINEMATICS)
       sf_int%x = x(1)
       sf_int%xb= xb1
       sf_int%E  = &
            energy (interaction_get_momentum (sf_int%interaction_t, 1))
    case (SF_FAILED_KINEMATICS)
       sf_int%x = 0
       sf_int%xb= 0
       f = 0
    end select
  end subroutine epa_complete_kinematics

\eatline
\nwendcode{}\nwbegindocs{1102}Compute inverse kinematics.  Here, we start with the $x$ array and
compute the ``input'' $r$ values and the Jacobian $f$.  After this, we
can set momenta by the same formula as for normal kinematics.
\nwenddocs{}\nwbegincode{1103}\moddef{SF epa: epa: TBP}\plusendmoddef
  procedure :: inverse_kinematics => epa_inverse_kinematics
\nwendcode{}\nwbegincode{1104}\moddef{SF epa: procedures}\plusendmoddef
  subroutine epa_inverse_kinematics (sf_int, x, f, r, rb, map, set_momenta)
    class(epa_t), intent(inout) :: sf_int
    real(default), dimension(:), intent(in) :: x
    real(default), intent(out) :: f
    real(default), dimension(:), intent(out) :: r
    real(default), dimension(:), intent(out) :: rb
    logical, intent(in) :: map
    logical, intent(in), optional :: set_momenta
    real(default) :: lx, delta, sqrt_delta, c
    logical :: set_mom
    set_mom = .false.;  if (present (set_momenta))  set_mom = set_momenta
    if (map) then
       associate (data => sf_int%data)
         lx = log (x(1))
         sqrt_delta = data%log - 2 * lx
         delta = sqrt_delta ** 2
         c = (data%log ** 2 - delta) / 4
         r (1) = (c - data%c0) / data%dc
         rb(1) = (data%c1 - c) / data%dc
         f = x(1) * data%dc / sqrt_delta
       end associate
    else
       r (1) = x(1)
       rb(1) = 1 - x(1)
       if (sf_int%data%x_min < x(1) .and. x(1) < sf_int%data%x_max) then
          f = 1
       else
          f = 0
       end if
    end if
    if (size(r) == 3) then
       r (2:3) = x(2:3)
       rb(2:3) = 1 - x(2:3)
    end if
    if (set_mom) then
       call sf_int%split_momentum (x, sf_int%xb)
       select case (sf_int%status)
       case (SF_DONE_KINEMATICS)
          sf_int%x  = x(1)
          sf_int%xb = 1 - x(1)
          sf_int%E  = &
               energy (interaction_get_momentum (sf_int%interaction_t, 1))
       case (SF_FAILED_KINEMATICS)
          sf_int%x = 0
          f = 0
       end select
    end if
  end subroutine epa_inverse_kinematics

\eatline
\nwendcode{}\nwbegindocs{1105}Overriding the default method: we compute the {\tt{}x} array from the
momentum configuration.  In the specific case of EPA, we also set the
internally stored $x$ and $\bar x$ values, so they can be used in the
following routine.

Note: the extraction of $\bar x$ is not numerically safe, but it can't
be as long as the base {\tt{}recover{\char95}x} isn't.
\nwenddocs{}\nwbegincode{1106}\moddef{SF epa: epa: TBP}\plusendmoddef
  procedure :: recover_x => sf_epa_recover_x
\nwendcode{}\nwbegincode{1107}\moddef{SF epa: procedures}\plusendmoddef
  subroutine sf_epa_recover_x (sf_int, x, x_free)
    class(epa_t), intent(inout) :: sf_int
    real(default), dimension(:), intent(out) :: x
    real(default), intent(inout), optional :: x_free
    call sf_int%base_recover_x (x, x_free)
    sf_int%x  = x(1)
    sf_int%xb = 1 - x(1)
  end subroutine sf_epa_recover_x
  
\eatline
\nwendcode{}\nwbegindocs{1108}\nwdocspar
\subsection{EPA application}
For EPA, we can in principle compute kinematics and function value in
a single step. In order to be able to reweight events, kinematics and
structure function application are separated. This function works on a
single beam, assuming that the input momentum has been set.  We need
three random numbers as input: one for $x$, and two for the polar and
azimuthal angles. Alternatively, for the no-recoil case, we can skip
$p_T$ generation; in this case, we only need one.

For obtaining splitting kinematics, we rely on the assumption that all
in-particles are mass-degenerate (or there is only one), so the
generated $x$ values are identical.
\nwenddocs{}\nwbegincode{1109}\moddef{SF epa: epa: TBP}\plusendmoddef
  procedure :: apply => epa_apply
\nwendcode{}\nwbegincode{1110}\moddef{SF epa: procedures}\plusendmoddef
  subroutine epa_apply (sf_int, scale)
    class(epa_t), intent(inout) :: sf_int
    real(default), intent(in) :: scale
    real(default) :: x, xb, qminsq, qmaxsq, f, E
    associate (data => sf_int%data)
      x = sf_int%x
      xb= sf_int%xb
      E = sf_int%E
      qminsq = max (x ** 2 / xb * data%mass ** 2, data%q_min ** 2)
      qmaxsq = min (4 * xb * E ** 2, data%q_max ** 2)
      if (qminsq < qmaxsq) then
         f = data%a / x &
              * ((xb + x ** 2 / 2) * log (qmaxsq / qminsq) &
              - (1 - x / 2) ** 2 &
              * log ((x**2 + qmaxsq / E ** 2) / (x**2 + qminsq / E ** 2)) &
              - x ** 2 * data%mass ** 2 / qminsq * (1 - qminsq / qmaxsq))
      else
         f = 0
      end if
      call interaction_set_matrix_element &
           (sf_int%interaction_t, cmplx (f, kind=default) * sf_int%charge2)
    end associate
    sf_int%status = SF_EVALUATED
  end subroutine epa_apply

\eatline
\nwendcode{}\nwbegindocs{1111}\nwdocspar
\subsection{Unit tests}
\nwenddocs{}\nwbegincode{1112}\moddef{SF epa: public}\plusendmoddef
  public :: sf_epa_test
\nwendcode{}\nwbegincode{1113}\moddef{SF epa: tests}\endmoddef
  subroutine sf_epa_test (u, results)
    integer, intent(in) :: u
    type(test_results_t), intent(inout) :: results
  \LA{}SF epa: execute tests\RA{}
  end subroutine sf_epa_test
  
\eatline
\nwendcode{}\nwbegindocs{1114}\nwdocspar
\subsubsection{Test structure function data}
Construct and display a test structure function data object.
\nwenddocs{}\nwbegincode{1115}\moddef{SF epa: execute tests}\endmoddef
  call test (sf_epa_1, "sf_epa_1", &
       "structure function configuration", &
       u, results)
\nwendcode{}\nwbegincode{1116}\moddef{SF epa: tests}\plusendmoddef
  subroutine sf_epa_1 (u)
    integer, intent(in) :: u
    type(model_data_t), target :: model
    type(pdg_array_t) :: pdg_in
    type(pdg_array_t), dimension(1) :: pdg_out
    integer, dimension(:), allocatable :: pdg1
    class(sf_data_t), allocatable :: data
    
    write (u, "(A)")  "* Test output: sf_epa_1"
    write (u, "(A)")  "*   Purpose: initialize and display &
         &test structure function data"
    write (u, "(A)")
    
    write (u, "(A)")  "* Create empty data object"
    write (u, "(A)")

    call model%init_qed_test ()
    pdg_in = ELECTRON

    allocate (epa_data_t :: data)
    call data%write (u)

    write (u, "(A)")
    write (u, "(A)")  "* Initialize"
    write (u, "(A)")

    select type (data)
    type is (epa_data_t)
       call data%init (model, pdg_in, 1./137._default, 0.01_default, &
            10._default, 50._default, 0.000511_default, recoil = .false.)
    end select

    call data%write (u)

    write (u, "(A)")

    write (u, "(1x,A)")  "Outgoing particle codes:"
    call data%get_pdg_out (pdg_out)
    pdg1 = pdg_out(1)
    write (u, "(2x,99(1x,I0))")  pdg1
        
    call model%final ()
    
    write (u, "(A)")
    write (u, "(A)")  "* Test output end: sf_epa_1"

  end subroutine sf_epa_1

\eatline
\nwendcode{}\nwbegindocs{1117}\nwdocspar
\subsubsection{Test and probe structure function}
Construct and display a structure function object based on the EPA
structure function.
\nwenddocs{}\nwbegincode{1118}\moddef{SF epa: execute tests}\plusendmoddef
  call test (sf_epa_2, "sf_epa_2", &
       "structure function instance", &
       u, results)
\nwendcode{}\nwbegincode{1119}\moddef{SF epa: tests}\plusendmoddef
  subroutine sf_epa_2 (u)
    integer, intent(in) :: u
    type(model_data_t), target :: model
    type(flavor_t) :: flv
    type(pdg_array_t) :: pdg_in
    class(sf_data_t), allocatable, target :: data
    class(sf_int_t), allocatable :: sf_int
    type(vector4_t) :: k
    type(vector4_t), dimension(2) :: q
    real(default) :: E
    real(default), dimension(:), allocatable :: r, rb, x
    real(default) :: f
    
    write (u, "(A)")  "* Test output: sf_epa_2"
    write (u, "(A)")  "*   Purpose: initialize and fill &
         &test structure function object"
    write (u, "(A)")
    
    write (u, "(A)")  "* Initialize configuration data"
    write (u, "(A)")

    call model%init_qed_test ()
    call flv%init (ELECTRON, model)
    pdg_in = ELECTRON

    call reset_interaction_counter ()
    
    allocate (epa_data_t :: data)
    select type (data)
    type is (epa_data_t)
       call data%init (model, pdg_in, 1./137._default, 0.01_default, &
            10._default, 50._default, 0.000511_default, recoil = .false.)
    end select
       
    write (u, "(A)")  "* Initialize structure-function object"
    write (u, "(A)")
    
    call data%allocate_sf_int (sf_int)
    call sf_int%init (data)
    call sf_int%set_beam_index ([1])
    call sf_int%setup_constants ()

    write (u, "(A)")  "* Initialize incoming momentum with E=500"
    write (u, "(A)")
    E = 500
    k = vector4_moving (E, sqrt (E**2 - flv%get_mass ()**2), 3)
    call pacify (k, 1e-10_default)
    call vector4_write (k, u)
    call sf_int%seed_kinematics ([k])
    
    write (u, "(A)")
    write (u, "(A)")  "* Set kinematics for r=0.4, no EPA mapping, collinear"
    write (u, "(A)")
    
    allocate (r (data%get_n_par ()))
    allocate (rb(size (r)))
    allocate (x (size (r)))
    
    r = 0.4_default
    rb = 1 - r
    call sf_int%complete_kinematics (x, f, r, rb, map=.false.)
    
    write (u, "(A,9(1x,F10.7))")  "r =", r
    write (u, "(A,9(1x,F10.7))")  "rb=", rb
    write (u, "(A,9(1x,F10.7))")  "x =", x
    write (u, "(A,9(1x,F10.7))")  "f =", f
    
    write (u, "(A)")
    write (u, "(A)")  "* Recover x from momenta"
    write (u, "(A)")
    
    q = interaction_get_momenta (sf_int%interaction_t, outgoing=.true.)
    call sf_int%final ()
    deallocate (sf_int)
    
    call data%allocate_sf_int (sf_int)
    call sf_int%init (data)
    call sf_int%set_beam_index ([1])
    call sf_int%setup_constants ()
    
    call sf_int%seed_kinematics ([k])
    call interaction_set_momenta (sf_int%interaction_t, q, outgoing=.true.)
    call sf_int%recover_x (x)
    call sf_int%inverse_kinematics (x, f, r, rb, map=.false., &
         set_momenta=.true.)
    
    write (u, "(A,9(1x,F10.7))")  "r =", r
    write (u, "(A,9(1x,F10.7))")  "rb=", rb
    write (u, "(A,9(1x,F10.7))")  "x =", x
    write (u, "(A,9(1x,F10.7))")  "f =", f

    write (u, "(A)")
    write (u, "(A)")  "* Evaluate EPA structure function"
    write (u, "(A)")
    
    call sf_int%apply (scale = 100._default)
    call sf_int%write (u)
    
    write (u, "(A)")
    write (u, "(A)")  "* Cleanup"
    
    call sf_int%final ()
    call model%final ()
    
    write (u, "(A)")
    write (u, "(A)")  "* Test output end: sf_epa_2"

  end subroutine sf_epa_2

\eatline
\nwendcode{}\nwbegindocs{1120}\nwdocspar
\subsubsection{Standard mapping}
Construct and display a structure function object based on the EPA
structure function, applying the standard single-particle mapping.
\nwenddocs{}\nwbegincode{1121}\moddef{SF epa: execute tests}\plusendmoddef
  call test (sf_epa_3, "sf_epa_3", &
       "apply mapping", &
       u, results)
\nwendcode{}\nwbegincode{1122}\moddef{SF epa: tests}\plusendmoddef
  subroutine sf_epa_3 (u)
    integer, intent(in) :: u
    type(model_data_t), target :: model
    type(flavor_t) :: flv
    type(pdg_array_t) :: pdg_in
    class(sf_data_t), allocatable, target :: data
    class(sf_int_t), allocatable :: sf_int
    type(vector4_t) :: k
    type(vector4_t), dimension(2) :: q
    real(default) :: E
    real(default), dimension(:), allocatable :: r, rb, x
    real(default) :: f
    
    write (u, "(A)")  "* Test output: sf_epa_3"
    write (u, "(A)")  "*   Purpose: initialize and fill &
         &test structure function object"
    write (u, "(A)")
    
    write (u, "(A)")  "* Initialize configuration data"
    write (u, "(A)")

    call model%init_qed_test ()
    call flv%init (ELECTRON, model)
    pdg_in = ELECTRON

    call reset_interaction_counter ()
    
    allocate (epa_data_t :: data)
    select type (data)
    type is (epa_data_t)
       call data%init (model, pdg_in, 1./137._default, 0.01_default, &
            10._default, 50._default, 0.000511_default, recoil = .false.)
    end select
       
    write (u, "(A)")  "* Initialize structure-function object"
    write (u, "(A)")
    
    call data%allocate_sf_int (sf_int)
    call sf_int%init (data)
    call sf_int%set_beam_index ([1])
    call sf_int%setup_constants ()

    write (u, "(A)")  "* Initialize incoming momentum with E=500"
    write (u, "(A)")
    E = 500
    k = vector4_moving (E, sqrt (E**2 - flv%get_mass ()**2), 3)
    call pacify (k, 1e-10_default)
    call vector4_write (k, u)
    call sf_int%seed_kinematics ([k])
    
    write (u, "(A)")
    write (u, "(A)")  "* Set kinematics for r=0.4, with EPA mapping, collinear"
    write (u, "(A)")
    
    allocate (r (data%get_n_par ()))
    allocate (rb(size (r)))
    allocate (x (size (r)))
    
    r = 0.4_default
    rb = 1 - r
    call sf_int%complete_kinematics (x, f, r, rb, map=.true.)
    
    write (u, "(A,9(1x,F10.7))")  "r =", r
    write (u, "(A,9(1x,F10.7))")  "rb=", rb
    write (u, "(A,9(1x,F10.7))")  "x =", x
    write (u, "(A,9(1x,F10.7))")  "f =", f
    
    write (u, "(A)")
    write (u, "(A)")  "* Recover x from momenta"
    write (u, "(A)")
    
    q = interaction_get_momenta (sf_int%interaction_t, outgoing=.true.)
    call sf_int%final ()
    deallocate (sf_int)
    
    call data%allocate_sf_int (sf_int)
    call sf_int%init (data)
    call sf_int%set_beam_index ([1])
    call sf_int%setup_constants ()
    
    call sf_int%seed_kinematics ([k])
    call interaction_set_momenta (sf_int%interaction_t, q, outgoing=.true.)
    call sf_int%recover_x (x)
    call sf_int%inverse_kinematics (x, f, r, rb, map=.true., &
         set_momenta=.true.)
    
    write (u, "(A,9(1x,F10.7))")  "r =", r
    write (u, "(A,9(1x,F10.7))")  "rb=", rb
    write (u, "(A,9(1x,F10.7))")  "x =", x
    write (u, "(A,9(1x,F10.7))")  "f =", f

    write (u, "(A)")
    write (u, "(A)")  "* Evaluate EPA structure function"
    write (u, "(A)")
    
    call sf_int%apply (scale = 100._default)
    call sf_int%write (u)
    
    write (u, "(A)")
    write (u, "(A)")  "* Cleanup"
    
    call sf_int%final ()
    call model%final ()
    
    write (u, "(A)")
    write (u, "(A)")  "* Test output end: sf_epa_3"

  end subroutine sf_epa_3

\eatline
\nwendcode{}\nwbegindocs{1123}\nwdocspar
\subsubsection{Non-collinear case}
Construct and display a structure function object based on the EPA
structure function.
\nwenddocs{}\nwbegincode{1124}\moddef{SF epa: execute tests}\plusendmoddef
  call test (sf_epa_4, "sf_epa_4", &
       "non-collinear", &
       u, results)
\nwendcode{}\nwbegincode{1125}\moddef{SF epa: tests}\plusendmoddef
  subroutine sf_epa_4 (u)
    integer, intent(in) :: u
    type(model_data_t), target :: model
    type(flavor_t) :: flv
    type(pdg_array_t) :: pdg_in
    class(sf_data_t), allocatable, target :: data
    class(sf_int_t), allocatable :: sf_int
    type(vector4_t) :: k
    type(vector4_t), dimension(2) :: q
    real(default) :: E, m
    real(default), dimension(:), allocatable :: r, rb, x
    real(default) :: f
    
    write (u, "(A)")  "* Test output: sf_epa_4"
    write (u, "(A)")  "*   Purpose: initialize and fill &
         &test structure function object"
    write (u, "(A)")
    
    write (u, "(A)")  "* Initialize configuration data"
    write (u, "(A)")

    call model%init_qed_test ()
    call flv%init (ELECTRON, model)
    pdg_in = ELECTRON

    call reset_interaction_counter ()

    allocate (epa_data_t :: data)
    select type (data)
    type is (epa_data_t)
       call data%init (model, pdg_in, 1./137._default, 0.01_default, &
            10._default, 50._default, 5.0_default, recoil = .true.)
    end select    

    write (u, "(A)")  "* Initialize structure-function object"
    write (u, "(A)")
        
    call data%allocate_sf_int (sf_int)
    call sf_int%init (data)
    call sf_int%set_beam_index ([1])
    call sf_int%setup_constants ()

    write (u, "(A)")  "* Initialize incoming momentum with E=500, me = 5 GeV"
    write (u, "(A)")
    E = 500
    m = 5
    k = vector4_moving (E, sqrt (E**2 - m**2), 3)
    call pacify (k, 1e-10_default)
    call vector4_write (k, u)
    call sf_int%seed_kinematics ([k])
        
    write (u, "(A)")
    write (u, "(A)")  "* Set kinematics for r=0.5/0.5/0.25, with EPA mapping, "
    write (u, "(A)")  "          non-coll., keeping energy, me = 5 GeV"
    write (u, "(A)")
    
    allocate (r (data%get_n_par ()))
    allocate (rb(size (r)))
    allocate (x (size (r)))
    
    r = [0.5_default, 0.5_default, 0.25_default]
    rb = 1 - r
    sf_int%on_shell_mode = KEEP_ENERGY    
    call sf_int%complete_kinematics (x, f, r, rb, map=.true.)
    call interaction_pacify_momenta (sf_int%interaction_t, 1e-10_default)

    write (u, "(A,9(1x,F10.7))")  "r =", r
    write (u, "(A,9(1x,F10.7))")  "rb=", rb
    write (u, "(A,9(1x,F10.7))")  "x =", x
    write (u, "(A,9(1x,F10.7))")  "f =", f

    write (u, "(A)")
    write (u, "(A)")  "* Recover x and r from momenta"
    write (u, "(A)")
    
    q = interaction_get_momenta (sf_int%interaction_t, outgoing=.true.)
    call sf_int%final ()
    deallocate (sf_int)
    
    call data%allocate_sf_int (sf_int)
    call sf_int%init (data)
    call sf_int%set_beam_index ([1])
    call sf_int%setup_constants ()
    
    call sf_int%seed_kinematics ([k])
    call interaction_set_momenta (sf_int%interaction_t, q, outgoing=.true.)
    call sf_int%recover_x (x)
    call sf_int%inverse_kinematics (x, f, r, rb, map=.true., &
         set_momenta=.true.)    
    call interaction_pacify_momenta (sf_int%interaction_t, 1e-10_default)    
    
    write (u, "(A,9(1x,F10.7))")  "r =", r
    write (u, "(A,9(1x,F10.7))")  "rb=", rb
    write (u, "(A,9(1x,F10.7))")  "x =", x
    write (u, "(A,9(1x,F10.7))")  "f =", f

    write (u, "(A)")
    write (u, "(A)")  "* Evaluate EPA structure function"
    write (u, "(A)")
    
    call sf_int%apply (scale = 100._default)
    call sf_int%write (u, testflag = .true.)
    
    write (u, "(A)")
    write (u, "(A)")  "* Cleanup"
    
    call sf_int%final ()
    call model%final ()
    
    write (u, "(A)")
    write (u, "(A)")  "* Test output end: sf_epa_4"

  end subroutine sf_epa_4

\eatline
\nwendcode{}\nwbegindocs{1126}\nwdocspar
\subsubsection{Structure function for multiple flavors}
Construct and display a structure function object based on the EPA
structure function.  The incoming state has multiple particles with
non-uniform charge.
\nwenddocs{}\nwbegincode{1127}\moddef{SF epa: execute tests}\plusendmoddef
  call test (sf_epa_5, "sf_epa_5", &
       "multiple flavors", &
       u, results)
\nwendcode{}\nwbegincode{1128}\moddef{SF epa: tests}\plusendmoddef
  subroutine sf_epa_5 (u)
    integer, intent(in) :: u
    type(model_data_t), target :: model
    type(flavor_t) :: flv
    type(pdg_array_t) :: pdg_in
    class(sf_data_t), allocatable, target :: data
    class(sf_int_t), allocatable :: sf_int
    type(vector4_t) :: k
    real(default) :: E
    real(default), dimension(:), allocatable :: r, rb, x
    real(default) :: f
    
    write (u, "(A)")  "* Test output: sf_epa_5"
    write (u, "(A)")  "*   Purpose: initialize and fill &
         &test structure function object"
    write (u, "(A)")
    
    write (u, "(A)")  "* Initialize configuration data"
    write (u, "(A)")

    call model%init_sm_test ()
    call flv%init (1, model)
    pdg_in = [1, 2, -1, -2]

    call reset_interaction_counter ()
    
    allocate (epa_data_t :: data)
    select type (data)
    type is (epa_data_t)
       call data%init (model, pdg_in, 1./137._default, 0.01_default, &
            10._default, 50._default, 0.000511_default, recoil = .false.)
       call data%check ()
    end select
       
    write (u, "(A)")  "* Initialize structure-function object"
    write (u, "(A)")
    
    call data%allocate_sf_int (sf_int)
    call sf_int%init (data)
    call sf_int%set_beam_index ([1])
    call sf_int%setup_constants ()

    write (u, "(A)")  "* Initialize incoming momentum with E=500"
    write (u, "(A)")
    E = 500
    k = vector4_moving (E, sqrt (E**2 - flv%get_mass ()**2), 3)
    call pacify (k, 1e-10_default)
    call vector4_write (k, u)
    call sf_int%seed_kinematics ([k])
    
    write (u, "(A)")
    write (u, "(A)")  "* Set kinematics for r=0.4, no EPA mapping, collinear"
    write (u, "(A)")
    
    allocate (r (data%get_n_par ()))
    allocate (rb(size (r)))
    allocate (x (size (r)))
    
    r = 0.4_default
    rb = 1 - r
    call sf_int%complete_kinematics (x, f, r, rb, map=.false.)
    
    write (u, "(A,9(1x,F10.7))")  "r =", r
    write (u, "(A,9(1x,F10.7))")  "rb=", rb
    write (u, "(A,9(1x,F10.7))")  "x =", x
    write (u, "(A,9(1x,F10.7))")  "f =", f
    
    write (u, "(A)")
    write (u, "(A)")  "* Evaluate EPA structure function"
    write (u, "(A)")
    
    call sf_int%apply (scale = 100._default)
    call sf_int%write (u)
    
    write (u, "(A)")
    write (u, "(A)")  "* Cleanup"
    
    call sf_int%final ()
    call model%final ()
    
    write (u, "(A)")
    write (u, "(A)")  "* Test output end: sf_epa_5"

  end subroutine sf_epa_5

\eatline
\nwendcode{}\nwbegindocs{1129}\nwdocspar
\clearpage
%------------------------------------------------------------------------
\section{EWA}

\nwenddocs{}\nwbegincode{1130}\moddef{\code{}sf{\char95}ewa.f90\edoc{}}\endmoddef
\LA{}File header\RA{}

module sf_ewa

\LA{}Use kinds\RA{}
\LA{}Use strings\RA{}
  use io_units
  use constants, only: pi
  use format_defs, only: FMT_17, FMT_19
  use unit_tests
  use diagnostics
  use physics_defs, only: W_BOSON, Z_BOSON
  use lorentz
  use pdg_arrays
  use model_data
  use flavors
  use colors
  use quantum_numbers
  use state_matrices
  use polarizations
  use interactions
  use sf_aux
  use sf_base

\LA{}Standard module head\RA{}

\LA{}SF ewa: public\RA{}

\LA{}SF ewa: parameters\RA{}

\LA{}SF ewa: types\RA{}

contains

\LA{}SF ewa: procedures\RA{}

\LA{}SF ewa: tests\RA{}

end module sf_ewa
\eatline
\nwendcode{}\nwbegindocs{1131}\nwdocspar
\subsection{Physics}
The EWA structure function for a $Z$ or $W$ inside a fermion (lepton
or quark) depends on the vector-boson polarization.  We distinguish
transversal ($\pm$) and longitudinal ($0$) polarization.
\begin{align}
  F_{+}(x) &= \frac{1}{16\pi^2}\,\frac{(v-a)^2 + (v+a)^2\bar x^2}{x}
  \left[
  \ln\left(\frac{p_{\perp,\textrm{max}}^2 + \bar x M^2}{\bar x M^2}\right)
  -
  \frac{p_{\perp,\textrm{max}}^2}{p_{\perp,\textrm{max}}^2 + \bar x M^2}
  \right]
  \\
  F_{-}(x) &= \frac{1}{16\pi^2}\,\frac{(v+a)^2 + (v-a)^2\bar x^2}{x}
  \left[
  \ln\left(\frac{p_{\perp,\textrm{max}}^2 + \bar x M^2}{\bar x M^2}\right)
  -
  \frac{p_{\perp,\textrm{max}}^2}{p_{\perp,\textrm{max}}^2 + \bar x M^2}
  \right]
  \\
  F_0(x) &= \frac{v^2+a^2}{8\pi^2}\,\frac{2\bar x}{x}\,
  \frac{p_{\perp,\textrm{max}}^2}{p_{\perp,\textrm{max}}^2 + \bar x M^2}
\end{align}
where $p_{\perp,\textrm{max}}$ is the cutoff in transversal momentum, $M$ is
the vector-boson mass, $v$ and $a$ are the vector and axial-vector
couplings, and $\bar x\equiv 1-x$.  Note that the longitudinal
structure function is finite for large cutoff, while the transversal
structure function is logarithmically divergent.

The maximal transverse momentum is given by the kinematical limit, it is
\begin{equation}
p_{\perp,\textrm{max}} = \bar x \sqrt{s}/2.
\end{equation}
The vector and axial couplings for a fermion branching into a $W$ are
\begin{align}
  v_W &= \frac{g}{2\sqrt 2},
& a_W &= \frac{g}{2\sqrt 2}.
\end{align}
For $Z$ emission, this is replaced by
\begin{align}
  v_Z &= \frac{g}{2\cos\theta_w}\left(t_3 - 2q\sin^2\theta_w\right),
& a_Z &= \frac{g}{2\cos\theta_w}t_3,
\end{align}
where $t_3=\pm\frac12$ is the fermion isospin, and $q$ its charge.

For an initial antifermion, the signs of the axial couplings are
inverted.  Note that a common sign change of $v$ and $a$ is
irrelevant.

%% Differentiating with respect to the cutoff, we get structure functions
%% \begin{align}
%%   f_{W,\pm}(x,p_T) &= \frac{g^2}{16\pi^2}\,
%%   \frac{1+\bar x^2}{x}
%%   \frac{p_\perp}{p_\perp^2 + \bar x M^2}
%%   \\
%%   f_{W,0}(x,p_T) &= \frac{g^2}{16\pi^2}\,
%%   \frac{2\bar x}{x}\,
%%   \frac{p_\perp \bar xM^2}{(p_\perp^2 + \bar x M^2)^2}
%%   \\
%%   F_{Z,\pm}(x,p_T) &= \frac{g^2}{16\pi^2\cos\theta_w^2}
%%   \left[(t_3^f-2q^2\sin\theta_w^2)^2 + (t_3^f)^2\right]\,
%%   \frac{1+\bar x^2}{x}
%%   \frac{p_\perp}{p_\perp^2 + \bar x M^2}
%%   \\
%%   F_{Z,0}(x,p_T) &= \frac{g^2}{16\pi^2\cos\theta_w^2}\,
%%   \left[(t_3^f-2q^2\sin\theta_w^2)^2 + (t_3^f)^2\right]\,
%%   \frac{2\bar x}{x}\,
%%   \frac{p_\perp \bar xM^2}{(p_\perp^2 + \bar x M^2)^2}
%% \end{align}
%% Here, $t_3^f$ is the $SU(2)_L$ quantum number of the fermion
%% $(\pm\frac12)$, and $q^f$ is the fermion charge in units of the
%% positron charge.

The EWA depends on the parameters $g$, $\sin^2\theta_w$, $M_W$, and
$M_Z$.  These can all be taken from the SM input, and the prefactors
are calculated from those and the incoming particle type.

Since these structure functions have a $1/x$ singularity (which is not
really relevant in practice, however, since the vector boson mass is
finite), we map this singularity allowing for nontrivial $x$ bounds:
\begin{equation}
  x = \exp(\bar r\ln x_0 + r\ln x_1)
\end{equation}
such that
\begin{equation}
  \int_{x_0}^{x_1}\frac{dx}{x} = (\ln x_1 - \ln x_0)\int_0^1 dr.
\end{equation}

As a user parameter, we have the cutoff $p_{\perp,\textrm{max}}$.
The divergence $1/x$ also requires a $x_0$ cutoff; and for
completeness we introduce a corresponding $x_1$.  Physically, the
minimal sensible value of $x$ is $M^2/s$, although the approximation
loses its value already at higher $x$ values.


\subsection{The EWA data block}
The EWA parameters are: $p_{T,\rm max}$, $c_V$, $c_A$, and
$m$. Instead of $m$ we can use the incoming particle PDG code as
input; from this we can deduce the mass and charges. In the
initialization phase it is not yet determined whether a $W$ or a $Z$
is radiated, hence we set the vector and axial-vector couplings equal
to the common prefactors $g/2 = e/2/\sin\theta_W$. 

In principle, for EWA it would make sense to allow the user to also
set the upper bound for $x$, $x_{\rm max}$, but we fix it to one here.
\nwenddocs{}\nwbegincode{1132}\moddef{SF ewa: public}\endmoddef
  public :: ewa_data_t
\nwendcode{}\nwbegincode{1133}\moddef{SF ewa: types}\endmoddef
  type, extends(sf_data_t) :: ewa_data_t
     private
     class(model_data_t), pointer :: model => null ()
     type(flavor_t), dimension(:), allocatable :: flv_in
     type(flavor_t), dimension(:), allocatable :: flv_out
     real(default) :: pt_max
     real(default) :: sqrts
     real(default) :: x_min
     real(default) :: x_max
     real(default) :: mass
     real(default) :: m_out
     real(default) :: q_min
     real(default) :: cv
     real(default) :: ca
     real(default) :: costhw
     real(default) :: sinthw
     real(default) :: mW
     real(default) :: mZ
     real(default) :: coeff
     logical :: mass_set = .false.
     logical :: keep_momentum
     logical :: keep_energy     
     integer :: id = 0 
     integer :: error = NONE
   contains
   \LA{}SF ewa: ewa data: TBP\RA{}  
  end type ewa_data_t

\eatline
\nwendcode{}\nwbegindocs{1134}Error codes
\nwenddocs{}\nwbegincode{1135}\moddef{SF ewa: parameters}\endmoddef
  integer, parameter :: NONE = 0
  integer, parameter :: ZERO_QMIN = 1
  integer, parameter :: Q_MAX_TOO_SMALL = 2
  integer, parameter :: ZERO_XMIN = 3
  integer, parameter :: MASS_MIX = 4
  integer, parameter :: ZERO_SW = 5
  integer, parameter :: ISOSPIN_MIX = 6
  integer, parameter :: WRONG_PRT = 7
  integer, parameter :: MASS_MIX_OUT = 8
  integer, parameter :: NO_EWA = 9
\nwendcode{}\nwbegincode{1136}\moddef{SF ewa: ewa data: TBP}\endmoddef
  procedure :: init => ewa_data_init
\nwendcode{}\nwbegincode{1137}\moddef{SF ewa: procedures}\endmoddef
  subroutine ewa_data_init (data, model, pdg_in, x_min, pt_max, &
        sqrts, keep_momentum, keep_energy, mass)
    class(ewa_data_t), intent(inout) :: data
    class(model_data_t), intent(in), target :: model
    type(pdg_array_t), intent(in) :: pdg_in
    real(default), intent(in) :: x_min, pt_max, sqrts
    logical, intent(in) :: keep_momentum, keep_energy
    real(default), intent(in), optional :: mass
    real(default) :: g, ee
    integer :: n_flv, i
    data%model => model
    if (.not. any (pdg_in .match. &
         [1,2,3,4,5,6,11,13,15,-1,-2,-3,-4,-5,-6,-11,-13,-15])) then
       data%error = WRONG_PRT;  return
    end if   
    n_flv = pdg_array_get_length (pdg_in)
    allocate (data%flv_in (n_flv))
    allocate (data%flv_out(n_flv))
    do i = 1, n_flv
       call data%flv_in(i)%init (pdg_array_get (pdg_in, i), model)
    end do
    data%pt_max = pt_max
    data%sqrts = sqrts
    data%x_min = x_min
    data%x_max = 1
    if (data%x_min == 0) then
       data%error = ZERO_XMIN;  return
    end if    
    select case (char (data%model%get_name ()))
    case ("QCD","QED","Test")
       data%error = NO_EWA;  return
    end select
    ee = data%model%get_real (var_str ("ee"))
    data%sinthw = data%model%get_real (var_str ("sw"))    
    data%costhw = data%model%get_real (var_str ("cw"))        
    data%mZ = data%model%get_real (var_str ("mZ"))
    data%mW = data%model%get_real (var_str ("mW"))    
    if (data%sinthw /= 0) then
       g = ee / data%sinthw
    else
       data%error = ZERO_SW;  return
    end if
    data%cv = g / 2._default
    data%ca = g / 2._default   
    data%coeff = 1._default / (8._default * PI**2)
    data%keep_momentum = keep_momentum
    data%keep_energy = keep_energy
    if (present (mass)) then
       data%mass = mass
       data%m_out = mass
       data%mass_set = .true.
    else
       data%mass = data%flv_in(1)%get_mass ()
       if (any (data%flv_in%get_mass () /= data%mass)) then
          data%error = MASS_MIX;  return
       end if
    end if
  end subroutine ewa_data_init

\eatline
\nwendcode{}\nwbegindocs{1138}Set the vector boson ID for distinguishing $W$ and $Z$ bosons.
\nwenddocs{}\nwbegincode{1139}\moddef{SF ewa: ewa data: TBP}\plusendmoddef
  procedure :: set_id => ewa_set_id
\nwendcode{}\nwbegincode{1140}\moddef{SF ewa: procedures}\plusendmoddef
  subroutine ewa_set_id (data, id)
    class(ewa_data_t), intent(inout) :: data
    integer, intent(in) :: id
    integer :: i, isospin, pdg
    if (.not. allocated (data%flv_in)) &
         call msg_bug ("EWA: incoming particles not set")
    data%id = id
    select case (data%id)
    case (23)
       data%m_out = data%mass
       data%flv_out = data%flv_in
    case (24)
       do i = 1, size (data%flv_in)
          pdg = data%flv_in(i)%get_pdg () 
          isospin = data%flv_in(i)%get_isospin_type ()
          if (isospin > 0) then            
             !!! up-type quark or neutrinos
             if (data%flv_in(i)%is_antiparticle ()) then
                call data%flv_out(i)%init (pdg + 1, data%model)
             else
                call data%flv_out(i)%init (pdg - 1, data%model)
             end if
          else
             !!! down-type quark or lepton
             if (data%flv_in(i)%is_antiparticle ()) then
                call data%flv_out(i)%init (pdg - 1, data%model)
             else
                call data%flv_out(i)%init (pdg + 1, data%model)
             end if
          end if
       end do
       if (.not. data%mass_set) then
          data%m_out = data%flv_out(1)%get_mass ()
          if (any (data%flv_out%get_mass () /= data%m_out)) then
             data%error = MASS_MIX_OUT;  return
          end if
       end if
    end select
  end subroutine ewa_set_id 

\eatline
\nwendcode{}\nwbegindocs{1141}Handle error conditions.  Should always be done after
initialization, unless we are sure everything is ok.
\nwenddocs{}\nwbegincode{1142}\moddef{SF ewa: ewa data: TBP}\plusendmoddef
  procedure :: check => ewa_data_check
\nwendcode{}\nwbegincode{1143}\moddef{SF ewa: procedures}\plusendmoddef
  subroutine ewa_data_check (data)
    class(ewa_data_t), intent(in) :: data
    select case (data%error)
    case (WRONG_PRT)
       call msg_fatal ("EWA structure function only accessible for " &
            // "SM quarks and leptons.")
    case (NO_EWA)
       call msg_fatal ("EWA structure function not available for model " &
            // char (data%model%get_name ()))
    case (ZERO_SW)
       call msg_fatal ("EWA: Vanishing value of sin(theta_w)")
    case (ZERO_QMIN)
       call msg_fatal ("EWA: Particle mass is zero")
    case (Q_MAX_TOO_SMALL)
       call msg_fatal ("EWA: Particle mass exceeds Qmax")
    case (ZERO_XMIN)
       call msg_fatal ("EWA: x_min must be larger than zero")
    case (MASS_MIX)
       call msg_fatal ("EWA: incoming particle masses must be uniform")
    case (MASS_MIX_OUT)
       call msg_fatal ("EWA: outgoing particle masses must be uniform")
    case (ISOSPIN_MIX)
       call msg_fatal ("EWA: incoming particle isospins must be uniform")
    end select
  end subroutine ewa_data_check

\eatline
\nwendcode{}\nwbegindocs{1144}Output
\nwenddocs{}\nwbegincode{1145}\moddef{SF ewa: ewa data: TBP}\plusendmoddef
  procedure :: write => ewa_data_write
\nwendcode{}\nwbegincode{1146}\moddef{SF ewa: procedures}\plusendmoddef
  subroutine ewa_data_write (data, unit, verbose) 
    class(ewa_data_t), intent(in) :: data
    integer, intent(in), optional :: unit
    logical, intent(in), optional :: verbose
    integer :: u, i
    u = given_output_unit (unit);  if (u < 0)  return
    write (u, "(1x,A)") "EWA data:"
    if (allocated (data%flv_in) .and. allocated (data%flv_out)) then
       write (u, "(3x,A)", advance="no") "  flavor(in)  =  "
       do i = 1, size (data%flv_in)
          if (i > 1)  write (u, "(',',1x)", advance="no")
          call data%flv_in(i)%write (u)
       end do
       write (u, *)
       write (u, "(3x,A)", advance="no") "  flavor(out) =  "
       do i = 1, size (data%flv_out)
          if (i > 1)  write (u, "(',',1x)", advance="no")
          call data%flv_out(i)%write (u)
       end do
       write (u, *)
       write (u, "(3x,A," // FMT_19 // ")") "  x_min     = ", data%x_min
       write (u, "(3x,A," // FMT_19 // ")") "  x_max     = ", data%x_max
       write (u, "(3x,A," // FMT_19 // ")") "  pt_max    = ", data%pt_max
       write (u, "(3x,A," // FMT_19 // ")") "  sqrts     = ", data%sqrts    
       write (u, "(3x,A," // FMT_19 // ")") "  mass      = ", data%mass
       write (u, "(3x,A," // FMT_19 // ")") "  cv        = ", data%cv
       write (u, "(3x,A," // FMT_19 // ")") "  ca        = ", data%ca
       write (u, "(3x,A," // FMT_19 // ")") "  coeff     = ", data%coeff
       write (u, "(3x,A," // FMT_19 // ")") "  costhw    = ", data%costhw
       write (u, "(3x,A," // FMT_19 // ")") "  sinthw    = ", data%sinthw    
       write (u, "(3x,A," // FMT_19 // ")") "  mZ        = ", data%mZ
       write (u, "(3x,A," // FMT_19 // ")") "  mW        = ", data%mW
       write (u, "(3x,A,L2)")      "  keep_mom. = ", data%keep_momentum
       write (u, "(3x,A,L2)")      "  keep_en.  = ", data%keep_energy 
       write (u, "(3x,A,I2)")      "  PDG (VB)  = ", data%id
    else
       write (u, "(3x,A)") "[undefined]"
    end if       
  end subroutine ewa_data_write

\eatline
\nwendcode{}\nwbegindocs{1147}The number of parameters is one for collinear splitting, in case one
of the options {\tt{}keep{\char95}energy} or {\tt{}keep{\char95}momentum} is set, we take
the recoil into account.
\nwenddocs{}\nwbegincode{1148}\moddef{SF ewa: ewa data: TBP}\plusendmoddef
  procedure :: get_n_par => ewa_data_get_n_par
\nwendcode{}\nwbegincode{1149}\moddef{SF ewa: procedures}\plusendmoddef
  function ewa_data_get_n_par (data) result (n)
    class(ewa_data_t), intent(in) :: data
    integer :: n
    if (data%keep_energy .or. data%keep_momentum) then
       n = 3
    else
       n = 1
    end if
  end function ewa_data_get_n_par
  
\eatline
\nwendcode{}\nwbegindocs{1150}Return the outgoing particles PDG codes. This depends, whether this
is a charged-current or neutral-current interaction.
\nwenddocs{}\nwbegincode{1151}\moddef{SF ewa: ewa data: TBP}\plusendmoddef
  procedure :: get_pdg_out => ewa_data_get_pdg_out
\nwendcode{}\nwbegincode{1152}\moddef{SF ewa: procedures}\plusendmoddef
  subroutine ewa_data_get_pdg_out (data, pdg_out)
    class(ewa_data_t), intent(in) :: data
    type(pdg_array_t), dimension(:), intent(inout) :: pdg_out
    integer, dimension(:), allocatable :: pdg1
    integer :: i, n_flv
    if (allocated (data%flv_out)) then
       n_flv = size (data%flv_out)
    else
       n_flv = 0
    end if
    allocate (pdg1 (n_flv))
    do i = 1, n_flv
       pdg1(i) = data%flv_out(i)%get_pdg ()
    end do
    pdg_out(1) = pdg1
  end subroutine ewa_data_get_pdg_out
  
\eatline
\nwendcode{}\nwbegindocs{1153}Allocate the interaction record.
\nwenddocs{}\nwbegincode{1154}\moddef{SF ewa: ewa data: TBP}\plusendmoddef
  procedure :: allocate_sf_int => ewa_data_allocate_sf_int
\nwendcode{}\nwbegincode{1155}\moddef{SF ewa: procedures}\plusendmoddef
  subroutine ewa_data_allocate_sf_int (data, sf_int)
    class(ewa_data_t), intent(in) :: data
    class(sf_int_t), intent(inout), allocatable :: sf_int
    allocate (ewa_t :: sf_int)
  end subroutine ewa_data_allocate_sf_int
  
\eatline
\nwendcode{}\nwbegindocs{1156}\nwdocspar
\subsection{The EWA object}
The {\tt{}ewa{\char95}t} data type is a $1\to 2$ interaction.  We should be able
to handle several flavors in parallel, since EWA is not necessarily
applied immediately after beam collision:  $W/Z$ bosons may be radiated
from quarks.  In that case, the partons are massless and $q_{\rm min}$
applies instead, so we do not need to generate several kinematical
configurations in parallel.

The particles are ordered as (incoming, radiated, W/Z), where the
W/Z initiates the hard interaction.

In the case of EPA, we generated an unpolarized photon and transferred
initial polarization to the radiated parton.  Color is transferred in
the same way. I do not know whether the same can/should be done for
EWA, as the structure functions depend on the W/Z polarization. If we
are having $Z$ bosons, both up- and down-type fermions can
participate. Otherwise, with a $W^+$ an up-type fermion is transferred
to a down-type fermion, and the other way round.
\nwenddocs{}\nwbegincode{1157}\moddef{SF ewa: types}\plusendmoddef
  type, extends (sf_int_t) :: ewa_t
     type(ewa_data_t), pointer :: data => null ()
     real(default) :: x  = 0
     real(default) :: xb = 0
     integer :: n_me = 0
     real(default), dimension(:), allocatable :: cv
     real(default), dimension(:), allocatable :: ca
   contains
   \LA{}SF ewa: ewa: TBP\RA{}
  end type ewa_t 
  
\eatline
\nwendcode{}\nwbegindocs{1158}Type string: has to be here, but there is no string variable on which EWA
depends. Hence, a dummy routine.
\nwenddocs{}\nwbegincode{1159}\moddef{SF ewa: ewa: TBP}\endmoddef
  procedure :: type_string => ewa_type_string
\nwendcode{}\nwbegincode{1160}\moddef{SF ewa: procedures}\plusendmoddef
  function ewa_type_string (object) result (string)
    class(ewa_t), intent(in) :: object
    type(string_t) :: string
    if (associated (object%data)) then
       string = "EWA: equivalent W/Z approx." 
    else
       string = "EWA: [undefined]"
    end if
  end function ewa_type_string
  
\eatline
\nwendcode{}\nwbegindocs{1161}Output.  Call the interaction routine after displaying the configuration.
\nwenddocs{}\nwbegincode{1162}\moddef{SF ewa: ewa: TBP}\plusendmoddef
  procedure :: write => ewa_write
\nwendcode{}\nwbegincode{1163}\moddef{SF ewa: procedures}\plusendmoddef
  subroutine ewa_write (object, unit, testflag)
    class(ewa_t), intent(in) :: object
    integer, intent(in), optional :: unit
    logical, intent(in), optional :: testflag
    integer :: u
    u = given_output_unit (unit)
    if (associated (object%data)) then
       call object%data%write (u)
       if (object%status >= SF_DONE_KINEMATICS) then
          write (u, "(1x,A)")  "SF parameters:"
          write (u, "(3x,A," // FMT_17 // ")")  "x =", object%x
       end if
       call object%base_write (u, testflag)
    else
       write (u, "(1x,A)")  "EWA data: [undefined]"
    end if    
  end subroutine ewa_write
    
\eatline
\nwendcode{}\nwbegindocs{1164}The current implementation requires uniform isospin for all incoming
particles, therefore we need to probe only the first one.
\nwenddocs{}\nwbegincode{1165}\moddef{SF ewa: ewa: TBP}\plusendmoddef
  procedure :: init => ewa_init
\nwendcode{}\nwbegincode{1166}\moddef{SF ewa: procedures}\plusendmoddef
  subroutine ewa_init (sf_int, data)
    class(ewa_t), intent(out) :: sf_int
    class(sf_data_t), intent(in), target :: data
    type(quantum_numbers_mask_t), dimension(3) :: mask
    integer, dimension(3) :: hel_lock
    type(polarization_t) :: pol
    type(quantum_numbers_t), dimension(1) :: qn_fc, qn_hel, qn_fc_fin
    type(flavor_t) :: flv_z, flv_wp, flv_wm
    type(color_t) :: col0
    type(quantum_numbers_t) :: qn_z, qn_wp, qn_wm, qn, qn_rad, qn_w
    type(state_iterator_t) :: it_hel
    integer :: i, isospin
    select type (data)
    type is (ewa_data_t)   
       mask = quantum_numbers_mask (.false., .false., &
            mask_h = [.false., .false., .true.])
       hel_lock = [2, 1, 0]
       call col0%init ()
       select case (data%id)
       case (23)
          !!! Z boson, flavor is not changing    
          call sf_int%base_init (mask, [data%mass**2], [data%mass**2], &
               [data%mZ**2], hel_lock = hel_lock)
          sf_int%data => data          
          call flv_z%init (Z_BOSON, data%model)
          call qn_z%init (flv_z, col0)
          do i = 1, size (data%flv_in)
             call polarization_init_generic (pol, data%flv_in(i))
             call qn_fc(1)%init ( &
                  flv = data%flv_in(i), &
                  col = color_from_flavor (data%flv_in(i), 1))
             call it_hel%init (pol%state)
             do while (it_hel%is_valid ())
                qn_hel = it_hel%get_quantum_numbers ()
                qn = qn_hel(1) .merge. qn_fc(1)
                qn_rad = qn
                call qn_rad%tag_radiated ()
                call interaction_add_state &
                     (sf_int%interaction_t, [qn, qn_rad, qn_z])
                call it_hel%advance ()
             end do
             call polarization_final (pol)
          end do
       case (24)    
          call sf_int%base_init (mask, [data%mass**2], [data%m_out**2], &
               [data%mW**2], hel_lock = hel_lock)
          sf_int%data => data                       
          call flv_wp%init (W_BOSON, data%model)
          call flv_wm%init (- W_BOSON, data%model)
          call qn_wp%init (flv_wp, col0)
          call qn_wm%init (flv_wm, col0)
          do i = 1, size (data%flv_in)
             isospin = data%flv_in(i)%get_isospin_type ()
             if (isospin > 0) then            
                !!! up-type quark or neutrinos
                if (data%flv_in(i)%is_antiparticle ()) then
                   qn_w = qn_wm
                else
                   qn_w = qn_wp
                end if
             else
                !!! down-type quark or lepton
                if (data%flv_in(i)%is_antiparticle ()) then
                   qn_w = qn_wp
                else
                   qn_w = qn_wm
                end if
             end if
             call polarization_init_generic (pol, data%flv_in(i))
             call qn_fc(1)%init ( &
                  flv = data%flv_in(i), &
                  col = color_from_flavor (data%flv_in(i), 1))
             call qn_fc_fin(1)%init ( &
                  flv = data%flv_out(i), &
                  col = color_from_flavor (data%flv_out(i), 1))
             call it_hel%init (pol%state)
             do while (it_hel%is_valid ())
                qn_hel = it_hel%get_quantum_numbers ()
                qn = qn_hel(1) .merge. qn_fc(1)
                qn_rad = qn_hel(1) .merge. qn_fc_fin(1)
                call qn_rad%tag_radiated ()
                call interaction_add_state &
                     (sf_int%interaction_t, [qn, qn_rad, qn_w])
                call it_hel%advance ()
             end do
             call polarization_final (pol)    
          end do
       case default
          call msg_fatal ("EWA initialization failed: wrong particle type.")
       end select
       call interaction_freeze (sf_int%interaction_t)
       if (data%keep_momentum) then
          if (data%keep_energy) then
             call msg_fatal ("EWA: momentum and energy" // &
                  "cannot be simultaneously conserved.")
          else          
             sf_int%on_shell_mode = KEEP_MOMENTUM
          end if
       else 
          if (data%keep_energy) then
             sf_int%on_shell_mode = KEEP_ENERGY
          end if
       end if
       call sf_int%set_incoming ([1])
       call sf_int%set_radiated ([2])
       call sf_int%set_outgoing ([3])
    end select
  end subroutine ewa_init
    
\eatline
\nwendcode{}\nwbegindocs{1167}Prepare the coupling arrays.  This is separate from the previous routine since
the state matrix may be helicity-contracted.
\nwenddocs{}\nwbegincode{1168}\moddef{SF ewa: ewa: TBP}\plusendmoddef
  procedure :: setup_constants => ewa_setup_constants
\nwendcode{}\nwbegincode{1169}\moddef{SF ewa: procedures}\plusendmoddef
  subroutine ewa_setup_constants (sf_int)
    class(ewa_t), intent(inout) :: sf_int
    type(state_iterator_t) :: it
    type(flavor_t) :: flv
    real(default) :: q, t3
    integer :: i
    sf_int%n_me = interaction_get_n_matrix_elements (sf_int%interaction_t) 
    allocate (sf_int%cv (sf_int%n_me))
    allocate (sf_int%ca (sf_int%n_me))
    associate (data => sf_int%data)
      select case (data%id)
      case (23)
         call it%init (interaction_get_state_matrix_ptr (sf_int%interaction_t))
         do while (it%is_valid ())
            i = it%get_me_index ()
            flv = it%get_flavor (1)
            q = flv%get_charge ()
            t3 = flv%get_isospin ()
            if (flv%is_antiparticle ()) then
               sf_int%cv(i) = - data%cv &
                    * (t3 - 2._default * q * data%sinthw**2) / data%costhw
               sf_int%ca(i) = data%ca *  t3 / data%costhw    
            else          
               sf_int%cv(i) = data%cv &
                    * (t3 - 2._default * q * data%sinthw**2) / data%costhw
               sf_int%ca(i) = data%ca *  t3 / data%costhw    
            end if
            call it%advance ()
         end do
      case (24)
         call it%init (interaction_get_state_matrix_ptr (sf_int%interaction_t))
         do while (it%is_valid ())
            i = it%get_me_index ()
            flv = it%get_flavor (1)
            if (flv%is_antiparticle ()) then
               sf_int%cv(i) = data%cv / sqrt(2._default)
               sf_int%ca(i) = - data%ca / sqrt(2._default)
            else          
               sf_int%cv(i) = data%cv / sqrt(2._default)
               sf_int%ca(i) = data%ca / sqrt(2._default)
            end if
            call it%advance ()
         end do
      end select
    end associate
    sf_int%status = SF_INITIAL
  end subroutine ewa_setup_constants
  
\eatline
\nwendcode{}\nwbegindocs{1170}\nwdocspar
\subsection{Kinematics}
Set kinematics.  The EWA structure function allows for a
straightforward mapping of the unit interval.  So, to leading order,
the structure function value is unity, but the $x$ value is
transformed.  Higher orders affect the function value.
If {\tt{}map} is unset, the $r$ and $x$ values coincide, and the Jacobian
$f(r)$ is trivial. 

If {\tt{}map} is set, the exponential mapping for the $1/x$ singularity
discussed above is applied. 
\nwenddocs{}\nwbegincode{1171}\moddef{SF ewa: ewa: TBP}\plusendmoddef
  procedure :: complete_kinematics => ewa_complete_kinematics
\nwendcode{}\nwbegincode{1172}\moddef{SF ewa: procedures}\plusendmoddef
  subroutine ewa_complete_kinematics (sf_int, x, f, r, rb, map)
    class(ewa_t), intent(inout) :: sf_int
    real(default), dimension(:), intent(out) :: x
    real(default), intent(out) :: f
    real(default), dimension(:), intent(in) :: r
    real(default), dimension(:), intent(in) :: rb
    logical, intent(in) :: map
    real(default) :: xb1, e_1
    real(default) :: x0, x1, lx0, lx1, lx
    e_1 = energy (interaction_get_momentum (sf_int%interaction_t, 1))
    if (sf_int%data%keep_momentum .or. sf_int%data%keep_energy) then
       select case (sf_int%data%id)
       case (23)
          x0 = max (sf_int%data%x_min, sf_int%data%mz / e_1)
       case (24)
          x0 = max (sf_int%data%x_min, sf_int%data%mw / e_1)
       end select
    else 
       x0 = sf_int%data%x_min
    end if
    x1 = sf_int%data%x_max
    if ( x0 >= x1) then
       f = 0
       sf_int%status = SF_FAILED_KINEMATICS
       return
    end if
    if (map) then
       lx0 = log (x0)
       lx1 = log (x1)
       lx = lx1 * r(1) + lx0 * rb(1)
       x(1) = exp(lx)       
       f = x(1) * (lx1 - lx0)
    else       
       x(1) = r(1)
       if (x0 < x(1) .and. x(1) < x1) then
          f = 1
       else
          sf_int%status = SF_FAILED_KINEMATICS
          f = 0
          return
       end if
    end if
    xb1 = 1 - x(1)
    if (size(x) == 3)  x(2:3) = r(2:3)
    call sf_int%split_momentum (x, xb1)
    select case (sf_int%status)
    case (SF_DONE_KINEMATICS)
       sf_int%x  = x(1)
       sf_int%xb = xb1
    case (SF_FAILED_KINEMATICS)
       sf_int%x  = 0
       sf_int%xb = 0
       f = 0
    end select
  end subroutine ewa_complete_kinematics

\eatline
\nwendcode{}\nwbegindocs{1173}Compute inverse kinematics.  Here, we start with the $x$ array and
compute the ``input'' $r$ values and the Jacobian $f$.  After this, we
can set momenta by the same formula as for normal kinematics.
\nwenddocs{}\nwbegincode{1174}\moddef{SF ewa: ewa: TBP}\plusendmoddef
  procedure :: inverse_kinematics => ewa_inverse_kinematics
\nwendcode{}\nwbegincode{1175}\moddef{SF ewa: procedures}\plusendmoddef
  subroutine ewa_inverse_kinematics (sf_int, x, f, r, rb, map, set_momenta)
    class(ewa_t), intent(inout) :: sf_int
    real(default), dimension(:), intent(in) :: x
    real(default), intent(out) :: f
    real(default), dimension(:), intent(out) :: r
    real(default), dimension(:), intent(out) :: rb
    logical, intent(in) :: map
    logical, intent(in), optional :: set_momenta
    real(default) :: x0, x1, lx0, lx1, lx, e_1
    logical :: set_mom
    set_mom = .false.;  if (present (set_momenta))  set_mom = set_momenta
    e_1 = energy (interaction_get_momentum (sf_int%interaction_t, 1))    
    if (sf_int%data%keep_momentum .or. sf_int%data%keep_energy) then
       select case (sf_int%data%id)
       case (23)
          x0 = max (sf_int%data%x_min, sf_int%data%mz / e_1)
       case (24)
          x0 = max (sf_int%data%x_min, sf_int%data%mw / e_1)
       end select
    else 
       x0 = sf_int%data%x_min
    end if
    x1 = sf_int%data%x_max
    if (map) then
       lx0 = log (x0)
       lx1 = log (x1)
       lx = log (x(1))
       r(1)  = (lx - lx0) / (lx1 - lx0)
       rb(1) = (lx1 - lx) / (lx1 - lx0)
       f = x(1) * (lx1 - lx0)
    else
       r (1) = x(1)
       rb(1) = 1 - x(1)
       if (x0 < x(1) .and. x(1) < x1) then
          f = 1
       else
          f = 0
       end if
    end if
    if (size(r) == 3) then 
       r (2:3) = x(2:3)
       rb(2:3) = 1 - x(2:3)
    end if
    if (set_mom) then
       call sf_int%split_momentum (x, sf_int%xb)
       select case (sf_int%status)
       case (SF_DONE_KINEMATICS)
          sf_int%x  = x(1)
          sf_int%xb = 1 - x(1)
       case (SF_FAILED_KINEMATICS)
          sf_int%x = 0
          f = 0
       end select
    end if
  end subroutine ewa_inverse_kinematics

\eatline
\nwendcode{}\nwbegindocs{1176}\nwdocspar
\subsection{EWA application}
For EWA, we can compute kinematics and function value in a single
step.  This function works on a single beam, assuming that the input
momentum has been set.  We need four random numbers as input: one for
$x$, one for $Q^2$, and two for the polar and azimuthal angles.
Alternatively, we can skip $p_T$ generation; in this case, we only
need one.

For obtaining splitting kinematics, we rely on the assumption that all
in-particles are mass-degenerate (or there is only one), so the
generated $x$ values are identical.
\nwenddocs{}\nwbegincode{1177}\moddef{SF ewa: ewa: TBP}\plusendmoddef
  procedure :: apply => ewa_apply
\nwendcode{}\nwbegincode{1178}\moddef{SF ewa: procedures}\plusendmoddef
  subroutine ewa_apply (sf_int, scale)
    class(ewa_t), intent(inout) :: sf_int
    real(default), intent(in) :: scale
    real(default) :: x, xb, pt2, c1, c2
    real(default) :: cv, ca
    real(default) :: f, fm, fp, fL
    integer :: i
    associate (data => sf_int%data)     
      x  = sf_int%x      
      xb = sf_int%xb
      pt2 = min ((data%pt_max)**2, (xb * data%sqrts / 2)**2)
      select case (data%id)
      case (23)
         !!! Z boson structure function
         c1 = log (1 + pt2 / (xb * (data%mZ)**2))
         c2 = 1 / (1 + (xb * (data%mZ)**2) / pt2)
      case (24)
         !!! W boson structure function
         c1 = log (1 + pt2 / (xb * (data%mW)**2))
         c2 = 1 / (1 + (xb * (data%mW)**2) / pt2)
      end select
      do i = 1, sf_int%n_me
         cv = sf_int%cv(i)
         ca = sf_int%ca(i)
         fm = data%coeff * &
              ((cv + ca)**2 + ((cv - ca) * xb)**2) * (c1 - c2) / (2 * x)
         fp = data%coeff * &
              ((cv - ca)**2 + ((cv + ca) * xb)**2) * (c1 - c2) / (2 * x)
         fL = data%coeff * &
              (cv**2 + ca**2) * (2 * xb / x) * c2       
         f = fp + fm + fL
         if (f /= 0) then
            fp = fp / f
            fm = fm / f
            fL = fL / f      
         end if
         call interaction_set_matrix_element &
              (sf_int%interaction_t, i, cmplx (f, kind=default))
      end do
    end associate
    sf_int%status = SF_EVALUATED
  end subroutine ewa_apply

\eatline
\nwendcode{}\nwbegindocs{1179}\nwdocspar
\subsection{Unit tests}
\nwenddocs{}\nwbegincode{1180}\moddef{SF ewa: public}\plusendmoddef
  public :: sf_ewa_test
\nwendcode{}\nwbegincode{1181}\moddef{SF ewa: tests}\endmoddef
  subroutine sf_ewa_test (u, results)
    integer, intent(in) :: u
    type(test_results_t), intent(inout) :: results
  \LA{}SF ewa: execute tests\RA{}
  end subroutine sf_ewa_test
  
\eatline
\nwendcode{}\nwbegindocs{1182}\nwdocspar
\subsubsection{Test structure function data}
Construct and display a test structure function data object.
\nwenddocs{}\nwbegincode{1183}\moddef{SF ewa: execute tests}\endmoddef
  call test (sf_ewa_1, "sf_ewa_1", &
       "structure function configuration", &
       u, results)
\nwendcode{}\nwbegincode{1184}\moddef{SF ewa: tests}\plusendmoddef
  subroutine sf_ewa_1 (u)
    integer, intent(in) :: u
    type(model_data_t), target :: model
    type(pdg_array_t) :: pdg_in
    type(pdg_array_t), dimension(1) :: pdg_out
    integer, dimension(:), allocatable :: pdg1
    class(sf_data_t), allocatable :: data
    
    write (u, "(A)")  "* Test output: sf_ewa_1"
    write (u, "(A)")  "*   Purpose: initialize and display &
         &test structure function data"
    write (u, "(A)")
    
    write (u, "(A)")  "* Create empty data object"
    write (u, "(A)")

    call model%init_sm_test ()
    pdg_in = 2

    allocate (ewa_data_t :: data)
    call data%write (u)

    write (u, "(A)")
    write (u, "(A)")  "* Initialize for Z boson"
    write (u, "(A)")

    select type (data)
    type is (ewa_data_t)
       call data%init (model, pdg_in, 0.01_default, &
            500._default, 5000._default, .false., .false.)
       call data%set_id (23)
    end select

    call data%write (u)

    write (u, "(A)")

    write (u, "(1x,A)")  "Outgoing particle codes:"
    call data%get_pdg_out (pdg_out)
    pdg1 = pdg_out(1)
    write (u, "(2x,99(1x,I0))")  pdg1
        
    write (u, "(A)")
    write (u, "(A)")  "* Initialize for W boson"
    write (u, "(A)")
    
    deallocate (data)
    allocate (ewa_data_t :: data)
    select type (data)
    type is (ewa_data_t)
       call data%init (model, pdg_in, 0.01_default, &
            500._default, 5000._default, .false., .false.)
       call data%set_id (24)
    end select

    call data%write (u)

    write (u, "(A)")

    write (u, "(1x,A)")  "Outgoing particle codes:"
    call data%get_pdg_out (pdg_out)
    pdg1 = pdg_out(1)
    write (u, "(2x,99(1x,I0))")  pdg1
            
    call model%final ()
    
    write (u, "(A)")
    write (u, "(A)")  "* Test output end: sf_ewa_1"

  end subroutine sf_ewa_1

\eatline
\nwendcode{}\nwbegindocs{1185}\nwdocspar
\subsubsection{Test and probe structure function}
Construct and display a structure function object based on the EWA
structure function.
\nwenddocs{}\nwbegincode{1186}\moddef{SF ewa: execute tests}\plusendmoddef
  call test (sf_ewa_2, "sf_ewa_2", &
       "structure function instance", &
       u, results)
\nwendcode{}\nwbegincode{1187}\moddef{SF ewa: tests}\plusendmoddef
  subroutine sf_ewa_2 (u)
    integer, intent(in) :: u
    type(model_data_t), target :: model
    type(flavor_t) :: flv
    type(pdg_array_t) :: pdg_in
    class(sf_data_t), allocatable, target :: data
    class(sf_int_t), allocatable :: sf_int
    type(vector4_t) :: k
    type(vector4_t), dimension(2) :: q
    real(default) :: E
    real(default), dimension(:), allocatable :: r, rb, x
    real(default) :: f
    
    write (u, "(A)")  "* Test output: sf_ewa_2"
    write (u, "(A)")  "*   Purpose: initialize and fill &
         &test structure function object"
    write (u, "(A)")
    
    write (u, "(A)")  "* Initialize configuration data"
    write (u, "(A)")

    call model%init_sm_test ()
    call flv%init (2, model)
    pdg_in = 2

    call reset_interaction_counter ()
    
    allocate (ewa_data_t :: data)
    select type (data)
    type is (ewa_data_t)
       call data%init (model, pdg_in, 0.01_default, &
            500._default, 3000._default, .false., .false.)
       call data%set_id (24)
    end select
       
    write (u, "(A)")  "* Initialize structure-function object"
    write (u, "(A)")
    
    call data%allocate_sf_int (sf_int)
    call sf_int%init (data)
    call sf_int%set_beam_index ([1])
    call sf_int%setup_constants ()
    
    call sf_int%write (u)

    write (u, "(A)")
    write (u, "(A)")  "* Initialize incoming momentum with E=1500"
    write (u, "(A)")
    E = 1500
    k = vector4_moving (E, sqrt (E**2 - flv%get_mass ()**2), 3)
    call pacify (k, 1e-10_default)
    call vector4_write (k, u)
    call sf_int%seed_kinematics ([k])
    
    write (u, "(A)")
    write (u, "(A)")  "* Set kinematics for r=0.4, no EWA mapping, collinear"
    write (u, "(A)")
    
    allocate (r (data%get_n_par ()))
    allocate (rb(size (r)))
    allocate (x (size (r)))
    
    r = 0.4_default
    rb = 1 - r
    call sf_int%complete_kinematics (x, f, r, rb, map=.false.)
    
    write (u, "(A,9(1x,F10.7))")  "r =", r
    write (u, "(A,9(1x,F10.7))")  "rb=", rb
    write (u, "(A,9(1x,F10.7))")  "x =", x
    write (u, "(A,9(1x,F10.7))")  "f =", f
    
    write (u, "(A)")
    write (u, "(A)")  "* Recover x from momenta"
    write (u, "(A)")
    
    q = interaction_get_momenta (sf_int%interaction_t, outgoing=.true.)
    call sf_int%final ()
    deallocate (sf_int)
    
    call data%allocate_sf_int (sf_int)
    call sf_int%init (data)
    call sf_int%set_beam_index ([1])
    call sf_int%setup_constants ()
    
    call sf_int%seed_kinematics ([k])
    call interaction_set_momenta (sf_int%interaction_t, q, outgoing=.true.)
    call sf_int%recover_x (x)
    call sf_int%inverse_kinematics (x, f, r, rb, map=.false., &
         set_momenta=.true.)
    
    write (u, "(A,9(1x,F10.7))")  "r =", r
    write (u, "(A,9(1x,F10.7))")  "rb=", rb
    write (u, "(A,9(1x,F10.7))")  "x =", x
    write (u, "(A,9(1x,F10.7))")  "f =", f    
    
    write (u, "(A)")
    write (u, "(A)")  "* Evaluate EWA structure function"
    write (u, "(A)")
    
    call sf_int%apply (scale = 100._default)
    call sf_int%write (u)
        
    write (u, "(A)")
    write (u, "(A)")  "* Cleanup"
    
    call sf_int%final ()
    call model%final ()

    write (u, "(A)")
    write (u, "(A)")  "* Test output end: sf_ewa_2"
    
  end subroutine sf_ewa_2
  
\eatline
\nwendcode{}\nwbegindocs{1188}\nwdocspar
\subsubsection{Standard mapping}
Construct and display a structure function object based on the EWA
structure function, applying the standard single-particle mapping.
\nwenddocs{}\nwbegincode{1189}\moddef{SF ewa: execute tests}\plusendmoddef
  call test (sf_ewa_3, "sf_ewa_3", &
       "apply mapping", &
       u, results)
\nwendcode{}\nwbegincode{1190}\moddef{SF ewa: tests}\plusendmoddef
  subroutine sf_ewa_3 (u)
    integer, intent(in) :: u
    type(model_data_t), target :: model
    type(flavor_t) :: flv
    type(pdg_array_t) :: pdg_in
    class(sf_data_t), allocatable, target :: data
    class(sf_int_t), allocatable :: sf_int
    type(vector4_t) :: k
    type(vector4_t), dimension(2) :: q
    real(default) :: E
    real(default), dimension(:), allocatable :: r, rb, x
    real(default) :: f
    
    write (u, "(A)")  "* Test output: sf_ewa_3"
    write (u, "(A)")  "*   Purpose: initialize and fill &
         &test structure function object"
    write (u, "(A)")
    
    write (u, "(A)")  "* Initialize configuration data"
    write (u, "(A)")

    call model%init_sm_test ()
    call flv%init (2, model)
    pdg_in = 2

    call reset_interaction_counter ()
    
    allocate (ewa_data_t :: data)
    select type (data)
    type is (ewa_data_t)
       call data%init (model, pdg_in, 0.01_default, &
            500._default, 3000._default, .false., .false.)
       call data%set_id (24)
    end select
    
    write (u, "(A)")  "* Initialize structure-function object"
    write (u, "(A)")    
    
    call data%allocate_sf_int (sf_int)
    call sf_int%init (data)
    call sf_int%set_beam_index ([1])
    call sf_int%setup_constants ()
    
    call sf_int%write (u)

    write (u, "(A)")
    write (u, "(A)")  "* Initialize incoming momentum with E=1500"
    write (u, "(A)")
    E = 1500
    k = vector4_moving (E, sqrt (E**2 - flv%get_mass ()**2), 3)
    call pacify (k, 1e-10_default)
    call vector4_write (k, u)
    call sf_int%seed_kinematics ([k])

    write (u, "(A)")
    write (u, "(A)")  "* Set kinematics for r=0.4, with EWA mapping, collinear"
    write (u, "(A)")
    
    allocate (r (data%get_n_par ()))
    allocate (rb(size (r)))
    allocate (x (size (r)))
    
    r = 0.4_default
    rb = 1 - r
    call sf_int%complete_kinematics (x, f, r, rb, map=.true.)
    
    write (u, "(A,9(1x,F10.7))")  "r =", r
    write (u, "(A,9(1x,F10.7))")  "rb=", rb
    write (u, "(A,9(1x,F10.7))")  "x =", x
    write (u, "(A,9(1x,F10.7))")  "f =", f
    
    write (u, "(A)")
    write (u, "(A)")  "* Recover x from momenta"
    write (u, "(A)")
    
    q = interaction_get_momenta (sf_int%interaction_t, outgoing=.true.)
    call sf_int%final ()
    deallocate (sf_int)
    
    call data%allocate_sf_int (sf_int)
    call sf_int%init (data)
    call sf_int%set_beam_index ([1])
    call sf_int%setup_constants ()
    
    call sf_int%seed_kinematics ([k])
    call interaction_set_momenta (sf_int%interaction_t, q, outgoing=.true.)
    call sf_int%recover_x (x)
    call sf_int%inverse_kinematics (x, f, r, rb, map=.true., &
         set_momenta=.true.)
    
    write (u, "(A,9(1x,F10.7))")  "r =", r
    write (u, "(A,9(1x,F10.7))")  "rb=", rb
    write (u, "(A,9(1x,F10.7))")  "x =", x
    write (u, "(A,9(1x,F10.7))")  "f =", f

    write (u, "(A)")
    write (u, "(A)")  "* Evaluate EWA structure function"
    write (u, "(A)")
    
    call sf_int%apply (scale = 100._default)
    call sf_int%write (u)
    
    write (u, "(A)")
    write (u, "(A)")  "* Cleanup"
    
    call sf_int%final ()
    call model%final ()
    
    write (u, "(A)")
    write (u, "(A)")  "* Test output end: sf_ewa_3"

  end subroutine sf_ewa_3

\eatline
\nwendcode{}\nwbegindocs{1191}\nwdocspar
\subsubsection{Non-collinear case}
Construct and display a structure function object based on the EPA
structure function.
\nwenddocs{}\nwbegincode{1192}\moddef{SF ewa: execute tests}\plusendmoddef
  call test (sf_ewa_4, "sf_ewa_4", &
       "non-collinear", &
       u, results)
\nwendcode{}\nwbegincode{1193}\moddef{SF ewa: tests}\plusendmoddef
  subroutine sf_ewa_4 (u)
    integer, intent(in) :: u
    type(model_data_t), target :: model
    type(flavor_t) :: flv
    type(pdg_array_t) :: pdg_in
    class(sf_data_t), allocatable, target :: data
    class(sf_int_t), allocatable :: sf_int
    type(vector4_t) :: k
    type(vector4_t), dimension(2) :: q
    real(default) :: E
    real(default), dimension(:), allocatable :: r, rb, x
    real(default) :: f
    
    write (u, "(A)")  "* Test output: sf_ewa_4"
    write (u, "(A)")  "*   Purpose: initialize and fill &
         &test structure function object"
    write (u, "(A)")
    
    write (u, "(A)")  "* Initialize configuration data"
    write (u, "(A)")

    call modeL%init_sm_test ()
    call flv%init (2, model)
    pdg_in = 2

    call reset_interaction_counter ()

    allocate (ewa_data_t :: data)
    select type (data)
    type is (ewa_data_t)
       call data%init (model, pdg_in, 0.01_default, &
            500._default, 3000.0_default, .false., .true.)
       call data%set_id (24)
    end select    

    write (u, "(A)")  "* Initialize structure-function object"
    write (u, "(A)")
        
    call data%allocate_sf_int (sf_int)
    call sf_int%init (data)
    call sf_int%set_beam_index ([1])
    call sf_int%setup_constants ()

    write (u, "(A)")  "* Initialize incoming momentum with E=1500"
    write (u, "(A)")
    E = 1500
    k = vector4_moving (E, sqrt (E**2 - flv%get_mass ()**2), 3)
    call pacify (k, 1e-10_default)
    call vector4_write (k, u)
    call sf_int%seed_kinematics ([k])        
    
    write (u, "(A)")
    write (u, "(A)")  "* Set kinematics for r=0.5/0.5/0.25, with EWA mapping, "
    write (u, "(A)")  "          non-coll., keeping energy"
    write (u, "(A)")
    
    allocate (r (data%get_n_par ()))
    allocate (rb(size (r)))
    allocate (x (size (r)))
    
    r = [0.5_default, 0.5_default, 0.25_default]
    rb = 1 - r
    sf_int%on_shell_mode = KEEP_ENERGY
    call sf_int%complete_kinematics (x, f, r, rb, map=.true.)
    call interaction_pacify_momenta (sf_int%interaction_t, 1e-10_default)
    
    write (u, "(A,9(1x,F10.7))")  "r =", r
    write (u, "(A,9(1x,F10.7))")  "rb=", rb
    write (u, "(A,9(1x,F10.7))")  "x =", x
    write (u, "(A,9(1x,F10.7))")  "f =", f
    
    write (u, "(A)")
    write (u, "(A)")  "* Recover x and r from momenta"
    write (u, "(A)")
    
    q = interaction_get_momenta (sf_int%interaction_t, outgoing=.true.)
    call sf_int%final ()
    deallocate (sf_int)
    
    call data%allocate_sf_int (sf_int)
    call sf_int%init (data)
    call sf_int%set_beam_index ([1])
    call sf_int%setup_constants ()
    
    call sf_int%seed_kinematics ([k])
    call interaction_set_momenta (sf_int%interaction_t, q, outgoing=.true.)
    call sf_int%recover_x (x)
    call sf_int%inverse_kinematics (x, f, r, rb, map=.true., &
         set_momenta=.true.)    
    call interaction_pacify_momenta (sf_int%interaction_t, 1e-10_default)
    
    write (u, "(A,9(1x,F10.7))")  "r =", r
    write (u, "(A,9(1x,F10.7))")  "rb=", rb
    write (u, "(A,9(1x,F10.7))")  "x =", x
    write (u, "(A,9(1x,F10.7))")  "f =", f
    
    write (u, "(A)")
    write (u, "(A)")  "* Evaluate EWA structure function"
    write (u, "(A)")
    
    call sf_int%apply (scale = 1500._default)
    call sf_int%write (u, testflag = .true.)    
    
    write (u, "(A)")
    write (u, "(A)")  "* Cleanup"
    
    call sf_int%final ()
    call model%final ()
    
    write (u, "(A)")
    write (u, "(A)")  "* Test output end: sf_ewa_4"

  end subroutine sf_ewa_4

\eatline
\nwendcode{}\nwbegindocs{1194}\nwdocspar
\subsubsection{Structure function for multiple flavors}
Construct and display a structure function object based on the EWA
structure function.  The incoming state has multiple particles with
non-uniform quantum numbers.
\nwenddocs{}\nwbegincode{1195}\moddef{SF ewa: execute tests}\plusendmoddef
  call test (sf_ewa_5, "sf_ewa_5", &
       "structure function instance", &
       u, results)
\nwendcode{}\nwbegincode{1196}\moddef{SF ewa: tests}\plusendmoddef
  subroutine sf_ewa_5 (u)
    integer, intent(in) :: u
    type(model_data_t), target :: model
    type(flavor_t) :: flv
    type(pdg_array_t) :: pdg_in
    class(sf_data_t), allocatable, target :: data
    class(sf_int_t), allocatable :: sf_int
    type(vector4_t) :: k
    real(default) :: E
    real(default), dimension(:), allocatable :: r, rb, x
    real(default) :: f
    
    write (u, "(A)")  "* Test output: sf_ewa_5"
    write (u, "(A)")  "*   Purpose: initialize and fill &
         &test structure function object"
    write (u, "(A)")
    
    write (u, "(A)")  "* Initialize configuration data"
    write (u, "(A)")

    call model%init_sm_test ()
    call flv%init (2, model)
    pdg_in = [1, 2, -1, -2]

    call reset_interaction_counter ()
    
    allocate (ewa_data_t :: data)
    select type (data)
    type is (ewa_data_t)
       call data%init (model, pdg_in, 0.01_default, &
            500._default, 3000._default, .false., .false.)
       call data%set_id (24)
    end select
       
    write (u, "(A)")  "* Initialize structure-function object"
    write (u, "(A)")
    
    call data%allocate_sf_int (sf_int)
    call sf_int%init (data)
    call sf_int%set_beam_index ([1])
    call sf_int%setup_constants ()
    
    call sf_int%write (u)

    write (u, "(A)")
    write (u, "(A)")  "* Initialize incoming momentum with E=1500"
    write (u, "(A)")
    E = 1500
    k = vector4_moving (E, sqrt (E**2 - flv%get_mass ()**2), 3)
    call pacify (k, 1e-10_default)
    call vector4_write (k, u)
    call sf_int%seed_kinematics ([k])
    
    write (u, "(A)")
    write (u, "(A)")  "* Set kinematics for r=0.4, no EWA mapping, collinear"
    write (u, "(A)")
    
    allocate (r (data%get_n_par ()))
    allocate (rb(size (r)))
    allocate (x (size (r)))
    
    r = 0.4_default
    rb = 1 - r
    call sf_int%complete_kinematics (x, f, r, rb, map=.false.)
    
    write (u, "(A,9(1x,F10.7))")  "r =", r
    write (u, "(A,9(1x,F10.7))")  "rb=", rb
    write (u, "(A,9(1x,F10.7))")  "x =", x
    write (u, "(A,9(1x,F10.7))")  "f =", f
    
    write (u, "(A)")
    write (u, "(A)")  "* Evaluate EWA structure function"
    write (u, "(A)")
    
    call sf_int%apply (scale = 100._default)
    call sf_int%write (u)
        
    write (u, "(A)")
    write (u, "(A)")  "* Cleanup"
    
    call sf_int%final ()
    call model%final ()

    write (u, "(A)")
    write (u, "(A)")  "* Test output end: sf_ewa_5"
    
  end subroutine sf_ewa_5
  
\eatline
\nwendcode{}\nwbegindocs{1197}\nwdocspar
\clearpage
%------------------------------------------------------------------------
\section{Energy-scan spectrum}

This spectrum is actually a trick that allows us to plot the c.m.\ energy
dependence of a cross section without scanning the input energy.  We
start with the observation that a spectrum $f(x)$, applied to one of
the incoming beams only, results in a cross section
\begin{equation}
  \sigma = \int dx\,f(x)\,\hat\sigma(xs).
\end{equation}
We want to compute the distribution of $E=\sqrt{\hat s}=\sqrt{xs}$, i.e.,
\begin{equation}
  \frac{d\sigma}{dE} = \frac{2\sqrt{x}}{\sqrt{s}}\,\frac{d\sigma}{dx}
  = \frac{2\sqrt{x}}{\sqrt{s}}\,f(x)\,\hat\sigma(xs),
\end{equation}
so if we set
\begin{equation}
  f(x) = \frac{\sqrt{s}}{2\sqrt{x}},
\end{equation}
we get the distribution
\begin{equation}
  \frac{d\sigma}{dE} = \hat\sigma(\hat s=E^2).
\end{equation}
We implement this as a spectrum with a single parameter $x$.  The
parameters for the individual beams are computed as $x_i=\sqrt{x}$, so
they are equal and the kinematics is always symmetric.
\nwenddocs{}\nwbegincode{1198}\moddef{\code{}sf{\char95}escan.f90\edoc{}}\endmoddef
\LA{}File header\RA{}

module sf_escan

\LA{}Use kinds\RA{}
\LA{}Use strings\RA{}
  use io_units
  use format_defs, only: FMT_12
  use unit_tests
  use diagnostics
  use lorentz
  use physics_defs, only: ELECTRON
  use pdg_arrays
  use model_data
  use flavors
  use helicities
  use colors
  use quantum_numbers
  use state_matrices
  use polarizations
  use interactions
  use sf_aux
  use sf_base
  
\LA{}Standard module head\RA{}

\LA{}SF escan: public\RA{}

\LA{}SF escan: types\RA{}

contains

\LA{}SF escan: procedures\RA{}

\LA{}SF escan: tests\RA{}

end module sf_escan
\eatline
\nwendcode{}\nwbegindocs{1199}\nwdocspar
\subsection{Data type}
The {\tt{}norm} is unity if the total cross section should be normalized
to one, and $\sqrt{s}$ if it should be normalized to the total
energy.  In the latter case, the differential distribution
$d\sigma/d\sqrt{\hat s}$ coincides with the partonic cross section
$\hat\sigma$ as a function of $\sqrt{\hat s}$.
\nwenddocs{}\nwbegincode{1200}\moddef{SF escan: public}\endmoddef
  public :: escan_data_t
\nwendcode{}\nwbegincode{1201}\moddef{SF escan: types}\endmoddef
  type, extends(sf_data_t) :: escan_data_t
     private
     type(flavor_t), dimension(:,:), allocatable :: flv_in
     integer, dimension(2) :: n_flv = 0
     real(default) :: norm = 1
   contains
   \LA{}SF escan: escan data: TBP\RA{}  
  end type escan_data_t

\nwendcode{}\nwbegindocs{1202} %def escan_data_t
\nwenddocs{}\nwbegincode{1203}\moddef{SF escan: escan data: TBP}\endmoddef
  procedure :: init => escan_data_init
\nwendcode{}\nwbegincode{1204}\moddef{SF escan: procedures}\endmoddef
  subroutine escan_data_init (data, model, pdg_in, norm)
    class(escan_data_t), intent(out) :: data
    class(model_data_t), intent(in), target :: model
    type(pdg_array_t), dimension(2), intent(in) :: pdg_in
    real(default), intent(in), optional :: norm
    real(default), dimension(2) :: m2
    integer :: i, j
    data%n_flv = pdg_array_get_length (pdg_in)
    allocate (data%flv_in (maxval (data%n_flv), 2))
    do i = 1, 2
       do j = 1, data%n_flv(i)
          call data%flv_in(j, i)%init (pdg_array_get (pdg_in(i), j), model)
       end do
    end do
    m2 = data%flv_in(1,:)%get_mass ()
    do i = 1, 2
       if (any (data%flv_in(1:data%n_flv(i),i)%get_mass () /= m2(i))) then
          call msg_fatal ("Energy scan: incoming particle mass must be uniform")
       end if
    end do
    if (present (norm))  data%norm = norm
  end subroutine escan_data_init

\eatline
\nwendcode{}\nwbegindocs{1205}Output
\nwenddocs{}\nwbegincode{1206}\moddef{SF escan: escan data: TBP}\plusendmoddef
  procedure :: write => escan_data_write
\nwendcode{}\nwbegincode{1207}\moddef{SF escan: procedures}\plusendmoddef
  subroutine escan_data_write (data, unit, verbose) 
    class(escan_data_t), intent(in) :: data
    integer, intent(in), optional :: unit
    logical, intent(in), optional :: verbose
    integer :: u, i, j
    u = given_output_unit (unit);  if (u < 0)  return
    write (u, "(1x,A)") "Energy-scan data:"
    write (u, "(3x,A)", advance="no")  "prt_in = "
    do i = 1, 2
       if (i > 1)  write (u, "(',',1x)", advance="no")
       do j = 1, data%n_flv(i)
          if (j > 1)  write (u, "(':')", advance="no")
          write (u, "(A)", advance="no")  char (data%flv_in(j,i)%get_name ())
       end do
    end do
    write (u, *)
    write (u, "(3x,A," // FMT_12 // ")")  "norm   =", data%norm
  end subroutine escan_data_write

\eatline
\nwendcode{}\nwbegindocs{1208}Kinematics is completely collinear, hence there is only one
parameter for a pair spectrum. 
\nwenddocs{}\nwbegincode{1209}\moddef{SF escan: escan data: TBP}\plusendmoddef
  procedure :: get_n_par => escan_data_get_n_par
\nwendcode{}\nwbegincode{1210}\moddef{SF escan: procedures}\plusendmoddef
  function escan_data_get_n_par (data) result (n)
    class(escan_data_t), intent(in) :: data
    integer :: n
    n = 1
  end function escan_data_get_n_par
  
\eatline
\nwendcode{}\nwbegindocs{1211}Return the outgoing particles PDG codes. This is always the same as
the incoming particle, where we use two indices for the two beams.
\nwenddocs{}\nwbegincode{1212}\moddef{SF escan: escan data: TBP}\plusendmoddef
  procedure :: get_pdg_out => escan_data_get_pdg_out
\nwendcode{}\nwbegincode{1213}\moddef{SF escan: procedures}\plusendmoddef
  subroutine escan_data_get_pdg_out (data, pdg_out)
    class(escan_data_t), intent(in) :: data
    type(pdg_array_t), dimension(:), intent(inout) :: pdg_out
    integer :: i, n
    n = 2
    do i = 1, n
       pdg_out(i) = data%flv_in(1:data%n_flv(i),i)%get_pdg ()
    end do
  end subroutine escan_data_get_pdg_out
  
\eatline
\nwendcode{}\nwbegindocs{1214}Allocate the interaction record.
\nwenddocs{}\nwbegincode{1215}\moddef{SF escan: escan data: TBP}\plusendmoddef
  procedure :: allocate_sf_int => escan_data_allocate_sf_int
\nwendcode{}\nwbegincode{1216}\moddef{SF escan: procedures}\plusendmoddef
  subroutine escan_data_allocate_sf_int (data, sf_int)
    class(escan_data_t), intent(in) :: data
    class(sf_int_t), intent(inout), allocatable :: sf_int
    allocate (escan_t :: sf_int)
  end subroutine escan_data_allocate_sf_int
  
\eatline
\nwendcode{}\nwbegindocs{1217}\nwdocspar
\subsection{The Energy-scan object}
This is a spectrum, not a radiation.  We create an interaction with
two incoming and two outgoing particles, flavor, color, and helicity
being carried through. $x$ nevertheless is only one-dimensional, as we
are always using only one beam parameter.
\nwenddocs{}\nwbegincode{1218}\moddef{SF escan: types}\plusendmoddef
  type, extends (sf_int_t) :: escan_t
     type(escan_data_t), pointer :: data => null ()
   contains
   \LA{}SF escan: escan: TBP\RA{}
  end type escan_t 
  
\eatline
\nwendcode{}\nwbegindocs{1219}Type string: for the energy scan this is just a dummy function.
\nwenddocs{}\nwbegincode{1220}\moddef{SF escan: escan: TBP}\endmoddef
  procedure :: type_string => escan_type_string
\nwendcode{}\nwbegincode{1221}\moddef{SF escan: procedures}\plusendmoddef
  function escan_type_string (object) result (string)
    class(escan_t), intent(in) :: object
    type(string_t) :: string
    if (associated (object%data)) then
       string = "Escan: energy scan" 
    else
       string = "Escan: [undefined]"
    end if
  end function escan_type_string
  
\eatline
\nwendcode{}\nwbegindocs{1222}Output.  Call the interaction routine after displaying the configuration.
\nwenddocs{}\nwbegincode{1223}\moddef{SF escan: escan: TBP}\plusendmoddef
  procedure :: write => escan_write
\nwendcode{}\nwbegincode{1224}\moddef{SF escan: procedures}\plusendmoddef
  subroutine escan_write (object, unit, testflag)
    class(escan_t), intent(in) :: object
    integer, intent(in), optional :: unit
    logical, intent(in), optional :: testflag
    integer :: u
    u = given_output_unit (unit)
    if (associated (object%data)) then
       call object%data%write (u)
       call object%base_write (u, testflag)
    else
       write (u, "(1x,A)")  "Energy scan data: [undefined]"
    end if
  end subroutine escan_write
    
\eatline
\nwendcode{}\nwbegindocs{1225}\nwdocspar
\nwenddocs{}\nwbegincode{1226}\moddef{SF escan: escan: TBP}\plusendmoddef
  procedure :: init => escan_init
\nwendcode{}\nwbegincode{1227}\moddef{SF escan: procedures}\plusendmoddef
  subroutine escan_init (sf_int, data)
    class(escan_t), intent(out) :: sf_int
    class(sf_data_t), intent(in), target :: data
    type(quantum_numbers_mask_t), dimension(4) :: mask
    integer, dimension(4) :: hel_lock
    real(default), dimension(2) :: m2
    real(default), dimension(0) :: mr2
    type(quantum_numbers_t), dimension(4) :: qn_fc, qn_hel, qn
    type(polarization_t) :: pol1, pol2
    type(state_iterator_t) :: it_hel1, it_hel2
    integer :: j1, j2
    select type (data)
    type is (escan_data_t)
       hel_lock = [3, 4, 1, 2]
       m2 = data%flv_in(1,:)%get_mass ()
       call sf_int%base_init (mask, m2, mr2, m2, hel_lock = hel_lock)
       sf_int%data => data       
       do j1 = 1, data%n_flv(1)
          call qn_fc(1)%init ( &
               flv = data%flv_in(j1,1), &
               col = color_from_flavor (data%flv_in(j1,1)))
          call qn_fc(3)%init ( &
               flv = data%flv_in(j1,1), &
               col = color_from_flavor (data%flv_in(j1,1)))
          call polarization_init_generic (pol1, data%flv_in(j1,1))
          do j2 = 1, data%n_flv(2)
             call qn_fc(2)%init ( &
                  flv = data%flv_in(j2,2), &
                  col = color_from_flavor (data%flv_in(j2,2)))
             call qn_fc(4)%init ( &
                  flv = data%flv_in(j2,2), &
                  col = color_from_flavor (data%flv_in(j2,2)))
             call polarization_init_generic (pol2, data%flv_in(j2,2))
             call it_hel1%init (pol1%state)
             do while (it_hel1%is_valid ())
                qn_hel(1:1) = it_hel1%get_quantum_numbers ()
                qn_hel(3:3) = it_hel1%get_quantum_numbers ()
                call it_hel2%init (pol2%state)
                do while (it_hel2%is_valid ())
                   qn_hel(2:2) = it_hel2%get_quantum_numbers ()
                   qn_hel(4:4) = it_hel2%get_quantum_numbers ()
                   qn = qn_hel .merge. qn_fc
                   call interaction_add_state (sf_int%interaction_t, qn)
                   call it_hel2%advance ()
                end do
                call it_hel1%advance ()
             end do
             call polarization_final (pol2)
          end do
          call polarization_final (pol1)
       end do
       call sf_int%set_incoming ([1,2])
       call sf_int%set_outgoing ([3,4])
       call interaction_freeze (sf_int%interaction_t)       
       sf_int%status = SF_INITIAL
    end select
  end subroutine escan_init
    
\eatline
\nwendcode{}\nwbegindocs{1228}\nwdocspar
\subsection{Kinematics}
Set kinematics.  We have a single parameter, but reduce both beams.
The {\tt{}map} flag is ignored.
\nwenddocs{}\nwbegincode{1229}\moddef{SF escan: escan: TBP}\plusendmoddef
  procedure :: complete_kinematics => escan_complete_kinematics
\nwendcode{}\nwbegincode{1230}\moddef{SF escan: procedures}\plusendmoddef
  subroutine escan_complete_kinematics (sf_int, x, f, r, rb, map)
    class(escan_t), intent(inout) :: sf_int
    real(default), dimension(:), intent(out) :: x
    real(default), intent(out) :: f
    real(default) :: sqrt_x
    real(default), dimension(:), intent(in) :: r
    real(default), dimension(:), intent(in) :: rb
    logical, intent(in) :: map
    x = r
    sqrt_x = sqrt (x(1))
    if (sqrt_x > 0) then
       f = 1 / (2 * sqrt_x)
    else
       f = 0
       sf_int%status = SF_FAILED_KINEMATICS
       return
    end if
    call sf_int%reduce_momenta ([sqrt_x, sqrt_x])
  end subroutine escan_complete_kinematics

\eatline
\nwendcode{}\nwbegindocs{1231}Recover $x$.  The base procedure should return two momentum
fractions for the two beams, while we have only one parameter.  This
is the product of the extracted momentum fractions.
\nwenddocs{}\nwbegincode{1232}\moddef{SF escan: escan: TBP}\plusendmoddef
  procedure :: recover_x => escan_recover_x
\nwendcode{}\nwbegincode{1233}\moddef{SF escan: procedures}\plusendmoddef
  subroutine escan_recover_x (sf_int, x, x_free)
    class(escan_t), intent(inout) :: sf_int
    real(default), dimension(:), intent(out) :: x
    real(default), intent(inout), optional :: x_free
    real(default), dimension(2) :: xi
    call sf_int%base_recover_x (xi, x_free)
    x = product (xi)
  end subroutine escan_recover_x
  
\eatline
\nwendcode{}\nwbegindocs{1234}Compute inverse kinematics.
\nwenddocs{}\nwbegincode{1235}\moddef{SF escan: escan: TBP}\plusendmoddef
  procedure :: inverse_kinematics => escan_inverse_kinematics
\nwendcode{}\nwbegincode{1236}\moddef{SF escan: procedures}\plusendmoddef
  subroutine escan_inverse_kinematics (sf_int, x, f, r, rb, map, set_momenta)
    class(escan_t), intent(inout) :: sf_int
    real(default), dimension(:), intent(in) :: x
    real(default), intent(out) :: f
    real(default), dimension(:), intent(out) :: r
    real(default), dimension(:), intent(out) :: rb
    logical, intent(in) :: map
    logical, intent(in), optional :: set_momenta
    real(default) :: sqrt_x
    logical :: set_mom
    set_mom = .false.;  if (present (set_momenta))  set_mom = set_momenta
    sqrt_x = sqrt (x(1))
    if (sqrt_x > 0) then
       f = 1 / (2 * sqrt_x)
    else
       f = 0
       sf_int%status = SF_FAILED_KINEMATICS
       return
    end if
    r = x
    rb = 1 - r
    if (set_mom) then
       call sf_int%reduce_momenta ([sqrt_x, sqrt_x])
    end if
  end subroutine escan_inverse_kinematics

\eatline
\nwendcode{}\nwbegindocs{1237}\nwdocspar
\subsection{Energy scan application}
Here, we insert the predefined norm.
\nwenddocs{}\nwbegincode{1238}\moddef{SF escan: escan: TBP}\plusendmoddef
  procedure :: apply => escan_apply
\nwendcode{}\nwbegincode{1239}\moddef{SF escan: procedures}\plusendmoddef
  subroutine escan_apply (sf_int, scale)
    class(escan_t), intent(inout) :: sf_int
    real(default), intent(in) :: scale
    real(default) :: f
    associate (data => sf_int%data)
      f = data%norm
    end associate
    call interaction_set_matrix_element (sf_int%interaction_t, &
         cmplx (f, kind=default))    
    sf_int%status = SF_EVALUATED
  end subroutine escan_apply

\eatline
\nwendcode{}\nwbegindocs{1240}\nwdocspar
\subsection{Unit tests}
\nwenddocs{}\nwbegincode{1241}\moddef{SF escan: public}\plusendmoddef
  public :: sf_escan_test
\nwendcode{}\nwbegincode{1242}\moddef{SF escan: tests}\endmoddef
  subroutine sf_escan_test (u, results)
    integer, intent(in) :: u
    type(test_results_t), intent(inout) :: results
  \LA{}SF escan: execute tests\RA{}
  end subroutine sf_escan_test
  
\eatline
\nwendcode{}\nwbegindocs{1243}\nwdocspar
\subsubsection{Test structure function data}
Construct and display a test structure function data object.
\nwenddocs{}\nwbegincode{1244}\moddef{SF escan: execute tests}\endmoddef
  call test (sf_escan_1, "sf_escan_1", &
       "structure function configuration", &
       u, results)
\nwendcode{}\nwbegincode{1245}\moddef{SF escan: tests}\plusendmoddef
  subroutine sf_escan_1 (u)
    integer, intent(in) :: u
    type(model_data_t), target :: model
    type(pdg_array_t), dimension(2) :: pdg_in
    type(pdg_array_t), dimension(2) :: pdg_out
    integer, dimension(:), allocatable :: pdg1, pdg2
    class(sf_data_t), allocatable :: data
    
    write (u, "(A)")  "* Test output: sf_escan_1"
    write (u, "(A)")  "*   Purpose: initialize and display &
         &energy-scan structure function data"
    write (u, "(A)")
    
    call model%init_qed_test ()
    pdg_in(1) = ELECTRON
    pdg_in(2) = -ELECTRON

    allocate (escan_data_t :: data)
    select type (data)
    type is (escan_data_t)
       call data%init (model, pdg_in, norm = 2._default)
    end select

    call data%write (u)

    write (u, "(A)")

    write (u, "(1x,A)")  "Outgoing particle codes:"
    call data%get_pdg_out (pdg_out)
    pdg1 = pdg_out(1)
    pdg2 = pdg_out(2)
    write (u, "(2x,99(1x,I0))")  pdg1, pdg2

    call model%final ()
    
    write (u, "(A)")
    write (u, "(A)")  "* Test output end: sf_escan_1"

  end subroutine sf_escan_1

\eatline
\nwendcode{}\nwbegindocs{1246}g@
\subsubsection{Probe the structure-function object}
Active the beam event reader, generate an event.
\nwenddocs{}\nwbegincode{1247}\moddef{SF escan: execute tests}\plusendmoddef
  call test (sf_escan_2, "sf_escan_2", &
       "generate event", &
       u, results)
\nwendcode{}\nwbegincode{1248}\moddef{SF escan: tests}\plusendmoddef
  subroutine sf_escan_2 (u)
    integer, intent(in) :: u
    type(model_data_t), target :: model
    type(flavor_t), dimension(2) :: flv
    type(pdg_array_t), dimension(2) :: pdg_in
    class(sf_data_t), allocatable, target :: data
    class(sf_int_t), allocatable :: sf_int
    type(vector4_t) :: k1, k2
    real(default) :: E
    real(default), dimension(:), allocatable :: r, rb, x
    real(default) :: x_free, f
    
    write (u, "(A)")  "* Test output: sf_escan_2"
    write (u, "(A)")  "*   Purpose: initialize and display &
         &beam-events structure function data"
    write (u, "(A)")
    
    call model%init_qed_test ()
    call flv(1)%init (ELECTRON, model)
    call flv(2)%init (-ELECTRON, model)
    pdg_in(1) = ELECTRON
    pdg_in(2) = -ELECTRON

    call reset_interaction_counter ()
    
    allocate (escan_data_t :: data)
    select type (data)
    type is (escan_data_t)
       call data%init (model, pdg_in)
    end select
       
    write (u, "(A)")  "* Initialize structure-function object"
    write (u, "(A)")
    
    call data%allocate_sf_int (sf_int)
    call sf_int%init (data)
    call sf_int%set_beam_index ([1,2])

    write (u, "(A)")  "* Initialize incoming momentum with E=500"
    write (u, "(A)")
    E = 250
    k1 = vector4_moving (E, sqrt (E**2 - flv(1)%get_mass ()**2), 3)
    k2 = vector4_moving (E,-sqrt (E**2 - flv(2)%get_mass ()**2), 3)
    call vector4_write (k1, u)
    call vector4_write (k2, u)
    call sf_int%seed_kinematics ([k1, k2])

    write (u, "(A)")
    write (u, "(A)")  "* Set dummy parameters and generate x"
    write (u, "(A)")

    allocate (r (data%get_n_par ()))
    allocate (rb(size (r)))
    allocate (x (size (r)))

    r  = 0.8
    rb = 1 - r
    x_free = 1

    call sf_int%complete_kinematics (x, f, r, rb, map=.false.)

    write (u, "(A,9(1x,F10.7))")  "r =", r
    write (u, "(A,9(1x,F10.7))")  "rb=", rb
    write (u, "(A,9(1x,F10.7))")  "x =", x
    write (u, "(A,9(1x,F10.7))")  "f =", f
    write (u, "(A,9(1x,F10.7))")  "xf=", x_free

    write (u, "(A)")
    write (u, "(A)")  "* Inverse kinematics"
    write (u, "(A)")

    call sf_int%recover_x (x, x_free)
    call sf_int%inverse_kinematics (x, f, r, rb, map=.false.)

    write (u, "(A,9(1x,F10.7))")  "r =", r
    write (u, "(A,9(1x,F10.7))")  "rb=", rb
    write (u, "(A,9(1x,F10.7))")  "x =", x
    write (u, "(A,9(1x,F10.7))")  "f =", f
    write (u, "(A,9(1x,F10.7))")  "xf=", x_free

    write (u, "(A)")
    write (u, "(A)")  "* Evaluate"
    write (u, "(A)")

    call sf_int%apply (scale = 0._default)
    call sf_int%write (u)

    write (u, "(A)")
    write (u, "(A)")  "* Cleanup"

    call sf_int%final ()
    call model%final ()
    
    write (u, "(A)")
    write (u, "(A)")  "* Test output end: sf_escan_2"

  end subroutine sf_escan_2

\eatline
\nwendcode{}\nwbegindocs{1249}%%%%%%%%%%%%%%%%%%%%%%%%%%%%%%%%%%%%%%%%%%%%%%%%%%%%%%%%%%%%%%%%%%%%%%%%
\section{Using beam event data}

Instead of an analytic beam description, beam data may be provided in
form of an event file.  In its most simple form, the event file
contains pairs of $x$ values, relative to nominal beam energies.  More
advanced formats may include polarization, etc.  The current
implementation carries beam polarization through, if specified.

The code is very similar to the energy scan described above.

However, we must include a file-handle manager for the beam-event
files.  Two different processes may access a given beam-event file at
the same time (i.e., serially but alternating).  Accessing an open
file from two different units is non-standard and not supported by all
compilers.  Therefore, we keep a global registry of open files,
associated units, and reference counts.  The {\tt{}beam{\char95}events{\char95}t} objects
act as proxies to this registry.
\nwenddocs{}\nwbegincode{1250}\moddef{\code{}sf{\char95}beam{\char95}events.f90\edoc{}}\endmoddef
\LA{}File header\RA{}

module sf_beam_events

\LA{}Use kinds\RA{}
\LA{}Use strings\RA{}
  use io_units
  use file_registries
  use unit_tests
  use diagnostics
  use physics_defs, only: ELECTRON
  use lorentz
  use pdg_arrays
  use model_data
  use flavors
  use helicities
  use colors
  use quantum_numbers
  use state_matrices
  use polarizations
  use interactions
  use sf_aux
  use sf_base

\LA{}Standard module head\RA{}

\LA{}SF beam events: public\RA{}

\LA{}SF beam events: types\RA{}

\LA{}SF beam events: variables\RA{}

contains

\LA{}SF beam events: procedures\RA{}

\LA{}SF beam events: tests\RA{}

end module sf_beam_events
\eatline
\nwendcode{}\nwbegindocs{1251}\nwdocspar
\subsection{The beam-data file registry}
We manage data files via the {\tt{}file{\char95}registries} module.  To this end,
we keep the registry as a private module variable here.
\nwenddocs{}\nwbegincode{1252}\moddef{SF beam events: variables}\endmoddef
  type(file_registry_t), save :: beam_file_registry

\eatline
\nwendcode{}\nwbegindocs{1253}\nwdocspar
\subsection{Data type}
\nwenddocs{}\nwbegincode{1254}\moddef{SF beam events: public}\endmoddef
  public :: beam_events_data_t
\nwendcode{}\nwbegincode{1255}\moddef{SF beam events: types}\endmoddef
  type, extends(sf_data_t) :: beam_events_data_t
     private
     type(flavor_t), dimension(2) :: flv_in
     type(string_t) :: dir
     type(string_t) :: file
     type(string_t) :: fqn
     integer :: unit = 0
     logical :: warn_eof = .true.
   contains
   \LA{}SF beam events: beam events data: TBP\RA{}
  end type beam_events_data_t

\nwendcode{}\nwbegindocs{1256} %def beam_events_data_t
\nwenddocs{}\nwbegincode{1257}\moddef{SF beam events: beam events data: TBP}\endmoddef
  procedure :: init => beam_events_data_init
\nwendcode{}\nwbegincode{1258}\moddef{SF beam events: procedures}\endmoddef
  subroutine beam_events_data_init (data, model, pdg_in, dir, file, warn_eof)
    class(beam_events_data_t), intent(out) :: data
    class(model_data_t), intent(in), target :: model
    type(pdg_array_t), dimension(2), intent(in) :: pdg_in
    type(string_t), intent(in) :: dir
    type(string_t), intent(in) :: file
    logical, intent(in), optional :: warn_eof
    if (any (pdg_array_get_length (pdg_in) /= 1)) then
       call msg_fatal ("Beam events: incoming beam particles must be unique")
    end if
    call data%flv_in(1)%init (pdg_array_get (pdg_in(1), 1), model)
    call data%flv_in(2)%init (pdg_array_get (pdg_in(2), 1), model)
    data%dir = dir
    data%file = file
    if (present (warn_eof))  data%warn_eof = warn_eof
  end subroutine beam_events_data_init

\eatline
\nwendcode{}\nwbegindocs{1259}Return true since this spectrum is always in generator mode.
\nwenddocs{}\nwbegincode{1260}\moddef{SF beam events: beam events data: TBP}\plusendmoddef
  procedure :: is_generator => beam_events_data_is_generator
\nwendcode{}\nwbegincode{1261}\moddef{SF beam events: procedures}\plusendmoddef
  function beam_events_data_is_generator (data) result (flag)
    class(beam_events_data_t), intent(in) :: data
    logical :: flag
    flag = .true.
  end function beam_events_data_is_generator
  
\eatline
\nwendcode{}\nwbegindocs{1262}The number of parameters is two.  They are free parameters.
\nwenddocs{}\nwbegincode{1263}\moddef{SF beam events: beam events data: TBP}\plusendmoddef
  procedure :: get_n_par => beam_events_data_get_n_par
\nwendcode{}\nwbegincode{1264}\moddef{SF beam events: procedures}\plusendmoddef
  function beam_events_data_get_n_par (data) result (n)
    class(beam_events_data_t), intent(in) :: data
    integer :: n
    n = 2
  end function beam_events_data_get_n_par
  
\eatline
\nwendcode{}\nwbegincode{1265}\moddef{SF beam events: beam events data: TBP}\plusendmoddef
  procedure :: get_pdg_out => beam_events_data_get_pdg_out
\nwendcode{}\nwbegincode{1266}\moddef{SF beam events: procedures}\plusendmoddef
  subroutine beam_events_data_get_pdg_out (data, pdg_out)
    class(beam_events_data_t), intent(in) :: data
    type(pdg_array_t), dimension(:), intent(inout) :: pdg_out
    integer :: i, n
    n = 2
    do i = 1, n
       pdg_out(i) = data%flv_in(i)%get_pdg ()
    end do
  end subroutine beam_events_data_get_pdg_out
  
\eatline
\nwendcode{}\nwbegindocs{1267}Allocate the interaction record.
\nwenddocs{}\nwbegincode{1268}\moddef{SF beam events: beam events data: TBP}\plusendmoddef
  procedure :: allocate_sf_int => beam_events_data_allocate_sf_int
\nwendcode{}\nwbegincode{1269}\moddef{SF beam events: procedures}\plusendmoddef
  subroutine beam_events_data_allocate_sf_int (data, sf_int)
    class(beam_events_data_t), intent(in) :: data
    class(sf_int_t), intent(inout), allocatable :: sf_int
    allocate (beam_events_t :: sf_int)
  end subroutine beam_events_data_allocate_sf_int
  
\eatline
\nwendcode{}\nwbegindocs{1270}Output
\nwenddocs{}\nwbegincode{1271}\moddef{SF beam events: beam events data: TBP}\plusendmoddef
  procedure :: write => beam_events_data_write
\nwendcode{}\nwbegincode{1272}\moddef{SF beam events: procedures}\plusendmoddef
  subroutine beam_events_data_write (data, unit, verbose) 
    class(beam_events_data_t), intent(in) :: data
    integer, intent(in), optional :: unit
    logical, intent(in), optional :: verbose        
    integer :: u
    u = given_output_unit (unit);  if (u < 0)  return
    write (u, "(1x,A)") "Beam-event file data:"
    write (u, "(3x,A,A,A,A)") "prt_in = ", &
         char (data%flv_in(1)%get_name ()), &
         ", ", char (data%flv_in(2)%get_name ())    
    write (u, "(3x,A,A,A)") "file   = '", char (data%file), "'"
    write (u, "(3x,A,I0)")  "unit   = ", data%unit
    write (u, "(3x,A,L1)")  "warn   = ", data%warn_eof
  end subroutine beam_events_data_write

\eatline
\nwendcode{}\nwbegindocs{1273}The data file needs to be opened and closed explicitly.  The
open/close message is communicated to the file handle registry, which
does the actual work.

We determine first whether to look in the local directory or in the
given system directory.
\nwenddocs{}\nwbegincode{1274}\moddef{SF beam events: beam events data: TBP}\plusendmoddef
  procedure :: open => beam_events_data_open
  procedure :: close => beam_events_data_close
\nwendcode{}\nwbegincode{1275}\moddef{SF beam events: procedures}\plusendmoddef
  subroutine beam_events_data_open (data)
    class(beam_events_data_t), intent(inout) :: data
    logical :: exist
    if (data%unit == 0) then
       data%fqn = data%file
       if (data%fqn == "") &
            call msg_fatal ("Beam events: $beam_events_file is not set")
       inquire (file = char (data%fqn), exist = exist)
       if (.not. exist) then
          data%fqn = data%dir // "/" // data%file
          inquire (file = char (data%fqn), exist = exist)
          if (.not. exist) then
             data%fqn = ""
             call msg_fatal ("Beam events: file '" &
                  // char (data%file) // "' not found")
             return
          end if
       end if
       call msg_message ("Beam events: reading from file '" &
            // char (data%file) // "'")
       call beam_file_registry%open (data%fqn, data%unit)
    else
       call msg_bug ("Beam events: file '" &
         // char (data%file) // "' is already open")
    end if
  end subroutine beam_events_data_open

  subroutine beam_events_data_close (data)
    class(beam_events_data_t), intent(inout) :: data
    if (data%unit /= 0) then
       call beam_file_registry%close (data%fqn)
       call msg_message ("Beam events: closed file '" &
         // char (data%file) // "'")
       data%unit = 0
    end if
  end subroutine beam_events_data_close

\eatline
\nwendcode{}\nwbegindocs{1276}\nwdocspar
\subsection{The beam events object}
Flavor and polarization carried through, no radiated particles.
\nwenddocs{}\nwbegincode{1277}\moddef{SF beam events: types}\plusendmoddef
  type, extends (sf_int_t) :: beam_events_t
     type(beam_events_data_t), pointer :: data => null ()
     integer :: count = 0
   contains
   \LA{}SF beam events: beam events: TBP\RA{}
  end type beam_events_t 
  
\eatline
\nwendcode{}\nwbegindocs{1278}Type string: show beam events file.
\nwenddocs{}\nwbegincode{1279}\moddef{SF beam events: beam events: TBP}\endmoddef
  procedure :: type_string => beam_events_type_string
\nwendcode{}\nwbegincode{1280}\moddef{SF beam events: procedures}\plusendmoddef
  function beam_events_type_string (object) result (string)
    class(beam_events_t), intent(in) :: object
    type(string_t) :: string
    if (associated (object%data)) then
       string = "Beam events: " // object%data%file 
    else
       string = "Beam events: [undefined]"
    end if
  end function beam_events_type_string
  
\eatline
\nwendcode{}\nwbegindocs{1281}Output.  Call the interaction routine after displaying the configuration.
\nwenddocs{}\nwbegincode{1282}\moddef{SF beam events: beam events: TBP}\plusendmoddef
  procedure :: write => beam_events_write
\nwendcode{}\nwbegincode{1283}\moddef{SF beam events: procedures}\plusendmoddef
  subroutine beam_events_write (object, unit, testflag)
    class(beam_events_t), intent(in) :: object
    integer, intent(in), optional :: unit
    logical, intent(in), optional :: testflag
    integer :: u
    u = given_output_unit (unit)
    if (associated (object%data)) then
       call object%data%write (u)
       call object%base_write (u, testflag)
    else
       write (u, "(1x,A)")  "Beam events data: [undefined]"
    end if
  end subroutine beam_events_write
    
\eatline
\nwendcode{}\nwbegindocs{1284}\nwdocspar
\nwenddocs{}\nwbegincode{1285}\moddef{SF beam events: beam events: TBP}\plusendmoddef
  procedure :: init => beam_events_init
\nwendcode{}\nwbegincode{1286}\moddef{SF beam events: procedures}\plusendmoddef
  subroutine beam_events_init (sf_int, data)
    class(beam_events_t), intent(out) :: sf_int
    class(sf_data_t), intent(in), target :: data
    real(default), dimension(2) :: m2
    real(default), dimension(0) :: mr2
    type(quantum_numbers_mask_t), dimension(4) :: mask
    integer, dimension(4) :: hel_lock
    type(quantum_numbers_t), dimension(4) :: qn_fc, qn_hel, qn
    type(polarization_t) :: pol1, pol2
    type(state_iterator_t) :: it_hel1, it_hel2
    integer :: i
    select type (data)
    type is (beam_events_data_t)
       m2 = data%flv_in%get_mass () ** 2
       hel_lock = [3, 4, 1, 2]
       mask = quantum_numbers_mask (.false., .false., .false.)
       call sf_int%base_init (mask, m2, mr2, m2, hel_lock = hel_lock)
       sf_int%data => data
       do i = 1, 2
          call qn_fc(i)%init ( &
               flv = data%flv_in(i), &
               col = color_from_flavor (data%flv_in(i)))
          call qn_fc(i+2)%init ( &
               flv = data%flv_in(i), &
               col = color_from_flavor (data%flv_in(i)))
       end do
       call polarization_init_generic (pol1, data%flv_in(1))
       call it_hel1%init (pol1%state)
       do while (it_hel1%is_valid ())
          qn_hel(1:1) = it_hel1%get_quantum_numbers ()
          qn_hel(3:3) = it_hel1%get_quantum_numbers ()
          call polarization_init_generic (pol2, data%flv_in(2))
          call it_hel2%init (pol2%state)
          do while (it_hel2%is_valid ())
             qn_hel(2:2) = it_hel2%get_quantum_numbers ()
             qn_hel(4:4) = it_hel2%get_quantum_numbers ()
             qn = qn_hel .merge. qn_fc
             call interaction_add_state (sf_int%interaction_t, qn)
             call it_hel2%advance ()
          end do
          call polarization_final (pol2)
          call it_hel1%advance ()
       end do
       call polarization_final (pol2)
       call interaction_freeze (sf_int%interaction_t)
       call sf_int%set_incoming ([1,2])
       call sf_int%set_outgoing ([3,4])
       call sf_int%data%open ()
       sf_int%status = SF_INITIAL
    end select
  end subroutine beam_events_init
    
\eatline
\nwendcode{}\nwbegindocs{1287}This spectrum type needs a finalizer, which closes the data file.
\nwenddocs{}\nwbegincode{1288}\moddef{SF beam events: beam events: TBP}\plusendmoddef
  procedure :: final => sf_beam_events_final
\nwendcode{}\nwbegincode{1289}\moddef{SF beam events: procedures}\plusendmoddef
  subroutine sf_beam_events_final (object)
    class(beam_events_t), intent(inout) :: object
    call object%data%close ()
    call interaction_final (object%interaction_t)
  end subroutine sf_beam_events_final

\eatline
\nwendcode{}\nwbegindocs{1290}\nwdocspar
\subsection{Kinematics}
Refer to the {\tt{}data} component.
\nwenddocs{}\nwbegincode{1291}\moddef{SF beam events: beam events: TBP}\plusendmoddef
  procedure :: is_generator => beam_events_is_generator
\nwendcode{}\nwbegincode{1292}\moddef{SF beam events: procedures}\plusendmoddef
  function beam_events_is_generator (sf_int) result (flag)
    class(beam_events_t), intent(in) :: sf_int
    logical :: flag
    flag = sf_int%data%is_generator ()
  end function beam_events_is_generator
  
\eatline
\nwendcode{}\nwbegindocs{1293}Generate free parameters.  We read them from file.
\nwenddocs{}\nwbegincode{1294}\moddef{SF beam events: beam events: TBP}\plusendmoddef
  procedure :: generate_free => beam_events_generate_free
\nwendcode{}\nwbegincode{1295}\moddef{SF beam events: procedures}\plusendmoddef
  recursive subroutine beam_events_generate_free (sf_int, r, rb,  x_free)
    class(beam_events_t), intent(inout) :: sf_int
    real(default), dimension(:), intent(out) :: r, rb
    real(default), intent(inout) :: x_free
    integer :: iostat
    associate (data => sf_int%data)
      if (data%unit /= 0) then
         read (data%unit, fmt=*, iostat=iostat)  r
         if (iostat > 0) then
            write (msg_buffer, "(A,I0,A)") &
                 "Beam events: I/O error after reading ", sf_int%count, &
                 " events"
            call msg_fatal ()
         else if (iostat < 0) then
            if (sf_int%count == 0) then
               call msg_fatal ("Beam events: file is empty")
            else if (sf_int%data%warn_eof) then
               write (msg_buffer, "(A,I0,A)") &
                    "Beam events: End of file after reading ", sf_int%count, &
                    " events, rewinding"
               call msg_warning ()
            end if
            rewind (data%unit)
            sf_int%count = 0
            call sf_int%generate_free (r, rb, x_free)
         else
            sf_int%count = sf_int%count + 1
            rb = 1 - r
            x_free = x_free * product (r)
         end if
      else
         call msg_bug ("Beam events: file is not open for reading")
      end if
    end associate
  end subroutine beam_events_generate_free
    
\eatline
\nwendcode{}\nwbegindocs{1296}Set kinematics.  Trivial transfer since this is a pure generator.

If {\tt{}map} is set, we are asked to provide an efficient mapping.
For the test case, we set $x=r^2$ and consequently $f(r)=2r$.
\nwenddocs{}\nwbegincode{1297}\moddef{SF beam events: beam events: TBP}\plusendmoddef
  procedure :: complete_kinematics => beam_events_complete_kinematics
\nwendcode{}\nwbegincode{1298}\moddef{SF beam events: procedures}\plusendmoddef
  subroutine beam_events_complete_kinematics (sf_int, x, f, r, rb, map)
    class(beam_events_t), intent(inout) :: sf_int
    real(default), dimension(:), intent(out) :: x
    real(default), intent(out) :: f
    real(default), dimension(:), intent(in) :: r
    real(default), dimension(:), intent(in) :: rb
    logical, intent(in) :: map
    if (map) then
       call msg_fatal ("Beam events: map flag not supported")
    else
       x = r
       f = 1
    end if
    call sf_int%reduce_momenta (x)
  end subroutine beam_events_complete_kinematics

\eatline
\nwendcode{}\nwbegindocs{1299}Compute inverse kinematics.  Trivial in this case.
\nwenddocs{}\nwbegincode{1300}\moddef{SF beam events: beam events: TBP}\plusendmoddef
  procedure :: inverse_kinematics => beam_events_inverse_kinematics
\nwendcode{}\nwbegincode{1301}\moddef{SF beam events: procedures}\plusendmoddef
  subroutine beam_events_inverse_kinematics &
       (sf_int, x, f, r, rb, map, set_momenta)
    class(beam_events_t), intent(inout) :: sf_int
    real(default), dimension(:), intent(in) :: x
    real(default), intent(out) :: f
    real(default), dimension(:), intent(out) :: r
    real(default), dimension(:), intent(out) :: rb
    logical, intent(in) :: map
    logical, intent(in), optional :: set_momenta
    logical :: set_mom
    set_mom = .false.;  if (present (set_momenta))  set_mom = set_momenta
    if (map) then
       call msg_fatal ("Beam events: map flag not supported")
    else
       r = x
       f = 1
    end if
    rb = 1 - r
    if (set_mom) then
       call sf_int%reduce_momenta (x)
    end if
  end subroutine beam_events_inverse_kinematics

\eatline
\nwendcode{}\nwbegindocs{1302}\nwdocspar
\subsection{Beam events application}
Trivial, just set the unit weight.
\nwenddocs{}\nwbegincode{1303}\moddef{SF beam events: beam events: TBP}\plusendmoddef
  procedure :: apply => beam_events_apply
\nwendcode{}\nwbegincode{1304}\moddef{SF beam events: procedures}\plusendmoddef
  subroutine beam_events_apply (sf_int, scale)
    class(beam_events_t), intent(inout) :: sf_int
    real(default), intent(in) :: scale
    real(default) :: f
    f = 1
    call interaction_set_matrix_element (sf_int%interaction_t, &
         cmplx (f, kind=default))        
    sf_int%status = SF_EVALUATED
  end subroutine beam_events_apply

\eatline
\nwendcode{}\nwbegindocs{1305}\nwdocspar
\subsection{Unit tests}
\nwenddocs{}\nwbegincode{1306}\moddef{SF beam events: public}\plusendmoddef
  public :: sf_beam_events_test
\nwendcode{}\nwbegincode{1307}\moddef{SF beam events: tests}\endmoddef
  subroutine sf_beam_events_test (u, results)
    integer, intent(in) :: u
    type(test_results_t), intent(inout) :: results
  \LA{}SF beam events: execute tests\RA{}
  end subroutine sf_beam_events_test
  
\eatline
\nwendcode{}\nwbegindocs{1308}\nwdocspar
\subsubsection{Test structure function data}
Construct and display a test structure function data object.
\nwenddocs{}\nwbegincode{1309}\moddef{SF beam events: execute tests}\endmoddef
  call test (sf_beam_events_1, "sf_beam_events_1", &
       "structure function configuration", &
       u, results)
\nwendcode{}\nwbegincode{1310}\moddef{SF beam events: tests}\plusendmoddef
  subroutine sf_beam_events_1 (u)
    integer, intent(in) :: u
    type(model_data_t), target :: model
    type(pdg_array_t), dimension(2) :: pdg_in
    type(pdg_array_t), dimension(2) :: pdg_out
    integer, dimension(:), allocatable :: pdg1, pdg2
    class(sf_data_t), allocatable :: data
    
    write (u, "(A)")  "* Test output: sf_beam_events_1"
    write (u, "(A)")  "*   Purpose: initialize and display &
         &beam-events structure function data"
    write (u, "(A)")
    
    call model%init_qed_test ()
    pdg_in(1) = ELECTRON
    pdg_in(2) = -ELECTRON

    allocate (beam_events_data_t :: data)
    select type (data)
    type is (beam_events_data_t)
       call data%init (model, pdg_in, var_str (""), var_str ("beam_events.dat"))
    end select

    call data%write (u)

    write (u, "(A)")

    write (u, "(1x,A)")  "Outgoing particle codes:"
    call data%get_pdg_out (pdg_out)
    pdg1 = pdg_out(1)
    pdg2 = pdg_out(2)
    write (u, "(2x,99(1x,I0))")  pdg1, pdg2

    call model%final ()
    
    write (u, "(A)")
    write (u, "(A)")  "* Test output end: sf_beam_events_1"

  end subroutine sf_beam_events_1

\eatline
\nwendcode{}\nwbegindocs{1311}\nwdocspar
\subsubsection{Probe the structure-function object}
Active the beam event reader, generate an event.
\nwenddocs{}\nwbegincode{1312}\moddef{SF beam events: execute tests}\plusendmoddef
  call test (sf_beam_events_2, "sf_beam_events_2", &
       "generate event", &
       u, results)
\nwendcode{}\nwbegincode{1313}\moddef{SF beam events: tests}\plusendmoddef
  subroutine sf_beam_events_2 (u)
    integer, intent(in) :: u
    type(model_data_t), target :: model
    type(flavor_t), dimension(2) :: flv
    type(pdg_array_t), dimension(2) :: pdg_in
    class(sf_data_t), allocatable, target :: data
    class(sf_int_t), allocatable :: sf_int
    type(vector4_t) :: k1, k2
    real(default) :: E
    real(default), dimension(:), allocatable :: r, rb, x
    real(default) :: x_free, f
    integer :: i
    
    write (u, "(A)")  "* Test output: sf_beam_events_2"
    write (u, "(A)")  "*   Purpose: initialize and display &
         &beam-events structure function data"
    write (u, "(A)")
    
    call model%init_qed_test ()
    call flv(1)%init (ELECTRON, model)
    call flv(2)%init (-ELECTRON, model)
    pdg_in(1) = ELECTRON
    pdg_in(2) = -ELECTRON

    call reset_interaction_counter ()
    
    allocate (beam_events_data_t :: data)
    select type (data)
    type is (beam_events_data_t)
       call data%init (model, pdg_in, &
            var_str (""), var_str ("test_beam_events.dat"))
    end select
       
    write (u, "(A)")  "* Initialize structure-function object"
    write (u, "(A)")
    
    call data%allocate_sf_int (sf_int)
    call sf_int%init (data)
    call sf_int%set_beam_index ([1,2])

    write (u, "(A)")  "* Initialize incoming momentum with E=500"
    write (u, "(A)")
    E = 250
    k1 = vector4_moving (E, sqrt (E**2 - flv(1)%get_mass ()**2), 3)
    k2 = vector4_moving (E,-sqrt (E**2 - flv(2)%get_mass ()**2), 3)
    call vector4_write (k1, u)
    call vector4_write (k2, u)
    call sf_int%seed_kinematics ([k1, k2])

    write (u, "(A)")
    write (u, "(A)")  "* Set dummy parameters and generate x."
    write (u, "(A)")

    allocate (r (data%get_n_par ()))
    allocate (rb(size (r)))
    allocate (x (size (r)))

    r  = 0
    rb = 0
    x_free = 1
    call sf_int%generate_free (r, rb, x_free)
    call sf_int%complete_kinematics (x, f, r, rb, map=.false.)

    write (u, "(A,9(1x,F10.7))")  "r =", r
    write (u, "(A,9(1x,F10.7))")  "rb=", rb
    write (u, "(A,9(1x,F10.7))")  "x =", x
    write (u, "(A,9(1x,F10.7))")  "f =", f
    write (u, "(A,9(1x,F10.7))")  "xf=", x_free
    select type (sf_int)
    type is (beam_events_t)     
       write (u, "(A,1x,I0)")  "count =", sf_int%count
    end select

    write (u, "(A)")
    write (u, "(A)")  "* Evaluate"
    write (u, "(A)")

    call sf_int%apply (scale = 0._default)
    call sf_int%write (u)

    write (u, "(A)")
    write (u, "(A)")  "* Generate more events, rewind"
    write (u, "(A)")

    select type (sf_int)
    type is (beam_events_t)     
       do i = 1, 3
          call sf_int%generate_free (r, rb, x_free)
          write (u, "(A,9(1x,F10.7))")  "r =", r
          write (u, "(A,1x,I0)")  "count =", sf_int%count
       end do
    end select

    write (u, "(A)")
    write (u, "(A)")  "* Cleanup"

    call sf_int%final ()
    call model%final ()
    
    write (u, "(A)")
    write (u, "(A)")  "* Test output end: sf_beam_events_2"

  end subroutine sf_beam_events_2

\eatline
\nwendcode{}\nwbegindocs{1314}\nwdocspar
\subsubsection{Check the file handle registry}
Open and close some files, checking the registry contents.
\nwenddocs{}\nwbegincode{1315}\moddef{SF beam events: execute tests}\plusendmoddef
  call test (sf_beam_events_3, "sf_beam_events_3", &
       "check registry", &
       u, results)
\nwendcode{}\nwbegincode{1316}\moddef{SF beam events: tests}\plusendmoddef
  subroutine sf_beam_events_3 (u)
    integer, intent(in) :: u
    integer :: u1
    
    write (u, "(A)")  "* Test output: sf_beam_events_2"
    write (u, "(A)")  "*   Purpose: check file handle registry"
    write (u, "(A)")

    write (u, "(A)")  "* Create some empty files"
    write (u, "(A)")

    u1 = free_unit ()
    open (u1, file = "sf_beam_events_f1.tmp", action="write", status="new")
    close (u1)
    open (u1, file = "sf_beam_events_f2.tmp", action="write", status="new")
    close (u1)
    open (u1, file = "sf_beam_events_f3.tmp", action="write", status="new")
    close (u1)
    
    write (u, "(A)")  "* Empty registry"
    write (u, "(A)")
    
    call beam_file_registry%write (u)
    
    write (u, "(A)")
    write (u, "(A)")  "* Insert three entries"
    write (u, "(A)")
    
    call beam_file_registry%open (var_str ("sf_beam_events_f3.tmp"))
    call beam_file_registry%open (var_str ("sf_beam_events_f2.tmp"))
    call beam_file_registry%open (var_str ("sf_beam_events_f1.tmp"))
    call beam_file_registry%write (u)

    write (u, "(A)")
    write (u, "(A)")  "* Open a second channel"
    write (u, "(A)")
    
    call beam_file_registry%open (var_str ("sf_beam_events_f2.tmp"))
    call beam_file_registry%write (u)

    write (u, "(A)")
    write (u, "(A)")  "* Close second entry twice"
    write (u, "(A)")
    
    call beam_file_registry%close (var_str ("sf_beam_events_f2.tmp"))
    call beam_file_registry%close (var_str ("sf_beam_events_f2.tmp"))
    call beam_file_registry%write (u)

    write (u, "(A)")
    write (u, "(A)")  "* Close last entry"
    write (u, "(A)")
    
    call beam_file_registry%close (var_str ("sf_beam_events_f3.tmp"))
    call beam_file_registry%write (u)

    write (u, "(A)")
    write (u, "(A)")  "* Close remaining entry"
    write (u, "(A)")
    
    call beam_file_registry%close (var_str ("sf_beam_events_f1.tmp"))
    call beam_file_registry%write (u)

    write (u, "(A)")
    write (u, "(A)")  "* Cleanup"

    open (u1, file = "sf_beam_events_f1.tmp", action="write")
    close (u1, status = "delete")
    open (u1, file = "sf_beam_events_f2.tmp", action="write")
    close (u1, status = "delete")
    open (u1, file = "sf_beam_events_f3.tmp", action="write")
    close (u1, status = "delete")
    
    write (u, "(A)")
    write (u, "(A)")  "* Test output end: sf_beam_events_3"

  end subroutine sf_beam_events_3

\eatline
\nwendcode{}\nwbegindocs{1317}\nwdocspar
\clearpage
%------------------------------------------------------------------------
\section{Lepton collider beamstrahlung: CIRCE1}

\nwenddocs{}\nwbegincode{1318}\moddef{\code{}sf{\char95}circe1.f90\edoc{}}\endmoddef
\LA{}File header\RA{}

module sf_circe1

\LA{}Use kinds\RA{}
  use kinds, only: double
\LA{}Use strings\RA{}
  use io_units
  use format_defs, only: FMT_17, FMT_19
  use unit_tests
  use diagnostics
  use physics_defs, only: UNDEFINED, ELECTRON, PHOTON
  use lorentz
  use rng_base
  use pdg_arrays
  use model_data
  use flavors
  use colors
  use quantum_numbers
  use state_matrices
  use polarizations
  use interactions
  use sf_mappings
  use sf_aux
  use sf_base
  use circe1, circe1_rng_t => rng_type !NODEP!

\LA{}Standard module head\RA{}

\LA{}SF circe1: public\RA{}

\LA{}SF circe1: types\RA{}

contains

\LA{}SF circe1: procedures\RA{}

\LA{}SF circe1: tests\RA{}

end module sf_circe1
\eatline
\nwendcode{}\nwbegindocs{1319}\nwdocspar
\subsection{Physics} 
Beamstrahlung is applied before ISR.  The {\tt{}CIRCE1} implementation has 
a single structure function for both beams (which makes sense since it 
has to be switched on or off for both beams simultaneously). 
Nevertheless it is factorized: 
 
The functional form in the {\tt{}CIRCE1} parameterization is defined for 
electrons or photons 
\begin{equation} 
  f(x) = \alpha\,x^\beta\,(1-x)^\gamma 
\end{equation} 
for $x<1-\epsilon$ (resp.\ $x>\epsilon$ in the photon case).  In the 
remaining interval, the standard form is zero, with a delta 
singularity at $x=1$ (resp.\ $x=0$).  Equivalently, the delta part may be 
distributed uniformly among this interval.  This latter form is 
implemented in the {\tt{}kirke} version of the {\tt{}CIRCE1} subroutines, and 
is used here. 
 
The parameter {\tt{}circe1{\char92}{\char95}eps} sets the peak mapping of the {\tt{}CIRCE1}
structure function. Its default value is $10^{-5}$. 
The other parameters are the parameterization version and revision 
number, the accelerator type, and the $\sqrt{s}$ value used by 
{\tt{}CIRCE1}. The chattiness can also be set. 
 
Since the energy is distributed in a narrow region around unity (for 
electrons) or zero (for photons), it is advantageous to map the 
interval first.  The mapping is controlled by the parameter
{\tt{}circe1{\char92}{\char95}epsilon} which is taken from the {\tt{}CIRCE1} 
internal data structure. 
 
The $\sqrt{s}$ value, if not explicitly set, is taken from the
process data.  Note that interpolating $\sqrt{s}$ is not recommended;
one should rather choose one of the distinct values known to {\tt{}CIRCE1}.
 
\subsection{The CIRCE1 data block} 
The CIRCE1 parameters are: The incoming flavors, the flags whether the photon
or the lepton is the parton in the hard interaction, the flags for the
generation mode (generator/mapping/no mapping), the mapping parameter
$\epsilon$, $\sqrt{s}$ and several steering parameters: {\tt{}ver},
{\tt{}rev}, {\tt{}acc}, {\tt{}chat}. 

In generator mode, the $x$ values are actually discarded and a random number
generator is used instead.
\nwenddocs{}\nwbegincode{1320}\moddef{SF circe1: public}\endmoddef
  public :: circe1_data_t 
\nwendcode{}\nwbegincode{1321}\moddef{SF circe1: types}\endmoddef
  type, extends (sf_data_t) :: circe1_data_t 
     private 
     class(model_data_t), pointer :: model => null () 
     type(flavor_t), dimension(2) :: flv_in 
     integer, dimension(2) :: pdg_in
     real(default), dimension(2) :: m_in = 0
     logical, dimension(2) :: photon = .false.
     logical :: generate = .false.
     class(rng_factory_t), allocatable :: rng_factory
     real(default) :: sqrts = 0
     real(default) :: eps = 0
     integer :: ver = 0 
     integer :: rev = 0      
     character(6) :: acc = "?"
     integer :: chat = 0    
   contains  
     \LA{}SF circe1: circe1 data: TBP\RA{}     
  end type circe1_data_t 
 
\eatline
\nwendcode{}\nwbegindocs{1322}\nwdocspar
\nwenddocs{}\nwbegincode{1323}\moddef{SF circe1: circe1 data: TBP}\endmoddef
  procedure :: init => circe1_data_init 
\nwendcode{}\nwbegincode{1324}\moddef{SF circe1: procedures}\endmoddef
  subroutine circe1_data_init &
       (data, model, pdg_in, sqrts, eps, out_photon, &
        ver, rev, acc, chat)
    class(circe1_data_t), intent(out) :: data 
    class(model_data_t), intent(in), target :: model
    type(pdg_array_t), dimension(2), intent(in) :: pdg_in
    real(default), intent(in) :: sqrts 
    real(default), intent(in) :: eps
    logical, dimension(2), intent(in) :: out_photon
    character(*), intent(in) :: acc
    integer, intent(in) :: ver, rev, chat 
    data%model => model 
    if (any (pdg_array_get_length (pdg_in) /= 1)) then
       call msg_fatal ("CIRCE1: incoming beam particles must be unique")
    end if
    call data%flv_in(1)%init (pdg_array_get (pdg_in(1), 1), model)
    call data%flv_in(2)%init (pdg_array_get (pdg_in(2), 1), model)
    data%pdg_in = data%flv_in%get_pdg ()
    data%m_in = data%flv_in%get_mass ()
    data%sqrts = sqrts
    data%eps = eps
    data%photon = out_photon
    data%ver = ver
    data%rev = rev
    data%acc = acc 
    data%chat = chat 
    call data%check ()
    call circex (0.d0, 0.d0, dble (data%sqrts), &
         data%acc, data%ver, data%rev, data%chat)
  end subroutine circe1_data_init 
 
\eatline
\nwendcode{}\nwbegindocs{1325}Activate the generator mode.  We import a RNG factory into the data
type, which can then spawn RNG generator objects.
\nwenddocs{}\nwbegincode{1326}\moddef{SF circe1: circe1 data: TBP}\plusendmoddef
  procedure :: set_generator_mode => circe1_data_set_generator_mode
\nwendcode{}\nwbegincode{1327}\moddef{SF circe1: procedures}\plusendmoddef
  subroutine circe1_data_set_generator_mode (data, rng_factory)
    class(circe1_data_t), intent(inout) :: data
    class(rng_factory_t), intent(inout), allocatable :: rng_factory
    data%generate = .true.
    call move_alloc (from = rng_factory, to = data%rng_factory)
  end subroutine circe1_data_set_generator_mode
  
\eatline
\nwendcode{}\nwbegindocs{1328}Handle error conditions.
\nwenddocs{}\nwbegincode{1329}\moddef{SF circe1: circe1 data: TBP}\plusendmoddef
  procedure :: check => circe1_data_check 
\nwendcode{}\nwbegincode{1330}\moddef{SF circe1: procedures}\plusendmoddef
  subroutine circe1_data_check (data) 
    class(circe1_data_t), intent(in) :: data 
    type(flavor_t) :: flv_electron, flv_photon
    call flv_electron%init (ELECTRON, data%model)
    call flv_photon%init (PHOTON, data%model)
    if (.not. flv_electron%is_defined () &
         .or. .not. flv_photon%is_defined ()) then
       call msg_fatal ("CIRCE1: model must contain photon and electron")
    end if
    if (any (abs (data%pdg_in) /= ELECTRON) &
         .or. (data%pdg_in(1) /= - data%pdg_in(2))) then
       call msg_fatal ("CIRCE1: applicable only for e+e- or e-e+ collisions")
    end if
    if (data%eps <= 0) then
       call msg_error ("CIRCE1: circe1_eps = 0: integration will &
            &miss x=1 peak")
    end if
  end subroutine circe1_data_check 
 
\eatline
\nwendcode{}\nwbegindocs{1331}Output 
\nwenddocs{}\nwbegincode{1332}\moddef{SF circe1: circe1 data: TBP}\plusendmoddef
  procedure :: write => circe1_data_write 
\nwendcode{}\nwbegincode{1333}\moddef{SF circe1: procedures}\plusendmoddef
  subroutine circe1_data_write (data, unit, verbose)
    class(circe1_data_t), intent(in) :: data 
    integer, intent(in), optional :: unit 
    logical, intent(in), optional :: verbose
    integer :: u
    u = given_output_unit (unit);  if (u < 0)  return 
    write (u, "(1x,A)") "CIRCE1 data:" 
    write (u, "(3x,A,2(1x,A))") "prt_in   =", &
         char (data%flv_in(1)%get_name ()), &
         char (data%flv_in(2)%get_name ())
    write (u, "(3x,A,2(1x,L1))")  "photon   =", data%photon
    write (u, "(3x,A,L1)")        "generate = ", data%generate
    write (u, "(3x,A,2(1x," // FMT_19 // "))") "m_in     =", data%m_in
    write (u, "(3x,A," // FMT_19 // ")") "sqrts    = ", data%sqrts 
    write (u, "(3x,A," // FMT_19 // ")") "eps      = ", data%eps
    write (u, "(3x,A,I0)") "ver      = ", data%ver 
    write (u, "(3x,A,I0)") "rev      = ", data%rev 
    write (u, "(3x,A,A)")  "acc      = ", data%acc 
    write (u, "(3x,A,I0)") "chat     = ", data%chat 
    if (data%generate)  call data%rng_factory%write (u)
  end subroutine circe1_data_write 
 
\eatline
\nwendcode{}\nwbegindocs{1334}Return true if this structure function is in generator mode.  In
that case, all parameters are free, otherwise bound.  (We do not
support mixed cases.)  Default is: no generator.
\nwenddocs{}\nwbegincode{1335}\moddef{SF circe1: circe1 data: TBP}\plusendmoddef
  procedure :: is_generator => circe1_data_is_generator
\nwendcode{}\nwbegincode{1336}\moddef{SF circe1: procedures}\plusendmoddef
  function circe1_data_is_generator (data) result (flag)
    class(circe1_data_t), intent(in) :: data
    logical :: flag
    flag = data%generate
  end function circe1_data_is_generator
  
\eatline
\nwendcode{}\nwbegindocs{1337}The number of parameters is two, collinear splitting for the two beams. 
\nwenddocs{}\nwbegincode{1338}\moddef{SF circe1: circe1 data: TBP}\plusendmoddef
  procedure :: get_n_par => circe1_data_get_n_par
\nwendcode{}\nwbegincode{1339}\moddef{SF circe1: procedures}\plusendmoddef
  function circe1_data_get_n_par (data) result (n)
    class(circe1_data_t), intent(in) :: data
    integer :: n
    n = 2
  end function circe1_data_get_n_par
  
\eatline
\nwendcode{}\nwbegindocs{1340}Return the outgoing particles PDG codes. This is either the incoming 
particle (if a photon is radiated), or the photon if that is the particle
of the hard interaction. The latter is determined via the {\tt{}photon}
flag. There are two entries for the two beams.
\nwenddocs{}\nwbegincode{1341}\moddef{SF circe1: circe1 data: TBP}\plusendmoddef
  procedure :: get_pdg_out => circe1_data_get_pdg_out
\nwendcode{}\nwbegincode{1342}\moddef{SF circe1: procedures}\plusendmoddef
  subroutine circe1_data_get_pdg_out (data, pdg_out)
    class(circe1_data_t), intent(in) :: data
    type(pdg_array_t), dimension(:), intent(inout) :: pdg_out
    integer :: i, n
    n = 2
    do i = 1, n
       if (data%photon(i)) then
          pdg_out(i) = PHOTON
       else
          pdg_out(i) = data%pdg_in(i)
       end if
    end do
  end subroutine circe1_data_get_pdg_out
  
\eatline
\nwendcode{}\nwbegindocs{1343}This variant is not inherited, it returns integers.
\nwenddocs{}\nwbegincode{1344}\moddef{SF circe1: circe1 data: TBP}\plusendmoddef
  procedure :: get_pdg_int => circe1_data_get_pdg_int
\nwendcode{}\nwbegincode{1345}\moddef{SF circe1: procedures}\plusendmoddef
  function circe1_data_get_pdg_int (data) result (pdg)
    class(circe1_data_t), intent(in) :: data
    integer, dimension(2) :: pdg
    integer :: i
    do i = 1, 2
       if (data%photon(i)) then
          pdg(i) = PHOTON
       else
          pdg(i) = data%pdg_in(i)
       end if
    end do
  end function circe1_data_get_pdg_int
  
\eatline
\nwendcode{}\nwbegindocs{1346}Allocate the interaction record.
\nwenddocs{}\nwbegincode{1347}\moddef{SF circe1: circe1 data: TBP}\plusendmoddef
  procedure :: allocate_sf_int => circe1_data_allocate_sf_int
\nwendcode{}\nwbegincode{1348}\moddef{SF circe1: procedures}\plusendmoddef
  subroutine circe1_data_allocate_sf_int (data, sf_int)
    class(circe1_data_t), intent(in) :: data
    class(sf_int_t), intent(inout), allocatable :: sf_int
    allocate (circe1_t :: sf_int)
  end subroutine circe1_data_allocate_sf_int
  
\eatline
\nwendcode{}\nwbegindocs{1349}\nwdocspar
\subsection{Random Number Generator for CIRCE}
The CIRCE implementation now supports a generic random-number
generator object that allows for a local state as a component.  To
support this, we must extend the abstract type provided by CIRCE and
delegate the generator call to the (also abstract) RNG used by WHIZARD.
\nwenddocs{}\nwbegincode{1350}\moddef{SF circe1: types}\plusendmoddef
  type, extends (circe1_rng_t) :: rng_obj_t
     class(rng_t), allocatable :: rng
   contains
     procedure :: generate => rng_obj_generate
  end type rng_obj_t
  
\eatline
\nwendcode{}\nwbegincode{1351}\moddef{SF circe1: procedures}\plusendmoddef
  subroutine rng_obj_generate (rng_obj, u)
    class(rng_obj_t), intent(inout) :: rng_obj
    real(double), intent(out) :: u
    real(default) :: x
    call rng_obj%rng%generate (x)
    u = x
  end subroutine rng_obj_generate

\eatline
\nwendcode{}\nwbegindocs{1352}\nwdocspar
\subsection{The CIRCE1 object} 
This is a $2\to 4$ interaction, where, depending on the parameters, any two of
the four outgoing particles are connected to the hard interactions, the others
are radiated.  Knowing that all particles are colorless, we do not have to
deal with color.

The flavors are sorted such that the first two particles are the incoming
leptons, the next two are the radiated particles, and the last two are the
partons initiating the hard interaction.

CIRCE1 does not support polarized beams explicitly.  For simplicity, we
nevertheless carry beam polarization through to the outgoing electrons and
make the photons unpolarized.
\nwenddocs{}\nwbegincode{1353}\moddef{SF circe1: types}\plusendmoddef
  type, extends (sf_int_t) :: circe1_t
     type(circe1_data_t), pointer :: data => null ()
     real(default), dimension(2) :: x = 0
     real(default) :: f = 0
     logical, dimension(2) :: continuum = .true.
     logical, dimension(2) :: peak = .true.
     type(rng_obj_t) :: rng_obj
   contains
   \LA{}SF circe1: circe1: TBP\RA{}
  end type circe1_t 
  
\eatline
\nwendcode{}\nwbegindocs{1354}Type string: has to be here, but there is no string variable on which CIRCE1
depends. Hence, a dummy routine.
\nwenddocs{}\nwbegincode{1355}\moddef{SF circe1: circe1: TBP}\endmoddef
  procedure :: type_string => circe1_type_string
\nwendcode{}\nwbegincode{1356}\moddef{SF circe1: procedures}\plusendmoddef
  function circe1_type_string (object) result (string)
    class(circe1_t), intent(in) :: object
    type(string_t) :: string
    if (associated (object%data)) then
       string = "CIRCE1: beamstrahlung" 
    else
       string = "CIRCE1: [undefined]"
    end if
  end function circe1_type_string
  
\eatline
\nwendcode{}\nwbegindocs{1357}Output.  Call the interaction routine after displaying the configuration.
\nwenddocs{}\nwbegincode{1358}\moddef{SF circe1: circe1: TBP}\plusendmoddef
  procedure :: write => circe1_write
\nwendcode{}\nwbegincode{1359}\moddef{SF circe1: procedures}\plusendmoddef
  subroutine circe1_write (object, unit, testflag)
    class(circe1_t), intent(in) :: object
    integer, intent(in), optional :: unit
    logical, intent(in), optional :: testflag
    integer :: u
    u = given_output_unit (unit)
    if (associated (object%data)) then
       call object%data%write (u)
       if (object%data%generate)  call object%rng_obj%rng%write (u)
       if (object%status >= SF_DONE_KINEMATICS) then
          write (u, "(3x,A,2(1x," // FMT_17 // "))")  "x =", object%x
          if (object%status >= SF_FAILED_EVALUATION) then
             write (u, "(3x,A,1x," // FMT_17 // ")")  "f =", object%f
          end if
       end if
       call object%base_write (u, testflag)
    else
       write (u, "(1x,A)")  "CIRCE1 data: [undefined]"
    end if
  end subroutine circe1_write
    
\eatline
\nwendcode{}\nwbegindocs{1360}\nwdocspar
\nwenddocs{}\nwbegincode{1361}\moddef{SF circe1: circe1: TBP}\plusendmoddef
  procedure :: init => circe1_init
\nwendcode{}\nwbegincode{1362}\moddef{SF circe1: procedures}\plusendmoddef
  subroutine circe1_init (sf_int, data) 
    class(circe1_t), intent(out) :: sf_int 
    class(sf_data_t), intent(in), target :: data
    logical, dimension(6) :: mask_h
    type(quantum_numbers_mask_t), dimension(6) :: mask 
    integer, dimension(6) :: hel_lock 
    type(polarization_t) :: pol1, pol2
    type(quantum_numbers_t), dimension(1) :: qn_fc1, qn_hel1, qn_fc2, qn_hel2
    type(flavor_t) :: flv_photon
    type(color_t) :: col0
    real(default), dimension(2) :: mi2, mr2, mo2
    type(quantum_numbers_t) :: qn_photon, qn1, qn2
    type(quantum_numbers_t), dimension(6) :: qn
    type(state_iterator_t) :: it_hel1, it_hel2
    hel_lock = 0
    mask_h = .false.
    select type (data)
    type is (circe1_data_t)
       mi2 = data%m_in**2
       if (data%photon(1)) then
          hel_lock(1) = 3;  hel_lock(3) = 1;  mask_h(5) = .true.
          mr2(1) = mi2(1)
          mo2(1) = 0._default
       else
          hel_lock(1) = 5;  hel_lock(5) = 1;  mask_h(3) = .true.
          mr2(1) = 0._default
          mo2(1) = mi2(1)
       end if
       if (data%photon(2)) then
          hel_lock(2) = 4;  hel_lock(4) = 2;  mask_h(6) = .true.
          mr2(2) = mi2(2)
          mo2(2) = 0._default
       else
          hel_lock(2) = 6;  hel_lock(6) = 2;  mask_h(4) = .true.
          mr2(2) = 0._default
          mo2(2) = mi2(2)
       end if
       mask = quantum_numbers_mask (.false., .false., mask_h)
       call sf_int%base_init (mask, mi2, mr2, mo2, &
            hel_lock = hel_lock)
       sf_int%data => data
       call flv_photon%init (PHOTON, data%model)
       call col0%init ()
       call qn_photon%init (flv_photon, col0)
       call polarization_init_generic (pol1, data%flv_in(1))
       call qn_fc1(1)%init (flv = data%flv_in(1), col = col0)
       call polarization_init_generic (pol2, data%flv_in(2))
       call qn_fc2(1)%init (flv = data%flv_in(2), col = col0)
       call it_hel1%init (pol1%state)
       
       do while (it_hel1%is_valid ()) 
          qn_hel1 = it_hel1%get_quantum_numbers ()
          qn1 = qn_hel1(1) .merge. qn_fc1(1) 
          qn(1) = qn1
          if (data%photon(1)) then
             qn(3) = qn1;  qn(5) = qn_photon
          else
             qn(3) = qn_photon;  qn(5) = qn1
          end if
          call it_hel2%init (pol2%state) 
          do while (it_hel2%is_valid ()) 
             qn_hel2 = it_hel2%get_quantum_numbers () 
             qn2 = qn_hel2(1) .merge. qn_fc2(1) 
             qn(2) = qn2
             if (data%photon(2)) then
                qn(4) = qn2;  qn(6) = qn_photon
             else
                qn(4) = qn_photon;  qn(6) = qn2
             end if
             call qn(3:4)%tag_radiated ()
             call interaction_add_state (sf_int%interaction_t, qn)
             call it_hel2%advance () 
          end do
          call it_hel1%advance ()
       end do
       call polarization_final (pol1)
       call polarization_final (pol2)
       call interaction_freeze (sf_int%interaction_t) 
       call sf_int%set_incoming ([1,2])
       call sf_int%set_radiated ([3,4])
       call sf_int%set_outgoing ([5,6])
       sf_int%status = SF_INITIAL
    end select
    if (sf_int%data%generate) then
       call sf_int%data%rng_factory%make (sf_int%rng_obj%rng)
    end if
  end subroutine circe1_init

\eatline
\nwendcode{}\nwbegindocs{1363}\nwdocspar
\subsection{Kinematics}
Refer to the {\tt{}data} component.
\nwenddocs{}\nwbegincode{1364}\moddef{SF circe1: circe1: TBP}\plusendmoddef
  procedure :: is_generator => circe1_is_generator
\nwendcode{}\nwbegincode{1365}\moddef{SF circe1: procedures}\plusendmoddef
  function circe1_is_generator (sf_int) result (flag)
    class(circe1_t), intent(in) :: sf_int
    logical :: flag
    flag = sf_int%data%is_generator ()
  end function circe1_is_generator
  
\eatline
\nwendcode{}\nwbegindocs{1366}Generate free parameters, if generator mode is on.  Otherwise, the
parameters will be discarded.
\nwenddocs{}\nwbegincode{1367}\moddef{SF circe1: circe1: TBP}\plusendmoddef
  procedure :: generate_free => circe1_generate_free
\nwendcode{}\nwbegincode{1368}\moddef{SF circe1: procedures}\plusendmoddef
  subroutine circe1_generate_free (sf_int, r, rb,  x_free)
    class(circe1_t), intent(inout) :: sf_int
    real(default), dimension(:), intent(out) :: r, rb
    real(default), intent(inout) :: x_free

    if (sf_int%data%generate) then
       call circe_generate (r, sf_int%data%get_pdg_int (), sf_int%rng_obj)
       rb = 1 - r
       x_free = x_free * product (r)
    else
       r = 0
       rb= 1
    end if
  end subroutine circe1_generate_free
    
\eatline
\nwendcode{}\nwbegindocs{1369}Generator mode: depending on the particle codes, call one of the
available {\tt{}girce} generators.  Illegal particle code combinations
should have been caught during data initialization.
\nwenddocs{}\nwbegincode{1370}\moddef{SF circe1: procedures}\plusendmoddef
  subroutine circe_generate (x, pdg, rng_obj)
    real(default), dimension(2), intent(out) :: x
    integer, dimension(2), intent(in) :: pdg
    class(rng_obj_t), intent(inout) :: rng_obj
    real(double) :: xc1, xc2
    select case (abs (pdg(1)))
    case (ELECTRON)
       select case (abs (pdg(2)))
       case (ELECTRON)
          call gircee (xc1, xc2, rng_obj = rng_obj)
       case (PHOTON)
          call girceg (xc1, xc2, rng_obj = rng_obj)
       end select
    case (PHOTON)
       select case (abs (pdg(2)))
       case (ELECTRON)
          call girceg (xc2, xc1, rng_obj = rng_obj)
       case (PHOTON)
          call gircgg (xc1, xc2, rng_obj = rng_obj)
       end select
    end select
    x = [xc1, xc2]
  end subroutine circe_generate

\eatline
\nwendcode{}\nwbegindocs{1371}Set kinematics.  The $r$ values (either from integration or from the
generator call above) are copied to $x$ unchanged, and $f$ is unity.
We store the $x$ values, so we can use them for the evaluation later.
\nwenddocs{}\nwbegincode{1372}\moddef{SF circe1: circe1: TBP}\plusendmoddef
  procedure :: complete_kinematics => circe1_complete_kinematics
\nwendcode{}\nwbegincode{1373}\moddef{SF circe1: procedures}\plusendmoddef
  subroutine circe1_complete_kinematics (sf_int, x, f, r, rb, map)
    class(circe1_t), intent(inout) :: sf_int
    real(default), dimension(:), intent(out) :: x
    real(default), intent(out) :: f
    real(default), dimension(:), intent(in) :: r
    real(default), dimension(:), intent(in) :: rb
    logical, intent(in) :: map
    real(default), dimension(2) :: xb1
    x = r
    sf_int%x = x
    f = 1
    xb1 = 1 - x
    call sf_int%split_momenta (x, xb1)
    select case (sf_int%status)
    case (SF_FAILED_KINEMATICS);  f = 0
    end select
  end subroutine circe1_complete_kinematics

\eatline
\nwendcode{}\nwbegindocs{1374}Compute inverse kinematics.  In generator mode, the $r$ values are
meaningless, but we copy them anyway.
\nwenddocs{}\nwbegincode{1375}\moddef{SF circe1: circe1: TBP}\plusendmoddef
  procedure :: inverse_kinematics => circe1_inverse_kinematics
\nwendcode{}\nwbegincode{1376}\moddef{SF circe1: procedures}\plusendmoddef
  subroutine circe1_inverse_kinematics (sf_int, x, f, r, rb, map, set_momenta)
    class(circe1_t), intent(inout) :: sf_int
    real(default), dimension(:), intent(in) :: x
    real(default), intent(out) :: f
    real(default), dimension(:), intent(out) :: r
    real(default), dimension(:), intent(out) :: rb
    logical, intent(in) :: map
    logical, intent(in), optional :: set_momenta
    real(default), dimension(2) :: xb1
    logical :: set_mom
    set_mom = .false.;  if (present (set_momenta))  set_mom = set_momenta
    r = x
    rb = 1 - x
    sf_int%x = x
    f = 1
    if (set_mom) then
       xb1 = 1 - x
       call sf_int%split_momenta (x, xb1)
       select case (sf_int%status)
       case (SF_FAILED_KINEMATICS);  f = 0
       end select
    end if
  end subroutine circe1_inverse_kinematics

\eatline
\nwendcode{}\nwbegindocs{1377} 
\subsection{CIRCE1 application} 
CIRCE is applied for the two beams at once.  We can safely assume that no
structure functions are applied before this, so the incoming particles are
on-shell electrons/positrons.

The scale is ignored.
\nwenddocs{}\nwbegincode{1378}\moddef{SF circe1: circe1: TBP}\plusendmoddef
  procedure :: apply => circe1_apply
\nwendcode{}\nwbegincode{1379}\moddef{SF circe1: procedures}\plusendmoddef
  subroutine circe1_apply (sf_int, scale)
    class(circe1_t), intent(inout) :: sf_int
    real(default), intent(in) :: scale
    real(default), dimension(2) :: xb
    real(double), dimension(2) :: xc
    real(double), parameter :: one = 1
    associate (data => sf_int%data)
      xc = sf_int%x
      xb = 1 - sf_int%x
      if (data%generate) then
         sf_int%f = 1
      else
         sf_int%f = 0
         if (all (sf_int%continuum)) then
            sf_int%f = circe (xc(1), xc(2), data%pdg_in(1), data%pdg_in(2))
         end if
         if (sf_int%continuum(2) .and. sf_int%peak(1)) then
            sf_int%f = sf_int%f &
                 + circe (one, xc(2), data%pdg_in(1), data%pdg_in(2)) &
                 * peak (xb(1), data%eps)
         end if
         if (sf_int%continuum(1) .and. sf_int%peak(2)) then
            sf_int%f = sf_int%f &
                 + circe (xc(1), one, data%pdg_in(1), data%pdg_in(2)) &
                 * peak (xb(2), data%eps)
         end if
         if (all (sf_int%peak)) then
            sf_int%f = sf_int%f &
                 + circe (one, one, data%pdg_in(1), data%pdg_in(2)) &
                 * peak (xb(1), data%eps) * peak (xb(2), data%eps)
         end if
      end if
    end associate
    call interaction_set_matrix_element &
        (sf_int%interaction_t, cmplx (sf_int%f, kind=default)) 
    sf_int%status = SF_EVALUATED
  end subroutine circe1_apply
 
\eatline
\nwendcode{}\nwbegindocs{1380}This is a smeared delta peak at zero, as an endpoint singularity.
We choose an exponentially decreasing function, starting at zero, with
integral (from $0$ to $1$) $1-e^{-1/\epsilon}$.  For small $\epsilon$,
this reduces to one.
\nwenddocs{}\nwbegincode{1381}\moddef{SF circe1: procedures}\plusendmoddef
  function peak (x, eps) result (f)
    real(default), intent(in) :: x, eps
    real(default) :: f
    f = exp (-x / eps) / eps
  end function peak
  
\eatline
\nwendcode{}\nwbegindocs{1382}\nwdocspar
\subsection{Unit tests}
\nwenddocs{}\nwbegincode{1383}\moddef{SF circe1: public}\plusendmoddef
  public :: sf_circe1_test
\nwendcode{}\nwbegincode{1384}\moddef{SF circe1: tests}\endmoddef
  subroutine sf_circe1_test (u, results)
    integer, intent(in) :: u
    type(test_results_t), intent(inout) :: results
  \LA{}SF circe1: execute tests\RA{}
  end subroutine sf_circe1_test
  
\eatline
\nwendcode{}\nwbegindocs{1385}\nwdocspar
\subsubsection{Test structure function data}
Construct and display a test structure function data object.
\nwenddocs{}\nwbegincode{1386}\moddef{SF circe1: execute tests}\endmoddef
  call test (sf_circe1_1, "sf_circe1_1", &
       "structure function configuration", &
       u, results)
\nwendcode{}\nwbegincode{1387}\moddef{SF circe1: tests}\plusendmoddef
  subroutine sf_circe1_1 (u)
    integer, intent(in) :: u
    type(model_data_t), target :: model
    type(pdg_array_t), dimension(2) :: pdg_in
    type(pdg_array_t), dimension(2) :: pdg_out
    integer, dimension(:), allocatable :: pdg1, pdg2
    class(sf_data_t), allocatable :: data
    
    write (u, "(A)")  "* Test output: sf_circe1_1"
    write (u, "(A)")  "*   Purpose: initialize and display &
         &CIRCE structure function data"
    write (u, "(A)")
    
    write (u, "(A)")  "* Create empty data object"
    write (u, "(A)")

    call model%init_qed_test ()
    pdg_in(1) = ELECTRON
    pdg_in(2) = -ELECTRON

    allocate (circe1_data_t :: data)
    call data%write (u)

    write (u, "(A)")
    write (u, "(A)")  "* Initialize"
    write (u, "(A)")

    select type (data)
    type is (circe1_data_t)
       call data%init (model, pdg_in, &
            sqrts = 500._default, &
            eps = 1e-6_default, &
            out_photon = [.false., .false.], &
            ver = 0, &
            rev = 0, &
            acc = "SBAND", &
            chat = 0)
    end select

    call data%write (u)

    write (u, "(A)")

    write (u, "(1x,A)")  "Outgoing particle codes:"
    call data%get_pdg_out (pdg_out)
    pdg1 = pdg_out(1)
    pdg2 = pdg_out(2)
    write (u, "(2x,99(1x,I0))")  pdg1, pdg2

    call model%final ()
    
    write (u, "(A)")
    write (u, "(A)")  "* Test output end: sf_circe1_1"

  end subroutine sf_circe1_1

\eatline
\nwendcode{}\nwbegindocs{1388}\nwdocspar
\subsubsection{Test and probe structure function}
Construct and display a structure function object based on the PDF builtin
structure function.
\nwenddocs{}\nwbegincode{1389}\moddef{SF circe1: execute tests}\plusendmoddef
  call test (sf_circe1_2, "sf_circe1_2", &
       "structure function instance", &
       u, results)
\nwendcode{}\nwbegincode{1390}\moddef{SF circe1: tests}\plusendmoddef
  subroutine sf_circe1_2 (u)
    integer, intent(in) :: u
    type(model_data_t), target :: model
    type(flavor_t), dimension(2) :: flv
    type(pdg_array_t), dimension(2) :: pdg_in
    class(sf_data_t), allocatable, target :: data
    class(sf_int_t), allocatable :: sf_int
    type(vector4_t) :: k1, k2
    type(vector4_t), dimension(4) :: q
    real(default) :: E
    real(default), dimension(:), allocatable :: r, rb, x
    real(default) :: f
    
    write (u, "(A)")  "* Test output: sf_circe1_2"
    write (u, "(A)")  "*   Purpose: initialize and fill &
         &circe1 structure function object"
    write (u, "(A)")
    
    write (u, "(A)")  "* Initialize configuration data"
    write (u, "(A)")

    call model%init_qed_test ()
    call flv(1)%init (ELECTRON, model)
    call flv(2)%init (-ELECTRON, model)
    pdg_in(1) = ELECTRON
    pdg_in(2) = -ELECTRON

    call reset_interaction_counter ()
    
    allocate (circe1_data_t :: data)
    select type (data)
    type is (circe1_data_t)
       call data%init (model, pdg_in, &
            sqrts = 500._default, &
            eps = 1e-6_default, &
            out_photon = [.false., .false.], &
            ver = 0, &
            rev = 0, &
            acc = "SBAND", &
            chat = 0)
    end select
       
    write (u, "(A)")  "* Initialize structure-function object"
    write (u, "(A)")
    
    call data%allocate_sf_int (sf_int)
    call sf_int%init (data)
    call sf_int%set_beam_index ([1,2])
    
    call sf_int%write (u)

    write (u, "(A)")
    write (u, "(A)")  "* Initialize incoming momentum with E=500"
    write (u, "(A)")
    E = 250
    k1 = vector4_moving (E, sqrt (E**2 - flv(1)%get_mass ()**2), 3)
    k2 = vector4_moving (E,-sqrt (E**2 - flv(2)%get_mass ()**2), 3)
    call vector4_write (k1, u)
    call vector4_write (k2, u)
    call sf_int%seed_kinematics ([k1, k2])

    write (u, "(A)")
    write (u, "(A)")  "* Set kinematics for x=0.95,0.85."
    write (u, "(A)")

    allocate (r (data%get_n_par ()))
    allocate (rb(size (r)))
    allocate (x (size (r)))

    r = [0.9_default, 0.8_default]
    call sf_int%complete_kinematics (x, f, r, rb, map=.false.)
    call sf_int%write (u)

    write (u, "(A)")
    write (u, "(A,9(1x,F10.7))")  "x =", x
    write (u, "(A,9(1x,F10.7))")  "f =", f

    write (u, "(A)")
    write (u, "(A)")  "* Recover x from momenta"
    write (u, "(A)")

    q = interaction_get_momenta (sf_int%interaction_t, outgoing=.true.)
    call sf_int%final ()
    deallocate (sf_int)

    call data%allocate_sf_int (sf_int)
    call sf_int%init (data)
    call sf_int%set_beam_index ([1, 2])

    call sf_int%seed_kinematics ([k1, k2])
    call interaction_set_momenta (sf_int%interaction_t, q, outgoing=.true.)
    call sf_int%recover_x (x)

    write (u, "(A,9(1x,F10.7))")  "x =", x

    write (u, "(A)")
    write (u, "(A)")  "* Evaluate"
    write (u, "(A)")

    call sf_int%complete_kinematics (x, f, r, rb, map=.false.) 
    call sf_int%apply (scale = 0._default)
    call sf_int%write (u)

    write (u, "(A)")
    write (u, "(A)")  "* Cleanup"

    call sf_int%final ()
    call model%final ()
    
    write (u, "(A)")
    write (u, "(A)")  "* Test output end: sf_circe1_2"

  end subroutine sf_circe1_2

\eatline
\nwendcode{}\nwbegindocs{1391}\nwdocspar
\subsubsection{Generator mode}
Construct and evaluate a structure function object in generator mode.
\nwenddocs{}\nwbegincode{1392}\moddef{SF circe1: execute tests}\plusendmoddef
  call test (sf_circe1_3, "sf_circe1_3", &
       "generator mode", &
       u, results)
\nwendcode{}\nwbegincode{1393}\moddef{SF circe1: tests}\plusendmoddef
  subroutine sf_circe1_3 (u)
    integer, intent(in) :: u
    type(model_data_t), target :: model
    type(flavor_t), dimension(2) :: flv
    type(pdg_array_t), dimension(2) :: pdg_in
    class(sf_data_t), allocatable, target :: data
    class(rng_factory_t), allocatable :: rng_factory
    class(sf_int_t), allocatable :: sf_int
    type(vector4_t) :: k1, k2
    real(default) :: E
    real(default), dimension(:), allocatable :: r, rb, x
    real(default) :: f, x_free
    
    write (u, "(A)")  "* Test output: sf_circe1_3"
    write (u, "(A)")  "*   Purpose: initialize and fill &
         &circe1 structure function object"
    write (u, "(A)")
    
    write (u, "(A)")  "* Initialize configuration data"
    write (u, "(A)")

    call model%init_qed_test ()
    call flv(1)%init (ELECTRON, model)
    call flv(2)%init (-ELECTRON, model)
    pdg_in(1) = ELECTRON
    pdg_in(2) = -ELECTRON

    call reset_interaction_counter ()
    
    allocate (circe1_data_t :: data)
    allocate (rng_test_factory_t :: rng_factory)
    select type (data)
    type is (circe1_data_t)
       call data%init (model, pdg_in, &
            sqrts = 500._default, &
            eps = 1e-6_default, &
            out_photon = [.false., .false.], &
            ver = 0, &
            rev = 0, &
            acc = "SBAND", &
            chat = 0)
       call data%set_generator_mode (rng_factory)
    end select
       
    write (u, "(A)")  "* Initialize structure-function object"
    write (u, "(A)")
    
    call data%allocate_sf_int (sf_int)
    call sf_int%init (data)
    call sf_int%set_beam_index ([1,2])
    select type (sf_int)
    type is (circe1_t)
       call sf_int%rng_obj%rng%init (3)
    end select

    write (u, "(A)")  "* Initialize incoming momentum with E=500"
    write (u, "(A)")
    E = 250
    k1 = vector4_moving (E, sqrt (E**2 - flv(1)%get_mass ()**2), 3)
    k2 = vector4_moving (E,-sqrt (E**2 - flv(2)%get_mass ()**2), 3)
    call vector4_write (k1, u)
    call vector4_write (k2, u)
    call sf_int%seed_kinematics ([k1, k2])

    write (u, "(A)")
    write (u, "(A)")  "* Generate x"
    write (u, "(A)")

    allocate (r (data%get_n_par ()))
    allocate (rb(size (r)))
    allocate (x (size (r)))

    r  = 0
    rb = 0
    x_free = 1
    call sf_int%generate_free (r, rb, x_free)
    call sf_int%complete_kinematics (x, f, r, rb, map=.false.)

    write (u, "(A,9(1x,F10.7))")  "x =", x
    write (u, "(A,9(1x,F10.7))")  "f =", f
    write (u, "(A,9(1x,F10.7))")  "xf=", x_free

    write (u, "(A)")
    write (u, "(A)")  "* Evaluate"
    write (u, "(A)")

    call sf_int%apply (scale = 0._default)
    call sf_int%write (u)

    write (u, "(A)")
    write (u, "(A)")  "* Cleanup"

    call sf_int%final ()
    call model%final ()
    
    write (u, "(A)")
    write (u, "(A)")  "* Test output end: sf_circe1_3"

  end subroutine sf_circe1_3

\eatline
\nwendcode{}\nwbegindocs{1394}\nwdocspar
\clearpage
%------------------------------------------------------------------------
\section{Lepton Collider Beamstrahlung and Photon collider: CIRCE2} 

\nwenddocs{}\nwbegincode{1395}\moddef{\code{}sf{\char95}circe2.f90\edoc{}}\endmoddef
\LA{}File header\RA{}

module sf_circe2

\LA{}Use kinds\RA{}
  use kinds, only: double
\LA{}Use strings\RA{}
  use io_units
  use format_defs, only: FMT_19
  use unit_tests
  use diagnostics
  use os_interface
  use physics_defs, only: UNDEFINED, PHOTON, ELECTRON
  use lorentz
  use rng_base
  use selectors
  use pdg_arrays
  use model_data
  use flavors
  use colors
  use helicities
  use quantum_numbers
  use state_matrices
  use polarizations
  use interactions
  use sf_aux
  use sf_base
  use circe2, circe2_rng_t => rng_type !NODEP!

\LA{}Standard module head\RA{}

\LA{}SF circe2: public\RA{}

\LA{}SF circe2: types\RA{}

contains

\LA{}SF circe2: procedures\RA{}

\LA{}SF circe2: tests\RA{}

end module sf_circe2
\eatline
\nwendcode{}\nwbegindocs{1396}\nwdocspar
\subsection{Physics}
{\tt{}CIRCE2} describes photon spectra
Beamstrahlung is applied before ISR.  The {\tt{}CIRCE2} implementation has
a single structure function for both beams (which makes sense since it
has to be switched on or off for both beams simultaneously).


\subsection{The CIRCE2 data block}
The CIRCE2 parameters are: file and collider specification, incoming
(= outgoing) particles.  The luminosity is returned by {\tt{}circe2{\char95}luminosity}.
\nwenddocs{}\nwbegincode{1397}\moddef{SF circe2: public}\endmoddef
  public :: circe2_data_t
\nwendcode{}\nwbegincode{1398}\moddef{SF circe2: types}\endmoddef
  type, extends (sf_data_t) :: circe2_data_t
     private
     class(model_data_t), pointer :: model => null () 
     type(flavor_t), dimension(2) :: flv_in
     integer, dimension(2) :: pdg_in
     real(default) :: sqrts = 0
     logical :: polarized = .false.
     class(rng_factory_t), allocatable :: rng_factory
     type(string_t) :: filename
     type(string_t) :: file 
     type(string_t) :: design 
     real(default) :: lumi = 0
     real(default), dimension(4) :: lumi_hel_frac = 0
     integer, dimension(0:4) :: h1 = [0, -1, -1, 1, 1]
     integer, dimension(0:4) :: h2 = [0, -1,  1,-1, 1]
     integer :: error = 1
   contains
     \LA{}SF circe2: circe2 data: TBP\RA{}
  end type circe2_data_t

\eatline
\nwendcode{}\nwbegincode{1399}\moddef{SF circe2: types}\plusendmoddef
  type(circe2_state) :: circe2_global_state
  
\nwendcode{}\nwbegindocs{1400}\nwdocspar
\nwenddocs{}\nwbegincode{1401}\moddef{SF circe2: circe2 data: TBP}\endmoddef
  procedure :: init => circe2_data_init
\nwendcode{}\nwbegincode{1402}\moddef{SF circe2: procedures}\endmoddef
  subroutine circe2_data_init &
       (data, os_data, model, pdg_in, sqrts, polarized, file, design)
    class(circe2_data_t), intent(out) :: data
    type(os_data_t), intent(in) :: os_data
    class(model_data_t), intent(in), target :: model
    type(pdg_array_t), dimension(2), intent(in) :: pdg_in
    real(default), intent(in) :: sqrts
    logical, intent(in) :: polarized
    type(string_t), intent(in) :: file, design
    integer :: h
    data%model => model
    if (any (pdg_array_get_length (pdg_in) /= 1)) then
       call msg_fatal ("CIRCE2: incoming beam particles must be unique")
    end if
    call data%flv_in(1)%init (pdg_array_get (pdg_in(1), 1), model)
    call data%flv_in(2)%init (pdg_array_get (pdg_in(2), 1), model)
    data%pdg_in = data%flv_in%get_pdg ()
    data%sqrts = sqrts
    data%polarized = polarized
    data%filename = file
    data%design = design
    call data%check_file (os_data)
    call circe2_load (circe2_global_state, trim (char(data%file)), &
            trim (char(data%design)), data%sqrts, data%error)
    data%lumi = circe2_luminosity (circe2_global_state, data%pdg_in, [0, 0])
    call data%check ()
    if (data%polarized) then
       do h = 1, 4
          data%lumi_hel_frac(h) = &
               circe2_luminosity (circe2_global_state, data%pdg_in, &
                                  [data%h1(h), data%h2(h)]) &
               / data%lumi
       end do
    end if
  end subroutine circe2_data_init

\eatline
\nwendcode{}\nwbegindocs{1403}Activate the generator mode.  We import a RNG factory into the data
type, which can then spawn RNG generator objects.
\nwenddocs{}\nwbegincode{1404}\moddef{SF circe2: circe2 data: TBP}\plusendmoddef
  procedure :: set_generator_mode => circe2_data_set_generator_mode
\nwendcode{}\nwbegincode{1405}\moddef{SF circe2: procedures}\plusendmoddef
  subroutine circe2_data_set_generator_mode (data, rng_factory)
    class(circe2_data_t), intent(inout) :: data
    class(rng_factory_t), intent(inout), allocatable :: rng_factory
    call move_alloc (from = rng_factory, to = data%rng_factory)
  end subroutine circe2_data_set_generator_mode
  
\eatline
\nwendcode{}\nwbegindocs{1406}Check whether the requested data file is in the system directory or
in the current directory.
\nwenddocs{}\nwbegincode{1407}\moddef{SF circe2: circe2 data: TBP}\plusendmoddef
  procedure :: check_file => circe2_check_file
\nwendcode{}\nwbegincode{1408}\moddef{SF circe2: procedures}\plusendmoddef
  subroutine circe2_check_file (data, os_data)
    class(circe2_data_t), intent(inout) :: data
    type(os_data_t), intent(in) :: os_data
    logical :: exist
    type(string_t) :: file
    file = data%filename
    if (file == "") &
         call msg_fatal ("CIRCE2: $circe2_file is not set")    
    inquire (file = char (file), exist = exist)
    if (exist) then
       data%file = file
    else
       file = os_data%whizard_circe2path // "/" // data%filename
       inquire (file = char (file), exist = exist)
       if (exist) then
          data%file = file
       else
          call msg_fatal ("CIRCE2: data file '" // char (data%filename) &
               // "' not found")
       end if
    end if
  end subroutine circe2_check_file
    
\eatline
\nwendcode{}\nwbegindocs{1409}Handle error conditions.
\nwenddocs{}\nwbegincode{1410}\moddef{SF circe2: circe2 data: TBP}\plusendmoddef
  procedure :: check => circe2_data_check 
\nwendcode{}\nwbegincode{1411}\moddef{SF circe2: procedures}\plusendmoddef
  subroutine circe2_data_check (data) 
    class(circe2_data_t), intent(in) :: data 
    type(flavor_t) :: flv_photon, flv_electron
    call flv_photon%init (PHOTON, data%model)
    if (.not. flv_photon%is_defined ()) then
       call msg_fatal ("CIRCE2: model must contain photon")
    end if
    call flv_electron%init (ELECTRON, data%model)
    if (.not. flv_electron%is_defined ()) then
       call msg_fatal ("CIRCE2: model must contain electron")
    end if
    if (any (abs (data%pdg_in) /= PHOTON .and. abs (data%pdg_in) /= ELECTRON)) &
         then
       call msg_fatal ("CIRCE2: applicable only for e+e- or photon collisions")
    end if
    select case (data%error)
    case (-1)
       call msg_fatal ("CIRCE2: data file not found.")
    case (-2)
       call msg_fatal ("CIRCE2: beam parameters do not match data file.")
    case (-3)
       call msg_fatal ("CIRCE2: invalid format of data file.")
    case (-4)
       call msg_fatal ("CIRCE2: data file too large.")
    end select
    if (data%lumi == 0) then
       call msg_fatal ("CIRCE2: luminosity vanishes for specified beams.")
    end if
  end subroutine circe2_data_check
  
\eatline
\nwendcode{}\nwbegindocs{1412}Output
\nwenddocs{}\nwbegincode{1413}\moddef{SF circe2: circe2 data: TBP}\plusendmoddef
  procedure :: write => circe2_data_write
\nwendcode{}\nwbegincode{1414}\moddef{SF circe2: procedures}\plusendmoddef
  subroutine circe2_data_write (data, unit, verbose) 
    class(circe2_data_t), intent(in) :: data
    integer, intent(in), optional :: unit
    logical, intent(in), optional :: verbose
    integer :: u, h
    u = given_output_unit (unit)
    write (u, "(1x,A)") "CIRCE2 data:"
    write (u, "(3x,A,A)")       "file   = ", char(data%filename)
    write (u, "(3x,A,A)")       "design = ", char(data%design)
    write (u, "(3x,A," // FMT_19 // ")") "sqrts  = ", data%sqrts
    write (u, "(3x,A,A,A,A)")   "prt_in = ", &
         char (data%flv_in(1)%get_name ()), &
         ", ", char (data%flv_in(2)%get_name ())    
    write (u, "(3x,A,L1)")      "polarized  = ", data%polarized
    write (u, "(3x,A," // FMT_19 // ")") "luminosity = ", data%lumi
    if (data%polarized) then
       do h = 1, 4
          write (u, "(6x,'(',I2,1x,I2,')',1x,'=',1x)", advance="no") &
               data%h1(h), data%h2(h)
          write (u, "(6x, " // FMT_19 // ")")  data%lumi_hel_frac(h)
       end do
    end if
    call data%rng_factory%write (u)
  end subroutine circe2_data_write
  
\eatline
\nwendcode{}\nwbegindocs{1415}This is always in generator mode.
\nwenddocs{}\nwbegincode{1416}\moddef{SF circe2: circe2 data: TBP}\plusendmoddef
  procedure :: is_generator => circe2_data_is_generator
\nwendcode{}\nwbegincode{1417}\moddef{SF circe2: procedures}\plusendmoddef
  function circe2_data_is_generator (data) result (flag)
    class(circe2_data_t), intent(in) :: data
    logical :: flag
    flag = .true.
  end function circe2_data_is_generator
  
\eatline
\nwendcode{}\nwbegindocs{1418}The number of parameters is two, collinear splitting for
the two beams.
\nwenddocs{}\nwbegincode{1419}\moddef{SF circe2: circe2 data: TBP}\plusendmoddef
  procedure :: get_n_par => circe2_data_get_n_par
\nwendcode{}\nwbegincode{1420}\moddef{SF circe2: procedures}\plusendmoddef
  function circe2_data_get_n_par (data) result (n)
    class(circe2_data_t), intent(in) :: data
    integer :: n
    n = 2
  end function circe2_data_get_n_par
  
\eatline
\nwendcode{}\nwbegindocs{1421}Return the outgoing particles PDG codes.  They are equal to the
incoming ones.
\nwenddocs{}\nwbegincode{1422}\moddef{SF circe2: circe2 data: TBP}\plusendmoddef
  procedure :: get_pdg_out => circe2_data_get_pdg_out
\nwendcode{}\nwbegincode{1423}\moddef{SF circe2: procedures}\plusendmoddef
  subroutine circe2_data_get_pdg_out (data, pdg_out)
    class(circe2_data_t), intent(in) :: data
    type(pdg_array_t), dimension(:), intent(inout) :: pdg_out
    integer :: i, n
    n = 2
    do i = 1, n
       pdg_out(i) = data%pdg_in(i)
    end do
  end subroutine circe2_data_get_pdg_out
  
\eatline
\nwendcode{}\nwbegindocs{1424}Allocate the interaction record.
\nwenddocs{}\nwbegincode{1425}\moddef{SF circe2: circe2 data: TBP}\plusendmoddef
  procedure :: allocate_sf_int => circe2_data_allocate_sf_int
\nwendcode{}\nwbegincode{1426}\moddef{SF circe2: procedures}\plusendmoddef
  subroutine circe2_data_allocate_sf_int (data, sf_int)
    class(circe2_data_t), intent(in) :: data
    class(sf_int_t), intent(inout), allocatable :: sf_int
    allocate (circe2_t :: sf_int)
  end subroutine circe2_data_allocate_sf_int
  
\eatline
\nwendcode{}\nwbegindocs{1427}\nwdocspar
\subsection{Random Number Generator for CIRCE}
The CIRCE implementation now supports a generic random-number
generator object that allows for a local state as a component.  To
support this, we must extend the abstract type provided by CIRCE and
delegate the generator call to the (also abstract) RNG used by WHIZARD.
\nwenddocs{}\nwbegincode{1428}\moddef{SF circe2: types}\plusendmoddef
  type, extends (circe2_rng_t) :: rng_obj_t
     class(rng_t), allocatable :: rng
   contains
     procedure :: generate => rng_obj_generate
  end type rng_obj_t
  
\eatline
\nwendcode{}\nwbegincode{1429}\moddef{SF circe2: procedures}\plusendmoddef
  subroutine rng_obj_generate (rng_obj, u)
    class(rng_obj_t), intent(inout) :: rng_obj
    real(default), intent(out) :: u
    real(default) :: x
    call rng_obj%rng%generate (x)
    u = x
  end subroutine rng_obj_generate

\eatline
\nwendcode{}\nwbegindocs{1430}\nwdocspar
\subsection{The CIRCE2 object}
For CIRCE2 spectra it does not make sense to describe the state matrix
as a radiation interaction, even if photons originate from laser
backscattering.  Instead, it is a $2\to 2$ interaction where the
incoming particles are identical to the outgoing ones.

The current implementation of CIRCE2 does support polarization and
classical correlations, but no entanglement, so the density matrix of
the outgoing particles is diagonal.  The incoming particles are
unpolarized (user-defined polarization for beams is meaningless, since
polarization is described by the data file).  The outgoing particles
are polarized or polarization-averaged, depending on user request.

When assigning matrix elements, we scan the previously initialized
state matrix.  For each entry, we extract helicity and call the
structure function.  In the unpolarized case, the helicity is
undefined and replaced by value zero.  In the polarized case, there
are four entries.  If the generator is used, only one entry is nonzero
in each call.  Which one, is determined by comparing with a previously
(randomly, distributed by relative luminosity) selected pair of
helicities.
\nwenddocs{}\nwbegincode{1431}\moddef{SF circe2: types}\plusendmoddef
  type, extends (sf_int_t) :: circe2_t
     type(circe2_data_t), pointer :: data => null ()
     type(rng_obj_t) :: rng_obj
     type(selector_t) :: selector
     integer :: h_sel = 0
   contains
   \LA{}SF circe2: circe2: TBP\RA{}
  end type circe2_t 
  
\eatline
\nwendcode{}\nwbegindocs{1432}Type string: show file and design of {\tt{}CIRCE2} structure function.
\nwenddocs{}\nwbegincode{1433}\moddef{SF circe2: circe2: TBP}\endmoddef
  procedure :: type_string => circe2_type_string
\nwendcode{}\nwbegincode{1434}\moddef{SF circe2: procedures}\plusendmoddef
  function circe2_type_string (object) result (string)
    class(circe2_t), intent(in) :: object
    type(string_t) :: string
    if (associated (object%data)) then
       string = "CIRCE2: " // object%data%design
    else
       string = "CIRCE2: [undefined]"
    end if
  end function circe2_type_string
  
\eatline
\nwendcode{}\nwbegindocs{1435}\nwdocspar
\nwenddocs{}\nwbegindocs{1436}Output.  Call the interaction routine after displaying the configuration.
\nwenddocs{}\nwbegincode{1437}\moddef{SF circe2: circe2: TBP}\plusendmoddef
  procedure :: write => circe2_write
\nwendcode{}\nwbegincode{1438}\moddef{SF circe2: procedures}\plusendmoddef
  subroutine circe2_write (object, unit, testflag)
    class(circe2_t), intent(in) :: object
    integer, intent(in), optional :: unit
    logical, intent(in), optional :: testflag
    integer :: u
    u = given_output_unit (unit)
    if (associated (object%data)) then
       call object%data%write (u)
       call object%base_write (u, testflag)
    else
       write (u, "(1x,A)")  "CIRCE2 data: [undefined]"
    end if
  end subroutine circe2_write
    
\eatline
\nwendcode{}\nwbegindocs{1439}\nwdocspar
\nwenddocs{}\nwbegincode{1440}\moddef{SF circe2: circe2: TBP}\plusendmoddef
  procedure :: init => circe2_init
\nwendcode{}\nwbegincode{1441}\moddef{SF circe2: procedures}\plusendmoddef
  subroutine circe2_init (sf_int, data)
    class(circe2_t), intent(out) :: sf_int
    class(sf_data_t), intent(in), target :: data
    logical, dimension(4) :: mask_h
    real(default), dimension(0) :: null_array
    type(quantum_numbers_mask_t), dimension(4) :: mask
    type(quantum_numbers_t), dimension(4) :: qn
    type(helicity_t) :: hel
    type(color_t) :: col0
    integer :: h
    select type (data)
    type is (circe2_data_t)
       mask_h(1:2) = .true.
       mask_h(3:4) = .not. data%polarized       
       mask = quantum_numbers_mask (.false., .false., mask_h)
       call sf_int%base_init (mask, [0._default, 0._default], &
            null_array, [0._default, 0._default])    
       sf_int%data => data              
       if (data%polarized) then
          call sf_int%selector%init (data%lumi_hel_frac)
       end if
       call col0%init ()
       call qn(1)%init (flv = data%flv_in(1), col = col0)
       call qn(2)%init (flv = data%flv_in(2), col = col0)
       if (data%polarized) then
          do h = 1, 4
             call hel%init (data%h1(h))
             call qn(3)%init &
                  (flv = data%flv_in(1), col = col0, hel = hel)
             call hel%init (data%h2(h))
             call qn(4)%init &
                  (flv = data%flv_in(2), col = col0, hel = hel)
             call interaction_add_state (sf_int%interaction_t, qn)
          end do
       else
          call qn(3)%init (flv = data%flv_in(1), col = col0)
          call qn(4)%init (flv = data%flv_in(2), col = col0)
          call qn(3:4)%tag_radiated ()
          call interaction_add_state (sf_int%interaction_t, qn)
       end if
       call interaction_freeze (sf_int%interaction_t)
       call sf_int%set_incoming ([1,2])
       call sf_int%set_outgoing ([3,4])
       call sf_int%data%rng_factory%make (sf_int%rng_obj%rng)
       sf_int%status = SF_INITIAL
    end select
  end subroutine circe2_init

\eatline
\nwendcode{}\nwbegindocs{1442}\nwdocspar
\subsection{Kinematics}
Refer to the {\tt{}data} component.
\nwenddocs{}\nwbegincode{1443}\moddef{SF circe2: circe2: TBP}\plusendmoddef
  procedure :: is_generator => circe2_is_generator
\nwendcode{}\nwbegincode{1444}\moddef{SF circe2: procedures}\plusendmoddef
  function circe2_is_generator (sf_int) result (flag)
    class(circe2_t), intent(in) :: sf_int
    logical :: flag
    flag = sf_int%data%is_generator ()
  end function circe2_is_generator
  
\eatline
\nwendcode{}\nwbegindocs{1445}Generate free parameters.  We first select a helicity, which we have
to store, then generate $x$ values for that helicity.
\nwenddocs{}\nwbegincode{1446}\moddef{SF circe2: circe2: TBP}\plusendmoddef
  procedure :: generate_free => circe2_generate_whizard_free
\nwendcode{}\nwbegincode{1447}\moddef{SF circe2: procedures}\plusendmoddef
  subroutine circe2_generate_whizard_free (sf_int, r, rb, x_free)
    class(circe2_t), intent(inout) :: sf_int
    real(default), dimension(:), intent(out) :: r, rb
    real(default), intent(inout) :: x_free
    integer :: h_sel
    if (sf_int%data%polarized) then
       call sf_int%selector%generate (sf_int%rng_obj%rng, h_sel)
    else
       h_sel = 0
    end if
    sf_int%h_sel = h_sel
    call circe2_generate_whizard (r, sf_int%data%pdg_in, &
         [sf_int%data%h1(h_sel), sf_int%data%h2(h_sel)], &
         sf_int%rng_obj)
    rb = 1 - r
    x_free = x_free * product (r)
  end subroutine circe2_generate_whizard_free
    
\eatline
\nwendcode{}\nwbegindocs{1448}Generator mode: call the CIRCE2 generator for the given particles
and helicities.  (For unpolarized generation, helicities are zero.)
\nwenddocs{}\nwbegincode{1449}\moddef{SF circe2: procedures}\plusendmoddef
  subroutine circe2_generate_whizard (x, pdg, hel, rng_obj)
    real(default), dimension(2), intent(out) :: x
    integer, dimension(2), intent(in) :: pdg
    integer, dimension(2), intent(in) :: hel
    class(rng_obj_t), intent(inout) :: rng_obj
    call circe2_generate (circe2_global_state, rng_obj, x, pdg, hel)
  end subroutine circe2_generate_whizard

\eatline
\nwendcode{}\nwbegindocs{1450}Set kinematics.  Trivial here.
\nwenddocs{}\nwbegincode{1451}\moddef{SF circe2: circe2: TBP}\plusendmoddef
  procedure :: complete_kinematics => circe2_complete_kinematics
\nwendcode{}\nwbegincode{1452}\moddef{SF circe2: procedures}\plusendmoddef
  subroutine circe2_complete_kinematics (sf_int, x, f, r, rb, map)
    class(circe2_t), intent(inout) :: sf_int
    real(default), dimension(:), intent(out) :: x
    real(default), intent(out) :: f
    real(default), dimension(:), intent(in) :: r
    real(default), dimension(:), intent(in) :: rb
    logical, intent(in) :: map
    if (map) then
       call msg_fatal ("CIRCE2: map flag not supported")
    else
       x = r
       f = 1
    end if
    call sf_int%reduce_momenta (x)
  end subroutine circe2_complete_kinematics

\eatline
\nwendcode{}\nwbegindocs{1453}Compute inverse kinematics.
\nwenddocs{}\nwbegincode{1454}\moddef{SF circe2: circe2: TBP}\plusendmoddef
  procedure :: inverse_kinematics => circe2_inverse_kinematics
\nwendcode{}\nwbegincode{1455}\moddef{SF circe2: procedures}\plusendmoddef
  subroutine circe2_inverse_kinematics (sf_int, x, f, r, rb, map, set_momenta)
    class(circe2_t), intent(inout) :: sf_int
    real(default), dimension(:), intent(in) :: x
    real(default), intent(out) :: f
    real(default), dimension(:), intent(out) :: r
    real(default), dimension(:), intent(out) :: rb
    logical, intent(in) :: map
    logical, intent(in), optional :: set_momenta
    logical :: set_mom
    set_mom = .false.;  if (present (set_momenta))  set_mom = set_momenta
    if (map) then
       call msg_fatal ("CIRCE2: map flag not supported")
    else
       r = x
       rb= 1 - r
       f = 1
    end if
    if (set_mom) then
       call sf_int%reduce_momenta (x)
    end if
  end subroutine circe2_inverse_kinematics

\eatline
\nwendcode{}\nwbegindocs{1456}\nwdocspar
\subsection{CIRCE2 application}
This function works on both beams.  In polarized mode, we set only the
selected helicity.  In unpolarized mode,
the interaction has only one entry, and the factor is unity.
\nwenddocs{}\nwbegincode{1457}\moddef{SF circe2: circe2: TBP}\plusendmoddef
  procedure :: apply => circe2_apply
\nwendcode{}\nwbegincode{1458}\moddef{SF circe2: procedures}\plusendmoddef
  subroutine circe2_apply (sf_int, scale)
    class(circe2_t), intent(inout) :: sf_int
    real(default), intent(in) :: scale    
    complex(default) :: f
    integer :: h
    associate (data => sf_int%data)
      f = 1
      if (data%polarized) then
         h = sf_int%h_sel
      else
         h = 1
      end if
      call interaction_set_matrix_element (sf_int%interaction_t, h, f)
    end associate
    sf_int%status = SF_EVALUATED
  end subroutine circe2_apply

\eatline
\nwendcode{}\nwbegindocs{1459}\nwdocspar
\subsection{Unit tests}
\nwenddocs{}\nwbegincode{1460}\moddef{SF circe2: public}\plusendmoddef
  public :: sf_circe2_test
\nwendcode{}\nwbegincode{1461}\moddef{SF circe2: tests}\endmoddef
  subroutine sf_circe2_test (u, results)
    integer, intent(in) :: u
    type(test_results_t), intent(inout) :: results
  \LA{}SF circe2: execute tests\RA{}
  end subroutine sf_circe2_test
  
\eatline
\nwendcode{}\nwbegindocs{1462}\nwdocspar
\subsubsection{Test structure function data}
Construct and display a test structure function data object.
\nwenddocs{}\nwbegincode{1463}\moddef{SF circe2: execute tests}\endmoddef
  call test (sf_circe2_1, "sf_circe2_1", &
       "structure function configuration", &
       u, results)
\nwendcode{}\nwbegincode{1464}\moddef{SF circe2: tests}\plusendmoddef
  subroutine sf_circe2_1 (u)
    integer, intent(in) :: u
    type(os_data_t) :: os_data
    type(model_data_t), target :: model
    type(pdg_array_t), dimension(2) :: pdg_in
    type(pdg_array_t), dimension(2) :: pdg_out
    integer, dimension(:), allocatable :: pdg1, pdg2
    class(sf_data_t), allocatable :: data
    class(rng_factory_t), allocatable :: rng_factory
    
    write (u, "(A)")  "* Test output: sf_circe2_1"
    write (u, "(A)")  "*   Purpose: initialize and display &
         &CIRCE structure function data"
    write (u, "(A)")
    
    write (u, "(A)")  "* Create empty data object"
    write (u, "(A)")

    call os_data_init (os_data)
    call model%init_qed_test ()
    pdg_in(1) = PHOTON
    pdg_in(2) = PHOTON

    allocate (circe2_data_t :: data)
    allocate (rng_test_factory_t :: rng_factory)

    write (u, "(A)")
    write (u, "(A)")  "* Initialize (unpolarized)"
    write (u, "(A)")

    select type (data)
    type is (circe2_data_t)
       call data%init (os_data, model, pdg_in, &
            sqrts = 500._default, &
            polarized = .false., &
            file = var_str ("teslagg_500_polavg.circe"), &
            design = var_str ("TESLA/GG"))
       call data%set_generator_mode (rng_factory)
    end select

    call data%write (u)

    write (u, "(A)")

    write (u, "(1x,A)")  "Outgoing particle codes:"
    call data%get_pdg_out (pdg_out)
    pdg1 = pdg_out(1)
    pdg2 = pdg_out(2)
    write (u, "(2x,99(1x,I0))")  pdg1, pdg2

    write (u, "(A)")
    write (u, "(A)")  "* Initialize (polarized)"
    write (u, "(A)")

    allocate (rng_test_factory_t :: rng_factory)

    select type (data)
    type is (circe2_data_t)
       call data%init (os_data, model, pdg_in, &
            sqrts = 500._default, &
            polarized = .true., &
            file = var_str ("teslagg_500.circe"), &
            design = var_str ("TESLA/GG"))
       call data%set_generator_mode (rng_factory)
    end select

    call data%write (u)

    call model%final ()
    
    write (u, "(A)")
    write (u, "(A)")  "* Test output end: sf_circe2_1"

  end subroutine sf_circe2_1

\eatline
\nwendcode{}\nwbegindocs{1465}\nwdocspar
\subsubsection{Generator mode, unpolarized}
Construct and evaluate a structure function object in generator mode.
\nwenddocs{}\nwbegincode{1466}\moddef{SF circe2: execute tests}\plusendmoddef
  call test (sf_circe2_2, "sf_circe2_2", &
       "generator, unpolarized", &
       u, results)
\nwendcode{}\nwbegincode{1467}\moddef{SF circe2: tests}\plusendmoddef
  subroutine sf_circe2_2 (u)
    integer, intent(in) :: u
    type(os_data_t) :: os_data
    type(model_data_t), target :: model
    type(flavor_t), dimension(2) :: flv
    type(pdg_array_t), dimension(2) :: pdg_in
    class(sf_data_t), allocatable, target :: data
    class(rng_factory_t), allocatable :: rng_factory
    class(sf_int_t), allocatable :: sf_int
    type(vector4_t) :: k1, k2
    real(default) :: E
    real(default), dimension(:), allocatable :: r, rb, x
    real(default) :: f, x_free
    
    write (u, "(A)")  "* Test output: sf_circe2_2"
    write (u, "(A)")  "*   Purpose: initialize and fill &
         &circe2 structure function object"
    write (u, "(A)")
    
    write (u, "(A)")  "* Initialize configuration data"
    write (u, "(A)")

    call os_data_init (os_data)
    call model%init_qed_test ()
    call flv(1)%init (PHOTON, model)
    call flv(2)%init (PHOTON, model)
    pdg_in(1) = PHOTON
    pdg_in(2) = PHOTON

    call reset_interaction_counter ()
    
    allocate (circe2_data_t :: data)
    allocate (rng_test_factory_t :: rng_factory)
    select type (data)
    type is (circe2_data_t)
       call data%init (os_data, model, pdg_in, &
            sqrts = 500._default, &
            polarized = .false., &
            file = var_str ("teslagg_500_polavg.circe"), &
            design = var_str ("TESLA/GG"))
       call data%set_generator_mode (rng_factory)
    end select
       
    write (u, "(A)")  "* Initialize structure-function object"
    write (u, "(A)")
    
    call data%allocate_sf_int (sf_int)
    call sf_int%init (data)
    call sf_int%set_beam_index ([1,2])
    select type (sf_int)
    type is (circe2_t)
       call sf_int%rng_obj%rng%init (3)
    end select

    write (u, "(A)")  "* Initialize incoming momentum with E=500"
    write (u, "(A)")
    E = 250
    k1 = vector4_moving (E, sqrt (E**2 - flv(1)%get_mass ()**2), 3)
    k2 = vector4_moving (E,-sqrt (E**2 - flv(2)%get_mass ()**2), 3)
    call vector4_write (k1, u)
    call vector4_write (k2, u)
    call sf_int%seed_kinematics ([k1, k2])

    write (u, "(A)")
    write (u, "(A)")  "* Generate x"
    write (u, "(A)")

    allocate (r (data%get_n_par ()))
    allocate (rb(size (r)))
    allocate (x (size (r)))

    r  = 0
    rb = 0
    x_free = 1
    call sf_int%generate_free (r, rb, x_free)
    call sf_int%complete_kinematics (x, f, r, rb, map=.false.)

    write (u, "(A,9(1x,F10.7))")  "x =", x
    write (u, "(A,9(1x,F10.7))")  "f =", f
    write (u, "(A,9(1x,F10.7))")  "xf=", x_free

    write (u, "(A)")
    write (u, "(A)")  "* Evaluate"
    write (u, "(A)")

    call sf_int%apply (scale = 0._default)
    call sf_int%write (u)

    write (u, "(A)")
    write (u, "(A)")  "* Cleanup"

    call sf_int%final ()
    call model%final ()
    
    write (u, "(A)")
    write (u, "(A)")  "* Test output end: sf_circe2_2"

  end subroutine sf_circe2_2

\eatline
\nwendcode{}\nwbegindocs{1468}\nwdocspar
\subsubsection{Generator mode, polarized}
Construct and evaluate a structure function object in generator mode.
\nwenddocs{}\nwbegincode{1469}\moddef{SF circe2: execute tests}\plusendmoddef
  call test (sf_circe2_3, "sf_circe2_3", &
       "generator, polarized", &
       u, results)
\nwendcode{}\nwbegincode{1470}\moddef{SF circe2: tests}\plusendmoddef
  subroutine sf_circe2_3 (u)
    integer, intent(in) :: u
    type(os_data_t) :: os_data
    type(model_data_t), target :: model
    type(flavor_t), dimension(2) :: flv
    type(pdg_array_t), dimension(2) :: pdg_in
    class(sf_data_t), allocatable, target :: data
    class(rng_factory_t), allocatable :: rng_factory
    class(sf_int_t), allocatable :: sf_int
    type(vector4_t) :: k1, k2
    real(default) :: E
    real(default), dimension(:), allocatable :: r, rb, x
    real(default) :: f, x_free
    
    write (u, "(A)")  "* Test output: sf_circe2_3"
    write (u, "(A)")  "*   Purpose: initialize and fill &
         &circe2 structure function object"
    write (u, "(A)")
    
    write (u, "(A)")  "* Initialize configuration data"
    write (u, "(A)")

    call os_data_init (os_data)
    call model%init_qed_test ()
    call flv(1)%init (PHOTON, model)
    call flv(2)%init (PHOTON, model)
    pdg_in(1) = PHOTON
    pdg_in(2) = PHOTON

    call reset_interaction_counter ()
    
    allocate (circe2_data_t :: data)
    allocate (rng_test_factory_t :: rng_factory)
    select type (data)
    type is (circe2_data_t)
       call data%init (os_data, model, pdg_in, &
            sqrts = 500._default, &
            polarized = .true., &
            file = var_str ("teslagg_500.circe"), &
            design = var_str ("TESLA/GG"))
       call data%set_generator_mode (rng_factory)
    end select
       
    write (u, "(A)")  "* Initialize structure-function object"
    write (u, "(A)")
    
    call data%allocate_sf_int (sf_int)
    call sf_int%init (data)
    call sf_int%set_beam_index ([1,2])
    select type (sf_int)
    type is (circe2_t)
       call sf_int%rng_obj%rng%init (3)
    end select

    write (u, "(A)")  "* Initialize incoming momentum with E=500"
    write (u, "(A)")
    E = 250
    k1 = vector4_moving (E, sqrt (E**2 - flv(1)%get_mass ()**2), 3)
    k2 = vector4_moving (E,-sqrt (E**2 - flv(2)%get_mass ()**2), 3)
    call vector4_write (k1, u)
    call vector4_write (k2, u)
    call sf_int%seed_kinematics ([k1, k2])

    write (u, "(A)")
    write (u, "(A)")  "* Generate x"
    write (u, "(A)")

    allocate (r (data%get_n_par ()))
    allocate (rb(size (r)))
    allocate (x (size (r)))

    r  = 0
    rb = 0
    x_free = 1
    call sf_int%generate_free (r, rb, x_free)
    call sf_int%complete_kinematics (x, f, r, rb, map=.false.)

    write (u, "(A,9(1x,F10.7))")  "x =", x
    write (u, "(A,9(1x,F10.7))")  "f =", f
    write (u, "(A,9(1x,F10.7))")  "xf=", x_free

    write (u, "(A)")
    write (u, "(A)")  "* Evaluate"
    write (u, "(A)")

    call sf_int%apply (scale = 0._default)
    call sf_int%write (u)

    write (u, "(A)")
    write (u, "(A)")  "* Cleanup"

    call sf_int%final ()
    call model%final ()
    
    write (u, "(A)")
    write (u, "(A)")  "* Test output end: sf_circe2_3"

  end subroutine sf_circe2_3

\eatline
\nwendcode{}\nwbegindocs{1471}\nwdocspar
\clearpage
%------------------------------------------------------------------------
\section{HOPPET interface}
Interface to the HOPPET wrapper necessary to perform
the LO vs. NLO matching of processes containing an initial
b quark.
\nwenddocs{}\nwbegincode{1472}\moddef{\code{}hoppet{\char95}interface.f90\edoc{}}\endmoddef
\LA{}File header\RA{}

module hoppet_interface
  use lhapdf !NODEP!     
  
\LA{}Standard module head\RA{}

  public :: hoppet_init, hoppet_eval

contains

  subroutine hoppet_init (pdf_builtin, pdf, pdf_id)
    logical, intent(in) :: pdf_builtin
    type(lhapdf_pdf_t), intent(inout), optional :: pdf
    integer, intent(in), optional :: pdf_id
    external InitForWhizard
    call InitForWhizard (pdf_builtin, pdf, pdf_id)
  end subroutine hoppet_init

  subroutine hoppet_eval (x, q, f)
    double precision, intent(in)  :: x, q
    double precision, intent(out) :: f(-6:6)
    external EvalForWhizard
    call EvalForWhizard (x, q, f)
  end subroutine hoppet_eval

end module hoppet_interface
\eatline
\nwendcode{}\nwbegindocs{1473}\nwdocspar
\clearpage
%------------------------------------------------------------------------
\section{Builtin PDF sets}
For convenience in order not to depend on the external package LHAPDF,
we ship some PDFs with WHIZARD. 
\nwenddocs{}\nwbegindocs{1474}\nwdocspar
\subsection{The module}
\nwenddocs{}\nwbegincode{1475}\moddef{\code{}sf{\char95}pdf{\char95}builtin.f90\edoc{}}\endmoddef
\LA{}File header\RA{}

module sf_pdf_builtin

\LA{}Use kinds\RA{}
  use kinds, only: double
\LA{}Use strings\RA{}
  use io_units
  use format_defs, only: FMT_17
  use unit_tests
  use diagnostics
  use os_interface
  use physics_defs, only: PROTON, PHOTON, GLUON
  use physics_defs, only: HADRON_REMNANT_SINGLET
  use physics_defs, only: HADRON_REMNANT_TRIPLET
  use physics_defs, only: HADRON_REMNANT_OCTET
  use sm_qcd
  use lorentz
  use pdg_arrays
  use model_data
  use flavors
  use colors
  use quantum_numbers
  use state_matrices
  use polarizations
  use interactions
  use sf_aux
  use sf_base
  use pdf_builtin !NODEP!
  use hoppet_interface

\LA{}Standard module head\RA{}

\LA{}SF pdf builtin: public\RA{}

\LA{}SF pdf builtin: types\RA{}

\LA{}SF pdf builtin: parameters\RA{}

contains

\LA{}SF pdf builtin: procedures\RA{}

\LA{}SF pdf builtin: tests\RA{}

end module sf_pdf_builtin
\eatline
\nwendcode{}\nwbegindocs{1476}\nwdocspar
\subsection{Codes for default PDF sets} 
\nwenddocs{}\nwbegincode{1477}\moddef{SF pdf builtin: parameters}\endmoddef
  character(*), parameter :: PDF_BUILTIN_DEFAULT_PROTON = "CTEQ6L"
  character(*), parameter :: PDF_BUILTIN_DEFAULT_PION   = "NONE"
  character(*), parameter :: PDF_BUILTIN_DEFAULT_PHOTON = "MRST2004QEDp"
\eatline
\nwendcode{}\nwbegindocs{1478}\nwdocspar
\subsection{The PDF builtin data block} 
The data block holds the incoming flavor (which has to be proton,
pion, or photon), the corresponding pointer to the global access data
(1, 2, or 3), the flag {\tt{}invert} which is set for an antiproton, the
bounds as returned by LHAPDF for the specified set, and a mask that
determines which partons will be actually in use.
\nwenddocs{}\nwbegincode{1479}\moddef{SF pdf builtin: public}\endmoddef
  public :: pdf_builtin_data_t
\nwendcode{}\nwbegincode{1480}\moddef{SF pdf builtin: types}\endmoddef
  type, extends (sf_data_t) :: pdf_builtin_data_t
     private
     integer :: id = -1
     type (string_t) :: name
     class(model_data_t), pointer :: model => null ()
     type(flavor_t) :: flv_in
     logical :: invert
     logical :: has_photon
     logical :: photon
     logical, dimension(-6:6) :: mask
     logical :: mask_photon
     logical :: hoppet_b_matching = .false.
   contains
   \LA{}SF pdf builtin: pdf builtin data: TBP\RA{}
  end type pdf_builtin_data_t

\eatline
\nwendcode{}\nwbegindocs{1481}Generate PDF data and initialize the requested set. Pion and photon PDFs
are disabled at the moment until we ship appropiate structure functions.
needed.
\nwenddocs{}\nwbegincode{1482}\moddef{SF pdf builtin: pdf builtin data: TBP}\endmoddef
  procedure :: init => pdf_builtin_data_init
\nwendcode{}\nwbegincode{1483}\moddef{SF pdf builtin: procedures}\endmoddef
  subroutine pdf_builtin_data_init (data, &
       model, pdg_in, name, path, hoppet_b_matching)
    class(pdf_builtin_data_t), intent(out) :: data
    class(model_data_t), intent(in), target :: model
    type(pdg_array_t), intent(in) :: pdg_in
    type(string_t), intent(in) :: name
    type(string_t), intent(in) :: path
    logical, intent(in), optional :: hoppet_b_matching
    data%model => model
    if (pdg_array_get_length (pdg_in) /= 1) &
         call msg_fatal ("PDF: incoming particle must be unique")
    call data%flv_in%init (pdg_array_get (pdg_in, 1), model)
    data%mask = .true.
    data%mask_photon = .true.
    select case (pdg_array_get (pdg_in, 1))
    case (PROTON)
       data%name = var_str (PDF_BUILTIN_DEFAULT_PROTON)
       data%invert = .false.
       data%photon = .false.
    case (-PROTON)
       data%name = var_str (PDF_BUILTIN_DEFAULT_PROTON)
       data%invert = .true.
       data%photon = .false.
       ! case (PIPLUS)
       !    data%name = var_str (PDF_BUILTIN_DEFAULT_PION)
       !    data%invert = .false.
       !    data%photon = .false.
       ! case (-PIPLUS)
       !    data%name = var_str (PDF_BUILTIN_DEFAULT_PION)
       !    data%invert = .true.
       !    data%photon = .false.
       ! case (PHOTON)
       !    data%name = var_str (PDF_BUILTIN_DEFAULT_PHOTON)
       !    data%invert = .false.
       !    data%photon = .true.
    case default
       call msg_fatal ("PDF: " &
            // "incoming particle must either proton or antiproton.")
       return
    end select
    data%name = name
    data%id = pdf_get_id (data%name)
    if (data%id < 0) call msg_fatal ("unknown PDF set " // char (data%name))
    data%has_photon = pdf_provides_photon (data%id)
    if (present (hoppet_b_matching))  data%hoppet_b_matching = hoppet_b_matching
    call pdf_init (data%id, path)
    if (data%hoppet_b_matching)  call hoppet_init (.true., pdf_id = data%id)
  end subroutine pdf_builtin_data_init

\eatline
\nwendcode{}\nwbegindocs{1484}Enable/disable partons explicitly.  If a mask entry is true,
applying the PDF will generate the corresponding flavor on output.
\nwenddocs{}\nwbegincode{1485}\moddef{SF pdf builtin: pdf builtin data: TBP}\plusendmoddef
  procedure :: set_mask => pdf_builtin_data_set_mask
\nwendcode{}\nwbegincode{1486}\moddef{SF pdf builtin: procedures}\plusendmoddef
  subroutine pdf_builtin_data_set_mask (data, mask)
    class(pdf_builtin_data_t), intent(inout) :: data
    logical, dimension(-6:6), intent(in) :: mask
    data%mask = mask
  end subroutine pdf_builtin_data_set_mask

\eatline
\nwendcode{}\nwbegindocs{1487}Output.
\nwenddocs{}\nwbegincode{1488}\moddef{SF pdf builtin: pdf builtin data: TBP}\plusendmoddef
  procedure :: write => pdf_builtin_data_write
\nwendcode{}\nwbegincode{1489}\moddef{SF pdf builtin: procedures}\plusendmoddef
  subroutine pdf_builtin_data_write (data, unit, verbose)
    class(pdf_builtin_data_t), intent(in) :: data
    integer, intent(in), optional :: unit
    logical, intent(in), optional :: verbose    
    integer :: u
    u = given_output_unit (unit);  if (u < 0)  return
    write (u, "(1x,A)")  "PDF builtin data:"
    if (data%id < 0) then
       write (u, "(3x,A)") "[undefined]"
       return
    end if
    write (u, "(3x,A)", advance="no") "flavor       = "
    call data%flv_in%write (u);  write (u, *)
    write (u, "(3x,A,A)")   "name         = ", char (data%name)
    write (u, "(3x,A,L1)")  "invert       = ", data%invert
    write (u, "(3x,A,L1)")  "has photon   = ", data%has_photon
    write (u, "(3x,A,6(1x,L1),1x,A,1x,L1,1x,A,6(1x,L1))") &
         "mask         =", &
         data%mask(-6:-1), "*", data%mask(0), "*", data%mask(1:6)
    write (u, "(3x,A,L1)") "photon mask  = ", data%mask_photon
    write (u, "(3x,A,L1)") "hoppet_b     = ", data%hoppet_b_matching    
  end subroutine pdf_builtin_data_write

\eatline
\nwendcode{}\nwbegindocs{1490}The number of parameters is one.  We do not generate transverse momentum.
\nwenddocs{}\nwbegincode{1491}\moddef{SF pdf builtin: pdf builtin data: TBP}\plusendmoddef
  procedure :: get_n_par => pdf_builtin_data_get_n_par
\nwendcode{}\nwbegincode{1492}\moddef{SF pdf builtin: procedures}\plusendmoddef
  function pdf_builtin_data_get_n_par (data) result (n)
    class(pdf_builtin_data_t), intent(in) :: data
    integer :: n
    n = 1
  end function pdf_builtin_data_get_n_par
  
\eatline
\nwendcode{}\nwbegindocs{1493}Return the outgoing particle PDG codes.  This is based on the mask.
\nwenddocs{}\nwbegincode{1494}\moddef{SF pdf builtin: pdf builtin data: TBP}\plusendmoddef
  procedure :: get_pdg_out => pdf_builtin_data_get_pdg_out
\nwendcode{}\nwbegincode{1495}\moddef{SF pdf builtin: procedures}\plusendmoddef
  subroutine pdf_builtin_data_get_pdg_out (data, pdg_out)
    class(pdf_builtin_data_t), intent(in) :: data
    type(pdg_array_t), dimension(:), intent(inout) :: pdg_out
    integer, dimension(:), allocatable :: pdg1
    integer :: n, np, i
    n = count (data%mask)
    np = 0;  if (data%has_photon .and. data%mask_photon)  np = 1
    allocate (pdg1 (n + np))
    pdg1(1:n) = pack ([(i, i = -6, 6)], data%mask)
    if (np == 1)  pdg1(n+np) = PHOTON
    pdg_out(1) = pdg1
  end subroutine pdf_builtin_data_get_pdg_out
  
\eatline
\nwendcode{}\nwbegindocs{1496}Allocate the interaction record.
\nwenddocs{}\nwbegincode{1497}\moddef{SF pdf builtin: pdf builtin data: TBP}\plusendmoddef
  procedure :: allocate_sf_int => pdf_builtin_data_allocate_sf_int
\nwendcode{}\nwbegincode{1498}\moddef{SF pdf builtin: procedures}\plusendmoddef
  subroutine pdf_builtin_data_allocate_sf_int (data, sf_int)
    class(pdf_builtin_data_t), intent(in) :: data
    class(sf_int_t), intent(inout), allocatable :: sf_int
    allocate (pdf_builtin_t :: sf_int)
  end subroutine pdf_builtin_data_allocate_sf_int
  
\eatline
\nwendcode{}\nwbegindocs{1499}Return the numerical PDF set index.
\nwenddocs{}\nwbegincode{1500}\moddef{SF pdf builtin: pdf builtin data: TBP}\plusendmoddef
  procedure :: get_pdf_set => pdf_builtin_data_get_pdf_set
\nwendcode{}\nwbegincode{1501}\moddef{SF pdf builtin: procedures}\plusendmoddef
  function pdf_builtin_data_get_pdf_set (data) result (pdf_set)
    class(pdf_builtin_data_t), intent(in) :: data
    integer :: pdf_set
    pdf_set = data%id
  end function pdf_builtin_data_get_pdf_set
  
\eatline
\nwendcode{}\nwbegindocs{1502}\nwdocspar
\subsection{The PDF object}
The PDF $1\to 2$ interaction which describes
the splitting of an (anti)proton into a parton and a beam remnant.  We
stay in the strict forward-splitting limit, but allow some invariant
mass for the beam remnant such that the outgoing parton is exactly
massless.  For a real event, we would replace this by a parton
cascade, where the outgoing partons have virtuality as dictated by
parton-shower kinematics, and transverse momentum is generated.

The PDF application is a $1\to 2$ splitting process, where the
particles are ordered as (hadron, remnant, parton).

Polarization is ignored completely.  The beam particle is colorless,
while partons and beam remnant carry color.  The remnant gets a
special flavor code.
\nwenddocs{}\nwbegincode{1503}\moddef{SF pdf builtin: types}\plusendmoddef
  type, extends (sf_int_t) :: pdf_builtin_t
     type(pdf_builtin_data_t), pointer :: data => null ()
     real(default) :: x = 0
     real(default) :: q = 0
   contains
   \LA{}SF pdf builtin: pdf builtin: TBP\RA{}
  end type pdf_builtin_t
  
\eatline
\nwendcode{}\nwbegindocs{1504}Type string: display the chosen PDF set.
\nwenddocs{}\nwbegincode{1505}\moddef{SF pdf builtin: pdf builtin: TBP}\endmoddef
  procedure :: type_string => pdf_builtin_type_string
\nwendcode{}\nwbegincode{1506}\moddef{SF pdf builtin: procedures}\plusendmoddef
  function pdf_builtin_type_string (object) result (string)
    class(pdf_builtin_t), intent(in) :: object
    type(string_t) :: string
    if (associated (object%data)) then
       string = "PDF builtin: " // object%data%name
    else
       string = "PDF builtin: [undefined]"
    end if
  end function pdf_builtin_type_string
  
\eatline
\nwendcode{}\nwbegindocs{1507}Output.  Call the interaction routine after displaying the configuration.
\nwenddocs{}\nwbegincode{1508}\moddef{SF pdf builtin: pdf builtin: TBP}\plusendmoddef
  procedure :: write => pdf_builtin_write
\nwendcode{}\nwbegincode{1509}\moddef{SF pdf builtin: procedures}\plusendmoddef
  subroutine pdf_builtin_write (object, unit, testflag)
    class(pdf_builtin_t), intent(in) :: object
    integer, intent(in), optional :: unit
    logical, intent(in), optional :: testflag
    integer :: u
    u = given_output_unit (unit)
    if (associated (object%data)) then
       call object%data%write (u)
       if (object%status >= SF_DONE_KINEMATICS) then
          write (u, "(1x,A)")  "SF parameters:"
          write (u, "(3x,A," // FMT_17 // ")")  "x =", object%x
          if (object%status >= SF_FAILED_EVALUATION) then
             write (u, "(3x,A," // FMT_17 // ")")  "Q =", object%q
          end if
       end if
       call object%base_write (u, testflag)
    else
       write (u, "(1x,A)")  "PDF builtin data: [undefined]"
    end if
  end subroutine pdf_builtin_write
    
\eatline
\nwendcode{}\nwbegindocs{1510}Initialize.  We know that {\tt{}data} will be of concrete type
{\tt{}sf{\char95}test{\char95}data{\char95}t}, but we have to cast this explicitly.

For this implementation, we set the incoming and outgoing masses equal
to the physical particle mass, but keep the radiated mass zero.

Optionally, we can provide minimum and maximum values for the momentum
transfer.
\nwenddocs{}\nwbegincode{1511}\moddef{SF pdf builtin: pdf builtin: TBP}\plusendmoddef
  procedure :: init => pdf_builtin_init
\nwendcode{}\nwbegincode{1512}\moddef{SF pdf builtin: procedures}\plusendmoddef
  subroutine pdf_builtin_init (sf_int, data)
    class(pdf_builtin_t), intent(out) :: sf_int
    class(sf_data_t), intent(in), target :: data
    type(quantum_numbers_mask_t), dimension(3) :: mask
    type(flavor_t) :: flv, flv_remnant
    type(color_t) :: col0
    type(quantum_numbers_t), dimension(3) :: qn
    integer :: i
    select type (data)
    type is (pdf_builtin_data_t)
       mask = quantum_numbers_mask (.false., .false., .true.)
       call col0%init ()
       call sf_int%base_init (mask, [0._default], [0._default], [0._default])
       sf_int%data => data
       do i = -6, 6
          if (data%mask(i)) then
             call qn(1)%init (data%flv_in, col = col0)
             if (i == 0) then
                call flv%init (GLUON, data%model)
                call flv_remnant%init (HADRON_REMNANT_OCTET, data%model)
             else
                call flv%init (i, data%model)
                call flv_remnant%init &
                     (sign (HADRON_REMNANT_TRIPLET, -i), data%model)
             end if
             call qn(2)%init ( &
                  flv = flv_remnant, col = color_from_flavor (flv_remnant, 1))
             call qn(2)%tag_radiated ()
             call qn(3)%init ( &
                  flv = flv, col = color_from_flavor (flv, 1, reverse=.true.))
             call interaction_add_state (sf_int%interaction_t, qn)
          end if
       end do
       if (data%has_photon .and. data%mask_photon) then
          call flv%init (PHOTON, data%model)
          call flv_remnant%init (HADRON_REMNANT_SINGLET, data%model)
          call qn(2)%init (flv = flv_remnant, &
               col = color_from_flavor (flv_remnant, 1))
          call qn(2)%tag_radiated ()
          call qn(3)%init (flv = flv, &
               col = color_from_flavor (flv, 1, reverse = .true.))
          call interaction_add_state (sf_int%interaction_t, qn)
       end if
       call interaction_freeze (sf_int%interaction_t)
       call sf_int%set_incoming ([1])
       call sf_int%set_radiated ([2])
       call sf_int%set_outgoing ([3])
       sf_int%status = SF_INITIAL
    end select
  end subroutine pdf_builtin_init

\eatline
\nwendcode{}\nwbegindocs{1513}\nwdocspar
\subsection{Kinematics}
Set kinematics.  If {\tt{}map} is unset, the $r$ and $x$ values
coincide, and the Jacobian $f(r)$ is trivial.

If {\tt{}map} is set, we are asked to provide an efficient mapping.
For the test case, we set $x=r^2$ and consequently $f(r)=2r$.
\nwenddocs{}\nwbegincode{1514}\moddef{SF pdf builtin: pdf builtin: TBP}\plusendmoddef
  procedure :: complete_kinematics => pdf_builtin_complete_kinematics
\nwendcode{}\nwbegincode{1515}\moddef{SF pdf builtin: procedures}\plusendmoddef
  subroutine pdf_builtin_complete_kinematics (sf_int, x, f, r, rb, map)
    class(pdf_builtin_t), intent(inout) :: sf_int
    real(default), dimension(:), intent(out) :: x
    real(default), intent(out) :: f
    real(default), dimension(:), intent(in) :: r
    real(default), dimension(:), intent(in) :: rb
    logical, intent(in) :: map
    real(default) :: xb1
    if (map) then
       call msg_fatal ("PDF builtin: map flag not supported")
    else
       x(1) = r(1)
       f = 1
    end if
    xb1 = 1 - x(1)
    call sf_int%split_momentum (x, xb1)
    select case (sf_int%status)
    case (SF_DONE_KINEMATICS)
       sf_int%x = x(1)
    case (SF_FAILED_KINEMATICS)
       sf_int%x = 0
       f = 0
    end select
  end subroutine pdf_builtin_complete_kinematics

\eatline
\nwendcode{}\nwbegindocs{1516}Compute inverse kinematics.  Here, we start with the $x$ array and
compute the ``input'' $r$ values and the Jacobian $f$.  After this, we
can set momenta by the same formula as for normal kinematics.
\nwenddocs{}\nwbegincode{1517}\moddef{SF pdf builtin: pdf builtin: TBP}\plusendmoddef
  procedure :: inverse_kinematics => pdf_builtin_inverse_kinematics
\nwendcode{}\nwbegincode{1518}\moddef{SF pdf builtin: procedures}\plusendmoddef
  subroutine pdf_builtin_inverse_kinematics (sf_int, x, f, r, rb, map, set_momenta)
    class(pdf_builtin_t), intent(inout) :: sf_int
    real(default), dimension(:), intent(in) :: x
    real(default), intent(out) :: f
    real(default), dimension(:), intent(out) :: r
    real(default), dimension(:), intent(out) :: rb
    logical, intent(in) :: map
    logical, intent(in), optional :: set_momenta
    real(default) :: xb1
    logical :: set_mom
    set_mom = .false.;  if (present (set_momenta))  set_mom = set_momenta
    if (map) then
       call msg_fatal ("PDF builtin: map flag not supported")
    else
       r(1) = x(1)
       f = 1
    end if
    xb1 = 1 - x(1)
    rb = 1 - r
    if (set_mom) then
       call sf_int%split_momentum (x, xb1)
       select case (sf_int%status)
       case (SF_DONE_KINEMATICS)
          sf_int%x = x(1)
       case (SF_FAILED_KINEMATICS)
          sf_int%x = 0
          f = 0
       end select
    end if
  end subroutine pdf_builtin_inverse_kinematics

\eatline
\nwendcode{}\nwbegindocs{1519}\nwdocspar
\subsection{Structure function}
Once the scale is also known, we can actually call the PDF and
set the values.  Contrary to LHAPDF, the wrapper already takes care of
adjusting to the $x$ and $Q$ bounds.  Account for the Jacobian.
\nwenddocs{}\nwbegincode{1520}\moddef{SF pdf builtin: pdf builtin: TBP}\plusendmoddef
  procedure :: apply => pdf_builtin_apply
\nwendcode{}\nwbegincode{1521}\moddef{SF pdf builtin: procedures}\plusendmoddef
  subroutine pdf_builtin_apply (sf_int, scale)
    class(pdf_builtin_t), intent(inout) :: sf_int
    real(default), intent(in) :: scale
    real(default), dimension(-6:6) :: ff
    real(double), dimension(-6:6) :: ff_dbl
    real(default) :: x, fph
    real(double) :: xx, qq
    complex(default), dimension(:), allocatable :: fc
    associate (data => sf_int%data)
      sf_int%q = scale
      x = sf_int%x
      xx = x
      qq = scale
      if (data%invert) then
         if (data%has_photon) then
            call pdf_evolve (data%id, x, scale, ff(6:-6:-1), fph)
         else
            if (data%hoppet_b_matching) then
               call hoppet_eval (xx, qq, ff_dbl(6:-6:-1))
               ff = ff_dbl
            else
               call pdf_evolve (data%id, x, scale, ff(6:-6:-1))
            end if
         end if
      else
         if (data%has_photon) then
            call pdf_evolve (data%id, x, scale, ff, fph)
         else
            if (data%hoppet_b_matching) then
               call hoppet_eval (xx, qq, ff_dbl)
               ff = ff_dbl
            else
               call pdf_evolve (data%id, x, scale, ff)               
            end if
         end if
      end if
      if (data%has_photon) then
         allocate (fc (count ([data%mask, data%mask_photon])))
         fc = max (pack ([ff, fph], &
              [data%mask, data%mask_photon]), 0._default)
      else
         allocate (fc (count (data%mask)))
         fc = max (pack (ff, data%mask), 0._default)
      end if
    end associate
    call interaction_set_matrix_element (sf_int%interaction_t, fc)
    sf_int%status = SF_EVALUATED
  end subroutine pdf_builtin_apply

\eatline
\nwendcode{}\nwbegindocs{1522}\nwdocspar
\subsection{Strong Coupling}
Since the PDF codes provide a function for computing the running
$\alpha_s$ value, we make this available as an implementation of the
abstract {\tt{}alpha{\char95}qcd{\char95}t} type, which is used for matrix element evaluation.
\nwenddocs{}\nwbegincode{1523}\moddef{SF pdf builtin: public}\plusendmoddef
  public :: alpha_qcd_pdf_builtin_t
\nwendcode{}\nwbegincode{1524}\moddef{SF pdf builtin: types}\plusendmoddef
  type, extends (alpha_qcd_t) :: alpha_qcd_pdf_builtin_t
     type(string_t) :: pdfset_name
     integer :: pdfset_id = -1
   contains
   \LA{}SF pdf builtin: alpha qcd: TBP\RA{}
  end type alpha_qcd_pdf_builtin_t
  
\eatline
\nwendcode{}\nwbegindocs{1525}Output.
\nwenddocs{}\nwbegincode{1526}\moddef{SF pdf builtin: alpha qcd: TBP}\endmoddef
  procedure :: write => alpha_qcd_pdf_builtin_write
\nwendcode{}\nwbegincode{1527}\moddef{SF pdf builtin: procedures}\plusendmoddef
  subroutine alpha_qcd_pdf_builtin_write (object, unit)
    class(alpha_qcd_pdf_builtin_t), intent(in) :: object
    integer, intent(in), optional :: unit
    integer :: u
    u = given_output_unit (unit)
    write (u, "(3x,A)")  "QCD parameters (pdf_builtin):"
    write (u, "(5x,A,A)")  "PDF set = ", char (object%pdfset_name)
    write (u, "(5x,A,I0)") "PDF ID  = ", object%pdfset_id
  end subroutine alpha_qcd_pdf_builtin_write
  
\eatline
\nwendcode{}\nwbegindocs{1528}Calculation: the numeric ID selects the correct PDF set, which must
be properly initialized.
\nwenddocs{}\nwbegincode{1529}\moddef{SF pdf builtin: alpha qcd: TBP}\plusendmoddef
  procedure :: get => alpha_qcd_pdf_builtin_get
\nwendcode{}\nwbegincode{1530}\moddef{SF pdf builtin: procedures}\plusendmoddef
  function alpha_qcd_pdf_builtin_get (alpha_qcd, scale) result (alpha)
    class(alpha_qcd_pdf_builtin_t), intent(in) :: alpha_qcd
    real(default), intent(in) :: scale
    real(default) :: alpha
    alpha = pdf_alphas (alpha_qcd%pdfset_id, scale)
  end function alpha_qcd_pdf_builtin_get
  
\eatline
\nwendcode{}\nwbegindocs{1531}\nwdocspar
Initialization.  We need to access the global initialization status.
\nwenddocs{}\nwbegincode{1532}\moddef{SF pdf builtin: alpha qcd: TBP}\plusendmoddef
  procedure :: init => alpha_qcd_pdf_builtin_init
\nwendcode{}\nwbegincode{1533}\moddef{SF pdf builtin: procedures}\plusendmoddef
  subroutine alpha_qcd_pdf_builtin_init (alpha_qcd, name, path)
    class(alpha_qcd_pdf_builtin_t), intent(out) :: alpha_qcd
    type(string_t), intent(in) :: name
    type(string_t), intent(in) :: path
    alpha_qcd%pdfset_name = name
    alpha_qcd%pdfset_id = pdf_get_id (name)
    if (alpha_qcd%pdfset_id < 0) &
         call msg_fatal ("QCD parameter initialization: PDF set " &
         // char (name) // " is unknown")
    call pdf_init (alpha_qcd%pdfset_id, path)
  end subroutine alpha_qcd_pdf_builtin_init
    
\eatline
\nwendcode{}\nwbegindocs{1534}\nwdocspar
\subsection{Unit tests}
\nwenddocs{}\nwbegincode{1535}\moddef{SF pdf builtin: public}\plusendmoddef
  public :: sf_pdf_builtin_test
\nwendcode{}\nwbegincode{1536}\moddef{SF pdf builtin: tests}\endmoddef
  subroutine sf_pdf_builtin_test (u, results)
    integer, intent(in) :: u
    type(test_results_t), intent(inout) :: results
  \LA{}SF pdf builtin: execute tests\RA{}
  end subroutine sf_pdf_builtin_test
  
\eatline
\nwendcode{}\nwbegindocs{1537}\nwdocspar
\subsubsection{Test structure function data}
Construct and display a test structure function data object.
\nwenddocs{}\nwbegincode{1538}\moddef{SF pdf builtin: execute tests}\endmoddef
  call test (sf_pdf_builtin_1, "sf_pdf_builtin_1", &
       "structure function configuration", &
       u, results)
\nwendcode{}\nwbegincode{1539}\moddef{SF pdf builtin: tests}\plusendmoddef
  subroutine sf_pdf_builtin_1 (u)
    integer, intent(in) :: u
    type(os_data_t) :: os_data
    type(model_data_t), target :: model
    type(pdg_array_t) :: pdg_in
    type(pdg_array_t), dimension(1) :: pdg_out
    integer, dimension(:), allocatable :: pdg1
    class(sf_data_t), allocatable :: data
    type(string_t) :: name
    
    write (u, "(A)")  "* Test output: sf_pdf_builtin_1"
    write (u, "(A)")  "*   Purpose: initialize and display &
         &test structure function data"
    write (u, "(A)")
    
    write (u, "(A)")  "* Create empty data object"
    write (u, "(A)")

    call os_data_init (os_data)

    call model%init_sm_test ()
    pdg_in = PROTON

    allocate (pdf_builtin_data_t :: data)
    call data%write (u)

    write (u, "(A)")
    write (u, "(A)")  "* Initialize"
    write (u, "(A)")

    name = "CTEQ6L"

    select type (data)
    type is (pdf_builtin_data_t)
       call data%init (model, pdg_in, name, &
            os_data%pdf_builtin_datapath)
    end select

    call data%write (u)

    write (u, "(A)")

    write (u, "(1x,A)")  "Outgoing particle codes:"
    call data%get_pdg_out (pdg_out)
    pdg1 = pdg_out(1)
    write (u, "(2x,99(1x,I0))")  pdg1

    call model%final ()
    
    write (u, "(A)")
    write (u, "(A)")  "* Test output end: sf_pdf_builtin_1"

  end subroutine sf_pdf_builtin_1

\eatline
\nwendcode{}\nwbegindocs{1540}\nwdocspar
\subsubsection{Test and probe structure function}
Construct and display a structure function object based on the PDF builtin
structure function.
\nwenddocs{}\nwbegincode{1541}\moddef{SF pdf builtin: execute tests}\plusendmoddef
  call test (sf_pdf_builtin_2, "sf_pdf_builtin_2", &
       "structure function instance", &
       u, results)
\nwendcode{}\nwbegincode{1542}\moddef{SF pdf builtin: tests}\plusendmoddef
  subroutine sf_pdf_builtin_2 (u)
    integer, intent(in) :: u
    type(os_data_t) :: os_data
    type(model_data_t), target :: model
    type(flavor_t) :: flv
    type(pdg_array_t) :: pdg_in
    class(sf_data_t), allocatable, target :: data
    class(sf_int_t), allocatable :: sf_int
    type(string_t) :: name
    type(vector4_t) :: k
    type(vector4_t), dimension(2) :: q
    real(default) :: E
    real(default), dimension(:), allocatable :: r, rb, x
    real(default) :: f
    
    write (u, "(A)")  "* Test output: sf_pdf_builtin_2"
    write (u, "(A)")  "*   Purpose: initialize and fill &
         &test structure function object"
    write (u, "(A)")
    
    write (u, "(A)")  "* Initialize configuration data"
    write (u, "(A)")

    call os_data_init (os_data)
    call model%init_sm_test ()
    call flv%init (PROTON, model)
    pdg_in = PROTON

    call reset_interaction_counter ()
    
    name = "CTEQ6L"

    allocate (pdf_builtin_data_t :: data)
    select type (data)
    type is (pdf_builtin_data_t)
       call data%init (model, pdg_in, name, &
            os_data%pdf_builtin_datapath)
    end select
       
    write (u, "(A)")  "* Initialize structure-function object"
    write (u, "(A)")
    
    call data%allocate_sf_int (sf_int)
    call sf_int%init (data)
    call sf_int%set_beam_index ([1])
    
    call sf_int%write (u)
    
    write (u, "(A)")
    write (u, "(A)")  "* Initialize incoming momentum with E=500"
    write (u, "(A)")
    E = 500
    k = vector4_moving (E, sqrt (E**2 - flv%get_mass ()**2), 3)
    call vector4_write (k, u)
    call sf_int%seed_kinematics ([k])

    write (u, "(A)")
    write (u, "(A)")  "* Set kinematics for x=0.5"
    write (u, "(A)")

    allocate (r (data%get_n_par ()))
    allocate (rb(size (r)))
    allocate (x (size (r)))

    r = 0.5_default
    rb = 1 - r
    call sf_int%complete_kinematics (x, f, r, rb, map=.false.)
    call sf_int%write (u)

    write (u, "(A)")
    write (u, "(A,9(1x,F10.7))")  "x =", x
    write (u, "(A,9(1x,F10.7))")  "f =", f

    write (u, "(A)")
    write (u, "(A)")  "* Recover x from momenta"
    write (u, "(A)")

    q = interaction_get_momenta (sf_int%interaction_t, outgoing=.true.)
    call sf_int%final ()
    deallocate (sf_int)

    call data%allocate_sf_int (sf_int)
    call sf_int%init (data)
    call sf_int%set_beam_index ([1])

    call sf_int%seed_kinematics ([k])
    call interaction_set_momenta (sf_int%interaction_t, q, outgoing=.true.)
    call sf_int%recover_x (x)

    write (u, "(A,9(1x,F10.7))")  "x =", x

    write (u, "(A)")
    write (u, "(A)")  "* Evaluate for Q = 100 GeV"
    write (u, "(A)")

    call sf_int%complete_kinematics (x, f, r, rb, map=.false.) 
    call sf_int%apply (scale = 100._default)
    call sf_int%write (u)


    write (u, "(A)")
    write (u, "(A)")  "* Cleanup"

    call sf_int%final ()
    call model%final ()
    
    write (u, "(A)")
    write (u, "(A)")  "* Test output end: sf_pdf_builtin_2"

  end subroutine sf_pdf_builtin_2

\eatline
\nwendcode{}\nwbegindocs{1543}\nwdocspar
\subsubsection{Strong Coupling}
Test $\alpha_s$ as an implementation of the {\tt{}alpha{\char95}qcd{\char95}t} abstract
type.
\nwenddocs{}\nwbegincode{1544}\moddef{SF pdf builtin: execute tests}\plusendmoddef
  call test (sf_pdf_builtin_3, "sf_pdf_builtin_3", &
       "running alpha_s", &
       u, results)
\nwendcode{}\nwbegincode{1545}\moddef{SF pdf builtin: tests}\plusendmoddef
  subroutine sf_pdf_builtin_3 (u)
    integer, intent(in) :: u
    type(os_data_t) :: os_data
    type(qcd_t) :: qcd
    type(string_t) :: name
    
    write (u, "(A)")  "* Test output: sf_pdf_builtin_3"
    write (u, "(A)")  "*   Purpose: initialize and evaluate alpha_s"
    write (u, "(A)")
    
    write (u, "(A)")  "* Initialize configuration data"
    write (u, "(A)")

    call os_data_init (os_data)

    name = "CTEQ6L"
       
    write (u, "(A)")  "* Initialize qcd object"
    write (u, "(A)")
    
    allocate (alpha_qcd_pdf_builtin_t :: qcd%alpha)
    select type (alpha => qcd%alpha)
    type is (alpha_qcd_pdf_builtin_t)
       call alpha%init (name, os_data%pdf_builtin_datapath)
    end select
    call qcd%write (u)

    write (u, "(A)")
    write (u, "(A)")  "* Evaluate for Q = 100"
    write (u, "(A)")
    
    write (u, "(1x,A,F8.5)")  "alpha = ", qcd%alpha%get (100._default)

    write (u, "(A)")
    write (u, "(A)")  "* Cleanup"

    write (u, "(A)")
    write (u, "(A)")  "* Test output end: sf_pdf_builtin_3"

  end subroutine sf_pdf_builtin_3

\eatline
\nwendcode{}\nwbegindocs{1546}\nwdocspar
\clearpage
%------------------------------------------------------------------------
\section{LHAPDF}
Parton distribution functions (PDFs) are available via an interface to
the LHAPDF standard library.
\nwenddocs{}\nwbegindocs{1547}\nwdocspar
\subsection{The module}
\nwenddocs{}\nwbegincode{1548}\moddef{\code{}sf{\char95}lhapdf.f90\edoc{}}\endmoddef
\LA{}File header\RA{}

module sf_lhapdf

\LA{}Use kinds\RA{}
\LA{}Use strings\RA{}
  use format_defs, only: FMT_17, FMT_19
  use io_units
  use unit_tests
  use system_dependencies, only: LHAPDF_PDFSETS_PATH
  use system_dependencies, only: LHAPDF5_AVAILABLE
  use system_dependencies, only: LHAPDF6_AVAILABLE  
  use diagnostics
  use physics_defs, only: PROTON, PHOTON, PIPLUS, GLUON
  use physics_defs, only: HADRON_REMNANT_SINGLET
  use physics_defs, only: HADRON_REMNANT_TRIPLET
  use physics_defs, only: HADRON_REMNANT_OCTET
  use lorentz
  use sm_qcd
  use pdg_arrays
  use model_data
  use flavors
  use colors
  use quantum_numbers
  use state_matrices
  use polarizations
  use interactions
  use sf_aux
  use sf_base
  use lhapdf !NODEP!
  use hoppet_interface

\LA{}Standard module head\RA{}

\LA{}SF lhapdf: public\RA{}

\LA{}SF lhapdf: types\RA{}

\LA{}SF lhapdf: parameters\RA{}

\LA{}SF lhapdf: variables\RA{}

\LA{}SF lhapdf: interfaces\RA{}

contains

\LA{}SF lhapdf: procedures\RA{}

\LA{}SF lhapdf: tests\RA{}

end module sf_lhapdf
\eatline
\nwendcode{}\nwbegindocs{1549}\nwdocspar
\subsection{Codes for default PDF sets}
The default PDF for protons set is chosen to be CTEQ6ll (LO fit with
LO $\alpha_s$).
\nwenddocs{}\nwbegincode{1550}\moddef{SF lhapdf: parameters}\endmoddef
  character(*), parameter :: LHAPDF5_DEFAULT_PROTON = "cteq6ll.LHpdf"
  character(*), parameter :: LHAPDF5_DEFAULT_PION   = "ABFKWPI.LHgrid"
  character(*), parameter :: LHAPDF5_DEFAULT_PHOTON = "GSG960.LHgrid"
  character(*), parameter :: LHAPDF6_DEFAULT_PROTON = "CT10"
\eatline
\eatline
\nwendcode{}\nwbegindocs{1551}\nwdocspar
\subsection{LHAPDF library interface}
Here we specify explicit interfaces for all LHAPDF routines that we
use below.
\nwenddocs{}\nwbegincode{1552}\moddef{SF lhapdf: interfaces}\endmoddef
  interface
     subroutine InitPDFsetM (set, file)
       integer, intent(in) :: set
       character(*), intent(in) :: file
     end subroutine InitPDFsetM
  end interface

\eatline
\nwendcode{}\nwbegincode{1553}\moddef{SF lhapdf: interfaces}\plusendmoddef
  interface
     subroutine InitPDFM (set, mem)
       integer, intent(in) :: set, mem
     end subroutine InitPDFM
  end interface

\eatline
\nwendcode{}\nwbegincode{1554}\moddef{SF lhapdf: interfaces}\plusendmoddef
  interface
     subroutine numberPDFM (set, n_members)
       integer, intent(in) :: set
       integer, intent(out) :: n_members
     end subroutine numberPDFM
  end interface

\eatline
\nwendcode{}\nwbegincode{1555}\moddef{SF lhapdf: interfaces}\plusendmoddef
  interface
     subroutine evolvePDFM (set, x, q, ff)
       integer, intent(in) :: set
       double precision, intent(in) :: x, q
       double precision, dimension(-6:6), intent(out) :: ff
     end subroutine evolvePDFM
  end interface

\eatline
\nwendcode{}\nwbegincode{1556}\moddef{SF lhapdf: interfaces}\plusendmoddef
  interface
     subroutine evolvePDFphotonM (set, x, q, ff, fphot)
       integer, intent(in) :: set
       double precision, intent(in) :: x, q
       double precision, dimension(-6:6), intent(out) :: ff
       double precision, intent(out) :: fphot
     end subroutine evolvePDFphotonM
  end interface

\eatline
\nwendcode{}\nwbegincode{1557}\moddef{SF lhapdf: interfaces}\plusendmoddef
  interface
     subroutine evolvePDFpM (set, x, q, s, scheme, ff)
       integer, intent(in) :: set
       double precision, intent(in) :: x, q, s
       integer, intent(in) :: scheme
       double precision, dimension(-6:6), intent(out) :: ff
     end subroutine evolvePDFpM
  end interface

\eatline
\nwendcode{}\nwbegincode{1558}\moddef{SF lhapdf: interfaces}\plusendmoddef
  interface
     subroutine GetXminM (set, mem, xmin)
       integer, intent(in) :: set, mem
       double precision, intent(out) :: xmin
     end subroutine GetXminM
  end interface

\eatline
\nwendcode{}\nwbegincode{1559}\moddef{SF lhapdf: interfaces}\plusendmoddef
  interface
     subroutine GetXmaxM (set, mem, xmax)
       integer, intent(in) :: set, mem
       double precision, intent(out) :: xmax
     end subroutine GetXmaxM
  end interface

\eatline
\nwendcode{}\nwbegincode{1560}\moddef{SF lhapdf: interfaces}\plusendmoddef
  interface
     subroutine GetQ2minM (set, mem, q2min)
       integer, intent(in) :: set, mem
       double precision, intent(out) :: q2min
     end subroutine GetQ2minM
  end interface

\eatline
\nwendcode{}\nwbegincode{1561}\moddef{SF lhapdf: interfaces}\plusendmoddef
  interface
     subroutine GetQ2maxM (set, mem, q2max)
       integer, intent(in) :: set, mem
       double precision, intent(out) :: q2max
     end subroutine GetQ2maxM
  end interface

\eatline
\nwendcode{}\nwbegincode{1562}\moddef{SF lhapdf: interfaces}\plusendmoddef
  interface
     function has_photon () result(flag)
        logical :: flag
     end function has_photon
  end interface

\eatline
\nwendcode{}\nwbegindocs{1563}\nwdocspar
\subsection{The LHAPDF status} 
This type holds the initialization status of the LHAPDF system. Entry
1 is for proton PDFs, entry 2 for pion PDFs, entry 3 for photon PDFs.

Since it is connected to the external LHAPDF library, this is a truly
global object.  We implement it as a a private module variable.  To
access it from elsewhere, the caller has to create and initialize an
object of type {\tt{}lhapdf{\char95}status{\char95}t}, which acts as a proxy.
\nwenddocs{}\nwbegincode{1564}\moddef{SF lhapdf: types}\endmoddef
  type :: lhapdf_global_status_t
     private
     logical, dimension(3) :: initialized = .false.
  end type lhapdf_global_status_t

\eatline
\nwendcode{}\nwbegincode{1565}\moddef{SF lhapdf: variables}\endmoddef
  type(lhapdf_global_status_t) :: lhapdf_global_status
\eatline
\nwendcode{}\nwbegincode{1566}\moddef{SF lhapdf: procedures}\endmoddef
  function lhapdf_global_status_is_initialized (set) result (flag)
    logical :: flag
    integer, intent(in), optional :: set
    if (present (set)) then
       select case (set)
       case (1:3);    flag = lhapdf_global_status%initialized(set)
       case default;  flag = .false.
       end select
    else
       flag = any (lhapdf_global_status%initialized)
    end if
  end function lhapdf_global_status_is_initialized

\eatline
\nwendcode{}\nwbegincode{1567}\moddef{SF lhapdf: procedures}\plusendmoddef
  subroutine lhapdf_global_status_set_initialized (set)
    integer, intent(in) :: set
    lhapdf_global_status%initialized(set) = .true.
  end subroutine lhapdf_global_status_set_initialized

\eatline
\nwendcode{}\nwbegindocs{1568}This is the only public procedure, it tells the system to forget
about previous initialization, allowing for changing the chosen PDF
set.  Note that such a feature works only if the global program flow
is serial, so no two distinct sets are accessed simultaneously.  But
this applies to LHAPDF anyway.
\nwenddocs{}\nwbegincode{1569}\moddef{SF lhapdf: public}\endmoddef
  public :: lhapdf_global_reset
\nwendcode{}\nwbegincode{1570}\moddef{SF lhapdf: procedures}\plusendmoddef
  subroutine lhapdf_global_reset ()
    lhapdf_global_status%initialized = .false.
  end subroutine lhapdf_global_reset

\eatline
\nwendcode{}\nwbegindocs{1571}\nwdocspar
\subsection{LHAPDF initialization}
Before using LHAPDF, we have to initialize it with a particular data
set and member.  This applies not just if we use structure functions,
but also if we just use an $\alpha_s$ formula.  The integer {\tt{}set}
should be $1$ for proton, $2$ for pion, and $3$ for photon, but this
is just convention.

It appears as if LHAPDF does not allow for multiple data sets being
used concurrently (?), so multi-threaded usage with different sets
(e.g., a scan) is excluded.  The current setup with a global flag that
indicates initialization is fine as long as Whizard itself is run in
serial mode at the Sindarin level.  If we introduce multithreading in
any form from Sindarin, we have to rethink the implementation of the
LHAPDF interface.  (The same considerations apply to builtin PDFs.)

If the particular set has already been initialized, do nothing.  This
implies that whenever we want to change the setup for a particular
set, we have to reset the LHAPDF status.
{\tt{}lhapdf{\char95}initialize} has an obvious name clash with {\tt{}lhapdf{\char95}init}, 
the reason it works for {\tt{}pdf{\char95}builtin} is that there things are 
outsourced to a separate module (inc. {\tt{}lhapdf{\char95}status} etc.).
\nwenddocs{}\nwbegincode{1572}\moddef{SF lhapdf: public}\plusendmoddef
  public :: lhapdf_initialize
\nwendcode{}\nwbegincode{1573}\moddef{SF lhapdf: procedures}\plusendmoddef
  subroutine lhapdf_initialize (set, prefix, file, member, pdf, b_match)
    integer, intent(in) :: set
    type(string_t), intent(inout) :: prefix
    type(string_t), intent(inout) :: file
    type(lhapdf_pdf_t), intent(inout), optional :: pdf
    integer, intent(inout) :: member
    logical, intent(in), optional :: b_match
    if (prefix == "")  prefix = LHAPDF_PDFSETS_PATH
    if (LHAPDF5_AVAILABLE) then
       if (lhapdf_global_status_is_initialized (set))  return
       if (file == "") then
          select case (set)
          case (1);  file = LHAPDF5_DEFAULT_PROTON
          case (2);  file = LHAPDF5_DEFAULT_PION
          case (3);  file = LHAPDF5_DEFAULT_PHOTON
          end select
       end if
       if (data_file_exists (prefix // "/" // file)) then
          call InitPDFsetM (set, char (prefix // "/" // file))
       else
          call msg_fatal ("LHAPDF: Data file '" &
               // char (file) // "' not found in '" // char (prefix) // "'.")
          return
       end if
       if (.not. dataset_member_exists (set, member)) then
          call msg_error (" LHAPDF: Chosen member does not exist for set '" &
               // char (file) // "', using default.")
          member = 0
       end if
       call InitPDFM (set, member)
    else if (LHAPDF6_AVAILABLE) then
       if (lhapdf_global_status_is_initialized (set) .and. &
            pdf%is_associated ()) return
       if (file == "") then
          select case (set)
          case (1);  file = LHAPDF6_DEFAULT_PROTON
          case (2);  
             call msg_fatal ("LHAPDF6: no pion PDFs supported")
          case (3);  
             call msg_fatal ("LHAPDF6: no photon PDFs supported")
          end select
       end if
       if (data_file_exists (prefix // "/" // file // "/" // file // ".info")) then          
          call pdf%init (char (file), member)
       else
          call msg_fatal ("LHAPDF: Data file '" &
               // char (file) // "' not found in '" // char (prefix) // "'.")
          return
       end if        
    end if
    if (present (b_match)) then
       if (b_match) then
          if (LHAPDF5_AVAILABLE) then
             call hoppet_init (.false.)
          else if (LHAPDF6_AVAILABLE) then
             call hoppet_init (.false., pdf)
          end if
       end if
    end if
    call lhapdf_global_status_set_initialized (set)
  contains
    function data_file_exists (fq_name) result (exist)
      type(string_t), intent(in) :: fq_name
      logical :: exist
      inquire (file = char(fq_name), exist = exist)
    end function data_file_exists
    function dataset_member_exists (set, member) result (exist)
      integer, intent(in) :: set, member
      logical :: exist
      integer :: n_members
      call numberPDFM (set, n_members)
      exist = member >= 0 .and. member <= n_members
    end function dataset_member_exists
  end subroutine lhapdf_initialize

\eatline
\nwendcode{}\nwbegindocs{1574}\nwdocspar
\subsection{Kinematics}
Set kinematics.  If {\tt{}map} is unset, the $r$ and $x$ values
coincide, and the Jacobian $f(r)$ is trivial.

If {\tt{}map} is set, we are asked to provide an efficient mapping.
For the test case, we set $x=r^2$ and consequently $f(r)=2r$.
\nwenddocs{}\nwbegincode{1575}\moddef{SF lhapdf: lhapdf: TBP}\endmoddef
  procedure :: complete_kinematics => lhapdf_complete_kinematics
\nwendcode{}\nwbegincode{1576}\moddef{SF lhapdf: procedures}\plusendmoddef
  subroutine lhapdf_complete_kinematics (sf_int, x, f, r, rb, map)
    class(lhapdf_t), intent(inout) :: sf_int
    real(default), dimension(:), intent(out) :: x
    real(default), intent(out) :: f
    real(default), dimension(:), intent(in) :: r
    real(default), dimension(:), intent(in) :: rb
    logical, intent(in) :: map
    real(default) :: xb1
    if (map) then
       call msg_fatal ("LHAPDF: map flag not supported")
    else
       x(1) = r(1)
       f = 1
    end if
    xb1 = 1 - x(1)
    call sf_int%split_momentum (x, xb1)
    select case (sf_int%status)
    case (SF_DONE_KINEMATICS)
       sf_int%x = x(1)
    case (SF_FAILED_KINEMATICS)
       sf_int%x = 0
       f = 0
    end select
  end subroutine lhapdf_complete_kinematics

\eatline
\nwendcode{}\nwbegindocs{1577}\nwdocspar
\nwenddocs{}\nwbegindocs{1578}Compute inverse kinematics.  Here, we start with the $x$ array and
compute the ``input'' $r$ values and the Jacobian $f$.  After this, we
can set momenta by the same formula as for normal kinematics.
\nwenddocs{}\nwbegincode{1579}\moddef{SF lhapdf: lhapdf: TBP}\plusendmoddef
  procedure :: inverse_kinematics => lhapdf_inverse_kinematics
\nwendcode{}\nwbegincode{1580}\moddef{SF lhapdf: procedures}\plusendmoddef
  subroutine lhapdf_inverse_kinematics (sf_int, x, f, r, rb, map, set_momenta)
    class(lhapdf_t), intent(inout) :: sf_int
    real(default), dimension(:), intent(in) :: x
    real(default), intent(out) :: f
    real(default), dimension(:), intent(out) :: r
    real(default), dimension(:), intent(out) :: rb
    logical, intent(in) :: map
    logical, intent(in), optional :: set_momenta
    real(default) :: xb1
    logical :: set_mom
    set_mom = .false.;  if (present (set_momenta))  set_mom = set_momenta
    if (map) then
       call msg_fatal ("LHAPDF: map flag not supported")
    else
       r(1) = x(1)
       f = 1
    end if
    xb1 = 1 - x(1)
    rb = 1 - r
    if (set_mom) then
       call sf_int%split_momentum (x, xb1)
       select case (sf_int%status)
       case (SF_DONE_KINEMATICS)
          sf_int%x = x(1)
       case (SF_FAILED_KINEMATICS)
          sf_int%x = 0
          f = 0
       end select
    end if
  end subroutine lhapdf_inverse_kinematics

\eatline
\nwendcode{}\nwbegindocs{1581}\nwdocspar
\subsection{The LHAPDF data block} 
The data block holds the incoming flavor (which has to be proton,
pion, or photon), the corresponding pointer to the global access data
(1, 2, or 3), the flag {\tt{}invert} which is set for an antiproton, the
bounds as returned by LHAPDF for the specified set, and a mask that
determines which partons will be actually in use.
\nwenddocs{}\nwbegincode{1582}\moddef{SF lhapdf: public}\plusendmoddef
  public :: lhapdf_data_t
\nwendcode{}\nwbegincode{1583}\moddef{SF lhapdf: types}\plusendmoddef
  type, extends (sf_data_t) :: lhapdf_data_t
     private
     type(string_t) :: prefix
     type(string_t) :: file
     type(lhapdf_pdf_t) :: pdf 
     integer :: member = 0
     class(model_data_t), pointer :: model => null ()
     type(flavor_t) :: flv_in
     integer :: set = 0
     logical :: invert = .false.
     logical :: photon = .false.
     logical :: has_photon = .false.
     integer :: photon_scheme = 0
     real(default) :: xmin = 0, xmax = 0
     real(default) :: qmin = 0, qmax = 0
     logical, dimension(-6:6) :: mask = .true.
     logical :: mask_photon = .true.
     logical :: hoppet_b_matching = .false.     
   contains
     \LA{}SF lhapdf: lhapdf data: TBP\RA{}
  end type lhapdf_data_t

\eatline
\nwendcode{}\nwbegindocs{1584}Generate PDF data.  This is provided as a function, but it has the
side-effect of initializing the requested PDF set.  A finalizer is not
needed.

The library uses double precision, so since the default precision may be
extended or quadruple, we use auxiliary variables for type casting.
\nwenddocs{}\nwbegincode{1585}\moddef{SF lhapdf: lhapdf data: TBP}\endmoddef
  procedure :: init => lhapdf_data_init
\nwendcode{}\nwbegincode{1586}\moddef{SF lhapdf: procedures}\plusendmoddef
  subroutine lhapdf_data_init &
       (data, model, pdg_in, prefix, file, member, photon_scheme, &
            hoppet_b_matching)
    class(lhapdf_data_t), intent(out) :: data
    class(model_data_t), intent(in), target :: model
    type(pdg_array_t), intent(in) :: pdg_in
    type(string_t), intent(in), optional :: prefix, file
    integer, intent(in), optional :: member
    integer, intent(in), optional :: photon_scheme
    logical, intent(in), optional :: hoppet_b_matching
    double precision :: xmin, xmax, q2min, q2max
    external :: InitPDFsetM, InitPDFM, numberPDFM
    external :: GetXminM, GetXmaxM, GetQ2minM, GetQ2maxM
    if (.not. LHAPDF5_AVAILABLE .and. .not. LHAPDF6_AVAILABLE) then
       call msg_fatal ("LHAPDF requested but library is not linked")
       return
    end if
    data%model => model
    if (pdg_array_get_length (pdg_in) /= 1) &
         call msg_fatal ("PDF: incoming particle must be unique")
    call data%flv_in%init (pdg_array_get (pdg_in, 1), model)
    select case (pdg_array_get (pdg_in, 1))
    case (PROTON)
       data%set = 1
    case (-PROTON)
       data%set = 1
       data%invert = .true.
    case (PIPLUS)
       data%set = 2
    case (-PIPLUS)
       data%set = 2
       data%invert = .true.
    case (PHOTON)
       data%set = 3
       data%photon = .true.
       if (present (photon_scheme))  data%photon_scheme = photon_scheme
    case default
       call msg_fatal (" LHAPDF: " &
            // "incoming particle must be (anti)proton, pion, or photon.")
       return
    end select
    if (present (prefix)) then
       data%prefix = prefix
    else
       data%prefix = ""
    end if
    if (present (file)) then
       data%file = file
    else
       data%file = ""
    end if
    if (present (hoppet_b_matching))  data%hoppet_b_matching = hoppet_b_matching    
    if (LHAPDF5_AVAILABLE) then
       call lhapdf_initialize &
            (data%set, data%prefix, data%file, data%member, &
            b_match = data%hoppet_b_matching) 
       call GetXminM (data%set, data%member, xmin)
       call GetXmaxM (data%set, data%member, xmax)
       call GetQ2minM (data%set, data%member, q2min)
       call GetQ2maxM (data%set, data%member, q2max)
       data%xmin = xmin
       data%xmax = xmax
       data%qmin = sqrt (q2min)
       data%qmax = sqrt (q2max)       
       data%has_photon = has_photon ()       
    else if (LHAPDF6_AVAILABLE) then
       call lhapdf_initialize &
            (data%set, data%prefix, data%file, data%member, &
            data%pdf, data%hoppet_b_matching)              
       data%xmin = data%pdf%getxmin ()
       data%xmax = data%pdf%getxmax ()
       data%qmin = sqrt(data%pdf%getq2min ())
       data%qmax = sqrt(data%pdf%getq2max ())
       data%has_photon = data%pdf%has_photon ()       
    end if
  end subroutine lhapdf_data_init

\eatline
\nwendcode{}\nwbegindocs{1587}Enable/disable partons explicitly.  If a mask entry is true,
applying the PDF will generate the corresponding flavor on output.
\nwenddocs{}\nwbegincode{1588}\moddef{LHAPDF: lhapdf data: TBP}\endmoddef
  procedure :: set_mask => lhapdf_data_set_mask
\nwendcode{}\nwbegincode{1589}\moddef{LHAPDF: procedures}\endmoddef
  subroutine lhapdf_data_set_mask (data, mask)
    class(lhapdf_data_t), intent(inout) :: data
    logical, dimension(-6:6), intent(in) :: mask
    data%mask = mask
  end subroutine lhapdf_data_set_mask

\eatline
\nwendcode{}\nwbegindocs{1590}Return the public part of the data set.
\nwenddocs{}\nwbegincode{1591}\moddef{LHAPDF: public}\endmoddef
  public :: lhapdf_data_get_public_info
\nwendcode{}\nwbegincode{1592}\moddef{LHAPDF: procedures}\plusendmoddef
  subroutine lhapdf_data_get_public_info &
       (data, lhapdf_dir, lhapdf_file, lhapdf_member)
    type(lhapdf_data_t), intent(in) :: data
    type(string_t), intent(out) :: lhapdf_dir, lhapdf_file
    integer, intent(out) :: lhapdf_member
    lhapdf_dir = data%prefix
    lhapdf_file = data%file
    lhapdf_member = data%member
  end subroutine lhapdf_data_get_public_info

\eatline
\nwendcode{}\nwbegindocs{1593}Return the number of the member of the data set.
\nwenddocs{}\nwbegincode{1594}\moddef{LHAPDF: public}\plusendmoddef
  public :: lhapdf_data_get_set
\nwendcode{}\nwbegincode{1595}\moddef{LHAPDF: procedures}\plusendmoddef
  function lhapdf_data_get_set(data) result(set)
    type(lhapdf_data_t), intent(in) :: data
    integer :: set
    set = data%set
  end function lhapdf_data_get_set

\eatline
\nwendcode{}\nwbegindocs{1596}Output
\nwenddocs{}\nwbegincode{1597}\moddef{SF lhapdf: lhapdf data: TBP}\plusendmoddef
  procedure :: write => lhapdf_data_write
\nwendcode{}\nwbegincode{1598}\moddef{SF lhapdf: procedures}\plusendmoddef
  subroutine lhapdf_data_write (data, unit, verbose) 
    class(lhapdf_data_t), intent(in) :: data
    integer, intent(in), optional :: unit
    logical, intent(in), optional :: verbose
    logical :: verb
    integer :: u
    if (present (verbose)) then
       verb = verbose
    else
       verb = .false.
    end if       
    u = given_output_unit (unit);  if (u < 0)  return
    write (u, "(1x,A)") "LHAPDF data:"
    if (data%set /= 0) then
       write (u, "(3x,A)", advance="no") "flavor       =  "
       call data%flv_in%write (u);  write (u, *)
       if (verb) then
          write (u, "(3x,A,A)")       "  prefix       =  ", char (data%prefix)
       else
          write (u, "(3x,A,A)")       "  prefix       = ", &
               " <empty (non-verbose version)>"
       end if
       write (u, "(3x,A,A)")       "  file         =  ", char (data%file)
       write (u, "(3x,A,I3)")      "  member       = ", data%member
       write (u, "(3x,A," // FMT_19 // ")") "  x(min)       = ", data%xmin
       write (u, "(3x,A," // FMT_19 // ")") "  x(max)       = ", data%xmax
       write (u, "(3x,A," // FMT_19 // ")") "  Q(min)       = ", data%qmin
       write (u, "(3x,A," // FMT_19 // ")") "  Q(max)       = ", data%qmax
       write (u, "(3x,A,L1)")      "  invert       =  ", data%invert
       if (data%photon)  write (u, "(3x,A,I3)") &
            "  IP2 (scheme) = ", data%photon_scheme
          write (u, "(3x,A,6(1x,L1),1x,A,1x,L1,1x,A,6(1x,L1))") &
               "  mask         = ", &
               data%mask(-6:-1), "*", data%mask(0), "*", data%mask(1:6)
          write (u, "(3x,A,L1)")   "  photon mask  =  ", data%mask_photon
       if (data%set == 1)  write (u, "(3x,A,L1)") &
            "  hoppet_b     =  ", data%hoppet_b_matching
    else
       write (u, "(3x,A)") "[undefined]"       
    end if
  end subroutine lhapdf_data_write
\eatline
\nwendcode{}\nwbegindocs{1599}The number of parameters is one.  We do not generate transverse momentum.
\nwenddocs{}\nwbegincode{1600}\moddef{SF lhapdf: lhapdf data: TBP}\plusendmoddef
  procedure :: get_n_par => lhapdf_data_get_n_par
\nwendcode{}\nwbegincode{1601}\moddef{SF lhapdf: procedures}\plusendmoddef
  function lhapdf_data_get_n_par (data) result (n)
    class(lhapdf_data_t), intent(in) :: data
    integer :: n
    n = 1
  end function lhapdf_data_get_n_par
  
\eatline
\nwendcode{}\nwbegindocs{1602}Return the outgoing particle PDG codes.  This is based on the mask.
\nwenddocs{}\nwbegincode{1603}\moddef{SF lhapdf: lhapdf data: TBP}\plusendmoddef
  procedure :: get_pdg_out => lhapdf_data_get_pdg_out
\nwendcode{}\nwbegincode{1604}\moddef{SF lhapdf: procedures}\plusendmoddef
  subroutine lhapdf_data_get_pdg_out (data, pdg_out)
    class(lhapdf_data_t), intent(in) :: data
    type(pdg_array_t), dimension(:), intent(inout) :: pdg_out
    integer, dimension(:), allocatable :: pdg1
    integer :: n, np, i
    n = count (data%mask)
    np = 0;  if (data%has_photon .and. data%mask_photon)  np = 1
    allocate (pdg1 (n + np))
    pdg1(1:n) = pack ([(i, i = -6, 6)], data%mask)
    if (np == 1)  pdg1(n+np) = PHOTON
    pdg_out(1) = pdg1
  end subroutine lhapdf_data_get_pdg_out
  
\eatline
\nwendcode{}\nwbegindocs{1605}Allocate the interaction record.
\nwenddocs{}\nwbegincode{1606}\moddef{SF lhapdf: lhapdf data: TBP}\plusendmoddef
  procedure :: allocate_sf_int => lhapdf_data_allocate_sf_int
\nwendcode{}\nwbegincode{1607}\moddef{SF lhapdf: procedures}\plusendmoddef
  subroutine lhapdf_data_allocate_sf_int (data, sf_int)
    class(lhapdf_data_t), intent(in) :: data
    class(sf_int_t), intent(inout), allocatable :: sf_int
    allocate (lhapdf_t :: sf_int)
  end subroutine lhapdf_data_allocate_sf_int
  
\eatline
\nwendcode{}\nwbegindocs{1608}Return the numerical PDF set index.
\nwenddocs{}\nwbegincode{1609}\moddef{SF lhapdf: lhapdf data: TBP}\plusendmoddef
  procedure :: get_pdf_set => lhapdf_data_get_pdf_set
\nwendcode{}\nwbegincode{1610}\moddef{SF lhapdf: procedures}\plusendmoddef
  function lhapdf_data_get_pdf_set (data) result (pdf_set)
    class(lhapdf_data_t), intent(in) :: data
    integer :: pdf_set
    pdf_set = data%set
  end function lhapdf_data_get_pdf_set
  
\eatline
\nwendcode{}\nwbegindocs{1611}\nwdocspar
\subsection{The LHAPDF object}
The {\tt{}lhapdf{\char95}t} data type is a $1\to 2$ interaction which describes
the splitting of an (anti)proton into a parton and a beam remnant.  We
stay in the strict forward-splitting limit, but allow some invariant
mass for the beam remnant such that the outgoing parton is exactly
massless.  For a real event, we would replace this by a parton
cascade, where the outgoing partons have virtuality as dictated by
parton-shower kinematics, and transverse momentum is generated.

This is the LHAPDF object which holds input data together with the
interaction.  We also store the $x$ momentum fraction and the scale,
since kinematics and function value are requested at different times.

The PDF application is a $1\to 2$ splitting process, where the
particles are ordered as (hadron, remnant, parton).

Polarization is ignored completely.  The beam particle is colorless,
while partons and beam remnant carry color.  The remnant gets a
special flavor code.
\nwenddocs{}\nwbegincode{1612}\moddef{SF lhapdf: types}\plusendmoddef
  type, extends (sf_int_t) :: lhapdf_t
     type(lhapdf_data_t), pointer :: data => null ()
     real(default) :: x = 0
     real(default) :: q = 0
     real(default) :: s = 0
   contains
   \LA{}SF lhapdf: lhapdf: TBP\RA{}
  end type lhapdf_t 
  
\eatline
\nwendcode{}\nwbegindocs{1613}Type string: display the chosen PDF set.
\nwenddocs{}\nwbegincode{1614}\moddef{SF lhapdf: lhapdf: TBP}\plusendmoddef
  procedure :: type_string => lhapdf_type_string
\nwendcode{}\nwbegincode{1615}\moddef{SF lhapdf: procedures}\plusendmoddef
  function lhapdf_type_string (object) result (string)
    class(lhapdf_t), intent(in) :: object
    type(string_t) :: string
    if (associated (object%data)) then
       string = "LHAPDF: " // object%data%file
    else
       string = "LHAPDF: [undefined]"
    end if
  end function lhapdf_type_string
  
\eatline
\nwendcode{}\nwbegindocs{1616}Output.  Call the interaction routine after displaying the configuration.
\nwenddocs{}\nwbegincode{1617}\moddef{SF lhapdf: lhapdf: TBP}\plusendmoddef
  procedure :: write => lhapdf_write
\nwendcode{}\nwbegincode{1618}\moddef{SF lhapdf: procedures}\plusendmoddef
  subroutine lhapdf_write (object, unit, testflag)
    class(lhapdf_t), intent(in) :: object
    integer, intent(in), optional :: unit
    logical, intent(in), optional :: testflag
    integer :: u
    u = given_output_unit (unit)
    if (associated (object%data)) then
       call object%data%write (u)
       if (object%status >= SF_DONE_KINEMATICS) then
          write (u, "(1x,A)")  "SF parameters:"
          write (u, "(3x,A," // FMT_17 // ")")  "x =", object%x
          if (object%status >= SF_FAILED_EVALUATION) then
             write (u, "(3x,A," // FMT_17 // ")")  "Q =", object%q
          end if
       end if
       call object%base_write (u, testflag)
    else
       write (u, "(1x,A)")  "LHAPDF data: [undefined]"
    end if
  end subroutine lhapdf_write
    
\eatline
\nwendcode{}\nwbegindocs{1619}Initialize.  We know that {\tt{}data} will be of concrete type
{\tt{}sf{\char95}lhapdf{\char95}data{\char95}t}, but we have to cast this explicitly.

For this implementation, we set the incoming and outgoing masses equal
to the physical particle mass, but keep the radiated mass zero.
\nwenddocs{}\nwbegincode{1620}\moddef{SF lhapdf: lhapdf: TBP}\plusendmoddef
  procedure :: init => lhapdf_init
\nwendcode{}\nwbegincode{1621}\moddef{SF lhapdf: procedures}\plusendmoddef
  subroutine lhapdf_init (sf_int, data)
    class(lhapdf_t), intent(out) :: sf_int
    class(sf_data_t), intent(in), target :: data
    type(quantum_numbers_mask_t), dimension(3) :: mask
    type(flavor_t) :: flv, flv_remnant
    type(color_t) :: col0
    type(quantum_numbers_t), dimension(3) :: qn
    integer :: i
    select type (data)
    type is (lhapdf_data_t)
       mask = quantum_numbers_mask (.false., .false., .true.)
       call col0%init ()
       call sf_int%base_init (mask, [0._default], [0._default], [0._default])
       sf_int%data => data
       do i = -6, 6
          if (data%mask(i)) then
             call qn(1)%init (data%flv_in, col = col0)
             if (i == 0) then                
                call flv%init (GLUON, data%model)
                call flv_remnant%init (HADRON_REMNANT_OCTET, data%model)
             else
                call flv%init (i, data%model)
                call flv_remnant%init &
                     (sign (HADRON_REMNANT_TRIPLET, -i), data%model)
             end if
             call qn(2)%init ( &
                  flv = flv_remnant, col = color_from_flavor (flv_remnant, 1))
             call qn(2)%tag_radiated ()
             call qn(3)%init ( &
                  flv = flv, col = color_from_flavor (flv, 1, reverse=.true.))
             call interaction_add_state (sf_int%interaction_t, qn)
          end if
       end do
       if (data%has_photon .and. data%mask_photon) then
          call flv%init (PHOTON, data%model)
          call flv_remnant%init (HADRON_REMNANT_SINGLET, data%model)
          call qn(2)%init (flv = flv_remnant, &
               col = color_from_flavor (flv_remnant, 1))
          call qn(2)%tag_radiated ()
          call qn(3)%init (flv = flv, &
               col = color_from_flavor (flv, 1, reverse=.true.))
          call interaction_add_state (sf_int%interaction_t, qn)
       end if
       call interaction_freeze (sf_int%interaction_t)
       call sf_int%set_incoming ([1])
       call sf_int%set_radiated ([2])
       call sf_int%set_outgoing ([3])
       sf_int%status = SF_INITIAL
    end select
  end subroutine lhapdf_init

\eatline
\nwendcode{}\nwbegindocs{1622}\nwdocspar
\subsection{Structure function}
We have to cast the LHAPDF arguments to/from double precision (possibly
from/to extended/quadruple precision), if necessary.  Furthermore,
some structure functions can yield negative results (sea quarks close
to $x=1$).  We set these unphysical values to zero.
\nwenddocs{}\nwbegincode{1623}\moddef{SF lhapdf: lhapdf: TBP}\plusendmoddef
  procedure :: apply => lhapdf_apply
\nwendcode{}\nwbegincode{1624}\moddef{SF lhapdf: procedures}\plusendmoddef
  subroutine lhapdf_apply (sf_int, scale)
    class(lhapdf_t), intent(inout) :: sf_int
    real(default), intent(in) :: scale
    real(default) :: x, s
    double precision :: xx, qq, ss    
    double precision, dimension(-6:6) :: ff
    double precision :: fphot
    complex(default), dimension(:), allocatable :: fc
    external :: evolvePDFM, evolvePDFpM
    associate (data => sf_int%data)
      sf_int%q = scale
      x = sf_int%x
      s = sf_int%s
      qq = min (data% qmax, scale)
      qq = max (data% qmin, qq)
      if (.not. data% photon) then
         xx = x
         if (data% invert) then
            if (data%has_photon) then
               if (LHAPDF5_AVAILABLE) then
                  call evolvePDFphotonM &
                       (data% set, xx, qq, ff(6:-6:-1), fphot)
               else if (LHAPDF6_AVAILABLE) then
                  call data%pdf%evolve_pdfphotonm &
                       (xx, qq, ff(6:-6:-1), fphot)
               end if
            else
               if (data%hoppet_b_matching) then
                  call hoppet_eval (xx, qq, ff(6:-6:-1)) 
               else
                  if (LHAPDF5_AVAILABLE) then
                     call evolvePDFM (data% set, xx, qq, ff(6:-6:-1))
                  else if (LHAPDF6_AVAILABLE) then
                     call data%pdf%evolve_pdfm (xx, qq, ff(6:-6:-1))
                  end if
               end if
            end if
         else
            if (data%has_photon) then
               if (LHAPDF5_AVAILABLE) then
                  call evolvePDFphotonM (data% set, xx, qq, ff, fphot)
               else if (LHAPDF6_AVAILABLE) then
                  call data%pdf%evolve_pdfphotonm (xx, qq, ff, fphot)
               end if
            else
               if (data%hoppet_b_matching) then
                  call hoppet_eval (xx, qq, ff) 
               else
                  if (LHAPDF5_AVAILABLE) then
                     call evolvePDFM (data% set, xx, qq, ff)
                  else if (LHAPDF6_AVAILABLE) then
                     call data%pdf%evolve_pdfm (xx, qq, ff)
                  end if
               end if
            end if
         end if
      else
         ss = s
         if (LHAPDF5_AVAILABLE) then
            call evolvePDFpM (data% set, xx, qq, &
                 ss, data% photon_scheme, ff)
         else if (LHAPDF6_AVAILABLE) then
            call data%pdf%evolve_pdfpm (xx, qq, ss, &
                 data%photon_scheme, ff)
         end if
      end if
      if (data%has_photon) then
         allocate (fc (count ([data%mask, data%mask_photon])))
         fc = max (pack ([ff, fphot] / x, &
              [data% mask, data%mask_photon]), 0._default)
      else
         allocate (fc (count (data%mask)))
         fc = max (pack (ff / x, data%mask), 0._default)
      end if
    end associate
    call interaction_set_matrix_element (sf_int%interaction_t, fc)
    sf_int%status = SF_EVALUATED
  end subroutine lhapdf_apply
  
\eatline
\nwendcode{}\nwbegindocs{1625}\nwdocspar
\subsection{Strong Coupling}
Since the PDF codes provide a function for computing the running
$\alpha_s$ value, we make this available as an implementation of the
abstract {\tt{}alpha{\char95}qcd{\char95}t} type, which is used for matrix element evaluation.
\nwenddocs{}\nwbegincode{1626}\moddef{SF lhapdf: public}\plusendmoddef
  public :: alpha_qcd_lhapdf_t
\nwendcode{}\nwbegincode{1627}\moddef{SF lhapdf: types}\plusendmoddef
  type, extends (alpha_qcd_t) :: alpha_qcd_lhapdf_t
     type(string_t) :: pdfset_dir
     type(string_t) :: pdfset_file
     integer :: pdfset_member = -1
     type(lhapdf_pdf_t) :: pdf
   contains
   \LA{}SF lhapdf: alpha qcd: TBP\RA{}
  end type alpha_qcd_lhapdf_t
  
\eatline
\nwendcode{}\nwbegindocs{1628}Output. As in earlier versions we leave the LHAPDF path out.
\nwenddocs{}\nwbegincode{1629}\moddef{SF lhapdf: alpha qcd: TBP}\endmoddef
  procedure :: write => alpha_qcd_lhapdf_write
\nwendcode{}\nwbegincode{1630}\moddef{SF lhapdf: procedures}\plusendmoddef
  subroutine alpha_qcd_lhapdf_write (object, unit)
    class(alpha_qcd_lhapdf_t), intent(in) :: object
    integer, intent(in), optional :: unit
    integer :: u
    u = given_output_unit (unit)
    write (u, "(3x,A)")  "QCD parameters (lhapdf):"
    write (u, "(5x,A,A)")  "PDF set    = ", char (object%pdfset_file)
    write (u, "(5x,A,I0)") "PDF member = ", object%pdfset_member
  end subroutine alpha_qcd_lhapdf_write
  
\eatline
\nwendcode{}\nwbegindocs{1631}Calculation: the numeric member ID selects the correct PDF set, which must
be properly initialized.
\nwenddocs{}\nwbegincode{1632}\moddef{SF lhapdf: interfaces}\plusendmoddef
  interface
     double precision function alphasPDF (Q)
       double precision, intent(in) :: Q
     end function alphasPDF
  end interface

\eatline
\nwendcode{}\nwbegindocs{1633}\nwdocspar
\nwenddocs{}\nwbegincode{1634}\moddef{SF lhapdf: alpha qcd: TBP}\plusendmoddef
  procedure :: get => alpha_qcd_lhapdf_get
\nwendcode{}\nwbegincode{1635}\moddef{SF lhapdf: procedures}\plusendmoddef
  function alpha_qcd_lhapdf_get (alpha_qcd, scale) result (alpha)
    class(alpha_qcd_lhapdf_t), intent(in) :: alpha_qcd
    real(default), intent(in) :: scale
    real(default) :: alpha
    if (LHAPDF5_AVAILABLE) then
       alpha = alphasPDF (dble (scale))
    else if (LHAPDF6_AVAILABLE) then
       alpha = alpha_qcd%pdf%alphas_pdf (dble (scale))
    end if
  end function alpha_qcd_lhapdf_get
  
\eatline
\nwendcode{}\nwbegindocs{1636}\nwdocspar
Initialization.  We need to access the (quasi-global) initialization status.
\nwenddocs{}\nwbegincode{1637}\moddef{SF lhapdf: alpha qcd: TBP}\plusendmoddef
  procedure :: init => alpha_qcd_lhapdf_init
\nwendcode{}\nwbegincode{1638}\moddef{SF lhapdf: procedures}\plusendmoddef
  subroutine alpha_qcd_lhapdf_init (alpha_qcd, file, member, path)
    class(alpha_qcd_lhapdf_t), intent(out) :: alpha_qcd
    type(string_t), intent(inout) :: file
    integer, intent(inout) :: member
    type(string_t), intent(inout) :: path
    alpha_qcd%pdfset_file = file
    alpha_qcd%pdfset_member = member
    if (alpha_qcd%pdfset_member < 0) &
         call msg_fatal ("QCD parameter initialization: PDF set " &
         // char (file) // " is unknown")
    if (LHAPDF5_AVAILABLE) then
       call lhapdf_initialize (1, path, file, member)
    else if (LHAPDF6_AVAILABLE) then
       call lhapdf_initialize &
            (1, path, file, member, alpha_qcd%pdf)
    end if
  end subroutine alpha_qcd_lhapdf_init
    
\eatline
\nwendcode{}\nwbegindocs{1639}\nwdocspar
\subsection{Unit tests}
\nwenddocs{}\nwbegincode{1640}\moddef{SF lhapdf: public}\plusendmoddef
  public :: sf_lhapdf_test
\nwendcode{}\nwbegincode{1641}\moddef{SF lhapdf: tests}\endmoddef
  subroutine sf_lhapdf_test (u, results)
    integer, intent(in) :: u
    type(test_results_t), intent(inout) :: results
  \LA{}SF lhapdf: execute tests\RA{}
  end subroutine sf_lhapdf_test
  
\eatline
\nwendcode{}\nwbegindocs{1642}\nwdocspar
\subsubsection{Test structure function data}
Construct and display a test structure function data object.
\nwenddocs{}\nwbegincode{1643}\moddef{SF lhapdf: execute tests}\endmoddef
  if (LHAPDF5_AVAILABLE) then  
     call test (sf_lhapdf_1, "sf_lhapdf5_1", &
          "structure function configuration", &
          u, results)
  else if (LHAPDF6_AVAILABLE) then
     call test (sf_lhapdf_1, "sf_lhapdf6_1", &
          "structure function configuration", &
          u, results)
  end if
\nwendcode{}\nwbegincode{1644}\moddef{SF lhapdf: tests}\plusendmoddef
  subroutine sf_lhapdf_1 (u)
    integer, intent(in) :: u
    type(model_data_t), target :: model
    type(pdg_array_t) :: pdg_in
    type(pdg_array_t), dimension(1) :: pdg_out
    integer, dimension(:), allocatable :: pdg1
    class(sf_data_t), allocatable :: data
    
    write (u, "(A)")  "* Test output: sf_lhapdf_1"
    write (u, "(A)")  "*   Purpose: initialize and display &
         &test structure function data"
    write (u, "(A)")
    
    write (u, "(A)")  "* Create empty data object"
    write (u, "(A)")

    call model%init_sm_test ()
    pdg_in = PROTON

    allocate (lhapdf_data_t :: data)
    call data%write (u)

    write (u, "(A)")
    write (u, "(A)")  "* Initialize"
    write (u, "(A)")

    select type (data)
    type is (lhapdf_data_t)
       call data%init (model, pdg_in)
    end select

    call data%write (u)

    write (u, "(A)")

    write (u, "(1x,A)")  "Outgoing particle codes:"
    call data%get_pdg_out (pdg_out)
    pdg1 = pdg_out(1)
    write (u, "(2x,99(1x,I0))")  pdg1

    call model%final ()
    
    write (u, "(A)")
    write (u, "(A)")  "* Test output end: sf_lhapdf_1"

  end subroutine sf_lhapdf_1

\eatline
\nwendcode{}\nwbegindocs{1645}\nwdocspar
\subsubsection{Test and probe structure function}
Construct and display a structure function object based on the PDF builtin
structure function.
\nwenddocs{}\nwbegincode{1646}\moddef{SF lhapdf: execute tests}\plusendmoddef
  if (LHAPDF5_AVAILABLE) then
     call test (sf_lhapdf_2, "sf_lhapdf5_2", &
          "structure function instance", &
          u, results)
  else if (LHAPDF6_AVAILABLE) then
     call test (sf_lhapdf_2, "sf_lhapdf6_2", &
          "structure function instance", &
          u, results)     
  end if
\nwendcode{}\nwbegincode{1647}\moddef{SF lhapdf: tests}\plusendmoddef
  subroutine sf_lhapdf_2 (u)
    integer, intent(in) :: u
    type(model_data_t), target :: model
    type(flavor_t) :: flv
    type(pdg_array_t) :: pdg_in
    class(sf_data_t), allocatable, target :: data
    class(sf_int_t), allocatable :: sf_int
    type(vector4_t) :: k
    type(vector4_t), dimension(2) :: q
    real(default) :: E
    real(default), dimension(:), allocatable :: r, rb, x
    real(default) :: f
    
    write (u, "(A)")  "* Test output: sf_lhapdf_2"
    write (u, "(A)")  "*   Purpose: initialize and fill &
         &test structure function object"
    write (u, "(A)")
    
    write (u, "(A)")  "* Initialize configuration data"
    write (u, "(A)")

    call model%init_sm_test ()
    call flv%init (PROTON, model)
    pdg_in = PROTON
    call lhapdf_global_reset ()
    
    call reset_interaction_counter ()
    
    allocate (lhapdf_data_t :: data)
    select type (data)
    type is (lhapdf_data_t)
       call data%init (model, pdg_in)
    end select
       
    write (u, "(A)")  "* Initialize structure-function object"
    write (u, "(A)")
    
    call data%allocate_sf_int (sf_int)
    call sf_int%init (data)
    call sf_int%set_beam_index ([1])
    
    call sf_int%write (u)

    write (u, "(A)")
    write (u, "(A)")  "* Initialize incoming momentum with E=500"
    write (u, "(A)")
    E = 500
    k = vector4_moving (E, sqrt (E**2 - flv%get_mass ()**2), 3)
    call vector4_write (k, u)
    call sf_int%seed_kinematics ([k])

    write (u, "(A)")
    write (u, "(A)")  "* Set kinematics for x=0.5"
    write (u, "(A)")

    allocate (r (data%get_n_par ()))
    allocate (rb(size (r)))
    allocate (x (size (r)))

    r = 0.5_default
    rb = 1 - r
    call sf_int%complete_kinematics (x, f, r, rb, map=.false.)
    call sf_int%write (u)

    write (u, "(A)")
    write (u, "(A,9(1x,F10.7))")  "x =", x
    write (u, "(A,9(1x,F10.7))")  "f =", f

    write (u, "(A)")
    write (u, "(A)")  "* Recover x from momenta"
    write (u, "(A)")

    q = interaction_get_momenta (sf_int%interaction_t, outgoing=.true.)
    call sf_int%final ()
    deallocate (sf_int)

    call data%allocate_sf_int (sf_int)
    call sf_int%init (data)
    call sf_int%set_beam_index ([1])

    call sf_int%seed_kinematics ([k])
    call interaction_set_momenta (sf_int%interaction_t, q, outgoing=.true.)
    call sf_int%recover_x (x)

    write (u, "(A,9(1x,F10.7))")  "x =", x

    write (u, "(A)")
    write (u, "(A)")  "* Evaluate for Q = 100 GeV"
    write (u, "(A)")

    call sf_int%complete_kinematics (x, f, r, rb, map=.false.) 
    call sf_int%apply (scale = 100._default)
    call sf_int%write (u, testflag = .true.)


    write (u, "(A)")
    write (u, "(A)")  "* Cleanup"

    call sf_int%final ()
    call model%final ()
    
    write (u, "(A)")
    write (u, "(A)")  "* Test output end: sf_lhapdf_2"

  end subroutine sf_lhapdf_2

\eatline
\nwendcode{}\nwbegindocs{1648}\nwdocspar
\subsubsection{Strong Coupling}
Test $\alpha_s$ as an implementation of the {\tt{}alpha{\char95}qcd{\char95}t} abstract
type.
\nwenddocs{}\nwbegincode{1649}\moddef{SF lhapdf: execute tests}\plusendmoddef
  if (LHAPDF5_AVAILABLE) then
     call test (sf_lhapdf_3, "sf_lhapdf5_3", &
          "running alpha_s", &
          u, results)
  else if (LHAPDF6_AVAILABLE) then
     call test (sf_lhapdf_3, "sf_lhapdf6_3", &
          "running alpha_s", &
          u, results)     
  end if
\nwendcode{}\nwbegincode{1650}\moddef{SF lhapdf: tests}\plusendmoddef
  subroutine sf_lhapdf_3 (u)
    integer, intent(in) :: u
    type(qcd_t) :: qcd
    type(string_t) :: name, path
    integer :: member
    
    write (u, "(A)")  "* Test output: sf_lhapdf_3"
    write (u, "(A)")  "*   Purpose: initialize and evaluate alpha_s"
    write (u, "(A)")
    
    write (u, "(A)")  "* Initialize configuration data"
    write (u, "(A)")

    call lhapdf_global_reset ()

    if (LHAPDF5_AVAILABLE) then
       name = "cteq6ll.LHpdf"
       member = 1
       path = ""
    else if (LHAPDF6_AVAILABLE) then
       name = "CT10"
       member = 1
       path = ""
    end if
       
    write (u, "(A)")  "* Initialize qcd object"
    write (u, "(A)")
    
    allocate (alpha_qcd_lhapdf_t :: qcd%alpha)
    select type (alpha => qcd%alpha)
    type is (alpha_qcd_lhapdf_t)
       call alpha%init (name, member, path)
    end select
    call qcd%write (u)

    write (u, "(A)")
    write (u, "(A)")  "* Evaluate for Q = 100"
    write (u, "(A)")
    
    write (u, "(1x,A,F8.5)")  "alpha = ", qcd%alpha%get (100._default)

    write (u, "(A)")
    write (u, "(A)")  "* Cleanup"

    write (u, "(A)")
    write (u, "(A)")  "* Test output end: sf_lhapdf_3"

  end subroutine sf_lhapdf_3

\eatline
\nwendcode{}\nwfilename{../../../src/me_methods/me_methods.nw}\nwbegindocs{0}% -*- ess-noweb-default-code-mode: f90-mode; noweb-default-code-mode: f90-mode; -*- 
% WHIZARD code as NOWEB source: matrix elements and process libraries

\nwenddocs{}\nwbegindocs{1}\nwdocspar
%%%%%%%%%%%%%%%%%%%%%%%%%%%%%%%%%%%%%%%%%%%%%%%%%%%%%%%%%%%%%%%%%%%%%%%%
\chapter{Interface for Matrix Element Objects}

These modules manage internal and, in particular, external
matrix-element code.
\begin{description}
\item[prc\_core]
  We define the abstract {\tt{}prc{\char95}core{\char95}t} type which handles all specific
  features of kinematics matrix-element evaluation that depend on a particular
  class of processes.  This abstract type supplements the
  {\tt{}prc{\char95}core{\char95}def{\char95}t} type and related types in another module.
  Together, they provide a complete set of matrix-element handlers
  that are implemented in the concrete types below.
\end{description}

These are the implementations:
\begin{description}
\item[prc\_template\_me]
  Implements matrix-element code without actual content (trivial
  value), but full-fledged interface.  This can be used for injecting
  user-defined matrix-element code.
\item[prc\_omega]
  Matrix elements calculated by \oMega\ are the default for WHIZARD.
  Here, we provide all necessary support.
\end{description}


%%%%%%%%%%%%%%%%%%%%%%%%%%%%%%%%%%%%%%%%%%%%%%%%%%%%%%%%%%%%%%%%%%%%%%%%
\section{Abstract process core}
In this module we provide abstract data types for process classes.  Each
process class represents a set of processes which are handled by a common
``method'', e.g., by the \oMega\ matrix-element generator.  The process class
is also able to select a particular implementation for the phase-space and
integration modules.

For a complete implementation of a process class, we have to
provide extensions of the following abstract types:
\begin{description}
\item[prc\_core\_def\_t] process and matrix-element configuration
\item[prc\_writer\_t] (optional) writing external matrix-element code
\item[prc\_driver\_t] accessing the matrix element (internal or external)
\item[prc\_core\_t] evaluating kinematics and matrix element.  The process
  core also selects phase-space and integrator implementations as appropriate
  for the process class and configuration.
\end{description}

In the actual process-handling data structures, each process component
contains an instance of such a process class as its core.  This allows us to
keep the {\tt{}processes} module below, which supervises matrix-element
evaluation, integration, and event generation, free of any reference to
concrete implementations (for the process class, phase space, and
integrator).

There are no unit tests, these are deferred to the {\tt{}processes} module.
\nwenddocs{}\nwbegincode{2}\moddef{\code{}prc{\char95}core.f90\edoc{}}\endmoddef
\LA{}File header\RA{}
module prc_core
  
\LA{}Use kinds\RA{}
\LA{}Use strings\RA{}
  use io_units
  use diagnostics
  use os_interface
  use lorentz
  use model_data
  use model_testbed
  use interactions

  use process_constants
  use prc_core_def
  use process_libraries
  use sf_base

\LA{}Standard module head\RA{}

\LA{}Prc core: public\RA{}

\LA{}Prc core: types\RA{}

\LA{}Prc core: variables\RA{}

\LA{}Prc core: interfaces\RA{}

contains

\LA{}Prc core: procedures\RA{}

end module prc_core
\eatline
\nwendcode{}\nwbegindocs{3}\nwdocspar
\subsection{The process core}
The process core is of abstract data type.  Different types of matrix
element will be represented by different implementations.
\nwenddocs{}\nwbegincode{4}\moddef{Prc core: public}\endmoddef
  public :: prc_core_t
\nwendcode{}\nwbegincode{5}\moddef{Prc core: types}\endmoddef
  type, abstract :: prc_core_t
     class(prc_core_def_t), pointer :: def => null ()
     logical :: data_known = .false.
     type(process_constants_t) :: data
     class(prc_core_driver_t), allocatable :: driver
     logical :: use_color_factors = .false.
     integer :: nc = 3
   contains
   \LA{}Prc core: process core: TBP\RA{}
  end type prc_core_t
  
\eatline
\nwendcode{}\nwbegindocs{6}In any case there must be an output routine.
\nwenddocs{}\nwbegincode{7}\moddef{Prc core: process core: TBP}\endmoddef
  procedure(prc_core_write), deferred :: write
\nwendcode{}\nwbegincode{8}\moddef{Prc core: interfaces}\endmoddef
  abstract interface
     subroutine prc_core_write (object, unit)
       import
       class(prc_core_t), intent(in) :: object
       integer, intent(in), optional :: unit
     end subroutine prc_core_write
  end interface
  
\eatline
\nwendcode{}\nwbegindocs{9}For initialization, we assign a pointer to the process entry in the
relevant library.  This allows us to access all process functions via
the implementation of {\tt{}prc{\char95}core{\char95}t}.

We declare the {\tt{}object} as {\tt{}intent(inout)}, since just after
allocation it may be useful to store some extra data in the object,
which we can then use in the actual initialization.  This applies to
extensions of {\tt{}prc{\char95}core} which override the {\tt{}init} method.
\nwenddocs{}\nwbegincode{10}\moddef{Prc core: process core: TBP}\plusendmoddef
  procedure :: init => prc_core_init
  procedure :: base_init => prc_core_init
\nwendcode{}\nwbegincode{11}\moddef{Prc core: procedures}\endmoddef
  subroutine prc_core_init (object, def, lib, id, i_component)
    class(prc_core_t), intent(inout) :: object
    class(prc_core_def_t), intent(in), target :: def
    type(process_library_t), intent(in), target :: lib
    type(string_t), intent(in) :: id
    integer, intent(in) :: i_component
    object%def => def
    call lib%connect_process (id, i_component, object%data, object%driver)    
    object%data_known = .true.
  end subroutine prc_core_init
  
\eatline
\nwendcode{}\nwbegindocs{12}Return true if the matrix element generation was successful.  This can be
tested by looking at the number of generated flavor states, which should be
nonzero.
\nwenddocs{}\nwbegincode{13}\moddef{Prc core: process core: TBP}\plusendmoddef
  procedure :: has_matrix_element => prc_core_has_matrix_element
\nwendcode{}\nwbegincode{14}\moddef{Prc core: procedures}\plusendmoddef
  function prc_core_has_matrix_element (object) result (flag)
    class(prc_core_t), intent(in) :: object
    logical :: flag
    flag = object%data%n_flv /= 0
  end function prc_core_has_matrix_element

\eatline
\nwendcode{}\nwbegindocs{15}Return true if a MC dataset should be attached to this process
component.  False if it shares the dataset with another component.
\nwenddocs{}\nwbegincode{16}\moddef{Prc core: process core: TBP}\plusendmoddef
  procedure(prc_core_get_flag), deferred :: needs_mcset
\nwendcode{}\nwbegincode{17}\moddef{Prc core: interfaces}\plusendmoddef
  abstract interface
     function prc_core_get_flag (object) result (flag)
       import
       class(prc_core_t), intent(in) :: object
       logical :: flag
     end function prc_core_get_flag
  end interface
  
\eatline
\nwendcode{}\nwbegindocs{18}Return an integer number
\nwenddocs{}\nwbegincode{19}\moddef{Prc core: interfaces}\plusendmoddef
  abstract interface
     function prc_core_get_integer (object) result (i)
       import
       class(prc_core_t), intent(in) :: object
       integer :: i
     end function prc_core_get_integer
  end interface

\eatline
\nwendcode{}\nwbegindocs{20}Return the number of distinct terms requested by
this process component.
\nwenddocs{}\nwbegincode{21}\moddef{Prc core: process core: TBP}\plusendmoddef
  procedure(prc_core_get_integer), deferred :: get_n_terms
\nwendcode{}\nwbegindocs{22}\nwdocspar
Tell whether a particular combination of flavor/helicity/color state
is allowed for the matrix element.
\nwenddocs{}\nwbegincode{23}\moddef{Prc core: process core: TBP}\plusendmoddef
  procedure(prc_core_is_allowed), deferred :: is_allowed
\nwendcode{}\nwbegincode{24}\moddef{Prc core: interfaces}\plusendmoddef
  abstract interface
     function prc_core_is_allowed (object, i_term, f, h, c) result (flag)
       import
       class(prc_core_t), intent(in) :: object
       integer, intent(in) :: i_term, f, h, c
       logical :: flag
     end function prc_core_is_allowed
  end interface

\eatline
\nwendcode{}\nwbegindocs{25}Set the constant process data for a specific term.  By default,
these are the constants stored inside the object, ignoring the term
index.  Type extensions may override this and provide term-specific data.
\nwenddocs{}\nwbegincode{26}\moddef{Prc core: process core: TBP}\plusendmoddef
  procedure :: get_constants => prc_core_get_constants
\nwendcode{}\nwbegincode{27}\moddef{Prc core: procedures}\plusendmoddef
  subroutine prc_core_get_constants (object, data, i_term)
    class(prc_core_t), intent(in) :: object
    type(process_constants_t), intent(out) :: data
    integer, intent(in) :: i_term
    data = object%data
  end subroutine prc_core_get_constants
  
\eatline
\nwendcode{}\nwbegindocs{28}The strong coupling is not among the process constants.  The default
implementation is to return a negative number, which indicates that $\alpha_s$
is not available.  This may be overridden by an implementation that provides
an (event-specific) value.  The value can be stored in the {\tt{}tmp} workspace.
\nwenddocs{}\nwbegincode{29}\moddef{Prc core: process core: TBP}\plusendmoddef
  procedure :: get_alpha_s => prc_core_get_alpha_s
\nwendcode{}\nwbegincode{30}\moddef{Prc core: procedures}\plusendmoddef
  function prc_core_get_alpha_s (object, tmp) result (alpha)
    class(prc_core_t), intent(in) :: object
    class(workspace_t), intent(in), allocatable :: tmp
    real(default) :: alpha
    alpha = -1
  end function prc_core_get_alpha_s
  
\eatline
\nwendcode{}\nwbegindocs{31}Allocate the workspace associated to a process component.  The default is
that there is no workspace, so we do nothing.  A type extension may override
this and allocate a workspace object of appropriate type, which can be used in
further calculations.

In any case, the {\tt{}intent(out)} attribute deletes any previously allocated
workspace.
\nwenddocs{}\nwbegincode{32}\moddef{Prc core: process core: TBP}\plusendmoddef
  procedure :: allocate_workspace => prc_core_ignore_workspace
\nwendcode{}\nwbegincode{33}\moddef{Prc core: procedures}\plusendmoddef
  subroutine prc_core_ignore_workspace (object, tmp)
    class(prc_core_t), intent(in) :: object
    class(workspace_t), intent(inout), allocatable :: tmp
  end subroutine prc_core_ignore_workspace

\eatline
\nwendcode{}\nwbegindocs{34}Initialize the structure-function instance that corresponds to a process
component.  In ordinary cases, this amounts to a straightforward copy of the
given template, but the process core may also choose to modify the setup,
and it may access its own workspace.

The {\tt{}sf{\char95}chain{\char95}instance} is the object that we want to prepare.  The
{\tt{}sf{\char95}chain} argument is a template for its structure.  For the
initialization, we also need the number of channels {\tt{}n{\char95}channel}.

Note: crash with nagfor 5.3.1 if {\tt{}sf{\char95}chain{\char95}instance} is declared
{\tt{}intent(out)}, as would be more appropriate.
\nwenddocs{}\nwbegincode{35}\moddef{Prc core: process core: TBP}\plusendmoddef
  procedure :: init_sf_chain => prc_core_init_sf_chain
\nwendcode{}\nwbegincode{36}\moddef{Prc core: procedures}\plusendmoddef
  subroutine prc_core_init_sf_chain &
       (object, sf_chain_instance, sf_chain, n_channel, tmp)
    class(prc_core_t), intent(in) :: object
    type(sf_chain_instance_t), intent(inout), target :: sf_chain_instance
    type(sf_chain_t), intent(in), target :: sf_chain
    integer, intent(in) :: n_channel
    class(workspace_t), intent(inout), allocatable :: tmp
    call sf_chain_instance%init (sf_chain, n_channel)
  end subroutine prc_core_init_sf_chain
  
\eatline
\nwendcode{}\nwbegindocs{37}Compute the momenta in the hard interaction, taking the seed
kinematics as input.  The {\tt{}i{\char95}term} index tells us which term we want
to compute.  (The standard method is to just transfer the momenta to the hard
interaction.)  
\nwenddocs{}\nwbegincode{38}\moddef{Prc core: process core: TBP}\plusendmoddef
  procedure(prc_core_compute_hard_kinematics), deferred :: &
       compute_hard_kinematics
\nwendcode{}\nwbegincode{39}\moddef{Prc core: interfaces}\plusendmoddef
  abstract interface
     subroutine prc_core_compute_hard_kinematics &
          (object, p_seed, i_term, int_hard, tmp)
       import
       class(prc_core_t), intent(in) :: object
       type(vector4_t), dimension(:), intent(in) :: p_seed
       integer, intent(in) :: i_term
       type(interaction_t), intent(inout) :: int_hard
       class(workspace_t), intent(inout), allocatable :: tmp
     end subroutine prc_core_compute_hard_kinematics
  end interface

\eatline
\nwendcode{}\nwbegindocs{40}Compute the momenta in the effective interaction, taking the hard
kinematics as input.  (This is called only if parton recombination is to be
applied for the process variant.)
\nwenddocs{}\nwbegincode{41}\moddef{Prc core: process core: TBP}\plusendmoddef
  procedure(prc_core_compute_eff_kinematics), deferred :: &
       compute_eff_kinematics
\nwendcode{}\nwbegincode{42}\moddef{Prc core: interfaces}\plusendmoddef
  abstract interface
     subroutine prc_core_compute_eff_kinematics &
          (object, i_term, int_hard, int_eff, tmp)
       import
       class(prc_core_t), intent(in) :: object
       integer, intent(in) :: i_term
       type(interaction_t), intent(in) :: int_hard
       type(interaction_t), intent(inout) :: int_eff
       class(workspace_t), intent(inout), allocatable :: tmp
     end subroutine prc_core_compute_eff_kinematics
  end interface

\eatline
\nwendcode{}\nwbegindocs{43}Recover the missing pieces.  We know the incoming momenta of the
{\tt{}p{\char95}seed} array and the outgoing momenta of the {\tt{}int{\char95}eff}
interaction.  We have to recover the outgoing momenta of {\tt{}p{\char95}seed}
and the incoming momenta of {\tt{}int{\char95}eff}.

(The trivial case is that these are identical.)

Furthermore, if {\tt{}int{\char95}hard} and {\tt{}int{\char95}eff} are not aliased, we have
to set the momenta there.  In the trivial case, {\tt{}int{\char95}eff} is a
pointer to {\tt{}int{\char95}hard}, so {\tt{}int{\char95}hard} should not be touched at all.
\nwenddocs{}\nwbegincode{44}\moddef{Prc core: process core: TBP}\plusendmoddef
  procedure(prc_core_recover_kinematics), deferred :: &
       recover_kinematics
\nwendcode{}\nwbegincode{45}\moddef{Prc core: interfaces}\plusendmoddef
  abstract interface
     subroutine prc_core_recover_kinematics &
          (object, p_seed, int_hard, int_eff, tmp)
       import
       class(prc_core_t), intent(in) :: object
       type(vector4_t), dimension(:), intent(inout) :: p_seed
       type(interaction_t), intent(inout) :: int_hard
       type(interaction_t), intent(inout) :: int_eff
       class(workspace_t), intent(inout), allocatable :: tmp
     end subroutine prc_core_recover_kinematics
  end interface

\eatline
\nwendcode{}\nwbegindocs{46}The process core must implement this function.  Here, {\tt{}j} is the index
of the particular term we want to compute.  The amplitude may depend on the
factorization and renormalization scales.

The {\tt{}tmp} (workspace) argument may be used if it is provided by the caller.
Otherwise, the routine should compute the result directly.
\nwenddocs{}\nwbegincode{47}\moddef{Prc core: process core: TBP}\plusendmoddef
  procedure(prc_core_compute_amplitude), deferred :: compute_amplitude
\nwendcode{}\nwbegincode{48}\moddef{Prc core: interfaces}\plusendmoddef
  abstract interface
     function prc_core_compute_amplitude &
          (object, j, p, f, h, c, fac_scale, ren_scale, alpha_qcd_forced, tmp) &
          result (amp)
       import
       class(prc_core_t), intent(in) :: object
       integer, intent(in) :: j
       type(vector4_t), dimension(:), intent(in) :: p
       integer, intent(in) :: f, h, c
       real(default), intent(in) :: fac_scale, ren_scale
       real(default), intent(in), allocatable :: alpha_qcd_forced
       class(workspace_t), intent(inout), allocatable, optional :: tmp
       complex(default) :: amp
     end function prc_core_compute_amplitude
  end interface
  
\eatline
\nwendcode{}\nwbegindocs{49}\nwdocspar
\subsection{Storage for intermediate results}

The abstract {\tt{}workspace{\char95}t} type allows process cores to set up temporary
workspace.  The object is an extra argument for each of the individual
calculations between kinematics setup and matrix-element evaluation.
\nwenddocs{}\nwbegincode{50}\moddef{Prc core: public}\plusendmoddef
  public :: workspace_t
\nwendcode{}\nwbegincode{51}\moddef{Prc core: types}\plusendmoddef
  type, abstract :: workspace_t
   contains
     procedure(workspace_write), deferred :: write     
     procedure(workspace_reset_new_kinematics), deferred :: reset_new_kinematics
  end type workspace_t
  
\eatline
\nwendcode{}\nwbegindocs{52}For debugging, we should at least have an output routine.
\nwenddocs{}\nwbegincode{53}\moddef{Prc core: interfaces}\plusendmoddef
  abstract interface
     subroutine workspace_write (object, unit)
       import
       class(workspace_t), intent(in) :: object
       integer, intent(in), optional :: unit
     end subroutine workspace_write
  end interface

\eatline
\nwendcode{}\nwbegindocs{54}This is used during the NLO calculation, see there for more information.
\nwenddocs{}\nwbegincode{55}\moddef{Prc core: interfaces}\plusendmoddef
  abstract interface
    subroutine workspace_reset_new_kinematics (object)
      import
      class(workspace_t), intent(inout) :: object
    end subroutine workspace_reset_new_kinematics
  end interface

\eatline
\nwendcode{}\nwbegindocs{56}\nwdocspar
\subsection{Helicity selection data}
This is intended for use with \oMega, but may also be made available to other
process methods.  We set thresholds for counting the times a specific
helicity amplitude is zero.  When the threshold is reached, we skip this
amplitude in subsequent calls.

For initializing the helicity counters, we need an object that holds the two
parameters, the threshold (large real number) and the cutoff (integer).

A helicity value suppressed by more than {\tt{}threshold} (a value which
multiplies {\tt{}epsilon}, to be compared with the average of the current
amplitude, default is $10^{10}$) is treated as zero.  A matrix element is
assumed to be zero and not called again if it has been zero {\tt{}cutoff} times.
\nwenddocs{}\nwbegincode{57}\moddef{Prc core: public}\plusendmoddef
  public :: helicity_selection_t
\nwendcode{}\nwbegincode{58}\moddef{Prc core: types}\plusendmoddef
  type :: helicity_selection_t
     logical :: active = .false.
     real(default) :: threshold = 0
     integer :: cutoff = 0
   contains
   \LA{}Prc core: helicity selection: TBP\RA{}
  end type helicity_selection_t
     
\eatline
\nwendcode{}\nwbegindocs{59}Output.  If the selection is inactive, print nothing.
\nwenddocs{}\nwbegincode{60}\moddef{Prc core: helicity selection: TBP}\endmoddef
  procedure :: write => helicity_selection_write
\nwendcode{}\nwbegincode{61}\moddef{Prc core: procedures}\plusendmoddef
  subroutine helicity_selection_write (object, unit)
    class(helicity_selection_t), intent(in) :: object
    integer, intent(in), optional :: unit
    integer :: u
    u = given_output_unit (unit)
    if (object%active) then
       write (u, "(3x,A)")  "Helicity selection data:"
       write (u, "(5x,A,ES17.10)") &
            "threshold =", object%threshold
       write (u, "(5x,A,I0)") &
            "cutoff    = ", object%cutoff
    end if
  end subroutine helicity_selection_write
    
\eatline
\nwendcode{}\nwbegindocs{62}\nwdocspar
%%%%%%%%%%%%%%%%%%%%%%%%%%%%%%%%%%%%%%%%%%%%%%%%%%%%%%%%%%%%%%%%%%%%%%%%
\section{Template matrix elements}

Here, we provide template matrix elements that are in structure 
very similar to \oMega\ matrix elements, but do not need its
infrastructure. Per default, the matrix elements are flat, i.e.
they have the constant value one. Analogous to the \oMega\
implementation, this section implements the interface
to the templates (via the makefile) and the driver.
\nwenddocs{}\nwbegincode{63}\moddef{\code{}prc{\char95}template{\char95}me.f90\edoc{}}\endmoddef
\LA{}File header\RA{}

module prc_template_me
  
  use iso_c_binding !NODEP!
  use kinds
  use io_units
\LA{}Use strings\RA{}
  use unit_tests
  use system_defs, only: TAB
  use diagnostics
  use os_interface
  use lorentz
  use flavors
  use sm_qcd
  use interactions
  use model_data
  use model_testbed

  use process_constants
  use prclib_interfaces
  use prc_core_def
  use particle_specifiers
  use process_libraries
  use prc_core

\LA{}Standard module head\RA{}

\LA{}Template matrix elements: public\RA{}

\LA{}Template matrix elements: types\RA{}

\LA{}Template matrix elements: interfaces\RA{}

contains
  
\LA{}Template matrix elements: procedures\RA{}

\LA{}Template matrix elements: tests\RA{}

end module prc_template_me
\eatline
\nwendcode{}\nwbegindocs{64}\nwdocspar
\subsection{Process definition}
For the process definition we implement an extension of the
{\tt{}prc{\char95}core{\char95}def{\char95}t} abstract type.
\nwenddocs{}\nwbegincode{65}\moddef{Template matrix elements: public}\endmoddef
  public :: template_me_def_t
\nwendcode{}\nwbegincode{66}\moddef{Template matrix elements: types}\endmoddef
  type, extends (prc_core_def_t) :: template_me_def_t
   contains
   \LA{}Template matrix elements: template ME def: TBP\RA{}
  end type template_me_def_t
  
\eatline
\nwendcode{}\nwbegincode{67}\moddef{Template matrix elements: template ME def: TBP}\endmoddef
  procedure, nopass :: type_string => template_me_def_type_string
\nwendcode{}\nwbegincode{68}\moddef{Template matrix elements: procedures}\endmoddef
  function template_me_def_type_string () result (string)
    type(string_t) :: string
    string = "template"
  end function template_me_def_type_string

\eatline
\nwendcode{}\nwbegindocs{69}Initialization: allocate the writer for the template matrix element.
Also set any data for this process that the writer needs.
\nwenddocs{}\nwbegincode{70}\moddef{Template matrix elements: template ME def: TBP}\plusendmoddef
  procedure :: init => template_me_def_init
\nwendcode{}\nwbegincode{71}\moddef{Template matrix elements: procedures}\plusendmoddef
  subroutine template_me_def_init &
       (object, model, prt_in, prt_out, unity)
    class(template_me_def_t), intent(out) :: object
    class(model_data_t), intent(in), target :: model
    type(string_t), dimension(:), intent(in) :: prt_in
    type(string_t), dimension(:), intent(in) :: prt_out
    logical, intent(in) :: unity
    allocate (template_me_writer_t :: object%writer)
    select type (writer => object%writer)
    type is (template_me_writer_t)
       call writer%init (model, prt_in, prt_out, unity)
    end select
  end subroutine template_me_def_init

\eatline
\nwendcode{}\nwbegindocs{72}Write/read process- and method-specific data.
\nwenddocs{}\nwbegincode{73}\moddef{Template matrix elements: template ME def: TBP}\plusendmoddef
  procedure :: write => template_me_def_write
\nwendcode{}\nwbegincode{74}\moddef{Template matrix elements: procedures}\plusendmoddef
  subroutine template_me_def_write (object, unit)
    class(template_me_def_t), intent(in) :: object
    integer, intent(in) :: unit
    select type (writer => object%writer)
    type is (template_me_writer_t)
       call writer%write (unit)
    end select
  end subroutine template_me_def_write
  
\eatline
\nwendcode{}\nwbegindocs{75}\nwdocspar
\nwenddocs{}\nwbegincode{76}\moddef{Template matrix elements: template ME def: TBP}\plusendmoddef
  procedure :: read => template_me_def_read
\nwendcode{}\nwbegincode{77}\moddef{Template matrix elements: procedures}\plusendmoddef
  subroutine template_me_def_read (object, unit)
    class(template_me_def_t), intent(out) :: object
    integer, intent(in) :: unit
    call msg_bug &
         ("WHIZARD template process definition: input not supported (yet)")
  end subroutine template_me_def_read
  
\eatline
\nwendcode{}\nwbegindocs{78}Allocate the driver for template matrix elements.
\nwenddocs{}\nwbegincode{79}\moddef{Template matrix elements: template ME def: TBP}\plusendmoddef
  procedure :: allocate_driver => template_me_def_allocate_driver
\nwendcode{}\nwbegincode{80}\moddef{Template matrix elements: procedures}\plusendmoddef
  subroutine template_me_def_allocate_driver (object, driver, basename)
    class(template_me_def_t), intent(in) :: object
    class(prc_core_driver_t), intent(out), allocatable :: driver
    type(string_t), intent(in) :: basename
    allocate (template_me_driver_t :: driver)
  end subroutine template_me_def_allocate_driver
  
\eatline
\nwendcode{}\nwbegindocs{81}We need code:
\nwenddocs{}\nwbegincode{82}\moddef{Template matrix elements: template ME def: TBP}\plusendmoddef
  procedure, nopass :: needs_code => template_me_def_needs_code
\nwendcode{}\nwbegincode{83}\moddef{Template matrix elements: procedures}\plusendmoddef
  function template_me_def_needs_code () result (flag)
    logical :: flag
    flag = .true.
  end function template_me_def_needs_code
  
\eatline
\nwendcode{}\nwbegindocs{84}These are the features that a template matrix element provides.
\nwenddocs{}\nwbegincode{85}\moddef{Template matrix elements: template ME def: TBP}\plusendmoddef
  procedure, nopass :: get_features => template_me_def_get_features
\nwendcode{}\nwbegincode{86}\moddef{Template matrix elements: procedures}\plusendmoddef
  subroutine template_me_def_get_features (features)
    type(string_t), dimension(:), allocatable, intent(out) :: features
    allocate (features (6))
    features = [ &
         var_str ("init"), &
         var_str ("update_alpha_s"), &
         var_str ("is_allowed"), &
         var_str ("new_event"), &
         var_str ("get_amplitude")]
  end subroutine template_me_def_get_features

\eatline
\nwendcode{}\nwbegindocs{87}The interface of the specific features.
\nwenddocs{}\nwbegincode{88}\moddef{Template matrix elements: interfaces}\endmoddef
  abstract interface
     subroutine init_t (par) bind(C)
       import
       real(c_default_float), dimension(*), intent(in) :: par
     end subroutine init_t
  end interface
  
  abstract interface
     subroutine update_alpha_s_t (alpha_s) bind(C)
       import
       real(c_default_float), intent(in) :: alpha_s
     end subroutine update_alpha_s_t
  end interface
  
  abstract interface
     subroutine is_allowed_t (flv, hel, col, flag) bind(C)
       import
       integer(c_int), intent(in) :: flv, hel, col
       logical(c_bool), intent(out) :: flag
     end subroutine is_allowed_t
  end interface

  abstract interface
     subroutine new_event_t (p) bind(C)
       import
       real(c_default_float), dimension(0:3,*), intent(in) :: p
     end subroutine new_event_t
  end interface
  
  abstract interface
     subroutine get_amplitude_t (flv, hel, col, amp) bind(C)
       import
       integer(c_int), intent(in) :: flv, hel, col
       complex(c_default_complex), intent(out):: amp
     end subroutine get_amplitude_t
  end interface

\eatline
\eatline
\nwendcode{}\nwbegindocs{89}Connect the template matrix element features with the process driver.
\nwenddocs{}\nwbegincode{90}\moddef{Template matrix elements: template ME def: TBP}\plusendmoddef
  procedure :: connect => template_me_def_connect
\nwendcode{}\nwbegincode{91}\moddef{Template matrix elements: procedures}\plusendmoddef
  subroutine template_me_def_connect (def, lib_driver, i, proc_driver)
    class(template_me_def_t), intent(in) :: def
    class(prclib_driver_t), intent(in) :: lib_driver
    integer, intent(in) :: i
    class(prc_core_driver_t), intent(inout) :: proc_driver
    integer(c_int) :: pid, fid
    type(c_funptr) :: fptr
    select type (proc_driver)
    type is  (template_me_driver_t)
       pid = i
       fid = 1
       call lib_driver%get_fptr (pid, fid, fptr)
       call c_f_procpointer (fptr, proc_driver%init)
       fid = 2
       call lib_driver%get_fptr (pid, fid, fptr)
       call c_f_procpointer (fptr, proc_driver%update_alpha_s)
       fid = 3
       call lib_driver%get_fptr (pid, fid, fptr)
       call c_f_procpointer (fptr, proc_driver%is_allowed)
       fid = 4
       call lib_driver%get_fptr (pid, fid, fptr)
       call c_f_procpointer (fptr, proc_driver%new_event)
       fid = 5
       call lib_driver%get_fptr (pid, fid, fptr)
       call c_f_procpointer (fptr, proc_driver%get_amplitude)
    end select
  end subroutine template_me_def_connect

\eatline
\nwendcode{}\nwbegindocs{92}\nwdocspar
\subsection{The Template Matrix element writer}
Unlike \oMega, the template matrix element is directly written by the main
\whizard\ program, so there will be no entry in the makefile for
calling an external program. The template matrix element writer is
responsible for writing interfaces and wrappers.
\nwenddocs{}\nwbegincode{93}\moddef{Template matrix elements: types}\plusendmoddef
  type, extends (prc_writer_f_module_t) :: template_me_writer_t
     class(model_data_t), pointer :: model => null ()
     type(string_t) :: model_name
     logical :: unity
     type(string_t), dimension(:), allocatable :: prt_in
     type(string_t), dimension(:), allocatable :: prt_out
     integer :: n_in
     integer :: n_out
     integer :: n_tot
   contains
   \LA{}Template matrix elements: template ME writer: TBP\RA{}
  end type template_me_writer_t

\eatline
\nwendcode{}\nwbegindocs{94}The reported type is the same as for the {\tt{}template{\char95}me{\char95}def{\char95}t} type.
\nwenddocs{}\nwbegincode{95}\moddef{Template matrix elements: template ME writer: TBP}\endmoddef
  procedure, nopass :: type_name => template_me_writer_type_name
\nwendcode{}\nwbegincode{96}\moddef{Template matrix elements: procedures}\plusendmoddef
  function template_me_writer_type_name () result (string)
    type(string_t) :: string
    string = "template"
  end function template_me_writer_type_name

\eatline
\nwendcode{}\nwbegindocs{97}Taking into account the prefix for template ME module names.
\nwenddocs{}\nwbegincode{98}\moddef{Template matrix elements: template ME writer: TBP}\plusendmoddef
  procedure, nopass :: get_module_name => template_me_writer_get_module_name
\nwendcode{}\nwbegincode{99}\moddef{Template matrix elements: procedures}\plusendmoddef
  function template_me_writer_get_module_name (id) result (name)
    type(string_t) :: name
    type(string_t), intent(in) :: id
    name = "tpr_" // id
  end function template_me_writer_get_module_name

\eatline
\nwendcode{}\nwbegindocs{100}Output.  This is called by {\tt{}template{\char95}me{\char95}def{\char95}write}.
\nwenddocs{}\nwbegincode{101}\moddef{Template matrix elements: template ME writer: TBP}\plusendmoddef
  procedure :: write => template_me_writer_write
\nwendcode{}\nwbegincode{102}\moddef{Template matrix elements: procedures}\plusendmoddef
  subroutine template_me_writer_write (object, unit)
    class(template_me_writer_t), intent(in) :: object
    integer, intent(in) :: unit
    integer :: i, j
    write (unit, "(5x,A,I0)") "# incoming part. = ", object%n_in
    write (unit, "(7x,A)", advance="no") &
                              "   Initial state: "  
    do i = 1, object%n_in - 1
       write (unit, "(1x,A)", advance="no") char (object%prt_in(i))
    end do
    write (unit, "(1x,A)") char (object%prt_in(object%n_in))    
    write (unit, "(5x,A,I0)") "# outgoing part. = ", object%n_out    
    write (unit, "(7x,A)", advance="no") &
                              "   Final state:   "      
    do j = 1, object%n_out - 1
       write (unit, "(1x,A)", advance="no") char (object%prt_out(i))
    end do    
    write (unit, "(1x,A)") char (object%prt_out(object%n_out))    
    write (unit, "(5x,A,I0)") "# part. (total) = ", object%n_tot
  end subroutine template_me_writer_write

\eatline
\nwendcode{}\nwbegindocs{103}Initialize with process data.
\nwenddocs{}\nwbegincode{104}\moddef{Template matrix elements: template ME writer: TBP}\plusendmoddef
  procedure :: init => template_me_writer_init
\nwendcode{}\nwbegincode{105}\moddef{Template matrix elements: procedures}\plusendmoddef
  subroutine template_me_writer_init (writer, model, &
       prt_in, prt_out, unity)
    class(template_me_writer_t), intent(out) :: writer
    class(model_data_t), intent(in), target :: model
    type(string_t), dimension(:), intent(in) :: prt_in
    type(string_t), dimension(:), intent(in) :: prt_out
    logical, intent(in) :: unity
    writer%model => model
    writer%model_name = model%get_name ()
    writer%n_in = size (prt_in)
    writer%n_out = size (prt_out)
    writer%n_tot = size (prt_in) + size (prt_out)
    allocate (writer%prt_in (size (prt_in)), source = prt_in)
    allocate (writer%prt_out (size (prt_out)), source = prt_out)    
    writer%unity = unity    
  end subroutine template_me_writer_init

\eatline
\nwendcode{}\nwbegindocs{106}The makefile is the driver file for the test matrix elements.
\nwenddocs{}\nwbegincode{107}\moddef{Template matrix elements: template ME writer: TBP}\plusendmoddef
  procedure :: write_makefile_code => template_me_write_makefile_code
\nwendcode{}\nwbegincode{108}\moddef{Template matrix elements: procedures}\plusendmoddef
  subroutine template_me_write_makefile_code (writer, unit, id, os_data, testflag)
    class(template_me_writer_t), intent(in) :: writer
    integer, intent(in) :: unit
    type(string_t), intent(in) :: id
    type(os_data_t), intent(in) :: os_data
    logical, intent(in), optional :: testflag
    write (unit, "(5A)")  "SOURCES += ", char (id), ".f90"
    write (unit, "(5A)")  "OBJECTS += ", char (id), ".lo"
    write (unit, "(5A)")  "clean-", char (id), ":"
    write (unit, "(5A)")  TAB, "rm -f tpr_", char (id), ".mod"
    write (unit, "(5A)")  TAB, "rm -f ", char (id), ".lo"
    write (unit, "(5A)")  "CLEAN_SOURCES += ", char (id), ".f90"    
    write (unit, "(5A)")  "CLEAN_OBJECTS += tpr_", char (id), ".mod"       
    write (unit, "(5A)")  "CLEAN_OBJECTS += ", char (id), ".lo"
    write (unit, "(5A)")  char (id), ".lo: ", char (id), ".f90"
    write (unit, "(5A)")  TAB, "$(LTFCOMPILE) $<"
  end subroutine template_me_write_makefile_code

\eatline
\nwendcode{}\nwbegindocs{109}The source is written by this routine.
\nwenddocs{}\nwbegincode{110}\moddef{Template matrix elements: template ME writer: TBP}\plusendmoddef
  procedure :: write_source_code => template_me_write_source_code
\nwendcode{}\nwbegincode{111}\moddef{Template matrix elements: procedures}\plusendmoddef
  subroutine template_me_write_source_code (writer, id)
    class(template_me_writer_t), intent(in) :: writer
    type(string_t), intent(in) :: id
    integer, dimension(writer%n_in) :: prt_in, mult_in, col_in
    type(flavor_t), dimension(1:writer%n_in) :: flv_in    
    integer, dimension(writer%n_out) :: prt_out, mult_out
    integer, dimension(writer%n_tot) :: prt, mult
    integer, dimension(:,:), allocatable :: sxxx
    integer :: dummy, status
    type(flavor_t), dimension(1:writer%n_out) :: flv_out    
    type(string_t) :: proc_str, comment_str
    integer :: u, i, j
    integer :: hel, hel_in, hel_out, fac, factor, col_fac
    type(string_t) :: filename
    comment_str = ""
    do i = 1, writer%n_in
       comment_str = comment_str // writer%prt_in(i) // " " 
    end do   
    do j = 1, writer%n_out
       comment_str = comment_str // writer%prt_out(j) // " " 
    end do       
    do i = 1, writer%n_in
       prt_in(i) = writer%model%get_pdg (writer%prt_in(i))
       call flv_in(i)%init (prt_in(i), writer%model)
       mult_in(i) = flv_in(i)%get_multiplicity ()
       col_in(i) = abs (flv_in(i)%get_color_type ())
       mult(i) = mult_in(i)
       end do
    do j = 1, writer%n_out
       prt_out(j) = writer%model%get_pdg (writer%prt_out(j))    
       call flv_out(j)%init (prt_out(j), writer%model)       
       mult_out(j) = flv_out(j)%get_multiplicity ()       
       mult(writer%n_in + j) = mult_out(j)
       end do
    prt(1:writer%n_in) = prt_in(1:writer%n_in)
    prt(writer%n_in+1:writer%n_tot) = prt_out(1:writer%n_out)
    proc_str = converter (prt)
    hel_in = product (mult_in)
    hel_out = product (mult_out)
    col_fac = product (col_in)
    hel = hel_in * hel_out
    fac = hel
    dummy = 1
    factor = 1
    if (writer%n_out >= 3) then
       do i = 3, writer%n_out
          factor = factor * (i - 2) * (i - 1)
       end do
    end if    
    factor = factor * col_fac
    allocate (sxxx(1:hel,1:writer%n_tot))
    call create_spin_table (dummy,hel,fac,mult,sxxx)
    call msg_message ("Writing test matrix element for process '" &
         // char (id) // "'")
    filename = id // ".f90"
    u = free_unit ()
    open (unit=u, file=char(filename), action="write")
    write (u, "(A)") "! File generated automatically by WHIZARD"   
    write (u, "(A)") "!                                        "
    write (u, "(A)") "! Note that irresp. of what you demanded WHIZARD"
    write (u, "(A)") "! treats this as colorless process       "    
    write (u, "(A)") "!                                        "
    write (u, "(A)") "module tpr_" // char(id)
    write (u, "(A)") "                                         "
    write (u, "(A)") "  use kinds"    
    write (u, "(A)") "  use omega_color, OCF => omega_color_factor"        
    write (u, "(A)") "                                         "
    write (u, "(A)") "  implicit none"        
    write (u, "(A)") "  private"            
    write (u, "(A)") "                                         "    
    write (u, "(A)") "  public :: md5sum"        
    write (u, "(A)") "  public :: number_particles_in, number_particles_out"       
    write (u, "(A)") "  public :: number_spin_states, spin_states"
    write (u, "(A)") "  public :: number_flavor_states, flavor_states"
    write (u, "(A)") "  public :: number_color_flows, color_flows"
    write (u, "(A)") "  public :: number_color_indices, number_color_factors, &"
    write (u, "(A)") "     color_factors, color_sum, openmp_supported"
    write (u, "(A)") "  public :: init, final, update_alpha_s"  
    write (u, "(A)") "                                         "    
    write (u, "(A)") "  public :: new_event, is_allowed, get_amplitude"        
    write (u, "(A)") "       "      
    write (u, "(A)") "  real(default), parameter :: &"
    write (u, "(A)") "       & conv = 0.38937966e12_default"
    write (u, "(A)") "       "
    write (u, "(A)") "  real(default), parameter :: &"
    write (u, "(A)") "       & pi = 3.1415926535897932384626433832795028841972_default"
    write (u, "(A)") "       "
    write (u, "(A)") "  real(default), parameter :: &"
    if (writer%unity) then
       write (u, "(A)") "                   & const = 1"
    else 
       write (u, "(A,1x,I0,A)") "       & const = (16 * pi / conv) * " &
          // "(16 * pi**2)**(", writer%n_out, "-2) " 
    end if
    write (u, "(A)") "       "    
    write (u, "(A,1x,I0)") "  integer, parameter, private :: n_prt =  ", &
       writer%n_tot
    write (u, "(A,1x,I0)") "  integer, parameter, private :: n_in = ", &
       writer%n_in
    write (u, "(A,1x,I0)") "  integer, parameter, private :: n_out = ", &
       writer%n_out
    write (u, "(A)") "  integer, parameter, private :: n_cflow = 1"
    write (u, "(A)") "  integer, parameter, private :: n_cindex = 2"
    write (u, "(A)") "  !!! We ignore tensor products and take only one flavor state."
    write (u, "(A)") "  integer, parameter, private :: n_flv = 1"
    write (u, "(A,1x,I0)") "  integer, parameter, private :: n_hel = ", hel
    write (u, "(A)") "                                           "
    write (u, "(A)") "  logical, parameter, private :: T = .true."
    write (u, "(A)") "  logical, parameter, private :: F = .false."
    write (u, "(A)") "                                           "    
    do i = 1, hel
       write (u, "(A)") "  integer, dimension(n_prt), parameter, private :: &"
       write (u, "(A)") "    " // s_conv(i) // " = [ " // &
            char(converter(sxxx(i,1:writer%n_tot))) // " ]"
    end do 
    write (u, "(A)") "  integer, dimension(n_prt,n_hel), parameter, private :: table_spin_states = &"
    write (u, "(A)") "    reshape ( [ & "
    do i = 1, hel-1
       write (u, "(A)") "                 " // s_conv(i) // ", & " 
    end do 
    write (u, "(A)") "                 " // s_conv(hel) // " & "     
    write (u, "(A)") "              ], [ n_prt, n_hel ] )"
    write (u, "(A)") "                                                 "
    write (u, "(A)") "  integer, dimension(n_prt), parameter, private :: &"
    write (u, "(A)") "    f0001 = [ " // char(proc_str) // " ]   !  " // char(comment_str)
    write (u, "(A)") "  integer, dimension(n_prt,n_flv), parameter, private :: table_flavor_states = &"
    write (u, "(A)") "    reshape ( [ f0001 ], [ n_prt, n_flv ] )"
    write (u, "(A)") "                                                 " 
    write (u, "(A)") "  integer, dimension(n_cindex, n_prt), parameter, private :: &"
    write (u, "(A)") "    c0001 = reshape ( [ " // char (dummy_colorizer (flv_in)) // " " // &
      (repeat ("0,0, ", writer%n_out-1)) // "0,0 ], " // " [ n_cindex, n_prt ] )"
    write (u, "(A)") "  integer, dimension(n_cindex, n_prt, n_cflow), parameter, private :: &"
    write (u, "(A)") "  table_color_flows = reshape ( [ c0001 ], [ n_cindex, n_prt, n_cflow ] )"
    write (u, "(A)") "                                           "   
    write (u, "(A)") "  logical, dimension(n_prt), parameter, private :: & "
    write (u, "(A)") "    g0001 = [ "  // (repeat ("F, ", writer%n_tot-1)) // "F ] "
    write (u, "(A)") "  logical, dimension(n_prt, n_cflow), parameter, private " &
         // ":: table_ghost_flags = &"
    write (u, "(A)") "    reshape ( [ g0001 ], [ n_prt, n_cflow ] )"
    write (u, "(A)") "                                           "   
    write (u, "(A)") "  integer, parameter, private :: n_cfactors = 1"
    write (u, "(A)") "  type(OCF), dimension(n_cfactors), parameter, private :: &"
    write (u, "(A)") "    table_color_factors = [  OCF(1,1,+1._default) ]"
    write (u, "(A)") "                                           "   
    write (u, "(A)") "  logical, dimension(n_flv), parameter, private :: a0001 = [ T ]"   
    write (u, "(A)") "  logical, dimension(n_flv, n_cflow), parameter, private :: &"   
    write (u, "(A)") "    flv_col_is_allowed = reshape ( [ a0001 ], [ n_flv, n_cflow ] )"   
    write (u, "(A)") "                                           "   
    write (u, "(A)") "  complex(default), dimension (n_flv, n_hel, n_cflow), private, save :: amp"    
    write (u, "(A)") "                                           "
    write (u, "(A)") "  logical, dimension(n_hel), private, save :: hel_is_allowed = T"
    write (u, "(A)") "                                           "
    write (u, "(A)") "contains"          
    write (u, "(A)") "                                           "      
    write (u, "(A)") "  pure function md5sum ()"
    write (u, "(A)") "    character(len=32) :: md5sum"    
    write (u, "(A)") "    ! DON'T EVEN THINK of modifying the following line!"        
    write (u, "(A)") "    md5sum = """ // writer%md5sum // """"
    write (u, "(A)") "  end function md5sum"
    write (u, "(A)") "                                           "          
    write (u, "(A)") "  subroutine init (par)"
    write (u, "(A)") "    real(default), dimension(*), intent(in) :: par"    
    write (u, "(A)") "  end subroutine init"    
    write (u, "(A)") "                                           " 
    write (u, "(A)") "  subroutine final ()" 
    write (u, "(A)") "  end subroutine final" 
    write (u, "(A)") "                                           " 
    write (u, "(A)") "  subroutine update_alpha_s (alpha_s)" 
    write (u, "(A)") "    real(default), intent(in) :: alpha_s"        
    write (u, "(A)") "  end subroutine update_alpha_s" 
    write (u, "(A)") "                                           " 
    write (u, "(A)") "  pure function number_particles_in () result (n)"
    write (u, "(A)") "    integer :: n"    
    write (u, "(A)") "    n = n_in"
    write (u, "(A)") "  end function number_particles_in"
    write (u, "(A)") "                                           "              
    write (u, "(A)") "  pure function number_particles_out () result (n)"
    write (u, "(A)") "    integer :: n"    
    write (u, "(A)") "    n = n_out"
    write (u, "(A)") "  end function number_particles_out"
    write (u, "(A)") "                                           "                  
    write (u, "(A)") "  pure function number_spin_states () result (n)"
    write (u, "(A)") "    integer :: n"    
    write (u, "(A)") "    n = size (table_spin_states, dim=2)"
    write (u, "(A)") "  end function number_spin_states"
    write (u, "(A)") "                                           "                      
    write (u, "(A)") "  pure subroutine spin_states (a)"
    write (u, "(A)") "    integer, dimension(:,:), intent(out) :: a"    
    write (u, "(A)") "    a = table_spin_states"
    write (u, "(A)") "  end subroutine spin_states"    
    write (u, "(A)") "                                           "                          
    write (u, "(A)") "  pure function number_flavor_states () result (n)"
    write (u, "(A)") "    integer :: n"    
    write (u, "(A)") "    n = 1"
    write (u, "(A)") "  end function number_flavor_states"
    write (u, "(A)") "                                           "                      
    write (u, "(A)") "  pure subroutine flavor_states (a)"
    write (u, "(A)") "    integer, dimension(:,:), intent(out) :: a"    
    write (u, "(A)") "    a = table_flavor_states"
    write (u, "(A)") "  end subroutine flavor_states"
    write (u, "(A)") "                                           "                          
    write (u, "(A)") "  pure function number_color_indices () result (n)"
    write (u, "(A)") "    integer :: n"    
    write (u, "(A)") "    n = size(table_color_flows, dim=1)"
    write (u, "(A)") "  end function number_color_indices"
    write (u, "(A)") "                                           "                          
    write (u, "(A)") "  pure subroutine color_factors (cf)"
    write (u, "(A)") "    type(OCF), dimension(:), intent(out) :: cf"    
    write (u, "(A)") "    cf = table_color_factors"
    write (u, "(A)") "  end subroutine color_factors"
    write (u, "(A)") "                                           "                              
    !pure unless OpenMP
    !write (u, "(A)") "  pure function color_sum (flv, hel) result (amp2)"
    write (u, "(A)") "  function color_sum (flv, hel) result (amp2)"
    write (u, "(A)") "    integer, intent(in) :: flv, hel"
    write (u, "(A)") "    real(kind=default) :: amp2"
    write (u, "(A)") "    amp2 = real (omega_color_sum (flv, hel, amp, table_color_factors))"
    write (u, "(A)") "  end function color_sum"
    write (u, "(A)") "                                           "       
    write (u, "(A)") "  pure function number_color_flows () result (n)"
    write (u, "(A)") "    integer :: n"    
    write (u, "(A)") "    n = size (table_color_flows, dim=3)"
    write (u, "(A)") "  end function number_color_flows"
    write (u, "(A)") "                                           "                                  
    write (u, "(A)") "  pure subroutine color_flows (a, g)"
    write (u, "(A)") "    integer, dimension(:,:,:), intent(out) :: a"
    write (u, "(A)") "    logical, dimension(:,:), intent(out) :: g"
    write (u, "(A)") "    a = table_color_flows"
    write (u, "(A)") "    g = table_ghost_flags"
    write (u, "(A)") "  end subroutine color_flows"    
    write (u, "(A)") "                                           "                              
    write (u, "(A)") "  pure function number_color_factors () result (n)"
    write (u, "(A)") "    integer :: n"    
    write (u, "(A)") "    n = size (table_color_factors)"
    write (u, "(A)") "  end function number_color_factors"
    write (u, "(A)") "                                           "                                  
    write (u, "(A)") "  pure function openmp_supported () result (status)"
    write (u, "(A)") "    logical :: status"
    write (u, "(A)") "    status = .false."
    write (u, "(A)") "  end function openmp_supported"
    write (u, "(A)") "                                           "                                  
    write (u, "(A)") "  subroutine new_event (p)"
    write (u, "(A)") "    real(default), dimension(0:3,*), intent(in) :: p"    
    write (u, "(A)") "    call calculate_amplitudes (amp, p)"        
    write (u, "(A)") "  end subroutine new_event"
    write (u, "(A)") "                                           "              
    write (u, "(A)") "  pure function is_allowed (flv, hel, col) result (yorn)"                  
    write (u, "(A)") "    logical :: yorn"                  
    write (u, "(A)") "    integer, intent(in) :: flv, hel, col"                  
    write (u, "(A)") "    yorn = hel_is_allowed(hel) .and. flv_col_is_allowed(flv,col)"                  
    write (u, "(A)") "  end function is_allowed"                     
    write (u, "(A)") "                                           "                  
    write (u, "(A)") "  pure function get_amplitude (flv, hel, col) result (amp_result)"
    write (u, "(A)") "    complex(default) :: amp_result"    
    write (u, "(A)") "    integer, intent(in) :: flv, hel, col"            
    write (u, "(A)") "    amp_result = amp (flv, hel, col)"        
    write (u, "(A)") "  end function get_amplitude"
    write (u, "(A)") "                                           "                  
    write (u, "(A)") "  pure subroutine calculate_amplitudes (amp, k)"
    write (u, "(A)") "    complex(default), dimension(:,:,:), intent(out) :: amp"    
    write (u, "(A)") "    real(default), dimension(0:3,*), intent(in) :: k"    
    write (u, "(A)") "    real(default) :: fac"        
    write (u, "(A)") "    integer :: i"            
    write (u, "(A)") "    ! We give all helicities the same weight!"            
    if (writer%unity) then 
       write (u, "(A,1x,I0,1x,A)") "    fac = ", col_fac
       write (u, "(A)") "    amp = const * sqrt(fac)"
    else
       write (u, "(A,1x,I0,1x,A)") "    fac = ", factor 
       write (u, "(A)") "    amp = sqrt((2 * (k(0,1)*k(0,2) &"
       write (u, "(A,1x,I0,A)") "         - dot_product (k(1:,1), k(1:,2)))) ** (3-", &
                                  writer%n_out, ")) * sqrt(const * fac)"
    end if                                  
    write (u, "(A,1x,I0,A)") "    amp = amp / sqrt(", hel_out, "._default)"
    write (u, "(A)") "  end subroutine calculate_amplitudes"
    write (u, "(A)") "                                           "                  
    write (u, "(A)") "end module tpr_" // char(id)    
    close (u, iostat=status)
    deallocate (sxxx)
  contains
    function s_conv (num) result (chrt)
      integer, intent(in) :: num
      character(len=10) :: chrt
      write (chrt, "(I10)") num
      chrt = trim(adjustl(chrt))
      if (num < 10) then
         chrt = "s000" // chrt
      else if (num < 100) then
         chrt = "s00" // chrt
      else if (num < 1000) then 
         chrt = "s0" // chrt     
      else
         chrt = "s" // chrt            
      end if             
    end function s_conv
    function converter (flv) result (str)
      integer, dimension(:), intent(in) :: flv
      type(string_t) :: str
      character(len=150), dimension(size(flv)) :: chrt
      integer :: i
      str = ""
      do i = 1, size(flv) - 1
         write (chrt(i), "(I10)") flv(i)
         str = str // var_str(trim(adjustl(chrt(i)))) // ", "
      end do    
      write (chrt(size(flv)), "(I10)") flv(size(flv))
      str = str // trim(adjustl(chrt(size(flv))))
    end function converter
    integer function sj (j,m)
      integer, intent(in) :: j, m
      if (((j == 1) .and. (m == 1)) .or. &
          ((j == 2) .and. (m == 2)) .or. &
          ((j == 3) .and. (m == 3)) .or. &
          ((j == 4) .and. (m == 3)) .or. &
          ((j == 5) .and. (m == 4))) then
         sj = 1
      else if (((j == 2) .and. (m == 1)) .or. &
          ((j == 3) .and. (m == 1)) .or. &         
          ((j == 4) .and. (m == 2)) .or. &
          ((j == 5) .and. (m == 2))) then
         sj = -1
      else if (((j == 3) .and. (m == 2)) .or. &
          ((j == 5) .and. (m == 3))) then
         sj = 0
      else if (((j == 4) .and. (m == 1)) .or. &
          ((j == 5) .and. (m == 1))) then         
         sj = -2
      else if (((j == 4) .and. (m == 4)) .or. &
          ((j == 5) .and. (m == 5))) then         
         sj = 2
      else
         call msg_fatal ("template_me_write_source_code: Wrong spin type")
      end if
    end function sj    
    recursive subroutine create_spin_table (index, nhel, fac, mult, inta)
      integer, intent(inout) :: index, fac
      integer, intent(in) :: nhel
      integer, dimension(:), intent(in) :: mult
      integer, dimension(nhel,size(mult)), intent(out) :: inta    
      integer :: j
      if (index > size(mult)) return
      fac = fac / mult(index)
      do j = 1, nhel 
         inta(j,index) = sj (mult(index),mod(((j-1)/fac),mult(index))+1)
      end do   
      index = index + 1
      call create_spin_table (index, nhel, fac, mult, inta)
    end subroutine create_spin_table  
    function dummy_colorizer (flv) result (str)
      type(flavor_t), dimension(:), intent(in) :: flv
      type(string_t) :: str
      integer :: i, k
      str = ""
      k = 0
      do i = 1, size(flv)
         k = k + 1
         select case (flv(i)%get_color_type ())
         case (1,-1)
            str = str // "0,0, "
         case (3)
            str = str // int2string(k) // ",0, "
         case (-3)
            str = str // "0," // int2string(-k) // ", "
         case (8)
            str = str // int2string(k) // "," // int2string(-k-1) // ", "
            k = k + 1
         case default
            call msg_error ("Color type not supported.")
         end select
      end do    
      str = adjustl(trim(str))
    end function dummy_colorizer    
  end subroutine template_me_write_source_code    

\eatline
\nwendcode{}\nwbegindocs{112}Return the name of a procedure that implements a given feature, as
it is provided by the template matrix-element code.  Template ME names
are chosen completely in analogy to the \oMega\ matrix element
conventions. 
\nwenddocs{}\nwbegincode{113}\moddef{Template matrix elements: template ME writer: TBP}\plusendmoddef
  procedure, nopass :: get_procname => template_me_writer_get_procname
\nwendcode{}\nwbegincode{114}\moddef{Template matrix elements: procedures}\plusendmoddef
  function template_me_writer_get_procname (feature) result (name)
    type(string_t) :: name
    type(string_t), intent(in) :: feature
    select case (char (feature))
    case ("n_in");   name = "number_particles_in"
    case ("n_out");  name = "number_particles_out"
    case ("n_flv");  name = "number_flavor_states"
    case ("n_hel");  name = "number_spin_states"
    case ("n_col");  name = "number_color_flows"
    case ("n_cin");  name = "number_color_indices"
    case ("n_cf");   name = "number_color_factors"
    case ("flv_state");  name = "flavor_states"
    case ("hel_state");  name = "spin_states"
    case ("col_state");  name = "color_flows"
    case default
       name = feature
    end select
  end function template_me_writer_get_procname
  
\eatline
\nwendcode{}\nwbegindocs{115}The interfaces for the template-specific features.
\nwenddocs{}\nwbegincode{116}\moddef{Template matrix elements: template ME writer: TBP}\plusendmoddef
  procedure :: write_interface => template_me_write_interface
\nwendcode{}\nwbegincode{117}\moddef{Template matrix elements: procedures}\plusendmoddef
  subroutine template_me_write_interface (writer, unit, id, feature)
    class(template_me_writer_t), intent(in) :: writer
    integer, intent(in) :: unit
    type(string_t), intent(in) :: id
    type(string_t), intent(in) :: feature
    type(string_t) :: name
    name = writer%get_c_procname (id, feature)
    write (unit, "(2x,9A)")  "interface"
    select case (char (feature))
    case ("init")
       write (unit, "(5x,9A)")  "subroutine ", char (name), " (par) bind(C)"
       write (unit, "(7x,9A)")  "import"
       write (unit, "(7x,9A)")  "real(c_default_float), dimension(*), &
            &intent(in) :: par"
       write (unit, "(5x,9A)")  "end subroutine ", char (name)
    case ("update_alpha_s")
       write (unit, "(5x,9A)")  "subroutine ", char (name), " (alpha_s) bind(C)"
       write (unit, "(7x,9A)")  "import"
       write (unit, "(7x,9A)")  "real(c_default_float), intent(in) :: alpha_s"
       write (unit, "(5x,9A)")  "end subroutine ", char (name)
    case ("is_allowed")
       write (unit, "(5x,9A)")  "subroutine ", char (name), " &
            &(flv, hel, col, flag) bind(C)"
       write (unit, "(7x,9A)")  "import"
       write (unit, "(7x,9A)")  "integer(c_int), intent(in) :: flv, hel, col"
       write (unit, "(7x,9A)")  "logical(c_bool), intent(out) :: flag"    
       write (unit, "(5x,9A)")  "end subroutine ", char (name)
    case ("new_event")
       write (unit, "(5x,9A)")  "subroutine ", char (name), " (p) bind(C)"
       write (unit, "(7x,9A)")  "import"
       write (unit, "(7x,9A)")  "real(c_default_float), dimension(0:3,*), &
            &intent(in) :: p"
       write (unit, "(5x,9A)")  "end subroutine ", char (name)
    case ("get_amplitude")
       write (unit, "(5x,9A)")  "subroutine ", char (name), " &
            &(flv, hel, col, amp) bind(C)"
       write (unit, "(7x,9A)")  "import"
       write (unit, "(7x,9A)")  "integer(c_int), intent(in) :: flv, hel, col"
       write (unit, "(7x,9A)")  "complex(c_default_complex), intent(out) &
            &:: amp"    
       write (unit, "(5x,9A)")  "end subroutine ", char (name)
    end select
    write (unit, "(2x,9A)")  "end interface"
  end subroutine template_me_write_interface

\eatline
\nwendcode{}\nwbegindocs{118}The wrappers have to take into account conversion between C and
Fortran data types.

NOTE: The case {\tt{}c{\char95}default{\char95}float} $\neq$ {\tt{}default} is not yet covered.
\nwenddocs{}\nwbegincode{119}\moddef{Template matrix elements: template ME writer: TBP}\plusendmoddef
  procedure :: write_wrapper => template_me_write_wrapper
\nwendcode{}\nwbegincode{120}\moddef{Template matrix elements: procedures}\plusendmoddef
  subroutine template_me_write_wrapper (writer, unit, id, feature)
    class(template_me_writer_t), intent(in) :: writer
    integer, intent(in) :: unit
    type(string_t), intent(in) :: id, feature
    type(string_t) :: name
    name = writer%get_c_procname (id, feature)
    write (unit, *)
    select case (char (feature))
    case ("init")
       write (unit, "(9A)")  "subroutine ", char (name), " (par) bind(C)"
       write (unit, "(2x,9A)")  "use iso_c_binding"
       write (unit, "(2x,9A)")  "use kinds"
       write (unit, "(2x,9A)")  "use tpr_", char (id)
       write (unit, "(2x,9A)")  "real(c_default_float), dimension(*), &
            &intent(in) :: par"
       if (c_default_float == default) then
          write (unit, "(2x,9A)")  "call ", char (feature), " (par)"
       end if
       write (unit, "(9A)")  "end subroutine ", char (name)
    case ("update_alpha_s")
       write (unit, "(9A)")  "subroutine ", char (name), " (alpha_s) bind(C)"
       write (unit, "(2x,9A)")  "use iso_c_binding"
       write (unit, "(2x,9A)")  "use kinds"
       write (unit, "(2x,9A)")  "use tpr_", char (id)
       if (c_default_float == default) then
          write (unit, "(2x,9A)")  "real(c_default_float), intent(in) &
               &:: alpha_s"
          write (unit, "(2x,9A)")  "call ", char (feature), " (alpha_s)"
       end if
       write (unit, "(9A)")  "end subroutine ", char (name)
    case ("is_allowed")
       write (unit, "(9A)")  "subroutine ", char (name), &
            " (flv, hel, col, flag) bind(C)"
       write (unit, "(2x,9A)")  "use iso_c_binding"
       write (unit, "(2x,9A)")  "use kinds"
       write (unit, "(2x,9A)")  "use tpr_", char (id)
       write (unit, "(2x,9A)")  "integer(c_int), intent(in) :: flv, hel, col"
       write (unit, "(2x,9A)")  "logical(c_bool), intent(out) :: flag"    
       write (unit, "(2x,9A)")  "flag = ", char (feature), &
            " (int (flv), int (hel), int (col))"
       write (unit, "(9A)")  "end subroutine ", char (name)
    case ("new_event")
       write (unit, "(9A)")  "subroutine ", char (name), " (p) bind(C)"
       write (unit, "(2x,9A)")  "use iso_c_binding"
       write (unit, "(2x,9A)")  "use kinds"
       write (unit, "(2x,9A)")  "use tpr_", char (id)
       if (c_default_float == default) then
          write (unit, "(2x,9A)")  "real(c_default_float), dimension(0:3,*), &
               &intent(in) :: p"
          write (unit, "(2x,9A)")  "call ", char (feature), " (p)"
       end if
       write (unit, "(9A)")  "end subroutine ", char (name)
    case ("get_amplitude")
       write (unit, "(9A)")  "subroutine ", char (name), &
            " (flv, hel, col, amp) bind(C)"
       write (unit, "(2x,9A)")  "use iso_c_binding"
       write (unit, "(2x,9A)")  "use kinds"
       write (unit, "(2x,9A)")  "use tpr_", char (id)
       write (unit, "(2x,9A)")  "integer(c_int), intent(in) :: flv, hel, col"
       write (unit, "(2x,9A)")  "complex(c_default_complex), intent(out) &
            &:: amp"    
       write (unit, "(2x,9A)")  "amp = ", char (feature), &
            " (int (flv), int (hel), int (col))"
       write (unit, "(9A)")  "end subroutine ", char (name)
    end select
  end subroutine template_me_write_wrapper

\eatline
\nwendcode{}\nwbegindocs{121}\nwdocspar
\subsection{Driver}
\nwenddocs{}\nwbegincode{122}\moddef{Template matrix elements: types}\plusendmoddef
  type, extends (prc_core_driver_t) :: template_me_driver_t
     procedure(init_t), nopass, pointer :: &
          init => null ()
     procedure(update_alpha_s_t), nopass, pointer :: &
          update_alpha_s => null ()
     procedure(is_allowed_t), nopass, pointer :: &
          is_allowed => null ()
     procedure(new_event_t), nopass, pointer :: &
          new_event => null ()
     procedure(get_amplitude_t), nopass, pointer :: &
          get_amplitude => null ()
   contains
   \LA{}Template matrix elements: template ME driver: TBP\RA{}
  end type template_me_driver_t

\eatline
\nwendcode{}\nwbegindocs{123}The reported type is the same as for the {\tt{}template{\char95}me{\char95}def{\char95}t} type.
\nwenddocs{}\nwbegincode{124}\moddef{Template matrix elements: template ME driver: TBP}\endmoddef
  procedure, nopass :: type_name => template_me_driver_type_name
\nwendcode{}\nwbegincode{125}\moddef{Template matrix elements: procedures}\plusendmoddef
  function template_me_driver_type_name () result (string)
    type(string_t) :: string
    string = "template"
  end function template_me_driver_type_name

\eatline
\nwendcode{}\nwbegindocs{126}\nwdocspar
\subsection{High-level process definition}
This procedure wraps the details filling a process-component
definition entry as appropriate for an template matrix element.

NOTE: For calling the {\tt{}import{\char95}component} method, we must explicitly
address the {\tt{}process{\char95}def{\char95}t} parent object.  The natural way to call
the method of the extended type triggers a bug in gfortran 4.6.  The
string array arguments {\tt{}prt{\char95}in} and {\tt{}prt{\char95}out} become corrupted and
cause a segfault.
\nwenddocs{}\nwbegincode{127}\moddef{Template matrix elements: public}\plusendmoddef
  public :: template_me_make_process_component
\nwendcode{}\nwbegincode{128}\moddef{Template matrix elements: procedures}\plusendmoddef
  subroutine template_me_make_process_component (entry, component_index, &
         model, model_name, prt_in, prt_out, unity)
    class(process_def_entry_t), intent(inout) :: entry
    integer, intent(in) :: component_index
    type(string_t), intent(in) :: model_name
    class(model_data_t), intent(in), target :: model
    type(string_t), dimension(:), intent(in) :: prt_in
    type(string_t), dimension(:), intent(in) :: prt_out
    logical, intent(in) :: unity
    class(prc_core_def_t), allocatable :: def
    allocate (template_me_def_t :: def)
    select type (def)
    type is (template_me_def_t)
       call def%init (model, prt_in, prt_out, unity)
    end select
    call entry%process_def_t%import_component (component_index, &
         n_out = size (prt_out), &
         prt_in  = new_prt_spec (prt_in), &
         prt_out = new_prt_spec (prt_out), &
         method = var_str ("template"), &
         variant = def)
  end subroutine template_me_make_process_component
    
\eatline
\nwendcode{}\nwbegindocs{129}\nwdocspar
\subsection{The {\tt{}prc{\char95}template{\char95}me{\char95}t} wrapper}
This is an instance of the generic {\tt{}prc{\char95}core{\char95}t} object.  It contains a
pointer to the process definition ({\tt{}template{\char95}me{\char95}def{\char95}t}), a data component
({\tt{}process{\char95}constants{\char95}t}), and the matrix-element driver
({\tt{}template{\char95}me{\char95}driver{\char95}t}).
\nwenddocs{}\nwbegincode{130}\moddef{Template matrix elements: public}\plusendmoddef
  public :: prc_template_me_t
\nwendcode{}\nwbegincode{131}\moddef{Template matrix elements: types}\plusendmoddef
  type, extends (prc_core_t) :: prc_template_me_t
     real(default), dimension(:), allocatable :: par
   contains
   \LA{}Template matrix elements: prc template ME: TBP\RA{}
  end type prc_template_me_t
  
\eatline
\nwendcode{}\nwbegindocs{132}The workspace associated to a {\tt{}prc{\char95}template{\char95}me{\char95}t} object contains a single flag.
The flag is used to suppress re-evaluating the matrix element for each
quantum-number combination, after the first amplitude belonging to a given
kinematics has been computed.

We can also store the value of a running coupling once it has been calculated
for an event.  The default value is negative, which indicates an undefined
value in this context.
\nwenddocs{}\nwbegincode{133}\moddef{Template matrix elements: types}\plusendmoddef
  type, extends (workspace_t) :: template_me_state_t
     logical :: new_kinematics = .true.
     real(default) :: alpha_qcd = -1
   contains
     procedure :: write => template_me_state_write
     procedure :: reset_new_kinematics => template_me_state_reset_new_kinematics
  end type template_me_state_t
  
\eatline
\nwendcode{}\nwbegincode{134}\moddef{Template matrix elements: procedures}\plusendmoddef
  subroutine template_me_state_write (object, unit)
    class(template_me_state_t), intent(in) :: object
    integer, intent(in), optional :: unit
    integer :: u
    u = given_output_unit (unit)
    write (u, "(3x,A,L1)")  "Template ME state: new kinematics = ", &
         object%new_kinematics
  end subroutine template_me_state_write
  
\eatline
\nwendcode{}\nwbegindocs{135}\nwdocspar
\nwenddocs{}\nwbegincode{136}\moddef{Template matrix elements: procedures}\plusendmoddef
  subroutine template_me_state_reset_new_kinematics (object)
    class(template_me_state_t), intent(inout) :: object
  end subroutine template_me_state_reset_new_kinematics

\nwendcode{}\nwbegindocs{137}\nwdocspar
\nwenddocs{}\nwbegindocs{138}Allocate the workspace with the above specific type.
\nwenddocs{}\nwbegincode{139}\moddef{Template matrix elements: prc template ME: TBP}\endmoddef
  procedure :: allocate_workspace => prc_template_me_allocate_workspace
\nwendcode{}\nwbegincode{140}\moddef{Template matrix elements: procedures}\plusendmoddef
  subroutine prc_template_me_allocate_workspace (object, tmp)
    class(prc_template_me_t), intent(in) :: object
    class(workspace_t), intent(inout), allocatable :: tmp
    allocate (template_me_state_t :: tmp)
  end subroutine prc_template_me_allocate_workspace
  
\eatline
\nwendcode{}\nwbegindocs{141}The following procedures are inherited from the base type as deferred, thus
must be implemented.  The corresponding unit tests are skipped here; the
procedures are tested when called from the {\tt{}processes} module.

Output: print just the ID of the associated matrix element.  Then display any
stored parameters.
\nwenddocs{}\nwbegincode{142}\moddef{Template matrix elements: prc template ME: TBP}\plusendmoddef
  procedure :: write => prc_template_me_write
\nwendcode{}\nwbegincode{143}\moddef{Template matrix elements: procedures}\plusendmoddef
  subroutine prc_template_me_write (object, unit)
    class(prc_template_me_t), intent(in) :: object
    integer, intent(in), optional :: unit
    integer :: u, i
    u = given_output_unit (unit)
    write (u, "(3x,A)", advance="no")  "Template process core:"
    if (object%data_known) then
       write (u, "(1x,A)")  char (object%data%id)
    else
       write (u, "(1x,A)")  "[undefined]"
    end if
    if (allocated (object%par)) then
       write (u, "(3x,A)")  "Parameter array:"
       do i = 1, size (object%par)
          write (u, "(5x,I0,1x,ES17.10)")  i, object%par(i)
       end do
    end if
  end subroutine prc_template_me_write
  
\eatline
\nwendcode{}\nwbegindocs{144}Temporarily store the parameter array inside the {\tt{}prc{\char95}template{\char95}me}
object, so we can use it later during the actual initialization. 
\nwenddocs{}\nwbegincode{145}\moddef{Template matrix elements: prc template ME: TBP}\plusendmoddef
  procedure :: set_parameters => prc_template_me_set_parameters
\nwendcode{}\nwbegincode{146}\moddef{Template matrix elements: procedures}\plusendmoddef
  subroutine prc_template_me_set_parameters (prc_template_me, model)
    class(prc_template_me_t), intent(inout) :: prc_template_me
    class(model_data_t), intent(in), target, optional :: model
    if (present (model)) then
       if (.not. allocated (prc_template_me%par)) &
            allocate (prc_template_me%par (model%get_n_real ()))
       call model%real_parameters_to_c_array (prc_template_me%par)
    end if
  end subroutine prc_template_me_set_parameters
  
\eatline
\nwendcode{}\nwbegindocs{147}To fully initialize the process core, we perform base
initialization, then initialize the external matrix element code.

This procedure overrides the {\tt{}init} method of the base type, which
we nevertheless can access via its binding {\tt{}base{\char95}init}.  When done, we
have an allocated driver.  The driver will call the {\tt{}init} procedure
for the external matrix element, and thus transfer the parameter set to
where it finally belongs.

\nwenddocs{}\nwbegincode{148}\moddef{Template matrix elements: prc template ME: TBP}\plusendmoddef
  procedure :: init => prc_template_me_init
\nwendcode{}\nwbegincode{149}\moddef{Template matrix elements: procedures}\plusendmoddef
  subroutine prc_template_me_init (object, def, lib, id, i_component)
    class(prc_template_me_t), intent(inout) :: object
    class(prc_core_def_t), intent(in), target :: def
    type(process_library_t), intent(in), target :: lib
    type(string_t), intent(in) :: id
    integer, intent(in) :: i_component
    call object%base_init (def, lib, id, i_component)
    call object%activate_parameters ()
  end subroutine prc_template_me_init
    
\eatline
\nwendcode{}\nwbegindocs{150}Activate the stored parameters by transferring them to the external
matrix element.  
\nwenddocs{}\nwbegincode{151}\moddef{Template matrix elements: prc template ME: TBP}\plusendmoddef
  procedure :: activate_parameters => prc_template_me_activate_parameters
\nwendcode{}\nwbegincode{152}\moddef{Template matrix elements: procedures}\plusendmoddef
  subroutine prc_template_me_activate_parameters (object)
    class (prc_template_me_t), intent(inout) :: object
    if (allocated (object%driver)) then
       if (allocated (object%par)) then
          select type (driver => object%driver)
          type is (template_me_driver_t)
             if (associated (driver%init))  call driver%init (object%par)
          end select
       else
          call msg_bug ("prc_template_me_activate: parameter set is not allocated")
       end if
    else
       call msg_bug ("prc_template_me_activate: driver is not allocated")
    end if
  end subroutine prc_template_me_activate_parameters
    
\eatline
\nwendcode{}\nwbegindocs{153}The template ME process is an independent process which needs its own Monte
Carlo parameter set for integration.
\nwenddocs{}\nwbegincode{154}\moddef{Template matrix elements: prc template ME: TBP}\plusendmoddef
  procedure :: needs_mcset => prc_template_me_needs_mcset
\nwendcode{}\nwbegincode{155}\moddef{Template matrix elements: procedures}\plusendmoddef
  function prc_template_me_needs_mcset (object) result (flag)
    class(prc_template_me_t), intent(in) :: object
    logical :: flag
    flag = .true.
  end function prc_template_me_needs_mcset
  
\eatline
\nwendcode{}\nwbegindocs{156}There is only one term for this tree-level process.
\nwenddocs{}\nwbegincode{157}\moddef{Template matrix elements: prc template ME: TBP}\plusendmoddef
  procedure :: get_n_terms => prc_template_me_get_n_terms
\nwendcode{}\nwbegincode{158}\moddef{Template matrix elements: procedures}\plusendmoddef
  function prc_template_me_get_n_terms (object) result (n)
    class(prc_template_me_t), intent(in) :: object
    integer :: n
    n = 1
  end function prc_template_me_get_n_terms

\eatline
\nwendcode{}\nwbegindocs{159}Tell whether a particular combination of flavor, helicity, color is
allowed.  Here we have to consult the matrix-element driver.
\nwenddocs{}\nwbegincode{160}\moddef{Template matrix elements: prc template ME: TBP}\plusendmoddef
  procedure :: is_allowed => prc_template_me_is_allowed
\nwendcode{}\nwbegincode{161}\moddef{Template matrix elements: procedures}\plusendmoddef
 function prc_template_me_is_allowed (object, i_term, f, h, c) result (flag)
    class(prc_template_me_t), intent(in) :: object
    integer, intent(in) :: i_term, f, h, c
    logical :: flag
    logical(c_bool) :: cflag
    select type (driver => object%driver)
    type is (template_me_driver_t)
       call driver%is_allowed (f, h, c, cflag)
       flag = cflag
    end select
  end function prc_template_me_is_allowed
 
\eatline
\nwendcode{}\nwbegindocs{162}Transfer the generated momenta directly to the hard interaction in
the (only) term.  We assume that everything has been set up correctly,
so the array fits.

%We reset the {\tt{}new{\char95}kinematics} flag, so that the next call to
%{\tt{}compute{\char95}amplitude} will evaluate the matrix element.
\nwenddocs{}\nwbegincode{163}\moddef{Template matrix elements: prc template ME: TBP}\plusendmoddef
  procedure :: compute_hard_kinematics => &
       prc_template_me_compute_hard_kinematics
\nwendcode{}\nwbegincode{164}\moddef{Template matrix elements: procedures}\plusendmoddef
  subroutine prc_template_me_compute_hard_kinematics &
       (object, p_seed, i_term, int_hard, tmp)
    class(prc_template_me_t), intent(in) :: object
    type(vector4_t), dimension(:), intent(in) :: p_seed
    integer, intent(in) :: i_term
    type(interaction_t), intent(inout) :: int_hard
    class(workspace_t), intent(inout), allocatable :: tmp
    call interaction_set_momenta (int_hard, p_seed)
!     if (allocated (tmp)) then
!        select type (tmp)
!        type is (template_me_state_t);  tmp%new_kinematics = .true.
!        end select
!     end if
  end subroutine prc_template_me_compute_hard_kinematics
  
\eatline
\nwendcode{}\nwbegindocs{165}This procedure is not called for {\tt{}prc{\char95}template{\char95}me{\char95}t}, just a placeholder.
\nwenddocs{}\nwbegincode{166}\moddef{Template matrix elements: prc template ME: TBP}\plusendmoddef
  procedure :: compute_eff_kinematics => &
       prc_template_me_compute_eff_kinematics
\nwendcode{}\nwbegincode{167}\moddef{Template matrix elements: procedures}\plusendmoddef
  subroutine prc_template_me_compute_eff_kinematics &
       (object, i_term, int_hard, int_eff, tmp)
    class(prc_template_me_t), intent(in) :: object
    integer, intent(in) :: i_term
    type(interaction_t), intent(in) :: int_hard
    type(interaction_t), intent(inout) :: int_eff
    class(workspace_t), intent(inout), allocatable :: tmp
  end subroutine prc_template_me_compute_eff_kinematics
  
\eatline
\nwendcode{}\nwbegindocs{168}Recover the momenta within the hard interaction.  For tree-level
processes, this is a trivial copy of the incoming seed and outgoing
effective momenta.  The effective interaction is a pointer to the hard
interaction, so the latter is ignored.
\nwenddocs{}\nwbegincode{169}\moddef{Template matrix elements: prc template ME: TBP}\plusendmoddef
  procedure :: recover_kinematics => prc_template_me_recover_kinematics
\nwendcode{}\nwbegincode{170}\moddef{Template matrix elements: tests}\endmoddef
  subroutine prc_template_me_recover_kinematics &
       (object, p_seed, int_hard, int_eff, tmp)
    class(prc_template_me_t), intent(in) :: object
    type(vector4_t), dimension(:), intent(inout) :: p_seed
    type(interaction_t), intent(inout) :: int_hard
    type(interaction_t), intent(inout) :: int_eff
    class(workspace_t), intent(inout), allocatable :: tmp
    integer :: n_in
    n_in = interaction_get_n_in (int_eff)
    call interaction_set_momenta (int_eff, p_seed(1:n_in), outgoing = .false.)
    p_seed(n_in+1:) = interaction_get_momenta (int_eff, outgoing = .true.)
  end subroutine prc_template_me_recover_kinematics
    
\eatline
\nwendcode{}\nwbegindocs{171}Compute the amplitude.  For the tree-level process, we can ignore the scale
settings.  The term index {\tt{}j} is also irrelevant.

We first call {\tt{}new{\char95}event} for the given momenta (which we must unpack), then
retrieve the amplitude value for the given quantum numbers.

If the {\tt{}tmp} status flag is present, we can make sure that we call
{\tt{}new{\char95}event} only once for a given kinematics.  After the first call, we
unset the {\tt{}new{\char95}kinematics} flag.
\nwenddocs{}\nwbegincode{172}\moddef{Template matrix elements: prc template ME: TBP}\plusendmoddef
  procedure :: compute_amplitude => prc_template_me_compute_amplitude
\nwendcode{}\nwbegincode{173}\moddef{Template matrix elements: procedures}\plusendmoddef
  function prc_template_me_compute_amplitude &
       (object, j, p, f, h, c, fac_scale, ren_scale, alpha_qcd_forced, tmp) &
       result (amp)
    class(prc_template_me_t), intent(in) :: object
    integer, intent(in) :: j
    type(vector4_t), dimension(:), intent(in) :: p
    integer, intent(in) :: f, h, c
    real(default), intent(in) :: fac_scale, ren_scale
    real(default), intent(in), allocatable :: alpha_qcd_forced
    class(workspace_t), intent(inout), allocatable, optional :: tmp
    complex(default) :: amp
    integer :: n_tot, i
    real(c_default_float), dimension(:,:), allocatable :: parray
    complex(c_default_complex) :: camp
    logical :: new_event
    select type (driver => object%driver)
    type is (template_me_driver_t)
       new_event = .true.
       if (present (tmp)) then
          if (allocated (tmp)) then
             select type (tmp)
             type is (template_me_state_t)
                new_event = tmp%new_kinematics
                tmp%new_kinematics = .false.
             end select
          end if
       end if
       if (new_event) then
          n_tot = object%data%n_in + object%data%n_out
          allocate (parray (0:3, n_tot))
          forall (i = 1:n_tot)  parray(:,i) = vector4_get_components (p(i))
          call driver%new_event (parray)
       end if
       if (object%is_allowed (1, f, h, c)) then
          call driver%get_amplitude &
               (int (f, c_int), int (h, c_int), int (c, c_int), camp)
          amp = camp
       else
          amp = 0
       end if
    end select
  end function prc_template_me_compute_amplitude
    
\eatline
\nwendcode{}\nwbegindocs{174}We do not overwrite the {\tt{}prc{\char95}core{\char95}t} routine for $\alpha_s$. 

\subsection{Test}
This is the master for calling self-test procedures. There is a
separate test for testing \oMega\ diagram generation as this depends on
a working analysis setup.
\nwenddocs{}\nwbegincode{175}\moddef{Template matrix elements: public}\plusendmoddef
  public :: prc_template_me_test
\nwendcode{}\nwbegincode{176}\moddef{Template matrix elements: tests}\plusendmoddef
  subroutine prc_template_me_test (u, results)
    integer, intent(in) :: u
    type(test_results_t), intent(inout) :: results
  \LA{}Template matrix elements: execute tests\RA{}
end subroutine prc_template_me_test

\eatline
\nwendcode{}\nwbegindocs{177}\nwdocspar
\subsubsection{Generate, compile and load a simple process matrix element}
The process is $e^+ e^- \to \mu^+\mu^-$ for vanishing masses and
$e=0.3$.  We initialize the process, build the library, and compute a
particular matrix element for momenta of unit energy and right-angle
scattering.  The matrix element, as it happens, is equal to $e^2$.
(Note that are no conversion factors applied, so this result is
exact.)

For {\tt{}GNU\ make}, {\tt{}makeflags} is set to {\tt{}-j1}.  This eliminates a
potential clash with a {\tt{}-j<n>} flag if this test is called from a
parallel make.
\nwenddocs{}\nwbegincode{178}\moddef{Template matrix elements: execute tests}\endmoddef
  call test (prc_template_me_1, "prc_template_me_1", &
       "build and load simple template process", &
       u, results)
\nwendcode{}\nwbegincode{179}\moddef{Template matrix elements: tests}\plusendmoddef
  subroutine prc_template_me_1 (u)
    integer, intent(in) :: u
    type(process_library_t) :: lib
    class(prc_core_def_t), allocatable :: def
    type(process_def_entry_t), pointer :: entry
    type(os_data_t) :: os_data
    class(model_data_t), pointer :: model
    type(string_t) :: model_name
    type(string_t), dimension(:), allocatable :: prt_in, prt_out
    type(process_constants_t) :: data
    class(prc_core_driver_t), allocatable :: driver
    integer, parameter :: cdf = c_default_float
    integer, parameter :: ci = c_int
    real(cdf), dimension(4) :: par
    real(cdf), dimension(0:3,4) :: p
    logical(c_bool) :: flag
    complex(c_default_complex) :: amp
    integer :: i
    
    write (u, "(A)")  "* Test output: prc_template_me_1"
    write (u, "(A)")  "*   Purpose: create a template matrix element,"
    write (u, "(A)")  "*            normalized to give unit integral,"    
    write (u, "(A)")  "*            build a library, link, load, and &
         &access the matrix element"
    write (u, "(A)")

    write (u, "(A)")  "* Initialize a process library with one entry"
    write (u, "(A)")
    call lib%init (var_str ("template_me1"))
    call os_data_init (os_data)

    model_name = "QED"
    call prepare_model (model, model_name)

    allocate (prt_in (2), prt_out (2))
    prt_in = [var_str ("e+"), var_str ("e-")]
    prt_out = [var_str ("m+"), var_str ("m-")]
    
    allocate (template_me_def_t :: def)
    select type (def)
    type is (template_me_def_t)
       call def%init (model, prt_in, prt_out, unity = .false.)
    end select
    allocate (entry)
    call entry%init (var_str ("template_me1_a"), model_name = model_name, &
         n_in = 2, n_components = 1)
    call entry%import_component (1, n_out = size (prt_out), &
         prt_in  = new_prt_spec (prt_in), &
         prt_out = new_prt_spec (prt_out), &
         method  = var_str ("template"), &
         variant = def)
    call lib%append (entry)
    
    write (u, "(A)")  "* Configure library"
    write (u, "(A)")
    call lib%configure (os_data)
    
    write (u, "(A)")  "* Write makefile"
    write (u, "(A)")
    call lib%write_makefile (os_data, force = .true.)

    write (u, "(A)")  "* Clean any left-over files"
    write (u, "(A)")
    call lib%clean (os_data, distclean = .false.)

    write (u, "(A)")  "* Write driver"
    write (u, "(A)")
    call lib%write_driver (force = .true.)

    write (u, "(A)")  "* Write process source code, compile, link, load"
    write (u, "(A)")
    call lib%load (os_data)

    call lib%write (u, libpath = .false.)
    
    write (u, "(A)")
    write (u, "(A)")  "* Probe library API:"
    write (u, "(A)")
       
    write (u, "(1x,A,L1)")  "is active                 = ", &
         lib%is_active ()
    write (u, "(1x,A,I0)")  "n_processes               = ", &
         lib%get_n_processes ()

    write (u, "(A)")
    write (u, "(A)")  "* Constants of template_me1_a_i1:"
    write (u, "(A)")

    call lib%connect_process (var_str ("template_me1_a"), 1, data, driver)

    write (u, "(1x,A,A)")  "component ID     = ", char (data%id)
    write (u, "(1x,A,A)")  "model name       = ", char (data%model_name)
    write (u, "(1x,A,A,A)")  "md5sum           = '", data%md5sum, "'"
    write (u, "(1x,A,I0)") "n_in  = ", data%n_in
    write (u, "(1x,A,I0)") "n_out = ", data%n_out
    write (u, "(1x,A,I0)") "n_flv = ", data%n_flv
    write (u, "(1x,A,I0)") "n_hel = ", data%n_hel
    write (u, "(1x,A,I0)") "n_col = ", data%n_col
    write (u, "(1x,A,I0)") "n_cin = ", data%n_cin
    write (u, "(1x,A,I0)") "n_cf  = ", data%n_cf
    write (u, "(1x,A,10(1x,I0))") "flv state =", data%flv_state
    write (u, "(1x,A,10(1x,I2))") "hel state =", data%hel_state(:,1)
    do i = 2, 16
       write (u, "(12x,4(1x,I2))")  data%hel_state(:,i)
    end do
    write (u, "(1x,A,10(1x,I0))") "col state =", data%col_state
    write (u, "(1x,A,10(1x,L1))") "ghost flag =", data%ghost_flag
    write (u, "(1x,A,10(1x,F5.3))") "color factors =", data%color_factors
    write (u, "(1x,A,10(1x,I0))") "cf index =", data%cf_index

    write (u, "(A)")
    write (u, "(A)")  "* Set parameters for template_me1_a and initialize:"
    write (u, "(A)")

    par = [0.3_cdf, 0.0_cdf, 0.0_cdf, 0.0_cdf]
    write (u, "(2x,A,F6.4)")  "ee   = ", par(1)
    write (u, "(2x,A,F6.4)")  "me   = ", par(2)
    write (u, "(2x,A,F6.4)")  "mmu  = ", par(3)
    write (u, "(2x,A,F6.4)")  "mtau = ", par(4)

    write (u, "(A)")
    write (u, "(A)")  "* Set kinematics:"
    write (u, "(A)")
    
    p = reshape ([ &
         1.0_cdf, 0.0_cdf, 0.0_cdf, 1.0_cdf, &
         1.0_cdf, 0.0_cdf, 0.0_cdf,-1.0_cdf, &
         1.0_cdf, 1.0_cdf, 0.0_cdf, 0.0_cdf, &
         1.0_cdf,-1.0_cdf, 0.0_cdf, 0.0_cdf &
         ], [4,4])
    do i = 1, 4
       write (u, "(2x,A,I0,A,4(1x,F7.4))")  "p", i, " =", p(:,i)
    end do

    select type (driver)
    type is (template_me_driver_t)
       call driver%init (par)

       call driver%new_event (p)

       write (u, "(A)")
       write (u, "(A)")  "* Compute matrix element:"
       write (u, "(A)")

       call driver%is_allowed (1_ci, 6_ci, 1_ci, flag)
       write (u, "(1x,A,L1)") "is_allowed (1, 6, 1) = ", flag
       
       call driver%get_amplitude (1_ci, 6_ci, 1_ci, amp)
       write (u, "(1x,A,1x,E11.4)") "|amp (1, 6, 1)| =", abs (amp)
    end select

    call lib%final ()
    call cleanup_model (model)
    
    write (u, "(A)")
    write (u, "(A)")  "* Test output end: prc_template_me_1"
    
  end subroutine prc_template_me_1
  
\eatline
\nwendcode{}\nwbegindocs{180}\nwdocspar
\nwenddocs{}\nwbegincode{181}\moddef{Template matrix elements: execute tests}\plusendmoddef
  call test (prc_template_me_2, "prc_template_me_2", &
       "build and load simple template_unity process", &
       u, results)
\nwendcode{}\nwbegincode{182}\moddef{Template matrix elements: tests}\plusendmoddef
  subroutine prc_template_me_2 (u)
    integer, intent(in) :: u
    type(process_library_t) :: lib
    class(prc_core_def_t), allocatable :: def
    type(process_def_entry_t), pointer :: entry
    type(os_data_t) :: os_data
    class(model_data_t), pointer :: model
    type(string_t) :: model_name
    type(string_t), dimension(:), allocatable :: prt_in, prt_out
    type(process_constants_t) :: data
    class(prc_core_driver_t), allocatable :: driver
    integer, parameter :: cdf = c_default_float
    integer, parameter :: ci = c_int
    real(cdf), dimension(4) :: par
    real(cdf), dimension(0:3,4) :: p
    logical(c_bool) :: flag
    complex(c_default_complex) :: amp
    integer :: i
    
    write (u, "(A)")  "* Test output: prc_template_me_1"
    write (u, "(A)")  "*   Purpose: create a template matrix element,"
    write (u, "(A)")  "*            being identical to unity,"    
    write (u, "(A)")  "*            build a library, link, load, and &
         &access the matrix element"
    write (u, "(A)")

    write (u, "(A)")  "* Initialize a process library with one entry"
    write (u, "(A)")
    call lib%init (var_str ("template_me2"))
    call os_data_init (os_data)
    
    model_name = "QED"
    call prepare_model (model, model_name)

    allocate (prt_in (2), prt_out (2))
    prt_in = [var_str ("e+"), var_str ("e-")]
    prt_out = [var_str ("m+"), var_str ("m-")]
    
    allocate (template_me_def_t :: def)
    select type (def)
    type is (template_me_def_t)
       call def%init (model, prt_in, prt_out, unity = .true.)
    end select
    allocate (entry)
    call entry%init (var_str ("template_me2_a"), model_name = model_name, &
         n_in = 2, n_components = 1)
    call entry%import_component (1, n_out = size (prt_out), &
         prt_in  = new_prt_spec (prt_in), &
         prt_out = new_prt_spec (prt_out), &
         method  = var_str ("template_unity"), &
         variant = def)
    call lib%append (entry)
    
    write (u, "(A)")  "* Configure library"
    write (u, "(A)")
    call lib%configure (os_data)
    
    write (u, "(A)")  "* Write makefile"
    write (u, "(A)")
    call lib%write_makefile (os_data, force = .true.)

    write (u, "(A)")  "* Clean any left-over files"
    write (u, "(A)")
    call lib%clean (os_data, distclean = .false.)

    write (u, "(A)")  "* Write driver"
    write (u, "(A)")
    call lib%write_driver (force = .true.)

    write (u, "(A)")  "* Write process source code, compile, link, load"
    write (u, "(A)")
    call lib%load (os_data)

    call lib%write (u, libpath = .false.)
    
    write (u, "(A)")
    write (u, "(A)")  "* Probe library API:"
    write (u, "(A)")
       
    write (u, "(1x,A,L1)")  "is active                 = ", &
         lib%is_active ()
    write (u, "(1x,A,I0)")  "n_processes               = ", &
         lib%get_n_processes ()

    write (u, "(A)")
    write (u, "(A)")  "* Constants of template_me2_a_i1:"
    write (u, "(A)")

    call lib%connect_process (var_str ("template_me2_a"), 1, data, driver)

    write (u, "(1x,A,A)")  "component ID     = ", char (data%id)
    write (u, "(1x,A,A)")  "model name       = ", char (data%model_name)
    write (u, "(1x,A,A,A)")  "md5sum           = '", data%md5sum, "'"
    write (u, "(1x,A,I0)") "n_in  = ", data%n_in
    write (u, "(1x,A,I0)") "n_out = ", data%n_out
    write (u, "(1x,A,I0)") "n_flv = ", data%n_flv
    write (u, "(1x,A,I0)") "n_hel = ", data%n_hel
    write (u, "(1x,A,I0)") "n_col = ", data%n_col
    write (u, "(1x,A,I0)") "n_cin = ", data%n_cin
    write (u, "(1x,A,I0)") "n_cf  = ", data%n_cf
    write (u, "(1x,A,10(1x,I0))") "flv state =", data%flv_state
    write (u, "(1x,A,10(1x,I2))") "hel state =", data%hel_state(:,1)
    do i = 2, 16
       write (u, "(12x,4(1x,I2))")  data%hel_state(:,i)
    end do
    write (u, "(1x,A,10(1x,I0))") "col state =", data%col_state
    write (u, "(1x,A,10(1x,L1))") "ghost flag =", data%ghost_flag
    write (u, "(1x,A,10(1x,F5.3))") "color factors =", data%color_factors
    write (u, "(1x,A,10(1x,I0))") "cf index =", data%cf_index

    write (u, "(A)")
    write (u, "(A)")  "* Set parameters for template_me2_a and initialize:"
    write (u, "(A)")

    par = [0.3_cdf, 0.0_cdf, 0.0_cdf, 0.0_cdf]
    write (u, "(2x,A,F6.4)")  "ee   = ", par(1)
    write (u, "(2x,A,F6.4)")  "me   = ", par(2)
    write (u, "(2x,A,F6.4)")  "mmu  = ", par(3)
    write (u, "(2x,A,F6.4)")  "mtau = ", par(4)

    write (u, "(A)")
    write (u, "(A)")  "* Set kinematics:"
    write (u, "(A)")
    
    p = reshape ([ &
         1.0_cdf, 0.0_cdf, 0.0_cdf, 1.0_cdf, &
         1.0_cdf, 0.0_cdf, 0.0_cdf,-1.0_cdf, &
         1.0_cdf, 1.0_cdf, 0.0_cdf, 0.0_cdf, &
         1.0_cdf,-1.0_cdf, 0.0_cdf, 0.0_cdf &
         ], [4,4])
    do i = 1, 4
       write (u, "(2x,A,I0,A,4(1x,F7.4))")  "p", i, " =", p(:,i)
    end do

    select type (driver)
    type is (template_me_driver_t)
       call driver%init (par)

       call driver%new_event (p)

       write (u, "(A)")
       write (u, "(A)")  "* Compute matrix element:"
       write (u, "(A)")

       call driver%is_allowed (1_ci, 6_ci, 1_ci, flag)
       write (u, "(1x,A,L1)") "is_allowed (1, 6, 1) = ", flag
       
       call driver%get_amplitude (1_ci, 6_ci, 1_ci, amp)
       write (u, "(1x,A,1x,E11.4)") "|amp (1, 6, 1)| =", abs (amp)
    end select

    call lib%final ()
    call cleanup_model (model)
    
    write (u, "(A)")
    write (u, "(A)")  "* Test output end: prc_template_me_2"
    
  end subroutine prc_template_me_2
  
\eatline
\nwendcode{}\nwbegindocs{183}\nwdocspar
%%%%%%%%%%%%%%%%%%%%%%%%%%%%%%%%%%%%%%%%%%%%%%%%%%%%%%%%%%%%%%%%%%%%%%%%
\section{\oMega\ Interface}

The standard method for process computation with \whizard\ is the
\oMega\ matrix element generator.

This section implements the interface to the code generator (via
the makefile) and the driver for the features provided by the \oMega\
matrix element.

There are actually two different methods steered by this interface, the
traditional one which delivers compiled Fortran code, while the \oMega\ 
virtual machine (OVM) produces bytecode with look-up tables.
\nwenddocs{}\nwbegincode{184}\moddef{\code{}prc{\char95}omega.f90\edoc{}}\endmoddef
\LA{}File header\RA{}

module prc_omega
  
  use, intrinsic :: iso_c_binding !NODEP!
  
  use kinds
\LA{}Use strings\RA{}
  use io_units
  use file_utils, only: delete_file
  use unit_tests
  use system_defs, only: TAB
  use diagnostics
  use os_interface
  use lorentz
  use sm_qcd
  use interactions
  use model_data
  use var_base
  use model_testbed

  use process_constants
  use prclib_interfaces
  use prc_core_def
  use particle_specifiers
  use process_libraries
  use prc_core

\LA{}Standard module head\RA{}

\LA{}Omega interface: public\RA{}

\LA{}Omega interface: types\RA{}

\LA{}Omega interface: interfaces\RA{}

contains
  
\LA{}Omega interface: procedures\RA{}

\LA{}Omega interface: tests\RA{}

end module prc_omega
\eatline
\nwendcode{}\nwbegindocs{185}\nwdocspar
\subsection{Process definition}
For the process definition we implement an extension of the
{\tt{}prc{\char95}core{\char95}def{\char95}t} abstract type.
\nwenddocs{}\nwbegincode{186}\moddef{Omega interface: public}\endmoddef
  public :: omega_def_t
\nwendcode{}\nwbegincode{187}\moddef{Omega interface: types}\endmoddef
  type, abstract, extends (prc_core_def_t) :: omega_def_t
   contains
   \LA{}Omega interface: omega def: TBP\RA{}
  end type omega_def_t
  
\eatline
\nwendcode{}\nwbegindocs{188}\nwdocspar
\nwenddocs{}\nwbegincode{189}\moddef{Omega interface: public}\plusendmoddef
  public :: omega_omega_def_t
\nwendcode{}\nwbegincode{190}\moddef{Omega interface: types}\plusendmoddef
  type, extends (omega_def_t) :: omega_omega_def_t
   contains
   \LA{}Omega interface: omega omega def: TBP\RA{}
  end type omega_omega_def_t
  
\eatline
\nwendcode{}\nwbegindocs{191}\nwdocspar
\nwenddocs{}\nwbegincode{192}\moddef{Omega interface: public}\plusendmoddef
  public :: omega_ovm_def_t
\nwendcode{}\nwbegincode{193}\moddef{Omega interface: types}\plusendmoddef
  type, extends (omega_def_t) :: omega_ovm_def_t
   contains
   \LA{}Omega interface: omega ovm def: TBP\RA{}
  end type omega_ovm_def_t
  
\eatline
\nwendcode{}\nwbegindocs{194}\nwdocspar
\nwenddocs{}\nwbegincode{195}\moddef{Omega interface: omega omega def: TBP}\endmoddef
  procedure, nopass :: type_string => omega_omega_def_type_string
\nwendcode{}\nwbegincode{196}\moddef{Omega interface: procedures}\endmoddef
  function omega_omega_def_type_string () result (string)
    type(string_t) :: string
    string = "omega"
  end function omega_omega_def_type_string

\eatline
\nwendcode{}\nwbegindocs{197}\nwdocspar
\nwenddocs{}\nwbegincode{198}\moddef{Omega interface: omega ovm def: TBP}\endmoddef
  procedure, nopass :: type_string => omega_ovm_def_type_string
\nwendcode{}\nwbegincode{199}\moddef{Omega interface: procedures}\plusendmoddef
  function omega_ovm_def_type_string () result (string)
    type(string_t) :: string
    string = "ovm"
  end function omega_ovm_def_type_string

\eatline
\nwendcode{}\nwbegindocs{200}Initialization: allocate the writer for the \oMega\ matrix element.
Also set any data for this process that the writer needs.
\nwenddocs{}\nwbegincode{201}\moddef{Omega interface: omega def: TBP}\endmoddef
  procedure :: init => omega_def_init
\nwendcode{}\nwbegincode{202}\moddef{Omega interface: procedures}\plusendmoddef
  subroutine omega_def_init (object, model_name, prt_in, prt_out, &
       restrictions, openmp_support, report_progress, extra_options, &
       diags, diags_color)
    class(omega_def_t), intent(out) :: object
    type(string_t), intent(in) :: model_name
    type(string_t), dimension(:), intent(in) :: prt_in
    type(string_t), dimension(:), intent(in) :: prt_out
    type(string_t), intent(in), optional :: restrictions
    logical, intent(in), optional :: openmp_support
    logical, intent(in), optional :: report_progress
    logical, intent(in), optional :: diags, diags_color
    type(string_t), intent(in), optional :: extra_options
    select type (object)
    type is (omega_omega_def_t)
       allocate (omega_omega_writer_t :: object%writer)
       select type (writer => object%writer)
       type is (omega_omega_writer_t)
          call writer%init (model_name, prt_in, prt_out, &
               restrictions, openmp_support, report_progress, &
               extra_options, diags, diags_color)
       end select
    type is (omega_ovm_def_t)
       allocate (omega_ovm_writer_t :: object%writer)
       select type (writer => object%writer)
       type is (omega_ovm_writer_t)
          call writer%init (model_name, prt_in, prt_out, &
               restrictions, openmp_support, report_progress, &
               extra_options, diags, diags_color)
       end select       
    end select
  end subroutine omega_def_init

\eatline
\nwendcode{}\nwbegindocs{203}Write/read process- and method-specific data.
\nwenddocs{}\nwbegincode{204}\moddef{Omega interface: omega def: TBP}\plusendmoddef
  procedure :: write => omega_def_write
\nwendcode{}\nwbegincode{205}\moddef{Omega interface: procedures}\plusendmoddef
  subroutine omega_def_write (object, unit)
    class(omega_def_t), intent(in) :: object
    integer, intent(in) :: unit
    select type (writer => object%writer)
    type is (omega_omega_writer_t)
       call writer%write (unit)
    type is (omega_ovm_writer_t)
       call writer%write (unit)
    end select
  end subroutine omega_def_write
  
\eatline
\nwendcode{}\nwbegindocs{206}\nwdocspar
\nwenddocs{}\nwbegincode{207}\moddef{Omega interface: omega def: TBP}\plusendmoddef
  procedure :: read => omega_def_read
\nwendcode{}\nwbegincode{208}\moddef{Omega interface: procedures}\plusendmoddef
  subroutine omega_def_read (object, unit)
    class(omega_def_t), intent(out) :: object
    integer, intent(in) :: unit
    call msg_bug ("O'Mega process definition: input not supported yet")
  end subroutine omega_def_read
  
\eatline
\nwendcode{}\nwbegindocs{209}Allocate the driver for \oMega matrix elements.
\nwenddocs{}\nwbegincode{210}\moddef{Omega interface: omega def: TBP}\plusendmoddef
  procedure :: allocate_driver => omega_def_allocate_driver
\nwendcode{}\nwbegincode{211}\moddef{Omega interface: procedures}\plusendmoddef
  subroutine omega_def_allocate_driver (object, driver, basename)
    class(omega_def_t), intent(in) :: object
    class(prc_core_driver_t), intent(out), allocatable :: driver
    type(string_t), intent(in) :: basename
    allocate (omega_driver_t :: driver)
  end subroutine omega_def_allocate_driver
  
\eatline
\nwendcode{}\nwbegindocs{212}We need code:
\nwenddocs{}\nwbegincode{213}\moddef{Omega interface: omega def: TBP}\plusendmoddef
  procedure, nopass :: needs_code => omega_def_needs_code
\nwendcode{}\nwbegincode{214}\moddef{Omega interface: procedures}\plusendmoddef
  function omega_def_needs_code () result (flag)
    logical :: flag
    flag = .true.
  end function omega_def_needs_code
  
\eatline
\nwendcode{}\nwbegindocs{215}These are the features that an \oMega\ matrix element provides.
\nwenddocs{}\nwbegincode{216}\moddef{Omega interface: omega def: TBP}\plusendmoddef
  procedure, nopass :: get_features => omega_def_get_features
\nwendcode{}\nwbegincode{217}\moddef{Omega interface: procedures}\plusendmoddef
  subroutine omega_def_get_features (features)
    type(string_t), dimension(:), allocatable, intent(out) :: features
    allocate (features (6))
    features = [ &
         var_str ("init"), &
         var_str ("update_alpha_s"), &
         var_str ("reset_helicity_selection"), &
         var_str ("is_allowed"), &
         var_str ("new_event"), &
         var_str ("get_amplitude")]
  end subroutine omega_def_get_features

\eatline
\nwendcode{}\nwbegindocs{218}The interface of the specific features.
\nwenddocs{}\nwbegincode{219}\moddef{Omega interface: interfaces}\endmoddef
  abstract interface
     subroutine init_t (par) bind(C)
       import
       real(c_default_float), dimension(*), intent(in) :: par
     end subroutine init_t
  end interface
  
  abstract interface
     subroutine update_alpha_s_t (alpha_s) bind(C)
       import
       real(c_default_float), intent(in) :: alpha_s
     end subroutine update_alpha_s_t
  end interface
  
  abstract interface
     subroutine reset_helicity_selection_t (threshold, cutoff) bind(C)
       import
       real(c_default_float), intent(in) :: threshold
       integer(c_int), intent(in) :: cutoff
     end subroutine reset_helicity_selection_t
  end interface

  abstract interface
     subroutine is_allowed_t (flv, hel, col, flag) bind(C)
       import
       integer(c_int), intent(in) :: flv, hel, col
       logical(c_bool), intent(out) :: flag
     end subroutine is_allowed_t
  end interface

  abstract interface
     subroutine new_event_t (p) bind(C)
       import
       real(c_default_float), dimension(0:3,*), intent(in) :: p
     end subroutine new_event_t
  end interface
  
  abstract interface
     subroutine get_amplitude_t (flv, hel, col, amp) bind(C)
       import
       integer(c_int), intent(in) :: flv, hel, col
       complex(c_default_complex), intent(out):: amp
     end subroutine get_amplitude_t
  end interface

\eatline
\eatline
\nwendcode{}\nwbegindocs{220}Connect the \oMega\ features with the process driver.
\nwenddocs{}\nwbegincode{221}\moddef{Omega interface: omega def: TBP}\plusendmoddef
  procedure :: connect => omega_def_connect
\nwendcode{}\nwbegincode{222}\moddef{Omega interface: procedures}\plusendmoddef
  subroutine omega_def_connect (def, lib_driver, i, proc_driver)
    class(omega_def_t), intent(in) :: def
    class(prclib_driver_t), intent(in) :: lib_driver
    integer, intent(in) :: i
    class(prc_core_driver_t), intent(inout) :: proc_driver
    integer(c_int) :: pid, fid
    type(c_funptr) :: fptr
    select type (proc_driver)
    type is  (omega_driver_t)
       pid = i
       fid = 1
       call lib_driver%get_fptr (pid, fid, fptr)
       call c_f_procpointer (fptr, proc_driver%init)
       fid = 2
       call lib_driver%get_fptr (pid, fid, fptr)
       call c_f_procpointer (fptr, proc_driver%update_alpha_s)
       fid = 3
       call lib_driver%get_fptr (pid, fid, fptr)
       call c_f_procpointer (fptr, proc_driver%reset_helicity_selection)
       fid = 4
       call lib_driver%get_fptr (pid, fid, fptr)
       call c_f_procpointer (fptr, proc_driver%is_allowed)
       fid = 5
       call lib_driver%get_fptr (pid, fid, fptr)
       call c_f_procpointer (fptr, proc_driver%new_event)
       fid = 6
       call lib_driver%get_fptr (pid, fid, fptr)
       call c_f_procpointer (fptr, proc_driver%get_amplitude)
    end select
  end subroutine omega_def_connect

\eatline
\nwendcode{}\nwbegindocs{223}\nwdocspar
\subsection{The \oMega\ writer}
The \oMega\ writer is responsible for inserting the appropriate lines
in the makefile that call \oMega, and for writing interfaces and
wrappers.
\nwenddocs{}\nwbegincode{224}\moddef{Omega interface: types}\plusendmoddef
  type, extends (prc_writer_f_module_t), abstract :: omega_writer_t
     type(string_t) :: model_name
     type(string_t) :: process_mode
     type(string_t) :: process_string
     type(string_t) :: restrictions
     logical :: openmp_support = .false.
     logical :: report_progress = .false.
     logical :: diags = .false.
     logical :: diags_color = .false.
     type(string_t) :: extra_options
   contains
   \LA{}Omega interface: omega writer: TBP\RA{}
  end type omega_writer_t

\eatline
\nwendcode{}\nwbegindocs{225}\nwdocspar
\nwenddocs{}\nwbegincode{226}\moddef{Omega interface: types}\plusendmoddef
  type, extends (omega_writer_t) :: omega_omega_writer_t
   contains
   \LA{}Omega interface: omega omega writer: TBP\RA{}
  end type omega_omega_writer_t

\eatline
\nwendcode{}\nwbegindocs{227}\nwdocspar
\nwenddocs{}\nwbegincode{228}\moddef{Omega interface: types}\plusendmoddef
  type, extends (omega_writer_t) :: omega_ovm_writer_t
   contains
   \LA{}Omega interface: omega ovm writer: TBP\RA{}
  end type omega_ovm_writer_t

\eatline
\nwendcode{}\nwbegincode{229}\moddef{Omega interface: omega omega writer: TBP}\endmoddef
  procedure, nopass :: type_name => omega_omega_writer_type_name
\nwendcode{}\nwbegincode{230}\moddef{Omega interface: procedures}\plusendmoddef
  function omega_omega_writer_type_name () result (string)
    type(string_t) :: string
    string = "omega"
  end function omega_omega_writer_type_name

\eatline
\nwendcode{}\nwbegindocs{231}\nwdocspar
\nwenddocs{}\nwbegincode{232}\moddef{Omega interface: omega ovm writer: TBP}\endmoddef
  procedure, nopass :: type_name => omega_ovm_writer_type_name
\nwendcode{}\nwbegincode{233}\moddef{Omega interface: procedures}\plusendmoddef
  function omega_ovm_writer_type_name () result (string)
    type(string_t) :: string
    string = "ovm"
  end function omega_ovm_writer_type_name

\eatline
\nwendcode{}\nwbegindocs{234}\nwdocspar
\nwenddocs{}\nwbegindocs{235}Taking into account the prefix for \oMega\ module names.
\nwenddocs{}\nwbegincode{236}\moddef{Omega interface: omega writer: TBP}\endmoddef
  procedure, nopass :: get_module_name => omega_writer_get_module_name
\nwendcode{}\nwbegincode{237}\moddef{Omega interface: procedures}\plusendmoddef
  function omega_writer_get_module_name (id) result (name)
    type(string_t) :: name
    type(string_t), intent(in) :: id
    name = "opr_" // id
  end function omega_writer_get_module_name

\eatline
\nwendcode{}\nwbegindocs{238}Output.  This is called by {\tt{}omega{\char95}def{\char95}write}.
\nwenddocs{}\nwbegincode{239}\moddef{Omega interface: omega writer: TBP}\plusendmoddef
  procedure :: write => omega_writer_write
\nwendcode{}\nwbegincode{240}\moddef{Omega interface: procedures}\plusendmoddef
  subroutine omega_writer_write (object, unit)
    class(omega_writer_t), intent(in) :: object
    integer, intent(in) :: unit
    write (unit, "(5x,A,A)")  "Model name        = ", &
         '"' // char (object%model_name) // '"'
    write (unit, "(5x,A,A)")  "Mode string       = ", &
         '"' // char (object%process_mode) // '"'
    write (unit, "(5x,A,A)")  "Process string    = ", &
         '"' // char (object%process_string) // '"'
    write (unit, "(5x,A,A)")  "Restrictions      = ", &
         '"' // char (object%restrictions) // '"'
    write (unit, "(5x,A,L1)")  "OpenMP support    = ", object%openmp_support
    write (unit, "(5x,A,L1)")  "Report progress   = ", object%report_progress
    write (unit, "(5x,A,A)")  "Extra options     = ", &
         '"' // char (object%extra_options) // '"'
    write (unit, "(5x,A,L1)")  "Write diagrams    = ", object%diags    
    write (unit, "(5x,A,L1)")  "Write color diag. = ", object%diags_color
  end subroutine omega_writer_write

\eatline
\nwendcode{}\nwbegindocs{241}Initialize with process data.
\nwenddocs{}\nwbegincode{242}\moddef{Omega interface: omega writer: TBP}\plusendmoddef
  procedure :: init => omega_writer_init
\nwendcode{}\nwbegincode{243}\moddef{Omega interface: procedures}\plusendmoddef
  subroutine omega_writer_init (writer, model_name, prt_in, prt_out, &
       restrictions, openmp_support, report_progress, extra_options, &
       diags, diags_color)
    class(omega_writer_t), intent(out) :: writer
    type(string_t), intent(in) :: model_name
    type(string_t), dimension(:), intent(in) :: prt_in
    type(string_t), dimension(:), intent(in) :: prt_out
    type(string_t), intent(in), optional :: restrictions
    logical, intent(in), optional :: openmp_support
    logical, intent(in), optional :: report_progress
    logical, intent(in), optional :: diags, diags_color    
    type(string_t), intent(in), optional :: extra_options
    integer :: i
    writer%model_name = model_name
    if (present (restrictions)) then
       writer%restrictions = restrictions
    else
       writer%restrictions = ""
    end if
    if (present (openmp_support))  writer%openmp_support = openmp_support
    if (present (report_progress))  writer%report_progress = report_progress
    if (present (diags))  writer%diags = diags
    if (present (diags_color))  writer%diags_color = diags_color
    if (present (extra_options)) then
       writer%extra_options = " " // extra_options
    else
       writer%extra_options = ""
    end if
    select case (size (prt_in))
    case (1);  writer%process_mode = " -decay"
    case (2);  writer%process_mode = " -scatter"
    end select
    associate (s => writer%process_string)
      s = " '"
      do i = 1, size (prt_in)
         if (i > 1)  s = s // " "
         s = s // prt_in(i)
      end do
      s = s // " ->"
      do i = 1, size (prt_out)
         s = s // " " // prt_out(i)
      end do
      s = s // "'"
    end associate
  end subroutine omega_writer_init

\eatline
\nwendcode{}\nwbegindocs{244}The makefile implements the actual \oMega\ call. For old \LaTeX\
distributions, we filter out the hyperref options for \oMega\
diagrams, at least in the testsuite.
\nwenddocs{}\nwbegincode{245}\moddef{Omega interface: omega writer: TBP}\plusendmoddef
  procedure :: write_makefile_code => omega_write_makefile_code
\nwendcode{}\nwbegincode{246}\moddef{Omega interface: procedures}\plusendmoddef
  subroutine omega_write_makefile_code (writer, unit, id, os_data, testflag)
    class(omega_writer_t), intent(in) :: writer
    integer, intent(in) :: unit
    type(string_t), intent(in) :: id
    type(os_data_t), intent(in) :: os_data
    logical, intent(in), optional :: testflag
    type(string_t) :: omega_binary, omega_path
    type(string_t) :: restrictions_string
    type(string_t) :: openmp_string
    type(string_t) :: kmatrix_string
    type(string_t) :: progress_string    
    type(string_t) :: diagrams_string
    logical :: escape_hyperref
    escape_hyperref = .false.
    if (present (testflag))  escape_hyperref = testflag
    select type (writer)
    type is (omega_omega_writer_t)       
       omega_binary = "omega_" // writer%model_name // ".opt"
    type is (omega_ovm_writer_t)
       select case (char (writer%model_name))
       case ("SM", "SM_CKM", "SM_Higgs", "2HDM", "2HDM_CKM", &
             "HSExt", "QED", "QCD", "Zprime")
       case default
          call msg_fatal ("The model " // char (writer%model_name) &
               // " is not available for the O'Mega VM.")
       end select
       omega_binary = "omega_" // writer%model_name // "_VM.opt"
    end select
    omega_path = os_data%whizard_omega_binpath // "/" // omega_binary
    if (writer%restrictions /= "") then
       restrictions_string = " -cascade '" // writer%restrictions // "'"
    else
       restrictions_string = ""
    end if
    if (writer%openmp_support) then
       openmp_string = " -target:openmp"
    else
       openmp_string = ""
    end if
    if (writer%report_progress) then
       progress_string = " -fusion:progress"
    else
       progress_string = ""
    end if
    if (writer%diags) then
       if (writer%diags_color) then
          diagrams_string = " -diagrams:C " // char(id) // &
               "_diags -diagrams_LaTeX"
       else
          diagrams_string = " -diagrams " // char(id) // &
               "_diags -diagrams_LaTeX"
       end if
    else 
       if (writer%diags_color) then
          diagrams_string = " -diagrams:c " // char(id) // &
               "_diags -diagrams_LaTeX"        
       else
          diagrams_string = ""
       end if
    end if
    select case (char (writer%model_name))
    case ("SM_rx", "SSC", "NoH_rx", "AltH")
       kmatrix_string = " -target:kmatrix_write"
    case default
       kmatrix_string = ""
    end select
    write (unit, "(5A)")  "SOURCES += ", char (id), ".f90"
    select type (writer)
    type is (omega_ovm_writer_t)
       write (unit, "(5A)")  "SOURCES += ", char (id), ".hbc"
    end select
    if (writer%diags .or. writer%diags_color) then
       write (unit, "(5A)")  "TEX_SOURCES += ", char (id), "_diags.tex"    
       if (os_data%event_analysis_pdf) then
          write (unit, "(5A)")  "TEX_OBJECTS += ", char (id), "_diags.pdf"
       else
          write (unit, "(5A)")  "TEX_OBJECTS += ", char (id), "_diags.ps"
       end if
    end if
    write (unit, "(5A)")  "OBJECTS += ", char (id), ".lo"    
    select type (writer)
    type is (omega_omega_writer_t)       
       write (unit, "(5A)")  char (id), ".f90:"       
       write (unit, "(99A)")  TAB, char (omega_path), &
            " -o ", char (id), ".f90", &
            " -target:whizard", &
            " -target:parameter_module parameters_", char (writer%model_name), &
            " -target:module opr_", char (id), &
            " -target:md5sum '", writer%md5sum, "'", &
            char (openmp_string), &
            char (progress_string), &
            char (kmatrix_string), &
            char (writer%process_mode), char (writer%process_string), &
            char (restrictions_string), char (diagrams_string), &
            char (writer%extra_options)
    type is (omega_ovm_writer_t)
       write (unit, "(5A)")  char (id), ".hbc:"
       write (unit, "(99A)")  TAB, char (omega_path), &
            " -o ", char (id), ".hbc", &
            char (progress_string), &
            char (writer%process_mode), char (writer%process_string), &
            char (restrictions_string), char (diagrams_string), &
            char (writer%extra_options)
       write (unit, "(5A)")  char (id), ".f90:"
       write (unit, "(99A)")  TAB, char (omega_path), &
            " -o ", char (id), ".f90 -params", &
            " -target:whizard ", &
            " -target:bytecode_file ", char (id), ".hbc", &
            " -target:wrapper_module opr_", char (id), &
            " -target:parameter_module_external parameters_", &           
            char (writer%model_name), &
            " -target:md5sum '", writer%md5sum, "'", &
            char (openmp_string)
    end select
    if (writer%diags .or. writer%diags_color) &
       write (unit, "(5A)")  char (id), "_diags.tex: ", char (id), ".f90"
    write (unit, "(5A)")  "clean-", char (id), ":"
    write (unit, "(5A)")  TAB, "rm -f ", char (id), ".f90"
    write (unit, "(5A)")  TAB, "rm -f opr_", char (id), ".mod"
    write (unit, "(5A)")  TAB, "rm -f ", char (id), ".lo"
    write (unit, "(5A)")  "CLEAN_SOURCES += ", char (id), ".f90"    
    select type (writer)
    type is (omega_ovm_writer_t)
       write (unit, "(5A)")  "CLEAN_SOURCES += ", char (id), ".hbc"
    end select
    if (writer%diags .or. writer%diags_color) then
       write (unit, "(5A)")  "CLEAN_SOURCES += ", char (id), "_diags.tex"
    end if
    write (unit, "(5A)")  "CLEAN_OBJECTS += opr_", char (id), ".mod"       
    write (unit, "(5A)")  "CLEAN_OBJECTS += ", char (id), ".lo"
    if (writer%diags .or. writer%diags_color) then    
       write (unit, "(5A)")  "CLEAN_OBJECTS += ", char (id), "_diags.aux"  
       write (unit, "(5A)")  "CLEAN_OBJECTS += ", char (id), "_diags.log"         
       write (unit, "(5A)")  "CLEAN_OBJECTS += ", char (id), "_diags.dvi"                
       write (unit, "(5A)")  "CLEAN_OBJECTS += ", char (id), "_diags.toc"                       
       write (unit, "(5A)")  "CLEAN_OBJECTS += ", char (id), "_diags.out"       
       write (unit, "(5A)")  "CLEAN_OBJECTS += ", char (id), "_diags-fmf.[1-9]"       
       write (unit, "(5A)")  "CLEAN_OBJECTS += ", char (id), "_diags-fmf.[1-9][0-9]"    
       write (unit, "(5A)")  "CLEAN_OBJECTS += ", char (id), "_diags-fmf.[1-9][0-9][0-9]"   
       write (unit, "(5A)")  "CLEAN_OBJECTS += ", char (id), "_diags-fmf.t[1-9]"       
       write (unit, "(5A)")  "CLEAN_OBJECTS += ", char (id), "_diags-fmf.t[1-9][0-9]"
       write (unit, "(5A)")  "CLEAN_OBJECTS += ", char (id), "_diags-fmf.t[1-9][0-9][0-9]"
       write (unit, "(5A)")  "CLEAN_OBJECTS += ", char (id), "_diags-fmf.mp"
       write (unit, "(5A)")  "CLEAN_OBJECTS += ", char (id), "_diags-fmf.log"       
       write (unit, "(5A)")  "CLEAN_OBJECTS += ", char (id), "_diags.dvi"              
       write (unit, "(5A)")  "CLEAN_OBJECTS += ", char (id), "_diags.ps"                     
       if (os_data%event_analysis_pdf) &
            write (unit, "(5A)")  "CLEAN_OBJECTS += ", char (id), "_diags.pdf" 
    end if
    write (unit, "(5A)")  char (id), ".lo: ", char (id), ".f90"
    write (unit, "(5A)")  TAB, "$(LTFCOMPILE) $<"
    if (writer%diags .or. writer%diags_color) then
       if (os_data%event_analysis_ps) then
          if (os_data%event_analysis_pdf) then
             write (unit, "(5A)")  char (id), "_diags.pdf: ", char (id), "_diags.tex"
          else
             write (unit, "(5A)")  char (id), "_diags.ps: ", char (id), "_diags.tex"
          end if
          if (escape_hyperref) then
             write (unit, "(5A)")  TAB, "-cat ", char (id), "_diags.tex | \\" 
             write (unit, "(5A)")  TAB, "   sed -e" // &
                "'s/\\\\usepackage\\[colorlinks\\]\{hyperref\}.*/%\\\\usepackage" // &
                "\\[colorlinks\\]\{hyperref\}/' > \\"
             write (unit, "(5A)")  TAB, "   ", char (id), "_diags.tex.tmp"
             write (unit, "(5A)")  TAB, "mv -f ", char (id), "_diags.tex.tmp \\"
             write (unit, "(5A)")  TAB, "   ", char (id), "_diags.tex"
          end if
          write (unit, "(5A)")  TAB, "-TEXINPUTS=$(TEX_FLAGS) $(LATEX) " // &
               char (id) // "_diags.tex"
          write (unit, "(5A)")  TAB, "MPINPUTS=$(MP_FLAGS) $(MPOST) " // &
               char (id) // "_diags-fmf.mp"
          write (unit, "(5A)")  TAB, "TEXINPUTS=$(TEX_FLAGS) $(LATEX) " // &
               char (id) // "_diags.tex"  
          write (unit, "(5A)")  TAB, "$(DVIPS) -o " // char (id) // "_diags.ps " // &
               char (id) // "_diags.dvi"
          if (os_data%event_analysis_pdf) then
             write (unit, "(5A)")  TAB, "$(PS2PDF) " // char (id) // "_diags.ps"
          end if
       end if
    end if
  end subroutine omega_write_makefile_code

\eatline
\nwendcode{}\nwbegindocs{247}The source is written by the makefile, so nothing to do here.
\nwenddocs{}\nwbegincode{248}\moddef{Omega interface: omega writer: TBP}\plusendmoddef
  procedure :: write_source_code => omega_write_source_code
\nwendcode{}\nwbegincode{249}\moddef{Omega interface: procedures}\plusendmoddef
  subroutine omega_write_source_code (writer, id)
    class(omega_writer_t), intent(in) :: writer
    type(string_t), intent(in) :: id
  end subroutine omega_write_source_code

\eatline
\nwendcode{}\nwbegindocs{250}Return the name of a procedure that implements a given feature, as
it is provided by the external matrix-element code.  \oMega\ names
some procedures differently, therefore we translate here and override
the binding of the base type.
\nwenddocs{}\nwbegincode{251}\moddef{Omega interface: omega writer: TBP}\plusendmoddef
  procedure, nopass :: get_procname => omega_writer_get_procname
\nwendcode{}\nwbegincode{252}\moddef{Omega interface: procedures}\plusendmoddef
  function omega_writer_get_procname (feature) result (name)
    type(string_t) :: name
    type(string_t), intent(in) :: feature
    select case (char (feature))
    case ("n_in");   name = "number_particles_in"
    case ("n_out");  name = "number_particles_out"
    case ("n_flv");  name = "number_flavor_states"
    case ("n_hel");  name = "number_spin_states"
    case ("n_col");  name = "number_color_flows"
    case ("n_cin");  name = "number_color_indices"
    case ("n_cf");   name = "number_color_factors"
    case ("flv_state");  name = "flavor_states"
    case ("hel_state");  name = "spin_states"
    case ("col_state");  name = "color_flows"
    case default
       name = feature
    end select
  end function omega_writer_get_procname
  
\eatline
\nwendcode{}\nwbegindocs{253}The interfaces for the \oMega-specific features.
\nwenddocs{}\nwbegincode{254}\moddef{Omega interface: omega writer: TBP}\plusendmoddef
  procedure :: write_interface => omega_write_interface
\nwendcode{}\nwbegincode{255}\moddef{Omega interface: procedures}\plusendmoddef
  subroutine omega_write_interface (writer, unit, id, feature)
    class(omega_writer_t), intent(in) :: writer
    integer, intent(in) :: unit
    type(string_t), intent(in) :: id
    type(string_t), intent(in) :: feature
    type(string_t) :: name
    name = writer%get_c_procname (id, feature)
    write (unit, "(2x,9A)")  "interface"
    select case (char (feature))
    case ("init")
       write (unit, "(5x,9A)")  "subroutine ", char (name), " (par) bind(C)"
       write (unit, "(7x,9A)")  "import"
       write (unit, "(7x,9A)")  "real(c_default_float), dimension(*), &
            &intent(in) :: par"
       write (unit, "(5x,9A)")  "end subroutine ", char (name)
    case ("update_alpha_s")
       write (unit, "(5x,9A)")  "subroutine ", char (name), " (alpha_s) bind(C)"
       write (unit, "(7x,9A)")  "import"
       write (unit, "(7x,9A)")  "real(c_default_float), intent(in) :: alpha_s"
       write (unit, "(5x,9A)")  "end subroutine ", char (name)
    case ("reset_helicity_selection")
       write (unit, "(5x,9A)")  "subroutine ", char (name), " &
            &(threshold, cutoff) bind(C)"
       write (unit, "(7x,9A)")  "import"
       write (unit, "(7x,9A)")  "real(c_default_float), intent(in) :: threshold"
       write (unit, "(7x,9A)")  "integer(c_int), intent(in) :: cutoff"
       write (unit, "(5x,9A)")  "end subroutine ", char (name)
    case ("is_allowed")
       write (unit, "(5x,9A)")  "subroutine ", char (name), " &
            &(flv, hel, col, flag) bind(C)"
       write (unit, "(7x,9A)")  "import"
       write (unit, "(7x,9A)")  "integer(c_int), intent(in) :: flv, hel, col"
       write (unit, "(7x,9A)")  "logical(c_bool), intent(out) :: flag"    
       write (unit, "(5x,9A)")  "end subroutine ", char (name)
    case ("new_event")
       write (unit, "(5x,9A)")  "subroutine ", char (name), " (p) bind(C)"
       write (unit, "(7x,9A)")  "import"
       write (unit, "(7x,9A)")  "real(c_default_float), dimension(0:3,*), &
            &intent(in) :: p"
       write (unit, "(5x,9A)")  "end subroutine ", char (name)
    case ("get_amplitude")
       write (unit, "(5x,9A)")  "subroutine ", char (name), " &
            &(flv, hel, col, amp) bind(C)"
       write (unit, "(7x,9A)")  "import"
       write (unit, "(7x,9A)")  "integer(c_int), intent(in) :: flv, hel, col"
       write (unit, "(7x,9A)")  "complex(c_default_complex), intent(out) &
            &:: amp"    
       write (unit, "(5x,9A)")  "end subroutine ", char (name)
    end select
    write (unit, "(2x,9A)")  "end interface"
  end subroutine omega_write_interface

\eatline
\nwendcode{}\nwbegindocs{256}The wrappers have to take into account conversion between C and
Fortran data types.

NOTE: The case {\tt{}c{\char95}default{\char95}float} $\neq$ {\tt{}default} is not yet covered.
\nwenddocs{}\nwbegincode{257}\moddef{Omega interface: omega writer: TBP}\plusendmoddef
  procedure :: write_wrapper => omega_write_wrapper
\nwendcode{}\nwbegincode{258}\moddef{Omega interface: procedures}\plusendmoddef
  subroutine omega_write_wrapper (writer, unit, id, feature)
    class(omega_writer_t), intent(in) :: writer
    integer, intent(in) :: unit
    type(string_t), intent(in) :: id, feature
    type(string_t) :: name
    name = writer%get_c_procname (id, feature)
    write (unit, *)
    select case (char (feature))
    case ("init")
       write (unit, "(9A)")  "subroutine ", char (name), " (par) bind(C)"
       write (unit, "(2x,9A)")  "use iso_c_binding"
       write (unit, "(2x,9A)")  "use kinds"
       write (unit, "(2x,9A)")  "use opr_", char (id)
       write (unit, "(2x,9A)")  "real(c_default_float), dimension(*), &
            &intent(in) :: par"
       if (c_default_float == default) then
          write (unit, "(2x,9A)")  "call ", char (feature), " (par)"
       end if
       write (unit, "(9A)")  "end subroutine ", char (name)
    case ("update_alpha_s")
       write (unit, "(9A)")  "subroutine ", char (name), " (alpha_s) bind(C)"
       write (unit, "(2x,9A)")  "use iso_c_binding"
       write (unit, "(2x,9A)")  "use kinds"
       write (unit, "(2x,9A)")  "use opr_", char (id)
       if (c_default_float == default) then
          write (unit, "(2x,9A)")  "real(c_default_float), intent(in) &
               &:: alpha_s"
          write (unit, "(2x,9A)")  "call ", char (feature), " (alpha_s)"
       end if
       write (unit, "(9A)")  "end subroutine ", char (name)
    case ("reset_helicity_selection")
       write (unit, "(9A)")  "subroutine ", char (name), &
            " (threshold, cutoff) bind(C)"
       write (unit, "(2x,9A)")  "use iso_c_binding"
       write (unit, "(2x,9A)")  "use kinds"
       write (unit, "(2x,9A)")  "use opr_", char (id)
       if (c_default_float == default) then
          write (unit, "(2x,9A)")  "real(c_default_float), intent(in) &
               &:: threshold"
          write (unit, "(2x,9A)")  "integer(c_int), intent(in) :: cutoff"
          write (unit, "(2x,9A)")  "call ", char (feature), &
               " (threshold, int (cutoff))"
       end if
       write (unit, "(9A)")  "end subroutine ", char (name)
    case ("is_allowed")
       write (unit, "(9A)")  "subroutine ", char (name), &
            " (flv, hel, col, flag) bind(C)"
       write (unit, "(2x,9A)")  "use iso_c_binding"
       write (unit, "(2x,9A)")  "use kinds"
       write (unit, "(2x,9A)")  "use opr_", char (id)
       write (unit, "(2x,9A)")  "integer(c_int), intent(in) :: flv, hel, col"
       write (unit, "(2x,9A)")  "logical(c_bool), intent(out) :: flag"    
       write (unit, "(2x,9A)")  "flag = ", char (feature), &
            " (int (flv), int (hel), int (col))"
       write (unit, "(9A)")  "end subroutine ", char (name)
    case ("new_event")
       write (unit, "(9A)")  "subroutine ", char (name), " (p) bind(C)"
       write (unit, "(2x,9A)")  "use iso_c_binding"
       write (unit, "(2x,9A)")  "use kinds"
       write (unit, "(2x,9A)")  "use opr_", char (id)
       if (c_default_float == default) then
          write (unit, "(2x,9A)")  "real(c_default_float), dimension(0:3,*), &
               &intent(in) :: p"
          write (unit, "(2x,9A)")  "call ", char (feature), " (p)"
       end if
       write (unit, "(9A)")  "end subroutine ", char (name)
    case ("get_amplitude")
       write (unit, "(9A)")  "subroutine ", char (name), &
            " (flv, hel, col, amp) bind(C)"
       write (unit, "(2x,9A)")  "use iso_c_binding"
       write (unit, "(2x,9A)")  "use kinds"
       write (unit, "(2x,9A)")  "use opr_", char (id)
       write (unit, "(2x,9A)")  "integer(c_int), intent(in) :: flv, hel, col"
       write (unit, "(2x,9A)")  "complex(c_default_complex), intent(out) &
            &:: amp"    
       write (unit, "(2x,9A)")  "amp = ", char (feature), &
            " (int (flv), int (hel), int (col))"
       write (unit, "(9A)")  "end subroutine ", char (name)
    end select
  end subroutine omega_write_wrapper

\eatline
\nwendcode{}\nwbegindocs{259}\nwdocspar
\subsection{Driver}
\nwenddocs{}\nwbegincode{260}\moddef{Omega interface: types}\plusendmoddef
  type, extends (prc_core_driver_t) :: omega_driver_t
     procedure(init_t), nopass, pointer :: &
          init => null ()
     procedure(update_alpha_s_t), nopass, pointer :: &
          update_alpha_s => null ()
     procedure(reset_helicity_selection_t), nopass, pointer :: &
          reset_helicity_selection => null ()
     procedure(is_allowed_t), nopass, pointer :: &
          is_allowed => null ()
     procedure(new_event_t), nopass, pointer :: &
          new_event => null ()
     procedure(get_amplitude_t), nopass, pointer :: &
          get_amplitude => null ()
   contains
   \LA{}Omega interface: omega driver: TBP\RA{}
  end type omega_driver_t

\eatline
\nwendcode{}\nwbegindocs{261}The reported type is the same as for the {\tt{}omega{\char95}def{\char95}t} type.
\nwenddocs{}\nwbegincode{262}\moddef{Omega interface: omega driver: TBP}\endmoddef
  procedure, nopass :: type_name => omega_driver_type_name
\nwendcode{}\nwbegincode{263}\moddef{Omega interface: procedures}\plusendmoddef
  function omega_driver_type_name () result (string)
    type(string_t) :: string
    string = "omega"
  end function omega_driver_type_name

\eatline
\nwendcode{}\nwbegindocs{264}\nwdocspar
\subsection{High-level process definition}
This procedure wraps the details filling a process-component
definition entry as appropriate for an
\oMega\ matrix element.

NOTE: For calling the {\tt{}import{\char95}component} method, we must explicitly
address the {\tt{}process{\char95}def{\char95}t} parent object.  The natural way to call
the method of the extended type triggers a bug in gfortran 4.6.  The
string array arguments {\tt{}prt{\char95}in} and {\tt{}prt{\char95}out} become corrupted and
cause a segfault.
\nwenddocs{}\nwbegincode{265}\moddef{Omega interface: public}\plusendmoddef
  public :: omega_make_process_component
\nwendcode{}\nwbegincode{266}\moddef{Omega interface: procedures}\plusendmoddef
  subroutine omega_make_process_component (entry, component_index, &
         model_name, prt_in, prt_out, restrictions, openmp_support, &
         report_progress, extra_options, diags, diags_color)
    class(process_def_entry_t), intent(inout) :: entry
    integer, intent(in) :: component_index
    type(string_t), intent(in) :: model_name
    type(string_t), dimension(:), intent(in) :: prt_in
    type(string_t), dimension(:), intent(in) :: prt_out
    type(string_t), intent(in), optional :: restrictions
    logical, intent(in), optional :: openmp_support
    logical, intent(in), optional :: report_progress
    logical, intent(in), optional :: diags, diags_color
    type(string_t), intent(in), optional :: extra_options
    class(prc_core_def_t), allocatable :: def
    allocate (omega_omega_def_t :: def)
    select type (def)
    type is (omega_omega_def_t)
       call def%init (model_name, prt_in, prt_out, &
            restrictions, openmp_support, report_progress, &
            extra_options, diags, diags_color)
    end select
    call entry%process_def_t%import_component (component_index, &
         n_out = size (prt_out), &
         prt_in  = new_prt_spec (prt_in), &
         prt_out = new_prt_spec (prt_out), &
         method = var_str ("omega"), &
         variant = def)
  end subroutine omega_make_process_component
    
\eatline
\nwendcode{}\nwbegindocs{267}\nwdocspar
\subsection{The {\tt{}prc{\char95}omega{\char95}t} wrapper}
This is an instance of the generic {\tt{}prc{\char95}core{\char95}t} object.  It contains a
pointer to the process definition ({\tt{}omega{\char95}def{\char95}t}), a data component
({\tt{}process{\char95}constants{\char95}t}), and the matrix-element driver
({\tt{}omega{\char95}driver{\char95}t}).
\nwenddocs{}\nwbegincode{268}\moddef{Omega interface: public}\plusendmoddef
  public :: prc_omega_t
\nwendcode{}\nwbegincode{269}\moddef{Omega interface: types}\plusendmoddef
  type, extends (prc_core_t) :: prc_omega_t
     real(default), dimension(:), allocatable :: par
     type(helicity_selection_t) :: helicity_selection
     type(qcd_t) :: qcd
   contains
   \LA{}Omega interface: prc omega: TBP\RA{}
  end type prc_omega_t
  
\eatline
\nwendcode{}\nwbegindocs{270}The workspace associated to a {\tt{}prc{\char95}omega{\char95}t} object contains a single flag.
The flag is used to suppress re-evaluating the matrix element for each
quantum-number combination, after the first amplitude belonging to a given
kinematics has been computed.

We can also store the value of a running coupling once it has been calculated
for an event.  The default value is negative, which indicates an undefined
value in this context.
\nwenddocs{}\nwbegincode{271}\moddef{Omega interface: types}\plusendmoddef
  type, extends (workspace_t) :: omega_state_t
     logical :: new_kinematics = .true.
     real(default) :: alpha_qcd = -1
   contains
     procedure :: write => omega_state_write
     procedure :: reset_new_kinematics => omega_state_reset_new_kinematics
  end type omega_state_t
  
\eatline
\nwendcode{}\nwbegincode{272}\moddef{Omega interface: procedures}\plusendmoddef
  subroutine omega_state_write (object, unit)
    class(omega_state_t), intent(in) :: object
    integer, intent(in), optional :: unit
    integer :: u
    u = given_output_unit (unit)
    write (u, "(3x,A,L1)")  "O'Mega state: new kinematics = ", &
         object%new_kinematics
  end subroutine omega_state_write
  
\eatline
\nwendcode{}\nwbegincode{273}\moddef{Omega interface: procedures}\plusendmoddef
  subroutine omega_state_reset_new_kinematics (object)
    class(omega_state_t), intent(inout) :: object
    object%new_kinematics = .true.
  end subroutine omega_state_reset_new_kinematics

\eatline
\nwendcode{}\nwbegindocs{274}Allocate the workspace with the above specific type.
\nwenddocs{}\nwbegincode{275}\moddef{Omega interface: prc omega: TBP}\endmoddef
  procedure :: allocate_workspace => prc_omega_allocate_workspace
\nwendcode{}\nwbegincode{276}\moddef{Omega interface: procedures}\plusendmoddef
  subroutine prc_omega_allocate_workspace (object, tmp)
    class(prc_omega_t), intent(in) :: object
    class(workspace_t), intent(inout), allocatable :: tmp
    allocate (omega_state_t :: tmp)
  end subroutine prc_omega_allocate_workspace
  
\eatline
\nwendcode{}\nwbegindocs{277}The following procedures are inherited from the base type as deferred, thus
must be implemented.  The corresponding unit tests are skipped here; the
procedures are tested when called from the {\tt{}processes} module.

Output: print just the ID of the associated matrix element.  Then display any
stored parameters and the helicity selection data.  (The latter are printed
only if active.)
\nwenddocs{}\nwbegincode{278}\moddef{Omega interface: prc omega: TBP}\plusendmoddef
  procedure :: write => prc_omega_write
\nwendcode{}\nwbegincode{279}\moddef{Omega interface: procedures}\plusendmoddef
  subroutine prc_omega_write (object, unit)
    class(prc_omega_t), intent(in) :: object
    integer, intent(in), optional :: unit
    integer :: u, i
    u = given_output_unit (unit)
    write (u, "(3x,A)", advance="no")  "O'Mega process core:"
    if (object%data_known) then
       write (u, "(1x,A)")  char (object%data%id)
    else
       write (u, "(1x,A)")  "[undefined]"
    end if
    if (allocated (object%par)) then
       write (u, "(3x,A)")  "Parameter array:"
       do i = 1, size (object%par)
          write (u, "(5x,I0,1x,ES17.10)")  i, object%par(i)
       end do
    end if
    call object%helicity_selection%write (u)
    call object%qcd%write (u)
  end subroutine prc_omega_write
  
\eatline
\nwendcode{}\nwbegindocs{280}Temporarily store the parameter array inside the {\tt{}prc{\char95}omega}
object, so we can use it later during the actual initialization.  Also
store threshold and cutoff for helicity selection.
\nwenddocs{}\nwbegincode{281}\moddef{Omega interface: prc omega: TBP}\plusendmoddef
  procedure :: set_parameters => prc_omega_set_parameters
\nwendcode{}\nwbegincode{282}\moddef{Omega interface: procedures}\plusendmoddef
  subroutine prc_omega_set_parameters (prc_omega, model, &
       helicity_selection, qcd, use_color_factors)
    class(prc_omega_t), intent(inout) :: prc_omega
    class(model_data_t), intent(in), target, optional :: model
    type(helicity_selection_t), intent(in), optional :: helicity_selection
    type(qcd_t), intent(in), optional :: qcd
    logical, intent(in), optional :: use_color_factors
    if (present (model)) then
       if (.not. allocated (prc_omega%par)) &
            allocate (prc_omega%par (model%get_n_real ()))
       call model%real_parameters_to_c_array (prc_omega%par)
    end if
    if (present (helicity_selection)) then
       prc_omega%helicity_selection = helicity_selection
    end if
    if (present (qcd)) then
       prc_omega%qcd = qcd
    end if
    if (present (use_color_factors)) then
       prc_omega%use_color_factors = use_color_factors
    end if
  end subroutine prc_omega_set_parameters
  
\eatline
\nwendcode{}\nwbegindocs{283}To fully initialize the process core, we perform base
initialization, then initialize the external matrix element code.

This procedure overrides the {\tt{}init} method of the base type, which
we nevertheless can access via its binding {\tt{}base{\char95}init}.  When done, we
have an allocated driver.  The driver will call the {\tt{}init} procedure
for the external matrix element, and thus transfer the parameter set to
where it finally belongs.

If requested, we initialize the helicity selction counter.
\nwenddocs{}\nwbegincode{284}\moddef{Omega interface: prc omega: TBP}\plusendmoddef
  procedure :: init => prc_omega_init
\nwendcode{}\nwbegincode{285}\moddef{Omega interface: procedures}\plusendmoddef
  subroutine prc_omega_init (object, def, lib, id, i_component)
    class(prc_omega_t), intent(inout) :: object
    class(prc_core_def_t), intent(in), target :: def
    type(process_library_t), intent(in), target :: lib
    type(string_t), intent(in) :: id
    integer, intent(in) :: i_component
    call object%base_init (def, lib, id, i_component)
    call object%activate_parameters ()
  end subroutine prc_omega_init
    
\eatline
\nwendcode{}\nwbegindocs{286}Activate the stored parameters by transferring them to the external
matrix element.  Also reset the helicity selection, if requested.
\nwenddocs{}\nwbegincode{287}\moddef{Omega interface: prc omega: TBP}\plusendmoddef
  procedure :: activate_parameters => prc_omega_activate_parameters
\nwendcode{}\nwbegincode{288}\moddef{Omega interface: procedures}\plusendmoddef
  subroutine prc_omega_activate_parameters (object)
    class (prc_omega_t), intent(inout) :: object
    if (allocated (object%driver)) then
       if (allocated (object%par)) then
          select type (driver => object%driver)
          type is (omega_driver_t)
             if (associated (driver%init))  call driver%init (object%par)
          end select
       else
          call msg_bug ("prc_omega_activate: parameter set is not allocated")
       end if
       call object%reset_helicity_selection ()
    else
       call msg_bug ("prc_omega_activate: driver is not allocated")
    end if
  end subroutine prc_omega_activate_parameters
    
\eatline
\nwendcode{}\nwbegindocs{289}The \oMega\ process is an independent process which needs its own Monte
Carlo parameter set for integration.
\nwenddocs{}\nwbegincode{290}\moddef{Omega interface: prc omega: TBP}\plusendmoddef
  procedure :: needs_mcset => prc_omega_needs_mcset
\nwendcode{}\nwbegincode{291}\moddef{Omega interface: procedures}\plusendmoddef
  function prc_omega_needs_mcset (object) result (flag)
    class(prc_omega_t), intent(in) :: object
    logical :: flag
    flag = .true.
  end function prc_omega_needs_mcset
  
\eatline
\nwendcode{}\nwbegindocs{292}There is only one term for this tree-level process.
\nwenddocs{}\nwbegincode{293}\moddef{Omega interface: prc omega: TBP}\plusendmoddef
  procedure :: get_n_terms => prc_omega_get_n_terms
\nwendcode{}\nwbegincode{294}\moddef{Omega interface: procedures}\plusendmoddef
  function prc_omega_get_n_terms (object) result (n)
    class(prc_omega_t), intent(in) :: object
    integer :: n
    n = 1
  end function prc_omega_get_n_terms

\eatline
\nwendcode{}\nwbegindocs{295}Tell whether a particular combination of flavor, helicity, color is
allowed.  Here we have to consult the matrix-element driver.
\nwenddocs{}\nwbegincode{296}\moddef{Omega interface: prc omega: TBP}\plusendmoddef
  procedure :: is_allowed => prc_omega_is_allowed
\nwendcode{}\nwbegincode{297}\moddef{Omega interface: procedures}\plusendmoddef
 function prc_omega_is_allowed (object, i_term, f, h, c) result (flag)
    class(prc_omega_t), intent(in) :: object
    integer, intent(in) :: i_term, f, h, c
    logical :: flag
    logical(c_bool) :: cflag
    select type (driver => object%driver)
    type is (omega_driver_t)
       call driver%is_allowed (f, h, c, cflag)
       flag = cflag
    end select
  end function prc_omega_is_allowed
 
\eatline
\nwendcode{}\nwbegindocs{298}Transfer the generated momenta directly to the hard interaction in
the (only) term.  We assume that everything has been set up correctly,
so the array fits.

We don't reset the {\tt{}new{\char95}kinematics} flag here.  This has to be done
explicitly by the caller ({\tt{}reset{\char95}new{\char95}kinematics}) when a new kinematics
configuration is to be considered.
\nwenddocs{}\nwbegincode{299}\moddef{Omega interface: prc omega: TBP}\plusendmoddef
  procedure :: compute_hard_kinematics => prc_omega_compute_hard_kinematics
\nwendcode{}\nwbegincode{300}\moddef{Omega interface: procedures}\plusendmoddef
  subroutine prc_omega_compute_hard_kinematics &
       (object, p_seed, i_term, int_hard, tmp)
    class(prc_omega_t), intent(in) :: object
    type(vector4_t), dimension(:), intent(in) :: p_seed
    integer, intent(in) :: i_term
    type(interaction_t), intent(inout) :: int_hard
    class(workspace_t), intent(inout), allocatable :: tmp
    call interaction_set_momenta (int_hard, p_seed)
  end subroutine prc_omega_compute_hard_kinematics
  
\eatline
\nwendcode{}\nwbegindocs{301}This procedure is not called for {\tt{}prc{\char95}omega{\char95}t}, just a placeholder.
\nwenddocs{}\nwbegincode{302}\moddef{Omega interface: prc omega: TBP}\plusendmoddef
  procedure :: compute_eff_kinematics => prc_omega_compute_eff_kinematics
\nwendcode{}\nwbegincode{303}\moddef{Omega interface: procedures}\plusendmoddef
  subroutine prc_omega_compute_eff_kinematics &
       (object, i_term, int_hard, int_eff, tmp)
    class(prc_omega_t), intent(in) :: object
    integer, intent(in) :: i_term
    type(interaction_t), intent(in) :: int_hard
    type(interaction_t), intent(inout) :: int_eff
    class(workspace_t), intent(inout), allocatable :: tmp
  end subroutine prc_omega_compute_eff_kinematics
  
\eatline
\nwendcode{}\nwbegindocs{304}Recover the momenta within the hard interaction.  For tree-level
processes, this is a trivial copy of the incoming seed and outgoing
effective momenta.  The effective interaction is a pointer to the hard
interaction, so the latter is ignored.
\nwenddocs{}\nwbegincode{305}\moddef{Omega interface: prc omega: TBP}\plusendmoddef
  procedure :: recover_kinematics => prc_omega_recover_kinematics
\nwendcode{}\nwbegincode{306}\moddef{Omega interface: tests}\endmoddef
  subroutine prc_omega_recover_kinematics &
       (object, p_seed, int_hard, int_eff, tmp)
    class(prc_omega_t), intent(in) :: object
    type(vector4_t), dimension(:), intent(inout) :: p_seed
    type(interaction_t), intent(inout) :: int_hard
    type(interaction_t), intent(inout) :: int_eff
    class(workspace_t), intent(inout), allocatable :: tmp
    integer :: n_in
    n_in = interaction_get_n_in (int_eff)
    call interaction_set_momenta (int_eff, p_seed(1:n_in), outgoing = .false.)
    p_seed(n_in+1:) = interaction_get_momenta (int_eff, outgoing = .true.)
  end subroutine prc_omega_recover_kinematics
    
\eatline
\nwendcode{}\nwbegindocs{307}Reset the helicity selection counters and start counting zero
helicities.  We assume that the {\tt{}helicity{\char95}selection} object is allocated.
Otherwise, reset and switch off helicity counting.

In the test routine, the driver is allocated but the driver methods are not.
Therefore, guard against a disassociated method.
\nwenddocs{}\nwbegincode{308}\moddef{Omega interface: prc omega: TBP}\plusendmoddef
  procedure :: reset_helicity_selection => prc_omega_reset_helicity_selection
\nwendcode{}\nwbegincode{309}\moddef{Omega interface: procedures}\plusendmoddef
  subroutine prc_omega_reset_helicity_selection (object)
    class(prc_omega_t), intent(inout) :: object
    select type (driver => object%driver)
    type is (omega_driver_t)
       if (associated (driver%reset_helicity_selection)) then
          if (object%helicity_selection%active) then
             call driver%reset_helicity_selection &
                  (real (object%helicity_selection%threshold, &
                  c_default_float), &
                  int (object%helicity_selection%cutoff, c_int))
          else
             call driver%reset_helicity_selection &
                  (0._c_default_float, 0_c_int)
          end if
       end if
    end select
  end subroutine prc_omega_reset_helicity_selection
  
\eatline
\nwendcode{}\nwbegindocs{310}Compute the amplitude.  For the tree-level process, we can ignore the scale
settings.  The term index {\tt{}j} is also irrelevant.

We first call {\tt{}new{\char95}event} for the given momenta (which we must unpack), then
retrieve the amplitude value for the given quantum numbers.

If the {\tt{}tmp} status flag is present, we can make sure that we call
{\tt{}new{\char95}event} only once for a given kinematics.  After the first call, we
unset the {\tt{}new{\char95}kinematics} flag.

The core objects computes the appropriate $\alpha_s$ value via the {\tt{}qcd}
subobject, taking into account the provided {\tt{}fac{\char95}scale} value.  However, if
the extra parameter {\tt{}alpha{\char95}qcd{\char95}forced} is allocated, it overrides this
setting.

The {\tt{}is{\char95}allowed} query is not redundant, since the status may change during
the run if helicities are switched off.
\nwenddocs{}\nwbegincode{311}\moddef{Omega interface: prc omega: TBP}\plusendmoddef
  procedure :: compute_amplitude => prc_omega_compute_amplitude
\nwendcode{}\nwbegincode{312}\moddef{Omega interface: procedures}\plusendmoddef
  function prc_omega_compute_amplitude &
       (object, j, p, f, h, c, fac_scale, ren_scale, alpha_qcd_forced, tmp) &
       result (amp)
    class(prc_omega_t), intent(in) :: object
    integer, intent(in) :: j
    type(vector4_t), dimension(:), intent(in) :: p
    integer, intent(in) :: f, h, c
    real(default), intent(in) :: fac_scale, ren_scale
    real(default), intent(in), allocatable :: alpha_qcd_forced
    class(workspace_t), intent(inout), allocatable, optional :: tmp
    real(default) :: alpha_qcd
    complex(default) :: amp
    integer :: n_tot, i
    real(c_default_float), dimension(:,:), allocatable :: parray
    complex(c_default_complex) :: camp
    logical :: new_event
    select type (driver => object%driver)
    type is (omega_driver_t)
       new_event = .true.
       if (present (tmp)) then
          if (allocated (tmp)) then
             select type (tmp)
             type is (omega_state_t)
                new_event = tmp%new_kinematics
                tmp%new_kinematics = .false.
             end select
          end if
       end if
       if (new_event) then
          if (allocated (object%qcd%alpha)) then
             if (allocated (alpha_qcd_forced)) then
                alpha_qcd = alpha_qcd_forced
             else
                alpha_qcd = object%qcd%alpha%get (fac_scale)
             end if
             call driver%update_alpha_s (alpha_qcd)
             if (present (tmp)) then
                if (allocated (tmp)) then
                   select type (tmp)
                   type is (omega_state_t)
                      tmp%alpha_qcd = alpha_qcd
                   end select
                end if
             end if
          end if
          n_tot = object%data%n_in + object%data%n_out
          allocate (parray (0:3, n_tot))
          do i = 1, n_tot
             parray(:,i) = vector4_get_components (p(i))
          end do
          call driver%new_event (parray)
       end if
       if (object%is_allowed (1, f, h, c)) then
          call driver%get_amplitude &
               (int (f, c_int), int (h, c_int), int (c, c_int), camp)
          amp = camp
       else
          amp = 0
       end if
    end select
  end function prc_omega_compute_amplitude
    
\eatline
\nwendcode{}\nwbegindocs{313}After the amplitude has been computed, we may read off the current value of
$\alpha_s$.  This works only if $\alpha_s$ varies, and if the workspace
{\tt{}tmp} is present which stores this value.
\nwenddocs{}\nwbegincode{314}\moddef{Omega interface: prc omega: TBP}\plusendmoddef
  procedure :: get_alpha_s => prc_omega_get_alpha_s
\nwendcode{}\nwbegincode{315}\moddef{Omega interface: procedures}\plusendmoddef
  function prc_omega_get_alpha_s (object, tmp) result (alpha)
    class(prc_omega_t), intent(in) :: object
    class(workspace_t), intent(in), allocatable :: tmp
    real(default) :: alpha
    alpha = -1
    if (allocated (object%qcd%alpha) .and. allocated (tmp)) then
       select type (tmp)
       type is (omega_state_t)
          alpha = tmp%alpha_qcd
       end select
    end if
  end function prc_omega_get_alpha_s
  
\eatline
\nwendcode{}\nwbegindocs{316}\nwdocspar
\subsection{Test}
This is the master for calling self-test procedures. There is a
separate test for testing \oMega\ diagram generation as this depends on
a working analysis setup.
\nwenddocs{}\nwbegincode{317}\moddef{Omega interface: public}\plusendmoddef
  public :: prc_omega_test
\nwendcode{}\nwbegincode{318}\moddef{Omega interface: tests}\plusendmoddef
  subroutine prc_omega_test (u, results)
    integer, intent(in) :: u
    type(test_results_t), intent(inout) :: results
  \LA{}Omega interface: execute tests\RA{}
end subroutine prc_omega_test

\eatline
\nwendcode{}\nwbegindocs{319}\nwdocspar
\nwenddocs{}\nwbegincode{320}\moddef{Omega interface: public}\plusendmoddef
  public :: prc_omega_diags_test
\nwendcode{}\nwbegincode{321}\moddef{Omega interface: tests}\plusendmoddef
  subroutine prc_omega_diags_test (u, results)
    integer, intent(in) :: u
    type(test_results_t), intent(inout) :: results
  \LA{}Omega interface: execute diags tests\RA{}
end subroutine prc_omega_diags_test

\eatline
\nwendcode{}\nwbegindocs{322}\nwdocspar
\subsubsection{Generate, compile and load a simple process matrix element}
The process is $e^+ e^- \to \mu^+\mu^-$ for vanishing masses and
$e=0.3$.  We initialize the process, build the library, and compute a
particular matrix element for momenta of unit energy and right-angle
scattering.  The matrix element, as it happens, is equal to $e^2$.
(Note that are no conversion factors applied, so this result is
exact.)

For {\tt{}GNU\ make}, {\tt{}makeflags} is set to {\tt{}-j1}.  This eliminates a
potential clash with a {\tt{}-j<n>} flag if this test is called from a
parallel make.
\nwenddocs{}\nwbegincode{323}\moddef{Omega interface: execute tests}\endmoddef
  call test (prc_omega_1, "prc_omega_1", &
       "build and load simple OMega process", &
       u, results)
\nwendcode{}\nwbegincode{324}\moddef{Omega interface: tests}\plusendmoddef
  subroutine prc_omega_1 (u)
    integer, intent(in) :: u
    type(process_library_t) :: lib
    class(prc_core_def_t), allocatable :: def
    type(process_def_entry_t), pointer :: entry
    type(os_data_t) :: os_data
    type(string_t) :: model_name
    type(string_t), dimension(:), allocatable :: prt_in, prt_out
    type(process_constants_t) :: data
    class(prc_core_driver_t), allocatable :: driver
    integer, parameter :: cdf = c_default_float
    integer, parameter :: ci = c_int
    real(cdf), dimension(4) :: par
    real(cdf), dimension(0:3,4) :: p
    logical(c_bool) :: flag
    complex(c_default_complex) :: amp
    integer :: i
    
    write (u, "(A)")  "* Test output: prc_omega_1"
    write (u, "(A)")  "*   Purpose: create a simple process with OMega"
    write (u, "(A)")  "*            build a library, link, load, and &
         &access the matrix element"
    write (u, "(A)")

    write (u, "(A)")  "* Initialize a process library with one entry"
    write (u, "(A)")
    call lib%init (var_str ("omega1"))
    call os_data_init (os_data)

    model_name = "QED"
    allocate (prt_in (2), prt_out (2))
    prt_in = [var_str ("e+"), var_str ("e-")]
    prt_out = [var_str ("m+"), var_str ("m-")]
    
    allocate (omega_omega_def_t :: def)
    select type (def)
    type is (omega_omega_def_t)
       call def%init (model_name, prt_in, prt_out)
    end select
    allocate (entry)
    call entry%init (var_str ("omega1_a"), model_name = model_name, &
         n_in = 2, n_components = 1)
    call entry%import_component (1, n_out = size (prt_out), &
         prt_in  = new_prt_spec (prt_in), &
         prt_out = new_prt_spec (prt_out), &
         method  = var_str ("omega"), &
         variant = def)
    call lib%append (entry)
    
    write (u, "(A)")  "* Configure library"
    write (u, "(A)")
    call lib%configure (os_data)
    
    write (u, "(A)")  "* Write makefile"
    write (u, "(A)")
    call lib%write_makefile (os_data, force = .true.)

    write (u, "(A)")  "* Clean any left-over files"
    write (u, "(A)")
    call lib%clean (os_data, distclean = .false.)

    write (u, "(A)")  "* Write driver"
    write (u, "(A)")
    call lib%write_driver (force = .true.)

    write (u, "(A)")  "* Write process source code, compile, link, load"
    write (u, "(A)")
    call lib%load (os_data)

    call lib%write (u, libpath = .false.)
    
    write (u, "(A)")
    write (u, "(A)")  "* Probe library API:"
    write (u, "(A)")
       
    write (u, "(1x,A,L1)")  "is active                 = ", &
         lib%is_active ()
    write (u, "(1x,A,I0)")  "n_processes               = ", &
         lib%get_n_processes ()

    write (u, "(A)")
    write (u, "(A)")  "* Constants of omega1_a_i1:"
    write (u, "(A)")

    call lib%connect_process (var_str ("omega1_a"), 1, data, driver)

    write (u, "(1x,A,A)")  "component ID     = ", char (data%id)
    write (u, "(1x,A,A)")  "model name       = ", char (data%model_name)
    write (u, "(1x,A,A,A)")  "md5sum           = '", data%md5sum, "'"
    write (u, "(1x,A,L1)") "openmp supported = ", data%openmp_supported
    write (u, "(1x,A,I0)") "n_in  = ", data%n_in
    write (u, "(1x,A,I0)") "n_out = ", data%n_out
    write (u, "(1x,A,I0)") "n_flv = ", data%n_flv
    write (u, "(1x,A,I0)") "n_hel = ", data%n_hel
    write (u, "(1x,A,I0)") "n_col = ", data%n_col
    write (u, "(1x,A,I0)") "n_cin = ", data%n_cin
    write (u, "(1x,A,I0)") "n_cf  = ", data%n_cf
    write (u, "(1x,A,10(1x,I0))") "flv state =", data%flv_state
    write (u, "(1x,A,10(1x,I2))") "hel state =", data%hel_state(:,1)
    do i = 2, 16
       write (u, "(12x,4(1x,I2))")  data%hel_state(:,i)
    end do
    write (u, "(1x,A,10(1x,I0))") "col state =", data%col_state
    write (u, "(1x,A,10(1x,L1))") "ghost flag =", data%ghost_flag
    write (u, "(1x,A,10(1x,F5.3))") "color factors =", data%color_factors
    write (u, "(1x,A,10(1x,I0))") "cf index =", data%cf_index

    write (u, "(A)")
    write (u, "(A)")  "* Set parameters for omega1_a and initialize:"
    write (u, "(A)")

    par = [0.3_cdf, 0.0_cdf, 0.0_cdf, 0.0_cdf]
    write (u, "(2x,A,F6.4)")  "ee   = ", par(1)
    write (u, "(2x,A,F6.4)")  "me   = ", par(2)
    write (u, "(2x,A,F6.4)")  "mmu  = ", par(3)
    write (u, "(2x,A,F6.4)")  "mtau = ", par(4)

    write (u, "(A)")
    write (u, "(A)")  "* Set kinematics:"
    write (u, "(A)")
    
    p = reshape ([ &
         1.0_cdf, 0.0_cdf, 0.0_cdf, 1.0_cdf, &
         1.0_cdf, 0.0_cdf, 0.0_cdf,-1.0_cdf, &
         1.0_cdf, 1.0_cdf, 0.0_cdf, 0.0_cdf, &
         1.0_cdf,-1.0_cdf, 0.0_cdf, 0.0_cdf &
         ], [4,4])
    do i = 1, 4
       write (u, "(2x,A,I0,A,4(1x,F7.4))")  "p", i, " =", p(:,i)
    end do

    select type (driver)
    type is (omega_driver_t)
       call driver%init (par)

       call driver%new_event (p)

       write (u, "(A)")
       write (u, "(A)")  "* Compute matrix element:"
       write (u, "(A)")

       call driver%is_allowed (1_ci, 6_ci, 1_ci, flag)
       write (u, "(1x,A,L1)") "is_allowed (1, 6, 1) = ", flag
       
       call driver%get_amplitude (1_ci, 6_ci, 1_ci, amp)
       write (u, "(1x,A,1x,E11.4)") "|amp (1, 6, 1)| =", abs (amp)
    end select

    call lib%final ()
    
    write (u, "(A)")
    write (u, "(A)")  "* Test output end: prc_omega_1"
    
  end subroutine prc_omega_1
  
\eatline
\nwendcode{}\nwbegindocs{325}\nwdocspar
\subsubsection{Check {\tt{}prc{\char95}omega{\char95}t} wrapper and options}
The process is $e^- e^+ \to e^- e^+$ for vanishing masses and
$e=0.3$.  We build the library using the high-level procedure
{\tt{}omega{\char95}make{\char95}process{\char95}component} and the ``black box''
{\tt{}prc{\char95}omega{\char95}t} object.  Two variants with different settings for
restrictions and OpenMP.

For {\tt{}GNU\ make}, {\tt{}makeflags} is set to {\tt{}-j1}.  This eliminates a
potential clash with a {\tt{}-j<n>} flag if this test is called from a
parallel make.
\nwenddocs{}\nwbegincode{326}\moddef{Omega interface: execute tests}\plusendmoddef
  call test (prc_omega_2, "prc_omega_2", &
       "OMega option passing", &
       u, results)
\nwendcode{}\nwbegincode{327}\moddef{Omega interface: tests}\plusendmoddef
  subroutine prc_omega_2 (u)
    integer, intent(in) :: u
    type(process_library_t), target :: lib
    type(process_def_entry_t), pointer :: entry
    type(os_data_t) :: os_data
    type(string_t) :: model_name
    class(model_data_t), pointer :: model
    class(vars_t), pointer :: vars
    type(string_t), dimension(:), allocatable :: prt_in, prt_out
    type(string_t) :: restrictions
    type(process_component_def_t), pointer :: config
    type(prc_omega_t) :: prc1, prc2
    type(process_constants_t) :: data
    integer, parameter :: cdf = c_default_float
    integer, parameter :: ci = c_int
    real(cdf), dimension(:), allocatable :: par
    real(cdf), dimension(0:3,4) :: p
    complex(c_default_complex) :: amp
    integer :: i
    logical :: exist
    
    write (u, "(A)")  "* Test output: prc_omega_2"
    write (u, "(A)")  "*   Purpose: create simple processes with OMega"
    write (u, "(A)")  "*            use the prc_omega wrapper for this"
    write (u, "(A)")  "*            and check OMega options"
    write (u, "(A)")

    write (u, "(A)")  "* Initialize a process library with two entries, &
         &different options."
    write (u, "(A)")  "* (1) e- e+ -> e- e+   &
         &(all diagrams, no OpenMP, report progress)"
    write (u, "(A)")  "* (2) e- e+ -> e- e+   &
         &(s-channel only, with OpenMP, report progress to file)"

    call lib%init (var_str ("omega2"))
    call os_data_init (os_data)

    model_name = "QED"
    call prepare_model (model, model_name, vars)

    allocate (prt_in (2), prt_out (2))
    prt_in = [var_str ("e-"), var_str ("e+")]
    prt_out = prt_in
    restrictions = "3+4~A"

    allocate (entry)
    call entry%init (var_str ("omega2_a"), &
         model, n_in = 2, n_components = 2)

    call omega_make_process_component (entry, 1, &
         model_name, prt_in, prt_out, &
         report_progress=.true.)
    call omega_make_process_component (entry, 2, &
         model_name, prt_in, prt_out, &
         restrictions=restrictions, openmp_support=.true., &
         extra_options=var_str ("-fusion:progress_file omega2.log"))

    call lib%append (entry)
    
    write (u, "(A)")
    write (u, "(A)")  "* Remove left-over file"
    write (u, "(A)")

    call delete_file ("omega2.log")
    inquire (file="omega2.log", exist=exist)
    write (u, "(1x,A,L1)")  "omega2.log exists = ", exist

    write (u, "(A)")
    write (u, "(A)")  "* Build and load library"

    call lib%configure (os_data)
    call lib%write_makefile (os_data, force = .true.)
    call lib%clean (os_data, distclean = .false.)
    call lib%write_driver (force = .true.)
    call lib%load (os_data)
    
    write (u, "(A)")
    write (u, "(A)")  "* Check extra output of OMega"
    write (u, "(A)")

    inquire (file="omega2.log", exist=exist)
    write (u, "(1x,A,L1)")  "omega2.log exists = ", exist

    write (u, "(A)")
    write (u, "(A)")  "* Probe library API:"
    write (u, "(A)")
       
    write (u, "(1x,A,L1)")  "is active                 = ", &
         lib%is_active ()
    write (u, "(1x,A,I0)")  "n_processes               = ", &
         lib%get_n_processes ()

    write (u, "(A)")
    write (u, "(A)")  "* Set parameters for omega2_a and initialize:"
    write (u, "(A)")

    call vars%set_rval (var_str ("ee"), 0.3_default)
    call vars%set_rval (var_str ("me"), 0._default)
    call vars%set_rval (var_str ("mmu"), 0._default)
    call vars%set_rval (var_str ("mtau"), 0._default)
    allocate (par (model%get_n_real ()))
    call model%real_parameters_to_c_array (par)

    write (u, "(2x,A,F6.4)")  "ee   = ", par(1)
    write (u, "(2x,A,F6.4)")  "me   = ", par(2)
    write (u, "(2x,A,F6.4)")  "mmu  = ", par(3)
    write (u, "(2x,A,F6.4)")  "mtau = ", par(4)

    call prc1%set_parameters (model)
    call prc2%set_parameters (model)

    write (u, "(A)")
    write (u, "(A)")  "* Constants of omega2_a_i1:"
    write (u, "(A)")

    config => lib%get_component_def_ptr (var_str ("omega2_a"), 1)
    call prc1%init (config%get_core_def_ptr (), &
         lib, var_str ("omega2_a"), 1)
    call prc1%get_constants (data, 1)

    write (u, "(1x,A,A)")  "component ID     = ", &
         char (data%id)
    write (u, "(1x,A,L1)") "openmp supported = ", &
         data%openmp_supported
    write (u, "(1x,A,A,A)") "model name       = '", &
         char (data%model_name), "'"

    write (u, "(A)")
    write (u, "(A)")  "* Constants of omega2_a_i2:"
    write (u, "(A)")

    config => lib%get_component_def_ptr (var_str ("omega2_a"), 2)
    call prc2%init (config%get_core_def_ptr (), &
         lib, var_str ("omega2_a"), 2)
    call prc2%get_constants (data, 1)

    write (u, "(1x,A,A)")  "component ID     = ", &
         char (data%id)
    write (u, "(1x,A,L1)") "openmp supported = ", &
         data%openmp_supported
    write (u, "(1x,A,A,A)") "model name       = '", &
         char (data%model_name), "'"

    write (u, "(A)")
    write (u, "(A)")  "* Set kinematics:"
    write (u, "(A)")
    
    p = reshape ([ &
         1.0_cdf, 0.0_cdf, 0.0_cdf, 1.0_cdf, &
         1.0_cdf, 0.0_cdf, 0.0_cdf,-1.0_cdf, &
         1.0_cdf, 1.0_cdf, 0.0_cdf, 0.0_cdf, &
         1.0_cdf,-1.0_cdf, 0.0_cdf, 0.0_cdf &
         ], [4,4])
    do i = 1, 4
       write (u, "(2x,A,I0,A,4(1x,F7.4))")  "p", i, " =", p(:,i)
    end do

    write (u, "(A)")
    write (u, "(A)")  "* Compute matrix element:"
    write (u, "(A)")

    select type (driver => prc1%driver)
    type is (omega_driver_t)
       call driver%new_event (p)
       call driver%get_amplitude (1_ci, 6_ci, 1_ci, amp)
       write (u, "(2x,A,1x,E11.4)") "(1) |amp (1, 6, 1)| =", abs (amp)
    end select

    select type (driver => prc2%driver)
    type is (omega_driver_t)
       call driver%new_event (p)
       call driver%get_amplitude (1_ci, 6_ci, 1_ci, amp)
       write (u, "(2x,A,1x,E11.4)") "(2) |amp (1, 6, 1)| =", abs (amp)
    end select

    write (u, "(A)")
    write (u, "(A)")  "* Set kinematics:"
    write (u, "(A)")
    
    p = reshape ([ &
         1.0_cdf, 0.0_cdf, 0.0_cdf, 1.0_cdf, &
         1.0_cdf, 0.0_cdf, 0.0_cdf,-1.0_cdf, &
         1.0_cdf, sqrt(0.5_cdf), 0.0_cdf, sqrt(0.5_cdf), &
         1.0_cdf,-sqrt(0.5_cdf), 0.0_cdf,-sqrt(0.5_cdf) &
         ], [4,4])
    do i = 1, 4
       write (u, "(2x,A,I0,A,4(1x,F7.4))")  "p", i, " =", p(:,i)
    end do

    write (u, "(A)")
    write (u, "(A)")  "* Compute matrix element:"
    write (u, "(A)")

    select type (driver => prc1%driver)
    type is (omega_driver_t)
       call driver%new_event (p)
       call driver%get_amplitude (1_ci, 6_ci, 1_ci, amp)
       write (u, "(2x,A,1x,E11.4)") "(1) |amp (1, 6, 1)| =", abs (amp)
    end select

    select type (driver => prc2%driver)
    type is (omega_driver_t)
       call driver%new_event (p)
       call driver%get_amplitude (1_ci, 6_ci, 1_ci, amp)
       write (u, "(2x,A,1x,E11.4)") "(2) |amp (1, 6, 1)| =", abs (amp)
    end select

    call lib%final ()
    call cleanup_model (model)

    write (u, "(A)")
    write (u, "(A)")  "* Test output end: prc_omega_2"
    
  end subroutine prc_omega_2
  
\eatline
\nwendcode{}\nwbegindocs{328}\nwdocspar
\subsubsection{Check helicity selection}
The process is $e^- e^+ \to e^- e^+$ for vanishing masses.  We call
the matrix element several times to verify the switching off of
irrelevant helicities.
\nwenddocs{}\nwbegincode{329}\moddef{Omega interface: execute tests}\plusendmoddef
  call test (prc_omega_3, "prc_omega_3", &
       "helicity selection", &
       u, results)
\nwendcode{}\nwbegincode{330}\moddef{Omega interface: tests}\plusendmoddef
  subroutine prc_omega_3 (u)
    integer, intent(in) :: u
    type(process_library_t), target :: lib
    type(process_def_entry_t), pointer :: entry
    type(os_data_t) :: os_data
    type(string_t) :: model_name
    class(model_data_t), pointer :: model
    class(vars_t), pointer :: vars => null ()
    type(string_t), dimension(:), allocatable :: prt_in, prt_out
    type(process_component_def_t), pointer :: config
    type(prc_omega_t) :: prc1
    type(process_constants_t) :: data
    integer, parameter :: cdf = c_default_float
    real(cdf), dimension(:), allocatable :: par
    real(cdf), dimension(0:3,4) :: p
    type(helicity_selection_t) :: helicity_selection
    integer :: i, h
    
    write (u, "(A)")  "* Test output: prc_omega_3"
    write (u, "(A)")  "*   Purpose: create simple process with OMega"
    write (u, "(A)")  "*            and check helicity selection"
    write (u, "(A)")

    write (u, "(A)")  "* Initialize a process library."
    write (u, "(A)")  "* (1) e- e+ -> e- e+   (all diagrams, no OpenMP)"

    call lib%init (var_str ("omega3"))
    call os_data_init (os_data)

    model_name = "QED"
    call prepare_model (model, model_name, vars)

    allocate (prt_in (2), prt_out (2))
    prt_in = [var_str ("e-"), var_str ("e+")]
    prt_out = prt_in

    allocate (entry)
    call entry%init (var_str ("omega3_a"), &
         model, n_in = 2, n_components = 1)

    call omega_make_process_component (entry, 1, &
         model_name, prt_in, prt_out)
    call lib%append (entry)
    
    write (u, "(A)")
    write (u, "(A)")  "* Build and load library"

    call lib%configure (os_data)
    call lib%write_makefile (os_data, force = .true.)
    call lib%clean (os_data, distclean = .false.)
    call lib%write_driver (force = .true.)
    call lib%load (os_data)
    
    write (u, "(A)")
    write (u, "(A)")  "* Probe library API:"
    write (u, "(A)")
       
    write (u, "(1x,A,L1)")  "is active                 = ", &
         lib%is_active ()
    write (u, "(1x,A,I0)")  "n_processes               = ", &
         lib%get_n_processes ()

    write (u, "(A)")
    write (u, "(A)")  "* Set parameters for omega3_a and initialize:"
    write (u, "(A)")

    call vars%set_rval (var_str ("ee"), 0.3_default)
    call vars%set_rval (var_str ("me"), 0._default)
    call vars%set_rval (var_str ("mmu"), 0._default)
    call vars%set_rval (var_str ("mtau"), 0._default)
    allocate (par (model%get_n_real ()))
    call model%real_parameters_to_c_array (par)

    write (u, "(2x,A,F6.4)")  "ee   = ", par(1)
    write (u, "(2x,A,F6.4)")  "me   = ", par(2)
    write (u, "(2x,A,F6.4)")  "mmu  = ", par(3)
    write (u, "(2x,A,F6.4)")  "mtau = ", par(4)

    call prc1%set_parameters (model, helicity_selection)

    write (u, "(A)")
    write (u, "(A)")  "* Helicity states of omega3_a_i1:"
    write (u, "(A)")

    config => lib%get_component_def_ptr (var_str ("omega3_a"), 1)
    call prc1%init (config%get_core_def_ptr (), &
         lib, var_str ("omega3_a"), 1)
    call prc1%get_constants (data, 1)

    do i = 1, data%n_hel
       write (u, "(3x,I2,':',4(1x,I2))") i, data%hel_state(:,i)
    end do

    write (u, "(A)")
    write (u, "(A)")  "* Initially allowed helicities:"
    write (u, "(A)")
    
    write (u, "(4x,16(1x,I2))")  [(h, h = 1, data%n_hel)]
    write (u, "(4x)", advance = "no")
    do h = 1, data%n_hel
       write (u, "(2x,L1)", advance = "no")  prc1%is_allowed (1, 1, h, 1)
    end do
    write (u, "(A)")
    
    write (u, "(A)")
    write (u, "(A)")  "* Reset helicity selection (cutoff = 4)"
    write (u, "(A)")

    helicity_selection%active = .true.
    helicity_selection%threshold = 1e10_default
    helicity_selection%cutoff = 4
    call helicity_selection%write (u)
    
    call prc1%set_parameters (model, helicity_selection)
    call prc1%reset_helicity_selection ()

    write (u, "(A)")
    write (u, "(A)")  "* Allowed helicities:"
    write (u, "(A)")
    
    write (u, "(4x,16(1x,I2))")  [(h, h = 1, data%n_hel)]
    write (u, "(4x)", advance = "no")
    do h = 1, data%n_hel
       write (u, "(2x,L1)", advance = "no")  prc1%is_allowed (1, 1, h, 1)
    end do
    write (u, "(A)")
    
    write (u, "(A)")
    write (u, "(A)")  "* Set kinematics:"
    write (u, "(A)")
    
    p = reshape ([ &
         1.0_cdf, 0.0_cdf, 0.0_cdf, 1.0_cdf, &
         1.0_cdf, 0.0_cdf, 0.0_cdf,-1.0_cdf, &
         1.0_cdf, 1.0_cdf, 0.0_cdf, 0.0_cdf, &
         1.0_cdf,-1.0_cdf, 0.0_cdf, 0.0_cdf &
         ], [4,4])
    do i = 1, 4
       write (u, "(2x,A,I0,A,4(1x,F7.4))")  "p", i, " =", p(:,i)
    end do

    write (u, "(A)")
    write (u, "(A)")  "* Compute scattering matrix 5 times"
    write (u, "(A)")

    write (u, "(4x,16(1x,I2))")  [(h, h = 1, data%n_hel)]

    select type (driver => prc1%driver)
    type is (omega_driver_t)
       do i = 1, 5
          call driver%new_event (p)
          write (u, "(2x,I2)", advance = "no")  i
          do h = 1, data%n_hel
             write (u, "(2x,L1)", advance = "no")  prc1%is_allowed (1, 1, h, 1)
          end do
          write (u, "(A)")
       end do
    end select

    write (u, "(A)")
    write (u, "(A)")  "* Reset helicity selection again"
    write (u, "(A)")

    call prc1%activate_parameters ()

    write (u, "(A)")  "* Allowed helicities:"
    write (u, "(A)")
    
    write (u, "(4x,16(1x,I2))")  [(h, h = 1, data%n_hel)]
    write (u, "(4x)", advance = "no")
    do h = 1, data%n_hel
       write (u, "(2x,L1)", advance = "no")  prc1%is_allowed (1, 1, h, 1)
    end do
    write (u, "(A)")
    
    call lib%final ()
    call cleanup_model (model)
    
    write (u, "(A)")
    write (u, "(A)")  "* Test output end: prc_omega_3"
    
  end subroutine prc_omega_3
  
\eatline
\nwendcode{}\nwbegindocs{331}\nwdocspar
\subsubsection{QCD coupling}
The process is $u\bar u \to d\bar d$ for vanishing masses.  We compute
the amplitude for a fixed configuration once, then reset $\alpha_s$,
then compute again.

For {\tt{}GNU\ make}, {\tt{}makeflags} is set to {\tt{}-j1}.  This eliminates a
potential clash with a {\tt{}-j<n>} flag if this test is called from a
parallel make.
\nwenddocs{}\nwbegincode{332}\moddef{Omega interface: execute tests}\plusendmoddef
  call test (prc_omega_4, "prc_omega_4", &
       "update QCD alpha", &
       u, results)
\nwendcode{}\nwbegincode{333}\moddef{Omega interface: tests}\plusendmoddef
  subroutine prc_omega_4 (u)
    integer, intent(in) :: u
    type(process_library_t) :: lib
    class(prc_core_def_t), allocatable :: def
    type(process_def_entry_t), pointer :: entry
    type(os_data_t) :: os_data
    type(string_t) :: model_name
    type(string_t), dimension(:), allocatable :: prt_in, prt_out
    type(process_constants_t) :: data
    class(prc_core_driver_t), allocatable :: driver
    integer, parameter :: cdf = c_default_float
    integer, parameter :: ci = c_int
    real(cdf), dimension(6) :: par
    real(cdf), dimension(0:3,4) :: p
    logical(c_bool) :: flag
    complex(c_default_complex) :: amp
    integer :: i
    real(cdf) :: alpha_s
    
    write (u, "(A)")  "* Test output: prc_omega_4"
    write (u, "(A)")  "*   Purpose: create a QCD process with OMega"
    write (u, "(A)")  "*            and check alpha_s dependence"
    write (u, "(A)")

    write (u, "(A)")  "* Initialize a process library with one entry"
    write (u, "(A)")
    call lib%init (var_str ("prc_omega_4_lib"))
    call os_data_init (os_data)

    model_name = "QCD"
    allocate (prt_in (2), prt_out (2))
    prt_in = [var_str ("u"), var_str ("ubar")]
    prt_out = [var_str ("d"), var_str ("dbar")]
    
    allocate (omega_omega_def_t :: def)
    select type (def)
    type is (omega_omega_def_t)
       call def%init (model_name, prt_in, prt_out)
    end select
    allocate (entry)
    call entry%init (var_str ("prc_omega_4_p"), model_name = model_name, &
         n_in = 2, n_components = 1)
    call entry%import_component (1, n_out = size (prt_out), &
         prt_in  = new_prt_spec (prt_in), &
         prt_out = new_prt_spec (prt_out), &
         method  = var_str ("omega"), &
         variant = def)
    call lib%append (entry)
    
    write (u, "(A)")  "* Configure and compile process"
    write (u, "(A)")
    call lib%configure (os_data)
    call lib%write_makefile (os_data, force = .true.)
    call lib%clean (os_data, distclean = .false.)
    call lib%write_driver (force = .true.)
    call lib%load (os_data)
    
    write (u, "(A)")  "* Probe library API:"
    write (u, "(A)")
       
    write (u, "(1x,A,L1)")  "is active = ", lib%is_active ()

    write (u, "(A)")
    write (u, "(A)")  "* Set parameters:"
    write (u, "(A)")

    alpha_s = 0.1178_cdf
    
    par = [alpha_s, 0._cdf, 0._cdf, 0._cdf, 173.1_cdf, 1.523_cdf]
    write (u, "(2x,A,F8.4)")  "alpha_s = ", par(1)
    write (u, "(2x,A,F8.4)")  "ms      = ", par(2)
    write (u, "(2x,A,F8.4)")  "mc      = ", par(3)
    write (u, "(2x,A,F8.4)")  "mb      = ", par(4)
    write (u, "(2x,A,F8.4)")  "mtop    = ", par(5)
    write (u, "(2x,A,F8.4)")  "wtop    = ", par(6)

    write (u, "(A)")
    write (u, "(A)")  "* Set kinematics:"
    write (u, "(A)")
    
    p = reshape ([ &
         100.0_cdf, 0.0_cdf, 0.0_cdf, 100.0_cdf, &
         100.0_cdf, 0.0_cdf, 0.0_cdf,-100.0_cdf, &
         100.0_cdf, 100.0_cdf, 0.0_cdf, 0.0_cdf, &
         100.0_cdf,-100.0_cdf, 0.0_cdf, 0.0_cdf &
         ], [4,4])
    do i = 1, 4
       write (u, "(2x,A,I0,A,4(1x,F7.1))")  "p", i, " =", p(:,i)
    end do

    call lib%connect_process (var_str ("prc_omega_4_p"), 1, data, driver)

    select type (driver)
    type is (omega_driver_t)
       call driver%init (par)

       write (u, "(A)")
       write (u, "(A)")  "* Compute matrix element:"
       write (u, "(A)")

       call driver%new_event (p)

       call driver%is_allowed (1_ci, 6_ci, 1_ci, flag)
       write (u, "(1x,A,L1)") "is_allowed (1, 6, 1) = ", flag
       
       call driver%get_amplitude (1_ci, 6_ci, 1_ci, amp)
       write (u, "(1x,A,1x,E11.4)") "|amp (1, 6, 1)| =", abs (amp)

       write (u, "(A)")
       write (u, "(A)")  "* Double alpha_s and compute matrix element again:"
       write (u, "(A)")

       call driver%update_alpha_s (2 * alpha_s)
       call driver%new_event (p)

       call driver%is_allowed (1_ci, 6_ci, 1_ci, flag)
       write (u, "(1x,A,L1)") "is_allowed (1, 6, 1) = ", flag
       
       call driver%get_amplitude (1_ci, 6_ci, 1_ci, amp)
       write (u, "(1x,A,1x,E11.4)") "|amp (1, 6, 1)| =", abs (amp)
    end select

    call lib%final ()
    
    write (u, "(A)")
    write (u, "(A)")  "* Test output end: prc_omega_4"
    
  end subroutine prc_omega_4
  
\eatline
\nwendcode{}\nwbegindocs{334}\nwdocspar
\subsubsection{Amplitude and QCD coupling}
The same process as before.  Here, we initialize with a running $\alpha_s$
coupling and compute twice with different scales.  We use the high-level
method {\tt{}compute{\char95}amplitude}.
\nwenddocs{}\nwbegincode{335}\moddef{Omega interface: execute tests}\plusendmoddef
  call test (prc_omega_5, "prc_omega_5", &
       "running QCD alpha", &
       u, results)
\nwendcode{}\nwbegincode{336}\moddef{Omega interface: tests}\plusendmoddef
  subroutine prc_omega_5 (u)
    integer, intent(in) :: u
    type(process_library_t) :: lib
    class(prc_core_def_t), allocatable :: def
    type(process_component_def_t), pointer :: cdef_ptr
    class(prc_core_def_t), pointer :: def_ptr
    type(process_def_entry_t), pointer :: entry
    type(os_data_t) :: os_data
    class(model_data_t), pointer :: model
    type(string_t) :: model_name
    type(string_t), dimension(:), allocatable :: prt_in, prt_out
    type(qcd_t) :: qcd
    class(prc_core_t), allocatable :: core
    class(workspace_t), allocatable :: tmp
    type(vector4_t), dimension(4) :: p
    complex(default) :: amp
    real(default) :: fac_scale
    real(default), allocatable :: alpha_qcd_forced
    integer :: i
    
    write (u, "(A)")  "* Test output: prc_omega_5"
    write (u, "(A)")  "*   Purpose: create a QCD process with OMega"
    write (u, "(A)")  "*            and check alpha_s dependence"
    write (u, "(A)")

    write (u, "(A)")  "* Initialize a process library with one entry"
    write (u, "(A)")
    call lib%init (var_str ("prc_omega_5_lib"))
    call os_data_init (os_data)

    model_name = "QCD"
    call prepare_model (model, model_name)

    allocate (prt_in (2), prt_out (2))
    prt_in = [var_str ("u"), var_str ("ubar")]
    prt_out = [var_str ("d"), var_str ("dbar")]
    
    allocate (omega_omega_def_t :: def)
    select type (def)
    type is (omega_omega_def_t)
       call def%init (model_name, prt_in, prt_out)
    end select
    allocate (entry)
    call entry%init (var_str ("prc_omega_5_p"), model_name = model_name, &
         n_in = 2, n_components = 1)
    call entry%import_component (1, n_out = size (prt_out), &
         prt_in  = new_prt_spec (prt_in), &
         prt_out = new_prt_spec (prt_out), &
         method  = var_str ("omega"), &
         variant = def)
    call lib%append (entry)
    
    write (u, "(A)")  "* Configure and compile process"
    write (u, "(A)")
    call lib%configure (os_data)
    call lib%write_makefile (os_data, force = .true.)
    call lib%clean (os_data, distclean = .false.)
    call lib%write_driver (force = .true.)
    call lib%load (os_data)
    
    write (u, "(A)")  "* Probe library API"
    write (u, "(A)")
       
    write (u, "(1x,A,L1)")  "is active = ", lib%is_active ()

    write (u, "(A)")
    write (u, "(A)")  "* Set kinematics"
    write (u, "(A)")
    
    p(1) = vector4_moving (100._default, 100._default, 3)
    p(2) = vector4_moving (100._default,-100._default, 3)
    p(3) = vector4_moving (100._default, 100._default, 1)
    p(4) = vector4_moving (100._default,-100._default, 1)
    do i = 1, 4
       call vector4_write (p(i), u)
    end do

    write (u, "(A)")
    write (u, "(A)")  "* Setup QCD data"
    write (u, "(A)")
    
    allocate (alpha_qcd_from_scale_t :: qcd%alpha)
    
    write (u, "(A)")  "* Setup process core"
    write (u, "(A)")
    
    allocate (prc_omega_t :: core)
    cdef_ptr => lib%get_component_def_ptr (var_str ("prc_omega_5_p"), 1)
    def_ptr => cdef_ptr%get_core_def_ptr ()

    select type (core)
    type is (prc_omega_t)
       call core%allocate_workspace (tmp)
       call core%set_parameters (model, qcd = qcd)
       call core%init (def_ptr, lib, var_str ("prc_omega_5_p"), 1)
       call core%write (u)

       write (u, "(A)")
       write (u, "(A)")  "* Compute matrix element"
       write (u, "(A)")

       fac_scale = 100
       write (u, "(1x,A,F4.0)")  "factorization scale = ", fac_scale

       amp = core%compute_amplitude &
            (1, p, 1, 6, 1, fac_scale, 100._default, alpha_qcd_forced)

       write (u, "(1x,A,1x,E11.4)") "|amp (1, 6, 1)| =", abs (amp)

       write (u, "(A)")
       write (u, "(A)")  "* Modify factorization scale and &
            &compute matrix element again"
       write (u, "(A)")

       fac_scale = 200
       write (u, "(1x,A,F4.0)")  "factorization scale = ", fac_scale

       amp = core%compute_amplitude &
            (1, p, 1, 6, 1, fac_scale, 100._default, alpha_qcd_forced)

       write (u, "(1x,A,1x,E11.4)") "|amp (1, 6, 1)| =", abs (amp)

       write (u, "(A)")
       write (u, "(A)")  "* Set alpha(QCD) directly and &
            &compute matrix element again"
       write (u, "(A)")

       allocate (alpha_qcd_forced, source = 0.1_default)
       write (u, "(1x,A,F6.4)")  "alpha_qcd = ", alpha_qcd_forced

       amp = core%compute_amplitude &
            (1, p, 1, 6, 1, fac_scale, 100._default, alpha_qcd_forced)

       write (u, "(1x,A,1x,E11.4)") "|amp (1, 6, 1)| =", abs (amp)

    end select

    call lib%final ()
    call cleanup_model (model)
    
    write (u, "(A)")
    write (u, "(A)")  "* Test output end: prc_omega_5"
    
  end subroutine prc_omega_5
  
\eatline
\nwendcode{}\nwbegindocs{337}\nwdocspar
\subsubsection{Generate matrix element diagrams}
The same process as before. No amplitude is computed here, instead we just
generate Feynman (and color flow) diagrams, and check whether PS and PDF
files have been generated. This test is only run if event analysis is 
possible.
\nwenddocs{}\nwbegincode{338}\moddef{Omega interface: execute diags tests}\endmoddef
  call test (prc_omega_diags_1, "prc_omega_diags_1", &
       "generate Feynman diagrams", &
       u, results)
\nwendcode{}\nwbegincode{339}\moddef{Omega interface: tests}\plusendmoddef
  subroutine prc_omega_diags_1 (u)
    integer, intent(in) :: u
    type(process_library_t) :: lib
    class(prc_core_def_t), allocatable :: def
    type(process_def_entry_t), pointer :: entry
    type(os_data_t) :: os_data
    type(string_t) :: model_name
    type(string_t), dimension(:), allocatable :: prt_in, prt_out
    type(string_t) :: diags_file, pdf_file, ps_file
    logical :: exist, exist_pdf, exist_ps
    integer :: iostat, u_diags
    character(128) :: buffer    
    
    write (u, "(A)")  "* Test output: prc_omega_diags_1"
    write (u, "(A)")  "*   Purpose: generate Feynman diagrams"
    write (u, "(A)")

    write (u, "(A)")  "* Initialize a process library with one entry"
    write (u, "(A)")
    call lib%init (var_str ("prc_omega_diags_1_lib"))
    call os_data_init (os_data)

    model_name = "SM"

    allocate (prt_in (2), prt_out (2))
    prt_in = [var_str ("u"), var_str ("ubar")]
    prt_out = [var_str ("d"), var_str ("dbar")]
    
    allocate (omega_omega_def_t :: def)
    select type (def)
    type is (omega_omega_def_t)
       call def%init (model_name, prt_in, prt_out, &
            diags = .true., diags_color = .true.)
    end select
    allocate (entry)
    call entry%init (var_str ("prc_omega_diags_1_p"), model_name = model_name, &
         n_in = 2, n_components = 1)
    call entry%import_component (1, n_out = size (prt_out), &
         prt_in  = new_prt_spec (prt_in), &
         prt_out = new_prt_spec (prt_out), &
         method  = var_str ("omega"), &
         variant = def)
    call lib%append (entry)
    
    write (u, "(A)")  "* Configure and compile process"
    write (u, "(A)")  "    and generate diagrams"
    write (u, "(A)")
    call lib%configure (os_data)
    call lib%write_makefile (os_data, force = .true., testflag = .true.)
    call lib%clean (os_data, distclean = .false.)
    call lib%write_driver (force = .true.)
    call lib%load (os_data)
    
    write (u, "(A)")  "* Probe library API"
    write (u, "(A)")
       
    write (u, "(1x,A,L1)")  "is active = ", lib%is_active ()

    write (u, "(A)")  "* Check produced diagram files"
    write (u, "(A)")        

    diags_file = "prc_omega_diags_1_p_i1_diags.tex"
    ps_file  = "prc_omega_diags_1_p_i1_diags.ps"
    pdf_file = "prc_omega_diags_1_p_i1_diags.pdf"    
    inquire (file = char (diags_file), exist = exist)
    if (exist) then
       u_diags = free_unit ()
       open (u_diags, file = char (diags_file), action = "read", status = "old")
       iostat = 0
       do while (iostat == 0)
          read (u_diags, "(A)", iostat = iostat)  buffer
          if (iostat == 0)  write (u, "(A)")  trim (buffer)
       end do
       close (u_diags)
    else
       write (u, "(A)")  "[Feynman diagrams LaTeX file is missing]"
    end if
    inquire (file = char (ps_file), exist = exist_ps)
    if (exist_ps) then
       write (u, "(A)")  "[Feynman diagrams Postscript file exists and is nonempty]"
    else
       write (u, "(A)")  "[Feynman diagrams Postscript file is missing/non-regular]"
    end if
    inquire (file = char (pdf_file), exist = exist_pdf)
    if (exist_pdf) then
       write (u, "(A)")  "[Feynman diagrams PDF file exists and is nonempty]"
    else
       write (u, "(A)")  "[Feynman diagrams PDF file is missing/non-regular]"
    end if               
    
    write (u, "(A)")
    write (u, "(A)")  "* Cleanup"
    write (u, "(A)")    
    
    call lib%final ()

    write (u, "(A)")
    write (u, "(A)")  "* Test output end: prc_omega_diags_1"
    
  end subroutine prc_omega_diags_1
  
\eatline
\nwendcode{}\nwfilename{../../../src/events/events.nw}\nwbegindocs{0}%% -*- ess-noweb-default-code-mode: f90-mode; noweb-default-code-mode: f90-mode; -*- 
% WHIZARD code as NOWEB source: event handling objects

\nwenddocs{}\nwbegindocs{1}\nwdocspar
%%%%%%%%%%%%%%%%%%%%%%%%%%%%%%%%%%%%%%%%%%%%%%%%%%%%%%%%%%%%%%%%%%%%%%%% 
\chapter{Generic Event Handling}

Event records allow the MC to communicate with the outside world.  The
event record should exhibit the observable contents of a physical
event.  We should be able to read and write events.  The actual
implementation of the event need not be defined yet, for that purpose.

We have the following basic modules:
\begin{description}
\item[event\_base]
  Abstract base type for event records.  The base type contains a
  reference to a {\tt{}particle{\char95}set{\char95}t} object as the event core, and it
  holds some data that we should always expect, such as the squared
  matrix element and event weight.
\item[eio\_data]
  Transparent container for the metadata of an event sample.
\item[eio\_base]
  Abstract base type for event-record input and output.  The
  implementations of this base type represent specific event I/O
  formats. 
\end{description}

These are the implementation modules:
\begin{description}
\item[eio\_checkpoints]
  Auxiliary output format.  The only purpose is to provide screen
  diagnostics during event output.
\item[eio\_weights]
  Print some event summary data, no details.  The main use if for
  testing purposes.
\item[hep\_common]
  Implements traditional HEP common blocks that are (still) used by
  some of the event I/O formats below.
\item[hepmc\_interface]
  Access particle objects of the HepMC package.  Functional only if this
  package is linked.    
\item[lcio\_interface]
  Access objects of the LCIO package.  Functional only if this
  package is linked.    
\item[hep\_events]
  Interface between the event record and the common blocks.
\item[eio\_ascii]
  Collection of event output formats that write ASCII files.
\item[eio\_lhef]
  LHEF for input and output.
\item[eio\_stdhep]
  Support for the StdHEP format (binary, machine-independent).
\item[eio\_hepmc]
  Support for the HepMC format (C++).
\item[eio\_lcio]
  Support for the LCIO format (C++).
\end{description}

%%%%%%%%%%%%%%%%%%%%%%%%%%%%%%%%%%%%%%%%%%%%%%%%%%%%%%%%%%%%%%%%%%%%%%%%
\section{Generic Event Handling}
We introduce events first in form of an abstract type, together with
some utilities.  Abstract events can be used by other modules, in
particular event I/O, without introducing an explicit dependency on
the event implementation.
\nwenddocs{}\nwbegincode{2}\moddef{\code{}event{\char95}base.f90\edoc{}}\endmoddef
\LA{}File header\RA{}

module event_base
  
  use kinds
\LA{}Use strings\RA{}
  use string_utils, only: lower_case
  use diagnostics
  use model_data
  use particles

\LA{}Standard module head\RA{}

\LA{}Event base: public\RA{}

\LA{}Event base: parameters\RA{}

\LA{}Event base: types\RA{}

\LA{}Event base: interfaces\RA{}

contains
  
\LA{}Event base: procedures\RA{}

end module event_base
\eatline
\nwendcode{}\nwbegindocs{3}\nwdocspar
\subsection{generic event type}
\nwenddocs{}\nwbegincode{4}\moddef{Event base: public}\endmoddef
  public :: generic_event_t
\nwendcode{}\nwbegincode{5}\moddef{Event base: types}\endmoddef
  type, abstract :: generic_event_t
     private
     logical :: particle_set_is_valid = .false.
     type(particle_set_t), pointer :: particle_set => null ()
     logical :: sqme_ref_known = .false.
     real(default) :: sqme_ref = 0
     logical :: sqme_prc_known = .false.
     real(default) :: sqme_prc = 0
     logical :: weight_ref_known = .false.
     real(default) :: weight_ref = 0
     logical :: weight_prc_known = .false.
     real(default) :: weight_prc = 0
     logical :: excess_prc_known = .false.
     real(default) :: excess_prc = 0
     integer :: n_alt = 0
     logical :: sqme_alt_known = .false.
     real(default), dimension(:), allocatable :: sqme_alt
     logical :: weight_alt_known = .false.
     real(default), dimension(:), allocatable :: weight_alt
   contains
   \LA{}Event base: generic event: TBP\RA{}
  end type generic_event_t
  
\eatline
\nwendcode{}\nwbegindocs{6}\nwdocspar
\subsection{Initialization}
This determines the number of alternate weights and sqme values.
\nwenddocs{}\nwbegincode{7}\moddef{Event base: generic event: TBP}\endmoddef
  procedure :: base_init => generic_event_init
\nwendcode{}\nwbegincode{8}\moddef{Event base: procedures}\endmoddef
  subroutine generic_event_init (event, n_alt)
    class(generic_event_t), intent(out) :: event
    integer, intent(in) :: n_alt
    event%n_alt = n_alt
    allocate (event%sqme_alt (n_alt))
    allocate (event%weight_alt (n_alt))
  end subroutine generic_event_init
  
\eatline
\nwendcode{}\nwbegindocs{9}\nwdocspar
\subsection{Access particle set}
The particle set is the core of the event.  We allow access to it via
a pointer, and we maintain the information whether the particle set
is valid, i.e., has been filled with meaningful data.
\nwenddocs{}\nwbegincode{10}\moddef{Event base: generic event: TBP}\plusendmoddef
  procedure :: has_valid_particle_set => generic_event_has_valid_particle_set
  procedure :: accept_particle_set => generic_event_accept_particle_set
  procedure :: discard_particle_set => generic_event_discard_particle_set
\nwendcode{}\nwbegincode{11}\moddef{Event base: procedures}\plusendmoddef
  function generic_event_has_valid_particle_set (event) result (flag)
    class(generic_event_t), intent(in) :: event
    logical :: flag
    flag = event%particle_set_is_valid
  end function generic_event_has_valid_particle_set
  
  subroutine generic_event_accept_particle_set (event)
    class(generic_event_t), intent(inout) :: event
    event%particle_set_is_valid = .true.
  end subroutine generic_event_accept_particle_set
  
  subroutine generic_event_discard_particle_set (event)
    class(generic_event_t), intent(inout) :: event
    event%particle_set_is_valid = .false.
  end subroutine generic_event_discard_particle_set
  
\eatline
\eatline
\eatline
\nwendcode{}\nwbegindocs{12}\nwdocspar
These procedures deal with the particle set directly.  Return the pointer:
\nwenddocs{}\nwbegincode{13}\moddef{Event base: generic event: TBP}\plusendmoddef
  procedure :: get_particle_set_ptr => generic_event_get_particle_set_ptr
\nwendcode{}\nwbegincode{14}\moddef{Event base: procedures}\plusendmoddef
  function generic_event_get_particle_set_ptr (event) result (ptr)
    class(generic_event_t), intent(in) :: event
    type(particle_set_t), pointer :: ptr
    ptr => event%particle_set
  end function generic_event_get_particle_set_ptr
  
\eatline
\nwendcode{}\nwbegindocs{15}\nwdocspar
Let it point to some existing particle set:
\nwenddocs{}\nwbegincode{16}\moddef{Event base: generic event: TBP}\plusendmoddef
  procedure :: link_particle_set => generic_event_link_particle_set
\nwendcode{}\nwbegincode{17}\moddef{Event base: procedures}\plusendmoddef
  subroutine generic_event_link_particle_set (event, particle_set)
    class(generic_event_t), intent(inout) :: event
    type(particle_set_t), intent(in), target :: particle_set
    event%particle_set => particle_set
    call event%accept_particle_set ()
  end subroutine generic_event_link_particle_set

\eatline
\nwendcode{}\nwbegindocs{18}\nwdocspar
\subsection{Access sqme and weight}
There are several incarnations: the current value, a reference value,
alternate values.
\nwenddocs{}\nwbegincode{19}\moddef{Event base: generic event: TBP}\plusendmoddef
  procedure :: sqme_prc_is_known => generic_event_sqme_prc_is_known
  procedure :: sqme_ref_is_known => generic_event_sqme_ref_is_known
  procedure :: sqme_alt_is_known => generic_event_sqme_alt_is_known
  procedure :: weight_prc_is_known => generic_event_weight_prc_is_known
  procedure :: weight_ref_is_known => generic_event_weight_ref_is_known
  procedure :: weight_alt_is_known => generic_event_weight_alt_is_known
  procedure :: excess_prc_is_known => generic_event_excess_prc_is_known
\nwendcode{}\nwbegincode{20}\moddef{Event base: procedures}\plusendmoddef
  function generic_event_sqme_prc_is_known (event) result (flag)
    class(generic_event_t), intent(in) :: event
    logical :: flag
    flag = event%sqme_prc_known
  end function generic_event_sqme_prc_is_known

  function generic_event_sqme_ref_is_known (event) result (flag)
    class(generic_event_t), intent(in) :: event
    logical :: flag
    flag = event%sqme_ref_known
  end function generic_event_sqme_ref_is_known

  function generic_event_sqme_alt_is_known (event) result (flag)
    class(generic_event_t), intent(in) :: event
    logical :: flag
    flag = event%sqme_alt_known
  end function generic_event_sqme_alt_is_known

  function generic_event_weight_prc_is_known (event) result (flag)
    class(generic_event_t), intent(in) :: event
    logical :: flag
    flag = event%weight_prc_known
  end function generic_event_weight_prc_is_known

  function generic_event_weight_ref_is_known (event) result (flag)
    class(generic_event_t), intent(in) :: event
    logical :: flag
    flag = event%weight_ref_known
  end function generic_event_weight_ref_is_known

  function generic_event_weight_alt_is_known (event) result (flag)
    class(generic_event_t), intent(in) :: event
    logical :: flag
    flag = event%weight_alt_known
  end function generic_event_weight_alt_is_known

  function generic_event_excess_prc_is_known (event) result (flag)
    class(generic_event_t), intent(in) :: event
    logical :: flag
    flag = event%excess_prc_known
  end function generic_event_excess_prc_is_known

\eatline
\eatline
\eatline
\eatline
\eatline
\eatline
\eatline
\nwendcode{}\nwbegindocs{21}\nwdocspar
\nwenddocs{}\nwbegincode{22}\moddef{Event base: generic event: TBP}\plusendmoddef
  procedure :: get_n_alt => generic_event_get_n_alt
\nwendcode{}\nwbegincode{23}\moddef{Event base: procedures}\plusendmoddef
  function generic_event_get_n_alt (event) result (n)
    class(generic_event_t), intent(in) :: event
    integer :: n
    n = event%n_alt
  end function generic_event_get_n_alt
  
\eatline
\nwendcode{}\nwbegindocs{24}\nwdocspar
\nwenddocs{}\nwbegincode{25}\moddef{Event base: generic event: TBP}\plusendmoddef
  procedure :: get_sqme_prc => generic_event_get_sqme_prc
  procedure :: get_sqme_ref => generic_event_get_sqme_ref
  generic :: get_sqme_alt => &
       generic_event_get_sqme_alt_0, generic_event_get_sqme_alt_1
  procedure :: generic_event_get_sqme_alt_0
  procedure :: generic_event_get_sqme_alt_1
  procedure :: get_weight_prc => generic_event_get_weight_prc
  procedure :: get_weight_ref => generic_event_get_weight_ref
  generic :: get_weight_alt => &
       generic_event_get_weight_alt_0, generic_event_get_weight_alt_1
  procedure :: generic_event_get_weight_alt_0
  procedure :: generic_event_get_weight_alt_1
  procedure :: get_excess_prc => generic_event_get_excess_prc
\nwendcode{}\nwbegincode{26}\moddef{Event base: procedures}\plusendmoddef
  function generic_event_get_sqme_prc (event) result (sqme)
    class(generic_event_t), intent(in) :: event
    real(default) :: sqme
    if (event%sqme_prc_known) then
       sqme = event%sqme_prc
    else
       sqme = 0
    end if
  end function generic_event_get_sqme_prc

  function generic_event_get_sqme_ref (event) result (sqme)
    class(generic_event_t), intent(in) :: event
    real(default) :: sqme
    if (event%sqme_ref_known) then
       sqme = event%sqme_ref
    else
       sqme = 0
    end if
  end function generic_event_get_sqme_ref

  function generic_event_get_sqme_alt_0 (event, i) result (sqme)
    class(generic_event_t), intent(in) :: event
    integer, intent(in) :: i
    real(default) :: sqme
    if (event%sqme_alt_known) then
       sqme = event%sqme_alt(i)
    else
       sqme = 0
    end if
  end function generic_event_get_sqme_alt_0

  function generic_event_get_sqme_alt_1 (event) result (sqme)
    class(generic_event_t), intent(in) :: event
    real(default), dimension(event%n_alt) :: sqme
    sqme = event%sqme_alt
  end function generic_event_get_sqme_alt_1

  function generic_event_get_weight_prc (event) result (weight)
    class(generic_event_t), intent(in) :: event
    real(default) :: weight
    if (event%weight_prc_known) then
       weight = event%weight_prc
    else
       weight = 0
    end if
  end function generic_event_get_weight_prc

  function generic_event_get_weight_ref (event) result (weight)
    class(generic_event_t), intent(in) :: event
    real(default) :: weight
    if (event%weight_ref_known) then
       weight = event%weight_ref
    else
       weight = 0
    end if
  end function generic_event_get_weight_ref

  function generic_event_get_weight_alt_0 (event, i) result (weight)
    class(generic_event_t), intent(in) :: event
    integer, intent(in) :: i
    real(default) :: weight
    if (event%weight_alt_known) then
       weight = event%weight_alt(i)
    else
       weight = 0
    end if
  end function generic_event_get_weight_alt_0

  function generic_event_get_weight_alt_1 (event) result (weight)
    class(generic_event_t), intent(in) :: event
    real(default), dimension(event%n_alt) :: weight
    weight = event%weight_alt
  end function generic_event_get_weight_alt_1

  function generic_event_get_excess_prc (event) result (excess)
    class(generic_event_t), intent(in) :: event
    real(default) :: excess
    if (event%excess_prc_known) then
       excess = event%excess_prc
    else
       excess = 0
    end if
  end function generic_event_get_excess_prc

\eatline
\eatline
\eatline
\eatline
\eatline
\eatline
\eatline
\nwendcode{}\nwbegindocs{27}\nwdocspar
\nwenddocs{}\nwbegincode{28}\moddef{Event base: generic event: TBP}\plusendmoddef
  procedure :: set_sqme_prc => generic_event_set_sqme_prc
  procedure :: set_sqme_ref => generic_event_set_sqme_ref
  procedure :: set_sqme_alt => generic_event_set_sqme_alt
  procedure :: set_weight_prc => generic_event_set_weight_prc
  procedure :: set_weight_ref => generic_event_set_weight_ref
  procedure :: set_weight_alt => generic_event_set_weight_alt
  procedure :: set_excess_prc => generic_event_set_excess_prc
\nwendcode{}\nwbegincode{29}\moddef{Event base: procedures}\plusendmoddef
  subroutine generic_event_set_sqme_prc (event, sqme)
    class(generic_event_t), intent(inout) :: event
    real(default), intent(in) :: sqme
    event%sqme_prc = sqme
    event%sqme_prc_known = .true.
  end subroutine generic_event_set_sqme_prc

  subroutine generic_event_set_sqme_ref (event, sqme)
    class(generic_event_t), intent(inout) :: event
    real(default), intent(in) :: sqme
    event%sqme_ref = sqme
    event%sqme_ref_known = .true.
  end subroutine generic_event_set_sqme_ref

  subroutine generic_event_set_sqme_alt (event, sqme)
    class(generic_event_t), intent(inout) :: event
    real(default), dimension(:), intent(in) :: sqme
    event%sqme_alt = sqme
    event%sqme_alt_known = .true.
  end subroutine generic_event_set_sqme_alt

  subroutine generic_event_set_weight_prc (event, weight)
    class(generic_event_t), intent(inout) :: event
    real(default), intent(in) :: weight
    event%weight_prc = weight
    event%weight_prc_known = .true.
  end subroutine generic_event_set_weight_prc

  subroutine generic_event_set_weight_ref (event, weight)
    class(generic_event_t), intent(inout) :: event
    real(default), intent(in) :: weight
    event%weight_ref = weight
    event%weight_ref_known = .true.
  end subroutine generic_event_set_weight_ref

  subroutine generic_event_set_weight_alt (event, weight)
    class(generic_event_t), intent(inout) :: event
    real(default), dimension(:), intent(in) :: weight
    event%weight_alt = weight
    event%weight_alt_known = .true.
  end subroutine generic_event_set_weight_alt

  subroutine generic_event_set_excess_prc (event, excess)
    class(generic_event_t), intent(inout) :: event
    real(default), intent(in) :: excess
    event%excess_prc = excess
    event%excess_prc_known = .true.
  end subroutine generic_event_set_excess_prc

\eatline
\eatline
\eatline
\eatline
\eatline
\eatline
\eatline
\nwendcode{}\nwbegindocs{30}\nwdocspar
Set the appropriate entry directly.
\nwenddocs{}\nwbegincode{31}\moddef{Event base: generic event: TBP}\plusendmoddef
  procedure :: set => generic_event_set
\nwendcode{}\nwbegincode{32}\moddef{Event base: procedures}\plusendmoddef
  subroutine generic_event_set (event, &
       weight_ref, weight_prc, weight_alt, &
       excess_prc, &
       sqme_ref, sqme_prc, sqme_alt)
    class(generic_event_t), intent(inout) :: event
    real(default), intent(in), optional :: weight_ref, weight_prc
    real(default), intent(in), optional :: sqme_ref, sqme_prc
    real(default), dimension(:), intent(in), optional :: sqme_alt, weight_alt
    real(default), intent(in), optional :: excess_prc
    if (present (sqme_prc)) then
       call event%set_sqme_prc (sqme_prc)
    end if
    if (present (sqme_ref)) then
       call event%set_sqme_ref (sqme_ref)
    end if 
    if (present (sqme_alt)) then
       call event%set_sqme_alt (sqme_alt)
    end if
    if (present (weight_prc)) then
       call event%set_weight_prc (weight_prc)
    end if
    if (present (weight_ref)) then
       call event%set_weight_ref (weight_ref)
    end if
    if (present (weight_alt)) then
       call event%set_weight_alt (weight_alt)
    end if
    if (present (excess_prc)) then
       call event%set_excess_prc (excess_prc)
    end if
  end subroutine generic_event_set
  
\eatline
\nwendcode{}\nwbegindocs{33}\nwdocspar
\subsection{Pure Virtual Methods}
These procedures can only implemented in the concrete implementation.

Output (verbose, depending on parameters).
\nwenddocs{}\nwbegincode{34}\moddef{Event base: generic event: TBP}\plusendmoddef
  procedure (generic_event_write), deferred :: write
\nwendcode{}\nwbegincode{35}\moddef{Event base: interfaces}\endmoddef
  abstract interface
     subroutine generic_event_write (object, unit, &
          show_process, show_transforms, &
          show_decay, verbose, testflag)
       import
       class(generic_event_t), intent(in) :: object
       integer, intent(in), optional :: unit
       logical, intent(in), optional :: show_process
       logical, intent(in), optional :: show_transforms
       logical, intent(in), optional :: show_decay
       logical, intent(in), optional :: verbose
       logical, intent(in), optional :: testflag
     end subroutine generic_event_write
  end interface

\eatline
\nwendcode{}\nwbegindocs{36}\nwdocspar
Generate an event, based on a selector index {\tt{}i{\char95}mci}, and optionally on an
extra set of random numbers {\tt{}r}.  For the main bunch of random numbers that
the generator needs, the event object should contain its own generator.
\nwenddocs{}\nwbegincode{37}\moddef{Event base: generic event: TBP}\plusendmoddef
  procedure (generic_event_generate), deferred :: generate
\nwendcode{}\nwbegincode{38}\moddef{Event base: interfaces}\plusendmoddef
  abstract interface
     subroutine generic_event_generate (event, i_mci, r)
       import
       class(generic_event_t), intent(inout) :: event
       integer, intent(in) :: i_mci
       real(default), dimension(:), intent(in), optional :: r
     end subroutine generic_event_generate
  end interface
     
\eatline
\nwendcode{}\nwbegindocs{39}\nwdocspar
Alternative : inject a particle set that is supposed to represent the hard
process. 
How this determines the event, is dependent on the event structure,
therefore this is a deferred method.
\nwenddocs{}\nwbegincode{40}\moddef{Event base: generic event: TBP}\plusendmoddef
  procedure (generic_event_set_hard_particle_set), deferred :: &
       set_hard_particle_set 
\nwendcode{}\nwbegincode{41}\moddef{Event base: interfaces}\plusendmoddef
  abstract interface
     subroutine generic_event_set_hard_particle_set (event, particle_set)
       import
       class(generic_event_t), intent(inout) :: event
       type(particle_set_t), intent(in) :: particle_set
     end subroutine generic_event_set_hard_particle_set
  end interface

\eatline
\nwendcode{}\nwbegindocs{42}Evaluate any expressions associated with the event.  No argument needed.
\nwenddocs{}\nwbegincode{43}\moddef{Event base: generic event: TBP}\plusendmoddef
  procedure (generic_event_handler), deferred :: evaluate_expressions
\nwendcode{}\nwbegincode{44}\moddef{Event base: interfaces}\plusendmoddef
  abstract interface
     subroutine generic_event_handler (event)
       import
       class(generic_event_t), intent(inout) :: event
     end subroutine generic_event_handler
  end interface
  
\eatline
\nwendcode{}\nwbegindocs{45}\nwdocspar
Select internal parameters
\nwenddocs{}\nwbegincode{46}\moddef{Event base: generic event: TBP}\plusendmoddef
  procedure (generic_event_select), deferred :: select
\nwendcode{}\nwbegincode{47}\moddef{Event base: interfaces}\plusendmoddef
  abstract interface
     subroutine generic_event_select (event,  i_mci, i_term, channel)
       import
       class(generic_event_t), intent(inout) :: event
       integer, intent(in) :: i_mci, i_term, channel
     end subroutine generic_event_select
  end interface
  
\eatline
\nwendcode{}\nwbegindocs{48}Return a pointer to the model for the currently active process.
\nwenddocs{}\nwbegincode{49}\moddef{Event base: generic event: TBP}\plusendmoddef
  procedure (generic_event_get_model_ptr), deferred :: get_model_ptr
\nwendcode{}\nwbegincode{50}\moddef{Event base: interfaces}\plusendmoddef
  abstract interface
     function generic_event_get_model_ptr (event) result (model)
       import
       class(generic_event_t), intent(in) :: event
       class(model_data_t), pointer :: model
     end function generic_event_get_model_ptr
  end interface
  
\eatline
\nwendcode{}\nwbegindocs{51}Return data used by external event formats.
\nwenddocs{}\nwbegincode{52}\moddef{Event base: generic event: TBP}\plusendmoddef
  procedure (generic_event_get_index), deferred :: get_index
  procedure (generic_event_get_fac_scale), deferred :: get_fac_scale
  procedure (generic_event_get_alpha_s), deferred :: get_alpha_s
\nwendcode{}\nwbegincode{53}\moddef{Event base: interfaces}\plusendmoddef
  abstract interface
     function generic_event_get_index (event) result (index)
       import
       class(generic_event_t), intent(in) :: event
       integer :: index
     end function generic_event_get_index
  end interface
    
  abstract interface
     function generic_event_get_fac_scale (event) result (fac_scale)
       import
       class(generic_event_t), intent(in) :: event
       real(default) :: fac_scale
     end function generic_event_get_fac_scale
  end interface
    
  abstract interface
     function generic_event_get_alpha_s (event) result (alpha_s)
       import
       class(generic_event_t), intent(in) :: event
       real(default) :: alpha_s
     end function generic_event_get_alpha_s
  end interface

\eatline
\eatline
\eatline
\nwendcode{}\nwbegindocs{54}Set data used by external event formats.
\nwenddocs{}\nwbegincode{55}\moddef{Event base: generic event: TBP}\plusendmoddef
  procedure (generic_event_set_alpha_qcd_forced), deferred :: &
       set_alpha_qcd_forced
  procedure (generic_event_set_scale_forced), deferred :: &
       set_scale_forced
\nwendcode{}\nwbegincode{56}\moddef{Event base: interfaces}\plusendmoddef
  abstract interface
     subroutine generic_event_set_alpha_qcd_forced (event, alpha_qcd)
       import
       class(generic_event_t), intent(inout) :: event
       real(default), intent(in) :: alpha_qcd
     end subroutine generic_event_set_alpha_qcd_forced
  end interface

  abstract interface
     subroutine generic_event_set_scale_forced (event, scale)
       import
       class(generic_event_t), intent(inout) :: event
       real(default), intent(in) :: scale
     end subroutine generic_event_set_scale_forced
  end interface

\eatline
\eatline
\nwendcode{}\nwbegindocs{57}\nwdocspar
\subsection{Utilities}
Applying this, current event contents are marked as incomplete but
are not deleted.  In particular, the initialization is kept.
\nwenddocs{}\nwbegincode{58}\moddef{Event base: generic event: TBP}\plusendmoddef
  procedure :: reset => generic_event_reset
  procedure :: base_reset => generic_event_reset
\nwendcode{}\nwbegincode{59}\moddef{Event base: procedures}\plusendmoddef
  subroutine generic_event_reset (event)
    class(generic_event_t), intent(inout) :: event
    call event%discard_particle_set ()
    event%sqme_ref_known = .false.
    event%sqme_prc_known = .false.
    event%sqme_alt_known = .false.
    event%weight_ref_known = .false.
    event%weight_prc_known = .false.
    event%weight_alt_known = .false.
    event%excess_prc_known = .false.
  end subroutine generic_event_reset
  
\eatline
\nwendcode{}\nwbegindocs{60}Pacify particle set.
\nwenddocs{}\nwbegincode{61}\moddef{Event base: generic event: TBP}\plusendmoddef
  procedure :: pacify_particle_set => generic_event_pacify_particle_set
\nwendcode{}\nwbegincode{62}\moddef{Event base: procedures}\plusendmoddef
  subroutine generic_event_pacify_particle_set (event)
    class(generic_event_t), intent(inout) :: event
    if (event%has_valid_particle_set ())  call pacify (event%particle_set)
  end subroutine generic_event_pacify_particle_set

\eatline
\nwendcode{}\nwbegindocs{63}\nwdocspar
\subsection{Event normalization}
The parameters for event normalization.  For unweighted events,
{\tt{}NORM{\char95}UNIT} is intended as default, while for weighted events, it is
{\tt{}NORM{\char95}SIGMA}.

Note: the unit test for this is in {\tt{}eio{\char95}data{\char95}2} below.
\nwenddocs{}\nwbegincode{64}\moddef{Event base: parameters}\endmoddef
  integer, parameter, public :: NORM_UNDEFINED = 0
  integer, parameter, public :: NORM_UNIT = 1
  integer, parameter, public :: NORM_N_EVT = 2
  integer, parameter, public :: NORM_SIGMA = 3
  integer, parameter, public :: NORM_S_N = 4

\eatline
\nwendcode{}\nwbegindocs{65}These functions translate between the user representation and the
internal one.
\nwenddocs{}\nwbegincode{66}\moddef{Event base: public}\plusendmoddef
  public :: event_normalization_mode
  public :: event_normalization_string
\nwendcode{}\nwbegincode{67}\moddef{Event base: procedures}\plusendmoddef
  function event_normalization_mode (string, unweighted) result (mode)
    integer :: mode
    type(string_t), intent(in) :: string
    logical, intent(in) :: unweighted
    select case (lower_case (char (string)))
    case ("auto")
       if (unweighted) then
          mode = NORM_UNIT
       else
          mode = NORM_SIGMA
       end if
    case ("1")
       mode = NORM_UNIT
    case ("1/n")
       mode = NORM_N_EVT
    case ("sigma")
       mode = NORM_SIGMA
    case ("sigma/n")
       mode = NORM_S_N
    case default
       call msg_fatal ("Event normalization: unknown value '" &
            // char (string) // "'")
    end select
  end function event_normalization_mode
  
  function event_normalization_string (norm_mode) result (string)
    integer, intent(in) :: norm_mode
    type(string_t) :: string
    select case (norm_mode)
    case (NORM_UNDEFINED); string = "[undefined]"
    case (NORM_UNIT);      string = "'1'"
    case (NORM_N_EVT);     string = "'1/n'"
    case (NORM_SIGMA);     string = "'sigma'"
    case (NORM_S_N);       string = "'sigma/n'"
    case default;          string = "???"
    end select
  end function event_normalization_string
  
\eatline
\eatline
\nwendcode{}\nwbegindocs{68}We place this here as a generic helper, so we can update event
weights whenever we need, not just in connection with an event sample
data object.
\nwenddocs{}\nwbegincode{69}\moddef{Event base: public}\plusendmoddef
  public :: event_normalization_update
\nwendcode{}\nwbegincode{70}\moddef{Event base: procedures}\plusendmoddef
  subroutine event_normalization_update (weight, sigma, n, mode_new, mode_old)
    real(default), intent(inout) :: weight
    real(default), intent(in) :: sigma
    integer, intent(in) :: n
    integer, intent(in) :: mode_new, mode_old
    if (mode_new /= mode_old) then
       if (sigma > 0 .and. n > 0) then
          weight = weight / factor (mode_old) * factor (mode_new)
       else
          call msg_fatal ("Event normalization update: null sample")
       end if
    end if
  contains
    function factor (mode)
      real(default) :: factor
      integer, intent(in) :: mode
      select case (mode)
      case (NORM_UNIT);   factor = 1._default
      case (NORM_N_EVT);  factor = 1._default / n
      case (NORM_SIGMA);  factor = sigma
      case (NORM_S_N);    factor = sigma / n
      case default
         call msg_fatal ("Event normalization update: undefined mode")
      end select
    end function factor
  end subroutine event_normalization_update
  
\eatline
\nwendcode{}\nwbegindocs{71}\nwdocspar
%%%%%%%%%%%%%%%%%%%%%%%%%%%%%%%%%%%%%%%%%%%%%%%%%%%%%%%%%%%%%%%%%%%%%%%%
\section{Event Sample Data}
We define a simple and transparent container for (meta)data that are
associated with an event sample.
\nwenddocs{}\nwbegincode{72}\moddef{\code{}eio{\char95}data.f90\edoc{}}\endmoddef
\LA{}File header\RA{}

module eio_data
  
  use kinds
  use io_units
\LA{}Use strings\RA{}
  use unit_tests
  use diagnostics

  use event_base

\LA{}Standard module head\RA{}

\LA{}EIO data: public\RA{}

\LA{}EIO data: types\RA{}

contains
  
\LA{}EIO data: procedures\RA{}

\LA{}EIO data: tests\RA{}

end module eio_data
\eatline
\nwendcode{}\nwbegindocs{73}\nwdocspar
\subsection{Event Sample Data}
These are data that apply to an event sample as a whole.  They are
given in an easily portable form (no fancy structure) and are used for
initializing event formats.

There are two MD5 sums here.  {\tt{}md5sum{\char95}proc} depends only on the
definition of the contributing processes.  A sample with matching
checksum can be rescanned with modified model parameters, beam
structure etc, to recalculate observables.  {\tt{}md5sum{\char95}config} includes
all relevant data.  Rescanning a sample with matching checksum will
produce identical observables.  (A third checksum might be added which
depends on the event sample itself.  This is not needed, so far.)

If alternate weights are part of the event sample ({\tt{}n{\char95}alt} nonzero),
there is a configuration MD5 sum for each of them.
\nwenddocs{}\nwbegincode{74}\moddef{EIO data: public}\endmoddef
  public :: event_sample_data_t
\nwendcode{}\nwbegincode{75}\moddef{EIO data: types}\endmoddef
  type :: event_sample_data_t
     character(32) :: md5sum_prc = ""
     character(32) :: md5sum_cfg = ""
     logical :: unweighted = .true.
     logical :: negative_weights = .false.
     integer :: norm_mode = NORM_UNDEFINED
     integer :: n_beam = 0
     integer, dimension(2) :: pdg_beam = 0
     real(default), dimension(2) :: energy_beam = 0
     integer :: n_proc = 0
     integer :: n_evt = 0
     integer :: split_n_evt = 0
     integer :: split_index = 0
     real(default) :: total_cross_section = 0
     integer, dimension(:), allocatable :: proc_num_id
     integer :: n_alt = 0
     character(32), dimension(:), allocatable :: md5sum_alt
     real(default), dimension(:), allocatable :: cross_section
     real(default), dimension(:), allocatable :: error
   contains
   \LA{}EIO data: event sample data: TBP\RA{}
  end type event_sample_data_t
  
\eatline
\nwendcode{}\nwbegindocs{76}Initialize: allocate for the number of processes
\nwenddocs{}\nwbegincode{77}\moddef{EIO data: event sample data: TBP}\endmoddef
  procedure :: init => event_sample_data_init
\nwendcode{}\nwbegincode{78}\moddef{EIO data: procedures}\endmoddef
  subroutine event_sample_data_init (data, n_proc, n_alt)
    class(event_sample_data_t), intent(out) :: data
    integer, intent(in) :: n_proc
    integer, intent(in), optional :: n_alt
    data%n_proc = n_proc
    allocate (data%proc_num_id (n_proc), source = 0)
    allocate (data%cross_section (n_proc), source = 0._default)
    allocate (data%error (n_proc), source = 0._default)
    if (present (n_alt)) then
       data%n_alt = n_alt
       allocate (data%md5sum_alt (n_alt))
       data%md5sum_alt = ""
    end if
  end subroutine event_sample_data_init
  
\eatline
\nwendcode{}\nwbegindocs{79}Output.
\nwenddocs{}\nwbegincode{80}\moddef{EIO data: event sample data: TBP}\plusendmoddef
  procedure :: write => event_sample_data_write
\nwendcode{}\nwbegincode{81}\moddef{EIO data: procedures}\plusendmoddef
  subroutine event_sample_data_write (data, unit)
    class(event_sample_data_t), intent(in) :: data
    integer, intent(in), optional :: unit
    integer :: u, i
    u = given_output_unit (unit)
    write (u, "(1x,A)")  "Event sample properties:"
    write (u, "(3x,A,A,A)")  "MD5 sum (proc)   = '", data%md5sum_prc, "'"
    write (u, "(3x,A,A,A)")  "MD5 sum (config) = '", data%md5sum_cfg, "'"
    write (u, "(3x,A,L1)")  "unweighted       = ", data%unweighted
    write (u, "(3x,A,L1)")  "negative weights = ", data%negative_weights
    write (u, "(3x,A,A)")   "normalization    = ", &
         char (event_normalization_string (data%norm_mode))
    write (u, "(3x,A,I0)")  "number of beams  = ", data%n_beam
    write (u, "(5x,A,2(1x,I19))")  "PDG    = ", &
         data%pdg_beam(:data%n_beam)
    write (u, "(5x,A,2(1x,ES19.12))")  "Energy = ", &
         data%energy_beam(:data%n_beam)
    if (data%n_evt > 0) then
       write (u, "(3x,A,I0)")  "number of events = ", data%n_evt
    end if
    if (data%total_cross_section /= 0) then
       write (u, "(3x,A,ES19.12)")  "total cross sec. = ", &
            data%total_cross_section
    end if
    write (u, "(3x,A,I0)")  "num of processes = ", data%n_proc
    do i = 1, data%n_proc
       write (u, "(3x,A,I0)")  "Process #", data%proc_num_id (i)
       select case (data%n_beam)
       case (1)
          write (u, "(5x,A,ES19.12)")  "Width = ", data%cross_section(i)
       case (2)
          write (u, "(5x,A,ES19.12)")  "CSec  = ", data%cross_section(i)
       end select
       write (u, "(5x,A,ES19.12)")  "Error = ", data%error(i)
    end do
    if (data%n_alt > 0) then
       write (u, "(3x,A,I0)")  "num of alt wgt   = ", data%n_alt
       do i = 1, data%n_alt
          write (u, "(5x,A,A,A,1x,I0)")  "MD5 sum (cfg)  = '", &
               data%md5sum_alt(i), "'", i
       end do
    end if
  end subroutine event_sample_data_write
    
\eatline
\nwendcode{}\nwbegindocs{82}\nwdocspar
\subsection{Unit tests}
\nwenddocs{}\nwbegincode{83}\moddef{EIO data: public}\plusendmoddef
  public :: eio_data_test
\nwendcode{}\nwbegincode{84}\moddef{EIO data: tests}\endmoddef
  subroutine eio_data_test (u, results)
    integer, intent(in) :: u
    type(test_results_t), intent(inout) :: results
  \LA{}EIO data: execute tests\RA{}
  end subroutine eio_data_test
  
\eatline
\nwendcode{}\nwbegindocs{85}\nwdocspar
\subsubsection{Event Sample Data}
Print the contents of a sample data block.
\nwenddocs{}\nwbegincode{86}\moddef{EIO data: execute tests}\endmoddef
  call test (eio_data_1, "eio_data_1", &
       "event sample data", &
       u, results)
\nwendcode{}\nwbegincode{87}\moddef{EIO data: tests}\plusendmoddef
  subroutine eio_data_1 (u)
    integer, intent(in) :: u
    type(event_sample_data_t) :: data

    write (u, "(A)")  "* Test output: eio_data_1"
    write (u, "(A)")  "*   Purpose:  display event sample data"
    write (u, "(A)")

    write (u, "(A)")  "* Decay process, one component"
    write (u, "(A)")
 
    call data%init (1, 1)
    data%n_beam = 1
    data%pdg_beam(1) = 25
    data%energy_beam(1) = 125

    data%norm_mode = NORM_UNIT
    
    data%proc_num_id = [42]
    data%cross_section = [1.23e-4_default]
    data%error = 5e-6_default
    
    data%md5sum_prc = "abcdefghijklmnopabcdefghijklmnop"
    data%md5sum_cfg = "12345678901234561234567890123456"
    data%md5sum_alt(1) = "uuuuuuuuuuuuuuuuuuuuuuuuuuuuuuuu"
    
    call data%write (u)

    write (u, "(A)")
    write (u, "(A)")  "* Scattering process, two components"
    write (u, "(A)")
    
    call data%init (2)
    data%n_beam = 2
    data%pdg_beam = [2212, -2212]
    data%energy_beam = [8._default, 10._default]
    
    data%norm_mode = NORM_SIGMA
    
    data%proc_num_id = [12, 34]
    data%cross_section = [100._default, 88._default]
    data%error = [1._default, 0.1_default]
    
    call data%write (u)
    
    write (u, "(A)")
    write (u, "(A)")  "* Test output end: eio_data_1"
    
  end subroutine eio_data_1
  
\eatline
\nwendcode{}\nwbegindocs{88}\nwdocspar
\subsubsection{Event Normalization}
Check the functions for translating modes and updating weights.
\nwenddocs{}\nwbegincode{89}\moddef{EIO data: execute tests}\plusendmoddef
  call test (eio_data_2, "eio_data_2", &
       "event normalization", &
       u, results)
\nwendcode{}\nwbegincode{90}\moddef{EIO data: tests}\plusendmoddef
  subroutine eio_data_2 (u)
    integer, intent(in) :: u
    type(string_t) :: s
    logical :: unweighted
    real(default) :: w, w0, sigma
    integer :: n

    write (u, "(A)")  "* Test output: eio_data_2"
    write (u, "(A)")  "*   Purpose:  handle event normalization"
    write (u, "(A)")

    write (u, "(A)")  "* Normalization strings"
    write (u, "(A)")

    s = "auto"
    unweighted = .true.
    write (u, "(1x,A,1x,L1,1x,A)")  char (s), unweighted, &
         char (event_normalization_string &
         (event_normalization_mode (s, unweighted)))
    s = "AUTO"
    unweighted = .false.
    write (u, "(1x,A,1x,L1,1x,A)")  char (s), unweighted, &
         char (event_normalization_string &
         (event_normalization_mode (s, unweighted)))

    unweighted = .true.
    
    s = "1"
    write (u, "(2(1x,A))") char (s), char (event_normalization_string &
         (event_normalization_mode (s, unweighted)))
    s = "1/n"
    write (u, "(2(1x,A))") char (s), char (event_normalization_string &
         (event_normalization_mode (s, unweighted)))
    s = "Sigma"
    write (u, "(2(1x,A))") char (s), char (event_normalization_string &
         (event_normalization_mode (s, unweighted)))
    s = "sigma/N"
    write (u, "(2(1x,A))") char (s), char (event_normalization_string &
         (event_normalization_mode (s, unweighted)))

    write (u, "(A)")
    write (u, "(A)")  "* Normalization update"
    write (u, "(A)")
    
    sigma = 5
    n = 2

    w0 = 1

    w = w0
    call event_normalization_update (w, sigma, n, NORM_UNIT, NORM_UNIT)
    write (u, "(2(F6.3))")  w0, w
    w = w0
    call event_normalization_update (w, sigma, n, NORM_N_EVT, NORM_UNIT)
    write (u, "(2(F6.3))")  w0, w
    w = w0
    call event_normalization_update (w, sigma, n, NORM_SIGMA, NORM_UNIT)
    write (u, "(2(F6.3))")  w0, w
    w = w0
    call event_normalization_update (w, sigma, n, NORM_S_N, NORM_UNIT)
    write (u, "(2(F6.3))")  w0, w

    write (u, *)
    
    w0 = 0.5

    w = w0
    call event_normalization_update (w, sigma, n, NORM_UNIT, NORM_N_EVT)
    write (u, "(2(F6.3))")  w0, w
    w = w0
    call event_normalization_update (w, sigma, n, NORM_N_EVT, NORM_N_EVT)
    write (u, "(2(F6.3))")  w0, w
    w = w0
    call event_normalization_update (w, sigma, n, NORM_SIGMA, NORM_N_EVT)
    write (u, "(2(F6.3))")  w0, w
    w = w0
    call event_normalization_update (w, sigma, n, NORM_S_N, NORM_N_EVT)
    write (u, "(2(F6.3))")  w0, w
    
    write (u, *)
    
    w0 = 5.0

    w = w0
    call event_normalization_update (w, sigma, n, NORM_UNIT, NORM_SIGMA)
    write (u, "(2(F6.3))")  w0, w
    w = w0
    call event_normalization_update (w, sigma, n, NORM_N_EVT, NORM_SIGMA)
    write (u, "(2(F6.3))")  w0, w
    w = w0
    call event_normalization_update (w, sigma, n, NORM_SIGMA, NORM_SIGMA)
    write (u, "(2(F6.3))")  w0, w
    w = w0
    call event_normalization_update (w, sigma, n, NORM_S_N, NORM_SIGMA)
    write (u, "(2(F6.3))")  w0, w
    
    write (u, *)
    
    w0 = 2.5

    w = w0
    call event_normalization_update (w, sigma, n, NORM_UNIT, NORM_S_N)
    write (u, "(2(F6.3))")  w0, w
    w = w0
    call event_normalization_update (w, sigma, n, NORM_N_EVT, NORM_S_N)
    write (u, "(2(F6.3))")  w0, w
    w = w0
    call event_normalization_update (w, sigma, n, NORM_SIGMA, NORM_S_N)
    write (u, "(2(F6.3))")  w0, w
    w = w0
    call event_normalization_update (w, sigma, n, NORM_S_N, NORM_S_N)
    write (u, "(2(F6.3))")  w0, w
    
    write (u, "(A)")
    write (u, "(A)")  "* Test output end: eio_data_2"
    
  end subroutine eio_data_2
  
\eatline
\nwendcode{}\nwbegindocs{91}\nwdocspar
%%%%%%%%%%%%%%%%%%%%%%%%%%%%%%%%%%%%%%%%%%%%%%%%%%%%%%%%%%%%%%%%%%%%%%%%
\section{Abstract I/O Handler}

This module defines an abstract object for event I/O and the
associated methods.  

There are {\tt{}output} and {\tt{}input} methods which
write or read a single event from/to the I/O stream, respectively.
The I/O stream itself may be a file, a common block, or an externally
linked structure, depending on the concrete implementation.

A {\tt{}write} method prints the current content of the
implementation-dependent event record in human-readable form.

The {\tt{}init{\char95}in}/{\tt{}init{\char95}out} and {\tt{}final} prepare
and finalize the I/O stream, respectively.  There is also a
{\tt{}switch{\char95}inout} method which turns an input stream into an output
stream where events can be appended.

Optionally, output files can be split in chunks of well-defined size.  The
{\tt{}split{\char95}out} method takes care of this.
\nwenddocs{}\nwbegincode{92}\moddef{\code{}eio{\char95}base.f90\edoc{}}\endmoddef
\LA{}File header\RA{}

module eio_base
  
  use kinds
  use io_units
\LA{}Use strings\RA{}
  use unit_tests
  use diagnostics
  use lorentz
  use model_data
  use particles
  use event_base
  use eio_data

\LA{}Standard module head\RA{}

\LA{}EIO base: public\RA{}

\LA{}EIO base: types\RA{}

\LA{}EIO base: interfaces\RA{}

\LA{}EIO base: test types\RA{}

contains
  
\LA{}EIO base: procedures\RA{}

\LA{}EIO base: tests\RA{}

end module eio_base
\eatline
\nwendcode{}\nwbegindocs{93}\nwdocspar
\subsection{Type}
We can assume that most implementations will need the file extension as a
fixed string and, if they support file splitting, the current file index.

The fallback model is useful for implementations that are able to read
unknown files which may contain hadrons etc., not in the current
hard-interaction model.
\nwenddocs{}\nwbegincode{94}\moddef{EIO base: public}\endmoddef
  public :: eio_t
\nwendcode{}\nwbegincode{95}\moddef{EIO base: types}\endmoddef
  type, abstract :: eio_t
     type(string_t) :: sample
     type(string_t) :: extension
     type(string_t) :: filename
     logical :: split = .false.
     integer :: split_n_evt = 0
     integer :: split_index = 0
     class(model_data_t), pointer :: fallback_model => null ()
   contains
   \LA{}EIO base: eio: TBP\RA{}
  end type eio_t
  
\eatline
\nwendcode{}\nwbegindocs{96}Write to screen.  If possible, this should display the contents of the
current event, i.e., the last one that was written or read.
\nwenddocs{}\nwbegincode{97}\moddef{EIO base: eio: TBP}\endmoddef
  procedure (eio_write), deferred :: write
\nwendcode{}\nwbegincode{98}\moddef{EIO base: interfaces}\endmoddef
  abstract interface
     subroutine eio_write (object, unit)
       import
       class(eio_t), intent(in) :: object
       integer, intent(in), optional :: unit
     end subroutine eio_write
  end interface
  
\eatline
\nwendcode{}\nwbegindocs{99}Finalize.  This should write/read footer data and close input/output
channels.
\nwenddocs{}\nwbegincode{100}\moddef{EIO base: eio: TBP}\plusendmoddef
  procedure (eio_final), deferred :: final
\nwendcode{}\nwbegincode{101}\moddef{EIO base: interfaces}\plusendmoddef
  abstract interface
     subroutine eio_final (object)
       import
       class(eio_t), intent(inout) :: object
     end subroutine eio_final
  end interface
  
\eatline
\nwendcode{}\nwbegindocs{102}Determine splitting parameters from the event sample data.
\nwenddocs{}\nwbegincode{103}\moddef{EIO base: eio: TBP}\plusendmoddef
  procedure :: set_splitting => eio_set_splitting
\nwendcode{}\nwbegincode{104}\moddef{EIO base: procedures}\endmoddef
  subroutine eio_set_splitting (eio, data)
    class(eio_t), intent(inout) :: eio
    type(event_sample_data_t), intent(in) :: data
    eio%split = data%split_n_evt > 0
    if (eio%split) then
       eio%split_n_evt = data%split_n_evt
       eio%split_index = data%split_index
    end if
  end subroutine eio_set_splitting
    
\eatline
\nwendcode{}\nwbegindocs{105}Generate a filename, taking a possible split index into account.
\nwenddocs{}\nwbegincode{106}\moddef{EIO base: eio: TBP}\plusendmoddef
  procedure :: set_filename => eio_set_filename
\nwendcode{}\nwbegincode{107}\moddef{EIO base: procedures}\plusendmoddef
  subroutine eio_set_filename (eio)
    class(eio_t), intent(inout) :: eio
    character(32) :: buffer
    if (eio%split) then
       write (buffer, "(I0,'.')")  eio%split_index
       eio%filename = eio%sample // "." // trim (buffer) // eio%extension
    else
       eio%filename = eio%sample // "." // eio%extension
    end if    
  end subroutine eio_set_filename

\eatline
\nwendcode{}\nwbegindocs{108}Set the fallback model.
\nwenddocs{}\nwbegincode{109}\moddef{EIO base: eio: TBP}\plusendmoddef
  procedure :: set_fallback_model => eio_set_fallback_model
\nwendcode{}\nwbegincode{110}\moddef{EIO base: procedures}\plusendmoddef
  subroutine eio_set_fallback_model (eio, model)
    class(eio_t), intent(inout) :: eio
    class(model_data_t), intent(in), target :: model
    eio%fallback_model => model
  end subroutine eio_set_fallback_model
  
\eatline
\nwendcode{}\nwbegindocs{111}Initialize for output.  We provide process names.  This should
open an event file if appropriate and write header data.  Some methods
may require event sample data.
\nwenddocs{}\nwbegincode{112}\moddef{EIO base: eio: TBP}\plusendmoddef
  procedure (eio_init_out), deferred :: init_out
\nwendcode{}\nwbegincode{113}\moddef{EIO base: interfaces}\plusendmoddef
  abstract interface
     subroutine eio_init_out (eio, sample, data, success, extension)
       import
       class(eio_t), intent(inout) :: eio
       type(string_t), intent(in) :: sample
       type(event_sample_data_t), intent(in), optional :: data
       logical, intent(out), optional :: success
       type(string_t), intent(in), optional :: extension       
     end subroutine eio_init_out
  end interface

\eatline
\nwendcode{}\nwbegindocs{114}Initialize for input.  We provide process names.  This should open an event
file if appropriate and read header data.  The {\tt{}md5sum} can be used to check
the integrity of the configuration, it it provides a checksum to compare with.
In case the extension has changed the extension is also given as an argument.

The {\tt{}data} argument is {\tt{}intent(inout)}: we may read part of it and
keep other parts and/or check them against the data in the file.
\nwenddocs{}\nwbegincode{115}\moddef{EIO base: eio: TBP}\plusendmoddef
  procedure (eio_init_in), deferred :: init_in
\nwendcode{}\nwbegincode{116}\moddef{EIO base: interfaces}\plusendmoddef
  abstract interface
     subroutine eio_init_in (eio, sample, data, success, extension)
       import
       class(eio_t), intent(inout) :: eio
       type(string_t), intent(in) :: sample
       type(event_sample_data_t), intent(inout), optional :: data
       logical, intent(out), optional :: success
       type(string_t), intent(in), optional :: extension
     end subroutine eio_init_in
  end interface

\eatline
\nwendcode{}\nwbegindocs{117}Re-initialize for output.  This should change the status of any event file
from input to output and position it for appending new events.
\nwenddocs{}\nwbegincode{118}\moddef{EIO base: eio: TBP}\plusendmoddef
  procedure (eio_switch_inout), deferred :: switch_inout
\nwendcode{}\nwbegincode{119}\moddef{EIO base: interfaces}\plusendmoddef
  abstract interface
     subroutine eio_switch_inout (eio, success)
       import
       class(eio_t), intent(inout) :: eio
       logical, intent(out), optional :: success
     end subroutine eio_switch_inout
  end interface

\eatline
\nwendcode{}\nwbegindocs{120}This is similar: split the output, i.e., close the current file and open a
new one.  The default implementation does nothing.  For the feature to work,
an implementation must override this.
\nwenddocs{}\nwbegincode{121}\moddef{EIO base: eio: TBP}\plusendmoddef
  procedure :: split_out => eio_split_out
\nwendcode{}\nwbegincode{122}\moddef{EIO base: procedures}\plusendmoddef
  subroutine eio_split_out (eio)
    class(eio_t), intent(inout) :: eio
  end subroutine eio_split_out

\eatline
\nwendcode{}\nwbegindocs{123}\nwdocspar
\nwenddocs{}\nwbegindocs{124}Output an event.  All data can be taken from the {\tt{}event} record.
The index {\tt{}i{\char95}prc} identifies the process among the processes that
are contained in the current sample.  The {\tt{}reading} flag, if present,
indicates that the event was read from file, not generated.

The {\tt{}passed} flag tells us that this event has passed the selection
criteria.  Depending on the event format, we may choose to skip events
that have not passed.
\nwenddocs{}\nwbegincode{125}\moddef{EIO base: eio: TBP}\plusendmoddef
  procedure (eio_output), deferred :: output
\nwendcode{}\nwbegincode{126}\moddef{EIO base: interfaces}\plusendmoddef
  abstract interface
     subroutine eio_output (eio, event, i_prc, reading, passed, pacify)
       import
       class(eio_t), intent(inout) :: eio
       class(generic_event_t), intent(in), target :: event
       integer, intent(in) :: i_prc
       logical, intent(in), optional :: reading, passed, pacify
     end subroutine eio_output
  end interface
  
\eatline
\nwendcode{}\nwbegindocs{127}Input an event.  This should fill all event data that cannot be inferred
from the associated process.

The input is broken down into two parts.  First we read the {\tt{}i{\char95}prc}
index.  So we know which process to expect in the subsequent event.
If we have reached end of file, we also will know.
Then, we read the event itself.

The parameter {\tt{}iostat} is supposed to be set as the Fortran standard
requires, negative for EOF and positive for error.
\nwenddocs{}\nwbegincode{128}\moddef{EIO base: eio: TBP}\plusendmoddef
  procedure (eio_input_i_prc), deferred :: input_i_prc
  procedure (eio_input_event), deferred :: input_event
\nwendcode{}\nwbegincode{129}\moddef{EIO base: interfaces}\plusendmoddef
  abstract interface
     subroutine eio_input_i_prc (eio, i_prc, iostat)
       import
       class(eio_t), intent(inout) :: eio
       integer, intent(out) :: i_prc
       integer, intent(out) :: iostat
     end subroutine eio_input_i_prc
  end interface
  
  abstract interface
     subroutine eio_input_event (eio, event, iostat)
       import
       class(eio_t), intent(inout) :: eio
       class(generic_event_t), intent(inout), target :: event
       integer, intent(out) :: iostat
     end subroutine eio_input_event
  end interface
  
\eatline
\nwendcode{}\nwbegindocs{130}\nwdocspar
\subsection{Unit tests}
\nwenddocs{}\nwbegincode{131}\moddef{EIO base: public}\plusendmoddef
  public :: eio_base_test
\nwendcode{}\nwbegincode{132}\moddef{EIO base: tests}\endmoddef
  subroutine eio_base_test (u, results)
    integer, intent(in) :: u
    type(test_results_t), intent(inout) :: results
  \LA{}EIO base: execute tests\RA{}
  end subroutine eio_base_test
  
\eatline
\nwendcode{}\nwbegindocs{133}The caller has to provide procedures that prepare and cleanup the test
environment.  They depend on modules that are not available here.
\nwenddocs{}\nwbegincode{134}\moddef{EIO base: interfaces}\plusendmoddef
  abstract interface
     subroutine eio_prepare_event (event, unweighted, n_alt)
       import
       class(generic_event_t), intent(inout), pointer :: event
       logical, intent(in), optional :: unweighted
       integer, intent(in), optional :: n_alt
     end subroutine eio_prepare_event
  end interface
  
  abstract interface
     subroutine eio_cleanup_event (event)
       import
       class(generic_event_t), intent(inout), pointer :: event
     end subroutine eio_cleanup_event
  end interface
  
\nwendcode{}\nwbegindocs{135}We store pointers to the test-environment handlers as module variables.
This allows us to call them from the test routines themselves, which don't
allow for extra arguments.
\nwenddocs{}\nwbegincode{136}\moddef{EIO base: public}\plusendmoddef
  public :: eio_prepare_test, eio_cleanup_test
\nwendcode{}\nwbegincode{137}\moddef{EIO base: test types}\endmoddef
  procedure(eio_prepare_event), pointer :: eio_prepare_test => null ()
  procedure(eio_cleanup_event), pointer :: eio_cleanup_test => null ()
  
\eatline
\nwendcode{}\nwbegindocs{138}Similarly, for the fallback (hadron) model that some eio tests require:
\nwenddocs{}\nwbegincode{139}\moddef{EIO base: interfaces}\plusendmoddef
  abstract interface
     subroutine eio_prepare_model (model)
       import
       class(model_data_t), intent(inout), pointer :: model
     end subroutine eio_prepare_model
  end interface
  
  abstract interface
     subroutine eio_cleanup_model (model)
       import
       class(model_data_t), intent(inout), pointer :: model
     end subroutine eio_cleanup_model
  end interface
  
\nwendcode{}\nwbegincode{140}\moddef{EIO base: public}\plusendmoddef
  public :: eio_prepare_fallback_model, eio_cleanup_fallback_model
\nwendcode{}\nwbegincode{141}\moddef{EIO base: test types}\plusendmoddef
  procedure(eio_prepare_model), pointer :: eio_prepare_fallback_model => null ()
  procedure(eio_cleanup_model), pointer :: eio_cleanup_fallback_model => null ()
  
\eatline
\nwendcode{}\nwbegindocs{142}\nwdocspar
\subsubsection{Test type for event I/O}
The contents simulate the contents of an external file.  We have the
{\tt{}sample} string as the file name and the array of momenta
{\tt{}event{\char95}p} as the list of events.  The 
second index is the event index.  The {\tt{}event{\char95}i} component is the pointer
to the current event, {\tt{}event{\char95}n} is the total number of stored events.
\nwenddocs{}\nwbegincode{143}\moddef{EIO base: test types}\plusendmoddef
  type, extends (eio_t) :: eio_test_t
     integer :: event_n = 0
     integer :: event_i = 0
     integer :: i_prc = 0
     type(vector4_t), dimension(:,:), allocatable :: event_p
   contains
   \LA{}EIO base: eio test: TBP\RA{}
  end type eio_test_t

\eatline
\nwendcode{}\nwbegindocs{144}Write to screen.  Pretend that this is an actual event format.
\nwenddocs{}\nwbegincode{145}\moddef{EIO base: eio test: TBP}\endmoddef
  procedure :: write => eio_test_write
\nwendcode{}\nwbegincode{146}\moddef{EIO base: procedures}\plusendmoddef
  subroutine eio_test_write (object, unit)
    class(eio_test_t), intent(in) :: object
    integer, intent(in), optional :: unit
    integer :: u, i
    u = given_output_unit (unit)
    write (u, "(1x,A)")  "Test event stream"
    if (object%event_i /= 0) then
       write (u, "(1x,A,I0,A)")  "Event #", object%event_i, ":"
       do i = 1, size (object%event_p, 1)
          call vector4_write (object%event_p(i, object%event_i), u)
       end do
    end if
  end subroutine eio_test_write
    
\eatline
\nwendcode{}\nwbegindocs{147}Finalizer.  For the test case, we just reset the event count,
but keep the stored ``events''.  For the real implementations, the events
would be stored on an external medium, so we would delete the object
contents.
\nwenddocs{}\nwbegincode{148}\moddef{EIO base: eio test: TBP}\plusendmoddef
  procedure :: final => eio_test_final
\nwendcode{}\nwbegincode{149}\moddef{EIO base: procedures}\plusendmoddef
  subroutine eio_test_final (object)
    class(eio_test_t), intent(inout) :: object
    object%event_i = 0
  end subroutine eio_test_final
    
\eatline
\nwendcode{}\nwbegindocs{150}Initialization: We store the process IDs and the energy from the beam-data
object.  We also allocate the momenta (i.e., the simulated event record) for a
fixed maximum size of 10 events, 2 momenta each.  There is only a single
process.
\nwenddocs{}\nwbegincode{151}\moddef{EIO base: eio test: TBP}\plusendmoddef
  procedure :: init_out => eio_test_init_out
\nwendcode{}\nwbegincode{152}\moddef{EIO base: procedures}\plusendmoddef
  subroutine eio_test_init_out (eio, sample, data, success, extension)
    class(eio_test_t), intent(inout) :: eio
    type(string_t), intent(in) :: sample
    type(event_sample_data_t), intent(in), optional :: data
    logical, intent(out), optional :: success
    type(string_t), intent(in), optional :: extension
    eio%sample = sample
    eio%event_n = 0
    eio%event_i = 0
    allocate (eio%event_p (2, 10))
    if (present (success))  success = .true.
  end subroutine eio_test_init_out
  
\eatline
\nwendcode{}\nwbegindocs{153}Initialization for input.  Nothing to do for the test type.
\nwenddocs{}\nwbegincode{154}\moddef{EIO base: eio test: TBP}\plusendmoddef
  procedure :: init_in => eio_test_init_in
\nwendcode{}\nwbegincode{155}\moddef{EIO base: procedures}\plusendmoddef
  subroutine eio_test_init_in (eio, sample, data, success, extension)
    class(eio_test_t), intent(inout) :: eio
    type(string_t), intent(in) :: sample
    type(event_sample_data_t), intent(inout), optional :: data
    logical, intent(out), optional :: success
    type(string_t), intent(in), optional :: extension
    if (present (success))  success = .true.
  end subroutine eio_test_init_in
  
\eatline
\nwendcode{}\nwbegindocs{156}Switch from output to input.  Again, nothing to do for the test type.
\nwenddocs{}\nwbegincode{157}\moddef{EIO base: eio test: TBP}\plusendmoddef
  procedure :: switch_inout => eio_test_switch_inout
\nwendcode{}\nwbegincode{158}\moddef{EIO base: procedures}\plusendmoddef
  subroutine eio_test_switch_inout (eio, success)
    class(eio_test_t), intent(inout) :: eio
    logical, intent(out), optional :: success
    if (present (success))  success = .true.
  end subroutine eio_test_switch_inout
  
\eatline
\nwendcode{}\nwbegindocs{159}Output.  Increment the event counter and store the momenta of the current
event.
\nwenddocs{}\nwbegincode{160}\moddef{EIO base: eio test: TBP}\plusendmoddef
  procedure :: output => eio_test_output
\nwendcode{}\nwbegincode{161}\moddef{EIO base: procedures}\plusendmoddef
  subroutine eio_test_output (eio, event, i_prc, reading, passed, pacify)
    class(eio_test_t), intent(inout) :: eio
    class(generic_event_t), intent(in), target :: event
    logical, intent(in), optional :: reading, passed, pacify
    integer, intent(in) :: i_prc
    type(particle_set_t), pointer :: pset
    type(particle_t) :: prt
    eio%event_n = eio%event_n + 1
    eio%event_i = eio%event_n
    eio%i_prc = i_prc
    pset => event%get_particle_set_ptr ()
    prt = pset%get_particle (3)
    eio%event_p(1, eio%event_i) = prt%get_momentum ()
    prt = pset%get_particle (4)
    eio%event_p(2, eio%event_i) = prt%get_momentum ()
  end subroutine eio_test_output

\eatline
\nwendcode{}\nwbegindocs{162}Input.  Increment the event counter and retrieve the momenta of the current
event.  For the test case, we do not actually modify the current event.
\nwenddocs{}\nwbegincode{163}\moddef{EIO base: eio test: TBP}\plusendmoddef
  procedure :: input_i_prc => eio_test_input_i_prc
  procedure :: input_event => eio_test_input_event
\nwendcode{}\nwbegincode{164}\moddef{EIO base: procedures}\plusendmoddef
  subroutine eio_test_input_i_prc (eio, i_prc, iostat)
    class(eio_test_t), intent(inout) :: eio
    integer, intent(out) :: i_prc
    integer, intent(out) :: iostat
    i_prc = eio%i_prc
    iostat = 0
  end subroutine eio_test_input_i_prc

  subroutine eio_test_input_event (eio, event, iostat)
    class(eio_test_t), intent(inout) :: eio
    class(generic_event_t), intent(inout), target :: event
    integer, intent(out) :: iostat
    eio%event_i = eio%event_i + 1
    iostat = 0
  end subroutine eio_test_input_event

\eatline
\eatline
\nwendcode{}\nwbegindocs{165}\nwdocspar
\subsubsection{Test I/O methods}
\nwenddocs{}\nwbegincode{166}\moddef{EIO base: execute tests}\endmoddef
  call test (eio_base_1, "eio_base_1", &
       "read and write event contents", &
       u, results)
\nwendcode{}\nwbegincode{167}\moddef{EIO base: tests}\plusendmoddef
  subroutine eio_base_1 (u)
    integer, intent(in) :: u
    class(generic_event_t), pointer :: event
    class(eio_t), allocatable :: eio
    integer :: i_prc,  iostat
    type(string_t) :: sample

    write (u, "(A)")  "* Test output: eio_base_1"
    write (u, "(A)")  "*   Purpose: generate and read/write an event"
    write (u, "(A)")

    write (u, "(A)")  "* Initialize test process"

    call eio_prepare_test (event, unweighted = .false.)
    
    write (u, "(A)")
    write (u, "(A)")  "* Generate and write an event"
    write (u, "(A)")
 
    sample = "eio_test1"
 
    allocate (eio_test_t :: eio)
    
    call eio%init_out (sample)
    call event%generate (1, [0._default, 0._default])
    call eio%output (event, 42)
    call eio%write (u)
    call eio%final ()
    
    write (u, "(A)")
    write (u, "(A)")  "* Re-read the event"
    write (u, "(A)")
    
    call eio%init_in (sample)
    call eio%input_i_prc (i_prc, iostat)
    call eio%input_event (event, iostat)
    call eio%write (u)
    write (u, "(A)")
    write (u, "(1x,A,I0)")  "i = ", i_prc
    
    write (u, "(A)")
    write (u, "(A)")  "* Generate and append another event"
    write (u, "(A)")
    
    call eio%switch_inout ()
    call event%generate (1, [0._default, 0._default])
    call eio%output (event, 5)
    call eio%write (u)
    call eio%final ()
    
    write (u, "(A)")
    write (u, "(A)")  "* Re-read both events"
    write (u, "(A)")
    
    call eio%init_in (sample)
    call eio%input_i_prc (i_prc, iostat)
    call eio%input_event (event, iostat)
    call eio%input_i_prc (i_prc, iostat)
    call eio%input_event (event, iostat)
    call eio%write (u)
    write (u, "(A)")
    write (u, "(1x,A,I0)")  "i = ", i_prc
    
    write (u, "(A)")
    write (u, "(A)")  "* Cleanup"
 
    call eio%final ()
    deallocate (eio)
 
    call eio_cleanup_test (event)

    write (u, "(A)")
    write (u, "(A)")  "* Test output end: eio_base_1"
    
  end subroutine eio_base_1
  
\eatline
\nwendcode{}\nwbegindocs{168}\nwdocspar
%%%%%%%%%%%%%%%%%%%%%%%%%%%%%%%%%%%%%%%%%%%%%%%%%%%%%%%%%%%%%%%%%%%%%%%% 
\section{Event Generation Checkpoints}
This is an output-only format.  Its only use is to write screen
messages every $n$ events, to inform the user about progress.
\nwenddocs{}\nwbegincode{169}\moddef{\code{}eio{\char95}checkpoints.f90\edoc{}}\endmoddef
\LA{}File header\RA{}

module eio_checkpoints
  
  use kinds
  use io_units
\LA{}Use strings\RA{}
  use unit_tests
  use diagnostics

  use cputime
  use model_data
  use event_base
  use eio_data
  use eio_base

\LA{}Standard module head\RA{}

\LA{}EIO checkpoints: public\RA{}

\LA{}EIO checkpoints: parameters\RA{}

\LA{}EIO checkpoints: types\RA{}

contains
  
\LA{}EIO checkpoints: procedures\RA{}

\LA{}EIO checkpoints: tests\RA{}

end module eio_checkpoints
\eatline
\nwendcode{}\nwbegindocs{170}\nwdocspar
\subsection{Type}
\nwenddocs{}\nwbegincode{171}\moddef{EIO checkpoints: public}\endmoddef
  public :: eio_checkpoints_t
\nwendcode{}\nwbegincode{172}\moddef{EIO checkpoints: types}\endmoddef
  type, extends (eio_t) :: eio_checkpoints_t
     logical :: active = .false.
     logical :: running = .false.
     integer :: val = 0
     integer :: n_events = 0
     integer :: n_read = 0
     integer :: i_evt = 0
     logical :: blank = .false.
     type(timer_t) :: timer
   contains
   \LA{}EIO checkpoints: eio checkpoints: TBP\RA{}
  end type eio_checkpoints_t
  
\eatline
\nwendcode{}\nwbegindocs{173}\nwdocspar
\subsection{Specific Methods}
Set parameters that are specifically used with LHEF.
\nwenddocs{}\nwbegincode{174}\moddef{EIO checkpoints: eio checkpoints: TBP}\endmoddef
  procedure :: set_parameters => eio_checkpoints_set_parameters
\nwendcode{}\nwbegincode{175}\moddef{EIO checkpoints: procedures}\endmoddef
  subroutine eio_checkpoints_set_parameters (eio, checkpoint, blank)
    class(eio_checkpoints_t), intent(inout) :: eio
    integer, intent(in) :: checkpoint
    logical, intent(in), optional :: blank
    eio%val = checkpoint
    if (present (blank))  eio%blank = blank
  end subroutine eio_checkpoints_set_parameters
  
\eatline
\nwendcode{}\nwbegindocs{176}\nwdocspar
\subsection{Common Methods}
Output.  This is not the actual event format, but a readable account
of the current status.
\nwenddocs{}\nwbegincode{177}\moddef{EIO checkpoints: eio checkpoints: TBP}\plusendmoddef
  procedure :: write => eio_checkpoints_write
\nwendcode{}\nwbegincode{178}\moddef{EIO checkpoints: procedures}\plusendmoddef
  subroutine eio_checkpoints_write (object, unit)
    class(eio_checkpoints_t), intent(in) :: object
    integer, intent(in), optional :: unit
    integer :: u
    u = given_output_unit (unit)
    if (object%active) then
       write (u, "(1x,A)")  "Event-sample checkpoints:  active"
       write (u, "(3x,A,I0)")  "interval  = ", object%val
       write (u, "(3x,A,I0)")  "n_events  = ", object%n_events
       write (u, "(3x,A,I0)")  "n_read    = ", object%n_read
       write (u, "(3x,A,I0)")  "n_current = ", object%i_evt
       write (u, "(3x,A,L1)")  "blanking  = ", object%blank
       call object%timer%write (u)
    else
       write (u, "(1x,A)")  "Event-sample checkpoints:  off"
    end if
  end subroutine eio_checkpoints_write
  
\eatline
\nwendcode{}\nwbegindocs{179}Finalizer: trivial.
\nwenddocs{}\nwbegincode{180}\moddef{EIO checkpoints: eio checkpoints: TBP}\plusendmoddef
  procedure :: final => eio_checkpoints_final
\nwendcode{}\nwbegincode{181}\moddef{EIO checkpoints: procedures}\plusendmoddef
  subroutine eio_checkpoints_final (object)
    class(eio_checkpoints_t), intent(inout) :: object
    object%active = .false.
  end subroutine eio_checkpoints_final
  
\eatline
\nwendcode{}\nwbegindocs{182}Activate checkpointing for event generation or writing.
\nwenddocs{}\nwbegincode{183}\moddef{EIO checkpoints: eio checkpoints: TBP}\plusendmoddef
  procedure :: init_out => eio_checkpoints_init_out
\nwendcode{}\nwbegincode{184}\moddef{EIO checkpoints: procedures}\plusendmoddef
  subroutine eio_checkpoints_init_out (eio, sample, data, success, extension)
    class(eio_checkpoints_t), intent(inout) :: eio
    type(string_t), intent(in) :: sample
    type(string_t), intent(in), optional :: extension
    type(event_sample_data_t), intent(in), optional :: data
    logical, intent(out), optional :: success
    if (present (data)) then
       if (eio%val > 0) then
          eio%active = .true.
          eio%i_evt = 0
          eio%n_read = 0
          eio%n_events = data%n_evt
       end if
    end if
    if (present (success))  success = .true.
  end subroutine eio_checkpoints_init_out
    
\eatline
\nwendcode{}\nwbegindocs{185}No checkpointing for event reading.
\nwenddocs{}\nwbegincode{186}\moddef{EIO checkpoints: eio checkpoints: TBP}\plusendmoddef
  procedure :: init_in => eio_checkpoints_init_in
\nwendcode{}\nwbegincode{187}\moddef{EIO checkpoints: procedures}\plusendmoddef
  subroutine eio_checkpoints_init_in (eio, sample, data, success, extension)
    class(eio_checkpoints_t), intent(inout) :: eio
    type(string_t), intent(in) :: sample
    type(string_t), intent(in), optional :: extension
    type(event_sample_data_t), intent(inout), optional :: data
    logical, intent(out), optional :: success
    call msg_bug ("Event checkpoints: event input not supported")
    if (present (success))  success = .false.
  end subroutine eio_checkpoints_init_in
    
\eatline
\nwendcode{}\nwbegindocs{188}Switch from input to output: also not supported.
\nwenddocs{}\nwbegincode{189}\moddef{EIO checkpoints: eio checkpoints: TBP}\plusendmoddef
  procedure :: switch_inout => eio_checkpoints_switch_inout
\nwendcode{}\nwbegincode{190}\moddef{EIO checkpoints: procedures}\plusendmoddef
  subroutine eio_checkpoints_switch_inout (eio, success)
    class(eio_checkpoints_t), intent(inout) :: eio
    logical, intent(out), optional :: success
    call msg_bug ("Event checkpoints: in-out switch not supported")
    if (present (success))  success = .false.
  end subroutine eio_checkpoints_switch_inout
  
\eatline
\nwendcode{}\nwbegindocs{191}Checkpoints: display progress for the current event, if applicable.
\nwenddocs{}\nwbegincode{192}\moddef{EIO checkpoints: eio checkpoints: TBP}\plusendmoddef
  procedure :: output => eio_checkpoints_output
\nwendcode{}\nwbegincode{193}\moddef{EIO checkpoints: procedures}\plusendmoddef
  subroutine eio_checkpoints_output (eio, event, i_prc, reading, passed, pacify)
    class(eio_checkpoints_t), intent(inout) :: eio
    class(generic_event_t), intent(in), target :: event
    integer, intent(in) :: i_prc
    logical, intent(in), optional :: reading, passed, pacify
    logical :: rd
    rd = .false.;  if (present (reading))  rd = reading
    if (eio%active) then
       if (.not. eio%running)  call eio%startup ()
       if (eio%running) then
          eio%i_evt = eio%i_evt + 1
          if (rd) then
             eio%n_read = eio%n_read + 1
          else if (mod (eio%i_evt, eio%val) == 0) then
             call eio%message (eio%blank)
          end if
          if (eio%i_evt == eio%n_events)  call eio%shutdown ()
       end if
    end if
  end subroutine eio_checkpoints_output

\eatline
\nwendcode{}\nwbegindocs{194}When the first event is called, we have to initialize the screen output.
\nwenddocs{}\nwbegincode{195}\moddef{EIO checkpoints: eio checkpoints: TBP}\plusendmoddef
  procedure :: startup => eio_checkpoints_startup
\nwendcode{}\nwbegincode{196}\moddef{EIO checkpoints: procedures}\plusendmoddef
  subroutine eio_checkpoints_startup (eio)
    class(eio_checkpoints_t), intent(inout) :: eio
    if (eio%active .and. eio%i_evt < eio%n_events) then
       call msg_message ("")
       call msg_message (checkpoint_bar)
       call msg_message (checkpoint_head)
       call msg_message (checkpoint_bar)
       write (msg_buffer, checkpoint_fmt) 0., 0, eio%n_events - eio%i_evt, "???"
       call msg_message ()
       eio%running = .true.
       call eio%timer%start ()
    end if
  end subroutine eio_checkpoints_startup
  
\eatline
\nwendcode{}\nwbegindocs{197}This message is printed at every checkpoint.
\nwenddocs{}\nwbegincode{198}\moddef{EIO checkpoints: eio checkpoints: TBP}\plusendmoddef
  procedure :: message => eio_checkpoints_message
\nwendcode{}\nwbegincode{199}\moddef{EIO checkpoints: procedures}\plusendmoddef
  subroutine eio_checkpoints_message (eio, testflag)
    class(eio_checkpoints_t), intent(inout) :: eio
    logical, intent(in), optional :: testflag
    real :: t
    type(time_t) :: time_remaining
    type(string_t) :: time_string
    call eio%timer%stop ()
    t = eio%timer
    call eio%timer%restart ()
    time_remaining = &
         nint (t / (eio%i_evt - eio%n_read) * (eio%n_events - eio%i_evt))
    time_string = time_remaining%to_string_ms (blank = testflag)
    write (msg_buffer, checkpoint_fmt) &
         100 * (eio%i_evt - eio%n_read) / real (eio%n_events - eio%n_read), &
         eio%i_evt - eio%n_read, &
         eio%n_events - eio%i_evt, &
         char (time_string)
    call msg_message ()
  end subroutine eio_checkpoints_message

\eatline
\nwendcode{}\nwbegindocs{200}When the last event is called, wrap up.
\nwenddocs{}\nwbegincode{201}\moddef{EIO checkpoints: eio checkpoints: TBP}\plusendmoddef
  procedure :: shutdown => eio_checkpoints_shutdown
\nwendcode{}\nwbegincode{202}\moddef{EIO checkpoints: procedures}\plusendmoddef
  subroutine eio_checkpoints_shutdown (eio)
    class(eio_checkpoints_t), intent(inout) :: eio
    if (mod (eio%i_evt, eio%val) /= 0) then
       write (msg_buffer, checkpoint_fmt) &
            100., eio%i_evt - eio%n_read, 0, "0m:00s"
       call msg_message ()
    end if
    call msg_message (checkpoint_bar)
    call msg_message ("")
    eio%running = .false.
  end subroutine eio_checkpoints_shutdown

\eatline
\nwendcode{}\nwbegincode{203}\moddef{EIO checkpoints: eio checkpoints: TBP}\plusendmoddef
  procedure :: input_i_prc => eio_checkpoints_input_i_prc
  procedure :: input_event => eio_checkpoints_input_event
\nwendcode{}\nwbegincode{204}\moddef{EIO checkpoints: procedures}\plusendmoddef
  subroutine eio_checkpoints_input_i_prc (eio, i_prc, iostat)
    class(eio_checkpoints_t), intent(inout) :: eio
    integer, intent(out) :: i_prc
    integer, intent(out) :: iostat
    call msg_bug ("Event checkpoints: event input not supported")
    i_prc = 0
    iostat = 1
  end subroutine eio_checkpoints_input_i_prc

  subroutine eio_checkpoints_input_event (eio, event, iostat)
    class(eio_checkpoints_t), intent(inout) :: eio
    class(generic_event_t), intent(inout), target :: event
    integer, intent(out) :: iostat
    call msg_bug ("Event checkpoints: event input not supported")
    iostat = 1
  end subroutine eio_checkpoints_input_event

\eatline
\eatline
\nwendcode{}\nwbegindocs{205}\nwdocspar
\subsection{Message header}
\nwenddocs{}\nwbegincode{206}\moddef{EIO checkpoints: parameters}\endmoddef
  character(*), parameter :: &
     checkpoint_head = "| % complete | events generated | events remaining &
     &| time remaining"
  character(*), parameter :: &
     checkpoint_bar  = "|==================================================&
     &=================|"
  character(*), parameter :: &
     checkpoint_fmt  = "('   ',F5.1,T16,I9,T35,I9,T58,A)"
\eatline
\eatline
\eatline
\eatline
\nwendcode{}\nwbegindocs{207}\nwdocspar
\subsection{Unit tests}
\nwenddocs{}\nwbegincode{208}\moddef{EIO checkpoints: public}\plusendmoddef
  public :: eio_checkpoints_test
\nwendcode{}\nwbegincode{209}\moddef{EIO checkpoints: tests}\endmoddef
  subroutine eio_checkpoints_test (u, results)
    integer, intent(in) :: u
    type(test_results_t), intent(inout) :: results
  \LA{}EIO checkpoints: execute tests\RA{}
  end subroutine eio_checkpoints_test
  
\eatline
\nwendcode{}\nwbegindocs{210}\nwdocspar
\subsubsection{Test I/O methods}
We test the implementation of all I/O methods.
\nwenddocs{}\nwbegincode{211}\moddef{EIO checkpoints: execute tests}\endmoddef
  call test (eio_checkpoints_1, "eio_checkpoints_1", &
       "read and write event contents", &
       u, results)
\nwendcode{}\nwbegincode{212}\moddef{EIO checkpoints: tests}\plusendmoddef
  subroutine eio_checkpoints_1 (u)
    integer, intent(in) :: u
    class(generic_event_t), pointer :: event
    class(eio_t), allocatable :: eio
    type(event_sample_data_t) :: data
    type(string_t) :: sample
    integer :: i, n_events

    write (u, "(A)")  "* Test output: eio_checkpoints_1"
    write (u, "(A)")  "*   Purpose: generate a number of events &
         &with screen output"
    write (u, "(A)")

    write (u, "(A)")  "* Initialize test process"
 
    call eio_prepare_test (event)

    write (u, "(A)")
    write (u, "(A)")  "* Generate events"
    write (u, "(A)")
 
    sample = "eio_checkpoints_1"
 
    allocate (eio_checkpoints_t :: eio)

    n_events = 10
    call data%init (1, 0)
    data%n_evt = n_events

    select type (eio)
    type is (eio_checkpoints_t)
       call eio%set_parameters (checkpoint = 4)
    end select

    call eio%init_out (sample, data)

    do i = 1, n_events
       call event%generate (1, [0._default, 0._default])
       call eio%output (event, i_prc = 0)
    end do
    
    write (u, "(A)")  "* Checkpointing status"
    write (u, "(A)")

    call eio%write (u)
    call eio%final ()

    write (u, "(A)")
    write (u, "(A)")  "* Cleanup"
 
    call eio_cleanup_test (event)

    write (u, "(A)")
    write (u, "(A)")  "* Test output end: eio_checkpoints_1"
    
  end subroutine eio_checkpoints_1
  
\eatline
\nwendcode{}\nwbegindocs{213}\nwdocspar
%%%%%%%%%%%%%%%%%%%%%%%%%%%%%%%%%%%%%%%%%%%%%%%%%%%%%%%%%%%%%%%%%%%%%%%% 
\section{Event Weight Output}
This is an output-only format.  For each event, we print the indices
that identify process, process part (MCI group), and term.  As
numerical information we print the squared matrix element (trace) and
the event weight.
\nwenddocs{}\nwbegincode{214}\moddef{\code{}eio{\char95}weights.f90\edoc{}}\endmoddef
\LA{}File header\RA{}

module eio_weights
  
  use kinds
  use io_units
\LA{}Use strings\RA{}
  use unit_tests
  use diagnostics

  use lorentz
  use model_data
  use particles
  use event_base
  use eio_data
  use eio_base

\LA{}Standard module head\RA{}

\LA{}EIO weights: public\RA{}

\LA{}EIO weights: types\RA{}

contains
  
\LA{}EIO weights: procedures\RA{}

\LA{}EIO weights: tests\RA{}

end module eio_weights
\eatline
\nwendcode{}\nwbegindocs{215}\nwdocspar
\subsection{Type}
\nwenddocs{}\nwbegincode{216}\moddef{EIO weights: public}\endmoddef
  public :: eio_weights_t
\nwendcode{}\nwbegincode{217}\moddef{EIO weights: types}\endmoddef
  type, extends (eio_t) :: eio_weights_t
     logical :: writing = .false.
     integer :: unit = 0
     logical :: pacify = .false.
   contains
   \LA{}EIO weights: eio weights: TBP\RA{}
  end type eio_weights_t
  
\eatline
\nwendcode{}\nwbegindocs{218}\nwdocspar
\subsection{Specific Methods}
Set pacify flags.
\nwenddocs{}\nwbegincode{219}\moddef{EIO weights: eio weights: TBP}\endmoddef
  procedure :: set_parameters => eio_weights_set_parameters
\nwendcode{}\nwbegincode{220}\moddef{EIO weights: procedures}\endmoddef
  subroutine eio_weights_set_parameters (eio, pacify)
    class(eio_weights_t), intent(inout) :: eio
    logical, intent(in), optional :: pacify
    if (present (pacify))  eio%pacify = pacify    
    eio%extension = "weights.dat"
  end subroutine eio_weights_set_parameters
  
\eatline
\nwendcode{}\nwbegindocs{221}\nwdocspar
\subsection{Common Methods}
\nwenddocs{}\nwbegindocs{222}Output.  This is not the actual event format, but a readable account
of the current object status.
\nwenddocs{}\nwbegincode{223}\moddef{EIO weights: eio weights: TBP}\plusendmoddef
  procedure :: write => eio_weights_write
\nwendcode{}\nwbegincode{224}\moddef{EIO weights: procedures}\plusendmoddef
  subroutine eio_weights_write (object, unit)
    class(eio_weights_t), intent(in) :: object
    integer, intent(in), optional :: unit
    integer :: u
    u = given_output_unit (unit)
    write (u, "(1x,A)")  "Weight stream:"
    if (object%writing) then
       write (u, "(3x,A,A)")  "Writing to file   = ", char (object%filename)
       write (u, "(3x,A,L1)") "Reduced I/O prec. = ", object%pacify 
    else
       write (u, "(3x,A)")  "[closed]"
    end if
  end subroutine eio_weights_write
  
\eatline
\nwendcode{}\nwbegindocs{225}Finalizer: close any open file.
\nwenddocs{}\nwbegincode{226}\moddef{EIO weights: eio weights: TBP}\plusendmoddef
  procedure :: final => eio_weights_final
\nwendcode{}\nwbegincode{227}\moddef{EIO weights: procedures}\plusendmoddef
  subroutine eio_weights_final (object)
    class(eio_weights_t), intent(inout) :: object
    if (object%writing) then
       write (msg_buffer, "(A,A,A)")  "Events: closing weight stream file '", &
            char (object%filename), "'"
       call msg_message ()
       close (object%unit)
       object%writing = .false.
    end if
  end subroutine eio_weights_final
  
\eatline
\nwendcode{}\nwbegindocs{228}Initialize event writing.
\nwenddocs{}\nwbegincode{229}\moddef{EIO weights: eio weights: TBP}\plusendmoddef
  procedure :: init_out => eio_weights_init_out
\nwendcode{}\nwbegincode{230}\moddef{EIO weights: procedures}\plusendmoddef
  subroutine eio_weights_init_out (eio, sample, data, success, extension)
    class(eio_weights_t), intent(inout) :: eio
    type(string_t), intent(in) :: sample
    type(string_t), intent(in), optional :: extension
    type(event_sample_data_t), intent(in), optional :: data
    logical, intent(out), optional :: success
    if (present(extension)) then
       eio%extension = extension
    else 
       eio%extension = "weights.dat"
    end if
    eio%filename = sample // "." // eio%extension
    eio%unit = free_unit ()
    write (msg_buffer, "(A,A,A)")  "Events: writing to weight stream file '", &
         char (eio%filename), "'"
    call msg_message ()
    eio%writing = .true.
    open (eio%unit, file = char (eio%filename), &
         action = "write", status = "replace")
    if (present (success))  success = .true.
  end subroutine eio_weights_init_out
    
\eatline
\nwendcode{}\nwbegindocs{231}Initialize event reading.
\nwenddocs{}\nwbegincode{232}\moddef{EIO weights: eio weights: TBP}\plusendmoddef
  procedure :: init_in => eio_weights_init_in
\nwendcode{}\nwbegincode{233}\moddef{EIO weights: procedures}\plusendmoddef
  subroutine eio_weights_init_in (eio, sample, data, success, extension)
    class(eio_weights_t), intent(inout) :: eio
    type(string_t), intent(in) :: sample
    type(string_t), intent(in), optional :: extension
    type(event_sample_data_t), intent(inout), optional :: data
    logical, intent(out), optional :: success
    call msg_bug ("Weight stream: event input not supported")
    if (present (success))  success = .false.
  end subroutine eio_weights_init_in
    
\eatline
\nwendcode{}\nwbegindocs{234}Switch from input to output: reopen the file for reading.
\nwenddocs{}\nwbegincode{235}\moddef{EIO weights: eio weights: TBP}\plusendmoddef
  procedure :: switch_inout => eio_weights_switch_inout
\nwendcode{}\nwbegincode{236}\moddef{EIO weights: procedures}\plusendmoddef
  subroutine eio_weights_switch_inout (eio, success)
    class(eio_weights_t), intent(inout) :: eio
    logical, intent(out), optional :: success
    call msg_bug ("Weight stream: in-out switch not supported")
    if (present (success))  success = .false.
  end subroutine eio_weights_switch_inout
  
\eatline
\nwendcode{}\nwbegindocs{237}Output an event.  Write first the event indices, then weight and two
values of the squared matrix element: {\tt{}sqme{\char95}ref} is the value stored
in the event record, and {\tt{}sqme{\char95}prc} is the one stored in the process
instance.  (They can differ: when recalculating, the former is read
from file and the latter is the result of the new calculation.)

For the alternative entries, the {\tt{}sqme} value is always obtained by
a new calculation, and thus qualifies as {\tt{}sqme{\char95}prc}.
\nwenddocs{}\nwbegincode{238}\moddef{EIO weights: eio weights: TBP}\plusendmoddef
  procedure :: output => eio_weights_output
\nwendcode{}\nwbegincode{239}\moddef{EIO weights: procedures}\plusendmoddef
  subroutine eio_weights_output (eio, event, i_prc, reading, passed, pacify)
    class(eio_weights_t), intent(inout) :: eio
    class(generic_event_t), intent(in), target :: event
    integer, intent(in) :: i_prc
    logical, intent(in), optional :: reading, passed, pacify
    integer :: n_alt, i
    real(default) :: weight, sqme_ref, sqme_prc
    if (eio%writing) then
       weight = event%get_weight_prc ()
       sqme_ref = event%get_sqme_ref ()
       sqme_prc = event%get_sqme_prc ()
       n_alt = event%get_n_alt ()
1      format (I0,3(1x,ES17.10),3(1x,I0))
2      format (I0,3(1x,ES15.8),3(1x,I0))       
       if (eio%pacify) then
          write (eio%unit, 2)  0, weight, sqme_prc, sqme_ref, &
               i_prc 
       else
          write (eio%unit, 1)  0, weight, sqme_prc, sqme_ref, &
               i_prc 
       end if
       do i = 1, n_alt
          weight = event%get_weight_alt(i)
          sqme_prc = event%get_sqme_alt(i)
          if (eio%pacify) then
             write (eio%unit, 2)  i, weight, sqme_prc
          else
             write (eio%unit, 1)  i, weight, sqme_prc          
          end if
       end do
    else
       call eio%write ()
       call msg_fatal ("Weight stream file is not open for writing")
    end if
  end subroutine eio_weights_output

\eatline
\nwendcode{}\nwbegindocs{240}Input an event.
\nwenddocs{}\nwbegincode{241}\moddef{EIO weights: eio weights: TBP}\plusendmoddef
  procedure :: input_i_prc => eio_weights_input_i_prc
  procedure :: input_event => eio_weights_input_event
\nwendcode{}\nwbegincode{242}\moddef{EIO weights: procedures}\plusendmoddef
  subroutine eio_weights_input_i_prc (eio, i_prc, iostat)
    class(eio_weights_t), intent(inout) :: eio
    integer, intent(out) :: i_prc
    integer, intent(out) :: iostat
    call msg_bug ("Weight stream: event input not supported")
    i_prc = 0
    iostat = 1
  end subroutine eio_weights_input_i_prc

  subroutine eio_weights_input_event (eio, event, iostat)
    class(eio_weights_t), intent(inout) :: eio
    class(generic_event_t), intent(inout), target :: event
    integer, intent(out) :: iostat
    call msg_bug ("Weight stream: event input not supported")
    iostat = 1
  end subroutine eio_weights_input_event

\eatline
\eatline
\nwendcode{}\nwbegindocs{243}\nwdocspar
\subsection{Unit tests}
\nwenddocs{}\nwbegincode{244}\moddef{EIO weights: public}\plusendmoddef
  public :: eio_weights_test
\nwendcode{}\nwbegincode{245}\moddef{EIO weights: tests}\endmoddef
  subroutine eio_weights_test (u, results)
    integer, intent(in) :: u
    type(test_results_t), intent(inout) :: results
  \LA{}EIO weights: execute tests\RA{}
  end subroutine eio_weights_test
  
\eatline
\nwendcode{}\nwbegindocs{246}\nwdocspar
\subsubsection{Test I/O methods}
We test the implementation of all I/O methods.
\nwenddocs{}\nwbegincode{247}\moddef{EIO weights: execute tests}\endmoddef
  call test (eio_weights_1, "eio_weights_1", &
       "read and write event contents", &
       u, results)
\nwendcode{}\nwbegincode{248}\moddef{EIO weights: tests}\plusendmoddef
  subroutine eio_weights_1 (u)
    integer, intent(in) :: u
    class(generic_event_t), pointer :: event
    class(eio_t), allocatable :: eio
    type(string_t) :: sample
    integer :: u_file
    character(80) :: buffer

    write (u, "(A)")  "* Test output: eio_weights_1"
    write (u, "(A)")  "*   Purpose: generate an event and write weight to file"
    write (u, "(A)")

    write (u, "(A)")  "* Initialize test process"
 
    call eio_prepare_test (event, unweighted = .false.)
    
    write (u, "(A)")
    write (u, "(A)")  "* Generate and write an event"
    write (u, "(A)")
 
    sample = "eio_weights_1"
 
    allocate (eio_weights_t :: eio)
    
    call eio%init_out (sample)
    call event%generate (1, [0._default, 0._default])

    call eio%output (event, i_prc = 42)
    call eio%write (u)
    call eio%final ()

    write (u, "(A)")
    write (u, "(A)")  "* File contents: &
         &(weight, sqme(evt), sqme(prc), i_prc, i_mci, i_term)"
    write (u, "(A)")

    u_file = free_unit ()
    open (u_file, file = "eio_weights_1.weights.dat", &
         action = "read", status = "old")
    read (u_file, "(A)")  buffer
    write (u, "(A)") trim (buffer)
    close (u_file)
    
    write (u, "(A)")
    write (u, "(A)")  "* Cleanup"
 
    call eio_cleanup_test (event) 

    write (u, "(A)")
    write (u, "(A)")  "* Test output end: eio_weights_1"
  end subroutine eio_weights_1
  
\eatline
\nwendcode{}\nwbegindocs{249}\nwdocspar
\subsubsection{Test I/O methods}
We test the implementation of all I/O methods.
\nwenddocs{}\nwbegincode{250}\moddef{EIO weights: execute tests}\plusendmoddef
  call test (eio_weights_2, "eio_weights_2", &
       "multiple weights", &
       u, results)
\nwendcode{}\nwbegincode{251}\moddef{EIO weights: tests}\plusendmoddef
  subroutine eio_weights_2 (u)
    integer, intent(in) :: u
    class(generic_event_t), pointer :: event
    class(eio_t), allocatable :: eio
    type(string_t) :: sample
    integer :: u_file, i
    character(80) :: buffer

    write (u, "(A)")  "* Test output: eio_weights_2"
    write (u, "(A)")  "*   Purpose: generate an event and write weight to file"
    write (u, "(A)")

    write (u, "(A)")  "* Initialize test process"
 
    call eio_prepare_test (event, unweighted = .false., n_alt = 2)

    write (u, "(A)")
    write (u, "(A)")  "* Generate and write an event"
    write (u, "(A)")
 
    sample = "eio_weights_2"
 
    allocate (eio_weights_t :: eio)
    
    call eio%init_out (sample)
    select type (eio)
    type is (eio_weights_t)
       call eio%set_parameters (pacify = .true.)
    end select
    call event%generate (1, [0._default, 0._default])
    call event%set (sqme_alt = [2._default, 3._default])
    call event%set (weight_alt = &
         [2 * event%get_weight_prc (), 3 * event%get_weight_prc ()])

    call eio%output (event, i_prc = 42)
    call eio%write (u)
    call eio%final ()

    write (u, "(A)") 
    write (u, "(A)")  "* File contents: &
         &(weight, sqme(evt), sqme(prc), i_prc, i_mci, i_term)"
    write (u, "(A)")

    u_file = free_unit ()
    open (u_file, file = "eio_weights_2.weights.dat", &
         action = "read", status = "old")
    do i = 1, 3
       read (u_file, "(A)")  buffer
       write (u, "(A)") trim (buffer)
    end do
    close (u_file)
    
    write (u, "(A)")
    write (u, "(A)")  "* Cleanup"
 
    call eio_cleanup_test (event)

    write (u, "(A)")
    write (u, "(A)")  "* Test output end: eio_weights_2"
    
  end subroutine eio_weights_2
  
\eatline
\nwendcode{}\nwbegindocs{252}\nwdocspar
%%%%%%%%%%%%%%%%%%%%%%%%%%%%%%%%%%%%%%%%%%%%%%%%%%%%%%%%%%%%%%%%%%%%%%%%
\section{ASCII File Formats}
Here, we implement several ASCII file formats. It is possible to
switch between them using flags.
\nwenddocs{}\nwbegincode{253}\moddef{\code{}eio{\char95}ascii.f90\edoc{}}\endmoddef
\LA{}File header\RA{}

module eio_ascii
  
  use kinds
  use io_units
\LA{}Use strings\RA{}
  use unit_tests
  use diagnostics

  use lorentz
  use model_data
  use particles
  use event_base
  use hep_common
  use hep_events
  use eio_data
  use eio_base

\LA{}Standard module head\RA{}

\LA{}EIO ascii: public\RA{}

\LA{}EIO ascii: types\RA{}

contains
  
\LA{}EIO ascii: procedures\RA{}

\LA{}EIO ascii: tests\RA{}

end module eio_ascii
\eatline
\nwendcode{}\nwbegindocs{254}\nwdocspar
\subsection{Type}
\nwenddocs{}\nwbegincode{255}\moddef{EIO ascii: public}\endmoddef
  public :: eio_ascii_t
\nwendcode{}\nwbegincode{256}\moddef{EIO ascii: types}\endmoddef
  type, abstract, extends (eio_t) :: eio_ascii_t
     logical :: writing = .false.
     integer :: unit = 0
     logical :: keep_beams = .false.     
     logical :: keep_remnants = .true.
   contains
   \LA{}EIO ascii: eio ascii: TBP\RA{}
  end type eio_ascii_t
  
\eatline
\nwendcode{}\nwbegindocs{257}\nwdocspar
\nwenddocs{}\nwbegincode{258}\moddef{EIO ascii: public}\plusendmoddef
  public :: eio_ascii_ascii_t
\nwendcode{}\nwbegincode{259}\moddef{EIO ascii: types}\plusendmoddef
  type, extends (eio_ascii_t) :: eio_ascii_ascii_t
  end type eio_ascii_ascii_t
  
\eatline
\nwendcode{}\nwbegindocs{260}\nwdocspar
\nwenddocs{}\nwbegincode{261}\moddef{EIO ascii: public}\plusendmoddef
  public :: eio_ascii_athena_t
\nwendcode{}\nwbegincode{262}\moddef{EIO ascii: types}\plusendmoddef
  type, extends (eio_ascii_t) :: eio_ascii_athena_t
  end type eio_ascii_athena_t
  
\eatline
\nwendcode{}\nwbegindocs{263}The debug format has a few options that can be controlled by
Sindarin variables.
\nwenddocs{}\nwbegincode{264}\moddef{EIO ascii: public}\plusendmoddef
  public :: eio_ascii_debug_t
\nwendcode{}\nwbegincode{265}\moddef{EIO ascii: types}\plusendmoddef
  type, extends (eio_ascii_t) :: eio_ascii_debug_t
     logical :: show_process = .true.
     logical :: show_transforms = .true.
     logical :: show_decay = .true.
     logical :: verbose = .true.
  end type eio_ascii_debug_t
  
\eatline
\nwendcode{}\nwbegindocs{266}\nwdocspar
\nwenddocs{}\nwbegincode{267}\moddef{EIO ascii: public}\plusendmoddef
  public :: eio_ascii_hepevt_t
\nwendcode{}\nwbegincode{268}\moddef{EIO ascii: types}\plusendmoddef
  type, extends (eio_ascii_t) :: eio_ascii_hepevt_t
  end type eio_ascii_hepevt_t
  
\eatline
\nwendcode{}\nwbegindocs{269}\nwdocspar
\nwenddocs{}\nwbegincode{270}\moddef{EIO ascii: public}\plusendmoddef
  public :: eio_ascii_hepevt_verb_t
\nwendcode{}\nwbegincode{271}\moddef{EIO ascii: types}\plusendmoddef
  type, extends (eio_ascii_t) :: eio_ascii_hepevt_verb_t
  end type eio_ascii_hepevt_verb_t
  
\eatline
\nwendcode{}\nwbegindocs{272}\nwdocspar
\nwenddocs{}\nwbegincode{273}\moddef{EIO ascii: public}\plusendmoddef
  public :: eio_ascii_lha_t
\nwendcode{}\nwbegincode{274}\moddef{EIO ascii: types}\plusendmoddef
  type, extends (eio_ascii_t) :: eio_ascii_lha_t
  end type eio_ascii_lha_t
  
\eatline
\nwendcode{}\nwbegindocs{275}\nwdocspar
\nwenddocs{}\nwbegincode{276}\moddef{EIO ascii: public}\plusendmoddef
  public :: eio_ascii_lha_verb_t
\nwendcode{}\nwbegincode{277}\moddef{EIO ascii: types}\plusendmoddef
  type, extends (eio_ascii_t) :: eio_ascii_lha_verb_t
  end type eio_ascii_lha_verb_t
   
\eatline
\nwendcode{}\nwbegindocs{278}\nwdocspar
\nwenddocs{}\nwbegincode{279}\moddef{EIO ascii: public}\plusendmoddef
  public :: eio_ascii_long_t
\nwendcode{}\nwbegincode{280}\moddef{EIO ascii: types}\plusendmoddef
  type, extends (eio_ascii_t) :: eio_ascii_long_t
  end type eio_ascii_long_t
  
\eatline
\nwendcode{}\nwbegindocs{281}\nwdocspar
\nwenddocs{}\nwbegincode{282}\moddef{EIO ascii: public}\plusendmoddef
  public :: eio_ascii_mokka_t
\nwendcode{}\nwbegincode{283}\moddef{EIO ascii: types}\plusendmoddef
  type, extends (eio_ascii_t) :: eio_ascii_mokka_t
  end type eio_ascii_mokka_t
  
\eatline
\nwendcode{}\nwbegindocs{284}\nwdocspar
\nwenddocs{}\nwbegincode{285}\moddef{EIO ascii: public}\plusendmoddef
  public :: eio_ascii_short_t
\nwendcode{}\nwbegincode{286}\moddef{EIO ascii: types}\plusendmoddef
  type, extends (eio_ascii_t) :: eio_ascii_short_t
  end type eio_ascii_short_t
  
\eatline
\nwendcode{}\nwbegindocs{287}\nwdocspar
\subsection{Specific Methods}
Set parameters that are specifically used with ASCII file formats.  In
particular, this is the file extension.
\nwenddocs{}\nwbegincode{288}\moddef{EIO ascii: eio ascii: TBP}\endmoddef
  procedure :: set_parameters => eio_ascii_set_parameters
\nwendcode{}\nwbegincode{289}\moddef{EIO ascii: procedures}\endmoddef
  subroutine eio_ascii_set_parameters (eio, &
       keep_beams, keep_remnants, extension, &
       show_process, show_transforms, show_decay, verbose)
    class(eio_ascii_t), intent(inout) :: eio
    logical, intent(in), optional :: keep_beams
    logical, intent(in), optional :: keep_remnants
    type(string_t), intent(in), optional :: extension
    logical, intent(in), optional :: show_process, show_transforms, show_decay
    logical, intent(in), optional :: verbose
    if (present (keep_beams))  eio%keep_beams = keep_beams
    if (present (keep_remnants))  eio%keep_remnants = keep_remnants
    if (present (extension)) then
       eio%extension = extension
    else
       select type (eio)
       type is (eio_ascii_ascii_t)
          eio%extension = "evt"
       type is (eio_ascii_athena_t)
          eio%extension = "athena.evt"
       type is (eio_ascii_debug_t)
          eio%extension = "debug"
       type is (eio_ascii_hepevt_t)
          eio%extension = "hepevt"
       type is (eio_ascii_hepevt_verb_t)
          eio%extension = "hepevt.verb"          
       type is (eio_ascii_lha_t)
          eio%extension = "lha"
       type is (eio_ascii_lha_verb_t)
          eio%extension = "lha.verb"          
       type is (eio_ascii_long_t)
          eio%extension = "long.evt"
       type is (eio_ascii_mokka_t)
          eio%extension = "mokka.evt"
       type is (eio_ascii_short_t)
          eio%extension = "short.evt"
       end select
    end if
    select type (eio)
    type is (eio_ascii_debug_t)
       if (present (show_process))  eio%show_process = show_process
       if (present (show_transforms))  eio%show_transforms = show_transforms
       if (present (show_decay))  eio%show_decay = show_decay
       if (present (verbose))  eio%verbose = verbose
    end select
  end subroutine eio_ascii_set_parameters
  
\eatline
\nwendcode{}\nwbegindocs{290}\nwdocspar
\subsection{Common Methods}
Output.  This is not the actual event format, but a readable account
of the current object status.
\nwenddocs{}\nwbegincode{291}\moddef{EIO ascii: eio ascii: TBP}\plusendmoddef
  procedure :: write => eio_ascii_write
\nwendcode{}\nwbegincode{292}\moddef{EIO ascii: procedures}\plusendmoddef
  subroutine eio_ascii_write (object, unit)
    class(eio_ascii_t), intent(in) :: object
    integer, intent(in), optional :: unit
    integer :: u
    u = given_output_unit (unit)
    select type (object)
    type is (eio_ascii_ascii_t)
       write (u, "(1x,A)")  "ASCII event stream (default format):"
    type is (eio_ascii_athena_t)
       write (u, "(1x,A)")  "ASCII event stream (ATHENA format):"
    type is (eio_ascii_debug_t)
       write (u, "(1x,A)")  "ASCII event stream (Debugging format):"
    type is (eio_ascii_hepevt_t)
       write (u, "(1x,A)")  "ASCII event stream (HEPEVT format):"
    type is (eio_ascii_hepevt_verb_t)
       write (u, "(1x,A)")  "ASCII event stream (verbose HEPEVT format):"
    type is (eio_ascii_lha_t)
       write (u, "(1x,A)")  "ASCII event stream (LHA format):"
    type is (eio_ascii_lha_verb_t)
       write (u, "(1x,A)")  "ASCII event stream (verbose LHA format):"
    type is (eio_ascii_long_t)
       write (u, "(1x,A)")  "ASCII event stream (long format):"
    type is (eio_ascii_mokka_t)
       write (u, "(1x,A)")  "ASCII event stream (MOKKA format):"
    type is (eio_ascii_short_t)
       write (u, "(1x,A)")  "ASCII event stream (short format):"
    end select
    if (object%writing) then
       write (u, "(3x,A,A)")  "Writing to file   = ", char (object%filename)
    else
       write (u, "(3x,A)")  "[closed]"
    end if
    write (u, "(3x,A,L1)")    "Keep beams        = ", object%keep_beams
    write (u, "(3x,A,L1)")    "Keep remnants     = ", object%keep_remnants
    select type (object)
    type is (eio_ascii_debug_t)
       write (u, "(3x,A,L1)")    "Show process      = ", object%show_process
       write (u, "(3x,A,L1)")    "Show transforms   = ", object%show_transforms
       write (u, "(3x,A,L1)")    "Show decay tree   = ", object%show_decay
       write (u, "(3x,A,L1)")    "Verbose output    = ", object%verbose
    end select
  end subroutine eio_ascii_write
  
\eatline
\nwendcode{}\nwbegindocs{293}Finalizer: close any open file.
\nwenddocs{}\nwbegincode{294}\moddef{EIO ascii: eio ascii: TBP}\plusendmoddef
  procedure :: final => eio_ascii_final
\nwendcode{}\nwbegincode{295}\moddef{EIO ascii: procedures}\plusendmoddef
  subroutine eio_ascii_final (object)
    class(eio_ascii_t), intent(inout) :: object
    if (object%writing) then
       write (msg_buffer, "(A,A,A)")  "Events: closing ASCII file '", &
            char (object%filename), "'"
       call msg_message ()
       close (object%unit)
       object%writing = .false.
    end if
  end subroutine eio_ascii_final
  
\eatline
\nwendcode{}\nwbegindocs{296}Initialize event writing.

Check weight normalization.  This applies to all ASCII-type files that
use the HEPRUP common block.  We can't allow normalization conventions
that are not covered by the HEPRUP definition.
\nwenddocs{}\nwbegincode{297}\moddef{EIO ascii: eio ascii: TBP}\plusendmoddef
  procedure :: init_out => eio_ascii_init_out
\nwendcode{}\nwbegincode{298}\moddef{EIO ascii: procedures}\plusendmoddef
  subroutine eio_ascii_init_out (eio, sample, data, success, extension)
    class(eio_ascii_t), intent(inout) :: eio
    type(string_t), intent(in) :: sample
    type(string_t), intent(in), optional :: extension
    type(event_sample_data_t), intent(in), optional :: data
    logical, intent(out), optional :: success
    integer :: i
    if (.not. present (data)) &
         call msg_bug ("ASCII initialization: missing data")
    if (data%n_beam /= 2) &
         call msg_fatal ("ASCII: defined for scattering processes only")
    eio%sample = sample
    call eio%check_normalization (data)
    call eio%set_splitting (data)
    call eio%set_filename ()
    eio%unit = free_unit ()
    write (msg_buffer, "(A,A,A)")  "Events: writing to ASCII file '", &
         char (eio%filename), "'"
    call msg_message ()
    eio%writing = .true.
    open (eio%unit, file = char (eio%filename), &
         action = "write", status = "replace")
    select type (eio)
    type is (eio_ascii_lha_t)
       call heprup_init &
            (data%pdg_beam, &
            data%energy_beam, &
            n_processes = data%n_proc, &
            unweighted = data%unweighted, &
            negative_weights = data%negative_weights)           
       do i = 1, data%n_proc
          call heprup_set_process_parameters (i = i, &
               process_id = data%proc_num_id(i), &
               cross_section = data%cross_section(i), &
               error = data%error(i))
       end do
       call heprup_write_ascii (eio%unit)    
    type is (eio_ascii_lha_verb_t)
       call heprup_init &
            (data%pdg_beam, &
            data%energy_beam, &
            n_processes = data%n_proc, &
            unweighted = data%unweighted, &
            negative_weights = data%negative_weights)           
       do i = 1, data%n_proc
          call heprup_set_process_parameters (i = i, &
               process_id = data%proc_num_id(i), &
               cross_section = data%cross_section(i), &
               error = data%error(i))
       end do
       call heprup_write_verbose (eio%unit)        
    end select
    if (present (success))  success = .true.
  end subroutine eio_ascii_init_out
    
\eatline
\nwendcode{}\nwbegindocs{299}Some event properties do not go well with some output formats.  In
particular, many formats require unweighted events.
\nwenddocs{}\nwbegincode{300}\moddef{EIO ascii: eio ascii: TBP}\plusendmoddef
  procedure :: check_normalization => eio_ascii_check_normalization
\nwendcode{}\nwbegincode{301}\moddef{EIO ascii: procedures}\plusendmoddef
  subroutine eio_ascii_check_normalization (eio, data)
    class(eio_ascii_t), intent(in) :: eio
    type(event_sample_data_t), intent(in) :: data
    if (data%unweighted) then
    else
       select type (eio)
       type is (eio_ascii_athena_t);  call msg_fatal &
            ("Event output (Athena format): events must be unweighted.")
       type is (eio_ascii_hepevt_t);  call msg_fatal &
            ("Event output (HEPEVT format): events must be unweighted.")
       type is (eio_ascii_hepevt_verb_t);  call msg_fatal &
            ("Event output (HEPEVT format): events must be unweighted.")
       end select
       select case (data%norm_mode)
       case (NORM_SIGMA)
       case default
          select type (eio)
          type is (eio_ascii_lha_t)
             call msg_fatal &
                  ("Event output (LHA): normalization for weighted events &
                  &must be 'sigma'")
          type is (eio_ascii_lha_verb_t)
             call msg_fatal &
                  ("Event output (LHA): normalization for weighted events &
                  &must be 'sigma'")
          end select
       end select
    end if
  end subroutine eio_ascii_check_normalization
  
\eatline
\nwendcode{}\nwbegindocs{302}Initialize event reading.
\nwenddocs{}\nwbegincode{303}\moddef{EIO ascii: eio ascii: TBP}\plusendmoddef
  procedure :: init_in => eio_ascii_init_in
\nwendcode{}\nwbegincode{304}\moddef{EIO ascii: procedures}\plusendmoddef
  subroutine eio_ascii_init_in (eio, sample, data, success, extension)
    class(eio_ascii_t), intent(inout) :: eio
    type(string_t), intent(in) :: sample
    type(string_t), intent(in), optional :: extension
    type(event_sample_data_t), intent(inout), optional :: data
    logical, intent(out), optional :: success
    call msg_bug ("ASCII: event input not supported")
    if (present (success))  success = .false.
  end subroutine eio_ascii_init_in
    
\eatline
\nwendcode{}\nwbegindocs{305}Switch from input to output: reopen the file for reading.
\nwenddocs{}\nwbegincode{306}\moddef{EIO ascii: eio ascii: TBP}\plusendmoddef
  procedure :: switch_inout => eio_ascii_switch_inout
\nwendcode{}\nwbegincode{307}\moddef{EIO ascii: procedures}\plusendmoddef
  subroutine eio_ascii_switch_inout (eio, success)
    class(eio_ascii_t), intent(inout) :: eio
    logical, intent(out), optional :: success
    call msg_bug ("ASCII: in-out switch not supported")
    if (present (success))  success = .false.
  end subroutine eio_ascii_switch_inout
  
\eatline
\nwendcode{}\nwbegindocs{308}Split event file: increment the counter, close the current file, open a new
one.  If the file needs a header, repeat it for the new file.  (We assume that
the common block contents are still intact.)
\nwenddocs{}\nwbegincode{309}\moddef{EIO ascii: eio ascii: TBP}\plusendmoddef
  procedure :: split_out => eio_ascii_split_out
\nwendcode{}\nwbegincode{310}\moddef{EIO ascii: procedures}\plusendmoddef
  subroutine eio_ascii_split_out (eio)
    class(eio_ascii_t), intent(inout) :: eio
    if (eio%split) then
       eio%split_index = eio%split_index + 1
       call eio%set_filename ()
       write (msg_buffer, "(A,A,A)")  "Events: writing to ASCII file '", &
            char (eio%filename), "'"
       call msg_message ()
       close (eio%unit)
       open (eio%unit, file = char (eio%filename), &
            action = "write", status = "replace")
       select type (eio)
       type is (eio_ascii_lha_t)
          call heprup_write_ascii (eio%unit)    
       type is (eio_ascii_lha_verb_t)
          call heprup_write_verbose (eio%unit)        
       end select
    end if
  end subroutine eio_ascii_split_out
  
\eatline
\nwendcode{}\nwbegindocs{311}Output an event.  Write first the event indices, then weight and
squared matrix element, then the particle set.
\nwenddocs{}\nwbegincode{312}\moddef{EIO ascii: eio ascii: TBP}\plusendmoddef
  procedure :: output => eio_ascii_output
\nwendcode{}\nwbegincode{313}\moddef{EIO ascii: procedures}\plusendmoddef
  subroutine eio_ascii_output (eio, event, i_prc, reading, passed, pacify)
    class(eio_ascii_t), intent(inout) :: eio
    class(generic_event_t), intent(in), target :: event
    integer, intent(in) :: i_prc
    logical, intent(in), optional :: reading, passed, pacify
    if (present (passed)) then
       if (.not. passed)  return
    end if
    if (eio%writing) then
       select type (eio)
       type is (eio_ascii_lha_t)
          call hepeup_from_event (event, &
               process_index = i_prc, &
               keep_beams = eio%keep_beams, &
               keep_remnants = eio%keep_remnants)
          call hepeup_write_lha (eio%unit)
       type is (eio_ascii_lha_verb_t)
          call hepeup_from_event (event, &
               process_index = i_prc, &
               keep_beams = eio%keep_beams, &
               keep_remnants = eio%keep_remnants)
          call hepeup_write_verbose (eio%unit)          
       type is (eio_ascii_ascii_t)
          call event%write (eio%unit, &
               show_process = .false., &
               show_transforms = .false., &
               show_decay = .false., &
               verbose = .false., testflag = pacify)
       type is (eio_ascii_athena_t)
          call hepevt_from_event (event, &
               i_evt = event%get_index (), &          
               keep_beams = eio%keep_beams, &
               keep_remnants = eio%keep_remnants)
          call hepevt_write_athena (eio%unit)                    
       type is (eio_ascii_debug_t)
          call event%write (eio%unit, &
               show_process = eio%show_process, &
               show_transforms = eio%show_transforms, &
               show_decay = eio%show_decay, &
               verbose = eio%verbose)
       type is (eio_ascii_hepevt_t)
          call hepevt_from_event (event, &
               i_evt = event%get_index (), &                         
               keep_beams = eio%keep_beams, &
               keep_remnants = eio%keep_remnants)
          call hepevt_write_hepevt (eio%unit)                              
       type is (eio_ascii_hepevt_verb_t)
          call hepevt_from_event (event, &
               i_evt = event%get_index (), &                         
               keep_beams = eio%keep_beams, &
               keep_remnants = eio%keep_remnants)
          call hepevt_write_verbose (eio%unit)
       type is (eio_ascii_long_t)
          call hepevt_from_event (event, &
               i_evt = event%get_index (), & 
               keep_beams = eio%keep_beams, &
               keep_remnants = eio%keep_remnants)
          call hepevt_write_ascii (eio%unit, .true.)                           
       type is (eio_ascii_mokka_t)
          call hepevt_from_event (event, &
               i_evt = event%get_index (), &                         
               keep_beams = eio%keep_beams, &
               keep_remnants = eio%keep_remnants)
          call hepevt_write_mokka (eio%unit)                              
       type is (eio_ascii_short_t)
          call hepevt_from_event (event, &
               i_evt = event%get_index (), &  
               keep_beams = eio%keep_beams, &
               keep_remnants = eio%keep_remnants)
          call hepevt_write_ascii (eio%unit, .false.)                    
       end select       
    else
       call eio%write ()
       call msg_fatal ("ASCII file is not open for writing")
    end if
  end subroutine eio_ascii_output

\eatline
\nwendcode{}\nwbegindocs{314}Input an event.
\nwenddocs{}\nwbegincode{315}\moddef{EIO ascii: eio ascii: TBP}\plusendmoddef
  procedure :: input_i_prc => eio_ascii_input_i_prc
  procedure :: input_event => eio_ascii_input_event
\nwendcode{}\nwbegincode{316}\moddef{EIO ascii: procedures}\plusendmoddef
  subroutine eio_ascii_input_i_prc (eio, i_prc, iostat)
    class(eio_ascii_t), intent(inout) :: eio
    integer, intent(out) :: i_prc
    integer, intent(out) :: iostat
    call msg_bug ("ASCII: event input not supported")
    i_prc = 0
    iostat = 1
  end subroutine eio_ascii_input_i_prc

  subroutine eio_ascii_input_event (eio, event, iostat)
    class(eio_ascii_t), intent(inout) :: eio
    class(generic_event_t), intent(inout), target :: event
    integer, intent(out) :: iostat
    call msg_bug ("ASCII: event input not supported")
    iostat = 1
  end subroutine eio_ascii_input_event

\eatline
\eatline
\nwendcode{}\nwbegindocs{317}\nwdocspar
\subsection{Unit tests}
\nwenddocs{}\nwbegincode{318}\moddef{EIO ascii: public}\plusendmoddef
  public :: eio_ascii_test
\nwendcode{}\nwbegincode{319}\moddef{EIO ascii: tests}\endmoddef
  subroutine eio_ascii_test (u, results)
    integer, intent(in) :: u
    type(test_results_t), intent(inout) :: results
  \LA{}EIO ascii: execute tests\RA{}
  end subroutine eio_ascii_test
  
\eatline
\nwendcode{}\nwbegindocs{320}\nwdocspar
\subsubsection{Test I/O methods}
We test the implementation of all I/O methods, method {\tt{}ascii}:
\nwenddocs{}\nwbegincode{321}\moddef{EIO ascii: execute tests}\endmoddef
  call test (eio_ascii_1, "eio_ascii_1", &
       "read and write event contents, format [ascii]", &
       u, results)
\nwendcode{}\nwbegincode{322}\moddef{EIO ascii: tests}\plusendmoddef
  subroutine eio_ascii_1 (u)
    integer, intent(in) :: u
    class(generic_event_t), pointer :: event
    type(event_sample_data_t) :: data
    class(eio_t), allocatable :: eio
    type(string_t) :: sample
    integer :: u_file, iostat
    character(80) :: buffer

    write (u, "(A)")  "* Test output: eio_ascii_1"
    write (u, "(A)")  "*   Purpose: generate an event in ASCII ascii format"
    write (u, "(A)")  "*      and write weight to file"
    write (u, "(A)")

    write (u, "(A)")  "* Initialize test process"
 
    call eio_prepare_test (event, unweighted = .false.)

    call data%init (1)
    data%n_evt = 1
    data%n_beam = 2
    data%pdg_beam = 25
    data%energy_beam = 500
    data%proc_num_id = [42]
    data%cross_section(1) = 100
    data%error(1) = 1
    data%total_cross_section = sum (data%cross_section)

    write (u, "(A)")
    write (u, "(A)")  "* Generate and write an event"
    write (u, "(A)")
 
    sample = "eio_ascii_1"
 
    allocate (eio_ascii_ascii_t :: eio)
    
    select type (eio)
    class is (eio_ascii_t);  call eio%set_parameters ()
    end select
    call eio%init_out (sample, data)
    call event%generate (1, [0._default, 0._default])
    call event%evaluate_expressions ()

    call eio%output (event, i_prc = 1)
    call eio%write (u)
    call eio%final ()

    write (u, "(A)")
    write (u, "(A)")  "* File contents:"
    write (u, "(A)")

    u_file = free_unit ()
    open (u_file, file = char (sample // ".evt"), &
         action = "read", status = "old")
    do
       read (u_file, "(A)", iostat = iostat)  buffer
       if (buffer(1:21) == "  <generator_version>")  buffer = "[...]"
       if (iostat /= 0)  exit
       write (u, "(A)") trim (buffer)
    end do
    close (u_file)
    
    write (u, "(A)")
    write (u, "(A)")  "* Reset data"
    write (u, "(A)")
 
    deallocate (eio)
    allocate (eio_ascii_ascii_t :: eio)
    
    select type (eio)
    type is (eio_ascii_ascii_t)
       call eio%set_parameters (keep_beams = .true.)
    end select
    call eio%write (u)

    write (u, "(A)")
    write (u, "(A)")  "* Cleanup"
 
    call eio_cleanup_test (event)

    write (u, "(A)")
    write (u, "(A)")  "* Test output end: eio_ascii_1"
    
  end subroutine eio_ascii_1
  
\eatline
\nwendcode{}\nwbegindocs{323}\nwdocspar
We test the implementation of all I/O methods, method {\tt{}athena}:
\nwenddocs{}\nwbegincode{324}\moddef{EIO ascii: execute tests}\plusendmoddef
  call test (eio_ascii_2, "eio_ascii_2", &
       "read and write event contents, format [athena]", &
       u, results)
\nwendcode{}\nwbegincode{325}\moddef{EIO ascii: tests}\plusendmoddef
  subroutine eio_ascii_2 (u)
    integer, intent(in) :: u
    class(generic_event_t), pointer :: event
    type(event_sample_data_t) :: data
    class(eio_t), allocatable :: eio
    type(string_t) :: sample
    integer :: u_file, iostat
    character(80) :: buffer

    write (u, "(A)")  "* Test output: eio_ascii_2"
    write (u, "(A)")  "*   Purpose: generate an event in ASCII athena format"
    write (u, "(A)")  "*      and write weight to file"
    write (u, "(A)")

    write (u, "(A)")  "* Initialize test process"
 
    call eio_prepare_test (event, unweighted = .false.)
    
    call data%init (1)
    data%n_evt = 1
    data%n_beam = 2
    data%pdg_beam = 25
    data%energy_beam = 500
    data%proc_num_id = [42]
    data%cross_section(1) = 100
    data%error(1) = 1
    data%total_cross_section = sum (data%cross_section)

    write (u, "(A)")
    write (u, "(A)")  "* Generate and write an event"
    write (u, "(A)")
 
    sample = "eio_ascii_2"
 
    allocate (eio_ascii_athena_t :: eio)
    
    select type (eio)
    class is (eio_ascii_t);  call eio%set_parameters ()
    end select
    call eio%init_out (sample, data)
    call event%generate (1, [0._default, 0._default])
    call event%evaluate_expressions ()

    call eio%output (event, i_prc = 1)
    call eio%write (u)
    call eio%final ()

    write (u, "(A)")
    write (u, "(A)")  "* File contents:"
    write (u, "(A)")

    u_file = free_unit ()
    open (u_file, file = char(sample // ".athena.evt"), &
         action = "read", status = "old")
    do
       read (u_file, "(A)", iostat = iostat)  buffer
       if (buffer(1:21) == "  <generator_version>")  buffer = "[...]"
       if (iostat /= 0)  exit
       write (u, "(A)") trim (buffer)
    end do
    close (u_file)
    
    write (u, "(A)")
    write (u, "(A)")  "* Reset data"
    write (u, "(A)")
 
    deallocate (eio)
    allocate (eio_ascii_athena_t :: eio)
    
    select type (eio)
    type is (eio_ascii_athena_t)
       call eio%set_parameters (keep_beams = .true.)
    end select
    call eio%write (u)

    write (u, "(A)")
    write (u, "(A)")  "* Cleanup"
 
    call eio_cleanup_test (event)

    write (u, "(A)")
    write (u, "(A)")  "* Test output end: eio_ascii_2"
    
  end subroutine eio_ascii_2
  
\eatline
\nwendcode{}\nwbegindocs{326}\nwdocspar
We test the implementation of all I/O methods, method {\tt{}debug}:
\nwenddocs{}\nwbegincode{327}\moddef{EIO ascii: execute tests}\plusendmoddef
  call test (eio_ascii_3, "eio_ascii_3", &
       "read and write event contents, format [debug]", &
       u, results)
\nwendcode{}\nwbegincode{328}\moddef{EIO ascii: tests}\plusendmoddef
  subroutine eio_ascii_3 (u)
    integer, intent(in) :: u
    class(generic_event_t), pointer :: event
    type(event_sample_data_t) :: data
    class(eio_t), allocatable :: eio
    type(string_t) :: sample
    integer :: u_file, iostat
    character(80) :: buffer

    write (u, "(A)")  "* Test output: eio_ascii_3"
    write (u, "(A)")  "*   Purpose: generate an event in ASCII debug format"
    write (u, "(A)")  "*      and write weight to file"
    write (u, "(A)")

    write (u, "(A)")  "* Initialize test process"
 
    call eio_prepare_test (event, unweighted = .false.)
    
    call data%init (1)
    data%n_evt = 1
    data%n_beam = 2
    data%pdg_beam = 25
    data%energy_beam = 500
    data%proc_num_id = [42]
    data%cross_section(1) = 100
    data%error(1) = 1
    data%total_cross_section = sum (data%cross_section)

    write (u, "(A)")
    write (u, "(A)")  "* Generate and write an event"
    write (u, "(A)")
 
    sample = "eio_ascii_3"
 
    allocate (eio_ascii_debug_t :: eio)
    
    select type (eio)
    class is (eio_ascii_t);  call eio%set_parameters ()
    end select
    call eio%init_out (sample, data)
    call event%generate (1, [0._default, 0._default])
    call event%evaluate_expressions ()

    call eio%output (event, i_prc = 1)
    call eio%write (u)
    call eio%final ()

    write (u, "(A)")
    write (u, "(A)")  "* File contents:"
    write (u, "(A)")

    u_file = free_unit ()
    open (u_file, file = char (sample // ".debug"), &
         action = "read", status = "old")
    do
       read (u_file, "(A)", iostat = iostat)  buffer
       if (buffer(1:21) == "  <generator_version>")  buffer = "[...]"
       if (iostat /= 0)  exit
       write (u, "(A)") trim (buffer)
    end do
    close (u_file)
    
    write (u, "(A)")
    write (u, "(A)")  "* Reset data"
    write (u, "(A)")
 
    deallocate (eio)
    allocate (eio_ascii_debug_t :: eio)
    
    select type (eio)
    type is (eio_ascii_debug_t)
       call eio%set_parameters (keep_beams = .true.)
    end select
    call eio%write (u)

    write (u, "(A)")
    write (u, "(A)")  "* Cleanup"
 
    call eio_cleanup_test (event)

    write (u, "(A)")
    write (u, "(A)")  "* Test output end: eio_ascii_3"
    
  end subroutine eio_ascii_3
  
\eatline
\nwendcode{}\nwbegindocs{329}\nwdocspar
We test the implementation of all I/O methods, method {\tt{}hepevt}:
\nwenddocs{}\nwbegincode{330}\moddef{EIO ascii: execute tests}\plusendmoddef
  call test (eio_ascii_4, "eio_ascii_4", &
       "read and write event contents, format [hepevt]", &
       u, results)
\nwendcode{}\nwbegincode{331}\moddef{EIO ascii: tests}\plusendmoddef
  subroutine eio_ascii_4 (u)
    integer, intent(in) :: u
    class(generic_event_t), pointer :: event
    type(event_sample_data_t) :: data
    class(eio_t), allocatable :: eio
    type(string_t) :: sample
    integer :: u_file, iostat
    character(80) :: buffer

    write (u, "(A)")  "* Test output: eio_ascii_4"
    write (u, "(A)")  "*   Purpose: generate an event in ASCII hepevt format"
    write (u, "(A)")  "*      and write weight to file"
    write (u, "(A)")

    write (u, "(A)")  "* Initialize test process"
 
    call eio_prepare_test (event, unweighted = .false.)
    
    call data%init (1)
    data%n_evt = 1
    data%n_beam = 2
    data%pdg_beam = 25
    data%energy_beam = 500
    data%proc_num_id = [42]
    data%cross_section(1) = 100
    data%error(1) = 1
    data%total_cross_section = sum (data%cross_section)

    write (u, "(A)")
    write (u, "(A)")  "* Generate and write an event"
    write (u, "(A)")
 
    sample = "eio_ascii_4"
 
    allocate (eio_ascii_hepevt_t :: eio)
    
    select type (eio)
    class is (eio_ascii_t);  call eio%set_parameters ()
    end select
    call eio%init_out (sample, data)
    call event%generate (1, [0._default, 0._default])
    call event%evaluate_expressions ()

    call eio%output (event, i_prc = 1)
    call eio%write (u)
    call eio%final ()

    write (u, "(A)")
    write (u, "(A)")  "* File contents:"
    write (u, "(A)")

    u_file = free_unit ()
    open (u_file, file = char (sample // ".hepevt"), &
         action = "read", status = "old")
    do
       read (u_file, "(A)", iostat = iostat)  buffer
       if (buffer(1:21) == "  <generator_version>")  buffer = "[...]"
       if (iostat /= 0)  exit
       write (u, "(A)") trim (buffer)
    end do
    close (u_file)
    
    write (u, "(A)")
    write (u, "(A)")  "* Reset data"
    write (u, "(A)")
 
    deallocate (eio)
    allocate (eio_ascii_hepevt_t :: eio)
    
    select type (eio)
    type is (eio_ascii_hepevt_t)
       call eio%set_parameters (keep_beams = .true.)
    end select
    call eio%write (u)

    write (u, "(A)")
    write (u, "(A)")  "* Cleanup"
 
    call eio_cleanup_test (event)

    write (u, "(A)")
    write (u, "(A)")  "* Test output end: eio_ascii_4"
    
  end subroutine eio_ascii_4
  
\eatline
\nwendcode{}\nwbegindocs{332}\nwdocspar
We test the implementation of all I/O methods, method {\tt{}lha} (old LHA):
\nwenddocs{}\nwbegincode{333}\moddef{EIO ascii: execute tests}\plusendmoddef
  call test (eio_ascii_5, "eio_ascii_5", &
       "read and write event contents, format [lha]", &
       u, results)
\nwendcode{}\nwbegincode{334}\moddef{EIO ascii: tests}\plusendmoddef
  subroutine eio_ascii_5 (u)
    integer, intent(in) :: u
    class(generic_event_t), pointer :: event
    type(event_sample_data_t) :: data
    class(eio_t), allocatable :: eio
    type(string_t) :: sample
    integer :: u_file, iostat
    character(80) :: buffer

    write (u, "(A)")  "* Test output: eio_ascii_5"
    write (u, "(A)")  "*   Purpose: generate an event in ASCII LHA format"
    write (u, "(A)")  "*      and write weight to file"
    write (u, "(A)")

    write (u, "(A)")  "* Initialize test process"
 
    call eio_prepare_test (event, unweighted = .false.)
    
    call data%init (1)
    data%n_evt = 1
    data%n_beam = 2
    data%pdg_beam = 25
    data%energy_beam = 500
    data%proc_num_id = [42]
    data%cross_section(1) = 100
    data%error(1) = 1
    data%total_cross_section = sum (data%cross_section)

    write (u, "(A)")
    write (u, "(A)")  "* Generate and write an event"
    write (u, "(A)")
 
    sample = "eio_ascii_5"
 
    allocate (eio_ascii_lha_t :: eio)
    
    select type (eio)
    class is (eio_ascii_t);  call eio%set_parameters ()
    end select
    call eio%init_out (sample, data)
    call event%generate (1, [0._default, 0._default])
    call event%evaluate_expressions ()

    call eio%output (event, i_prc = 1)
    call eio%write (u)
    call eio%final ()

    write (u, "(A)")
    write (u, "(A)")  "* File contents:"
    write (u, "(A)")

    u_file = free_unit ()
    open (u_file, file = char (sample // ".lha"), &
         action = "read", status = "old")
    do
       read (u_file, "(A)", iostat = iostat)  buffer
       if (buffer(1:21) == "  <generator_version>")  buffer = "[...]"
       if (iostat /= 0)  exit
       write (u, "(A)") trim (buffer)
    end do
    close (u_file)
    
    write (u, "(A)")
    write (u, "(A)")  "* Reset data"
    write (u, "(A)")
 
    deallocate (eio)
    allocate (eio_ascii_lha_t :: eio)
    
    select type (eio)
    type is (eio_ascii_lha_t)
       call eio%set_parameters (keep_beams = .true.)
    end select
    call eio%write (u)

    write (u, "(A)")
    write (u, "(A)")  "* Cleanup"
 
    call eio_cleanup_test (event)

    write (u, "(A)")
    write (u, "(A)")  "* Test output end: eio_ascii_5"
    
  end subroutine eio_ascii_5
  
\eatline
\nwendcode{}\nwbegindocs{335}\nwdocspar
We test the implementation of all I/O methods, method {\tt{}long}:
\nwenddocs{}\nwbegincode{336}\moddef{EIO ascii: execute tests}\plusendmoddef
  call test (eio_ascii_6, "eio_ascii_6", &
       "read and write event contents, format [long]", &
       u, results)
\nwendcode{}\nwbegincode{337}\moddef{EIO ascii: tests}\plusendmoddef
  subroutine eio_ascii_6 (u)
    integer, intent(in) :: u
    class(generic_event_t), pointer :: event
    type(event_sample_data_t) :: data
    class(eio_t), allocatable :: eio
    type(string_t) :: sample
    integer :: u_file, iostat
    character(80) :: buffer

    write (u, "(A)")  "* Test output: eio_ascii_6"
    write (u, "(A)")  "*   Purpose: generate an event in ASCII long format"
    write (u, "(A)")  "*      and write weight to file"
    write (u, "(A)")

    write (u, "(A)")  "* Initialize test process"
 
    call eio_prepare_test (event, unweighted = .false.)
    
    call data%init (1)
    data%n_evt = 1
    data%n_beam = 2
    data%pdg_beam = 25
    data%energy_beam = 500
    data%proc_num_id = [42]
    data%cross_section(1) = 100
    data%error(1) = 1
    data%total_cross_section = sum (data%cross_section)

    write (u, "(A)")
    write (u, "(A)")  "* Generate and write an event"
    write (u, "(A)")
 
    sample = "eio_ascii_6"
 
    allocate (eio_ascii_long_t :: eio)
    
    select type (eio)
    class is (eio_ascii_t);  call eio%set_parameters ()
    end select
    call eio%init_out (sample, data)
    call event%generate (1, [0._default, 0._default])
    call event%evaluate_expressions ()

    call eio%output (event, i_prc = 1)
    call eio%write (u)
    call eio%final ()

    write (u, "(A)")
    write (u, "(A)")  "* File contents:"
    write (u, "(A)")

    u_file = free_unit ()
    open (u_file, file = char (sample // ".long.evt"), &
         action = "read", status = "old")
    do
       read (u_file, "(A)", iostat = iostat)  buffer
       if (buffer(1:21) == "  <generator_version>")  buffer = "[...]"
       if (iostat /= 0)  exit
       write (u, "(A)") trim (buffer)
    end do
    close (u_file)
    
    write (u, "(A)")
    write (u, "(A)")  "* Reset data"
    write (u, "(A)")
 
    deallocate (eio)
    allocate (eio_ascii_long_t :: eio)
    
    select type (eio)
    type is (eio_ascii_long_t)
       call eio%set_parameters (keep_beams = .true.)
    end select
    call eio%write (u)

    write (u, "(A)")
    write (u, "(A)")  "* Cleanup"
 
    call eio_cleanup_test (event)

    write (u, "(A)")
    write (u, "(A)")  "* Test output end: eio_ascii_6"
    
  end subroutine eio_ascii_6
  
\eatline
\nwendcode{}\nwbegindocs{338}\nwdocspar
We test the implementation of all I/O methods, method {\tt{}mokka}:
\nwenddocs{}\nwbegincode{339}\moddef{EIO ascii: execute tests}\plusendmoddef
  call test (eio_ascii_7, "eio_ascii_7", &
       "read and write event contents, format [mokka]", &
       u, results)
\nwendcode{}\nwbegincode{340}\moddef{EIO ascii: tests}\plusendmoddef
  subroutine eio_ascii_7 (u)
    integer, intent(in) :: u
    type(model_data_t), target :: model
    class(generic_event_t), pointer :: event
    type(event_sample_data_t) :: data
    class(eio_t), allocatable :: eio
    type(string_t) :: sample
    integer :: u_file, iostat
    character(80) :: buffer

    write (u, "(A)")  "* Test output: eio_ascii_7"
    write (u, "(A)")  "*   Purpose: generate an event in ASCII mokka format"
    write (u, "(A)")  "*      and write weight to file"
    write (u, "(A)")

    write (u, "(A)")  "* Initialize test process"
 
    call eio_prepare_test (event, unweighted = .false.)
    
    call data%init (1)
    data%n_evt = 1
    data%n_beam = 2
    data%pdg_beam = 25
    data%energy_beam = 500
    data%proc_num_id = [42]
    data%cross_section(1) = 100
    data%error(1) = 1
    data%total_cross_section = sum (data%cross_section)

    write (u, "(A)")
    write (u, "(A)")  "* Generate and write an event"
    write (u, "(A)")
 
    sample = "eio_ascii_7"
 
    allocate (eio_ascii_mokka_t :: eio)
    
    select type (eio)
    class is (eio_ascii_t);  call eio%set_parameters ()
    end select
    call eio%init_out (sample, data)
    call event%generate (1, [0._default, 0._default])
    call event%evaluate_expressions ()

    call eio%output (event, i_prc = 1)
    call eio%write (u)
    call eio%final ()

    write (u, "(A)")
    write (u, "(A)")  "* File contents:"
    write (u, "(A)")

    u_file = free_unit ()
    open (u_file, file = char (sample // ".mokka.evt"), &
         action = "read", status = "old")
    do
       read (u_file, "(A)", iostat = iostat)  buffer
       if (buffer(1:21) == "  <generator_version>")  buffer = "[...]"
       if (iostat /= 0)  exit
       write (u, "(A)") trim (buffer)
    end do
    close (u_file)
    
    write (u, "(A)")
    write (u, "(A)")  "* Reset data"
    write (u, "(A)")
 
    deallocate (eio)
    allocate (eio_ascii_mokka_t :: eio)
    
    select type (eio)
    type is (eio_ascii_mokka_t)
       call eio%set_parameters (keep_beams = .true.)
    end select
    call eio%write (u)

    write (u, "(A)")
    write (u, "(A)")  "* Cleanup"

    call eio_cleanup_test (event)

    write (u, "(A)")
    write (u, "(A)")  "* Test output end: eio_ascii_7"
    
  end subroutine eio_ascii_7
  
\eatline
\nwendcode{}\nwbegindocs{341}\nwdocspar
We test the implementation of all I/O methods, method {\tt{}short}:
\nwenddocs{}\nwbegincode{342}\moddef{EIO ascii: execute tests}\plusendmoddef
  call test (eio_ascii_8, "eio_ascii_8", &
       "read and write event contents, format [short]", &
       u, results)
\nwendcode{}\nwbegincode{343}\moddef{EIO ascii: tests}\plusendmoddef
  subroutine eio_ascii_8 (u)
    integer, intent(in) :: u
    class(generic_event_t), pointer :: event
    type(event_sample_data_t) :: data
    class(eio_t), allocatable :: eio
    type(string_t) :: sample
    integer :: u_file, iostat
    character(80) :: buffer

    write (u, "(A)")  "* Test output: eio_ascii_8"
    write (u, "(A)")  "*   Purpose: generate an event in ASCII short format"
    write (u, "(A)")  "*      and write weight to file"
    write (u, "(A)")

    write (u, "(A)")  "* Initialize test process"
 
    call eio_prepare_test (event, unweighted = .false.)
    
    call data%init (1)
    data%n_evt = 1
    data%n_beam = 2
    data%pdg_beam = 25
    data%energy_beam = 500
    data%proc_num_id = [42]
    data%cross_section(1) = 100
    data%error(1) = 1
    data%total_cross_section = sum (data%cross_section)

    write (u, "(A)")
    write (u, "(A)")  "* Generate and write an event"
    write (u, "(A)")
 
    sample = "eio_ascii_8"
 
    allocate (eio_ascii_short_t :: eio)
    
    select type (eio)
    class is (eio_ascii_t);  call eio%set_parameters ()
    end select
    call eio%init_out (sample, data)
    call event%generate (1, [0._default, 0._default])
    call event%evaluate_expressions ()

    call eio%output (event, i_prc = 1)
    call eio%write (u)
    call eio%final ()

    write (u, "(A)")
    write (u, "(A)")  "* File contents:"
    write (u, "(A)")

    u_file = free_unit ()
    open (u_file, file = char (sample // ".short.evt"), &
         action = "read", status = "old")
    do
       read (u_file, "(A)", iostat = iostat)  buffer
       if (buffer(1:21) == "  <generator_version>")  buffer = "[...]"
       if (iostat /= 0)  exit
       write (u, "(A)") trim (buffer)
    end do
    close (u_file)
    
    write (u, "(A)")
    write (u, "(A)")  "* Reset data"
    write (u, "(A)")
 
    deallocate (eio)
    allocate (eio_ascii_short_t :: eio)
    
    select type (eio)
    type is (eio_ascii_short_t)
       call eio%set_parameters (keep_beams = .true.)
    end select
    call eio%write (u)

    write (u, "(A)")
    write (u, "(A)")  "* Cleanup"
 
    call eio_cleanup_test (event)

    write (u, "(A)")
    write (u, "(A)")  "* Test output end: eio_ascii_8"
    
  end subroutine eio_ascii_8
  
\eatline
\nwendcode{}\nwbegindocs{344}\nwdocspar
We test the implementation of all I/O methods, method {\tt{}lha}  (old
LHA) in verbose version:
\nwenddocs{}\nwbegincode{345}\moddef{EIO ascii: execute tests}\plusendmoddef
  call test (eio_ascii_9, "eio_ascii_9", &
       "read and write event contents, format [lha_verb]", &
       u, results)
\nwendcode{}\nwbegincode{346}\moddef{EIO ascii: tests}\plusendmoddef
  subroutine eio_ascii_9 (u)
    integer, intent(in) :: u
    class(generic_event_t), pointer :: event
    type(event_sample_data_t) :: data
    class(eio_t), allocatable :: eio
    type(string_t) :: sample
    integer :: u_file, iostat
    character(80) :: buffer

    write (u, "(A)")  "* Test output: eio_ascii_9"
    write (u, "(A)")  "*   Purpose: generate an event in ASCII LHA verbose format"
    write (u, "(A)")  "*      and write weight to file"
    write (u, "(A)")

    write (u, "(A)")  "* Initialize test process"
 
    call eio_prepare_test (event, unweighted = .false.)
    
    call data%init (1)
    data%n_evt = 1
    data%n_beam = 2
    data%pdg_beam = 25
    data%energy_beam = 500
    data%proc_num_id = [42]
    data%cross_section(1) = 100
    data%error(1) = 1
    data%total_cross_section = sum (data%cross_section)

    write (u, "(A)")
    write (u, "(A)")  "* Generate and write an event"
    write (u, "(A)")
 
    sample = "eio_ascii_9"
 
    allocate (eio_ascii_lha_verb_t :: eio)
    
    select type (eio)
    class is (eio_ascii_t);  call eio%set_parameters ()
    end select
    call eio%init_out (sample, data)
    call event%generate (1, [0._default, 0._default])
    call event%evaluate_expressions ()

    call eio%output (event, i_prc = 1)
    call eio%write (u)
    call eio%final ()

    write (u, "(A)")
    write (u, "(A)")  "* File contents:"
    write (u, "(A)")

    u_file = free_unit ()
    open (u_file, file = char (sample // ".lha.verb"), &
         action = "read", status = "old")
    do
       read (u_file, "(A)", iostat = iostat)  buffer
       if (buffer(1:21) == "  <generator_version>")  buffer = "[...]"
       if (iostat /= 0)  exit
       write (u, "(A)") trim (buffer)
    end do
    close (u_file)
    
    write (u, "(A)")
    write (u, "(A)")  "* Reset data"
    write (u, "(A)")
 
    deallocate (eio)
    allocate (eio_ascii_lha_verb_t :: eio)
    
    select type (eio)
    type is (eio_ascii_lha_verb_t)
       call eio%set_parameters (keep_beams = .true.)
    end select
    call eio%write (u)

    write (u, "(A)")
    write (u, "(A)")  "* Cleanup"
 
    call eio_cleanup_test (event)

    write (u, "(A)")
    write (u, "(A)")  "* Test output end: eio_ascii_9"
    
  end subroutine eio_ascii_9
  
\eatline
\nwendcode{}\nwbegindocs{347}\nwdocspar
We test the implementation of all I/O methods, method {\tt{}hepevt{\char95}verb}:
\nwenddocs{}\nwbegincode{348}\moddef{EIO ascii: execute tests}\plusendmoddef
  call test (eio_ascii_10, "eio_ascii_10", &
       "read and write event contents, format [hepevt_verb]", &
       u, results)
\nwendcode{}\nwbegincode{349}\moddef{EIO ascii: tests}\plusendmoddef
  subroutine eio_ascii_10 (u)
    integer, intent(in) :: u
    class(generic_event_t), pointer :: event
    type(event_sample_data_t) :: data
    class(eio_t), allocatable :: eio
    type(string_t) :: sample
    integer :: u_file, iostat
    character(80) :: buffer

    write (u, "(A)")  "* Test output: eio_ascii_10"
    write (u, "(A)")  "*   Purpose: generate an event in ASCII hepevt verbose format"
    write (u, "(A)")  "*      and write weight to file"
    write (u, "(A)")

    write (u, "(A)")  "* Initialize test process"
 
    call eio_prepare_test (event, unweighted = .false.)
    
    call data%init (1)
    data%n_evt = 1
    data%n_beam = 2
    data%pdg_beam = 25
    data%energy_beam = 500
    data%proc_num_id = [42]
    data%cross_section(1) = 100
    data%error(1) = 1
    data%total_cross_section = sum (data%cross_section)

    write (u, "(A)")
    write (u, "(A)")  "* Generate and write an event"
    write (u, "(A)")
 
    sample = "eio_ascii_10"
 
    allocate (eio_ascii_hepevt_verb_t :: eio)
    
    select type (eio)
    class is (eio_ascii_t);  call eio%set_parameters ()
    end select
    call eio%init_out (sample, data)
    call event%generate (1, [0._default, 0._default])
    call event%evaluate_expressions ()

    call eio%output (event, i_prc = 1)
    call eio%write (u)
    call eio%final ()

    write (u, "(A)")
    write (u, "(A)")  "* File contents:"
    write (u, "(A)")

    u_file = free_unit ()
    open (u_file, file = char (sample // ".hepevt.verb"), &
         action = "read", status = "old")
    do
       read (u_file, "(A)", iostat = iostat)  buffer
       if (buffer(1:21) == "  <generator_version>")  buffer = "[...]"
       if (iostat /= 0)  exit
       write (u, "(A)") trim (buffer)
    end do
    close (u_file)
    
    write (u, "(A)")
    write (u, "(A)")  "* Reset data"
    write (u, "(A)")
 
    deallocate (eio)
    allocate (eio_ascii_hepevt_verb_t :: eio)
    
    select type (eio)
    type is (eio_ascii_hepevt_verb_t)
       call eio%set_parameters (keep_beams = .true.)
    end select
    call eio%write (u)

    write (u, "(A)")
    write (u, "(A)")  "* Cleanup"
 
    call eio_cleanup_test (event)

    write (u, "(A)")
    write (u, "(A)")  "* Test output end: eio_ascii_10"
    
  end subroutine eio_ascii_10
  
\eatline
\nwendcode{}\nwbegindocs{350}\nwdocspar
%%%%%%%%%%%%%%%%%%%%%%%%%%%%%%%%%%%%%%%%%%%%%%%%%%%%%%%%%%%%%%%%%%%%%%%%
\section{HEP Common Blocks}

Long ago, to transfer data between programs one had to set up a common
block and link both programs as libraries to the main executable.  The
HEP community standardizes several of those common blocks.

The modern way of data exchange uses data files with standard
formats.  However, the LHEF standard data format derives from a common
block (actually, two).

\whizard\ used to support those common blocks, and LHEF was
implemented via writing/reading blocks.  We still keep this
convention, but intend to eliminate common blocks (or any other static
storage) from the workflow in the future.  This will gain flexibility
towards concurrent running of program images.

We encapsulate everything here in a module.  The module holds the
variables which are part of the common block.  To access the common
block variables, we just have to {\tt{}use} this module.  (They are
nevertheless in the common block, since external software may access
it in this way.)

Note: This code is taken essentially unchanged from \whizard\ 2.1 and
does not (yet) provide unit tests.
\nwenddocs{}\nwbegincode{351}\moddef{\code{}hep{\char95}common.f90\edoc{}}\endmoddef
\LA{}File header\RA{}

module hep_common
  
  use kinds
  use io_units
\LA{}Use strings\RA{}
  use diagnostics
  use physics_defs, only: HADRON_REMNANT
  use physics_defs, only: HADRON_REMNANT_SINGLET
  use physics_defs, only: HADRON_REMNANT_TRIPLET
  use physics_defs, only: HADRON_REMNANT_OCTET
  use lorentz
  use flavors
  use colors
  use polarizations
  use model_data
  use particles
  use subevents

\LA{}Standard module head\RA{}

\LA{}HEP common: public\RA{}

\LA{}HEP common: interfaces\RA{}

\LA{}HEP common: parameters\RA{}

\LA{}HEP common: variables\RA{}

\LA{}HEP common: common blocks\RA{}

contains
  
\LA{}HEP common: procedures\RA{}

end module hep_common
\eatline
\nwendcode{}\nwbegindocs{352}\nwdocspar
\subsection{Event characteristics}
The maximal number of particles in an event record.
\nwenddocs{}\nwbegincode{353}\moddef{HEP common: parameters}\endmoddef
  integer, parameter, public :: MAXNUP = 500
\eatline
\nwendcode{}\nwbegindocs{354}The number of particles in this event.
\nwenddocs{}\nwbegincode{355}\moddef{HEP common: variables}\endmoddef
  integer, public :: NUP
\eatline
\nwendcode{}\nwbegindocs{356}The process ID for this event.
\nwenddocs{}\nwbegincode{357}\moddef{HEP common: variables}\plusendmoddef
  integer, public :: IDPRUP
\eatline
\nwendcode{}\nwbegindocs{358}The weight of this event ($\pm 1$ for unweighted events).
\nwenddocs{}\nwbegincode{359}\moddef{HEP common: variables}\plusendmoddef
  double precision, public :: XWGTUP
\eatline
\nwendcode{}\nwbegindocs{360}The factorization scale that is used for PDF calculation ($-1$ if
undefined).
\nwenddocs{}\nwbegincode{361}\moddef{HEP common: variables}\plusendmoddef
  double precision, public :: SCALUP
\eatline
\nwendcode{}\nwbegindocs{362}The QED and QCD couplings $\alpha$ used for this event ($-1$ if
undefined).
\nwenddocs{}\nwbegincode{363}\moddef{HEP common: variables}\plusendmoddef
  double precision, public :: AQEDUP
  double precision, public :: AQCDUP
\eatline
\nwendcode{}\nwbegindocs{364}\nwdocspar

\subsection{Particle characteristics}
The PDG code:
\nwenddocs{}\nwbegincode{365}\moddef{HEP common: variables}\plusendmoddef
  integer, dimension(MAXNUP) :: IDUP
\eatline
\nwendcode{}\nwbegindocs{366}The status code.  Incoming: $-1$, outgoing: $+1$.  Intermediate
t-channel propagator: $-2$ (currently not used by WHIZARD).
Intermediate resonance whose mass should be preserved: $2$.
Intermediate resonance for documentation: $3$ (currently not used).
Beam particles: $-9$.
\nwenddocs{}\nwbegincode{367}\moddef{HEP common: variables}\plusendmoddef
  integer, dimension(MAXNUP), public :: ISTUP
\eatline
\nwendcode{}\nwbegindocs{368}Index of first and last mother.
\nwenddocs{}\nwbegincode{369}\moddef{HEP common: variables}\plusendmoddef
  integer, dimension(2,MAXNUP), public :: MOTHUP
\eatline
\nwendcode{}\nwbegindocs{370}Color line index of the color and anticolor entry for the particle.
The standard recommends using large numbers; we start from MAXNUP+1. 
\nwenddocs{}\nwbegincode{371}\moddef{HEP common: variables}\plusendmoddef
  integer, dimension(2,MAXNUP), public :: ICOLUP
\eatline
\nwendcode{}\nwbegindocs{372}Momentum, energy, and invariant mass: $(p_x,p_y,p_z,E,M)$.  For
space-like particles, $M$ is the negative square root of the absolute
value of the invariant mass.
\nwenddocs{}\nwbegincode{373}\moddef{HEP common: variables}\plusendmoddef
  double precision, dimension(5,MAXNUP), public :: PUP
\eatline
\nwendcode{}\nwbegindocs{374}Invariant lifetime (distance) from production to decay in mm.
\nwenddocs{}\nwbegincode{375}\moddef{HEP common: variables}\plusendmoddef
  double precision, dimension(MAXNUP), public :: VTIMUP
\eatline
\nwendcode{}\nwbegindocs{376}Cosine of the angle between the spin-vector and a particle and the
3-momentum of its mother, given in the lab frame.  If
undefined/unpolarized: $9$.
\nwenddocs{}\nwbegincode{377}\moddef{HEP common: variables}\plusendmoddef
  double precision, dimension(MAXNUP), public :: SPINUP
\eatline
\nwendcode{}\nwbegindocs{378}\nwdocspar
\subsection{The HEPRUP common block}
This common block is filled once per run.

\subsubsection{Run characteristics}
The maximal number of different processes.
\nwenddocs{}\nwbegincode{379}\moddef{HEP common: parameters}\plusendmoddef
  integer, parameter, public :: MAXPUP = 100
\eatline
\nwendcode{}\nwbegindocs{380}The beam PDG codes.
\nwenddocs{}\nwbegincode{381}\moddef{HEP common: variables}\plusendmoddef
  integer, dimension(2), public :: IDBMUP
\eatline
\nwendcode{}\nwbegindocs{382}The beam energies in GeV.
\nwenddocs{}\nwbegincode{383}\moddef{HEP common: variables}\plusendmoddef
  double precision, dimension(2), public :: EBMUP
\eatline
\nwendcode{}\nwbegindocs{384}The PDF group and set for the two beams.  (Undefined: use $-1$;
LHAPDF: use group = $0$).
\nwenddocs{}\nwbegincode{385}\moddef{HEP common: variables}\plusendmoddef
  integer, dimension(2), public :: PDFGUP
  integer, dimension(2), public :: PDFSUP
\eatline
\nwendcode{}\nwbegindocs{386}The (re)weighting model.  1: events are weighted, the shower
generator (SHG) selects processes according to the maximum weight (in
pb) and unweights events.  2: events are weighted, the SHG selects
processes according to their cross section (in pb) and unweights
events.  3: events are unweighted and simply run through the SHG.  4:
events are weighted, and the SHG keeps the weight.  Negative numbers:
negative weights are allowed (and are reweighted to $\pm 1$ by the
SHG, if allowed).

\whizard\ only supports modes 3 and 4, as the SHG is not given control
over process selection.  This is consistent with writing events to
file, for offline showering.
\nwenddocs{}\nwbegincode{387}\moddef{HEP common: variables}\plusendmoddef
  integer, public :: IDWTUP
\eatline
\nwendcode{}\nwbegindocs{388}The number of different processes.
\nwenddocs{}\nwbegincode{389}\moddef{HEP common: variables}\plusendmoddef
  integer, public :: NPRUP
\eatline
\nwendcode{}\nwbegindocs{390}\nwdocspar
\subsubsection{Process characteristics}
Cross section and error in pb.  (Cross section is needed only for
${\tt{}IDWTUP} = 2$, so here both values are given for informational
purposes only.)
\nwenddocs{}\nwbegincode{391}\moddef{HEP common: variables}\plusendmoddef
  double precision, dimension(MAXPUP), public :: XSECUP
  double precision, dimension(MAXPUP), public :: XERRUP
\eatline
\nwendcode{}\nwbegindocs{392}Maximum weight, i.e., the maximum value that {\tt{}XWGTUP} can take.
Also unused for the supported weighting models. It is $\pm 1$ for
unweighted events.
\nwenddocs{}\nwbegincode{393}\moddef{HEP common: variables}\plusendmoddef
  double precision, dimension(MAXPUP), public :: XMAXUP
\eatline
\nwendcode{}\nwbegindocs{394}Internal ID of the selected process, matches {\tt{}IDPRUP} below.
\nwenddocs{}\nwbegincode{395}\moddef{HEP common: variables}\plusendmoddef
  integer, dimension(MAXPUP), public :: LPRUP
\eatline
\nwendcode{}\nwbegindocs{396}\nwdocspar
\subsubsection{The common block}
\nwenddocs{}\nwbegincode{397}\moddef{HEP common: common blocks}\endmoddef
  common /HEPRUP/ &
       IDBMUP, EBMUP, PDFGUP, PDFSUP, IDWTUP, NPRUP, &
       XSECUP, XERRUP, XMAXUP, LPRUP
  save /HEPRUP/

\eatline
\nwendcode{}\nwbegindocs{398}Fill the run characteristics of the common block.  The
initialization sets the beam properties, number of processes, and
weighting model.
\nwenddocs{}\nwbegincode{399}\moddef{HEP common: public}\endmoddef
  public :: heprup_init
\nwendcode{}\nwbegincode{400}\moddef{HEP common: procedures}\endmoddef
  subroutine heprup_init &
       (beam_pdg, beam_energy, n_processes, unweighted, negative_weights)
    integer, dimension(2), intent(in) :: beam_pdg
    real(default), dimension(2), intent(in) :: beam_energy
    integer, intent(in) :: n_processes
    logical, intent(in) :: unweighted
    logical, intent(in) :: negative_weights
    IDBMUP = beam_pdg
    EBMUP = beam_energy
    PDFGUP = -1
    PDFSUP = -1
    if (unweighted) then
       IDWTUP = 3
    else
       IDWTUP = 4
    end if
    if (negative_weights)  IDWTUP = - IDWTUP
    NPRUP = n_processes
  end subroutine heprup_init

\eatline
\nwendcode{}\nwbegindocs{401}Extract parameters from the common block.  We leave it to the caller
to specify which parameters it actually needs.

{\tt{}PDFGUP} and {\tt{}PDFSUP} are not extracted.  {\tt{}IDWTUP=1,2} are not
supported by \whizard, but correspond to weighted events.
\nwenddocs{}\nwbegincode{402}\moddef{HEP common: public}\plusendmoddef
  public :: heprup_get_run_parameters
\nwendcode{}\nwbegincode{403}\moddef{HEP common: procedures}\plusendmoddef
  subroutine heprup_get_run_parameters &
       (beam_pdg, beam_energy, n_processes, unweighted, negative_weights)
    integer, dimension(2), intent(out), optional :: beam_pdg
    real(default), dimension(2), intent(out), optional :: beam_energy
    integer, intent(out), optional :: n_processes
    logical, intent(out), optional :: unweighted
    logical, intent(out), optional :: negative_weights
    if (present (beam_pdg))  beam_pdg = IDBMUP
    if (present (beam_energy))  beam_energy = EBMUP
    if (present (n_processes))  n_processes = NPRUP
    if (present (unweighted)) then
       select case (abs (IDWTUP))
       case (3)
          unweighted = .true.
       case (4)
          unweighted = .false.
       case (1,2)  !!! not supported by WHIZARD
          unweighted = .false.
       case default
          call msg_fatal ("HEPRUP: unsupported IDWTUP value")
       end select
    end if
    if (present (negative_weights)) then
       negative_weights = IDWTUP < 0
    end if
  end subroutine heprup_get_run_parameters

\eatline
\nwendcode{}\nwbegindocs{404}Specify PDF set info.  Since we support only LHAPDF, the group entry
is zero.
\nwenddocs{}\nwbegincode{405}\moddef{HEP common: public}\plusendmoddef
  public :: heprup_set_lhapdf_id
\nwendcode{}\nwbegincode{406}\moddef{HEP common: procedures}\plusendmoddef
  subroutine heprup_set_lhapdf_id (i_beam, pdf_id)
    integer, intent(in) :: i_beam, pdf_id
    PDFGUP(i_beam) = 0
    PDFSUP(i_beam) = pdf_id
  end subroutine heprup_set_lhapdf_id

\eatline
\nwendcode{}\nwbegindocs{407}Fill the characteristics for a particular process.  Only the process
ID is mandatory.  Note that \whizard\ computes cross sections in fb,
so we have to rescale to pb.  The maximum weight is meaningless for
unweighted events.
\nwenddocs{}\nwbegincode{408}\moddef{HEP common: public}\plusendmoddef
  public :: heprup_set_process_parameters
\nwendcode{}\nwbegincode{409}\moddef{HEP common: procedures}\plusendmoddef
  subroutine heprup_set_process_parameters &
       (i, process_id, cross_section, error, max_weight)
    integer, intent(in) :: i, process_id
    real(default), intent(in), optional :: cross_section, error, max_weight
    real(default), parameter :: pb_per_fb = 1.e-3_default
    LPRUP(i) = process_id
    if (present (cross_section)) then
       XSECUP(i) = cross_section * pb_per_fb
    else
       XSECUP(i) = 0
    end if
    if (present (error)) then
       XERRUP(i) = error * pb_per_fb
    else
       XERRUP(i) = 0
    end if
    select case (IDWTUP)
    case (3);  XMAXUP(i) = 1
    case (4)
       if (present (max_weight)) then
          XMAXUP(i) = max_weight * pb_per_fb
       else
          XMAXUP(i) = 0
       end if
    end select
  end subroutine heprup_set_process_parameters

\eatline
\nwendcode{}\nwbegindocs{410}Extract the process parameters, as far as needed.
\nwenddocs{}\nwbegincode{411}\moddef{HEP common: public}\plusendmoddef
  public :: heprup_get_process_parameters
\nwendcode{}\nwbegincode{412}\moddef{HEP common: procedures}\plusendmoddef
  subroutine heprup_get_process_parameters  &
       (i, process_id, cross_section, error, max_weight)
    integer, intent(in) :: i
    integer, intent(out), optional :: process_id
    real(default), intent(out), optional :: cross_section, error, max_weight
    real(default), parameter :: pb_per_fb = 1.e-3_default
    if (present (process_id))  process_id = LPRUP(i)
    if (present (cross_section)) then
       cross_section = XSECUP(i) / pb_per_fb
    end if
    if (present (error)) then
       error = XERRUP(i) / pb_per_fb
    end if
    if (present (max_weight)) then
       select case (IDWTUP)
       case (3)
          max_weight = 1
       case (4)
          max_weight = XMAXUP(i) / pb_per_fb
       case (1,2)   !!! not supported by WHIZARD
          max_weight = 0
       case default
          call msg_fatal ("HEPRUP: unsupported IDWTUP value")
       end select
    end if
  end subroutine heprup_get_process_parameters

\eatline
\nwendcode{}\nwbegindocs{413}\nwdocspar
\subsection{Run parameter output (verbose)}
This is a verbose output of the HEPRUP block.
\nwenddocs{}\nwbegincode{414}\moddef{HEP common: public}\plusendmoddef
  public :: heprup_write_verbose
\nwendcode{}\nwbegincode{415}\moddef{HEP common: procedures}\plusendmoddef
  subroutine heprup_write_verbose (unit)
    integer, intent(in), optional :: unit
    integer :: u, i
    u = given_output_unit (unit);  if (u < 0)  return
    write (u, "(A)")  "HEPRUP Common Block"
    write (u, "(3x,A6,' = ',I9,3x,1x,I9,3x,8x,A)")  "IDBMUP", IDBMUP, &
         "PDG code of beams"
    write (u, "(3x,A6,' = ',G12.5,1x,G12.5,8x,A)")  "EBMUP ", EBMUP, &
         "Energy of beams in GeV"
    write (u, "(3x,A6,' = ',I9,3x,1x,I9,3x,8x,A)")  "PDFGUP", PDFGUP, &
         "PDF author group [-1 = undefined]"
    write (u, "(3x,A6,' = ',I9,3x,1x,I9,3x,8x,A)")  "PDFSUP", PDFSUP, &
         "PDF set ID       [-1 = undefined]"
    write (u, "(3x,A6,' = ',I9,3x,1x,9x,3x,8x,A)")  "IDWTUP", IDWTUP, &
         "LHA code for event weight mode"
    write (u, "(3x,A6,' = ',I9,3x,1x,9x,3x,8x,A)")  "NPRUP ", NPRUP, &
         "Number of user subprocesses"
    do i = 1, NPRUP
       write (u, "(1x,A,I0)")  "Subprocess #", i
       write (u, "(3x,A6,' = ',ES12.5,1x,12x,8x,A)")  "XSECUP", XSECUP(i), &
            "Cross section in pb"
       write (u, "(3x,A6,' = ',ES12.5,1x,12x,8x,A)")  "XERRUP", XERRUP(i), &
            "Cross section error in pb"
       write (u, "(3x,A6,' = ',ES12.5,1x,12x,8x,A)")  "XMAXUP", XMAXUP(i), &
            "Maximum event weight (cf. IDWTUP)"
       write (u, "(3x,A6,' = ',I9,3x,1x,12x,8x,A)")  "LPRUP ", LPRUP(i), &
            "Subprocess ID"
    end do
  end subroutine heprup_write_verbose

\eatline
\nwendcode{}\nwbegindocs{416}\nwdocspar
\subsection{Run parameter output (other formats)}
This routine writes the initialization block according to the LHEF
standard.  It uses the current contents of the HEPRUP block.
\nwenddocs{}\nwbegincode{417}\moddef{HEP common: public}\plusendmoddef
  public :: heprup_write_lhef
\nwendcode{}\nwbegincode{418}\moddef{HEP common: procedures}\plusendmoddef
  subroutine heprup_write_lhef (unit)
    integer, intent(in), optional :: unit
    integer :: u, i
    u = given_output_unit (unit);  if (u < 0)  return
    write (u, "(2(1x,I0),2(1x,ES17.10),6(1x,I0))") &
         IDBMUP, EBMUP, PDFGUP, PDFSUP, IDWTUP, NPRUP
    do i = 1, NPRUP
       write (u, "(3(1x,ES17.10),1x,I0)") &
            XSECUP(i), XERRUP(i), XMAXUP(i), LPRUP(i)
    end do
  end subroutine heprup_write_lhef
  
\eatline
\nwendcode{}\nwbegindocs{419}\nwdocspar
This routine is a complete dummy at the moment. It uses the current
contents of the HEPRUP block. At the end, it should depend on certain
input flags for the different ASCII event formats.
\nwenddocs{}\nwbegincode{420}\moddef{HEP common: public}\plusendmoddef
  public :: heprup_write_ascii
\nwendcode{}\nwbegincode{421}\moddef{HEP common: procedures}\plusendmoddef
  subroutine heprup_write_ascii (unit)
    integer, intent(in), optional :: unit
    integer :: u, i
    u = given_output_unit (unit);  if (u < 0)  return
    write (u, "(2(1x,I0),2(1x,ES17.10),6(1x,I0))") &
         IDBMUP, EBMUP, PDFGUP, PDFSUP, IDWTUP, NPRUP
    do i = 1, NPRUP
       write (u, "(3(1x,ES17.10),1x,I0)") &
            XSECUP(i), XERRUP(i), XMAXUP(i), LPRUP(i)
    end do
  end subroutine heprup_write_ascii
  
\eatline
\nwendcode{}\nwbegindocs{422}\nwdocspar
\subsubsection{Run parameter input (LHEF)}
In a LHEF file, the parameters are written in correct order on
separate lines, but we should not count on the precise format.
List-directed input should just work.
\nwenddocs{}\nwbegincode{423}\moddef{HEP common: public}\plusendmoddef
  public :: heprup_read_lhef
\nwendcode{}\nwbegincode{424}\moddef{HEP common: procedures}\plusendmoddef
  subroutine heprup_read_lhef (u)
    integer, intent(in) :: u
    integer :: i
    read (u, *) &
         IDBMUP, EBMUP, PDFGUP, PDFSUP, IDWTUP, NPRUP
    do i = 1, NPRUP
       read (u, *) &
            XSECUP(i), XERRUP(i), XMAXUP(i), LPRUP(i)
    end do
  end subroutine heprup_read_lhef
    
\eatline
\nwendcode{}\nwbegindocs{425}\nwdocspar
\subsection{The HEPEUP common block}
\nwenddocs{}\nwbegincode{426}\moddef{HEP common: common blocks}\plusendmoddef
  common /HEPEUP/ &
       NUP, IDPRUP, XWGTUP, SCALUP, AQEDUP, AQCDUP, &
       IDUP, ISTUP, MOTHUP, ICOLUP, PUP, VTIMUP, SPINUP
  save /HEPEUP/

\eatline
\nwendcode{}\nwbegindocs{427}\nwdocspar
\subsubsection{Initialization}
Fill the event characteristics of the common block.  The
initialization sets only the number of particles and initializes the
rest with default values.  The other routine sets the optional
parameters.
\nwenddocs{}\nwbegincode{428}\moddef{HEP common: public}\plusendmoddef
  public :: hepeup_init
  public :: hepeup_set_event_parameters
\nwendcode{}\nwbegincode{429}\moddef{HEP common: procedures}\plusendmoddef
  subroutine hepeup_init (n_tot)
    integer, intent(in) :: n_tot
    NUP = n_tot
    IDPRUP = 0
    XWGTUP = 1
    SCALUP = -1
    AQEDUP = -1
    AQCDUP = -1
  end subroutine hepeup_init

  subroutine hepeup_set_event_parameters &
       (proc_id, weight, scale, alpha_qed, alpha_qcd)
    integer, intent(in), optional :: proc_id
    real(default), intent(in), optional :: weight, scale, alpha_qed, alpha_qcd
    if (present (proc_id))   IDPRUP = proc_id
    if (present (weight))    XWGTUP = weight
    if (present (scale))     SCALUP = scale
    if (present (alpha_qed)) AQEDUP = alpha_qed
    if (present (alpha_qcd)) AQCDUP = alpha_qcd
  end subroutine hepeup_set_event_parameters

\eatline
\nwendcode{}\nwbegindocs{430}\nwdocspar
Extract event information.  The caller determines the parameters.
\nwenddocs{}\nwbegincode{431}\moddef{HEP common: public}\plusendmoddef
  public :: hepeup_get_event_parameters
\nwendcode{}\nwbegincode{432}\moddef{HEP common: procedures}\plusendmoddef
  subroutine hepeup_get_event_parameters &
       (proc_id, weight, scale, alpha_qed, alpha_qcd)
    integer, intent(out), optional :: proc_id
    real(default), intent(out), optional :: weight, scale, alpha_qed, alpha_qcd
    if (present (proc_id))   proc_id   = IDPRUP
    if (present (weight))    weight    = XWGTUP
    if (present (scale))     scale     = SCALUP
    if (present (alpha_qed)) alpha_qed = AQEDUP
    if (present (alpha_qcd)) alpha_qcd = AQCDUP
  end subroutine hepeup_get_event_parameters

\eatline
\nwendcode{}\nwbegindocs{433}\nwdocspar
\subsubsection{Particle data}
Below we need the particle status codes which are actually defined
in the {\tt{}subevents} module.

Set the entry for a specific particle.  All parameters are set with
the exception of lifetime and spin, where default values are stored.
\nwenddocs{}\nwbegincode{434}\moddef{HEP common: public}\plusendmoddef
  public :: hepeup_set_particle
\nwendcode{}\nwbegincode{435}\moddef{HEP common: procedures}\plusendmoddef
  subroutine hepeup_set_particle (i, pdg, status, parent, col, p, m2)
    integer, intent(in) :: i
    integer, intent(in) :: pdg, status
    integer, dimension(:), intent(in) :: parent
    type(vector4_t), intent(in) :: p
    integer, dimension(2), intent(in) :: col
    real(default), intent(in) :: m2
    if (i > MAXNUP) then
       call msg_error (arr=[ &
            var_str ("Too many particles in HEPEUP common block. " // &
                            "If this happened "), &
            var_str ("during event output, your events will be " // &
                            "invalid; please consider "), &
            var_str ("switching to a modern event format like HEPMC. " // &
                            "If you are not "), &
            var_str ("using an old, HEPEUP based format and " // &
                            "nevertheless get this error,"), &
            var_str ("please notify the WHIZARD developers,") ])
       return
    end if
    IDUP(i) = pdg
    select case (status)
    case (PRT_BEAM);         ISTUP(i) = -9
    case (PRT_INCOMING);     ISTUP(i) = -1
    case (PRT_BEAM_REMNANT); ISTUP(i) =  3
    case (PRT_OUTGOING);     ISTUP(i) =  1
    case (PRT_RESONANT);     ISTUP(i) =  2
    case (PRT_VIRTUAL);      ISTUP(i) =  3
    case default;            ISTUP(i) =  0
    end select
    select case (size (parent))
    case (0);      MOTHUP(:,i) = 0
    case (1);      MOTHUP(1,i) = parent(1); MOTHUP(2,i) = 0
    case default;  MOTHUP(:,i) = [ parent(1), parent(size (parent)) ]
    end select
    if (col(1) > 0) then
       ICOLUP(1,i) = 500 + col(1)
    else
       ICOLUP(1,i) = 0
    end if
    if (col(2) > 0) then
       ICOLUP(2,i) = 500 + col(2)
    else
       ICOLUP(2,i) = 0
    end if
    PUP(1:3,i) = vector3_get_components (space_part (p))
    PUP(4,i) = energy (p)
    PUP(5,i) = sign (sqrt (abs (m2)), m2)
    VTIMUP(i) = 0
    SPINUP(i) = 9
  end subroutine hepeup_set_particle

\eatline
\nwendcode{}\nwbegindocs{436}Set the lifetime, actually $c\tau$ measured im mm, where $\tau$ is
the invariant lifetime.
\nwenddocs{}\nwbegincode{437}\moddef{HEP common: public}\plusendmoddef
  public :: hepeup_set_particle_lifetime
\nwendcode{}\nwbegincode{438}\moddef{HEP common: procedures}\plusendmoddef
  subroutine hepeup_set_particle_lifetime (i, lifetime)
    integer, intent(in) :: i
    real(default), intent(in) :: lifetime
    VTIMUP(i) = lifetime
  end subroutine hepeup_set_particle_lifetime

\eatline
\nwendcode{}\nwbegindocs{439}Set the particle spin entry.  We need the cosine of the angle of the
spin axis with respect to the three-momentum of the parent particle.

If the particle has a full polarization density matrix given, we need
the particle momentum and polarization as well as the mother-particle
momentum.  The polarization is transformed into a spin vector (which
is sensible only for spin-1/2 or massless particles), which then is
transformed into the lab frame (by a rotation of the 3-axis to the
particle momentum axis).  Finally, we compute the scalar product of
this vector with the mother-particle three-momentum.

This puts severe restrictions on the applicability of this definition,
and Lorentz invariance is lost.  Unfortunately, the Les Houches Accord
requires this computation.
\nwenddocs{}\nwbegincode{440}\moddef{HEP common: public}\plusendmoddef
  public :: hepeup_set_particle_spin
\nwendcode{}\nwbegincode{441}\moddef{HEP common: interfaces}\endmoddef
  interface hepeup_set_particle_spin
     module procedure hepeup_set_particle_spin_pol
  end interface
\nwendcode{}\nwbegincode{442}\moddef{HEP common: procedures}\plusendmoddef
  subroutine hepeup_set_particle_spin_pol (i, p, pol, p_mother)
    integer, intent(in) :: i
    type(vector4_t), intent(in) :: p
    type(polarization_t), intent(in) :: pol
    type(vector4_t), intent(in) :: p_mother
    type(vector3_t) :: s3, p3
    type(vector4_t) :: s4
    s3 = vector3_moving (polarization_get_axis (pol))
    p3 = space_part (p)
    s4 = rotation_to_2nd (3, p3) * vector4_moving (0._default, s3)
    SPINUP(i) = enclosed_angle_ct (s4, p_mother)
  end subroutine hepeup_set_particle_spin_pol
    
\eatline
\nwendcode{}\nwbegindocs{443}\nwdocspar
Extract particle data.  The caller decides which ones to retrieve.

Status codes: beam remnants share the status code with virtual particles.
However, for the purpose of WHIZARD we should identify them.  We
use the PDG code for this.
\nwenddocs{}\nwbegincode{444}\moddef{HEP common: public}\plusendmoddef
  public :: hepeup_get_particle
\nwendcode{}\nwbegincode{445}\moddef{HEP common: procedures}\plusendmoddef
  subroutine hepeup_get_particle (i, pdg, status, parent, col, p, m2)
    integer, intent(in) :: i
    integer, intent(out), optional :: pdg, status
    integer, dimension(:), intent(out), optional :: parent
    type(vector4_t), intent(out), optional :: p
    integer, dimension(2), intent(out), optional :: col
    real(default), dimension(5,MAXNUP) :: pup_def
    real(default), intent(out), optional :: m2
    if (present (pdg))  pdg = IDUP(i)
    if (present (status)) then
       select case (ISTUP(i))
       case (-9);  status = PRT_BEAM
       case (-1);  status = PRT_INCOMING
       case (1);   status = PRT_OUTGOING
       case (2);   status = PRT_RESONANT
       case (3);
          select case (abs (IDUP(i)))
          case (HADRON_REMNANT, HADRON_REMNANT_SINGLET, &
               HADRON_REMNANT_TRIPLET, HADRON_REMNANT_OCTET)
             status = PRT_BEAM_REMNANT
          case default
             status = PRT_VIRTUAL
          end select
       case default
          status = PRT_UNDEFINED
       end select
    end if
    if (present (parent)) then
       select case (size (parent))
       case (0)
       case (1);    parent(1) = MOTHUP(1,i)
       case (2);    parent = MOTHUP(:,i)
       end select
    end if
    if (present (col)) then
       col = ICOLUP(:,i)
    end if
    if (present (p)) then
       pup_def = PUP
       p = vector4_moving (pup_def(4,i), vector3_moving (pup_def(1:3,i)))
    end if
    if (present (m2)) then
       m2 = sign (PUP(5,i) ** 2, PUP(5,i))
    end if
  end subroutine hepeup_get_particle

\eatline
\nwendcode{}\nwbegindocs{446}\nwdocspar
\subsection{The HEPEVT common block}

For the LEP Monte Carlos, a standard common block has been proposed
in AKV89.  We strongly recommend its use.  (The description is an
abbreviated transcription of AKV89, Vol. 3, pp. 327-330).

{\tt{}NMXHEP} is the maximum number of entries:
\nwenddocs{}\nwbegincode{447}\moddef{HEP common: variables}\plusendmoddef
  integer, parameter :: NMXHEP = 4000

\eatline
\nwendcode{}\nwbegindocs{448}{\tt{}NEVHEP} is normally the event number, but may take special
values as follows:

   0   the program does not keep track of event numbers.
  -1   a special initialization record.
  -2   a special final record.
\nwenddocs{}\nwbegincode{449}\moddef{HEP common: variables}\plusendmoddef
  integer :: NEVHEP

\eatline
\nwendcode{}\nwbegindocs{450} {\tt{}NHEP} holds the number of entries for this event.
\nwenddocs{}\nwbegincode{451}\moddef{HEP common: variables}\plusendmoddef
  integer :: NHEP

\eatline
\nwendcode{}\nwbegindocs{452}The entry {\tt{}ISTHEP(N)} gives the status code for the {\tt{}N}th entry,
with the following semantics:
   0       a null entry.
   1       an existing entry, which has not decayed or fragmented.
   2       a decayed or fragmented entry, which is retained for
           event history information.
   3       documentation line.
   4- 10   reserved for future standards.
  11-200   at the disposal of each model builder.
 201-      at the disposal of users.
\nwenddocs{}\nwbegincode{453}\moddef{HEP common: variables}\plusendmoddef
  integer, dimension(NMXHEP) :: ISTHEP

\eatline
\nwendcode{}\nwbegindocs{454}\nwdocspar
The Particle Data Group has proposed standard particle codes,
which are to be stored in {\tt{}IDHEP(N)}.
\nwenddocs{}\nwbegincode{455}\moddef{HEP common: variables}\plusendmoddef
  integer, dimension(NMXHEP) :: IDHEP

\eatline
\nwendcode{}\nwbegindocs{456}{\tt{}JMOHEP(1,N)} points to the mother of the {\tt{}N}th entry, if any.
It is set to zero for initial entries.
{\tt{}JMOHEP(2,N)} points to the second mother, if any.
\nwenddocs{}\nwbegincode{457}\moddef{HEP common: variables}\plusendmoddef
  integer, dimension(2, NMXHEP) :: JMOHEP

\eatline
\nwendcode{}\nwbegindocs{458}{\tt{}JDAHEP(1,N)} and {\tt{}JDAHEP(2,N)} point to the first and last daughter 
of the {\tt{}N}th entry, if any.  These are zero for entries which have not
yet decayed.  The other daughters are stored in between these two.
\nwenddocs{}\nwbegincode{459}\moddef{HEP common: variables}\plusendmoddef
  integer, dimension(2, NMXHEP) :: JDAHEP

\eatline
\nwendcode{}\nwbegindocs{460}In {\tt{}PHEP} we store the momentum of the particle, more specifically
this means that {\tt{}PHEP(1,N)}, {\tt{}PHEP(2,N)}, and {\tt{}PHEP(3,N)} contain the
momentum in the $x$, $y$, and $z$ direction (as defined by the machine
people), measured in GeV/c.  {\tt{}PHEP(4,N)} contains the energy in GeV
and {\tt{}PHEP(5,N)} the mass in GeV$/c^2$.  The latter may be negative for
spacelike partons.
\nwenddocs{}\nwbegincode{461}\moddef{HEP common: variables}\plusendmoddef
  double precision, dimension(5, NMXHEP) :: PHEP
  
\eatline
\nwendcode{}\nwbegindocs{462}Finally {\tt{}VHEP} is the place to store the position of the production
vertex.  {\tt{}VHEP(1,N)}, {\tt{}VHEP(2,N)}, and {\tt{}VHEP(3,N)} contain the $x$, $y$,
and $z$ coordinate (as defined by the machine people), measured in mm.
{\tt{}VHEP(4,N)} contains the production time in mm/c.
\nwenddocs{}\nwbegincode{463}\moddef{HEP common: variables}\plusendmoddef
  double precision, dimension(4, NMXHEP) :: VHEP
  
\eatline
\nwendcode{}\nwbegindocs{464}As an amendment to the proposed standard common block HEPEVT, we
also have a polarisation common block HEPSPN, as described in
AKV89.  {\tt{}SHEP(1,N)}, {\tt{}SHEP(2,N)}, and {\tt{}SHEP(3,N)} give the $x$, $y$, and $z$
component of the spinvector $s$ of a fermion in the fermions restframe.

Furthermore, we add the polarization of the corresponding outgoing
particles:
\nwenddocs{}\nwbegincode{465}\moddef{HEP common: variables}\plusendmoddef
  integer, dimension(NMXHEP) :: hepevt_pol

\eatline
\nwendcode{}\nwbegindocs{466}\nwdocspar
By convention, {\tt{}SHEP(4,N)} is always 1. All this is taken from StdHep
4.06 manual and written using Fortran90 conventions.
\nwenddocs{}\nwbegincode{467}\moddef{HEP common: common blocks}\plusendmoddef
  common /HEPEVT/ &
       NEVHEP, NHEP, ISTHEP, IDHEP, &
       JMOHEP, JDAHEP, PHEP, VHEP
  save /HEPEVT/
  
\eatline
\nwendcode{}\nwbegindocs{468}Here we store HEPEVT parameters of the WHIZARD 1 realization which
are not part of the HEPEVT common block. 
\nwenddocs{}\nwbegincode{469}\moddef{HEP common: variables}\plusendmoddef
  integer :: hepevt_n_out, hepevt_n_remnants

\eatline
\nwendcode{}\nwbegindocs{470} 
\nwenddocs{}\nwbegincode{471}\moddef{HEP common: variables}\plusendmoddef
  double precision :: hepevt_weight, hepevt_function_value
  double precision :: hepevt_function_ratio
  
\eatline
\nwendcode{}\nwbegindocs{472}Filling HEPEVT: If the event count is not provided, set {\tt{}NEVHEP}
to zero.  If the event count is {\tt{}-1} or {\tt{}-2}, the record
corresponds to initialization and finalization, and the event is
irrelevant.

Note that the event count may be larger than $2^{31}$ (2 GEvents).  In
that case, cut off the upper bits since {\tt{}NEVHEP} is probably limited
to default integer. 
\nwenddocs{}\nwbegincode{473}\moddef{HEP common: public}\plusendmoddef
  public :: hepevt_init
  public :: hepevt_set_event_parameters
\nwendcode{}\nwbegincode{474}\moddef{HEP common: procedures}\plusendmoddef
  subroutine hepevt_init (n_tot, n_out)
    integer, intent(in) :: n_tot, n_out
    NHEP              = n_tot
    NEVHEP            = 0
    hepevt_n_out      = n_out
    hepevt_n_remnants = 0
    hepevt_weight     = 1
    hepevt_function_value = 0
    hepevt_function_ratio = 1
  end subroutine hepevt_init
  
  subroutine hepevt_set_event_parameters &
       (weight, function_value, function_ratio, i_evt)
    integer, intent(in), optional :: i_evt
    real(default), intent(in), optional :: weight, function_value, &
       function_ratio
    if (present (i_evt)) NEVHEP = i_evt
    if (present (weight)) hepevt_weight = weight
    if (present (function_value)) hepevt_function_value = &
         function_value
    if (present (function_ratio)) hepevt_function_ratio = &
         function_ratio
  end subroutine hepevt_set_event_parameters

\eatline
\nwendcode{}\nwbegindocs{475}Set the entry for a specific particle.  All parameters are set with
the exception of lifetime and spin, where default values are stored.
\nwenddocs{}\nwbegincode{476}\moddef{HEP common: public}\plusendmoddef
  public :: hepevt_set_particle
\nwendcode{}\nwbegincode{477}\moddef{HEP common: procedures}\plusendmoddef
  subroutine hepevt_set_particle (i, pdg, status, parent, child, p, m2, hel)
    integer, intent(in) :: i
    integer, intent(in) :: pdg, status
    integer, dimension(:), intent(in) :: parent
    integer, dimension(:), intent(in) :: child
    type(vector4_t), intent(in) :: p
    real(default), intent(in) :: m2
    integer, intent(in) :: hel
    IDHEP(i) = pdg
    select case (status)
      case (PRT_BEAM);      ISTHEP(i) = 2
      case (PRT_INCOMING);  ISTHEP(i) = 2
      case (PRT_OUTGOING);  ISTHEP(i) = 1
      case (PRT_VIRTUAL);   ISTHEP(i) = 2
      case (PRT_RESONANT);  ISTHEP(i) = 2
      case default;         ISTHEP(i) = 0
    end select
    select case (size (parent))
    case (0);      JMOHEP(:,i) = 0
    case (1);      JMOHEP(1,i) = parent(1); JMOHEP(2,i) = 0
    case default;  JMOHEP(:,i) = [ parent(1), parent(size (parent)) ]
    end select
    select case (size (child))
    case (0);      JDAHEP(:,i) = 0
    case (1);      JDAHEP(:,i) = child(1)
    case default;  JDAHEP(:,i) = [ child(1), child(size (child)) ]
    end select
    PHEP(1:3,i) = vector3_get_components (space_part (p))
    PHEP(4,i) = energy (p)
    PHEP(5,i) = sign (sqrt (abs (m2)), m2)
    VHEP(1:4,i) = 0
    hepevt_pol(i) = hel
  end subroutine hepevt_set_particle

\eatline
\nwendcode{}\nwbegindocs{478}\nwdocspar
\subsection{Event output}
This is a verbose output of the HEPEVT block.
\nwenddocs{}\nwbegincode{479}\moddef{HEP common: public}\plusendmoddef
  public :: hepevt_write_verbose
\nwendcode{}\nwbegincode{480}\moddef{HEP common: procedures}\plusendmoddef
  subroutine hepevt_write_verbose (unit)
    integer, intent(in), optional :: unit
    integer :: u, i
    u = given_output_unit (unit);  if (u < 0)  return
    write (u, "(A)")  "HEPEVT Common Block"
    write (u, "(3x,A6,' = ',I9,3x,1x,20x,A)")  "NEVHEP", NEVHEP, &
         "Event number"
    write (u, "(3x,A6,' = ',I9,3x,1x,20x,A)")  "NHEP  ", NHEP, &
         "Number of particles in event"
    do i = 1, NHEP
       write (u, "(1x,A,I0)")  "Particle #", i
       write (u, "(3x,A6,' = ',I9,3x,1x,20x,A)", advance="no") &
            "ISTHEP", ISTHEP(i), "Status code: "
       select case (ISTHEP(i))
       case ( 0);  write (u, "(A)")  "null entry"
       case ( 1);  write (u, "(A)")  "outgoing"
       case ( 2);  write (u, "(A)")  "decayed"
       case ( 3);  write (u, "(A)")  "documentation"
       case (4:10);  write (u, "(A)")  "[unspecified]"
       case (11:200);  write (u, "(A)")  "[model-specific]"
       case (201:);  write (u, "(A)")  "[user-defined]"
       case default;  write (u, "(A)")  "[undefined]"
       end select
       write (u, "(3x,A6,' = ',I9,3x,1x,20x,A)")  "IDHEP ", IDHEP(i), &
            "PDG code of particle"
       write (u, "(3x,A6,' = ',I9,3x,1x,I9,3x,8x,A)")  "JMOHEP", JMOHEP(:,i), &
            "Index of first/second mother"
       write (u, "(3x,A6,' = ',I9,3x,1x,I9,3x,8x,A)")  "JDAHEP", JDAHEP(:,i), &
            "Index of first/last daughter"
       write (u, "(3x,A6,' = ',ES12.5,1x,ES12.5,8x,A)")  "PHEP12", &
            PHEP(1:2,i), "Transversal momentum (x/y) in GeV"
       write (u, "(3x,A6,' = ',ES12.5,1x,12x,8x,A)")  "PHEP3 ", PHEP(3,i), &
            "Longitudinal momentum (z) in GeV"
       write (u, "(3x,A6,' = ',ES12.5,1x,12x,8x,A)")  "PHEP4 ", PHEP(4,i), &
            "Energy in GeV"
       write (u, "(3x,A6,' = ',ES12.5,1x,12x,8x,A)")  "PHEP5 ", PHEP(5,i), &
            "Invariant mass in GeV"
       write (u, "(3x,A6,' = ',ES12.5,1x,ES12.5,8x,A)")  "VHEP12", VHEP(1:2,i), &
            "Transversal displacement (xy) in mm"
       write (u, "(3x,A6,' = ',ES12.5,1x,12x,8x,A)")  "VHEP3 ", VHEP(3,i), &
            "Longitudinal displacement (z) in mm"
       write (u, "(3x,A6,' = ',ES12.5,1x,12x,8x,A)")  "VHEP4 ", VHEP(4,i), &
            "Production time in mm"
    end do
  end subroutine hepevt_write_verbose

\eatline
\nwendcode{}\nwbegindocs{481}\nwdocspar
This is a verbose output of the HEPEUP block.
\nwenddocs{}\nwbegincode{482}\moddef{HEP common: public}\plusendmoddef
  public :: hepeup_write_verbose
\nwendcode{}\nwbegincode{483}\moddef{HEP common: procedures}\plusendmoddef
  subroutine hepeup_write_verbose (unit)
    integer, intent(in), optional :: unit
    integer :: u, i
    u = given_output_unit (unit);  if (u < 0)  return
    write (u, "(A)")  "HEPEUP Common Block"
    write (u, "(3x,A6,' = ',I9,3x,1x,20x,A)")  "NUP   ", NUP, &
         "Number of particles in event"
    write (u, "(3x,A6,' = ',I9,3x,1x,20x,A)")  "IDPRUP", IDPRUP, &
         "Subprocess ID"
    write (u, "(3x,A6,' = ',ES12.5,1x,20x,A)")  "XWGTUP", XWGTUP, &
         "Event weight"
    write (u, "(3x,A6,' = ',ES12.5,1x,20x,A)")  "SCALUP", SCALUP, &
         "Event energy scale in GeV"
    write (u, "(3x,A6,' = ',ES12.5,1x,20x,A)")  "AQEDUP", AQEDUP, &
         "QED coupling [-1 = undefined]"
    write (u, "(3x,A6,' = ',ES12.5,1x,20x,A)")  "AQCDUP", AQCDUP, &
         "QCD coupling [-1 = undefined]"
    do i = 1, NUP
       write (u, "(1x,A,I0)")  "Particle #", i
       write (u, "(3x,A6,' = ',I9,3x,1x,20x,A)")  "IDUP  ", IDUP(i), &
            "PDG code of particle"
       write (u, "(3x,A6,' = ',I9,3x,1x,20x,A)", advance="no") &
            "ISTUP ", ISTUP(i), "Status code: "
       select case (ISTUP(i))
       case (-1);  write (u, "(A)")  "incoming"
       case ( 1);  write (u, "(A)")  "outgoing"
       case (-2);  write (u, "(A)")  "spacelike"
       case ( 2);  write (u, "(A)")  "resonance"
       case ( 3);  write (u, "(A)")  "resonance (doc)"
       case (-9);  write (u, "(A)")  "beam"
       case default;  write (u, "(A)")  "[undefined]"
       end select
       write (u, "(3x,A6,' = ',I9,3x,1x,I9,3x,8x,A)")  "MOTHUP", MOTHUP(:,i), &
            "Index of first/last mother"
       write (u, "(3x,A6,' = ',I9,3x,1x,I9,3x,8x,A)")  "ICOLUP", ICOLUP(:,i), &
            "Color/anticolor flow index"
       write (u, "(3x,A6,' = ',ES12.5,1x,ES12.5,8x,A)")  "PUP1/2", PUP(1:2,i), &
            "Transversal momentum (x/y) in GeV"
       write (u, "(3x,A6,' = ',ES12.5,1x,12x,8x,A)")  "PUP3  ", PUP(3,i), &
            "Longitudinal momentum (z) in GeV"
       write (u, "(3x,A6,' = ',ES12.5,1x,12x,8x,A)")  "PUP4  ", PUP(4,i), &
            "Energy in GeV"
       write (u, "(3x,A6,' = ',ES12.5,1x,12x,8x,A)")  "PUP5  ", PUP(5,i), &
            "Invariant mass in GeV"
       write (u, "(3x,A6,' = ',ES12.5,1x,12x,8x,A)")  "VTIMUP", VTIMUP(i), &
            "Invariant lifetime in mm"
       write (u, "(3x,A6,' = ',ES12.5,1x,12x,8x,A)")  "SPINUP", SPINUP(i), &
            "cos(spin angle) [9 = undefined]"
    end do
  end subroutine hepeup_write_verbose

\eatline
\nwendcode{}\nwbegindocs{484}\nwdocspar
\subsection{Event output in various formats}
This routine writes event output according to the LHEF standard.  It
uses the current contents of the HEPEUP block.
\nwenddocs{}\nwbegincode{485}\moddef{HEP common: public}\plusendmoddef
  public :: hepeup_write_lhef
  public :: hepeup_write_lha
\nwendcode{}\nwbegincode{486}\moddef{HEP common: procedures}\plusendmoddef
  subroutine hepeup_write_lhef (unit)
    integer, intent(in), optional :: unit
    integer :: u, i
    u = given_output_unit (unit);  if (u < 0)  return
    write (u, "(2(1x,I0),4(1x,ES17.10))") &
         NUP, IDPRUP, XWGTUP, SCALUP, AQEDUP, AQCDUP
    do i = 1, NUP
       write (u, "(6(1x,I0),7(1x,ES17.10))") &
            IDUP(i), ISTUP(i), MOTHUP(:,i), ICOLUP(:,i), &
            PUP(:,i), VTIMUP(i), SPINUP(i)
    end do
  end subroutine hepeup_write_lhef

  subroutine hepeup_write_lha (unit)
    integer, intent(in), optional :: unit
    integer :: u, i
    integer, dimension(MAXNUP) :: spin_up
    spin_up = SPINUP
    u = given_output_unit (unit);  if (u < 0)  return
    write (u, "(2(1x,I5),1x,ES17.10,3(1x,ES13.6))") &
         NUP, IDPRUP, XWGTUP, SCALUP, AQEDUP, AQCDUP
    write (u, "(500(1x,I5))") IDUP(:NUP)
    write (u, "(500(1x,I5))") MOTHUP(1,:NUP)
    write (u, "(500(1x,I5))") MOTHUP(2,:NUP)
    write (u, "(500(1x,I5))") ICOLUP(1,:NUP)
    write (u, "(500(1x,I5))") ICOLUP(2,:NUP)
    write (u, "(500(1x,I5))") ISTUP(:NUP)
    write (u, "(500(1x,I5))") spin_up(:NUP)
    do i = 1, NUP
            write (u, "(1x,I5,4(1x,ES17.10))") i, PUP([ 4,1,2,3 ], i)
    end do

  end subroutine hepeup_write_lha  

\eatline
\nwendcode{}\nwbegindocs{487}This routine writes event output according to the HEPEVT standard.  It
uses the current contents of the HEPEVT block and some additional
parameters according to the standard in WHIZARD 1. For the long ASCII
format, the value of the sample function (i.e. the product of squared
matrix element, structure functions and phase space factor is printed out).
The option of reweighting matrix elements with respect to some
reference cross section is not implemented in WHIZARD 2 for this event
format, therefore the second entry in the long ASCII format (the
function ratio) is always one. The ATHENA format is an implementation
of the HEPEVT format that is readable by the ATLAS ATHENA software
framework. It is very similar to the WHIZARD 1 HEPEVT format, except
that it contains an event counter, a particle counter inside the
event, and has the HEPEVT {\tt{}ISTHEP} status before the PDG code. The
MOKKA format is a special ASCII format  that contains the information
to be parsed to the MOKKA LC fast simulation software.
\nwenddocs{}\nwbegincode{488}\moddef{HEP common: public}\plusendmoddef
  public :: hepevt_write_hepevt
  public :: hepevt_write_ascii
  public :: hepevt_write_athena
  public :: hepevt_write_mokka  
\nwendcode{}\nwbegincode{489}\moddef{HEP common: procedures}\plusendmoddef
  subroutine hepevt_write_hepevt (unit)
    integer, intent(in), optional :: unit
    integer :: u, i
    u = given_output_unit (unit);  if (u < 0)  return
    write (u, "(3(1x,I0),(1x,ES17.10))") &
         NHEP, hepevt_n_out, hepevt_n_remnants, hepevt_weight
    do i = 1, NHEP
       write (u, "(7(1x,I0))") &
            ISTHEP(i), IDHEP(i), JMOHEP(:,i), JDAHEP(:,i), hepevt_pol(i)
       write (u, "(5(1x,ES17.10))") PHEP(:,i)
       write (u, "(5(1x,ES17.10))") VHEP(:,i), 0.d0
    end do
  end subroutine hepevt_write_hepevt
  
  subroutine hepevt_write_ascii (unit, long)
    integer, intent(in), optional :: unit
    logical, intent(in) :: long   
    integer :: u, i
    u = given_output_unit (unit);  if (u < 0)  return
    write (u, "(3(1x,I0),(1x,ES17.10))") &
         NHEP, hepevt_n_out, hepevt_n_remnants, hepevt_weight
    do i = 1, NHEP
       if (ISTHEP(i) /= 1)  cycle
       write (u, "(2(1x,I0))") IDHEP(i), hepevt_pol(i)
       write (u, "(5(1x,ES17.10))") PHEP(:,i)
    end do
    if (long) then 
       write (u, "(2(1x,ES17.10))") &
            hepevt_function_value, hepevt_function_ratio
    end if
  end subroutine hepevt_write_ascii
  
  subroutine hepevt_write_athena (unit, i_evt)
    integer, intent(in), optional :: unit, i_evt
    integer :: u, i, num_event
    num_event = 0
    if (present (i_evt)) num_event = i_evt
    u = given_output_unit (unit);  if (u < 0)  return
    write (u, "(2(1x,I0))") num_event, NHEP
    do i = 1, NHEP
       write (u, "(7(1x,I0))") &
            i, ISTHEP(i), IDHEP(i), JMOHEP(:,i), JDAHEP(:,i)
       write (u, "(5(1x,ES17.10))") PHEP(:,i)
       write (u, "(5(1x,ES17.10))") VHEP(1:4,i)
    end do
  end subroutine hepevt_write_athena
  
  subroutine hepevt_write_mokka (unit)
    integer, intent(in), optional :: unit
    integer :: u, i
    u = given_output_unit (unit);  if (u < 0)  return
    write (u, "(3(1x,I0),(1x,ES17.10))") &
         NHEP, hepevt_n_out, hepevt_n_remnants, hepevt_weight
    do i = 1, NHEP
       write (u, "(4(1x,I0),4(1x,ES17.10))") &
            ISTHEP(i), IDHEP(i), JDAHEP(1,i), JDAHEP(2,i), &
            PHEP(1:3,i), PHEP(5,i)
    end do
  end subroutine hepevt_write_mokka
  
\eatline
\eatline
\nwendcode{}\nwbegindocs{490}\nwdocspar
\subsection{Event input in various formats}
This routine writes event output according to the LHEF standard.  It
uses the current contents of the HEPEUP block.
\nwenddocs{}\nwbegincode{491}\moddef{HEP common: public}\plusendmoddef
  public :: hepeup_read_lhef
\nwendcode{}\nwbegincode{492}\moddef{HEP common: procedures}\plusendmoddef
  subroutine hepeup_read_lhef (u)
    integer, intent(in) :: u
    integer :: i
    read (u, *) &
         NUP, IDPRUP, XWGTUP, SCALUP, AQEDUP, AQCDUP
    do i = 1, NUP
       read (u, *) &
            IDUP(i), ISTUP(i), MOTHUP(:,i), ICOLUP(:,i), &
            PUP(:,i), VTIMUP(i), SPINUP(i)
    end do
  end subroutine hepeup_read_lhef

\eatline
\nwendcode{}\nwbegindocs{493}\nwdocspar
\subsection{Data Transfer: particle sets}
The \whizard\ format for handling particle data in events is
{\tt{}particle{\char95}set{\char95}t}.  We have to interface this to the common blocks.

We first create a new particle set that contains only the particles
that are supported by the LHEF format.  These are: beam, incoming,
resonant, outgoing.  We drop particles with unknown, virtual or
beam-remnant status.

From this set we fill the common block.  Event information such as
process ID and weight is not transferred here; this has to be done by
the caller.  The spin information is set only if the particle has a
unique mother, and if its polarization is fully defined.
\nwenddocs{}\nwbegincode{494}\moddef{HEP common: public}\plusendmoddef
  public :: hepeup_from_particle_set
\nwendcode{}\nwbegincode{495}\moddef{HEP common: procedures}\plusendmoddef
  subroutine hepeup_from_particle_set (particle_set, keep_beams, keep_remnants)
    type(particle_set_t), intent(in) :: particle_set
    logical, intent(in), optional :: keep_beams
    logical, intent(in), optional :: keep_remnants
    type(particle_set_t), target :: pset_hepevt
    integer :: i, n_parents, status, n_tot
    integer, dimension(1) :: i_mother
    logical :: activate_remnants
    activate_remnants = .true.
    if (present (keep_remnants))  activate_remnants = keep_remnants
    call particle_set%to_hepevt_form (pset_hepevt, keep_beams)
    n_tot = pset_hepevt%get_n_tot ()
    call hepeup_init (n_tot)
    do i = 1, n_tot
       associate (prt => pset_hepevt%prt(i))
         status = prt%get_status ()
         if (activate_remnants &
              .and. status == PRT_BEAM_REMNANT &
              .and. prt%get_n_children () == 0) &
              status = PRT_OUTGOING
         call hepeup_set_particle (i, &
              prt%get_pdg (), &
              status, &
              prt%get_parents (), &
              prt%get_color (), &
              prt%get_momentum (), &
              prt%get_p2 ())
         n_parents = prt%get_n_parents ()
         if (n_parents == 1) then
            i_mother = prt%get_parents ()
            select case (prt%get_polarization_status ())
            case (PRT_GENERIC_POLARIZATION)
               call hepeup_set_particle_spin (i, &
                    prt%get_momentum (), &
                    prt%get_polarization (), &
                    pset_hepevt%prt(i_mother(1))%get_momentum ())
            end select
         end if
       end associate
    end do
    call pset_hepevt%final ()
  end subroutine hepeup_from_particle_set

\eatline
\nwendcode{}\nwbegindocs{496}Input.  The particle set should be allocated properly, but we
replace the particle contents.

If there are no beam particles in the event, we try to reconstruct beam
particles and beam remnants.  We assume for simplicity that the beam
particles, if any, are the first two particles.  If they are absent, the first
two particles should be the incoming partons.
\nwenddocs{}\nwbegincode{497}\moddef{HEP common: public}\plusendmoddef
  public :: hepeup_to_particle_set
\nwendcode{}\nwbegincode{498}\moddef{HEP common: procedures}\plusendmoddef
  subroutine hepeup_to_particle_set &
       (particle_set, recover_beams, model, alt_model)
    type(particle_set_t), intent(inout), target :: particle_set
    logical, intent(in), optional :: recover_beams
    class(model_data_t), intent(in), target :: model, alt_model
    type(particle_t), dimension(:), allocatable :: prt
    integer, dimension(2) :: parent
    integer, dimension(:), allocatable :: child
    integer :: i, j, k, pdg, status
    type(flavor_t) :: flv
    type(color_t) :: col
    integer, dimension(2) :: c
    type(vector4_t) :: p
    real(default) :: p2
    logical :: reconstruct
    integer :: off
    if (present (recover_beams)) then
       reconstruct = recover_beams .and. .not. all (ISTUP(1:2) == PRT_BEAM)
    else
       reconstruct = .false.
    end if
    if (reconstruct) then
       off = 4
    else
       off = 0
    end if
    allocate (prt (NUP + off), child (NUP + off))
    do i = 1, NUP
       k = i + off
       call hepeup_get_particle (i, pdg, status, col = c, p = p, m2 = p2)
       call flv%init (pdg, model, alt_model)
       call prt(k)%set_flavor (flv)
       call prt(k)%reset_status (status)
       call col%init (c)
       call prt(k)%set_color (col)
       call prt(k)%set_momentum (p, p2)
       where (MOTHUP(:,i) /= 0)
          parent = MOTHUP(:,i) + off
       elsewhere
          parent = 0
       end where
       call prt(k)%set_parents (parent)
       child = [(j, j = 1 + off, NUP + off)]
       where (MOTHUP(1,:NUP) /= i .and. MOTHUP(2,:NUP) /= i)  child = 0
       call prt(k)%set_children (child)
    end do
    if (reconstruct) then
       do k = 1, 2
          call prt(k)%reset_status (PRT_BEAM)
          call prt(k)%set_children ([k+2,k+4])
       end do
       do k = 3, 4
          call prt(k)%reset_status (PRT_BEAM_REMNANT)
          call prt(k)%set_parents ([k-2])
       end do
       do k = 5, 6
          call prt(k)%set_parents ([k-4])
       end do
    end if
    call particle_set%replace (prt)
  end subroutine hepeup_to_particle_set
  
\eatline
\nwendcode{}\nwbegindocs{499}\nwdocspar
The HEPEVT common block is quite similar, but does contain less
information, e.g. no color flows (it was LEP time). The spin
information is set only if the particle has a unique mother, and if
its polarization is fully defined. 
\nwenddocs{}\nwbegincode{500}\moddef{HEP common: public}\plusendmoddef
  public :: hepevt_from_particle_set
\nwendcode{}\nwbegincode{501}\moddef{HEP common: procedures}\plusendmoddef
  subroutine hepevt_from_particle_set &
       (particle_set, keep_beams, keep_remnants)
    type(particle_set_t), intent(in) :: particle_set
    type(particle_set_t), target :: pset_hepevt
    logical, intent(in), optional :: keep_beams
    logical, intent(in), optional :: keep_remnants
    integer :: i, status, n_tot
    logical :: activate_remnants
    activate_remnants = .true.
    if (present (keep_remnants))  activate_remnants = keep_remnants
    call particle_set%to_hepevt_form (pset_hepevt, keep_beams)
    n_tot = pset_hepevt%get_n_tot ()
    call hepevt_init (n_tot, pset_hepevt%get_n_out ())
    do i = 1, n_tot
       associate (prt => pset_hepevt%prt(i))
         status = prt%get_status ()
         if (activate_remnants &
              .and. status == PRT_BEAM_REMNANT &
              .and. prt%get_n_children () == 0) &
              status = PRT_OUTGOING
         call hepevt_set_particle (i, &
              prt%get_pdg (), status, &
              prt%get_parents (), &
              prt%get_children (), &
              prt%get_momentum (), &
              prt%get_p2 (), &
              prt%get_helicity ())
       end associate
    end do
    call pset_hepevt%final ()
  end subroutine hepevt_from_particle_set

\eatline
\nwendcode{}\nwbegindocs{502}\nwdocspar
%%%%%%%%%%%%%%%%%%%%%%%%%%%%%%%%%%%%%%%%%%%%%%%%%%%%%%%%%%%%%%%%%%%%%%%%
\section{HepMC events}
This section provides the interface to the HepMC C++ library for handling
Monte-Carlo events.

Each C++ class of HepMC that we use is mirrored by a Fortran type,
which contains as its only component the C pointer to the C++ object.

Each C++ method of HepMC that we use has a C wrapper function.  This
function takes a pointer to the host object as its first argument.
Further arguments are either C pointers, or in the case of simple
types (integer, real), interoperable C/Fortran objects.

The C wrapper functions have explicit interfaces in the Fortran
module.  They are called by Fortran wrapper procedures.  These are
treated as methods of the corresponding Fortran type.
\nwenddocs{}\nwbegincode{503}\moddef{\code{}hepmc{\char95}interface.f90\edoc{}}\endmoddef
\LA{}File header\RA{}

module hepmc_interface

  use, intrinsic :: iso_c_binding !NODEP!
  
\LA{}Use kinds\RA{}
\LA{}Use strings\RA{}
  use io_units
  use constants
  use lorentz
  use unit_tests
  use model_data
  use flavors
  use colors
  use helicities
  use quantum_numbers
  use polarizations

\LA{}Standard module head\RA{}

\LA{}HepMC interface: public\RA{}

\LA{}HepMC interface: types\RA{}

\LA{}HepMC interface: interfaces\RA{}

contains

\LA{}HepMC interface: procedures\RA{}

\LA{}HepMC interface: tests\RA{}

end module hepmc_interface
\eatline
\nwendcode{}\nwbegindocs{504}\nwdocspar
\subsection{Interface check}
This function can be called in order to verify that we are using the
actual HepMC library, and not the dummy version.
\nwenddocs{}\nwbegincode{505}\moddef{HepMC interface: interfaces}\endmoddef
  interface
     logical(c_bool) function hepmc_available () bind(C)
       import
     end function hepmc_available
  end interface
\nwendcode{}\nwbegincode{506}\moddef{HepMC interface: public}\endmoddef
  public :: hepmc_is_available
\nwendcode{}\nwbegincode{507}\moddef{HepMC interface: procedures}\endmoddef
  function hepmc_is_available () result (flag)
    logical :: flag
    flag = hepmc_available ()
  end function hepmc_is_available

\eatline
\nwendcode{}\nwbegindocs{508}\nwdocspar
\subsection{FourVector}
The C version of four-vectors is often transferred by value, and the
associated procedures are all inlined.  The wrapper needs to transfer
by reference, so we create FourVector objects on the heap which have
to be deleted explicitly.  The input is a {\tt{}vector4{\char95}t} or
{\tt{}vector3{\char95}t} object from the {\tt{}lorentz} module.
\nwenddocs{}\nwbegincode{509}\moddef{HepMC interface: public}\plusendmoddef
  public :: hepmc_four_vector_t
\nwendcode{}\nwbegincode{510}\moddef{HepMC interface: types}\endmoddef
  type :: hepmc_four_vector_t
     private
     type(c_ptr) :: obj
  end type hepmc_four_vector_t

\eatline
\nwendcode{}\nwbegindocs{511}In the C constructor, the zero-component (fourth argument) is
optional; if missing, it is set to zero.  The Fortran version has
initializer form and takes either a three-vector or a four-vector.
A further version extracts the four-vector from a HepMC particle
object.
\nwenddocs{}\nwbegincode{512}\moddef{HepMC interface: interfaces}\plusendmoddef
  interface
     type(c_ptr) function new_four_vector_xyz (x, y, z) bind(C)
       import
       real(c_double), value :: x, y, z
     end function new_four_vector_xyz
  end interface
  interface
     type(c_ptr) function new_four_vector_xyzt (x, y, z, t) bind(C)
       import
       real(c_double), value :: x, y, z, t
     end function new_four_vector_xyzt
  end interface
\eatline
\nwendcode{}\nwbegincode{513}\moddef{HepMC interface: public}\plusendmoddef
  public :: hepmc_four_vector_init
\nwendcode{}\nwbegincode{514}\moddef{HepMC interface: interfaces}\plusendmoddef
  interface hepmc_four_vector_init
     module procedure hepmc_four_vector_init_v4
     module procedure hepmc_four_vector_init_v3
     module procedure hepmc_four_vector_init_hepmc_prt
  end interface
\nwendcode{}\nwbegincode{515}\moddef{HepMC interface: procedures}\plusendmoddef
  subroutine hepmc_four_vector_init_v4 (pp, p)
    type(hepmc_four_vector_t), intent(out) :: pp
    type(vector4_t), intent(in) :: p
    real(default), dimension(0:3) :: pa
    pa = vector4_get_components (p)
    pp%obj = new_four_vector_xyzt &
         (real (pa(1), c_double), &
          real (pa(2), c_double), &
          real (pa(3), c_double), &
          real (pa(0), c_double))
  end subroutine hepmc_four_vector_init_v4
  
  subroutine hepmc_four_vector_init_v3 (pp, p)
    type(hepmc_four_vector_t), intent(out) :: pp
    type(vector3_t), intent(in) :: p
    real(default), dimension(3) :: pa
    pa = vector3_get_components (p)
    pp%obj = new_four_vector_xyz &
         (real (pa(1), c_double), &
          real (pa(2), c_double), &
          real (pa(3), c_double))
  end subroutine hepmc_four_vector_init_v3
  
  subroutine hepmc_four_vector_init_hepmc_prt (pp, prt)
    type(hepmc_four_vector_t), intent(out) :: pp
    type(hepmc_particle_t), intent(in) :: prt
    pp%obj = gen_particle_momentum (prt%obj)
  end subroutine hepmc_four_vector_init_hepmc_prt
  
\eatline
\nwendcode{}\nwbegindocs{516}Here, the destructor is explicitly needed.
\nwenddocs{}\nwbegincode{517}\moddef{HepMC interface: interfaces}\plusendmoddef
  interface
     subroutine four_vector_delete (p_obj) bind(C)
       import
       type(c_ptr), value :: p_obj
     end subroutine four_vector_delete
  end interface
\eatline
\nwendcode{}\nwbegincode{518}\moddef{HepMC interface: public}\plusendmoddef
  public :: hepmc_four_vector_final
\nwendcode{}\nwbegincode{519}\moddef{HepMC interface: procedures}\plusendmoddef
  subroutine hepmc_four_vector_final (p)
    type(hepmc_four_vector_t), intent(inout) :: p
    call four_vector_delete (p%obj)
  end subroutine hepmc_four_vector_final

\eatline
\nwendcode{}\nwbegindocs{520}Convert to a Lorentz vector.
\nwenddocs{}\nwbegincode{521}\moddef{HepMC interface: interfaces}\plusendmoddef
  interface
     function four_vector_px (p_obj) result (px) bind(C)
       import
       real(c_double) :: px
       type(c_ptr), value :: p_obj
     end function four_vector_px
  end interface
  interface
     function four_vector_py (p_obj) result (py) bind(C)
       import
       real(c_double) :: py
       type(c_ptr), value :: p_obj
     end function four_vector_py
  end interface
  interface
     function four_vector_pz (p_obj) result (pz) bind(C)
       import
       real(c_double) :: pz
       type(c_ptr), value :: p_obj
     end function four_vector_pz
  end interface
  interface
     function four_vector_e (p_obj) result (e) bind(C)
       import
       real(c_double) :: e
       type(c_ptr), value :: p_obj
     end function four_vector_e
  end interface
\eatline
\nwendcode{}\nwbegincode{522}\moddef{HepMC interface: public}\plusendmoddef
  public :: hepmc_four_vector_to_vector4
\nwendcode{}\nwbegincode{523}\moddef{HepMC interface: procedures}\plusendmoddef
  subroutine hepmc_four_vector_to_vector4 (pp, p)
    type(hepmc_four_vector_t), intent(in) :: pp
    type(vector4_t), intent(out) :: p
    real(default) :: E
    real(default), dimension(3) :: p3
    E = four_vector_e (pp%obj)
    p3(1) = four_vector_px (pp%obj)
    p3(2) = four_vector_py (pp%obj)
    p3(3) = four_vector_pz (pp%obj)
    p = vector4_moving (E, vector3_moving (p3))
  end subroutine hepmc_four_vector_to_vector4

\eatline
\nwendcode{}\nwbegindocs{524}\nwdocspar
\subsection{Polarization}
Polarization objects are temporarily used for assigning particle
polarization.  We add a flag {\tt{}polarized}.  If this is false, the
polarization is not set and should not be transferred to
{\tt{}hepmc{\char95}particle} objects.
\nwenddocs{}\nwbegincode{525}\moddef{HepMC interface: public}\plusendmoddef
  public :: hepmc_polarization_t
\nwendcode{}\nwbegincode{526}\moddef{HepMC interface: types}\plusendmoddef
  type :: hepmc_polarization_t
     private
     logical :: polarized = .false.
     type(c_ptr) :: obj
  end type hepmc_polarization_t

\eatline
\nwendcode{}\nwbegindocs{527}Constructor.  The C wrapper takes polar and azimuthal angle as
arguments.  The Fortran version allows for either a complete
polarization density matrix, or for a definite (diagonal) helicity.

\emph{HepMC does not allow to specify the degree of polarization,
  therefore we have to map it to either 0 or 1.  We choose 0 for
  polarization less than $0.5$ and 1 for polarization greater than
  $0.5$.  Even this simplification works only for spin-1/2 and for
  massless particles; massive vector bosons cannot be treated this
  way.  In particular, zero helicity is always translated as
  unpolarized.}

\emph{For massive vector bosons, we arbitrarily choose the convention
  that the longitudinal (zero) helicity state is mapped to the theta
  angle $\pi/2$.  This works under the condition that helicity is
  projected onto one of the basis states.}
\nwenddocs{}\nwbegincode{528}\moddef{HepMC interface: interfaces}\plusendmoddef
  interface
     type(c_ptr) function new_polarization (theta, phi) bind(C)
       import
       real(c_double), value :: theta, phi
     end function new_polarization
  end interface
\eatline
\nwendcode{}\nwbegincode{529}\moddef{HepMC interface: public}\plusendmoddef
  public :: hepmc_polarization_init
\nwendcode{}\nwbegincode{530}\moddef{HepMC interface: interfaces}\plusendmoddef
  interface hepmc_polarization_init
     module procedure hepmc_polarization_init_pol
     module procedure hepmc_polarization_init_hel
     module procedure hepmc_polarization_init_int
  end interface
\nwendcode{}\nwbegincode{531}\moddef{HepMC interface: procedures}\plusendmoddef
  subroutine hepmc_polarization_init_pol (hpol, pol)
    type(hepmc_polarization_t), intent(out) :: hpol
    type(polarization_t), intent(in) :: pol
    real(default) :: r, theta, phi
    if (polarization_is_polarized (pol)) then
       call polarization_to_angles (pol, r, theta, phi)
       if (r >= 0.5) then
          hpol%polarized = .true.
          hpol%obj = new_polarization &
               (real (theta, c_double), real (phi, c_double))
       end if
    end if
  end subroutine hepmc_polarization_init_pol

  subroutine hepmc_polarization_init_hel (hpol, hel)
    type(hepmc_polarization_t), intent(out) :: hpol
    type(helicity_t), intent(in) :: hel
    integer, dimension(2) :: h
    if (hel%is_defined ()) then
       h = hel%to_pair ()
       select case (h(1))
       case (1:)
          hpol%polarized = .true.
          hpol%obj = new_polarization (0._c_double, 0._c_double)
       case (:-1)
          hpol%polarized = .true.
          hpol%obj = new_polarization (real (pi, c_double), 0._c_double)
       case (0)
          hpol%polarized = .true.
          hpol%obj = new_polarization (real (pi/2, c_double), 0._c_double)
       end select
    end if
  end subroutine hepmc_polarization_init_hel

  subroutine hepmc_polarization_init_int (hpol, hel)
    type(hepmc_polarization_t), intent(out) :: hpol
    integer, intent(in) :: hel
    select case (hel)
    case (1:)
       hpol%polarized = .true.
       hpol%obj = new_polarization (0._c_double, 0._c_double)
    case (:-1)
       hpol%polarized = .true.
       hpol%obj = new_polarization (real (pi, c_double), 0._c_double)
    case (0)
       hpol%polarized = .true.
       hpol%obj = new_polarization (real (pi/2, c_double), 0._c_double)
    end select
  end subroutine hepmc_polarization_init_int

\eatline
\nwendcode{}\nwbegindocs{532}Destructor.  The C object is deallocated only if the {\tt{}polarized}
flag is set.
\nwenddocs{}\nwbegincode{533}\moddef{HepMC interface: interfaces}\plusendmoddef
  interface
     subroutine polarization_delete (pol_obj) bind(C)
       import
       type(c_ptr), value :: pol_obj
     end subroutine polarization_delete
  end interface
\eatline
\nwendcode{}\nwbegincode{534}\moddef{HepMC interface: public}\plusendmoddef
  public :: hepmc_polarization_final
\nwendcode{}\nwbegincode{535}\moddef{HepMC interface: procedures}\plusendmoddef
  subroutine hepmc_polarization_final (hpol)
    type(hepmc_polarization_t), intent(inout) :: hpol
    if (hpol%polarized)  call polarization_delete (hpol%obj)
  end subroutine hepmc_polarization_final

\eatline
\nwendcode{}\nwbegindocs{536}Recover polarization from HepMC polarization object (with the
abovementioned deficiencies).
\nwenddocs{}\nwbegincode{537}\moddef{HepMC interface: interfaces}\plusendmoddef
  interface
     function polarization_theta (pol_obj) result (theta) bind(C)
       import
       real(c_double) :: theta
       type(c_ptr), value :: pol_obj
     end function polarization_theta
  end interface
  interface
     function polarization_phi (pol_obj) result (phi) bind(C)
       import
       real(c_double) :: phi
       type(c_ptr), value :: pol_obj
     end function polarization_phi
  end interface
\eatline
\nwendcode{}\nwbegincode{538}\moddef{HepMC interface: public}\plusendmoddef
  public :: hepmc_polarization_to_pol
\nwendcode{}\nwbegincode{539}\moddef{HepMC interface: procedures}\plusendmoddef
  subroutine hepmc_polarization_to_pol (hpol, flv, pol)
    type(hepmc_polarization_t), intent(in) :: hpol
    type(flavor_t), intent(in) :: flv
    type(polarization_t), intent(out) :: pol
    real(default) :: theta, phi
    theta = polarization_theta (hpol%obj)
    phi = polarization_phi (hpol%obj)
    call polarization_init_angles (pol, flv, 1._default, theta, phi)
  end subroutine hepmc_polarization_to_pol

\eatline
\nwendcode{}\nwbegindocs{540}Recover helicity.  Here, $\phi$ is ignored and only the sign of
$\cos\theta$ is relevant, mapped to positive/negative helicity.
\nwenddocs{}\nwbegincode{541}\moddef{HepMC interface: public}\plusendmoddef
  public :: hepmc_polarization_to_hel
\nwendcode{}\nwbegincode{542}\moddef{HepMC interface: procedures}\plusendmoddef
  subroutine hepmc_polarization_to_hel (hpol, flv, hel)
    type(hepmc_polarization_t), intent(in) :: hpol
    type(flavor_t), intent(in) :: flv
    type(helicity_t), intent(out) :: hel
    real(default) :: theta
    integer :: hmax
    theta = polarization_theta (hpol%obj)
    hmax = flv%get_spin_type () / 2
    call hel%init (sign (hmax, nint (cos (theta))))
  end subroutine hepmc_polarization_to_hel

\eatline
\nwendcode{}\nwbegindocs{543}\nwdocspar
\subsection{GenParticle}
Particle objects have the obvious meaning.
\nwenddocs{}\nwbegincode{544}\moddef{HepMC interface: public}\plusendmoddef
  public :: hepmc_particle_t
\nwendcode{}\nwbegincode{545}\moddef{HepMC interface: types}\plusendmoddef
  type :: hepmc_particle_t
     private
     type(c_ptr) :: obj
  end type hepmc_particle_t

\eatline
\nwendcode{}\nwbegindocs{546}Constructor.  The C version takes a FourVector object, which in the
Fortran wrapper is created on the fly from a {\tt{}vector4} Lorentz
vector.

No destructor is needed as long as all particles are entered into
vertex containers.
\nwenddocs{}\nwbegincode{547}\moddef{HepMC interface: interfaces}\plusendmoddef
  interface
     type(c_ptr) function new_gen_particle (prt_obj, pdg_id, status) bind(C)
       import
       type(c_ptr), value :: prt_obj
       integer(c_int), value :: pdg_id, status
     end function new_gen_particle
  end interface
\eatline
\nwendcode{}\nwbegincode{548}\moddef{HepMC interface: public}\plusendmoddef
  public :: hepmc_particle_init
\nwendcode{}\nwbegincode{549}\moddef{HepMC interface: procedures}\plusendmoddef
  subroutine hepmc_particle_init (prt, p, pdg, status)
    type(hepmc_particle_t), intent(out) :: prt
    type(vector4_t), intent(in) :: p
    integer, intent(in) :: pdg, status
    type(hepmc_four_vector_t) :: pp
    call hepmc_four_vector_init (pp, p)
    prt%obj = new_gen_particle (pp%obj, int (pdg, c_int), int (status, c_int))
    call hepmc_four_vector_final (pp)
  end subroutine hepmc_particle_init

\eatline
\nwendcode{}\nwbegindocs{550}Set the particle color flow.
\nwenddocs{}\nwbegincode{551}\moddef{HepMC interface: interfaces}\plusendmoddef
  interface
     subroutine gen_particle_set_flow (prt_obj, code_index, code) bind(C)
       import
       type(c_ptr), value :: prt_obj
       integer(c_int), value :: code_index, code
     end subroutine gen_particle_set_flow
  end interface
\eatline
\nwendcode{}\nwbegindocs{552}Set the particle color.  Either from a {\tt{}color{\char95}t} object or
directly from a pair of integers.
\nwenddocs{}\nwbegincode{553}\moddef{HepMC interface: interfaces}\plusendmoddef
  interface hepmc_particle_set_color
     module procedure hepmc_particle_set_color_col
     module procedure hepmc_particle_set_color_int
  end interface hepmc_particle_set_color
\nwendcode{}\nwbegincode{554}\moddef{HepMC interface: public}\plusendmoddef
  public :: hepmc_particle_set_color
\nwendcode{}\nwbegincode{555}\moddef{HepMC interface: procedures}\plusendmoddef
  subroutine hepmc_particle_set_color_col (prt, col)
    type(hepmc_particle_t), intent(inout) :: prt
    type(color_t), intent(in) :: col
    integer(c_int) :: c
    c = col%get_col ()
    if (c /= 0)  call gen_particle_set_flow (prt%obj, 1_c_int, c)
    c = col%get_acl ()
    if (c /= 0)  call gen_particle_set_flow (prt%obj, 2_c_int, c)
  end subroutine hepmc_particle_set_color_col

  subroutine hepmc_particle_set_color_int (prt, col)
    type(hepmc_particle_t), intent(inout) :: prt
    integer, dimension(2), intent(in) :: col
    integer(c_int) :: c
    c = col(1)
    if (c /= 0)  call gen_particle_set_flow (prt%obj, 1_c_int, c)
    c = col(2)
    if (c /= 0)  call gen_particle_set_flow (prt%obj, 2_c_int, c)
  end subroutine hepmc_particle_set_color_int

\eatline
\nwendcode{}\nwbegindocs{556}Set the particle polarization.  For the restrictions on particle
polarization in HepMC, see above {\tt{}hepmc{\char95}polarization{\char95}init}.
\nwenddocs{}\nwbegincode{557}\moddef{HepMC interface: interfaces}\plusendmoddef
  interface
     subroutine gen_particle_set_polarization (prt_obj, pol_obj) bind(C)
       import
       type(c_ptr), value :: prt_obj, pol_obj
     end subroutine gen_particle_set_polarization
  end interface
\eatline
\nwendcode{}\nwbegincode{558}\moddef{HepMC interface: public}\plusendmoddef
  public :: hepmc_particle_set_polarization
\nwendcode{}\nwbegincode{559}\moddef{HepMC interface: interfaces}\plusendmoddef
  interface hepmc_particle_set_polarization
     module procedure hepmc_particle_set_polarization_pol
     module procedure hepmc_particle_set_polarization_hel
     module procedure hepmc_particle_set_polarization_int
  end interface
\nwendcode{}\nwbegincode{560}\moddef{HepMC interface: procedures}\plusendmoddef
  subroutine hepmc_particle_set_polarization_pol (prt, pol)
    type(hepmc_particle_t), intent(inout) :: prt
    type(polarization_t), intent(in) :: pol
    type(hepmc_polarization_t) :: hpol
    call hepmc_polarization_init (hpol, pol)
    if (hpol%polarized)  call gen_particle_set_polarization (prt%obj, hpol%obj)
    call hepmc_polarization_final (hpol)
  end subroutine hepmc_particle_set_polarization_pol

  subroutine hepmc_particle_set_polarization_hel (prt, hel)
    type(hepmc_particle_t), intent(inout) :: prt
    type(helicity_t), intent(in) :: hel
    type(hepmc_polarization_t) :: hpol
    call hepmc_polarization_init (hpol, hel)
    if (hpol%polarized)  call gen_particle_set_polarization (prt%obj, hpol%obj)
    call hepmc_polarization_final (hpol)
  end subroutine hepmc_particle_set_polarization_hel

  subroutine hepmc_particle_set_polarization_int (prt, hel)
    type(hepmc_particle_t), intent(inout) :: prt
    integer, intent(in) :: hel
    type(hepmc_polarization_t) :: hpol
    call hepmc_polarization_init (hpol, hel)
    if (hpol%polarized)  call gen_particle_set_polarization (prt%obj, hpol%obj)
    call hepmc_polarization_final (hpol)
  end subroutine hepmc_particle_set_polarization_int

\eatline
\nwendcode{}\nwbegindocs{561}Return the HepMC barcode (unique integer ID) of the particle.
\nwenddocs{}\nwbegincode{562}\moddef{HepMC interface: interfaces}\plusendmoddef
  interface
     function gen_particle_barcode (prt_obj) result (barcode) bind(C)
       import
       integer(c_int) :: barcode
       type(c_ptr), value :: prt_obj
     end function gen_particle_barcode
  end interface
\eatline
\nwendcode{}\nwbegincode{563}\moddef{HepMC interface: public}\plusendmoddef
  public :: hepmc_particle_get_barcode
\nwendcode{}\nwbegincode{564}\moddef{HepMC interface: procedures}\plusendmoddef
  function hepmc_particle_get_barcode (prt) result (barcode)
    integer :: barcode
    type(hepmc_particle_t), intent(in) :: prt
    barcode = gen_particle_barcode (prt%obj)
  end function hepmc_particle_get_barcode

\eatline
\nwendcode{}\nwbegindocs{565}Return the four-vector component of the particle object as a {\tt{}vector4{\char95}t} Lorentz vector.
\nwenddocs{}\nwbegincode{566}\moddef{HepMC interface: interfaces}\plusendmoddef
  interface
     type(c_ptr) function gen_particle_momentum (prt_obj) bind(C)
       import
       type(c_ptr), value :: prt_obj
     end function gen_particle_momentum
  end interface
\eatline
\nwendcode{}\nwbegincode{567}\moddef{HepMC interface: public}\plusendmoddef
  public :: hepmc_particle_get_momentum
\nwendcode{}\nwbegincode{568}\moddef{HepMC interface: procedures}\plusendmoddef
  function hepmc_particle_get_momentum (prt) result (p)
    type(vector4_t) :: p
    type(hepmc_particle_t), intent(in) :: prt
    type(hepmc_four_vector_t) :: pp
    call hepmc_four_vector_init (pp, prt)
    call hepmc_four_vector_to_vector4 (pp, p)
    call hepmc_four_vector_final (pp)
  end function hepmc_particle_get_momentum

\eatline
\nwendcode{}\nwbegindocs{569}Return the invariant mass squared of the particle object.  HepMC
stores the signed invariant mass (no squaring).
\nwenddocs{}\nwbegincode{570}\moddef{HepMC interface: interfaces}\plusendmoddef
  interface
     function gen_particle_generated_mass (prt_obj) result (mass) bind(C)
       import
       real(c_double) :: mass
       type(c_ptr), value :: prt_obj
     end function gen_particle_generated_mass
  end interface
\eatline
\nwendcode{}\nwbegincode{571}\moddef{HepMC interface: public}\plusendmoddef
  public :: hepmc_particle_get_mass_squared
\nwendcode{}\nwbegincode{572}\moddef{HepMC interface: procedures}\plusendmoddef
  function hepmc_particle_get_mass_squared (prt) result (m2)
    real(default) :: m2
    type(hepmc_particle_t), intent(in) :: prt
    real(default) :: m
    m = gen_particle_generated_mass (prt%obj)
    m2 = sign (m**2, m)
  end function hepmc_particle_get_mass_squared

\eatline
\nwendcode{}\nwbegindocs{573}Return the PDG ID:
\nwenddocs{}\nwbegincode{574}\moddef{HepMC interface: interfaces}\plusendmoddef
  interface
     function gen_particle_pdg_id (prt_obj) result (pdg_id) bind(C)
       import
       integer(c_int) :: pdg_id
       type(c_ptr), value :: prt_obj
     end function gen_particle_pdg_id
  end interface
\eatline
\nwendcode{}\nwbegincode{575}\moddef{HepMC interface: public}\plusendmoddef
  public :: hepmc_particle_get_pdg
\nwendcode{}\nwbegincode{576}\moddef{HepMC interface: procedures}\plusendmoddef
  function hepmc_particle_get_pdg (prt) result (pdg)
    integer :: pdg
    type(hepmc_particle_t), intent(in) :: prt
    pdg = gen_particle_pdg_id (prt%obj)
  end function hepmc_particle_get_pdg

\eatline
\nwendcode{}\nwbegindocs{577}Return the status code:
\nwenddocs{}\nwbegincode{578}\moddef{HepMC interface: interfaces}\plusendmoddef
  interface
     function gen_particle_status (prt_obj) result (status) bind(C)
       import
       integer(c_int) :: status
       type(c_ptr), value :: prt_obj
     end function gen_particle_status
  end interface
\eatline
\nwendcode{}\nwbegincode{579}\moddef{HepMC interface: public}\plusendmoddef
  public :: hepmc_particle_get_status
\nwendcode{}\nwbegincode{580}\moddef{HepMC interface: procedures}\plusendmoddef
  function hepmc_particle_get_status (prt) result (status)
    integer :: status
    type(hepmc_particle_t), intent(in) :: prt
    status = gen_particle_status (prt%obj)
  end function hepmc_particle_get_status

\eatline
\nwendcode{}\nwbegincode{581}\moddef{HepMC interface: interfaces}\plusendmoddef
  interface
     function gen_particle_is_beam (prt_obj) result (is_beam) bind(C)
       import
       logical(c_bool) :: is_beam
       type(c_ptr), value :: prt_obj
     end function gen_particle_is_beam
  end interface
\eatline
\nwendcode{}\nwbegindocs{582}Determine whether a particle is a beam particle. 
\nwenddocs{}\nwbegincode{583}\moddef{HepMC interface: public}\plusendmoddef
  public :: hepmc_particle_is_beam
\nwendcode{}\nwbegincode{584}\moddef{HepMC interface: procedures}\plusendmoddef
  function hepmc_particle_is_beam (prt) result (is_beam)
    logical :: is_beam
    type(hepmc_particle_t), intent(in) :: prt
    is_beam = gen_particle_is_beam (prt%obj)
  end function hepmc_particle_is_beam
    
\eatline
\nwendcode{}\nwbegindocs{585}Return the production/decay vertex (as a pointer, no finalization
necessary).
\nwenddocs{}\nwbegincode{586}\moddef{HepMC interface: interfaces}\plusendmoddef
  interface
     type(c_ptr) function gen_particle_production_vertex (prt_obj) bind(C)
       import
       type(c_ptr), value :: prt_obj
     end function gen_particle_production_vertex
  end interface
  interface
     type(c_ptr) function gen_particle_end_vertex (prt_obj) bind(C)
       import
       type(c_ptr), value :: prt_obj
     end function gen_particle_end_vertex
  end interface
\eatline
\nwendcode{}\nwbegincode{587}\moddef{HepMC interface: public}\plusendmoddef
  public :: hepmc_particle_get_production_vertex
  public :: hepmc_particle_get_decay_vertex
\nwendcode{}\nwbegincode{588}\moddef{HepMC interface: procedures}\plusendmoddef
  function hepmc_particle_get_production_vertex (prt) result (v)
    type(hepmc_vertex_t) :: v
    type(hepmc_particle_t), intent(in) :: prt
    v%obj = gen_particle_production_vertex (prt%obj)
  end function hepmc_particle_get_production_vertex

  function hepmc_particle_get_decay_vertex (prt) result (v)
    type(hepmc_vertex_t) :: v
    type(hepmc_particle_t), intent(in) :: prt
    v%obj = gen_particle_end_vertex (prt%obj)
  end function hepmc_particle_get_decay_vertex

\eatline
\nwendcode{}\nwbegindocs{589}Return the number of parents/children.
\nwenddocs{}\nwbegincode{590}\moddef{HepMC interface: public}\plusendmoddef
  public :: hepmc_particle_get_n_parents
  public :: hepmc_particle_get_n_children
\nwendcode{}\nwbegincode{591}\moddef{HepMC interface: procedures}\plusendmoddef
  function hepmc_particle_get_n_parents (prt) result (n_parents)
    integer :: n_parents
    type(hepmc_particle_t), intent(in) :: prt
    type(hepmc_vertex_t) :: v
    v = hepmc_particle_get_production_vertex (prt)
    if (hepmc_vertex_is_valid (v)) then
       n_parents = hepmc_vertex_get_n_in (v)
    else
       n_parents = 0
    end if
  end function hepmc_particle_get_n_parents

  function hepmc_particle_get_n_children (prt) result (n_children)
    integer :: n_children
    type(hepmc_particle_t), intent(in) :: prt
    type(hepmc_vertex_t) :: v
    v = hepmc_particle_get_decay_vertex (prt)
    if (hepmc_vertex_is_valid (v)) then
       n_children = hepmc_vertex_get_n_out (v)
    else
       n_children = 0
    end if
  end function hepmc_particle_get_n_children

\eatline
\eatline
\nwendcode{}\nwbegindocs{592}Convenience function: Return the array of parent particles for a
given HepMC particle.  The contents are HepMC barcodes that still have
to be mapped to the particle indices.
\nwenddocs{}\nwbegincode{593}\moddef{HepMC interface: public}\plusendmoddef
  public :: hepmc_particle_get_parent_barcodes
  public :: hepmc_particle_get_child_barcodes
\nwendcode{}\nwbegincode{594}\moddef{HepMC interface: procedures}\plusendmoddef
  function hepmc_particle_get_parent_barcodes (prt) result (parent_barcode)
    type(hepmc_particle_t), intent(in) :: prt
    integer, dimension(:), allocatable :: parent_barcode
    type(hepmc_vertex_t) :: v
    type(hepmc_vertex_particle_in_iterator_t) :: it
    integer :: i
    v = hepmc_particle_get_production_vertex (prt)
    if (hepmc_vertex_is_valid (v)) then
       allocate (parent_barcode (hepmc_vertex_get_n_in (v)))
       if (size (parent_barcode) /= 0) then
          call hepmc_vertex_particle_in_iterator_init (it, v)
          do i = 1, size (parent_barcode)
             parent_barcode(i) = hepmc_particle_get_barcode &
                  (hepmc_vertex_particle_in_iterator_get (it))
             call hepmc_vertex_particle_in_iterator_advance (it)
          end do
          call hepmc_vertex_particle_in_iterator_final (it)
       end if
    else
       allocate (parent_barcode (0))
    end if
  end function hepmc_particle_get_parent_barcodes

  function hepmc_particle_get_child_barcodes (prt) result (child_barcode)
    type(hepmc_particle_t), intent(in) :: prt
    integer, dimension(:), allocatable :: child_barcode
    type(hepmc_vertex_t) :: v
    type(hepmc_vertex_particle_out_iterator_t) :: it
    integer :: i
    v = hepmc_particle_get_decay_vertex (prt)
    if (hepmc_vertex_is_valid (v)) then
       allocate (child_barcode (hepmc_vertex_get_n_out (v)))
       call hepmc_vertex_particle_out_iterator_init (it, v)
       if (size (child_barcode) /= 0) then
          do i = 1, size (child_barcode)
             child_barcode(i) = hepmc_particle_get_barcode &
                  (hepmc_vertex_particle_out_iterator_get (it))
             call hepmc_vertex_particle_out_iterator_advance (it)
          end do
          call hepmc_vertex_particle_out_iterator_final (it)
       end if
    else
       allocate (child_barcode (0))
    end if
  end function hepmc_particle_get_child_barcodes

\eatline
\nwendcode{}\nwbegindocs{595}Return the polarization (assuming that the particle is completely
polarized).  Note that the generated polarization object needs finalization.
\nwenddocs{}\nwbegincode{596}\moddef{HepMC interface: interfaces}\plusendmoddef
  interface
     type(c_ptr) function gen_particle_polarization (prt_obj) bind(C)
       import
       type(c_ptr), value :: prt_obj
     end function gen_particle_polarization
  end interface
\eatline
\nwendcode{}\nwbegincode{597}\moddef{HepMC interface: public}\plusendmoddef
  public :: hepmc_particle_get_polarization
\nwendcode{}\nwbegincode{598}\moddef{HepMC interface: procedures}\plusendmoddef
  function hepmc_particle_get_polarization (prt) result (pol)
    type(hepmc_polarization_t) :: pol
    type(hepmc_particle_t), intent(in) :: prt
    pol%obj = gen_particle_polarization (prt%obj)
  end function hepmc_particle_get_polarization

\eatline
\nwendcode{}\nwbegindocs{599}Return the particle color as a two-dimensional array (color, anticolor).
\nwenddocs{}\nwbegincode{600}\moddef{HepMC interface: interfaces}\plusendmoddef
  interface
     function gen_particle_flow (prt_obj, code_index) result (code) bind(C)
       import
       integer(c_int) :: code
       type(c_ptr), value :: prt_obj
       integer(c_int), value :: code_index
     end function gen_particle_flow
  end interface
\eatline
\nwendcode{}\nwbegincode{601}\moddef{HepMC interface: public}\plusendmoddef
  public :: hepmc_particle_get_color
\nwendcode{}\nwbegincode{602}\moddef{HepMC interface: procedures}\plusendmoddef
  function hepmc_particle_get_color (prt) result (col)
    integer, dimension(2) :: col
    type(hepmc_particle_t), intent(in) :: prt
    col(1) = gen_particle_flow (prt%obj, 1)
    col(2) = - gen_particle_flow (prt%obj, 2)
  end function hepmc_particle_get_color

\eatline
\nwendcode{}\nwbegindocs{603}\nwdocspar
\subsection{GenVertex}
Vertices are made of particles (incoming and outgoing).
\nwenddocs{}\nwbegincode{604}\moddef{HepMC interface: public}\plusendmoddef
  public :: hepmc_vertex_t
\nwendcode{}\nwbegincode{605}\moddef{HepMC interface: types}\plusendmoddef
  type :: hepmc_vertex_t
     private
     type(c_ptr) :: obj
  end type hepmc_vertex_t

\eatline
\nwendcode{}\nwbegindocs{606}Constructor.  Two versions, one plain, one with the position in
space and time (measured in mm) as argument.  The Fortran version has
initializer form, and the vertex position is an optional argument.

A destructor is unnecessary as long as all vertices are entered into
an event container.
\nwenddocs{}\nwbegincode{607}\moddef{HepMC interface: interfaces}\plusendmoddef
  interface
     type(c_ptr) function new_gen_vertex () bind(C)
       import
     end function new_gen_vertex
  end interface
  interface
     type(c_ptr) function new_gen_vertex_pos (prt_obj) bind(C)
       import
       type(c_ptr), value :: prt_obj
     end function new_gen_vertex_pos
  end interface
\eatline
\nwendcode{}\nwbegincode{608}\moddef{HepMC interface: public}\plusendmoddef
  public :: hepmc_vertex_init
\nwendcode{}\nwbegincode{609}\moddef{HepMC interface: procedures}\plusendmoddef
  subroutine hepmc_vertex_init (v, x)
    type(hepmc_vertex_t), intent(out) :: v
    type(vector4_t), intent(in), optional :: x
    type(hepmc_four_vector_t) :: pos
    if (present (x)) then
       call hepmc_four_vector_init (pos, x)
       v%obj = new_gen_vertex_pos (pos%obj)
       call hepmc_four_vector_final (pos)
    else
       v%obj = new_gen_vertex ()
    end if
  end subroutine hepmc_vertex_init

\eatline
\nwendcode{}\nwbegindocs{610}Return true if the vertex pointer is non-null:
\nwenddocs{}\nwbegincode{611}\moddef{HepMC interface: interfaces}\plusendmoddef
  interface
     function gen_vertex_is_valid (v_obj) result (flag) bind(C)
       import
       logical(c_bool) :: flag
       type(c_ptr), value :: v_obj
     end function gen_vertex_is_valid
  end interface
\eatline
\nwendcode{}\nwbegincode{612}\moddef{HepMC interface: public}\plusendmoddef
  public :: hepmc_vertex_is_valid
\nwendcode{}\nwbegincode{613}\moddef{HepMC interface: procedures}\plusendmoddef
  function hepmc_vertex_is_valid (v) result (flag)
    logical :: flag
    type(hepmc_vertex_t), intent(in) :: v
    flag = gen_vertex_is_valid (v%obj)
  end function hepmc_vertex_is_valid

\eatline
\nwendcode{}\nwbegindocs{614}Add a particle to a vertex, incoming or outgoing.
\nwenddocs{}\nwbegincode{615}\moddef{HepMC interface: interfaces}\plusendmoddef
  interface
     subroutine gen_vertex_add_particle_in (v_obj, prt_obj) bind(C)
       import
       type(c_ptr), value :: v_obj, prt_obj
     end subroutine gen_vertex_add_particle_in
  end interface
  interface
     subroutine gen_vertex_add_particle_out (v_obj, prt_obj) bind(C)
       import
       type(c_ptr), value :: v_obj, prt_obj
     end subroutine gen_vertex_add_particle_out
  end interface
\nwendcode{}\nwbegincode{616}\moddef{HepMC interface: public}\plusendmoddef
  public :: hepmc_vertex_add_particle_in
  public :: hepmc_vertex_add_particle_out
\eatline
\nwendcode{}\nwbegincode{617}\moddef{HepMC interface: procedures}\plusendmoddef
  subroutine hepmc_vertex_add_particle_in (v, prt)
    type(hepmc_vertex_t), intent(inout) :: v
    type(hepmc_particle_t), intent(in) :: prt
    call gen_vertex_add_particle_in (v%obj, prt%obj)
  end subroutine hepmc_vertex_add_particle_in

  subroutine hepmc_vertex_add_particle_out (v, prt)
    type(hepmc_vertex_t), intent(inout) :: v
    type(hepmc_particle_t), intent(in) :: prt
    call gen_vertex_add_particle_out (v%obj, prt%obj)
  end subroutine hepmc_vertex_add_particle_out

\eatline
\nwendcode{}\nwbegindocs{618}Return the number of incoming/outgoing particles.
\nwenddocs{}\nwbegincode{619}\moddef{HepMC interface: interfaces}\plusendmoddef
  interface
     function gen_vertex_particles_in_size (v_obj) result (size) bind(C)
       import
       integer(c_int) :: size
       type(c_ptr), value :: v_obj
     end function gen_vertex_particles_in_size
  end interface
  interface
     function gen_vertex_particles_out_size (v_obj) result (size) bind(C)
       import
       integer(c_int) :: size
       type(c_ptr), value :: v_obj
     end function gen_vertex_particles_out_size
  end interface
\eatline
\nwendcode{}\nwbegincode{620}\moddef{HepMC interface: public}\plusendmoddef
  public :: hepmc_vertex_get_n_in
  public :: hepmc_vertex_get_n_out
\nwendcode{}\nwbegincode{621}\moddef{HepMC interface: procedures}\plusendmoddef
  function hepmc_vertex_get_n_in (v) result (n_in)
    integer :: n_in
    type(hepmc_vertex_t), intent(in) :: v
    n_in = gen_vertex_particles_in_size (v%obj)
  end function hepmc_vertex_get_n_in

  function hepmc_vertex_get_n_out (v) result (n_out)
    integer :: n_out
    type(hepmc_vertex_t), intent(in) :: v
    n_out = gen_vertex_particles_out_size (v%obj)
  end function hepmc_vertex_get_n_out

\eatline
\nwendcode{}\nwbegindocs{622}\nwdocspar
\subsection{Vertex-particle-in iterator}
This iterator iterates over all incoming particles in an vertex.  We store a
pointer to the vertex in addition to the iterator.  This allows for
simple end checking.

The iterator is actually a constant iterator; it can only read.
\nwenddocs{}\nwbegincode{623}\moddef{HepMC interface: public}\plusendmoddef
  public :: hepmc_vertex_particle_in_iterator_t
\nwendcode{}\nwbegincode{624}\moddef{HepMC interface: types}\plusendmoddef
  type :: hepmc_vertex_particle_in_iterator_t
     private
     type(c_ptr) :: obj
     type(c_ptr) :: v_obj
  end type hepmc_vertex_particle_in_iterator_t

\eatline
\nwendcode{}\nwbegindocs{625}Constructor.  The iterator is initialized at the first particle in
the vertex.
\nwenddocs{}\nwbegincode{626}\moddef{HepMC interface: interfaces}\plusendmoddef
  interface
     type(c_ptr) function &
          new_vertex_particles_in_const_iterator (v_obj) bind(C)
       import
       type(c_ptr), value :: v_obj
     end function new_vertex_particles_in_const_iterator
  end interface
\eatline
\nwendcode{}\nwbegincode{627}\moddef{HepMC interface: public}\plusendmoddef
  public :: hepmc_vertex_particle_in_iterator_init
\nwendcode{}\nwbegincode{628}\moddef{HepMC interface: procedures}\plusendmoddef
  subroutine hepmc_vertex_particle_in_iterator_init (it, v)
    type(hepmc_vertex_particle_in_iterator_t), intent(out) :: it
    type(hepmc_vertex_t), intent(in) :: v
    it%obj = new_vertex_particles_in_const_iterator (v%obj)
    it%v_obj = v%obj
  end subroutine hepmc_vertex_particle_in_iterator_init

\eatline
\nwendcode{}\nwbegindocs{629}Destructor.  Necessary because the iterator is allocated on the
heap.
\nwenddocs{}\nwbegincode{630}\moddef{HepMC interface: interfaces}\plusendmoddef
  interface
     subroutine vertex_particles_in_const_iterator_delete (it_obj) bind(C)
       import
       type(c_ptr), value :: it_obj
     end subroutine vertex_particles_in_const_iterator_delete
  end interface
\eatline
\nwendcode{}\nwbegincode{631}\moddef{HepMC interface: public}\plusendmoddef
  public :: hepmc_vertex_particle_in_iterator_final
\nwendcode{}\nwbegincode{632}\moddef{HepMC interface: procedures}\plusendmoddef
  subroutine hepmc_vertex_particle_in_iterator_final (it)
    type(hepmc_vertex_particle_in_iterator_t), intent(inout) :: it
    call vertex_particles_in_const_iterator_delete (it%obj)
  end subroutine hepmc_vertex_particle_in_iterator_final

\eatline
\nwendcode{}\nwbegindocs{633}Increment
\nwenddocs{}\nwbegincode{634}\moddef{HepMC interface: interfaces}\plusendmoddef
  interface
     subroutine vertex_particles_in_const_iterator_advance (it_obj) bind(C)
       import
       type(c_ptr), value :: it_obj
     end subroutine vertex_particles_in_const_iterator_advance
  end interface
\eatline
\nwendcode{}\nwbegincode{635}\moddef{HepMC interface: public}\plusendmoddef
  public :: hepmc_vertex_particle_in_iterator_advance
\nwendcode{}\nwbegincode{636}\moddef{HepMC interface: procedures}\plusendmoddef
  subroutine hepmc_vertex_particle_in_iterator_advance (it)
    type(hepmc_vertex_particle_in_iterator_t), intent(inout) :: it
    call vertex_particles_in_const_iterator_advance (it%obj)
  end subroutine hepmc_vertex_particle_in_iterator_advance

\eatline
\nwendcode{}\nwbegindocs{637}Reset to the beginning
\nwenddocs{}\nwbegincode{638}\moddef{HepMC interface: interfaces}\plusendmoddef
  interface
     subroutine vertex_particles_in_const_iterator_reset &
          (it_obj, v_obj) bind(C)
       import
       type(c_ptr), value :: it_obj, v_obj
     end subroutine vertex_particles_in_const_iterator_reset
  end interface
\eatline
\nwendcode{}\nwbegincode{639}\moddef{HepMC interface: public}\plusendmoddef
  public :: hepmc_vertex_particle_in_iterator_reset
\nwendcode{}\nwbegincode{640}\moddef{HepMC interface: procedures}\plusendmoddef
  subroutine hepmc_vertex_particle_in_iterator_reset (it)
    type(hepmc_vertex_particle_in_iterator_t), intent(inout) :: it
    call vertex_particles_in_const_iterator_reset (it%obj, it%v_obj)
  end subroutine hepmc_vertex_particle_in_iterator_reset

\eatline
\nwendcode{}\nwbegindocs{641}Test: return true as long as we are not past the end.
\nwenddocs{}\nwbegincode{642}\moddef{HepMC interface: interfaces}\plusendmoddef
  interface
     function vertex_particles_in_const_iterator_is_valid &
          (it_obj, v_obj) result (flag) bind(C)
       import
       logical(c_bool) :: flag
       type(c_ptr), value :: it_obj, v_obj
     end function vertex_particles_in_const_iterator_is_valid
  end interface
\eatline
\nwendcode{}\nwbegincode{643}\moddef{HepMC interface: public}\plusendmoddef
  public :: hepmc_vertex_particle_in_iterator_is_valid
\nwendcode{}\nwbegincode{644}\moddef{HepMC interface: procedures}\plusendmoddef
  function hepmc_vertex_particle_in_iterator_is_valid (it) result (flag)
    logical :: flag
    type(hepmc_vertex_particle_in_iterator_t), intent(in) :: it
    flag = vertex_particles_in_const_iterator_is_valid (it%obj, it%v_obj)
  end function hepmc_vertex_particle_in_iterator_is_valid

\eatline
\nwendcode{}\nwbegindocs{645}Return the particle pointed to by the iterator.  (The particle
object should not be finalized, since it contains merely a pointer to
the particle which is owned by the vertex.)
\nwenddocs{}\nwbegincode{646}\moddef{HepMC interface: interfaces}\plusendmoddef
  interface
     type(c_ptr) function &
          vertex_particles_in_const_iterator_get (it_obj) bind(C)
       import
       type(c_ptr), value :: it_obj
     end function vertex_particles_in_const_iterator_get
  end interface
\eatline
\nwendcode{}\nwbegincode{647}\moddef{HepMC interface: public}\plusendmoddef
  public :: hepmc_vertex_particle_in_iterator_get
\nwendcode{}\nwbegincode{648}\moddef{HepMC interface: procedures}\plusendmoddef
  function hepmc_vertex_particle_in_iterator_get (it) result (prt)
    type(hepmc_particle_t) :: prt
    type(hepmc_vertex_particle_in_iterator_t), intent(in) :: it
    prt%obj = vertex_particles_in_const_iterator_get (it%obj)
  end function hepmc_vertex_particle_in_iterator_get

\eatline
\nwendcode{}\nwbegindocs{649}\nwdocspar
\subsection{Vertex-particle-out iterator}
This iterator iterates over all incoming particles in an vertex.  We store a
pointer to the vertex in addition to the iterator.  This allows for
simple end checking.

The iterator is actually a constant iterator; it can only read.
\nwenddocs{}\nwbegincode{650}\moddef{HepMC interface: public}\plusendmoddef
  public :: hepmc_vertex_particle_out_iterator_t
\nwendcode{}\nwbegincode{651}\moddef{HepMC interface: types}\plusendmoddef
  type :: hepmc_vertex_particle_out_iterator_t
     private
     type(c_ptr) :: obj
     type(c_ptr) :: v_obj
  end type hepmc_vertex_particle_out_iterator_t

\eatline
\nwendcode{}\nwbegindocs{652}Constructor.  The iterator is initialized at the first particle in
the vertex.
\nwenddocs{}\nwbegincode{653}\moddef{HepMC interface: interfaces}\plusendmoddef
  interface
     type(c_ptr) function &
          new_vertex_particles_out_const_iterator (v_obj) bind(C)
       import
       type(c_ptr), value :: v_obj
     end function new_vertex_particles_out_const_iterator
  end interface
\eatline
\nwendcode{}\nwbegincode{654}\moddef{HepMC interface: public}\plusendmoddef
  public :: hepmc_vertex_particle_out_iterator_init
\nwendcode{}\nwbegincode{655}\moddef{HepMC interface: procedures}\plusendmoddef
  subroutine hepmc_vertex_particle_out_iterator_init (it, v)
    type(hepmc_vertex_particle_out_iterator_t), intent(out) :: it
    type(hepmc_vertex_t), intent(in) :: v
    it%obj = new_vertex_particles_out_const_iterator (v%obj)
    it%v_obj = v%obj
  end subroutine hepmc_vertex_particle_out_iterator_init

\eatline
\nwendcode{}\nwbegindocs{656}Destructor.  Necessary because the iterator is allocated on the
heap.
\nwenddocs{}\nwbegincode{657}\moddef{HepMC interface: interfaces}\plusendmoddef
  interface
     subroutine vertex_particles_out_const_iterator_delete (it_obj) bind(C)
       import
       type(c_ptr), value :: it_obj
     end subroutine vertex_particles_out_const_iterator_delete
  end interface
\eatline
\nwendcode{}\nwbegincode{658}\moddef{HepMC interface: public}\plusendmoddef
  public :: hepmc_vertex_particle_out_iterator_final
\nwendcode{}\nwbegincode{659}\moddef{HepMC interface: procedures}\plusendmoddef
  subroutine hepmc_vertex_particle_out_iterator_final (it)
    type(hepmc_vertex_particle_out_iterator_t), intent(inout) :: it
    call vertex_particles_out_const_iterator_delete (it%obj)
  end subroutine hepmc_vertex_particle_out_iterator_final

\eatline
\nwendcode{}\nwbegindocs{660}Increment
\nwenddocs{}\nwbegincode{661}\moddef{HepMC interface: interfaces}\plusendmoddef
  interface
     subroutine vertex_particles_out_const_iterator_advance (it_obj) bind(C)
       import
       type(c_ptr), value :: it_obj
     end subroutine vertex_particles_out_const_iterator_advance
  end interface
\eatline
\nwendcode{}\nwbegincode{662}\moddef{HepMC interface: public}\plusendmoddef
  public :: hepmc_vertex_particle_out_iterator_advance
\nwendcode{}\nwbegincode{663}\moddef{HepMC interface: procedures}\plusendmoddef
  subroutine hepmc_vertex_particle_out_iterator_advance (it)
    type(hepmc_vertex_particle_out_iterator_t), intent(inout) :: it
    call vertex_particles_out_const_iterator_advance (it%obj)
  end subroutine hepmc_vertex_particle_out_iterator_advance

\eatline
\nwendcode{}\nwbegindocs{664}Reset to the beginning
\nwenddocs{}\nwbegincode{665}\moddef{HepMC interface: interfaces}\plusendmoddef
  interface
     subroutine vertex_particles_out_const_iterator_reset &
          (it_obj, v_obj) bind(C)
       import
       type(c_ptr), value :: it_obj, v_obj
     end subroutine vertex_particles_out_const_iterator_reset
  end interface
\eatline
\nwendcode{}\nwbegincode{666}\moddef{HepMC interface: public}\plusendmoddef
  public :: hepmc_vertex_particle_out_iterator_reset
\nwendcode{}\nwbegincode{667}\moddef{HepMC interface: procedures}\plusendmoddef
  subroutine hepmc_vertex_particle_out_iterator_reset (it)
    type(hepmc_vertex_particle_out_iterator_t), intent(inout) :: it
    call vertex_particles_out_const_iterator_reset (it%obj, it%v_obj)
  end subroutine hepmc_vertex_particle_out_iterator_reset

\eatline
\nwendcode{}\nwbegindocs{668}Test: return true as long as we are not past the end.
\nwenddocs{}\nwbegincode{669}\moddef{HepMC interface: interfaces}\plusendmoddef
  interface
     function vertex_particles_out_const_iterator_is_valid &
          (it_obj, v_obj) result (flag) bind(C)
       import
       logical(c_bool) :: flag
       type(c_ptr), value :: it_obj, v_obj
     end function vertex_particles_out_const_iterator_is_valid
  end interface
\eatline
\nwendcode{}\nwbegincode{670}\moddef{HepMC interface: public}\plusendmoddef
  public :: hepmc_vertex_particle_out_iterator_is_valid
\nwendcode{}\nwbegincode{671}\moddef{HepMC interface: procedures}\plusendmoddef
  function hepmc_vertex_particle_out_iterator_is_valid (it) result (flag)
    logical :: flag
    type(hepmc_vertex_particle_out_iterator_t), intent(in) :: it
    flag = vertex_particles_out_const_iterator_is_valid (it%obj, it%v_obj)
  end function hepmc_vertex_particle_out_iterator_is_valid

\eatline
\nwendcode{}\nwbegindocs{672}Return the particle pointed to by the iterator.  (The particle
object should not be finalized, since it contains merely a pointer to
the particle which is owned by the vertex.)
\nwenddocs{}\nwbegincode{673}\moddef{HepMC interface: interfaces}\plusendmoddef
  interface
     type(c_ptr) function &
          vertex_particles_out_const_iterator_get (it_obj) bind(C)
       import
       type(c_ptr), value :: it_obj
     end function vertex_particles_out_const_iterator_get
  end interface
\eatline
\nwendcode{}\nwbegincode{674}\moddef{HepMC interface: public}\plusendmoddef
  public :: hepmc_vertex_particle_out_iterator_get
\nwendcode{}\nwbegincode{675}\moddef{HepMC interface: procedures}\plusendmoddef
  function hepmc_vertex_particle_out_iterator_get (it) result (prt)
    type(hepmc_particle_t) :: prt
    type(hepmc_vertex_particle_out_iterator_t), intent(in) :: it
    prt%obj = vertex_particles_out_const_iterator_get (it%obj)
  end function hepmc_vertex_particle_out_iterator_get

\eatline
\nwendcode{}\nwbegindocs{676}\nwdocspar
\subsection{GenEvent}
The main object of HepMC is a GenEvent.  The object is filled by
GenVertex objects, which in turn contain GenParticle objects.
\nwenddocs{}\nwbegincode{677}\moddef{HepMC interface: public}\plusendmoddef
  public :: hepmc_event_t
\nwendcode{}\nwbegincode{678}\moddef{HepMC interface: types}\plusendmoddef
  type :: hepmc_event_t
     private
     type(c_ptr) :: obj
  end type hepmc_event_t

\eatline
\nwendcode{}\nwbegindocs{679}Constructor.  Arguments are process ID (integer) and event ID
(integer).

The Fortran version has initializer form.
\nwenddocs{}\nwbegincode{680}\moddef{HepMC interface: interfaces}\plusendmoddef
  interface
     type(c_ptr) function new_gen_event (proc_id, event_id) bind(C)
       import
       integer(c_int), value :: proc_id, event_id
     end function new_gen_event
  end interface
\eatline
\nwendcode{}\nwbegincode{681}\moddef{HepMC interface: public}\plusendmoddef
  public :: hepmc_event_init
\nwendcode{}\nwbegincode{682}\moddef{HepMC interface: procedures}\plusendmoddef
  subroutine hepmc_event_init (evt, proc_id, event_id)
    type(hepmc_event_t), intent(out) :: evt
    integer, intent(in), optional :: proc_id, event_id
    integer(c_int) :: pid, eid
    pid = 0;  if (present (proc_id))  pid = proc_id
    eid = 0;  if (present (event_id)) eid = event_id
    evt%obj = new_gen_event (pid, eid)
  end subroutine hepmc_event_init

\eatline
\nwendcode{}\nwbegindocs{683}Destructor.
\nwenddocs{}\nwbegincode{684}\moddef{HepMC interface: interfaces}\plusendmoddef
  interface
     subroutine gen_event_delete (evt_obj) bind(C)
       import
       type(c_ptr), value :: evt_obj
     end subroutine gen_event_delete
  end interface
\eatline
\nwendcode{}\nwbegincode{685}\moddef{HepMC interface: public}\plusendmoddef
  public :: hepmc_event_final
\nwendcode{}\nwbegincode{686}\moddef{HepMC interface: procedures}\plusendmoddef
  subroutine hepmc_event_final (evt)
    type(hepmc_event_t), intent(inout) :: evt
    call gen_event_delete (evt%obj)
  end subroutine hepmc_event_final

\eatline
\nwendcode{}\nwbegindocs{687}Screen output.  Printing to file is possible in principle (using a
C++ output channel), by allowing an argument.  Printing to an open
Fortran unit is obviously not possible.
\nwenddocs{}\nwbegincode{688}\moddef{HepMC interface: interfaces}\plusendmoddef
  interface
     subroutine gen_event_print (evt_obj) bind(C)
       import
       type(c_ptr), value :: evt_obj
     end subroutine gen_event_print
  end interface
\eatline
\nwendcode{}\nwbegincode{689}\moddef{HepMC interface: public}\plusendmoddef
  public :: hepmc_event_print
\nwendcode{}\nwbegincode{690}\moddef{HepMC interface: procedures}\plusendmoddef
  subroutine hepmc_event_print (evt)
    type(hepmc_event_t), intent(in) :: evt
    call gen_event_print (evt%obj)
  end subroutine hepmc_event_print
    
\eatline
\nwendcode{}\nwbegindocs{691}Get the event number.
\nwenddocs{}\nwbegincode{692}\moddef{HepMC interface: interfaces}\plusendmoddef
  interface
     integer(c_int) function gen_event_event_number (evt_obj) bind(C)
       use iso_c_binding !NODEP!
       type(c_ptr), value :: evt_obj
     end function gen_event_event_number
  end interface
\eatline
\nwendcode{}\nwbegincode{693}\moddef{HepMC interface: public}\plusendmoddef
  public :: hepmc_event_get_event_index
\nwendcode{}\nwbegincode{694}\moddef{HepMC interface: procedures}\plusendmoddef
  function hepmc_event_get_event_index (evt) result (i_proc)
    integer :: i_proc
    type(hepmc_event_t), intent(in) :: evt
    i_proc = gen_event_event_number (evt%obj)
  end function hepmc_event_get_event_index

\eatline
\nwendcode{}\nwbegindocs{695}Set the numeric signal process ID
\nwenddocs{}\nwbegincode{696}\moddef{HepMC interface: interfaces}\plusendmoddef
  interface
     subroutine gen_event_set_signal_process_id (evt_obj, proc_id) bind(C)
       import
       type(c_ptr), value :: evt_obj
       integer(c_int), value :: proc_id
     end subroutine gen_event_set_signal_process_id
  end interface
\eatline
\nwendcode{}\nwbegincode{697}\moddef{HepMC interface: public}\plusendmoddef
  public :: hepmc_event_set_process_id
\nwendcode{}\nwbegincode{698}\moddef{HepMC interface: procedures}\plusendmoddef
  subroutine hepmc_event_set_process_id (evt, proc)
    type(hepmc_event_t), intent(in) :: evt
    integer, intent(in) :: proc
    integer(c_int) :: i_proc
    i_proc = proc
    call gen_event_set_signal_process_id (evt%obj, i_proc)
  end subroutine hepmc_event_set_process_id

\eatline
\nwendcode{}\nwbegindocs{699}Get the numeric signal process ID
\nwenddocs{}\nwbegincode{700}\moddef{HepMC interface: interfaces}\plusendmoddef
  interface
     integer(c_int) function gen_event_signal_process_id (evt_obj) bind(C)
       import
       type(c_ptr), value :: evt_obj
     end function gen_event_signal_process_id
  end interface
\eatline
\nwendcode{}\nwbegincode{701}\moddef{HepMC interface: public}\plusendmoddef
  public :: hepmc_event_get_process_id
\nwendcode{}\nwbegincode{702}\moddef{HepMC interface: procedures}\plusendmoddef
  function hepmc_event_get_process_id (evt) result (i_proc)
    integer :: i_proc
    type(hepmc_event_t), intent(in) :: evt
    i_proc = gen_event_signal_process_id (evt%obj)
  end function hepmc_event_get_process_id

\eatline
\nwendcode{}\nwbegindocs{703}Set the event energy scale
\nwenddocs{}\nwbegincode{704}\moddef{HepMC interface: interfaces}\plusendmoddef
  interface
     subroutine gen_event_set_event_scale (evt_obj, scale) bind(C)
       import
       type(c_ptr), value :: evt_obj
       real(c_double), value :: scale
     end subroutine gen_event_set_event_scale
  end interface
\eatline
\nwendcode{}\nwbegincode{705}\moddef{HepMC interface: public}\plusendmoddef
  public :: hepmc_event_set_scale
\nwendcode{}\nwbegincode{706}\moddef{HepMC interface: procedures}\plusendmoddef
  subroutine hepmc_event_set_scale (evt, scale)
    type(hepmc_event_t), intent(in) :: evt
    real(default), intent(in) :: scale
    real(c_double) :: cscale
    cscale = scale
    call gen_event_set_event_scale (evt%obj, cscale)
  end subroutine hepmc_event_set_scale

\eatline
\nwendcode{}\nwbegindocs{707}Get the event energy scale
\nwenddocs{}\nwbegincode{708}\moddef{HepMC interface: interfaces}\plusendmoddef
  interface
     real(c_double) function gen_event_event_scale (evt_obj) bind(C)
       import
       type(c_ptr), value :: evt_obj
     end function gen_event_event_scale
  end interface
\eatline
\nwendcode{}\nwbegincode{709}\moddef{HepMC interface: public}\plusendmoddef
  public :: hepmc_event_get_scale
\nwendcode{}\nwbegincode{710}\moddef{HepMC interface: procedures}\plusendmoddef
  function hepmc_event_get_scale (evt) result (scale)
    real(default) :: scale
    type(hepmc_event_t), intent(in) :: evt
    scale = gen_event_event_scale (evt%obj)
  end function hepmc_event_get_scale

\eatline
\nwendcode{}\nwbegindocs{711}Set the value of $\alpha_{\rm QCD}$.
\nwenddocs{}\nwbegincode{712}\moddef{HepMC interface: interfaces}\plusendmoddef
  interface
     subroutine gen_event_set_alpha_qcd (evt_obj, a) bind(C)
       import
       type(c_ptr), value :: evt_obj
       real(c_double), value :: a
     end subroutine gen_event_set_alpha_qcd
  end interface
\eatline
\nwendcode{}\nwbegincode{713}\moddef{HepMC interface: public}\plusendmoddef
  public :: hepmc_event_set_alpha_qcd
\nwendcode{}\nwbegincode{714}\moddef{HepMC interface: procedures}\plusendmoddef
  subroutine hepmc_event_set_alpha_qcd (evt, alpha)
    type(hepmc_event_t), intent(in) :: evt
    real(default), intent(in) :: alpha
    real(c_double) :: a
    a = alpha
    call gen_event_set_alpha_qcd (evt%obj, a)
  end subroutine hepmc_event_set_alpha_qcd

\eatline
\nwendcode{}\nwbegindocs{715}Get the value of $\alpha_{\rm QCD}$.
\nwenddocs{}\nwbegincode{716}\moddef{HepMC interface: interfaces}\plusendmoddef
  interface
     real(c_double) function gen_event_alpha_qcd (evt_obj) bind(C)
       import
       type(c_ptr), value :: evt_obj
     end function gen_event_alpha_qcd
  end interface
\eatline
\nwendcode{}\nwbegincode{717}\moddef{HepMC interface: public}\plusendmoddef
  public :: hepmc_event_get_alpha_qcd
\nwendcode{}\nwbegincode{718}\moddef{HepMC interface: procedures}\plusendmoddef
  function hepmc_event_get_alpha_qcd (evt) result (alpha)
    real(default) :: alpha
    type(hepmc_event_t), intent(in) :: evt
    alpha = gen_event_alpha_qcd (evt%obj)
  end function hepmc_event_get_alpha_qcd

\eatline
\nwendcode{}\nwbegindocs{719}Set the value of $\alpha_{\rm QED}$.
\nwenddocs{}\nwbegincode{720}\moddef{HepMC interface: interfaces}\plusendmoddef
  interface
     subroutine gen_event_set_alpha_qed (evt_obj, a) bind(C)
       import
       type(c_ptr), value :: evt_obj
       real(c_double), value :: a
     end subroutine gen_event_set_alpha_qed
  end interface
\eatline
\nwendcode{}\nwbegincode{721}\moddef{HepMC interface: public}\plusendmoddef
  public :: hepmc_event_set_alpha_qed
\nwendcode{}\nwbegincode{722}\moddef{HepMC interface: procedures}\plusendmoddef
  subroutine hepmc_event_set_alpha_qed (evt, alpha)
    type(hepmc_event_t), intent(in) :: evt
    real(default), intent(in) :: alpha
    real(c_double) :: a
    a = alpha
    call gen_event_set_alpha_qed (evt%obj, a)
  end subroutine hepmc_event_set_alpha_qed

\eatline
\nwendcode{}\nwbegindocs{723}Get the value of $\alpha_{\rm QED}$.
\nwenddocs{}\nwbegincode{724}\moddef{HepMC interface: interfaces}\plusendmoddef
  interface
     real(c_double) function gen_event_alpha_qed (evt_obj) bind(C)
       import
       type(c_ptr), value :: evt_obj
     end function gen_event_alpha_qed
  end interface
\eatline
\nwendcode{}\nwbegincode{725}\moddef{HepMC interface: public}\plusendmoddef
  public :: hepmc_event_get_alpha_qed
\nwendcode{}\nwbegincode{726}\moddef{HepMC interface: procedures}\plusendmoddef
  function hepmc_event_get_alpha_qed (evt) result (alpha)
    real(default) :: alpha
    type(hepmc_event_t), intent(in) :: evt
    alpha = gen_event_alpha_qed (evt%obj)
  end function hepmc_event_get_alpha_qed

\eatline
\nwendcode{}\nwbegindocs{727}Clear a weight value to the end of the weight container.
\nwenddocs{}\nwbegincode{728}\moddef{HepMC interface: interfaces}\plusendmoddef
  interface
     subroutine gen_event_clear_weights (evt_obj) bind(C)
       use iso_c_binding !NODEP!
       type(c_ptr), value :: evt_obj
     end subroutine gen_event_clear_weights
  end interface
\eatline
\nwendcode{}\nwbegincode{729}\moddef{HepMC interface: public}\plusendmoddef
  public :: hepmc_event_clear_weights
\nwendcode{}\nwbegincode{730}\moddef{HepMC interface: procedures}\plusendmoddef
  subroutine hepmc_event_clear_weights (evt)
    type(hepmc_event_t), intent(in) :: evt
    call gen_event_clear_weights (evt%obj)
  end subroutine hepmc_event_clear_weights

\eatline
\nwendcode{}\nwbegindocs{731}Add a weight value to the end of the weight container.
\nwenddocs{}\nwbegincode{732}\moddef{HepMC interface: interfaces}\plusendmoddef
  interface
     subroutine gen_event_add_weight (evt_obj, w) bind(C)
       use iso_c_binding !NODEP!
       type(c_ptr), value :: evt_obj
       real(c_double), value :: w
     end subroutine gen_event_add_weight
  end interface
\eatline
\nwendcode{}\nwbegincode{733}\moddef{HepMC interface: public}\plusendmoddef
  public :: hepmc_event_add_weight
\nwendcode{}\nwbegincode{734}\moddef{HepMC interface: procedures}\plusendmoddef
  subroutine hepmc_event_add_weight (evt, weight)
    type(hepmc_event_t), intent(in) :: evt
    real(default), intent(in) :: weight
    real(c_double) :: w
    w = weight
    call gen_event_add_weight (evt%obj, w)
  end subroutine hepmc_event_add_weight

\eatline
\nwendcode{}\nwbegindocs{735}Get the size of the weight container (the number of valid elements).
\nwenddocs{}\nwbegincode{736}\moddef{HepMC interface: interfaces}\plusendmoddef
  interface
     integer(c_int) function gen_event_weights_size (evt_obj) bind(C)
       use iso_c_binding !NODEP!
       type(c_ptr), value :: evt_obj
     end function gen_event_weights_size
  end interface
\eatline
\nwendcode{}\nwbegincode{737}\moddef{HepMC interface: public}\plusendmoddef
  public :: hepmc_event_get_weights_size
\nwendcode{}\nwbegincode{738}\moddef{HepMC interface: procedures}\plusendmoddef
  function hepmc_event_get_weights_size (evt) result (n)
    integer :: n
    type(hepmc_event_t), intent(in) :: evt
    n = gen_event_weights_size (evt%obj)
  end function hepmc_event_get_weights_size

\eatline
\nwendcode{}\nwbegindocs{739}Get the value of the weight with index {\tt{}i}.  (Count from 1, while C counts
from zero.)
\nwenddocs{}\nwbegincode{740}\moddef{HepMC interface: interfaces}\plusendmoddef
  interface
     real(c_double) function gen_event_weight (evt_obj, i) bind(C)
       use iso_c_binding !NODEP!
       type(c_ptr), value :: evt_obj
       integer(c_int), value :: i
     end function gen_event_weight
  end interface
\eatline
\nwendcode{}\nwbegincode{741}\moddef{HepMC interface: public}\plusendmoddef
  public :: hepmc_event_get_weight
\nwendcode{}\nwbegincode{742}\moddef{HepMC interface: procedures}\plusendmoddef
  function hepmc_event_get_weight (evt, index) result (weight)
    real(default) :: weight
    type(hepmc_event_t), intent(in) :: evt
    integer, intent(in) :: index
    integer(c_int) :: i
    i = index - 1
    weight = gen_event_weight (evt%obj, i)
  end function hepmc_event_get_weight

\eatline
\nwendcode{}\nwbegindocs{743}Add a vertex to the event container.
\nwenddocs{}\nwbegincode{744}\moddef{HepMC interface: interfaces}\plusendmoddef
  interface
     subroutine gen_event_add_vertex (evt_obj, v_obj) bind(C)
       import
       type(c_ptr), value :: evt_obj
       type(c_ptr), value :: v_obj
     end subroutine gen_event_add_vertex
  end interface
\eatline
\nwendcode{}\nwbegincode{745}\moddef{HepMC interface: public}\plusendmoddef
  public :: hepmc_event_add_vertex
\nwendcode{}\nwbegincode{746}\moddef{HepMC interface: procedures}\plusendmoddef
  subroutine hepmc_event_add_vertex (evt, v)
    type(hepmc_event_t), intent(inout) :: evt
    type(hepmc_vertex_t), intent(in) :: v
    call gen_event_add_vertex (evt%obj, v%obj)
  end subroutine hepmc_event_add_vertex

\eatline
\nwendcode{}\nwbegindocs{747}Mark a particular vertex as the signal process (hard interaction).
\nwenddocs{}\nwbegincode{748}\moddef{HepMC interface: interfaces}\plusendmoddef
  interface
     subroutine gen_event_set_signal_process_vertex (evt_obj, v_obj) bind(C)
       import
       type(c_ptr), value :: evt_obj
       type(c_ptr), value :: v_obj
     end subroutine gen_event_set_signal_process_vertex
  end interface
\eatline
\nwendcode{}\nwbegincode{749}\moddef{HepMC interface: public}\plusendmoddef
  public :: hepmc_event_set_signal_process_vertex
\nwendcode{}\nwbegincode{750}\moddef{HepMC interface: procedures}\plusendmoddef
  subroutine hepmc_event_set_signal_process_vertex (evt, v)
    type(hepmc_event_t), intent(inout) :: evt
    type(hepmc_vertex_t), intent(in) :: v
    call gen_event_set_signal_process_vertex (evt%obj, v%obj)
  end subroutine hepmc_event_set_signal_process_vertex

\eatline
\nwendcode{}\nwbegindocs{751}Return the the signal process (hard interaction).
\nwenddocs{}\nwbegincode{752}\moddef{HepMC interface: interfaces}\plusendmoddef
  interface
     function gen_event_get_signal_process_vertex (evt_obj) &
          result (v_obj) bind(C)
       import
       type(c_ptr), value :: evt_obj
       type(c_ptr) :: v_obj
     end function gen_event_get_signal_process_vertex
  end interface
\eatline
\nwendcode{}\nwbegincode{753}\moddef{HepMC interface: public}\plusendmoddef
  public :: hepmc_event_get_signal_process_vertex
\nwendcode{}\nwbegincode{754}\moddef{HepMC interface: procedures}\plusendmoddef
  function hepmc_event_get_signal_process_vertex (evt) result (v)
    type(hepmc_event_t), intent(in) :: evt
    type(hepmc_vertex_t) :: v
    v%obj = gen_event_get_signal_process_vertex (evt%obj)
  end function hepmc_event_get_signal_process_vertex

\eatline
\nwendcode{}\nwbegindocs{755}Set the beam particles explicitly.
\nwenddocs{}\nwbegincode{756}\moddef{HepMC interface: public}\plusendmoddef
  public :: hepmc_event_set_beam_particles
\nwendcode{}\nwbegincode{757}\moddef{HepMC interface: procedures}\plusendmoddef
  subroutine hepmc_event_set_beam_particles (evt, prt1, prt2)
    type(hepmc_event_t), intent(inout) :: evt
    type(hepmc_particle_t), intent(in) :: prt1, prt2
    logical(c_bool) :: flag
    flag = gen_event_set_beam_particles (evt%obj, prt1%obj, prt2%obj)
  end subroutine hepmc_event_set_beam_particles

\eatline
\nwendcode{}\nwbegindocs{758}The C function returns a boolean which we do not use.
\nwenddocs{}\nwbegincode{759}\moddef{HepMC interface: interfaces}\plusendmoddef
  interface
     logical(c_bool) function gen_event_set_beam_particles &
          (evt_obj, prt1_obj, prt2_obj) bind(C)
       import
       type(c_ptr), value :: evt_obj, prt1_obj, prt2_obj
     end function gen_event_set_beam_particles
  end interface

\eatline
\nwendcode{}\nwbegindocs{760}Set the cross section and error explicitly.  Note that HepMC uses
pb, while WHIZARD uses fb.
\nwenddocs{}\nwbegincode{761}\moddef{HepMC interface: public}\plusendmoddef
  public :: hepmc_event_set_cross_section
\nwendcode{}\nwbegincode{762}\moddef{HepMC interface: procedures}\plusendmoddef
  subroutine hepmc_event_set_cross_section (evt, xsec, xsec_err)
    type(hepmc_event_t), intent(inout) :: evt
    real(default), intent(in) :: xsec, xsec_err
    call gen_event_set_cross_section &
         (evt%obj, &
         real (xsec * 1e-3_default, c_double), &
         real (xsec_err * 1e-3_default, c_double))
  end subroutine hepmc_event_set_cross_section

\eatline
\nwendcode{}\nwbegindocs{763}The C function returns a boolean which we do not use.
\nwenddocs{}\nwbegincode{764}\moddef{HepMC interface: interfaces}\plusendmoddef
  interface
     subroutine gen_event_set_cross_section (evt_obj, xs, xs_err) bind(C)
       import
       type(c_ptr), value :: evt_obj
       real(c_double), value :: xs, xs_err
     end subroutine gen_event_set_cross_section
  end interface

\eatline
\nwendcode{}\nwbegindocs{765}\nwdocspar
\subsection{Event-particle iterator}
This iterator iterates over all particles in an event.  We store a
pointer to the event in addition to the iterator.  This allows for
simple end checking.

The iterator is actually a constant iterator; it can only read.
\nwenddocs{}\nwbegincode{766}\moddef{HepMC interface: public}\plusendmoddef
  public :: hepmc_event_particle_iterator_t
\nwendcode{}\nwbegincode{767}\moddef{HepMC interface: types}\plusendmoddef
  type :: hepmc_event_particle_iterator_t
     private
     type(c_ptr) :: obj
     type(c_ptr) :: evt_obj
  end type hepmc_event_particle_iterator_t

\eatline
\nwendcode{}\nwbegindocs{768}Constructor.  The iterator is initialized at the first particle in
the event.
\nwenddocs{}\nwbegincode{769}\moddef{HepMC interface: interfaces}\plusendmoddef
  interface
     type(c_ptr) function new_event_particle_const_iterator (evt_obj) bind(C)
       import
       type(c_ptr), value :: evt_obj
     end function new_event_particle_const_iterator
  end interface
\eatline
\nwendcode{}\nwbegincode{770}\moddef{HepMC interface: public}\plusendmoddef
  public :: hepmc_event_particle_iterator_init
\nwendcode{}\nwbegincode{771}\moddef{HepMC interface: procedures}\plusendmoddef
  subroutine hepmc_event_particle_iterator_init (it, evt)
    type(hepmc_event_particle_iterator_t), intent(out) :: it
    type(hepmc_event_t), intent(in) :: evt
    it%obj = new_event_particle_const_iterator (evt%obj)
    it%evt_obj = evt%obj
  end subroutine hepmc_event_particle_iterator_init

\eatline
\nwendcode{}\nwbegindocs{772}Destructor.  Necessary because the iterator is allocated on the
heap.
\nwenddocs{}\nwbegincode{773}\moddef{HepMC interface: interfaces}\plusendmoddef
  interface
     subroutine event_particle_const_iterator_delete (it_obj) bind(C)
       import
       type(c_ptr), value :: it_obj
     end subroutine event_particle_const_iterator_delete
  end interface
\eatline
\nwendcode{}\nwbegincode{774}\moddef{HepMC interface: public}\plusendmoddef
  public :: hepmc_event_particle_iterator_final
\nwendcode{}\nwbegincode{775}\moddef{HepMC interface: procedures}\plusendmoddef
  subroutine hepmc_event_particle_iterator_final (it)
    type(hepmc_event_particle_iterator_t), intent(inout) :: it
    call event_particle_const_iterator_delete (it%obj)
  end subroutine hepmc_event_particle_iterator_final

\eatline
\nwendcode{}\nwbegindocs{776}Increment
\nwenddocs{}\nwbegincode{777}\moddef{HepMC interface: interfaces}\plusendmoddef
  interface
     subroutine event_particle_const_iterator_advance (it_obj) bind(C)
       import
       type(c_ptr), value :: it_obj
     end subroutine event_particle_const_iterator_advance
  end interface
\eatline
\nwendcode{}\nwbegincode{778}\moddef{HepMC interface: public}\plusendmoddef
  public :: hepmc_event_particle_iterator_advance
\nwendcode{}\nwbegincode{779}\moddef{HepMC interface: procedures}\plusendmoddef
  subroutine hepmc_event_particle_iterator_advance (it)
    type(hepmc_event_particle_iterator_t), intent(inout) :: it
    call event_particle_const_iterator_advance (it%obj)
  end subroutine hepmc_event_particle_iterator_advance

\eatline
\nwendcode{}\nwbegindocs{780}Reset to the beginning
\nwenddocs{}\nwbegincode{781}\moddef{HepMC interface: interfaces}\plusendmoddef
  interface
     subroutine event_particle_const_iterator_reset (it_obj, evt_obj) bind(C)
       import
       type(c_ptr), value :: it_obj, evt_obj
     end subroutine event_particle_const_iterator_reset
  end interface
\eatline
\nwendcode{}\nwbegincode{782}\moddef{HepMC interface: public}\plusendmoddef
  public :: hepmc_event_particle_iterator_reset
\nwendcode{}\nwbegincode{783}\moddef{HepMC interface: procedures}\plusendmoddef
  subroutine hepmc_event_particle_iterator_reset (it)
    type(hepmc_event_particle_iterator_t), intent(inout) :: it
    call event_particle_const_iterator_reset (it%obj, it%evt_obj)
  end subroutine hepmc_event_particle_iterator_reset

\eatline
\nwendcode{}\nwbegindocs{784}Test: return true as long as we are not past the end.
\nwenddocs{}\nwbegincode{785}\moddef{HepMC interface: interfaces}\plusendmoddef
  interface
     function event_particle_const_iterator_is_valid &
          (it_obj, evt_obj) result (flag) bind(C)
       import
       logical(c_bool) :: flag
       type(c_ptr), value :: it_obj, evt_obj
     end function event_particle_const_iterator_is_valid
  end interface
\eatline
\nwendcode{}\nwbegincode{786}\moddef{HepMC interface: public}\plusendmoddef
  public :: hepmc_event_particle_iterator_is_valid
\nwendcode{}\nwbegincode{787}\moddef{HepMC interface: procedures}\plusendmoddef
  function hepmc_event_particle_iterator_is_valid (it) result (flag)
    logical :: flag
    type(hepmc_event_particle_iterator_t), intent(in) :: it
    flag = event_particle_const_iterator_is_valid (it%obj, it%evt_obj)
  end function hepmc_event_particle_iterator_is_valid

\eatline
\nwendcode{}\nwbegindocs{788}Return the particle pointed to by the iterator.  (The particle
object should not be finalized, since it contains merely a pointer to
the particle which is owned by the vertex.)
\nwenddocs{}\nwbegincode{789}\moddef{HepMC interface: interfaces}\plusendmoddef
  interface
     type(c_ptr) function event_particle_const_iterator_get (it_obj) bind(C)
       import
       type(c_ptr), value :: it_obj
     end function event_particle_const_iterator_get
  end interface
\eatline
\nwendcode{}\nwbegincode{790}\moddef{HepMC interface: public}\plusendmoddef
  public :: hepmc_event_particle_iterator_get
\nwendcode{}\nwbegincode{791}\moddef{HepMC interface: procedures}\plusendmoddef
  function hepmc_event_particle_iterator_get (it) result (prt)
    type(hepmc_particle_t) :: prt
    type(hepmc_event_particle_iterator_t), intent(in) :: it
    prt%obj = event_particle_const_iterator_get (it%obj)
  end function hepmc_event_particle_iterator_get

\eatline
\nwendcode{}\nwbegindocs{792}\nwdocspar
\subsection{I/O streams}
There is a specific I/O stream type for handling the output of
GenEvent objects (i.e., Monte Carlo event samples) to file.  Opening
the file is done by the constructor, closing by the destructor.
\nwenddocs{}\nwbegincode{793}\moddef{HepMC interface: public}\plusendmoddef
  public :: hepmc_iostream_t
\nwendcode{}\nwbegincode{794}\moddef{HepMC interface: types}\plusendmoddef
  type :: hepmc_iostream_t
     private
     type(c_ptr) :: obj
  end type hepmc_iostream_t

\eatline
\nwendcode{}\nwbegindocs{795}Constructor for an output stream associated to a file.
\nwenddocs{}\nwbegincode{796}\moddef{HepMC interface: interfaces}\plusendmoddef
  interface
     type(c_ptr) function new_io_gen_event_out (filename) bind(C)
       import
       character(c_char), dimension(*), intent(in) :: filename
     end function new_io_gen_event_out
  end interface
\eatline
\nwendcode{}\nwbegincode{797}\moddef{HepMC interface: public}\plusendmoddef
  public :: hepmc_iostream_open_out
\nwendcode{}\nwbegincode{798}\moddef{HepMC interface: procedures}\plusendmoddef
  subroutine hepmc_iostream_open_out (iostream, filename)
    type(hepmc_iostream_t), intent(out) :: iostream
    type(string_t), intent(in) :: filename
    iostream%obj = new_io_gen_event_out (char (filename) // c_null_char)
  end subroutine hepmc_iostream_open_out

\eatline
\nwendcode{}\nwbegindocs{799}Constructor for an input stream associated to a file.
\nwenddocs{}\nwbegincode{800}\moddef{HepMC interface: interfaces}\plusendmoddef
  interface
     type(c_ptr) function new_io_gen_event_in (filename) bind(C)
       import
       character(c_char), dimension(*), intent(in) :: filename
     end function new_io_gen_event_in
  end interface
\eatline
\nwendcode{}\nwbegincode{801}\moddef{HepMC interface: public}\plusendmoddef
  public :: hepmc_iostream_open_in
\nwendcode{}\nwbegincode{802}\moddef{HepMC interface: procedures}\plusendmoddef
  subroutine hepmc_iostream_open_in (iostream, filename)
    type(hepmc_iostream_t), intent(out) :: iostream
    type(string_t), intent(in) :: filename
    iostream%obj = new_io_gen_event_in (char (filename) // c_null_char)
  end subroutine hepmc_iostream_open_in

\eatline
\nwendcode{}\nwbegindocs{803}Destructor:
\nwenddocs{}\nwbegincode{804}\moddef{HepMC interface: interfaces}\plusendmoddef
  interface
     subroutine io_gen_event_delete (io_obj) bind(C)
       import
       type(c_ptr), value :: io_obj
     end subroutine io_gen_event_delete
  end interface
\eatline
\nwendcode{}\nwbegincode{805}\moddef{HepMC interface: public}\plusendmoddef
  public :: hepmc_iostream_close
\nwendcode{}\nwbegincode{806}\moddef{HepMC interface: procedures}\plusendmoddef
  subroutine hepmc_iostream_close (iostream)
    type(hepmc_iostream_t), intent(inout) :: iostream
    call io_gen_event_delete (iostream%obj)
  end subroutine hepmc_iostream_close

\eatline
\nwendcode{}\nwbegindocs{807}Write a single event to the I/O stream.
\nwenddocs{}\nwbegincode{808}\moddef{HepMC interface: interfaces}\plusendmoddef
  interface
     subroutine io_gen_event_write_event (io_obj, evt_obj) bind(C)
       import
       type(c_ptr), value :: io_obj, evt_obj
     end subroutine io_gen_event_write_event
  end interface
\eatline
\nwendcode{}\nwbegincode{809}\moddef{HepMC interface: public}\plusendmoddef
  public :: hepmc_iostream_write_event
\nwendcode{}\nwbegincode{810}\moddef{HepMC interface: procedures}\plusendmoddef
  subroutine hepmc_iostream_write_event (iostream, evt)
    type(hepmc_iostream_t), intent(inout) :: iostream
    type(hepmc_event_t), intent(in) :: evt
    call io_gen_event_write_event (iostream%obj, evt%obj)
  end subroutine hepmc_iostream_write_event

\eatline
\nwendcode{}\nwbegindocs{811}Read a single event from the I/O stream.  Return true if successful.
\nwenddocs{}\nwbegincode{812}\moddef{HepMC interface: interfaces}\plusendmoddef
  interface
     logical(c_bool) function io_gen_event_read_event (io_obj, evt_obj) bind(C)
       import
       type(c_ptr), value :: io_obj, evt_obj
     end function io_gen_event_read_event
  end interface
\eatline
\nwendcode{}\nwbegincode{813}\moddef{HepMC interface: public}\plusendmoddef
  public :: hepmc_iostream_read_event
\nwendcode{}\nwbegincode{814}\moddef{HepMC interface: procedures}\plusendmoddef
  subroutine hepmc_iostream_read_event (iostream, evt, ok)
    type(hepmc_iostream_t), intent(inout) :: iostream
    type(hepmc_event_t), intent(inout) :: evt
    logical, intent(out) :: ok
    ok = io_gen_event_read_event (iostream%obj, evt%obj)
  end subroutine hepmc_iostream_read_event

\eatline
\nwendcode{}\nwbegindocs{815}\nwdocspar
\subsection{Unit tests}
This test example is an abridged version from the build-from-scratch
example in the HepMC distribution.  We create two vertices for $p\to
q$ PDF splitting, then a vertex for a $qq\to W^-g$ hard-interaction
process, and finally a vertex for $W^-\to qq$ decay.  The setup is for
LHC kinematics.

Extending the original example, we set color flow for the incoming
quarks and polarization for the outgoing photon.  For the latter, we
have to define a particle-data object for the photon, so a flavor
object can be correctly initialized.
\nwenddocs{}\nwbegincode{816}\moddef{HepMC interface: public}\plusendmoddef
  public :: hepmc_test
\nwendcode{}\nwbegincode{817}\moddef{HepMC interface: procedures}\plusendmoddef
  subroutine hepmc_test (u, results)
    integer, intent(in) :: u
    type(test_results_t), intent(inout) :: results
  \LA{}HepMC interface: execute tests\RA{}
  end subroutine hepmc_test
  
\nwendcode{}\nwbegindocs{818} %def hepmc_test
\nwenddocs{}\nwbegindocs{819}\nwdocspar
\nwenddocs{}\nwbegincode{820}\moddef{HepMC interface: execute tests}\endmoddef
  call test (hepmc_interface_1, "hepmc_interface_1", &
       "check HepMC interface", &
       u, results)
\nwendcode{}\nwbegincode{821}\moddef{HepMC interface: tests}\endmoddef
  subroutine hepmc_interface_1 (u)
    use physics_defs, only: VECTOR
    use model_data, only: field_data_t
    integer, intent(in) :: u
    integer :: u_file, iostat
    type(hepmc_event_t) :: evt
    type(hepmc_vertex_t) :: v1, v2, v3, v4
    type(hepmc_particle_t) :: prt1, prt2, prt3, prt4, prt5, prt6, prt7, prt8
    type(hepmc_iostream_t) :: iostream
    type(flavor_t) :: flv
    type(color_t) :: col
    type(polarization_t) :: pol
    type(field_data_t), target :: photon_data
    character(80) :: buffer    

    write (u, "(A)")  "* Test output: HepMC interface"
    write (u, "(A)")  "*   Purpose: test HepMC interface"
    write (u, "(A)")      
    
    write (u, "(A)")  "* Initialization"
    write (u, "(A)")
    
    ! Initialize a photon flavor object and some polarization
    call photon_data%init (var_str ("PHOTON"), 22)
    call photon_data%set (spin_type=VECTOR)
    call photon_data%freeze ()
    call flv%init (photon_data)
    call polarization_init_angles &
         (pol, flv, 0.6_default, 1._default, 0.5_default)

    ! Event initialization
    call hepmc_event_init (evt, 20, 1)

    write (u, "(A)")  "* p -> q splitting"
    write (u, "(A)")
            
    ! $p\\to q$ splittings
    call hepmc_vertex_init (v1)
    call hepmc_event_add_vertex (evt, v1)
    call hepmc_vertex_init (v2)
    call hepmc_event_add_vertex (evt, v2)
    call particle_init (prt1, &
         0._default, 0._default, 7000._default, 7000._default, &
         2212, 3)
    call hepmc_vertex_add_particle_in (v1, prt1)
    call particle_init (prt2, &
         0._default, 0._default,-7000._default, 7000._default, &
         2212, 3)
    call hepmc_vertex_add_particle_in (v2, prt2)
    call particle_init (prt3, &
         .750_default, -1.569_default, 32.191_default, 32.238_default, &
         1, 3)
    call color_init_from_array (col, [501])
    call hepmc_particle_set_color (prt3, col)
    call hepmc_vertex_add_particle_out (v1, prt3)
    call particle_init (prt4, &
         -3.047_default, -19._default, -54.629_default, 57.920_default, &
         -2, 3)
    call color_init_from_array (col, [-501])
    call hepmc_particle_set_color (prt4, col)
    call hepmc_vertex_add_particle_out (v2, prt4)

    write (u, "(A)")  "* Hard interaction"
    write (u, "(A)")
        
    ! Hard interaction
    call hepmc_vertex_init (v3)
    call hepmc_event_add_vertex (evt, v3)
    call hepmc_vertex_add_particle_in (v3, prt3)
    call hepmc_vertex_add_particle_in (v3, prt4)
    call particle_init (prt6, &
         -3.813_default, 0.113_default, -1.833_default, 4.233_default, &
         22, 1)
    call hepmc_particle_set_polarization (prt6, pol)
    call hepmc_vertex_add_particle_out (v3, prt6)
    call particle_init (prt5, &
         1.517_default, -20.68_default, -20.605_default, 85.925_default, &
         -24, 3)
    call hepmc_vertex_add_particle_out (v3, prt5)
    call hepmc_event_set_signal_process_vertex (evt, v3)
    
    ! $W^-$ decay
    call vertex_init_pos (v4, &
         0.12_default, -0.3_default, 0.05_default, 0.004_default)
    call hepmc_event_add_vertex (evt, v4)
    call hepmc_vertex_add_particle_in (v4, prt5)
    call particle_init (prt7, &
         -2.445_default, 28.816_default, 6.082_default, 29.552_default, &
         1, 1)
    call hepmc_vertex_add_particle_out (v4, prt7)
    call particle_init (prt8, &
         3.962_default, -49.498_default, -26.687_default, 56.373_default, &
         -2, 1)
    call hepmc_vertex_add_particle_out (v4, prt8)
    
    ! Event output
    call hepmc_event_print (evt)
    write (u, "(A)")  "Writing to file 'hepmc_test.hepmc'"
    write (u, "(A)")
    
    call hepmc_iostream_open_out (iostream , var_str ("hepmc_test.hepmc"))
    call hepmc_iostream_write_event (iostream, evt)
    call hepmc_iostream_close (iostream)

    write (u, "(A)")  "Writing completed"
    
    write (u, "(A)")
    write (u, "(A)")  "* File contents:"
    write (u, "(A)")

    u_file = free_unit ()
    open (u_file, file = "hepmc_test.hepmc", &
         action = "read", status = "old")
    do
       read (u_file, "(A)", iostat = iostat)  buffer
       if (buffer(1:14) == "HepMC::Version")  buffer = "[...]"
       if (iostat /= 0)  exit
       write (u, "(A)") trim (buffer)
    end do
    close (u_file)
    
    write (u, "(A)")
    write (u, "(A)")  "* Cleanup"
    write (u, "(A)")
        
    ! Wrapup
    call polarization_final (pol)
    call hepmc_event_final (evt)

    write (u, "(A)")
    write (u, "(A)")  "* Test output end: hepmc_interface_1"        
    
  contains

    subroutine vertex_init_pos (v, x, y, z, t)
      type(hepmc_vertex_t), intent(out) :: v
      real(default), intent(in) :: x, y, z, t
      type(vector4_t) :: xx
      xx = vector4_moving (t, vector3_moving ([x, y, z]))
      call hepmc_vertex_init (v, xx)
    end subroutine vertex_init_pos

    subroutine particle_init (prt, px, py, pz, E, pdg, status)
      type(hepmc_particle_t), intent(out) :: prt
      real(default), intent(in) :: px, py, pz, E
      integer, intent(in) :: pdg, status
      type(vector4_t) :: p
      p = vector4_moving (E, vector3_moving ([px, py, pz]))
      call hepmc_particle_init (prt, p, pdg, status)
    end subroutine particle_init

  end subroutine hepmc_interface_1

\eatline
\nwendcode{}\nwbegindocs{822}\nwdocspar
%%%%%%%%%%%%%%%%%%%%%%%%%%%%%%%%%%%%%%%%%%%%%%%%%%%%%%%%%%%%%%%%%%%%%%%%
\section{LCIO events}
This section provides the interface to the LCIO C++ library for handling
Monte-Carlo events.

Each C++ class of LCIO that we use is mirrored by a Fortran type,
which contains as its only component the C pointer to the C++ object.

Each C++ method of LCIO that we use has a C wrapper function.  This
function takes a pointer to the host object as its first argument.
Further arguments are either C pointers, or in the case of simple
types (integer, real), interoperable C/Fortran objects.

The C wrapper functions have explicit interfaces in the Fortran
module.  They are called by Fortran wrapper procedures.  These are
treated as methods of the corresponding Fortran type.
\nwenddocs{}\nwbegincode{823}\moddef{\code{}lcio{\char95}interface.f90\edoc{}}\endmoddef
\LA{}File header\RA{}

module lcio_interface

  use, intrinsic :: iso_c_binding !NODEP!
  
\LA{}Use kinds\RA{}
\LA{}Use strings\RA{}
  use io_units
  use constants
  use lorentz
  use unit_tests
  use diagnostics
  use model_data
  use flavors
  use colors
  use helicities
  use quantum_numbers
  use polarizations

\LA{}Standard module head\RA{}

\LA{}LCIO interface: public\RA{}

\LA{}LCIO interface: types\RA{}

\LA{}LCIO interface: interfaces\RA{}

contains

\LA{}LCIO interface: procedures\RA{}

\LA{}LCIO interface: tests\RA{}

end module lcio_interface
\eatline
\nwendcode{}\nwbegindocs{824}\nwdocspar
\subsection{Interface check}
This function can be called in order to verify that we are using the
actual LCIO library, and not the dummy version.
\nwenddocs{}\nwbegincode{825}\moddef{LCIO interface: interfaces}\endmoddef
  interface
     logical(c_bool) function lcio_available () bind(C)
       import
     end function lcio_available
  end interface
\nwendcode{}\nwbegincode{826}\moddef{LCIO interface: public}\endmoddef
  public :: lcio_is_available
\nwendcode{}\nwbegincode{827}\moddef{LCIO interface: procedures}\endmoddef
  function lcio_is_available () result (flag)
    logical :: flag
    flag = lcio_available ()
  end function lcio_is_available

\eatline
\nwendcode{}\nwbegindocs{828}\nwdocspar
\subsection{LCIO Run Header}
This is a type for the run header of the LCIO file.
\nwenddocs{}\nwbegincode{829}\moddef{LCIO interface: public}\plusendmoddef
  public :: lcio_run_header_t
\nwendcode{}\nwbegincode{830}\moddef{LCIO interface: types}\endmoddef
  type :: lcio_run_header_t
     private
     type(c_ptr) :: obj
  end type lcio_run_header_t

\eatline
\nwendcode{}\nwbegindocs{831}The Fortran version has initializer form.
\nwenddocs{}\nwbegincode{832}\moddef{LCIO interface: interfaces}\plusendmoddef
  interface
     type(c_ptr) function new_lcio_run_header (proc_id) bind(C)
       import
       integer(c_int), value :: proc_id
     end function new_lcio_run_header
  end interface
\eatline
\nwendcode{}\nwbegincode{833}\moddef{LCIO interface: interfaces}\plusendmoddef
  interface
     subroutine run_header_set_simstring &
          (runhdr_obj, simstring) bind(C)
       import
       type(c_ptr), value :: runhdr_obj
       character(c_char), dimension(*), intent(in) :: simstring
     end subroutine run_header_set_simstring
  end interface
\eatline
\nwendcode{}\nwbegincode{834}\moddef{LCIO interface: public}\plusendmoddef
  public :: lcio_run_header_init
\nwendcode{}\nwbegincode{835}\moddef{LCIO interface: procedures}\plusendmoddef
  subroutine lcio_run_header_init (runhdr, proc_id, run_id)
    type(lcio_run_header_t), intent(out) :: runhdr
    integer, intent(in), optional :: proc_id, run_id
    integer(c_int) :: pid, rid
    pid = 0;  if (present (proc_id))  pid = proc_id
    rid = 0; if (present (run_id))  rid = run_id
    runhdr%obj = new_lcio_run_header (rid)
    call run_header_set_simstring (runhdr%obj, &
         "WHIZARD version:" // "\LA{}Version\RA{}")
  end subroutine lcio_run_header_init

\eatline
\nwendcode{}\nwbegindocs{836}\nwdocspar
\nwenddocs{}\nwbegincode{837}\moddef{LCIO interface: interfaces}\plusendmoddef
  interface
     subroutine write_run_header (lcwrt_obj, runhdr_obj) bind(C)
       import
       type(c_ptr), value :: lcwrt_obj
       type(c_ptr), value :: runhdr_obj
     end subroutine write_run_header
  end interface
\eatline
\nwendcode{}\nwbegincode{838}\moddef{LCIO interface: public}\plusendmoddef
  public :: lcio_run_header_write
\nwendcode{}\nwbegincode{839}\moddef{LCIO interface: procedures}\plusendmoddef
  subroutine lcio_run_header_write (wrt, hdr)
    type(lcio_writer_t), intent(inout) :: wrt
    type(lcio_run_header_t), intent(inout) :: hdr
    call write_run_header (wrt%obj, hdr%obj)
  end subroutine lcio_run_header_write

\eatline
\nwendcode{}\nwbegindocs{840}\nwdocspar
\subsection{LCIO Event and LC Collection}
The main object of LCIO is a LCEventImpl.  The object is filled by
MCParticle objects, which are set as LCCollection.
\nwenddocs{}\nwbegincode{841}\moddef{LCIO interface: public}\plusendmoddef
  public :: lccollection_t
\nwendcode{}\nwbegincode{842}\moddef{LCIO interface: types}\plusendmoddef
  type :: lccollection_t
     private
     type(c_ptr) :: obj
  end type lccollection_t

\eatline
\nwendcode{}\nwbegindocs{843}Initializer.
\nwenddocs{}\nwbegincode{844}\moddef{LCIO interface: interfaces}\plusendmoddef
  interface
     type(c_ptr) function new_lccollection () bind(C)
       import
     end function new_lccollection
  end interface
\eatline
\nwendcode{}\nwbegincode{845}\moddef{LCIO interface: public}\plusendmoddef
  public :: lcio_event_t
\nwendcode{}\nwbegincode{846}\moddef{LCIO interface: types}\plusendmoddef
  type :: lcio_event_t
     private
     type(c_ptr) :: obj
     type(lccollection_t) :: lccoll
  end type lcio_event_t

\eatline
\nwendcode{}\nwbegindocs{847}Constructor.  Arguments are process ID (integer) and event ID
(integer).

The Fortran version has initializer form. 
\nwenddocs{}\nwbegincode{848}\moddef{LCIO interface: interfaces}\plusendmoddef
  interface
     type(c_ptr) function new_lcio_event (proc_id, event_id, run_id) bind(C)
       import
       integer(c_int), value :: proc_id, event_id, run_id
     end function new_lcio_event
  end interface
\eatline
\nwendcode{}\nwbegindocs{849}\nwdocspar
\nwenddocs{}\nwbegincode{850}\moddef{LCIO interface: public}\plusendmoddef
  public :: lcio_event_init
\nwendcode{}\nwbegincode{851}\moddef{LCIO interface: procedures}\plusendmoddef
  subroutine lcio_event_init (evt, proc_id, event_id, run_id)
    type(lcio_event_t), intent(out) :: evt
    integer, intent(in), optional :: proc_id, event_id, run_id
    integer(c_int) :: pid, eid, rid
    pid = 0;  if (present (proc_id))  pid = proc_id
    eid = 0;  if (present (event_id)) eid = event_id
    rid = 0;  if (present (run_id))   rid = run_id
    evt%obj = new_lcio_event (pid, eid, rid)
    evt%lccoll%obj = new_lccollection ()
  end subroutine lcio_event_init

\eatline
\nwendcode{}\nwbegindocs{852}Destructor.
\nwenddocs{}\nwbegincode{853}\moddef{LCIO interface: interfaces}\plusendmoddef
  interface
     subroutine lcio_event_delete (evt_obj) bind(C)
       import
       type(c_ptr), value :: evt_obj
     end subroutine lcio_event_delete
  end interface
\eatline
\nwendcode{}\nwbegindocs{854}Show event on screen.
\nwenddocs{}\nwbegincode{855}\moddef{LCIO interface: interfaces}\plusendmoddef
  interface
     subroutine dump_lcio_event (evt_obj) bind(C)
       import
       type(c_ptr), value :: evt_obj
     end subroutine dump_lcio_event
  end interface
\eatline
\nwendcode{}\nwbegincode{856}\moddef{LCIO interface: public}\plusendmoddef
  public :: show_lcio_event
\nwendcode{}\nwbegincode{857}\moddef{LCIO interface: procedures}\plusendmoddef
  subroutine show_lcio_event (evt)
    type(lcio_event_t), intent(in) :: evt
    if (c_associated (evt%obj)) then
       call dump_lcio_event (evt%obj)
    else
       call msg_error ("LCIO event is not allocated.")
    end if
  end subroutine show_lcio_event

\eatline
\nwendcode{}\nwbegindocs{858}Put a single event to file.
\nwenddocs{}\nwbegincode{859}\moddef{LCIO interface: interfaces}\plusendmoddef
  interface
     subroutine lcio_event_to_file (evt_obj, filename) bind(C)
       import
       type(c_ptr), value :: evt_obj
       character(c_char), dimension(*), intent(in) :: filename       
     end subroutine lcio_event_to_file
  end interface
\eatline
\nwendcode{}\nwbegincode{860}\moddef{LCIO interface: public}\plusendmoddef
  public :: write_lcio_event
\nwendcode{}\nwbegincode{861}\moddef{LCIO interface: procedures}\plusendmoddef
  subroutine write_lcio_event (evt, filename)
    type(lcio_event_t), intent(in) :: evt
    type(string_t), intent(in) :: filename
    call lcio_event_to_file (evt%obj, char (filename) // c_null_char)
  end subroutine write_lcio_event

\eatline
\nwendcode{}\nwbegindocs{862}\nwdocspar
\nwenddocs{}\nwbegincode{863}\moddef{LCIO interface: public}\plusendmoddef
  public :: lcio_event_final
\nwendcode{}\nwbegincode{864}\moddef{LCIO interface: procedures}\plusendmoddef
  subroutine lcio_event_final (evt)
    type(lcio_event_t), intent(inout) :: evt
    call lcio_event_delete (evt%obj)
  end subroutine lcio_event_final

\eatline
\nwendcode{}\nwbegindocs{865}\nwdocspar
\nwenddocs{}\nwbegincode{866}\moddef{LCIO interface: interfaces}\plusendmoddef
  interface
     subroutine lcio_set_alpha_qcd (evt_obj, alphas) bind(C)
       import
       type(c_ptr), value :: evt_obj
       real(c_double), value :: alphas
     end subroutine lcio_set_alpha_qcd
  end interface
  interface 
     subroutine lcio_set_scale (evt_obj, scale) bind(C)
       import
       type(c_ptr), value :: evt_obj
       real(c_double), value :: scale
     end subroutine lcio_set_scale
  end interface  
\eatline
\nwendcode{}\nwbegindocs{867}\nwdocspar
\nwenddocs{}\nwbegincode{868}\moddef{LCIO interface: public}\plusendmoddef
  public :: lcio_event_set_alpha_qcd
\nwendcode{}\nwbegincode{869}\moddef{LCIO interface: procedures}\plusendmoddef
  subroutine lcio_event_set_alpha_qcd (evt, alphas)
    type(lcio_event_t), intent(inout) :: evt
    real(default), intent(in) :: alphas
    call lcio_set_alpha_qcd (evt%obj, real (alphas, c_double))
  end subroutine lcio_event_set_alpha_qcd

\eatline
\nwendcode{}\nwbegindocs{870}\nwdocspar
\nwenddocs{}\nwbegincode{871}\moddef{LCIO interface: public}\plusendmoddef
  public :: lcio_event_set_scale
\nwendcode{}\nwbegincode{872}\moddef{LCIO interface: procedures}\plusendmoddef
  subroutine lcio_event_set_scale (evt, scale) 
    type(lcio_event_t), intent(inout) :: evt
    real(default), intent(in) :: scale
    call lcio_set_scale (evt%obj, real (scale, c_double))
  end subroutine lcio_event_set_scale

\eatline
\nwendcode{}\nwbegindocs{873}\nwdocspar
\nwenddocs{}\nwbegincode{874}\moddef{LCIO interface: interfaces}\plusendmoddef
  interface
     subroutine lcio_event_add_collection &
          (evt_obj, lccoll_obj) bind(C)
       import
       type(c_ptr), value :: evt_obj, lccoll_obj
     end subroutine lcio_event_add_collection
  end interface
\eatline
\nwendcode{}\nwbegincode{875}\moddef{LCIO interface: public}\plusendmoddef
  public :: lcio_event_add_coll
\nwendcode{}\nwbegincode{876}\moddef{LCIO interface: procedures}\plusendmoddef
  subroutine lcio_event_add_coll (evt) 
    type(lcio_event_t), intent(inout) :: evt
    call lcio_event_add_collection (evt%obj, &
         evt%lccoll%obj)
  end subroutine lcio_event_add_coll

\eatline
\nwendcode{}\nwbegindocs{877}\nwdocspar
\subsection{LCIO Particle}
Particle objects have the obvious meaning.
\nwenddocs{}\nwbegincode{878}\moddef{LCIO interface: public}\plusendmoddef
  public :: lcio_particle_t
\nwendcode{}\nwbegincode{879}\moddef{LCIO interface: types}\plusendmoddef
  type :: lcio_particle_t
     private
     type(c_ptr) :: obj
  end type lcio_particle_t

\eatline
\nwendcode{}\nwbegindocs{880}Constructor. 
\nwenddocs{}\nwbegincode{881}\moddef{LCIO interface: interfaces}\plusendmoddef
  interface
     type(c_ptr) function new_lcio_particle &
          (px, py, pz, pdg_id, mass, status) bind(C)
       import
       integer(c_int), value :: pdg_id, status
       real(c_double), value :: px, py, pz, mass
     end function new_lcio_particle
  end interface
\eatline
\nwendcode{}\nwbegindocs{882}\nwdocspar
\nwenddocs{}\nwbegincode{883}\moddef{LCIO interface: interfaces}\plusendmoddef
  interface
     subroutine add_particle_to_collection &
          (prt_obj, lccoll_obj) bind(C)
       import
       type(c_ptr), value :: prt_obj, lccoll_obj
     end subroutine add_particle_to_collection
  end interface
\eatline
\nwendcode{}\nwbegincode{884}\moddef{LCIO interface: public}\plusendmoddef
  public :: lcio_particle_add_to_evt_coll
\nwendcode{}\nwbegincode{885}\moddef{LCIO interface: procedures}\plusendmoddef
  subroutine lcio_particle_add_to_evt_coll &
       (lprt, evt)
    type(lcio_particle_t), intent(in) :: lprt
    type(lcio_event_t), intent(inout) :: evt
    call add_particle_to_collection (lprt%obj, evt%lccoll%obj)
  end subroutine lcio_particle_add_to_evt_coll

\eatline
\nwendcode{}\nwbegindocs{886}\nwdocspar
\nwenddocs{}\nwbegincode{887}\moddef{LCIO interface: public}\plusendmoddef
  public :: lcio_particle_init
\nwendcode{}\nwbegincode{888}\moddef{LCIO interface: procedures}\plusendmoddef
  subroutine lcio_particle_init (prt, p, pdg, status)
    type(lcio_particle_t), intent(out) :: prt
    type(vector4_t), intent(in) :: p
    real(default) :: mass
    real(default) :: px, py, pz
    integer, intent(in) :: pdg, status
    px = vector4_get_component (p, 1)
    py = vector4_get_component (p, 2)
    pz = vector4_get_component (p, 3)
    mass = p**1
    prt%obj = new_lcio_particle (real (px, c_double), real (py, c_double), &
         real (pz, c_double), int (pdg, c_int), &
         real (mass, c_double), int (status, c_int))
  end subroutine lcio_particle_init

\eatline
\nwendcode{}\nwbegindocs{889}Set the particle color flow.
\nwenddocs{}\nwbegincode{890}\moddef{LCIO interface: interfaces}\plusendmoddef
  interface
     subroutine lcio_set_color_flow (prt_obj, col1, col2) bind(C)
       import
       type(c_ptr), value :: prt_obj
       integer(c_int), value :: col1, col2
     end subroutine lcio_set_color_flow
  end interface
\eatline
\nwendcode{}\nwbegindocs{891}Set the particle color.  Either from a {\tt{}color{\char95}t} object or
directly from a pair of integers.
\nwenddocs{}\nwbegincode{892}\moddef{LCIO interface: interfaces}\plusendmoddef
  interface lcio_particle_set_color
     module procedure lcio_particle_set_color_col
     module procedure lcio_particle_set_color_int
  end interface lcio_particle_set_color
\nwendcode{}\nwbegincode{893}\moddef{LCIO interface: public}\plusendmoddef
  public :: lcio_particle_set_color
\nwendcode{}\nwbegincode{894}\moddef{LCIO interface: procedures}\plusendmoddef
  subroutine lcio_particle_set_color_col (prt, col)
    type(lcio_particle_t), intent(inout) :: prt
    type(color_t), intent(in) :: col
    integer(c_int), dimension(2) :: c
    c(1) = col%get_col ()
    c(2) = col%get_acl ()    
    if (c(1) /= 0 .or. c(2) /= 0)  then
       call lcio_set_color_flow (prt%obj, c(1), c(2))
    end  if
  end subroutine lcio_particle_set_color_col

  subroutine lcio_particle_set_color_int (prt, col)
    type(lcio_particle_t), intent(inout) :: prt
    integer, dimension(2), intent(in) :: col
    integer(c_int), dimension(2) :: c
    c = col
    if (c(1) /= 0 .or. c(2) /= 0) then
       call lcio_set_color_flow (prt%obj, c(1), c(2))
    end if
  end subroutine lcio_particle_set_color_int

\eatline
\nwendcode{}\nwbegindocs{895}Return the particle color as a two-dimensional array (color, anticolor).
\nwenddocs{}\nwbegincode{896}\moddef{LCIO interface: interfaces}\plusendmoddef
  interface
     integer(c_int) function lcio_particle_flow (prt_obj, col_index) bind(C)
       use iso_c_binding !NODEP!
       type(c_ptr), value :: prt_obj
       integer(c_int), value :: col_index
     end function lcio_particle_flow
  end interface
\eatline
\nwendcode{}\nwbegincode{897}\moddef{LCIO interface: public}\plusendmoddef
  public :: lcio_particle_get_flow
\nwendcode{}\nwbegincode{898}\moddef{LCIO interface: procedures}\plusendmoddef
  function lcio_particle_get_flow (prt) result (col)
    integer, dimension(2) :: col
    type(lcio_particle_t), intent(in) :: prt
    col(1) = lcio_particle_flow (prt%obj, 0_c_int)
    col(2) = - lcio_particle_flow (prt%obj, 1_c_int)    
  end function lcio_particle_get_flow

\eatline
\nwendcode{}\nwbegindocs{899}Return the four-momentum of a LCIO particle.
\nwenddocs{}\nwbegincode{900}\moddef{LCIO interface: interfaces}\plusendmoddef
  interface
     real(c_double) function lcio_three_momentum (prt_obj, p_index) bind(C)
       use iso_c_binding !NODEP!
       type(c_ptr), value :: prt_obj
       integer(c_int), value :: p_index
     end function lcio_three_momentum
  end interface
\eatline
\nwendcode{}\nwbegincode{901}\moddef{LCIO interface: interfaces}\plusendmoddef
  interface
     real(c_double) function lcio_energy (prt_obj) bind(C)
       import
       type(c_ptr), intent(in), value :: prt_obj
     end function lcio_energy
  end interface
\eatline
\nwendcode{}\nwbegincode{902}\moddef{LCIO interface: public}\plusendmoddef
  public :: lcio_particle_get_momentum
\nwendcode{}\nwbegincode{903}\moddef{LCIO interface: procedures}\plusendmoddef
  function lcio_particle_get_momentum (prt) result (p)
    type(vector4_t) :: p
    type(lcio_particle_t), intent(in) :: prt
    real(default) :: E, px, py, pz
    E = lcio_energy (prt%obj)
    px = lcio_three_momentum (prt%obj, 0_c_int)
    py = lcio_three_momentum (prt%obj, 1_c_int)
    pz = lcio_three_momentum (prt%obj, 2_c_int)
    p = vector4_moving ( E, vector3_moving ([ px, py, pz ]))
  end function lcio_particle_get_momentum

\eatline
\nwendcode{}\nwbegindocs{904}Return the invariant mass squared of the particle object.  LCIO
stores the signed invariant mass (no squaring).
\nwenddocs{}\nwbegincode{905}\moddef{LCIO interface: interfaces}\plusendmoddef
  interface
     function lcio_mass (prt_obj) result (mass) bind(C)
       import
       real(c_double) :: mass
       type(c_ptr), value :: prt_obj
     end function lcio_mass
  end interface
\eatline
\nwendcode{}\nwbegincode{906}\moddef{LCIO interface: public}\plusendmoddef
  public :: lcio_particle_get_mass_squared
\nwendcode{}\nwbegincode{907}\moddef{LCIO interface: procedures}\plusendmoddef
  function lcio_particle_get_mass_squared (prt) result (m2)
    real(default) :: m2
    type(lcio_particle_t), intent(in) :: prt
    real(default) :: m
    m = lcio_mass (prt%obj)
    m2 = sign (m**2, m)
  end function lcio_particle_get_mass_squared

\eatline
\nwendcode{}\nwbegindocs{908}\nwdocspar
\subsection{Polarization}
For polarization there is a three-component float entry foreseen in the LCIO 
format. Completely generic density matrices can in principle be attached to 
events as float vectors added to {\tt{}LCCollection} of the {\tt{}LCEvent}. This is 
not yet implemented currently. Here, we restrict ourselves to the same 
implementation as in HepMC format: we use two entries as the polarization 
angles, while the first entry gives the degree of polarization (something 
not specified in the HepMC format).
\emph{For massive vector bosons, we arbitrarily choose the convention
  that the longitudinal (zero) helicity state is mapped to the theta
  angle $\pi/2$.  This works under the condition that helicity is
  projected onto one of the basis states.}

\nwenddocs{}\nwbegincode{909}\moddef{LCIO interface: interfaces}\plusendmoddef
  interface
     subroutine lcio_particle_set_spin (prt_obj, s1, s2, s3) bind(C)
       import
       type(c_ptr), value :: prt_obj
       real(c_double), value :: s1, s2, s3
     end subroutine lcio_particle_set_spin
  end interface
\eatline
\nwendcode{}\nwbegindocs{910}\nwdocspar
\nwenddocs{}\nwbegincode{911}\moddef{LCIO interface: public}\plusendmoddef
  public :: lcio_polarization_init
\nwendcode{}\nwbegincode{912}\moddef{LCIO interface: interfaces}\plusendmoddef
  interface lcio_polarization_init
     module procedure lcio_polarization_init_pol
     module procedure lcio_polarization_init_hel
     module procedure lcio_polarization_init_int
  end interface
\nwendcode{}\nwbegincode{913}\moddef{LCIO interface: procedures}\plusendmoddef
  subroutine lcio_polarization_init_pol (prt, pol)
    type(lcio_particle_t), intent(out) :: prt
    type(polarization_t), intent(in) :: pol
    real(default) :: r, theta, phi
    if (polarization_is_polarized (pol)) then
       call polarization_to_angles (pol, r, theta, phi)
       call lcio_particle_set_spin (prt%obj, &
            real(r, c_double), real (theta, c_double), real (phi, c_double))
    end if
  end subroutine lcio_polarization_init_pol

  subroutine lcio_polarization_init_hel (prt, hel)
    type(lcio_particle_t), intent(out) :: prt
    type(helicity_t), intent(in) :: hel
    integer, dimension(2) :: h
    if (hel%is_defined ()) then
       h = hel%to_pair ()
       select case (h(1))
       case (1:)
          call lcio_particle_set_spin (prt%obj, 1._c_double, &
               0._c_double, 0._c_double)
       case (:-1)
          call lcio_particle_set_spin (prt%obj, 1._c_double, &
               real (pi, c_double), 0._c_double)
       case (0)
          call lcio_particle_set_spin (prt%obj, 1._c_double, &
               real (pi/2, c_double), 0._c_double)
       end select
    end if
  end subroutine lcio_polarization_init_hel

  subroutine lcio_polarization_init_int (prt, hel)
    type(lcio_particle_t), intent(out) :: prt
    integer, intent(in) :: hel
    select case (hel)
    case (1:)
       call lcio_particle_set_spin (prt%obj, 1._c_double, &
            0._c_double, 0._c_double)
    case (:-1)
       call lcio_particle_set_spin (prt%obj, 1._c_double, &
            real (pi, c_double), 0._c_double)
    case (0)
       call lcio_particle_set_spin (prt%obj, 1._c_double, &
            real (pi/2, c_double), 0._c_double)
    end select
  end subroutine lcio_polarization_init_int

\eatline
\nwendcode{}\nwbegindocs{914}Recover polarization from LCIO particle (with the
abovementioned deficiencies).
\nwenddocs{}\nwbegincode{915}\moddef{LCIO interface: interfaces}\plusendmoddef
  interface
     function lcio_polarization_degree (prt_obj) result (degree) bind(C)
       import
       real(c_double) :: degree
       type(c_ptr), value :: prt_obj
     end function lcio_polarization_degree
  end interface  
  interface
     function lcio_polarization_theta (prt_obj) result (theta) bind(C)
       import
       real(c_double) :: theta
       type(c_ptr), value :: prt_obj
     end function lcio_polarization_theta
  end interface
  interface
     function lcio_polarization_phi (prt_obj) result (phi) bind(C)
       import
       real(c_double) :: phi
       type(c_ptr), value :: prt_obj
     end function lcio_polarization_phi
  end interface
\eatline
\nwendcode{}\nwbegincode{916}\moddef{LCIO interface: public}\plusendmoddef
  public :: lcio_particle_to_pol
\nwendcode{}\nwbegincode{917}\moddef{LCIO interface: procedures}\plusendmoddef
  subroutine lcio_particle_to_pol (prt, flv, pol)
    type(lcio_particle_t), intent(in) :: prt
    type(flavor_t), intent(in) :: flv
    type(polarization_t), intent(out) :: pol
    real(default) :: degree, theta, phi
    degree = lcio_polarization_degree (prt%obj)
    theta = lcio_polarization_theta (prt%obj)
    phi = lcio_polarization_phi (prt%obj)
    call polarization_init_angles (pol, flv, degree, theta, phi)
  end subroutine lcio_particle_to_pol

\eatline
\nwendcode{}\nwbegindocs{918}Recover helicity.  Here, $\phi$ and {\tt{}degree} is ignored and only the sign of 
$\cos\theta$ is relevant, mapped to positive/negative helicity.
\nwenddocs{}\nwbegincode{919}\moddef{LCIO interface: public}\plusendmoddef
  public :: lcio_particle_to_hel
\nwendcode{}\nwbegincode{920}\moddef{LCIO interface: procedures}\plusendmoddef
  subroutine lcio_particle_to_hel (prt, flv, hel)
    type(lcio_particle_t), intent(in) :: prt
    type(flavor_t), intent(in) :: flv
    type(helicity_t), intent(out) :: hel
    real(default) :: theta
    integer :: hmax
    theta = lcio_polarization_theta (prt%obj)
    hmax = flv%get_spin_type () / 2
    call hel%init (sign (hmax, nint (cos (theta))))
  end subroutine lcio_particle_to_hel

\eatline
\nwendcode{}\nwbegindocs{921}\nwdocspar
\nwenddocs{}\nwbegincode{922}\moddef{LCIO interface: interfaces}\plusendmoddef
  interface
     subroutine lcio_particle_add_parent (prt_obj1, prt_obj2) bind(C)
       import
       type(c_ptr), value :: prt_obj1, prt_obj2
     end subroutine lcio_particle_add_parent
  end interface  
\eatline
\nwendcode{}\nwbegincode{923}\moddef{LCIO interface: public}\plusendmoddef
  public :: lcio_particle_set_parent
\nwendcode{}\nwbegincode{924}\moddef{LCIO interface: procedures}\plusendmoddef
  subroutine lcio_particle_set_parent (daughter, parent) 
    type(lcio_particle_t), intent(inout) :: daughter, parent
    call lcio_particle_add_parent (daughter%obj, parent%obj)
  end subroutine lcio_particle_set_parent

\eatline
\nwendcode{}\nwbegindocs{925}\nwdocspar
\nwenddocs{}\nwbegincode{926}\moddef{LCIO interface: interfaces}\plusendmoddef
  interface
     integer(c_int) function lcio_particle_get_generator_status &
          (prt_obj) bind(C)
       import
       type(c_ptr), value :: prt_obj
     end function lcio_particle_get_generator_status
  end interface
\eatline
\nwendcode{}\nwbegincode{927}\moddef{LCIO interface: public}\plusendmoddef
  public :: lcio_particle_get_status  
\nwendcode{}\nwbegincode{928}\moddef{LCIO interface: procedures}\plusendmoddef
  function lcio_particle_get_status (lptr) result (status)
    integer :: status
    type(lcio_particle_t), intent(in) :: lptr
    status = lcio_particle_get_generator_status (lptr%obj)
  end function lcio_particle_get_status
  
\eatline
\nwendcode{}\nwbegindocs{929}Getting the PDG code.
\nwenddocs{}\nwbegincode{930}\moddef{LCIO interface: interfaces}\plusendmoddef
  interface
     integer(c_int) function lcio_particle_get_pdg_code (prt_obj) bind(C)
       import
       type(c_ptr), value :: prt_obj
     end function lcio_particle_get_pdg_code
  end interface
\eatline
\nwendcode{}\nwbegindocs{931}\nwdocspar
\nwenddocs{}\nwbegincode{932}\moddef{LCIO interface: public}\plusendmoddef
  public :: lcio_particle_get_pdg  
\nwendcode{}\nwbegincode{933}\moddef{LCIO interface: procedures}\plusendmoddef
  function lcio_particle_get_pdg (lptr) result (pdg)
    integer :: pdg
    type(lcio_particle_t), intent(in) :: lptr
    pdg = lcio_particle_get_pdg_code (lptr%obj)
  end function lcio_particle_get_pdg
  
\eatline
\nwendcode{}\nwbegindocs{934}Obtaining the number of parents and daughters of an LCIO particle.
\nwenddocs{}\nwbegincode{935}\moddef{LCIO interface: interfaces}\plusendmoddef
  interface
     integer(c_int) function lcio_n_parents (prt_obj) bind(C)
       import
       type(c_ptr), value :: prt_obj
     end function lcio_n_parents
  end interface
\eatline
\nwendcode{}\nwbegindocs{936}\nwdocspar
\nwenddocs{}\nwbegincode{937}\moddef{LCIO interface: interfaces}\plusendmoddef
  interface
     integer(c_int) function lcio_n_daughters (prt_obj) bind(C)
       import
       type(c_ptr), value :: prt_obj
     end function lcio_n_daughters
  end interface
\eatline
\nwendcode{}\nwbegindocs{938}\nwdocspar
\nwenddocs{}\nwbegincode{939}\moddef{LCIO interface: public}\plusendmoddef
  public :: lcio_particle_get_n_parents 
\nwendcode{}\nwbegincode{940}\moddef{LCIO interface: procedures}\plusendmoddef
  function lcio_particle_get_n_parents (lptr) result (n_parents)
    integer :: n_parents
    type(lcio_particle_t), intent(in) :: lptr
    n_parents = lcio_n_parents (lptr%obj)
  end function lcio_particle_get_n_parents
  
\eatline
\nwendcode{}\nwbegindocs{941}\nwdocspar
\nwenddocs{}\nwbegincode{942}\moddef{LCIO interface: public}\plusendmoddef
  public :: lcio_particle_get_n_children
\nwendcode{}\nwbegincode{943}\moddef{LCIO interface: procedures}\plusendmoddef
  function lcio_particle_get_n_children (lptr) result (n_children)
    integer :: n_children
    type(lcio_particle_t), intent(in) :: lptr
    n_children = lcio_n_daughters (lptr%obj)
  end function lcio_particle_get_n_children
  
\eatline
\nwendcode{}\nwbegindocs{944}This provides access from the LCIO event {\tt{}lcio{\char95}event{\char95}t} to the array entries 
of the parent and daughter arrays of the LCIO particles.
\nwenddocs{}\nwbegincode{945}\moddef{LCIO interface: interfaces}\plusendmoddef
  interface
     integer(c_int) function lcio_event_parent_k &
          (evt_obj, num_part, k_parent) bind (C)
       use iso_c_binding !NODEP!
       type(c_ptr), value :: evt_obj
       integer(c_int), value :: num_part, k_parent
     end function lcio_event_parent_k
  end interface
\eatline
\nwendcode{}\nwbegincode{946}\moddef{LCIO interface: interfaces}\plusendmoddef
  interface
     integer(c_int) function lcio_event_daughter_k &
          (evt_obj, num_part, k_daughter) bind (C)
       use iso_c_binding !NODEP!
       type(c_ptr), value :: evt_obj
       integer(c_int), value :: num_part, k_daughter
     end function lcio_event_daughter_k
  end interface
\eatline
\nwendcode{}\nwbegindocs{947}\nwdocspar
\nwenddocs{}\nwbegincode{948}\moddef{LCIO interface: public}\plusendmoddef
  public :: lcio_get_n_parents
\nwendcode{}\nwbegincode{949}\moddef{LCIO interface: procedures}\plusendmoddef
  function lcio_get_n_parents (evt, num_part, k_parent) result (index_parent)
    type(lcio_event_t), intent(in) :: evt
    integer, intent(in) :: num_part, k_parent
    integer :: index_parent
    index_parent = lcio_event_parent_k (evt%obj, int (num_part, c_int), &
         int (k_parent, c_int))
  end function lcio_get_n_parents

\eatline
\nwendcode{}\nwbegindocs{950}\nwdocspar
\nwenddocs{}\nwbegincode{951}\moddef{LCIO interface: public}\plusendmoddef
  public :: lcio_get_n_children
\nwendcode{}\nwbegincode{952}\moddef{LCIO interface: procedures}\plusendmoddef
  function lcio_get_n_children (evt, num_part, k_daughter) result (index_daughter)
    type(lcio_event_t), intent(in) :: evt
    integer, intent(in) :: num_part, k_daughter
    integer :: index_daughter
    index_daughter = lcio_event_daughter_k (evt%obj, int (num_part, c_int), &
         int (k_daughter, c_int))
  end function lcio_get_n_children

\eatline
\nwendcode{}\nwbegindocs{953}\nwdocspar
\subsection{LCIO Writer type}
There is a specific LCIO Writer type for handling the output of
LCEventImpl objects (i.e., Monte Carlo event samples) to file.  Opening
the file is done by the constructor, closing by the destructor.
\nwenddocs{}\nwbegincode{954}\moddef{LCIO interface: public}\plusendmoddef
  public :: lcio_writer_t
\nwendcode{}\nwbegincode{955}\moddef{LCIO interface: types}\plusendmoddef
  type :: lcio_writer_t
     private
     type(c_ptr) :: obj
  end type lcio_writer_t

\eatline
\nwendcode{}\nwbegindocs{956}Constructor for an output associated to a file.
\nwenddocs{}\nwbegincode{957}\moddef{LCIO interface: interfaces}\plusendmoddef
  interface
     type(c_ptr) function open_lcio_writer_new (filename, complevel) bind(C)
       import
       character(c_char), dimension(*), intent(in) :: filename
       integer(c_int), intent(in) :: complevel
     end function open_lcio_writer_new
  end interface
\eatline
\nwendcode{}\nwbegincode{958}\moddef{LCIO interface: public}\plusendmoddef
  public :: lcio_writer_open_out
\nwendcode{}\nwbegincode{959}\moddef{LCIO interface: procedures}\plusendmoddef
  subroutine lcio_writer_open_out (lcio_writer, filename)
    type(lcio_writer_t), intent(out) :: lcio_writer
    type(string_t), intent(in) :: filename
    lcio_writer%obj = open_lcio_writer_new (char (filename) // &
         c_null_char, 9_c_int)
  end subroutine lcio_writer_open_out

\eatline
\nwendcode{}\nwbegindocs{960}Destructor:
\nwenddocs{}\nwbegincode{961}\moddef{LCIO interface: interfaces}\plusendmoddef
  interface
     subroutine lcio_writer_delete (io_obj) bind(C)
       import
       type(c_ptr), value :: io_obj
     end subroutine lcio_writer_delete
  end interface
\eatline
\nwendcode{}\nwbegincode{962}\moddef{LCIO interface: public}\plusendmoddef
  public :: lcio_writer_close
\nwendcode{}\nwbegincode{963}\moddef{LCIO interface: procedures}\plusendmoddef
  subroutine lcio_writer_close (lciowriter)
    type(lcio_writer_t), intent(inout) :: lciowriter
    call lcio_writer_delete (lciowriter%obj)
  end subroutine lcio_writer_close

\eatline
\nwendcode{}\nwbegindocs{964}Write a single event to the LCIO writer.
\nwenddocs{}\nwbegincode{965}\moddef{LCIO interface: interfaces}\plusendmoddef
  interface
     subroutine lcio_write_event (io_obj, evt_obj) bind(C)
       import
       type(c_ptr), value :: io_obj, evt_obj
     end subroutine lcio_write_event
  end interface
\eatline
\nwendcode{}\nwbegincode{966}\moddef{LCIO interface: public}\plusendmoddef
  public :: lcio_event_write
\nwendcode{}\nwbegincode{967}\moddef{LCIO interface: procedures}\plusendmoddef
  subroutine lcio_event_write (wrt, evt)
    type(lcio_writer_t), intent(inout) :: wrt
    type(lcio_event_t), intent(in) :: evt
    call lcio_write_event (wrt%obj, evt%obj)
  end subroutine lcio_event_write

\eatline
\nwendcode{}\nwbegindocs{968}\nwdocspar
\subsection{LCIO Reader type}
There is a specific LCIO Reader type for handling the input of
LCEventImpl objects (i.e., Monte Carlo event samples) from file.  Opening
the file is done by the constructor, closing by the destructor.
\nwenddocs{}\nwbegincode{969}\moddef{LCIO interface: public}\plusendmoddef
  public :: lcio_reader_t
\nwendcode{}\nwbegincode{970}\moddef{LCIO interface: types}\plusendmoddef
  type :: lcio_reader_t
     private
     type(c_ptr) :: obj
  end type lcio_reader_t

\eatline
\nwendcode{}\nwbegindocs{971}Constructor for an output associated to a file.
\nwenddocs{}\nwbegincode{972}\moddef{LCIO interface: interfaces}\plusendmoddef
  interface
     type(c_ptr) function open_lcio_reader (filename) bind(C)
       import
       character(c_char), dimension(*), intent(in) :: filename
     end function open_lcio_reader
  end interface
\eatline
\nwendcode{}\nwbegincode{973}\moddef{LCIO interface: public}\plusendmoddef
  public :: lcio_open_file
\nwendcode{}\nwbegincode{974}\moddef{LCIO interface: procedures}\plusendmoddef
  subroutine lcio_open_file (lcio_reader, filename)
    type(lcio_reader_t), intent(out) :: lcio_reader
    type(string_t), intent(in) :: filename
    lcio_reader%obj = open_lcio_reader (char (filename) // c_null_char)
  end subroutine lcio_open_file

\eatline
\nwendcode{}\nwbegindocs{975}Destructor:
\nwenddocs{}\nwbegincode{976}\moddef{LCIO interface: interfaces}\plusendmoddef
  interface
     subroutine lcio_reader_delete (io_obj) bind(C)
       import
       type(c_ptr), value :: io_obj
     end subroutine lcio_reader_delete
  end interface
\eatline
\nwendcode{}\nwbegincode{977}\moddef{LCIO interface: public}\plusendmoddef
  public :: lcio_reader_close
\nwendcode{}\nwbegincode{978}\moddef{LCIO interface: procedures}\plusendmoddef
  subroutine lcio_reader_close (lcioreader)
    type(lcio_reader_t), intent(inout) :: lcioreader
    call lcio_reader_delete (lcioreader%obj)
  end subroutine lcio_reader_close

\eatline
\nwendcode{}\nwbegindocs{979}\nwdocspar
\nwenddocs{}\nwbegindocs{980}Read a single event from the event file.  Return true if successful.
\nwenddocs{}\nwbegincode{981}\moddef{LCIO interface: interfaces}\plusendmoddef
  interface
     type(c_ptr) function read_lcio_event (io_obj) bind(C)
       import
       type(c_ptr), value :: io_obj
     end function read_lcio_event
  end interface
\eatline
\nwendcode{}\nwbegincode{982}\moddef{LCIO interface: public}\plusendmoddef
  public :: lcio_read_event
\nwendcode{}\nwbegincode{983}\moddef{LCIO interface: procedures}\plusendmoddef
  subroutine lcio_read_event (lcrdr, evt, ok)
    type(lcio_reader_t), intent(inout) :: lcrdr
    type(lcio_event_t), intent(out) :: evt
    logical, intent(out) :: ok
    evt%obj = read_lcio_event (lcrdr%obj)
    ok = c_associated (evt%obj)
  end subroutine lcio_read_event

\eatline
\nwendcode{}\nwbegindocs{984}Extract the process ID. This is stored (at the moment abusively) in the 
RUN ID as well as in an additional event parameter.
\nwenddocs{}\nwbegincode{985}\moddef{LCIO interface: interfaces}\plusendmoddef
  interface
     integer(c_int) function lcio_event_signal_process_id (evt_obj) bind(C)
       import
       type(c_ptr), value :: evt_obj
     end function lcio_event_signal_process_id
  end interface
\eatline
\nwendcode{}\nwbegincode{986}\moddef{LCIO interface: public}\plusendmoddef
  public :: lcio_event_get_process_id
\nwendcode{}\nwbegincode{987}\moddef{LCIO interface: procedures}\plusendmoddef
  function lcio_event_get_process_id (evt) result (i_proc)
    integer :: i_proc
    type(lcio_event_t), intent(in) :: evt
    i_proc = lcio_event_signal_process_id (evt%obj)
  end function lcio_event_get_process_id

\eatline
\nwendcode{}\nwbegindocs{988}Number of particles in an LCIO event.
\nwenddocs{}\nwbegincode{989}\moddef{LCIO interface: interfaces}\plusendmoddef
  interface
     integer(c_int) function lcio_event_get_n_particles (evt_obj) bind(C)
       import
       type(c_ptr), value :: evt_obj
     end function lcio_event_get_n_particles
  end interface
\eatline
\nwendcode{}\nwbegincode{990}\moddef{LCIO interface:}\endmoddef
\nwendcode{}\nwbegindocs{991}\nwdocspar
\nwenddocs{}\nwbegincode{992}\moddef{LCIO interface: public}\plusendmoddef
  public :: lcio_event_get_n_tot
\nwendcode{}\nwbegincode{993}\moddef{LCIO interface: procedures}\plusendmoddef
  function lcio_event_get_n_tot (evt) result (n_tot)
    integer :: n_tot
    type(lcio_event_t), intent(in) :: evt
    n_tot = lcio_event_get_n_particles (evt%obj)
  end function lcio_event_get_n_tot

\eatline
\nwendcode{}\nwbegindocs{994}Extracting $\alpha_s$ and the scale.
\nwenddocs{}\nwbegincode{995}\moddef{LCIO interface: interfaces}\plusendmoddef
  interface
     function lcio_event_get_alpha_qcd (evt_obj) result (as) bind(C)
       import
       real(c_double) :: as
       type(c_ptr), value :: evt_obj
     end function lcio_event_get_alpha_qcd
  end interface
  interface
     function lcio_event_get_scale (evt_obj) result (scale) bind(C)
       import 
       real(c_double) :: scale
       type(c_ptr), value :: evt_obj
     end function lcio_event_get_scale
  end interface
\eatline
\nwendcode{}\nwbegindocs{996}\nwdocspar
\nwenddocs{}\nwbegincode{997}\moddef{LCIO interface: public}\plusendmoddef
  public :: lcio_event_get_alphas
\nwendcode{}\nwbegincode{998}\moddef{LCIO interface: procedures}\plusendmoddef
  function lcio_event_get_alphas (evt) result (as)
    type(lcio_event_t), intent(in) :: evt
    real(default) :: as
    as = lcio_event_get_alpha_qcd (evt%obj)
  end function lcio_event_get_alphas

\eatline
\nwendcode{}\nwbegindocs{999}\nwdocspar
\nwenddocs{}\nwbegincode{1000}\moddef{LCIO interface: public}\plusendmoddef
  public :: lcio_event_get_scaleval  
\nwendcode{}\nwbegincode{1001}\moddef{LCIO interface: procedures}\plusendmoddef
  function lcio_event_get_scaleval (evt) result (scale)
    type(lcio_event_t), intent(in) :: evt
    real(default) :: scale
    scale = lcio_event_get_scale (evt%obj)
  end function lcio_event_get_scaleval

\eatline
\nwendcode{}\nwbegindocs{1002}Extracting particles by index from an LCIO event.
\nwenddocs{}\nwbegincode{1003}\moddef{LCIO interface: interfaces}\plusendmoddef
  interface
     type(c_ptr) function lcio_event_particle_k (evt_obj, k) bind(C)
       import
       type(c_ptr), value :: evt_obj
       integer(c_int), value :: k
     end function lcio_event_particle_k
  end interface
\eatline
\nwendcode{}\nwbegindocs{1004}\nwdocspar
\nwenddocs{}\nwbegincode{1005}\moddef{LCIO interface: public}\plusendmoddef
  public :: lcio_event_get_particle
\nwendcode{}\nwbegincode{1006}\moddef{LCIO interface: procedures}\plusendmoddef
  function lcio_event_get_particle (evt, n) result (prt)
    type(lcio_event_t), intent(in) :: evt
    integer, intent(in) :: n
    type(lcio_particle_t) :: prt
    prt%obj = lcio_event_particle_k (evt%obj, int (n, c_int))
  end function lcio_event_get_particle
\eatline
\nwendcode{}\nwbegindocs{1007}\nwdocspar
\subsection{Unit tests}
This test example is an abridged version from the build-from-scratch
example in the HepMC distribution, translated to LCIO. 
There are three different branchings for $p\to q$ PDF splitting, then a
branching for a $qq\to W^-g$ hard-interaction process, and finally a 
branching for $W^-\to qq$ decay.  The setup is for LHC kinematics.

Extending the original example, we set color flow for the incoming
quarks and polarization for the outgoing photon.  For the latter, we
have to define a particle-data object for the photon, so a flavor
object can be correctly initialized.
\nwenddocs{}\nwbegincode{1008}\moddef{LCIO interface: public}\plusendmoddef
  public :: lcio_test
\nwendcode{}\nwbegincode{1009}\moddef{LCIO interface: procedures}\plusendmoddef
  subroutine lcio_test (u, results)
    integer, intent(in) :: u
    type(test_results_t), intent(inout) :: results
  \LA{}LCIO interface: execute tests\RA{}
  end subroutine lcio_test
  
\nwendcode{}\nwbegindocs{1010} %def lcio_test
\nwenddocs{}\nwbegindocs{1011}\nwdocspar
\nwenddocs{}\nwbegincode{1012}\moddef{LCIO interface: execute tests}\endmoddef
  call test (lcio_interface_1, "lcio_interface_1", &
       "check LCIO interface", &
       u, results)
\nwendcode{}\nwbegincode{1013}\moddef{LCIO interface: tests}\endmoddef
  subroutine lcio_interface_1 (u)
    use physics_defs, only: VECTOR
    use model_data, only: field_data_t
    integer, intent(in) :: u
    integer :: u_file, iostat
    type(lcio_event_t) :: evt
    type(lcio_particle_t) :: prt1, prt2, prt3, prt4, prt5, prt6, prt7, prt8
    type(flavor_t) :: flv
    type(color_t) :: col
    type(polarization_t) :: pol
    type(field_data_t), target :: photon_data
    character(220) :: buffer    

    write (u, "(A)")  "* Test output: LCIO interface"
    write (u, "(A)")  "*   Purpose: test LCIO interface"
    write (u, "(A)")      
    
    write (u, "(A)")  "* Initialization"
    write (u, "(A)")
    
    ! Initialize a photon flavor object and some polarization
    call photon_data%init (var_str ("PHOTON"), 22)
    call photon_data%set (spin_type=VECTOR)
    call photon_data%freeze ()
    call flv%init (photon_data)
    call polarization_init_angles &
         (pol, flv, 0.6_default, 1._default, 0.5_default)

    ! Event initialization
    call lcio_event_init (evt, 20, 1, 42)

    write (u, "(A)")  "* p -> q splitting"
    write (u, "(A)")
            
    ! $p\\to q$ splittings
    call particle_init (prt1, &
         0._default, 0._default, 7000._default, 7000._default, &
         2212, 3)
    call particle_init (prt2, &
         0._default, 0._default,-7000._default, 7000._default, &
         2212, 3)
    call particle_init (prt3, &
          .750_default, -1.569_default, 32.191_default, 32.238_default, &
          1, 3)
    call color_init_from_array (col, [501])
    call lcio_particle_set_color (prt3, col)
    call lcio_particle_set_parent (prt3, prt1)
    call lcio_particle_set_parent (prt3, prt2)
    call particle_init (prt4, &
         -3.047_default, -19._default, -54.629_default, 57.920_default, &
         -2, 3)
    call color_init_from_array (col, [-501])
    call lcio_particle_set_color (prt4, col)
    call lcio_particle_set_parent (prt4, prt1)
    call lcio_particle_set_parent (prt4, prt2)
    
    write (u, "(A)")  "* Hard interaction"
    write (u, "(A)")
        
    ! Hard interaction
    call particle_init (prt6, &
         -3.813_default, 0.113_default, -1.833_default, 4.233_default, &
         22, 1)
    call lcio_polarization_init (prt6, pol)
    call particle_init (prt5, &
         1.517_default, -20.68_default, -20.605_default, 85.925_default, &
         -24, 3)
    call lcio_particle_set_parent (prt5, prt3)
    call lcio_particle_set_parent (prt5, prt4)
    call lcio_particle_set_parent (prt6, prt3)
    call lcio_particle_set_parent (prt6, prt4)    
    
    ! $W^-$ decay
    call particle_init (prt7, &
         -2.445_default, 28.816_default, 6.082_default, 29.552_default, &
         1, 1)
    call particle_init (prt8, &
         3.962_default, -49.498_default, -26.687_default, 56.373_default, &
         -2, 1)
    call lcio_particle_set_parent (prt7, prt5)
    call lcio_particle_set_parent (prt8, prt5)
    call lcio_particle_add_to_evt_coll (prt1, evt)
    call lcio_particle_add_to_evt_coll (prt2, evt)
    call lcio_particle_add_to_evt_coll (prt3, evt)
    call lcio_particle_add_to_evt_coll (prt4, evt)
    call lcio_particle_add_to_evt_coll (prt5, evt)
    call lcio_particle_add_to_evt_coll (prt6, evt)
    call lcio_particle_add_to_evt_coll (prt7, evt)
    call lcio_particle_add_to_evt_coll (prt8, evt)    
    call lcio_event_add_coll (evt)
    
    ! Event output
    write (u, "(A)")  "Writing in ASCII form to file 'lcio_test.slcio'"
    write (u, "(A)")
    
    call write_lcio_event (evt, var_str ("lcio_test.slcio"))
    
    write (u, "(A)")  "Writing completed"
    
    write (u, "(A)")
    write (u, "(A)")  "* File contents:"
    write (u, "(A)")
    
    u_file = free_unit ()
    open (u_file, file = "lcio_test.slcio", &
         action = "read", status = "old")
    do
       read (u_file, "(A)", iostat = iostat)  buffer
       if (trim (buffer) == "")  cycle
       if (buffer(1:12) == " - timestamp")  buffer = "[...]"
       if (buffer(1:6) == " date:")  buffer = "[...]"
       if (iostat /= 0)  exit
       write (u, "(A)") trim (buffer)
    end do
    close (u_file)
    
    write (u, "(A)")
    write (u, "(A)")  "* Cleanup"
    write (u, "(A)")
        
    ! Wrapup
    call polarization_final (pol)
    call lcio_event_final (evt)
    
    write (u, "(A)")
    write (u, "(A)")  "* Test output end: lcio_interface_1"
    
  contains

    subroutine particle_init (prt, px, py, pz, E, pdg, status)
      type(lcio_particle_t), intent(out) :: prt
      real(default), intent(in) :: px, py, pz, E
      integer, intent(in) :: pdg, status
      type(vector4_t) :: p
      p = vector4_moving (E, vector3_moving ([px, py, pz]))
      call lcio_particle_init (prt, p, pdg, status)
    end subroutine particle_init

  end subroutine lcio_interface_1

\eatline
\nwendcode{}\nwbegindocs{1014}\nwdocspar
%%%%%%%%%%%%%%%%%%%%%%%%%%%%%%%%%%%%%%%%%%%%%%%%%%%%%%%%%%%%%%%%%%%%%%%%
\section{HEP Common and Events}

This is a separate module that manages data exchange between the common blocks
and {\tt{}event{\char95}t} objects. We separate this from the previous module in order 
to avoid a circular module dependency. It also contains the functions
necessary for communication between {\tt{}hepmc{\char95}event{\char95}t} and
{\tt{}event{\char95}t} or {\tt{}lcio{\char95}event{\char95}t} and {\tt{}event{\char95}t} as well as
{\tt{}particle{\char95}set{\char95}t} and {\tt{}particle{\char95}t} objects.  
\nwenddocs{}\nwbegincode{1015}\moddef{\code{}hep{\char95}events.f90\edoc{}}\endmoddef
\LA{}File header\RA{}

module hep_events
  
  use kinds
  use io_units
\LA{}Use strings\RA{}  
  use diagnostics
  use lorentz
  use unit_tests
  
  use flavors
  use colors
  use helicities
  use quantum_numbers
  use interactions
  use state_matrices
  use evaluators
  use polarizations
  use model_data
  use subevents
  use particles
  use hep_common
  use hepmc_interface
  use lcio_interface
  use event_base

\LA{}Standard module head\RA{}

\LA{}HEP events: public\RA{}

contains
  
\LA{}HEP events: procedures\RA{}

\LA{}HEP events: tests\RA{}

end module hep_events
\eatline
\nwendcode{}\nwbegindocs{1016}\nwdocspar
\subsection{Data Transfer: events}
Fill the HEPEUP block, given a \whizard\ event object.
\nwenddocs{}\nwbegincode{1017}\moddef{HEP events: public}\endmoddef
  public :: hepeup_from_event
\nwendcode{}\nwbegincode{1018}\moddef{HEP events: procedures}\endmoddef
  subroutine hepeup_from_event &
       (event, keep_beams, keep_remnants, process_index)
    class(generic_event_t), intent(in), target :: event
    logical, intent(in), optional :: keep_beams
    logical, intent(in), optional :: keep_remnants
    integer, intent(in), optional :: process_index
    type(particle_set_t), pointer :: particle_set
    real(default) :: scale, alpha_qcd
    if (event%has_valid_particle_set ()) then
       particle_set => event%get_particle_set_ptr ()
       call hepeup_from_particle_set (particle_set, keep_beams, keep_remnants)
       if (present (process_index)) &
            call hepeup_set_event_parameters (proc_id = process_index)
       scale = event%get_fac_scale ()
       if (scale /= 0) &
            call hepeup_set_event_parameters (scale = scale)
       alpha_qcd = event%get_alpha_s ()       
       if (alpha_qcd /= 0) &
            call hepeup_set_event_parameters (alpha_qcd = alpha_qcd)
       if (event%weight_prc_is_known ()) then
          call hepeup_set_event_parameters (weight = event%get_weight_prc ())
       else
          call msg_bug ("HEPEUP: process weight is unknown")
       end if
    else
       call msg_bug ("HEPEUP: event incomplete")
    end if
  end subroutine hepeup_from_event

\eatline
\nwendcode{}\nwbegindocs{1019}Reverse.  

Note: The current implementation sets the particle set of the hard
process and is therefore not useful if the event on file is dressed.
This should be reconsidered.

Note: setting of scale or alpha is not yet supported by the
{\tt{}event{\char95}t} object.  Ticket \#628.
\nwenddocs{}\nwbegincode{1020}\moddef{HEP events: public}\plusendmoddef
  public :: hepeup_to_event
\nwendcode{}\nwbegincode{1021}\moddef{HEP events: procedures}\plusendmoddef
  subroutine hepeup_to_event &
       (event, fallback_model, process_index, recover_beams, &
       use_alpha_s, use_scale)
    class(generic_event_t), intent(inout), target :: event
    class(model_data_t), intent(in), target :: fallback_model
    integer, intent(out), optional :: process_index
    logical, intent(in), optional :: recover_beams
    logical, intent(in), optional :: use_alpha_s
    logical, intent(in), optional :: use_scale
    class(model_data_t), pointer :: model
    real(default) :: weight, scale, alpha_qcd
    type(particle_set_t) :: particle_set
    model => event%get_model_ptr ()
    call hepeup_to_particle_set &
         (particle_set, recover_beams, model, fallback_model)
    call event%set_hard_particle_set (particle_set)
    call particle_set%final ()
    if (present (process_index)) then
       call hepeup_get_event_parameters (proc_id = process_index)
    end if
    call hepeup_get_event_parameters (weight = weight, &
         scale = scale, alpha_qcd = alpha_qcd)
    call event%set_weight_ref (weight)
    if (present (use_alpha_s)) then
       if (use_alpha_s .and. alpha_qcd > 0) &
            call event%set_alpha_qcd_forced (alpha_qcd)
    end if
    if (present (use_scale)) then
       if (use_scale .and. scale > 0) &
            call event%set_scale_forced (scale)
    end if
  end subroutine hepeup_to_event

\eatline
\nwendcode{}\nwbegindocs{1022}Fill the HEPEVT (event) common block:
\nwenddocs{}\nwbegincode{1023}\moddef{HEP events: public}\plusendmoddef
  public :: hepevt_from_event
\nwendcode{}\nwbegincode{1024}\moddef{HEP events: procedures}\plusendmoddef
  subroutine hepevt_from_event (event, i_evt, keep_beams, keep_remnants)
    class(generic_event_t), intent(in), target :: event
    integer, intent(in), optional :: i_evt
    logical, intent(in), optional :: keep_beams  
    logical, intent(in), optional :: keep_remnants
    type(particle_set_t), pointer :: particle_set
    if (event%has_valid_particle_set ()) then
       particle_set => event%get_particle_set_ptr ()
       call hepevt_from_particle_set (particle_set, keep_beams, keep_remnants)
       if (event%weight_prc_is_known () .and. event%sqme_prc_is_known ()) then
          call hepevt_set_event_parameters ( &
               weight = event%get_weight_prc (), &
               function_value = event%get_sqme_prc ())
       else
          call msg_bug ("HEPEVT: event weight and/or sqme unknown")
       end if
       if (present (i_evt)) &
            call hepevt_set_event_parameters (i_evt = i_evt)
    else
       call msg_bug ("HEPEVT: event incomplete")
    end if
  end subroutine hepevt_from_event

\eatline
\nwendcode{}\nwbegindocs{1025}\subsubsection{HepMC format}
The master output function fills a HepMC GenEvent object that is
already initialized, but has no vertices in it.

We first set up the vertex lists and enter the vertices into the HepMC
event.  Then, we assign first all incoming particles and then all
outgoing particles to their associated vertices.  Particles which have
neither parent nor children entries (this should not happen) are
dropped.

Finally, we insert the beam particles.  If there are none, use the incoming
particles instead.
\nwenddocs{}\nwbegindocs{1026}Transform a particle into a {\tt{}hepmc{\char95}particle} object, including
color and polarization.  The HepMC status is equivalent to the HEPEVT
status, in particular: 0 = null entry, 1 = physical particle, 2 =
decayed/fragmented SM hadron, tau or muon, 3 = other unphysical particle 
entry, 4 = incoming particles, 11 = intermediate resonance such as squarks.
The use of 11 for intermediate resonances is as done by HERWIG, see 
http://herwig.hepforge.org/trac/wiki/FaQs.
\nwenddocs{}\nwbegincode{1027}\moddef{HEP events: procedures}\plusendmoddef
  subroutine particle_to_hepmc (prt, hprt)
    type(particle_t), intent(in) :: prt
    type(hepmc_particle_t), intent(out) :: hprt
    integer :: hepmc_status
    select case (prt%get_status ())
    case (PRT_UNDEFINED)
       hepmc_status = 0
    case (PRT_OUTGOING)
       hepmc_status = 1
    case (PRT_BEAM)
       hepmc_status = 4
    case (PRT_RESONANT)
       if (abs(prt%get_pdg()) == 13 .or. &
            abs(prt%get_pdg()) == 15) then
          hepmc_status = 2
       else
          hepmc_status = 11
       end if
    case default
       hepmc_status = 3
    end select
    call hepmc_particle_init (hprt, &
         prt%get_momentum (), prt%get_pdg (), &
         hepmc_status)
    call hepmc_particle_set_color (hprt, prt%get_color ())
    select case (prt%get_polarization_status ())
    case (PRT_DEFINITE_HELICITY)
       call hepmc_particle_set_polarization (hprt, &
            prt%get_helicity ())
    case (PRT_GENERIC_POLARIZATION)
       call hepmc_particle_set_polarization (hprt, &
            prt%get_polarization ())
    end select
  end subroutine particle_to_hepmc

\eatline
\nwendcode{}\nwbegindocs{1028}\nwdocspar
\nwenddocs{}\nwbegincode{1029}\moddef{HEP events: public}\plusendmoddef
  public :: hepmc_event_from_particle_set 
\nwendcode{}\nwbegincode{1030}\moddef{HEP events: procedures}\plusendmoddef
  subroutine hepmc_event_from_particle_set (evt, particle_set)
    type(hepmc_event_t), intent(inout) :: evt
    type(particle_set_t), intent(in) :: particle_set
    type(hepmc_vertex_t), dimension(:), allocatable :: v
    type(hepmc_particle_t), dimension(:), allocatable :: hprt
    type(hepmc_particle_t), dimension(2) :: hbeam
    logical, dimension(:), allocatable :: is_beam
    integer, dimension(:), allocatable :: v_from, v_to
    integer :: n_vertices, n_tot, i
    n_tot = particle_set%get_n_tot ()
    allocate (v_from (n_tot), v_to (n_tot))
    call particle_set%assign_vertices (v_from, v_to, n_vertices)
    allocate (v (n_vertices))
    do i = 1, n_vertices
       call hepmc_vertex_init (v(i))
       call hepmc_event_add_vertex (evt, v(i))
    end do
    allocate (hprt (n_tot))
    do i = 1, n_tot
       if (v_to(i) /= 0 .or. v_from(i) /= 0) then
          call particle_to_hepmc (particle_set%prt(i), hprt(i))
       end if
    end do
    allocate (is_beam (n_tot))
    is_beam = particle_set%prt(1:n_tot)%get_status () == PRT_BEAM
    if (.not. any (is_beam)) then
       is_beam = particle_set%prt(1:n_tot)%get_status () == PRT_INCOMING
    end if
    if (count (is_beam) == 2) then
       hbeam = pack (hprt, is_beam)
       call hepmc_event_set_beam_particles (evt, hbeam(1), hbeam(2))
    end if
    do i = 1, n_tot
       if (v_to(i) /= 0) then
          call hepmc_vertex_add_particle_in (v(v_to(i)), hprt(i))
       end if
    end do
    do i = 1, n_tot
       if (v_from(i) /= 0) then
          call hepmc_vertex_add_particle_out (v(v_from(i)), hprt(i))
       end if
    end do
    FIND_SIGNAL_PROCESS: do i = 1, n_tot
       if (particle_set%prt(i)%get_status () == PRT_INCOMING) then
          call hepmc_event_set_signal_process_vertex (evt, v(v_to(i)))
          exit FIND_SIGNAL_PROCESS
       end if
    end do FIND_SIGNAL_PROCESS
  end subroutine hepmc_event_from_particle_set
  
\eatline
\nwendcode{}\nwbegindocs{1031}Initialize a particle from a HepMC particle object.  The model is
necessary for making a fully qualified flavor component.  We have the
additional flag {\tt{}polarized} which tells whether the polarization
information should be interpreted or ignored, and the lookup array of
barcodes.  Note that the lookup array is searched linearly, a possible
bottleneck for large particle arrays.  If necessary, the barcode array
could be replaced by a hash table.
\nwenddocs{}\nwbegincode{1032}\moddef{HEP events: procedures}\plusendmoddef
  subroutine particle_from_hepmc_particle &
       (prt, hprt, model, polarization, barcode)
    type(particle_t), intent(out) :: prt
    type(hepmc_particle_t), intent(in) :: hprt
    type(model_data_t), intent(in), target :: model
    integer, intent(in) :: polarization
    integer, dimension(:), intent(in) :: barcode
    type(hepmc_polarization_t) :: hpol
    type(flavor_t) :: flv
    type(color_t) :: col
    type(helicity_t) :: hel
    type(polarization_t) :: pol
    integer :: n_parents, n_children
    integer, dimension(:), allocatable :: &
         parent_barcode, child_barcode, parent, child
    integer :: i
    select case (hepmc_particle_get_status (hprt))
    case (1);  call prt%set_status (PRT_OUTGOING)
    case (2);  call prt%set_status (PRT_RESONANT)
    case (3);  call prt%set_status (PRT_VIRTUAL)
    end select
    if (hepmc_particle_is_beam (hprt)) call prt%set_status (PRT_BEAM)
    call flv%init (hepmc_particle_get_pdg (hprt), model)
    call col%init (hepmc_particle_get_color (hprt))
    call prt%set_flavor (flv)   
    call prt%set_color (col)
    call prt%set_polarization (polarization)
    select case (polarization)
    case (PRT_DEFINITE_HELICITY)
       hpol = hepmc_particle_get_polarization (hprt)
       call hepmc_polarization_to_hel (hpol, prt%get_flv (), hel)
       call prt%set_helicity (hel)
       call hepmc_polarization_final (hpol)
    case (PRT_GENERIC_POLARIZATION)
       hpol = hepmc_particle_get_polarization (hprt)
       call hepmc_polarization_to_pol (hpol, prt%get_flv (), pol)
       call prt%set_pol (pol)
       call hepmc_polarization_final (hpol)
    end select
    call prt%set_momentum (hepmc_particle_get_momentum (hprt), &
         hepmc_particle_get_mass_squared (hprt))
    n_parents  = hepmc_particle_get_n_parents  (hprt)
    n_children = hepmc_particle_get_n_children (hprt)
    allocate (parent_barcode (n_parents),  parent (n_parents))
    allocate (child_barcode  (n_children), child  (n_children))
    parent_barcode = hepmc_particle_get_parent_barcodes (hprt)
    child_barcode  = hepmc_particle_get_child_barcodes  (hprt)
    do i = 1, size (barcode)
       where (parent_barcode == barcode(i))  parent = i
       where (child_barcode  == barcode(i))  child  = i
    end do
    call prt%set_parents (parent)
    call prt%set_children (child)
    if (prt%get_status () == PRT_VIRTUAL .and. n_parents == 0) &
         call prt%set_status (PRT_INCOMING)
  end subroutine particle_from_hepmc_particle

\eatline
\nwendcode{}\nwbegindocs{1033}If a particle set is initialized from a HepMC event record, we have
to specify the treatment of polarization (unpolarized or density
matrix) which is common to all particles.  Correlated polarization
information is not available.

There is some complication in reconstructing incoming particles and
beam remnants.  First of all, they all will be tagged as virtual.  We
then define an incoming particle as 
\nwenddocs{}\nwbegincode{1034}\moddef{HEP events: procedures}\plusendmoddef
  subroutine hepmc_event_to_particle_set &
       (particle_set, evt, model, fallback_model, polarization)
    type(particle_set_t), intent(inout), target :: particle_set
    type(hepmc_event_t), intent(in) :: evt
    class(model_data_t), intent(in), target :: model, fallback_model
    integer, intent(in) :: polarization
    type(hepmc_event_particle_iterator_t) :: it
    type(hepmc_vertex_t) :: v
    type(hepmc_vertex_particle_in_iterator_t) :: v_it
    type(hepmc_particle_t) :: prt
    integer, dimension(:), allocatable :: barcode
    integer :: n_tot, i, j, bc
    logical :: has_parents, has_children
    n_tot = 0
    call hepmc_event_particle_iterator_init (it, evt)
    do while (hepmc_event_particle_iterator_is_valid (it))
       n_tot = n_tot + 1
       call hepmc_event_particle_iterator_advance (it)
    end do
    allocate (barcode (n_tot))
    call hepmc_event_particle_iterator_reset (it)
    do i = 1, n_tot
       barcode(i) = hepmc_particle_get_barcode &
            (hepmc_event_particle_iterator_get (it))
       call hepmc_event_particle_iterator_advance (it)
    end do
    allocate (particle_set%prt (n_tot))
    call hepmc_event_particle_iterator_reset (it)
    do i = 1, n_tot
       prt = hepmc_event_particle_iterator_get (it)
       call particle_from_hepmc_particle (particle_set%prt(i), &
            prt, model, polarization, barcode)
       call hepmc_event_particle_iterator_advance (it)
    end do
    call hepmc_event_particle_iterator_final (it)
    v = hepmc_event_get_signal_process_vertex (evt)
    if (hepmc_vertex_is_valid (v)) then
       call hepmc_vertex_particle_in_iterator_init (v_it, v)
       do while (hepmc_vertex_particle_in_iterator_is_valid (v_it))
          prt = hepmc_vertex_particle_in_iterator_get (v_it)
          bc = hepmc_particle_get_barcode &
               (hepmc_vertex_particle_in_iterator_get (v_it))
          do i = 1, size(barcode)
             if (bc == barcode(i))  &
               call particle_set%prt(i)%set_status (PRT_INCOMING)
          end do
          call hepmc_vertex_particle_in_iterator_advance (v_it)
       end do
       call hepmc_vertex_particle_in_iterator_final (v_it)
    end if
    do i = 1, n_tot
       if (particle_set%prt(i)%get_status () == PRT_VIRTUAL &
            .and. particle_set%prt(i)%get_n_children () == 0) &
            call particle_set%prt(i)%set_status (PRT_OUTGOING)
    end do
    particle_set%n_tot = n_tot
    particle_set%n_beam  = &
         count (particle_set%prt%get_status () == PRT_BEAM)
    particle_set%n_in  = &
         count (particle_set%prt%get_status () == PRT_INCOMING)
    particle_set%n_out = &
         count (particle_set%prt%get_status () == PRT_OUTGOING)
    particle_set%n_vir = &
         particle_set%n_tot - particle_set%n_in - particle_set%n_out
  end subroutine hepmc_event_to_particle_set

\eatline
\nwendcode{}\nwbegindocs{1035}Fill a WHIZARD event from a HepMC event record.  Recall that
HepMC events are always unweighted, so the weight is unity.
\nwenddocs{}\nwbegincode{1036}\moddef{HEP events: public}\plusendmoddef
  public :: hepmc_to_event
\nwendcode{}\nwbegincode{1037}\moddef{HEP events: procedures}\plusendmoddef
  subroutine hepmc_to_event &
       (event, hepmc_event, fallback_model, process_index, recover_beams, &
       use_alpha_s, use_scale)
    class(generic_event_t), intent(inout), target :: event
    type(hepmc_event_t), intent(inout) :: hepmc_event
    class(model_data_t), intent(in), target :: fallback_model
    integer, intent(out), optional :: process_index
    logical, intent(in), optional :: recover_beams
    logical, intent(in), optional :: use_alpha_s
    logical, intent(in), optional :: use_scale
    class(model_data_t), pointer :: model
    real(default) :: weight, scale, alpha_qcd
    type(particle_set_t) :: particle_set
    model => event%get_model_ptr ()
    call hepmc_event_to_particle_set (particle_set, &
         hepmc_event, model, fallback_model, PRT_DEFINITE_HELICITY)
    call event%set_hard_particle_set (particle_set)
    call particle_set%final ()
    call event%set_weight_ref (1._default)
    alpha_qcd = hepmc_event_get_alpha_qcd (hepmc_event)
    scale = hepmc_event_get_scale (hepmc_event)
    if (present (use_alpha_s)) then
       if (use_alpha_s .and. alpha_qcd > 0) &
            call event%set_alpha_qcd_forced (alpha_qcd)
    end if
    if (present (use_scale)) then
       if (use_scale .and. scale > 0) &
            call event%set_scale_forced (scale)
    end if
  end subroutine hepmc_to_event
  
\eatline
\nwendcode{}\nwbegindocs{1038}\nwdocspar
\subsubsection{LCIO event format}
The master output function fills a LCIO event object that is
already initialized, but has no particles in it.

In contrast to HepMC in LCIO there are no vertices (except for tracker 
and other detector specifications). So we assign first all incoming 
particles and then all outgoing particles to LCIO particle types.  
Particles which have neither parent nor children entries (this 
should not happen) are dropped. Finally, we insert the beam particles.  
If there are none, use the incoming particles instead.

Transform a particle into a {\tt{}lcio{\char95}particle} object, including
color and polarization.  The LCIO status is equivalent to the HepMC
status, in particular: 0 = null entry, 1 = physical particle, 2 =
decayed/fragmented SM hadron, tau or muon, 3 = other unphysical particle 
entry, 4 = incoming particles, 11 = intermediate resonance such as squarks.
The use of 11 for intermediate resonances is as done by HERWIG, see 
http://herwig.hepforge.org/trac/wiki/FaQs.

A beam-remnant particle (e.g., ISR photon) that has no children is
tagged as outgoing, otherwise unphysical.
\nwenddocs{}\nwbegincode{1039}\moddef{HEP events: public}\plusendmoddef
  public :: particle_to_lcio
\nwendcode{}\nwbegincode{1040}\moddef{HEP events: procedures}\plusendmoddef
  subroutine particle_to_lcio (prt, lprt)
    type(particle_t), intent(in) :: prt
    type(lcio_particle_t), intent(out) :: lprt
    integer :: lcio_status
    select case (prt%get_status ())
    case (PRT_UNDEFINED)
       lcio_status = 0
    case (PRT_OUTGOING)
       lcio_status = 1
    case (PRT_BEAM_REMNANT)
       if (prt%get_n_children () == 0) then
          lcio_status = 1
       else
          lcio_status = 3
       end if
    case (PRT_BEAM)
       lcio_status = 4
    case (PRT_RESONANT)
       if (abs (prt%get_pdg ()) == 13 .or. &
            abs (prt%get_pdg ()) == 15) then
          lcio_status = 2
       else
          lcio_status = 11
       end if
    case default
       lcio_status = 3
    end select
    call lcio_particle_init (lprt, &
         prt%get_momentum (), &
         prt%get_pdg (), &
         lcio_status)
    call lcio_particle_set_color (lprt, prt%get_color ())
    select case (prt%get_polarization_status ())
    case (PRT_DEFINITE_HELICITY)
       call lcio_polarization_init (lprt, prt%get_helicity ())
    case (PRT_GENERIC_POLARIZATION)
       call lcio_polarization_init (lprt, prt%get_polarization ())
    end select
  end subroutine particle_to_lcio

\eatline
\nwendcode{}\nwbegindocs{1041}\nwdocspar
\nwenddocs{}\nwbegindocs{1042}Initialize a particle from a LCIO particle object.  The model is
necessary for making a fully qualified flavor component.
\nwenddocs{}\nwbegincode{1043}\moddef{HEP events: public}\plusendmoddef
  public :: particle_from_lcio_particle
\nwendcode{}\nwbegincode{1044}\moddef{HEP events: procedures}\plusendmoddef
  subroutine particle_from_lcio_particle &
     (prt, lprt, model, daughters, parents, polarization)
    type(particle_t), intent(out) :: prt
    type(lcio_particle_t), intent(in) :: lprt
    type(model_data_t), intent(in), target :: model
    integer, dimension(:), intent(in) :: daughters, parents    
    type(flavor_t) :: flv
    type(color_t) :: col
    type(helicity_t) :: hel
    type(polarization_t) :: pol
    integer, intent(in) :: polarization
    integer :: i
    select case (lcio_particle_get_status (lprt))
    case (1);  call prt%set_status (PRT_OUTGOING)
    case (2);  call prt%set_status (PRT_RESONANT)
    case (3);  call prt%set_status (PRT_VIRTUAL)
    end select
    call flv%init (lcio_particle_get_pdg (lprt), model)
    call col%init (lcio_particle_get_flow (lprt))    
    if (flv%is_beam_remnant ())  call prt%set_status (PRT_BEAM_REMNANT)
    call prt%set_flavor (flv)
    call prt%set_color (col)
    call prt%set_polarization (polarization)
    select case (polarization)
    case (PRT_DEFINITE_HELICITY)
       call lcio_particle_to_hel (lprt, prt%get_flv (), hel)
       call prt%set_helicity (hel)
    case (PRT_GENERIC_POLARIZATION)
       call lcio_particle_to_pol (lprt, prt%get_flv (), pol)
       call prt%set_pol (pol)
    end select
    call prt%set_momentum (lcio_particle_get_momentum (lprt), &
         lcio_particle_get_mass_squared (lprt))
    call prt%set_parents (parents)
    call prt%set_children (daughters)
    if (prt%get_status () == PRT_VIRTUAL .and. size(parents) == 0) &
         call prt%set_status (PRT_INCOMING)
  end subroutine particle_from_lcio_particle

\eatline
\nwendcode{}\nwbegindocs{1045}\nwdocspar
\nwenddocs{}\nwbegincode{1046}\moddef{HEP events: public}\plusendmoddef
  public :: lcio_event_from_particle_set 
\nwendcode{}\nwbegincode{1047}\moddef{HEP events: procedures}\plusendmoddef
  subroutine lcio_event_from_particle_set (evt, particle_set)
    type(lcio_event_t), intent(inout) :: evt
    type(particle_set_t), intent(in) :: particle_set
    type(lcio_particle_t), dimension(:), allocatable :: lprt
    type(lcio_particle_t), dimension(2) :: lbeam
    integer, dimension(:), allocatable :: parent, child    
    logical, dimension(:), allocatable :: is_beam
    integer :: n_tot, i, j, n_parents
    n_tot = particle_set%n_tot
    allocate (lprt (n_tot))
    do i = 1, n_tot
       call particle_to_lcio (particle_set%prt(i), lprt(i))
       n_parents = particle_set%prt(i)%get_n_parents ()
       if (n_parents /= 0) then
          allocate (parent (n_parents))
          parent = particle_set%prt(i)%get_parents ()
          do j = 1, n_parents
             call lcio_particle_set_parent (lprt(i), lprt(parent(j)))
          end do
          deallocate (parent)
       end if
       call lcio_particle_add_to_evt_coll (lprt(i), evt)
    end do
    call lcio_event_add_coll (evt)
  end subroutine lcio_event_from_particle_set
  
\eatline
\nwendcode{}\nwbegindocs{1048}If a particle set is initialized from a LCIO event record, we have
to specify the treatment of polarization (unpolarized or density
matrix) which is common to all particles.  Correlated polarization
information is not available.
\nwenddocs{}\nwbegincode{1049}\moddef{HEP events: public}\plusendmoddef
  public :: lcio_event_to_particle_set
\nwendcode{}\nwbegincode{1050}\moddef{HEP events: procedures}\plusendmoddef
  subroutine lcio_event_to_particle_set &
       (particle_set, evt, model, fallback_model, polarization)
    type(particle_set_t), intent(inout), target :: particle_set
    type(lcio_event_t), intent(in) :: evt
    class(model_data_t), intent(in), target :: model, fallback_model
    integer, intent(in) :: polarization
    type(lcio_particle_t) :: prt
    integer, dimension(:), allocatable :: parents, daughters
    integer :: n_tot, i, j, n_parents, n_children
    n_tot = lcio_event_get_n_tot (evt)
    allocate (particle_set%prt (n_tot))
    do i = 1, n_tot
       prt = lcio_event_get_particle (evt, i-1)
       n_parents = lcio_particle_get_n_parents (prt)
       n_children = lcio_particle_get_n_children (prt) 
       allocate (daughters (n_children))
       allocate (parents (n_parents))
       if (n_children > 0) then
          do j = 1, n_children
             daughters(j) = lcio_get_n_children (evt,i,j)          
          end do
       end if
       if (n_parents > 0) then
          do j = 1, n_parents
             parents(j) = lcio_get_n_parents (evt,i,j)
          end do
       end if
       call particle_from_lcio_particle (particle_set%prt(i), prt, model, &
            daughters, parents, polarization)
       deallocate (daughters, parents)
    end do
    do i = 1, n_tot
       if (particle_set%prt(i)%get_status () == PRT_VIRTUAL) then
          CHECK_BEAM: do j = 1, particle_set%prt(i)%get_n_parents ()
             if (particle_set%prt(j)%get_status () == PRT_BEAM) &
                  call particle_set%prt(i)%set_status (PRT_INCOMING)
             exit CHECK_BEAM
          end do CHECK_BEAM
       end if
    end do
    particle_set%n_tot = n_tot
    particle_set%n_beam  = &
         count (particle_set%prt%get_status () == PRT_BEAM)
    particle_set%n_in  = &
         count (particle_set%prt%get_status () == PRT_INCOMING)
    particle_set%n_out = &
         count (particle_set%prt%get_status () == PRT_OUTGOING)
    particle_set%n_vir = &
         particle_set%n_tot - particle_set%n_in - particle_set%n_out
  end subroutine lcio_event_to_particle_set

\eatline
\nwendcode{}\nwbegindocs{1051}\nwdocspar
\nwenddocs{}\nwbegincode{1052}\moddef{HEP events: public}\plusendmoddef
  public :: lcio_to_event
\nwendcode{}\nwbegincode{1053}\moddef{HEP events: procedures}\plusendmoddef
  subroutine lcio_to_event &
       (event, lcio_event, fallback_model, process_index, recover_beams, &
       use_alpha_s, use_scale)
    class(generic_event_t), intent(inout), target :: event
    type(lcio_event_t), intent(inout) :: lcio_event
    class(model_data_t), intent(in), target :: fallback_model
    integer, intent(out), optional :: process_index
    logical, intent(in), optional :: recover_beams
    logical, intent(in), optional :: use_alpha_s
    logical, intent(in), optional :: use_scale
    class(model_data_t), pointer :: model
    real(default) :: weight, scale, alpha_qcd
    type(particle_set_t) :: particle_set
    model => event%get_model_ptr ()
    call lcio_event_to_particle_set (particle_set, &
         lcio_event, model, fallback_model, PRT_DEFINITE_HELICITY)
    call event%set_hard_particle_set (particle_set)
    call particle_set%final ()
    alpha_qcd = lcio_event_get_alphas (lcio_event)
    scale = lcio_event_get_scaleval (lcio_event)
    if (present (use_alpha_s)) then
       if (use_alpha_s .and. alpha_qcd > 0) &
            call event%set_alpha_qcd_forced (alpha_qcd)
    end if
    if (present (use_scale)) then
       if (use_scale .and. scale > 0) &
            call event%set_scale_forced (scale)
    end if
  end subroutine lcio_to_event
  
\eatline
\nwendcode{}\nwbegindocs{1054}\nwdocspar
\subsection{Unit tests}
Set up a chain of production and decay and factorize the result into
particles.  The process is $d\bar d \to Z \to q\bar q$.
\nwenddocs{}\nwbegincode{1055}\moddef{HEP events: public}\plusendmoddef
  public :: hep_events_test
\nwendcode{}\nwbegincode{1056}\moddef{HEP events: procedures}\plusendmoddef
  subroutine hep_events_test (u, results)
    integer, intent(in) :: u
    type(test_results_t), intent(inout) :: results
  \LA{}HEP events: execute tests\RA{}  
  end subroutine hep_events_test

\nwendcode{}\nwbegindocs{1057} %def particles_test  
\nwenddocs{}\nwbegindocs{1058}If {\tt{}HepMC} is available, check the routines via {\tt{}HepMC}.
\nwenddocs{}\nwbegincode{1059}\moddef{HEP events: execute tests}\endmoddef
  if (hepmc_is_available ()) then
     call test (hep_events_1, "hep_events_1", &
          "check HepMC event routines", &
          u, results)
  end if
\nwendcode{}\nwbegincode{1060}\moddef{HEP events: tests}\endmoddef
  subroutine hep_events_1 (u)
    use os_interface
    integer, intent(in) :: u
    type(model_data_t), target :: model
    type(flavor_t), dimension(3) :: flv
    type(color_t), dimension(3) :: col
    type(helicity_t), dimension(3) :: hel
    type(quantum_numbers_t), dimension(3) :: qn
    type(vector4_t), dimension(3) :: p
    type(interaction_t), target :: int1, int2
    type(quantum_numbers_mask_t) :: qn_mask_conn, qn_rest
    type(evaluator_t), target :: eval
    type(interaction_t), pointer :: int
    type(particle_set_t) :: particle_set1, particle_set2
    type(particle_set_t) :: particle_set3, particle_set4
    type(subevt_t) :: subevt
    logical :: ok
    integer :: unit, iostat      
    type(hepmc_event_t) :: hepmc_event
    type(hepmc_iostream_t) :: iostream    

    write (u, "(A)")  "* Test output: HEP events"
    write (u, "(A)")  "*   Purpose: test HepMC event routines"
    write (u, "(A)")      

    write (u, "(A)")  "* Reading model file"

    call model%init_sm_test ()

    write (u, "(A)")
    write (u, "(A)")  "* Initializing production process"

    call interaction_init (int1, 2, 0, 1, set_relations=.true.)
    call flv%init ([1, -1, 23], model)
    call col%init_col_acl ([0, 0, 0], [0, 0, 0])
    call hel(3)%init ( 1, 1)
    call qn%init (flv, col, hel)
    call interaction_add_state (int1, qn, value=(0.25_default, 0._default))
    call hel(3)%init ( 1,-1)
    call qn%init (flv, col, hel)
    call interaction_add_state (int1, qn, value=(0._default, 0.25_default))
    call hel(3)%init (-1, 1)
    call qn%init (flv, col, hel)
    call interaction_add_state (int1, qn, value=(0._default,-0.25_default))
    call hel(3)%init (-1,-1)
    call qn%init (flv, col, hel)
    call interaction_add_state (int1, qn, value=(0.25_default, 0._default))
    call hel(3)%init ( 0, 0)
    call qn%init (flv, col, hel)
    call interaction_add_state (int1, qn, value=(0.5_default, 0._default))
    call interaction_freeze (int1)
    p(1) = vector4_moving (45._default, 45._default, 3)
    p(2) = vector4_moving (45._default,-45._default, 3)
    p(3) = p(1) + p(2)
    call interaction_set_momenta (int1, p)

    write (u, "(A)")
    write (u, "(A)")  "* Setup decay process"

    call interaction_init (int2, 1, 0, 2, set_relations=.true.)
    call flv%init ([23, 1, -1], model)
    call col%init_col_acl ([0, 501, 0], [0, 0, 501])
    call hel%init ([1, 1, 1], [1, 1, 1])
    call qn%init (flv, col, hel)
    call interaction_add_state (int2, qn, value=(1._default, 0._default))
    call hel%init ([1, 1, 1], [-1,-1,-1])
    call qn%init (flv, col, hel)
    call interaction_add_state (int2, qn, value=(0._default, 0.1_default))
    call hel%init ([-1,-1,-1], [1, 1, 1])
    call qn%init (flv, col, hel)
    call interaction_add_state (int2, qn, value=(0._default,-0.1_default))
    call hel%init ([-1,-1,-1], [-1,-1,-1])
    call qn%init (flv, col, hel)
    call interaction_add_state (int2, qn, value=(1._default, 0._default))
    call hel%init ([0, 1,-1], [0, 1,-1])
    call qn%init (flv, col, hel)
    call interaction_add_state (int2, qn, value=(4._default, 0._default))
    call hel%init ([0,-1, 1], [0, 1,-1])
    call qn%init (flv, col, hel)
    call interaction_add_state (int2, qn, value=(2._default, 0._default))
    call hel%init ([0, 1,-1], [0,-1, 1])
    call qn%init (flv, col, hel)
    call interaction_add_state (int2, qn, value=(2._default, 0._default))
    call hel%init ([0,-1, 1], [0,-1, 1])
    call qn%init (flv, col, hel)
    call interaction_add_state (int2, qn, value=(4._default, 0._default))
    call flv%init ([23, 2, -2], model)
    call hel%init ([0, 1,-1], [0, 1,-1])
    call qn%init (flv, col, hel)
    call interaction_add_state (int2, qn, value=(0.5_default, 0._default))
    call hel%init ([0,-1, 1], [0,-1, 1])
    call qn%init (flv, col, hel)
    call interaction_add_state (int2, qn, value=(0.5_default, 0._default))
    call interaction_freeze (int2)
    p(2) = vector4_moving (45._default, 45._default, 2)
    p(3) = vector4_moving (45._default,-45._default, 2)
    call interaction_set_momenta (int2, p)
    call interaction_set_source_link (int2, 1, int1, 3)
    call interaction_write (int1, u)
    call interaction_write (int2, u)

    write (u, "(A)")
    write (u, "(A)")  "* Concatenate production and decay"

    call evaluator_init_product (eval, int1, int2, qn_mask_conn, &
         connections_are_resonant=.true.)
    call evaluator_receive_momenta (eval)
    call eval%evaluate ()
    call eval%write (u)
    
    write (u, "(A)")
    write (u, "(A)")  "* Factorize as subevent (complete, polarized)"
    write (u, "(A)")
    
    int => evaluator_get_int_ptr (eval)
    call particle_set1%init &
         (ok, int, int, FM_FACTOR_HELICITY, &
          [0.2_default, 0.2_default], .false., .true.)
    call particle_set1%write (u)

    write (u, "(A)")
    write (u, "(A)")  "* Factorize as subevent (in/out only, selected helicity)"
    write (u, "(A)")
    
    int => evaluator_get_int_ptr (eval)
    call particle_set2%init &
         (ok, int, int, FM_SELECT_HELICITY, &
          [0.9_default, 0.9_default], .false., .false.)
    call particle_set2%write (u)
    call particle_set2%final ()

    write (u, "(A)")
    write (u, "(A)")  "* Factorize as subevent (complete, selected helicity)"
    write (u, "(A)") 
    
    int => evaluator_get_int_ptr (eval)
    call particle_set2%init &
         (ok, int, int, FM_SELECT_HELICITY, &
          [0.7_default, 0.7_default], .false., .true.)
    call particle_set2%write (u)      

    write (u, "(A)")
    write (u, "(A)")  "* Transfer particle_set to HepMC, print, and output to"
    write (u, "(A)")  "        hep_events.hepmc.dat"
    write (u, "(A)")
    
    call hepmc_event_init (hepmc_event, 11, 127)
    call hepmc_event_from_particle_set (hepmc_event, particle_set2)
    call hepmc_event_print (hepmc_event)
    call hepmc_iostream_open_out &
         (iostream , var_str ("hep_events.hepmc.dat"))
    call hepmc_iostream_write_event (iostream, hepmc_event)
    call hepmc_iostream_close (iostream)

    write (u, "(A)")
    write (u, "(A)")  "* Recover from HepMC file"
    write (u, "(A)")
    
    call particle_set2%final ()
    call hepmc_event_final (hepmc_event)
    call hepmc_event_init (hepmc_event)
    call hepmc_iostream_open_in &
         (iostream , var_str ("hep_events.hepmc.dat"))
    call hepmc_iostream_read_event (iostream, hepmc_event, ok)
    call hepmc_iostream_close (iostream)
    call hepmc_event_to_particle_set (particle_set2, &
         hepmc_event, model, model, PRT_DEFINITE_HELICITY)
    call particle_set2%write (u)   

    write (u, "(A)")
    write (u, "(A)")  "* Cleanup"
    
    call particle_set1%final ()
    call particle_set2%final ()
    call evaluator_final (eval)
    call interaction_final (int1)
    call interaction_final (int2)
    call hepmc_event_final (hepmc_event)            
    call model%final ()
       
    write (u, "(A)")
    write (u, "(A)")  "* Test output end: particles_2"

  end subroutine hep_events_1
  
\nwendcode{}\nwbegindocs{1061}\nwdocspar
%
%%%%%%%%%%%%%%%%%%%%%%%%%%%%%%%%%%%%%%%%%%%%%%%%%%%%%%%%%%%%%%%%%%%%%%%% 
\section{LHEF Input/Output}
The LHEF event record is standardized.  It is an ASCII format.  We try
our best at using it for both input and output.
\nwenddocs{}\nwbegincode{1062}\moddef{\code{}eio{\char95}lhef.f90\edoc{}}\endmoddef
\LA{}File header\RA{}

module eio_lhef
  
  use kinds
  use io_units
\LA{}Use strings\RA{}
  use string_utils
  use unit_tests
  use diagnostics
  use os_interface

  use xml
  use lorentz
  use model_data
  use particles
  use event_base
  use eio_data
  use eio_base
  use hep_common
  use hep_events

\LA{}Standard module head\RA{}

\LA{}EIO LHEF: public\RA{}

\LA{}EIO LHEF: types\RA{}

\LA{}EIO LHEF: interfaces\RA{}

contains
  
\LA{}EIO LHEF: procedures\RA{}

\LA{}EIO LHEF: tests\RA{}

end module eio_lhef
\eatline
\nwendcode{}\nwbegindocs{1063}\nwdocspar
\subsection{Type}
With sufficient confidence that it will always be three characters, we
can store the version string with a default value.
\nwenddocs{}\nwbegincode{1064}\moddef{EIO LHEF: public}\endmoddef
  public :: eio_lhef_t
\nwendcode{}\nwbegincode{1065}\moddef{EIO LHEF: types}\endmoddef
  type, extends (eio_t) :: eio_lhef_t
     logical :: writing = .false.
     logical :: reading = .false.
     integer :: unit = 0
     type(event_sample_data_t) :: data
     type(cstream_t) :: cstream
     character(3) :: version = "1.0"
     logical :: keep_beams = .false.
     logical :: keep_remnants = .true.
     logical :: recover_beams = .true.
     logical :: unweighted = .true.
     logical :: write_sqme_ref = .false.
     logical :: write_sqme_prc = .false.
     logical :: write_sqme_alt = .false.
     logical :: use_alpha_s_from_file = .false.
     logical :: use_scale_from_file = .false.
     integer :: n_alt = 0
     integer, dimension(:), allocatable :: proc_num_id
     integer :: i_weight_sqme = 0
     type(xml_tag_t) :: tag_lhef, tag_head, tag_init, tag_event
     type(xml_tag_t), allocatable :: tag_gen_n, tag_gen_v
     type(xml_tag_t), allocatable :: tag_generator, tag_xsecinfo
     type(xml_tag_t), allocatable :: tag_sqme_ref, tag_sqme_prc
     type(xml_tag_t), dimension(:), allocatable :: tag_sqme_alt, tag_wgts_alt
     type(xml_tag_t), allocatable :: tag_weight, tag_weightinfo, tag_weights
   contains
   \LA{}EIO LHEF: eio lhef: TBP\RA{}
  end type eio_lhef_t
  
\eatline
\nwendcode{}\nwbegindocs{1066}\nwdocspar
\subsection{Specific Methods}
Set parameters that are specifically used with LHEF.
\nwenddocs{}\nwbegincode{1067}\moddef{EIO LHEF: eio lhef: TBP}\endmoddef
  procedure :: set_parameters => eio_lhef_set_parameters
\nwendcode{}\nwbegincode{1068}\moddef{EIO LHEF: procedures}\endmoddef
  subroutine eio_lhef_set_parameters (eio, &
       keep_beams, keep_remnants, recover_beams, &
       use_alpha_s_from_file, use_scale_from_file, &
       version, extension, write_sqme_ref, write_sqme_prc, write_sqme_alt)
    class(eio_lhef_t), intent(inout) :: eio
    logical, intent(in), optional :: keep_beams
    logical, intent(in), optional :: keep_remnants
    logical, intent(in), optional :: recover_beams
    logical, intent(in), optional :: use_alpha_s_from_file
    logical, intent(in), optional :: use_scale_from_file
    character(*), intent(in), optional :: version
    type(string_t), intent(in), optional :: extension
    logical, intent(in), optional :: write_sqme_ref
    logical, intent(in), optional :: write_sqme_prc
    logical, intent(in), optional :: write_sqme_alt
    if (present (keep_beams))  eio%keep_beams = keep_beams
    if (present (keep_remnants))  eio%keep_remnants = keep_remnants
    if (present (recover_beams))  eio%recover_beams = recover_beams
    if (present (use_alpha_s_from_file)) &
         eio%use_alpha_s_from_file = use_alpha_s_from_file
    if (present (use_scale_from_file)) &
         eio%use_scale_from_file = use_scale_from_file
    if (present (version)) then
       select case (version)
       case ("1.0", "2.0", "3.0")
          eio%version = version
       case default
          call msg_error ("LHEF version " // version &
               // " is not supported.  Inserting 2.0")
          eio%version = "2.0"
       end select
    end if
    if (present (extension)) then
       eio%extension = extension
    else
       eio%extension = "lhe"
    end if
    if (present (write_sqme_ref))  eio%write_sqme_ref = write_sqme_ref
    if (present (write_sqme_prc))  eio%write_sqme_prc = write_sqme_prc
    if (present (write_sqme_alt))  eio%write_sqme_alt = write_sqme_alt
  end subroutine eio_lhef_set_parameters
  
\eatline
\nwendcode{}\nwbegindocs{1069}\nwdocspar
\subsection{Common Methods}
Output.  This is not the actual event format, but a readable account
of the current object status.
\nwenddocs{}\nwbegincode{1070}\moddef{EIO LHEF: eio lhef: TBP}\plusendmoddef
  procedure :: write => eio_lhef_write
\nwendcode{}\nwbegincode{1071}\moddef{EIO LHEF: procedures}\plusendmoddef
  subroutine eio_lhef_write (object, unit)
    class(eio_lhef_t), intent(in) :: object
    integer, intent(in), optional :: unit
    integer :: u, i
    u = given_output_unit (unit)
    write (u, "(1x,A)")  "LHEF event stream:"
    if (object%writing) then
       write (u, "(3x,A,A)")  "Writing to file   = ", char (object%filename)
    else if (object%reading) then
       write (u, "(3x,A,A)")  "Reading from file = ", char (object%filename)
    else
       write (u, "(3x,A)")  "[closed]"
    end if
    write (u, "(3x,A,L1)")    "Keep beams        = ", object%keep_beams
    write (u, "(3x,A,L1)")    "Keep remnants     = ", object%keep_remnants
    write (u, "(3x,A,L1)")    "Recover beams     = ", object%recover_beams
    write (u, "(3x,A,L1)")    "Alpha_s from file = ", &
         object%use_alpha_s_from_file
    write (u, "(3x,A,L1)")    "Scale from file   = ", &
         object%use_scale_from_file
    write (u, "(3x,A,A)")     "Version           = ", object%version
    write (u, "(3x,A,A,A)")     "File extension    = '", &
         char (object%extension), "'"
    if (allocated (object%proc_num_id)) then
       write (u, "(3x,A)")  "Numerical process IDs:"
       do i = 1, size (object%proc_num_id)
          write (u, "(5x,I0,': ',I0)")  i, object%proc_num_id(i)
       end do
    end if
  end subroutine eio_lhef_write
  
\eatline
\nwendcode{}\nwbegindocs{1072}Finalizer: close any open file.
\nwenddocs{}\nwbegincode{1073}\moddef{EIO LHEF: eio lhef: TBP}\plusendmoddef
  procedure :: final => eio_lhef_final
\nwendcode{}\nwbegincode{1074}\moddef{EIO LHEF: procedures}\plusendmoddef
  subroutine eio_lhef_final (object)
    class(eio_lhef_t), intent(inout) :: object
    if (allocated (object%proc_num_id))  deallocate (object%proc_num_id)
    if (object%writing) then
       write (msg_buffer, "(A,A,A)")  "Events: closing LHEF file '", &
            char (object%filename), "'"
       call msg_message ()
       call object%write_footer ()
       close (object%unit)
       object%writing = .false.
    else if (object%reading) then
       write (msg_buffer, "(A,A,A)")  "Events: closing LHEF file '", &
            char (object%filename), "'"
       call msg_message ()
       call object%cstream%final ()
       close (object%unit)
       object%reading = .false.
    end if
  end subroutine eio_lhef_final
  
\eatline
\nwendcode{}\nwbegindocs{1075}Common initialization for input and output.
\nwenddocs{}\nwbegincode{1076}\moddef{EIO LHEF: eio lhef: TBP}\plusendmoddef
  procedure :: common_init => eio_lhef_common_init
\nwendcode{}\nwbegincode{1077}\moddef{EIO LHEF: procedures}\plusendmoddef
  subroutine eio_lhef_common_init (eio, sample, data, extension)
    class(eio_lhef_t), intent(inout) :: eio
    type(string_t), intent(in) :: sample
    type(string_t), intent(in), optional :: extension
    type(event_sample_data_t), intent(in), optional :: data
    if (.not. present (data)) &
         call msg_bug ("LHEF initialization: missing data")
    eio%data = data
    if (data%n_beam /= 2) &
         call msg_fatal ("LHEF: defined for scattering processes only")
    eio%unweighted = data%unweighted
    if (eio%unweighted) then
       select case (data%norm_mode)
       case (NORM_UNIT)
       case default;  call msg_fatal &
            ("LHEF: normalization for unweighted events must be '1'")
       end select
    else
       select case (data%norm_mode)
       case (NORM_SIGMA)
       case default;  call msg_fatal &
            ("LHEF: normalization for weighted events must be 'sigma'")
       end select
    end if
    eio%n_alt = data%n_alt
    eio%sample = sample
    if (present (extension)) then
       eio%extension = extension
    end if
    call eio%set_filename ()
    eio%unit = free_unit ()
    call eio%init_tags (data)
    allocate (eio%proc_num_id (data%n_proc), source = data%proc_num_id)
  end subroutine eio_lhef_common_init
  
\eatline
\nwendcode{}\nwbegindocs{1078}Initialize the tag objects.  Some tags depend on the LHEF
version.  In particular, the tags that in LHEF 2.0 identify
individual weights by name in each event block, in LHEF 3.0 are
replaced by info tags in the init block and a single \texttt{weights}
tag in the event block.  The name attributes of those tags
are specific for \whizard.
\nwenddocs{}\nwbegincode{1079}\moddef{EIO LHEF: eio lhef: TBP}\plusendmoddef
  procedure :: init_tags => eio_lhef_init_tags
\nwendcode{}\nwbegincode{1080}\moddef{EIO LHEF: procedures}\plusendmoddef
  subroutine eio_lhef_init_tags (eio, data)
    class(eio_lhef_t), intent(inout) :: eio
    type(event_sample_data_t), intent(in) :: data
    real(default), parameter :: pb_per_fb = 1.e-3_default
    integer :: i
    call eio%tag_lhef%init ( &
         var_str ("LesHouchesEvents"), &
         [xml_attribute (var_str ("version"), var_str (eio%version))], &
         .true.)
    call eio%tag_head%init ( &
         var_str ("header"), &
         .true.)
    call eio%tag_init%init ( &
         var_str ("init"), &
         .true.)
    call eio%tag_event%init (var_str ("event"), &
         .true.)
    select case (eio%version)
    case ("1.0")
       allocate (eio%tag_gen_n)
       call eio%tag_gen_n%init ( &
            var_str ("generator_name"), &
            .true.)
       allocate (eio%tag_gen_v)
       call eio%tag_gen_v%init ( &
            var_str ("generator_version"), &
            .true.)
    end select
    select case (eio%version)
    case ("2.0", "3.0")
       allocate (eio%tag_generator)
       call eio%tag_generator%init ( &
            var_str ("generator"), &
            [xml_attribute (var_str ("version"), var_str ("\LA{}Version\RA{}"))], &
            .true.)
       allocate (eio%tag_xsecinfo)
       call eio%tag_xsecinfo%init ( &
            var_str ("xsecinfo"), &
            [xml_attribute (var_str ("neve"), str (data%n_evt)), &
             xml_attribute (var_str ("totxsec"), &
                            str (data%total_cross_section * pb_per_fb))])
    end select
    select case (eio%version)
    case ("2.0")
       allocate (eio%tag_weight)
       call eio%tag_weight%init (var_str ("weight"), &
            [xml_attribute (var_str ("name"))])
       if (eio%write_sqme_ref) then
          allocate (eio%tag_sqme_ref)
          call eio%tag_sqme_ref%init (var_str ("weight"), &
               [xml_attribute (var_str ("name"), var_str ("sqme_ref"))], &
               .true.)
       end if
       if (eio%write_sqme_prc) then
          allocate (eio%tag_sqme_prc)
          call eio%tag_sqme_prc%init (var_str ("weight"), &
               [xml_attribute (var_str ("name"), var_str ("sqme_prc"))], &
               .true.)
       end if
       if (eio%n_alt > 0) then
          if (eio%write_sqme_alt) then
             allocate (eio%tag_sqme_alt (1))
             call eio%tag_sqme_alt(1)%init (var_str ("weight"), &
                  [xml_attribute (var_str ("name"), var_str ("sqme_alt"))], &
                  .true.)
          end if
          allocate (eio%tag_wgts_alt (1))
          call eio%tag_wgts_alt(1)%init (var_str ("weight"), &
               [xml_attribute (var_str ("name"), var_str ("wgts_alt"))], &
               .true.)
       end if
    case ("3.0")
       if (eio%write_sqme_ref) then
          allocate (eio%tag_sqme_ref)
          call eio%tag_sqme_ref%init (var_str ("weightinfo"), &
               [xml_attribute (var_str ("name"), var_str ("sqme_ref"))])
       end if
       if (eio%write_sqme_prc) then
          allocate (eio%tag_sqme_prc)
          call eio%tag_sqme_prc%init (var_str ("weightinfo"), &
               [xml_attribute (var_str ("name"), var_str ("sqme_prc"))])
       end if
       if (eio%n_alt > 0) then
          if (eio%write_sqme_alt) then
             allocate (eio%tag_sqme_alt (eio%n_alt))
             do i = 1, eio%n_alt
                call eio%tag_sqme_alt(i)%init (var_str ("weightinfo"), &
                     [xml_attribute (var_str ("name"), &
                                     var_str ("sqme_alt") // str (i))])
             end do
          end if
          allocate (eio%tag_wgts_alt (eio%n_alt))
          do i = 1, eio%n_alt
             call eio%tag_wgts_alt(i)%init (var_str ("weightinfo"), &
                  [xml_attribute (var_str ("name"), &
                                  var_str ("wgts_alt") // str (i))])
          end do
       end if
       allocate (eio%tag_weightinfo)
       call eio%tag_weightinfo%init (var_str ("weightinfo"), &
            [xml_attribute (var_str ("name"))])
       allocate (eio%tag_weights)
       call eio%tag_weights%init (var_str ("weights"), .true.)
    end select
  end subroutine eio_lhef_init_tags
  
\eatline
\nwendcode{}\nwbegindocs{1081}Initialize event writing.
\nwenddocs{}\nwbegincode{1082}\moddef{EIO LHEF: eio lhef: TBP}\plusendmoddef
  procedure :: init_out => eio_lhef_init_out
\nwendcode{}\nwbegincode{1083}\moddef{EIO LHEF: procedures}\plusendmoddef
  subroutine eio_lhef_init_out (eio, sample, data, success, extension)
    class(eio_lhef_t), intent(inout) :: eio
    type(string_t), intent(in) :: sample
    type(string_t), intent(in), optional :: extension
    type(event_sample_data_t), intent(in), optional :: data
    logical, intent(out), optional :: success
    integer :: u, i
    call eio%set_splitting (data)
    call eio%common_init (sample, data, extension)
    write (msg_buffer, "(A,A,A)")  "Events: writing to LHEF file '", &
         char (eio%filename), "'"
    call msg_message ()
    eio%writing = .true.
    u = eio%unit
    open (u, file = char (eio%filename), &
         action = "write", status = "replace")
    call eio%write_header ()
    call heprup_init &
         (data%pdg_beam, &
          data%energy_beam, &
          n_processes = data%n_proc, &
          unweighted = data%unweighted, &
          negative_weights = data%negative_weights)             
    do i = 1, data%n_proc
       call heprup_set_process_parameters (i = i, &
            process_id = data%proc_num_id(i), &
            cross_section = data%cross_section(i), &
            error = data%error(i))
    end do
    call eio%tag_init%write (u);  write (u, *)
    call heprup_write_lhef (u)
    select case (eio%version)
    case ("2.0");  call eio%write_init_20 (data)
    case ("3.0");  call eio%write_init_30 (data)
    end select
    call eio%tag_init%close (u);  write (u, *)
    if (present (success))  success = .true.
  end subroutine eio_lhef_init_out
    
\eatline
\nwendcode{}\nwbegindocs{1084}Initialize event reading.  First read the LHEF tag and version, then
read the header and skip over its contents, then read the init block.
(We require the opening and closing tags of the init block to be placed
on separate lines without extra stuff.)

For input, we do not (yet?) support split event files.
\nwenddocs{}\nwbegincode{1085}\moddef{EIO LHEF: eio lhef: TBP}\plusendmoddef
  procedure :: init_in => eio_lhef_init_in
\nwendcode{}\nwbegincode{1086}\moddef{EIO LHEF: procedures}\plusendmoddef
  subroutine eio_lhef_init_in (eio, sample, data, success, extension)
    class(eio_lhef_t), intent(inout) :: eio
    type(string_t), intent(in) :: sample
    type(string_t), intent(in), optional :: extension
    type(event_sample_data_t), intent(inout), optional :: data
    logical, intent(out), optional :: success
    logical :: exist, ok, closing
    type(event_sample_data_t) :: data_file
    type(string_t) :: string
    integer :: u
    eio%split = .false.
    call eio%common_init (sample, data, extension)
    write (msg_buffer, "(A,A,A)")  "Events: reading from LHEF file '", &
         char (eio%filename), "'"
    call msg_message ()
    inquire (file = char (eio%filename), exist = exist)
    if (.not. exist)  call msg_fatal ("Events: LHEF file not found.")
    eio%reading = .true.
    u = eio%unit
    open (u, file = char (eio%filename), &
         action = "read", status = "old")
    call eio%cstream%init (u)
    call eio%read_header ()
    call eio%tag_init%read (eio%cstream, ok)
    if (.not. ok)  call err_init
    select case (eio%version)
    case ("1.0");  call eio%read_init_10 (data_file)
       call eio%tag_init%read_content (eio%cstream, string, closing)
       if (string /= "" .or. .not. closing)  call err_init
    case ("2.0");  call eio%read_init_20 (data_file)
    case ("3.0");  call eio%read_init_30 (data_file)
    end select
    call eio%merge_data (data, data_file)
    if (present (success))  success = .true.

  contains
    
    subroutine err_init
      call msg_fatal ("LHEF: syntax error in init tag")
    end subroutine err_init
      
  end subroutine eio_lhef_init_in
    
\eatline
\nwendcode{}\nwbegindocs{1087}Merge event sample data: we can check the data in the file against
our assumptions and set or reset parameters.
\nwenddocs{}\nwbegincode{1088}\moddef{EIO LHEF: eio lhef: TBP}\plusendmoddef
  procedure :: merge_data => eio_merge_data
\nwendcode{}\nwbegincode{1089}\moddef{EIO LHEF: procedures}\plusendmoddef
  subroutine eio_merge_data (eio, data, data_file)
    class(eio_lhef_t), intent(inout) :: eio
    type(event_sample_data_t), intent(inout) :: data
    type(event_sample_data_t), intent(in) :: data_file
    real, parameter :: tolerance = 1000 * epsilon (1._default)
    if (data%unweighted .neqv. data_file%unweighted)  call err_weights
    if (data%negative_weights .neqv. data_file%negative_weights) &
         call err_weights
    if (data%norm_mode /= data_file%norm_mode)  call err_norm
    if (data%n_beam /= data_file%n_beam)  call err_beams
    if (any (data%pdg_beam /= data_file%pdg_beam))  call err_beams
    if (any (abs ((data%energy_beam - data_file%energy_beam)) &
         > (data%energy_beam + data_file%energy_beam) * tolerance)) &
         call err_beams
    if (data%n_proc /= data_file%n_proc)  call err_proc
    if (any (data%proc_num_id /= data_file%proc_num_id))  call err_proc
    where (data%cross_section == 0)
       data%cross_section = data_file%cross_section
       data%error = data_file%error
    end where
    data%total_cross_section = sum (data%cross_section)
    if (data_file%n_evt > 0) then
       if (data%n_evt > 0 .and. data_file%n_evt /= data%n_evt)  call err_n_evt
       data%n_evt = data_file%n_evt
    end if
  contains
    subroutine err_weights
      call msg_fatal ("LHEF: mismatch in event weight properties")
    end subroutine err_weights
    subroutine err_norm
      call msg_fatal ("LHEF: mismatch in event normalization")
    end subroutine err_norm
    subroutine err_beams
      call msg_fatal ("LHEF: mismatch in beam properties")
    end subroutine err_beams
    subroutine err_proc
      call msg_fatal ("LHEF: mismatch in process definitions")
    end subroutine err_proc
    subroutine err_n_evt
      call msg_error ("LHEF: mismatch in specified number of events (ignored)")
    end subroutine err_n_evt
  end subroutine eio_merge_data
    
\eatline
\nwendcode{}\nwbegindocs{1090}Switch from input to output: reopen the file for reading.
\nwenddocs{}\nwbegincode{1091}\moddef{EIO LHEF: eio lhef: TBP}\plusendmoddef
  procedure :: switch_inout => eio_lhef_switch_inout
\nwendcode{}\nwbegincode{1092}\moddef{EIO LHEF: procedures}\plusendmoddef
  subroutine eio_lhef_switch_inout (eio, success)
    class(eio_lhef_t), intent(inout) :: eio
    logical, intent(out), optional :: success
    call msg_bug ("LHEF: in-out switch not supported")
    if (present (success))  success = .false.
  end subroutine eio_lhef_switch_inout
  
\eatline
\nwendcode{}\nwbegindocs{1093}Split event file: increment the counter, close the current file, open a new
one.  If the file needs a header, repeat it for the new file.  (We assume that
the common block contents are still intact.)
\nwenddocs{}\nwbegincode{1094}\moddef{EIO LHEF: eio lhef: TBP}\plusendmoddef
  procedure :: split_out => eio_lhef_split_out
\nwendcode{}\nwbegincode{1095}\moddef{EIO LHEF: procedures}\plusendmoddef
  subroutine eio_lhef_split_out (eio)
    class(eio_lhef_t), intent(inout) :: eio
    integer :: u
    if (eio%split) then
       eio%split_index = eio%split_index + 1
       call eio%set_filename ()
       write (msg_buffer, "(A,A,A)")  "Events: writing to LHEF file '", &
            char (eio%filename), "'"
       call msg_message ()
       call eio%write_footer ()
       u = eio%unit
       close (u)
       open (u, file = char (eio%filename), &
            action = "write", status = "replace")
       call eio%write_header ()
       call eio%tag_init%write (u);  write (u, *)
       call heprup_write_lhef (u)
       select case (eio%version)
       case ("2.0");  call eio%write_init_20 (eio%data)
       case ("3.0");  call eio%write_init_30 (eio%data)
       end select
       call eio%tag_init%close (u);  write (u, *)
    end if
  end subroutine eio_lhef_split_out
  
\eatline
\nwendcode{}\nwbegindocs{1096}Output an event.  Write first the event indices, then weight and
squared matrix element, then the particle set.
\nwenddocs{}\nwbegincode{1097}\moddef{EIO LHEF: eio lhef: TBP}\plusendmoddef
  procedure :: output => eio_lhef_output
\nwendcode{}\nwbegincode{1098}\moddef{EIO LHEF: procedures}\plusendmoddef
  subroutine eio_lhef_output (eio, event, i_prc, reading, passed, pacify)
    class(eio_lhef_t), intent(inout) :: eio
    class(generic_event_t), intent(in), target :: event
    integer, intent(in) :: i_prc
    logical, intent(in), optional :: reading, passed, pacify
    integer :: u
    u = given_output_unit (eio%unit);  if (u < 0)  return
    if (present (passed)) then
       if (.not. passed)  return
    end if
    if (eio%writing) then
       call hepeup_from_event (event, &
            process_index = eio%proc_num_id (i_prc), &
            keep_beams = eio%keep_beams, &
            keep_remnants = eio%keep_remnants)
       write (u, '(A)') "<event>"
       call hepeup_write_lhef (eio%unit)
       select case (eio%version)
       case ("2.0");  call eio%write_event_20 (event)
       case ("3.0");  call eio%write_event_30 (event)
       end select
       write (u, '(A)') "</event>"
    else
       call eio%write ()
       call msg_fatal ("LHEF file is not open for writing")
    end if
  end subroutine eio_lhef_output

\eatline
\nwendcode{}\nwbegindocs{1099}Input an event.  Upon input of {\tt{}i{\char95}prc}, we can just read in the
whole HEPEUP common block.  These data are known to come first.  The
{\tt{}i{\char95}prc} value can be deduced from the IDPRUP value by a table
lookup.

Reading the common block bypasses the {\tt{}cstream} which accesses the
input unit.  This is consistent with the LHEF specification.  After
the common-block data have been swallowed, we can resume reading from
stream.

We don't catch actual I/O errors.  However, we return a negative value in
{\tt{}iostat} if we reached the terminating {\tt{}</LesHouchesEvents>} tag.
\nwenddocs{}\nwbegincode{1100}\moddef{EIO LHEF: eio lhef: TBP}\plusendmoddef
  procedure :: input_i_prc => eio_lhef_input_i_prc
\nwendcode{}\nwbegincode{1101}\moddef{EIO LHEF: procedures}\plusendmoddef
  subroutine eio_lhef_input_i_prc (eio, i_prc, iostat)
    class(eio_lhef_t), intent(inout) :: eio
    integer, intent(out) :: i_prc
    integer, intent(out) :: iostat
    integer :: i, proc_num_id
    type(string_t) :: s
    logical :: ok
    iostat = 0
    call eio%tag_lhef%read_content (eio%cstream, s, ok)
    if (ok) then
       if (s == "") then
          iostat = -1
       else
          call err_close
       end if
       return
    else
       call eio%cstream%revert_record (s)
    end if
    call eio%tag_event%read (eio%cstream, ok)
    if (.not. ok) then
       call err_evt1
       return
    end if
    call hepeup_read_lhef (eio%unit)
    call hepeup_get_event_parameters (proc_id = proc_num_id)
    i_prc = 0
    FIND_I_PRC: do i = 1, size (eio%proc_num_id)
       if (eio%proc_num_id(i) == proc_num_id) then
          i_prc = i
          exit FIND_I_PRC
       end if
    end do FIND_I_PRC
    if (i_prc == 0)  call err_index
  contains
    subroutine err_close
      call msg_error ("LHEF: reading events: syntax error in closing tag")
      iostat = 1
    end subroutine
    subroutine err_evt1
      call msg_error ("LHEF: reading events: invalid event tag, &
           &aborting read")
      iostat = 2
    end subroutine err_evt1
    subroutine err_index
      call msg_error ("LHEF: reading events: undefined process ID " &
           // char (str (proc_num_id)) // ", aborting read")
      iostat = 3
    end subroutine err_index
  end subroutine eio_lhef_input_i_prc

\eatline
\nwendcode{}\nwbegindocs{1102}Since we have already read the event information from file, this
input routine can transfer the common-block contents to the event
record.  Also, we read any further information in the event record.

Since LHEF doesn't give this information, we must assume that the MCI
group, term, and channel can all be safely set to 1.  This works if
there is only one MCI group and term.  The channel doesn't matter for
the matrix element.
\nwenddocs{}\nwbegincode{1103}\moddef{EIO LHEF: eio lhef: TBP}\plusendmoddef
  procedure :: input_event => eio_lhef_input_event
\nwendcode{}\nwbegincode{1104}\moddef{EIO LHEF: procedures}\plusendmoddef
  subroutine eio_lhef_input_event (eio, event, iostat)
    class(eio_lhef_t), intent(inout) :: eio
    class(generic_event_t), intent(inout), target :: event
    integer, intent(out) :: iostat
    type(particle_set_t), pointer :: pset
    type(string_t) :: s
    logical :: closing
    iostat = 0
    call event%reset ()
    call event%select (1, 1, 1)
    call hepeup_to_event (event, eio%fallback_model, &
         recover_beams = eio%recover_beams, &
         use_alpha_s = eio%use_alpha_s_from_file, &
         use_scale = eio%use_scale_from_file)
    select case (eio%version)
    case ("1.0")
       call eio%tag_event%read_content (eio%cstream, s, closing = closing)
       if (s /= "" .or. .not. closing)  call err_evt2
    case ("2.0");  call eio%read_event_20 (event)
    case ("3.0");  call eio%read_event_30 (event)
    end select
  contains
    subroutine err_evt2
      call msg_error ("LHEF: reading events: syntax error in event record, &
           &aborting read")
      iostat = 2
    end subroutine err_evt2

  end subroutine eio_lhef_input_event

\eatline
\nwendcode{}\nwbegindocs{1105}\nwdocspar
\subsection{Les Houches Event File: header/footer}
These two routines write the header and footer for the Les Houches
Event File format (LHEF).

The current version writes no information except for the generator
name and version (v.1.0 only).
\nwenddocs{}\nwbegincode{1106}\moddef{EIO LHEF: eio lhef: TBP}\plusendmoddef
  procedure :: write_header => eio_lhef_write_header
  procedure :: write_footer => eio_lhef_write_footer
\nwendcode{}\nwbegincode{1107}\moddef{EIO LHEF: procedures}\plusendmoddef
  subroutine eio_lhef_write_header (eio)
    class(eio_lhef_t), intent(in) :: eio
    integer :: u
    u = given_output_unit (eio%unit);  if (u < 0)  return
    call eio%tag_lhef%write (u);  write (u, *)
    call eio%tag_head%write (u);  write (u, *)
    select case (eio%version)
    case ("1.0")
       write (u, "(2x)", advance = "no")
       call eio%tag_gen_n%write (var_str ("WHIZARD"), u)
       write (u, *)
       write (u, "(2x)", advance = "no")
       call eio%tag_gen_v%write (var_str ("\LA{}Version\RA{}"), u)
       write (u, *)
    end select
    call eio%tag_head%close (u);  write (u, *)
  end subroutine eio_lhef_write_header

  subroutine eio_lhef_write_footer (eio)
    class(eio_lhef_t), intent(in) :: eio
    integer :: u
    u = given_output_unit (eio%unit);  if (u < 0)  return
    call eio%tag_lhef%close (u)
  end subroutine eio_lhef_write_footer

\eatline
\nwendcode{}\nwbegindocs{1108}Reading the header just means finding the tags and ignoring any
contents.  When done, we should stand just after the header tag.
\nwenddocs{}\nwbegincode{1109}\moddef{EIO LHEF: eio lhef: TBP}\plusendmoddef
  procedure :: read_header => eio_lhef_read_header
\nwendcode{}\nwbegincode{1110}\moddef{EIO LHEF: procedures}\plusendmoddef
  subroutine eio_lhef_read_header (eio)
    class(eio_lhef_t), intent(inout) :: eio
    logical :: success, closing
    type(string_t) :: content
    call eio%tag_lhef%read (eio%cstream, success)
    if (.not. success .or. .not. eio%tag_lhef%has_content)  call err_lhef
    if (eio%tag_lhef%get_attribute (1) /= eio%version)  call err_version
    call eio%tag_head%read (eio%cstream, success)
    if (.not. success)  call err_header
    if (eio%tag_head%has_content) then
       SKIP_HEADER_CONTENT: do
          call eio%tag_head%read_content (eio%cstream, content, closing)
          if (closing)  exit SKIP_HEADER_CONTENT
       end do SKIP_HEADER_CONTENT
    end if
  contains
    subroutine err_lhef
      call msg_fatal ("LHEF: LesHouchesEvents tag absent or corrupted")
    end subroutine err_lhef
    subroutine err_header
      call msg_fatal ("LHEF: header tag absent or corrupted")
    end subroutine err_header
    subroutine err_version
       call msg_error ("LHEF: version mismatch: expected " &
            // eio%version // ", found " &
            // char (eio%tag_lhef%get_attribute (1)))
    end subroutine err_version
  end subroutine eio_lhef_read_header

\eatline
\nwendcode{}\nwbegindocs{1111}\nwdocspar
\subsection{Version-Specific Code: 1.0}
In version 1.0, the init tag contains just HEPRUP data.  While a
{\tt{}cstream} is connected to the input unit, we bypass it temporarily
for the purpose of reading the HEPRUP contents.  This is consistent
with the LHEF standard.

This routine does not read the closing tag of the init block.
\nwenddocs{}\nwbegincode{1112}\moddef{EIO LHEF: eio lhef: TBP}\plusendmoddef
  procedure :: read_init_10 => eio_lhef_read_init_10
\nwendcode{}\nwbegincode{1113}\moddef{EIO LHEF: procedures}\plusendmoddef
  subroutine eio_lhef_read_init_10 (eio, data)
    class(eio_lhef_t), intent(in) :: eio
    type(event_sample_data_t), intent(out) :: data
    integer :: n_proc, i
    call heprup_read_lhef (eio%unit)
    call heprup_get_run_parameters (n_processes = n_proc)
    call data%init (n_proc)
    data%n_beam = 2
    call heprup_get_run_parameters ( &
         unweighted = data%unweighted, &
         negative_weights = data%negative_weights, &
         beam_pdg = data%pdg_beam, &
         beam_energy = data%energy_beam)
    if (data%unweighted) then
       data%norm_mode = NORM_UNIT
    else
       data%norm_mode = NORM_SIGMA
    end if
    do i = 1, n_proc
       call heprup_get_process_parameters (i, &
            process_id = data%proc_num_id(i), &
            cross_section = data%cross_section(i), &
            error = data%error(i))
    end do
  end subroutine eio_lhef_read_init_10
  
\eatline
\nwendcode{}\nwbegindocs{1114}\nwdocspar
\subsection{Version-Specific Code: 2.0}
This is the init information for the 2.0 format, after the HEPRUP
data.  We have the following tags:
\begin{itemize}
\item \texttt{generator}  Generator name and version.
\item \texttt{xsecinfo}  Cross section and weights data.  We have the
  total cross section and number of events (assuming that the event
  file is intact), but information on minimum and maximum weights is
  not available before the file is complete.  We just write the
  mandatory tags.  (Note that the default values of the other tags
  describe a uniform unit weight, but we can determine most values
  only after the sample is complete.)
\item \texttt{cutsinfo}  This optional tag is too specific to represent the
  possibilities of WHIZARD, so we skip it.
\item \texttt{procinfo}  This optional tag is useful for giving
  details of NLO calculations.  Skipped.
\item \texttt{mergetype}  Optional, also not applicable.
\end{itemize}
\nwenddocs{}\nwbegincode{1115}\moddef{EIO LHEF: eio lhef: TBP}\plusendmoddef
  procedure :: write_init_20 => eio_lhef_write_init_20
\nwendcode{}\nwbegincode{1116}\moddef{EIO LHEF: procedures}\plusendmoddef
  subroutine eio_lhef_write_init_20 (eio, data)
    class(eio_lhef_t), intent(in) :: eio
    type(event_sample_data_t), intent(in) :: data
    integer :: u
    u = eio%unit
    call eio%tag_generator%write (u)
    write (u, "(A)", advance="no")  "WHIZARD"
    call eio%tag_generator%close (u);  write (u, *)
    call eio%tag_xsecinfo%write (u);  write (u, *)
  end subroutine eio_lhef_write_init_20
    
\eatline
\nwendcode{}\nwbegindocs{1117}When reading the init block, we first call the 1.0 routine that
fills HEPRUP.  Then we consider the possible tags.  Only the
\texttt{generator} and \texttt{xsecinfo} tags are of interest.  We
skip everything else except for the closing tag.
\nwenddocs{}\nwbegincode{1118}\moddef{EIO LHEF: eio lhef: TBP}\plusendmoddef
  procedure :: read_init_20 => eio_lhef_read_init_20
\nwendcode{}\nwbegincode{1119}\moddef{EIO LHEF: procedures}\plusendmoddef
  subroutine eio_lhef_read_init_20 (eio, data)
    class(eio_lhef_t), intent(inout) :: eio
    type(event_sample_data_t), intent(out) :: data
    real(default), parameter :: pb_per_fb = 1.e-3_default
    type(string_t) :: content
    logical :: found, closing
    call eio_lhef_read_init_10 (eio, data)
    SCAN_INIT_TAGS: do
       call eio%tag_generator%read (eio%cstream, found)
       if (found) then
          if (.not. eio%tag_generator%has_content)  call err_generator
          call eio%tag_generator%read_content (eio%cstream, content, closing)
          call msg_message ("LHEF: Event file has been generated by " &
               // char (content) // " " &
               // char (eio%tag_generator%get_attribute (1)))
          cycle SCAN_INIT_TAGS
       end if
       call eio%tag_xsecinfo%read (eio%cstream, found)
       if (found) then
          if (eio%tag_xsecinfo%has_content)  call err_xsecinfo
          cycle SCAN_INIT_TAGS
       end if
       call eio%tag_init%read_content (eio%cstream, content, closing)
       if (closing) then
          if (content /= "")  call err_init
          exit SCAN_INIT_TAGS
       end if
    end do SCAN_INIT_TAGS
    data%n_evt = &
         read_ival (eio%tag_xsecinfo%get_attribute (1))
    data%total_cross_section = &
         read_rval (eio%tag_xsecinfo%get_attribute (2)) / pb_per_fb
  contains
    subroutine err_generator
      call msg_fatal ("LHEF: invalid generator tag")
    end subroutine err_generator
    subroutine err_xsecinfo
      call msg_fatal ("LHEF: invalid xsecinfo tag")
    end subroutine err_xsecinfo
    subroutine err_init
      call msg_fatal ("LHEF: syntax error after init tag")
    end subroutine err_init
  end subroutine eio_lhef_read_init_20
  
\eatline
\nwendcode{}\nwbegindocs{1120}This is additional event-specific information for the 2.0 format,
after the HEPEUP data.  We can specify weights, starting from the
master weight and adding alternative weights.  The alternative weights
are collected in a common tag.
\nwenddocs{}\nwbegincode{1121}\moddef{EIO LHEF: eio lhef: TBP}\plusendmoddef
  procedure :: write_event_20 => eio_lhef_write_event_20
\nwendcode{}\nwbegincode{1122}\moddef{EIO LHEF: procedures}\plusendmoddef
  subroutine eio_lhef_write_event_20 (eio, event)
    class(eio_lhef_t), intent(in) :: eio
    class(generic_event_t), intent(in) :: event
    type(string_t) :: s
    integer :: i, u
    u = eio%unit
    if (eio%write_sqme_ref) then
       s = str (event%get_sqme_ref ())
       call eio%tag_sqme_ref%write (s, u);  write (u, *)
    end if
    if (eio%write_sqme_prc) then
       s = str (event%get_sqme_prc ())
       call eio%tag_sqme_prc%write (s, u);  write (u, *)
    end if
    if (eio%n_alt > 0) then
       if (eio%write_sqme_alt) then
          s = str (event%get_sqme_alt(1))
          do i = 2, eio%n_alt
             s = s // " " // str (event%get_sqme_alt(i));  write (u, *)
          end do
          call eio%tag_sqme_alt(1)%write (s, u)
       end if
       s = str (event%get_weight_alt(1))
       do i = 2, eio%n_alt
          s = s // " " // str (event%get_weight_alt(i));  write (u, *)
       end do
       call eio%tag_wgts_alt(1)%write (s, u)
    end if
  end subroutine eio_lhef_write_event_20
    
\eatline
\nwendcode{}\nwbegindocs{1123}Read extra event data.  If there is a weight entry labeled {\tt{}sqme{\char95}prc}, we
take this as the squared matrix-element value (the new
\emph{reference} value {\tt{}sqme{\char95}ref}). Other tags, including
tags written by the above writer, are skipped.
\nwenddocs{}\nwbegincode{1124}\moddef{EIO LHEF: eio lhef: TBP}\plusendmoddef
  procedure :: read_event_20 => eio_lhef_read_event_20
\nwendcode{}\nwbegincode{1125}\moddef{EIO LHEF: procedures}\plusendmoddef
  subroutine eio_lhef_read_event_20 (eio, event)
    class(eio_lhef_t), intent(inout) :: eio
    class(generic_event_t), intent(inout) :: event
    type(string_t) :: content
    logical :: found, closing
    SCAN_EVENT_TAGS: do
       call eio%tag_weight%read (eio%cstream, found)
       if (found) then
          if (.not. eio%tag_weight%has_content)  call err_weight
          call eio%tag_weight%read_content (eio%cstream, content, closing)
          if (.not. closing)  call err_weight
          if (eio%tag_weight%get_attribute (1) == "sqme_prc") then
             call event%set_sqme_ref (read_rval (content))
          end if
          cycle SCAN_EVENT_TAGS
       end if
       call eio%tag_event%read_content (eio%cstream, content, closing)
       if (closing) then
          if (content /= "")  call err_event
          exit SCAN_EVENT_TAGS
       end if
    end do SCAN_EVENT_TAGS
  contains
    subroutine err_weight
      call msg_fatal ("LHEF: invalid weight tag in event record")
    end subroutine err_weight
    subroutine err_event
      call msg_fatal ("LHEF: syntax error after event tag")
    end subroutine err_event
  end subroutine eio_lhef_read_event_20
    
\eatline
\nwendcode{}\nwbegindocs{1126}\nwdocspar
\subsection{Version-Specific Code: 3.0}
This is the init information for the 3.0 format, after the HEPRUP
data.  We have the following tags:
\begin{itemize}
\item \texttt{generator}  Generator name and version.
\item \texttt{xsecinfo}  Cross section and weights data.  We have the
  total cross section and number of events (assuming that the event
  file is intact), but information on minimum and maximum weights is
  not available before the file is complete.  We just write the
  mandatory tags.  (Note that the default values of the other tags
  describe a uniform unit weight, but we can determine most values
  only after the sample is complete.)
\item \texttt{cutsinfo}  This optional tag is too specific to represent the
  possibilities of WHIZARD, so we skip it.
\item \texttt{procinfo}  This optional tag is useful for giving
  details of NLO calculations.  Skipped.
\item \texttt{weightinfo}  Determine the meaning of optional weights, whose
  values are given in the event record.
\end{itemize}
\nwenddocs{}\nwbegincode{1127}\moddef{EIO LHEF: eio lhef: TBP}\plusendmoddef
  procedure :: write_init_30 => eio_lhef_write_init_30
\nwendcode{}\nwbegincode{1128}\moddef{EIO LHEF: procedures}\plusendmoddef
  subroutine eio_lhef_write_init_30 (eio, data)
    class(eio_lhef_t), intent(in) :: eio
    type(event_sample_data_t), intent(in) :: data
    real(default), parameter :: pb_per_fb = 1.e-3_default
    integer :: u, i
    u = given_output_unit (eio%unit)
    call eio%tag_generator%write (u)
    write (u, "(A)", advance="no")  "WHIZARD"
    call eiO%tag_generator%close (u);  write (u, *)
    call eio%tag_xsecinfo%write (u);  write (u, *)
    if (eio%write_sqme_ref) then
       call eio%tag_sqme_ref%write (u);  write (u, *)
    end if
    if (eio%write_sqme_prc) then
       call eio%tag_sqme_prc%write (u);  write (u, *)
    end if
    if (eio%write_sqme_alt) then
       do i = 1, eio%n_alt
          call eio%tag_sqme_alt(i)%write (u);  write (u, *)
       end do
    end if
    do i = 1, eio%n_alt
       call eio%tag_wgts_alt(i)%write (u);  write (u, *)
    end do
  end subroutine eio_lhef_write_init_30
    
\eatline
\nwendcode{}\nwbegindocs{1129}When reading the init block, we first call the 1.0 routine that
fills HEPRUP.  Then we consider the possible tags.  Only the
\texttt{generator} and \texttt{xsecinfo} tags are of interest.  We
skip everything else except for the closing tag.
\nwenddocs{}\nwbegincode{1130}\moddef{EIO LHEF: eio lhef: TBP}\plusendmoddef
  procedure :: read_init_30 => eio_lhef_read_init_30
\nwendcode{}\nwbegincode{1131}\moddef{EIO LHEF: procedures}\plusendmoddef
  subroutine eio_lhef_read_init_30 (eio, data)
    class(eio_lhef_t), intent(inout) :: eio
    type(event_sample_data_t), intent(out) :: data
    real(default), parameter :: pb_per_fb = 1.e-3_default
    type(string_t) :: content
    logical :: found, closing
    integer :: n_weightinfo
    call eio_lhef_read_init_10 (eio, data)
    n_weightinfo = 0
    eio%i_weight_sqme = 0
    SCAN_INIT_TAGS: do
       call eio%tag_generator%read (eio%cstream, found)
       if (found) then
          if (.not. eio%tag_generator%has_content)  call err_generator
          call eio%tag_generator%read_content (eio%cstream, content, closing)
          call msg_message ("LHEF: Event file has been generated by " &
               // char (content) // " " &
               // char (eio%tag_generator%get_attribute (1)))
          cycle SCAN_INIT_TAGS
       end if
       call eio%tag_xsecinfo%read (eio%cstream, found)
       if (found) then
          if (eio%tag_xsecinfo%has_content)  call err_xsecinfo
          cycle SCAN_INIT_TAGS
       end if
       call eio%tag_weightinfo%read (eio%cstream, found)
       if (found) then
          if (eio%tag_weightinfo%has_content)  call err_xsecinfo
          n_weightinfo = n_weightinfo + 1
          if (eio%tag_weightinfo%get_attribute (1) == "sqme_prc") then
             eio%i_weight_sqme = n_weightinfo
          end if
          cycle SCAN_INIT_TAGS
       end if
       call eio%tag_init%read_content (eio%cstream, content, closing)
       if (closing) then
          if (content /= "")  call err_init
          exit SCAN_INIT_TAGS
       end if
    end do SCAN_INIT_TAGS
    data%n_evt = &
         read_ival (eio%tag_xsecinfo%get_attribute (1))
    data%total_cross_section = &
         read_rval (eio%tag_xsecinfo%get_attribute (2)) / pb_per_fb
  contains
    subroutine err_generator
      call msg_fatal ("LHEF: invalid generator tag")
    end subroutine err_generator
    subroutine err_xsecinfo
      call msg_fatal ("LHEF: invalid xsecinfo tag")
    end subroutine err_xsecinfo
    subroutine err_init
      call msg_fatal ("LHEF: syntax error after init tag")
    end subroutine err_init
  end subroutine eio_lhef_read_init_30
  
\eatline
\nwendcode{}\nwbegindocs{1132}This is additional event-specific information for the 3.0 format,
after the HEPEUP data.  We can specify weights, starting from the
master weight and adding alternative weights.  The weight tags are
already allocated, so we just have to transfer the weight values to
strings, assemble them and write them to file.  All weights are
collected in a single tag.

Note: If efficiency turns out to be an issue, we may revert to
traditional character buffer writing.  However, we need to know the
maximum length.
\nwenddocs{}\nwbegincode{1133}\moddef{EIO LHEF: eio lhef: TBP}\plusendmoddef
  procedure :: write_event_30 => eio_lhef_write_event_30
\nwendcode{}\nwbegincode{1134}\moddef{EIO LHEF: procedures}\plusendmoddef
  subroutine eio_lhef_write_event_30 (eio, event)
    class(eio_lhef_t), intent(in) :: eio
    class(generic_event_t), intent(in) :: event
    type(string_t) :: s
    integer :: u, i
    u = eio%unit
    s = ""
    if (eio%write_sqme_ref) then
       s = s // str (event%get_sqme_ref ()) // " "
    end if
    if (eio%write_sqme_prc) then
       s = s // str (event%get_sqme_prc ()) // " "
    end if
    if (eio%n_alt > 0) then
       if (eio%write_sqme_alt) then
          s = s // str (event%get_sqme_alt(1)) // " "
          do i = 2, eio%n_alt
             s = s // str (event%get_sqme_alt(i)) // " "
          end do
       end if
       s = s // str (event%get_weight_alt(1)) // " "
       do i = 2, eio%n_alt
          s = s // str (event%get_weight_alt(i)) // " "
       end do
    end if
    if (len_trim (s) > 0) then
       call eio%tag_weights%write (trim (s), u);  write (u, *)
    end if
  end subroutine eio_lhef_write_event_30
    
\eatline
\nwendcode{}\nwbegindocs{1135}Read extra event data.  If there is a {\tt{}weights} tag and if there
was a {\tt{}weightinfo} entry labeled {\tt{}sqme{\char95}prc}, we extract the
corresponding entry from the weights string and store this as the
event's  squared matrix-element value.  Other tags, including
tags written by the above writer, are skipped.
\nwenddocs{}\nwbegincode{1136}\moddef{EIO LHEF: eio lhef: TBP}\plusendmoddef
  procedure :: read_event_30 => eio_lhef_read_event_30
\nwendcode{}\nwbegincode{1137}\moddef{EIO LHEF: procedures}\plusendmoddef
  subroutine eio_lhef_read_event_30 (eio, event)
    class(eio_lhef_t), intent(inout) :: eio
    class(generic_event_t), intent(inout) :: event
    type(string_t) :: content, string
    logical :: found, closing
    integer :: i
    SCAN_EVENT_TAGS: do
       call eio%tag_weights%read (eio%cstream, found)
       if (found) then
          if (.not. eio%tag_weights%has_content)  call err_weights
          call eio%tag_weights%read_content (eio%cstream, content, closing)
          if (.not. closing)  call err_weights
          if (eio%i_weight_sqme > 0) then
             SCAN_WEIGHTS: do i = 1, eio%i_weight_sqme
                call split (content, string, " ")
                content = adjustl (content)
                if (i == eio%i_weight_sqme) then
                   call event%set_sqme_ref (read_rval (string))
                   exit SCAN_WEIGHTS
                end if
             end do SCAN_WEIGHTS
          end if
          cycle SCAN_EVENT_TAGS
       end if
       call eio%tag_event%read_content (eio%cstream, content, closing)
       if (closing) then
          if (content /= "")  call err_event
          exit SCAN_EVENT_TAGS
       end if
    end do SCAN_EVENT_TAGS
  contains
    subroutine err_weights
      call msg_fatal ("LHEF: invalid weights tag in event record")
    end subroutine err_weights
    subroutine err_event
      call msg_fatal ("LHEF: syntax error after event tag")
    end subroutine err_event
  end subroutine eio_lhef_read_event_30
    
\eatline
\nwendcode{}\nwbegindocs{1138}\nwdocspar
\subsection{Unit tests}
\nwenddocs{}\nwbegincode{1139}\moddef{EIO LHEF: public}\plusendmoddef
  public :: eio_lhef_test
\nwendcode{}\nwbegincode{1140}\moddef{EIO LHEF: tests}\endmoddef
  subroutine eio_lhef_test (u, results)
    integer, intent(in) :: u
    type(test_results_t), intent(inout) :: results
  \LA{}EIO LHEF: execute tests\RA{}
  end subroutine eio_lhef_test
  
\eatline
\nwendcode{}\nwbegindocs{1141}\nwdocspar
\subsubsection{Version 1.0 Output}
We test the implementation of all I/O methods.  We start with output
according to version 1.0.
\nwenddocs{}\nwbegincode{1142}\moddef{EIO LHEF: execute tests}\endmoddef
  call test (eio_lhef_1, "eio_lhef_1", &
       "write version 1.0", &
       u, results)
\nwendcode{}\nwbegincode{1143}\moddef{EIO LHEF: tests}\plusendmoddef
  subroutine eio_lhef_1 (u)
    integer, intent(in) :: u
    class(generic_event_t), pointer :: event
    type(event_sample_data_t) :: data
    class(eio_t), allocatable :: eio
    type(string_t) :: sample
    integer :: u_file, iostat
    character(80) :: buffer

    write (u, "(A)")  "* Test output: eio_lhef_1"
    write (u, "(A)")  "*   Purpose: generate an event and write weight to file"
    write (u, "(A)")

    write (u, "(A)")  "* Initialize test process"
 
    call eio_prepare_test (event, unweighted = .false.)

    call data%init (1)
    data%n_evt = 1
    data%n_beam = 2
    data%unweighted = .true.
    data%norm_mode = NORM_UNIT
    data%pdg_beam = 25
    data%energy_beam = 500
    data%proc_num_id = [42]
    data%cross_section(1) = 100
    data%error(1) = 1
    data%total_cross_section = sum (data%cross_section)

    write (u, "(A)")
    write (u, "(A)")  "* Generate and write an event"
    write (u, "(A)")
 
    sample = "eio_lhef_1"
    
    allocate (eio_lhef_t :: eio)
    select type (eio)
    type is (eio_lhef_t)
       call eio%set_parameters ()
    end select
    
    call eio%init_out (sample, data)
    call event%generate (1, [0._default, 0._default])

    call eio%output (event, i_prc = 1)
    call eio%write (u)
    call eio%final ()

    write (u, "(A)")
    write (u, "(A)")  "* File contents:"
    write (u, "(A)")

    u_file = free_unit ()
    open (u_file, file = char (sample // "." // eio%extension), &
         action = "read", status = "old")
    do
       read (u_file, "(A)", iostat = iostat)  buffer
       if (buffer(1:21) == "  <generator_version>")  buffer = "[...]"
       if (iostat /= 0)  exit
       write (u, "(A)") trim (buffer)
    end do
    close (u_file)
    
    write (u, "(A)")
    write (u, "(A)")  "* Reset data"
    write (u, "(A)")
 
    deallocate (eio)
    allocate (eio_lhef_t :: eio)
    select type (eio)
    type is (eio_lhef_t)
       call eio%set_parameters ()
    end select
    
    select type (eio)
    type is (eio_lhef_t)
       call eio%set_parameters (keep_beams = .true.)
    end select
    call eio%write (u)

    write (u, "(A)")
    write (u, "(A)")  "* Cleanup"
 
    call eio_cleanup_test (event)
    
    write (u, "(A)")
    write (u, "(A)")  "* Test output end: eio_lhef_1"
    
  end subroutine eio_lhef_1
  
\eatline
\nwendcode{}\nwbegindocs{1144}\nwdocspar
\subsubsection{Version 2.0 Output}
Version 2.0 has added a lot of options to the LHEF format.  We
implement some of them.
\nwenddocs{}\nwbegincode{1145}\moddef{EIO LHEF: execute tests}\plusendmoddef
  call test (eio_lhef_2, "eio_lhef_2", &
       "write version 2.0", &
       u, results)
\nwendcode{}\nwbegincode{1146}\moddef{EIO LHEF: tests}\plusendmoddef
  subroutine eio_lhef_2 (u)
    integer, intent(in) :: u
    type(model_data_t), target :: model
    class(generic_event_t), pointer :: event
    type(event_sample_data_t) :: data
    class(eio_t), allocatable :: eio
    type(string_t) :: sample
    integer :: u_file, iostat
    character(80) :: buffer

    write (u, "(A)")  "* Test output: eio_lhef_2"
    write (u, "(A)")  "*   Purpose: generate an event and write weight to file"
    write (u, "(A)")

    write (u, "(A)")  "* Initialize test process"
 
    call eio_prepare_test (event, unweighted = .false.)

    call data%init (1)
    data%unweighted = .false.
    data%norm_mode = NORM_SIGMA
    data%n_evt = 1
    data%n_beam = 2
    data%pdg_beam = 25
    data%energy_beam = 500
    data%proc_num_id = [42]
    data%cross_section(1) = 100
    data%error(1) = 1
    data%total_cross_section = sum (data%cross_section)

    write (u, "(A)")
    write (u, "(A)")  "* Generate and write an event"
    write (u, "(A)")
 
    sample = "eio_lhef_2"
    
    allocate (eio_lhef_t :: eio)
    select type (eio)
    type is (eio_lhef_t)
       call eio%set_parameters (version = "2.0", write_sqme_prc = .true.)
    end select
    
    call eio%init_out (sample, data)
    call event%generate (1, [0._default, 0._default])


    call eio%output (event, i_prc = 1)
    call eio%write (u)
    call eio%final ()

    write (u, "(A)")
    write (u, "(A)")  "* File contents:"
    write (u, "(A)")

    u_file = free_unit ()
    open (u_file, file = char (sample // "." // eio%extension), &
         action = "read", status = "old")
    do
       read (u_file, "(A)", iostat = iostat)  buffer
       if (buffer(1:10) == "<generator")  buffer = "[...]"
       if (iostat /= 0)  exit
       write (u, "(A)") trim (buffer)
    end do
    close (u_file)
    
    write (u, "(A)")
    write (u, "(A)")  "* Cleanup"
 
    call eio_cleanup_test (event)

    write (u, "(A)")
    write (u, "(A)")  "* Test output end: eio_lhef_2"
    
  end subroutine eio_lhef_2
  
\eatline
\nwendcode{}\nwbegindocs{1147}\nwdocspar
\subsubsection{Version 3.0 Output}
Version 3.0 is an update which removes some tags (which we didn't use anyway)
and suggests a new treatment of weights.
\nwenddocs{}\nwbegincode{1148}\moddef{EIO LHEF: execute tests}\plusendmoddef
  call test (eio_lhef_3, "eio_lhef_3", &
       "write version 3.0", &
       u, results)
\nwendcode{}\nwbegincode{1149}\moddef{EIO LHEF: tests}\plusendmoddef
  subroutine eio_lhef_3 (u)
    integer, intent(in) :: u
    class(generic_event_t), pointer :: event
    type(event_sample_data_t) :: data
    class(eio_t), allocatable :: eio
    type(string_t) :: sample
    integer :: u_file, iostat
    character(80) :: buffer

    write (u, "(A)")  "* Test output: eio_lhef_3"
    write (u, "(A)")  "*   Purpose: generate an event and write weight to file"
    write (u, "(A)")

    write (u, "(A)")  "* Initialize test process"
 
    call eio_prepare_test (event, unweighted = .false.)

    call data%init (1)
    data%unweighted = .false.
    data%norm_mode = NORM_SIGMA
    data%n_evt = 1
    data%n_beam = 2
    data%pdg_beam = 25
    data%energy_beam = 500
    data%proc_num_id = [42]
    data%cross_section(1) = 100
    data%error(1) = 1
    data%total_cross_section = sum (data%cross_section)

    write (u, "(A)")
    write (u, "(A)")  "* Generate and write an event"
    write (u, "(A)")
 
    sample = "eio_lhef_3"
    
    allocate (eio_lhef_t :: eio)
    select type (eio)
    type is (eio_lhef_t)
       call eio%set_parameters (version = "3.0", write_sqme_prc = .true.)
    end select
    
    call eio%init_out (sample, data)
    call event%generate (1, [0._default, 0._default])


    call eio%output (event, i_prc = 1)
    call eio%write (u)
    call eio%final ()

    write (u, "(A)")
    write (u, "(A)")  "* File contents:"
    write (u, "(A)")

    u_file = free_unit ()
    open (u_file, file = char (sample // ".lhe"), &
         action = "read", status = "old")
    do
       read (u_file, "(A)", iostat = iostat)  buffer
       if (buffer(1:10) == "<generator")  buffer = "[...]"
       if (iostat /= 0)  exit
       write (u, "(A)") trim (buffer)
    end do
    close (u_file)
    
    write (u, "(A)")
    write (u, "(A)")  "* Cleanup"
 
    call eio_cleanup_test (event)

    write (u, "(A)")
    write (u, "(A)")  "* Test output end: eio_lhef_3"
    
  end subroutine eio_lhef_3
  
\eatline
\nwendcode{}\nwbegindocs{1150}\nwdocspar
\subsubsection{Version 1.0 Input}
Check input of a version-1.0 conforming LHEF file.
\nwenddocs{}\nwbegincode{1151}\moddef{EIO LHEF: execute tests}\plusendmoddef
  call test (eio_lhef_4, "eio_lhef_4", &
       "read version 1.0", &
       u, results)
\nwendcode{}\nwbegincode{1152}\moddef{EIO LHEF: tests}\plusendmoddef
  subroutine eio_lhef_4 (u)
    integer, intent(in) :: u
    class(model_data_t), pointer :: fallback_model
    class(generic_event_t), pointer :: event
    type(event_sample_data_t) :: data
    class(eio_t), allocatable :: eio
    type(string_t) :: sample
    integer :: u_file, iostat, i_prc

    write (u, "(A)")  "* Test output: eio_lhef_4"
    write (u, "(A)")  "*   Purpose: read a LHEF 1.0 file"
    write (u, "(A)")

    write (u, "(A)")  "* Write a LHEF data file"
    write (u, "(A)")
 
    u_file = free_unit ()
    sample = "eio_lhef_4"
    open (u_file, file = char (sample // ".lhe"), &
         status = "replace", action = "readwrite")
    
    write (u_file, "(A)")  '<LesHouchesEvents version="1.0">'
    write (u_file, "(A)")  '<header>'
    write (u_file, "(A)")  '  <arbitrary_tag opt="foo">content</arbitrary_tag>'
    write (u_file, "(A)")  '  Text'
    write (u_file, "(A)")  '  <another_tag />'
    write (u_file, "(A)")  '</header>'
    write (u_file, "(A)")  '<init>'
    write (u_file, "(A)")  ' 25 25  5.0000000000E+02  5.0000000000E+02 &
         & -1 -1 -1 -1 3 1'
    write (u_file, "(A)")  '  1.0000000000E-01  1.0000000000E-03 &
         & 1.0000000000E+00 42'
    write (u_file, "(A)")  '</init>'
    write (u_file, "(A)")  '<event>'
    write (u_file, "(A)")  ' 4 42  3.0574068604E+08  1.0000000000E+03 &
         & -1.0000000000E+00 -1.0000000000E+00'
    write (u_file, "(A)")  ' 25 -1 0 0 0 0  0.0000000000E+00  0.0000000000E+00 &
         & 4.8412291828E+02  5.0000000000E+02  1.2500000000E+02 &
         & 0.0000000000E+00  9.0000000000E+00'
    write (u_file, "(A)")  ' 25 -1 0 0 0 0  0.0000000000E+00  0.0000000000E+00 &
         &-4.8412291828E+02  5.0000000000E+02  1.2500000000E+02 &
         & 0.0000000000E+00  9.0000000000E+00'
    write (u_file, "(A)")  ' 25 1 1 2 0 0 -1.4960220911E+02 -4.6042825611E+02 &
         & 0.0000000000E+00  5.0000000000E+02  1.2500000000E+02 &
         & 0.0000000000E+00  9.0000000000E+00'
    write (u_file, "(A)")  ' 25 1 1 2 0 0  1.4960220911E+02  4.6042825611E+02 &
         & 0.0000000000E+00  5.0000000000E+02  1.2500000000E+02 &
         & 0.0000000000E+00  9.0000000000E+00'
    write (u_file, "(A)")  '</event>'
    write (u_file, "(A)")  '</LesHouchesEvents>'
    close (u_file)
    

    write (u, "(A)")  "* Initialize test process"
    write (u, "(A)")

    call eio_prepare_fallback_model (fallback_model)
    call eio_prepare_test (event, unweighted = .false.)

    allocate (eio_lhef_t :: eio)
    select type (eio)
    type is (eio_lhef_t)
       call eio%set_parameters (recover_beams = .false.)
    end select
    call eio%set_fallback_model (fallback_model)
    
    call data%init (1)
    data%n_beam = 2
    data%unweighted = .true.
    data%norm_mode = NORM_UNIT
    data%pdg_beam = 25
    data%energy_beam = 500
    data%proc_num_id = [42]
    call data%write (u)
    write (u, *)

    write (u, "(A)")  "* Initialize and read header"
    write (u, "(A)")

    call eio%init_in (sample, data)
    call eio%write (u)
    
    write (u, *)
    
    select type (eio)
    type is (eio_lhef_t)
       call eio%tag_lhef%write (u);  write (u, *)
    end select

    write (u, *)
    call data%write (u)

    write (u, "(A)")
    write (u, "(A)")  "* Read event"
    write (u, "(A)")

    call eio%input_i_prc (i_prc, iostat)
    
    select type (eio)
    type is (eio_lhef_t)
       write (u, "(A,I0,A,I0)")  "Found process #", i_prc, &
            " with ID = ", eio%proc_num_id(i_prc)
    end select
   
    call eio%input_event (event, iostat)

    call event%write (u)

    write (u, "(A)")
    write (u, "(A)")  "* Read closing"
    write (u, "(A)")

    call eio%input_i_prc (i_prc, iostat)
    write (u, "(A,I0)")  "iostat = ", iostat
    
    write (u, "(A)")
    write (u, "(A)")  "* Cleanup"
 
    call eio%final ()

    call eio_cleanup_test (event)
    call eio_cleanup_fallback_model (fallback_model)

    write (u, "(A)")
    write (u, "(A)")  "* Test output end: eio_lhef_4"
    
  end subroutine eio_lhef_4
  
\eatline
\nwendcode{}\nwbegindocs{1153}\nwdocspar
\subsubsection{Version 2.0 Input}
Check input of a version-2.0 conforming LHEF file.
\nwenddocs{}\nwbegincode{1154}\moddef{EIO LHEF: execute tests}\plusendmoddef
  call test (eio_lhef_5, "eio_lhef_5", &
       "read version 2.0", &
       u, results)
\nwendcode{}\nwbegincode{1155}\moddef{EIO LHEF: tests}\plusendmoddef
  subroutine eio_lhef_5 (u)
    integer, intent(in) :: u
    class(model_data_t), pointer :: fallback_model
    class(generic_event_t), pointer :: event
    type(event_sample_data_t) :: data
    class(eio_t), allocatable :: eio
    type(string_t) :: sample
    integer :: u_file, iostat, i_prc

    write (u, "(A)")  "* Test output: eio_lhef_5"
    write (u, "(A)")  "*   Purpose: read a LHEF 2.0 file"
    write (u, "(A)")

    write (u, "(A)")  "* Write a LHEF data file"
    write (u, "(A)")
 
    u_file = free_unit ()
    sample = "eio_lhef_5"
    open (u_file, file = char (sample // ".lhe"), &
         status = "replace", action = "readwrite")
    
    write (u_file, "(A)")  '<LesHouchesEvents version="2.0">'
    write (u_file, "(A)")  '<header>'
    write (u_file, "(A)")  '</header>'
    write (u_file, "(A)")  '<init>'
    write (u_file, "(A)")  ' 25 25  5.0000000000E+02  5.0000000000E+02 &
         &-1 -1 -1 -1 4 1'
    write (u_file, "(A)")  '  1.0000000000E-01  1.0000000000E-03 &
         & 0.0000000000E+00 42'
    write (u_file, "(A)")  '<generator version="2.2.3">WHIZARD&
         &</generator>'
    write (u_file, "(A)")  '<xsecinfo neve="1" totxsec="1.0000000000E-01" />'
    write (u_file, "(A)")  '</init>'
    write (u_file, "(A)")  '<event>'
    write (u_file, "(A)")  ' 4 42  3.0574068604E+08  1.0000000000E+03 &
         &-1.0000000000E+00 -1.0000000000E+00'
    write (u_file, "(A)")  ' 25 -1 0 0 0 0  0.0000000000E+00 &
         & 0.0000000000E+00  4.8412291828E+02  5.0000000000E+02 &
         & 1.2500000000E+02  0.0000000000E+00  9.0000000000E+00'
    write (u_file, "(A)")  ' 25 -1 0 0 0 0  0.0000000000E+00 &
         & 0.0000000000E+00 -4.8412291828E+02  5.0000000000E+02 &
         & 1.2500000000E+02  0.0000000000E+00  9.0000000000E+00'
    write (u_file, "(A)")  ' 25 1 1 2 0 0 -1.4960220911E+02 &
         &-4.6042825611E+02  0.0000000000E+00  5.0000000000E+02 &
         & 1.2500000000E+02  0.0000000000E+00  9.0000000000E+00'
    write (u_file, "(A)")  ' 25 1 1 2 0 0  1.4960220911E+02 &
         & 4.6042825611E+02  0.0000000000E+00  5.0000000000E+02 &
         & 1.2500000000E+02  0.0000000000E+00  9.0000000000E+00'
    write (u_file, "(A)")  '<weight name="sqme_prc">1.0000000000E+00</weight>'
    write (u_file, "(A)")  '</event>'
    write (u_file, "(A)")  '</LesHouchesEvents>'
    close (u_file)

    write (u, "(A)")  "* Initialize test process"
    write (u, "(A)")

    call eio_prepare_fallback_model (fallback_model)
    call eio_prepare_test (event, unweighted = .false.)

    allocate (eio_lhef_t :: eio)
    select type (eio)
    type is (eio_lhef_t)
       call eio%set_parameters (version = "2.0", recover_beams = .false.)
    end select
    call eio%set_fallback_model (fallback_model)
    
    call data%init (1)
    data%unweighted = .false.
    data%norm_mode = NORM_SIGMA
    data%n_beam = 2
    data%pdg_beam = 25
    data%energy_beam = 500
    data%proc_num_id = [42]
    call data%write (u)
    write (u, *)

    write (u, "(A)")  "* Initialize and read header"
    write (u, "(A)")

    call eio%init_in (sample, data)
    call eio%write (u)
    
    write (u, *)
    
    select type (eio)
    type is (eio_lhef_t)
       call eio%tag_lhef%write (u);  write (u, *)
    end select

    write (u, *)
    call data%write (u)

    write (u, "(A)")
    write (u, "(A)")  "* Read event"
    write (u, "(A)")

    call eio%input_i_prc (i_prc, iostat)
    
    select type (eio)
    type is (eio_lhef_t)
       write (u, "(A,I0,A,I0)")  "Found process #", i_prc, &
            " with ID = ", eio%proc_num_id(i_prc)
    end select
   
    call eio%input_event (event, iostat)

    call event%write (u)

    write (u, "(A)")
    write (u, "(A)")  "* Read closing"
    write (u, "(A)")

    call eio%input_i_prc (i_prc, iostat)
    write (u, "(A,I0)")  "iostat = ", iostat
    
    write (u, "(A)")
    write (u, "(A)")  "* Cleanup"
 
    call eio%final ()

    call eio_cleanup_test (event)
    call eio_cleanup_fallback_model (fallback_model)
    
    write (u, "(A)")
    write (u, "(A)")  "* Test output end: eio_lhef_5"
    
  end subroutine eio_lhef_5
  
\eatline
\nwendcode{}\nwbegindocs{1156}\nwdocspar
\subsubsection{Version 3.0 Input}
Check input of a version-3.0 conforming LHEF file.
\nwenddocs{}\nwbegincode{1157}\moddef{EIO LHEF: execute tests}\plusendmoddef
  call test (eio_lhef_6, "eio_lhef_6", &
       "read version 3.0", &
       u, results)
\nwendcode{}\nwbegincode{1158}\moddef{EIO LHEF: tests}\plusendmoddef
  subroutine eio_lhef_6 (u)
    integer, intent(in) :: u
    class(model_data_t), pointer :: fallback_model
    class(generic_event_t), pointer :: event
    type(event_sample_data_t) :: data
    class(eio_t), allocatable :: eio
    type(string_t) :: sample
    integer :: u_file, iostat, i_prc

    write (u, "(A)")  "* Test output: eio_lhef_6"
    write (u, "(A)")  "*   Purpose: read a LHEF 3.0 file"
    write (u, "(A)")

    write (u, "(A)")  "* Write a LHEF data file"
    write (u, "(A)")
 
    u_file = free_unit ()
    sample = "eio_lhef_6"
    open (u_file, file = char (sample // ".lhe"), &
         status = "replace", action = "readwrite")
    
    write (u_file, "(A)")  '<LesHouchesEvents version="3.0">'
    write (u_file, "(A)")  '<header>'
    write (u_file, "(A)")  '</header>'
    write (u_file, "(A)")  '<init>'
    write (u_file, "(A)")  ' 25 25  5.0000000000E+02  5.0000000000E+02 &
         &-1 -1 -1 -1 4 1'
    write (u_file, "(A)")  '  1.0000000000E-01  1.0000000000E-03 &
         & 0.0000000000E+00 42'
    write (u_file, "(A)")  '<generator version="2.2.3">WHIZARD&
         &</generator>'
    write (u_file, "(A)")  '<xsecinfo neve="1" totxsec="1.0000000000E-01" />'
    write (u_file, "(A)")  '<weightinfo name="sqme_prc" />'
    write (u_file, "(A)")  '</init>'
    write (u_file, "(A)")  '<event>'
    write (u_file, "(A)")  ' 4 42  3.0574068604E+08  1.0000000000E+03 &
         &-1.0000000000E+00 -1.0000000000E+00'
    write (u_file, "(A)")  ' 25 -1 0 0 0 0  0.0000000000E+00 &
         & 0.0000000000E+00  4.8412291828E+02  5.0000000000E+02 &
         & 1.2500000000E+02  0.0000000000E+00  9.0000000000E+00'
    write (u_file, "(A)")  ' 25 -1 0 0 0 0  0.0000000000E+00 &
         & 0.0000000000E+00 -4.8412291828E+02  5.0000000000E+02 &
         & 1.2500000000E+02  0.0000000000E+00  9.0000000000E+00'
    write (u_file, "(A)")  ' 25 1 1 2 0 0 -1.4960220911E+02 &
         &-4.6042825611E+02  0.0000000000E+00  5.0000000000E+02 &
         & 1.2500000000E+02  0.0000000000E+00  9.0000000000E+00'
    write (u_file, "(A)")  ' 25 1 1 2 0 0  1.4960220911E+02 &
         & 4.6042825611E+02  0.0000000000E+00  5.0000000000E+02 &
         & 1.2500000000E+02  0.0000000000E+00  9.0000000000E+00'
    write (u_file, "(A)")  '<weights>1.0000000000E+00</weights>'
    write (u_file, "(A)")  '</event>'
    write (u_file, "(A)")  '</LesHouchesEvents>'
    close (u_file)

    write (u, "(A)")  "* Initialize test process"
    write (u, "(A)")

    call eio_prepare_fallback_model (fallback_model)
    call eio_prepare_test (event, unweighted = .false.)

    allocate (eio_lhef_t :: eio)
    select type (eio)
    type is (eio_lhef_t)
       call eio%set_parameters (version = "3.0", recover_beams = .false.)
    end select
    call eio%set_fallback_model (fallback_model)
    
    call data%init (1)
    data%unweighted = .false.
    data%norm_mode = NORM_SIGMA
    data%n_beam = 2
    data%pdg_beam = 25
    data%energy_beam = 500
    data%proc_num_id = [42]
    call data%write (u)
    write (u, *)

    write (u, "(A)")  "* Initialize and read header"
    write (u, "(A)")

    call eio%init_in (sample, data)
    call eio%write (u)
    
    write (u, *)
    
    select type (eio)
    type is (eio_lhef_t)
       call eio%tag_lhef%write (u);  write (u, *)
    end select

    write (u, *)
    call data%write (u)

    write (u, "(A)")
    write (u, "(A)")  "* Read event"
    write (u, "(A)")

    call eio%input_i_prc (i_prc, iostat)
    
    select type (eio)
    type is (eio_lhef_t)
       write (u, "(A,I0,A,I0)")  "Found process #", i_prc, &
            " with ID = ", eio%proc_num_id(i_prc)
    end select
   
    call eio%input_event (event, iostat)

    call event%write (u)

    write (u, "(A)")
    write (u, "(A)")  "* Read closing"
    write (u, "(A)")

    call eio%input_i_prc (i_prc, iostat)
    write (u, "(A,I0)")  "iostat = ", iostat
    
    write (u, "(A)")
    write (u, "(A)")  "* Cleanup"
 
    call eio%final ()

    call eio_cleanup_test (event)
    call eio_cleanup_fallback_model (fallback_model)

    write (u, "(A)")
    write (u, "(A)")  "* Test output end: eio_lhef_6"
    
  end subroutine eio_lhef_6
  
\eatline
\nwendcode{}\nwbegindocs{1159}\nwdocspar
%%%%%%%%%%%%%%%%%%%%%%%%%%%%%%%%%%%%%%%%%%%%%%%%%%%%%%%%%%%%%%%%%%%%%%%%
\section{STDHEP File Formats}
Here, we implement the two existing STDHEP file formats, one based on the
HEPRUP/HEPEUP common blocks, the other based on the HEPEVT common block.
The second one is actually the standard STDHEP format.
\nwenddocs{}\nwbegincode{1160}\moddef{\code{}eio{\char95}stdhep.f90\edoc{}}\endmoddef
\LA{}File header\RA{}

module eio_stdhep
  
  use kinds
  use io_units
\LA{}Use strings\RA{}
  use string_utils
  use unit_tests
  use diagnostics

  use lorentz
  use model_data
  use particles
  use event_base
  use hep_common
  use hep_events
  use xdr_stdhep
  use eio_data
  use eio_base

\LA{}Standard module head\RA{}

\LA{}EIO stdhep: public\RA{}

\LA{}EIO stdhep: types\RA{}

\LA{}EIO stdhep: variables\RA{}

contains
  
\LA{}EIO stdhep: procedures\RA{}

\LA{}EIO stdhep: tests\RA{}

end module eio_stdhep
\eatline
\nwendcode{}\nwbegindocs{1161}\nwdocspar
\subsection{Type}
\nwenddocs{}\nwbegincode{1162}\moddef{EIO stdhep: public}\endmoddef
  public :: eio_stdhep_t
\nwendcode{}\nwbegincode{1163}\moddef{EIO stdhep: types}\endmoddef
  type, abstract, extends (eio_t) :: eio_stdhep_t
     logical :: writing = .false.
     logical :: reading = .false.
     integer :: unit = 0
     logical :: keep_beams = .false.     
     logical :: keep_remnants = .true.
     logical :: recover_beams = .false.
     logical :: use_alpha_s_from_file = .false.
     logical :: use_scale_from_file = .false.
     integer, dimension(:), allocatable :: proc_num_id     
     integer(i64) :: n_events_expected = 0
   contains
   \LA{}EIO stdhep: eio stdhep: TBP\RA{}
  end type eio_stdhep_t
  
\eatline
\nwendcode{}\nwbegindocs{1164}\nwdocspar
\nwenddocs{}\nwbegincode{1165}\moddef{EIO stdhep: public}\plusendmoddef
  public :: eio_stdhep_hepevt_t
\nwendcode{}\nwbegincode{1166}\moddef{EIO stdhep: types}\plusendmoddef
  type, extends (eio_stdhep_t) :: eio_stdhep_hepevt_t
  end type eio_stdhep_hepevt_t
  
\eatline
\nwendcode{}\nwbegindocs{1167}\nwdocspar
\nwenddocs{}\nwbegincode{1168}\moddef{EIO stdhep: public}\plusendmoddef
  public :: eio_stdhep_hepeup_t
\nwendcode{}\nwbegincode{1169}\moddef{EIO stdhep: types}\plusendmoddef
  type, extends (eio_stdhep_t) :: eio_stdhep_hepeup_t
  end type eio_stdhep_hepeup_t
  
\eatline
\nwendcode{}\nwbegindocs{1170}\nwdocspar
\subsection{Specific Methods}
Set parameters that are specifically used with STDHEP file formats.
\nwenddocs{}\nwbegincode{1171}\moddef{EIO stdhep: eio stdhep: TBP}\endmoddef
  procedure :: set_parameters => eio_stdhep_set_parameters
\nwendcode{}\nwbegincode{1172}\moddef{EIO stdhep: procedures}\endmoddef
  subroutine eio_stdhep_set_parameters (eio, &
       keep_beams, keep_remnants, recover_beams, use_alpha_s_from_file, &
       use_scale_from_file, extension)
    class(eio_stdhep_t), intent(inout) :: eio
    logical, intent(in), optional :: keep_beams
    logical, intent(in), optional :: keep_remnants
    logical, intent(in), optional :: recover_beams
    logical, intent(in), optional :: use_alpha_s_from_file
    logical, intent(in), optional :: use_scale_from_file
    type(string_t), intent(in), optional :: extension
    if (present (keep_beams))  eio%keep_beams = keep_beams
    if (present (keep_remnants))  eio%keep_remnants = keep_remnants
    if (present (recover_beams))  eio%recover_beams = recover_beams
    if (present (use_alpha_s_from_file)) &
         eio%use_alpha_s_from_file = use_alpha_s_from_file
    if (present (use_scale_from_file))  &
         eio%use_scale_from_file = use_scale_from_file
    if (present (extension)) then
       eio%extension = extension
    else
       select type (eio)
       type is (eio_stdhep_hepevt_t)
          eio%extension = "hep"
       type is (eio_stdhep_hepeup_t)
          eio%extension = "up.hep"
       end select
    end if
  end subroutine eio_stdhep_set_parameters
  
\eatline
\nwendcode{}\nwbegindocs{1173}\nwdocspar
\subsection{Common Methods}
Output.  This is not the actual event format, but a readable account
of the current object status.
\nwenddocs{}\nwbegincode{1174}\moddef{EIO stdhep: eio stdhep: TBP}\plusendmoddef
  procedure :: write => eio_stdhep_write
\nwendcode{}\nwbegincode{1175}\moddef{EIO stdhep: procedures}\plusendmoddef
  subroutine eio_stdhep_write (object, unit)
    class(eio_stdhep_t), intent(in) :: object
    integer, intent(in), optional :: unit
    integer :: u, i
    u = given_output_unit (unit)
    write (u, "(1x,A)")  "STDHEP event stream:"
    if (object%writing) then
       write (u, "(3x,A,A)")  "Writing to file   = ", char (object%filename)
    else if (object%reading) then
       write (u, "(3x,A,A)")  "Reading from file = ", char (object%filename)
    else
       write (u, "(3x,A)")  "[closed]"
    end if
    write (u, "(3x,A,L1)")    "Keep beams        = ", object%keep_beams
    write (u, "(3x,A,L1)")    "Keep remnants     = ", object%keep_remnants
    write (u, "(3x,A,L1)")    "Recover beams     = ", object%recover_beams
    write (u, "(3x,A,L1)")    "Alpha_s from file = ", &
         object%use_alpha_s_from_file
    write (u, "(3x,A,L1)")    "Scale from file   = ", &
         object%use_scale_from_file    
    if (allocated (object%proc_num_id)) then
       write (u, "(3x,A)")  "Numerical process IDs:"
       do i = 1, size (object%proc_num_id)
          write (u, "(5x,I0,': ',I0)")  i, object%proc_num_id(i)
       end do
    end if      
  end subroutine eio_stdhep_write
  
\eatline
\nwendcode{}\nwbegindocs{1176}Finalizer: close any open file.
\nwenddocs{}\nwbegincode{1177}\moddef{EIO stdhep: eio stdhep: TBP}\plusendmoddef
  procedure :: final => eio_stdhep_final
\nwendcode{}\nwbegincode{1178}\moddef{EIO stdhep: procedures}\plusendmoddef
  subroutine eio_stdhep_final (object)
    class(eio_stdhep_t), intent(inout) :: object
    if (allocated (object%proc_num_id))  deallocate (object%proc_num_id)
    if (object%writing) then
       write (msg_buffer, "(A,A,A)")  "Events: closing STDHEP file '", &
            char (object%filename), "'"
       call msg_message ()
       call stdhep_write (200)
       call stdhep_end ()
       object%writing = .false.
    else if (object%reading) then
       write (msg_buffer, "(A,A,A)")  "Events: closing STDHEP file '", &
            char (object%filename), "'"
       call msg_message ()
       object%reading = .false.
    end if
  end subroutine eio_stdhep_final
  
\eatline
\nwendcode{}\nwbegindocs{1179}Common initialization for input and output.
\nwenddocs{}\nwbegincode{1180}\moddef{EIO stdhep: eio stdhep: TBP}\plusendmoddef
  procedure :: common_init => eio_stdhep_common_init
\nwendcode{}\nwbegincode{1181}\moddef{EIO stdhep: procedures}\plusendmoddef
  subroutine eio_stdhep_common_init (eio, sample, data, extension)
    class(eio_stdhep_t), intent(inout) :: eio
    type(string_t), intent(in) :: sample
    type(string_t), intent(in), optional :: extension
    type(event_sample_data_t), intent(in), optional :: data
    if (.not. present (data)) &
         call msg_bug ("STDHEP initialization: missing data")
    if (data%n_beam /= 2) &
         call msg_fatal ("STDHEP: defined for scattering processes only")
    if (present (extension)) then
       eio%extension = extension
    end if
    eio%sample = sample    
    call eio%set_filename ()    
    eio%unit = free_unit () 
    allocate (eio%proc_num_id (data%n_proc), source = data%proc_num_id)
  end subroutine eio_stdhep_common_init

\eatline
\nwendcode{}\nwbegindocs{1182}Split event file: increment the counter, close the current file, open a new
one.  If the file needs a header, repeat it for the new file.  (We assume that
the common block contents are still intact.)
\nwenddocs{}\nwbegincode{1183}\moddef{EIO stdhep: eio stdhep: TBP}\plusendmoddef
  procedure :: split_out => eio_stdhep_split_out
\nwendcode{}\nwbegincode{1184}\moddef{EIO stdhep: procedures}\plusendmoddef
  subroutine eio_stdhep_split_out (eio)
    class(eio_stdhep_t), intent(inout) :: eio
    if (eio%split) then
       eio%split_index = eio%split_index + 1
       call eio%set_filename ()
       write (msg_buffer, "(A,A,A)")  "Events: writing to STDHEP file '", &
            char (eio%filename), "'"
       call msg_message ()
       call stdhep_write (200)
       call stdhep_end ()
       select type (eio)
       type is (eio_stdhep_hepeup_t)
          call stdhep_init_out (char (eio%filename), &
               "WHIZARD \LA{}Version\RA{}", eio%n_events_expected)
          call stdhep_write (100)
          call stdhep_write (STDHEP_HEPRUP)
       type is (eio_stdhep_hepevt_t)
          call stdhep_init_out (char (eio%filename), &
               "WHIZARD \LA{}Version\RA{}", eio%n_events_expected) 
          call stdhep_write (100)
       end select
    end if
  end subroutine eio_stdhep_split_out
  
\eatline
\nwendcode{}\nwbegindocs{1185}Initialize event writing.
\nwenddocs{}\nwbegincode{1186}\moddef{EIO stdhep: eio stdhep: TBP}\plusendmoddef
  procedure :: init_out => eio_stdhep_init_out
\nwendcode{}\nwbegincode{1187}\moddef{EIO stdhep: procedures}\plusendmoddef
  subroutine eio_stdhep_init_out (eio, sample, data, success, extension)
    class(eio_stdhep_t), intent(inout) :: eio
    type(string_t), intent(in) :: sample
    type(string_t), intent(in), optional :: extension
    type(event_sample_data_t), intent(in), optional :: data
    logical, intent(out), optional :: success
    integer :: i
    if (.not. present (data)) &
         call msg_bug ("STDHEP initialization: missing data")        
    call eio%set_splitting (data)    
    call eio%common_init (sample, data, extension)
    eio%n_events_expected = data%n_evt
    write (msg_buffer, "(A,A,A)")  "Events: writing to STDHEP file '", &
         char (eio%filename), "'"
    call msg_message ()
    eio%writing = .true.
    select type (eio)
    type is (eio_stdhep_hepeup_t)
       call heprup_init &
            (data%pdg_beam, &
            data%energy_beam, &
            n_processes = data%n_proc, &
            unweighted = data%unweighted, &
            negative_weights = data%negative_weights)           
       do i = 1, data%n_proc
          call heprup_set_process_parameters (i = i, &
               process_id = data%proc_num_id(i), &
               cross_section = data%cross_section(i), &
               error = data%error(i))          
       end do
       call stdhep_init_out (char (eio%filename), &
            "WHIZARD \LA{}Version\RA{}", eio%n_events_expected)
       call stdhep_write (100)
       call stdhep_write (STDHEP_HEPRUP)
    type is (eio_stdhep_hepevt_t)
       call stdhep_init_out (char (eio%filename), &
            "WHIZARD \LA{}Version\RA{}", eio%n_events_expected) 
       call stdhep_write (100)
    end select    
    if (present (success))  success = .true.
  end subroutine eio_stdhep_init_out
    
\eatline
\nwendcode{}\nwbegindocs{1188}Initialize event reading.
\nwenddocs{}\nwbegincode{1189}\moddef{EIO stdhep: eio stdhep: TBP}\plusendmoddef
  procedure :: init_in => eio_stdhep_init_in
\nwendcode{}\nwbegincode{1190}\moddef{EIO stdhep: procedures}\plusendmoddef
  subroutine eio_stdhep_init_in (eio, sample, data, success, extension)
    class(eio_stdhep_t), intent(inout) :: eio
    type(string_t), intent(in) :: sample
    type(string_t), intent(in), optional :: extension
    type(event_sample_data_t), intent(inout), optional :: data
    logical, intent(out), optional :: success
    integer :: ilbl, lok
    logical :: exist
    call eio%common_init (sample, data, extension)
    write (msg_buffer, "(A,A,A)")  "Events: reading from STDHEP file '", &
         char (eio%filename), "'"
    call msg_message ()
    inquire (file = char (eio%filename), exist = exist)
    if (.not. exist)  call msg_fatal ("Events: STDHEP file not found.")
    eio%reading = .true.
    call stdhep_init_in (char (eio%filename), eio%n_events_expected)
    call stdhep_read (ilbl, lok)
    if (lok /= 0) then
       call stdhep_end () 
       write (msg_buffer, "(A)")  "Events: STDHEP file appears to" // &
            " be empty."
       call msg_message ()      
    end if
    if (ilbl == 100) then
       write (msg_buffer, "(A)")  "Events: reading in STDHEP events"
       call msg_message ()
    end if 
    if (present (success))  success = .false.
  end subroutine eio_stdhep_init_in
    
\eatline
\nwendcode{}\nwbegindocs{1191}Switch from input to output: reopen the file for reading.
\nwenddocs{}\nwbegincode{1192}\moddef{EIO stdhep: eio stdhep: TBP}\plusendmoddef
  procedure :: switch_inout => eio_stdhep_switch_inout
\nwendcode{}\nwbegincode{1193}\moddef{EIO stdhep: procedures}\plusendmoddef
  subroutine eio_stdhep_switch_inout (eio, success)
    class(eio_stdhep_t), intent(inout) :: eio
    logical, intent(out), optional :: success
    call msg_bug ("STDHEP: in-out switch not supported")
    if (present (success))  success = .false.
  end subroutine eio_stdhep_switch_inout
  
\eatline
\nwendcode{}\nwbegindocs{1194}Output an event.  Write first the event indices, then weight and
squared matrix element, then the particle set.
\nwenddocs{}\nwbegincode{1195}\moddef{EIO stdhep: eio stdhep: TBP}\plusendmoddef
  procedure :: output => eio_stdhep_output
\nwendcode{}\nwbegincode{1196}\moddef{EIO stdhep: procedures}\plusendmoddef
  subroutine eio_stdhep_output (eio, event, i_prc, reading, passed, pacify)
    class(eio_stdhep_t), intent(inout) :: eio
    class(generic_event_t), intent(in), target :: event
    integer, intent(in) :: i_prc
    logical, intent(in), optional :: reading, passed, pacify
    if (present (passed)) then
       if (.not. passed)  return
    end if
    if (eio%writing) then
       select type (eio)
       type is (eio_stdhep_hepeup_t)
          call hepeup_from_event (event, &
               process_index = eio%proc_num_id (i_prc), &
               keep_beams = eio%keep_beams, &
               keep_remnants = eio%keep_remnants)
          call stdhep_write (STDHEP_HEPEUP)
       type is (eio_stdhep_hepevt_t)
          call hepevt_from_event (event, &
               i_evt = event%get_index (), &                         
               keep_beams = eio%keep_beams, &
               keep_remnants = eio%keep_remnants)
          call stdhep_write (STDHEP_HEPEVT)                              
       end select       
    else
       call eio%write ()
       call msg_fatal ("STDHEP file is not open for writing")
    end if
  end subroutine eio_stdhep_output

\eatline
\nwendcode{}\nwbegindocs{1197}Input an event. We do not allow to read in STDHEP files written via
the HEPEVT common block as there is no control on the process ID.
\nwenddocs{}\nwbegincode{1198}\moddef{EIO stdhep: eio stdhep: TBP}\plusendmoddef
  procedure :: input_i_prc => eio_stdhep_input_i_prc
  procedure :: input_event => eio_stdhep_input_event
\nwendcode{}\nwbegincode{1199}\moddef{EIO stdhep: procedures}\plusendmoddef
  subroutine eio_stdhep_input_i_prc (eio, i_prc, iostat)
    class(eio_stdhep_t), intent(inout) :: eio
    integer, intent(out) :: i_prc
    integer, intent(out) :: iostat
    integer :: i, ilbl, proc_num_id
    iostat = 0
    select type (eio)
    type is (eio_stdhep_hepevt_t)
       if (size (eio%proc_num_id) > 1) then
          call msg_fatal ("Events: only single processes allowed " // &
               "with the STDHEP HEPEVT format.")
       else
          proc_num_id = eio%proc_num_id (1)
          call stdhep_read (ilbl, lok)
       end if
    type is (eio_stdhep_hepeup_t)
       call stdhep_read (ilbl, lok)
       if (lok /= 0)  call msg_error ("Events: STDHEP appears to be " // &
            "empty or corrupted.")
       if (ilbl == 12) then
          call stdhep_read (ilbl, lok)
       end if
       if (ilbl == 11) then
          proc_num_id = IDPRUP
       end if       
    end select
    FIND_I_PRC: do i = 1, size (eio%proc_num_id)
       if (eio%proc_num_id(i) == proc_num_id) then
          i_prc = i
          exit FIND_I_PRC
       end if
    end do FIND_I_PRC
    if (i_prc == 0)  call err_index
  contains
    subroutine err_index
      call msg_error ("STDHEP: reading events: undefined process ID " &
           // char (str (proc_num_id)) // ", aborting read")
      iostat = 1
    end subroutine err_index    
  end subroutine eio_stdhep_input_i_prc

  subroutine eio_stdhep_input_event (eio, event, iostat)
    class(eio_stdhep_t), intent(inout) :: eio
    class(generic_event_t), intent(inout), target :: event
    integer, intent(out) :: iostat
    call event%reset ()
    call event%select (1, 1, 1)
    call hepeup_to_event (event, eio%fallback_model, &
         recover_beams = eio%recover_beams, &
         use_alpha_s = eio%use_alpha_s_from_file, &
         use_scale = eio%use_scale_from_file)    
  end subroutine eio_stdhep_input_event

\eatline
\eatline
\nwendcode{}\nwbegindocs{1200}STDHEP speficic routines.
\nwenddocs{}\nwbegincode{1201}\moddef{EIO stdhep: public}\plusendmoddef
  public :: stdhep_init_out
  public :: stdhep_init_in
  public :: stdhep_write 
  public :: stdhep_end  
\nwendcode{}\nwbegincode{1202}\moddef{EIO stdhep: procedures}\plusendmoddef
  subroutine stdhep_init_out (file, title, nevt)
    character(len=*), intent(in) :: file, title
    integer(i64), intent(in) :: nevt
    integer(i32) :: nevt32
    external stdxwinit, stdxwrt
    nevt32 = min (nevt, int (huge (1_i32), i64))
    call stdxwinit (file, title, nevt32, istr, lok)    
  end subroutine stdhep_init_out

  subroutine stdhep_init_in (file, nevt)
    character(len=*), intent(in) :: file
    integer(i64), intent(out) :: nevt
    integer(i32) :: nevt32
    external stdxrinit, stdxrd
    call stdxrinit (file, nevt32, istr, lok)
    if (lok /= 0)  call msg_fatal ("STDHEP: error in reading file '" // &
         file // "'.")
    nevt = int (nevt32, i64)
  end subroutine stdhep_init_in
  
  subroutine stdhep_write (ilbl)
    integer, intent(in) :: ilbl
    external stdxwrt
    call stdxwrt (ilbl, istr, lok)
  end subroutine stdhep_write

  subroutine stdhep_read (ilbl, lok)
    integer, intent(out) :: ilbl, lok
    external stdxrd
    call stdxrd (ilbl, istr, lok)
    if (lok /= 0)  return
  end subroutine stdhep_read
  
  subroutine stdhep_end
    external stdxend
    call stdxend (istr)
  end subroutine stdhep_end  
  
\nwendcode{}\nwbegindocs{1203} %def stdhep_init stdhep_read stdhep_write stdhep_end
\nwenddocs{}\nwbegindocs{1204}\nwdocspar
\subsection{Variables}
\nwenddocs{}\nwbegincode{1205}\moddef{EIO stdhep: variables}\endmoddef
  integer, save :: istr, lok
  integer, parameter :: &
       STDHEP_HEPEVT = 1, STDHEP_HEPEUP = 11, STDHEP_HEPRUP = 12
\nwendcode{}\nwbegindocs{1206}\nwdocspar
\nwenddocs{}\nwbegindocs{1207}\nwdocspar
\subsection{Unit tests}
\nwenddocs{}\nwbegincode{1208}\moddef{EIO stdhep: public}\plusendmoddef
  public :: eio_stdhep_test
\nwendcode{}\nwbegincode{1209}\moddef{EIO stdhep: tests}\endmoddef
  subroutine eio_stdhep_test (u, results)
    integer, intent(in) :: u
    type(test_results_t), intent(inout) :: results
  \LA{}EIO stdhep: execute tests\RA{}
  end subroutine eio_stdhep_test
  
\eatline
\nwendcode{}\nwbegindocs{1210}\nwdocspar
\subsubsection{Test I/O methods}
We test the implementation of the STDHEP HEPEVT I/O method:
\nwenddocs{}\nwbegincode{1211}\moddef{EIO stdhep: execute tests}\endmoddef
  call test (eio_stdhep_1, "eio_stdhep_1", &
       "read and write event contents, format [stdhep]", &
       u, results)
\nwendcode{}\nwbegincode{1212}\moddef{EIO stdhep: tests}\plusendmoddef
  subroutine eio_stdhep_1 (u)
    integer, intent(in) :: u
    class(generic_event_t), pointer :: event
    type(event_sample_data_t) :: data
    class(eio_t), allocatable :: eio
    type(string_t) :: sample
    integer :: u_file, iostat
    character(215) :: buffer

    write (u, "(A)")  "* Test output: eio_stdhep_1"
    write (u, "(A)")  "*   Purpose: generate an event in STDHEP HEPEVT format"
    write (u, "(A)")  "*      and write weight to file"
    write (u, "(A)")

    write (u, "(A)")  "* Initialize test process"
 
    call eio_prepare_test (event)
    
    call data%init (1)
    data%n_evt = 1
    data%n_beam = 2
    data%pdg_beam = 25
    data%energy_beam = 500
    data%proc_num_id = [42]
    data%cross_section(1) = 100
    data%error(1) = 1
    data%total_cross_section = sum (data%cross_section)

    write (u, "(A)")
    write (u, "(A)")  "* Generate and write an event"
    write (u, "(A)")
 
    sample = "eio_stdhep_1"
 
    allocate (eio_stdhep_hepevt_t :: eio)
    select type (eio)
    type is (eio_stdhep_hepevt_t)
       call eio%set_parameters ()
    end select
    
    call eio%init_out (sample, data)
    call event%generate (1, [0._default, 0._default])
    call event%evaluate_expressions ()
    call event%pacify_particle_set ()
    
    call eio%output (event, i_prc = 1)
    call eio%write (u)
    call eio%final ()
    
    write (u, "(A)") 
    write (u, "(A)")  "* Write STDHEP file contents to ASCII file"
    write (u, "(A)")

    call write_stdhep_event &
         (sample // ".hep", var_str ("test_1.hep"), 1)
        
    write (u, "(A)") 
    write (u, "(A)")  "* Read in ASCII contents of STDHEP file"
    write (u, "(A)")    
    
    u_file = free_unit ()
    open (u_file, file = "test_1.hep", &
         action = "read", status = "old")
    do
       read (u_file, "(A)", iostat = iostat)  buffer
       if (iostat /= 0)  exit
       if (trim (buffer) == "")  cycle
       if (buffer(1:18) == "    total blocks: ")  &
            buffer = "    total blocks: [...]"
       if (buffer(1:25) == "           title: WHIZARD")  &
            buffer = "           title: WHIZARD [version]"
       if (buffer(1:17) == "            date:")  &
            buffer = "            date: [...]"
       if (buffer(1:17) == "    closing date:")  &
            buffer = "    closing date: [...]"       
       write (u, "(A)") trim (buffer)
    end do
    close (u_file)    
        
    write (u, "(A)")
    write (u, "(A)")  "* Reset data"
    write (u, "(A)")
 
    deallocate (eio)
    allocate (eio_stdhep_hepevt_t :: eio)
    
    select type (eio)
    type is (eio_stdhep_hepevt_t)
       call eio%set_parameters (keep_beams = .true.)
    end select
    call eio%write (u)

    write (u, "(A)")
    write (u, "(A)")  "* Cleanup"
 
    call eio_cleanup_test (event)

    write (u, "(A)")
    write (u, "(A)")  "* Test output end: eio_stdhep_1"
    
  end subroutine eio_stdhep_1
  
\eatline
\nwendcode{}\nwbegindocs{1213}\nwdocspar
We test the implementation of the STDHEP HEPEUP I/O method:
\nwenddocs{}\nwbegincode{1214}\moddef{EIO stdhep: execute tests}\plusendmoddef
  call test (eio_stdhep_2, "eio_stdhep_2", &
       "read and write event contents, format [stdhep]", &
       u, results)
\nwendcode{}\nwbegincode{1215}\moddef{EIO stdhep: tests}\plusendmoddef
  subroutine eio_stdhep_2 (u)
    integer, intent(in) :: u
    class(generic_event_t), pointer :: event
    type(event_sample_data_t) :: data
    class(model_data_t), pointer :: fallback_model
    class(eio_t), allocatable :: eio
    type(string_t) :: sample
    integer :: u_file, iostat
    character(215) :: buffer

    write (u, "(A)")  "* Test output: eio_stdhep_2"
    write (u, "(A)")  "*   Purpose: generate an event in STDHEP HEPEUP format"
    write (u, "(A)")  "*      and write weight to file"
    write (u, "(A)")

    write (u, "(A)")  "* Initialize test process"
 
    call eio_prepare_fallback_model (fallback_model)
    call eio_prepare_test (event, unweighted = .false.)
    
    call data%init (1)
    data%n_evt = 1
    data%n_beam = 2
    data%pdg_beam = 25
    data%energy_beam = 500
    data%proc_num_id = [42]
    data%cross_section(1) = 100
    data%error(1) = 1
    data%total_cross_section = sum (data%cross_section)

    write (u, "(A)")
    write (u, "(A)")  "* Generate and write an event"
    write (u, "(A)")
 
    sample = "eio_stdhep_2"
 
    allocate (eio_stdhep_hepeup_t :: eio)
    select type (eio)
    type is (eio_stdhep_hepeup_t)
       call eio%set_parameters ()
    end select  
    call eio%set_fallback_model (fallback_model)
    
    call eio%init_out (sample, data)
    call event%generate (1, [0._default, 0._default])
    call event%evaluate_expressions ()

    call eio%output (event, i_prc = 1)
    call eio%write (u)
    call eio%final ()

    write (u, "(A)") 
    write (u, "(A)")  "* Write STDHEP file contents to ASCII file"
    write (u, "(A)")

    call write_stdhep_event & 
         (sample // ".up.hep", var_str ("test_2.hep"), 2)
        
    write (u, "(A)") 
    write (u, "(A)")  "* Read in ASCII contents of STDHEP file"
    write (u, "(A)")    
    
    u_file = free_unit ()
    open (u_file, file = "test_2.hep", &
         action = "read", status = "old")
    do
       read (u_file, "(A)", iostat = iostat)  buffer
       if (iostat /= 0)  exit
       if (trim (buffer) == "")  cycle
       if (buffer(1:18) == "    total blocks: ")  &
            buffer = "    total blocks: [...]"       
       if (buffer(1:25) == "           title: WHIZARD")  &
            buffer = "           title: WHIZARD [version]"
       if (buffer(1:17) == "            date:")  &
            buffer = "            date: [...]"
       if (buffer(1:17) == "    closing date:")  &
            buffer = "    closing date: [...]"       
       write (u, "(A)") trim (buffer)
    end do
    close (u_file)    

    write (u, "(A)")
    write (u, "(A)")  "* Reset data"
    write (u, "(A)")
 
    deallocate (eio)
    allocate (eio_stdhep_hepeup_t :: eio)
    
    select type (eio)
    type is (eio_stdhep_hepeup_t)
       call eio%set_parameters (keep_beams = .true.)
    end select
    call eio%write (u)

    write (u, "(A)")
    write (u, "(A)")  "* Cleanup"
 
    call eio_cleanup_test (event)
    call eio_cleanup_fallback_model (fallback_model)

    write (u, "(A)")
    write (u, "(A)")  "* Test output end: eio_stdhep_2"
    
  end subroutine eio_stdhep_2
  
\eatline
\nwendcode{}\nwbegindocs{1216}\nwdocspar
Check input from a StdHep file, HEPEVT block.
\nwenddocs{}\nwbegincode{1217}\moddef{EIO stdhep: execute tests}\plusendmoddef
  call test (eio_stdhep_3, "eio_stdhep_3", &
       "read StdHep file, HEPEVT block", &
       u, results)
\nwendcode{}\nwbegincode{1218}\moddef{EIO stdhep: tests}\plusendmoddef
  subroutine eio_stdhep_3 (u)
    integer, intent(in) :: u
    class(model_data_t), pointer :: fallback_model
    class(generic_event_t), pointer :: event
    type(event_sample_data_t) :: data
    class(eio_t), allocatable :: eio
    type(string_t) :: sample
    integer :: u_file, iostat, i_prc

    write (u, "(A)")  "* Test output: eio_stdhep_3"
    write (u, "(A)")  "*   Purpose: read a StdHep file, HEPEVT block"
    write (u, "(A)")

    write (u, "(A)")  "* Write a StdHep data file, HEPEVT block"
    write (u, "(A)")
 
    call eio_prepare_fallback_model (fallback_model)    
    call eio_prepare_test (event)
    
    call data%init (1)
    data%n_evt = 1
    data%n_beam = 2
    data%pdg_beam = 25
    data%energy_beam = 500
    data%proc_num_id = [42]
    data%cross_section(1) = 100
    data%error(1) = 1
    data%total_cross_section = sum (data%cross_section)

    write (u, "(A)")
    write (u, "(A)")  "* Generate and write an event"
    write (u, "(A)")
 
    sample = "eio_stdhep_3"
 
    allocate (eio_stdhep_hepevt_t :: eio)
    select type (eio)
    type is (eio_stdhep_hepevt_t)
       call eio%set_parameters ()
    end select  
    call eio%set_fallback_model (fallback_model)
    
    call eio%init_out (sample, data)
    call event%generate (1, [0._default, 0._default])
    call event%evaluate_expressions ()

    call eio%output (event, i_prc = 1)
    call eio%write (u)
    call eio%final ()

    call eio_cleanup_test (event)
    call eio_cleanup_fallback_model (fallback_model)    
    deallocate (eio)
    
    write (u, "(A)")  "* Initialize test process"
    write (u, "(A)")

    call eio_prepare_fallback_model (fallback_model)
    call eio_prepare_test (event, unweighted = .false.)

    allocate (eio_stdhep_hepevt_t :: eio)
    select type (eio)
    type is (eio_stdhep_hepevt_t)
       call eio%set_parameters (recover_beams = .false.)
    end select
    call eio%set_fallback_model (fallback_model)
    
    call data%init (1)
    data%n_beam = 2
    data%unweighted = .true.
    data%norm_mode = NORM_UNIT
    data%pdg_beam = 25
    data%energy_beam = 500
    data%proc_num_id = [42]
    call data%write (u)
    write (u, *)

    write (u, "(A)")  "* Initialize"
    write (u, "(A)")

    call eio%init_in (sample, data)
    call eio%write (u)
    
    write (u, "(A)")
    write (u, "(A)")  "* Read event"
    write (u, "(A)")

    call eio%input_i_prc (i_prc, iostat)
    
    select type (eio)
    type is (eio_stdhep_hepevt_t)
       write (u, "(A,I0,A,I0)")  "Found process #", i_prc, &
            " with ID = ", eio%proc_num_id(i_prc)
    end select

    call eio%input_event (event, iostat)

    call event%write (u)

    write (u, "(A)")
    write (u, "(A)")  "* Read closing"
    write (u, "(A)")

    call eio%input_i_prc (i_prc, iostat)
    write (u, "(A,I0)")  "iostat = ", iostat
    
    write (u, "(A)")
    write (u, "(A)")  "* Cleanup"
 
    call eio%final ()

    call eio_cleanup_test (event)
    call eio_cleanup_fallback_model (fallback_model)

    write (u, "(A)")
    write (u, "(A)")  "* Test output end: eio_stdhep_3"
    
  end subroutine eio_stdhep_3
  
\eatline
\nwendcode{}\nwbegindocs{1219}\nwdocspar
Check input from a StdHep file, HEPEVT block.
\nwenddocs{}\nwbegincode{1220}\moddef{EIO stdhep: execute tests}\plusendmoddef
  call test (eio_stdhep_4, "eio_stdhep_4", &
       "read StdHep file, HEPRUP/HEPEUP block", &
       u, results)
\nwendcode{}\nwbegincode{1221}\moddef{EIO stdhep: tests}\plusendmoddef
  subroutine eio_stdhep_4 (u)
    integer, intent(in) :: u
    class(model_data_t), pointer :: fallback_model
    class(generic_event_t), pointer :: event
    type(event_sample_data_t) :: data
    class(eio_t), allocatable :: eio
    type(string_t) :: sample
    integer :: u_file, iostat, i_prc

    write (u, "(A)")  "* Test output: eio_stdhep_3"
    write (u, "(A)")  "*   Purpose: read a StdHep file, HEPRUP/HEPEUP block"
    write (u, "(A)")

    write (u, "(A)")  "* Write a StdHep data file, HEPRUP/HEPEUP block"
    write (u, "(A)")
 
    call eio_prepare_fallback_model (fallback_model)    
    call eio_prepare_test (event)
    
    call data%init (1)
    data%n_evt = 1
    data%n_beam = 2
    data%pdg_beam = 25
    data%energy_beam = 500
    data%proc_num_id = [42]
    data%cross_section(1) = 100
    data%error(1) = 1
    data%total_cross_section = sum (data%cross_section)

    write (u, "(A)")
    write (u, "(A)")  "* Generate and write an event, HEPEUP/HEPRUP"
    write (u, "(A)")
 
    sample = "eio_stdhep_4"
 
    allocate (eio_stdhep_hepeup_t :: eio)
    select type (eio)
    type is (eio_stdhep_hepeup_t)
       call eio%set_parameters ()
    end select  
    call eio%set_fallback_model (fallback_model)
    
    call eio%init_out (sample, data)
    call event%generate (1, [0._default, 0._default])
    call event%evaluate_expressions ()
    call event%pacify_particle_set ()    

    call eio%output (event, i_prc = 1)
    call eio%write (u)
    call eio%final ()

    call eio_cleanup_test (event)
    call eio_cleanup_fallback_model (fallback_model)    
    deallocate (eio)
    
    write (u, "(A)")  "* Initialize test process"
    write (u, "(A)")

    call eio_prepare_fallback_model (fallback_model)
    call eio_prepare_test (event, unweighted = .false.)

    allocate (eio_stdhep_hepeup_t :: eio)
    select type (eio)
    type is (eio_stdhep_hepeup_t)
       call eio%set_parameters (recover_beams = .false.)
    end select
    call eio%set_fallback_model (fallback_model)
    
    call data%init (1)
    data%n_beam = 2
    data%unweighted = .true.
    data%norm_mode = NORM_UNIT
    data%pdg_beam = 25
    data%energy_beam = 500
    data%proc_num_id = [42]
    call data%write (u)
    write (u, *)

    write (u, "(A)")  "* Initialize"
    write (u, "(A)")

    call eio%init_in (sample, data)
    call eio%write (u)
    
    write (u, "(A)")
    write (u, "(A)")  "* Read event"
    write (u, "(A)")

    call eio%input_i_prc (i_prc, iostat)
    
    select type (eio)
    type is (eio_stdhep_hepeup_t)
       write (u, "(A,I0,A,I0)")  "Found process #", i_prc, &
            " with ID = ", eio%proc_num_id(i_prc)
    end select

    call eio%input_event (event, iostat)

    call event%write (u)

    write (u, "(A)")
    write (u, "(A)")  "* Read closing"
    write (u, "(A)")

    call eio%input_i_prc (i_prc, iostat)
    write (u, "(A,I0)")  "iostat = ", iostat
    
    write (u, "(A)")
    write (u, "(A)")  "* Cleanup"
 
    call eio%final ()

    call eio_cleanup_test (event)
    call eio_cleanup_fallback_model (fallback_model)

    write (u, "(A)")
    write (u, "(A)")  "* Test output end: eio_stdhep_4"
    
  end subroutine eio_stdhep_4
  
\eatline
\nwendcode{}\nwbegindocs{1222}\nwdocspar
%%%%%%%%%%%%%%%%%%%%%%%%%%%%%%%%%%%%%%%%%%%%%%%%%%%%%%%%%%%%%%%%%%%%%%%%
\section{HepMC Output}
The HepMC event record is standardized.  It is an ASCII format.  We try
our best at using it for both input and output.
\nwenddocs{}\nwbegincode{1223}\moddef{\code{}eio{\char95}hepmc.f90\edoc{}}\endmoddef
\LA{}File header\RA{}

module eio_hepmc
  
  use kinds
  use io_units
\LA{}Use strings\RA{}
  use string_utils
  use unit_tests
  use diagnostics
  use os_interface
  use lorentz
  use model_data
  use particles
  use subevents
  use event_base
  use hep_events
  use eio_data
  use eio_base
  use hepmc_interface

\LA{}Standard module head\RA{}

\LA{}EIO HepMC: public\RA{}

\LA{}EIO HepMC: types\RA{}

\LA{}EIO HepMC: interfaces\RA{}

contains
  
\LA{}EIO HepMC: procedures\RA{}

\LA{}EIO HepMC: tests\RA{}

end module eio_hepmc
\eatline
\nwendcode{}\nwbegindocs{1224}\nwdocspar
\subsection{Type}
A type {\tt{}hepmc{\char95}event} is introduced as container to store HepMC event
data, particularly for splitting the reading into read out of the process
index and the proper event data.

Note: the {\tt{}keep{\char95}beams} flag is not supported.  Beams will always
be written.
\nwenddocs{}\nwbegincode{1225}\moddef{EIO HepMC: public}\endmoddef
  public :: eio_hepmc_t
\nwendcode{}\nwbegincode{1226}\moddef{EIO HepMC: types}\endmoddef
  type, extends (eio_t) :: eio_hepmc_t
     logical :: writing = .false.
     logical :: reading = .false.
     ! logical :: keep_beams = .false.
     logical :: recover_beams = .false.
     logical :: use_alpha_s_from_file = .false.
     logical :: use_scale_from_file = .false.
     type(hepmc_iostream_t) :: iostream
     type(hepmc_event_t) :: hepmc_event
     integer, dimension(:), allocatable :: proc_num_id
   contains
   \LA{}EIO HepMC: eio hepmc: TBP\RA{}
  end type eio_hepmc_t
  
\eatline
\nwendcode{}\nwbegindocs{1227}\nwdocspar
\subsection{Specific Methods}
Set parameters that are specifically used with HepMC.
\nwenddocs{}\nwbegincode{1228}\moddef{EIO HepMC: eio hepmc: TBP}\endmoddef
  procedure :: set_parameters => eio_hepmc_set_parameters
\nwendcode{}\nwbegincode{1229}\moddef{EIO HepMC: procedures}\endmoddef
  ! subroutine eio_hepmc_set_parameters (eio, keep_beams, recover_beams, extension)
  subroutine eio_hepmc_set_parameters &
       (eio, recover_beams, use_alpha_s_from_file, use_scale_from_file, &
       extension)
    class(eio_hepmc_t), intent(inout) :: eio
    logical, intent(in), optional :: recover_beams 
    logical, intent(in), optional :: use_alpha_s_from_file
    logical, intent(in), optional :: use_scale_from_file
    type(string_t), intent(in), optional :: extension    
    if (present (recover_beams))  eio%recover_beams = recover_beams
    if (present (use_alpha_s_from_file)) &
         eio%use_alpha_s_from_file = use_alpha_s_from_file
    if (present (use_scale_from_file)) &
         eio%use_scale_from_file = use_scale_from_file
    if (present (extension)) then
       eio%extension = extension
    else
       eio%extension = "hepmc"
    end if
  end subroutine eio_hepmc_set_parameters
  
\eatline
\nwendcode{}\nwbegindocs{1230}\nwdocspar
\subsection{Common Methods}
Output.  This is not the actual event format, but a readable account
of the current object status.
\nwenddocs{}\nwbegincode{1231}\moddef{EIO HepMC: eio hepmc: TBP}\plusendmoddef
  procedure :: write => eio_hepmc_write
\nwendcode{}\nwbegincode{1232}\moddef{EIO HepMC: procedures}\plusendmoddef
  subroutine eio_hepmc_write (object, unit)
    class(eio_hepmc_t), intent(in) :: object
    integer, intent(in), optional :: unit
    integer :: u, i
    u = given_output_unit (unit)
    write (u, "(1x,A)")  "HepMC event stream:"
    if (object%writing) then
       write (u, "(3x,A,A)")  "Writing to file   = ", char (object%filename)
    else if (object%reading) then
       write (u, "(3x,A,A)")  "Reading from file = ", char (object%filename)
    else
       write (u, "(3x,A)")  "[closed]"
    end if
    write (u, "(3x,A,L1)")    "Recover beams     = ", object%recover_beams
    write (u, "(3x,A,L1)")    "Alpha_s from file = ", &
         object%use_alpha_s_from_file
    write (u, "(3x,A,L1)")    "Scale from file   = ", &
         object%use_scale_from_file
    write (u, "(3x,A,A,A)")     "File extension    = '", &
         char (object%extension), "'"
    if (allocated (object%proc_num_id)) then
       write (u, "(3x,A)")  "Numerical process IDs:"
       do i = 1, size (object%proc_num_id)
          write (u, "(5x,I0,': ',I0)")  i, object%proc_num_id(i)
       end do
    end if    
  end subroutine eio_hepmc_write
  
\eatline
\nwendcode{}\nwbegindocs{1233}Finalizer: close any open file.
\nwenddocs{}\nwbegincode{1234}\moddef{EIO HepMC: eio hepmc: TBP}\plusendmoddef
  procedure :: final => eio_hepmc_final
\nwendcode{}\nwbegincode{1235}\moddef{EIO HepMC: procedures}\plusendmoddef
  subroutine eio_hepmc_final (object)
    class(eio_hepmc_t), intent(inout) :: object
    if (allocated (object%proc_num_id))  deallocate (object%proc_num_id)
    if (object%writing) then
       write (msg_buffer, "(A,A,A)")  "Events: closing HepMC file '", &
            char (object%filename), "'"
       call msg_message ()
       call hepmc_iostream_close (object%iostream)
       object%writing = .false.
    else if (object%reading) then
       write (msg_buffer, "(A,A,A)")  "Events: closing HepMC file '", &
            char (object%filename), "'"
       call msg_message ()
       call hepmc_iostream_close (object%iostream)
       object%reading = .false.
    end if
  end subroutine eio_hepmc_final
  
\eatline
\nwendcode{}\nwbegindocs{1236}Split event file: increment the counter, close the current file, open a new
one.  If the file needs a header, repeat it for the new file.
\nwenddocs{}\nwbegincode{1237}\moddef{EIO HepMC: eio hepmc: TBP}\plusendmoddef
  procedure :: split_out => eio_hepmc_split_out
\nwendcode{}\nwbegincode{1238}\moddef{EIO HepMC: procedures}\plusendmoddef
  subroutine eio_hepmc_split_out (eio)
    class(eio_hepmc_t), intent(inout) :: eio
    if (eio%split) then
       eio%split_index = eio%split_index + 1
       call eio%set_filename ()
       write (msg_buffer, "(A,A,A)")  "Events: writing to HepMC file '", &
            char (eio%filename), "'"
       call msg_message ()
       call hepmc_iostream_close (eio%iostream)
       call hepmc_iostream_open_out (eio%iostream, eio%filename)
    end if
  end subroutine eio_hepmc_split_out
  
\eatline
\nwendcode{}\nwbegindocs{1239}Common initialization for input and output.
\nwenddocs{}\nwbegincode{1240}\moddef{EIO HepMC: eio hepmc: TBP}\plusendmoddef
  procedure :: common_init => eio_hepmc_common_init
\nwendcode{}\nwbegincode{1241}\moddef{EIO HepMC: procedures}\plusendmoddef
  subroutine eio_hepmc_common_init (eio, sample, data, extension)
    class(eio_hepmc_t), intent(inout) :: eio
    type(string_t), intent(in) :: sample
    type(string_t), intent(in), optional :: extension
    type(event_sample_data_t), intent(in), optional :: data
    if (.not. present (data)) &
         call msg_bug ("HepMC initialization: missing data")
    if (data%n_beam /= 2) &
         call msg_fatal ("HepMC: defined for scattering processes only")    
    if (data%unweighted) then
       select case (data%norm_mode)
       case (NORM_UNIT)
       case default; call msg_fatal &
            ("HepMC: normalization for unweighted events must be '1'")
       end select
    else
       call msg_fatal ("HepMC: events must be unweighted")    
    end if
    eio%sample = sample    
    if (present (extension)) then
       eio%extension = extension
    end if
    call eio%set_filename ()
    allocate (eio%proc_num_id (data%n_proc), source = data%proc_num_id)
  end subroutine eio_hepmc_common_init
  
\eatline
\nwendcode{}\nwbegindocs{1242}Initialize event writing.
\nwenddocs{}\nwbegincode{1243}\moddef{EIO HepMC: eio hepmc: TBP}\plusendmoddef
  procedure :: init_out => eio_hepmc_init_out
\nwendcode{}\nwbegincode{1244}\moddef{EIO HepMC: procedures}\plusendmoddef
  subroutine eio_hepmc_init_out (eio, sample, data, success, extension)
    class(eio_hepmc_t), intent(inout) :: eio
    type(string_t), intent(in) :: sample
    type(string_t), intent(in), optional :: extension
    type(event_sample_data_t), intent(in), optional :: data
    logical, intent(out), optional :: success
    call eio%set_splitting (data)    
    call eio%common_init (sample, data, extension)
    write (msg_buffer, "(A,A,A)")  "Events: writing to HepMC file '", &
         char (eio%filename), "'"
    call msg_message ()
    eio%writing = .true.
    call hepmc_iostream_open_out (eio%iostream, eio%filename)
    if (present (success))  success = .true.
  end subroutine eio_hepmc_init_out
    
\eatline
\nwendcode{}\nwbegindocs{1245}Initialize event reading. For input, we do not (yet) support split
event files. 
\nwenddocs{}\nwbegincode{1246}\moddef{EIO HepMC: eio hepmc: TBP}\plusendmoddef
  procedure :: init_in => eio_hepmc_init_in
\nwendcode{}\nwbegincode{1247}\moddef{EIO HepMC: procedures}\plusendmoddef
  subroutine eio_hepmc_init_in (eio, sample, data, success, extension)
    class(eio_hepmc_t), intent(inout) :: eio
    type(string_t), intent(in) :: sample
    type(string_t), intent(in), optional :: extension
    type(event_sample_data_t), intent(inout), optional :: data
    logical, intent(out), optional :: success
    logical :: exist
    eio%split = .false.
    call eio%common_init (sample, data, extension)
    write (msg_buffer, "(A,A,A)")  "Events: reading from HepMC file '", &
         char (eio%filename), "'"
    call msg_message ()
    inquire (file = char (eio%filename), exist = exist)
    if (.not. exist)  call msg_fatal ("Events: HepMC file not found.")
    eio%reading = .true.
    call hepmc_iostream_open_in (eio%iostream, eio%filename)
    if (present (success))  success = .true.
  end subroutine eio_hepmc_init_in
    
\eatline
\nwendcode{}\nwbegindocs{1248}Switch from input to output: reopen the file for reading.
\nwenddocs{}\nwbegincode{1249}\moddef{EIO HepMC: eio hepmc: TBP}\plusendmoddef
  procedure :: switch_inout => eio_hepmc_switch_inout
\nwendcode{}\nwbegincode{1250}\moddef{EIO HepMC: procedures}\plusendmoddef
  subroutine eio_hepmc_switch_inout (eio, success)
    class(eio_hepmc_t), intent(inout) :: eio
    logical, intent(out), optional :: success
    call msg_bug ("HepMC: in-out switch not supported")
    if (present (success))  success = .false.
  end subroutine eio_hepmc_switch_inout
  
\eatline
\nwendcode{}\nwbegindocs{1251}Output an event to the allocated HepMC output stream.
\nwenddocs{}\nwbegincode{1252}\moddef{EIO HepMC: eio hepmc: TBP}\plusendmoddef
  procedure :: output => eio_hepmc_output
\nwendcode{}\nwbegincode{1253}\moddef{EIO HepMC: procedures}\plusendmoddef
  subroutine eio_hepmc_output (eio, event, i_prc, reading, passed, pacify)
    class(eio_hepmc_t), intent(inout) :: eio
    class(generic_event_t), intent(in), target :: event
    integer, intent(in) :: i_prc
    logical, intent(in), optional :: reading, passed, pacify
    type(particle_set_t), pointer :: pset_ptr
    if (present (passed)) then
       if (.not. passed)  return
    end if
    if (eio%writing) then
       pset_ptr => event%get_particle_set_ptr ()
       call hepmc_event_init (eio%hepmc_event, &
            proc_id = eio%proc_num_id (i_prc), &
            event_id = event%get_index ())
       call hepmc_event_from_particle_set (eio%hepmc_event, pset_ptr)
       call hepmc_event_set_scale (eio%hepmc_event, event%get_fac_scale ())
       call hepmc_event_set_alpha_qcd (eio%hepmc_event, event%get_alpha_s ())
       call hepmc_iostream_write_event (eio%iostream, eio%hepmc_event)
       call hepmc_event_final (eio%hepmc_event)
    else
       call eio%write ()
       call msg_fatal ("HepMC file is not open for writing")
    end if
  end subroutine eio_hepmc_output

\eatline
\nwendcode{}\nwbegindocs{1254}Input an event.
\nwenddocs{}\nwbegincode{1255}\moddef{EIO HepMC: eio hepmc: TBP}\plusendmoddef
  procedure :: input_i_prc => eio_hepmc_input_i_prc
  procedure :: input_event => eio_hepmc_input_event
\nwendcode{}\nwbegincode{1256}\moddef{EIO HepMC: procedures}\plusendmoddef
  subroutine eio_hepmc_input_i_prc (eio, i_prc, iostat)
    class(eio_hepmc_t), intent(inout) :: eio
    integer, intent(out) :: i_prc
    integer, intent(out) :: iostat
    logical :: ok
    integer :: i, proc_num_id
    iostat = 0
    call hepmc_event_init (eio%hepmc_event)
    call hepmc_iostream_read_event (eio%iostream, eio%hepmc_event, ok)
    proc_num_id = hepmc_event_get_process_id (eio%hepmc_event)
    if (.not. ok) then 
       iostat = -1 
       return
    end if
    i_prc = 0
    FIND_I_PRC: do i = 1, size (eio%proc_num_id)
       if (eio%proc_num_id(i) == proc_num_id) then
          i_prc = i
          exit FIND_I_PRC
       end if
    end do FIND_I_PRC
    if (i_prc == 0)  call err_index
  contains
    subroutine err_index
      call msg_error ("HepMC: reading events: undefined process ID " &
           // char (str (proc_num_id)) // ", aborting read")
      iostat = 1
    end subroutine err_index
  end subroutine eio_hepmc_input_i_prc

  subroutine eio_hepmc_input_event (eio, event, iostat)
    class(eio_hepmc_t), intent(inout) :: eio
    class(generic_event_t), intent(inout), target :: event
    integer, intent(out) :: iostat
    type(particle_set_t), pointer :: pset
    logical :: ok    
    iostat = 0
    call event%reset ()
    call event%select (1, 1, 1)
    call hepmc_to_event (event, eio%hepmc_event, eio%fallback_model, &
         recover_beams = eio%recover_beams, &
         use_alpha_s = eio%use_alpha_s_from_file, &
         use_scale = eio%use_scale_from_file) 
    call hepmc_event_final (eio%hepmc_event)
  end subroutine eio_hepmc_input_event

\eatline
\eatline
\nwendcode{}\nwbegindocs{1257}\nwdocspar
\subsection{Unit tests}
\nwenddocs{}\nwbegincode{1258}\moddef{EIO HepMC: public}\plusendmoddef
  public :: eio_hepmc_test
\nwendcode{}\nwbegincode{1259}\moddef{EIO HepMC: tests}\endmoddef
  subroutine eio_hepmc_test (u, results)
    integer, intent(in) :: u
    type(test_results_t), intent(inout) :: results
  \LA{}EIO HepMC: execute tests\RA{}
  end subroutine eio_hepmc_test
  
\eatline
\nwendcode{}\nwbegindocs{1260}\nwdocspar
\subsubsection{Test I/O methods}
We test the implementation of all I/O methods.
\nwenddocs{}\nwbegincode{1261}\moddef{EIO HepMC: execute tests}\endmoddef
  call test (eio_hepmc_1, "eio_hepmc_1", &
       "write event contents", &
       u, results)
\nwendcode{}\nwbegincode{1262}\moddef{EIO HepMC: tests}\plusendmoddef
  subroutine eio_hepmc_1 (u)
    integer, intent(in) :: u
    class(generic_event_t), pointer :: event
    type(event_sample_data_t) :: data
    class(eio_t), allocatable :: eio
    type(string_t) :: sample
    integer :: u_file, iostat
    character(116) :: buffer

    write (u, "(A)")  "* Test output: eio_hepmc_1"
    write (u, "(A)")  "*   Purpose: write a HepMC file"
    write (u, "(A)")

    write (u, "(A)")  "* Initialize test process"
 
    call eio_prepare_test (event, unweighted=.false.)

    call data%init (1)
    data%n_beam = 2
    data%unweighted = .true.
    data%norm_mode = NORM_UNIT
    data%pdg_beam = 25
    data%energy_beam = 500
    data%proc_num_id = [42]
    data%cross_section(1) = 100
    data%error(1) = 1
    data%total_cross_section = sum (data%cross_section)

    write (u, "(A)")
    write (u, "(A)")  "* Generate and write an event"
    write (u, "(A)")
 
    sample = "eio_hepmc_1"
 
    allocate (eio_hepmc_t :: eio)
    select type (eio)
    type is (eio_hepmc_t)
       call eio%set_parameters ()
    end select
    
    call eio%init_out (sample, data)
    call event%generate (1, [0._default, 0._default])

    call eio%output (event, i_prc = 1)
    call eio%write (u)
    call eio%final ()

    write (u, "(A)")
    write (u, "(A)")  "* File contents (blanking out last two digits):"
    write (u, "(A)")

    u_file = free_unit ()
    open (u_file, file = char (sample // ".hepmc"), &
         action = "read", status = "old")
    do
       read (u_file, "(A)", iostat = iostat)  buffer
       if (iostat /= 0)  exit
       if (trim (buffer) == "")  cycle
       if (buffer(1:14) == "HepMC::Version")  cycle
       if (buffer(1:10) == "P 10001 25") &
            call buffer_blanker (buffer, 32, 55, 78)
       if (buffer(1:10) == "P 10002 25") &
            call buffer_blanker (buffer, 33, 56, 79)
       if (buffer(1:10) == "P 10003 25") &
            call buffer_blanker (buffer, 29, 53, 78, 101)
       if (buffer(1:10) == "P 10004 25") &
            call buffer_blanker (buffer, 28, 51, 76, 99)       
       write (u, "(A)") trim (buffer)
    end do
    close (u_file)
    
    write (u, "(A)")
    write (u, "(A)")  "* Reset data"
    write (u, "(A)")
 
    deallocate (eio)
    allocate (eio_hepmc_t :: eio)
    
    select type (eio)
    type is (eio_hepmc_t)
       call eio%set_parameters ()
    end select
    call eio%write (u)

    write (u, "(A)")
    write (u, "(A)")  "* Cleanup"
 
    call eio_cleanup_test (event)

    write (u, "(A)")
    write (u, "(A)")  "* Test output end: eio_hepmc_1"
    
  contains
    
    subroutine buffer_blanker (buf, pos1, pos2, pos3, pos4)
      character(len=*), intent(inout) :: buf
      integer, intent(in) :: pos1, pos2, pos3
      integer, intent(in), optional :: pos4
      type(string_t) :: line
      line = var_str (trim (buf))
      line = replace (line, pos1, "XX")
      line = replace (line, pos2, "XX")
      line = replace (line, pos3, "XX")
      if (present (pos4)) then
         line = replace (line, pos4, "XX")
      end if
      line = replace (line, "4999999999999", "5000000000000")
      buf = char (line)
    end subroutine buffer_blanker
    
  end subroutine eio_hepmc_1
  
\eatline
\nwendcode{}\nwbegindocs{1263}Test also the reading of HepMC events.
\nwenddocs{}\nwbegincode{1264}\moddef{EIO HepMC: execute tests}\plusendmoddef
  call test (eio_hepmc_2, "eio_hepmc_2", &
       "read event contents", &
       u, results)
\nwendcode{}\nwbegincode{1265}\moddef{EIO HepMC: tests}\plusendmoddef
  subroutine eio_hepmc_2 (u)
    integer, intent(in) :: u
    class(model_data_t), pointer :: fallback_model
    class(generic_event_t), pointer :: event
    type(event_sample_data_t) :: data
    class(eio_t), allocatable :: eio
    type(string_t) :: sample
    integer :: u_file, iostat, i_prc

    write (u, "(A)")  "* Test output: eio_hepmc_2"
    write (u, "(A)")  "*   Purpose: read a HepMC event"
    write (u, "(A)")

    write (u, "(A)")  "* Write a HepMC data file"
    write (u, "(A)")
    
    u_file = free_unit ()
    sample = "eio_hepmc_2"
    open (u_file, file = char (sample // ".hepmc"), &
         status = "replace", action = "readwrite")
    
    write (u_file, "(A)")  "HepMC::Version 2.06.09"
    write (u_file, "(A)")  "HepMC::IO_GenEvent-START_EVENT_LISTING"
    write (u_file, "(A)")  "E 0 -1 -1.0000000000000000e+00 &
         &-1.0000000000000000e+00 &
         &-1.0000000000000000e+00 42 0 1 10001 10002 0 0"
    write (u_file, "(A)")  "U GEV MM"
    write (u_file, "(A)")  "V -1 0 0 0 0 0 2 2 0"
    write (u_file, "(A)")  "P 10001 25 0 0 4.8412291827592713e+02 &
         &5.0000000000000000e+02 &
         &1.2499999999999989e+02 3 0 0 -1 0"
    write (u_file, "(A)")  "P 10002 25 0 0 -4.8412291827592713e+02 &
         &5.0000000000000000e+02 &
         &1.2499999999999989e+02 3 0 0 -1 0"
    write (u_file, "(A)")  "P 10003 25 -1.4960220911365536e+02 &
         &-4.6042825611414656e+02 &
         &0 5.0000000000000000e+02 1.2500000000000000e+02 1 0 0 0 0"
    write (u_file, "(A)")  "P 10004 25 1.4960220911365536e+02 &
         &4.6042825611414656e+02 &
         &0 5.0000000000000000e+02 1.2500000000000000e+02 1 0 0 0 0"
    write (u_file, "(A)")  "HepMC::IO_GenEvent-END_EVENT_LISTING"
    close (u_file)

    write (u, "(A)")  "* Initialize test process" 
    write (u, "(A)")
    
    call eio_prepare_fallback_model (fallback_model)
    call eio_prepare_test (event, unweighted=.false.)

    allocate (eio_hepmc_t :: eio)
    select type (eio)
    type is (eio_hepmc_t)
       call eio%set_parameters (recover_beams = .false.)
    end select            
    call eio%set_fallback_model (fallback_model)
    
    call data%init (1)
    data%n_beam = 2
    data%unweighted = .true.
    data%norm_mode = NORM_UNIT
    data%pdg_beam = 25
    data%energy_beam = 500
    data%proc_num_id = [42]
    call data%write (u)

    write (u, "(A)")    
    write (u, "(A)")  "* Initialize"
    write (u, "(A)")
             
    call eio%init_in (sample, data)
    call eio%write (u)    
    
    write (u, "(A)")
    write (u, "(A)")  "* Read event"
    write (u, "(A)")
 
    call eio%input_i_prc (i_prc, iostat)

    select type (eio)
    type is (eio_hepmc_t)
       write (u, "(A,I0,A,I0)")  "Found process #", i_prc, &
            " with ID = ", eio%proc_num_id(i_prc)
    end select
    
    call eio%input_event (event, iostat)
    
    call event%write (u)
        
    write (u, "(A)")
    write (u, "(A)")  "* Read closing"
    write (u, "(A)")
    
    call eio%input_i_prc (i_prc, iostat)
    write (u, "(A,I0)")  "iostat = ", iostat

    write (u, "(A)")
    write (u, "(A)")  "* Cleanup"

    call eio%final ()

    call eio_cleanup_test (event)
    call eio_cleanup_fallback_model (fallback_model)

    write (u, "(A)")
    write (u, "(A)")  "* Test output end: eio_hepmc_2"
    
  end subroutine eio_hepmc_2
  
\eatline
\nwendcode{}\nwbegindocs{1266}\nwdocspar
%%%%%%%%%%%%%%%%%%%%%%%%%%%%%%%%%%%%%%%%%%%%%%%%%%%%%%%%%%%%%%%%%%%%%%%%
\section{LCIO Output}
The LCIO event record is standardized for the use with Linear $e^+e^-$
colliders. It is a binary event format.  We try our best at using it
for both input and output. 
\nwenddocs{}\nwbegincode{1267}\moddef{\code{}eio{\char95}lcio.f90\edoc{}}\endmoddef
\LA{}File header\RA{}

module eio_lcio
  
  use kinds
  use io_units
\LA{}Use strings\RA{}
  use string_utils
  use unit_tests
  use diagnostics
  use os_interface
  use lorentz
  use model_data
  use particles
  use subevents
  use event_base
  use hep_events
  use eio_data
  use eio_base
  use lcio_interface

\LA{}Standard module head\RA{}

\LA{}EIO LCIO: public\RA{}

\LA{}EIO LCIO: types\RA{}

\LA{}EIO LCIO: interfaces\RA{}

contains
  
\LA{}EIO LCIO: procedures\RA{}

\LA{}EIO LCIO: tests\RA{}

end module eio_lcio
\eatline
\nwendcode{}\nwbegindocs{1268}\nwdocspar
\subsection{Type}
A type {\tt{}lcio{\char95}event} is introduced as container to store LCIO event
data, particularly for splitting the reading into read out of the process
index and the proper event data.

Note: the {\tt{}keep{\char95}beams} flag is not supported.
\nwenddocs{}\nwbegincode{1269}\moddef{EIO LCIO: public}\endmoddef
  public :: eio_lcio_t
\nwendcode{}\nwbegincode{1270}\moddef{EIO LCIO: types}\endmoddef
  type, extends (eio_t) :: eio_lcio_t
     logical :: writing = .false.
     logical :: reading = .false.
     logical :: recover_beams = .false.
     logical :: use_alpha_s_from_file = .false.
     logical :: use_scale_from_file = .false.
     type(lcio_writer_t) :: lcio_writer
     type(lcio_reader_t) :: lcio_reader
     type(lcio_run_header_t) :: lcio_run_hdr
     type(lcio_event_t) :: lcio_event
     integer, dimension(:), allocatable :: proc_num_id
   contains
   \LA{}EIO LCIO: eio lcio: TBP\RA{}
  end type eio_lcio_t
  
\eatline
\nwendcode{}\nwbegindocs{1271}\nwdocspar
\subsection{Specific Methods}
Set parameters that are specifically used with LCIO.
\nwenddocs{}\nwbegincode{1272}\moddef{EIO LCIO: eio lcio: TBP}\endmoddef
  procedure :: set_parameters => eio_lcio_set_parameters
\nwendcode{}\nwbegincode{1273}\moddef{EIO LCIO: procedures}\endmoddef
  subroutine eio_lcio_set_parameters &
       (eio, recover_beams, use_alpha_s_from_file, use_scale_from_file, &
       extension)
    class(eio_lcio_t), intent(inout) :: eio
    logical, intent(in), optional :: recover_beams 
    logical, intent(in), optional :: use_alpha_s_from_file
    logical, intent(in), optional :: use_scale_from_file
    type(string_t), intent(in), optional :: extension    
    if (present (recover_beams))  eio%recover_beams = recover_beams
    if (present (use_alpha_s_from_file)) &
         eio%use_alpha_s_from_file = use_alpha_s_from_file
    if (present (use_scale_from_file)) &
         eio%use_scale_from_file = use_scale_from_file
    if (present (extension)) then
       eio%extension = extension
    else
       eio%extension = "slcio"
    end if
  end subroutine eio_lcio_set_parameters
  
\eatline
\nwendcode{}\nwbegindocs{1274}\nwdocspar
\subsection{Common Methods}
Output.  This is not the actual event format, but a readable account
of the current object status.
\nwenddocs{}\nwbegincode{1275}\moddef{EIO LCIO: eio lcio: TBP}\plusendmoddef
  procedure :: write => eio_lcio_write
\nwendcode{}\nwbegincode{1276}\moddef{EIO LCIO: procedures}\plusendmoddef
  subroutine eio_lcio_write (object, unit)
    class(eio_lcio_t), intent(in) :: object
    integer, intent(in), optional :: unit
    integer :: u, i
    u = given_output_unit (unit)
    write (u, "(1x,A)")  "LCIO event stream:"
    if (object%writing) then
       write (u, "(3x,A,A)")  "Writing to file   = ", char (object%filename)
    else if (object%reading) then
       write (u, "(3x,A,A)")  "Reading from file = ", char (object%filename)
    else
       write (u, "(3x,A)")  "[closed]"
    end if
    write (u, "(3x,A,L1)")    "Recover beams     = ", object%recover_beams
    write (u, "(3x,A,L1)")    "Alpha_s from file = ", &
         object%use_alpha_s_from_file
    write (u, "(3x,A,L1)")    "Scale from file   = ", &
         object%use_scale_from_file
    write (u, "(3x,A,A,A)")     "File extension    = '", &
         char (object%extension), "'"
    if (allocated (object%proc_num_id)) then
       write (u, "(3x,A)")  "Numerical process IDs:"
       do i = 1, size (object%proc_num_id)
          write (u, "(5x,I0,': ',I0)")  i, object%proc_num_id(i)
       end do
    end if    
  end subroutine eio_lcio_write
  
\eatline
\nwendcode{}\nwbegindocs{1277}Finalizer: close any open file.
\nwenddocs{}\nwbegincode{1278}\moddef{EIO LCIO: eio lcio: TBP}\plusendmoddef
  procedure :: final => eio_lcio_final
\nwendcode{}\nwbegincode{1279}\moddef{EIO LCIO: procedures}\plusendmoddef
  subroutine eio_lcio_final (object)
    class(eio_lcio_t), intent(inout) :: object
    if (allocated (object%proc_num_id))  deallocate (object%proc_num_id)
    if (object%writing) then
       write (msg_buffer, "(A,A,A)")  "Events: closing LCIO file '", &
            char (object%filename), "'"
       call msg_message ()
       call lcio_writer_close (object%lcio_writer)
       object%writing = .false.
    else if (object%reading) then
       write (msg_buffer, "(A,A,A)")  "Events: closing LCIO file '", &
            char (object%filename), "'"
       call msg_message ()
       call lcio_reader_close (object%lcio_reader)
       object%reading = .false.
    end if
  end subroutine eio_lcio_final
  
\eatline
\nwendcode{}\nwbegindocs{1280}Split event file: increment the counter, close the current file, open a new
one.  If the file needs a header, repeat it for the new file.
\nwenddocs{}\nwbegincode{1281}\moddef{EIO LCIO: eio lcio: TBP}\plusendmoddef
  procedure :: split_out => eio_lcio_split_out
\nwendcode{}\nwbegincode{1282}\moddef{EIO LCIO: procedures}\plusendmoddef
  subroutine eio_lcio_split_out (eio)
    class(eio_lcio_t), intent(inout) :: eio
    if (eio%split) then
       eio%split_index = eio%split_index + 1
       call eio%set_filename ()
       write (msg_buffer, "(A,A,A)")  "Events: writing to LCIO file '", &
            char (eio%filename), "'"
       call msg_message ()
       call lcio_writer_close (eio%lcio_writer)
       call lcio_writer_open_out (eio%lcio_writer, eio%filename)       
    end if
  end subroutine eio_lcio_split_out
  
\eatline
\nwendcode{}\nwbegindocs{1283}Common initialization for input and output.
\nwenddocs{}\nwbegincode{1284}\moddef{EIO LCIO: eio lcio: TBP}\plusendmoddef
  procedure :: common_init => eio_lcio_common_init
\nwendcode{}\nwbegincode{1285}\moddef{EIO LCIO: procedures}\plusendmoddef
  subroutine eio_lcio_common_init (eio, sample, data, extension)
    class(eio_lcio_t), intent(inout) :: eio
    type(string_t), intent(in) :: sample
    type(string_t), intent(in), optional :: extension
    type(event_sample_data_t), intent(in), optional :: data
    if (.not. present (data)) &
         call msg_bug ("LCIO initialization: missing data")
    if (data%n_beam /= 2) &
         call msg_fatal ("LCIO: defined for scattering processes only")    
    if (data%unweighted) then
       select case (data%norm_mode)
       case (NORM_UNIT)
       case default; call msg_fatal &
            ("LCIO: normalization for unweighted events must be '1'")
       end select
    else
       call msg_fatal ("LCIO: events must be unweighted")    
    end if
    eio%sample = sample    
    if (present (extension)) then
       eio%extension = extension
    end if
    call eio%set_filename ()
    allocate (eio%proc_num_id (data%n_proc), source = data%proc_num_id)
  end subroutine eio_lcio_common_init
  
\eatline
\nwendcode{}\nwbegindocs{1286}Initialize event writing.
\nwenddocs{}\nwbegincode{1287}\moddef{EIO LCIO: eio lcio: TBP}\plusendmoddef
  procedure :: init_out => eio_lcio_init_out
\nwendcode{}\nwbegincode{1288}\moddef{EIO LCIO: procedures}\plusendmoddef
  subroutine eio_lcio_init_out (eio, sample, data, success, extension)
    class(eio_lcio_t), intent(inout) :: eio
    type(string_t), intent(in) :: sample
    type(string_t), intent(in), optional :: extension
    type(event_sample_data_t), intent(in), optional :: data
    logical, intent(out), optional :: success
    call eio%set_splitting (data)    
    call eio%common_init (sample, data, extension)
    write (msg_buffer, "(A,A,A)")  "Events: writing to LCIO file '", &
         char (eio%filename), "'"
    call msg_message ()
    eio%writing = .true.
    call lcio_writer_open_out (eio%lcio_writer, eio%filename)
    call lcio_run_header_init (eio%lcio_run_hdr)
    call lcio_run_header_write (eio%lcio_writer, eio%lcio_run_hdr)
    if (present (success))  success = .true.
  end subroutine eio_lcio_init_out
    
\eatline
\nwendcode{}\nwbegindocs{1289}Initialize event reading. For input, we do not (yet) support split
event files. 
\nwenddocs{}\nwbegincode{1290}\moddef{EIO LCIO: eio lcio: TBP}\plusendmoddef
  procedure :: init_in => eio_lcio_init_in
\nwendcode{}\nwbegincode{1291}\moddef{EIO LCIO: procedures}\plusendmoddef
  subroutine eio_lcio_init_in (eio, sample, data, success, extension)
    class(eio_lcio_t), intent(inout) :: eio
    type(string_t), intent(in) :: sample
    type(string_t), intent(in), optional :: extension
    type(event_sample_data_t), intent(inout), optional :: data
    logical, intent(out), optional :: success
    logical :: exist
    eio%split = .false.
    call eio%common_init (sample, data, extension)
    write (msg_buffer, "(A,A,A)")  "Events: reading from LCIO file '", &
         char (eio%filename), "'"
    call msg_message ()
    inquire (file = char (eio%filename), exist = exist)
    if (.not. exist)  call msg_fatal ("Events: LCIO file not found.")
    eio%reading = .true.
    call lcio_open_file (eio%lcio_reader, eio%filename)
    if (present (success))  success = .true.
  end subroutine eio_lcio_init_in
    
\eatline
\nwendcode{}\nwbegindocs{1292}Switch from input to output: reopen the file for reading.
\nwenddocs{}\nwbegincode{1293}\moddef{EIO LCIO: eio lcio: TBP}\plusendmoddef
  procedure :: switch_inout => eio_lcio_switch_inout
\nwendcode{}\nwbegincode{1294}\moddef{EIO LCIO: procedures}\plusendmoddef
  subroutine eio_lcio_switch_inout (eio, success)
    class(eio_lcio_t), intent(inout) :: eio
    logical, intent(out), optional :: success
    call msg_bug ("LCIO: in-out switch not supported")
    if (present (success))  success = .false.
  end subroutine eio_lcio_switch_inout
  
\eatline
\nwendcode{}\nwbegindocs{1295}Output an event to the allocated LCIO writer.
\nwenddocs{}\nwbegincode{1296}\moddef{EIO LCIO: eio lcio: TBP}\plusendmoddef
  procedure :: output => eio_lcio_output
\nwendcode{}\nwbegincode{1297}\moddef{EIO LCIO: procedures}\plusendmoddef
  subroutine eio_lcio_output (eio, event, i_prc, reading, passed, pacify)
    class(eio_lcio_t), intent(inout) :: eio
    class(generic_event_t), intent(in), target :: event
    integer, intent(in) :: i_prc
    logical, intent(in), optional :: reading, passed, pacify
    type(particle_set_t), pointer :: pset_ptr
    if (present (passed)) then
       if (.not. passed)  return
    end if
    if (eio%writing) then
       pset_ptr => event%get_particle_set_ptr ()
       call lcio_event_init (eio%lcio_event, &
             proc_id = eio%proc_num_id (i_prc), &
             event_id = event%get_index ())
       call lcio_event_from_particle_set (eio%lcio_event, pset_ptr)
       call lcio_event_set_scale (eio%lcio_event, event%get_fac_scale ())
       call lcio_event_set_alpha_qcd (eio%lcio_event, event%get_alpha_s ())
       call lcio_event_write (eio%lcio_writer, eio%lcio_event)
       call lcio_event_final (eio%lcio_event)
    else
       call eio%write ()
       call msg_fatal ("LCIO file is not open for writing")
    end if
  end subroutine eio_lcio_output

\eatline
\nwendcode{}\nwbegindocs{1298}Input an event.
\nwenddocs{}\nwbegincode{1299}\moddef{EIO LCIO: eio lcio: TBP}\plusendmoddef
  procedure :: input_i_prc => eio_lcio_input_i_prc
  procedure :: input_event => eio_lcio_input_event
\nwendcode{}\nwbegincode{1300}\moddef{EIO LCIO: procedures}\plusendmoddef
  subroutine eio_lcio_input_i_prc (eio, i_prc, iostat)
    class(eio_lcio_t), intent(inout) :: eio
    integer, intent(out) :: i_prc
    integer, intent(out) :: iostat
    logical :: ok
    integer :: i, proc_num_id
    iostat = 0
    call lcio_read_event (eio%lcio_reader, eio%lcio_event, ok)
    if (.not. ok) then 
       iostat = -1 
       return
    end if
    proc_num_id = lcio_event_get_process_id (eio%lcio_event)    
    i_prc = 0
    FIND_I_PRC: do i = 1, size (eio%proc_num_id)
       if (eio%proc_num_id(i) == proc_num_id) then
          i_prc = i
          exit FIND_I_PRC
       end if
    end do FIND_I_PRC
    if (i_prc == 0)  call err_index
  contains
    subroutine err_index
      call msg_error ("LCIO: reading events: undefined process ID " &
           // char (str (proc_num_id)) // ", aborting read")
      iostat = 1
    end subroutine err_index
  end subroutine eio_lcio_input_i_prc

  subroutine eio_lcio_input_event (eio, event, iostat)
    class(eio_lcio_t), intent(inout) :: eio
    class(generic_event_t), intent(inout), target :: event
    integer, intent(out) :: iostat
    type(particle_set_t), pointer :: pset
    logical :: ok    
    iostat = 0
    call event%reset ()
    call event%select (1, 1, 1)
    call lcio_to_event (event, eio%lcio_event, eio%fallback_model, &
         recover_beams = eio%recover_beams, &
         use_alpha_s = eio%use_alpha_s_from_file, &
         use_scale = eio%use_scale_from_file) 
    call lcio_event_final (eio%lcio_event)
  end subroutine eio_lcio_input_event

\eatline
\eatline
\nwendcode{}\nwbegindocs{1301}\nwdocspar
\subsection{Unit tests}
\nwenddocs{}\nwbegincode{1302}\moddef{EIO LCIO: public}\plusendmoddef
  public :: eio_lcio_test
\nwendcode{}\nwbegincode{1303}\moddef{EIO LCIO: tests}\endmoddef
  subroutine eio_lcio_test (u, results)
    integer, intent(in) :: u
    type(test_results_t), intent(inout) :: results
  \LA{}EIO LCIO: execute tests\RA{}
  end subroutine eio_lcio_test
  
\eatline
\nwendcode{}\nwbegindocs{1304}\nwdocspar
\subsubsection{Test I/O methods}
We test the implementation of all I/O methods.
\nwenddocs{}\nwbegincode{1305}\moddef{EIO LCIO: execute tests}\endmoddef
  call test (eio_lcio_1, "eio_lcio_1", &
       "write event contents", &
       u, results)
\nwendcode{}\nwbegincode{1306}\moddef{EIO LCIO: tests}\plusendmoddef
  subroutine eio_lcio_1 (u)
    integer, intent(in) :: u
    class(generic_event_t), pointer :: event
    type(event_sample_data_t) :: data
    class(eio_t), allocatable :: eio
    type(particle_set_t), pointer :: pset_ptr    
    type(string_t) :: sample
    integer :: u_file, iostat
    character(215) :: buffer

    write (u, "(A)")  "* Test output: eio_lcio_1"
    write (u, "(A)")  "*   Purpose: write a LCIO file"
    write (u, "(A)")

    write (u, "(A)")  "* Initialize test process"
 
    call eio_prepare_test (event)

    call data%init (1)
    data%n_beam = 2
    data%unweighted = .true.
    data%norm_mode = NORM_UNIT
    data%pdg_beam = 25
    data%energy_beam = 500
    data%proc_num_id = [42]
    data%cross_section(1) = 100
    data%error(1) = 1
    data%total_cross_section = sum (data%cross_section)

    write (u, "(A)")
    write (u, "(A)")  "* Generate and write an event"
    write (u, "(A)")
 
    sample = "eio_lcio_1"
 
    allocate (eio_lcio_t :: eio)
    select type (eio)
    type is (eio_lcio_t)
       call eio%set_parameters ()
    end select

    call eio%init_out (sample, data)
    
    call event%generate (1, [0._default, 0._default])
    call event%pacify_particle_set ()
    
    call eio%output (event, i_prc = 1)    
    call eio%write (u)
    call eio%final ()

    write (u, "(A)")
    write (u, "(A)")  "* Reset data"
    write (u, "(A)")
 
    deallocate (eio)
    allocate (eio_lcio_t :: eio)
    
    select type (eio)
    type is (eio_lcio_t)
       call eio%set_parameters ()
    end select
    call eio%write (u)

    write (u, "(A)") 
    write (u, "(A)")  "* Write LCIO file contents to ASCII file"
    write (u, "(A)")

    select type (eio)    
    type is (eio_lcio_t)       
       call lcio_event_init (eio%lcio_event, &
            proc_id = 42, &
            event_id = event%get_index ())
       pset_ptr => event%get_particle_set_ptr ()
       call lcio_event_from_particle_set &
            (eio%lcio_event,  pset_ptr)
       call write_lcio_event (eio%lcio_event, var_str ("test_file.slcio"))
       call lcio_event_final (eio%lcio_event)       
    end select
    
    write (u, "(A)") 
    write (u, "(A)")  "* Read in ASCII contents of LCIO file"
    write (u, "(A)")    
    
    u_file = free_unit ()
    open (u_file, file = "test_file.slcio", &
         action = "read", status = "old")
    do
       read (u_file, "(A)", iostat = iostat)  buffer
       if (iostat /= 0)  exit
       if (trim (buffer) == "")  cycle
       if (buffer(1:12) == " - timestamp")  cycle
       if (buffer(1:6) == " date:")  cycle       
       write (u, "(A)") trim (buffer)
    end do
    close (u_file)    
    
    write (u, "(A)")
    write (u, "(A)")  "* Cleanup"
 
    call eio_cleanup_test (event)

    write (u, "(A)")
    write (u, "(A)")  "* Test output end: eio_lcio_1"

  end subroutine eio_lcio_1
  
\eatline
\nwendcode{}\nwbegindocs{1307}Test also the reading of LCIO events.
\nwenddocs{}\nwbegincode{1308}\moddef{EIO LCIO: execute tests}\plusendmoddef
  call test (eio_lcio_2, "eio_lcio_2", &
       "read event contents", &
       u, results)
\nwendcode{}\nwbegincode{1309}\moddef{EIO LCIO: tests}\plusendmoddef
  subroutine eio_lcio_2 (u)
    integer, intent(in) :: u
    class(model_data_t), pointer :: fallback_model
    class(generic_event_t), pointer :: event
    type(event_sample_data_t) :: data
    class(eio_t), allocatable :: eio
    type(string_t) :: sample
    integer :: u_file, iostat, i_prc

    write (u, "(A)")  "* Test output: eio_lcio_2"
    write (u, "(A)")  "*   Purpose: read a LCIO event"
    write (u, "(A)")
    
    write (u, "(A)")  "* Initialize test process" 
    
    call eio_prepare_fallback_model (fallback_model)
    call eio_prepare_test (event)
    
    call data%init (1)
    data%n_beam = 2
    data%unweighted = .true.
    data%norm_mode = NORM_UNIT
    data%pdg_beam = 25
    data%energy_beam = 500
    data%proc_num_id = [42]
    data%cross_section(1) = 100
    data%error(1) = 1
    data%total_cross_section = sum (data%cross_section)    

    write (u, "(A)")
    write (u, "(A)")  "* Generate and write an event"
    write (u, "(A)")
 
    sample = "eio_lcio_2"
 
    allocate (eio_lcio_t :: eio)
    select type (eio)
    type is (eio_lcio_t)
       call eio%set_parameters (recover_beams = .false.)
    end select            
    call eio%set_fallback_model (fallback_model)
    
    call eio%init_out (sample, data)
    call event%generate (1, [0._default, 0._default])
    call event%evaluate_expressions ()
    call event%pacify_particle_set ()

    call eio%output (event, i_prc = 1)
    call eio%write (u)
    call eio%final ()
    deallocate (eio)
    
    write (u, "(A)")    
    write (u, "(A)")  "* Initialize"
    write (u, "(A)")
             
    allocate (eio_lcio_t :: eio)
    select type (eio)
    type is (eio_lcio_t)
       call eio%set_parameters (recover_beams = .false.)
    end select
    call eio%set_fallback_model (fallback_model)
    
    call data%init (1)
    data%n_beam = 2
    data%unweighted = .true.
    data%norm_mode = NORM_UNIT
    data%pdg_beam = 25
    data%energy_beam = 500
    data%proc_num_id = [42]
    call data%write (u)
    write (u, *)
    
    write (u, "(A)")  "* Initialize"
    write (u, "(A)")
    
    call eio%init_in (sample, data)
    call eio%write (u)        

    write (u, "(A)")
    write (u, "(A)")  "* Read event"
    write (u, "(A)")
 
    call eio%input_i_prc (i_prc, iostat)

    select type (eio)
    type is (eio_lcio_t)
       write (u, "(A,I0,A,I0)")  "Found process #", i_prc, &
            " with ID = ", eio%proc_num_id(i_prc)
    end select
    
    call eio%input_event (event, iostat)    
    call event%write (u)
    
    write (u, "(A)")
    write (u, "(A)")  "* Read closing"
    write (u, "(A)")
    
    call eio%input_i_prc (i_prc, iostat)
    write (u, "(A,I0)")  "iostat = ", iostat

    write (u, "(A)")
    write (u, "(A)")  "* Cleanup"

    call eio%final ()

    call eio_cleanup_test (event)
    call eio_cleanup_fallback_model (fallback_model)

    write (u, "(A)")
    write (u, "(A)")  "* Test output end: eio_lcio_2"
    
  end subroutine eio_lcio_2
  
\eatline
\nwendcode{}\nwfilename{../../../src/phase_space/phase_space.nw}\nwbegindocs{0}%% -*- ess-noweb-default-code-mode: f90-mode; noweb-default-code-mode: f90-mode; -*- 
% WHIZARD code as NOWEB source: phase space
%%%%%%%%%%%%%%%%%%%%%%%%%%%%%%%%%%%%%%%%%%%%%%%%%%%%%%%%%%%%%%%%%%%%%%%%
\chapter{Phase Space}

The abstract representation of a type that parameterizes phase space,
with methods for construction and evaluation.
\begin{description}
\item[phs\_base]
  Abstract phase-space representation.
\end{description}

A simple implementation:
\begin{description}
\item[phs\_single]
  Parameterize the phase space of a single particle, i.e., the solid
  angle.  This is useful only for very restricted problems, but it
  avoids the complexity of a generic approach in those trivial cases.
\end{description}

The standard implementation is called \emph{wood} phase space.  It
consists of several auxiliary modules and the actual implementation
module. 
\begin{description}
\item[mappings] 
  Generate invariant masses and decay angles from given
  random numbers (or the inverse operation).  Each mapping pertains to a
  particular node in a phase-space tree.  Different mappings account for
  uniform distributions, resonances, zero-mass behavior, and so on.
\item[phs\_trees] 
  Phase space parameterizations for scattering
  processes are defined recursively as if there was an initial particle
  decaying.  This module sets up a representation in terms of abstract
  trees, where each node gets a unique binary number.  Each tree is
  stored as an array of branches, where integers indicate the
  connections.  This emulates pointers in a transparent way.  Real
  pointers would also be possible, but seem to be less efficient for
  this particular case.
\item[phs\_forests] 
  The type defined by this module collects the
  decay trees corresponding to a given process and the applicable
  mappings.  To set this up, a file is read which is either written by
  the user or by the \textbf{cascades} module functions.  The module
  also contains the routines that evaluate phase space, i.e., generate
  momenta from random numbers and back.
\item[cascades] 
  This module is a pseudo Feynman diagram generator with the
  particular purpose of finding the phase space parameterizations best
  suited for a given process.  It uses a model file to set up the
  possible vertices, generates all possible diagrams, identifies
  resonances and singularities, and simplifies the list by merging
  equivalent diagrams and dropping irrelevant ones.  This process can be
  controlled at several points by user-defined parameters.  Note that it
  depends on the particular values of particle masses, so it cannot be
  done before reading the input file.
\item[phs\_wood]
  Make the functionality available in form of an implementation of the
  abstract phase-space type.
\end{description}

\clearpage
%%%%%%%%%%%%%%%%%%%%%%%%%%%%%%%%%%%%%%%%%%%%%%%%%%%%%%%%%%%%%%%%%%%%%%%%
\section{Abstract phase-space module}

In this module we define an abstract base type (and a trivial test
implementation) for multi-channel phase-space parameterizations.
\nwenddocs{}\nwbegincode{1}\moddef{\code{}phs{\char95}base.f90\edoc{}}\endmoddef
\LA{}File header\RA{}

module phs_base

\LA{}Use kinds\RA{}
\LA{}Use strings\RA{}
  use io_units
  use constants, only: TWOPI, TWOPI4
  use format_defs, only: FMT_19
  use unit_tests
  use diagnostics
  use md5
  use physics_defs
  use lorentz
  use model_data
  use flavors
  use sf_mappings
  use sf_base
  use process_constants

\LA{}Standard module head\RA{}

\LA{}PHS base: public\RA{}

\LA{}PHS base: types\RA{}

\LA{}PHS base: interfaces\RA{}

\LA{}PHS base: test types\RA{}

contains

\LA{}PHS base: procedures\RA{}

\LA{}PHS base: tests\RA{}

end module phs_base
\eatline
\nwendcode{}\nwbegindocs{2}\nwdocspar
\subsection{Phase-space channels}
The kinematics configuration may generate multiple parameterizations of phase
space.  Some of those have specific properties, such as a resonance in the s
channel.

\subsubsection{Channel properties}
This is the abstract type for the channel properties.  We need them as
a data transfer container, so everything is public and transparent.
\nwenddocs{}\nwbegincode{3}\moddef{PHS base: public}\endmoddef
  public :: channel_prop_t
\nwendcode{}\nwbegincode{4}\moddef{PHS base: types}\endmoddef
  type, abstract :: channel_prop_t
   contains
     procedure (channel_prop_to_string), deferred :: to_string
     generic :: operator (==) => is_equal
     procedure (channel_eq), deferred :: is_equal
  end type channel_prop_t

\eatline
\nwendcode{}\nwbegincode{5}\moddef{PHS base: interfaces}\endmoddef
  abstract interface
     function channel_prop_to_string (object) result (string)
       import
       class(channel_prop_t), intent(in) :: object
       type(string_t) :: string
     end function channel_prop_to_string
  end interface

\eatline
\nwendcode{}\nwbegincode{6}\moddef{PHS base: interfaces}\plusendmoddef
  abstract interface
     function channel_eq (prop1, prop2) result (flag)
       import
       class(channel_prop_t), intent(in) :: prop1, prop2
       logical :: flag
     end function channel_eq
  end interface

\eatline
\nwendcode{}\nwbegindocs{7}\nwdocspar
Here is a resonance as a channel property.  Mass and width are stored
here in physical units.
\nwenddocs{}\nwbegincode{8}\moddef{PHS base: public}\plusendmoddef
  public :: resonance_t
\nwendcode{}\nwbegincode{9}\moddef{PHS base: types}\plusendmoddef
  type, extends (channel_prop_t) :: resonance_t
     real(default) :: mass = 0
     real(default) :: width = 0
   contains
     procedure :: to_string => resonance_to_string
     procedure :: is_equal => resonance_is_equal
  end type resonance_t
  
\eatline
\nwendcode{}\nwbegindocs{10}Print mass and width.
\nwenddocs{}\nwbegincode{11}\moddef{PHS base: procedures}\endmoddef
  function resonance_to_string (object) result (string)
    class(resonance_t), intent(in) :: object
    type(string_t) :: string
    character(32) :: buffer
    string = "resonant: m ="
    write (buffer, "(" // FMT_19 // ")")  object%mass
    string = string // trim (buffer) // " GeV, w ="
    write (buffer, "(" // FMT_19 // ")")  object%width
    string = string // trim (buffer) // " GeV"
  end function resonance_to_string

\eatline
\nwendcode{}\nwbegindocs{12}Equality.
\nwenddocs{}\nwbegincode{13}\moddef{PHS base: procedures}\plusendmoddef
  function resonance_is_equal (prop1, prop2) result (flag)
    class(resonance_t), intent(in) :: prop1
    class(channel_prop_t), intent(in) :: prop2
    logical :: flag
    select type (prop2)
    type is (resonance_t)
       flag = prop1%mass == prop2%mass .and. prop1%width == prop2%width
    class default
       flag = .false.
    end select
  end function resonance_is_equal
  
\eatline
\nwendcode{}\nwbegindocs{14}\nwdocspar
This is the limiting case of a resonance, namely an on-shell particle.
We just store the mass in physical units.
\nwenddocs{}\nwbegincode{15}\moddef{PHS base: public}\plusendmoddef
  public :: on_shell_t
\nwendcode{}\nwbegincode{16}\moddef{PHS base: types}\plusendmoddef
  type, extends (channel_prop_t) :: on_shell_t
     real(default) :: mass = 0
   contains
     procedure :: to_string => on_shell_to_string
     procedure :: is_equal => on_shell_is_equal
  end type on_shell_t
  
\eatline
\nwendcode{}\nwbegindocs{17}Print mass and width.
\nwenddocs{}\nwbegincode{18}\moddef{PHS base: procedures}\plusendmoddef
  function on_shell_to_string (object) result (string)
    class(on_shell_t), intent(in) :: object
    type(string_t) :: string
    character(32) :: buffer
    string = "on shell: m ="
    write (buffer, "(" // FMT_19 // ")")  object%mass
    string = string // trim (buffer) // " GeV"
  end function on_shell_to_string

\eatline
\nwendcode{}\nwbegindocs{19}Equality.
\nwenddocs{}\nwbegincode{20}\moddef{PHS base: procedures}\plusendmoddef
  function on_shell_is_equal (prop1, prop2) result (flag)
    class(on_shell_t), intent(in) :: prop1
    class(channel_prop_t), intent(in) :: prop2
    logical :: flag
    select type (prop2)
    type is (on_shell_t)
       flag = prop1%mass == prop2%mass
    class default
       flag = .false.
    end select
  end function on_shell_is_equal
  
\eatline
\nwendcode{}\nwbegindocs{21}\nwdocspar
\subsubsection{Channel equivalences}
This type describes an equivalence.  The current channel is equivalent
to channel {\tt{}c}.  The equivalence involves a permutation {\tt{}perm} of
integration dimensions and, within each integration dimension, a
mapping {\tt{}mode}.
\nwenddocs{}\nwbegincode{22}\moddef{PHS base: types}\plusendmoddef
  type :: phs_equivalence_t
     integer :: c = 0
     integer, dimension(:), allocatable :: perm
     integer, dimension(:), allocatable :: mode
   contains
   \LA{}PHS base: phs equivalence: TBP\RA{}
  end type phs_equivalence_t
  
\eatline
\nwendcode{}\nwbegindocs{23}\nwdocspar
The mapping modes are
\nwenddocs{}\nwbegincode{24}\moddef{PHS base: types}\plusendmoddef
  integer, parameter, public :: &
       EQ_IDENTITY = 0, EQ_INVERT = 1, EQ_SYMMETRIC = 2, EQ_INVARIANT = 3

\eatline
\nwendcode{}\nwbegindocs{25}In particular, if a channel is equivalent to itself in the
{\tt{}EQ{\char95}SYMMETRIC} mode, the integrand can be assumed to be symmetric
w.r.t.\ a reflection $x\to 1 - x$ of the correponding integration variable.

These are the associated tags, for output:
\nwenddocs{}\nwbegincode{26}\moddef{PHS base: types}\plusendmoddef
  character, dimension(0:3), parameter :: TAG = ["+", "-", ":", "x"]

\eatline
\nwendcode{}\nwbegindocs{27}Write an equivalence.
\nwenddocs{}\nwbegincode{28}\moddef{PHS base: phs equivalence: TBP}\endmoddef
  procedure :: write => phs_equivalence_write
\nwendcode{}\nwbegincode{29}\moddef{PHS base: procedures}\plusendmoddef
  subroutine phs_equivalence_write (object, unit)
    class(phs_equivalence_t), intent(in) :: object
    integer, intent(in), optional :: unit
    integer :: u, j
    u = given_output_unit (unit)
    write (u, "(5x,'=',1x,I0,1x)", advance = "no")  object%c
    if (allocated (object%perm)) then
       write (u, "(A)", advance = "no")  "("
       do j = 1, size (object%perm)
          if (j > 1)  write (u, "(1x)", advance = "no")
          write (u, "(I0,A1)", advance = "no") &
               object%perm(j), TAG(object%mode(j))
       end do
       write (u, "(A)")  ")"
    else
       write (u, "(A)")
    end if
  end subroutine phs_equivalence_write
  
\eatline
\nwendcode{}\nwbegindocs{30}Initialize an equivalence.  This allocates the {\tt{}perm} and {\tt{}mode}
arrays with equal size.
\nwenddocs{}\nwbegincode{31}\moddef{PHS base: phs equivalence: TBP}\plusendmoddef
  procedure :: init => phs_equivalence_init
\nwendcode{}\nwbegincode{32}\moddef{PHS base: procedures}\plusendmoddef
  subroutine phs_equivalence_init (eq, n_dim)
    class(phs_equivalence_t), intent(out) :: eq
    integer, intent(in) :: n_dim
    allocate (eq%perm (n_dim), source = 0)
    allocate (eq%mode (n_dim), source = EQ_IDENTITY)
  end subroutine phs_equivalence_init
  
\eatline
\nwendcode{}\nwbegindocs{33}\nwdocspar
\subsubsection{Channel objects}
The channel entry holds (optionally) specific properties.

{\tt{}sf{\char95}channel} is the structure-function channel that corresponds to this
phase-space channel.  The structure-function channel may be set up with a
specific mapping that depends on the phase-space channel properties.  (The
default setting is to leave the properties empty.)
\nwenddocs{}\nwbegincode{34}\moddef{PHS base: public}\plusendmoddef
  public :: phs_channel_t
\nwendcode{}\nwbegincode{35}\moddef{PHS base: types}\plusendmoddef
  type :: phs_channel_t
     class(channel_prop_t), allocatable :: prop
     integer :: sf_channel = 1
     type(phs_equivalence_t), dimension(:), allocatable :: eq
   contains
   \LA{}PHS base: phs channel: TBP\RA{}
  end type phs_channel_t
  
\eatline
\nwendcode{}\nwbegindocs{36}Output.
\nwenddocs{}\nwbegincode{37}\moddef{PHS base: phs channel: TBP}\endmoddef
  procedure :: write => phs_channel_write
\nwendcode{}\nwbegincode{38}\moddef{PHS base: procedures}\plusendmoddef
  subroutine phs_channel_write (object, unit)
    class(phs_channel_t), intent(in) :: object
    integer, intent(in), optional :: unit
    integer :: u, j
    u = given_output_unit (unit)
    write (u, "(1x,I0)", advance="no") object%sf_channel
    if (allocated (object%prop)) then
       write (u, "(1x,A)")  char (object%prop%to_string ())
    else
       write (u, *)
    end if
    if (allocated (object%eq)) then
       do j = 1, size (object%eq)
          call object%eq(j)%write (u)
       end do
    end if
  end subroutine phs_channel_write
    
\eatline
\nwendcode{}\nwbegindocs{39}Identify the channel with an s-channel resonance.
\nwenddocs{}\nwbegincode{40}\moddef{PHS base: phs channel: TBP}\plusendmoddef
  procedure :: set_resonant => channel_set_resonant
\nwendcode{}\nwbegincode{41}\moddef{PHS base: procedures}\plusendmoddef
  subroutine channel_set_resonant (channel, mass, width)
    class(phs_channel_t), intent(inout) :: channel
    real(default), intent(in) :: mass, width
    allocate (resonance_t :: channel%prop)
    select type (prop => channel%prop)
    type is (resonance_t)
       prop%mass = mass
       prop%width = width
    end select
  end subroutine channel_set_resonant
  
\eatline
\nwendcode{}\nwbegindocs{42}Identify the channel with an on-shell particle.
\nwenddocs{}\nwbegincode{43}\moddef{PHS base: phs channel: TBP}\plusendmoddef
  procedure :: set_on_shell => channel_set_on_shell
\nwendcode{}\nwbegincode{44}\moddef{PHS base: procedures}\plusendmoddef
  subroutine channel_set_on_shell (channel, mass)
    class(phs_channel_t), intent(inout) :: channel
    real(default), intent(in) :: mass
    allocate (on_shell_t :: channel%prop)
    select type (prop => channel%prop)
    type is (on_shell_t)
       prop%mass = mass
    end select
  end subroutine channel_set_on_shell
  
\eatline
\nwendcode{}\nwbegindocs{45}\nwdocspar
\subsection{Property collection}
We can set up a list of all distinct channel properties for a given
set of channels.
\nwenddocs{}\nwbegincode{46}\moddef{PHS base: public}\plusendmoddef
  public :: phs_channel_collection_t
\nwendcode{}\nwbegincode{47}\moddef{PHS base: types}\plusendmoddef
  type :: prop_entry_t
     integer :: i = 0
     class(channel_prop_t), allocatable :: prop
     type(prop_entry_t), pointer :: next => null ()
  end type prop_entry_t
  
  type :: phs_channel_collection_t
     integer :: n = 0
     type(prop_entry_t), pointer :: first => null ()
   contains
   \LA{}PHS base: phs channel collection: TBP\RA{}
  end type phs_channel_collection_t

\eatline
\eatline
\nwendcode{}\nwbegindocs{48}Finalizer for the list.
\nwenddocs{}\nwbegincode{49}\moddef{PHS base: phs channel collection: TBP}\endmoddef
  procedure :: final => phs_channel_collection_final
\nwendcode{}\nwbegincode{50}\moddef{PHS base: procedures}\plusendmoddef
  subroutine phs_channel_collection_final (object)
    class(phs_channel_collection_t), intent(inout) :: object
    type(prop_entry_t), pointer :: entry
    do while (associated (object%first))
       entry => object%first
       object%first => entry%next
       deallocate (entry)
    end do
  end subroutine phs_channel_collection_final

\eatline
\nwendcode{}\nwbegindocs{51}Output.

Note: eliminating the {\tt{}string} auxiliary triggers an ICE in gfortran 4.7.2.
\nwenddocs{}\nwbegincode{52}\moddef{PHS base: phs channel collection: TBP}\plusendmoddef
  procedure :: write => phs_channel_collection_write
\nwendcode{}\nwbegincode{53}\moddef{PHS base: procedures}\plusendmoddef
  subroutine phs_channel_collection_write (object, unit)
    class(phs_channel_collection_t), intent(in) :: object
    integer, intent(in), optional :: unit
    type(prop_entry_t), pointer :: entry
    type(string_t) :: string
    integer :: u
    u = given_output_unit (unit)
    entry => object%first
    do while (associated (entry))
       if (allocated (entry%prop)) then
          string = entry%prop%to_string ()
          write (u, "(1x,I0,1x,A)")  entry%i, char (string)
       else
          write (u, "(1x,I0)")  entry%i
       end if
       entry => entry%next
    end do
  end subroutine phs_channel_collection_write
  
\eatline
\nwendcode{}\nwbegindocs{54}Push a new property to the stack if it is not yet included.
Simultaneously, set the {\tt{}sf{\char95}channel} entry in the phase-space
channel object to the index of the matching entry, or the new entry if
there was no match.
\nwenddocs{}\nwbegincode{55}\moddef{PHS base: phs channel collection: TBP}\plusendmoddef
  procedure :: push => phs_channel_collection_push
\nwendcode{}\nwbegincode{56}\moddef{PHS base: procedures}\plusendmoddef
  subroutine phs_channel_collection_push (coll, channel)
    class(phs_channel_collection_t), intent(inout) :: coll
    type(phs_channel_t), intent(inout) :: channel
    type(prop_entry_t), pointer :: entry, new
    if (associated (coll%first)) then
       entry => coll%first
       do
          if (allocated (entry%prop)) then
             if (allocated (channel%prop)) then
                if (entry%prop == channel%prop) then
                   channel%sf_channel = entry%i
                   return
                end if
             end if
          else if (.not. allocated (channel%prop)) then
             channel%sf_channel = entry%i
             return
          end if
          if (associated (entry%next)) then
             entry => entry%next
          else
             exit
          end if
       end do
       allocate (new)
       entry%next => new
    else
       allocate (new)
       coll%first => new
    end if
    coll%n = coll%n + 1
    new%i = coll%n
    channel%sf_channel = new%i
    if (allocated (channel%prop)) then
       allocate (new%prop, source = channel%prop)
    end if
  end subroutine phs_channel_collection_push
    
\eatline
\nwendcode{}\nwbegindocs{57}Return the number of collected distinct channels.
\nwenddocs{}\nwbegincode{58}\moddef{PHS base: phs channel collection: TBP}\plusendmoddef
  procedure :: get_n => phs_channel_collection_get_n
\nwendcode{}\nwbegincode{59}\moddef{PHS base: procedures}\plusendmoddef
  function phs_channel_collection_get_n (coll) result (n)
    class(phs_channel_collection_t), intent(in) :: coll
    integer :: n
    n = coll%n
  end function phs_channel_collection_get_n
  
\eatline
\nwendcode{}\nwbegindocs{60}Return a specific channel (property object).
\nwenddocs{}\nwbegincode{61}\moddef{PHS base: phs channel collection: TBP}\plusendmoddef
  procedure :: get_entry => phs_channel_collection_get_entry
\nwendcode{}\nwbegincode{62}\moddef{PHS base: procedures}\plusendmoddef
  subroutine phs_channel_collection_get_entry (coll, i, prop)
    class(phs_channel_collection_t), intent(in) :: coll
    integer, intent(in) :: i
    class(channel_prop_t), intent(out), allocatable :: prop
    type(prop_entry_t), pointer :: entry
    integer :: k
    if (i > 0 .and. i <= coll%n) then
       entry => coll%first
       do k = 2, i
          entry => entry%next
       end do
       if (allocated (entry%prop)) then
          if (allocated (prop))  deallocate (prop)
          allocate (prop, source = entry%prop)
       end if
    else
       call msg_bug ("PHS channel collection: get entry: illegal index")
    end if
  end subroutine phs_channel_collection_get_entry
    
\eatline
\nwendcode{}\nwbegindocs{63}\nwdocspar
\subsection{Kinematics configuration}
Here, we store the universal information that is specifically relevant
for phase-space generation.  It is a subset of the process data,
supplemented by basic information on phase-space parameterization
channels.

A concrete implementation will contain more data, that describe the
phase space in detail.

MD5 sums: the phase space setup depends on the process, it depends on
the model parameters (the masses, that is), and on the configuration
parameters.  (It doesn't depend on the QCD setup.)
\nwenddocs{}\nwbegincode{64}\moddef{PHS base: public}\plusendmoddef
  public :: phs_config_t
\nwendcode{}\nwbegincode{65}\moddef{PHS base: types}\plusendmoddef
  type, abstract :: phs_config_t
     ! private
     type(string_t) :: id
     integer :: n_in = 0
     integer :: n_out = 0
     integer :: n_tot = 0
     integer :: n_state = 0
     integer :: n_par = 0
     integer :: n_channel = 0
     real(default) :: sqrts = 0
     logical :: sqrts_fixed = .true.
     logical :: cm_frame = .true.
     logical :: azimuthal_dependence = .false.
     integer, dimension(:), allocatable :: dim_flat
     logical :: provides_equivalences = .false.
     logical :: provides_chains = .false.
     logical :: vis_channels = .false.
     integer, dimension(:), allocatable :: chain
     class(model_data_t), pointer :: model => null ()
     type(flavor_t), dimension(:,:), allocatable :: flv
     type(phs_channel_t), dimension(:), allocatable :: channel
     character(32) :: md5sum_process = ""
     character(32) :: md5sum_model_par = ""
     character(32) :: md5sum_phs_config = ""
     type(string_t) :: nlo_type
   contains
   \LA{}PHS base: phs config: TBP\RA{}
  end type phs_config_t
  
\eatline
\nwendcode{}\nwbegindocs{66}Finalizer, deferred.
\nwenddocs{}\nwbegincode{67}\moddef{PHS base: phs config: TBP}\endmoddef
  procedure (phs_config_final), deferred :: final
\nwendcode{}\nwbegincode{68}\moddef{PHS base: interfaces}\plusendmoddef
  abstract interface
     subroutine phs_config_final (object)
       import
       class(phs_config_t), intent(inout) :: object
     end subroutine phs_config_final
  end interface

\eatline
\nwendcode{}\nwbegindocs{69}Output.  We provide an implementation for the output of the base-type
contents and an interface for the actual write method.
\nwenddocs{}\nwbegincode{70}\moddef{PHS base: phs config: TBP}\plusendmoddef
  procedure (phs_config_write), deferred :: write
  procedure :: base_write => phs_config_write
\nwendcode{}\nwbegincode{71}\moddef{PHS base: procedures}\plusendmoddef
  subroutine phs_config_write (object, unit)
    class(phs_config_t), intent(in) :: object
    integer, intent(in), optional :: unit
    integer :: u, i, j
    integer :: n_tot_flv
    n_tot_flv = object%n_tot
    u = given_output_unit (unit)
    write (u, "(3x,A,A,A)") "ID        = '", char (object%id), "'"
    write (u, "(3x,A,I0)")  "n_in      = ", object%n_in
    write (u, "(3x,A,I0)")  "n_out     = ", object%n_out
    write (u, "(3x,A,I0)")  "n_tot     = ", object%n_tot
    write (u, "(3x,A,I0)")  "n_state   = ", object%n_state
    write (u, "(3x,A,I0)")  "n_par     = ", object%n_par
    write (u, "(3x,A,I0)")  "n_channel = ", object%n_channel
    write (u, "(3x,A," // FMT_19 // ")")  "sqrts     = ", object%sqrts
    write (u, "(3x,A,L1)")  "s_fixed   = ", object%sqrts_fixed
    write (u, "(3x,A,L1)")  "cm_frame  = ", object%cm_frame
    write (u, "(3x,A,L1)")  "azim.dep. = ", object%azimuthal_dependence
    if (allocated (object%dim_flat)) then
       write (u, "(3x,A,I0)")  "flat dim. = ", object%dim_flat
    end if
    write (u, "(1x,A)")  "Flavor combinations:"
    do i = 1, object%n_state
       write (u, "(3x,I0,':')", advance="no")  i
!       do j = 1, object%n_tot
       do j = 1, n_tot_flv
          write (u, "(1x,A)", advance="no")  char (object%flv(j,i)%get_name ())
       end do
       write (u, *)
    end do
    if (allocated (object%channel)) then
       write (u, "(1x,A)")  "Phase-space / structure-function channels:"
       do i = 1, object%n_channel
          write (u, "(3x,I0,':')", advance="no") i
          call object%channel(i)%write (u)
       end do
    end if
    if (object%md5sum_process /= "") then
       write (u, "(3x,A,A,A)") "MD5 sum (process)    = '", &
            object%md5sum_process, "'"
    end if
    if (object%md5sum_model_par /= "") then
       write (u, "(3x,A,A,A)") "MD5 sum (model par)  = '", &
            object%md5sum_model_par, "'"
    end if
    if (object%md5sum_phs_config /= "") then
       write (u, "(3x,A,A,A)") "MD5 sum (phs config) = '", &
            object%md5sum_phs_config, "'"
    end if
  end subroutine phs_config_write

\eatline
\nwendcode{}\nwbegindocs{72}Similarly, a basic initializer and an interface.  The model pointer is taken
as an argument; we may verify that this has the expected model name.

The intent is {\tt{}inout}.  We want to be able to set parameters in advance.
\nwenddocs{}\nwbegincode{73}\moddef{PHS base: phs config: TBP}\plusendmoddef
  procedure :: init => phs_config_init
\nwendcode{}\nwbegincode{74}\moddef{PHS base: procedures}\plusendmoddef
  subroutine phs_config_init (phs_config, data, model)
    class(phs_config_t), intent(inout) :: phs_config
    type(process_constants_t), intent(in) :: data
    class(model_data_t), intent(in), target :: model
    integer :: i, j
    phs_config%id = data%id
    phs_config%n_in  = data%n_in
    phs_config%n_out = data%n_out
    phs_config%n_tot = data%n_in + data%n_out
    phs_config%n_state = data%n_flv
    if (data%model_name == model%get_name ()) then
       phs_config%model => model
    else
       call msg_bug ("phs_config_init: model name mismatch")
    end if
    allocate (phs_config%flv (phs_config%n_tot, phs_config%n_state))    
    do i = 1, phs_config%n_state
       do j = 1, phs_config%n_tot
          call phs_config%flv(j,i)%init (data%flv_state(j,i), &
               phs_config%model)
       end do
    end do
    phs_config%md5sum_process = data%md5sum
  end subroutine phs_config_init
    
\eatline
\nwendcode{}\nwbegindocs{75}This procedure should complete the phase-space configuration.  We
need the {\tt{}sqrts} value as overall scale, which is known only after
the beams have been defined.  The procedure should determine the number of
channels, their properties (if any), and allocate and fill the {\tt{}channel}
array accordingly.
\nwenddocs{}\nwbegincode{76}\moddef{PHS base: phs config: TBP}\plusendmoddef
  procedure (phs_config_configure), deferred :: configure
\nwendcode{}\nwbegincode{77}\moddef{PHS base: interfaces}\plusendmoddef
  abstract interface
     subroutine phs_config_configure (phs_config, sqrts, &
          sqrts_fixed, cm_frame, azimuthal_dependence, rebuild, ignore_mismatch, &
          nlo_type)
       import
       class(phs_config_t), intent(inout) :: phs_config
       real(default), intent(in) :: sqrts
       logical, intent(in), optional :: sqrts_fixed
       logical, intent(in), optional :: cm_frame
       logical, intent(in), optional :: azimuthal_dependence
       logical, intent(in), optional :: rebuild
       logical, intent(in), optional :: ignore_mismatch
       type(string_t), intent(inout), optional :: nlo_type
     end subroutine phs_config_configure
  end interface

\eatline
\nwendcode{}\nwbegindocs{78}Manually assign structure-function channel indices to the phase-space
channel objects.  (Used by a test routine.)
\nwenddocs{}\nwbegincode{79}\moddef{PHS base: phs config: TBP}\plusendmoddef
  procedure :: set_sf_channel => phs_config_set_sf_channel
\nwendcode{}\nwbegincode{80}\moddef{PHS base: procedures}\plusendmoddef
  subroutine phs_config_set_sf_channel (phs_config, sf_channel)
    class(phs_config_t), intent(inout) :: phs_config
    integer, dimension(:), intent(in) :: sf_channel
    phs_config%channel%sf_channel = sf_channel
  end subroutine phs_config_set_sf_channel
  
\eatline
\nwendcode{}\nwbegindocs{81}Collect new channels not yet in the collection from this phase-space
configuration object.  At the same time, assign structure-function channels.
\nwenddocs{}\nwbegincode{82}\moddef{PHS base: phs config: TBP}\plusendmoddef
  procedure :: collect_channels => phs_config_collect_channels
\nwendcode{}\nwbegincode{83}\moddef{PHS base: procedures}\plusendmoddef
  subroutine phs_config_collect_channels (phs_config, coll)
    class(phs_config_t), intent(inout) :: phs_config
    type(phs_channel_collection_t), intent(inout) :: coll
    integer :: c
    do c = 1, phs_config%n_channel
       call coll%push (phs_config%channel(c))
    end do
  end subroutine phs_config_collect_channels
    
\eatline
\nwendcode{}\nwbegindocs{84}Compute the MD5 sum.  We abuse the {\tt{}write} method.  In
type implementations, {\tt{}write} should only display information that is
relevant for the MD5 sum.  The data include the process MD5 sum which is taken
from the process constants, and the MD5 sum of the model parameters.  This may
change, so it is computed here.
\nwenddocs{}\nwbegincode{85}\moddef{PHS base: phs config: TBP}\plusendmoddef
  procedure :: compute_md5sum => phs_config_compute_md5sum
\nwendcode{}\nwbegincode{86}\moddef{PHS base: procedures}\plusendmoddef
  subroutine phs_config_compute_md5sum (phs_config)
    class(phs_config_t), intent(inout) :: phs_config
    integer :: u
    phs_config%md5sum_model_par = phs_config%model%get_parameters_md5sum ()
    phs_config%md5sum_phs_config = ""
    u = free_unit ()
    open (u, status = "scratch", action = "readwrite")
    call phs_config%write (u)
    rewind (u)
    phs_config%md5sum_phs_config = md5sum (u)
    close (u)
  end subroutine phs_config_compute_md5sum
    
\eatline
\nwendcode{}\nwbegindocs{87}Print an informative message after phase-space configuration.
\nwenddocs{}\nwbegincode{88}\moddef{PHS base: phs config: TBP}\plusendmoddef
  procedure (phs_startup_message), deferred :: startup_message
  procedure :: base_startup_message => phs_startup_message
\nwendcode{}\nwbegincode{89}\moddef{PHS base: procedures}\plusendmoddef
  subroutine phs_startup_message (phs_config, unit)
    class(phs_config_t), intent(in) :: phs_config
    integer, intent(in), optional :: unit
    write (msg_buffer, "(A,3(1x,I0,1x,A))") &
         "Phase space:", &
         phs_config%n_channel, "channels,", &
         phs_config%n_par, "dimensions"
    call msg_message (unit = unit)
  end subroutine phs_startup_message

\eatline
\nwendcode{}\nwbegindocs{90}This procedure should be implemented such that the phase-space
configuration object allocates a phase-space instance of matching type.
\nwenddocs{}\nwbegincode{91}\moddef{PHS base: phs config: TBP}\plusendmoddef
  procedure (phs_config_allocate_instance), nopass, deferred :: &
       allocate_instance
\nwendcode{}\nwbegincode{92}\moddef{PHS base: interfaces}\plusendmoddef
  abstract interface
     subroutine phs_config_allocate_instance (phs)
       import
       class(phs_t), intent(inout), pointer :: phs
     end subroutine phs_config_allocate_instance
  end interface

\eatline
\nwendcode{}\nwbegindocs{93}\nwdocspar
\subsection{Extract data}
Return the number of MC input parameters.
\nwenddocs{}\nwbegincode{94}\moddef{PHS base: phs config: TBP}\plusendmoddef
  procedure :: get_n_par => phs_config_get_n_par
\nwendcode{}\nwbegincode{95}\moddef{PHS base: procedures}\plusendmoddef
  function phs_config_get_n_par (phs_config) result (n)
    class(phs_config_t), intent(in) :: phs_config
    integer :: n
    n = phs_config%n_par
  end function phs_config_get_n_par
  
\eatline
\nwendcode{}\nwbegindocs{96}Return dimensions (parameter indices) for which the phase-space
dimension is flat, so integration and event generation can be simplified.
\nwenddocs{}\nwbegincode{97}\moddef{PHS base: phs config: TBP}\plusendmoddef
  procedure :: get_flat_dimensions => phs_config_get_flat_dimensions
\nwendcode{}\nwbegincode{98}\moddef{PHS base: procedures}\plusendmoddef
  function phs_config_get_flat_dimensions (phs_config) result (dim_flat)
    class(phs_config_t), intent(in) :: phs_config
    integer, dimension(:), allocatable :: dim_flat
    if (allocated (phs_config%dim_flat)) then
       allocate (dim_flat (size (phs_config%dim_flat)))
       dim_flat = phs_config%dim_flat
    else
       allocate (dim_flat (0))
    end if
  end function phs_config_get_flat_dimensions
  
\eatline
\nwendcode{}\nwbegindocs{99}Return the number of phase-space channels.
\nwenddocs{}\nwbegincode{100}\moddef{PHS base: phs config: TBP}\plusendmoddef
  procedure :: get_n_channel => phs_config_get_n_channel
\nwendcode{}\nwbegincode{101}\moddef{PHS base: procedures}\plusendmoddef
  function phs_config_get_n_channel (phs_config) result (n)
    class(phs_config_t), intent(in) :: phs_config
    integer :: n
    n = phs_config%n_channel
  end function phs_config_get_n_channel
  
\eatline
\nwendcode{}\nwbegindocs{102}Return the structure-function channel that corresponds to the
phase-space channel {\tt{}c}.  If the channel array is not allocated (which
happens if there is no structure function), return zero.
\nwenddocs{}\nwbegincode{103}\moddef{PHS base: phs config: TBP}\plusendmoddef
  procedure :: get_sf_channel => phs_config_get_sf_channel
\nwendcode{}\nwbegincode{104}\moddef{PHS base: procedures}\plusendmoddef
  function phs_config_get_sf_channel (phs_config, c) result (c_sf)
    class(phs_config_t), intent(in) :: phs_config
    integer, intent(in) :: c
    integer :: c_sf
    if (allocated (phs_config%channel)) then
       c_sf = phs_config%channel(c)%sf_channel
    else
       c_sf = 0
    end if
  end function phs_config_get_sf_channel
  
\eatline
\nwendcode{}\nwbegindocs{105}Return the mass(es) of the incoming particle(s).  We take the first flavor
combination in the array, assuming that masses must be degenerate among
flavors.
\nwenddocs{}\nwbegincode{106}\moddef{PHS base: phs config: TBP}\plusendmoddef
  procedure :: get_masses_in => phs_config_get_masses_in
\nwendcode{}\nwbegincode{107}\moddef{PHS base: procedures}\plusendmoddef
  subroutine phs_config_get_masses_in (phs_config, m)
    class(phs_config_t), intent(in) :: phs_config
    real(default), dimension(:), intent(out) :: m
    integer :: i
    do i = 1, phs_config%n_in
       m(i) = phs_config%flv(i,1)%get_mass ()
    end do
  end subroutine phs_config_get_masses_in
  
\eatline
\nwendcode{}\nwbegindocs{108}Return the MD5 sum of the configuration.
\nwenddocs{}\nwbegincode{109}\moddef{PHS base: phs config: TBP}\plusendmoddef
  procedure :: get_md5sum => phs_config_get_md5sum
\nwendcode{}\nwbegincode{110}\moddef{PHS base: procedures}\plusendmoddef
  function phs_config_get_md5sum (phs_config) result (md5sum)
    class(phs_config_t), intent(in) :: phs_config
    character(32) :: md5sum
    md5sum = phs_config%md5sum_phs_config
  end function phs_config_get_md5sum
  
\eatline
\nwendcode{}\nwbegindocs{111}\nwdocspar
\subsection{Phase-space point instance}
The {\tt{}phs{\char95}t} object holds the workspace for phase-space generation.
In the base object, we have the MC input parameters {\tt{}r} and the
Jacobian factor {\tt{}f}, for each channel, and the incoming and outgoing
momenta.

Note: The {\tt{}active{\char95}channel} array is not used yet, all elements are
initialized with {\tt{}.true.}.  It should be touched by the integrator if it
decides to drop irrelevant channels.
\nwenddocs{}\nwbegincode{112}\moddef{PHS base: public}\plusendmoddef
  public :: phs_t
\nwendcode{}\nwbegincode{113}\moddef{PHS base: types}\plusendmoddef
  type, abstract :: phs_t
     class(phs_config_t), pointer :: config => null ()
     logical :: r_defined = .false.
     integer :: selected_channel = 0
     logical, dimension(:), allocatable :: active_channel
     real(default), dimension(:,:), allocatable :: r
     real(default), dimension(:), allocatable :: f
     real(default), dimension(:), allocatable :: m_in
     real(default), dimension(:), allocatable :: m_out
     real(default) :: flux = 0
     real(default) :: volume = 0
     type(lorentz_transformation_t) :: lt_cm_to_lab
     logical :: p_defined = .false.
     real(default) :: sqrts_hat = 0
     type(vector4_t), dimension(:), allocatable :: p
     logical :: q_defined = .false.
     type(vector4_t), dimension(:), allocatable :: q
   contains
   \LA{}PHS base: phs: TBP\RA{}
  end type phs_t
     
\eatline
\nwendcode{}\nwbegindocs{114}Output.  Since phase space may get complicated, we include a
{\tt{}verbose} option for the abstract {\tt{}write} procedure.
\nwenddocs{}\nwbegincode{115}\moddef{PHS base: phs: TBP}\endmoddef
  procedure (phs_write), deferred :: write
\nwendcode{}\nwbegincode{116}\moddef{PHS base: interfaces}\plusendmoddef
  abstract interface
     subroutine phs_write (object, unit, verbose)
       import
       class(phs_t), intent(in) :: object
       integer, intent(in), optional :: unit
       logical, intent(in), optional :: verbose
     end subroutine phs_write
  end interface
  
\eatline
\nwendcode{}\nwbegindocs{117}This procedure can be called to print the contents of the base type.
\nwenddocs{}\nwbegincode{118}\moddef{PHS base: phs: TBP}\plusendmoddef
  procedure :: base_write => phs_base_write
\nwendcode{}\nwbegincode{119}\moddef{PHS base: procedures}\plusendmoddef
  subroutine phs_base_write (object, unit)
    class(phs_t), intent(in) :: object
    integer, intent(in), optional :: unit
    integer :: u, c, i
    u = given_output_unit (unit)
    write (u, "(1x,A)", advance="no")  "Partonic phase space: parameters"
    if (object%r_defined) then
       write (u, *)
    else
       write (u, "(1x,A)")  "[undefined]"
    end if
    write (u, "(3x,A,999(1x," // FMT_19 // "))") "m_in    =", object%m_in
    write (u, "(3x,A,999(1x," // FMT_19 // "))") "m_out   =", object%m_out
    write (u, "(3x,A," // FMT_19 // ")")  "Flux   = ", object%flux
    write (u, "(3x,A," // FMT_19 // ")")  "Volume = ", object%volume
    if (allocated (object%f)) then
       do c = 1, size (object%r, 2)
          write (u, "(1x,A,I0,A)", advance="no")  "Channel #", c, ":"
          if (c == object%selected_channel) then
             write (u, "(1x,A)")  "[selected]"
          else
             write (u, *)
          end if
          write (u, "(3x,A)", advance="no")  "r ="
          do i = 1, size (object%r, 1)
             write (u, "(1x,F9.7)", advance="no")  object%r(i,c)
          end do
          write (u, *)
          write (u, "(3x,A,1x,ES13.7)")  "f =", object%f(c)
       end do
    end if
    write (u, "(1x,A)")  "Partonic phase space: momenta"
    if (object%p_defined) then
       write (u, "(3x,A," // FMT_19 // ")")  "sqrts  = ", object%sqrts_hat
    end if
    write (u, "(1x,A)", advance="no")  "Incoming:"
    if (object%p_defined) then
       write (u, *)
    else
       write (u, "(1x,A)")  "[undefined]"
    end if
    if (allocated (object%p)) then
       do i = 1, size (object%p)
          call vector4_write (object%p(i), u)
       end do
    end if
    write (u, "(1x,A)", advance="no")  "Outgoing:"
    if (object%q_defined) then
       write (u, *)
    else
       write (u, "(1x,A)")  "[undefined]"
    end if
    if (allocated (object%q)) then
       do i = 1, size (object%q)
          call vector4_write (object%q(i), u)
       end do
    end if
    if (object%p_defined .and. .not. object%config%cm_frame) then
       write (u, "(1x,A)")  "Transformation c.m -> lab frame"
       call lorentz_transformation_write (object%lt_cm_to_lab, u)
    end if
  end subroutine phs_base_write
          
\eatline
\nwendcode{}\nwbegindocs{120}Finalizer.  The base type does not need it, but extensions may.
\nwenddocs{}\nwbegincode{121}\moddef{PHS base: phs: TBP}\plusendmoddef
  procedure (phs_final), deferred :: final
\nwendcode{}\nwbegincode{122}\moddef{PHS base: interfaces}\plusendmoddef
  abstract interface
     subroutine phs_final (object)
       import
       class(phs_t), intent(inout) :: object
     end subroutine phs_final
  end interface
  
\eatline
\nwendcode{}\nwbegindocs{123}Initializer.  Everything should be contained in the {\tt{}process{\char95}data}
configuration object, so we can require a universal interface.
\nwenddocs{}\nwbegincode{124}\moddef{PHS base: phs: TBP}\plusendmoddef
  procedure (phs_init), deferred :: init
\nwendcode{}\nwbegincode{125}\moddef{PHS base: interfaces}\plusendmoddef
  abstract interface
     subroutine phs_init (phs, phs_config)
       import
       class(phs_t), intent(out) :: phs
       class(phs_config_t), intent(in), target :: phs_config
     end subroutine phs_init
  end interface
  
\eatline
\nwendcode{}\nwbegindocs{126}The base version will just allocate the arrays.  It should be called
at the beginning of the implementation of {\tt{}phs{\char95}init}.
\nwenddocs{}\nwbegincode{127}\moddef{PHS base: phs: TBP}\plusendmoddef
  procedure :: base_init => phs_base_init
\nwendcode{}\nwbegincode{128}\moddef{PHS base: procedures}\plusendmoddef
  subroutine phs_base_init (phs, phs_config)
    class(phs_t), intent(out) :: phs
    class(phs_config_t), intent(in), target :: phs_config
    phs%config => phs_config
    allocate (phs%active_channel (phs%config%n_channel))
    phs%active_channel = .true.
    allocate (phs%r (phs%config%n_par, phs%config%n_channel));  phs%r = 0
    allocate (phs%f (phs%config%n_channel));                    phs%f = 0
    allocate (phs%p (phs%config%n_in))
    allocate (phs%m_in  (phs%config%n_in), &
         source = phs_config%flv(:phs_config%n_in, 1)%get_mass ())
    allocate (phs%q (phs%config%n_out))
    allocate (phs%m_out (phs%config%n_out), &
         source = phs_config%flv(phs_config%n_in+1:, 1)%get_mass ())
    call phs%compute_flux ()
  end subroutine phs_base_init

\eatline
\nwendcode{}\nwbegindocs{129}Manually select a channel.
\nwenddocs{}\nwbegincode{130}\moddef{PHS base: phs: TBP}\plusendmoddef
  procedure :: select_channel => phs_base_select_channel
\nwendcode{}\nwbegincode{131}\moddef{PHS base: procedures}\plusendmoddef
  subroutine phs_base_select_channel (phs, channel)
    class(phs_t), intent(inout) :: phs
    integer, intent(in), optional :: channel
    if (present (channel)) then
       phs%selected_channel = channel
    else
       phs%selected_channel = 0
    end if
  end subroutine phs_base_select_channel
  
\eatline
\nwendcode{}\nwbegindocs{132}Set incoming momenta.  Assume that array shapes match.  If
requested, compute the Lorentz transformation from the c.m.\ to the
lab frame and apply that transformation to the incoming momenta.

In the c.m.\ frame, the sum of three-momenta is zero.  In a scattering
process, the $z$ axis is the direction of the first beam, the second
beam is along the negative $z$ axis.  The transformation from the
c.m.\ to the lab frame is a rotation from the $z$ axis to the boost
axis followed by a boost, such that the c.m.\ momenta are transformed
into the lab-frame momenta.  In a decay process, we just boost along
the flight direction, without rotation.
\nwenddocs{}\nwbegincode{133}\moddef{PHS base: phs: TBP}\plusendmoddef
  procedure :: set_incoming_momenta => phs_set_incoming_momenta
\nwendcode{}\nwbegincode{134}\moddef{PHS base: procedures}\plusendmoddef
  subroutine phs_set_incoming_momenta (phs, p)
    class(phs_t), intent(inout) :: phs
    type(vector4_t), dimension(:), intent(in) :: p
    type(vector4_t) :: p0, p1
    type(lorentz_transformation_t) :: lt0
    phs%p = p
    if (phs%config%cm_frame) then
       phs%sqrts_hat = phs%config%sqrts
       phs%p = p
    else
       p0 = sum (p)
       if (phs%config%sqrts_fixed) then
          phs%sqrts_hat = phs%config%sqrts
       else
          phs%sqrts_hat = p0 ** 1
       end if
       lt0 = boost (p0, phs%sqrts_hat)
       select case (phs%config%n_in)
       case (1)
          phs%lt_cm_to_lab = lt0
       case (2)
          p1 = inverse (lt0) * p(1)
          phs%lt_cm_to_lab = lt0 * rotation_to_2nd (3, space_part (p1))
       end select
       phs%p = inverse (phs%lt_cm_to_lab) * p
    end if
    phs%p_defined = .true.
  end subroutine phs_set_incoming_momenta
  
\eatline
\nwendcode{}\nwbegindocs{135}Set outgoing momenta.  Assume that array shapes match.  The incoming
momenta must be known, so can apply the Lorentz transformation from
c.m.\ to lab (inverse) to the momenta.
\nwenddocs{}\nwbegincode{136}\moddef{PHS base: phs: TBP}\plusendmoddef
  procedure :: set_outgoing_momenta => phs_set_outgoing_momenta
\nwendcode{}\nwbegincode{137}\moddef{PHS base: procedures}\plusendmoddef
  subroutine phs_set_outgoing_momenta (phs, q)
    class(phs_t), intent(inout) :: phs
    type(vector4_t), dimension(:), intent(in) :: q
    if (phs%p_defined) then
       if (phs%config%cm_frame) then
          phs%q = q
       else
          phs%q = inverse (phs%lt_cm_to_lab) * q
       end if
       phs%q_defined = .true.
    end if
  end subroutine phs_set_outgoing_momenta
  
\eatline
\nwendcode{}\nwbegindocs{138}Return outgoing momenta.  Apply the c.m.\ to lab transformation if
necessary.
\nwenddocs{}\nwbegincode{139}\moddef{PHS base: phs: TBP}\plusendmoddef
  procedure :: get_outgoing_momenta => phs_get_outgoing_momenta
\nwendcode{}\nwbegincode{140}\moddef{PHS base: procedures}\plusendmoddef
  subroutine phs_get_outgoing_momenta (phs, q)
    class(phs_t), intent(in) :: phs
    type(vector4_t), dimension(:), intent(out) :: q
    if (phs%p_defined .and. phs%q_defined) then
       if (phs%config%cm_frame) then
          q = phs%q
       else
          q = phs%lt_cm_to_lab * phs%q
       end if
    else
       q = vector4_null
    end if
  end subroutine phs_get_outgoing_momenta
  
\eatline
\nwendcode{}\nwbegindocs{141}Return the input parameter array for a channel.
\nwenddocs{}\nwbegincode{142}\moddef{PHS base: phs: TBP}\plusendmoddef
  procedure :: get_mcpar => phs_get_mcpar
\nwendcode{}\nwbegincode{143}\moddef{PHS base: procedures}\plusendmoddef
  subroutine phs_get_mcpar (phs, c, r)
    class(phs_t), intent(in) :: phs
    integer, intent(in) :: c
    real(default), dimension(:), intent(out) :: r
    if (phs%r_defined) then
       r = phs%r(:,c)
    else
       r = 0
    end if
  end subroutine phs_get_mcpar
  
\eatline
\nwendcode{}\nwbegindocs{144}Return the Jacobian factor for a channel.
\nwenddocs{}\nwbegincode{145}\moddef{PHS base: phs: TBP}\plusendmoddef
  procedure :: get_f => phs_get_f
\nwendcode{}\nwbegincode{146}\moddef{PHS base: procedures}\plusendmoddef
  function phs_get_f (phs, c) result (f)
    class(phs_t), intent(in) :: phs
    integer, intent(in) :: c
    real(default) :: f
    if (phs%r_defined) then
       f = phs%f(c)
    else
       f = 0
    end if
  end function phs_get_f
    
\eatline
\nwendcode{}\nwbegindocs{147}Return the overall factor, which is the product of the flux factor for the
incoming partons and the phase-space volume for the outgoing partons.
\nwenddocs{}\nwbegincode{148}\moddef{PHS base: phs: TBP}\plusendmoddef
  procedure :: get_overall_factor => phs_get_overall_factor
\nwendcode{}\nwbegincode{149}\moddef{PHS base: procedures}\plusendmoddef
  function phs_get_overall_factor (phs) result (f)
    class(phs_t), intent(in) :: phs
    real(default) :: f
    f = phs%flux * phs%volume
  end function phs_get_overall_factor
  
\eatline
\nwendcode{}\nwbegindocs{150}Compute flux factor.  We do this during initialization (when the
incoming momenta {\tt{}p} are undefined), unless {\tt{}sqrts} is variable.  We do
this again once for each phase-space point, but then we skip the calculation
if {\tt{}sqrts} is fixed.

There are three different flux factors.
\begin{enumerate}
\item
  For a decaying massive particle, the factor is
  \begin{equation}
    f = (2\pi)^4 / (2M)
  \end{equation}
\item
  For a $2\to n$ scattering process with $n>1$, the factor is
  \begin{equation}
    f = (2\pi)^4 / (2\sqrt{\lambda})
  \end{equation}
  where for massless incoming particles, $\sqrt{\lambda} = s$.
\item For a $2\to 1$ on-shell production process, the factor includes
  an extra $1/(2\pi)^3$ factor and a $1/m^2$ factor from the
  phase-space delta function $\delta (x_1x_2 - m^2/s)$, which
  originate from the one-particle phase space that we integrate out.
  \begin{equation}
    f = 2\pi / (2s m^2)
  \end{equation}
  The delta function is handled by the structure-function
  parameterization.
\end{enumerate}
\nwenddocs{}\nwbegincode{151}\moddef{PHS base: phs: TBP}\plusendmoddef
  procedure :: compute_flux => phs_compute_flux
\nwendcode{}\nwbegincode{152}\moddef{PHS base: procedures}\plusendmoddef
  subroutine phs_compute_flux (phs)
    class(phs_t), intent(inout) :: phs
    real(default) :: s_hat, lda
    select case (phs%config%n_in)
    case (1)
       if (.not. phs%p_defined) then
          phs%flux = twopi4 / (2 * phs%m_in(1))
       end if
    case (2)
       if (phs%p_defined) then
          if (phs%config%sqrts_fixed) then
             return
          else
             s_hat = sum (phs%p) ** 2
          end if
       else
          if (phs%config%sqrts_fixed) then
             s_hat = phs%config%sqrts ** 2
          else
             return
          end if
       end if
       select case (phs%config%n_out)
       case (2:)
          lda = lambda (s_hat, phs%m_in(1) ** 2, phs%m_in(2) ** 2)
          if (lda > 0) then
             phs%flux = conv * twopi4 / (2 * sqrt (lda))
          else
             phs%flux = 0
          end if
       case (1)
          phs%flux = conv * twopi &
               / (2 * phs%config%sqrts ** 2 * phs%m_out(1) ** 2)
       case default
          phs%flux = 0
       end select
    end select
  end subroutine phs_compute_flux

\eatline
\nwendcode{}\nwbegindocs{153}Evaluate the phase-space point for a particular channel and compute momenta,
Jacobian, and phase-space volume. This is, of course, deferred to
the implementation.
\nwenddocs{}\nwbegincode{154}\moddef{PHS base: phs: TBP}\plusendmoddef
  procedure (phs_evaluate_selected_channel), deferred :: &
       evaluate_selected_channel
\nwendcode{}\nwbegincode{155}\moddef{PHS base: interfaces}\plusendmoddef
  abstract interface
     subroutine phs_evaluate_selected_channel (phs, c_in, r_in)
       import
       class(phs_t), intent(inout) :: phs
       integer, intent(in) :: c_in
       real(default), dimension(:), intent(in) :: r_in
     end subroutine phs_evaluate_selected_channel
  end interface
  
\eatline
\nwendcode{}\nwbegindocs{156}Compute the inverse mappings to completely fill the {\tt{}r} and {\tt{}f} arrays,
for the non-selected channels.
\nwenddocs{}\nwbegincode{157}\moddef{PHS base: phs: TBP}\plusendmoddef
  procedure (phs_evaluate_other_channels), deferred :: &
       evaluate_other_channels
\nwendcode{}\nwbegincode{158}\moddef{PHS base: interfaces}\plusendmoddef
  abstract interface
     subroutine phs_evaluate_other_channels (phs, c_in)
       import
       class(phs_t), intent(inout) :: phs
       integer, intent(in) :: c_in
     end subroutine phs_evaluate_other_channels
  end interface
  
\eatline
\nwendcode{}\nwbegindocs{159}Inverse evaluation.  If all momenta are known, we compute the
inverse mappings to fill the {\tt{}r} and {\tt{}f} arrays.
\nwenddocs{}\nwbegincode{160}\moddef{PHS base: phs: TBP}\plusendmoddef
  procedure (phs_inverse), deferred :: inverse
\nwendcode{}\nwbegincode{161}\moddef{PHS base: interfaces}\plusendmoddef
  abstract interface
     subroutine phs_inverse (phs)
       import
       class(phs_t), intent(inout) :: phs
     end subroutine phs_inverse
  end interface
  
\eatline
\nwendcode{}\nwbegindocs{162}\nwdocspar
\subsection{Auxiliary stuff}
The {\tt{}pacify} subroutine, which is provided by the Lorentz module,
has the purpose of setting numbers to zero which are (by comparing
with a {\tt{}tolerance} parameter) considered equivalent with zero.  This
is useful for numerical checks.
\nwenddocs{}\nwbegincode{163}\moddef{PHS base: public}\plusendmoddef
  public :: pacify
\nwendcode{}\nwbegincode{164}\moddef{PHS base: interfaces}\plusendmoddef
  interface pacify
     module procedure pacify_phs
  end interface pacify
  
\nwendcode{}\nwbegincode{165}\moddef{PHS base: procedures}\plusendmoddef
  subroutine pacify_phs (phs)
    class(phs_t), intent(inout) :: phs
    if (phs%p_defined) then
       call pacify (phs%p, 30 * epsilon (1._default) * phs%config%sqrts)
       call pacify (phs%lt_cm_to_lab, 30 * epsilon (1._default))
    end if
    if (phs%q_defined) then
       call pacify (phs%q, 30 * epsilon (1._default) * phs%config%sqrts)
    end if
  end subroutine pacify_phs
    
\eatline
\nwendcode{}\nwbegindocs{166}\nwdocspar
\subsection{Unit tests}
\nwenddocs{}\nwbegincode{167}\moddef{PHS base: public}\plusendmoddef
  public :: phs_base_test
\nwendcode{}\nwbegincode{168}\moddef{PHS base: tests}\endmoddef
  subroutine phs_base_test (u, results)
    integer, intent(in) :: u
    type(test_results_t), intent(inout) :: results
  \LA{}PHS base: execute tests\RA{}
  end subroutine phs_base_test
  
\eatline
\nwendcode{}\nwbegindocs{169}\nwdocspar
\subsubsection{Test process data}
We provide a procedure that initializes a test case for the process
constants.  This set of process data contains just the minimal
contents that we need for the phase space.  The rest is left
uninitialized.
\nwenddocs{}\nwbegincode{170}\moddef{PHS base: public}\plusendmoddef
  public :: init_test_process_data
\nwendcode{}\nwbegincode{171}\moddef{PHS base: tests}\plusendmoddef
  subroutine init_test_process_data (id, data)
    type(process_constants_t), intent(out) :: data 
    type(string_t), intent(in), optional :: id
    if (present (id)) then
       data%id = id
    else
       data%id = "testproc"
    end if
    data%model_name = "Test"
    data%n_in = 2
    data%n_out = 2
    data%n_flv = 1
    allocate (data%flv_state (data%n_in + data%n_out, data%n_flv))
    data%flv_state = 25
  end subroutine init_test_process_data

\eatline
\nwendcode{}\nwbegindocs{172}This is the variant for a decay process.
\nwenddocs{}\nwbegincode{173}\moddef{PHS base: public}\plusendmoddef
  public :: init_test_decay_data
\nwendcode{}\nwbegincode{174}\moddef{PHS base: tests}\plusendmoddef
  subroutine init_test_decay_data (id, data)
    type(process_constants_t), intent(out) :: data 
    type(string_t), intent(in), optional :: id
    if (present (id)) then
       data%id = id
    else
       data%id = "testproc"
    end if
    data%model_name = "Test"
    data%n_in = 1
    data%n_out = 2
    data%n_flv = 1
    allocate (data%flv_state (data%n_in + data%n_out, data%n_flv))
    data%flv_state(:,1) = [25, 6, -6]
  end subroutine init_test_decay_data

\eatline
\nwendcode{}\nwbegindocs{175}\nwdocspar
\subsubsection{Test kinematics configuration}
This is a trivial implementation of the {\tt{}phs{\char95}config{\char95}t} configuration object.
\nwenddocs{}\nwbegincode{176}\moddef{PHS base: public}\plusendmoddef
  public :: phs_test_config_t
\nwendcode{}\nwbegincode{177}\moddef{PHS base: test types}\endmoddef
  type, extends (phs_config_t) :: phs_test_config_t
     logical :: create_equivalences = .false.
   contains
     procedure :: final => phs_test_config_final
     procedure :: write => phs_test_config_write
     procedure :: configure => phs_test_config_configure
     procedure :: startup_message => phs_test_config_startup_message
     procedure, nopass :: allocate_instance => phs_test_config_allocate_instance
  end type phs_test_config_t

\eatline
\nwendcode{}\nwbegindocs{178}The finalizer is empty.
\nwenddocs{}\nwbegincode{179}\moddef{PHS base: tests}\plusendmoddef
  subroutine phs_test_config_final (object)
    class(phs_test_config_t), intent(inout) :: object
  end subroutine phs_test_config_final

\eatline
\nwendcode{}\nwbegindocs{180}The {\tt{}cm{\char95}frame} parameter is not tested here; we defer this to the
{\tt{}phs{\char95}single} implementation.
\nwenddocs{}\nwbegincode{181}\moddef{PHS base: tests}\plusendmoddef
  subroutine phs_test_config_write (object, unit)
    class(phs_test_config_t), intent(in) :: object
    integer, intent(in), optional :: unit
    integer :: u
    u = given_output_unit (unit)
    write (u, "(1x,A)")  "Partonic phase-space configuration:"
    call object%base_write (unit)
  end subroutine phs_test_config_write
  
  subroutine phs_test_config_configure (phs_config, sqrts, &
       sqrts_fixed, cm_frame, azimuthal_dependence, rebuild, ignore_mismatch, &
       nlo_type)
    class(phs_test_config_t), intent(inout) :: phs_config
    real(default), intent(in) :: sqrts
    logical, intent(in), optional :: sqrts_fixed
    logical, intent(in), optional :: cm_frame
    logical, intent(in), optional :: azimuthal_dependence
    logical, intent(in), optional :: rebuild
    logical, intent(in), optional :: ignore_mismatch
    type(string_t), intent(inout), optional :: nlo_type
    phs_config%n_channel = 2
    phs_config%n_par = 2
    phs_config%sqrts = sqrts
    if (.not. present (nlo_type)) &
      phs_config%nlo_type = 'Born'
    if (present (sqrts_fixed)) then
       phs_config%sqrts_fixed = sqrts_fixed
    end if
    if (present (cm_frame)) then
       phs_config%cm_frame = cm_frame
    end if
    if (present (azimuthal_dependence)) then
       phs_config%azimuthal_dependence = azimuthal_dependence
    end if
    if (allocated (phs_config%channel))  deallocate (phs_config%channel)
    allocate (phs_config%channel (phs_config%n_channel))
    if (phs_config%create_equivalences) then
       call setup_test_equivalences (phs_config)
       call setup_test_channel_props (phs_config)
    end if
    call phs_config%compute_md5sum ()
  end subroutine phs_test_config_configure

\eatline
\eatline
\nwendcode{}\nwbegindocs{182}If requested, we make up an arbitrary set of equivalences.
\nwenddocs{}\nwbegincode{183}\moddef{PHS base: tests}\plusendmoddef
  subroutine setup_test_equivalences (phs_config)
    class(phs_test_config_t), intent(inout) :: phs_config
    integer :: i
    associate (channel => phs_config%channel(1))
      allocate (channel%eq (2))
      do i = 1, size (channel%eq)
         call channel%eq(i)%init (phs_config%n_par)
      end do
      associate (eq => channel%eq(1))
        eq%c = 1;  eq%perm = [1, 2];  eq%mode = [EQ_IDENTITY, EQ_SYMMETRIC]
      end associate
      associate (eq => channel%eq(2))
        eq%c = 2;  eq%perm = [2, 1];  eq%mode = [EQ_INVARIANT, EQ_IDENTITY]
      end associate
    end associate
  end subroutine setup_test_equivalences
      
\eatline
\nwendcode{}\nwbegindocs{184}Ditto, for channel properties.
\nwenddocs{}\nwbegincode{185}\moddef{PHS base: tests}\plusendmoddef
  subroutine setup_test_channel_props (phs_config)
    class(phs_test_config_t), intent(inout) :: phs_config
    associate (channel => phs_config%channel(2))
      call channel%set_resonant (140._default, 3.1415_default)
    end associate
  end subroutine setup_test_channel_props
      
\eatline
\nwendcode{}\nwbegindocs{186}Startup message
\nwenddocs{}\nwbegincode{187}\moddef{PHS base: tests}\plusendmoddef
  subroutine phs_test_config_startup_message (phs_config, unit)
    class(phs_test_config_t), intent(in) :: phs_config
    integer, intent(in), optional :: unit
    call phs_config%base_startup_message (unit)
    write (msg_buffer, "(A)") "Phase space: Test"
    call msg_message (unit = unit)
  end subroutine phs_test_config_startup_message
  
\eatline
\nwendcode{}\nwbegindocs{188}The instance type that matches {\tt{}phs{\char95}test{\char95}config{\char95}t} is {\tt{}phs{\char95}test{\char95}t}.
\nwenddocs{}\nwbegincode{189}\moddef{PHS base: tests}\plusendmoddef
  subroutine phs_test_config_allocate_instance (phs)
    class(phs_t), intent(inout), pointer :: phs
    allocate (phs_test_t :: phs)
  end subroutine phs_test_config_allocate_instance
  
\eatline
\nwendcode{}\nwbegindocs{190}\nwdocspar
\subsubsection{Test kinematics implementation}
This implementation of kinematics generates a simple two-particle
configuration from the incoming momenta.  The incoming momenta must be
in the c.m.\ system, all masses equal.

There are two channels: one generates $\cos\theta$ and $\phi$
uniformly, in the other channel we map the $r_1$ parameter which
belongs to $\cos\theta$.

We should store the mass parameter that we need.
\nwenddocs{}\nwbegincode{191}\moddef{PHS base: public}\plusendmoddef
  public :: phs_test_t
\nwendcode{}\nwbegincode{192}\moddef{PHS base: test types}\plusendmoddef
  type, extends (phs_t) :: phs_test_t
     real(default) :: m = 0
     real(default), dimension(:), allocatable :: x
   contains
   \LA{}PHS base: phs test: TBP\RA{}
  end type phs_test_t
  
\eatline
\nwendcode{}\nwbegindocs{193}Output.  The specific data are displayed only if {\tt{}verbose} is set.
\nwenddocs{}\nwbegincode{194}\moddef{PHS base: phs test: TBP}\endmoddef
  procedure :: write => phs_test_write
\nwendcode{}\nwbegincode{195}\moddef{PHS base: tests}\plusendmoddef
  subroutine phs_test_write (object, unit, verbose)
    class(phs_test_t), intent(in) :: object
    integer, intent(in), optional :: unit
    logical, intent(in), optional :: verbose
    integer :: u
    logical :: verb
    u = given_output_unit (unit)
    verb = .false.;  if (present (verbose))  verb = verbose
    if (verb) then
       write (u, "(1x,A)")  "Partonic phase space: data"
       write (u, "(3x,A," // FMT_19 // ")")  "m = ", object%m
    end if
    call object%base_write (u)
  end subroutine phs_test_write
    
\eatline
\nwendcode{}\nwbegindocs{196}The finalizer is empty.
\nwenddocs{}\nwbegincode{197}\moddef{PHS base: phs test: TBP}\plusendmoddef
  procedure :: final => phs_test_final
\nwendcode{}\nwbegincode{198}\moddef{PHS base: tests}\plusendmoddef
  subroutine phs_test_final (object)
    class(phs_test_t), intent(inout) :: object
  end subroutine phs_test_final
  
\eatline
\nwendcode{}\nwbegindocs{199}Initialization: set the mass value.
\nwenddocs{}\nwbegincode{200}\moddef{PHS base: phs test: TBP}\plusendmoddef
  procedure :: init => phs_test_init
\nwendcode{}\nwbegincode{201}\moddef{PHS base: tests}\plusendmoddef
  subroutine phs_test_init (phs, phs_config)
    class(phs_test_t), intent(out) :: phs
    class(phs_config_t), intent(in), target :: phs_config
    call phs%base_init (phs_config)
    phs%m = phs%config%flv(1,1)%get_mass ()
    allocate (phs%x (phs_config%n_par), source = 0._default)
  end subroutine phs_test_init
  
\eatline
\nwendcode{}\nwbegindocs{202}Evaluation.  In channel 1, we uniformly generate $\cos\theta$ and
$\phi$, with Jacobian normalized to one.  In channel 2, we prepend a
mapping $r_1 \to r_1^(1/3)$ with Jacobian $f=3r_1^2$.

The component {\tt{}x} is allocated in the first subroutine, used and deallocated
in the second one.
\nwenddocs{}\nwbegincode{203}\moddef{PHS base: phs test: TBP}\plusendmoddef
  procedure :: evaluate_selected_channel => phs_test_evaluate_selected_channel
  procedure :: evaluate_other_channels => phs_test_evaluate_other_channels
\nwendcode{}\nwbegincode{204}\moddef{PHS base: tests}\plusendmoddef
  subroutine phs_test_evaluate_selected_channel (phs, c_in, r_in)
    class(phs_test_t), intent(inout) :: phs
    integer, intent(in) :: c_in
    real(default), intent(in), dimension(:) :: r_in
    if (phs%p_defined) then
       call phs%select_channel (c_in)
       phs%r(:,c_in) = r_in
       select case (c_in)
       case (1)
          phs%x = r_in
       case (2)
          phs%x(1) = r_in(1) ** (1 / 3._default)
          phs%x(2) = r_in(2)
       end select
       call compute_kinematics_solid_angle (phs%p, phs%q, phs%x) 
       phs%volume = 1
       phs%q_defined = .true.
    end if
  end subroutine phs_test_evaluate_selected_channel
  
  subroutine phs_test_evaluate_other_channels (phs, c_in)
    class(phs_test_t), intent(inout) :: phs
    integer, intent(in) :: c_in
    integer :: c, n_channel
    if (phs%p_defined) then
       n_channel = phs%config%n_channel
       do c = 1, n_channel
          if (c /= c_in) then
             call inverse_kinematics_solid_angle (phs%p, phs%q, phs%x)
             select case (c)
             case (1)
                phs%r(:,c) = phs%x
             case (2)
                phs%r(1,c) = phs%x(1) ** 3
                phs%r(2,c) = phs%x(2)
             end select
          end if
       end do
       phs%f(1) = 1
       if (phs%r(1,2) /= 0) then
          phs%f(2) = 1 / (3 * phs%r(1,2) ** (2/3._default))
       else
          phs%f(2) = 0
       end if
       phs%r_defined = .true.
    end if
  end subroutine phs_test_evaluate_other_channels
  
\eatline
\eatline
\nwendcode{}\nwbegindocs{205}Inverse evaluation.
\nwenddocs{}\nwbegincode{206}\moddef{PHS base: phs test: TBP}\plusendmoddef
  procedure :: inverse => phs_test_inverse
\nwendcode{}\nwbegincode{207}\moddef{PHS base: tests}\plusendmoddef
  subroutine phs_test_inverse (phs)
    class(phs_test_t), intent(inout) :: phs
    integer :: c, n_channel
    real(default), dimension(:), allocatable :: x
    if (phs%p_defined .and. phs%q_defined) then
       call phs%select_channel ()
       n_channel = phs%config%n_channel
       allocate (x (phs%config%n_par))
       do c = 1, n_channel
          call inverse_kinematics_solid_angle (phs%p, phs%q, x)
          select case (c)
          case (1)
             phs%r(:,c) = x
          case (2)
             phs%r(1,c) = x(1) ** 3
             phs%r(2,c) = x(2)
          end select
       end do
       phs%f(1) = 1
       if (phs%r(1,2) /= 0) then
          phs%f(2) = 1 / (3 * phs%r(1,2) ** (2/3._default))
       else
          phs%f(2) = 0
       end if
       phs%volume = 1
       phs%r_defined = .true.
    end if
  end subroutine phs_test_inverse
  
\eatline
\nwendcode{}\nwbegindocs{208}\nwdocspar
\subsubsection{Uniform angular distribution}
These procedures implement the uniform angular distribution, generated
from two parameters $x_1$ and $x_2$:
\begin{equation}
  \cos\theta = 1 - 2x_1, \qquad \phi = 2\pi x_2
\end{equation}
We generate a rotation (Lorentz transformation) which rotates the
positive $z$ axis into this point on the unit sphere.  This rotation
is applied to the {\tt{}p} momenta, which are assumed to be
back-to-back, on-shell, and with the correct mass.

We do not compute a Jacobian (constant).  The uniform distribution is
assumed to be normalized.
\nwenddocs{}\nwbegincode{209}\moddef{PHS base: public}\plusendmoddef
  public :: compute_kinematics_solid_angle
\nwendcode{}\nwbegincode{210}\moddef{PHS base: tests}\plusendmoddef
  subroutine compute_kinematics_solid_angle (p, q, x)
    type(vector4_t), dimension(2), intent(in) :: p
    type(vector4_t), dimension(2), intent(out) :: q
    real(default), dimension(2), intent(in) :: x
    real(default) :: ct, st, phi
    type(lorentz_transformation_t) :: rot
    integer :: i
    ct = 1 - 2*x(1)
    st = sqrt (1 - ct**2)
    phi = twopi * x(2)
    rot = rotation (phi, 3) * rotation (ct, st, 2)
    do i = 1, 2
       q(i) = rot * p(i)
    end do
  end subroutine compute_kinematics_solid_angle
  
\eatline
\nwendcode{}\nwbegindocs{211}This is the inverse transformation.  We assume that the outgoing
momenta are rotated versions of the incoming momenta, back-to-back.
Thus, we determine the angles from $q(1)$ alone.  {\tt{}p} is unused.
\nwenddocs{}\nwbegincode{212}\moddef{PHS base: public}\plusendmoddef
  public :: inverse_kinematics_solid_angle
\nwendcode{}\nwbegincode{213}\moddef{PHS base: tests}\plusendmoddef
  subroutine inverse_kinematics_solid_angle (p, q, x)
    type(vector4_t), dimension(:), intent(in) :: p
    type(vector4_t), dimension(2), intent(in) :: q
    real(default), dimension(2), intent(out) :: x
    real(default) :: ct, phi
    ct = polar_angle_ct (q(1))
    phi = azimuthal_angle (q(1))
    x(1) = (1 - ct) / 2
    x(2) = phi / twopi
  end subroutine inverse_kinematics_solid_angle

\eatline
\nwendcode{}\nwbegindocs{214}\nwdocspar
\subsubsection{Phase-space configuration data}
Construct and display a test phase-space configuration object.
\nwenddocs{}\nwbegincode{215}\moddef{PHS base: execute tests}\endmoddef
  call test (phs_base_1, "phs_base_1", &
       "phase-space configuration", &
       u, results)
\nwendcode{}\nwbegincode{216}\moddef{PHS base: tests}\plusendmoddef
  subroutine phs_base_1 (u)
    integer, intent(in) :: u
    type(model_data_t), target :: model
    type(process_constants_t) :: process_data
    class(phs_config_t), allocatable :: phs_data
    
    write (u, "(A)")  "* Test output: phs_base_1"
    write (u, "(A)")  "*   Purpose: initialize and display &
         &test phase-space configuration data"
    write (u, "(A)")
    
    call model%init_test ()

    write (u, "(A)")  "* Initialize a process and a matching &
         &phase-space configuration"
    write (u, "(A)")

    call init_test_process_data (var_str ("phs_base_1"), process_data)

    allocate (phs_test_config_t :: phs_data)
    call phs_data%init (process_data, model)
       
    call phs_data%write (u)
    
    call phs_data%final ()
    call model%final ()
    
    write (u, "(A)")
    write (u, "(A)")  "* Test output end: phs_base_1"

  end subroutine phs_base_1

\eatline
\nwendcode{}\nwbegindocs{217}\nwdocspar
\subsubsection{Phase space evaluation}
Compute kinematics for given parameters, also invert the calculation.
\nwenddocs{}\nwbegincode{218}\moddef{PHS base: execute tests}\plusendmoddef
  call test (phs_base_2, "phs_base_2", &
       "phase-space evaluation", &
       u, results)
\nwendcode{}\nwbegincode{219}\moddef{PHS base: tests}\plusendmoddef
  subroutine phs_base_2 (u)
    integer, intent(in) :: u
    type(model_data_t), target :: model
    type(flavor_t) :: flv
    type(process_constants_t) :: process_data
    real(default) :: sqrts, E
    class(phs_config_t), allocatable, target :: phs_data
    class(phs_t), pointer :: phs => null ()
    type(vector4_t), dimension(2) :: p, q
    
    write (u, "(A)")  "* Test output: phs_base_2"
    write (u, "(A)")  "*   Purpose: test simple two-channel phase space"
    write (u, "(A)")
    
    call model%init_test ()
    call flv%init (25, model)

    write (u, "(A)")  "* Initialize a process and a matching &
         &phase-space configuration"
    write (u, "(A)")

    call init_test_process_data (var_str ("phs_base_2"), process_data)

    allocate (phs_test_config_t :: phs_data)
    call phs_data%init (process_data, model)

    sqrts = 1000._default
    call phs_data%configure (sqrts)

    call phs_data%write (u)
    
    write (u, "(A)")
    write (u, "(A)")  "* Initialize the phase-space instance"
    write (u, "(A)")

    call phs_data%allocate_instance (phs)
    select type (phs)
    type is (phs_test_t)
       call phs%init (phs_data)
    end select
       
    call phs%write (u, verbose=.true.)
    
    write (u, "(A)")
    write (u, "(A)")  "* Set incoming momenta"
    write (u, "(A)")

    E = sqrts / 2
    p(1) = vector4_moving (E, sqrt (E**2 - flv%get_mass ()**2), 3)
    p(2) = vector4_moving (E,-sqrt (E**2 - flv%get_mass ()**2), 3)

    call phs%set_incoming_momenta (p)
    call phs%compute_flux ()
    call phs%write (u)
    
    write (u, "(A)")
    write (u, "(A)")  "* Compute phase-space point in channel 1 &
         &for x = 0.5, 0.125"
    write (u, "(A)")

    call phs%evaluate_selected_channel (1, [0.5_default, 0.125_default])
    call phs%evaluate_other_channels (1)
    call phs%write (u)
    
    write (u, "(A)")
    write (u, "(A)")  "* Compute phase-space point in channel 2 &
         &for x = 0.125, 0.125"
    write (u, "(A)")

    call phs%evaluate_selected_channel (2, [0.125_default, 0.125_default])
    call phs%evaluate_other_channels (2)
    call phs%write (u)
    
    write (u, "(A)")
    write (u, "(A)")  "* Inverse kinematics"
    write (u, "(A)")

    call phs%get_outgoing_momenta (q)
    deallocate (phs)
    call phs_data%allocate_instance (phs)
    call phs%init (phs_data)
       
    sqrts = 1000._default
    select type (phs_data)
    type is (phs_test_config_t)
       call phs_data%configure (sqrts)
    end select

    call phs%set_incoming_momenta (p)
    call phs%compute_flux ()
    call phs%set_outgoing_momenta (q)
    
    call phs%inverse ()
    call phs%write (u)
    
    call phs%final ()
    deallocate (phs)
    
    call phs_data%final ()
    call model%final ()
    
    write (u, "(A)")
    write (u, "(A)")  "* Test output end: phs_base_2"

  end subroutine phs_base_2

\eatline
\nwendcode{}\nwbegindocs{220}\nwdocspar
\subsubsection{Phase-space equivalences}
Construct a test phase-space configuration which contains channel
equivalences.
\nwenddocs{}\nwbegincode{221}\moddef{PHS base: execute tests}\plusendmoddef
  call test (phs_base_3, "phs_base_3", &
       "channel equivalences", &
       u, results)
\nwendcode{}\nwbegincode{222}\moddef{PHS base: tests}\plusendmoddef
  subroutine phs_base_3 (u)
    integer, intent(in) :: u
    type(model_data_t), target :: model
    type(process_constants_t) :: process_data
    class(phs_config_t), allocatable :: phs_data
    
    write (u, "(A)")  "* Test output: phs_base_3"
    write (u, "(A)")  "*   Purpose: construct phase-space configuration data &
         &with equivalences"
    write (u, "(A)")
    
    call model%init_test ()

    write (u, "(A)")  "* Initialize a process and a matching &
         &phase-space configuration"
    write (u, "(A)")

    call init_test_process_data (var_str ("phs_base_3"), process_data)

    allocate (phs_test_config_t :: phs_data)
    call phs_data%init (process_data, model)
    select type (phs_data)
    type is (phs_test_config_t)
       phs_data%create_equivalences = .true.
    end select
       
    call phs_data%configure (1000._default)
    call phs_data%write (u)
    
    call phs_data%final ()
    call model%final ()
    
    write (u, "(A)")
    write (u, "(A)")  "* Test output end: phs_base_3"

  end subroutine phs_base_3

\eatline
\nwendcode{}\nwbegindocs{223}\nwdocspar
\subsubsection{MD5 sum checks}
Construct a test phase-space configuration, compute and compare MD5 sums.
\nwenddocs{}\nwbegincode{224}\moddef{PHS base: execute tests}\plusendmoddef
  call test (phs_base_4, "phs_base_4", &
       "MD5 sum", &
       u, results)
\nwendcode{}\nwbegincode{225}\moddef{PHS base: tests}\plusendmoddef
  subroutine phs_base_4 (u)
    integer, intent(in) :: u
    type(model_data_t), target :: model
    type(process_constants_t) :: process_data
    class(phs_config_t), allocatable :: phs_data
    
    write (u, "(A)")  "* Test output: phs_base_4"
    write (u, "(A)")  "*   Purpose: compute and compare MD5 sums"
    write (u, "(A)")
    
    call model%init_test ()

    write (u, "(A)")  "* Model parameters"
    write (u, "(A)")

    call model%write (unit = u, &
         show_parameters = .true., &
         show_particles = .false., show_vertices = .false.)

    write (u, "(A)")
    write (u, "(A)")  "* Initialize a process and a matching &
         &phase-space configuration"
    write (u, "(A)")

    call init_test_process_data (var_str ("phs_base_4"), process_data)
    process_data%md5sum = "test_process_data_m6sum_12345678"

    allocate (phs_test_config_t :: phs_data)
    call phs_data%init (process_data, model)
    
    call phs_data%compute_md5sum ()
    call phs_data%write (u)
    
    write (u, "(A)")
    write (u, "(A)")  "* Modify model parameter"
    write (u, "(A)")

    call model%set_par (var_str ("ms"), 100._default)
    call model%write (show_parameters = .true., &
         show_particles = .false., show_vertices = .false.)

    write (u, "(A)")
    write (u, "(A)")  "* PHS configuration"
    write (u, "(A)")

    call phs_data%compute_md5sum ()
    call phs_data%write (u)

    call phs_data%final ()
    call model%final ()
    
    write (u, "(A)")
    write (u, "(A)")  "* Test output end: phs_base_4"

  end subroutine phs_base_4

\eatline
\nwendcode{}\nwbegindocs{226}\nwdocspar
\subsubsection{Phase-space channel collection}
Set up an array of various phase-space channels and collect them in a list.
\nwenddocs{}\nwbegincode{227}\moddef{PHS base: execute tests}\plusendmoddef
  call test (phs_base_5, "phs_base_5", &
       "channel collection", &
       u, results)
\nwendcode{}\nwbegincode{228}\moddef{PHS base: tests}\plusendmoddef
  subroutine phs_base_5 (u)
    integer, intent(in) :: u
    type(phs_channel_t), dimension(:), allocatable :: channel
    type(phs_channel_collection_t) :: coll
    integer :: i, n
    
    write (u, "(A)")  "* Test output: phs_base_5"
    write (u, "(A)")  "*   Purpose: collect channel properties"
    write (u, "(A)")
    
    write (u, "(A)")  "* Set up an array of channels"
    write (u, "(A)")

    n = 6

    allocate (channel (n))
    call channel(2)%set_resonant (75._default, 3._default)
    call channel(4)%set_resonant (130._default, 1._default)
    call channel(5)%set_resonant (75._default, 3._default)
    call channel(6)%set_on_shell (33._default)

    do i = 1, n
       write (u, "(1x,I0)", advance="no")  i
       call channel(i)%write (u)
    end do
    
    write (u, "(A)")
    write (u, "(A)")  "* Collect distinct properties"
    write (u, "(A)")
    
    do i = 1, n
       call coll%push (channel(i))
    end do

    write (u, "(1x,A,I0)")  "n = ", coll%get_n ()
    write (u, "(A)")

    call coll%write (u)

    write (u, "(A)")
    write (u, "(A)")  "* Channel array with collection index assigned"
    write (u, "(A)")

    do i = 1, n
       write (u, "(1x,I0)", advance="no")  i
       call channel(i)%write (u)
    end do
    
    write (u, "(A)")
    write (u, "(A)")  "* Cleanup"
    
    call coll%final ()

    write (u, "(A)")
    write (u, "(A)")  "* Test output end: phs_base_5"

  end subroutine phs_base_5

\eatline
\nwendcode{}\nwbegindocs{229}\nwdocspar
%%%%%%%%%%%%%%%%%%%%%%%%%%%%%%%%%%%%%%%%%%%%%%%%%%%%%%%%%%%%%%%%%%%%%%%%
\clearpage
\section{Single-particle phase space}

This module implements the phase space for a single particle, i.e., the solid
angle, in a straightforward parameterization with a single channel.  The
phase-space implementation may be used either for $1\to 2$ decays or for $2\to
2$ scattering processes, so the number of incoming particles is the only free
parameter in the configuration.  In the latter case, we should restrict its
use to non-resonant s-channel processes, because there is no mapping of the
scattering angle.

(We might extend this later to account for generic $2\to 2$ situations, e.g.,
account for a Coulomb singularity or detect an s-channel resonance structure
that requires matching structure-function mappings.)

This is derived from the {\tt{}phs{\char95}test} implementation in the
{\tt{}phs{\char95}base} module above, even more simplified, but intended for
actual use.
\nwenddocs{}\nwbegincode{230}\moddef{\code{}phs{\char95}single.f90\edoc{}}\endmoddef
\LA{}File header\RA{}

module phs_single

\LA{}Use kinds\RA{}
\LA{}Use strings\RA{}
  use io_units
  use constants
  use unit_tests
  use diagnostics
  use os_interface
  use lorentz
  use model_data
  use flavors
  use process_constants
  use sf_mappings
  use sf_base
  use phs_base
  
\LA{}Standard module head\RA{}

\LA{}PHS single: public\RA{}

\LA{}PHS single: types\RA{}

contains

\LA{}PHS single: procedures\RA{}

\LA{}PHS single: tests\RA{}

end module phs_single
\eatline
\nwendcode{}\nwbegindocs{231}\nwdocspar
\subsection{Configuration}
\nwenddocs{}\nwbegincode{232}\moddef{PHS single: public}\endmoddef
  public :: phs_single_config_t
\nwendcode{}\nwbegincode{233}\moddef{PHS single: types}\endmoddef
  type, extends (phs_config_t) :: phs_single_config_t  
   contains
   \LA{}PHS single: phs single config: TBP\RA{}
  end type phs_single_config_t

\eatline
\nwendcode{}\nwbegindocs{234}The finalizer is empty.
\nwenddocs{}\nwbegincode{235}\moddef{PHS single: phs single config: TBP}\endmoddef
  procedure :: final => phs_single_config_final
\nwendcode{}\nwbegincode{236}\moddef{PHS single: procedures}\endmoddef
  subroutine phs_single_config_final (object)
    class(phs_single_config_t), intent(inout) :: object
  end subroutine phs_single_config_final

\eatline
\nwendcode{}\nwbegindocs{237}Output.
\nwenddocs{}\nwbegincode{238}\moddef{PHS single: phs single config: TBP}\plusendmoddef
  procedure :: write => phs_single_config_write
\nwendcode{}\nwbegincode{239}\moddef{PHS single: procedures}\plusendmoddef
  subroutine phs_single_config_write (object, unit)
    class(phs_single_config_t), intent(in) :: object
    integer, intent(in), optional :: unit
    integer :: u
    u = given_output_unit (unit)
    write (u, "(1x,A)")  "Partonic phase-space configuration (single-particle):"
    call object%base_write (unit)
  end subroutine phs_single_config_write
  
\eatline
\nwendcode{}\nwbegindocs{240}Configuration: there is only one channel and two parameters.  The
second parameter is the azimuthal angle, which may be a flat dimension.
\nwenddocs{}\nwbegincode{241}\moddef{PHS single: phs single config: TBP}\plusendmoddef
  procedure :: configure => phs_single_config_configure
\nwendcode{}\nwbegincode{242}\moddef{PHS single: procedures}\plusendmoddef
  subroutine phs_single_config_configure (phs_config, sqrts, &
       sqrts_fixed, cm_frame, azimuthal_dependence, rebuild, ignore_mismatch, &
       nlo_type)
    class(phs_single_config_t), intent(inout) :: phs_config
    real(default), intent(in) :: sqrts
    logical, intent(in), optional :: sqrts_fixed
    logical, intent(in), optional :: cm_frame
    logical, intent(in), optional :: azimuthal_dependence
    logical, intent(in), optional :: rebuild
    logical, intent(in), optional :: ignore_mismatch
    type(string_t), intent(inout), optional :: nlo_type
    if (.not. present (nlo_type)) &
      phs_config%nlo_type = 'Born'
    if (phs_config%n_out == 2) then
       phs_config%n_channel = 1
       phs_config%n_par = 2
       phs_config%sqrts = sqrts
       if (present (sqrts_fixed))  phs_config%sqrts_fixed = sqrts_fixed
       if (present (cm_frame))  phs_config%cm_frame = cm_frame
       if (present (azimuthal_dependence)) then
          phs_config%azimuthal_dependence = azimuthal_dependence
          if (.not. azimuthal_dependence) then
             allocate (phs_config%dim_flat (1))
             phs_config%dim_flat(1) = 2
          end if
       end if
       if (allocated (phs_config%channel))  deallocate (phs_config%channel)
       allocate (phs_config%channel (1))
       call phs_config%compute_md5sum ()
    else
       call msg_fatal ("Single-particle phase space requires n_out = 2")
    end if
  end subroutine phs_single_config_configure
  
\eatline
\nwendcode{}\nwbegindocs{243}Startup message, after configuration is complete.
\nwenddocs{}\nwbegincode{244}\moddef{PHS single: phs single config: TBP}\plusendmoddef
  procedure :: startup_message => phs_single_config_startup_message
\nwendcode{}\nwbegincode{245}\moddef{PHS single: procedures}\plusendmoddef
  subroutine phs_single_config_startup_message (phs_config, unit)
    class(phs_single_config_t), intent(in) :: phs_config
    integer, intent(in), optional :: unit
    call phs_config%base_startup_message (unit)
    write (msg_buffer, "(A,2(1x,I0,1x,A))") &
         "Phase space: single-particle"
    call msg_message (unit = unit)
  end subroutine phs_single_config_startup_message
    
\eatline
\nwendcode{}\nwbegindocs{246}Allocate an instance: the actual phase-space object.
\nwenddocs{}\nwbegincode{247}\moddef{PHS single: phs single config: TBP}\plusendmoddef
  procedure, nopass :: allocate_instance => phs_single_config_allocate_instance
\nwendcode{}\nwbegincode{248}\moddef{PHS single: procedures}\plusendmoddef
  subroutine phs_single_config_allocate_instance (phs)
    class(phs_t), intent(inout), pointer :: phs
    allocate (phs_single_t :: phs)
  end subroutine phs_single_config_allocate_instance
  
\eatline
\nwendcode{}\nwbegindocs{249}\nwdocspar
\subsection{Kinematics implementation}
We generate $\cos\theta$ and $\phi$ uniformly, covering the solid angle.

Note: The incoming momenta must be in the c.m. system.
\nwenddocs{}\nwbegincode{250}\moddef{PHS single: public}\plusendmoddef
  public :: phs_single_t
\nwendcode{}\nwbegincode{251}\moddef{PHS single: types}\plusendmoddef
  type, extends (phs_t) :: phs_single_t
   contains
   \LA{}PHS single: phs single: TBP\RA{}
  end type phs_single_t
  
\eatline
\nwendcode{}\nwbegindocs{252}Output.  The {\tt{}verbose} setting is irrelevant, we just display the contents
of the base object.
\nwenddocs{}\nwbegincode{253}\moddef{PHS single: phs single: TBP}\endmoddef
  procedure :: write => phs_single_write
\nwendcode{}\nwbegincode{254}\moddef{PHS single: procedures}\plusendmoddef
  subroutine phs_single_write (object, unit, verbose)
    class(phs_single_t), intent(in) :: object
    integer, intent(in), optional :: unit
    logical, intent(in), optional :: verbose
    integer :: u
    u = given_output_unit (unit)
    call object%base_write (u)
  end subroutine phs_single_write
    
\eatline
\nwendcode{}\nwbegindocs{255}The finalizer is empty.
\nwenddocs{}\nwbegincode{256}\moddef{PHS single: phs single: TBP}\plusendmoddef
  procedure :: final => phs_single_final
\nwendcode{}\nwbegincode{257}\moddef{PHS single: procedures}\plusendmoddef
  subroutine phs_single_final (object)
    class(phs_single_t), intent(inout) :: object
  end subroutine phs_single_final
  
\eatline
\nwendcode{}\nwbegindocs{258}Initialization.  We allocate arrays ({\tt{}base{\char95}init}) and adjust the
phase-space volume.  The massless two-particle phase space volume is
\begin{equation}
  \Phi_2 = \frac{1}{4(2\pi)^5} = 2.55294034614 \times 10^{-5}
\end{equation}
For a decay with nonvanishing masses ($m_3$, $m_4$), there is a correction
factor 
\begin{equation}
  \Phi_2(m) / \Phi_2(0) = \frac{1}{\hat s}
  \lambda^{1/2}(\hat s, m_3^2, m_4^2).
\end{equation}
For a scattering process with nonvanishing masses, the correction
factor is
\begin{equation}
  \Phi_2(m) / \Phi_2(0) = \frac{1}{\hat s ^ 2}
  \lambda^{1/2}(\hat s, m_1^2, m_2^2)\,
  \lambda^{1/2}(\hat s, m_3^2, m_4^2).
\end{equation}
If the energy is fixed, this is constant.  Otherwise, we have to account for
varying $\hat s$.
\nwenddocs{}\nwbegincode{259}\moddef{PHS single: phs single: TBP}\plusendmoddef
  procedure :: init => phs_single_init
\nwendcode{}\nwbegincode{260}\moddef{PHS single: procedures}\plusendmoddef
  subroutine phs_single_init (phs, phs_config)
    class(phs_single_t), intent(out) :: phs
    class(phs_config_t), intent(in), target :: phs_config
    call phs%base_init (phs_config)
    phs%volume = 1 / (4 * twopi5)
    call phs%compute_factor ()
  end subroutine phs_single_init
  
\eatline
\nwendcode{}\nwbegindocs{261}Compute the correction factor for nonzero masses.  We do this during
initialization (when the incoming momenta {\tt{}p} are undefined), unless
{\tt{}sqrts} is variable.  We do this again once for each phase-space point, but
then we skip the calculation if {\tt{}sqrts} is fixed.
\nwenddocs{}\nwbegincode{262}\moddef{PHS single: phs single: TBP}\plusendmoddef
  procedure :: compute_factor => phs_single_compute_factor
\nwendcode{}\nwbegincode{263}\moddef{PHS single: procedures}\plusendmoddef
  subroutine phs_single_compute_factor (phs)
    class(phs_single_t), intent(inout) :: phs
    real(default) :: s_hat
    select case (phs%config%n_in)
    case (1)
       if (.not. phs%p_defined) then
          if (sum (phs%m_out) < phs%m_in(1)) then
             s_hat = phs%m_in(1) ** 2
             phs%f(1) = 1 / s_hat &
                  * sqrt (lambda (s_hat, phs%m_out(1)**2, phs%m_out(2)**2))
          else
             print *, "m_in  = ", phs%m_in
             print *, "m_out = ", phs%m_out
             call msg_fatal ("Decay is kinematically forbidden")
          end if
       end if
    case (2)
       if (phs%config%sqrts_fixed) then
          if (phs%p_defined)  return
          s_hat = phs%config%sqrts ** 2
       else
          if (.not. phs%p_defined)  return
          s_hat = sum (phs%p) ** 2
       end if
       if (sum (phs%m_in)**2 < s_hat .and. sum (phs%m_out)**2 < s_hat) then
          phs%f(1) = 1 / s_hat * &
               ( lambda (s_hat, phs%m_in (1)**2, phs%m_in (2)**2)   &
               * lambda (s_hat, phs%m_out(1)**2, phs%m_out(2)**2) ) &
               ** 0.25_default
       else
          phs%f(1) = 0
       end if
    end select
  end subroutine phs_single_compute_factor
    
\eatline
\nwendcode{}\nwbegindocs{264}Evaluation.  We uniformly generate $\cos\theta$ and
$\phi$, with Jacobian normalized to one.

There is only a single channel, so the second subroutine does nothing.

Note: the current implementation works for elastic scattering only.
\nwenddocs{}\nwbegincode{265}\moddef{PHS single: phs single: TBP}\plusendmoddef
  procedure :: evaluate_selected_channel => phs_single_evaluate_selected_channel
  procedure :: evaluate_other_channels => phs_single_evaluate_other_channels
\nwendcode{}\nwbegincode{266}\moddef{PHS single: procedures}\plusendmoddef
  subroutine phs_single_evaluate_selected_channel (phs, c_in, r_in)
    class(phs_single_t), intent(inout) :: phs
    integer, intent(in) :: c_in
    real(default), intent(in), dimension(:) :: r_in
    !!! Catching a gfortran bogus warning
    type(vector4_t), dimension(2) :: p_dum
    if (phs%p_defined) then
       call phs%select_channel (c_in)
       phs%r(:,c_in) = r_in
       select case (phs%config%n_in)
       case (2)
          if (all (phs%m_in == phs%m_out)) then
             call compute_kinematics_solid_angle (phs%p, phs%q, r_in) 
          else
             call msg_bug ("PHS single: inelastic scattering not implemented")
          end if
       case (1)
          !!! Catching a gfortran bogus warning
          !!! call compute_kinematics_solid_angle (phs%decay_p (), phs%q, x)
          p_dum = phs%decay_p ()
          call compute_kinematics_solid_angle (p_dum, phs%q, r_in)
       end select
       call phs%compute_factor ()
       phs%q_defined = .true.
       phs%r_defined = .true.
    end if
  end subroutine phs_single_evaluate_selected_channel
  
  subroutine phs_single_evaluate_other_channels (phs, c_in)
    class(phs_single_t), intent(inout) :: phs
    integer, intent(in) :: c_in
  end subroutine phs_single_evaluate_other_channels
  
\eatline
\eatline
\nwendcode{}\nwbegindocs{267}Auxiliary: split a decaying particle at rest into the decay products,
aligned along the $z$ axis.
\nwenddocs{}\nwbegincode{268}\moddef{PHS single: phs single: TBP}\plusendmoddef
  procedure :: decay_p => phs_single_decay_p
\nwendcode{}\nwbegincode{269}\moddef{PHS single: procedures}\plusendmoddef
  function phs_single_decay_p (phs) result (p)
    class(phs_single_t), intent(in) :: phs
    type(vector4_t), dimension(2) :: p
    real(default) :: k
    real(default), dimension(2) :: E
    k = sqrt (lambda (phs%m_in(1) ** 2, phs%m_out(1) ** 2, phs%m_out(2) ** 2)) &
         / (2 * phs%m_in(1))
    E = sqrt (phs%m_out ** 2 + k ** 2)
    p(1) = vector4_moving (E(1), k, 3)
    p(2) = vector4_moving (E(2),-k, 3)
  end function phs_single_decay_p
  
\eatline
\nwendcode{}\nwbegindocs{270}Inverse evaluation.
\nwenddocs{}\nwbegincode{271}\moddef{PHS single: phs single: TBP}\plusendmoddef
  procedure :: inverse => phs_single_inverse
\nwendcode{}\nwbegincode{272}\moddef{PHS single: procedures}\plusendmoddef
  subroutine phs_single_inverse (phs)
    class(phs_single_t), intent(inout) :: phs
    real(default), dimension(:), allocatable :: x
    if (phs%p_defined .and. phs%q_defined) then
       call phs%select_channel ()
       allocate (x (phs%config%n_par))
       call inverse_kinematics_solid_angle (phs%p, phs%q, x)
       phs%r(:,1) = x
       call phs%compute_factor ()
       phs%r_defined = .true.
    end if
  end subroutine phs_single_inverse
  
\eatline
\nwendcode{}\nwbegindocs{273}\nwdocspar
\subsection{Unit tests}
\nwenddocs{}\nwbegincode{274}\moddef{PHS single: public}\plusendmoddef
  public :: phs_single_test
\nwendcode{}\nwbegincode{275}\moddef{PHS single: tests}\endmoddef
  subroutine phs_single_test (u, results)
    integer, intent(in) :: u
    type(test_results_t), intent(inout) :: results
  \LA{}PHS single: execute tests\RA{}
  end subroutine phs_single_test
  
\eatline
\nwendcode{}\nwbegindocs{276}\nwdocspar
\subsubsection{Phase-space configuration data}
Construct and display a test phase-space configuration object.  Also
check the {\tt{}azimuthal{\char95}dependence} flag.
\nwenddocs{}\nwbegincode{277}\moddef{PHS single: execute tests}\endmoddef
  call test (phs_single_1, "phs_single_1", &
       "phase-space configuration", &
       u, results)
\nwendcode{}\nwbegincode{278}\moddef{PHS single: tests}\plusendmoddef
  subroutine phs_single_1 (u)
    integer, intent(in) :: u
    type(model_data_t), target :: model
    type(process_constants_t) :: process_data
    class(phs_config_t), allocatable :: phs_data
    real(default) :: sqrts
    
    write (u, "(A)")  "* Test output: phs_single_1"
    write (u, "(A)")  "*   Purpose: initialize and display &
         &phase-space configuration data"
    write (u, "(A)")
    
    call model%init_test ()

    write (u, "(A)")  "* Initialize a process and a matching &
         &phase-space configuration"
    write (u, "(A)")

    call init_test_process_data (var_str ("phs_single_1"), process_data)

    allocate (phs_single_config_t :: phs_data)
    call phs_data%init (process_data, model)

    sqrts = 1000._default
    call phs_data%configure (sqrts, azimuthal_dependence=.false.)
       
    call phs_data%write (u)
  
    call phs_data%final ()
    call model%final ()
    
    write (u, "(A)")
    write (u, "(A)")  "* Test output end: phs_single_1"

  end subroutine phs_single_1

\eatline
\nwendcode{}\nwbegindocs{279}\nwdocspar
\subsubsection{Phase space evaluation}
Compute kinematics for given parameters, also invert the calculation.
\nwenddocs{}\nwbegincode{280}\moddef{PHS single: execute tests}\plusendmoddef
  call test (phs_single_2, "phs_single_2", &
       "phase-space evaluation", &
       u, results)
\nwendcode{}\nwbegincode{281}\moddef{PHS single: tests}\plusendmoddef
  subroutine phs_single_2 (u)
    integer, intent(in) :: u
    type(model_data_t), target :: model
    type(flavor_t) :: flv
    type(process_constants_t) :: process_data
    real(default) :: sqrts, E
    class(phs_config_t), allocatable, target :: phs_data
    class(phs_t), pointer :: phs => null ()
    type(vector4_t), dimension(2) :: p, q
    
    write (u, "(A)")  "* Test output: phs_single_2"
    write (u, "(A)")  "*   Purpose: test simple two-channel phase space"
    write (u, "(A)")
    
    call model%init_test ()
    call flv%init (25, model)

    write (u, "(A)")  "* Initialize a process and a matching &
         &phase-space configuration"
    write (u, "(A)")

    call init_test_process_data (var_str ("phs_single_2"), process_data)

    allocate (phs_single_config_t :: phs_data)
    call phs_data%init (process_data, model)

    sqrts = 1000._default
    call phs_data%configure (sqrts)

    call phs_data%write (u)
    
    write (u, "(A)")
    write (u, "(A)")  "* Initialize the phase-space instance"
    write (u, "(A)")

    call phs_data%allocate_instance (phs)
    call phs%init (phs_data)
       
    call phs%write (u, verbose=.true.)
    
    write (u, "(A)")
    write (u, "(A)")  "* Set incoming momenta"
    write (u, "(A)")

    E = sqrts / 2
    p(1) = vector4_moving (E, sqrt (E**2 - flv%get_mass ()**2), 3)
    p(2) = vector4_moving (E,-sqrt (E**2 - flv%get_mass ()**2), 3)

    call phs%set_incoming_momenta (p)
    call phs%compute_flux ()
    call phs%write (u)
    
    write (u, "(A)")
    write (u, "(A)")  "* Compute phase-space point &
         &for x = 0.5, 0.125"
    write (u, "(A)")

    call phs%evaluate_selected_channel (1, [0.5_default, 0.125_default])
    call phs%evaluate_other_channels (1)
    call phs%write (u)
    
    write (u, "(A)")
    write (u, "(A)")  "* Inverse kinematics"
    write (u, "(A)")

    call phs%get_outgoing_momenta (q)
    deallocate (phs)
    call phs_data%allocate_instance (phs)
    call phs%init (phs_data)
       
    sqrts = 1000._default
    call phs_data%configure (sqrts)

    call phs%set_incoming_momenta (p)
    call phs%compute_flux ()
    call phs%set_outgoing_momenta (q)
    
    call phs%inverse ()
    call phs%write (u)
    
    call phs%final ()
    deallocate (phs)
    
    call phs_data%final ()
    call model%final ()
    
    write (u, "(A)")
    write (u, "(A)")  "* Test output end: phs_single_2"

  end subroutine phs_single_2

\eatline
\nwendcode{}\nwbegindocs{282}\nwdocspar
\subsubsection{Phase space for non-c.m. system}
Compute kinematics for given parameters, also invert the calculation.
Since this will involve cancellations, we call {\tt{}pacify} to eliminate
numerical noise.
\nwenddocs{}\nwbegincode{283}\moddef{PHS single: execute tests}\plusendmoddef
  call test (phs_single_3, "phs_single_3", &
       "phase-space evaluation in lab frame", &
       u, results)
\nwendcode{}\nwbegincode{284}\moddef{PHS single: tests}\plusendmoddef
  subroutine phs_single_3 (u)
    integer, intent(in) :: u
    type(model_data_t), target :: model
    type(flavor_t) :: flv
    type(process_constants_t) :: process_data
    real(default) :: sqrts, E
    class(phs_config_t), allocatable, target :: phs_data
    class(phs_t), pointer :: phs => null ()
    type(vector4_t), dimension(2) :: p, q
    type(lorentz_transformation_t) :: lt
    
    write (u, "(A)")  "* Test output: phs_single_3"
    write (u, "(A)")  "*   Purpose: test simple two-channel phase space"
    write (u, "(A)")  "*            without c.m. kinematics assumption"
    write (u, "(A)")
    
    call model%init_test ()
    call flv%init (25, model)

    write (u, "(A)")  "* Initialize a process and a matching &
         &phase-space configuration"
    write (u, "(A)")

    call init_test_process_data (var_str ("phs_single_3"), process_data)

    allocate (phs_single_config_t :: phs_data)
    call phs_data%init (process_data, model)

    sqrts = 1000._default
    call phs_data%configure (sqrts, cm_frame=.false., sqrts_fixed=.false.)

    call phs_data%write (u)
    
    write (u, "(A)")
    write (u, "(A)")  "* Initialize the phase-space instance"
    write (u, "(A)")

    call phs_data%allocate_instance (phs)
    call phs%init (phs_data)
       
    call phs%write (u, verbose=.true.)
    
    write (u, "(A)")
    write (u, "(A)")  "* Set incoming momenta in lab system"
    write (u, "(A)")

    lt = boost (0.1_default, 1) * boost (0.3_default, 3)
    
    E = sqrts / 2
    p(1) = lt * vector4_moving (E, sqrt (E**2 - flv%get_mass ()**2), 3)
    p(2) = lt * vector4_moving (E,-sqrt (E**2 - flv%get_mass ()**2), 3)

    call vector4_write (p(1), u)
    call vector4_write (p(2), u)

    write (u, "(A)")
    write (u, "(A)")  "* Compute phase-space point &
         &for x = 0.5, 0.125"
    write (u, "(A)")

    call phs%set_incoming_momenta (p)
    call phs%compute_flux ()
    
    call phs%evaluate_selected_channel (1, [0.5_default, 0.125_default])
    call phs%evaluate_other_channels (1)
    call pacify (phs)
    call phs%write (u)
    
    write (u, "(A)")
    write (u, "(A)")  "* Extract outgoing momenta in lab system"
    write (u, "(A)")
    
    call phs%get_outgoing_momenta (q)
    call vector4_write (q(1), u)
    call vector4_write (q(2), u)

    write (u, "(A)")
    write (u, "(A)")  "* Inverse kinematics"
    write (u, "(A)")

    deallocate (phs)
    call phs_data%allocate_instance (phs)
    call phs%init (phs_data)
       
    sqrts = 1000._default
    call phs_data%configure (sqrts)

    call phs%set_incoming_momenta (p)
    call phs%compute_flux ()
    call phs%set_outgoing_momenta (q)
    
    call phs%inverse ()
    call pacify (phs)
    call phs%write (u)
    
    call phs%final ()
    deallocate (phs)
    
    call phs_data%final ()
    call model%final ()
    
    write (u, "(A)")
    write (u, "(A)")  "* Test output end: phs_single_3"

  end subroutine phs_single_3

\eatline
\nwendcode{}\nwbegindocs{285}\nwdocspar
\subsubsection{Decay Phase space evaluation}
Compute kinematics for given parameters, also invert the calculation.  This
time, implement a decay process.
\nwenddocs{}\nwbegincode{286}\moddef{PHS single: execute tests}\plusendmoddef
  call test (phs_single_4, "phs_single_4", &
       "decay phase-space evaluation", &
       u, results)
\nwendcode{}\nwbegincode{287}\moddef{PHS single: tests}\plusendmoddef
  subroutine phs_single_4 (u)
    integer, intent(in) :: u
    type(model_data_t), target :: model
    type(flavor_t) :: flv
    type(process_constants_t) :: process_data
    class(phs_config_t), allocatable, target :: phs_data
    class(phs_t), pointer :: phs => null ()
    type(vector4_t), dimension(1) :: p
    type(vector4_t), dimension(2) :: q
    
    write (u, "(A)")  "* Test output: phs_single_4"
    write (u, "(A)")  "*   Purpose: test simple two-channel phase space"
    write (u, "(A)")
    
    call model%init_test ()

    call model%set_par (var_str ("ff"), 0.4_default)
    call model%set_par (var_str ("mf"), &
         model%get_real (var_str ("ff")) * model%get_real (var_str ("ms")))
    call flv%init (25, model)

    write (u, "(A)")  "* Initialize a decay and a matching &
         &phase-space configuration"
    write (u, "(A)")

    call init_test_decay_data (var_str ("phs_single_4"), process_data)

    allocate (phs_single_config_t :: phs_data)
    call phs_data%init (process_data, model)

    call phs_data%configure (flv%get_mass ())

    call phs_data%write (u)
    
    write (u, "(A)")
    write (u, "(A)")  "* Initialize the phase-space instance"
    write (u, "(A)")

    call phs_data%allocate_instance (phs)
    call phs%init (phs_data)
       
    call phs%write (u, verbose=.true.)
    
    write (u, "(A)")
    write (u, "(A)")  "* Set incoming momenta"
    write (u, "(A)")

    p(1) = vector4_at_rest (flv%get_mass ())

    call phs%set_incoming_momenta (p)
    call phs%compute_flux ()
    call phs%write (u)
    
    write (u, "(A)")
    write (u, "(A)")  "* Compute phase-space point &
         &for x = 0.5, 0.125"
    write (u, "(A)")

    call phs%evaluate_selected_channel (1, [0.5_default, 0.125_default])
    call phs%evaluate_other_channels (1)
    call phs%write (u)
    
    write (u, "(A)")
    write (u, "(A)")  "* Inverse kinematics"
    write (u, "(A)")

    call phs%get_outgoing_momenta (q)
    deallocate (phs)
    call phs_data%allocate_instance (phs)
    call phs%init (phs_data)
       
    call phs_data%configure (flv%get_mass ())

    call phs%set_incoming_momenta (p)
    call phs%compute_flux ()
    call phs%set_outgoing_momenta (q)
    
    call phs%inverse ()
    call phs%write (u)
    
    call phs%final ()
    deallocate (phs)
    
    call phs_data%final ()
    call model%final ()
    
    write (u, "(A)")
    write (u, "(A)")  "* Test output end: phs_single_4"

  end subroutine phs_single_4

\eatline
\nwendcode{}\nwbegindocs{288}\nwdocspar
%%%%%%%%%%%%%%%%%%%%%%%%%%%%%%%%%%%%%%%%%%%%%%%%%%%%%%%%%%%%%%%%%%%%%%%%
\clearpage
\section{Mappings}

Mappings are objects that encode the transformation of the interval
$(0,1)$ to a physical variable $m^2$ or $\cos\theta$ (and back), as it
is used in the phase space parameterization.  The mapping objects
contain fixed parameters, the associated methods implement the mapping
and inverse mapping operations, including the computation of the
Jacobian (phase space factor).
\nwenddocs{}\nwbegincode{289}\moddef{\code{}mappings.f90\edoc{}}\endmoddef
\LA{}File header\RA{}

module mappings

\LA{}Use kinds\RA{}
  use kinds, only: TC
\LA{}Use strings\RA{}
  use io_units
  use constants, only: pi
  use format_defs, only: FMT_19
  use diagnostics
  use md5
  use model_data
  use flavors

\LA{}Standard module head\RA{}

\LA{}Mappings: public\RA{}

\LA{}Mappings: parameters\RA{}

\LA{}Mappings: types\RA{}

\LA{}Mappings: interfaces\RA{}

contains

\LA{}Mappings: procedures\RA{}

end module mappings
\eatline
\nwendcode{}\nwbegindocs{290}\nwdocspar
\subsection{Default parameters}
This type holds the default parameters, needed for setting the scale
in cases where no mass parameter is available.  The contents are public.
\nwenddocs{}\nwbegincode{291}\moddef{Mappings: public}\endmoddef
  public :: mapping_defaults_t
\nwendcode{}\nwbegincode{292}\moddef{Mappings: types}\endmoddef
  type :: mapping_defaults_t
     real(default) :: energy_scale = 10
     real(default) :: invariant_mass_scale = 10
     real(default) :: momentum_transfer_scale = 10
     logical :: step_mapping = .true.
     logical :: step_mapping_exp = .true.
     logical :: enable_s_mapping = .false.
   contains
   \LA{}Mappings: mapping defaults: TBP\RA{}
  end type mapping_defaults_t

\eatline
\nwendcode{}\nwbegindocs{293}Output.
\nwenddocs{}\nwbegincode{294}\moddef{Mappings: mapping defaults: TBP}\endmoddef
  procedure :: write => mapping_defaults_write
\nwendcode{}\nwbegincode{295}\moddef{Mappings: procedures}\endmoddef
  subroutine mapping_defaults_write (object, unit)
    class(mapping_defaults_t), intent(in) :: object
    integer, intent(in), optional :: unit
    integer :: u
    u = given_output_unit (unit)
    write (u, "(3x,A," // FMT_19 // ")") "energy scale  = ", &
         object%energy_scale
    write (u, "(3x,A," // FMT_19 // ")") "mass scale    = ", &
         object%invariant_mass_scale
    write (u, "(3x,A," // FMT_19 // ")") "q scale       = ", &
         object%momentum_transfer_scale
    write (u, "(3x,A,L1)") "step mapping     = ", &
         object%step_mapping
    write (u, "(3x,A,L1)") "step exp. mode   = ", &
         object%step_mapping_exp
    write (u, "(3x,A,L1)") "allow s mapping  = ", &
         object%enable_s_mapping
  end subroutine mapping_defaults_write    

\eatline
\nwendcode{}\nwbegindocs{296}\nwdocspar
\nwenddocs{}\nwbegincode{297}\moddef{Mappings: public}\plusendmoddef
  public :: mapping_defaults_md5sum
\nwendcode{}\nwbegincode{298}\moddef{Mappings: procedures}\plusendmoddef
  function mapping_defaults_md5sum (mapping_defaults) result (md5sum_map)
    character(32) :: md5sum_map
    type(mapping_defaults_t), intent(in) :: mapping_defaults
    integer :: u
    u = free_unit ()
    open (u, status = "scratch")
    write (u, *)  mapping_defaults%energy_scale
    write (u, *)  mapping_defaults%invariant_mass_scale
    write (u, *)  mapping_defaults%momentum_transfer_scale
    write (u, *)  mapping_defaults%step_mapping
    write (u, *)  mapping_defaults%step_mapping_exp
    write (u, *)  mapping_defaults%enable_s_mapping
    rewind (u)
    md5sum_map = md5sum (u)
    close (u)
  end function mapping_defaults_md5sum

\eatline
\nwendcode{}\nwbegindocs{299}\nwdocspar
\subsection{The Mapping type}
Each mapping has a type (e.g., s-channel, infrared), a binary code
(redundant, but useful for debugging), and a reference particle.  The
flavor code of this particle is stored for bookkeeping reasons, what
matters are the mass and width of this particle.  Furthermore,
depending on the type, various mapping parameters can be set and used.

The parameters {\tt{}a1} to {\tt{}a3} (for $m^2$ mappings) and {\tt{}b1} to
{\tt{}b3} (for $\cos\theta$ mappings) are values that are stored once to
speed up the calculation, if {\tt{}variable{\char95}limits} is false.  The exact
meaning of these parameters depends on the mapping type.  The limits
are fixed if there is a fixed c.m. energy.
\nwenddocs{}\nwbegincode{300}\moddef{Mappings: public}\plusendmoddef
  public :: mapping_t
\nwendcode{}\nwbegincode{301}\moddef{Mappings: types}\plusendmoddef
  type :: mapping_t
     private
     integer :: type = NO_MAPPING
     integer(TC) :: bincode
     type(flavor_t) :: flv
     real(default) :: mass = 0
     real(default) :: width = 0
     logical :: a_unknown = .true.
     real(default) :: a1 = 0
     real(default) :: a2 = 0
     real(default) :: a3 = 0
     logical :: b_unknown = .true.
     real(default) :: b1 = 0
     real(default) :: b2 = 0
     real(default) :: b3 = 0
     logical :: variable_limits = .true.
  end type mapping_t
  
\eatline
\nwendcode{}\nwbegindocs{302}The valid mapping types.  The extra type {\tt{}STEP{\char95}MAPPING} is used
only internally.
\nwenddocs{}\nwbegincode{303}\moddef{Mappings: parameters}\endmoddef
\LA{}Mapping modes\RA{}
\nwendcode{}\nwbegindocs{304}\nwdocspar
\subsection{Screen output}
Do not write empty mappings.
\nwenddocs{}\nwbegincode{305}\moddef{Mappings: public}\plusendmoddef
  public :: mapping_write
\nwendcode{}\nwbegincode{306}\moddef{Mappings: procedures}\plusendmoddef
  subroutine mapping_write (map, unit, verbose)
    type(mapping_t), intent(in) :: map
    integer, intent(in), optional :: unit
    logical, intent(in), optional :: verbose
    integer :: u
    character(len=9) :: str
    u = given_output_unit (unit);  if (u < 0)  return
    select case(map%type)
    case(S_CHANNEL); str = "s_channel"
    case(COLLINEAR); str = "collinear"
    case(INFRARED);  str = "infrared "
    case(RADIATION); str = "radiation"
    case(T_CHANNEL); str = "t_channel"
    case(U_CHANNEL); str = "u_channel"
    case(STEP_MAPPING_E);  str = "step_exp"
    case(STEP_MAPPING_H);  str = "step_hyp"
    case(ON_SHELL);  str = "on_shell"
    case default;    str = "????????"
    end select
    if (map%type /= NO_MAPPING) then
       write (u, '(1x,A,I4,A)') &
            "Branch #", map%bincode, ":  " // &
            "Mapping (" // str // ") for particle " // &
            '"' // char (map%flv%get_name ()) // '"'
       if (present (verbose)) then
          if (verbose) then
             select case (map%type)
             case (S_CHANNEL, RADIATION, STEP_MAPPING_E, STEP_MAPPING_H)
                write (u, "(1x,A,3(" // FMT_19 // "))")  &
                     "  m/w    = ", map%mass, map%width
             case default
                write (u, "(1x,A,3(" // FMT_19 // "))")  &
                     "  m      = ", map%mass
             end select
             select case (map%type)
             case (S_CHANNEL, T_CHANNEL, U_CHANNEL, &
                  STEP_MAPPING_E, STEP_MAPPING_H, &
                  COLLINEAR, INFRARED, RADIATION)
                write (u, "(1x,A,3(" // FMT_19 // "))")  &
                     "  a1/2/3 = ", map%a1, map%a2, map%a3
             end select
             select case (map%type)
             case (T_CHANNEL, U_CHANNEL, COLLINEAR)
                write (u, "(1x,A,3(" // FMT_19 // "))")  &
                     "  b1/2/3 = ", map%b1, map%b2, map%b3
             end select
          end if
       end if
    end if
  end subroutine mapping_write

\eatline
\nwendcode{}\nwbegindocs{307}\nwdocspar
\subsection{Define a mapping}
The initialization routine sets the mapping type and the particle
(binary code and flavor code) for which the mapping applies (e.g., a
$Z$ resonance in branch \#3).  We only need the absolute value of the
flavor code.
\nwenddocs{}\nwbegincode{308}\moddef{Mappings: public}\plusendmoddef
  public :: mapping_init
\nwendcode{}\nwbegincode{309}\moddef{Mappings: procedures}\plusendmoddef
  subroutine mapping_init (mapping, bincode, type, f, model)
    type(mapping_t), intent(inout) :: mapping
    integer(TC), intent(in) :: bincode
    type(string_t), intent(in) :: type
    integer, intent(in), optional :: f
    class(model_data_t), intent(in), optional, target :: model
    mapping%bincode = bincode
    select case (char (type))
    case ("s_channel");  mapping%type = S_CHANNEL
    case ("collinear");  mapping%type = COLLINEAR
    case ("infrared");   mapping%type = INFRARED
    case ("radiation");  mapping%type = RADIATION
    case ("t_channel");  mapping%type = T_CHANNEL
    case ("u_channel");  mapping%type = U_CHANNEL
    case ("step_exp");  mapping%type = STEP_MAPPING_E
    case ("step_hyp");  mapping%type = STEP_MAPPING_H
    case ("on_shell");  mapping%type = ON_SHELL
    case default
       call msg_bug ("Mappings: encountered undefined mapping key '" &
            // char (type) // "'")
    end select
    if (present (f) .and. present (model)) &
         call mapping%flv%init (abs (f), model)
  end subroutine mapping_init

\eatline
\nwendcode{}\nwbegindocs{310}This sets the actual mass and width, using a parameter set.  Since
the auxiliary parameters will only be determined when the mapping is
first called, they are marked as unknown.
\nwenddocs{}\nwbegincode{311}\moddef{Mappings: public}\plusendmoddef
  public :: mapping_set_parameters
\nwendcode{}\nwbegincode{312}\moddef{Mappings: procedures}\plusendmoddef
  subroutine mapping_set_parameters (map, mapping_defaults, variable_limits)
    type(mapping_t), intent(inout) :: map
    type(mapping_defaults_t), intent(in) :: mapping_defaults
    logical, intent(in) :: variable_limits
    if (map%type /= NO_MAPPING) then
       map%mass  = map%flv%get_mass ()
       map%width = map%flv%get_width ()
       map%variable_limits = variable_limits
       map%a_unknown = .true.
       map%b_unknown = .true.
       select case (map%type)
       case (S_CHANNEL)
          if (map%mass <= 0) then
             call mapping_write (map)
             call msg_fatal &
                  & (" S-channel resonance must have positive mass")
          else if (map%width <= 0) then
             call mapping_write (map)
             call msg_fatal &
                  & (" S-channel resonance must have positive width")
          end if
       case (RADIATION)
          map%width = max (map%width, mapping_defaults%energy_scale)
       case (INFRARED, COLLINEAR)
          map%mass = max (map%mass, mapping_defaults%invariant_mass_scale)
       case (T_CHANNEL, U_CHANNEL)
          map%mass = max (map%mass, mapping_defaults%momentum_transfer_scale)
       end select
    end if
  end subroutine mapping_set_parameters

\eatline
\nwendcode{}\nwbegindocs{313}For a step mapping the mass and width are set directly, instead of
being determined from the flavor parameter (which is meaningless
here).  They correspond to the effective upper bound of phase space
due to a resonance, as opposed to the absolute upper bound.
\nwenddocs{}\nwbegincode{314}\moddef{Mappings: public}\plusendmoddef
  public :: mapping_set_step_mapping_parameters
\nwendcode{}\nwbegincode{315}\moddef{Mappings: procedures}\plusendmoddef
  subroutine mapping_set_step_mapping_parameters (map, &
       mass, width, variable_limits)
    type(mapping_t), intent(inout) :: map
    real(default), intent(in) :: mass, width
    logical, intent(in) :: variable_limits
    select case (map%type)
    case (STEP_MAPPING_E, STEP_MAPPING_H)
       map%variable_limits = variable_limits
       map%a_unknown = .true.
       map%b_unknown = .true.
       map%mass = mass
       map%width = width
    end select
  end subroutine mapping_set_step_mapping_parameters

\eatline
\nwendcode{}\nwbegindocs{316}\nwdocspar
\subsection{Retrieve contents}
Return true if there is any / an s-channel mapping.
\nwenddocs{}\nwbegincode{317}\moddef{Mappings: public}\plusendmoddef
  public :: mapping_is_set
  public :: mapping_is_s_channel
  public :: mapping_is_on_shell
\nwendcode{}\nwbegincode{318}\moddef{Mappings: procedures}\plusendmoddef
  function mapping_is_set (mapping) result (flag)
    logical :: flag
    type(mapping_t), intent(in) :: mapping
    flag = mapping%type /= NO_MAPPING
  end function mapping_is_set

  function mapping_is_s_channel (mapping) result (flag)
    logical :: flag
    type(mapping_t), intent(in) :: mapping
    flag = mapping%type == S_CHANNEL
  end function mapping_is_s_channel

  function mapping_is_on_shell (mapping) result (flag)
    logical :: flag
    type(mapping_t), intent(in) :: mapping
    flag = mapping%type == ON_SHELL
  end function mapping_is_on_shell

\eatline
\eatline
\eatline
\nwendcode{}\nwbegindocs{319}Return stored mass and width, respectively.
\nwenddocs{}\nwbegincode{320}\moddef{Mappings: public}\plusendmoddef
  public :: mapping_get_mass
  public :: mapping_get_width
\nwendcode{}\nwbegincode{321}\moddef{Mappings: procedures}\plusendmoddef
  function mapping_get_mass (mapping) result (mass)
    real(default) :: mass
    type(mapping_t), intent(in) :: mapping
    mass = mapping%mass
  end function mapping_get_mass

  function mapping_get_width (mapping) result (width)
    real(default) :: width
    type(mapping_t), intent(in) :: mapping
    width = mapping%width
  end function mapping_get_width

\eatline
\eatline
\nwendcode{}\nwbegindocs{322}\nwdocspar
\subsection{Compare mappings}
Equality for single mappings and arrays
\nwenddocs{}\nwbegincode{323}\moddef{Mappings: public}\plusendmoddef
  public :: operator(==)
\nwendcode{}\nwbegincode{324}\moddef{Mappings: interfaces}\endmoddef
  interface operator(==)
     module procedure mapping_equal
  end interface
\nwendcode{}\nwbegincode{325}\moddef{Mappings: procedures}\plusendmoddef
  function mapping_equal (m1, m2) result (equal)
    type(mapping_t), intent(in) :: m1, m2
    logical :: equal
    if (m1%type == m2%type) then
       select case (m1%type)
       case (NO_MAPPING)
          equal = .true.
       case (S_CHANNEL, RADIATION, STEP_MAPPING_E, STEP_MAPPING_H)
          equal = (m1%mass == m2%mass) .and. (m1%width == m2%width)
       case default
          equal = (m1%mass == m2%mass)
       end select
    else
       equal = .false.
    end if
  end function mapping_equal

\eatline
\nwendcode{}\nwbegindocs{326}\nwdocspar
\subsection{Mappings of the invariant mass}
Inserting an $x$ value between 0 and 1, we want to compute the
corresponding invariant mass $m^2(x)$ and the jacobian, aka phase
space factor $f(x)$.  We also need the reverse operation.

In general, the phase space factor $f$ is defined by
\begin{equation}
  \frac{1}{s}\int_{m^2_{\textrm{min}}}^{m^2_{\textrm{max}}} dm^2\,g(m^2)
  = \int_0^1 dx\,\frac{1}{s}\,\frac{dm^2}{dx}\,g(m^2(x))
  = \int_0^1 dx\,f(x)\,g(x),
\end{equation}
where thus
\begin{equation}
  f(x) = \frac{1}{s}\,\frac{dm^2}{dx}.
\end{equation}
With this mapping, a function of the form
\begin{equation}
  g(m^2) = c\frac{dx(m^2)}{dm^2}
\end{equation}
is mapped to a constant:
\begin{equation}
  \frac{1}{s}\int_{m^2_{\textrm{min}}}^{m^2_{\textrm{max}}} dm^2\,g(m^2)
  = \int_0^1 dx\,f(x)\,g(m^2(x)) = \int_0^1 dx\,\frac{c}{s}.
\end{equation}

Here is the mapping routine.  Input are the available energy
squared {\tt{}s}, the limits for $m^2$, and the $x$ value.  Output are
the $m^2$ value and the phase space factor $f$.
\nwenddocs{}\nwbegincode{327}\moddef{Mappings: public}\plusendmoddef
  public :: mapping_compute_msq_from_x
\nwendcode{}\nwbegincode{328}\moddef{Mappings: procedures}\plusendmoddef
  subroutine mapping_compute_msq_from_x (map, s, msq_min, msq_max, msq, f, x)
    type(mapping_t), intent(inout) :: map
    real(default), intent(in) :: s, msq_min, msq_max
    real(default), intent(out) :: msq, f
    real(default), intent(in) :: x
    real(default) :: z, msq0, msq1, tmp
    integer :: type
    type = map%type
    if (s == 0) &
         call msg_fatal (" Applying msq mapping for zero energy")
  \LA{}Modify mapping type if necessary\RA{}
    select case(type)
    case (NO_MAPPING)
     \LA{}Constants for trivial msq mapping\RA{}
     \LA{}Apply trivial msq mapping\RA{}
    case (S_CHANNEL)
     \LA{}Constants for s-channel resonance mapping\RA{}
     \LA{}Apply s-channel resonance mapping\RA{}
    case (COLLINEAR, INFRARED, RADIATION)
     \LA{}Constants for s-channel pole mapping\RA{}
     \LA{}Apply s-channel pole mapping\RA{}
    case (T_CHANNEL, U_CHANNEL)
     \LA{}Constants for t-channel pole mapping\RA{}
     \LA{}Apply t-channel pole mapping\RA{}
    case (STEP_MAPPING_E)
     \LA{}Constants for exponential step mapping\RA{}
     \LA{}Apply exponential step mapping\RA{}
    case (STEP_MAPPING_H)
     \LA{}Constants for hyperbolic step mapping\RA{}
     \LA{}Apply hyperbolic step mapping\RA{}
    case default
       call msg_fatal ( " Attempt to apply undefined msq mapping")
    end select
  end subroutine mapping_compute_msq_from_x

\eatline
\nwendcode{}\nwbegindocs{329}The inverse mapping
\nwenddocs{}\nwbegincode{330}\moddef{Mappings: public}\plusendmoddef
  public :: mapping_compute_x_from_msq
\nwendcode{}\nwbegincode{331}\moddef{Mappings: procedures}\plusendmoddef
  subroutine mapping_compute_x_from_msq (map, s, msq_min, msq_max, msq, f, x)
    type(mapping_t), intent(inout) :: map
    real(default), intent(in) :: s, msq_min, msq_max
    real(default), intent(in) :: msq
    real(default), intent(out) :: f, x
    real(default) :: msq0, msq1, tmp, z
    integer :: type
    type = map%type
    if (s == 0) &
         call msg_fatal (" Applying inverse msq mapping for zero energy")
  \LA{}Modify mapping type if necessary\RA{}
    select case (type)
    case (NO_MAPPING)
     \LA{}Constants for trivial msq mapping\RA{}
     \LA{}Apply inverse trivial msq mapping\RA{}
    case (S_CHANNEL)
     \LA{}Constants for s-channel resonance mapping\RA{}
     \LA{}Apply inverse s-channel resonance mapping\RA{}
    case (COLLINEAR, INFRARED, RADIATION)
     \LA{}Constants for s-channel pole mapping\RA{}
     \LA{}Apply inverse s-channel pole mapping\RA{}
    case (T_CHANNEL, U_CHANNEL)
     \LA{}Constants for t-channel pole mapping\RA{}
     \LA{}Apply inverse t-channel pole mapping\RA{}
    case (STEP_MAPPING_E)
     \LA{}Constants for exponential step mapping\RA{}
     \LA{}Apply inverse exponential step mapping\RA{}
    case (STEP_MAPPING_H)
     \LA{}Constants for hyperbolic step mapping\RA{}
     \LA{}Apply inverse hyperbolic step mapping\RA{}
    case default
       call msg_fatal ( " Attempt to apply undefined msq mapping")
    end select
  end subroutine mapping_compute_x_from_msq

\eatline
\nwendcode{}\nwbegindocs{332}\nwdocspar
\subsubsection{Trivial mapping}
We simply map the boundaries of the interval $(m_{\textrm{min}},
m_{\textrm{max}})$ to $(0,1)$:
\begin{equation}
  m^2 = (1-x) m_{\textrm{min}}^2 + x m_{\textrm{max}}^2;
\end{equation}
the inverse is
\begin{equation}
  x = \frac{m^2 - m_{\textrm{min}}^2}{m_{\textrm{max}}^2- m_{\textrm{min}}^2}.
\end{equation}
Hence
\begin{equation}
  f(x) = \frac{m_{\textrm{max}}^2 - m_{\textrm{min}}^2}{s},
\end{equation}
and we have, as required,
\begin{equation}
  f(x)\,\frac{dx}{dm^2} = \frac{1}{s}.
\end{equation}

We store the constant parameters the first time the mapping is called
-- or, if limits vary, recompute them each time.
\nwenddocs{}\nwbegincode{333}\moddef{Constants for trivial msq mapping}\endmoddef
  if (map%variable_limits .or. map%a_unknown) then
     map%a1 = 0
     map%a2 = msq_max - msq_min
     map%a3 = map%a2 / s
     map%a_unknown = .false.
  end if
\nwendcode{}\nwbegincode{334}\moddef{Apply trivial msq mapping}\endmoddef
  msq = (1-x) * msq_min + x * msq_max
  f = map%a3
\nwendcode{}\nwbegincode{335}\moddef{Apply inverse trivial msq mapping}\endmoddef
  if (map%a2 /= 0) then
     x = (msq - msq_min) / map%a2
  else
     x = 0
  end if
  f = map%a3
\nwendcode{}\nwbegindocs{336}Resonance or step mapping does not make much sense if the resonance mass is
outside the kinematical bounds.  If this is the case, revert to
{\tt{}NO{\char95}MAPPING}.  This is possible even if the kinematical bounds vary
from event to event.
\nwenddocs{}\nwbegincode{337}\moddef{Modify mapping type if necessary}\endmoddef
  select case (type)
  case (S_CHANNEL, STEP_MAPPING_E, STEP_MAPPING_H)
     msq0 = map%mass**2
     if (msq0 < msq_min .or. msq0 > msq_max)  type = NO_MAPPING
  end select
\nwendcode{}\nwbegindocs{338}\nwdocspar
\subsubsection{Breit-Wigner mapping}
A Breit-Wigner resonance with mass $M$ and width $\Gamma$ is flattened
by the following mapping:

This mapping does not make much sense if the resonance mass is too low.
If this is the case, revert to {\tt{}NO{\char95}MAPPING}.  There is a tricky
point with this if the mass is too high: {\tt{}msq{\char95}max} is not a
constant if structure functions are around.  However, switching the
type depending on the overall energy does not change the integral, it
is just another branching point.
\begin{equation}
  m^2 = M(M+t\Gamma),
\end{equation}
where
\begin{equation}
  t = \tan\left[(1-x)\arctan\frac{m^2_{\textrm{min}} - M^2}{M\Gamma}
                + x  \arctan\frac{m^2_{\textrm{max}} - M^2}{M\Gamma}\right].
\end{equation}
The inverse:
\begin{equation}
  x = \frac{  \arctan\frac{m^2 - M^2}{M\Gamma}
            - \arctan\frac{m^2_{\textrm{min}} - M^2}{M\Gamma}}
           {  \arctan\frac{m^2_{\textrm{max}} - M^2}{M\Gamma}
            - \arctan\frac{m^2_{\textrm{min}} - M^2}{M\Gamma}}
\end{equation}
The phase-space factor of this transformation is
\begin{equation}
  f(x) = \frac{M\Gamma}{s}\left(
           \arctan\frac{m^2_{\textrm{max}} - M^2}{M\Gamma}
         - \arctan\frac{m^2_{\textrm{min}} - M^2}{M\Gamma}\right)
         (1 + t^2).
\end{equation}
This maps any function proportional to
\begin{equation}
  g(m^2) = \frac{M\Gamma}{(m^2-M^2)^2 + M^2\Gamma^2}
\end{equation}
to a constant times $1/s$.
\nwenddocs{}\nwbegincode{339}\moddef{Constants for s-channel resonance mapping}\endmoddef
  if (map%variable_limits .or. map%a_unknown) then
     msq0 = map%mass ** 2
     map%a1 = atan ((msq_min - msq0) / (map%mass * map%width))
     map%a2 = atan ((msq_max - msq0) / (map%mass * map%width))
     map%a3 = (map%a2 - map%a1) * (map%mass * map%width) / s 
     map%a_unknown = .false.
  end if
\nwendcode{}\nwbegincode{340}\moddef{Apply s-channel resonance mapping}\endmoddef
  z = (1-x) * map%a1 + x * map%a2
  if (-pi/2 < z .and. z < pi/2) then
     tmp = tan (z)
     msq = map%mass * (map%mass + map%width * tmp)
     f = map%a3 * (1 + tmp**2) 
  else
     msq = 0
     f = 0
  end if
\nwendcode{}\nwbegincode{341}\moddef{Apply inverse s-channel resonance mapping}\endmoddef
  tmp = (msq - msq0) / (map%mass * map%width)
  x = (atan (tmp) - map%a1) / (map%a2 - map%a1)
  f = map%a3 * (1 + tmp**2)
\nwendcode{}\nwbegindocs{342}\nwdocspar
\subsubsection{Mapping for massless splittings}
This mapping accounts for approximately scale-invariant behavior where
$\ln M^2$ is evenly distributed.
\begin{equation}
  m^2 = m_{\textrm{min}}^2 + M^2\left(\exp(xL)-1\right)
\end{equation}
where
\begin{equation}
  L = \ln\left(\frac{m_{\textrm{max}}^2 - m_{\textrm{min}}^2}{M^2} + 1\right).
\end{equation}
The inverse:
\begin{equation}
  x = \frac1L\ln\left(\frac{m^2-m_{\textrm{min}}^2}{M^2} + 1\right)
\end{equation}
The constant $M$ is a characteristic scale.  Above this scale
($m^2-m_{\textrm{min}}^2 \gg M^2$), this mapping behaves like
$x\propto\ln m^2$, while below the scale it reverts to a linear
mapping.

The phase-space factor is
\begin{equation}
  f(x) = \frac{M^2}{s}\,\exp(xL)\,L.
\end{equation}
A function proportional to
\begin{equation}
  g(m^2) = \frac{1}{(m^2-m_{\textrm{min}}^2) + M^2}
\end{equation}
is mapped to a constant, i.e., a simple pole near $m_{\textrm{min}}$
with a regulator mass $M$.

This type of mapping is useful for massless collinear and infrared
singularities, where the scale is stored as the mass parameter.  In
the radiation case (IR radiation off massive particle), the heavy
particle width is the characteristic scale.
\nwenddocs{}\nwbegincode{343}\moddef{Constants for s-channel pole mapping}\endmoddef
  if (map%variable_limits .or. map%a_unknown) then
     if (type == RADIATION) then
        msq0 = map%width**2
     else
        msq0 = map%mass**2
     end if
     map%a1 = msq0
     map%a2 = log ((msq_max - msq_min) / msq0 + 1)
     map%a3 = map%a2 / s
     map%a_unknown = .false.
  end if
\nwendcode{}\nwbegincode{344}\moddef{Apply s-channel pole mapping}\endmoddef
  msq1 = map%a1 * exp (x * map%a2)
  msq = msq1 - map%a1 + msq_min 
  f = map%a3 * msq1
\nwendcode{}\nwbegincode{345}\moddef{Apply inverse s-channel pole mapping}\endmoddef
  msq1 = msq - msq_min + map%a1
  x = log (msq1 / map%a1) / map%a2
  f = map%a3 * msq1
\nwendcode{}\nwbegindocs{346}\nwdocspar
\subsubsection{Mapping for t-channel poles}
This is also approximately scale-invariant, and we use the same type
of mapping as before.  However, we map $1/x$ singularities at both
ends of the interval; again, the mapping becomes linear when the
distance is less than $M^2$:
\begin{equation}
  m^2 = 
  \begin{cases}
    m_{\textrm{min}}^2 + M^2\left(\exp(xL)-1\right)
    &
    \text{for $0 < x < \frac12$}
    \\
    m_{\textrm{max}}^2 - M^2\left(\exp((1-x)L)-1\right)
    &
    \text{for $\frac12 \leq x < 1$}
  \end{cases}
\end{equation}
where
\begin{equation}
  L = 2\ln\left(\frac{m_{\textrm{max}}^2 - m_{\textrm{min}}^2}{2M^2} 
                + 1\right).
\end{equation}
The inverse:
\begin{equation}
  x = 
  \begin{cases}
    \frac1L\ln\left(\frac{m^2-m_{\textrm{min}}^2}{M^2} + 1\right)
    &
    \text{for $m^2 < (m_{\textrm{max}}^2 - m_{\textrm{min}}^2)/2$}
    \\
    1 - \frac1L\ln\left(\frac{m_{\textrm{max}}-m^2}{M^2} + 1\right)
    &
    \text{for $m^2 \geq (m_{\textrm{max}}^2 - m_{\textrm{min}}^2)/2$}
  \end{cases}
\end{equation}
The phase-space factor is
\begin{equation}
  f(x) = 
  \begin{cases}
    \frac{M^2}{s}\,\exp(xL)\,L.
    &
    \text{for $0 < x < \frac12$}
    \\
    \frac{M^2}{s}\,\exp((1-x)L)\,L.
    &
    \text{for $\frac12 \leq x < 1$}
  \end{cases}
\end{equation}
A (continuous) function proportional to
\begin{equation}
  g(m^2) =
  \begin{cases}
    1/(m^2-m_{\textrm{min}}^2) + M^2)
    &
    \text{for $m^2 < (m_{\textrm{max}}^2 - m_{\textrm{min}}^2)/2$}
    \\
    1/((m_{\textrm{max}}^2 - m^2) + M^2)
    &
    \text{for $m^2 \leq (m_{\textrm{max}}^2 - m_{\textrm{min}}^2)/2$}
  \end{cases}
\end{equation}
is mapped to a constant by this mapping, i.e., poles near both ends of
the interval.
\nwenddocs{}\nwbegincode{347}\moddef{Constants for t-channel pole mapping}\endmoddef
  if (map%variable_limits .or. map%a_unknown) then
     msq0 = map%mass**2
     map%a1 = msq0
     map%a2 = 2 * log ((msq_max - msq_min)/(2*msq0) + 1)
     map%a3 = map%a2 / s
     map%a_unknown = .false.
  end if
\nwendcode{}\nwbegincode{348}\moddef{Apply t-channel pole mapping}\endmoddef
  if (x < .5_default) then
     msq1 = map%a1 * exp (x * map%a2)
     msq = msq1 - map%a1 + msq_min
  else
     msq1 = map%a1 * exp ((1-x) * map%a2)
     msq = -(msq1 - map%a1) + msq_max
  end if
  f = map%a3 * msq1
\nwendcode{}\nwbegincode{349}\moddef{Apply inverse t-channel pole mapping}\endmoddef
  if (msq < (msq_max + msq_min)/2) then
     msq1 = msq - msq_min + map%a1
     x = log (msq1/map%a1) / map%a2
  else
     msq1 = msq_max - msq + map%a1
     x = 1 - log (msq1/map%a1) / map%a2
  end if
  f = map%a3 * msq1
\nwendcode{}\nwbegindocs{350}\nwdocspar
\subsection{Step mapping}
Step mapping is useful when the allowed range for a squared-mass
variable is large, but only a fraction at the lower end is populated
because the particle in question is an (off-shell) decay product of a
narrow resonance.  I.e., if the resonance was forced to be on-shell,
the upper end of the range would be the resonance mass, minus the
effective (real or resonance) mass of the particle(s) in the sibling
branch of the decay.

The edge of this phase space section has a width which is determined
by the width of the parent, plus the width of the sibling branch.  (The
widths might be added in quadrature, but this precision is probably
not important.)

\subsubsection{Fermi function}

A possible mapping is derived from the Fermi function which has
precisely this behavior.  The Fermi function is given by
\begin{equation}
  f(x) = \frac{1}{1 + \exp\frac{x-\mu}{\gamma}}
\end{equation}
where $x$ is taken as the invariant mass squared, $\mu$ is the
invariant mass squared of the edge, and $\gamma$ is the effective
width which is given by the widths of the parent and the sibling
branch.  (Widths might be added in quadrature, but we do not require
this level of precision.)
\begin{align}
  x &= \frac{m^2 - m_{\text{min}}^2}{\Delta m^2}
\\
  \mu &= 
  \frac{m_{\text{max,eff}}^2 - m_{\text{min}}^2}
          {\Delta m^2}
\\
  \gamma &= \frac{2m_{\text{max,eff}}\Gamma}{\Delta m^2}
\end{align}
with
\begin{equation}
  \Delta m^2 = m_{\text{max}}^2 - m_{\text{min}}^2
\end{equation}
$m^2$ is thus given by
\begin{equation}
  m^2(x) = xm_{\text{max}}^2 + (1-x)m_{\text{min}}^2
\end{equation}
For the mapping, we compute the integral $g(x)$ of the Fermi function,
normalized such that $g(0)=0$ and $g(1)=1$.  We introduce the abbreviations
\begin{align}
  \alpha &= 1 - \gamma\ln\frac{1 + \beta e^{1/\gamma}}{1 + \beta}
\\
  \beta  &= e^{- \mu/\gamma}
\end{align}
and obtain
\begin{equation}
  g(x) = \frac{1}{\alpha}
         \left(x - \gamma\ln\frac{1 + \beta e^{x/\gamma}}
                                 {1 + \beta}\right)
\end{equation}
The actual mapping is the inverse function $h(y) = g^{-1}(y)$,
\begin{equation}
  h(y) = -\gamma\ln\left(e^{-\alpha y/\gamma}(1 + \beta) - \beta\right)
\end{equation}
The Jacobian is
\begin{equation}
  \frac{dh}{dy} = \alpha\left(1 - e^{\alpha y/\gamma}
                                  \frac{\beta}{1 + \beta}\right)^{-1}
\end{equation}
which is equal to $1/(dg/dx)$, namely
\begin{equation}
  \frac{dg}{dx} = \frac{1}{\alpha}\,\frac{1}{1 + \beta e^{x/\gamma}}
\end{equation}
The final result is
\begin{align}
  \int_{m_{\text{min}}^2}^{m_{\text{max}}^2} dm^2\,F(m^2)
  &= \Delta m^2\int_0^1\,dx\,F(m^2(x))
\\
  &= \Delta m^2\int_0^1\,dy\,F(m^2(h(y)))\,\frac{dh}{dy}
\end{align}
Here is the implementation.  We fill {\tt{}a1}, {\tt{}a2}, {\tt{}a3} with
$\alpha,\beta,\gamma$, respectively.
\nwenddocs{}\nwbegincode{351}\moddef{Constants for exponential step mapping}\endmoddef
  if (map%variable_limits .or. map%a_unknown) then
     map%a3 = max (2 * map%mass * map%width / (msq_max - msq_min), 0.01_default)
     map%a2 = exp (- (map%mass**2 - msq_min) / (msq_max - msq_min) &
                     / map%a3)
     map%a1 = 1 - map%a3 * log ((1 + map%a2 * exp (1 / map%a3)) / (1 + map%a2))
  end if
\nwendcode{}\nwbegincode{352}\moddef{Apply exponential step mapping}\endmoddef
  tmp = exp (- x * map%a1 / map%a3) * (1 + map%a2)
  z = - map%a3 * log (tmp - map%a2)
  msq  = z * msq_max + (1 - z) * msq_min
  f = map%a1 / (1 - map%a2 / tmp) * (msq_max - msq_min) / s
\nwendcode{}\nwbegincode{353}\moddef{Apply inverse exponential step mapping}\endmoddef
  z = (msq - msq_min) / (msq_max - msq_min)
  tmp = 1 + map%a2 * exp (z / map%a3)
  x = (z - map%a3 * log (tmp / (1 + map%a2))) &
      / map%a1
  f = map%a1 * tmp * (msq_max - msq_min) / s
\nwendcode{}\nwbegindocs{354}\nwdocspar
\subsubsection{Hyperbolic mapping}

The Fermi function has the drawback that it decreases exponentially.
It might be preferable to take a function with a power-law decrease,
such that the high-mass region is not completely depopulated.

Here, we start with the actual mapping which we take as
\begin{equation}
  h(y) = \frac{b}{a-y} - \frac{b}{a} + \mu y
\end{equation}
with the abbreviation
\begin{equation}
  a = \frac12\left(1 + \sqrt{1 + \frac{4b}{1-\mu}}\right)
\end{equation}
This is a hyperbola in the $xy$ plane.  The derivative is
\begin{equation}
  \frac{dh}{dy} = \frac{b}{(a-y)^2} + \mu
\end{equation}
The constants correspond to
\begin{align}
  \mu &= 
  \frac{m_{\text{max,eff}}^2 - m_{\text{min}}^2}
          {\Delta m^2}
\\
  b &= \frac{1}{\mu}\left(\frac{2m_{\text{max,eff}}\Gamma}{\Delta m^2}\right)^2
\end{align}
The inverse function is the solution of a quadratic equation,
\begin{equation}
  g(x) = \frac{1}{2}
         \left[\left(a + \frac{x}{\mu} + \frac{b}{a\mu}\right) 
               - \sqrt{\left(a-\frac{x}{\mu}\right)^2
                       + 2\frac{b}{a\mu}\left(a + \frac{x}{\mu}\right)
                       + \left(\frac{b}{a\mu}\right)^2}\right]
\end{equation}
The constants $a_{1,2,3}$ are identified with $a,b,\mu$.
\nwenddocs{}\nwbegincode{355}\moddef{Constants for hyperbolic step mapping}\endmoddef
  if (map%variable_limits .or. map%a_unknown) then
     map%a3 = (map%mass**2 - msq_min) / (msq_max - msq_min)
     map%a2 = max ((2 * map%mass * map%width / (msq_max - msq_min))**2 &
                   / map%a3, 1e-6_default)
     map%a1 = (1 + sqrt (1 + 4 * map%a2 / (1 - map%a3))) / 2
  end if
\nwendcode{}\nwbegincode{356}\moddef{Apply hyperbolic step mapping}\endmoddef
  z = map%a2 / (map%a1 - x) - map%a2 / map%a1 + map%a3 * x
  msq = z * msq_max + (1 - z) * msq_min
  f = (map%a2 / (map%a1 - x)**2 + map%a3) * (msq_max - msq_min) / s
\nwendcode{}\nwbegincode{357}\moddef{Apply inverse hyperbolic step mapping}\endmoddef
  z = (msq - msq_min) / (msq_max - msq_min)
  tmp = map%a2 / (map%a1 * map%a3)
  x = ((map%a1 + z / map%a3 + tmp) &
       - sqrt ((map%a1 - z / map%a3)**2 + 2 * tmp * (map%a1 + z / map%a3) &
               + tmp**2)) / 2
  f = (map%a2 / (map%a1 - x)**2 + map%a3) * (msq_max - msq_min) / s
\nwendcode{}\nwbegindocs{358}\nwdocspar
\subsection{Mappings of the polar angle}
The other type of singularity, a simple pole just outside the
integration region, can occur in the integration over $\cos\theta$.
This applies to exchange of massless (or light) particles.

Double poles (Coulomb scattering) are also possible, but only in
certain cases.  These are also handled by the single-pole mapping.

The mapping is analogous to the previous $m^2$ pole mapping, but with
a different normalization and notation of variables:
\begin{equation}
  \frac12\int_{-1}^1 d\cos\theta\,g(\theta)
  = \int_0^1 dx\,\frac{d\cos\theta}{dx}\,g(\theta(x))
  = \int_0^1 dx\,f(x)\,g(x),
\end{equation}
where thus
\begin{equation}
  f(x) = \frac12\,\frac{d\cos\theta}{dx}.
\end{equation}
With this mapping, a function of the form
\begin{equation}
  g(\theta) = c\frac{dx(\cos\theta)}{d\cos\theta}
\end{equation}
is mapped to a constant:
\begin{equation}
  \int_{-1}^1 d\cos\theta\,g(\theta)
  = \int_0^1 dx\,f(x)\,g(\theta(x)) = \int_0^1 dx\,c.
\end{equation}
\nwenddocs{}\nwbegincode{359}\moddef{Mappings: public}\plusendmoddef
  public :: mapping_compute_ct_from_x
\nwendcode{}\nwbegincode{360}\moddef{Mappings: procedures}\plusendmoddef
  subroutine mapping_compute_ct_from_x (map, s, ct, st, f, x)
    type(mapping_t), intent(inout) :: map
    real(default), intent(in) :: s
    real(default), intent(out) :: ct, st, f
    real(default), intent(in) :: x
    real(default) :: tmp, ct1
    select case (map%type)
    case (NO_MAPPING, S_CHANNEL, INFRARED, RADIATION, &
         STEP_MAPPING_E, STEP_MAPPING_H)
     \LA{}Apply trivial ct mapping\RA{}
    case (T_CHANNEL, U_CHANNEL, COLLINEAR)
     \LA{}Constants for ct pole mapping\RA{}
     \LA{}Apply ct pole mapping\RA{}
    case default
       call msg_fatal (" Attempt to apply undefined ct mapping")
    end select
  end subroutine mapping_compute_ct_from_x

\eatline
\nwendcode{}\nwbegincode{361}\moddef{Mappings: public}\plusendmoddef
  public :: mapping_compute_x_from_ct
\nwendcode{}\nwbegincode{362}\moddef{Mappings: procedures}\plusendmoddef
  subroutine mapping_compute_x_from_ct (map, s, ct, f, x)
    type(mapping_t), intent(inout) :: map
    real(default), intent(in) :: s
    real(default), intent(in) :: ct
    real(default), intent(out) :: f, x
    real(default) :: ct1
    select case (map%type)
    case (NO_MAPPING, S_CHANNEL, INFRARED, RADIATION, &
         STEP_MAPPING_E, STEP_MAPPING_H)
     \LA{}Apply inverse trivial ct mapping\RA{}
    case (T_CHANNEL, U_CHANNEL, COLLINEAR)
     \LA{}Constants for ct pole mapping\RA{}
     \LA{}Apply inverse ct pole mapping\RA{}
    case default
       call msg_fatal (" Attempt to apply undefined inverse ct mapping")
    end select
  end subroutine mapping_compute_x_from_ct

\eatline
\nwendcode{}\nwbegindocs{363}\nwdocspar
\subsubsection{Trivial mapping}
This is just the mapping of the interval $(-1,1)$ to $(0,1)$:
\begin{equation}
  \cos\theta = -1 + 2x
\end{equation}
and
\begin{equation}
  f(x) = 1
\end{equation}
with the inverse
\begin{equation}
  x = \frac{1+\cos\theta}{2}
\end{equation}
\nwenddocs{}\nwbegincode{364}\moddef{Apply trivial ct mapping}\endmoddef
  tmp = 2 * (1-x)
  ct = 1 - tmp
  st = sqrt (tmp * (2-tmp))
  f = 1
\nwendcode{}\nwbegincode{365}\moddef{Apply inverse trivial ct mapping}\endmoddef
  x = (ct + 1) / 2
  f = 1
\nwendcode{}\nwbegindocs{366}\nwdocspar
\subsubsection{Pole mapping}
As above for $m^2$, we simultaneously map poles at both ends of the
$\cos\theta$ interval.  The formulae are completely analogous:
\begin{equation}
  \cos\theta =
  \begin{cases}
    \frac{M^2}{s}\left[\exp(xL)-1\right] - 1
    &
    \text{for $x<\frac12$}
    \\
    -\frac{M^2}{s}\left[\exp((1-x)L)-1\right] + 1
    &
    \text{for $x\geq\frac12$}
  \end{cases}
\end{equation}
where
\begin{equation}
  L = 2\ln\frac{M^2+s}{M^2}.
\end{equation}
Inverse:
\begin{equation}
  x = 
  \begin{cases}
    \frac{1}{2L}\ln\frac{1 + \cos\theta + M^2/s}{M^2/s}
    &
    \text{for $\cos\theta < 0$}
    \\
    1 - \frac{1}{2L}\ln\frac{1 - \cos\theta + M^2/s}{M^2/s}
    &
    \text{for $\cos\theta \geq 0$}
  \end{cases}
\end{equation}
The phase-space factor:
\begin{equation}
  f(x) = 
  \begin{cases}
    \frac{M^2}{s}\exp(xL)\,L
    &
    \text{for $x<\frac12$}
    \\
    \frac{M^2}{s}\exp((1-x)L)\,L
    &
    \text{for $x\geq\frac12$}
  \end{cases}
\end{equation}
\nwenddocs{}\nwbegincode{367}\moddef{Constants for ct pole mapping}\endmoddef
  if (map%variable_limits .or. map%b_unknown) then
     map%b1 = map%mass**2 / s
     map%b2 = log ((map%b1 + 1) / map%b1)
     map%b3 = 0
     map%b_unknown = .false.
  end if
\nwendcode{}\nwbegincode{368}\moddef{Apply ct pole mapping}\endmoddef
  if (x < .5_default) then
     ct1 = map%b1 * exp (2 * x * map%b2)
     ct = ct1 - map%b1 - 1
  else
     ct1 = map%b1 * exp (2 * (1-x) * map%b2)
     ct = -(ct1 - map%b1) + 1
  end if
  if (ct >= -1 .and. ct <= 1) then
     st = sqrt (1 - ct**2)
     f = ct1 * map%b2
  else
     ct = 1;  st = 0;  f = 0
  end if
\nwendcode{}\nwbegincode{369}\moddef{Apply inverse ct pole mapping}\endmoddef
  if (ct < 0) then
     ct1 = ct + map%b1 + 1
     x = log (ct1 / map%b1) / (2 * map%b2)
  else
     ct1 = -ct + map%b1 + 1
     x = 1 - log (ct1 / map%b1) / (2 * map%b2)
  end if
  f = ct1 * map%b2
\nwendcode{}\nwbegindocs{370}\nwdocspar
%%%%%%%%%%%%%%%%%%%%%%%%%%%%%%%%%%%%%%%%%%%%%%%%%%%%%%%%%%%%%%%%%%%%%%%%
\clearpage
\section{Phase-space trees}
The phase space evaluation is organized in terms of trees, where each
branch corresponds to three integrations: $m^2$, $\cos\theta$, and
$\phi$.  The complete tree thus makes up a specific parameterization
of the multidimensional phase-space integral.  For the multi-channel
integration, the phase-space tree is a single channel.

The trees imply mappings of formal Feynman tree graphs into arrays of
integer numbers: Each branch, corresponding to a particular line in
the graph, is assigned an integer code $c$ (with kind value {\tt{}TC} =
tree code).

In this integer, each bit determines whether a particular external
momentum flows through the line.  The external branches therefore have
codes $1,2,4,8,\ldots$.  An internal branch has those bits ORed
corresponding to the momenta flowing through it.  For example, a
branch with momentum $p_1+p_4$ has code $2^0+2^3=1+8=9$.

There is a two-fold ambiguity: Momentum conservation implies that the
branch with code
\begin{equation}
  c_0 = \sum_{i=1}^{n(\rm{ext})} 2^{i-1}
\end{equation}
i.e. the branch with momentum $p_1+p_2+\ldots p_n$ has momentum zero,
which is equivalent to tree code $0$ by definition.  Correspondingly,
\begin{equation}
  c \quad\textrm{and}\quad c_0 - c = c\;\textrm{XOR}\;c_0
\end{equation}
are equivalent.  E.g., if there are five externals with codes
$c=1,2,4,8,16$, then $c=9$ and $\bar c=31-9=22$ are equivalent.

This ambiguity may be used to assign a direction to the line: If all
momenta are understood as outgoing, $c=9$ in the example above means
$p_1+p_4$, but $c=22$ means $p_2+p_3+p_5 = -(p_1+p_4)$.

Here we make use of the ambiguity in a slightly different way.  First,
the initial particles are singled out as those externals with the
highest bits, the IN-bits.  (Here: $8$ and $16$ for a $2\to 3$
scattering process, $16$ only for a $1\to 4$ decay.)  Then we invert
those codes where all IN-bits are set.  For a decay process this maps
each tree of an equivalence class onto a unique representative (that one
with the smallest integer codes).  For a scattering process we proceed
further:

The ambiguity remains in all branches where only one IN-bit is set,
including the initial particles.  If there are only externals with
this property, we have an $s$-channel graph which we leave as it is.
In all other cases, an internal with only one IN-bit is a $t$-channel
line, which for phase space integration should be associated with one
of the initial momenta as a reference axis.  We take that one whose
bit is set in the current tree code.  (E.g., for branch $c=9$ we use
the initial particle $c=8$ as reference axis, whereas for the same
branch we would take $c=16$ if it had been assigned $\bar c=31-9=22$
as tree code.) Thus, different ways of coding the same $t$-channel
graph imply different phase space parameterizations.

$s$-channel graphs have a unique parameterization.  The same sets of
parameterizations are used for $t$-channel graphs, except for the
reference frames of their angular parts.  We map each
$t$-channel graph onto an $s$-channel graph as follows:

Working in ascending order, for each $t$-channel line (whose code has
exactly one IN-bit set) the attached initial line is flipped upstream,
while the outgoing line is flipped downstream.  (This works only if
$t$-channel graphs are always parameterized beginning at their outer
vertices, which we require as a restriction.)  After all possible
flips have been applied, we have an $s$-channel graph.  We only have
to remember the initial particle a vertex was originally attached to.
\nwenddocs{}\nwbegincode{371}\moddef{\code{}phs{\char95}trees.f90\edoc{}}\endmoddef
\LA{}File header\RA{}

module phs_trees

\LA{}Use kinds\RA{}
  use kinds, only: TC
\LA{}Use strings\RA{}
  use io_units
  use constants, only: twopi, twopi2, twopi5
  use format_defs, only: FMT_19
  use diagnostics
  use lorentz
  use permutations, only: permutation_t, permutation_size
  use permutations, only: permutation_init, permutation_find
  use permutations, only: tc_decay_level, tc_permute
  use model_data
  use flavors
  use mappings

\LA{}Standard module head\RA{}

\LA{}PHS trees: public\RA{}

\LA{}PHS trees: types\RA{}

contains

\LA{}PHS trees: procedures\RA{}

end module phs_trees
\eatline
\nwendcode{}\nwbegindocs{372}\nwdocspar
\subsection{Particles}
We define a particle type which contains only four-momentum and
invariant mass squared, and a flag that tells whether the momentum is
filled or not.
\nwenddocs{}\nwbegincode{373}\moddef{PHS trees: public}\endmoddef
  public :: phs_prt_t
\nwendcode{}\nwbegincode{374}\moddef{PHS trees: types}\endmoddef
  type :: phs_prt_t
     private
     logical :: defined = .false.
     type(vector4_t) :: p
     real(default) :: p2
  end type phs_prt_t

\eatline
\nwendcode{}\nwbegindocs{375}Set contents:
\nwenddocs{}\nwbegincode{376}\moddef{PHS trees: public}\plusendmoddef
  public :: phs_prt_set_defined
  public :: phs_prt_set_undefined
  public :: phs_prt_set_momentum
  public :: phs_prt_set_msq
\nwendcode{}\nwbegincode{377}\moddef{PHS trees: procedures}\endmoddef
  elemental subroutine phs_prt_set_defined (prt)
    type(phs_prt_t), intent(inout) :: prt
    prt%defined = .true.
  end subroutine phs_prt_set_defined

  elemental subroutine phs_prt_set_undefined (prt)
    type(phs_prt_t), intent(inout) :: prt
    prt%defined = .false.
  end subroutine phs_prt_set_undefined

  elemental subroutine phs_prt_set_momentum (prt, p)
    type(phs_prt_t), intent(inout) :: prt
    type(vector4_t), intent(in) :: p
    prt%p = p
  end subroutine phs_prt_set_momentum

  elemental subroutine phs_prt_set_msq (prt, p2)
    type(phs_prt_t), intent(inout) :: prt
    real(default), intent(in) :: p2
    prt%p2 = p2
  end subroutine phs_prt_set_msq

\eatline
\nwendcode{}\nwbegindocs{378}Access methods:
\nwenddocs{}\nwbegincode{379}\moddef{PHS trees: public}\plusendmoddef
  public :: phs_prt_is_defined
  public :: phs_prt_get_momentum
  public :: phs_prt_get_msq
\nwendcode{}\nwbegincode{380}\moddef{PHS trees: procedures}\plusendmoddef
  elemental function phs_prt_is_defined (prt) result (defined)
    logical :: defined
    type(phs_prt_t), intent(in) :: prt
    defined = prt%defined
  end function phs_prt_is_defined

  elemental function phs_prt_get_momentum (prt) result (p)
    type(vector4_t) :: p
    type(phs_prt_t), intent(in) :: prt
    p = prt%p
  end function phs_prt_get_momentum

  elemental function phs_prt_get_msq (prt) result (p2)
    real(default) :: p2
    type(phs_prt_t), intent(in) :: prt
    p2 = prt%p2
  end function phs_prt_get_msq

\eatline
\nwendcode{}\nwbegindocs{381}Addition of momenta (invariant mass square is computed).
\nwenddocs{}\nwbegincode{382}\moddef{PHS trees: public}\plusendmoddef
  public :: phs_prt_combine
\nwendcode{}\nwbegincode{383}\moddef{PHS trees: procedures}\plusendmoddef
  elemental subroutine phs_prt_combine (prt, prt1, prt2)
    type(phs_prt_t), intent(inout) :: prt
    type(phs_prt_t), intent(in) :: prt1, prt2
    prt%defined = .true.
    prt%p = prt1%p + prt2%p
    prt%p2 = prt%p ** 2
    call phs_prt_check (prt)
  end subroutine phs_prt_combine

\eatline
\nwendcode{}\nwbegindocs{384}Output
\nwenddocs{}\nwbegincode{385}\moddef{PHS trees: public}\plusendmoddef
  public :: phs_prt_write
\nwendcode{}\nwbegincode{386}\moddef{PHS trees: procedures}\plusendmoddef
  subroutine phs_prt_write (prt, unit)
    type(phs_prt_t), intent(in) :: prt
    integer, intent(in), optional :: unit
    integer :: u
    u = given_output_unit (unit);  if (u < 0)  return
    if (prt%defined) then
       call vector4_write (prt%p, u)
       write (u, "(1x,A,1x," // FMT_19 // ")") "T = ", prt%p2
    else
       write (u, "(3x,A)") "[undefined]"
    end if
  end subroutine phs_prt_write

\eatline
\nwendcode{}\nwbegincode{387}\moddef{PHS trees: public}\plusendmoddef
  public :: phs_prt_check
\nwendcode{}\nwbegincode{388}\moddef{PHS trees: procedures}\plusendmoddef
  elemental subroutine phs_prt_check (prt)
    type(phs_prt_t), intent(inout) :: prt
    real(default), parameter :: eps = 1E-9
    if (prt%p2 < 0._default) then
       prt%p2 = 0._default
    end if
  end subroutine phs_prt_check

\eatline
\nwendcode{}\nwbegindocs{389}\nwdocspar
\subsection{The phase-space tree type}
\subsubsection{Definition}
In the concrete implementation, each branch $c$ may have two
\emph{daughters} $c_1$ and $c_2$ such that $c_1+c_2=c$, a
\emph{sibling} $c_s$ and a \emph{mother} $c_m$ such that $c+c_s =
c_m$, and a \emph{friend} which is kept during flips, such that it can
indicate a fixed reference frame.  Absent entries are set $c=0$.

First, declare the branch type.  There is some need to have this
public. Give initializations for all components, so no {\tt{}init}
routine is necessary.  The branch has some information about the
associated coordinates and about connections.
\nwenddocs{}\nwbegincode{390}\moddef{PHS trees: types}\plusendmoddef
  type :: phs_branch_t
     private
     logical :: set = .false.
     logical :: inverted_decay = .false.
     logical :: inverted_axis = .false.
     integer(TC) :: mother = 0
     integer(TC) :: sibling = 0
     integer(TC) :: friend = 0
     integer(TC) :: origin = 0
     integer(TC), dimension(2) :: daughter = 0
     integer :: firstborn = 0
     logical :: has_children = .false.
     logical :: has_friend = .false.
     logical :: is_real = .false.
  end type phs_branch_t

\eatline
\nwendcode{}\nwbegindocs{391}The tree type: No initialization, this is done by
{\tt{}phs{\char95}tree{\char95}init}.  In addition to the branch array which 

The branches are collected in an array which holds all possible
branches, of which only a few are set.  After flips have been applied,
the branch $c_M=\sum_{i=1}^{n({\rm fin})}2^{i-1}$ must be there,
indicating the mother of all decay products.  In addition, we should
check for consistency at the beginning.

{\tt{}n{\char95}branches} is the number of those actually set.  {\tt{}n{\char95}externals}
defines the number of significant bit, and {\tt{}mask} is a code where all
bits are set.  Analogous: {\tt{}n{\char95}in} and {\tt{}mask{\char95}in} for the incoming
particles.

The {\tt{}mapping} array contains the mappings associated to the branches
(corresponding indices).  The array {\tt{}mass{\char95}sum} contains the sum of
the real masses of the external final-state particles associated to
the branch.  During phase-space evaluation, this determines the
boundaries.
\nwenddocs{}\nwbegincode{392}\moddef{PHS trees: public}\plusendmoddef
  public :: phs_tree_t
\nwendcode{}\nwbegincode{393}\moddef{PHS trees: types}\plusendmoddef
  type :: phs_tree_t
     private
     integer :: n_branches, n_externals, n_in, n_msq, n_angles
     integer(TC) :: n_branches_tot, n_branches_out
     integer(TC) :: mask, mask_in, mask_out
     type(phs_branch_t), dimension(:), allocatable :: branch
     type(mapping_t), dimension(:), allocatable :: mapping
     real(default), dimension(:), allocatable :: mass_sum
     real(default), dimension(:), allocatable :: effective_mass
     real(default), dimension(:), allocatable :: effective_width
     logical :: real_phsp = .false.
     integer(TC) :: emitter = 0
     integer, dimension(:), allocatable :: momentum_link
   contains
   \LA{}PHS trees: phs tree: TBP\RA{}
  end type phs_tree_t

\eatline
\nwendcode{}\nwbegindocs{394}The maximum number of external particles that can be represented is
related to the bit size of the integer that stores binary codes.  With
the default integer of 32 bit on common machines, this is more than
enough space.  If {\tt{}TC} is actually the default integer kind, there
is no need to keep it separate, but doing so marks this as a
special type of integer.  So, just state that the maximum number is 32:
\nwenddocs{}\nwbegincode{395}\moddef{Limits: public parameters}\endmoddef
  integer, parameter, public :: MAX_EXTERNAL = 32
\eatline
\nwendcode{}\nwbegindocs{396}\nwdocspar
\subsubsection{Constructor and destructor}
Allocate memory for a phase-space tree with given number of externals and
incoming.  The number of allocated branches can easily become large,
but appears manageable for realistic cases, e.g., for {\tt{}n{\char95}in=2} and
{\tt{}n{\char95}out=8} we get $2^{10}-1=1023$.
\nwenddocs{}\nwbegincode{397}\moddef{PHS trees: public}\plusendmoddef
  public :: phs_tree_init
  public :: phs_tree_final
\nwendcode{}\nwbegindocs{398}Here we set the masks for incoming and for all externals.
\nwenddocs{}\nwbegincode{399}\moddef{PHS trees: procedures}\plusendmoddef
  elemental subroutine phs_tree_init (tree, n_in, n_out, n_masses, n_angles)
    type(phs_tree_t), intent(inout) :: tree
    integer, intent(in) :: n_in, n_out, n_masses, n_angles
    integer(TC) :: i
    tree%n_externals = n_in + n_out
    tree%n_branches_tot = 2**(n_in+n_out) - 1
    tree%n_branches_out = 2**n_out - 1
    tree%mask = 0
    do i = 0, n_in + n_out - 1
       tree%mask = ibset (tree%mask, i)
    end do
    tree%n_in = n_in
    tree%mask_in = 0
    do i = n_out, n_in + n_out - 1
       tree%mask_in = ibset (tree%mask_in, i)
    end do
    tree%mask_out = ieor (tree%mask, tree%mask_in)
    tree%n_msq = n_masses
    tree%n_angles = n_angles
    allocate (tree%branch (tree%n_branches_tot))
    tree%n_branches  = 0
    allocate (tree%mapping (tree%n_branches_out))
    allocate (tree%mass_sum (tree%n_branches_out))
    allocate (tree%effective_mass (tree%n_branches_out))
    allocate (tree%effective_width (tree%n_branches_out))
  end subroutine phs_tree_init

  elemental subroutine phs_tree_final (tree)
    type(phs_tree_t), intent(inout) :: tree
    deallocate (tree%branch)
    deallocate (tree%mapping)
    deallocate (tree%mass_sum)
    deallocate (tree%effective_mass)
    deallocate (tree%effective_width)
  end subroutine phs_tree_final

\eatline
\nwendcode{}\nwbegindocs{400}\nwdocspar
\subsubsection{Screen output}
Write only the branches that are set:
\nwenddocs{}\nwbegincode{401}\moddef{PHS trees: public}\plusendmoddef
  public :: phs_tree_write
\nwendcode{}\nwbegincode{402}\moddef{PHS trees: procedures}\plusendmoddef
  subroutine phs_tree_write (tree, unit)
    type(phs_tree_t), intent(in) :: tree
    integer, intent(in), optional :: unit
    integer :: u
    integer(TC) :: k
    u = given_output_unit (unit);  if (u < 0)  return
    write (u, '(3X,A,1x,I0,5X,A,I3)') &
         'External:', tree%n_externals, 'Mask:', tree%mask
    write (u, '(3X,A,1x,I0,5X,A,I3)') &
         'Incoming:', tree%n_in, 'Mask:', tree%mask_in
    write (u, '(3X,A,1x,I0,5X,A,I3)') &
         'Branches:', tree%n_branches
    do k = size (tree%branch), 1, -1
       if (tree%branch(k)%set) &
            call phs_branch_write (tree%branch(k), unit=unit, kval=k)
    end do
    do k = 1, size (tree%mapping)
       call mapping_write (tree%mapping (k), unit, verbose=.true.)
    end do
    write (u, "(3x,A)") "Arrays: mass_sum, effective_mass, effective_width"
    do k = 1, size (tree%mass_sum)
       if (tree%branch(k)%set) then
          write (u, "(5x,I0,3(2x," // FMT_19 // "))") k, tree%mass_sum(k), &
               tree%effective_mass(k), tree%effective_width(k)
       end if
    end do
  end subroutine phs_tree_write

  subroutine phs_branch_write (b, unit, kval)
    type(phs_branch_t), intent(in) :: b
    integer, intent(in), optional :: unit
    integer(TC), intent(in), optional :: kval
    integer :: u
    integer(TC) :: k
    character(len=6) :: tmp
    character(len=1) :: firstborn(2), sign_decay, sign_axis
    integer :: i
    u = given_output_unit (unit);  if (u < 0)  return
    k = 0;  if (present (kval))  k = kval
    if (b%origin /= 0) then
       write(tmp, '(A,I4,A)') '(', b%origin, ')'
    else
       tmp = ' '
    end if
    do i=1, 2
       if (b%firstborn == i) then
          firstborn(i) = "*"
       else
          firstborn(i) = " "
       end if
    end do
    if (b%inverted_decay) then
       sign_decay = "-"
    else
       sign_decay = "+"
    end if
    if (b%inverted_axis) then
       sign_axis = "-"
    else
       sign_axis = "+"
    end if
    if (b%has_children) then
       if (b%has_friend) then
          write(u,'(4X,A1,I0,3x,A,1X,A,I0,A1,1x,I0,A1,1X,A1,1X,A,1x,I0)') &
               &   '*', k, tmp, &
               &   'Daughters: ', &
               &   b%daughter(1), firstborn(1), &
               &   b%daughter(2), firstborn(2), sign_decay, &
               &   'Friend:    ', b%friend
       else
          write(u,'(4X,A1,I0,3x,A,1X,A,I0,A1,1x,I0,A1,1X,A1,1X,A)') &
               &   '*', k, tmp, &
               &   'Daughters: ', &
               &   b%daughter(1), firstborn(1), &
               &   b%daughter(2), firstborn(2), sign_decay, &
               &   '(axis '//sign_axis//')'
       end if
    else
       write(u,'(5X,I0)') k
    end if
  end subroutine phs_branch_write

\eatline
\nwendcode{}\nwbegindocs{403}\nwdocspar
\subsection{PHS tree setup}
\subsubsection{Transformation into an array of branch codes and back}
Assume that the tree/array has been created before with the
appropriate length and is empty.  
\nwenddocs{}\nwbegincode{404}\moddef{PHS trees: public}\plusendmoddef
  public :: phs_tree_from_array
\nwendcode{}\nwbegincode{405}\moddef{PHS trees: procedures}\plusendmoddef
  subroutine phs_tree_from_array (tree, a)
    type(phs_tree_t), intent(inout) :: tree
    integer(TC), dimension(:), intent(in) :: a
    integer :: i
    integer(TC) :: k
  \LA{}Set branches from array \code{}a\edoc{}\RA{}
  \LA{}Set external branches if necessary\RA{}
  \LA{}Check number of branches\RA{}
  \LA{}Determine the connections\RA{}
  contains
  \LA{}Subroutine: set relatives\RA{}
  end subroutine phs_tree_from_array

\eatline
\nwendcode{}\nwbegindocs{406}First, set all branches specified by the user.  If all IN-bits
are set, we invert the branch code.  
\nwenddocs{}\nwbegincode{407}\moddef{Set branches from array \code{}a\edoc{}}\endmoddef
  do i=1, size(a)
     k = a(i)
     if (iand(k, tree%mask_in) == tree%mask_in)  k = ieor(tree%mask, k)
     tree%branch(k)%set = .true.
     tree%n_branches = tree%n_branches+1
  end do
\nwendcode{}\nwbegindocs{408}The external branches are understood, so set them now if not yet
done.  In all cases ensure that the representative with one bit set is
used, except for decays where the in-particle is represented by all
OUT-bits set instead.
\nwenddocs{}\nwbegincode{409}\moddef{Set external branches if necessary}\endmoddef
  do i=0, tree%n_externals-1
     k = ibset(0,i)
     if (iand(k, tree%mask_in) == tree%mask_in)  k = ieor(tree%mask, k)
     if (tree%branch(ieor(tree%mask, k))%set) then
        tree%branch(ieor(tree%mask, k))%set = .false.
        tree%branch(k)%set = .true.
     else if (.not.tree%branch(k)%set) then
        tree%branch(k)%set = .true.
        tree%n_branches = tree%n_branches+1
     end if
  end do
\nwendcode{}\nwbegindocs{410}Now the number of branches set can be checked.  Here we assume that
the tree is binary.  For three externals there are three branches in
total, and for each additional external branch we get another internal
one.
\nwenddocs{}\nwbegincode{411}\moddef{Check number of branches}\endmoddef
  if (tree%n_branches /= tree%n_externals*2-3) then
     call phs_tree_write (tree)
     call msg_bug &
          & (" Wrong number of branches set in phase space tree")
  end if
\nwendcode{}\nwbegindocs{412}For all branches that are set, except for the externals, we try to
find the daughter branches:
\nwenddocs{}\nwbegincode{413}\moddef{Determine the connections}\endmoddef
  do k=1, size (tree%branch)
     if (tree%branch(k)%set .and. tc_decay_level (k) /= 1) then
        call branch_set_relatives(k)
     end if
  end do
\nwendcode{}\nwbegindocs{414}To this end, we scan all codes less than the current code, whether
we can find two branches which are set and which together give the
current code.  After that, the tree may still not be connected, but at
least we know if a branch does not have daughters: This indicates some
inconsistency.

The algorithm ensures that, at this stage, the first daughter has a
smaller code value than the second one.
\nwenddocs{}\nwbegincode{415}\moddef{Subroutine: set relatives}\endmoddef
  subroutine branch_set_relatives (k)
    integer(TC), intent(in) :: k
    integer(TC) :: m,n
    do m=1, k-1
       if(iand(k,m)==m) then
          n = ieor(k,m)
          if ( tree%branch(m)%set .and. tree%branch(n)%set ) then
             tree%branch(k)%daughter(1) = m;  tree%branch(k)%daughter(2) = n
             tree%branch(m)%mother      = k;  tree%branch(n)%mother      = k
             tree%branch(m)%sibling     = n;  tree%branch(n)%sibling     = m
             tree%branch(k)%has_children = .true.
             return
          end if
       end if
    end do
    call phs_tree_write (tree)
    call msg_bug &
         & (" Missing daughter branch(es) in phase space tree")
  end subroutine branch_set_relatives

\nwendcode{}\nwbegindocs{416}The inverse: this is trivial, fortunately.
\nwenddocs{}\nwbegindocs{417}\nwdocspar
\subsubsection{Flip $t$-channel into $s$-channel}
Flipping the tree is done upwards, beginning from the decay products.
First we select a $t$-channel branch {\tt{}k}: one which is set, which
does have an IN-bit, and which is not an external particle.

Next, we determine the adjacent in-particle (called the 'friend' {\tt{}f}
here, since it will provide the reference axis for the angular
integration).  In addition, we look for the 'mother' and 'sibling' of
this particle.  If the latter field is empty, we select the (unique)
other out-particle which has no mother, calling the internal
subroutine {\tt{}find{\char95}orphan}.  

The flip is done as follows: We assume that the first daughter {\tt{}d}
is an $s$-channel line, which is true if the daughters are sorted.
This will stay the first daughter.  The second one is a $t$-channel
line; it is exchanged with the 'sibling' {\tt{}s}.  The new line which
replaces the branch {\tt{}k} is just the sum of {\tt{}s} and {\tt{}d}.  In
addition, we have to rearrange the relatives of {\tt{}s} and {\tt{}d}, as
well of {\tt{}f}.

Finally, we flip 'sibling' and 'friend' and set the new $s$-channel
branch {\tt{}n} which replaces the $t$-channel branch {\tt{}k}.  After this
is complete, we are ready to execute another flip.

[Although the friend is not needed for the final flip, since it would
be an initial particle anyway, we need to know whether we have $t$- or
$u$-channel.]
\nwenddocs{}\nwbegincode{418}\moddef{PHS trees: public}\plusendmoddef
  public :: phs_tree_flip_t_to_s_channel
\nwendcode{}\nwbegincode{419}\moddef{PHS trees: procedures}\plusendmoddef
  subroutine phs_tree_flip_t_to_s_channel (tree)
    type(phs_tree_t), intent(inout) :: tree
    integer(TC) :: k, f, m, n, d, s
    if (tree%n_in == 2) then
       FLIP: do k=3, tree%mask-1
          if (.not. tree%branch(k)%set) cycle FLIP
          f = iand(k,tree%mask_in)
          if (f==0 .or. f==k) cycle FLIP
          m = tree%branch(k)%mother
          s = tree%branch(k)%sibling
          if (s==0) call find_orphan(s)
          d = tree%branch(k)%daughter(1)
          n = ior(d,s)  
          tree%branch(k)%set = .false.
          tree%branch(n)%set = .true.
          tree%branch(n)%origin = k
          tree%branch(n)%daughter(1) = d; tree%branch(d)%mother  = n
          tree%branch(n)%daughter(2) = s; tree%branch(s)%mother  = n
          tree%branch(n)%has_children = .true.
          tree%branch(d)%sibling = s;  tree%branch(s)%sibling = d
          tree%branch(n)%sibling = f;  tree%branch(f)%sibling = n
          tree%branch(n)%mother      = m
          tree%branch(f)%mother      = m
          if (m/=0) then
             tree%branch(m)%daughter(1) = n
             tree%branch(m)%daughter(2) = f
          end if
          tree%branch(n)%friend = f
          tree%branch(n)%has_friend = .true.
          tree%branch(n)%firstborn = 2
       end do FLIP
    end if
  contains
    subroutine find_orphan(s)
      integer(TC) :: s
      do s=1, tree%mask_out
         if (tree%branch(s)%set .and. tree%branch(s)%mother==0) return
      end do
      call phs_tree_write (tree)
      call msg_bug (" Can't flip phase space tree to channel")
    end subroutine find_orphan
  end subroutine phs_tree_flip_t_to_s_channel

\eatline
\nwendcode{}\nwbegindocs{420}After the tree has been flipped, one may need to determine what has
become of a particular $t$-channel branch.  This function gives the
bincode of the flipped tree.  If the original bincode does not contain
IN-bits, we leave it as it is.
\nwenddocs{}\nwbegincode{421}\moddef{PHS trees: procedures}\plusendmoddef
  function tc_flipped (tree, kt) result (ks)
    type(phs_tree_t), intent(in) :: tree
    integer(TC), intent(in) :: kt
    integer(TC) :: ks
    if (iand (kt, tree%mask_in) == 0) then
       ks = kt
    else
       ks = tree%branch(iand (kt, tree%mask_out))%mother
    end if
  end function tc_flipped

\eatline
\nwendcode{}\nwbegindocs{422}Scan a tree and make sure that the first daughter has always a
smaller code than the second one.  Furthermore, delete any {\tt{}friend}
entry in the root branch -- this branching has the incoming particle
direction as axis anyway.  Keep track of reordering by updating
{\tt{}inverted{\char95}axis}, {\tt{}inverted{\char95}decay} and {\tt{}firstborn}.
\nwenddocs{}\nwbegincode{423}\moddef{PHS trees: public}\plusendmoddef
  public :: phs_tree_canonicalize
\nwendcode{}\nwbegincode{424}\moddef{PHS trees: procedures}\plusendmoddef
  subroutine phs_tree_canonicalize (tree)
    type(phs_tree_t), intent(inout) :: tree
    integer :: n_out
    integer(TC) :: k_out
    call branch_canonicalize (tree%branch(tree%mask_out))
    n_out = tree%n_externals - tree%n_in
    k_out = tree%mask_out
    if (tree%branch(k_out)%has_friend &
         & .and. tree%branch(k_out)%friend == ibset (0, n_out)) then
       tree%branch(k_out)%inverted_axis = .not.tree%branch(k_out)%inverted_axis
    end if
    tree%branch(k_out)%has_friend = .false.
    tree%branch(k_out)%friend = 0
  contains
    recursive subroutine branch_canonicalize (b)
      type(phs_branch_t), intent(inout) :: b
      integer(TC) :: d1, d2
      if (b%has_children) then
         d1 = b%daughter(1)
         d2 = b%daughter(2)
         if (d1 > d2) then
            b%daughter(1) = d2
            b%daughter(2) = d1
            b%inverted_decay = .not.b%inverted_decay
            if (b%firstborn /= 0)  b%firstborn = 3 - b%firstborn
         end if
         call branch_canonicalize (tree%branch(b%daughter(1)))
         call branch_canonicalize (tree%branch(b%daughter(2)))
      end if
    end subroutine branch_canonicalize
  end subroutine phs_tree_canonicalize

\eatline
\nwendcode{}\nwbegindocs{425}\nwdocspar
\subsubsection{Mappings}
Initialize a mapping for the current tree.  This is done while reading
from file, so the mapping parameters are read, but applied to the
flipped tree.  Thus, the size of the array of mappings is given by the
number of outgoing particles only.
\nwenddocs{}\nwbegincode{426}\moddef{PHS trees: public}\plusendmoddef
  public :: phs_tree_init_mapping
\nwendcode{}\nwbegincode{427}\moddef{PHS trees: procedures}\plusendmoddef
  subroutine phs_tree_init_mapping (tree, k, type, pdg, model)
    type(phs_tree_t), intent(inout) :: tree
    integer(TC), intent(in) :: k
    type(string_t), intent(in) :: type
    integer, intent(in) :: pdg
    class(model_data_t), intent(in), target :: model
    integer(TC) :: kk
    kk = tc_flipped (tree, k)
    call mapping_init (tree%mapping(kk), kk, type, pdg, model)
  end subroutine phs_tree_init_mapping

\eatline
\nwendcode{}\nwbegindocs{428}Set the physical parameters for the mapping, using a specific
parameter set.  Also set the mass sum array.
\nwenddocs{}\nwbegincode{429}\moddef{PHS trees: public}\plusendmoddef
  public :: phs_tree_set_mapping_parameters
\nwendcode{}\nwbegincode{430}\moddef{PHS trees: procedures}\plusendmoddef
  subroutine phs_tree_set_mapping_parameters &
       (tree, mapping_defaults, variable_limits)
    type(phs_tree_t), intent(inout) :: tree
    type(mapping_defaults_t), intent(in) :: mapping_defaults
    logical, intent(in) :: variable_limits
    integer(TC) :: k
    do k = 1, tree%n_branches_out
       call mapping_set_parameters &
            (tree%mapping(k), mapping_defaults, variable_limits) 
    end do
  end subroutine phs_tree_set_mapping_parameters

\eatline
\nwendcode{}\nwbegindocs{431}Return the mapping for the sum of all outgoing particles.  This
should either be no mapping or a global s-channel mapping.
\nwenddocs{}\nwbegincode{432}\moddef{PHS trees: public}\plusendmoddef
  public :: phs_tree_assign_s_mapping
\nwendcode{}\nwbegincode{433}\moddef{PHS trees: procedures}\plusendmoddef
  subroutine phs_tree_assign_s_mapping (tree, mapping)
    type(phs_tree_t), intent(in) :: tree
    type(mapping_t), intent(out) :: mapping
    mapping = tree%mapping(tree%mask_out)
  end subroutine phs_tree_assign_s_mapping

\eatline
\nwendcode{}\nwbegindocs{434}\nwdocspar
\subsubsection{Kinematics}
Fill the mass sum array, starting from the external particles and
working down to the tree root.  For each bincode {\tt{}k} we scan the
bits in {\tt{}k}; if only one is set, we take the physical mass of the
corresponding external particle; if more then one is set, we sum up
the two masses (which we know have already been set).
\nwenddocs{}\nwbegincode{435}\moddef{PHS trees: public}\plusendmoddef
  public :: phs_tree_set_mass_sum
\nwendcode{}\nwbegincode{436}\moddef{PHS trees: procedures}\plusendmoddef
  subroutine phs_tree_set_mass_sum (tree, flv)
    type(phs_tree_t), intent(inout) :: tree
    type(flavor_t), dimension(:), intent(in) :: flv
    integer(TC) :: k
    integer :: i
    tree%mass_sum = 0
    do k = 1, tree%n_branches_out
       do i = 0, size (flv) - 1
          if (btest(k,i)) then
             if (ibclr(k,i) == 0) then
                tree%mass_sum(k) = flv(i+1)%get_mass ()
             else
                tree%mass_sum(k) = &
                     tree%mass_sum(ibclr(k,i)) + tree%mass_sum(ibset(0,i))
             end if
             exit
          end if
       end do
    end do
  end subroutine phs_tree_set_mass_sum
  
\eatline
\nwendcode{}\nwbegindocs{437}Set the effective masses and widths.  For each non-resonant branch
in a tree, the effective mass is equal to the sum of the effective
masses of the children (and analogous for the width).  External
particles have their real mass and width zero.  For resonant branches,
we insert mass and width from the corresponding mapping.

This routine has {\tt{}phs{\char95}tree{\char95}set{\char95}mass{\char95}sum} and
{\tt{}phs{\char95}tree{\char95}set{\char95}mapping{\char95}parameters} as prerequisites.
\nwenddocs{}\nwbegincode{438}\moddef{PHS trees: public}\plusendmoddef
  public :: phs_tree_set_effective_masses
\nwendcode{}\nwbegincode{439}\moddef{PHS trees: procedures}\plusendmoddef
  subroutine phs_tree_set_effective_masses (tree)
    type(phs_tree_t), intent(inout) :: tree
    tree%effective_mass = 0
    tree%effective_width = 0
    call set_masses_x (tree%mask_out)
  contains
    recursive subroutine set_masses_x (k)
      integer(TC), intent(in) :: k
      integer(TC) :: k1, k2
      if (tree%branch(k)%has_children) then
         k1 = tree%branch(k)%daughter(1)
         k2 = tree%branch(k)%daughter(2)
         call set_masses_x (k1)
         call set_masses_x (k2)
         if (mapping_is_s_channel (tree%mapping(k))) then
            tree%effective_mass(k) = mapping_get_mass (tree%mapping(k))
            tree%effective_width(k) = mapping_get_width (tree%mapping(k))
         else
            tree%effective_mass(k) = &
                 tree%effective_mass(k1) + tree%effective_mass(k2)
            tree%effective_width(k) = &
                 tree%effective_width(k1) + tree%effective_width(k2)
         end if
      else
         tree%effective_mass(k) = tree%mass_sum(k)
      end if
    end subroutine set_masses_x
  end subroutine phs_tree_set_effective_masses

\eatline
\nwendcode{}\nwbegindocs{440}Define step mappings, recursively, for the decay products of all
intermediate resonances.  Step mappings account for the fact that a
branch may originate from a resonance, which almost replaces the
upper limit on the possible invariant mass.  The step
mapping implements a smooth cutoff that interpolates between the
resonance and the real kinematic limit.  The mapping width determines
the sharpness of the cutoff.

Step mappings are inserted only for branches that are not mapped
otherwise.

At each branch, we record the mass that is effectively available for
phase space, by taking the previous limit and subtracting the
effective mass of the sibling branch.  Widths are added, not subtracted.

If we encounter a resonance decay, we discard the previous limit and
replace it by the mass and width of the resonance, also subtracting
the sibling branch.

Initially, the limit is zero, so it becomes negative at any branch.  Only
if there is a resonance, the limit becomes positive.  Whenever the
limit is positive, and the current branch decays, we activate a step
mapping for the current branch.

As a result, step mappings are implemented for all internal lines that
originate from an intermediate resonance decay.

The flag {\tt{}variable{\char95}limits} applies to the ultimate limit from the
available energy, not to the intermediate resonances whose masses are
always fixed.

This routine requires {\tt{}phs{\char95}tree{\char95}set{\char95}effective{\char95}masses}
\nwenddocs{}\nwbegincode{441}\moddef{PHS trees: public}\plusendmoddef
  public :: phs_tree_set_step_mappings
\nwendcode{}\nwbegincode{442}\moddef{PHS trees: procedures}\plusendmoddef
  subroutine phs_tree_set_step_mappings (tree, exp_type, variable_limits)
    type(phs_tree_t), intent(inout) :: tree
    logical, intent(in) :: exp_type
    logical, intent(in) :: variable_limits
    type(string_t) :: map_str
    integer(TC) :: k
    if (exp_type) then
       map_str = "step_exp"
    else
       map_str = "step_hyp"
    end if
    k = tree%mask_out
    call set_step_mappings_x (k, 0._default, 0._default)
  contains
    recursive subroutine set_step_mappings_x (k, m_limit, w_limit)
      integer(TC), intent(in) :: k
      real(default), intent(in) :: m_limit, w_limit
      integer(TC), dimension(2) :: kk
      real(default), dimension(2) :: m, w
      if (tree%branch(k)%has_children) then
         if (m_limit > 0) then
            if (.not. mapping_is_set (tree%mapping(k))) then
               call mapping_init (tree%mapping(k), k, map_str)
               call mapping_set_step_mapping_parameters (tree%mapping(k), &
                    m_limit, w_limit, &
                    variable_limits)
            end if
         end if
         kk = tree%branch(k)%daughter
         m = tree%effective_mass(kk)
         w = tree%effective_width(kk)
         if (mapping_is_s_channel (tree%mapping(k))) then
            call set_step_mappings_x (kk(1), &
                 mapping_get_mass (tree%mapping(k)) - m(2), &
                 mapping_get_width (tree%mapping(k)) + w(2))
            call set_step_mappings_x (kk(2), &
                 mapping_get_mass (tree%mapping(k)) - m(1), &
                 mapping_get_width (tree%mapping(k)) + w(1))
         else if (m_limit > 0) then
            call set_step_mappings_x (kk(1), &
                 m_limit - m(2), &
                 w_limit + w(2))
            call set_step_mappings_x (kk(2), &
                 m_limit - m(1), &
                 w_limit + w(1))
         else
            call set_step_mappings_x (kk(1), &
                 - m(2), &
                 + w(2))
            call set_step_mappings_x (kk(2), &
                 - m(1), &
                 + w(1))
         end if
      end if
    end subroutine set_step_mappings_x
  end subroutine phs_tree_set_step_mappings

\eatline
\nwendcode{}\nwbegindocs{443}\nwdocspar
\subsubsection{Structural comparison}
This function allows to check whether one tree is the permutation of
another one.  The permutation is applied to the second tree in the
argument list.  We do not make up a temporary permuted tree, but
compare the two trees directly.  The branches are scanned recursively,
where for each daughter we check the friend and the mapping as well.
Once a discrepancy is found, the recursion is exited immediately.
\nwenddocs{}\nwbegincode{444}\moddef{PHS trees: public}\plusendmoddef
  public :: phs_tree_equivalent
\nwendcode{}\nwbegincode{445}\moddef{PHS trees: procedures}\plusendmoddef
  function phs_tree_equivalent (t1, t2, perm) result (is_equal)
    type(phs_tree_t), intent(in) :: t1, t2
    type(permutation_t), intent(in) :: perm
    logical :: equal, is_equal
    integer(TC) :: k1, k2, mask_in
    k1 = t1%mask_out
    k2 = t2%mask_out
    mask_in = t1%mask_in
    equal = .true.
    call check (t1%branch(k1), t2%branch(k2), k1, k2)
    is_equal = equal
  contains
    recursive subroutine check (b1, b2, k1, k2)
      type(phs_branch_t), intent(in) :: b1, b2
      integer(TC), intent(in) :: k1, k2
      integer(TC), dimension(2) :: d1, d2, pd2
      integer :: i
      if (.not.b1%has_friend .and. .not.b2%has_friend) then
         equal = .true.
      else if (b1%has_friend .and. b2%has_friend) then
         equal = (b1%friend == tc_permute (b2%friend, perm, mask_in))
      end if
      if (equal) then
         if (b1%has_children .and. b2%has_children) then
            d1 = b1%daughter
            d2 = b2%daughter
            do i=1, 2
               pd2(i) = tc_permute (d2(i), perm, mask_in)
            end do
            if (d1(1)==pd2(1) .and. d1(2)==pd2(2)) then
               equal = (b1%firstborn == b2%firstborn)
               if (equal) call check &
                    &     (t1%branch(d1(1)), t2%branch(d2(1)), d1(1), d2(1))
               if (equal) call check &
                    &     (t1%branch(d1(2)), t2%branch(d2(2)), d1(2), d2(2))
            else if (d1(1)==pd2(2) .and. d1(2)==pd2(1)) then
               equal = ( (b1%firstborn == 0 .and. b2%firstborn == 0) &
                    &   .or. (b1%firstborn == 3 - b2%firstborn) )
               if (equal) call check &
                    &     (t1%branch(d1(1)), t2%branch(d2(2)), d1(1), d2(2))
               if (equal) call check &
                    &     (t1%branch(d1(2)), t2%branch(d2(1)), d1(2), d2(1))
            else
               equal = .false.
            end if
         end if
      end if
      if (equal) then
         equal = (t1%mapping(k1) == t2%mapping(k2))
      end if
    end subroutine check
  end function phs_tree_equivalent

\eatline
\nwendcode{}\nwbegindocs{446}Scan two decay trees and determine the correspondence of mass
variables, i.e., the permutation that transfers the ordered list of
mass variables belonging to the second tree into the first one.  Mass
variables are assigned beginning from branches and ending at the root.
\nwenddocs{}\nwbegincode{447}\moddef{PHS trees: public}\plusendmoddef
  public :: phs_tree_find_msq_permutation
\nwendcode{}\nwbegincode{448}\moddef{PHS trees: procedures}\plusendmoddef
  subroutine phs_tree_find_msq_permutation (tree1, tree2, perm2, msq_perm)
    type(phs_tree_t), intent(in) :: tree1, tree2
    type(permutation_t), intent(in) :: perm2
    type(permutation_t), intent(out) :: msq_perm
    type(permutation_t) :: perm1
    integer(TC) :: mask_in, root
    integer(TC), dimension(:), allocatable :: index1, index2
    integer :: i
    allocate (index1 (tree1%n_msq), index2 (tree2%n_msq))
    call permutation_init (perm1, permutation_size (perm2))
    mask_in = tree1%mask_in
    root = tree1%mask_out
    i = 0
    call tree_scan (tree1, root, perm1, index1)
    i = 0
    call tree_scan (tree2, root, perm2, index2)
    call permutation_find (msq_perm, index1, index2)
  contains
    recursive subroutine tree_scan (tree, k, perm, index)
      type(phs_tree_t), intent(in) :: tree
      integer(TC), intent(in) :: k
      type(permutation_t), intent(in) :: perm
      integer, dimension(:), intent(inout) :: index
      if (tree%branch(k)%has_children) then
         call tree_scan (tree, tree%branch(k)%daughter(1), perm, index)
         call tree_scan (tree, tree%branch(k)%daughter(2), perm, index)
         i = i + 1
         if (i <= size (index))  index(i) = tc_permute (k, perm, mask_in)
      end if
    end subroutine tree_scan
  end subroutine phs_tree_find_msq_permutation

\eatline
\nwendcode{}\nwbegincode{449}\moddef{PHS trees: public}\plusendmoddef
  public :: phs_tree_find_angle_permutation
\nwendcode{}\nwbegincode{450}\moddef{PHS trees: procedures}\plusendmoddef
  subroutine phs_tree_find_angle_permutation &
       (tree1, tree2, perm2, angle_perm, sig2)
    type(phs_tree_t), intent(in) :: tree1, tree2
    type(permutation_t), intent(in) :: perm2
    type(permutation_t), intent(out) :: angle_perm
    logical, dimension(:), allocatable, intent(out) :: sig2
    type(permutation_t) :: perm1
    integer(TC) :: mask_in, root
    integer(TC), dimension(:), allocatable :: index1, index2
    logical, dimension(:), allocatable :: sig1
    integer :: i
    allocate (index1 (tree1%n_angles), index2 (tree2%n_angles))
    allocate (sig1 (tree1%n_angles), sig2 (tree2%n_angles))
    call permutation_init (perm1, permutation_size (perm2))
    mask_in = tree1%mask_in
    root = tree1%mask_out
    i = 0
    call tree_scan (tree1, root, perm1, index1, sig1)
    i = 0
    call tree_scan (tree2, root, perm2, index2, sig2)
    call permutation_find (angle_perm, index1, index2)
  contains
    recursive subroutine tree_scan (tree, k, perm, index, sig)
      type(phs_tree_t), intent(in) :: tree
      integer(TC), intent(in) :: k
      type(permutation_t), intent(in) :: perm
      integer, dimension(:), intent(inout) :: index
      logical, dimension(:), intent(inout) :: sig
      integer(TC) :: k1, k2, kp
      logical :: s
      if (tree%branch(k)%has_children) then
         k1 = tree%branch(k)%daughter(1)
         k2 = tree%branch(k)%daughter(2)
         s = (tc_permute(k1, perm, mask_in) < tc_permute(k2, perm, mask_in))
         kp = tc_permute (k, perm, mask_in)
         i = i + 1
         index(i) = kp
         sig(i) = s
         i = i + 1
         index(i) = - kp
         sig(i) = s
         call tree_scan (tree, k1, perm, index, sig)
         call tree_scan (tree, k2, perm, index, sig)
      end if
    end subroutine tree_scan
  end subroutine phs_tree_find_angle_permutation

\eatline
\nwendcode{}\nwbegindocs{451}\nwdocspar
\subsection{Phase-space evaluation} 
\subsubsection{Phase-space volume}
We compute the phase-space volume recursively, following the same path
as for computing other kinematical variables.  However, the volume
depends just on $\sqrt{\hat s}$, not on the momentum configuration.

Note: counting branches, we may replace this by a simple formula.
\nwenddocs{}\nwbegincode{452}\moddef{PHS trees: public}\plusendmoddef
  public :: phs_tree_compute_volume
\nwendcode{}\nwbegincode{453}\moddef{PHS trees: procedures}\plusendmoddef
  subroutine phs_tree_compute_volume (tree, sqrts, volume)
    type(phs_tree_t), intent(in) :: tree
    real(default), intent(in) :: sqrts
    real(default), intent(out) :: volume
    integer(TC) :: k
    k  = tree%mask_out
    if (tree%branch(k)%has_children) then
       call compute_volume_x (tree%branch(k), k, volume, .true.)
    else
       volume = 1
    end if
  contains
    recursive subroutine compute_volume_x (b, k, volume, initial)
      type(phs_branch_t), intent(in) :: b
      integer(TC), intent(in) :: k
      real(default), intent(out) :: volume
      logical, intent(in) :: initial
      integer(TC) :: k1, k2
      real(default) :: v1, v2
      k1 = b%daughter(1);  k2 = b%daughter(2)
      if (tree%branch(k1)%has_children) then
         call compute_volume_x (tree%branch(k1), k1, v1, .false.)
      else
         v1 = 1
      end if
      if (tree%branch(k2)%has_children) then
         call compute_volume_x (tree%branch(k2), k2, v2, .false.)
      else
         v2 = 1
      end if
      if (initial) then
         volume = v1 * v2 / (4 * twopi5)
      else
         volume = v1 * v2 * sqrts**2 / (4 * twopi2)
      end if
    end subroutine compute_volume_x
  end subroutine phs_tree_compute_volume

\eatline
\nwendcode{}\nwbegindocs{454}\nwdocspar
\subsubsection{Determine momenta}
This is done in two steps: First the masses are determined.  This step
may fail, in which case {\tt{}ok} is set to false.  If successful, we
generate angles and the actual momenta.  The array {\tt{}decay{\char95}p} serves
for transferring the individual three-momenta of the daughter
particles in their mother rest frame from the mass generation to the
momentum generation step.
\nwenddocs{}\nwbegincode{455}\moddef{PHS trees: public}\plusendmoddef
  public :: phs_tree_compute_momenta_from_x
\nwendcode{}\nwbegincode{456}\moddef{PHS trees: procedures}\plusendmoddef
  subroutine phs_tree_compute_momenta_from_x &
       (tree, prt, factor, volume, sqrts, x, ok)
    type(phs_tree_t), intent(inout) :: tree
    type(phs_prt_t), dimension(:), intent(inout) :: prt
    real(default), intent(out) :: factor, volume
    real(default), intent(in) :: sqrts
    real(default), dimension(:), intent(in) :: x
    logical, intent(out) :: ok
    real(default), dimension(tree%mask_out) :: decay_p
    integer :: n1, n2
    integer :: n_out
    if (tree%real_phsp) then
      n_out = tree%n_externals - tree%n_in - 1
      n1 = max (n_out-2, 0)
      n2 = n1 + max (2*n_out, 0)
    else
      n1 = tree%n_msq
      n2 = n1 + tree%n_angles
    end if
    call phs_tree_set_msq &
         (tree, prt, factor, volume, decay_p, sqrts, x(1:n1), ok)
    if (ok) call phs_tree_set_angles &
         (tree, prt, factor, decay_p, sqrts, x(n1+1:n2))
  end subroutine phs_tree_compute_momenta_from_x

\eatline
\nwendcode{}\nwbegindocs{457}Mass generation is done recursively.  The {\tt{}ok} flag causes the
filled tree to be discarded if set to {\tt{}.false.}.  This happens if a
three-momentum turns out to be imaginary, indicating impossible
kinematics.  The index {\tt{}ix} tells us how far we have used up the
input array {\tt{}x}.
\nwenddocs{}\nwbegincode{458}\moddef{PHS trees: procedures}\plusendmoddef
  subroutine phs_tree_set_msq &
       (tree, prt, factor, volume, decay_p, sqrts, x, ok)
    type(phs_tree_t), intent(inout) :: tree
    type(phs_prt_t), dimension(:), intent(inout) :: prt
    real(default), intent(out) :: factor, volume
    real(default), dimension(:), intent(out) :: decay_p
    real(default), intent(in) :: sqrts
    real(default), dimension(:), intent(in) :: x
    logical, intent(out) :: ok
    integer :: ix
    integer(TC) :: k
    real(default) :: m_tot
    ok =.true.
    ix = 1
    k  = tree%mask_out
    m_tot = tree%mass_sum(k)
    decay_p(k) = 0.
    if (m_tot < sqrts .or. k == 1) then
       if (tree%branch(k)%has_children) then
          call set_msq_x (tree%branch(k), k, factor, volume, .true.)
       else
          factor = 1
          volume = 1
       end if
    else
       ok = .false.
    end if
  contains
    recursive subroutine set_msq_x (b, k, factor, volume, initial)
      type(phs_branch_t), intent(in) :: b
      integer(TC), intent(in) :: k
      real(default), intent(out) :: factor, volume
      logical, intent(in) :: initial
      real(default) :: msq, m, m_min, m_max, m1, m2, msq1, msq2, lda, rlda
      integer(TC) :: k1, k2
      real(default) :: f1, f2, v1, v2
      k1 = b%daughter(1);  k2 = b%daughter(2)
      if (tree%branch(k1)%has_children) then
         call set_msq_x (tree%branch(k1), k1, f1, v1, .false.)
         if (.not.ok) return
      else
         f1 = 1;  v1 = 1
      end if
      if (tree%branch(k2)%has_children) then
         call set_msq_x (tree%branch(k2), k2, f2, v2, .false.)
         if (.not.ok) return
      else
         f2 = 1;  v2 = 1
      end if
      m_min = tree%mass_sum(k)
      if (initial) then
         msq = sqrts**2
         m = sqrts
         m_max = sqrts
         factor = f1 * f2
         volume = v1 * v2 / (4 * twopi5)
      else
         m_max = sqrts - m_tot + m_min
         call mapping_compute_msq_from_x &
              (tree%mapping(k), sqrts**2, m_min**2, m_max**2, msq, factor, &
               x(ix)); ix = ix + 1
         if (msq >= 0) then
            m = sqrt (msq)
            factor = f1 * f2 * factor
            volume = v1 * v2 * sqrts**2 / (4 * twopi2)
            call phs_prt_set_msq (prt(k), msq)
            call phs_prt_set_defined (prt(k))
         else
            ok = .false.
         end if
      end if
      if (ok) then
         msq1 = phs_prt_get_msq (prt(k1));  m1 = sqrt (msq1)
         msq2 = phs_prt_get_msq (prt(k2));  m2 = sqrt (msq2)
         lda = lambda (msq, msq1, msq2)
         if (lda > 0 .and. m > m1 + m2 .and. m <= m_max) then
            rlda = sqrt (lda)
            decay_p(k1) = rlda / (2*m)
            decay_p(k2) = - decay_p(k1)
            factor = rlda / msq * factor
         else
            ok = .false.
         end if
      end if
    end subroutine set_msq_x

  end subroutine phs_tree_set_msq

\eatline
\nwendcode{}\nwbegindocs{459}\nwdocspar
The heart of phase space generation: Now we have the invariant masses,
let us generate angles.  At each branch, we take a Lorentz
transformation and augment it by a boost to the current particle
rest frame, and by rotations $\phi$ and $\theta$ around the $z$ and
$y$ axis, respectively.  This transformation is passed down to the
daughter particles, if present.
\nwenddocs{}\nwbegincode{460}\moddef{PHS trees: procedures}\plusendmoddef
  subroutine phs_tree_set_angles (tree, prt, factor, decay_p, sqrts, x)
    type(phs_tree_t), intent(inout) :: tree
    type(phs_prt_t), dimension(:), intent(inout) :: prt
    real(default), intent(inout) :: factor
    real(default), dimension(:), intent(in) :: decay_p
    real(default), intent(in) :: sqrts
    real(default), dimension(:), intent(in) :: x
    integer :: ix
    integer(TC) :: k
    ix = 1
    k  = tree%mask_out
    call set_angles_x (tree%branch(k), k)
  contains
    recursive subroutine set_angles_x (b, k, L0)
      type(phs_branch_t), intent(in) :: b
      integer(TC), intent(in) :: k
      type(lorentz_transformation_t), intent(in), optional :: L0
      real(default) :: m, msq, ct, st, phi, f, E, p, bg
      type(lorentz_transformation_t) :: L, LL
      integer(TC) :: k1, k2
      type(vector3_t) :: axis
      p = decay_p(k)
      msq = phs_prt_get_msq (prt(k));  m = sqrt (msq)
      E = sqrt (msq + p**2)
      if (present (L0)) then
         call phs_prt_set_momentum (prt(k), L0 * vector4_moving (E,p,3))
      else
         call phs_prt_set_momentum (prt(k), vector4_moving (E,p,3))
      end if
      call phs_prt_set_defined (prt(k))
      if (b%has_children) then
         k1 = b%daughter(1)
         k2 = b%daughter(2)
         if (m > 0) then
            bg = p / m
         else
            bg = 0
         end if
         phi = x(ix) * twopi;  ix = ix + 1
         call mapping_compute_ct_from_x &
              (tree%mapping(k), sqrts**2, ct, st, f, x(ix));  ix = ix + 1
         factor = factor * f
         if (.not. b%has_friend) then
            L = LT_compose_r2_r3_b3 (ct, st, cos(phi), sin(phi), bg)
            !!! The function above is equivalent to:
            ! L = boost (bg,3) * rotation (phi,3) * rotation (ct,st,2)
         else
            LL = boost (-bg,3);  if (present (L0))  LL = LL * inverse(L0)
            axis = space_part ( &
                 LL * phs_prt_get_momentum (prt(tree%branch(k)%friend)) )
            L = boost(bg,3) * rotation_to_2nd (vector3_canonical(3), axis) &
                 * LT_compose_r2_r3_b3 (ct, st, cos(phi), sin(phi), 0._default)
         end if
         if (present (L0))  L = L0 * L
         call set_angles_x (tree%branch(k1), k1, L)
         call set_angles_x (tree%branch(k2), k2, L)
      end if
    end subroutine set_angles_x

  end subroutine phs_tree_set_angles

\eatline
\nwendcode{}\nwbegindocs{461}\nwdocspar
\subsubsection{Recover random numbers}
For the other channels we want to compute the random numbers that
would have generated the momenta that we already know.
\nwenddocs{}\nwbegincode{462}\moddef{PHS trees: public}\plusendmoddef
  public :: phs_tree_compute_x_from_momenta
\nwendcode{}\nwbegincode{463}\moddef{PHS trees: procedures}\plusendmoddef
  subroutine phs_tree_compute_x_from_momenta (tree, prt, factor, sqrts, x)
    type(phs_tree_t), intent(inout) :: tree
    type(phs_prt_t), dimension(:), intent(in) :: prt
    real(default), intent(out) :: factor
    real(default), intent(in) :: sqrts
    real(default), dimension(:), intent(inout) :: x
    real(default), dimension(tree%mask_out) :: decay_p
    integer :: n1, n2
    n1 = tree%n_msq
    n2 = n1 + tree%n_angles
    call phs_tree_get_msq &
         (tree, prt, factor, decay_p, sqrts, x(1:n1))
    call phs_tree_get_angles &
         (tree, prt, factor, decay_p, sqrts, x(n1+1:n2))
  end subroutine phs_tree_compute_x_from_momenta

\eatline
\nwendcode{}\nwbegindocs{464}The inverse operation follows exactly the same steps.  The tree is
{\tt{}inout} because it contains mappings whose parameters can be reset
when the mapping is applied.
\nwenddocs{}\nwbegincode{465}\moddef{PHS trees: procedures}\plusendmoddef
  subroutine phs_tree_get_msq (tree, prt, factor, decay_p, sqrts, x)
    type(phs_tree_t), intent(inout) :: tree
    type(phs_prt_t), dimension(:), intent(in) :: prt
    real(default), intent(out) :: factor
    real(default), dimension(:), intent(out) :: decay_p
    real(default), intent(in) :: sqrts
    real(default), dimension(:), intent(inout) :: x
    integer :: ix
    integer(TC) :: k
    real(default) :: m_tot
    ix = 1
    k  = tree%mask_out
    m_tot = tree%mass_sum(k)
    decay_p(k) = 0.
    if (tree%branch(k)%has_children) then
       call get_msq_x (tree%branch(k), k, factor, .true.)
    else
       factor = 1
    end if
  contains
    recursive subroutine get_msq_x (b, k, factor, initial)
      type(phs_branch_t), intent(in) :: b
      integer(TC), intent(in) :: k
      real(default), intent(out) :: factor
      logical, intent(in) :: initial
      real(default) :: msq, m, m_min, m_max, msq1, msq2, lda, rlda
      integer(TC) :: k1, k2
      real(default) :: f1, f2
      k1 = b%daughter(1);  k2 = b%daughter(2)
      if (tree%branch(k1)%has_children) then
         call get_msq_x (tree%branch(k1), k1, f1, .false.)
      else
         f1 = 1
      end if
      if (tree%branch(k2)%has_children) then
         call get_msq_x (tree%branch(k2), k2, f2, .false.)
      else
         f2 = 1
      end if
      m_min = tree%mass_sum(k)
      m_max = sqrts - m_tot + m_min
      msq = phs_prt_get_msq (prt(k));  m = sqrt (msq)
      if (initial) then
         factor = f1 * f2
      else
         call mapping_compute_x_from_msq &
              (tree%mapping(k), sqrts**2, m_min**2, m_max**2, msq, factor, &
               x(ix));  ix = ix + 1
         factor = f1 * f2 * factor
      end if
      msq1 = phs_prt_get_msq (prt(k1))
      msq2 = phs_prt_get_msq (prt(k2))  
      lda = lambda (msq, msq1, msq2)
      if (lda > 0) then
         rlda = sqrt (lda)
         decay_p(k1) = rlda / (2 * m)
         decay_p(k2) = - decay_p(k1)
         factor = rlda / msq * factor
      else
         decay_p(k1) = 0
         decay_p(k2) = 0
         factor = 0
      end if
    end subroutine get_msq_x

  end subroutine phs_tree_get_msq

\eatline
\nwendcode{}\nwbegindocs{466}This subroutine is the most time-critical part of the whole
program.  Therefore, we do not exactly parallel the angle generation
routine above but make sure that things get evaluated only if they are
really needed, at the expense of readability.  Particularly important
is to have as few multiplications of Lorentz transformations as
possible.
\nwenddocs{}\nwbegincode{467}\moddef{PHS trees: procedures}\plusendmoddef
  subroutine phs_tree_get_angles (tree, prt, factor, decay_p, sqrts, x)
    type(phs_tree_t), intent(inout) :: tree
    type(phs_prt_t), dimension(:), intent(in) :: prt
    real(default), intent(inout) :: factor
    real(default), dimension(:), intent(in) :: decay_p
    real(default), intent(in) :: sqrts
    real(default), dimension(:), intent(out) :: x
    integer :: ix
    integer(TC) :: k
    ix = 1
    k  = tree%mask_out
    if (tree%branch(k)%has_children) then
       call get_angles_x (tree%branch(k), k)
    end if
  contains
    recursive subroutine get_angles_x (b, k, ct0, st0, phi0, L0)
      type(phs_branch_t), intent(in) :: b
      integer(TC), intent(in) :: k
      real(default), intent(in), optional :: ct0, st0, phi0
      type(lorentz_transformation_t), intent(in), optional :: L0
      real(default) :: cp0, sp0, m, msq, ct, st, phi, bg, f
      type(lorentz_transformation_t) :: L, LL
      type(vector4_t) :: p1, pf
      type(vector3_t) :: n, axis
      integer(TC) :: k1, k2, kf
      logical :: has_friend, need_L
      k1 = b%daughter(1)
      k2 = b%daughter(2)
      kf = b%friend
      has_friend = b%has_friend
      if (present(L0)) then
         p1 = L0 * phs_prt_get_momentum (prt(k1))
         if (has_friend)  pf = L0 * phs_prt_get_momentum (prt(kf))
      else
         p1 = phs_prt_get_momentum (prt(k1))
         if (has_friend)  pf = phs_prt_get_momentum (prt(kf))
      end if
      if (present(phi0)) then
         cp0 = cos (phi0)
         sp0 = sin (phi0)
      end if
      msq = phs_prt_get_msq (prt(k));  m = sqrt (msq)
      if (m > 0) then
         bg = decay_p(k) / m
      else
         bg = 0
      end if
      if (has_friend) then
         if (present (phi0)) then
            axis = axis_from_p_r3_r2_b3 (pf, cp0, -sp0, ct0, -st0, -bg)
            LL = rotation_to_2nd (axis, vector3_canonical (3)) &
                 * LT_compose_r3_r2_b3 (cp0, -sp0, ct0, -st0, -bg)
         else
            axis = axis_from_p_b3 (pf, -bg)
            LL = rotation_to_2nd (axis, vector3_canonical(3))
            if (bg /= 0)  LL = LL * boost(-bg, 3)
         end if
         n = space_part (LL * p1)
      else if (present (phi0)) then
         n = axis_from_p_r3_r2_b3 (p1, cp0, -sp0, ct0, -st0, -bg)
      else
         n = axis_from_p_b3 (p1, -bg)
      end if
      phi = azimuthal_angle (n)
      x(ix) = phi / twopi;  ix = ix + 1
      ct = polar_angle_ct (n)
      st = sqrt (1 - ct**2)
      call mapping_compute_x_from_ct (tree%mapping(k), sqrts**2, ct, f, &
           x(ix));  ix = ix + 1
      factor = factor * f
      if (tree%branch(k1)%has_children .or. tree%branch(k2)%has_children) then
         need_L = .true.
         if (has_friend) then
            if (present (L0)) then
               L = LL * L0
            else
               L = LL
            end if
         else if (present (L0)) then
            L = LT_compose_r3_r2_b3 (cp0, -sp0, ct0, -st0, -bg) * L0
         else if (present (phi0)) then
            L = LT_compose_r3_r2_b3 (cp0, -sp0, ct0, -st0, -bg)
         else if (bg /= 0) then
            L = boost(-bg, 3)
         else
            need_L = .false.
         end if
         if (need_L) then
            if (tree%branch(k1)%has_children) &
                 call get_angles_x (tree%branch(k1), k1, ct, st, phi, L)
            if (tree%branch(k2)%has_children) &
                 call get_angles_x (tree%branch(k2), k2, ct, st, phi, L)
         else
            if (tree%branch(k1)%has_children) &
                 call get_angles_x (tree%branch(k1), k1, ct, st, phi)
            if (tree%branch(k2)%has_children) &
                 call get_angles_x (tree%branch(k2), k2, ct, st, phi)
         end if
      end if
    end subroutine get_angles_x
  end subroutine phs_tree_get_angles

\eatline
\nwendcode{}\nwbegindocs{468}\nwdocspar
\subsubsection{Auxiliary stuff}
This calculates all momenta that are not yet known by summing up
daughter particle momenta.  The external particles must be known.
Only composite particles not yet known are calculated.
\nwenddocs{}\nwbegincode{469}\moddef{PHS trees: public}\plusendmoddef
  public :: phs_tree_combine_particles
\nwendcode{}\nwbegincode{470}\moddef{PHS trees: procedures}\plusendmoddef
  subroutine phs_tree_combine_particles (tree, prt)
    type(phs_tree_t), intent(in) :: tree
    type(phs_prt_t), dimension(:), intent(inout) :: prt
    call combine_particles_x (tree%mask_out)
  contains
    recursive subroutine combine_particles_x (k)
      integer(TC), intent(in) :: k
      integer :: k1, k2
      if (tree%branch(k)%has_children) then
         k1 = tree%branch(k)%daughter(1);  k2 = tree%branch(k)%daughter(2)
         call combine_particles_x (k1)
         call combine_particles_x (k2)
         if (.not. prt(k)%defined) then
            call phs_prt_combine (prt(k), prt(k1), prt(k2))
         end if
      end if
    end subroutine combine_particles_x
  end subroutine phs_tree_combine_particles

\eatline
\nwendcode{}\nwbegindocs{471}The previous routine is to be evaluated at runtime.  Instead of
scanning trees, we can as well set up a multiplication table.  This is
generated here.  Note that the table is {\tt{}intent(out)}.
\nwenddocs{}\nwbegincode{472}\moddef{PHS trees: public}\plusendmoddef
  public :: phs_tree_setup_prt_combinations
\nwendcode{}\nwbegincode{473}\moddef{PHS trees: procedures}\plusendmoddef
  subroutine phs_tree_setup_prt_combinations (tree, comb)
    type(phs_tree_t), intent(in) :: tree
    integer, dimension(:,:), intent(out) :: comb
    comb = 0
    call setup_prt_combinations_x (tree%mask_out)
  contains
    recursive subroutine setup_prt_combinations_x (k)
      integer(TC), intent(in) :: k
      integer, dimension(2) :: kk
      if (tree%branch(k)%has_children) then
         kk = tree%branch(k)%daughter
         call setup_prt_combinations_x (kk(1))
         call setup_prt_combinations_x (kk(2))
         comb(:,k) = kk
      end if
    end subroutine setup_prt_combinations_x
  end subroutine phs_tree_setup_prt_combinations
      
\eatline
\nwendcode{}\nwbegindocs{474}\nwdocspar
\nwenddocs{}\nwbegincode{475}\moddef{PHS trees: public}\plusendmoddef
  public :: phs_tree_reshuffle_mappings
\nwendcode{}\nwbegincode{476}\moddef{PHS trees: procedures}\plusendmoddef
  subroutine phs_tree_reshuffle_mappings (tree)
   type(phs_tree_t), intent(inout) :: tree
   integer(TC) :: k0, k_old, k_new, k2
   integer :: i
   type(mapping_t) :: mapping_tmp
   real(default) :: mass_tmp
   do i = 1, size (tree%momentum_link)
     if (i /= tree%momentum_link (i)) then
       k_old = 2**(i-tree%n_in-1)
       k_new = 2**(tree%momentum_link(i)-tree%n_in-1)
       k0 = tree%branch(k_old)%mother
       k2 = k_new + tree%branch(k_old)%sibling
       mapping_tmp = tree%mapping(k0)
       mass_tmp = tree%mass_sum(k0)
       tree%mapping(k0) = tree%mapping(k2)
       tree%mapping(k2) = mapping_tmp
       tree%mass_sum(k0) = tree%mass_sum(k2)
       tree%mass_sum(k2) = mass_tmp
     end if
   end do
  end subroutine phs_tree_reshuffle_mappings

\eatline
\nwendcode{}\nwbegindocs{477}\nwdocspar
\nwenddocs{}\nwbegincode{478}\moddef{PHS trees: public}\plusendmoddef
  public :: phs_tree_set_momentum_links
\nwendcode{}\nwbegincode{479}\moddef{PHS trees: procedures}\plusendmoddef
  subroutine phs_tree_set_momentum_links (tree, list)
    type(phs_tree_t), intent(inout) :: tree
    integer, dimension(:), allocatable :: list
    tree%momentum_link = list
  end subroutine phs_tree_set_momentum_links

\eatline
\nwendcode{}\nwbegindocs{480}\nwdocspar
\clearpage
%%%%%%%%%%%%%%%%%%%%%%%%%%%%%%%%%%%%%%%%%%%%%%%%%%%%%%%%%%%%%%%%%%%%%%%%
\section{The phase-space forest}
Simply stated, a phase-space forest is a collection of phase-space
trees.  More precisely, a {\tt{}phs{\char95}forest} object contains all
parameterizations of phase space that \whizard\ will use for a single
hard process, prepared in the form of {\tt{}phs{\char95}tree} objects.  This is
suitable for evaluation by the \vamp\ integration package:  each
parameterization (tree) is a valid channel in the multi-channel
adaptive integration, and each variable in a tree corresponds to an
integration dimension, defined by an appropriate mapping of the
$(0,1)$ interval to the allowed range of the integration variable.

The trees are grouped in groves.  The trees (integration channels)
within a grove share a common weight, assuming that they are related
by some approximate symmetry.

Trees/channels that are related by an exact symmetry are connected by
an array of equivalences; each equivalence object holds the data that
relate one channel to another.

The phase-space setup, i.e., the detailed structure of trees and
forest, are read from a file.  Therefore, this module also contains
the syntax definition and the parser needed for interpreting this
file.
\nwenddocs{}\nwbegincode{481}\moddef{\code{}phs{\char95}forests.f90\edoc{}}\endmoddef
\LA{}File header\RA{}

module phs_forests

\LA{}Use kinds\RA{}
  use kinds, only: TC
\LA{}Use strings\RA{}
  use io_units
  use format_defs, only: FMT_12, FMT_19
  use diagnostics
  use lorentz
  use unit_tests
  use permutations
  use ifiles
  use syntax_rules
  use lexers
  use parser
  use model_data
  use model_data
  use flavors
  use interactions

  use phs_base
  use mappings
  use phs_trees

\LA{}Standard module head\RA{}

\LA{}PHS forests: public\RA{}

\LA{}PHS forests: types\RA{}

\LA{}PHS forests: interfaces\RA{}

\LA{}PHS forests: variables\RA{}

contains

\LA{}PHS forests: procedures\RA{}

\LA{}PHS forests: tests\RA{}

end module phs_forests
\eatline
\nwendcode{}\nwbegindocs{482}\nwdocspar
\subsection{Phase-space setup parameters}
This transparent container holds the parameters that the algorithm
needs for phase-space setup, with reasonable defaults.

The threshold mass (for considering a particle as effectively
massless) is specified separately for s- and t-channel. The default is
to treat $W$ and $Z$ bosons as massive in the s-channel, but as
massless in the t-channel.  The $b$-quark is treated always massless,
the $t$-quark always massive.
\nwenddocs{}\nwbegincode{483}\moddef{PHS forests: public}\endmoddef
  public :: phs_parameters_t
\nwendcode{}\nwbegincode{484}\moddef{PHS forests: types}\endmoddef
  type :: phs_parameters_t
     real(default) :: sqrts = 0
     real(default) :: m_threshold_s = 50._default
     real(default) :: m_threshold_t = 100._default
     integer :: off_shell = 1
     integer :: t_channel = 2
     logical :: keep_nonresonant = .true.
  end type phs_parameters_t

\eatline
\nwendcode{}\nwbegindocs{485}Write phase-space parameters to file.
\nwenddocs{}\nwbegincode{486}\moddef{PHS forests: public}\plusendmoddef
  public :: phs_parameters_write
\nwendcode{}\nwbegincode{487}\moddef{PHS forests: procedures}\endmoddef
  subroutine phs_parameters_write (phs_par, unit)
    type(phs_parameters_t), intent(in) :: phs_par
    integer, intent(in), optional :: unit
    integer :: u
    u = given_output_unit (unit)
    write (u, "(3x,A," // FMT_19 // ")") "sqrts         = ", phs_par%sqrts
    write (u, "(3x,A," // FMT_19 // ")") "m_threshold_s = ", phs_par%m_threshold_s
    write (u, "(3x,A," // FMT_19 // ")") "m_threshold_t = ", phs_par%m_threshold_t
    write (u, "(3x,A,I0)") "off_shell = ", phs_par%off_shell
    write (u, "(3x,A,I0)") "t_channel = ", phs_par%t_channel
    write (u, "(3x,A,L1)") "keep_nonresonant = ", phs_par%keep_nonresonant
  end subroutine phs_parameters_write

\eatline
\nwendcode{}\nwbegindocs{488}Read phase-space parameters from file.
\nwenddocs{}\nwbegincode{489}\moddef{PHS forests: public}\plusendmoddef
  public :: phs_parameters_read
\nwendcode{}\nwbegincode{490}\moddef{PHS forests: procedures}\plusendmoddef
  subroutine phs_parameters_read (phs_par, unit)
    type(phs_parameters_t), intent(out) :: phs_par
    integer, intent(in) :: unit
    character(20) :: dummy
    character :: equals
    read (unit, *)  dummy, equals, phs_par%sqrts
    read (unit, *)  dummy, equals, phs_par%m_threshold_s
    read (unit, *)  dummy, equals, phs_par%m_threshold_t
    read (unit, *)  dummy, equals, phs_par%off_shell
    read (unit, *)  dummy, equals, phs_par%t_channel
    read (unit, *)  dummy, equals, phs_par%keep_nonresonant
  end subroutine phs_parameters_read

\eatline
\nwendcode{}\nwbegindocs{491}Comparison.
\nwenddocs{}\nwbegincode{492}\moddef{PHS forests: interfaces}\endmoddef
  interface operator(==)
     module procedure phs_parameters_eq
  end interface
  interface operator(/=)
     module procedure phs_parameters_ne
  end interface
\nwendcode{}\nwbegincode{493}\moddef{PHS forests: procedures}\plusendmoddef
  function phs_parameters_eq (phs_par1, phs_par2) result (equal)
    logical :: equal
    type(phs_parameters_t), intent(in) :: phs_par1, phs_par2
    equal = phs_par1%sqrts == phs_par2%sqrts &
         .and. phs_par1%m_threshold_s == phs_par2%m_threshold_s &
         .and. phs_par1%m_threshold_t == phs_par2%m_threshold_t &
         .and. phs_par1%off_shell == phs_par2%off_shell &
         .and. phs_par1%t_channel == phs_par2%t_channel &
         .and.(phs_par1%keep_nonresonant .eqv. phs_par2%keep_nonresonant)
  end function phs_parameters_eq

  function phs_parameters_ne (phs_par1, phs_par2) result (ne)
    logical :: ne
    type(phs_parameters_t), intent(in) :: phs_par1, phs_par2
    ne = phs_par1%sqrts /= phs_par2%sqrts &
         .or. phs_par1%m_threshold_s /= phs_par2%m_threshold_s &
         .or. phs_par1%m_threshold_t /= phs_par2%m_threshold_t &
         .or. phs_par1%off_shell /= phs_par2%off_shell &
         .or. phs_par1%t_channel /= phs_par2%t_channel &
         .or.(phs_par1%keep_nonresonant .neqv. phs_par2%keep_nonresonant)
  end function phs_parameters_ne

\eatline
\nwendcode{}\nwbegindocs{494}\nwdocspar
\subsection{Equivalences}
This type holds information about equivalences between phase-space
trees.  We make a linked list, where each node contains the two
trees which are equivalent and the corresponding permutation of
external particles.  Two more arrays are to be filled: The permutation
of mass variables and the permutation of angular variables, where the
signature indicates a necessary exchange of daughter branches.
\nwenddocs{}\nwbegincode{495}\moddef{PHS forests: types}\plusendmoddef
  type :: equivalence_t
     private
     integer :: left, right
     type(permutation_t) :: perm
     type(permutation_t) :: msq_perm, angle_perm
     logical, dimension(:), allocatable :: angle_sig
     type(equivalence_t), pointer :: next => null ()
  end type equivalence_t

\eatline
\nwendcode{}\nwbegincode{496}\moddef{PHS forests: types}\plusendmoddef
  type :: equivalence_list_t
     private
     integer :: length = 0
     type(equivalence_t), pointer :: first => null ()
     type(equivalence_t), pointer :: last => null ()
  end type equivalence_list_t

\eatline
\nwendcode{}\nwbegindocs{497}Append an equivalence to the list
\nwenddocs{}\nwbegincode{498}\moddef{PHS forests: procedures}\plusendmoddef
  subroutine equivalence_list_add (eql, left, right, perm)
    type(equivalence_list_t), intent(inout) :: eql
    integer, intent(in) :: left, right
    type(permutation_t), intent(in) :: perm
    type(equivalence_t), pointer :: eq
    allocate (eq)
    eq%left = left
    eq%right = right
    eq%perm = perm
    if (associated (eql%last)) then
       eql%last%next => eq
    else
       eql%first => eq
    end if
    eql%last => eq
    eql%length = eql%length + 1
  end subroutine equivalence_list_add

\eatline
\nwendcode{}\nwbegindocs{499}Delete the list contents.  Has to be pure because it is called from
an elemental subroutine.
\nwenddocs{}\nwbegincode{500}\moddef{PHS forests: procedures}\plusendmoddef
  pure subroutine equivalence_list_final (eql)
    type(equivalence_list_t), intent(inout) :: eql
    type(equivalence_t), pointer :: eq
    do while (associated (eql%first))
       eq => eql%first
       eql%first => eql%first%next
       deallocate (eq)
    end do
    eql%last => null ()
    eql%length = 0
  end subroutine equivalence_list_final

\eatline
\nwendcode{}\nwbegindocs{501}Make a deep copy of the equivalence list.  This allows for deep
copies of groves and forests.
\nwenddocs{}\nwbegincode{502}\moddef{PHS forests: interfaces}\plusendmoddef
  interface assignment(=)
     module procedure equivalence_list_assign
  end interface

\nwendcode{}\nwbegincode{503}\moddef{PHS forests: procedures}\plusendmoddef
  subroutine equivalence_list_assign (eql_out, eql_in)
    type(equivalence_list_t), intent(out) :: eql_out
    type(equivalence_list_t), intent(in) :: eql_in
    type(equivalence_t), pointer :: eq, eq_copy
    eq => eql_in%first
    do while (associated (eq))
       allocate (eq_copy)
       eq_copy = eq
       eq_copy%next => null ()
       if (associated (eql_out%first)) then
          eql_out%last%next => eq_copy
       else
          eql_out%first => eq_copy
       end if
       eql_out%last => eq_copy
       eq => eq%next
    end do
  end subroutine equivalence_list_assign

\eatline
\nwendcode{}\nwbegindocs{504}The number of list entries
\nwenddocs{}\nwbegincode{505}\moddef{PHS forests: procedures}\plusendmoddef
  elemental function equivalence_list_length (eql) result (length)
    integer :: length
    type(equivalence_list_t), intent(in) :: eql
    length = eql%length
  end function equivalence_list_length

\eatline
\nwendcode{}\nwbegindocs{506}Recursively write the equivalences list
\nwenddocs{}\nwbegincode{507}\moddef{PHS forests: procedures}\plusendmoddef
  subroutine equivalence_list_write (eql, unit)
    type(equivalence_list_t), intent(in) :: eql
    integer, intent(in), optional :: unit
    integer :: u
    u = given_output_unit (unit);  if (u < 0)  return
    if (associated (eql%first)) then
       call equivalence_write_rec (eql%first, u)
    else
       write (u, *) " [empty]"
    end if
  contains
    recursive subroutine equivalence_write_rec (eq, u)
      type(equivalence_t), intent(in) :: eq
      integer, intent(in) :: u
      integer :: i
      write (u, "(3x,A,1x,I0,1x,I0,2x,A)", advance="no") &
           "Equivalence:", eq%left, eq%right, "Final state permutation:"
      call permutation_write (eq%perm, u)
      write (u, "(1x,12x,1x,A,1x)", advance="no") &
           "       msq permutation:  "
      call permutation_write (eq%msq_perm, u)
      write (u, "(1x,12x,1x,A,1x)", advance="no") &
           "       angle permutation:"
      call permutation_write (eq%angle_perm, u)
      write (u, "(1x,12x,1x,26x)", advance="no")
      do i = 1, size (eq%angle_sig)
         if (eq%angle_sig(i)) then
            write (u, "(1x,A)", advance="no") "+"
         else
            write (u, "(1x,A)", advance="no") "-"
         end if
      end do
      write (u, *) 
      if (associated (eq%next))  call equivalence_write_rec (eq%next, u)
    end subroutine equivalence_write_rec
  end subroutine equivalence_list_write

\eatline
\nwendcode{}\nwbegindocs{508}\nwdocspar
\subsection{Groves}
A grove is a group of trees (phase-space channels) that share a common
weight in the integration.  Within a grove, channels can be declared
equivalent, so they also share their integration grids (up to
symmetries).  The grove contains a list of equivalences.  The
{\tt{}tree{\char95}count{\char95}offset} is the total number of trees of the preceding
groves; when the trees are counted per forest (integration channels),
the offset has to be added to all tree indices.
\nwenddocs{}\nwbegincode{509}\moddef{PHS forests: types}\plusendmoddef
  type :: phs_grove_t
     private
     integer :: tree_count_offset
     type(phs_tree_t), dimension(:), allocatable :: tree
     type(equivalence_list_t) :: equivalence_list
  end type phs_grove_t

\eatline
\nwendcode{}\nwbegindocs{510}Call {\tt{}phs{\char95}tree{\char95}init} which is also elemental:
\nwenddocs{}\nwbegincode{511}\moddef{PHS forests: procedures}\plusendmoddef
  elemental subroutine phs_grove_init &
       (grove, n_trees, n_in, n_out, n_masses, n_angles)
    type(phs_grove_t), intent(inout) :: grove
    integer, intent(in) :: n_trees, n_in, n_out, n_masses, n_angles
    grove%tree_count_offset = 0
    allocate (grove%tree (n_trees))
    call phs_tree_init (grove%tree, n_in, n_out, n_masses, n_angles)
  end subroutine phs_grove_init

\eatline
\nwendcode{}\nwbegindocs{512}The trees do not have pointer components, thus no call to
{\tt{}phs{\char95}tree{\char95}final}:
\nwenddocs{}\nwbegincode{513}\moddef{PHS forests: procedures}\plusendmoddef
  elemental subroutine phs_grove_final (grove)
    type(phs_grove_t), intent(inout) :: grove
    deallocate (grove%tree)
    call equivalence_list_final (grove%equivalence_list)
  end subroutine phs_grove_final

\eatline
\nwendcode{}\nwbegindocs{514}Deep copy.
\nwenddocs{}\nwbegincode{515}\moddef{PHS forests: interfaces}\plusendmoddef
  interface assignment(=)
     module procedure phs_grove_assign0
     module procedure phs_grove_assign1
  end interface

\nwendcode{}\nwbegincode{516}\moddef{PHS forests: procedures}\plusendmoddef
  subroutine phs_grove_assign0 (grove_out, grove_in)
    type(phs_grove_t), intent(out) :: grove_out
    type(phs_grove_t), intent(in) :: grove_in
    grove_out%tree_count_offset = grove_in%tree_count_offset
    if (allocated (grove_in%tree)) then
       allocate (grove_out%tree (size (grove_in%tree)))
       grove_out%tree = grove_in%tree
    end if
    grove_out%equivalence_list = grove_in%equivalence_list
  end subroutine phs_grove_assign0

  subroutine phs_grove_assign1 (grove_out, grove_in)
    type(phs_grove_t), dimension(:), intent(out) :: grove_out
    type(phs_grove_t), dimension(:), intent(in) :: grove_in
    integer :: i
    do i = 1, size (grove_in)
       call phs_grove_assign0 (grove_out(i), grove_in(i))
    end do
  end subroutine phs_grove_assign1

\eatline
\nwendcode{}\nwbegindocs{517}Get the global (s-channel) mappings.  Implemented as a subroutine
which returns an array (slice).
\nwenddocs{}\nwbegincode{518}\moddef{PHS forests: procedures}\plusendmoddef
  subroutine phs_grove_assign_s_mappings (grove, mapping)
    type(phs_grove_t), intent(in) :: grove
    type(mapping_t), dimension(:), intent(out) :: mapping
    integer :: i
    if (size (mapping) == size (grove%tree)) then
       do i = 1, size (mapping)
          call phs_tree_assign_s_mapping (grove%tree(i), mapping(i))
       end do
    else
       call msg_bug ("phs_grove_assign_s_mappings: array size mismatch")
    end if
  end subroutine phs_grove_assign_s_mappings

\eatline
\nwendcode{}\nwbegindocs{519}\nwdocspar
\subsection{The forest type}
This is a collection of trees and associated particles.  In a given
tree, each branch code corresponds to a particle in the {\tt{}prt} array.
Furthermore, we have an array of mass sums which is independent of the
decay tree and of the particular event.  The mappings directly
correspond to the decay trees, and the decay groves collect the trees
in classes.  The permutation list consists of all permutations of
outgoing particles that map the decay forest onto itself.

The particle codes {\tt{}flv} (one for each external particle) are needed
for determining masses and such.  The trees and associated information
are collected in the {\tt{}grove} array, together with a lookup table
that associates tree indices to groves.  Finally, the {\tt{}prt} array
serves as workspace for phase-space evaluation.

The {\tt{}prt{\char95}combination} is a list of index pairs, namely the particle
momenta pairs that need to be combined in order to provide all
momentum combinations that the phase-space trees need to know.
\nwenddocs{}\nwbegincode{520}\moddef{PHS forests: public}\plusendmoddef
  public :: phs_forest_t
\nwendcode{}\nwbegincode{521}\moddef{PHS forests: types}\plusendmoddef
  type :: phs_forest_t
     private
     integer :: n_in, n_out, n_tot
     integer :: n_masses, n_angles, n_dimensions
     integer :: n_trees, n_equivalences
     type(flavor_t), dimension(:), allocatable :: flv
     type(phs_grove_t), dimension(:), allocatable :: grove
     integer, dimension(:), allocatable :: grove_lookup
     type(phs_prt_t), dimension(:), allocatable :: prt_in
     type(phs_prt_t), dimension(:), allocatable :: prt_out
     type(phs_prt_t), dimension(:), allocatable :: prt
     integer(TC), dimension(:,:), allocatable :: prt_combination
     type(mapping_t), dimension(:), allocatable :: s_mapping
   contains
   \LA{}PHS forests: phs forest: TBP\RA{}
  end type phs_forest_t

\eatline
\nwendcode{}\nwbegindocs{522}\nwdocspar
The initialization merely allocates memory.  We have to know how many
trees there are in each grove, so we can initialize everything.  The
number of groves is the size of the {\tt{}n{\char95}tree} array.

In the {\tt{}grove{\char95}lookup} table we store the grove index that belongs to
each absolute tree index.  The difference between the absolute index
and the relative (to the grove) index is stored, for each grove, as
{\tt{}tree{\char95}count{\char95}offset}.

The particle array is allocated according to the total number of
branches each tree has, but not filled.
\nwenddocs{}\nwbegincode{523}\moddef{PHS forests: public}\plusendmoddef
  public :: phs_forest_init
\nwendcode{}\nwbegincode{524}\moddef{PHS forests: procedures}\plusendmoddef
  subroutine phs_forest_init (forest, n_tree, n_in, n_out)
    type(phs_forest_t), intent(inout) :: forest
    integer, dimension(:), intent(in) :: n_tree
    integer, intent(in) :: n_in, n_out
    integer :: g, count, k_root
    forest%n_in = n_in
    forest%n_out = n_out
    forest%n_tot = n_in + n_out
    forest%n_masses = max (n_out - 2, 0)
    forest%n_angles = max (2*n_out - 2, 0)
    forest%n_dimensions = forest%n_masses + forest%n_angles
    forest%n_trees = sum (n_tree)
    forest%n_equivalences = 0
    allocate (forest%grove (size (n_tree)))
    call phs_grove_init &
         (forest%grove, n_tree, n_in, n_out, forest%n_masses, &
          forest%n_angles)
    allocate (forest%grove_lookup (forest%n_trees))
    count = 0
    do g = 1, size (forest%grove)
       forest%grove(g)%tree_count_offset = count
       forest%grove_lookup (count+1:count+n_tree(g)) = g
       count = count + n_tree(g)
    end do
    allocate (forest%prt_in  (n_in))
    allocate (forest%prt_out (forest%n_out))
    k_root = 2**forest%n_tot - 1
    allocate (forest%prt (k_root))
    allocate (forest%prt_combination (2, k_root))
    allocate (forest%s_mapping (forest%n_trees))
  end subroutine phs_forest_init

\eatline
\nwendcode{}\nwbegindocs{525}Assign the global (s-channel) mappings.
\nwenddocs{}\nwbegincode{526}\moddef{PHS forests: public}\plusendmoddef
  public :: phs_forest_set_s_mappings
\nwendcode{}\nwbegincode{527}\moddef{PHS forests: procedures}\plusendmoddef
  subroutine phs_forest_set_s_mappings (forest)
    type(phs_forest_t), intent(inout) :: forest
    integer :: g, i0, i1, n
    do g = 1, size (forest%grove)
       call phs_forest_get_grove_bounds (forest, g, i0, i1, n)
       call phs_grove_assign_s_mappings &
            (forest%grove(g), forest%s_mapping(i0:i1))
    end do
  end subroutine phs_forest_set_s_mappings

\eatline
\nwendcode{}\nwbegindocs{528}The grove finalizer is called because it contains the equivalence list:
\nwenddocs{}\nwbegincode{529}\moddef{PHS forests: public}\plusendmoddef
  public :: phs_forest_final
\nwendcode{}\nwbegincode{530}\moddef{PHS forests: procedures}\plusendmoddef
  subroutine phs_forest_final (forest)
    type(phs_forest_t), intent(inout) :: forest
    if (allocated (forest%grove)) then
       call phs_grove_final (forest%grove)
       deallocate (forest%grove)
    end if
    if (allocated (forest%grove_lookup))  deallocate (forest%grove_lookup)
    if (allocated (forest%prt))  deallocate (forest%prt)
    if (allocated (forest%s_mapping))  deallocate (forest%s_mapping)
  end subroutine phs_forest_final

\eatline
\nwendcode{}\nwbegindocs{531}\nwdocspar
\subsection{Screen output}
Write the particles that are non-null, then the trees which point to
them:
\nwenddocs{}\nwbegincode{532}\moddef{PHS forests: public}\plusendmoddef
  public :: phs_forest_write
\nwendcode{}\nwbegincode{533}\moddef{PHS forests: phs forest: TBP}\endmoddef
  procedure :: write => phs_forest_write
\nwendcode{}\nwbegincode{534}\moddef{PHS forests: procedures}\plusendmoddef
  subroutine phs_forest_write (forest, unit)
    class(phs_forest_t), intent(in) :: forest
    integer, intent(in), optional :: unit
    integer :: u
    integer :: i, g, k
    u = given_output_unit (unit);  if (u < 0)  return
    write (u, "(1x,A)") "Phase space forest:"
    write (u, "(3x,A,I0)") "n_in  = ", forest%n_in
    write (u, "(3x,A,I0)") "n_out = ", forest%n_out
    write (u, "(3x,A,I0)") "n_tot = ", forest%n_tot
    write (u, "(3x,A,I0)") "n_masses = ", forest%n_masses
    write (u, "(3x,A,I0)") "n_angles = ", forest%n_angles
    write (u, "(3x,A,I0)") "n_dim    = ", forest%n_dimensions
    write (u, "(3x,A,I0)") "n_trees  = ", forest%n_trees
    write (u, "(3x,A,I0)") "n_equiv  = ", forest%n_equivalences
    write (u, "(3x,A)", advance="no") "flavors  ="
    if (allocated (forest%flv)) then
       do i = 1, size (forest%flv)
          write (u, "(1x,I0)", advance="no")  forest%flv(i)%get_pdg ()
       end do
       write (u, "(A)")
    else
       write (u, "(1x,A)") "[empty]"
    end if
    write (u, "(1x,A)") "Particle combinations:"
    if (allocated (forest%prt_combination)) then
       do k = 1, size (forest%prt_combination, 2)
          if (forest%prt_combination(1, k) /= 0) then
             write (u, "(3x,I0,1x,'<=',1x,I0,1x,'+',1x,I0)") &
                  k, forest%prt_combination(:,k)
          end if
       end do
    else
       write (u, "(3x,A)") "  [empty]"
    end if
    write (u, "(1x,A)") "Groves and trees:"
    if (allocated (forest%grove)) then
       do g = 1, size (forest%grove)
          write (u, "(3x,A,1x,I0)") "Grove    ", g
          call phs_grove_write (forest%grove(g), unit)
       end do
    else
       write (u, "(3x,A)") "  [empty]"
    end if
    write (u, "(1x,A,I0)") "Total number of equivalences: ", &
         forest%n_equivalences
    write (u, "(A)")
    write (u, "(1x,A)") "Global s-channel mappings:"
    if (allocated (forest%s_mapping)) then
       do i = 1, size (forest%s_mapping)
          associate (mapping => forest%s_mapping(i))
            if (mapping_is_s_channel (mapping) &
                 .or. mapping_is_on_shell (mapping)) then
               write (u, "(1x,I0,':',1x)", advance="no")  i
               call mapping_write (forest%s_mapping(i), unit)
            end if
          end associate
       end do
    else
       write (u, "(3x,A)") "  [empty]"
    end if
    write (u, "(A)")
    write (u, "(1x,A)") "Incoming particles:"
    if (allocated (forest%prt_in)) then
       if (any (phs_prt_is_defined (forest%prt_in))) then
          do i = 1, size (forest%prt_in)
             if (phs_prt_is_defined (forest%prt_in(i))) then
                write (u, "(1x,A,1x,I0)")  "Particle", i
                call phs_prt_write (forest%prt_in(i), u)
             end if
          end do
       else
          write (u, "(3x,A)")  "[all undefined]"
       end if
    else
       write (u, "(3x,A)")  "  [empty]"
    end if
    write (u, "(A)")
    write (u, "(1x,A)") "Outgoing particles:"
    if (allocated (forest%prt_out)) then
       if (any (phs_prt_is_defined (forest%prt_out))) then
          do i = 1, size (forest%prt_out)
             if (phs_prt_is_defined (forest%prt_out(i))) then
                write (u, "(1x,A,1x,I0)")  "Particle", i
                call phs_prt_write (forest%prt_out(i), u)
             end if
          end do
       else
          write (u, "(3x,A)")  "[all undefined]"
       end if
    else
       write (u, "(1x,A)")  "  [empty]"
    end if
    write (u, "(A)")
    write (u, "(1x,A)") "Tree particles:"
    if (allocated (forest%prt)) then
       if (any (phs_prt_is_defined (forest%prt))) then
          do i = 1, size (forest%prt)
             if (phs_prt_is_defined (forest%prt(i))) then
                write (u, "(1x,A,1x,I0)")  "Particle", i
                call phs_prt_write (forest%prt(i), u)
             end if
          end do
       else
          write (u, "(3x,A)")  "[all undefined]"
       end if
    else
       write (u, "(3x,A)")  "  [empty]"
    end if
  end subroutine phs_forest_write

  subroutine phs_grove_write (grove, unit)
    type(phs_grove_t), intent(in) :: grove
    integer, intent(in), optional :: unit
    integer :: u
    integer :: t
    u = given_output_unit (unit);  if (u < 0)  return
    do t = 1, size (grove%tree)
       write (u, "(3x,A,I0)") "Tree      ", t
       call phs_tree_write (grove%tree(t), unit)
    end do
    write (u, "(1x,A)") "Equivalence list:"
    call equivalence_list_write (grove%equivalence_list, unit)
  end subroutine phs_grove_write

\eatline
\nwendcode{}\nwbegindocs{535}Deep copy.
\nwenddocs{}\nwbegincode{536}\moddef{PHS forests: public}\plusendmoddef
  public :: assignment(=)
\nwendcode{}\nwbegincode{537}\moddef{PHS forests: interfaces}\plusendmoddef
  interface assignment(=)
     module procedure phs_forest_assign
  end interface

\nwendcode{}\nwbegincode{538}\moddef{PHS forests: procedures}\plusendmoddef
  subroutine phs_forest_assign (forest_out, forest_in)
    type(phs_forest_t), intent(out) :: forest_out
    type(phs_forest_t), intent(in) :: forest_in
    forest_out%n_in  = forest_in%n_in
    forest_out%n_out = forest_in%n_out
    forest_out%n_tot = forest_in%n_tot
    forest_out%n_masses = forest_in%n_masses
    forest_out%n_angles = forest_in%n_angles
    forest_out%n_dimensions  = forest_in%n_dimensions
    forest_out%n_trees  = forest_in%n_trees
    forest_out%n_equivalences  = forest_in%n_equivalences
    if (allocated (forest_in%flv)) then
       allocate (forest_out%flv (size (forest_in%flv)))
       forest_out%flv = forest_in%flv
    end if
    if (allocated (forest_in%grove)) then
       allocate (forest_out%grove (size (forest_in%grove)))
       forest_out%grove = forest_in%grove
    end if
    if (allocated (forest_in%grove_lookup)) then
       allocate (forest_out%grove_lookup (size (forest_in%grove_lookup)))
       forest_out%grove_lookup = forest_in%grove_lookup
    end if
    if (allocated (forest_in%prt_in)) then
       allocate (forest_out%prt_in (size (forest_in%prt_in)))
       forest_out%prt_in = forest_in%prt_in
    end if
    if (allocated (forest_in%prt_out)) then
       allocate (forest_out%prt_out (size (forest_in%prt_out)))
       forest_out%prt_out = forest_in%prt_out
    end if
    if (allocated (forest_in%prt)) then
       allocate (forest_out%prt (size (forest_in%prt)))
       forest_out%prt = forest_in%prt
    end if
    if (allocated (forest_in%s_mapping)) then
       allocate (forest_out%s_mapping (size (forest_in%s_mapping)))
       forest_out%s_mapping = forest_in%s_mapping
    end if
    if (allocated (forest_in%prt_combination)) then
       allocate (forest_out%prt_combination &
            (2, size (forest_in%prt_combination, 2)))
       forest_out%prt_combination = forest_in%prt_combination
    end if
  end subroutine phs_forest_assign

\eatline
\nwendcode{}\nwbegindocs{539}\nwdocspar
\subsection{Accessing contents}
Get the number of integration parameters
\nwenddocs{}\nwbegincode{540}\moddef{PHS forests: public}\plusendmoddef
  public :: phs_forest_get_n_parameters
\nwendcode{}\nwbegincode{541}\moddef{PHS forests: procedures}\plusendmoddef
  function phs_forest_get_n_parameters (forest) result (n)
    integer :: n
    type(phs_forest_t), intent(in) :: forest
    n = forest%n_dimensions
  end function phs_forest_get_n_parameters

\eatline
\nwendcode{}\nwbegindocs{542}Get the number of integration channels
\nwenddocs{}\nwbegincode{543}\moddef{PHS forests: public}\plusendmoddef
  public :: phs_forest_get_n_channels
\nwendcode{}\nwbegincode{544}\moddef{PHS forests: procedures}\plusendmoddef
  function phs_forest_get_n_channels (forest) result (n)
    integer :: n
    type(phs_forest_t), intent(in) :: forest
    n = forest%n_trees
  end function phs_forest_get_n_channels

\eatline
\nwendcode{}\nwbegindocs{545}Get the number of groves
\nwenddocs{}\nwbegincode{546}\moddef{PHS forests: public}\plusendmoddef
  public :: phs_forest_get_n_groves
\nwendcode{}\nwbegincode{547}\moddef{PHS forests: procedures}\plusendmoddef
  function phs_forest_get_n_groves (forest) result (n)
    integer :: n
    type(phs_forest_t), intent(in) :: forest
    n = size (forest%grove)
  end function phs_forest_get_n_groves

\eatline
\nwendcode{}\nwbegindocs{548}Get the index bounds for a specific grove.
\nwenddocs{}\nwbegincode{549}\moddef{PHS forests: public}\plusendmoddef
  public :: phs_forest_get_grove_bounds
\nwendcode{}\nwbegincode{550}\moddef{PHS forests: procedures}\plusendmoddef
  subroutine phs_forest_get_grove_bounds (forest, g, i0, i1, n)
    type(phs_forest_t), intent(in) :: forest
    integer, intent(in) :: g
    integer, intent(out) :: i0, i1, n
    n = size (forest%grove(g)%tree)
    i0 = forest%grove(g)%tree_count_offset + 1
    i1 = forest%grove(g)%tree_count_offset + n
  end subroutine phs_forest_get_grove_bounds

\eatline
\nwendcode{}\nwbegindocs{551}Get the number of equivalences
\nwenddocs{}\nwbegincode{552}\moddef{PHS forests: public}\plusendmoddef
  public :: phs_forest_get_n_equivalences
\nwendcode{}\nwbegincode{553}\moddef{PHS forests: procedures}\plusendmoddef
  function phs_forest_get_n_equivalences (forest) result (n)
    integer :: n
    type(phs_forest_t), intent(in) :: forest
    n = forest%n_equivalences
  end function phs_forest_get_n_equivalences

\eatline
\nwendcode{}\nwbegindocs{554}Return true if a particular channel has a global (s-channel)
mapping; also return the resonance mass and width for this mapping.
\nwenddocs{}\nwbegincode{555}\moddef{PHS forests: public}\plusendmoddef
  public :: phs_forest_get_s_mapping
  public :: phs_forest_get_on_shell
\nwendcode{}\nwbegincode{556}\moddef{PHS forests: procedures}\plusendmoddef
  subroutine phs_forest_get_s_mapping (forest, channel, flag, mass, width)
    type(phs_forest_t), intent(in) :: forest
    integer, intent(in) :: channel
    logical, intent(out) :: flag
    real(default), intent(out) :: mass, width
    flag = mapping_is_s_channel (forest%s_mapping(channel))
    if (flag) then
       mass = mapping_get_mass (forest%s_mapping(channel))
       width = mapping_get_width (forest%s_mapping(channel))
    else
       mass = 0
       width = 0
    end if
  end subroutine phs_forest_get_s_mapping

  subroutine phs_forest_get_on_shell (forest, channel, flag, mass)
    type(phs_forest_t), intent(in) :: forest
    integer, intent(in) :: channel
    logical, intent(out) :: flag
    real(default), intent(out) :: mass
    flag = mapping_is_on_shell (forest%s_mapping(channel))
    if (flag) then
       mass = mapping_get_mass (forest%s_mapping(channel))
    else
       mass = 0
    end if
  end subroutine phs_forest_get_on_shell

\eatline
\eatline
\nwendcode{}\nwbegindocs{557}\nwdocspar
\subsection{Read the phase space setup from file}
The phase space setup is stored in a file.  The file may be generated
by the {\tt{}cascades} module below, or by other means.  This file has to
be read and parsed to create the PHS forest as the internal
phase-space representation.

Create lexer and syntax:
\nwenddocs{}\nwbegincode{558}\moddef{PHS forests: procedures}\plusendmoddef
  subroutine define_phs_forest_syntax (ifile)
    type(ifile_t) :: ifile
    call ifile_append (ifile, "SEQ phase_space_list = process_phase_space*")
    call ifile_append (ifile, "SEQ process_phase_space = " &
         // "process_def process_header phase_space")
    call ifile_append (ifile, "SEQ process_def = process process_list")
    call ifile_append (ifile, "KEY process")
    call ifile_append (ifile, "LIS process_list = process_tag*")
    call ifile_append (ifile, "IDE process_tag")
    call ifile_append (ifile, "SEQ process_header = " &
         // "md5sum_process = md5sum " &
         // "md5sum_model_par = md5sum " &
         // "md5sum_phs_config = md5sum " &
         // "sqrts = real " &
         // "m_threshold_s = real " &
         // "m_threshold_t = real " &
         // "off_shell = integer " &
         // "t_channel = integer " &
         // "keep_nonresonant = logical")
    call ifile_append (ifile, "KEY '='")
    call ifile_append (ifile, "KEY md5sum_process")
    call ifile_append (ifile, "KEY md5sum_model_par")
    call ifile_append (ifile, "KEY md5sum_phs_config")
    call ifile_append (ifile, "KEY sqrts")
    call ifile_append (ifile, "KEY m_threshold_s")
    call ifile_append (ifile, "KEY m_threshold_t")
    call ifile_append (ifile, "KEY off_shell")
    call ifile_append (ifile, "KEY t_channel")
    call ifile_append (ifile, "KEY keep_nonresonant")
    call ifile_append (ifile, "QUO md5sum = '""' ... '""'")
    call ifile_append (ifile, "REA real")
    call ifile_append (ifile, "INT integer")
    call ifile_append (ifile, "IDE logical")
    call ifile_append (ifile, "SEQ phase_space = grove_def+")
    call ifile_append (ifile, "SEQ grove_def = grove tree_def+")
    call ifile_append (ifile, "KEY grove")
    call ifile_append (ifile, "SEQ tree_def = tree bincodes mapping*")
    call ifile_append (ifile, "KEY tree")
    call ifile_append (ifile, "SEQ bincodes = bincode*")
    call ifile_append (ifile, "INT bincode")
    call ifile_append (ifile, "SEQ mapping = map bincode channel pdg")
    call ifile_append (ifile, "KEY map")
    call ifile_append (ifile, "ALT channel = &
         &s_channel | t_channel | u_channel | &
         &collinear | infrared | radiation | on_shell")
    call ifile_append (ifile, "KEY s_channel")
    ! call ifile_append (ifile, "KEY t_channel")   !!! Key already exists
    call ifile_append (ifile, "KEY u_channel")
    call ifile_append (ifile, "KEY collinear")
    call ifile_append (ifile, "KEY infrared")
    call ifile_append (ifile, "KEY radiation")
    call ifile_append (ifile, "KEY on_shell")
    call ifile_append (ifile, "INT pdg")
  end subroutine define_phs_forest_syntax

\eatline
\nwendcode{}\nwbegindocs{559}The model-file syntax and lexer are fixed, therefore stored as
module variables:
\nwenddocs{}\nwbegincode{560}\moddef{PHS forests: variables}\endmoddef
  type(syntax_t), target, save :: syntax_phs_forest

\eatline
\nwendcode{}\nwbegincode{561}\moddef{PHS forests: public}\plusendmoddef
  public :: syntax_phs_forest_init
\nwendcode{}\nwbegincode{562}\moddef{PHS forests: procedures}\plusendmoddef
  subroutine syntax_phs_forest_init ()
    type(ifile_t) :: ifile
    call define_phs_forest_syntax (ifile)
    call syntax_init (syntax_phs_forest, ifile)
    call ifile_final (ifile)
  end subroutine syntax_phs_forest_init

\eatline
\nwendcode{}\nwbegincode{563}\moddef{PHS forests: procedures}\plusendmoddef
  subroutine lexer_init_phs_forest (lexer)
    type(lexer_t), intent(out) :: lexer
    call lexer_init (lexer, &
         comment_chars = "#!", &
         quote_chars = '"', &
         quote_match = '"', &
         single_chars = "", &
         special_class = ["="] , &
         keyword_list = syntax_get_keyword_list_ptr (syntax_phs_forest))
  end subroutine lexer_init_phs_forest

\eatline
\nwendcode{}\nwbegincode{564}\moddef{PHS forests: public}\plusendmoddef
  public :: syntax_phs_forest_final
\nwendcode{}\nwbegincode{565}\moddef{PHS forests: procedures}\plusendmoddef
  subroutine syntax_phs_forest_final ()
    call syntax_final (syntax_phs_forest)
  end subroutine syntax_phs_forest_final

\eatline
\nwendcode{}\nwbegincode{566}\moddef{PHS forests: public}\plusendmoddef
  public :: syntax_phs_forest_write
\nwendcode{}\nwbegincode{567}\moddef{PHS forests: procedures}\plusendmoddef
  subroutine syntax_phs_forest_write (unit)
    integer, intent(in), optional :: unit
    call syntax_write (syntax_phs_forest, unit)
  end subroutine syntax_phs_forest_write

\eatline
\nwendcode{}\nwbegindocs{568}The concrete parser and interpreter.  Generate an input stream for
the external {\tt{}unit}, read the parse tree (with given {\tt{}syntax} and
{\tt{}lexer}) from this stream, and transfer the contents of the parse
tree to the PHS {\tt{}forest}.

We look for the matching {\tt{}process} tag, count groves and trees for
initializing the {\tt{}forest}, and fill the trees.

If the optional parameters are set, compare the parameters stored in
the file to those.  Set {\tt{}match} true if everything agrees.
\nwenddocs{}\nwbegincode{569}\moddef{PHS forests: public}\plusendmoddef
  public :: phs_forest_read
\nwendcode{}\nwbegincode{570}\moddef{PHS forests: interfaces}\plusendmoddef
  interface phs_forest_read
     module procedure phs_forest_read_file
     module procedure phs_forest_read_unit
     module procedure phs_forest_read_parse_tree
  end interface

\nwendcode{}\nwbegincode{571}\moddef{PHS forests: procedures}\plusendmoddef
  subroutine phs_forest_read_file &
       (forest, filename, process_id, n_in, n_out, model, found, &
        md5sum_process, md5sum_model_par, &
        md5sum_phs_config, phs_par, match)
    type(phs_forest_t), intent(out) :: forest
    type(string_t), intent(in) :: filename
    type(string_t), intent(in) :: process_id
    integer, intent(in) :: n_in, n_out
    class(model_data_t), intent(in), target :: model
    logical, intent(out) :: found
    character(32), intent(in), optional :: &
         md5sum_process, md5sum_model_par, md5sum_phs_config
    type(phs_parameters_t), intent(in), optional :: phs_par
    logical, intent(out), optional :: match
    type(parse_tree_t), target :: parse_tree
    type(stream_t), target :: stream
    type(lexer_t) :: lexer
    call lexer_init_phs_forest (lexer)
    call stream_init (stream, char (filename))
    call lexer_assign_stream (lexer, stream)
    call parse_tree_init (parse_tree, syntax_phs_forest, lexer)
    call phs_forest_read (forest, parse_tree, &
         process_id, n_in, n_out, model, found, &
         md5sum_process, md5sum_model_par, md5sum_phs_config, phs_par, match)
    call stream_final (stream)
    call lexer_final (lexer)
    call parse_tree_final (parse_tree)
  end subroutine phs_forest_read_file

  subroutine phs_forest_read_unit &
       (forest, unit, process_id, n_in, n_out, model, found, &
        md5sum_process, md5sum_model_par, md5sum_phs_config, &
        phs_par, match)
    type(phs_forest_t), intent(out) :: forest
    integer, intent(in) :: unit
    type(string_t), intent(in) :: process_id
    integer, intent(in) :: n_in, n_out
    class(model_data_t), intent(in), target :: model
    logical, intent(out) :: found
    character(32), intent(in), optional :: &
         md5sum_process, md5sum_model_par, md5sum_phs_config
    type(phs_parameters_t), intent(in), optional :: phs_par
    logical, intent(out), optional :: match
    type(parse_tree_t), target :: parse_tree
    type(stream_t), target :: stream
    type(lexer_t) :: lexer
    call lexer_init_phs_forest (lexer)
    call stream_init (stream, unit)
    call lexer_assign_stream (lexer, stream)
    call parse_tree_init (parse_tree, syntax_phs_forest, lexer)
    call phs_forest_read (forest, parse_tree, &
         process_id, n_in, n_out, model, found, &
         md5sum_process, md5sum_model_par, md5sum_phs_config, &
         phs_par, match)
    call stream_final (stream)
    call lexer_final (lexer)
    call parse_tree_final (parse_tree)
  end subroutine phs_forest_read_unit

  subroutine phs_forest_read_parse_tree &
       (forest, parse_tree, process_id, n_in, n_out, model, found, &
        md5sum_process, md5sum_model_par, md5sum_phs_config, &
        phs_par, match)
    type(phs_forest_t), intent(out) :: forest
    type(parse_tree_t), intent(in), target :: parse_tree
    type(string_t), intent(in) :: process_id
    integer, intent(in) :: n_in, n_out
    class(model_data_t), intent(in), target :: model
    logical, intent(out) :: found
    character(32), intent(in), optional :: &
         md5sum_process, md5sum_model_par, md5sum_phs_config
    type(phs_parameters_t), intent(in), optional :: phs_par
    logical, intent(out), optional :: match
    type(parse_node_t), pointer :: node_header, node_phs, node_grove
    integer :: n_grove, g
    integer, dimension(:), allocatable :: n_tree
    integer :: t
    logical :: real_phsp_work
    node_header => parse_tree_get_process_ptr (parse_tree, process_id)
    found = associated (node_header);  if (.not. found)  return
    if (present (match)) then
       call phs_forest_check_input (node_header, &
            md5sum_process, md5sum_model_par, md5sum_phs_config, phs_par, match)
       if (.not. match)  return
    end if
    node_phs => parse_node_get_next_ptr (node_header)
    n_grove = parse_node_get_n_sub (node_phs)
    allocate (n_tree (n_grove))
    do g = 1, n_grove
       node_grove => parse_node_get_sub_ptr (node_phs, g)
       n_tree(g) = parse_node_get_n_sub (node_grove) - 1
    end do
    call phs_forest_init (forest, n_tree, n_in, n_out)
    do g = 1, n_grove
       node_grove => parse_node_get_sub_ptr (node_phs, g)
       do t = 1, n_tree(g)
          call phs_tree_set (forest%grove(g)%tree(t), &
               parse_node_get_sub_ptr (node_grove, t+1), model)
       end do
    end do
  end subroutine phs_forest_read_parse_tree

\eatline
\nwendcode{}\nwbegindocs{572}Check the input for consistency.  If any MD5 sum or phase-space
parameter disagrees, the phase-space file cannot be used.  The MD5
sum checks are skipped if the stored MD5 sum is empty.
\nwenddocs{}\nwbegincode{573}\moddef{PHS forests: procedures}\plusendmoddef
  subroutine phs_forest_check_input (pn_header, &
       md5sum_process, md5sum_model_par, md5sum_phs_config, phs_par, match)
    type(parse_node_t), intent(in), target :: pn_header
    character(32), intent(in) :: &
         md5sum_process, md5sum_model_par, md5sum_phs_config
    type(phs_parameters_t), intent(in), optional :: phs_par
    logical, intent(out) :: match
    type(parse_node_t), pointer :: pn_md5sum, pn_rval, pn_ival, pn_lval
    character(32) :: md5sum
    type(phs_parameters_t) :: phs_par_old
    character(1) :: lstr
    pn_md5sum => parse_node_get_sub_ptr (pn_header, 3)
    md5sum = parse_node_get_string (pn_md5sum)
    if (md5sum /= "" .and. md5sum /= md5sum_process) then
       call msg_message ("Phase space: discarding old configuration &
            &(process changed)")
       match = .false.;  return
    end if
    pn_md5sum => parse_node_get_next_ptr (pn_md5sum, 3)
    md5sum = parse_node_get_string (pn_md5sum)
    if (md5sum /= "" .and. md5sum /= md5sum_model_par) then
       call msg_message ("Phase space: discarding old configuration &
            &(model parameters changed)")
       match = .false.;  return
    end if
    pn_md5sum => parse_node_get_next_ptr (pn_md5sum, 3)
    md5sum = parse_node_get_string (pn_md5sum)
    if (md5sum /= "" .and. md5sum /= md5sum_phs_config) then
       call msg_message ("Phase space: discarding old configuration &
            &(configuration parameters changed)")
       match = .false.;  return
    end if
    if (present (phs_par)) then
       pn_rval => parse_node_get_next_ptr (pn_md5sum, 3)
       phs_par_old%sqrts = parse_node_get_real (pn_rval)
       pn_rval => parse_node_get_next_ptr (pn_rval, 3)
       phs_par_old%m_threshold_s = parse_node_get_real (pn_rval)
       pn_rval => parse_node_get_next_ptr (pn_rval, 3)
       phs_par_old%m_threshold_t = parse_node_get_real (pn_rval)
       pn_ival => parse_node_get_next_ptr (pn_rval, 3)
       phs_par_old%off_shell = parse_node_get_integer (pn_ival)
       pn_ival => parse_node_get_next_ptr (pn_ival, 3)
       phs_par_old%t_channel = parse_node_get_integer (pn_ival)
       pn_lval => parse_node_get_next_ptr (pn_ival, 3)
       lstr = parse_node_get_string (pn_lval)
       read (lstr, "(L1)")  phs_par_old%keep_nonresonant
       if (phs_par_old /= phs_par) then
          call msg_message &
               ("Phase space: discarding old configuration &
               &(configuration parameters changed)")
          match = .false.;  return
       end if
    end if
    match = .true.
  end subroutine phs_forest_check_input

\eatline
\nwendcode{}\nwbegindocs{574}Initialize a specific tree in the forest, using the contents of the
'tree' node.  First, count the bincodes, allocate an array and read
them in, and make the tree.  Each $t$-channel tree is flipped to
$s$-channel.  Then, find mappings and initialize them.
\nwenddocs{}\nwbegincode{575}\moddef{PHS forests: procedures}\plusendmoddef
  subroutine phs_tree_set (tree, node, model)
    type(phs_tree_t), intent(inout) :: tree
    type(parse_node_t), intent(in), target :: node
    class(model_data_t), intent(in), target :: model
    type(parse_node_t), pointer :: node_bincodes, node_mapping
    integer :: n_bincodes, offset
    integer(TC), dimension(:), allocatable :: bincode
    integer :: b, n_mappings, m
    integer(TC) :: k
    type(string_t) :: type
    integer :: pdg
    node_bincodes => parse_node_get_sub_ptr (node, 2)
    if (associated (node_bincodes)) then
       select case (char (parse_node_get_rule_key (node_bincodes)))
       case ("bincodes")
          n_bincodes = parse_node_get_n_sub (node_bincodes)
          offset = 2
       case default
          n_bincodes = 0
          offset = 1
       end select
    else
       n_bincodes = 0
       offset = 2
    end if
    allocate (bincode (n_bincodes))
    do b = 1, n_bincodes
       bincode(b) = parse_node_get_integer &
            (parse_node_get_sub_ptr (node_bincodes, b))
    end do
    call phs_tree_from_array (tree, bincode)
    call phs_tree_flip_t_to_s_channel (tree)
    call phs_tree_canonicalize (tree)
    n_mappings = parse_node_get_n_sub (node) - offset
    do m = 1, n_mappings
       node_mapping => parse_node_get_sub_ptr (node, m + offset)
       k = parse_node_get_integer &
            (parse_node_get_sub_ptr (node_mapping, 2))
       type = parse_node_get_key &
            (parse_node_get_sub_ptr (node_mapping, 3))
       pdg = parse_node_get_integer &
            (parse_node_get_sub_ptr (node_mapping, 4))
       call phs_tree_init_mapping (tree, k, type, pdg, model)
    end do
  end subroutine phs_tree_set

\eatline
\nwendcode{}\nwbegindocs{576}\nwdocspar
\subsection{Preparation}
The trees that we read from file do not carry flavor information.
This is set separately:

The flavor list must be unique for a unique set of masses; if a given
particle can have different flavor, the mass must be degenerate, so we
can choose one of the possible flavor combinations.
\nwenddocs{}\nwbegincode{577}\moddef{PHS forests: public}\plusendmoddef
  public :: phs_forest_set_flavors
\nwendcode{}\nwbegincode{578}\moddef{PHS forests: procedures}\plusendmoddef
  subroutine phs_forest_set_flavors (forest, flv, reshuffle, flv_extra)
    type(phs_forest_t), intent(inout) :: forest
    type(flavor_t), dimension(:), intent(in) :: flv
    integer, intent(in), dimension(:), allocatable, optional :: reshuffle
    type(flavor_t), intent(in), optional :: flv_extra
    integer :: i, n_flv0
    if (present (reshuffle) .and. present (flv_extra)) then
       n_flv0 = size (flv)
       do i = 1, n_flv0
          if (reshuffle(i) <= n_flv0) then
             forest%flv(i) = flv (reshuffle(i))
          else
             forest%flv(i) = flv_extra
          end if
       end do
    else
       allocate (forest%flv (size (flv)))
       forest%flv = flv
    end if
  end subroutine phs_forest_set_flavors

\eatline
\nwendcode{}\nwbegindocs{579}\nwdocspar
\nwenddocs{}\nwbegincode{580}\moddef{PHS forests: public}\plusendmoddef
  public :: phs_forest_set_momentum_links
\nwendcode{}\nwbegincode{581}\moddef{PHS forests: procedures}\plusendmoddef
  subroutine phs_forest_set_momentum_links (forest, list)
    type(phs_forest_t), intent(inout) :: forest
    integer, intent(in), dimension(:), allocatable :: list
    integer :: g, t
    do g = 1, size (forest%grove)
      do t = 1, size (forest%grove(g)%tree)
        associate (tree => forest%grove(g)%tree(t))
          call phs_tree_set_momentum_links (tree, list)
!!!          call phs_tree_reshuffle_mappings (tree)
        end associate
      end do
    end do
  end subroutine phs_forest_set_momentum_links

\eatline
\nwendcode{}\nwbegindocs{582}Once the parameter set is fixed, the masses and the widths of the
particles are known and the {\tt{}mass{\char95}sum} arrays as well as the mapping
parameters can be computed.  Note that order is important: we first
compute the mass sums, then the ordinary mappings.  The resonances
obtained here determine the effective masses, which in turn are used
to implement step mappings for resonance decay products that are not
mapped otherwise.
\nwenddocs{}\nwbegincode{583}\moddef{PHS forests: public}\plusendmoddef
  public :: phs_forest_set_parameters
\nwendcode{}\nwbegincode{584}\moddef{PHS forests: procedures}\plusendmoddef
  subroutine phs_forest_set_parameters &
       (forest, mapping_defaults, variable_limits)
    type(phs_forest_t), intent(inout) :: forest
    type(mapping_defaults_t), intent(in) :: mapping_defaults
    logical, intent(in) :: variable_limits
    integer :: g, t
    do g = 1, size (forest%grove)
       do t = 1, size (forest%grove(g)%tree)
          call phs_tree_set_mass_sum &
               (forest%grove(g)%tree(t), forest%flv(forest%n_in+1:))
          call phs_tree_set_mapping_parameters (forest%grove(g)%tree(t), &
               mapping_defaults, variable_limits)
          call phs_tree_set_effective_masses (forest%grove(g)%tree(t))
          if (mapping_defaults%step_mapping) then
             call phs_tree_set_step_mappings (forest%grove(g)%tree(t), &
                  mapping_defaults%step_mapping_exp, variable_limits)
          end if
       end do
    end do
  end subroutine phs_forest_set_parameters

\eatline
\nwendcode{}\nwbegindocs{585}Generate the particle combination table.  Scan all trees and merge
their individual combination tables.  At the end, valid entries are
non-zero, and they indicate the indices of a pair of particles to be
combined to a new particle.  If a particle is accessible by more than
one tree (this is usual), only keep the first possibility.
\nwenddocs{}\nwbegincode{586}\moddef{PHS forests: public}\plusendmoddef
  public :: phs_forest_setup_prt_combinations
\nwendcode{}\nwbegincode{587}\moddef{PHS forests: procedures}\plusendmoddef
  subroutine phs_forest_setup_prt_combinations (forest)
    type(phs_forest_t), intent(inout) :: forest
    integer :: g, t
    integer, dimension(:,:), allocatable :: tree_prt_combination
    forest%prt_combination = 0
    allocate (tree_prt_combination (2, size (forest%prt_combination, 2)))
    do g = 1, size (forest%grove)
       do t = 1, size (forest%grove(g)%tree)
          call phs_tree_setup_prt_combinations &
               (forest%grove(g)%tree(t), tree_prt_combination)
          where (tree_prt_combination /= 0 .and. forest%prt_combination == 0)
             forest%prt_combination = tree_prt_combination
          end where
       end do
    end do
  end subroutine phs_forest_setup_prt_combinations

\eatline
\nwendcode{}\nwbegindocs{588}\nwdocspar
\subsection{Accessing the particle arrays}
Set the incoming particles from the contents of an interaction.
\nwenddocs{}\nwbegincode{589}\moddef{PHS forests: public}\plusendmoddef
  public :: phs_forest_set_prt_in
\nwendcode{}\nwbegincode{590}\moddef{PHS forests: interfaces}\plusendmoddef
  interface phs_forest_set_prt_in
     module procedure phs_forest_set_prt_in_int, phs_forest_set_prt_in_mom
  end interface phs_forest_set_prt_in
\nwendcode{}\nwbegincode{591}\moddef{PHS forests: procedures}\plusendmoddef
  subroutine phs_forest_set_prt_in_int (forest, int, lt_cm_to_lab)
    type(phs_forest_t), intent(inout) :: forest
    type(interaction_t), intent(in) :: int
    type(lorentz_transformation_t), intent(in), optional :: lt_cm_to_lab
    if (present (lt_cm_to_lab)) then
       call phs_prt_set_momentum (forest%prt_in, &
            inverse (lt_cm_to_lab) * &
            interaction_get_momenta (int, outgoing=.false.))
    else
       call phs_prt_set_momentum (forest%prt_in, &
            interaction_get_momenta (int, outgoing=.false.))
    end if
    associate (m_in => forest%flv(:forest%n_in)%get_mass ())
      call phs_prt_set_msq (forest%prt_in, m_in ** 2)
    end associate
    call phs_prt_set_defined (forest%prt_in)
  end subroutine phs_forest_set_prt_in_int

  subroutine phs_forest_set_prt_in_mom (forest, mom, lt_cm_to_lab)
    type(phs_forest_t), intent(inout) :: forest
    type(vector4_t), dimension(size (forest%prt_in)), intent(in) :: mom
    type(lorentz_transformation_t), intent(in), optional :: lt_cm_to_lab
    if (present (lt_cm_to_lab)) then
       call phs_prt_set_momentum (forest%prt_in, &
            inverse (lt_cm_to_lab) * mom)
    else
       call phs_prt_set_momentum (forest%prt_in, mom)
    end if
    associate (m_in => forest%flv(:forest%n_in)%get_mass ())
      call phs_prt_set_msq (forest%prt_in, m_in ** 2)
    end associate
    call phs_prt_set_defined (forest%prt_in)
  end subroutine phs_forest_set_prt_in_mom

\eatline
\nwendcode{}\nwbegindocs{592}Set the outgoing particles from the contents of an interaction.
\nwenddocs{}\nwbegincode{593}\moddef{PHS forests: public}\plusendmoddef
  public :: phs_forest_set_prt_out
\nwendcode{}\nwbegincode{594}\moddef{PHS forests: interfaces}\plusendmoddef
  interface phs_forest_set_prt_out
     module procedure phs_forest_set_prt_out_int, phs_forest_set_prt_out_mom
  end interface phs_forest_set_prt_out
\nwendcode{}\nwbegincode{595}\moddef{PHS forests: procedures}\plusendmoddef
  subroutine phs_forest_set_prt_out_int (forest, int, lt_cm_to_lab)
    type(phs_forest_t), intent(inout) :: forest
    type(interaction_t), intent(in) :: int
    type(lorentz_transformation_t), intent(in), optional :: lt_cm_to_lab
    if (present (lt_cm_to_lab)) then
       call phs_prt_set_momentum (forest%prt_out, &
            inverse (lt_cm_to_lab) * &
            interaction_get_momenta (int, outgoing=.true.))
    else
       call phs_prt_set_momentum (forest%prt_out, &
            interaction_get_momenta (int, outgoing=.true.))
    end if
    associate (m_out => forest%flv(forest%n_in+1:)%get_mass ())
      call phs_prt_set_msq (forest%prt_out, m_out ** 2)
    end associate
    call phs_prt_set_defined (forest%prt_out)
  end subroutine phs_forest_set_prt_out_int

  subroutine phs_forest_set_prt_out_mom (forest, mom, lt_cm_to_lab)
    type(phs_forest_t), intent(inout) :: forest
    type(vector4_t), dimension(size (forest%prt_out)), intent(in) :: mom
    type(lorentz_transformation_t), intent(in), optional :: lt_cm_to_lab
    if (present (lt_cm_to_lab)) then
       call phs_prt_set_momentum (forest%prt_out, &
            inverse (lt_cm_to_lab) * mom)
    else
       call phs_prt_set_momentum (forest%prt_out, mom)
    end if
    associate (m_out => forest%flv(forest%n_in+1:)%get_mass ())
      call phs_prt_set_msq (forest%prt_out, m_out ** 2)
    end associate
    call phs_prt_set_defined (forest%prt_out)
  end subroutine phs_forest_set_prt_out_mom

\eatline
\nwendcode{}\nwbegindocs{596}Combine particles as described by the particle combination table.
Particle momentum sums will be calculated only if the resulting
particle is contained in at least one of the trees in the current
forest.  The others are kept undefined.
\nwenddocs{}\nwbegincode{597}\moddef{PHS forests: public}\plusendmoddef
  public :: phs_forest_combine_particles
\nwendcode{}\nwbegincode{598}\moddef{PHS forests: procedures}\plusendmoddef
  subroutine phs_forest_combine_particles (forest)
    type(phs_forest_t), intent(inout) :: forest
    integer :: k
    integer, dimension(2) :: kk
    do k = 1, size (forest%prt_combination, 2)
       kk = forest%prt_combination(:,k)
       if (kk(1) /= 0) then
          call phs_prt_combine (forest%prt(k), &
               forest%prt(kk(1)), forest%prt(kk(2)))
       end if
    end do
  end subroutine phs_forest_combine_particles

\eatline
\nwendcode{}\nwbegindocs{599}Extract the outgoing particles and insert into an interaction.
\nwenddocs{}\nwbegincode{600}\moddef{PHS forests: public}\plusendmoddef
  public :: phs_forest_get_prt_out
\nwendcode{}\nwbegincode{601}\moddef{PHS forests: procedures}\plusendmoddef
  subroutine phs_forest_get_prt_out (forest, int, lt_cm_to_lab)
    type(phs_forest_t), intent(in) :: forest
    type(interaction_t), intent(inout) :: int
    type(lorentz_transformation_t), intent(in), optional :: lt_cm_to_lab
    if (present (lt_cm_to_lab)) then
       call interaction_set_momenta (int, &
            lt_cm_to_lab * &
            phs_prt_get_momentum (forest%prt_out), outgoing=.true.)
    else
       call interaction_set_momenta (int, &
            phs_prt_get_momentum (forest%prt_out), outgoing=.true.)
    end if
  end subroutine phs_forest_get_prt_out

\eatline
\nwendcode{}\nwbegindocs{602}Extract the outgoing particle momenta
\nwenddocs{}\nwbegincode{603}\moddef{PHS forests: public}\plusendmoddef
  public :: phs_forest_get_momenta_out
\nwendcode{}\nwbegincode{604}\moddef{PHS forests: procedures}\plusendmoddef
  function phs_forest_get_momenta_out (forest, lt_cm_to_lab) result (p)
    type(phs_forest_t), intent(in) :: forest
    type(lorentz_transformation_t), intent(in), optional :: lt_cm_to_lab
    type(vector4_t), dimension(size (forest%prt_out)) :: p
    p = phs_prt_get_momentum (forest%prt_out)
    if (present (lt_cm_to_lab)) p = p * lt_cm_to_lab
  end function phs_forest_get_momenta_out
  
\eatline
\nwendcode{}\nwbegindocs{605}\nwdocspar
\subsection{Find equivalences among phase-space trees}
Scan phase space for equivalences.  We generate the complete set of
unique permutations for the given list of outgoing particles, and use
this for scanning equivalences within each grove.
\nwenddocs{}\nwbegindocs{606}We scan all pairs of trees, using all permutations.  This implies
that trivial equivalences are included, and equivalences between
different trees are recorded twice.  This is intentional.
\nwenddocs{}\nwbegincode{607}\moddef{PHS forests: procedures}\plusendmoddef
  subroutine phs_grove_set_equivalences (grove, perm_array)
    type(phs_grove_t), intent(inout) :: grove
    type(permutation_t), dimension(:), intent(in) :: perm_array
    type(equivalence_t), pointer :: eq
    integer :: t1, t2, i
    do t1 = 1, size (grove%tree)
       do t2 = 1, size (grove%tree)
          SCAN_PERM: do i = 1, size (perm_array)
             if (phs_tree_equivalent &
                  (grove%tree(t1), grove%tree(t2), perm_array(i))) then
                call equivalence_list_add &
                     (grove%equivalence_list, t1, t2, perm_array(i))
                eq => grove%equivalence_list%last
                call phs_tree_find_msq_permutation &
                     (grove%tree(t1), grove%tree(t2), eq%perm, &
                      eq%msq_perm)
                call phs_tree_find_angle_permutation &
                     (grove%tree(t1), grove%tree(t2), eq%perm, &
                      eq%angle_perm, eq%angle_sig)
             end if
          end do SCAN_PERM
       end do
    end do
  end subroutine phs_grove_set_equivalences

\eatline
\nwendcode{}\nwbegincode{608}\moddef{PHS forests: public}\plusendmoddef
  public :: phs_forest_set_equivalences
\nwendcode{}\nwbegincode{609}\moddef{PHS forests: procedures}\plusendmoddef
  subroutine phs_forest_set_equivalences (forest)
    type(phs_forest_t), intent(inout) :: forest
    type(permutation_t), dimension(:), allocatable :: perm_array
    integer :: i
    call permutation_array_make &
         (perm_array, forest%flv(forest%n_in+1:)%get_pdg ())
    do i = 1, size (forest%grove)
       call phs_grove_set_equivalences (forest%grove(i), perm_array)
    end do
    forest%n_equivalences = sum (forest%grove%equivalence_list%length)
  end subroutine phs_forest_set_equivalences

\eatline
\nwendcode{}\nwbegindocs{610}\nwdocspar
\subsection{Interface for channel equivalences}
Here, we store the equivalence list in the appropriate containers that
the {\tt{}phs{\char95}base} module provides.  There is one separate list for each
channel.
\nwenddocs{}\nwbegincode{611}\moddef{PHS forests: public}\plusendmoddef
  public :: phs_forest_get_equivalences 
\nwendcode{}\nwbegincode{612}\moddef{PHS forests: procedures}\plusendmoddef
  subroutine phs_forest_get_equivalences (forest, channel, azimuthal_dependence)
    type(phs_forest_t), intent(in) :: forest
    type(phs_channel_t), dimension(:), intent(out) :: channel
    logical, intent(in) :: azimuthal_dependence
    integer :: n_masses, n_angles
    integer :: mode_azimuthal_angle
    integer, dimension(:), allocatable :: n_eq
    type(equivalence_t), pointer :: eq
    integer, dimension(:), allocatable :: perm, mode
    integer :: g, c, j, left, right
    n_masses = forest%n_masses
    n_angles = forest%n_angles
    allocate (n_eq (forest%n_trees), source = 0)
    allocate (perm (forest%n_dimensions))
    allocate (mode (forest%n_dimensions), source = EQ_IDENTITY)
    do g = 1, size (forest%grove)
       eq => forest%grove(g)%equivalence_list%first
       do while (associated (eq))
          left = eq%left + forest%grove(g)%tree_count_offset
          n_eq(left) = n_eq(left) + 1
          eq => eq%next
       end do
    end do
    do c = 1, size (channel)
       allocate (channel(c)%eq (n_eq(c)))
       do j = 1, n_eq(c)
          call channel(c)%eq(j)%init (forest%n_dimensions)
       end do
    end do
    n_eq = 0
    if (azimuthal_dependence) then
       mode_azimuthal_angle = EQ_IDENTITY
    else
       mode_azimuthal_angle = EQ_INVARIANT
    end if
    do g = 1, size (forest%grove)
       eq => forest%grove(g)%equivalence_list%first
       do while (associated (eq))
          left = eq%left + forest%grove(g)%tree_count_offset
          right = eq%right + forest%grove(g)%tree_count_offset
          do j = 1, n_masses
             perm(j) = permute (j, eq%msq_perm)
             mode(j) = EQ_IDENTITY
          end do
          do j = 1, n_angles
             perm(n_masses+j) = n_masses + permute (j, eq%angle_perm)
             if (j == 1) then
                mode(n_masses+j) = mode_azimuthal_angle   ! first az. angle
             else if (mod(j,2) == 1) then
                mode(n_masses+j) = EQ_SYMMETRIC          ! other az. angles
             else if (eq%angle_sig(j)) then
                mode(n_masses+j) = EQ_IDENTITY           ! polar angle +
             else
                mode(n_masses+j) = EQ_INVERT             ! polar angle -
             end if
          end do
          n_eq(left) = n_eq(left) + 1
          associate (eq_cur => channel(left)%eq(n_eq(left)))
            eq_cur%c = right
            eq_cur%perm = perm
            eq_cur%mode = mode
          end associate
          eq => eq%next
       end do
    end do
  end subroutine phs_forest_get_equivalences

\eatline
\nwendcode{}\nwbegindocs{613}\nwdocspar
\subsection{Phase-space evaluation}
Given one row of the {\tt{}x} parameter array and the corresponding
channel index, compute first all relevant momenta and then recover the
remainder of the {\tt{}x} array, the Jacobians {\tt{}phs{\char95}factor}, and the
phase-space {\tt{}volume}.

The output argument {\tt{}ok} indicates whether this was successful.
\nwenddocs{}\nwbegincode{614}\moddef{PHS forests: public}\plusendmoddef
  public :: phs_forest_evaluate_selected_channel
\nwendcode{}\nwbegincode{615}\moddef{PHS forests: procedures}\plusendmoddef
  subroutine phs_forest_evaluate_selected_channel &
       (forest, channel, active, sqrts, x, phs_factor, volume, ok)
    type(phs_forest_t), intent(inout) :: forest
    integer, intent(in) :: channel
    logical, dimension(:), intent(in) :: active
    real(default), intent(in) :: sqrts
    real(default), dimension(:,:), intent(inout) :: x
    real(default), dimension(:), intent(out) :: phs_factor
    real(default), intent(out) :: volume
    logical, intent(out) :: ok
    integer :: g, t
    integer(TC) :: k, k_root, k_in

    g = forest%grove_lookup (channel)
    t = channel - forest%grove(g)%tree_count_offset
    call phs_prt_set_undefined (forest%prt)
    call phs_prt_set_undefined (forest%prt_out)
    k_in = forest%n_tot
    
    do k = 1,forest%n_in
       forest%prt(ibset(0,k_in-k)) = forest%prt_in(k)
    end do

    do k = 1, forest%n_out
       call phs_prt_set_msq (forest%prt(ibset(0,k-1)), &
            forest%flv(forest%n_in+k)%get_mass () ** 2)
    end do


    k_root = 2**forest%n_out - 1
    select case (forest%n_in)
    case (1)
       forest%prt(k_root) = forest%prt_in(1)
    case (2)
       call phs_prt_combine &
            (forest%prt(k_root), forest%prt_in(1), forest%prt_in(2))
    end select
    call phs_tree_compute_momenta_from_x (forest%grove(g)%tree(t), &
         forest%prt, phs_factor(channel), volume, sqrts, x(:,channel), ok)
    if (ok) then
       do k = 1, forest%n_out
          forest%prt_out(k) = forest%prt(ibset(0,k-1))
       end do
    end if
  end subroutine phs_forest_evaluate_selected_channel

\eatline
\nwendcode{}\nwbegindocs{616}The remainder: recover $x$ values for all channels except for the current 
channel.

NOTE: OpenMP not used for the first loop.  {\tt{}combine{\char95}particles} is not a
channel-local operation.
\nwenddocs{}\nwbegincode{617}\moddef{PHS forests: public}\plusendmoddef
  public :: phs_forest_evaluate_other_channels
\nwendcode{}\nwbegincode{618}\moddef{PHS forests: procedures}\plusendmoddef
  subroutine phs_forest_evaluate_other_channels &
       (forest, channel, active, sqrts, x, phs_factor, combine)
    type(phs_forest_t), intent(inout) :: forest
    integer, intent(in) :: channel
    logical, dimension(:), intent(in) :: active
    real(default), intent(in) :: sqrts
    real(default), dimension(:,:), intent(inout) :: x
    real(default), dimension(:), intent(inout) :: phs_factor
    logical, intent(in) :: combine
    integer :: g, t, ch, n_channel

    g = forest%grove_lookup (channel)
    t = channel - forest%grove(g)%tree_count_offset

    n_channel = forest%n_trees
    if (combine) then
       do ch = 1, n_channel
          if (ch == channel)  cycle
          if (active(ch)) then
             g = forest%grove_lookup(ch)
             t = ch - forest%grove(g)%tree_count_offset
             call phs_tree_combine_particles &
                  (forest%grove(g)%tree(t), forest%prt)
          end if
       end do
    end if

    !OMP PARALLEL PRIVATE (g,t,ch) SHARED(active,forest,sqrts,x,channel)
    !OMP DO SCHEDULE(STATIC)
    do ch = 1, n_channel
       if (ch == channel)  cycle
       if (active(ch)) then
          g = forest%grove_lookup(ch)
          t = ch - forest%grove(g)%tree_count_offset
          call phs_tree_compute_x_from_momenta &
               (forest%grove(g)%tree(t), &
               forest%prt, phs_factor(ch), sqrts, x(:,ch))
       end if
    end do
    !OMP END DO
    !OMP END PARALLEL

  end subroutine phs_forest_evaluate_other_channels

\eatline
\nwendcode{}\nwbegindocs{619}The complement: recover one row of the {\tt{}x} array and the
associated Jacobian entry, corresponding to
{\tt{}channel}, from incoming and outgoing momenta.  Also compute the
phase-space volume.
\nwenddocs{}\nwbegincode{620}\moddef{PHS forests: public}\plusendmoddef
  public :: phs_forest_recover_channel
\nwendcode{}\nwbegincode{621}\moddef{PHS forests: procedures}\plusendmoddef
  subroutine phs_forest_recover_channel &
       (forest, channel, sqrts, x, phs_factor, volume)
    type(phs_forest_t), intent(inout) :: forest
    integer, intent(in) :: channel
    real(default), intent(in) :: sqrts
    real(default), dimension(:,:), intent(inout) :: x
    real(default), dimension(:), intent(inout) :: phs_factor
    real(default), intent(out) :: volume
    integer :: g, t
    integer(TC) :: k, k_in
    g = forest%grove_lookup (channel)
    t = channel - forest%grove(g)%tree_count_offset
    call phs_prt_set_undefined (forest%prt)
    k_in = forest%n_tot
    forall (k = 1:forest%n_in)
       forest%prt(ibset(0,k_in-k)) = forest%prt_in(k)
    end forall
    forall (k = 1:forest%n_out)
       forest%prt(ibset(0,k-1)) = forest%prt_out(k)
    end forall
    call phs_forest_combine_particles (forest)
    call phs_tree_compute_volume &
         (forest%grove(g)%tree(t), sqrts, volume)
    call phs_tree_compute_x_from_momenta &
         (forest%grove(g)%tree(t), &
         forest%prt, phs_factor(channel), sqrts, x(:,channel))
  end subroutine phs_forest_recover_channel

\eatline
\nwendcode{}\nwbegindocs{622}\nwdocspar
\subsection{Unit tests}
Write a possible phase-space file for a $2\to 3$ process and make the
corresponding forest, print the forest. Choose some in-particle
momenta and a random-number array and evaluate out-particles and
phase-space factors. 
\nwenddocs{}\nwbegincode{623}\moddef{PHS forests: public}\plusendmoddef
  public :: phs_forest_test
\nwendcode{}\nwbegincode{624}\moddef{PHS forests: procedures}\plusendmoddef
  subroutine phs_forest_test (u, results)
    integer, intent(in) :: u
    type(test_results_t), intent(inout) :: results
  \LA{}PHS forests: execute tests\RA{}  
  end subroutine phs_forest_test
  
\nwendcode{}\nwbegindocs{625} %def phs_forest_test
\nwenddocs{}\nwbegincode{626}\moddef{PHS forests: execute tests}\endmoddef
  call test (phs_forest_1, "phs_forest_1", &
       "check phs forest setup", &
       u, results)
\nwendcode{}\nwbegincode{627}\moddef{PHS forests: tests}\endmoddef
  subroutine phs_forest_1 (u)
    use os_interface
    integer, intent(in) :: u
    type(phs_forest_t) :: forest
    type(phs_channel_t), dimension(:), allocatable :: channel
    type(model_data_t), target :: model
    type(string_t) :: process_id
    type(flavor_t), dimension(5) :: flv
    type(string_t) :: filename
    type(interaction_t) :: int
    integer, parameter :: unit_fix = 20
    type(mapping_defaults_t) :: mapping_defaults
    logical :: found_process, ok
    integer :: n_channel, ch, i
    logical, dimension(4) :: active = .true.
    real(default) :: sqrts = 1000
    real(default), dimension(5,4) :: x
    real(default), dimension(4) :: factor
    real(default) :: volume

    write (u, "(A)")  "* Test output: PHS forest"
    write (u, "(A)")  "*   Purpose: test PHS forest routines"
    write (u, "(A)")      
    
    write (u, "(A)")  "* Reading model file"
    
    call model%init_sm_test ()

    write (u, "(A)")
    write (u, "(A)")  "* Create phase-space file 'phs_forest_test.phs'"
    write (u, "(A)")
    
    call flv%init ([11, -11, 11, -11, 22], model)
    open (file="phs_forest_test.phs", unit=unit_fix, action="write")
    write (unit_fix, *) "process foo"
    write (unit_fix, *) 'md5sum_process    = "6ABA33BC2927925D0F073B1C1170780A"'
    write (unit_fix, *) 'md5sum_model_par  = "1A0B151EE6E2DEB92D880320355A3EAB"'
    write (unit_fix, *) 'md5sum_phs_config = "B6A8877058809A8BDD54753CDAB83ACE"'
    write (unit_fix, *) "sqrts         =    100.00000000000000"     
    write (unit_fix, *) "m_threshold_s =    50.000000000000000"     
    write (unit_fix, *) "m_threshold_t =    100.00000000000000"     
    write (unit_fix, *) "off_shell =            2"
    write (unit_fix, *) "t_channel =            6"
    write (unit_fix, *) "keep_nonresonant =  F"
    write (unit_fix, *) ""
    write (unit_fix, *) "  grove"
    write (unit_fix, *) "    tree 3 7"
    write (unit_fix, *) "      map 3 s_channel 23"
    write (unit_fix, *) "    tree 5 7"
    write (unit_fix, *) "    tree 6 7"
    write (unit_fix, *) "  grove"
    write (unit_fix, *) "    tree 9 11"
    write (unit_fix, *) "      map 9 t_channel 22"
    close (unit_fix)

    write (u, "(A)")
    write (u, "(A)")  "* Read phase-space file 'phs_forest_test.phs'"

    call syntax_phs_forest_init ()
    process_id = "foo"
    filename = "phs_forest_test.phs"
    call phs_forest_read &
         (forest, filename, process_id, 2, 3, model, found_process)

    write (u, "(A)")
    write (u, "(A)")  "* Set parameters, flavors, equiv, momenta"
    write (u, "(A)") 
    
    call phs_forest_set_flavors (forest, flv)
    call phs_forest_set_parameters (forest, mapping_defaults, .false.)
    call phs_forest_setup_prt_combinations (forest)
    call phs_forest_set_equivalences (forest)
    call interaction_init (int, 2, 0, 3)
    call interaction_set_momentum (int, &
         vector4_moving (500._default, 500._default, 3), 1)
    call interaction_set_momentum (int, &
         vector4_moving (500._default,-500._default, 3), 2)
    call phs_forest_set_prt_in (forest, int)
    n_channel = 2
    x = 0
    x(:,n_channel) = [0.3, 0.4, 0.1, 0.9, 0.6]
    write (u, "(A)")  "   Input values:"
    write (u, "(3x,5(1x," // FMT_12 // "))")  x(:,n_channel)

    write (u, "(A)")
    write (u, "(A)")  "* Evaluating phase space"

    call phs_forest_evaluate_selected_channel (forest, &
         n_channel, active, sqrts, x, factor, volume, ok)
    call phs_forest_evaluate_other_channels (forest, &
         n_channel, active, sqrts, x, factor, combine=.true.)
    call phs_forest_get_prt_out (forest, int)
    write (u, "(A)")  "   Output values:"
    do ch = 1, 4
       write (u, "(3x,5(1x," // FMT_12 // "))")  x(:,ch)
    end do
    call interaction_write (int, u)
    write (u, "(A)")  "   Factors:"
    write (u, "(3x,5(1x," // FMT_12 // "))")  factor
    write (u, "(A)")  "   Volume:"
    write (u, "(3x,5(1x," // FMT_12 // "))")  volume
    call phs_forest_write (forest, u)

    write (u, "(A)")
    write (u, "(A)")  "* Compute equivalences"

    n_channel = 4
    allocate (channel (n_channel))
    call phs_forest_get_equivalences (forest, &
         channel, .true.)
    do i = 1, n_channel
       write (u, "(1x,I0,':')", advance = "no")  ch       
       call channel(i)%write (u)
    end do
    
    write (u, "(A)")
    write (u, "(A)")  "* Cleanup"

    call model%final ()
    call phs_forest_final (forest)
    call syntax_phs_forest_final ()
        
    write (u, "(A)")
    write (u, "(A)")  "* Test output end: phs_forest_1"    

  end subroutine phs_forest_1

\eatline
\nwendcode{}\nwbegindocs{628}\nwdocspar
%%%%%%%%%%%%%%%%%%%%%%%%%%%%%%%%%%%%%%%%%%%%%%%%%%%%%%%%%%%%%%%%%%%%%%%%
\section{Finding phase space parameterizations}
If the phase space configuration is not found in the appropriate file,
we should generate one.

The idea is to construct all Feynman diagrams subject to certain
constraints which eliminate everything that is probably irrelevant for
the integration.  These Feynman diagrams (cascades) are grouped in
groves by finding equivalence classes related by symmetry and ordered
with respect to their importance (resonances).  Finally, the result
(or part of it) is written to file and used for the integration.

This module may eventually disappear and be replaced by CAML code.
In particular, we need here a set of Feynman rules (vertices with
particle codes, but not the factors).  Thus, the module works for the
Standard Model only.

Note that this module is stand-alone, it communicates to the main
program only via the generated ASCII phase-space configuration file.
\nwenddocs{}\nwbegincode{629}\moddef{\code{}cascades.f90\edoc{}}\endmoddef
\LA{}File header\RA{}

module cascades

\LA{}Use kinds\RA{}
  use kinds, only: TC, i8, i32
\LA{}Use strings\RA{}
  use io_units
  use format_defs, only: FMT_12, FMT_19
  use unit_tests
  use diagnostics
  use hashes
  use sorting
  use physics_defs, only: SCALAR, SPINOR, VECTOR, VECTORSPINOR, TENSOR
  use physics_defs, only: UNDEFINED
  use model_data
  use flavors
  use phs_forests

\LA{}Standard module head\RA{}

\LA{}Cascades: public\RA{}

\LA{}Cascades: parameters\RA{}

\LA{}Cascades: types\RA{}

\LA{}Cascades: interfaces\RA{}

contains

\LA{}Cascades: procedures\RA{}

\LA{}Cascades: tests\RA{}

end module cascades
\eatline
\nwendcode{}\nwbegindocs{630}\nwdocspar
\subsection{The mapping modes}
The valid mapping modes, to be used below.  We will make use of the convention
that mappings of internal particles have a positive value.  Only for positive
values, the flavor code is propagated when combining cascades.
\nwenddocs{}\nwbegincode{631}\moddef{Mapping modes}\endmoddef
  integer, parameter :: &
       & EXTERNAL_PRT = -1, &
       & NO_MAPPING = 0, S_CHANNEL = 1, T_CHANNEL =  2, U_CHANNEL = 3, &
       & RADIATION = 4, COLLINEAR = 5, INFRARED = 6, &
       & STEP_MAPPING_E = 11, STEP_MAPPING_H = 12, &
       & ON_SHELL = 99
\eatline
\eatline
\eatline
\eatline
\eatline
\nwendcode{}\nwbegincode{632}\moddef{Cascades: parameters}\endmoddef
\LA{}Mapping modes\RA{}
\nwendcode{}\nwbegindocs{633}\nwdocspar
\subsection{The cascade type}
A cascade is essentially the same as a decay tree (both definitions
may be merged in a later version).  It contains a linked tree of
nodes, each of which representing an internal particle.  In contrast
to decay trees, each node has a definite particle code.  These nodes
need not be modified, therefore we can use pointers and do not have to
copy them.  Thus, physically each cascades has only a single node, the
mother particle.  However, to be able to compare trees quickly, we
store in addition an array of binary codes which is always sorted in
ascending order.  This is accompanied by a corresponding list of
particle codes.  The index is the location of the corresponding
cascade in the cascade set, this may be used to access the daughters
directly.

The real mass is the particle mass belonging to the particle code.
The minimal mass is the sum of the real masses of all its daughters;
this is the kinematical cutoff.  The effective mass may be zero if the
particle mass is below a certain threshold; it may be the real mass if
the particle is resonant; or it may be some other value.

The logical {\tt{}t{\char95}channel} is set if this a $t$-channel line, while
{\tt{}initial} is true only for an initial particle.  Note that both
initial particles are also {\tt{}t{\char95}channel} by definition, and that they
are distinguished by the direction of the tree:  One of them decays
and is the root of the tree, while the other one is one of the leaves.

The cascade is a list of nodes (particles) which are linked via the
{\tt{}daughter} entries.  The node is the mother particle of
the decay cascade.  Much of the information in the nodes is repeated
in arrays, to be accessible more easily.  The arrays will be kept
sorted by binary codes.

The counter {\tt{}n{\char95}off{\char95}shell} is increased for each internal line that
is neither resonant nor log-enhanced.  It is set to zero if the
current line is resonant, since this implies on-shell particle production
and subsequent decay.

The counter {\tt{}n{\char95}t{\char95}channel} is non-negative once an initial particle
is included in the tree: then, it counts the number of $t$-channel lines.

The {\tt{}multiplicity} is the number of branchings to follow until all
daughters are on-shell.  A resonant or non-decaying particle has
multiplicity one.  Merging nodes, the multiplicities add unless the
mother is a resonance.  An initial or final node has multiplicity
zero.

The arrays correspond to the subnode tree {\tt{}tree} of the current
cascade.  PDG codes are stored only for those positions which are
resonant, with the exception of the last entry, i.e., the current node.
Other positions, in particular external legs, are assigned undefined
PDG code.

A cascade is uniquely identified by its tree, the tree of PDG codes,
and the tree of mappings.  The tree of resonances is kept only to mask
the PDG tree as described above.
\nwenddocs{}\nwbegincode{634}\moddef{Cascades: types}\endmoddef
  type :: cascade_t
     private
     ! counters
     integer :: index = 0
     integer :: grove = 0
     ! status
     logical :: active = .false.
     logical :: complete = .false.
     logical :: incoming = .false.
     ! this node
     integer(TC) :: bincode = 0
     type(flavor_t) :: flv
     integer :: pdg = UNDEFINED
     logical :: is_vector = .false.
     real(default) :: m_min = 0
     real(default) :: m_rea = 0
     real(default) :: m_eff = 0
     integer :: mapping = NO_MAPPING
     logical :: on_shell = .false.
     logical :: resonant = .false.
     logical :: log_enhanced = .false.
     logical :: t_channel = .false.
     ! global tree properties
     integer :: multiplicity = 0
     integer :: internal = 0
     integer :: n_off_shell = 0
     integer :: n_resonances = 0
     integer :: n_log_enhanced = 0
     integer :: n_t_channel = 0
     integer :: res_hash = 0
     ! the sub-node tree
     integer :: depth = 0
     integer(TC), dimension(:), allocatable :: tree
     integer, dimension(:), allocatable :: tree_pdg
     integer, dimension(:), allocatable :: tree_mapping
     logical, dimension(:), allocatable :: tree_resonant
     ! branch connections
     logical :: has_children = .false.
     type(cascade_t), pointer :: daughter1 => null ()
     type(cascade_t), pointer :: daughter2 => null ()
     type(cascade_t), pointer :: mother => null ()
     ! next in list
     type(cascade_t), pointer :: next => null ()
  end type cascade_t

\eatline
\nwendcode{}\nwbegincode{635}\moddef{Cascades: procedures}\endmoddef
  subroutine cascade_init (cascade, depth)
    type(cascade_t), intent(out) :: cascade
    integer, intent(in) :: depth
    integer, save :: index = 0
    index = cascade_index ()
    cascade%index = index
    cascade%depth = depth
    cascade%active = .true.
    allocate (cascade%tree (depth))
    allocate (cascade%tree_pdg (depth))
    allocate (cascade%tree_mapping (depth))
    allocate (cascade%tree_resonant (depth))
  end subroutine cascade_init
\eatline
\nwendcode{}\nwbegindocs{636}Keep and increment a global index
\nwenddocs{}\nwbegincode{637}\moddef{Cascades: procedures}\plusendmoddef
  function cascade_index (seed) result (index)
    integer :: index
    integer, intent(in), optional :: seed
    integer, save :: i = 0
    if (present (seed))  i = seed
    i = i + 1
    index = i
  end function cascade_index

\eatline
\nwendcode{}\nwbegindocs{638}We need three versions of writing cascades.  This goes to the
phase-space file:
\nwenddocs{}\nwbegincode{639}\moddef{Cascades: procedures}\plusendmoddef
  subroutine cascade_write_file_format (cascade, model, unit)
    type(cascade_t), intent(in) :: cascade
    class(model_data_t), intent(in), target :: model
    integer, intent(in), optional :: unit
    type(flavor_t) :: flv
    integer :: u, i
1   format(3x,A,1x,40(1x,I4))
2   format(3x,A,1x,I3,1x,A,1x,I7,1x,'!',1x,A)
    u = given_output_unit (unit);  if (u < 0)  return
    call write_reduced (cascade%tree, u)
    write (u, "(A)")
    do i = 1, cascade%depth
       call flv%init (cascade%tree_pdg(i), model)
       select case (cascade%tree_mapping(i))
       case (NO_MAPPING, EXTERNAL_PRT)
       case (S_CHANNEL)
          write(u,2) 'map', &
               cascade%tree(i), 's_channel', abs (cascade%tree_pdg(i)), &
               char (flv%get_name ())
       case (T_CHANNEL)
          write(u,2) 'map', &
               cascade%tree(i), 't_channel', abs (cascade%tree_pdg(i)), &
               char (flv%get_name ())
       case (U_CHANNEL)
          write(u,2) 'map', &
               cascade%tree(i), 'u_channel', abs (cascade%tree_pdg(i)), &
               char (flv%get_name ())
       case (RADIATION)
          write(u,2) 'map', &
               cascade%tree(i), 'radiation', abs (cascade%tree_pdg(i)), &
               char (flv%get_name ())
       case (COLLINEAR)
          write(u,2) 'map', &
               cascade%tree(i), 'collinear', abs (cascade%tree_pdg(i)), &
               char (flv%get_name ())
       case (INFRARED)
          write(u,2) 'map', &
               cascade%tree(i), 'infrared ',  abs (cascade%tree_pdg(i)), &
               char (flv%get_name ())
       case (ON_SHELL)
          write(u,2) 'map', &
               cascade%tree(i), 'on_shell ', abs (cascade%tree_pdg(i)), &
               char (flv%get_name ())
       case default
          call msg_bug (" Impossible mapping mode encountered")
       end select
    end do
  contains
    subroutine write_reduced (array, unit)
      integer(TC), dimension(:), intent(in) :: array
      integer, intent(in) :: unit
      integer :: i
      write (u, "(3x,A,1x)", advance="no")  "tree"
      do i = 1, size (array)
         if (decay_level (array(i)) > 1) then
            write (u, "(1x,I0)", advance="no")  array(i)
         end if
      end do
    end subroutine write_reduced

    elemental function decay_level (k) result (l)
      integer(TC), intent(in) :: k
      integer :: l
      integer :: i
      l = 0
      do i = 0, bit_size(k) - 1
         if (btest(k,i)) l = l + 1
      end do
    end function decay_level
    subroutine start_comment (u)
      integer, intent(in) :: u
      write(u, '(1x,A)', advance='no') '!'
    end subroutine start_comment
  end subroutine cascade_write_file_format

\eatline
\nwendcode{}\nwbegindocs{640}This creates metapost source for graphical display:
\nwenddocs{}\nwbegincode{641}\moddef{Cascades: procedures}\plusendmoddef
  subroutine cascade_write_graph_format (cascade, count, unit)
    type(cascade_t), intent(in) :: cascade
    integer, intent(in) :: count
    integer, intent(in), optional :: unit
    integer :: u
    integer(TC) :: mask
    type(string_t) :: left_str, right_str
    u = given_output_unit (unit);  if (u < 0)  return
    mask = 2**((cascade%depth+3)/2) - 1
    left_str = ""
    right_str = ""
    write (u, '(A)') "\\begin\{minipage\}\{105pt\}"
    write (u, '(A)') "\\vspace\{30pt\}"
    write (u, '(A)') "\\begin\{center\}"
    write (u, '(A)') "\\begin\{fmfgraph*\}(55,55)"
    call graph_write (cascade, mask)
    write (u, '(A)') "\\fmfleft\{" // char (extract (left_str, 2)) // "\}"
    write (u, '(A)') "\\fmfright\{" // char (extract (right_str, 2)) // "\}"
    write (u, '(A)') "\\end\{fmfgraph*\}\\\\"
    write (u, '(A,I5,A)') "\\fbox\{$", count, "$\}"
    write (u, '(A)') "\\end\{center\}"
    write (u, '(A)') "\\end\{minipage\}"
    write (u, '(A)') "%"
  contains
    recursive subroutine graph_write (cascade, mask, reverse)
      type(cascade_t), intent(in) :: cascade
      integer(TC), intent(in) :: mask
      logical, intent(in), optional :: reverse
      type(flavor_t) :: anti
      logical :: rev
      rev = .false.;  if (present(reverse))  rev = reverse
      if (cascade%has_children) then
         if (.not.rev) then
            call vertex_write (cascade, cascade%daughter1, mask)
            call vertex_write (cascade, cascade%daughter2, mask)
         else
            call vertex_write (cascade, cascade%daughter2, mask, .true.)
            call vertex_write (cascade, cascade%daughter1, mask, .true.)
         end if
         if (cascade%complete) then
            call vertex_write (cascade, cascade%mother, mask, .true.)
            write (u, '(A,I0,A)') "\\fmfv\{d.shape=square\}\{v0\}"
         end if
      else
         if (cascade%incoming) then
            anti = cascade%flv%anti ()
            call external_write (cascade%bincode, anti%get_tex_name (), &
                 left_str)
         else
            call external_write (cascade%bincode, cascade%flv%get_tex_name (), &
                 right_str)
         end if
      end if
    end subroutine graph_write
    recursive subroutine vertex_write (cascade, daughter, mask, reverse)
      type(cascade_t), intent(in) :: cascade, daughter
      integer(TC), intent(in) :: mask
      logical, intent(in), optional :: reverse
      integer :: bincode
      if (cascade%complete) then
         bincode = 0
      else
         bincode = cascade%bincode
      end if
      call graph_write (daughter, mask, reverse)
      if (daughter%has_children) then
         call line_write (bincode, daughter%bincode, daughter%flv, &
              mapping=daughter%mapping)
      else
         call line_write (bincode, daughter%bincode, daughter%flv)
      end if
    end subroutine vertex_write
    subroutine line_write (i1, i2, flv, mapping)
      integer(TC), intent(in) :: i1, i2
      type(flavor_t), intent(in) :: flv
      integer, intent(in), optional :: mapping
      integer :: k1, k2
      type(string_t) :: prt_type
      select case (flv%get_spin_type ())
      case (SCALAR);       prt_type = "plain"
      case (SPINOR);       prt_type = "fermion"
      case (VECTOR);       prt_type = "boson"
      case (VECTORSPINOR); prt_type = "fermion"
      case (TENSOR);       prt_type = "dbl_wiggly"
      case default;        prt_type = "dashes"
      end select
      if (flv%is_antiparticle ()) then
         k1 = i2;  k2 = i1
      else
         k1 = i1;  k2 = i2
      end if
      if (present (mapping)) then
         select case (mapping)
         case (S_CHANNEL)
            write (u, '(A,I0,A,I0,A)') "\\fmf\{" // char (prt_type) // &
                 & ",f=blue,lab=\\sm\\blue$" // &
                 & char (flv%get_tex_name ()) // "$\}" // &
                 & "\{v", k1, ",v", k2, "\}"
         case (T_CHANNEL, U_CHANNEL)
            write (u, '(A,I0,A,I0,A)') "\\fmf\{" // char (prt_type) // &
                 & ",f=cyan,lab=\\sm\\cyan$" // &
                 & char (flv%get_tex_name ()) // "$\}" // &
                 & "\{v", k1, ",v", k2, "\}"
         case (RADIATION)
            write (u, '(A,I0,A,I0,A)') "\\fmf\{" // char (prt_type) // &
                 & ",f=green,lab=\\sm\\green$" // &
                 & char (flv%get_tex_name ()) // "$\}" // &
                 & "\{v", k1, ",v", k2, "\}"
         case (COLLINEAR)
            write (u, '(A,I0,A,I0,A)') "\\fmf\{" // char (prt_type) // &
                 & ",f=magenta,lab=\\sm\\magenta$" // &
                 & char (flv%get_tex_name ()) // "$\}" // &
                 & "\{v", k1, ",v", k2, "\}"
         case (INFRARED)
            write (u, '(A,I0,A,I0,A)') "\\fmf\{" // char (prt_type) // &
                 & ",f=red,lab=\\sm\\red$" // &
                 & char (flv%get_tex_name ()) // "$\}" // &
                 & "\{v", k1, ",v", k2, "\}"
         case default
            write (u, '(A,I0,A,I0,A)') "\\fmf\{" // char (prt_type) // &
                 & ",f=black\}" // &
                 & "\{v", k1, ",v", k2, "\}"
         end select
      else
         write (u, '(A,I0,A,I0,A)') "\\fmf\{" // char (prt_type) // &
                 & "\}" // &
                 & "\{v", k1, ",v", k2, "\}"
      end if
    end subroutine line_write
    subroutine external_write (bincode, name, ext_str)
      integer(TC), intent(in) :: bincode
      type(string_t), intent(in) :: name
      type(string_t), intent(inout) :: ext_str
      character(len=20) :: str
      write (str, '(A2,I0)') ",v", bincode
      ext_str = ext_str // trim (str)
      write (u, '(A,I0,A,I0,A)') "\\fmflabel\{\\sm$" &
        // char (name) &
        // "\\,(", bincode, ")" &
        // "$\}\{v", bincode, "\}"
    end subroutine external_write
  end subroutine cascade_write_graph_format

\eatline
\nwendcode{}\nwbegindocs{642}This is for screen/debugging output:
\nwenddocs{}\nwbegincode{643}\moddef{Cascades: procedures}\plusendmoddef
  subroutine cascade_write (cascade, unit)
    type(cascade_t), intent(in) :: cascade
    integer, intent(in), optional :: unit
    integer :: u
    character(9) :: depth
    u = given_output_unit (unit);  if (u < 0)  return
    write (u, "(A,(1x,I7))") 'Cascade #', cascade%index
    write (u, "(A,(1x,I7))") '  Grove:       #', cascade%grove
    write (u, "(A,3(1x,L1))") '  act/cmp/inc:  ', &
         cascade%active, cascade%complete, cascade%incoming
    write (u, "(A,I0)") '  Bincode:      ', cascade%bincode
    write (u, "(A)", advance="no") '  Flavor:       '
    call cascade%flv%write (unit)
    write (u, "(A,I9)") '  Active flavor:', cascade%pdg
    write (u, "(A,L1)") '  Is vector:    ', cascade%is_vector
    write (u, "(A,3(1x," // FMT_19 // "))") '  Mass (m/r/e): ', &
         cascade%m_min, cascade%m_rea, cascade%m_eff
    write (u, "(A,I1)") '  Mapping:      ', cascade%mapping
    write (u, "(A,3(1x,L1))") '  res/log/tch:  ', &
         cascade%resonant, cascade%log_enhanced, cascade%t_channel
    write (u, "(A,(1x,I7))") '  Multiplicity: ', cascade%multiplicity
    write (u, "(A,2(1x,I7))") '  n intern/off: ', &
         cascade%internal, cascade%n_off_shell
    write (u, "(A,3(1x,I7))") '  n res/log/tch:', &
         cascade%n_resonances, cascade%n_log_enhanced, cascade%n_t_channel
    write (u, "(A,I7)") '  Depth:        ', cascade%depth
    write (depth, "(I7)") cascade%depth
    write (u, "(A," // depth // "(1x,I7))") &
       '  Tree:         ', cascade%tree
    write (u, "(A," // depth // "(1x,I7))") &
       '  Tree(PDG):    ', cascade%tree_pdg
    write (u, "(A," // depth // "(1x,I7))") &
       '  Tree(mapping):', cascade%tree_mapping
    write (u, "(A," // depth // "(1x,L1))") &
       '  Tree(res):    ', cascade%tree_resonant
    if (cascade%has_children) then
       write (u, "(A,I7,1x,I7)") '  Daughter1/2:  ', &
            cascade%daughter1%index, cascade%daughter2%index
    end if
    if (associated (cascade%mother)) then
       write (u, "(A,I7)") '  Mother:       ', cascade%mother%index
    end if
  end subroutine cascade_write

\eatline
\nwendcode{}\nwbegindocs{644}\nwdocspar
\subsection{Creating new cascades}
This initializes a single-particle cascade (external, final state).
The PDG entry in the tree is set undefined because the cascade is not
resonant.  However, the flavor entry is set, so the cascade flavor
is identified nevertheless.
\nwenddocs{}\nwbegincode{645}\moddef{Cascades: procedures}\plusendmoddef
  subroutine cascade_init_outgoing (cascade, flv, pos, m_thr)
    type(cascade_t), intent(out) :: cascade
    type(flavor_t), intent(in) :: flv
    integer, intent(in) :: pos
    real(default), intent(in) :: m_thr
    call cascade_init (cascade, 1)
    cascade%bincode = ibset (0_TC, pos-1)
    cascade%flv = flv
    cascade%pdg = abs (cascade%flv%get_pdg ())
    cascade%is_vector = flv%get_spin_type () == VECTOR
    cascade%m_min = flv%get_mass ()
    cascade%m_rea = cascade%m_min
    if (cascade%m_rea >= m_thr) then
       cascade%m_eff = cascade%m_rea
    end if
    cascade%on_shell = .true.
    cascade%multiplicity = 1
    cascade%tree(1) = cascade%bincode
    cascade%tree_pdg(1) = cascade%pdg
    cascade%tree_mapping(1) = EXTERNAL_PRT
    cascade%tree_resonant(1) = .false.
  end subroutine cascade_init_outgoing

\eatline
\nwendcode{}\nwbegindocs{646}The same for an incoming line:
\nwenddocs{}\nwbegincode{647}\moddef{Cascades: procedures}\plusendmoddef
  subroutine cascade_init_incoming (cascade, flv, pos, m_thr)
    type(cascade_t), intent(out) :: cascade
    type(flavor_t), intent(in) :: flv
    integer, intent(in) :: pos
    real(default), intent(in) :: m_thr
    call cascade_init (cascade, 1)
    cascade%incoming = .true.
    cascade%bincode = ibset (0_TC, pos-1)
    cascade%flv = flv%anti ()
    cascade%pdg = abs (flv%get_pdg ())
    cascade%is_vector = flv%get_spin_type () == VECTOR
    cascade%m_min = flv%get_mass ()
    cascade%m_rea = cascade%m_min
    if (cascade%m_rea >= m_thr) then
       cascade%m_eff = cascade%m_rea
    end if
    cascade%on_shell = .true.
    cascade%n_t_channel = 0
    cascade%n_off_shell = 0
    cascade%tree(1) = cascade%bincode
    cascade%tree_pdg(1) = cascade%pdg
    cascade%tree_mapping(1) = EXTERNAL_PRT
    cascade%tree_resonant(1) = .false.
  end subroutine cascade_init_incoming

\eatline
\nwendcode{}\nwbegindocs{648}\nwdocspar
\subsection{Tools}
This function returns true if the two cascades share no common
external particle.  This is a requirement for joining them.
\nwenddocs{}\nwbegincode{649}\moddef{Cascades: interfaces}\endmoddef
  interface operator(.disjunct.)
     module procedure cascade_disjunct
  end interface

\nwendcode{}\nwbegincode{650}\moddef{Cascades: procedures}\plusendmoddef
  function cascade_disjunct (cascade1, cascade2) result (flag)
    logical :: flag
    type(cascade_t), intent(in) :: cascade1, cascade2
    flag = iand (cascade1%bincode, cascade2%bincode) == 0
  end function cascade_disjunct

\eatline
\eatline
\nwendcode{}\nwbegindocs{651}Compute a hash code for the resonance pattern of a cascade.  We count the
number of times each particle appears as a resonance.

We pack the PDG codes of the resonances in two arrays (s-channel and 
t-channel), sort them both, concatenate the results, transfer to
{\tt{}i8} integers, and compute the hash code from this byte stream.
\nwenddocs{}\nwbegincode{652}\moddef{Cascades: procedures}\plusendmoddef
  subroutine cascade_assign_resonance_hash (cascade)
    type(cascade_t), intent(inout) :: cascade
    integer(i8), dimension(1) :: mold
    cascade%res_hash = hash (transfer &
         (concat (sort (pack (cascade%tree_pdg, &
                              cascade%tree_resonant)), &
                  sort (pack (cascade%tree_pdg, &
                              cascade%tree_mapping == T_CHANNEL .or. &
                              cascade%tree_mapping == U_CHANNEL))), &
          mold))
  end subroutine cascade_assign_resonance_hash

\eatline
\nwendcode{}\nwbegindocs{653}\nwdocspar
\subsection{Hash entries for cascades}
We will set up a hash array which contains keys of and pointers to
cascades.  We hold a list of cascade (pointers) within each bucket.
This is not for collision resolution, but for keeping similar, but
unequal cascades together.
\nwenddocs{}\nwbegincode{654}\moddef{Cascades: types}\plusendmoddef
  type :: cascade_p
     type(cascade_t), pointer :: cascade => null ()
     type(cascade_p), pointer :: next => null ()
  end type cascade_p

\eatline
\nwendcode{}\nwbegindocs{655}Here is the bucket or hash entry type:
\nwenddocs{}\nwbegincode{656}\moddef{Cascades: types}\plusendmoddef
  type :: hash_entry_t
     integer(i32) :: hashval = 0
     integer(i8), dimension(:), allocatable :: key
     type(cascade_p), pointer :: first => null ()
     type(cascade_p), pointer :: last => null ()
  end type hash_entry_t

\eatline
\nwendcode{}\nwbegindocs{657}Finalize: just deallocate the list; the contents are just pointers.
\nwenddocs{}\nwbegincode{658}\moddef{Cascades: procedures}\plusendmoddef
  subroutine hash_entry_final (hash_entry)
    type(hash_entry_t), intent(inout) :: hash_entry
    type(cascade_p), pointer :: current
    do while (associated (hash_entry%first))
       current => hash_entry%first
       hash_entry%first => current%next
       deallocate (current)
    end do
  end subroutine hash_entry_final

\eatline
\nwendcode{}\nwbegindocs{659}Output: concise format for debugging, just list cascade indices.
\nwenddocs{}\nwbegincode{660}\moddef{Cascades: procedures}\plusendmoddef
  subroutine hash_entry_write (hash_entry, unit)
    type(hash_entry_t), intent(in) :: hash_entry
    integer, intent(in), optional :: unit
    type(cascade_p), pointer :: current
    integer :: u, i
    u = given_output_unit (unit);  if (u < 0)  return
    write (u, "(1x,A)", advance="no")  "Entry:"
    do i = 1, size (hash_entry%key)
       write (u, "(1x,I0)", advance="no")  hash_entry%key(i)
    end do
    write (u, "(1x,A)", advance="no")  "->"
    current => hash_entry%first
    do while (associated (current))
       write (u, "(1x,I7)", advance="no") current%cascade%index
       current => current%next
    end do
    write (u, *)
  end subroutine hash_entry_write

\eatline
\nwendcode{}\nwbegindocs{661}This function adds a cascade pointer to the bucket.  If {\tt{}ok} is
present, check first if it is already there and return failure if yes.
If {\tt{}cascade{\char95}ptr} is also present, set it to the current cascade if
successful.  If not, set it to the cascade that is already there.
\nwenddocs{}\nwbegincode{662}\moddef{Cascades: procedures}\plusendmoddef
  subroutine hash_entry_add_cascade_ptr (hash_entry, cascade, ok, cascade_ptr)
    type(hash_entry_t), intent(inout) :: hash_entry
    type(cascade_t), intent(in), target :: cascade
    logical, intent(out), optional :: ok
    type(cascade_t), optional, pointer :: cascade_ptr
    type(cascade_p), pointer :: current
    if (present (ok)) then
       call hash_entry_check_cascade (hash_entry, cascade, ok, cascade_ptr)
       if (.not. ok)  return
    end if
    allocate (current)
    current%cascade => cascade
    if (associated (hash_entry%last)) then
       hash_entry%last%next => current
    else 
       hash_entry%first => current
    end if
    hash_entry%last => current
  end subroutine hash_entry_add_cascade_ptr

\eatline
\nwendcode{}\nwbegindocs{663}This function checks whether a cascade is already in the bucket.
For incomplete cascades, we look for an exact match. It should suffice
to verify the tree, the PDG codes, and the mapping modes.  This is the
information that is written to the phase space file.

For complete cascades, we ignore the PDG code at positions with
mappings infrared, collinear, or t/u-channel.  Thus a cascade which is
distinguished only by PDG code at such places, is flagged existent.
If the convention is followed that light particles come before heavier
ones (in the model definition), this ensures that the lightest
particle is kept in the appropriate place, corresponding to the
strongest peak.

For external cascades (incoming/outgoing) we take the PDG code into
account even though it is zeroed in the PDG-code tree.
\nwenddocs{}\nwbegincode{664}\moddef{Cascades: procedures}\plusendmoddef
  subroutine hash_entry_check_cascade (hash_entry, cascade, ok, cascade_ptr)
    type(hash_entry_t), intent(in), target :: hash_entry
    type(cascade_t), intent(in), target :: cascade
    logical, intent(out) :: ok
    type(cascade_t), optional, pointer :: cascade_ptr
    type(cascade_p), pointer :: current
    integer, dimension(:), allocatable :: tree_pdg
    ok = .true.
    allocate (tree_pdg (size (cascade%tree_pdg)))
    if (cascade%complete) then
       where (cascade%tree_mapping == INFRARED .or. &
            cascade%tree_mapping == COLLINEAR .or. &
            cascade%tree_mapping == T_CHANNEL .or. &
            cascade%tree_mapping == U_CHANNEL)
          tree_pdg = 0
       elsewhere
          tree_pdg = cascade%tree_pdg
       end where
    else
       tree_pdg = cascade%tree_pdg
    end if
    current => hash_entry%first
    do while (associated (current))
       if (current%cascade%depth == cascade%depth) then
          if (all (current%cascade%tree == cascade%tree)) then
             if (all (current%cascade%tree_mapping == cascade%tree_mapping)) &
                  then
                if (all (current%cascade%tree_pdg .match. tree_pdg)) then
                   if (present (cascade_ptr))  cascade_ptr => current%cascade
                   ok = .false.;  return
                end if
             end if
          end if
       end if
       current => current%next
    end do
    if (present (cascade_ptr))  cascade_ptr => cascade
  end subroutine hash_entry_check_cascade

\eatline
\nwendcode{}\nwbegindocs{665}For PDG codes, we specify that the undefined code matches any code.
This is already defined for flavor objects, but here we need it for
the codes themselves.
\nwenddocs{}\nwbegincode{666}\moddef{Cascades: interfaces}\plusendmoddef
  interface operator(.match.)
     module procedure pdg_match
  end interface
\nwendcode{}\nwbegincode{667}\moddef{Cascades: procedures}\plusendmoddef
  elemental function pdg_match (pdg1, pdg2) result (flag)
    logical :: flag
    integer(TC), intent(in) :: pdg1, pdg2
    select case (pdg1)
    case (0)
       flag = .true.
    case default
       select case (pdg2)
       case (0)
          flag = .true.
       case default
          flag = pdg1 == pdg2
       end select
    end select
  end function pdg_match

\eatline
\nwendcode{}\nwbegindocs{668}\nwdocspar
\subsection{The cascade set}
The cascade set will later be transformed into the decay forest.  It
is set up as a linked list.  In addition to the usual {\tt{}first} and
{\tt{}last} pointers, there is a {\tt{}first{\char95}t} pointer which points to the
first t-channel cascade (after all s-channel cascades), and a
{\tt{}first{\char95}k} pointer which points to the first final cascade (with a
keystone).

As an auxiliary device, the object contains a hash array with
associated parameters where an additional pointer is stored for each
cascade.  The keys are made from the relevant cascade data.  This hash
is used for fast detection (and thus avoidance) of double entries in
the cascade list.
\nwenddocs{}\nwbegincode{669}\moddef{Cascades: public}\endmoddef
  public :: cascade_set_t
\nwendcode{}\nwbegincode{670}\moddef{Cascades: types}\plusendmoddef
  type :: cascade_set_t
     private
     class(model_data_t), pointer :: model
     integer :: n_in, n_out, n_tot
     type(flavor_t), dimension(:,:), allocatable :: flv
     integer :: depth_out, depth_tot
     real(default) :: sqrts = 0
     real(default) :: m_threshold_s = 0
     real(default) :: m_threshold_t = 0
     integer :: off_shell = 0
     integer :: t_channel = 0
     logical :: keep_nonresonant
     integer :: n_groves = 0
     ! The cascade list
     type(cascade_t), pointer :: first => null ()
     type(cascade_t), pointer :: last => null ()
     type(cascade_t), pointer :: first_t => null ()
     type(cascade_t), pointer :: first_k => null ()
     ! The hashtable
     integer :: n_entries = 0
     real :: fill_ratio = 0
     integer :: n_entries_max = 0
     integer(i32) :: mask = 0
     logical :: fatal_beam_decay = .true.
     type(hash_entry_t), dimension(:), allocatable :: entry
  end type cascade_set_t

\eatline
\nwendcode{}\nwbegindocs{671}Return true if there are cascades which are active and complete, so
the phase space file would be nonempty.
\nwenddocs{}\nwbegincode{672}\moddef{Cascades: public}\plusendmoddef
  public :: cascade_set_is_valid
\nwendcode{}\nwbegincode{673}\moddef{Cascades: procedures}\plusendmoddef
  function cascade_set_is_valid (cascade_set) result (flag)
    logical :: flag
    type(cascade_set_t), intent(in) :: cascade_set
    type(cascade_t), pointer :: cascade
    flag = .false.
    cascade => cascade_set%first_k
    do while (associated (cascade))
       if (cascade%active .and. cascade%complete) then
          flag = .true.
          return
       end if
       cascade => cascade%next
    end do
  end function cascade_set_is_valid

\eatline
\nwendcode{}\nwbegindocs{674}The initializer sets up the hash table with some initial size
guessed by looking at the number of external particles.  We choose 256
for 3 external particles and a factor of 4 for each additional
particle, limited at $2^{30}$=1G.

Note: the explicit initialization loop might be avoided (ELEMENTAL),
but a bug in nagfor 5.3.2 prevents this.
\nwenddocs{}\nwbegincode{675}\moddef{Cascades: parameters}\plusendmoddef
  real, parameter, public :: CASCADE_SET_FILL_RATIO = 0.1
\nwendcode{}\nwbegincode{676}\moddef{Cascades: procedures}\plusendmoddef
  subroutine cascade_set_init (cascade_set, model, n_in, n_out, phs_par, &
        fatal_beam_decay, flv)
    type(cascade_set_t), intent(out) :: cascade_set
    class(model_data_t), intent(in), target :: model
    integer, intent(in) :: n_in, n_out
    type(phs_parameters_t), intent(in) :: phs_par
    logical, intent(in) :: fatal_beam_decay
    type(flavor_t), dimension(:,:), intent(in), optional :: flv
    integer :: size_guess
    integer :: i, j
    cascade_set%model => model
    cascade_set%n_in = n_in
    cascade_set%n_out = n_out
    cascade_set%n_tot = n_in + n_out
    if (present (flv)) then
       allocate (cascade_set%flv (size (flv, 1), size (flv, 2)))
       do i = 1, size (flv, 2)
          do j = 1, size (flv, 1)
             call cascade_set%flv(j,i)%init (flv(j,i)%get_pdg (), model)
          end do
       end do
    end if
    select case (n_in)
    case (1);  cascade_set%depth_out = 2 * n_out - 3
    case (2);  cascade_set%depth_out = 2 * n_out - 1
    end select
    cascade_set%depth_tot = 2 * cascade_set%n_tot - 3
    cascade_set%sqrts = phs_par%sqrts
    cascade_set%m_threshold_s = phs_par%m_threshold_s
    cascade_set%m_threshold_t = phs_par%m_threshold_t
    cascade_set%off_shell = phs_par%off_shell
    cascade_set%t_channel = phs_par%t_channel
    cascade_set%keep_nonresonant = phs_par%keep_nonresonant
    cascade_set%fill_ratio = CASCADE_SET_FILL_RATIO
    size_guess = ishft (256, min (2 * (cascade_set%n_tot - 3), 22))
    cascade_set%n_entries_max = size_guess * cascade_set%fill_ratio
    cascade_set%mask = size_guess - 1
    allocate (cascade_set%entry (0:cascade_set%mask))
    cascade_set%fatal_beam_decay = fatal_beam_decay
  end subroutine cascade_set_init
    
\eatline
\nwendcode{}\nwbegindocs{677}The finalizer has to delete both the hash and the list.
\nwenddocs{}\nwbegincode{678}\moddef{Cascades: public}\plusendmoddef
  public :: cascade_set_final
\nwendcode{}\nwbegincode{679}\moddef{Cascades: procedures}\plusendmoddef
  subroutine cascade_set_final (cascade_set)
    type(cascade_set_t), intent(inout), target :: cascade_set
    type(cascade_t), pointer :: current
    integer :: i
    if (allocated (cascade_set%entry)) then
       do i = 0, cascade_set%mask
          call hash_entry_final (cascade_set%entry(i))
       end do
       deallocate (cascade_set%entry)
    end if
    do while (associated (cascade_set%first))
       current => cascade_set%first
       cascade_set%first => cascade_set%first%next
       deallocate (current)
    end do
  end subroutine cascade_set_final

\eatline
\nwendcode{}\nwbegindocs{680}Write the process in ASCII format, in columns that are headed by the
corresponding bincode.
\nwenddocs{}\nwbegincode{681}\moddef{Cascades: public}\plusendmoddef
  public :: cascade_set_write_process_bincode_format
\nwendcode{}\nwbegincode{682}\moddef{Cascades: procedures}\plusendmoddef
  subroutine cascade_set_write_process_bincode_format (cascade_set, unit)
    type(cascade_set_t), intent(in), target :: cascade_set
    integer, intent(in), optional :: unit
    integer, dimension(:), allocatable :: bincode, field_width
    integer :: n_in, n_out, n_tot, n_flv
    integer :: u, f, i, bc
    character(20) :: str
    type(string_t) :: fmt_head
    type(string_t), dimension(:), allocatable :: fmt_proc
    u = given_output_unit (unit);  if (u < 0)  return
    if (.not. allocated (cascade_set%flv)) return
    write (u, "('!',1x,A)")  "List of subprocesses with particle bincodes:"
    n_in  = cascade_set%n_in
    n_out = cascade_set%n_out
    n_tot = cascade_set%n_tot
    n_flv = size (cascade_set%flv, 2)
    allocate (bincode (n_tot), field_width (n_tot), fmt_proc (n_tot))
    bc = 1
    do i = 1, n_out
       bincode(n_in + i) = bc
       bc = 2 * bc
    end do
    do i = n_in, 1, -1
       bincode(i) = bc
       bc = 2 * bc
    end do       
    do i = 1, n_tot
       write (str, "(I0)")  bincode(i)
       field_width(i) = len_trim (str)
       do f = 1, n_flv
          field_width(i) = max (field_width(i), &
               len (cascade_set%flv(i,f)%get_name ()))
       end do
    end do
    fmt_head = "('!'"
    do i = 1, n_tot
       fmt_head = fmt_head // ",1x,"
       fmt_proc(i) = "(1x,"
       write (str, "(I0)")  field_width(i)
       fmt_head = fmt_head // "I" // trim(str)
       fmt_proc(i) = fmt_proc(i) // "A" // trim(str)
       if (i == n_in) then
          fmt_head = fmt_head // ",1x,'  '"
       end if
    end do
    fmt_proc = fmt_proc // ")"
    fmt_head = fmt_head // ")"
    write (u, char (fmt_head))  bincode
    do f = 1, n_flv
       write (u, "('!')", advance="no")
       do i = 1, n_tot
          write (u, char (fmt_proc(i)), advance="no") &
               char (cascade_set%flv(i,f)%get_name ())
          if (i == n_in)  write (u, "(1x,'=>')", advance="no")
       end do
       write (u, *)
    end do
    write (u, char (fmt_head))  bincode
  end subroutine cascade_set_write_process_bincode_format

\eatline
\nwendcode{}\nwbegindocs{683}Write the process as a \LaTeX\ expression.
\nwenddocs{}\nwbegincode{684}\moddef{Cascades: procedures}\plusendmoddef
  subroutine cascade_set_write_process_tex_format (cascade_set, unit)
    type(cascade_set_t), intent(in), target :: cascade_set
    integer, intent(in), optional :: unit
    integer :: u, f, i
    u = given_output_unit (unit);  if (u < 0)  return
    if (.not. allocated (cascade_set%flv)) return
    write (u, "(A)")  "\\begin\{align*\}"
    do f = 1, size (cascade_set%flv, 2)
       do i = 1, cascade_set%n_in
          if (i > 1)  write (u, "(A)", advance="no") "\\quad "
          write (u, "(A)", advance="no") &
               char (cascade_set%flv(i,f)%get_tex_name ())
       end do
       write (u, "(A)", advance="no")  "\\quad &\\to\\quad "
       do i = cascade_set%n_in + 1, cascade_set%n_tot
          if (i > cascade_set%n_in + 1)  write (u, "(A)", advance="no") "\\quad "
          write (u, "(A)", advance="no") &
               char (cascade_set%flv(i,f)%get_tex_name ())
       end do
       if (f < size (cascade_set%flv, 2)) then
          write (u, "(A)")  "\\\\"
       else
          write (u, "(A)")  ""
       end if
    end do       
    write (u, "(A)")  "\\end\{align*\}"
  end subroutine cascade_set_write_process_tex_format

\eatline
\nwendcode{}\nwbegindocs{685}Three output routines: phase-space file, graph source code, and
screen output.

This version generates the phase space file.  It deals only with
complete cascades.
\nwenddocs{}\nwbegincode{686}\moddef{Cascades: public}\plusendmoddef
  public :: cascade_set_write_file_format
\nwendcode{}\nwbegincode{687}\moddef{Cascades: procedures}\plusendmoddef
  subroutine cascade_set_write_file_format (cascade_set, unit)
    type(cascade_set_t), intent(in), target :: cascade_set
    integer, intent(in), optional :: unit
    type(cascade_t), pointer :: cascade
    integer :: u, grove, count
    logical :: first_in_grove
    u = given_output_unit (unit);  if (u < 0)  return
    count = 0
    do grove = 1, cascade_set%n_groves
       first_in_grove = .true.
       cascade => cascade_set%first_k
       do while (associated (cascade))
          if (cascade%active .and. cascade%complete) then
             if (cascade%grove == grove) then
                if (first_in_grove) then
                   first_in_grove = .false.
                   write (u, *)
                   write (u, "(1x,'!',1x,A,1x,I0,A)", advance='no') &
                      'Multiplicity =', cascade%multiplicity, ","
                   select case (cascade%n_resonances)
                   case (0)
                      write (u, '(1x,A)', advance='no') 'no resonances, '
                   case (1)
                      write (u, '(1x,A)', advance='no') '1 resonance,  '
                   case default
                      write (u, '(1x,I0,1x,A)', advance='no') &
                           cascade%n_resonances, 'resonances, '
                   end select
                   write (u, '(1x,I0,1x,A)', advance='no') &
                        cascade%n_log_enhanced, 'logs, '
                   write (u, '(1x,I0,1x,A)', advance='no') &
                        cascade%n_off_shell, 'off-shell, '
                   select case (cascade%n_t_channel)
                   case (0);  write (u, '(1x,A)') 's-channel graph'
                   case (1);  write (u, '(1x,A)') '1 t-channel line'
                   case default
                      write(u,'(1x,I0,1x,A)') &
                           cascade%n_t_channel, 't-channel lines'
                   end select
                   write (u, '(1x,A,I0)') 'grove #', grove
                end if
                count = count + 1
                write (u, "(1x,'!',1x,A,I0)")  "Channel #", count
                call cascade_write_file_format (cascade, cascade_set%model, u)
             end if
          end if
          cascade => cascade%next
       end do
    end do
  end subroutine cascade_set_write_file_format

\eatline
\nwendcode{}\nwbegindocs{688}This is the graph output format, the driver-file
\nwenddocs{}\nwbegincode{689}\moddef{Cascades: public}\plusendmoddef
  public :: cascade_set_write_graph_format
\nwendcode{}\nwbegincode{690}\moddef{Cascades: procedures}\plusendmoddef
  subroutine cascade_set_write_graph_format &
      (cascade_set, filename, process_id, unit)
    type(cascade_set_t), intent(in), target :: cascade_set
    type(string_t), intent(in) :: filename, process_id
    integer, intent(in), optional :: unit
    type(cascade_t), pointer :: cascade
    integer :: u, grove, count, pgcount
    logical :: first_in_grove
    u = given_output_unit (unit);  if (u < 0)  return
    write (u, '(A)') "\\documentclass[10pt]\{article\}"
    write (u, '(A)') "\\usepackage\{amsmath\}"
    write (u, '(A)') "\\usepackage\{feynmp\}"
    write (u, '(A)') "\\usepackage\{url\}"
    write (u, '(A)') "\\usepackage\{color\}"
    write (u, *)
    write (u, '(A)') "\\textwidth 18.5cm"
    write (u, '(A)') "\\evensidemargin -1.5cm"
    write (u, '(A)') "\\oddsidemargin -1.5cm"
    write (u, *)
    write (u, '(A)') "\\newcommand\{\\blue\}\{\\color\{blue\}\}"
    write (u, '(A)') "\\newcommand\{\\green\}\{\\color\{green\}\}"
    write (u, '(A)') "\\newcommand\{\\red\}\{\\color\{red\}\}"
    write (u, '(A)') "\\newcommand\{\\magenta\}\{\\color\{magenta\}\}"
    write (u, '(A)') "\\newcommand\{\\cyan\}\{\\color\{cyan\}\}"
    write (u, '(A)') "\\newcommand\{\\sm\}\{\\footnotesize\}"
    write (u, '(A)') "\\setlength\{\\parindent\}\{0pt\}"
    write (u, '(A)') "\\setlength\{\\parsep\}\{20pt\}"
    write (u, *)
    write (u, '(A)') "\\begin\{document\}"
    write (u, '(A)') "\\begin\{fmffile\}\{" // char (filename) // "\}"
    write (u, '(A)') "\\fmfcmd\{color magenta; magenta = red + blue;\}"
    write (u, '(A)') "\\fmfcmd\{color cyan; cyan = green + blue;\}"
    write (u, '(A)') "\\begin\{fmfshrink\}\{0.5\}"
    write (u, '(A)') "\\begin\{flushleft\}"
    write (u, *)
    write (u, '(A)') "\\noindent" // &
         & "\\textbf\{\\large\\texttt\{WHIZARD\} phase space channels\}" // &
         & "\\hfill\\today"
    write (u, *)
    write (u, '(A)') "\\vspace\{10pt\}"
    write (u, '(A)') "\\noindent" // &
         & "\\textbf\{Process:\} \\url\{" // char (process_id) // "\}"
    call cascade_set_write_process_tex_format (cascade_set, u)
    write (u, *)
    write (u, '(A)') "\\noindent" // &
         & "\\textbf\{Note:\} These are pseudo Feynman graphs that " 
    write (u, '(A)') "visualize phase-space parameterizations " // &
         & "(``integration channels'').  "
    write (u, '(A)') "They do \\emph\{not\} indicate Feynman graphs used for the " // &
         & "matrix element."
    write (u, *)
    write (u, '(A)') "\\textbf\{Color code:\} " // &
         & "\{\\blue resonance,\} " // &
         & "\{\\cyan t-channel,\} " // &
         & "\{\\green radiation,\} "
    write (u, '(A)') "\{\\red infrared,\} " // &
         & "\{\\magenta collinear,\} " // &
         & "external/off-shell"
    write (u, *)
    write (u, '(A)') "\\noindent" // &
         & "\\textbf\{Black square:\} Keystone, indicates ordering of " // &
         & "phase space parameters."
    write (u, *)
    write (u, '(A)') "\\vspace\{-20pt\}"
    count = 0
    pgcount = 0
    do grove = 1, cascade_set%n_groves
       first_in_grove = .true.
       cascade => cascade_set%first
       do while (associated (cascade))
          if (cascade%active .and. cascade%complete) then
             if (cascade%grove == grove) then
                if (first_in_grove) then
                   first_in_grove = .false.
                   write (u, *)
                   write (u, '(A)') "\\vspace\{20pt\}"
                   write (u, '(A)') "\\begin\{tabular\}\{l\}"
                   write (u, '(A,I5,A)') &
                        & "\\fbox\{\\bf Grove \\boldmath$", grove, "$\} \\\\[10pt]"
                   write (u, '(A,I1,A)') "Multiplicity: ", &
                        cascade%multiplicity, "\\\\"
                   write (u, '(A,I1,A)') "Resonances:   ", &
                        cascade%n_resonances, "\\\\"
                   write (u, '(A,I1,A)') "Log-enhanced: ", &
                        cascade%n_log_enhanced, "\\\\"
                   write (u, '(A,I1,A)') "Off-shell:    ", &
                        cascade%n_off_shell, "\\\\"
                   write (u, '(A,I1,A)') "t-channel:    ", &
                        cascade%n_t_channel, ""
                   write (u, '(A)') "\\end\{tabular\}"
                end if
                count = count + 1
                call cascade_write_graph_format (cascade, count, unit)
                if (pgcount >= 250) then
                   write (u, '(A)') "\\clearpage"
                   pgcount = 0
                end if
             end if
          end if
          cascade => cascade%next
       end do
    end do
    write (u, '(A)') "\\end\{flushleft\}"
    write (u, '(A)') "\\end\{fmfshrink\}"
    write (u, '(A)') "\\end\{fmffile\}"
    write (u, '(A)') "\\end\{document\}"
 end subroutine cascade_set_write_graph_format

\eatline
\nwendcode{}\nwbegindocs{691}This is for screen output and debugging:
\nwenddocs{}\nwbegincode{692}\moddef{Cascades: public}\plusendmoddef
  public :: cascade_set_write
\nwendcode{}\nwbegincode{693}\moddef{Cascades: procedures}\plusendmoddef
  subroutine cascade_set_write (cascade_set, unit, active_only, complete_only)
    type(cascade_set_t), intent(in), target :: cascade_set
    integer, intent(in), optional :: unit
    logical, intent(in), optional :: active_only, complete_only
    logical :: active, complete
    type(cascade_t), pointer :: cascade
    integer :: u, i
    u = given_output_unit (unit);  if (u < 0)  return
    active = .true.;  if (present (active_only))  active = active_only
    complete = .false.;  if (present (complete_only))  complete = complete_only
    write (u, "(A)") "Cascade set:"
    write (u, "(3x,A)", advance="no")  "Model:"
    if (associated (cascade_set%model)) then
       write (u, "(1x,A)") char (cascade_set%model%get_name ())
    else
       write (u, "(1x,A)") "[none]"
    end if
    write (u, "(3x,A)", advance="no")  "n_in/out/tot  ="
    write (u, "(3(1x,I7))")  &
         cascade_set%n_in, cascade_set%n_out, cascade_set%n_tot
    write (u, "(3x,A)", advance="no")  "depth_out/tot ="
    write (u, "(2(1x,I7))")  cascade_set%depth_out, cascade_set%depth_tot
    write (u, "(3x,A)", advance="no")  "mass thr(s/t) ="
    write (u, "(2(1x," // FMT_19 // "))")  &
         cascade_set%m_threshold_s, cascade_set%m_threshold_t
    write (u, "(3x,A)", advance="no")  "off shell     ="
    write (u, "(1x,I7)")  cascade_set%off_shell
    write (u, "(3x,A)", advance="no")  "keep_nonreson ="
    write (u, "(1x,L1)")  cascade_set%keep_nonresonant
    write (u, "(3x,A)", advance="no")  "n_groves      ="
    write (u, "(1x,I7)")  cascade_set%n_groves
    write (u, "(A)")
    write (u, "(A)") "Cascade list:"
    if (associated (cascade_set%first)) then
       cascade => cascade_set%first
       do while (associated (cascade))
          if (active .and. .not. cascade%active)  cycle
          if (complete .and. .not. cascade%complete)  cycle
          call cascade_write (cascade, unit)
          cascade => cascade%next
       end do
    else
       write (u, "(A)") "[empty]"
    end if
    write (u, "(A)") "Hash array"
    write (u, "(3x,A)", advance="no")  "n_entries     ="
    write (u, "(1x,I7)")  cascade_set%n_entries
    write (u, "(3x,A)", advance="no")  "fill_ratio    ="
    write (u, "(1x," // FMT_12 // ")")  cascade_set%fill_ratio
    write (u, "(3x,A)", advance="no")  "n_entries_max ="
    write (u, "(1x,I7)")  cascade_set%n_entries_max
    write (u, "(3x,A)", advance="no")  "mask          ="
    write (u, "(1x,I0)")  cascade_set%mask
    do i = 0, ubound (cascade_set%entry, 1)
       if (allocated (cascade_set%entry(i)%key)) then
          write (u, "(1x,I7)") i
          call hash_entry_write (cascade_set%entry(i), u)
       end if
    end do
  end subroutine cascade_set_write

\eatline
\nwendcode{}\nwbegindocs{694}\nwdocspar
\subsection{Adding cascades}
Add a deep copy of a cascade to the set.  The copy has all content of the
original, but the pointers are nullified.  We do not care whether insertion
was successful or not.  The pointer argument, if present, is assigned to the
input cascade, or to the hash entry if it is already present.

The procedure is recursive: any daughter or mother entries are also
deep-copied and added to the cascade set before the current copy is added.
\nwenddocs{}\nwbegincode{695}\moddef{Cascades: procedures}\plusendmoddef
  recursive subroutine cascade_set_add_copy &
       (cascade_set, cascade_in, cascade_ptr)
    type(cascade_set_t), intent(inout), target :: cascade_set
    type(cascade_t), intent(in) :: cascade_in
    type(cascade_t), optional, pointer :: cascade_ptr
    type(cascade_t), pointer :: cascade
    logical :: ok
    allocate (cascade)
    cascade = cascade_in
    if (associated (cascade_in%daughter1))  call cascade_set_add_copy &
         (cascade_set, cascade_in%daughter1, cascade%daughter1)
    if (associated (cascade_in%daughter2))  call cascade_set_add_copy &
         (cascade_set, cascade_in%daughter2, cascade%daughter2)
    if (associated (cascade_in%mother))  call cascade_set_add_copy &
         (cascade_set, cascade_in%mother, cascade%mother)
    cascade%next => null ()
    call cascade_set_add (cascade_set, cascade, ok, cascade_ptr)
    if (.not. ok)  deallocate (cascade)
  end subroutine cascade_set_add_copy

\eatline
\nwendcode{}\nwbegindocs{696}Add a cascade to the set.  This does not deep-copy.  We first try to insert
it in the hash array.  If successful, add it to the list.  Failure indicates
that it is already present, and we drop it.

The hash key is built solely from the tree array, so neither particle
codes nor resonances count, just topology.

Technically, hash and list receive only pointers, so the cascade can
be considered as being in either of both.  We treat it as part of the
list.
\nwenddocs{}\nwbegincode{697}\moddef{Cascades: procedures}\plusendmoddef
  subroutine cascade_set_add (cascade_set, cascade, ok, cascade_ptr)
    type(cascade_set_t), intent(inout), target :: cascade_set
    type(cascade_t), intent(in), target :: cascade
    logical, intent(out) :: ok
    type(cascade_t), optional, pointer :: cascade_ptr
    integer(i8), dimension(1) :: mold
    call cascade_set_hash_insert &
         (cascade_set, transfer (cascade%tree, mold), cascade, ok, cascade_ptr)
    if (ok)  call cascade_set_list_add (cascade_set, cascade)
  end subroutine cascade_set_add

\eatline
\nwendcode{}\nwbegindocs{698}Add a new cascade to the list:
\nwenddocs{}\nwbegincode{699}\moddef{Cascades: procedures}\plusendmoddef
  subroutine cascade_set_list_add (cascade_set, cascade)
    type(cascade_set_t), intent(inout) :: cascade_set
    type(cascade_t), intent(in), target :: cascade
    if (associated (cascade_set%last)) then
       cascade_set%last%next => cascade
    else
       cascade_set%first => cascade
    end if
    cascade_set%last => cascade
  end subroutine cascade_set_list_add

\eatline
\nwendcode{}\nwbegindocs{700}Add a cascade entry to the hash array:
\nwenddocs{}\nwbegincode{701}\moddef{Cascades: procedures}\plusendmoddef
  subroutine cascade_set_hash_insert &
       (cascade_set, key, cascade, ok, cascade_ptr)
    type(cascade_set_t), intent(inout), target :: cascade_set
    integer(i8), dimension(:), intent(in) :: key
    type(cascade_t), intent(in), target :: cascade
    logical, intent(out) :: ok
    type(cascade_t), optional, pointer :: cascade_ptr
    integer(i32) :: h
    if (cascade_set%n_entries >= cascade_set%n_entries_max) &
         call cascade_set_hash_expand (cascade_set)
    h = hash (key)
    call cascade_set_hash_insert_rec &
         (cascade_set, h, h, key, cascade, ok, cascade_ptr)
  end subroutine cascade_set_hash_insert

\eatline
\nwendcode{}\nwbegindocs{702}Double the hashtable size when necesssary:
\nwenddocs{}\nwbegincode{703}\moddef{Cascades: procedures}\plusendmoddef
  subroutine cascade_set_hash_expand (cascade_set)
    type(cascade_set_t), intent(inout), target :: cascade_set
    type(hash_entry_t), dimension(:), allocatable, target :: table_tmp
    type(cascade_p), pointer :: current
    integer :: i, s
    allocate (table_tmp (0:cascade_set%mask))
    table_tmp = cascade_set%entry
    deallocate (cascade_set%entry)
    s = 2 * size (table_tmp)
    cascade_set%n_entries = 0
    cascade_set%n_entries_max = s * cascade_set%fill_ratio
    cascade_set%mask = s - 1
    allocate (cascade_set%entry (0:cascade_set%mask))
    do i = 0, ubound (table_tmp, 1)
       current => table_tmp(i)%first
       do while (associated (current))
          call cascade_set_hash_insert_rec &
               (cascade_set, table_tmp(i)%hashval, table_tmp(i)%hashval, &
                table_tmp(i)%key, current%cascade)
          current => current%next
       end do
    end do
  end subroutine cascade_set_hash_expand

\eatline
\nwendcode{}\nwbegindocs{704}Insert the cascade at the bucket determined by the hash value.  If
the bucket is filled, check first for a collision (unequal keys).  In
that case, choose the following bucket and repeat.  Otherwise, add the
cascade to the bucket.

If the bucket is empty, record the hash value, allocate and store the
key, and then add the cascade to the bucket.

If {\tt{}ok} is present, before insertion we check whether the cascade is
already stored, and return failure if yes.
\nwenddocs{}\nwbegincode{705}\moddef{Cascades: procedures}\plusendmoddef
  recursive subroutine cascade_set_hash_insert_rec &
       (cascade_set, h, hashval, key, cascade, ok, cascade_ptr)
    type(cascade_set_t), intent(inout) :: cascade_set
    integer(i32), intent(in) :: h, hashval
    integer(i8), dimension(:), intent(in) :: key
    type(cascade_t), intent(in), target :: cascade
    logical, intent(out), optional :: ok
    type(cascade_t), optional, pointer :: cascade_ptr
    integer(i32) :: i
    i = iand (h, cascade_set%mask)
    if (allocated (cascade_set%entry(i)%key)) then
       if (size (cascade_set%entry(i)%key) /= size (key)) then
          call cascade_set_hash_insert_rec &
               (cascade_set, h + 1, hashval, key, cascade, ok, cascade_ptr)
       else if (any (cascade_set%entry(i)%key /= key)) then
          call cascade_set_hash_insert_rec &
               (cascade_set, h + 1, hashval, key, cascade, ok, cascade_ptr)
       else
          call hash_entry_add_cascade_ptr &
               (cascade_set%entry(i), cascade, ok, cascade_ptr)
       end if
    else
       cascade_set%entry(i)%hashval = hashval
       allocate (cascade_set%entry(i)%key (size (key)))
       cascade_set%entry(i)%key = key
       call hash_entry_add_cascade_ptr &
            (cascade_set%entry(i), cascade, ok, cascade_ptr)
       cascade_set%n_entries = cascade_set%n_entries + 1
    end if
  end subroutine cascade_set_hash_insert_rec

\eatline
\nwendcode{}\nwbegindocs{706}\nwdocspar
\subsection{External particles}
We want to initialize the cascade set with the outgoing particles.  In
case of multiple processes, initial cascades are prepared for all of
them.  The hash array check ensures that no particle appears more than
once at the same place.
\nwenddocs{}\nwbegincode{707}\moddef{Cascades: interfaces}\plusendmoddef
  interface cascade_set_add_outgoing
     module procedure cascade_set_add_outgoing1
     module procedure cascade_set_add_outgoing2
  end interface
 
\nwendcode{}\nwbegincode{708}\moddef{Cascades: procedures}\plusendmoddef
  subroutine cascade_set_add_outgoing2 (cascade_set, flv)
    type(cascade_set_t), intent(inout), target :: cascade_set
    type(flavor_t), dimension(:,:), intent(in) :: flv
    integer :: pos, prc, n_out, n_prc
    type(cascade_t), pointer :: cascade
    logical :: ok
    n_out = size (flv, dim=1)
    n_prc = size (flv, dim=2)
    do prc = 1, n_prc
       do pos = 1, n_out
          allocate (cascade)
          call cascade_init_outgoing &
               (cascade, flv(pos,prc), pos, cascade_set%m_threshold_s)
          call cascade_set_add (cascade_set, cascade, ok)
          if (.not. ok) then
             deallocate (cascade)
          end if
       end do
    end do
  end subroutine cascade_set_add_outgoing2

  subroutine cascade_set_add_outgoing1 (cascade_set, flv)
    type(cascade_set_t), intent(inout), target :: cascade_set
    type(flavor_t), dimension(:), intent(in) :: flv
    integer :: pos, n_out
    type(cascade_t), pointer :: cascade
    logical :: ok
    n_out = size (flv, dim=1)
    do pos = 1, n_out
       allocate (cascade)
       call cascade_init_outgoing &
            (cascade, flv(pos), pos, cascade_set%m_threshold_s)
       call cascade_set_add (cascade_set, cascade, ok)
       if (.not. ok) then
          deallocate (cascade)
       end if
    end do
  end subroutine cascade_set_add_outgoing1

\eatline
\nwendcode{}\nwbegindocs{709}The incoming particles are added one at a time.  Nevertheless, we
may have several processes which are looped over.  At the first
opportunity, we set the pointer {\tt{}first{\char95}t} in the cascade set which
should point to the first t-channel cascade.

Return the indices of the first and last cascade generated.
\nwenddocs{}\nwbegincode{710}\moddef{Cascades: interfaces}\plusendmoddef
  interface cascade_set_add_incoming
     module procedure cascade_set_add_incoming0
     module procedure cascade_set_add_incoming1
  end interface
 
\nwendcode{}\nwbegincode{711}\moddef{Cascades: procedures}\plusendmoddef
  subroutine cascade_set_add_incoming1 (cascade_set, n1, n2, pos, flv)
    type(cascade_set_t), intent(inout), target :: cascade_set
    integer, intent(out) :: n1, n2
    integer, intent(in) :: pos
    type(flavor_t), dimension(:), intent(in) :: flv
    integer :: prc, n_prc
    type(cascade_t), pointer :: cascade
    logical :: ok
    n1 = 0
    n2 = 0
    n_prc = size (flv)
    do prc = 1, n_prc
       allocate (cascade)
       call cascade_init_incoming &
            (cascade, flv(prc), pos, cascade_set%m_threshold_t)
       call cascade_set_add (cascade_set, cascade, ok)
       if (ok) then
          if (n1 == 0)  n1 = cascade%index
          n2 = cascade%index
          if (.not. associated (cascade_set%first_t)) then
             cascade_set%first_t => cascade
          end if
       else
          deallocate (cascade)
       end if
    end do
  end subroutine cascade_set_add_incoming1

  subroutine cascade_set_add_incoming0 (cascade_set, n1, n2, pos, flv)
    type(cascade_set_t), intent(inout), target :: cascade_set
    integer, intent(out) :: n1, n2
    integer, intent(in) :: pos
    type(flavor_t), intent(in) :: flv
    type(cascade_t), pointer :: cascade
    logical :: ok
    n1 = 0
    n2 = 0
    allocate (cascade)
    call cascade_init_incoming &
         (cascade, flv, pos, cascade_set%m_threshold_t)
    call cascade_set_add (cascade_set, cascade, ok)
    if (ok) then
       if (n1 == 0)  n1 = cascade%index
       n2 = cascade%index
       if (.not. associated (cascade_set%first_t)) then
          cascade_set%first_t => cascade
       end if
    else
       deallocate (cascade)
    end if
  end subroutine cascade_set_add_incoming0

\eatline
\nwendcode{}\nwbegindocs{712}\nwdocspar
\subsection{Cascade combination I: flavor assignment}
We have two disjunct cascades, now use the vertex table to determine
the possible flavors of the combination cascade.  For each
possibility, try to generate a new cascade.  The total cascade depth
has to be one less than the limit, because this is reached by setting
the keystone.
\nwenddocs{}\nwbegincode{713}\moddef{Cascades: procedures}\plusendmoddef
  subroutine cascade_match_pair (cascade_set, cascade1, cascade2, s_channel)
    type(cascade_set_t), intent(inout), target :: cascade_set
    type(cascade_t), intent(in), target :: cascade1, cascade2
    logical, intent(in) :: s_channel
    integer, dimension(:), allocatable :: pdg3
    integer :: i, depth_max
    type(flavor_t) :: flv
    if (s_channel) then
       depth_max = cascade_set%depth_out
    else
       depth_max = cascade_set%depth_tot
    end if
    if (cascade1%depth + cascade2%depth < depth_max) then
       call cascade_set%model%match_vertex ( &
            cascade1%flv%get_pdg (), &
            cascade2%flv%get_pdg (), &
            pdg3)
       do i = 1, size (pdg3)
          call flv%init (pdg3(i), cascade_set%model)
          if (s_channel) then
             call cascade_combine_s (cascade_set, cascade1, cascade2, flv)
          else
             call cascade_combine_t (cascade_set, cascade1, cascade2, flv)
          end if
       end do
       deallocate (pdg3)
    end if
  end subroutine cascade_match_pair

\eatline
\nwendcode{}\nwbegindocs{714}The triplet version takes a third cascade, and we check whether this
triplet has a matching vertex in the database.  If yes, we make a
keystone cascade.
\nwenddocs{}\nwbegincode{715}\moddef{Cascades: procedures}\plusendmoddef
  subroutine cascade_match_triplet &
       (cascade_set, cascade1, cascade2, cascade3, s_channel)
    type(cascade_set_t), intent(inout), target :: cascade_set
    type(cascade_t), intent(in), target :: cascade1, cascade2, cascade3
    logical, intent(in) :: s_channel
    integer :: depth_max
    depth_max = cascade_set%depth_tot
    if (cascade1%depth + cascade2%depth + cascade3%depth == depth_max) then
       if (cascade_set%model%check_vertex ( &
            cascade1%flv%get_pdg (), &
            cascade2%flv%get_pdg (), &
            cascade3%flv%get_pdg ())) then
          call cascade_combine_keystone &
               (cascade_set, cascade1, cascade2, cascade3, s_channel)
       end if
    end if
  end subroutine cascade_match_triplet

\eatline
\nwendcode{}\nwbegindocs{716}\nwdocspar
\subsection{Cascade combination II: kinematics setup and check}
Having three matching flavors, we start constructing the combination
cascade.  We look at the mass hierarchies and determine whether the
cascade is to be kept.  In passing we set mapping modes, resonance
properties and such.

If successful, the cascade is finalized.  For a resonant cascade, we
prepare in addition a copy without the resonance.
\nwenddocs{}\nwbegincode{717}\moddef{Cascades: procedures}\plusendmoddef
  subroutine cascade_combine_s (cascade_set, cascade1, cascade2, flv)
    type(cascade_set_t), intent(inout), target :: cascade_set
    type(cascade_t), intent(in), target :: cascade1, cascade2
    type(flavor_t), intent(in) :: flv
    type(cascade_t), pointer :: cascade3, cascade4
    logical :: keep
    keep = .false.
    allocate (cascade3)
    call cascade_init (cascade3, cascade1%depth + cascade2%depth + 1)
    cascade3%bincode = ior (cascade1%bincode, cascade2%bincode)
    cascade3%flv = flv%anti ()
    cascade3%pdg = abs (cascade3%flv%get_pdg ())
    cascade3%is_vector = flv%get_spin_type () == VECTOR
    cascade3%m_min = cascade1%m_min + cascade2%m_min
    cascade3%m_rea = flv%get_mass ()
    if (cascade3%m_rea > cascade_set%m_threshold_s) then
       cascade3%m_eff = cascade3%m_rea
    end if
    ! Potentially resonant cases [sqrts = m_rea for on-shell decay]
    if (cascade3%m_rea > cascade3%m_min &
         .and. cascade3%m_rea <= cascade_set%sqrts) then
       if (flv%get_width () /= 0) then
          if (cascade1%on_shell .or. cascade2%on_shell) then
             keep = .true.
             cascade3%mapping = S_CHANNEL
             cascade3%resonant = .true.
          end if
       else
          call warn_decay (flv)
       end if
    ! Collinear and IR singular cases
    else if (cascade3%m_rea < cascade_set%sqrts) then
       ! Massless splitting
       if (cascade1%m_eff == 0 .and. cascade2%m_eff == 0 &
            .and. cascade3%depth <= 3) then
          keep = .true.
          cascade3%log_enhanced = .true.
          if (cascade3%is_vector) then
             if (cascade1%is_vector .and. cascade2%is_vector) then
                cascade3%mapping = COLLINEAR   ! three-vector-vertex
             else
                cascade3%mapping = INFRARED    ! vector splitting into matter
             end if
          else
             if (cascade1%is_vector .or. cascade2%is_vector) then
                cascade3%mapping = COLLINEAR   ! vector radiation off matter
             else
                cascade3%mapping = INFRARED    ! scalar radiation/splitting
             end if
          end if
       ! IR radiation off massive particle
       else if (cascade3%m_eff > 0 .and. cascade1%m_eff > 0 &
            .and. cascade2%m_eff == 0 &
            .and. (cascade1%on_shell .or. cascade1%mapping == RADIATION) &
            .and. abs (cascade3%m_eff - cascade1%m_eff) &
                       < cascade_set%m_threshold_s) &
            then
          keep = .true.
          cascade3%log_enhanced = .true.
          cascade3%mapping = RADIATION
       else if (cascade3%m_eff > 0 .and. cascade2%m_eff > 0 &
            .and. cascade1%m_eff == 0 &
            .and. (cascade2%on_shell .or. cascade2%mapping == RADIATION) &
            .and. abs (cascade3%m_eff - cascade2%m_eff) &
                  < cascade_set%m_threshold_s) &
            then
          keep = .true.
          cascade3%log_enhanced = .true.
          cascade3%mapping = RADIATION
       end if
    end if
    ! Non-singular cases, including failed resonances
    if (.not. keep) then
       ! Two on-shell particles from a virtual mother
       if (cascade1%on_shell .or. cascade2%on_shell) then
          keep = .true.
          cascade3%m_eff = max (cascade3%m_min, &
                                cascade1%m_eff + cascade2%m_eff)
          if (cascade3%m_eff < cascade_set%m_threshold_s) then
             cascade3%m_eff = 0
          end if
       end if
    end if
    ! Complete and register the cascade (two in case of resonance)
    if (keep) then
       cascade3%on_shell = cascade3%resonant .or. cascade3%log_enhanced
       if (cascade3%resonant) then
          cascade3%pdg = abs (cascade3%flv%get_pdg ())
          if (cascade_set%keep_nonresonant) then
             allocate (cascade4)
             cascade4 = cascade3
             cascade4%index = cascade_index ()
             cascade4%pdg = UNDEFINED
             cascade4%mapping = NO_MAPPING
             cascade4%resonant = .false.
             cascade4%on_shell = .false.
          end if      
          cascade3%m_min = cascade3%m_rea
          call cascade_fusion (cascade_set, cascade1, cascade2, cascade3)
          if (cascade_set%keep_nonresonant) then
             call cascade_fusion (cascade_set, cascade1, cascade2, cascade4)
          end if
       else
          call cascade_fusion (cascade_set, cascade1, cascade2, cascade3)
       end if
    else
       deallocate (cascade3)
    end if
  contains
    subroutine warn_decay (flv)
      type(flavor_t), intent(in) :: flv
      integer :: i
      integer, dimension(MAX_WARN_RESONANCE), save :: warned_code = 0
      LOOP_WARNED: do i = 1, MAX_WARN_RESONANCE
         if (warned_code(i) == 0) then
            warned_code(i) = flv%get_pdg ()
            write (msg_buffer, "(A)") &
                 & " Intermediate decay of zero-width particle " &
                 & // char (flv%get_name ()) &
                 & // " may be possible."
            call msg_warning
            exit LOOP_WARNED
         else if (warned_code(i) == flv%get_pdg ()) then
            exit LOOP_WARNED
         end if
      end do LOOP_WARNED
    end subroutine warn_decay
  end subroutine cascade_combine_s

\eatline
\nwendcode{}\nwbegincode{718}\moddef{Cascades: parameters}\plusendmoddef
  integer, parameter, public :: MAX_WARN_RESONANCE = 50
\eatline
\nwendcode{}\nwbegindocs{719}This is the t-channel version.  {\tt{}cascade1} is t-channel and
contains the seed, {\tt{}cascade2} is s-channel.  We check for
kinematically allowed beam decay (which is a fatal error), or massless
splitting / soft radiation.  The cascade is kept in all remaining
cases and submitted for registration.
\nwenddocs{}\nwbegincode{720}\moddef{Cascades: procedures}\plusendmoddef
  subroutine cascade_combine_t (cascade_set, cascade1, cascade2, flv)
    type(cascade_set_t), intent(inout), target :: cascade_set
    type(cascade_t), intent(in), target :: cascade1, cascade2
    type(flavor_t), intent(in) :: flv
    type(cascade_t), pointer :: cascade3
    allocate (cascade3)
    call cascade_init (cascade3, cascade1%depth + cascade2%depth + 1)
    cascade3%bincode = ior (cascade1%bincode, cascade2%bincode)
    cascade3%flv = flv%anti ()
    cascade3%pdg = abs (cascade3%flv%get_pdg ())
    cascade3%is_vector = flv%get_spin_type () == VECTOR
    if (cascade1%incoming) then
       cascade3%m_min = cascade2%m_min
    else
       cascade3%m_min = cascade1%m_min + cascade2%m_min
    end if
    cascade3%m_rea = flv%get_mass ()
    if (cascade3%m_rea > cascade_set%m_threshold_t) then
       cascade3%m_eff = max (cascade3%m_rea, cascade2%m_eff)
    else if (cascade2%m_eff > cascade_set%m_threshold_t) then
       cascade3%m_eff = cascade2%m_eff
    else
       cascade3%m_eff = 0
    end if
    ! Allowed decay of beam particle
    if (cascade1%incoming &
         .and. cascade1%m_rea > cascade2%m_rea + cascade3%m_rea) then
         call beam_decay (cascade_set%fatal_beam_decay)
    ! Massless splitting
    else if (cascade1%m_eff == 0 &
         .and. cascade2%m_eff < cascade_set%m_threshold_t &
         .and. cascade3%m_eff == 0) then
       cascade3%mapping = U_CHANNEL
       cascade3%log_enhanced = .true.
    ! IR radiation off massive particle
    else if (cascade1%m_eff /= 0 .and. cascade2%m_eff == 0 &
         .and. cascade3%m_eff /= 0 &
         .and. (cascade1%on_shell .or. cascade1%mapping == RADIATION) &
         .and. abs (cascade1%m_eff - cascade3%m_eff) &
               < cascade_set%m_threshold_t) &
         then
       cascade3%pdg = abs (flv%get_pdg ())
       cascade3%log_enhanced = .true.
       cascade3%mapping = RADIATION
    end if
    cascade3%t_channel = .true.
    call cascade_fusion (cascade_set, cascade1, cascade2, cascade3)
  contains
    subroutine beam_decay (fatal_beam_decay)
      logical, intent(in) :: fatal_beam_decay
      write (msg_buffer, "(1x,A,1x,'->',1x,A,1x,A)") &
           char (cascade1%flv%get_name ()), &
           char (cascade3%flv%get_name ()), &
           char (cascade2%flv%get_name ())
      call msg_message
      write (msg_buffer, "(1x,'mass(',A,') =',1x,E17.10)") &
           char (cascade1%flv%get_name ()), cascade1%m_rea
      call msg_message
      write (msg_buffer, "(1x,'mass(',A,') =',1x,E17.10)") &
           char (cascade3%flv%get_name ()), cascade3%m_rea
      call msg_message
      write (msg_buffer, "(1x,'mass(',A,') =',1x,E17.10)") &
           char (cascade2%flv%get_name ()), cascade2%m_rea
      call msg_message
      if (fatal_beam_decay) then
         call msg_fatal (" Phase space: Initial beam particle can decay")
      else
         call msg_warning (" Phase space: Initial beam particle can decay")
      end if 
    end subroutine beam_decay
  end subroutine cascade_combine_t

\eatline
\nwendcode{}\nwbegindocs{721}Here we complete a decay cascade.  The third input is the
single-particle cascade for the initial particle.  There is no
resonance or mapping assignment.  The only condition for keeping the
cascade is the mass sum of the final state, which must be less than
the available energy.

Two modifications are necessary for scattering cascades: a pure
s-channel diagram (cascade1 is the incoming particle) do not have a
logarithmic mapping at top-level.  And in a t-channel diagram, the
last line exchanged is mapped t-channel, not u-channel.  Finally, we
can encounter the case of a $2\to 1$ process, where cascade1 is
incoming, and cascade2 is the outgoing particle.  In all three cases
we register a new cascade with the modified mapping.
\nwenddocs{}\nwbegincode{722}\moddef{Cascades: procedures}\plusendmoddef
  subroutine cascade_combine_keystone &
       (cascade_set, cascade1, cascade2, cascade3, s_channel)
    type(cascade_set_t), intent(inout), target :: cascade_set
    type(cascade_t), intent(in), target :: cascade1, cascade2, cascade3
    logical, intent(in) :: s_channel
    type(cascade_t), pointer :: cascade4, cascade0
    logical :: keep, ok
    keep = .false.
    allocate (cascade4)
    call cascade_init &
         (cascade4, cascade1%depth + cascade2%depth + cascade3%depth)
    cascade4%complete = .true.
    if (s_channel) then
       cascade4%bincode = ior (cascade1%bincode, cascade2%bincode)
    else
       cascade4%bincode = cascade3%bincode
    end if
    cascade4%flv = cascade3%flv
    cascade4%pdg = cascade3%pdg
    cascade4%mapping = EXTERNAL_PRT
    cascade4%is_vector = cascade3%is_vector
    cascade4%m_min = cascade1%m_min + cascade2%m_min
    cascade4%m_rea = cascade3%m_rea
    cascade4%m_eff = cascade3%m_rea
    if (cascade4%m_min < cascade_set%sqrts) then
       keep = .true.
    end if
    if (keep) then
       if (cascade1%incoming .and. cascade2%log_enhanced) then
          allocate (cascade0)
          cascade0 = cascade2
          cascade0%next => null ()
          cascade0%index = cascade_index ()
          cascade0%mapping = NO_MAPPING
          cascade0%log_enhanced = .false.
          cascade0%n_log_enhanced = cascade0%n_log_enhanced - 1
          cascade0%tree_mapping(cascade0%depth) = NO_MAPPING
          call cascade_keystone &
               (cascade_set, cascade1, cascade0, cascade3, cascade4, ok)
          if (ok) then
             call cascade_set_add (cascade_set, cascade0, ok)
          else
             deallocate (cascade0)
          end if
       else if (cascade1%t_channel .and. cascade1%mapping == U_CHANNEL) then
          allocate (cascade0)
          cascade0 = cascade1
          cascade0%next => null ()
          cascade0%index = cascade_index ()
          cascade0%mapping = T_CHANNEL
          cascade0%tree_mapping(cascade0%depth) = T_CHANNEL
          call cascade_keystone &
               (cascade_set, cascade0, cascade2, cascade3, cascade4, ok)
          if (ok) then
             call cascade_set_add (cascade_set, cascade0, ok)
          else
             deallocate (cascade0)
          end if
       else if (cascade1%incoming .and. cascade2%depth == 1) then
          allocate (cascade0)
          cascade0 = cascade2
          cascade0%next => null ()
          cascade0%index = cascade_index ()
          cascade0%mapping = ON_SHELL
          cascade0%tree_mapping(cascade0%depth) = ON_SHELL
          call cascade_keystone &
               (cascade_set, cascade1, cascade0, cascade3, cascade4, ok)
          if (ok) then
             call cascade_set_add (cascade_set, cascade0, ok)
          else
             deallocate (cascade0)
          end if
       else
          call cascade_keystone &
               (cascade_set, cascade1, cascade2, cascade3, cascade4, ok)
       end if
    else
       deallocate (cascade4)
    end if
  end subroutine cascade_combine_keystone

\eatline
\nwendcode{}\nwbegindocs{723}\nwdocspar
\subsection{Cascade combination III: node connections and tree fusion}
Here we assign global tree properties.  If the allowed number of
off-shell lines is exceeded, discard the new cascade.  Otherwise,
assign the trees, sort them, and assign connections.  Finally, append
the cascade to the list.  This may fail (because in the hash array
there is already an equivalent cascade).  On failure, discard the
cascade.
\nwenddocs{}\nwbegincode{724}\moddef{Cascades: procedures}\plusendmoddef
  subroutine cascade_fusion (cascade_set, cascade1, cascade2, cascade3)
    type(cascade_set_t), intent(inout), target :: cascade_set
    type(cascade_t), intent(in), target :: cascade1, cascade2
    type(cascade_t), pointer :: cascade3
    integer :: i1, i2, i3, i4
    logical :: ok
    cascade3%internal = (cascade3%depth - 3) / 2
    if (cascade3%resonant) then
       cascade3%multiplicity = 1
       cascade3%n_resonances = &
            cascade1%n_resonances + cascade2%n_resonances + 1
    else
       cascade3%multiplicity = cascade1%multiplicity + cascade2%multiplicity
       cascade3%n_resonances = cascade1%n_resonances + cascade2%n_resonances
    end if
    if (cascade3%log_enhanced) then
       cascade3%n_log_enhanced = &
            cascade1%n_log_enhanced + cascade2%n_log_enhanced + 1
    else
       cascade3%n_log_enhanced = &
            cascade1%n_log_enhanced + cascade2%n_log_enhanced
    end if
    if (cascade3%resonant) then
       cascade3%n_off_shell = 0    
    else if (cascade3%log_enhanced) then
       cascade3%n_off_shell = cascade1%n_off_shell + cascade2%n_off_shell
    else
       cascade3%n_off_shell = cascade1%n_off_shell + cascade2%n_off_shell + 1
    end if
    if (cascade3%t_channel) then
       cascade3%n_t_channel = cascade1%n_t_channel + 1
    end if
    if (cascade3%n_off_shell > cascade_set%off_shell) then
       deallocate (cascade3)
    else if (cascade3%n_t_channel > cascade_set%t_channel) then
       deallocate (cascade3)
    else
       i1 = cascade1%depth
       i2 = i1 + 1
       i3 = i1 + cascade2%depth
       i4 = cascade3%depth
       cascade3%tree(:i1) = cascade1%tree
       where (cascade1%tree_mapping > NO_MAPPING)
          cascade3%tree_pdg(:i1) = cascade1%tree_pdg
       elsewhere
          cascade3%tree_pdg(:i1) = UNDEFINED
       end where
       cascade3%tree_mapping(:i1) = cascade1%tree_mapping
       cascade3%tree_resonant(:i1) = cascade1%tree_resonant
       cascade3%tree(i2:i3) = cascade2%tree
       where (cascade2%tree_mapping > NO_MAPPING)
          cascade3%tree_pdg(i2:i3) = cascade2%tree_pdg
       elsewhere
          cascade3%tree_pdg(i2:i3) = UNDEFINED
       end where
       cascade3%tree_mapping(i2:i3) = cascade2%tree_mapping
       cascade3%tree_resonant(i2:i3) = cascade2%tree_resonant
       cascade3%tree(i4) = cascade3%bincode
       cascade3%tree_pdg(i4) = cascade3%pdg
       cascade3%tree_mapping(i4) = cascade3%mapping
       cascade3%tree_resonant(i4) = cascade3%resonant
       call tree_sort (cascade3%tree, &
            cascade3%tree_pdg, cascade3%tree_mapping, cascade3%tree_resonant)
       cascade3%has_children = .true.
       cascade3%daughter1 => cascade1
       cascade3%daughter2 => cascade2
       call cascade_set_add (cascade_set, cascade3, ok)
       if (.not. ok)  deallocate (cascade3)
    end if
  end subroutine cascade_fusion

\eatline
\nwendcode{}\nwbegindocs{725}Here we combine a cascade pair with an incoming particle, i.e., we
set a keystone.  Otherwise, this is similar.  On the first
opportunity, we set the {\tt{}first{\char95}k} pointer in the cascade set.
\nwenddocs{}\nwbegincode{726}\moddef{Cascades: procedures}\plusendmoddef
  subroutine cascade_keystone &
       (cascade_set, cascade1, cascade2, cascade3, cascade4, ok)
    type(cascade_set_t), intent(inout), target :: cascade_set
    type(cascade_t), intent(in), target :: cascade1, cascade2, cascade3
    type(cascade_t), pointer :: cascade4
    logical, intent(out) :: ok
    integer :: i1, i2, i3, i4
    cascade4%internal = (cascade4%depth - 3) / 2
    cascade4%multiplicity = cascade1%multiplicity + cascade2%multiplicity
    cascade4%n_resonances = cascade1%n_resonances + cascade2%n_resonances
    cascade4%n_off_shell = cascade1%n_off_shell + cascade2%n_off_shell
    cascade4%n_log_enhanced = &
            cascade1%n_log_enhanced + cascade2%n_log_enhanced
    cascade4%n_t_channel = cascade1%n_t_channel + cascade2%n_t_channel
    if (cascade4%n_off_shell > cascade_set%off_shell) then
       deallocate (cascade4)
       ok = .false.
    else if (cascade4%n_t_channel > cascade_set%t_channel) then
       deallocate (cascade4)
       ok = .false.
    else
       i1 = cascade1%depth
       i2 = i1 + 1
       i3 = i1 + cascade2%depth
       i4 = cascade4%depth
       cascade4%tree(:i1) = cascade1%tree
       where (cascade1%tree_mapping > NO_MAPPING)
          cascade4%tree_pdg(:i1) = cascade1%tree_pdg
       elsewhere
          cascade4%tree_pdg(:i1) = UNDEFINED
       end where
       cascade4%tree_mapping(:i1) = cascade1%tree_mapping
       cascade4%tree_resonant(:i1) = cascade1%tree_resonant
       cascade4%tree(i2:i3) = cascade2%tree
       where (cascade2%tree_mapping > NO_MAPPING)
          cascade4%tree_pdg(i2:i3) = cascade2%tree_pdg
       elsewhere
          cascade4%tree_pdg(i2:i3) = UNDEFINED
       end where
       cascade4%tree_mapping(i2:i3) = cascade2%tree_mapping
       cascade4%tree_resonant(i2:i3) = cascade2%tree_resonant
       cascade4%tree(i4) = cascade4%bincode
       cascade4%tree_pdg(i4) = UNDEFINED
       cascade4%tree_mapping(i4) = cascade4%mapping
       cascade4%tree_resonant(i4) = .false.
       call tree_sort (cascade4%tree, &
            cascade4%tree_pdg, cascade4%tree_mapping, cascade4%tree_resonant)
       cascade4%has_children = .true.
       cascade4%daughter1 => cascade1
       cascade4%daughter2 => cascade2
       cascade4%mother => cascade3
       call cascade_set_add (cascade_set, cascade4, ok)
       if (ok) then
          if (.not. associated (cascade_set%first_k)) then
             cascade_set%first_k => cascade4
          end if
       else
          deallocate (cascade4)
       end if
    end if
  end subroutine cascade_keystone

\eatline
\nwendcode{}\nwbegindocs{727}\nwdocspar
Sort a tree (array of binary codes) and particle code array
simultaneously, by ascending binary codes.  A convenient method is to
use the {\tt{}maxloc} function iteratively, to find and remove the
largest entry in the tree array one by one.
\nwenddocs{}\nwbegincode{728}\moddef{Cascades: procedures}\plusendmoddef
  subroutine tree_sort (tree, pdg, mapping, resonant)
    integer(TC), dimension(:), intent(inout) :: tree
    integer, dimension(:), intent(inout) :: pdg, mapping
    logical, dimension(:), intent(inout) :: resonant
    integer(TC), dimension(size(tree)) :: tree_tmp
    integer, dimension(size(pdg)) :: pdg_tmp, mapping_tmp
    logical, dimension(size(resonant)) :: resonant_tmp
    integer, dimension(1) :: pos
    integer :: i
    tree_tmp = tree
    pdg_tmp = pdg
    mapping_tmp = mapping
    resonant_tmp = resonant
    do i = size(tree),1,-1
       pos = maxloc (tree_tmp)
       tree(i) = tree_tmp (pos(1))
       pdg(i) = pdg_tmp (pos(1))
       mapping(i) = mapping_tmp (pos(1))
       resonant(i) = resonant_tmp (pos(1))
       tree_tmp(pos(1)) = 0
    end do
  end subroutine tree_sort

\eatline
\nwendcode{}\nwbegindocs{729}\nwdocspar
\subsection{Cascade set generation}
These procedures loop over cascades and build up the cascade set.  After each
iteration of the innermost loop, we set a breakpoint.

s-channel: We use a nested scan to combine all cascades with all other
cascades.
\nwenddocs{}\nwbegincode{730}\moddef{Cascades: procedures}\plusendmoddef
  subroutine cascade_set_generate_s (cascade_set)
    type(cascade_set_t), intent(inout), target :: cascade_set
    type(cascade_t), pointer :: cascade1, cascade2
    cascade1 => cascade_set%first
    LOOP1: do while (associated (cascade1))
       cascade2 => cascade_set%first
       LOOP2: do while (associated (cascade2))
          if (cascade2%index >= cascade1%index)  exit LOOP2
          if (cascade1 .disjunct. cascade2) then
             call cascade_match_pair (cascade_set, cascade1, cascade2, .true.)
          end if
          call terminate_now_if_signal ()
          cascade2 => cascade2%next
       end do LOOP2
       cascade1 => cascade1%next
    end do LOOP1
  end subroutine cascade_set_generate_s

\eatline
\nwendcode{}\nwbegindocs{731}The t-channel cascades are directed and have a seed (one of the
incoming particles) and a target (the other one).  We loop over all
possible seeds and targets.  Inside this, we loop over all t-channel
cascades ({\tt{}cascade1}) and s-channel cascades ({\tt{}cascade2}) and try
to combine them.
\nwenddocs{}\nwbegincode{732}\moddef{Cascades: procedures}\plusendmoddef
  subroutine cascade_set_generate_t (cascade_set, pos_seed, pos_target)
    type(cascade_set_t), intent(inout), target :: cascade_set
    integer, intent(in) :: pos_seed, pos_target
    type(cascade_t), pointer :: cascade_seed, cascade_target
    type(cascade_t), pointer :: cascade1, cascade2
    integer(TC) :: bc_seed, bc_target
    bc_seed = ibset (0_TC, pos_seed-1)
    bc_target = ibset (0_TC, pos_target-1)
    cascade_seed => cascade_set%first_t
    LOOP_SEED: do while (associated (cascade_seed))
       if (cascade_seed%bincode == bc_seed) then
          cascade_target => cascade_set%first_t
          LOOP_TARGET: do while (associated (cascade_target))
             if (cascade_target%bincode == bc_target) then
                cascade1 => cascade_set%first_t
                LOOP_T: do while (associated (cascade1))
                   if ((cascade1 .disjunct. cascade_target) &
                        .and. .not. (cascade1 .disjunct. cascade_seed)) then
                      cascade2 => cascade_set%first
                      LOOP_S: do while (associated (cascade2))
                         if ((cascade2 .disjunct. cascade_target) &
                              .and. (cascade2 .disjunct. cascade1)) then
                            call cascade_match_pair &
                                 (cascade_set, cascade1, cascade2, .false.)
                         end if
                         call terminate_now_if_signal ()
                         cascade2 => cascade2%next
                      end do LOOP_S
                   end if
                   call terminate_now_if_signal ()
                   cascade1 => cascade1%next
                end do LOOP_T
             end if
             call terminate_now_if_signal ()
             cascade_target => cascade_target%next
          end do LOOP_TARGET
       end if
       call terminate_now_if_signal ()
       cascade_seed => cascade_seed%next
    end do LOOP_SEED
  end subroutine cascade_set_generate_t
                   
\eatline
\nwendcode{}\nwbegindocs{733}This part completes the phase space for decay processes.  It is
similar to s-channel cascade generation, but combines two cascade with
the particular cascade of the incoming particle.  This particular
cascade is expected to be pointed at by {\tt{}first{\char95}t}.
\nwenddocs{}\nwbegincode{734}\moddef{Cascades: procedures}\plusendmoddef
  subroutine cascade_set_generate_decay (cascade_set)
    type(cascade_set_t), intent(inout), target :: cascade_set
    type(cascade_t), pointer :: cascade1, cascade2
    type(cascade_t), pointer :: cascade_in
    cascade_in => cascade_set%first_t
    cascade1 => cascade_set%first
    do while (associated (cascade1))
       if (cascade1 .disjunct. cascade_in) then
          cascade2 => cascade1%next
          do while (associated (cascade2))
             if ((cascade2 .disjunct. cascade1) &
                  .and. (cascade2 .disjunct. cascade_in)) then
                call cascade_match_triplet (cascade_set, &
                     cascade1, cascade2, cascade_in, .true.)
             end if
             call terminate_now_if_signal ()
             cascade2 => cascade2%next
          end do
       end if
       call terminate_now_if_signal ()
       cascade1 => cascade1%next
    end do
  end subroutine cascade_set_generate_decay

\eatline
\nwendcode{}\nwbegindocs{735}This part completes the phase space for scattering processes.  We
combine a t-channel cascade (containing the seed) with a s-channel
cascade and the target.
\nwenddocs{}\nwbegincode{736}\moddef{Cascades: procedures}\plusendmoddef
  subroutine cascade_set_generate_scattering &
       (cascade_set, ns1, ns2, nt1, nt2, pos_seed, pos_target)
    type(cascade_set_t), intent(inout), target :: cascade_set
    integer, intent(in) :: pos_seed, pos_target
    integer, intent(in) :: ns1, ns2, nt1, nt2
    type(cascade_t), pointer :: cascade_seed, cascade_target
    type(cascade_t), pointer :: cascade1, cascade2
    integer(TC) :: bc_seed, bc_target
    bc_seed = ibset (0_TC, pos_seed-1)
    bc_target = ibset (0_TC, pos_target-1)
    cascade_seed => cascade_set%first_t
    LOOP_SEED: do while (associated (cascade_seed))
       if (cascade_seed%index < ns1) then
          cascade_seed => cascade_seed%next
          cycle LOOP_SEED
       else if (cascade_seed%index > ns2) then
          exit LOOP_SEED
       else if (cascade_seed%bincode == bc_seed) then
          cascade_target => cascade_set%first_t
          LOOP_TARGET: do while (associated (cascade_target))
             if (cascade_target%index < nt1) then
                cascade_target => cascade_target%next
                cycle LOOP_TARGET
             else if (cascade_target%index > nt2) then
                exit LOOP_TARGET
             else if (cascade_target%bincode == bc_target) then
                cascade1 => cascade_set%first_t
                LOOP_T: do while (associated (cascade1))
                   if ((cascade1 .disjunct. cascade_target) &
                        .and. .not. (cascade1 .disjunct. cascade_seed)) then
                      cascade2 => cascade_set%first
                      LOOP_S: do while (associated (cascade2))
                         if ((cascade2 .disjunct. cascade_target) &
                              .and. (cascade2 .disjunct. cascade1)) then
                            call cascade_match_triplet (cascade_set, &
                                 cascade1, cascade2, cascade_target, .false.)
                         end if
                         call terminate_now_if_signal ()
                         cascade2 => cascade2%next
                      end do LOOP_S
                   end if
                   call terminate_now_if_signal ()
                   cascade1 => cascade1%next
                end do LOOP_T
             end if
             call terminate_now_if_signal ()
             cascade_target => cascade_target%next
          end do LOOP_TARGET
       end if
       call terminate_now_if_signal ()
       cascade_seed => cascade_seed%next
    end do LOOP_SEED
  end subroutine cascade_set_generate_scattering
    
\eatline
\nwendcode{}\nwbegindocs{737}\nwdocspar
\subsection{Groves}
Before assigning groves, assign hashcodes to the resonance patterns, so they
can easily be compared.
\nwenddocs{}\nwbegincode{738}\moddef{Cascades: procedures}\plusendmoddef
  subroutine cascade_set_assign_resonance_hash (cascade_set)
    type(cascade_set_t), intent(inout) :: cascade_set
    type(cascade_t), pointer :: cascade
    cascade => cascade_set%first_k
    do while (associated (cascade))
       call cascade_assign_resonance_hash (cascade)
       cascade => cascade%next
    end do
  end subroutine cascade_set_assign_resonance_hash

\eatline
\nwendcode{}\nwbegindocs{739}After all cascades are recorded, we group the complete cascades in
groves.  A grove consists of cascades with identical multiplicity,
number of resonances, log-enhanced, t-channel lines, and resonance flavors.
\nwenddocs{}\nwbegincode{740}\moddef{Cascades: procedures}\plusendmoddef
  subroutine cascade_set_assign_groves (cascade_set)
    type(cascade_set_t), intent(inout), target :: cascade_set
    type(cascade_t), pointer :: cascade1, cascade2
    integer :: multiplicity
    integer :: n_resonances, n_log_enhanced, n_t_channel, n_off_shell
    integer :: res_hash
    integer :: grove
    grove = 0
    cascade1 => cascade_set%first_k
    do while (associated (cascade1))
       if (cascade1%active .and. cascade1%complete &
            .and. cascade1%grove == 0) then
          grove = grove + 1
          cascade1%grove = grove
          multiplicity = cascade1%multiplicity
          n_resonances = cascade1%n_resonances
          n_log_enhanced = cascade1%n_log_enhanced
          n_off_shell = cascade1%n_off_shell
          n_t_channel = cascade1%n_t_channel
          res_hash = cascade1%res_hash
          cascade2 => cascade1%next
          do while (associated (cascade2))
             if (cascade2%grove == 0) then
                if (cascade2%multiplicity == multiplicity &
                     .and. cascade2%n_resonances == n_resonances &
                     .and. cascade2%n_log_enhanced == n_log_enhanced &
                     .and. cascade2%n_off_shell == n_off_shell &
                     .and. cascade2%n_t_channel == n_t_channel &
                     .and. cascade2%res_hash == res_hash) then
                   cascade2%grove = grove
                end if
             end if
             call terminate_now_if_signal ()
             cascade2 => cascade2%next
          end do
       end if
       call terminate_now_if_signal ()
       cascade1 => cascade1%next
    end do
    cascade_set%n_groves = grove
  end subroutine cascade_set_assign_groves

\eatline
\nwendcode{}\nwbegindocs{741}\nwdocspar
\subsection{Generate the phase space file}
Generate a complete phase space configuration.  

For each flavor assignment: First, all s-channel
graphs that can be built up from the outgoing particles.  Then we
distinguish (1) decay, where we complete the s-channel graphs by
connecting to the input line, and (2) scattering, where we now
generate t-channel graphs by introducing an incoming particle, and
complete this by connecting to the other incoming particle.

After all cascade sets have been generated, merge them into a common set.
This eliminates redunancies between flavor assignments.
\nwenddocs{}\nwbegincode{742}\moddef{Cascades: public}\plusendmoddef
  public :: cascade_set_generate
\nwendcode{}\nwbegincode{743}\moddef{Cascades: procedures}\plusendmoddef
  subroutine cascade_set_generate &
       (cascade_set, model, n_in, n_out, flv, phs_par, fatal_beam_decay)
    type(cascade_set_t), intent(out) :: cascade_set
    class(model_data_t), intent(in), target :: model
    integer, intent(in) :: n_in, n_out
    type(flavor_t), dimension(:,:), intent(in) :: flv
    type(phs_parameters_t), intent(in) :: phs_par
    logical, intent(in) :: fatal_beam_decay
    type(cascade_set_t), dimension(:), allocatable :: cset
    type(cascade_t), pointer :: cascade
    integer :: i
    if (phase_space_vanishes (phs_par%sqrts, n_in, flv))  return
    call cascade_set_init (cascade_set, model, n_in, n_out, phs_par, &
       fatal_beam_decay, flv)
    allocate (cset (size (flv, 2)))
    do i = 1, size (cset)
       call cascade_set_generate_single (cset(i), &
            model, n_in, n_out, flv(:,i), phs_par, fatal_beam_decay)
       cascade => cset(i)%first_k
       do while (associated (cascade))
          if (cascade%active .and. cascade%complete) then
             call cascade_set_add_copy (cascade_set, cascade)
          end if
          cascade => cascade%next
       end do
       call cascade_set_final (cset(i))
    end do           
    cascade_set%first_k => cascade_set%first
    call cascade_set_assign_resonance_hash (cascade_set)
    call cascade_set_assign_groves (cascade_set)
  end subroutine cascade_set_generate

\eatline
\nwendcode{}\nwbegindocs{744}This generates phase space for a single channel, without assigning groves.
\nwenddocs{}\nwbegincode{745}\moddef{Cascades: procedures}\plusendmoddef
  subroutine cascade_set_generate_single (cascade_set, &
      model, n_in, n_out, flv, phs_par, fatal_beam_decay)
    type(cascade_set_t), intent(out) :: cascade_set
    class(model_data_t), intent(in), target :: model
    integer, intent(in) :: n_in, n_out
    type(flavor_t), dimension(:), intent(in) :: flv
    type(phs_parameters_t), intent(in) :: phs_par
    logical, intent(in) :: fatal_beam_decay
    integer :: n11, n12, n21, n22
    call cascade_set_init (cascade_set, model, n_in, n_out, phs_par, &
       fatal_beam_decay)
    call cascade_set_add_outgoing (cascade_set, flv(n_in+1:))
    call cascade_set_generate_s (cascade_set)
    select case (n_in)
    case(1)
       call cascade_set_add_incoming &
            (cascade_set, n11, n12, n_out + 1, flv(1))
       call cascade_set_generate_decay (cascade_set)
    case(2)
       call cascade_set_add_incoming &
            (cascade_set, n11, n12, n_out + 1, flv(2))
       call cascade_set_add_incoming &
            (cascade_set, n21, n22, n_out + 2, flv(1))
       call cascade_set_generate_t (cascade_set, n_out + 1, n_out + 2)
       call cascade_set_generate_t (cascade_set, n_out + 2, n_out + 1)
       call cascade_set_generate_scattering &
            (cascade_set, n11, n12, n21, n22, n_out + 1, n_out + 2)
       call cascade_set_generate_scattering &
            (cascade_set, n21, n22, n11, n12, n_out + 2, n_out + 1)
    end select
  end subroutine cascade_set_generate_single

\eatline
\nwendcode{}\nwbegindocs{746}Sanity check: Before anything else is done, check if there could
possibly be any phase space.
\nwenddocs{}\nwbegincode{747}\moddef{Cascades: procedures}\plusendmoddef
  function phase_space_vanishes (sqrts, n_in, flv) result (flag)
    logical :: flag
    real(default), intent(in) :: sqrts
    integer, intent(in) :: n_in
    type(flavor_t), dimension(:,:), intent(in) :: flv
    real(default), dimension(:,:), allocatable :: mass
    real(default), dimension(:), allocatable :: mass_in, mass_out
    integer :: n_prt, n_flv
    flag = .false.
    if (sqrts <= 0) then
       call msg_error ("Phase space vanishes (sqrts must be positive)")
       flag = .true.;  return
    end if
    n_prt = size (flv, 1)
    n_flv = size (flv, 2)
    allocate (mass (n_prt, n_flv), mass_in (n_flv), mass_out (n_flv))
    mass = flv%get_mass ()
    mass_in = sum (mass(:n_in,:), 1)
    mass_out = sum (mass(n_in+1:,:), 1)
    if (any (mass_in > sqrts)) then
       call msg_error ("Mass sum of incoming particles " &
            // "is more than available energy")
       flag = .true.;  return
    end if
    if (any (mass_out > sqrts)) then
       call msg_error ("Mass sum of outgoing particles " &
            // "is more than available energy")
       flag = .true.;  return
    end if
  end function phase_space_vanishes
       
\eatline
\nwendcode{}\nwbegindocs{748}\nwdocspar
\subsection{Unit tests}
\nwenddocs{}\nwbegincode{749}\moddef{Cascades: public}\plusendmoddef
  public :: cascade_test
\nwendcode{}\nwbegincode{750}\moddef{Cascades: procedures}\plusendmoddef
  subroutine cascade_test (u, results)
    integer, intent(in) :: u
    type(test_results_t), intent(inout) :: results
  \LA{}Cascades: execute tests\RA{}
  end subroutine cascade_test

\eatline
\nwendcode{}\nwbegindocs{751}Checking the basic setup up of the phase space cascade parameterizations. 
\nwenddocs{}\nwbegincode{752}\moddef{Cascades: execute tests}\endmoddef
  call test (cascade_1, "cascade_1", &
       "check cascade setup", &
       u, results)
\nwendcode{}\nwbegincode{753}\moddef{Cascades: tests}\endmoddef
  subroutine cascade_1 (u)  
    integer, intent(in) :: u
    type(model_data_t), target :: model
    type(flavor_t), dimension(5,2) :: flv
    type(cascade_set_t) :: cascade_set
    type(phs_parameters_t) :: phs_par

    write (u, "(A)")  "* Test output: Cascades"
    write (u, "(A)")  "*   Purpose: test cascade phase space functions"
    write (u, "(A)")  
    
    write (u, "(A)")  "* Initializing"
    write (u, "(A)")    
    
    call model%init_sm_test ()

    call flv(1,1)%init ( 2, model)
    call flv(2,1)%init (-2, model)
    call flv(3,1)%init ( 1, model)
    call flv(4,1)%init (-1, model)
    call flv(5,1)%init (21, model)
    call flv(1,2)%init ( 2, model)
    call flv(2,2)%init (-2, model)
    call flv(3,2)%init ( 2, model)
    call flv(4,2)%init (-2, model)
    call flv(5,2)%init (21, model)
    phs_par%sqrts = 1000._default
    phs_par%off_shell = 2
    
    write (u, "(A)")
    write (u, "(A)")  "* Generating the cascades"
    write (u, "(A)")
    
    call cascade_set_generate (cascade_set, model, 2, 3, flv, phs_par,.true.)
    call cascade_set_write (cascade_set, u)
    call cascade_set_write_file_format (cascade_set, u)

    write (u, "(A)")  "* Cleanup"
    write (u, "(A)")
    
    call cascade_set_final (cascade_set)
    call model%final ()
    
    write (u, *)
    write (u, "(A)")  "* Test output end: cascade_1"
        
  end subroutine cascade_1

\eatline
\nwendcode{}\nwbegindocs{754}\nwdocspar
%%%%%%%%%%%%%%%%%%%%%%%%%%%%%%%%%%%%%%%%%%%%%%%%%%%%%%%%%%%%%%%%%%%%%%%%
\section{WOOD phase space}

This is the module that interfaces the {\tt{}phs{\char95}forests} phase-space
treatment and the {\tt{}cascades} module for generating phase-space
channels.  As an extension of the {\tt{}phs{\char95}base} abstract type,
the phase-space configuration and instance implement the standard API.

(Currently, this is the only generic phase-space implementation of
\whizard.  For trivial two-particle phase space, there is
{\tt{}phs{\char95}wood} as an alternative.)
\nwenddocs{}\nwbegincode{755}\moddef{\code{}phs{\char95}wood.f90\edoc{}}\endmoddef
\LA{}File header\RA{}

module phs_wood

\LA{}Use kinds\RA{}
\LA{}Use strings\RA{}
  use io_units
  use constants
  use unit_tests
  use diagnostics
  use os_interface
  use md5
  use lorentz
  use model_data
  use flavors
  use process_constants
  use sf_mappings
  use sf_base
  use phs_base
  use mappings
  use phs_forests
  use cascades

\LA{}Standard module head\RA{}

\LA{}PHS wood: public\RA{}

\LA{}PHS wood: types\RA{}

contains

\LA{}PHS wood: procedures\RA{}

\LA{}PHS wood: tests\RA{}

end module phs_wood
\eatline
\nwendcode{}\nwbegindocs{756}\nwdocspar
\subsection{Configuration}
\nwenddocs{}\nwbegincode{757}\moddef{PHS wood: public}\endmoddef
  public :: phs_wood_config_t
\nwendcode{}\nwbegincode{758}\moddef{PHS wood: types}\endmoddef
  type, extends (phs_config_t) :: phs_wood_config_t
     character(32) :: md5sum_forest = ""
     integer :: io_unit = 0
     logical :: io_unit_keep_open = .false.
     logical :: use_equivalences = .false.
     logical :: fatal_beam_decay = .true.
     type(mapping_defaults_t) :: mapping_defaults
     type(phs_parameters_t) :: par
     type(string_t) :: run_id 
     type(cascade_set_t), allocatable :: cascade_set
     type(phs_forest_t) :: forest
     type(os_data_t) :: os_data
     logical :: extended_phs = .false.
   contains
   \LA{}PHS wood: phs wood config: TBP\RA{}
  end type phs_wood_config_t
  
\eatline
\nwendcode{}\nwbegindocs{759}Finalizer.  We should delete the cascade set and the forest subobject.

Also close the I/O unit, just in case.  (We assume that {\tt{}io{\char95}unit} is
not standard input/output.)
\nwenddocs{}\nwbegincode{760}\moddef{PHS wood: phs wood config: TBP}\endmoddef
  procedure :: final => phs_wood_config_final
\nwendcode{}\nwbegincode{761}\moddef{PHS wood: procedures}\endmoddef
  subroutine phs_wood_config_final (object)
    class(phs_wood_config_t), intent(inout) :: object
    logical :: opened
    if (object%io_unit /= 0) then
       inquire (unit = object%io_unit, opened = opened)
       if (opened)  close (object%io_unit)
    end if
    call object%clear_phase_space ()
    call phs_forest_final (object%forest)
  end subroutine phs_wood_config_final
  
\eatline
\nwendcode{}\nwbegindocs{762}\nwdocspar
\nwenddocs{}\nwbegincode{763}\moddef{PHS wood: phs wood config: TBP}\plusendmoddef
  procedure :: increase_n_par => phs_wood_config_increase_n_par
\nwendcode{}\nwbegincode{764}\moddef{PHS wood: procedures}\plusendmoddef
  subroutine phs_wood_config_increase_n_par (phs_config)
    class(phs_wood_config_t), intent(inout) :: phs_config
    phs_config%n_par = phs_config%n_par + 3
  end subroutine phs_wood_config_increase_n_par

\eatline
\nwendcode{}\nwbegindocs{765}\nwdocspar
\nwenddocs{}\nwbegincode{766}\moddef{PHS wood: phs wood config: TBP}\plusendmoddef
  procedure :: set_extended_phs => phs_wood_config_set_extended_phs
\nwendcode{}\nwbegincode{767}\moddef{PHS wood: procedures}\plusendmoddef
  subroutine phs_wood_config_set_extended_phs (phs_config)
    class(phs_wood_config_t), intent(inout) :: phs_config
    phs_config%extended_phs = .true.
  end subroutine phs_wood_config_set_extended_phs

\nwendcode{}\nwbegindocs{768}\nwdocspar
\nwenddocs{}\nwbegindocs{769}Output.  The contents of the PHS forest are not printed explicitly.
\nwenddocs{}\nwbegincode{770}\moddef{PHS wood: phs wood config: TBP}\plusendmoddef
  procedure :: write => phs_wood_config_write
\nwendcode{}\nwbegincode{771}\moddef{PHS wood: procedures}\plusendmoddef
  subroutine phs_wood_config_write (object, unit)
    class(phs_wood_config_t), intent(in) :: object
    integer, intent(in), optional :: unit
    integer :: u
    u = given_output_unit (unit)
    write (u, "(1x,A)") &
         "Partonic phase-space configuration (phase-space forest):"
    call object%base_write (unit)
    write (u, "(1x,A)")    "Phase-space configuration parameters:"
    call phs_parameters_write (object%par, u)
    call object%mapping_defaults%write (u)
    write (u, "(3x,A,A,A)")  "Run ID: '", char (object%run_id), "'"
  end subroutine phs_wood_config_write
  
\eatline
\nwendcode{}\nwbegindocs{772}Print the PHS forest contents.
\nwenddocs{}\nwbegincode{773}\moddef{PHS wood: phs wood config: TBP}\plusendmoddef
  procedure :: write_forest => phs_wood_config_write_forest
\nwendcode{}\nwbegincode{774}\moddef{PHS wood: procedures}\plusendmoddef
  subroutine phs_wood_config_write_forest (object, unit)
    class(phs_wood_config_t), intent(in) :: object
    integer, intent(in), optional :: unit
    integer :: u
    u = given_output_unit (unit)
    call phs_forest_write (object%forest, u)
  end subroutine phs_wood_config_write_forest
  
\eatline
\nwendcode{}\nwbegindocs{775}Set the phase-space parameters that the configuration generator requests.
\nwenddocs{}\nwbegincode{776}\moddef{PHS wood: phs wood config: TBP}\plusendmoddef
  procedure :: set_parameters => phs_wood_config_set_parameters
\nwendcode{}\nwbegincode{777}\moddef{PHS wood: procedures}\plusendmoddef
  subroutine phs_wood_config_set_parameters (phs_config, par)
    class(phs_wood_config_t), intent(inout) :: phs_config
    type(phs_parameters_t), intent(in) :: par
    phs_config%par = par
  end subroutine phs_wood_config_set_parameters

\eatline
\nwendcode{}\nwbegindocs{778}Enable the generation of channel equivalences (when calling {\tt{}configure}).
\nwenddocs{}\nwbegincode{779}\moddef{PHS wood: phs wood config: TBP}\plusendmoddef
  procedure :: enable_equivalences => phs_wood_config_enable_equivalences
\nwendcode{}\nwbegincode{780}\moddef{PHS wood: procedures}\plusendmoddef
  subroutine phs_wood_config_enable_equivalences (phs_config)
    class(phs_wood_config_t), intent(inout) :: phs_config
    phs_config%use_equivalences = .true.
  end subroutine phs_wood_config_enable_equivalences
  
\eatline
\nwendcode{}\nwbegindocs{781}Set the phase-space mapping parameters that the configuration generator
requests.g
\nwenddocs{}\nwbegincode{782}\moddef{PHS wood: phs wood config: TBP}\plusendmoddef
  procedure :: set_mapping_defaults => phs_wood_config_set_mapping_defaults
\nwendcode{}\nwbegincode{783}\moddef{PHS wood: procedures}\plusendmoddef
  subroutine phs_wood_config_set_mapping_defaults (phs_config, mapping_defaults)
    class(phs_wood_config_t), intent(inout) :: phs_config
    type(mapping_defaults_t), intent(in) :: mapping_defaults
    phs_config%mapping_defaults = mapping_defaults
  end subroutine phs_wood_config_set_mapping_defaults

\eatline
\nwendcode{}\nwbegindocs{784}Define the input stream for the phase-space file as an open logical unit.
The unit must be connected.
\nwenddocs{}\nwbegincode{785}\moddef{PHS wood: phs wood config: TBP}\plusendmoddef
  procedure :: set_input => phs_wood_config_set_input
\nwendcode{}\nwbegincode{786}\moddef{PHS wood: procedures}\plusendmoddef
  subroutine phs_wood_config_set_input (phs_config, unit)
    class(phs_wood_config_t), intent(inout) :: phs_config
    integer, intent(in) :: unit
    phs_config%io_unit = unit
    rewind (unit)
  end subroutine phs_wood_config_set_input
  
\eatline
\nwendcode{}\nwbegindocs{787}\nwdocspar
\subsection{Phase-space generation}
This subroutine generates a phase space configuration using the
{\tt{}cascades} module.   Note that this may take time, and the
{\tt{}cascade{\char95}set} subobject may consume a large amount of memory.
\nwenddocs{}\nwbegincode{788}\moddef{PHS wood: phs wood config: TBP}\plusendmoddef
  procedure :: generate_phase_space => phs_wood_config_generate_phase_space
\nwendcode{}\nwbegincode{789}\moddef{PHS wood: procedures}\plusendmoddef
  subroutine phs_wood_config_generate_phase_space (phs_config)
    class(phs_wood_config_t), intent(inout) :: phs_config
    integer :: off_shell, extra_off_shell
    call msg_message ("Phase space: generating configuration ...")
    off_shell = phs_config%par%off_shell
    allocate (phs_config%cascade_set)
    do extra_off_shell = 0, max (phs_config%n_tot - 3, 0)
       phs_config%par%off_shell = off_shell + extra_off_shell
       call cascade_set_generate (phs_config%cascade_set, &
            phs_config%model, phs_config%n_in, phs_config%n_out, &
            phs_config%flv, &
            phs_config%par, phs_config%fatal_beam_decay)
       if (cascade_set_is_valid (phs_config%cascade_set)) then
          exit
       else
          call msg_message ("Phase space: ... failed.  &
               &Increasing phs_off_shell ...")
       end if
    end do
    if (cascade_set_is_valid (phs_config%cascade_set)) then
       call msg_message ("Phase space: ... success.")
    else
       call msg_fatal ("Phase-space: generation failed")
    end if
  end subroutine phs_wood_config_generate_phase_space
    
\eatline
\nwendcode{}\nwbegindocs{790}Using the generated phase-space configuration, write an appropriate
phase-space file to the stored (or explicitly specified) I/O unit.
\nwenddocs{}\nwbegincode{791}\moddef{PHS wood: phs wood config: TBP}\plusendmoddef
  procedure :: write_phase_space => phs_wood_config_write_phase_space
\nwendcode{}\nwbegincode{792}\moddef{PHS wood: procedures}\plusendmoddef
  subroutine phs_wood_config_write_phase_space (phs_config, &
       filename_vis, unit)
    class(phs_wood_config_t), intent(in) :: phs_config
    integer, intent(in), optional :: unit
    type(string_t), intent(in), optional :: filename_vis
    type(string_t) :: setenv_tex, setenv_mp, pipe, pipe_dvi
    integer :: u, unit_tex, unit_dev, status   
    if (allocated (phs_config%cascade_set)) then
       if (present (unit)) then
          u = unit
       else
          u = phs_config%io_unit
       end if
       write (u, "(1x,A,A)") "process ", char (phs_config%id)
       write (u, "(A)")
       call cascade_set_write_process_bincode_format (phs_config%cascade_set, u)
       write (u, "(A)")
       write (u, "(3x,A,A,A32,A)") "md5sum_process    = ", &
            '"', phs_config%md5sum_process, '"'
       write (u, "(3x,A,A,A32,A)") "md5sum_model_par  = ", &
            '"', phs_config%md5sum_model_par, '"'
       write (u, "(3x,A,A,A32,A)") "md5sum_phs_config = ", &
            '"', phs_config%md5sum_phs_config, '"'
       call phs_parameters_write (phs_config%par, u)
       call cascade_set_write_file_format (phs_config%cascade_set, u)
       if (phs_config%vis_channels) then 
          unit_tex = free_unit ()
          open (unit=unit_tex, file=char(filename_vis // ".tex"), &
               action="write", status="replace")      
          call cascade_set_write_graph_format (phs_config%cascade_set, &
               filename_vis // "-graphs", phs_config%id, unit_tex)
          close (unit_tex)      
          call msg_message ("Phase space: visualizing channels in file " & 
               // char(trim(filename_vis)) // "...")
          if (phs_config%os_data%event_analysis_ps) then
             BLOCK: do
                unit_dev = free_unit ()
                open (file = "/dev/null", unit = unit_dev, &
                     action = "write", iostat = status)
                if (status /= 0) then
                   pipe = ""
                   pipe_dvi = ""
                else
                   pipe = " > /dev/null"
                   pipe_dvi = " 2>/dev/null 1>/dev/null"
                end if
                close (unit_dev)
                if (phs_config%os_data%whizard_texpath /= "") then
                   setenv_tex = "TEXINPUTS=" // &
                        phs_config%os_data%whizard_texpath // ":$TEXINPUTS "
                   setenv_mp = "MPINPUTS=" // &
                        phs_config%os_data%whizard_texpath // ":$MPINPUTS "
                else
                   setenv_tex = ""
                   setenv_mp = ""
                end if
                call os_system_call (setenv_tex // &
                     phs_config%os_data%latex // " " // &
                     filename_vis // ".tex " // pipe, status)
                if (status /= 0)  exit BLOCK
                if (phs_config%os_data%mpost /= "") then
                   call os_system_call (setenv_mp // &
                        phs_config%os_data%mpost // " " // &
                        filename_vis // "-graphs.mp" // pipe, status)
                else 
                   call msg_fatal ("Could not use MetaPOST.")
                end if
                if (status /= 0)  exit BLOCK
                call os_system_call (setenv_tex // &
                     phs_config%os_data%latex // " " // &
                     filename_vis // ".tex" // pipe, status)
                if (status /= 0)  exit BLOCK
                call os_system_call &
                     (phs_config%os_data%dvips // " -o " // filename_vis &
                     // ".ps " // filename_vis // ".dvi" // pipe_dvi, status)
                if (status /= 0)  exit BLOCK
                if (phs_config%os_data%event_analysis_pdf) then
                   call os_system_call (phs_config%os_data%ps2pdf // " " // &
                        filename_vis // ".ps", status)
                   if (status /= 0)  exit BLOCK
                end if
                exit BLOCK
             end do BLOCK
             if (status /= 0) then
                call msg_error ("Unable to compile analysis output file")
             end if
          end if
       end if
    else
       call msg_fatal ("Phase-space configuration: &
            &no phase space object generated")
    end if
  end subroutine phs_wood_config_write_phase_space
       
\eatline
\nwendcode{}\nwbegindocs{793}Clear the phase-space configuration.  This is useful since the
object may become \emph{really} large.
\nwenddocs{}\nwbegincode{794}\moddef{PHS wood: phs wood config: TBP}\plusendmoddef
  procedure :: clear_phase_space => phs_wood_config_clear_phase_space
\nwendcode{}\nwbegincode{795}\moddef{PHS wood: procedures}\plusendmoddef
  subroutine phs_wood_config_clear_phase_space (phs_config)
    class(phs_wood_config_t), intent(inout) :: phs_config
    if (allocated (phs_config%cascade_set)) then
       call cascade_set_final (phs_config%cascade_set)
       deallocate (phs_config%cascade_set)
    end if
  end subroutine phs_wood_config_clear_phase_space
  
\eatline
\nwendcode{}\nwbegindocs{796}\nwdocspar
\subsection{Phase-space configuration}
We read the phase-space configuration from the stored I/O unit.  If
this is not set, we assume that we have to generate a phase space
configuration.  When done, we open a scratch file and write the
configuration.

If {\tt{}rebuild} is set, we should trash any existing phase space file
and build a new one.  Otherwise, we try to use an old one, which we
check for existence and integrity.  If {\tt{}ignore{\char95}mismatch} is set, we
reuse an existing file even if it does not match the current setup.
\nwenddocs{}\nwbegincode{797}\moddef{PHS wood: phs wood config: TBP}\plusendmoddef
  procedure :: configure => phs_wood_config_configure
\nwendcode{}\nwbegincode{798}\moddef{PHS wood: procedures}\plusendmoddef
  subroutine phs_wood_config_configure (phs_config, sqrts, &
       sqrts_fixed, cm_frame, azimuthal_dependence, rebuild, ignore_mismatch, &
       nlo_type)
    class(phs_wood_config_t), intent(inout) :: phs_config
    real(default), intent(in) :: sqrts
    logical, intent(in), optional :: sqrts_fixed
    logical, intent(in), optional :: cm_frame
    logical, intent(in), optional :: azimuthal_dependence
    logical, intent(in), optional :: rebuild
    logical, intent(in), optional :: ignore_mismatch
    type(string_t), intent(inout), optional :: nlo_type
    type(string_t) :: filename, filename_vis
    logical :: variable_limits
    logical :: ok, exist, found, check, match, rebuild_phs
    integer :: g, c0, c1, n
    if (present (nlo_type)) then
      phs_config%nlo_type = nlo_type
    else
      phs_config%nlo_type = 'Born'
    end if
    phs_config%sqrts = sqrts
    phs_config%par%sqrts = sqrts
    if (present (sqrts_fixed)) &
         phs_config%sqrts_fixed = sqrts_fixed
    if (present (cm_frame)) &
         phs_config%cm_frame = cm_frame
    if (present (azimuthal_dependence)) &
         phs_config%azimuthal_dependence = azimuthal_dependence
    if (present (rebuild)) then
       rebuild_phs = rebuild
    else
       rebuild_phs = .true.
    end if
    if (present (ignore_mismatch)) then
       check = .not. ignore_mismatch
       if (ignore_mismatch) &
            call msg_warning ("Reading phs file: MD5 sum check disabled")
    else
       check = .true.
    end if
    phs_config%md5sum_forest = ""
    call phs_config%compute_md5sum ()
    if (phs_config%io_unit == 0) then
       if (phs_config%run_id /= "") then
          filename = phs_config%id // "." // phs_config%run_id // ".phs"
          filename_vis = phs_config%id // "." // phs_config%run_id // "_phs"
       else
          filename = phs_config%id // ".phs"
          filename_vis = phs_config%id // "_phs"          
       end if
       if (.not. rebuild_phs) then
          if (check) then
             call phs_config%read_phs_file (exist, found, match)
             rebuild_phs = .not. (exist .and. found .and. match)
          else
             call phs_config%read_phs_file (exist, found)
             rebuild_phs = .not. (exist .and. found)
          end if
       end if
       if (rebuild_phs) then
          call phs_config%generate_phase_space ()
          phs_config%io_unit = free_unit ()
          if (phs_config%id /= "") then
             call msg_message ("Phase space: writing configuration file '" &
                  // char (filename) // "'")
             open (phs_config%io_unit, file = char (filename), &
                  status = "replace", action = "readwrite")
          else
             open (phs_config%io_unit, status = "scratch", action = "readwrite")
          end if
          call phs_config%write_phase_space (filename_vis)
          rewind (phs_config%io_unit)
       else
          call msg_message ("Phase space: keeping configuration file '" &
               // char (filename) // "'")
       end if
    end if
    if (phs_config%io_unit == 0) then
       ok = .true.
    else
      call phs_forest_read (phs_config%forest, phs_config%io_unit, &
            phs_config%id, phs_config%n_in, phs_config%n_out, &
            phs_config%model, ok)
       if (.not. phs_config%io_unit_keep_open) then
          close (phs_config%io_unit)
          phs_config%io_unit = 0
       end if
    end if
    if (ok) then
       call phs_forest_set_flavors (phs_config%forest, phs_config%flv(:,1))
       variable_limits = .not. phs_config%cm_frame
       call phs_forest_set_parameters &
            (phs_config%forest, phs_config%mapping_defaults, variable_limits)
       call phs_forest_setup_prt_combinations (phs_config%forest)
       phs_config%n_channel = phs_forest_get_n_channels (phs_config%forest)
       phs_config%n_par = phs_forest_get_n_parameters (phs_config%forest)
       allocate (phs_config%channel (phs_config%n_channel))
       if (phs_config%use_equivalences) then
          call phs_forest_set_equivalences (phs_config%forest)
          call phs_forest_get_equivalences (phs_config%forest, &
               phs_config%channel, phs_config%azimuthal_dependence)
          phs_config%provides_equivalences = .true.
       end if
       call phs_forest_set_s_mappings (phs_config%forest)
       call phs_config%record_on_shell ()
       if (phs_config%mapping_defaults%enable_s_mapping) then
          call phs_config%record_s_mappings ()
       end if
       allocate (phs_config%chain (phs_config%n_channel), source = 0)
       do g = 1, phs_forest_get_n_groves (phs_config%forest)
          call phs_forest_get_grove_bounds (phs_config%forest, g, c0, c1, n)
          phs_config%chain (c0:c1) = g
       end do
       phs_config%provides_chains = .true.
       call phs_config%compute_md5sum_forest ()
    else
       write (msg_buffer, "(A,A,A)") &
            "Phase space: process '", &
            char (phs_config%id), "' not found in configuration file"
       call msg_fatal ()
    end if
  end subroutine phs_wood_config_configure
  
\eatline
\nwendcode{}\nwbegindocs{799}\nwdocspar
\nwenddocs{}\nwbegincode{800}\moddef{PHS wood: phs wood config: TBP}\plusendmoddef
  procedure :: set_extra_parameters => phs_wood_config_set_extra_parameters
\nwendcode{}\nwbegincode{801}\moddef{PHS wood: procedures}\plusendmoddef
  subroutine phs_wood_config_set_extra_parameters (config)
    class(phs_wood_config_t), intent(inout) :: config
    config%n_par = config%n_par + 3
  end subroutine phs_wood_config_set_extra_parameters

\eatline
\nwendcode{}\nwbegindocs{802}\nwdocspar
\nwenddocs{}\nwbegincode{803}\moddef{PHS wood: phs wood config: TBP}\plusendmoddef
  procedure :: reshuffle_flavors => phs_wood_config_reshuffle_flavors
\nwendcode{}\nwbegincode{804}\moddef{PHS wood: procedures}\plusendmoddef
  subroutine phs_wood_config_reshuffle_flavors (phs_config, reshuffle, flv_extra)
    class(phs_wood_config_t), intent(inout) :: phs_config
    integer, intent(in), dimension(:), allocatable :: reshuffle
    type(flavor_t), intent(in) :: flv_extra
    call phs_forest_set_flavors (phs_config%forest, phs_config%flv(:,1), reshuffle, flv_extra)
  end subroutine phs_wood_config_reshuffle_flavors

\eatline
\nwendcode{}\nwbegindocs{805}\nwdocspar
\nwenddocs{}\nwbegincode{806}\moddef{PHS wood: phs wood config: TBP}\plusendmoddef
  procedure :: set_momentum_links => phs_wood_config_set_momentum_links
\nwendcode{}\nwbegincode{807}\moddef{PHS wood: procedures}\plusendmoddef
  subroutine phs_wood_config_set_momentum_links (phs_config, reshuffle)
    class(phs_wood_config_t), intent(inout) :: phs_config
    integer, intent(in), dimension(:), allocatable :: reshuffle
    call phs_forest_set_momentum_links (phs_config%forest, reshuffle)
  end subroutine phs_wood_config_set_momentum_links

\eatline
\nwendcode{}\nwbegindocs{808}Identify resonances which are marked by s-channel mappings for the
whole phase space and report them to the channel array.
\nwenddocs{}\nwbegincode{809}\moddef{PHS wood: phs wood config: TBP}\plusendmoddef
  procedure :: record_s_mappings => phs_wood_config_record_s_mappings
\nwendcode{}\nwbegincode{810}\moddef{PHS wood: procedures}\plusendmoddef
  subroutine phs_wood_config_record_s_mappings (phs_config)
    class(phs_wood_config_t), intent(inout) :: phs_config
    logical :: flag
    real(default) :: mass, width
    integer :: c
    do c = 1, phs_config%n_channel
       call phs_forest_get_s_mapping (phs_config%forest, c, flag, mass, width)
       if (flag) then
          if (mass == 0) then
             call msg_fatal ("Phase space: s-channel resonance " &
                  // " has zero mass")
          end if
          if (width == 0) then
             call msg_fatal ("Phase space: s-channel resonance " &
                  // " has zero width")
          end if
          call phs_config%channel(c)%set_resonant (mass, width)
       end if
    end do
  end subroutine phs_wood_config_record_s_mappings

\eatline
\nwendcode{}\nwbegindocs{811}Identify on-shell mappings for the whole phase space and report them
to the channel array.
\nwenddocs{}\nwbegincode{812}\moddef{PHS wood: phs wood config: TBP}\plusendmoddef
  procedure :: record_on_shell => phs_wood_config_record_on_shell
\nwendcode{}\nwbegincode{813}\moddef{PHS wood: procedures}\plusendmoddef
  subroutine phs_wood_config_record_on_shell (phs_config)
    class(phs_wood_config_t), intent(inout) :: phs_config
    logical :: flag
    real(default) :: mass
    integer :: c
    do c = 1, phs_config%n_channel
       call phs_forest_get_on_shell (phs_config%forest, c, flag, mass)
       if (flag) then
          call phs_config%channel(c)%set_on_shell (mass)
       end if
    end do
  end subroutine phs_wood_config_record_on_shell

\eatline
\nwendcode{}\nwbegindocs{814}The MD5 sum of the forest is computed in addition to the MD5 sum of
the configuration.  The reason is that the forest may depend on a
user-provided external file.  On the other hand, this MD5 sum encodes
all information that is relevant for further processing.  Therefore,
the {\tt{}get{\char95}md5sum} method returns this result, once it is available.
\nwenddocs{}\nwbegincode{815}\moddef{PHS wood: phs wood config: TBP}\plusendmoddef
  procedure :: compute_md5sum_forest => phs_wood_config_compute_md5sum_forest
\nwendcode{}\nwbegincode{816}\moddef{PHS wood: procedures}\plusendmoddef
  subroutine phs_wood_config_compute_md5sum_forest (phs_config)
    class(phs_wood_config_t), intent(inout) :: phs_config
    integer :: u
    u = free_unit ()
    open (u, status = "scratch", action = "readwrite")
    call phs_config%write_forest (u)
    rewind (u)
    phs_config%md5sum_forest = md5sum (u)
    close (u)
  end subroutine phs_wood_config_compute_md5sum_forest
  
\eatline
\nwendcode{}\nwbegindocs{817}Return the most relevant MD5 sum.  This overrides the method of the
base type.
\nwenddocs{}\nwbegincode{818}\moddef{PHS wood: phs wood config: TBP}\plusendmoddef
  procedure :: get_md5sum => phs_wood_config_get_md5sum
\nwendcode{}\nwbegincode{819}\moddef{PHS wood: procedures}\plusendmoddef
  function phs_wood_config_get_md5sum (phs_config) result (md5sum)
    class(phs_wood_config_t), intent(in) :: phs_config
    character(32) :: md5sum
    if (phs_config%md5sum_forest /= "") then
       md5sum = phs_config%md5sum_forest
    else
       md5sum = phs_config%md5sum_phs_config
    end if
  end function phs_wood_config_get_md5sum
 
\eatline
\nwendcode{}\nwbegindocs{820}Check whether a phase-space configuration for the current process exists.
We look for the phase-space file that should correspond to the current
process.  If we find it, we check the MD5 sums stored in the file against the
MD5 sums in the current configuration (if required).

If successful, read the PHS file.
\nwenddocs{}\nwbegincode{821}\moddef{PHS wood: phs wood config: TBP}\plusendmoddef
  procedure :: read_phs_file => phs_wood_read_phs_file
\nwendcode{}\nwbegincode{822}\moddef{PHS wood: procedures}\plusendmoddef
  subroutine phs_wood_read_phs_file (phs_config, exist, found, match)
    class(phs_wood_config_t), intent(inout) :: phs_config
    logical, intent(out) :: exist
    logical, intent(out) :: found
    logical, intent(out), optional :: match
    type(string_t) :: filename
    integer :: u
    filename = phs_config%id // ".phs"
    inquire (file = char (filename), exist = exist)
    if (exist) then
       u = free_unit ()
       open (u, file = char (filename), action = "read", status = "old")
       call phs_forest_read (phs_config%forest, u, &
            phs_config%id, phs_config%n_in, phs_config%n_out, &
            phs_config%model, found, &
            phs_config%md5sum_process, &
            phs_config%md5sum_model_par, &
            phs_config%md5sum_phs_config, &
            match = match)
       close (u)
    else
       found = .false.
       if (present (match))  match = .false.
    end if
  end subroutine phs_wood_read_phs_file
    
\eatline
\nwendcode{}\nwbegindocs{823}Startup message, after configuration is complete.
\nwenddocs{}\nwbegincode{824}\moddef{PHS wood: phs wood config: TBP}\plusendmoddef
  procedure :: startup_message => phs_wood_config_startup_message
\nwendcode{}\nwbegincode{825}\moddef{PHS wood: procedures}\plusendmoddef
  subroutine phs_wood_config_startup_message (phs_config, unit)
    class(phs_wood_config_t), intent(in) :: phs_config
    integer, intent(in), optional :: unit
    integer :: n_groves, n_eq
    n_groves = phs_forest_get_n_groves (phs_config%forest)
    n_eq = phs_forest_get_n_equivalences (phs_config%forest)
    call phs_config%base_startup_message (unit)
    if (phs_config%n_channel == 1) then
       write (msg_buffer, "(A,2(I0,A))") &
            "Phase space: found ", phs_config%n_channel, &
            " channel, collected in ", n_groves, &
            " grove."
    else if (n_groves == 1) then
       write (msg_buffer, "(A,2(I0,A))") &
            "Phase space: found ", phs_config%n_channel, &
            " channels, collected in ", n_groves, &
            " grove." 
       else
       write (msg_buffer, "(A,2(I0,A))") &
            "Phase space: found ", phs_config%n_channel, &
            " channels, collected in ", &
            phs_forest_get_n_groves (phs_config%forest), &
            " groves."
    end if
    call msg_message (unit = unit)    
    if (phs_config%use_equivalences) then
       if (n_eq == 1) then 
          write (msg_buffer, "(A,I0,A)") &
               "Phase space: Using ", n_eq, &
               " equivalence between channels."
       else
          write (msg_buffer, "(A,I0,A)") &
               "Phase space: Using ", n_eq, &
               " equivalences between channels."          
       end if
    else
       write (msg_buffer, "(A)") &
            "Phase space: no equivalences between channels used."
    end if
    call msg_message (unit = unit)
    write (msg_buffer, "(A,2(1x,I0,1x,A))") &
         "Phase space: wood"
    call msg_message (unit = unit)
  end subroutine phs_wood_config_startup_message
    
\eatline
\nwendcode{}\nwbegindocs{826}Allocate an instance: the actual phase-space object.
\nwenddocs{}\nwbegincode{827}\moddef{PHS wood: phs wood config: TBP}\plusendmoddef
  procedure, nopass :: allocate_instance => phs_wood_config_allocate_instance
\nwendcode{}\nwbegincode{828}\moddef{PHS wood: procedures}\plusendmoddef
  subroutine phs_wood_config_allocate_instance (phs)
    class(phs_t), intent(inout), pointer :: phs
    allocate (phs_wood_t :: phs)
  end subroutine phs_wood_config_allocate_instance
  
\eatline
\nwendcode{}\nwbegindocs{829}\nwdocspar
\subsection{Kinematics implementation}
We generate $\cos\theta$ and $\phi$ uniformly, covering the solid angle.
\nwenddocs{}\nwbegincode{830}\moddef{PHS wood: public}\plusendmoddef
  public :: phs_wood_t
\nwendcode{}\nwbegincode{831}\moddef{PHS wood: types}\plusendmoddef
  type, extends (phs_t) :: phs_wood_t
     real(default) :: sqrts = 0
     type(phs_forest_t) :: forest
     real(default), dimension(3) :: r_real
     integer :: n_r_born = 0
   contains
   \LA{}PHS wood: phs wood: TBP\RA{}
  end type phs_wood_t
  
\eatline
\nwendcode{}\nwbegindocs{832}Output.  The {\tt{}verbose} setting is irrelevant, we just display the contents
of the base object.
\nwenddocs{}\nwbegincode{833}\moddef{PHS wood: phs wood: TBP}\endmoddef
  procedure :: write => phs_wood_write
\nwendcode{}\nwbegincode{834}\moddef{PHS wood: procedures}\plusendmoddef
  subroutine phs_wood_write (object, unit, verbose)
    class(phs_wood_t), intent(in) :: object
    integer, intent(in), optional :: unit
    logical, intent(in), optional :: verbose
    integer :: u
    u = given_output_unit (unit)
    call object%base_write (u)
  end subroutine phs_wood_write
    
\eatline
\nwendcode{}\nwbegindocs{835}Write the forest separately.
\nwenddocs{}\nwbegincode{836}\moddef{PHS wood: phs wood: TBP}\plusendmoddef
  procedure :: write_forest => phs_wood_write_forest
\nwendcode{}\nwbegincode{837}\moddef{PHS wood: procedures}\plusendmoddef
  subroutine phs_wood_write_forest (object, unit)
    class(phs_wood_t), intent(in) :: object
    integer, intent(in), optional :: unit
    integer :: u
    u = given_output_unit (unit)
    call phs_forest_write (object%forest, u)
  end subroutine phs_wood_write_forest
  
\eatline
\nwendcode{}\nwbegindocs{838}Finalizer.
\nwenddocs{}\nwbegincode{839}\moddef{PHS wood: phs wood: TBP}\plusendmoddef
  procedure :: final => phs_wood_final
\nwendcode{}\nwbegincode{840}\moddef{PHS wood: procedures}\plusendmoddef
  subroutine phs_wood_final (object)
    class(phs_wood_t), intent(inout) :: object
    call phs_forest_final (object%forest)
  end subroutine phs_wood_final
  
\eatline
\nwendcode{}\nwbegindocs{841}Initialization.  We allocate arrays ({\tt{}base{\char95}init}) and adjust the
phase-space volume.  The two-particle phase space volume is
\begin{equation}
  \Phi_2 = \frac{1}{4(2\pi)^5} = 2.55294034614 \times 10^{-5}
\end{equation}
independent of the particle masses.
\nwenddocs{}\nwbegincode{842}\moddef{PHS wood: phs wood: TBP}\plusendmoddef
  procedure :: init => phs_wood_init
\nwendcode{}\nwbegincode{843}\moddef{PHS wood: procedures}\plusendmoddef
  subroutine phs_wood_init (phs, phs_config)
    class(phs_wood_t), intent(out) :: phs
    class(phs_config_t), intent(in), target :: phs_config
    call phs%base_init (phs_config)
    select type (phs_config)
    type is (phs_wood_config_t)
       phs%forest = phs_config%forest
       if (phs_config%extended_phs) &
          phs%n_r_born = phs_config%n_par - 3
    end select
  end subroutine phs_wood_init
  
\eatline
\nwendcode{}\nwbegindocs{844}\nwdocspar
\subsection{Evaluation}
We compute the outgoing momenta from the incoming momenta and
the input parameter set {\tt{}r{\char95}in} in channel {\tt{}r{\char95}in}.  We also compute the
{\tt{}r} parameters and Jacobians {\tt{}f} for all other channels.

We do \emph{not} need to a apply a transformation from/to the c.m.\ frame,
because in {\tt{}phs{\char95}base} the momenta are already boosted to the c.m.\ frame
before assigning them in the {\tt{}phs} object, and inversely boosted when
extracting them.
\nwenddocs{}\nwbegincode{845}\moddef{PHS wood: phs wood: TBP}\plusendmoddef
  procedure :: evaluate_selected_channel => phs_wood_evaluate_selected_channel
  procedure :: evaluate_other_channels => phs_wood_evaluate_other_channels
\nwendcode{}\nwbegincode{846}\moddef{PHS wood: procedures}\plusendmoddef
  subroutine phs_wood_evaluate_selected_channel (phs, c_in, r_in)
    class(phs_wood_t), intent(inout) :: phs
    integer, intent(in) :: c_in
    real(default), intent(in), dimension(:) :: r_in
    logical :: ok
    phs%q_defined = .false.
    if (phs%p_defined) then
       call phs_forest_set_prt_in (phs%forest, phs%p)
       phs%r(:,c_in) = r_in
       call phs_forest_evaluate_selected_channel (phs%forest, &
            c_in, phs%active_channel, &
            phs%sqrts_hat, phs%r, phs%f, phs%volume, ok)
       select type (config => phs%config)
       type is (phs_wood_config_t)
          if (config%extended_phs) then
             if (phs%n_r_born > 0) then
                phs%r_real = r_in (phs%n_r_born+1:phs%n_r_born+3)
             else
                call msg_fatal ("n_r_born should be larger than 0!")
             end if
          end if
       end select
       if (ok) then
          phs%q = phs_forest_get_momenta_out (phs%forest)
          phs%q_defined = .true.
       end if
    end if
  end subroutine phs_wood_evaluate_selected_channel
  
  subroutine phs_wood_evaluate_other_channels (phs, c_in)
    class(phs_wood_t), intent(inout) :: phs
    integer, intent(in) :: c_in
    integer :: c
    if (phs%q_defined) then
       call phs_forest_evaluate_other_channels (phs%forest, &
            c_in, phs%active_channel, &
            phs%sqrts_hat, phs%r, phs%f, combine=.true.)
       select type (config => phs%config)
       type is (phs_wood_config_t)
          if (config%extended_phs) then
             if (phs%n_r_born > 0) then
                do c = 1, size (phs%r, 2)
                   phs%r(phs%n_r_born+1:phs%n_r_born+3,c) = phs%r_real
                end do
             else
                phs%r_defined = .false.
             end if
          end if
       end select
       phs%r_defined = .true.
    end if
  end subroutine phs_wood_evaluate_other_channels
  
\eatline
\eatline
\nwendcode{}\nwbegindocs{847}Inverse evaluation.
\nwenddocs{}\nwbegincode{848}\moddef{PHS wood: phs wood: TBP}\plusendmoddef
  procedure :: inverse => phs_wood_inverse
\nwendcode{}\nwbegincode{849}\moddef{PHS wood: procedures}\plusendmoddef
  subroutine phs_wood_inverse (phs)
    class(phs_wood_t), intent(inout) :: phs
    if (phs%p_defined .and. phs%q_defined) then
       call phs_forest_set_prt_in (phs%forest, phs%p)
       call phs_forest_set_prt_out (phs%forest, phs%q)
       call phs_forest_recover_channel (phs%forest, &
            1, &
            phs%sqrts_hat, phs%r, phs%f, phs%volume)
       call phs_forest_evaluate_other_channels (phs%forest, &
            1, phs%active_channel, &
            phs%sqrts_hat, phs%r, phs%f, combine=.false.)
       phs%r_defined = .true.
    end if
  end subroutine phs_wood_inverse
  
\eatline
\nwendcode{}\nwbegindocs{850}\nwdocspar
\subsection{Unit tests}
\nwenddocs{}\nwbegincode{851}\moddef{PHS wood: public}\plusendmoddef
  public :: phs_wood_test
\nwendcode{}\nwbegincode{852}\moddef{PHS wood: tests}\endmoddef
  subroutine phs_wood_test (u, results)
    integer, intent(in) :: u
    type(test_results_t), intent(inout) :: results
  \LA{}PHS wood: execute tests\RA{}
  end subroutine phs_wood_test
  
\eatline
\nwendcode{}\nwbegincode{853}\moddef{PHS wood: public}\plusendmoddef
  public :: phs_wood_vis_test
\nwendcode{}\nwbegincode{854}\moddef{PHS wood: tests}\plusendmoddef
  subroutine phs_wood_vis_test (u, results)
    integer, intent(in) :: u
    type(test_results_t), intent(inout) :: results
  \LA{}PHS wood: execute vis tests\RA{}
  end subroutine phs_wood_vis_test
  
\eatline
\nwendcode{}\nwbegindocs{855}\nwdocspar
\subsubsection{Phase-space configuration data}
Construct and display a test phase-space configuration object.  Also
check the {\tt{}azimuthal{\char95}dependence} flag.

This auxiliary routine writes a phase-space configuration file to unit
{\tt{}u{\char95}phs}.
\nwenddocs{}\nwbegincode{856}\moddef{PHS wood: public}\plusendmoddef
  public :: write_test_phs_file
\nwendcode{}\nwbegincode{857}\moddef{PHS wood: tests}\plusendmoddef
  subroutine write_test_phs_file (u_phs, procname)
    integer, intent(in) :: u_phs
    type(string_t), intent(in), optional :: procname
    if (present (procname)) then
       write (u_phs, "(A,A)")  "process ", char (procname)
    else
       write (u_phs, "(A)")  "process testproc"
    end if
    write (u_phs, "(A,A)")  "   md5sum_process    = ", '""'
    write (u_phs, "(A,A)")  "   md5sum_model_par  = ", '""'
    write (u_phs, "(A,A)")  "   md5sum_phs_config = ", '""'
    write (u_phs, "(A)")  "   sqrts         = 1000"
    write (u_phs, "(A)")  "   m_threshold_s =   50"    
    write (u_phs, "(A)")  "   m_threshold_t =  100"    
    write (u_phs, "(A)")  "   off_shell = 2"
    write (u_phs, "(A)")  "   t_channel = 6"
    write (u_phs, "(A)")  "   keep_nonresonant = T"
    write (u_phs, "(A)")  "  grove #1"
    write (u_phs, "(A)")  "    tree 3"
  end subroutine write_test_phs_file

\eatline
\nwendcode{}\nwbegindocs{858}\nwdocspar
\nwenddocs{}\nwbegincode{859}\moddef{PHS wood: execute tests}\endmoddef
  call test (phs_wood_1, "phs_wood_1", &
       "phase-space configuration", &
       u, results)
\nwendcode{}\nwbegincode{860}\moddef{PHS wood: tests}\plusendmoddef
  subroutine phs_wood_1 (u)
    integer, intent(in) :: u
    type(model_data_t), target :: model
    type(process_constants_t) :: process_data
    class(phs_config_t), allocatable :: phs_data
    type(mapping_defaults_t) :: mapping_defaults
    real(default) :: sqrts
    integer :: u_phs, iostat
    character(32) :: buffer
    
    write (u, "(A)")  "* Test output: phs_wood_1"
    write (u, "(A)")  "*   Purpose: initialize and display &
         &phase-space configuration data"
    write (u, "(A)")
    
    call model%init_test ()

    call syntax_phs_forest_init ()
    
    write (u, "(A)")  "* Initialize a process"
    write (u, "(A)")

    call init_test_process_data (var_str ("phs_wood_1"), process_data)

    write (u, "(A)")  "* Create a scratch phase-space file"
    write (u, "(A)")

    u_phs = free_unit ()
    open (u_phs, status = "scratch", action = "readwrite")
    call write_test_phs_file (u_phs, var_str ("phs_wood_1"))
    rewind (u_phs)
    do
       read (u_phs, "(A)", iostat = iostat)  buffer
       if (iostat /= 0)  exit
       write (u, "(A)") trim (buffer)
    end do

    write (u, "(A)")
    write (u, "(A)")  "* Setup phase-space configuration object"
    write (u, "(A)")

    mapping_defaults%step_mapping = .false.

    allocate (phs_wood_config_t :: phs_data)
    call phs_data%init (process_data, model)
    select type (phs_data)
    type is (phs_wood_config_t)
       call phs_data%set_input (u_phs)
       call phs_data%set_mapping_defaults (mapping_defaults)
    end select

    sqrts = 1000._default
    call phs_data%configure (sqrts)
       
    call phs_data%write (u)
    write (u, "(A)")

    select type (phs_data)
    type is (phs_wood_config_t)
       call phs_data%write_forest (u)
    end select
    
    write (u, "(A)")
    write (u, "(A)")  "* Cleanup"

    close (u_phs)
    call phs_data%final ()
    call model%final ()
    
    write (u, "(A)")
    write (u, "(A)")  "* Test output end: phs_wood_1"

  end subroutine phs_wood_1

\eatline
\nwendcode{}\nwbegindocs{861}\nwdocspar
\subsubsection{Phase space evaluation}
Compute kinematics for given parameters, also invert the calculation.
\nwenddocs{}\nwbegincode{862}\moddef{PHS wood: execute tests}\plusendmoddef
  call test (phs_wood_2, "phs_wood_2", &
       "phase-space evaluation", &
       u, results)
\nwendcode{}\nwbegincode{863}\moddef{PHS wood: tests}\plusendmoddef
  subroutine phs_wood_2 (u)
    integer, intent(in) :: u
    type(model_data_t), target :: model
    type(flavor_t) :: flv
    type(process_constants_t) :: process_data
    real(default) :: sqrts, E
    class(phs_config_t), allocatable, target :: phs_data
    class(phs_t), pointer :: phs => null ()
    type(vector4_t), dimension(2) :: p, q
    integer :: u_phs
    
    write (u, "(A)")  "* Test output: phs_wood_2"
    write (u, "(A)")  "*   Purpose: test simple single-channel phase space"
    write (u, "(A)")
    
    call model%init_test ()
    call flv%init (25, model)

    write (u, "(A)")  "* Initialize a process and a matching &
         &phase-space configuration"
    write (u, "(A)")

    call init_test_process_data (var_str ("phs_wood_2"), process_data)
    u_phs = free_unit ()
    open (u_phs, status = "scratch", action = "readwrite")
    call write_test_phs_file (u_phs, var_str ("phs_wood_2"))
    rewind (u_phs)

    allocate (phs_wood_config_t :: phs_data)
    call phs_data%init (process_data, model)
    select type (phs_data)
    type is (phs_wood_config_t)
       call phs_data%set_input (u_phs)
    end select

    sqrts = 1000._default
    call phs_data%configure (sqrts)

    call phs_data%write (u)
    
    write (u, "(A)")
    write (u, "(A)")  "* Initialize the phase-space instance"
    write (u, "(A)")

    call phs_data%allocate_instance (phs)
    call phs%init (phs_data)
       
    call phs%write (u, verbose=.true.)
    
    write (u, "(A)")
    write (u, "(A)")  "* Set incoming momenta"
    write (u, "(A)")

    E = sqrts / 2
    p(1) = vector4_moving (E, sqrt (E**2 - flv%get_mass ()**2), 3)
    p(2) = vector4_moving (E,-sqrt (E**2 - flv%get_mass ()**2), 3)

    call phs%set_incoming_momenta (p)
    call phs%compute_flux ()
    call phs%write (u)
    
    write (u, "(A)")
    write (u, "(A)")  "* Compute phase-space point &
         &for x = 0.125, 0.5"
    write (u, "(A)")

    call phs%evaluate_selected_channel (1, [0.125_default, 0.5_default])
    call phs%evaluate_other_channels (1)
    call phs%write (u)
    write (u, "(A)")
    select type (phs)
    type is (phs_wood_t)
       call phs%write_forest (u)
    end select
    
    write (u, "(A)")
    write (u, "(A)")  "* Inverse kinematics"
    write (u, "(A)")

    call phs%get_outgoing_momenta (q)
    call phs%final ()
    deallocate (phs)
    
    call phs_data%allocate_instance (phs)
    call phs%init (phs_data)
       
    call phs%set_incoming_momenta (p)
    call phs%compute_flux ()
    call phs%set_outgoing_momenta (q)
    
    call phs%inverse ()
    call phs%write (u)
    write (u, "(A)")
    select type (phs)
    type is (phs_wood_t)
       call phs%write_forest (u)
    end select
    
    call phs%final ()
    deallocate (phs)
    
    close (u_phs)
    call phs_data%final ()
    call model%final ()
    
    write (u, "(A)")
    write (u, "(A)")  "* Test output end: phs_wood_2"

  end subroutine phs_wood_2

\eatline
\nwendcode{}\nwbegindocs{864}\nwdocspar
\subsubsection{Phase-space generation}
Generate phase space for a simple process.
\nwenddocs{}\nwbegincode{865}\moddef{PHS wood: execute tests}\plusendmoddef
  call test (phs_wood_3, "phs_wood_3", &
       "phase-space generation", &
       u, results)
\nwendcode{}\nwbegincode{866}\moddef{PHS wood: tests}\plusendmoddef
  subroutine phs_wood_3 (u)
    integer, intent(in) :: u
    type(model_data_t), target :: model
    type(process_constants_t) :: process_data
    type(phs_parameters_t) :: phs_par
    class(phs_config_t), allocatable :: phs_data
    integer :: iostat
    character(80) :: buffer
   
    write (u, "(A)")  "* Test output: phs_wood_3"
    write (u, "(A)")  "*   Purpose: generate a phase-space configuration"
    write (u, "(A)")
    
    call model%init_test ()

    call syntax_phs_forest_init ()
    
    write (u, "(A)")  "* Initialize a process and phase-space parameters"
    write (u, "(A)")

    call init_test_process_data (var_str ("phs_wood_3"), process_data)
    allocate (phs_wood_config_t :: phs_data)
    call phs_data%init (process_data, model)

    phs_par%sqrts = 1000
    select type (phs_data)
    type is (phs_wood_config_t)
       call phs_data%set_parameters (phs_par)
       phs_data%io_unit_keep_open = .true.
    end select

    write (u, "(A)")
    write (u, "(A)")  "* Generate a scratch phase-space file"
    write (u, "(A)")

    call phs_data%configure (phs_par%sqrts)

    select type (phs_data)
    type is (phs_wood_config_t)
       rewind (phs_data%io_unit)
       do
          read (phs_data%io_unit, "(A)", iostat = iostat)  buffer
          if (iostat /= 0)  exit
          write (u, "(A)") trim (buffer)
       end do
    end select
    
    write (u, "(A)")
    write (u, "(A)")  "* Cleanup"

    call phs_data%final ()
    call model%final ()
    
    write (u, "(A)")
    write (u, "(A)")  "* Test output end: phs_wood_3"

  end subroutine phs_wood_3

\eatline
\nwendcode{}\nwbegindocs{867}\nwdocspar
\subsubsection{Nontrivial process}
Generate phase space for a $2\to 3$ process.
\nwenddocs{}\nwbegincode{868}\moddef{PHS wood: execute tests}\plusendmoddef
  call test (phs_wood_4, "phs_wood_4", &
       "nontrivial process", &
       u, results)
\nwendcode{}\nwbegincode{869}\moddef{PHS wood: tests}\plusendmoddef
  subroutine phs_wood_4 (u)
    integer, intent(in) :: u
    type(model_data_t), target :: model
    type(process_constants_t) :: process_data
    type(phs_parameters_t) :: phs_par
    class(phs_config_t), allocatable, target :: phs_data
    integer :: iostat
    character(80) :: buffer
    class(phs_t), pointer :: phs => null ()
    real(default) :: E, pL
    type(vector4_t), dimension(2) :: p
    type(vector4_t), dimension(3) :: q
   
    write (u, "(A)")  "* Test output: phs_wood_4"
    write (u, "(A)")  "*   Purpose: generate a phase-space configuration"
    write (u, "(A)")
    
    call model%init_test ()

    call syntax_phs_forest_init ()
    
    write (u, "(A)")  "* Initialize a process and phase-space parameters"
    write (u, "(A)")

    process_data%id = "phs_wood_4"
    process_data%model_name = "Test"
    process_data%n_in = 2
    process_data%n_out = 3
    process_data%n_flv = 1
    allocate (process_data%flv_state (process_data%n_in + process_data%n_out, &
         process_data%n_flv))
    process_data%flv_state(:,1) = [25, 25, 25, 6, -6]

    allocate (phs_wood_config_t :: phs_data)
    call phs_data%init (process_data, model)

    phs_par%sqrts = 1000
    select type (phs_data)
    type is (phs_wood_config_t)
       call phs_data%set_parameters (phs_par)
       phs_data%io_unit_keep_open = .true.
    end select

    write (u, "(A)")
    write (u, "(A)")  "* Generate a scratch phase-space file"
    write (u, "(A)")

    call phs_data%configure (phs_par%sqrts)

    select type (phs_data)
    type is (phs_wood_config_t)
       rewind (phs_data%io_unit)
       do
          read (phs_data%io_unit, "(A)", iostat = iostat)  buffer
          if (iostat /= 0)  exit
          write (u, "(A)") trim (buffer)
       end do
    end select
    
    write (u, "(A)")
    write (u, "(A)")  "* Initialize the phase-space instance"
    write (u, "(A)")

    call phs_data%allocate_instance (phs)
    call phs%init (phs_data)
       
    write (u, "(A)")  "* Set incoming momenta"
    write (u, "(A)")

    select type (phs_data)
    type is (phs_wood_config_t)
       E = phs_data%sqrts / 2
       pL = sqrt (E**2 - phs_data%flv(1,1)%get_mass ()**2)
    end select
    p(1) = vector4_moving (E, pL, 3)
    p(2) = vector4_moving (E, -pL, 3)

    call phs%set_incoming_momenta (p)
    call phs%compute_flux ()
    
    write (u, "(A)")  "* Compute phase-space point &
         &for x = 0.1, 0.2, 0.3, 0.4, 0.5"
    write (u, "(A)")

    call phs%evaluate_selected_channel (1, &
         [0.1_default, 0.2_default, 0.3_default, 0.4_default, 0.5_default])
    call phs%evaluate_other_channels (1)
    call phs%write (u)
    
    write (u, "(A)")
    write (u, "(A)")  "* Inverse kinematics"
    write (u, "(A)")

    call phs%get_outgoing_momenta (q)
    call phs%final ()
    deallocate (phs)
    
    call phs_data%allocate_instance (phs)
    call phs%init (phs_data)
       
    call phs%set_incoming_momenta (p)
    call phs%compute_flux ()
    call phs%set_outgoing_momenta (q)
    
    call phs%inverse ()
    call phs%write (u)
    
    write (u, "(A)")
    write (u, "(A)")  "* Cleanup"

    call phs%final ()
    deallocate (phs)

    call phs_data%final ()
    call model%final ()
    
    write (u, "(A)")
    write (u, "(A)")  "* Test output end: phs_wood_4"

  end subroutine phs_wood_4

\eatline
\nwendcode{}\nwbegindocs{870}\nwdocspar
\subsubsection{Equivalences}
Generate phase space for a simple process, including channel equivalences.
\nwenddocs{}\nwbegincode{871}\moddef{PHS wood: execute tests}\plusendmoddef
  call test (phs_wood_5, "phs_wood_5", &
       "equivalences", &
       u, results)
\nwendcode{}\nwbegincode{872}\moddef{PHS wood: tests}\plusendmoddef
  subroutine phs_wood_5 (u)
    integer, intent(in) :: u
    type(model_data_t), target :: model
    type(process_constants_t) :: process_data
    type(phs_parameters_t) :: phs_par
    class(phs_config_t), allocatable :: phs_data
   
    write (u, "(A)")  "* Test output: phs_wood_5"
    write (u, "(A)")  "*   Purpose: generate a phase-space configuration"
    write (u, "(A)")
    
    call model%init_test ()

    call syntax_phs_forest_init ()
    
    write (u, "(A)")  "* Initialize a process and phase-space parameters"
    write (u, "(A)")

    call init_test_process_data (var_str ("phs_wood_5"), process_data)
    allocate (phs_wood_config_t :: phs_data)
    call phs_data%init (process_data, model)

    phs_par%sqrts = 1000
    select type (phs_data)
    type is (phs_wood_config_t)
       call phs_data%set_parameters (phs_par)
       call phs_data%enable_equivalences ()
    end select

    write (u, "(A)")
    write (u, "(A)")  "* Generate a scratch phase-space file"
    write (u, "(A)")

    call phs_data%configure (phs_par%sqrts)
    call phs_data%write (u)
    write (u, "(A)")

    select type (phs_data)
    type is (phs_wood_config_t)
       call phs_data%write_forest (u)
    end select

    write (u, "(A)")
    write (u, "(A)")  "* Cleanup"

    call phs_data%final ()
    call model%final ()
    
    write (u, "(A)")
    write (u, "(A)")  "* Test output end: phs_wood_5"

  end subroutine phs_wood_5

\eatline
\nwendcode{}\nwbegindocs{873}\nwdocspar
\subsubsection{MD5 sum checks}
Generate phase space for a simple process.  Repeat this with and without
parameter change.
\nwenddocs{}\nwbegincode{874}\moddef{PHS wood: execute tests}\plusendmoddef
  call test (phs_wood_6, "phs_wood_6", &
       "phase-space generation", &
       u, results)
\nwendcode{}\nwbegincode{875}\moddef{PHS wood: tests}\plusendmoddef
  subroutine phs_wood_6 (u)
    integer, intent(in) :: u
    type(model_data_t), target :: model
    type(process_constants_t) :: process_data
    type(phs_parameters_t) :: phs_par
    class(phs_config_t), allocatable :: phs_data
    logical :: exist, found, match
    integer :: u_phs
    character(*), parameter :: filename = "phs_wood_6_p.phs"
   
    write (u, "(A)")  "* Test output: phs_wood_6"
    write (u, "(A)")  "*   Purpose: generate and check  phase-space file"
    write (u, "(A)")
    
    call model%init_test ()

    call syntax_phs_forest_init ()
    
    write (u, "(A)")  "* Initialize a process and phase-space parameters"
    write (u, "(A)")

    call init_test_process_data (var_str ("phs_wood_6"), process_data)
    process_data%id = "phs_wood_6_p"
    process_data%md5sum = "1234567890abcdef1234567890abcdef"
    allocate (phs_wood_config_t :: phs_data)
    call phs_data%init (process_data, model)
    
    phs_par%sqrts = 1000
    select type (phs_data)
    type is (phs_wood_config_t)
       call phs_data%set_parameters (phs_par)
    end select

    write (u, "(A)")  "* Remove previous phs file, if any"
    write (u, "(A)")

    inquire (file = filename, exist = exist)
    if (exist) then
       u_phs = free_unit ()
       open (u_phs, file = filename, action = "write")
       close (u_phs, status = "delete")
    end if

    write (u, "(A)")  "* Check phase-space file (should fail)"
    write (u, "(A)")
    
    select type (phs_data)
    type is (phs_wood_config_t)
       call phs_data%read_phs_file (exist, found, match)
       write (u, "(1x,A,L1)")  "exist = ", exist
       write (u, "(1x,A,L1)")  "found = ", found
       write (u, "(1x,A,L1)")  "match = ", match
    end select

    write (u, "(A)")
    write (u, "(A)")  "* Generate a phase-space file"
    write (u, "(A)")

    call phs_data%configure (phs_par%sqrts)

    write (u, "(1x,A,A,A)")  "MD5 sum (process)    = '", &
         phs_data%md5sum_process, "'"
    write (u, "(1x,A,A,A)")  "MD5 sum (model par)  = '", &
         phs_data%md5sum_model_par, "'"
    write (u, "(1x,A,A,A)")  "MD5 sum (phs config) = '", &
         phs_data%md5sum_phs_config, "'"

    write (u, "(A)")
    write (u, "(A)")  "* Check MD5 sum"
    write (u, "(A)")

    call phs_data%final ()
    deallocate (phs_data)
    allocate (phs_wood_config_t :: phs_data)
    call phs_data%init (process_data, model)
    phs_par%sqrts = 1000
    select type (phs_data)
    type is (phs_wood_config_t)
       call phs_data%set_parameters (phs_par)
       phs_data%sqrts = phs_par%sqrts
       phs_data%par%sqrts = phs_par%sqrts
    end select
    call phs_data%compute_md5sum ()

    write (u, "(1x,A,A,A)")  "MD5 sum (process)    = '", &
         phs_data%md5sum_process, "'"
    write (u, "(1x,A,A,A)")  "MD5 sum (model par)  = '", &
         phs_data%md5sum_model_par, "'"
    write (u, "(1x,A,A,A)")  "MD5 sum (phs config) = '", &
         phs_data%md5sum_phs_config, "'"

    select type (phs_data)
    type is (phs_wood_config_t)
       call phs_data%read_phs_file (exist, found, match)
       write (u, "(1x,A,L1)")  "exist = ", exist
       write (u, "(1x,A,L1)")  "found = ", found
       write (u, "(1x,A,L1)")  "match = ", match
    end select

    write (u, "(A)")
    write (u, "(A)")  "* Modify sqrts and check MD5 sum"
    write (u, "(A)")

    call phs_data%final ()
    deallocate (phs_data)
    allocate (phs_wood_config_t :: phs_data)
    call phs_data%init (process_data, model)
    phs_par%sqrts = 500
    select type (phs_data)
    type is (phs_wood_config_t)
       call phs_data%set_parameters (phs_par)
       phs_data%sqrts = phs_par%sqrts
       phs_data%par%sqrts = phs_par%sqrts
    end select
    call phs_data%compute_md5sum ()

    write (u, "(1x,A,A,A)")  "MD5 sum (process)    = '", &
         phs_data%md5sum_process, "'"
    write (u, "(1x,A,A,A)")  "MD5 sum (model par)  = '", &
         phs_data%md5sum_model_par, "'"
    write (u, "(1x,A,A,A)")  "MD5 sum (phs config) = '", &
         phs_data%md5sum_phs_config, "'"

    select type (phs_data)
    type is (phs_wood_config_t)
       call phs_data%read_phs_file (exist, found, match)
       write (u, "(1x,A,L1)")  "exist = ", exist
       write (u, "(1x,A,L1)")  "found = ", found
       write (u, "(1x,A,L1)")  "match = ", match
    end select

    write (u, "(A)")
    write (u, "(A)")  "* Modify process and check MD5 sum"
    write (u, "(A)")

    call phs_data%final ()
    deallocate (phs_data)
    process_data%md5sum = "77777777777777777777777777777777"
    allocate (phs_wood_config_t :: phs_data)
    call phs_data%init (process_data, model)
    phs_par%sqrts = 1000
    select type (phs_data)
    type is (phs_wood_config_t)
       call phs_data%set_parameters (phs_par)
       phs_data%sqrts = phs_par%sqrts
       phs_data%par%sqrts = phs_par%sqrts
    end select
    call phs_data%compute_md5sum ()

    write (u, "(1x,A,A,A)")  "MD5 sum (process)    = '", &
         phs_data%md5sum_process, "'"
    write (u, "(1x,A,A,A)")  "MD5 sum (model par)  = '", &
         phs_data%md5sum_model_par, "'"
    write (u, "(1x,A,A,A)")  "MD5 sum (phs config) = '", &
         phs_data%md5sum_phs_config, "'"

    select type (phs_data)
    type is (phs_wood_config_t)
       call phs_data%read_phs_file (exist, found, match)
       write (u, "(1x,A,L1)")  "exist = ", exist
       write (u, "(1x,A,L1)")  "found = ", found
       write (u, "(1x,A,L1)")  "match = ", match
    end select

    write (u, "(A)")
    write (u, "(A)")  "* Modify phs parameter and check MD5 sum"
    write (u, "(A)")

    call phs_data%final ()
    deallocate (phs_data)
    allocate (phs_wood_config_t :: phs_data)
    process_data%md5sum = "1234567890abcdef1234567890abcdef"
    call phs_data%init (process_data, model)
    phs_par%sqrts = 1000
    phs_par%off_shell = 17
    select type (phs_data)
    type is (phs_wood_config_t)
       call phs_data%set_parameters (phs_par)
       phs_data%sqrts = phs_par%sqrts
       phs_data%par%sqrts = phs_par%sqrts
    end select
    call phs_data%compute_md5sum ()

    write (u, "(1x,A,A,A)")  "MD5 sum (process)    = '", &
         phs_data%md5sum_process, "'"
    write (u, "(1x,A,A,A)")  "MD5 sum (model par)  = '", &
         phs_data%md5sum_model_par, "'"
    write (u, "(1x,A,A,A)")  "MD5 sum (phs config) = '", &
         phs_data%md5sum_phs_config, "'"

    select type (phs_data)
    type is (phs_wood_config_t)
       call phs_data%read_phs_file (exist, found, match)
       write (u, "(1x,A,L1)")  "exist = ", exist
       write (u, "(1x,A,L1)")  "found = ", found
       write (u, "(1x,A,L1)")  "match = ", match
    end select

    write (u, "(A)")
    write (u, "(A)")  "* Modify model parameter and check MD5 sum"
    write (u, "(A)")

    call phs_data%final ()
    deallocate (phs_data)
    allocate (phs_wood_config_t :: phs_data)
    call model%set_par (var_str ("ms"), 100._default)
    call phs_data%init (process_data, model)
    phs_par%sqrts = 1000
    phs_par%off_shell = 1
    select type (phs_data)
    type is (phs_wood_config_t)
       call phs_data%set_parameters (phs_par)
       phs_data%sqrts = phs_par%sqrts
       phs_data%par%sqrts = phs_par%sqrts
    end select
    call phs_data%compute_md5sum ()

    write (u, "(1x,A,A,A)")  "MD5 sum (process)    = '", &
         phs_data%md5sum_process, "'"
    write (u, "(1x,A,A,A)")  "MD5 sum (model par)  = '", &
         phs_data%md5sum_model_par, "'"
    write (u, "(1x,A,A,A)")  "MD5 sum (phs config) = '", &
         phs_data%md5sum_phs_config, "'"

    select type (phs_data)
    type is (phs_wood_config_t)
       call phs_data%read_phs_file (exist, found, match)
       write (u, "(1x,A,L1)")  "exist = ", exist
       write (u, "(1x,A,L1)")  "found = ", found
       write (u, "(1x,A,L1)")  "match = ", match
    end select

    write (u, "(A)")
    write (u, "(A)")  "* Cleanup"

    call phs_data%final ()
    call model%final ()
    
    write (u, "(A)")
    write (u, "(A)")  "* Test output end: phs_wood_6"

  end subroutine phs_wood_6

\eatline
\nwendcode{}\nwbegindocs{876}\nwdocspar
\nwenddocs{}\nwbegincode{877}\moddef{PHS wood: execute vis tests}\endmoddef
  call test (phs_wood_vis_1, "phs_wood_vis_1", &
       "visualizing phase space channels", &
       u, results)
\nwendcode{}\nwbegincode{878}\moddef{PHS wood: tests}\plusendmoddef
  subroutine phs_wood_vis_1 (u)
    integer, intent(in) :: u
    type(os_data_t) :: os_data
    type(model_data_t), target :: model
    type(process_constants_t) :: process_data
    class(phs_config_t), allocatable :: phs_data
    type(mapping_defaults_t) :: mapping_defaults
    type(string_t) :: vis_file, pdf_file, ps_file
    real(default) :: sqrts
    logical :: exist, exist_pdf, exist_ps
    integer :: u_phs, iostat, u_vis
    character(95) :: buffer
    
    write (u, "(A)")  "* Test output: phs_wood_vis_1"
    write (u, "(A)")  "*   Purpose: visualizing the &
         &phase-space configuration"
    write (u, "(A)")
    
    call os_data_init (os_data)
    call model%init_test ()

    call syntax_phs_forest_init ()
    
    write (u, "(A)")  "* Initialize a process"
    write (u, "(A)")

    call init_test_process_data (var_str ("phs_wood_vis_1"), process_data)

    write (u, "(A)")  "* Create a scratch phase-space file"
    write (u, "(A)")

    u_phs = free_unit ()
    open (u_phs, status = "scratch", action = "readwrite")
    call write_test_phs_file (u_phs, var_str ("phs_wood_vis_1"))
    rewind (u_phs)
    do
       read (u_phs, "(A)", iostat = iostat)  buffer
       if (iostat /= 0)  exit
       write (u, "(A)") trim (buffer)
    end do

    write (u, "(A)")
    write (u, "(A)")  "* Setup phase-space configuration object"
    write (u, "(A)")

    mapping_defaults%step_mapping = .false.

    allocate (phs_wood_config_t :: phs_data)
    call phs_data%init (process_data, model)
    select type (phs_data)
    type is (phs_wood_config_t)
       call phs_data%set_input (u_phs)
       call phs_data%set_mapping_defaults (mapping_defaults)
       phs_data%os_data = os_data
       phs_data%io_unit = 0
       phs_data%io_unit_keep_open = .true.
       phs_data%vis_channels = .true.
    end select

    sqrts = 1000._default
    call phs_data%configure (sqrts)
       
    call phs_data%write (u)
    write (u, "(A)")

    select type (phs_data)
    type is (phs_wood_config_t)
       call phs_data%write_forest (u)
    end select
    
    vis_file = "phs_wood_vis_1_phs.tex"
    ps_file  = "phs_wood_vis_1_phs.ps"
    pdf_file = "phs_wood_vis_1_phs.pdf"    
    inquire (file = char (vis_file), exist = exist)
    if (exist) then
       u_vis = free_unit ()
       open (u_vis, file = char (vis_file), action = "read", status = "old")
       iostat = 0
       do while (iostat == 0)
          read (u_vis, "(A)", iostat = iostat)  buffer
          if (iostat == 0)  write (u, "(A)")  trim (buffer)
       end do
       close (u_vis)
    else
       write (u, "(A)")  "[Visualize LaTeX file is missing]"
    end if
    inquire (file = char (ps_file), exist = exist_ps)
    if (exist_ps) then
       write (u, "(A)")  "[Visualize Postscript file exists and is nonempty]"
    else
       write (u, "(A)")  "[Visualize Postscript file is missing/non-regular]"
    end if
    inquire (file = char (pdf_file), exist = exist_pdf)
    if (exist_pdf) then
       write (u, "(A)")  "[Visualize PDF file exists and is nonempty]"
    else
       write (u, "(A)")  "[Visualize PDF file is missing/non-regular]"
    end if        
    
    write (u, "(A)")
    write (u, "(A)")  "* Cleanup"

    close (u_phs)
    call phs_data%final ()
    call model%final ()
    
    write (u, "(A)")
    write (u, "(A)")  "* Test output end: phs_wood_vis_1"

  end subroutine phs_wood_vis_1

\eatline
\nwendcode{}\nwfilename{../../../src/mci/mci.nw}\nwbegindocs{0}%% -*- ess-noweb-default-code-mode: f90-mode; noweb-default-code-mode: f90-mode; -*- 
% WHIZARD code as NOWEB source: integration and event generation
%%%%%%%%%%%%%%%%%%%%%%%%%%%%%%%%%%%%%%%%%%%%%%%%%%%%%%%%%%%%%%%%%%%%%%%%
\chapter{Multi-Channel Integration}

The abstract representation of multi-channel
Monte Carlo algorithms for integration and event generation.
\begin{description}
\item[Module {\tt{}mci{\char95}base}:]
  The abstract types and their methods.  It provides a test integrator
  that is referenced in later unit tests.
\item[iterations]
  Container for defining integration call and pass settings.
\item[integration\_results]
  This module handles results from integrating processes.  It records passes
  and iterations, calculates statistical averages, and provides the user
  output of integration results.
\end{description}

These are the implementations:
\begin{description}
\item[Module {\tt{}mci{\char95}midpoint}:]
  A simple integrator that uses the midpoint rule to sample the
  integrand uniformly over the unit hypercube.  There is only one
  integration channel, so this can be matched only to single-channel
  phase space.
\item[Module {\tt{}mci{\char95}vamp}:]
  Interface for the VAMP package.
\end{description}


%%%%%%%%%%%%%%%%%%%%%%%%%%%%%%%%%%%%%%%%%%%%%%%%%%%%%%%%%%%%%%%%%%%%%%%%
\clearpage
\section{Generic Integrator}
This module provides a multi-channel integrator (MCI) base type, a
corresponding configuration type, and methods for integration and event
generation. 

\nwenddocs{}\nwbegincode{1}\moddef{\code{}mci{\char95}base.f90\edoc{}}\endmoddef
\LA{}File header\RA{}

module mci_base

  use kinds
\LA{}Use strings\RA{}
  use io_units
  use format_utils, only: pac_fmt
  use format_defs, only: FMT_14, FMT_17
  use unit_tests
  use diagnostics
  use cputime
  use phs_base
  use rng_base

\LA{}Standard module head\RA{}

\LA{}MCI base: public\RA{}

\LA{}MCI base: types\RA{}

\LA{}MCI base: interfaces\RA{}

\LA{}MCI base: test types\RA{}

contains
  
\LA{}MCI base: procedures\RA{}

\LA{}MCI base: tests\RA{}

end module mci_base
\eatline
\nwendcode{}\nwbegindocs{2}\nwdocspar
\subsection{MCI: integrator}
The MCI object contains the methods for integration and event generation.
For the actual work and data storage, it spawns an MCI instance object.

The base object contains the number of integration dimensions and the number
of channels as configuration data.  Further configuration data are stored in
the concrete extensions.

The MCI sum contains all relevant information about the integrand.  It can be
used for comparing the current configuration against a previous one.  If they
match, we can skip an actual integration.  (Implemented only for the VAMP
version.) 

There is a random-number generator (its state with associated methods)
available as {\tt{}rng}.  It may or may not be used for integration.  It
will be used for event generation.
\nwenddocs{}\nwbegincode{3}\moddef{MCI base: public}\endmoddef
  public :: mci_t
\nwendcode{}\nwbegincode{4}\moddef{MCI base: types}\endmoddef
  type, abstract :: mci_t
     integer :: n_dim = 0
     integer :: n_channel = 0
     integer :: n_chain = 0
     integer, dimension(:), allocatable :: chain
     real(default), dimension(:), allocatable :: chain_weights
     character(32) :: md5sum = ""
     logical :: integral_known = .false.
     logical :: error_known = .false.
     logical :: efficiency_known = .false.
     real(default) :: integral = 0
     real(default) :: error = 0
     real(default) :: efficiency = 0
     logical :: use_timer = .false.
     type(timer_t) :: timer
     class(rng_t), allocatable :: rng
   contains
   \LA{}MCI base: mci: TBP\RA{}
  end type mci_t
  
\eatline
\nwendcode{}\nwbegindocs{5}Finalizer: the random-number generator may need one.
\nwenddocs{}\nwbegincode{6}\moddef{MCI base: mci: TBP}\endmoddef
  procedure :: base_final => mci_final
  procedure (mci_final), deferred :: final
\nwendcode{}\nwbegincode{7}\moddef{MCI base: procedures}\endmoddef
  subroutine mci_final (object)
    class(mci_t), intent(inout) :: object
    if (allocated (object%rng))  call object%rng%final ()
  end subroutine mci_final

\eatline
\nwendcode{}\nwbegindocs{8}Output: basic and extended output.
\nwenddocs{}\nwbegincode{9}\moddef{MCI base: mci: TBP}\plusendmoddef
  procedure :: base_write => mci_write
  procedure (mci_write), deferred :: write
\nwendcode{}\nwbegincode{10}\moddef{MCI base: procedures}\plusendmoddef
  subroutine mci_write (object, unit, pacify, md5sum_version)
    class(mci_t), intent(in) :: object
    integer, intent(in), optional :: unit
    logical, intent(in), optional :: pacify
    logical, intent(in), optional :: md5sum_version
    logical :: md5sum_ver
    integer :: u, i, j
    character(len=7) :: fmt 
    call pac_fmt (fmt, FMT_17, FMT_14, pacify)
    u = given_output_unit (unit)
    md5sum_ver = .false.
    if (present (md5sum_version))  md5sum_ver = md5sum_version
    if (object%use_timer .and. .not. md5sum_ver) then
       write (u, "(2x)", advance="no")
       call object%timer%write (u)
    end if
    if (object%integral_known) then
       write (u, "(3x,A," // fmt // ")") &
            "Integral             = ", object%integral
    end if
    if (object%error_known) then
       write (u, "(3x,A," // fmt // ")") &
            "Error                = ", object%error
    end if
    if (object%efficiency_known) then
       write (u, "(3x,A," // fmt // ")")  &
            "Efficiency           = ", object%efficiency
    end if
    write (u, "(3x,A,I0)")  "Number of channels   = ", object%n_channel
    write (u, "(3x,A,I0)")  "Number of dimensions = ", object%n_dim
    if (object%n_chain > 0) then
       write (u, "(3x,A,I0)")  "Number of chains     = ", object%n_chain
       write (u, "(3x,A)")  "Chains:"
       do i = 1, object%n_chain
          write (u, "(5x,I0,':')", advance = "no")  i
          do j = 1, object%n_channel
             if (object%chain(j) == i) &
                  write (u, "(1x,I0)", advance = "no")  j
          end do
          write (u, "(A)")
       end do
    end if
  end subroutine mci_write

\eatline
\nwendcode{}\nwbegindocs{11}Print an informative message when starting integration.
\nwenddocs{}\nwbegincode{12}\moddef{MCI base: mci: TBP}\plusendmoddef
  procedure (mci_startup_message), deferred :: startup_message
  procedure :: base_startup_message => mci_startup_message
\nwendcode{}\nwbegincode{13}\moddef{MCI base: procedures}\plusendmoddef
  subroutine mci_startup_message (mci, unit, n_calls)
    class(mci_t), intent(in) :: mci
    integer, intent(in), optional :: unit, n_calls
    if (mci%n_chain > 0) then
       write (msg_buffer, "(A,3(1x,I0,1x,A))") &
            "Integrator:", mci%n_chain, "chains,", &
            mci%n_channel, "channels,", &
            mci%n_dim, "dimensions"
    else
       write (msg_buffer, "(A,3(1x,I0,1x,A))") &
            "Integrator:", &
            mci%n_channel, "channels,", &
            mci%n_dim, "dimensions"
    end if
    call msg_message (unit = unit)
  end subroutine mci_startup_message

\eatline
\nwendcode{}\nwbegindocs{14}Record the index of the MCI object within a process.  For
multi-component processes with more than one integrator, the
integrator should know about its own index, so file names can be
unique, etc.  The default implementation does nothing, however.
\nwenddocs{}\nwbegincode{15}\moddef{MCI base: mci: TBP}\plusendmoddef
  procedure :: record_index => mci_record_index
\nwendcode{}\nwbegincode{16}\moddef{MCI base: procedures}\plusendmoddef
  subroutine mci_record_index (mci, i_mci)
    class(mci_t), intent(inout) :: mci
    integer, intent(in) :: i_mci
  end subroutine mci_record_index

\eatline
\nwendcode{}\nwbegindocs{17}There is no Initializer for the abstract type, but a generic setter
for the number of channels and dimensions.  We make two aliases
available, to be able to override it.
\nwenddocs{}\nwbegincode{18}\moddef{MCI base: mci: TBP}\plusendmoddef
  procedure :: set_dimensions => mci_set_dimensions
  procedure :: base_set_dimensions => mci_set_dimensions
\nwendcode{}\nwbegincode{19}\moddef{MCI base: procedures}\plusendmoddef
  subroutine mci_set_dimensions (mci, n_dim, n_channel)
    class(mci_t), intent(inout) :: mci
    integer, intent(in) :: n_dim
    integer, intent(in) :: n_channel
    mci%n_dim = n_dim
    mci%n_channel = n_channel
  end subroutine mci_set_dimensions
  
\eatline
\nwendcode{}\nwbegindocs{20}Declare particular dimensions as flat.  This information can be used
to simplify integration.  When generating events, the flat dimensions
should be sampled with uniform and uncorrelated distribution.  It
depends on the integrator what to do with that information.
\nwenddocs{}\nwbegincode{21}\moddef{MCI base: mci: TBP}\plusendmoddef
  procedure (mci_declare_flat_dimensions), deferred :: declare_flat_dimensions
\nwendcode{}\nwbegincode{22}\moddef{MCI base: interfaces}\endmoddef
  abstract interface
     subroutine mci_declare_flat_dimensions (mci, dim_flat)
       import
       class(mci_t), intent(inout) :: mci
       integer, dimension(:), intent(in) :: dim_flat
     end subroutine mci_declare_flat_dimensions
  end interface
  
\eatline
\nwendcode{}\nwbegindocs{23}Declare particular channels as equivalent, possibly allowing for
permutations or reflections of dimensions.  We use the information
stored in the {\tt{}phs{\char95}channel{\char95}t} object array that the phase-space module
provides.

(We do not test this here, deferring the unit test to the {\tt{}mci{\char95}vamp}
implementation where we actually use this feature.)
\nwenddocs{}\nwbegincode{24}\moddef{MCI base: mci: TBP}\plusendmoddef
  procedure (mci_declare_equivalences), deferred :: declare_equivalences
\nwendcode{}\nwbegincode{25}\moddef{MCI base: interfaces}\plusendmoddef
  abstract interface
     subroutine mci_declare_equivalences (mci, channel, dim_offset)
       import
       class(mci_t), intent(inout) :: mci
       type(phs_channel_t), dimension(:), intent(in) :: channel
       integer, intent(in) :: dim_offset
     end subroutine mci_declare_equivalences
  end interface
  
\eatline
\nwendcode{}\nwbegindocs{26}Declare particular channels as chained together.  The implementation may use
this array for keeping their weights equal to each other, etc.

The chain array is an array sized by the number of channels.  For each
channel, there is an integer entry that indicates the correponding
chains.  The total number of chains is the maximum value of this
entry.
\nwenddocs{}\nwbegincode{27}\moddef{MCI base: mci: TBP}\plusendmoddef
  procedure :: declare_chains => mci_declare_chains
\nwendcode{}\nwbegincode{28}\moddef{MCI base: procedures}\plusendmoddef
  subroutine mci_declare_chains (mci, chain)
    class(mci_t), intent(inout) :: mci
    integer, dimension(:), intent(in) :: chain
    allocate (mci%chain (size (chain)))
    mci%n_chain = maxval (chain)
    allocate (mci%chain_weights (mci%n_chain), source = 0._default)
    mci%chain = chain
  end subroutine mci_declare_chains
  
\eatline
\nwendcode{}\nwbegindocs{29}Collect channel weights according to chains and store them in the
{\tt{}chain{\char95}weights} for output.  We sum up the weights for all channels that
share the same {\tt{}chain} index and store the results in the {\tt{}chain{\char95}weights}
array.
\nwenddocs{}\nwbegincode{30}\moddef{MCI base: mci: TBP}\plusendmoddef
  procedure :: collect_chain_weights => mci_collect_chain_weights
\nwendcode{}\nwbegincode{31}\moddef{MCI base: procedures}\plusendmoddef
  subroutine mci_collect_chain_weights (mci, weight)
    class(mci_t), intent(inout) :: mci
    real(default), dimension(:), intent(in) :: weight
    integer :: i, c
    if (allocated (mci%chain)) then
       mci%chain_weights = 0
       do i = 1, size (mci%chain)
          c = mci%chain(i)
          mci%chain_weights(c) = mci%chain_weights(c) + weight(i)
       end do
    end if
  end subroutine mci_collect_chain_weights
    
\eatline
\nwendcode{}\nwbegindocs{32}Check if there are chains.
\nwenddocs{}\nwbegincode{33}\moddef{MCI base: mci: TBP}\plusendmoddef
  procedure :: has_chains => mci_has_chains
\nwendcode{}\nwbegincode{34}\moddef{MCI base: procedures}\plusendmoddef
  function mci_has_chains (mci) result (flag)
    class(mci_t), intent(in) :: mci
    logical :: flag
    flag = allocated (mci%chain)
  end function mci_has_chains

\eatline
\nwendcode{}\nwbegindocs{35}Output of the chain weights, kept separate from the main {\tt{}write} method.

[The formatting will work as long as the number of chains is less than
  $10^{10}$\ldots]
\nwenddocs{}\nwbegincode{36}\moddef{MCI base: mci: TBP}\plusendmoddef
  procedure :: write_chain_weights => mci_write_chain_weights
\nwendcode{}\nwbegincode{37}\moddef{MCI base: procedures}\plusendmoddef
  subroutine mci_write_chain_weights (mci, unit)
    class(mci_t), intent(in) :: mci
    integer, intent(in), optional :: unit
    integer :: u, i, n, n_digits
    character(4) :: ifmt
    u = given_output_unit (unit)
    if (allocated (mci%chain_weights)) then
       write (u, "(1x,A)")  "Weights of channel chains (groves):"
       n_digits = 0
       n = size (mci%chain_weights)
       do while (n > 0)
          n = n / 10
          n_digits = n_digits + 1
       end do
       write (ifmt, "(A1,I1)") "I", n_digits
       do i = 1, size (mci%chain_weights)
          write (u, "(3x," // ifmt // ",F13.10)")  i, mci%chain_weights(i)
       end do
    end if
  end subroutine mci_write_chain_weights
  
\eatline
\nwendcode{}\nwbegindocs{38}Set the MD5 sum, independent of initialization.
\nwenddocs{}\nwbegincode{39}\moddef{MCI base: mci: TBP}\plusendmoddef
  procedure :: set_md5sum => mci_set_md5sum
\nwendcode{}\nwbegincode{40}\moddef{MCI base: procedures}\plusendmoddef
  subroutine mci_set_md5sum (mci, md5sum)
    class(mci_t), intent(inout) :: mci
    character(32), intent(in) :: md5sum
    mci%md5sum = md5sum
  end subroutine mci_set_md5sum
  
\eatline
\nwendcode{}\nwbegindocs{41}Initialize a new integration pass.  This is not necessarily
meaningful, so we provide an empty base method.  The {\tt{}mci{\char95}vamp}
implementation overrides this.
\nwenddocs{}\nwbegincode{42}\moddef{MCI base: mci: TBP}\plusendmoddef
  procedure :: add_pass => mci_add_pass
\nwendcode{}\nwbegincode{43}\moddef{MCI base: procedures}\plusendmoddef
  subroutine mci_add_pass (mci, adapt_grids, adapt_weights, final)
    class(mci_t), intent(inout) :: mci
    logical, intent(in), optional :: adapt_grids
    logical, intent(in), optional :: adapt_weights
    logical, intent(in), optional :: final
  end subroutine mci_add_pass
    
\eatline
\nwendcode{}\nwbegindocs{44}Allocate an instance with matching type.  This must be deferred.
\nwenddocs{}\nwbegincode{45}\moddef{MCI base: mci: TBP}\plusendmoddef
  procedure (mci_allocate_instance), deferred :: allocate_instance
\nwendcode{}\nwbegincode{46}\moddef{MCI base: interfaces}\plusendmoddef
  abstract interface
     subroutine mci_allocate_instance (mci, mci_instance)
       import
       class(mci_t), intent(in) :: mci
       class(mci_instance_t), intent(out), pointer :: mci_instance
     end subroutine mci_allocate_instance
  end interface
       
\eatline
\nwendcode{}\nwbegindocs{47}Import a random-number generator.  We transfer the allocation of an
existing generator state into the object.  The generator state may
already be initialized, or we can reset it by its {\tt{}init}
method.
\nwenddocs{}\nwbegincode{48}\moddef{MCI base: mci: TBP}\plusendmoddef
  procedure :: import_rng => mci_import_rng
\nwendcode{}\nwbegincode{49}\moddef{MCI base: procedures}\plusendmoddef
  subroutine mci_import_rng (mci, rng)
    class(mci_t), intent(inout) :: mci
    class(rng_t), intent(inout), allocatable :: rng
    call move_alloc (rng, mci%rng)
  end subroutine mci_import_rng
  
\eatline
\nwendcode{}\nwbegindocs{50}Activate or deactivate the timer.
\nwenddocs{}\nwbegincode{51}\moddef{MCI base: mci: TBP}\plusendmoddef
  procedure :: set_timer => mci_set_timer
\nwendcode{}\nwbegincode{52}\moddef{MCI base: procedures}\plusendmoddef
  subroutine mci_set_timer (mci, active)
    class(mci_t), intent(inout) :: mci
    logical, intent(in) :: active
    mci%use_timer = active
  end subroutine mci_set_timer
    
\eatline
\nwendcode{}\nwbegindocs{53}Start and stop signal for the timer, if active.  The elapsed time
can then be retrieved from the MCI record.
\nwenddocs{}\nwbegincode{54}\moddef{MCI base: mci: TBP}\plusendmoddef
  procedure :: start_timer => mci_start_timer
  procedure :: stop_timer => mci_stop_timer
\nwendcode{}\nwbegincode{55}\moddef{MCI base: procedures}\plusendmoddef
  subroutine mci_start_timer (mci)
    class(mci_t), intent(inout) :: mci
    if (mci%use_timer)  call mci%timer%start ()
  end subroutine mci_start_timer
  
  subroutine mci_stop_timer (mci)
    class(mci_t), intent(inout) :: mci
    if (mci%use_timer)  call mci%timer%stop ()
  end subroutine mci_stop_timer
  
\eatline
\eatline
\nwendcode{}\nwbegindocs{56}Sampler test.  Evaluate the sampler a given number of times.  Results are
discarded, so we don't need the MCI instance which would record them.

The evaluation channel is iterated, and the {\tt{}x} parameters are randomly
chosen.
\nwenddocs{}\nwbegincode{57}\moddef{MCI base: mci: TBP}\plusendmoddef
  procedure :: sampler_test => mci_sampler_test
\nwendcode{}\nwbegincode{58}\moddef{MCI base: procedures}\plusendmoddef
  subroutine mci_sampler_test (mci, sampler, n_calls)
    class(mci_t), intent(inout) :: mci
    class(mci_sampler_t), intent(inout), target :: sampler
    integer, intent(in) :: n_calls
    real(default), dimension(:), allocatable :: x_in, f
    real(default), dimension(:,:), allocatable :: x_out
    real(default) :: val
    integer :: i, c
    allocate (x_in (mci%n_dim))
    allocate (f (mci%n_channel))
    allocate (x_out (mci%n_dim, mci%n_channel))
    do i = 1, n_calls
       c = mod (i, mci%n_channel) + 1
       call mci%rng%generate_array (x_in)
       call sampler%evaluate (c, x_in, val, x_out, f)
    end do
  end subroutine mci_sampler_test
  
\eatline
\nwendcode{}\nwbegindocs{59}Integrate: this depends on the implementation. We foresee a pacify
flag to take care of small numerical noise on different platforms.
\nwenddocs{}\nwbegincode{60}\moddef{MCI base: mci: TBP}\plusendmoddef
  procedure (mci_integrate), deferred :: integrate
\nwendcode{}\nwbegincode{61}\moddef{MCI base: interfaces}\plusendmoddef
  abstract interface
     subroutine mci_integrate (mci, instance, sampler, &
          n_it, n_calls, results, pacify)
       import
       class(mci_t), intent(inout) :: mci
       class(mci_instance_t), intent(inout) :: instance
       class(mci_sampler_t), intent(inout) :: sampler
       integer, intent(in) :: n_it
       integer, intent(in) :: n_calls
       logical, intent(in), optional :: pacify
       class(mci_results_t), intent(inout), optional :: results
     end subroutine mci_integrate
  end interface
  
\eatline
\nwendcode{}\nwbegindocs{62}Event generation.  Depending on the implementation,
event generation may or may not require a previous integration pass.

Instead of a black-box {\tt{}simulate} method, we require an initializer,
a finalizer, and procedures for generating a single event.  This
allows us to interface simulation event by event from the outside, and
it facilitates the further processing of an event after successful
generation.  For integration, this is not necessary.

The initializer has {\tt{}intent(inout)} for the {\tt{}mci} passed object.  The
reason is that the initializer can read integration results and grids from
file, where the results can modify the {\tt{}mci} record.
\nwenddocs{}\nwbegincode{63}\moddef{MCI base: mci: TBP}\plusendmoddef
  procedure (mci_prepare_simulation), deferred :: prepare_simulation
\eatline
\nwendcode{}\nwbegincode{64}\moddef{MCI base: interfaces}\plusendmoddef
  abstract interface
     subroutine mci_prepare_simulation (mci)
       import
       class(mci_t), intent(inout) :: mci
     end subroutine mci_prepare_simulation
  end interface
  
\eatline
\nwendcode{}\nwbegindocs{65}\nwdocspar
The generated event will reside in in the {\tt{}instance} object (overall
results and weight) and in the {\tt{}sampler} object (detailed data).  In
the real application, we can subsequently call methods of the
{\tt{}sampler} in order to further process the generated event.

The {\tt{}target} attributes are required by the VAMP implementation,
which uses pointers to refer to the instance and sampler objects from
within the integration function.
\nwenddocs{}\nwbegincode{66}\moddef{MCI base: mci: TBP}\plusendmoddef
  procedure (mci_generate), deferred :: generate_weighted_event
  procedure (mci_generate), deferred :: generate_unweighted_event
\eatline
\eatline
\nwendcode{}\nwbegincode{67}\moddef{MCI base: interfaces}\plusendmoddef
  abstract interface
     subroutine mci_generate (mci, instance, sampler)
       import
       class(mci_t), intent(inout) :: mci
       class(mci_instance_t), intent(inout), target :: instance
       class(mci_sampler_t), intent(inout), target :: sampler
     end subroutine mci_generate
  end interface
  
\eatline
\nwendcode{}\nwbegindocs{68}This is analogous, but we rebuild the event from the information
stored in {\tt{}state} instead of generating it.

Note: currently unused outside of tests, might be deleted later.
\nwenddocs{}\nwbegincode{69}\moddef{MCI base: mci: TBP}\plusendmoddef
  procedure (mci_rebuild), deferred :: rebuild_event
\nwendcode{}\nwbegincode{70}\moddef{MCI base: interfaces}\plusendmoddef
  abstract interface
     subroutine mci_rebuild (mci, instance, sampler, state)
       import
       class(mci_t), intent(inout) :: mci
       class(mci_instance_t), intent(inout) :: instance
       class(mci_sampler_t), intent(inout) :: sampler
       class(mci_state_t), intent(in) :: state
     end subroutine mci_rebuild
  end interface
  
\eatline
\nwendcode{}\nwbegindocs{71}Return the value of the integral, error, efficiency, and time per call.
\nwenddocs{}\nwbegincode{72}\moddef{MCI base: mci: TBP}\plusendmoddef
  procedure :: get_integral => mci_get_integral
  procedure :: get_error => mci_get_error
  procedure :: get_efficiency => mci_get_efficiency
  procedure :: get_time => mci_get_time
\nwendcode{}\nwbegincode{73}\moddef{MCI base: procedures}\plusendmoddef
  function mci_get_integral (mci) result (integral)
    class(mci_t), intent(in) :: mci
    real(default) :: integral
    if (mci%integral_known) then
       integral = mci%integral
    else
       call msg_bug ("The integral is unknown. This is presumably a" // &
            "WHIZARD bug.")
    end if
  end function mci_get_integral
  
  function mci_get_error (mci) result (error)
    class(mci_t), intent(in) :: mci
    real(default) :: error
    if (mci%error_known) then
       error = mci%error
    else
       error = 0
    end if
  end function mci_get_error
  
  function mci_get_efficiency (mci) result (efficiency)
    class(mci_t), intent(in) :: mci
    real(default) :: efficiency
    if (mci%efficiency_known) then
       efficiency = mci%efficiency
    else
       efficiency = 0
    end if
  end function mci_get_efficiency
  
  function mci_get_time (mci) result (time)
    class(mci_t), intent(in) :: mci
    real(default) :: time
    if (mci%use_timer) then
       time = mci%timer
    else
       time = 0
    end if
  end function mci_get_time
  
\eatline
\eatline
\eatline
\eatline
\nwendcode{}\nwbegindocs{74}Return the MD5 sum of the configuration.  This may be overridden in
an extension, to return a different MD5 sum.
\nwenddocs{}\nwbegincode{75}\moddef{MCI base: mci: TBP}\plusendmoddef
  procedure :: get_md5sum => mci_get_md5sum
\nwendcode{}\nwbegincode{76}\moddef{MCI base: procedures}\plusendmoddef
  function mci_get_md5sum (mci) result (md5sum)
    class(mci_t), intent(in) :: mci
    character(32) :: md5sum
    md5sum = mci%md5sum
  end function mci_get_md5sum
  
\eatline
\nwendcode{}\nwbegindocs{77}\nwdocspar
\subsection{MCI instance}
The base type contains an array of channel weights.  The value {\tt{}mci{\char95}weight}
is the combined MCI weight that corresponds to a particular sampling point.

For convenience, we also store the {\tt{}x} and Jacobian values for this sampling
point.
\nwenddocs{}\nwbegincode{78}\moddef{MCI base: public}\plusendmoddef
  public :: mci_instance_t
\nwendcode{}\nwbegincode{79}\moddef{MCI base: types}\plusendmoddef
  type, abstract :: mci_instance_t
     logical :: valid = .false.
     real(default), dimension(:), allocatable :: w
     real(default), dimension(:), allocatable :: f
     real(default), dimension(:,:), allocatable :: x
     integer :: selected_channel = 0
     real(default) :: mci_weight = 0
     real(default) :: integrand  = 0
   contains
   \LA{}MCI base: mci instance: TBP\RA{}
  end type mci_instance_t
  
\eatline
\nwendcode{}\nwbegindocs{80}Output: deferred
\nwenddocs{}\nwbegincode{81}\moddef{MCI base: mci instance: TBP}\endmoddef
  procedure (mci_instance_write), deferred :: write
\nwendcode{}\nwbegincode{82}\moddef{MCI base: interfaces}\plusendmoddef
  abstract interface
     subroutine mci_instance_write (object, unit, pacify)
       import
       class(mci_instance_t), intent(in) :: object
       integer, intent(in), optional :: unit
       logical, intent(in), optional :: pacify
     end subroutine mci_instance_write
  end interface
    
\eatline
\nwendcode{}\nwbegindocs{83}A finalizer, just in case.
\nwenddocs{}\nwbegincode{84}\moddef{MCI base: mci instance: TBP}\plusendmoddef
  procedure (mci_instance_final), deferred :: final
\nwendcode{}\nwbegincode{85}\moddef{MCI base: interfaces}\plusendmoddef
  abstract interface
     subroutine mci_instance_final (object)
       import
       class(mci_instance_t), intent(inout) :: object
     end subroutine mci_instance_final
  end interface
  
\eatline
\nwendcode{}\nwbegindocs{86}Init: basic initializer for the arrays, otherwise deferred.  Assigning
the {\tt{}mci} object is also deferred, because it depends on the concrete type.

The weights are initialized with an uniform normalized value.
\nwenddocs{}\nwbegincode{87}\moddef{MCI base: mci instance: TBP}\plusendmoddef
  procedure (mci_instance_base_init), deferred :: init
  procedure :: base_init => mci_instance_base_init
\nwendcode{}\nwbegincode{88}\moddef{MCI base: procedures}\plusendmoddef
  subroutine mci_instance_base_init (mci_instance, mci)
    class(mci_instance_t), intent(out) :: mci_instance
    class(mci_t), intent(in), target :: mci
    allocate (mci_instance%w (mci%n_channel))
    allocate (mci_instance%f (mci%n_channel))
    allocate (mci_instance%x (mci%n_dim, mci%n_channel))
    if (size (mci_instance%w) /= 0) then
       mci_instance%w = 1._default / size (mci_instance%w)
    end if
    mci_instance%f = 0
    mci_instance%x = 0
  end subroutine mci_instance_base_init
    
\eatline
\nwendcode{}\nwbegindocs{89}Compute the overall weight factor for a configuration of $x$ values and
Jacobians $f$.  The $x$ values come in {\tt{}n{\char95}channel} rows with {\tt{}n{\char95}dim}
entries each.  The $f$ factors constitute an array with {\tt{}n{\char95}channel} entries.

We assume that the $x$ and $f$ arrays are already stored inside the MC
instance.  The result is also stored there.
\nwenddocs{}\nwbegincode{90}\moddef{MCI base: mci instance: TBP}\plusendmoddef
  procedure (mci_instance_compute_weight), deferred :: compute_weight
\nwendcode{}\nwbegincode{91}\moddef{MCI base: interfaces}\plusendmoddef
  abstract interface
     subroutine mci_instance_compute_weight (mci, c)
       import
       class(mci_instance_t), intent(inout) :: mci
       integer, intent(in) :: c
     end subroutine mci_instance_compute_weight
  end interface
    
\eatline
\nwendcode{}\nwbegindocs{92}Record the integrand as returned by the sampler.  Depending on the
implementation, this may merely copy the value, or do more complicated things.

We may need the MCI weight for the actual computations, so this should
be called after the previous routine.
\nwenddocs{}\nwbegincode{93}\moddef{MCI base: mci instance: TBP}\plusendmoddef
  procedure (mci_instance_record_integrand), deferred :: record_integrand
\nwendcode{}\nwbegincode{94}\moddef{MCI base: interfaces}\plusendmoddef
  abstract interface
     subroutine mci_instance_record_integrand (mci, integrand)
       import
       class(mci_instance_t), intent(inout) :: mci
       real(default), intent(in) :: integrand
     end subroutine mci_instance_record_integrand
  end interface
  
\eatline
\nwendcode{}\nwbegindocs{95}Sample a point directly: evaluate the sampler, then compute the weight and
the weighted integrand.  Finally, record the integrand within the MCI instance.

If a signal (interrupt) was raised recently, we abort the calculation before
entering the sampler.  Thus, a previous calculation will have
completed and any data are already recorded, but any new point can be
discarded.  If the {\tt{}abort} flag is present, we may delay the interrupt, so
we can do some cleanup.
\nwenddocs{}\nwbegincode{96}\moddef{MCI base: mci instance: TBP}\plusendmoddef
  procedure :: evaluate => mci_instance_evaluate
\nwendcode{}\nwbegincode{97}\moddef{MCI base: procedures}\plusendmoddef
  subroutine mci_instance_evaluate (mci, sampler, c, x)
    class(mci_instance_t), intent(inout) :: mci
    class(mci_sampler_t), intent(inout) :: sampler
    integer, intent(in) :: c
    real(default), dimension(:), intent(in) :: x
    real(default) :: val
    call sampler%evaluate (c, x, val, mci%x, mci%f)
    mci%valid = sampler%is_valid ()
    if (mci%valid) then
       call mci%compute_weight (c)
       call mci%record_integrand (val)
    end if
  end subroutine mci_instance_evaluate
    
\eatline
\nwendcode{}\nwbegindocs{98}Initiate and terminate simulation.  In contrast to integration, we implement
these as methods of the process instance, since the {\tt{}mci} configuration
object is unchanged.

The safety factor reduces the acceptance probability for unweighted
events.  The implementation of this feature depends on the concrete type.
\nwenddocs{}\nwbegincode{99}\moddef{MCI base: mci instance: TBP}\plusendmoddef
  procedure (mci_instance_init_simulation), deferred :: init_simulation
  procedure (mci_instance_final_simulation), deferred :: final_simulation
\nwendcode{}\nwbegincode{100}\moddef{MCI base: interfaces}\plusendmoddef
  abstract interface
     subroutine mci_instance_init_simulation (instance, safety_factor)
       import
       class(mci_instance_t), intent(inout) :: instance
       real(default), intent(in), optional :: safety_factor
     end subroutine mci_instance_init_simulation
  end interface
  
  abstract interface
     subroutine mci_instance_final_simulation (instance)
       import
       class(mci_instance_t), intent(inout) :: instance
     end subroutine mci_instance_final_simulation
  end interface
  
\eatline
\nwendcode{}\nwbegindocs{101}Assuming that the sampler is in a completely defined state, just
extract the data that {\tt{}evaluate} would compute.  Also record the integrand.
\nwenddocs{}\nwbegincode{102}\moddef{MCI base: mci instance: TBP}\plusendmoddef
  procedure :: fetch => mci_instance_fetch
\nwendcode{}\nwbegincode{103}\moddef{MCI base: procedures}\plusendmoddef
  subroutine mci_instance_fetch (mci, sampler, c)
    class(mci_instance_t), intent(inout) :: mci
    class(mci_sampler_t), intent(in) :: sampler
    integer, intent(in) :: c
    real(default) :: val
    mci%valid = sampler%is_valid ()
    if (mci%valid) then
       call sampler%fetch (val, mci%x, mci%f)
       call mci%compute_weight (c)
       call mci%record_integrand (val)
    end if
  end subroutine mci_instance_fetch

\eatline
\nwendcode{}\nwbegindocs{104}The value, i.e., the weighted integrand, is the integrand (which
should be taken as-is from the sampler) multiplied by the MCI weight.
\nwenddocs{}\nwbegincode{105}\moddef{MCI base: mci instance: TBP}\plusendmoddef
  procedure :: get_value => mci_instance_get_value
\nwendcode{}\nwbegincode{106}\moddef{MCI base: procedures}\plusendmoddef
  function mci_instance_get_value (mci) result (value)
    class(mci_instance_t), intent(in) :: mci
    real(default) :: value
    if (mci%valid) then
       value = mci%integrand * mci%mci_weight
    else
       value = 0
    end if
  end function mci_instance_get_value

\eatline
\nwendcode{}\nwbegindocs{107}This is an extra routine.  By default, the event weight is equal to
the value returned by the previous routine.  However, if we select a
channel for event generation not just based on the channel weights,
the event weight has to account for this bias, so the event weight
that applies to event generation is different.  In that case, we
should override the default routine.
\nwenddocs{}\nwbegincode{108}\moddef{MCI base: mci instance: TBP}\plusendmoddef
  procedure :: get_event_weight => mci_instance_get_value
\eatline
\nwendcode{}\nwbegindocs{109}Excess weight can occur during unweighted event generation, if the
assumed maximum value of the integrand is too small.  This excess
should be normalized in the same way as the event weight above (which
for unweighted events becomes unity).
\nwenddocs{}\nwbegincode{110}\moddef{MCI base: mci instance: TBP}\plusendmoddef
  procedure (mci_instance_get_event_excess), deferred :: get_event_excess
\nwendcode{}\nwbegincode{111}\moddef{MCI base: interfaces}\plusendmoddef
  abstract interface
     function mci_instance_get_event_excess (mci) result (excess)
       import
       class(mci_instance_t), intent(in) :: mci
       real(default) :: excess
     end function mci_instance_get_event_excess
  end interface
  
\eatline
\nwendcode{}\nwbegindocs{112}\nwdocspar
\subsection{MCI state}
This object can hold the relevant information that allows us to
reconstruct the MCI instance without re-evaluating the sampler completely.

We store the {\tt{}x{\char95}in} MC input parameter set, which coincides with the
section of the complete {\tt{}x} array that belongs to a particular
channel.  We also store the MC function value.  When we want to
reconstruct the state, we can use the input array to recover the
complete {\tt{}x} and {\tt{}f} arrays (i.e., the kinematics), but do not
need to recompute the MC function value (the dynamics).

The {\tt{}mci{\char95}state{\char95}t} may be extended, to allow storing/recalling more
information.  In that case, we would override the type-bound
procedures.  However, the base type is also a concrete type and
self-contained.
\nwenddocs{}\nwbegincode{113}\moddef{MCI base: public}\plusendmoddef
  public :: mci_state_t
\nwendcode{}\nwbegincode{114}\moddef{MCI base: types}\plusendmoddef
  type :: mci_state_t
     integer :: selected_channel = 0
     real(default), dimension(:), allocatable :: x_in
     real(default) :: val
   contains
   \LA{}MCI base: mci state: TBP\RA{}
  end type mci_state_t
  
\eatline
\nwendcode{}\nwbegindocs{115}Output:
\nwenddocs{}\nwbegincode{116}\moddef{MCI base: mci state: TBP}\endmoddef
  procedure :: write => mci_state_write
\nwendcode{}\nwbegincode{117}\moddef{MCI base: procedures}\plusendmoddef
  subroutine mci_state_write (object, unit)
    class(mci_state_t), intent(in) :: object
    integer, intent(in), optional :: unit
    integer :: u
    u = given_output_unit (unit)
    write (u, "(1x,A)")  "MCI state:"
    write (u, "(3x,A,I0)")  "Channel   = ", object%selected_channel
    write (u, "(3x,A,999(1x,F12.10))")  "x (in)    =", object%x_in
    write (u, "(3x,A,ES19.12)")  "Integrand = ", object%val
  end subroutine mci_state_write
  
\eatline
\nwendcode{}\nwbegindocs{118}To store the object, we take the relevant section of the {\tt{}x}
array.  The channel used for storing data is taken from the
{\tt{}instance} object, but it could be arbitrary in principle.
\nwenddocs{}\nwbegincode{119}\moddef{MCI base: mci instance: TBP}\plusendmoddef
  procedure :: store => mci_instance_store
\nwendcode{}\nwbegincode{120}\moddef{MCI base: procedures}\plusendmoddef
  subroutine mci_instance_store (mci, state)
    class(mci_instance_t), intent(in) :: mci
    class(mci_state_t), intent(out) :: state
    state%selected_channel = mci%selected_channel
    allocate (state%x_in (size (mci%x, 1)))
    state%x_in = mci%x(:,mci%selected_channel)
    state%val = mci%integrand
  end subroutine mci_instance_store
    
\eatline
\nwendcode{}\nwbegindocs{121}Recalling the state, we must consult the sampler in order to fully
reconstruct the {\tt{}x} and {\tt{}f} arrays.  The integrand value is known,
and we also give it to the sampler, bypassing evaluation.

The final steps are equivalent to the {\tt{}evaluate} method above.
\nwenddocs{}\nwbegincode{122}\moddef{MCI base: mci instance: TBP}\plusendmoddef
  procedure :: recall => mci_instance_recall
\nwendcode{}\nwbegincode{123}\moddef{MCI base: procedures}\plusendmoddef
  subroutine mci_instance_recall (mci, sampler, state)
    class(mci_instance_t), intent(inout) :: mci
    class(mci_sampler_t), intent(inout) :: sampler
    class(mci_state_t), intent(in) :: state
    if (size (state%x_in) == size (mci%x, 1) &
         .and. state%selected_channel <= size (mci%x, 2)) then
       call sampler%rebuild (state%selected_channel, &
            state%x_in, state%val, mci%x, mci%f)
       call mci%compute_weight (state%selected_channel)
       call mci%record_integrand (state%val)
    else
       call msg_fatal ("Recalling event: mismatch in channel or dimension")
    end if
  end subroutine mci_instance_recall
    
\eatline
\nwendcode{}\nwbegindocs{124}\nwdocspar
\subsection{MCI sampler}
A sampler is an object that implements a multi-channel parameterization of the
unit hypercube.  Specifically, it is able to compute, given a channel and a
set of $x$ MC parameter values, a the complete set of $x$ values and
associated Jacobian factors $f$ for all channels.

Furthermore, the sampler should return a single real value, the integrand, for
the given point in the hypercube.

It must implement a method {\tt{}evaluate} for performing the above
computations.
\nwenddocs{}\nwbegincode{125}\moddef{MCI base: public}\plusendmoddef
  public :: mci_sampler_t
\nwendcode{}\nwbegincode{126}\moddef{MCI base: types}\plusendmoddef
  type, abstract :: mci_sampler_t
   contains
   \LA{}MCI base: mci sampler: TBP\RA{}
  end type mci_sampler_t

\eatline
\nwendcode{}\nwbegindocs{127}Output, deferred to the implementation.
\nwenddocs{}\nwbegincode{128}\moddef{MCI base: mci sampler: TBP}\endmoddef
  procedure (mci_sampler_write), deferred :: write
\nwendcode{}\nwbegincode{129}\moddef{MCI base: interfaces}\plusendmoddef
  abstract interface
     subroutine mci_sampler_write (object, unit, testflag)
       import
       class(mci_sampler_t), intent(in) :: object
       integer, intent(in), optional :: unit
       logical, intent(in), optional :: testflag
     end subroutine mci_sampler_write
  end interface
  
\eatline
\nwendcode{}\nwbegindocs{130}The evaluation routine.  Input is the channel index {\tt{}c} and the
one-dimensional parameter array {\tt{}x{\char95}in}.  Output are the integrand value
{\tt{}val}, the two-dimensional parameter array {\tt{}x} and the Jacobian array
{\tt{}f}.
\nwenddocs{}\nwbegincode{131}\moddef{MCI base: mci sampler: TBP}\plusendmoddef
  procedure (mci_sampler_evaluate), deferred :: evaluate
\nwendcode{}\nwbegincode{132}\moddef{MCI base: interfaces}\plusendmoddef
  abstract interface
     subroutine mci_sampler_evaluate (sampler, c, x_in, val, x, f)
       import
       class(mci_sampler_t), intent(inout) :: sampler
       integer, intent(in) :: c
       real(default), dimension(:), intent(in) :: x_in
       real(default), intent(out) :: val
       real(default), dimension(:,:), intent(out) :: x
       real(default), dimension(:), intent(out) :: f
     end subroutine mci_sampler_evaluate
  end interface

\eatline
\nwendcode{}\nwbegindocs{133}Query the validity of the sampling point.  Can be called after
{\tt{}evaluate}.
\nwenddocs{}\nwbegincode{134}\moddef{MCI base: mci sampler: TBP}\plusendmoddef
  procedure (mci_sampler_is_valid), deferred :: is_valid
\nwendcode{}\nwbegincode{135}\moddef{MCI base: interfaces}\plusendmoddef
  abstract interface
     function mci_sampler_is_valid (sampler) result (valid)
       import
       class(mci_sampler_t), intent(in) :: sampler
       logical :: valid
     end function mci_sampler_is_valid
  end interface

\eatline
\nwendcode{}\nwbegindocs{136}The shortcut.  Again, the channel index {\tt{}c} and the parameter
array {\tt{}x{\char95}in} are input.  However, we also provide the integrand
value {\tt{}val}, and we just require that the complete parameter array
{\tt{}x} and Jacobian array {\tt{}f} are recovered.
\nwenddocs{}\nwbegincode{137}\moddef{MCI base: mci sampler: TBP}\plusendmoddef
  procedure (mci_sampler_rebuild), deferred :: rebuild
\nwendcode{}\nwbegincode{138}\moddef{MCI base: interfaces}\plusendmoddef
  abstract interface
     subroutine mci_sampler_rebuild (sampler, c, x_in, val, x, f)
       import
       class(mci_sampler_t), intent(inout) :: sampler
       integer, intent(in) :: c
       real(default), dimension(:), intent(in) :: x_in
       real(default), intent(in) :: val
       real(default), dimension(:,:), intent(out) :: x
       real(default), dimension(:), intent(out) :: f
     end subroutine mci_sampler_rebuild
  end interface
  
\eatline
\nwendcode{}\nwbegindocs{139}This routine should extract the important data from a sampler that
has been filled by other means.  We fetch the integrand value {\tt{}val},
the two-dimensional parameter array {\tt{}x} and the Jacobian array {\tt{}f}.
\nwenddocs{}\nwbegincode{140}\moddef{MCI base: mci sampler: TBP}\plusendmoddef
  procedure (mci_sampler_fetch), deferred :: fetch
\nwendcode{}\nwbegincode{141}\moddef{MCI base: interfaces}\plusendmoddef
  abstract interface
     subroutine mci_sampler_fetch (sampler, val, x, f)
       import
       class(mci_sampler_t), intent(in) :: sampler
       real(default), intent(out) :: val
       real(default), dimension(:,:), intent(out) :: x
       real(default), dimension(:), intent(out) :: f
     end subroutine mci_sampler_fetch
  end interface
  
\eatline
\nwendcode{}\nwbegindocs{142}\nwdocspar
\subsection{Results record}
This is an abstract type which allows us to implement callback: each
integration results can optionally be recorded to an instance of this
object.  The actual object may store a new result, average results,
etc.  It may also display a result on-line or otherwise, whenever
the {\tt{}record} method is called.
\nwenddocs{}\nwbegincode{143}\moddef{MCI base: public}\plusendmoddef
  public :: mci_results_t
\nwendcode{}\nwbegincode{144}\moddef{MCI base: types}\plusendmoddef
  type, abstract :: mci_results_t
   contains
   \LA{}MCI base: mci results: TBP\RA{}
  end type mci_results_t
  
\eatline
\nwendcode{}\nwbegindocs{145}The output routine is deferred.  We provide an extra {\tt{}verbose}
flag, which could serve any purpose.
\nwenddocs{}\nwbegincode{146}\moddef{MCI base: mci results: TBP}\endmoddef
  procedure (mci_results_write), deferred :: write
\nwendcode{}\nwbegincode{147}\moddef{MCI base: interfaces}\plusendmoddef
  abstract interface
     subroutine mci_results_write (object, unit, verbose, suppress)
       import
       class(mci_results_t), intent(in) :: object
       integer, intent(in), optional :: unit
       logical, intent(in), optional :: verbose, suppress
     end subroutine mci_results_write
  end interface
  
\eatline
\nwendcode{}\nwbegindocs{148}This is the {\tt{}record} method, which can be called directly from the
integrator.
\nwenddocs{}\nwbegincode{149}\moddef{MCI base: mci results: TBP}\plusendmoddef
  procedure (mci_results_record), deferred :: record
\nwendcode{}\nwbegincode{150}\moddef{MCI base: interfaces}\plusendmoddef
  abstract interface
     subroutine mci_results_record (object, n_it, &
          n_calls, integral, error, efficiency, chain_weights, suppress)
       import
       class(mci_results_t), intent(inout) :: object
       integer, intent(in) :: n_it
       integer, intent(in) :: n_calls
       real(default), intent(in) :: integral
       real(default), intent(in) :: error
       real(default), intent(in) :: efficiency
       real(default), dimension(:), intent(in), optional :: chain_weights
       logical, intent(in), optional :: suppress
     end subroutine mci_results_record
  end interface

\eatline
\nwendcode{}\nwbegindocs{151}\nwdocspar
\subsection{Unit tests}
\nwenddocs{}\nwbegincode{152}\moddef{MCI base: public}\plusendmoddef
  public :: mci_base_test
\nwendcode{}\nwbegincode{153}\moddef{MCI base: tests}\endmoddef
  subroutine mci_base_test (u, results)
    integer, intent(in) :: u
    type(test_results_t), intent(inout) :: results
  \LA{}MCI base: execute tests\RA{}
  end subroutine mci_base_test
  
\eatline
\nwendcode{}\nwbegindocs{154}\nwdocspar
\subsubsection{Test implementation of the configuration type}
The concrete type contains the number of requested calls and the integral
result, to be determined.

The {\tt{}max{\char95}factor} entry is set for the actual test integration, where the
integrand is not unity but some other constant value.  This value should be
set here, such that the actual maximum of the integrand is known when
vetoing unweighted events.
\nwenddocs{}\nwbegincode{155}\moddef{MCI base: public}\plusendmoddef
  public :: mci_test_t
\nwendcode{}\nwbegincode{156}\moddef{MCI base: test types}\endmoddef
  type, extends (mci_t) :: mci_test_t
     integer :: divisions = 0
     integer :: tries = 0
     real(default) :: max_factor = 1
   contains
     procedure :: final => mci_test_final
     procedure :: write => mci_test_write
     procedure :: startup_message => mci_test_startup_message
     procedure :: declare_flat_dimensions => mci_test_ignore_flat_dimensions
     procedure :: declare_equivalences => mci_test_ignore_equivalences
     procedure :: set_divisions => mci_test_set_divisions
     procedure :: set_max_factor => mci_test_set_max_factor
     procedure :: allocate_instance => mci_test_allocate_instance
     procedure :: integrate => mci_test_integrate
     procedure :: prepare_simulation => mci_test_ignore_prepare_simulation
     procedure :: generate_weighted_event => mci_test_generate_weighted_event
     procedure :: generate_unweighted_event => &
          mci_test_generate_unweighted_event
     procedure :: rebuild_event => mci_test_rebuild_event
  end type mci_test_t
  
\eatline
\nwendcode{}\nwbegindocs{157}Finalizer: base version is sufficient
\nwenddocs{}\nwbegincode{158}\moddef{MCI base: tests}\plusendmoddef
  subroutine mci_test_final (object)
    class(mci_test_t), intent(inout) :: object
    call object%base_final ()
  end subroutine mci_test_final
  
\eatline
\nwendcode{}\nwbegindocs{159}Output: trivial
\nwenddocs{}\nwbegincode{160}\moddef{MCI base: tests}\plusendmoddef
  subroutine mci_test_write (object, unit, pacify, md5sum_version)
    class(mci_test_t), intent(in) :: object
    integer, intent(in), optional :: unit
    logical, intent(in), optional :: pacify
    logical, intent(in), optional :: md5sum_version
    integer :: u
    u = given_output_unit (unit)
    write (u, "(1x,A)")  "Test integrator:"
    call object%base_write (u, pacify, md5sum_version)
    if (object%divisions /= 0) then
       write (u, "(3x,A,I0)")  "Number of divisions  = ", object%divisions
    end if
    if (allocated (object%rng))  call object%rng%write (u)
  end subroutine mci_test_write
  
\eatline
\nwendcode{}\nwbegindocs{161}Short version.
\nwenddocs{}\nwbegincode{162}\moddef{MCI base: tests}\plusendmoddef
  subroutine mci_test_startup_message (mci, unit, n_calls)
    class(mci_test_t), intent(in) :: mci
    integer, intent(in), optional :: unit, n_calls
    call mci%base_startup_message (unit = unit, n_calls = n_calls)
    write (msg_buffer, "(A,1x,I0,1x,A)") &
         "Integrator: Test:", mci%divisions, "divisions"
    call msg_message (unit = unit)
  end subroutine mci_test_startup_message
  
\eatline
\nwendcode{}\nwbegindocs{163}This is a no-op for the test integrator.
\nwenddocs{}\nwbegincode{164}\moddef{MCI base: tests}\plusendmoddef
  subroutine mci_test_ignore_flat_dimensions (mci, dim_flat)
    class(mci_test_t), intent(inout) :: mci
    integer, dimension(:), intent(in) :: dim_flat
  end subroutine mci_test_ignore_flat_dimensions
  
\eatline
\nwendcode{}\nwbegindocs{165}Ditto.
\nwenddocs{}\nwbegincode{166}\moddef{MCI base: tests}\plusendmoddef
  subroutine mci_test_ignore_equivalences (mci, channel, dim_offset)
    class(mci_test_t), intent(inout) :: mci
    type(phs_channel_t), dimension(:), intent(in) :: channel
    integer, intent(in) :: dim_offset
  end subroutine mci_test_ignore_equivalences
  
\eatline
\nwendcode{}\nwbegindocs{167}Set the number of divisions to a nonzero value.
\nwenddocs{}\nwbegincode{168}\moddef{MCI base: tests}\plusendmoddef
  subroutine mci_test_set_divisions (object, divisions)
    class(mci_test_t), intent(inout) :: object
    integer, intent(in) :: divisions
    object%divisions = divisions
  end subroutine mci_test_set_divisions
  
\eatline
\nwendcode{}\nwbegindocs{169}Set the maximum factor (default is 1).
\nwenddocs{}\nwbegincode{170}\moddef{MCI base: tests}\plusendmoddef
  subroutine mci_test_set_max_factor (object, max_factor)
    class(mci_test_t), intent(inout) :: object
    real(default), intent(in) :: max_factor
    object%max_factor = max_factor
  end subroutine mci_test_set_max_factor
  
\eatline
\nwendcode{}\nwbegindocs{171}Allocate instance with matching type.
\nwenddocs{}\nwbegincode{172}\moddef{MCI base: tests}\plusendmoddef
  subroutine mci_test_allocate_instance (mci, mci_instance)
    class(mci_test_t), intent(in) :: mci
    class(mci_instance_t), intent(out), pointer :: mci_instance
    allocate (mci_test_instance_t :: mci_instance)
  end subroutine mci_test_allocate_instance
  
\eatline
\nwendcode{}\nwbegindocs{173}Integrate: sample at the midpoints of uniform bits and add the results.
We implement this for one and for two dimensions.  In the latter case,
we scan over two channels and multiply with the channel weights.

The arguments {\tt{}n{\char95}it} and {\tt{}n{\char95}calls} are ignored in this implementations.

The test integrator does not set error or efficiency, so those will
remain undefined.
\nwenddocs{}\nwbegincode{174}\moddef{MCI base: tests}\plusendmoddef
  subroutine mci_test_integrate (mci, instance, sampler, &
       n_it, n_calls, results, pacify)
    class(mci_test_t), intent(inout) :: mci
    class(mci_instance_t), intent(inout) :: instance
    class(mci_sampler_t), intent(inout) :: sampler
    integer, intent(in) :: n_it
    integer, intent(in) :: n_calls
    logical, intent(in), optional :: pacify
    class(mci_results_t), intent(inout), optional :: results
    real(default), dimension(:), allocatable :: integral
    real(default), dimension(:), allocatable :: x
    integer :: i, j, c
    select type (instance)
    type is (mci_test_instance_t)
       allocate (integral (mci%n_channel))
       integral = 0
       allocate (x (mci%n_dim))
       select case (mci%n_dim)
       case (1)
          do c = 1, mci%n_channel
             do i = 1, mci%divisions
                x(1) = (i - 0.5_default) / mci%divisions
                call instance%evaluate (sampler, c, x)
                integral(c) = integral(c) + instance%get_value ()
             end do
          end do
          mci%integral = dot_product (instance%w, integral) &
               / mci%divisions
          mci%integral_known = .true.
       case (2)
          do c = 1, mci%n_channel
             do i = 1, mci%divisions
                x(1) = (i - 0.5_default) / mci%divisions
                do j = 1, mci%divisions
                   x(2) = (j - 0.5_default) / mci%divisions
                   call instance%evaluate (sampler, c, x)
                   integral(c) = integral(c) + instance%get_value ()
                end do
             end do
          end do
          mci%integral = dot_product (instance%w, integral) &
               / mci%divisions / mci%divisions
          mci%integral_known = .true.
       end select
       if (present (results)) then
          call results%record (n_it, n_calls, &
               mci%integral, mci%error, &
               efficiency = 0._default)
       end if
    end select
  end subroutine mci_test_integrate
  
\eatline
\nwendcode{}\nwbegindocs{175}Simulation initializer and finalizer: nothing to do here.
\nwenddocs{}\nwbegincode{176}\moddef{MCI base: tests}\plusendmoddef
  subroutine mci_test_ignore_prepare_simulation (mci)
    class(mci_test_t), intent(inout) :: mci
  end subroutine mci_test_ignore_prepare_simulation
  
\eatline
\nwendcode{}\nwbegindocs{177}Event generator.  We use mock random numbers for first selecting the
channel and then setting the $x$ values.  The results reside in the
state of {\tt{}instance} and {\tt{}sampler}.
\nwenddocs{}\nwbegincode{178}\moddef{MCI base: tests}\plusendmoddef
  subroutine mci_test_generate_weighted_event (mci, instance, sampler)
    class(mci_test_t), intent(inout) :: mci
    class(mci_instance_t), intent(inout), target :: instance
    class(mci_sampler_t), intent(inout), target :: sampler
    real(default) :: r
    real(default), dimension(:), allocatable :: x
    integer :: c
    select type (instance)
    type is (mci_test_instance_t)
       allocate (x (mci%n_dim))
       select case (mci%n_channel)
       case (1)
          c = 1
          call mci%rng%generate (x(1))
       case (2)
          call mci%rng%generate (r)
          if (r < instance%w(1)) then
             c = 1
          else
             c = 2
          end if
          call mci%rng%generate (x)
       end select
       call instance%evaluate (sampler, c, x)
    end select
  end subroutine mci_test_generate_weighted_event
       
\eatline
\nwendcode{}\nwbegindocs{179}For unweighted events, we generate weighted events and apply a
simple rejection step to the relative event weight, until an event passes.

(This might result in an endless loop if we happen to be in sync with
the mock random generator cycle.  Therefore, limit the number of tries.)
\nwenddocs{}\nwbegincode{180}\moddef{MCI base: tests}\plusendmoddef
  subroutine mci_test_generate_unweighted_event (mci, instance, sampler)
    class(mci_test_t), intent(inout) :: mci
    class(mci_instance_t), intent(inout), target :: instance
    class(mci_sampler_t), intent(inout), target :: sampler
    real(default) :: r
    integer :: i 
    select type (instance)
    type is (mci_test_instance_t)
       mci%tries = 0
       do i = 1, 10
          call mci%generate_weighted_event (instance, sampler)
          mci%tries = mci%tries + 1
          call mci%rng%generate (r)
          if (r < instance%rel_value)  exit
       end do
    end select
  end subroutine mci_test_generate_unweighted_event
    
\eatline
\nwendcode{}\nwbegindocs{181}Here, we rebuild the event from the state without consulting the rng.
\nwenddocs{}\nwbegincode{182}\moddef{MCI base: tests}\plusendmoddef
  subroutine mci_test_rebuild_event (mci, instance, sampler, state)
    class(mci_test_t), intent(inout) :: mci
    class(mci_instance_t), intent(inout) :: instance
    class(mci_sampler_t), intent(inout) :: sampler
    class(mci_state_t), intent(in) :: state
    select type (instance)
    type is (mci_test_instance_t)
       call instance%recall (sampler, state)
    end select
  end subroutine mci_test_rebuild_event
    
\eatline
\nwendcode{}\nwbegindocs{183}\nwdocspar
\subsubsection{Instance of the test MCI type}
This instance type simulates the VAMP approach.  We implement the VAMP
multi-channel formula, but keep the channel-specific probability
functions $g_i$ smooth and fixed.  We also keep the weights fixed.

The setup is as follows: we have $n$ mappings of the unit hypercube
\begin{equation}
  x = x (x^{(k)}) \qquad \text{where $x=(x_1,\ldots)$}.
\end{equation}
The Jacobian factors are the determinants
\begin{equation}
  f^{(k)}(x^{(k)}) = \left|\frac{\partial x}{\partial x^{(k)}}\right|
\end{equation}
We introduce arbitrary probability functions
\begin{equation}
  g^{(k)}(x^{(k)}) \qquad 
  \text{with}\quad \int dx^{(k)} g^{(k)}(x^{(k)}) = 1
\end{equation}
and weights
\begin{equation}
  w_k \qquad \text{with}\quad \sum_k w_k = 1
\end{equation}
and construct the joint probability function
\begin{equation}
  g(x) = \sum_k w_k\frac{g^{(k)}(x^{(k)}(x))}{f^{(k)}(x^{(k)}(x))}
\end{equation}
which also satisfies
\begin{equation}
  \int g(x)\,dx = 1.
\end{equation}
The algorithm implements a resolution of unity as follows
\begin{align}
  1 &= \int dx = \int\frac{g(x)}{g(x)} dx
  \nonumber\\
  &= \sum w_k \int \frac{g^{(k)}(x^{(k)}(x))}{f^{(k)}(x^{(k)}(x))}
  \,\frac{dx}{g(x)}
  \nonumber\\
  &= \sum w_k \int g^{(k)}(x^{(k)}) \frac{dx^{(k)}}{g(x(x^{(k)}))}
\end{align}
where each of the integrals in the sum is evaluated using the
channel-specific variables $x^{(k)}$.

We provide two examples: (1) trivial with one channel, one dimension,
and all functions unity and (2) two channels and two dimensions with
\begin{align}
  x (x^{(1)}) &= (x^{(1)}_1, x^{(1)}_2)
  \nonumber\\
  x (x^{(2)}) &= (x^{(2)}_1{}^2, x^{(2)}_2)
\end{align}
hence
\begin{align}
  f^{(1)}&\equiv 1,
  &f^{(2)}(x^{(2)}) &= 2x^{(2)}_1
\end{align}
The probability functions are
\begin{align}
  g^{(1)}&\equiv 1,
  &g^{(2)}(x^{(2)}) = 2 x^{(2)}_2
\end{align}

In the concrete implementation of the integrator instance we store
values for the channel probabilities $g_i$ and the accumulated
probability $g$.

We also store the result (product of integrand and MCI weight), the
expected maximum for the result in each channel.
\nwenddocs{}\nwbegincode{184}\moddef{MCI base: public}\plusendmoddef
  public :: mci_test_instance_t
\nwendcode{}\nwbegincode{185}\moddef{MCI base: test types}\plusendmoddef
  type, extends (mci_instance_t) :: mci_test_instance_t
     type(mci_test_t), pointer :: mci => null ()
     real(default) :: g = 0
     real(default), dimension(:), allocatable :: gi
     real(default) :: value = 0
     real(default) :: rel_value = 0
     real(default), dimension(:), allocatable :: max
   contains
     procedure :: write => mci_test_instance_write
     procedure :: final => mci_test_instance_final
     procedure :: init => mci_test_instance_init
     procedure :: compute_weight => mci_test_instance_compute_weight
     procedure :: record_integrand => mci_test_instance_record_integrand
     procedure :: init_simulation => mci_test_instance_init_simulation
     procedure :: final_simulation => mci_test_instance_final_simulation
     procedure :: get_event_excess => mci_test_instance_get_event_excess
  end type mci_test_instance_t

\eatline
\nwendcode{}\nwbegindocs{186}Output: trivial
\nwenddocs{}\nwbegincode{187}\moddef{MCI base: tests}\plusendmoddef
  subroutine mci_test_instance_write (object, unit, pacify)
    class(mci_test_instance_t), intent(in) :: object
    integer, intent(in), optional :: unit
    logical, intent(in), optional :: pacify
    integer :: u, c
    u = given_output_unit (unit)
    write (u, "(1x,A,ES13.7)") "Result value = ", object%value
    write (u, "(1x,A,ES13.7)") "Rel. weight  = ", object%rel_value
    write (u, "(1x,A,ES13.7)") "Integrand    = ", object%integrand
    write (u, "(1x,A,ES13.7)") "MCI weight   = ", object%mci_weight
    write (u, "(3x,A,I0)")  "c = ", object%selected_channel
    write (u, "(3x,A,ES13.7)") "g = ", object%g
    write (u, "(1x,A)")  "Channel parameters:"
    do c = 1, object%mci%n_channel
       write (u, "(1x,I0,A,4(1x,ES13.7))")  c, ": w/f/g/m =", &
            object%w(c), object%f(c), object%gi(c), object%max(c)
       write (u, "(4x,A,9(1x,F9.7))")  "x =", object%x(:,c)
    end do
  end subroutine mci_test_instance_write
  
\eatline
\nwendcode{}\nwbegindocs{188}The finalizer is empty.
\nwenddocs{}\nwbegincode{189}\moddef{MCI base: tests}\plusendmoddef
  subroutine mci_test_instance_final (object)
    class(mci_test_instance_t), intent(inout) :: object
  end subroutine mci_test_instance_final
  
\eatline
\nwendcode{}\nwbegindocs{190}Initializer.  We make use of the analytical result that the maximum of
the weighted integrand, in each channel, is equal to $1$
(one-dimensional case) and $2$ (two-dimensional case), respectively.
\nwenddocs{}\nwbegincode{191}\moddef{MCI base: tests}\plusendmoddef
  subroutine mci_test_instance_init (mci_instance, mci)
    class(mci_test_instance_t), intent(out) :: mci_instance
    class(mci_t), intent(in), target :: mci
    call mci_instance%base_init (mci)
    select type (mci)
    type is (mci_test_t)
       mci_instance%mci => mci
    end select
    allocate (mci_instance%gi (mci%n_channel))
    mci_instance%gi = 0
    allocate (mci_instance%max (mci%n_channel))
    select case (mci%n_channel)
    case (1)
       mci_instance%max = 1._default
    case (2)
       mci_instance%max = 2._default
    end select
  end subroutine mci_test_instance_init
    
\eatline
\nwendcode{}\nwbegindocs{192}Compute weight: we implement the VAMP multi-channel formula.  The channel
probabilities {\tt{}gi} are predefined functions.
\nwenddocs{}\nwbegincode{193}\moddef{MCI base: tests}\plusendmoddef
  subroutine mci_test_instance_compute_weight (mci, c)
    class(mci_test_instance_t), intent(inout) :: mci
    integer, intent(in) :: c
    integer :: i
    mci%selected_channel = c
    select case (mci%mci%n_dim)
    case (1)
       mci%gi(1) = 1
    case (2)
       mci%gi(1) = 1
       mci%gi(2) = 2 * mci%x(2,2)
    end select
    mci%g = 0
    do i = 1, mci%mci%n_channel
       mci%g = mci%g + mci%w(i) * mci%gi(i) / mci%f(i)
    end do
    mci%mci_weight = mci%gi(c) / mci%g
  end subroutine mci_test_instance_compute_weight
    
\eatline
\nwendcode{}\nwbegindocs{194}Record the integrand.  Apply the Jacobian weight to get the absolute value.
Divide by the channel maximum and by any overall factor to get the value
relative to the maximum.
\nwenddocs{}\nwbegincode{195}\moddef{MCI base: tests}\plusendmoddef
  subroutine mci_test_instance_record_integrand (mci, integrand)
    class(mci_test_instance_t), intent(inout) :: mci
    real(default), intent(in) :: integrand
    mci%integrand = integrand
    mci%value = mci%integrand * mci%mci_weight
    mci%rel_value = mci%value / mci%max(mci%selected_channel) &
         / mci%mci%max_factor
  end subroutine mci_test_instance_record_integrand
  
\eatline
\nwendcode{}\nwbegindocs{196}Nothing to do here.
\nwenddocs{}\nwbegincode{197}\moddef{MCI base: tests}\plusendmoddef
  subroutine mci_test_instance_init_simulation (instance, safety_factor)
    class(mci_test_instance_t), intent(inout) :: instance
    real(default), intent(in), optional :: safety_factor
  end subroutine mci_test_instance_init_simulation
  
  subroutine mci_test_instance_final_simulation (instance)
    class(mci_test_instance_t), intent(inout) :: instance
  end subroutine mci_test_instance_final_simulation
  
\eatline
\eatline
\nwendcode{}\nwbegindocs{198}Return always zero.
\nwenddocs{}\nwbegincode{199}\moddef{MCI base: tests}\plusendmoddef
  function mci_test_instance_get_event_excess (mci) result (excess)
    class(mci_test_instance_t), intent(in) :: mci
    real(default) :: excess
    excess = 0
  end function mci_test_instance_get_event_excess
  
\eatline
\nwendcode{}\nwbegindocs{200}\nwdocspar
\subsubsection{Test sampler}
The test sampler implements a fixed configuration, either trivial
(one-channel, one-dimension), or slightly nontrivial (two-channel,
two-dimension).  In the second channel, the first parameter is mapped
according to $x_1 = x^{(2)}_1{}^2$, so we have $f^{(2)}(x^{(2)}) =
2x^{(2)}_1$.

For display purposes, we store the return values inside the object.  This is
not strictly necessary.
\nwenddocs{}\nwbegincode{201}\moddef{MCI base: test types}\plusendmoddef
  type, extends (mci_sampler_t) :: test_sampler_t
     real(default) :: integrand = 0
     integer :: selected_channel = 0
     real(default), dimension(:,:), allocatable :: x
     real(default), dimension(:), allocatable :: f
   contains
     procedure :: init => test_sampler_init
     procedure :: write => test_sampler_write
     procedure :: compute => test_sampler_compute
     procedure :: is_valid => test_sampler_is_valid
     procedure :: evaluate => test_sampler_evaluate
     procedure :: rebuild => test_sampler_rebuild
     procedure :: fetch => test_sampler_fetch
  end type test_sampler_t
  
\eatline
\nwendcode{}\nwbegincode{202}\moddef{MCI base: tests}\plusendmoddef
  subroutine test_sampler_init (sampler, n)
    class(test_sampler_t), intent(out) :: sampler
    integer, intent(in) :: n
    allocate (sampler%x (n, n))
    allocate (sampler%f (n))
  end subroutine test_sampler_init
  
\eatline
\nwendcode{}\nwbegindocs{203}Output
\nwenddocs{}\nwbegincode{204}\moddef{MCI base: tests}\plusendmoddef
  subroutine test_sampler_write (object, unit, testflag)
    class(test_sampler_t), intent(in) :: object
    integer, intent(in), optional :: unit
    logical, intent(in), optional :: testflag
    integer :: u, c
    u = given_output_unit (unit)
    write (u, "(1x,A)")  "Test sampler:"
    write (u, "(3x,A,ES13.7)")  "Integrand = ", object%integrand
    write (u, "(3x,A,I0)")      "Channel   = ", object%selected_channel
    do c = 1, size (object%f)
       write (u, "(1x,I0,':',1x,A,ES13.7)") c, "f = ", object%f(c)
       write (u, "(4x,A,9(1x,F9.7))") "x =", object%x(:,c)
    end do
  end subroutine test_sampler_write
  
\eatline
\nwendcode{}\nwbegindocs{205}Compute $x$ and Jacobians, given the input parameter array.  This is called
both by {\tt{}evaluate} and {\tt{}rebuild}.
\nwenddocs{}\nwbegincode{206}\moddef{MCI base: tests}\plusendmoddef
  subroutine test_sampler_compute (sampler, c, x_in)
    class(test_sampler_t), intent(inout) :: sampler
    integer, intent(in) :: c
    real(default), dimension(:), intent(in) :: x_in
    sampler%selected_channel = c
    select case (size (sampler%f))
    case (1)
       sampler%x(:,1) = x_in
       sampler%f = 1
    case (2)
       select case (c)
       case (1)
          sampler%x(:,1) = x_in
          sampler%x(1,2) = sqrt (x_in(1))
          sampler%x(2,2) = x_in(2)
       case (2)
          sampler%x(1,1) = x_in(1) ** 2
          sampler%x(2,1) = x_in(2)
          sampler%x(:,2) = x_in
       end select
       sampler%f(1) = 1
       sampler%f(2) = 2 * sampler%x(1,2)
    end select
  end subroutine test_sampler_compute

\eatline
\nwendcode{}\nwbegindocs{207}The point is always valid.
\nwenddocs{}\nwbegincode{208}\moddef{MCI base: tests}\plusendmoddef
  function test_sampler_is_valid (sampler) result (valid)
    class(test_sampler_t), intent(in) :: sampler
    logical :: valid
    valid = .true.
  end function test_sampler_is_valid
  
\eatline
\nwendcode{}\nwbegindocs{209}The integrand is always equal to 1.
\nwenddocs{}\nwbegincode{210}\moddef{MCI base: tests}\plusendmoddef
  subroutine test_sampler_evaluate (sampler, c, x_in, val, x, f)
    class(test_sampler_t), intent(inout) :: sampler
    integer, intent(in) :: c
    real(default), dimension(:), intent(in) :: x_in
    real(default), intent(out) :: val
    real(default), dimension(:,:), intent(out) :: x
    real(default), dimension(:), intent(out) :: f
    call sampler%compute (c, x_in)
    sampler%integrand = 1
    val = sampler%integrand
    x = sampler%x
    f = sampler%f
  end subroutine test_sampler_evaluate

\eatline
\nwendcode{}\nwbegindocs{211}Construct kinematics from the input $x$ array.  Set the integrand
instead of evaluating it.
\nwenddocs{}\nwbegincode{212}\moddef{MCI base: tests}\plusendmoddef
  subroutine test_sampler_rebuild (sampler, c, x_in, val, x, f)
    class(test_sampler_t), intent(inout) :: sampler
    integer, intent(in) :: c
    real(default), dimension(:), intent(in) :: x_in
    real(default), intent(in) :: val
    real(default), dimension(:,:), intent(out) :: x
    real(default), dimension(:), intent(out) :: f
    call sampler%compute (c, x_in)
    sampler%integrand = val
    x = sampler%x
    f = sampler%f
  end subroutine test_sampler_rebuild

\eatline
\nwendcode{}\nwbegindocs{213}Recall contents.
\nwenddocs{}\nwbegincode{214}\moddef{MCI base: tests}\plusendmoddef
  subroutine test_sampler_fetch (sampler, val, x, f)
    class(test_sampler_t), intent(in) :: sampler
    real(default), intent(out) :: val
    real(default), dimension(:,:), intent(out) :: x
    real(default), dimension(:), intent(out) :: f
    val = sampler%integrand
    x = sampler%x
    f = sampler%f
  end subroutine test_sampler_fetch

\eatline
\nwendcode{}\nwbegindocs{215}\nwdocspar
\subsubsection{Test results object}
This mock object just stores and displays the current result.
\nwenddocs{}\nwbegincode{216}\moddef{MCI base: test types}\plusendmoddef
  type, extends (mci_results_t) :: mci_test_results_t
     integer :: n_it = 0
     integer :: n_calls = 0
     real(default) :: integral = 0
     real(default) :: error = 0
     real(default) :: efficiency = 0
   contains
   \LA{}MCI base: mci test results: TBP\RA{}
  end type mci_test_results_t
  
\eatline
\nwendcode{}\nwbegindocs{217}Output.
\nwenddocs{}\nwbegincode{218}\moddef{MCI base: mci test results: TBP}\endmoddef
  procedure :: write => mci_test_results_write
\nwendcode{}\nwbegincode{219}\moddef{MCI base: tests}\plusendmoddef
  subroutine mci_test_results_write (object, unit, verbose, suppress)
    class(mci_test_results_t), intent(in) :: object
    integer, intent(in), optional :: unit
    logical, intent(in), optional :: verbose, suppress
    integer :: u
    u = given_output_unit (unit)
    write (u, "(3x,A,1x,I0)") "Iterations = ", object%n_it
    write (u, "(3x,A,1x,I0)") "Calls      = ", object%n_calls
    write (u, "(3x,A,1x,F12.10)")  "Integral   = ", object%integral
    write (u, "(3x,A,1x,F12.10)")  "Error      = ", object%error
    write (u, "(3x,A,1x,F12.10)")  "Efficiency = ", object%efficiency
  end subroutine mci_test_results_write
  
\eatline
\nwendcode{}\nwbegindocs{220}Record result.
\nwenddocs{}\nwbegincode{221}\moddef{MCI base: mci test results: TBP}\plusendmoddef
  procedure :: record => mci_test_results_record
\nwendcode{}\nwbegincode{222}\moddef{MCI base: tests}\plusendmoddef
  subroutine mci_test_results_record (object, n_it, n_calls, &
       integral, error, efficiency, chain_weights, suppress)
    class(mci_test_results_t), intent(inout) :: object
    integer, intent(in) :: n_it
    integer, intent(in) :: n_calls
    real(default), intent(in) :: integral
    real(default), intent(in) :: error
    real(default), intent(in) :: efficiency
    real(default), dimension(:), intent(in), optional :: chain_weights
    logical, intent(in), optional :: suppress
    object%n_it = n_it
    object%n_calls = n_calls
    object%integral = integral
    object%error = error
    object%efficiency = efficiency
  end subroutine mci_test_results_record

\eatline
\nwendcode{}\nwbegindocs{223}\nwdocspar
\subsubsection{Integrator configuration data}
Construct and display a test integrator configuration object.
\nwenddocs{}\nwbegincode{224}\moddef{MCI base: execute tests}\endmoddef
  call test (mci_base_1, "mci_base_1", &
       "integrator configuration", &
       u, results)
\nwendcode{}\nwbegincode{225}\moddef{MCI base: tests}\plusendmoddef
  subroutine mci_base_1 (u)
    integer, intent(in) :: u
    class(mci_t), allocatable, target :: mci
    class(mci_instance_t), pointer :: mci_instance => null ()
    class(mci_sampler_t), allocatable :: sampler

    real(default) :: integrand
    
    write (u, "(A)")  "* Test output: mci_base_1"
    write (u, "(A)")  "*   Purpose: initialize and display &
         &test integrator"
    write (u, "(A)")
    
    write (u, "(A)")  "* Initialize integrator"
    write (u, "(A)")

    allocate (mci_test_t :: mci)
    call mci%set_dimensions (2, 2)
    
    call mci%write (u)
    
    write (u, "(A)")
    write (u, "(A)")  "* Initialize instance"
    write (u, "(A)")
    
    call mci%allocate_instance (mci_instance)
    call mci_instance%init (mci)
    
    write (u, "(A)")  "* Initialize test sampler"
    write (u, "(A)")
    
    allocate (test_sampler_t :: sampler)
    select type (sampler)
    type is (test_sampler_t)
       call sampler%init (2)
    end select

    write (u, "(A)")  "* Evaluate sampler for given point and channel"
    write (u, "(A)")
    
    call sampler%evaluate (1, [0.25_default, 0.8_default], &
         integrand, mci_instance%x, mci_instance%f)

    call sampler%write (u)
    
    write (u, "(A)")
    write (u, "(A)")  "* Compute MCI weight"
    write (u, "(A)")
    
    call mci_instance%compute_weight (1)
    call mci_instance%write (u)
    
    write (u, "(A)")
    write (u, "(A)")  "* Get integrand and compute weight for another point"
    write (u, "(A)")
    
    call mci_instance%evaluate (sampler, 2, [0.5_default, 0.6_default])
    call mci_instance%write (u)
    
    write (u, "(A)")
    write (u, "(A)")  "* Recall results, again"
    write (u, "(A)")
    
    call mci_instance%final ()
    deallocate (mci_instance)
    
    call mci%allocate_instance (mci_instance)
    call mci_instance%init (mci)

    call mci_instance%fetch (sampler, 2)
    call mci_instance%write (u)
    
    write (u, "(A)")
    write (u, "(A)")  "* Retrieve value"
    write (u, "(A)")
    
    write (u, "(1x,A,ES13.7)")  "Weighted integrand = ", &
         mci_instance%get_value ()

    call mci_instance%final ()
    call mci%final ()

    write (u, "(A)")
    write (u, "(A)")  "* Test output end: mci_base_1"

  end subroutine mci_base_1

\eatline
\nwendcode{}\nwbegindocs{226}\nwdocspar
\subsubsection{Trivial integral}
Use the MCI approach to compute a trivial one-dimensional integral.
\nwenddocs{}\nwbegincode{227}\moddef{MCI base: execute tests}\plusendmoddef
  call test (mci_base_2, "mci_base_2", &
       "integration", &
       u, results)
\nwendcode{}\nwbegincode{228}\moddef{MCI base: tests}\plusendmoddef
  subroutine mci_base_2 (u)
    integer, intent(in) :: u
    class(mci_t), allocatable, target :: mci
    class(mci_instance_t), pointer :: mci_instance => null ()
    class(mci_sampler_t), allocatable :: sampler
    
    write (u, "(A)")  "* Test output: mci_base_2"
    write (u, "(A)")  "*   Purpose: perform a test integral"
    write (u, "(A)")
    
    write (u, "(A)")  "* Initialize integrator"
    write (u, "(A)")

    allocate (mci_test_t :: mci)
    call mci%set_dimensions (1, 1)
    select type (mci)
    type is (mci_test_t)
       call mci%set_divisions (10)
    end select
    
    call mci%write (u)
    
    write (u, "(A)")
    write (u, "(A)")  "* Initialize instance"
    write (u, "(A)")
    
    call mci%allocate_instance (mci_instance)
    call mci_instance%init (mci)
    
    write (u, "(A)")  "* Initialize test sampler"
    write (u, "(A)")
    
    allocate (test_sampler_t :: sampler)
    select type (sampler)
    type is (test_sampler_t)
       call sampler%init (1)
    end select
    
    write (u, "(A)")  "* Integrate"
    write (u, "(A)")
    
    call mci%integrate (mci_instance, sampler, 0, 0)
    
    call mci%write (u)
    
    call mci_instance%final ()
    call mci%final ()

    write (u, "(A)")
    write (u, "(A)")  "* Test output end: mci_base_2"

  end subroutine mci_base_2

\eatline
\nwendcode{}\nwbegindocs{229}\nwdocspar
\subsubsection{Nontrivial integral}
Use the MCI approach to compute a simple two-dimensional integral with
two channels.
\nwenddocs{}\nwbegincode{230}\moddef{MCI base: execute tests}\plusendmoddef
  call test (mci_base_3, "mci_base_3", &
       "integration (two channels)", &
       u, results)
\nwendcode{}\nwbegincode{231}\moddef{MCI base: tests}\plusendmoddef
  subroutine mci_base_3 (u)
    integer, intent(in) :: u
    class(mci_t), allocatable, target :: mci
    class(mci_instance_t), pointer :: mci_instance => null ()
    class(mci_sampler_t), allocatable :: sampler
    
    write (u, "(A)")  "* Test output: mci_base_3"
    write (u, "(A)")  "*   Purpose: perform a nontrivial test integral"
    write (u, "(A)")
    
    write (u, "(A)")  "* Initialize integrator"
    write (u, "(A)")

    allocate (mci_test_t :: mci)
    call mci%set_dimensions (2, 2)
    select type (mci)
    type is (mci_test_t)
       call mci%set_divisions (10)
    end select

    write (u, "(A)")  "* Initialize instance"
    write (u, "(A)")
    
    call mci%allocate_instance (mci_instance)
    call mci_instance%init (mci)
    
    write (u, "(A)")  "* Initialize test sampler"
    write (u, "(A)")
    
    allocate (test_sampler_t :: sampler)
    select type (sampler)
    type is (test_sampler_t)
       call sampler%init (2)
    end select
    
    write (u, "(A)")  "* Integrate"
    write (u, "(A)")
    
    call mci%integrate (mci_instance, sampler, 0, 0)
    call mci%write (u)
    
    write (u, "(A)")
    write (u, "(A)")  "* Integrate with higher resolution"
    write (u, "(A)")

    select type (mci)
    type is (mci_test_t)
       call mci%set_divisions (100)
    end select

    call mci%integrate (mci_instance, sampler, 0, 0)
    call mci%write (u)

    call mci_instance%final ()
    call mci%final ()

    write (u, "(A)")
    write (u, "(A)")  "* Test output end: mci_base_3"

  end subroutine mci_base_3

\eatline
\nwendcode{}\nwbegindocs{232}\nwdocspar
\subsubsection{Event generation}
We generate ``random'' events, one weighted and one unweighted.  The
test implementation does not require an integration pass, we can
generate events immediately.
\nwenddocs{}\nwbegincode{233}\moddef{MCI base: execute tests}\plusendmoddef
  call test (mci_base_4, "mci_base_4", &
       "event generation (two channels)", &
       u, results)
\nwendcode{}\nwbegincode{234}\moddef{MCI base: tests}\plusendmoddef
  subroutine mci_base_4 (u)
    integer, intent(in) :: u
    class(mci_t), allocatable, target :: mci
    class(mci_instance_t), pointer :: mci_instance => null ()
    class(mci_sampler_t), allocatable :: sampler
    class(rng_t), allocatable :: rng
    
    write (u, "(A)")  "* Test output: mci_base_4"
    write (u, "(A)")  "*   Purpose: generate events"
    write (u, "(A)")
    
    write (u, "(A)")  "* Initialize integrator, instance, sampler"
    write (u, "(A)")

    allocate (mci_test_t :: mci)
    call mci%set_dimensions (2, 2)
    select type (mci)
    type is (mci_test_t)
       call mci%set_divisions (10)
    end select
    
    call mci%allocate_instance (mci_instance)
    call mci_instance%init (mci)
    
    allocate (test_sampler_t :: sampler)
    select type (sampler)
    type is (test_sampler_t)
       call sampler%init (2)
    end select
    
    allocate (rng_test_t :: rng)
    call mci%import_rng (rng)
    
    write (u, "(A)")  "* Generate weighted event"
    write (u, "(A)")
    
    call mci%generate_weighted_event (mci_instance, sampler)

    call sampler%write (u)
    write (u, *)
    call mci_instance%write (u)
    
    write (u, "(A)")
    write (u, "(A)")  "* Generate unweighted event"
    write (u, "(A)")
    
    call mci%generate_unweighted_event (mci_instance, sampler)

    select type (mci)
    type is (mci_test_t)
       write (u, "(A,I0)")  " Success in try ", mci%tries
       write (u, "(A)")
    end select
    
    call sampler%write (u)
    write (u, *)
    call mci_instance%write (u)
    
    call mci_instance%final ()
    call mci%final ()

    write (u, "(A)")
    write (u, "(A)")  "* Test output end: mci_base_4"

  end subroutine mci_base_4

\eatline
\nwendcode{}\nwbegindocs{235}\nwdocspar
\subsubsection{Store and recall data}
We generate an event and store the relevant data, i.e., the input
parameters and the result value for a particular channel.  Then we use
those data to recover the event, as far as the MCI record is concerned.
\nwenddocs{}\nwbegincode{236}\moddef{MCI base: execute tests}\plusendmoddef
  call test (mci_base_5, "mci_base_5", &
       "store and recall", &
       u, results)
\nwendcode{}\nwbegincode{237}\moddef{MCI base: tests}\plusendmoddef
  subroutine mci_base_5 (u)
    integer, intent(in) :: u
    class(mci_t), allocatable, target :: mci
    class(mci_instance_t), pointer :: mci_instance => null ()
    class(mci_sampler_t), allocatable :: sampler
    class(rng_t), allocatable :: rng
    class(mci_state_t), allocatable :: state
    
    write (u, "(A)")  "* Test output: mci_base_5"
    write (u, "(A)")  "*   Purpose: store and recall an event"
    write (u, "(A)")
    
    write (u, "(A)")  "* Initialize integrator, instance, sampler"
    write (u, "(A)")

    allocate (mci_test_t :: mci)
    call mci%set_dimensions (2, 2)
    select type (mci)
    type is (mci_test_t)
       call mci%set_divisions (10)
    end select
    
    call mci%allocate_instance (mci_instance)
    call mci_instance%init (mci)
    
    allocate (test_sampler_t :: sampler)
    select type (sampler)
    type is (test_sampler_t)
       call sampler%init (2)
    end select
    
    allocate (rng_test_t :: rng)
    call mci%import_rng (rng)
    
    write (u, "(A)")  "* Generate weighted event"
    write (u, "(A)")
    
    call mci%generate_weighted_event (mci_instance, sampler)

    call sampler%write (u)
    write (u, *)
    call mci_instance%write (u)
    
    write (u, "(A)")
    write (u, "(A)")  "* Store data"
    write (u, "(A)")
    
    allocate (state)
    call mci_instance%store (state)
    call mci_instance%final ()
    deallocate (mci_instance)
    
    call state%write (u)
    
    write (u, "(A)")
    write (u, "(A)")  "* Recall data and rebuild event"
    write (u, "(A)")
    
    call mci%allocate_instance (mci_instance)
    call mci_instance%init (mci)
    call mci%rebuild_event (mci_instance, sampler, state)

    call sampler%write (u)
    write (u, *)
    call mci_instance%write (u)
    
    call mci_instance%final ()
    call mci%final ()

    write (u, "(A)")
    write (u, "(A)")  "* Test output end: mci_base_5"

  end subroutine mci_base_5

\eatline
\nwendcode{}\nwbegindocs{238}\nwdocspar
\subsubsection{Chained channels}
Chain channels together.  In the base configuration, this just fills entries
in an extra array (each channel may belong to a chain).  In type
implementations, this will be used for grouping equivalent channels by keeping
their weights equal.
\nwenddocs{}\nwbegincode{239}\moddef{MCI base: execute tests}\plusendmoddef
  call test (mci_base_6, "mci_base_6", &
       "chained channels", &
       u, results)
\nwendcode{}\nwbegincode{240}\moddef{MCI base: tests}\plusendmoddef
  subroutine mci_base_6 (u)
    integer, intent(in) :: u
    class(mci_t), allocatable, target :: mci
    
    write (u, "(A)")  "* Test output: mci_base_6"
    write (u, "(A)")  "*   Purpose: initialize and display &
         &test integrator with chains"
    write (u, "(A)")
    
    write (u, "(A)")  "* Initialize integrator"
    write (u, "(A)")

    allocate (mci_test_t :: mci)
    call mci%set_dimensions (1, 5)
    
    write (u, "(A)")  "* Introduce chains"
    write (u, "(A)")
    
    call mci%declare_chains ([1, 2, 2, 1, 2])

    call mci%write (u)
    
    call mci%final ()

    write (u, "(A)")
    write (u, "(A)")  "* Test output end: mci_base_6"

  end subroutine mci_base_6

\eatline
\nwendcode{}\nwbegindocs{241}\nwdocspar
\subsubsection{Recording results}
Compute a simple two-dimensional integral and record the result.
\nwenddocs{}\nwbegincode{242}\moddef{MCI base: execute tests}\plusendmoddef
  call test (mci_base_7, "mci_base_7", &
       "recording results", &
       u, results)
\nwendcode{}\nwbegincode{243}\moddef{MCI base: tests}\plusendmoddef
  subroutine mci_base_7 (u)
    integer, intent(in) :: u
    class(mci_t), allocatable, target :: mci
    class(mci_instance_t), pointer :: mci_instance => null ()
    class(mci_sampler_t), allocatable :: sampler
    class(mci_results_t), allocatable :: results
    
    write (u, "(A)")  "* Test output: mci_base_7"
    write (u, "(A)")  "*   Purpose: perform a nontrivial test integral &
         &and record results"
    write (u, "(A)")
    
    write (u, "(A)")  "* Initialize integrator"
    write (u, "(A)")

    allocate (mci_test_t :: mci)
    call mci%set_dimensions (2, 2)
    select type (mci)
    type is (mci_test_t)
       call mci%set_divisions (10)
    end select

    write (u, "(A)")  "* Initialize instance"
    write (u, "(A)")
    
    call mci%allocate_instance (mci_instance)
    call mci_instance%init (mci)
    
    write (u, "(A)")  "* Initialize test sampler"
    write (u, "(A)")
    
    allocate (test_sampler_t :: sampler)
    select type (sampler)
    type is (test_sampler_t)
       call sampler%init (2)
    end select
    
    allocate (mci_test_results_t :: results)

    write (u, "(A)")  "* Integrate"
    write (u, "(A)")
    
    call mci%integrate (mci_instance, sampler, 1, 1000, results)
    call mci%write (u)
    
    write (u, "(A)")
    write (u, "(A)")  "* Display results"
    write (u, "(A)")

    call results%write (u)
    
    call mci_instance%final ()
    call mci%final ()

    write (u, "(A)")
    write (u, "(A)")  "* Test output end: mci_base_7"

  end subroutine mci_base_7

\eatline
\nwendcode{}\nwbegindocs{244}\nwdocspar
\subsubsection{Timer}
Simple checks for the embedded timer.
\nwenddocs{}\nwbegincode{245}\moddef{MCI base: execute tests}\plusendmoddef
  call test (mci_base_8, "mci_base_8", &
       "timer", &
       u, results)
\nwendcode{}\nwbegincode{246}\moddef{MCI base: tests}\plusendmoddef
  subroutine mci_base_8 (u)
    integer, intent(in) :: u
    class(mci_t), allocatable, target :: mci
    
    write (u, "(A)")  "* Test output: mci_base_8"
    write (u, "(A)")  "*   Purpose: check timer availability"
    write (u, "(A)")
    
    write (u, "(A)")  "* Initialize integrator with timer"
    write (u, "(A)")

    allocate (mci_test_t :: mci)
    call mci%set_dimensions (2, 2)
    select type (mci)
    type is (mci_test_t)
       call mci%set_divisions (10)
    end select

    call mci%set_timer (active = .true.)
    call mci%write (u)
    
    write (u, "(A)")
    write (u, "(A)")  "* Start timer"
    write (u, "(A)")

    call mci%start_timer ()
    call mci%write (u)

    write (u, "(A)")
    write (u, "(A)")  "* Stop timer"
    write (u, "(A)")

    call mci%stop_timer ()
    call mci%write (u)

    write (u, "(A)")
    write (u, "(A)")  "* Readout"
    write (u, "(A)")

    write (u, "(1x,A,F6.3)")  "Time = ", mci%get_time ()
    
    write (u, "(A)")
    write (u, "(A)")  "* Deactivate timer"
    write (u, "(A)")

    call mci%set_timer (active = .false.)
    call mci%write (u)

    write (u, "(A)")
    write (u, "(A)")  "* Cleanup"

    call mci%final ()

    write (u, "(A)")
    write (u, "(A)")  "* Test output end: mci_base_8"

  end subroutine mci_base_8

\eatline
\nwendcode{}\nwbegindocs{247}\nwdocspar
%%%%%%%%%%%%%%%%%%%%%%%%%%%%%%%%%%%%%%%%%%%%%%%%%%%%%%%%%%%%%%%%%%%%%%%%
\section{Iterations}
This module defines a container for the list of iterations and calls, to be
submitted to integration.
\nwenddocs{}\nwbegincode{248}\moddef{\code{}iterations.f90\edoc{}}\endmoddef
\LA{}File header\RA{}

module iterations

\LA{}Use strings\RA{}
  use io_units
  use unit_tests
  use diagnostics

\LA{}Standard module head\RA{}

\LA{}Iterations: public\RA{}

\LA{}Iterations: types\RA{}

contains

\LA{}Iterations: procedures\RA{}

\LA{}Iterations: tests\RA{}

end module iterations
\eatline
\nwendcode{}\nwbegindocs{249}\nwdocspar
\subsection{The iterations list}
Each integration pass has a number of iterations and a number of calls
per iteration.  The last pass produces the end result; the previous
passes are used for adaptation.

The flags {\tt{}adapt{\char95}grid} and {\tt{}adapt{\char95}weight} are used only if
{\tt{}custom{\char95}adaptation} is set.  Otherwise, default settings are used
that depend on the integration pass.
\nwenddocs{}\nwbegincode{250}\moddef{Iterations: types}\endmoddef
  type :: iterations_spec_t
     private
     integer :: n_it = 0
     integer :: n_calls = 0
     logical :: custom_adaptation = .false.
     logical :: adapt_grids = .false.
     logical :: adapt_weights = .false.
  end type iterations_spec_t

\eatline
\nwendcode{}\nwbegindocs{251}We build up a list of iterations.
\nwenddocs{}\nwbegincode{252}\moddef{Iterations: public}\endmoddef
  public :: iterations_list_t
\nwendcode{}\nwbegincode{253}\moddef{Iterations: types}\plusendmoddef
  type :: iterations_list_t
     private
     integer :: n_pass = 0
     type(iterations_spec_t), dimension(:), allocatable :: pass
   contains
   \LA{}Iterations: iterations list: TBP\RA{}
  end type iterations_list_t
     
\eatline
\nwendcode{}\nwbegindocs{254}Initialize an iterations list.  For each pass, we have to specify
the number of iterations and calls.  We may provide the adaption
conventions explicitly, either as character codes or as logicals.  

For passes where the adaptation conventions are not specified, we use
the following default setting: adapt weights and grids for all passes
except the last one.
\nwenddocs{}\nwbegincode{255}\moddef{Iterations: iterations list: TBP}\endmoddef
  procedure :: init => iterations_list_init
\nwendcode{}\nwbegincode{256}\moddef{Iterations: procedures}\endmoddef
  subroutine iterations_list_init &
       (it_list, n_it, n_calls, adapt, adapt_code, adapt_grids, adapt_weights)
    class(iterations_list_t), intent(inout) :: it_list
    integer, dimension(:), intent(in) :: n_it, n_calls
    logical, dimension(:), intent(in), optional :: adapt
    type(string_t), dimension(:), intent(in), optional :: adapt_code
    logical, dimension(:), intent(in), optional :: adapt_grids, adapt_weights
    integer :: i
    it_list%n_pass = size (n_it)
    if (allocated (it_list%pass)) deallocate (it_list%pass)    
    allocate (it_list%pass (it_list%n_pass))
    it_list%pass%n_it = n_it
    it_list%pass%n_calls = n_calls
    if (present (adapt)) then
       it_list%pass%custom_adaptation = adapt
       do i = 1, it_list%n_pass
          if (adapt(i)) then
             if (verify (adapt_code(i), "wg") /= 0) then
                call msg_error ("iteration specification: " &
                     // "adaptation code letters must be 'w' or 'g'")
             end if
             it_list%pass(i)%adapt_grids = scan (adapt_code(i), "g") /= 0
             it_list%pass(i)%adapt_weights = scan (adapt_code(i), "w") /= 0
          end if
       end do
    else if (present (adapt_grids) .and. present (adapt_weights)) then
       it_list%pass%custom_adaptation = .true.
       it_list%pass%adapt_grids = adapt_grids
       it_list%pass%adapt_weights = adapt_weights
    end if
    do i = 1, it_list%n_pass - 1
       if (.not. it_list%pass(i)%custom_adaptation) then
          it_list%pass(i)%adapt_grids = .true.
          it_list%pass(i)%adapt_weights = .true.
       end if
    end do
  end subroutine iterations_list_init

\eatline
\nwendcode{}\nwbegincode{257}\moddef{Iterations: iterations list: TBP}\plusendmoddef
  procedure :: clear => iterations_list_clear
\nwendcode{}\nwbegincode{258}\moddef{Iterations: procedures}\plusendmoddef
  subroutine iterations_list_clear (it_list)
    class(iterations_list_t), intent(inout) :: it_list
    it_list%n_pass = 0
    deallocate (it_list%pass)
  end subroutine iterations_list_clear

\eatline
\nwendcode{}\nwbegindocs{259}Write the list of iterations.
\nwenddocs{}\nwbegincode{260}\moddef{Iterations: iterations list: TBP}\plusendmoddef
  procedure :: write => iterations_list_write
\nwendcode{}\nwbegincode{261}\moddef{Iterations: procedures}\plusendmoddef
  subroutine iterations_list_write (it_list, unit)
    class(iterations_list_t), intent(in) :: it_list
    integer, intent(in), optional :: unit
    integer :: u
    u = given_output_unit (unit)
    write (u, "(A)")  char (it_list%to_string ())
  end subroutine iterations_list_write

\eatline
\nwendcode{}\nwbegindocs{262}The output as a single-line string.
\nwenddocs{}\nwbegincode{263}\moddef{Iterations: iterations list: TBP}\plusendmoddef
  procedure :: to_string => iterations_list_to_string
\nwendcode{}\nwbegincode{264}\moddef{Iterations: procedures}\plusendmoddef
  function iterations_list_to_string (it_list) result (buffer)
    class(iterations_list_t), intent(in) :: it_list
    type(string_t) :: buffer
    character(30) :: ibuf
    integer :: i
    buffer = "iterations = "
    if (it_list%n_pass > 0) then
       do i = 1, it_list%n_pass
          if (i > 1)  buffer = buffer // ", "
          write (ibuf, "(I0,':',I0)") &
               it_list%pass(i)%n_it, it_list%pass(i)%n_calls
          buffer = buffer // trim (ibuf)
          if (it_list%pass(i)%custom_adaptation &
               .or. it_list%pass(i)%adapt_grids &
               .or. it_list%pass(i)%adapt_weights) then
             buffer = buffer // ':"'
             if (it_list%pass(i)%adapt_grids)  buffer = buffer // "g"
             if (it_list%pass(i)%adapt_weights)  buffer = buffer // "w"
             buffer = buffer // '"'
          end if          
       end do
    else
       buffer = buffer // "[undefined]"
    end if
  end function iterations_list_to_string
    
\eatline
\nwendcode{}\nwbegindocs{265}\nwdocspar
\subsection{Tools}
Return the total number of passes.
\nwenddocs{}\nwbegincode{266}\moddef{Iterations: iterations list: TBP}\plusendmoddef
  procedure :: get_n_pass => iterations_list_get_n_pass
\nwendcode{}\nwbegincode{267}\moddef{Iterations: procedures}\plusendmoddef
  function iterations_list_get_n_pass (it_list) result (n_pass)
    class(iterations_list_t), intent(in) :: it_list
    integer :: n_pass
    n_pass = it_list%n_pass
  end function iterations_list_get_n_pass

\eatline
\nwendcode{}\nwbegindocs{268}Return the number of calls for a specific pass.
\nwenddocs{}\nwbegincode{269}\moddef{Iterations: iterations list: TBP}\plusendmoddef
  procedure :: get_n_calls => iterations_list_get_n_calls
\nwendcode{}\nwbegincode{270}\moddef{Iterations: procedures}\plusendmoddef
  function iterations_list_get_n_calls (it_list, pass) result (n_calls)
    class(iterations_list_t), intent(in) :: it_list
    integer :: n_calls
    integer, intent(in) :: pass
    if (pass <= it_list%n_pass) then
       n_calls = it_list%pass(pass)%n_calls
    else
       n_calls = 0
    end if
  end function iterations_list_get_n_calls

\eatline
\nwendcode{}\nwbegindocs{271}Get the adaptation mode (automatic/custom) and, for custom adaptation, the
flags for a specific pass.
\nwenddocs{}\nwbegincode{272}\moddef{Iterations: iterations list: TBP}\plusendmoddef
  procedure :: adapt_grids => iterations_list_adapt_grids
  procedure :: adapt_weights => iterations_list_adapt_weights
\nwendcode{}\nwbegincode{273}\moddef{Iterations: procedures}\plusendmoddef
  function iterations_list_adapt_grids (it_list, pass) result (flag)
    logical :: flag
    class(iterations_list_t), intent(in) :: it_list
    integer, intent(in) :: pass
    if (pass <= it_list%n_pass) then
       flag = it_list%pass(pass)%adapt_grids
    else
       flag = .false.
    end if
  end function iterations_list_adapt_grids

  function iterations_list_adapt_weights (it_list, pass) result (flag)
    logical :: flag
    class(iterations_list_t), intent(in) :: it_list
    integer, intent(in) :: pass
    if (pass <= it_list%n_pass) then
       flag = it_list%pass(pass)%adapt_weights
    else
       flag = .false.
    end if
  end function iterations_list_adapt_weights

\eatline
\eatline
\eatline
\nwendcode{}\nwbegindocs{274}Return the total number of iterations / the iterations for a specific pass.
\nwenddocs{}\nwbegincode{275}\moddef{Iterations: iterations list: TBP}\plusendmoddef
  procedure :: get_n_it => iterations_list_get_n_it
\nwendcode{}\nwbegincode{276}\moddef{Iterations: procedures}\plusendmoddef
  function iterations_list_get_n_it (it_list, pass) result (n_it)
    class(iterations_list_t), intent(in) :: it_list
    integer :: n_it
    integer, intent(in) :: pass
    if (pass <= it_list%n_pass) then
       n_it = it_list%pass(pass)%n_it
    else
       n_it = 0
    end if
  end function iterations_list_get_n_it

\eatline
\nwendcode{}\nwbegindocs{277}\nwdocspar
\subsection{Test}
This is the master for calling self-test procedures.
\nwenddocs{}\nwbegincode{278}\moddef{Iterations: public}\plusendmoddef
  public :: iterations_test
\nwendcode{}\nwbegincode{279}\moddef{Iterations: tests}\endmoddef
  subroutine iterations_test (u, results)
    integer, intent(in) :: u
    type(test_results_t), intent(inout) :: results
  \LA{}Iterations: execute tests\RA{}
end subroutine iterations_test

\eatline
\nwendcode{}\nwbegindocs{280}\nwdocspar
\subsubsection{Empty list}
\nwenddocs{}\nwbegincode{281}\moddef{Iterations: execute tests}\endmoddef
  call test (iterations_1, "iterations_1", &
       "empty iterations list", &
       u, results)
\nwendcode{}\nwbegincode{282}\moddef{Iterations: tests}\plusendmoddef
  subroutine iterations_1 (u)
    integer, intent(in) :: u
    type(iterations_list_t) :: it_list
    
    write (u, "(A)")  "* Test output: iterations_1"
    write (u, "(A)")  "*   Purpose: display empty iterations list"
    write (u, "(A)")

    call it_list%write (u)
    
    write (u, "(A)")
    write (u, "(A)")  "* Test output end: iterations_1"
    
  end subroutine iterations_1
  
\eatline
\nwendcode{}\nwbegindocs{283}\nwdocspar
\subsubsection{Fill list}
\nwenddocs{}\nwbegincode{284}\moddef{Iterations: execute tests}\plusendmoddef
  call test (iterations_2, "iterations_2", &
       "create iterations list", &
       u, results)
\nwendcode{}\nwbegincode{285}\moddef{Iterations: tests}\plusendmoddef
  subroutine iterations_2 (u)
    integer, intent(in) :: u
    type(iterations_list_t) :: it_list
    
    write (u, "(A)")  "* Test output: iterations_2"
    write (u, "(A)")  "*   Purpose: fill and display iterations list"
    write (u, "(A)")

    write (u, "(A)")  "* Minimal setup (2 passes)"
    write (u, "(A)")

    call it_list%init ([2, 4], [5000, 20000])

    call it_list%write (u)
    call it_list%clear ()
    
    write (u, "(A)")
    write (u, "(A)")  "* Setup with flags (3 passes)"
    write (u, "(A)")

    call it_list%init ([2, 4, 5], [5000, 20000, 400], &
         [.false., .true., .true.], &
         [var_str (""), var_str ("g"), var_str ("wg")])

    call it_list%write (u)
    
    write (u, "(A)")
    write (u, "(A)")  "* Extract data"
    write (u, "(A)")
    
    write (u, "(A,I0)")  "n_pass = ", it_list%get_n_pass ()
    write (u, "(A)")
    write (u, "(A,I0)")  "n_calls(2) = ", it_list%get_n_calls (2)
    write (u, "(A)")
    write (u, "(A,I0)")  "n_it(3) = ", it_list%get_n_it (3)

    write (u, "(A)")
    write (u, "(A)")  "* Test output end: iterations_2"
    
  end subroutine iterations_2
  
\eatline
\nwendcode{}\nwbegindocs{286}\nwdocspar
%%%%%%%%%%%%%%%%%%%%%%%%%%%%%%%%%%%%%%%%%%%%%%%%%%%%%%%%%%%%%%%%%%%%%%%%
\section{Integration results}

We record integration results and errors in a dedicated type.  This
allows us to do further statistics such as weighted average,
chi-squared, grouping by integration passes, etc.

Note WHIZARD 2.2.0: This code is taken from the previous {\tt{}processes}
module essentially unchanged and converted into a separate module.  It
lacks an overhaul and, in particular, self-tests.
\nwenddocs{}\nwbegincode{287}\moddef{\code{}integration{\char95}results.f90\edoc{}}\endmoddef
module integration_results
  
\LA{}Use kinds\RA{}
\LA{}Use strings\RA{}
  use io_units
  use format_utils, only: mp_format, pac_fmt
  use format_defs, only: FMT_10, FMT_14
  use diagnostics
  use md5
  use os_interface
  use mci_base
  
\LA{}Standard module head\RA{}

\LA{}Integration results: public\RA{}

\LA{}Integration results: parameters\RA{}

\LA{}Integration results: types\RA{}

contains
  
\LA{}Integration results: procedures\RA{}

end module integration_results
\eatline
\nwendcode{}\nwbegindocs{288}\nwdocspar
\subsection{Integration results entry}
This object collects the results of an integration pass and makes them
available to the outside.

The results object has to distinguish the process type:

We store the process type, the index of the integration pass and the
absolute iteration index, the number of iterations contained in this
result (for averages), and the integral (cross section or partial
width), error estimate, efficiency.

For intermediate results, we set a flag if this result is an
improvement w.r.t. previous ones.

The process type indicates decay or scattering.  Dummy entries
(skipped iterations) have a process type of {\tt{}PRC{\char95}UNKNOWN}.
\nwenddocs{}\nwbegincode{289}\moddef{Integration results: types}\endmoddef
  type :: integration_entry_t
     private
     integer :: process_type = PRC_UNKNOWN
     integer :: pass = 0
     integer :: it = 0
     integer :: n_it = 0
     integer :: n_calls = 0
     logical :: improved = .false.
     real(default) :: integral = 0
     real(default) :: error = 0
     real(default) :: efficiency = 0
     real(default) :: chi2 = 0
     real(default), dimension(:), allocatable :: chain_weights
  end type integration_entry_t

\eatline
\nwendcode{}\nwbegindocs{290}\nwdocspar
The possible values of the type indicator:
\nwenddocs{}\nwbegincode{291}\moddef{Integration results: parameters}\endmoddef
  integer, parameter, public :: PRC_UNKNOWN = 0
  integer, parameter, public :: PRC_DECAY = 1
  integer, parameter, public :: PRC_SCATTERING = 2

\eatline
\nwendcode{}\nwbegindocs{292}Initialize with all relevant data
\nwenddocs{}\nwbegincode{293}\moddef{Integration results: procedures}\endmoddef
  subroutine integration_entry_init (entry, &
       process_type, pass, it, n_it, n_calls, improved, &
       integral, error, efficiency, chi2, chain_weights)
    type(integration_entry_t), intent(out) :: entry
    integer, intent(in) :: process_type, pass, it, n_it, n_calls
    logical, intent(in) :: improved
    real(default), intent(in) :: integral, error, efficiency
    real(default), intent(in), optional :: chi2
    real(default), dimension(:), intent(in), optional :: chain_weights
    entry%process_type = process_type
    entry%pass = pass
    entry%it = it
    entry%n_it = n_it
    entry%n_calls = n_calls
    entry%improved = improved
    entry%integral = integral
    entry%error = error
    entry%efficiency = efficiency
    if (present (chi2)) &
         entry%chi2 = chi2
    if (present (chain_weights)) then
       allocate (entry%chain_weights (size (chain_weights)))
       entry%chain_weights = chain_weights
    end if
  end subroutine integration_entry_init

\eatline
\nwendcode{}\nwbegindocs{294}Access values, some of them computed on demand:
\nwenddocs{}\nwbegincode{295}\moddef{Integration results: procedures}\plusendmoddef
  elemental function integration_entry_get_pass (entry) result (n)
    integer :: n
    type(integration_entry_t), intent(in) :: entry
    n = entry%pass
  end function integration_entry_get_pass

  elemental function integration_entry_get_n_calls (entry) result (n)
    integer :: n
    type(integration_entry_t), intent(in) :: entry
    n = entry%n_calls
  end function integration_entry_get_n_calls

  elemental function integration_entry_get_integral (entry) result (int)
    real(default) :: int
    type(integration_entry_t), intent(in) :: entry
    int = entry%integral
  end function integration_entry_get_integral

  elemental function integration_entry_get_error (entry) result (err)
    real(default) :: err
    type(integration_entry_t), intent(in) :: entry
    err = entry%error
  end function integration_entry_get_error

  elemental function integration_entry_get_relative_error (entry) result (err)
    real(default) :: err
    type(integration_entry_t), intent(in) :: entry
    if (entry%integral /= 0) then
       err = entry%error / entry%integral
    else
       err = 0
    end if
  end function integration_entry_get_relative_error

  elemental function integration_entry_get_accuracy (entry) result (acc)
    real(default) :: acc
    type(integration_entry_t), intent(in) :: entry
    acc = accuracy (entry%integral, entry%error, entry%n_calls)
  end function integration_entry_get_accuracy

  elemental function accuracy (integral, error, n_calls) result (acc)
    real(default) :: acc
    real(default), intent(in) :: integral, error
    integer, intent(in) :: n_calls
    if (integral /= 0) then
       acc = error / integral * sqrt (real (n_calls, default))
    else
       acc = 0
    end if
  end function accuracy

  elemental function integration_entry_get_efficiency (entry) result (eff)
    real(default) :: eff
    type(integration_entry_t), intent(in) :: entry
    eff = entry%efficiency
  end function integration_entry_get_efficiency

  elemental function integration_entry_get_chi2 (entry) result (chi2)
    real(default) :: chi2
    type(integration_entry_t), intent(in) :: entry
    chi2 = entry%chi2
  end function integration_entry_get_chi2

  elemental function integration_entry_has_improved (entry) result (flag)
    logical :: flag
    type(integration_entry_t), intent(in) :: entry
    flag = entry%improved
  end function integration_entry_has_improved

  elemental function integration_entry_get_n_groves (entry) result (n_groves)
    integer :: n_groves
    type(integration_entry_t), intent(in) :: entry
    if (allocated (entry%chain_weights)) then
       n_groves = size (entry%chain_weights, 1)
    else
       n_groves = 0
    end if
  end function integration_entry_get_n_groves

\eatline
\eatline
\eatline
\eatline
\eatline
\eatline
\eatline
\eatline
\eatline
\eatline
\nwendcode{}\nwbegindocs{296}Output. This writes the header line for the result account below:
\nwenddocs{}\nwbegincode{297}\moddef{Integration results: procedures}\plusendmoddef
  subroutine write_header (process_type, unit, logfile)
    integer, intent(in) :: process_type
    integer, intent(in), optional :: unit
    logical, intent(in), optional :: logfile
    character(5) :: phys_unit
    integer :: u
    u = given_output_unit (unit);  if (u < 0)  return
    select case (process_type)
    case (PRC_DECAY);      phys_unit = "[GeV]"
    case (PRC_SCATTERING); phys_unit = "[fb] "
    case default
       phys_unit = ""
    end select
    write (msg_buffer, "(A)") &
         "It      Calls  Integral" // phys_unit // &
         " Error" // phys_unit // &
         "  Err[%]    Acc  Eff[%]   Chi2 N[It] |"
    call msg_message (unit=u, logfile=logfile)
  end subroutine write_header
       
\eatline
\nwendcode{}\nwbegindocs{298}This writes a separator for result display:
\nwenddocs{}\nwbegincode{299}\moddef{Integration results: procedures}\plusendmoddef
  subroutine write_hline (unit)
    integer, intent(in), optional :: unit
    integer :: u
    u = given_output_unit (unit);  if (u < 0)  return
    write (u, "(A)")  "|" // (repeat ("-", 77)) // "|"
    flush (u)
  end subroutine write_hline
  
  subroutine write_dline (unit)
    integer, intent(in), optional :: unit
    integer :: u
    u = given_output_unit (unit);  if (u < 0)  return
    write (u, "(A)")  "|" // (repeat ("=", 77)) // "|"
    flush (u)
  end subroutine write_dline
  
\eatline
\eatline
\nwendcode{}\nwbegindocs{300}This writes the standard result account into one screen line.  The
verbose version uses multiple lines and prints the unabridged values.
Dummy entries are not written.
\nwenddocs{}\nwbegincode{301}\moddef{Integration results: procedures}\plusendmoddef
  subroutine integration_entry_write (entry, unit, verbose, suppress)
    type(integration_entry_t), intent(in) :: entry
    integer, intent(in), optional :: unit
    logical, intent(in), optional :: verbose
    logical, intent(in), optional :: suppress   
    integer :: u
    character(1) :: star
    character(12) :: fmt
    character(7) :: fmt2
    logical :: verb, supp
    u = given_output_unit (unit);  if (u < 0)  return
    verb = .false.;  if (present (verbose))  verb = verbose
    supp = .false.;  if (present (suppress)) supp = suppress 
    if (verb)  then
       write (u, *)  "process_type = ", entry%process_type
       write (u, *)  "        pass = ", entry%pass
       write (u, *)  "          it = ", entry%it
       write (u, *)  "        n_it = ", entry%n_it
       write (u, *)  "     n_calls = ", entry%n_calls
       write (u, *)  "    improved = ", entry%improved
       write (u, *)  "    integral = ", entry%integral
       write (u, *)  "       error = ", entry%error
       write (u, *)  "  efficiency = ", entry%efficiency
       write (u, *)  "        chi2 = ", entry%chi2
       if (allocated (entry%chain_weights)) then
          write (u, *)  "    n_groves = ", size (entry%chain_weights)
          write (u, *)  "chain_weights = ", entry%chain_weights
       else
          write (u, *)  "    n_groves = 0"
       end if
    else if (entry%process_type /= PRC_UNKNOWN) then
       if (entry%improved .and. .not. supp) then
          star = "*"
       else
          star = " "
       end if
       call pac_fmt (fmt, FMT_14, "3x," // FMT_10 // ",1x", suppress)
       call pac_fmt (fmt2, "1x,F6.2", "2x,F5.1", suppress)
       if (entry%n_it /= 1) then
          write (u, "(1x,I3,1x,I10,1x," // fmt // ",1x,ES9.2,1x,F7.2," // &
            "1x,F7.2,A1," // fmt2 // ",1x,F7.2,1x,I3)") &
               entry%it, &
               entry%n_calls, &
               entry%integral, &
               abs(entry%error), &
               abs(integration_entry_get_relative_error (entry)) * 100, &
               abs(integration_entry_get_accuracy (entry)), &
               star, &
               entry%efficiency * 100, &
               entry%chi2, &
               entry%n_it
       else
          write (u, "(1x,I3,1x,I10,1x," // fmt // ",1x,ES9.2,1x,F7.2," // &
            "1x,F7.2,A1," // fmt2 // ",1x,F7.2,1x,I3)") &          
               entry%it, &
               entry%n_calls, &
               entry%integral, &
               abs(entry%error), &
               abs(integration_entry_get_relative_error (entry)) * 100, &
               abs(integration_entry_get_accuracy (entry)), &
               star, &
               entry%efficiency * 100
       end if
    end if
    flush (u)
  end subroutine integration_entry_write

\eatline
\nwendcode{}\nwbegindocs{302}Read the entry, assuming it has been written in verbose format.
\nwenddocs{}\nwbegincode{303}\moddef{Integration results: procedures}\plusendmoddef
  subroutine integration_entry_read (entry, unit)
    type(integration_entry_t), intent(out) :: entry
    integer, intent(in) :: unit
    character(30) :: dummy
    character :: equals
    integer :: n_groves
    read (unit, *)  dummy, equals, entry%process_type
    read (unit, *)  dummy, equals, entry%pass
    read (unit, *)  dummy, equals, entry%it
    read (unit, *)  dummy, equals, entry%n_it
    read (unit, *)  dummy, equals, entry%n_calls
    read (unit, *)  dummy, equals, entry%improved
    read (unit, *)  dummy, equals, entry%integral
    read (unit, *)  dummy, equals, entry%error
    read (unit, *)  dummy, equals, entry%efficiency
    read (unit, *)  dummy, equals, entry%chi2
    read (unit, *)  dummy, equals, n_groves
    if (n_groves /= 0) then
       allocate (entry%chain_weights (n_groves))
       read (unit, *)  dummy, equals, entry%chain_weights
    end if
  end subroutine integration_entry_read
    
\eatline
\nwendcode{}\nwbegindocs{304}Write an account of the channel weights, accumulated by groves.
\nwenddocs{}\nwbegincode{305}\moddef{Integration results: procedures}\plusendmoddef
  subroutine integration_entry_write_chain_weights (entry, unit)
    type(integration_entry_t), intent(in) :: entry
    integer, intent(in), optional :: unit
    integer :: u, i
    u = given_output_unit (unit);  if (u < 0)  return    
    if (allocated (entry%chain_weights)) then
       do i = 1, size (entry%chain_weights)
          write (u, "(1x,I3)", advance="no")  nint (entry%chain_weights(i) * 100)
       end do
       write (u, *)
    end if
  end subroutine integration_entry_write_chain_weights

\eatline
\nwendcode{}\nwbegindocs{306}Compute the average for all entries in the specified integration
pass.  The integrals are weighted w.r.t.\ their individual errors.

The quoted error of the result is the expected error, computed from
the weighted average of the given individual errors.

This should be compared to the actual distribution of the results,
from which we also can compute an error estimate if there is more than
one iteration.  The ratio of the distribution error and the averaged
error, is the $\chi^2$ value.

All error distributions are assumed Gaussian, of course.  The $\chi^2$
value is a partial check for this assumption.  If it is significantly
greater than unity, there is something wrong with the individual errors.

The efficiency returned is the one of the last entry in the
integration pass. 

If any error vanishes, averaging by this algorithm would fail.  In this case,
we simply average the entries and use the deviations from this average (if
any) to estimate the error.
\nwenddocs{}\nwbegincode{307}\moddef{Integration results: procedures}\plusendmoddef
  function compute_average (entry, pass) result (result)
    type(integration_entry_t) :: result
    type(integration_entry_t), dimension(:), intent(in) :: entry
    integer, intent(in) :: pass
    integer :: i
    logical, dimension(size(entry)) :: mask
    real(default), dimension(size(entry)) :: ivar
    real(default) :: sum_ivar, variance
    result%process_type = entry(1)%process_type
    result%pass = pass
    mask = entry%pass == pass .and. entry%process_type /= PRC_UNKNOWN
    result%it = maxval (entry%it, mask)
    result%n_it = count (mask)
    result%n_calls = sum (entry%n_calls, mask)
    if (.not. any (mask .and. entry%error == 0)) then
       where (mask)
          ivar = 1 / entry%error ** 2
       elsewhere
          ivar = 0
       end where
       sum_ivar = sum (ivar, mask)
       if (sum_ivar /= 0) then
          variance = 1 / sum_ivar
       else
          variance = 0
       end if
       result%integral = sum (entry%integral * ivar, mask) * variance
       if (result%n_it > 1) then
          result%chi2 = &
               sum ((entry%integral - result%integral)**2 * ivar, mask) &
               / (result%n_it - 1)
       end if
    else if (result%n_it /= 0) then
       result%integral = sum (entry%integral, mask) / result%n_it
       if (result%n_it > 1) then
          variance = &
               sum ((entry%integral - result%integral)**2, mask) &
               / (result%n_it - 1)
          if (result%integral /= 0) then
             if (abs (variance / result%integral) &
                  < 100 * epsilon (1._default)) then
                variance = 0
             end if
          end if
          result%chi2 = variance / result%n_it
       else
          variance = 0
       end if
    end if
    result%error = sqrt (variance)
    do i = size (entry), 1, -1
       if (mask(i)) then
          result%efficiency = entry(i)%efficiency
          exit
       end if
    end do
  end function compute_average

\eatline
\nwendcode{}\nwbegindocs{308}\nwdocspar
\subsection{Combined integration results}
We collect a list of results which grows during the execution of the
program.  This is implemented as an array which grows if necessary; so
we can easily compute averages.

We implement this as an extension of the {\tt{}mci{\char95}results{\char95}t} which is
defined in {\tt{}mci{\char95}base} as an abstract type.  We thus decouple the
implementation of the integrator from the implementation of the
results display, but nevertheless can record intermediate results
during integration.  This implies that the present extension
implements a {\tt{}record} method.
\nwenddocs{}\nwbegincode{309}\moddef{Integration results: public}\endmoddef
  public :: integration_results_t
\nwendcode{}\nwbegincode{310}\moddef{Integration results: types}\plusendmoddef
  type, extends (mci_results_t) :: integration_results_t
     private
     integer :: process_type = PRC_UNKNOWN
     integer :: current_pass = 0
     integer :: n_pass = 0
     integer :: n_it = 0
     logical :: screen = .false.
     integer :: unit = 0
     real(default) :: error_threshold = 0
     type(integration_entry_t), dimension(:), allocatable :: entry
     type(integration_entry_t), dimension(:), allocatable :: average
   contains
   \LA{}Integration results: integration results: TBP\RA{}
  end type integration_results_t

\eatline
\nwendcode{}\nwbegindocs{311}The array is extended in chunks of 10 entries.
\nwenddocs{}\nwbegincode{312}\moddef{Integration results: parameters}\plusendmoddef
  integer, parameter :: RESULTS_CHUNK_SIZE = 10

\eatline
\nwendcode{}\nwbegindocs{313}The standard does not require to explicitly initialize the integers;
however, some gfortran version has a bug here and misses the default
initialization in the type definition.
\nwenddocs{}\nwbegincode{314}\moddef{Integration results: integration results: TBP}\endmoddef
  procedure :: init => integration_results_init
\nwendcode{}\nwbegincode{315}\moddef{Integration results: procedures}\plusendmoddef
  subroutine integration_results_init (results, process_type)
    class(integration_results_t), intent(out) :: results
    integer, intent(in) :: process_type
    results%process_type = process_type
    results%n_pass = 0
    results%n_it = 0
    allocate (results%entry (RESULTS_CHUNK_SIZE))
    allocate (results%average (RESULTS_CHUNK_SIZE))
  end subroutine integration_results_init

\eatline
\nwendcode{}\nwbegindocs{316}Set addititional parameters: the {\tt{}error{\char95}threshold} declares that any error
value (in absolute numbers) smaller than this is to be considered zero.
\nwenddocs{}\nwbegincode{317}\moddef{Integration results: integration results: TBP}\plusendmoddef
  procedure :: set_error_threshold => integration_results_set_error_threshold
\nwendcode{}\nwbegincode{318}\moddef{Integration results: procedures}\plusendmoddef
  subroutine integration_results_set_error_threshold (results, error_threshold)
    class(integration_results_t), intent(inout) :: results
    real(default), intent(in) :: error_threshold
    results%error_threshold = error_threshold
  end subroutine integration_results_set_error_threshold
  
\eatline
\nwendcode{}\nwbegindocs{319}Output (ASCII format).  The {\tt{}verbose} format is used for writing
the header in grid files.
\nwenddocs{}\nwbegincode{320}\moddef{Integration results: integration results: TBP}\plusendmoddef
  procedure :: write => integration_results_write
\nwendcode{}\nwbegincode{321}\moddef{Integration results: procedures}\plusendmoddef
  subroutine integration_results_write (object, unit, verbose, suppress)
    class(integration_results_t), intent(in) :: object
    integer, intent(in), optional :: unit
    logical, intent(in), optional :: verbose
    logical, intent(in), optional :: suppress
    logical :: verb
    integer :: u, n
    u = given_output_unit (unit);  if (u < 0)  return
    verb = .false.;  if (present (verbose))  verb = verbose
    if (.not. verb) then
       call write_dline (unit)
       if (object%n_it /= 0) then
          call write_header (object%entry(1)%process_type, unit, &
               logfile=.false.)
          call write_dline (unit)
          do n = 1, object%n_it
             if (n > 1) then
                if (object%entry(n)%pass /= object%entry(n-1)%pass) then
                   call write_hline (unit)
                   call integration_entry_write &
                        (object%average(object%entry(n-1)%pass), &
                           unit, suppress = suppress)
                   call write_hline (unit)
                end if
             end if
             call integration_entry_write (object%entry(n), unit, &
                     suppress = suppress)
          end do
          call write_hline(unit)
          call integration_entry_write (object%average(object%n_pass), &
                   unit, suppress = suppress)
       else
          call msg_message ("[WHIZARD integration results: empty]", unit)
       end if
       call write_dline (unit)
    else
       write (u, *)  "begin(integration_results)"
       write (u, *)  "  n_pass = ", object%n_pass
       write (u, *)  "    n_it = ", object%n_it
       if (object%n_it > 0) then
          write (u, *)  "begin(integration_pass)"
          do n = 1, object%n_it
             if (n > 1) then
                if (object%entry(n)%pass /= object%entry(n-1)%pass) then
                   write (u, *)  "end(integration_pass)"
                   write (u, *)  "begin(integration_pass)"
                end if
             end if
             write (u, *)  "begin(iteration)"
             call integration_entry_write (object%entry(n), unit, &
                      verbose = verb, suppress = suppress)
             write (u, *)  "end(iteration)"
          end do
          write (u, *)  "end(integration_pass)"
       end if
       write (u, *)  "end(integration_results)"
    end if
    flush (u)
  end subroutine integration_results_write

\eatline
\nwendcode{}\nwbegindocs{322}During integration, we do not want to print all results at once, but
each intermediate result as soon as we get it.  Thus, the previous procedure
is chopped in pieces.  First piece: store the
output unit and a flag whether we want to print to standard output as
well.  Then write the header if the results are still empty, i.e.,
before integration has started.  The second piece writes a single
result to the saved output channels.  We call this from the {\tt{}record}
method, which can be called from the integrator directly.  The third
piece writes the average result, once a pass has been completed.  The
fourth piece writes a footer (if any), assuming that this is the final result.
\nwenddocs{}\nwbegincode{323}\moddef{Integration results: integration results: TBP}\plusendmoddef
  procedure :: display_init => integration_results_display_init
  procedure :: display_current => integration_results_display_current
  procedure :: display_pass => integration_results_display_pass
  procedure :: display_final => integration_results_display_final
\nwendcode{}\nwbegincode{324}\moddef{Integration results: procedures}\plusendmoddef
  subroutine integration_results_display_init &
       (results, process_type, screen, unit)
    class(integration_results_t), intent(inout) :: results
    integer, intent(in) :: process_type
    logical, intent(in) :: screen
    integer, intent(in), optional :: unit
    integer :: u
    if (present (unit))  results%unit = unit
    u = given_output_unit ()
    results%screen = screen
    if (results%n_it == 0) then
       if (results%screen) then
          call write_dline (u)
          call write_header (process_type, u, &
               logfile=.false.)
          call write_dline (u)
       end if
       if (results%unit /= 0) then
          call write_dline (results%unit)
          call write_header (process_type, results%unit, &
               logfile=.false.)
          call write_dline (results%unit)
       end if
    else
       if (results%screen) then
          call write_hline (u)
       end if
       if (results%unit /= 0) then
          call write_hline (results%unit)
       end if
    end if
  end subroutine integration_results_display_init
  
  subroutine integration_results_display_current (results, pacify)
    class(integration_results_t), intent(in) :: results
    integer :: u
    logical, intent(in), optional :: pacify
    u = given_output_unit ()
    if (results%screen) then
       call integration_entry_write (results%entry(results%n_it), u, &
            suppress = pacify)
    end if
    if (results%unit /= 0) then
       call integration_entry_write (results%entry(results%n_it), &
            results%unit, suppress = pacify)
    end if
  end subroutine integration_results_display_current

  subroutine integration_results_display_pass (results, pacify)
    class(integration_results_t), intent(in) :: results
    logical, intent(in), optional :: pacify
    integer :: u
    u = given_output_unit ()
    if (results%screen) then
       call write_hline (u)
       call integration_entry_write &
            (results%average(results%entry(results%n_it)%pass), &
                u, suppress = pacify)
    end if
    if (results%unit /= 0) then
       call write_hline (results%unit)
       call integration_entry_write &
            (results%average(results%entry(results%n_it)%pass), &
                results%unit, suppress = pacify)
    end if
  end subroutine integration_results_display_pass

  subroutine integration_results_display_final (results)
    class(integration_results_t), intent(inout) :: results
    integer :: u
    u = given_output_unit ()
    if (results%screen) then
       call write_dline (u)
    end if
    if (results%unit /= 0) then
       call write_dline (results%unit)
    end if
    results%screen = .false.
    results%unit = 0
  end subroutine integration_results_display_final

\eatline
\eatline
\eatline
\nwendcode{}\nwbegindocs{325}Write a concise table of chain weights, i.e., the channel history where
channels are collected by chains.
\nwenddocs{}\nwbegincode{326}\moddef{Integration results: integration results: TBP}\plusendmoddef
  procedure :: write_chain_weights => &
       integration_results_write_chain_weights 
\nwendcode{}\nwbegincode{327}\moddef{Integration results: procedures}\plusendmoddef
  subroutine integration_results_write_chain_weights (results, unit)
    class(integration_results_t), intent(in) :: results
    integer, intent(in), optional :: unit
    integer :: u, i, n
    u = given_output_unit (unit);  if (u < 0)  return
    if (allocated (results%entry(1)%chain_weights) .and. results%n_it /= 0) then
       call msg_message ("Phase-space chain (grove) weight history: " &
            // "(numbers in %)", unit)
       write (u, "(A9)", advance="no")  "| chain |"
          do i = 1, integration_entry_get_n_groves (results%entry(1))
          write (u, "(1x,I3)", advance="no")  i
       end do
       write (u, *)
       call write_dline (unit)
       do n = 1, results%n_it
          if (n > 1) then
             if (results%entry(n)%pass /= results%entry(n-1)%pass) then
                call write_hline (unit)
             end if
          end if
          write (u, "(1x,I6,1x,A1)", advance="no")  n, "|"
          call integration_entry_write_chain_weights (results%entry(n), unit)
       end do
       flush (u)
       call write_dline(unit)
    end if
  end subroutine integration_results_write_chain_weights

\eatline
\nwendcode{}\nwbegindocs{328}Read the list from file.  The file must be written using the
{\tt{}verbose} option of the writing routine.
\nwenddocs{}\nwbegincode{329}\moddef{Integration results: procedures}\plusendmoddef
  subroutine integration_results_read (results, unit)
    type(integration_results_t), intent(out) :: results
    integer, intent(in) :: unit
    character(80) :: buffer
    character :: equals
    integer :: pass, it
    read (unit, *)  buffer
    if (trim (adjustl (buffer)) /= "begin(integration_results)") then
       call read_err ();  return
    end if
    read (unit, *)  buffer, equals, results%n_pass
    read (unit, *)  buffer, equals, results%n_it
    allocate (results%entry (results%n_it + RESULTS_CHUNK_SIZE))
    allocate (results%average (results%n_it + RESULTS_CHUNK_SIZE))
    it = 0
    do pass = 1, results%n_pass
       read (unit, *)  buffer
       if (trim (adjustl (buffer)) /= "begin(integration_pass)") then
          call read_err ();  return
       end if
       READ_ENTRIES: do
          read (unit, *)  buffer
          if (trim (adjustl (buffer)) /= "begin(iteration)") then
             exit READ_ENTRIES
          end if
          it = it + 1
          call integration_entry_read (results%entry(it), unit)
          read (unit, *)  buffer
          if (trim (adjustl (buffer)) /= "end(iteration)") then
             call read_err (); return
          end if
       end do READ_ENTRIES
       if (trim (adjustl (buffer)) /= "end(integration_pass)") then
          call read_err (); return
       end if
       results%average(pass) = compute_average (results%entry, pass)
    end do
    read (unit, *)  buffer
    if (trim (adjustl (buffer)) /= "end(integration_results)") then
       call read_err (); return
    end if
  contains
    subroutine read_err ()
      call msg_fatal ("Reading integration results from file: syntax error")
    end subroutine read_err
  end subroutine integration_results_read

\eatline
\nwendcode{}\nwbegindocs{330}Check integration results for consistency.  We compare against an
array of pass indices and call numbers.  If there is a difference, up
to the number of iterations done so far, we return failure.  Dummy
entries (where {\tt{}pass} = 0) are ignored.
\nwenddocs{}\nwbegincode{331}\moddef{Integration results: procedures}\plusendmoddef
  function integration_results_iterations_are_consistent &
       (results, pass, n_calls) result (flag)
    logical :: flag
    type(integration_results_t), intent(in) :: results
    integer, dimension(:), intent(in) :: pass, n_calls
    integer :: n_it
    n_it = results%n_it
    flag = size (pass) >= n_it .and. size (n_calls) >= n_it
    if (flag) then
       flag = all (results%entry(:n_it)%pass == pass(:n_it) &
                   .and. &
                   (results%entry(:n_it)%n_calls == n_calls(:n_it) &
                    .or. &
                    results%entry(:n_it)%process_type == PRC_UNKNOWN))
    end if
  end function integration_results_iterations_are_consistent

\eatline
\nwendcode{}\nwbegindocs{332}Discard all results starting from the specified iteration.
\nwenddocs{}\nwbegincode{333}\moddef{Integration results: procedures}\plusendmoddef
  subroutine integration_results_discard (results, it)
    type(integration_results_t), intent(inout) :: results
    integer, intent(in) :: it
    if (it <= results%n_it) then
       select case (it)
       case (:1)
          results%n_it = 0
          results%n_pass = 0
          results%current_pass = 0
       case default
          results%n_it = it - 1
          results%n_pass = maxval (results%entry(1:results%n_it)%pass)
          results%current_pass = results%n_pass
       end select
    end if
  end subroutine integration_results_discard

\eatline
\nwendcode{}\nwbegindocs{334}Expand the list of entries if the limit has been reached:
\nwenddocs{}\nwbegincode{335}\moddef{Integration results: integration results: TBP}\plusendmoddef
  procedure :: expand => integration_results_expand
\nwendcode{}\nwbegincode{336}\moddef{Integration results: procedures}\plusendmoddef
  subroutine integration_results_expand (results)
    class(integration_results_t), intent(inout) :: results
    type(integration_entry_t), dimension(:), allocatable :: entry_tmp
    if (results%n_it == size (results%entry)) then
       allocate (entry_tmp (results%n_it))
       entry_tmp = results%entry
       deallocate (results%entry)
       allocate (results%entry (results%n_it + RESULTS_CHUNK_SIZE))
       results%entry(:results%n_it) = entry_tmp
       deallocate (entry_tmp)
    end if
    if (results%n_pass == size (results%average)) then
       allocate (entry_tmp (results%n_pass))
       entry_tmp = results%average
       deallocate (results%average)
       allocate (results%average (results%n_it + RESULTS_CHUNK_SIZE))
       results%average(:results%n_pass) = entry_tmp
       deallocate (entry_tmp)
    end if
  end subroutine integration_results_expand

\eatline
\nwendcode{}\nwbegindocs{337}Increment the {\tt{}current{\char95}pass} counter.  Can be done before
integration; after integration, the recording method may use the value
of this counter to define the entry.
\nwenddocs{}\nwbegincode{338}\moddef{Integration results: integration results: TBP}\plusendmoddef
  procedure :: new_pass => integration_results_new_pass
\nwendcode{}\nwbegincode{339}\moddef{Integration results: procedures}\plusendmoddef
  subroutine integration_results_new_pass (results)
    class(integration_results_t), intent(inout) :: results
    results%current_pass = results%current_pass + 1
  end subroutine integration_results_new_pass
  
\eatline
\nwendcode{}\nwbegindocs{340}Append a new entry to the list and, if appropriate, compute the average.
\nwenddocs{}\nwbegincode{341}\moddef{Integration results: integration results: TBP}\plusendmoddef
  procedure :: append_entry => integration_results_append_entry
\nwendcode{}\nwbegincode{342}\moddef{Integration results: procedures}\plusendmoddef
  subroutine integration_results_append_entry (results, entry)
    class(integration_results_t), intent(inout) :: results
    type(integration_entry_t), intent(in), optional :: entry
    if (results%n_it == 0) then
       results%n_it = 1
       results%n_pass = 1
    else
       call results%expand ()
       if (present (entry)) then
          if (entry%pass /= results%entry(results%n_it)%pass) &
               results%n_pass = results%n_pass + 1
       end if
       results%n_it = results%n_it + 1
    end if
    if (present (entry)) then
       results%entry(results%n_it) = entry
       results%average(results%n_pass) = &
            compute_average (results%entry, entry%pass)
    end if
  end subroutine integration_results_append_entry

\eatline
\nwendcode{}\nwbegindocs{343}Enter results into the results list.  For the error value, we may compare
them with a given threshold.  This guards against numerical noise, if the
exact error would be zero.
\nwenddocs{}\nwbegincode{344}\moddef{Integration results: integration results: TBP}\plusendmoddef
  procedure :: append => integration_results_append
\nwendcode{}\nwbegincode{345}\moddef{Integration results: procedures}\plusendmoddef
  subroutine integration_results_append (results, &
       n_it, n_calls, &
       integral, error, efficiency, &      
       chain_weights)
    class(integration_results_t), intent(inout) :: results
    integer, intent(in) :: n_it, n_calls
    real(default), intent(in) :: integral, error, efficiency
    real(default), dimension(:), intent(in), optional :: chain_weights
    logical :: improved
    type(integration_entry_t) :: entry
    real(default) :: err_checked
    if (results%n_it /= 0) then
       improved = abs(accuracy (integral, error, n_calls)) &
            < abs(integration_entry_get_accuracy (results%entry(results%n_it)))
    else
       improved = .true.
    end if
    if (abs (error) >= results%error_threshold) then
       err_checked = error
    else
       err_checked = 0
    end if
    call integration_entry_init (entry, &
         results%process_type, results%current_pass, &
         results%n_it+1, n_it, n_calls, improved, & 
         integral, err_checked, efficiency, &
         chain_weights=chain_weights)
    call results%append_entry (entry)
  end subroutine integration_results_append
         
\eatline
\nwendcode{}\nwbegindocs{346}Enter an empty result into the results list.
\nwenddocs{}\nwbegincode{347}\moddef{Integration results: public}\plusendmoddef
  public :: integration_results_append_null
\nwendcode{}\nwbegincode{348}\moddef{Integration results: procedures}\plusendmoddef
  subroutine integration_results_append_null (results, pass, n_it)
    type(integration_results_t), intent(inout) :: results
    integer, intent(in) :: pass, n_it
    type(integration_entry_t) :: entry
    call integration_entry_init (entry, &
         PRC_UNKNOWN, results%current_pass, n_it, 1, 0, .false., &
         0._default, 0._default, 0._default)
    call results%append_entry (entry)
  end subroutine integration_results_append_null
         
\eatline
\nwendcode{}\nwbegindocs{349}Record an integration pass executed by an {\tt{}mci} integrator
object.

There is a tolerance below we treat an error (relative to the
integral) as zero.
\nwenddocs{}\nwbegincode{350}\moddef{Integration results: parameters}\plusendmoddef
  real(default), parameter, public :: INTEGRATION_ERROR_TOLERANCE = 1e-10
\eatline
\nwendcode{}\nwbegindocs{351}\nwdocspar
\nwenddocs{}\nwbegincode{352}\moddef{Integration results: integration results: TBP}\plusendmoddef
  procedure :: record => integration_results_record
\nwendcode{}\nwbegincode{353}\moddef{Integration results: procedures}\plusendmoddef
  subroutine integration_results_record &
       (object, n_it, n_calls, integral, error, efficiency, &
        chain_weights, suppress)
    class(integration_results_t), intent(inout) :: object
    integer, intent(in) :: n_it, n_calls
    real(default), intent(in) :: integral, error, efficiency
    real(default), dimension(:), intent(in), optional :: chain_weights    
    real(default) :: err
    logical, intent(in), optional :: suppress
    
    if (abs (error) >= abs (integral) * INTEGRATION_ERROR_TOLERANCE) then
       err = error
    else
       err = 0
    end if
    call object%append (n_it, n_calls, integral, err, efficiency, chain_weights)
    call object%display_current (suppress)
  end subroutine integration_results_record
    
\eatline
\nwendcode{}\nwbegindocs{354}\nwdocspar
\subsection{Access results}
Return true if the results object has entries.
\nwenddocs{}\nwbegincode{355}\moddef{Integration results: integration results: TBP}\plusendmoddef
  procedure :: exist => integration_results_exist
\nwendcode{}\nwbegincode{356}\moddef{Integration results: procedures}\plusendmoddef
  function integration_results_exist (results) result (flag)
    logical :: flag
    class(integration_results_t), intent(in) :: results
    flag = results%n_pass > 0
  end function integration_results_exist

\eatline
\nwendcode{}\nwbegindocs{357}Retrieve information from the results record.  If {\tt{}last} is set and
true, take the last iteration.  If {\tt{}it} is set instead, take this
iteration.  If {\tt{}pass} is set, take this average.  If none is set,
take the final average.

If the result would be invalid, the entry is not assigned.  Due to
default initialization, this returns a null entry.
\nwenddocs{}\nwbegincode{358}\moddef{Integration results: integration results: TBP}\plusendmoddef
  procedure :: get_entry => results_get_entry
\nwendcode{}\nwbegincode{359}\moddef{Integration results: procedures}\plusendmoddef
  function results_get_entry (results, last, it, pass) result (entry)
    class(integration_results_t), intent(in) :: results
    type(integration_entry_t) :: entry
    logical, intent(in), optional :: last
    integer, intent(in), optional :: it, pass
    if (present (last)) then
       if (allocated (results%entry) .and. results%n_it > 0) then
          entry = results%entry(results%n_it)
       else
          call error ()
       end if
    else if (present (it)) then
       if (allocated (results%entry) .and. it > 0 .and. it <= results%n_it) then
          entry = results%entry(it)
       else
          call error ()
       end if
    else if (present (pass)) then
       if (allocated (results%average) &
            .and. pass > 0 .and. pass <= results%n_pass) then
          entry = results%average (pass)
       else
          call error ()
       end if
    else
       if (allocated (results%average) .and. results%n_pass > 0) then
          entry = results%average (results%n_pass)
       else
          call error ()
       end if
    end if
  contains
    subroutine error ()
      call msg_fatal ("Requested integration result is not available")
    end subroutine error
  end function results_get_entry
  
\eatline
\nwendcode{}\nwbegindocs{360}The individual procedures.  The {\tt{}results} record should have the
{\tt{}target} attribute, but only locally within the function.
\nwenddocs{}\nwbegincode{361}\moddef{Integration results: integration results: TBP}\plusendmoddef
  procedure :: get_n_calls => integration_results_get_n_calls
  procedure :: get_integral => integration_results_get_integral
  procedure :: get_error => integration_results_get_error
  procedure :: get_accuracy => integration_results_get_accuracy
  procedure :: get_chi2 => integration_results_get_chi2
  procedure :: get_efficiency => integration_results_get_efficiency
\nwendcode{}\nwbegincode{362}\moddef{Integration results: procedures}\plusendmoddef
  function integration_results_get_n_calls (results, last, it, pass) &
       result (n_calls)
    class(integration_results_t), intent(in), target :: results
    integer :: n_calls
    logical, intent(in), optional :: last
    integer, intent(in), optional :: it, pass
    n_calls = integration_entry_get_n_calls &
         (results%get_entry (last, it, pass))
  end function integration_results_get_n_calls

  function integration_results_get_integral (results, last, it, pass) &
       result (integral)
    class(integration_results_t), intent(in), target :: results
    real(default) :: integral
    logical, intent(in), optional :: last
    integer, intent(in), optional :: it, pass
    integral = integration_entry_get_integral &
         (results%get_entry (last, it, pass))
  end function integration_results_get_integral

  function integration_results_get_error (results, last, it, pass) &
       result (error)
    class(integration_results_t), intent(in), target :: results
    real(default) :: error
    logical, intent(in), optional :: last
    integer, intent(in), optional :: it, pass
    error = integration_entry_get_error &
         (results%get_entry (last, it, pass))
  end function integration_results_get_error

  function integration_results_get_accuracy (results, last, it, pass) &
       result (accuracy)
    class(integration_results_t), intent(in), target :: results
    real(default) :: accuracy
    logical, intent(in), optional :: last
    integer, intent(in), optional :: it, pass
    accuracy = integration_entry_get_accuracy &
         (results%get_entry (last, it, pass))
  end function integration_results_get_accuracy

  function integration_results_get_chi2 (results, last, it, pass) &
       result (chi2)
    class(integration_results_t), intent(in), target :: results
    real(default) :: chi2
    logical, intent(in), optional :: last
    integer, intent(in), optional :: it, pass
    chi2 = integration_entry_get_chi2 &
         (results%get_entry (last, it, pass))
  end function integration_results_get_chi2

  function integration_results_get_efficiency (results, last, it, pass) &
       result (efficiency)
    class(integration_results_t), intent(in), target :: results
    real(default) :: efficiency
    logical, intent(in), optional :: last
    integer, intent(in), optional :: it, pass
    efficiency = integration_entry_get_efficiency &
         (results%get_entry (last, it, pass))
  end function integration_results_get_efficiency

\eatline
\eatline
\eatline
\eatline
\eatline
\eatline
\nwendcode{}\nwbegindocs{363}Return the last pass index and the index of the last iteration
\emph{within} the last pass.  The third routine returns the absolute
index of the last iteration.
\nwenddocs{}\nwbegincode{364}\moddef{Integration results: procedures}\plusendmoddef
  function integration_results_get_current_pass (results) result (pass)
    integer :: pass
    type(integration_results_t), intent(in) :: results
    pass = results%n_pass
  end function integration_results_get_current_pass

  function integration_results_get_current_it (results) result (it)
    integer :: it
    type(integration_results_t), intent(in) :: results
    if (allocated (results%entry)) then
       it = count (results%entry(1:results%n_it)%pass == results%n_pass)
    else
       it = 0
    end if
  end function integration_results_get_current_it

  function integration_results_get_last_it (results) result (it)
    integer :: it
    type(integration_results_t), intent(in) :: results
    it = results%n_it
  end function integration_results_get_last_it

\eatline
\eatline
\eatline
\nwendcode{}\nwbegindocs{365}Return the index of the best iteration (lowest accuracy value)
within the current pass.  If none qualifies, return zero.
\nwenddocs{}\nwbegincode{366}\moddef{Integration results: procedures}\plusendmoddef
  function integration_results_get_best_it (results) result (it)
    integer :: it
    type(integration_results_t), intent(in) :: results
    integer :: i
    real(default) :: acc, acc_best
    acc_best = -1
    it = 0
    do i = 1, results%n_it
       if (results%entry(i)%pass == results%n_pass) then
          acc = integration_entry_get_accuracy (results%entry(i))
          if (acc_best < 0 .or. acc <= acc_best) then
             acc_best = acc
             it = i
          end if
       end if
    end do
  end function integration_results_get_best_it

\eatline
\nwendcode{}\nwbegindocs{367}Compute the MD5 sum by printing everything and checksumming the
resulting file.
\nwenddocs{}\nwbegincode{368}\moddef{Integration results: procedures}\plusendmoddef
  function integration_results_get_md5sum (results) result (md5sum_results)
    character(32) :: md5sum_results
    type(integration_results_t), intent(in) :: results
    integer :: u
    u = free_unit ()
    open (unit = u, status = "scratch", action = "readwrite")
    call integration_results_write (results, u, verbose=.true.)
    rewind (u)
    md5sum_results = md5sum (u)
    close (u)
  end function integration_results_get_md5sum

\eatline
\nwendcode{}\nwbegindocs{369}\nwdocspar
This is (ab)used to suppress numerical noise when integrating constant
matrix elements. 
\nwenddocs{}\nwbegincode{370}\moddef{Integration results: integration results: TBP}\plusendmoddef
  procedure :: pacify => integration_results_pacify
\nwendcode{}\nwbegincode{371}\moddef{Integration results: procedures}\plusendmoddef
  subroutine integration_results_pacify (results, efficiency_reset)
    class(integration_results_t), intent(inout) :: results
    logical, intent(in), optional :: efficiency_reset
    integer :: i 
    logical :: reset    
    reset = .false.
    if (present (efficiency_reset))  reset = efficiency_reset
    if (allocated (results%entry)) then
       do i = 1, size (results%entry)
          call pacify (results%entry(i)%error, &
               results%entry(i)%integral * 1.E-9_default)
          if (reset)  results%entry(i)%efficiency = 1
       end do
    end if    
    if (allocated (results%average)) then
       do i = 1, size (results%average)
          call pacify (results%average(i)%error, &
               results%average(i)%integral * 1.E-9_default)
          if (reset)  results%average(i)%efficiency = 1
       end do
    end if    
  end subroutine integration_results_pacify
        
\eatline
\nwendcode{}\nwbegindocs{372}\nwdocspar
\nwenddocs{}\nwbegincode{373}\moddef{Integration results: integration results: TBP}\plusendmoddef
  procedure :: record_correction => integration_results_record_correction
\nwendcode{}\nwbegincode{374}\moddef{Integration results: procedures}\plusendmoddef
  subroutine integration_results_record_correction (object, corr, err)
    class(integration_results_t), intent(inout) :: object
    real(default), intent(in) :: corr, err
    integer :: u
    u = given_output_unit ()
    if (object%screen) then
      call write_hline (u)
      call msg_message ("NLO Results: O(alpha_s+1)/O(alpha_s)")
      write(u,'(1X,F6.4,A4,F7.5)') corr, ' +- ', err
    end if
  end subroutine integration_results_record_correction

\eatline
\nwendcode{}\nwbegindocs{375}\nwdocspar
\subsection{Results display}
Write a driver file for history visualization.

The ratio of $y$ range over $y$ value must not become too small, otherwise
we run into an arithmetic overflow in GAMELAN.  2\% appears to be safe.
\nwenddocs{}\nwbegincode{376}\moddef{Integration results: parameters}\plusendmoddef
  real, parameter, public :: GML_MIN_RANGE_RATIO = 0.02
\nwendcode{}\nwbegincode{377}\moddef{Integration results: public}\plusendmoddef
  public :: integration_results_write_driver
\nwendcode{}\nwbegincode{378}\moddef{Integration results: procedures}\plusendmoddef
  subroutine integration_results_write_driver (results, filename, eff_reset)
    type(integration_results_t), intent(inout) :: results
    type(string_t), intent(in) :: filename
    logical, intent(in), optional :: eff_reset
    type(string_t) :: file_tex
    integer :: unit
    integer :: n, i, n_pass, pass
    integer, dimension(:), allocatable :: ipass
    real(default) :: ymin, ymax, yavg, ydif, y0, y1
    logical :: reset 
    file_tex = filename // ".tex"
    unit = free_unit ()
    open (unit=unit, file=char(file_tex), action="write", status="replace")
    reset = .false.; if (present (eff_reset))  reset = eff_reset
    n = results%n_it
    n_pass = results%n_pass
    allocate (ipass (results%n_pass))
    ipass(1) = 0
    pass = 2
    do i = 1, n-1
       if (integration_entry_get_pass (results%entry(i)) &
           /= integration_entry_get_pass (results%entry(i+1))) then
          ipass(pass) = i
          pass = pass + 1
       end if
    end do
    ymin = minval (integration_entry_get_integral (results%entry(:n)) &
                   - integration_entry_get_error (results%entry(:n)))
    ymax = maxval (integration_entry_get_integral (results%entry(:n)) &
                   + integration_entry_get_error (results%entry(:n)))
    yavg = (ymax + ymin) / 2
    ydif = (ymax - ymin)
    if (ydif * 1.5 > GML_MIN_RANGE_RATIO * yavg) then
       y0 = yavg - ydif * 0.75
       y1 = yavg + ydif * 0.75
    else
       y0 = yavg * (1 - GML_MIN_RANGE_RATIO / 2)
       y1 = yavg * (1 + GML_MIN_RANGE_RATIO / 2)
    end if
    write (unit, "(A)") "\\documentclass\{article\}"
    write (unit, "(A)") "\\usepackage\{a4wide\}"
    write (unit, "(A)") "\\usepackage\{gamelan\}"
    write (unit, "(A)") "\\usepackage\{amsmath\}"
    write (unit, "(A)") ""
    write (unit, "(A)") "\\begin\{document\}"
    write (unit, "(A)") "\\begin\{gmlfile\}"
    write (unit, "(A)") "\\section*\{Integration Results Display\}"
    write (unit, "(A)") ""
    write (unit, "(A)") "Process: \\verb|" // char (filename) // "|"
    write (unit, "(A)") ""
    write (unit, "(A)") "\\vspace*\{2\\baselineskip\}"
    write (unit, "(A)") "\\unitlength 1mm"
    write (unit, "(A)") "\\begin\{gmlcode\}"
    write (unit, "(A)") "  picture sym;  sym = fshape (circle scaled 1mm)();"
    write (unit, "(A)") "  color col.band;  col.band = 0.9white;"
    write (unit, "(A)") "  color col.eband;  col.eband = 0.98white;"
    write (unit, "(A)") "\\end\{gmlcode\}"
    write (unit, "(A)") "\\begin\{gmlgraph*\}(130,180)[history]"
    write (unit, "(A)") "  setup (linear, linear);"
    write (unit, "(A,I0,A)") "  history.n_pass = ", n_pass, ";"
    write (unit, "(A,I0,A)") "  history.n_it   = ", n, ";"
    write (unit, "(A,A,A)")  "  history.y0 = #""", char (mp_format (y0)), """;"
    write (unit, "(A,A,A)")  "  history.y1 = #""", char (mp_format (y1)), """;"
    write (unit, "(A)") &
         "  graphrange (#0.5, history.y0), (#(n+0.5), history.y1);"
    do pass = 1, n_pass
       write (unit, "(A,I0,A,I0,A)") &
            "  history.pass[", pass, "] = ", ipass(pass), ";"
       write (unit, "(A,I0,A,A,A)") &
            "  history.avg[", pass, "] = #""", &
            char (mp_format &
               (integration_entry_get_integral (results%average(pass)))), &
            """;"
       write (unit, "(A,I0,A,A,A)") &
            "  history.err[", pass, "] = #""", &
            char (mp_format &
               (integration_entry_get_error (results%average(pass)))), &
            """;"
       write (unit, "(A,I0,A,A,A)") &
            "  history.chi[", pass, "] = #""", &
            char (mp_format &
               (integration_entry_get_chi2 (results%average(pass)))), &
            """;"
    end do
    write (unit, "(A,I0,A,I0,A)") &
         "  history.pass[", n_pass + 1, "] = ", n, ";"
    write (unit, "(A)")  "  for i = 1 upto history.n_pass:"
    write (unit, "(A)")  "    if history.chi[i] greater one:"
    write (unit, "(A)")  "    fill plot ("
    write (unit, "(A)")  &
         "      (#(history.pass[i]  +.5), " &
         // "history.avg[i] minus history.err[i] times history.chi[i]),"
    write (unit, "(A)")  &
         "      (#(history.pass[i+1]+.5), " &
         // "history.avg[i] minus history.err[i] times history.chi[i]),"
    write (unit, "(A)")  &
         "      (#(history.pass[i+1]+.5), " &
         // "history.avg[i] plus history.err[i] times history.chi[i]),"
    write (unit, "(A)")  &
         "      (#(history.pass[i]  +.5), " &
         // "history.avg[i] plus history.err[i] times history.chi[i])"
    write (unit, "(A)")  "    ) withcolor col.eband fi;"
    write (unit, "(A)")  "    fill plot ("
    write (unit, "(A)")  &
         "      (#(history.pass[i]  +.5), history.avg[i] minus history.err[i]),"
    write (unit, "(A)")  &
         "      (#(history.pass[i+1]+.5), history.avg[i] minus history.err[i]),"
    write (unit, "(A)")  &
         "      (#(history.pass[i+1]+.5), history.avg[i] plus history.err[i]),"
    write (unit, "(A)")  &
         "      (#(history.pass[i]  +.5), history.avg[i] plus history.err[i])"
    write (unit, "(A)")  "    ) withcolor col.band;"
    write (unit, "(A)")  "    draw plot ("
    write (unit, "(A)")  &
         "      (#(history.pass[i]  +.5), history.avg[i]),"
    write (unit, "(A)")  &
         "      (#(history.pass[i+1]+.5), history.avg[i])"
    write (unit, "(A)")  "      ) dashed evenly;"
    write (unit, "(A)")  "  endfor"
    write (unit, "(A)")  "  for i = 1 upto history.n_pass + 1:"
    write (unit, "(A)")  "    draw plot ("
    write (unit, "(A)")  &
         "      (#(history.pass[i]+.5), history.y0),"
    write (unit, "(A)")  &
         "      (#(history.pass[i]+.5), history.y1)"
    write (unit, "(A)")  "      ) dashed withdots;"
    write (unit, "(A)")  "  endfor"
    do i = 1, n
       write (unit, "(A,I0,A,A,A,A,A)") "  plot (history) (#", &
          i, ", #""", &
          char (mp_format (integration_entry_get_integral (results%entry(i)))),&
          """) vbar #""", &
          char (mp_format (integration_entry_get_error (results%entry(i)))), &
          """;"
    end do
    write (unit, "(A)") "  draw piecewise from (history) " &
      // "withsymbol sym;"
    write (unit, "(A)") "  fullgrid.lr (5,20);"
    write (unit, "(A)") "  standardgrid.bt (n);"
    write (unit, "(A)")  "  begingmleps ""Whizard-Logo.eps"";"
    write (unit, "(A)")  "    base := (120*unitlength,170*unitlength);"
    write (unit, "(A)")  "    height := 9.6*unitlength;"
    write (unit, "(A)")  "    width := 11.2*unitlength;"
    write (unit, "(A)")  "  endgmleps;"    
    write (unit, "(A)") "\\end\{gmlgraph*\}"
    write (unit, "(A)") "\\end\{gmlfile\}"
    write (unit, "(A)") "\\clearpage"
    write (unit, "(A)") "\\begin\{verbatim\}"
    if (reset) then
      call results%pacify (reset)
    end if      
    call integration_results_write (results, unit)
    write (unit, "(A)") "\\end\{verbatim\}"
    write (unit, "(A)") "\\end\{document\}"
    close (unit)
  end subroutine integration_results_write_driver

\eatline
\nwendcode{}\nwbegindocs{379}Call \LaTeX\ and Metapost for the history driver file, and convert to PS and
PDF.
\nwenddocs{}\nwbegincode{380}\moddef{Integration results: public}\plusendmoddef
  public :: integration_results_compile_driver
\nwendcode{}\nwbegincode{381}\moddef{Integration results: procedures}\plusendmoddef
  subroutine integration_results_compile_driver (results, filename, os_data)
    type(integration_results_t), intent(in) :: results
    type(string_t), intent(in) :: filename
    type(os_data_t), intent(in) :: os_data
    integer :: unit_dev, status
    type(string_t) :: file_tex, file_dvi, file_ps, file_pdf, file_mp
    type(string_t) :: setenv_tex, setenv_mp, pipe, pipe_dvi
    if (.not. os_data%event_analysis) then
       call msg_warning ("Skipping integration history display " &
           // "because latex or mpost is not available")
       return
    end if
    file_tex = filename // ".tex"
    file_dvi = filename // ".dvi"
    file_ps = filename // ".ps"
    file_pdf = filename // ".pdf"
    file_mp = filename // ".mp"
    call msg_message ("Creating integration history display "& 
         // char (file_ps) // " and " // char (file_pdf))
    BLOCK: do
       unit_dev = free_unit ()
       open (file = "/dev/null", unit = unit_dev, &
              action = "write", iostat = status)
       if (status /= 0) then
          pipe = ""
          pipe_dvi = ""
       else
          pipe = " > /dev/null"
          pipe_dvi = " 2>/dev/null 1>/dev/null"
       end if
       close (unit_dev)
       if (os_data%whizard_texpath /= "") then
          setenv_tex = &
               "TEXINPUTS=" // os_data%whizard_texpath // ":$TEXINPUTS "
          setenv_mp = &
               "MPINPUTS=" // os_data%whizard_texpath // ":$MPINPUTS "
       else
          setenv_tex = ""
          setenv_mp = ""
       end if
       call os_system_call (setenv_tex // os_data%latex // " " // &
            file_tex // pipe, status)
       if (status /= 0)  exit BLOCK
       if (os_data%gml /= "") then
          call os_system_call (setenv_mp // os_data%gml // " " // &
               file_mp // pipe, status)
       else 
          call msg_error ("Could not use GAMELAN/MetaPOST.")
          exit BLOCK
       end if
       if (status /= 0)  exit BLOCK
       call os_system_call (setenv_tex // os_data%latex // " " // &
             file_tex // pipe, status)
       if (status /= 0)  exit BLOCK
       if (os_data%event_analysis_ps) then
          call os_system_call (os_data%dvips // " " // &
             file_dvi // pipe_dvi, status)
          if (status /= 0)  exit BLOCK
       else
          call msg_warning ("Skipping PostScript generation because dvips " &
               // "is not available")
          exit BLOCK
       end if
       if (os_data%event_analysis_pdf) then
          call os_system_call (os_data%ps2pdf // " " // &
                  file_ps, status)
          if (status /= 0)  exit BLOCK
       else
          call msg_warning ("Skipping PDF generation because ps2pdf " &
               // "is not available")
          exit BLOCK
       end if
       exit BLOCK
    end do BLOCK
    if (status /= 0) then
       call msg_error ("Unable to compile integration history display")
    end if
  end subroutine integration_results_compile_driver

\eatline
\nwendcode{}\nwbegindocs{382}\nwdocspar
\clearpage
%%%%%%%%%%%%%%%%%%%%%%%%%%%%%%%%%%%%%%%%%%%%%%%%%%%%%%%%%%%%%%%%%%%%%%%%
\section{Simple midpoint integration}
This is a most simple implementation of an integrator.  The algorithm
is the straightforward multi-dimensional midpoint rule, i.e., the
integration hypercube is binned uniformly, the integrand is evaluated
at the midpoints of each bin, and the result is the average.  The
binning is equivalent for all integration dimensions.

This rule is accurate to the order $h^2$, where $h$ is the bin width.
Given that $h=N^{-1/d}$, where $d$ is the integration dimension and
$N$ is the total number of sampling points, we get a relative error
of order $N^{-2/d}$.  This is superior to MC integration if $d<4$, and
equivalent if $d=4$.  It is not worse than higher-order formulas
(such as Gauss integration) if the integrand is not smooth, e.g., if
it contains cuts.

The integrator is specifically single-channel.  However, we do not
limit the dimension.
\nwenddocs{}\nwbegincode{383}\moddef{\code{}mci{\char95}midpoint.f90\edoc{}}\endmoddef
\LA{}File header\RA{}

module mci_midpoint

  use kinds
\LA{}Use strings\RA{}
  use io_units
  use unit_tests
  use diagnostics

  use phs_base
  use rng_base
  use mci_base
  
\LA{}Standard module head\RA{}

\LA{}MCI midpoint: public\RA{}

\LA{}MCI midpoint: types\RA{}

\LA{}MCI midpoint: test types\RA{}

contains
  
\LA{}MCI midpoint: procedures\RA{}

\LA{}MCI midpoint: tests\RA{}

end module mci_midpoint
\eatline
\nwendcode{}\nwbegindocs{384}\nwdocspar
\subsection{Integrator}
The object contains the methods for integration and event generation.
For the actual work and data storage, it spawns an instance object.

After an integration pass, we update the {\tt{}max} parameter to indicate
the maximum absolute value of the integrand that the integrator
encountered.  This is required for event generation.
\nwenddocs{}\nwbegincode{385}\moddef{MCI midpoint: public}\endmoddef
  public :: mci_midpoint_t
\nwendcode{}\nwbegincode{386}\moddef{MCI midpoint: types}\endmoddef
  type, extends (mci_t) :: mci_midpoint_t
     integer :: n_dim_binned = 0
     logical, dimension(:), allocatable :: dim_is_binned
     logical :: calls_known = .false.
     integer :: n_calls = 0
     integer :: n_calls_pos = 0
     integer :: n_calls_nul = 0
     integer :: n_calls_neg = 0
     real(default) :: integral_pos = 0
     real(default) :: integral_neg = 0
     integer, dimension(:), allocatable :: n_bin
     logical :: max_known = .false.
     real(default) :: max = 0
     real(default) :: min = 0
     real(default) :: max_abs = 0
     real(default) :: min_abs = 0
   contains
   \LA{}MCI midpoint: mci midpoint: TBP\RA{}
end type mci_midpoint_t
  
\eatline
\nwendcode{}\nwbegindocs{387}Finalizer: base version is sufficient
\nwenddocs{}\nwbegincode{388}\moddef{MCI midpoint: mci midpoint: TBP}\endmoddef
  procedure :: final => mci_midpoint_final
\nwendcode{}\nwbegincode{389}\moddef{MCI midpoint: procedures}\endmoddef
  subroutine mci_midpoint_final (object)
    class(mci_midpoint_t), intent(inout) :: object
    call object%base_final ()
  end subroutine mci_midpoint_final
  
\eatline
\nwendcode{}\nwbegindocs{390}Output.
\nwenddocs{}\nwbegincode{391}\moddef{MCI midpoint: mci midpoint: TBP}\plusendmoddef
  procedure :: write => mci_midpoint_write
\nwendcode{}\nwbegincode{392}\moddef{MCI midpoint: procedures}\plusendmoddef
  subroutine mci_midpoint_write (object, unit, pacify, md5sum_version)
    class(mci_midpoint_t), intent(in) :: object
    integer, intent(in), optional :: unit
    logical, intent(in), optional :: pacify
    logical, intent(in), optional :: md5sum_version
    integer :: u, i
    u = given_output_unit (unit)
    write (u, "(1x,A)") "Single-channel midpoint rule integrator:"
    call object%base_write (u, pacify, md5sum_version)
    if (object%n_dim_binned < object%n_dim) then
       write (u, "(3x,A,99(1x,I0))")  "Flat dimensions      =", &
            pack ([(i, i = 1, object%n_dim)], mask = .not. object%dim_is_binned)
       write (u, "(3x,A,I0)")  "Number of binned dim = ", object%n_dim_binned
    end if
    if (object%calls_known) then
       write (u, "(3x,A,99(1x,I0))")  "Number of bins       =", object%n_bin
       write (u, "(3x,A,I0)")  "Number of calls      = ", object%n_calls
       if (object%n_calls_pos /= object%n_calls) then
          write (u, "(3x,A,I0)")  "  positive value     = ", object%n_calls_pos
          write (u, "(3x,A,I0)")  "  zero value         = ", object%n_calls_nul
          write (u, "(3x,A,I0)")  "  negative value     = ", object%n_calls_neg
          write (u, "(3x,A,ES17.10)") &
               "Integral (pos. part) = ", object%integral_pos
          write (u, "(3x,A,ES17.10)") &
               "Integral (neg. part) = ", object%integral_neg
       end if
    end if
    if (object%max_known) then
       write (u, "(3x,A,ES17.10)")  "Maximum of integrand = ", object%max
       write (u, "(3x,A,ES17.10)")  "Minimum of integrand = ", object%min
       if (object%min /= object%min_abs) then
          write (u, "(3x,A,ES17.10)")  "Maximum (abs. value) = ", object%max_abs
          write (u, "(3x,A,ES17.10)")  "Minimum (abs. value) = ", object%min_abs
       end if
    end if
    if (allocated (object%rng))  call object%rng%write (u)
  end subroutine mci_midpoint_write
    
\eatline
\nwendcode{}\nwbegindocs{393}Startup message: short version.
\nwenddocs{}\nwbegincode{394}\moddef{MCI midpoint: mci midpoint: TBP}\plusendmoddef
  procedure :: startup_message => mci_midpoint_startup_message
\nwendcode{}\nwbegincode{395}\moddef{MCI midpoint: procedures}\plusendmoddef
  subroutine mci_midpoint_startup_message (mci, unit, n_calls)
    class(mci_midpoint_t), intent(in) :: mci
    integer, intent(in), optional :: unit, n_calls
    call mci%base_startup_message (unit = unit, n_calls = n_calls)
    if (mci%n_dim_binned < mci%n_dim) then
       write (msg_buffer, "(A,2(1x,I0,1x,A))") &
            "Integrator: Midpoint rule:", &
            mci%n_dim_binned, "binned dimensions"
    else
       write (msg_buffer, "(A,2(1x,I0,1x,A))") &
            "Integrator: Midpoint rule"
    end if
    call msg_message (unit = unit)
  end subroutine mci_midpoint_startup_message
    
\eatline
\nwendcode{}\nwbegindocs{396}The number of channels must be one.
\nwenddocs{}\nwbegincode{397}\moddef{MCI midpoint: mci midpoint: TBP}\plusendmoddef
  procedure :: set_dimensions => mci_midpoint_set_dimensions
\nwendcode{}\nwbegincode{398}\moddef{MCI midpoint: procedures}\plusendmoddef
  subroutine mci_midpoint_set_dimensions (mci, n_dim, n_channel)
    class(mci_midpoint_t), intent(inout) :: mci
    integer, intent(in) :: n_dim
    integer, intent(in) :: n_channel
    if (n_channel == 1) then
       mci%n_channel = n_channel
       mci%n_dim = n_dim
       allocate (mci%dim_is_binned (mci%n_dim))
       mci%dim_is_binned = .true.
       mci%n_dim_binned = count (mci%dim_is_binned)
       allocate (mci%n_bin (mci%n_dim))
       mci%n_bin = 0
    else
       call msg_fatal ("Attempt to initialize single-channel integrator &
            &for multiple channels")
    end if
  end subroutine mci_midpoint_set_dimensions
  
\eatline
\nwendcode{}\nwbegindocs{399}Declare particular dimensions as flat.  These dimensions will not be binned.
\nwenddocs{}\nwbegincode{400}\moddef{MCI midpoint: mci midpoint: TBP}\plusendmoddef
  procedure :: declare_flat_dimensions => mci_midpoint_declare_flat_dimensions
\nwendcode{}\nwbegincode{401}\moddef{MCI midpoint: procedures}\plusendmoddef
  subroutine mci_midpoint_declare_flat_dimensions (mci, dim_flat)
    class(mci_midpoint_t), intent(inout) :: mci
    integer, dimension(:), intent(in) :: dim_flat
    integer :: d
    mci%n_dim_binned = mci%n_dim - size (dim_flat)
    do d = 1, size (dim_flat)
       mci%dim_is_binned(dim_flat(d)) = .false.
    end do
    mci%n_dim_binned = count (mci%dim_is_binned)
  end subroutine mci_midpoint_declare_flat_dimensions
  
\eatline
\nwendcode{}\nwbegindocs{402}Declare particular channels as equivalent.  This has no effect.
\nwenddocs{}\nwbegincode{403}\moddef{MCI midpoint: mci midpoint: TBP}\plusendmoddef
  procedure :: declare_equivalences => mci_midpoint_ignore_equivalences
\nwendcode{}\nwbegincode{404}\moddef{MCI midpoint: procedures}\plusendmoddef
  subroutine mci_midpoint_ignore_equivalences (mci, channel, dim_offset)
    class(mci_midpoint_t), intent(inout) :: mci
    type(phs_channel_t), dimension(:), intent(in) :: channel
    integer, intent(in) :: dim_offset
  end subroutine mci_midpoint_ignore_equivalences
  
\eatline
\nwendcode{}\nwbegindocs{405}Allocate instance with matching type.
\nwenddocs{}\nwbegincode{406}\moddef{MCI midpoint: mci midpoint: TBP}\plusendmoddef
  procedure :: allocate_instance => mci_midpoint_allocate_instance
\nwendcode{}\nwbegincode{407}\moddef{MCI midpoint: procedures}\plusendmoddef
  subroutine mci_midpoint_allocate_instance (mci, mci_instance)
    class(mci_midpoint_t), intent(in) :: mci
    class(mci_instance_t), intent(out), pointer :: mci_instance
    allocate (mci_midpoint_instance_t :: mci_instance)
  end subroutine mci_midpoint_allocate_instance
  
\eatline
\nwendcode{}\nwbegindocs{408}Integrate.  The number of dimensions is arbitrary.  We make sure
that the number of calls is evenly distributed among the dimensions.
The actual number of calls will typically be smaller than the
requested number, but never smaller than 1.

The sampling over a variable number of dimensions implies a variable
number of nested loops.  We implement this by a recursive subroutine,
one loop in each recursion level.

The number of iterations {\tt{}n{\char95}it} is ignored.  Also, the error is set
to zero in the current implementation.

With this integrator, we allow the calculation to abort immediately when
forced by a signal.  There is no state that we can save, hence we do not catch
an interrupt.
\nwenddocs{}\nwbegincode{409}\moddef{MCI midpoint: mci midpoint: TBP}\plusendmoddef
  procedure :: integrate => mci_midpoint_integrate
\nwendcode{}\nwbegincode{410}\moddef{MCI midpoint: procedures}\plusendmoddef
  subroutine mci_midpoint_integrate (mci, instance, sampler, n_it, n_calls, &
       results, pacify)
    class(mci_midpoint_t), intent(inout) :: mci
    class(mci_instance_t), intent(inout) :: instance
    class(mci_sampler_t), intent(inout) :: sampler
    integer, intent(in) :: n_it
    integer, intent(in) :: n_calls
    logical, intent(in), optional :: pacify
    class(mci_results_t), intent(inout), optional :: results
    real(default), dimension(:), allocatable :: x
    real(default) :: integral, integral_pos, integral_neg
    integer :: n_bin
    select type (instance)
    type is (mci_midpoint_instance_t)
       allocate (x (mci%n_dim))
       integral = 0
       integral_pos = 0
       integral_neg = 0
       select case (mci%n_dim_binned)
       case (1)
          n_bin = n_calls
       case (2:)
          n_bin = max (int (n_calls ** (1. / mci%n_dim_binned)), 1)
       end select
       where (mci%dim_is_binned)
          mci%n_bin = n_bin
       elsewhere
          mci%n_bin = 1
       end where
       mci%n_calls = product (mci%n_bin)
       mci%n_calls_pos = 0
       mci%n_calls_nul = 0
       mci%n_calls_neg = 0
       mci%calls_known = .true.
       call sample_dim (mci%n_dim)
       mci%integral = integral / mci%n_calls
       mci%integral_pos = integral_pos / mci%n_calls
       mci%integral_neg = integral_neg / mci%n_calls
       mci%integral_known = .true.
       call instance%set_max ()
       if (present (results)) then
          call results%record (1, mci%n_calls, &
               mci%integral, mci%error, mci%efficiency)
       end if
    end select
  contains
    recursive subroutine sample_dim (d)
      integer, intent(in) :: d
      integer :: i
      real(default) :: value
      do i = 1, mci%n_bin(d)
         x(d) = (i - 0.5_default) / mci%n_bin(d)
         if (d > 1) then
            call sample_dim (d - 1)
         else
            if (signal_is_pending ())  return
            call instance%evaluate (sampler, 1, x)
            value = instance%get_value ()
            if (value > 0) then
               mci%n_calls_pos = mci%n_calls_pos + 1
               integral = integral + value
               integral_pos = integral_pos + value
            else if (value == 0) then
               mci%n_calls_nul = mci%n_calls_nul + 1
            else
               mci%n_calls_neg = mci%n_calls_neg + 1
               integral = integral + value
               integral_neg = integral_neg + value
            end if
         end if
      end do
    end subroutine sample_dim
  end subroutine mci_midpoint_integrate

\eatline
\nwendcode{}\nwbegindocs{411}Simulation initializer and finalizer: nothing to do here.
\nwenddocs{}\nwbegincode{412}\moddef{MCI midpoint: mci midpoint: TBP}\plusendmoddef
  procedure :: prepare_simulation => mci_midpoint_ignore_prepare_simulation
\nwendcode{}\nwbegincode{413}\moddef{MCI midpoint: procedures}\plusendmoddef
  subroutine mci_midpoint_ignore_prepare_simulation (mci)
    class(mci_midpoint_t), intent(inout) :: mci
  end subroutine mci_midpoint_ignore_prepare_simulation
  
\eatline
\nwendcode{}\nwbegindocs{414}Generate weighted event.
\nwenddocs{}\nwbegincode{415}\moddef{MCI midpoint: mci midpoint: TBP}\plusendmoddef
  procedure :: generate_weighted_event => mci_midpoint_generate_weighted_event
\nwendcode{}\nwbegincode{416}\moddef{MCI midpoint: procedures}\plusendmoddef
  subroutine mci_midpoint_generate_weighted_event (mci, instance, sampler)
    class(mci_midpoint_t), intent(inout) :: mci
    class(mci_instance_t), intent(inout), target :: instance
    class(mci_sampler_t), intent(inout), target :: sampler
    real(default), dimension(mci%n_dim) :: x
    select type (instance)
    type is (mci_midpoint_instance_t)
       call mci%rng%generate (x)
       call instance%evaluate (sampler, 1, x)
       instance%excess_weight = 0
    end select
  end subroutine mci_midpoint_generate_weighted_event
       
\eatline
\nwendcode{}\nwbegindocs{417}For unweighted events, we generate weighted events and apply a
simple rejection step to the relative event weight, until an event
passes.  

Note that we use the {\tt{}max{\char95}abs} value stored in the configuration
record, not the one stored in the instance.  The latter may change
during event generation.  After an event generation pass is over, we
may update the value for a subsequent pass.
\nwenddocs{}\nwbegincode{418}\moddef{MCI midpoint: mci midpoint: TBP}\plusendmoddef
  procedure :: generate_unweighted_event => &
       mci_midpoint_generate_unweighted_event
\nwendcode{}\nwbegincode{419}\moddef{MCI midpoint: procedures}\plusendmoddef
  subroutine mci_midpoint_generate_unweighted_event (mci, instance, sampler)
    class(mci_midpoint_t), intent(inout) :: mci
    class(mci_instance_t), intent(inout), target :: instance
    class(mci_sampler_t), intent(inout), target :: sampler
    real(default) :: x, norm, int
    select type (instance)
    type is (mci_midpoint_instance_t)
       if (mci%max_known .and. mci%max_abs > 0) then
          norm = abs (mci%max_abs * instance%safety_factor)
          REJECTION: do
             call mci%generate_weighted_event (instance, sampler)
             if (sampler%is_valid ()) then
                call mci%rng%generate (x)
                int = abs (instance%integrand)
                if (x * norm <= int) then
                   if (norm > 0 .and. norm < int) then
                      instance%excess_weight = int / norm - 1
                   end if
                   exit REJECTION
                end if
             end if
             if (signal_is_pending ())  return
          end do REJECTION
       else
          call msg_fatal ("Unweighted event generation: &
               &maximum of integrand is zero or unknown")
       end if
    end select
  end subroutine mci_midpoint_generate_unweighted_event
    
\eatline
\nwendcode{}\nwbegindocs{420}Rebuild an event, using the {\tt{}state} input.
\nwenddocs{}\nwbegincode{421}\moddef{MCI midpoint: mci midpoint: TBP}\plusendmoddef
  procedure :: rebuild_event => mci_midpoint_rebuild_event
\nwendcode{}\nwbegincode{422}\moddef{MCI midpoint: procedures}\plusendmoddef
  subroutine mci_midpoint_rebuild_event (mci, instance, sampler, state)
    class(mci_midpoint_t), intent(inout) :: mci
    class(mci_instance_t), intent(inout) :: instance
    class(mci_sampler_t), intent(inout) :: sampler
    class(mci_state_t), intent(in) :: state
    select type (instance)
    type is (mci_midpoint_instance_t)
       call instance%recall (sampler, state)
    end select
  end subroutine mci_midpoint_rebuild_event
       
\eatline
\nwendcode{}\nwbegindocs{423}\nwdocspar
\subsection{Integrator instance}
Covering the case of flat dimensions, we store a complete {\tt{}x} array.  This
is filled when generating events.
\nwenddocs{}\nwbegincode{424}\moddef{MCI midpoint: public}\plusendmoddef
  public :: mci_midpoint_instance_t
\nwendcode{}\nwbegincode{425}\moddef{MCI midpoint: types}\plusendmoddef
  type, extends (mci_instance_t) :: mci_midpoint_instance_t
     type(mci_midpoint_t), pointer :: mci => null ()
     logical :: max_known = .false.
     real(default) :: max = 0
     real(default) :: min = 0
     real(default) :: max_abs = 0
     real(default) :: min_abs = 0
     real(default) :: safety_factor = 1
     real(default) :: excess_weight = 0
   contains
   \LA{}MCI midpoint: mci midpoint instance: TBP\RA{}
  end type mci_midpoint_instance_t
  
\eatline
\nwendcode{}\nwbegindocs{426}Output.
\nwenddocs{}\nwbegincode{427}\moddef{MCI midpoint: mci midpoint instance: TBP}\endmoddef
  procedure :: write => mci_midpoint_instance_write
\nwendcode{}\nwbegincode{428}\moddef{MCI midpoint: procedures}\plusendmoddef
  subroutine mci_midpoint_instance_write (object, unit, pacify)
    class(mci_midpoint_instance_t), intent(in) :: object
    integer, intent(in), optional :: unit
    logical, intent(in), optional :: pacify
    integer :: u
    u = given_output_unit (unit)
    write (u, "(1x,A,9(1x,F12.10))")  "x =", object%x(:,1)
    write (u, "(1x,A,ES19.12)") "Integrand = ", object%integrand
    write (u, "(1x,A,ES19.12)") "Weight    = ", object%mci_weight
    if (object%safety_factor /= 1) then
       write (u, "(1x,A,ES19.12)") "Safety f  = ", object%safety_factor
    end if
    if (object%excess_weight /= 0) then
       write (u, "(1x,A,ES19.12)") "Excess    = ", object%excess_weight
    end if
    if (object%max_known) then
       write (u, "(1x,A,ES19.12)") "Maximum   = ", object%max
       write (u, "(1x,A,ES19.12)") "Minimum   = ", object%min
       if (object%min /= object%min_abs) then
          write (u, "(1x,A,ES19.12)") "Max.(abs) = ", object%max_abs
          write (u, "(1x,A,ES19.12)") "Min.(abs) = ", object%min_abs
       end if
    end if
  end subroutine mci_midpoint_instance_write
  
\eatline
\nwendcode{}\nwbegindocs{429}The finalizer is empty.
\nwenddocs{}\nwbegincode{430}\moddef{MCI midpoint: mci midpoint instance: TBP}\plusendmoddef
  procedure :: final => mci_midpoint_instance_final
\nwendcode{}\nwbegincode{431}\moddef{MCI midpoint: tests}\endmoddef
  subroutine mci_midpoint_instance_final (object)
    class(mci_midpoint_instance_t), intent(inout) :: object
  end subroutine mci_midpoint_instance_final
  
\eatline
\nwendcode{}\nwbegindocs{432}Initializer.
\nwenddocs{}\nwbegincode{433}\moddef{MCI midpoint: mci midpoint instance: TBP}\plusendmoddef
  procedure :: init => mci_midpoint_instance_init
\nwendcode{}\nwbegincode{434}\moddef{MCI midpoint: procedures}\plusendmoddef
  subroutine mci_midpoint_instance_init (mci_instance, mci)
    class(mci_midpoint_instance_t), intent(out) :: mci_instance
    class(mci_t), intent(in), target :: mci
    call mci_instance%base_init (mci)
    select type (mci)
    type is (mci_midpoint_t)
       mci_instance%mci => mci
       call mci_instance%get_max ()
       mci_instance%selected_channel = 1
    end select
  end subroutine mci_midpoint_instance_init
    
\eatline
\nwendcode{}\nwbegindocs{435}Copy the stored extrema of the integrand in the instance record.
\nwenddocs{}\nwbegincode{436}\moddef{MCI midpoint: mci midpoint instance: TBP}\plusendmoddef
  procedure :: get_max => mci_midpoint_instance_get_max
\nwendcode{}\nwbegincode{437}\moddef{MCI midpoint: procedures}\plusendmoddef
  subroutine mci_midpoint_instance_get_max (instance)
    class(mci_midpoint_instance_t), intent(inout) :: instance
    associate (mci => instance%mci)
      if (mci%max_known) then
         instance%max_known = .true.
         instance%max = mci%max
         instance%min = mci%min
         instance%max_abs = mci%max_abs
         instance%min_abs = mci%min_abs
      end if
    end associate
  end subroutine mci_midpoint_instance_get_max
  
\eatline
\nwendcode{}\nwbegindocs{438}Reverse operations: recall the extrema, but only if they are wider
than the extrema already stored in the configuration.  Also recalculate the
efficiency value.
\nwenddocs{}\nwbegincode{439}\moddef{MCI midpoint: mci midpoint instance: TBP}\plusendmoddef
  procedure :: set_max => mci_midpoint_instance_set_max
\nwendcode{}\nwbegincode{440}\moddef{MCI midpoint: procedures}\plusendmoddef
  subroutine mci_midpoint_instance_set_max (instance)
    class(mci_midpoint_instance_t), intent(inout) :: instance
    associate (mci => instance%mci)
      if (instance%max_known) then
         if (mci%max_known) then
            mci%max = max (mci%max, instance%max)
            mci%min = min (mci%min, instance%min)
            mci%max_abs = max (mci%max_abs, instance%max_abs)
            mci%min_abs = min (mci%min_abs, instance%min_abs)
         else
            mci%max = instance%max
            mci%min = instance%min
            mci%max_abs = instance%max_abs
            mci%min_abs = instance%min_abs
            mci%max_known = .true.
         end if
         if (mci%max_abs /= 0) then
            if (mci%integral == mci%integral_pos) then
               mci%efficiency = mci%integral / mci%max_abs
               mci%efficiency_known = .true.
            else if (mci%n_calls /= 0) then
               mci%efficiency = &
                    (mci%n_calls_pos * mci%integral_pos &
                    - mci%n_calls_neg * mci%integral_neg) &
                    / mci%n_calls / mci%max_abs
               mci%efficiency_known = .true.
            end if
         end if
      end if
    end associate
  end subroutine mci_midpoint_instance_set_max
  
\eatline
\nwendcode{}\nwbegindocs{441}The weight is the Jacobian of the mapping for the only channel.
\nwenddocs{}\nwbegincode{442}\moddef{MCI midpoint: mci midpoint instance: TBP}\plusendmoddef
  procedure :: compute_weight => mci_midpoint_instance_compute_weight
\nwendcode{}\nwbegincode{443}\moddef{MCI midpoint: procedures}\plusendmoddef
  subroutine mci_midpoint_instance_compute_weight (mci, c)
    class(mci_midpoint_instance_t), intent(inout) :: mci
    integer, intent(in) :: c
    select case (c)
    case (1)
       mci%mci_weight = mci%f(1)
    case default
       call msg_fatal ("MCI midpoint integrator: only single channel supported")
    end select
  end subroutine mci_midpoint_instance_compute_weight
    
\eatline
\nwendcode{}\nwbegindocs{444}Record the integrand.  Update stored values for maximum and minimum.
\nwenddocs{}\nwbegincode{445}\moddef{MCI midpoint: mci midpoint instance: TBP}\plusendmoddef
  procedure :: record_integrand => mci_midpoint_instance_record_integrand
\nwendcode{}\nwbegincode{446}\moddef{MCI midpoint: procedures}\plusendmoddef
  subroutine mci_midpoint_instance_record_integrand (mci, integrand)
    class(mci_midpoint_instance_t), intent(inout) :: mci
    real(default), intent(in) :: integrand
    mci%integrand = integrand
    if (mci%max_known) then
       mci%max = max (mci%max, integrand)
       mci%min = min (mci%min, integrand)
       mci%max_abs = max (mci%max_abs, abs (integrand))
       mci%min_abs = min (mci%min_abs, abs (integrand))
    else
       mci%max = integrand
       mci%min = integrand
       mci%max_abs = abs (integrand)
       mci%min_abs = abs (integrand)
       mci%max_known = .true.
    end if
  end subroutine mci_midpoint_instance_record_integrand
  
\eatline
\nwendcode{}\nwbegindocs{447}We store the safety factor, otherwise nothing to do here.
\nwenddocs{}\nwbegincode{448}\moddef{MCI midpoint: mci midpoint instance: TBP}\plusendmoddef
  procedure :: init_simulation => mci_midpoint_instance_init_simulation
  procedure :: final_simulation => mci_midpoint_instance_final_simulation
\nwendcode{}\nwbegincode{449}\moddef{MCI midpoint: procedures}\plusendmoddef
  subroutine mci_midpoint_instance_init_simulation (instance, safety_factor)
    class(mci_midpoint_instance_t), intent(inout) :: instance
    real(default), intent(in), optional :: safety_factor
    if (present (safety_factor))  instance%safety_factor = safety_factor
  end subroutine mci_midpoint_instance_init_simulation
  
  subroutine mci_midpoint_instance_final_simulation (instance)
    class(mci_midpoint_instance_t), intent(inout) :: instance
  end subroutine mci_midpoint_instance_final_simulation
  
\eatline
\eatline
\nwendcode{}\nwbegindocs{450}Return excess weight for the current event.
\nwenddocs{}\nwbegincode{451}\moddef{MCI midpoint: mci midpoint instance: TBP}\plusendmoddef
  procedure :: get_event_excess => mci_midpoint_instance_get_event_excess
\nwendcode{}\nwbegincode{452}\moddef{MCI midpoint: procedures}\plusendmoddef
  function mci_midpoint_instance_get_event_excess (mci) result (excess)
    class(mci_midpoint_instance_t), intent(in) :: mci
    real(default) :: excess
    excess = mci%excess_weight
  end function mci_midpoint_instance_get_event_excess
  
\eatline
\nwendcode{}\nwbegindocs{453}\nwdocspar
\subsection{Unit tests}
\nwenddocs{}\nwbegincode{454}\moddef{MCI midpoint: public}\plusendmoddef
  public :: mci_midpoint_test
\nwendcode{}\nwbegincode{455}\moddef{MCI midpoint: tests}\plusendmoddef
  subroutine mci_midpoint_test (u, results)
    integer, intent(in) :: u
    type(test_results_t), intent(inout) :: results
  \LA{}MCI midpoint: execute tests\RA{}
  end subroutine mci_midpoint_test
  
\eatline
\nwendcode{}\nwbegindocs{456}\nwdocspar
\subsubsection{Test sampler}
A test sampler object should implement a function with known integral that
we can use to check the integrator.

This is the function $f(x) = 3 x^2$ with integral $\int_0^1
f(x)\,dx=1$ and maximum $f(1)=3$.  If the integration dimension is
greater than one, the function is extended as a constant in the other
dimension(s).

Mimicking the behavior of a process object, we store the argument
and result inside the sampler, so we can {\tt{}fetch} results.
\nwenddocs{}\nwbegincode{457}\moddef{MCI midpoint: test types}\endmoddef
  type, extends (mci_sampler_t) :: test_sampler_1_t
     real(default), dimension(:), allocatable :: x
     real(default) :: val
   contains
   \LA{}MCI midpoint: test sampler 1: TBP\RA{}
  end type test_sampler_1_t

\eatline
\nwendcode{}\nwbegindocs{458}Output: There is nothing stored inside, so just print an informative line.
\nwenddocs{}\nwbegincode{459}\moddef{MCI midpoint: test sampler 1: TBP}\endmoddef
  procedure :: write => test_sampler_1_write
\nwendcode{}\nwbegincode{460}\moddef{MCI midpoint: procedures}\plusendmoddef
  subroutine test_sampler_1_write (object, unit, testflag)
    class(test_sampler_1_t), intent(in) :: object
    integer, intent(in), optional :: unit
    logical, intent(in), optional :: testflag
    integer :: u
    u = given_output_unit (unit)
    write (u, "(1x,A)") "Test sampler: f(x) = 3 x^2"
  end subroutine test_sampler_1_write
  
\eatline
\nwendcode{}\nwbegindocs{461}Evaluation: compute the function value.  The output $x$ parameter
(only one channel) is identical to the input $x$, and the Jacobian is 1.
\nwenddocs{}\nwbegincode{462}\moddef{MCI midpoint: test sampler 1: TBP}\plusendmoddef
  procedure :: evaluate => test_sampler_1_evaluate
\nwendcode{}\nwbegincode{463}\moddef{MCI midpoint: procedures}\plusendmoddef
  subroutine test_sampler_1_evaluate (sampler, c, x_in, val, x, f)
    class(test_sampler_1_t), intent(inout) :: sampler
    integer, intent(in) :: c
    real(default), dimension(:), intent(in) :: x_in
    real(default), intent(out) :: val
    real(default), dimension(:,:), intent(out) :: x
    real(default), dimension(:), intent(out) :: f
    if (allocated (sampler%x))  deallocate (sampler%x)
    allocate (sampler%x (size (x_in)))
    sampler%x = x_in
    sampler%val = 3 * x_in(1) ** 2
    call sampler%fetch (val, x, f)
  end subroutine test_sampler_1_evaluate

\eatline
\nwendcode{}\nwbegindocs{464}The point is always valid.
\nwenddocs{}\nwbegincode{465}\moddef{MCI midpoint: test sampler 1: TBP}\plusendmoddef
  procedure :: is_valid => test_sampler_1_is_valid
\nwendcode{}\nwbegincode{466}\moddef{MCI midpoint: tests}\plusendmoddef
  function test_sampler_1_is_valid (sampler) result (valid)
    class(test_sampler_1_t), intent(in) :: sampler
    logical :: valid
    valid = .true.
  end function test_sampler_1_is_valid
  
\eatline
\nwendcode{}\nwbegindocs{467}Rebuild: compute all but the function value.
\nwenddocs{}\nwbegincode{468}\moddef{MCI midpoint: test sampler 1: TBP}\plusendmoddef
  procedure :: rebuild => test_sampler_1_rebuild
\nwendcode{}\nwbegincode{469}\moddef{MCI midpoint: procedures}\plusendmoddef
  subroutine test_sampler_1_rebuild (sampler, c, x_in, val, x, f)
    class(test_sampler_1_t), intent(inout) :: sampler
    integer, intent(in) :: c
    real(default), dimension(:), intent(in) :: x_in
    real(default), intent(in) :: val
    real(default), dimension(:,:), intent(out) :: x
    real(default), dimension(:), intent(out) :: f
    if (allocated (sampler%x))  deallocate (sampler%x)
    allocate (sampler%x (size (x_in)))
    sampler%x = x_in
    sampler%val = val
    x(:,1) = sampler%x
    f = 1
  end subroutine test_sampler_1_rebuild

\eatline
\nwendcode{}\nwbegindocs{470}Extract the results.
\nwenddocs{}\nwbegincode{471}\moddef{MCI midpoint: test sampler 1: TBP}\plusendmoddef
  procedure :: fetch => test_sampler_1_fetch
\nwendcode{}\nwbegincode{472}\moddef{MCI midpoint: procedures}\plusendmoddef
  subroutine test_sampler_1_fetch (sampler, val, x, f)
    class(test_sampler_1_t), intent(in) :: sampler
    real(default), intent(out) :: val
    real(default), dimension(:,:), intent(out) :: x
    real(default), dimension(:), intent(out) :: f
    val = sampler%val
    x(:,1) = sampler%x
    f = 1
  end subroutine test_sampler_1_fetch
    
\eatline
\nwendcode{}\nwbegindocs{473}\nwdocspar
This is the function $f(x) = 3 x^2 + 2 y$ with integral $\int_0^1
f(x,y)\,dx\,dy=2$ and maximum $f(1)=5$.
\nwenddocs{}\nwbegincode{474}\moddef{MCI midpoint: test types}\plusendmoddef
  type, extends (mci_sampler_t) :: test_sampler_2_t
     real(default) :: val
     real(default), dimension(2) :: x
   contains
   \LA{}MCI midpoint: test sampler 2: TBP\RA{}
  end type test_sampler_2_t

\eatline
\nwendcode{}\nwbegindocs{475}Output: There is nothing stored inside, so just print an informative line.
\nwenddocs{}\nwbegincode{476}\moddef{MCI midpoint: test sampler 2: TBP}\endmoddef
  procedure :: write => test_sampler_2_write
\nwendcode{}\nwbegincode{477}\moddef{MCI midpoint: procedures}\plusendmoddef
  subroutine test_sampler_2_write (object, unit, testflag)
    class(test_sampler_2_t), intent(in) :: object
    integer, intent(in), optional :: unit
    logical, intent(in), optional :: testflag
    integer :: u
    u = given_output_unit (unit)
    write (u, "(1x,A)") "Test sampler: f(x) = 3 x^2 + 2 y"
  end subroutine test_sampler_2_write
  
\eatline
\nwendcode{}\nwbegindocs{478}Evaluate: compute the function value.  The output $x$ parameter
(only one channel) is identical to the input $x$, and the Jacobian is 1.
\nwenddocs{}\nwbegincode{479}\moddef{MCI midpoint: test sampler 2: TBP}\plusendmoddef
  procedure :: evaluate => test_sampler_2_evaluate
\nwendcode{}\nwbegincode{480}\moddef{MCI midpoint: procedures}\plusendmoddef
  subroutine test_sampler_2_evaluate (sampler, c, x_in, val, x, f)
    class(test_sampler_2_t), intent(inout) :: sampler
    integer, intent(in) :: c
    real(default), dimension(:), intent(in) :: x_in
    real(default), intent(out) :: val
    real(default), dimension(:,:), intent(out) :: x
    real(default), dimension(:), intent(out) :: f
    sampler%x = x_in
    sampler%val = 3 * x_in(1) ** 2 + 2 * x_in(2)
    call sampler%fetch (val, x, f)
  end subroutine test_sampler_2_evaluate

\eatline
\nwendcode{}\nwbegindocs{481}The point is always valid.
\nwenddocs{}\nwbegincode{482}\moddef{MCI midpoint: test sampler 2: TBP}\plusendmoddef
  procedure :: is_valid => test_sampler_2_is_valid
\nwendcode{}\nwbegincode{483}\moddef{MCI midpoint: tests}\plusendmoddef
  function test_sampler_2_is_valid (sampler) result (valid)
    class(test_sampler_2_t), intent(in) :: sampler
    logical :: valid
    valid = .true.
  end function test_sampler_2_is_valid
  
\eatline
\nwendcode{}\nwbegindocs{484}Rebuild: compute all but the function value.
\nwenddocs{}\nwbegincode{485}\moddef{MCI midpoint: test sampler 2: TBP}\plusendmoddef
  procedure :: rebuild => test_sampler_2_rebuild
\nwendcode{}\nwbegincode{486}\moddef{MCI midpoint: procedures}\plusendmoddef
  subroutine test_sampler_2_rebuild (sampler, c, x_in, val, x, f)
    class(test_sampler_2_t), intent(inout) :: sampler
    integer, intent(in) :: c
    real(default), dimension(:), intent(in) :: x_in
    real(default), intent(in) :: val
    real(default), dimension(:,:), intent(out) :: x
    real(default), dimension(:), intent(out) :: f
    sampler%x = x_in
    sampler%val = val
    x(:,1) = sampler%x
    f = 1
  end subroutine test_sampler_2_rebuild

\eatline
\nwendcode{}\nwbegincode{487}\moddef{MCI midpoint: test sampler 2: TBP}\plusendmoddef
  procedure :: fetch => test_sampler_2_fetch
\nwendcode{}\nwbegincode{488}\moddef{MCI midpoint: procedures}\plusendmoddef
  subroutine test_sampler_2_fetch (sampler, val, x, f)
    class(test_sampler_2_t), intent(in) :: sampler
    real(default), intent(out) :: val
    real(default), dimension(:,:), intent(out) :: x
    real(default), dimension(:), intent(out) :: f
    val = sampler%val
    x(:,1) = sampler%x
    f = 1
  end subroutine test_sampler_2_fetch
  
\eatline
\nwendcode{}\nwbegindocs{489}\nwdocspar
This is the function $f(x) = (1 - 3 x^2)\,\theta(x-1/2)$ with integral
$\int_0^1 f(x)\,dx=-3/8$, minimum $f(1)=-2$ and maximum $f(1/2)=1/4$.
If the integration dimension is greater than one, the function is
extended as a constant in the other dimension(s).
\nwenddocs{}\nwbegincode{490}\moddef{MCI midpoint: test types}\plusendmoddef
  type, extends (mci_sampler_t) :: test_sampler_4_t
     real(default) :: val
     real(default), dimension(:), allocatable :: x
   contains
   \LA{}MCI midpoint: test sampler 4: TBP\RA{}
  end type test_sampler_4_t

\eatline
\nwendcode{}\nwbegindocs{491}Output: There is nothing stored inside, so just print an informative line.
\nwenddocs{}\nwbegincode{492}\moddef{MCI midpoint: test sampler 4: TBP}\endmoddef
  procedure :: write => test_sampler_4_write
\nwendcode{}\nwbegincode{493}\moddef{MCI midpoint: procedures}\plusendmoddef
  subroutine test_sampler_4_write (object, unit, testflag)
    class(test_sampler_4_t), intent(in) :: object
    integer, intent(in), optional :: unit
    logical, intent(in), optional :: testflag
    integer :: u
    u = given_output_unit (unit)
    write (u, "(1x,A)") "Test sampler: f(x) = 1 - 3 x^2"
  end subroutine test_sampler_4_write
  
\eatline
\nwendcode{}\nwbegindocs{494}Evaluation: compute the function value.  The output $x$ parameter
(only one channel) is identical to the input $x$, and the Jacobian is 1.
\nwenddocs{}\nwbegincode{495}\moddef{MCI midpoint: test sampler 4: TBP}\plusendmoddef
  procedure :: evaluate => test_sampler_4_evaluate
\nwendcode{}\nwbegincode{496}\moddef{MCI midpoint: procedures}\plusendmoddef
  subroutine test_sampler_4_evaluate (sampler, c, x_in, val, x, f)
    class(test_sampler_4_t), intent(inout) :: sampler
    integer, intent(in) :: c
    real(default), dimension(:), intent(in) :: x_in
    real(default), intent(out) :: val
    real(default), dimension(:,:), intent(out) :: x
    real(default), dimension(:), intent(out) :: f
    if (x_in(1) >= .5_default) then
       sampler%val = 1 - 3 * x_in(1) ** 2
    else
       sampler%val = 0
    end if
    if (.not. allocated (sampler%x))  allocate (sampler%x (size (x_in)))
    sampler%x = x_in
    call sampler%fetch (val, x, f)
  end subroutine test_sampler_4_evaluate

\eatline
\nwendcode{}\nwbegindocs{497}The point is always valid.
\nwenddocs{}\nwbegincode{498}\moddef{MCI midpoint: test sampler 4: TBP}\plusendmoddef
  procedure :: is_valid => test_sampler_4_is_valid
\nwendcode{}\nwbegincode{499}\moddef{MCI midpoint: tests}\plusendmoddef
  function test_sampler_4_is_valid (sampler) result (valid)
    class(test_sampler_4_t), intent(in) :: sampler
    logical :: valid
    valid = .true.
  end function test_sampler_4_is_valid
  
\eatline
\nwendcode{}\nwbegindocs{500}Rebuild: compute all but the function value.
\nwenddocs{}\nwbegincode{501}\moddef{MCI midpoint: test sampler 4: TBP}\plusendmoddef
  procedure :: rebuild => test_sampler_4_rebuild
\nwendcode{}\nwbegincode{502}\moddef{MCI midpoint: procedures}\plusendmoddef
  subroutine test_sampler_4_rebuild (sampler, c, x_in, val, x, f)
    class(test_sampler_4_t), intent(inout) :: sampler
    integer, intent(in) :: c
    real(default), dimension(:), intent(in) :: x_in
    real(default), intent(in) :: val
    real(default), dimension(:,:), intent(out) :: x
    real(default), dimension(:), intent(out) :: f
    sampler%x = x_in
    sampler%val = val
    x(:,1) = sampler%x
    f = 1
  end subroutine test_sampler_4_rebuild

\eatline
\nwendcode{}\nwbegincode{503}\moddef{MCI midpoint: test sampler 4: TBP}\plusendmoddef
  procedure :: fetch => test_sampler_4_fetch
\nwendcode{}\nwbegincode{504}\moddef{MCI midpoint: procedures}\plusendmoddef
  subroutine test_sampler_4_fetch (sampler, val, x, f)
    class(test_sampler_4_t), intent(in) :: sampler
    real(default), intent(out) :: val
    real(default), dimension(:,:), intent(out) :: x
    real(default), dimension(:), intent(out) :: f
    val = sampler%val
    x(:,1) = sampler%x
    f = 1
  end subroutine test_sampler_4_fetch
    
\eatline
\nwendcode{}\nwbegindocs{505}\nwdocspar
\subsubsection{One-dimensional integration}
Construct an integrator and use it for a one-dimensional sampler.
\nwenddocs{}\nwbegincode{506}\moddef{MCI midpoint: execute tests}\endmoddef
  call test (mci_midpoint_1, "mci_midpoint_1", &
       "one-dimensional integral", &
       u, results)
\nwendcode{}\nwbegincode{507}\moddef{MCI midpoint: tests}\plusendmoddef
  subroutine mci_midpoint_1 (u)
    integer, intent(in) :: u
    class(mci_t), allocatable, target :: mci
    class(mci_instance_t), pointer :: mci_instance => null ()
    class(mci_sampler_t), allocatable :: sampler
    
    write (u, "(A)")  "* Test output: mci_midpoint_1"
    write (u, "(A)")  "*   Purpose: integrate function in one dimension"
    write (u, "(A)")
    
    write (u, "(A)")  "* Initialize integrator"
    write (u, "(A)")

    allocate (mci_midpoint_t :: mci)
    call mci%set_dimensions (1, 1)
    
    call mci%write (u)
    
    write (u, "(A)")
    write (u, "(A)")  "* Initialize instance"
    write (u, "(A)")
    
    call mci%allocate_instance (mci_instance)
    call mci_instance%init (mci)
    
    write (u, "(A)")  "* Initialize test sampler"
    write (u, "(A)")
    
    allocate (test_sampler_1_t :: sampler)
    call sampler%write (u)
     
    write (u, "(A)")
    write (u, "(A)")  "* Evaluate for x = 0.8"
    write (u, "(A)")
    
    call mci_instance%evaluate (sampler, 1, [0.8_default])
    call mci_instance%write (u)
    
    write (u, "(A)")
    write (u, "(A)")  "* Evaluate for x = 0.7"
    write (u, "(A)")
    
    call mci_instance%evaluate (sampler, 1, [0.7_default])
    call mci_instance%write (u)
    
    write (u, "(A)")
    write (u, "(A)")  "* Evaluate for x = 0.9"
    write (u, "(A)")
    
    call mci_instance%evaluate (sampler, 1, [0.9_default])
    call mci_instance%write (u)
    
    write (u, "(A)")
    write (u, "(A)")  "* Integrate with n_calls = 1000"
    write (u, "(A)")
    
    call mci%integrate (mci_instance, sampler, 1, 1000)
    call mci%write (u)

    call mci_instance%final ()
    call mci%final ()

    write (u, "(A)")
    write (u, "(A)")  "* Test output end: mci_midpoint_1"

  end subroutine mci_midpoint_1

\eatline
\nwendcode{}\nwbegindocs{508}\nwdocspar
\subsubsection{Two-dimensional integration}
Construct an integrator and use it for a two-dimensional sampler.
\nwenddocs{}\nwbegincode{509}\moddef{MCI midpoint: execute tests}\plusendmoddef
  call test (mci_midpoint_2, "mci_midpoint_2", &
       "two-dimensional integral", &
       u, results)
\nwendcode{}\nwbegincode{510}\moddef{MCI midpoint: tests}\plusendmoddef
  subroutine mci_midpoint_2 (u)
    integer, intent(in) :: u
    class(mci_t), allocatable, target :: mci
    class(mci_instance_t), pointer :: mci_instance => null ()
    class(mci_sampler_t), allocatable :: sampler
    
    write (u, "(A)")  "* Test output: mci_midpoint_2"
    write (u, "(A)")  "*   Purpose: integrate function in two dimensions"
    write (u, "(A)")
    
    write (u, "(A)")  "* Initialize integrator"
    write (u, "(A)")

    allocate (mci_midpoint_t :: mci)
    call mci%set_dimensions (2, 1)
    
    call mci%write (u)
    
    write (u, "(A)")
    write (u, "(A)")  "* Initialize instance"
    write (u, "(A)")
    
    call mci%allocate_instance (mci_instance)
    call mci_instance%init (mci)
    
    write (u, "(A)")  "* Initialize test sampler"
    write (u, "(A)")
    
    allocate (test_sampler_2_t :: sampler)
    call sampler%write (u)
     
    write (u, "(A)")
    write (u, "(A)")  "* Evaluate for x = 0.8, y = 0.2"
    write (u, "(A)")
    
    call mci_instance%evaluate (sampler, 1, [0.8_default, 0.2_default])
    call mci_instance%write (u)
    
    write (u, "(A)")
    write (u, "(A)")  "* Integrate with n_calls = 1000"
    write (u, "(A)")
    
    call mci%integrate (mci_instance, sampler, 1, 1000)
    call mci%write (u)

    call mci_instance%final ()
    call mci%final ()

    write (u, "(A)")
    write (u, "(A)")  "* Test output end: mci_midpoint_2"

  end subroutine mci_midpoint_2

\eatline
\nwendcode{}\nwbegindocs{511}\nwdocspar
\subsubsection{Two-dimensional integration with flat dimension}
Construct an integrator and use it for a two-dimensional sampler,
where the function is constant in the second dimension.
\nwenddocs{}\nwbegincode{512}\moddef{MCI midpoint: execute tests}\plusendmoddef
  call test (mci_midpoint_3, "mci_midpoint_3", &
       "two-dimensional integral with flat dimension", &
       u, results)
\nwendcode{}\nwbegincode{513}\moddef{MCI midpoint: tests}\plusendmoddef
  subroutine mci_midpoint_3 (u)
    integer, intent(in) :: u
    class(mci_t), allocatable, target :: mci
    class(mci_instance_t), pointer :: mci_instance => null ()
    class(mci_sampler_t), allocatable :: sampler
    
    write (u, "(A)")  "* Test output: mci_midpoint_3"
    write (u, "(A)")  "*   Purpose: integrate function with one flat dimension"
    write (u, "(A)")
    
    write (u, "(A)")  "* Initialize integrator"
    write (u, "(A)")

    allocate (mci_midpoint_t :: mci)
    select type (mci)
    type is (mci_midpoint_t)
       call mci%set_dimensions (2, 1)
       call mci%declare_flat_dimensions ([2])
    end select
    
    call mci%write (u)
    
    write (u, "(A)")
    write (u, "(A)")  "* Initialize instance"
    write (u, "(A)")
    
    call mci%allocate_instance (mci_instance)
    call mci_instance%init (mci)
    
    write (u, "(A)")  "* Initialize test sampler"
    write (u, "(A)")
    
    allocate (test_sampler_1_t :: sampler)
    call sampler%write (u)
     
    write (u, "(A)")
    write (u, "(A)")  "* Evaluate for x = 0.8, y = 0.2"
    write (u, "(A)")
    
    call mci_instance%evaluate (sampler, 1, [0.8_default, 0.2_default])
    call mci_instance%write (u)
    
    write (u, "(A)")
    write (u, "(A)")  "* Integrate with n_calls = 1000"
    write (u, "(A)")
    
    call mci%integrate (mci_instance, sampler, 1, 1000)
    call mci%write (u)

    call mci_instance%final ()
    call mci%final ()

    write (u, "(A)")
    write (u, "(A)")  "* Test output end: mci_midpoint_3"

  end subroutine mci_midpoint_3

\eatline
\nwendcode{}\nwbegindocs{514}\nwdocspar
\subsubsection{Integrand with sign flip}
Construct an integrator and use it for a one-dimensional sampler.
\nwenddocs{}\nwbegincode{515}\moddef{MCI midpoint: execute tests}\plusendmoddef
  call test (mci_midpoint_4, "mci_midpoint_4", &
       "integrand with sign flip", &
       u, results)
\nwendcode{}\nwbegincode{516}\moddef{MCI midpoint: tests}\plusendmoddef
  subroutine mci_midpoint_4 (u)
    integer, intent(in) :: u
    class(mci_t), allocatable, target :: mci
    class(mci_instance_t), pointer :: mci_instance => null ()
    class(mci_sampler_t), allocatable :: sampler
    
    write (u, "(A)")  "* Test output: mci_midpoint_4"
    write (u, "(A)")  "*   Purpose: integrate function with sign flip &
         &in one dimension"
    write (u, "(A)")
    
    write (u, "(A)")  "* Initialize integrator"
    write (u, "(A)")

    allocate (mci_midpoint_t :: mci)
    call mci%set_dimensions (1, 1)
    
    call mci%write (u)
    
    write (u, "(A)")
    write (u, "(A)")  "* Initialize instance"
    write (u, "(A)")
    
    call mci%allocate_instance (mci_instance)
    call mci_instance%init (mci)
    
    write (u, "(A)")  "* Initialize test sampler"
    write (u, "(A)")
    
    allocate (test_sampler_4_t :: sampler)
    call sampler%write (u)
     
    write (u, "(A)")
    write (u, "(A)")  "* Evaluate for x = 0.8"
    write (u, "(A)")
    
    call mci_instance%evaluate (sampler, 1, [0.8_default])
    call mci_instance%write (u)
    
    write (u, "(A)")
    write (u, "(A)")  "* Integrate with n_calls = 1000"
    write (u, "(A)")
    
    call mci%integrate (mci_instance, sampler, 1, 1000)
    call mci%write (u)

    call mci_instance%final ()
    call mci%final ()

    write (u, "(A)")
    write (u, "(A)")  "* Test output end: mci_midpoint_4"

  end subroutine mci_midpoint_4

\eatline
\nwendcode{}\nwbegindocs{517}\nwdocspar
\subsubsection{Weighted events}
Generate weighted events.  Without rejection, we do not need to know maxima
and minima, so we can start generating events immediately.  We have two
dimensions.
\nwenddocs{}\nwbegincode{518}\moddef{MCI midpoint: execute tests}\plusendmoddef
  call test (mci_midpoint_5, "mci_midpoint_5", &
       "weighted events", &
       u, results)
\nwendcode{}\nwbegincode{519}\moddef{MCI midpoint: tests}\plusendmoddef
  subroutine mci_midpoint_5 (u)
    integer, intent(in) :: u
    class(mci_t), allocatable, target :: mci
    class(mci_instance_t), pointer :: mci_instance => null ()
    class(mci_sampler_t), allocatable :: sampler
    class(rng_t), allocatable :: rng
    class(mci_state_t), allocatable :: state

    write (u, "(A)")  "* Test output: mci_midpoint_5"
    write (u, "(A)")  "*   Purpose: generate weighted events"
    write (u, "(A)")
    
    write (u, "(A)")  "* Initialize integrator"
    write (u, "(A)")

    allocate (mci_midpoint_t :: mci)
    call mci%set_dimensions (2, 1)
    
    call mci%write (u)
    
    write (u, "(A)")
    write (u, "(A)")  "* Initialize instance"
    write (u, "(A)")
    
    call mci%allocate_instance (mci_instance)
    call mci_instance%init (mci)
    
    write (u, "(A)")  "* Initialize test sampler"
    write (u, "(A)")
    
    allocate (test_sampler_2_t :: sampler)

    write (u, "(A)")  "* Initialize random-number generator"
    write (u, "(A)")
    
    allocate (rng_test_t :: rng)
    call rng%init ()
    call mci%import_rng (rng)
    
    write (u, "(A)")  "* Generate weighted event"
    write (u, "(A)")
    
    call mci%generate_weighted_event (mci_instance, sampler)
    call mci_instance%write (u)
    
    write (u, "(A)")
    write (u, "(A)")  "* Generate weighted event"
    write (u, "(A)")
    
    call mci%generate_weighted_event (mci_instance, sampler)
    call mci_instance%write (u)
    
    write (u, "(A)")
    write (u, "(A)")  "* Store data"
    write (u, "(A)")
    
    allocate (state)
    call mci_instance%store (state)
    call mci_instance%final ()
    deallocate (mci_instance)
    
    call state%write (u)
    
    write (u, "(A)")
    write (u, "(A)")  "* Recall data and rebuild event"
    write (u, "(A)")
    
    call mci%allocate_instance (mci_instance)
    call mci_instance%init (mci)
    call mci%rebuild_event (mci_instance, sampler, state)

    call mci_instance%write (u)
    
    write (u, "(A)")
    write (u, "(A)")  "* Cleanup"

    call mci_instance%final ()
    deallocate (mci_instance)
    call mci%final ()

    write (u, "(A)")
    write (u, "(A)")  "* Test output end: mci_midpoint_5"

  end subroutine mci_midpoint_5
    
\eatline
\nwendcode{}\nwbegindocs{520}\nwdocspar
\subsubsection{Unweighted events}
Generate unweighted events.  The integrand has a sign flip in it.
\nwenddocs{}\nwbegincode{521}\moddef{MCI midpoint: execute tests}\plusendmoddef
  call test (mci_midpoint_6, "mci_midpoint_6", &
       "unweighted events", &
       u, results)
\nwendcode{}\nwbegincode{522}\moddef{MCI midpoint: tests}\plusendmoddef
  subroutine mci_midpoint_6 (u)
    integer, intent(in) :: u
    class(mci_t), allocatable, target :: mci
    class(mci_instance_t), pointer :: mci_instance => null ()
    class(mci_sampler_t), allocatable :: sampler
    class(rng_t), allocatable :: rng

    write (u, "(A)")  "* Test output: mci_midpoint_6"
    write (u, "(A)")  "*   Purpose: generate unweighted events"
    write (u, "(A)")
    
    write (u, "(A)")  "* Initialize integrator"
    write (u, "(A)")

    allocate (mci_midpoint_t :: mci)
    call mci%set_dimensions (1, 1)
    
    write (u, "(A)")  "* Initialize instance"
    write (u, "(A)")
    
    call mci%allocate_instance (mci_instance)
    call mci_instance%init (mci)
    
    write (u, "(A)")  "* Initialize test sampler"
    write (u, "(A)")
    
    allocate (test_sampler_4_t :: sampler)

    write (u, "(A)")  "* Initialize random-number generator"
    write (u, "(A)")
    
    allocate (rng_test_t :: rng)
    call rng%init ()
    call mci%import_rng (rng)
    
    write (u, "(A)")  "* Integrate (determine maximum of integrand"
    write (u, "(A)")
    call mci%integrate (mci_instance, sampler, 1, 1000)
    call mci%write (u)
    
    write (u, "(A)")
    write (u, "(A)")  "* Generate unweighted event"
    write (u, "(A)")
    
    call mci%generate_unweighted_event (mci_instance, sampler)
    call mci_instance%write (u)
    
    write (u, "(A)")
    write (u, "(A)")  "* Cleanup"

    call mci_instance%final ()
    deallocate (mci_instance)
    call mci%final ()

    write (u, "(A)")
    write (u, "(A)")  "* Test output end: mci_midpoint_6"

  end subroutine mci_midpoint_6
    
\eatline
\nwendcode{}\nwbegindocs{523}\nwdocspar
\subsubsection{Excess weight}
Generate unweighted events.  With only 2 points for integration, the
maximum of the integrand is too low, and we produce excess weight.
\nwenddocs{}\nwbegincode{524}\moddef{MCI midpoint: execute tests}\plusendmoddef
  call test (mci_midpoint_7, "mci_midpoint_7", &
       "excess weight", &
       u, results)
\nwendcode{}\nwbegincode{525}\moddef{MCI midpoint: tests}\plusendmoddef
  subroutine mci_midpoint_7 (u)
    integer, intent(in) :: u
    class(mci_t), allocatable, target :: mci
    class(mci_instance_t), pointer :: mci_instance => null ()
    class(mci_sampler_t), allocatable :: sampler
    class(rng_t), allocatable :: rng

    write (u, "(A)")  "* Test output: mci_midpoint_7"
    write (u, "(A)")  "*   Purpose: generate unweighted event &
         &with excess weight"
    write (u, "(A)")
    
    write (u, "(A)")  "* Initialize integrator"
    write (u, "(A)")

    allocate (mci_midpoint_t :: mci)
    call mci%set_dimensions (1, 1)
    
    write (u, "(A)")  "* Initialize instance"
    write (u, "(A)")
    
    call mci%allocate_instance (mci_instance)
    call mci_instance%init (mci)
    
    write (u, "(A)")  "* Initialize test sampler"
    write (u, "(A)")
    
    allocate (test_sampler_4_t :: sampler)

    write (u, "(A)")  "* Initialize random-number generator"
    write (u, "(A)")
    
    allocate (rng_test_t :: rng)
    call rng%init ()
    call mci%import_rng (rng)
    
    write (u, "(A)")  "* Integrate (determine maximum of integrand"
    write (u, "(A)")
    call mci%integrate (mci_instance, sampler, 1, 2)
    call mci%write (u)
    
    write (u, "(A)")
    write (u, "(A)")  "* Generate unweighted event"
    write (u, "(A)")
    
    call mci_instance%init_simulation ()
    call mci%generate_unweighted_event (mci_instance, sampler)
    call mci_instance%write (u)
    
    write (u, "(A)")
    write (u, "(A)")  "* Use getter methods"
    write (u, "(A)")
    
    write (u, "(1x,A,1x,ES19.12)")  "weight =", mci_instance%get_event_weight ()
    write (u, "(1x,A,1x,ES19.12)")  "excess =", mci_instance%get_event_excess ()

    write (u, "(A)")
    write (u, "(A)")  "* Apply safety factor"
    write (u, "(A)")

    call mci_instance%init_simulation (safety_factor = 2.1_default)

    write (u, "(A)")  "* Generate unweighted event"
    write (u, "(A)")
    
    call mci%generate_unweighted_event (mci_instance, sampler)
    call mci_instance%write (u)
    
    write (u, "(A)")
    write (u, "(A)")  "* Use getter methods"
    write (u, "(A)")
    
    write (u, "(1x,A,1x,ES19.12)")  "weight =", mci_instance%get_event_weight ()
    write (u, "(1x,A,1x,ES19.12)")  "excess =", mci_instance%get_event_excess ()

    write (u, "(A)")
    write (u, "(A)")  "* Cleanup"

    call mci_instance%final ()
    deallocate (mci_instance)
    call mci%final ()

    write (u, "(A)")
    write (u, "(A)")  "* Test output end: mci_midpoint_7"

  end subroutine mci_midpoint_7
    
\eatline
\nwendcode{}\nwbegindocs{526}\nwdocspar
\clearpage
%%%%%%%%%%%%%%%%%%%%%%%%%%%%%%%%%%%%%%%%%%%%%%%%%%%%%%%%%%%%%%%%%%%%%%%%
\section{\vamp\ interface}
The standard method for integration is \vamp: the multi-channel
version of the VEGAS algorithm.  Each parameterization (channel) of
the hypercube is binned in each dimension.  The binning is equally
equidistant, but an iteration of the integration procedure,
the binning is updated for each dimension, according to the variance
distribution of the integrand, summed over all other dimension.  In
the next iteration, the binning approximates (hopefully) follows the
integrand more closely, and the accuracy of the result is increased.
Furthermore, the relative weight of the individual channels is also updated
after an iteration.

The bin distribution is denoted as the grid for a channel, which we
can write to file and reuse later.

In our implementation we specify the generic \vamp\ algorithm more
tightly: the number of bins is equal for all dimensions, the initial
weights are all equal.  The user controls whether to update bins
and/or weights after each iteration.  The integration is organized in
passes, each one consisting of several iterations with a common number
of calls to the integrand.  The first passes are intended as warmup,
so the results are displayed but otherwise discarded.  In the final
pass, the integration estimates for the individual iterations are
averaged for the final result.
\nwenddocs{}\nwbegincode{527}\moddef{\code{}mci{\char95}vamp.f90\edoc{}}\endmoddef
\LA{}File header\RA{}

module mci_vamp

  use kinds
\LA{}Use strings\RA{}
  use io_units
  use constants
  use format_utils, only: pac_fmt
  use format_defs, only: FMT_12, FMT_14, FMT_17, FMT_19
  use unit_tests
  use diagnostics
  use md5

  use phs_base
  use rng_base
  use rng_tao
  use mci_base
  
  use vamp !NODEP!
  use exceptions !NODEP!
  
\LA{}Standard module head\RA{}

\LA{}MCI vamp: public\RA{}

\LA{}MCI vamp: types\RA{}

\LA{}MCI vamp: interfaces\RA{}

\LA{}MCI vamp: test types\RA{}

contains
  
\LA{}MCI vamp: procedures\RA{}

\LA{}MCI vamp: tests\RA{}

end module mci_vamp
\eatline
\nwendcode{}\nwbegindocs{528}\nwdocspar
\subsection{Grid parameters}
This is a transparent container.  It holds the parameters that are
stored in grid files, and are checked when grid files are read.
\nwenddocs{}\nwbegincode{529}\moddef{MCI vamp: public}\endmoddef
  public :: grid_parameters_t
\nwendcode{}\nwbegincode{530}\moddef{MCI vamp: types}\endmoddef
  type :: grid_parameters_t
     integer :: threshold_calls = 0
     integer :: min_calls_per_channel = 10
     integer :: min_calls_per_bin = 10
     integer :: min_bins = 3
     integer :: max_bins = 20
     logical :: stratified = .true.
     logical :: use_vamp_equivalences = .true.
     real(default) :: channel_weights_power = 0.25_default
     real(default) :: accuracy_goal = 0
     real(default) :: error_goal = 0
     real(default) :: rel_error_goal = 0
   contains
   \LA{}MCI vamp: grid parameters: TBP\RA{}
  end type grid_parameters_t

\eatline
\nwendcode{}\nwbegindocs{531}I/O:
\nwenddocs{}\nwbegincode{532}\moddef{MCI vamp: grid parameters: TBP}\endmoddef
  procedure :: write => grid_parameters_write
\nwendcode{}\nwbegincode{533}\moddef{MCI vamp: procedures}\endmoddef
  subroutine grid_parameters_write (object, unit)
    class(grid_parameters_t), intent(in) :: object
    integer, intent(in), optional :: unit
    integer :: u
    u = given_output_unit (unit)
    write (u, "(3x,A,I0)") "threshold_calls       = ", &
         object%threshold_calls 
    write (u, "(3x,A,I0)") "min_calls_per_channel = ", &
         object%min_calls_per_channel
    write (u, "(3x,A,I0)") "min_calls_per_bin     = ", &
         object%min_calls_per_bin
    write (u, "(3x,A,I0)") "min_bins              = ", &
         object%min_bins
    write (u, "(3x,A,I0)") "max_bins              = ", &
         object%max_bins
    write (u, "(3x,A,L1)") "stratified            = ", &
         object%stratified
    write (u, "(3x,A,L1)") "use_vamp_equivalences = ", &
         object%use_vamp_equivalences
    write (u, "(3x,A,F10.7)") "channel_weights_power = ", &
         object%channel_weights_power
    if (object%accuracy_goal > 0) then
       write (u, "(3x,A,F10.7)") "accuracy_goal         = ", &
            object%accuracy_goal
    end if
    if (object%error_goal > 0) then
       write (u, "(3x,A,F10.7)") "error_goal            = ", &
            object%error_goal
    end if
    if (object%rel_error_goal > 0) then
       write (u, "(3x,A,F10.7)") "rel_error_goal        = ", &
            object%rel_error_goal
    end if
  end subroutine grid_parameters_write

\eatline
\nwendcode{}\nwbegindocs{534}\nwdocspar
\subsection{History parameters}
The history parameters are also stored in a transparent container.
This is not a part of the grid definition, and should not be included
in the MD5 sum.
\nwenddocs{}\nwbegincode{535}\moddef{MCI vamp: public}\plusendmoddef
  public :: history_parameters_t
\nwendcode{}\nwbegincode{536}\moddef{MCI vamp: types}\plusendmoddef
  type :: history_parameters_t
     logical :: global = .true.
     logical :: global_verbose = .false.
     logical :: channel = .false.
     logical :: channel_verbose = .false.
   contains
   \LA{}MCI vamp: history parameters: TBP\RA{}
  end type history_parameters_t

\eatline
\nwendcode{}\nwbegindocs{537}I/O:
\nwenddocs{}\nwbegincode{538}\moddef{MCI vamp: history parameters: TBP}\endmoddef
  procedure :: write => history_parameters_write
\nwendcode{}\nwbegincode{539}\moddef{MCI vamp: procedures}\plusendmoddef
  subroutine history_parameters_write (object, unit)
    class(history_parameters_t), intent(in) :: object
    integer, intent(in), optional :: unit
    integer :: u
    u = given_output_unit (unit)
    write (u, "(3x,A,L1)") "history(global)       = ", object%global
    write (u, "(3x,A,L1)") "history(global) verb. = ", object%global_verbose
    write (u, "(3x,A,L1)") "history(channels)     = ", object%channel
    write (u, "(3x,A,L1)") "history(chann.) verb. = ", object%channel_verbose
  end subroutine history_parameters_write

\eatline
\nwendcode{}\nwbegindocs{540}\nwdocspar
\subsection{Integration pass}
We store the parameters for each integration pass in a linked list.
\nwenddocs{}\nwbegincode{541}\moddef{MCI vamp: types}\plusendmoddef
  type :: pass_t
     integer :: i_pass = 0
     integer :: i_first_it = 0
     integer :: n_it = 0
     integer :: n_calls = 0
     integer :: n_bins = 0
     logical :: adapt_grids = .false.
     logical :: adapt_weights = .false.
     logical :: is_final_pass = .false.
     logical :: integral_defined = .false.
     integer, dimension(:), allocatable :: calls
     real(default), dimension(:), allocatable :: integral
     real(default), dimension(:), allocatable :: error
     real(default), dimension(:), allocatable :: efficiency
     type(vamp_history), dimension(:), allocatable :: v_history
     type(vamp_history), dimension(:,:), allocatable :: v_histories
     type(pass_t), pointer :: next => null ()
   contains
   \LA{}MCI vamp: pass: TBP\RA{}
  end type pass_t
     
\eatline
\nwendcode{}\nwbegindocs{542}Finalizer.  The VAMP histories contain a pointer array.
\nwenddocs{}\nwbegincode{543}\moddef{MCI vamp: pass: TBP}\endmoddef
  procedure :: final => pass_final
\nwendcode{}\nwbegincode{544}\moddef{MCI vamp: procedures}\plusendmoddef
  subroutine pass_final (object)
    class(pass_t), intent(inout) :: object
    if (allocated (object%v_history)) then
       call vamp_delete_history (object%v_history)
    end if
    if (allocated (object%v_histories)) then
       call vamp_delete_history (object%v_histories)
    end if
  end subroutine pass_final
  
\eatline
\nwendcode{}\nwbegindocs{545}Output.  Note that the precision of the numerical values should match the
precision for comparing output from file with data.
\nwenddocs{}\nwbegincode{546}\moddef{MCI vamp: pass: TBP}\plusendmoddef
  procedure :: write => pass_write
\nwendcode{}\nwbegincode{547}\moddef{MCI vamp: procedures}\plusendmoddef
  subroutine pass_write (object, unit, pacify)
    class(pass_t), intent(in) :: object
    integer, intent(in) :: unit
    logical, intent(in), optional :: pacify
    integer :: u, i
    character(len=7) :: fmt
    call pac_fmt (fmt, FMT_17, FMT_14, pacify)
    u = given_output_unit (unit)
    write (u, "(3x,A,I0)")  "n_it          = ", object%n_it
    write (u, "(3x,A,I0)")  "n_calls       = ", object%n_calls
    write (u, "(3x,A,I0)")  "n_bins        = ", object%n_bins
    write (u, "(3x,A,L1)")  "adapt grids   = ", object%adapt_grids
    write (u, "(3x,A,L1)")  "adapt weights = ", object%adapt_weights
    if (object%integral_defined) then
       write (u, "(3x,A)")  "Results:  [it, calls, integral, error, efficiency]"
       do i = 1, object%n_it
          write (u, "(5x,I0,1x,I0,3(1x," // fmt // "))") &
               i, object%calls(i), object%integral(i), object%error(i), &
               object%efficiency(i)
       end do
    else
       write (u, "(3x,A)")  "Results: [undefined]"
    end if
  end subroutine pass_write
  
\eatline
\nwendcode{}\nwbegindocs{548}Read and reconstruct the pass.
\nwenddocs{}\nwbegincode{549}\moddef{MCI vamp: pass: TBP}\plusendmoddef
  procedure :: read => pass_read
\nwendcode{}\nwbegincode{550}\moddef{MCI vamp: procedures}\plusendmoddef
  subroutine pass_read (object, u, n_pass, n_it)
    class(pass_t), intent(out) :: object
    integer, intent(in) :: u, n_pass, n_it
    integer :: i, j
    character(80) :: buffer
    object%i_pass = n_pass + 1
    object%i_first_it = n_it + 1
    call read_ival (u, object%n_it)
    call read_ival (u, object%n_calls)
    call read_ival (u, object%n_bins)
    call read_lval (u, object%adapt_grids)
    call read_lval (u, object%adapt_weights)
    allocate (object%calls (object%n_it), source = 0)
    allocate (object%integral (object%n_it), source = 0._default)
    allocate (object%error (object%n_it), source = 0._default)
    allocate (object%efficiency (object%n_it), source = 0._default)
    read (u, "(A)")  buffer
    select case (trim (adjustl (buffer)))
    case ("Results:  [it, calls, integral, error, efficiency]")
       do i = 1, object%n_it
          read (u, *) &
               j, object%calls(i), object%integral(i), object%error(i), &
               object%efficiency(i)
       end do
       object%integral_defined = .true.
    case ("Results: [undefined]")
       object%integral_defined = .false.
    case default
       call msg_fatal ("Reading integration pass: corrupted file")
    end select
  end subroutine pass_read
  
\eatline
\nwendcode{}\nwbegindocs{551}Write the VAMP history for this pass.  (The subroutine writes the
whole array at once.)
\nwenddocs{}\nwbegincode{552}\moddef{MCI vamp: pass: TBP}\plusendmoddef
  procedure :: write_history => pass_write_history
\nwendcode{}\nwbegincode{553}\moddef{MCI vamp: procedures}\plusendmoddef
  subroutine pass_write_history (pass, unit)
    class(pass_t), intent(in) :: pass
    integer, intent(in), optional :: unit
    integer :: u
    u = given_output_unit (unit)
    if (allocated (pass%v_history)) then
       call vamp_write_history (u, pass%v_history)
    else
       write (u, "(1x,A)")  "Global history: [undefined]"
    end if
    if (allocated (pass%v_histories)) then
       write (u, "(1x,A)")  "Channel histories:"
       call vamp_write_history (u, pass%v_histories)
    else
       write (u, "(1x,A)")  "Channel histories: [undefined]"
    end if
  end subroutine pass_write_history
    
\eatline
\nwendcode{}\nwbegindocs{554}Given a number of calls and iterations, compute remaining data.
\nwenddocs{}\nwbegincode{555}\moddef{MCI vamp: pass: TBP}\plusendmoddef
  procedure :: configure => pass_configure
\nwendcode{}\nwbegincode{556}\moddef{MCI vamp: procedures}\plusendmoddef
  subroutine pass_configure (pass, n_it, n_calls, min_calls, &
       min_bins, max_bins, min_channel_calls)
    class(pass_t), intent(inout) :: pass
    integer, intent(in) :: n_it, n_calls, min_channel_calls
    integer, intent(in) :: min_calls, min_bins, max_bins
    pass%n_it = n_it    
    if (min_calls /= 0) then
       pass%n_bins =  max (min_bins, &
            min (n_calls / min_calls, max_bins))
    else
       pass%n_bins = max_bins
    end if
    pass%n_calls = max (n_calls, max (min_calls, min_channel_calls))
    if (pass%n_calls /= n_calls) then
       write (msg_buffer, "(A,I0)")  "VAMP: too few calls, resetting " &
            // "n_calls to ", pass%n_calls
       call msg_warning ()
    end if
    allocate (pass%calls (n_it), source = 0)
    allocate (pass%integral (n_it), source = 0._default)
    allocate (pass%error (n_it), source = 0._default)
    allocate (pass%efficiency (n_it), source = 0._default)
  end subroutine pass_configure
  
\eatline
\nwendcode{}\nwbegindocs{557}Allocate the VAMP history and give options.  We assume that the
{\tt{}configure} routine above has been executed, so the number of
iterations is known.
\nwenddocs{}\nwbegincode{558}\moddef{MCI vamp: pass: TBP}\plusendmoddef
  procedure :: configure_history => pass_configure_history
\nwendcode{}\nwbegincode{559}\moddef{MCI vamp: procedures}\plusendmoddef
  subroutine pass_configure_history (pass, n_channels, par)
    class(pass_t), intent(inout) :: pass
    integer, intent(in) :: n_channels
    type(history_parameters_t), intent(in) :: par
    if (par%global) then
       allocate (pass%v_history (pass%n_it))
       call vamp_create_history (pass%v_history, &
            verbose = par%global_verbose)
    end if
    if (par%channel) then
       allocate (pass%v_histories (pass%n_it, n_channels))
       call vamp_create_history (pass%v_histories, &
            verbose = par%channel_verbose)
    end if
  end subroutine pass_configure_history
  
\eatline
\nwendcode{}\nwbegindocs{560}Given two pass objects, compare them.  All parameters must match.  Where
integrations are done in both (number of calls nonzero), the results must be
equal (up to numerical noise).

The allocated array sizes might be different, but should match up to the
common {\tt{}n{\char95}it} value.
\nwenddocs{}\nwbegincode{561}\moddef{MCI vamp: interfaces}\endmoddef
  interface operator (.matches.)
     module procedure pass_matches
  end interface operator (.matches.)
\nwendcode{}\nwbegincode{562}\moddef{MCI vamp: procedures}\plusendmoddef
  function pass_matches (pass, ref) result (ok)
    type(pass_t), intent(in) :: pass, ref
    integer :: n
    logical :: ok
    ok = .true.
    if (ok)  ok = pass%i_pass == ref%i_pass
    if (ok)  ok = pass%i_first_it == ref%i_first_it
    if (ok)  ok = pass%n_it == ref%n_it
    if (ok)  ok = pass%n_calls == ref%n_calls
    if (ok)  ok = pass%n_bins == ref%n_bins
    if (ok)  ok = pass%adapt_grids .eqv. ref%adapt_grids
    if (ok)  ok = pass%adapt_weights .eqv. ref%adapt_weights
    if (ok)  ok = pass%integral_defined .eqv. ref%integral_defined
    if (pass%integral_defined) then
       n = pass%n_it
       if (ok)  ok = all (pass%calls(:n) == ref%calls(:n))
       if (ok)  ok = all (pass%integral(:n) .matches. ref%integral(:n))
       if (ok)  ok = all (pass%error(:n) .matches. ref%error(:n))
       if (ok)  ok = all (pass%efficiency(:n) .matches. ref%efficiency(:n))
    end if
  end function pass_matches
    
\eatline
\nwendcode{}\nwbegindocs{563}Update a pass object, given a reference.  The parameters must match, except
for the {\tt{}n{\char95}it} entry.  The number of complete iterations must be less or
equal to the reference, and the number of complete iterations in the reference
must be no larger than {\tt{}n{\char95}it}.  Where results are present in both passes,
they must match.  Where results are present in the reference only, the pass is
updated accordingly.
\nwenddocs{}\nwbegincode{564}\moddef{MCI vamp: pass: TBP}\plusendmoddef
  procedure :: update => pass_update
\nwendcode{}\nwbegincode{565}\moddef{MCI vamp: procedures}\plusendmoddef
  subroutine pass_update (pass, ref, ok)
    class(pass_t), intent(inout) :: pass
    type(pass_t), intent(in) :: ref
    logical, intent(out) :: ok
    integer :: n, n_ref
    ok = .true.
    if (ok)  ok = pass%i_pass == ref%i_pass
    if (ok)  ok = pass%i_first_it == ref%i_first_it
    if (ok)  ok = pass%n_calls == ref%n_calls
    if (ok)  ok = pass%n_bins == ref%n_bins
    if (ok)  ok = pass%adapt_grids .eqv. ref%adapt_grids
    if (ok)  ok = pass%adapt_weights .eqv. ref%adapt_weights
    if (ok) then
       if (ref%integral_defined) then
          if (.not. allocated (pass%calls)) then
             allocate (pass%calls (pass%n_it), source = 0)
             allocate (pass%integral (pass%n_it), source = 0._default)
             allocate (pass%error (pass%n_it), source = 0._default)
             allocate (pass%efficiency (pass%n_it), source = 0._default)
          end if
          n = count (pass%calls /= 0)
          n_ref = count (ref%calls /= 0)
          ok = n <= n_ref .and. n_ref <= pass%n_it
          if (ok)  ok = all (pass%calls(:n) == ref%calls(:n))
          if (ok)  ok = all (pass%integral(:n) .matches. ref%integral(:n))
          if (ok)  ok = all (pass%error(:n) .matches. ref%error(:n))
          if (ok)  ok = all (pass%efficiency(:n) .matches. ref%efficiency(:n))
          if (ok) then
             pass%calls(n+1:n_ref) = ref%calls(n+1:n_ref)
             pass%integral(n+1:n_ref) = ref%integral(n+1:n_ref)
             pass%error(n+1:n_ref) = ref%error(n+1:n_ref)
             pass%efficiency(n+1:n_ref) = ref%efficiency(n+1:n_ref)
             pass%integral_defined = any (pass%calls /= 0)
          end if
       end if
    end if
  end subroutine pass_update

\eatline
\nwendcode{}\nwbegindocs{566}Match two real numbers: they are equal up to a tolerance, which is
$10^{-8}$, matching the number of digits that are output by {\tt{}pass{\char95}write}.
In particular, if one number is exactly zero, the other one must also be zero.
\nwenddocs{}\nwbegincode{567}\moddef{MCI vamp: interfaces}\plusendmoddef
  interface operator (.matches.)
     module procedure real_matches
  end interface operator (.matches.)
\nwendcode{}\nwbegincode{568}\moddef{MCI vamp: procedures}\plusendmoddef
  elemental function real_matches (x, y) result (ok)
    real(default), intent(in) :: x, y
    logical :: ok
    real(default), parameter :: tolerance = 1.e-8_default
    ok = abs (x - y) <= tolerance * max (abs (x), abs (y))
  end function real_matches
  
\eatline
\nwendcode{}\nwbegindocs{569}Return the index of the most recent complete integration.  If there is none,
return zero.
\nwenddocs{}\nwbegincode{570}\moddef{MCI vamp: pass: TBP}\plusendmoddef
  procedure :: get_integration_index => pass_get_integration_index
\nwendcode{}\nwbegincode{571}\moddef{MCI vamp: procedures}\plusendmoddef
  function pass_get_integration_index (pass) result (n)
    class (pass_t), intent(in) :: pass
    integer :: n
    integer :: i
    n = 0
    if (allocated (pass%calls)) then
       do i = 1, pass%n_it
          if (pass%calls(i) == 0)  exit
          n = i
       end do
    end if
  end function pass_get_integration_index

\eatline
\nwendcode{}\nwbegindocs{572}Return the most recent integral and error, if available.
\nwenddocs{}\nwbegincode{573}\moddef{MCI vamp: pass: TBP}\plusendmoddef
  procedure :: get_calls => pass_get_calls
  procedure :: get_integral => pass_get_integral
  procedure :: get_error => pass_get_error
  procedure :: get_efficiency => pass_get_efficiency
\nwendcode{}\nwbegincode{574}\moddef{MCI vamp: procedures}\plusendmoddef
  function pass_get_calls (pass) result (calls)
    class(pass_t), intent(in) :: pass
    integer :: calls
    integer :: n
    n = pass%get_integration_index ()
    if (n /= 0) then
       calls = pass%calls(n)
    else
       calls = 0
    end if
  end function pass_get_calls

  function pass_get_integral (pass) result (integral)
    class(pass_t), intent(in) :: pass
    real(default) :: integral
    integer :: n
    n = pass%get_integration_index ()
    if (n /= 0) then
       integral = pass%integral(n)
    else
       integral = 0
    end if
  end function pass_get_integral

  function pass_get_error (pass) result (error)
    class(pass_t), intent(in) :: pass
    real(default) :: error
    integer :: n
    n = pass%get_integration_index ()
    if (n /= 0) then
       error = pass%error(n)
    else
       error = 0
    end if
  end function pass_get_error

  function pass_get_efficiency (pass) result (efficiency)
    class(pass_t), intent(in) :: pass
    real(default) :: efficiency
    integer :: n
    n = pass%get_integration_index ()
    if (n /= 0) then
       efficiency = pass%efficiency(n)
    else
       efficiency = 0
    end if
  end function pass_get_efficiency

\eatline
\eatline
\eatline
\nwendcode{}\nwbegindocs{575}\nwdocspar
\subsection{Integrator}
\nwenddocs{}\nwbegincode{576}\moddef{MCI vamp: public}\plusendmoddef
  public :: mci_vamp_t
\nwendcode{}\nwbegincode{577}\moddef{MCI vamp: types}\plusendmoddef
  type, extends (mci_t) :: mci_vamp_t
     logical, dimension(:), allocatable :: dim_is_flat
     type(grid_parameters_t) :: grid_par
     type(history_parameters_t) :: history_par
     integer :: min_calls = 0
     type(pass_t), pointer :: first_pass => null ()
     type(pass_t), pointer :: current_pass => null ()
     type(vamp_equivalences_t) :: equivalences
     logical :: rebuild = .true.
     logical :: check_grid_file = .true.
     logical :: grid_filename_set = .false.
     logical :: negative_weights = .false.
     logical :: verbose = .false.
     type(string_t) :: grid_filename
     character(32) :: md5sum_adapted = ""
   contains
   \LA{}MCI vamp: mci vamp: TBP\RA{}
  end type mci_vamp_t
  
\eatline
\nwendcode{}\nwbegindocs{578}Reset: delete integration-pass entries.
\nwenddocs{}\nwbegincode{579}\moddef{MCI vamp: mci vamp: TBP}\endmoddef
  procedure :: reset => mci_vamp_reset
\nwendcode{}\nwbegincode{580}\moddef{MCI vamp: procedures}\plusendmoddef
  subroutine mci_vamp_reset (object)
    class(mci_vamp_t), intent(inout) :: object
    type(pass_t), pointer :: current_pass
    do while (associated (object%first_pass))
       current_pass => object%first_pass
       object%first_pass => current_pass%next
       call current_pass%final ()
       deallocate (current_pass)
    end do
    object%current_pass => null ()
  end subroutine mci_vamp_reset
  
\eatline
\nwendcode{}\nwbegindocs{581}Finalizer: reset and finalize the equivalences list.
\nwenddocs{}\nwbegincode{582}\moddef{MCI vamp: mci vamp: TBP}\plusendmoddef
  procedure :: final => mci_vamp_final
\nwendcode{}\nwbegincode{583}\moddef{MCI vamp: procedures}\plusendmoddef
  subroutine mci_vamp_final (object)
    class(mci_vamp_t), intent(inout) :: object
    call object%reset ()
    call vamp_equivalences_final (object%equivalences)
    call object%base_final ()
  end subroutine mci_vamp_final
  
\eatline
\nwendcode{}\nwbegindocs{584}Output.  Do not output the grids themselves, this may result in tons
of data.
\nwenddocs{}\nwbegincode{585}\moddef{MCI vamp: mci vamp: TBP}\plusendmoddef
  procedure :: write => mci_vamp_write
\nwendcode{}\nwbegincode{586}\moddef{MCI vamp: procedures}\plusendmoddef
  subroutine mci_vamp_write (object, unit, pacify, md5sum_version)
    class(mci_vamp_t), intent(in) :: object
    integer, intent(in), optional :: unit
    logical, intent(in), optional :: pacify
    logical, intent(in), optional :: md5sum_version
    type(pass_t), pointer :: current_pass
    integer :: u, i
    u = given_output_unit (unit)
    write (u, "(1x,A)")  "VAMP integrator:"
    call object%base_write (u, pacify, md5sum_version)
    if (allocated (object%dim_is_flat)) then
       write (u, "(3x,A,999(1x,I0))")  "Flat dimensions    =", &
            pack ([(i, i = 1, object%n_dim)], object%dim_is_flat)
    end if
    write (u, "(1x,A)")  "Grid parameters:"
    call object%grid_par%write (u)
    write (u, "(3x,A,I0)") "min_calls             = ", object%min_calls
    write (u, "(3x,A,L1)") "negative weights      = ", &
         object%negative_weights
    write (u, "(3x,A,L1)") "verbose               = ", &
         object%verbose  
    if (object%grid_par%use_vamp_equivalences) then
       call vamp_equivalences_write (object%equivalences, u)
    end if
    current_pass => object%first_pass
    do while (associated (current_pass))
       write (u, "(1x,A,I0,A)")  "Integration pass:"
       call current_pass%write (u, pacify)
       current_pass => current_pass%next
    end do
    if (object%md5sum_adapted /= "") then
       write (u, "(1x,A,A,A)")  "MD5 sum (including results) = '", &
            object%md5sum_adapted, "'"
    end if
  end subroutine mci_vamp_write
  
\eatline
\nwendcode{}\nwbegindocs{587}Write the history parameters.
\nwenddocs{}\nwbegincode{588}\moddef{MCI vamp: mci vamp: TBP}\plusendmoddef
  procedure :: write_history_parameters => mci_vamp_write_history_parameters
\nwendcode{}\nwbegincode{589}\moddef{MCI vamp: procedures}\plusendmoddef
  subroutine mci_vamp_write_history_parameters (mci, unit)
    class(mci_vamp_t), intent(in) :: mci
    integer, intent(in), optional :: unit
    integer :: u
    u = given_output_unit (unit)
    write (u, "(1x,A)")  "VAMP history parameters:"
    call mci%history_par%write (unit)
  end subroutine mci_vamp_write_history_parameters

\eatline
\nwendcode{}\nwbegindocs{590}Write the history, iterating over passes.  We keep this separate
from the generic {\tt{}write} routine.
\nwenddocs{}\nwbegincode{591}\moddef{MCI vamp: mci vamp: TBP}\plusendmoddef
  procedure :: write_history => mci_vamp_write_history
\nwendcode{}\nwbegincode{592}\moddef{MCI vamp: procedures}\plusendmoddef
  subroutine mci_vamp_write_history (mci, unit)
    class(mci_vamp_t), intent(in) :: mci
    integer, intent(in), optional :: unit
    type(pass_t), pointer :: current_pass
    integer :: i_pass
    integer :: u
    u = given_output_unit (unit)
    if (associated (mci%first_pass)) then
       write (u, "(1x,A)")  "VAMP history (global):"
       i_pass = 0
       current_pass => mci%first_pass
       do while (associated (current_pass))
          i_pass = i_pass + 1
          write (u, "(1x,A,I0,':')")  "Pass #", i_pass
          call current_pass%write_history (u)
          current_pass => current_pass%next
       end do
    end if
  end subroutine mci_vamp_write_history
  
\eatline
\nwendcode{}\nwbegindocs{593}Compute the MD5 sum, including the configuration MD5 sum and the
printout, which incorporates the current results.
\nwenddocs{}\nwbegincode{594}\moddef{MCI vamp: mci vamp: TBP}\plusendmoddef
  procedure :: compute_md5sum => mci_vamp_compute_md5sum
\nwendcode{}\nwbegincode{595}\moddef{MCI vamp: procedures}\plusendmoddef
  subroutine mci_vamp_compute_md5sum (mci, pacify)
    class(mci_vamp_t), intent(inout) :: mci
    logical, intent(in), optional :: pacify
    integer :: u
    mci%md5sum_adapted = ""
    u = free_unit ()
    open (u, status = "scratch", action = "readwrite")
    write (u, "(A)")  mci%md5sum
    call mci%write (u, pacify, md5sum_version = .true.)
    rewind (u)
    mci%md5sum_adapted = md5sum (u)
    close (u)
  end subroutine mci_vamp_compute_md5sum
    
\eatline
\nwendcode{}\nwbegindocs{596}Return the MD5 sum: If available, return the adapted one.
\nwenddocs{}\nwbegincode{597}\moddef{MCI vamp: mci vamp: TBP}\plusendmoddef
  procedure :: get_md5sum => mci_vamp_get_md5sum
\nwendcode{}\nwbegincode{598}\moddef{MCI vamp: procedures}\plusendmoddef
  function mci_vamp_get_md5sum (mci) result (md5sum)
    class(mci_vamp_t), intent(in) :: mci
    character(32) :: md5sum
    if (mci%md5sum_adapted /= "") then
       md5sum = mci%md5sum_adapted
    else
       md5sum = mci%md5sum
    end if
  end function mci_vamp_get_md5sum
  
\eatline
\nwendcode{}\nwbegindocs{599}Startup message: short version.
\nwenddocs{}\nwbegincode{600}\moddef{MCI vamp: mci vamp: TBP}\plusendmoddef
  procedure :: startup_message => mci_vamp_startup_message
\nwendcode{}\nwbegincode{601}\moddef{MCI vamp: procedures}\plusendmoddef
  subroutine mci_vamp_startup_message (mci, unit, n_calls)
    class(mci_vamp_t), intent(in) :: mci
    integer, intent(in), optional :: unit, n_calls
    integer :: num_calls, n_bins
    if (present (n_calls)) then
       num_calls = n_calls 
    else
       num_calls = 0
    end if
    if (mci%min_calls /= 0) then
       n_bins =  max (mci%grid_par%min_bins, &
            min (num_calls / mci%min_calls, &
            mci%grid_par%max_bins))
    else
       n_bins = mci%grid_par%max_bins  
    end if    
    call mci%base_startup_message (unit = unit, n_calls = n_calls)
    if (mci%grid_par%use_vamp_equivalences) then
       write (msg_buffer, "(A,2(1x,I0,1x,A))") &
            "Integrator: Using VAMP channel equivalences"    
       call msg_message (unit = unit)
    end if
    write (msg_buffer, "(A,2(1x,I0,1x,A),L1)") &
         "Integrator:", num_calls, &
         "initial calls,", n_bins, & 
         "bins, stratified = ", &
         mci%grid_par%stratified
    call msg_message (unit = unit)
    write (msg_buffer, "(A,2(1x,I0,1x,A))") &
         "Integrator: VAMP"
    call msg_message (unit = unit)
  end subroutine mci_vamp_startup_message
    
\eatline
\nwendcode{}\nwbegindocs{602}Set the MCI index (necessary for processes with multiple components).
We append the index to the grid filename, just before the final dotted
suffix.
\nwenddocs{}\nwbegincode{603}\moddef{MCI vamp: mci vamp: TBP}\plusendmoddef
  procedure :: record_index => mci_vamp_record_index
\nwendcode{}\nwbegincode{604}\moddef{MCI vamp: procedures}\plusendmoddef
  subroutine mci_vamp_record_index (mci, i_mci)
    class(mci_vamp_t), intent(inout) :: mci
    integer, intent(in) :: i_mci
    type(string_t) :: basename, suffix
    character(32) :: buffer
    if (mci%grid_filename_set) then
       basename = mci%grid_filename
       call split (basename, suffix, ".", back=.true.)
       write (buffer, "(I0)")  i_mci
       if (basename /= "") then
          mci%grid_filename = basename // "_m" // trim (buffer) // "." // suffix
       else
          mci%grid_filename = suffix // "_m" // trim (buffer) // ".vg"
       end if
    end if
  end subroutine mci_vamp_record_index

\eatline
\nwendcode{}\nwbegindocs{605}Set the grid parameters.
\nwenddocs{}\nwbegincode{606}\moddef{MCI vamp: mci vamp: TBP}\plusendmoddef
  procedure :: set_grid_parameters => mci_vamp_set_grid_parameters
\nwendcode{}\nwbegincode{607}\moddef{MCI vamp: procedures}\plusendmoddef
  subroutine mci_vamp_set_grid_parameters (mci, grid_par)
    class(mci_vamp_t), intent(inout) :: mci
    type(grid_parameters_t), intent(in) :: grid_par
    mci%grid_par = grid_par
    mci%min_calls = grid_par%min_calls_per_bin * mci%n_channel
  end subroutine mci_vamp_set_grid_parameters
  
\eatline
\nwendcode{}\nwbegindocs{608}Set the history parameters.
\nwenddocs{}\nwbegincode{609}\moddef{MCI vamp: mci vamp: TBP}\plusendmoddef
  procedure :: set_history_parameters => mci_vamp_set_history_parameters
\nwendcode{}\nwbegincode{610}\moddef{MCI vamp: procedures}\plusendmoddef
  subroutine mci_vamp_set_history_parameters (mci, history_par)
    class(mci_vamp_t), intent(inout) :: mci
    type(history_parameters_t), intent(in) :: history_par
    mci%history_par = history_par
  end subroutine mci_vamp_set_history_parameters
  
\eatline
\nwendcode{}\nwbegindocs{611}Set the rebuild flag, also the flag for checking the grid file.
\nwenddocs{}\nwbegincode{612}\moddef{MCI vamp: mci vamp: TBP}\plusendmoddef
  procedure :: set_rebuild_flag => mci_vamp_set_rebuild_flag
\nwendcode{}\nwbegincode{613}\moddef{MCI vamp: procedures}\plusendmoddef
  subroutine mci_vamp_set_rebuild_flag (mci, rebuild, check_grid_file)
    class(mci_vamp_t), intent(inout) :: mci
    logical, intent(in) :: rebuild
    logical, intent(in) :: check_grid_file
    mci%rebuild = rebuild
    mci%check_grid_file = check_grid_file
  end subroutine mci_vamp_set_rebuild_flag
  
\eatline
\nwendcode{}\nwbegindocs{614}Set the filename.
\nwenddocs{}\nwbegincode{615}\moddef{MCI vamp: mci vamp: TBP}\plusendmoddef
  procedure :: set_grid_filename => mci_vamp_set_grid_filename
\nwendcode{}\nwbegincode{616}\moddef{MCI vamp: procedures}\plusendmoddef
  subroutine mci_vamp_set_grid_filename (mci, name, run_id)
    class(mci_vamp_t), intent(inout) :: mci
    type(string_t), intent(in) :: name
    type(string_t), intent(in), optional :: run_id
    if (present (run_id)) then
       mci%grid_filename = name // "." // run_id // ".vg"
    else
       mci%grid_filename = name // ".vg"
    end if
    mci%grid_filename_set = .true.
  end subroutine mci_vamp_set_grid_filename
  
\eatline
\nwendcode{}\nwbegindocs{617}Declare particular dimensions as flat.
\nwenddocs{}\nwbegincode{618}\moddef{MCI vamp: mci vamp: TBP}\plusendmoddef
  procedure :: declare_flat_dimensions => mci_vamp_declare_flat_dimensions
\nwendcode{}\nwbegincode{619}\moddef{MCI vamp: procedures}\plusendmoddef
  subroutine mci_vamp_declare_flat_dimensions (mci, dim_flat)
    class(mci_vamp_t), intent(inout) :: mci
    integer, dimension(:), intent(in) :: dim_flat
    integer :: d
    allocate (mci%dim_is_flat (mci%n_dim), source = .false.)
    do d = 1, size (dim_flat)
       mci%dim_is_flat(dim_flat(d)) = .true.
    end do
  end subroutine mci_vamp_declare_flat_dimensions
  
\eatline
\nwendcode{}\nwbegindocs{620}Declare equivalences.  We have an array of channel equivalences,
provided by the phase-space module.  Here, we translate this into the
{\tt{}vamp{\char95}equivalences} array.
\nwenddocs{}\nwbegincode{621}\moddef{MCI vamp: mci vamp: TBP}\plusendmoddef
  procedure :: declare_equivalences => mci_vamp_declare_equivalences
\nwendcode{}\nwbegincode{622}\moddef{MCI vamp: procedures}\plusendmoddef
  subroutine mci_vamp_declare_equivalences (mci, channel, dim_offset)
    class(mci_vamp_t), intent(inout) :: mci
    type(phs_channel_t), dimension(:), intent(in) :: channel
    integer, intent(in) :: dim_offset
    integer, dimension(:), allocatable :: perm, mode
    integer :: n_channels, n_dim, n_equivalences
    integer :: c, i, j, left, right
    n_channels = mci%n_channel
    n_dim = mci%n_dim
    n_equivalences = 0
    do c = 1, n_channels
       n_equivalences = n_equivalences + size (channel(c)%eq)
    end do
    call vamp_equivalences_init (mci%equivalences, &
         n_equivalences, n_channels, n_dim)
    allocate (perm (n_dim))
    allocate (mode (n_dim))
    perm(1:dim_offset) = [(i, i = 1, dim_offset)]
    mode(1:dim_offset) = VEQ_IDENTITY
    c = 1
    j = 0
    do i = 1, n_equivalences
       if (j < size (channel(c)%eq)) then
          j = j + 1
       else
          c = c + 1
          j = 1
       end if
       associate (eq => channel(c)%eq(j))
         left = c
         right = eq%c
         perm(dim_offset+1:) = eq%perm + dim_offset
         mode(dim_offset+1:) = eq%mode
         call vamp_equivalence_set (mci%equivalences, &
              i, left, right, perm, mode)
       end associate
    end do
    call vamp_equivalences_complete (mci%equivalences)
  end subroutine mci_vamp_declare_equivalences

\eatline
\nwendcode{}\nwbegindocs{623}Allocate instance with matching type.
\nwenddocs{}\nwbegincode{624}\moddef{MCI vamp: mci vamp: TBP}\plusendmoddef
  procedure :: allocate_instance => mci_vamp_allocate_instance
\nwendcode{}\nwbegincode{625}\moddef{MCI vamp: procedures}\plusendmoddef
  subroutine mci_vamp_allocate_instance (mci, mci_instance)
    class(mci_vamp_t), intent(in) :: mci
    class(mci_instance_t), intent(out), pointer :: mci_instance
    allocate (mci_vamp_instance_t :: mci_instance)
  end subroutine mci_vamp_allocate_instance
  
\eatline
\nwendcode{}\nwbegindocs{626}Allocate a new integration pass.  We can preset everything that does not
depend on the number of iterations and calls.  This is postponed to
the {\tt{}integrate} method.

In the final pass, we do not check accuracy goal etc., since we can assume
that the user wants to perform and average all iterations in this pass.
\nwenddocs{}\nwbegincode{627}\moddef{MCI vamp: mci vamp: TBP}\plusendmoddef
  procedure :: add_pass => mci_vamp_add_pass
\nwendcode{}\nwbegincode{628}\moddef{MCI vamp: procedures}\plusendmoddef
  subroutine mci_vamp_add_pass (mci, adapt_grids, adapt_weights, final)
    class(mci_vamp_t), intent(inout) :: mci
    logical, intent(in), optional :: adapt_grids, adapt_weights, final
    integer :: i_pass, i_it
    type(pass_t), pointer :: new
    allocate (new)
    if (associated (mci%current_pass)) then
       i_pass = mci%current_pass%i_pass + 1
       i_it   = mci%current_pass%i_first_it + mci%current_pass%n_it
       mci%current_pass%next => new
    else
       i_pass = 1
       i_it = 1
       mci%first_pass => new
    end if
    mci%current_pass => new
    new%i_pass = i_pass
    new%i_first_it = i_it
    if (present (adapt_grids)) then
       new%adapt_grids = adapt_grids
    else
       new%adapt_grids = .false.
    end if
    if (present (adapt_weights)) then
       new%adapt_weights = adapt_weights
    else
       new%adapt_weights = .false.
    end if
    if (present (final)) then
       new%is_final_pass = final
    else
       new%is_final_pass = .false.
    end if
  end subroutine mci_vamp_add_pass
  
\eatline
\nwendcode{}\nwbegindocs{629}Update the list of integration passes.  All passes except for the last one
must match exactly.  For the last one, integration results are updated.  The
reference output may contain extra passes, these are ignored.
\nwenddocs{}\nwbegincode{630}\moddef{MCI vamp: mci vamp: TBP}\plusendmoddef
  procedure :: update_from_ref => mci_vamp_update_from_ref
\nwendcode{}\nwbegincode{631}\moddef{MCI vamp: procedures}\plusendmoddef
  subroutine mci_vamp_update_from_ref (mci, mci_ref, success)
    class(mci_vamp_t), intent(inout) :: mci
    class(mci_t), intent(in) :: mci_ref
    logical, intent(out) :: success
    type(pass_t), pointer :: current_pass, ref_pass
    select type (mci_ref)
    type is (mci_vamp_t)
       current_pass => mci%first_pass
       ref_pass => mci_ref%first_pass
       success = .true.
       do while (success .and. associated (current_pass))
          if (associated (ref_pass)) then
             if (associated (current_pass%next)) then
                success = current_pass .matches. ref_pass
             else
                call current_pass%update (ref_pass, success)
                if (current_pass%integral_defined) then
                   mci%integral = current_pass%get_integral ()
                   mci%error = current_pass%get_error ()
                   mci%efficiency = current_pass%get_efficiency ()
                   mci%integral_known = .true.
                   mci%error_known = .true.
                   mci%efficiency_known = .true.
                end if
             end if
             current_pass => current_pass%next
             ref_pass => ref_pass%next
          else
             success = .false.
          end if
       end do
    end select
  end subroutine mci_vamp_update_from_ref
  
\eatline
\nwendcode{}\nwbegindocs{632}Update the MCI record (i.e., the integration passes) by reading from input
stream.  The stream should contain a {\tt{}write} output from a previous run.  We
first check the MD5 sum of the configuration parameters.  If that matches, we
proceed directly to the stored integration passes.  If successful, we may
continue to read the file; the position will be after a blank line that
must follow the MCI record.
\nwenddocs{}\nwbegincode{633}\moddef{MCI vamp: mci vamp: TBP}\plusendmoddef
  procedure :: update => mci_vamp_update
\nwendcode{}\nwbegincode{634}\moddef{MCI vamp: procedures}\plusendmoddef
  subroutine mci_vamp_update (mci, u, success)
    class(mci_vamp_t), intent(inout) :: mci
    integer, intent(in) :: u
    logical, intent(out) :: success
    character(80) :: buffer
    character(32) :: md5sum_file
    type(mci_vamp_t) :: mci_file
    integer :: n_pass, n_it
    call read_sval (u, md5sum_file)
    if (mci%check_grid_file) then
       success = md5sum_file == mci%md5sum
    else
       success = .true.
    end if
    if (success) then
       read (u, *)
       read (u, "(A)")  buffer
       if (trim (adjustl (buffer)) == "VAMP integrator:") then
          n_pass = 0
          n_it = 0
          do
             read (u, "(A)")  buffer
             select case (trim (adjustl (buffer)))
             case ("")
                exit
             case ("Integration pass:")
                call mci_file%add_pass ()
                call mci_file%current_pass%read (u, n_pass, n_it)
                n_pass = n_pass + 1
                n_it = n_it + mci_file%current_pass%n_it
             end select
          end do
          call mci%update_from_ref (mci_file, success)
          call mci_file%final ()
       else
          call msg_fatal ("VAMP: reading grid file: corrupted data")
       end if
    end if
  end subroutine mci_vamp_update
  
\eatline
\nwendcode{}\nwbegindocs{635}Read / write grids from / to file.
\nwenddocs{}\nwbegincode{636}\moddef{MCI vamp: mci vamp: TBP}\plusendmoddef
  procedure :: write_grids => mci_vamp_write_grids
  procedure :: read_grids_header => mci_vamp_read_grids_header
  procedure :: read_grids_data => mci_vamp_read_grids_data
  procedure :: read_grids => mci_vamp_read_grids
\nwendcode{}\nwbegincode{637}\moddef{MCI vamp: procedures}\plusendmoddef
  subroutine mci_vamp_write_grids (mci, instance)
    class(mci_vamp_t), intent(in) :: mci
    class(mci_instance_t), intent(inout) :: instance
    integer :: u
    select type (instance)
    type is (mci_vamp_instance_t)
       if (mci%grid_filename_set) then
          if (instance%grids_defined) then
             u = free_unit ()
             open (u, file = char (mci%grid_filename), &
                  action = "write", status = "replace")
             write (u, "(1x,A,A,A)")  "MD5sum = '", mci%md5sum, "'"
             write (u, *)
             call mci%write (u)
             write (u, *)
             write (u, "(1x,A)")  "VAMP grids:"
             call vamp_write_grids (instance%grids, u, &
                  write_integrals = .true.)
             close (u)
          else
             call msg_bug ("VAMP: write grids: grids undefined")
          end if
       else
          call msg_bug ("VAMP: write grids: filename undefined")
       end if
    end select
  end subroutine mci_vamp_write_grids

  subroutine mci_vamp_read_grids_header (mci, success)
    class(mci_vamp_t), intent(inout) :: mci
    logical, intent(out) :: success
    logical :: exist
    integer :: u
    success = .false.
    if (mci%grid_filename_set) then
       inquire (file = char (mci%grid_filename), exist = exist)
       if (exist) then
          u = free_unit ()
          open (u, file = char (mci%grid_filename), &
               action = "read", status = "old")
          call mci%update (u, success)
          close (u)
          if (.not. success) then
             write (msg_buffer, "(A,A,A)") &
                  "VAMP: parameter mismatch, discarding grid file '", &
                  char (mci%grid_filename), "'"
             call msg_message ()
          end if
       end if
    else
       call msg_bug ("VAMP: read grids: filename undefined")
    end if
  end subroutine mci_vamp_read_grids_header
  
  subroutine mci_vamp_read_grids_data (mci, instance, read_integrals)
    class(mci_vamp_t), intent(in) :: mci
    class(mci_instance_t), intent(inout) :: instance
    logical, intent(in), optional :: read_integrals
    integer :: u
    character(80) :: buffer
    select type (instance)
    type is (mci_vamp_instance_t)
       if (.not. instance%grids_defined) then
          u = free_unit ()
          open (u, file = char (mci%grid_filename), &
               action = "read", status = "old") 
          do
             read (u, "(A)")  buffer
             if (trim (adjustl (buffer)) == "VAMP grids:")  exit
          end do
          call vamp_read_grids (instance%grids, u, read_integrals)
          close (u)
          instance%grids_defined = .true.
       else
          call msg_bug ("VAMP: read grids: grids already defined")
       end if
    end select
  end subroutine mci_vamp_read_grids_data
  
  subroutine mci_vamp_read_grids (mci, instance, success)
    class(mci_vamp_t), intent(inout) :: mci
    class(mci_instance_t), intent(inout) :: instance
    logical, intent(out) :: success
    logical :: exist
    integer :: u
    character(80) :: buffer
    select type (instance)
    type is (mci_vamp_instance_t)
       success = .false.
       if (mci%grid_filename_set) then
          if (.not. instance%grids_defined) then
             inquire (file = char (mci%grid_filename), exist = exist)
             if (exist) then
                u = free_unit ()
                open (u, file = char (mci%grid_filename), &
                     action = "read", status = "old")
                call mci%update (u, success)
                if (success) then
                   read (u, "(A)")  buffer
                   if (trim (adjustl (buffer)) == "VAMP grids:") then
                      call vamp_read_grids (instance%grids, u)
                   else
                      call msg_fatal ("VAMP: reading grid file: &
                           &corrupted grid data")
                   end if
                else
                   write (msg_buffer, "(A,A,A)") &
                        "VAMP: parameter mismatch, discarding grid file '", &
                        char (mci%grid_filename), "'"
                   call msg_message ()
                end if
                close (u)
                instance%grids_defined = success
             end if
          else
             call msg_bug ("VAMP: read grids: grids already defined")
          end if
       else
          call msg_bug ("VAMP: read grids: filename undefined")
       end if
    end select
  end subroutine mci_vamp_read_grids
  
\eatline
\eatline
\eatline
\eatline
\nwendcode{}\nwbegindocs{638}Auxiliary: Read real, integer, string value.  We search for an equals sign,
the value must follow.
\nwenddocs{}\nwbegincode{639}\moddef{MCI vamp: procedures}\plusendmoddef
  subroutine read_rval (u, rval)
    integer, intent(in) :: u
    real(default), intent(out) :: rval
    character(80) :: buffer
    read (u, "(A)")  buffer
    buffer = adjustl (buffer(scan (buffer, "=") + 1:))
    read (buffer, *)  rval
  end subroutine read_rval
    
  subroutine read_ival (u, ival)
    integer, intent(in) :: u
    integer, intent(out) :: ival
    character(80) :: buffer
    read (u, "(A)")  buffer
    buffer = adjustl (buffer(scan (buffer, "=") + 1:))
    read (buffer, *)  ival
  end subroutine read_ival
    
  subroutine read_sval (u, sval)
    integer, intent(in) :: u
    character(*), intent(out) :: sval
    character(80) :: buffer
    read (u, "(A)")  buffer
    buffer = adjustl (buffer(scan (buffer, "=") + 1:))
    read (buffer, *)  sval
  end subroutine read_sval
    
  subroutine read_lval (u, lval)
    integer, intent(in) :: u
    logical, intent(out) :: lval
    character(80) :: buffer
    read (u, "(A)")  buffer
    buffer = adjustl (buffer(scan (buffer, "=") + 1:))
    read (buffer, *)  lval
  end subroutine read_lval
    
\eatline
\nwendcode{}\nwbegindocs{640}Integrate.  Perform a new integration pass (possibly reusing
previous results), which may consist of several iterations.

Note: we record the integral once per iteration.  The integral stored
in the {\tt{}mci} record itself is the last integral of the current
iteration, no averaging done.  The {\tt{}results} record may average results.

Note: recording the efficiency is not supported yet.
\nwenddocs{}\nwbegincode{641}\moddef{MCI vamp: mci vamp: TBP}\plusendmoddef
  procedure :: integrate => mci_vamp_integrate
\nwendcode{}\nwbegincode{642}\moddef{MCI vamp: procedures}\plusendmoddef
  subroutine mci_vamp_integrate (mci, instance, sampler, &
       n_it, n_calls, results, pacify)
    class(mci_vamp_t), intent(inout) :: mci
    class(mci_instance_t), intent(inout) :: instance
    class(mci_sampler_t), intent(inout) :: sampler
    integer, intent(in) :: n_it
    integer, intent(in) :: n_calls
    class(mci_results_t), intent(inout), optional :: results
    logical, intent(in), optional :: pacify
    integer :: it
    logical :: reshape, from_file, success
    select type (instance)
    type is (mci_vamp_instance_t)
       if (associated (mci%current_pass)) then
          mci%current_pass%integral_defined = .false.
          call mci%current_pass%configure (n_it, n_calls, &
               mci%min_calls, mci%grid_par%min_bins, &
               mci%grid_par%max_bins, &
               mci%grid_par%min_calls_per_channel * mci%n_channel)
          call mci%current_pass%configure_history &
               (mci%n_channel, mci%history_par)
          instance%pass_complete = .false.
          instance%it_complete = .false.
          call instance%new_pass (reshape)
          if (.not. instance%grids_defined .or. instance%grids_from_file) then
             if (mci%grid_filename_set .and. .not. mci%rebuild) then
                call mci%read_grids_header (success)
                from_file = success
                if (.not. instance%grids_defined .and. success) then
                   call mci%read_grids_data (instance)
                end if
             else
                from_file = .false.
             end if
          else
             from_file = .false.
          end if
          if (from_file) then
             if (.not. mci%check_grid_file) &
                  call msg_warning ("Reading grid file: MD5 sum check disabled")
             call msg_message ("VAMP: " &
                  // "using grids and results from file '" &
                  // char (mci%grid_filename) // "'")
          else if (.not. instance%grids_defined) then
             call instance%create_grids ()
          end if
          do it = 1, instance%n_it
             if (signal_is_pending ())  return
             instance%grids_from_file = from_file .and. &
                  it <= mci%current_pass%get_integration_index ()
             if (.not. instance%grids_from_file) then
                instance%it_complete = .false.
                call instance%adapt_grids ()
                if (signal_is_pending ())  return
                call instance%adapt_weights ()
                if (signal_is_pending ())  return
                call instance%discard_integrals (reshape)
                if (mci%grid_par%use_vamp_equivalences) then
                   call instance%sample_grids (mci%rng, sampler, &
                        mci%equivalences)
                else
                   call instance%sample_grids (mci%rng, sampler)
                end if
                if (signal_is_pending ())  return
                instance%it_complete = .true.
                if (instance%integral /= 0) then
                   mci%current_pass%calls(it) = instance%calls
                   mci%current_pass%integral(it) = instance%integral
                   if (abs (instance%error / instance%integral) &
                        > epsilon (1._default)) then
                      mci%current_pass%error(it) = instance%error
                   end if
                   mci%current_pass%efficiency(it) = instance%efficiency
                end if
                mci%current_pass%integral_defined = .true.
             end if
             if (present (results)) then
                if (mci%has_chains ()) then
                   call mci%collect_chain_weights (instance%w)
                   call results%record (1, &
                        n_calls    = mci%current_pass%calls(it), &
                        integral   = mci%current_pass%integral(it), &
                        error      = mci%current_pass%error(it), &
                        efficiency = mci%current_pass%efficiency(it), &
                        chain_weights = mci%chain_weights, &
                        suppress = pacify)
                else
                   call results%record (1, &
                        n_calls    = mci%current_pass%calls(it), &
                        integral   = mci%current_pass%integral(it), &
                        error      = mci%current_pass%error(it), &
                        efficiency = mci%current_pass%efficiency(it), &
                        suppress = pacify)
                end if
             end if
             if (.not. instance%grids_from_file &
                  .and. mci%grid_filename_set) then
                call mci%write_grids (instance)
             end if
             call instance%allow_adaptation ()
             reshape = .false.
             if (.not. mci%current_pass%is_final_pass) then
                call mci%check_goals (it, success)
                if (success)  exit
             end if
          end do
          if (signal_is_pending ())  return
          instance%pass_complete = .true.
          mci%integral = mci%current_pass%get_integral()
          mci%error = mci%current_pass%get_error()
          mci%efficiency = mci%current_pass%get_efficiency()
          mci%integral_known = .true.
          mci%error_known = .true.
          mci%efficiency_known = .true.
          call mci%compute_md5sum (pacify)
       else
          call msg_bug ("MCI integrate: current_pass object not allocated")
       end if
    end select
  end subroutine mci_vamp_integrate

\eatline
\nwendcode{}\nwbegindocs{643}Check whether we are already finished with this pass.
\nwenddocs{}\nwbegincode{644}\moddef{MCI vamp: mci vamp: TBP}\plusendmoddef
  procedure :: check_goals => mci_vamp_check_goals
\nwendcode{}\nwbegincode{645}\moddef{MCI vamp: procedures}\plusendmoddef
  subroutine mci_vamp_check_goals (mci, it, success)
    class(mci_vamp_t), intent(inout) :: mci
    integer, intent(in) :: it
    logical, intent(out) :: success
    success = .false.
    if (mci%error_reached (it)) then
       mci%current_pass%n_it = it
       call msg_message ("VAMP: error goal reached; &
            &skipping iterations")
       success = .true.
       return
    end if
    if (mci%rel_error_reached (it)) then
       mci%current_pass%n_it = it
       call msg_message ("VAMP: relative error goal reached; &
            &skipping iterations")
       success = .true.
       return
    end if
    if (mci%accuracy_reached (it)) then
       mci%current_pass%n_it = it
       call msg_message ("VAMP: accuracy goal reached; &
            &skipping iterations")
       success = .true.
       return
    end if
  end subroutine mci_vamp_check_goals
    
\eatline
\nwendcode{}\nwbegindocs{646}Return true if the error, relative error, or accuracy goal has been reached,
if any.
\nwenddocs{}\nwbegincode{647}\moddef{MCI vamp: mci vamp: TBP}\plusendmoddef
  procedure :: error_reached => mci_vamp_error_reached
  procedure :: rel_error_reached => mci_vamp_rel_error_reached
  procedure :: accuracy_reached => mci_vamp_accuracy_reached
\nwendcode{}\nwbegincode{648}\moddef{MCI vamp: procedures}\plusendmoddef
  function mci_vamp_error_reached (mci, it) result (flag)
    class(mci_vamp_t), intent(in) :: mci
    integer, intent(in) :: it
    logical :: flag
    real(default) :: error_goal, error
    error_goal = mci%grid_par%error_goal
    if (error_goal > 0) then
       associate (pass => mci%current_pass)
         if (pass%integral_defined) then
            error = abs (pass%error(it))
            flag = error < error_goal
         else
            flag = .false.
         end if
       end associate
    else
       flag = .false.
    end if
  end function mci_vamp_error_reached
  
  function mci_vamp_rel_error_reached (mci, it) result (flag)
    class(mci_vamp_t), intent(in) :: mci
    integer, intent(in) :: it
    logical :: flag
    real(default) :: rel_error_goal, rel_error
    rel_error_goal = mci%grid_par%rel_error_goal
    if (rel_error_goal > 0) then
       associate (pass => mci%current_pass)
         if (pass%integral_defined) then
            if (pass%integral(it) /= 0) then
               rel_error = abs (pass%error(it) / pass%integral(it))
               flag = rel_error < rel_error_goal
            else
               flag = .true.
            end if
         else
            flag = .false.
         end if
       end associate
    else
       flag = .false.
    end if
  end function mci_vamp_rel_error_reached
  
  function mci_vamp_accuracy_reached (mci, it) result (flag)
    class(mci_vamp_t), intent(in) :: mci
    integer, intent(in) :: it
    logical :: flag
    real(default) :: accuracy_goal, accuracy
    accuracy_goal = mci%grid_par%accuracy_goal
    if (accuracy_goal > 0) then
       associate (pass => mci%current_pass)
         if (pass%integral_defined) then
            if (pass%integral(it) /= 0) then
               accuracy = abs (pass%error(it) / pass%integral(it)) &
                    * sqrt (real (pass%calls(it), default))
               flag = accuracy < accuracy_goal
            else
               flag = .true.
            end if
         else
            flag = .false.
         end if
       end associate
    else
       flag = .false.
    end if
  end function mci_vamp_accuracy_reached
  
\eatline
\eatline
\eatline
\nwendcode{}\nwbegindocs{649}Prepare an event generation pass.  Should be called before a sequence of
events is generated, then we should call the corresponding finalizer.

The pass-specific data of the previous integration pass are retained,
but we reset the number of iterations and calls to zero.  The latter
now counts the number of events (calls to the sampling function, actually).
\nwenddocs{}\nwbegincode{650}\moddef{MCI vamp: mci vamp: TBP}\plusendmoddef
  procedure :: prepare_simulation => mci_vamp_prepare_simulation
\nwendcode{}\nwbegincode{651}\moddef{MCI vamp: procedures}\plusendmoddef
  subroutine mci_vamp_prepare_simulation (mci)
    class(mci_vamp_t), intent(inout) :: mci
    logical :: success
    if (mci%grid_filename_set) then
       call mci%read_grids_header (success)
       call mci%compute_md5sum ()
       if (.not. success) then
          call msg_fatal ("Simulate: " &
               // "reading integration grids from file '" &
               // char (mci%grid_filename) // "' failed")
       end if
    else
       call msg_bug ("VAMP: simulation: no grids, no grid filename")
    end if
  end subroutine mci_vamp_prepare_simulation
  
\eatline
\nwendcode{}\nwbegindocs{652}Generate weighted event.  Note that the event weight
({\tt{}vamp{\char95}weight}) is not just the MCI weight.  {\tt{}vamp{\char95}next{\char95}event}
selects a channel based on the channel weights multiplied by the
(previously recorded) maximum integrand value of the channel.  The
MCI weight is renormalized accordingly, to cancel this effect on the
result.
\nwenddocs{}\nwbegincode{653}\moddef{MCI vamp: mci vamp: TBP}\plusendmoddef
  procedure :: generate_weighted_event => mci_vamp_generate_weighted_event
\nwendcode{}\nwbegincode{654}\moddef{MCI vamp: procedures}\plusendmoddef
  subroutine mci_vamp_generate_weighted_event (mci, instance, sampler)
    class(mci_vamp_t), intent(inout) :: mci
    class(mci_instance_t), intent(inout), target :: instance
    class(mci_sampler_t), intent(inout), target :: sampler
    class(vamp_data_t), allocatable :: data
    type(exception) :: vamp_exception
    select type (instance)
    type is (mci_vamp_instance_t)
       instance%vamp_weight_set = .false.
       allocate (mci_workspace_t :: data)
       select type (data)
       type is (mci_workspace_t)
          data%sampler => sampler
          data%instance => instance
       end select
       select type (rng => mci%rng)
       type is (rng_tao_t)
          if (instance%grids_defined) then
             call vamp_next_event ( &
                  instance%vamp_x, &
                  rng%state, &
                  instance%grids, &
                  vamp_sampling_function, &
                  data, &
                  phi = phi_trivial, &
                  weight = instance%vamp_weight, &
                  exc = vamp_exception)
             call handle_vamp_exception (vamp_exception, mci%verbose)
             instance%vamp_excess = 0
             instance%vamp_weight_set = .true.
          else
             call msg_bug ("VAMP: generate event: grids undefined")
          end if
       class default
         call msg_fatal ("VAMP event generation: &
               &random-number generator must be TAO")
       end select
    end select
  end subroutine mci_vamp_generate_weighted_event
       
\eatline
\nwendcode{}\nwbegindocs{655}Generate unweighted event.
\nwenddocs{}\nwbegincode{656}\moddef{MCI vamp: mci vamp: TBP}\plusendmoddef
  procedure :: generate_unweighted_event => &
       mci_vamp_generate_unweighted_event
\nwendcode{}\nwbegincode{657}\moddef{MCI vamp: procedures}\plusendmoddef
  subroutine mci_vamp_generate_unweighted_event (mci, instance, sampler)
    class(mci_vamp_t), intent(inout) :: mci
    class(mci_instance_t), intent(inout), target :: instance
    class(mci_sampler_t), intent(inout), target :: sampler
    class(vamp_data_t), allocatable :: data
    type(exception) :: vamp_exception
    select type (instance)
    type is (mci_vamp_instance_t)
       instance%vamp_weight_set = .false.
       allocate (mci_workspace_t :: data)
       select type (data)
       type is (mci_workspace_t)
          data%sampler => sampler
          data%instance => instance
       end select
       select type (rng => mci%rng)
       type is (rng_tao_t)
          if (instance%grids_defined) then
             REJECTION: do
                call vamp_next_event ( &
                     instance%vamp_x, &
                     rng%state, &
                     instance%grids, &
                     vamp_sampling_function, &
                     data, &
                     phi = phi_trivial, &
                     excess = instance%vamp_excess, &
                     exc = vamp_exception)
                if (signal_is_pending ())  return
                if (sampler%is_valid ())  exit REJECTION
             end do REJECTION
             call handle_vamp_exception (vamp_exception, mci%verbose)
             instance%vamp_weight = 1
             instance%vamp_weight_set = .true.
          else
             call msg_bug ("VAMP: generate event: grids undefined")
          end if
       class default
         call msg_fatal ("VAMP event generation: &
               &random-number generator must be TAO")
       end select
    end select
  end subroutine mci_vamp_generate_unweighted_event
    
\eatline
\nwendcode{}\nwbegindocs{658}Rebuild an event, using the {\tt{}state} input.

Note: This feature is currently unused.
\nwenddocs{}\nwbegincode{659}\moddef{MCI vamp: mci vamp: TBP}\plusendmoddef
  procedure :: rebuild_event => mci_vamp_rebuild_event
\nwendcode{}\nwbegincode{660}\moddef{MCI vamp: procedures}\plusendmoddef
  subroutine mci_vamp_rebuild_event (mci, instance, sampler, state)
    class(mci_vamp_t), intent(inout) :: mci
    class(mci_instance_t), intent(inout) :: instance
    class(mci_sampler_t), intent(inout) :: sampler
    class(mci_state_t), intent(in) :: state
    call msg_bug ("MCI vamp rebuild event not implemented yet")
  end subroutine mci_vamp_rebuild_event
       
\eatline
\nwendcode{}\nwbegindocs{661}\nwdocspar
\nwenddocs{}\nwbegincode{662}\moddef{MCI vamp: mci vamp: TBP}\plusendmoddef
  procedure :: pacify => mci_vamp_pacify
\nwendcode{}\nwbegincode{663}\moddef{MCI vamp: procedures}\plusendmoddef
  subroutine mci_vamp_pacify (object, efficiency_reset, error_reset)
    class(mci_vamp_t), intent(inout) :: object
    logical, intent(in), optional :: efficiency_reset, error_reset
    logical :: err_reset
    type(pass_t), pointer :: current_pass
    err_reset = .false.
    if (present (error_reset))  err_reset = error_reset
    current_pass => object%first_pass    
    do while (associated (current_pass))
       if (allocated (current_pass%error) .and. err_reset) then
          current_pass%error = 0
       end if
       if (allocated (current_pass%efficiency) .and. err_reset) then
          current_pass%efficiency = 1
       end if
       current_pass => current_pass%next
    end do
  end subroutine mci_vamp_pacify
    
\nwendcode{}\nwbegindocs{664} %def mci_vamp_pacify
\nwenddocs{}\nwbegindocs{665}\nwdocspar
\subsection{Sampler as Workspace}
In the full setup, the sampling function requires the process instance
object as workspace.  We implement this by (i) implementing the
process instance as a type extension of the abstract {\tt{}sampler{\char95}t}
object used by the MCI implementation and (ii) providing such an
object as an extra argument to the sampling function that VAMP can
call.  To minimize cross-package dependencies, we use an abstract type
{\tt{}vamp{\char95}workspace} that VAMP declares and extend this by including a
pointer to the {\tt{}sampler} and {\tt{}instance} objects.  In the body of
the sampling function, we dereference this pointer and can then work
with the contents.
\nwenddocs{}\nwbegincode{666}\moddef{MCI vamp: types}\plusendmoddef
  type, extends (vamp_data_t) :: mci_workspace_t
     class(mci_sampler_t), pointer :: sampler => null ()
     class(mci_vamp_instance_t), pointer :: instance => null ()
  end type mci_workspace_t
  
\eatline
\nwendcode{}\nwbegindocs{667}\nwdocspar
\subsection{Integrator instance}
The history entries should point to the corresponding history entry in
the {\tt{}pass{\char95}t} object.  If there is none, we may allocate a local
history, which is then just transient.
\nwenddocs{}\nwbegincode{668}\moddef{MCI vamp: public}\plusendmoddef
  public :: mci_vamp_instance_t
\nwendcode{}\nwbegincode{669}\moddef{MCI vamp: types}\plusendmoddef
  type, extends (mci_instance_t) :: mci_vamp_instance_t
     type(mci_vamp_t), pointer :: mci => null ()
     logical :: grids_defined = .false.
     logical :: grids_from_file = .false.
     integer :: n_it = 0
     integer :: it = 0
     logical :: pass_complete = .false.
     integer :: n_calls = 0
     integer :: calls = 0
     logical :: it_complete = .false.
     logical :: enable_adapt_grids = .false.
     logical :: enable_adapt_weights = .false.
     logical :: allow_adapt_grids = .false.
     logical :: allow_adapt_weights = .false.
     logical :: negative_weights = .false.
     integer :: n_adapt_grids = 0
     integer :: n_adapt_weights = 0
     logical :: generating_events = .false.
     real(default) :: safety_factor = 1
     type(vamp_grids) :: grids
     real(default) :: g = 0
     real(default), dimension(:), allocatable :: gi
     real(default) :: integral = 0
     real(default) :: error = 0
     real(default) :: efficiency = 0
     real(default), dimension(:), allocatable :: vamp_x
     logical :: vamp_weight_set = .false.
     real(default) :: vamp_weight = 0
     real(default) :: vamp_excess = 0
     logical :: allocate_global_history = .false.
     type(vamp_history), dimension(:), pointer :: v_history => null ()
     logical :: allocate_channel_history = .false.
     type(vamp_history), dimension(:,:), pointer :: v_histories => null ()
   contains
   \LA{}MCI vamp: mci vamp instance: TBP\RA{}
  end type mci_vamp_instance_t
  
\eatline
\nwendcode{}\nwbegindocs{670}Output.
\nwenddocs{}\nwbegincode{671}\moddef{MCI vamp: mci vamp instance: TBP}\endmoddef
  procedure :: write => mci_vamp_instance_write
\nwendcode{}\nwbegincode{672}\moddef{MCI vamp: procedures}\plusendmoddef
  subroutine mci_vamp_instance_write (object, unit, pacify)
    class(mci_vamp_instance_t), intent(in) :: object
    integer, intent(in), optional :: unit
    logical, intent(in), optional :: pacify    
    integer :: u, i
    character(len=7) :: fmt 
    call pac_fmt (fmt, FMT_17, FMT_14, pacify)
    u = given_output_unit (unit)
    write (u, "(3x,A," // FMT_19 // ")") "Integrand = ", object%integrand
    write (u, "(3x,A," // FMT_19 // ")") "Weight    = ", object%mci_weight
    if (object%vamp_weight_set) then
       write (u, "(3x,A," // FMT_19 // ")") "VAMP wgt  = ", object%vamp_weight
       if (object%vamp_excess /= 0) then
          write (u, "(3x,A," // FMT_19 // ")") "VAMP exc  = ", &
               object%vamp_excess
       end if
    end if
    write (u, "(3x,A,L1)")  "adapt grids   = ", object%enable_adapt_grids
    write (u, "(3x,A,L1)")  "adapt weights = ", object%enable_adapt_weights
    if (object%grids_defined) then
       if (object%grids_from_file) then
          write (u, "(3x,A)")  "VAMP grids: read from file"
       else
          write (u, "(3x,A)")  "VAMP grids: defined"
       end if
    else
       write (u, "(3x,A)")  "VAMP grids: [undefined]"
    end if
    write (u, "(3x,A,I0)")  "n_it          = ", object%n_it
    write (u, "(3x,A,I0)")  "it            = ", object%it
    write (u, "(3x,A,L1)")  "pass complete = ", object%it_complete
    write (u, "(3x,A,I0)")  "n_calls       = ", object%n_calls
    write (u, "(3x,A,I0)")  "calls         = ", object%calls
    write (u, "(3x,A,L1)")  "it complete   = ", object%it_complete
    write (u, "(3x,A,I0)")  "n adapt.(g)   = ", object%n_adapt_grids
    write (u, "(3x,A,I0)")  "n adapt.(w)   = ", object%n_adapt_weights
    write (u, "(3x,A,L1)")  "gen. events   = ", object%generating_events
    write (u, "(3x,A,L1)")  "neg. weights  = ", object%negative_weights
    if (object%safety_factor /= 1)  write &
          (u, "(3x,A," // fmt // ")")  "safety f = ", object%safety_factor
    write (u, "(3x,A," // fmt // ")")  "integral = ", object%integral
    write (u, "(3x,A," // fmt // ")")  "error    = ", object%error
    write (u, "(3x,A," // fmt // ")")  "eff.     = ", object%efficiency
    write (u, "(3x,A)")  "weights:"
    do i = 1, size (object%w)
       write (u, "(5x,I0,1x," // FMT_12 // ")")  i, object%w(i)
    end do
  end subroutine mci_vamp_instance_write
  
\eatline
\nwendcode{}\nwbegindocs{673}Write the grids to the specified unit.
\nwenddocs{}\nwbegincode{674}\moddef{MCI vamp: mci vamp instance: TBP}\plusendmoddef
  procedure :: write_grids => mci_vamp_instance_write_grids
\nwendcode{}\nwbegincode{675}\moddef{MCI vamp: procedures}\plusendmoddef
  subroutine mci_vamp_instance_write_grids (object, unit)
    class(mci_vamp_instance_t), intent(in) :: object
    integer, intent(in), optional :: unit
    integer :: u
    u = given_output_unit (unit)
    if (object%grids_defined) then
       call vamp_write_grids (object%grids, u, write_integrals = .true.)
    end if
  end subroutine mci_vamp_instance_write_grids
  
\eatline
\nwendcode{}\nwbegindocs{676}Finalizer: the history arrays are pointer arrays and need finalization.
\nwenddocs{}\nwbegincode{677}\moddef{MCI vamp: mci vamp instance: TBP}\plusendmoddef
  procedure :: final => mci_vamp_instance_final
\nwendcode{}\nwbegincode{678}\moddef{MCI vamp: tests}\endmoddef
  subroutine mci_vamp_instance_final (object)
    class(mci_vamp_instance_t), intent(inout) :: object
    if (object%allocate_global_history) then
       if (associated (object%v_history)) then
          call vamp_delete_history (object%v_history)
          deallocate (object%v_history)
       end if
    end if
    if (object%allocate_channel_history) then
       if (associated (object%v_histories)) then
          call vamp_delete_history (object%v_histories)
          deallocate (object%v_histories)
       end if
    end if
    if (object%grids_defined) then
       call vamp_delete_grids (object%grids)
       object%grids_defined = .false.       
    end if
  end subroutine mci_vamp_instance_final
  
\eatline
\nwendcode{}\nwbegindocs{679}Initializer.
\nwenddocs{}\nwbegincode{680}\moddef{MCI vamp: mci vamp instance: TBP}\plusendmoddef
  procedure :: init => mci_vamp_instance_init
\nwendcode{}\nwbegincode{681}\moddef{MCI vamp: procedures}\plusendmoddef
  subroutine mci_vamp_instance_init (mci_instance, mci)
    class(mci_vamp_instance_t), intent(out) :: mci_instance
    class(mci_t), intent(in), target :: mci
    call mci_instance%base_init (mci)
    select type (mci)
    type is (mci_vamp_t)
       mci_instance%mci => mci
       allocate (mci_instance%gi (mci%n_channel))
       mci_instance%allocate_global_history = .not. mci%history_par%global
       mci_instance%allocate_channel_history = .not. mci%history_par%channel
       mci_instance%negative_weights = mci%negative_weights
    end select
  end subroutine mci_vamp_instance_init
    
\eatline
\nwendcode{}\nwbegindocs{682}Prepare a new integration pass: write the pass-specific settings to
the {\tt{}instance} object.  This should be called initially, together
with the {\tt{}create{\char95}grids} procedure, and whenever we start a new
integration pass.

Set {\tt{}reshape} if the number of calls is different than previously (unless it
was zero, indicating the first pass).

We link VAMP histories to the allocated histories in the current pass
object, so the recorded results are persistent.  However, if there are
no histories present there, we allocate them locally.  In that case,
the histories will disappear together with the MCI instance object.
\nwenddocs{}\nwbegincode{683}\moddef{MCI vamp: mci vamp instance: TBP}\plusendmoddef
  procedure :: new_pass => mci_vamp_instance_new_pass
\nwendcode{}\nwbegincode{684}\moddef{MCI vamp: procedures}\plusendmoddef
  subroutine mci_vamp_instance_new_pass (instance, reshape)
    class(mci_vamp_instance_t), intent(inout) :: instance
    logical, intent(out) :: reshape
    type(pass_t), pointer :: current
    associate (mci => instance%mci)
      current => mci%current_pass
      instance%n_it = current%n_it
      if (instance%n_calls == 0) then
         reshape = .false.
         instance%n_calls = current%n_calls
      else if (instance%n_calls == current%n_calls) then
         reshape = .false.
      else
         reshape = .true.
         instance%n_calls = current%n_calls
      end if
      instance%it = 0
      instance%calls = 0
      instance%enable_adapt_grids = current%adapt_grids
      instance%enable_adapt_weights = current%adapt_weights
      instance%generating_events = .false.
      if (instance%allocate_global_history) then
         if (associated (instance%v_history)) then
            call vamp_delete_history (instance%v_history)
            deallocate (instance%v_history)
         end if
         allocate (instance%v_history (instance%n_it))
         call vamp_create_history (instance%v_history, verbose = .false.)
      else
         instance%v_history => current%v_history
      end if
      if (instance%allocate_channel_history) then
         if (associated (instance%v_histories)) then
            call vamp_delete_history (instance%v_histories)
            deallocate (instance%v_histories)
         end if
         allocate (instance%v_histories (instance%n_it, mci%n_channel))
         call vamp_create_history (instance%v_histories, verbose = .false.)
      else
         instance%v_histories => current%v_histories
      end if
    end associate
  end subroutine mci_vamp_instance_new_pass
  
\eatline
\nwendcode{}\nwbegindocs{685}\nwdocspar
Create a grid set within the {\tt{}instance} object, using the data of
the current integration pass.  Also reset counters that track this
grid set.
\nwenddocs{}\nwbegincode{686}\moddef{MCI vamp: mci vamp instance: TBP}\plusendmoddef
  procedure :: create_grids => mci_vamp_instance_create_grids
\nwendcode{}\nwbegincode{687}\moddef{MCI vamp: procedures}\plusendmoddef
  subroutine mci_vamp_instance_create_grids (instance)
    class(mci_vamp_instance_t), intent(inout) :: instance
    type (pass_t), pointer :: current
    integer, dimension(:), allocatable :: num_div
    real(default), dimension(:,:), allocatable :: region
    associate (mci => instance%mci)
      current => mci%current_pass
      allocate (num_div (mci%n_dim))
      allocate (region (2, mci%n_dim))
      region(1,:) = 0
      region(2,:) = 1
      num_div = current%n_bins
      instance%n_adapt_grids = 0
      instance%n_adapt_weights = 0
      if (.not. instance%grids_defined) then
         call vamp_create_grids (instance%grids, &
              region, &
              current%n_calls, &
              weights = instance%w, &
              num_div = num_div, &
              stratified = mci%grid_par%stratified)
         instance%grids_defined = .true.
      else
         call msg_bug ("VAMP: create grids: grids already defined")
      end if
    end associate
  end subroutine mci_vamp_instance_create_grids

\eatline
\nwendcode{}\nwbegindocs{688}Reset a grid set, so we can start a fresh integration pass.  In
effect, we delete results of previous integrations, but keep the grid
shapes, weights, and variance arrays, so adaptation is still possible.
The grids are prepared for a specific number of calls (per iteration)
and sampling mode (stratified/importance).

The {\tt{}vamp{\char95}discard{\char95}integrals} implementation will reshape the grids
only if the argument {\tt{}num{\char95}calls} is present.
\nwenddocs{}\nwbegincode{689}\moddef{MCI vamp: mci vamp instance: TBP}\plusendmoddef
  procedure :: discard_integrals => mci_vamp_instance_discard_integrals
\nwendcode{}\nwbegincode{690}\moddef{MCI vamp: procedures}\plusendmoddef
  subroutine mci_vamp_instance_discard_integrals (instance, reshape)
    class(mci_vamp_instance_t), intent(inout) :: instance
    logical, intent(in) :: reshape
    instance%calls = 0
    instance%integral = 0
    instance%error = 0
    instance%efficiency = 0
    associate (mci => instance%mci)
      if (instance%grids_defined) then
         if (mci%grid_par%use_vamp_equivalences) then
            if (reshape) then
               call vamp_discard_integrals (instance%grids, &
                    num_calls = instance%n_calls, &
                    stratified = mci%grid_par%stratified, &
                    eq = mci%equivalences)
            else
               call vamp_discard_integrals (instance%grids, &
                    stratified = mci%grid_par%stratified, &
                    eq = mci%equivalences)
            end if
         else
            if (reshape) then
               call vamp_discard_integrals (instance%grids, &
                    num_calls = instance%n_calls, &
                    stratified = mci%grid_par%stratified)
            else
               call vamp_discard_integrals (instance%grids, &
                    stratified = mci%grid_par%stratified)
            end if
         end if
      else
         call msg_bug ("VAMP: discard integrals: grids undefined")
      end if
    end associate
  end subroutine mci_vamp_instance_discard_integrals
    
\eatline
\nwendcode{}\nwbegindocs{691}After grids are created (with equidistant binning and equal weight),
adaptation is redundant.  Therefore, we should allow it only after a
complete integration step has been performed, calling this.
\nwenddocs{}\nwbegincode{692}\moddef{MCI vamp: mci vamp instance: TBP}\plusendmoddef
  procedure :: allow_adaptation => mci_vamp_instance_allow_adaptation
\nwendcode{}\nwbegincode{693}\moddef{MCI vamp: procedures}\plusendmoddef
  subroutine mci_vamp_instance_allow_adaptation (instance)
    class(mci_vamp_instance_t), intent(inout) :: instance
    instance%allow_adapt_grids = .true.
    instance%allow_adapt_weights = .true.
  end subroutine mci_vamp_instance_allow_adaptation

\eatline
\nwendcode{}\nwbegindocs{694}Adapt grids.  
\nwenddocs{}\nwbegincode{695}\moddef{MCI vamp: mci vamp instance: TBP}\plusendmoddef
  procedure :: adapt_grids => mci_vamp_instance_adapt_grids
\nwendcode{}\nwbegincode{696}\moddef{MCI vamp: procedures}\plusendmoddef
  subroutine mci_vamp_instance_adapt_grids (instance)
    class(mci_vamp_instance_t), intent(inout) :: instance
    if (instance%enable_adapt_grids .and. instance%allow_adapt_grids) then
       if (instance%grids_defined) then
          call vamp_refine_grids (instance%grids)
          instance%n_adapt_grids = instance%n_adapt_grids + 1
      else
         call msg_bug ("VAMP: adapt grids: grids undefined")
      end if
    end if
  end subroutine mci_vamp_instance_adapt_grids
  
\eatline
\nwendcode{}\nwbegindocs{697}Adapt weights.  Use the variance array returned by \vamp\ for
recalculating the weight array.  The parameter
{\tt{}channel{\char95}weights{\char95}power} dampens fluctuations.

If the number of calls in a given channel falls below a user-defined threshold,
the weight is not lowered further but kept at this threshold.  The other
channel weights are reduced accordingly.
\nwenddocs{}\nwbegincode{698}\moddef{MCI vamp: mci vamp instance: TBP}\plusendmoddef
  procedure :: adapt_weights => mci_vamp_instance_adapt_weights
\nwendcode{}\nwbegincode{699}\moddef{MCI vamp: procedures}\plusendmoddef
  subroutine mci_vamp_instance_adapt_weights (instance)
    class(mci_vamp_instance_t), intent(inout) :: instance
    real(default) :: w_sum, w_avg_ch, sum_w_underflow, w_min
    integer :: n_ch, ch, n_underflow
    logical, dimension(:), allocatable :: mask, underflow
    type(exception) :: vamp_exception
    if (instance%enable_adapt_weights .and. instance%allow_adapt_weights) then
       associate (mci => instance%mci)
         if (instance%grids_defined) then
            instance%w = instance%grids%weights &
                 * vamp_get_variance (instance%grids%grids) &
                 ** mci%grid_par%channel_weights_power
            w_sum = sum (instance%w)
            if (w_sum /= 0) then
               instance%w = instance%w / w_sum
               if (mci%n_chain /= 0) then
                  allocate (mask (mci%n_channel))
                  do ch = 1, mci%n_chain
                     mask = mci%chain == ch
                     n_ch = count (mask)
                     if (n_ch /= 0) then
                        w_avg_ch = sum (instance%w, mask) / n_ch
                        where (mask)  instance%w = w_avg_ch
                     end if
                  end do
               end if
               if (mci%grid_par%threshold_calls /= 0) then
                  w_min = &
                       real (mci%grid_par%threshold_calls, default) &
                       / instance%n_calls
                  allocate (underflow (mci%n_channel))
                  underflow = instance%w /= 0 .and. abs (instance%w) < w_min
                  n_underflow = count (underflow)
                  sum_w_underflow = sum (instance%w, mask=underflow)
                  if (sum_w_underflow /= 1) then
                     where (underflow)
                        instance%w = w_min
                     elsewhere
                        instance%w = instance%w &
                             * (1 - n_underflow * w_min) / (1 - sum_w_underflow)
                     end where
                  end if
               end if
            end if
            call vamp_update_weights (instance%grids, instance%w, &
                 exc = vamp_exception)
            call handle_vamp_exception (vamp_exception, mci%verbose)
         else
            call msg_bug ("VAMP: adapt weights: grids undefined")
         end if
       end associate
       instance%n_adapt_weights = instance%n_adapt_weights + 1
    end if
  end subroutine mci_vamp_instance_adapt_weights
  
\eatline
\nwendcode{}\nwbegindocs{700}Integration: sample the VAMP grids.  The number of calls etc. are
already stored inside the grids.  We provide the random-number
generator, the sampling function, and a link to the workspace object,
which happens to contain a pointer to the sampler object.  The sampler
object thus becomes the workspace of the sampling function.

Note: in the current implementation, the random-number generator must
be the TAO generator.  This explicit dependence should be removed from
the VAMP implementation.
\nwenddocs{}\nwbegincode{701}\moddef{MCI vamp: mci vamp instance: TBP}\plusendmoddef
  procedure :: sample_grids => mci_vamp_instance_sample_grids
\nwendcode{}\nwbegincode{702}\moddef{MCI vamp: procedures}\plusendmoddef
  subroutine mci_vamp_instance_sample_grids (instance, rng, sampler, eq)
    class(mci_vamp_instance_t), intent(inout), target :: instance
    class(rng_t), intent(inout) :: rng
    class(mci_sampler_t), intent(inout), target :: sampler
    type(vamp_equivalences_t), intent(in), optional :: eq
    class(vamp_data_t), allocatable :: data
    type(exception) :: vamp_exception
    allocate (mci_workspace_t :: data)
    select type (data)
    type is (mci_workspace_t)
       data%sampler => sampler
       data%instance => instance
    end select
    select type (rng)
    type is (rng_tao_t)
       instance%it = instance%it + 1
       instance%calls = 0
       if (instance%grids_defined) then
          call vamp_sample_grids ( &
               rng%state, &
               instance%grids, &
               vamp_sampling_function, &
               data, &
               1, &
               eq = eq, &
               history = instance%v_history(instance%it:), &
               histories = instance%v_histories(instance%it:,:), &
               integral = instance%integral, &
               std_dev = instance%error, &
               exc = vamp_exception, &
               negative_weights = instance%negative_weights)
          call handle_vamp_exception (vamp_exception, instance%mci%verbose)
          instance%efficiency = instance%get_efficiency ()
       else
          call msg_bug ("VAMP: sample grids: grids undefined")
       end if
    class default
       call msg_fatal ("VAMP integration: random-number generator must be TAO")
    end select
  end subroutine mci_vamp_instance_sample_grids

\eatline
\nwendcode{}\nwbegindocs{703}\nwdocspar
Compute the reweighting efficiency for the current grids, suitable
averaged over all active channels.
\nwenddocs{}\nwbegincode{704}\moddef{MCI vamp: mci vamp instance: TBP}\plusendmoddef
  procedure :: get_efficiency_array => mci_vamp_instance_get_efficiency_array
  procedure :: get_efficiency => mci_vamp_instance_get_efficiency
\nwendcode{}\nwbegincode{705}\moddef{MCI vamp: procedures}\plusendmoddef
  function mci_vamp_instance_get_efficiency_array (mci) result (efficiency)
    class(mci_vamp_instance_t), intent(in) :: mci
    real(default), dimension(:), allocatable :: efficiency
    allocate (efficiency (mci%mci%n_channel))
    where (mci%grids%grids%f_max /= 0)
       efficiency = mci%grids%grids%mu(1) / abs (mci%grids%grids%f_max)
    elsewhere
       efficiency = 0
    end where
  end function mci_vamp_instance_get_efficiency_array

  function mci_vamp_instance_get_efficiency (mci) result (efficiency)
    class(mci_vamp_instance_t), intent(in) :: mci
    real(default) :: efficiency
    real(default), dimension(:), allocatable :: weight
    real(default) :: norm
    allocate (weight (mci%mci%n_channel))
    weight = mci%grids%weights * abs (mci%grids%grids%f_max)
    norm = sum (weight)
    if (norm /= 0) then
       efficiency = dot_product (mci%get_efficiency_array (), weight) / norm
    else
       efficiency = 1
    end if
  end function mci_vamp_instance_get_efficiency

\eatline
\eatline
\nwendcode{}\nwbegindocs{706}Prepare an event generation pass.  Should be called before a sequence of
events is generated, then we should call the corresponding finalizer.

The pass-specific data of the previous integration pass are retained,
but we reset the number of iterations and calls to zero.  The latter
now counts the number of events (calls to the sampling function, actually).
\nwenddocs{}\nwbegincode{707}\moddef{MCI vamp: mci vamp instance: TBP}\plusendmoddef
  procedure :: init_simulation => mci_vamp_instance_init_simulation
\nwendcode{}\nwbegincode{708}\moddef{MCI vamp: procedures}\plusendmoddef
  subroutine mci_vamp_instance_init_simulation (instance, safety_factor)
    class(mci_vamp_instance_t), intent(inout) :: instance
    real(default), intent(in), optional :: safety_factor
    associate (mci => instance%mci)
      allocate (instance%vamp_x (mci%n_dim))
      instance%it = 0
      instance%calls = 0
      instance%generating_events = .true.
      if (present (safety_factor))  instance%safety_factor = safety_factor
      if (.not. instance%grids_defined) then
         if (mci%grid_filename_set) then
            if (.not. mci%check_grid_file) &
                 call msg_warning ("Reading grid file: MD5 sum check disabled")
            call msg_message ("Simulate: " &
                 // "using integration grids from file '" &
                 // char (mci%grid_filename) // "'")
            call mci%read_grids_data (instance)
            if (instance%safety_factor /= 1) then
               write (msg_buffer, "(A,ES10.3,A)")  "Simulate: &
                    &applying safety factor", instance%safety_factor, &
                    " to event rejection"
               call msg_message ()
               instance%grids%grids%f_max = &
                    instance%grids%grids%f_max * instance%safety_factor
            end if
         else
            call msg_bug ("VAMP: simulation: no grids, no grid filename")
         end if
      end if
    end associate
  end subroutine mci_vamp_instance_init_simulation
  
\eatline
\nwendcode{}\nwbegindocs{709}Finalize an event generation pass.  Should be called before a sequence of
events is generated, then we should call the corresponding finalizer.
\nwenddocs{}\nwbegincode{710}\moddef{MCI vamp: mci vamp instance: TBP}\plusendmoddef
  procedure :: final_simulation => mci_vamp_instance_final_simulation
\nwendcode{}\nwbegincode{711}\moddef{MCI vamp: procedures}\plusendmoddef
  subroutine mci_vamp_instance_final_simulation (instance)
    class(mci_vamp_instance_t), intent(inout) :: instance
    if (allocated (instance%vamp_x))  deallocate (instance%vamp_x)
  end subroutine mci_vamp_instance_final_simulation
  
\eatline
\nwendcode{}\nwbegindocs{712}\nwdocspar
\subsection{Sampling function}
The VAMP sampling function has a well-defined interface which we have
to implement.  The {\tt{}data} argument allows us to pass pointers to the
{\tt{}sampler} and {\tt{}instance} objects, so we can access configuration
data and fill point-dependent contents within these objects.

The {\tt{}weights} and {\tt{}channel} argument must be present in the call.

Note: we would normally declare the {\tt{}instance} pointer with the
concrete type, or just use the {\tt{}data} component directly.
Unfortunately, gfortran 4.6 forgets the inherited base-type methods in
that case.

Note: this is the place where we must look for external signals, i.e.,
interrupt from the OS.  We would like to raise a \vamp\ exception which is then
caught by {\tt{}vamp{\char95}sample{\char95}grids} as the caller, so it dumps its current state
and returns (with the signal still pending).  \whizard\ will then terminate
gracefully.  Of course, VAMP should be able to resume from the dump.  

In the current implementation, we handle the exception in place and terminate
immediately.  The incomplete current integration pass is lost.
\nwenddocs{}\nwbegincode{713}\moddef{MCI vamp: procedures}\plusendmoddef
  function vamp_sampling_function &
       (xi, data, weights, channel, grids) result (f)
    real(default) :: f
    real(default), dimension(:), intent(in) :: xi
    class(vamp_data_t), intent(in) :: data
    real(default), dimension(:), intent(in), optional :: weights
    integer, intent(in), optional :: channel
    type(vamp_grid), dimension(:), intent(in), optional :: grids
    type(exception) :: exc
    class(mci_instance_t), pointer :: instance
    logical :: verbose
    character(*), parameter :: FN = "WHIZARD sampling function"
    select type (data)
    type is (mci_workspace_t)
       instance => data%instance
       select type (instance)
       class is (mci_vamp_instance_t)
          instance%calls = instance%calls + 1
          verbose = instance%mci%verbose
       end select
       call instance%evaluate (data%sampler, channel, xi)
       if (signal_is_pending ()) then
          call raise_exception (exc, EXC_FATAL, FN, "signal received")
          call handle_vamp_exception (exc, verbose)
          call terminate_now_if_signal ()
       end if
       f = instance%get_value ()
    end select
  end function vamp_sampling_function

\eatline
\nwendcode{}\nwbegindocs{714}This is supposed to be the mapping between integration channels.
The VAMP event generating procedures technically require it, but it is
meaningless in our setup where all transformations happen inside the
sampler object.  So, this implementation is trivial:
\nwenddocs{}\nwbegincode{715}\moddef{MCI vamp: procedures}\plusendmoddef
  pure function phi_trivial (xi, channel_dummy) result (x)
    real(default), dimension(:), intent(in) :: xi
    integer, intent(in) :: channel_dummy
    real(default), dimension(size(xi)) :: x
    x = xi
  end function phi_trivial

\eatline
\nwendcode{}\nwbegindocs{716}\nwdocspar
\subsection{Integrator instance: evaluation}
Here, we compute the multi-channel reweighting factor for the current
channel, that accounts for the Jacobians of the transformations
from/to all other channels.

The computation of the VAMP probabilities may consume considerable
time, therefore we enable parallel evaluation. (Collecting the
contributions to {\tt{}mci{\char37}g} is a reduction, which we should also
implement via OpenMP.)
\nwenddocs{}\nwbegincode{717}\moddef{MCI vamp: mci vamp instance: TBP}\plusendmoddef
  procedure :: compute_weight => mci_vamp_instance_compute_weight
\nwendcode{}\nwbegincode{718}\moddef{MCI vamp: procedures}\plusendmoddef
  subroutine mci_vamp_instance_compute_weight (mci, c)
    class(mci_vamp_instance_t), intent(inout) :: mci
    integer, intent(in) :: c
    integer :: i
    mci%selected_channel = c
    !$OMP PARALLEL PRIVATE(i) SHARED(mci)
    !$OMP DO
    do i = 1, mci%mci%n_channel
       if (mci%w(i) /= 0) then
          mci%gi(i) = vamp_probability (mci%grids%grids(i), mci%x(:,i))
       else
          mci%gi(i) = 0
       end if
    end do
    !$OMP END DO
    !$OMP END PARALLEL
    mci%g = 0
    if (mci%gi(c) /= 0) then
       do i = 1, mci%mci%n_channel
          if (mci%w(i) /= 0 .and. mci%f(i) /= 0) then
             mci%g = mci%g + mci%w(i) * mci%gi(i) / mci%f(i)
          end if
       end do
    end if
    if (mci%g /= 0) then
       mci%mci_weight = mci%gi(c) / mci%g
    else
       mci%mci_weight = 0
    end if
  end subroutine mci_vamp_instance_compute_weight
    
\eatline
\nwendcode{}\nwbegindocs{719}Record the integrand.
\nwenddocs{}\nwbegincode{720}\moddef{MCI vamp: mci vamp instance: TBP}\plusendmoddef
  procedure :: record_integrand => mci_vamp_instance_record_integrand
\nwendcode{}\nwbegincode{721}\moddef{MCI vamp: procedures}\plusendmoddef
  subroutine mci_vamp_instance_record_integrand (mci, integrand)
    class(mci_vamp_instance_t), intent(inout) :: mci
    real(default), intent(in) :: integrand
    mci%integrand = integrand
  end subroutine mci_vamp_instance_record_integrand
  
\eatline
\nwendcode{}\nwbegindocs{722}Get the event weight.  The default routine returns the same value that
we would use for integration.  This is correct if we select the integration
channel according to the channel weight.  {\tt{}vamp{\char95}next{\char95}event} does
differently, so we should rather rely on the weight that VAMP
returns.  This is the value stored in {\tt{}vamp{\char95}weight}.  We override
the default TBP accordingly.
\nwenddocs{}\nwbegincode{723}\moddef{MCI vamp: mci vamp instance: TBP}\plusendmoddef
  procedure :: get_event_weight => mci_vamp_instance_get_event_weight
  procedure :: get_event_excess => mci_vamp_instance_get_event_excess
\nwendcode{}\nwbegincode{724}\moddef{MCI vamp: procedures}\plusendmoddef
  function mci_vamp_instance_get_event_weight (mci) result (value)
    class(mci_vamp_instance_t), intent(in) :: mci
    real(default) :: value
    if (mci%vamp_weight_set) then
       value = mci%vamp_weight
    else
       call msg_bug ("VAMP: attempt to read undefined event weight")
    end if
  end function mci_vamp_instance_get_event_weight
   
  function mci_vamp_instance_get_event_excess (mci) result (value)
    class(mci_vamp_instance_t), intent(in) :: mci
    real(default) :: value
    if (mci%vamp_weight_set) then
       value = mci%vamp_excess
    else
       call msg_bug ("VAMP: attempt to read undefined event excess weight")
    end if
  end function mci_vamp_instance_get_event_excess
   
\eatline
\nwendcode{}\nwbegindocs{725}\nwdocspar
\subsection{VAMP exceptions}
A VAMP routine may have raised an exception.  Turn this into a WHIZARD
error message.

An external signal could raise a fatal exception, but this should be delayed and
handled by the correct termination routine.
\nwenddocs{}\nwbegincode{726}\moddef{MCI vamp: procedures}\plusendmoddef
  subroutine handle_vamp_exception (exc, verbose)
    type(exception), intent(in) :: exc
    logical, intent(in) :: verbose
    integer :: exc_level
    if (verbose) then
       exc_level = EXC_INFO
    else
       exc_level = EXC_ERROR
    end if
    if (exc%level >= exc_level) then
       write (msg_buffer, "(A,':',1x,A)")  trim (exc%origin), trim (exc%message)
       select case (exc%level)
       case (EXC_INFO);  call msg_message ()
       case (EXC_WARN);  call msg_warning ()
       case (EXC_ERROR); call msg_error ()
       case (EXC_FATAL)
          if (signal_is_pending ()) then
             call msg_message ()
          else
             call msg_fatal ()
          end if
       end select
    end if
  end subroutine handle_vamp_exception
  
\eatline
\nwendcode{}\nwbegindocs{727}\nwdocspar
\subsection{Unit tests}
\nwenddocs{}\nwbegincode{728}\moddef{MCI vamp: public}\plusendmoddef
  public :: mci_vamp_test
\nwendcode{}\nwbegincode{729}\moddef{MCI vamp: tests}\plusendmoddef
  subroutine mci_vamp_test (u, results)
    integer, intent(in) :: u
    type(test_results_t), intent(inout) :: results
  \LA{}MCI vamp: execute tests\RA{}
  end subroutine mci_vamp_test
  
\eatline
\nwendcode{}\nwbegindocs{730}\nwdocspar
\subsubsection{Test sampler}
A test sampler object should implement a function with known integral that
we can use to check the integrator.

In mode {\tt{}1}, the function is $f(x) = 3 x^2$ with integral $\int_0^1
f(x)\,dx=1$ and maximum $f(1)=3$.  If the integration dimension is
greater than one, the function is extended as a constant in the other
dimension(s).

In mode {\tt{}2}, the function is $11 x^{10}$, also with integral $1$.
\nwenddocs{}\nwbegincode{731}\moddef{MCI vamp: test types}\endmoddef
  type, extends (mci_sampler_t) :: test_sampler_1_t
     real(default), dimension(:), allocatable :: x
     real(default) :: val
     integer :: mode = 1
   contains
   \LA{}MCI vamp: test sampler 1: TBP\RA{}
  end type test_sampler_1_t

\eatline
\nwendcode{}\nwbegindocs{732}Output: There is nothing stored inside, so just print an informative line.
\nwenddocs{}\nwbegincode{733}\moddef{MCI vamp: test sampler 1: TBP}\endmoddef
  procedure :: write => test_sampler_1_write
\nwendcode{}\nwbegincode{734}\moddef{MCI vamp: tests}\plusendmoddef
  subroutine test_sampler_1_write (object, unit, testflag)
    class(test_sampler_1_t), intent(in) :: object
    integer, intent(in), optional :: unit
    logical, intent(in), optional :: testflag
    integer :: u
    u = given_output_unit (unit)
    select case (object%mode)
    case (1)
       write (u, "(1x,A)") "Test sampler: f(x) = 3 x^2"
    case (2)
       write (u, "(1x,A)") "Test sampler: f(x) = 11 x^10"
    case (3)
       write (u, "(1x,A)") "Test sampler: f(x) = 11 x^10 * 2 * cos^2 (2 pi y)"
    end select
  end subroutine test_sampler_1_write
  
\eatline
\nwendcode{}\nwbegindocs{735}Evaluation: compute the function value.  The output $x$ parameter
(only one channel) is identical to the input $x$, and the Jacobian is 1.
\nwenddocs{}\nwbegincode{736}\moddef{MCI vamp: test sampler 1: TBP}\plusendmoddef
  procedure :: evaluate => test_sampler_1_evaluate
\nwendcode{}\nwbegincode{737}\moddef{MCI vamp: tests}\plusendmoddef
  subroutine test_sampler_1_evaluate (sampler, c, x_in, val, x, f)
    class(test_sampler_1_t), intent(inout) :: sampler
    integer, intent(in) :: c
    real(default), dimension(:), intent(in) :: x_in
    real(default), intent(out) :: val
    real(default), dimension(:,:), intent(out) :: x
    real(default), dimension(:), intent(out) :: f
    if (allocated (sampler%x))  deallocate (sampler%x)
    allocate (sampler%x (size (x_in)))
    sampler%x = x_in
    select case (sampler%mode)
    case (1)
       sampler%val = 3 * x_in(1) ** 2
    case (2)
       sampler%val = 11 * x_in(1) ** 10
    case (3)
       sampler%val = 11 * x_in(1) ** 10 * 2 * cos (twopi * x_in(2)) ** 2
    end select
    call sampler%fetch (val, x, f)
  end subroutine test_sampler_1_evaluate

\eatline
\nwendcode{}\nwbegindocs{738}The point is always valid.
\nwenddocs{}\nwbegincode{739}\moddef{MCI vamp: test sampler 1: TBP}\plusendmoddef
  procedure :: is_valid => test_sampler_1_is_valid
\nwendcode{}\nwbegincode{740}\moddef{MCI vamp: tests}\plusendmoddef
  function test_sampler_1_is_valid (sampler) result (valid)
    class(test_sampler_1_t), intent(in) :: sampler
    logical :: valid
    valid = .true.
  end function test_sampler_1_is_valid
  
\eatline
\nwendcode{}\nwbegindocs{741}Rebuild: compute all but the function value.
\nwenddocs{}\nwbegincode{742}\moddef{MCI vamp: test sampler 1: TBP}\plusendmoddef
  procedure :: rebuild => test_sampler_1_rebuild
\nwendcode{}\nwbegincode{743}\moddef{MCI vamp: tests}\plusendmoddef
  subroutine test_sampler_1_rebuild (sampler, c, x_in, val, x, f)
    class(test_sampler_1_t), intent(inout) :: sampler
    integer, intent(in) :: c
    real(default), dimension(:), intent(in) :: x_in
    real(default), intent(in) :: val
    real(default), dimension(:,:), intent(out) :: x
    real(default), dimension(:), intent(out) :: f
    if (allocated (sampler%x))  deallocate (sampler%x)
    allocate (sampler%x (size (x_in)))
    sampler%x = x_in
    sampler%val = val
    x(:,1) = sampler%x
    f = 1
  end subroutine test_sampler_1_rebuild

\eatline
\nwendcode{}\nwbegindocs{744}Extract the results.
\nwenddocs{}\nwbegincode{745}\moddef{MCI vamp: test sampler 1: TBP}\plusendmoddef
  procedure :: fetch => test_sampler_1_fetch
\nwendcode{}\nwbegincode{746}\moddef{MCI vamp: tests}\plusendmoddef
  subroutine test_sampler_1_fetch (sampler, val, x, f)
    class(test_sampler_1_t), intent(in) :: sampler
    real(default), intent(out) :: val
    real(default), dimension(:,:), intent(out) :: x
    real(default), dimension(:), intent(out) :: f
    val = sampler%val
    x(:,1) = sampler%x
    f = 1
  end subroutine test_sampler_1_fetch
    
\eatline
\nwendcode{}\nwbegindocs{747}\nwdocspar
\subsubsection{Two-channel, two dimension test sampler}
This sampler implements the function
\begin{equation}
  f(x, y) = 4\sin^2(\pi x)\sin^2(\pi y) + 2\sin^2(\pi v)
\end{equation}
where
\begin{align}
  x &= u^v        &u &= xy
\\
  y &= u^{(1-v)}  &v &= \frac12\left(1 + \frac{\log(x/y)}{\log xy}\right)
\end{align}
Each term contributes $1$ to the integral.  The first term in the function is
peaked along a cross aligned to the coordinates $x$ and $y$, while the second
term is peaked along the diagonal $x=y$.

The Jacobian is
\begin{equation}
  \frac{\partial(x,y)}{\partial(u,v)} = |\log u|
\end{equation}
\nwenddocs{}\nwbegincode{748}\moddef{MCI vamp: test types}\plusendmoddef
  type, extends (mci_sampler_t) :: test_sampler_2_t
     real(default), dimension(:,:), allocatable :: x
     real(default), dimension(:), allocatable :: f
     real(default) :: val
   contains
   \LA{}MCI vamp: test sampler 2: TBP\RA{}
  end type test_sampler_2_t

\eatline
\nwendcode{}\nwbegindocs{749}Output: There is nothing stored inside, so just print an informative line.
\nwenddocs{}\nwbegincode{750}\moddef{MCI vamp: test sampler 2: TBP}\endmoddef
  procedure :: write => test_sampler_2_write
\nwendcode{}\nwbegincode{751}\moddef{MCI vamp: tests}\plusendmoddef
  subroutine test_sampler_2_write (object, unit, testflag)
    class(test_sampler_2_t), intent(in) :: object
    integer, intent(in), optional :: unit
    logical, intent(in), optional :: testflag
    integer :: u
    u = given_output_unit (unit)
    write (u, "(1x,A)") "Two-channel test sampler 2"
  end subroutine test_sampler_2_write
  
\eatline
\nwendcode{}\nwbegindocs{752}Kinematics: compute $x$ and Jacobians, given the input parameter array.
\nwenddocs{}\nwbegincode{753}\moddef{MCI vamp: test sampler 2: TBP}\plusendmoddef
  procedure :: compute => test_sampler_2_compute
\nwendcode{}\nwbegincode{754}\moddef{MCI vamp: tests}\plusendmoddef
  subroutine test_sampler_2_compute (sampler, c, x_in)
    class(test_sampler_2_t), intent(inout) :: sampler
    integer, intent(in) :: c
    real(default), dimension(:), intent(in) :: x_in
    real(default) :: xx, yy, uu, vv
    if (.not. allocated (sampler%x)) &
         allocate (sampler%x (size (x_in), 2))
    if (.not. allocated (sampler%f)) &
         allocate (sampler%f (2))
    select case (c)
    case (1)
       xx = x_in(1)
       yy = x_in(2)
       uu = xx * yy
       vv = (1 + log (xx/yy) / log (xx*yy)) / 2
    case (2)
       uu = x_in(1)
       vv = x_in(2)
       xx = uu ** vv
       yy = uu ** (1 - vv)
    end select
    sampler%val = (2 * sin (pi * xx) * sin (pi * yy)) ** 2 &
         + 2 * sin (pi * vv) ** 2
    sampler%f(1) = 1
    sampler%f(2) = abs (log (uu))
    sampler%x(:,1) = [xx, yy]
    sampler%x(:,2) = [uu, vv]
  end subroutine test_sampler_2_compute

\eatline
\nwendcode{}\nwbegindocs{755}Evaluation: compute the function value.  The output $x$ parameter
(only one channel) is identical to the input $x$, and the Jacobian is 1.
\nwenddocs{}\nwbegincode{756}\moddef{MCI vamp: test sampler 2: TBP}\plusendmoddef
  procedure :: evaluate => test_sampler_2_evaluate
\nwendcode{}\nwbegincode{757}\moddef{MCI vamp: tests}\plusendmoddef
  subroutine test_sampler_2_evaluate (sampler, c, x_in, val, x, f)
    class(test_sampler_2_t), intent(inout) :: sampler
    integer, intent(in) :: c
    real(default), dimension(:), intent(in) :: x_in
    real(default), intent(out) :: val
    real(default), dimension(:,:), intent(out) :: x
    real(default), dimension(:), intent(out) :: f
    call sampler%compute (c, x_in)
    call sampler%fetch (val, x, f)
  end subroutine test_sampler_2_evaluate

\eatline
\nwendcode{}\nwbegindocs{758}The point is always valid.
\nwenddocs{}\nwbegincode{759}\moddef{MCI vamp: test sampler 2: TBP}\plusendmoddef
  procedure :: is_valid => test_sampler_2_is_valid
\nwendcode{}\nwbegincode{760}\moddef{MCI vamp: tests}\plusendmoddef
  function test_sampler_2_is_valid (sampler) result (valid)
    class(test_sampler_2_t), intent(in) :: sampler
    logical :: valid
    valid = .true.
  end function test_sampler_2_is_valid
  
\eatline
\nwendcode{}\nwbegindocs{761}Rebuild: compute all but the function value.
\nwenddocs{}\nwbegincode{762}\moddef{MCI vamp: test sampler 2: TBP}\plusendmoddef
  procedure :: rebuild => test_sampler_2_rebuild
\nwendcode{}\nwbegincode{763}\moddef{MCI vamp: tests}\plusendmoddef
  subroutine test_sampler_2_rebuild (sampler, c, x_in, val, x, f)
    class(test_sampler_2_t), intent(inout) :: sampler
    integer, intent(in) :: c
    real(default), dimension(:), intent(in) :: x_in
    real(default), intent(in) :: val
    real(default), dimension(:,:), intent(out) :: x
    real(default), dimension(:), intent(out) :: f
    call sampler%compute (c, x_in)
    x = sampler%x
    f = sampler%f
  end subroutine test_sampler_2_rebuild

\eatline
\nwendcode{}\nwbegindocs{764}Extract the results.
\nwenddocs{}\nwbegincode{765}\moddef{MCI vamp: test sampler 2: TBP}\plusendmoddef
  procedure :: fetch => test_sampler_2_fetch
\nwendcode{}\nwbegincode{766}\moddef{MCI vamp: tests}\plusendmoddef
  subroutine test_sampler_2_fetch (sampler, val, x, f)
    class(test_sampler_2_t), intent(in) :: sampler
    real(default), intent(out) :: val
    real(default), dimension(:,:), intent(out) :: x
    real(default), dimension(:), intent(out) :: f
    val = sampler%val
    x = sampler%x
    f = sampler%f
  end subroutine test_sampler_2_fetch
    
\eatline
\nwendcode{}\nwbegindocs{767}\nwdocspar
\subsubsection{Two-channel, one dimension test sampler}
This sampler implements the function
\begin{equation}
  f(x, y) = a * 5 x^4 + b * 5 (1-x)^4
\end{equation}
Each term contributes $1$ to the integral, multiplied by $a$ or $b$,
respectively.  The first term is peaked at $x=1$, the second one at $x=0$..

We implement the two mappings
\begin{equation}
  x = u^{1/5} \quad\text{and}\quad x = 1 - v^{1/5},
\end{equation}
with Jacobians
\begin{equation}
  \frac{\partial(x)}{\partial(u)} = u^{-4/5}/5 \quad\text{and}\quad v^{-4/5}/5,
\end{equation}
respectively.  The first mapping concentrates points near $x=1$, the
second one near $x=0$.
\nwenddocs{}\nwbegincode{768}\moddef{MCI vamp: test types}\plusendmoddef
  type, extends (mci_sampler_t) :: test_sampler_3_t
     real(default), dimension(:,:), allocatable :: x
     real(default), dimension(:), allocatable :: f
     real(default) :: val
     real(default) :: a = 1
     real(default) :: b = 1
   contains
   \LA{}MCI vamp: test sampler 3: TBP\RA{}
  end type test_sampler_3_t

\eatline
\nwendcode{}\nwbegindocs{769}Output: display $a$ and $b$
\nwenddocs{}\nwbegincode{770}\moddef{MCI vamp: test sampler 3: TBP}\endmoddef
  procedure :: write => test_sampler_3_write
\nwendcode{}\nwbegincode{771}\moddef{MCI vamp: tests}\plusendmoddef
  subroutine test_sampler_3_write (object, unit, testflag)
    class(test_sampler_3_t), intent(in) :: object
    integer, intent(in), optional :: unit
    logical, intent(in), optional :: testflag
    integer :: u
    u = given_output_unit (unit)
    write (u, "(1x,A)") "Two-channel test sampler 3"
    write (u, "(3x,A,F5.2)")  "a = ", object%a
    write (u, "(3x,A,F5.2)")  "b = ", object%b
  end subroutine test_sampler_3_write
  
\eatline
\nwendcode{}\nwbegindocs{772}Kinematics: compute $x$ and Jacobians, given the input parameter array.
\nwenddocs{}\nwbegincode{773}\moddef{MCI vamp: test sampler 3: TBP}\plusendmoddef
  procedure :: compute => test_sampler_3_compute
\nwendcode{}\nwbegincode{774}\moddef{MCI vamp: tests}\plusendmoddef
  subroutine test_sampler_3_compute (sampler, c, x_in)
    class(test_sampler_3_t), intent(inout) :: sampler
    integer, intent(in) :: c
    real(default), dimension(:), intent(in) :: x_in
    real(default) :: u, v, xx
    if (.not. allocated (sampler%x)) &
         allocate (sampler%x (size (x_in), 2))
    if (.not. allocated (sampler%f)) &
         allocate (sampler%f (2))
    select case (c)
    case (1)
       u = x_in(1)
       xx = u ** 0.2_default
       v = (1 - xx) ** 5._default
    case (2)
       v = x_in(1)
       xx = 1 - v ** 0.2_default
       u = xx ** 5._default
    end select
    sampler%val = sampler%a * 5 * xx ** 4 + sampler%b * 5 * (1 - xx) ** 4
    sampler%f(1) = 0.2_default * u ** (-0.8_default)
    sampler%f(2) = 0.2_default * v ** (-0.8_default)
    sampler%x(:,1) = [u]
    sampler%x(:,2) = [v]
  end subroutine test_sampler_3_compute

\eatline
\nwendcode{}\nwbegindocs{775}Evaluation: compute the function value.  The output $x$ parameter
(only one channel) is identical to the input $x$, and the Jacobian is 1.
\nwenddocs{}\nwbegincode{776}\moddef{MCI vamp: test sampler 3: TBP}\plusendmoddef
  procedure :: evaluate => test_sampler_3_evaluate
\nwendcode{}\nwbegincode{777}\moddef{MCI vamp: tests}\plusendmoddef
  subroutine test_sampler_3_evaluate (sampler, c, x_in, val, x, f)
    class(test_sampler_3_t), intent(inout) :: sampler
    integer, intent(in) :: c
    real(default), dimension(:), intent(in) :: x_in
    real(default), intent(out) :: val
    real(default), dimension(:,:), intent(out) :: x
    real(default), dimension(:), intent(out) :: f
    call sampler%compute (c, x_in)
    call sampler%fetch (val, x, f)
  end subroutine test_sampler_3_evaluate

\eatline
\nwendcode{}\nwbegindocs{778}The point is always valid.
\nwenddocs{}\nwbegincode{779}\moddef{MCI vamp: test sampler 3: TBP}\plusendmoddef
  procedure :: is_valid => test_sampler_3_is_valid
\nwendcode{}\nwbegincode{780}\moddef{MCI vamp: tests}\plusendmoddef
  function test_sampler_3_is_valid (sampler) result (valid)
    class(test_sampler_3_t), intent(in) :: sampler
    logical :: valid
    valid = .true.
  end function test_sampler_3_is_valid
  
\eatline
\nwendcode{}\nwbegindocs{781}Rebuild: compute all but the function value.
\nwenddocs{}\nwbegincode{782}\moddef{MCI vamp: test sampler 3: TBP}\plusendmoddef
  procedure :: rebuild => test_sampler_3_rebuild
\nwendcode{}\nwbegincode{783}\moddef{MCI vamp: tests}\plusendmoddef
  subroutine test_sampler_3_rebuild (sampler, c, x_in, val, x, f)
    class(test_sampler_3_t), intent(inout) :: sampler
    integer, intent(in) :: c
    real(default), dimension(:), intent(in) :: x_in
    real(default), intent(in) :: val
    real(default), dimension(:,:), intent(out) :: x
    real(default), dimension(:), intent(out) :: f
    call sampler%compute (c, x_in)
    x = sampler%x
    f = sampler%f
  end subroutine test_sampler_3_rebuild

\eatline
\nwendcode{}\nwbegindocs{784}Extract the results.
\nwenddocs{}\nwbegincode{785}\moddef{MCI vamp: test sampler 3: TBP}\plusendmoddef
  procedure :: fetch => test_sampler_3_fetch
\nwendcode{}\nwbegincode{786}\moddef{MCI vamp: tests}\plusendmoddef
  subroutine test_sampler_3_fetch (sampler, val, x, f)
    class(test_sampler_3_t), intent(in) :: sampler
    real(default), intent(out) :: val
    real(default), dimension(:,:), intent(out) :: x
    real(default), dimension(:), intent(out) :: f
    val = sampler%val
    x = sampler%x
    f = sampler%f
  end subroutine test_sampler_3_fetch
    
\eatline
\nwendcode{}\nwbegindocs{787}\nwdocspar
\subsubsection{One-dimensional integration}
Construct an integrator and use it for a one-dimensional sampler.

Note: We would like to check the precise contents of the grid
allocated during integration, but the output format for reals is very
long (for good reasons), so the last digits in the grid content
display are numerical noise.  So, we just check the integration
results.
\nwenddocs{}\nwbegincode{788}\moddef{MCI vamp: execute tests}\endmoddef
  call test (mci_vamp_1, "mci_vamp_1", &
       "one-dimensional integral", &
       u, results)
\nwendcode{}\nwbegincode{789}\moddef{MCI vamp: tests}\plusendmoddef
  subroutine mci_vamp_1 (u)
    integer, intent(in) :: u
    type(grid_parameters_t) :: grid_par
    class(mci_t), allocatable, target :: mci
    class(mci_instance_t), pointer :: mci_instance => null ()
    class(mci_sampler_t), allocatable :: sampler
    class(rng_t), allocatable :: rng
    
    write (u, "(A)")  "* Test output: mci_vamp_1"
    write (u, "(A)")  "*   Purpose: integrate function in one dimension &
         &(single channel)"
    
    write (u, "(A)")
    write (u, "(A)")  "* Initialize integrator"
    write (u, "(A)")

    allocate (mci_vamp_t :: mci)
    call mci%set_dimensions (1, 1)
    select type (mci)
    type is (mci_vamp_t)
       grid_par%use_vamp_equivalences = .false.
       call mci%set_grid_parameters (grid_par)
    end select
    
    allocate (rng_tao_t :: rng)
    call rng%init ()
    call mci%import_rng (rng)

    call mci%write (u)
    
    write (u, "(A)")
    write (u, "(A)")  "* Initialize instance"
    write (u, "(A)")
    
    call mci%allocate_instance (mci_instance)
    call mci_instance%init (mci)
    
    write (u, "(A)")  "* Initialize test sampler"
    write (u, "(A)")
    
    allocate (test_sampler_1_t :: sampler)
    call sampler%write (u)

    write (u, "(A)")
    write (u, "(A)")  "* Integrate with n_calls = 1000"
    write (u, "(A)")  "   (lower precision to avoid"
    write (u, "(A)")  "      numerical noise)"
    write (u, "(A)")
    
    select type (mci)
    type is (mci_vamp_t)
       call mci%add_pass ()
    end select
    call mci%integrate (mci_instance, sampler, 1, 1000, pacify = .true.)
    call mci%write (u, .true.)
 
    write (u, "(A)")
    write (u, "(A)")  "* Contents of mci_instance:"
    write (u, "(A)")
    
    call mci_instance%write (u, .true.)

    write (u, "(A)")
    write (u, "(A)")  "* Cleanup"

    call mci_instance%final ()
    call mci%final ()

    write (u, "(A)")
    write (u, "(A)")  "* Test output end: mci_vamp_1"

  end subroutine mci_vamp_1

\eatline
\nwendcode{}\nwbegindocs{790}\nwdocspar
\subsubsection{Multiple iterations}
Construct an integrator and use it for a one-dimensional sampler.
Integrate with five iterations without grid adaptation.
\nwenddocs{}\nwbegincode{791}\moddef{MCI vamp: execute tests}\plusendmoddef
  call test (mci_vamp_2, "mci_vamp_2", &
       "multiple iterations", &
       u, results)
\nwendcode{}\nwbegincode{792}\moddef{MCI vamp: tests}\plusendmoddef
  subroutine mci_vamp_2 (u)
    integer, intent(in) :: u
    type(grid_parameters_t) :: grid_par
    class(mci_t), allocatable, target :: mci
    class(mci_instance_t), pointer :: mci_instance => null ()
    class(mci_sampler_t), allocatable :: sampler
    class(rng_t), allocatable :: rng
    
    write (u, "(A)")  "* Test output: mci_vamp_2"
    write (u, "(A)")  "*   Purpose: integrate function in one dimension &
         &(single channel)" 
    
    write (u, "(A)")
    write (u, "(A)")  "* Initialize integrator, sampler, instance"
    write (u, "(A)")

    allocate (mci_vamp_t :: mci)
    call mci%set_dimensions (1, 1)
    select type (mci)
    type is (mci_vamp_t)
       grid_par%use_vamp_equivalences = .false.
       call mci%set_grid_parameters (grid_par)
    end select
    
    allocate (rng_tao_t :: rng)
    call rng%init ()
    call mci%import_rng (rng)

    call mci%allocate_instance (mci_instance)
    call mci_instance%init (mci)
    
    allocate (test_sampler_1_t :: sampler)
    select type (sampler)
    type is (test_sampler_1_t)
       sampler%mode = 2
    end select
    call sampler%write (u)

    write (u, "(A)")
    write (u, "(A)")  "* Integrate with n_it = 3 and n_calls = 100"
    write (u, "(A)")
    
    select type (mci)
    type is (mci_vamp_t)
       call mci%add_pass (adapt_grids = .false.)
    end select
    call mci%integrate (mci_instance, sampler, 3, 100)
    call mci%write (u)
 
    write (u, "(A)")
    write (u, "(A)")  "* Contents of mci_instance:"
    write (u, "(A)")
    
    call mci_instance%write (u)
    
    write (u, "(A)")
    write (u, "(A)")  "* Cleanup"

    call mci_instance%final ()
    call mci%final ()

    write (u, "(A)")
    write (u, "(A)")  "* Test output end: mci_vamp_2"

  end subroutine mci_vamp_2

\eatline
\nwendcode{}\nwbegindocs{793}\nwdocspar
\subsubsection{Grid adaptation}
Construct an integrator and use it for a one-dimensional sampler.
Integrate with three iterations and in-between grid adaptations.
\nwenddocs{}\nwbegincode{794}\moddef{MCI vamp: execute tests}\plusendmoddef
  call test (mci_vamp_3, "mci_vamp_3", &
       "grid adaptation", &
       u, results)
\nwendcode{}\nwbegincode{795}\moddef{MCI vamp: tests}\plusendmoddef
  subroutine mci_vamp_3 (u)
    integer, intent(in) :: u
    type(grid_parameters_t) :: grid_par
    class(mci_t), allocatable, target :: mci
    class(mci_instance_t), pointer :: mci_instance => null ()
    class(mci_sampler_t), allocatable :: sampler
    class(rng_t), allocatable :: rng
    
    write (u, "(A)")  "* Test output: mci_vamp_3"
    write (u, "(A)")  "*   Purpose: integrate function in one dimension &
         &(single channel)" 
    write (u, "(A)")  "*            and adapt grid"
    
    write (u, "(A)")
    write (u, "(A)")  "* Initialize integrator, sampler, instance"
    write (u, "(A)")

    allocate (mci_vamp_t :: mci)
    call mci%set_dimensions (1, 1)
    select type (mci)
    type is (mci_vamp_t)
       grid_par%use_vamp_equivalences = .false.
       call mci%set_grid_parameters (grid_par)
    end select
    
    allocate (rng_tao_t :: rng)
    call rng%init ()
    call mci%import_rng (rng)

    call mci%allocate_instance (mci_instance)
    call mci_instance%init (mci)
    
    allocate (test_sampler_1_t :: sampler)
    select type (sampler)
    type is (test_sampler_1_t)
       sampler%mode = 2
    end select
    call sampler%write (u)

    write (u, "(A)")
    write (u, "(A)")  "* Integrate with n_it = 3 and n_calls = 100"
    write (u, "(A)")
    
    select type (mci)
    type is (mci_vamp_t)
       call mci%add_pass (adapt_grids = .true.)
    end select
    call mci%integrate (mci_instance, sampler, 3, 100)
    call mci%write (u)
 
    write (u, "(A)")
    write (u, "(A)")  "* Contents of mci_instance:"
    write (u, "(A)")
    
    call mci_instance%write (u)
    
    write (u, "(A)")
    write (u, "(A)")  "* Cleanup"

    call mci_instance%final ()
    call mci%final ()

    write (u, "(A)")
    write (u, "(A)")  "* Test output end: mci_vamp_3"

  end subroutine mci_vamp_3

\eatline
\nwendcode{}\nwbegindocs{796}\nwdocspar
\subsubsection{Two-dimensional integral}
Construct an integrator and use it for a two-dimensional sampler.
Integrate with three iterations and in-between grid adaptations.
\nwenddocs{}\nwbegincode{797}\moddef{MCI vamp: execute tests}\plusendmoddef
  call test (mci_vamp_4, "mci_vamp_4", &
       "two-dimensional integration", &
       u, results)
\nwendcode{}\nwbegincode{798}\moddef{MCI vamp: tests}\plusendmoddef
  subroutine mci_vamp_4 (u)
    integer, intent(in) :: u
    type(grid_parameters_t) :: grid_par
    class(mci_t), allocatable, target :: mci
    class(mci_instance_t), pointer :: mci_instance => null ()
    class(mci_sampler_t), allocatable :: sampler
    class(rng_t), allocatable :: rng
    
    write (u, "(A)")  "* Test output: mci_vamp_4"
    write (u, "(A)")  "*   Purpose: integrate function in two dimensions &
         &(single channel)" 
    write (u, "(A)")  "*            and adapt grid"
    
    write (u, "(A)")
    write (u, "(A)")  "* Initialize integrator, sampler, instance"
    write (u, "(A)")

    allocate (mci_vamp_t :: mci)
    call mci%set_dimensions (2, 1)
    select type (mci)
    type is (mci_vamp_t)
       grid_par%use_vamp_equivalences = .false.
       call mci%set_grid_parameters (grid_par)
    end select
    
    allocate (rng_tao_t :: rng)
    call rng%init ()
    call mci%import_rng (rng)

    call mci%allocate_instance (mci_instance)
    call mci_instance%init (mci)
    
    allocate (test_sampler_1_t :: sampler)
    select type (sampler)
    type is (test_sampler_1_t)
       sampler%mode = 3
    end select
    call sampler%write (u)

    write (u, "(A)")
    write (u, "(A)")  "* Integrate with n_it = 3 and n_calls = 1000"
    write (u, "(A)")
    
    select type (mci)
    type is (mci_vamp_t)
       call mci%add_pass (adapt_grids = .true.)
    end select
    call mci%integrate (mci_instance, sampler, 3, 1000)
    call mci%write (u)
 
    write (u, "(A)")
    write (u, "(A)")  "* Contents of mci_instance:"
    write (u, "(A)")
    
    call mci_instance%write (u)
    
    write (u, "(A)")
    write (u, "(A)")  "* Cleanup"

    call mci_instance%final ()
    call mci%final ()

    write (u, "(A)")
    write (u, "(A)")  "* Test output end: mci_vamp_4"

  end subroutine mci_vamp_4

\eatline
\nwendcode{}\nwbegindocs{799}\nwdocspar
\subsubsection{Two-channel integral}
Construct an integrator and use it for a two-dimensional sampler with two
channels.

Integrate with three iterations and in-between grid adaptations.
\nwenddocs{}\nwbegincode{800}\moddef{MCI vamp: execute tests}\plusendmoddef
  call test (mci_vamp_5, "mci_vamp_5", &
       "two-dimensional integration", &
       u, results)
\nwendcode{}\nwbegincode{801}\moddef{MCI vamp: tests}\plusendmoddef
  subroutine mci_vamp_5 (u)
    integer, intent(in) :: u
    type(grid_parameters_t) :: grid_par
    class(mci_t), allocatable, target :: mci
    class(mci_instance_t), pointer :: mci_instance => null ()
    class(mci_sampler_t), allocatable :: sampler
    class(rng_t), allocatable :: rng
    
    write (u, "(A)")  "* Test output: mci_vamp_5"
    write (u, "(A)")  "*   Purpose: integrate function in two dimensions &
         &(two channels)" 
    write (u, "(A)")  "*            and adapt grid"
    
    write (u, "(A)")
    write (u, "(A)")  "* Initialize integrator, sampler, instance"
    write (u, "(A)")

    allocate (mci_vamp_t :: mci)
    call mci%set_dimensions (2, 2)
    select type (mci)
    type is (mci_vamp_t)
       grid_par%stratified = .false.
       grid_par%use_vamp_equivalences = .false.
       call mci%set_grid_parameters (grid_par)
    end select
    
    allocate (rng_tao_t :: rng)
    call rng%init ()
    call mci%import_rng (rng)

    call mci%allocate_instance (mci_instance)
    call mci_instance%init (mci)
    
    allocate (test_sampler_2_t :: sampler)
    call sampler%write (u)

    write (u, "(A)")
    write (u, "(A)")  "* Integrate with n_it = 3 and n_calls = 1000"
    write (u, "(A)")
    
    select type (mci)
    type is (mci_vamp_t)
       call mci%add_pass (adapt_grids = .true.)
    end select
    call mci%integrate (mci_instance, sampler, 3, 1000)
    call mci%write (u)
 
    write (u, "(A)")
    write (u, "(A)")  "* Contents of mci_instance:"
    write (u, "(A)")
    
    call mci_instance%write (u)
    
    write (u, "(A)")
    write (u, "(A)")  "* Cleanup"

    call mci_instance%final ()
    call mci%final ()

    write (u, "(A)")
    write (u, "(A)")  "* Test output end: mci_vamp_5"

  end subroutine mci_vamp_5

\eatline
\nwendcode{}\nwbegindocs{802}\nwdocspar
\subsubsection{Weight adaptation}
Construct an integrator and use it for a one-dimensional sampler with two
channels.

Integrate with three iterations and in-between weight adaptations.
\nwenddocs{}\nwbegincode{803}\moddef{MCI vamp: execute tests}\plusendmoddef
  call test (mci_vamp_6, "mci_vamp_6", &
       "weight adaptation", &
       u, results)
\nwendcode{}\nwbegincode{804}\moddef{MCI vamp: tests}\plusendmoddef
  subroutine mci_vamp_6 (u)
    integer, intent(in) :: u
    type(grid_parameters_t) :: grid_par
    class(mci_t), allocatable, target :: mci
    class(mci_instance_t), pointer :: mci_instance => null ()
    class(mci_sampler_t), allocatable :: sampler
    class(rng_t), allocatable :: rng
    
    write (u, "(A)")  "* Test output: mci_vamp_6"
    write (u, "(A)")  "*   Purpose: integrate function in one dimension &
         &(two channels)" 
    write (u, "(A)")  "*            and adapt weights"
    
    write (u, "(A)")
    write (u, "(A)")  "* Initialize integrator, sampler, instance"
    write (u, "(A)")

    allocate (mci_vamp_t :: mci)
    call mci%set_dimensions (1, 2)
    select type (mci)
    type is (mci_vamp_t)
       grid_par%stratified = .false.
       grid_par%use_vamp_equivalences = .false.
       call mci%set_grid_parameters (grid_par)
    end select
    
    allocate (rng_tao_t :: rng)
    call rng%init ()
    call mci%import_rng (rng)

    call mci%allocate_instance (mci_instance)
    call mci_instance%init (mci)
    
    allocate (test_sampler_3_t :: sampler)
    select type (sampler)
    type is (test_sampler_3_t)
       sampler%a = 0.9_default
       sampler%b = 0.1_default
    end select
    call sampler%write (u)

    write (u, "(A)")
    write (u, "(A)")  "* Integrate with n_it = 3 and n_calls = 1000"
    write (u, "(A)")
    
    select type (mci)
    type is (mci_vamp_t)
       call mci%add_pass (adapt_weights = .true.)
    end select
    call mci%integrate (mci_instance, sampler, 3, 1000)
    call mci%write (u)
 
    write (u, "(A)")
    write (u, "(A)")  "* Contents of mci_instance:"
    write (u, "(A)")
    
    call mci_instance%write (u)
    
    write (u, "(A)")
    write (u, "(A)")  "* Cleanup"

    call mci_instance%final ()
    call mci%final ()
    deallocate (mci_instance)
    deallocate (mci)

    write (u, "(A)")
    write (u, "(A)")  "* Re-initialize with chained channels"
    write (u, "(A)")

    allocate (mci_vamp_t :: mci)
    call mci%set_dimensions (1, 2)
    call mci%declare_chains ([1,1])
    select type (mci)
    type is (mci_vamp_t)
       grid_par%stratified = .false.
       grid_par%use_vamp_equivalences = .false.
       call mci%set_grid_parameters (grid_par)
    end select
    
    allocate (rng_tao_t :: rng)
    call rng%init ()
    call mci%import_rng (rng)

    call mci%allocate_instance (mci_instance)
    call mci_instance%init (mci)
    
    write (u, "(A)")  "* Integrate with n_it = 3 and n_calls = 1000"
    write (u, "(A)")
    
    select type (mci)
    type is (mci_vamp_t)
       call mci%add_pass (adapt_weights = .true.)
    end select
    call mci%integrate (mci_instance, sampler, 3, 1000)
    call mci%write (u)
 
    write (u, "(A)")
    write (u, "(A)")  "* Contents of mci_instance:"
    write (u, "(A)")
    
    call mci_instance%write (u)
    
    write (u, "(A)")
    write (u, "(A)")  "* Cleanup"

    call mci_instance%final ()
    call mci%final ()

    write (u, "(A)")
    write (u, "(A)")  "* Test output end: mci_vamp_6"

  end subroutine mci_vamp_6

\eatline
\nwendcode{}\nwbegindocs{805}\nwdocspar
\subsubsection{Equivalences}
Construct an integrator and use it for a one-dimensional sampler with two
channels.

Integrate with three iterations and in-between grid adaptations.
Apply an equivalence between the two channels, so the binning of the
two channels is forced to coincide.  Compare this with the behavior
without equivalences.
\nwenddocs{}\nwbegincode{806}\moddef{MCI vamp: execute tests}\plusendmoddef
  call test (mci_vamp_7, "mci_vamp_7", &
       "use channel equivalences", &
       u, results)
\nwendcode{}\nwbegincode{807}\moddef{MCI vamp: tests}\plusendmoddef
  subroutine mci_vamp_7 (u)
    integer, intent(in) :: u
    type(grid_parameters_t) :: grid_par
    class(mci_t), allocatable, target :: mci
    class(mci_instance_t), pointer :: mci_instance => null ()
    class(mci_sampler_t), allocatable :: sampler
    type(phs_channel_t), dimension(:), allocatable :: channel
    class(rng_t), allocatable :: rng
    real(default), dimension(:,:), allocatable :: x
    integer :: u_grid, iostat, i, div, ch
    character(16) :: buffer
    
    write (u, "(A)")  "* Test output: mci_vamp_7"
    write (u, "(A)")  "*   Purpose: check effect of channel equivalences"
    
    write (u, "(A)")
    write (u, "(A)")  "* Initialize integrator, sampler, instance"
    write (u, "(A)")

    allocate (mci_vamp_t :: mci)
    call mci%set_dimensions (1, 2)
    select type (mci)
    type is (mci_vamp_t)
       grid_par%stratified = .false.
       grid_par%use_vamp_equivalences = .false.
       call mci%set_grid_parameters (grid_par)
    end select
    
    allocate (rng_tao_t :: rng)
    call rng%init ()
    call mci%import_rng (rng)

    call mci%allocate_instance (mci_instance)
    call mci_instance%init (mci)
    
    allocate (test_sampler_3_t :: sampler)
    select type (sampler)
    type is (test_sampler_3_t)
       sampler%a = 0.7_default
       sampler%b = 0.3_default
    end select
    call sampler%write (u)

    write (u, "(A)")
    write (u, "(A)")  "* Integrate with n_it = 2 and n_calls = 1000, &
         &adapt grids"
    write (u, "(A)")
    
    select type (mci)
    type is (mci_vamp_t)
       call mci%add_pass (adapt_grids = .true.)
    end select
    call mci%integrate (mci_instance, sampler, 2, 1000)

    call mci%write (u)
    
    write (u, "(A)")
    write (u, "(A)") "* Write grids and extract binning"
    write (u, "(A)")

    u_grid = free_unit ()
    open (u_grid, status = "scratch", action = "readwrite")
    select type (mci_instance)
    type is (mci_vamp_instance_t)
       call vamp_write_grids (mci_instance%grids, u_grid)
    end select
    rewind (u_grid)
    allocate (x (0:20, 2))
    do div = 1, 2
       FIND_BINS1: do
          read (u_grid, "(A)")  buffer
          if (trim (adjustl (buffer)) == "begin d%x") then
             do
                read (u_grid, *, iostat = iostat)  i, x(i,div)
                if (iostat /= 0)  exit FIND_BINS1
             end do
          end if
       end do FIND_BINS1
    end do
    close (u_grid)

    write (u, "(1x,A,L1)")  "Equal binning in both channels = ", &
         all (x(:,1) == x(:,2))
    deallocate (x)
    
    write (u, "(A)")
    write (u, "(A)")  "* Cleanup"

    call mci_instance%final ()
    call mci%final ()
    deallocate (mci_instance)
    deallocate (mci)

    write (u, "(A)")
    write (u, "(A)")  "* Re-initialize integrator, instance"
    write (u, "(A)")

    allocate (mci_vamp_t :: mci)
    call mci%set_dimensions (1, 2)
    select type (mci)
    type is (mci_vamp_t)
       grid_par%stratified = .false.
       grid_par%use_vamp_equivalences = .true.
       call mci%set_grid_parameters (grid_par)
    end select
    
    write (u, "(A)")  "* Define equivalences"
    write (u, "(A)")

    allocate (channel (2))
    do ch = 1, 2
       allocate (channel(ch)%eq (2))
       do i = 1, 2
          associate (eq => channel(ch)%eq(i))
            call eq%init (1)
            eq%c = i
            eq%perm = [1]
            eq%mode = [0]
          end associate
       end do
       write (u, "(1x,I0,':')", advance = "no")  ch
       call channel(ch)%write (u)
    end do
    call mci%declare_equivalences (channel, dim_offset = 0)

    allocate (rng_tao_t :: rng)
    call rng%init ()
    call mci%import_rng (rng)

    call mci%allocate_instance (mci_instance)
    call mci_instance%init (mci)
    
    write (u, "(A)")
    write (u, "(A)")  "* Integrate with n_it = 2 and n_calls = 1000, &
         &adapt grids"
    write (u, "(A)")
    
    select type (mci)
    type is (mci_vamp_t)
       call mci%add_pass (adapt_grids = .true.)
    end select
    call mci%integrate (mci_instance, sampler, 2, 1000)
 
    call mci%write (u)
    
    write (u, "(A)")
    write (u, "(A)") "* Write grids and extract binning"
    write (u, "(A)")

    u_grid = free_unit ()
    open (u_grid, status = "scratch", action = "readwrite")
    select type (mci_instance)
    type is (mci_vamp_instance_t)
       call vamp_write_grids (mci_instance%grids, u_grid)
    end select
    rewind (u_grid)
    allocate (x (0:20, 2))
    do div = 1, 2
       FIND_BINS2: do
          read (u_grid, "(A)")  buffer
          if (trim (adjustl (buffer)) == "begin d%x") then
             do
                read (u_grid, *, iostat = iostat)  i, x(i,div)
                if (iostat /= 0)  exit FIND_BINS2
             end do
          end if
       end do FIND_BINS2
    end do
    close (u_grid)

    write (u, "(1x,A,L1)")  "Equal binning in both channels = ", &
         all (x(:,1) == x(:,2))
    deallocate (x)
    

    write (u, "(A)")
    write (u, "(A)")  "* Cleanup"

    call mci_instance%final ()
    call mci%final ()

    write (u, "(A)")
    write (u, "(A)")  "* Test output end: mci_vamp_7"

  end subroutine mci_vamp_7

\eatline
\nwendcode{}\nwbegindocs{808}\nwdocspar
\subsubsection{Multiple passes}
Integrate with three passes and different settings for weight and grid
adaptation.
\nwenddocs{}\nwbegincode{809}\moddef{MCI vamp: execute tests}\plusendmoddef
  call test (mci_vamp_8, "mci_vamp_8", &
       "integration passes", &
       u, results)
\nwendcode{}\nwbegincode{810}\moddef{MCI vamp: tests}\plusendmoddef
  subroutine mci_vamp_8 (u)
    integer, intent(in) :: u
    type(grid_parameters_t) :: grid_par
    class(mci_t), allocatable, target :: mci
    class(mci_instance_t), pointer :: mci_instance => null ()
    class(mci_sampler_t), allocatable :: sampler
    class(rng_t), allocatable :: rng
    
    write (u, "(A)")  "* Test output: mci_vamp_8"
    write (u, "(A)")  "*   Purpose: integrate function in one dimension &
         &(two channels)" 
    write (u, "(A)")  "*            in three passes"
    
    write (u, "(A)")
    write (u, "(A)")  "* Initialize integrator, sampler, instance"
    write (u, "(A)")

    allocate (mci_vamp_t :: mci)
    call mci%set_dimensions (1, 2)
    select type (mci)
    type is (mci_vamp_t)
       grid_par%stratified = .false.
       grid_par%use_vamp_equivalences = .false.
       call mci%set_grid_parameters (grid_par)
    end select
    
    allocate (rng_tao_t :: rng)
    call rng%init ()
    call mci%import_rng (rng)

    call mci%allocate_instance (mci_instance)
    call mci_instance%init (mci)
    
    allocate (test_sampler_3_t :: sampler)
    select type (sampler)
    type is (test_sampler_3_t)
       sampler%a = 0.9_default
       sampler%b = 0.1_default
    end select
    call sampler%write (u)

    write (u, "(A)")
    write (u, "(A)")  "* Integrate with grid and weight adaptation"
    write (u, "(A)")
    
    select type (mci)
    type is (mci_vamp_t)
       call mci%add_pass (adapt_grids = .true., adapt_weights = .true.)
    end select
    call mci%integrate (mci_instance, sampler, 3, 1000)
    call mci%write (u)
 
    write (u, "(A)")
    write (u, "(A)")  "* Contents of mci_instance:"
    write (u, "(A)")
    
    call mci_instance%write (u)
    
    write (u, "(A)")
    write (u, "(A)")  "* Integrate with grid adaptation"
    write (u, "(A)")
    
    select type (mci)
    type is (mci_vamp_t)
       call mci%add_pass (adapt_grids = .true.)
    end select
    call mci%integrate (mci_instance, sampler, 3, 1000)
    call mci%write (u)
 
    write (u, "(A)")
    write (u, "(A)")  "* Contents of mci_instance:"
    write (u, "(A)")
    
    call mci_instance%write (u)
    
    write (u, "(A)")
    write (u, "(A)")  "* Integrate without adaptation"
    write (u, "(A)")
    
    select type (mci)
    type is (mci_vamp_t)
       call mci%add_pass ()
    end select
    call mci%integrate (mci_instance, sampler, 3, 1000)
    call mci%write (u)
 
    write (u, "(A)")
    write (u, "(A)")  "* Contents of mci_instance:"
    write (u, "(A)")
    
    call mci_instance%write (u)
    
    write (u, "(A)")
    write (u, "(A)")  "* Cleanup"

    call mci_instance%final ()
    call mci%final ()

    write (u, "(A)")
    write (u, "(A)")  "* Test output end: mci_vamp_8"

  end subroutine mci_vamp_8

\eatline
\nwendcode{}\nwbegindocs{811}\nwdocspar
\subsubsection{Weighted events}
Construct an integrator and use it for a two-dimensional sampler with two
channels.  Integrate and generate a weighted event.
\nwenddocs{}\nwbegincode{812}\moddef{MCI vamp: execute tests}\plusendmoddef
  call test (mci_vamp_9, "mci_vamp_9", &
       "weighted event", &
       u, results)
\nwendcode{}\nwbegincode{813}\moddef{MCI vamp: tests}\plusendmoddef
  subroutine mci_vamp_9 (u)
    integer, intent(in) :: u
    type(grid_parameters_t) :: grid_par
    class(mci_t), allocatable, target :: mci
    class(mci_instance_t), pointer :: mci_instance => null ()
    class(mci_sampler_t), allocatable :: sampler
    class(rng_t), allocatable :: rng
    
    write (u, "(A)")  "* Test output: mci_vamp_9"
    write (u, "(A)")  "*   Purpose: integrate function in two dimensions &
         &(two channels)" 
    write (u, "(A)")  "*            and generate a weighted event"
    
    write (u, "(A)")
    write (u, "(A)")  "* Initialize integrator, sampler, instance"
    write (u, "(A)")

    allocate (mci_vamp_t :: mci)
    call mci%set_dimensions (2, 2)
    select type (mci)
    type is (mci_vamp_t)
       grid_par%stratified = .false.
       grid_par%use_vamp_equivalences = .false.
       call mci%set_grid_parameters (grid_par)
    end select
    
    allocate (rng_tao_t :: rng)
    call rng%init ()
    call mci%import_rng (rng)

    call mci%allocate_instance (mci_instance)
    call mci_instance%init (mci)
    
    allocate (test_sampler_2_t :: sampler)
    call sampler%write (u)

    write (u, "(A)")
    write (u, "(A)")  "* Integrate with n_it = 3 and n_calls = 1000"
    write (u, "(A)")
    
    call mci%add_pass ()
    call mci%integrate (mci_instance, sampler, 1, 1000)
    call mci%write (u)
 
    write (u, "(A)")
    write (u, "(A)")  "* Generate a weighted event"
    write (u, "(A)")

    call mci_instance%init_simulation ()
    call mci%generate_weighted_event (mci_instance, sampler)

    write (u, "(1x,A)")  "MCI instance:"
    call mci_instance%write (u)
    
    write (u, "(A)")
    write (u, "(A)")  "* Cleanup"

    call mci_instance%final_simulation ()
    call mci_instance%final ()
    call mci%final ()

    write (u, "(A)")
    write (u, "(A)")  "* Test output end: mci_vamp_9"

  end subroutine mci_vamp_9

\eatline
\nwendcode{}\nwbegindocs{814}\nwdocspar
\subsubsection{Grids I/O}
Construct an integrator and allocate grids.  Write grids to file, read
them in again and compare.
\nwenddocs{}\nwbegincode{815}\moddef{MCI vamp: execute tests}\plusendmoddef
  call test (mci_vamp_10, "mci_vamp_10", &
       "grids I/O", &
       u, results)
\nwendcode{}\nwbegincode{816}\moddef{MCI vamp: tests}\plusendmoddef
  subroutine mci_vamp_10 (u)
    integer, intent(in) :: u
    type(grid_parameters_t) :: grid_par
    class(mci_t), allocatable, target :: mci
    class(mci_instance_t), pointer :: mci_instance => null ()
    class(mci_sampler_t), allocatable :: sampler
    class(rng_t), allocatable :: rng
    type(string_t) :: file1, file2
    character(80) :: buffer1, buffer2
    integer :: u1, u2, iostat1, iostat2
    logical :: equal, success
   
    write (u, "(A)")  "* Test output: mci_vamp_10"
    write (u, "(A)")  "*   Purpose: write and read VAMP grids"
    
    write (u, "(A)")
    write (u, "(A)")  "* Initialize integrator, sampler, instance"
    write (u, "(A)")

    allocate (mci_vamp_t :: mci)
    call mci%set_dimensions (2, 2)
    select type (mci)
    type is (mci_vamp_t)
       grid_par%stratified = .false.
       grid_par%use_vamp_equivalences = .false.
       call mci%set_grid_parameters (grid_par)
    end select
    
    allocate (rng_tao_t :: rng)
    call rng%init ()
    call mci%import_rng (rng)

    mci%md5sum = "1234567890abcdef1234567890abcdef"

    call mci%allocate_instance (mci_instance)
    call mci_instance%init (mci)
    
    allocate (test_sampler_2_t :: sampler)
    call sampler%write (u)

    write (u, "(A)")
    write (u, "(A)")  "* Integrate with n_it = 3 and n_calls = 1000"
    write (u, "(A)")
    
    call mci%add_pass ()
    call mci%integrate (mci_instance, sampler, 1, 1000)
 
    write (u, "(A)")  "* Write grids to file"
    write (u, "(A)")
 
    file1 = "mci_vamp_10.1"
    select type (mci)
    type is (mci_vamp_t)
       call mci%set_grid_filename (file1)
       call mci%write_grids (mci_instance)
    end select

    call mci_instance%final ()
    call mci%final ()
    deallocate (mci)
    
    write (u, "(A)")  "* Read grids from file"
    write (u, "(A)")

    allocate (mci_vamp_t :: mci)
    call mci%set_dimensions (2, 2)
    select type (mci)
    type is (mci_vamp_t)
       call mci%set_grid_parameters (grid_par)
    end select
    
    allocate (rng_tao_t :: rng)
    call rng%init ()
    call mci%import_rng (rng)

    mci%md5sum = "1234567890abcdef1234567890abcdef"

    call mci%allocate_instance (mci_instance)
    call mci_instance%init (mci)

    select type (mci)
    type is (mci_vamp_t)
       call mci%set_grid_filename (file1)
       call mci%add_pass ()
       call mci%current_pass%configure (1, 1000, &
            mci%min_calls, &
            mci%grid_par%min_bins, mci%grid_par%max_bins, &
            mci%grid_par%min_calls_per_channel * mci%n_channel)
       call mci%read_grids_header (success)
       call mci%compute_md5sum ()
       call mci%read_grids_data (mci_instance, read_integrals = .true.)
    end select
    write (u, "(1x,A,L1)")  "success = ", success
    
    write (u, "(A)")
    write (u, "(A)")  "* Write grids again"
    write (u, "(A)")

    file2 = "mci_vamp_10.2"
    select type (mci)
    type is (mci_vamp_t)
       call mci%set_grid_filename (file2)
       call mci%write_grids (mci_instance)
    end select

    u1 = free_unit ()
    open (u1, file = char (file1) // ".vg", action = "read", status = "old")
    u2 = free_unit ()
    open (u2, file = char (file2) // ".vg", action = "read", status = "old")

    equal = .true.
    iostat1 = 0
    iostat2 = 0
    do while (equal .and. iostat1 == 0 .and. iostat2 == 0)
       read (u1, "(A)", iostat = iostat1)  buffer1
       read (u2, "(A)", iostat = iostat2)  buffer2
       equal = buffer1 == buffer2 .and. iostat1 == iostat2
    end do
    close (u1)
    close (u2)
    
    if (equal) then
       write (u, "(1x,A)")  "Success: grid files are identical"
    else
       write (u, "(1x,A)")  "Failure: grid files differ"
    end if

    write (u, "(A)")
    write (u, "(A)")  "* Cleanup"
    
    call mci_instance%final ()
    call mci%final ()

    write (u, "(A)")
    write (u, "(A)")  "* Test output end: mci_vamp_10"

  end subroutine mci_vamp_10

\eatline
\nwendcode{}\nwbegindocs{817}\nwdocspar
\subsubsection{Weighted events}
Construct an integrator and use it for a two-dimensional sampler with two
channels.  Integrate, write grids, and generate a weighted event using
the grids from file.
\nwenddocs{}\nwbegincode{818}\moddef{MCI vamp: execute tests}\plusendmoddef
 call test (mci_vamp_11, "mci_vamp_11", &
      "weighted events with grid I/O", &
      u, results)
\nwendcode{}\nwbegincode{819}\moddef{MCI vamp: tests}\plusendmoddef
  subroutine mci_vamp_11 (u)
    integer, intent(in) :: u
    type(grid_parameters_t) :: grid_par
    class(mci_t), allocatable, target :: mci
    class(mci_instance_t), pointer :: mci_instance => null ()
    class(mci_sampler_t), allocatable :: sampler
    class(rng_t), allocatable :: rng
    
    write (u, "(A)")  "* Test output: mci_vamp_11"
    write (u, "(A)")  "*   Purpose: integrate function in two dimensions &
         &(two channels)" 
    write (u, "(A)")  "*            and generate a weighted event"
    
    write (u, "(A)")
    write (u, "(A)")  "* Initialize integrator, sampler, instance"
    write (u, "(A)")

    allocate (mci_vamp_t :: mci)
    call mci%set_dimensions (2, 2)
    select type (mci)
    type is (mci_vamp_t)
       grid_par%stratified = .false.
       grid_par%use_vamp_equivalences = .false.
       call mci%set_grid_parameters (grid_par)
       call mci%set_grid_filename (var_str ("mci_vamp_11"))
    end select
    
    allocate (rng_tao_t :: rng)
    call rng%init ()
    call mci%import_rng (rng)

    call mci%allocate_instance (mci_instance)
    call mci_instance%init (mci)
    
    allocate (test_sampler_2_t :: sampler)

    write (u, "(A)")  "* Integrate with n_it = 3 and n_calls = 1000"
    write (u, "(A)")
    
    call mci%add_pass ()
    call mci%integrate (mci_instance, sampler, 1, 1000)
 
    write (u, "(A)")  "* Reset instance"
    write (u, "(A)")

    call mci_instance%final ()
    call mci%allocate_instance (mci_instance)
    call mci_instance%init (mci)

    write (u, "(A)")  "* Generate a weighted event"
    write (u, "(A)")

    call mci_instance%init_simulation ()
    call mci%generate_weighted_event (mci_instance, sampler)

    write (u, "(A)")  "* Cleanup"

    call mci_instance%final_simulation ()
    call mci_instance%final ()
    call mci%final ()

    write (u, "(A)")
    write (u, "(A)")  "* Test output end: mci_vamp_11"

  end subroutine mci_vamp_11

\eatline
\nwendcode{}\nwbegindocs{820}\nwdocspar
\subsubsection{Weighted events}
Construct an integrator and use it for a two-dimensional sampler with two
channels.  Integrate, write grids, and generate a weighted event using
the grids from file.
\nwenddocs{}\nwbegincode{821}\moddef{MCI vamp: execute tests}\plusendmoddef
 call test (mci_vamp_12, "mci_vamp_12", &
      "unweighted events with grid I/O", &
      u, results)
\nwendcode{}\nwbegincode{822}\moddef{MCI vamp: tests}\plusendmoddef
  subroutine mci_vamp_12 (u)
    integer, intent(in) :: u
    type(grid_parameters_t) :: grid_par
    class(mci_t), allocatable, target :: mci
    class(mci_instance_t), pointer :: mci_instance => null ()
    class(mci_sampler_t), allocatable :: sampler
    class(rng_t), allocatable :: rng
    
    write (u, "(A)")  "* Test output: mci_vamp_12"
    write (u, "(A)")  "*   Purpose: integrate function in two dimensions &
         &(two channels)" 
    write (u, "(A)")  "*            and generate an unweighted event"
    
    write (u, "(A)")
    write (u, "(A)")  "* Initialize integrator, sampler, instance"
    write (u, "(A)")

    allocate (mci_vamp_t :: mci)
    call mci%set_dimensions (2, 2)
    select type (mci)
    type is (mci_vamp_t)
       grid_par%stratified = .false.
       grid_par%use_vamp_equivalences = .false.
       call mci%set_grid_parameters (grid_par)
       call mci%set_grid_filename (var_str ("mci_vamp_12"))
    end select
    
    allocate (rng_tao_t :: rng)
    call rng%init ()
    call mci%import_rng (rng)

    call mci%allocate_instance (mci_instance)
    call mci_instance%init (mci)
    
    allocate (test_sampler_2_t :: sampler)

    write (u, "(A)")  "* Integrate with n_it = 3 and n_calls = 1000"
    write (u, "(A)")
    
    call mci%add_pass ()
    call mci%integrate (mci_instance, sampler, 1, 1000)
 
    write (u, "(A)")  "* Reset instance"
    write (u, "(A)")

    call mci_instance%final ()
    call mci%allocate_instance (mci_instance)
    call mci_instance%init (mci)

    write (u, "(A)")  "* Generate an unweighted event"
    write (u, "(A)")

    call mci_instance%init_simulation ()
    call mci%generate_unweighted_event (mci_instance, sampler)

    write (u, "(1x,A)")  "MCI instance:"
    call mci_instance%write (u)
    
    write (u, "(A)")
    write (u, "(A)")  "* Cleanup"

    call mci_instance%final_simulation ()
    call mci_instance%final ()
    call mci%final ()

    write (u, "(A)")
    write (u, "(A)")  "* Test output end: mci_vamp_12"

  end subroutine mci_vamp_12

\eatline
\nwendcode{}\nwbegindocs{823}\nwdocspar
\subsubsection{Update integration results}
Compare two {\tt{}mci} objects; match the two and update the first if
successful.
\nwenddocs{}\nwbegincode{824}\moddef{MCI vamp: execute tests}\plusendmoddef
  call test (mci_vamp_13, "mci_vamp_13", &
       "updating integration results", &
       u, results)
\nwendcode{}\nwbegincode{825}\moddef{MCI vamp: tests}\plusendmoddef
  subroutine mci_vamp_13 (u)
    integer, intent(in) :: u
    type(grid_parameters_t) :: grid_par
    class(mci_t), allocatable, target :: mci, mci_ref
    logical :: success
    
    write (u, "(A)")  "* Test output: mci_vamp_13"
    write (u, "(A)")  "*   Purpose: match and update integrators"
    
    write (u, "(A)")
    write (u, "(A)")  "* Initialize integrator with no passes"
    write (u, "(A)")

    allocate (mci_vamp_t :: mci)
    call mci%set_dimensions (2, 2)
    select type (mci)
    type is (mci_vamp_t)
       grid_par%stratified = .false.
       grid_par%use_vamp_equivalences = .false.
       call mci%set_grid_parameters (grid_par)
    end select
    call mci%write (u)
    
    write (u, "(A)")
    write (u, "(A)")  "* Initialize reference"
    write (u, "(A)")

    allocate (mci_vamp_t :: mci_ref)
    call mci_ref%set_dimensions (2, 2)
    select type (mci_ref)
    type is (mci_vamp_t)
       call mci_ref%set_grid_parameters (grid_par)
    end select
    
    select type (mci_ref)
    type is (mci_vamp_t)
       call mci_ref%add_pass (adapt_grids = .true.)
       call mci_ref%current_pass%configure (2, 1000, 0, 1, 5, 0)
       mci_ref%current_pass%calls = [77, 77]
       mci_ref%current_pass%integral = [1.23_default, 3.45_default]
       mci_ref%current_pass%error = [0.23_default, 0.45_default]
       mci_ref%current_pass%efficiency = [0.1_default, 0.6_default]
       mci_ref%current_pass%integral_defined = .true.

       call mci_ref%add_pass ()
       call mci_ref%current_pass%configure (2, 2000, 0, 1, 7, 0)
       mci_ref%current_pass%calls = [99, 0]
       mci_ref%current_pass%integral = [7.89_default, 0._default]
       mci_ref%current_pass%error = [0.89_default, 0._default]
       mci_ref%current_pass%efficiency = [0.86_default, 0._default]
       mci_ref%current_pass%integral_defined = .true.
    end select

    call mci_ref%write (u)
    
    write (u, "(A)")
    write (u, "(A)")  "* Update integrator (no-op, should succeed)"
    write (u, "(A)")

    select type (mci)
    type is (mci_vamp_t)
       call mci%update_from_ref (mci_ref, success)
    end select

    write (u, "(1x,A,L1)")  "success = ", success
    write (u, "(A)")
    call mci%write (u)
    
    write (u, "(A)")
    write (u, "(A)")  "* Add pass to integrator"
    write (u, "(A)")

    select type (mci)
    type is (mci_vamp_t)
       call mci%add_pass (adapt_grids = .true.)
       call mci%current_pass%configure (2, 1000, 0, 1, 5, 0)
       mci%current_pass%calls = [77, 77]
       mci%current_pass%integral = [1.23_default, 3.45_default]
       mci%current_pass%error = [0.23_default, 0.45_default]
       mci%current_pass%efficiency = [0.1_default, 0.6_default]
       mci%current_pass%integral_defined = .true.
    end select

    write (u, "(A)")  "* Update integrator (no-op, should succeed)"
    write (u, "(A)")

    select type (mci)
    type is (mci_vamp_t)
       call mci%update_from_ref (mci_ref, success)
    end select

    write (u, "(1x,A,L1)")  "success = ", success
    write (u, "(A)")
    call mci%write (u)
    
    write (u, "(A)")
    write (u, "(A)")  "* Add pass to integrator, wrong parameters"
    write (u, "(A)")

    select type (mci)
    type is (mci_vamp_t)
       call mci%add_pass ()
       call mci%current_pass%configure (2, 1000, 0, 1, 7, 0)
    end select

    write (u, "(A)")  "* Update integrator (should fail)"
    write (u, "(A)")

    select type (mci)
    type is (mci_vamp_t)
       call mci%update_from_ref (mci_ref, success)
    end select

    write (u, "(1x,A,L1)")  "success = ", success
    write (u, "(A)")
    call mci%write (u)
    
    write (u, "(A)")
    write (u, "(A)")  "* Reset and add passes to integrator"
    write (u, "(A)")

    select type (mci)
    type is (mci_vamp_t)
       call mci%reset ()
       call mci%add_pass (adapt_grids = .true.)
       call mci%current_pass%configure (2, 1000, 0, 1, 5, 0)
       mci%current_pass%calls = [77, 77]
       mci%current_pass%integral = [1.23_default, 3.45_default]
       mci%current_pass%error = [0.23_default, 0.45_default]
       mci%current_pass%efficiency = [0.1_default, 0.6_default]
       mci%current_pass%integral_defined = .true.
       
       call mci%add_pass ()
       call mci%current_pass%configure (2, 2000, 0, 1, 7, 0)
    end select

    write (u, "(A)")  "* Update integrator (should succeed)"
    write (u, "(A)")

    select type (mci)
    type is (mci_vamp_t)
       call mci%update_from_ref (mci_ref, success)
    end select

    write (u, "(1x,A,L1)")  "success = ", success
    write (u, "(A)")
    call mci%write (u)
    
    write (u, "(A)")
    write (u, "(A)")  "* Update again (no-op, should succeed)"
    write (u, "(A)")

    select type (mci)
    type is (mci_vamp_t)
       call mci%update_from_ref (mci_ref, success)
    end select

    write (u, "(1x,A,L1)")  "success = ", success
    write (u, "(A)")
    call mci%write (u)
    
    write (u, "(A)")
    write (u, "(A)")  "* Add extra result to integrator"
    write (u, "(A)")

    select type (mci)
    type is (mci_vamp_t)
       mci%current_pass%calls(2) = 1234
    end select

    write (u, "(A)")  "* Update integrator (should fail)"
    write (u, "(A)")

    select type (mci)
    type is (mci_vamp_t)
       call mci%update_from_ref (mci_ref, success)
    end select

    write (u, "(1x,A,L1)")  "success = ", success
    write (u, "(A)")
    call mci%write (u)
    
    write (u, "(A)")
    write (u, "(A)")  "* Cleanup"

    call mci%final ()
    call mci_ref%final ()

    write (u, "(A)")
    write (u, "(A)")  "* Test output end: mci_vamp_13"

  end subroutine mci_vamp_13

\eatline
\nwendcode{}\nwbegindocs{826}\nwdocspar
\subsubsection{Accuracy Goal}
Integrate with multiple iterations.  Skip iterations once an accuracy goal has
been reached.
\nwenddocs{}\nwbegincode{827}\moddef{MCI vamp: execute tests}\plusendmoddef
  call test (mci_vamp_14, "mci_vamp_14", &
       "accuracy goal", &
       u, results)
\nwendcode{}\nwbegincode{828}\moddef{MCI vamp: tests}\plusendmoddef
  subroutine mci_vamp_14 (u)
    integer, intent(in) :: u
    type(grid_parameters_t) :: grid_par
    class(mci_t), allocatable, target :: mci
    class(mci_instance_t), pointer :: mci_instance => null ()
    class(mci_sampler_t), allocatable :: sampler
    class(rng_t), allocatable :: rng
    
    write (u, "(A)")  "* Test output: mci_vamp_14"
    write (u, "(A)")  "*   Purpose: integrate function in one dimension &
         &(single channel)" 
    write (u, "(A)")  "*            and check accuracy goal"
    
    write (u, "(A)")
    write (u, "(A)")  "* Initialize integrator, sampler, instance"
    write (u, "(A)")

    allocate (mci_vamp_t :: mci)
    call mci%set_dimensions (1, 1)
    select type (mci)
    type is (mci_vamp_t)
       grid_par%use_vamp_equivalences = .false.
       grid_par%accuracy_goal = 5E-2_default
       call mci%set_grid_parameters (grid_par)
    end select
    
    allocate (rng_tao_t :: rng)
    call rng%init ()
    call mci%import_rng (rng)

    call mci%allocate_instance (mci_instance)
    call mci_instance%init (mci)
    
    allocate (test_sampler_1_t :: sampler)
    select type (sampler)
    type is (test_sampler_1_t)
       sampler%mode = 2
    end select
    call sampler%write (u)

    write (u, "(A)")
    write (u, "(A)")  "* Integrate with n_it = 5 and n_calls = 100"
    write (u, "(A)")
    
    select type (mci)
    type is (mci_vamp_t)
       call mci%add_pass (adapt_grids = .true.)
    end select
    call mci%integrate (mci_instance, sampler, 5, 100)
    call mci%write (u)
 
    write (u, "(A)")
    write (u, "(A)")  "* Cleanup"

    call mci_instance%final ()
    call mci%final ()

    write (u, "(A)")
    write (u, "(A)")  "* Test output end: mci_vamp_14"

  end subroutine mci_vamp_14

\eatline
\nwendcode{}\nwbegindocs{829}\nwdocspar
\subsubsection{VAMP history}
Integrate with three passes and different settings for weight and grid
adaptation.  Then show the VAMP history.
\nwenddocs{}\nwbegincode{830}\moddef{MCI vamp: execute tests}\plusendmoddef
  call test (mci_vamp_15, "mci_vamp_15", &
       "VAMP history", &
       u, results)
\nwendcode{}\nwbegincode{831}\moddef{MCI vamp: tests}\plusendmoddef
  subroutine mci_vamp_15 (u)
    integer, intent(in) :: u
    type(grid_parameters_t) :: grid_par
    type(history_parameters_t) :: history_par
    class(mci_t), allocatable, target :: mci
    class(mci_instance_t), pointer :: mci_instance => null ()
    class(mci_sampler_t), allocatable :: sampler
    class(rng_t), allocatable :: rng
    
    write (u, "(A)")  "* Test output: mci_vamp_15"
    write (u, "(A)")  "*   Purpose: integrate function in one dimension &
         &(two channels)" 
    write (u, "(A)")  "*            in three passes, show history"
    
    write (u, "(A)")
    write (u, "(A)")  "* Initialize integrator, sampler, instance"
    write (u, "(A)")

    history_par%channel = .true.

    allocate (mci_vamp_t :: mci)
    call mci%set_dimensions (1, 2)
    select type (mci)
    type is (mci_vamp_t)
       grid_par%stratified = .false.
       grid_par%use_vamp_equivalences = .false.
       call mci%set_grid_parameters (grid_par)
       call mci%set_history_parameters (history_par)
    end select
    
    allocate (rng_tao_t :: rng)
    call rng%init ()
    call mci%import_rng (rng)

    call mci%allocate_instance (mci_instance)
    call mci_instance%init (mci)
    
    allocate (test_sampler_3_t :: sampler)
    select type (sampler)
    type is (test_sampler_3_t)
       sampler%a = 0.9_default
       sampler%b = 0.1_default
    end select
    call sampler%write (u)

    write (u, "(A)")
    write (u, "(A)")  "* Pass 1: grid and weight adaptation"
    
    select type (mci)
    type is (mci_vamp_t)
       call mci%add_pass (adapt_grids = .true., adapt_weights = .true.)
    end select
    call mci%integrate (mci_instance, sampler, 3, 1000)
 
    write (u, "(A)")
    write (u, "(A)")  "* Pass 2: grid adaptation"
    
    select type (mci)
    type is (mci_vamp_t)
       call mci%add_pass (adapt_grids = .true.)
    end select
    call mci%integrate (mci_instance, sampler, 3, 1000)
 
    write (u, "(A)")
    write (u, "(A)")  "* Pass 3: without adaptation"
    
    select type (mci)
    type is (mci_vamp_t)
       call mci%add_pass ()
    end select
    call mci%integrate (mci_instance, sampler, 3, 1000)

    write (u, "(A)")
    write (u, "(A)")  "* Contents of MCI record, with history"
    write (u, "(A)")
    
    call mci%write (u)
    select type (mci)
    type is (mci_vamp_t)
       call mci%write_history (u)
    end select
 
    write (u, "(A)")
    write (u, "(A)")  "* Cleanup"

    call mci_instance%final ()
    call mci%final ()

    write (u, "(A)")
    write (u, "(A)")  "* Test output end: mci_vamp_15"

  end subroutine mci_vamp_15

\eatline
\nwendcode{}\nwfilename{../../../src/shower/shower.nw}\nwbegindocs{0}% -*- ess-noweb-default-code-mode: f90-mode; noweb-default-code-mode: f90-mode; -*-
% WHIZARD shower code as NOWEB source
%%%%%%%%%%%%%%%%%%%%%%%%%%%%%%%%%%%%%%%%%%%%%%%%%%%%%%%%%%%%%%%%%%%%%%%%
\chapter{Parton shower and interface to hadronization}

This is the code for the \whizard\ QCD parton shower for final state
radiation (FSR) and initial state radiation (ISR) as well as the interface
to the \pythia\ module for showering and hadronization.

\section{Basics of the shower}
\nwenddocs{}\nwbegincode{1}\moddef{\code{}shower{\char95}base.f90\edoc{}}\endmoddef
\LA{}File header\RA{}

module shower_base

\LA{}Use kinds\RA{}
  use constants
  use diagnostics
  use rng_base
  use physics_defs
  use sm_physics, only: running_as_lam

\LA{}Standard module head\RA{}

\LA{}Shower base: public\RA{}

\LA{}Shower base: parameters\RA{}

contains

\LA{}Shower base: procedures\RA{}

end module shower_base
\eatline
\nwendcode{}\nwbegindocs{2}\nwdocspar
\subsection{PDF Parameters}
The shower has to know about the PDF set used for the parent process,
if any.  We support the following:
\nwenddocs{}\nwbegincode{3}\moddef{Shower base: parameters}\endmoddef
  integer, parameter, public :: STRF_NONE = 0
  integer, parameter, public :: STRF_LHAPDF6 = 1
  integer, parameter, public :: STRF_LHAPDF5 = 2
  integer, parameter, public :: STRF_PDF_BUILTIN = 3

\eatline
\nwendcode{}\nwbegindocs{4}\nwdocspar
\subsection{Additional parameters}
{\tt{}D{\char95}PRINT} decides whether to print out additional information while
{\tt{}ASSERT} does additional sanity checks that should be disabled for
production.
\nwenddocs{}\nwbegincode{5}\moddef{Shower base: parameters}\plusendmoddef
  logical, parameter, public :: D_PRINT = .false.
  logical, parameter, public :: ASSERT = .false.
\eatline
\nwendcode{}\nwbegindocs{6}\nwdocspar
These parameters are the cut-off scale $t_{\text{cut}}$, given in
GeV$^2$ ({\tt{}PARJ(82)}), the cut-off scale for the $P_t^2$-ordered
shower in GeV$^2$, and the two shower parameters {\tt{}PARP(72)} and
{\tt{}PARP(61)}, respectively.
\nwenddocs{}\nwbegincode{7}\moddef{Shower base: parameters}\plusendmoddef
  real(default), public :: D_Min_t = one
  real(default), public :: D_min_scale = 0.5_default
  real(default), public :: D_Lambda_fsr = 0.29_default
  real(default), public :: D_Lambda_isr = 0.29_default
\eatline
\eatline
\nwendcode{}\nwbegindocs{8}\nwdocspar
The parameter {\tt{}MSTJ(45)} gives the maximum number of flavours in
gluon decay to quarks:
\nwenddocs{}\nwbegincode{9}\moddef{Shower base: parameters}\plusendmoddef
  integer, public :: D_Nf = 5
\nwendcode{}\nwbegindocs{10} %def D_Nf
\nwenddocs{}\nwbegindocs{11}\nwdocspar
These two parameters decide whether to use constant or running
$alpha_s$ in the form of the function $D\_{alpha_s} (t)$ for the FSR
and ISR ({\tt{}MSTJ(44)}, {\tt{}MSTP(64)}), respectively. The next is the
parameter {\tt{}PARU(111)}, which sets the value for constant $\alpha_s$,
and the flag whether to use $P_t$-ordered ISR.
\nwenddocs{}\nwbegincode{12}\moddef{Shower base: parameters}\plusendmoddef
  logical, public :: D_running_alpha_s_fsr = .true.
  logical, public :: D_running_alpha_s_isr = .true.
  real(default), public :: D_constalpha_s = 0.20_default
  logical, public :: alpha_s_fudged = .true.
  logical, public :: isr_pt_ordered = .false.
\eatline
\eatline
\eatline
\eatline
\nwendcode{}\nwbegindocs{13}Treating either $u$ and $d$, or all quarks except $t$ as massless:
\nwenddocs{}\nwbegincode{14}\moddef{Shower base: parameters}\plusendmoddef
  logical, public :: treat_light_quarks_massless = .true.
  logical, public :: treat_duscb_quarks_massless = .false.
\eatline
\eatline
\nwendcode{}\nwbegindocs{15}\nwdocspar
The width and the cutoff of the Gaussian primordial $k_t$
distribution, {\tt{}PARP(91)} and {\tt{}PARP(93)}, in GeV:
\nwenddocs{}\nwbegincode{16}\moddef{Shower base: parameters}\plusendmoddef
  real(default), public :: primordial_kt_width = 1.5_default
  real(default), public :: primordial_kt_cutoff = five
\eatline
\eatline
\nwendcode{}\nwbegindocs{17}Temporary parameters for the $P_t$-ordered shower:
\nwenddocs{}\nwbegincode{18}\moddef{Shower base: parameters}\plusendmoddef
  real(default), public :: scalefactor1 = 0.02_default
  real(default), public :: scalefactor2 = 0.02_default
\eatline
\nwendcode{}\nwbegindocs{19}\nwdocspar
\nwenddocs{}\nwbegincode{20}\moddef{Shower base: parameters}\plusendmoddef
    class(rng_t), allocatable, public :: rng
\eatline
\nwendcode{}\nwbegincode{21}\moddef{Shower base: public}\endmoddef
  public :: D_alpha_s_isr
  public :: D_alpha_s_fsr
\nwendcode{}\nwbegincode{22}\moddef{Shower base: procedures}\endmoddef
  function D_alpha_s_isr (tin) result (alpha_s)
    real(default), intent(in) :: tin
    real(default) :: t
    real(default) :: alpha_s
    if (alpha_s_fudged) then
       t = max (max (0.1_default * D_Min_t, &
                     1.1_default * D_Lambda_isr**2), abs(tin))
    else
       t = abs(tin)
    end if
    if (D_running_alpha_s_isr) then
       alpha_s = running_as_lam (number_of_flavors(t), sqrt(t), D_Lambda_isr, 0)
    else
       alpha_s = D_constalpha_s
    end if
  end function D_alpha_s_isr

  function D_alpha_s_fsr (tin) result (alpha_s)
    real(default), intent(in) :: tin
    real(default) :: t
    real(default) :: alpha_s
    if (alpha_s_fudged) then
       t = max (max (0.1_default * D_Min_t, &
                     1.1_default * D_Lambda_isr**2), abs(tin))
    else
       t = abs(tin)
    end if
    if (D_running_alpha_s_fsr) then
       alpha_s = running_as_lam (number_of_flavors(t), sqrt(t), D_Lambda_fsr, 0)
    else
       alpha_s = D_constalpha_s
    end if
  end function D_alpha_s_fsr

\eatline
\nwendcode{}\nwbegindocs{23}Mass and mass squared selection functions. All masses are in
GeV. Light quarks are assumed to be ordered, $m_1 < m_2 < m_3 \ldots$,
and they get current masses, not elementary ones. Mesons and baryons
other than proton and neutron are needed as beam-remnants. Particles
with PDG number zero are taken massless, as well as proper beam
remnants and any other particles.
\nwenddocs{}\nwbegincode{24}\moddef{Shower base: public}\plusendmoddef
  public :: mass_type
  public :: mass_squared_type
\nwendcode{}\nwbegincode{25}\moddef{Shower base: procedures}\plusendmoddef
  elemental function mass_type (type) result (mass)
    integer, intent(in) :: type
    real(default) :: mass
    mass = sqrt (mass_squared_type (type))
  end function mass_type

  elemental function mass_squared_type (type) result (mass2)
    integer, intent(in) :: type
    real(default) :: mass2

    select case (abs (type))
    case (1,2)
       if (treat_light_quarks_massless .or. &
            treat_duscb_quarks_massless) then
          mass2 = zero
       else
          mass2 = 0.330_default**2
       end if
    case (3)
       if (treat_duscb_quarks_massless) then
          mass2 = zero
       else
          mass2 = 0.500_default**2
       end if
    case (4)
       if (treat_duscb_quarks_massless) then
          mass2 = zero
       else
          mass2 = 1.500_default**2
       end if
    case (5)
       if (treat_duscb_quarks_massless) then
          mass2 = zero
       else
          mass2 = 4.800_default**2
       end if
    case (6)
       mass2 = 175.00_default**2
    case (GLUON)
       mass2 = zero
    case (NEUTRON)
       mass2 = 0.939565_default**2
    case (PROTON)
       mass2 = 0.93827_default**2
    case (DPLUS)
       mass2 = 1.86960_default**2
    case (D0)
       mass2 = 1.86483_default**2
    case (B0)
       mass2 = 5.27950_default**2
    case (BPLUS)
       mass2 = 5.27917_default**2
    case (DELTAPLUSPLUS)
       mass2 = 1.232_default**2
    case (SIGMA0)
       mass2 = 1.192642_default**2
    case (SIGMAPLUS)
       mass2 = 1.18937_default**2
    case (SIGMACPLUS)
       mass2 = 2.4529_default**2
    case (SIGMACPLUSPLUS)
       mass2 = 2.45402_default**2
    case (SIGMAB0)
       mass2 = 5.8152_default**2
    case (SIGMABPLUS)
       mass2 = 5.8078_default**2
    case (UNDEFINED)
       mass2 = zero
    case (BEAM_REMNANT)
       mass2 = zero ! don't know how to handle the beamremnant
    case default !others not implemented
       mass2 = zero
    end select
  end function mass_squared_type

\eatline
\nwendcode{}\nwbegindocs{26}The number of flavours active at a certain scale (virtuality) $t$.
\nwenddocs{}\nwbegincode{27}\moddef{Shower base: public}\plusendmoddef
  public :: number_of_flavors
\nwendcode{}\nwbegincode{28}\moddef{Shower base: procedures}\plusendmoddef
  elemental function number_of_flavors (t) result (nr)
    real(default), intent(in) :: t
    real(default) :: nr
    integer :: i
    nr = 0
    if (t < D_Min_t) return   ! arbitrary cut off
    ! TODO: do i = 1, min (max (3, D_Nf), 6)
    do i = 1, min (3, D_Nf)
    !!! to do: take heavier quarks(-> cuts on allowed costheta in g->qq)
    !!!        into account
       if ((four * mass_squared_type (i) + D_Min_t) < t ) then
          nr = i
       else
          exit
       end if
    end do
  end function number_of_flavors

\eatline
\nwendcode{}\nwbegindocs{29}Methods to set parameters, once there should be a better interface.
\nwenddocs{}\nwbegincode{30}\moddef{Shower base: public}\plusendmoddef
  public :: shower_set_D_min_t
  public :: shower_set_D_lambda_fsr
  public :: shower_set_D_lambda_isr
  public :: shower_set_D_Nf
  public :: shower_set_D_running_alpha_s_fsr
  public :: shower_set_D_running_alpha_s_isr
  public :: shower_set_D_constantalpha_s
  public :: shower_set_isr_pt_ordered
  public :: shower_set_primordial_kt_width
  public :: shower_set_primordial_kt_cutoff
  public :: shower_set_rng
  public :: shower_get_rng
\nwendcode{}\nwbegincode{31}\moddef{Shower base: procedures}\plusendmoddef
  subroutine shower_set_D_Min_t (input)
    real(default) :: input
    D_Min_t = input
  end subroutine shower_set_D_Min_t

  subroutine shower_set_D_Lambda_fsr (input)
    real(default) :: input
    D_Lambda_fsr = input
  end subroutine shower_set_D_Lambda_fsr

  subroutine shower_set_D_Lambda_isr (input)
    real(default) :: input
    D_Lambda_isr = input
  end subroutine shower_set_D_Lambda_isr

  subroutine shower_set_D_Nf (input)
    integer :: input
    D_Nf = input
  end subroutine shower_set_D_Nf

  subroutine shower_set_D_running_alpha_s_fsr (input)
    logical :: input
    D_running_alpha_s_fsr = input
  end subroutine shower_set_D_running_alpha_s_fsr

  subroutine shower_set_D_running_alpha_s_isr (input)
    logical :: input
    D_running_alpha_s_isr = input
  end subroutine shower_set_D_running_alpha_s_isr

  subroutine shower_set_D_constantalpha_s (input)
    real(default) :: input
    D_constalpha_s = input
  end subroutine shower_set_D_constantalpha_s

  subroutine shower_set_isr_pt_ordered (input)
    logical :: input
    isr_pt_ordered = input
  end subroutine shower_set_isr_pt_ordered

  subroutine shower_set_primordial_kt_width (input)
    real(default) :: input
    primordial_kt_width = input
  end subroutine shower_set_primordial_kt_width

  subroutine shower_set_primordial_kt_cutoff (input)
    real(default) :: input
    primordial_kt_cutoff = input
  end subroutine shower_set_primordial_kt_cutoff

  subroutine shower_set_rng (rng_inc)
    class(rng_t), intent(inout), allocatable :: rng_inc
    call move_alloc (from = rng_inc, to = rng)
  end subroutine shower_set_rng

  subroutine shower_get_rng (rng_out)
    class(rng_t), intent(inout), allocatable :: rng_out
    call move_alloc (from = rng, to = rng_out)
  end subroutine shower_get_rng

\eatline
\eatline
\eatline
\eatline
\eatline
\eatline
\eatline
\eatline
\eatline
\nwendcode{}\nwbegindocs{32}\nwdocspar
%%%%%%%%%%%%%%%%%%%%%%%%%%%%%%%%%%%%%%%%%%%%%%%%%%%%%%%%%%%%%%%%%%%%%%%%
\section{Parton module for the shower}
\nwenddocs{}\nwbegincode{33}\moddef{\code{}shower{\char95}partons.f90\edoc{}}\endmoddef
\LA{}File header\RA{}

module shower_partons

\LA{}Use kinds\RA{}
  use io_units
  use constants
  use system_defs, only: TAB
  use diagnostics
  use physics_defs
  use lorentz
  use sm_physics
  use shower_base

\LA{}Standard module head\RA{}

\LA{}Shower partons: public\RA{}

\LA{}Shower partons: types\RA{}

contains

\LA{}Shower partons: procedures\RA{}

end module shower_partons
\eatline
\nwendcode{}\nwbegindocs{34}\nwdocspar
\subsection{The basic type defintions}

The type {\tt{}parton{\char95}t} defines a parton for the shower. The {\tt{}x} value
of the parton is only needed for spacelike showers. The pointer
{\tt{}initial} is only needed for partons in initial showers, it points to
the hadron the parton is coming from. An auxiliary value for the
$P_t$-ordered ISR is {\tt{}aux{\char95}pt}. Then, there are two auxiliary entries
for the clustering of CKKW pseudo weights and CKKW matching,
{\tt{}ckkwlabel} and {\tt{}ckkwscale}.
\nwenddocs{}\nwbegincode{35}\moddef{Shower partons: public}\endmoddef
  public :: parton_t
\nwendcode{}\nwbegincode{36}\moddef{Shower partons: types}\endmoddef
  type :: parton_t
     integer :: nr = 0
     integer :: type = 0
     type(vector4_t) :: momentum = vector4_null
     real(default) :: t  = zero
     real(default) :: scale = zero
     real(default) :: z = zero
     real(default) :: costheta = zero
     real(default) :: x = zero
     logical :: simulated = .false.
     logical :: belongstoFSR = .true.
     logical :: belongstointeraction = .false.
     type(parton_t), pointer :: parent => null ()
     type(parton_t), pointer :: child1 => null ()
     type(parton_t), pointer :: child2 => null ()
     type(parton_t), pointer :: initial => null ()
     integer :: c1 = 0, c2 = 0
     integer :: aux_pt = 0
     integer :: ckkwlabel = 0
     real(default) :: ckkwscale = zero
     integer :: ckkwtype = -1
     integer :: interactionnr = 0
  end type parton_t

\eatline
\nwendcode{}\nwbegindocs{37}\nwdocspar
\nwenddocs{}\nwbegincode{38}\moddef{Shower partons: public}\plusendmoddef
  public :: parton_pointer_t
\nwendcode{}\nwbegincode{39}\moddef{Shower partons: types}\plusendmoddef
  type :: parton_pointer_t
     type(parton_t), pointer :: p => null ()
  end type parton_pointer_t

\eatline
\nwendcode{}\nwbegindocs{40}\nwdocspar
\subsection{Routines}
\nwenddocs{}\nwbegincode{41}\moddef{Shower partons: public}\plusendmoddef
  public :: parton_copy
\nwendcode{}\nwbegincode{42}\moddef{Shower partons: procedures}\endmoddef
  subroutine parton_copy (prt1, prt2)
    type(parton_t), intent(in) :: prt1
    type(parton_t), intent(out) :: prt2
    prt2%nr = prt1%nr
    prt2%type = prt1%type
    prt2%momentum = prt1%momentum
    prt2%t = prt1%t
    prt2%scale = prt1%scale
    prt2%z = prt1%z
    prt2%costheta = prt1%costheta
    prt2%x = prt1%x
    prt2%simulated = prt1%simulated
    prt2%belongstoFSR = prt1%belongstoFSR
    prt2%belongstointeraction = prt1%belongstointeraction
    prt2%interactionnr = prt1%interactionnr
    if (associated (prt1%parent))  prt2%parent  => prt1%parent
    if (associated (prt1%child1))  prt2%child1  => prt1%child1
    if (associated (prt1%child2))  prt2%child2  => prt1%child2
    if (associated (prt1%initial)) prt2%initial => prt1%initial
    prt2%c1 = prt1%c1
    prt2%c2 = prt1%c2
    prt2%aux_pt = prt1%aux_pt
  end subroutine parton_copy

\nwendcode{}\nwbegindocs{43} %def parton_copy
\nwenddocs{}\nwbegindocs{44}This returns the angle between the daughters assuming them to be
massless.
\nwenddocs{}\nwbegincode{45}\moddef{Shower partons: public}\plusendmoddef
  public :: parton_get_costheta
\nwendcode{}\nwbegincode{46}\moddef{Shower partons: procedures}\plusendmoddef
  elemental function parton_get_costheta (prt) result (costheta)
    type(parton_t), intent(in) :: prt
    real(default) :: costheta
    real(default) :: denom
    denom = two * prt%z * (one - prt%z) * prt%momentum%p(0)**2
    if (denom > eps0) then
       costheta = one - prt%t / denom
    else
       costheta = - one
    end if
  end function parton_get_costheta

\eatline
\nwendcode{}\nwbegindocs{47}The same for massive daughters.
\nwenddocs{}\nwbegincode{48}\moddef{Shower partons: public}\plusendmoddef
  public :: parton_get_costheta_correct
\nwendcode{}\nwbegindocs{49}\nwdocspar
\nwenddocs{}\nwbegincode{50}\moddef{Shower partons: procedures}\plusendmoddef
  elemental function parton_get_costheta_correct (prt) result (costheta)
    type(parton_t), intent(in) :: prt
    real(default) :: costheta, sqrt12
    if (parton_is_branched (prt)) then
       if (prt%child1%simulated .and. &
           prt%child2%simulated) then
          sqrt12 = sqrt (max (zero, (prt%z)**2 * prt%momentum%p(0)**2 &
                                    - prt%child1%t)) * &
                   sqrt (max (zero, (one - prt%z)**2 * prt%momentum%p(0)**2 &
                                    - prt%child2%t))
          if (sqrt12 > eps0) then
             costheta = (prt%t - prt%child1%t - prt%child2%t - &
                  two * prt%z * (one - prt%z) * prt%momentum%p(0)**2) / &
                  (- two * sqrt12)
             return
          end if
       end if
    end if
    costheta = parton_get_costheta (prt)
  end function parton_get_costheta_correct

\eatline
\nwendcode{}\nwbegindocs{51}This function returns the angle between the momentum vectors of the
parton and first daughter.  This is only used for debugging.
\nwenddocs{}\nwbegincode{52}\moddef{Shower partons: public}\plusendmoddef
  public :: parton_get_costheta_motherfirst
\nwendcode{}\nwbegincode{53}\moddef{Shower partons: procedures}\plusendmoddef
  elemental function parton_get_costheta_motherfirst (prt) result (costheta)
    real(default) :: costheta
    type(parton_t), intent(in) :: prt
    if (parton_is_branched (prt)) then
       if ((prt%child1%simulated .or. &
            parton_is_final (prt%child1) .or. &
            parton_is_branched (prt%child1)) .and. &
            (prt%child2%simulated .or. &
             parton_is_final (prt%child2) .or. &
             parton_is_branched (prt%child2))) then
          costheta = enclosed_angle_ct (prt%momentum, prt%child1%momentum)
          return
       end if
    end if
    costheta = - two
  end function parton_get_costheta_motherfirst

\eatline
\nwendcode{}\nwbegindocs{54}Get the parton velocities.
\nwenddocs{}\nwbegincode{55}\moddef{Shower partons: public}\plusendmoddef
  public :: parton_get_beta
\nwendcode{}\nwbegindocs{56}\nwdocspar
\nwenddocs{}\nwbegincode{57}\moddef{Shower partons: procedures}\plusendmoddef
  pure function get_beta (t,E) result (beta)
    real(default), intent(in) :: t,E
    real(default) :: beta
    beta = sqrt (max (tiny_07, one - t /(E**2)))
  end function get_beta

  elemental function parton_get_beta (prt) result (beta)
    type(parton_t), intent(in) :: prt
    real(default) :: beta
    beta = sqrt (max (tiny_07, one - prt%t / prt%momentum%p(0)**2))
  end function parton_get_beta

\eatline
\nwendcode{}\nwbegindocs{58}Write routine.
\nwenddocs{}\nwbegincode{59}\moddef{Shower partons: public}\plusendmoddef
  public :: parton_write
\nwendcode{}\nwbegincode{60}\moddef{Shower partons: procedures}\plusendmoddef
  subroutine parton_write (prt, unit)
    type(parton_t), intent(in) :: prt
    integer, intent(in), optional :: unit
    integer :: u
    u = given_output_unit (unit); if (u < 0) return

    write (u, "(1x,7A)") "Shower parton <nr>", TAB, "<type>", TAB // TAB, &
         "<parent>", TAB, "<mom(0:3)>"
    write (u, "(2x,I5,3A)", advance = "no")  prt%nr, TAB, TAB, TAB
    if (parton_is_final (prt)) then
       write (u, "(1x,I5,1x,A)", advance = "no") prt%type, TAB // TAB
    else
       write (u, "('[',I5,']',A)", advance = "no") prt%type, TAB // TAB
    end if
    if (associated (prt%parent)) then
       write (u, "(I5,A)", advance = "no") prt%parent%nr, TAB // TAB
    else
       write (u, "(5x,2A)", advance = "no") TAB, TAB
    end if
    write (u, "(4(ES12.5,A))") prt%momentum%p(0), TAB, &
                               prt%momentum%p(1), TAB, &
                               prt%momentum%p(2), TAB, &
                               prt%momentum%p(3)
    write (u, "(1x,9A)") "<p4square>", TAB // TAB, "<t>", TAB // TAB, &
         "<scale>", TAB // TAB, "<c1>", TAB, "<c2>"
    write (u, "(1x,3(ES12.5,A))", advance = "no") &
         prt%momentum ** 2, TAB // TAB, prt%t, TAB, prt%scale, TAB
    write (u, "(2(I4,A))") prt%c1, TAB, prt%c2, TAB
    if (parton_is_branched (prt)) then
       if (prt%belongstoFSR) then
          write (u, "(1x,9A)") "costheta(prt)", TAB, &
               "costheta_correct(prt)", TAB, &
               "prt%costheta", TAB, "prt%z", TAB, &
               "costheta_motherfirst(prt)"
          write (u, "(1X,5(ES12.5,A))") &
               parton_get_costheta (prt), TAB, &
               parton_get_costheta_correct (prt), TAB // TAB, &
               prt%costheta, TAB, prt%z, TAB, &
               parton_get_costheta_motherfirst (prt), TAB
       else
          write (u, "(1x,9A)") "prt%z", TAB, "prt%x", TAB, &
               "costheta_correct(prt)", TAB, &
               "prt%costheta", TAB, &
               "costheta_motherfirst(prt)"
          write (u, "(1X,5(ES12.5,A))") &
               prt%z, TAB, prt%x, TAB, &
               parton_get_costheta_correct (prt), TAB, &
               prt%costheta, TAB, &
               parton_get_costheta_motherfirst (prt), TAB
       end if
    else
       if (prt%belongstoFSR) then
          write (u, "(1X,A)") "not branched."
       else
          write (u, "(1X,A,ES12.5)") "not branched. x = ",  prt%x
       end if
    end if
    write (u, "(A)", advance = "no") " Parton"
    if (prt%belongstoFSR) then
       write (u, "(A)", advance = "no")  " is FSR,"
    else
       if (associated (prt%initial)) then
          write (u, "(A,I1)", advance = "no")  " from hadron,", prt%initial%nr
       else
          write (u, "(A)", advance = "no")  ""
       end if
    end if
    if (parton_is_final (prt)) then
       write (u, "(A)", advance = "no")  " is final,"
    else
       write (u, "(A)", advance = "no")  ""
    end if
    if (prt%simulated) then
       write (u, "(A)", advance = "no")  " is simulated,"
    else
       write (u, "(A)", advance = "no")  ""
    end if
    if (associated (prt%child1) .and. associated (prt%child2)) then
       write (u, "(A,2(I5),A)", advance = "no") &
            " has children: ", prt%child1%nr, prt%child2%nr, ","
    else if (associated (prt%child1)) then
       write (u, "(A,1(I5),A)", advance = "no") &
            " has one child: ", prt%child1%nr, ", "
    end if
    if (prt%belongstointeraction) then
       write (u, "(A,I2)") " belongs to interaction ", &
            prt%interactionnr
    else
       write (u, "(A,I2)") " does not belong to interaction ", &
            prt%interactionnr
    end if
    write (u,"(A)")  TAB
  end subroutine parton_write

\eatline
\nwendcode{}\nwbegindocs{61}\nwdocspar
\nwenddocs{}\nwbegincode{62}\moddef{Shower partons: public}\plusendmoddef
  public :: parton_is_final
\nwendcode{}\nwbegincode{63}\moddef{Shower partons: procedures}\plusendmoddef
  elemental function parton_is_final (prt) result (is_final)
    type(parton_t), intent(in) :: prt
    logical :: is_final
    is_final = .false.
    if (prt%belongstoFSR) then
       is_final = .not. associated (prt%child1) .and. &
            (.not. prt%belongstointeraction .or. &
            (prt%belongstointeraction .and. prt%simulated))
    end if
  end function parton_is_final

\eatline
\nwendcode{}\nwbegindocs{64}\nwdocspar
\nwenddocs{}\nwbegincode{65}\moddef{Shower partons: public}\plusendmoddef
  public :: parton_is_branched
\nwendcode{}\nwbegincode{66}\moddef{Shower partons: procedures}\plusendmoddef
  elemental function parton_is_branched (prt) result (is_branched)
    type(parton_t), intent(in) :: prt
    logical :: is_branched
    is_branched = associated (prt%child1) .and. associated (prt%child2)
  end function parton_is_branched

\eatline
\nwendcode{}\nwbegindocs{67}\nwdocspar
\nwenddocs{}\nwbegincode{68}\moddef{Shower partons: public}\plusendmoddef
  public :: parton_set_simulated
\nwendcode{}\nwbegincode{69}\moddef{Shower partons: procedures}\plusendmoddef
  pure subroutine parton_set_simulated (prt, sim)
    type(parton_t), intent(inout) :: prt
    logical, intent(in), optional :: sim
    if (present (sim)) then
       prt%simulated = sim
    else
       prt%simulated = .true.
    end if
  end subroutine parton_set_simulated

\eatline
\nwendcode{}\nwbegindocs{70}\nwdocspar
\nwenddocs{}\nwbegincode{71}\moddef{Shower partons: public}\plusendmoddef
  public :: parton_set_parent
\nwendcode{}\nwbegincode{72}\moddef{Shower partons: procedures}\plusendmoddef
  subroutine parton_set_parent (prt, parent)
    type(parton_t), intent(inout) :: prt
    type(parton_t), intent(in) , target :: parent
    prt%parent => parent
  end subroutine parton_set_parent

\eatline
\nwendcode{}\nwbegindocs{73}\nwdocspar
\nwenddocs{}\nwbegincode{74}\moddef{Shower partons: public}\plusendmoddef
  public :: parton_get_parent
\nwendcode{}\nwbegincode{75}\moddef{Shower partons: procedures}\plusendmoddef
  function parton_get_parent (prt) result (parent)
    type(parton_t), intent(in) :: prt
    type(parton_t), pointer :: parent
    parent => prt%parent
  end function parton_get_parent

\eatline
\nwendcode{}\nwbegindocs{76}\nwdocspar
\nwenddocs{}\nwbegincode{77}\moddef{Shower partons: public}\plusendmoddef
  public :: parton_set_initial
\nwendcode{}\nwbegincode{78}\moddef{Shower partons: procedures}\plusendmoddef
  subroutine parton_set_initial (prt, initial)
    type(parton_t), intent(inout) :: prt
    type(parton_t), intent(in) , target :: initial
    prt%initial => initial
  end subroutine parton_set_initial

\eatline
\nwendcode{}\nwbegindocs{79}\nwdocspar
\nwenddocs{}\nwbegincode{80}\moddef{Shower partons: public}\plusendmoddef
  public :: parton_get_initial
\nwendcode{}\nwbegincode{81}\moddef{Shower partons: procedures}\plusendmoddef
  function parton_get_initial (prt) result (initial)
    type(parton_t), intent(in) :: prt
    type(parton_t), pointer :: initial
    initial => prt%initial
  end function parton_get_initial

\eatline
\nwendcode{}\nwbegindocs{82}\nwdocspar
\nwenddocs{}\nwbegincode{83}\moddef{Shower partons: public}\plusendmoddef
  public :: parton_set_child
\nwendcode{}\nwbegincode{84}\moddef{Shower partons: procedures}\plusendmoddef
  subroutine parton_set_child (prt, child, i)
    type(parton_t), intent(inout) :: prt
    type(parton_t), intent(in), target :: child
    integer, intent(in) ::  i
    if (i == 1) then
       prt%child1 => child
    else
       prt%child2 => child
    end if
  end subroutine parton_set_child

\eatline
\nwendcode{}\nwbegindocs{85}\nwdocspar
\nwenddocs{}\nwbegincode{86}\moddef{Shower partons: public}\plusendmoddef
  public :: parton_get_child
\nwendcode{}\nwbegincode{87}\moddef{Shower partons: procedures}\plusendmoddef
  function parton_get_child (prt, i) result (child)
    type(parton_t), pointer :: child
    type(parton_t), intent(in) :: prt
    integer, intent(in) :: i
    child => null ()
    if (i == 1) then
       child => prt%child1
    else
       child => prt%child2
    end if
  end function parton_get_child

\eatline
\nwendcode{}\nwbegindocs{88}\nwdocspar
\nwenddocs{}\nwbegincode{89}\moddef{Shower partons: public}\plusendmoddef
  public :: parton_is_quark
\nwendcode{}\nwbegindocs{90}\nwdocspar
\nwenddocs{}\nwbegincode{91}\moddef{Shower partons: procedures}\plusendmoddef
  elemental function parton_is_quark (prt) result (is_quark)
    type(parton_t), intent(in) ::prt
    logical :: is_quark
    is_quark = abs (prt%type) <= 6 .and. prt%type /= 0
  end function parton_is_quark

\eatline
\nwendcode{}\nwbegindocs{92}\nwdocspar
\nwenddocs{}\nwbegincode{93}\moddef{Shower partons: public}\plusendmoddef
  public :: parton_is_squark
\nwendcode{}\nwbegincode{94}\moddef{Shower partons: procedures}\plusendmoddef
  elemental function parton_is_squark (prt) result (is_squark)
    type(parton_t), intent(in) ::prt
    logical :: is_squark
    is_squark = ((abs(prt%type) >= 1000001) .and. (abs(prt%type) <= 1000006)) &
             .or. ((abs(prt%type) >= 2000001) .and. (abs(prt%type) <= 2000006))
  end function parton_is_squark

\eatline
\nwendcode{}\nwbegindocs{95}\nwdocspar
\nwenddocs{}\nwbegincode{96}\moddef{Shower partons: public}\plusendmoddef
  public :: parton_is_gluon
\nwendcode{}\nwbegincode{97}\moddef{Shower partons: procedures}\plusendmoddef
  elemental function parton_is_gluon (prt) result (is_gluon)
    type(parton_t), intent(in) :: prt
    logical :: is_gluon
    is_gluon = prt%type == GLUON .or. prt%type == 9
  end function parton_is_gluon

\eatline
\nwendcode{}\nwbegincode{98}\moddef{Shower partons: public}\plusendmoddef
  public :: parton_is_gluino
\nwendcode{}\nwbegincode{99}\moddef{Shower partons: procedures}\plusendmoddef
  elemental function parton_is_gluino (prt) result (is_gluino)
    type(parton_t), intent(in) :: prt
    logical :: is_gluino
    is_gluino = prt%type == 1000021
  end function parton_is_gluino

\eatline
\nwendcode{}\nwbegindocs{100}Only the proton is implemented yet.
\nwenddocs{}\nwbegincode{101}\moddef{Shower partons: public}\plusendmoddef
  public :: parton_is_hadron
\nwendcode{}\nwbegincode{102}\moddef{Shower partons: procedures}\plusendmoddef
  elemental function parton_is_hadron (prt) result (is_hadron)
    type(parton_t), intent(in) :: prt
    logical :: is_hadron
    is_hadron = abs (prt%type) == PROTON
  end function parton_is_hadron

\eatline
\nwendcode{}\nwbegindocs{103}TODO: SUSY partons.
\nwenddocs{}\nwbegincode{104}\moddef{Shower partons: public}\plusendmoddef
  public :: parton_is_colored
\nwendcode{}\nwbegincode{105}\moddef{Shower partons: procedures}\plusendmoddef
  elemental function parton_is_colored (prt) result (is_colored)
    type(parton_t), intent(in) ::prt
    logical :: is_colored
    is_colored = parton_is_quark (prt) .or. parton_is_gluon (prt)
  end function parton_is_colored

\eatline
\nwendcode{}\nwbegindocs{106}\nwdocspar
\nwenddocs{}\nwbegincode{107}\moddef{Shower partons: public}\plusendmoddef
  public :: parton_mass
\nwendcode{}\nwbegincode{108}\moddef{Shower partons: procedures}\plusendmoddef
  function parton_mass (prt) result (mass)
    type(parton_t), intent(in) :: prt
    real(default) :: mass
    mass = mass_type (prt%type)
  end function parton_mass

\eatline
\nwendcode{}\nwbegindocs{109}\nwdocspar
\nwenddocs{}\nwbegincode{110}\moddef{Shower partons: public}\plusendmoddef
  public :: parton_mass_squared
\nwendcode{}\nwbegincode{111}\moddef{Shower partons: procedures}\plusendmoddef
  function parton_mass_squared (prt) result (mass_squared)
    type(parton_t), intent(in) :: prt
    real(default) :: mass_squared
    mass_squared = mass_squared_type (prt%type)
  end function parton_mass_squared

\eatline
\nwendcode{}\nwbegindocs{112}\nwdocspar
\nwenddocs{}\nwbegincode{113}\moddef{Shower partons: public}\plusendmoddef
  public :: P_prt_to_child1
\nwendcode{}\nwbegincode{114}\moddef{Shower partons: procedures}\plusendmoddef
  function P_prt_to_child1 (prt) result (retvalue)
    type(parton_t), intent(in) :: prt
    real(default) :: retvalue
    retvalue = zero
    if (parton_is_gluon (prt)) then
       if (parton_is_quark (prt%child1)) then
          retvalue = P_gqq (prt%z)
       else if (parton_is_gluon (prt%child1)) then
          retvalue = P_ggg (prt%z) + P_ggg (one - prt%z)
       end if
    else if (parton_is_quark (prt)) then
       if (parton_is_quark (prt%child1)) then
          retvalue = P_qqg (prt%z)
       else if (parton_is_gluon (prt%child1)) then
          retvalue = P_qqg (one - prt%z)
       end if
    end if
  end function P_prt_to_child1

\eatline
\nwendcode{}\nwbegindocs{115}This function returns whether the kinematics of the branching of
parton {\tt{}prt} into its daughters are allowed or not.
\nwenddocs{}\nwbegincode{116}\moddef{Shower partons: public}\plusendmoddef
  public :: thetabar
\nwendcode{}\nwbegincode{117}\moddef{Shower partons: procedures}\plusendmoddef
  function thetabar (prt, recoiler, isr_ang, E3out) result (retvalue)
    type(parton_t), intent(inout) :: prt
    type(parton_t), intent(in) :: recoiler
    real(default), intent(out), optional :: E3out
    logical, intent(in) :: isr_ang
    logical :: retvalue
    real(default) :: ctheta, cthetachild1
    real(default) p1, p4, p3, E3, shat

    shat = (prt%child1%momentum + recoiler%momentum)**2
    E3 = 0.5_default * (shat / prt%z -recoiler%t + prt%child1%t - &
                        parton_mass_squared (prt%child2)) / sqrt(shat)
    if (present (E3out)) then
       E3out = E3
    end if
    !!! absolute values of momenta in a 3 -> 1 + 4 branching
    p3 = sqrt (E3**2 - prt%t)
    p1 = sqrt (prt%child1%momentum%p(0)**2 - prt%child1%t)
    p4 = sqrt (max (zero, (E3 - prt%child1%momentum%p(0))**2 &
                          - prt%child2%t))
    if (p3 > zero) then
       retvalue = ((p1 + p4 >= p3) .and. (p3 >= abs(p1 - p4)) )
       if (retvalue .and. isr_ang) then
          !!! check angular ordering
          if (associated (prt%child1)) then
             if (associated (prt%child1%child2)) then
                ctheta = (E3**2 - p1**2 - p4**2 +prt%t) / (two * p1 * p4)
                cthetachild1 = (prt%child1%momentum%p(0)**2 - &
                     space_part (prt%child1%child1%momentum)**2 &
                     - space_part (prt%child1%child2%momentum)**2 + prt%child1%t) &
                     / (two * space_part (prt%child1%child1%momentum)**1 * &
                              space_part (prt%child1%child2%momentum)**1)
                retvalue = (ctheta > cthetachild1)
             end if
          end if
       end if
    else
       retvalue = .false.
    end if
  end function thetabar

\eatline
\nwendcode{}\nwbegindocs{118}\nwdocspar
\nwenddocs{}\nwbegincode{119}\moddef{Shower partons: public}\plusendmoddef
  public :: parton_apply_z
\nwendcode{}\nwbegincode{120}\moddef{Shower partons: procedures}\plusendmoddef
  recursive subroutine parton_apply_z(prt, newz)
    type(parton_t), intent(inout) :: prt
    real(default), intent(in) :: newz
    if (D_PRINT) print *, "D: old z = ", prt%z , " new z = ", newz
    prt%z = newz
    if (associated (prt%child1) .and. associated (prt%child2)) then
       prt%child1%momentum%p(0) = newz * prt%momentum%p(0)
       call parton_apply_z (prt%child1, prt%child1%z)
       prt%child2%momentum%p(0) = (one - newz) * prt%momentum%p(0)
       call parton_apply_z (prt%child2, prt%child2%z)
    end if
  end subroutine parton_apply_z

\eatline
\nwendcode{}\nwbegindocs{121}\nwdocspar
\nwenddocs{}\nwbegincode{122}\moddef{Shower partons: public}\plusendmoddef
  public :: parton_apply_costheta
\nwendcode{}\nwbegincode{123}\moddef{Shower partons: procedures}\plusendmoddef
  recursive subroutine parton_apply_costheta (prt)
    type(parton_t), intent(inout) :: prt
    prt%z = 0.5_default * (one + parton_get_beta (prt) * prt%costheta)
    if (associated (prt%child1) .and. associated (prt%child2) ) then
       if (prt%child1%simulated .and. &
           prt%child2%simulated) then
          prt%z = 0.5_default * (one + (prt%child1%t - prt%child2%t) / &
               prt%t + parton_get_beta (prt) * prt%costheta * &
                sqrt((prt%t - prt%child1%t - prt%child2%t)**2 - &
                4 * prt%child1%t * prt%child2%t) / prt%t)
          if (prt%type /= INTERNAL) then
             prt%child1%momentum%p(0) = prt%z * prt%momentum%p(0)
             prt%child2%momentum%p(0) = (one - prt%z) * prt%momentum%p(0)
          end if
          call parton_generate_ps (prt)
          call parton_apply_costheta (prt%child1)
          call parton_apply_costheta (prt%child2)
       end if
    end if
  end subroutine parton_apply_costheta

\eatline
\nwendcode{}\nwbegindocs{124}\nwdocspar
\nwenddocs{}\nwbegincode{125}\moddef{Shower partons: public}\plusendmoddef
  public :: parton_apply_lorentztrafo
\nwendcode{}\nwbegincode{126}\moddef{Shower partons: procedures}\plusendmoddef
  subroutine parton_apply_lorentztrafo (prt, L)
    type(parton_t), intent(inout) :: prt
    type(lorentz_transformation_t), intent(in) :: L
    prt%momentum = L * prt%momentum
  end subroutine parton_apply_lorentztrafo

\eatline
\nwendcode{}\nwbegindocs{127}\nwdocspar
\nwenddocs{}\nwbegincode{128}\moddef{Shower partons: public}\plusendmoddef
  public :: parton_apply_lorentztrafo_recursive
\nwendcode{}\nwbegincode{129}\moddef{Shower partons: procedures}\plusendmoddef
  recursive subroutine parton_apply_lorentztrafo_recursive (prt, L)
    type(parton_t), intent(inout) :: prt
    type(lorentz_transformation_t) ,intent(in) :: L
    if (prt%type /= PROTON .and. prt%type /= BEAM_REMNANT) then
       !!! don't boost hadrons and beam-remnants
       call parton_apply_lorentztrafo (prt, L)
    end if
    if (associated (prt%child1) .and. associated (prt%child2)) then
       if ((space_part_norm (prt%child1%momentum) < eps0) .and. &
           (space_part_norm (prt%child2%momentum) < eps0) .and. &
           (.not. prt%child1%belongstointeraction) .and. &
           (.not. prt%child2%belongstointeraction)) then
          !!! don't boost unevolved timelike partons
       else
          call parton_apply_lorentztrafo_recursive (prt%child1, L)
          call parton_apply_lorentztrafo_recursive (prt%child2, L)
       end if
    else
       if (associated (prt%child1)) then
          call parton_apply_lorentztrafo_recursive (prt%child1, L)
       end if
       if (associated (prt%child2)) then
          call parton_apply_lorentztrafo_recursive (prt%child2, L)
       end if
    end if
  end subroutine parton_apply_lorentztrafo_recursive

\eatline
\nwendcode{}\nwbegindocs{130}This takes the three-momentum of a parton and generates
three-momenta of its children.
\nwenddocs{}\nwbegincode{131}\moddef{Shower partons: public}\plusendmoddef
  public :: parton_generate_ps
\nwendcode{}\nwbegincode{132}\moddef{Shower partons: procedures}\plusendmoddef
  subroutine parton_generate_ps (prt)
    type(parton_t), intent(inout) :: prt
    real(default), dimension(1:3, 1:3) :: directions
    integer i,j
    real(default) :: scproduct, pabs, p1abs, p2abs, x, ptabs, phi
    real(default), dimension(1:3) :: momentum
    type(vector3_t) :: pchild1_direction
    type(lorentz_transformation_t) :: L, rotation
    if (D_PRINT) print *, "D: parton_generate_ps for parton " , prt%nr
    if (ASSERT) then
       if (.not. (associated (prt%child1) .and. associated (prt%child2))) then
          print *, "no children for generate_ps"
          stop 1
       end if
    end if
    !!! test if parton is a virtual parton from the imagined parton shower history
    if (prt%type == INTERNAL) then
       L = inverse (boost (prt%momentum, sqrt(prt%t)))
       !!! boost to restframe of mother
       call parton_apply_lorentztrafo (prt, L)
       call parton_apply_lorentztrafo (prt%child1, L)
       call parton_apply_lorentztrafo (prt%child2, L)
       !!! Store child1's momenta
       pchild1_direction = direction (space_part (prt%child1%momentum))
       !!! Redistribute energy
       prt%child1%momentum%p(0) = (prt%momentum%p(0)**2 - &
            prt%child2%t + prt%child1%t) / (two * prt%momentum%p(0))
       prt%child2%momentum%p(0) = prt%momentum%p(0) - &
            prt%child1%momentum%p(0)

       ! rescale momenta and set momenta to be along z-axis
       prt%child1%momentum = vector4_moving (prt%child1%momentum%p(0), &
            vector3_canonical(3) * &
            sqrt(prt%child1%momentum%p(0)**2 - prt%child1%t))
       prt%child2%momentum = vector4_moving (prt%child2%momentum%p(0), &
            - vector3_canonical(3) * &
            sqrt(prt%child2%momentum%p(0)**2 - prt%child2%t))

       !!! rotate so that total momentum is along former total momentum
       rotation = rotation_to_2nd (space_part (prt%child1%momentum), &
            pchild1_direction)
       call parton_apply_lorentztrafo (prt%child1, rotation)
       call parton_apply_lorentztrafo (prt%child2, rotation)

       L = inverse (L)             !!! inverse of the boost to restframe of mother
       call parton_apply_lorentztrafo (prt, L)
       call parton_apply_lorentztrafo (prt%child1, L)
       call parton_apply_lorentztrafo (prt%child2, L)
    else
       !!! directions(1,:) -> direction of the parent parton
       if (space_part_norm (prt%momentum) < eps0) return
       directions(1,1:3) = prt%momentum%p(1:3) / space_part_norm (prt%momentum)
       !!! directions(2,:) and directions(3,:) -> two random directions
       !!!   perpendicular to the direction of the parent parton
        do j = 2, 3
           call rng%generate(directions(j,:))
        end do
       do i = 2, 3
          scproduct = zero
          do j = 1, i - 1
             scproduct = directions(i,1) * directions(j,1) + &
                directions(i,2) * directions(j,2) + directions(i,3) * directions(j,3)
             directions(i,1) = directions(i,1) - directions(j,1) * scproduct
             directions(i,2) = directions(i,2) - directions(j,2) * scproduct
             directions(i,3) = directions(i,3) - directions(j,3) * scproduct
          end do
          scproduct = directions(i,1)**2 + directions(i,2)**2 + directions(i,3)**2
          do j = 1, 3
             directions(i,j) = directions(i,j) / sqrt(scproduct)
          end do
       end do
       \LA{}Enforce right-handed system for \code{}directions\edoc{}\RA{}

       pabs = space_part_norm (prt%momentum)
       if ((prt%child1%momentum%p(0)**2 - prt%child1%t < 0) .or. &
           (prt%child2%momentum%p(0)**2 - prt%child2%t < 0)) then
          if (D_PRINT) print *, "D: generate_ps error at E^2 < t"
          return
       end if
       p1abs = sqrt (prt%child1%momentum%p(0)**2 - prt%child1%t)
       p2abs = sqrt (prt%child2%momentum%p(0)**2 - prt%child2%t)
       x = (pabs**2 + p1abs**2 - p2abs**2) / (two * pabs)
       if (pabs > p1abs + p2abs .or. &
            pabs < abs(p1abs - p2abs)) then
          if (D_PRINT) then
             print *, "D: parton_generate_ps Dreiecksungleichung error for parton ", &
                     prt%nr, " ", space_part_norm (prt%momentum)," ",p1abs," ",p2abs
             call parton_write (prt)
             call parton_write (prt%child1)
             call parton_write (prt%child2)
          end if
          return
       end if
       !!! Due to numerical problems transverse momentum could be imaginary ->
       !!!     set transverse momentum to zero
       ptabs = sqrt (max (p1abs * p1abs - x * x, zero))
       call rng%generate (phi)
       phi = twopi * phi
       do i = 1, 3
          momentum(i) = x * directions(1,i) + ptabs * &
                (cos(phi) * directions(2,i) + sin(phi) * directions(3,i))
       end do
       prt%child1%momentum%p(1:3) = momentum(1:3)
       do i = 1, 3
          momentum(i) = (space_part_norm (prt%momentum) - x) * directions(1,i) - &
               ptabs * (cos(phi) * directions(2,i) + sin(phi) * directions(3,i))
       end do
       prt%child2%momentum%p(1:3) = momentum(1:3)
    end if
  end subroutine parton_generate_ps

\eatline
\nwendcode{}\nwbegindocs{133}\nwdocspar
\nwenddocs{}\nwbegincode{134}\moddef{Enforce right-handed system for \code{}directions\edoc{}}\endmoddef
if ((directions(1,1) * (directions(2,2) * directions(3,3) - &
     directions(2,3) * directions(3,2)) + &
     directions(1,2) * (directions(2,3) * directions(3,1) - &
     directions(2,1) * directions(3,3)) + &
     directions(1,3) * (directions(2,1) * directions(3,2) - &
     directions(2,2) * directions(3,1))) < 0) then
   directions(3,:) = - directions(3,:)
end if
\nwendcode{}\nwbegindocs{135}\nwdocspar
This routine is similar to {\tt{}parton{\char95}generate{\char95}ps}, but now for the
ISR. It takes the three-momentum of a parton's first child as fixed and
generates the two remaining three-momenta.
\nwenddocs{}\nwbegincode{136}\moddef{Shower partons: public}\plusendmoddef
  public :: parton_generate_ps_ini
\nwendcode{}\nwbegincode{137}\moddef{Shower partons: procedures}\plusendmoddef
  subroutine parton_generate_ps_ini(prt)
    type(parton_t), intent(inout) :: prt
    real(default), dimension(1:3, 1:3) :: directions
    integer :: i,j
    real(default) :: scproduct, pabs, p1abs, p2abs, x, ptabs, phi
    real(default), dimension(1:3) :: momentum
    if (D_PRINT) print *, "D: parton_generate_ps_ini: for parton " , prt%nr
    if (ASSERT) then
       if (.not. (associated (prt%child1) .and. associated (prt%child2))) then
          print *, "no children for generate_ps"
          stop 1
       end if
    end if

    if (parton_is_hadron(prt) .eqv. .false.) then
       !!! generate ps for normal partons
       do i = 1, 3
          directions(1,i) = prt%child1%momentum%p(i) / &
               space_part_norm(prt%child1%momentum)
       end do
       do j = 2, 3
          call rng%generate (directions(j,:))
       end do
       do i = 2, 3
          scproduct = zero
          do j = 1, i - 1
             scproduct = directions(i,1) * directions(j,1) + &
                  directions(i,2) * directions(j,2) + directions(i,3) * directions(j,3)
             directions(i,1) = directions(i,1) - directions(j,1) * scproduct
             directions(i,2) = directions(i,2) - directions(j,2) * scproduct
             directions(i,3) = directions(i,3) - directions(j,3) * scproduct
          end do
          scproduct = directions(i,1)**2 + directions(i,2)**2 + directions(i,3)**2
          do j = 1, 3
             directions(i,j) = directions(i,j) / sqrt(scproduct)
          end do
       end do
       \LA{}Enforce right-handed system for \code{}directions\edoc{}\RA{}

       pabs = space_part_norm (prt%child1%momentum)
       p1abs = sqrt (prt%momentum%p(0)**2 - prt%t)
       p2abs = sqrt (max(zero, prt%child2%momentum%p(0)**2 - &
            prt%child2%t))

       x = (pabs**2 + p1abs**2 - p2abs**2) / (two * pabs)
       if (ASSERT) then
          if (pabs > p1abs + p2abs .or. pabs < abs(p1abs - p2abs)) then
             print *, "error at generate_ps, Dreiecksungleichung for parton ", &
                  prt%nr, " ", pabs," ",p1abs," ",p2abs
             call parton_write (prt)
             call parton_write (prt%child1)
             call parton_write (prt%child2)
             stop 1
          end if
       end if
       if (D_PRINT) print *, "D: parton_generate_ps_ini: x = ", x
       ptabs = sqrt (p1abs * p1abs - x**2)
       call rng%generate (phi)
       phi = twopi * phi
       do i = 1,3
          momentum(i) = x * directions(1,i) + ptabs * (cos(phi) * &
               directions(2,i) + sin(phi) * directions(3,i))
       end do
       prt%momentum%p(1:3) = momentum
       do i = 1, 3
          momentum(i) = (x - pabs) * directions(1,i) + ptabs * (cos(phi) * &
               directions(2,i) + sin(phi) * directions(3,i))
       end do
       prt%child2%momentum%p(1:3) = momentum(1:3)
    else
       !!! for first partons just set beam remnants momentum
       prt%child2%momentum = prt%momentum - prt%child1%momentum
    end if
  end subroutine parton_generate_ps_ini

\eatline
\nwendcode{}\nwbegindocs{138}\nwdocspar
\subsection{The analytic FSR}
\nwenddocs{}\nwbegincode{139}\moddef{Shower partons: public}\plusendmoddef
  public :: parton_next_t_ana
\nwendcode{}\nwbegincode{140}\moddef{Shower partons: procedures}\plusendmoddef
  subroutine parton_next_t_ana (prt)
    type(parton_t), intent(inout) :: prt
    integer :: gtoqq
    real(default) :: integral, random

    if (signal_is_pending ()) return
    if (D_PRINT) then
       print *, "D: parton_next_t_ana: for parton " , prt%nr
    end if

    ! check if branchings are possible at all
    if (min (prt%t, prt%momentum%p(0)**2) < &
         parton_mass_squared(prt) + D_Min_t) then
       prt%t = parton_mass_squared (prt)
       call parton_set_simulated (prt)
       return
    end if

    integral = zero
    call rng%generate (random)
    do
       if (signal_is_pending ()) return
       call parton_simulate_stept (prt, integral, random, gtoqq, .false.)
       if (prt%simulated) then
          if (parton_is_gluon (prt)) then
             !!! Abusing the x-variable to store the information to which
             !!! quark flavour the gluon branches (if any)
             prt%x = one * gtoqq + 0.1_default
             !!! x = gtoqq + 0.1 -> int(x) will be the quark flavour or
             !!! zero for g -> gg
          end if
          exit
       end if
    end do
  end subroutine parton_next_t_ana

\eatline
\nwendcode{}\nwbegindocs{141}The shower is actually sensitive to how close we go to the one here.
\nwenddocs{}\nwbegincode{142}\moddef{Shower partons: procedures}\plusendmoddef
  pure function cmax (prt, tt) result (cma)
    type(parton_t), intent(in) :: prt
    real(default), intent(in), optional :: tt
    real(default) :: cma
    real(default) :: t, cost

    if (present(tt)) then
       t = tt
    else
       t = prt%t
    end if

    if (associated (prt%parent)) then
       cost = parton_get_costheta (prt%parent)
       cma = min (0.99999_default, sqrt( max(zero, one - t/ &
              (parton_get_beta(prt) * prt%momentum%p(0))**2 * &
              (one + cost) / (one - cost))))
    else
       cma = 0.99999_default
    end if
  end function cmax

\eatline
\nwendcode{}\nwbegindocs{143}Simulation routine. The variable {\tt{}lookatsister} takes constraints
from the sister parton into account, if not given it is assumed
{\tt{}.true.}. {\tt{}a} and {\tt{}x} are three-dimensional arrays for values
used for the integration.
\nwenddocs{}\nwbegincode{144}\moddef{Shower partons: public}\plusendmoddef
  public :: parton_simulate_stept
\nwendcode{}\nwbegincode{145}\moddef{Shower partons: procedures}\plusendmoddef
  subroutine parton_simulate_stept &
       (prt, integral, random, gtoqq, lookatsister)
    type(parton_t), intent(inout) :: prt
    real(default), intent(inout) :: integral
    real(default), intent(inout) :: random
    integer, intent(out) :: gtoqq
    logical, intent(in), optional :: lookatsister

    type(parton_t), pointer :: sister
    real(default) :: tstep, tmin, oldt
    real(default) :: c, cstep
    real(default), dimension(3) :: z, P
    real(default) :: to_integral
    real(default) :: a11,a12,a13,a21,a22,a23
    real(default) :: cmax_t
    real(default) :: temprand
    real(default), dimension(3) :: a, x

    ! higher values -> faster but coarser
    real(default), parameter :: tstepfactor = 0.02_default
    real(default), parameter :: tstepmin = 0.5_default
    real(default), parameter :: cstepfactor = 0.8_default
    real(default), parameter :: cstepmin = 0.03_default

    if (signal_is_pending ()) return
    gtoqq = 111 ! illegal value
    call parton_set_simulated (prt, .false.)

    \LA{}Set \code{}sister\edoc{} if \code{}lookatsister\edoc{} is true or not given\RA{}

    tmin = D_Min_t + parton_mass_squared (prt)
    if (parton_is_quark(prt)) then
       to_integral = three *pi * log(one / random)
    else if (parton_is_gluon(prt)) then
       to_integral = four *pi * log(one / random)
    else
       prt%t = parton_mass_squared (prt)
       call parton_set_simulated (prt)
       return
    end if

    if (associated (sister)) then
       if (sqrt(prt%t) > sqrt(prt%parent%t) - &
            sqrt(parton_mass_squared (sister))) then
          prt%t = (sqrt (prt%parent%t) - sqrt (parton_mass_squared (sister)))**2
       end if
    end if
    if (prt%t > prt%momentum%p(0)**2) then
        prt%t = prt%momentum%p(0)**2
    end if

    if (prt%t <= tmin) then
       prt%t = parton_mass_squared (prt)
       call parton_set_simulated (prt)
       return
    end if

    ! simulate the branchings between prt%t and prt%t - tstep
    tstep = max(tstepfactor * (prt%t - 0.9_default * tmin), tstepmin)
    cmax_t = cmax(prt)
    c = - cmax_t ! take highest t -> minimal constraint
    cstep = max(cstepfactor * (one - abs(c)), cstepmin)
    ! get values at border of "previous" bin -> to be used in first bin
    z(3) = 0.5_default + 0.5_default * get_beta (prt%t - &
         0.5_default * tstep, prt%momentum%p(0)) * c
    if (parton_is_gluon (prt)) then
       P(3) = P_ggg (z(3)) + P_gqq (z(3)) * number_of_flavors (prt%t)
    else
       P(3) = P_qqg (z(3))
    end if
    a(3) = D_alpha_s_fsr (z(3) * (one - z(3)) * prt%t) * P(3) / &
         (prt%t - 0.5_default * tstep)

    do while (c < cmax_t .and. (integral < to_integral))
       if (signal_is_pending ()) return
       cmax_t = cmax (prt)
       cstep = max (cstepfactor * (one - abs(c)**2), cstepmin)
       if (c + cstep > cmax_t) then
          cstep = cmax_t - c
       end if
       if (cstep < 1E-9_default) then
          !!! reject too small bins
          exit
       end if
       z(1) = z(3)
       z(2) = 0.5_default + 0.5_default * get_beta &
            (prt%t - 0.5_default * tstep, prt%momentum%p(0)) * &
            (c + 0.5_default * cstep)
       z(3) = 0.5_default + 0.5_default * get_beta &
            (prt%t - 0.5_default * tstep, prt%momentum%p(0)) * (c + cstep)
       P(1) = P(3)
       if (parton_is_gluon (prt)) then
          P(2) = P_ggg(z(2)) + P_gqq(z(2)) * number_of_flavors (prt%t)
          P(3) = P_ggg(z(3)) + P_gqq(z(3)) * number_of_flavors (prt%t)
       else
          P(2) = P_qqg(z(2))
          P(3) = P_qqg(z(3))
       end if
       ! get values at borders of the intgral and in the middle
       a(1) = a(3)
       a(2) = D_alpha_s_fsr(z(2) * (one - z(2)) * prt%t) * P(2) / &
            (prt%t - 0.5_default * tstep)
       a(3) = D_alpha_s_fsr(z(3) * (one - z(3)) * prt%t) * P(3) / &
            (prt%t - 0.5_default * tstep)

       ! fit x(1) + x(2)/(1 + c) + x(3)/(1 - c) to these values !! a little tricky
       a11 = (one+c+0.5_default*cstep) * (one-c-0.5_default*cstep) - &
             (one-c) * (one+c+0.5_default*cstep)
       a12 = (one-c-0.5_default*cstep) - (one+c+0.5_default*cstep) * &
             (one-c) / (one+c)
       a13 = a(2) * (one+c+0.5_default*cstep) * (one-c-0.5_default*cstep) - &
             a(1) * (one-c) * (one+c+0.5_default*cstep)
       a21 = (one+c+cstep) * (one-c-cstep) - (one+c+cstep) * (one-c)
       a22 = (one-c-cstep) - (one+c+cstep) * (one-c) / (one+c)
       a23 = a(3) * (one+c+cstep) * (one-c-cstep) - &
            a(1) * (one-c) * (one+c+cstep)

       x(2) = (a23 - a21 * a13 / a11) / (a22 - a12 * a21 / a11)
       x(1) = (a13 - a12 * x(2)) / a11
       x(3) = a(1) * (one - c) - x(1) * (one - c) - x(2) * (one - c) / (one + c)

       integral = integral + tstep * (x(1) * cstep + x(2) * &
            log((one + c + cstep) / (one + c)) - x(3) * &
            log((one - c - cstep) / (one - c)))

       if (integral > to_integral) then
          oldt = prt%t
          call rng%generate (temprand)
          prt%t = prt%t - temprand * tstep
          call rng%generate (temprand)
          prt%costheta = c + (0.5_default - temprand) * cstep
          call parton_set_simulated (prt)

          if (prt%t < D_Min_t + parton_mass_squared(prt)) then
             prt%t = parton_mass_squared (prt)
          end if
          if (abs(prt%costheta) > cmax_t) then
             ! reject branching due to violation of costheta-limits
             call rng%generate (random)
             if (parton_is_quark (prt)) then
                to_integral = three * pi * log(one / random)
             else if (parton_is_gluon(prt)) then
                to_integral = four * pi * log(one / random)
             end if
             integral = zero
             prt%t = oldt
             call parton_set_simulated (prt, .false.)
          end if
          if (parton_is_gluon (prt)) then
             ! decide between g->gg and g->qqbar splitting
             z(1) = 0.5_default + 0.5_default * prt%costheta
             call rng%generate (temprand)
             if (P_ggg(z(1)) > temprand * (P_ggg (z(1)) + P_gqq (z(1)) * &
                  number_of_flavors(prt%t))) then
                gtoqq = 0
             else
                call rng%generate (temprand)
                gtoqq = 1 + int (temprand * number_of_flavors (prt%t))
             end if
          end if
       else
          c = c + cstep
       end if
       cmax_t = cmax (prt)
    end do
    if (integral <= to_integral) then
       prt%t = prt%t - tstep
       if (prt%t < D_Min_t + parton_mass_squared (prt)) then
          prt%t = parton_mass_squared (prt)
          call parton_set_simulated(prt)
       end if
    end if
  end subroutine parton_simulate_stept

\eatline
\nwendcode{}\nwbegindocs{146}\nwdocspar
\nwenddocs{}\nwbegincode{147}\moddef{Set \code{}sister\edoc{} if \code{}lookatsister\edoc{} is true or not given}\endmoddef
sister => null()
SET_SISTER: do
   if (present (lookatsister)) then
      if (.not. lookatsister) then
         exit SET_SISTER
      end if
   end if
   if (prt%nr == prt%parent%child1%nr) then
      sister => prt%parent%child2
   else
      sister => prt%parent%child1
   end if
   exit SET_SISTER
end do SET_SISTER
\nwendcode{}\nwbegindocs{148}\nwdocspar
\nwenddocs{}\nwbegindocs{149}From the whole ISR algorithm all functionality has been moved to
{\tt{}shower{\char95}core.f90}. Only {\tt{}maxzz} remains here, because more than
one module needs to access it.
\nwenddocs{}\nwbegincode{150}\moddef{Shower partons: public}\plusendmoddef
  public :: maxzz
\nwendcode{}\nwbegincode{151}\moddef{Shower partons: procedures}\plusendmoddef
  function maxzz (shat, s, maxz_isr, minenergy_timelike) result (maxz)
    real(default), intent(in) :: shat, s, minenergy_timelike, maxz_isr
    real(default) :: maxz
    maxz = min (maxz_isr, one - (two * minenergy_timelike * sqrt(shat)) / s)
  end function maxzz

\eatline
\nwendcode{}\nwbegindocs{152}\nwdocspar
%%%%%%%%%%%%%%%%%%%%%%%%%%%%%%%%%%%%%%%%%%%%%%%%%%%%%%%%%%%%%%%%%%%%%%%%
\section{Main shower module}
\nwenddocs{}\nwbegincode{153}\moddef{\code{}shower{\char95}core.f90\edoc{}}\endmoddef
\LA{}File header\RA{}

module shower_core

\LA{}Use kinds with double\RA{}
  use io_units
  use constants
  use unit_tests, only: vanishes, nearly_equal
  use system_defs, only: TAB
  use diagnostics
  use physics_defs
  use lorentz
  use sm_physics
  use pdf_builtin !NODEP!
  use lhapdf !NODEP!
  use shower_base
  use shower_partons
  use ckkw_pseudo_weights

\LA{}Standard module head\RA{}

\LA{}Shower core: public\RA{}

\LA{}Shower core: parameters\RA{}

\LA{}Shower core: types\RA{}

\LA{}Shower core: interfaces\RA{}

contains

\LA{}Shower core: procedures\RA{}

end module shower_core
\eatline
\nwendcode{}\nwbegindocs{154}\nwdocspar
\nwenddocs{}\nwbegincode{155}\moddef{Shower core: interfaces}\endmoddef
  interface
     subroutine evolvePDFM (set, x, q, ff)
       integer, intent(in) :: set
       double precision, intent(in) :: x, q
       double precision, dimension(-6:6), intent(out) :: ff
     end subroutine evolvePDFM
  end interface

\eatline
\nwendcode{}\nwbegindocs{156}\nwdocspar
\nwenddocs{}\nwbegincode{157}\moddef{Shower core: public}\endmoddef
  public :: shower_interaction_t
\nwendcode{}\nwbegincode{158}\moddef{Shower core: types}\endmoddef
  type :: shower_interaction_t
     type(parton_pointer_t), dimension(:), allocatable :: partons
  end type shower_interaction_t

  type :: shower_interaction_pointer_t
     type(shower_interaction_t), pointer :: i => null ()
  end type shower_interaction_pointer_t

\eatline
\eatline
\nwendcode{}\nwbegindocs{159}The main type of {\tt{}shower{\char95}core}: It contains many parameter of the
shower, namely {\tt{}maxz{\char95}isr} ({\tt{}PARP(66)}), the minimal energy
of the emitted timelike parton in the ISR {\tt{}minenergy{\char95}timelike}
({\tt{}PARP(65)}), the factor for the first scale, as
well as the factor by which the integral in the Sudakhov factor is
suppressed for the respective first scale in ISR. Higher values
here correspond  to higher starting scales. Also, there is a flag
whether ISR is angular ordered ({\tt{}MSTP(62)}). The flag corresponding
to {\tt{}MSTP(63)} sets emitted timelike partons in the spacelike
shower on shell.
\nwenddocs{}\nwbegincode{160}\moddef{Shower core: public}\plusendmoddef
  public :: shower_t
\nwendcode{}\nwbegincode{161}\moddef{Shower core: types}\plusendmoddef
  type :: shower_t
     type(shower_interaction_pointer_t), dimension(:), allocatable :: &
          interactions
     type(parton_pointer_t), dimension(:), allocatable :: partons
     type(lhapdf_pdf_t), pointer :: pdf => null ()
     integer :: next_free_nr
     integer :: next_color_nr
     logical :: valid
     integer :: pdf_set = 0
     integer :: pdf_type = STRF_NONE
     logical :: isr_angular_ordered = .true.
     real(default) :: maxz_isr = 0.999_default
     real(default) :: minenergy_timelike = one
     real(default) :: tscalefactor_isr = one
     real(default) :: first_integral_suppression_factor = one
     logical :: isr_only_onshell_emitted_partons = .true.
     real(default) :: xmin = 0, xmax = 0
     real(default) :: qmin = 0, qmax = 0
   contains
   \LA{}Shower core: shower: TBP\RA{}
  end type shower_t

\eatline
\nwendcode{}\nwbegindocs{162}The parameters of the shower module:
\nwenddocs{}\nwbegincode{163}\moddef{Shower core: parameters}\endmoddef
  real(default), save :: alphasxpdfmax = 12._default

\eatline
\nwendcode{}\nwbegindocs{164}\nwdocspar
\nwenddocs{}\nwbegincode{165}\moddef{Shower core: shower: TBP}\endmoddef
  procedure :: add_interaction_2ton => shower_add_interaction_2ton
\nwendcode{}\nwbegincode{166}\moddef{Shower core: procedures}\endmoddef
  subroutine shower_add_interaction_2ton (shower, partons)
    class(shower_t), intent(inout) :: shower
    type(parton_pointer_t), intent(in), dimension(:), allocatable :: partons
    type(ckkw_pseudo_shower_weights_t) :: ckkw_pseudo_weights
    call shower%add_interaction_2ton_CKKW (partons, ckkw_pseudo_weights)
  end subroutine shower_add_interaction_2ton

\eatline
\nwendcode{}\nwbegindocs{167}In this routine, {\tt{}y} and {\tt{}ymin} are the jet measures, {\tt{}w} and
{\tt{}wmax} are weights, {\tt{}s} is the kinematic energy squared of the
interaction. The flag {\tt{}isr{\char95}is{\char95}possible} checks whether the initial
parton is set, lepton-hadron collisions are not implemented (yet). As
a workaround: as WHIZARD can treat partons as massless, there might be
partons with $E < m$: if such a parton is found, quarks will be
treated massless.
\nwenddocs{}\nwbegincode{168}\moddef{Shower core: shower: TBP}\plusendmoddef
  procedure :: add_interaction_2ton_CKKW => shower_add_interaction_2ton_CKKW
\nwendcode{}\nwbegincode{169}\moddef{Shower core: procedures}\plusendmoddef
  subroutine shower_add_interaction_2ton_CKKW &
       (shower, partons, ckkw_pseudo_weights)
    class(shower_t), intent(inout) :: shower
    type(parton_pointer_t), intent(in), dimension(:), allocatable :: partons
    type(ckkw_pseudo_shower_weights_t), intent(in) :: ckkw_pseudo_weights

    integer :: n_partons, n_out, n_partons_shower
    integer :: i, j, imin, jmin
    real(default) :: y, ymin
    real(default) :: w, wmax
    real(default) :: random, sum
    type(parton_pointer_t), dimension(:), allocatable :: new_partons
    type(parton_t), pointer :: prt
    integer :: n_int
    type(shower_interaction_pointer_t), dimension(:), allocatable :: temp
    type(vector4_t) :: prtmomentum, childmomentum
    logical :: isr_is_possible
    type(lorentz_transformation_t) :: L

    if (signal_is_pending ()) return
    n_partons = size (partons)
    n_out = n_partons - 2
    if (n_out < 2) then
       call msg_bug &
            ("Shower core: trying to add a 2-> (something<2) interaction")
    end if

    isr_is_possible = associated (partons(1)%p%initial) .and. &
         associated (partons(2)%p%initial)

    if (associated (partons(1)%p%initial) .and. &
         parton_is_quark (partons(1)%p)) then
       if (partons(1)%p%momentum%p(0) < &
            two * parton_mass (partons(1)%p)) then
          if (abs(partons(1)%p%type) < 2) then
             treat_light_quarks_massless = .true.
          else
             treat_duscb_quarks_massless = .true.
          end if
       end if
    end if
    if (associated (partons(2)%p%initial) .and. &
         parton_is_quark (partons(2)%p)) then
       if (partons(2)%p%momentum%p(0) < &
            two * parton_mass (partons(2)%p)) then
          if (abs(partons(2)%p%type) < 2) then
             treat_light_quarks_massless = .true.
          else
             treat_duscb_quarks_massless = .true.
          end if
       end if
    end if

    \LA{}Add a new interaction to \code{}shower{\char37}interactions\edoc{}\RA{}

    if (associated (shower%interactions(n_int)%i%partons(1)%p%initial))  &
         call parton_set_simulated &
         (shower%interactions(n_int)%i%partons(1)%p%initial)
    if (associated (shower%interactions(n_int)%i%partons(2)%p%initial))  &
         call parton_set_simulated &
         (shower%interactions(n_int)%i%partons(2)%p%initial)
    if (isr_is_possible) then
       !!! boost to the CMFrame of the incoming partons
       L = boost (-(shower%interactions(n_int)%i%partons(1)%p%momentum + &
            shower%interactions(n_int)%i%partons(2)%p%momentum), &
            (shower%interactions(n_int)%i%partons(1)%p%momentum + &
            shower%interactions(n_int)%i%partons(2)%p%momentum)**1 )
       do i = 1, n_partons
          call parton_apply_lorentztrafo &
               (shower%interactions(n_int)%i%partons(i)%p, L)
       end do
    end if
    do i = 1, size (partons)
       if (signal_is_pending ()) return
       !!! ensure that partons are marked as belonging to the hard interaction
       shower%interactions(n_int)%i%partons(i)%p%belongstointeraction &
            = .true.
       !!! ensure that incoming partons are marked as belonging to ISR
       shower%interactions(n_int)%i%partons(i)%p%belongstoFSR = i > 2

       shower%interactions(n_int)%i%partons(i)%p%interactionnr &
            = n_int
       !!! include a 2^(i - 1) number as a label for the ckkw clustering
       shower%interactions(n_int)%i%partons(i)%p%ckkwlabel = 2**(i - 1)
    end do

    \LA{}Add partons from \code{}shower{\char37}interactions\edoc{} to \code{}shower{\char37}partons\edoc{}\RA{}

    if (isr_is_possible) then
       if (isr_pt_ordered) then
          call shower_prepare_for_simulate_isr_pt &
               (shower, shower%interactions(size (shower%interactions))%i)
       else
          call shower_prepare_for_simulate_isr_ana_test &
               (shower, shower%interactions(n_int)%i%partons(1)%p, &
               shower%interactions(n_int)%i%partons(2)%p)
       end if
    end if

    !!! generate pseudo parton shower history and add all partons to
    !!!          shower%partons-array
    !!! TODO initial -> initial + final branchings ??
    allocate (new_partons(1:(n_partons - 2)))
    do i = 1, size (new_partons)
       nullify (new_partons(i)%p)
    end do
    do i = 1, size (new_partons)
       new_partons(i)%p => shower%interactions(n_int)%i%partons(i + 2)%p
    end do
    imin = 0
    jmin = 0

    if (allocated (ckkw_pseudo_weights%weights)) then
       \LA{}Perform clustering using the CKKW weights\RA{}
    else
       \LA{}Perform clustering in the usual way\RA{}
    end if

    !!! add all partons to the shower
    \LA{}Count number of associated \code{}shower{\char37}partons(i){\char37}p\edoc{}\RA{}

    !!! set the FSR starting scale for all partons
    do i = 1, size (new_partons)
       !!! the imaginary mother is the only parton remaining in new_partons
       if (.not. associated (new_partons(i)%p)) cycle

       call set_starting_scale (new_partons(i)%p, &
            get_starting_scale (new_partons(i)%p))
       exit
    end do

    deallocate (new_partons)

  contains

    \LA{}Procedures of \code{}shower{\char95}add{\char95}interaction{\char95}2ton{\char95}CKKW\edoc{}\RA{}
  end subroutine shower_add_interaction_2ton_CKKW

\eatline
\nwendcode{}\nwbegindocs{170}\nwdocspar
\nwenddocs{}\nwbegincode{171}\moddef{Add a new interaction to \code{}shower{\char37}interactions\edoc{}}\endmoddef
if (allocated (shower%interactions)) then
   n_int = size (shower%interactions) + 1
else
   n_int = 1
end if
allocate (temp (1:n_int))
do i = 1, n_int - 1
   allocate (temp(i)%i)
   temp(i)%i = shower%interactions(i)%i
end do
allocate (temp(n_int)%i)
allocate (temp(n_int)%i%partons(1:n_partons))
do i = 1, n_partons
   allocate (temp(n_int)%i%partons(i)%p)
   call parton_copy (partons(i)%p, temp(n_int)%i%partons(i)%p)
end do
if (allocated (shower%interactions)) deallocate(shower%interactions)
allocate (shower%interactions(1:n_int))
do i = 1, n_int
   shower%interactions(i)%i => temp(i)%i
end do
deallocate (temp)
\nwendcode{}\nwbegindocs{172}\nwdocspar
\nwenddocs{}\nwbegincode{173}\moddef{Add partons from \code{}shower{\char37}interactions\edoc{} to \code{}shower{\char37}partons\edoc{}}\endmoddef
if (allocated (shower%partons)) then
   allocate (new_partons(1:size(shower%partons) + &
        size(shower%interactions(n_int)%i%partons)))
   do i = 1, size (shower%partons)
      new_partons(i)%p => shower%partons(i)%p
   end do
   do i = 1, size (shower%interactions(n_int)%i%partons)
      new_partons(size(shower%partons) + i)%p => &
           shower%interactions(n_int)%i%partons(i)%p
   end do
   deallocate (shower%partons)
else
   allocate (new_partons(1:size(shower%interactions(n_int)%i%partons)))
   do i = 1, size (partons)
      new_partons(i)%p => shower%interactions(n_int)%i%partons(i)%p
   end do
end if
allocate (shower%partons(1:size (new_partons)))
do i = 1, size (new_partons)
   shower%partons(i)%p => new_partons(i)%p
end do
deallocate (new_partons)
\nwendcode{}\nwbegindocs{174}\nwdocspar
\nwenddocs{}\nwbegincode{175}\moddef{Perform clustering using the CKKW weights}\endmoddef
CKKW_CLUSTERING: do
   !!! search for the combination with the highest weight
   wmax = zero
   CKKW_OUTER: do i = 1, size (new_partons)
      CKKW_INNER: do j = i + 1, size (new_partons)
         if (.not. associated (new_partons(i)%p)) cycle
         if (.not. associated (new_partons(j)%p)) cycle
         w = ckkw_pseudo_weights%weights(new_partons(i)%p%ckkwlabel + &
              new_partons(j)%p%ckkwlabel)
         if (w > wmax .or. vanishes(wmax)) then
            wmax = w
            imin = i
            jmin = j
         end if
      end do CKKW_INNER
   end do CKKW_OUTER
   if (wmax > zero) then
      call shower%add_parent (new_partons(imin)%p)
      call parton_set_child (new_partons(imin)%p%parent, &
           new_partons(jmin)%p, 2)
      call parton_set_parent (new_partons(jmin)%p, &
           new_partons(imin)%p%parent)
      prt => new_partons(imin)%p%parent
      prt%nr = shower_get_next_free_nr (shower)
      prt%type = INTERNAL

      prt%momentum = new_partons(imin)%p%momentum + &
           new_partons(jmin)%p%momentum
      prt%t = prt%momentum**2

      !!! auxilliary values for the ckkw matching
      !!! for now, randomly choose the type of the intermediate
      prt%ckkwlabel = new_partons(imin)%p%ckkwlabel + &
           new_partons(jmin)%p%ckkwlabel
      sum = zero
      call rng%generate (random)
      CKKW_TYPE: do i = 0, 4
         if (sum + &
              ckkw_pseudo_weights%weights_by_type(prt%ckkwlabel, i) > &
              random * ckkw_pseudo_weights%weights(prt%ckkwlabel) ) then
            prt%ckkwtype = i
            exit ckkw_type
         end if
         sum = sum + &
              ckkw_pseudo_weights%weights_by_type(prt%ckkwlabel, i)
      end do CKKW_TYPE

      !!! TODO -> calculate costheta and store it for
      !!!        later use in generate_ps

      if (space_part_norm(prt%momentum) > tiny_10) then
         prtmomentum = prt%momentum
         childmomentum = prt%child1%momentum
         prtmomentum = boost (-parton_get_beta(prt) / &
              sqrt (one - &
              (parton_get_beta(prt))**2), space_part (prt%momentum) / &
              space_part_norm(prt%momentum)) * prtmomentum
         childmomentum = boost (-parton_get_beta(prt) / &
              sqrt(one - &
              (parton_get_beta(prt))**2), space_part (prt%momentum) / &
              space_part_norm(prt%momentum)) * childmomentum
         prt%costheta = enclosed_angle_ct(prtmomentum, childmomentum)
      else
         prt%costheta = - one
      end if

      prt%belongstointeraction = .true.
      prt%belongstoFSR = &
           new_partons(imin)%p%belongstoFSR .and. &
           new_partons(jmin)%p%belongstoFSR

      nullify (new_partons(imin)%p)
      nullify (new_partons(jmin)%p)
      new_partons(imin)%p => prt
   else
      exit CKKW_CLUSTERING
   end if
end do CKKW_CLUSTERING
\nwendcode{}\nwbegindocs{176}\nwdocspar
\nwenddocs{}\nwbegincode{177}\moddef{Perform clustering in the usual way}\endmoddef
CLUSTERING: do
   !!! search for the partons to be clustered together
   ymin = zero
   OUTER: do i = 1, size (new_partons)
      INNER: do j = i + 1, size (new_partons)
         !!! calculate the jet measure
         if (.not.associated (new_partons(i)%p)) cycle INNER
         if (.not.associated (new_partons(j)%p)) cycle INNER
         !if (.not. shower_clustering_allowed &
              !(shower, new_partons, i,j)) &
              !cycle inner
         !!! Durham jet-measure ! don't care about constants
         y = min (new_partons(i)%p%momentum%p(0), &
              new_partons(j)%p%momentum%p(0)) * &
              (one - enclosed_angle_ct &
              (new_partons(i)%p%momentum, &
              new_partons(j)%p%momentum))
         if (y < ymin .or. vanishes(ymin)) then
            ymin = y
            imin = i
            jmin = j
         end if
      end do INNER
   end do OUTER
   if (ymin > zero) then
      call shower%add_parent (new_partons(imin)%p)
      call parton_set_child &
           (new_partons(imin)%p%parent, new_partons(jmin)%p, 2)
      call parton_set_parent &
           (new_partons(jmin)%p, new_partons(imin)%p%parent)
      prt => new_partons(imin)%p%parent
      prt%nr = shower_get_next_free_nr (shower)
      prt%type = INTERNAL

      prt%momentum = new_partons(imin)%p%momentum + &
           new_partons(jmin)%p%momentum
      prt%t = prt%momentum**2
      !!! TODO -> calculate costheta and store it for
      !!!       later use in generate_ps

      if (space_part_norm(prt%momentum) > tiny_10) then
         prtmomentum = prt%momentum
         childmomentum = prt%child1%momentum
         prtmomentum = boost (-parton_get_beta(prt) / sqrt(one - &
              (parton_get_beta(prt))**2), space_part(prt%momentum) / &
              space_part_norm(prt%momentum)) * prtmomentum
         childmomentum = boost (-parton_get_beta(prt) / &
              sqrt(one - &
              (parton_get_beta(prt))**2), space_part(prt%momentum) / &
              space_part_norm(prt%momentum)) * childmomentum
         prt%costheta = enclosed_angle_ct(prtmomentum, childmomentum)
      else
         prt%costheta = - one
      end if

      prt%belongstointeraction = .true.
      nullify (new_partons(imin)%p)
      nullify (new_partons(jmin)%p)
      new_partons(imin)%p => prt
   else
      exit CLUSTERING
   end if
end do CLUSTERING
\nwendcode{}\nwbegindocs{178}\nwdocspar
\nwenddocs{}\nwbegincode{179}\moddef{Count number of associated \code{}shower{\char37}partons(i){\char37}p\edoc{}}\endmoddef
n_partons_shower = 0
if (allocated (shower%partons)) then
   do i = 1, size (shower%partons)
      if (associated (shower%partons(i)%p)) &
           n_partons_shower = n_partons_shower + 1
   end do
end if
\nwendcode{}\nwbegindocs{180}\nwdocspar
\nwenddocs{}\nwbegincode{181}\moddef{Procedures of \code{}shower{\char95}add{\char95}interaction{\char95}2ton{\char95}CKKW\edoc{}}\endmoddef
recursive subroutine transfer_pointers (destiny, start, prt)
  type(parton_pointer_t), dimension(:), allocatable :: destiny
  integer, intent(inout) :: start
  type(parton_t), pointer :: prt
  destiny(start)%p => prt
  start = start + 1
  if (associated (prt%child1)) then
     call transfer_pointers (destiny, start, prt%child1)
  end if
  if (associated (prt%child2)) then
     call transfer_pointers (destiny, start, prt%child2)
  end if
end subroutine transfer_pointers

\nwendcode{}\nwbegindocs{182}\nwdocspar
\nwenddocs{}\nwbegincode{183}\moddef{Procedures of \code{}shower{\char95}add{\char95}interaction{\char95}2ton{\char95}CKKW\edoc{}}\plusendmoddef
recursive function get_starting_scale (prt) result (scale)
  type(parton_t), pointer :: prt
  real(default) :: scale
  scale = huge (scale)
  if (associated (prt%child1) .and. associated (prt%child2)) then
     scale = min(scale, prt%t)
  end if
  if (associated (prt%child1)) then
     scale = min (scale, get_starting_scale (prt%child1))
  end if
  if (associated (prt%child2)) then
     scale = min (scale, get_starting_scale (prt%child2))
  end if
end function get_starting_scale

\nwendcode{}\nwbegindocs{184}\nwdocspar
\nwenddocs{}\nwbegincode{185}\moddef{Procedures of \code{}shower{\char95}add{\char95}interaction{\char95}2ton{\char95}CKKW\edoc{}}\plusendmoddef
recursive subroutine set_starting_scale (prt, scale)
  type(parton_t), pointer :: prt
  real(default) :: scale
  if (prt%type /= INTERNAL) then
     if (scale > D_Min_t + parton_mass_squared (prt)) then
        prt%t = scale
     else
        prt%t = parton_mass_squared (prt)
        call parton_set_simulated (prt)
     end if
  end if
  if (associated (prt%child1)) then
     call set_starting_scale (prt%child1, scale)
  end if
  if (associated (prt%child2)) then
     call set_starting_scale (prt%child2, scale)
  end if
end subroutine set_starting_scale

\nwendcode{}\nwbegindocs{186}\nwdocspar
\nwenddocs{}\nwbegincode{187}\moddef{Shower core: shower: TBP}\plusendmoddef
  procedure :: simulate_no_isr_shower => shower_simulate_no_isr_shower
\nwendcode{}\nwbegincode{188}\moddef{Shower core: procedures}\plusendmoddef
  subroutine shower_simulate_no_isr_shower (shower)
    class(shower_t), intent(inout) :: shower
    integer :: i, j
    type(parton_t), pointer :: prt
    do i = 1, size (shower%interactions)
       do j = 1, 2
          prt => shower%interactions(i)%i%partons(j)%p
          if (associated (prt%initial)) then
             !!! for virtuality ordered: remove unneeded partons
             if (associated (prt%parent)) then
                if (.not. parton_is_hadron(prt%parent)) then
                   if (associated (prt%parent%parent)) then
                      if (.not. parton_is_hadron (prt%parent)) then
                         call shower_remove_parton_from_partons &
                              (shower, prt%parent%parent)
                      end if
                   end if
                   call shower_remove_parton_from_partons &
                        (shower, prt%parent)
                end if
             end if
             call parton_set_parent (prt, prt%initial)
             call parton_set_child (prt%initial, prt, 1)
             if (associated (prt%initial%child2)) then
                call shower_remove_parton_from_partons &
                     (shower,prt%initial%child2)
                deallocate (prt%initial%child2)
             end if
             call shower%add_child (prt%initial, 2)
          end if
       end do
    end do
  end subroutine shower_simulate_no_isr_shower

\eatline
\nwendcode{}\nwbegindocs{189}\nwdocspar
\nwenddocs{}\nwbegincode{190}\moddef{Shower core: shower: TBP}\plusendmoddef
  procedure :: simulate_no_fsr_shower => shower_simulate_no_fsr_shower
\nwendcode{}\nwbegincode{191}\moddef{Shower core: procedures}\plusendmoddef
  subroutine shower_simulate_no_fsr_shower (shower)
    class(shower_t), intent(inout) :: shower
    integer :: i, j
    type(parton_t), pointer :: prt
    do i = 1, size (shower%interactions)
       do j = 3, size (shower%interactions(i)%i%partons)
          prt => shower%interactions(i)%i%partons(j)%p
          call parton_set_simulated (prt)
          prt%scale = zero
          prt%t = parton_mass_squared (prt)
       end do
    end do
  end subroutine shower_simulate_no_fsr_shower

\eatline
\nwendcode{}\nwbegindocs{192}\nwdocspar
\nwenddocs{}\nwbegincode{193}\moddef{Shower core: procedures}\plusendmoddef
  subroutine swap_pointers (prtp1, prtp2)
    type(parton_pointer_t), intent(inout) :: prtp1, prtp2
    type(parton_pointer_t) :: prtptemp
    prtptemp%p => prtp1%p
    prtp1%p => prtp2%p
    prtp2%p => prtptemp%p
  end subroutine swap_pointers

\eatline
\nwendcode{}\nwbegindocs{194}This removes emitted timelike partons.
\nwenddocs{}\nwbegincode{195}\moddef{Shower core: procedures}\plusendmoddef
  recursive subroutine shower_remove_parton_from_partons (shower, prt)
    type(shower_t), intent(inout) :: shower
    type(parton_t), pointer :: prt
    integer :: i
    if (.not. prt%belongstoFSR .and. associated (prt%child2)) then
       call shower_remove_parton_from_partons_recursive (shower, prt%child2)
    end if
    do i = 1, size (shower%partons)
       if (associated (shower%partons(i)%p, prt)) then
          shower%partons(i)%p => null()
          exit
       end if
       if (ASSERT) then
          if (i == size (shower%partons)) then
             print *, "shower_remove_parton_from_partons: parton to be removed not found"
             stop 1
          end if
       end if
    end do
  end subroutine shower_remove_parton_from_partons

\eatline
\nwendcode{}\nwbegindocs{196}This removes the parton {\tt{}prt} and all its children.
\nwenddocs{}\nwbegincode{197}\moddef{Shower core: procedures}\plusendmoddef
  recursive subroutine shower_remove_parton_from_partons_recursive (shower, prt)
    type(shower_t), intent(inout) :: shower
    type(parton_t), pointer :: prt
    if (associated (prt%child1)) then
       call shower_remove_parton_from_partons_recursive (shower, prt%child1)
       deallocate (prt%child1)
    end if
    if (associated (prt%child2)) then
       call shower_remove_parton_from_partons_recursive (shower, prt%child2)
       deallocate (prt%child2)
    end if
    call shower_remove_parton_from_partons (shower, prt)
  end subroutine shower_remove_parton_from_partons_recursive

\eatline
\nwendcode{}\nwbegindocs{198}\nwdocspar
\nwenddocs{}\nwbegincode{199}\moddef{Shower core: shower: TBP}\plusendmoddef
  procedure :: sort_partons => shower_sort_partons
\nwendcode{}\nwbegincode{200}\moddef{Shower core: procedures}\plusendmoddef
  subroutine shower_sort_partons (shower)
    class(shower_t), intent(inout) :: shower
    integer :: i, j, maxsort, size_partons
    logical :: changed
    if (D_PRINT) print *, "D: shower_sort_partons"
    if (.not. allocated (shower%partons)) return
    size_partons = size (shower%partons)
    maxsort = 0
    do i = 1, size_partons
       if (associated (shower%partons(i)%p))  maxsort = i
    end do
    if (signal_is_pending ()) return
    size_partons = size (shower%partons)
    if (size_partons <= 1) return
    do i = 1, maxsort
       if (.not. associated (shower%partons(i)%p)) cycle
       if (.not. isr_pt_ordered) then
          !!! set unsimulated ISR partons to be "typeless" to prevent
          !!!     influences from "wrong" masses
          if (.not. shower%partons(i)%p%belongstoFSR .and. &
              .not. shower%partons(i)%p%simulated .and. &
              .not. shower%partons(i)%p%belongstointeraction) then
             shower%partons(i)%p%type = 0
          end if
       end if
    end do
    if (signal_is_pending ()) return
    !!! Just a Bubblesort
    !!! Different algorithms needed for t-ordered and pt^2-ordered shower
    !!! Pt-ordered:
    if (isr_pt_ordered) then
       OUTERDO_PT: do i = 1, maxsort - 1
          changed = .false.
          INNERDO_PT: do j = 1, maxsort - i
             if (.not. associated (shower%partons(j + 1)%p)) cycle
             if (.not. associated (shower%partons(j)%p)) then
                !!! change if j + 1 ist assoaciated and j is not
                call swap_pointers (shower%partons(j), shower%partons(j + 1))
                changed = .true.
             else if (shower%partons(j)%p%scale < &
                  shower%partons(j + 1)%p%scale) then
                call swap_pointers (shower%partons(j), shower%partons(j + 1))
                changed = .true.
             else if (nearly_equal(shower%partons(j)%p%scale, &
                  shower%partons(j + 1)%p%scale)) then
                if (shower%partons(j)%p%nr > shower%partons(j + 1)%p%nr) then
                   call swap_pointers (shower%partons(j), shower%partons(j + 1))
                   changed = .true.
                end if
             end if
          end do INNERDO_PT
          if (.not. changed) exit OUTERDO_PT
       end do outerdo_pt
    !!! |t|-ordered
    else
       OUTERDO_T: do i = 1, maxsort - 1
          changed = .false.
          INNERDO_T: do j = 1, maxsort - i
             if (.not. associated (shower%partons(j + 1)%p)) cycle
             if (.not. associated (shower%partons(j)%p)) then
                !!! change if j+1 is associated and j isn't
                call swap_pointers (shower%partons(j), shower%partons(j + 1))
                changed = .true.
             else if (.not. shower%partons(j)%p%belongstointeraction .and. &
                  shower%partons(j + 1)%p%belongstointeraction) then
                !!! move partons belonging to the interaction to the front
                call swap_pointers (shower%partons(j), shower%partons(j + 1))
                changed = .true.
             else if (.not. shower%partons(j)%p%belongstointeraction .and. &
                  .not. shower%partons(j + 1)%p%belongstointeraction ) then
                if (abs (shower%partons(j)%p%t) - parton_mass_squared &
                     (shower%partons(j)%p) < &
                     abs(shower%partons(j + 1)%p%t) - parton_mass_squared &
                     (shower%partons(j + 1)%p)) then
                   call swap_pointers (shower%partons(j), shower%partons(j + 1))
                   changed = .true.
                else
                   if (nearly_equal(abs (shower%partons(j)%p%t) - &
                        parton_mass_squared (shower%partons(j)%p), &
                        abs(shower%partons(j + 1)%p%t) - parton_mass_squared &
                        (shower%partons(j + 1)%p))) then
                      if (shower%partons(j)%p%nr > &
                           shower%partons(j + 1)%p%nr) then
                         call swap_pointers (shower%partons(j), &
                              shower%partons(j + 1))
                         changed = .true.
                      end if
                   end if
                end if
             end if
          end do INNERDO_T
          if (.not. changed) exit OUTERDO_T
       end do OUTERDO_T
    end if
    if (signal_is_pending ()) return
    if (D_PRINT) print *, "D: shower_sort_partons: finished"
  end subroutine shower_sort_partons

\eatline
\nwendcode{}\nwbegindocs{201}Creation and finalization.
\nwenddocs{}\nwbegincode{202}\moddef{Shower core: shower: TBP}\plusendmoddef
  procedure :: create => shower_create
\nwendcode{}\nwbegincode{203}\moddef{Shower core: procedures}\plusendmoddef
  subroutine shower_create (shower, xmin, xmax, qmin, qmax, pdf)
    class(shower_t), intent(inout), target :: shower
    real(double), intent(in) :: xmin, xmax, qmin, qmax
    type(lhapdf_pdf_t), intent(inout), target, optional :: pdf
    shower%next_free_nr = 1
    shower%next_color_nr = 1
    if (ASSERT) then
       if (allocated (shower%interactions)) then
          call msg_bug ("Shower: creating new shower while old one " // &
               "is still associated (interactions)")
       end if
       if (allocated (shower%partons)) then
          call msg_bug ("Shower: creating new shower while old one " // &
               "is still associated (partons)")
       end if
    end if
    treat_light_quarks_massless = .true.
    treat_duscb_quarks_massless = .false.
    shower%valid = .true.
    shower%xmin = xmin
    shower%xmax = xmax
    shower%qmin = qmin
    shower%qmax = qmax
    if (shower%pdf_type == STRF_LHAPDF6 .and. present (pdf)) then
       shower%pdf => pdf
       call lhapdf_transfer_pointer (pdf, shower%pdf)
    end if
  end subroutine shower_create

\eatline
\nwendcode{}\nwbegindocs{204}Deallocate the interaction pointers.
\nwenddocs{}\nwbegincode{205}\moddef{Shower core: shower: TBP}\plusendmoddef
  procedure :: final => shower_final
\nwendcode{}\nwbegincode{206}\moddef{Shower core: procedures}\plusendmoddef
  subroutine shower_final (shower)
    class(shower_t), intent(inout) :: shower
    integer :: i
    if (.not. allocated (shower%interactions))  return
    do i = 1, size (shower%interactions)
       if (allocated (shower%interactions(i)%i%partons)) &
            deallocate (shower%interactions(i)%i%partons)
       deallocate (shower%interactions(i)%i)
    end do
    deallocate (shower%interactions)
    deallocate (shower%partons)
    if (associated (shower%pdf)) then
       shower%pdf => null ()
    end if
  end subroutine shower_final

\eatline
\nwendcode{}\nwbegindocs{207}Bookkeeping functions.
\nwenddocs{}\nwbegincode{208}\moddef{Shower core: shower: TBP}\plusendmoddef
  procedure :: get_next_free_nr => shower_get_next_free_nr
\nwendcode{}\nwbegincode{209}\moddef{Shower core: procedures}\plusendmoddef
  function shower_get_next_free_nr (shower) result (next_number)
    class(shower_t), intent(inout) :: shower
    integer :: next_number
    next_number = shower%next_free_nr
    shower%next_free_nr = shower%next_free_nr + 1
  end function shower_get_next_free_nr

\eatline
\nwendcode{}\nwbegindocs{210}\nwdocspar
\nwenddocs{}\nwbegincode{211}\moddef{Shower core: shower: TBP}\plusendmoddef
  procedure :: set_next_color_nr => shower_set_next_color_nr
\nwendcode{}\nwbegincode{212}\moddef{Shower core: procedures}\plusendmoddef
  subroutine shower_set_next_color_nr (shower, index)
    class(shower_t), intent(inout) :: shower
    integer, intent(in) :: index
    if (index < shower%next_color_nr) then
       call msg_error ("in shower_set_next_color_nr")
    else
       shower%next_color_nr = max(shower%next_color_nr, index)
    end if
  end subroutine shower_set_next_color_nr

\eatline
\nwendcode{}\nwbegindocs{213}\nwdocspar
\nwenddocs{}\nwbegincode{214}\moddef{Shower core: shower: TBP}\plusendmoddef
  procedure :: get_next_color_nr => shower_get_next_color_nr
\nwendcode{}\nwbegincode{215}\moddef{Shower core: procedures}\plusendmoddef
  function shower_get_next_color_nr (shower) result (next_color)
    class(shower_t), intent(inout) :: shower
    integer :: next_color
    next_color = shower%next_color_nr
    shower%next_color_nr = shower%next_color_nr + 1
  end function shower_get_next_color_nr

\eatline
\nwendcode{}\nwbegindocs{216}\nwdocspar
\nwenddocs{}\nwbegincode{217}\moddef{Shower core: procedures}\plusendmoddef
  subroutine shower_enlarge_partons_array (shower, custom_length)
    type(shower_t), intent(inout) :: shower
    integer, intent(in), optional :: custom_length
    integer :: i, length, oldlength
    type(parton_pointer_t), dimension(:), allocatable :: tmp_partons
    if (D_PRINT) print *, "D: shower_enlarge_partons_array"
    if (present(custom_length)) then
       length = custom_length
    else
       length = 10
    end if
    if (ASSERT) then
       if (length < 1) then
          call msg_bug ("Shower: no parton_pointers added in shower%partons")
       end if
    end if
    if (allocated (shower%partons)) then
       oldlength = size (shower%partons)
       allocate (tmp_partons(1:oldlength))
       do i = 1, oldlength
          tmp_partons(i)%p => shower%partons(i)%p
       end do
       deallocate (shower%partons)
    else
       oldlength = 0
    end if
    allocate (shower%partons(1:oldlength + length))
    do i = 1, oldlength
       shower%partons(i)%p => tmp_partons(i)%p
    end do
    do i = oldlength + 1, oldlength + length
       shower%partons(i)%p => null()
    end do
  end subroutine shower_enlarge_partons_array

\eatline
\nwendcode{}\nwbegindocs{218}\nwdocspar
\nwenddocs{}\nwbegincode{219}\moddef{Shower core: shower: TBP}\plusendmoddef
  procedure :: add_child => shower_add_child
\nwendcode{}\nwbegincode{220}\moddef{Shower core: procedures}\plusendmoddef
  subroutine shower_add_child (shower, prt, child)
    class(shower_t), intent(inout) :: shower
    type(parton_t), pointer :: prt
    integer, intent(in) :: child
    integer :: i, lastfree
    type(parton_pointer_t) :: newprt
    if (child /= 1 .and. child /= 2) then
       call msg_bug ("Shower: Adding child in nonexisting place")
    end if
    allocate (newprt%p)
    newprt%p%nr = shower%get_next_free_nr ()
    !!! add new parton as child
    if (child == 1) then
       prt%child1 => newprt%p
    else
       prt%child2 => newprt%p
    end if
    newprt%p%parent => prt
    newprt%p%interactionnr = prt%interactionnr
    !!! add new parton to shower%partons list
    if (associated (shower%partons (size(shower%partons))%p)) then
       call shower_enlarge_partons_array (shower)
    end if
    !!! find last free pointer and let it point to the new parton
    lastfree = 0
    do i = size (shower%partons), 1, -1
       if (.not. associated (shower%partons(i)%p)) then
          lastfree = i
       end if
    end do
    if (lastfree == 0) then
       call msg_bug ("Shower: no free pointers found")
    end if
    shower%partons(lastfree)%p => newprt%p
  end subroutine shower_add_child

\eatline
\nwendcode{}\nwbegindocs{221}\nwdocspar
\nwenddocs{}\nwbegincode{222}\moddef{Shower core: shower: TBP}\plusendmoddef
  procedure :: add_parent => shower_add_parent
\nwendcode{}\nwbegincode{223}\moddef{Shower core: procedures}\plusendmoddef
  subroutine shower_add_parent (shower, prt)
    class(shower_t), intent(inout) :: shower
    type(parton_t), intent(inout), target :: prt
    integer :: i, lastfree
    type(parton_pointer_t) :: newprt
    if (D_PRINT) print *, "D: shower_add_parent: for parton ", prt%nr
    allocate (newprt%p)
    newprt%p%nr = shower%get_next_free_nr ()
    !!! add new parton as parent
    newprt%p%child1 => prt
    prt%parent => newprt%p
    newprt%p%interactionnr = prt%interactionnr
    !!! add new parton to shower%partons list
    if (.not. allocated (shower%partons) .or. &
         associated (shower%partons(size(shower%partons))%p)) then
       call shower_enlarge_partons_array (shower)
    end if
    !!! find last free pointer and let it point to the new parton
    lastfree = 0
    do i = size(shower%partons), 1, -1
       if (.not. associated (shower%partons(i)%p)) then
          lastfree = i
       end if
    end do
    if (lastfree == 0) then
       call msg_bug ("Shower: no free pointers found")
    end if
    shower%partons(lastfree)%p => newprt%p
  end subroutine shower_add_parent

\eatline
\nwendcode{}\nwbegindocs{224}For debugging:
\nwenddocs{}\nwbegincode{225}\moddef{Shower core: procedures}\plusendmoddef
  pure function shower_get_total_momentum (shower) result (mom)
    type(shower_t), intent(in) :: shower
    type(vector4_t) :: mom
    integer :: i
    if (.not. allocated (shower%partons)) return
    mom = vector4_null
    do i = 1, size (shower%partons)
       if (.not. associated (shower%partons(i)%p)) cycle
       if (parton_is_final (shower%partons(i)%p)) then
          mom = mom + shower%partons(i)%p%momentum
       end if
    end do
  end function shower_get_total_momentum

\eatline
\nwendcode{}\nwbegindocs{226}Count the number of partons by going through {\tt{}shower{\char37}partons}
whereby you can require a minimum energy {\tt{}mine} and specify whether to
{\tt{}include{\char95}remnants}, which is done if not given.
\nwenddocs{}\nwbegincode{227}\moddef{Shower core: procedures}\plusendmoddef
  function shower_get_nr_of_partons (shower, mine, include_remnants) &
       result (nr)
    type(shower_t), intent(in) :: shower
    real(default), intent(in), optional :: mine
    logical, intent(in), optional :: include_remnants
    integer :: nr
    integer :: i
    type(parton_t), pointer :: prt
    nr = 0
    do i = 1, size (shower%partons)
       prt => shower%partons(i)%p
       if (.not. associated (prt)) cycle
       if (.not. parton_is_final (prt)) cycle
       if (prt%type == BEAM_REMNANT) then
          if (present (include_remnants)) then
             if (.not. include_remnants) cycle
          end if
       end if
       if (present(mine)) then
          if (prt%momentum%p(0) < mine) cycle
       end if
       nr = nr + 1
    end do
  end function shower_get_nr_of_partons

\eatline
\nwendcode{}\nwbegindocs{228}\nwdocspar
\nwenddocs{}\nwbegincode{229}\moddef{Shower core: procedures}\plusendmoddef
  function shower_get_nr_of_final_colored_ME_partons (shower) result (nr)
    type(shower_t), intent(in) :: shower
    integer :: nr
    integer :: i, j
    type(parton_t), pointer :: prt
    nr = 0
    do i = 1, size (shower%interactions)
       do j = 1, size (shower%interactions(i)%i%partons)
          prt => shower%interactions(i)%i%partons(j)%p
          if (.not. associated (prt)) cycle
          if (.not. parton_is_colored (prt)) cycle
          if (prt%belongstointeraction .and. prt%belongstoFSR .and. &
               (prt%type /= INTERNAL)) then
             nr = nr +1
          end if
       end do
    end do
  end function shower_get_nr_of_final_colored_ME_partons

\eatline
\nwendcode{}\nwbegindocs{230}\nwdocspar
\nwenddocs{}\nwbegincode{231}\moddef{Shower core: shower: TBP}\plusendmoddef
  procedure :: get_final_colored_ME_partons => &
       shower_get_final_colored_ME_partons
\nwendcode{}\nwbegincode{232}\moddef{Shower core: procedures}\plusendmoddef
  subroutine shower_get_final_colored_ME_partons (shower, partons)
    class(shower_t), intent(in) :: shower
    type(parton_pointer_t), dimension(:), allocatable, intent(inout) :: partons
    integer :: i, j, index, s
    type(parton_t), pointer :: prt
    if (allocated(partons))  deallocate(partons)
    s = shower_get_nr_of_final_colored_ME_partons (shower)
    if (s == 0) return
    allocate (partons(1:s))
    index = 0
    do i = 1, size (shower%interactions)
       do j = 1, size (shower%interactions(i)%i%partons)
          prt => shower%interactions(i)%i%partons(j)%p
          if (.not. associated (prt)) cycle
          if (.not. parton_is_colored (prt)) cycle
          if (prt%belongstointeraction .and. prt%belongstoFSR .and. &
               (prt%type /= INTERNAL)) then
             index = index + 1
             partons(index)%p => prt
          end if
       end do
    end do
  end subroutine shower_get_final_colored_ME_partons

\eatline
\nwendcode{}\nwbegindocs{233}\nwdocspar
\nwenddocs{}\nwbegincode{234}\moddef{Shower core: procedures}\plusendmoddef
  recursive function interaction_fsr_is_finished_for_parton(prt) result (finished)
    type(parton_t), intent(in) :: prt
    logical :: finished
    if (prt%belongstoFSR) then
       !!! FSR partons
       if (associated (prt%child1)) then
          finished = interaction_fsr_is_finished_for_parton (prt%child1) &
               .and. interaction_fsr_is_finished_for_parton (prt%child2)
       else
          finished = prt%t <= parton_mass_squared(prt)
       end if
    else
       !!! search for emitted timelike partons in ISR shower
       if (.not. associated (prt%initial)) then
          !!! no inital -> no ISR
          finished = .true.
       else if (.not. associated (prt%parent)) then
          finished = .false.
       else
          if (.not. parton_is_hadron (prt%parent)) then
             if (associated (prt%child2)) then
                finished = interaction_fsr_is_finished_for_parton (prt%parent) .and. &
                     interaction_fsr_is_finished_for_parton (prt%child2)
             else
                finished = interaction_fsr_is_finished_for_parton (prt%parent)
             end if
          else
             if (associated (prt%child2)) then
                finished = interaction_fsr_is_finished_for_parton (prt%child2)
             else
                !!! only second partons can come here -> if that happens FSR
                !!!     evolution is not existing
                finished = .true.
             end if
          end if
       end if
    end if
  end function interaction_fsr_is_finished_for_parton

\eatline
\nwendcode{}\nwbegindocs{235}\nwdocspar
\nwenddocs{}\nwbegincode{236}\moddef{Shower core: procedures}\plusendmoddef
  function interaction_fsr_is_finished (interaction) result (finished)
    type(shower_interaction_t), intent(in) :: interaction
    logical :: finished
    integer :: i
    finished = .true.
    if (.not. allocated (interaction%partons)) return
    do i = 1, size (interaction%partons)
       if (.not. interaction_fsr_is_finished_for_parton &
            (interaction%partons(i)%p)) then
          finished = .false.
          exit
       end if
    end do
  end function interaction_fsr_is_finished

\eatline
\nwendcode{}\nwbegindocs{237}\nwdocspar
\nwenddocs{}\nwbegincode{238}\moddef{Shower core: public}\plusendmoddef
  public :: shower_interaction_get_s
\nwendcode{}\nwbegincode{239}\moddef{Shower core: procedures}\plusendmoddef
  function shower_interaction_get_s (interaction) result (s)
    type(shower_interaction_t), intent(in) :: interaction
    real(default) :: s
    s = (interaction%partons(1)%p%initial%momentum + &
         interaction%partons(2)%p%initial%momentum)**2
  end function shower_interaction_get_s

\eatline
\nwendcode{}\nwbegindocs{240}\nwdocspar
\nwenddocs{}\nwbegincode{241}\moddef{Shower core: procedures}\plusendmoddef
  function shower_fsr_is_finished (shower) result (finished)
    type(shower_t), intent(in) :: shower
    logical :: finished
    integer :: i
    finished = .true.
    if (.not.allocated (shower%interactions)) return
    do i = 1, size(shower%interactions)
       if (.not. interaction_fsr_is_finished (shower%interactions(i)%i)) then
          finished = .false.
          exit
       end if
    end do
  end function shower_fsr_is_finished

\eatline
\nwendcode{}\nwbegindocs{242}\nwdocspar
\nwenddocs{}\nwbegincode{243}\moddef{Shower core: procedures}\plusendmoddef
  function shower_isr_is_finished (shower) result (finished)
    type(shower_t), intent(in) :: shower
    logical :: finished
    integer :: i
    type(parton_t), pointer :: prt
    finished = .true.
    if (.not.allocated (shower%partons)) return
    do i = 1, size (shower%partons)
       if (.not. associated (shower%partons(i)%p)) cycle
       prt => shower%partons(i)%p
       if (isr_pt_ordered) then
          if (.not. prt%belongstoFSR .and. .not. prt%simulated &
              .and. prt%scale > zero) then
             finished = .false.
             exit
          end if
       else
          if (.not. prt%belongstoFSR .and. .not. prt%simulated &
              .and. prt%t < zero) then
             finished = .false.
             exit
          end if
       end if
    end do
  end function shower_isr_is_finished

\eatline
\nwendcode{}\nwbegindocs{244}\nwdocspar
\nwenddocs{}\nwbegincode{245}\moddef{Shower core: procedures}\plusendmoddef
  subroutine interaction_find_partons_nearest_to_hadron (interaction, prt1, prt2)
    type(shower_interaction_t), intent(in) :: interaction
    type(parton_t), pointer :: prt1, prt2
    prt1 => null ()
    prt2 => null ()
    prt1 => interaction%partons(1)%p
    do
       if (associated (prt1%parent)) then
          if (parton_is_hadron (prt1%parent)) then
             exit
          else if ((.not. isr_pt_ordered .and. .not. prt1%parent%simulated) &
               .or. (isr_pt_ordered .and. .not. prt1%simulated)) then
             exit
          else
             prt1 => prt1%parent
          end if
       else
          exit
       end if
    end do
    prt2 => interaction%partons(2)%p
    do
       if (associated (prt2%parent)) then
          if (parton_is_hadron (prt2%parent)) then
             exit
          else if ((.not. isr_pt_ordered .and. .not. prt2%parent%simulated) &
               .or. (isr_pt_ordered .and. .not. prt2%simulated)) then
             exit
          else
             prt2 => prt2%parent
          end if
       else
          exit
       end if
    end do
  end subroutine interaction_find_partons_nearest_to_hadron

\eatline
\nwendcode{}\nwbegindocs{246}\nwdocspar
\nwenddocs{}\nwbegincode{247}\moddef{Shower core: shower: TBP}\plusendmoddef
  procedure :: update_beamremnants => shower_update_beamremnants
\nwendcode{}\nwbegincode{248}\moddef{Shower core: procedures}\plusendmoddef
  subroutine shower_update_beamremnants (shower)
    class(shower_t), intent(inout) :: shower
    type(parton_t), pointer :: hadron, remnant
    integer :: i
    real(default) :: random
    !!! only proton in first interaction !!?
    !!! currently only first beam-remnant will be updated
    do i = 1,2
       if (associated (shower%interactions(1)%i%partons(i)%p%initial)) then
          hadron => shower%interactions(1)%i%partons(i)%p%initial
       else
          cycle
       end if
       remnant => hadron%child2
       if (associated (remnant)) then
          remnant%momentum = hadron%momentum - hadron%child1%momentum
       end if
       !!! generate flavor of the beam-remnant if beam was proton
       if (abs (hadron%type) == PROTON .and. associated (hadron%child1)) then
          if (parton_is_quark (hadron%child1)) then
             !!! decide if valence (u,d) or sea quark (s,c,b)
             if ((abs (hadron%child1%type) <= 2) .and. &
                  (hadron%type * hadron%child1%type > zero)) then
                !!! valence quark
                if (abs (hadron%child1%type) == 1) then
                   !!! if d then remaining diquark is uu_1
                   remnant%type = sign (UU1, hadron%type)
                else
                   call rng%generate (random)
                   !!! if u then remaining diquark is ud_0 or ud_1
                   if (random < 0.75_default) then
                      remnant%type = sign (UD0, hadron%type)
                   else
                      remnant%type = sign (UD1, hadron%type)
                   end if
                end if
                remnant%c1 = hadron%child1%c2
                remnant%c2 = hadron%child1%c1
             else if ((hadron%type * hadron%child1%type) < zero) then
                !!! antiquark
                if (.not. associated (remnant%child1)) then
                   call shower%add_child (remnant, 1)
                end if
                if (.not. associated (remnant%child2)) then
                   call shower%add_child (remnant, 2)
                end if
                call rng%generate (random)
                if (random < 0.6666_default) then
                   !!! 2/3 into udq + u
                   if (abs (hadron%child1%type) == 1) then
                      remnant%child1%type = sign (NEUTRON, hadron%type)
                   else if (abs (hadron%child1%type) == 2) then
                      remnant%child1%type = sign (PROTON, hadron%type)
                   else if (abs (hadron%child1%type) == 3) then
                      remnant%child1%type = sign (SIGMA0, hadron%type)
                   else if (abs (hadron%child1%type) == 4) then
                      remnant%child1%type = sign (SIGMACPLUS, hadron%type)
                   else if (abs (hadron%child1%type) == 5) then
                      remnant%child1%type = sign (SIGMAB0, hadron%type)
                   end if
                   remnant%child2%type = sign (2, hadron%type)
                else
                   !!! 1/3 into uuq + d
                   if (abs (hadron%child1%type) == 1) then
                      remnant%child1%type = sign (PROTON, hadron%type)
                   else if (abs (hadron%child1%type) == 2) then
                      remnant%child1%type = sign (DELTAPLUSPLUS, hadron%type)
                   else if (abs (hadron%child1%type) == 3) then
                      remnant%child1%type = sign (SIGMAPLUS, hadron%type)
                   else if (abs (hadron%child1%type) == 4) then
                      remnant%child1%type = sign (SIGMACPLUSPLUS, hadron%type)
                   else if (abs (hadron%child1%type) == 5) then
                      remnant%child1%type = sign (SIGMABPLUS, hadron%type)
                   end if
                   remnant%child2%type = sign (1, hadron%type)
                end if
                remnant%c1 = hadron%child1%c2
                remnant%c2 = hadron%child1%c1
                remnant%child1%c1 = 0
                remnant%child1%c2 = 0
                remnant%child2%c1 = remnant%c1
                remnant%child2%c2 = remnant%c2
             else
                !!! sea quark
                if (.not. associated (remnant%child1)) then
                   call shower%add_child (remnant, 1)
                end if
                if (.not. associated (remnant%child2)) then
                   call shower%add_child (remnant, 2)
                end if
                call rng%generate (random)
                if (random < 0.5_default) then
                   !!! 1/2 into usbar + ud_0
                   if (abs (hadron%child1%type) == 3) then
                      remnant%child1%type = sign (KPLUS, hadron%type)
                   else if (abs (hadron%child1%type) == 4) then
                      remnant%child1%type = sign (D0, hadron%type)
                   else if (abs (hadron%child1%type) == 5) then
                      remnant%child1%type = sign (BPLUS, hadron%type)
                   end if
                   remnant%child2%type = sign (UD0, hadron%type)
                else if (random < 0.6666_default) then
                   !!! 1/6 into usbar + ud_1
                   if (abs (hadron%child1%type) == 3) then
                      remnant%child1%type = sign (KPLUS, hadron%type)
                   else if (abs (hadron%child1%type) == 4) then
                      remnant%child1%type = sign (D0, hadron%type)
                   else if (abs (hadron%child1%type) == 5) then
                      remnant%child1%type = sign (BPLUS, hadron%type)
                   end if
                   remnant%child2%type = sign (UD1, hadron%type)
                else
                   !!! 1/3 into dsbar + uu_1
                   if (abs (hadron%child1%type) == 3) then
                      remnant%child1%type = sign (K0, hadron%type)
                   else if (abs (hadron%child1%type) == 4) then
                      remnant%child1%type = sign (DPLUS, hadron%type)
                   else if (abs (hadron%child1%type) == 5) then
                      remnant%child1%type = sign (B0, hadron%type)
                   end if
                   remnant%child2%type = sign (UU1, hadron%type)
                end if
                remnant%c1 = hadron%child1%c2
                remnant%c2 = hadron%child1%c1
                remnant%child2%c1 = remnant%c1
                remnant%child2%c2 = remnant%c2
             end if
          else if (parton_is_gluon (hadron%child1)) then
             if (.not.associated (remnant%child1)) then
                call shower%add_child (remnant, 1)
             end if
             if (.not.associated (remnant%child2)) then
                call shower%add_child (remnant, 2)
             end if
             call rng%generate (random)
             if (random < 0.5_default) then
                !!! 1/2 into u + ud_0
                remnant%child1%type = sign (2, hadron%type)
                remnant%child2%type = sign (UD0, hadron%type)
             else if (random < 0.6666_default) then
                !!! 1/6 into u + ud_1
                remnant%child1%type = sign (2, hadron%type)
                remnant%child2%type = sign (UD1, hadron%type)
             else
                !!! 1/3 into d + uu_1
                remnant%child1%type = sign (1, hadron%type)
                remnant%child2%type = sign (UU1, hadron%type)
             end if
             remnant%c1 = hadron%child1%c2
             remnant%c2 = hadron%child1%c1
             if (sign (1,hadron%type) > 0) then
                remnant%child1%c1 = remnant%c1
                remnant%child2%c2 = remnant%c2
             else
                remnant%child1%c2 = remnant%c2
                remnant%child2%c1 = remnant%c1
             end if
          end if
          if (associated (remnant%child1)) then
             !!! don't care about on-shellness for now
             remnant%child1%momentum = 0.5_default * remnant%momentum
             remnant%child2%momentum = 0.5_default * remnant%momentum
             !!! but care about on-shellness for baryons
             if (mod (remnant%child1%type, 100) >= 10) then
                !!! check if the third quark is set -> meson or baryon
                remnant%child1%t = parton_mass_squared(remnant%child1)
                remnant%child1%momentum = [remnant%child1%momentum%p(0), &
                     (remnant%child1%momentum%p(1:3) / &
                      remnant%child1%momentum%p(1:3)**1) * &
                     sqrt (remnant%child1%momentum%p(0)**2 - remnant%child1%t)]
                remnant%child2%momentum = remnant%momentum &
                     - remnant%child1%momentum
             end if
          end if
       end if
    end do
  end subroutine shower_update_beamremnants

\eatline
\nwendcode{}\nwbegindocs{249}\nwdocspar
\nwenddocs{}\nwbegincode{250}\moddef{Shower core: procedures}\plusendmoddef
  subroutine interaction_apply_lorentztrafo (interaction, L)
    type(shower_interaction_t), intent(inout) :: interaction
    type(lorentz_transformation_t), intent(in) :: L
    type(parton_t), pointer :: prt
    integer :: i
    !!! ISR part
    do i = 1,2
       prt => interaction%partons(i)%p
       !!! loop over ancestors
       MOTHERS: do
          !!! boost parton
          call parton_apply_lorentztrafo (prt, L)
          if (associated (prt%child2)) then
             !!! boost emitted timelike parton (and daughters)
             call parton_apply_lorentztrafo_recursive (prt%child2, L)
          end if
          if (associated (prt%parent)) then
             if (.not. parton_is_hadron (prt%parent)) then
                prt => prt%parent
             else
                exit
             end if
          else
             exit
          end if
       end do MOTHERS
    end do
    !!! FSR part
    if (associated (interaction%partons(3)%p%parent)) then
       !!! pseudo Parton-Shower histora has been generated -> find
       !!!     mother and go on from there recursively
       prt => interaction%partons(3)%p
       do while (associated (prt%parent))
          prt => prt%parent
       end do
       call parton_apply_lorentztrafo_recursive (prt, L)
    else
       do i = 3, size (interaction%partons)
          call parton_apply_lorentztrafo (interaction%partons(i)%p, L)
       end do
    end if
  end subroutine interaction_apply_lorentztrafo

\nwendcode{}\nwbegindocs{251} %def interaction_apply_lorentztrafo
\nwenddocs{}\nwbegindocs{252}\nwdocspar
\nwenddocs{}\nwbegincode{253}\moddef{Shower core: procedures}\plusendmoddef
  subroutine shower_apply_lorentztrafo (shower, L)
    type(shower_t), intent(inout) :: shower
    type(lorentz_transformation_t), intent(in) :: L
    integer :: i
    do i = 1, size (shower%interactions)
       call interaction_apply_lorentztrafo (shower%interactions(i)%i, L)
    end do
  end subroutine shower_apply_lorentztrafo

\eatline
\nwendcode{}\nwbegindocs{254}This boosts partons belonging to the interaction to the
center-of-mass frame of its partons nearest to the hadron.
\nwenddocs{}\nwbegincode{255}\moddef{Shower core: procedures}\plusendmoddef
  subroutine interaction_boost_to_CMframe (interaction)
    type(shower_interaction_t), intent(inout) :: interaction
    type(vector4_t) :: beta
    type(parton_t), pointer :: prt1, prt2
    call interaction_find_partons_nearest_to_hadron (interaction, prt1, prt2)
    beta = prt1%momentum + prt2%momentum
    beta = beta / beta%p(0)
    if (ASSERT) then
       if (beta**2 > one) then
          call msg_error ("Shower: boost to CM frame: beta > 1")
          return
       end if
    end if
    if (space_part(beta)**2 > tiny_13) then
       call interaction_apply_lorentztrafo(interaction, &
            boost(space_part(beta)**1 / &
              sqrt (one - space_part(beta)**2), -direction(beta)))
    end if
  end subroutine interaction_boost_to_CMframe

\eatline
\nwendcode{}\nwbegindocs{256}This boosts every interaction to the center-of-mass-frame of its
partons nearest to the hadron.
\nwenddocs{}\nwbegincode{257}\moddef{Shower core: procedures}\plusendmoddef
  subroutine shower_boost_to_CMframe (shower)
    type(shower_t), intent(inout) :: shower
    integer :: i
    do i = 1, size (shower%interactions)
       call interaction_boost_to_CMframe (shower%interactions(i)%i)
    end do
    call shower%update_beamremnants ()
  end subroutine shower_boost_to_CMframe

\eatline
\nwendcode{}\nwbegindocs{258}This boost all partons so that initial partons have their assigned
$x$-value.
\nwenddocs{}\nwbegincode{259}\moddef{Shower core: shower: TBP}\plusendmoddef
  procedure :: boost_to_labframe => shower_boost_to_labframe
\nwendcode{}\nwbegincode{260}\moddef{Shower core: procedures}\plusendmoddef
  subroutine shower_boost_to_labframe (shower)
    class(shower_t), intent(inout) :: shower
    integer :: i
    do i = 1, size (shower%interactions)
       call interaction_boost_to_labframe (shower%interactions(i)%i)
    end do
  end subroutine shower_boost_to_labframe

\eatline
\nwendcode{}\nwbegindocs{261}This boosts all partons so that initial partons have their
assigned $x$-value.
\nwenddocs{}\nwbegincode{262}\moddef{Shower core: procedures}\plusendmoddef
  subroutine interaction_boost_to_labframe (interaction)
    type(shower_interaction_t), intent(inout) :: interaction
    type(parton_t), pointer :: prt1, prt2
    type(vector3_t) :: beta
    call interaction_find_partons_nearest_to_hadron (interaction, prt1, prt2)
    if (.not. associated (prt1%initial) .or. .not. associated (prt2%initial)) then
       return
    end if
    !!! transform partons to overall labframe.
    beta = vector3_canonical(3) * &
         ((prt1%x * prt2%momentum%p(0) - &
           prt2%x * prt1%momentum%p(0)) / &
          (prt1%x * prt2%momentum%p(3) - &
           prt2%x * prt1%momentum%p(3)))
    if (beta**1 > tiny_10) &
         call interaction_apply_lorentztrafo (interaction, &
             boost (beta**1 / sqrt(one - beta**2), -direction(beta)))
  end subroutine interaction_boost_to_labframe

\eatline
\nwendcode{}\nwbegindocs{263}Only rotate to z if inital hadrons are given (and they are assumed
to be aligned along the z-axis).
\nwenddocs{}\nwbegincode{264}\moddef{Shower core: procedures}\plusendmoddef
  subroutine interaction_rotate_to_z (interaction)
    type(shower_interaction_t), intent(inout) :: interaction
    type(parton_t), pointer :: prt1, prt2
    call interaction_find_partons_nearest_to_hadron (interaction, prt1, prt2)
    if (associated (prt1%initial)) then
       call interaction_apply_lorentztrafo (interaction, &
            rotation_to_2nd (space_part (prt1%momentum), &
            vector3_canonical(3) * sign (one, &
            prt1%initial%momentum%p(3))))
    end if
  end subroutine interaction_rotate_to_z

\eatline
\nwendcode{}\nwbegindocs{265}Rotate initial partons to lie along $\pm z$ axis.
\nwenddocs{}\nwbegincode{266}\moddef{Shower core: procedures}\plusendmoddef
  subroutine shower_rotate_to_z (shower)
    type(shower_t), intent(inout) :: shower
    integer :: i
    do i = 1, size (shower%interactions)
       call interaction_rotate_to_z (shower%interactions(i)%i)
    end do
    call shower%update_beamremnants ()
  end subroutine shower_rotate_to_z

\eatline
\nwendcode{}\nwbegindocs{267}\nwdocspar
\nwenddocs{}\nwbegincode{268}\moddef{Shower core: procedures}\plusendmoddef
  subroutine interaction_generate_primordial_kt (interaction)
    type(shower_interaction_t), intent(inout) :: interaction
    type(parton_t), pointer :: had1, had2
    type(vector4_t) :: momenta(2)
    type(vector3_t) :: beta
    real(default) :: pt (2), phi(2)
    real(default) :: shat
    real(default) :: btheta, bphi
    integer :: i
    if (vanishes(primordial_kt_width))  return
    !!! Return if there are no initials, electron-hadron collision not implemented
    if (.not. associated (interaction%partons(1)%p%initial) .or. &
        .not. associated (interaction%partons(2)%p%initial)) then
       return
    end if
    had1 => interaction%partons(1)%p%initial
    had2 => interaction%partons(2)%p%initial
    !!! copy momenta and energy
    momenta(1) = had1%child1%momentum
    momenta(2) = had2%child1%momentum
    GENERATE_PT_PHI: do i = 1, 2
       !!! generate transverse momentum and phi
       GENERATE_PT: do
          call rng%generate (pt (i))
          pt(i) = primordial_kt_width * sqrt(-log(pt(i)))
          if (pt(i) < primordial_kt_cutoff) exit
       end do GENERATE_PT
       call rng%generate (phi (i))
       phi(i) = twopi * phi(i)
    end do GENERATE_PT_PHI
    !!! adjust momenta
    shat = (momenta(1) + momenta(2))**2
    momenta(1) = [momenta(1)%p(0), &
                  pt(1) * cos(phi(1)), &
                  pt(1) * sin(phi(1)), &
                  momenta(1)%p(3)]
    momenta(2) = [momenta(2)%p(0), &
                  pt(2) * cos(phi(2)), &
                  pt(2) * sin(phi(2)), &
                  momenta(2)%p(3)]
    beta = [momenta(1)%p(1) + momenta(2)%p(1), &
            momenta(1)%p(2) + momenta(2)%p(2), zero] / sqrt(shat)
    momenta(1) = boost (beta**1 / sqrt(one - beta**2), -direction(beta)) * momenta(1)
    bphi = azimuthal_angle (momenta(1))
    btheta = polar_angle (momenta(1))
    call interaction_apply_lorentztrafo (interaction, &
         rotation (cos(bphi), sin(bphi), 3) * rotation(cos(btheta), &
         sin(btheta), 2) * rotation(cos(-bphi), sin(-bphi), 3))
    call interaction_apply_lorentztrafo (interaction, &
         boost (beta**1 / sqrt(one - beta**2), -direction(beta)))
  end subroutine interaction_generate_primordial_kt

\eatline
\nwendcode{}\nwbegindocs{269}\nwdocspar
\nwenddocs{}\nwbegincode{270}\moddef{Shower core: shower: TBP}\plusendmoddef
  procedure :: generate_primordial_kt => shower_generate_primordial_kt
\nwendcode{}\nwbegincode{271}\moddef{Shower core: procedures}\plusendmoddef
  subroutine shower_generate_primordial_kt (shower)
    class(shower_t), intent(inout) :: shower
    integer :: i
    do i = 1, size (shower%interactions)
       call interaction_generate_primordial_kt (shower%interactions(i)%i)
    end do
    call shower%update_beamremnants ()
  end subroutine shower_generate_primordial_kt

\eatline
\nwendcode{}\nwbegindocs{272}Output.
\nwenddocs{}\nwbegincode{273}\moddef{Shower core: procedures}\plusendmoddef
  subroutine interaction_write (interaction, unit)
    type(shower_interaction_t), intent(in) :: interaction
    integer, intent(in), optional :: unit
    integer :: i, u
    u = given_output_unit (unit); if (u < 0) return
    if (associated (interaction%partons(1)%p)) then
       if (associated (interaction%partons(1)%p%initial)) &
            call parton_write (interaction%partons(1)%p%initial, u)
    end if
    if (associated (interaction%partons(2)%p)) then
       if (associated (interaction%partons(2)%p%initial)) &
            call parton_write (interaction%partons(2)%p%initial, u)
    end if
    if (allocated (interaction%partons)) then
       do i = 1, size (interaction%partons)
          call parton_write (interaction%partons(i)%p, u)
       end do
    end if
    write (u, "(A)")
  end subroutine interaction_write

\eatline
\nwendcode{}\nwbegindocs{274}\nwdocspar
\nwenddocs{}\nwbegincode{275}\moddef{Shower core: shower: TBP}\plusendmoddef
  procedure :: write => shower_write
\nwendcode{}\nwbegincode{276}\moddef{Shower core: procedures}\plusendmoddef
  subroutine shower_write (shower, unit)
    class(shower_t), intent(in) :: shower
    integer, intent(in), optional :: unit
    integer :: i, u
    u = given_output_unit (unit); if (u < 0) return
    write (u, "(1x,A)") "------------------------------"
    write (u, "(1x,A)") "WHIZARD internal parton shower"
    write (u, "(1x,A)") "------------------------------"
    write (u, "(3x,A,I0)") "PDF set  = ", shower%pdf_set
    write (u, "(3x,A,I0)") "PDF type = ", shower%pdf_type
    write (u, "(3x,A,L1)") "ISR: angular ordered = ", &
         shower%isr_angular_ordered
    write (u, "(3x,A,ES19.12)") "ISR: maxz_isr = ", shower%maxz_isr
    write (u, "(3x,A,ES19.12)") "ISR: min. energy/timelike emission = ", &
         shower%minenergy_timelike
    write (u, "(3x,A,ES19.12)") "ISR: first scale = ", shower%tscalefactor_isr
    write (u, "(3x,A,ES19.12)") "ISR: 1st integral suppression factor = ", &
         shower%first_integral_suppression_factor
    write (u, "(3x,A,L1)") "ISR: partons only onshell = ", &
         shower%isr_only_onshell_emitted_partons
    if (size (shower%interactions) > 0) then
       write (u, "(3x,A)") "Interactions: "
       do i = 1, size (shower%interactions)
          write (u, "(4x,A,I0)") "Interaction number ", i
          if (.not. associated (shower%interactions(i)%i)) then
             call msg_fatal ("Shower: missing interaction in shower")
          end if
          call interaction_write (shower%interactions(i)%i, u)
       end do
    else
       write (u, "(3x,A)") "[no interactions in shower]"
    end if
    write (u, "(A)")
    if (allocated (shower%partons)) then
       write (u, "(5x,A)")  "Partons:"
       do i = 1, size (shower%partons)
          if (associated (shower%partons(i)%p)) then
             call parton_write(shower%partons(i)%p, u)
             if (i < size (shower%partons)) then
                if (associated (shower%partons(i + 1)%p)) then
                   if (shower%partons(i)%p%belongstointeraction .and. &
                        .not. shower%partons(i + 1)%p%belongstointeraction) then
                      write (u, "(A)") &
                           "-------------------------------------------------------"
                   end if
                end if
             end if
          end if
       end do
    else
       write (u, "(5x,A)")  "[no partons in shower]"
    end if
    write (u, "(4x,A)")  "Total Momentum: "
    call vector4_write (shower_get_total_momentum (shower))
    write (u, "(1x,A,L1)") "ISR finished: ", shower_isr_is_finished (shower)
    write (u, "(1x,A,L1)") "FSR finished: ", shower_fsr_is_finished (shower)
  end subroutine shower_write

\eatline
\nwendcode{}\nwbegindocs{277}\nwdocspar
\nwenddocs{}\nwbegincode{278}\moddef{Shower core: shower: TBP}\plusendmoddef
  procedure :: write_lhef => shower_write_lhef
\nwendcode{}\nwbegincode{279}\moddef{Shower core: procedures}\plusendmoddef
  subroutine shower_write_lhef (shower, unit)
    class(shower_t), intent(in) :: shower
    integer, intent(in), optional :: unit
    integer :: u
    integer :: i
    integer :: c1, c2
    u = given_output_unit (unit);  if (u < 0)  return
    write(u,'(A)') '<LesHouchesEvents version="1.0">'
    write(u,'(A)') '<-- not a complete lhe file - just one event -->'
    write(u,'(A)') '<event>'
    write(u, *) 2 + shower_get_nr_of_partons (shower), 1, 1.0, 1.0, 1.0, 1.0
    !!! write incoming partons
    do i = 1, 2
       if (abs (shower%partons(i)%p%type) < 1000) then
          c1 = 0
          c2 = 0
          if (parton_is_colored (shower%partons(i)%p)) then
             if (shower%partons(i)%p%c1 /= 0) c1 = 500 + shower%partons(i)%p%c1
             if (shower%partons(i)%p%c2 /= 0) c2 = 500 + shower%partons(i)%p%c2
          end if
          write(u,*) shower%partons(i)%p%type, -1, 0, 0, c1, c2, &
               shower%partons(i)%p%momentum%p(1), &
               shower%partons(i)%p%momentum%p(2), &
               shower%partons(i)%p%momentum%p(3), &
               shower%partons(i)%p%momentum%p(0), &
               shower%partons(i)%p%momentum**2, zero, 9.0
       else
          write(u,*) shower%partons(i)%p%type, -9, 0, 0, 0, 0, &
               shower%partons(i)%p%momentum%p(1), &
               shower%partons(i)%p%momentum%p(2), &
               shower%partons(i)%p%momentum%p(3), &
               shower%partons(i)%p%momentum%p(0), &
               shower%partons(i)%p%momentum**2, zero, 9.0
       end if
    end do
    !!! write outgoing partons
    do i = 3, size (shower%partons)
       if (.not. associated (shower%partons(i)%p)) cycle
       if (.not. parton_is_final (shower%partons(i)%p)) cycle
       c1 = 0
       c2 = 0
       if (parton_is_colored (shower%partons(i)%p)) then
          if (shower%partons(i)%p%c1 .ne. 0) c1 = 500 + shower%partons(i)%p%c1
          if (shower%partons(i)%p%c2 .ne. 0) c2 = 500 + shower%partons(i)%p%c2
       end if
       write(u,*) shower%partons(i)%p%type, 1, 1, 2, c1, c2, &
            shower%partons(i)%p%momentum%p(1), &
            shower%partons(i)%p%momentum%p(2), &
            shower%partons(i)%p%momentum%p(3), &
            shower%partons(i)%p%momentum%p(0), &
            shower%partons(i)%p%momentum**2, zero, 9.0
    end do
    write(u,'(A)') '</event>'
    write(u,'(A)') '</LesHouchesEvents>'
  end subroutine shower_write_lhef

\eatline
\nwendcode{}\nwbegindocs{280}\nwdocspar
\nwenddocs{}\nwbegincode{281}\moddef{Shower core: procedures}\plusendmoddef
  subroutine shower_replace_parent_by_hadron (shower, prt)
    type(shower_t), intent(inout) :: shower
    type(parton_t), intent(inout), target :: prt
    type(parton_t), pointer :: remnant => null()
    if (associated (prt%parent)) then
       call shower_remove_parton_from_partons (shower, prt%parent)
       deallocate (prt%parent)
    end if
    if (.not. associated (prt%initial%child2)) then
       call shower%add_child (prt%initial, 2)
    end if
    prt%parent => prt%initial
    prt%parent%child1 => prt
    ! make other child to be a beam-remnant
    remnant => prt%initial%child2
    remnant%type = BEAM_REMNANT
    remnant%momentum = prt%parent%momentum - prt%momentum
    remnant%x = one - prt%x
    remnant%parent => prt%initial
    remnant%t = zero
  end subroutine shower_replace_parent_by_hadron

\eatline
\nwendcode{}\nwbegindocs{282}\nwdocspar
\nwenddocs{}\nwbegincode{283}\moddef{Shower core: procedures}\plusendmoddef
  subroutine shower_get_first_ISR_scale_for_parton (shower, prt, tmax)
    type(shower_t), intent(inout), target :: shower
    type(parton_t), intent(inout), target :: prt
    real(default), intent(in), optional :: tmax
    real(default) :: t, tstep, random, integral, temp1
    real(default) :: temprand
    if (present(tmax)) then
       t = max (max (-shower%tscalefactor_isr * prt%momentum%p(0)**2, &
            -abs(tmax)), prt%t)
    else
       t = max (-shower%tscalefactor_isr * prt%momentum%p(0)**2, prt%t)
    end if
    call rng%generate (random)
    random = -twopi * log(random)
    !!! compare Integral and log(random) instead of random and exp(-Integral)
    random = random / shower%first_integral_suppression_factor
    integral = zero
    call parton_set_simulated (prt, .false.)
    do
       call rng%generate (temprand)
       tstep = max (abs (0.01_default * t) * temprand, 0.1_default * D_Min_t)
       if (t + 0.5_default * tstep > - D_Min_t) then
          prt%t = parton_mass_squared (prt)
          call parton_set_simulated (prt)
          exit
       end if
       prt%t = t + 0.5_default * tstep
       temp1 = integral_over_z_simple (prt, (random - integral) / tstep)
       integral = integral + tstep * temp1
       if (integral > random) then
          prt%t = t + 0.5_default * tstep
          exit
       end if
       t = t + tstep
    end do
    if (prt%t > - D_Min_t) then
       call shower_replace_parent_by_hadron (shower, prt)
    end if

  contains

    function integral_over_z_simple (prt, final) result (integral)
      type(parton_t), intent(inout) :: prt
      real(default), intent(in) :: final
      real(default) :: integral

      real(default), parameter :: zstepfactor = one
      real(default), parameter :: zstepmin = 0.0001_default
      real(default) :: z, zstep, minz, maxz
      real(default) :: pdfsum
      integer :: quark

      integral = zero
      if (D_PRINT) then
         print *, "D: integral_over_z_simple: t = ", prt%t
      end if
      minz = prt%x
      ! maxz = maxzz(shat, s, shower%maxz_isr, shower%minenergy_timelike)
      maxz = shower%maxz_isr
      z = minz
      !!! TODO -> Adapt zstep to structure of divergencies
      if (parton_is_gluon (prt%child1)) then
         !!! gluon coming from g->gg
         do
            call rng%generate (temprand)
            zstep = max(zstepmin, temprand * zstepfactor * z * (one - z))
            zstep = min(zstep, maxz - z)
            integral = integral + zstep * (D_alpha_s_isr ((one - &
                 (z + 0.5_default * zstep)) * abs(prt%t)) / (abs(prt%t))) * &
                 P_ggg (z + 0.5_default * zstep) * shower%get_pdf (prt%initial%type, &
                 prt%x / (z + 0.5_default * zstep), abs(prt%t), GLUON)
            if (integral > final) then
               exit
            end if
            z = z + zstep
            if (z >= maxz) then
               exit
            end if
         end do
         !!! gluon coming from q->qg  ! correctly implemented yet?
         if (integral < final) then
            z = minz
            do
               call rng%generate (temprand)
               zstep = max(zstepmin, temprand * zstepfactor * z * (one - z))
               zstep = min(zstep, maxz - z)
               pdfsum = zero
               do quark = -D_Nf, D_Nf
                  if (quark == 0) cycle
                  pdfsum = pdfsum + shower%get_pdf (prt%initial%type, &
                       prt%x / (z + 0.5_default * zstep), abs(prt%t), quark)
               end do
               integral = integral + zstep * (D_alpha_s_isr &
                    ((z + 0.5_default * zstep) * abs(prt%t)) / (abs(prt%t))) * &
                        P_qqg (one - (z + 0.5_default * zstep)) * pdfsum
               if (integral > final) then
                  exit
               end if
               z = z + zstep
               if (z >= maxz) then
                  exit
               end if
            end do
         end if
      else if (parton_is_quark (prt%child1)) then
         !!! quark coming from q->qg
         do
            call rng%generate(temprand)
            zstep = max(zstepmin, temprand * zstepfactor * z * (one - z))
            zstep = min(zstep, maxz - z)
            integral = integral + zstep * (D_alpha_s_isr ((one - &
                 (z + 0.5_default * zstep)) * abs(prt%t)) / (abs(prt%t))) * &
                 P_qqg (z + 0.5_default * zstep) * shower%get_pdf (prt%initial%type, &
                 prt%x / (z + 0.5_default * zstep), abs(prt%t), prt%type)
            if (integral > final) then
               exit
            end if
            z = z + zstep
            if (z >= maxz) then
               exit
            end if
         end do
         !!! quark coming from g->qqbar
         if (integral < final) then
            z = minz
            do
               call rng%generate (temprand)
               zstep = max(zstepmin, temprand * zstepfactor * z*(one - z))
               zstep = min(zstep, maxz - z)
               integral = integral + zstep * (D_alpha_s_isr &
                    ((one - (z + 0.5_default * zstep)) * abs(prt%t)) / (abs(prt%t))) * &
                    P_gqq (z + 0.5_default * zstep) * shower%get_pdf (prt%initial%type, &
                    prt%x / (z + 0.5_default * zstep), abs(prt%t), GLUON)
               if (integral > final) then
                  exit
               end if
               z = z + zstep
               if (z >= maxz) then
                  exit
               end if
            end do
         end if

      end if
      integral = integral / shower%get_pdf (prt%initial%type, prt%x, &
           abs(prt%t), prt%type)
    end function integral_over_z_simple

  end subroutine shower_get_first_ISR_scale_for_parton

\eatline
\nwendcode{}\nwbegindocs{284}\nwdocspar
\nwenddocs{}\nwbegincode{285}\moddef{Shower core: procedures}\plusendmoddef
  subroutine shower_prepare_for_simulate_isr_pt (shower, interaction)
    type(shower_t), intent(inout) :: shower
    type(shower_interaction_t), intent(inout) :: interaction
    real(default) :: s
    s = (interaction%partons(1)%p%momentum + &
         interaction%partons(2)%p%momentum)**2
    interaction%partons(1)%p%scale = shower%tscalefactor_isr * 0.25_default * s
    interaction%partons(2)%p%scale = shower%tscalefactor_isr * 0.25_default * s
  end subroutine shower_prepare_for_simulate_isr_pt

\eatline
\nwendcode{}\nwbegindocs{286}\nwdocspar
\nwenddocs{}\nwbegincode{287}\moddef{Shower core: procedures}\plusendmoddef
  subroutine shower_prepare_for_simulate_isr_ana_test (shower, prt1, prt2)
    type(shower_t), intent(inout) :: shower
    type(parton_t), intent(inout), target :: prt1, prt2
    type(parton_t), pointer :: prt, prta, prtb
    real(default) ::  scale, factor, E
    integer :: i
    if (.not. associated (prt1%initial) .or. .not. associated (prt2%initial)) then
       return
    end if
    scale = - (prt1%momentum + prt2%momentum) ** 2
    call parton_set_simulated (prt1)
    call parton_set_simulated (prt2)
    call shower%add_parent (prt1)
    call shower%add_parent (prt2)
    factor = sqrt (energy (prt1%momentum)**2 - scale) / &
         space_part_norm(prt1%momentum)
    prt1%parent%type = prt1%type
    prt1%parent%z = one
    prt1%parent%momentum = prt1%momentum
    prt1%parent%t = scale
    prt1%parent%x = prt1%x
    prt1%parent%initial => prt1%initial
    prt1%parent%belongstoFSR = .false.
    prt1%parent%c1 = prt1%c1
    prt1%parent%c2 = prt1%c2

    prt2%parent%type= prt2%type
    prt2%parent%z = one
    prt2%parent%momentum = prt2%momentum
    prt2%parent%t = scale
    prt2%parent%x = prt2%x
    prt2%parent%initial => prt2%initial
    prt2%parent%belongstoFSR = .false.
    prt2%parent%c1 = prt2%c1
    prt2%parent%c2 = prt2%c2

    do
       call shower_get_first_ISR_scale_for_parton (shower, prt1%parent)
       call shower_get_first_ISR_scale_for_parton (shower, prt2%parent)

       !!! redistribute energy among first partons
       prta => prt1%parent
       prtb => prt2%parent

       E = energy (prt1%momentum + prt2%momentum)
       prta%momentum%p(0) = (E**2 - prtb%t + prta%t) / (two * E)
       prtb%momentum%p(0) = E - prta%momentum%p(0)

       exit
    end do

    call parton_set_simulated (prt1%parent)
    call parton_set_simulated (prt2%parent)
    !!! rescale momenta
    do i = 1, 2
       if (i == 1) then
          prt => prt1%parent
       else
          prt => prt2%parent
       end if
       factor = sqrt (energy (prt%momentum)**2 - prt%t) &
            / space_part_norm (prt%momentum)
       prt%momentum = vector4_moving (energy (prt%momentum), &
            factor * space_part (prt%momentum))
    end do

    if (prt1%parent%t < zero) then
       call shower%add_parent (prt1%parent)
       prt1%parent%parent%momentum = prt1%parent%momentum
       prt1%parent%parent%t = prt1%parent%t
       prt1%parent%parent%x = prt1%parent%x
       prt1%parent%parent%initial => prt1%parent%initial
       prt1%parent%parent%belongstoFSR = .false.
       call shower%add_child (prt1%parent%parent, 2)
    end if

    if (prt2%parent%t < zero) then
       call shower%add_parent (prt2%parent)
       prt2%parent%parent%momentum = prt2%parent%momentum
       prt2%parent%parent%t = prt2%parent%t
       prt2%parent%parent%x = prt2%parent%x
       prt2%parent%parent%initial => prt2%parent%initial
       prt2%parent%parent%belongstoFSR = .false.
       call shower%add_child (prt2%parent%parent, 2)
    end if

  end subroutine shower_prepare_for_simulate_isr_ana_test

\eatline
\nwendcode{}\nwbegindocs{288}\nwdocspar
\nwenddocs{}\nwbegincode{289}\moddef{Shower core: procedures}\plusendmoddef
  subroutine shower_add_children_of_emitted_timelike_parton (shower, prt)
    type(shower_t), intent(inout) :: shower
    type(parton_t), pointer :: prt

    if (prt%t > parton_mass_squared (prt) + D_Min_t) then
       if (parton_is_quark (prt)) then
          !!! q -> qg
          call shower%add_child (prt, 1)
          prt%child1%type = prt%type
          prt%child1%momentum%p(0) = prt%z * prt%momentum%p(0)
          prt%child1%t = prt%t
          call shower%add_child (prt, 2)
          prt%child2%type = GLUON
          prt%child2%momentum%p(0) = (one - prt%z) * prt%momentum%p(0)
          prt%child2%t = prt%t
       else
          if (int (prt%x) > 0) then
             call shower%add_child (prt, 1)
             prt%child1%type = int (prt%x)
             prt%child1%momentum%p(0) = prt%z * prt%momentum%p(0)
             prt%child1%t = prt%t
             call shower%add_child (prt, 2)
             prt%child2%type = -int (prt%x)
             prt%child2%momentum%p(0) = (one - prt%z) * prt%momentum%p(0)
             prt%child2%t= prt%t
          else
             call shower%add_child (prt, 1)
             prt%child1%type = GLUON
             prt%child1%momentum%p(0) = prt%z * prt%momentum%p(0)
             prt%child1%t = prt%t
             call shower%add_child (prt, 2)
             prt%child2%type = GLUON
             prt%child2%momentum%p(0) = (one - prt%z) * prt%momentum%p(0)
             prt%child2%t = prt%t
          end if
       end if
    end if
  end subroutine shower_add_children_of_emitted_timelike_parton

\eatline
\nwendcode{}\nwbegindocs{290}\nwdocspar
\nwenddocs{}\nwbegincode{291}\moddef{Shower core: procedures}\plusendmoddef
  subroutine shower_simulate_children_ana (shower,prt)
    type(shower_t), intent(inout) :: shower
    type(parton_t), intent(inout) :: prt
    real(default), dimension(1:2) :: t, random, integral
    integer, dimension(1:2) :: gtoqq
    integer :: daughter
    type(parton_t), pointer :: daughterprt
    integer :: n_loop

    if (signal_is_pending ()) return
    gtoqq = 0

    if (D_PRINT) print *, "D: shower_simulate_children_ana: for parton " , prt%nr

    if (.not. associated (prt%child1) .or. .not. associated (prt%child2)) then
       call msg_error ("Shower: error in simulate_children_ana: no children.")
       return
    end if

    \LA{}Set beam-remnants and internal partons as simulated\RA{}

    integral = zero

    !!! impose constraints by angular ordering -> cf. (26) of Gaining analytic control
    !!! check if no branchings are possible
    if (.not. prt%child1%simulated) then
       prt%child1%t = min (prt%child1%t, &
            0.5_default * prt%child1%momentum%p(0)**2 * (one - &
            parton_get_costheta (prt)))
       if (min (prt%child1%t, prt%child1%momentum%p(0)**2) < &
            parton_mass_squared (prt%child1) + D_Min_t) then
          prt%child1%t = parton_mass_squared (prt%child1)
          call parton_set_simulated (prt%child1)
       end if
    end if
    if (.not. prt%child2%simulated) then
       prt%child2%t = min (prt%child2%t, &
            0.5_default * prt%child2%momentum%p(0)**2 * (one - &
            parton_get_costheta (prt)))
       if (min (prt%child2%t, prt%child2%momentum%p(0)**2) < &
            parton_mass_squared (prt%child2) + D_Min_t) then
          prt%child2%t = parton_mass_squared (prt%child2)
          call parton_set_simulated (prt%child2)
       end if
    end if

    call rng%generate (random)

    n_loop = 0
    do
       if (signal_is_pending ()) return
       n_loop = n_loop + 1
       if (n_loop > 900) then
          !!! try with massless quarks
          treat_duscb_quarks_massless = .true.
       end if
       if (n_loop > 1000) then
          call shower%write ()
          call msg_message (" simulate_children_ana failed for parton ", prt%nr)
          call msg_error ("BUG: too many loops in simulate_children_ana (?)")
          shower%valid = .false.
          return
       end if

       t(1) = prt%child1%t
       t(2) = prt%child2%t

       !!! check if a branching in the range t(i) to t(i) - tstep(i) occurs

       !!! check for child1
       if (.not. prt%child1%simulated) then
          call parton_simulate_stept &
               (prt%child1, integral(1), random(1), gtoqq(1))
       end if
       !!! check for child2
       if (.not. prt%child2%simulated) then
          call parton_simulate_stept &
               (prt%child2, integral(2), random(2), gtoqq(2))
       end if

       if (prt%child1%simulated .and. &
           prt%child2%simulated) then
          if (sqrt (prt%t) <= sqrt (prt%child1%t) + sqrt (prt%child2%t)) then
            \LA{}Repeat the simulation for the parton with the lower virtuality\RA{}
          else
             exit
          end if
       end if
    end do

    call parton_apply_costheta (prt)

    \LA{}Add children\RA{}
    call shower_parton_update_color_connections (shower, prt)
  end subroutine shower_simulate_children_ana

\eatline
\nwendcode{}\nwbegindocs{292}\nwdocspar
\nwenddocs{}\nwbegincode{293}\moddef{Set beam-remnants and internal partons as simulated}\endmoddef
if (HADRON_REMNANT <= abs (prt%type) .and. abs (prt%type) <= HADRON_REMNANT_OCTET) then
   !!! prt is beam-remnant
   call parton_set_simulated (prt)
   return
end if

!!! check if partons are "internal" -> fixed scale
if (prt%child1%type == INTERNAL) then
   call parton_set_simulated (prt%child1)
end if
if (prt%child2%type == INTERNAL) then
   call parton_set_simulated (prt%child2)
end if
\nwendcode{}\nwbegindocs{294}\nwdocspar
\nwenddocs{}\nwbegincode{295}\moddef{Repeat the simulation for the parton with the lower virtuality}\endmoddef
!!! virtuality : t - m**2 (assuming it's not fixed)
if (prt%child1%type == INTERNAL .and. prt%child2%type == INTERNAL) then
   call msg_fatal &
        ("Shower: both partons fixed, but momentum not conserved")
else if (prt%child1%type == INTERNAL) then
   !!! reset child2
   call parton_set_simulated (prt%child2, .false.)
   prt%child2%t = min (prt%child1%t, (sqrt (prt%t) - &
        sqrt (prt%child1%t))**2)
   integral(2) = zero
   call rng%generate (random(2))
else if (prt%child2%type == INTERNAL) then
   ! reset child1
   call parton_set_simulated (prt%child1, .false.)
   prt%child1%t = min (prt%child2%t, (sqrt (prt%t) - &
        sqrt (prt%child2%t))**2)
   integral(1) = zero
   call rng%generate (random(1))
else if (prt%child1%t - parton_mass_squared (prt%child1) > &
     prt%child2%t - parton_mass_squared (prt%child2)) then
   !!! reset child2
   call parton_set_simulated (prt%child2, .false.)
   prt%child2%t = min (prt%child1%t, (sqrt (prt%t) - &
        sqrt (prt%child1%t))**2)
   integral(2) = zero
   call rng%generate (random(2))
else
   !!! reset child1 ! TODO choose child according to their t
   call parton_set_simulated (prt%child1, .false.)
   prt%child1%t = min (prt%child2%t, (sqrt (prt%t) - &
        sqrt (prt%child2%t))**2)
   integral(1) = zero
   call rng%generate (random(1))
end if
\nwendcode{}\nwbegindocs{296}\nwdocspar
\nwenddocs{}\nwbegincode{297}\moddef{Add children}\endmoddef
do daughter = 1, 2
   if (signal_is_pending ()) return
   if (daughter == 1) then
      daughterprt => prt%child1
   else
      daughterprt => prt%child2
   end if
   if (daughterprt%t < parton_mass_squared (daughterprt) + D_Min_t) then
      cycle
   end if
   if (.not. (parton_is_quark (daughterprt) .or. &
        parton_is_gluon (daughterprt))) then
      cycle
   end if
   if (parton_is_quark (daughterprt)) then
      !!! q -> qg
      call shower%add_child (daughterprt, 1)
      daughterprt%child1%type = daughterprt%type
      daughterprt%child1%momentum%p(0) = daughterprt%z * &
           daughterprt%momentum%p(0)
      daughterprt%child1%t = daughterprt%t
      call shower%add_child (daughterprt, 2)
      daughterprt%child2%type = GLUON
      daughterprt%child2%momentum%p(0) = (one - daughterprt%z) * &
           daughterprt%momentum%p(0)
      daughterprt%child2%t = daughterprt%t
   else if (parton_is_gluon (daughterprt)) then
      if (gtoqq(daughter) > 0) then
         call shower%add_child (daughterprt, 1)
         daughterprt%child1%type = gtoqq (daughter)
         daughterprt%child1%momentum%p(0) = &
              daughterprt%z * daughterprt%momentum%p(0)
         daughterprt%child1%t = daughterprt%t
         call shower%add_child (daughterprt, 2)
         daughterprt%child2%type = - gtoqq (daughter)
         daughterprt%child2%momentum%p(0) = (one - &
              daughterprt%z) * daughterprt%momentum%p(0)
         daughterprt%child2%t = daughterprt%t
      else
         call shower%add_child (daughterprt, 1)
         daughterprt%child1%type = GLUON
         daughterprt%child1%momentum%p(0) = &
              daughterprt%z * daughterprt%momentum%p(0)
         daughterprt%child1%t = daughterprt%t
         call shower%add_child (daughterprt, 2)
         daughterprt%child2%type = GLUON
         daughterprt%child2%momentum%p(0) = (one - &
              daughterprt%z) * daughterprt%momentum%p(0)
         daughterprt%child2%t = daughterprt%t
      end if
   end if
end do
\nwendcode{}\nwbegindocs{298}\nwdocspar
\nwenddocs{}\nwbegindocs{299}The recoiler is {\tt{}otherprt}. Instead of the random number and the
exponential of the integral, we compare the logarithm of the random
number and the integral.
\nwenddocs{}\nwbegincode{300}\moddef{Shower core: procedures}\plusendmoddef
  subroutine shower_isr_step_pt (shower, prt)
    type(shower_t), intent(inout) :: shower
    type(parton_t), target, intent(inout) :: prt
    type(parton_t), pointer :: otherprt

    real(default) :: scale, scalestep
    real(default) :: integral, random, factor
    real(default) :: temprand1, temprand2

    otherprt => shower_find_recoiler (shower, prt)

    scale = prt%scale
    call rng%generate (temprand1)
    call rng%generate (temprand2)
    scalestep = max (abs (scalefactor1 * scale) * temprand1, &
         scalefactor2 * temprand2 * D_Min_scale)
    call rng%generate (random)
    random = - twopi * log(random)
    integral = zero

    if (scale - 0.5_default * scalestep < D_Min_scale) then
       !!! close enough to cut-off scale -> ignore
       prt%scale = zero
       prt%t = parton_mass_squared (prt)
       call parton_set_simulated (prt)
    else
       prt%scale = scale - 0.5_default * scalestep
       factor = scalestep * (D_alpha_s_isr (prt%scale) / (prt%scale * &
            shower%get_pdf (prt%initial%type, prt%x, prt%scale, prt%type)))
       integral = integral + factor * integral_over_z_isr_pt &
            (prt, otherprt, (random - integral) / factor)
       if (integral > random) then
          !!! prt%scale set above and prt%z set in integral_over_z_isr_pt
          call parton_set_simulated (prt)
          prt%t = - prt%scale / (one - prt%z)
       else
          prt%scale = scale - scalestep
       end if
    end if

  contains

    function integral_over_z_isr_pt (prt, otherprt, final) &
         result (integral)
      type(parton_t), intent(inout) :: prt, otherprt
      real(default), intent(in) :: final
      real(default) :: integral
      real(default) :: mbr, r
      real(default) :: zmin, zmax, z, zstep
      integer :: n_bin
      integer, parameter :: n_total_bins = 100
      real(default) :: quarkpdfsum
      real(default) :: temprand
      integer :: quark

      quarkpdfsum = zero
      if (D_PRINT) then
         print *, "D: integral_over_z_isr_pt: for scale = ", prt%scale
      end if

      integral = zero
      mbr = (prt%momentum + otherprt%momentum)**1
      zmin = prt%x
      zmax = min (one - (sqrt (prt%scale) / mbr) * &
           (sqrt(one + 0.25_default * prt%scale / mbr**2) - &
           0.25_default * sqrt(prt%scale) / mbr), shower%maxz_isr)
      zstep = (zmax - zmin) / n_total_bins

      if (ASSERT) then
         if (zmin > zmax) then
            call msg_bug(" error in integral_over_z_isr_pt: zmin > zmax ")
            integral = zero
         end if
      end if

      !!! divide the range [zmin:zmax] in n_total_bins
      BINS: do n_bin = 1, n_total_bins
         z = zmin + zstep * (n_bin - 0.5_default)
         !!! z-value in the middle of the bin

         if (parton_is_gluon (prt)) then
            QUARKS: do quark = -D_Nf, D_Nf
               if (quark == 0) cycle quarks
               quarkpdfsum = quarkpdfsum + shower%get_pdf &
                    (prt%initial%type, prt%x / z, prt%scale, quark)
            end do QUARKS
            !!! g -> gg or q -> gq
            integral = integral + (zstep / z) * ((P_ggg (z) + &
                 P_ggg (one - z)) * shower%get_pdf (prt%initial%type, &
                 prt%x / z, prt%scale, GLUON) + P_qqg (one - z) * quarkpdfsum)
         else if (parton_is_quark (prt)) then
            !!! q -> qg or g -> qq
            integral = integral + (zstep / z) * ( P_qqg (z) * &
                 shower%get_pdf (prt%initial%type, prt%x / z, prt%scale, prt%type) + &
                 P_gqq(z) * shower%get_pdf (prt%initial%type, prt%x / z, prt%scale, GLUON))
         else
            ! call msg_fatal ("Bug neither quark nor gluon in" &
            !           // " integral_over_z_isr_pt")
         end if
         if (integral > final) then
            prt%z = z
            call rng%generate (temprand)
            !!! decide type of father partons
              if (parton_is_gluon (prt)) then
                 if (temprand > (P_qqg (one - z) * quarkpdfsum) / &
                      ((P_ggg (z) + P_ggg (one - z)) * shower%get_pdf &
                      (prt%initial%type, prt%x / z, prt%scale, GLUON) &
                      + P_qqg (one - z) * quarkpdfsum)) then
                    !!! gluon => gluon + gluon
                    prt%aux_pt = GLUON
                 else
                    !!! quark => quark + gluon
                    !!! decide which quark flavor the parent is
                    r = temprand * quarkpdfsum
                    WHICH_QUARK: do quark = -D_Nf, D_Nf
                       if (quark == 0) cycle WHICH_QUARK
                       if (r > quarkpdfsum - shower%get_pdf (prt%initial%type, &
                            prt%x / z, prt%scale, quark)) then
                          prt%aux_pt = quark
                          exit WHICH_QUARK
                       else
                          quarkpdfsum = quarkpdfsum - shower%get_pdf &
                               (prt%initial%type, prt%x / z, prt%scale, quark)
                       end if
                    end do WHICH_QUARK
                 end if

              else if (parton_is_quark (prt)) then
                 if (temprand > (P_qqg (z) * shower%get_pdf (prt%initial%type, &
                      prt%x / z, prt%scale, prt%type)) / &
                      (P_qqg (z) * shower%get_pdf (prt%initial%type, prt%x / z, &
                      prt%scale, prt%type) + &
                      P_gqq (z) * shower%get_pdf (prt%initial%type, prt%x / z, &
                      prt%scale, GLUON))) then
                    !!! gluon => quark + antiquark
                    prt%aux_pt = GLUON
                 else
                    !!! quark => quark + gluon
                    prt%aux_pt = prt%type
                 end if
              end if
              exit BINS
           end if
        end do BINS
      end function integral_over_z_isr_pt
    end subroutine shower_isr_step_pt

\eatline
\nwendcode{}\nwbegindocs{301}This function returns a pointer to the parton with the next ISR
branching, while FSR branchings are ignored.
\nwenddocs{}\nwbegincode{302}\moddef{Shower core: shower: TBP}\plusendmoddef
  procedure :: generate_next_isr_branching_veto => &
       shower_generate_next_isr_branching_veto
\nwendcode{}\nwbegincode{303}\moddef{Shower core: procedures}\plusendmoddef
  function shower_generate_next_isr_branching_veto &
       (shower) result (next_brancher)
    class(shower_t), intent(inout) :: shower
    type(parton_pointer_t) :: next_brancher
    integer :: i
    type(parton_t),  pointer :: prt
    real(default) :: random
    !!! pointers to branchable partons
    type(parton_pointer_t), dimension(:), allocatable :: partons
    integer :: n_partons
    real(default) :: weight
    real(default) :: temp1, temp2, temp3, E3

    if (signal_is_pending ()) return

    if (isr_pt_ordered) then
       next_brancher = shower%generate_next_isr_branching ()
       return
    end if
    next_brancher%p => null()
    !!! branchable partons
    n_partons = 0
    do i = 1,size (shower%partons)
       prt => shower%partons(i)%p
       if (signal_is_pending ()) return
       if (.not. associated (prt)) cycle
       if (prt%belongstoFSR) cycle
       if (parton_is_final (prt)) cycle
       if (.not. prt%belongstoFSR .and. prt%simulated) cycle
       n_partons = n_partons + 1
    end do
    if (n_partons == 0) then
       return
    end if
    allocate (partons(1:n_partons))
    n_partons = 1
    do i = 1, size (shower%partons)
       prt => shower%partons(i)%p
       if (.not. associated (prt)) cycle
       if (prt%belongstoFSR) cycle
       if (parton_is_final (prt)) cycle
       if (.not. prt%belongstoFSR .and. prt%simulated) cycle
       if (signal_is_pending ()) return
       partons(n_partons)%p => shower%partons(i)%p
       n_partons = n_partons + 1
    end do
    !!! generate initial trial scales
    do i = 1, size (partons)
       if (signal_is_pending ()) return
       call generate_next_trial_scale (partons(i)%p)
    end do

    do
       !!! search for parton with the highest trial scale
       prt => partons(1)%p
       do i = 1, size (partons)
          if (prt%t >= zero) cycle
          if (abs (partons(i)%p%t) > abs (prt%t)) then
             prt => partons(i)%p
          end if
       end do

       if (prt%t >= zero) then
          next_brancher%p => null()
          exit
       end if
       !!! generate trial z and type of mother prt
       call generate_trial_z_and_typ (prt)

       !!! weight with pdf and alpha_s
       temp1 = (D_alpha_s_isr ((one - prt%z) * abs(prt%t)) &
            / sqrt (alphasxpdfmax))
       temp2 = shower%get_xpdf (prt%initial%type, prt%x, prt%t, &
            prt%type) / sqrt (alphasxpdfmax)
       temp3 = shower%get_xpdf (prt%initial%type, prt%child1%x, prt%child1%t, &
            prt%child1%type) / &
            shower%get_xpdf (prt%initial%type, prt%child1%x, prt%t, prt%child1%type)
       if (temp1 * temp2 * temp3 > one) then
          print *, "weights:", temp1, temp2, temp3
       end if
       weight = (D_alpha_s_isr ((one - prt%z) * abs(prt%t))) * &
            shower%get_xpdf (prt%initial%type, prt%x, prt%t, prt%type) * &
            shower%get_xpdf (prt%initial%type, prt%child1%x, prt%child1%t, &
            prt%child1%type) / &
            shower%get_xpdf (prt%initial%type, prt%child1%x, prt%t, prt%child1%type)
       if (weight > alphasxpdfmax) then
          print *, "Setting alphasxpdfmax from ", alphasxpdfmax, " to ", weight
          alphasxpdfmax = weight
       end if
       weight = weight / alphasxpdfmax
       call rng%generate (random)
       if (weight < random) then
          !!! discard branching
          call generate_next_trial_scale (prt)
          cycle
       end if
       !!! branching accepted so far
       !!! generate emitted parton
       prt%child2%t = abs(prt%t)
       prt%child2%momentum%p(0) = sqrt (abs(prt%t))
       if (shower%isr_only_onshell_emitted_partons) then
          prt%child2%t = parton_mass_squared(prt%child2)
       else
          call parton_next_t_ana (prt%child2)
       end if

       if (thetabar (prt, shower_find_recoiler (shower, prt), &
            shower%isr_angular_ordered, E3)) then
          prt%momentum%p(0) = E3
          prt%child2%momentum%p(0) = E3 - prt%child1%momentum%p(0)

          !!! found branching
          call parton_generate_ps_ini (prt)
          next_brancher%p => prt
          call parton_set_simulated (prt)
          exit
       else
          call generate_next_trial_scale (prt)
          cycle
       end if
    end do
    if (.not. associated (next_brancher%p)) then
       !!! no further branching found -> all partons emitted by hadron
       print *, "--all partons emitted by hadrons---"
       do i = 1, size(partons)
          call shower_replace_parent_by_hadron (shower, partons(i)%p%child1)
       end do
    end if
    !!! some bookkeeping
    call shower%sort_partons ()
    ! call shower_boost_to_CMframe(shower)         ! really necessary?
    ! call shower_rotate_to_z(shower)              ! really necessary?
  contains

    subroutine generate_next_trial_scale(prt)
      type(parton_t), pointer, intent(inout) :: prt
      real(default) :: random, F
      real(default) :: zmax = 0.99_default !! ??
      call rng%generate (random)
      F = one   !!! TODO
      F = alphasxpdfmax / (two * pi)
      if (parton_is_quark (prt%child1)) then
         F = F * (integral_over_P_gqq (prt%child1%x, zmax) + &
                  integral_over_P_qqg (prt%child1%x, zmax))
      else if (parton_is_gluon (prt%child1)) then
         F = F * (integral_over_P_ggg (prt%child1%x, zmax) + &
              two * D_Nf * &
              integral_over_P_qqg (one - zmax, one - prt%child1%x))
      else
         print *, "neither quark nor gluon in generate_next_trial_scale"
      end if
      F = F / shower%get_xpdf (prt%child1%initial%type, prt%child1%x, &
           prt%child1%t, prt%child1%type)
      prt%t = prt%t * random**(one / F)
      if (abs (prt%t) - parton_mass_squared (prt) < D_Min_t) then
         prt%t = parton_mass_squared (prt)
      end if
    end subroutine generate_next_trial_scale

    subroutine generate_trial_z_and_typ(prt)
      type(parton_t), pointer, intent(inout) :: prt
      real(default) :: random
      real(default) :: z, zstep, zmin, integral
      real(default) :: zmax = 0.99_default !! ??
      ! print *, "generate next_z"
      call rng%generate (random)
      integral = zero
      !!! decide which branching a->bc occurs
      if (parton_is_quark (prt%child1)) then
         if (random < integral_over_P_qqg (prt%child1%x, zmax) / &
              (integral_over_P_qqg (prt%child1%x, zmax) + &
               integral_over_P_gqq(prt%child1%x, zmax))) then
            prt%type = prt%child1%type
            prt%child2%type = GLUON
            integral = integral_over_P_qqg (prt%child1%x, zmax)
         else
            prt%type = GLUON
            prt%child2%type = - prt%child1%type
            integral = integral_over_P_gqq (prt%child1%x, zmax)
         end if
      else if (parton_is_gluon (prt%child1)) then
         if (random < integral_over_P_ggg (prt%child1%x, zmax) / &
              (integral_over_P_ggg (prt%child1%x, zmax) + two * D_Nf * &
              integral_over_P_qqg (one - zmax, &
              one - prt%child1%x))) then
            prt%type = GLUON
            prt%child2%type = GLUON
            integral = integral_over_P_ggg (prt%child1%x, zmax)
         else
            call rng%generate (random)
            prt%type = 1 + floor(random * D_Nf)
            call rng%generate (random)
            if (random > 0.5_default) prt%type = - prt%type
            prt%child2%type = prt%type
            integral = integral_over_P_qqg (one - zmax, &
                 one - prt%child1%x)
         end if
      else
         print *, "neither quark nor gluon in generate_next_trial_scale"
      end if
      !!! generate the z-value
      !!! z between prt%child1%x and zmax
      ! prt%z = one - random * (one - prt%child1%x)      ! TODO

      call rng%generate (random)
      zmin = prt%child1%x
      zstep = max(0.1_default, 0.5_default * (zmax - zmin))
      z = zmin
      if (zmin > zmax) then
         print *, " zmin = ", zmin, " zmax = ", zmax
         call msg_fatal ("Shower: zmin greater than zmax")
      end if
      !!! procedure pointers would be helpful here
      if (parton_is_quark (prt) .and. parton_is_quark (prt%child1)) then
         do
            zstep = min(zstep, 0.5_default * (zmax - z))
            if (abs(zstep) < 0.00001) exit
            if (integral_over_P_qqg (zmin, z) < random * integral) then
               if (integral_over_P_qqg (zmin, min(z + zstep, zmax)) &
                    < random * integral) then
                  z = min (z + zstep, zmax)
                  cycle
               else
                  zstep = zstep * 0.5_default
                  cycle
               end if
            end if
         end do
      else if (parton_is_quark (prt) .and. parton_is_gluon (prt%child1)) then
         do
            zstep = min(zstep, 0.5_default * (zmax - z))
            if (abs(zstep) < 0.00001) exit
            if (integral_over_P_qqg (zmin, z) < random * integral) then
               if (integral_over_P_qqg (zmin, min(z + zstep, zmax)) &
                    < random * integral) then
                  z = min(z + zstep, zmax)
                  cycle
               else
                  zstep = zstep * 0.5_default
                  cycle
               end if
            end if
         end do
      else if (parton_is_gluon (prt) .and. parton_is_quark (prt%child1)) then
         do
            zstep = min(zstep, 0.5_default * (zmax - z))
            if (abs (zstep) < 0.00001) exit
            if (integral_over_P_gqq (zmin, z) < random * integral) then
               if (integral_over_P_gqq (zmin, min(z + zstep, zmax)) &
                    < random * integral) then
                  z = min (z + zstep, zmax)
                  cycle
               else
                  zstep = zstep * 0.5_default
                  cycle
               end if
            end if
         end do
      else if (parton_is_gluon (prt) .and. parton_is_gluon (prt%child1)) then
         do
            zstep = min(zstep, 0.5_default * (zmax - z))
            if (abs (zstep) < 0.00001) exit
            if (integral_over_P_ggg(zmin, z) < random * integral) then
               if (integral_over_P_ggg (zmin, min(z + zstep, zmax)) &
                    < random * integral) then
                  z = min(z + zstep, zmax)
                  cycle
               else
                  zstep = zstep * 0.5_default
                  cycle
               end if
            end if
         end do
      else
      end if
      prt%z = z
      prt%x = prt%child1%x / prt%z
    end subroutine generate_trial_z_and_typ
  end function shower_generate_next_isr_branching_veto

\eatline
\nwendcode{}\nwbegindocs{304}\nwdocspar
\nwenddocs{}\nwbegincode{305}\moddef{Shower core: procedures}\plusendmoddef
  function shower_find_recoiler (shower, prt) result(recoiler)
    type(shower_t), intent(inout) :: shower
    type(parton_t), intent(inout), target :: prt
    type(parton_t), pointer :: recoiler
    type(parton_t), pointer :: otherprt1, otherprt2
    integer :: n_int
    otherprt1 => null()
    otherprt2 => null()
    DO_INTERACTIONS: do n_int = 1, size(shower%interactions)
       otherprt1=>shower%interactions(n_int)%i%partons(1)%p
       otherprt2=>shower%interactions(n_int)%i%partons(2)%p
       PARTON1: do
          if (associated (otherprt1%parent)) then
             if (.not. parton_is_hadron (otherprt1%parent) .and. &
                  otherprt1%parent%simulated) then
                otherprt1 => otherprt1%parent
                if (associated (otherprt1, prt)) then
                   exit PARTON1
                end if
             else
                exit PARTON1
             end if
          else
             exit PARTON1
          end if
       end do PARTON1
       PARTON2: do
          if (associated (otherprt2%parent)) then
             if (.not. parton_is_hadron (otherprt2%parent) .and. &
                  otherprt2%parent%simulated) then
                otherprt2 => otherprt2%parent
                if (associated (otherprt2, prt)) then
                   exit PARTON2
                end if
             else
                exit PARTON2
             end if
          else
             exit PARTON2
          end if
       end do PARTON2

       if (associated (otherprt1, prt) .or. associated (otherprt2, prt)) then
          exit DO_INTERACTIONS
       end if
       if (associated (otherprt1%parent, prt) .or. &
            associated (otherprt2%parent, prt)) then
          exit DO_INTERACTIONS
       end if
    end do DO_INTERACTIONS

    recoiler => null()
    if (associated (otherprt1%parent, prt)) then
       recoiler => otherprt2
    else if (associated (otherprt2%parent, prt)) then
       recoiler => otherprt1
    else if (associated (otherprt1, prt)) then
       recoiler => otherprt2
    else if (associated (otherprt2, prt)) then
       recoiler => otherprt1
    else
       call shower%write ()
       call parton_write (prt)
       call msg_error ("Shower: no otherparton found")
    end if
  end function shower_find_recoiler

\eatline
\nwendcode{}\nwbegindocs{306}\nwdocspar
\nwenddocs{}\nwbegincode{307}\moddef{Shower core: procedures}\plusendmoddef
  subroutine shower_isr_step (shower, prt)
    type(shower_t), intent(inout) :: shower
    type(parton_t), target, intent(inout) :: prt
    type(parton_t), pointer :: otherprt => null()
    real(default) :: t, tstep
    real(default) :: integral, random
    real(default) :: temprand1, temprand2
    otherprt => shower_find_recoiler(shower, prt)
    ! if (.not. otherprt%child1%belongstointeraction) then
    !    otherprt => otherprt%child1
    ! end if

    if (signal_is_pending ()) return
    t = max(prt%t, prt%child1%t)
    call rng%generate (random)
    ! compare Integral and log(random) instead of random and exp(-Integral)
    random = - twopi * log(random)
    integral = zero
    call rng%generate (temprand1)
    call rng%generate (temprand2)
    tstep = max (abs (0.02_default * t) * temprand1, &
         0.02_default * temprand2 * D_Min_t)
    if (t + 0.5_default * tstep > - D_Min_t) then
       prt%t = parton_mass_squared (prt)
       call parton_set_simulated (prt)
    else
       prt%t = t + 0.5_default * tstep
       integral = integral + tstep * &
            integral_over_z_isr (prt, otherprt,(random - integral) / tstep)
       if (integral > random) then
          prt%t = t + 0.5_default * tstep
          prt%x = prt%child1%x / prt%z
          call parton_set_simulated (prt)
       else
          prt%t = t + tstep
       end if
    end if

  contains

    function integral_over_z_isr (prt, otherprt, final) result (integral)
      type(parton_t), intent(inout) :: prt, otherprt
      real(default), intent(in) :: final
      real(default) integral
      real(default) :: minz, maxz, z, shat,s
      integer :: quark

      !!! calculate shat -> s of parton-parton system
      shat= (otherprt%momentum + prt%child1%momentum)**2
      !!! calculate s -> s of hadron-hadron system
      s= (otherprt%initial%momentum + prt%initial%momentum)**2
      integral = zero
      minz = prt%child1%x
      maxz = maxzz (shat, s, shower%maxz_isr, shower%minenergy_timelike)

      !!! for gluon
      if (parton_is_gluon (prt%child1)) then
         !!! 1: g->gg
         prt%type = GLUON
         prt%child2%type = GLUON
           z = minz
           prt%child2%t = abs(prt%t)
           call integral_over_z_part_isr &
                (prt, otherprt, shat, minz, maxz, integral, final)
           if (integral > final) then
              return
           end if
           !!! 2: q->gq
           do quark = -D_Nf, D_Nf
              if (quark == 0) cycle
              prt%type = quark
              prt%child2%type = quark
              z = minz
              prt%child2%t = abs(prt%t)
              call integral_over_z_part_isr &
                   (prt, otherprt, shat, minz, maxz, integral, final)
              if (integral > final) then
                 return
              end if
           end do
        else if (parton_is_quark (prt%child1)) then
           !!! 1: q->qg
           prt%type = prt%child1%type
           prt%child2%type = GLUON
           z = minz
           prt%child2%t = abs(prt%t)
           call integral_over_z_part_isr &
                (prt,otherprt, shat, minz, maxz, integral, final)
           if (integral > final) then
              return
           end if
           !!! 2: g->qqbar
           prt%type = GLUON
           prt%child2%type = -prt%child1%type
           z = minz
           prt%child2%t = abs(prt%t)
           call integral_over_z_part_isr &
                (prt,otherprt, shat, minz, maxz, integral, final)
           if (integral > final) then
              return
           end if
        end if
      end function integral_over_z_isr

      subroutine integral_over_z_part_isr &
           (prt, otherprt, shat ,minz, maxz, retvalue, final)
        type(parton_t), intent(inout) :: prt, otherprt
        real(default), intent(in) :: shat, minz, maxz, final
        real(default), intent(inout) :: retvalue
        real(default) :: z, zstep
        real(default) :: r1,r3,s1,s3
        real(default) :: pdf_divisor
        real(default) :: temprand
        real(default), parameter :: zstepfactor = 0.1_default
        real(default), parameter :: zstepmin = 0.0001_default
        if (D_PRINT) print *, "D: integral_over_z_part_isr"
        pdf_divisor = shower%get_pdf &
             (prt%initial%type, prt%child1%x, prt%t, prt%child1%type)
        z = minz
        s1 = shat + abs(otherprt%t) + abs(prt%child1%t)
        r1 = sqrt (s1**2 - four * abs(otherprt%t * prt%child1%t))
        ZLOOP: do
           if (signal_is_pending ()) return
           if (z >= maxz) then
              exit
           end if
           call rng%generate (temprand)
           if (parton_is_gluon (prt%child1)) then
              if (parton_is_gluon (prt)) then
                 !!! g-> gg -> divergencies at z->0 and z->1
                 zstep = max(zstepmin, temprand * zstepfactor * z * (one - z))
              else
                 !!! q-> gq -> divergencies at z->0
                 zstep = max(zstepmin, temprand * zstepfactor * (one - z))
              end if
           else
              if (parton_is_gluon (prt)) then
                 !!! g-> qqbar -> no divergencies
                 zstep = max(zstepmin, temprand * zstepfactor)
              else
                 !!! q-> qg -> divergencies at z->1
                 zstep = max(zstepmin, temprand * zstepfactor * (one - z))
              end if
           end if
           zstep = min(zstep, maxz - z)
           prt%z = z + 0.5_default * zstep
           s3 = shat / prt%z + abs(otherprt%t) + abs(prt%t)
           r3 = sqrt (s3**2 - four * abs(otherprt%t * prt%t))
           !!! TODO: WHY is this if needed?
           if (abs(otherprt%t) > eps0) then
              prt%child2%t = min ((s1 * s3 - r1 * r3) / &
                   (two * abs(otherprt%t)) - abs(prt%child1%t) - &
                   abs(prt%t), abs(prt%child1%t))
           else
              prt%child2%t = abs(prt%child1%t)
           end if
           do
              prt%child2%momentum%p(0) = sqrt (abs(prt%child2%t))
              if (shower%isr_only_onshell_emitted_partons) then
                 prt%child2%t = parton_mass_squared (prt%child2)
              else
                 call parton_next_t_ana (prt%child2)
              end if
              !!! take limits by recoiler into account
              prt%momentum%p(0) = (shat / prt%z + &
                   abs(otherprt%t) - abs(prt%child1%t) - &
                   prt%child2%t) / (two * sqrt(shat))
              prt%child2%momentum%p(0) = &
                   prt%momentum%p(0) - prt%child1%momentum%p(0)
              !!! check if E and t of prt%child2 are consistent
              if (prt%child2%momentum%p(0)**2 < prt%child2%t &
                   .and. prt%child2%t > parton_mass_squared (prt%child2)) then
                 !!! E is too small to have p_T^2 = E^2 - t > 0
                 !!!      -> cycle to find another solution
                 cycle
              else
                 !!! E is big enough -> exit
                 exit
              end if
           end do
           if (thetabar (prt, otherprt, shower%isr_angular_ordered) &
                .and. pdf_divisor > zero &
                .and. prt%child2%momentum%p(0) > zero) then
              retvalue = retvalue + (zstep / prt%z) * &
                   (D_alpha_s_isr ((one - prt%z) * prt%t) * &
                   P_prt_to_child1 (prt) * &
                   shower%get_pdf (prt%initial%type, prt%child1%x / prt%z, &
                   prt%t, prt%type)) / (abs(prt%t) * pdf_divisor)
           end if
           if (retvalue > final) then
              exit
           else
              z = z + zstep
           end if
        end do ZLOOP
      end subroutine integral_over_z_part_isr
    end subroutine shower_isr_step

\eatline
\nwendcode{}\nwbegindocs{308}This returns a pointer to the parton with the next ISR branching, again
FSR branchings are ignored.
\nwenddocs{}\nwbegincode{309}\moddef{Shower core: shower: TBP}\plusendmoddef
  procedure :: generate_next_isr_branching => &
       shower_generate_next_isr_branching
\nwendcode{}\nwbegincode{310}\moddef{Shower core: procedures}\plusendmoddef
  function shower_generate_next_isr_branching &
       (shower) result (next_brancher)
    class(shower_t), intent(inout) :: shower
    type(parton_pointer_t) :: next_brancher
    integer i, index
    type(parton_t),  pointer :: prt
    real(default) :: maxscale
    next_brancher%p => null()
    do
       if (signal_is_pending ()) return
       if (shower_isr_is_finished (shower)) exit
       !!! find mother with highest |t| or pt to be simulated
       index = 0
       maxscale = zero
       call shower%sort_partons ()
       do i = 1,size (shower%partons)
          prt => shower%partons(i)%p
          if (.not. associated (prt)) cycle
          if (.not. isr_pt_ordered) then
             if (prt%belongstointeraction) cycle
          end if
          if (prt%belongstoFSR) cycle
          if (parton_is_final (prt)) cycle
          if (.not. prt%belongstoFSR .and. prt%simulated) cycle
          index = i
          exit
       end do
       if (ASSERT) then
          if (index == 0) then
             call msg_fatal(" no branchable partons found")
             return
          end if
       end if

       prt => shower%partons(index)%p

       !!! ISR simulation
       if (isr_pt_ordered) then
          call shower_isr_step_pt (shower, prt)
       else
          call shower_isr_step (shower, prt)
       end if
       if (prt%simulated) then
          if (prt%t < zero) then
             next_brancher%p => prt
             if (.not. isr_pt_ordered) call parton_generate_ps_ini (prt)
             exit
          else
             if (.not. isr_pt_ordered) then
                call shower_replace_parent_by_hadron (shower, prt%child1)
             else
                call shower_replace_parent_by_hadron (shower, prt)
             end if
          end if
       end if
    end do

    !!! some bookkeeping
    call shower%sort_partons ()
    call shower_boost_to_CMframe (shower)         !!! really necessary?
    call shower_rotate_to_z (shower)              !!! really necessary?
    ! print *, "shower_generate_next_isr_branching finished"
    end function shower_generate_next_isr_branching

\eatline
\nwendcode{}\nwbegindocs{311}This is a loop which searches for all emitted and branched partons.
\nwenddocs{}\nwbegincode{312}\moddef{Shower core: shower: TBP}\plusendmoddef
  procedure :: generate_fsr_for_isr_partons => &
       shower_generate_fsr_for_partons_emitted_in_ISR
\nwendcode{}\nwbegincode{313}\moddef{Shower core: procedures}\plusendmoddef
  subroutine shower_generate_fsr_for_partons_emitted_in_ISR (shower)
    class(shower_t), intent(inout) :: shower
    integer :: n_int, i
    type(parton_t), pointer :: prt
    if (shower%isr_only_onshell_emitted_partons) return
    if (D_PRINT) print *, "D: shower_generate_fsr_for_partons_emitted_in_ISR"
    INTERACTIONS_LOOP: do n_int = 1, size (shower%interactions)
       INCOMING_PARTONS_LOOP: do i = 1, 2
          if (signal_is_pending ()) return
          prt => shower%interactions(n_int)%i%partons(i)%p
          PARENT_PARTONS_LOOP: do
             if (associated (prt%parent)) then
                if (.not. parton_is_hadron (prt%parent)) then
                   prt => prt%parent
                else
                   exit
                end if
             else
                exit
             end if
             if (associated (prt%child2)) then
                if (parton_is_branched (prt%child2)) then
                   call shower_parton_generate_fsr (shower, prt%child2)
                end if
             else
                ! call msg_fatal ("Shower: no child2 associated?")
             end if
          end do PARENT_PARTONS_LOOP
       end do INCOMING_PARTONS_LOOP
    end do INTERACTIONS_LOOP
  end subroutine shower_generate_fsr_for_partons_emitted_in_ISR

\eatline
\nwendcode{}\nwbegindocs{314}This executes the branching generated by
{\tt{}shower{\char95}generate{\char95}next{\char95}isr{\char95}branching}, that means it generates the
flavours, momenta, etc.
\nwenddocs{}\nwbegincode{315}\moddef{Shower core: shower: TBP}\plusendmoddef
  procedure :: execute_next_isr_branching => shower_execute_next_isr_branching
\nwendcode{}\nwbegincode{316}\moddef{Shower core: procedures}\plusendmoddef
  subroutine shower_execute_next_isr_branching (shower, prtp)
    class(shower_t), intent(inout) :: shower
    type(parton_pointer_t), intent(inout) :: prtp
    type(parton_t), pointer :: prt, otherprt
    type(parton_t), pointer :: prta, prtb, prtc, prtr
    real(default) :: mar, mbr
    real(default) :: phirand
    ! print *, "shower_execute_next_isr_branching"
    if (.not. associated (prtp%p)) then
       call msg_fatal ("Shower: prtp not associated")
    end if

    prt => prtp%p

    if ((.not. isr_pt_ordered .and. prt%t > - D_Min_t) .or. &
         (isr_pt_ordered .and. prt%scale < D_Min_scale)) then
       call msg_error ("Shower: no branching to be executed.")
    end if

    otherprt => shower_find_recoiler (shower, prt)
    if (isr_pt_ordered) then
       !!! get the recoiler
       otherprt => shower_find_recoiler (shower, prt)
       if (associated (otherprt%parent)) then
          !!! Why only for pt ordered
          if (.not. parton_is_hadron (otherprt%parent) .and. &
               isr_pt_ordered) otherprt => otherprt%parent
       end if
       if (.not. associated (prt%parent)) then
          call shower%add_parent (prt)
       end if
       prt%parent%belongstoFSR = .false.
       if (.not. associated (prt%parent%child2)) then
          call shower%add_child (prt%parent, 2)
       end if

       prta => prt%parent          !!! new parton a with branching a->bc
       prtb => prt                 !!! former parton
       prtc => prt%parent%child2   !!! emitted parton
       prtr => otherprt            !!! recoiler

       mbr = (prtb%momentum + prtr%momentum)**1
       mar = mbr / sqrt(prt%z)

       !!! 1. assume you are in the restframe
       !!! 2. rotate by random phi
       call rng%generate (phirand)
       phirand = twopi * phirand
       call shower_apply_lorentztrafo (shower, &
            rotation(cos(phirand), sin(phirand),vector3_canonical(3)))
       !!! 3. Put the b off-shell
       !!! and
       !!! 4. construct the massless a
       !!! and the parton (eventually emitted by a)

       !!! generate the flavour of the parent (prta)
       if (prtb%aux_pt /= 0) prta%type = prtb%aux_pt
       if (parton_is_quark (prtb)) then
          if (prta%type == prtb%type) then
             !!! (anti)-quark -> (anti-)quark + gluon
             prta%type = prtb%type   ! quarks have same flavour
             prtc%type = GLUON        ! emitted gluon
          else
             !!! gluon -> quark + antiquark
             prta%type = GLUON
             prtc%type = - prtb%type
          end if
       else if (parton_is_gluon (prtb)) then
          prta%type = GLUON
          prtc%type = GLUON
       else
          ! STOP "Bug in shower_execute_next_branching: neither quark nor gluon"
       end if

       prta%initial => prtb%initial
       prta%belongstoFSR = .false.
       prta%scale = prtb%scale
       prta%x = prtb%x / prtb%z

       prtb%momentum = vector4_moving ((mbr**2 + prtb%t) / (two * mbr), &
            vector3_canonical(3) * &
            sign ((mbr**2 - prtb%t) / (two * mbr), &
            prtb%momentum%p(3)))
       prtr%momentum = vector4_moving ((mbr**2 - prtb%t) / (two * mbr), &
            vector3_canonical(3) * &
            sign( (mbr**2 - prtb%t) / (two * mbr), &
            prtr%momentum%p(3)))

       prta%momentum = vector4_moving ((0.5_default / mbr) * &
            ((mbr**2 / prtb%z) + prtb%t - parton_mass_squared(prtc)), &
            vector3_null)
       prta%momentum = vector4_moving (prta%momentum%p(0), &
            vector3_canonical(3) * &
            (0.5_default / prtb%momentum%p(3)) * &
            ((mbr**2 / prtb%z) - two &
            * prtr%momentum%p(0) * prta%momentum%p(0) ) )
       if (prta%momentum%p(0)**2 - prta%momentum%p(3)**2 - parton_mass_squared (prtc) &
            > zero) then
          !!! This SHOULD be always fulfilled???
          prta%momentum = vector4_moving (prta%momentum%p(0), &
               vector3_moving([sqrt (prta%momentum%p(0)**2 - &
               prta%momentum%p(3)**2 - &
               parton_mass_squared (prtc)), zero, &
               prta%momentum%p(3)]))
       end if
       prtc%momentum = prta%momentum - prtb%momentum

       !!! 5. rotate to have a along z-axis
       call shower_boost_to_CMframe (shower)
       call shower_rotate_to_z (shower)
       !!! 6. rotate back in phi
       call shower_apply_lorentztrafo (shower, rotation &
            (cos(-phirand), sin(-phirand), vector3_canonical(3)))
    else
       if (prt%child2%t > parton_mass_squared(prt%child2)) then
          call shower_add_children_of_emitted_timelike_parton &
               (shower, prt%child2)
          call parton_set_simulated (prt%child2)
       end if

       call shower%add_parent (prt)
       call shower%add_child (prt%parent, 2)

       prt%parent%momentum = prt%momentum
       prt%parent%t = prt%t
       prt%parent%x = prt%x
       prt%parent%initial => prt%initial
       prt%parent%belongstoFSR = .false.

       prta => prt
       prtb => prt%child1
       prtc => prt%child2
    end if
    if (signal_is_pending ()) return
    if (isr_pt_ordered) then
       call parton_generate_ps_ini (prt%parent)
    else
       call parton_generate_ps_ini (prt)
    end if

    !!! add color connections
    if (parton_is_quark (prtb)) then

       if (prta%type == prtb%type) then
          if (prtb%type > 0) then
             !!! quark -> quark + gluon
             prtc%c2 = prtb%c1
             prtc%c1 = shower%get_next_color_nr ()
             prta%c1 = prtc%c1
          else
             !!! antiquark -> antiquark + gluon
             prtc%c1 = prtb%c2
             prtc%c2 = shower%get_next_color_nr ()
             prta%c2 = prtc%c2
          end if
       else
          !!! gluon -> quark + antiquark
          if (prtb%type > 0) then
             !!! gluon -> quark + antiquark
             prta%c1 = prtb%c1
             prtc%c1 = 0
             prtc%c2 = shower%get_next_color_nr ()
             prta%c2 = prtc%c2
          else
             !!! gluon -> antiquark + quark
             prta%c2 = prtb%c2
             prtc%c1 = shower%get_next_color_nr ()
             prtc%c2 = 0
             prta%c1 = prtc%c1
            end if
         end if
      else if (parton_is_gluon (prtb)) then
         if (parton_is_gluon (prta)) then
            !!! g -> gg
            prtc%c2 = prtb%c1
            prtc%c1 = shower%get_next_color_nr ()
            prta%c1 = prtc%c1
            prta%c2 = prtb%c2
         else if (parton_is_quark (prta)) then
            if (prta%type > 0) then
               prta%c1 = prtb%c1
               prta%c2 = 0
               prtc%c1 = prtb%c2
               prtc%c2 = 0
            else
               prta%c1 = 0
               prta%c2 = prtb%c2
               prtc%c1 = 0
               prtc%c2 = prtb%c1
            end if
         end if
      end if

      call shower%sort_partons ()
      call shower_boost_to_CMframe (shower)
      call shower_rotate_to_z (shower)

    end subroutine shower_execute_next_isr_branching

\eatline
\nwendcode{}\nwbegindocs{317}\nwdocspar
\nwenddocs{}\nwbegincode{318}\moddef{Shower core: procedures}\plusendmoddef
  subroutine shower_remove_parents_and_stuff (shower, prt)
    type(shower_t), intent(inout) :: shower
    type(parton_t), intent(inout), target :: prt
    type(parton_t), pointer :: actprt, nextprt
    nextprt => prt%parent
    actprt => null()
    !!! remove children of emitted timelike parton
    if (associated (prt%child2)) then
       if (associated (prt%child2%child1)) then
          call shower_remove_parton_from_partons_recursive &
               (shower, prt%child2%child1)
       end if
       prt%child2%child1 => null()
       if (associated (prt%child2%child2)) then
          call shower_remove_parton_from_partons_recursive &
               (shower, prt%child2%child2)
       end if
       prt%child2%child2 => null()
    end if
    do
       actprt => nextprt
       if (.not. associated (actprt)) then
          exit
       else if (parton_is_hadron (actprt)) then
          !!! remove beam-remnant
          call shower_remove_parton_from_partons (shower, actprt%child2)
          exit
       end if
       if (associated (actprt%parent)) then
          nextprt => actprt%parent
       else
          nextprt => null()
       end if
       call shower_remove_parton_from_partons_recursive &
            (shower, actprt%child2)
       call shower_remove_parton_from_partons (shower, actprt)

    end do
    prt%parent=>null()

  end subroutine shower_remove_parents_and_stuff

\eatline
\nwendcode{}\nwbegindocs{319}\nwdocspar
\nwenddocs{}\nwbegincode{320}\moddef{Shower core: shower: TBP}\plusendmoddef
  procedure :: get_ISR_scale => shower_get_ISR_scale
\nwendcode{}\nwbegincode{321}\moddef{Shower core: procedures}\plusendmoddef
  function shower_get_ISR_scale (shower) result (scale)
    class(shower_t), intent(in) :: shower
    real(default) :: scale
    type(parton_t), pointer :: prt1, prt2
    integer :: i
    scale = zero
    do i = 1, size (shower%interactions)
       call interaction_find_partons_nearest_to_hadron &
            (shower%interactions(i)%i, prt1, prt2)
       if (.not. prt1%simulated .and. abs(prt1%scale) > scale) &
            scale = abs(prt1%scale)
       if (.not. prt1%simulated .and. abs(prt2%scale) > scale) &
            scale = abs(prt2%scale)
    end do
  end function shower_get_ISR_scale

\eatline
\nwendcode{}\nwbegindocs{322}\nwdocspar
\nwenddocs{}\nwbegincode{323}\moddef{Shower core: shower: TBP}\plusendmoddef
  procedure :: set_max_isr_scale => shower_set_max_isr_scale
\nwendcode{}\nwbegincode{324}\moddef{Shower core: procedures}\plusendmoddef
  subroutine shower_set_max_isr_scale (shower, newscale)
    class(shower_t), intent(inout) :: shower
    real(default), intent(in) :: newscale
    real(default) :: scale
    type(parton_t), pointer :: prt
    integer :: i,j
    print *, "shower_set_max_isr_scale", newscale
    call shower%write ()

    if (isr_pt_ordered) then
       scale = newscale
    else
       scale = -abs(newscale)
    end if

    INTERACTIONS: do i = 1, size (shower%interactions)
       PARTONS: do j = 1, 2
          prt => shower%interactions(i)%i%partons(j)%p
          do
             if (.not. isr_pt_ordered) then
                if (prt%belongstointeraction) prt => prt%parent
             end if
               if (prt%t < scale) then
                  if (associated (prt%parent)) then
                     prt => prt%parent
                  else
                     exit   !!! unresolved prt found
                  end if
               else
                  exit   !!! prt with scale above newscale found
               end if
            end do
            if (.not. isr_pt_ordered) then
               if (prt%child1%belongstointeraction .or. &
                    parton_is_hadron (prt)) then
                  !!! don't reset scales of "first" spacelike partons
                  !!!    in virtuality ordered shower or hadrons
                  cycle
               end if
            else
               if (parton_is_hadron (prt)) then
                  !!! don't reset scales of hadrons
                  cycle
               end if
            end if
            if (isr_pt_ordered) then
               prt%scale = scale
            else
               prt%t = scale
            end if
            call parton_set_simulated (prt, .false.)
            call shower_remove_parents_and_stuff (shower, prt)
         end do PARTONS
      end do INTERACTIONS

      call shower%write ()
      print *, "shower_set_max_isr_scale finished"
    end subroutine shower_set_max_isr_scale

\eatline
\nwendcode{}\nwbegindocs{325}\nwdocspar
\nwenddocs{}\nwbegincode{326}\moddef{Shower core: shower: TBP}\plusendmoddef
  procedure :: interaction_generate_fsr_2ton => &
       shower_interaction_generate_fsr_2ton
\nwendcode{}\nwbegindocs{327}\nwdocspar
\nwenddocs{}\nwbegincode{328}\moddef{Shower core: procedures}\plusendmoddef
  subroutine shower_interaction_generate_fsr_2ton (shower, interaction)
    class(shower_t), intent(inout) :: shower
    type(shower_interaction_t), intent(inout) :: interaction
    type(parton_t), pointer :: prt
    prt => interaction%partons(3)%p
    do
       if (.not. associated (prt%parent)) exit
       prt => prt%parent
    end do
    call shower_parton_generate_fsr (shower, prt)
    call shower_parton_update_color_connections (shower, prt)
  end subroutine shower_interaction_generate_fsr_2ton

\eatline
\nwendcode{}\nwbegindocs{329} Perform the FSR for one parton, it is assumed, that the parton
already branched.  Hence, its children are to be simulated.  This
procedure is intended for branched FSR-partons emitted in the ISR.
\nwenddocs{}\nwbegincode{330}\moddef{Shower core: procedures}\plusendmoddef
  subroutine shower_parton_generate_fsr (shower, prt)
    type(shower_t), intent(inout) :: shower
    type(parton_t), intent(inout), target :: prt
    type(parton_pointer_t), dimension(:), allocatable :: partons
    if (signal_is_pending ()) return
    if (ASSERT) then
       if (.not. parton_is_branched (prt)) then
          print *, " error in shower_parton_generate_fsr:", &
               "parton not branched"
          return
       end if
       if (prt%child1%simulated .or. &
           prt%child2%simulated) then
          print *, " error in shower_parton_generate_fsr:", &
               "children already simulated for parton ", prt%nr
          call shower%write ()
          return
       end if
    end if

    allocate (partons(1:1))
    partons(1)%p => prt
    call shower_parton_pointer_array_generate_fsr (shower, partons)

  end subroutine shower_parton_generate_fsr

\eatline
\nwendcode{}\nwbegindocs{331}\nwdocspar
\nwenddocs{}\nwbegincode{332}\moddef{Shower core: procedures}\plusendmoddef
    recursive subroutine shower_parton_pointer_array_generate_fsr &
         (shower, partons)
      type(shower_t), intent(inout) :: shower
      type(parton_pointer_t), dimension(:), allocatable, intent(inout) :: &
           partons
      type(parton_pointer_t), dimension(:), allocatable :: partons_new
      integer :: i, size_partons, size_partons_new
      size_partons = size (partons)
      if (signal_is_pending ()) return
      if (size_partons == 0) return
      !!! sort partons -> necessary ?? -> probably not
      !!! simulate highest/first parton
      call shower_simulate_children_ana (shower, partons(1)%p)
      !!! check for new daughters to be included in new_partons
      size_partons_new = size_partons - 1   !!! partons(1) not needed anymore
      if (parton_is_branched (partons(1)%p%child1)) &
           size_partons_new = size_partons_new + 1
      if (parton_is_branched (partons(1)%p%child2)) &
           size_partons_new = size_partons_new + 1

      allocate (partons_new(1:size_partons_new))

      if (size_partons > 1) then
         do i = 2, size_partons
            partons_new(i - 1)%p => partons(i)%p
         end do
      end if
      if (parton_is_branched (partons(1)%p%child1)) &
           partons_new(size_partons)%p => partons(1)%p%child1
      if (parton_is_branched (partons(1)%p%child2)) then
         !!! check if child1 is already included
         if (size_partons_new == size_partons) then
            partons_new(size_partons)%p => partons(1)%p%child2
         else if (size_partons_new == size_partons + 1) then
            partons_new(size_partons + 1)%p => partons(1)%p%child2
         else
            call msg_fatal ("Shower: wrong sizes in" &
                 // "shower_parton_pointer_array_generate_fsr")
         end if
      end if
      deallocate(partons)

      call shower_parton_pointer_array_generate_fsr (shower, partons_new)

    end subroutine shower_parton_pointer_array_generate_fsr

\eatline
\nwendcode{}\nwbegindocs{333}\nwdocspar
\nwenddocs{}\nwbegincode{334}\moddef{Shower core: procedures}\plusendmoddef
  recursive subroutine shower_parton_update_color_connections &
       (shower, prt)
    type(shower_t), intent(inout) :: shower
    type(parton_t), intent(inout) :: prt
    real(default) :: temprand
    if (.not. associated (prt%child1) .or. &
        .not. associated (prt%child2)) return

    if (signal_is_pending ()) return
    if (parton_is_gluon (prt)) then
       if (parton_is_quark (prt%child1)) then
          !!! give the quark the colorpartner and the antiquark
          !!!     the anticolorpartner
          if (prt%child1%type > 0) then
             !!! child1 is quark, child2 is antiquark
             prt%child1%c1 = prt%c1
             prt%child2%c2 = prt%c2
          else
             !!! child1 is antiquark, child2 is quark
             prt%child1%c2 = prt%c2
             prt%child2%c1 = prt%c1
          end if
       else
          !!! g -> gg splitting -> random choosing of partners
          call rng%generate (temprand)
          if (temprand > 0.5_default) then
             prt%child1%c1 = prt%c1
             prt%child1%c2 = shower%get_next_color_nr ()
             prt%child2%c1 = prt%child1%c2
             prt%child2%c2 = prt%c2
          else
             prt%child1%c2 = prt%c2
             prt%child2%c1 = prt%c1
             prt%child2%c2 = shower%get_next_color_nr ()
             prt%child1%c1 = prt%child2%c2
          end if
       end if
    else if (parton_is_quark (prt)) then
       if (parton_is_quark (prt%child1)) then
          if (prt%child1%type > 0) then
             !!! q -> q + g
             prt%child2%c1 = prt%c1
             prt%child2%c2 = shower%get_next_color_nr ()
             prt%child1%c1 = prt%child2%c2
          else
             !!! qbar -> qbar + g
             prt%child2%c2 = prt%c2
             prt%child2%c1 = shower%get_next_color_nr ()
             prt%child1%c2 = prt%child2%c1
          end if
       else
          if (prt%child2%type > 0) then
             !!! q -> g + q
             prt%child1%c1 = prt%c1
             prt%child1%c2 = shower%get_next_color_nr ()
             prt%child2%c1 = prt%child1%c2
          else
             !!! qbar -> g + qbar
             prt%child1%c2 = prt%c2
             prt%child1%c1 = shower%get_next_color_nr ()
             prt%child2%c2 = prt%child1%c1
          end if
       end if
    end if

    call shower_parton_update_color_connections (shower, prt%child1)
    call shower_parton_update_color_connections (shower, prt%child2)
  end subroutine shower_parton_update_color_connections

\eatline
\nwendcode{}\nwbegindocs{335}The next two routines are for PDFs. Wrapper function to return
parton densities.
\nwenddocs{}\nwbegincode{336}\moddef{Shower core: shower: TBP}\plusendmoddef
  procedure :: get_pdf => shower_get_pdf
\nwendcode{}\nwbegincode{337}\moddef{Shower core: procedures}\plusendmoddef
  function shower_get_pdf (shower, mother, x, Q2, daughter) result (pdf)
    \LA{}get pdf\RA{}
    if (x > eps0) then
       pdf = pdf / x
    end if
  end function shower_get_pdf

\eatline
\nwendcode{}\nwbegincode{338}\moddef{Shower core: shower: TBP}\plusendmoddef
  procedure :: get_xpdf => shower_get_xpdf
\nwendcode{}\nwbegincode{339}\moddef{Shower core: procedures}\plusendmoddef
  function shower_get_xpdf (shower, mother, x, Q2, daughter) result (pdf)
    \LA{}get pdf\RA{}
  end function shower_get_xpdf

\eatline
\nwendcode{}\nwbegindocs{340}\nwdocspar
\nwenddocs{}\nwbegincode{341}\moddef{get pdf}\endmoddef
class(shower_t), intent(inout), target :: shower
integer, intent(in) :: mother, daughter
real(default), intent(in) :: x, Q2
real(default) :: pdf
real(double), save :: f(-6:6) = 0._double
real(double), save :: lastx, lastQ2 = 0._double
if (ASSERT) then
   if (abs (mother) /= PROTON) then
      pdf = zero
      print *, "mother = ", mother
      call msg_fatal ("Shower: pdf only implemented for (anti-)proton")
   end if
   if (.not. (abs (daughter) >= 1 .and. abs (daughter) <= 6 .or. &
              daughter == GLUON)) then
      pdf = zero
      print *, "daughter = ", daughter
      call msg_fatal ("Shower: error in pdf, unknown daughter")
   end if
end if
if (x > zero .and. x < one) then
   if ((dble(Q2) - lastQ2) > eps0 .or. (dble(x) - lastx) > eps0) then
      call shower%pdf_func &
           (shower%pdf_set, dble(x), sqrt (abs (dble(Q2))), f)
   end if
   if (abs (daughter) >= 1 .and. abs (daughter) <= 6) then
      pdf = max (f(daughter * sign (1,mother)), tiny_10)
   else
      pdf = max (f(0), tiny_10)
   end if
else
   pdf = zero
end if
lastQ2 = dble(Q2)
lastx  = dble(x)
\nwendcode{}\nwbegindocs{342}\nwdocspar
\nwenddocs{}\nwbegincode{343}\moddef{Shower core: shower: TBP}\plusendmoddef
  procedure :: pdf_func => shower_pdf_func
\nwendcode{}\nwbegincode{344}\moddef{Shower core: procedures}\plusendmoddef
  subroutine shower_pdf_func (shower, set, x, q, f)
    class(shower_t), intent(inout) :: shower
    integer, intent(in) :: set
    real(double) :: x, q
    real(double), dimension(-6:6), intent(out) :: f
    select case (shower%pdf_type)
    case (STRF_PDF_BUILTIN)
       call pdf_evolve_LHAPDF (set, x, q, f)
    case (STRF_LHAPDF6)
       q = min (shower%qmax, q)
       q = max (shower%qmin, q)
       call shower%pdf%evolve_pdfm (x, q, f)
    case (STRF_LHAPDF5)
       q = min (shower%qmax, q)
       q = max (shower%qmin, q)
       call evolvePDFM (set, x, q, f)
    case default
       call msg_fatal ("Shower PDF function: unknown PDF method.")
    end select
  end subroutine shower_pdf_func
\nwendcode{}\nwbegindocs{345}\nwdocspar
%%%%%%%%%%%%%%%%%%%%%%%%%%%%%%%%%%%%%%%%%%%%%%%%%%%%%%%%%%%%%%%%%%%%%%%%
\section{Talking to PYTHIA}
\nwenddocs{}\nwbegincode{346}\moddef{\code{}shower{\char95}topythia.f90\edoc{}}\endmoddef
\LA{}File header\RA{}

module shower_topythia

\LA{}Use kinds\RA{}
  use physics_defs
  use lorentz
  use shower_base
  use shower_partons
  use shower_core

\LA{}Standard module head\RA{}

\LA{}Shower2pythia: public\RA{}

contains

\LA{}Shower2pythia: procedures\RA{}

  end module shower_topythia

\eatline
\nwendcode{}\nwbegindocs{347}\nwdocspar
\nwenddocs{}\nwbegincode{348}\moddef{Shower2pythia: public}\endmoddef
  public :: shower_converttopythia
\nwendcode{}\nwbegindocs{349}\nwdocspar
\nwenddocs{}\nwbegincode{350}\moddef{Shower2pythia: procedures}\endmoddef
  subroutine shower_converttopythia (shower)
    IMPLICIT DOUBLE PRECISION(A-H, O-Z)
    IMPLICIT INTEGER(I-N)
    !!!    C...  Commonblocks.
    COMMON/PYJETS/N,NPAD,K(4000,5),P(4000,5),V(4000,5)
    COMMON/PYDAT1/MSTU(200),PARU(200),MSTJ(200),PARJ(200)
    COMMON/PYDAT2/KCHG(500,4),PMAS(500,4),PARF(2000),VCKM(4,4)
    COMMON/PYDAT3/MDCY(500,3),MDME(8000,2),BRAT(8000),KFDP(8000,5)
    COMMON/PYSUBS/MSEL,MSELPD,MSUB(500),KFIN(2,-40:40),CKIN(200)
    COMMON/PYPARS/MSTP(200),PARP(200),MSTI(200),PARI(200)
    SAVE /PYJETS/,/PYDAT1/,/PYDAT2/,/PYDAT3/,/PYSUBS/,/PYPARS/

    type(shower_t), intent(in) :: shower
    type(parton_t), pointer :: pp, ppparent
    integer :: i, j, nz

    !!! currently only works for one interaction

    do i = 1, 2
       !!! get history of the event
       pp => shower%interactions(1)%i%partons(i)%p
       !!! add these partons to the event record
       if (associated (pp%initial)) then
          !!! add hadrons
          K(i,1) = GLUON
          K(i,2) = pp%initial%type
          K(i,3) = 0
          P(i,1) = pp%initial%momentum%p(1)
          P(i,2) = pp%initial%momentum%p(2)
          P(i,3) = pp%initial%momentum%p(3)
          P(i,4) = pp%initial%momentum%p(0)
          P(I,5) = pp%initial%momentum**2
          !!! add partons emitted by the hadron
          ppparent => pp
          do while (associated (ppparent%parent))
             if (parton_is_hadron (ppparent%parent)) then
                exit
             else
                ppparent => ppparent%parent
             end if
          end do
          K(i+2,1) = GLUON
          K(i+2,2) = ppparent%type
          K(i+2,3) = i
          P(i+2,1) = ppparent%momentum%p(1)
          P(i+2,2) = ppparent%momentum%p(2)
          P(i+2,3) = ppparent%momentum%p(3)
          P(i+2,4) = ppparent%momentum%p(0)
          P(I+2,5) = ppparent%momentum**2
          !!! add partons in the initial state of the ME
          K(i+4,1) = GLUON
          K(i+4,2) = pp%type
          K(i+4,3) = i
          P(i+4,1) = pp%momentum%p(1)
          P(i+4,2) = pp%momentum%p(2)
          P(i+4,3) = pp%momentum%p(3)
          P(i+4,4) = pp%momentum%p(0)
          P(I+4,5) = pp%momentum**2
       else
          !!! for e+e- without ISR all entries are the same
          K(i,1) = GLUON
          K(i,2) = pp%type
          K(i,3) = 0
          P(i,1) = pp%momentum%p(1)
          P(i,2) = pp%momentum%p(2)
          P(i,3) = pp%momentum%p(3)
          P(i,4) = pp%momentum%p(0)
          P(I,5) = pp%momentum**2
          do j = 1, 5
             P(i+2,j) = P(1,j)
             K(i+2,j) = K(1,j)
             K(i+2,3) = i
             P(i+4,j) = P(1,j)
             K(i+4,j) = K(1,j)
             K(i+4,3) = i
          end do
          P(i+4,5) = 0.
       end if
    end do
    N = 6
    !!! create intermediate (fake) Z-Boson
    K(7,1) = GLUON
    K(7,2) = 23
    K(7,3) = 0
    P(7,1) = P(5,1) + P(6,1)
    P(7,2) = P(5,2) + P(6,2)
    P(7,3) = P(5,3) + P(6,3)
    P(7,4) = P(5,4) + P(6,4)
    P(7,5) = P(7,4)**2 - P(7,3)**2 - P(7,2)**2 - P(7,1)**2
    N = 7
    !!! include partons in the final state of the hard matrix element
    do i = 1, size (shower%interactions(1)%i%partons) - 2
       !!! get partons that are in the final state of the hard matrix element
       pp => shower%interactions(1)%i%partons(2+i)%p
       !!! add these partons to the event record
       K(7+I,1) = GLUON
       K(7+I,2) = pp%type
       K(7+I,3) = 7
       P(7+I,1) = pp%momentum%p(1)
       P(7+I,2) = pp%momentum%p(2)
       P(7+I,3) = pp%momentum%p(3)
       P(7+I,4) = pp%momentum%p(0)
       P(7+I,5) = P(7+I,4)**2 - P(7+I,3)**2 - P(7+I,2)**2 - P(7+I,1)**2
       N = 7 + I
    end do
    !!! include "Z" (again)
    N = N + 1
    K(N,1) = 11
    K(N,2) = 23
    K(N,3) = 7
    P(N,1) = P(7,1)
    P(N,2) = P(7,2)
    P(N,3) = P(7,3)
    P(N,4) = P(7,4)
    P(N,5) = P(7,5)
    nz = N
    !!! include partons from the final state of the parton shower
    call shower_transfer_final_partons_to_pythia (shower, 8)
    !!! set "children" of "Z"
    K(nz,4) = 11
    K(nz,5) = N
    !!! mark spacers
    MSTU(73) = N
    MSTU(74) = N

    !!! be sure to remove the next partons (=first obsolete partons)
    K(N+1,1) = 0
    K(N+1,2) = 0
    K(N+1,3) = 0
    K(N+2,1) = 0
    K(N+2,2) = 0
    K(N+2,3) = 0
    K(N+3,1) = 0
    K(N+3,2) = 0
    K(N+3,3) = 0
    !!! otherwise they might be interpreted as thrust information
  end subroutine shower_converttopythia

\eatline
\nwendcode{}\nwbegindocs{351}\nwdocspar
\nwenddocs{}\nwbegincode{352}\moddef{Shower2pythia: procedures}\plusendmoddef
  subroutine shower_transfer_final_partons_to_pythia (shower, first)
    IMPLICIT DOUBLE PRECISION(A-H, O-Z)
    IMPLICIT INTEGER(I-N)
    !    C...  Commonblocks.
    COMMON/PYJETS/N,NPAD,K(4000,5),P(4000,5),V(4000,5)
    COMMON/PYDAT1/MSTU(200),PARU(200),MSTJ(200),PARJ(200)
    COMMON/PYDAT2/KCHG(500,4),PMAS(500,4),PARF(2000),VCKM(4,4)
    COMMON/PYDAT3/MDCY(500,3),MDME(8000,2),BRAT(8000),KFDP(8000,5)
    COMMON/PYSUBS/MSEL,MSELPD,MSUB(500),KFIN(2,-40:40),CKIN(200)
    COMMON/PYPARS/MSTP(200),PARP(200),MSTI(200),PARI(200)
    SAVE /PYJETS/,/PYDAT1/,/PYDAT2/,/PYDAT3/,/PYSUBS/,/PYPARS/

    type(shower_t), intent(in) :: shower
    integer, intent(in) :: first
    type(parton_t), pointer :: prt
    integer :: i, j, n_finals
    type(parton_t), dimension(:), allocatable :: final_partons
    type(parton_t) :: temp_parton
    integer :: minindex, maxindex

    prt => null()

    !!! get total number of final partons
    n_finals = 0
    do i = 1, size (shower%partons)
       if (.not. associated (shower%partons(i)%p)) cycle
       prt => shower%partons(i)%p
       if (.not. prt%belongstoFSR) cycle
       if (associated (prt%child1)) cycle
       n_finals = n_finals + 1
    end do

    allocate (final_partons(1:n_finals))
    j = 1
    do i = 1, size (shower%partons)
       if (.not. associated (shower%partons(i)%p)) cycle
       prt => shower%partons(i)%p
       if (.not. prt%belongstoFSR) cycle
       if (associated (prt%child1)) cycle
       final_partons(j) = shower%partons(i)%p
       j = j + 1
    end do

    !!! move quark to front as beginning of color string
    minindex = 1
    maxindex = size (final_partons)
    FIND_Q: do i = minindex, maxindex
       if (final_partons(i)%type >= 1 .and. final_partons(i)%type <= 6) then
          temp_parton = final_partons(minindex)
          final_partons(minindex) = final_partons(i)
          final_partons(i) = temp_parton
          exit FIND_Q
       end if
    end do FIND_Q

    !!! sort so that connected partons are next to each other, don't care about zeros
    do i = 1, size (final_partons)
       !!! ensure that final_partnons begins with a color (not an anticolor)
       if (final_partons(i)%c1 > 0 .and. final_partons(i)%c2 == 0) then
          if (i == 1) then
             exit
          else
             temp_parton = final_partons(1)
             final_partons(1) = final_partons(i)
             final_partons(i) = temp_parton
             exit
          end if
       end if
    end do

    do i = 1, size (final_partons) - 1
       !!! search for color partner and move it to i + 1
       PARTNERS: do j = i + 1, size (final_partons)
          if (final_partons(j)%c2 == final_partons(i)%c1) exit PARTNERS
       end do PARTNERS
       if (j > size (final_partons)) then
          print *, "no color connected parton found" !WRONG???
          print *, "particle: ", final_partons(i)%nr, " index: ", &
               final_partons(i)%c1
          exit
       end if
       temp_parton = final_partons(i + 1)
       final_partons(i + 1) = final_partons(j)
       final_partons(j) = temp_parton
    end do

    !!! transfering partons
    do i = 1, size (final_partons)
       prt = final_partons(i)
       N = N + 1
       K(N,1) = 2
       if (prt%c1 == 0) K(N,1) = 1       !!! end of color string
       K(N,2) = prt%type
       K(N,3) = first
       K(N,4) = 0
       K(N,5) = 0
       P(N,1) = prt%momentum%p(1)
       P(N,2) = prt%momentum%p(2)
       P(N,3) = prt%momentum%p(3)
       P(N,4) = prt%momentum%p(0)
       P(N,5) = prt%momentum**2
    end do
    deallocate (final_partons)
  end subroutine shower_transfer_final_partons_to_pythia

\eatline
\nwendcode{}\nwbegindocs{353}\nwdocspar
%%%%%%%%%%%%%%%%%%%%%%%%%%%%%%%%%%%%%%%%%%%%%%%%%%%%%%%%%%%%%%%%%%%%%%%%
\nwenddocs{}\nwbegincode{354}\moddef{\code{}pythia6{\char95}up.f\edoc{}}\endmoddef
C*********************************************************************
C*********************************************************************
C*                                                                  **
C*                                                       Mar 2011   **
C*                                                                  **
C*                       The Lund Monte Carlo                       **
C*                                                                  **
C*                        PYTHIA version 6.4                        **
C*                                                                  **
C*                        Torbjorn Sjostrand                        **
C*                 Department of Theoretical Physics                **
C*                         Lund University                          **
C*               Solvegatan 14A, S-223 62 Lund, Sweden              **
C*                    E-mail torbjorn@thep.lu.se                    **
C*                                                                  **
C*                  SUSY and Technicolor parts by                   **
C*                         Stephen Mrenna                           **
C*                       Computing Division                         **
C*            Generators and Detector Simulation Group              **
C*              Fermi National Accelerator Laboratory               **
C*                 MS 234, Batavia, IL  60510, USA                  **
C*                   phone + 1 - 630 - 840 - 2556                   **
C*                      E-mail mrenna@fnal.gov                      **
C*                                                                  **
C*         New multiple interactions and more SUSY parts by         **
C*                          Peter Skands                            **
C*               CERN/PH, CH-1211 Geneva, Switzerland               **
C*                    phone +41 - 22 - 767 2447                     **
C*                   E-mail peter.skands@cern.ch                    **
C*                                                                  **
C*         Several parts are written by Hans-Uno Bengtsson          **
C*          PYSHOW is written together with Mats Bengtsson          **
C*               PYMAEL is written by Emanuel Norrbin               **
C*     advanced popcorn baryon production written by Patrik Eden    **
C*    code for virtual photons mainly written by Christer Friberg   **
C*    code for low-mass strings mainly written by Emanuel Norrbin   **
C*        Bose-Einstein code mainly written by Leif Lonnblad        **
C*      CTEQ  parton distributions are by the CTEQ collaboration    **
C*      GRV 94 parton distributions are by Glueck, Reya and Vogt    **
C*   SaS photon parton distributions together with Gerhard Schuler  **
C*     g + g and q + qbar -> t + tbar + H code by Zoltan Kunszt     **
C*         MSSM Higgs mass calculation code by M. Carena,           **
C*           J.R. Espinosa, M. Quiros and C.E.M. Wagner             **
C*  UED implementation by M. Elkacimi, D. Goujdami, H. Przysiezniak **
C*         PYGAUS adapted from CERN library (K.S. Kolbig)           **
C*        NRQCD/colour octet production of onium by S. Wolf         **
C*                                                                  **
C*   The latest program version and documentation is found on WWW   **
C*            http://www.thep.lu.se/~torbjorn/Pythia.html           **
C*                                                                  **
C*              Copyright Torbjorn Sjostrand, Lund 2010             **
C*                                                                  **
C*********************************************************************
C*********************************************************************

C...UPINIT
C...Dummy routine, to be replaced by a user implementing external
C...processes. Is supposed to fill the HEPRUP commonblock with info
C...on incoming beams and allowed processes.

C...New example: handles a standard Les Houches Events File.

      SUBROUTINE UPINIT

C...Double precision and integer declarations.
      IMPLICIT DOUBLE PRECISION(A-H, O-Z)
      IMPLICIT INTEGER(I-N)

C...PYTHIA commonblock: only used to provide read unit MSTP(161).
      COMMON/PYPARS/MSTP(200),PARP(200),MSTI(200),PARI(200)
      SAVE /PYPARS/

C...User process initialization commonblock.
      INTEGER MAXPUP
      PARAMETER (MAXPUP=100)
      INTEGER IDBMUP,PDFGUP,PDFSUP,IDWTUP,NPRUP,LPRUP
      DOUBLE PRECISION EBMUP,XSECUP,XERRUP,XMAXUP
      COMMON/HEPRUP/IDBMUP(2),EBMUP(2),PDFGUP(2),PDFSUP(2),
     &IDWTUP,NPRUP,XSECUP(MAXPUP),XERRUP(MAXPUP),XMAXUP(MAXPUP),
     &LPRUP(MAXPUP)
      SAVE /HEPRUP/

C...Lines to read in assumed never longer than 200 characters.
      PARAMETER (MAXLEN=200)
      CHARACTER*(MAXLEN) STRING

C...Format for reading lines.
      CHARACTER(len=6) STRFMT
      STRFMT='(A000)'
      WRITE(STRFMT(3:5),'(I3)') MAXLEN

C...Loop until finds line beginning with "<init>" or "<init ".
  100 READ(MSTP(161),STRFMT,END=130,ERR=130) STRING
      IBEG=0
  110 IBEG=IBEG+1
C...Allow indentation.
      IF(STRING(IBEG:IBEG).EQ.' '.AND.IBEG.LT.MAXLEN-5) GOTO 110
      IF(STRING(IBEG:IBEG+5).NE.'<init>'.AND.
     &STRING(IBEG:IBEG+5).NE.'<init ') GOTO 100

C...Read first line of initialization info.
      READ(MSTP(161),*,END=130,ERR=130) IDBMUP(1),IDBMUP(2),EBMUP(1),
     &EBMUP(2),PDFGUP(1),PDFGUP(2),PDFSUP(1),PDFSUP(2),IDWTUP,NPRUP

C...Read NPRUP subsequent lines with information on each process.
      DO 120 IPR=1,NPRUP
        READ(MSTP(161),*,END=130,ERR=130) XSECUP(IPR),XERRUP(IPR),
     &  XMAXUP(IPR),LPRUP(IPR)
  120 CONTINUE
      RETURN

C...Error exit: give up if initalization does not work.
  130 WRITE(*,*) ' Failed to read LHEF initialization information.'
      WRITE(*,*) ' Event generation will be stopped.'
      CALL PYSTOP(12)

      RETURN
      END

C...Old example: handles a simple Pythia 6.4 initialization file.

c      SUBROUTINE UPINIT

C...Double precision and integer declarations.
c      IMPLICIT DOUBLE PRECISION(A-H, O-Z)
c      IMPLICIT INTEGER(I-N)

C...Commonblocks.
c      COMMON/PYDAT1/MSTU(200),PARU(200),MSTJ(200),PARJ(200)
c      COMMON/PYPARS/MSTP(200),PARP(200),MSTI(200),PARI(200)
c      SAVE /PYDAT1/,/PYPARS/

C...User process initialization commonblock.
c      INTEGER MAXPUP
c      PARAMETER (MAXPUP=100)
c      INTEGER IDBMUP,PDFGUP,PDFSUP,IDWTUP,NPRUP,LPRUP
c      DOUBLE PRECISION EBMUP,XSECUP,XERRUP,XMAXUP
c      COMMON/HEPRUP/IDBMUP(2),EBMUP(2),PDFGUP(2),PDFSUP(2),
c     &IDWTUP,NPRUP,XSECUP(MAXPUP),XERRUP(MAXPUP),XMAXUP(MAXPUP),
c     &LPRUP(MAXPUP)
c      SAVE /HEPRUP/

C...Read info from file.
c      IF(MSTP(161).GT.0) THEN
c        READ(MSTP(161),*,END=110,ERR=110) IDBMUP(1),IDBMUP(2),EBMUP(1),
c     &  EBMUP(2),PDFGUP(1),PDFGUP(2),PDFSUP(1),PDFSUP(2),IDWTUP,NPRUP
c        DO 100 IPR=1,NPRUP
c          READ(MSTP(161),*,END=110,ERR=110) XSECUP(IPR),XERRUP(IPR),
c     &    XMAXUP(IPR),LPRUP(IPR)
c  100   CONTINUE
c        RETURN
C...Error or prematurely reached end of file.
c  110   WRITE(MSTU(11),5000)
c        STOP

C...Else not implemented.
c      ELSE
c        WRITE(MSTU(11),5100)
c        STOP
c      ENDIF

C...Format for error printout.
c 5000 FORMAT(1X,'Error: UPINIT routine failed to read information'/
c     &1X,'Execution stopped!')
c 5100 FORMAT(1X,'Error: You have not implemented UPINIT routine'/
c     &1X,'Dummy routine in PYTHIA file called instead.'/
c     &1X,'Execution stopped!')

c      RETURN
c      END

C*********************************************************************

C...UPEVNT
C...Dummy routine, to be replaced by a user implementing external
C...processes. Depending on cross section model chosen, it either has
C...to generate a process of the type IDPRUP requested, or pick a type
C...itself and generate this event. The event is to be stored in the
C...HEPEUP commonblock, including (often) an event weight.

C...New example: handles a standard Les Houches Events File.

      SUBROUTINE UPEVNT

C...Double precision and integer declarations.
      IMPLICIT DOUBLE PRECISION(A-H, O-Z)
      IMPLICIT INTEGER(I-N)

C...PYTHIA commonblock: only used to provide read unit MSTP(162).
      COMMON/PYPARS/MSTP(200),PARP(200),MSTI(200),PARI(200)
      SAVE /PYPARS/

C...Added by WHIZARD
      COMMON/PYDAT1/MSTU(200),PARU(200),MSTJ(200),PARJ(200)
      SAVE/PYDAT1/

C...User process event common block.
      INTEGER MAXNUP
      PARAMETER (MAXNUP=500)
      INTEGER NUP,IDPRUP,IDUP,ISTUP,MOTHUP,ICOLUP
      DOUBLE PRECISION XWGTUP,SCALUP,AQEDUP,AQCDUP,PUP,VTIMUP,SPINUP
      COMMON/HEPEUP/NUP,IDPRUP,XWGTUP,SCALUP,AQEDUP,AQCDUP,IDUP(MAXNUP),
     &ISTUP(MAXNUP),MOTHUP(2,MAXNUP),ICOLUP(2,MAXNUP),PUP(5,MAXNUP),
     &VTIMUP(MAXNUP),SPINUP(MAXNUP)
      SAVE /HEPEUP/

C...Lines to read in assumed never longer than 200 characters.
      PARAMETER (MAXLEN=200)
      CHARACTER*(MAXLEN) STRING

C...Format for reading lines.
      CHARACTER(len=6) STRFMT
      STRFMT='(A000)'
      WRITE(STRFMT(3:5),'(I3)') MAXLEN

C...Loop until finds line beginning with "<event>" or "<event ".
  100 READ(MSTP(162),STRFMT,END=130,ERR=130) STRING
      IBEG=0
  110 IBEG=IBEG+1
C...Allow indentation.
      IF(STRING(IBEG:IBEG).EQ.' '.AND.IBEG.LT.MAXLEN-6) GOTO 110
      IF(STRING(IBEG:IBEG+6).NE.'<event>'.AND.
     &STRING(IBEG:IBEG+6).NE.'<event ') GOTO 100

C...Read first line of event info.
      READ(MSTP(162),*,END=130,ERR=130) NUP,IDPRUP,XWGTUP,SCALUP,
     &AQEDUP,AQCDUP

C...Read NUP subsequent lines with information on each particle.
      DO 120 I=1,NUP
        READ(MSTP(162),*,END=130,ERR=130) IDUP(I),ISTUP(I),
     &  MOTHUP(1,I),MOTHUP(2,I),ICOLUP(1,I),ICOLUP(2,I),
     &  (PUP(J,I),J=1,5),VTIMUP(I),SPINUP(I)
  120 CONTINUE
      RETURN

C...Error exit, typically when no more events.
  130 CONTINUE
C      WRITE(*,*) ' Failed to read LHEF event information.'
C      WRITE(*,*) ' Will assume end of file has been reached.'
      NUP=0
      MSTI(51)=1
C...Added by WHIZARD, mark these failed events
      MSTU(23)=1

      RETURN
      END

C...Old example: handles a simple Pythia 6.4 event file.

c      SUBROUTINE UPEVNT

C...Double precision and integer declarations.
c      IMPLICIT DOUBLE PRECISION(A-H, O-Z)
c      IMPLICIT INTEGER(I-N)

C...Commonblocks.
c      COMMON/PYDAT1/MSTU(200),PARU(200),MSTJ(200),PARJ(200)
c      COMMON/PYPARS/MSTP(200),PARP(200),MSTI(200),PARI(200)
c      SAVE /PYDAT1/,/PYPARS/

C...User process event common block.
c      INTEGER MAXNUP
c      PARAMETER (MAXNUP=500)
c      INTEGER NUP,IDPRUP,IDUP,ISTUP,MOTHUP,ICOLUP
c      DOUBLE PRECISION XWGTUP,SCALUP,AQEDUP,AQCDUP,PUP,VTIMUP,SPINUP
c      COMMON/HEPEUP/NUP,IDPRUP,XWGTUP,SCALUP,AQEDUP,AQCDUP,IDUP(MAXNUP),
c     &ISTUP(MAXNUP),MOTHUP(2,MAXNUP),ICOLUP(2,MAXNUP),PUP(5,MAXNUP),
c     &VTIMUP(MAXNUP),SPINUP(MAXNUP)
c      SAVE /HEPEUP/

C...Read info from file.
c      IF(MSTP(162).GT.0) THEN
c        READ(MSTP(162),*,END=110,ERR=110) NUP,IDPRUP,XWGTUP,SCALUP,
c     &  AQEDUP,AQCDUP
c        DO 100 I=1,NUP
c          READ(MSTP(162),*,END=110,ERR=110) IDUP(I),ISTUP(I),
c     &    MOTHUP(1,I),MOTHUP(2,I),ICOLUP(1,I),ICOLUP(2,I),
c     &    (PUP(J,I),J=1,5),VTIMUP(I),SPINUP(I)
c  100   CONTINUE
c        RETURN
C...Special when reached end of file or other error.
c  110   NUP=0

C...Else not implemented.
c      ELSE
c        WRITE(MSTU(11),5000)
c        STOP
c      ENDIF

C...Format for error printout.
c 5000 FORMAT(1X,'Error: You have not implemented UPEVNT routine'/
c     &1X,'Dummy routine in PYTHIA file called instead.'/
c     &1X,'Execution stopped!')

c      RETURN
c      END

C*********************************************************************

C...UPVETO
C...Dummy routine, to be replaced by user, to veto event generation
C...on the parton level, after parton showers but before multiple
C...interactions, beam remnants and hadronization is added.
C...If resonances like W, Z, top, Higgs and SUSY particles are handed
C...undecayed from UPEVNT, or are generated by PYTHIA, they will also
C...be undecayed at this stage; if decayed their decay products will
C...have been allowed to shower.

C...All partons at the end of the shower phase are stored in the
C...HEPEVT commonblock. The interesting information is
C...NHEP = the number of such partons, in entries 1 <= i <= NHEP,
C...IDHEP(I) = the particle ID code according to PDG conventions,
C...PHEP(J,I) = the (p_x, p_y, p_z, E, m) of the particle.
C...All ISTHEP entries are 1, while the rest is zeroed.

C...The user decision is to be conveyed by the IVETO value.
C...IVETO = 0 : retain current event and generate in full;
C...      = 1 : abort generation of current event and move to next.

      SUBROUTINE UPVETO(IVETO)

C...HEPEVT commonblock.
      PARAMETER (NMXHEP=4000)
      COMMON/HEPEVT/NEVHEP,NHEP,ISTHEP(NMXHEP),IDHEP(NMXHEP),
     &JMOHEP(2,NMXHEP),JDAHEP(2,NMXHEP),PHEP(5,NMXHEP),VHEP(4,NMXHEP)
      DOUBLE PRECISION PHEP,VHEP
      SAVE /HEPEVT/

C...Next few lines allow you to see what info PYVETO extracted from
C...the full event record for the first two events.
C...Delete if you don't want it.
      DATA NLIST/0/
      SAVE NLIST
      IF(NLIST.LE.2) THEN
        WRITE(*,*) ' Full event record at time of UPVETO call:'
        CALL PYLIST(1)
        WRITE(*,*) ' Part of event record made available to UPVETO:'
        CALL PYLIST(5)
        NLIST=NLIST+1
      ENDIF

C...Make decision here.
      IVETO = 0

      RETURN
      END

\eatline
\nwendcode{}\nwbegindocs{355}\nwdocspar
%%%%%%%%%%%%%%%%%%%%%%%%%%%%%%%%%%%%%%%%%%%%%%%%%%%%%%%%%%%%%%%%%%%%%%%%
\section{Matching}
\subsection{MLM matching}
\nwenddocs{}\nwbegincode{356}\moddef{\code{}mlm{\char95}matching.f90\edoc{}}\endmoddef
\LA{}File header\RA{}

module mlm_matching

\LA{}Use kinds with double\RA{}
  use io_units
  use constants
  use format_utils, only: write_separator
  use diagnostics
  use file_utils
  use lorentz

\LA{}Standard module head\RA{}

\LA{}MLM matching: public\RA{}

\LA{}MLM matching: types\RA{}

contains

\LA{}MLM matching: procedures\RA{}

end module mlm_matching
\eatline
\nwendcode{}\nwbegindocs{357}This is a container for the (colored) parton momenta as well as the
jet momenta.
\nwenddocs{}\nwbegincode{358}\moddef{MLM matching: public}\endmoddef
  public :: mlm_matching_data_t
\nwendcode{}\nwbegincode{359}\moddef{MLM matching: types}\endmoddef
  type :: mlm_matching_data_t
     logical :: is_hadron_collision = .false.
     type(vector4_t), dimension(:), allocatable, public :: P_ME
     type(vector4_t), dimension(:), allocatable, public :: P_PS
     type(vector4_t), dimension(:), allocatable, private :: JETS_ME
     type(vector4_t), dimension(:), allocatable, private :: JETS_PS
  end type mlm_matching_data_t

\eatline
\nwendcode{}\nwbegindocs{360}\nwdocspar
\nwenddocs{}\nwbegincode{361}\moddef{MLM matching: public}\plusendmoddef
  public :: mlm_matching_settings_t
\nwendcode{}\nwbegincode{362}\moddef{MLM matching: types}\plusendmoddef
  type :: mlm_matching_settings_t
     real(default) :: mlm_Qcut_ME = one
     real(default) :: mlm_Qcut_PS = one
     real(default) :: mlm_ptmin, mlm_etamax, mlm_Rmin, mlm_Emin
     real(default) :: mlm_ETclusfactor = 0.2_default
     real(default) :: mlm_ETclusminE = five
     real(default) :: mlm_etaclusfactor = one
     real(default) :: mlm_Rclusfactor = one
     real(default) :: mlm_Eclusfactor = one
     integer :: kt_imode_hadronic = 4313
     integer :: kt_imode_leptonic = 1111
     integer :: mlm_nmaxMEjets = 0
  end type mlm_matching_settings_t

\eatline
\nwendcode{}\nwbegindocs{363}\nwdocspar
\nwenddocs{}\nwbegincode{364}\moddef{MLM matching: public}\plusendmoddef
  public :: mlm_matching_settings_write
\nwendcode{}\nwbegincode{365}\moddef{MLM matching: procedures}\endmoddef
  subroutine mlm_matching_settings_write (settings, unit)
    type(mlm_matching_settings_t), intent(in) :: settings
    integer, intent(in), optional :: unit
    integer :: u
    u = given_output_unit (unit);  if (u < 0)  return
    write (u, "(3x,A,ES19.12)") &
         "mlm_Qcut_ME                  = ", settings%mlm_Qcut_ME
    write (u, "(3x,A,ES19.12)") &
         "mlm_Qcut_PS                  = ", settings%mlm_Qcut_PS
    write (u, "(3x,A,ES19.12)") &
         "mlm_ptmin                    = ", settings%mlm_ptmin
    write (u, "(3x,A,ES19.12)") &
         "mlm_etamax                   = ", settings%mlm_etamax
    write (u, "(3x,A,ES19.12)") &
         "mlm_Rmin                     = ", settings%mlm_Rmin
    write (u, "(3x,A,ES19.12)") &
         "mlm_Emin                     = ", settings%mlm_Emin
    write (u, "(3x,A,1x,I0)") &
         "mlm_nmaxMEjets               = ", settings%mlm_nmaxMEjets
    write (u, "(3x,A,ES19.12)") &
         "mlm_ETclusfactor  (D=0.2)    = ", settings%mlm_ETclusfactor
    write (u, "(3x,A,ES19.12)") &
         "mlm_ETclusminE    (D=5.0)    = ", settings%mlm_ETclusminE
    write (u, "(3x,A,ES19.12)") &
         "mlm_etaclusfactor (D=1.0)    = ", settings%mlm_etaClusfactor
    write (u, "(3x,A,ES19.12)") &
         "mlm_Rclusfactor   (D=1.0)    = ", settings%mlm_RClusfactor
    write (u, "(3x,A,ES19.12)") &
         "mlm_Eclusfactor   (D=1.0)    = ", settings%mlm_EClusfactor
  end subroutine mlm_matching_settings_write

\eatline
\nwendcode{}\nwbegindocs{366}\nwdocspar
\nwenddocs{}\nwbegincode{367}\moddef{MLM matching: public}\plusendmoddef
  public :: mlm_matching_data_write
\nwendcode{}\nwbegincode{368}\moddef{MLM matching: procedures}\plusendmoddef
  subroutine mlm_matching_data_write (data, unit)
    type(mlm_matching_data_t), intent(in) :: data
    integer, intent(in), optional :: unit
    integer :: i
    integer :: u
    u = given_output_unit (unit);  if (u < 0)  return

    write (u, "(1x,A)") "MLM matching data:"
    write (u, "(3x,A)") "Momenta of ME partons:"
    if (allocated (data%P_ME)) then
       do i = 1, size (data%P_ME)
          write (u, "(4x)", advance = "no")
          call vector4_write (data%P_ME(i), unit = u)
       end do
    else
       write (u, "(5x,A)")  "[empty]"
    end if
    call write_separator (u)
    write (u, "(3x,A)")  "Momenta of ME jets:"
    if (allocated (data%JETS_ME)) then
       do i = 1, size (data%JETS_ME)
          write (u, "(4x)", advance = "no")
          call vector4_write (data%JETS_ME(i), unit = u)
       end do
    else
       write (u, "(5x,A)")  "[empty]"
    end if
    call write_separator (u)
    write(u, "(3x,A)")  "Momenta of shower partons:"
    if (allocated (data%P_PS)) then
       do i = 1, size (data%P_PS)
          write (u, "(4x)", advance = "no")
          call vector4_write (data%P_PS(i), unit = u)
       end do
    else
       write (u, "(5x,A)")  "[empty]"
    end if
    call write_separator (u)
    write (u, "(3x,A)")  "Momenta of shower jets:"
    if (allocated (data%JETS_PS)) then
       do i = 1, size (data%JETS_PS)
          write (u, "(4x)", advance = "no")
          call vector4_write (data%JETS_PS(i), unit = u)
       end do
    else
       write (u, "(5x,A)")  "[empty]"
    end if
    call write_separator (u)
  end subroutine mlm_matching_data_write

\eatline
\nwendcode{}\nwbegindocs{369}\nwdocspar
\nwenddocs{}\nwbegincode{370}\moddef{MLM matching: public}\plusendmoddef
  public :: mlm_matching_data_final
\nwendcode{}\nwbegincode{371}\moddef{MLM matching: procedures}\plusendmoddef
  subroutine mlm_matching_data_final (data)
    type(mlm_matching_data_t), intent(inout) :: data
    ! if (allocated (data%P_ME))  deallocate(data%P_ME)
    if (allocated (data%P_PS))  deallocate (data%P_PS)
    if (allocated (data%JETS_ME))  deallocate(data%JETS_ME)
    if (allocated (data%JETS_PS))  deallocate(data%JETS_PS)
  end subroutine mlm_matching_data_final

\eatline
\nwendcode{}\nwbegindocs{372}\nwdocspar
\nwenddocs{}\nwbegincode{373}\moddef{MLM matching: public}\plusendmoddef
  public :: mlm_matching_apply
\nwendcode{}\nwbegincode{374}\moddef{MLM matching: procedures}\plusendmoddef
  subroutine mlm_matching_apply (data, settings, vetoed)
    type(mlm_matching_data_t), intent(inout) :: data
    type(mlm_matching_settings_t), intent(in) :: settings
    logical, intent(out) :: vetoed

    integer :: i, j
    integer :: n_jets_ME, n_jets_PS, n_jets_PS_atycut
    real(double) :: ycut
    real(double), dimension(:, :), allocatable :: PP
    real(double), dimension(:), allocatable :: Y
    real(double), dimension(:,:), allocatable :: P_JETS
    real(double), dimension(:,:), allocatable :: P_ME
    integer, dimension(:), allocatable :: JET
    integer :: NJET, NSUB
    integer :: imode
!!! TODO: (bcn 2014-03-26) Why is ECUT hard coded to 1?
!!! It is the denominator of the KT measure. Candidate for removal
    real(double) :: ECUT = 1._double
    integer :: ip1,ip2

    ! KTCLUS COMMON BLOCK
    INTEGER NMAX,NUM,HIST
    PARAMETER (NMAX=512)
    DOUBLE PRECISION P,KT,KTP,KTS,ETOT,RSQ,KTLAST
    COMMON /KTCOMM/ETOT,RSQ,P(9,NMAX),KTP(NMAX,NMAX),KTS(NMAX), &
         KT(NMAX),KTLAST(NMAX),HIST(NMAX),NUM

    vetoed = .true.
    if (signal_is_pending ())  return

    \LA{}Set \code{}n{\char95}jets{\char95}ME/PS\edoc{} from \code{}data\edoc{} (or equal zero)\RA{}

    \LA{}Jet clustering for partons after matrix element\RA{}

    \LA{}Jet clustering for partons after shower\RA{}

    \LA{}Veto: too many or not enough jets after PS\RA{}

    \LA{}Cluster ME jets with PS jets one at a time\RA{}

    vetoed = .false.
999 continue
  end subroutine mlm_matching_apply

\eatline
\nwendcode{}\nwbegindocs{375}\nwdocspar
\nwenddocs{}\nwbegincode{376}\moddef{Set \code{}n{\char95}jets{\char95}ME/PS\edoc{} from \code{}data\edoc{} (or equal zero)}\endmoddef
if (allocated (data%P_ME)) then
   ! print *, "number of partons after ME: ", size(data%P_ME)
   n_jets_ME = size (data%P_ME)
else
   n_jets_ME = 0
end if
if (allocated (data%p_PS)) then
   ! print *, "number of partons after PS: ", size(data%p_PS)
   n_jets_PS = size (data%p_PS)
else
   n_jets_PS = 0
end if
\nwendcode{}\nwbegindocs{377}\nwdocspar
\nwenddocs{}\nwbegincode{378}\moddef{Jet clustering for partons after matrix element}\endmoddef
if (n_jets_ME > 0) then
   ycut = (settings%mlm_ptmin)**2
   allocate (PP(1:4, 1:N_jets_ME))
   do i = 1, n_jets_ME
      PP(1:3,i) = data%p_ME(i)%p(1:3)
      PP(4,i) = data%p_ME(i)%p(0)
   end do

   \LA{}Set \code{}imode\edoc{} for lepton or hadron collisions\RA{}

   allocate (P_ME(1:4,1:n_jets_ME))
   allocate (JET(1:n_jets_ME))
   allocate (Y(1:n_jets_ME))

   if (signal_is_pending ())  return
   call KTCLUR (imode, PP, n_jets_ME, &
        dble (settings%mlm_Rclusfactor * settings%mlm_Rmin), ECUT, y, *999)
   call KTRECO (1, PP, n_jets_ME, ECUT, ycut, ycut, P_ME, JET, &
        NJET, NSUB, *999)

   n_jets_ME = NJET
   if (NJET > 0) then
      allocate (data%JETS_ME (1:NJET))
      do i = 1, NJET
         data%JETS_ME(i) = vector4_moving (REAL(P_ME(4,i), default), &
              vector3_moving([REAL(P_ME(1,i), default), &
              REAL(P_ME(2,i), default), REAL(P_ME(3,i), default)]))
      end do
   end if
   deallocate (P_ME)
   deallocate (JET)
   deallocate (Y)
   deallocate (PP)
end if
\nwendcode{}\nwbegindocs{379}\nwdocspar
\nwenddocs{}\nwbegincode{380}\moddef{Jet clustering for partons after shower}\endmoddef
if (n_jets_PS > 0) then
   ycut = (settings%mlm_ptmin + max (settings%mlm_ETclusminE, &
        settings%mlm_ETclusfactor * settings%mlm_ptmin))**2
   allocate (PP(1:4, 1:n_jets_PS))
   do i = 1, n_jets_PS
      PP(1:3,i) = data%p_PS(i)%p(1:3)
      PP(4,i) = data%p_PS(i)%p(0)
   end do

   \LA{}Set \code{}imode\edoc{} for lepton or hadron collisions\RA{}

   allocate (P_JETS(1:4,1:n_jets_PS))
   allocate (JET(1:n_jets_PS))
   allocate (Y(1:n_jets_PS))

   if (signal_is_pending ()) return
   call KTCLUR (imode, PP, n_jets_PS, &
        dble (settings%mlm_Rclusfactor * settings%mlm_Rmin), &
        ECUT, y, *999)
   call KTRECO (1, PP, n_jets_PS, ECUT, ycut, ycut, P_JETS, JET, &
        NJET, NSUB, *999)
   n_jets_PS_atycut = NJET
   if (n_jets_ME == settings%mlm_nmaxMEjets .and. NJET > 0) then
      ! print *, " resetting ycut to ", Y(settings%mlm_nmaxMEjets)
      ycut = y(settings%mlm_nmaxMEjets)
      call KTRECO (1, PP, n_jets_PS, ECUT, ycut, ycut, P_JETS, JET, &
           NJET, NSUB, *999)
   end if

   ! !Sample of code for a FastJet interface
   ! palg = 1d0         ! 1.0d0 = kt, 0.0d0 = Cam/Aachen, -1.0d0 = anti-kt
   ! R = 0.7_double     ! radius parameter
   ! f = 0.75_double    ! overlap threshold
   ! !call fastjetppgenkt(PP,n,R,palg,P_JETS,NJET)   ! KT-Algorithm
   ! !call fastjetsiscone(PP,n,R,f,P_JETS,NJET)      ! SiSCone-Algorithm

   if (NJET > 0) then
      allocate (data%JETS_PS(1:NJET))
      do i = 1, NJET
         data%JETS_PS(i) = vector4_moving (REAL(P_JETS(4,i), default), &
              vector3_moving([REAL(P_JETS(1,i), default), &
              REAL(P_JETS(2,i), default), REAL(P_JETS(3,i), default)]))
      end do
   end if

   deallocate (P_JETS)
   deallocate (JET)
   deallocate (Y)
else
   n_jets_PS_atycut = 0
end if
\nwendcode{}\nwbegindocs{381}\nwdocspar
\nwenddocs{}\nwbegincode{382}\moddef{Set \code{}imode\edoc{} for lepton or hadron collisions}\endmoddef
if (data%is_hadron_collision) then
   imode = settings%kt_imode_hadronic
else
   imode = settings%kt_imode_leptonic
end if
\nwendcode{}\nwbegindocs{383}\nwdocspar
\nwenddocs{}\nwbegincode{384}\moddef{Veto: too many or not enough jets after PS}\endmoddef
if (n_jets_PS_atycut < n_jets_ME) then
   ! print *, "DISCARDING: Not enough PS jets: ", n_jets_PS_atycut
   return
end if
if (n_jets_PS_atycut > n_jets_ME .and. n_jets_ME /= settings%mlm_nmaxMEjets) then
   ! print *, "DISCARDING: Too many PS jets: ", n_jets_PS_atycut
   return
end if
\nwendcode{}\nwbegindocs{385}\nwdocspar
\nwenddocs{}\nwbegincode{386}\moddef{Cluster ME jets with PS jets one at a time}\endmoddef
if (allocated(data%JETS_PS)) then
   ! print *, "number of jets after PS: ", size(data%JETS_PS)
   n_jets_PS = size (data%JETS_PS)
else
   n_jets_PS = 0
end if
if (n_jets_ME > 0 .and. n_jets_PS > 0) then
   n_jets_PS = size (data%JETS_PS)
   if (allocated (PP))  deallocate(PP)
   allocate (PP(1:4, 1:n_jets_PS + 1))
   do i = 1, n_jets_PS
      if (signal_is_pending ()) return
      PP(1:3,i) = data%JETS_PS(i)%p(1:3)
      PP(4,i) = data%JETS_PS(i)%p(0)
   end do
   if (allocated (Y))  deallocate(Y)
   allocate (Y(1:n_jets_PS + 1))
   y = zero
   do i = 1, n_jets_ME
      PP(1:3,n_jets_PS + 2 - i) = data%JETS_ME(i)%p(1:3)
      PP(4,n_jets_PS + 2 - i) = data%JETS_ME(i)%p(0)
      !!! This makes more sense than hardcoding
      ! call KTCLUS (4313, PP, (n_jets_PS + 2 - i), 1.0_double, Y, *999)
      call KTCLUR (imode, PP, (n_jets_PS + 2 - i), &
        dble (settings%mlm_Rclusfactor * settings%mlm_Rmin), &
        ECUT, y, *999)
      if (0.99 * y(n_jets_PS + 1 - (i - 1)).gt.ycut) then
         ! print *, "DISCARDING: Jet ", i, " not clusterd"
         return
      end if
      !!! search for and remove PS jet clustered with ME Jet
      ip1 = HIST(n_jets_PS + 2 - i) / NMAX
      ip2 = mod(hist(n_jets_PS + 2 - i), NMAX)
      if ((ip2 /= n_jets_PS + 2 - i) .or. (ip1 <= 0)) then
         ! print *, "DISCARDING: Jet ", i, " not clustered ", ip1, ip2, &
         !      hist(n_jets_PS + 2 - i)
         return
      else
         ! print *, "PARTON clustered", ip1, ip2, hist(n_jets_PS + 2 - i)
         PP(:,IP1) = zero
         do j = IP1, n_jets_PS - i
            PP(:, j) = PP(:,j + 1)
         end do
      end if
   end do
end if
\nwendcode{}\nwbegindocs{387}\nwdocspar
\nwenddocs{}\nwbegincode{388}\moddef{MLM matching: procedures}\plusendmoddef
!C-----------------------------------------------------------------------
!C-----------------------------------------------------------------------
!C-----------------------------------------------------------------------
!C     KTCLUS: written by Mike Seymour, July 1992.
!C     Last modified November 2000.
!C     Please send comments or suggestions to Mike.Seymour@rl.ac.uk
!C
!C     This is a general-purpose kt clustering package.
!C     It can handle ee, ep and pp collisions.
!C     It is loosely based on the program of Siggi Bethke.
!C
!C     The time taken (on a 10MIP machine) is (0.2microsec)*N**3
!C     where N is the number of particles.
!C     Over 90 percent of this time is used in subroutine KTPMIN, which
!C     simply finds the minimum member of a one-dimensional array.
!C     It is well worth thinking about optimization: on the SPARCstation
!C     a factor of two increase was obtained simply by increasing the
!C     optimization level from its default value.
!C
!C     The approach is to separate the different stages of analysis.
!C     KTCLUS does all the clustering and records a merging history.
!C     It returns a simple list of the y values at which each merging
!C     occured. Then the following routines can be called to give extra
!C     information on the most recently analysed event.
!C     KTCLUR is identical but includes an R parameter, see below.
!C     KTYCUT gives the number of jets at each given YCUT value.
!C     KTYSUB gives the number of sub-jets at each given YCUT value.
!C     KTBEAM gives same info as KTCLUS but only for merges with the beam
!C     KTJOIN gives same info as KTCLUS but for merges of sub-jets.
!C     KTRECO reconstructs the jet momenta at a given value of YCUT.
!C     It also gives information on which jets at scale YCUT belong to
!C     which macro-jets at scale YMAC, for studying sub-jet properties.
!C     KTINCL reconstructs the jet momenta according to the inclusive jet
!C     definition of Ellis and Soper.
!C     KTISUB, KTIJOI and KTIREC are like KTYSUB, KTJOIN and KTRECO,
!C     except that they only apply to one inclusive jet at a time,
!C     with the pt of that jet automatically used for ECUT.
!C     KTWICH gives a list of which particles ended up in which jets.
!C     KTWCHS gives the same thing, but only for subjets.
!C     Note that the numbering of jets used by these two routines is
!C     guaranteed to be the same as that used by KTRECO.
!C
!C     The collision type and analysis type are indicated by the first
!C     argument of KTCLUS. IMODE=<TYPE><ANGLE><MONO><RECOM> where
!C     TYPE:  1=>ee, 2=>ep with p in -z direction, 3=>pe, 4=>pp
!C     ANGLE: 1=>angular kt def., 2=>DeltaR, 3=>f(DeltaEta,DeltaPhi)
!C            where f()=2(cosh(eta)-cos(phi)) is the QCD emission metric
!C     MONO:  1=>derive relative pseudoparticle angles from jets
!C            2=>monotonic definitions of relative angles
!C     RECOM: 1=>E recombination scheme, 2=>pt scheme, 3=>pt**2 scheme
!C
!C     There are also abbreviated forms for the most common combinations:
!C     IMODE=1 => E scheme in e+e-                              (=1111)
!C           2 => E scheme in ep                                (=2111)
!C           3 => E scheme in pe                                (=3111)
!C           4 => E scheme in pp                                (=4111)
!C           5 => covariant E scheme in pp                      (=4211)
!C           6 => covariant pt-scheme in pp                     (=4212)
!C           7 => covariant monotonic pt**2-scheme in pp        (=4223)
!C
!C     KTRECO no longer needs to reconstruct the momenta according to the
!C     same recombination scheme in which they were clustered. Its first
!C     argument gives the scheme, taking the same values as RECOM above.
!C
!C     Note that unlike previous versions, all variables which hold y
!C     values have been named in a consistent way:
!C     Y()  is the output scale at which jets were merged,
!C     YCUT is the input scale at which jets should be counted, and
!C          jet-momenta reconstructed etc,
!C     YMAC is the input macro-jet scale, used in determining whether
!C          or not each jet is a sub-jet.
!C     The original scheme defined in our papers is equivalent to always
!C     setting YMAC=1.
!C     Whenever a YCUT or YMAC variable is used, it is rounded down
!C     infinitesimally, so that for example, setting YCUT=Y(2) refers
!C     to the scale where the event is 2-jet, even if rounding errors
!C     have shifted its value slightly.
!C
!C     An R parameter can be used in hadron-hadron collisions by
!C     calling KTCLUR instead of KTCLUS.  This is as suggested by
!C     Ellis and Soper, but implemented slightly differently,
!C     as in M.H. Seymour, LU TP 94/2 (submitted to Nucl. Phys. B.).
!C     R**2 multiplies the single Kt everywhere it is used.
!C     Calling KTCLUR with R=1 is identical to calling KTCLUS.
!C     R plays a similar role to the jet radius in a cone-type algorithm,
!C     but is scaled up by about 40% (ie R=0.7 in a cone algorithm is
!C     similar to this algorithm with R=1).
!C     Note that R.EQ.1 must be used for the e+e- and ep versions,
!C     and is strongly recommended for the hadron-hadron version.
!C     However, R values smaller than 1 have been found to be useful for
!C     certain applications, particularly the mass reconstruction of
!C     highly-boosted colour-singlets such as high-pt hadronic Ws,
!C     as in M.H. Seymour, LU TP 93/8 (to appear in Z. Phys. C.).
!C     Situations in which R<1 is useful are likely to also be those in
!C     which the inclusive reconstruction method is more useful.
!C
!C     Also included is a set of routines for doing Lorentz boosts:
!C     KTLBST finds the boost matrix to/from the cm frame of a 4-vector
!C     KTRROT finds the rotation matrix from one vector to another
!C     KTMMUL multiplies together two matrices
!C     KTVMUL multiplies a vector by a matrix
!C     KTINVT inverts a transformation matrix (nb NOT a general 4 by 4)
!C     KTFRAM boosts a list of vectors between two arbitrary frames
!C     KTBREI boosts a list of vectors between the lab and Breit frames
!C     KTHADR boosts a list of vectors between the lab and hadronic cmf
!C       The last two need the momenta in the +z direction of the lepton
!C       and hadron beams, and the 4-momentum of the outgoing lepton.
!C
!C     The main reference is:
!C       S. Catani, Yu.L. Dokshitzer, M.H. Seymour and B.R. Webber,
!C         Nucl.Phys.B406(1993)187.
!C     The ep version was proposed in:
!C       S. Catani, Yu.L. Dokshitzer and B.R. Webber,
!C         Phys.Lett.285B(1992)291.
!C     The inclusive reconstruction method was proposed in:
!C       S.D. Ellis and D.E. Soper,
!C         Phys.Rev.D48(1993)3160.
!C
!C-----------------------------------------------------------------------
!C-----------------------------------------------------------------------
!C-----------------------------------------------------------------------
\nwendcode{}\nwbegincode{389}\moddef{MLM matching: procedures}\plusendmoddef
  SUBROUTINE KTCLUR(IMODE,PP,NN,R,ECUT,Y,*)
    IMPLICIT NONE
!C---DO CLUSTER ANALYSIS OF PARTICLES IN PP
!C
!C   IMODE   = INPUT  : DESCRIBED ABOVE
!C   PP(I,J) = INPUT  : 4-MOMENTUM OF Jth PARTICLE: I=1,4 => PX,PY,PZ,E
!C   NN      = INPUT  : NUMBER OF PARTICLES
!C   R       = INPUT  : ELLIS AND SOPER'S R PARAMETER, SEE ABOVE.
!C   ECUT    = INPUT  : DENOMINATOR OF KT MEASURE. IF ZERO, ETOT IS USED
!C   Y(J)    = OUTPUT : VALUE OF Y FOR WHICH EVENT CHANGES FROM BEING
!C                        J JET TO J-1 JET
!C   LAST ARGUMENT IS LABEL TO JUMP TO IF FOR ANY REASON THE EVENT
!C   COULD NOT BE PROCESSED (MOST LIKELY DUE TO TOO MANY PARTICLES)
!C
!C   NOTE THAT THE MOMENTA ARE DECLARED DOUBLE PRECISION,
!C   AND ALL OTHER FLOATING POINT VARIABLES ARE DECLARED DOUBLE PRECISION
!C
    INTEGER NMAX,IM,IMODE,TYPE,ANGL,MONO,RECO,N,I,J,NN, &
         IMIN,JMIN,KMIN,NUM,HIST,INJET,IABBR,NABBR
    PARAMETER (NMAX=512,NABBR=7)
    DOUBLE PRECISION PP(4,*)
    integer :: u
!CHANGE    DOUBLE PRECISION R,ECUT,Y(*),P,KT,ETOT,RSQ,KTP,KTS,KTPAIR,KTSING, &
    DOUBLE PRECISION R,ECUT,Y(*),P,KT,ETOT,RSQ,KTP,KTS, &
         KTMIN,ETSQ,KTLAST,KTMAX,KTTMP
    LOGICAL FIRST
    CHARACTER TITLE(4,4)*10
!C---KT RECORDS THE KT**2 OF EACH MERGING.
!C---KTLAST RECORDS FOR EACH MERGING, THE HIGHEST ECUT**2 FOR WHICH THE
!C   RESULT IS NOT MERGED WITH THE BEAM (COULD BE LARGER THAN THE
!C   KT**2 AT WHICH IT WAS MERGED IF THE KT VALUES ARE NOT MONOTONIC).
!C   THIS MAY SOUND POINTLESS, BUT ITS USEFUL FOR DETERMINING WHETHER
!C   SUB-JETS SURVIVED TO SCALE Y=YMAC OR NOT.
!C---HIST RECORDS MERGING HISTORY:
!C   N=>DELETED TRACK N, M*NMAX+N=>MERGED TRACKS M AND N (M<N).
    COMMON /KTCOMM/ETOT,RSQ,P(9,NMAX),KTP(NMAX,NMAX),KTS(NMAX), &
         KT(NMAX),KTLAST(NMAX),HIST(NMAX),NUM
    DIMENSION INJET(NMAX),IABBR(NABBR)
    DATA FIRST,TITLE,IABBR/.TRUE., &
         'e+e-      ','ep        ','pe        ','pp        ', &
         'angle     ','DeltaR    ','f(DeltaR) ','**********', &
         'no        ','yes       ','**********','**********', &
         'E         ','Pt        ','Pt**2     ','**********', &
         1111,2111,3111,4111,4211,4212,4223/
!C---CHECK INPUT
    IM=IMODE
    IF (IM.GE.1.AND.IM.LE.NABBR) IM=IABBR(IM)
    TYPE=MOD(IM/1000,10)
    ANGL=MOD(IM/100 ,10)
    MONO=MOD(IM/10  ,10)
    RECO=MOD(IM     ,10)
    IF (NN.GT.NMAX) CALL KTWARN('KT-MAX',100,*999)
    IF (NN.LT.1) CALL KTWARN('KT-LT1',100,*999)
    IF (NN.LT.2.AND.TYPE.EQ.1) CALL KTWARN('KT-LT2',100,*999)
    IF (TYPE.LT.1.OR.TYPE.GT.4.OR.ANGL.LT.1.OR.ANGL.GT.4.OR. &
         MONO.LT.1.OR.MONO.GT.2.OR.RECO.LT.1.OR.RECO.GT.3) CALL KTWARN('KTCLUS',101,*999)
    u = given_output_unit ()
    IF (FIRST) THEN
       WRITE (u,'(/,1X,54(''*'')/A)') &
            ' KTCLUS: written by Mike Seymour, July 1992.'
       WRITE (u,'(A)') &
            ' Last modified November 2000.'
       WRITE (u,'(A)') &
            ' Please send comments or suggestions to Mike.Seymour@rl.ac.uk'
       WRITE (u,'(/A,I2,2A)') &
            '       Collision type =',TYPE,' = ',TITLE(TYPE,1)
       WRITE (u,'(A,I2,2A)') &
            '     Angular variable =',ANGL,' = ',TITLE(ANGL,2)
       WRITE (u,'(A,I2,2A)') &
            ' Monotonic definition =',MONO,' = ',TITLE(MONO,3)
       WRITE (u,'(A,I2,2A)') &
            ' Recombination scheme =',RECO,' = ',TITLE(RECO,4)
       IF (R.NE.1) THEN
          WRITE (u,'(A,F5.2)') &
               '     Radius parameter =',R
          IF (TYPE.NE.4) WRITE (u,'(A)') &
               ' R.NE.1 is strongly discouraged for this collision type!'
       ENDIF
       WRITE (u,'(1X,54(''*'')/)')
       FIRST=.FALSE.
    ENDIF
!C---COPY PP TO P
    N=NN
    NUM=NN
    CALL KTCOPY(PP,N,P,(RECO.NE.1))
    ETOT=0
    DO I=1,N
       ETOT=ETOT+P(4,I)
    END DO
    IF (ETOT.EQ.0) CALL KTWARN('KTCLUS',102,*999)
    IF (ECUT.EQ.0) THEN
       ETSQ=1/ETOT**2
    ELSE
       ETSQ=1/ECUT**2
    ENDIF
    RSQ=R**2
!C---CALCULATE ALL PAIR KT's
    DO I=1,N-1
       DO J=I+1,N
          KTP(J,I)=-1
          KTP(I,J)=KTPAIR(ANGL,P(1,I),P(1,J),KTP(J,I))
       END DO
    END DO
!C---CALCULATE ALL SINGLE KT's
    DO I=1,N
       KTS(I)=KTSING(ANGL,TYPE,P(1,I))
    END DO
    KTMAX=0
!C---MAIN LOOP
300 CONTINUE
!C---FIND MINIMUM MEMBER OF KTP
    CALL KTPMIN(KTP,NMAX,N,IMIN,JMIN)
!C---FIND MINIMUM MEMBER OF KTS
    CALL KTSMIN(KTS,NMAX,N,KMIN)
!C---STORE Y VALUE OF TRANSITION FROM N TO N-1 JETS
    KTMIN=KTP(IMIN,JMIN)
    KTTMP=RSQ*KTS(KMIN)
    IF ((TYPE.GE.2.AND.TYPE.LE.4).AND. &
         (KTTMP.LE.KTMIN.OR.N.EQ.1)) &
         KTMIN=KTTMP
    KT(N)=KTMIN
    Y(N)=KT(N)*ETSQ
!C---IF MONO.GT.1, SEQUENCE IS SUPPOSED TO BE MONOTONIC, IF NOT, WARN
    IF (KTMIN.LT.KTMAX.AND.MONO.GT.1) CALL KTWARN('KTCLUS',1,*999)
    IF (KTMIN.GE.KTMAX) KTMAX=KTMIN
!C---IF LOWEST KT IS TO A BEAM, THROW IT AWAY AND MOVE LAST ENTRY UP
    IF (KTMIN.EQ.KTTMP) THEN
       CALL KTMOVE(P,KTP,KTS,NMAX,N,KMIN,1)
!C---UPDATE HISTORY AND CROSS-REFERENCES
       HIST(N)=KMIN
       INJET(N)=KMIN
       DO I=N,NN
          IF (INJET(I).EQ.KMIN) THEN
             KTLAST(I)=KTMAX
             INJET(I)=0
          ELSEIF (INJET(I).EQ.N) THEN
             INJET(I)=KMIN
          ENDIF
       END DO
!C---OTHERWISE MERGE JETS IMIN AND JMIN AND MOVE LAST ENTRY UP
    ELSE
       CALL KTMERG(P,KTP,KTS,NMAX,IMIN,JMIN,N,TYPE,ANGL,MONO,RECO)
       CALL KTMOVE(P,KTP,KTS,NMAX,N,JMIN,1)
!C---UPDATE HISTORY AND CROSS-REFERENCES
       HIST(N)=IMIN*NMAX+JMIN
       INJET(N)=IMIN
       DO I=N,NN
          IF (INJET(I).EQ.JMIN) THEN
             INJET(I)=IMIN
          ELSEIF (INJET(I).EQ.N) THEN
             INJET(I)=JMIN
          ENDIF
       END DO
    ENDIF
!C---THATS ALL THERE IS TO IT
    N=N-1
    IF (N.GT.1 .OR. N.GT.0.AND.(TYPE.GE.2.AND.TYPE.LE.4)) GOTO 300
    IF (N.EQ.1) THEN
       KT(N)=1D20
       Y(N)=KT(N)*ETSQ
    ENDIF
    RETURN
999 RETURN 1
  END SUBROUTINE KTCLUR
!C-----------------------------------------------------------------------
\nwendcode{}\nwbegincode{390}\moddef{MLM matching: procedures}\plusendmoddef
!C-----------------------------------------------------------------------
    SUBROUTINE KTRECO(RECO,PP,NN,ECUT,YCUT,YMAC,PJET,JET,NJET,NSUB,*)
      IMPLICIT NONE
!C---RECONSTRUCT KINEMATICS OF JET SYSTEM, WHICH HAS ALREADY BEEN
!C   ANALYSED BY KTCLUS. NOTE THAT NO CONSISTENCY CHECK IS MADE: USER
!C   IS TRUSTED TO USE THE SAME PP VALUES AS FOR KTCLUS
!C
!C   RECO     = INPUT : RECOMBINATION SCHEME (NEED NOT BE SAME AS KTCLUS)
!C   PP(I,J)  = INPUT : 4-MOMENTUM OF Jth PARTICLE: I=1,4 => PX,PY,PZ,E
!C   NN       = INPUT : NUMBER OF PARTICLES
!C   ECUT     = INPUT : DENOMINATOR OF KT MEASURE. IF ZERO, ETOT IS USED
!C   YCUT     = INPUT : Y VALUE AT WHICH TO RECONSTRUCT JET MOMENTA
!C   YMAC     = INPUT : Y VALUE USED TO DEFINE MACRO-JETS, TO DETERMINE
!C                        WHICH JETS ARE SUB-JETS
!C   PJET(I,J)=OUTPUT : 4-MOMENTUM OF Jth JET AT SCALE YCUT
!C   JET(J)   =OUTPUT : THE MACRO-JET WHICH CONTAINS THE Jth JET,
!C                        SET TO ZERO IF JET IS NOT A SUB-JET
!C   NJET     =OUTPUT : THE NUMBER OF JETS
!C   NSUB     =OUTPUT : THE NUMBER OF SUB-JETS (EQUAL TO THE NUMBER OF
!C                        NON-ZERO ENTRIES IN JET())
!C   LAST ARGUMENT IS LABEL TO JUMP TO IF FOR ANY REASON THE EVENT
!C   COULD NOT BE PROCESSED
!C
!C   NOTE THAT THE MOMENTA ARE DECLARED DOUBLE PRECISION,
!C   AND ALL OTHER FLOATING POINT VARIABLES ARE DECLARED DOUBLE PRECISION
!C
      INTEGER NMAX,RECO,NUM,N,NN,NJET,NSUB,JET(*),HIST,IMIN,JMIN,I,J
      PARAMETER (NMAX=512)
      DOUBLE PRECISION PP(4,*),PJET(4,*)
      DOUBLE PRECISION ECUT,P,KT,KTP,KTS,ETOT,RSQ,ETSQ,YCUT,YMAC,KTLAST, &
           ROUND
      PARAMETER (ROUND=0.99999D0)
      COMMON /KTCOMM/ETOT,RSQ,P(9,NMAX),KTP(NMAX,NMAX),KTS(NMAX), &
           KT(NMAX),KTLAST(NMAX),HIST(NMAX),NUM
!C---CHECK INPUT
      if (signal_is_pending ()) return
      IF (RECO.LT.1.OR.RECO.GT.3) THEN
         PRINT *,'RECO=',RECO
         CALL KTWARN('KTRECO',100,*999)
      ENDIF
!C---COPY PP TO P
      N=NN
      IF (NUM.NE.NN) CALL KTWARN('KTRECO',101,*999)
      CALL KTCOPY(PP,N,P,(RECO.NE.1))
      IF (ECUT.EQ.0) THEN
         ETSQ=1/ETOT**2
      ELSE
         ETSQ=1/ECUT**2
      ENDIF
!C---KEEP MERGING UNTIL YCUT
 100  IF (ETSQ*KT(N).LT.ROUND*YCUT) THEN
         IF (HIST(N).LE.NMAX) THEN
            CALL KTMOVE(P,KTP,KTS,NMAX,N,HIST(N),0)
         ELSE
            IMIN=HIST(N)/NMAX
            JMIN=HIST(N)-IMIN*NMAX
            CALL KTMERG(P,KTP,KTS,NMAX,IMIN,JMIN,N,0,0,0,RECO)
            CALL KTMOVE(P,KTP,KTS,NMAX,N,JMIN,0)
         ENDIF
         N=N-1
         IF (N.GT.0) GOTO 100
      ENDIF
!C---IF YCUT IS TOO LARGE THERE ARE NO JETS
      NJET=N
      NSUB=N
      IF (N.EQ.0) RETURN
!C---SET UP OUTPUT MOMENTA
      DO I=1,NJET
         IF (RECO.EQ.1) THEN
            DO J=1,4
               PJET(J,I)=P(J,I)
            END DO
         ELSE
            PJET(1,I)=P(6,I)*COS(P(8,I))
            PJET(2,I)=P(6,I)*SIN(P(8,I))
            PJET(3,I)=P(6,I)*SINH(P(7,I))
            PJET(4,I)=P(6,I)*COSH(P(7,I))
         ENDIF
         JET(I)=I
      END DO
!C---KEEP MERGING UNTIL YMAC TO FIND THE FATE OF EACH JET
300   IF (ETSQ*KT(N).LT.ROUND*YMAC) THEN
         IF (HIST(N).LE.NMAX) THEN
            IMIN=0
            JMIN=HIST(N)
            NSUB=NSUB-1
         ELSE
            IMIN=HIST(N)/NMAX
            JMIN=HIST(N)-IMIN*NMAX
            IF (ETSQ*KTLAST(N).LT.ROUND*YMAC) NSUB=NSUB-1
         ENDIF
         DO I=1,NJET
            IF (JET(I).EQ.JMIN) JET(I)=IMIN
            IF (JET(I).EQ.N) JET(I)=JMIN
         END DO
         N=N-1
         IF (N.GT.0) GOTO 300
      ENDIF
      RETURN
999   RETURN 1
    END SUBROUTINE KTRECO
!C-----------------------------------------------------------------------
\nwendcode{}\nwbegincode{391}\moddef{MLM matching: procedures}\plusendmoddef
!C-----------------------------------------------------------------------
    FUNCTION KTPAIR(ANGL,P,Q,ANGLE)
      IMPLICIT NONE
!C---CALCULATE LOCAL KT OF PAIR, USING ANGULAR SCHEME:
!C   1=>ANGULAR, 2=>DeltaR, 3=>f(DeltaEta,DeltaPhi)
!C   WHERE f(eta,phi)=2(COSH(eta)-COS(phi)) IS THE QCD EMISSION METRIC
!C---IF ANGLE<0, IT IS SET TO THE ANGULAR PART OF THE LOCAL KT ON RETURN
!C   IF ANGLE>0, IT IS USED INSTEAD OF THE ANGULAR PART OF THE LOCAL KT
      INTEGER ANGL
! CHANGE      DOUBLE PRECISION P(9),Q(9),KTPAIR,R,KTMDPI,ANGLE,ETA,PHI,ESQ
      DOUBLE PRECISION P(9),Q(9),KTPAIR,R,ANGLE,ETA,PHI,ESQ
!C---COMPONENTS OF MOMENTA ARE PX,PY,PZ,E,1/P,PT,ETA,PHI,PT**2
      R=ANGLE
      IF (ANGL.EQ.1) THEN
         IF (R.LE.0) R=2*(1-(P(1)*Q(1)+P(2)*Q(2)+P(3)*Q(3))*(P(5)*Q(5)))
         ESQ=MIN(P(4),Q(4))**2
      ELSEIF (ANGL.EQ.2.OR.ANGL.EQ.3) THEN
         IF (R.LE.0) THEN
            ETA=P(7)-Q(7)
            PHI=KTMDPI(P(8)-Q(8))
            IF (ANGL.EQ.2) THEN
               R=ETA**2+PHI**2
            ELSE
               R=2*(COSH(ETA)-COS(PHI))
            ENDIF
         ENDIF
         ESQ=MIN(P(9),Q(9))
      ELSEIF (ANGL.EQ.4) THEN
         ESQ=(1d0/(P(5)*Q(5))-P(1)*Q(1)-P(2)*Q(2)- &
              P(3)*Q(3))*2D0/(P(5)*Q(5))/(0.0001D0+1d0/P(5)+1d0/Q(5))**2
         R=1d0
      ELSE
         CALL KTWARN('KTPAIR',200,*999)
         STOP
      ENDIF
      KTPAIR=ESQ*R
      IF (ANGLE.LT.0) ANGLE=R
999 END FUNCTION KTPAIR
!C-----------------------------------------------------------------------
    FUNCTION KTSING(ANGL,TYPE,P)
      IMPLICIT NONE
!C---CALCULATE KT OF PARTICLE, USING ANGULAR SCHEME:
!C   1=>ANGULAR, 2=>DeltaR, 3=>f(DeltaEta,DeltaPhi)
!C---TYPE=1 FOR E+E-, 2 FOR EP, 3 FOR PE, 4 FOR PP
!C   FOR EP, PROTON DIRECTION IS DEFINED AS -Z
!C   FOR PE, PROTON DIRECTION IS DEFINED AS +Z
      INTEGER ANGL,TYPE
      DOUBLE PRECISION P(9),KTSING,COSTH,R,SMALL
      DATA SMALL/1D-4/
      IF (ANGL.EQ.1.OR.ANGL.EQ.4) THEN
         COSTH=P(3)*P(5)
         IF (TYPE.EQ.2) THEN
            COSTH=-COSTH
         ELSEIF (TYPE.EQ.4) THEN
            COSTH=ABS(COSTH)
         ELSEIF (TYPE.NE.1.AND.TYPE.NE.3) THEN
            CALL KTWARN('KTSING',200,*999)
            STOP
         ENDIF
         R=2*(1-COSTH)
!C---IF CLOSE TO BEAM, USE APPROX 2*(1-COS(THETA))=SIN**2(THETA)
         IF (R.LT.SMALL) R=(P(1)**2+P(2)**2)*P(5)**2
         KTSING=P(4)**2*R
      ELSEIF (ANGL.EQ.2.OR.ANGL.EQ.3) THEN
         KTSING=P(9)
      ELSE
         CALL KTWARN('KTSING',201,*999)
         STOP
      ENDIF
999 END FUNCTION KTSING
!C-----------------------------------------------------------------------
    SUBROUTINE KTPMIN(A,NMAX,N,IMIN,JMIN)
      IMPLICIT NONE
!C---FIND THE MINIMUM MEMBER OF A(NMAX,NMAX) WITH IMIN < JMIN <= N
      INTEGER NMAX,N,IMIN,JMIN,KMIN,I,J,K
!C---REMEMBER THAT A(X+(Y-1)*NMAX)=A(X,Y)
!C   THESE LOOPING VARIABLES ARE J=Y-2, I=X+(Y-1)*NMAX
      DOUBLE PRECISION A(*),AMIN
      K=1+NMAX
      KMIN=K
      AMIN=A(KMIN)
      DO J=0,N-2
         DO I=K,K+J
            IF (A(I).LT.AMIN) THEN
               KMIN=I
               AMIN=A(KMIN)
            ENDIF
         END DO
         K=K+NMAX
      END DO
      JMIN=KMIN/NMAX+1
      IMIN=KMIN-(JMIN-1)*NMAX
    END SUBROUTINE KTPMIN
!C-----------------------------------------------------------------------
    SUBROUTINE KTSMIN(A,NMAX,N,IMIN)
      IMPLICIT NONE
!C---FIND THE MINIMUM MEMBER OF A
      INTEGER N,NMAX,IMIN,I
      DOUBLE PRECISION A(NMAX)
      IMIN=1
      DO I=1,N
         IF (A(I).LT.A(IMIN)) IMIN=I
      END DO
    END SUBROUTINE KTSMIN
!C-----------------------------------------------------------------------
    SUBROUTINE KTCOPY(A,N,B,ONSHLL)
      IMPLICIT NONE
!C---COPY FROM A TO B. 5TH=1/(3-MTM), 6TH=PT, 7TH=ETA, 8TH=PHI, 9TH=PT**2
!C   IF ONSHLL IS .TRUE. PARTICLE ENTRIES ARE PUT ON-SHELL BY SETTING E=P
      INTEGER I,N
      DOUBLE PRECISION A(4,N)
      LOGICAL ONSHLL
      DOUBLE PRECISION B(9,N),ETAMAX,SINMIN,EPS
      DATA ETAMAX,SINMIN,EPS/10,0,1D-6/
!C---SINMIN GETS CALCULATED ON FIRST CALL
      IF (SINMIN.EQ.0) SINMIN=1/COSH(ETAMAX)
      DO I=1,N
         B(1,I)=A(1,I)
         B(2,I)=A(2,I)
         B(3,I)=A(3,I)
         B(4,I)=A(4,I)
         B(5,I)=SQRT(A(1,I)**2+A(2,I)**2+A(3,I)**2)
         IF (ONSHLL) B(4,I)=B(5,I)
         IF (B(5,I).EQ.0) B(5,I)=1D-10
         B(5,I)=1/B(5,I)
         B(9,I)=A(1,I)**2+A(2,I)**2
         B(6,I)=SQRT(B(9,I))
         B(7,I)=B(6,I)*B(5,I)
         IF (B(7,I).GT.SINMIN) THEN
            B(7,I)=A(4,I)**2-A(3,I)**2
            IF (B(7,I).LE.EPS*B(4,I)**2.OR.ONSHLL) B(7,I)=B(9,I)
            B(7,I)=LOG((B(4,I)+ABS(B(3,I)))**2/B(7,I))/2
         ELSE
            B(7,I)=ETAMAX+2
         ENDIF
         B(7,I)=SIGN(B(7,I),B(3,I))
         IF (A(1,I).EQ.0 .AND. A(2,I).EQ.0) THEN
            B(8,I)=0
         ELSE
            B(8,I)=ATAN2(A(2,I),A(1,I))
         ENDIF
      END DO
    END SUBROUTINE KTCOPY
!C-----------------------------------------------------------------------
    SUBROUTINE KTMERG(P,KTP,KTS,NMAX,I,J,N,TYPE,ANGL,MONO,RECO)
      IMPLICIT NONE
!C---MERGE THE Jth PARTICLE IN P INTO THE Ith PARTICLE
!C   J IS ASSUMED GREATER THAN I. P CONTAINS N PARTICLES BEFORE MERGING.
!C---ALSO RECALCULATING THE CORRESPONDING KTP AND KTS VALUES IF MONO.GT.0
!C   FROM THE RECOMBINED ANGULAR MEASURES IF MONO.GT.1
!C---NOTE THAT IF MONO.LE.0, TYPE AND ANGL ARE NOT USED
      INTEGER ANGL,RECO,TYPE,I,J,K,N,NMAX,MONO
      DOUBLE PRECISION P(9,NMAX),KTP(NMAX,NMAX),KTS(NMAX),PT,PTT, &
! CHANGE           KTMDPI,KTUP,PI,PJ,ANG,KTPAIR,KTSING,ETAMAX,EPS
           KTUP,PI,PJ,ANG,ETAMAX,EPS
      KTUP(I,J)=KTP(MAX(I,J),MIN(I,J))
      DATA ETAMAX,EPS/10,1D-6/
      IF (J.LE.I) CALL KTWARN('KTMERG',200,*999)
!C---COMBINE ANGULAR MEASURES IF NECESSARY
      IF (MONO.GT.1) THEN
         DO K=1,N
            IF (K.NE.I.AND.K.NE.J) THEN
               IF (RECO.EQ.1) THEN
                  PI=P(4,I)
                  PJ=P(4,J)
               ELSEIF (RECO.EQ.2) THEN
                  PI=P(6,I)
                  PJ=P(6,J)
               ELSEIF (RECO.EQ.3) THEN
                  PI=P(9,I)
                  PJ=P(9,J)
               ELSE
                  CALL KTWARN('KTMERG',201,*999)
                  STOP
               ENDIF
               IF (PI.EQ.0.AND.PJ.EQ.0) THEN
                  PI=1
                  PJ=1
               ENDIF
               KTP(MAX(I,K),MIN(I,K))= &
                    (PI*KTUP(I,K)+PJ*KTUP(J,K))/(PI+PJ)
            ENDIF
         END DO
      ENDIF
      IF (RECO.EQ.1) THEN
!C---VECTOR ADDITION
         P(1,I)=P(1,I)+P(1,J)
         P(2,I)=P(2,I)+P(2,J)
         P(3,I)=P(3,I)+P(3,J)
!c         P(4,I)=P(4,I)+P(4,J) ! JA
         P(5,I)=SQRT(P(1,I)**2+P(2,I)**2+P(3,I)**2)
         P(4,I)=P(5,I) ! JA (Massless scheme)
         IF (P(5,I).EQ.0) THEN
            P(5,I)=1
         ELSE
            P(5,I)=1/P(5,I)
         ENDIF
      ELSEIF (RECO.EQ.2) THEN
!C---PT WEIGHTED ETA-PHI ADDITION
         PT=P(6,I)+P(6,J)
         IF (PT.EQ.0) THEN
            PTT=1
         ELSE
            PTT=1/PT
         ENDIF
         P(7,I)=(P(6,I)*P(7,I)+P(6,J)*P(7,J))*PTT
         P(8,I)=KTMDPI(P(8,I)+P(6,J)*PTT*KTMDPI(P(8,J)-P(8,I)))
         P(6,I)=PT
         P(9,I)=PT**2
      ELSEIF (RECO.EQ.3) THEN
!C---PT**2 WEIGHTED ETA-PHI ADDITION
         PT=P(9,I)+P(9,J)
         IF (PT.EQ.0) THEN
            PTT=1
         ELSE
            PTT=1/PT
         ENDIF
         P(7,I)=(P(9,I)*P(7,I)+P(9,J)*P(7,J))*PTT
         P(8,I)=KTMDPI(P(8,I)+P(9,J)*PTT*KTMDPI(P(8,J)-P(8,I)))
         P(6,I)=P(6,I)+P(6,J)
         P(9,I)=P(6,I)**2
      ELSE
         CALL KTWARN('KTMERG',202,*999)
         STOP
      ENDIF
!C---IF MONO.GT.0 CALCULATE NEW KT MEASURES. IF MONO.GT.1 USE ANGULAR ONES.
      IF (MONO.LE.0) RETURN
!C---CONVERTING BETWEEN 4-MTM AND PT,ETA,PHI IF NECESSARY
      IF (ANGL.NE.1.AND.RECO.EQ.1) THEN
         P(9,I)=P(1,I)**2+P(2,I)**2
         P(7,I)=P(4,I)**2-P(3,I)**2
         IF (P(7,I).LE.EPS*P(4,I)**2) P(7,I)=P(9,I)
         IF (P(7,I).GT.0) THEN
            P(7,I)=LOG((P(4,I)+ABS(P(3,I)))**2/P(7,I))/2
            IF (P(7,I).GT.ETAMAX) P(7,I)=ETAMAX+2
         ELSE
            P(7,I)=ETAMAX+2
         ENDIF
         P(7,I)=SIGN(P(7,I),P(3,I))
         IF (P(1,I).NE.0.AND.P(2,I).NE.0) THEN
            P(8,I)=ATAN2(P(2,I),P(1,I))
         ELSE
            P(8,I)=0
         ENDIF
      ELSEIF (ANGL.EQ.1.AND.RECO.NE.1) THEN
         P(1,I)=P(6,I)*COS(P(8,I))
         P(2,I)=P(6,I)*SIN(P(8,I))
         P(3,I)=P(6,I)*SINH(P(7,I))
         P(4,I)=P(6,I)*COSH(P(7,I))
         IF (P(4,I).NE.0) THEN
            P(5,I)=1/P(4,I)
         ELSE
            P(5,I)=1
         ENDIF
      ENDIF
      ANG=0
      DO K=1,N
         IF (K.NE.I.AND.K.NE.J) THEN
            IF (MONO.GT.1) ANG=KTUP(I,K)
            KTP(MIN(I,K),MAX(I,K))= &
                 KTPAIR(ANGL,P(1,I),P(1,K),ANG)
         ENDIF
      END DO
      KTS(I)=KTSING(ANGL,TYPE,P(1,I))
999 END SUBROUTINE KTMERG
!C-----------------------------------------------------------------------
    SUBROUTINE KTMOVE(P,KTP,KTS,NMAX,N,J,IOPT)
      IMPLICIT NONE
!C---MOVE THE Nth PARTICLE IN P TO THE Jth POSITION
!C---ALSO MOVING KTP AND KTS IF IOPT.GT.0
      INTEGER I,J,N,NMAX,IOPT
      DOUBLE PRECISION P(9,NMAX),KTP(NMAX,NMAX),KTS(NMAX)
      DO I=1,9
         P(I,J)=P(I,N)
      END DO
      IF (IOPT.LE.0) RETURN
      DO I=1,J-1
         KTP(I,J)=KTP(I,N)
         KTP(J,I)=KTP(N,I)
      END DO
      DO I=J+1,N-1
         KTP(J,I)=KTP(I,N)
         KTP(I,J)=KTP(N,I)
      END DO
      KTS(J)=KTS(N)
    END SUBROUTINE KTMOVE
!C-----------------------------------------------------------------------
\nwendcode{}\nwbegincode{392}\moddef{MLM matching: procedures}\plusendmoddef
    FUNCTION KTMDPI(PHI)
      IMPLICIT NONE
!C---RETURNS PHI, MOVED ONTO THE RANGE [-PI,PI)
      DOUBLE PRECISION KTMDPI,PHI,PI,TWOPI,THRPI,EPS
      PARAMETER (PI=3.14159265358979324D0,TWOPI=6.28318530717958648D0, &
           THRPI=9.42477796076937972D0)
      PARAMETER (EPS=1D-15)
      KTMDPI=PHI
      IF (KTMDPI.LE.PI) THEN
        IF (KTMDPI.GT.-PI) THEN
          GOTO 100
        ELSEIF (KTMDPI.GT.-THRPI) THEN
          KTMDPI=KTMDPI+TWOPI
        ELSE
          KTMDPI=-MOD(PI-KTMDPI,TWOPI)+PI
        ENDIF
      ELSEIF (KTMDPI.LE.THRPI) THEN
        KTMDPI=KTMDPI-TWOPI
      ELSE
        KTMDPI=MOD(PI+KTMDPI,TWOPI)-PI
      ENDIF
 100  IF (ABS(KTMDPI).LT.EPS) KTMDPI=0
    END FUNCTION KTMDPI
!C-----------------------------------------------------------------------
    SUBROUTINE KTWARN(SUBRTN,ICODE,*)
!C     DEALS WITH ERRORS DURING EXECUTION
!C     SUBRTN = NAME OF CALLING SUBROUTINE
!C     ICODE  = ERROR CODE:    - 99 PRINT WARNING & CONTINUE
!C                          100-199 PRINT WARNING & JUMP
!C                          200-    PRINT WARNING & STOP DEAD
!C-----------------------------------------------------------------------
      INTEGER ICODE
      CHARACTER(len=6) SUBRTN
      WRITE (6,10) SUBRTN,ICODE
10    FORMAT(/' KTWARN CALLED FROM SUBPROGRAM ',A6,': CODE =',I4/)
      IF (ICODE.LT.100) RETURN
      IF (ICODE.LT.200) RETURN 1
      STOP
    END SUBROUTINE KTWARN
!C-----------------------------------------------------------------------
!C-----------------------------------------------------------------------
!C-----------------------------------------------------------------------

\eatline
\eatline
\eatline
\eatline
\eatline
\eatline
\eatline
\eatline
\nwendcode{}\nwbegindocs{393}\nwdocspar
%%%%%%%%%%%%%%%%%%%%%%%%%%%%%%%%%%%%%%%%%%%%%%%%%%%%%%%%%%%%%%%%%%%%%%%%
\subsection{CKKW matching}
\nwenddocs{}\nwbegincode{394}\moddef{\code{}ckkw{\char95}matching.f90\edoc{}}\endmoddef
\LA{}File header\RA{}

module ckkw_matching

\LA{}Use kinds with double\RA{}
  use io_units
  use constants
  use diagnostics
  use physics_defs
  use lorentz
  use shower_base
  use shower_partons
  use shower_core
  use ckkw_pseudo_weights

\LA{}Standard module head\RA{}

\LA{}CKKW matching: public\RA{}

\LA{}CKKW matching: types\RA{}

contains

\LA{}CKKW matching: procedures\RA{}

end module ckkw_matching
\eatline
\nwendcode{}\nwbegindocs{395}\nwdocspar
\nwenddocs{}\nwbegincode{396}\moddef{CKKW matching: public}\endmoddef
  public :: ckkw_matching_settings_t
\nwendcode{}\nwbegincode{397}\moddef{CKKW matching: types}\endmoddef
  type :: ckkw_matching_settings_t
     real(default) :: alphaS = 0.118_default
     real(default) :: Qmin = one
     integer :: n_max_jets = 0
  end type ckkw_matching_settings_t

\eatline
\nwendcode{}\nwbegindocs{398}\nwdocspar
\nwenddocs{}\nwbegincode{399}\moddef{CKKW matching: public}\plusendmoddef
  public :: ckkw_matching_apply
\nwendcode{}\nwbegindocs{400}\nwdocspar
\nwenddocs{}\nwbegincode{401}\moddef{CKKW matching: procedures}\endmoddef
  subroutine ckkw_matching_apply (shower, settings, weights, veto)
    type(shower_t), intent(inout) :: shower
    type(ckkw_matching_settings_t), intent(in) :: settings
    type(ckkw_pseudo_shower_weights_t), intent(in) :: weights
    logical, intent(out) :: veto

    real(default), dimension(:), allocatable :: scales
    real(double) :: weight, sf
    real(default) :: rand
    integer :: i

    if (signal_is_pending ()) return
    weight = one

    call shower%write ()

    !!! the pseudo parton shower is already simulated by shower_add_interaction
    !!! get the respective clustering scales
    allocate (scales(1:size(shower%partons)))
    do i = 1, size (shower%partons)
       if (.not. associated (shower%partons(i)%p)) cycle
       if (shower%partons(i)%p%type == INTERNAL) then
          scales(i) = two * min (shower%partons(i)%p%child1%momentum%p(0),  &
                                 shower%partons(i)%p%child2%momentum%p(0))**2 * &
               (1.0 - (space_part (shower%partons(i)%p%child1%momentum) * &
                space_part (shower%partons(i)%p%child2%momentum)) / &
               (space_part (shower%partons(i)%p%child1%momentum)**1 * &
                space_part (shower%partons(i)%p%child2%momentum)**1))
          scales(i) = sqrt (scales(i))
          shower%partons(i)%p%ckkwscale = scales(i)
          print *, scales(i)
       end if
    end do

    print *, " scales finished"
    !!! if (highest multiplicity) -> reweight with PDF(mu_F) / PDF(mu_cut)
    call shower%write ()

    !!! Reweight and possibly veto the whole event

    !!! calculate the relative alpha_S weight

    !! calculate the Sudakov weights for internal lines
    !! calculate the Sudakov weights for external lines
    do i = 1, size (shower%partons)
       if (signal_is_pending ()) return
       if (.not. associated (shower%partons(i)%p)) cycle
       if (shower%partons(i)%p%type == INTERNAL) then
          !!! get type
          !!! check that all particles involved are colored
          if ((parton_is_colored (shower%partons(i)%p) .or. &
               shower%partons(i)%p%ckkwtype > 0) .and. &
               (parton_is_colored (shower%partons(i)%p%child1) .or. &
               shower%partons(i)%p%child1%ckkwtype > 0) .and. &
               (parton_is_colored (shower%partons(i)%p%child1) .or. &
               shower%partons(i)%p%child1%ckkwtype > 0)) then
             print *, "reweight with alphaS(" , shower%partons(i)%p%ckkwscale, &
                  ") for particle ", shower%partons(i)%p%nr
             if (shower%partons(i)%p%belongstoFSR) then
                print *, "FSR"
                weight = weight * D_alpha_s_fsr (shower%partons(i)%p%ckkwscale**2) &
                     / settings%alphas
             else
                print *, "ISR"
                weight = weight * &
                     D_alpha_s_isr (shower%partons(i)%p%ckkwscale**2) &
                     / settings%alphas
             end if
          else
             print *, "no reweight with alphaS for ", shower%partons(i)%p%nr
          end if
          if (shower%partons(i)%p%child1%type == INTERNAL) then
             print *, "internal line from ", &
                  shower%partons(i)%p%child1%ckkwscale, &
                  " to ", shower%partons(i)%p%ckkwscale, &
                  " for type ", shower%partons(i)%p%child1%ckkwtype
             if (shower%partons(i)%p%child1%ckkwtype == 0) then
                sf = 1.0
             else if (shower%partons(i)%p%child1%ckkwtype == 1) then
                sf = SudakovQ (shower%partons(i)%p%child1%ckkwscale, &
                     shower%partons(i)%p%ckkwscale, .true.)
                print *, "SFQ = ", sf
             else if (shower%partons(i)%p%child1%ckkwtype == 2) then
                sf = SudakovG (shower%partons(i)%p%child1%ckkwscale, &
                     shower%partons(i)%p%ckkwscale, .true.)
                print *, "SFG = ", sf
             else
                print *, "SUSY not yet implemented"
             end if
             weight = weight * min (one, sf)
          else
             print *, "external line from ", settings%Qmin, &
                  shower%partons(i)%p%ckkwscale
             if (parton_is_quark (shower%partons(i)%p%child1)) then
                sf = SudakovQ (settings%Qmin, &
                     shower%partons(i)%p%ckkwscale, .true.)
                print *, "SFQ = ", sf
             else if (parton_is_gluon (shower%partons(i)%p%child1)) then
                sf = SudakovG (settings%Qmin, &
                     shower%partons(i)%p%ckkwscale, .true.)
                print *, "SFG = ", sf
             else
                print *, "not yet implemented (", &
                     shower%partons(i)%p%child2%type, ")"
                sf = one
             end if
             weight = weight * min (one, sf)
          end if
          if (shower%partons(i)%p%child2%type == INTERNAL) then
             print *, "internal line from ", shower%partons(i)%p%child2%ckkwscale, &
                  " to ", shower%partons(i)%p%ckkwscale, &
                  " for type ", shower%partons(i)%p%child2%ckkwtype
             if (shower%partons(i)%p%child2%ckkwtype == 0) then
                sf = 1.0
             else if (shower%partons(i)%p%child2%ckkwtype == 1) then
                sf = SudakovQ (shower%partons(i)%p%child2%ckkwscale, &
                     shower%partons(i)%p%ckkwscale, .true.)
                print *, "SFQ = ", sf
             else if (shower%partons(i)%p%child2%ckkwtype == 2) then
                sf = SudakovG (shower%partons(i)%p%child2%ckkwscale, &
                     shower%partons(i)%p%ckkwscale, .true.)
                print *, "SFG = ", sf
             else
                print *, "SUSY not yet implemented"
             end if
             weight = weight * min (one, sf)
          else
             print *, "external line from ", settings%Qmin, &
                  shower%partons(i)%p%ckkwscale
             if (parton_is_quark (shower%partons(i)%p%child2)) then
                sf = SudakovQ (settings%Qmin, &
                     shower%partons(i)%p%ckkwscale, .true.)
                print *, "SFQ = ", sf
             else if (parton_is_gluon (shower%partons(i)%p%child2)) then
                sf = SudakovG (settings%Qmin, &
                     shower%partons(i)%p%ckkwscale, .true.)
                print *, "SFG = ", sf
             else
                print *, "not yet implemented (", &
                     shower%partons(i)%p%child2%type, ")"
                sf = one
             end if
             weight = weight * min (one, sf)
          end if
       end if
    end do

    call rng%generate (rand)

    print *, "final weight: ", weight

    !!!!!!! WRONG
    veto = .false.
    veto = (rand > weight)
    if (veto) then
       return
    end if

    !!! finally perform the parton shower
    !!! veto emissions that are too hard

    deallocate (scales)
  end subroutine ckkw_matching_apply

\eatline
\nwendcode{}\nwbegindocs{402}\nwdocspar
\nwenddocs{}\nwbegincode{403}\moddef{CKKW matching: procedures}\plusendmoddef
  function GammaQ (smallq, largeq, fsr) result (gamma)
    real(default), intent(in) :: smallq, largeq
    logical, intent(in) :: fsr
    real(default) :: gamma
    gamma = (8._default / three) / (pi * smallq)
    gamma = gamma * (log(largeq / smallq) - 0.75)
    if (fsr) then
       gamma = gamma * D_alpha_s_fsr (smallq**2)
    else
       gamma = gamma * D_alpha_s_isr (smallq**2)
    end if
  end function GammaQ

\eatline
\nwendcode{}\nwbegindocs{404}\nwdocspar
\nwenddocs{}\nwbegincode{405}\moddef{CKKW matching: procedures}\plusendmoddef
  function GammaG(smallq, largeq, fsr) result(gamma)
    real(default), intent(in) :: smallq, largeq
    logical, intent(in) :: fsr
    real(default) :: gamma
    gamma = 6._default / (pi * smallq)
    gamma = gamma *( log(largeq / smallq) - 11.0 / 12.0)
    if (fsr) then
       gamma = gamma * D_alpha_s_fsr(smallq**2)
    else
       gamma = gamma * D_alpha_s_isr(smallq**2)
    end if
  end function GammaG

\eatline
\nwendcode{}\nwbegindocs{406}\nwdocspar
\nwenddocs{}\nwbegincode{407}\moddef{CKKW matching: procedures}\plusendmoddef
  function GammaF (smallq, fsr) result (gamma)
    real(default), intent(in) :: smallq
    logical, intent(in) :: fsr
    real(default) :: gamma
    gamma = number_of_flavors (smallq) / (three * pi * smallq)
    if (fsr) then
       gamma = gamma * D_alpha_s_fsr (smallq**2)
    else
       gamma = gamma * D_alpha_s_isr (smallq**2)
    end if
  end function GammaF

\eatline
\nwendcode{}\nwbegindocs{408}\nwdocspar
\nwenddocs{}\nwbegincode{409}\moddef{CKKW matching: procedures}\plusendmoddef
  function SudakovQ (Q1, Q, fsr) result (sf)
    real(default), intent(in) :: Q1, Q
    logical, intent(in) :: fsr
    real(default) :: sf
    real(default) :: integral
    integer, parameter :: NTRIES = 100
    integer :: i
    real(default) :: rand
    integral = zero
    do i = 1, NTRIES
       call rng%generate (rand)
       integral = integral + GammaQ (Q1 + rand * (Q - Q1), Q, fsr)
    end do
    integral = integral / NTRIES
    sf = exp (-integral)
  end function SudakovQ

\eatline
\nwendcode{}\nwbegindocs{410}\nwdocspar
\nwenddocs{}\nwbegincode{411}\moddef{CKKW matching: procedures}\plusendmoddef
  function SudakovG (Q1, Q, fsr) result (sf)
    real(default), intent(in) :: Q1, Q
    logical, intent(in) :: fsr
    real(default) :: sf
    real(default) :: integral
    integer, parameter :: NTRIES = 100
    integer :: i
    real(default) :: rand
    integral = zero
    do i = 1, NTRIES
       call rng%generate (rand)
       integral = integral + GammaG (Q1 + rand * (Q - Q1), Q, fsr) + &
            GammaF (Q1 +rand * (Q - Q1), fsr)
    end do
    integral = integral / NTRIES
    sf = exp (-integral)
  end function SudakovG

\eatline
\nwendcode{}\nwbegindocs{412}\nwdocspar
%%%%%%%%%%%%%%%%%%%%%%%%%%%%%%%%%%%%%%%%%%%%%%%%%%%%%%%%%%%%%%%%%%%%%%%%
\subsubsection{CKKW (pseudo) weights}

The type {\tt{}ckkw{\char95}pseudo{\char95}shower{\char95}weights{\char95}t} gives the (relative) weights
for different clusterings of the final particles, as given in Eq.~(2.7) of
hep-ph/0503281v1. Each particle has a binary labelling (power of 2)
(first particle = 1, second particle = 2, third particle = 4,
...). Each recombination therefore corresponds to an integer, that is
not a power of 2. Fur multiple subsequent recombinations, no different
weights for different sequences of clustering are stored. It is
assumed that the weight of a multiply recombined state is a
combination of the states with one fewer recombination and that these
states' contributions are proportional to their weights. For a $2->n$
event, the weights array thus has the size $2^{(2 + n) - 1}$. The
{\tt{}weights{\char95}by{\char95}type} array gives the weights depending on the type of
the particle, the first index is the same as for weights, the second
index gives the type of the new mother particle:
\begin{itemize}
  \item[0:] uncolored ($\gamma$, $Z$, $W$, Higgs)
  \item[1:] colored (quark)
  \item[2:] gluon
  \item[3:] squark
  \item[4:] gluino
\end{itemize}
{\tt{}alphaS} gives the value for $alpha_s$ used in the generation of the
matrix element. This is needed for the reweighting using the values
for a running $alpha_s$ at the scales of the clusterings.
\nwenddocs{}\nwbegincode{413}\moddef{\code{}ckkw{\char95}pseudo{\char95}weights.f90\edoc{}}\endmoddef
\LA{}File header\RA{}

module ckkw_pseudo_weights

\LA{}Use kinds\RA{}
  use io_units
  use constants

\LA{}Standard module head\RA{}

\LA{}CKKW pseudo weights: public\RA{}

\LA{}CKKW pseudo weights: types\RA{}

contains

\LA{}CKKW pseudo weights: procedures\RA{}

end module ckkw_pseudo_weights
\eatline
\nwendcode{}\nwbegindocs{414}\nwdocspar
\nwenddocs{}\nwbegincode{415}\moddef{CKKW pseudo weights: public}\endmoddef
  public :: ckkw_pseudo_shower_weights_t
\nwendcode{}\nwbegincode{416}\moddef{CKKW pseudo weights: types}\endmoddef
  type :: ckkw_pseudo_shower_weights_t
     real(default) :: alphaS
     real(default), dimension(:), allocatable :: weights
     real(default), dimension(:,:), allocatable :: weights_by_type
  end type ckkw_pseudo_shower_weights_t

\eatline
\nwendcode{}\nwbegindocs{417}\nwdocspar
\nwenddocs{}\nwbegincode{418}\moddef{CKKW pseudo weights: public}\plusendmoddef
  public :: ckkw_pseudo_shower_weights_init
\nwendcode{}\nwbegincode{419}\moddef{CKKW pseudo weights: procedures}\endmoddef
  subroutine ckkw_pseudo_shower_weights_init (weights)
    type(ckkw_pseudo_shower_weights_t), intent(out) :: weights
    weights%alphaS = zero
    if (allocated (weights%weights))  deallocate(weights%weights)
    if (allocated (weights%weights_by_type))  &
         deallocate(weights%weights_by_type)
  end subroutine ckkw_pseudo_shower_weights_init

\eatline
\nwendcode{}\nwbegindocs{420}\nwdocspar
\nwenddocs{}\nwbegincode{421}\moddef{CKKW pseudo weights: public}\plusendmoddef
  public :: ckkw_pseudo_shower_weights_write
\nwendcode{}\nwbegincode{422}\moddef{CKKW pseudo weights: procedures}\plusendmoddef
  subroutine ckkw_pseudo_shower_weights_write (weights, unit)
    type(ckkw_pseudo_shower_weights_t), intent(in) :: weights
    integer, intent(in), optional :: unit
    integer :: s, i, u
    u = given_output_unit (unit); if (u < 0) return
    s = size (weights%weights)
    write (u, "(1x,A)")  "CKKW (pseudo) shower weights: "
    do i = 1, s
       write (u, "(3x,I0,2(ES19.12))")  i, weights%weights(i), &
            weights%weights_by_type(i,:)
    end do
    write (u, "(3x,A,1x,I0)")  "alphaS =", weights%alphaS
  end subroutine ckkw_pseudo_shower_weights_write

\eatline
\nwendcode{}\nwbegindocs{423}\nwdocspar
%%%%%%%%%%%%%%%%%%%%%%%%%%%%%%%%%%%%%%%%%%%%%%%%%%%%%%%%%%%%%%%%%%%%%%%%
\section{Attic}
\nwenddocs{}\nwbegincode{424}\moddef{Procedures of \code{}shower{\char95}add{\char95}interaction{\char95}2ton{\char95}CKKW\edoc{} (not used)}\endmoddef
function shower_clustering_allowed (shower, partons, i, j) result (allowed)
  type(shower_t), intent(inout) :: shower
  logical :: allowed
  type(parton_pointer_t), intent(in), dimension(:), allocatable :: partons
  integer, intent(in) :: i, j
  !!! TODO implement checking if clustering is allowed, e.g.
  !!!     in e+e- -> qqg don't cluster the quarks together first
  allowed = .true.
end function shower_clustering_allowed

\nwendcode{}\nwbegindocs{425}\nwdocspar
\nwenddocs{}\nwbegincode{426}\moddef{Shower core: procedures (not used)}\endmoddef
  subroutine shower_get_final_partons (shower, partons, include_remnants)
    type(shower_t), intent(in) :: shower
    type(parton_pointer_t), dimension(:), allocatable, intent(inout) :: partons
    logical, intent(in), optional :: include_remnants
    integer :: i, j
    type(parton_t), pointer :: prt
    if (allocated (partons))  deallocate (partons)
    allocate (partons(1:shower_get_nr_of_partons(shower, &
         include_remnants = include_remnants)))
    j = 0
    do i = 1, size (shower%partons)
       prt => shower%partons(i)%p
       if (.not. associated (prt)) cycle
       if (.not. parton_is_final (prt)) cycle
       if (prt%type == BEAM_REMNANT) then
          if (present (include_remnants)) then
             if (.not. include_remnants) cycle
          end if
       end if
       j = j + 1
       partons(j)%p => prt
    end do
  end subroutine shower_get_final_partons

\eatline
\nwendcode{}\nwbegindocs{427}\nwdocspar
\nwenddocs{}\nwbegincode{428}\moddef{Shower core: procedures (not used)}\plusendmoddef
  function shower_interaction_get_shat (interaction) result (shat)
    type(shower_interaction_t), intent(in) :: interaction
    real(default) :: shat
    shat = (interaction%partons(1)%p%momentum + &
            interaction%partons(2)%p%momentum)**2
  end function shower_interaction_get_shat

\eatline
\nwendcode{}\nwbegindocs{429}\nwdocspar
\nwenddocs{}\nwbegincode{430}\moddef{Shower core: procedures (not used)}\plusendmoddef
  function shower_is_finished (shower) result (finished)
    type(shower_t), intent(in) :: shower
    logical :: finished
    finished = shower_isr_is_finished (shower) .and. &
               shower_fsr_is_finished(shower)
  end function shower_is_finished

\eatline
\nwendcode{}\nwbegindocs{431}\nwdocspar
\nwenddocs{}\nwbegincode{432}\moddef{Shower core: procedures (not used)}\plusendmoddef
  subroutine shower_prepare_for_simulate_isr_ana (shower, prt1, prt2)
    type(shower_t), intent(inout) :: shower
    type(parton_t), intent(inout), target :: prt1, prt2
    type(parton_t), pointer :: prt
    real(default) :: scale, factor
    real(default) :: oldscales(1:2)
    type(parton_pointer_t) :: temppp
    integer :: i

    if (signal_is_pending ()) return

    if (.not. associated (prt1%initial) .or. &
        .not. associated (prt2%initial)) then
       return
    end if

    scale = - (prt1%momentum + prt2%momentum) ** 2

    prt1%t = -shower%tscalefactor_isr * abs(scale)
    prt2%t = -shower%tscalefactor_isr * abs(scale)
    !!! rescale momenta
    do i = 1, 2
       if (i == 1) then
          prt => prt1
       else
          prt => prt2
       end if

       factor = sqrt (prt%momentum%p(0)**2 - prt%t) / &
            space_part_norm(prt%momentum)

       prt%momentum = vector4_moving (prt%momentum%p(0), &
            factor * space_part (prt%momentum))
    end do

    !!! ensure that belongstointeraction bits are set correctly
    prt1%belongstointeraction = .true.
    prt2%belongstointeraction = .true.

    call shower%add_parent (prt1)
    call shower%add_parent (prt2)

    call parton_set_simulated (prt1)
    prt1%parent%type = prt1%type
    prt1%parent%z = one
    prt1%parent%momentum = prt1%momentum
    prt1%parent%t = scale
    prt1%parent%x = prt1%x
    prt1%parent%initial => prt1%initial
    prt1%parent%belongstoFSR = .false.
    prt1%parent%c1 = prt1%c1
    prt1%parent%c2 = prt1%c2
    call shower%add_child (prt1%parent, 2)

    call parton_set_simulated (prt2)
    prt2%parent%type = prt2%type
    prt2%parent%z = one
    prt2%parent%momentum = prt2%momentum
    prt2%parent%t = scale
    prt2%parent%x = prt2%x
    prt2%parent%initial => prt2%initial
    prt2%parent%belongstoFSR = .false.
    prt2%parent%c1 = prt2%c1
    prt2%parent%c2 = prt2%c2
    call shower%add_child (prt2%parent, 2)

    FIRST_BRANCHINGS: do
       if (signal_is_pending ()) return
       oldscales(1) = prt1%parent%t
       oldscales(2) = prt2%parent%t
       if (abs(prt1%parent%t) > abs(prt2%parent%t)) then
          temppp%p => prt1%parent
       else
          temppp%p => prt2%parent
       end if
       if (.not. temppp%p%simulated .and. .not. parton_is_hadron &
            (temppp%p)) then
          call shower_isr_step (shower, temppp%p)
          if (temppp%p%simulated) then
             ! call parton_generate_ps_ini(prt2%parent)
             if (temppp%p%t < zero) then
                call shower%execute_next_isr_branching (temppp)
                ! call shower%write ()
             else
                call shower_replace_parent_by_hadron (shower, temppp%p%child1)
             end if
          end if
       end if

       if (oldscales(1) == prt1%parent%t .and. &
           oldscales(2) == prt2%parent%t) then
          exit FIRST_BRANCHINGS
       end if
    end do FIRST_BRANCHINGS

  end subroutine shower_prepare_for_simulate_isr_ana

\eatline
\nwendcode{}\nwbegindocs{433}\nwdocspar
\nwenddocs{}\nwbegincode{434}\moddef{Shower core: procedures (not used)}\plusendmoddef
  subroutine shower_generate_next_fsr_branchings (shower)
    type(shower_t), intent(inout) :: shower
    integer i, index
    type(parton_t),  pointer :: prt

    !!! find mother with highest t to be simulated
    index = 0
    do i = 1, size (shower%partons)
       prt => shower%partons(i)%p
       if (.not. prt%belongstoFSR) cycle
       if (prt%belongstointeraction) cycle
       if (associated(prt%child1) .and. associated(prt%child2)) then
          if (prt%child1%simulated .and. &
              prt%child2%simulated) cycle
       end if
       if (parton_is_final (prt)) cycle

         index = i
         exit
      end do

      if (index == 0) then
         call msg_message ("Shower: no branchable partons found")
         return
      end if

      prt => shower%partons(index)%p
      call shower_simulate_children_ana (shower, prt)

    end subroutine shower_generate_next_fsr_branchings

\eatline
\nwendcode{}\nwbegindocs{435}\nwdocspar
\nwenddocs{}\nwbegincode{436}\moddef{Shower2pythia: procedures (not used)}\endmoddef
  recursive subroutine shower_topythia_recursive_weighted (prt, mode, first)
    IMPLICIT DOUBLE PRECISION(A-H, O-Z)
    IMPLICIT INTEGER(I-N)
    !    C...  Commonblocks.
    COMMON/PYJETS/N,NPAD,K(4000,5),P(4000,5),V(4000,5)
    COMMON/PYDAT1/MSTU(200),PARU(200),MSTJ(200),PARJ(200)
    COMMON/PYDAT2/KCHG(500,4),PMAS(500,4),PARF(2000),VCKM(4,4)
    COMMON/PYDAT3/MDCY(500,3),MDME(8000,2),BRAT(8000),KFDP(8000,5)
    COMMON/PYSUBS/MSEL,MSELPD,MSUB(500),KFIN(2,-40:40),CKIN(200)
    COMMON/PYPARS/MSTP(200),PARP(200),MSTI(200),PARI(200)
    SAVE /PYJETS/,/PYDAT1/,/PYDAT2/,/PYDAT3/,/PYSUBS/,/PYPARS/

    type(parton_t), intent(in), target :: prt
    integer, intent(in) :: mode, first
    integer :: n_emissions
    type(parton_t), pointer :: tempprt, finalprt
    real(default), dimension(:), allocatable :: costhetas
    type(parton_pointer_t), dimension(:), allocatable :: emittedpartons
    integer :: i, max
    real(default) :: maxcostheta

    n_emissions = 0

    if (parton_is_final(prt)) then
       N = N + 1
       K(N,1) = 2
       if (parton_is_quark(prt)) then
          ! check if quark is end of a color connection
          if (prt%type < 0) then
             K(N,1) = 1
          end if
       end if
       K(N,2) = prt%type
       K(N,3) = first
       K(N,4) = 0
       K(N,5) = 0
       P(N,1) = prt%momentum%p(1)
       P(N,2) = prt%momentum%p(2)
       P(N,3) = prt%momentum%p(3)
       P(N,4) = prt%momentum%p(0)
       P(N,5) = prt%t
    else
       !!! search for following final partons
       if (parton_is_gluon (prt)) then
          if (parton_is_gluon(prt%child1)) then
             !!! g-> gg so sequence is unimportant
             call shower_topythia_recursive_weighted (prt%child1, 1, first)
             call shower_topythia_recursive_weighted (prt%child2, 1, first)
          else
             !!! g-> qqbar -> use antiquark first,
             !!!       so that color flow is given correctly
             if (prt%child1%type < 0) then
                call shower_topythia_recursive_weighted (prt%child1, 1, first)
                call shower_topythia_recursive_weighted (prt%child2, 2, first)
             else
                call shower_topythia_recursive_weighted (prt%child2, 2, first)
                call shower_topythia_recursive_weighted (prt%child1, 1, first)
             end if
          end if
       else
          !!! parton is quark
          !!! find the emitted gluons and order them by the emission angle
          n_emissions = 0
          tempprt = prt
          do
             !!! calculate how many emissions there are
             if (associated (tempprt%child1)) then
                tempprt => tempprt%child1
                n_emissions = n_emissions +1
                cycle
             else
                exit
             end if
          end do

          allocate (costhetas(1:n_emissions))
          allocate (emittedpartons(1:n_emissions))

          tempprt => prt
          n_emissions = 1
          do
             if (associated (tempprt%child1)) then
                costhetas(n_emissions) = parton_get_costheta_correct (tempprt)
                emittedpartons(n_emissions)%p => tempprt%child2
                n_emissions = n_emissions +1
                tempprt => tempprt%child1
                cycle
             else
                finalprt=>tempprt
                exit
             end if
          end do

          !!! if mode == 1 write quark first
          if (mode == 1) &
               call shower_topythia_recursive_weighted (finalprt, 1, first)

          !!! if mode == 2 write out gluons in recursive order <= replace
          !!!      costheta by 1- costheta
          if (mode == 2) then
             do i = 1, size (costhetas)
                costhetas(i) = one - costhetas(i)
             end do
          end if

          do
             max = 0
             maxcostheta = zero
             do i = 1, size (costhetas)
                if (costhetas(i) > maxcostheta) then
                   maxcostheta = costhetas(i)
                   max = i
                end if
             end do
             if (maxcostheta == zero) then
                exit
             end if

             call shower_topythia_recursive_weighted &
                  (emittedpartons(max)%p, mode, first)
             costhetas(max) = zero
          end do

          !!! if mode == 2 write quark last
          if (mode == 2) then
             call shower_topythia_recursive_weighted (finalprt, 2, first)
          end if

       end if
    end if
  end subroutine shower_topythia_recursive_weighted

\eatline
\nwendcode{}\nwbegindocs{437}\nwdocspar
\nwenddocs{}\nwbegincode{438}\moddef{MLM matching: procedures (not used)}\endmoddef
  SUBROUTINE KTCLUS(IMODE,PP,NN,ECUT,Y,*)
    IMPLICIT NONE
!C---DO CLUSTER ANALYSIS OF PARTICLES IN PP
!C
!C   IMODE   = INPUT  : DESCRIBED ABOVE
!C   PP(I,J) = INPUT  : 4-MOMENTUM OF Jth PARTICLE: I=1,4 => PX,PY,PZ,E
!C   NN      = INPUT  : NUMBER OF PARTICLES
!C   ECUT    = INPUT  : DENOMINATOR OF KT MEASURE. IF ZERO, ETOT IS USED
!C   Y(J)    = OUTPUT : VALUE OF Y FOR WHICH EVENT CHANGES FROM BEING
!C                        J JET TO J-1 JET
!C   LAST ARGUMENT IS LABEL TO JUMP TO IF FOR ANY REASON THE EVENT
!C   COULD NOT BE PROCESSED (MOST LIKELY DUE TO TOO MANY PARTICLES)
!C
!C   NOTE THAT THE MOMENTA ARE DECLARED DOUBLE PRECISION,
!C   AND ALL OTHER FLOATING POINT VARIABLES ARE DECLARED DOUBLE PRECISION
!C
    INTEGER IMODE,NN
    DOUBLE PRECISION PP(4,*)
    DOUBLE PRECISION ECUT,Y(*),ONE
    ONE=1
    CALL KTCLUR(IMODE,PP,NN,ONE,ECUT,Y,*999)
    RETURN
999 RETURN 1
  END SUBROUTINE KTCLUS
!C-----------------------------------------------------------------------
\nwendcode{}\nwbegincode{439}\moddef{MLM matching: procedures (not used)}\plusendmoddef
      SUBROUTINE KTYCUT(ECUT,NY,YCUT,NJET,*)
      IMPLICIT NONE
!C---COUNT THE NUMBER OF JETS AT EACH VALUE OF YCUT, FOR EVENT WHICH HAS
!C   ALREADY BEEN ANALYSED BY KTCLUS.
!C
!C   ECUT    = INPUT : DENOMINATOR OF KT MEASURE. IF ZERO, ETOT IS USED
!C   NY      = INPUT : NUMBER OF YCUT VALUES
!C   YCUT(J) = INPUT : Y VALUES AT WHICH NUMBERS OF JETS ARE COUNTED
!C   NJET(J) =OUTPUT : NUMBER OF JETS AT YCUT(J)
!C   LAST ARGUMENT IS LABEL TO JUMP TO IF FOR ANY REASON THE EVENT
!C   COULD NOT BE PROCESSED
!C
!C   NOTE THAT ALL FLOATING POINT VARIABLES ARE DECLARED DOUBLE PRECISION
!C
      INTEGER NY,NJET(NY),NMAX,HIST,I,J,NUM
      PARAMETER (NMAX=512)
      DOUBLE PRECISION YCUT(NY),ETOT,RSQ,P,KT,KTP,KTS,ETSQ,ECUT,KTLAST, &
           ROUND
      PARAMETER (ROUND=0.99999D0)
      COMMON /KTCOMM/ETOT,RSQ,P(9,NMAX),KTP(NMAX,NMAX),KTS(NMAX), &
           KT(NMAX),KTLAST(NMAX),HIST(NMAX),NUM
      IF (ETOT.EQ.0) CALL KTWARN('KTYCUT',100,*999)
      IF (ECUT.EQ.0) THEN
         ETSQ=1/ETOT**2
      ELSE
         ETSQ=1/ECUT**2
      ENDIF
      DO I=1,NY
         NJET(I)=0
      END DO
      DO I=NUM,1,-1
         DO J=1,NY
            IF (NJET(J).EQ.0.AND.KT(I)*ETSQ.GE.ROUND*YCUT(J)) NJET(J)=I
         END DO
      END DO
      RETURN
999   RETURN 1
    END SUBROUTINE KTYCUT
!C-----------------------------------------------------------------------
    SUBROUTINE KTYSUB(ECUT,NY,YCUT,YMAC,NSUB,*)
      IMPLICIT NONE
!C---COUNT THE NUMBER OF SUB-JETS AT EACH VALUE OF YCUT, FOR EVENT WHICH
!C   HAS ALREADY BEEN ANALYSED BY KTCLUS.
!C   REMEMBER THAT A SUB-JET IS DEFINED AS A JET AT Y=YCUT WHICH HAS NOT
!C   YET BEEN MERGED WITH THE BEAM AT Y=YMAC.
!C
!C   ECUT    = INPUT : DENOMINATOR OF KT MEASURE. IF ZERO, ETOT IS USED
!C   NY      = INPUT : NUMBER OF YCUT VALUES
!C   YCUT(J) = INPUT : Y VALUES AT WHICH NUMBERS OF SUB-JETS ARE COUNTED
!C   YMAC    = INPUT : Y VALUE USED TO DEFINE MACRO-JETS, TO DETERMINE
!C                       WHICH JETS ARE SUB-JETS
!C   NSUB(J) =OUTPUT : NUMBER OF SUB-JETS AT YCUT(J)
!C   LAST ARGUMENT IS LABEL TO JUMP TO IF FOR ANY REASON THE EVENT
!C   COULD NOT BE PROCESSED
!C
!C   NOTE THAT ALL FLOATING POINT VARIABLES ARE DECLARED DOUBLE PRECISION
!C
      INTEGER NY,NSUB(NY),NMAX,HIST,I,J,NUM
      PARAMETER (NMAX=512)
      DOUBLE PRECISION YCUT(NY),YMAC,ETOT,RSQ,P,KT,KTP,KTS,ETSQ,ECUT, &
           KTLAST,ROUND
      PARAMETER (ROUND=0.99999D0)
      COMMON /KTCOMM/ETOT,RSQ,P(9,NMAX),KTP(NMAX,NMAX),KTS(NMAX), &
           KT(NMAX),KTLAST(NMAX),HIST(NMAX),NUM
      IF (ETOT.EQ.0) CALL KTWARN('KTYSUB',100,*999)
      IF (ECUT.EQ.0) THEN
         ETSQ=1/ETOT**2
      ELSE
         ETSQ=1/ECUT**2
      ENDIF
      DO I=1,NY
         NSUB(I)=0
      END DO
      DO I=NUM,1,-1
         DO J=1,NY
            IF (NSUB(J).EQ.0.AND.KT(I)*ETSQ.GE.ROUND*YCUT(J)) NSUB(J)=I
            IF (NSUB(J).NE.0.AND.KTLAST(I)*ETSQ.LT.ROUND*YMAC) &
                 NSUB(J)=NSUB(J)-1
         END DO
      END DO
      RETURN
999   RETURN 1
    END SUBROUTINE KTYSUB
!C-----------------------------------------------------------------------
    SUBROUTINE KTBEAM(ECUT,Y,*)
      IMPLICIT NONE
!C---GIVE SAME INFORMATION AS LAST CALL TO KTCLUS EXCEPT THAT ONLY
!C   TRANSITIONS WHERE A JET WAS MERGED WITH THE BEAM JET ARE RECORDED
!C
!C   ECUT    = INPUT : DENOMINATOR OF KT MEASURE. IF ZERO, ETOT IS USED
!C   Y(J)    =OUTPUT : Y VALUE WHERE Jth HARDEST JET WAS MERGED WITH BEAM
!C   LAST ARGUMENT IS LABEL TO JUMP TO IF FOR ANY REASON THE EVENT
!C   COULD NOT BE PROCESSED
!C
!C   NOTE THAT ALL FLOATING POINT VARIABLES ARE DECLARED DOUBLE PRECISION
!C
      INTEGER NMAX,HIST,NUM,I,J
      PARAMETER (NMAX=512)
      DOUBLE PRECISION ETOT,RSQ,P,KT,KTP,KTS,ECUT,ETSQ,Y(*),KTLAST
      COMMON /KTCOMM/ETOT,RSQ,P(9,NMAX),KTP(NMAX,NMAX),KTS(NMAX), &
           KT(NMAX),KTLAST(NMAX),HIST(NMAX),NUM
      IF (ETOT.EQ.0) CALL KTWARN('KTBEAM',100,*999)
      IF (ECUT.EQ.0) THEN
         ETSQ=1/ETOT**2
      ELSE
         ETSQ=1/ECUT**2
      ENDIF
      J=1
      DO I=1,NUM
         IF (HIST(I).LE.NMAX) THEN
            Y(J)=ETSQ*KT(I)
            J=J+1
         ENDIF
      END DO
      DO I=J,NUM
         Y(I)=0
      END DO
      RETURN
999   RETURN 1
    END SUBROUTINE KTBEAM
!C-----------------------------------------------------------------------
    SUBROUTINE KTJOIN(ECUT,YMAC,Y,*)
      IMPLICIT NONE
!C---GIVE SAME INFORMATION AS LAST CALL TO KTCLUS EXCEPT THAT ONLY
!C   TRANSITIONS WHERE TWO SUB-JETS WERE JOINED ARE RECORDED
!C   REMEMBER THAT A SUB-JET IS DEFINED AS A JET AT Y=YCUT WHICH HAS NOT
!C   YET BEEN MERGED WITH THE BEAM AT Y=YMAC.
!C
!C   ECUT    = INPUT : DENOMINATOR OF KT MEASURE. IF ZERO, ETOT IS USED
!C   YMAC    = INPUT : VALUE OF Y USED TO DEFINE MACRO-JETS
!C   Y(J)    =OUTPUT : Y VALUE WHERE EVENT CHANGED FROM HAVING
!C                         N+J SUB-JETS TO HAVING N+J-1, WHERE N IS
!C                         THE NUMBER OF MACRO-JETS AT SCALE YMAC
!C   LAST ARGUMENT IS LABEL TO JUMP TO IF FOR ANY REASON THE EVENT
!C   COULD NOT BE PROCESSED
!C
!C   NOTE THAT ALL FLOATING POINT VARIABLES ARE DECLARED DOUBLE PRECISION
!C
      INTEGER NMAX,HIST,NUM,I,J
      PARAMETER (NMAX=512)
      DOUBLE PRECISION ETOT,RSQ,P,KT,KTP,KTS,ECUT,ETSQ,Y(*),YMAC,KTLAST, &
           ROUND
      PARAMETER (ROUND=0.99999D0)
      COMMON /KTCOMM/ETOT,RSQ,P(9,NMAX),KTP(NMAX,NMAX),KTS(NMAX), &
           KT(NMAX),KTLAST(NMAX),HIST(NMAX),NUM
      IF (ETOT.EQ.0) CALL KTWARN('KTJOIN',100,*999)
      IF (ECUT.EQ.0) THEN
         ETSQ=1/ETOT**2
      ELSE
         ETSQ=1/ECUT**2
      ENDIF
      J=1
      DO I=1,NUM
         IF (HIST(I).GT.NMAX.AND.ETSQ*KTLAST(I).GE.ROUND*YMAC) THEN
            Y(J)=ETSQ*KT(I)
            J=J+1
         ENDIF
      END DO
      DO I=J,NUM
         Y(I)=0
      END DO
      RETURN
999   RETURN 1
    END SUBROUTINE KTJOIN
\nwendcode{}\nwbegincode{440}\moddef{MLM matching: procedures (not used)}\plusendmoddef
    SUBROUTINE KTINCL(RECO,PP,NN,PJET,JET,NJET,*)
      IMPLICIT NONE
!C---RECONSTRUCT KINEMATICS OF JET SYSTEM, WHICH HAS ALREADY BEEN
!C   ANALYSED BY KTCLUS ACCORDING TO THE INCLUSIVE JET DEFINITION. NOTE
!C   THAT NO CONSISTENCY CHECK IS MADE: USER IS TRUSTED TO USE THE SAME
!C   PP VALUES AS FOR KTCLUS
!C
!C   RECO     = INPUT : RECOMBINATION SCHEME (NEED NOT BE SAME AS KTCLUS)
!C   PP(I,J)  = INPUT : 4-MOMENTUM OF Jth PARTICLE: I=1,4 => PX,PY,PZ,E
!C   NN       = INPUT : NUMBER OF PARTICLES
!C   PJET(I,J)=OUTPUT : 4-MOMENTUM OF Jth JET AT SCALE YCUT
!C   JET(J)   =OUTPUT : THE JET WHICH CONTAINS THE Jth PARTICLE
!C   NJET     =OUTPUT : THE NUMBER OF JETS
!C   LAST ARGUMENT IS LABEL TO JUMP TO IF FOR ANY REASON THE EVENT
!C   COULD NOT BE PROCESSED
!C
!C   NOTE THAT THE MOMENTA ARE DECLARED DOUBLE PRECISION,
!C   AND ALL OTHER FLOATING POINT VARIABLES ARE DECLARED DOUBLE PRECISION
!C
      INTEGER NMAX,RECO,NUM,N,NN,NJET,JET(*),HIST,IMIN,JMIN,I,J
      PARAMETER (NMAX=512)
      DOUBLE PRECISION PP(4,*),PJET(4,*)
      DOUBLE PRECISION P,KT,KTP,KTS,ETOT,RSQ,KTLAST
      COMMON /KTCOMM/ETOT,RSQ,P(9,NMAX),KTP(NMAX,NMAX),KTS(NMAX), &
           KT(NMAX),KTLAST(NMAX),HIST(NMAX),NUM
!C---CHECK INPUT
      IF (RECO.LT.1.OR.RECO.GT.3) CALL KTWARN('KTINCL',100,*999)
!C---COPY PP TO P
      N=NN
      IF (NUM.NE.NN) CALL KTWARN('KTINCL',101,*999)
      CALL KTCOPY(PP,N,P,(RECO.NE.1))
!C---INITIALLY EVERY PARTICLE IS IN ITS OWN JET
      DO I=1,NN
         JET(I)=I
      END DO
!C---KEEP MERGING TO THE BITTER END
      NJET=0
200   IF (N.GT.0) THEN
         IF (HIST(N).LE.NMAX) THEN
            IMIN=0
            JMIN=HIST(N)
            NJET=NJET+1
            IF (RECO.EQ.1) THEN
               DO J=1,4
                  PJET(J,NJET)=P(J,JMIN)
               END DO
            ELSE
               PJET(1,NJET)=P(6,JMIN)*COS(P(8,JMIN))
               PJET(2,NJET)=P(6,JMIN)*SIN(P(8,JMIN))
               PJET(3,NJET)=P(6,JMIN)*SINH(P(7,JMIN))
               PJET(4,NJET)=P(6,JMIN)*COSH(P(7,JMIN))
            ENDIF
            CALL KTMOVE(P,KTP,KTS,NMAX,N,JMIN,0)
         ELSE
            IMIN=HIST(N)/NMAX
            JMIN=HIST(N)-IMIN*NMAX
            CALL KTMERG(P,KTP,KTS,NMAX,IMIN,JMIN,N,0,0,0,RECO)
            CALL KTMOVE(P,KTP,KTS,NMAX,N,JMIN,0)
         ENDIF
         DO I=1,NN
            IF (JET(I).EQ.JMIN) JET(I)=IMIN
            IF (JET(I).EQ.N) JET(I)=JMIN
            IF (JET(I).EQ.0) JET(I)=-NJET
         END DO
         N=N-1
         GOTO 200
      ENDIF
!C---FINALLY EVERY PARTICLE MUST BE IN AN INCLUSIVE JET
      DO I=1,NN
!C---IF THERE ARE ANY UNASSIGNED PARTICLES SOMETHING MUST HAVE GONE WRONG
         IF (JET(I).GE.0) CALL KTWARN('KTINCL',102,*999)
         JET(I)=-JET(I)
      END DO
      RETURN
 999  RETURN 1
    END SUBROUTINE KTINCL
!C-----------------------------------------------------------------------
    SUBROUTINE KTISUB(N,NY,YCUT,NSUB,*)
      IMPLICIT NONE
!C---COUNT THE NUMBER OF SUB-JETS IN THE Nth INCLUSIVE JET OF AN EVENT
!C   THAT HAS ALREADY BEEN ANALYSED BY KTCLUS.
!C
!C   N       = INPUT : WHICH INCLUSIVE JET TO USE
!C   NY      = INPUT : NUMBER OF YCUT VALUES
!C   YCUT(J) = INPUT : Y VALUES AT WHICH NUMBERS OF SUB-JETS ARE COUNTED
!C   NSUB(J) =OUTPUT : NUMBER OF SUB-JETS AT YCUT(J)
!C   LAST ARGUMENT IS LABEL TO JUMP TO IF FOR ANY REASON THE EVENT
!C   COULD NOT BE PROCESSED
!C
!C   NOTE THAT ALL FLOATING POINT VARIABLES ARE DECLARED DOUBLE PRECISION
!C
      INTEGER N,NY,NSUB(NY),NMAX,HIST,I,J,NUM,NM
      PARAMETER (NMAX=512)
      DOUBLE PRECISION YCUT(NY),ETOT,RSQ,P,KT,KTP,KTS,KTLAST,ROUND,EPS
      PARAMETER (ROUND=0.99999D0)
      COMMON /KTCOMM/ETOT,RSQ,P(9,NMAX),KTP(NMAX,NMAX),KTS(NMAX), &
           KT(NMAX),KTLAST(NMAX),HIST(NMAX),NUM
      DATA EPS/1D-6/
      DO I=1,NY
         NSUB(I)=0
      END DO
!C---FIND WHICH MERGING CORRESPONDS TO THE NTH INCLUSIVE JET
      NM=0
      J=0
      DO I=NUM,1,-1
        IF (HIST(I).LE.NMAX) J=J+1
        IF (J.EQ.N) THEN
          NM=I
          GOTO 120
        ENDIF
     END DO
120  CONTINUE
!C---GIVE UP IF THERE ARE LESS THAN N INCLUSIVE JETS
      IF (NM.EQ.0) CALL KTWARN('KTISUB',100,*999)
      DO I=NUM,1,-1
         DO J=1,NY
            IF (NSUB(J).EQ.0.AND.RSQ*KT(I).GE.ROUND*YCUT(J)*KT(NM)) &
                 NSUB(J)=I
            IF (NSUB(J).NE.0.AND.ABS(KTLAST(I)-KTLAST(NM)).GT.EPS) &
                 NSUB(J)=NSUB(J)-1
         END DO
      END DO
      RETURN
999   RETURN 1
    END SUBROUTINE KTISUB
!C-----------------------------------------------------------------------
    SUBROUTINE KTIJOI(N,Y,*)
      IMPLICIT NONE
!C---GIVE SAME INFORMATION AS LAST CALL TO KTCLUS EXCEPT THAT ONLY
!C   MERGES OF TWO SUB-JETS INSIDE THE Nth INCLUSIVE JET ARE RECORDED
!C
!C   N       = INPUT : WHICH INCLUSIVE JET TO USE
!C   Y(J)    =OUTPUT : Y VALUE WHERE JET CHANGED FROM HAVING
!C                         J+1 SUB-JETS TO HAVING J
!C   LAST ARGUMENT IS LABEL TO JUMP TO IF FOR ANY REASON THE EVENT
!C   COULD NOT BE PROCESSED
!C
!C   NOTE THAT ALL FLOATING POINT VARIABLES ARE DECLARED DOUBLE PRECISION
!C
      INTEGER NMAX,HIST,NUM,I,J,N,NM
      PARAMETER (NMAX=512)
      DOUBLE PRECISION ETOT,RSQ,P,KT,KTP,KTS,Y(*),KTLAST,EPS
      COMMON /KTCOMM/ETOT,RSQ,P(9,NMAX),KTP(NMAX,NMAX),KTS(NMAX), &
           KT(NMAX),KTLAST(NMAX),HIST(NMAX),NUM
      DATA EPS/1D-6/
!C---FIND WHICH MERGING CORRESPONDS TO THE NTH INCLUSIVE JET
      NM=0
      J=0
      DO I=NUM,1,-1
        IF (HIST(I).LE.NMAX) J=J+1
        IF (J.EQ.N) THEN
          NM=I
          GOTO 105
        ENDIF
     END DO
105  CONTINUE
!C---GIVE UP IF THERE ARE LESS THAN N INCLUSIVE JETS
     IF (NM.EQ.0) CALL KTWARN('KTIJOI',100,*999)
     J=1
     DO I=1,NUM
        IF (HIST(I).GT.NMAX.AND.ABS(KTLAST(I)-KTLAST(NM)).LT.EPS) THEN
           Y(J)=RSQ*KT(I)/KT(NM)
           J=J+1
        ENDIF
     END DO
     DO I=J,NUM
        Y(I)=0
     END DO
      RETURN
999   RETURN 1
    END SUBROUTINE KTIJOI
!C-----------------------------------------------------------------------
    SUBROUTINE KTIREC(RECO,PP,NN,N,YCUT,PSUB,NSUB,*)
      IMPLICIT NONE
!C---RECONSTRUCT KINEMATICS OF SUB-JET SYSTEM IN THE Nth INCLUSIVE JET
!C   OF AN EVENT THAT HAS ALREADY BEEN ANALYSED BY KTCLUS
!C
!C   RECO     = INPUT : RECOMBINATION SCHEME (NEED NOT BE SAME AS KTCLUS)
!C   PP(I,J)  = INPUT : 4-MOMENTUM OF Jth PARTICLE: I=1,4 => PX,PY,PZ,E
!C   NN       = INPUT : NUMBER OF PARTICLES
!C   N        = INPUT : WHICH INCLUSIVE JET TO USE
!C   YCUT     = INPUT : Y VALUE AT WHICH TO RECONSTRUCT JET MOMENTA
!C   PSUB(I,J)=OUTPUT : 4-MOMENTUM OF Jth SUB-JET AT SCALE YCUT
!C   NSUB     =OUTPUT : THE NUMBER OF SUB-JETS
!C   LAST ARGUMENT IS LABEL TO JUMP TO IF FOR ANY REASON THE EVENT
!C   COULD NOT BE PROCESSED
!C
!C   NOTE THAT THE MOMENTA ARE DECLARED DOUBLE PRECISION,
!C   AND ALL OTHER FLOATING POINT VARIABLES ARE DECLARED DOUBLE PRECISION
!C
      INTEGER NMAX,RECO,NUM,NN,NJET,NSUB,JET,HIST,I,J,N,NM
      PARAMETER (NMAX=512)
      DOUBLE PRECISION PP(4,*),PSUB(4,*)
      DOUBLE PRECISION ECUT,P,KT,KTP,KTS,ETOT,RSQ,YCUT,YMAC,KTLAST
      COMMON /KTCOMM/ETOT,RSQ,P(9,NMAX),KTP(NMAX,NMAX),KTS(NMAX), &
           KT(NMAX),KTLAST(NMAX),HIST(NMAX),NUM
      DIMENSION JET(NMAX)
!C---FIND WHICH MERGING CORRESPONDS TO THE NTH INCLUSIVE JET
      NM=0
      J=0
      DO I=NUM,1,-1
         IF (HIST(I).LE.NMAX) J=J+1
         IF (J.EQ.N) THEN
            NM=I
            GOTO 110
         ENDIF
      END DO
110   CONTINUE
!C---GIVE UP IF THERE ARE LESS THAN N INCLUSIVE JETS
      IF (NM.EQ.0) CALL KTWARN('KTIREC',102,*999)
!C---RECONSTRUCT THE JETS AT THE APPROPRIATE SCALE
      ECUT=SQRT(KT(NM)/RSQ)
      YMAC=RSQ
      CALL KTRECO(RECO,PP,NN,ECUT,YCUT,YMAC,PSUB,JET,NJET,NSUB,*999)
!C---GET RID OF THE ONES THAT DO NOT END UP IN THE JET WE WANT
      NSUB=0
      DO I=1,NJET
         IF (JET(I).EQ.HIST(NM)) THEN
            NSUB=NSUB+1
            DO J=1,4
               PSUB(J,NSUB)=PSUB(J,I)
            END DO
         ENDIF
      END DO
      RETURN
999   RETURN 1
    END SUBROUTINE KTIREC
!C-----------------------------------------------------------------------
    SUBROUTINE KTWICH(ECUT,YCUT,JET,NJET,*)
      IMPLICIT NONE
!C---GIVE A LIST OF WHICH JET EACH ORIGINAL PARTICLE ENDED UP IN AT SCALE
!C   YCUT, TOGETHER WITH THE NUMBER OF JETS AT THAT SCALE.
!C
!C   ECUT     = INPUT : DENOMINATOR OF KT MEASURE. IF ZERO, ETOT IS USED
!C   YCUT     = INPUT : Y VALUE AT WHICH TO DEFINE JETS
!C   JET(J)   =OUTPUT : THE JET WHICH CONTAINS THE Jth PARTICLE,
!C                        SET TO ZERO IF IT WAS PUT INTO THE BEAM JETS
!C   NJET     =OUTPUT : THE NUMBER OF JETS AT SCALE YCUT (SO JET()
!C                        ENTRIES WILL BE IN THE RANGE 0 -> NJET)
!C   LAST ARGUMENT IS LABEL TO JUMP TO IF FOR ANY REASON THE EVENT
!C   COULD NOT BE PROCESSED
!C
!C   NOTE THAT ALL FLOATING POINT VARIABLES ARE DECLARED DOUBLE PRECISION
!C
      INTEGER JET(*),NJET,NTEMP
      DOUBLE PRECISION ECUT,YCUT
      CALL KTWCHS(ECUT,YCUT,YCUT,JET,NJET,NTEMP,*999)
      RETURN
 999  RETURN 1
    END SUBROUTINE KTWICH
!C-----------------------------------------------------------------------
    SUBROUTINE KTWCHS(ECUT,YCUT,YMAC,JET,NJET,NSUB,*)
      IMPLICIT NONE
!C---GIVE A LIST OF WHICH SUB-JET EACH ORIGINAL PARTICLE ENDED UP IN AT
!C   SCALE YCUT, WITH MACRO-JET SCALE YMAC, TOGETHER WITH THE NUMBER OF
!C   JETS AT SCALE YCUT AND THE NUMBER OF THEM WHICH ARE SUB-JETS.
!C
!C   ECUT     = INPUT : DENOMINATOR OF KT MEASURE. IF ZERO, ETOT IS USED
!C   YCUT     = INPUT : Y VALUE AT WHICH TO DEFINE JETS
!C   YMAC     = INPUT : Y VALUE AT WHICH TO DEFINE MACRO-JETS
!C   JET(J)   =OUTPUT : THE JET WHICH CONTAINS THE Jth PARTICLE,
!C                        SET TO ZERO IF IT WAS PUT INTO THE BEAM JETS
!C   NJET     =OUTPUT : THE NUMBER OF JETS AT SCALE YCUT (SO JET()
!C                        ENTRIES WILL BE IN THE RANGE 0 -> NJET)
!C   NSUB     =OUTPUT : THE NUMBER OF SUB-JETS AT SCALE YCUT, WITH
!C                        MACRO-JETS DEFINED AT SCALE YMAC (SO ONLY NSUB
!C                        OF THE JETS 1 -> NJET WILL APPEAR IN JET())
!C   LAST ARGUMENT IS LABEL TO JUMP TO IF FOR ANY REASON THE EVENT
!C   COULD NOT BE PROCESSED
!C
!C   NOTE THAT ALL FLOATING POINT VARIABLES ARE DECLARED DOUBLE PRECISION
!C
      INTEGER NMAX,JET(*),NJET,NSUB,HIST,NUM,I,J,JSUB
      PARAMETER (NMAX=512)
      DOUBLE PRECISION P1(4,NMAX),P2(4,NMAX)
      DOUBLE PRECISION ECUT,YCUT,YMAC,ZERO,ETOT,RSQ,P,KTP,KTS,KT,KTLAST
      COMMON /KTCOMM/ETOT,RSQ,P(9,NMAX),KTP(NMAX,NMAX),KTS(NMAX), &
           KT(NMAX),KTLAST(NMAX),HIST(NMAX),NUM
      DIMENSION JSUB(NMAX)
!C---THE MOMENTA HAVE TO BEEN GIVEN LEGAL VALUES,
!C   EVEN THOUGH THEY WILL NEVER BE USED
      DATA ((P1(J,I),I=1,NMAX),J=1,4),ZERO &
           /NMAX*1,NMAX*0,NMAX*0,NMAX*1,0/
!C---FIRST GET A LIST OF WHICH PARTICLE IS IN WHICH JET AT YCUT
      CALL KTRECO(1,P1,NUM,ECUT,ZERO,YCUT,P2,JET,NJET,NSUB,*999)
!C---THEN FIND OUT WHICH JETS ARE SUBJETS
      CALL KTRECO(1,P1,NUM,ECUT,YCUT,YMAC,P2,JSUB,NJET,NSUB,*999)
!C---AND MODIFY JET() ACCORDINGLY
      DO I=1,NUM
         IF (JET(I).NE.0) THEN
            IF (JSUB(JET(I)).EQ.0) JET(I)=0
         ENDIF
      END DO
      RETURN
999   RETURN 1
    END SUBROUTINE KTWCHS
!C-----------------------------------------------------------------------
    SUBROUTINE KTFRAM(IOPT,CMF,SIGN,Z,XZ,N,P,Q,*)
      IMPLICIT NONE
!C---BOOST PARTICLES IN P TO/FROM FRAME GIVEN BY CMF, Z, XZ.
!C---IN THIS FRAME CMZ IS STATIONARY,
!C                   Z IS ALONG THE (SIGN)Z-AXIS (SIGN=+ OR -)
!C                  XZ IS IN THE X-Z PLANE (WITH POSITIVE X COMPONENT)
!C---IF Z HAS LENGTH ZERO, OR SIGN=0, NO ROTATION IS PERFORMED
!C---IF XZ HAS ZERO COMPONENT PERPENDICULAR TO Z IN THAT FRAME,
!C   NO AZIMUTHAL ROTATION IS PERFORMED
!C
!C   IOPT    = INPUT  : 0=TO FRAME, 1=FROM FRAME
!C   CMF(I)  = INPUT  : 4-MOMENTUM WHICH IS STATIONARY IN THE FRAME
!C   SIGN    = INPUT  : DIRECTION OF Z IN THE FRAME, NOTE THAT
!C                        ONLY ITS SIGN IS USED, NOT ITS MAGNITUDE
!C   Z(I)    = INPUT  : 4-MOMENTUM WHICH LIES ON THE (SIGN)Z-AXIS
!C   XZ(I)   = INPUT  : 4-MOMENTUM WHICH LIES IN THE X-Z PLANE
!C   N       = INPUT  : NUMBER OF PARTICLES IN P
!C   P(I,J)  = INPUT  : 4-MOMENTUM OF JTH PARTICLE BEFORE
!C   Q(I,J)  = OUTPUT : 4-MOMENTUM OF JTH PARTICLE AFTER
!C   LAST ARGUMENT IS LABEL TO JUMP TO IF FOR ANY REASON THE EVENT
!C   COULD NOT BE PROCESSED
!C
!C   NOTE THAT ALL MOMENTA ARE DOUBLE PRECISION
!C
!C   NOTE THAT IT IS SAFE TO CALL WITH P=Q
!C
      INTEGER IOPT,I,N
      DOUBLE PRECISION CMF(4),SIGN,Z(4),XZ(4),P(4,N),Q(4,N), &
           R(4,4),NEW(4),OLD(4)
      IF (IOPT.LT.0.OR.IOPT.GT.1) CALL KTWARN('KTFRAM',200,*999)
!C---FIND BOOST TO GET THERE FROM LAB
      CALL KTUNIT(R)
      CALL KTLBST(0,R,CMF,*999)
!C---FIND ROTATION TO PUT BOOSTED Z ON THE (SIGN)Z AXIS
      IF (SIGN.NE.0) THEN
         CALL KTVMUL(R,Z,OLD)
         IF (OLD(1).NE.0.OR.OLD(2).NE.0.OR.OLD(3).NE.0) THEN
            NEW(1)=0
            NEW(2)=0
            NEW(3)=SIGN
            NEW(4)=ABS(SIGN)
            CALL KTRROT(R,OLD,NEW,*999)
!C---FIND ROTATION TO PUT BOOSTED AND ROTATED XZ INTO X-Z PLANE
            CALL KTVMUL(R,XZ,OLD)
            IF (OLD(1).NE.0.OR.OLD(2).NE.0) THEN
               NEW(1)=1
               NEW(2)=0
               NEW(3)=0
               NEW(4)=1
               OLD(3)=0
!C---NOTE THAT A POTENTIALLY AWKWARD SPECIAL CASE IS AVERTED, BECAUSE IF
!C   OLD AND NEW ARE EXACTLY BACK-TO-BACK, THE ROTATION AXIS IS UNDEFINED
!C   BUT IN THAT CASE KTRROT WILL USE THE Z AXIS, AS REQUIRED
               CALL KTRROT(R,OLD,NEW,*999)
            ENDIF
         ENDIF
      ENDIF
!C---INVERT THE TRANSFORMATION IF NECESSARY
      IF (IOPT.EQ.1) CALL KTINVT(R,R)
!C---APPLY THE RESULT TO ALL THE VECTORS
      DO I=1,N
         CALL KTVMUL(R,P(1,I),Q(1,I))
      END DO
      RETURN
999   RETURN 1
    END SUBROUTINE KTFRAM
!C-----------------------------------------------------------------------
    SUBROUTINE KTBREI(IOPT,PLEP,PHAD,POUT,N,P,Q,*)
      IMPLICIT NONE
!C---BOOST PARTICLES IN P TO/FROM BREIT FRAME
!C
!C   IOPT    = INPUT  : 0/2=TO BREIT FRAME, 1/3=FROM BREIT FRAME
!C                      0/1=NO AZIMUTHAL ROTATION AFTERWARDS
!C                      2/3=LEPTON PLANE ROTATED INTO THE X-Z PLANE
!C   PLEP    = INPUT  : MOMENTUM OF INCOMING LEPTON IN +Z DIRECTION
!C   PHAD    = INPUT  : MOMENTUM OF INCOMING HADRON IN +Z DIRECTION
!C   POUT(I) = INPUT  : 4-MOMENTUM OF OUTGOING LEPTON
!C   N       = INPUT  : NUMBER OF PARTICLES IN P
!C   P(I,J)  = INPUT  : 4-MOMENTUM OF JTH PARTICLE BEFORE
!C   Q(I,J)  = OUTPUT : 4-MOMENTUM OF JTH PARTICLE AFTER
!C   LAST ARGUMENT IS LABEL TO JUMP TO IF FOR ANY REASON THE EVENT
!C   COULD NOT BE PROCESSED (MOST LIKELY DUE TO PARTICLES HAVING SMALLER
!C   ENERGY THAN MOMENTUM)
!C
!C   NOTE THAT ALL MOMENTA ARE DOUBLE PRECISION
!C
!C   NOTE THAT IT IS SAFE TO CALL WITH P=Q
!C
      INTEGER IOPT,N
      DOUBLE PRECISION PLEP,PHAD,POUT(4),P(4,N),Q(4,N), &
           CMF(4),Z(4),XZ(4),DOT,QDQ
!C---CHECK INPUT
      IF (IOPT.LT.0.OR.IOPT.GT.3) CALL KTWARN('KTBREI',200,*999)
!C---FIND 4-MOMENTUM OF BREIT FRAME (TIMES AN ARBITRARY FACTOR)
      DOT=ABS(PHAD)*(ABS(PLEP)-POUT(4))-PHAD*(PLEP-POUT(3))
      QDQ=(ABS(PLEP)-POUT(4))**2-(PLEP-POUT(3))**2-POUT(2)**2-POUT(1)**2
      CMF(1)=DOT*(         -POUT(1))
      CMF(2)=DOT*(         -POUT(2))
      CMF(3)=DOT*(    PLEP -POUT(3))-QDQ*    PHAD
      CMF(4)=DOT*(ABS(PLEP)-POUT(4))-QDQ*ABS(PHAD)
!C---FIND ROTATION TO PUT INCOMING HADRON BACK ON Z-AXIS
      Z(1)=0
      Z(2)=0
      Z(3)=PHAD
      Z(4)=ABS(PHAD)
      XZ(1)=0
      XZ(2)=0
      XZ(3)=0
      XZ(4)=0
!C---DO THE BOOST
      IF (IOPT.LE.1) THEN
         CALL KTFRAM(IOPT,CMF,PHAD,Z,XZ,N,P,Q,*999)
      ELSE
         CALL KTFRAM(IOPT-2,CMF,PHAD,Z,POUT,N,P,Q,*999)
      ENDIF
      RETURN
999   RETURN 1
    END SUBROUTINE KTBREI
!C-----------------------------------------------------------------------
    SUBROUTINE KTHADR(IOPT,PLEP,PHAD,POUT,N,P,Q,*)
      IMPLICIT NONE
!C---BOOST PARTICLES IN P TO/FROM HADRONIC CMF
!C
!C   ARGUMENTS ARE EXACTLY AS FOR KTBREI
!C
!C   NOTE THAT ALL MOMENTA ARE DOUBLE PRECISION
!C
!C   NOTE THAT IT IS SAFE TO CALL WITH P=Q
!C
      INTEGER IOPT,N
      DOUBLE PRECISION PLEP,PHAD,POUT(4),P(4,N),Q(4,N), &
           CMF(4),Z(4),XZ(4)
!C---CHECK INPUT
      IF (IOPT.LT.0.OR.IOPT.GT.3) CALL KTWARN('KTHADR',200,*999)
!C---FIND 4-MOMENTUM OF HADRONIC CMF
      CMF(1)=         -POUT(1)
      CMF(2)=         -POUT(2)
      CMF(3)=    PLEP -POUT(3)+    PHAD
      CMF(4)=ABS(PLEP)-POUT(4)+ABS(PHAD)
!C---FIND ROTATION TO PUT INCOMING HADRON BACK ON Z-AXIS
      Z(1)=0
      Z(2)=0
      Z(3)=PHAD
      Z(4)=ABS(PHAD)
      XZ(1)=0
      XZ(2)=0
      XZ(3)=0
      XZ(4)=0
!C---DO THE BOOST
      IF (IOPT.LE.1) THEN
         CALL KTFRAM(IOPT,CMF,PHAD,Z,XZ,N,P,Q,*999)
      ELSE
         CALL KTFRAM(IOPT-2,CMF,PHAD,Z,POUT,N,P,Q,*999)
      ENDIF
      RETURN
999   RETURN 1
    END SUBROUTINE KTHADR
\nwendcode{}\nwbegincode{441}\moddef{MLM matching: procedures (not used)}\plusendmoddef
    SUBROUTINE KTUNIT(R)
      IMPLICIT NONE
!C   SET R EQUAL TO THE 4 BY 4 IDENTITY MATRIX
      DOUBLE PRECISION R(4,4)
      INTEGER I,J
      DO I=1,4
         DO J=1,4
            R(I,J)=0
            IF (I.EQ.J) R(I,J)=1
         END DO
      END DO
    END SUBROUTINE KTUNIT
!C-----------------------------------------------------------------------
    SUBROUTINE KTLBST(IOPT,R,A,*)
      IMPLICIT NONE
!C   PREMULTIPLY R BY THE 4 BY 4 MATRIX TO
!C   LORENTZ BOOST TO/FROM THE CM FRAME OF A
!C   IOPT=0 => TO
!C   IOPT=1 => FROM
!C
!C   LAST ARGUMENT IS LABEL TO JUMP TO IF A IS NOT TIME-LIKE
!C
      INTEGER IOPT,I,J
      DOUBLE PRECISION R(4,4),A(4),B(4),C(4,4),M
      DO I=1,4
         B(I)=A(I)
      END DO
      M=B(4)**2-B(1)**2-B(2)**2-B(3)**2
      IF (M.LE.0) CALL KTWARN('KTLBST',100,*999)
      M=SQRT(M)
      B(4)=B(4)+M
      M=1/(M*B(4))
      IF (IOPT.EQ.0) THEN
        B(4)=-B(4)
      ELSEIF (IOPT.NE.1) THEN
        CALL KTWARN('KTLBST',200,*999)
        STOP
      ENDIF
      DO I=1,4
         DO J=1,4
            C(I,J)=B(I)*B(J)*M
            IF (I.EQ.J) C(I,J)=C(I,J)+1
         END DO
      END DO
      C(4,4)=C(4,4)-2
      CALL KTMMUL(C,R,R)
      RETURN
999   RETURN 1
    END SUBROUTINE KTLBST
!C-----------------------------------------------------------------------
    SUBROUTINE KTRROT(R,A,B,*)
      IMPLICIT NONE
!C   PREMULTIPLY R BY THE 4 BY 4 MATRIX TO
!C   ROTATE FROM VECTOR A TO VECTOR B BY THE SHORTEST ROUTE
!C   IF THEY ARE EXACTLY BACK-TO-BACK, THE ROTATION AXIS IS THE VECTOR
!C   WHICH IS PERPENDICULAR TO THEM AND THE X AXIS, UNLESS THEY ARE
!C   PERPENDICULAR TO THE Y AXIS, WHEN IT IS THE VECTOR WHICH IS
!C   PERPENDICULAR TO THEM AND THE Y AXIS.
!C   NOTE THAT THESE CONDITIONS GUARANTEE THAT IF BOTH ARE PERPENDICULAR
!C   TO THE Z AXIS, IT WILL BE USED AS THE ROTATION AXIS.
!C
!C   LAST ARGUMENT IS LABEL TO JUMP TO IF EITHER HAS LENGTH ZERO
!C
      DOUBLE PRECISION R(4,4),M(4,4),A(4),B(4),C(4),D(4),AL,BL,CL,DL,EPS
!C---SQRT(2*EPS) IS THE ANGLE IN RADIANS OF THE SMALLEST ALLOWED ROTATION
!C   NOTE THAT IF YOU CONVERT THIS PROGRAM TO SINGLE PRECISION, YOU WILL
!C   NEED TO INCREASE EPS TO AROUND 0.5E-4
      PARAMETER (EPS=0.5D-6)
      AL=A(1)**2+A(2)**2+A(3)**2
      BL=B(1)**2+B(2)**2+B(3)**2
      IF (AL.LE.0.OR.BL.LE.0) CALL KTWARN('KTRROT',100,*999)
      AL=1/SQRT(AL)
      BL=1/SQRT(BL)
      CL=(A(1)*B(1)+A(2)*B(2)+A(3)*B(3))*AL*BL
!C---IF THEY ARE COLLINEAR, DON'T NEED TO DO ANYTHING
      IF (CL.GE.1-EPS) THEN
         RETURN
!C---IF THEY ARE BACK-TO-BACK, USE THE AXIS PERP TO THEM AND X AXIS
      ELSEIF (CL.LE.-1+EPS) THEN
         IF (ABS(B(2)).GT.EPS) THEN
            C(1)= 0
            C(2)=-B(3)
            C(3)= B(2)
!C---UNLESS THEY ARE PERPENDICULAR TO THE Y AXIS,
         ELSE
            C(1)= B(3)
            C(2)= 0
            C(3)=-B(1)
         ENDIF
!C---OTHERWISE FIND ROTATION AXIS
      ELSE
         C(1)=A(2)*B(3)-A(3)*B(2)
         C(2)=A(3)*B(1)-A(1)*B(3)
         C(3)=A(1)*B(2)-A(2)*B(1)
      ENDIF
      CL=C(1)**2+C(2)**2+C(3)**2
      IF (CL.LE.0) CALL KTWARN('KTRROT',101,*999)
      CL=1/SQRT(CL)
!C---FIND ROTATION TO INTERMEDIATE AXES FROM A
      D(1)=A(2)*C(3)-A(3)*C(2)
      D(2)=A(3)*C(1)-A(1)*C(3)
      D(3)=A(1)*C(2)-A(2)*C(1)
      DL=AL*CL
      M(1,1)=A(1)*AL
      M(1,2)=A(2)*AL
      M(1,3)=A(3)*AL
      M(1,4)=0
      M(2,1)=C(1)*CL
      M(2,2)=C(2)*CL
      M(2,3)=C(3)*CL
      M(2,4)=0
      M(3,1)=D(1)*DL
      M(3,2)=D(2)*DL
      M(3,3)=D(3)*DL
      M(3,4)=0
      M(4,1)=0
      M(4,2)=0
      M(4,3)=0
      M(4,4)=1
      CALL KTMMUL(M,R,R)
!C---AND ROTATION FROM INTERMEDIATE AXES TO B
      D(1)=B(2)*C(3)-B(3)*C(2)
      D(2)=B(3)*C(1)-B(1)*C(3)
      D(3)=B(1)*C(2)-B(2)*C(1)
      DL=BL*CL
      M(1,1)=B(1)*BL
      M(2,1)=B(2)*BL
      M(3,1)=B(3)*BL
      M(1,2)=C(1)*CL
      M(2,2)=C(2)*CL
      M(3,2)=C(3)*CL
      M(1,3)=D(1)*DL
      M(2,3)=D(2)*DL
      M(3,3)=D(3)*DL
      CALL KTMMUL(M,R,R)
      RETURN
 999  RETURN 1
    END SUBROUTINE KTRROT
!C-----------------------------------------------------------------------
    SUBROUTINE KTVMUL(M,A,B)
      IMPLICIT NONE
!C   4 BY 4 MATRIX TIMES 4 VECTOR: B=M*A.
!C   ALL ARE DOUBLE PRECISION
!C   IT IS SAFE TO CALL WITH B=A
!C   FIRST SUBSCRIPT=ROWS, SECOND=COLUMNS
      DOUBLE PRECISION M(4,4),A(4),B(4),C(4)
      INTEGER I,J
      DO I=1,4
         C(I)=0
         DO J=1,4
            C(I)=C(I)+M(I,J)*A(J)
         END DO
      END DO
      DO I=1,4
         B(I)=C(I)
      END DO
    END SUBROUTINE KTVMUL
!C-----------------------------------------------------------------------
    SUBROUTINE KTMMUL(A,B,C)
      IMPLICIT NONE
!C   4 BY 4 MATRIX MULTIPLICATION: C=A*B.
!C   ALL ARE DOUBLE PRECISION
!C   IT IS SAFE TO CALL WITH C=A OR B.
!C   FIRST SUBSCRIPT=ROWS, SECOND=COLUMNS
      DOUBLE PRECISION A(4,4),B(4,4),C(4,4),D(4,4)
      INTEGER I,J,K
      DO I=1,4
         DO J=1,4
            D(I,J)=0
            DO K=1,4
               D(I,J)=D(I,J)+A(I,K)*B(K,J)
            END DO
         END DO
      END DO
      DO I=1,4
         DO J=1,4
            C(I,J)=D(I,J)
         END DO
      END DO
    END SUBROUTINE KTMMUL
!C-----------------------------------------------------------------------
    SUBROUTINE KTINVT(A,B)
      IMPLICIT NONE
!C---INVERT TRANSFORMATION MATRIX A
!C
!C   A = INPUT  : 4 BY 4 TRANSFORMATION MATRIX
!C   B = OUTPUT : INVERTED TRANSFORMATION MATRIX
!C
!C   IF A IS NOT A TRANSFORMATION MATRIX YOU WILL GET STRANGE RESULTS
!C
!C   NOTE THAT IT IS SAFE TO CALL WITH A=B
!C
      DOUBLE PRECISION A(4,4),B(4,4),C(4,4)
      INTEGER I,J
!C---TRANSPOSE
      DO I=1,4
         DO J=1,4
            C(I,J)=A(J,I)
         END DO
      END DO
!C---NEGATE ENERGY-MOMENTUM MIXING TERMS
      DO I=1,3
         C(4,I)=-C(4,I)
         C(I,4)=-C(I,4)
      END DO
!C---OUTPUT
      DO I=1,4
         DO J=1,4
            B(I,J)=C(I,J)
         END DO
      END DO
    END SUBROUTINE KTINVT
!C-----------------------------------------------------------------------
\nwendcode{}\nwbegincode{442}\moddef{Shower2pythia: procedures (not used)}\plusendmoddef
  recursive subroutine shower_topythia_recursive (prt, mode,first)
    IMPLICIT DOUBLE PRECISION(A-H, O-Z)
    IMPLICIT INTEGER(I-N)
    !    C...  Commonblocks.
    COMMON/PYJETS/N,NPAD,K(4000,5),P(4000,5),V(4000,5)
    COMMON/PYDAT1/MSTU(200),PARU(200),MSTJ(200),PARJ(200)
    COMMON/PYDAT2/KCHG(500,4),PMAS(500,4),PARF(2000),VCKM(4,4)
    COMMON/PYDAT3/MDCY(500,3),MDME(8000,2),BRAT(8000),KFDP(8000,5)
    COMMON/PYSUBS/MSEL,MSELPD,MSUB(500),KFIN(2,-40:40),CKIN(200)
    COMMON/PYPARS/MSTP(200),PARP(200),MSTI(200),PARI(200)
    SAVE /PYJETS/,/PYDAT1/,/PYDAT2/,/PYDAT3/,/PYSUBS/,/PYPARS/

    type(parton_t), intent(in) :: prt
    integer, intent(in) :: mode, first

    if (parton_is_final (prt)) then
       N = N + 1
       K(N,1) = 2
       if (parton_is_quark(prt)) then
          !!! check if quark is end of a color connection
          if (prt%type < 0) then
             K(N,1) = 1
          end if
       end if
       K(N,2) = prt%type
       K(N,3) = first
       K(N,4) = 0
       K(N,5) = 0
       P(N,1) = prt%momentum%p(1)
       P(N,2) = prt%momentum%p(2)
       P(N,3) = prt%momentum%p(3)
       P(N,4) = prt%momentum%p(0)
       P(N,5) = prt%t
    else
       !!! search for following final partons
       if (parton_is_gluon (prt)) then
          if (parton_is_gluon (prt%child1)) then
             !!! g-> gg so sequence is unimportant
             call shower_topythia_recursive (prt%child1, 1, first)
             call shower_topythia_recursive (prt%child2, 1, first)
          else
             !!! g-> qqbar -> use antiquark first, so that color flow
             !!!     is given correctly
             if (prt%child1%type < 0) then
                call shower_topythia_recursive (prt%child1, 1, first)
                call shower_topythia_recursive (prt%child2, 2, first)
             else
                call shower_topythia_recursive (prt%child2, 2, first)
                call shower_topythia_recursive (prt%child1, 1, first)
             end if
          endif
       else
          if (mode == 1) then
             call shower_topythia_recursive (prt%child1, 1, first)
             call shower_topythia_recursive (prt%child2, 1, first)
          else
             call shower_topythia_recursive (prt%child2, 2, first)
             call shower_topythia_recursive (prt%child1, 2, first)
          end if
       end if
    end if
  end subroutine shower_topythia_recursive

\eatline
\nwendcode{}\nwbegindocs{443}\nwdocspar
\nwenddocs{}\nwfilename{../../../src/muli/muli.nw}\nwbegindocs{0}% -*- ess-noweb-default-code-mode: f90-mode; noweb-default-code-mode: f90-mode; -*-
% WHIZARD multiple interactions code as NOWEB source
%%%%%%%%%%%%%%%%%%%%%%%%%%%%%%%%%%%%%%%%%%%%%%%%%%%%%%%%%%%%%%%%%%%%%%%%
\chapter{Multiple Interactions (MPI) Code}

This is the code for the \whizard\ module for multiple interactions (MPI)
in hadron collisions. It also provides the interleaved shower with together
with the shower module.

%%%%%%%%%%%%%%%%%%%%%%%%%%%%%%%%%%%%%%%%%%%%%%%%%%%%%%%%%%%%%%%%%%%%%%%%
\section{Basic types}

This file contains the module {\tt{}muli{\char95}basic} which is included by all
other modules of the MPI code. It's main purpose is serialization and
deserialization, but it also contains some bitmodel parameters and
some minor definitions common to all  modules and types.

Serialization is implemented in three layers:

\begin{enumerate}
\item {\bf I/O layer}: streamfile access and exact, retreivable and
  compiler independent representation of all intrinsic types
\item {\bf Marking layer}: storing/restoring pointer association and
    better human readability
\item {\bf Derived type layer}: abstract type as parent for all
    serializable derived types
\end{enumerate}

\paragraph{The I/O Layer}

The I/O layer is implemented by {\tt{}page{\char95}ring{\char95}t} and its type bound 
procedures. We chose not to use the standard formatted sequential I/O for
several reasons:
\begin{itemize}
\item Sequential I/O is line orientated, but serialization is not.
\item Formatted I/O of floating point numbers is inexact. There were problems
  in reading numbers like {\tt{}1+epsilon(1)} or {\tt{}huge(1)}
  including arithmetic exeptions and wrong results.
\item Formatted reading is slow. This does matter, when you read files
  of $\mathcal{O}(100)$ MB.
\end{itemize}

{\tt{}page{\char95}ring{\char95}t} is a ring buffer of default kind characters holding
a region of the addressed file. All read and write procedures use the
ring buffer, not the file. {\tt{}page{\char95}ring} automatically reads more 
characters from the file when needed and writes and discards parts of
the ring buffer when not any longer needed.

\paragraph{Marking layer}

Marking is done by {\tt{}marking{\char95}t} and its type bound procedures. Files 
written by {\tt{}marking{\char95}t} are regular XML files, so you can use your
favorite XML editor to examine or edit serialized contents. The main
purpose of this layer is to recover the type of polymorphic entities
and to restore association of pointers, but it also assigns names to
all contents, so it is much easier to read generated files.

Type recovery is done by a reference list. You must push a representive 
object to this reference list for each type that might get resolved before 
starting deserialization. {\tt{}marker{\char95}t} will care for comparing these
representives to the contents of the files.

Association restoration is done by a heap list. It is technically equal
to the reference list, but holds all targets that have been processed so
far. Generation of this list is handled by {\tt{}marking{\char95}t}, so you
dont't have to care for this aspect at all. Up to the present it is
not possible to restore association, when a non-pointer target is
serialized after its associated pointer is serialized. 
There is no trivial solution and this case does not appear here,
so we will not take care of this.

\paragraph{Derived type layer}

Each instance that shall become serialized must extend
{\tt{}ser{\char95}class{\char95}t}. Essential for type recovery is the virtual
type bound procedure {\tt{}get{\char95}type}. Each non-virtual type shall override
{\tt{}get{\char95}type} and return the actual name of its type in lower-case letters.
Each type which adds new, non-redundant components shall override
{\tt{}write{\char95}to{\char95}marker} and {\tt{}read{\char95}from{\char95}marker}. These type-bound
procedures define, what contents get serialized. While the marker
cares about tagging the type and association of the instance, every
instance still has to define what to serialize. The rule is to mark
the begin of its contents, then its parents procedure, then mark all
non-redundant components, then mark the end of its contents.
Finally, each serializable type shall override {\tt{}print{\char95}to{\char95}unit}. This
procedure is called for an arbitrary human-readable output. It is
quite similar to {\tt{}write{\char95}to{\char95}ring}, but without strict format and
ignoring machine-readability. 

{\tt{}ser{\char95}class{\char95}t} has strictly speaking two layers. {\tt{}write{\char95}to{\char95}marker} and
{\tt{}read{\char95}from{\char95}marker} are only for internal usage. Serialization and
deserialization are triggered by the TBPs serialize and deserialize. These
procedures take care of initialization and finalization of the marker. A
serializable type should override these procedures to push a representive
of itself and any other references to the reference list of it's marker
before (de)serialization and to pop them from the list afterwards.

\nwenddocs{}\nwbegincode{1}\moddef{\code{}muli{\char95}base.f90\edoc{}}\endmoddef
\LA{}File header\RA{}

module muli_base
  use, intrinsic :: iso_fortran_env
\LA{}Use kinds with double\RA{}
  use kinds, only: i64
\LA{}Use strings\RA{}
  use constants
  use io_units
  use diagnostics

\LA{}Standard module head\RA{}

\LA{}Muli base: variables\RA{}

\LA{}Muli base: public\RA{}

\LA{}Muli base: types\RA{}

\LA{}Muli base: interfaces\RA{}

contains

\LA{}Muli base: procedures\RA{}  

end module muli_base

\eatline
\nwendcode{}\nwbegindocs{2}These are the bitmodel parameters.
\nwenddocs{}\nwbegincode{3}\moddef{Muli base: variables}\endmoddef
  integer, public, parameter :: dik = i64
  integer(dik), public, parameter :: i_one = int(1, kind=dik)
  integer(dik), public, parameter :: i_zero = int(0, kind=dik)
\eatline
\nwendcode{}\nwbegindocs{4}These are the serialization parameters.
\nwenddocs{}\nwbegincode{5}\moddef{Muli base: variables}\plusendmoddef
  integer(dik), public, parameter :: serialize_page_size = 1024
  integer(dik), public, parameter :: serialize_ok = 0000
  integer(dik), public, parameter :: serialize_syntax_error = 1001
  integer(dik), public, parameter :: serialize_wrong_tag = 1002
  integer(dik), public, parameter :: serialize_wrong_id = 1003
  integer(dik), public, parameter :: serialize_wrong_type = 1004
  integer(dik), public, parameter :: serialize_wrong_name = 1005
  integer(dik), public, parameter :: serialize_no_target = 1006
  integer(dik), public, parameter :: serialize_no_pointer = 1007
  integer(dik), public, parameter :: serialize_wrong_action = 1008
  integer(dik), public, parameter :: serialize_unexpected_content = 1009
  integer(dik), public, parameter :: serialize_null = 1010
  integer(dik), public, parameter :: serialize_nothing = 1011
  logical, public, parameter :: serialize_default_indent = .true.
  logical, public, parameter :: serialize_default_line_break = .true.
  logical, public, parameter :: serialize_default_asynchronous = .false.
\eatline
\eatline
\eatline
\eatline
\eatline
\eatline
\eatline
\eatline
\nwendcode{}\nwbegindocs{6}\nwdocspar
And some private variables:  
\nwenddocs{}\nwbegincode{7}\moddef{Muli base: variables}\plusendmoddef
  integer(dik) :: last_id = 0
  character(len=*), parameter :: serialize_integer_characters = &
       "-0123456789"

\eatline
\nwendcode{}\nwbegincode{8}\moddef{Muli base: interfaces}\endmoddef
  abstract interface
     subroutine ser_write_if (this, marker, status)
       import ser_class_t
       import marker_t
       import dik
       class(ser_class_t), intent(in) :: this
       class(marker_t), intent(inout) :: marker
       integer(dik), intent(out) :: status
     end subroutine ser_write_if
  end interface
  
  abstract interface
     subroutine ser_read_if (this, marker, status)
       import ser_class_t
       import marker_t
       import dik
       class(ser_class_t), intent(out) :: this
       class(marker_t), intent(inout) :: marker
       integer(dik), intent(out) :: status
     end subroutine ser_read_if
  end interface
  
  abstract interface
     subroutine ser_unit (this, unit, parents, components, peers)
       import ser_class_t
       import dik
       class(ser_class_t), intent(in) :: this
       integer,intent(in) :: unit
       integer(dik), intent(in) :: parents,components,peers
     end subroutine ser_unit
  end interface
  
  abstract interface
     pure subroutine ser_type (type)
       character(:), allocatable, intent(out) :: type
     end subroutine ser_type
  end interface  

\eatline
\nwendcode{}\nwbegindocs{9}\nwdocspar
\nwenddocs{}\nwbegincode{10}\moddef{Muli base: interfaces}\plusendmoddef
  abstract interface
     elemental function measure_int (this)
       import 
       class(measure_class_t), intent(in) :: this
       real(default) :: measure_int
     end function measure_int
  end interface  
  
\eatline
\nwendcode{}\nwbegindocs{11}\nwdocspar
\nwenddocs{}\nwbegincode{12}\moddef{Muli base: public}\endmoddef
  public :: operator(<)
\nwendcode{}\nwbegincode{13}\moddef{Muli base: interfaces}\plusendmoddef
  interface operator(<)
     module procedure measurable_less_measurable
     module procedure measurable_less_default
  end interface operator(<)
\nwendcode{}\nwbegincode{14}\moddef{Muli base: public}\plusendmoddef
  public :: operator(<=)
\nwendcode{}\nwbegincode{15}\moddef{Muli base: interfaces}\plusendmoddef
  interface operator(<=)
     module procedure measurable_less_or_equal_measurable
     module procedure measurable_less_or_equal_default
  end interface operator(<=)
\nwendcode{}\nwbegincode{16}\moddef{Muli base: public}\plusendmoddef
  public :: operator(==)
\nwendcode{}\nwbegincode{17}\moddef{Muli base: interfaces}\plusendmoddef
  interface operator(==)
     module procedure measurable_equal_measurable
     module procedure measurable_equal_default
  end interface operator(==)
\nwendcode{}\nwbegincode{18}\moddef{Muli base: public}\plusendmoddef
  public :: operator(>=)
\nwendcode{}\nwbegincode{19}\moddef{Muli base: interfaces}\plusendmoddef
  interface operator(>=)
     module procedure measurable_equal_or_greater_measurable
     module procedure measurable_equal_or_greater_default
  end interface operator(>=)
\nwendcode{}\nwbegincode{20}\moddef{Muli base: public}\plusendmoddef
  public :: operator(>)
\nwendcode{}\nwbegincode{21}\moddef{Muli base: interfaces}\plusendmoddef
  interface operator(>)
     module procedure measurable_greater_measurable
     module procedure measurable_greater_default
  end interface operator(>)

\eatline
\nwendcode{}\nwbegindocs{22}\nwdocspar
\nwenddocs{}\nwbegincode{23}\moddef{Muli base: interfaces}\plusendmoddef
  interface page_ring_position_is_before
     module procedure page_ring_position_is_before_int_pos
     module procedure page_ring_position_is_before_pos_pos
     module procedure page_ring_position_is_before_pos_int
  end interface 
  
\nwendcode{}\nwbegindocs{24}\nwdocspar
\nwenddocs{}\nwbegincode{25}\moddef{Muli base: public}\plusendmoddef
  public :: ser_class_t
\nwendcode{}\nwbegincode{26}\moddef{Muli base: types}\endmoddef
  type, abstract :: ser_class_t
   contains
  \LA{}Muli base: serial class: TBP\RA{}
  end type ser_class_t
  
\eatline
\nwendcode{}\nwbegincode{27}\moddef{Muli base: serial class: TBP}\endmoddef
  procedure(ser_write_if), deferred :: write_to_marker
\nwendcode{}\nwbegindocs{28}\nwdocspar
This is a dummy procedure. Usually, you do not need to deserialize
targets, so by implementing this dummy we don't force all descendants
to override this procedure. Then again this is the only way to read
targets from markers. 
\nwenddocs{}\nwbegincode{29}\moddef{Muli base: serial class: TBP}\plusendmoddef
  procedure(ser_read_if), deferred :: read_from_marker  
\nwendcode{}\nwbegincode{30}\moddef{Muli base: procedures}\endmoddef
  subroutine serializable_read_target_from_marker (this, marker, status)
    class(ser_class_t), target, intent(out) :: this
    class(marker_t), intent(inout) :: marker
    integer(dik), intent(out) :: status
    write (output_unit, "(A)") "serializable_read_target_from_marker:"
    write (output_unit, "(A)") "This is a dummy procedure. Usually, this " &
         // "message indicates a missing overridden " &
         // "read_target_from_marker TPB for "
    call this%write_type (output_unit)
    write (output_unit, "(A)") ""
    call this%read_from_marker (marker, status)
  end subroutine serializable_read_target_from_marker
  
\eatline
\nwendcode{}\nwbegindocs{31}\nwdocspar
\nwenddocs{}\nwbegincode{32}\moddef{Muli base: serial class: TBP}\plusendmoddef
  procedure(ser_unit), deferred :: print_to_unit
\nwendcode{}\nwbegincode{33}\moddef{Muli base: procedures}\plusendmoddef
  subroutine serializable_serialize_to_unit (this, unit, name)
    class(ser_class_t), intent(in) :: this
    integer, intent(in) :: unit
    character (len=*), intent(in) :: name
    logical :: opened
    character(32) :: file
    !!!    gfortran bug
    !!!    character::stream
    character::write
    type(marker_t)::marker
    !    inquire(unit=unit,opened=opened,stream=stream,write=write)
    inquire (unit=unit,opened=opened,write=write)
    if (opened) then
       !!! if(stream=="Y")then
          if(write=="Y")then
             print *,"dummy: serializable_serialize_to_unit"
             stop
          else
             print *,"serializable_serialize_to_unit: cannot write to read-only unit."
          end if
          !!!       else
          !!!          print *,"serializable_serialize_to_unit: access kind of unit is not 'stream'."
          !!!       end if
    else
       call msg_error ("serializable_serialize_to_unit: file is not opened.")
    end if
  end subroutine serializable_serialize_to_unit
  
\eatline
\nwendcode{}\nwbegindocs{34}\nwdocspar
\nwenddocs{}\nwbegincode{35}\moddef{Muli base: serial class: TBP}\plusendmoddef
  procedure(ser_type), nopass, deferred :: get_type  
\nwendcode{}\nwbegindocs{36}\nwdocspar
\nwenddocs{}\nwbegincode{37}\moddef{Muli base: serial class: TBP}\plusendmoddef
  procedure, nopass :: verify_type => serializable_verify_type  
\nwendcode{}\nwbegincode{38}\moddef{Muli base: procedures}\plusendmoddef
  elemental function serializable_verify_type (type) result (match)
    character(*), intent(in) :: type
    logical :: match
    match = type == "ser_class_t"
  end function serializable_verify_type
  
\eatline
\nwendcode{}\nwbegindocs{39}\nwdocspar
\nwenddocs{}\nwbegincode{40}\moddef{Muli base: serial class: TBP}\plusendmoddef
  procedure :: read_target_from_marker => &
       serializable_read_target_from_marker
\nwendcode{}\nwbegindocs{41}\nwdocspar
\nwenddocs{}\nwbegincode{42}\moddef{Muli base: serial class: TBP}\plusendmoddef
  procedure :: write_type => serializable_write_type
\nwendcode{}\nwbegincode{43}\moddef{Muli base: procedures}\plusendmoddef
  subroutine serializable_write_type (this, unit)
    class(ser_class_t), intent(in) :: this
    integer,intent(in) :: unit
    character(:), allocatable :: this_type
    call this%get_type (this_type)
    write (unit, "(A)", advance="no") this_type
  end subroutine serializable_write_type
  
\eatline
\nwendcode{}\nwbegindocs{44}\nwdocspar
\nwenddocs{}\nwbegincode{45}\moddef{Muli base: serial class: TBP}\plusendmoddef
  procedure :: print => serializable_print
\nwendcode{}\nwbegincode{46}\moddef{Muli base: procedures}\plusendmoddef
  recursive subroutine serializable_print &
       (this, parents, components, peers, unit)
    class(ser_class_t), intent(in) :: this
    integer(dik), intent(in) :: parents, components, peers
    integer, intent(in), optional :: unit
    integer :: u
    u = given_output_unit (unit)
    write (u, "(A)")
    write (u, "(A)", advance="no")  "Instance of type: "
    call this%write_type (u)
    write (u, "(A)")
    call this%print_to_unit (u, parents, components, peers)
  end subroutine serializable_print
  
\eatline
\nwendcode{}\nwbegindocs{47}\nwdocspar
\nwenddocs{}\nwbegincode{48}\moddef{Muli base: serial class: TBP}\plusendmoddef
  procedure :: print_error => serializable_print_error
\nwendcode{}\nwbegincode{49}\moddef{Muli base: procedures}\plusendmoddef
  recursive subroutine serializable_print_error (this)
    class(ser_class_t), intent(in) :: this
    call this%print_to_unit (error_unit, i_zero, i_zero, i_zero)
  end subroutine serializable_print_error
 
\eatline
\nwendcode{}\nwbegindocs{50}\nwdocspar
\nwenddocs{}\nwbegincode{51}\moddef{Muli base: serial class: TBP}\plusendmoddef
  procedure :: print_all => serializable_print_all
\nwendcode{}\nwbegincode{52}\moddef{Muli base: procedures}\plusendmoddef
  recursive subroutine serializable_print_all (this, unit)
    class(ser_class_t), intent(in) :: this
    integer, intent(in), optional :: unit
    integer :: u
    u = given_output_unit (unit)
    write (u, "(A)")
    write (u, "(A)", advance="no")  "Instance of type: "
    call this%write_type (u)
    write (u, "(A)")
    call this%print_to_unit (u, huge(i_one), huge(i_one), huge(i_one))
  end subroutine serializable_print_all
  
\eatline
\nwendcode{}\nwbegindocs{53}\nwdocspar
\nwenddocs{}\nwbegincode{54}\moddef{Muli base: serial class: TBP}\plusendmoddef
  procedure :: print_little => serializable_print_little
\nwendcode{}\nwbegincode{55}\moddef{Muli base: procedures}\plusendmoddef
  recursive subroutine serializable_print_little (this, unit)
    class(ser_class_t), intent(in) :: this
    integer, intent(in), optional :: unit
    integer :: u
    u = given_output_unit (u)
    write(u, "(A)")
    write(u, "(A)", advance="no")  "Instance of type: "
    call this%write_type (u)
    write(u, "(A)")
    call this%print_to_unit (u, i_zero, i_zero, i_zero)
  end subroutine serializable_print_little
  
\eatline
\nwendcode{}\nwbegindocs{56}\nwdocspar
\nwenddocs{}\nwbegincode{57}\moddef{Muli base: serial class: TBP}\plusendmoddef
  procedure :: print_parents => serializable_print_parents
\nwendcode{}\nwbegincode{58}\moddef{Muli base: procedures}\plusendmoddef
  recursive subroutine serializable_print_parents (this)
    class(ser_class_t), intent(in) :: this
    write(output_unit, "(A)")
    write(output_unit, "(A)", advance="no")  "Instance of type: "
    call this%write_type (output_unit)
    write (output_unit, "(A)")
    call this%print_to_unit (output_unit, huge(i_one), i_zero, i_zero)
  end subroutine serializable_print_parents
  
\eatline
\nwendcode{}\nwbegindocs{59}\nwdocspar
\nwenddocs{}\nwbegincode{60}\moddef{Muli base: serial class: TBP}\plusendmoddef
  procedure :: print_components => serializable_print_components
\nwendcode{}\nwbegincode{61}\moddef{Muli base: procedures}\plusendmoddef
  recursive subroutine serializable_print_components(this)
    class(ser_class_t), intent(in) :: this
    write (output_unit, "(A)")
    write (output_unit, "(A)", advance="no")  "Instance of type: "
    call this%write_type (output_unit)
    write(output_unit, "(A)")
    call this%print_to_unit (output_unit, i_zero, huge(i_one), i_zero)
  end subroutine serializable_print_components
  
\eatline
\nwendcode{}\nwbegindocs{62}\nwdocspar
\nwenddocs{}\nwbegincode{63}\moddef{Muli base: serial class: TBP}\plusendmoddef
  procedure :: print_peers => serializable_print_peers
\nwendcode{}\nwbegincode{64}\moddef{Muli base: procedures}\plusendmoddef
  recursive subroutine serializable_print_peers (this)
    class(ser_class_t), intent(in) :: this
    write (output_unit, "(A)")
    write (output_unit, "(A)", advance="no")  "Instance of type: "
    call this%write_type (output_unit)
    write (output_unit, "(A)")
    call this%print_to_unit (output_unit, i_zero, i_zero, huge(i_one))
  end subroutine serializable_print_peers
  
\eatline
\nwendcode{}\nwbegindocs{65}\nwdocspar
\nwenddocs{}\nwbegincode{66}\moddef{Muli base: serial class: TBP}\plusendmoddef
  procedure :: serialize_to_file => serializable_serialize_to_file
\nwendcode{}\nwbegincode{67}\moddef{Muli base: procedures}\plusendmoddef
  subroutine serializable_serialize_to_file (this, name, file)
    class(ser_class_t), intent(in) :: this
    character(len=*), intent(in) :: file, name
    type(marker_t) :: marker
    call marker%open_for_write_access (file)
    write (output_unit, "(A,A)")  &
         "Serializable_serialize_to_file: writing xml preamble to ", file
    call marker%activate_next_page ()
    call marker%push ('<?xml version="1.0"?>')
    call marker%mark_begin (tag="file", name = file)
    flush(marker%unit)
    call this%serialize_to_marker (marker, name)
    call marker%mark_end ("file")
    call marker%close ()
    call marker%finalize ()
  end subroutine serializable_serialize_to_file
  
\eatline
\nwendcode{}\nwbegindocs{68}\nwdocspar
\nwenddocs{}\nwbegincode{69}\moddef{Muli base: serial class: TBP}\plusendmoddef
  procedure :: serialize_to_unit => serializable_serialize_to_unit
\nwendcode{}\nwbegindocs{70}\nwdocspar
\nwenddocs{}\nwbegincode{71}\moddef{Muli base: serial class: TBP}\plusendmoddef
  procedure :: serialize_to_marker => serializable_serialize_to_marker
\nwendcode{}\nwbegincode{72}\moddef{Muli base: procedures}\plusendmoddef
  recursive subroutine serializable_serialize_to_marker (this, marker, name)
    class(ser_class_t), intent(in) :: this
    class(marker_t), intent(inout) :: marker
    character(len=*), intent(in) :: name
    if (marker%action == 1) then
       call marker%mark_instance (this, name)
    else       
       call msg_error ("serializable_serialize_to_marker: Marker is " &
            // "not ready for write access.")
    end if
  end subroutine serializable_serialize_to_marker
  
\eatline
\nwendcode{}\nwbegindocs{73}\nwdocspar
\nwenddocs{}\nwbegincode{74}\moddef{Muli base: serial class: TBP}\plusendmoddef
  procedure :: deserialize_from_file => serializable_deserialize_from_file
\nwendcode{}\nwbegincode{75}\moddef{Muli base: procedures}\plusendmoddef
  subroutine serializable_deserialize_from_file (this, name, file)
    class(ser_class_t), intent(out) :: this
    character(*), intent(in) :: name, file
    type(marker_t) :: marker
    integer(dik), dimension(2) :: p1, p2
    call marker%open_for_read_access (file, "</file>")
    marker%eof_int = huge(i_one)
    marker%eof_pos = page_ring_position (marker%eof_int)
    call marker%read_page ()
    call marker%find ('<?', skip=2, proceed=.true., pos=p1)
    call marker%find ('?>', skip=3, proceed=.false., pos=p2)
    if ((p1(2) <= 0) .or. (p2(2) <= 0)) then
       call msg_error ("serializable_deserialize_from_file: no " &
            // "version substring found.")
    end if
    call marker%set_position (p2)
    call marker%find ('<file ', skip=4, proceed=.true., pos=p1)
    call marker%find ('>', skip=1, proceed=.false., pos=p2)
    if((p1(2)>0) .and. (p2(2)>0))then
       call marker%push_position (p2)
       call marker%find ('name="', skip=4, proceed=.true., pos=p1)
       call marker%find ('"', skip=1, proceed=.false., pos=p2)
       call marker%pop_position ()
    else
       call msg_error ("serializable_deserialize_from_file: no file " &
            // "header found.")
    end if
    call this%deserialize_from_marker (name, marker)
    call marker%close ()
    call marker%finalize ()
  end subroutine serializable_deserialize_from_file
  
\eatline
\nwendcode{}\nwbegindocs{76}\nwdocspar
\nwenddocs{}\nwbegincode{77}\moddef{Muli base: serial class: TBP}\plusendmoddef
  procedure :: deserialize_from_unit => &
       serializable_deserialize_from_unit
\nwendcode{}\nwbegincode{78}\moddef{Muli base: procedures}\plusendmoddef
  subroutine serializable_deserialize_from_unit (this, unit, name)
    class(ser_class_t), intent(inout) :: this
    integer, intent(in) :: unit
    character(len=*), intent(in) :: name
    logical::opened
    !!!    gfortran bug
    !!!    character::stream
    character::read
    type(marker_t)::marker
    !!!    inquire(unit=unit,opened=opened,stream=stream,read=read)
    inquire(unit=unit,opened=opened,read=read)
    if(opened)then
       !!!       if(stream=="Y")then
          if(read=="Y")then
             print *,"dummy: serializable_serialize_from_unit"
             stop
          else
             print *,"serializable_serialize_from_unit: cannot write from read-only unit."
          end if
          !!!       else
          !!!          print *,"serializable_serialize_from_unit: access kind of unit is not 'stream'."
          !!!       end if
    else
       print *,"serializable_serialize_from_unit: file is not opened."
    end if
  end subroutine serializable_deserialize_from_unit
  
\eatline
\nwendcode{}\nwbegindocs{79}This needs to be made public, and not only be present as a TBP.
\nwenddocs{}\nwbegincode{80}\moddef{Muli base: public}\plusendmoddef
  public :: serializable_deserialize_from_marker
\nwendcode{}\nwbegincode{81}\moddef{Muli base: serial class: TBP}\plusendmoddef
  procedure :: deserialize_from_marker => &
       serializable_deserialize_from_marker
\nwendcode{}\nwbegincode{82}\moddef{Muli base: procedures}\plusendmoddef
  subroutine serializable_deserialize_from_marker (this, name, marker)
    class(ser_class_t), intent(out) :: this
    character(*), intent(in) :: name
    class(marker_t), intent(inout) :: marker
    integer(dik) :: status
    if (marker%action == 2) then
       call marker%pick_instance (name, this, status)
    else       
       call msg_error ("serializable_deserialize_from_marker: Marker is " &
            // "not ready for read access.")
    end if
  end subroutine serializable_deserialize_from_marker
  
\eatline
\nwendcode{}\nwbegindocs{83}\nwdocspar
\nwenddocs{}\nwbegincode{84}\moddef{Muli base: serial class: TBP}\plusendmoddef
  generic :: serialize => serialize_to_file, serialize_to_unit, &
       serialize_to_marker
\nwendcode{}\nwbegindocs{85}\nwdocspar
\nwenddocs{}\nwbegincode{86}\moddef{Muli base: serial class: TBP}\plusendmoddef
  generic :: deserialize => deserialize_from_file, & 
       deserialize_from_unit, deserialize_from_marker
\nwendcode{}\nwbegindocs{87}\nwdocspar
\nwenddocs{}\nwbegincode{88}\moddef{Muli base: public}\plusendmoddef
  public :: serialize_print_peer_pointer
\nwendcode{}\nwbegincode{89}\moddef{Muli base: procedures}\plusendmoddef
  recursive subroutine serialize_print_peer_pointer &
       (ser, unit, parents, components, peers, name)
    class(ser_class_t), pointer, intent(in) :: ser
    integer, intent(in) :: unit
    integer(dik) :: parents, components, peers
    character(len=*), intent(in) :: name
    if (associated (ser)) then
       write (unit,*) name, " is associated."
       if (peers>0) then
          write (unit,*) "Printing components of ", name
          call ser%print_to_unit (unit, parents, components, peers - i_one)
       else
          write (unit,*) "Skipping components of ", name
       end if
    else
       write (unit,*) name, " is not associated."
    end if
  end subroutine serialize_print_peer_pointer
  
\eatline
\nwendcode{}\nwbegindocs{90}\nwdocspar
\nwenddocs{}\nwbegincode{91}\moddef{Muli base: public}\plusendmoddef
  public :: serialize_print_comp_pointer
\nwendcode{}\nwbegincode{92}\moddef{Muli base: procedures}\plusendmoddef
  recursive subroutine serialize_print_comp_pointer &
       (ser, unit, parents, components, peers, name)
    class(ser_class_t), pointer, intent(in) :: ser
    integer, intent(in) :: unit
    integer(dik), intent(in) :: parents, components, peers
    character(len=*), intent(in) :: name
    if (associated (ser)) then
       write (unit,*)  name," is associated."
       if (components > 0) then
          write (unit,*) "Printing components of ", name
          call ser%print_to_unit (unit, parents, components - i_one, peers)
       else
          write (unit,*) "Skipping components of ", name
       end if
    else
       write (unit,*) name," is not associated."
    end if
  end subroutine serialize_print_comp_pointer
  
\eatline
\nwendcode{}\nwbegindocs{93}\nwdocspar
\nwenddocs{}\nwbegincode{94}\moddef{Muli base: public}\plusendmoddef
  public :: serialize_print_allocatable
\nwendcode{}\nwbegincode{95}\moddef{Muli base: procedures}\plusendmoddef
  subroutine serialize_print_allocatable &
       (ser, unit, parents, components, peers, name)
    class(ser_class_t), allocatable, intent(in) :: ser
    integer, intent(in) :: unit
    integer(dik), intent(in) :: parents, components, peers
    character(len=*), intent(in) :: name
    if (allocated (ser)) then
       write (unit,*) name, " is allocated."
       if (components > 0) then
          write (unit,*) "Printing components of ",name
          call ser%print_to_unit (unit, parents, components-1, peers)
       else
          write (unit,*) "Skipping components of ",name
       end if
    else
       write (unit,*) name," is not allocated."
    end if
  end subroutine serialize_print_allocatable

\eatline
\nwendcode{}\nwbegindocs{96}\nwdocspar
\nwenddocs{}\nwbegincode{97}\moddef{Muli base: public}\plusendmoddef
  public :: measure_class_t
\nwendcode{}\nwbegincode{98}\moddef{Muli base: types}\plusendmoddef
  type, abstract, extends (ser_class_t) :: measure_class_t
    contains
      procedure(measure_int), public, deferred :: measure
  end type measure_class_t

\eatline
\nwendcode{}\nwbegindocs{99}\nwdocspar
\nwenddocs{}\nwbegincode{100}\moddef{Muli base: public}\plusendmoddef
  public :: identified_t
\nwendcode{}\nwbegincode{101}\moddef{Muli base: types}\plusendmoddef
  type, extends (ser_class_t) :: identified_t 
     private
     integer(dik) :: id
     type(string_t) :: name
   contains
   \LA{}Muli base: identified type: TBP\RA{}
  end type identified_t
  
\eatline
\nwendcode{}\nwbegindocs{102}\nwdocspar
\nwenddocs{}\nwbegincode{103}\moddef{Muli base: identified type: TBP}\endmoddef
  procedure :: base_write_to_marker => identified_write_to_marker
  procedure :: write_to_marker => identified_write_to_marker
\nwendcode{}\nwbegincode{104}\moddef{Muli base: procedures}\plusendmoddef
  subroutine identified_write_to_marker (this, marker, status)
    class(identified_t), intent(in) :: this
    class(marker_t), intent(inout) :: marker
    integer(dik), intent(out) :: status
    integer(dik) :: id
    id = this%get_id ()
    call marker%mark_begin ("identified_t")
    call marker%mark ("name", this%get_name ())
    call marker%mark ("id", id)
    call marker%mark_end ("identified_t")
  end subroutine identified_write_to_marker
  
\eatline
\nwendcode{}\nwbegindocs{105}\nwdocspar
\nwenddocs{}\nwbegincode{106}\moddef{Muli base: identified type: TBP}\plusendmoddef
  procedure :: base_read_from_marker => identified_read_from_marker
  procedure :: read_from_marker => identified_read_from_marker
\nwendcode{}\nwbegincode{107}\moddef{Muli base: procedures}\plusendmoddef
  subroutine identified_read_from_marker (this, marker, status)
    class(identified_t), intent(out) :: this
    class(marker_t), intent(inout) :: marker
    integer(dik), intent(out) :: status
    character(:), allocatable :: name
    call marker%pick_begin ("identified_t", status=status)
    call marker%pick ("name", name, status)
    call marker%pick ("id", this%id, status)
    call marker%pick_end ("identified_t", status=status)
    this%name = name
  end subroutine identified_read_from_marker
  
\eatline
\nwendcode{}\nwbegindocs{108}\nwdocspar
\nwenddocs{}\nwbegincode{109}\moddef{Muli base: identified type: TBP}\plusendmoddef
  procedure :: base_print_to_unit => identified_print_to_unit
  procedure :: print_to_unit => identified_print_to_unit  
\nwendcode{}\nwbegincode{110}\moddef{Muli base: procedures}\plusendmoddef
  subroutine identified_print_to_unit (this, unit, parents, components, peers)
    class(identified_t), intent(in) :: this
    integer, intent(in) :: unit
    integer(dik), intent(in) :: parents, components, peers
    write (unit, "(A)")      "Components of identified_t:"
    write (unit, "(A,A)")    "Name:             ", this%get_name ()
    write (unit, "(A,I10)")  "ID:               ", this%get_id ()
  end subroutine identified_print_to_unit
    
\eatline
\nwendcode{}\nwbegindocs{111}\nwdocspar
\nwenddocs{}\nwbegincode{112}\moddef{Muli base: identified type: TBP}\plusendmoddef
  procedure, nopass :: get_type => identified_get_type     
\nwendcode{}\nwbegincode{113}\moddef{Muli base: procedures}\plusendmoddef
  pure subroutine identified_get_type (type)
    character(:), allocatable, intent(out) :: type
    allocate (type, source="identified_t")
  end subroutine identified_get_type
 
\eatline
\nwendcode{}\nwbegindocs{114}\nwdocspar
\nwenddocs{}\nwbegincode{115}\moddef{Muli base: identified type: TBP}\plusendmoddef
  procedure, nopass :: verify_type => identified_verify_type  
\nwendcode{}\nwbegincode{116}\moddef{Muli base: procedures}\plusendmoddef
  elemental logical function identified_verify_type (type)
    character(len=*), intent(in) ::type
    identified_verify_type = (type == "identified_t")
  end function identified_verify_type
  
\eatline
\nwendcode{}\nwbegindocs{117}\nwdocspar
\nwenddocs{}\nwbegincode{118}\moddef{Muli base: identified type: TBP}\plusendmoddef
  generic :: initialize => identified_initialize  
  procedure, private :: identified_initialize 
\nwendcode{}\nwbegincode{119}\moddef{Muli base: procedures}\plusendmoddef
  subroutine identified_initialize (this, id, name)
    class(identified_t), intent(out) :: this
    integer(dik), intent(in) :: id
    character(len=*), intent(in) :: name
    this%name = name
    this%id = id 
  end subroutine identified_initialize
  
\eatline
\nwendcode{}\nwbegindocs{120}\nwdocspar
\nwenddocs{}\nwbegincode{121}\moddef{Muli base: identified type: TBP}\plusendmoddef
  procedure :: get_id => identified_get_id  
\nwendcode{}\nwbegincode{122}\moddef{Muli base: procedures}\plusendmoddef
  elemental function identified_get_id (this) result(id)
    class(identified_t), intent(in) :: this
    integer(dik) :: id
    id = this%id
  end function identified_get_id

\eatline
\nwendcode{}\nwbegindocs{123}\nwdocspar
\nwenddocs{}\nwbegincode{124}\moddef{Muli base: identified type: TBP}\plusendmoddef
  procedure :: get_name => identified_get_name
\nwendcode{}\nwbegincode{125}\moddef{Muli base: procedures}\plusendmoddef
  pure function identified_get_name (this)
    class(identified_t), intent(in) :: this
    character(len (this%name)) :: identified_get_name
    identified_get_name = char (this%name)
  end function identified_get_name

\eatline
\nwendcode{}\nwbegindocs{126}\nwdocspar
\nwenddocs{}\nwbegincode{127}\moddef{Muli base: public}\plusendmoddef
  public :: unique_t
\nwendcode{}\nwbegincode{128}\moddef{Muli base: types}\plusendmoddef
  type, extends (identified_t) :: unique_t
     private
     integer(dik) :: unique_id
   contains
   \LA{}Muli base: unique type: TBP\RA{}  
  end type unique_t

\eatline
\nwendcode{}\nwbegindocs{129}\nwdocspar
\nwenddocs{}\nwbegincode{130}\moddef{Muli base: unique type: TBP}\endmoddef
  procedure, nopass :: get_type => unique_get_type  
\nwendcode{}\nwbegincode{131}\moddef{Muli base: procedures}\plusendmoddef
  pure subroutine unique_get_type (type)
    character(:), allocatable, intent(out) :: type
    allocate (type, source="unique_t")
  end subroutine unique_get_type
  
\eatline
\nwendcode{}\nwbegindocs{132}\nwdocspar
\nwenddocs{}\nwbegincode{133}\moddef{Muli base: unique type: TBP}\plusendmoddef
  procedure, nopass :: verify_type => unique_verify_type  
\nwendcode{}\nwbegincode{134}\moddef{Muli base: procedures}\plusendmoddef
  elemental logical function unique_verify_type (type)
    character(len=*), intent(in) :: type
    unique_verify_type = (type == "unique_t")
  end function unique_verify_type
  
\eatline
\nwendcode{}\nwbegindocs{135}\nwdocspar
\nwenddocs{}\nwbegincode{136}\moddef{Muli base: unique type: TBP}\plusendmoddef
  procedure :: write_to_marker => unique_write_to_marker  
\nwendcode{}\nwbegincode{137}\moddef{Muli base: procedures}\plusendmoddef
  subroutine unique_write_to_marker (this, marker, status)
    class(unique_t), intent(in) :: this
    class(marker_t), intent(inout) :: marker
    integer(dik), intent(out) :: status
    call marker%mark_begin ("unique_t")
    call identified_write_to_marker (this, marker, status)
    call marker%mark ("unique_id", this%get_unique_id ())
    call marker%mark_end ("unique_t")
  end subroutine unique_write_to_marker
  
\eatline
\nwendcode{}\nwbegindocs{138}\nwdocspar
\nwenddocs{}\nwbegincode{139}\moddef{Muli base: unique marker: TBP}\endmoddef
  procedure :: read_from_marker => unique_read_from_marker  
\nwendcode{}\nwbegincode{140}\moddef{Muli base: procedures}\plusendmoddef
  subroutine unique_read_from_marker (this, marker, status)
    class(unique_t), intent(out) :: this
    class(marker_t), intent(inout) :: marker
    integer(dik), intent(out) :: status
    call marker%pick_begin ("unique_t", status=status)    
    call identified_read_from_marker (this, marker, status)
    call marker%pick ("unique_id", this%unique_id, status)
    call marker%pick_end ("unique_t", status)
  end subroutine unique_read_from_marker
  
\eatline
\nwendcode{}\nwbegindocs{141}\nwdocspar
\nwenddocs{}\nwbegincode{142}\moddef{Muli base: unique type: TBP}\plusendmoddef
  procedure :: print_to_unit => unique_print_to_unit  
\nwendcode{}\nwbegincode{143}\moddef{Muli base: procedures}\plusendmoddef
  subroutine unique_print_to_unit (this, unit, parents, components, peers)
    class(unique_t), intent(in) :: this
    integer,intent(in) :: unit
    integer(dik), intent(in) :: parents, components, peers
    if (parents > 0) call identified_print_to_unit &
         (this, unit, parents-1, components, peers)
    write (unit, "(A,I10)")  "Unique ID:        ", this%get_unique_id ()
  end subroutine unique_print_to_unit
  
\eatline
\nwendcode{}\nwbegindocs{144}\nwdocspar
\nwenddocs{}\nwbegincode{145}\moddef{Muli base: unique type: TBP}\plusendmoddef
  procedure :: identified_initialize => unique_initialize  
\nwendcode{}\nwbegincode{146}\moddef{Muli base: procedures}\plusendmoddef
  subroutine unique_initialize(this,id,name)
    class(unique_t), intent(out) :: this
    integer(dik), intent(in) :: id
    character(len=*), intent(in) :: name
    call identified_initialize (this, id, name)
    last_id = last_id + 1
    this%unique_id = last_id
  end subroutine unique_initialize
  
\eatline
\nwendcode{}\nwbegindocs{147}\nwdocspar
\nwenddocs{}\nwbegincode{148}\moddef{Muli base: unique type: TBP}\plusendmoddef
  procedure :: get_unique_id => unique_get_unique_id  
\nwendcode{}\nwbegincode{149}\moddef{Muli base: procedures}\plusendmoddef
  pure function unique_get_unique_id (this)
    class(unique_t), intent(in) :: this
    integer(dik) :: unique_get_unique_id
    unique_get_unique_id = this%unique_id
  end function unique_get_unique_id
  
\eatline
\nwendcode{}\nwbegindocs{150}\nwdocspar
\nwenddocs{}\nwbegincode{151}\moddef{Muli base: types}\plusendmoddef
  type :: serializable_ref_type
     private
     integer(dik) :: id
     class(ser_class_t), pointer :: ref => null()
     class(serializable_ref_type), pointer :: next => null()
   contains
   \LA{}Muli base: serial ref: TBP\RA{}  
  end type serializable_ref_type

\eatline
\nwendcode{}\nwbegindocs{152}\nwdocspar
\nwenddocs{}\nwbegincode{153}\moddef{Muli base: serial ref: TBP}\endmoddef
  procedure :: finalize => serializable_ref_finalize    
\nwendcode{}\nwbegincode{154}\moddef{Muli base: procedures}\plusendmoddef
  subroutine serializable_ref_finalize (this)
    class(serializable_ref_type), intent(inout) :: this
    class(serializable_ref_type), pointer :: next
    do while (associated (this%next))
       next => this%next
       this%next => next%next
       nullify (next%ref)
       deallocate (next)
    end do
    if (associated (this%ref))  nullify (this%ref)
  end subroutine serializable_ref_finalize
  
\nwendcode{}\nwbegindocs{155} %def serializable_ref_finalize    
\nwenddocs{}\nwbegindocs{156}\nwdocspar
\nwenddocs{}\nwbegincode{157}\moddef{Muli base: types}\plusendmoddef
  type :: position_stack_t
     private
     integer(dik), dimension(2) :: position
     class(position_stack_t), pointer :: next => null()
   contains
   \LA{}Muli base: position stack: TBP\RA{}  
  end type position_stack_t

\eatline
\nwendcode{}\nwbegindocs{158}\nwdocspar
\nwenddocs{}\nwbegincode{159}\moddef{Muli base: position stack: TBP}\endmoddef
  generic :: push => push_head, push_given  
  procedure :: push_head => position_stack_push_head  
  procedure :: push_given => position_stack_push_given    
\nwendcode{}\nwbegincode{160}\moddef{Muli base: procedures}\plusendmoddef
  subroutine position_stack_push_head (this)
    class(position_stack_t) :: this
    class(position_stack_t), pointer :: new
    allocate (new)
    new%next => this%next
    new%position = this%position
    this%next => new
  end subroutine position_stack_push_head
  
\eatline
\nwendcode{}\nwbegindocs{161}\nwdocspar
\nwenddocs{}\nwbegincode{162}\moddef{Muli base: procedures}\plusendmoddef
  subroutine position_stack_push_given (this, position)
    class(position_stack_t) :: this
    integer(dik), dimension(2), intent(in) :: position
    class(position_stack_t), pointer:: new
    allocate (new)
    new%next => this%next
    new%position = position
    this%next => new
  end subroutine position_stack_push_given
  
\eatline
\nwendcode{}\nwbegindocs{163}\nwdocspar
\nwenddocs{}\nwbegincode{164}\moddef{Muli base: position stack: TBP}\plusendmoddef
  generic :: pop => position_stack_pop, position_stack_drop  
  procedure :: position_stack_pop
  procedure :: position_stack_drop  
\nwendcode{}\nwbegincode{165}\moddef{Muli base: procedures}\plusendmoddef
  subroutine position_stack_pop (this)
    class(position_stack_t) :: this
    class(position_stack_t), pointer :: old
    if (associated (this%next)) then
       old => this%next
       this%next => old%next
       this%position = old%position
       deallocate (old)
    end if
  end subroutine position_stack_pop

\eatline
\nwendcode{}\nwbegindocs{166}\nwdocspar
\nwenddocs{}\nwbegincode{167}\moddef{Muli base: procedures}\plusendmoddef
  subroutine position_stack_drop (this, position)
    class(position_stack_t) :: this
    integer(dik), dimension(2), intent(out) :: position
    class(position_stack_t), pointer :: old
    if (associated (this%next)) then
       old => this%next
       this%next => old%next
       position = old%position
       deallocate (old)
    else
       position= [0,0]
    end if
  end subroutine position_stack_drop
  
\eatline
\nwendcode{}\nwbegindocs{168}\nwdocspar
\nwenddocs{}\nwbegincode{169}\moddef{Muli base: position stack: TBP}\plusendmoddef
  procedure :: nth_position => position_stack_nth_position  
\nwendcode{}\nwbegincode{170}\moddef{Muli base: procedures}\plusendmoddef
  function position_stack_nth_position (this, n) result (position)
    class(position_stack_t), intent(in) :: this
    integer(dik), intent(in) :: n
    integer(dik), dimension(2) :: position
    class(position_stack_t), pointer :: tmp
    integer(dik) :: pos
    tmp => this%next
    pos = n
    do while (associated (tmp) .and. pos>0)
       tmp => tmp%next
       pos = pos - 1
    end do
    if (associated(tmp)) then
       position = tmp%position
    else
       position = [0,0]
    end if
  end function position_stack_nth_position
  
\eatline
\nwendcode{}\nwbegindocs{171}\nwdocspar
\nwenddocs{}\nwbegincode{172}\moddef{Muli base: position stack: TBP}\plusendmoddef
  procedure :: first => position_stack_first  
\nwendcode{}\nwbegincode{173}\moddef{Muli base: procedures}\plusendmoddef
  function position_stack_first(this) result(position)
    class(position_stack_t), intent(in) :: this
    integer(kind=dik), dimension(2) :: position, tmp_position
    class(position_stack_t), pointer :: tmp_stack
    tmp_position = this%position
    tmp_stack => this%next
    do while (associated (tmp_stack))
       if (page_ring_position_is_before (tmp_stack%position, tmp_position)) then
          tmp_position = tmp_stack%position
       end if
       tmp_stack => tmp_stack%next
    end do
  end function position_stack_first
  
\eatline
\nwendcode{}\nwbegindocs{174}\nwdocspar
\nwenddocs{}\nwbegincode{175}\moddef{Muli base: position stack: TBP}\plusendmoddef
  procedure :: last => position_stack_last  
\nwendcode{}\nwbegincode{176}\moddef{Muli base: procedures}\plusendmoddef
  function position_stack_last (this) result (position)
    class(position_stack_t), intent(in) :: this
    integer(dik), dimension(2) :: position, tmp_position
    class(position_stack_t), pointer :: tmp_stack
    tmp_position = this%position
    tmp_stack => this%next
    do while (associated (tmp_stack))
       if (page_ring_position_is_before (tmp_position, tmp_stack%position)) then
          tmp_position = tmp_stack%position
       end if
       tmp_stack => tmp_stack%next
    end do
  end function position_stack_last
    
\eatline
\nwendcode{}\nwbegindocs{177}\nwdocspar
\nwenddocs{}\nwbegincode{178}\moddef{Muli base: position stack: TBP}\plusendmoddef
  procedure :: range => position_stack_range  
\nwendcode{}\nwbegincode{179}\moddef{Muli base: procedures}\plusendmoddef
  pure function position_stack_range (this) result (position)
    class(position_stack_t), intent(in) :: this
    integer(dik), dimension(2) :: position
    class(position_stack_t), pointer :: tmp
  end function position_stack_range
    
\eatline
\nwendcode{}\nwbegindocs{180}\nwdocspar
\nwenddocs{}\nwbegincode{181}\moddef{Muli base: public}\plusendmoddef
  public :: page_ring_t
\nwendcode{}\nwbegincode{182}\moddef{Muli base: types}\plusendmoddef
  type :: page_ring_t
     private
     logical :: asynchronous = serialize_default_asynchronous
     logical :: eof_reached = .false.
     integer :: unit = -1
     integer(dik) :: ring_size = 2
     integer(dik) :: action = 0
     integer(dik) :: eof_int = -1
     integer(dik) :: out_unit = output_unit
     integer(dik) :: err_unit = error_unit
     integer(dik), dimension(2) :: active_pages = [0,-1]
     integer(dik), dimension(2) :: eof_pos = [-1,-1]
     type(string_t) :: eof_string
     type(position_stack_t) :: position_stack
     character(serialize_page_size), dimension(:), allocatable::ring     
   contains
   \LA{}Muli base: page ring: TBP\RA{}  
  end type page_ring_t

\eatline
\nwendcode{}\nwbegindocs{183}These are the {\tt{}page{\char95}ring{\char95}t} procedures, here for read access only:
\nwenddocs{}\nwbegincode{184}\moddef{Muli base: page ring: TBP}\endmoddef
  procedure :: open_for_read_access => page_ring_open_for_read_access
\nwendcode{}\nwbegincode{185}\moddef{Muli base: procedures}\plusendmoddef
  subroutine page_ring_open_for_read_access &
       (this, file, eof_string, asynchronous)
    class(page_ring_t), intent(inout) :: this
    character(*), intent(in) :: file, eof_string
    logical, intent(in), optional :: asynchronous
    logical :: exist
    this%eof_string = eof_string
    inquire (file=file, exist=exist)
    if (exist) then
       this%action = 2
    else
       call msg_error ("page_ring_open: File " // file // " is opened " &
            // "for read access but does not exist.")
    end if
    if (present (asynchronous)) this%asynchronous = asynchronous
    if (this%unit < 0)  call generate_unit (this%unit, 100, 1000)
    if (this%unit < 0) then
       call msg_error ("page_ring_open: No free unit found.")
    end if
    this%ring_size = 2
    call this%set_position ([i_zero,i_one])
    this%active_pages = [i_zero,-i_one]
    if (allocated (this%ring)) deallocate (this%ring)
    allocate (this%ring (i_zero:this%ring_size - i_one))
    if (this%asynchronous) then
       open (this%unit, file=file, access="stream", &
            action="read", asynchronous="yes", status="old")
    else
       open (this%unit, file=file, access="stream", action="read", &
            asynchronous="no", status="old")
    end if
    call this%read_page ()
  end subroutine page_ring_open_for_read_access
  
\eatline
\nwendcode{}\nwbegindocs{186}\nwdocspar
\nwenddocs{}\nwbegincode{187}\moddef{Muli base: page ring: TBP}\plusendmoddef
  procedure :: read_page => page_ring_read_page  
\nwendcode{}\nwbegincode{188}\moddef{Muli base: procedures}\plusendmoddef
  subroutine page_ring_read_page (this)
    class(page_ring_t), intent(inout) :: this
    integer(dik) :: iostat
    character(8) :: iomsg
    if (.not. this%eof_reached) then
       call this%activate_next_page ()
       read (this%unit, iostat=iostat) this%ring (this%last_index ())
       if (iostat == iostat_end) then
          this%eof_reached = .true.
          this%eof_pos(1) = this%last_page ()
          this%eof_pos(2) = index(this%ring(this%last_index()), &
               char(this%eof_string))
          this%eof_pos(2) = this%eof_pos(2) + len(this%eof_string) - 1
          this%eof_int = page_ring_ordinal(this%eof_pos)
       end if
    end if
  end subroutine page_ring_read_page
  
\eatline
\nwendcode{}\nwbegindocs{189}Those are the write access only type-bound procedures of
{\tt{}page{\char95}ring{\char95}t}:
\nwenddocs{}\nwbegincode{190}\moddef{Muli base: page ring: TBP}\plusendmoddef
  procedure :: open_for_write_access => page_ring_open_for_write_access
\nwendcode{}\nwbegincode{191}\moddef{Muli base: procedures}\plusendmoddef
  subroutine page_ring_open_for_write_access (this, file, asynchronous)
    class(page_ring_t), intent(inout) :: this
    character(*), intent(in) :: file
    logical, intent(in), optional :: asynchronous
    this%action = 1
    if (present (asynchronous)) this%asynchronous = asynchronous
    if (this%unit < 0)  call generate_unit (this%unit, 100, 1000)
    if (this%unit < 0) then
       call msg_error ("page_ring_open: No free unit found.")
    end if
    this%ring_size = 2 
    call this%set_position ([i_zero,i_one])
    this%active_pages = [i_zero,-i_one]
    if (allocated (this%ring))  deallocate (this%ring)
    allocate (this%ring (i_zero:this%ring_size-i_one))
    if (this%asynchronous) then
       open (this%unit, file=file, access="stream", action="write", &
            asynchronous="yes", status="replace")
    else
       open (this%unit, file=file, access="stream", action="write", &
            asynchronous="no",status="replace")
    end if
  end subroutine page_ring_open_for_write_access

\eatline
\nwendcode{}\nwbegindocs{192}\nwdocspar
\nwenddocs{}\nwbegincode{193}\moddef{Muli base: page ring: TBP}\plusendmoddef
  procedure :: flush => page_ring_flush
\nwendcode{}\nwbegincode{194}\moddef{Muli base: procedures}\plusendmoddef
  subroutine page_ring_flush (this)
    class(page_ring_t), intent(inout) :: this
    integer(dik) :: page
    do while (this%active_pages(1) < this%actual_page ())
       if (this%asynchronous) then
          write (this%unit, asynchronous="yes") &
               this%ring(mod(this%active_pages(1), this%ring_size))
       else
          write (this%unit, asynchronous="no") &
               this%ring(mod(this%active_pages(1), this%ring_size))
       end if
       this%active_pages(1) = this%active_pages(1) + 1
    end do
  end subroutine page_ring_flush

\eatline
\nwendcode{}\nwbegindocs{195} 
\nwenddocs{}\nwbegincode{196}\moddef{Muli base: page ring: TBP}\plusendmoddef
  procedure :: break => page_ring_break
\nwendcode{}\nwbegincode{197}\moddef{Muli base: procedures}\plusendmoddef
  subroutine page_ring_break(this)
    class(page_ring_t), intent(inout) :: this
    if (this%actual_page () >= this%active_pages(2)) &
         call this%activate_next_page ()
    call this%turn_page ()
  end subroutine page_ring_break
  
\eatline
\nwendcode{}\nwbegindocs{198}For comparisons:
\nwenddocs{}\nwbegincode{199}\moddef{Muli base: page ring: TBP}\plusendmoddef
  procedure :: str_equal => page_ring_str_equal  
\nwendcode{}\nwbegincode{200}\moddef{Muli base: procedures}\plusendmoddef
  pure logical function page_ring_str_equal (this, string, pos)
    class(page_ring_t), intent(in) :: this
    character(*), intent(in) :: string
    integer(dik), dimension(2,2), intent(in) :: pos
    page_ring_str_equal = string == this%substring (pos)
  end function page_ring_str_equal
  
\eatline
\nwendcode{}\nwbegindocs{201}Routines for searching:
\nwenddocs{}\nwbegincode{202}\moddef{Muli base: page ring: TBP}\plusendmoddef
  generic :: find => page_ring_find, page_ring_find_default
  procedure, private :: page_ring_find
  procedure, private :: page_ring_find_default  
\nwendcode{}\nwbegincode{203}\moddef{Muli base: procedures}\plusendmoddef
  recursive subroutine page_ring_find &
       (this, exp, start, limit, skip, proceed, pos)
    class(page_ring_t), intent(inout) :: this
    integer(dik), dimension(2), intent(in) :: start
    integer(dik), dimension(2), intent(in) :: limit
    character(*), intent(in) :: exp
    integer, intent(in) :: skip
    logical, intent(in) :: proceed
    integer(dik), dimension(2), intent(out) :: pos
    integer(dik) :: page, page2, ind
    page = this%ring_index (start(1))
    if (limit(1) == start(1)) then
       ind = index(this%ring(page) (start(2):limit(2)), exp)
       if (ind > 0) then
          select case (skip)
          case (1)
             pos= [start(1), start(2)+ind-2]
             if (pos(2) == 0) then
                pos(1) = pos(1) - 1
                pos(2) = serialize_page_size
             end if
          case (2)
             pos = [start(1), start(2)+ind-1]
          case (3)
             pos = [start(1), start(2)+ind+len(exp)-2]
          case (4)
             pos = [start(1),start(2)+ind+len(exp)-1]
             if (pos(1) == this%last_page())  call this%read_page ()
             if (pos(2) > serialize_page_size) then
                pos(1) = pos(1) + 1
                pos(2) = pos(2) - serialize_page_size
             end if
          end select
          if (proceed)  call this%set_position (pos)
       else
          Call msg_warning ("page_ring_find: limit reached.")
          pos = [-1, -1]
       end if
    else
       ind = index (this%ring(page) (start(2):), exp)
       if (ind > 0) then
          select case (skip)
          case (1)
             pos = [start(1), start(2)+ind-2]
             if (pos(2) == 0) then
                pos(1) = pos(1) - 1
                pos(2) = serialize_page_size
             end if
          case (2)
             pos = [start(1), start(2)+ind-1]
          case (3)
             pos = [start(1), start(2)+ind+len(exp)-2]
          case (4)
             pos = [start(1), start(2)+ind+len(exp)-1]
             if (pos(1) == this%last_page ())  call this%read_page ()
             if (pos(2) > serialize_page_size) then
                pos(1) = pos(1) + 1
                pos(2) = i_one
             end if
          end select
          if(proceed)call this%set_position(pos)
       else
          if (start(1) + 1 > this%active_pages (2)) then
             call this%read_page ()
             page = this%ring_index(start(1))
          end if
          page2 = this%ring_index(start(1)+1)
          ind = index(this%ring(page) (serialize_page_size - &
               len(exp)+1:)//this%ring(page2)(:len(exp)),exp)
          if (ind > 0) then
             select case (skip)
             case (1)
                pos = [start(1), serialize_page_size-len(exp)+ind-1]
             case (2)
                pos = [start(1), serialize_page_size-len(exp)+ind]
             case (3)
                pos = [start(1)+1, ind-1]
             case (4)
                pos = [start(1)+1, ind]
             end select
             if (pos(2) > serialize_page_size) then
                pos(1) = pos(1) + 1
                pos(2) = pos(2) - serialize_page_size
             else
                if (pos(2) < 0) then
                   pos(1) = pos(1) - 1
                   pos(2) = pos(2) + serialize_page_size
                end if
             end if
             if (proceed)  call this%set_position (pos)
          else
             if (proceed)  this%active_pages(1) = this%active_pages(2)
             call this%find (exp, [start(1) + i_one, i_one], &
                  limit, skip, proceed, pos)
          end if
       end if
    end if
  end subroutine page_ring_find

\eatline
\nwendcode{}\nwbegindocs{204}\nwdocspar
\nwenddocs{}\nwbegincode{205}\moddef{Muli base: procedures}\plusendmoddef
  subroutine page_ring_find_default (this, exp, skip, proceed, pos)
    class(page_ring_t), intent(inout) :: this
    character(*), intent(in), optional :: exp
    integer, intent(in) :: skip
    logical, intent(in) :: proceed
    integer(dik), dimension(2), intent(out) :: pos
    call this%find (exp, this%position_stack%position, this%eof_pos, &
         skip, proceed, pos)
  end subroutine page_ring_find_default
  
\eatline
\nwendcode{}\nwbegindocs{206}\nwdocspar
\nwenddocs{}\nwbegincode{207}\moddef{Muli base: page ring: TBP}\plusendmoddef
  procedure :: find_pure => page_ring_find_pure  
\nwendcode{}\nwbegincode{208}\moddef{Muli base: procedures}\plusendmoddef
  pure recursive function page_ring_find_pure &
       (this, exp, start, limit, skip) result (pos)
    class(page_ring_t),intent(in) :: this
    integer(dik), dimension(2), intent(in) :: start
    integer(dik), dimension(2), intent(in) :: limit
    character(*),intent(in) :: exp
    integer,optional,intent(in) :: skip
    integer(dik), dimension(2) :: pos
    integer(dik) :: page, page2, ind, actual_skip
    !!! Is the starting point before limit?
    if (start(1) <= limit(1)) then
       !!! Default skip is what you expect from the build-in index function
       if (present(skip)) then
          actual_skip = skip
       else
          actual_skip = 2
       end if
       page = mod(start(1), this%ring_size)
       !!! Does the scanning region end on the page?
       if (start(1) == limit(1)) then
          ind = index (this%ring (page) (start(2):limit(2)),exp)
       else
          ind = index (this%ring (page) (start(2):),exp)
       end if
       if (ind > 0) then
          !!! substring found on first page
          select case (actual_skip)
          case (1)
             pos = [start(1), start(2)+ind-2]
             if (pos(2) == 0) then
                pos(1) = pos(1) - 1
                pos(2) = serialize_page_size
             end if
          case (2)
             pos= [start(1), start(2)+ind-1]
          case (3)
             pos= [start(1), start(2)+ind+len(exp)-2]
          case (4)
             pos= [start(1), start(2)+ind+len(exp)-1]
             if (pos(2) > serialize_page_size) then
                pos(1) = pos(1) + 1
                pos(2) = pos(2) - serialize_page_size
             end if
          end select
       else
          !!! Substring not found on first page. Is the next page already read?
          if ((start(1) >= limit(1)) .or. &
              (start(1)+1 > this%active_pages(2))) then
             !!! Either the limit is reached or the next page is not ready.
             pos = [0, 0]
          else
             !!! The next page is available.
             page2 = mod(start(1)+1, this%ring_size)
             !!! We concatenate the edges. When l is the length of exp, 
             !!! then we want to concatenate the l-1 last characters of 
             !!! page one and the first l characters of page two.
             ! print *,"overlap: |",this%ring(page) &
             !    (serialize_page_size-len(exp)+2:)//this%ring(page2) &
             !    (:len(exp)),"|"
             ind = index (this%ring(page) (serialize_page_size - &
                  len(exp)+2:)//this%ring(page2) (:len(exp)),exp)
             if (ind > 0) then
                select case (actual_skip)
                case (1)
                   pos = [start(1), serialize_page_size-len(exp)+ind]
                case (2)
                   pos = [start(1), serialize_page_size-len(exp)+ind+1]
                case (3)
                   pos = [start(1)+1, ind]
                case (4)
                   pos = [start(1)+1, ind+1]
                end select
             else
                !!! EXP is not found in the overlap region. 
                !!! We recursively search the next pages.
                pos = this%find_pure (exp, [start(i_one) + i_one, i_one], &
                     limit, skip)
             end if
          end if
       end if
    else
       !!! Limit is before start
       pos = [0, 0]
    end if
  end function page_ring_find_pure
  
\eatline
\nwendcode{}\nwbegindocs{209}{\tt{}page{\char95}ring{\char95}t} routines for positioning:
\nwenddocs{}\nwbegincode{210}\moddef{Muli base: page ring: TBP}\plusendmoddef
  generic :: get_position => page_ring_get_position1, page_ring_get_position2  
  procedure, private :: page_ring_get_position1
  procedure, private :: page_ring_get_position2  
\nwendcode{}\nwbegincode{211}\moddef{Muli base: procedures}\plusendmoddef
  pure subroutine page_ring_get_position1 (this, pos)
    class(page_ring_t), intent(in) :: this
    integer(dik), intent(out) :: pos
    pos = page_ring_ordinal (this%position_stack%position)
  end subroutine page_ring_get_position1

\eatline
\nwendcode{}\nwbegindocs{212}\nwdocspar
\nwenddocs{}\nwbegincode{213}\moddef{Muli base: procedures}\plusendmoddef
  pure subroutine page_ring_get_position2 (this, pos)
    class(page_ring_t), intent(in) :: this
    integer(dik), dimension(2), intent(out) :: pos
    pos = this%position_stack%position
  end subroutine page_ring_get_position2

\eatline
\nwendcode{}\nwbegindocs{214} 
\nwenddocs{}\nwbegincode{215}\moddef{Muli base: page ring: TBP}\plusendmoddef
  generic :: pop_position => pop_actual_position, pop_given_position  
  procedure, private :: pop_actual_position => &
       page_ring_ring_pop_actual_position
  procedure, private :: pop_given_position => &
       page_ring_ring_pop_given_position  
\nwendcode{}\nwbegincode{216}\moddef{Muli base: procedures}\plusendmoddef
  subroutine page_ring_ring_pop_actual_position (this)
    class(page_ring_t), intent(inout) :: this
    call this%position_stack%pop ()
  end subroutine page_ring_ring_pop_actual_position

\eatline
\nwendcode{}\nwbegindocs{217}\nwdocspar
\nwenddocs{}\nwbegincode{218}\moddef{Muli base: procedures}\plusendmoddef
  subroutine page_ring_ring_pop_given_position (this, pos)
    class(page_ring_t), intent(inout) :: this
    integer(dik), dimension(2), intent(out) :: pos
    call this%position_stack%pop (pos)
  end subroutine page_ring_ring_pop_given_position

\eatline
\nwendcode{}\nwbegindocs{219}\nwdocspar
\nwenddocs{}\nwbegincode{220}\moddef{Muli base: page ring: TBP}\plusendmoddef
  generic :: push_position => push_actual_position, push_given_position  
  procedure, private :: push_actual_position => &
       page_ring_ring_push_actual_position
  procedure, private :: push_given_position => &
       page_ring_ring_push_given_position  
\nwendcode{}\nwbegincode{221}\moddef{Muli base: procedures}\plusendmoddef
  subroutine page_ring_ring_push_actual_position (this)
    class(page_ring_t), intent(inout) :: this
    call this%position_stack%push ()
  end subroutine page_ring_ring_push_actual_position

\eatline
\nwendcode{}\nwbegindocs{222}\nwdocspar
\nwenddocs{}\nwbegincode{223}\moddef{Muli base: procedures}\plusendmoddef
  subroutine page_ring_ring_push_given_position (this, pos)
    class(page_ring_t), intent(inout) :: this
    integer(dik), dimension(2), intent(in) :: pos
    call this%position_stack%push (pos)
  end subroutine page_ring_ring_push_given_position

\eatline
\nwendcode{}\nwbegindocs{224}\nwdocspar
\nwenddocs{}\nwbegincode{225}\moddef{Muli base: page ring: TBP}\plusendmoddef
  procedure :: set_position => page_ring_set_position  
\nwendcode{}\nwbegincode{226}\moddef{Muli base: procedures}\plusendmoddef
  subroutine page_ring_set_position (this, pos)
    class(page_ring_t), intent(inout) :: this
    integer(dik), dimension(2), intent(in) :: pos
    this%position_stack%position = pos
  end subroutine page_ring_set_position

\eatline
\nwendcode{}\nwbegindocs{227}\nwdocspar
\nwenddocs{}\nwbegincode{228}\moddef{Muli base: page ring: TBP}\plusendmoddef
  procedure :: turn_page => page_ring_turn_page 
\nwendcode{}\nwbegincode{229}\moddef{Muli base: procedures}\plusendmoddef
  subroutine page_ring_turn_page (this)
    class(page_ring_t), intent(inout) :: this
    this%position_stack%position(1) = this%position_stack%position(1) + 1
    this%position_stack%position(2) = 1
  end subroutine page_ring_turn_page

\eatline
\nwendcode{}\nwbegindocs{230}\nwdocspar
\nwenddocs{}\nwbegincode{231}\moddef{Muli base: page ring: TBP}\plusendmoddef
  procedure :: proceed => page_ring_proceed
\nwendcode{}\nwbegincode{232}\moddef{Muli base: procedures}\plusendmoddef
  subroutine page_ring_proceed (this, n, deactivate)
    class(page_ring_t), intent(inout) :: this
    integer(dik), intent(in) :: n
    logical, intent(in), optional :: deactivate
    integer(dik) :: offset
    offset = this%position_stack%position(2) + n
    do while (offset > serialize_page_size)
       if (this%position_stack%position(1) >= this%active_pages(2)) &
            call this%activate_next_page ()
       this%position_stack%position(1) = this%position_stack%position(1) + 1
       offset = offset - serialize_page_size
    end do
    this%position_stack%position(2) = offset
    if (present (deactivate)) then
       if (deactivate)this%active_pages(1) = this%actual_page ()
    end if
  end subroutine page_ring_proceed
  
\eatline
\nwendcode{}\nwbegindocs{233}These are the {\tt{}page{\char95}ring{\char95}t} routines for printing:
\nwenddocs{}\nwbegincode{234}\moddef{Muli base: page ring: TBP}\plusendmoddef
  procedure :: print_to_unit => page_ring_print_to_unit
\nwendcode{}\nwbegincode{235}\moddef{Muli base: procedures}\plusendmoddef
  subroutine page_ring_print_to_unit (this, unit, parents, components, peers)
    class(page_ring_t), intent(in) :: this
    integer, intent(in) :: unit
    integer(dik), intent(in) :: parents, components, peers
    write (unit, "(1x,A)")   "Components of page_ring_t: "
    write (unit, "(3x,A,A)")  "asynchronous: ", this%asynchronous
    write (unit, "(3x,A,L1)") "eof reached:  ", this%eof_reached
    write (unit, "(3x,A,I0)") "ring_size:    ", this%ring_size
    write (unit, "(3x,A,I0)") "unit:         ", this%unit
    write (unit, "(3x,A,I0)") "action:       ", this%action
    write (unit, "(3x,A,I0,I0)") &
         "position:     ", this%position_stack%position
    write (unit, "(3x,A,I0)") "active_pages: ", this%active_pages
    write (unit, "(3x,A,I0)") "file size:    ", this%eof_int
    write (unit, "(3x,A,I0,I0)")  "eof position: ", this%eof_pos
    write (unit, "(3x,A,A)")  "eof string:   ", char(this%eof_string)
    if (allocated (this%ring)) then
       write (unit, "(3x,A)") "Ring is allocated."
       if (components > 0)  call this%print_ring (unit)
    else
       write (unit, "(3x,A)")  "Ring is not allocated."
    end if
  end subroutine page_ring_print_to_unit
  
\eatline
\nwendcode{}\nwbegindocs{236}\nwdocspar
\nwenddocs{}\nwbegincode{237}\moddef{Muli base: page ring: TBP}\plusendmoddef
  procedure :: print_ring => page_ring_print_ring
\nwendcode{}\nwbegincode{238}\moddef{Muli base: procedures}\plusendmoddef
  subroutine page_ring_print_ring (this, unit)
    class(page_ring_t), intent(in) :: this
    integer, intent(in) :: unit
    integer(dik) :: n
    write (unit, "(1x,A)")  "Begin of page ring"
    do n = this%active_pages(1), this%active_pages(2)
       write (unit, "(3x,A,I0,A,A)") &
            "(", n, ")", this%ring (mod(n, this%ring_size))
    end do
    write (unit, "(1x,A)")  "End of page ring"
  end subroutine page_ring_print_ring
  
\eatline
\nwendcode{}\nwbegindocs{239}\nwdocspar
\nwenddocs{}\nwbegincode{240}\moddef{Muli base: page ring: TBP}\plusendmoddef
  procedure :: print_position => page_ring_print_position
\nwendcode{}\nwbegincode{241}\moddef{Muli base: procedures}\plusendmoddef
  subroutine page_ring_print_position(this)
    class(page_ring_t), intent(inout) :: this
    print *, this%actual_position(), &
         this%ring(this%actual_index()) (:this%actual_offset() - 1), "|", &
         this%ring(this%actual_index()) (this%actual_offset():)
  end subroutine page_ring_print_position
  
\eatline
\nwendcode{}\nwbegindocs{242}Here are the {\tt{}page{\char95}ring{\char95}t} routines for writing:
\nwenddocs{}\nwbegincode{243}\moddef{Muli base: page ring: TBP}\plusendmoddef
  procedure :: put => page_ring_put
\nwendcode{}\nwbegincode{244}\moddef{Muli base: procedures}\plusendmoddef
  subroutine page_ring_put (this)
    class(page_ring_t), intent(inout) :: this
  end subroutine page_ring_put

\eatline
\nwendcode{}\nwbegindocs{245}\nwdocspar
\nwenddocs{}\nwbegincode{246}\moddef{Muli base: page ring: TBP}\plusendmoddef
  generic :: push => push_string, push_integer, push_integer_dik, &
       push_real, push_integer_array, push_integer_array_dik, &
       push_real_array
\nwendcode{}\nwbegincode{247}\moddef{Muli base: page ring: TBP}\plusendmoddef
  procedure, private :: push_string => page_ring_push_string
\nwendcode{}\nwbegincode{248}\moddef{Muli base: procedures}\plusendmoddef
  recursive subroutine page_ring_push_string (this, string)
    class(page_ring_t), intent(inout) :: this
    character(*), intent(in) :: string
    integer(dik) :: cut, l
    l = len(string)
    if (l <= serialize_page_size-this%actual_offset()+1) then
       this%ring(this%actual_index()) &
            (this%actual_offset():this%actual_offset()+l-1)=string
       if (l == serialize_page_size-this%actual_offset()+1) then
          call this%break()
          call this%flush()
       else
          call this%proceed(l)
       end if       
    else
       cut = serialize_page_size-this%actual_offset() + 1
       call this%push_string(string(:cut))
       call this%push_string(string(cut+1:))
    end if
  end subroutine page_ring_push_string

\eatline
\nwendcode{}\nwbegindocs{249}\nwdocspar
\nwenddocs{}\nwbegincode{250}\moddef{Muli base: page ring: TBP}\plusendmoddef
  procedure, private :: push_integer => page_ring_push_integer
\nwendcode{}\nwbegincode{251}\moddef{Muli base: procedures}\plusendmoddef
  subroutine page_ring_push_integer (this, in)
    class(page_ring_t), intent(inout) :: this
    integer, intent(in) :: in
    call this%push_integer_dik (int(in,kind=dik))
  end subroutine page_ring_push_integer

\eatline
\nwendcode{}\nwbegindocs{252}\nwdocspar
\nwenddocs{}\nwbegincode{253}\moddef{Muli base: page ring: TBP}\plusendmoddef
  procedure, private :: push_integer_dik => page_ring_push_integer_dik
\nwendcode{}\nwbegincode{254}\moddef{Muli base: procedures}\plusendmoddef
  recursive subroutine page_ring_push_integer_dik (this, int)
    class(page_ring_t), intent(inout) :: this
    integer(dik), intent(in) :: int
    integer(dik) :: int1
    if (int < 0) then
       call this%push ("-")
       call this%push_integer_dik (-int)
    else
       if (int > 9)  call this%push (int/10)
       int1 = mod(int, 10*i_one)
       select case (int1)
       case (0)
          call this%push ("0")
       case (1)
          call this%push ("1")
       case (2)
          call this%push ("2")
       case (3)
          call this%push ("3")
       case (4)
          call this%push ("4")
       case (5)
          call this%push ("5")
       case (6)
          call this%push ("6")
       case (7)
          call this%push ("7")
       case (8)
          call this%push ("8")
       case (9)
          call this%push ("9")
       end select
    end if
  end subroutine page_ring_push_integer_dik
  
\eatline
\nwendcode{}\nwbegindocs{255}\nwdocspar
\nwenddocs{}\nwbegincode{256}\moddef{Muli base: page ring: TBP}\plusendmoddef
  procedure, private :: push_integer_array => page_ring_push_integer_array
\nwendcode{}\nwbegincode{257}\moddef{Muli base: procedures}\plusendmoddef
  subroutine page_ring_push_integer_array(this,int)
    class(page_ring_t), intent(inout) :: this
    integer, dimension(:), intent(in) :: int
    integer :: n
    do n = 1, size(int)
       call this%push (int(n))
       call this%push (" ")
    end do
  end subroutine page_ring_push_integer_array
  
\eatline
\nwendcode{}\nwbegindocs{258}\nwdocspar
\nwenddocs{}\nwbegincode{259}\moddef{Muli base: page ring: TBP}\plusendmoddef
  procedure, private :: push_integer_array_dik => &
       page_ring_push_integer_array_dik 
\nwendcode{}\nwbegincode{260}\moddef{Muli base: procedures}\plusendmoddef
  subroutine page_ring_push_integer_array_dik(this,int)
    class(page_ring_t), intent(inout) :: this
    integer(dik), dimension(:), intent(in) :: int
    integer(dik) :: n
    do n = 1, size(int)
       call this%push (int(n))
       call this%push (" ")
    end do
  end subroutine page_ring_push_integer_array_dik
    
\eatline
\nwendcode{}\nwbegindocs{261}\nwdocspar
\nwenddocs{}\nwbegincode{262}\moddef{Muli base: page ring: TBP}\plusendmoddef
  procedure, private :: push_real => page_ring_push_real
\nwendcode{}\nwbegincode{263}\moddef{Muli base: procedures}\plusendmoddef
  subroutine page_ring_push_real (this, dou)
    class(page_ring_t), intent(inout) :: this
    real(default), intent(in) :: dou
    integer(dik) :: f
    ! print *,"page_ring_push_real: ",dou
    if (dou == 0D0) then
       call this%push ("0")
    else
       f = int (scale (fraction(dou), digits(dou)), kind=dik)
       call this%push (digits(dou))
       call this%push (":")    
       call this%push (f)
       call this%push (":")
       call this%push (exponent(dou))
    end if
    call this%push (" ")
  end subroutine page_ring_push_real

\eatline
\nwendcode{}\nwbegindocs{264}\nwdocspar
\nwenddocs{}\nwbegincode{265}\moddef{Muli base: page ring: TBP}\plusendmoddef
  procedure, private :: push_real_array => page_ring_push_real_array  
\nwendcode{}\nwbegincode{266}\moddef{Muli base: procedures}\plusendmoddef
  subroutine page_ring_push_real_array (this, dou)
    class(page_ring_t), intent(inout) :: this
    real(default), dimension(:), intent(in) :: dou
    integer(dik) :: n
    do n=1, size(dou)
       call this%push (dou(n))
    end do
  end subroutine page_ring_push_real_array

\eatline
\nwendcode{}\nwbegindocs{267}\nwdocspar
\nwenddocs{}\nwbegincode{268}\moddef{Muli base: page ring: TBP}\plusendmoddef
  procedure :: get_character => page_ring_get_character  
\nwendcode{}\nwbegincode{269}\moddef{Muli base: procedures}\plusendmoddef
  elemental function  page_ring_get_character (this)
    class(page_ring_t), intent(in) :: this
    character :: page_ring_get_character
    page_ring_get_character = this%ring (this%actual_index()) &
         (this%actual_offset():this%actual_offset())
  end function page_ring_get_character

\eatline
\nwendcode{}\nwbegindocs{270}\nwdocspar
\nwenddocs{}\nwbegincode{271}\moddef{Muli base: page ring: TBP}\plusendmoddef
  procedure :: allocate_substring => page_ring_allocate_substring     
\nwendcode{}\nwbegincode{272}\moddef{Muli base: procedures}\plusendmoddef
  subroutine page_ring_allocate_substring (this, p1, p2, string)
    class(page_ring_t), intent(in) :: this
    integer(dik), dimension(2), intent(in) :: p1, p2
    character(:), allocatable, intent(out) :: string
    string = this%substring (p1, p2)
  end subroutine page_ring_allocate_substring
  
\eatline
\nwendcode{}\nwbegindocs{273}\nwdocspar
\nwenddocs{}\nwbegincode{274}\moddef{Muli base: page ring: TBP}\plusendmoddef
  procedure :: pop_character => page_ring_pop_character
\nwendcode{}\nwbegincode{275}\moddef{Muli base: procedures}\plusendmoddef
  subroutine page_ring_pop_character (this, c)
    class(page_ring_t), intent(inout) :: this
    character, intent(out) :: c
    c = this%ring (this%actual_index()) &
         (this%actual_offset():this%actual_offset())
    if (this%actual_offset () == serialize_page_size)  call this%read_page
    call this%proceed (i_one)
  end subroutine page_ring_pop_character
  
\eatline
\nwendcode{}\nwbegindocs{276}\nwdocspar
\nwenddocs{}\nwbegincode{277}\moddef{Muli base: page ring: TBP}\plusendmoddef
  procedure :: pop_by_keys => page_ring_pop_by_keys
\nwendcode{}\nwbegincode{278}\moddef{Muli base: procedures}\plusendmoddef
  subroutine page_ring_pop_by_keys (this, start, stop, inclusive, res)
    class(page_ring_t), intent(inout) :: this
    character(*), intent(in), optional :: start
    character(*), intent(in) :: stop
    logical, optional, intent(in) :: inclusive
    character(len=*), intent(out) :: res
    integer(dik), dimension(2) :: i1, i2
    if (inclusive) then
       call this%find (start, 2, .true., i1)
       call this%find (stop, 3, .false., i2)
    else
       call this%find (start, 4, .true., i1)
       call this%find (stop, 1, .false., i2)
    end if
    res = this%substring (i1, i2)
    call this%set_position (i2)
  end subroutine page_ring_pop_by_keys
  
\eatline
\nwendcode{}\nwbegindocs{279}\nwdocspar
\nwenddocs{}\nwbegincode{280}\moddef{Muli base: page ring: TBP}\plusendmoddef
  generic :: substring => page_ring_substring1, page_ring_substring2
  procedure, private :: page_ring_substring1
  procedure, private :: page_ring_substring2  
\nwendcode{}\nwbegindocs{281}\nwdocspar
\nwenddocs{}\nwbegincode{282}\moddef{Muli base: procedures}\plusendmoddef
  pure function page_ring_substring1 (this, i) result (res)
    class(page_ring_t), intent(in) :: this
    integer(dik), dimension(2,2), intent(in) :: i
    character(ring_position_metric1(i)) :: res
    integer(dik) :: page, pos
    if (i(1,1) == i(1,2)) then
       res = this%ring (mod(i(1,1), this%ring_size)) (i(2,1):i(2,2))
    else
       pos = serialize_page_size - i(2,1)
       res(1:pos+1) = this%ring (mod(i(1,1),this%ring_size)) (i(2,1):)
       do page = i(1,1) + 1, i(1,1) - 1
          res (pos+2:pos+2+serialize_page_size) = &
               this%ring (mod(page,this%ring_size))
          pos = pos + serialize_page_size
       end do
       res(pos+2:pos+1+i(2,2)) = &
            this%ring (mod(page,this%ring_size)) (1:i(2,2))
    end if
  end function page_ring_substring1
  
\eatline
\nwendcode{}\nwbegindocs{283}\nwdocspar
\nwenddocs{}\nwbegincode{284}\moddef{Muli base: procedures}\plusendmoddef
  pure function page_ring_substring2 (this, i1, i2) result (res)
    class(page_ring_t), intent(in) :: this
    integer(dik), dimension(2), intent(in) :: i1,i2
    character(ring_position_metric2(i1,i2)) :: res
    integer(dik) :: page, pos
    if (i1(1) == i2(1)) then
       res = this%ring(mod(i1(1),this%ring_size)) (i1(2):i2(2))
    else
       pos = serialize_page_size - i1(2)
       res(1:pos+1) = this%ring(mod(i1(1),this%ring_size)) (i1(2):)
       do page = i1(1)+1, i2(1)-1
          res(pos+2:pos+2+serialize_page_size) = &
               this%ring(mod(page, this%ring_size))
          pos = pos + serialize_page_size
       end do
       res(pos+2:pos+1+i2(2)) = this%ring(mod(page, this%ring_size)) (1:i2(2))
    end if
  end function page_ring_substring2
  
\eatline
\nwendcode{}\nwbegindocs{285}\nwdocspar
\nwenddocs{}\nwbegincode{286}\moddef{Muli base: page ring: TBP}\plusendmoddef
  generic :: substring_by_keys => page_ring_character_by_keys, &
       page_ring_positions_by_keys
  procedure, private :: page_ring_character_by_keys
  procedure, private :: page_ring_positions_by_keys
\nwendcode{}\nwbegincode{287}\moddef{Muli base: procedures}\plusendmoddef
  pure recursive subroutine page_ring_character_by_keys (this, exp1, &
       exp2, start, limit, inclusive, length, string)
    class(page_ring_t), intent(in) :: this
    character(*), intent(in) :: exp1, exp2
    integer(dik), dimension(2), intent(in) :: start, limit
    logical, optional, intent(in) :: inclusive
    integer(dik), intent(out), optional :: length
    character(:), allocatable, intent(out) :: string
    integer(dik), dimension(2,2) :: pos
    call this%substring_by_keys (exp1, exp2, start, limit, &
         inclusive, length, pos)
    string = this%substring (pos(:,1),pos(:,2))
  end subroutine page_ring_character_by_keys

\eatline
\nwendcode{}\nwbegindocs{288}\nwdocspar
\nwenddocs{}\nwbegincode{289}\moddef{Muli base: procedures}\plusendmoddef
  pure recursive subroutine page_ring_positions_by_keys (this, exp1, &
       exp2, start, limit, inclusive, length, pos)
    class(page_ring_t), intent(in) :: this
    character(*), intent(in) :: exp1, exp2
    integer(dik), dimension(2), intent(in) :: start, limit
    logical, optional, intent(in) :: inclusive
    integer(dik), intent(out), optional :: length
    integer(dik), dimension(2,2), intent(out) :: pos
    if (inclusive) then
       pos(1:2,1) = this%find_pure (exp1, start, limit, 2)
    else
       pos(1:2,1) = this%find_pure (exp1,start, limit, 4)
    end if
    ! print *,pos1
    if (present(length)) then
       length = 0
    end if
    if (pos(2,1) > 0) then
       if (inclusive) then
          pos(1:2,2) = this%find_pure (exp2, pos(1:2,1), limit, 3)
       else
          pos(1:2,2) = this%find_pure (exp2, pos(1:2,1), limit, 1)
       end if
       ! print *,pos2
       if (pos(2,2) > 0) then
          if (present (length)) then
             length = ring_position_metric1 (pos)
          end if
       end if
    end if
  end subroutine page_ring_positions_by_keys
  
\eatline
\nwendcode{}\nwbegindocs{290}\nwdocspar
\nwenddocs{}\nwbegincode{291}\moddef{Muli base: page ring: TBP}\plusendmoddef
  generic :: pop => pop_string, pop_integer, pop_integer_dik, &
       pop_real, pop_logical, pop_integer_array, &
       pop_integer_array_dik, pop_real_array
\nwendcode{}\nwbegincode{292}\moddef{Muli base: page ring: TBP}\plusendmoddef
  procedure, private :: pop_string => page_ring_pop_string 
\nwendcode{}\nwbegincode{293}\moddef{Muli base: procedures}\plusendmoddef
  recursive subroutine page_ring_pop_string (this, res)
    class(page_ring_t), intent(inout) :: this
    character(len=*), intent(out) :: res
    integer(dik) :: n, cut
    n = len(res)
    cut = serialize_page_size-this%actual_offset() + 1
    if (n <= cut) then
       res = this%ring (this%actual_index()) & 
            (this%actual_offset():this%actual_offset()+n)
       if (n == cut) then
          call this%read_page
       end if
       call this%proceed (n)
    else
       call this%pop (res(:cut))
       call this%pop (res(cut+1:))
    end if
  end subroutine page_ring_pop_string

\eatline
\nwendcode{}\nwbegindocs{294}\nwdocspar
\nwenddocs{}\nwbegincode{295}\moddef{Muli base: page ring: TBP}\plusendmoddef
  procedure, private :: pop_integer => page_ring_pop_integer
\nwendcode{}\nwbegincode{296}\moddef{Muli base: procedures}\plusendmoddef
  subroutine page_ring_pop_integer (this,in)
    class(page_ring_t), intent(inout) :: this
    integer, intent(out) :: in
    integer(dik) :: in_dik
    call this%pop (in_dik)
    in = int(in_dik)
  end subroutine page_ring_pop_integer
   
\eatline
\nwendcode{}\nwbegindocs{297}\nwdocspar
\nwenddocs{}\nwbegincode{298}\moddef{Muli base: page ring: TBP}\plusendmoddef
  procedure, private :: pop_integer_dik => page_ring_pop_integer_dik
\nwendcode{}\nwbegincode{299}\moddef{Muli base: procedures}\plusendmoddef
  subroutine page_ring_pop_integer_dik (this, int)
    class(page_ring_t), intent(inout) :: this
    integer(dik), intent(out) :: int
    integer(dik) :: int1
    integer(dik) :: sign
    character :: c
    int = 0
    sign = 1
    c = " "
    do while (scan (c, serialize_integer_characters) == 0)
       call this%pop_character (c)
    end do
    if (c == "-") then
       sign = -1
       call this%pop_character (c)
    end if
    do while (scan (c, serialize_integer_characters) > 0)
       int = int * 10
       select case (c)
       case ("1")
          int = int + 1
       case ("2")
          int = int + 2
       case ("3")
          int = int + 3
       case ("4")
          int = int + 4
       case ("5")
          int = int + 5
       case ("6")
          int = int + 6
       case ("7")
          int = int + 7
       case ("8")
          int = int + 8
       case ("9")
          int = int + 9
       end select
       call this%pop_character (c)
    end do
    int = int * sign
    if (c == "<")  call this%proceed (-i_one)
  end subroutine page_ring_pop_integer_dik
  
\eatline
\nwendcode{}\nwbegindocs{300}\nwdocspar
\nwenddocs{}\nwbegincode{301}\moddef{Muli base: page ring: TBP}\plusendmoddef
  procedure, private :: pop_integer_array => page_ring_pop_integer_array
\nwendcode{}\nwbegindocs{302}\nwdocspar
\nwenddocs{}\nwbegincode{303}\moddef{Muli base: procedures}\plusendmoddef
  subroutine page_ring_pop_integer_array (this, int)
    class(page_ring_t), intent(inout) :: this
    integer, dimension(:), intent(out) :: int
    integer :: n
    do n = 1, size(int)
       call this%pop (int(n))
    end do
  end subroutine page_ring_pop_integer_array
  
\eatline
\nwendcode{}\nwbegindocs{304}\nwdocspar
\nwenddocs{}\nwbegincode{305}\moddef{Muli base: page ring: TBP}\plusendmoddef
  procedure, private :: pop_integer_array_dik => &
       page_ring_pop_integer_array_dik
\nwendcode{}\nwbegindocs{306}\nwdocspar
\nwenddocs{}\nwbegincode{307}\moddef{Muli base: procedures}\plusendmoddef
   subroutine page_ring_pop_integer_array_dik (this, int)
    class(page_ring_t), intent(inout) :: this
    integer(dik), dimension(:), intent(out) :: int
    integer(dik) :: n
    do n = 1, size(int)
       call this%pop (int(n))
    end do
  end subroutine page_ring_pop_integer_array_dik
  
\eatline
\nwendcode{}\nwbegindocs{308}\nwdocspar
\nwenddocs{}\nwbegincode{309}\moddef{Muli base: page ring: TBP}\plusendmoddef
  procedure, private :: pop_logical => page_ring_pop_logical  
\nwendcode{}\nwbegincode{310}\moddef{Muli base: procedures}\plusendmoddef
  subroutine page_ring_pop_logical (this, l)
    class(page_ring_t), intent(inout) :: this
    logical, intent(out) :: l
    character(1) :: lc
    call this%pop (lc)
    do while (scan (lc,"tTfF") == 0)
       call this%pop (lc)
    end do
    read (lc, "(L1)") l
  end subroutine page_ring_pop_logical
  
\eatline
\nwendcode{}\nwbegindocs{311}\nwdocspar
\nwenddocs{}\nwbegincode{312}\moddef{Muli base: page ring: TBP}\plusendmoddef
  procedure, private :: pop_real => page_ring_pop_real
\nwendcode{}\nwbegindocs{313}\nwdocspar
\nwenddocs{}\nwbegincode{314}\moddef{Muli base: procedures}\plusendmoddef
  subroutine page_ring_pop_real (this, def, skip)
    class(page_ring_t), intent(inout) :: this
    real(default), intent(out) :: def
    logical, optional, intent(in) :: skip
    integer(dik) :: d, f, e
    call this%pop (d)
    if (d == i_zero) then
       def = zero
    else
       call this%pop (f)
       call this%pop (e)
       def = set_exponent (scale (real(f, kind=default), -d), e)
    end if
    if (present (skip)) then
       if (.not. skip)  call this%proceed (-i_one)
    end if
  end subroutine page_ring_pop_real
  
\eatline
\nwendcode{}\nwbegindocs{315}\nwdocspar
\nwenddocs{}\nwbegincode{316}\moddef{Muli base: page ring: TBP}\plusendmoddef
  procedure, private :: pop_real_array => page_ring_pop_real_array
\nwendcode{}\nwbegindocs{317}  
\nwenddocs{}\nwbegincode{318}\moddef{Muli base: procedures}\plusendmoddef
  subroutine page_ring_pop_real_array (this, def, skip)
    class(page_ring_t), intent(inout) :: this
    real(default), dimension(:), intent(out) :: def
    logical, optional, intent(in) :: skip
    integer(dik) :: n    
    call this%pop_real (def(1))
    do n = 2, size(def)
       call this%pop_real (def(n))
    end do
    if (present(skip)) then
       if (.not. skip)  call this%proceed (-i_one)
    end if
  end subroutine page_ring_pop_real_array
  
\eatline
\nwendcode{}\nwbegindocs{319}\nwdocspar
\nwenddocs{}\nwbegincode{320}\moddef{Muli base: page ring: TBP}\plusendmoddef
  procedure :: close => page_ring_close
\nwendcode{}\nwbegincode{321}\moddef{Muli base: procedures}\plusendmoddef
  subroutine page_ring_close (this)
    class(page_ring_t), intent(inout) :: this
    if (this%action == 1) then
       call this%flush ()
       ! call this%print_position()
       if (this%asynchronous) then
          write (this%unit, asynchronous="yes") &
               this%ring (this%actual_index()) (:this%actual_offset() - 1)
       else
          write (this%unit, asynchronous="no") &
               this%ring (this%actual_index()) (:this%actual_offset() - 1)
       end if
    end if
    close (this%unit)
  end subroutine page_ring_close
  
\eatline
\nwendcode{}\nwbegindocs{322}\nwdocspar
\nwenddocs{}\nwbegincode{323}\moddef{Muli base: page ring: TBP}\plusendmoddef
  procedure :: ring_index => page_ring_ring_index 
\nwendcode{}\nwbegincode{324}\moddef{Muli base: procedures}\plusendmoddef
  elemental integer(dik) function page_ring_ring_index (this, n)
    class(page_ring_t), intent(in) :: this
    integer(dik), intent(in) :: n
    page_ring_ring_index = mod(n, this%ring_size)
  end function page_ring_ring_index
  
\eatline
\nwendcode{}\nwbegindocs{325}\nwdocspar
\nwenddocs{}\nwbegincode{326}\moddef{Muli base: page ring: TBP}\plusendmoddef
  procedure, private :: activate_next_page => page_ring_activate_next_page  
\nwendcode{}\nwbegincode{327}\moddef{Muli base: procedures}\plusendmoddef
  subroutine page_ring_activate_next_page (this)
    class(page_ring_t), intent(inout) :: this
    if (this%active_pages(2) - this%active_pages(1) + 1 >= &
         this%ring_size)  call this%enlarge
    this%active_pages(2) = this%active_pages(2) + 1
  end subroutine page_ring_activate_next_page

\eatline
\nwendcode{}\nwbegindocs{328}\nwdocspar
\nwenddocs{}\nwbegincode{329}\moddef{Muli base: page ring: TBP}\plusendmoddef
  procedure, private :: enlarge => page_ring_enlarge  
\nwendcode{}\nwbegincode{330}\moddef{Muli base: procedures}\plusendmoddef
  subroutine page_ring_enlarge (this)
    class(page_ring_t), intent(inout) :: this
    character(serialize_page_size), dimension(:), allocatable :: tmp_ring
    integer(dik) :: n
    call move_alloc (this%ring, tmp_ring)
    allocate (this%ring(0:this%ring_size*2-1))
    do n = this%active_pages(1), this%active_pages(2)
       this%ring (mod(n,this%ring_size*2)) = tmp_ring (mod(n,this%ring_size))
    end do
    this%ring_size = this%ring_size * 2
  end subroutine page_ring_enlarge

\eatline
\nwendcode{}\nwbegindocs{331}These are specific implementations of generic procedures:
\nwenddocs{}\nwbegincode{332}\moddef{Muli base: page ring: TBP}\plusendmoddef
  procedure, private :: actual_index => page_ring_actual_index
\nwendcode{}\nwbegincode{333}\moddef{Muli base: procedures}\plusendmoddef
  elemental integer(dik) function page_ring_actual_index (this)
    class(page_ring_t), intent(in) :: this
    page_ring_actual_index = &
         mod (this%position_stack%position(1), this%ring_size)
  end function page_ring_actual_index

\eatline
\nwendcode{}\nwbegindocs{334}\nwdocspar
\nwenddocs{}\nwbegincode{335}\moddef{Muli base: page ring: TBP}\plusendmoddef
  procedure, private :: actual_page => page_ring_actual_page 
\nwendcode{}\nwbegincode{336}\moddef{Muli base: procedures}\plusendmoddef
  elemental integer(dik) function page_ring_actual_page (this)
    class(page_ring_t), intent(in) :: this
    page_ring_actual_page = this%position_stack%position(1)
  end function page_ring_actual_page
  
\eatline
\nwendcode{}\nwbegindocs{337}\nwdocspar
\nwenddocs{}\nwbegincode{338}\moddef{Muli base: page ring: TBP}\plusendmoddef
  procedure, private :: actual_offset => page_ring_actual_offset
\nwendcode{}\nwbegincode{339}\moddef{Muli base: procedures}\plusendmoddef
  elemental integer(kind=dik) function page_ring_actual_offset(this)
    class(page_ring_t),intent(in) :: this
    page_ring_actual_offset=this%position_stack%position(2)
  end function page_ring_actual_offset
  
\eatline
\nwendcode{}\nwbegindocs{340}\nwdocspar
\nwenddocs{}\nwbegincode{341}\moddef{Muli base: page ring: TBP}\plusendmoddef
  procedure, private :: actual_position => page_ring_actual_position
\nwendcode{}\nwbegincode{342}\moddef{Muli base: procedures}\plusendmoddef
  pure function page_ring_actual_position(this)
    class(page_ring_t), intent(in) :: this
    integer(dik), dimension(2) :: page_ring_actual_position
    page_ring_actual_position = this%position_stack%position
  end function page_ring_actual_position
   
\eatline
\nwendcode{}\nwbegindocs{343}\nwdocspar
\nwenddocs{}\nwbegincode{344}\moddef{Muli base: page ring: TBP}\plusendmoddef
  procedure, private :: first_index => page_ring_first_index
\nwendcode{}\nwbegincode{345}\moddef{Muli base: procedures}\plusendmoddef
  elemental integer(dik) function page_ring_first_index (this)
    class(page_ring_t), intent(in) :: this
    page_ring_first_index = mod(this%active_pages(1), this%ring_size)
  end function page_ring_first_index
  
\eatline
\nwendcode{}\nwbegindocs{346}  
\nwenddocs{}\nwbegincode{347}\moddef{Muli base: page ring: TBP}\plusendmoddef
  procedure, private :: first_page => page_ring_first_page
\nwendcode{}\nwbegincode{348}\moddef{Muli base: procedures}\plusendmoddef
  elemental integer(dik) function page_ring_first_page (this)
    class(page_ring_t), intent(in) :: this
    page_ring_first_page = this%active_pages(1)
  end function page_ring_first_page
  
\eatline
\nwendcode{}\nwbegindocs{349}  
\nwenddocs{}\nwbegincode{350}\moddef{Muli base: page ring: TBP}\plusendmoddef
  procedure, private :: last_index => page_ring_last_index
\nwendcode{}\nwbegincode{351}\moddef{Muli base: procedures}\plusendmoddef
  elemental integer(dik) function page_ring_last_index (this)
    class(page_ring_t), intent(in) :: this
    page_ring_last_index = mod(this%active_pages(2), this%ring_size)
  end function page_ring_last_index

\eatline
\nwendcode{}\nwbegindocs{352}  
\nwenddocs{}\nwbegincode{353}\moddef{Muli base: page ring: TBP}\plusendmoddef
  procedure, private :: last_page => page_ring_last_page
\nwendcode{}\nwbegincode{354}\moddef{Muli base: procedures}\plusendmoddef
  elemental integer(dik) function page_ring_last_page (this)
    class(page_ring_t), intent(in) :: this
    page_ring_last_page = this%active_pages(2)
  end function page_ring_last_page
    
\eatline
\nwendcode{}\nwbegindocs{355}\nwdocspar
\nwenddocs{}\nwbegincode{356}\moddef{Muli base: public}\plusendmoddef
  public :: marker_t
\nwendcode{}\nwbegincode{357}\moddef{Muli base: types}\plusendmoddef
  type, extends (page_ring_t) :: marker_t
     private
     integer(dik) :: indentation=0
     integer(dik) :: n_instances=0
     logical :: do_break=.true.
     logical :: do_indent=.false.
     class(serializable_ref_type),pointer :: heap=>null()
     class(serializable_ref_type),pointer :: references=>null()
   contains
   \LA{}Muli base: marker: TBP\RA{}  
  end type marker_t

\eatline
\nwendcode{}\nwbegindocs{358}\nwdocspar
\nwenddocs{}\nwbegincode{359}\moddef{Muli base: marker: TBP}\endmoddef
  procedure :: mark_begin => marker_mark_begin
\nwendcode{}\nwbegincode{360}\moddef{Muli base: procedures}\plusendmoddef
  subroutine marker_mark_begin (this, tag, type, name, target, pointer, shape)
    class(marker_t), intent(inout) :: this
    character(*), intent(in) :: tag
    character(*), intent(in), optional :: type, name
    integer(kind=dik), intent(in), optional :: target, pointer
    integer,intent(in), dimension(:), optional :: shape
    call this%indent ()
    call this%push ("<")
    call this%push (tag)
    if (present (type))  call this%push (' type="'//type//'"')
    if (present (name))  call this%push (' name="'//name//'"')
    if (present (target)) then
       call this%push (' target="')
       call this%push (target)
       call this%push ('"')
    end if
    if (present (pointer))then
       call this%push (' pointer="')
       call this%push (pointer)
       call this%push ('"')
    end if
    if (present (shape))then
       call this%push (' shape="')
       call this%push (shape)
       call this%push ('"')
    end if
    call this%push (">")
    this%indentation = this%indentation + 1
  end subroutine marker_mark_begin
  
\eatline
\nwendcode{}\nwbegindocs{361}\nwdocspar
\nwenddocs{}\nwbegincode{362}\moddef{Muli base: marker: TBP}\plusendmoddef
  procedure :: mark_instance_begin => marker_mark_instance_begin
\nwendcode{}\nwbegincode{363}\moddef{Muli base: procedures}\plusendmoddef
  subroutine marker_mark_instance_begin &
       (this, ser, name, target, pointer, shape)
    class(marker_t), intent(inout) :: this
    class(ser_class_t), intent(in) :: ser
    character(*), intent(in) :: name
    integer(dik), intent(in), optional :: target, pointer
    integer, dimension(:), intent(in), optional :: shape
    character(:), allocatable :: this_type
    call ser%get_type (this_type)
    call this%mark_begin ("ser", this_type, name, target, pointer, shape)
  end subroutine marker_mark_instance_begin
  
\eatline
\nwendcode{}\nwbegindocs{364}\nwdocspar
\nwenddocs{}\nwbegincode{365}\moddef{Muli base: marker: TBP}\plusendmoddef
  procedure :: mark_end => marker_mark_end
\nwendcode{}\nwbegincode{366}\moddef{Muli base: procedures}\plusendmoddef
  subroutine marker_mark_end (this, tag)
    class(marker_t), intent(inout) :: this
    character(*), intent(in), optional :: tag
    this%indentation = this%indentation - 1
    call this%indent ()
    if (present (tag)) then
       call this%push ("</"//tag//">")
    else
       call this%push ("</ser>")
    end if
  end subroutine marker_mark_end

\eatline
\nwendcode{}\nwbegindocs{367}\nwdocspar
\nwenddocs{}\nwbegincode{368}\moddef{Muli base: marker: TBP}\plusendmoddef
  procedure :: mark_instance_end => marker_mark_instance_end
\nwendcode{}\nwbegincode{369}\moddef{Muli base: procedures}\plusendmoddef
  subroutine marker_mark_instance_end (this)
    class(marker_t), intent(inout) :: this
    call this%mark_end ("ser")
  end subroutine marker_mark_instance_end

\eatline
\nwendcode{}\nwbegindocs{370}\nwdocspar
\nwenddocs{}\nwbegincode{371}\moddef{Muli base: marker: TBP}\plusendmoddef
  generic :: mark => mark_logical, &
       mark_integer, mark_integer_array, mark_integer_matrix, &
       mark_integer_dik, mark_integer_array_dik, mark_integer_matrix_dik, &
       mark_default, mark_default_array, mark_default_matrix, mark_string
  procedure, private :: mark_logical => marker_mark_logical
\nwendcode{}\nwbegincode{372}\moddef{Muli base: procedures}\plusendmoddef
  subroutine marker_mark_logical (this, name, content)
    class(marker_t), intent(inout) :: this
    character(*), intent(in) :: name
    logical, intent(in) :: content
    call this%indent ()
    call this%push ("<"//name//">")
    if (content) then
       call this%push ("T")
    else
       call this%push ("F")
    end if
    call this%push ("</"//name//">")
  end subroutine marker_mark_logical
  
\eatline
\nwendcode{}\nwbegindocs{373}\nwdocspar
\nwenddocs{}\nwbegincode{374}\moddef{Muli base: marker: TBP}\plusendmoddef
  procedure :: mark_integer => marker_mark_integer
\nwendcode{}\nwbegincode{375}\moddef{Muli base: procedures}\plusendmoddef
  subroutine marker_mark_integer (this, name, content)
    class(marker_t), intent(inout) :: this
    character(*), intent(in) :: name
    integer, intent(in) :: content
    call this%indent ()
    call this%push ("<"//name//">")
    call this%push (content)
    call this%push ("</"//name//">")
  end subroutine marker_mark_integer

\eatline
\nwendcode{}\nwbegindocs{376}\nwdocspar
\nwenddocs{}\nwbegincode{377}\moddef{Muli base: marker: TBP}\plusendmoddef
  procedure :: mark_integer_array => marker_mark_integer_array
\nwendcode{}\nwbegincode{378}\moddef{Muli base: procedures}\plusendmoddef
  subroutine marker_mark_integer_array (this, name, content)
    class(marker_t), intent(inout) :: this
    character(*), intent(in) :: name
    integer, dimension(:), intent(in) :: content
    call this%indent ()
    call this%push ("<"//name//">")
    call this%push (content)
    call this%push ("</"//name//">")
  end subroutine marker_mark_integer_array

\eatline
\nwendcode{}\nwbegindocs{379} 
\nwenddocs{}\nwbegincode{380}\moddef{Muli base: marker: TBP}\plusendmoddef
  procedure :: mark_integer_matrix => marker_mark_integer_matrix
\nwendcode{}\nwbegincode{381}\moddef{Muli base: procedures}\plusendmoddef
  
  subroutine marker_mark_integer_matrix (this, name, content)
    class(marker_t), intent(inout) :: this
    character(*), intent(in) :: name
    integer, dimension(:,:), intent(in) :: content
    integer :: n
    integer, dimension(2) :: s
    s= shape(content)
    call this%indent ()
    call this%push ("<"//name//">")
    do n = 1, s(2)
       call this%push (content(:,n))
       call this%push (" ")
    end do
    call this%push ("</"//name//">")
  end subroutine marker_mark_integer_matrix

\nwendcode{}\nwbegindocs{382} %def marker_mark_integer_matrix
\nwenddocs{}\nwbegindocs{383}\nwdocspar
\nwenddocs{}\nwbegincode{384}\moddef{Muli base: marker: TBP}\plusendmoddef
  procedure :: mark_integer_dik => marker_mark_integer_dik
\nwendcode{}\nwbegincode{385}\moddef{Muli base: procedures}\plusendmoddef
  subroutine marker_mark_integer_dik (this, name, content)
    class(marker_t), intent(inout) :: this
    character(*), intent(in) :: name
    integer(dik), intent(in) :: content
    call this%indent ()
    call this%push ("<"//name//">")
    call this%push (content)
    call this%push ("</"//name//">")
  end subroutine marker_mark_integer_dik

\eatline
\nwendcode{}\nwbegindocs{386} 
\nwenddocs{}\nwbegincode{387}\moddef{Muli base: marker: TBP}\plusendmoddef
  procedure :: mark_integer_array_dik => marker_mark_integer_array_dik
\nwendcode{}\nwbegincode{388}\moddef{Muli base: procedures}\plusendmoddef
  subroutine marker_mark_integer_array_dik (this, name, content)
    class(marker_t), intent(inout) :: this
    character(*), intent(in) :: name
    integer(dik), dimension(:), intent(in) :: content
    call this%indent ()
    call this%push ("<"//name//">")
    call this%push (content)
    call this%push ("</"//name//">")
  end subroutine marker_mark_integer_array_dik

\eatline
\nwendcode{}\nwbegindocs{389}\nwdocspar
\nwenddocs{}\nwbegincode{390}\moddef{Muli base: marker: TBP}\plusendmoddef
  procedure :: mark_integer_matrix_dik => marker_mark_integer_matrix_dik
\nwendcode{}\nwbegincode{391}\moddef{Muli base: procedures}\plusendmoddef
  subroutine marker_mark_integer_matrix_dik (this, name, content)
    class(marker_t), intent(inout) :: this
    character(*), intent(in) :: name
    integer(dik), dimension(:,:), intent(in) :: content
    integer :: n
    integer, dimension(2) :: s
    call this%indent ()
    call this%push ("<"//name//">")
    do n = 1, s(2)
       call this%push (content(:,n))
       call this%push (" ")
    end do
    call this%push ("</"//name//">")
  end subroutine marker_mark_integer_matrix_dik
  
\eatline
\nwendcode{}\nwbegindocs{392}\nwdocspar
\nwenddocs{}\nwbegincode{393}\moddef{Muli base: marker: TBP}\plusendmoddef
  procedure :: mark_default => marker_mark_default
\nwendcode{}\nwbegincode{394}\moddef{Muli base: procedures}\plusendmoddef
  subroutine marker_mark_default (this, name, content)
    class(marker_t), intent(inout) :: this
    character(*), intent(in) :: name
    real(default), intent(in) :: content
    call this%indent ()
    call this%push ("<"//name//">")
    call this%push (content)
    call this%push ("</"//name//">")
  end subroutine marker_mark_default

\eatline
\nwendcode{}\nwbegindocs{395}\nwdocspar
\nwenddocs{}\nwbegincode{396}\moddef{Muli base: marker: TBP}\plusendmoddef
  procedure :: mark_default_array => marker_mark_default_array
\nwendcode{}\nwbegincode{397}\moddef{Muli base: procedures}\plusendmoddef
  subroutine marker_mark_default_array (this, name, content)
    class(marker_t), intent(inout) :: this
    character(*), intent(in) :: name
    real(default), dimension(:), intent(in) :: content
    call this%indent ()
    call this%push ("<"//name//">")
    call this%push (content)
    call this%push ("</"//name//">")
  end subroutine marker_mark_default_array

\eatline
\nwendcode{}\nwbegindocs{398}\nwdocspar
\nwenddocs{}\nwbegincode{399}\moddef{Muli base: marker: TBP}\plusendmoddef
  procedure :: mark_default_matrix => marker_mark_default_matrix
\nwendcode{}\nwbegincode{400}\moddef{Muli base: procedures}\plusendmoddef
  subroutine marker_mark_default_matrix (this, name, content)
    class(marker_t), intent(inout) :: this
    character(*), intent(in) :: name
    real(default), dimension(:,:), intent(in) :: content
    integer :: n
    integer, dimension(2) :: s
    s = shape(content)
    call this%indent ()
    call this%push ("<"//name//">")
    do n = 1, s(2)
       call this%push (content(:,n))
       call this%push (" ")
    end do
    call this%push ("</"//name//">")
  end subroutine marker_mark_default_matrix

\nwendcode{}\nwbegindocs{401} %def marker_mark_default_matrix
\nwenddocs{}\nwbegindocs{402}\nwdocspar
\nwenddocs{}\nwbegincode{403}\moddef{Muli base: marker: TBP}\plusendmoddef
  procedure :: mark_string => marker_mark_string
\nwendcode{}\nwbegincode{404}\moddef{Muli base: procedures}\plusendmoddef
  subroutine marker_mark_string (this, name, content)
    class(marker_t), intent(inout) :: this
    character(*), intent(in) :: name, content
    call this%indent ()
    call this%push ("<"//name//">")
    call this%push (content)
    call this%push ("</"//name//">")
  end subroutine marker_mark_string

\eatline
\nwendcode{}\nwbegindocs{405}\nwdocspar
\nwenddocs{}\nwbegincode{406}\moddef{Muli base: marker: TBP}\plusendmoddef
  procedure :: mark_instance => marker_mark_instance
\nwendcode{}\nwbegincode{407}\moddef{Muli base: procedures}\plusendmoddef
  recursive subroutine marker_mark_instance (this, ser, name, target, pointer)
    class(marker_t), intent(inout) :: this
    class(ser_class_t), intent(in) :: ser
    character(len=*), intent(in) :: name
    integer(dik), intent(in), optional :: target, pointer
    integer(dik) :: status
    call this%mark_instance_begin (ser, name, target, pointer)
    call ser%write_to_marker (this, status)
    call this%mark_end ("ser")
  end subroutine marker_mark_instance

\eatline
\nwendcode{}\nwbegindocs{408}\nwdocspar
\nwenddocs{}\nwbegincode{409}\moddef{Muli base: marker: TBP}\plusendmoddef
  procedure :: mark_target => marker_mark_target
\nwendcode{}\nwbegincode{410}\moddef{Muli base: procedures}\plusendmoddef
  recursive subroutine marker_mark_target (this, name, ser)
    class(marker_t), intent(inout) :: this
    class(ser_class_t), target, intent(in) :: ser
    character(len=*), intent(in) :: name
    this%n_instances = this%n_instances + 1
    call this%push_heap (ser, this%n_instances)
    call this%mark_instance (ser, name, target = this%n_instances)
  end subroutine marker_mark_target

\eatline
\nwendcode{}\nwbegindocs{411}\nwdocspar
\nwenddocs{}\nwbegincode{412}\moddef{Muli base: marker: TBP}\plusendmoddef
  procedure :: mark_allocatable => marker_mark_allocatable
\nwendcode{}\nwbegincode{413}\moddef{Muli base: procedures}\plusendmoddef
  subroutine marker_mark_allocatable (this, name, ser)
    class(marker_t), intent(inout) :: this
    class(ser_class_t), allocatable, intent(in) :: ser
    character(len=*), intent(in) :: name
    if (allocated (ser)) then
       call this%mark_instance (ser, name)
    else
       call this%mark_null (name)
    end if
  end subroutine marker_mark_allocatable

\eatline
\nwendcode{}\nwbegindocs{414}\nwdocspar
\nwenddocs{}\nwbegincode{415}\moddef{Muli base: marker: TBP}\plusendmoddef
  procedure :: mark_pointer => marker_mark_pointer
\nwendcode{}\nwbegincode{416}\moddef{Muli base: procedures}\plusendmoddef
  recursive subroutine marker_mark_pointer (this, name, ser)
    class(marker_t), intent(inout) :: this
    class(ser_class_t), pointer, intent(in) :: ser
    character(len=*), intent(in) :: name
    character(:), allocatable :: type
    integer(dik) :: p
    if (associated (ser)) then
       call this%search_heap (ser, p)
       if (p > 0) then
          call ser%get_type (type)
          call this%push ('<ser type="')
          call this%push (type)
          call this%push ('" name="')
          call this%push (name)
          call this%push ('" pointer="')
          call this%push (p)
          call this%push ('"/>')
       else
          call this%mark_target (name, ser)
       end if
    else
       call this%mark_null (name)
    end if
  end subroutine marker_mark_pointer

\eatline
\nwendcode{}\nwbegindocs{417}\nwdocspar
\nwenddocs{}\nwbegincode{418}\moddef{Muli base: marker: TBP}\plusendmoddef
  procedure :: mark_null => marker_mark_null
\nwendcode{}\nwbegincode{419}\moddef{Muli base: procedures}\plusendmoddef
  subroutine marker_mark_null (this, name)
    class(marker_t), intent(inout) :: this
    character(*), intent(in) :: name
    call this%indent ()
    call this%push ('<ser type="null" name="')
    call this%push (name)
    call this%push ('"/>')
  end subroutine marker_mark_null

\eatline
\nwendcode{}\nwbegindocs{420}\nwdocspar
\nwenddocs{}\nwbegincode{421}\moddef{Muli base: marker: TBP}\plusendmoddef
  procedure :: mark_nothing => marker_mark_nothing
\nwendcode{}\nwbegincode{422}\moddef{Muli base: procedures}\plusendmoddef
  subroutine marker_mark_nothing (this, name)
    class(marker_t), intent(inout) :: this
    character(*), intent(in) :: name
    call this%indent ()
    call this%push ('<')
    call this%push (name)
    call this%push ('/>')
  end subroutine marker_mark_nothing

\eatline
\nwendcode{}\nwbegindocs{423}\nwdocspar
\nwenddocs{}\nwbegincode{424}\moddef{Muli base: marker: TBP}\plusendmoddef
  procedure :: mark_empty => marker_mark_empty
\nwendcode{}\nwbegincode{425}\moddef{Muli base: procedures}\plusendmoddef
  subroutine marker_mark_empty (this, tag, type, name, target, pointer, shape)
    class(marker_t), intent(inout) :: this
    character(*), intent(in) :: tag
    character(*), intent(in), optional :: type, name
    integer(dik), intent(in), optional :: target, pointer
    integer, dimension(:), intent(in), optional :: shape
    call this%push ("<")
    call this%push (tag)
    if (present (type))  call this%push (' type="'//type//'"')
    if (present (name))  call this%push (' name="'//name//'"')
    if (present (target)) then
       call this%push (' target="')
       call this%push (target)
       call this%push ('"')
    end if
    if (present (pointer)) then
       call this%push (' pointer="')
       call this%push (pointer)
       call this%push ('"')
    end if
    if (present (shape)) then
       call this%push (' shape="')
       call this%push (shape)
       call this%push ('"')
    end if
    call this%push ("/>")
  end subroutine marker_mark_empty
  
\eatline
\nwendcode{}\nwbegindocs{426}\nwdocspar
\nwenddocs{}\nwbegincode{427}\moddef{Muli base: marker: TBP}\plusendmoddef
  procedure :: pick_begin => marker_pick_begin
\nwendcode{}\nwbegincode{428}\moddef{Muli base: procedures}\plusendmoddef
  subroutine marker_pick_begin (this, tag, type, name, target, &
       pointer, shape, status)
    class(marker_t), intent(inout) :: this
    character(*), intent(in) :: tag
    integer(dik), dimension(2,2),intent(out),optional :: type,name
    integer(dik), intent(out), optional :: target, pointer
    integer, dimension(:), allocatable, optional, intent(out) :: shape
    integer(dik), intent(out) :: status
    integer(dik), dimension(2) :: p1, p2, p3
    integer(dik) :: l
    call this%find ("<", skip=4, proceed=.true., pos=p1)
    call this%find (">", skip=1, proceed=.false., pos=p2)
    p3 = this%find_pure (" ",p1,p2,skip=1)
    if (p3(2) > 0) then
       if (this%substring(p1, p3) == tag) then
          status = serialize_ok
          if (present (type)) then
             call this%substring_by_keys &
                  ('type="','"', p3, p2, .false., l, type)
             if (l <= 0) then
                call msg_error ("marker_pick_begin: No type found")
                status = serialize_wrong_type
             end if
          end if
          if (present (name)) then
             call this%substring_by_keys &
                  ('name="','"', p3, p2, .false., l, name)
             if (l <= 0) then
                call msg_error ("marker_pick_begin: No name found")
                status = serialize_wrong_name
                call this%print_position ()
                stop
             end if
          end if
          if (present (target)) then
             p1 = this%find_pure ('target="', p3, p2, 4)
             if (p1(2) > 0) then
                call this%set_position (p1)
                call this%pop (target)
             else
                target = -1
                status = serialize_ok
             end if
          end if
          if (present (pointer)) then
             p1=this%find_pure ('pointer="', p3, p2, 4)
             if (p1(2) > 0)then
                call this%set_position (p1)
                call this%pop (pointer)
             else
                pointer = -1
                status = serialize_ok
             end if
          end if
          if (present (shape)) then
             p1 = this%find_pure ('shape="', p3, p2, 4)
             if (p1(2) > 0) then
                call this%set_position (p1)
                call this%pop (shape)
             else
                status = serialize_ok
             end if
          end if
       else
          call msg_error ("marker_pick_begin: Wrong tag. Expected: " // &
               tag // " Found: " // this%substring(p1, p3))
          status = serialize_wrong_tag
          call this%print_position ()
       end if
    else
       if (this%substring(p1, p2) == tag) then
          status = serialize_ok
       else
          call msg_error ("marker_pick_begin: Wrong tag. Expected: " // &
               tag // " Found: " // this%substring(p1, p2))
          status = serialize_wrong_tag
       end if
    end if
    call this%set_position (p2)
    call this%proceed (i_one*2, .true.)
  end subroutine marker_pick_begin

\eatline
\nwendcode{}\nwbegindocs{429}\nwdocspar
\nwenddocs{}\nwbegincode{430}\moddef{Muli base: marker: TBP}\plusendmoddef
  procedure :: query_instance_begin => marker_query_instance_begin
\nwendcode{}\nwbegincode{431}\moddef{Muli base: procedures}\plusendmoddef
  subroutine marker_query_instance_begin &
       (this, type, name, target, pointer, shape,status)
    class(marker_t), intent(inout) :: this
    integer(dik), dimension(2,2), intent(out), optional :: type, name
    integer(dik), intent(out), optional :: target, pointer
    integer, dimension(:), allocatable, intent(out), optional :: shape
    integer(dik), intent(out) :: status
    call this%pick_begin ("ser", type, name, target, pointer, shape, status)
  end subroutine marker_query_instance_begin

\eatline
\nwendcode{}\nwbegindocs{432}\nwdocspar
\nwenddocs{}\nwbegincode{433}\moddef{Muli base: marker: TBP}\plusendmoddef
  procedure :: pick_instance_begin => marker_pick_instance_begin
\nwendcode{}\nwbegincode{434}\moddef{Muli base: procedures}\plusendmoddef
  subroutine marker_pick_instance_begin &
       (this, name, type, target, pointer, shape, status)
    class(marker_t), intent(inout) :: this
    character(*), intent(in) :: name
    integer(dik), dimension(2,2), intent(out), optional :: type
    integer(dik), intent(out), optional :: target,pointer
    integer, dimension(:), allocatable,intent(out), optional :: shape
    integer(dik), intent(out) :: status
    integer(dik), dimension(2,2) :: read_name
    call this%query_instance_begin &
         (type, read_name, target, pointer, shape, status)
    if (status == serialize_ok) then
       if (.not. this%str_equal (name, read_name)) &
            status = serialize_wrong_name
    end if
  end subroutine marker_pick_instance_begin

\eatline
\nwendcode{}\nwbegindocs{435}\nwdocspar
\nwenddocs{}\nwbegincode{436}\moddef{Muli base: marker: TBP}\plusendmoddef
  procedure :: pick_end => marker_pick_end
\nwendcode{}\nwbegincode{437}\moddef{Muli base: procedures}\plusendmoddef
  subroutine marker_pick_end (this, tag, status)
    class(marker_t), intent(inout) :: this
    character(*), intent(in) :: tag
    integer(dik), intent(out) :: status
    integer(dik), dimension(2) :: p1, p2
    call this%find ("</", skip=4, proceed=.true., pos=p1)
    call this%find (">", skip=1, proceed=.false., pos=p2)
    if (tag == this%substring (p1, p2)) then
       status = serialize_ok
    else
       call msg_error ("marker_pick_end: Wrong tag. Expected: " // tag &
            // " Found: " // this%substring (p1, p2))
       ! print *,"p1=",p1,"p2=",p2
       call this%print_position ()
    end if
    call this%set_position (p2)
    call this%proceed (i_one*2, .true.)
  end subroutine marker_pick_end

\eatline
\nwendcode{}\nwbegindocs{438}\nwdocspar
\nwenddocs{}\nwbegincode{439}\moddef{Muli base: marker: TBP}\plusendmoddef
  procedure :: pick_instance_end => marker_pick_instance_end
\nwendcode{}\nwbegincode{440}\moddef{Muli base: procedures}\plusendmoddef
  subroutine marker_pick_instance_end (this, status)
    class(marker_t), intent(inout) :: this
    integer(dik), intent(out) :: status
    call this%pick_end ("ser",status)
  end subroutine marker_pick_instance_end

\eatline
\nwendcode{}\nwbegindocs{441}\nwdocspar
\nwenddocs{}\nwbegincode{442}\moddef{Muli base: marker: TBP}\plusendmoddef
  procedure :: pick_instance => marker_pick_instance
\nwendcode{}\nwbegincode{443}\moddef{Muli base: procedures}\plusendmoddef
  subroutine marker_pick_instance (this, name, ser, status)
    class(marker_t), intent(inout) :: this
    class(ser_class_t), intent(out) :: ser
    character(*), intent(in) :: name
    integer(dik), intent(out) :: status
    integer(dik), dimension(2,2) :: type, r_name
    call this%pick_begin ("ser", type, r_name, status=status)
    if (status == serialize_ok) then
       if (ser%verify_type (this%substring(type))) then
          if (this%str_equal (name, r_name)) then
             call ser%read_from_marker (this, status)
             call this%pick_end ("ser", status)
          else
             call msg_error ("marker_pick_instance: Name mismatch")
             write (*,*) "Expected: ", name, " Found: ", r_name
             status = serialize_wrong_name
             call this%print_position
          end if
       else
          call msg_error ("marker_pick_instance: Type mismatch: ")
          write (*,*) type
          call ser%write_type (output_unit)
          write (*,*)
          status = serialize_wrong_type
          call this%print_position
       end if
    end if
  end subroutine marker_pick_instance

\eatline
\nwendcode{}\nwbegindocs{444}\nwdocspar
\nwenddocs{}\nwbegincode{445}\moddef{Muli base: marker: TBP}\plusendmoddef
  procedure :: pick_target => marker_pick_target
\nwendcode{}\nwbegincode{446}\moddef{Muli base: procedures}\plusendmoddef
  subroutine marker_pick_target (this, name, ser, status)
    class(marker_t), intent(inout) :: this
    class(ser_class_t), target, intent(out) :: ser
    character(*), intent(in) :: name
    integer(dik), intent(out) :: status
    integer(dik), dimension(2,2) :: type, r_name
    integer(dik) :: target
    call this%pick_begin ("ser", type, r_name, target, status=status)
    if (status == serialize_ok) then
       if (ser%verify_type (this%substring(type))) then
          if (this%str_equal (name, r_name)) then
             call ser%read_target_from_marker (this, status)
             if (target > 0)  call this%push_heap (ser, target)
          else
             call msg_error ("marker_pick_instance: Name mismatch: ")
             write (*,*) "Expected: ", name, " Found: ", r_name
             status = serialize_wrong_name
          end if
       else
          call msg_error ("marker_pick_instance: Type mismatch: ")
          write (*,*)  type
          status = serialize_wrong_type
       end if
    end if
    call this%pick_end ("ser", status)
  end subroutine marker_pick_target

\eatline
\nwendcode{}\nwbegindocs{447}\nwdocspar
\nwenddocs{}\nwbegincode{448}\moddef{Muli base: marker: TBP}\plusendmoddef
  procedure :: pick_allocatable => marker_pick_allocatable
\nwendcode{}\nwbegincode{449}\moddef{Muli base: procedures}\plusendmoddef
  subroutine marker_pick_allocatable (this, name, ser)
    class(marker_t), intent(inout) :: this
    character(*), intent(in) :: name
    class(ser_class_t), allocatable, intent(out) :: ser
    class(ser_class_t), pointer :: ref
    integer(dik),dimension(2,2) :: type, r_name
    integer(dik) :: status
    call this%pick_begin ("ser", type, r_name, status=status)
    if (status == serialize_ok) then
       if (ser%verify_type (this%substring(type))) then
          if (this%str_equal (name, r_name)) then
             call this%search_reference (type, ref)
             if (associated (ref)) then
                allocate (ser, source=ref)
                call ser%read_from_marker (this, status)
             else
                call msg_error ("marker_pick_allocatable:")
                write (*,*) "Type ", type, " not found on reference stack."
             end if
          else
             call msg_error ("marker_pick_instance: Name mismatch: ")
             write (*,*)  "Expected: ",name," Found: ",r_name
             status = serialize_wrong_name
          end if
       else
          call msg_error ("marker_pick_instance: Type mismatch: ")
          write (*,*)  type
          status = serialize_wrong_type
       end if
    end if
    call this%pick_end ("ser", status)
  end subroutine marker_pick_allocatable

\eatline
\nwendcode{}\nwbegindocs{450}\nwdocspar
\nwenddocs{}\nwbegincode{451}\moddef{Muli base: marker: TBP}\plusendmoddef
  procedure :: pick_pointer => marker_pick_pointer
\nwendcode{}\nwbegincode{452}\moddef{Muli base: procedures}\plusendmoddef
  recursive subroutine marker_pick_pointer (this, name, ser)
    class(marker_t), intent(inout) :: this
    character(*), intent(in) :: name
    class(ser_class_t), pointer, intent(out) :: ser
    class(ser_class_t), pointer :: ref
    integer(dik), dimension(2,2) :: type, r_name
    integer(dik) :: status, t, p
    nullify (ser)
    call this%pick_begin &
         ("ser", type, r_name, target=t, pointer=p, status=status)
    if (status == serialize_ok) then
       if (.not. this%str_equal ("null",type)) then
          if (p > 0) then
             call this%search_heap (p, ser)
          else
             call this%search_reference (type, ref)
             if (associated (ref))then
                allocate (ser, source=ref)
                call ser%read_target_from_marker (this, status)
                call this%pick_end ("ser", status)
                if (t > 0)  call this%push_heap (ser, t)
             else
                write (*,*) "marker_pick_pointer:&
                     & Type ",type," not found on reference stack."
             end if
          end if
       end if
    end if
  end subroutine marker_pick_pointer

\eatline
\nwendcode{}\nwbegindocs{453}\nwdocspar
\nwenddocs{}\nwbegincode{454}\moddef{Muli base: marker: TBP}\plusendmoddef
  generic :: pick => pick_logical, &
       pick_integer, pick_integer_array, pick_integer_matrix, &
       pick_integer_dik, pick_integer_array_dik, pick_integer_matrix_dik, &
       pick_default, pick_default_array, pick_default_matrix, pick_string    
  procedure :: pick_logical => marker_pick_logical
\nwendcode{}\nwbegincode{455}\moddef{Muli base: procedures}\plusendmoddef
  subroutine marker_pick_logical (this, name, content, status)
    class(marker_t), intent(inout) :: this
    character(*), intent(in) :: name
    logical, intent(out) :: content
    integer(dik), intent(out) :: status
    call this%pick_begin (name, status=status)
    if (status == serialize_ok) then
       call this%pop (content)
       call this%pick_end (name, status)
    end if
  end subroutine marker_pick_logical

\eatline
\nwendcode{}\nwbegindocs{456}\nwdocspar
\nwenddocs{}\nwbegincode{457}\moddef{Muli base: marker: TBP}\plusendmoddef
  procedure :: pick_integer => marker_pick_integer
\nwendcode{}\nwbegincode{458}\moddef{Muli base: procedures}\plusendmoddef
  subroutine marker_pick_integer (this, name, content, status)
    class(marker_t), intent(inout) :: this
    character(*), intent(in) :: name
    integer, intent(out) :: content
    integer(dik), intent(out) :: status
    call this%pick_begin (name, status=status)
    if (status == serialize_ok) then
       call this%pop (content)
       call this%pick_end (name, status)
    end if
  end subroutine marker_pick_integer

\eatline
\nwendcode{}\nwbegindocs{459} 
\nwenddocs{}\nwbegincode{460}\moddef{Muli base: marker: TBP}\plusendmoddef
  procedure :: pick_integer_array => marker_pick_integer_array
\nwendcode{}\nwbegincode{461}\moddef{Muli base: procedures}\plusendmoddef
  subroutine marker_pick_integer_array (this, name, content, status)
    class(marker_t), intent(inout) :: this
    character(*), intent(in) :: name
    integer, dimension(:), intent(out) :: content
    integer(dik), intent(out) :: status
    call this%pick_begin (name, status=status)
    if (status == serialize_ok) then
       call this%pop (content)
       call this%pick_end (name, status)
    end if
  end subroutine marker_pick_integer_array

\eatline
\nwendcode{}\nwbegindocs{462}\nwdocspar
\nwenddocs{}\nwbegincode{463}\moddef{Muli base: marker: TBP}\plusendmoddef
  procedure :: pick_integer_matrix => marker_pick_integer_matrix
\nwendcode{}\nwbegincode{464}\moddef{Muli base: procedures}\plusendmoddef
  subroutine marker_pick_integer_matrix (this, name, content, status)
    class(marker_t), intent(inout) :: this
    character(*), intent(in) :: name
    integer, dimension(:,:), intent(out) :: content
    integer(dik), intent(out) :: status
    integer :: n
    integer, dimension(2) :: s
    s = shape(content)
    call this%pick_begin (name, status=status)
    if (status == serialize_ok) then
       do n = 1, s(2)
          call this%pop (content(:,n))
       end do
       call this%pick_end (name, status)
    end if
  end subroutine marker_pick_integer_matrix

\eatline
\nwendcode{}\nwbegindocs{465}\nwdocspar
\nwenddocs{}\nwbegincode{466}\moddef{Muli base: marker: TBP}\plusendmoddef
  procedure :: pick_integer_dik => marker_pick_integer_dik
\nwendcode{}\nwbegincode{467}\moddef{Muli base: procedures}\plusendmoddef
  subroutine marker_pick_integer_dik (this, name, content, status)
    class(marker_t), intent(inout) :: this
    character(*), intent(in) :: name
    integer(dik), intent(out) :: content
    integer(dik), intent(out) :: status
    call this%pick_begin (name, status=status)
    if (status == serialize_ok) then
       call this%pop (content)
       call this%pick_end (name,status)
    end if
  end subroutine marker_pick_integer_dik

\eatline
\nwendcode{}\nwbegindocs{468}\nwdocspar
\nwenddocs{}\nwbegincode{469}\moddef{Muli base: marker: TBP}\plusendmoddef
  procedure :: pick_integer_array_dik => marker_pick_integer_array_dik
\nwendcode{}\nwbegincode{470}\moddef{Muli base: procedures}\plusendmoddef
  subroutine marker_pick_integer_array_dik (this, name, content, status)
    class(marker_t), intent(inout) :: this
    character(*), intent(in) :: name
    integer(dik), dimension(:), intent(out) :: content
    integer(dik), intent(out) :: status
    call this%pick_begin (name, status=status)
    if (status == serialize_ok) then
       call this%pop (content)
       call this%pick_end (name, status)
    end if
  end subroutine marker_pick_integer_array_dik

\eatline
\nwendcode{}\nwbegindocs{471}\nwdocspar
\nwenddocs{}\nwbegincode{472}\moddef{Muli base: marker: TBP}\plusendmoddef
  procedure :: pick_integer_matrix_dik => marker_pick_integer_matrix_dik  
\nwendcode{}\nwbegincode{473}\moddef{Muli base: procedures}\plusendmoddef
  subroutine marker_pick_integer_matrix_dik (this, name, content, status)
    class(marker_t), intent(inout) :: this
    character(*), intent(in) :: name
    integer(dik), dimension(:,:), intent(out) :: content
    integer(dik), intent(out) :: status
    integer :: n
    integer, dimension(2) :: s
    s = shape(content)
    call this%pick_begin (name, status=status)
    if (status == serialize_ok) then
       do n = 1, s(2)
          call this%pop (content(:,n))
       end do
       call this%pick_end (name,status)
    end if
  end subroutine marker_pick_integer_matrix_dik

\eatline
\nwendcode{}\nwbegindocs{474}\nwdocspar
\nwenddocs{}\nwbegincode{475}\moddef{Muli base: marker: TBP}\plusendmoddef
  procedure :: pick_default => marker_pick_default
\nwendcode{}\nwbegincode{476}\moddef{Muli base: procedures}\plusendmoddef
  subroutine marker_pick_default (this, name, content, status)
    class(marker_t), intent(inout) :: this
    character(*), intent(in) :: name
    real(default), intent(out) :: content
    integer(dik), intent(out) :: status
    call this%pick_begin (name, status=status)
    if (status == serialize_ok) then
       call this%pop (content)
       call this%pick_end (name,status)
    end if
  end subroutine marker_pick_default

\eatline
\nwendcode{}\nwbegindocs{477}\nwdocspar
\nwenddocs{}\nwbegincode{478}\moddef{Muli base: marker: TBP}\plusendmoddef
  procedure :: pick_default_array => marker_pick_default_array
\nwendcode{}\nwbegincode{479}\moddef{Muli base: procedures}\plusendmoddef
  subroutine marker_pick_default_array (this, name, content, status)
    class(marker_t), intent(inout) :: this
    character(*), intent(in) :: name
    real(default), dimension(:), intent(out) :: content
    integer(dik), intent(out) :: status
    call this%pick_begin (name, status=status)
    if (status == serialize_ok) then
       call this%pop (content)
       call this%pick_end (name, status)
    end if
  end subroutine marker_pick_default_array

\eatline
\nwendcode{}\nwbegindocs{480}\nwdocspar
\nwenddocs{}\nwbegincode{481}\moddef{Muli base: marker: TBP}\plusendmoddef
  procedure :: pick_default_matrix => marker_pick_default_matrix
\nwendcode{}\nwbegincode{482}\moddef{Muli base: procedures}\plusendmoddef
  subroutine marker_pick_default_matrix (this, name, content, status)
    class(marker_t), intent(inout) :: this
    character(*), intent(in) :: name
    real(default), dimension(:,:), intent(out) :: content
    integer(dik), intent(out) :: status
    integer :: n
    integer, dimension(2) :: s
    s = shape(content)
    call this%pick_begin (name, status=status)
    if (status == serialize_ok) then
       do n = 1, s(2)
          call this%pop (content(:,n))
       end do
       call this%pick_end (name,status)
    end if
  end subroutine marker_pick_default_matrix
  
\eatline
\nwendcode{}\nwbegindocs{483} 
\nwenddocs{}\nwbegincode{484}\moddef{Muli base: marker: TBP}\plusendmoddef
  procedure :: pick_string => marker_pick_string  
\nwendcode{}\nwbegincode{485}\moddef{Muli base: procedures}\plusendmoddef
  subroutine marker_pick_string (this, name, content, status)
    class(marker_t), intent(inout) :: this
    character(*), intent(in) :: name
    character(:), allocatable, intent(out) :: content
    integer(dik), intent(out) :: status
    call this%pick_begin (name, status=status)
    if (status == serialize_ok) then
       call this%pop (content)
       call this%pick_end (name, status)
    end if
  end subroutine marker_pick_string
    
\eatline
\nwendcode{}\nwbegindocs{486}\nwdocspar
\nwenddocs{}\nwbegincode{487}\moddef{Muli base: marker: TBP}\plusendmoddef
  procedure :: verify_nothing => marker_verify_nothing
\nwendcode{}\nwbegincode{488}\moddef{Muli base: procedures}\plusendmoddef
  subroutine marker_verify_nothing (this, name, status)
    class(marker_t), intent(inout) :: this
    character(*), intent(in) :: name
    integer(dik) ,intent(out) :: status
    integer(dik), dimension(2) :: p1, p2
    call this%find ("<", skip=4, proceed=.false., pos=p1)
    call this%find (">", 1, .false., p2)
    if (name//"/" == this%substring(p1, p2)) then
       status = serialize_nothing
       call this%set_position (p2)
       call this%proceed (i_one*3, .true.)
    else
       if (name == this%substring(p1, p2)) then
          status = serialize_ok
       else
          status = serialize_wrong_tag
       end if
    end if
  end subroutine marker_verify_nothing

\eatline
\nwendcode{}\nwbegindocs{489}\nwdocspar
\nwenddocs{}\nwbegincode{490}\moddef{Muli base: marker: TBP}\plusendmoddef
  procedure :: indent => marker_indent
\nwendcode{}\nwbegincode{491}\moddef{Muli base: procedures}\plusendmoddef
  subroutine marker_indent (this, step)
    class(marker_t), intent(inout) :: this
    integer(dik), optional :: step
    if (this%do_break)  call this%push (new_line(" "))
    if (this%do_indent) then
       if (present(step)) this%indentation = this%indentation + step
       call this%push (repeat(" ", this%indentation))
    end if
    this%active_pages(1) = this%actual_page()
  end subroutine marker_indent

\eatline
\nwendcode{}\nwbegindocs{492}\nwdocspar
\nwenddocs{}\nwbegincode{493}\moddef{Muli base: marker: TBP}\plusendmoddef
  procedure :: push_heap => marker_push_heap
\nwendcode{}\nwbegincode{494}\moddef{Muli base: procedures}\plusendmoddef
  subroutine marker_push_heap (this, ser, id)
    class(marker_t), intent(inout) :: this
    class(ser_class_t), target, intent(in) :: ser
    integer(dik), intent(in) :: id
    class(serializable_ref_type), pointer :: new_ref
    allocate (new_ref)
    new_ref%next => this%heap
    new_ref%ref => ser
    new_ref%id = id
    this%heap => new_ref
  end subroutine marker_push_heap
  
\eatline
\nwendcode{}\nwbegindocs{495}\nwdocspar
\nwenddocs{}\nwbegincode{496}\moddef{Muli base: marker: TBP}\plusendmoddef
  procedure :: pop_heap => marker_pop_heap
\nwendcode{}\nwbegincode{497}\moddef{Muli base: procedures}\plusendmoddef
  subroutine marker_pop_heap (this, ser)
    class(marker_t), intent(inout) :: this
    class(ser_class_t), pointer, intent(out) :: ser
    class(serializable_ref_type), pointer :: old_ref
    if (associated (this%heap)) then
       old_ref => this%heap
       ser => old_ref%ref
       this%heap => this%heap%next
       deallocate (old_ref)
    else
       call msg_error ("marker_pop_heap: heap_stack is not associated.")
    end if
  end subroutine marker_pop_heap

\eatline
\nwendcode{}\nwbegindocs{498}\nwdocspar
\nwenddocs{}\nwbegincode{499}\moddef{Muli base: marker: TBP}\plusendmoddef
  procedure :: push_reference => marker_push_reference
\nwendcode{}\nwbegincode{500}\moddef{Muli base: procedures}\plusendmoddef
  subroutine marker_push_reference (this, ser, id)
    class(marker_t), intent(inout) :: this
    class(ser_class_t), target, intent(in) :: ser
    integer(kind=dik), intent(in), optional :: id
    class(serializable_ref_type), pointer :: new_ref
    allocate (new_ref)
    new_ref%next => this%references
    new_ref%ref => ser
    if (present(id)) then
       new_ref%id = id
    else
       new_ref%id = -1
    end if
    this%references => new_ref
  end subroutine marker_push_reference
  
\eatline
\nwendcode{}\nwbegindocs{501}\nwdocspar
\nwenddocs{}\nwbegincode{502}\moddef{Muli base: marker: TBP}\plusendmoddef
  procedure :: pop_reference => marker_pop_reference
\nwendcode{}\nwbegincode{503}\moddef{Muli base: procedures}\plusendmoddef
  subroutine marker_pop_reference (this, ser)
    class(marker_t), intent(inout) :: this
    class(ser_class_t), pointer, intent(out) :: ser
    class(serializable_ref_type), pointer :: old_ref
    if (associated (this%references)) then
       old_ref => this%references
       ser => old_ref%ref
       this%references => this%references%next
       deallocate (old_ref)
    else
       call msg_error &
            ("marker_pop_reference: reference_stack is not associated.")
    end if
  end subroutine marker_pop_reference

\eatline
\nwendcode{}\nwbegindocs{504}\nwdocspar
\nwenddocs{}\nwbegincode{505}\moddef{Muli base: marker: TBP}\plusendmoddef
  procedure :: reset_references => marker_reset_references
\nwendcode{}\nwbegincode{506}\moddef{Muli base: procedures}\plusendmoddef
  subroutine marker_reset_references (this)
    class(marker_t), intent(inout) :: this
    if (associated (this%references)) then
       call this%references%finalize ()
       deallocate (this%references)
    end if
  end subroutine marker_reset_references
  
\eatline
\nwendcode{}\nwbegindocs{507}\nwdocspar
\nwenddocs{}\nwbegincode{508}\moddef{Muli base: marker: TBP}\plusendmoddef
  procedure :: search_reference => marker_search_reference
\nwendcode{}\nwbegincode{509}\moddef{Muli base: procedures}\plusendmoddef
  subroutine marker_search_reference (this, type, ser)
    class(marker_t), intent(in) :: this
    integer(dik), dimension(2,2), intent(in) :: type
    class(ser_class_t), pointer, intent(out) :: ser
    !!! !!! !!! NAG bug workaround
    class(ser_class_t), pointer :: tmp_ser 
    class(serializable_ref_type), pointer :: ref
    ref => this%references
    nullify (ser)
    do while (associated (ref))
       tmp_ser => ref%ref
       if (tmp_ser%verify_type (this%substring(type))) then
          ser => tmp_ser
          exit
       end if
       ref => ref%next
    end do
  end subroutine marker_search_reference

\nwendcode{}\nwbegindocs{510}%ref marker_search_reference 
\nwenddocs{}\nwbegindocs{511}\nwdocspar
\nwenddocs{}\nwbegincode{512}\moddef{Muli base: marker: TBP}\plusendmoddef
  procedure :: reset_heap => marker_reset_heap  
\nwendcode{}\nwbegincode{513}\moddef{Muli base: procedures}\plusendmoddef
  subroutine marker_reset_heap (this)
    class(marker_t), intent(inout) :: this
    if (associated (this%heap)) then
       call this%heap%finalize ()
       deallocate (this%heap)
    end if
  end subroutine marker_reset_heap

\eatline
\nwendcode{}\nwbegindocs{514}\nwdocspar
\nwenddocs{}\nwbegincode{515}\moddef{Muli base: marker: TBP}\plusendmoddef
  procedure :: finalize => marker_finalize
\nwendcode{}\nwbegincode{516}\moddef{Muli base: procedures}\plusendmoddef
  subroutine marker_finalize (this)
    class(marker_t), intent(inout) :: this
    call this%reset_heap ()
    call this%reset_references ()
  end subroutine marker_finalize
  
\eatline
\nwendcode{}\nwbegindocs{517}\nwdocspar
\nwenddocs{}\nwbegincode{518}\moddef{Muli base: marker: TBP}\plusendmoddef
  generic :: search_heap => search_heap_by_id, search_heap_by_ref
  procedure :: search_heap_by_id => marker_search_heap_by_id
  procedure :: search_heap_by_ref => marker_search_heap_by_ref  
\nwendcode{}\nwbegincode{519}\moddef{Muli base: procedures}\plusendmoddef
  subroutine marker_search_heap_by_ref (this, ref, id)
    class(marker_t), intent(in) :: this
    class(ser_class_t), pointer, intent(in) :: ref
    integer(dik), intent(out) :: id
    class(serializable_ref_type), pointer :: ref_p
    ref_p => this%heap
    id = 0
    do while (associated (ref_p))
       if (associated (ref, ref_p%ref)) then
          id = ref_p%id
          exit
       end if
       ref_p => ref_p%next
    end do
  end subroutine marker_search_heap_by_ref
  
\eatline
\nwendcode{}\nwbegindocs{520}\nwdocspar
\nwenddocs{}\nwbegincode{521}\moddef{Muli base: procedures}\plusendmoddef
  subroutine marker_search_heap_by_id (this, id, ser)
    class(marker_t), intent(in) :: this
    integer(dik), intent(in) :: id
    class(ser_class_t), pointer, intent(out) :: ser
    class(serializable_ref_type), pointer :: ref
    ref => this%heap
    do while (associated (ref))
       if (id == ref%id) then
          ser => ref%ref
          exit
       end if
       ref => ref%next
    end do
  end subroutine marker_search_heap_by_id
  
\eatline
\nwendcode{}\nwbegindocs{522}\nwdocspar
\nwenddocs{}\nwbegincode{523}\moddef{Muli base: procedures}\plusendmoddef
  elemental function measurable_less_measurable (mea1, mea2)
    class(measure_class_t), intent(in) :: mea1, mea2
    logical :: measurable_less_measurable
    measurable_less_measurable = mea1%measure() < mea2%measure()
  end function measurable_less_measurable
  
\eatline
\nwendcode{}\nwbegindocs{524}\nwdocspar
\nwenddocs{}\nwbegincode{525}\moddef{Muli base: procedures}\plusendmoddef
  elemental function measurable_less_default (mea1, def)
    class(measure_class_t), intent(in) :: mea1
    real(default), intent(in) :: def
    logical :: measurable_less_default
    measurable_less_default = mea1%measure() < def
  end function measurable_less_default
  
\eatline
\nwendcode{}\nwbegindocs{526}\nwdocspar
\nwenddocs{}\nwbegincode{527}\moddef{Muli base: procedures}\plusendmoddef
  elemental function measurable_less_or_equal_measurable (mea1, mea2)
    class(measure_class_t), intent(in) :: mea1, mea2
    logical :: measurable_less_or_equal_measurable
    measurable_less_or_equal_measurable = mea1%measure() <= mea2%measure()
  end function measurable_less_or_equal_measurable

\eatline
\nwendcode{}\nwbegindocs{528}\nwdocspar
\nwenddocs{}\nwbegincode{529}\moddef{Muli base: procedures}\plusendmoddef
  elemental function measurable_less_or_equal_default (mea1, def)
    class(measure_class_t), intent(in) :: mea1
    real(default), intent(in) :: def
    logical :: measurable_less_or_equal_default
    measurable_less_or_equal_default = mea1%measure() <= def
  end function measurable_less_or_equal_default

\eatline
\nwendcode{}\nwbegindocs{530}\nwdocspar
\nwenddocs{}\nwbegincode{531}\moddef{Muli base: procedures}\plusendmoddef
 elemental function measurable_equal_measurable (mea1, mea2)
    class(measure_class_t), intent(in) :: mea1, mea2
    logical :: measurable_equal_measurable
    measurable_equal_measurable = mea1%measure() == mea2%measure()
  end function measurable_equal_measurable

\eatline
\nwendcode{}\nwbegindocs{532}\nwdocspar
\nwenddocs{}\nwbegincode{533}\moddef{Muli base: procedures}\plusendmoddef
  elemental function measurable_equal_default (mea1, def)
    class(measure_class_t), intent(in) :: mea1
    real(default), intent(in) :: def
    logical :: measurable_equal_default
    measurable_equal_default = mea1%measure() == def
  end function measurable_equal_default

\eatline
\nwendcode{}\nwbegindocs{534}\nwdocspar
\nwenddocs{}\nwbegincode{535}\moddef{Muli base: procedures}\plusendmoddef
  elemental function measurable_equal_or_greater_measurable (mea1, mea2)
    class(measure_class_t), intent(in) :: mea1, mea2
    logical :: measurable_equal_or_greater_measurable
    measurable_equal_or_greater_measurable = mea1%measure() >= mea2%measure()
  end function measurable_equal_or_greater_measurable

\eatline
\nwendcode{}\nwbegindocs{536}\nwdocspar
\nwenddocs{}\nwbegincode{537}\moddef{Muli base: procedures}\plusendmoddef
  elemental function measurable_equal_or_greater_default (mea1, def)
    class(measure_class_t), intent(in) :: mea1
    real(default), intent(in) :: def
    logical :: measurable_equal_or_greater_default
    measurable_equal_or_greater_default = mea1%measure() >= def
  end function measurable_equal_or_greater_default
  
\eatline
\nwendcode{}\nwbegindocs{538}\nwdocspar
\nwenddocs{}\nwbegincode{539}\moddef{Muli base: procedures}\plusendmoddef
  elemental function measurable_greater_measurable (mea1, mea2)
    class(measure_class_t), intent(in) :: mea1,mea2
    logical :: measurable_greater_measurable
    measurable_greater_measurable = mea1%measure() > mea2%measure()
  end function measurable_greater_measurable

\eatline
\nwendcode{}\nwbegindocs{540}\nwdocspar
\nwenddocs{}\nwbegincode{541}\moddef{Muli base: procedures}\plusendmoddef
  elemental function measurable_greater_default (mea1, def)
    class(measure_class_t), intent(in) :: mea1
    real(default), intent(in) :: def
    logical :: measurable_greater_default
    measurable_greater_default = mea1%measure() > def
  end function measurable_greater_default

\eatline
\nwendcode{}\nwbegindocs{542}\nwdocspar
\nwenddocs{}\nwbegincode{543}\moddef{Muli base: procedures}\plusendmoddef
  pure function page_ring_position (n)
    integer(dik), intent(in) :: n
    integer(dik), dimension(2) :: page_ring_position
    page_ring_position(2) = mod(n, serialize_page_size)
    page_ring_position(1) = (n-page_ring_position(2)) / serialize_page_size
  end function page_ring_position

\eatline
\nwendcode{}\nwbegindocs{544}\nwdocspar
\nwenddocs{}\nwbegincode{545}\moddef{Muli base: procedures}\plusendmoddef
  pure integer(dik) function page_ring_ordinal (pos)
    integer(dik), dimension(2), intent(in) :: pos
    page_ring_ordinal = pos(1) * serialize_page_size + pos(2)
  end function page_ring_ordinal

\eatline
\nwendcode{}\nwbegindocs{546}\nwdocspar
\nwenddocs{}\nwbegincode{547}\moddef{Muli base: procedures}\plusendmoddef
  pure logical function page_ring_position_is_before_int_pos (m, n)
    integer(dik), intent(in) :: m
    integer(dik), dimension(2), intent(in) :: n
    if (m < page_ring_ordinal(n)) then
       page_ring_position_is_before_int_pos = .true.
    else
       page_ring_position_is_before_int_pos = .false.
    end if
  end function page_ring_position_is_before_int_pos

\eatline
\nwendcode{}\nwbegindocs{548}\nwdocspar
\nwenddocs{}\nwbegincode{549}\moddef{Muli base: procedures}\plusendmoddef
  pure logical function page_ring_position_is_before_pos_int (m, n)
    integer(dik), dimension(2), intent(in) :: m
    integer(dik), intent(in) :: n
    if (page_ring_ordinal(m) < n) then
       page_ring_position_is_before_pos_int = .true.
    else
       page_ring_position_is_before_pos_int = .false.
    end if
  end function page_ring_position_is_before_pos_int

\eatline
\nwendcode{}\nwbegindocs{550}\nwdocspar
\nwenddocs{}\nwbegincode{551}\moddef{Muli base: procedures}\plusendmoddef
  pure logical function page_ring_position_is_before_pos_pos (m, n)
    integer(dik), dimension(2), intent(in) :: m,n
    if (m(1) < n(1)) then
       page_ring_position_is_before_pos_pos = .true.
    else
       if (m(1) > n(1)) then
          page_ring_position_is_before_pos_pos = .false.
       else
          if (m(2) < n(2)) then
             page_ring_position_is_before_pos_pos = .true.
          else
             page_ring_position_is_before_pos_pos = .false.
          end if
       end if
    end if
  end function page_ring_position_is_before_pos_pos

\eatline
\nwendcode{}\nwbegindocs{552}\nwdocspar
\nwenddocs{}\nwbegincode{553}\moddef{Muli base: procedures}\plusendmoddef
  subroutine ring_position_increase (pos, n)
    integer(dik), dimension(2), intent(inout) :: pos
    integer(dik), intent(in) :: n
    pos = page_ring_position (page_ring_ordinal(pos) + n)
  end subroutine ring_position_increase

\eatline
\nwendcode{}\nwbegindocs{554}\nwdocspar
\nwenddocs{}\nwbegincode{555}\moddef{Muli base: procedures}\plusendmoddef
  pure integer(dik) function ring_position_metric1 (p)
    integer(dik), dimension(2,2), intent(in) :: p
    ring_position_metric1 = (p(1,2) - p(1,1)) * serialize_page_size + &
         p(2,2) - p(2,1) + 1
  end function ring_position_metric1
  
  pure integer(dik) function ring_position_metric2 (p1, p2)
    integer(dik), dimension(2), intent(in) :: p1, p2
    ring_position_metric2 = (p2(1) - p1(1)) * &
         serialize_page_size + p2(2) - p1(2) + 1
  end function ring_position_metric2
  
\eatline
\nwendcode{}\nwbegindocs{556}\nwdocspar
\nwenddocs{}\nwbegincode{557}\moddef{Muli base: public}\plusendmoddef
  public :: generate_unit
\nwendcode{}\nwbegincode{558}\moddef{Muli base: procedures}\plusendmoddef
  subroutine generate_unit (unit, min, max)
    integer, intent(out) :: unit
    integer, intent(in), optional :: min,max
    integer :: min_u, max_u
    logical :: is_open
    ! print *,"generate_unit"
    unit = -1
    if (present (min)) then
       min_u = min
    else
       min_u = 10
    end if
    if (present (max)) then
       max_u = max
    else
       max_u = huge (max_u)
    end if
    do unit = min_u, max_u
       !print *,"testing ",unit
       inquire (unit, opened = is_open)
       if (.not. is_open) then
          exit
       end if
    end do
  end subroutine generate_unit  
  
\eatline
\nwendcode{}\nwbegindocs{559}\nwdocspar
\nwenddocs{}\nwbegincode{560}\moddef{Muli base: public}\plusendmoddef
  public :: ilog2
\nwendcode{}\nwbegincode{561}\moddef{Muli base: procedures}\plusendmoddef
  subroutine ilog2 (int, exp, rem)
    integer,intent(in) :: int
    integer,intent(out) :: exp, rem
    integer :: count
    count = 2
    exp = 1
    do while (count < int)
       exp = exp + 1
       count = ishft(count, 1)
    end do
    if (count > int) then
       rem = (int - ishft(count, -1))
    else
       rem = 0
    end if    
  end subroutine ilog2
  
\eatline
\nwendcode{}\nwbegindocs{562}\nwdocspar
\nwenddocs{}\nwbegincode{563}\moddef{Muli base: public}\plusendmoddef
  public :: integer_with_leading_zeros
\nwendcode{}\nwbegincode{564}\moddef{Muli base: procedures}\plusendmoddef
  subroutine integer_with_leading_zeros (number, length, string)
    integer, intent(in) :: number, length
    character(len=*), intent(out) :: string
    integer :: zeros
    character::sign
    if (number == 0) then
       string = repeat("0", length)
    else
       if (number > 0) then
          zeros = length -floor(log10 (real(number))) - 1
          if (zeros < 0) then
             string = repeat("*", length)
          else
             write (string, fmt="(A,I0)")  repeat("0", zeros), number
          end if
       else
          zeros = length - floor (log10 (real (-number))) - 2
          if (zeros < 0) then
             string = repeat("*", length)
          else
             write (string, fmt="(A,A,I0)") "-", repeat("0", zeros), &
                  abs(number)
          end if
       end if
    end if
  end subroutine integer_with_leading_zeros

\eatline
\nwendcode{}\nwbegindocs{565} 
\nwenddocs{}\nwbegincode{566}\moddef{Muli base: procedures}\plusendmoddef
  pure logical function character_is_in (c, array)
    character, intent(in) :: c
    character, dimension(:), intent(in) :: array
    integer(dik) :: n
    character_is_in = .false.
    do n=1,size(array)
       if (c == array(n)) then
          character_is_in = .true.
          exit
       end if
    end do
  end function character_is_in

\eatline
\nwendcode{}\nwbegindocs{567}\nwdocspar
%%%%%%%%%%%%%%%%%%%%%%%%%%%%%%%%%%%%%%%%%%%%%%%%%%%%%%%%%%%%%%%%%%%%%%%%
\section{Transverse momenta}

This file contains the module {\tt{}muli{\char95}momentum}. Its purpose is to
store the actual value of the evolution parameter $p_t^2$ in a
convenient way. We use the normalized value $p_t^2 /
p_{t,\text{max}}^2$  for generating the next value of the scale, also
need the square root of both $p_t^2$ and $p_t^2 / p_{t,\text{max}}^2$
for other purposes. That's why I store all four combinations together
with $p_{t, \text{max}}$ in an array.  

\nwenddocs{}\nwbegincode{568}\moddef{\code{}muli{\char95}momentum.f90\edoc{}}\endmoddef
\LA{}File header\RA{}

module muli_momentum
  
\LA{}Use kinds\RA{}
  use constants
  use muli_base
  
\LA{}Standard module head\RA{}
  
\LA{}Muli mom: public\RA{}

\LA{}Muli mom: types\RA{}

\LA{}Muli mom: interfaces\RA{}

contains

\LA{}Muli mom: procedures\RA{}  
    
end module muli_momentum

\eatline
\nwendcode{}\nwbegindocs{569}\nwdocspar
\nwenddocs{}\nwbegincode{570}\moddef{Muli mom: public}\endmoddef
  public :: transverse_mom_t
\nwendcode{}\nwbegincode{571}\moddef{Muli mom: types}\endmoddef
  type, extends (ser_class_t) :: transverse_mom_t
     private
     real(default), dimension(0:4) :: momentum = [0, 0, 0, 0, 0]
   contains
   \LA{}Muli mom: trans mom: TBP\RA{}
  end type transverse_mom_t
  
\eatline
\nwendcode{}\nwbegindocs{572}\nwdocspar
\nwenddocs{}\nwbegincode{573}\moddef{Muli mom: trans mom: TBP}\endmoddef
  procedure :: mom_write_to_marker => transverse_mom_write_to_marker
  procedure :: write_to_marker => transverse_mom_write_to_marker
\nwendcode{}\nwbegincode{574}\moddef{Muli mom: procedures}\endmoddef
  subroutine transverse_mom_write_to_marker (this, marker, status)
    class(transverse_mom_t), intent(in) :: this
    class(marker_t), intent(inout) :: marker
    integer(dik), intent(out) :: status
    call marker%mark_begin ("transverse_mom_t")
    call marker%mark ("gev_momenta", this%momentum(0:1))
    call marker%mark_end ("transverse_mom_t")
  end subroutine transverse_mom_write_to_marker

\eatline
\nwendcode{}\nwbegindocs{575}\nwdocspar
\nwenddocs{}\nwbegincode{576}\moddef{Muli mom: trans mom: TBP}\plusendmoddef
  procedure :: mom_read_from_marker => transverse_mom_read_from_marker
  procedure :: read_from_marker => transverse_mom_read_from_marker
\nwendcode{}\nwbegincode{577}\moddef{Muli mom: procedures}\plusendmoddef
  subroutine transverse_mom_read_from_marker (this, marker, status)
    class(transverse_mom_t), intent(out) :: this
    class(marker_t), intent(inout) :: marker
    integer(dik), intent(out) :: status
    call marker%pick_begin ("transverse_mom_t", status=status)
    call marker%pick ("gev_momenta", this%momentum(0:1), status)
    this%momentum(2:4) = [ this%momentum(1)**2, &
                           this%momentum(1) / this%momentum(0), &
                           (this%momentum(1)/this%momentum(0))**2 ]
    call marker%pick_end ("transverse_mom_t", status=status)
  end subroutine transverse_mom_read_from_marker

\eatline
\nwendcode{}\nwbegindocs{578} 
\nwenddocs{}\nwbegincode{579}\moddef{Muli mom: trans mom: TBP}\plusendmoddef
  procedure :: mom_print_to_unit => transverse_mom_print_to_unit
  procedure :: print_to_unit => transverse_mom_print_to_unit
\nwendcode{}\nwbegincode{580}\moddef{Muli mom: procedures}\plusendmoddef
  subroutine transverse_mom_print_to_unit &
       (this, unit, parents, components, peers)
    class(transverse_mom_t), intent(in) :: this
    integer, intent(in) :: unit
    integer(dik), intent(in) :: parents, components, peers
    write (unit, "(1x,A)")  "Components of transverse_mom_t:"
    write (unit, "(3x,A)")  "Actual energy scale:"
    write (unit, "(A,E20.10)")  "Max scale (MeV)   :", this%momentum(0)
    write (unit, "(A,E20.10)")  "Scale (MeV)       :", this%momentum(1)
    write (unit, "(A,E20.10)")  "Scale^2 (MeV^2)   :", this%momentum(2)
    write (unit, "(A,E20.10)")  "Scale normalized  :", this%momentum(3)
    write (unit, "(A,E20.10)")  "Scale^2 normalized:", this%momentum(4)
  end subroutine transverse_mom_print_to_unit
    
\eatline
\nwendcode{}\nwbegindocs{581}\nwdocspar
\nwenddocs{}\nwbegincode{582}\moddef{Muli mom: trans mom: TBP}\plusendmoddef
  procedure, nopass :: get_type => transverse_mom_get_type    
\nwendcode{}\nwbegincode{583}\moddef{Muli mom: procedures}\plusendmoddef
  pure subroutine transverse_mom_get_type (type)
    character(:), allocatable, intent(out) :: type
    allocate (type, source="transverse_mom_t")
  end subroutine transverse_mom_get_type
  
\eatline
\nwendcode{}\nwbegindocs{584}\nwdocspar
\nwenddocs{}\nwbegincode{585}\moddef{Muli mom: trans mom: TBP}\plusendmoddef
  procedure :: get_gev_initial_cme => transverse_mom_get_gev_initial_cme
\nwendcode{}\nwbegincode{586}\moddef{Muli mom: procedures}\plusendmoddef
  elemental function transverse_mom_get_gev_initial_cme (this) result(scale)
    class(transverse_mom_t), intent(in) :: this
    real(default) :: scale
    scale = this%momentum(0) * 2D0
  end function transverse_mom_get_gev_initial_cme

\eatline
\nwendcode{}\nwbegindocs{587}\nwdocspar
\nwenddocs{}\nwbegincode{588}\moddef{Muli mom: trans mom: TBP}\plusendmoddef
  procedure :: get_gev_max_scale => transverse_mom_get_gev_max_scale
\nwendcode{}\nwbegincode{589}\moddef{Muli mom: procedures}\plusendmoddef
  elemental function transverse_mom_get_gev_max_scale (this) result (scale)
    class(transverse_mom_t), intent(in) :: this
    real(default) :: scale
    scale = this%momentum(0)
  end function transverse_mom_get_gev_max_scale

\eatline
\nwendcode{}\nwbegindocs{590}\nwdocspar
\nwenddocs{}\nwbegincode{591}\moddef{Muli mom: trans mom: TBP}\plusendmoddef
  procedure :: get_gev2_max_scale => transverse_mom_get_gev2_max_scale
\nwendcode{}\nwbegincode{592}\moddef{Muli mom: procedures}\plusendmoddef
  elemental function transverse_mom_get_gev2_max_scale (this) result (scale)
    class(transverse_mom_t), intent(in) :: this
    real(default) :: scale
    scale = this%momentum(0)**2
  end function transverse_mom_get_gev2_max_scale

\eatline
\nwendcode{}\nwbegindocs{593}\nwdocspar
\nwenddocs{}\nwbegincode{594}\moddef{Muli mom: trans mom: TBP}\plusendmoddef
  procedure :: get_gev_scale => transverse_mom_get_gev_scale
\nwendcode{}\nwbegincode{595}\moddef{Muli mom: procedures}\plusendmoddef
  elemental function transverse_mom_get_gev_scale(this) result(scale)
    class(transverse_mom_t), intent(in) :: this
    real(default) :: scale
    scale = this%momentum(1)
  end function transverse_mom_get_gev_scale

\eatline
\nwendcode{}\nwbegindocs{596}\nwdocspar
\nwenddocs{}\nwbegincode{597}\moddef{Muli mom: trans mom: TBP}\plusendmoddef
  procedure :: get_gev2_scale => transverse_mom_get_gev2_scale
\nwendcode{}\nwbegincode{598}\moddef{Muli mom: procedures}\plusendmoddef
  elemental function transverse_mom_get_gev2_scale (this) result (scale)
    class(transverse_mom_t), intent(in) :: this
    real(default) :: scale
    scale = this%momentum(2)
  end function transverse_mom_get_gev2_scale

\eatline
\nwendcode{}\nwbegindocs{599}\nwdocspar
\nwenddocs{}\nwbegincode{600}\moddef{Muli mom: trans mom: TBP}\plusendmoddef
  procedure :: get_unit_scale => transverse_mom_get_unit_scale
\nwendcode{}\nwbegincode{601}\moddef{Muli mom: procedures}\plusendmoddef
  pure function transverse_mom_get_unit_scale (this) result (scale)
    class(transverse_mom_t), intent(in) :: this
    real(default) :: scale
    scale = this%momentum(3)
  end function transverse_mom_get_unit_scale

\eatline
\nwendcode{}\nwbegindocs{602}\nwdocspar
\nwenddocs{}\nwbegincode{603}\moddef{Muli mom: trans mom: TBP}\plusendmoddef
  procedure :: get_unit2_scale => transverse_mom_get_unit2_scale
\nwendcode{}\nwbegincode{604}\moddef{Muli mom: procedures}\plusendmoddef
  pure function transverse_mom_get_unit2_scale (this) result (scale)
    class(transverse_mom_t), intent(in) :: this
    real(default) :: scale
    scale = this%momentum(4)
  end function transverse_mom_get_unit2_scale

\eatline
\nwendcode{}\nwbegindocs{605}\nwdocspar
\nwenddocs{}\nwbegincode{606}\moddef{Muli mom: trans mom: TBP}\plusendmoddef
  procedure :: set_gev_initial_cme => transverse_mom_set_gev_initial_cme
\nwendcode{}\nwbegincode{607}\moddef{Muli mom: procedures}\plusendmoddef
  subroutine transverse_mom_set_gev_initial_cme (this, new_gev_initial_cme)
    class(transverse_mom_t), intent(inout) :: this
    real(default), intent(in)  ::  new_gev_initial_cme
    this%momentum(0) = new_gev_initial_cme / 2D0
    this%momentum(3) = this%momentum(1) / this%momentum(0)
    this%momentum(4) = this%momentum(3)**2
  end subroutine transverse_mom_set_gev_initial_cme

\eatline
\nwendcode{}\nwbegindocs{608}\nwdocspar
\nwenddocs{}\nwbegincode{609}\moddef{Muli mom: trans mom: TBP}\plusendmoddef
  procedure :: set_gev_max_scale => transverse_mom_set_gev_max_scale
\nwendcode{}\nwbegincode{610}\moddef{Muli mom: procedures}\plusendmoddef
  subroutine transverse_mom_set_gev_max_scale (this, new_gev_max_scale)
    class(transverse_mom_t), intent(inout) :: this
    real(default), intent(in)  ::  new_gev_max_scale
    this%momentum(0) = new_gev_max_scale
    this%momentum(3) = this%momentum(1) / this%momentum(0)
    this%momentum(4) = this%momentum(3)**2
  end subroutine transverse_mom_set_gev_max_scale
  
\eatline
\nwendcode{}\nwbegindocs{611}\nwdocspar
\nwenddocs{}\nwbegincode{612}\moddef{Muli mom: trans mom: TBP}\plusendmoddef
  procedure :: set_gev2_max_scale => transverse_mom_set_gev2_max_scale
\nwendcode{}\nwbegincode{613}\moddef{Muli mom: procedures}\plusendmoddef
  subroutine transverse_mom_set_gev2_max_scale (this, new_gev2_max_scale)
    class(transverse_mom_t), intent(inout) :: this
    real(default), intent(in)  ::  new_gev2_max_scale
    this%momentum(0) = sqrt (new_gev2_max_scale)
    this%momentum(3) = this%momentum(1) / this%momentum(0)
    this%momentum(4) = this%momentum(3)**2
  end subroutine transverse_mom_set_gev2_max_scale

\eatline
\nwendcode{}\nwbegindocs{614}\nwdocspar
\nwenddocs{}\nwbegincode{615}\moddef{Muli mom: trans mom: TBP}\plusendmoddef
  procedure :: set_gev_scale => transverse_mom_set_gev_scale
\nwendcode{}\nwbegincode{616}\moddef{Muli mom: procedures}\plusendmoddef
  subroutine transverse_mom_set_gev_scale (this, new_gev_scale)
    class(transverse_mom_t), intent(inout) :: this
    real(default), intent(in)  ::  new_gev_scale
    this%momentum(1) = new_gev_scale
    this%momentum(2) = new_gev_scale**2
    this%momentum(3) = new_gev_scale / this%momentum(0)
    this%momentum(4) = this%momentum(3)**2
  end subroutine transverse_mom_set_gev_scale

\eatline
\nwendcode{}\nwbegindocs{617}\nwdocspar
\nwenddocs{}\nwbegincode{618}\moddef{Muli mom: trans mom: TBP}\plusendmoddef
  procedure :: set_gev2_scale => transverse_mom_set_gev2_scale
\nwendcode{}\nwbegincode{619}\moddef{Muli mom: procedures}\plusendmoddef
  subroutine transverse_mom_set_gev2_scale (this, new_gev2_scale)
    class(transverse_mom_t), intent(inout) :: this
    real(default), intent(in) :: new_gev2_scale
    this%momentum(1) = sqrt (new_gev2_scale)
    this%momentum(2) = new_gev2_scale
    this%momentum(3) = this%momentum(1) / this%momentum(0)
    this%momentum(4) = this%momentum(3)**2
  end subroutine transverse_mom_set_gev2_scale

\eatline
\nwendcode{}\nwbegindocs{620}\nwdocspar
\nwenddocs{}\nwbegincode{621}\moddef{Muli mom: trans mom: TBP}\plusendmoddef
  procedure :: set_unit_scale => transverse_mom_set_unit_scale
\nwendcode{}\nwbegincode{622}\moddef{Muli mom: procedures}\plusendmoddef
  subroutine transverse_mom_set_unit_scale (this, new_unit_scale)
    class(transverse_mom_t), intent(inout)::this
    real(default), intent(in) :: new_unit_scale
    this%momentum(1) = new_unit_scale * this%momentum(0)
    this%momentum(2) = this%momentum(1)**2
    this%momentum(3) = new_unit_scale
    this%momentum(4) = this%momentum(3)**2
  end subroutine transverse_mom_set_unit_scale

\eatline
\nwendcode{}\nwbegindocs{623}\nwdocspar
\nwenddocs{}\nwbegincode{624}\moddef{Muli mom: trans mom: TBP}\plusendmoddef
  procedure :: set_unit2_scale => transverse_mom_set_unit2_scale
\nwendcode{}\nwbegincode{625}\moddef{Muli mom: procedures}\plusendmoddef
  subroutine transverse_mom_set_unit2_scale (this, new_unit2_scale)
    class(transverse_mom_t), intent(inout)::this
    real(default), intent(in) :: new_unit2_scale
    this%momentum(3) = sqrt (new_unit2_scale)
    this%momentum(4) = new_unit2_scale
    this%momentum(1) = this%momentum(3) * this%momentum(0)
    this%momentum(2) = this%momentum(1)**2
  end subroutine transverse_mom_set_unit2_scale

\eatline
\nwendcode{}\nwbegindocs{626}\nwdocspar
\nwenddocs{}\nwbegincode{627}\moddef{Muli mom: trans mom: TBP}\plusendmoddef
  generic :: initialize => transverse_mom_initialize
  procedure :: transverse_mom_initialize  
\nwendcode{}\nwbegincode{628}\moddef{Muli mom: procedures}\plusendmoddef
  subroutine transverse_mom_initialize (this, gev2_s)
    class(transverse_mom_t), intent(out) :: this
    real(default), intent(in) :: gev2_s
    real(default) :: gev_s
    gev_s = sqrt (gev2_s)
    this%momentum = [gev_s/2, gev_s/2, gev2_s/4, one, one]
  end subroutine transverse_mom_initialize
  
\eatline
\nwendcode{}\nwbegindocs{629}\nwdocspar
\nwenddocs{}\nwbegincode{630}\moddef{Muli mom: public}\plusendmoddef
  public :: qcd_2_2_class
\nwendcode{}\nwbegincode{631}\moddef{Muli mom: types}\plusendmoddef
  type, extends (transverse_mom_t), abstract :: qcd_2_2_class
   contains
   \LA{}Muli mom: QCD22: TBP\RA{}  
  end type qcd_2_2_class

\eatline
\nwendcode{}\nwbegindocs{632}\nwdocspar
\nwenddocs{}\nwbegincode{633}\moddef{Muli mom: QCD22: TBP}\endmoddef
  procedure(qcd_get_int), deferred :: get_process_id
\nwendcode{}\nwbegincode{634}\moddef{Muli mom: QCD22: TBP}\plusendmoddef
  procedure(qcd_get_int), deferred :: get_integrand_id
\nwendcode{}\nwbegincode{635}\moddef{Muli mom: QCD22: TBP}\plusendmoddef
  procedure(qcd_get_int), deferred :: get_diagram_kind
\nwendcode{}\nwbegincode{636}\moddef{Muli mom: QCD22: TBP}\plusendmoddef
  procedure(qcd_get_int_4), deferred :: get_lha_flavors
\nwendcode{}\nwbegincode{637}\moddef{Muli mom: QCD22: TBP}\plusendmoddef
  procedure(qcd_get_int_4), deferred :: get_pdg_flavors
\nwendcode{}\nwbegincode{638}\moddef{Muli mom: QCD22: TBP}\plusendmoddef
  procedure(qcd_get_int_by_int), deferred :: get_parton_id
\nwendcode{}\nwbegincode{639}\moddef{Muli mom: QCD22: TBP}\plusendmoddef
  procedure(qcd_get_int_2), deferred :: get_parton_kinds
\nwendcode{}\nwbegincode{640}\moddef{Muli mom: QCD22: TBP}\plusendmoddef
  procedure(qcd_get_int_2), deferred :: get_pdf_int_kinds
\nwendcode{}\nwbegincode{641}\moddef{Muli mom: QCD22: TBP}\plusendmoddef
  procedure(qcd_get_real), deferred :: get_momentum_boost
\nwendcode{}\nwbegincode{642}\moddef{Muli mom: QCD22: TBP}\plusendmoddef
  ! procedure(qcd_get_real_3),deferred :: get_parton_in_momenta
\nwendcode{}\nwbegincode{643}\moddef{Muli mom: QCD22: TBP}\plusendmoddef
  procedure(qcd_get_real_2), deferred :: get_remnant_momentum_fractions
\nwendcode{}\nwbegincode{644}\moddef{Muli mom: QCD22: TBP}\plusendmoddef
  procedure(qcd_get_real_2), deferred :: get_total_momentum_fractions  
\nwendcode{}\nwbegincode{645}\moddef{Muli mom: interfaces}\endmoddef
  abstract interface
     subroutine qcd_none (this)
       import qcd_2_2_class
       class(qcd_2_2_class), target, intent(in) :: this
     end subroutine qcd_none
  end interface
\eatline
\nwendcode{}\nwbegindocs{646}\nwdocspar
\nwenddocs{}\nwbegincode{647}\moddef{Muli mom: interfaces}\plusendmoddef
  ! abstract interface
  !    subroutine qcd_get_beam (this, beam)
  !      import qcd_2_2_class
  !      import pp_remnant_class
  !      class(qcd_2_2_class),target, intent(in) :: this
  !      class(pp_remnant_class),pointer, intent(out) :: beam
  !     end subroutine qcd_get_beam
  ! end interface
\eatline
\nwendcode{}\nwbegindocs{648}\nwdocspar
\nwenddocs{}\nwbegincode{649}\moddef{Muli mom: interfaces}\plusendmoddef
  abstract interface
     elemental function qcd_get_real (this)
       import 
       class(qcd_2_2_class), intent(in) :: this
       real(default) :: qcd_get_real
     end function qcd_get_real
  end interface
\eatline
\nwendcode{}\nwbegindocs{650}\nwdocspar
\nwenddocs{}\nwbegincode{651}\moddef{Muli mom: interfaces}\plusendmoddef
  abstract interface
     pure function qcd_get_real_2 (this)
       import 
       class(qcd_2_2_class), intent(in) :: this
       real(default), dimension(2) :: qcd_get_real_2
     end function qcd_get_real_2
  end interface
\eatline
\nwendcode{}\nwbegindocs{652}\nwdocspar
\nwenddocs{}\nwbegincode{653}\moddef{Muli mom: interfaces}\plusendmoddef
  abstract interface
     pure function qcd_get_real_3 (this)
       import 
       class(qcd_2_2_class), intent(in) :: this
       real(default), dimension(3) :: qcd_get_real_3
     end function qcd_get_real_3
  end interface
\eatline
\nwendcode{}\nwbegindocs{654}\nwdocspar
\nwenddocs{}\nwbegincode{655}\moddef{Muli mom: interfaces}\plusendmoddef
  abstract interface
     elemental function qcd_get_int (this)
       import 
       class(qcd_2_2_class), intent(in) :: this
       integer :: qcd_get_int
     end function qcd_get_int
  end interface
\eatline
\nwendcode{}\nwbegindocs{656}\nwdocspar
\nwenddocs{}\nwbegincode{657}\moddef{Muli mom: interfaces}\plusendmoddef
  abstract interface
     pure function qcd_get_int_by_int (this, n)
       import 
       class(qcd_2_2_class), intent(in) :: this
       integer, intent(in) :: n
       integer :: qcd_get_int_by_int
     end function qcd_get_int_by_int
  end interface
\eatline
\nwendcode{}\nwbegindocs{658}\nwdocspar
\nwenddocs{}\nwbegincode{659}\moddef{Muli mom: interfaces}\plusendmoddef
  abstract interface
     pure function qcd_get_int_2 (this)
       import 
       class(qcd_2_2_class), intent(in) :: this
       integer, dimension(2) :: qcd_get_int_2
     end function qcd_get_int_2
  end interface
\eatline
\nwendcode{}\nwbegindocs{660}\nwdocspar
\nwenddocs{}\nwbegincode{661}\moddef{Muli mom: interfaces}\plusendmoddef
  abstract interface
     pure function qcd_get_int_4 (this)
       import 
       class(qcd_2_2_class), intent(in) :: this
       integer, dimension(4) :: qcd_get_int_4
     end function qcd_get_int_4
  end interface

\eatline
\nwendcode{}\nwbegindocs{662}\nwdocspar
%%%%%%%%%%%%%%%%%%%%%%%%%%%%%%%%%%%%%%%%%%%%%%%%%%%%%%%%%%%%%%%%%%%%%%%%
\section{Multi parton interactions}

This file contains the module {\tt{}muli{\char95}interactions}. The allowed
interactions and their cross sections are defined here. Additionaly,
some coordinate transformations which annihilate divergencies of the
cross sections are defined. Since the phase space border is
hyperbolic, this transformations are also hyperbolic. That's why all
interactions are named {\tt{}x{\char95}cart} for cartesian or {\tt{}x{\char95}hyp} for
hyperbolic to avoid confusion. 

\nwenddocs{}\nwbegincode{663}\moddef{\code{}muli{\char95}interactions.f90\edoc{}}\endmoddef
\LA{}File header\RA{}

module muli_interactions

\LA{}Use kinds with double\RA{}
  use constants
  use muli_momentum

\LA{}Standard module head\RA{}  

\LA{}Muli interactions: public\RA{}

\LA{}Muli interactions: variables\RA{}

\LA{}Muli interactions: interfaces\RA{}

contains
        
\LA{}Muli interactions: procedures\RA{}  
  
end module muli_interactions

\eatline
\nwendcode{}\nwbegindocs{664}\nwdocspar
\nwenddocs{}\nwbegincode{665}\moddef{Muli interactions: variables}\endmoddef
  character(len=2), dimension(-6:6), parameter :: integer_parton_names = &
       ["-6", "-5", "-4", "-3", "-2", "-1", "00", &
        "+1", "+2", "+3", "+4", "+5", "+6" ]
  character, dimension(-6:6),parameter :: traditional_parton_names = &
       ["T", "B", "C", "S", "U", "D", "g", "d", "u", "s", "c", "b", "t"]    
\eatline
\nwendcode{}\nwbegindocs{666}These are the phase space coefficients of the polynomial mappings,
the evolution variable is {\tt{}pt2s/(x1*x2)}.
\nwenddocs{}\nwbegincode{667}\moddef{Muli interactions: variables}\plusendmoddef
  real(default), dimension(1:4,1:5), parameter :: &
       phase_space_coefficients_in = reshape (source = &
       [ 6144, -4608,  +384,    0, &
         6144, -5120,  +384,    0, &
         6144, -2048,  +128, -576, &
        13824, -9600, +1056,    0, &
        31104,-19872, +2160, +486 ], shape=[4,5])
\eatline
\nwendcode{}\nwbegindocs{668}\nwdocspar
\nwenddocs{}\nwbegincode{669}\moddef{Muli interactions: variables}\plusendmoddef
  integer, parameter :: hadron_A_kind =  2212 
  integer, parameter :: hadron_B_kind = -2212 
  integer, dimension(4), parameter, public :: &
       parton_kind_of_int_kind = [1, 1, 2, 2]
  real(default), parameter :: b_sigma_tot_all = 100 !mb PDG
  real(default), parameter :: &
       b_sigma_tot_nd = 0.5*b_sigma_tot_all !!! PRD 49 n5 1994
  real(default), parameter, public :: &
       gev_cme_tot = 14000 ! total center of mass energie
  real(default), parameter :: gev2_cme_tot = gev_cme_tot**2 !!! s
  real(default), parameter :: gev_pt_max = gev_cme_tot/2D0
  real(default), parameter :: gev2_pt_max = gev2_cme_tot/4D0
  !model parameters
  real(default), parameter :: gev_pt_min = 8E-1_default
  real(default), parameter :: gev2_pt_min = gev_pt_min**2
  real(default), parameter :: pts_min = gev_pt_min / gev_pt_max
  real(default), parameter :: pts2_min = gev2_pt_min / gev2_pt_max
  real(default), parameter :: gev_p_t_0 = 2.0
  real(default), parameter :: gev2_p_t_0 = gev_p_t_0**2
  real(default), parameter :: norm_p_t_0 = gev_p_t_0 / gev_pt_max
  real(default), parameter :: norm2_p_t_0 = gev2_p_t_0 / gev2_pt_max
  !mathematical constants
  real(default), parameter, public :: euler = exp(one)
  !physical constants
  real(default), parameter :: gev2_mbarn = 0.389379304_default
  real(default), parameter :: const_pref = pi * gev2_mbarn / &
       (gev2_cme_tot * b_sigma_tot_nd)

\nwendcode{}\nwbegindocs{670}\nwdocspar
\nwenddocs{}\nwbegindocs{671} 
\nwenddocs{}\nwbegincode{672}\moddef{Muli interactions: variables}\plusendmoddef
  integer, parameter, public :: LHA_FLAVOR_AT = -6
  integer, parameter, public :: LHA_FLAVOR_AB = -5
  integer, parameter, public :: LHA_FLAVOR_AC = -4
  integer, parameter, public :: LHA_FLAVOR_AS = -3
  integer, parameter, public :: LHA_FLAVOR_AU = -2
  integer, parameter, public :: LHA_FLAVOR_AD = -1
  integer, parameter, public :: LHA_FLAVOR_G = 0
  integer, parameter, public :: LHA_FLAVOR_D = 1
  integer, parameter, public :: LHA_FLAVOR_U = 2
  integer, parameter, public :: LHA_FLAVOR_S = 3
  integer, parameter, public :: LHA_FLAVOR_C = 4
  integer, parameter, public :: LHA_FLAVOR_B = 5
  integer, parameter, public :: LHA_FLAVOR_T = 6

\eatline
\eatline
\eatline
\eatline
\nwendcode{}\nwbegindocs{673}\nwdocspar
\nwenddocs{}\nwbegincode{674}\moddef{Muli interactions: variables}\plusendmoddef
  integer, parameter, public :: PDG_FLAVOR_AT = -6
  integer, parameter, public :: PDG_FLAVOR_AB = -5
  integer, parameter, public :: PDG_FLAVOR_AC = -4
  integer, parameter, public :: PDG_FLAVOR_AS = -3
  integer, parameter, public :: PDG_FLAVOR_AU = -2
  integer, parameter, public :: PDG_FLAVOR_AD = -1
  integer, parameter, public :: PDG_FLAVOR_G = 21
  integer, parameter, public :: PDG_FLAVOR_D = 1
  integer, parameter, public :: PDG_FLAVOR_U = 2
  integer, parameter, public :: PDG_FLAVOR_S = 3
  integer, parameter, public :: PDG_FLAVOR_C = 4
  integer, parameter, public :: PDG_FLAVOR_B = 5
  integer, parameter, public :: PDG_FLAVOR_T = 6

\eatline
\eatline
\eatline
\eatline
\nwendcode{}\nwbegindocs{675}\nwdocspar
\nwenddocs{}\nwbegincode{676}\moddef{Muli interactions: variables}\plusendmoddef
  integer, parameter, public :: PARTON_SEA = 1
  integer, parameter, public :: PARTON_VALENCE = 2
  integer, parameter, public :: PARTON_SEA_AND_VALENCE = 3
  integer, parameter, public :: PARTON_TWIN = 4
  integer, parameter, public :: PARTON_SEA_AND_TWIN = 5
  integer, parameter, public :: PARTON_VALENCE_AND_TWIN = 6
  integer, parameter, public :: PARTON_ALL = 7

\eatline
\eatline
\nwendcode{}\nwbegindocs{677}\nwdocspar
\nwenddocs{}\nwbegincode{678}\moddef{Muli interactions: variables}\plusendmoddef
  integer, parameter, public :: PDF_UNDEFINED = 0
  integer, parameter, public :: PDF_GLUON = 1
  integer, parameter, public :: PDF_SEA = 2
  integer, parameter, public :: PDF_VALENCE_DOWN = 3
  integer, parameter, public :: PDF_VALENCE_UP = 4
  integer, parameter, public :: PDF_TWIN = 5  

\eatline
\eatline
\nwendcode{}\nwbegindocs{679}Evolution variable is {\tt{}pt2s/(x1*x2)}.
\nwenddocs{}\nwbegincode{680}\moddef{Muli interactions: variables}\plusendmoddef
  real(default), dimension(1:4,1:8),parameter :: &
       phase_space_coefficients_inout = reshape(source=[ &
       3072,  -2304,  +192,    0,  & 
       6144,  -5120,  +384,    0,  &
       0,         0,   192,  -96,  &
       3072,  -2048,  +192,  -96,  & 
       0,      2048, -2176, +576,  & 
       0,       288,  -306,  +81,  &
       6912,  -4800,  +528,    0,  &
       31104,-23328, +5832, -486], &
       shape=[4,8])  
  
\eatline
\nwendcode{}\nwbegindocs{681}\nwdocspar
\nwenddocs{}\nwbegincode{682}\moddef{Muli interactions: variables}\plusendmoddef
  integer, dimension(1:4,0:8), parameter :: inout_signatures = &
       reshape (source = [ &
        1, 1, 1, 1, &   !1a
       -1, 1,-1, 1, &   !1b
        1, 1, 1, 1, &   !2
        1,-1, 1,-1, &   !3
        1,-1, 1,-1, &   !4
        1,-1, 0, 0, &   !5
        0, 0, 1,-1, &   !6
        1, 0, 1, 0, &   !7
        0, 0, 0, 0  ], &
        shape = [4,9])
  
\eatline
\nwendcode{}\nwbegindocs{683}\nwdocspar
\nwenddocs{}\nwbegincode{684}\moddef{Muli interactions: variables}\plusendmoddef
  integer, dimension(6,-234:234), parameter, public :: &
       valid_processes = reshape( [ &
       -6,  -6,  -6,  -6,   2,   2, &   !-234
       -6,  -5,  -6,  -5,   1,   1, &   !-233
       -6,  -5,  -5,  -6,   1,   1, &   !-232
       -6,  -4,  -6,  -4,   1,   1, &   !-231
       -6,  -4,  -4,  -6,   1,   1, &   !-230
       -6,  -3,  -6,  -3,   1,   1, &   !-229
       -6,  -3,  -3,  -6,   1,   1, &   !-228
       -6,  -2,  -6,  -2,   1,   1, &   !-227
       -6,  -2,  -2,  -6,   1,   1, &   !-226
       -6,  -1,  -6,  -1,   1,   1, &   !-225
       -6,  -1,  -1,  -6,   1,   1, &   !-224
       -6,   0,  -6,   0,   4,   7, &   !-223
       -6,   0,   0,  -6,   4,   7, &   !-222
       -6,   1,  -6,   1,   1,   1, &   !-221
       -6,   1,   1,  -6,   1,   1, &   !-220
       -6,   2,  -6,   2,   1,   1, &   !-219
       -6,   2,   2,  -6,   1,   1, &   !-218
       -6,   3,  -6,   3,   1,   1, &   !-217
       -6,   3,   3,  -6,   1,   1, &   !-216
       -6,   4,  -6,   4,   1,   1, &   !-215
       -6,   4,   4,  -6,   1,   1, &   !-214
       -6,   5,  -6,   5,   1,   1, &   !-213
       -6,   5,   5,  -6,   1,   1, &   !-212
       -6,   6,  -6,   6,   3,   4, &   !-211
       -6,   6,  -5,   5,   3,   3, &   !-210
       -6,   6,  -4,   4,   3,   3, &   !-209
       -6,   6,  -3,   3,   3,   3, &   !-208
       -6,   6,  -2,   2,   3,   3, &   !-207
       -6,   6,  -1,   1,   3,   3, &   !-206
       -6,   6,   0,   0,   3,   5, &   !-205
       -6,   6,   1,  -1,   3,   3, &   !-204
       -6,   6,   2,  -2,   3,   3, &   !-203
       -6,   6,   3,  -3,   3,   3, &   !-202
       -6,   6,   4,  -4,   3,   3, &   !-201
       -6,   6,   5,  -5,   3,   3, &   !-200
       -6,   6,   6,  -6,   3,   4, &   !-199
       -5,  -6,  -6,  -5,   1,   1, &   !-198
       -5,  -6,  -5,  -6,   1,   1, &   !-197
       -5,  -5,  -5,  -5,   2,   2, &   !-196
       -5,  -4,  -5,  -4,   1,   1, &   !-195
       -5,  -4,  -4,  -5,   1,   1, &   !-194
       -5,  -3,  -5,  -3,   1,   1, &   !-193
       -5,  -3,  -3,  -5,   1,   1, &   !-192
       -5,  -2,  -5,  -2,   1,   1, &   !-191
       -5,  -2,  -2,  -5,   1,   1, &   !-190
       -5,  -1,  -5,  -1,   1,   1, &   !-189
       -5,  -1,  -1,  -5,   1,   1, &   !-188
       -5,   0,  -5,   0,   4,   7, &   !-187
       -5,   0,   0,  -5,   4,   7, &   !-186
       -5,   1,  -5,   1,   1,   1, &   !-185
       -5,   1,   1,  -5,   1,   1, &   !-184
       -5,   2,  -5,   2,   1,   1, &   !-183
       -5,   2,   2,  -5,   1,   1, &   !-182
       -5,   3,  -5,   3,   1,   1, &   !-181
       -5,   3,   3,  -5,   1,   1, &   !-180
       -5,   4,  -5,   4,   1,   1, &   !-179
       -5,   4,   4,  -5,   1,   1, &   !-178
       -5,   5,  -6,   6,   3,   3, &   !-177
       -5,   5,  -5,   5,   3,   4, &   !-176
       -5,   5,  -4,   4,   3,   3, &   !-175
       -5,   5,  -3,   3,   3,   3, &   !-174
       -5,   5,  -2,   2,   3,   3, &   !-173
       -5,   5,  -1,   1,   3,   3, &   !-172
       -5,   5,   0,   0,   3,   5, &   !-171
       -5,   5,   1,  -1,   3,   3, &   !-170
       -5,   5,   2,  -2,   3,   3, &   !-169
       -5,   5,   3,  -3,   3,   3, &   !-168
       -5,   5,   4,  -4,   3,   3, &   !-167
       -5,   5,   5,  -5,   3,   4, &   !-166
       -5,   5,   6,  -6,   3,   3, &   !-165
       -5,   6,  -5,   6,   1,   1, &   !-164
       -5,   6,   6,  -5,   1,   1, &   !-163
       -4,  -6,  -6,  -4,   1,   1, &   !-162
       -4,  -6,  -4,  -6,   1,   1, &   !-161
       -4,  -5,  -5,  -4,   1,   1, &   !-160
       -4,  -5,  -4,  -5,   1,   1, &   !-159
       -4,  -4,  -4,  -4,   2,   2, &   !-158
       -4,  -3,  -4,  -3,   1,   1, &   !-157
       -4,  -3,  -3,  -4,   1,   1, &   !-156
       -4,  -2,  -4,  -2,   1,   1, &   !-155
       -4,  -2,  -2,  -4,   1,   1, &   !-154
       -4,  -1,  -4,  -1,   1,   1, &   !-153
       -4,  -1,  -1,  -4,   1,   1, &   !-152
       -4,   0,  -4,   0,   4,   7, &   !-151
       -4,   0,   0,  -4,   4,   7, &   !-150
       -4,   1,  -4,   1,   1,   1, &   !-149
       -4,   1,   1,  -4,   1,   1, &   !-148
       -4,   2,  -4,   2,   1,   1, &   !-147
       -4,   2,   2,  -4,   1,   1, &   !-146
       -4,   3,  -4,   3,   1,   1, &   !-145
       -4,   3,   3,  -4,   1,   1, &   !-144
       -4,   4,  -6,   6,   3,   3, &   !-143
       -4,   4,  -5,   5,   3,   3, &   !-142
       -4,   4,  -4,   4,   3,   4, &   !-141
       -4,   4,  -3,   3,   3,   3, &   !-140
       -4,   4,  -2,   2,   3,   3, &   !-139
       -4,   4,  -1,   1,   3,   3, &   !-138
       -4,   4,   0,   0,   3,   5, &   !-137
       -4,   4,   1,  -1,   3,   3, &   !-136
       -4,   4,   2,  -2,   3,   3, &   !-135
       -4,   4,   3,  -3,   3,   3, &   !-134
       -4,   4,   4,  -4,   3,   4, &   !-133
       -4,   4,   5,  -5,   3,   3, &   !-132
       -4,   4,   6,  -6,   3,   3, &   !-131
       -4,   5,  -4,   5,   1,   1, &   !-130
       -4,   5,   5,  -4,   1,   1, &   !-129
       -4,   6,  -4,   6,   1,   1, &   !-128
       -4,   6,   6,  -4,   1,   1, &   !-127
       -3,  -6,  -6,  -3,   1,   1, &   !-126
       -3,  -6,  -3,  -6,   1,   1, &   !-125
       -3,  -5,  -5,  -3,   1,   1, &   !-124
       -3,  -5,  -3,  -5,   1,   1, &   !-123
       -3,  -4,  -4,  -3,   1,   1, &   !-122
       -3,  -4,  -3,  -4,   1,   1, &   !-121
       -3,  -3,  -3,  -3,   2,   2, &   !-120
       -3,  -2,  -3,  -2,   1,   1, &   !-119
       -3,  -2,  -2,  -3,   1,   1, &   !-118
       -3,  -1,  -3,  -1,   1,   1, &   !-117
       -3,  -1,  -1,  -3,   1,   1, &   !-116
       -3,   0,  -3,   0,   4,   7, &   !-115
       -3,   0,   0,  -3,   4,   7, &   !-114
       -3,   1,  -3,   1,   1,   1, &   !-113
       -3,   1,   1,  -3,   1,   1, &   !-112
       -3,   2,  -3,   2,   1,   1, &   !-111
       -3,   2,   2,  -3,   1,   1, &   !-110
       -3,   3,  -6,   6,   3,   3, &   !-109
       -3,   3,  -5,   5,   3,   3, &   !-108
       -3,   3,  -4,   4,   3,   3, &   !-107
       -3,   3,  -3,   3,   3,   4, &   !-106
       -3,   3,  -2,   2,   3,   3, &   !-105
       -3,   3,  -1,   1,   3,   3, &   !-104
       -3,   3,   0,   0,   3,   5, &   !-103
       -3,   3,   1,  -1,   3,   3, &   !-102
       -3,   3,   2,  -2,   3,   3, &   !-101
       -3,   3,   3,  -3,   3,   4, &   !-100
       -3,   3,   4,  -4,   3,   3, &   ! -99
       -3,   3,   5,  -5,   3,   3, &   ! -98
       -3,   3,   6,  -6,   3,   3, &   ! -97
       -3,   4,  -3,   4,   1,   1, &   ! -96
       -3,   4,   4,  -3,   1,   1, &   ! -95
       -3,   5,  -3,   5,   1,   1, &   ! -94
       -3,   5,   5,  -3,   1,   1, &   ! -93
       -3,   6,  -3,   6,   1,   1, &   ! -92
       -3,   6,   6,  -3,   1,   1, &   ! -91
       -2,  -6,  -6,  -2,   1,   1, &   ! -90
       -2,  -6,  -2,  -6,   1,   1, &   ! -89
       -2,  -5,  -5,  -2,   1,   1, &   ! -88
       -2,  -5,  -2,  -5,   1,   1, &   ! -87
       -2,  -4,  -4,  -2,   1,   1, &   ! -86
       -2,  -4,  -2,  -4,   1,   1, &   ! -85
       -2,  -3,  -3,  -2,   1,   1, &   ! -84
       -2,  -3,  -2,  -3,   1,   1, &   ! -83
       -2,  -2,  -2,  -2,   2,   2, &   ! -82
       -2,  -1,  -2,  -1,   1,   1, &   ! -81
       -2,  -1,  -1,  -2,   1,   1, &   ! -80
       -2,   0,  -2,   0,   4,   7, &   ! -79
       -2,   0,   0,  -2,   4,   7, &   ! -78
       -2,   1,  -2,   1,   1,   1, &   ! -77
       -2,   1,   1,  -2,   1,   1, &   ! -76
       -2,   2,  -6,   6,   3,   3, &   ! -75
       -2,   2,  -5,   5,   3,   3, &   ! -74
       -2,   2,  -4,   4,   3,   3, &   ! -73
       -2,   2,  -3,   3,   3,   3, &   ! -72
       -2,   2,  -2,   2,   3,   4, &   ! -71
       -2,   2,  -1,   1,   3,   3, &   ! -70
       -2,   2,   0,   0,   3,   5, &   ! -69
       -2,   2,   1,  -1,   3,   3, &   ! -68
       -2,   2,   2,  -2,   3,   4, &   ! -67
       -2,   2,   3,  -3,   3,   3, &   ! -66
       -2,   2,   4,  -4,   3,   3, &   ! -65
       -2,   2,   5,  -5,   3,   3, &   ! -64
       -2,   2,   6,  -6,   3,   3, &   ! -63
       -2,   3,  -2,   3,   1,   1, &   ! -62
       -2,   3,   3,  -2,   1,   1, &   ! -61
       -2,   4,  -2,   4,   1,   1, &   ! -60
       -2,   4,   4,  -2,   1,   1, &   ! -59
       -2,   5,  -2,   5,   1,   1, &   ! -58
       -2,   5,   5,  -2,   1,   1, &   ! -57
       -2,   6,  -2,   6,   1,   1, &   ! -56
       -2,   6,   6,  -2,   1,   1, &   ! -55
       -1,  -6,  -6,  -1,   1,   1, &   ! -54
       -1,  -6,  -1,  -6,   1,   1, &   ! -53
       -1,  -5,  -5,  -1,   1,   1, &   ! -52
       -1,  -5,  -1,  -5,   1,   1, &   ! -51
       -1,  -4,  -4,  -1,   1,   1, &   ! -50
       -1,  -4,  -1,  -4,   1,   1, &   ! -49
       -1,  -3,  -3,  -1,   1,   1, &   ! -48
       -1,  -3,  -1,  -3,   1,   1, &   ! -47
       -1,  -2,  -2,  -1,   1,   1, &   ! -46
       -1,  -2,  -1,  -2,   1,   1, &   ! -45
       -1,  -1,  -1,  -1,   2,   2, &   ! -44
       -1,   0,  -1,   0,   4,   7, &   ! -43
       -1,   0,   0,  -1,   4,   7, &   ! -42
       -1,   1,  -6,   6,   3,   3, &   ! -41
       -1,   1,  -5,   5,   3,   3, &   ! -40
       -1,   1,  -4,   4,   3,   3, &   ! -39
       -1,   1,  -3,   3,   3,   3, &   ! -38
       -1,   1,  -2,   2,   3,   3, &   ! -37
       -1,   1,  -1,   1,   3,   4, &   ! -36
       -1,   1,   0,   0,   3,   5, &   ! -35
       -1,   1,   1,  -1,   3,   4, &   ! -34
       -1,   1,   2,  -2,   3,   3, &   ! -33
       -1,   1,   3,  -3,   3,   3, &   ! -32
       -1,   1,   4,  -4,   3,   3, &   ! -31
       -1,   1,   5,  -5,   3,   3, &   ! -30
       -1,   1,   6,  -6,   3,   3, &   ! -29
       -1,   2,  -1,   2,   1,   1, &   ! -28
       -1,   2,   2,  -1,   1,   1, &   ! -27
       -1,   3,  -1,   3,   1,   1, &   ! -26
       -1,   3,   3,  -1,   1,   1, &   ! -25
       -1,   4,  -1,   4,   1,   1, &   ! -24
       -1,   4,   4,  -1,   1,   1, &   ! -23
       -1,   5,  -1,   5,   1,   1, &   ! -22
       -1,   5,   5,  -1,   1,   1, &   ! -21
       -1,   6,  -1,   6,   1,   1, &   ! -20
       -1,   6,   6,  -1,   1,   1, &   ! -19
        0,  -6,  -6,   0,   4,   7, &   ! -18
        0,  -6,   0,  -6,   4,   7, &   ! -17
        0,  -5,  -5,   0,   4,   7, &   ! -16
        0,  -5,   0,  -5,   4,   7, &   ! -15
        0,  -4,  -4,   0,   4,   7, &   ! -14
        0,  -4,   0,  -4,   4,   7, &   ! -13
        0,  -3,  -3,   0,   4,   7, &   ! -12
        0,  -3,   0,  -3,   4,   7, &   ! -11
        0,  -2,  -2,   0,   4,   7, &   ! -10
        0,  -2,   0,  -2,   4,   7, &   !  -9
        0,  -1,  -1,   0,   4,   7, &   !  -8
        0,  -1,   0,  -1,   4,   7, &   !  -7
        0,   0,  -6,   6,   5,   6, &   !  -6
        0,   0,  -5,   5,   5,   6, &   !  -5
        0,   0,  -4,   4,   5,   6, &   !  -4
        0,   0,  -3,   3,   5,   6, &   !  -3
        0,   0,  -2,   2,   5,   6, &   !  -2
        0,   0,  -1,   1,   5,   6, &   !  -1
        0,   0,   0,   0,   5,   8, &   !   0
        0,   0,   1,  -1,   5,   6, &   !   1
        0,   0,   2,  -2,   5,   6, &   !   2
        0,   0,   3,  -3,   5,   6, &   !   3
        0,   0,   4,  -4,   5,   6, &   !   4
        0,   0,   5,  -5,   5,   6, &   !   5
        0,   0,   6,  -6,   5,   6, &   !   6
        0,   1,   0,   1,   4,   7, &   !   7
        0,   1,   1,   0,   4,   7, &   !   8
        0,   2,   0,   2,   4,   7, &   !   9
        0,   2,   2,   0,   4,   7, &   !  10
        0,   3,   0,   3,   4,   7, &   !  11
        0,   3,   3,   0,   4,   7, &   !  12
        0,   4,   0,   4,   4,   7, &   !  13
        0,   4,   4,   0,   4,   7, &   !  14
        0,   5,   0,   5,   4,   7, &   !  15
        0,   5,   5,   0,   4,   7, &   !  16
        0,   6,   0,   6,   4,   7, &   !  17
        0,   6,   6,   0,   4,   7, &   !  18
        1,  -6,  -6,   1,   1,   1, &   !  19
        1,  -6,   1,  -6,   1,   1, &   !  20
        1,  -5,  -5,   1,   1,   1, &   !  21
        1,  -5,   1,  -5,   1,   1, &   !  22
        1,  -4,  -4,   1,   1,   1, &   !  23
        1,  -4,   1,  -4,   1,   1, &   !  24
        1,  -3,  -3,   1,   1,   1, &   !  25
        1,  -3,   1,  -3,   1,   1, &   !  26
        1,  -2,  -2,   1,   1,   1, &   !  27
        1,  -2,   1,  -2,   1,   1, &   !  28
        1,  -1,  -6,   6,   3,   3, &   !  29
        1,  -1,  -5,   5,   3,   3, &   !  30
        1,  -1,  -4,   4,   3,   3, &   !  31
        1,  -1,  -3,   3,   3,   3, &   !  32
        1,  -1,  -2,   2,   3,   3, &   !  33
        1,  -1,  -1,   1,   3,   4, &   !  34
        1,  -1,   0,   0,   3,   5, &   !  35
        1,  -1,   1,  -1,   3,   4, &   !  36
        1,  -1,   2,  -2,   3,   3, &   !  37
        1,  -1,   3,  -3,   3,   3, &   !  38
        1,  -1,   4,  -4,   3,   3, &   !  39
        1,  -1,   5,  -5,   3,   3, &   !  40
        1,  -1,   6,  -6,   3,   3, &   !  41
        1,   0,   0,   1,   4,   7, &   !  42
        1,   0,   1,   0,   4,   7, &   !  43
        1,   1,   1,   1,   2,   2, &   !  44
        1,   2,   1,   2,   1,   1, &   !  45
        1,   2,   2,   1,   1,   1, &   !  46
        1,   3,   1,   3,   1,   1, &   !  47
        1,   3,   3,   1,   1,   1, &   !  48
        1,   4,   1,   4,   1,   1, &   !  49
        1,   4,   4,   1,   1,   1, &   !  50
        1,   5,   1,   5,   1,   1, &   !  51
        1,   5,   5,   1,   1,   1, &   !  52
        1,   6,   1,   6,   1,   1, &   !  53
        1,   6,   6,   1,   1,   1, &   !  54
        2,  -6,  -6,   2,   1,   1, &   !  55
        2,  -6,   2,  -6,   1,   1, &   !  56
        2,  -5,  -5,   2,   1,   1, &   !  57
        2,  -5,   2,  -5,   1,   1, &   !  58
        2,  -4,  -4,   2,   1,   1, &   !  59
        2,  -4,   2,  -4,   1,   1, &   !  60
        2,  -3,  -3,   2,   1,   1, &   !  61
        2,  -3,   2,  -3,   1,   1, &   !  62
        2,  -2,  -6,   6,   3,   3, &   !  63
        2,  -2,  -5,   5,   3,   3, &   !  64
        2,  -2,  -4,   4,   3,   3, &   !  65
        2,  -2,  -3,   3,   3,   3, &   !  66
        2,  -2,  -2,   2,   3,   4, &   !  67
        2,  -2,  -1,   1,   3,   3, &   !  68
        2,  -2,   0,   0,   3,   5, &   !  69
        2,  -2,   1,  -1,   3,   3, &   !  70
        2,  -2,   2,  -2,   3,   4, &   !  71
        2,  -2,   3,  -3,   3,   3, &   !  72
        2,  -2,   4,  -4,   3,   3, &   !  73
        2,  -2,   5,  -5,   3,   3, &   !  74
        2,  -2,   6,  -6,   3,   3, &   !  75
        2,  -1,  -1,   2,   1,   1, &   !  76
        2,  -1,   2,  -1,   1,   1, &   !  77
        2,   0,   0,   2,   4,   7, &   !  78
        2,   0,   2,   0,   4,   7, &   !  79
        2,   1,   1,   2,   1,   1, &   !  80
        2,   1,   2,   1,   1,   1, &   !  81
        2,   2,   2,   2,   2,   2, &   !  82
        2,   3,   2,   3,   1,   1, &   !  83
        2,   3,   3,   2,   1,   1, &   !  84
        2,   4,   2,   4,   1,   1, &   !  85
        2,   4,   4,   2,   1,   1, &   !  86
        2,   5,   2,   5,   1,   1, &   !  87
        2,   5,   5,   2,   1,   1, &   !  88
        2,   6,   2,   6,   1,   1, &   !  89
        2,   6,   6,   2,   1,   1, &   !  90
        3,  -6,  -6,   3,   1,   1, &   !  91
        3,  -6,   3,  -6,   1,   1, &   !  92
        3,  -5,  -5,   3,   1,   1, &   !  93
        3,  -5,   3,  -5,   1,   1, &   !  94
        3,  -4,  -4,   3,   1,   1, &   !  95
        3,  -4,   3,  -4,   1,   1, &   !  96
        3,  -3,  -6,   6,   3,   3, &   !  97
        3,  -3,  -5,   5,   3,   3, &   !  98
        3,  -3,  -4,   4,   3,   3, &   !  99
        3,  -3,  -3,   3,   3,   4, &   ! 100
        3,  -3,  -2,   2,   3,   3, &   ! 101
        3,  -3,  -1,   1,   3,   3, &   ! 102
        3,  -3,   0,   0,   3,   5, &   ! 103
        3,  -3,   1,  -1,   3,   3, &   ! 104
        3,  -3,   2,  -2,   3,   3, &   ! 105
        3,  -3,   3,  -3,   3,   4, &   ! 106
        3,  -3,   4,  -4,   3,   3, &   ! 107
        3,  -3,   5,  -5,   3,   3, &   ! 108
        3,  -3,   6,  -6,   3,   3, &   ! 109
        3,  -2,  -2,   3,   1,   1, &   ! 110
        3,  -2,   3,  -2,   1,   1, &   ! 111
        3,  -1,  -1,   3,   1,   1, &   ! 112
        3,  -1,   3,  -1,   1,   1, &   ! 113
        3,   0,   0,   3,   4,   7, &   ! 114
        3,   0,   3,   0,   4,   7, &   ! 115
        3,   1,   1,   3,   1,   1, &   ! 116
        3,   1,   3,   1,   1,   1, &   ! 117
        3,   2,   2,   3,   1,   1, &   ! 118
        3,   2,   3,   2,   1,   1, &   ! 119
        3,   3,   3,   3,   2,   2, &   ! 120
        3,   4,   3,   4,   1,   1, &   ! 121
        3,   4,   4,   3,   1,   1, &   ! 122
        3,   5,   3,   5,   1,   1, &   ! 123
        3,   5,   5,   3,   1,   1, &   ! 124
        3,   6,   3,   6,   1,   1, &   ! 125
        3,   6,   6,   3,   1,   1, &   ! 126
        4,  -6,  -6,   4,   1,   1, &   ! 127
        4,  -6,   4,  -6,   1,   1, &   ! 128
        4,  -5,  -5,   4,   1,   1, &   ! 129
        4,  -5,   4,  -5,   1,   1, &   ! 130
        4,  -4,  -6,   6,   3,   3, &   ! 131
        4,  -4,  -5,   5,   3,   3, &   ! 132
        4,  -4,  -4,   4,   3,   4, &   ! 133
        4,  -4,  -3,   3,   3,   3, &   ! 134
        4,  -4,  -2,   2,   3,   3, &   ! 135
        4,  -4,  -1,   1,   3,   3, &   ! 136
        4,  -4,   0,   0,   3,   5, &   ! 137
        4,  -4,   1,  -1,   3,   3, &   ! 138
        4,  -4,   2,  -2,   3,   3, &   ! 139
        4,  -4,   3,  -3,   3,   3, &   ! 140
        4,  -4,   4,  -4,   3,   4, &   ! 141
        4,  -4,   5,  -5,   3,   3, &   ! 142
        4,  -4,   6,  -6,   3,   3, &   ! 143
        4,  -3,  -3,   4,   1,   1, &   ! 144
        4,  -3,   4,  -3,   1,   1, &   ! 145
        4,  -2,  -2,   4,   1,   1, &   ! 146
        4,  -2,   4,  -2,   1,   1, &   ! 147
        4,  -1,  -1,   4,   1,   1, &   ! 148
        4,  -1,   4,  -1,   1,   1, &   ! 149
        4,   0,   0,   4,   4,   7, &   ! 150
        4,   0,   4,   0,   4,   7, &   ! 151
        4,   1,   1,   4,   1,   1, &   ! 152
        4,   1,   4,   1,   1,   1, &   ! 153
        4,   2,   2,   4,   1,   1, &   ! 154
        4,   2,   4,   2,   1,   1, &   ! 155
        4,   3,   3,   4,   1,   1, &   ! 156
        4,   3,   4,   3,   1,   1, &   ! 157
        4,   4,   4,   4,   2,   2, &   ! 158
        4,   5,   4,   5,   1,   1, &   ! 159
        4,   5,   5,   4,   1,   1, &   ! 160
        4,   6,   4,   6,   1,   1, &   ! 161
        4,   6,   6,   4,   1,   1, &   ! 162
        5,  -6,  -6,   5,   1,   1, &   ! 163
        5,  -6,   5,  -6,   1,   1, &   ! 164
        5,  -5,  -6,   6,   3,   3, &   ! 165
        5,  -5,  -5,   5,   3,   4, &   ! 166
        5,  -5,  -4,   4,   3,   3, &   ! 167
        5,  -5,  -3,   3,   3,   3, &   ! 168
        5,  -5,  -2,   2,   3,   3, &   ! 169
        5,  -5,  -1,   1,   3,   3, &   ! 170
        5,  -5,   0,   0,   3,   5, &   ! 171
        5,  -5,   1,  -1,   3,   3, &   ! 172
        5,  -5,   2,  -2,   3,   3, &   ! 173
        5,  -5,   3,  -3,   3,   3, &   ! 174
        5,  -5,   4,  -4,   3,   3, &   ! 175
        5,  -5,   5,  -5,   3,   4, &   ! 176
        5,  -5,   6,  -6,   3,   3, &   ! 177
        5,  -4,  -4,   5,   1,   1, &   ! 178
        5,  -4,   5,  -4,   1,   1, &   ! 179
        5,  -3,  -3,   5,   1,   1, &   ! 180
        5,  -3,   5,  -3,   1,   1, &   ! 181
        5,  -2,  -2,   5,   1,   1, &   ! 182
        5,  -2,   5,  -2,   1,   1, &   ! 183
        5,  -1,  -1,   5,   1,   1, &   ! 184
        5,  -1,   5,  -1,   1,   1, &   ! 185
        5,   0,   0,   5,   4,   7, &   ! 186
        5,   0,   5,   0,   4,   7, &   ! 187
        5,   1,   1,   5,   1,   1, &   ! 188
        5,   1,   5,   1,   1,   1, &   ! 189
        5,   2,   2,   5,   1,   1, &   ! 190
        5,   2,   5,   2,   1,   1, &   ! 191
        5,   3,   3,   5,   1,   1, &   ! 192
        5,   3,   5,   3,   1,   1, &   ! 193
        5,   4,   4,   5,   1,   1, &   ! 194
        5,   4,   5,   4,   1,   1, &   ! 195
        5,   5,   5,   5,   2,   2, &   ! 196
        5,   6,   5,   6,   1,   1, &   ! 197
        5,   6,   6,   5,   1,   1, &   ! 198
        6,  -6,  -6,   6,   3,   4, &   ! 199
        6,  -6,  -5,   5,   3,   3, &   ! 200
        6,  -6,  -4,   4,   3,   3, &   ! 201
        6,  -6,  -3,   3,   3,   3, &   ! 202
        6,  -6,  -2,   2,   3,   3, &   ! 203
        6,  -6,  -1,   1,   3,   3, &   ! 204
        6,  -6,   0,   0,   3,   5, &   ! 205
        6,  -6,   1,  -1,   3,   3, &   ! 206
        6,  -6,   2,  -2,   3,   3, &   ! 207
        6,  -6,   3,  -3,   3,   3, &   ! 208
        6,  -6,   4,  -4,   3,   3, &   ! 209
        6,  -6,   5,  -5,   3,   3, &   ! 210
        6,  -6,   6,  -6,   3,   4, &   ! 211
        6,  -5,  -5,   6,   1,   1, &   ! 212
        6,  -5,   6,  -5,   1,   1, &   ! 213
        6,  -4,  -4,   6,   1,   1, &   ! 214
        6,  -4,   6,  -4,   1,   1, &   ! 215
        6,  -3,  -3,   6,   1,   1, &   ! 216
        6,  -3,   6,  -3,   1,   1, &   ! 217
        6,  -2,  -2,   6,   1,   1, &   ! 218
        6,  -2,   6,  -2,   1,   1, &   ! 219
        6,  -1,  -1,   6,   1,   1, &   ! 220
        6,  -1,   6,  -1,   1,   1, &   ! 221
        6,   0,   0,   6,   4,   7, &   ! 222
        6,   0,   6,   0,   4,   7, &   ! 223
        6,   1,   1,   6,   1,   1, &   ! 224
        6,   1,   6,   1,   1,   1, &   ! 225
        6,   2,   2,   6,   1,   1, &   ! 226
        6,   2,   6,   2,   1,   1, &   ! 227
        6,   3,   3,   6,   1,   1, &   ! 228
        6,   3,   6,   3,   1,   1, &   ! 229
        6,   4,   4,   6,   1,   1, &   ! 230
        6,   4,   6,   4,   1,   1, &   ! 231
        6,   5,   5,   6,   1,   1, &   ! 232
        6,   5,   6,   5,   1,   1, &   ! 233
        6,   6,   6,   6,   2,   2  &   ! 234
       ], [6, 469])
  
\nwendcode{}\nwbegindocs{685}% valid_processes
\nwenddocs{}\nwbegindocs{686}\nwdocspar
\nwenddocs{}\nwbegincode{687}\moddef{Muli interactions: variables}\plusendmoddef
  integer, dimension(2,0:16), parameter, public :: &
       double_pdf_kinds = reshape ( [ &
       0, 0, &
       1, 1, &
       1, 2, &
       1, 3, &
       1, 4, &
       2, 1, &
       2, 2, &
       2, 3, &
       2, 4, &
       3, 1, &
       3, 2, &
       3, 3, &
       3, 4, &
       4, 1, &
       4, 2, &
       4, 3, &
       4, 4], [2, 17])
  
\eatline
\nwendcode{}\nwbegindocs{688}\nwdocspar
\nwenddocs{}\nwbegincode{689}\moddef{Muli interactions: variables}\plusendmoddef
  integer, parameter, dimension(371), public :: int_all = [ &
         -6,  -5,  -4,  -3,  -2,  -1,   0,  1 ,   2, &
          3,   4,   5,   6, -14, -13, -12, -11, -10, &
          9,  -8,  -7,   7,   8,   9,  10,  11,  12, &
         13,  14,   7,   8,   9,  10,-151,-150,-115, &
       -114, -79, -78, -43, -42,  42,  43,  78,  79, & 
        114, 115, 150, 151,-158,-157,-156,-155,-154, &
       -153,-152,-149,-148,-147,-146,-145,-144,-143, &
       -142,-141,-140,-139,-138,-137,-136,-135,-134, &
       -133,-132,-131,-122,-121,-120,-119,-118,-117, &
       -116,-113,-112,-111,-110,-109,-108,-107,-106, &
       -105,-104,-103,-102,-101,-100, -99, -98, -97, &
        -96, -95, -86, -85, -84, -83, -82, -81, -80, &
        -77, -76, -75, -74, -73, -72, -71, -70, -69, &
        -68, -67, -66, -65, -64, -63, -62, -61, -60, &
        -59, -50, -49, -48, -47, -46, -45, -44, -41, &
        -40, -39, -38, -37, -36, -35, -34, -33, -32, &
        -31, -30, -29, -28, -27, -26, -25, -24, -23, &
         23,  24,  25,  26,  27,  28,  29,  30,  31, &
         32,  33,  34,  35,  36,  37,  38,  39,  40, &
         41,  44,  45,  46,  47,  48,  49,  50,  59, &
         60,  61,  62,  63,  64,  65,  66,  67,  68, &
         69,  70,  71,  72,  73,  74,  75,  76,  77, &
         80,  81,  82,  83,  84,  85,  86,  95,  96, &
         97,  98,  99, 100, 101, 102, 103, 104, 105, &
        106, 107, 108, 109, 110, 111, 112, 113, 116, &
        117, 118, 119, 120, 121, 122, 131, 132, 133, &
        134, 135, 136, 137, 138, 139, 140, 141, 142, &
        143, 144, 145, 146, 147, 148, 149, 152, 153, &
        154, 155, 156, 157, 158,-149,-148,-113,-112, &
        -77, -76, -41, -40, -39, -38, -37, -36, -35, &
        -34, -33, -32, -31, -30, -29,  44,  80,  81, &
        116, 117, 152, 153,-147,-146,-111,-110, -75, &
        -74, -73, -72, -71, -70, -69, -68, -67, -66, &
        -65, -64, -63, -28, -27,  45,  46,  82, 118, &
        119, 154, 155,  42,  43,  23,  24,  25,  26, &
         27,  28,  29,  30,  31,  32,  33,  34,  35, &
         36,  37,  38,  39,  40,  41,  44,  45,  46, &
         47,  48,  49,  50,  44,  45,  46,  78,  79, &
         59,  60,  61,  62,  63,  64,  65,  66,  67, &
         68,  69,  70,  71,  72,  73,  74,  75,  76, &
         77,  80,  81,  82,  83,  84,  85,  86,  80, &
         81,  82 ]
  
\eatline
\nwendcode{}\nwbegindocs{690}\nwdocspar
\nwenddocs{}\nwbegincode{691}\moddef{Muli interactions: variables}\plusendmoddef
  integer, parameter, dimension(16), public :: int_sizes_all = &
       [13, 16, 2, 2, 16, 208, 26, 26, 2, 26, 1, 2, 2, 26, 2, 1]
  
\eatline
\nwendcode{}\nwbegindocs{692}\nwdocspar
\nwenddocs{}\nwbegincode{693}\moddef{Muli interactions: variables}\plusendmoddef
  integer, parameter, dimension(3,0:8), public :: muli_flow_stats = &
       reshape( [ &
           1,  2,  4,   &
           3,  4,  4,   &
           5,  6,  8,   &
           7,  8,  4,   &
           9, 10,  8,   &
          11, 16, 16,   &
          17, 22, 16,   &
          23, 28, 16,   &
          29, 52, 96 ], &
          [3,9])
  
\eatline
\nwendcode{}\nwbegindocs{694}\nwdocspar
\nwenddocs{}\nwbegincode{695}\moddef{Muli interactions: variables}\plusendmoddef
  integer, parameter, dimension(0:4,52), public :: muli_flows = &
       reshape( [ &
       3, 0, 0, 1, 2, &   !1a
       1, 0, 0, 2, 1, &
       1, 2, 0, 0, 3, &   !1b
       3, 3, 0, 0, 2, &
       4, 0, 0, 1, 2, &   !2
       4, 0, 0, 2, 1, &
       3, 2, 0, 0, 3, &   !3
       1, 3, 0, 0, 2, &
       4, 2, 0, 0, 3, &   !4
       4, 3, 0, 0, 2, &
       4, 0, 1, 3, 4, &   !5
       4, 0, 1, 4, 3, &
       2, 0, 3, 1, 4, &
       2, 0, 4, 1, 3, &
       2, 0, 3, 4, 1, &
       2, 0, 4, 3, 1, &
       4, 1, 2, 4, 0, &    !6
       2, 1, 4, 2, 0, &
       4, 2, 1, 4, 0, &
       2, 4, 1, 2, 0, &
       2, 2, 4, 1, 0, &
       2, 4, 2, 1, 0, &
       2, 0, 1, 2, 4, &    !7
       2, 0, 1, 4, 2, &
       4, 0, 2, 1, 4, &
       4, 0, 4, 1, 2, &
       2, 0, 2, 4, 1, &
       2, 0, 4, 2, 1, &
       9, 1, 2, 3, 4, &    !8
       5, 1, 2, 4, 3, &
       5, 1, 3, 2, 4, &
       3, 1, 4, 2, 3, &
       3, 1, 3, 4, 2, &
       5, 1, 4, 3, 2, &
       5, 2, 1, 3, 4, &
       5, 2, 1, 4, 3, &
       3, 3, 1, 2, 4, &
       3, 4, 1, 2, 3, &
       3, 3, 1, 4, 2, &
       3, 4, 1, 3, 2, &
       3, 2, 3, 1, 4, &
       3, 2, 4, 1, 3, &
       5, 3, 2, 1, 4, &
       3, 4, 2, 1, 3, &
       5, 3, 4, 1, 2, &
       3, 4, 3, 1, 2, &
       3, 2, 3, 4, 1, &
       3, 2, 4, 3, 1, &
       3, 3, 2, 4, 1, &
       5, 4, 2, 3, 1, &
       3, 3, 4, 2, 1, &
       5, 4, 3, 2, 1], [5, 52])

\eatline
\nwendcode{}\nwbegindocs{696}This value, {\tt{}pts2{\char95}scale}, seems to be nowhere set in the code.
\nwenddocs{}\nwbegincode{697}\moddef{Muli interactions: variables}\plusendmoddef
  real(default) :: pts2_scale  
\nwendcode{}\nwbegindocs{698}\nwdocspar
\nwenddocs{}\nwbegincode{699}\moddef{Muli interactions: interfaces}\endmoddef
  abstract interface
     function trafo_in (in) 
       use kinds !NODEP!
       real(default), dimension(3) :: trafo_in
       real(default), dimension(3), intent(in) :: in
     end function trafo_in
  end interface
\eatline
\nwendcode{}\nwbegindocs{700}\nwdocspar
\nwenddocs{}\nwbegincode{701}\moddef{Muli interactions: interfaces}\plusendmoddef
  abstract interface
     pure function coord_scalar_in (hyp)
       use kinds !NODEP!
       real(default) :: coord_scalar_in
       real(double), dimension(3), intent(in) :: hyp
     end function coord_scalar_in
  end interface
\eatline
\nwendcode{}\nwbegindocs{702}\nwdocspar
\nwenddocs{}\nwbegincode{703}\moddef{Muli interactions: interfaces}\plusendmoddef
  abstract interface
     subroutine coord_hcd_in (hyp, cart, denom)
       use kinds !NODEP!
       real(default), dimension(3), intent(in) :: hyp
       real(default), dimension(3), intent(out) :: cart
       real(default), intent(out) :: denom
     end subroutine coord_hcd_in
  end interface
\eatline
\nwendcode{}\nwbegindocs{704}This is the interface to the routines {\tt{}alphasPDF} and
{\tt{}evolvePDF} from LHAPDF or internal PDFs which therefore need to be
explicitly in {\tt{}double} precision.
\nwenddocs{}\nwbegincode{705}\moddef{Muli interactions: interfaces}\plusendmoddef
  interface
     pure function alphaspdf (Q)
       use kinds !NODEP!
       real(double) :: alphaspdf
       real(double), intent(in) :: Q
     end function alphaspdf
  end interface
\eatline
\nwendcode{}\nwbegindocs{706}\nwdocspar
\nwenddocs{}\nwbegincode{707}\moddef{Muli interactions: interfaces}\plusendmoddef
  interface
     pure subroutine evolvepdf (x, q, f)
       use kinds !NODEP!
       real(double), intent(in) :: x, q
       real(double), intent(out), dimension(-6:6) :: f
     end subroutine evolvepdf
  end interface  
  
\eatline
\nwendcode{}\nwbegindocs{708}\nwdocspar
\nwenddocs{}\nwbegincode{709}\moddef{Muli interactions: public}\endmoddef
  public :: muli_get_state_transformations
\nwendcode{}\nwbegincode{710}\moddef{Muli interactions: procedures}\endmoddef
  pure function muli_get_state_transformations &
       (inout_kind, lha_flavors) result (transformations)
    integer, intent(in) :: inout_kind
    integer, dimension(4), intent(in) :: lha_flavors
    integer, dimension(4) :: signature
    logical, dimension(3) :: transformations
    where (lha_flavors > 0)
       signature = 1
    elsewhere (lha_flavors < 0)
       signature = -1
    elsewhere
       signature = 0
    end where
    ! print *,"inout_kind=",inout_kind
    ! print *,"lha_flavors=",lha_flavors
    ! print *,"signature",signature
    if ((sum(inout_signatures(1:2,inout_kind)) == sum(signature(1:2))) .and. &
        (sum(inout_signatures(3:4,inout_kind)) == sum(signature(3:4)))) then
           transformations(1) = .false.
        else
           transformations(1) = .true.
           signature = -signature
        end if
    if (all (inout_signatures(1:2,inout_kind) == signature(1:2))) then
       transformations(2) = .false.
    else
       transformations(2) = .true.
    end if
    if (all(inout_signatures(3:4,inout_kind) == signature(3:4))) then
       transformations(3) = .false.
    else
       transformations(3) = .true.
    end if
    ! print *,"signature",signature
    ! print *,"transformations=",transformations
  end function muli_get_state_transformations
  
\nwendcode{}\nwbegindocs{711}%muli_get_state_transformations
\nwenddocs{}\nwbegindocs{712}\nwdocspar
\nwenddocs{}\nwbegincode{713}\moddef{Muli interactions: public}\plusendmoddef
  public :: h_to_c_param
\nwendcode{}\nwbegincode{714}\moddef{Muli interactions: procedures}\plusendmoddef
  pure function h_to_c_param (hyp)
    real(default), dimension(3) :: h_to_c_param
    real(default), dimension(3), intent(in) :: hyp
    h_to_c_param = [sqrt (sqrt ((((hyp(1)**4) * (one-hyp(3))) + &
         hyp(3))**2 + (((hyp(2)-(5E-1_default))**3)*4)**2) - &
         ((hyp(2)-(5E-1_default))**3)*4), &
         sqrt (sqrt ((((hyp(1)**4)*(one-hyp(3))) + hyp(3))**2 + &
         (((hyp(2)-(5E-1_default))**3)*4)**2) + &
         ((hyp(2)-(5E-1_default))**3)*4), hyp(3)]
  end function h_to_c_param
  
\eatline
\nwendcode{}\nwbegindocs{715}\nwdocspar
\nwenddocs{}\nwbegincode{716}\moddef{Muli interactions: public}\plusendmoddef
  public :: c_to_h_param
\nwendcode{}\nwbegincode{717}\moddef{Muli interactions: procedures}\plusendmoddef
  pure function c_to_h_param (cart)
    real(default), dimension(3) :: c_to_h_param
    real(default), dimension(3), intent(in)::cart
    c_to_h_param= [ (((cart(1)*cart(2)) - cart(3)) / &
         (one - cart(3)))**(1/four), (one + sign(abs((cart(2)**2) - &
         (cart(1)**2))**(1/three), cart(2) - cart(1))) / two, cart(3) ]
  end function c_to_h_param
  
\eatline
\nwendcode{}\nwbegindocs{718}\nwdocspar
\nwenddocs{}\nwbegincode{719}\moddef{Muli interactions: public}\plusendmoddef
  public :: h_to_c_param_def
\nwendcode{}\nwbegincode{720}\moddef{Muli interactions: procedures}\plusendmoddef
  pure function h_to_c_param_def (hyp)
    real(default), dimension(3) :: h_to_c_param_def
    real(default), dimension(3), intent(in) :: hyp
    h_to_c_param_def = h_to_c_param ([hyp(1), hyp(2), pts2_scale])
  end function h_to_c_param_def
    
\eatline
\nwendcode{}\nwbegindocs{721}\nwdocspar
\nwenddocs{}\nwbegincode{722}\moddef{Muli interactions: public}\plusendmoddef
  public :: h_to_c_ort
\nwendcode{}\nwbegincode{723}\moddef{Muli interactions: procedures}\plusendmoddef
  pure function h_to_c_ort (hyp)
    real(default), dimension(3) :: h_to_c_ort
    real(default), dimension(3), intent(in) :: hyp
    h_to_c_ort = [sqrt (sqrt (((hyp(1) * (one - hyp(3))) + hyp(3))**2 + &
         (hyp(2) - (5E-1_default))**2) - (hyp(2) - (5E-1_default))), &
         sqrt (sqrt (((hyp(1) * (one - hyp(3))) + hyp(3))**2 + &
         (hyp(2)-(5E-1_default))**2) + (hyp(2) - (5E-1_default))), hyp(3)]
  end function h_to_c_ort
  
\eatline
\nwendcode{}\nwbegindocs{724}\nwdocspar
\nwenddocs{}\nwbegincode{725}\moddef{Muli interactions: public}\plusendmoddef
  public :: c_to_h_ort
\nwendcode{}\nwbegincode{726}\moddef{Muli interactions: procedures}\plusendmoddef
  pure function c_to_h_ort (cart)
    real(default), dimension(3) :: c_to_h_ort
    real(default), dimension(3), intent(in) :: cart
    c_to_h_ort = [ (cart(3) - (cart(1)*cart(2))) / (cart(3) - one), &
         (one - cart(1)**2 + cart(2)**2) / two, cart(3)]
  end function c_to_h_ort
  
\eatline
\nwendcode{}\nwbegindocs{727}\nwdocspar
\nwenddocs{}\nwbegincode{728}\moddef{Muli interactions: public}\plusendmoddef
  public :: h_to_c_ort_def
\nwendcode{}\nwbegincode{729}\moddef{Muli interactions: procedures}\plusendmoddef
  pure function h_to_c_ort_def (hyp)
    real(default), dimension(3) :: h_to_c_ort_def
    real(default), dimension(3), intent(in) :: hyp
    h_to_c_ort_def = h_to_c_ort ([hyp(1), hyp(2), pts2_scale])
  end function h_to_c_ort_def
  
\eatline
\nwendcode{}\nwbegindocs{730}\nwdocspar
\nwenddocs{}\nwbegincode{731}\moddef{Muli interactions: public}\plusendmoddef
  public :: c_to_h_ort_def
\nwendcode{}\nwbegincode{732}\moddef{Muli interactions: procedures}\plusendmoddef
  pure function c_to_h_ort_def (cart)
    real(default), dimension(3) :: c_to_h_ort_def
    real(default), dimension(3), intent(in) :: cart
    c_to_h_ort_def = c_to_h_ort ([ cart(1), cart(2), pts2_scale])
  end function c_to_h_ort_def
    
\eatline
\nwendcode{}\nwbegindocs{733}\nwdocspar
\nwenddocs{}\nwbegincode{734}\moddef{Muli interactions: public}\plusendmoddef
  public :: h_to_c_noparam
\nwendcode{}\nwbegincode{735}\moddef{Muli interactions: procedures}\plusendmoddef
  pure function h_to_c_noparam (hyp)
    real(default), dimension(2) :: h_to_c_noparam
    real(default), dimension(2), intent(in) :: hyp
    h_to_c_noparam = [sqrt (sqrt (hyp(1)**8 + (((hyp(2) - &
         (5E-1_default))**3)*4)**2) - ((hyp(2)-(5E-1_default))**3)*4), &
         sqrt (sqrt (hyp(1)**8 + (((hyp(2)-(5E-1_default))**3)*4)**2) + &
         ((hyp(2)-(5E-1_default))**3)*4)]
  end function h_to_c_noparam

\eatline
\nwendcode{}\nwbegindocs{736}\nwdocspar
\nwenddocs{}\nwbegincode{737}\moddef{Muli interactions: public}\plusendmoddef
  public :: c_to_h_noparam
\nwendcode{}\nwbegincode{738}\moddef{Muli interactions: procedures}\plusendmoddef
  pure function c_to_h_noparam (cart)
    real(default), dimension(2) :: c_to_h_noparam
    real(default), dimension(2), intent(in) :: cart
    c_to_h_noparam = [sqrt (sqrt (cart(1)*cart(2))), &
         (one + sign(abs((cart(2)**2) - (cart(1)**2))**(one/three), &
         cart(2)-cart(1)))/two]
  end function c_to_h_noparam

\eatline
\nwendcode{}\nwbegindocs{739}\nwdocspar
\nwenddocs{}\nwbegincode{740}\moddef{Muli interactions: public}\plusendmoddef
  public :: c_to_h_param_def
\nwendcode{}\nwbegincode{741}\moddef{Muli interactions: procedures}\plusendmoddef
  pure function c_to_h_param_def (cart)
    real(default), dimension(3) :: c_to_h_param_def
    real(default), dimension(3), intent(in) :: cart
    if (product (cart(1:2)) >= pts2_scale) then
       c_to_h_param_def = c_to_h_param ([cart(1), cart(2), pts2_scale])
    else
       c_to_h_param_def = [-one, -one, -one]
    end if
  end function c_to_h_param_def
  
\eatline
\nwendcode{}\nwbegindocs{742}\nwdocspar
\nwenddocs{}\nwbegincode{743}\moddef{Muli interactions: public}\plusendmoddef
  public :: h_to_c_smooth
\nwendcode{}\nwbegincode{744}\moddef{Muli interactions: procedures}\plusendmoddef
  pure function h_to_c_smooth (hyp)
    real(default), dimension(3) :: h_to_c_smooth
    real(default), dimension(3), intent(in) :: hyp
    real(default) :: h2
    h2 = (((hyp(2) - 5E-1_default)**3) * 4._default + hyp(2)-5E-1_default) &
         / two
    h_to_c_smooth = &
        [sqrt (sqrt((((hyp(1)**4)*(one-hyp(3)))+hyp(3))**2+h2**2) - h2), &
         sqrt (sqrt((((hyp(1)**4)*(one-hyp(3)))+hyp(3))**2+h2**2) + h2), &
         hyp(3)]
  end function h_to_c_smooth
  
\eatline
\nwendcode{}\nwbegindocs{745}\nwdocspar
\nwenddocs{}\nwbegincode{746}\moddef{Muli interactions: public}\plusendmoddef
  public :: c_to_h_smooth
\nwendcode{}\nwbegincode{747}\moddef{Muli interactions: procedures}\plusendmoddef
  pure function c_to_h_smooth (cart)
    real(default), dimension(3) :: c_to_h_smooth
    real(default), dimension(3), intent(in) :: cart
    c_to_h_smooth = &
         [((product (cart(1:2)) - cart(3)) / (one - cart(3)))**(1/four), &
          (three-three**(two/3) / (-9._default * cart(1)**2 + &
           9._default * cart(2)**2 + sqrt (three + 81._default * &
           (cart(1)**2 - cart(2)**2)**2))**(one/three)&
         + 3**(one/3)*(-9._default * cart(1)**2 + 9._default*cart(2)**2 &
         + sqrt(three + 81._default*(cart(1)**2&
         - cart(2)**2)**2))**(one/3))/6._default,cart(3)]
  end function c_to_h_smooth

\eatline
\nwendcode{}\nwbegindocs{748}\nwdocspar
\nwenddocs{}\nwbegincode{749}\moddef{Muli interactions: public}\plusendmoddef
  public :: h_to_c_smooth_def
\nwendcode{}\nwbegincode{750}\moddef{Muli interactions: procedures}\plusendmoddef
  pure function h_to_c_smooth_def (hyp)
    real(default), dimension(3) :: h_to_c_smooth_def
    real(default), dimension(3), intent(in) :: hyp
    h_to_c_smooth_def = h_to_c_smooth ([hyp(1), hyp(2), pts2_scale])
  end function h_to_c_smooth_def
  
\eatline
\nwendcode{}\nwbegindocs{751}\nwdocspar
\nwenddocs{}\nwbegincode{752}\moddef{Muli interactions: public}\plusendmoddef
\nwendcode{}\nwbegindocs{753}\nwdocspar
\nwenddocs{}\nwbegincode{754}\moddef{Muli interactions: procedures}\plusendmoddef
  pure function c_to_h_smooth_def (cart)
    real(default), dimension(3)::c_to_h_smooth_def
    real(default), dimension(3), intent(in) :: cart
    if (product (cart(1:2)) >= pts2_scale) then
       c_to_h_smooth_def = c_to_h_smooth ([cart(1), cart(2), pts2_scale])
    else
       c_to_h_smooth_def = [-one, -one, -one]
    end if
  end function c_to_h_smooth_def

\eatline
\nwendcode{}\nwbegindocs{755}\nwdocspar
\nwenddocs{}\nwbegincode{756}\moddef{Muli interactions: public}\plusendmoddef
  public :: voxel_h_to_c_ort
\nwendcode{}\nwbegincode{757}\moddef{Muli interactions: procedures}\plusendmoddef
  pure function voxel_h_to_c_ort (hyp)
    real(default) :: voxel_h_to_c_ort
    real(default), dimension(3), intent(in) :: hyp
    real(default) :: T, TH1
    T = one - hyp(3)
    TH1 = T * (one - hyp(1))
    voxel_h_to_c_ort = sqrt (T**2 / (five - four*(one-hyp(2))*hyp(2) - &    
         four*(two-TH1)*TH1))
  end function voxel_h_to_c_ort

\eatline
\nwendcode{}\nwbegindocs{758}\nwdocspar
\nwenddocs{}\nwbegincode{759}\moddef{Muli interactions: public}\plusendmoddef
  public :: voxel_c_to_h_ort
\nwendcode{}\nwbegincode{760}\moddef{Muli interactions: procedures}\plusendmoddef
  pure function voxel_c_to_h_ort(cart)
    real(default) :: voxel_c_to_h_ort
    real(default), dimension(3), intent(in) :: cart
    real(default) :: P
    P = product (cart(1:2))
    if (P > cart(3)) then
       voxel_c_to_h_ort = (cart(1)**2 + cart(2)**2) / (one -cart(3))
    else
       voxel_c_to_h_ort = zero
    end if
  end function voxel_c_to_h_ort

\eatline
\nwendcode{}\nwbegindocs{761}\nwdocspar
\nwenddocs{}\nwbegincode{762}\moddef{Muli interactions: public}\plusendmoddef
  public :: voxel_h_to_c_noparam
\nwendcode{}\nwbegincode{763}\moddef{Muli interactions: procedures}\plusendmoddef
  pure function voxel_h_to_c_noparam (hyp)
    real(default) :: voxel_h_to_c_noparam
    real(default), dimension(3), intent(in) :: hyp
    voxel_h_to_c_noparam = 12._default * sqrt ((hyp(1)**6 * &
         (one - two*hyp(2))**4) / (4*hyp(1)**8 + (one - two*hyp(2))**6))
  end function voxel_h_to_c_noparam

\eatline
\nwendcode{}\nwbegindocs{764}\nwdocspar
\nwenddocs{}\nwbegincode{765}\moddef{Muli interactions: public}\plusendmoddef
  public :: voxel_c_to_h_noparam
\nwendcode{}\nwbegincode{766}\moddef{Muli interactions: procedures}\plusendmoddef
  pure function voxel_c_to_h_noparam (cart)
    real(default) :: voxel_c_to_h_noparam
    real(default), dimension(3), intent(in) :: cart
    real(default) :: P
    voxel_c_to_h_noparam = (cart(1)**2 + cart(2)**2) / (12._default * &
         (cart(1)*cart(2))**(three/four) * &
         (cart(2)**2 + cart(1)**2)**(two/three))
  end function voxel_c_to_h_noparam

\eatline
\nwendcode{}\nwbegindocs{767}\nwdocspar
\nwenddocs{}\nwbegincode{768}\moddef{Muli interactions: public}\plusendmoddef
  public :: voxel_h_to_c_param
\nwendcode{}\nwbegincode{769}\moddef{Muli interactions: procedures}\plusendmoddef
  pure function voxel_h_to_c_param (hyp)
    real(default) :: voxel_h_to_c_param
    real(default), dimension(3), intent(in) :: hyp
    voxel_h_to_c_param = 12*Sqrt((hyp(1)**6 * &
         (one - 2._default*hyp(2))**4 * (hyp(3) - one)**2) / &
         ((one - two * hyp(2))**6 + four * &
         (hyp(3)-(hyp(1)**4*(hyp(3)-one)))**2))
  end function voxel_h_to_c_param

\eatline
\nwendcode{}\nwbegindocs{770}\nwdocspar
\nwenddocs{}\nwbegincode{771}\moddef{Muli interactions: public}\plusendmoddef
  public :: voxel_c_to_h_param
\nwendcode{}\nwbegincode{772}\moddef{Muli interactions: procedures}\plusendmoddef
  pure function voxel_c_to_h_param (cart)
    real(default)::voxel_c_to_h_param
    real(default), dimension(3), intent(in) :: cart
    real(default) :: P, T, CP, CM
    P = product (cart(1:2))
    if (P > cart(3)) then
       P = P - cart(3)
       CP = cart(1)**2 + cart(2)**2
       CM = abs(cart(2)**2 - cart(1)**2)
       T = 1 - cart(3)
       voxel_c_to_h_param = (Cp*sqrt(sqrt(P/T))) / (12*Cm**(two/three)*P)
    else
       voxel_c_to_h_param = zero
    end if
  end function voxel_c_to_h_param

\eatline
\nwendcode{}\nwbegindocs{773}\nwdocspar
\nwenddocs{}\nwbegincode{774}\moddef{Muli interactions: public}\plusendmoddef
  public :: voxel_h_to_c_smooth 
\nwendcode{}\nwbegindocs{775}\nwdocspar
\nwenddocs{}\nwbegincode{776}\moddef{Muli interactions: procedures}\plusendmoddef
  pure function voxel_h_to_c_smooth (hyp)
    real(default) :: voxel_h_to_c_smooth
    real(default), dimension(3), intent(in) :: hyp
    real(default) :: T
    T = one - hyp(3)
    voxel_h_to_c_smooth = 8._default * (hyp(1)**3 * (one + three * &
         (hyp(2) - one)*hyp(2))*T) / sqrt ((one - two*hyp(2) * (two + &
         hyp(2)*(two*hyp(2)-three)))**2 + &
         four * (one + (hyp(1)**4 - one)*T)**2) 
  end function voxel_h_to_c_smooth

\eatline
\nwendcode{}\nwbegindocs{777}\nwdocspar
\nwenddocs{}\nwbegincode{778}\moddef{Muli interactions: public}\plusendmoddef
  public :: voxel_c_to_h_smooth
\nwendcode{}\nwbegincode{779}\moddef{Muli interactions: procedures}\plusendmoddef
  pure function voxel_c_to_h_smooth (cart)
    real(default) :: voxel_c_to_h_smooth
    real(default), dimension(3), intent(in) :: cart
    real(default) :: P, S, T, CM, CP
    P = product (cart(1:2))
    if (P > cart(3)) then
       P = P - cart(3)
       CP = cart(1)**2 + cart(2)**2
       CM = cart(2)**2 - cart(1)**2
       T = 1 - cart(3)
       S = sqrt(three + 81._default*cm**2)
       voxel_c_to_h_smooth = (three**(one/three) * Cp*(three**(one/three) + &
            (9._default*Cm + S)**(two/three)) * sqrt (sqrt (P/T))) / &
            (four * P * S * (9._default * Cm + S)**(one/three))
    else
       voxel_c_to_h_smooth = zero
    end if
end function voxel_c_to_h_smooth

\eatline
\nwendcode{}\nwbegindocs{780}\nwdocspar
\nwenddocs{}\nwbegincode{781}\moddef{Muli interactions: public}\plusendmoddef
  public :: voxel_h_to_c_ort_def
\nwendcode{}\nwbegincode{782}\moddef{Muli interactions: procedures}\plusendmoddef
  pure function voxel_h_to_c_ort_def (hyp)
    real(default) :: voxel_h_to_c_ort_def
    real(default), dimension(3), intent(in) :: hyp
    voxel_h_to_c_ort_def = voxel_h_to_c_ort (hyp)
  end function voxel_h_to_c_ort_def

\eatline
\nwendcode{}\nwbegindocs{783}\nwdocspar
\nwenddocs{}\nwbegincode{784}\moddef{Muli interactions: public}\plusendmoddef
  public :: voxel_c_to_h_ort_def
\nwendcode{}\nwbegincode{785}\moddef{Muli interactions: procedures}\plusendmoddef
  pure function voxel_c_to_h_ort_def (cart)
    real(default) :: voxel_c_to_h_ort_def
    real(default), dimension(3), intent(in) :: cart
    voxel_c_to_h_ort_def = voxel_c_to_h_ort (cart)
  end function voxel_c_to_h_ort_def

\eatline
\nwendcode{}\nwbegindocs{786}\nwdocspar
\nwenddocs{}\nwbegincode{787}\moddef{Muli interactions: public}\plusendmoddef
  public :: voxel_h_to_c_param_def 
\nwendcode{}\nwbegincode{788}\moddef{Muli interactions: procedures}\plusendmoddef
  pure function voxel_h_to_c_param_def (hyp)
    real(default) :: voxel_h_to_c_param_def
    real(default), dimension(3), intent(in) :: hyp
    voxel_h_to_c_param_def = voxel_h_to_c_param (hyp)
  end function voxel_h_to_c_param_def

\eatline
\nwendcode{}\nwbegindocs{789}\nwdocspar
\nwenddocs{}\nwbegincode{790}\moddef{Muli interactions: public}\plusendmoddef
  public :: voxel_c_to_h_param_def
\nwendcode{}\nwbegincode{791}\moddef{Muli interactions: procedures}\plusendmoddef
  pure function voxel_c_to_h_param_def (cart)
    real(default) :: voxel_c_to_h_param_def
    real(default), dimension(3), intent(in) :: cart
    voxel_c_to_h_param_def = voxel_c_to_h_param (cart)
  end function voxel_c_to_h_param_def

\eatline
\nwendcode{}\nwbegindocs{792}\nwdocspar
\nwenddocs{}\nwbegincode{793}\moddef{Muli interactions: public}\plusendmoddef
  public :: voxel_h_to_c_smooth_def
\nwendcode{}\nwbegincode{794}\moddef{Muli interactions: procedures}\plusendmoddef
  pure function voxel_h_to_c_smooth_def (hyp)
    real(default) :: voxel_h_to_c_smooth_def
    real(default), dimension(3), intent(in) :: hyp
    voxel_h_to_c_smooth_def = voxel_h_to_c_smooth (hyp)
  end function voxel_h_to_c_smooth_def

\eatline
\nwendcode{}\nwbegindocs{795}\nwdocspar
\nwenddocs{}\nwbegincode{796}\moddef{Muli interactions: public}\plusendmoddef
  public :: voxel_c_to_h_smooth_def
\nwendcode{}\nwbegincode{797}\moddef{Muli interactions: procedures}\plusendmoddef
  pure function voxel_c_to_h_smooth_def (cart)
    real(default) :: voxel_c_to_h_smooth_def
    real(default), dimension(3), intent(in) :: cart
    voxel_c_to_h_smooth_def = voxel_c_to_h_smooth (cart)
  end function voxel_c_to_h_smooth_def

\eatline
\nwendcode{}\nwbegindocs{798}\nwdocspar
\nwenddocs{}\nwbegincode{799}\moddef{Muli interactions: public}\plusendmoddef
  public :: denom_cart
\nwendcode{}\nwbegincode{800}\moddef{Muli interactions: procedures}\plusendmoddef
  pure function denom_cart (cart)
    real(default) :: denom_cart
    real(default), dimension(3), intent(in) :: cart
    denom_cart = 1._default / (864._default * sqrt (cart(3)**3 * &
         (1._default - cart(3) / product(cart(1:2)))))
  end function denom_cart

\eatline
\nwendcode{}\nwbegindocs{801}\nwdocspar
\nwenddocs{}\nwbegincode{802}\moddef{Muli interactions: public}\plusendmoddef
  public :: denom_ort
\nwendcode{}\nwbegincode{803}\moddef{Muli interactions: procedures}\plusendmoddef
  pure function denom_ort (hyp)
    real(default) :: denom_ort
    real(default), dimension(3), intent(in) :: hyp
    real(default) :: Y, P
    Y = (one - two * hyp(2))**2
    P = one - hyp(3)
    if (hyp(1) > zero .and. hyp(3) > zero) then
       denom_ort = sqrt ((P + (-1 + Hyp(1))*P**2) / &
            (746496*hyp(1)*hyp(3)**3 * (4*(1 + (-1 + hyp(1))*P)**2 + Y)))
    else
       denom_ort = zero
    end if
  end function denom_ort
  
\eatline
\nwendcode{}\nwbegindocs{804}\nwdocspar
\nwenddocs{}\nwbegincode{805}\moddef{Muli interactions: public}\plusendmoddef
  public :: denom_param
\nwendcode{}\nwbegincode{806}\moddef{Muli interactions: procedures}\plusendmoddef
  pure function denom_param (hyp)
    real(default) :: denom_param
    real(default), dimension(3), intent(in) :: hyp
    real(default) :: X, Y, P
    X = hyp(1)**4
    Y = 1._default - 2._default * hyp(2)
    P = 1._default - hyp(3)
    if (hyp(3) > 0._default) then
       denom_param = sqrt ((P * (1+P*(X-1)) * Sqrt(X)*Y**4) / &
            (5184*(4*(1+P*(X-1))**2+Y**6)*hyp(3)**3))
    else
       denom_param = zero
    end if
  end function denom_param

\eatline
\nwendcode{}\nwbegindocs{807}\nwdocspar
\nwenddocs{}\nwbegincode{808}\moddef{Muli interactions: public}\plusendmoddef
  public :: denom_param_reg
\nwendcode{}\nwbegincode{809}\moddef{Muli interactions: procedures}\plusendmoddef
  pure function denom_param_reg (hyp)
    real(default) :: denom_param_reg
    real(default), dimension(3), intent(in) :: hyp
    real(default) :: X, Y, P
    X = hyp(1)**4
    Y = one - two * hyp(2)
    P = one - hyp(3)
    if (hyp(3) > zero) then
       denom_param_reg = sqrt ((P*(1+P*(X-1)) * Sqrt(X)*Y**4) / &
            (5184*(4*(1+P*(X-1))**2+Y**6) * (hyp(3) + norm2_p_t_0)**3))
    else
       denom_param_reg = zero
    end if
  end function denom_param_reg

\eatline
\nwendcode{}\nwbegindocs{810}\nwdocspar
\nwenddocs{}\nwbegincode{811}\moddef{Muli interactions: public}\plusendmoddef
  public :: denom_smooth
\nwendcode{}\nwbegincode{812}\moddef{Muli interactions: procedures}\plusendmoddef
  pure function denom_smooth (hyp)
    real(default) :: denom_smooth
    real(default), dimension(3), intent(in) :: hyp
    real(default) :: X, Y, P
    X = hyp(1)**2
    Y = (one - two * hyp(2))**2
    P = one - hyp(3)
    if (hyp(3) > zero) then
       denom_smooth = sqrt ((P * X * (one + P*(-one + X**2)) * &
            (1 + three*Y)**2)/(46656*hyp(3)**3 &
            *(16*(1 + P*(-1 + X**2))**2 + Y + 2*Y**2 + Y**3)))
    else
       denom_smooth = zero
    end if
  end function denom_smooth

\eatline
\nwendcode{}\nwbegindocs{813}\nwdocspar
\nwenddocs{}\nwbegincode{814}\moddef{Muli interactions: public}\plusendmoddef
  public :: denom_smooth_reg
\nwendcode{}\nwbegincode{815}\moddef{Muli interactions: procedures}\plusendmoddef
  pure function denom_smooth_reg (hyp)
    real(default) :: denom_smooth_reg
    real(default), dimension(3), intent(in) :: hyp
    real(default) :: X, Y, P
    X = hyp(1)**2
    Y = (one - two * hyp(2))**2
    P = one - hyp(3)
    if (hyp(3) > zero) then
       denom_smooth_reg = sqrt ((P * X * (1 + P*(-1 + X**2)) * &
            (1 + 3*Y)**2)/(46656*(hyp(3) + norm2_p_t_0)**3 * &
            (16 * (1 + P*(-1 + X**2))**2 + Y + 2*Y**2 + Y**3)))
    else
       denom_smooth_reg = zero
    end if
  end function denom_smooth_reg
  
\eatline
\nwendcode{}\nwbegindocs{816}\nwdocspar
\nwenddocs{}\nwbegincode{817}\moddef{Muli interactions: public}\plusendmoddef
  public :: denom_cart_save
\nwendcode{}\nwbegincode{818}\moddef{Muli interactions: procedures}\plusendmoddef
   pure function denom_cart_save (cart)
    real(default) :: denom_cart_save
    real(default), dimension(3), intent(in) :: cart
    if (product(cart(1:2)) > cart(3)) then
       denom_cart_save = denom_cart (cart)
    else
       denom_cart_save = zero
    end if
  end function denom_cart_save

\eatline
\nwendcode{}\nwbegindocs{819}\nwdocspar
\nwenddocs{}\nwbegincode{820}\moddef{Muli interactions: public}\plusendmoddef
  public :: denom_ort_save
\nwendcode{}\nwbegincode{821}\moddef{Muli interactions: procedures}\plusendmoddef
  pure function denom_ort_save (hyp)
    real(default) :: denom_ort_save
    real(default), dimension(3), intent(in) :: hyp
    real(default) :: Y, Z, W
    real(default), dimension(3) :: cart
    cart = h_to_c_ort (hyp)
    if (cart(1) > one .or. cart(2) > one) then
       denom_ort_save = zero
    else
       denom_ort_save = denom_ort (hyp)
    end if
  end function denom_ort_save

\eatline
\nwendcode{}\nwbegindocs{822}\nwdocspar
\nwenddocs{}\nwbegincode{823}\moddef{Muli interactions: public}\plusendmoddef
  public :: denom_param_save
\nwendcode{}\nwbegincode{824}\moddef{Muli interactions: procedures}\plusendmoddef
  pure function denom_param_save (hyp)
    real(default) :: denom_param_save
    real(default), dimension(3), intent(in) :: hyp
    real(default) :: Y, Z, W
    real(default), dimension(3) :: cart
    cart=h_to_c_param (hyp)
    if (cart(1) > one .or. cart(2) > one) then
       denom_param_save = zero
    else
       denom_param_save = denom_param (hyp)
    end if
  end function denom_param_save

\eatline
\nwendcode{}\nwbegindocs{825}\nwdocspar
\nwenddocs{}\nwbegincode{826}\moddef{Muli interactions: public}\plusendmoddef
  public :: denom_smooth_save
\nwendcode{}\nwbegincode{827}\moddef{Muli interactions: procedures}\plusendmoddef
  pure function denom_smooth_save (hyp)
    real(default) :: denom_smooth_save
    real(default), dimension(3), intent(in) :: hyp
    real(default) :: Y, Z, W
    real(default), dimension(3) :: cart
    cart = h_to_c_smooth (hyp)
    if (cart(1) > one .or. cart(2) > one) then
       denom_smooth_save = zero
    else
       denom_smooth_save = denom_smooth (hyp)
    end if
  end function denom_smooth_save

\eatline
\nwendcode{}\nwbegindocs{828}\nwdocspar
\nwenddocs{}\nwbegincode{829}\moddef{Muli interactions: public}\plusendmoddef
  public :: denom_cart_cuba_int
\nwendcode{}\nwbegincode{830}\moddef{Muli interactions: procedures}\plusendmoddef
  subroutine denom_cart_cuba_int (d_cart, cart, d_denom, denom, pt2s)
    real(default), dimension(3), intent(in) :: cart
    real(default), dimension(1), intent(out) :: denom
    real(default), intent(in) :: pt2s
    integer, intent(in) :: d_cart, d_denom
    denom(1) = denom_cart_save ([cart(1), cart(2), pt2s])
  end subroutine denom_cart_cuba_int

\eatline
\nwendcode{}\nwbegindocs{831}\nwdocspar
\nwenddocs{}\nwbegincode{832}\moddef{Muli interactions: public}\plusendmoddef
  public :: denom_ort_cuba_int
\nwendcode{}\nwbegincode{833}\moddef{Muli interactions: procedures}\plusendmoddef
  subroutine denom_ort_cuba_int (d_hyp, hyp, d_denom, denom, pt2s)
    real(default), dimension(3), intent(in) :: hyp
    real(default), dimension(1), intent(out) :: denom
    real(default), intent(in) :: pt2s
    integer, intent(in) :: d_hyp, d_denom
    denom(1) = denom_ort_save ([hyp(1), hyp(2), pt2s])
  end subroutine denom_ort_cuba_int

\eatline
\nwendcode{}\nwbegindocs{834}\nwdocspar
\nwenddocs{}\nwbegincode{835}\moddef{Muli interactions: public}\plusendmoddef
  public :: denom_param_cuba_int
\nwendcode{}\nwbegincode{836}\moddef{Muli interactions: procedures}\plusendmoddef
  subroutine denom_param_cuba_int (d_hyp, hyp, d_denom, denom, pt2s)
    real(default), dimension(3), intent(in) :: hyp
    real(default), dimension(1), intent(out) :: denom
    real(default), intent(in) :: pt2s
    integer, intent(in) :: d_hyp, d_denom
    denom(1) = denom_param_save ([hyp(1), hyp(2), pt2s])
  end subroutine denom_param_cuba_int

\eatline
\nwendcode{}\nwbegindocs{837}\nwdocspar
\nwenddocs{}\nwbegincode{838}\moddef{Muli interactions: public}\plusendmoddef
  public :: denom_smooth_cuba_int
\nwendcode{}\nwbegincode{839}\moddef{Muli interactions: procedures}\plusendmoddef
  subroutine denom_smooth_cuba_int (d_hyp, hyp, d_denom, denom, pt2s)
    real(default), dimension(3), intent(in) :: hyp
    real(default), dimension(1), intent(out) :: denom
    real(default), intent(in) :: pt2s
    integer, intent(in) :: d_hyp, d_denom
    denom(1) = denom_smooth_save ([hyp(1), hyp(2), pt2s])
  end subroutine denom_smooth_cuba_int

\eatline
\nwendcode{}\nwbegindocs{840}\nwdocspar
\nwenddocs{}\nwbegincode{841}\moddef{Muli interactions: public}\plusendmoddef
  public :: coordinates_hcd_cart
\nwendcode{}\nwbegincode{842}\moddef{Muli interactions: procedures}\plusendmoddef
  subroutine coordinates_hcd_cart (hyp, cart, denom)
    real(default), dimension(3), intent(in) :: hyp
    real(default), dimension(3), intent(out) :: cart
    real(default), intent(out) :: denom
    cart = hyp
    denom = denom_cart_save (cart)
  end subroutine coordinates_hcd_cart

\eatline
\nwendcode{}\nwbegindocs{843}\nwdocspar
\nwenddocs{}\nwbegincode{844}\moddef{Muli interactions: public}\plusendmoddef
  public :: coordinates_hcd_ort
\nwendcode{}\nwbegincode{845}\moddef{Muli interactions: procedures}\plusendmoddef
  subroutine coordinates_hcd_ort (hyp, cart, denom)
    real(default), dimension(3), intent(in) :: hyp
    real(default), dimension(3), intent(out) :: cart
    real(default), intent(out)::denom
    cart = h_to_c_ort (hyp)
    denom = denom_ort (hyp)
  end subroutine coordinates_hcd_ort

\eatline
\nwendcode{}\nwbegindocs{846}\nwdocspar
\nwenddocs{}\nwbegincode{847}\moddef{Muli interactions: public}\plusendmoddef
  public :: coordinates_hcd_param
\nwendcode{}\nwbegincode{848}\moddef{Muli interactions: procedures}\plusendmoddef
  subroutine coordinates_hcd_param (hyp, cart, denom)
    real(default), dimension(3), intent(in) :: hyp
    real(default), dimension(3), intent(out) :: cart
    real(default), intent(out) :: denom
    cart = h_to_c_param (hyp)
    denom = denom_param (hyp)
  end subroutine coordinates_hcd_param

\eatline
\nwendcode{}\nwbegindocs{849}\nwdocspar
\nwenddocs{}\nwbegincode{850}\moddef{Muli interactions: public}\plusendmoddef
  public :: coordinates_hcd_param_reg
\nwendcode{}\nwbegincode{851}\moddef{Muli interactions: procedures}\plusendmoddef
  subroutine coordinates_hcd_param_reg (hyp, cart, denom)
    real(default), dimension(3), intent(in) :: hyp
    real(default), dimension(3), intent(out) :: cart
    real(default), intent(out) :: denom
    cart = h_to_c_param (hyp)
    denom = denom_param_reg (hyp)
  end subroutine coordinates_hcd_param_reg

\eatline
\nwendcode{}\nwbegindocs{852}\nwdocspar
\nwenddocs{}\nwbegincode{853}\moddef{Muli interactions: public}\plusendmoddef
  public :: coordinates_hcd_smooth
\nwendcode{}\nwbegincode{854}\moddef{Muli interactions: procedures}\plusendmoddef
  subroutine coordinates_hcd_smooth (hyp, cart, denom)
    real(default), dimension(3), intent(in) :: hyp
    real(default), dimension(3), intent(out) :: cart
    real(default), intent(out) :: denom
    cart = h_to_c_smooth (hyp)
    denom = denom_smooth (hyp)
  end subroutine coordinates_hcd_smooth
  
\eatline
\nwendcode{}\nwbegindocs{855}\nwdocspar
\nwenddocs{}\nwbegincode{856}\moddef{Muli interactions: public}\plusendmoddef
  public :: coordinates_hcd_smooth_reg
\nwendcode{}\nwbegincode{857}\moddef{Muli interactions: procedures}\plusendmoddef
  subroutine coordinates_hcd_smooth_reg (hyp, cart, denom)
    real(default), dimension(3), intent(in) :: hyp
    real(default), dimension(3), intent(out) :: cart
    real(default), intent(out) :: denom
    cart = h_to_c_smooth (hyp)
    denom = denom_smooth_reg (hyp)
  end subroutine coordinates_hcd_smooth_reg

\eatline
\nwendcode{}\nwbegindocs{858}\nwdocspar
\nwenddocs{}\nwbegincode{859}\moddef{Muli interactions: public}\plusendmoddef
  public :: interactions_dddsigma_reg
\nwendcode{}\nwbegincode{860}\moddef{Muli interactions: procedures}\plusendmoddef
  pure subroutine interactions_dddsigma_reg &
       (process_id, double_pdf_id, hyp, cart, dddsigma)
    real(default), intent(out) :: dddsigma
    integer, intent(in) :: process_id, double_pdf_id
    real(default), dimension(3), intent(in) :: hyp
    real(default), dimension(3), intent(out) :: cart
    real(default) :: a, pt2shat, gev_pt, gev2_pt
    cart = h_to_c_param (hyp)
    a = product (cart(1:2))
    if (cart(1) <= 1D0 .and. cart(2) <= 1D0) then
       pt2shat = hyp(3) / a
       gev_pt = sqrt(hyp(3)) * gev_pt_max
       gev2_pt = hyp(3) * gev2_pt_max
       ! print *,process_id,pt2shat
       dddsigma = &
            const_pref &
            * alphasPDF (dble (sqrt (gev2_pt+gev2_p_t_0)))**2 &
            * ps_io_pol (process_id, pt2shat) &
            * pdf_in_in_kind &
               (process_id, double_pdf_id, cart(1), cart(2), gev_pt) &
            * denom_param_reg (hyp) / a
    else
       dddsigma = zero
    end if
  end subroutine interactions_dddsigma_reg
  
\eatline
\nwendcode{}\nwbegindocs{861}\nwdocspar
\nwenddocs{}\nwbegincode{862}\moddef{Muli interactions: public}\plusendmoddef
  public :: pdf_in_in_kind
\nwendcode{}\nwbegincode{863}\moddef{Muli interactions: procedures}\plusendmoddef
  pure function pdf_in_in_kind (process_id, double_pdf_id, c1, c2, gev_pt)
    real(default) :: pdf_in_in_kind
    real(default), intent(in) :: c1, c2, gev_pt
    integer, intent(in) :: process_id, double_pdf_id
    real(default) :: pdf1, pdf2
    call single_pdf (valid_processes(1, process_id), &
         double_pdf_kinds(1, double_pdf_id), c1, gev_pt, pdf1)
    call single_pdf (valid_processes(2, process_id), &
         double_pdf_kinds(2, double_pdf_id), c2, gev_pt, pdf2)
    pdf_in_in_kind = pdf1 * pdf2
  contains
    pure subroutine single_pdf (flavor, pdf_kind, c, gev_pt, pdf)
      integer, intent(in) :: flavor, pdf_kind
      real(default), intent(in) :: c, gev_pt
      real(default), intent(out) :: pdf
      real(double), dimension(-6:6) :: lha_pdf
      call evolvePDF (dble (c), dble (gev_pt), lha_pdf)
      select case (pdf_kind)
      case (1)
         pdf = lha_pdf (0)
      case (2)
         if (flavor==1 .or. flavor==2) then
            pdf = lha_pdf (-flavor)
         else
            pdf = lha_pdf (flavor)
         end if
      case (3)
         pdf = lha_pdf(1) - lha_pdf(-1)
      case (4)
         pdf = lha_pdf(2) - lha_pdf(-2)
      end select
    end subroutine single_pdf
  end function pdf_in_in_kind
  
\eatline
\nwendcode{}\nwbegindocs{864}\nwdocspar
\nwenddocs{}\nwbegincode{865}\moddef{Muli interactions: public}\plusendmoddef
  public :: ps_io_pol
\nwendcode{}\nwbegincode{866}\moddef{Muli interactions: procedures}\plusendmoddef
  elemental function ps_io_pol (process_io_id, pt2shat)
    real(default) :: ps_io_pol
    integer, intent(in) :: process_io_id
    real(default), intent(in) :: pt2shat
    ps_io_pol = dot_product([1._default, pt2shat, pt2shat**2, pt2shat**3], &
         phase_space_coefficients_inout (1:4, &
         valid_processes (6, process_io_id)))
  end function ps_io_pol

\eatline
\nwendcode{}\nwbegindocs{867}\nwdocspar
\nwenddocs{}\nwbegincode{868}\moddef{Muli interactions: public}\plusendmoddef
  public :: interactions_dddsigma
\nwendcode{}\nwbegincode{869}\moddef{Muli interactions: procedures}\plusendmoddef
  pure subroutine interactions_dddsigma &
       (process_id, double_pdf_id, hyp, cart, dddsigma)
    real(default), intent(out) :: dddsigma
    integer, intent(in) :: process_id, double_pdf_id
    real(default), dimension(3), intent(in) :: hyp
    real(default), dimension(3), intent(out) :: cart
    real(default) :: a, pt2shat, gev_pt
    cart = h_to_c_param (hyp)
    a = product (cart(1:2))
    if (cart(1) <= 1._default .and. cart(2) <= 1._default) then
       pt2shat = hyp(3) / a
       gev_pt = sqrt(hyp(3)) * gev_pt_max
       ! print *,process_id,pt2shat
       dddsigma = const_pref * &
            alphasPDF (dble (gev_pt))**2 * &
            ps_io_pol (process_id, pt2shat) * &
            pdf_in_in_kind &
               (process_id, double_pdf_id, cart(1), cart(2), gev_pt) * &
            denom_param(hyp) / a
    else
       dddsigma = zero
    end if
  end subroutine interactions_dddsigma

\eatline
\nwendcode{}\nwbegindocs{870}\nwdocspar
\nwenddocs{}\nwbegincode{871}\moddef{Muli interactions: public}\plusendmoddef
  public :: interactions_dddsigma_print
\nwendcode{}\nwbegincode{872}\moddef{Muli interactions: procedures}\plusendmoddef
  subroutine interactions_dddsigma_print &
       (process_id, double_pdf_id, hyp, cart, dddsigma)
    real(default), intent(out) :: dddsigma
    integer, intent(in) :: process_id, double_pdf_id
    real(default), dimension(3), intent(in) :: hyp
    real(default), dimension(3), intent(out) :: cart
    real(default) :: a, pt2shat, gev_pt
    cart = h_to_c_param (hyp)
    a = product (cart(1:2))
    if (cart(1) <= 1._default .and. cart(2) <= 1._default) then
       pt2shat = hyp(3) / a
       gev_pt=sqrt(hyp(3))*gev_pt_max
       ! print *,process_id,pt2shat
       dddsigma = const_pref * &
            ! alphasPDF(dble (gev_pt))**2 * &
            ps_io_pol (process_id, pt2shat) * &
            pdf_in_in_kind &
               (process_id, double_pdf_id, cart(1), cart(2), gev_pt) * &
            denom_param (hyp) / a
    else
       dddsigma = zero
    end if
    write(11, *) dddsigma, pt2shat, &
         pdf_in_in_kind (process_id, double_pdf_id, cart(1), cart(2), &
         gev_pt), ps_io_pol (process_id, pt2shat), const_pref, &
         denom_param(hyp), a
    flush(11)
  end subroutine interactions_dddsigma_print

\eatline
\nwendcode{}\nwbegindocs{873}\nwdocspar
\nwenddocs{}\nwbegincode{874}\moddef{Muli interactions: public}\plusendmoddef
  public :: interactions_dddsigma_cart 
\nwendcode{}\nwbegincode{875}\moddef{Muli interactions: procedures}\plusendmoddef
  pure subroutine interactions_dddsigma_cart &
       (process_id, double_pdf_id, cart, dddsigma)
    real(default), intent(out) :: dddsigma
    integer, intent(in) :: process_id, double_pdf_id
    real(default), dimension(3), intent(in) :: cart
    real(default) :: a, pt2shat, gev_pt
    a = product (cart(1:2))
    if (cart(1) <= one .and. cart(2) <= one) then
       pt2shat = cart(3) / a
       gev_pt = sqrt(cart(3)) * gev_pt_max
       ! print *,process_id,pt2shat
       dddsigma = const_pref * &
            alphasPDF (dble (gev_pt))**2 * &
            ps_io_pol (process_id, pt2shat) * &
            pdf_in_in_kind &
               (process_id, double_pdf_id, cart(1), cart(2), gev_pt) * &
            denom_cart (cart) / a
    else
       dddsigma = zero
    end if
  end subroutine interactions_dddsigma_cart

\eatline
\nwendcode{}\nwbegindocs{876}\nwdocspar
\nwenddocs{}\nwbegincode{877}\moddef{Muli interactions: public}\plusendmoddef
  public :: cuba_gg_me_smooth
\nwendcode{}\nwbegincode{878}\moddef{Muli interactions: procedures}\plusendmoddef
  subroutine cuba_gg_me_smooth (d_hyp, hyp, d_me, me, pt2s)
    integer, intent(in) :: d_hyp, d_me
    real(default), dimension(d_hyp), intent(in) :: hyp
    real(default), dimension(1), intent(out) :: me
    real(default), dimension(3) :: cart
    real(default), intent(in) :: pt2s
    real(default) :: p, p2
    if (d_hyp == 3) then
       p = hyp(3)
       p2 = hyp(3)**2
    else
       if (d_hyp == 2) then
          p = sqrt (pt2s)
          p2 = pt2s
       end if
    end if
    cart = h_to_c_smooth ([hyp(1), hyp(2), p2])
    if (p > pts_min .and. product (cart(1:2)) > p2) then
       me(1) = const_pref * &
            alphasPDF (dble (p*gev_pt_max))**2 * &
            ps_io_pol (109, p2) * &
            pdf_in_in_kind (109, 11, cart(1), cart(2), p2) * &
            denom_smooth ([hyp(1), hyp(2), p2]) / product (cart(1:2))
    else
       me(1) = zero
    end if
  end subroutine cuba_gg_me_smooth

\eatline
\nwendcode{}\nwbegindocs{879}\nwdocspar
\nwenddocs{}\nwbegincode{880}\moddef{Muli interactions: public}\plusendmoddef
  public :: cuba_gg_me_param
\nwendcode{}\nwbegincode{881}\moddef{Muli interactions: procedures}\plusendmoddef
  subroutine cuba_gg_me_param (d_hyp, hyp, d_me, me, pt2s)
    integer, intent(in)::d_hyp,d_me
    real(default), dimension(d_hyp), intent(in) :: hyp
    real(default), dimension(1), intent(out) :: me
    real(default), dimension(3) :: cart
    real(default), intent(in) :: pt2s
    real(default) :: p, p2
    if (d_hyp == 3) then
       p = hyp(3)
       p2 = hyp(3)**2
    else
       if (d_hyp == 2) then
          p = sqrt (pt2s)
          p2 = pt2s
       end if
    end if
    cart = h_to_c_param ([hyp(1), hyp(2), p2])
    if (p>pts_min .and. product (cart(1:2))>p2) then
       me(1) = const_pref * &
            alphasPDF(dble (p*gev_pt_max))**2 * &
            ps_io_pol (109, p2) * &
            pdf_in_in_kind (109, 11, cart(1), cart(2), p2) * &
            denom_param ([hyp(1), hyp(2), p2]) / product(cart(1:2))
    else
       me(1) = zero
    end if
  end subroutine cuba_gg_me_param

\eatline
\nwendcode{}\nwbegindocs{882}\nwdocspar
\nwenddocs{}\nwbegincode{883}\moddef{Muli interactions: public}\plusendmoddef
  public :: cuba_gg_me_ort
\nwendcode{}\nwbegincode{884}\moddef{Muli interactions: procedures}\plusendmoddef
  subroutine cuba_gg_me_ort (d_hyp, hyp, d_me, me, pt2s)
    integer, intent(in) :: d_hyp, d_me
    real(default), dimension(d_hyp), intent(in) :: hyp
    real(default), dimension(1), intent(out) :: me
    real(default), dimension(3) :: cart
    real(default), intent(in) :: pt2s
    real(default) :: p, p2
    if (d_hyp == 3) then
       p = hyp(3)
       p2 = hyp(3)**2
    else
       if (d_hyp == 2) then
          p = sqrt(pt2s)
          p2 = pt2s
       end if
    end if
    cart = h_to_c_ort ([hyp(1), cart(2), p2])
    if (p > pts_min .and. product (cart(1:2)) > p2) then
       me(1) = const_pref * &
            alphasPDF(dble (p*gev_pt_max))**2 * &
            ps_io_pol (109, p2) * &
            pdf_in_in_kind (109, 11, cart(1), cart(2), p2) * &
            denom_ort ([hyp(1), hyp(2), p2]) / product (cart(1:2))
    else
       me(1) = zero
    end if
  end subroutine cuba_gg_me_ort

\eatline
\nwendcode{}\nwbegindocs{885}\nwdocspar
\nwenddocs{}\nwbegincode{886}\moddef{Muli interactions: public}\plusendmoddef
  public :: cuba_gg_me_cart
\nwendcode{}\nwbegincode{887}\moddef{Muli interactions: procedures}\plusendmoddef
  subroutine cuba_gg_me_cart (d_cart, cart, d_me, me, pt2s)
    integer, intent(in) :: d_cart, d_me
    real(default), dimension(d_cart), intent(in) :: cart
    real(default), dimension(1), intent(out) :: me
    real(default), intent(in) :: pt2s
    real(default) :: a, p, p2
    if (d_cart == 3) then
       p = cart(3)
       p2 = cart(3)**2
    else
       if (d_cart == 2) then
          p = sqrt (pt2s)
          p2 = pt2s
       end if
    end if
    a = product (cart(1:2))
    if (p > pts_min .and. a > p2) then
       me(1) = const_pref * &
            alphasPDF (dble (p*gev_pt_max))**2 * &
            ps_io_pol (109, p2) * &
            pdf_in_in_kind (109, 11, cart(1), cart(2), p2) * &
            denom_cart ([cart(1), cart(2), p2]) / a
    else
       me(1) = zero
    end if
  end subroutine cuba_gg_me_cart

\eatline
\nwendcode{}\nwbegindocs{888}\nwdocspar
\nwenddocs{}\nwbegincode{889}\moddef{Muli interactions: public}\plusendmoddef
  public :: interactions_proton_proton_integrand_generic_17_reg
\nwendcode{}\nwbegincode{890}\moddef{Muli interactions: procedures}\plusendmoddef
  subroutine interactions_proton_proton_integrand_generic_17_reg & 
       (hyp_2, trafo, f, pt)
    real(default), dimension(2), intent(in) :: hyp_2
    procedure(coord_hcd_in) :: trafo
    real(default), dimension(17), intent(out) :: f
    class(transverse_mom_t), intent(in) :: pt
    real(default), dimension(3) :: cart, hyp_3
    real(default), dimension(5) :: psin
    real(double), dimension(-6:6) :: c_dble, d_dble
    real(default), dimension(-6:6) :: c, d
    real(default) :: gev_pt, gev2_pt, pts, pt2s, pt2shat, a, &
         pdf_seaquark_seaquark, pdf_seaquark_gluon, pdf_gluon_gluon, &
         pdf_up_seaquark, pdf_up_gluon, pdf_down_seaquark, pdf_down_gluon, &
         v1u, v1d, v2u, v2d, denom
        
    pts = pt%get_unit_scale()
    pt2s = pt%get_unit2_scale()
    gev_pt = pt%get_gev_scale()
    gev2_pt = pt%get_gev2_scale()

    hyp_3(1:2) = hyp_2
    hyp_3(3) = pt2s
    call trafo (hyp_3, cart, denom)
    a = product (cart(1:2))
    if (cart(1) <= one .and. cart(2) <= one .and. a > pt2s) then
       pt2shat = pt2s / a
       ! phase space polynom
       psin = matmul ([one, pt2shat, pt2shat**2, pt2shat**3], &
            phase_space_coefficients_in)
       ! pdf
       call evolvepdf (dble (cart(1)), dble (gev_pt), c_dble)
       call evolvepdf (dble (cart(2)), dble (gev_pt), d_dble)
       c = c_dble
       d = d_dble
       ! c = [1,1,1,1,1,1,1,1,1,1,1,1,1]*1D0
       ! d = c
       v1d = c(1) - c(-1)
       v1u = c(2) - c(-2)
       v2d = d(1) - d(-1)
       v2u = d(2) - d(-2)
       c(1) = c(-1)
       c(2) = c(-2)
       d(1) = d(-1)
       d(2) = d(-2)
       f(1) = zero
       !!! gluon_gluon
       f( 2) = (c(0)*d(0)) * psin(5)
            !!! type5
       !!! gluon_seaquark
       f( 3) = (c(0)*d(-4) + c(0)*d(-3) + c(0)*d(-2) + c(0)*d(-1) + &
                c(0)*d(1)  + c(0)*d(2)  + c(0)*d(3)  + c(0)*d(4)) * psin(4)
            !!! type4
       !!! gluon_down
       f( 4) = (c(0)*v2d) * psin(4)
            !!! type4
       !!! gluon_up
       f( 5) = (c(0)*v2u) * psin(4)
            !!! type4
       !!! seaquark_gluon
       f( 6) = (c(-4)*d(0) + c(-3)*d(0) + c(-2)*d(0) + c(-1)*d(0) + &
                c(1)*d(0)  + c(2)*d(0)  + c(3)*d(0)  + c(4)*d(0)) * psin(4)
            !!! type4
       !!! seaquark_seaquark
       f( 7) = &
            !!! type1
               (c(-4)*d(-3) + c(-4)*d(-2) + c(-4)*d(-1) + c(-4)*d( 1) + &
                c(-4)*d( 2) + c(-4)*d( 3) + c(-3)*d(-4) + c(-3)*d(-2) + &
                c(-3)*d(-1) + c(-3)*d( 1) + c(-3)*d( 2) + c(-3)*d( 4) + &
                c(-2)*d(-4) + c(-2)*d(-3) + c(-2)*d(-1) + c(-2)*d( 1) + &
                c(-2)*d( 3) + c(-2)*d( 4) + c(-1)*d(-4) + c(-1)*d(-3) + &
                c(-1)*d(-2) + c(-1)*d( 2) + c(-1)*d( 3) + c(-1)*d( 4) + &
                c( 1)*d(-4) + c( 1)*d(-3) + c( 1)*d(-2) + c( 1)*d( 2) + &
                c( 1)*d( 3) + c( 1)*d( 4) + c( 2)*d(-4) + c( 2)*d(-3) + &
                c( 2)*d(-1) + c( 2)*d( 1) + c( 2)*d( 3) + c( 2)*d( 4) + &
                c( 3)*d(-4) + c( 3)*d(-2) + c( 3)*d(-1) + c( 3)*d( 1) + &
                c( 3)*d( 2) + c( 3)*d( 4) + c( 4)*d(-3) + c( 4)*d(-2) + &
                c( 4)*d(-1) + c( 4)*d( 1) + c( 4)*d( 2) + c( 4)*d( 3)) * &
                psin(1) + &
            !!! type2
               (c(-4)*d(-4) + c(-3)*d(-3) + c(-2)*d(-2) + c(-1)*d(-1) + &
                c( 4)*d( 4) + c( 3)*d( 3) + c(2)*d( 2)  + c(1)*d( 1)) * &
                psin(2) + &
            !!! type3
               (c(-4)*d( 4) + c(-3)*d( 3) + c(-2)*d( 2) + c(-1)*d( 1) + &
                c( 4)*d(-4) + c( 3)*d(-3) + c(2)*d(-2)  + c(1)*d(-1)) * &
                psin(3)
       !!! seaquark_down
       f( 8) = &
            !!! type1
               (c(-4)*v2d + c(-3)*v2d + c(-2)*v2d + c( 2)*v2d + &
                c( 3)*v2d + c( 4)*v2d) * psin(1) + &
            !!! type2
                c( 1)*v2d * psin(2) + &
            !!! type3
                c(-1)*v2d * psin(3)
       !!! seaquark_up
       f( 9) = &
            !!! type1
               (c(-4)*v2u + c(-3)*v2u + c(-1)*v2u + c( 1)*v2u + &
                c( 3)*v2u + c( 4)*v2u) * psin(1) + &
            !!! type2
                c(2)*v2u * psin(2) + &
            !!! type3
                c(-2)*v2u * psin(3)
       !!! down_gluon
       f(10) = (v1d*d( 0)) * psin(4)
            !!! type4
       !!! down_seaquark
       f(11) = &
            !!! type1
               (v1d*d(-4) + v1d*d(-3) + v1d*d(-2) + v1d*d( 2) + &
                v1d*d( 3) + v1d*d( 4)) * psin(1) + &
            !!! type2
                v1d*d( 1) * psin(2) + &
            !!! type3
                v1d*d(-1) * psin(3)
       !!! down_down
       f(12) = v1d*v2d * psin(2)
       !!! down_up
       f(13) = v1d*v2u * psin(1)
       !!! up_gluon
       f(14) = (v1u*d(0)) * psin(4)
            !!! type4
       !!! up_seaquark
       f(15) = &
            !!! type1
               (v1u*d(-4) + v1u*d(-3) + v1u*d(-1) + v1u*d( 1) + &
                v1u*d( 3) + v1u*d( 4)) * psin(1) + &
            !!! type2
                v1u*d(2) * psin(2) + &
            !!! type3
                v1u*d(-2) * psin(3)
       !!! up_down
       f(16) = v1u * v2d * psin(1)
       !!! up_up
       f(17) = v1u * v2u * psin(2)
       f=f * const_pref &
           * alphasPDF (dble (sqrt(gev2_pt+gev2_p_t_0)))**2 &
           * denom / a
       ! print *, const_pref, alphasPDF(gev_pt)**2, denom_smooth (hyp), a
    else
       f = [zero, zero, zero, zero, zero, zero, zero, zero, zero, &
            zero, zero, zero, zero, zero, zero, zero, zero]
    end if
    ! print *, pt2shat, c(0)*d(0), psin(5), const_pref, &
    !     alphasPDF(gev_pt)**2, denom, a
  end subroutine interactions_proton_proton_integrand_generic_17_reg

\eatline
\nwendcode{}\nwbegindocs{891}\nwdocspar
\nwenddocs{}\nwbegincode{892}\moddef{Muli interactions: procedures}\plusendmoddef
  !  subroutine coordinates_proton_proton_integrand_cart_11 &
  !       (d_hyp, hyp_2, d_f, f)
  !    integer, intent(in) :: d_hyp, d_f
  !    real(default), dimension(2), intent(in) :: hyp_2
  !    real(default), dimension(11), intent(out) :: f
  !    call coordinates_proton_proton_integrand_generic_11 &
  !       (hyp_2, coordinates_hcd_cart, f)
  !    ! write (51,*) hyp_2, momentum_get_pts_scale(), f
  !  end subroutine coordinates_proton_proton_integrand_cart_11

\eatline
\nwendcode{}\nwbegindocs{893}\nwdocspar
\nwenddocs{}\nwbegincode{894}\moddef{Muli interactions: procedures}\plusendmoddef
  !  subroutine coordinates_proton_proton_integrand_ort_11 &
  !       (d_hyp, hyp_2, d_f, f)
  !    integer, intent(in) :: d_hyp, d_f
  !    real(default), dimension(2), intent(in) :: hyp_2
  !    real(default), dimension(11), intent(out) :: f
  !    call coordinates_proton_proton_integrand_generic_11 &
  !       (hyp_2, coordinates_hcd_ort, f)
  !    ! write (52,*) hyp_2, momentum_get_pts_scale(), f
  !  end subroutine coordinates_proton_proton_integrand_ort_11

\eatline
\nwendcode{}\nwbegindocs{895}\nwdocspar
\nwenddocs{}\nwbegincode{896}\moddef{Muli interactions: procedures}\plusendmoddef
  !  subroutine coordinates_proton_proton_integrand_param_11 &
  !       (d_hyp, hyp_2, d_f, f)
  !    integer, intent(in) :: d_hyp, d_f
  !    real(default), dimension(2), intent(in) :: hyp_2
  !    real(default), dimension(11), intent(out) :: f
  !    call coordinates_proton_proton_integrand_generic_11 &
  !       (hyp_2, coordinates_hcd_param, f)
  !    ! write(53,*) hyp_2, momentum_get_pts_scale(), f
  !  end subroutine coordinates_proton_proton_integrand_param_11

\eatline
\nwendcode{}\nwbegindocs{897}\nwdocspar
\nwenddocs{}\nwbegincode{898}\moddef{Muli interactions: procedures}\plusendmoddef
  !  subroutine coordinates_proton_proton_integrand_smooth_11 &
  !       (d_hyp, hyp_2, d_f, f)
  !    integer, intent(in)::d_hyp,d_f
  !    real(default), dimension(2), intent(in) :: hyp_2
  !    real(default), dimension(11), intent(out) :: f
  !    call coordinates_proton_proton_integrand_generic_11 (hyp_2, &
  !       coordinates_hcd_smooth, f)
  !    ! write (54,*) hyp_2, momentum_get_pts_scale(), f
  !  end subroutine coordinates_proton_proton_integrand_smooth_11

\eatline
\nwendcode{}\nwbegindocs{899}\nwdocspar
\nwenddocs{}\nwbegincode{900}\moddef{Muli interactions: public}\plusendmoddef
  public :: interactions_proton_proton_integrand_param_17_reg
\nwendcode{}\nwbegincode{901}\moddef{Muli interactions: procedures}\plusendmoddef
  subroutine interactions_proton_proton_integrand_param_17_reg &
       (d_hyp, hyp_2, d_f, f, pt)
    integer, intent(in) :: d_hyp, d_f
    real(default), dimension(2), intent(in) :: hyp_2
    real(default), dimension(17), intent(out) :: f
    class(transverse_mom_t), intent(in) :: pt
    call interactions_proton_proton_integrand_generic_17_reg &
         (hyp_2, coordinates_hcd_param_reg, f, pt)
    ! write (53,*)  hyp_2,momentum_get_pts_scale(),f
  end subroutine interactions_proton_proton_integrand_param_17_reg
  
\eatline
\nwendcode{}\nwbegindocs{902}\nwdocspar
\nwenddocs{}\nwbegincode{903}\moddef{Muli interactions: public}\plusendmoddef
  public :: interactions_proton_proton_integrand_smooth_17_reg
\nwendcode{}\nwbegincode{904}\moddef{Muli interactions: procedures}\plusendmoddef
  subroutine interactions_proton_proton_integrand_smooth_17_reg &
       (d_hyp, hyp_2, d_f, f, pt)
    integer, intent(in) :: d_hyp, d_f
    real(default), dimension(2), intent(in) :: hyp_2
    real(default), dimension(17), intent(out) :: f
    class(transverse_mom_t), intent(in) :: pt
    call interactions_proton_proton_integrand_generic_17_reg &
         (hyp_2, coordinates_hcd_smooth_reg, f, pt)
    ! write (53,*)hyp_2,momentum_get_pts_scale(), f
  end subroutine interactions_proton_proton_integrand_smooth_17_reg

\eatline
\nwendcode{}\nwbegindocs{905}\nwdocspar
%%%%%%%%%%%%%%%%%%%%%%%%%%%%%%%%%%%%%%%%%%%%%%%%%%%%%%%%%%%%%%%%%%%%%%%%
\section{VEGAS and CUBA integration routines}

This file contains the module {\tt{}muli{\char95}cuba}, a wrapper for the CUBA
integration library. Different algorithms and settings have been tried out 
for the integration, including VEGAS, and this wrapper has been mainly
written for that purpose.

\nwenddocs{}\nwbegincode{906}\moddef{\code{}muli{\char95}cuba.f90\edoc{}}\endmoddef
\LA{}File header\RA{}

module muli_cuba
  
\LA{}Use kinds\RA{}  
  use constants
  use diagnostics
  use muli_base
  use muli_momentum

\LA{}Standard module head\RA{}  

\LA{}Muli CUBA: public\RA{}

\LA{}Muli CUBA: variables\RA{}

\LA{}Muli CUBA: types\RA{}

\LA{}Muli CUBA: interfaces\RA{}

contains
  
\LA{}Muli CUBA: procedures\RA{}  

end module muli_cuba
\eatline
\nwendcode{}\nwbegindocs{907}\nwdocspar
\nwenddocs{}\nwbegincode{908}\moddef{Muli CUBA: variables}\endmoddef
  integer, parameter :: max_maxeval = huge(1)
  
\eatline
\nwendcode{}\nwbegindocs{909}\nwdocspar
\nwenddocs{}\nwbegincode{910}\moddef{Muli CUBA: public}\endmoddef
  public :: cuba_class
\nwendcode{}\nwbegincode{911}\moddef{Muli CUBA: types}\endmoddef
  type, extends (ser_class_t), abstract :: cuba_class  
     real(default) :: start_time = zero
     real(default) :: stop_time = zero
     real(default) :: run_time = zero
     integer :: dim_x = 2
     integer :: dim_f = 1
     type(transverse_mom_t) :: userdata
     real(default) :: eps_rel = 1.E-3_default
     real(default) :: eps_abs = 0._default
     integer :: flags = 0
     integer :: seed = 1
     integer :: min_eval = 0
     integer :: max_eval = max_maxeval
     integer :: neval = 0
     integer, public :: fail = -1
     integer :: nregions = 0
     real(default), dimension(:), allocatable :: integral
     real(default), dimension(:), allocatable :: error
     real(default), dimension(:), allocatable :: prob
     procedure(integrand_interface), nopass, pointer :: integrand
   contains
   \LA{}Muli CUBA: cuba class: TBP\RA{}
  end type cuba_class

\eatline
\nwendcode{}\nwbegindocs{912}\nwdocspar
\nwenddocs{}\nwbegincode{913}\moddef{Muli CUBA: types}\plusendmoddef
  type, extends (cuba_class) :: cuba_cuhre_t
     private
     integer :: key = 13
   contains
   \LA{}Muli CUBA: cuba cuhre: TBP\RA{}     
  end type cuba_cuhre_t

\eatline
\nwendcode{}\nwbegindocs{914}\nwdocspar
\nwenddocs{}\nwbegincode{915}\moddef{Muli CUBA: types}\plusendmoddef
  type, extends (cuba_class) :: cuba_suave_t
     private
     integer :: nnew = 10000    !1000
     integer :: flatness = 5   !50
   contains
   \LA{}Muli CUBA: cuba suave: TBP\RA{}  
  end type cuba_suave_t

\eatline
\nwendcode{}\nwbegindocs{916}\nwdocspar
\nwenddocs{}\nwbegincode{917}\moddef{Muli CUBA: public}\plusendmoddef
  public :: cuba_divonne_t
\nwendcode{}\nwbegincode{918}\moddef{Muli CUBA: types}\plusendmoddef
  type, extends (cuba_class) :: cuba_divonne_t
     private
     integer :: key1 = 13
     integer :: key2 = 13
     integer :: key3 = 13
     integer :: maxpass = 2
     real(default) :: border = zero
     real(default) :: maxchisq = 10._default
     real(default) :: mindeviation = .25_default
     integer :: ngiven = 0
     integer :: ldxgiven = 0
     ! real(default), dimension(ldxgiven,ngiven) :: &
     !    xgiven = reshape( source = [ 0.0,0.0 ], shape = [2,1])
     real(default), dimension(:,:), allocatable :: xgiven
     ! real(default), dimension(2) :: xgiven = [1E-1_default, 5E-1_default]
     integer :: nextra = 0
   contains
   \LA{}Muli CUBA: cuba divonne: TBP\RA{}  
  end type cuba_divonne_t

\eatline
\nwendcode{}\nwbegindocs{919}\nwdocspar
\nwenddocs{}\nwbegincode{920}\moddef{Muli CUBA: types}\plusendmoddef
  type, extends (cuba_class) :: cuba_vegas_t
     private
     integer :: nstart = 500
     integer :: nincrease = 1000
     integer :: nbatch = 1000
     integer :: gridno = 0
     character(len=8), pointer :: statefile => null()
   contains
   \LA{}Muli CUBA: cuba vegas: TBP\RA{}  
  end type cuba_vegas_t

\eatline
\nwendcode{}\nwbegindocs{921}\nwdocspar
\nwenddocs{}\nwbegincode{922}\moddef{Muli CUBA: cuba class: TBP}\endmoddef
  procedure :: write_to_marker => cuba_write_to_marker
\nwendcode{}\nwbegincode{923}\moddef{Muli CUBA: procedures}\endmoddef
  subroutine cuba_write_to_marker (this, marker, status)
    class(cuba_class), intent(in) :: this
    class(marker_t), intent(inout) :: marker
    integer(dik), intent(out) :: status
    call marker%mark_begin ("cuba_class")
    call marker%mark ("dim_x", this%dim_x)
    call marker%mark ("dim_f", this%dim_f)    
    call marker%mark ("eps_rel", this%eps_rel)
    call marker%mark ("eps_abs", this%eps_abs)
    call marker%mark ("flags", this%flags)    
    call marker%mark ("min_eval", this%min_eval)
    call marker%mark ("max_eval", this%max_eval)
    call marker%mark ("neval", this%neval)
    call marker%mark ("fail", this%fail)
    call marker%mark ("nregions", this%nregions)
    if (allocated (this%integral)) then
       call marker%mark ("integral", this%integral)
    else
       call marker%mark_null ("integral")
    end if
    if (allocated(this%error)) then
       call marker%mark ("error", this%error)
    else
       call marker%mark_null ("error")

    end if
    if (allocated (this%prob)) then
       call marker%mark ("prob", this%prob)
    else
       call marker%mark_null ("prob")
    end if    
    call marker%mark_null ("cuba_class")
  end subroutine cuba_write_to_marker
  
\eatline
\nwendcode{}\nwbegindocs{924}\nwdocspar
\nwenddocs{}\nwbegincode{925}\moddef{Muli CUBA: cuba class: TBP}\plusendmoddef
  procedure :: read_from_marker => cuba_read_from_marker  
\nwendcode{}\nwbegincode{926}\moddef{Muli CUBA: procedures}\plusendmoddef
  subroutine cuba_read_from_marker (this, marker, status)
    class(cuba_class), intent(out) :: this
    class(marker_t), intent(inout) :: marker
    integer(dik), intent(out)::status
    call marker%pick_begin ("CUBA_CLASS", status=status)
    call marker%pick ("dim_x", this%dim_x,status)
    call marker%pick ("dim_f", this%dim_f,status)
    call marker%pick ("eps_rel", this%eps_rel,status)
    call marker%pick ("eps_abs", this%eps_abs,status)
    call marker%pick ("flags", this%flags,status)
    call marker%pick ("min_eval", this%min_eval,status)
    call marker%pick ("max_eval", this%max_eval,status)
    call marker%pick ("neval", this%neval,status)
    call marker%pick ("fail", this%fail,status)
    call marker%pick ("nregions", this%nregions,status)
    call marker%verify_nothing("integral",status)
    if (allocated (this%integral))  deallocate (this%integral)
    if (status == serialize_ok) then
       allocate (this%integral (this%dim_f))
       call marker%pick ("integral", this%integral, status)
    end if
    call marker%verify_nothing ("error", status)
    if (allocated (this%error))  deallocate (this%error)
    if (status == serialize_ok) then
       allocate (this%error (this%dim_f))
       call marker%pick ("error", this%error, status)
    end if
    call marker%verify_nothing ("prob", status)
    if (allocated (this%prob))  deallocate (this%prob)
    if (status == serialize_ok) then
       allocate (this%prob (this%dim_f))
       call marker%pick ("prob", this%prob, status)
    end if
    call marker%pick_end ("cuba_class", status)
  end subroutine cuba_read_from_marker
  
\eatline
\nwendcode{}\nwbegindocs{927}\nwdocspar
\nwenddocs{}\nwbegincode{928}\moddef{Muli CUBA: cuba class: TBP}\plusendmoddef
  procedure :: print_to_unit => cuba_print_to_unit    
\nwendcode{}\nwbegincode{929}\moddef{Muli CUBA: procedures}\plusendmoddef
  subroutine cuba_print_to_unit (this, unit, parents, components, peers)
    class(cuba_class), intent(in) :: this
    integer, intent(in) :: unit
    integer(dik), intent(in) :: parents, components, peers
    character(11) :: n
    write (n, '("(",I2,"(E12.4))")')  this%dim_f
    write (unit, "(1x,A)")  "Components of cuba_class:"
    write (unit, "(3x,A)")  "Parameters:"
    write (unit, "(3x,A,I10)")   "dim_f:     ", this%dim_f
    write (unit, "(3x,A,I10)")   "dim_x:     ", this%dim_x
    call this%userdata%print_to_unit (unit, parents, components-1, peers)
    write (unit, "(3x,A,E10.4)") "eps_rel:   ", this%eps_rel
    write (unit, "(3x,A,E10.4)") "eps_abs:   ", this%eps_abs
    write (unit, "(3x,A,I10)")   "flags:     ", this%flags
    write (unit, "(3x,A,I10)")   "seed:      ", this%seed
    write (unit, "(3x,A,I10)")   "min_eval:  ", this%min_eval
    write (unit, "(3x,A,I10)")   "max_eval:  ", this%max_eval  
    write (unit, "(3x,A)")  "Results:"
    write (unit, "(3x,A,I10)")   "neval:     ", this%neval
    write (unit, "(3x,A,I10)")   "fail:      ", this%fail
    write (unit, "(3x,A)", advance="no") "integral:  "
    write (unit, fmt=n)  this%integral
    write (unit, "(3x,A)", advance="no")  "error:     "
    write (unit, fmt=n)  this%error
    write (unit, "(3x,A)", advance="no")  "prob:      "
    write (unit, fmt=n)  this%prob
    write (unit, "(3x,A,E10.4)")  "time:      ", &
         this%stop_time - this%start_time
    ! write(unit,'("time:      ",E10.4)') this%run_time
  end subroutine cuba_print_to_unit

\eatline
\nwendcode{}\nwbegindocs{930}\nwdocspar
\nwenddocs{}\nwbegincode{931}\moddef{Muli CUBA: cuba class: TBP}\plusendmoddef
  generic   :: get_integral => get_integral_array, get_integral_1 
  procedure :: get_integral_array => cuba_get_integral_array  
  procedure :: get_integral_1 => cuba_get_integral_1  
\nwendcode{}\nwbegincode{932}\moddef{Muli CUBA: procedures}\plusendmoddef
  subroutine cuba_get_integral_array (this, integral)
    class(cuba_class) :: this
    real(default), intent(out), dimension(:) :: integral
    integral = this%integral
  end subroutine cuba_get_integral_array

\eatline
\nwendcode{}\nwbegindocs{933}\nwdocspar
\nwenddocs{}\nwbegincode{934}\moddef{Muli CUBA: procedures}\plusendmoddef
  subroutine cuba_get_integral_1 (this, integral)
    class(cuba_class) :: this
    real(default), intent(out) :: integral
    integral = this%integral(1)
  end subroutine cuba_get_integral_1
  
\eatline
\nwendcode{}\nwbegindocs{935}\nwdocspar
\nwenddocs{}\nwbegincode{936}\moddef{Muli CUBA: cuba class: TBP}\plusendmoddef
  procedure :: copy_common => cuba_copy_common  
\nwendcode{}\nwbegincode{937}\moddef{Muli CUBA: procedures}\plusendmoddef
  subroutine cuba_copy_common (this, source)
    class(cuba_class), intent(out) :: this
    class(cuba_class), intent(in) :: source
    this%dim_x = source%dim_x
    this%dim_f = source%dim_f
    this%eps_rel = source%eps_rel
    this%eps_abs = source%eps_abs
    this%flags = source%flags
    this%min_eval = source%min_eval
    this%max_eval = source%max_eval
    call this%alloc()
  end subroutine cuba_copy_common
  
\eatline
\nwendcode{}\nwbegindocs{938}\nwdocspar
\nwenddocs{}\nwbegincode{939}\moddef{Muli CUBA: cuba class: TBP}\plusendmoddef
  procedure :: set_common => cuba_set_common  
\nwendcode{}\nwbegindocs{940}\nwdocspar
\nwenddocs{}\nwbegincode{941}\moddef{Muli CUBA: procedures}\plusendmoddef
  subroutine cuba_set_common (this, dim_x, dim_f, eps_rel, eps_abs, &
       flags, seed, min_eval, max_eval, integrand, userdata)
    class(cuba_class), intent(inout) :: this
    integer, intent(in), optional :: dim_x, dim_f, flags, min_eval, &
         max_eval, seed
    real(default), intent(in), optional :: eps_rel,eps_abs
    type(transverse_mom_t), intent(in), optional :: userdata
    procedure(integrand_interface), optional :: integrand
    if (present (dim_x)) then
       call this%set_dim_x (dim_x)
    end if
    if (present (dim_f)) then
       call this%set_dim_f (dim_f)
    end if
    if (present (flags)) then
       this%flags = flags
    end if
    if (present (seed)) then
       this%seed = seed
    end if    
    if (present (min_eval)) then
       this%min_eval = min_eval
    end if
    if (present (max_eval)) then
       if (max_eval < max_maxeval) then
          this%max_eval = max_eval
       else
          call msg_warning &
               ("cuba_set_common: Value of max_eval is too large.")
          this%max_eval = max_maxeval
       end if
    end if
    if (present (eps_rel)) then
       this%eps_rel = eps_rel
    end if
    if (present (eps_abs)) then
       this%eps_abs = eps_abs
    end if
    if (present (integrand))  this%integrand=>integrand
    if (present (userdata))  this%userdata=userdata
  end subroutine cuba_set_common

\eatline
\nwendcode{}\nwbegindocs{942}\nwdocspar
\nwenddocs{}\nwbegincode{943}\moddef{Muli CUBA: cuba class: TBP}\plusendmoddef
  procedure :: set_dim_f => cuba_set_dim_f  
\nwendcode{}\nwbegincode{944}\moddef{Muli CUBA: procedures}\plusendmoddef
  subroutine cuba_set_dim_f (this, new_dim_f)
    class(cuba_class) :: this
    integer, intent(in) :: new_dim_f
    if (new_dim_f > 0) then
       this%dim_f = new_dim_f
       call this%alloc_dim_f
    else
       call msg_error ("cuba_set_dim_f: New value for dim_f is " &
            // "negative. dim_f is not set.")
    end if
  end subroutine cuba_set_dim_f

\eatline
\nwendcode{}\nwbegindocs{945}\nwdocspar
\nwenddocs{}\nwbegincode{946}\moddef{Muli CUBA: cuba class: TBP}\plusendmoddef
  procedure :: set_dim_x => cuba_set_dim_x
\nwendcode{}\nwbegincode{947}\moddef{Muli CUBA: procedures}\plusendmoddef
  subroutine cuba_set_dim_x (this, new_dim_x)
    class(cuba_class) :: this
    integer, intent(in) :: new_dim_x
    if (new_dim_x > 0) then
       this%dim_x = new_dim_x
    else
       call msg_error ("cuba_set_dim_x: New value for dim_x is " &
            // "negative. dim_x is not set.")
    end if
  end subroutine cuba_set_dim_x

\eatline
\nwendcode{}\nwbegindocs{948}\nwdocspar
\nwenddocs{}\nwbegincode{949}\moddef{Muli CUBA: cuba class: TBP}\plusendmoddef
  procedure :: reset_timer => cuba_reset_timer  
\nwendcode{}\nwbegincode{950}\moddef{Muli CUBA: procedures}\plusendmoddef
  subroutine cuba_reset_timer (this)
    class(cuba_class), intent(inout) :: this
    this%start_time = zero
    this%stop_time = zero
    this%run_time = zero
  end subroutine cuba_reset_timer

\eatline
\nwendcode{}\nwbegindocs{951}\nwdocspar
\nwenddocs{}\nwbegincode{952}\moddef{Muli CUBA: cuba class: TBP}\plusendmoddef
  procedure :: integrate_with_timer => cuba_integrate_with_timer 
\nwendcode{}\nwbegincode{953}\moddef{Muli CUBA: procedures}\plusendmoddef
  subroutine cuba_integrate_with_timer (this, integrand)
    class(cuba_class), intent(inout) :: this
    procedure(integrand_interface) :: integrand
    call cpu_time (this%start_time)
    call this%integrate (integrand)
    call cpu_time (this%stop_time)
    this%run_time = this%run_time + this%stop_time - this%start_time
  end subroutine cuba_integrate_with_timer

\eatline
\nwendcode{}\nwbegindocs{954}\nwdocspar
\nwenddocs{}\nwbegincode{955}\moddef{Muli CUBA: cuba class: TBP}\plusendmoddef
  procedure :: integrate_associated => cuba_integrate_associated  
\nwendcode{}\nwbegincode{956}\moddef{Muli CUBA: procedures}\plusendmoddef
  subroutine cuba_integrate_associated (this)
    class(cuba_class), intent(inout) :: this
    call this%integrate_with_timer (this%integrand)
  end subroutine cuba_integrate_associated

\eatline
\nwendcode{}\nwbegindocs{957}\nwdocspar
\nwenddocs{}\nwbegincode{958}\moddef{Muli CUBA: cuba class: TBP}\plusendmoddef
  generic :: integrate => integrate_nd, integrate_userdata
  procedure(integrate_interface), deferred :: integrate_nd
  procedure(integrate_userdata_interface), deferred :: integrate_userdata
  procedure(cuba_copy_interface), deferred :: copy  
\nwendcode{}\nwbegindocs{959}\nwdocspar
\nwenddocs{}\nwbegincode{960}\moddef{Muli CUBA: cuba class: TBP}\plusendmoddef
  procedure :: dealloc_dim_f => cuba_dealloc_dim_f
\nwendcode{}\nwbegincode{961}\moddef{Muli CUBA: procedures}\plusendmoddef
  subroutine cuba_dealloc_dim_f(this)
    class(cuba_class) :: this
    ! print '("cuba_dealloc_dim_f...")'
    if (allocated (this%integral)) then
       deallocate (this%integral)
    end if
    if (allocated (this%error)) then
       deallocate (this%error)
    end if
    if (allocated (this%prob)) then
       deallocate (this%prob)
    end if
    ! print '("done")'
  end subroutine cuba_dealloc_dim_f

\eatline
\nwendcode{}\nwbegindocs{962} 
\nwenddocs{}\nwbegincode{963}\moddef{Muli CUBA: cuba class: TBP}\plusendmoddef
  procedure :: alloc_dim_f => cuba_alloc_dim_f
\nwendcode{}\nwbegincode{964}\moddef{Muli CUBA: procedures}\plusendmoddef
  subroutine cuba_alloc_dim_f (this)
    class(cuba_class) :: this
    call this%dealloc_dim_f ()
    allocate (this%integral (this%dim_f))
    allocate (this%error (this%dim_f))
    allocate (this%prob (this%dim_f))
  end subroutine cuba_alloc_dim_f

\eatline
\nwendcode{}\nwbegindocs{965}\nwdocspar
\nwenddocs{}\nwbegincode{966}\moddef{Muli CUBA: cuba class: TBP}\plusendmoddef
  procedure :: dealloc => cuba_dealloc
\nwendcode{}\nwbegincode{967}\moddef{Muli CUBA: procedures}\plusendmoddef
  subroutine cuba_dealloc (this)
    class(cuba_class) :: this
    call this%dealloc_dim_f
  end subroutine cuba_dealloc
  
\eatline
\nwendcode{}\nwbegindocs{968}\nwdocspar
\nwenddocs{}\nwbegincode{969}\moddef{Muli CUBA: cuba class: TBP}\plusendmoddef
  procedure :: alloc => cuba_alloc       
\nwendcode{}\nwbegincode{970}\moddef{Muli CUBA: procedures}\plusendmoddef
  subroutine cuba_alloc (this)
    class(cuba_class) :: this
    call this%alloc_dim_f
  end subroutine cuba_alloc
  
\eatline
\nwendcode{}\nwbegindocs{971}\nwdocspar
\nwenddocs{}\nwbegincode{972}\moddef{Muli CUBA: cuba vegas: TBP}\endmoddef
  procedure :: write_to_marker => cuba_vegas_write_to_marker  
\nwendcode{}\nwbegincode{973}\moddef{Muli CUBA: procedures}\plusendmoddef
  subroutine cuba_vegas_write_to_marker (this, marker, status)
    class(cuba_vegas_t), intent(in) :: this
    class(marker_t), intent(inout) :: marker
    integer(dik), intent(out) :: status
    call marker%mark_begin ("cuba_vegas_t")
    call cuba_write_to_marker (this, marker, status)
    call marker%mark("nstart", this%nstart)
    call marker%mark("nincrease", this%nincrease)
    call marker%mark_null ("cuba_vegas_t")
  end subroutine cuba_vegas_write_to_marker
  
\eatline
\nwendcode{}\nwbegindocs{974}\nwdocspar
\nwenddocs{}\nwbegincode{975}\moddef{Muli CUBA: cuba vegas: TBP}\plusendmoddef
  procedure :: read_from_marker => cuba_vegas_read_from_marker
\nwendcode{}\nwbegincode{976}\moddef{Muli CUBA: procedures}\plusendmoddef
  subroutine cuba_vegas_read_from_marker (this, marker, status)
    class(cuba_vegas_t), intent(out) :: this
    class(marker_t), intent(inout) :: marker
    integer(dik), intent(out) :: status
    call marker%pick_begin ("cuba_vegas_t", status=status)
    call cuba_read_from_marker (this, marker, status)
    call marker%pick ("nstart", this%nstart, status)
    call marker%pick ("nincrease", this%nincrease, status)
    call marker%pick_end ("cuba_vegas_t", status)
  end subroutine cuba_vegas_read_from_marker
  
\eatline
\nwendcode{}\nwbegindocs{977}\nwdocspar
\nwenddocs{}\nwbegincode{978}\moddef{Muli CUBA: cuba vegas: TBP}\plusendmoddef
  procedure :: print_to_unit => cuba_vegas_print_to_unit  
\nwendcode{}\nwbegincode{979}\moddef{Muli CUBA: procedures}\plusendmoddef
  subroutine cuba_vegas_print_to_unit(this,unit,parents,components,peers)
    class(cuba_vegas_t), intent(in) :: this
    INTEGER, INTENT(IN) :: unit
    integer(dik), intent(in)::parents,components,peers
    if (parents>0)call cuba_print_to_unit(this,unit,parents-1,components,peers)
    write (unit, "(1x,A)")       "Components of cuba_vegas_t:"
    write (unit, "(3x,A,I10)")   "nstart:    ", this%nstart
    write (unit, "(3x,A,I10)")   "nincrease: ", this%nincrease
    write (unit, "(3x,A,I10)")   "nbatch:    ", this%nbatch
    write (unit, "(3x,A,I10)")   "gridno:    ", this%gridno
    if (associated (this%statefile)) then
       write (unit, "(3x,A,A)")  "statefile: ", this%statefile
    else
       write (unit, "(3x,A)")    "statefile:    not associated"
    end if
  end subroutine cuba_vegas_print_to_unit
  
\eatline
\nwendcode{}\nwbegindocs{980}\nwdocspar
\nwenddocs{}\nwbegincode{981}\moddef{Muli CUBA: cuba vegas: TBP}\plusendmoddef
  procedure, nopass :: get_type => cuba_vegas_get_type    
\nwendcode{}\nwbegincode{982}\moddef{Muli CUBA: procedures}\plusendmoddef
  pure subroutine cuba_vegas_get_type (type)
    character(:), allocatable, intent(out) :: type
    allocate (type, source="cuba_vegas_t")
  end subroutine cuba_vegas_get_type
  
\eatline
\nwendcode{}\nwbegindocs{983}\nwdocspar
\nwenddocs{}\nwbegincode{984}\moddef{Muli CUBA: cuba vegas: TBP}\plusendmoddef
  procedure :: integrate_nd => integrate_vegas  
\nwendcode{}\nwbegincode{985}\moddef{Muli CUBA: procedures}\plusendmoddef
  subroutine integrate_vegas (this, integrand)
    class(cuba_vegas_t), intent(inout) :: this
    procedure(integrand_interface) :: integrand
    ! print '("vegas")'
    ! call vegas (this%dim_x, this%dim_f, integrand, this%userdata, &
    !      this%eps_rel, this%eps_abs, this%flags, this%seed, &
    !      this%min_eval, this%max_eval, this%nstart, this%nincrease, &
    !      this%nbatch, this%gridno, this%statefile, this%neval, &
    !      this%fail, this%integral, this%error, this%prob)
  end subroutine integrate_vegas
  
\eatline
\nwendcode{}\nwbegindocs{986}\nwdocspar
\nwenddocs{}\nwbegincode{987}\moddef{Muli CUBA: cuba vegas: TBP}\plusendmoddef
  procedure :: integrate_userdata => integrate_vegas_userdata  
\nwendcode{}\nwbegincode{988}\moddef{Muli CUBA: procedures}\plusendmoddef
  subroutine integrate_vegas_userdata (this, integrand, userdata)
    class(cuba_vegas_t), intent(inout) :: this
    procedure(integrand_interface) :: integrand
    class(transverse_mom_t), intent(in) :: userdata
    ! print '("vegas")'
    ! call vegas(this%dim_x, this%dim_f, integrand, userdata, this%eps_rel, &
    !      this%eps_abs, this%flags, this%seed, this%min_eval, this%max_eval, &
    !      this%nstart, this%nincrease, this%nbatch, this%gridno, &
    !      this%statefile, this%neval, this%fail, this%integral, &
    !      this%error, this%prob)
  end subroutine integrate_vegas_userdata
  
\eatline
\nwendcode{}\nwbegindocs{989}\nwdocspar
\nwenddocs{}\nwbegincode{990}\moddef{Muli CUBA: cuba vegas: TBP}\plusendmoddef
  procedure :: copy => cuba_vegas_copy  
\nwendcode{}\nwbegincode{991}\moddef{Muli CUBA: procedures}\plusendmoddef
  subroutine cuba_vegas_copy (this, source)
    class(cuba_vegas_t), intent(out) :: this
    class(cuba_class), intent(in) :: source
    select type (source)
    class is (cuba_vegas_t)
       call this%copy_common (source)
       this%nstart = source%nstart
       this%nincrease = source%nincrease
    class default
       call msg_error ("cuba_vegas_copy: type of source is not type " &
            // "compatible with cuba_vegas_t.")
    end select
  end subroutine cuba_vegas_copy
    
\eatline
\nwendcode{}\nwbegindocs{992}\nwdocspar
\nwenddocs{}\nwbegincode{993}\moddef{Muli CUBA: cuba vegas: TBP}\plusendmoddef
  procedure :: set_deferred => cuba_vegas_set_deferred  
\nwendcode{}\nwbegincode{994}\moddef{Muli CUBA: procedures}\plusendmoddef
  subroutine cuba_vegas_set_deferred &
       (this, n_start, n_increase, nbatch, gridno, statefile)
    class(cuba_vegas_t), intent(inout) :: this  
    integer, intent(in), optional :: n_start, n_increase, nbatch, gridno
    character(len=*), intent(in), target, optional :: statefile
    if (present (n_start))  this%nstart = n_start
    if (present (n_increase))  this%nincrease = n_increase
    if (present (nbatch))  this%nbatch = nbatch
    if (present (gridno))  this%gridno = gridno
    if (present (statefile))  this%statefile => statefile
  end subroutine cuba_vegas_set_deferred
  
\eatline
\nwendcode{}\nwbegindocs{995}\nwdocspar
\nwenddocs{}\nwbegincode{996}\moddef{Muli CUBA: cuba divonne: TBP}\endmoddef
  procedure :: write_to_marker => cuba_divonne_write_to_marker  
\nwendcode{}\nwbegincode{997}\moddef{Muli CUBA: procedures}\plusendmoddef
  subroutine cuba_divonne_write_to_marker (this, marker, status)
    class(cuba_divonne_t), intent(in) :: this
    class(marker_t), intent(inout) :: marker
    integer(dik), intent(out) :: status
    call marker%mark_begin ("cuba_divonne_t")
    call cuba_write_to_marker (this, marker, status)
    call marker%mark ("key1", this%key1)
    call marker%mark ("key2", this%key2)
    call marker%mark ("key3", this%key3)
    call marker%mark ("maxpass", this%maxpass)
    call marker%mark ("border", this%border)
    call marker%mark ("maxchisq", this%maxchisq)
    call marker%mark ("mindeviation", this%mindeviation)
    call marker%mark ("ngiven", this%ngiven)
    call marker%mark ("ldxgiven", this%ldxgiven)
    call marker%mark ("nextra", this%nextra)
    call marker%mark ("xgiven", this%xgiven)
    call marker%mark_null ("cuba_divonne_t")
  end subroutine cuba_divonne_write_to_marker
    
\eatline
\nwendcode{}\nwbegindocs{998}\nwdocspar
\nwenddocs{}\nwbegincode{999}\moddef{Muli CUBA: cuba divonne: TBP}\plusendmoddef
  procedure :: read_from_marker => cuba_divonne_read_from_marker  
\nwendcode{}\nwbegincode{1000}\moddef{Muli CUBA: procedures}\plusendmoddef
  subroutine cuba_divonne_read_from_marker (this, marker, status)
    class(cuba_divonne_t), intent(out) :: this
    class(marker_t), intent(inout) :: marker
    integer(dik), intent(out) :: status
    call marker%pick_begin ("cuba_divonne_t", status=status)
    call cuba_read_from_marker (this, marker, status)
    call marker%pick ("key1", this%key1, status)
    call marker%pick ("key2", this%key2, status)
    call marker%pick ("key3", this%key3, status)
    call marker%pick ("maxpass", this%maxpass, status)
    call marker%pick ("border", this%border, status)
    call marker%pick ("maxchisq", this%maxchisq, status)
    call marker%pick ("mindeviation", this%mindeviation, status)
    call marker%pick ("ngiven", this%ngiven, status)
    call marker%pick ("ldxgiven", this%ldxgiven, status)
    call marker%pick ("nextra", this%nextra, status)
    if (allocated (this%xgiven))  deallocate (this%xgiven)
    allocate (this%xgiven (this%ldxgiven, this%ngiven))
    call marker%pick ("xgiven", this%xgiven, status)
    call marker%pick_end ("cuba_divonne_t", status)
  end subroutine cuba_divonne_read_from_marker
    
\eatline
\nwendcode{}\nwbegindocs{1001}\nwdocspar
\nwenddocs{}\nwbegincode{1002}\moddef{Muli CUBA: cuba divonne: TBP}\plusendmoddef
  procedure :: print_to_unit => cuba_divonne_print_to_unit  
\nwendcode{}\nwbegincode{1003}\moddef{Muli CUBA: procedures}\plusendmoddef
  subroutine cuba_divonne_print_to_unit (this, unit, parents, components, peers)
    class(cuba_divonne_t), intent(in) :: this
    INTEGER, INTENT(IN) :: unit
    integer(dik), intent(in) :: parents, components, peers
    if (parents > 0)  &
         call cuba_print_to_unit (this, unit, parents-1, components, peers)
    write (unit, "(1x,A)")          "Components of cuba_divonne_t:"
    write (unit, "(3x,A,I10)")      "key1:        ", this%key1
    write (unit, "(3x,A,I10)")      "key2:        ", this%key2
    write (unit, "(3x,A,I10)")      "key3:        ", this%key3
    write (unit, "(3x,A,I10)")      "maxpass:     ", this%maxpass
    write (unit, "(3x,A,I10)")      "ngiven:      ", this%ngiven
    write (unit, "(3x,A,I10)")      "ldxgiven:    ", this%ldxgiven
    write (unit, "(3x,A,I10)")      "nextra:      ", this%nextra
    write (unit, "(3x,A,E10.4)")    "border:      ", this%border
    write (unit, "(3x,A,E10.4)")    "maxchisq:    ", this%maxchisq
    write (unit, "(3x,A,E10.4)")    "mindeviation:", this%mindeviation
    write (unit, "(3x,A,2(E10.4))") "xgiven:      ", this%xgiven
  end subroutine cuba_divonne_print_to_unit
  
\eatline
\nwendcode{}\nwbegindocs{1004}\nwdocspar
\nwenddocs{}\nwbegincode{1005}\moddef{Muli CUBA: cuba divonne: TBP}\plusendmoddef
  procedure, nopass :: get_type => cuba_divonne_get_type   
\nwendcode{}\nwbegincode{1006}\moddef{Muli CUBA: procedures}\plusendmoddef
  pure subroutine cuba_divonne_get_type (type)
    character(:), allocatable, intent(out) :: type
    allocate(type, source="cuba_divonne_t")
  end subroutine cuba_divonne_get_type
  
\eatline
\nwendcode{}\nwbegindocs{1007}\nwdocspar
\nwenddocs{}\nwbegincode{1008}\moddef{Muli CUBA: cuba divonne: TBP}\plusendmoddef
  procedure :: integrate_nd => integrate_divonne
\nwendcode{}\nwbegincode{1009}\moddef{Muli CUBA: procedures}\plusendmoddef
  subroutine integrate_divonne (this, integrand)
    class(cuba_divonne_t), intent(inout) :: this  
    procedure(integrand_interface) :: integrand
    ! call this%reset_output()
    ! print '("divonne")'
    ! call divonne(this%dim_x, this%dim_f, integrand, this%userdata, &
    !      this%eps_rel, this%eps_abs, this%flags, this%seed, this%min_eval, &
    !      this%max_eval, this%key1, this%key2, this%key3, this%maxpass, &
    !      this%border, this%maxchisq, this%mindeviation, this%ngiven, &
    !      this%ldxgiven, this%xgiven, this%nextra, &
    !      ! this%peakfinder, &
    !      0, this%nregions, this%neval, this%fail, this%integral, &
    !      this%error, this%prob)
  end subroutine integrate_divonne

\eatline
\nwendcode{}\nwbegindocs{1010}\nwdocspar
\nwenddocs{}\nwbegincode{1011}\moddef{Muli CUBA: cuba divonne: TBP}\plusendmoddef
  procedure :: integrate_userdata => integrate_divonne_userdata
\nwendcode{}\nwbegincode{1012}\moddef{Muli CUBA: procedures}\plusendmoddef
  subroutine integrate_divonne_userdata (this, integrand, userdata)
    class(cuba_divonne_t), intent(inout) :: this  
    procedure(integrand_interface) :: integrand
    class(transverse_mom_t), intent(in) :: userdata
    ! call this%reset_output()
    ! print '("divonne")'
    ! call divonne (this%dim_x, this%dim_f, integrand, userdata, &
    !      this%eps_rel, this%eps_abs, this%flags, this%seed, this%min_eval, &
    !      this%max_eval, this%key1, this%key2, this%key3, this%maxpass, &
    !      this%border, this%maxchisq, this%mindeviation, this%ngiven, &
    !      this%ldxgiven, this%xgiven, this%nextra, &
    !      ! this%peakfinder, &
    !      0, this%nregions, this%neval, this%fail, this%integral, &
    !      this%error, this%prob)
  end subroutine integrate_divonne_userdata

\eatline
\nwendcode{}\nwbegindocs{1013}\nwdocspar
\nwenddocs{}\nwbegincode{1014}\moddef{Muli CUBA: cuba divonne: TBP}\plusendmoddef
  procedure :: copy => cuba_divonne_copy
\nwendcode{}\nwbegincode{1015}\moddef{Muli CUBA: procedures}\plusendmoddef
  subroutine cuba_divonne_copy (this, source)
    class(cuba_divonne_t), intent(out) :: this
    class(cuba_class), intent(in) :: source
    select type (source)
    class is (cuba_divonne_t)
       call this%copy_common(source)
       call this%set_deferred (source%key1, source%key2, source%key3, &
            source%maxpass, source%border, source%maxchisq, &
            source%mindeviation, source%xgiven)
    class default
       call msg_error ("cuba_divonne_copy: type of source is not " &
            // "type compatible with cuba_divonne_t.")
    end select
  end subroutine cuba_divonne_copy
  
\eatline
\nwendcode{}\nwbegindocs{1016}\nwdocspar
\nwenddocs{}\nwbegincode{1017}\moddef{Muli CUBA: cuba divonne: TBP}\plusendmoddef
  procedure :: set_deferred => cuba_divonne_set_deferred  
\nwendcode{}\nwbegincode{1018}\moddef{Muli CUBA: procedures}\plusendmoddef
  subroutine cuba_divonne_set_deferred (this, key1, key2, key3, maxpass, &
       border, maxchisq, mindeviation, xgiven, xgiven_flat)
    class(cuba_divonne_t) :: this
    integer, optional, intent(in) :: key1, key2, key3, maxpass
    real(default), optional, intent(in) :: border, maxchisq, mindeviation
    real(default), dimension(:,:), optional, intent(in) :: xgiven
    real(default), dimension(:), optional, intent(in) :: xgiven_flat
    integer, dimension(2) :: s
    if (present (key1))  this%key1 = key1
    if (present (key2))  this%key2 = key2
    if (present (key3))  this%key3 = key3
    if (present (maxpass))  this%maxpass = maxpass
    if (present (border))  this%border = border
    if (present (maxchisq))  this%maxchisq = maxchisq
    if (present (mindeviation))  this%mindeviation = mindeviation
    if (present (xgiven)) then
       if (allocated (this%xgiven))  deallocate (this%xgiven)
       s = shape(xgiven)
       if (s(1) == this%dim_x) then
          allocate (this%xgiven (s(1), s(2)), source=xgiven)
          this%ldxgiven = s(1)
          this%ngiven = s(2)
       else
          call msg_error ("cuba_divonne_set_deferred: shape of xgiven " &
               // "is not [dim_x,:].")
          this%ngiven = 0
       end if
    end if
    if (present (xgiven_flat)) then
       if (allocated (this%xgiven))  deallocate (this%xgiven)
       if (mod(size(xgiven_flat), this%dim_x) == 0) then
          this%ngiven = size(xgiven_flat) / this%dim_x
          this%ldxgiven = this%dim_x
          allocate (this%xgiven (this%ldxgiven,this%ngiven))
          this%xgiven = reshape(xgiven_flat, [this%ldxgiven, this%ngiven])
       else
          call msg_error ("cuba_divonne_set_deferred: size of xgiven_flat " &
               // "is no multiple of dim_x.")
          this%ngiven = 0
       end if
    end if
  end subroutine cuba_divonne_set_deferred
  
\eatline
\nwendcode{}\nwbegindocs{1019}\nwdocspar
\nwenddocs{}\nwbegincode{1020}\moddef{Muli CUBA: cuba cuhre: TBP}\endmoddef
  procedure :: write_to_marker => cuba_cuhre_write_to_marker   
\nwendcode{}\nwbegincode{1021}\moddef{Muli CUBA: procedures}\plusendmoddef
  subroutine cuba_cuhre_write_to_marker (this, marker, status)
    class(cuba_cuhre_t), intent(in) :: this
    class(marker_t), intent(inout) :: marker
    integer(dik), intent(out) :: status
    call marker%mark_begin ("cuba_cuhre_t")
    call cuba_write_to_marker (this, marker, status)
    call marker%mark ("key", this%key)
    call marker%pick_end ("cuba_cuhre_t", status)
  end subroutine cuba_cuhre_write_to_marker
    
\eatline
\nwendcode{}\nwbegindocs{1022}\nwdocspar
\nwenddocs{}\nwbegincode{1023}\moddef{Muli CUBA: cuba cuhre: TBP}\plusendmoddef
  procedure :: read_from_marker => cuba_cuhre_read_from_marker  
\nwendcode{}\nwbegincode{1024}\moddef{Muli CUBA: procedures}\plusendmoddef
  subroutine cuba_cuhre_read_from_marker (this, marker, status)
    class(cuba_cuhre_t), intent(out) :: this
    class(marker_t), intent(inout) :: marker
    integer(dik), intent(out) :: status
    call marker%pick_begin ("cuba_cuhre_t", status=status)
    call cuba_read_from_marker (this, marker, status)
    call marker%pick ("key",this%key, status)
    call marker%pick_end ("cuba_cuhre_t", status)
  end subroutine cuba_cuhre_read_from_marker
  
\eatline
\nwendcode{}\nwbegindocs{1025}\nwdocspar
\nwenddocs{}\nwbegincode{1026}\moddef{Muli CUBA: cuba cuhre: TBP}\plusendmoddef
  procedure :: print_to_unit => cuba_cuhre_print_to_unit  
\nwendcode{}\nwbegincode{1027}\moddef{Muli CUBA: procedures}\plusendmoddef
  subroutine cuba_cuhre_print_to_unit (this, unit, parents, components, peers)
    class(cuba_cuhre_t), intent(in) :: this
    integer, intent(in) :: unit
    integer(dik), intent(in) :: parents, components, peers
    if (parents > 0)  &
         call cuba_print_to_unit (this, unit, parents-1, components, peers)
    write (unit, "(1x,A)")      "Components of cuba_cuhre_t:"
    write (unit, "(3x,A,I10)")  "key:       ", this%key
  end subroutine cuba_cuhre_print_to_unit  
  
\eatline
\nwendcode{}\nwbegindocs{1028}\nwdocspar
\nwenddocs{}\nwbegincode{1029}\moddef{Muli CUBA: cuba cuhre: TBP}\plusendmoddef
  procedure, nopass :: get_type => cuba_cuhre_get_type  
\nwendcode{}\nwbegincode{1030}\moddef{Muli CUBA: procedures}\plusendmoddef
  pure subroutine cuba_cuhre_get_type (type)
    character(:), allocatable, intent(out) :: type
    allocate (type, source="cuba_cuhre_t")
  end subroutine cuba_cuhre_get_type
  
\eatline
\nwendcode{}\nwbegindocs{1031}\nwdocspar
\nwenddocs{}\nwbegincode{1032}\moddef{Muli CUBA: cuba cuhre: TBP}\plusendmoddef
  procedure :: integrate_nd => integrate_cuhre
\nwendcode{}\nwbegincode{1033}\moddef{Muli CUBA: procedures}\plusendmoddef
  subroutine integrate_cuhre (this, integrand)
    class(cuba_cuhre_t), intent(inout) :: this  
    procedure(integrand_interface) :: integrand
    ! print '("cuhre")'
    ! call cuhre(this%dim_x, this%dim_f, integrand, this%userdata, &
    !      this%eps_rel, this%eps_abs, this%flags, &
    !      ! this%seed, &
    !      this%min_eval, this%max_eval, this%key, this%nregions, &
    !      this%neval, this%fail, this%integral, this%error, this%prob)
  end subroutine integrate_cuhre

\eatline
\nwendcode{}\nwbegindocs{1034} 
\nwenddocs{}\nwbegincode{1035}\moddef{Muli CUBA: cuba cuhre: TBP}\plusendmoddef
  procedure :: integrate_userdata => integrate_cuhre_userdata  
\nwendcode{}\nwbegincode{1036}\moddef{Muli CUBA: procedures}\plusendmoddef
  subroutine integrate_cuhre_userdata (this, integrand, userdata)
    class(cuba_cuhre_t), intent(inout) :: this
    procedure(integrand_interface) :: integrand
    class(transverse_mom_t), intent(in) :: userdata
    ! print '("cuhre")'
    ! call cuhre(this%dim_x, this%dim_f, integrand, userdata, this%eps_rel, &
    !      this%eps_abs, this%flags, &
    !      ! this%seed, &
    !      this%min_eval, this%max_eval, this%key, this%nregions, &
    !      this%neval, this%fail, this%integral, this%error, this%prob)
  end subroutine integrate_cuhre_userdata
  
\eatline
\nwendcode{}\nwbegindocs{1037}\nwdocspar
\nwenddocs{}\nwbegincode{1038}\moddef{Muli CUBA: cuba cuhre: TBP}\plusendmoddef
  procedure :: copy => cuba_cuhre_copy
\nwendcode{}\nwbegincode{1039}\moddef{Muli CUBA: procedures}\plusendmoddef
  subroutine cuba_cuhre_copy (this, source)
    class(cuba_cuhre_t), intent(out) :: this
    class(cuba_class), intent(in) :: source
    select type (source)
    class is (cuba_cuhre_t)
       call this%copy_common (source)
       this%key = source%key
    class default
       call msg_error ("cuba_cuhre_copy: type of source is not type " &
            // "compatible with cuba_cuhre_t.")
    end select
  end subroutine cuba_cuhre_copy

\eatline
\nwendcode{}\nwbegindocs{1040}\nwdocspar
\nwenddocs{}\nwbegincode{1041}\moddef{Muli CUBA: cuba cuhre: TBP}\plusendmoddef
  procedure :: set_deferred => cuba_cuhre_set_deferred 
\nwendcode{}\nwbegincode{1042}\moddef{Muli CUBA: procedures}\plusendmoddef
  subroutine cuba_cuhre_set_deferred (this, key)
    class(cuba_cuhre_t), intent(inout) :: this  
    integer, intent(in) :: key
    this%key = key
  end subroutine cuba_cuhre_set_deferred
  
\eatline
\nwendcode{}\nwbegindocs{1043}\nwdocspar
\nwenddocs{}\nwbegincode{1044}\moddef{Muli CUBA: cuba suave: TBP}\endmoddef
  procedure :: write_to_marker => cuba_suave_write_to_marker
\nwendcode{}\nwbegincode{1045}\moddef{Muli CUBA: procedures}\plusendmoddef
  subroutine cuba_suave_write_to_marker (this, marker, status)
    class(cuba_suave_t), intent(in) :: this
    class(marker_t), intent(inout) :: marker
    integer(dik), intent(out) :: status
    call marker%mark_begin ("cuba_suave_t")
    call cuba_write_to_marker (this, marker, status)
    call marker%mark ("nnew", this%nnew)
    call marker%mark ("flatness", this%flatness)
    call marker%mark_null ("cuba_suave_t")
  end subroutine cuba_suave_write_to_marker
  
\eatline
\nwendcode{}\nwbegindocs{1046}\nwdocspar
\nwenddocs{}\nwbegincode{1047}\moddef{Muli CUBA: cuba suave: TBP}\plusendmoddef
  procedure :: read_from_marker => cuba_suave_read_from_marker    
\nwendcode{}\nwbegincode{1048}\moddef{Muli CUBA: procedures}\plusendmoddef
  subroutine cuba_suave_read_from_marker (this, marker, status)
    class(cuba_suave_t), intent(out) :: this
    class(marker_t), intent(inout) :: marker
    integer(dik), intent(out) :: status
    call marker%pick_begin ("cuba_suave_t", status=status)
    call cuba_read_from_marker (this, marker, status)
    call marker%pick ("nnew", this%nnew, status)
    call marker%pick ("flatnes", this%flatness, status)
    call marker%pick_end ("cuba_suave_t", status)
  end subroutine cuba_suave_read_from_marker
  
\eatline
\nwendcode{}\nwbegindocs{1049}\nwdocspar
\nwenddocs{}\nwbegincode{1050}\moddef{Muli CUBA: cuba suave: TBP}\plusendmoddef
  procedure::print_to_unit=>cuba_suave_print_to_unit  
\nwendcode{}\nwbegincode{1051}\moddef{Muli CUBA: procedures}\plusendmoddef
  subroutine cuba_suave_print_to_unit (this, unit, parents, components, peers)
    class(cuba_suave_t), intent(in) :: this
    integer, intent(in) :: unit
    integer(dik), intent(in) :: parents, components, peers
    if (parents > 0) &
         call cuba_print_to_unit (this, unit, parents-1, components, peers)
    write (unit, "(1x,A)")     "Components of cuba_suave_t:"
    write (unit, "(3x,A,I10)") "nnew:      ", this%nnew
    write (unit, "(3x,A,I10)") "flatness:  ", this%flatness
  end subroutine cuba_suave_print_to_unit
  
\eatline
\nwendcode{}\nwbegindocs{1052}\nwdocspar
\nwenddocs{}\nwbegincode{1053}\moddef{Muli CUBA: cuba suave: TBP}\plusendmoddef
  procedure, nopass :: get_type => cuba_suave_get_type  
\nwendcode{}\nwbegincode{1054}\moddef{Muli CUBA: procedures}\plusendmoddef
  pure subroutine cuba_suave_get_type (type)
    character(:), allocatable, intent(out) :: type
    allocate (type, source="cuba_suave_t")
  end subroutine cuba_suave_get_type
  
\eatline
\nwendcode{}\nwbegindocs{1055}\nwdocspar
\nwenddocs{}\nwbegincode{1056}\moddef{Muli CUBA: cuba suave: TBP}\plusendmoddef
  procedure :: integrate_nd => integrate_suave
\nwendcode{}\nwbegincode{1057}\moddef{Muli CUBA: procedures}\plusendmoddef
  subroutine integrate_suave (this, integrand)
    class(cuba_suave_t), intent(inout) :: this  
    procedure(integrand_interface) :: integrand
    ! print '("suave")'
    ! call suave(this%dim_x, this%dim_f, integrand, this%userdata, &
    !      this%eps_rel, this%eps_abs, this%flags, this%seed, &
    !      this%min_eval, this%max_eval, this%nnew, this%flatness, &
    !      this%nregions, this%neval, this%fail, this%integral, &
    !      this%error, this%prob)
  end subroutine integrate_suave
  
\eatline
\nwendcode{}\nwbegindocs{1058} 
\nwenddocs{}\nwbegincode{1059}\moddef{Muli CUBA: cuba suave: TBP}\plusendmoddef
  procedure :: integrate_userdata => integrate_suave_userdata
\nwendcode{}\nwbegincode{1060}\moddef{Muli CUBA: procedures}\plusendmoddef
   subroutine integrate_suave_userdata (this, integrand, userdata)
    class(cuba_suave_t), intent(inout) :: this  
    procedure(integrand_interface) :: integrand
    class(transverse_mom_t), intent(in) :: userdata
    ! print '("suave")'
    ! call suave (this%dim_x, this%dim_f, integrand, userdata, &
    !      this%eps_rel, this%eps_abs, this%flags, this%seed, &
    !      this%min_eval, this%max_eval, this%nnew, this%flatness, &
    !      this%nregions, this%neval, this%fail, this%integral, &
    !      this%error, this%prob)
  end subroutine integrate_suave_userdata
  
\eatline
\nwendcode{}\nwbegindocs{1061}\nwdocspar
\nwenddocs{}\nwbegincode{1062}\moddef{Muli CUBA: cuba suave: TBP}\plusendmoddef
  procedure :: copy => cuba_suave_copy  
\nwendcode{}\nwbegincode{1063}\moddef{Muli CUBA: procedures}\plusendmoddef
  subroutine cuba_suave_copy (this, source)
    class(cuba_suave_t), intent(out) :: this
    class(cuba_class), intent(in) :: source
    select type (source)
    class is (cuba_suave_t)
       call this%copy_common (source)
       this%nnew = source%nnew
       this%flatness = source%flatness
    class default
       call msg_error ("cuba_suave_copy: type of source is not type " &
            // "compatible with cuba_suave_t.")
    end select
  end subroutine cuba_suave_copy

\eatline
\nwendcode{}\nwbegindocs{1064}\nwdocspar
\nwenddocs{}\nwbegincode{1065}\moddef{Muli CUBA: interfaces}\endmoddef
  interface
     subroutine integrand_interface (dim_x, x, dim_f, f,userdata)
     \LA{}Use kinds\RA{}  
       use muli_momentum
       integer, intent(in) :: dim_x, dim_f
       real(default), dimension(dim_x), intent(in) :: x
       real(default), dimension(dim_f), intent(out) :: f
       class(transverse_mom_t), intent(in) :: userdata
     end subroutine integrand_interface
  end interface
\eatline
\nwendcode{}\nwbegindocs{1066}\nwdocspar
\nwenddocs{}\nwbegincode{1067}\moddef{Muli CUBA: interfaces}\plusendmoddef
  interface
     subroutine cuba_copy_interface (this, source)
       import :: cuba_class
       class(cuba_class), intent(out) :: this
       class(cuba_class), intent(in) :: source
     end subroutine cuba_copy_interface
  end interface
\eatline
\nwendcode{}\nwbegindocs{1068}\nwdocspar
\nwenddocs{}\nwbegincode{1069}\moddef{Muli CUBA: interfaces}\plusendmoddef
  interface
     subroutine ca_plain (this)
       import :: cuba_class
       class(cuba_class) :: this 
     end subroutine ca_plain
  end interface
\eatline
\nwendcode{}\nwbegindocs{1070}\nwdocspar
\nwenddocs{}\nwbegincode{1071}\moddef{Muli CUBA: interfaces}\plusendmoddef
  interface 
     subroutine integrate_interface (this, integrand)
       import :: cuba_class
       class(cuba_class), intent(inout) :: this  
       interface
          subroutine integrand (dim_x, x, dim_f, f,userdata)
          \LA{}Use kinds\RA{}  
            use muli_momentum
            integer, intent(in) :: dim_x, dim_f
            real(default), dimension(dim_x), intent(in) :: x
            real(default), dimension(dim_f), intent(out) :: f
            class(transverse_mom_t), intent(in) :: userdata
          end subroutine integrand
       end interface
     end subroutine integrate_interface
  end interface
\eatline
\nwendcode{}\nwbegindocs{1072}\nwdocspar
\nwenddocs{}\nwbegincode{1073}\moddef{Muli CUBA: interfaces}\plusendmoddef
  interface
     subroutine integrate_userdata_interface (this, integrand,userdata)
       use muli_momentum
       import :: cuba_class
       class(cuba_class), intent(inout) :: this  
       interface
          subroutine integrand (dim_x, x, dim_f, f,userdata)
          \LA{}Use kinds\RA{}
            use muli_momentum
            integer, intent(in) :: dim_x, dim_f
            real(default), dimension(dim_x), intent(in) :: x
            real(default), dimension(dim_f), intent(out) :: f
            class(transverse_mom_t), intent(in) :: userdata
          end subroutine integrand
       end interface
       class(transverse_mom_t), intent(in) :: userdata
     end subroutine integrate_userdata_interface
  end interface
  
\eatline
\nwendcode{}\nwbegindocs{1074} 
%%%%%%%%%%%%%%%%%%%%%%%%%%%%%%%%%%%%%%%%%%%%%%%%%%%%%%%%%%%%%%%%%%%%%%%%
\section{Trapezoidal integration routines}

This file contains the module {\tt{}"muli{\char95}trapezium"}. The name is
derived from the trapezoidal integration rule. The purpose of this
module is to define a binary tree {\tt{}muli{\char95}trapezium{\char95}tree{\char95}t} which
holds a probability function in terms of trapezoidal segments. Its
leaves of type {\tt{}muli{\char95}trapezium{\char95}list{\char95}t} are connected to form a list,
so you can either walk the root function back and forth or pick a
certain segment in logarithmic time by walking down the tree.
All nodes extend {\tt{}muli{\char95}trapezium{\char95}t} which holds the actual
values. Those are the values of the density function, the integral
from this segment to the end of the integration area and a probability
function calculated from these values. All values are taken at the
upper "right" bound of the segment. Additionally the differences of
these values to the values of the left neighbor is stored. 

\nwenddocs{}\nwbegincode{1075}\moddef{\code{}muli{\char95}trapezium.f90\edoc{}}\endmoddef
\LA{}File header\RA{}

module muli_trapezium
  use, intrinsic :: iso_fortran_env, only: output_unit
\LA{}Use kinds\RA{}
  use constants
  use diagnostics
  use muli_base

\LA{}Standard module head\RA{}  
  
\LA{}Muli trapezium: variables\RA{}

\LA{}Muli trapezium: public\RA{}

\LA{}Muli trapezium: types\RA{}

\LA{}Muli trapezium: interfaces\RA{}

contains
  
\LA{}Muli trapezium: procedures\RA{}  
  
end module muli_trapezium
\eatline
\nwendcode{}\nwbegindocs{1076}\nwdocspar
\nwenddocs{}\nwbegincode{1077}\moddef{Muli trapezium: variables}\endmoddef
  integer, private, parameter :: value_dimension = 7
  integer, private, parameter :: r_value_index = 1
  integer, private, parameter :: d_value_index = 2
  integer, private, parameter :: r_integral_index = 3
  integer, private, parameter :: d_integral_index = 4
  integer, private, parameter :: r_probability_index = 5
  integer, private, parameter :: d_probability_index = 6
  integer, private, parameter :: error_index = 7  

\eatline
\eatline
\eatline
\nwendcode{}\nwbegindocs{1078}\nwdocspar
\nwenddocs{}\nwbegincode{1079}\moddef{Muli trapezium: interfaces}\endmoddef
  abstract interface
     subroutine muli_trapezium_append_interface (this, right)
       import muli_trapezium_node_class_t
       class(muli_trapezium_node_class_t), intent(inout), target :: this, right
     end subroutine muli_trapezium_append_interface
  end interface  
\eatline
\nwendcode{}\nwbegindocs{1080}\nwdocspar
\nwenddocs{}\nwbegincode{1081}\moddef{Muli trapezium: interfaces}\plusendmoddef
  abstract interface
     subroutine muli_trapezium_final_interface (this)
       import muli_trapezium_node_class_t
       class(muli_trapezium_node_class_t), intent(inout) :: this
     end subroutine muli_trapezium_final_interface
  end interface
  
\eatline
\nwendcode{}\nwbegindocs{1082}This is the base type {\tt{}muli{\char95}trapezium{\char95}t}. Its component {\tt{}values}
has a first index is in $\left\{ 0, \ldots, \text{dim}-1 \right\}$, while the
second index distingiushes between {\tt{}r{\char95}value}, {\tt{}d{\char95}value},
{\tt{}r{\char95}integral}, {\tt{}d{\char95}integral}, {\tt{}r{\char95}probability}, {\tt{}d{\char95}probability}.
\nwenddocs{}\nwbegincode{1083}\moddef{Muli trapezium: public}\endmoddef
  public :: muli_trapezium_t
\nwendcode{}\nwbegincode{1084}\moddef{Muli trapezium: types}\endmoddef
  type, extends (measure_class_t) :: muli_trapezium_t
     private
     integer::dim=0
     real(default) :: r_position = 0
     real(default) :: d_position = 0
     real(default) :: measure_comp = 0
     real(default), dimension(:,:), allocatable :: values
   contains
   \LA{}Muli trapezium: trapezium: TBP\RA{}  
  end type muli_trapezium_t  
  
\eatline
\nwendcode{}\nwbegindocs{1085}\nwdocspar
\nwenddocs{}\nwbegincode{1086}\moddef{Muli trapezium: public}\plusendmoddef
  public :: muli_trapezium_node_class_t 
\nwendcode{}\nwbegincode{1087}\moddef{Muli trapezium: types}\plusendmoddef
  type, extends (muli_trapezium_t), abstract :: muli_trapezium_node_class_t
     private
     class(muli_trapezium_node_class_t), pointer :: left => null()
     class(muli_trapezium_node_class_t), pointer :: right => null()
     ! real(default) :: criterion
   contains
   \LA{}Muli trapezium: node: TBP\RA{}
  end type muli_trapezium_node_class_t
  
\eatline
\nwendcode{}\nwbegindocs{1088}\nwdocspar
\nwenddocs{}\nwbegincode{1089}\moddef{Muli trapezium: public}\plusendmoddef
  public :: muli_trapezium_tree_t
\nwendcode{}\nwbegincode{1090}\moddef{Muli trapezium: types}\plusendmoddef
  type, extends(muli_trapezium_node_class_t) :: muli_trapezium_tree_t
     class(muli_trapezium_node_class_t), pointer :: down => null()
   contains
   \LA{}Muli trapezium: tree: TBP\RA{}
  end type muli_trapezium_tree_t
  
\eatline
\nwendcode{}\nwbegindocs{1091}\nwdocspar
\nwenddocs{}\nwbegincode{1092}\moddef{Muli trapezium: public}\plusendmoddef
  public :: muli_trapezium_list_t
\nwendcode{}\nwbegincode{1093}\moddef{Muli trapezium: types}\plusendmoddef
  type, extends (muli_trapezium_node_class_t) :: muli_trapezium_list_t
   contains
   \LA{}Muli trapezium: list: TBP\RA{}
  end type muli_trapezium_list_t
  
\eatline
\nwendcode{}\nwbegindocs{1094}\nwdocspar
\nwenddocs{}\nwbegincode{1095}\moddef{Muli trapezium: trapezium: TBP}\endmoddef
  procedure :: write_to_marker => muli_trapezium_write_to_marker
\nwendcode{}\nwbegincode{1096}\moddef{Muli trapezium: procedures}\endmoddef
    subroutine muli_trapezium_write_to_marker (this,marker,status)
    class(muli_trapezium_t), intent(in) :: this
    class(marker_t), intent(inout) :: marker
    integer(dik), intent(out) :: status
    integer::dim
    call marker%mark_begin ("muli_trapezium_t")
    call marker%mark ("dim", this%dim)
    call marker%mark ("r_position", this%r_position)
    call marker%mark ("d_position", this%d_position)
    if (allocated(this%values)) then
       call marker%mark ("values", this%values)
    else
       call marker%mark_null ("values")
    end if
    call marker%mark_end ("muli_trapezium_t")
  end subroutine muli_trapezium_write_to_marker
  
\eatline
\nwendcode{}\nwbegindocs{1097}\nwdocspar
\nwenddocs{}\nwbegincode{1098}\moddef{Muli trapezium: trapezium: TBP}\plusendmoddef
  procedure :: read_from_marker => muli_trapezium_read_from_marker  
\nwendcode{}\nwbegincode{1099}\moddef{Muli trapezium: procedures}\plusendmoddef
  subroutine muli_trapezium_read_from_marker (this,marker,status)
    class(muli_trapezium_t), intent(out) :: this
    class(marker_t), intent(inout) :: marker
    integer(dik), intent(out) :: status
    integer :: dim
    call marker%pick_begin ("muli_trapezium_t", status=status)
    call marker%pick ("dim", this%dim,status)
    call marker%pick ("r_position", this%r_position, status)
    call marker%pick ("d_position", this%d_position, status)
    if (allocated (this%values))  deallocate (this%values)
    call marker%verify_nothing ("values", status)
    if (status == serialize_ok) then
       allocate(this%values(0:this%dim-1,7))
       call marker%pick ("values",this%values, status)
    end if
    call marker%pick_end("muli_trapezium_t",status)
  end subroutine muli_trapezium_read_from_marker
  
\eatline
\nwendcode{}\nwbegindocs{1100}\nwdocspar
\nwenddocs{}\nwbegincode{1101}\moddef{Muli trapezium: trapezium: TBP}\plusendmoddef
  procedure :: print_to_unit => muli_trapezium_print_to_unit  
\nwendcode{}\nwbegincode{1102}\moddef{Muli trapezium: procedures}\plusendmoddef
  subroutine muli_trapezium_print_to_unit (this, unit, parents, components, peers)
    class(muli_trapezium_t), intent(in) :: this
    integer, intent(in) :: unit
    integer(dik), intent(in) :: parents, components, peers
    write (unit, "(1x,A)") "Components of muli_trapezium_t:"
    write (unit, fmt=*)"Dimension:        ",this%dim
    write (unit,fmt=*)"Right position:   ",this%r_position
    write (unit,fmt=*)"Position step:    ",this%d_position
    if (allocated(this%values)) then    
       if (components>0) then          
          write (unit,fmt=*)"Right values:     ",muli_trapezium_get_r_value_array(this)
          write (unit,fmt=*) "Value step:       ", this%get_d_value()
          write (unit,fmt=*)"Right integrals:  ",this%get_r_integral()
          write (unit,fmt=*)"Integral step:    ",this%get_d_integral()
          write (unit,fmt=*)"Right propabilities:",this%get_r_probability()
          write (unit,fmt=*)"Probability step: ",this%get_d_probability()
          write (unit,fmt=*)"Errors:           ",this%get_error()
       else
          write (unit, "(3x,A)") "Values are allocated."
       end if
    else
       write (unit, "(3x,A)") "Values are not allocated."
    end if
  end subroutine muli_trapezium_print_to_unit  

\eatline
\nwendcode{}\nwbegindocs{1103}\nwdocspar
\nwenddocs{}\nwbegincode{1104}\moddef{Muli trapezium: trapezium: TBP}\plusendmoddef
  procedure, nopass :: get_type => muli_trapezium_get_type  
\nwendcode{}\nwbegincode{1105}\moddef{Muli trapezium: procedures}\plusendmoddef
  pure subroutine muli_trapezium_get_type (type)
    character(:),allocatable, intent(out) :: type
    allocate (type, source="muli_trapezium_t")
  end subroutine muli_trapezium_get_type
  
\eatline
\nwendcode{}\nwbegindocs{1106}\nwdocspar
\nwenddocs{}\nwbegincode{1107}\moddef{Muli trapezium: trapezium: TBP}\plusendmoddef
  procedure, nopass :: verify_type => muli_trapezium_verify_type
\nwendcode{}\nwbegincode{1108}\moddef{Muli trapezium: procedures}\plusendmoddef
  elemental logical function muli_trapezium_verify_type (type) result (match)
    character(*), intent(in) :: type
    match = type == "muli_trapezium_t"
  end function muli_trapezium_verify_type
  
\eatline
\nwendcode{}\nwbegindocs{1109}\nwdocspar
\nwenddocs{}\nwbegincode{1110}\moddef{Muli trapezium: trapezium: TBP}\plusendmoddef
  procedure :: measure => muli_trapezium_measure
\nwendcode{}\nwbegincode{1111}\moddef{Muli trapezium: procedures}\plusendmoddef
  elemental function muli_trapezium_measure (this)
    class(muli_trapezium_t), intent(in) :: this
    real(default) :: muli_trapezium_measure
    muli_trapezium_measure = this%measure_comp
  end function muli_trapezium_measure
  
\eatline
\nwendcode{}\nwbegindocs{1112}\nwdocspar
\nwenddocs{}\nwbegincode{1113}\moddef{Muli trapezium: trapezium: TBP}\plusendmoddef
  procedure :: initialize => muli_trapezium_initialize  
\nwendcode{}\nwbegincode{1114}\moddef{Muli trapezium: procedures}\plusendmoddef
  subroutine muli_trapezium_initialize (this, dim, r_position, d_position)
    class(muli_trapezium_t), intent(inout) :: this
    integer, intent(in) :: dim
    real(default), intent(in) :: r_position, d_position
    integer :: dim1, dim2
    this%dim = dim
    this%r_position = r_position
    this%d_position = d_position
    if (allocated (this%values))  deallocate (this%values)
    allocate (this%values(0:dim-1,value_dimension))
    do dim2 = 1, value_dimension-1
       do dim1 = 0, dim-1
          this%values(dim1,dim2) = zero
       end do       
    end do
    do dim1 = 0, dim-1
       this%values(dim1, value_dimension) = huge(one)
    end do
    this%measure_comp = huge(one)
  end subroutine muli_trapezium_initialize

\eatline
\nwendcode{}\nwbegindocs{1115}\nwdocspar
\nwenddocs{}\nwbegincode{1116}\moddef{Muli trapezium: trapezium: TBP}\plusendmoddef
  procedure :: get_dimension => muli_trapezium_get_dimension  
\nwendcode{}\nwbegincode{1117}\moddef{Muli trapezium: procedures}\plusendmoddef
  elemental function muli_trapezium_get_dimension (this) result (dim)
    class(muli_trapezium_t), intent(in) :: this
    integer :: dim
    dim = this%dim
  end function muli_trapezium_get_dimension
  
\eatline
\nwendcode{}\nwbegindocs{1118}\nwdocspar
\nwenddocs{}\nwbegincode{1119}\moddef{Muli trapezium: trapezium: TBP}\plusendmoddef
  procedure :: get_l_position => muli_trapezium_get_l_position
\nwendcode{}\nwbegincode{1120}\moddef{Muli trapezium: procedures}\plusendmoddef
  pure function muli_trapezium_get_l_position (this) result (pos)
    class(muli_trapezium_t), intent(in) :: this
    real(default) :: pos
    pos = this%r_position - this%d_position
  end function muli_trapezium_get_l_position
  
\eatline
\nwendcode{}\nwbegindocs{1121}\nwdocspar
\nwenddocs{}\nwbegincode{1122}\moddef{Muli trapezium: trapezium: TBP}\plusendmoddef
  procedure :: get_r_position => muli_trapezium_get_r_position  
\nwendcode{}\nwbegincode{1123}\moddef{Muli trapezium: procedures}\plusendmoddef
  pure function muli_trapezium_get_r_position (this) result (pos)
    class(muli_trapezium_t), intent(in) :: this
    real(default) :: pos
    pos = this%r_position
  end function muli_trapezium_get_r_position
  
\eatline
\nwendcode{}\nwbegindocs{1124}\nwdocspar
\nwenddocs{}\nwbegincode{1125}\moddef{Muli trapezium: trapezium: TBP}\plusendmoddef
  procedure :: get_d_position => muli_trapezium_get_d_position  
\nwendcode{}\nwbegincode{1126}\moddef{Muli trapezium: procedures}\plusendmoddef
  pure function muli_trapezium_get_d_position (this) result (pos)
    class(muli_trapezium_t), intent(in) :: this
    real(default) :: pos
    pos = this%d_position
  end function muli_trapezium_get_d_position
  
\eatline
\nwendcode{}\nwbegindocs{1127}\nwdocspar
\nwenddocs{}\nwbegincode{1128}\moddef{Muli trapezium: trapezium: TBP}\plusendmoddef
  generic :: get_l_value => get_l_value_array, get_l_value_element
  procedure :: get_l_value_array => muli_trapezium_get_l_value_array
  procedure :: get_l_value_element => muli_trapezium_get_l_value_element    
\nwendcode{}\nwbegincode{1129}\moddef{Muli trapezium: procedures}\plusendmoddef
  pure function muli_trapezium_get_l_value_array (this) result (subarray)
    class(muli_trapezium_t), intent(in) :: this
    real(default), dimension(this%dim) :: subarray
    subarray = this%values(0:this%dim-1, r_value_index) - &
         this%values(0:this%dim-1, d_value_index)
  end function muli_trapezium_get_l_value_array
  
\eatline
\nwendcode{}\nwbegindocs{1130}\nwdocspar
\nwenddocs{}\nwbegincode{1131}\moddef{Muli trapezium: procedures}\plusendmoddef
  pure function muli_trapezium_get_l_value_element (this, set) result (element)
    class(muli_trapezium_t), intent(in) :: this
    integer, intent(in) :: set
    real(default) :: element
    element = this%values(set, r_value_index) - this%values(set, d_value_index)
  end function muli_trapezium_get_l_value_element
  
\eatline
\nwendcode{}\nwbegindocs{1132}\nwdocspar
\nwenddocs{}\nwbegincode{1133}\moddef{Muli trapezium: trapezium: TBP}\plusendmoddef
  generic :: get_r_value => get_r_value_array, get_r_value_element  
  procedure :: get_r_value_array => muli_trapezium_get_r_value_array
  procedure :: get_r_value_element => muli_trapezium_get_r_value_element  
\nwendcode{}\nwbegindocs{1134}\nwdocspar
\nwenddocs{}\nwbegincode{1135}\moddef{Muli trapezium: procedures}\plusendmoddef
  pure function muli_trapezium_get_r_value_element (this, set) result (element)
    class(muli_trapezium_t), intent(in) :: this
    integer, intent(in) :: set
    real(default) :: element
    element = this%values (set, r_value_index)
  end function muli_trapezium_get_r_value_element

\eatline
\nwendcode{}\nwbegindocs{1136}\nwdocspar
\nwenddocs{}\nwbegincode{1137}\moddef{Muli trapezium: procedures}\plusendmoddef
  pure function muli_trapezium_get_r_value_array (this) result (subarray)
    class(muli_trapezium_t), intent(in) :: this
    real(default), dimension(this%dim) :: subarray
    subarray = this%values(0:this%dim-1, r_value_index)
  end function muli_trapezium_get_r_value_array
  
\eatline
\nwendcode{}\nwbegindocs{1138}\nwdocspar
\nwenddocs{}\nwbegincode{1139}\moddef{Muli trapezium: trapezium: TBP}\plusendmoddef
  generic :: get_d_value => get_d_value_array, get_d_value_element  
  procedure :: get_d_value_array => muli_trapezium_get_d_value_array
  procedure :: get_d_value_element => muli_trapezium_get_d_value_element  
\nwendcode{}\nwbegincode{1140}\moddef{Muli trapezium: procedures}\plusendmoddef
  pure function muli_trapezium_get_d_value_element (this, set) result (element)
    class(muli_trapezium_t), intent(in) :: this
    integer, intent(in) :: set
    real(default) :: element
    element=this%values (set, d_value_index)
  end function muli_trapezium_get_d_value_element

\eatline
\nwendcode{}\nwbegindocs{1141}\nwdocspar
\nwenddocs{}\nwbegincode{1142}\moddef{Muli trapezium: procedures}\plusendmoddef
  pure function muli_trapezium_get_d_value_array (this) result (subarray)
    class(muli_trapezium_t), intent(in) :: this
    real(default), dimension(this%dim) :: subarray
    subarray = this%values(0:this%dim-1, d_value_index)
  end function muli_trapezium_get_d_value_array
  
\eatline
\nwendcode{}\nwbegindocs{1143}\nwdocspar
\nwenddocs{}\nwbegincode{1144}\moddef{Muli trapezium: trapezium: TBP}\plusendmoddef
  generic :: get_l_integral => get_l_integral_array, get_l_integral_element  
  procedure :: get_l_integral_array => muli_trapezium_get_l_integral_array
  procedure :: get_l_integral_element => muli_trapezium_get_l_integral_element  
\nwendcode{}\nwbegincode{1145}\moddef{Muli trapezium: procedures}\plusendmoddef
  pure function muli_trapezium_get_l_integral_element &
       (this, set) result (element)
    class(muli_trapezium_t), intent(in) :: this
    integer, intent(in) :: set
    real(default) :: element
    element = this%values (set, r_integral_index) - &
         this%values (set, d_integral_index)
  end function muli_trapezium_get_l_integral_element

\eatline
\nwendcode{}\nwbegindocs{1146}\nwdocspar
\nwenddocs{}\nwbegincode{1147}\moddef{Muli trapezium: procedures}\plusendmoddef
  pure function muli_trapezium_get_l_integral_array (this) result (subarray)
    class(muli_trapezium_t), intent(in) :: this
    real(default), dimension(this%dim) :: subarray
    subarray = this%values (0:this%dim-1, r_integral_index) - &
         this%values (0:this%dim-1, d_integral_index)
  end function muli_trapezium_get_l_integral_array

\eatline
\nwendcode{}\nwbegindocs{1148}\nwdocspar
\nwenddocs{}\nwbegincode{1149}\moddef{Muli trapezium: trapezium: TBP}\plusendmoddef
  generic :: get_r_integral => get_r_integral_array, get_r_integral_element  
  procedure :: get_r_integral_array => muli_trapezium_get_r_integral_array
  procedure :: get_r_integral_element => muli_trapezium_get_r_integral_element   
\nwendcode{}\nwbegincode{1150}\moddef{Muli trapezium: procedures}\plusendmoddef
  pure function muli_trapezium_get_r_integral_element (this, set) result (element)
    class(muli_trapezium_t), intent(in) :: this
    integer, intent(in) :: set
    real(default) :: element
    element = this%values (set, r_integral_index)
  end function muli_trapezium_get_r_integral_element

\eatline
\nwendcode{}\nwbegindocs{1151}\nwdocspar
\nwenddocs{}\nwbegincode{1152}\moddef{Muli trapezium: procedures}\plusendmoddef
  pure function muli_trapezium_get_r_integral_array (this) result (subarray)
    class(muli_trapezium_t), intent(in) :: this
    real(default), dimension(this%dim) :: subarray
    subarray = this%values (0:this%dim-1, r_integral_index)
  end function muli_trapezium_get_r_integral_array
  
\eatline
\nwendcode{}\nwbegindocs{1153}\nwdocspar
\nwenddocs{}\nwbegincode{1154}\moddef{Muli trapezium: trapezium: TBP}\plusendmoddef
  generic :: get_d_integral => get_d_integral_array, get_d_integral_element  
  procedure :: get_d_integral_array => muli_trapezium_get_d_integral_array
  procedure :: get_d_integral_element => muli_trapezium_get_d_integral_element
\nwendcode{}\nwbegincode{1155}\moddef{Muli trapezium: procedures}\plusendmoddef
  pure function muli_trapezium_get_d_integral_element & 
       (this, set) result (element)
    class(muli_trapezium_t), intent(in) :: this
    integer, intent(in) :: set
    real(default) :: element
    element = this%values (set, d_integral_index)
  end function muli_trapezium_get_d_integral_element

\eatline
\nwendcode{}\nwbegindocs{1156}\nwdocspar
\nwenddocs{}\nwbegincode{1157}\moddef{Muli trapezium: procedures}\plusendmoddef
  pure function muli_trapezium_get_d_integral_array (this) result (subarray)
    class(muli_trapezium_t), intent(in) :: this
    real(default), dimension(this%dim) :: subarray
    subarray = this%values (0:this%dim-1, d_integral_index)
  end function muli_trapezium_get_d_integral_array
  
\eatline
\nwendcode{}\nwbegindocs{1158}\nwdocspar
\nwenddocs{}\nwbegincode{1159}\moddef{Muli trapezium: trapezium: TBP}\plusendmoddef
  generic :: get_l_probability => &
       get_l_probability_array, get_l_probability_element  
  procedure :: get_l_probability_element => &
       muli_trapezium_get_l_probability_element
  procedure :: get_l_probability_array => &
       muli_trapezium_get_l_probability_array  
\nwendcode{}\nwbegindocs{1160}\nwdocspar
\nwenddocs{}\nwbegincode{1161}\moddef{Muli trapezium: procedures}\plusendmoddef
  pure function muli_trapezium_get_l_probability_element &
       (this, set) result (element)
    class(muli_trapezium_t), intent(in) :: this
    integer, intent(in) :: set
    real(default) :: element
    element = this%values (set, r_probability_index) - &
         this%values (set, d_probability_index)
  end function muli_trapezium_get_l_probability_element

\eatline
\nwendcode{}\nwbegindocs{1162}\nwdocspar
\nwenddocs{}\nwbegincode{1163}\moddef{Muli trapezium: procedures}\plusendmoddef
  pure function muli_trapezium_get_l_probability_array (this) result (subarray)
    class(muli_trapezium_t), intent(in) :: this
    real(default), dimension(this%dim) :: subarray
    subarray = this%values (0:this%dim-1, r_probability_index) - &
         this%values (0:this%dim-1, d_probability_index)
  end function muli_trapezium_get_l_probability_array

\eatline
\nwendcode{}\nwbegindocs{1164}\nwdocspar
\nwenddocs{}\nwbegincode{1165}\moddef{Muli trapezium: trapezium: TBP}\plusendmoddef
  generic :: get_r_probability => &
       get_r_probability_array, get_r_probability_element  
  procedure :: get_r_probability_element => &
       muli_trapezium_get_r_probability_element
  procedure :: get_r_probability_array => &
       muli_trapezium_get_r_probability_array  
\nwendcode{}\nwbegindocs{1166}\nwdocspar
\nwenddocs{}\nwbegincode{1167}\moddef{Muli trapezium: procedures}\plusendmoddef
  pure function muli_trapezium_get_r_probability_element &
       (this, set) result (element)
    class(muli_trapezium_t), intent(in) :: this
    integer, intent(in) :: set
    real(default) :: element
    element = this%values (set, r_probability_index)
  end function muli_trapezium_get_r_probability_element

\eatline
\nwendcode{}\nwbegindocs{1168}\nwdocspar
\nwenddocs{}\nwbegincode{1169}\moddef{Muli trapezium: procedures}\plusendmoddef
  pure function muli_trapezium_get_r_probability_array (this) result (subarray)
    class(muli_trapezium_t), intent(in) :: this
    real(default), dimension(this%dim) :: subarray
    subarray = this%values (0:this%dim-1, r_probability_index)
  end function muli_trapezium_get_r_probability_array
  
\eatline
\nwendcode{}\nwbegindocs{1170}\nwdocspar
\nwenddocs{}\nwbegincode{1171}\moddef{Muli trapezium: trapezium: TBP}\plusendmoddef
  generic :: get_d_probability => &
       get_d_probability_array, get_d_probability_element  
  procedure :: get_d_probability_element => &
       muli_trapezium_get_d_probability_element
  procedure :: get_d_probability_array => &
       muli_trapezium_get_d_probability_array  
\nwendcode{}\nwbegincode{1172}\moddef{Muli trapezium: procedures}\plusendmoddef
  pure function muli_trapezium_get_d_probability_array (this) result (subarray)
    class(muli_trapezium_t), intent(in) :: this
    real(default), dimension(this%dim) :: subarray
    subarray = this%values (0:this%dim-1, d_probability_index)
  end function muli_trapezium_get_d_probability_array

\eatline
\nwendcode{}\nwbegindocs{1173}\nwdocspar
\nwenddocs{}\nwbegincode{1174}\moddef{Muli trapezium: procedures}\plusendmoddef
  pure function muli_trapezium_get_d_probability_element &
       (this, set) result (element)
    class(muli_trapezium_t), intent(in) :: this
    integer, intent(in) :: set
    real(default) :: element
    element = this%values (set, d_probability_index)
  end function muli_trapezium_get_d_probability_element

\eatline
\nwendcode{}\nwbegindocs{1175}\nwdocspar
\nwenddocs{}\nwbegincode{1176}\moddef{Muli trapezium: trapezium: TBP}\plusendmoddef
  procedure :: get_error => muli_trapezium_get_error
\nwendcode{}\nwbegincode{1177}\moddef{Muli trapezium: procedures}\plusendmoddef
  pure function muli_trapezium_get_error_sum (this) result (error)
    class(muli_trapezium_t), intent(in) :: this
    real(default) :: error
    error = sum (this%values (0:this%dim-1, error_index))
  end function muli_trapezium_get_error_sum

\eatline
\nwendcode{}\nwbegindocs{1178}\nwdocspar
\nwenddocs{}\nwbegincode{1179}\moddef{Muli trapezium: trapezium: TBP}\plusendmoddef
  procedure :: get_error_sum => muli_trapezium_get_error_sum
\nwendcode{}\nwbegincode{1180}\moddef{Muli trapezium: procedures}\plusendmoddef
  pure function muli_trapezium_get_error (this) result (error)
    class(muli_trapezium_t), intent(in) :: this
    real(default), dimension(this%dim) :: error
    error = this%values (0:this%dim-1, error_index)
  end function muli_trapezium_get_error
  
\eatline
\nwendcode{}\nwbegindocs{1181}\nwdocspar
\nwenddocs{}\nwbegincode{1182}\moddef{Muli trapezium: trapezium: TBP}\plusendmoddef
  procedure :: get_integral_sum => muli_trapezium_get_integral_sum
\nwendcode{}\nwbegincode{1183}\moddef{Muli trapezium: procedures}\plusendmoddef
  pure function muli_trapezium_get_integral_sum (this) result (error)
    class(muli_trapezium_t), intent(in) :: this
    real(default) :: error
    error = sum (this%values (0:this%dim-1, d_integral_index))
  end function muli_trapezium_get_integral_sum
  
\eatline
\nwendcode{}\nwbegindocs{1184}\nwdocspar
\nwenddocs{}\nwbegincode{1185}\moddef{Muli trapezium: trapezium: TBP}\plusendmoddef
  procedure :: get_value_at_position => muli_trapezium_get_value_at_position  
\nwendcode{}\nwbegincode{1186}\moddef{Muli trapezium: procedures}\plusendmoddef
  subroutine muli_trapezium_get_value_at_position (this, pos, subarray)
    class(muli_trapezium_t), intent(in) :: this
    real(default), intent(in) :: pos
    real(default), dimension(this%dim), intent(out) :: subarray
    subarray = this%get_r_value_array() - this%get_d_value() * &
         this%d_position / (this%r_position-pos)
  end subroutine muli_trapezium_get_value_at_position

\eatline
\nwendcode{}\nwbegindocs{1187}\nwdocspar
\nwenddocs{}\nwbegincode{1188}\moddef{Muli trapezium: trapezium: TBP}\plusendmoddef
  procedure :: set_r_value => muli_trapezium_set_r_value
\nwendcode{}\nwbegincode{1189}\moddef{Muli trapezium: procedures}\plusendmoddef
  subroutine muli_trapezium_set_r_value (this, subarray)
    class(muli_trapezium_t), intent(inout) :: this
    real(default), intent(in), dimension(0:this%dim-1) :: subarray
    this%values(0:this%dim-1, r_value_index) = subarray
  end subroutine muli_trapezium_set_r_value

\eatline
\nwendcode{}\nwbegindocs{1190}\nwdocspar
\nwenddocs{}\nwbegincode{1191}\moddef{Muli trapezium: trapezium: TBP}\plusendmoddef
  procedure :: set_d_value => muli_trapezium_set_d_value
\nwendcode{}\nwbegincode{1192}\moddef{Muli trapezium: procedures}\plusendmoddef
  subroutine muli_trapezium_set_d_value (this, subarray)
    class(muli_trapezium_t), intent(inout) :: this
    real(default), intent(in), dimension(0:this%dim-1) :: subarray
    this%values(0:this%dim-1,d_value_index) = subarray
  end subroutine muli_trapezium_set_d_value

\eatline
\nwendcode{}\nwbegindocs{1193}\nwdocspar
\nwenddocs{}\nwbegincode{1194}\moddef{Muli trapezium: trapezium: TBP}\plusendmoddef
  procedure :: set_r_integral => muli_trapezium_set_r_integral
\nwendcode{}\nwbegincode{1195}\moddef{Muli trapezium: procedures}\plusendmoddef
  subroutine muli_trapezium_set_r_integral (this, subarray)
    class(muli_trapezium_t), intent(inout) :: this
    real(default), intent(in), dimension(0:this%dim-1) :: subarray
    this%values(0:this%dim-1,r_integral_index) = subarray
  end subroutine muli_trapezium_set_r_integral

\eatline
\nwendcode{}\nwbegindocs{1196}\nwdocspar
\nwenddocs{}\nwbegincode{1197}\moddef{Muli trapezium: trapezium: TBP}\plusendmoddef
  procedure :: set_d_integral => muli_trapezium_set_d_integral
\nwendcode{}\nwbegincode{1198}\moddef{Muli trapezium: procedures}\plusendmoddef
  subroutine muli_trapezium_set_d_integral (this, subarray)
    class(muli_trapezium_t), intent(inout) :: this
    real(default), intent(in), dimension(0:this%dim-1) :: subarray
    this%values (0:this%dim-1, d_integral_index) = subarray
  end subroutine muli_trapezium_set_d_integral

\eatline
\nwendcode{}\nwbegindocs{1199}\nwdocspar
\nwenddocs{}\nwbegincode{1200}\moddef{Muli trapezium: trapezium: TBP}\plusendmoddef
  procedure :: set_r_probability => muli_trapezium_set_r_probability
\nwendcode{}\nwbegincode{1201}\moddef{Muli trapezium: procedures}\plusendmoddef
  subroutine muli_trapezium_set_r_probability (this, subarray)
    class(muli_trapezium_t), intent(inout) :: this
    real(default), intent(in), dimension(0:this%dim-1) :: subarray
    this%values (0:this%dim-1,r_probability_index) = subarray
  end subroutine muli_trapezium_set_r_probability
  
\eatline
\nwendcode{}\nwbegindocs{1202}\nwdocspar
\nwenddocs{}\nwbegincode{1203}\moddef{Muli trapezium: trapezium: TBP}\plusendmoddef
  procedure :: set_d_probability => muli_trapezium_set_d_probability
\nwendcode{}\nwbegincode{1204}\moddef{Muli trapezium: procedures}\plusendmoddef
  subroutine muli_trapezium_set_d_probability (this, subarray)
    class(muli_trapezium_t), intent(inout) :: this
    real(default), intent(in), dimension(0:this%dim-1) :: subarray
    this%values (0:this%dim-1,d_probability_index) = subarray
  end subroutine muli_trapezium_set_d_probability
  
\eatline
\nwendcode{}\nwbegindocs{1205}\nwdocspar
\nwenddocs{}\nwbegincode{1206}\moddef{Muli trapezium: trapezium: TBP}\plusendmoddef
  procedure :: set_error => muli_trapezium_set_error  
\nwendcode{}\nwbegincode{1207}\moddef{Muli trapezium: procedures}\plusendmoddef
  subroutine muli_trapezium_set_error (this, subarray)
    class(muli_trapezium_t), intent(inout) :: this
    real(default), intent(in), dimension(0:this%dim-1) :: subarray
    this%values (0:this%dim-1, error_index) = subarray
    this%measure_comp = sum (subarray)
  end subroutine muli_trapezium_set_error
  
\eatline
\nwendcode{}\nwbegindocs{1208}\nwdocspar
\nwenddocs{}\nwbegincode{1209}\moddef{Muli trapezium: trapezium: TBP}\plusendmoddef
  procedure :: is_left_of => muli_trapezium_is_left_of
\nwendcode{}\nwbegincode{1210}\moddef{Muli trapezium: procedures}\plusendmoddef
  pure function muli_trapezium_is_left_of (this, that) result (is_left)
    logical :: is_left
    class(muli_trapezium_t), intent(in) :: this, that
    is_left = this%r_position <= that%r_position   !-that%d_position
    ! if (is_left.and.that%r_position < this%r_position) then
    !    print *,"!"
    !    STOP
    ! end if
  end function muli_trapezium_is_left_of

\nwendcode{}\nwbegindocs{1211} %def muli_trapezium_is_left_of  
\nwenddocs{}\nwbegindocs{1212}\nwdocspar
\nwenddocs{}\nwbegincode{1213}\moddef{Muli trapezium: trapezium: TBP}\plusendmoddef
  procedure :: includes => muli_trapezium_includes
\nwendcode{}\nwbegincode{1214}\moddef{Muli trapezium: procedures}\plusendmoddef
  elemental logical function muli_trapezium_includes &
       (this, dim, position, value, integral, probability) result (includes)
    class(muli_trapezium_t), intent(in) :: this
    integer, intent(in) :: dim
    real(default), intent(in),optional :: position, value, integral, probability
    includes = .true.
    if (present (position)) then
       if (this%get_l_position() > position .or. &
           position >= this%get_r_position())  includes = .false.
    end if
    if (present (value)) then
       if (this%get_l_value(dim) > value .or. value >= &
            this%get_r_value(dim))  includes = .false.
    end if
    if (present (integral)) then
       if (this%get_l_integral(dim) > integral .or. integral >= &
            this%get_r_integral(dim))  includes = .false.
    end if
    if (present (probability)) then
       if (this%get_l_probability(dim) > probability .or. &
            probability >= this%get_r_probability(dim))  includes = .false.
    end if    
  end function muli_trapezium_includes

\eatline
\nwendcode{}\nwbegindocs{1215}\nwdocspar
\nwenddocs{}\nwbegincode{1216}\moddef{Muli trapezium: trapezium: TBP}\plusendmoddef
  procedure :: to_node => muli_trapezium_to_node
\nwendcode{}\nwbegincode{1217}\moddef{Muli trapezium: procedures}\plusendmoddef
  subroutine muli_trapezium_to_node (this, value, list, tree)
    class(muli_trapezium_t), intent(in) :: this
    real(default), intent(in) :: value  
    ! class(muli_trapezium_node_class_t), optional, pointer, intent(out) :: node
    class(muli_trapezium_list_t), optional, pointer, intent(out) :: list
    class(muli_trapezium_tree_t), optional, pointer, intent(out) :: tree
    ! if (present (node)) then
    !    allocate (node)
    !    node%dim = this%dim
    !    node%r_position = this%r_position
    !    node%d_position = this%d_position
    !    allocate (node%values (this%dim, value_dimension), source=this%values)
    ! end if
    if (present (list)) then
       allocate (list)
       list%dim = this%dim
       list%r_position = this%r_position
       list%d_position = this%d_position
       allocate (list%values (0:this%dim-1, value_dimension), source=this%values)
    end if
    if (present (tree)) then
       allocate (tree)
       tree%dim = this%dim
       tree%r_position = this%r_position
       tree%d_position = this%d_position
       allocate (tree%values (0:this%dim-1, value_dimension), source=this%values)
    end if    
  end subroutine muli_trapezium_to_node

\eatline
\nwendcode{}\nwbegindocs{1218}\nwdocspar
\nwenddocs{}\nwbegincode{1219}\moddef{Muli trapezium: trapezium: TBP}\plusendmoddef
  procedure :: sum_up => muli_trapezium_sum_up  
\nwendcode{}\nwbegincode{1220}\moddef{Muli trapezium: procedures}\plusendmoddef
  subroutine muli_trapezium_sum_up (this)
    class(muli_trapezium_t), intent(inout) :: this
    integer :: i
    if (allocated (this%values)) then
       do i = 1, 7
          this%values (0,i) = sum (this%values (1:this%dim-1,i))
       end do
    end if
  end subroutine muli_trapezium_sum_up
  
\eatline
\nwendcode{}\nwbegindocs{1221}\nwdocspar
\nwenddocs{}\nwbegincode{1222}\moddef{Muli trapezium: trapezium: TBP}\plusendmoddef
  procedure :: approx_value => muli_trapezium_approx_value
\nwendcode{}\nwbegincode{1223}\moddef{Muli trapezium: procedures}\plusendmoddef
  pure function muli_trapezium_approx_value (this, x) result (val)
    ! returns the values at x
    class(muli_trapezium_t), intent(in) :: this
    real(default), dimension(this%dim) :: val
    real(default), intent(in) :: x
    val = this%get_r_value_array() + (x - this%r_position) * &
         this%get_d_value() / this%d_position
  end function muli_trapezium_approx_value

\eatline
\nwendcode{}\nwbegindocs{1224}This function returns the value at {\tt{}x}. 
\nwenddocs{}\nwbegincode{1225}\moddef{Muli trapezium: trapezium: TBP}\plusendmoddef
  procedure :: approx_value_n => muli_trapezium_approx_value_n
\nwendcode{}\nwbegincode{1226}\moddef{Muli trapezium: procedures}\plusendmoddef
  elemental function muli_trapezium_approx_value_n (this, x, n) result (val)
    class(muli_trapezium_t), intent(in) :: this
    real(default) :: val
    real(default), intent(in) :: x
    integer, intent(in) :: n
    val = this%get_r_value_element(n) + (x - this%r_position) * &
         this%get_d_value_element(n) / this%d_position
  end function muli_trapezium_approx_value_n

\eatline
\nwendcode{}\nwbegindocs{1227}This function returns the integral from {\tt{}x} to {\tt{}r{\char95}position}.
\nwenddocs{}\nwbegincode{1228}\moddef{Muli trapezium: trapezium: TBP}\plusendmoddef
  procedure :: approx_integral => muli_trapezium_approx_integral
\nwendcode{}\nwbegincode{1229}\moddef{Muli trapezium: procedures}\plusendmoddef
  pure function muli_trapezium_approx_integral (this, x)
    class(muli_trapezium_t), intent(in) :: this
    real(default), dimension(this%dim) :: muli_trapezium_approx_integral
    real(default), intent(in) :: x
    muli_trapezium_approx_integral = &
         ! this%get_r_integral()+&
         ! (this%r_position-x)*this%get_r_value()+&
         ! (x**2-this%r_position**2)*this%get_d_integral()/(this%d_position*2D0)
         this%get_r_integral() + &
         ((this%r_position - x) * &
         (-this%get_d_value() * (this%r_position - x) + 2 * &
         this%d_position*this%get_r_value_array())) / &
         (2 * this%d_position)
  end function muli_trapezium_approx_integral
  
\eatline
\nwendcode{}\nwbegindocs{1230}This function returns the integral from {\tt{}x} to {\tt{}r{\char95}position}.
\nwenddocs{}\nwbegincode{1231}\moddef{Muli trapezium: trapezium: TBP}\plusendmoddef
  procedure :: approx_integral_n => muli_trapezium_approx_integral_n
\nwendcode{}\nwbegincode{1232}\moddef{Muli trapezium: procedures}\plusendmoddef
  elemental function muli_trapezium_approx_integral_n (this, x, n) result (val)    
    class(muli_trapezium_t), intent(in) :: this
    real(default) :: val
    real(default), intent(in) :: x
    integer, intent(in) :: n
    val = this%get_r_integral_element (n) + ((this%r_position - x) * &
         (-this%get_d_value_element (n) * (this%r_position - x) + 2 * &
         this%d_position * this%get_r_value_element (n))) / &
         (2 * this%d_position)
  end function muli_trapezium_approx_integral_n
    
\eatline
\nwendcode{}\nwbegindocs{1233}This function returns the values at {\tt{}x}.
\nwenddocs{}\nwbegincode{1234}\moddef{Muli trapezium: trapezium: TBP}\plusendmoddef
  procedure :: approx_probability => muli_trapezium_approx_probability
\nwendcode{}\nwbegincode{1235}\moddef{Muli trapezium: procedures}\plusendmoddef
   pure function muli_trapezium_approx_probability (this, x) result (prop)
    class(muli_trapezium_t), intent(in) :: this
    real(default), dimension(this%dim) :: prop
    real(default), intent(in) :: x
    prop = exp (- this%approx_integral (x))
  end function muli_trapezium_approx_probability

\eatline
\nwendcode{}\nwbegindocs{1236} This function returns the integral from {\tt{}x} to {\tt{}r{\char95}position}.
\nwenddocs{}\nwbegincode{1237}\moddef{Muli trapezium: trapezium: TBP}\plusendmoddef
  procedure :: approx_probability_n => muli_trapezium_approx_probability_n
\nwendcode{}\nwbegincode{1238}\moddef{Muli trapezium: procedures}\plusendmoddef
  elemental function muli_trapezium_approx_probability_n (this, x, n) result (val)
    class(muli_trapezium_t), intent(in) :: this
    real(default) :: val
    real(default), intent(in) :: x
    integer, intent(in) :: n
    val = exp (- this%approx_integral_n (x, n))
  end function muli_trapezium_approx_probability_n
     
\eatline
\nwendcode{}\nwbegindocs{1239}\nwdocspar
\nwenddocs{}\nwbegincode{1240}\moddef{Muli trapezium: trapezium: TBP}\plusendmoddef
  procedure :: approx_position_by_integral => &
       muli_trapezium_approx_position_by_integral
\nwendcode{}\nwbegincode{1241}\moddef{Muli trapezium: procedures}\plusendmoddef
  elemental function muli_trapezium_approx_position_by_integral &
       (this, dim, int) result (val)
    class(muli_trapezium_t), intent(in) :: this
    real(default) :: val
    integer, intent(in) :: dim
    real(default), intent(in) :: int
    real(default) :: dpdv    
    dpdv = (this%d_position / this%values (dim,d_value_index))
    val = this%r_position - dpdv * (this%values (dim, r_value_index) - &
           sqrt (((this%values (dim, r_integral_index) - int) * two / dpdv) + &
           this%values (dim, r_value_index)**2))
  end function muli_trapezium_approx_position_by_integral
  
\eatline
\nwendcode{}\nwbegindocs{1242}\nwdocspar
\nwenddocs{}\nwbegincode{1243}\moddef{Muli trapezium: trapezium: TBP}\plusendmoddef
  ! procedure :: choose_partons => muli_trapezium_choose_partons
\nwendcode{}\nwbegindocs{1244}\nwdocspar
\nwenddocs{}\nwbegindocs{1245}\nwdocspar
\nwenddocs{}\nwbegincode{1246}\moddef{Muli trapezium: trapezium: TBP}\plusendmoddef
  procedure :: split => muli_trapezium_split
\nwendcode{}\nwbegincode{1247}\moddef{Muli trapezium: procedures}\plusendmoddef
  subroutine muli_trapezium_split (this, c_value, c_position, new_node)
    class(muli_trapezium_t), intent(inout) :: this
    real(default), intent(in) :: c_position
    real(default), intent(in), dimension(this%dim) :: c_value
    class(muli_trapezium_t), intent(out), pointer :: new_node
    real(default) :: ndpr, ndpl
    real(default), dimension(:), allocatable :: ov, edv
    ndpr = this%r_position - c_position
    ndpl = this%d_position - ndpr
    allocate (ov (0:this%dim-1), source=this%get_r_value_array() - ndpr * &
         this%get_d_value() / this%d_position)
    allocate (edv (0:this%dim-1), source=c_value-ov)
    allocate (new_node)
    call new_node%initialize (dim=this%dim, r_position=c_position, &
         d_position=ndpl)
    call new_node%set_r_value (c_value)
    call new_node%set_d_value (this%get_d_value() + &
         c_value-this%get_r_value_array())
    call new_node%set_d_integral (ndpl*(this%get_d_value() - &
         this%get_r_value_array() - c_value) / two)
    call new_node%set_error (abs((edv*ndpl) / two))
    ! new_node%measure_comp = sum (abs((edv*ndpl) / two))
    this%d_position = ndpr
    call this%set_d_value (this%get_r_value_array() - c_value)
    call this%set_d_integral (- (ndpr*(this%get_r_value_array() + c_value) / two))
    call this%set_error (abs(edv*ndpr / two))
    ! this%measure_comp = sum (abs(edv*ndpr / two))
    ! write (*, "(1x,A)") "muli_trapezium_split: new errors:"
    ! write (*, "(3x,ES14.7)")  this%get_error()
    ! write (*, "(3x,ES14.7)")  new_node%get_error()
    ! write (*, "(3x,11(ES20.10)")  new_node%get_d_integral()
    ! write (*, "(3x,11(ES20.10)")  this%get_d_integral()
  end subroutine muli_trapezium_split
  
\eatline
\nwendcode{}\nwbegindocs{1248}\nwdocspar
\nwenddocs{}\nwbegincode{1249}\moddef{Muli trapezium: trapezium: TBP}\plusendmoddef
  procedure :: update => muli_trapezium_update  
\nwendcode{}\nwbegincode{1250}\moddef{Muli trapezium: procedures}\plusendmoddef
  subroutine muli_trapezium_update (this)
    class(muli_trapezium_t), intent(inout) :: this
    real(default), dimension(:), allocatable :: integral
    allocate (integral (0:this%dim-1), source=this%get_d_integral())
    call this%set_d_integral (-this%d_position * (this%get_r_value_array() &
         - this%get_d_value() / 2))
    call this%set_error (abs (this%get_d_integral() - integral))
    ! write (*, "(3x,11(ES20.10)")  this%get_d_integral()
  end subroutine muli_trapezium_update
  
\eatline
\nwendcode{}\nwbegindocs{1251}\nwdocspar
\nwenddocs{}\nwbegincode{1252}\moddef{Muli trapezium: node: TBP}\endmoddef
  procedure :: deserialize_from_marker => &
       muli_trapezium_node_deserialize_from_marker  
\nwendcode{}\nwbegincode{1253}\moddef{Muli trapezium: procedures}\plusendmoddef
  subroutine muli_trapezium_node_deserialize_from_marker (this, name, marker)
    class(muli_trapezium_node_class_t), intent(out) :: this
    character(*), intent(in) :: name
    class(marker_t), intent(inout) :: marker
    integer(dik) :: status
    class(ser_class_t), pointer :: ser
    allocate (muli_trapezium_tree_t :: ser)
    call marker%push_reference (ser)
    allocate (muli_trapezium_list_t::ser)
    call marker%push_reference (ser)
    call serializable_deserialize_from_marker (this, name, marker)
    call marker%pop_reference (ser)
    deallocate (ser)
    call marker%pop_reference (ser)
    deallocate (ser)
  end subroutine muli_trapezium_node_deserialize_from_marker
  
\eatline
\nwendcode{}\nwbegindocs{1254}\nwdocspar
\nwenddocs{}\nwbegincode{1255}\moddef{Muli trapezium: node: TBP}\plusendmoddef
  procedure(muli_trapezium_append_interface), deferred :: append  
\nwendcode{}\nwbegindocs{1256}\nwdocspar
\nwenddocs{}\nwbegindocs{1257}\nwdocspar
\nwenddocs{}\nwbegincode{1258}\moddef{Muli trapezium: node: TBP}\plusendmoddef
  procedure(muli_trapezium_final_interface), deferred :: finalize
\nwendcode{}\nwbegindocs{1259}\nwdocspar
\nwenddocs{}\nwbegindocs{1260} 
\nwenddocs{}\nwbegincode{1261}\moddef{Muli trapezium: list: TBP}\endmoddef
  procedure :: append => muli_trapezium_list_append  
\nwendcode{}\nwbegincode{1262}\moddef{Muli trapezium: procedures}\plusendmoddef
  subroutine muli_trapezium_list_append (this, right)
    class(muli_trapezium_list_t), intent(inout), target :: this
    class(muli_trapezium_node_class_t), intent(inout), target :: right
    this%right => right
    right%left => this
  end subroutine muli_trapezium_list_append

\eatline
\nwendcode{}\nwbegindocs{1263}\nwdocspar
\nwenddocs{}\nwbegincode{1264}\moddef{Muli trapezium: node: TBP}\plusendmoddef
  procedure :: nullify => muli_trapezium_node_nullify
\nwendcode{}\nwbegincode{1265}\moddef{Muli trapezium: procedures}\plusendmoddef
  subroutine muli_trapezium_node_nullify (this)
    class(muli_trapezium_node_class_t), intent(out) :: this
    nullify (this%left)
    nullify (this%right)
  end subroutine muli_trapezium_node_nullify
  
\eatline
\nwendcode{}\nwbegindocs{1266}\nwdocspar
\nwenddocs{}\nwbegincode{1267}\moddef{Muli trapezium: node: TBP}\plusendmoddef
  procedure :: get_left => muli_trapezium_node_get_left
\nwendcode{}\nwbegincode{1268}\moddef{Muli trapezium: procedures}\plusendmoddef
  subroutine muli_trapezium_node_get_left (this, left)
    class(muli_trapezium_node_class_t), intent(in) :: this
    class(muli_trapezium_node_class_t), pointer, intent(out) :: left
    left => this%left
  end subroutine muli_trapezium_node_get_left
  
\eatline
\nwendcode{}\nwbegindocs{1269}\nwdocspar
\nwenddocs{}\nwbegincode{1270}\moddef{Muli trapezium: node: TBP}\plusendmoddef
  procedure :: get_right => muli_trapezium_node_get_right
\nwendcode{}\nwbegincode{1271}\moddef{Muli trapezium: procedures}\plusendmoddef
  subroutine muli_trapezium_node_get_right (this, right)
    class(muli_trapezium_node_class_t), intent(in) :: this
    class(muli_trapezium_node_class_t), pointer, intent(out) :: right
    right => this%right
  end subroutine muli_trapezium_node_get_right
  
\eatline
\nwendcode{}\nwbegindocs{1272}\nwdocspar
\nwenddocs{}\nwbegincode{1273}\moddef{Muli trapezium: node: TBP}\plusendmoddef
  procedure :: get_leftmost => muli_trapezium_node_get_leftmost
\nwendcode{}\nwbegincode{1274}\moddef{Muli trapezium: procedures}\plusendmoddef
  subroutine muli_trapezium_node_get_leftmost (this, node)
    class(muli_trapezium_node_class_t), intent(in) :: this
    class(muli_trapezium_node_class_t), pointer, intent(out) :: node
    if (associated (this%left)) then
       node => this%left
       do while (associated (node%left))
          node => node%left
       end do
    else
       nullify (node)
    end if
  end subroutine muli_trapezium_node_get_leftmost
  
\eatline
\nwendcode{}\nwbegindocs{1275}\nwdocspar
\nwenddocs{}\nwbegincode{1276}\moddef{Muli trapezium: node: TBP}\plusendmoddef
  procedure :: get_rightmost => muli_trapezium_node_get_rightmost
\nwendcode{}\nwbegincode{1277}\moddef{Muli trapezium: procedures}\plusendmoddef
  subroutine muli_trapezium_node_get_rightmost (this, right)
    class(muli_trapezium_node_class_t), intent(in) :: this
    class(muli_trapezium_node_class_t), pointer, intent(out) :: right
    if (associated (this%right)) then
       right => this%right
       do while (associated (right%right))
          right => right%right
       end do
    else
       nullify (right)
    end if
  end subroutine muli_trapezium_node_get_rightmost
  
\eatline
\nwendcode{}\nwbegindocs{1278}\nwdocspar
\nwenddocs{}\nwbegincode{1279}\moddef{Muli trapezium: node: TBP}\plusendmoddef
  generic :: decide => decide_by_value, decide_by_position  
  procedure :: decide_by_value => muli_trapezium_node_decide_by_value  
  procedure :: decide_by_position => muli_trapezium_node_decide_by_position
\nwendcode{}\nwbegincode{1280}\moddef{Muli trapezium: procedures}\plusendmoddef
  subroutine muli_trapezium_node_decide_by_value (this, value, dim, record, node)
    class(muli_trapezium_node_class_t), intent(in) :: this
    real(default), intent(in) :: value
    integer, intent(in) :: record, dim
    class(muli_trapezium_node_class_t), pointer, intent(out) :: node
    if (this%values (dim, record) > value) then
       node => this%left
    else
       node => this%right
    end if
  end subroutine muli_trapezium_node_decide_by_value

\eatline
\nwendcode{}\nwbegindocs{1281}\nwdocspar
\nwenddocs{}\nwbegincode{1282}\moddef{Muli trapezium: procedures}\plusendmoddef
  subroutine muli_trapezium_node_decide_by_position (this, position, node)
    class(muli_trapezium_node_class_t), intent(in) :: this
    real(default), intent(in) :: position
    class(muli_trapezium_node_class_t), pointer, intent(out) :: node
    if (this%r_position > position) then
       node => this%left
    else
       node => this%right
    end if
  end subroutine muli_trapezium_node_decide_by_position
  
\eatline
\nwendcode{}\nwbegindocs{1283} 
\nwenddocs{}\nwbegincode{1284}\moddef{Muli trapezium: node: TBP}\plusendmoddef
  procedure :: decide_decreasing => muli_trapezium_node_decide_decreasing
\nwendcode{}\nwbegincode{1285}\moddef{Muli trapezium: procedures}\plusendmoddef
  subroutine muli_trapezium_node_decide_decreasing &
       (this, value, dim, record, node)
    class(muli_trapezium_node_class_t), intent(in) :: this
    real(default), intent(in) :: value
    integer, intent(in) :: record, dim
    class(muli_trapezium_node_class_t), pointer, intent(out) :: node
    if (this%values (dim, record) <= value) then
       node => this%left
    else
       node => this%right
    end if
  end subroutine muli_trapezium_node_decide_decreasing
  
\eatline
\nwendcode{}\nwbegindocs{1286}\nwdocspar
\nwenddocs{}\nwbegincode{1287}\moddef{Muli trapezium: node: TBP}\plusendmoddef
  procedure :: to_tree => muli_trapezium_node_to_tree
\nwendcode{}\nwbegincode{1288}\moddef{Muli trapezium: procedures}\plusendmoddef
  subroutine muli_trapezium_node_to_tree (this, out_tree)
    class(muli_trapezium_node_class_t), target, intent(in) :: this
    class(muli_trapezium_tree_t), intent(out) :: out_tree
    out_tree%left => this%left
    out_tree%right => this%right
  end subroutine muli_trapezium_node_to_tree
  
\eatline
\nwendcode{}\nwbegindocs{1289}\nwdocspar
\nwenddocs{}\nwbegincode{1290}\moddef{Muli trapezium: node: TBP}\plusendmoddef
  procedure :: untangle => muli_trapezium_node_untangle
\nwendcode{}\nwbegincode{1291}\moddef{Muli trapezium: procedures}\plusendmoddef
  subroutine muli_trapezium_node_untangle(this)
    class(muli_trapezium_node_class_t), intent(inout), target :: this
    if (associated (this%left)) then
       if (associated (this%left%right, this)) then
          nullify (this%left%right)
          nullify (this%left)
       end if
    end if    
  end subroutine muli_trapezium_node_untangle

\eatline
\nwendcode{}\nwbegindocs{1292} 
\nwenddocs{}\nwbegincode{1293}\moddef{Muli trapezium: node: TBP}\plusendmoddef
  procedure :: apply => muli_trapezium_node_apply 
\nwendcode{}\nwbegincode{1294}\moddef{Muli trapezium: procedures}\plusendmoddef
  recursive subroutine muli_trapezium_node_apply(this,proc)
    class(muli_trapezium_node_class_t), intent(inout) :: this
    interface
       subroutine proc(this)
         import muli_trapezium_node_class_t
         class(muli_trapezium_node_class_t), intent(inout) :: this
       end subroutine proc
    end interface
    if (associated(this%right))call proc(this%right)
    if (associated(this%left))call proc(this%left)
    call proc(this)
  end subroutine muli_trapezium_node_apply
 
\eatline
\nwendcode{}\nwbegindocs{1295}\nwdocspar
\nwenddocs{}\nwbegincode{1296}\moddef{Muli trapezium: node: TBP}\plusendmoddef
  ! procedure :: copy => muli_trapezium_node_copy
  ! generic :: assignment(=) => copy
  ! procedure, deferred :: approx => muli_trapezium_node_approx
\nwendcode{}\nwbegindocs{1297}\nwdocspar
\nwenddocs{}\nwbegindocs{1298}\nwdocspar
\nwenddocs{}\nwbegincode{1299}\moddef{Muli trapezium: tree: TBP}\endmoddef
  procedure :: write_to_marker => muli_trapezium_tree_write_to_marker
\nwendcode{}\nwbegincode{1300}\moddef{Muli trapezium: procedures}\plusendmoddef
  subroutine muli_trapezium_tree_write_to_marker (this, marker, status)
    class(muli_trapezium_tree_t), intent(in) :: this
    class(marker_t), intent(inout) :: marker
    integer(dik), intent(out) :: status
    class(muli_trapezium_list_t), pointer :: list
    class(ser_class_t), pointer :: ser
    call marker%mark_begin ("muli_trapezium_tree_t")
    call this%get_left_list (list)
    ser => list
    call marker%mark_pointer ("list", ser)
    call marker%mark_end ("muli_trapezium_tree_t")
  end subroutine muli_trapezium_tree_write_to_marker

\eatline
\nwendcode{}\nwbegindocs{1301}\nwdocspar
\nwenddocs{}\nwbegincode{1302}\moddef{Muli trapezium: tree: TBP}\plusendmoddef
  procedure :: read_from_marker => muli_trapezium_tree_read_from_marker  
\nwendcode{}\nwbegincode{1303}\moddef{Muli trapezium: procedures}\plusendmoddef
  subroutine muli_trapezium_tree_read_from_marker (this, marker, status)
    class(muli_trapezium_tree_t), intent(out) :: this
    class(marker_t), intent(inout) :: marker
    integer(dik), intent(out) :: status
    class(ser_class_t), pointer :: ser
    call marker%pick_begin ("muli_trapezium_tree_t", status=status)
    call marker%pick_pointer ("list", ser)
    if (associated (ser)) then
       select type (ser)
       class is (muli_trapezium_list_t)
          call ser%to_tree (this)
       class default
          nullify (this%left)
          nullify (this%right)
          nullify (this%down)
       end select
    else
       nullify (this%left)
       nullify (this%right)
       nullify (this%down)
    end if
    call marker%pick_end ("muli_trapezium_tree_t", status)
  end subroutine muli_trapezium_tree_read_from_marker

\eatline
\nwendcode{}\nwbegindocs{1304}\nwdocspar
\nwenddocs{}\nwbegincode{1305}\moddef{Muli trapezium: tree: TBP}\plusendmoddef
  procedure :: print_to_unit => muli_trapezium_tree_print_to_unit    
\nwendcode{}\nwbegincode{1306}\moddef{Muli trapezium: procedures}\plusendmoddef
  recursive subroutine muli_trapezium_tree_print_to_unit &
       (this, unit, parents, components, peers)
    class(muli_trapezium_tree_t), intent(in) :: this
    integer, intent(in) :: unit
    integer(dik), intent(in) :: parents, components, peers
    class(ser_class_t), pointer :: ser
    if (parents > 0)  call muli_trapezium_print_to_unit &
         (this, unit, parents-1, components, peers)
    ser => this%down
    call serialize_print_peer_pointer (ser, unit, i_one, i_zero, i_one, "DOWN")
    if (associated (this%left)) then
       select type (sertmp => this%left)
       class is (muli_trapezium_list_t)
          ser => sertmp
          call serialize_print_peer_pointer &
               (ser, unit, parents, components, i_zero, "LEFT")
       class default
          call serialize_print_peer_pointer &
               (ser, unit, parents, components, peers, "LEFT")
       end select
    else
       write (unit, "(1x,A)")  "Left is not associated."
    end if
    if (associated (this%right)) then
       select type (sertmp => this%right)
       class is (muli_trapezium_list_t)
          ser => sertmp
          call serialize_print_peer_pointer &
               (ser, unit, parents, components, i_zero, "RIGHT")
       class default
          call serialize_print_peer_pointer &
               (ser, unit, parents, components, peers, "RIGHT")
       end select
    else
       write (unit, "(1x,A)")  "Right is not associated."
    end if
  end subroutine muli_trapezium_tree_print_to_unit
  
\eatline
\nwendcode{}\nwbegindocs{1307}\nwdocspar
\nwenddocs{}\nwbegincode{1308}\moddef{Muli trapezium: tree: TBP}\plusendmoddef
  procedure, nopass :: get_type => muli_trapezium_tree_get_type
\nwendcode{}\nwbegincode{1309}\moddef{Muli trapezium: procedures}\plusendmoddef
  pure subroutine muli_trapezium_tree_get_type (type)
    character(:),allocatable, intent(out) :: type
    allocate (type, source="muli_trapezium_tree_t")
  end subroutine muli_trapezium_tree_get_type
  
\eatline
\nwendcode{}\nwbegindocs{1310}\nwdocspar
\nwenddocs{}\nwbegincode{1311}\moddef{Muli trapezium: tree: TBP}\plusendmoddef
  procedure, nopass :: verify_type => muli_trapezium_tree_verify_type  
\nwendcode{}\nwbegincode{1312}\moddef{Muli trapezium: procedures}\plusendmoddef
  elemental logical function muli_trapezium_tree_verify_type (type) result (match)
    character(*), intent(in) :: type
    match = type == "muli_trapezium_tree_t"
  end function muli_trapezium_tree_verify_type
  
\eatline
\nwendcode{}\nwbegindocs{1313}\nwdocspar
\nwenddocs{}\nwbegincode{1314}\moddef{Muli trapezium: tree: TBP}\plusendmoddef
  procedure :: nullify => muli_trapezium_tree_nullify  
\nwendcode{}\nwbegincode{1315}\moddef{Muli trapezium: procedures}\plusendmoddef
  subroutine muli_trapezium_tree_nullify (this)
    class(muli_trapezium_tree_t), intent(out) :: this
    call muli_trapezium_node_nullify (this)
    nullify (this%down)
  end subroutine muli_trapezium_tree_nullify
    
\eatline
\nwendcode{}\nwbegindocs{1316}\nwdocspar
\nwenddocs{}\nwbegincode{1317}\moddef{Muli trapezium: tree: TBP}\plusendmoddef
  procedure :: finalize => muli_trapezium_tree_finalize
\nwendcode{}\nwbegincode{1318}\moddef{Muli trapezium: procedures}\plusendmoddef
  recursive subroutine muli_trapezium_tree_finalize (this)
    class(muli_trapezium_tree_t), intent(inout) :: this
    if (associated (this%right)) then
       call this%right%untangle ()
       call this%right%finalize ()
       deallocate (this%right)
    end if
    if (associated (this%left)) then
       call this%left%untangle ()
       call this%left%finalize ()
       deallocate (this%left)
    end if
    this%dim = 0 
  end subroutine muli_trapezium_tree_finalize
  
\eatline
\nwendcode{}\nwbegindocs{1319}\nwdocspar
\nwenddocs{}\nwbegincode{1320}\moddef{Muli trapezium: tree: TBP}\plusendmoddef
  procedure :: decide_by_value => muli_trapezium_tree_decide_by_value    
\nwendcode{}\nwbegincode{1321}\moddef{Muli trapezium: procedures}\plusendmoddef
  subroutine muli_trapezium_tree_decide_by_value (this, value, dim, record, node)
    class(muli_trapezium_tree_t), intent(in) :: this
    real(default), intent(in) :: value
    integer, intent(in) :: record, dim
    class(muli_trapezium_node_class_t), pointer, intent(out) :: node
    if (this%down%values (dim, record) > value) then
       node => this%left
    else
       node => this%right
    end if
  end subroutine muli_trapezium_tree_decide_by_value
  
\eatline
\nwendcode{}\nwbegindocs{1322}\nwdocspar
\nwenddocs{}\nwbegincode{1323}\moddef{Muli trapezium: tree: TBP}\plusendmoddef
  procedure :: decide_by_position => muli_trapezium_tree_decide_by_position
\nwendcode{}\nwbegincode{1324}\moddef{Muli trapezium: procedures}\plusendmoddef
  subroutine muli_trapezium_tree_decide_by_position (this, position, node)
    class(muli_trapezium_tree_t), intent(in) :: this
    real(default), intent(in) :: position
    class(muli_trapezium_node_class_t), pointer, intent(out) :: node
    if (this%down%r_position > position) then
       node => this%left
    else
       node => this%right
    end if
  end subroutine muli_trapezium_tree_decide_by_position
  
\eatline
\nwendcode{}\nwbegincode{1325}\moddef{Muli trapezium: tree: TBP}\plusendmoddef
  procedure :: decide_decreasing => muli_trapezium_tree_decide_decreasing
\nwendcode{}\nwbegincode{1326}\moddef{Muli trapezium: procedures}\plusendmoddef
  subroutine muli_trapezium_tree_decide_decreasing &
       (this, value, dim, record, node)
    class(muli_trapezium_tree_t), intent(in) :: this
    real(default), intent(in) :: value
    integer, intent(in) :: record, dim
    ! integer, save :: count=0
    class(muli_trapezium_node_class_t), pointer, intent(out) :: node
    ! count = count + 1
    if (this%down%values (dim, record) <= value) then
       ! print ('("Decide: value(",I2,",",I1,")=",E20.7," > ",E20.7, &
       !     ": go left.")'), dim, record, this%down%values(dim, record), value
       node => this%left
    else
       ! print ('("Decide: value(",I2,",",I1,")=",E20.7," <= ", &
       !   E20.7,": go right.")'), &
       !   dim, record, this%down%values(dim, record), value
       node => this%right
    end if
  end subroutine muli_trapezium_tree_decide_decreasing
  
\eatline
\nwendcode{}\nwbegindocs{1327}\nwdocspar
\nwenddocs{}\nwbegincode{1328}\moddef{Muli trapezium: tree: TBP}\plusendmoddef
  procedure :: get_left_list => muli_trapezium_tree_get_left_list
\nwendcode{}\nwbegincode{1329}\moddef{Muli trapezium: procedures}\plusendmoddef
  subroutine muli_trapezium_tree_get_left_list (this, list)
    class(muli_trapezium_tree_t), intent(in) :: this
    class(muli_trapezium_list_t), pointer, intent(out) :: list
    class(muli_trapezium_node_class_t), pointer::node
    call this%get_leftmost (node)
    if (associated (node)) then
       select type (node)
       class is (muli_trapezium_list_t)
          list => node
       class default
          nullify (list)
       end select
    else
       nullify (list)
    end if
  end subroutine muli_trapezium_tree_get_left_list

\eatline
\nwendcode{}\nwbegindocs{1330}\nwdocspar
\nwenddocs{}\nwbegincode{1331}\moddef{Muli trapezium: tree: TBP}\plusendmoddef
  procedure :: get_right_list => muli_trapezium_tree_get_right_list  
\nwendcode{}\nwbegincode{1332}\moddef{Muli trapezium: procedures}\plusendmoddef
  subroutine muli_trapezium_tree_get_right_list (this, list)
    class(muli_trapezium_tree_t), intent(in) :: this
    class(muli_trapezium_list_t), pointer, intent(out) :: list
    class(muli_trapezium_node_class_t), pointer::node
    call this%get_rightmost (node)
    if (associated (node)) then
       select type (node)
       class is (muli_trapezium_list_t)
          list => node
       class default
          nullify (list)
       end select
    else
       nullify (list) 
    end if
  end subroutine muli_trapezium_tree_get_right_list
  
\eatline
\nwendcode{}\nwbegindocs{1333}\nwdocspar
\nwenddocs{}\nwbegincode{1334}\moddef{Muli trapezium: tree: TBP}\plusendmoddef
  generic :: find => find_by_value, find_by_position  
  procedure :: find_by_value => muli_trapezium_tree_find_by_value
  procedure :: find_by_position => muli_trapezium_tree_find_by_position  
\nwendcode{}\nwbegincode{1335}\moddef{Muli trapezium: procedures}\plusendmoddef
  subroutine muli_trapezium_tree_find_by_value (this, value, dim, record, node)
    class(muli_trapezium_tree_t), intent(in), target :: this
    real(default), intent(in) :: value
    integer, intent(in) :: record, dim
    class(muli_trapezium_node_class_t), pointer, intent(out) :: node
    node => this
    do while (.not. allocated (node%values))
       call node%decide (value, dim, record, node)
    end do
  end subroutine muli_trapezium_tree_find_by_value

\eatline
\nwendcode{}\nwbegindocs{1336} 
\nwenddocs{}\nwbegincode{1337}\moddef{Muli trapezium: procedures}\plusendmoddef
  subroutine muli_trapezium_tree_find_by_position (this, position, node)
    class(muli_trapezium_tree_t), intent(in), target :: this
    real(default), intent(in) :: position
    class(muli_trapezium_node_class_t), pointer, intent(out) :: node
    node => this
    do while (.not. allocated (node%values))
       call node%decide (position, node)
    end do
  end subroutine muli_trapezium_tree_find_by_position

\eatline
\nwendcode{}\nwbegindocs{1338}\nwdocspar
\nwenddocs{}\nwbegincode{1339}\moddef{Muli trapezium: tree: TBP}\plusendmoddef
  procedure :: find_decreasing => muli_trapezium_tree_find_decreasing
\nwendcode{}\nwbegincode{1340}\moddef{Muli trapezium: procedures}\plusendmoddef
  subroutine muli_trapezium_tree_find_decreasing (this, value, dim, node)
    class(muli_trapezium_tree_t), intent(in), target :: this
    real(default), intent(in) :: value
    integer, intent(in) :: dim
    class(muli_trapezium_node_class_t), pointer, intent(out) :: node
    node => this
    do while (.not. allocated (node%values))
       call node%decide_decreasing (value, dim, r_integral_index, node)
    end do
  end subroutine muli_trapezium_tree_find_decreasing
  
\eatline
\nwendcode{}\nwbegindocs{1341}\nwdocspar
\nwenddocs{}\nwbegincode{1342}\moddef{Muli trapezium: tree: TBP}\plusendmoddef
  procedure :: approx_by_integral => muli_trapezium_tree_approx_by_integral
\nwendcode{}\nwbegincode{1343}\moddef{Muli trapezium: procedures}\plusendmoddef
  subroutine muli_trapezium_tree_approx_by_integral &
       (this, int, dim, in_range, position, value, integral, content)
    class(muli_trapezium_tree_t), intent(in), target :: this
    real(default), intent(in) :: int
    integer, intent(in) :: dim
    logical, intent(out) :: in_range
    class(muli_trapezium_node_class_t), pointer, intent(out), optional :: content
    real(default), intent(out), optional :: position, value, integral
    integer :: i
    real(default) :: DINT    !,l_prop,r_prop,d_prop
    real(default) :: RP, DP, RV, DV, RI    !FC = gfortran
    class(muli_trapezium_node_class_t), pointer :: node
    node => this
    do while (.not. allocated (node%values))
       call node%decide_decreasing(INT, dim, r_integral_index, node)
    end do
    if (   int<=node%values(dim,r_integral_index)-node%values(dim,d_integral_index)&
         &.and.&
         &int>=node%values(dim,r_integral_index)) then
       in_range=.true.
!       associate(&!FC = nagfor
!            &RP=>node%r_position,&!FC = nagfor
!            &DP=>node%d_position,&!FC = nagfor
!            &RV=>node%values(dim,r_value_index),&!FC = nagfor
!            &DV=>node%values(dim,d_value_index),&!FC = nagfor
!            &RI=>node%values(dim,r_integral_index))!FC = nagfor
         RP=node%r_position!FC = gfortran
         DP=node%d_position!FC = gfortran
         RV=node%values(dim,r_value_index)!FC = gfortran
         DV=node%values(dim,d_value_index)!FC = gfortran
         RI=node%values(dim,r_integral_index)!FC = gfortran
         if (present(position)) then
            DINT=(ri-int)*2D0*dv/dp
            position=rp-(dp/dv)*(rv-sqrt(dint+rv**2))
         end if
         if (present(value)) then
            value=Sqrt(dp*(-2*dv*int + 2*dv*ri + dp*rv**2))/dp
         end if
         if (present(integral)) then
            integral=int
         end if
         if (present(content)) then
            content=>node
         end if
!       end associate!FC = nagfor
    else
       in_range=.false.
    end if
  end subroutine muli_trapezium_tree_approx_by_integral

\eatline
\nwendcode{}\nwbegindocs{1344}\nwdocspar
\nwenddocs{}\nwbegincode{1345}\moddef{Muli trapezium: tree: TBP}\plusendmoddef
  procedure :: approx_by_probability => muli_trapezium_tree_approx_by_probability
\nwendcode{}\nwbegincode{1346}\moddef{Muli trapezium: procedures}\plusendmoddef
  subroutine muli_trapezium_tree_approx_by_probability &
       (this, prop, dim, in_range, position, value, integral, content)
    class(muli_trapezium_tree_t), intent(in), target :: this
    real(default), intent(in) :: prop
    integer, intent(in) :: dim
    logical, intent(out) :: in_range
    class(muli_trapezium_node_class_t), pointer, intent(out), optional :: content
    real(default), intent(out), optional :: position, value, integral
    integer :: i
    real(default) :: int
    class(muli_trapezium_node_class_t), pointer :: node
    if (zero < prop .and. prop < one) then
       node => this
       int = -log (prop)
       call muli_trapezium_tree_approx_by_integral &
            (this, int, dim, in_range, position, value, integral, content)
    else
       in_range = .false.
    end if
  end subroutine muli_trapezium_tree_approx_by_probability
  
\eatline
\nwendcode{}\nwbegindocs{1347}\nwdocspar
\nwenddocs{}\nwbegincode{1348}\moddef{Muli trapezium: tree: TBP}\plusendmoddef
  procedure :: to_tree => muli_trapezium_tree_to_tree  
\nwendcode{}\nwbegincode{1349}\moddef{Muli trapezium: procedures}\plusendmoddef
  subroutine muli_trapezium_tree_to_tree (this, out_tree)
    class(muli_trapezium_tree_t), target, intent(in) :: this
    class(muli_trapezium_tree_t), intent(out) :: out_tree
    out_tree%left => this%left
    out_tree%right => this%right
    out_tree%down => this%down
  end subroutine muli_trapezium_tree_to_tree
  
\eatline
\nwendcode{}\nwbegindocs{1350}\nwdocspar
\nwenddocs{}\nwbegincode{1351}\moddef{Muli trapezium: tree: TBP}\plusendmoddef
  procedure :: append => muli_trapezium_tree_append    
\nwendcode{}\nwbegincode{1352}\moddef{Muli trapezium: procedures}\plusendmoddef
  subroutine muli_trapezium_tree_append(this,right)
    class(muli_trapezium_tree_t), intent(inout), target :: this
    class(muli_trapezium_node_class_t), intent(inout), target :: right
    call msg_error ("muli_trapezium_tree_append: Not yet implemented.")
  end subroutine muli_trapezium_tree_append
  
\eatline
\nwendcode{}\nwbegindocs{1353}\nwdocspar
\nwenddocs{}\nwbegincode{1354}\moddef{Muli trapezium: tree: TBP}\plusendmoddef
  procedure :: gnuplot => muli_trapezium_tree_gnuplot
\nwendcode{}\nwbegincode{1355}\moddef{Muli trapezium: procedures}\plusendmoddef
  subroutine muli_trapezium_tree_gnuplot (this, dir)
    class(muli_trapezium_tree_t), intent(in) :: this
    character(len=*), intent(in) :: dir
    class(muli_trapezium_list_t), pointer :: list
    call this%get_left_list (list)
    call list%gnuplot (dir)
  end subroutine muli_trapezium_tree_gnuplot
  
\eatline
\nwendcode{}\nwbegindocs{1356}\nwdocspar
\nwenddocs{}\nwbegincode{1357}\moddef{Muli trapezium: list: TBP}\plusendmoddef
  procedure :: write_to_marker => muli_trapezium_list_write_to_marker
\nwendcode{}\nwbegincode{1358}\moddef{Muli trapezium: procedures}\plusendmoddef
  recursive subroutine muli_trapezium_list_write_to_marker (this, marker, status)
    class(muli_trapezium_list_t), intent(in) :: this
    class(marker_t), intent(inout) :: marker
    integer(dik), intent(out) :: status 
    class(ser_class_t), pointer :: ser
    call marker%mark_begin ("muli_trapezium_list_t")
    call muli_trapezium_write_to_marker (this, marker, status)   
    ser => this%right
    call marker%mark_pointer ("right", ser)
    call marker%mark_end ("muli_trapezium_list_t")
  end subroutine muli_trapezium_list_write_to_marker

\eatline
\nwendcode{}\nwbegindocs{1359}\nwdocspar
\nwenddocs{}\nwbegincode{1360}\moddef{Muli trapezium: list: TBP}\plusendmoddef
  procedure :: read_from_marker => muli_trapezium_list_read_from_marker  
\nwendcode{}\nwbegincode{1361}\moddef{Muli trapezium: procedures}\plusendmoddef
  recursive subroutine muli_trapezium_list_read_from_marker (this, marker, status)
    class(muli_trapezium_list_t), intent(out) :: this
    class(marker_t), intent(inout) :: marker
    integer(dik), intent(out) :: status
    call msg_warning ("muli_trapezium_list_read_from_marker: " // &
         "You cannot deserialize a list with this subroutine.")
    call msg_error ("Use muli_trapezium_list_read_target_from_marker instead.")
  end subroutine muli_trapezium_list_read_from_marker

\eatline
\nwendcode{}\nwbegindocs{1362}\nwdocspar
\nwenddocs{}\nwbegincode{1363}\moddef{Muli trapezium: list: TBP}\plusendmoddef
  procedure :: read_target_from_marker => &
       muli_trapezium_list_read_target_from_marker  
\nwendcode{}\nwbegincode{1364}\moddef{Muli trapezium: procedures}\plusendmoddef
  recursive subroutine muli_trapezium_list_read_target_from_marker &
       (this, marker, status)
    class(muli_trapezium_list_t), target, intent(out) :: this
    class(marker_t), intent(inout) :: marker
    integer(dik), intent(out) :: status
    class(ser_class_t), pointer :: ser
    call marker%pick_begin ("muli_trapezium_list_t", status=status)
    call muli_trapezium_read_from_marker (this, marker, status)
    call marker%pick_pointer ("right", ser)
    if (associated (ser)) then
       select type (ser)
       class is (muli_trapezium_list_t)
          this%right => ser
          ser%left => this
       class default
          nullify (this%right)
          call msg_error ("muli_trapezium_list_read_target_from_marker: " &
               // "Unexpected type for right component.")
       end select
    else
       nullify (this%right)
    end if
    call marker%pick_end ("muli_trapezium_list_t", status)
  end subroutine muli_trapezium_list_read_target_from_marker
  
\eatline
\nwendcode{}\nwbegindocs{1365}\nwdocspar
\nwenddocs{}\nwbegincode{1366}\moddef{Muli trapezium: list: TBP}\plusendmoddef
  procedure :: print_to_unit => muli_trapezium_list_print_to_unit
\nwendcode{}\nwbegincode{1367}\moddef{Muli trapezium: procedures}\plusendmoddef
  recursive subroutine muli_trapezium_list_print_to_unit &
       (this, unit, parents, components, peers)
    class(muli_trapezium_list_t), intent(in) :: this
    integer, intent(in) :: unit
    integer(dik), intent(in) :: parents, components, peers
    class(ser_class_t), pointer :: ser
    if (parents > 0)  call muli_trapezium_print_to_unit &
         (this, unit, parents-1, components, peers)
    ser => this%left
    call serialize_print_peer_pointer &
         (ser, unit, -i_one, -i_one, -i_one, "LEFT")
    ser => this%right
    call serialize_print_peer_pointer &
         (ser, unit, parents, components, peers, "RIGHT")
  end subroutine muli_trapezium_list_print_to_unit
    
\eatline
\nwendcode{}\nwbegindocs{1368}\nwdocspar
\nwenddocs{}\nwbegincode{1369}\moddef{Muli trapezium: list: TBP}\plusendmoddef
  procedure, nopass :: get_type => muli_trapezium_list_get_type
\nwendcode{}\nwbegincode{1370}\moddef{Muli trapezium: procedures}\plusendmoddef
  pure subroutine muli_trapezium_list_get_type (type)
    character(:), allocatable, intent(out) :: type
    allocate (type, source="muli_trapezium_list_t")
  end subroutine muli_trapezium_list_get_type
  
\eatline
\nwendcode{}\nwbegindocs{1371}\nwdocspar
\nwenddocs{}\nwbegincode{1372}\moddef{Muli trapezium: list: TBP}\plusendmoddef
  procedure, nopass :: verify_type => muli_trapezium_list_verify_type
\nwendcode{}\nwbegindocs{1373}\nwdocspar
\nwenddocs{}\nwbegincode{1374}\moddef{Muli trapezium: procedures}\plusendmoddef
  elemental logical function muli_trapezium_list_verify_type (type) result (match)
    character(*), intent(in) :: type
    match = type == "muli_trapezium_list_t"
  end function muli_trapezium_list_verify_type
  
\eatline
\nwendcode{}\nwbegindocs{1375}\nwdocspar
\nwenddocs{}\nwbegincode{1376}\moddef{Muli trapezium: list: TBP}\plusendmoddef
  procedure :: finalize => muli_trapezium_list_finalize
\nwendcode{}\nwbegindocs{1377}\nwdocspar
\nwenddocs{}\nwbegincode{1378}\moddef{Muli trapezium: procedures}\plusendmoddef
  recursive subroutine muli_trapezium_list_finalize(this)
    class(muli_trapezium_list_t), intent(inout)::this
    if (associated(this%right)) then
       call this%right%finalize()
       deallocate(this%right)
    end if
    this%dim=0
  end subroutine muli_trapezium_list_finalize
  
\eatline
\nwendcode{}\nwbegindocs{1379}\nwdocspar
\nwenddocs{}\nwbegincode{1380}\moddef{Muli trapezium: list: TBP}\plusendmoddef
  generic :: insert_right => insert_right_a   !, insert_right_b    
  procedure :: insert_right_a => muli_trapezium_list_insert_right_a
  ! procedure :: insert_right_b => muli_trapezium_list_insert_right_b  
\nwendcode{}\nwbegincode{1381}\moddef{Muli trapezium: procedures}\plusendmoddef
  subroutine muli_trapezium_list_insert_right_a (this, value, content, new_node)
    class(muli_trapezium_list_t), intent(inout), target :: this
    real(default), intent(in) :: value
    class(muli_trapezium_t), intent(in) :: content
    class(muli_trapezium_list_t), pointer, intent(out) :: new_node
    class(muli_trapezium_list_t), pointer :: tmp_list
    call content%to_node (value, list=tmp_list)
    if (associated (this%right)) then
       this%right%left => tmp_list
       tmp_list%right => this%right
    else
       nullify (tmp_list%right)
    end if
    this%right => tmp_list
    tmp_list%left => this
    new_node => tmp_list
  end subroutine muli_trapezium_list_insert_right_a
  
\eatline
\nwendcode{}\nwbegindocs{1382}\nwdocspar
\nwenddocs{}\nwbegincode{1383}\moddef{Muli trapezium: list: TBP}\plusendmoddef
  generic :: insert_left => insert_left_a   !, insert_left_b  
  procedure :: insert_left_a => muli_trapezium_list_insert_left_a
  ! procedure :: insert_left_b => muli_trapezium_list_insert_left_b  
\nwendcode{}\nwbegindocs{1384}\nwdocspar
\nwenddocs{}\nwbegincode{1385}\moddef{Muli trapezium: procedures}\plusendmoddef
  subroutine muli_trapezium_list_insert_left_a (this, value, content, new_node)
    class(muli_trapezium_list_t), intent(inout), target :: this
    real(default), intent(in) :: value
    class(muli_trapezium_t), intent(in) :: content
    class(muli_trapezium_list_t), pointer, intent(out) :: new_node
    call content%to_node (value, list=new_node)
    new_node%right => this
    if (associated (this%left)) then
       new_node%left => this%left
       this%left%right => new_node
    else
       nullify (new_node%left)
    end if
    this%left => new_node
  end subroutine muli_trapezium_list_insert_left_a

\eatline
\nwendcode{}\nwbegindocs{1386}\nwdocspar
\nwenddocs{}\nwbegincode{1387}\moddef{Muli trapezium: list: TBP}\plusendmoddef
  procedure :: to_tree => muli_trapezium_list_to_tree
\nwendcode{}\nwbegincode{1388}\moddef{Muli trapezium: procedures}\plusendmoddef
  subroutine muli_trapezium_list_to_tree (this, out_tree)
    class(muli_trapezium_list_t), target, intent(in) :: this
    class(muli_trapezium_tree_t), intent(out) :: out_tree
    type(muli_trapezium_tree_t),target :: do_list
    class(muli_trapezium_node_class_t),pointer :: this_entry,do_list_entry,node
    class(muli_trapezium_tree_t),pointer :: tree1,tree2
    integer :: ite,log,n_deep,n_leaves
    n_leaves=0
    this_entry => this
    count: do while(associated(this_entry))
       n_leaves=n_leaves+1
       this_entry=>this_entry%right
    end do count
    call ilog2(n_leaves,log,n_deep)
    this_entry => this
    do_list_entry => do_list
    deep: do ite=0,n_deep-1
       allocate(tree1)
       tree1%down=>this_entry%right
       allocate(tree2)
       tree2%down=>this_entry
       tree2%left=>this_entry
       tree2%right=>this_entry%right
       tree1%left=>tree2
       this_entry => this_entry%right%right
       do_list_entry%right=>tree1
       do_list_entry=>tree1
    end do deep
    rest: do while(associated(this_entry))
       allocate(tree1)
       tree1%down=>this_entry
       tree1%left=>this_entry
       do_list_entry%right => tree1
       do_list_entry => tree1
       this_entry => this_entry%right
       ite=ite+1
    end do rest
    tree: do while(ite>2)
       do_list_entry => do_list%right
       node=>do_list
       level: do while(associated(do_list_entry))
          node%right=>do_list_entry%right
          node=>do_list_entry%right
          do_list_entry%right=>node%left
          node%left=>do_list_entry
          do_list_entry=>node%right
          ite=ite-1
       end do level
    end do tree
    node=>do_list%right
    select type(node)
    type is (muli_trapezium_tree_t)
       call node%to_tree(out_tree)
    class default
       print *,"muli_trapezium_list_to_tree"
       print *,"unexpeted type for do_list%right"
    end select
    out_tree%right=>out_tree%right%left
    if (allocated(out_tree%values)) then
       deallocate(out_tree%values)
    end if
    deallocate(do_list%right%right)
    deallocate(do_list%right)
  end subroutine muli_trapezium_list_to_tree
  
\eatline
\nwendcode{}\nwbegindocs{1389}\nwdocspar
\nwenddocs{}\nwbegincode{1390}\moddef{Muli trapezium: list: TBP}\plusendmoddef
  procedure :: gnuplot => muli_trapezium_list_gnuplot  
\nwendcode{}\nwbegindocs{1391}\nwdocspar
\nwenddocs{}\nwbegincode{1392}\moddef{Muli trapezium: procedures}\plusendmoddef
  subroutine muli_trapezium_list_gnuplot (this, dir)
    class(muli_trapezium_list_t), intent(in), target :: this
    character(len=*), intent(in) :: dir
    character(len=*), parameter :: val_file = "/value.plot"
    character(len=*), parameter :: int_file = "/integral.plot"
    character(len=*), parameter :: err_file = "/integral_error.plot"
    character(len=*), parameter :: pro_file = "/probability.plot"
    character(len=*), parameter :: den_file = "/density.plot"
    character(len=*), parameter :: fmt = "(ES20.10)"
    class(muli_trapezium_node_class_t), pointer :: list
    integer :: val_unit, err_unit, int_unit, pro_unit, den_unit
    list => this
    call generate_unit (val_unit, 100, 1000)
    open (val_unit, file = dir // val_file)
    call generate_unit (int_unit, 100, 1000)
    open (int_unit, file = dir // int_file)
    call generate_unit (err_unit, 100, 1000)
    open (err_unit, file = dir // err_file)
    call generate_unit (pro_unit, 100, 1000)
    open (pro_unit, file = dir // pro_file)
    call generate_unit (den_unit, 100, 1000)
    open (den_unit, file = dir // den_file)
    do while (associated (list))
       ! print *,list%r_position,list%get_r_value()
       !!! !!! !!! gfortran 5.0.0 ICE       
       ! write (val_unit, fmt, advance="no")  list%r_position
       ! call write_array (val_unit, list%get_r_value_array(), fmt)
       ! write (int_unit,fmt,advance="no")  list%r_position
       ! call write_array (int_unit, list%get_r_integral(), fmt)
       ! write (err_unit, fmt, advance="no")  list%r_position
       ! call write_array (err_unit, list%get_error(), fmt)
       ! write (pro_unit, fmt, advance="no")  list%r_position
       ! call write_array (pro_unit, list%get_r_probability(), fmt)
       ! write (den_unit, fmt, advance="no")  list%r_position
       ! call write_array (den_unit, list%get_r_probability() * &
       !      list%get_r_value_array(), fmt)
       list => list%right
    end do
    close (val_unit)
    close (int_unit)
    close (err_unit)
    close (pro_unit)
    close (den_unit)
    contains
      subroutine write_array (unit, array, form)
        integer, intent(in) :: unit
        real(default), dimension(:), intent(in) :: array
        character(len=*), intent(in) :: form
        integer :: n
        do n = 1, size(array)
           write (unit, form, advance="no")  array(n)
           flush (unit)
        end do
        write (unit, *)
      end subroutine write_array
  end subroutine muli_trapezium_list_gnuplot

\nwendcode{}\nwbegindocs{1393} %def muli_trapezium_list_gnuplot  
\nwenddocs{}\nwbegindocs{1394}\nwdocspar
\nwenddocs{}\nwbegincode{1395}\moddef{Muli trapezium: list: TBP}\plusendmoddef
  procedure :: integrate => muli_trapezium_list_integrate  
\nwendcode{}\nwbegincode{1396}\moddef{Muli trapezium: procedures}\plusendmoddef
  subroutine muli_trapezium_list_integrate (this, integral_sum, error_sum)
    class(muli_trapezium_list_t), intent(in), target :: this
    real(default), intent(out) :: error_sum, integral_sum
    real(default), dimension(:), allocatable :: integral
    class(muli_trapezium_node_class_t), pointer :: node
    allocate (integral (0:this%dim-1))
    call this%get_rightmost (node)
    integral = 0._default
    integral_sum = 0._default
    error_sum = 0._default
    integrate: do while (associated (node))
       node%values(1,r_value_index) = sum(node%values(1:this%dim-1,r_value_index))
       node%values(1,d_value_index) = sum(node%values(1:this%dim-1,d_value_index))
       ! node%values (1, r_integral_index) = &
       !    sum (node%values (1:this%dim-1, r_integral_index))
       ! node%values (1, d_integral_index) = &
       !    sum (node%values (1:this%dim-1, d_integral_index))
       node%values(1, error_index) = sum (node%values(1:this%dim-1, error_index))
       error_sum = error_sum + node%values (1, error_index)
       !!! !!! !!! gfortran 5.0.0 ICE
       ! call node%set_d_integral (node%get_d_position() * &
       !      (node%get_d_value() / 2 - node%get_r_value_array ()))
       call node%set_r_probability (exp (-integral))
       call node%set_r_integral (integral)
       !!! !!! !!! gfortran 5.0.0 ICE       
       ! integral = integral - node%get_d_integral()
       ! call node%set_d_probability (node%get_r_probability() - exp(-integral))
       ! call muli_trapezium_write (node, output_unit)
       call node%get_left (node)
    end do integrate
    integral_sum = integral (1)
  end subroutine muli_trapezium_list_integrate
  
\eatline
\nwendcode{}\nwbegindocs{1397}\nwdocspar
\nwenddocs{}\nwbegincode{1398}\moddef{Muli trapezium: list: TBP}\plusendmoddef
  procedure :: check => muli_trapezium_list_check  
\nwendcode{}\nwbegincode{1399}\moddef{Muli trapezium: procedures}\plusendmoddef
  recursive subroutine muli_trapezium_list_check (this)
    class(muli_trapezium_list_t), intent(in),target :: this
    class(muli_trapezium_node_class_t), pointer :: tn, next
    real(default), parameter :: eps = 1E-10_default
    logical::test
    if (associated(this%right)) then
       next=>this%right
       test=(this%r_position.le.this%right%get_l_position()+eps)
       print *,"position check:  ",test
       if (.not.test) then
          call this%print_parents()
          call next%print_parents()
       end if
       select type (next)
       class is (muli_trapezium_list_t)
          tn=>this
          print *,"structure check: ",associated(tn,next%left)
          print *,"class check:    T"
          call next%check()
       class default
          print *,"class check:    F"
       end select
    else
       print *,"end of list at ",this%r_position
    end if
  end subroutine muli_trapezium_list_check
  
\eatline
\nwendcode{}\nwbegindocs{1400}\nwdocspar
\nwenddocs{}\nwbegincode{1401}\moddef{Muli trapezium: list: TBP}\plusendmoddef
  procedure :: apply => muli_trapezium_list_apply  
\nwendcode{}\nwbegincode{1402}\moddef{Muli trapezium: procedures}\plusendmoddef
  recursive subroutine muli_trapezium_list_apply (this, proc)
    class(muli_trapezium_list_t), intent(inout) :: this
    interface
       subroutine proc (this)
         import muli_trapezium_node_class_t
         class(muli_trapezium_node_class_t), intent(inout) :: this
       end subroutine proc
    end interface
    if (associated (this%right))call this%right%apply (proc)
    call proc (this)
  end subroutine muli_trapezium_list_apply
  
\eatline
\nwendcode{}\nwbegindocs{1403}\nwdocspar
\nwenddocs{}\nwbegincode{1404}\moddef{Muli trapezium: procedures}\plusendmoddef
!  subroutine muli_trapezium_list_insert_right_old &
!       (this, value, content, new_node)
!    class(muli_trapezium_list_t), intent(inout), target :: this
!    real(default), intent(in) :: value
!    class(muli_trapezium_t), intent(in) :: content
!    class(muli_trapezium_list_t), pointer, intent(out) :: new_node
!    call content%to_node (value, list=new_node)
!    new_node%left => this
!    if (associated (this%right)) then
!       new_node%right => this%right
!       this%right%left => new_node
!    else
!       nullify (new_node%right)
!    end if
!    this%right => new_node
!  end subroutine muli_trapezium_list_insert_right_old

\eatline
\nwendcode{}\nwbegindocs{1405}\nwdocspar
\nwenddocs{}\nwbegincode{1406}\moddef{Muli trapezium: procedures}\plusendmoddef
!  subroutine muli_trapezium_node_error_no_content (this)
!    class(muli_trapezium_node_class_t), intent(in) :: this
!!    print ("muli_trapezium_node: Trying to access unallocated content.")
!!    call this%print()
!  end subroutine muli_trapezium_node_error_no_content

\eatline
\nwendcode{}\nwbegindocs{1407}\nwdocspar
%%%%%%%%%%%%%%%%%%%%%%%%%%%%%%%%%%%%%%%%%%%%%%%%%%%%%%%%%%%%%%%%%%%%%%%%
\section{Fibonacci trees}

This file contains the module {\tt{}muli{\char95}fibonacci{\char95}tree}. A fibonacci
tree is a self-balancing binary tree. "Balanced" means that the depth
of the left successor may differ from the depth of the right successor
in any node by maximally one. The maximally unbalanced tree is maximally
unbalanced in every node. So the number of leaves of a maximally
unbalanced fibonacci tree of depth $n$ is precisely the $n$th fibonacci
number. 

We use this tree for the adaptive quadrature. In every iteration of the
algorithm, we have to pick the segment with largest integration error,
cut it into pieces and calculate the new integration error for the pieces.
A binary tree is obvioulsly a good choice for soring these segments. The 
problem is, that we always pick leaves from the same side of the tree. So 
the tree must decline at one side and grow at the other side. Sorting gets
less efficient and finally most of the overall CPU time gets wasted for
sorting. This tree outperforms a naive binary tree significantly in
this particular job.

\nwenddocs{}\nwbegincode{1408}\moddef{File header}\plusendmoddef
\nwendcode{}\nwbegincode{1409}\moddef{\code{}muli{\char95}fibonacci{\char95}tree.f90\edoc{}}\endmoddef

module muli_fibonacci_tree
  
\LA{}Use kinds\RA{}
  use diagnostics
  use muli_base

\LA{}Standard module head\RA{}  

\LA{}Muli fibonacci: public\RA{}

\LA{}Muli fibonacci: variables\RA{}

\LA{}Muli fibonacci: types\RA{}

contains

\LA{}Muli fibonacci: procedures\RA{}  
   
end module muli_fibonacci_tree

\eatline
\nwendcode{}\nwbegindocs{1410}\nwdocspar
\nwenddocs{}\nwbegincode{1411}\moddef{Muli fibonacci: variables}\endmoddef
  character(*), parameter :: no_par = "edge=\\noparent"
  character(*), parameter :: no_ret = "edge=\\noreturn"
  character(*), parameter :: no_kid = "edge=\\nochild"
  character(*), parameter :: le_kid = "edge=\\childofleave"
    
\eatline
\nwendcode{}\nwbegindocs{1412}\nwdocspar
\nwenddocs{}\nwbegincode{1413}\moddef{Muli fibonacci: public}\endmoddef
  public :: fibonacci_node_t
\nwendcode{}\nwbegincode{1414}\moddef{Muli fibonacci: types}\endmoddef
  type, extends (measure_class_t) :: fibonacci_node_t
     ! private
     class(fibonacci_node_t), pointer :: up => null()
     class(measure_class_t), pointer :: down => null()
     class(fibonacci_node_t), pointer :: left => null()
     class(fibonacci_node_t), pointer :: right => null()
     integer :: depth = 0
     ! real(default) :: value
   contains
   \LA{}Muli fibonacci: fibonacci node: TBP\RA{}
  end type fibonacci_node_t
  
\eatline
\nwendcode{}\nwbegindocs{1415}\nwdocspar
\nwenddocs{}\nwbegincode{1416}\moddef{Muli fibonacci: fibonacci node: TBP}\endmoddef
  procedure :: write_to_marker => fibonacci_node_write_to_marker
\nwendcode{}\nwbegincode{1417}\moddef{Muli fibonacci: procedures}\endmoddef
  recursive  subroutine fibonacci_node_write_to_marker (this, marker, status)
    class(fibonacci_node_t), intent(in) :: this
    class(marker_t), intent(inout) :: marker
    integer(dik), intent(out) :: status
    class(ser_class_t), pointer :: ser
    call marker%mark_begin ("fibonacci_node_t")
    ser => this%left
    call marker%mark_pointer ("left", ser)
    ser => this%right
    call marker%mark_pointer ("right", ser)
    ser => this%down
    call marker%mark_pointer ("down", ser)
    call marker%mark_end ("fibonacci_node_t")
  end subroutine fibonacci_node_write_to_marker

\eatline
\nwendcode{}\nwbegindocs{1418}\nwdocspar
\nwenddocs{}\nwbegincode{1419}\moddef{Muli fibonacci: fibonacci node: TBP}\plusendmoddef
  procedure :: read_from_marker => fibonacci_node_read_from_marker
\nwendcode{}\nwbegincode{1420}\moddef{Muli fibonacci: procedures}\plusendmoddef
  recursive subroutine fibonacci_node_read_from_marker (this, marker, status)
    class(fibonacci_node_t), intent(out) :: this
    class(marker_t), intent(inout) :: marker
    integer(dik), intent(out) :: status
    call msg_warning ("fibonacci_node_read_from_marker: You cannot " // &
         "deserialize a list with this subroutine.")
    call msg_error ("Use fibonacci_node_read_target_from_marker instead.")
  end subroutine fibonacci_node_read_from_marker

\eatline
\nwendcode{}\nwbegindocs{1421}\nwdocspar
\nwenddocs{}\nwbegincode{1422}\moddef{Muli fibonacci: fibonacci node: TBP}\plusendmoddef
  procedure :: read_target_from_marker => fibonacci_node_read_target_from_marker
\nwendcode{}\nwbegincode{1423}\moddef{Muli fibonacci: procedures}\plusendmoddef
  recursive subroutine fibonacci_node_read_target_from_marker &
       (this, marker, status)
    class(fibonacci_node_t), target, intent(out) :: this
    class(marker_t), intent(inout) :: marker
    integer(dik), intent(out) :: status
    class(ser_class_t), pointer :: ser
    call marker%pick_begin ("fibonacci_node_t", status=status)
    call marker%pick_pointer ("left", ser)
    if (status == 0) then
       select type (ser)
       class is (fibonacci_node_t)
          this%left => ser
          this%left%up => this
       end select
    end if
    call marker%pick_pointer ("right", ser)
    if (status == 0) then
       select type (ser)
       class is (fibonacci_node_t)
          this%right => ser
          this%right%up => this
       end select
    end if
    call marker%pick_pointer ("down", ser)
    if (status == 0) then
       select type (ser)
       class is (measure_class_t)
          this%down => ser
       end select
    end if
    call marker%pick_end ("fibonacci_node_t", status)
  end subroutine fibonacci_node_read_target_from_marker

\eatline
\nwendcode{}\nwbegindocs{1424}\nwdocspar
\nwenddocs{}\nwbegincode{1425}\moddef{Muli fibonacci: fibonacci node: TBP}\plusendmoddef
  procedure :: print_to_unit => fibonacci_node_print_to_unit  
\nwendcode{}\nwbegincode{1426}\moddef{Muli fibonacci: procedures}\plusendmoddef
  recursive subroutine fibonacci_node_print_to_unit &
       (this, unit, parents, components, peers)
    class(fibonacci_node_t), intent(in) :: this
    integer, intent(in) :: unit
    integer(dik), intent(in) :: parents, components, peers
    class(ser_class_t), pointer :: ser
    write (unit, "(1x,A)")        "Components of fibonacci_node_t:"
    write (unit, "(3x,A,I22)")    "Depth:   ", this%depth
    write (unit, "(3x,A,E23.16)") "Value:   ", this%measure ()
    ser => this%up
    call serialize_print_comp_pointer &
         (ser, unit, parents, -i_one, -i_one, "Up:     ")
    ser => this%left
    call serialize_print_peer_pointer &
         (ser, unit, parents, components, peers, "Left:   ")
    ser => this%right
    call serialize_print_peer_pointer &
         (ser, unit, parents, components, peers, "Right:  ")
  end subroutine fibonacci_node_print_to_unit  

\eatline
\nwendcode{}\nwbegindocs{1427}\nwdocspar
\nwenddocs{}\nwbegincode{1428}\moddef{Muli fibonacci: fibonacci node: TBP}\plusendmoddef
  procedure, nopass :: get_type => fibonacci_node_get_type  
\nwendcode{}\nwbegincode{1429}\moddef{Muli fibonacci: procedures}\plusendmoddef
  pure subroutine fibonacci_node_get_type (type)
    character(:), allocatable, intent(out) :: type
    allocate (type, source="fibonacci_node_t")
  end subroutine fibonacci_node_get_type
    
\eatline
\nwendcode{}\nwbegindocs{1430}\nwdocspar
\nwenddocs{}\nwbegincode{1431}\moddef{Muli fibonacci: fibonacci node: TBP}\plusendmoddef
  procedure :: deserialize_from_marker => fibonacci_node_deserialize_from_marker
\nwendcode{}\nwbegincode{1432}\moddef{Muli fibonacci: procedures}\plusendmoddef
  subroutine fibonacci_node_deserialize_from_marker (this, name, marker)
    class(fibonacci_node_t), intent(out) :: this
    character(*), intent(in) :: name
    class(marker_t), intent(inout) :: marker
    class(ser_class_t), pointer :: ser
    allocate (fibonacci_leave_t :: ser)
    call marker%push_reference (ser)
    allocate (fibonacci_node_t :: ser)
    call marker%push_reference (ser)
    call serializable_deserialize_from_marker (this, name, marker)
    call marker%pop_reference (ser)
    deallocate (ser)
    call marker%pop_reference (ser)
    deallocate (ser)    
  end subroutine fibonacci_node_deserialize_from_marker
  
\eatline
\nwendcode{}\nwbegindocs{1433}\nwdocspar
\nwenddocs{}\nwbegincode{1434}\moddef{Muli fibonacci: fibonacci node: TBP}\plusendmoddef
  procedure :: measure => fibonacci_node_measure
\nwendcode{}\nwbegincode{1435}\moddef{Muli fibonacci: procedures}\plusendmoddef
  elemental function fibonacci_node_measure (this)
    class(fibonacci_node_t), intent(in) :: this
    real(default) :: fibonacci_node_measure
    fibonacci_node_measure = this%down%measure ()
  end function fibonacci_node_measure
  
\eatline
\nwendcode{}\nwbegindocs{1436}\nwdocspar
\nwenddocs{}\nwbegincode{1437}\moddef{Muli fibonacci: fibonacci node: TBP}\plusendmoddef
  procedure :: deallocate_tree => fibonacci_node_deallocate_tree
\nwendcode{}\nwbegincode{1438}\moddef{Muli fibonacci: procedures}\plusendmoddef
  recursive subroutine fibonacci_node_deallocate_tree (this)
    class(fibonacci_node_t), intent(inout) :: this
    if (associated (this%left)) then
       call this%left%deallocate_tree ()
       deallocate (this%left)
    end if
    if (associated (this%right)) then
       call this%right%deallocate_tree ()
       deallocate (this%right)
    end if
    call this%set_depth (0)
  end subroutine fibonacci_node_deallocate_tree

\eatline
\nwendcode{}\nwbegindocs{1439} 
\nwenddocs{}\nwbegincode{1440}\moddef{Muli fibonacci: fibonacci node: TBP}\plusendmoddef
  procedure :: deallocate_all => fibonacci_node_deallocate_all  
\nwendcode{}\nwbegincode{1441}\moddef{Muli fibonacci: procedures}\plusendmoddef
  recursive subroutine fibonacci_node_deallocate_all (this)
    class(fibonacci_node_t), intent(inout) :: this
    if (associated (this%left)) then
       call this%left%deallocate_all ()
       deallocate (this%left)
    end if
    if (associated (this%right)) then
       call this%right%deallocate_all ()
       deallocate (this%right)
    end if
    call this%set_depth (0)
  end subroutine fibonacci_node_deallocate_all
  
\eatline
\nwendcode{}\nwbegindocs{1442}\nwdocspar
\nwenddocs{}\nwbegincode{1443}\moddef{Muli fibonacci: fibonacci node: TBP}\plusendmoddef
  procedure :: get_depth => fibonacci_node_get_depth
\nwendcode{}\nwbegincode{1444}\moddef{Muli fibonacci: procedures}\plusendmoddef
  elemental function fibonacci_node_get_depth (this)
    class(fibonacci_node_t), intent(in) :: this
    integer :: fibonacci_node_get_depth
    fibonacci_node_get_depth = this%depth
  end function fibonacci_node_get_depth
  
\eatline
\nwendcode{}\nwbegindocs{1445}\nwdocspar
\nwenddocs{}\nwbegincode{1446}\moddef{Muli fibonacci: fibonacci node: TBP}\plusendmoddef
  procedure :: count_leaves => fibonacci_node_count_leaves
\nwendcode{}\nwbegincode{1447}\moddef{Muli fibonacci: procedures}\plusendmoddef
  recursive subroutine fibonacci_node_count_leaves (this, n)
    class(fibonacci_node_t), intent(in) :: this
    integer, intent(out) :: n
    integer :: n1, n2
    if (associated (this%left) .and. associated (this%right)) then
       call fibonacci_node_count_leaves (this%left, n1)
       call fibonacci_node_count_leaves (this%right, n2)
       n = n1 + n2
    else
       n = 1
    end if
  end subroutine fibonacci_node_count_leaves

\eatline
\nwendcode{}\nwbegindocs{1448}\nwdocspar
\nwenddocs{}\nwbegincode{1449}\moddef{Muli fibonacci: fibonacci node: TBP}\plusendmoddef
  procedure,public,nopass :: is_leave => fibonacci_node_is_leave
\nwendcode{}\nwbegincode{1450}\moddef{Muli fibonacci: procedures}\plusendmoddef
  elemental function fibonacci_node_is_leave ()
    logical :: fibonacci_node_is_leave
    fibonacci_node_is_leave = .false.
  end function fibonacci_node_is_leave

\eatline
\nwendcode{}\nwbegindocs{1451}\nwdocspar
\nwenddocs{}\nwbegincode{1452}\moddef{Muli fibonacci: fibonacci node: TBP}\plusendmoddef
  procedure,public,nopass :: is_root => fibonacci_node_is_root
\nwendcode{}\nwbegincode{1453}\moddef{Muli fibonacci: procedures}\plusendmoddef
  elemental function fibonacci_node_is_root ()
    logical :: fibonacci_node_is_root
    fibonacci_node_is_root = .false.
  end function fibonacci_node_is_root

\eatline
\nwendcode{}\nwbegindocs{1454}\nwdocspar
\nwenddocs{}\nwbegincode{1455}\moddef{Muli fibonacci: fibonacci node: TBP}\plusendmoddef
  procedure,public,nopass :: is_inner => fibonacci_node_is_inner  
\nwendcode{}\nwbegincode{1456}\moddef{Muli fibonacci: procedures}\plusendmoddef
  elemental function fibonacci_node_is_inner ()
    logical :: fibonacci_node_is_inner
    fibonacci_node_is_inner = .true.
  end function fibonacci_node_is_inner
    
\eatline
\nwendcode{}\nwbegindocs{1457}\nwdocspar
\nwenddocs{}\nwbegincode{1458}\moddef{Muli fibonacci: fibonacci node: TBP}\plusendmoddef
  procedure :: write_association => fibonacci_node_write_association
\nwendcode{}\nwbegincode{1459}\moddef{Muli fibonacci: procedures}\plusendmoddef
  subroutine fibonacci_node_write_association (this, that)
    class(fibonacci_node_t), intent(in), target :: this
    class(fibonacci_node_t), intent(in), target :: that
    if (associated (that%left, this)) then
       write(*, "(A)")  "This is left child of that"
    end if
    if (associated (that%right, this)) then
       write(*, "(A)")  "This is right child of that"
    end if
    if (associated (that%up, this)) then
       write(*, "(A)")  "This is parent of that"
    end if    
    if (associated (this%left, that)) then
       write(*, "(A)")  "That is left child of this"
    end if
    if (associated (this%right, that)) then
       write(*, "(A)")  "That is right child of this"
    end if
    if (associated (this%up, that)) then
       write(*, "(A)")  "That is parent of this"
    end if
  end subroutine fibonacci_node_write_association
  
\eatline
\nwendcode{}\nwbegindocs{1460}\nwdocspar
\nwenddocs{}\nwbegincode{1461}\moddef{Muli fibonacci: fibonacci node: TBP}\plusendmoddef
  procedure :: write_contents => fibonacci_node_write_contents
\nwendcode{}\nwbegincode{1462}\moddef{Muli fibonacci: procedures}\plusendmoddef
  subroutine fibonacci_node_write_contents (this, unit)
    class(fibonacci_node_t), intent(in), target :: this
    integer, intent(in), optional :: unit
    call this%apply_to_leaves (fibonacci_leave_write_content, unit)
  end subroutine fibonacci_node_write_contents
  
\eatline
\nwendcode{}\nwbegindocs{1463}\nwdocspar
\nwenddocs{}\nwbegincode{1464}\moddef{Muli fibonacci: fibonacci node: TBP}\plusendmoddef
  procedure :: write_values => fibonacci_node_write_values
\nwendcode{}\nwbegincode{1465}\moddef{Muli fibonacci: procedures}\plusendmoddef
  subroutine fibonacci_node_write_values (this, unit)
    class(fibonacci_node_t), intent(in), target :: this
    integer, intent(in), optional :: unit
    call this%apply_to_leaves (fibonacci_leave_write_value, unit)
  end subroutine fibonacci_node_write_values
  
\eatline
\nwendcode{}\nwbegindocs{1466}\nwdocspar
\nwenddocs{}\nwbegincode{1467}\moddef{Muli fibonacci: fibonacci node: TBP}\plusendmoddef
  procedure :: write_leaves => fibonacci_node_write_leaves
\nwendcode{}\nwbegincode{1468}\moddef{Muli fibonacci: procedures}\plusendmoddef
  subroutine fibonacci_node_write_leaves (this, unit)
    class(fibonacci_node_t), intent(in), target :: this
    integer, intent(in),optional :: unit
    call this%apply_to_leaves (fibonacci_leave_write, unit)
  end subroutine fibonacci_node_write_leaves
    
\eatline
\nwendcode{}\nwbegindocs{1469}\nwdocspar
\nwenddocs{}\nwbegincode{1470}\moddef{Muli fibonacci: fibonacci node: TBP}\plusendmoddef
  ! procedure :: write => fibonacci_node_write_contents  
\nwendcode{}\nwbegindocs{1471}\nwdocspar
\nwenddocs{}\nwbegincode{1472}\moddef{Muli fibonacci: fibonacci node: TBP}\plusendmoddef
  procedure :: write_pstricks => fibonacci_node_write_pstricks  
\nwendcode{}\nwbegincode{1473}\moddef{Muli fibonacci: procedures}\plusendmoddef
  recursive subroutine fibonacci_node_write_pstricks (this, unitnr)
    class(fibonacci_node_t), intent(in), target :: this
    integer, intent(in) :: unitnr
    if (associated (this%up)) then
       if (associated (this%up%left, this) .neqv. &
          (associated (this%up%right, this))) then
          ! write (unitnr,'("\\begin\{psTree\}\{\\Toval\{$",i3,"$\}\}")') int(this%depth)
          write (unitnr, &
               '("\\begin\{psTree\}\{\\Toval\{\\node\{",i3,"\}\{",f9.3,"\}\}\}")') &
               int(this%depth), this%measure()
       else
          write (unitnr, &
               '("\\begin\{psTree\}\{\\Toval[",a,"]\{\\node\{",i3,"\}\{",f9.3,"\}\}\}")') &
               no_ret, int(this%depth), this%measure()
       end if
    else
       write (unitnr, &
            '("\\begin\{psTree\}\{\\Toval[",a,"]\{\\node\{",i3,"\}\{",f9.3,"\}\}\}")') &
            no_par, int(this%depth), this%measure()
    end if
    if (associated (this%left)) then
       call this%left%write_pstricks (unitnr)
    else
       write (unitnr,'("\\Tr[edge=brokenline]\{\}")')
    end if
    if (associated (this%right)) then
       call this%right%write_pstricks (unitnr)
    else
       write (unitnr, '("\\Tr[edge=brokenline]\{\}")')
    end if
    write (unitnr, '("\\end\{psTree\}")')
  end subroutine fibonacci_node_write_pstricks
  
\eatline
\nwendcode{}\nwbegindocs{1474}\nwdocspar
\nwenddocs{}\nwbegincode{1475}\moddef{Muli fibonacci: fibonacci node: TBP}\plusendmoddef
  procedure :: copy_node => fibonacci_node_copy_node
\nwendcode{}\nwbegincode{1476}\moddef{Muli fibonacci: procedures}\plusendmoddef
  subroutine fibonacci_node_copy_node (this, primitive)
    class(fibonacci_node_t), intent(out) :: this
    class(fibonacci_node_t), intent(in) :: primitive
    this%up => primitive%up
    this%left => primitive%left
    this%right => primitive%right
    this%depth = primitive%depth
    this%down => primitive%down
  end subroutine fibonacci_node_copy_node
  
\eatline
\nwendcode{}\nwbegindocs{1477}\nwdocspar
\nwenddocs{}\nwbegincode{1478}\moddef{Muli fibonacci: fibonacci node: TBP}\plusendmoddef
  procedure :: find_root => fibonacci_node_find_root  
\nwendcode{}\nwbegincode{1479}\moddef{Muli fibonacci: procedures}\plusendmoddef
  subroutine fibonacci_node_find_root (this, root)
    class(fibonacci_node_t), intent(in), target :: this
    class(fibonacci_root_t), pointer, intent(out) :: root
    class(fibonacci_node_t), pointer :: node
    node => this
    do while (associated (node%up))
       node => node%up
    end do
    select type (node)
    class is (fibonacci_root_t)
       root => node
    class default
       nullify (root)
       call msg_error ("fibonacci_node_find_root: root is not type " // &
            "compatible to fibonacci_root_t. Retured NULL().")
    end select
  end subroutine fibonacci_node_find_root
  
\eatline
\nwendcode{}\nwbegindocs{1480}\nwdocspar
\nwenddocs{}\nwbegincode{1481}\moddef{Muli fibonacci: fibonacci node: TBP}\plusendmoddef
  procedure :: find_leftmost => fibonacci_node_find_leftmost
\nwendcode{}\nwbegincode{1482}\moddef{Muli fibonacci: procedures}\plusendmoddef
  subroutine fibonacci_node_find_leftmost (this, leave)
    class(fibonacci_node_t), intent(in), target :: this
    class(fibonacci_leave_t), pointer, intent(out) :: leave
    class(fibonacci_node_t),  pointer :: node
    node => this
    do while (associated (node%left))
       node => node%left
    end do
    select type (node)
    class is (fibonacci_leave_t)
       leave => node
    class default
       leave => null()
    end select
  end subroutine fibonacci_node_find_leftmost
  
\eatline
\nwendcode{}\nwbegindocs{1483}\nwdocspar
\nwenddocs{}\nwbegincode{1484}\moddef{Muli fibonacci: fibonacci node: TBP}\plusendmoddef
  procedure :: find_rightmost => fibonacci_node_find_rightmost
\nwendcode{}\nwbegincode{1485}\moddef{Muli fibonacci: procedures}\plusendmoddef
  subroutine fibonacci_node_find_rightmost (this, leave)
    class(fibonacci_node_t), intent(in), target :: this
    class(fibonacci_leave_t), pointer, intent(out) :: leave
    class(fibonacci_node_t), pointer :: node
    node => this
    do while (associated (node%right))
       node => node%right
    end do
    select type (node)
    class is (fibonacci_leave_t)
       leave => node
    class default
       leave => null()
    end select
  end subroutine fibonacci_node_find_rightmost
  
\eatline
\nwendcode{}\nwbegindocs{1486}\nwdocspar
\nwenddocs{}\nwbegincode{1487}\moddef{Muli fibonacci: fibonacci node: TBP}\plusendmoddef
  procedure :: find => fibonacci_node_find
\nwendcode{}\nwbegincode{1488}\moddef{Muli fibonacci: procedures}\plusendmoddef
  subroutine fibonacci_node_find (this, value, leave)
    class(fibonacci_node_t), intent(in), target :: this
    real(default), intent(in) :: value
    class(fibonacci_leave_t), pointer, intent(out) :: leave
    class(fibonacci_node_t), pointer :: node
    node => this
    do
       if (node >= value) then
          if (associated (node%left)) then
             node => node%left
          else
             call msg_warning ("fibonacci_node_find: broken tree!")
             leave => null()
             return
          end if
       else
          if (associated (node%right)) then
             node => node%right
          else
             call msg_warning ("fibonacci_node_find: broken tree!")
             leave => null()
             return
          end if
       end if
       select type (node)
       class is (fibonacci_leave_t)
          leave => node
          exit
       end select
    end do
  end subroutine fibonacci_node_find
  
\eatline
\nwendcode{}\nwbegindocs{1489}\nwdocspar
\nwenddocs{}\nwbegincode{1490}\moddef{Muli fibonacci: fibonacci node: TBP}\plusendmoddef
  procedure :: find_left_leave => fibonacci_node_find_left_leave
\nwendcode{}\nwbegincode{1491}\moddef{Muli fibonacci: procedures}\plusendmoddef
  subroutine fibonacci_node_find_left_leave (this, leave)
    class(fibonacci_node_t), intent(in), target :: this
    class(fibonacci_node_t), pointer :: node
    class(fibonacci_leave_t), pointer, intent(out) :: leave
    nullify(leave)
    node => this
    do while (associated (node%up))
       if (associated (node%up%right, node)) then
          node => node%up%left
          do while (associated (node%right))
             node => node%right
          end do
          select type (node)
          class is (fibonacci_leave_t)
          leave => node
          end select
          exit
       end if
       node => node%up
    end do
  end subroutine fibonacci_node_find_left_leave
  
\eatline
\nwendcode{}\nwbegindocs{1492}\nwdocspar
\nwenddocs{}\nwbegincode{1493}\moddef{Muli fibonacci: fibonacci node: TBP}\plusendmoddef
  procedure :: find_right_leave => fibonacci_node_find_right_leave
\nwendcode{}\nwbegincode{1494}\moddef{Muli fibonacci: procedures}\plusendmoddef
  subroutine fibonacci_node_find_right_leave (this, leave)
    class(fibonacci_node_t), intent(in), target :: this
    class(fibonacci_node_t), pointer :: node
    class(fibonacci_leave_t), pointer, intent(out) :: leave
    nullify (leave)
    node => this
    do while (associated (node%up))
       if (associated (node%up%left, node)) then
          node => node%up%right
          do while (associated (node%left))
             node => node%left
          end do
          select type (node)
          class is (fibonacci_leave_t)
          leave => node
          end select
          exit
       end if
       node => node%up
    end do
  end subroutine fibonacci_node_find_right_leave

\eatline
\nwendcode{}\nwbegindocs{1495}\nwdocspar
\nwenddocs{}\nwbegincode{1496}\moddef{Muli fibonacci: fibonacci node: TBP}\plusendmoddef
  procedure :: apply_to_leaves => fibonacci_node_apply_to_leaves 
\nwendcode{}\nwbegincode{1497}\moddef{Muli fibonacci: procedures}\plusendmoddef
  recursive subroutine fibonacci_node_apply_to_leaves (node, func, unit)
    class(fibonacci_node_t), intent(in), target :: node
    interface
       subroutine func (this, unit)
         import fibonacci_leave_t
         class(fibonacci_leave_t), intent(in), target :: this
         integer, intent(in), optional :: unit
       end subroutine func
    end interface
    integer, intent(in), optional :: unit
    select type (node)
    class is (fibonacci_leave_t)
       call func (node, unit)
    class default 
       call node%left%apply_to_leaves (func, unit)
       call node%right%apply_to_leaves (func, unit)
    end select
  end subroutine fibonacci_node_apply_to_leaves
  
\eatline
\nwendcode{}\nwbegindocs{1498}\nwdocspar
\nwenddocs{}\nwbegincode{1499}\moddef{Muli fibonacci: fibonacci node: TBP}\plusendmoddef
  procedure :: apply_to_leaves_rl => fibonacci_node_apply_to_leaves_rl  
\nwendcode{}\nwbegincode{1500}\moddef{Muli fibonacci: procedures}\plusendmoddef
  recursive subroutine fibonacci_node_apply_to_leaves_rl (node, func, unit)
    class(fibonacci_node_t), intent(in), target :: node
    interface
       subroutine func (this, unit)
         import fibonacci_leave_t
         class(fibonacci_leave_t), intent(in), target :: this
         integer, intent(in), optional :: unit
       end subroutine func
    end interface
    integer, intent(in), optional :: unit
    select type (node)
    class is (fibonacci_leave_t)
       call func (node, unit)
    class default 
       call node%right%apply_to_leaves_rl (func, unit)
       call node%left%apply_to_leaves_rl (func, unit)
    end select
  end subroutine fibonacci_node_apply_to_leaves_rl
  
\eatline
\nwendcode{}\nwbegindocs{1501}\nwdocspar
\nwenddocs{}\nwbegincode{1502}\moddef{Muli fibonacci: fibonacci node: TBP}\plusendmoddef
  procedure :: set_depth => fibonacci_node_set_depth
\nwendcode{}\nwbegindocs{1503}\nwdocspar
\nwenddocs{}\nwbegincode{1504}\moddef{Muli fibonacci: procedures}\plusendmoddef
  subroutine fibonacci_node_set_depth (this, depth)
    class(fibonacci_node_t), intent(inout) :: this
    integer, intent(in) :: depth
    this%depth = depth
  end subroutine fibonacci_node_set_depth
  
\eatline
\nwendcode{}\nwbegindocs{1505}\nwdocspar
\nwenddocs{}\nwbegincode{1506}\moddef{Muli fibonacci: fibonacci node: TBP}\plusendmoddef
  procedure :: append_left => fibonacci_node_append_left
\nwendcode{}\nwbegincode{1507}\moddef{Muli fibonacci: procedures}\plusendmoddef
  subroutine fibonacci_node_append_left(this,new_branch)
    class(fibonacci_node_t),target :: this
    class(fibonacci_node_t),target :: new_branch
    this%left => new_branch
    new_branch%up => this
  end subroutine fibonacci_node_append_left
    
\eatline
\nwendcode{}\nwbegindocs{1508}\nwdocspar
\nwenddocs{}\nwbegincode{1509}\moddef{Muli fibonacci: fibonacci node: TBP}\plusendmoddef
  procedure :: append_right => fibonacci_node_append_right
\nwendcode{}\nwbegincode{1510}\moddef{Muli fibonacci: procedures}\plusendmoddef
  subroutine fibonacci_node_append_right (this, new_branch)
    class(fibonacci_node_t), intent(inout), target :: this
    class(fibonacci_node_t), target :: new_branch
    this%right => new_branch
    new_branch%up => this
  end subroutine fibonacci_node_append_right
  
\eatline
\nwendcode{}\nwbegindocs{1511}\nwdocspar
\nwenddocs{}\nwbegincode{1512}\moddef{Muli fibonacci: fibonacci node: TBP}\plusendmoddef
  procedure :: replace => fibonacci_node_replace
\nwendcode{}\nwbegincode{1513}\moddef{Muli fibonacci: procedures}\plusendmoddef
  subroutine fibonacci_node_replace (this, old_node)
    class(fibonacci_node_t), intent(inout), target :: this
    class(fibonacci_node_t), target :: old_node
    if (associated (old_node%up)) then
       if (old_node%is_left_child ()) then
          old_node%up%left => this
       else
          if (old_node%is_right_child ()) then
             old_node%up%right => this
          end if
       end if
       this%up => old_node%up
    else
       nullify (this%up)
    end if
  end subroutine fibonacci_node_replace
  
\eatline
\nwendcode{}\nwbegincode{1514}\moddef{Muli fibonacci: fibonacci node: TBP}\plusendmoddef
  procedure :: swap => fibonacci_node_swap_nodes
\nwendcode{}\nwbegincode{1515}\moddef{Muli fibonacci: procedures}\plusendmoddef
  subroutine fibonacci_node_swap_nodes (left, right)
    class(fibonacci_node_t), target, intent(inout) :: left, right
    class(fibonacci_node_t), pointer :: left_left, right_right
    class(measure_class_t), pointer::down
    ! swap branches
    left_left => left%left
    right_right => right%right
    left%left => right%right
    right%right => left_left
    ! repair up components
    right_right%up => left
    left_left%up => right
    ! repair down components
    down => left%down
    left%down => right%down
    right%down => down
  end subroutine fibonacci_node_swap_nodes

!  subroutine fibonacci_node_swap_nodes (this, that)
!    class(fibonacci_node_t),target :: this
!    class(fibonacci_node_t), pointer, intent(in) :: that
!    class(fibonacci_node_t), pointer :: par_i, par_a
!    par_i => this%up
!    par_a => that%up
!    if (associated (par_i%left, this)) then
!       par_i%left => that
!    else
!       par_i%right => that
!    end if
!    if (associated (par_a%left, that)) then
!       par_a%left => this
!    else
!       par_a%right => this
!    end if
!    this%up => par_a
!    that%up => par_i
!  end subroutine fibonacci_node_swap_nodes
  
\nwendcode{}\nwbegindocs{1516}\nwdocspar
\nwenddocs{}\nwbegincode{1517}\moddef{Muli fibonacci: fibonacci node: TBP}\plusendmoddef
  procedure :: flip => fibonacci_node_flip_children
\nwendcode{}\nwbegincode{1518}\moddef{Muli fibonacci: procedures}\plusendmoddef
  subroutine fibonacci_node_flip_children (this)
    class(fibonacci_node_t), intent(inout) :: this
    class(fibonacci_node_t), pointer :: child
    child => this%left
    this%left => this%right
    this%right => child
  end subroutine fibonacci_node_flip_children
  
\eatline
\nwendcode{}\nwbegindocs{1519}\nwdocspar
\nwenddocs{}\nwbegincode{1520}\moddef{Muli fibonacci: fibonacci node: TBP}\plusendmoddef
  procedure :: rip => fibonacci_node_rip
\nwendcode{}\nwbegincode{1521}\moddef{Muli fibonacci: procedures}\plusendmoddef
  subroutine fibonacci_node_rip (this)
    class(fibonacci_node_t), intent(inout), target :: this
    if (this%is_left_child ()) then
       nullify (this%up%left)
    end if
    if (this%is_right_child ()) then
       nullify (this%up%right)
    end if
    nullify (this%up)
  end subroutine fibonacci_node_rip
  
\eatline
\nwendcode{}\nwbegindocs{1522}\nwdocspar
\nwenddocs{}\nwbegincode{1523}\moddef{Muli fibonacci: fibonacci node: TBP}\plusendmoddef
  procedure :: remove_and_keep_parent => fibonacci_node_remove_and_keep_parent
\nwendcode{}\nwbegincode{1524}\moddef{Muli fibonacci: procedures}\plusendmoddef
  subroutine fibonacci_node_remove_and_keep_parent (this, pa)
    class(fibonacci_node_t), intent(inout), target :: this
    class(fibonacci_node_t), intent(out), pointer :: pa
    class(fibonacci_node_t), pointer :: twin
    if (.not. (this%is_root ())) then
       pa => this%up
       if (this%is_left_child ()) then         
          twin => pa%right
       else
          twin => pa%left
       end if
       twin%up => pa%up
       if (associated (twin%left)) then
          twin%left%up => pa
       end if
       if (associated (twin%right)) then
          twin%right%up => pa
       end if
       call pa%copy_node (twin)
       select type (pa)
       class is (fibonacci_root_t)
          call pa%set_leftmost ()
          call pa%set_rightmost ()
       end select
       if (associated (this%right)) then
          this%right%left => this%left
       end if
       if (associated (this%left)) then
          this%left%right => this%right
       end if
       nullify (this%left)
       nullify (this%right)
       nullify (this%up)
       deallocate (twin)
    else
       pa => this
    end if
  end subroutine fibonacci_node_remove_and_keep_parent
  
\eatline
\nwendcode{}\nwbegindocs{1525}\nwdocspar
\nwenddocs{}\nwbegincode{1526}\moddef{Muli fibonacci: fibonacci node: TBP}\plusendmoddef
  procedure :: remove_and_keep_twin => fibonacci_node_remove_and_keep_twin
\nwendcode{}\nwbegincode{1527}\moddef{Muli fibonacci: procedures}\plusendmoddef
  subroutine fibonacci_node_remove_and_keep_twin (this, twin)
    class(fibonacci_node_t), intent(inout), target :: this
    class(fibonacci_node_t), intent(out), pointer :: twin
    class(fibonacci_node_t), pointer :: pa
    if (.not. (this%is_root ())) then
       pa => this%up
       if (.not. pa%is_root ()) then
          if (this%is_left_child ()) then         
             twin => pa%right
          else
             twin => pa%left
          end if
          if (pa%is_left_child ()) then
             pa%up%left => twin
          else
             pa%up%right => twin
          end if
       end if
       twin%up => pa%up
       if (associated (this%right)) then
          this%right%left => this%left
       end if
       if (associated (this%left)) then
          this%left%right => this%right
       end if
       nullify (this%left)
       nullify (this%right)
       nullify (this%up)
       deallocate (pa)
    end if
  end subroutine fibonacci_node_remove_and_keep_twin
  
\eatline
\nwendcode{}\nwbegindocs{1528}\nwdocspar
\nwenddocs{}\nwbegincode{1529}\moddef{Muli fibonacci: fibonacci node: TBP}\plusendmoddef
  procedure :: rotate_left => fibonacci_node_rotate_left
\nwendcode{}\nwbegincode{1530}\moddef{Muli fibonacci: procedures}\plusendmoddef
  subroutine fibonacci_node_rotate_left (this)
    class(fibonacci_node_t), intent(inout), target :: this
    call this%swap (this%right)
    call this%right%flip ()
    call this%right%update_depth_unsave ()
    call this%flip ()
    ! value = this%value
    ! this%value = this%left%value
    ! this%left%value = value
  end subroutine fibonacci_node_rotate_left

\eatline
\nwendcode{}\nwbegindocs{1531} 
\nwenddocs{}\nwbegincode{1532}\moddef{Muli fibonacci: fibonacci node: TBP}\plusendmoddef
  procedure :: rotate_right => fibonacci_node_rotate_right
\nwendcode{}\nwbegincode{1533}\moddef{Muli fibonacci: procedures}\plusendmoddef
  subroutine fibonacci_node_rotate_right (this)
    class(fibonacci_node_t), intent(inout), target :: this
    call this%left%swap (this)
    call this%left%flip ()
    call this%left%update_depth_unsave ()
    call this%flip ()
    ! value = this%value
    ! this%value = this%right%value
    ! this%right%value = value
  end subroutine fibonacci_node_rotate_right  

\eatline
\nwendcode{}\nwbegindocs{1534}\nwdocspar
\nwenddocs{}\nwbegincode{1535}\moddef{Muli fibonacci: fibonacci node: TBP}\plusendmoddef
  procedure :: rotate => fibonacci_node_rotate
\nwendcode{}\nwbegincode{1536}\moddef{Muli fibonacci: procedures}\plusendmoddef
  subroutine fibonacci_node_rotate (this)
    class(fibonacci_node_t), intent(inout), target :: this
    if (this%is_left_short ()) then
       call this%rotate_left ()
    else
       if (this%is_right_short ()) then
          call this%rotate_right ()
       end if
    end if
  end subroutine fibonacci_node_rotate

\eatline
\nwendcode{}\nwbegindocs{1537}\nwdocspar
\nwenddocs{}\nwbegincode{1538}\moddef{Muli fibonacci: fibonacci node: TBP}\plusendmoddef
  procedure :: balance_node => fibonacci_node_balance_node
\nwendcode{}\nwbegincode{1539}\moddef{Muli fibonacci: procedures}\plusendmoddef
  subroutine fibonacci_node_balance_node (this, changed)
    class(fibonacci_node_t), intent(inout), target :: this
    logical, intent(out) :: changed
    changed = .false.
    if (this%is_left_too_short ()) then
       if (this%right%is_right_short ()) then
          call this%right%rotate_right
       end if
       call this%rotate_left ()
       changed = .true.
    else
       if (this%is_right_too_short ()) then
          if (this%left%is_left_short ()) then
             call this%left%rotate_left
          end if
          call this%rotate_right ()
          changed = .true.
       end if
    end if
  end subroutine fibonacci_node_balance_node

\eatline
\nwendcode{}\nwbegindocs{1540}\nwdocspar
\nwenddocs{}\nwbegincode{1541}\moddef{Muli fibonacci: fibonacci node: TBP}\plusendmoddef
  procedure :: update_depth_save => fibonacci_node_update_depth_save
\nwendcode{}\nwbegindocs{1542}\nwdocspar
\nwenddocs{}\nwbegincode{1543}\moddef{Muli fibonacci: procedures}\plusendmoddef
  subroutine fibonacci_node_update_depth_save (this, updated)
    class(fibonacci_node_t), intent(inout) :: this
    logical, intent(out) :: updated
    integer :: left, right, new_depth
    if (associated (this%left)) then
       left = this%left%depth + 1
    else
       left = -1
    end if
    if (associated (this%right)) then
       right = this%right%depth + 1
    else
       right = -1
    end if
    new_depth = max(left, right)
    if (this%depth == new_depth) then
       updated = .false.
    else
       this%depth = new_depth
       updated = .true.
    end if
  end subroutine fibonacci_node_update_depth_save

\eatline
\nwendcode{}\nwbegindocs{1544}\nwdocspar
\nwenddocs{}\nwbegincode{1545}\moddef{Muli fibonacci: fibonacci node: TBP}\plusendmoddef
  procedure :: update_depth_unsave => fibonacci_node_update_depth_unsave
\nwendcode{}\nwbegindocs{1546}\nwdocspar
\nwenddocs{}\nwbegincode{1547}\moddef{Muli fibonacci: procedures}\plusendmoddef
  subroutine fibonacci_node_update_depth_unsave (this)
    class(fibonacci_node_t), intent(inout) :: this
    this%depth = max (this%left%depth+1, this%right%depth+1)
  end subroutine fibonacci_node_update_depth_unsave
  
\eatline
\nwendcode{}\nwbegindocs{1548}\nwdocspar
\nwenddocs{}\nwbegincode{1549}\moddef{Muli fibonacci: fibonacci node: TBP}\plusendmoddef
  procedure :: repair => fibonacci_node_repair
\nwendcode{}\nwbegincode{1550}\moddef{Muli fibonacci: procedures}\plusendmoddef
  subroutine fibonacci_node_repair (this)
    class(fibonacci_node_t), intent(inout), target :: this
    class(fibonacci_node_t), pointer:: node
    logical :: new_depth, new_balance
    new_depth = .true.
    node => this
    do while ((new_depth .or. new_balance) .and. (associated (node)))
       call node%balance_node (new_balance)
       call node%update_depth_save (new_depth)
       node => node%up
    end do
  end subroutine fibonacci_node_repair

\eatline
\nwendcode{}\nwbegindocs{1551}\nwdocspar
\nwenddocs{}\nwbegincode{1552}\moddef{Muli fibonacci: fibonacci node: TBP}\plusendmoddef
  procedure :: is_left_short => fibonacci_node_is_left_short
\nwendcode{}\nwbegincode{1553}\moddef{Muli fibonacci: procedures}\plusendmoddef
  elemental logical function fibonacci_node_is_left_short(this)
    class(fibonacci_node_t), intent(in) :: this
    fibonacci_node_is_left_short = (this%left%depth < this%right%depth)
  end function fibonacci_node_is_left_short

\nwendcode{}\nwbegindocs{1554} %def fibonacci_node_is_left_short  
\nwenddocs{}\nwbegindocs{1555}\nwdocspar
\nwenddocs{}\nwbegincode{1556}\moddef{Muli fibonacci: fibonacci node: TBP}\plusendmoddef
  procedure :: is_right_short => fibonacci_node_is_right_short
\nwendcode{}\nwbegindocs{1557}\nwdocspar
\nwenddocs{}\nwbegincode{1558}\moddef{Muli fibonacci: procedures}\plusendmoddef
  elemental logical function fibonacci_node_is_right_short (this)
    class(fibonacci_node_t), intent(in) :: this
    fibonacci_node_is_right_short = (this%right%depth < this%left%depth)
  end function fibonacci_node_is_right_short
    
\eatline
\nwendcode{}\nwbegindocs{1559}\nwdocspar
\nwenddocs{}\nwbegincode{1560}\moddef{Muli fibonacci: fibonacci node: TBP}\plusendmoddef
  procedure :: is_unbalanced => fibonacci_node_is_unbalanced
\nwendcode{}\nwbegincode{1561}\moddef{Muli fibonacci: procedures}\plusendmoddef
  elemental logical function fibonacci_node_is_unbalanced (this)
    class(fibonacci_node_t), intent(in) :: this
    fibonacci_node_is_unbalanced = &
         (this%is_left_short () .or. this%is_right_short ())
  end function fibonacci_node_is_unbalanced
  
\eatline
\nwendcode{}\nwbegindocs{1562}\nwdocspar
\nwenddocs{}\nwbegincode{1563}\moddef{Muli fibonacci: fibonacci node: TBP}\plusendmoddef
  procedure :: is_left_too_short => fibonacci_node_is_left_too_short
\nwendcode{}\nwbegincode{1564}\moddef{Muli fibonacci: procedures}\plusendmoddef
  elemental logical function fibonacci_node_is_left_too_short (this)
    class(fibonacci_node_t), intent(in) :: this
    fibonacci_node_is_left_too_short = (this%left%depth+1 < this%right%depth)
  end function fibonacci_node_is_left_too_short

\eatline
\nwendcode{}\nwbegindocs{1565}\nwdocspar
\nwenddocs{}\nwbegincode{1566}\moddef{Muli fibonacci: fibonacci node: TBP}\plusendmoddef
  procedure :: is_right_too_short => fibonacci_node_is_right_too_short
\nwendcode{}\nwbegincode{1567}\moddef{Muli fibonacci: procedures}\plusendmoddef
  elemental logical function fibonacci_node_is_right_too_short (this)
    class(fibonacci_node_t), intent(in) :: this
    fibonacci_node_is_right_too_short = (this%right%depth+1 < this%left%depth)
  end function fibonacci_node_is_right_too_short

\nwendcode{}\nwbegindocs{1568} %def fibonacci_node_is_right_too_short  
\nwenddocs{}\nwbegindocs{1569}\nwdocspar
\nwenddocs{}\nwbegincode{1570}\moddef{Muli fibonacci: fibonacci node: TBP}\plusendmoddef
  procedure :: is_too_unbalanced => fibonacci_node_is_too_unbalanced
\nwendcode{}\nwbegincode{1571}\moddef{Muli fibonacci: procedures}\plusendmoddef
  elemental logical function fibonacci_node_is_too_unbalanced (this)
    class(fibonacci_node_t), intent(in) :: this
    fibonacci_node_is_too_unbalanced = &
         (this%is_left_too_short() .or. this%is_right_too_short())
  end function fibonacci_node_is_too_unbalanced
  
\eatline
\nwendcode{}\nwbegindocs{1572}\nwdocspar
\nwenddocs{}\nwbegincode{1573}\moddef{Muli fibonacci: fibonacci node: TBP}\plusendmoddef
  procedure :: is_left_child => fibonacci_node_is_left_child
\nwendcode{}\nwbegincode{1574}\moddef{Muli fibonacci: procedures}\plusendmoddef
  elemental logical function fibonacci_node_is_left_child (this)
    class(fibonacci_node_t), intent(in),target :: this
    fibonacci_node_is_left_child = associated (this%up%left, this)
  end function fibonacci_node_is_left_child
  
\nwendcode{}\nwbegindocs{1575} %def fibonacci_node_is_left_child  
\nwenddocs{}\nwbegindocs{1576}\nwdocspar
\nwenddocs{}\nwbegincode{1577}\moddef{Muli fibonacci: fibonacci node: TBP}\plusendmoddef
  procedure :: is_right_child => fibonacci_node_is_right_child  
\nwendcode{}\nwbegincode{1578}\moddef{Muli fibonacci: procedures}\plusendmoddef
  elemental logical function fibonacci_node_is_right_child (this)
    class(fibonacci_node_t), intent(in),target :: this
    fibonacci_node_is_right_child = associated (this%up%right, this)
  end function fibonacci_node_is_right_child
  
\eatline
\nwendcode{}\nwbegindocs{1579}\nwdocspar
\nwenddocs{}\nwbegincode{1580}\moddef{Muli fibonacci: fibonacci node: TBP}\plusendmoddef
  ! user
  ! node
  ! tree
  ! procedure :: balance
  ! procedure :: sort
  ! procedure :: merge
  ! procedure :: split  
\nwendcode{}\nwbegindocs{1581}\nwdocspar
\nwenddocs{}\nwbegincode{1582}\moddef{Muli fibonacci: public}\plusendmoddef
  public :: fibonacci_leave_t
\nwendcode{}\nwbegincode{1583}\moddef{Muli fibonacci: types}\plusendmoddef
  type, extends (fibonacci_node_t) :: fibonacci_leave_t
     ! class(measure_class_t), pointer :: content
  contains
  \LA{}Muli fibonacci: fibonacci leave: TBP\RA{}  
  end type fibonacci_leave_t
  
\eatline
\nwendcode{}\nwbegindocs{1584}\nwdocspar
\nwenddocs{}\nwbegincode{1585}\moddef{Muli fibonacci: fibonacci leave: TBP}\endmoddef
  ! procedure :: write_to_marker => fibonacci_leave_write_to_marker
  ! procedure :: read_from_marker => fibonacci_leave_read_from_marker
  procedure :: print_to_unit => fibonacci_leave_print_to_unit    
\nwendcode{}\nwbegincode{1586}\moddef{Muli fibonacci: procedures}\plusendmoddef
  subroutine fibonacci_leave_print_to_unit &
       (this, unit, parents, components, peers)
    class(fibonacci_leave_t), intent(in) :: this
    integer, intent(in) :: unit
    integer(dik), intent(in) :: parents, components, peers
    class(ser_class_t), pointer :: ser
    if (parents > 0)  call fibonacci_node_print_to_unit &
         (this, unit, parents-i_one, components, -i_one)
    write(unit, "(A)")  "Components of fibonacci_leave_t:"
    ser => this%down
    call serialize_print_comp_pointer &
         (ser, unit, parents, components, peers, "Content:")
  end subroutine fibonacci_leave_print_to_unit
  
\eatline
\nwendcode{}\nwbegindocs{1587}\nwdocspar
\nwenddocs{}\nwbegincode{1588}\moddef{Muli fibonacci: fibonacci leave: TBP}\plusendmoddef
  procedure, nopass :: get_type => fibonacci_leave_get_type
\nwendcode{}\nwbegincode{1589}\moddef{Muli fibonacci: procedures}\plusendmoddef
  pure subroutine fibonacci_leave_get_type (type)
    character(:), allocatable, intent(out) :: type
    allocate (type, source="fibonacci_leave_t")
  end subroutine fibonacci_leave_get_type

\eatline
\nwendcode{}\nwbegindocs{1590}\nwdocspar
\nwenddocs{}\nwbegincode{1591}\moddef{Muli fibonacci: fibonacci leave: TBP}\plusendmoddef
  procedure :: deallocate_all => fibonacci_leave_deallocate_all
\nwendcode{}\nwbegincode{1592}\moddef{Muli fibonacci: procedures}\plusendmoddef
  subroutine fibonacci_leave_deallocate_all (this)
    class(fibonacci_leave_t), intent(inout) :: this
    if (associated (this%down)) then
       deallocate (this%down)
    end if
  end subroutine fibonacci_leave_deallocate_all

\eatline
\nwendcode{}\nwbegindocs{1593}\nwdocspar
\nwenddocs{}\nwbegincode{1594}\moddef{Muli fibonacci: fibonacci leave: TBP}\plusendmoddef
  procedure :: pick => fibonacci_leave_pick
\nwendcode{}\nwbegindocs{1595}\nwdocspar
\nwenddocs{}\nwbegincode{1596}\moddef{Muli fibonacci: procedures}\plusendmoddef
  subroutine fibonacci_leave_pick (this)
    class(fibonacci_leave_t), target, intent(inout) :: this
    class(fibonacci_node_t), pointer :: other
    class(fibonacci_root_t), pointer :: root
    ! call this%up%print_parents()
    call this%find_root (root)
    if (associated (this%up, root)) then
       if (this%up%depth < 2) then
          call msg_error ("fibonacci_leave_pick: Cannot pick leave. " // &
               "Tree must have at least three leaves.")
       else
          call this%remove_and_keep_parent (other)
          call other%repair ()
       end if
    else
       call this%remove_and_keep_twin (other)
       call other%up%repair ()
    end if
    if (associated (root%leftmost, this))  call root%set_leftmost ()
    if (associated (root%rightmost, this))  call root%set_rightmost ()
  end subroutine fibonacci_leave_pick

\eatline
\nwendcode{}\nwbegindocs{1597}\nwdocspar
\nwenddocs{}\nwbegincode{1598}\moddef{Muli fibonacci: fibonacci leave: TBP}\plusendmoddef
  procedure :: get_left => fibonacci_leave_get_left
\nwendcode{}\nwbegindocs{1599}\nwdocspar
\nwenddocs{}\nwbegincode{1600}\moddef{Muli fibonacci: procedures}\plusendmoddef
  subroutine fibonacci_leave_get_left (this, leave)
    class(fibonacci_leave_t), intent(in) :: this
    class(fibonacci_leave_t), intent(out), pointer :: leave
    class(fibonacci_node_t), pointer :: node
    node => this%left
    select type (node)
    class is (fibonacci_leave_t)
       leave => node
    end select
  end subroutine fibonacci_leave_get_left
    
\eatline
\nwendcode{}\nwbegindocs{1601}\nwdocspar
\nwenddocs{}\nwbegincode{1602}\moddef{Muli fibonacci: fibonacci leave: TBP}\plusendmoddef
  procedure :: get_right => fibonacci_leave_get_right
\nwendcode{}\nwbegincode{1603}\moddef{Muli fibonacci: procedures}\plusendmoddef
  subroutine fibonacci_leave_get_right (this, leave)
    class(fibonacci_leave_t), intent(in) :: this
    class(fibonacci_leave_t), intent(out), pointer :: leave
    class(fibonacci_node_t), pointer :: node
    ! print *,"fibonacci_leave_get_right"
    ! call this%down%print_little
    if (associated (this%right)) then
       node => this%right
       ! call node%down%print_little
       select type (node)
       class is (fibonacci_leave_t)
          leave => node
       end select
    else
       ! print *,"no right leave"
       nullify (leave)
    end if
  end subroutine fibonacci_leave_get_right
  
\eatline
\nwendcode{}\nwbegindocs{1604}\nwdocspar
\nwenddocs{}\nwbegincode{1605}\moddef{Muli fibonacci: fibonacci leave: TBP}\plusendmoddef
  procedure :: write_pstricks => fibonacci_leave_write_pstricks
\nwendcode{}\nwbegincode{1606}\moddef{Muli fibonacci: procedures}\plusendmoddef
  subroutine fibonacci_leave_write_pstricks (this, unitnr)
    class(fibonacci_leave_t), intent(in), target :: this
    integer, intent(in) :: unitnr
    write (unitnr, "(A,I3,A,F9.3,A)")  &
         "\\begin\{psTree\}\{\\Toval[linecolor=green]\{\\node\{", this%depth, "\}\{", &
         this%measure(), "\}\}\}"
    if (associated (this%left)) then
       write (unitnr, "(A,A,A)")  "\\Tr[", le_kid, "]\{\}"
    end if
    if (associated (this%right)) then
       write (unitnr, "(A,A,A)")  "\\Tr[", le_kid, "]\{\}"
    end if
    write (unitnr, "(A)")  "\\end\{psTree\}"
  end subroutine fibonacci_leave_write_pstricks
  
\eatline
\nwendcode{}\nwbegindocs{1607}\nwdocspar
\nwenddocs{}\nwbegincode{1608}\moddef{Muli fibonacci: fibonacci leave: TBP}\plusendmoddef
  procedure :: copy_content => fibonacci_leave_copy_content
\nwendcode{}\nwbegincode{1609}\moddef{Muli fibonacci: procedures}\plusendmoddef
  subroutine fibonacci_leave_copy_content (this, content)
    class(fibonacci_leave_t) :: this
    class(measure_class_t), intent(in) :: content
    allocate (this%down, source=content)
  end subroutine fibonacci_leave_copy_content
  
\eatline
\nwendcode{}\nwbegindocs{1610}\nwdocspar
\nwenddocs{}\nwbegincode{1611}\moddef{Muli fibonacci: fibonacci leave: TBP}\plusendmoddef
  procedure :: set_content => fibonacci_leave_set_content
\nwendcode{}\nwbegincode{1612}\moddef{Muli fibonacci: procedures}\plusendmoddef
  subroutine fibonacci_leave_set_content (this, content)
    class(fibonacci_leave_t) :: this
    class(measure_class_t), target, intent(in) :: content
    this%down => content
  end subroutine fibonacci_leave_set_content
  
\eatline
\nwendcode{}\nwbegindocs{1613}\nwdocspar
\nwenddocs{}\nwbegincode{1614}\moddef{Muli fibonacci: fibonacci leave: TBP}\plusendmoddef
  procedure :: get_content => fibonacci_leave_get_content
\nwendcode{}\nwbegincode{1615}\moddef{Muli fibonacci: procedures}\plusendmoddef
  subroutine fibonacci_leave_get_content (this, content)
    class(fibonacci_leave_t), intent(in) :: this
    class(measure_class_t), pointer :: content
    content => this%down
  end subroutine fibonacci_leave_get_content
  
\eatline
\nwendcode{}\nwbegindocs{1616}\nwdocspar
\nwenddocs{}\nwbegincode{1617}\moddef{Muli fibonacci: fibonacci leave: TBP}\plusendmoddef
  procedure, nopass :: is_inner => fibonacci_leave_is_inner
\nwendcode{}\nwbegincode{1618}\moddef{Muli fibonacci: procedures}\plusendmoddef
  elemental logical function fibonacci_leave_is_inner ()
    fibonacci_leave_is_inner = .false.
  end function fibonacci_leave_is_inner

\eatline
\nwendcode{}\nwbegindocs{1619}\nwdocspar
\nwenddocs{}\nwbegincode{1620}\moddef{Muli fibonacci: fibonacci leave: TBP}\plusendmoddef
  procedure, nopass :: is_leave => fibonacci_leave_is_leave
\nwendcode{}\nwbegincode{1621}\moddef{Muli fibonacci: procedures}\plusendmoddef
  elemental logical function fibonacci_leave_is_leave ()
    fibonacci_leave_is_leave = .true.
  end function fibonacci_leave_is_leave
  
\eatline
\nwendcode{}\nwbegindocs{1622}\nwdocspar
\nwenddocs{}\nwbegincode{1623}\moddef{Muli fibonacci: fibonacci leave: TBP}\plusendmoddef
  procedure :: insert_leave_by_node => fibonacci_leave_insert_leave_by_node
\nwendcode{}\nwbegincode{1624}\moddef{Muli fibonacci: procedures}\plusendmoddef
  subroutine fibonacci_leave_insert_leave_by_node (this, new_leave)
    class(fibonacci_leave_t), target, intent(inout) :: this,new_leave
    class(fibonacci_node_t), pointer :: parent, new_node
    parent => this%up
    !print *, associated (this%left), associated (this%right)
    if (this < new_leave) then
       call fibonacci_node_spawn (new_node, this, new_leave, this%left, this%right)
       ! print *,"Repair! ",this%measure(),new_leave%measure()
    else
       call fibonacci_node_spawn (new_node, new_leave, this, this%left, this%right)
    end if
    if (associated (parent%left, this)) then
       call parent%append_left (new_node)
    else
       call parent%append_right (new_node)
    end if
    call parent%repair ()
  end subroutine fibonacci_leave_insert_leave_by_node

\eatline
\nwendcode{}\nwbegindocs{1625}\nwdocspar
\nwenddocs{}\nwbegincode{1626}\moddef{Muli fibonacci: fibonacci leave: TBP}\plusendmoddef
  procedure :: is_left_short => fibonacci_leave_is_left_short
\nwendcode{}\nwbegincode{1627}\moddef{Muli fibonacci: procedures}\plusendmoddef
  elemental logical function fibonacci_leave_is_left_short (this)
    class(fibonacci_leave_t), intent(in) :: this
    fibonacci_leave_is_left_short = .false.
  end function fibonacci_leave_is_left_short

\eatline
\nwendcode{}\nwbegindocs{1628}\nwdocspar
\nwenddocs{}\nwbegincode{1629}\moddef{Muli fibonacci: fibonacci leave: TBP}\plusendmoddef
  procedure :: is_right_short => fibonacci_leave_is_right_short
\nwendcode{}\nwbegincode{1630}\moddef{Muli fibonacci: procedures}\plusendmoddef
  elemental logical function fibonacci_leave_is_right_short (this)
    class(fibonacci_leave_t), intent(in) :: this
    fibonacci_leave_is_right_short = .false.
  end function fibonacci_leave_is_right_short

\eatline
\nwendcode{}\nwbegindocs{1631}\nwdocspar
\nwenddocs{}\nwbegincode{1632}\moddef{Muli fibonacci: fibonacci leave: TBP}\plusendmoddef
  procedure :: is_unbalanced => fibonacci_leave_is_unbalanced
\nwendcode{}\nwbegincode{1633}\moddef{Muli fibonacci: procedures}\plusendmoddef
  elemental logical function fibonacci_leave_is_unbalanced (this)
    class(fibonacci_leave_t), intent(in) :: this
    fibonacci_leave_is_unbalanced = .false.
  end function fibonacci_leave_is_unbalanced

\eatline
\nwendcode{}\nwbegindocs{1634}\nwdocspar
\nwenddocs{}\nwbegincode{1635}\moddef{Muli fibonacci: fibonacci leave: TBP}\plusendmoddef
  procedure :: is_left_too_short => fibonacci_leave_is_left_too_short
\nwendcode{}\nwbegincode{1636}\moddef{Muli fibonacci: procedures}\plusendmoddef
  elemental logical function fibonacci_leave_is_left_too_short (this)
    class(fibonacci_leave_t), intent(in) :: this
    fibonacci_leave_is_left_too_short = .false.
  end function fibonacci_leave_is_left_too_short

\eatline
\nwendcode{}\nwbegindocs{1637}\nwdocspar
\nwenddocs{}\nwbegincode{1638}\moddef{Muli fibonacci: fibonacci leave: TBP}\plusendmoddef
  procedure :: is_right_too_short => fibonacci_leave_is_right_too_short
\nwendcode{}\nwbegincode{1639}\moddef{Muli fibonacci: procedures}\plusendmoddef
  elemental logical function fibonacci_leave_is_right_too_short (this)
    class(fibonacci_leave_t), intent(in) :: this
    fibonacci_leave_is_right_too_short = .false.
  end function fibonacci_leave_is_right_too_short

\eatline
\nwendcode{}\nwbegindocs{1640}\nwdocspar
\nwenddocs{}\nwbegincode{1641}\moddef{Muli fibonacci: fibonacci leave: TBP}\plusendmoddef
  procedure :: is_too_unbalanced => fibonacci_leave_is_too_unbalanced  
\nwendcode{}\nwbegincode{1642}\moddef{Muli fibonacci: procedures}\plusendmoddef
  elemental logical function fibonacci_leave_is_too_unbalanced (this)
    class(fibonacci_leave_t), intent(in) :: this
    fibonacci_leave_is_too_unbalanced = .false.
  end function fibonacci_leave_is_too_unbalanced

\eatline
\nwendcode{}\nwbegindocs{1643}\nwdocspar
\nwenddocs{}\nwbegincode{1644}\moddef{Muli fibonacci: public}\plusendmoddef
  public :: fibonacci_root_t
\nwendcode{}\nwbegincode{1645}\moddef{Muli fibonacci: types}\plusendmoddef
  type, extends (fibonacci_node_t) :: fibonacci_root_t
     logical::is_valid_c=.false.
     class(fibonacci_leave_t),pointer :: leftmost => null()
     class(fibonacci_leave_t),pointer :: rightmost => null()
  contains
  \LA{}Muli fibonacci: fibonacci root: TBP\RA{}  
     procedure :: is_left_child => fibonacci_root_is_left_child
     procedure :: is_right_child => fibonacci_root_is_right_child
  end type fibonacci_root_t
  
\eatline
\nwendcode{}\nwbegindocs{1646}\nwdocspar
\nwenddocs{}\nwbegincode{1647}\moddef{Muli fibonacci: fibonacci root: TBP}\endmoddef
  procedure :: write_to_marker => fibonacci_root_write_to_marker
\nwendcode{}\nwbegincode{1648}\moddef{Muli fibonacci: procedures}\plusendmoddef
  subroutine fibonacci_root_write_to_marker (this, marker, status)
    class(fibonacci_root_t), intent(in) :: this
    class(marker_t), intent(inout) :: marker
    integer(dik), intent(out) :: status
    ! call marker%mark_begin ("fibonacci_root_t")
    call fibonacci_node_write_to_marker (this, marker, status)
    ! marker%mark_end ("fibonacci_root_t")
  end subroutine fibonacci_root_write_to_marker
  
\nwendcode{}\nwbegincode{1649}\moddef{Muli fibonacci: fibonacci root: TBP}\plusendmoddef
  procedure :: read_target_from_marker => fibonacci_root_read_target_from_marker
\nwendcode{}\nwbegincode{1650}\moddef{Muli fibonacci: procedures}\plusendmoddef
  subroutine fibonacci_root_read_target_from_marker (this, marker, status)
    class(fibonacci_root_t), target, intent(out) :: this
    class(marker_t), intent(inout) :: marker
    integer(dik), intent(out) :: status
    ! call marker%pick_begin ("fibonacci_root_t", status)
    call fibonacci_node_read_from_marker (this, marker, status)
    call this%find_leftmost (this%leftmost)
    call this%find_rightmost (this%rightmost)
    ! call marker%pick_end ("fibonacci_root_t", status)
  end subroutine fibonacci_root_read_target_from_marker

\eatline
\nwendcode{}\nwbegindocs{1651}\nwdocspar
\nwenddocs{}\nwbegincode{1652}\moddef{Muli fibonacci: fibonacci root: TBP}\plusendmoddef
  procedure :: print_to_unit => fibonacci_root_print_to_unit  
\nwendcode{}\nwbegincode{1653}\moddef{Muli fibonacci: procedures}\plusendmoddef
  subroutine fibonacci_root_print_to_unit (this, unit, parents, components, peers)
    class(fibonacci_root_t), intent(in) :: this
    integer, intent(in) :: unit
    integer(dik), intent(in) :: parents, components, peers
    class(ser_class_t), pointer :: ser
    if (parents > 0)  call fibonacci_node_print_to_unit &
         (this, unit, parents-1, components, peers)
    write (unit, "(A)")  "Components of fibonacci_root_t:"
    ser => this%leftmost
    call serialize_print_peer_pointer &
         (ser, unit, parents, components, min(peers, i_one), "Leftmost: ")
    ser => this%rightmost
    call serialize_print_peer_pointer &
         (ser, unit, parents, components, min(peers, i_one), "Rightmost:")
  end subroutine fibonacci_root_print_to_unit
  
\eatline
\nwendcode{}\nwbegindocs{1654}\nwdocspar
\nwenddocs{}\nwbegincode{1655}\moddef{Muli fibonacci: procedures}\plusendmoddef
  elemental logical function fibonacci_root_is_left_child (this)
    class(fibonacci_root_t),target, intent(in) :: this
    fibonacci_root_is_left_child = .false.
  end function fibonacci_root_is_left_child
  
\eatline
\nwendcode{}\nwbegindocs{1656}\nwdocspar
\nwenddocs{}\nwbegincode{1657}\moddef{Muli fibonacci: procedures}\plusendmoddef
  elemental logical function fibonacci_root_is_right_child (this)
    class(fibonacci_root_t),target, intent(in) :: this
    fibonacci_root_is_right_child = .false.
  end function fibonacci_root_is_right_child
  
\eatline
\nwendcode{}\nwbegindocs{1658}\nwdocspar
\nwenddocs{}\nwbegincode{1659}\moddef{Muli fibonacci: fibonacci root: TBP}\plusendmoddef
  procedure, nopass :: get_type => fibonacci_root_get_type  
\nwendcode{}\nwbegindocs{1660}\nwdocspar
\nwenddocs{}\nwbegincode{1661}\moddef{Muli fibonacci: procedures}\plusendmoddef
  pure subroutine fibonacci_root_get_type (type)
    character(:),allocatable, intent(out) :: type
    allocate (type, source="fibonacci_root_t")
  end subroutine fibonacci_root_get_type

\eatline
\nwendcode{}\nwbegindocs{1662}\nwdocspar
\nwenddocs{}\nwbegincode{1663}\moddef{Muli fibonacci: fibonacci root: TBP}\plusendmoddef
  procedure :: get_leftmost=>fibonacci_root_get_leftmost
\nwendcode{}\nwbegincode{1664}\moddef{Muli fibonacci: procedures}\plusendmoddef
  subroutine fibonacci_root_get_leftmost (this, leftmost)
    class(fibonacci_root_t), intent(in) :: this
    class(fibonacci_leave_t), pointer :: leftmost
    leftmost => this%leftmost
  end subroutine fibonacci_root_get_leftmost

\eatline
\nwendcode{}\nwbegindocs{1665}\nwdocspar
\nwenddocs{}\nwbegincode{1666}\moddef{Muli fibonacci: fibonacci root: TBP}\plusendmoddef
  procedure :: get_rightmost=>fibonacci_root_get_rightmost  
\nwendcode{}\nwbegincode{1667}\moddef{Muli fibonacci: procedures}\plusendmoddef
  subroutine fibonacci_root_get_rightmost (this, rightmost)
    class(fibonacci_root_t), intent(in) :: this
    class(fibonacci_leave_t),pointer :: rightmost
    rightmost => this%rightmost
  end subroutine fibonacci_root_get_rightmost
    
\eatline
\nwendcode{}\nwbegindocs{1668}\nwdocspar
\nwenddocs{}\nwbegincode{1669}\moddef{Muli fibonacci: fibonacci root: TBP}\plusendmoddef
  procedure, nopass :: is_root => fibonacci_root_is_root
\nwendcode{}\nwbegincode{1670}\moddef{Muli fibonacci: procedures}\plusendmoddef
  elemental function fibonacci_root_is_root ()
    logical::fibonacci_root_is_root
    fibonacci_root_is_root = .true.
  end function fibonacci_root_is_root
  
\eatline
\nwendcode{}\nwbegindocs{1671}\nwdocspar
\nwenddocs{}\nwbegincode{1672}\moddef{Muli fibonacci: fibonacci root: TBP}\plusendmoddef
  procedure, nopass :: is_inner => fibonacci_root_is_inner   
\nwendcode{}\nwbegincode{1673}\moddef{Muli fibonacci: procedures}\plusendmoddef
  elemental function fibonacci_root_is_inner ()
    logical::fibonacci_root_is_inner
    fibonacci_root_is_inner = .false.
  end function fibonacci_root_is_inner
  
\eatline
\nwendcode{}\nwbegindocs{1674}\nwdocspar
\nwenddocs{}\nwbegincode{1675}\moddef{Muli fibonacci: fibonacci root: TBP}\plusendmoddef
  procedure :: is_valid => fibonacci_root_is_valid
\nwendcode{}\nwbegincode{1676}\moddef{Muli fibonacci: procedures}\plusendmoddef
  elemental function fibonacci_root_is_valid (this)
    class(fibonacci_root_t), intent(in) :: this
    logical :: fibonacci_root_is_valid
    fibonacci_root_is_valid = this%is_valid_c
  end function fibonacci_root_is_valid
  
\eatline
\nwendcode{}\nwbegindocs{1677}\nwdocspar
\nwenddocs{}\nwbegincode{1678}\moddef{Muli fibonacci: fibonacci root: TBP}\plusendmoddef
  procedure :: count_leaves => fibonacci_root_count_leaves 
\nwendcode{}\nwbegincode{1679}\moddef{Muli fibonacci: procedures}\plusendmoddef
  subroutine fibonacci_root_count_leaves (this, n)
    class(fibonacci_root_t), intent(in) :: this
    integer, intent(out) :: n
    n = 0
    call fibonacci_node_count_leaves (this, n)
  end subroutine fibonacci_root_count_leaves
  
\eatline
\nwendcode{}\nwbegindocs{1680}\nwdocspar
\nwenddocs{}\nwbegincode{1681}\moddef{Muli fibonacci: fibonacci root: TBP}\plusendmoddef
  procedure :: write_pstricks => fibonacci_root_write_pstricks
\nwendcode{}\nwbegincode{1682}\moddef{Muli fibonacci: procedures}\plusendmoddef
  subroutine fibonacci_root_write_pstricks (this, unitnr)
    class(fibonacci_root_t), intent(in), target :: this
    integer, intent(in) :: unitnr
    logical :: is_opened
    character :: is_sequential, is_formatted, is_writeable
    print *,"pstricks"
    inquire (unitnr, opened=is_opened, sequential=is_sequential, &
         formatted=is_formatted, write=is_writeable)
    if (is_opened) then
       if (is_sequential == "Y" .and. is_formatted == "Y " &
            .and. is_writeable == "Y") then
          ! write (unitnr, "(A,I3,A)")  &
          !     "\\begin\{psTree\}\{\\Toval[linecolor=blue]\{$", int(this%depth), "$\}\}"
          write (unitnr, "(A,I3,A,F9.3,A)")  &
               "\\begin\{psTree\}\{\\Toval[linecolor=blue]\{\\node\{", this%depth, &
               "\}\{", this%measure(), "\}\}\}"
          if (associated (this%leftmost)) then
             call this%leftmost%write_pstricks (unitnr)
          else
             write (unitnr, "(A,A,A)") "\\Tr[", no_kid, "]\{\}"
          end if
          if (associated (this%left)) then
             call this%left%write_pstricks (unitnr)
          else
             write (unitnr, "(A,A,A)") "\\Tr[", no_kid, "]\{\}"
          end if
          if (associated (this%right)) then
             call this%right%write_pstricks (unitnr)
          else
             write (unitnr, "(A,A,A)") "\\Tr[", no_kid, "]\{\}"
          end if
          if (associated (this%rightmost)) then
             call this%rightmost%write_pstricks (unitnr)
          else
             write(unitnr,'("\\Tr[",a,"]\{\}")') no_kid
          end if
          write (unitnr, "(A)")  "\\end\{psTree\}"
          write (unitnr, "(A)")  "\\\\"
       else
          write (*, "(A,I2,A)") &
             "fibonacci_node_write_pstricks: Unit ", unitnr, &
               " is not opened properly."
          write (*, "(A)")  "No output is written to unit."
       end if
    else
       write (*, "(A,I2,A)") &
          "fibonacci_node_write_pstricks: Unit ", unitnr, &
            " is not opened."
       write (*, "(A)")  "No output is written to unit."
    end if
  end subroutine fibonacci_root_write_pstricks
  
\eatline
\nwendcode{}\nwbegindocs{1683}\nwdocspar
\nwenddocs{}\nwbegincode{1684}\moddef{Muli fibonacci: fibonacci root: TBP}\plusendmoddef
  procedure :: copy_root => fibonacci_root_copy_root
\nwendcode{}\nwbegincode{1685}\moddef{Muli fibonacci: procedures}\plusendmoddef
  subroutine fibonacci_root_copy_root (this, primitive)
    class(fibonacci_root_t), intent(out) :: this
    class(fibonacci_root_t), intent(in) :: primitive
    call fibonacci_node_copy_node (this, primitive)
    this%leftmost => primitive%leftmost
    this%rightmost => primitive%rightmost
  end subroutine fibonacci_root_copy_root
  
\eatline
\nwendcode{}\nwbegindocs{1686}\nwdocspar
\nwenddocs{}\nwbegincode{1687}\moddef{Muli fibonacci: fibonacci root: TBP}\plusendmoddef
  procedure :: push_by_content => fibonacci_root_push_by_content
\nwendcode{}\nwbegincode{1688}\moddef{Muli fibonacci: procedures}\plusendmoddef
  subroutine fibonacci_root_push_by_content (this, content)
    class(fibonacci_root_t), target, intent(inout) :: this
    class(measure_class_t), target, intent(in) :: content
    class(fibonacci_leave_t), pointer :: node
    ! print *,"fibonacci_root_push_by_content: ",content%measure()
    allocate (node)
    node%down => content
    call this%push_by_leave (node)
  end subroutine fibonacci_root_push_by_content
    
\eatline
\nwendcode{}\nwbegindocs{1689}This is a workaround for gfortran bug 44696. This subroutine is a
merge of {\tt{}fibonacci{\char95}tree{\char95}push{\char95}by{\char95}node}, {\tt{}fibonacci{\char95}node{\char95}find}, and 
{\tt{}fibonacci{\char95}leave{\char95}insert{\char95}leave{\char95}by{\char95}node}.
\nwenddocs{}\nwbegincode{1690}\moddef{Muli fibonacci: fibonacci root: TBP}\plusendmoddef
  procedure :: push_by_leave => fibonacci_root_push_by_leave
\nwendcode{}\nwbegincode{1691}\moddef{Muli fibonacci: procedures}\plusendmoddef
  subroutine fibonacci_root_push_by_leave (this, new_leave)
    class(fibonacci_root_t), target, intent(inout) :: this
    class(fibonacci_leave_t), pointer, intent(inout) :: new_leave
    class(fibonacci_leave_t), pointer :: old_leave
    class(fibonacci_node_t), pointer :: node, new_node, leave_c
    ! write (11, fmt=*)  "push by leave(", new_leave%measure(), ")\\\\" !PSTRICKS
    ! flush(11)  !PSTRICKS
    if (new_leave <= this%leftmost) then
       old_leave => this%leftmost
       this%leftmost => new_leave
       node => old_leave%up
       call fibonacci_node_spawn &
            (new_node, new_leave, old_leave, old_leave%left, old_leave%right)
       call node%append_left (new_node)
    else
       if (new_leave > this%rightmost) then
          old_leave => this%rightmost
          this%rightmost => new_leave
          node => old_leave%up
          call fibonacci_node_spawn &
               (new_node, old_leave, new_leave, old_leave%left, old_leave%right)
          call node%append_right (new_node)
       else
          node => this
          do
             if (new_leave <= node) then
                leave_c => node%left
                select type (leave_c)
                class is (fibonacci_leave_t)
                   if (new_leave <= leave_c) then
                      ! print *,"left left"
                      call fibonacci_node_spawn (new_node, new_leave, &
                           leave_c, leave_c%left, leave_c%right)
                   else
                      ! print *,"left right"
                      call fibonacci_node_spawn (new_node, leave_c, &
                           new_leave, leave_c%left, leave_c%right)
                   end if
                   call node%append_left (new_node)
                   exit
                class default
                   ! print *,"left"
                   node => node%left
                end select
             else
                leave_c => node%right
                select type (leave_c)
                class is (fibonacci_leave_t)          
                   if (new_leave <= leave_c) then
                      ! print *,"right left"
                      call fibonacci_node_spawn (new_node, new_leave, &
                           leave_c, leave_c%left, leave_c%right)
                   else
                      ! print *,"right right"
                      call fibonacci_node_spawn (new_node, leave_c, &
                           new_leave, leave_c%left, leave_c%right)
                   end if
                   call node%append_right (new_node)
                   exit
                class default
                   ! print *,"right"
                   node => node%right
                end select
             end if
          end do
       end if
    end if
    ! call this%write_pstricks(11) ! PSTRICKS
    ! flush(11) ! PSTRICKS
    ! write(11,fmt=*)"repair\\\\" ! PSTRICKS
    call node%repair ()
    ! call this%write_pstricks (11) !PSTRICKS
    ! flush(11) !PSTRICKS
    ! call node%update_value (right_value)
    ! call this%write_pstricks (11)
    ! print *, new_node%value, new_node%left%value, new_node%right%value
  end subroutine fibonacci_root_push_by_leave
  
\eatline
\nwendcode{}\nwbegindocs{1692}\nwdocspar
\nwenddocs{}\nwbegincode{1693}\moddef{Muli fibonacci: fibonacci root: TBP}\plusendmoddef
  procedure :: pop_left => fibonacci_root_pop_left
\nwendcode{}\nwbegincode{1694}\moddef{Muli fibonacci: procedures}\plusendmoddef
  subroutine fibonacci_root_pop_left (this, leave)
    class(fibonacci_root_t), intent(inout), target :: this
    class(fibonacci_leave_t), pointer, intent(out) :: leave
    class(fibonacci_node_t), pointer :: parent, grand
    ! write (11,fmt=*) "fibonacci root pop left\\\\"  ! PSTRICKS
    ! flush (11)   ! PSTRICKS
    leave => this%leftmost
    if (this%left%depth >= 1) then
       parent => leave%up
       grand => parent%up
       grand%left => parent%right
       parent%right%up => grand
       deallocate (parent)
       parent => grand%left
       if (.not. parent%is_leave ()) then
          parent => parent%left
       end if
       select type (parent)
       class is (fibonacci_leave_t)
          this%leftmost => parent
       class default
          call parent%print_all()          
          call msg_fatal ("fibonacci_root_pop_left: ERROR: leftmost is no leave.")
       end select
       ! call this%write_pstricks (11)   ! PSTRICKS
       ! flush (11)   ! PSTRICKS
       ! write (11,fmt=*)  "fibonacci node repair\\\\"   ! PSTRICKS
       ! flush (11)   ! PSTRICKS
       call grand%repair ()
    else
       if (this%left%depth == 0 .and. this%right%depth == 1) then
          parent => this%right
          parent%right%up => this
          parent%left%up => this
          this%left => parent%left
          this%right => parent%right
          this%depth = 1 
          deallocate (parent)
          parent => this%left
          select type (parent)
          class is (fibonacci_leave_t)
          this%leftmost => parent
          end select
          this%down => this%leftmost%down
       end if
    end if
    nullify (leave%right%left)
    nullify (leave%up)
    nullify (leave%right)
    nullify (this%leftmost%left)
    ! call this%write_pstricks (11)   ! PSTRICKS
    ! flush (11)   ! PSTRICKS
  end subroutine fibonacci_root_pop_left
  
\eatline
\nwendcode{}\nwbegindocs{1695}\nwdocspar
\nwenddocs{}\nwbegincode{1696}\moddef{Muli fibonacci: fibonacci root: TBP}\plusendmoddef
  procedure :: pop_right => fibonacci_root_pop_right
\nwendcode{}\nwbegincode{1697}\moddef{Muli fibonacci: procedures}\plusendmoddef
  subroutine fibonacci_root_pop_right (this, leave)
    class(fibonacci_root_t), intent(inout), target :: this
    class(fibonacci_leave_t), pointer, intent(out) :: leave
    class(fibonacci_node_t), pointer :: parent, grand
    leave => this%rightmost
    if (this%right%depth >= 1) then
       parent => leave%up
       grand => parent%up
       grand%right => parent%left
       parent%left%up => grand
       deallocate (parent)
       parent => grand%right
       if (.not. parent%is_leave ()) then
          parent => parent%right
       end if
       select type (parent)
       class is (fibonacci_leave_t)
          this%rightmost => parent
       class default
          call parent%print_all()          
          call msg_fatal ("fibonacci_root_pop_left: ERROR: leftmost is no leave.")
       end select
       call grand%repair ()
    else
       if (this%right%depth == 0 .and. this%left%depth == 1) then
          parent => this%left
          parent%left%up => this
          parent%right%up => this
          this%right => parent%right
          this%left => parent%left
          this%depth = 1
          deallocate (parent)
          parent => this%right
          select type (parent)
          class is (fibonacci_leave_t)
          this%rightmost => parent
          end select
          this%down => this%rightmost%down
       end if
    end if
  end subroutine fibonacci_root_pop_right
  
\eatline
\nwendcode{}\nwbegindocs{1698}\nwdocspar
\nwenddocs{}\nwbegincode{1699}\moddef{Muli fibonacci: fibonacci root: TBP}\plusendmoddef
  procedure :: list_to_tree => fibonacci_root_list_to_tree
\nwendcode{}\nwbegincode{1700}\moddef{Muli fibonacci: procedures}\plusendmoddef
  subroutine fibonacci_root_list_to_tree (this, n_leaves, leave_list_target)
    class(fibonacci_root_t), target, intent(inout) :: this
    integer, intent(in) :: n_leaves
    type(fibonacci_leave_list_t), target, intent(in) :: leave_list_target
    ! class(fibonacci_root_t), pointer, intent(out) :: tree
    integer :: depth, n_deep, n_merge
    class(fibonacci_node_t), pointer :: node
    class(fibonacci_leave_list_t), pointer :: leave_list
    class(fibonacci_leave_t), pointer :: content
    real(default) :: up_value
    leave_list => leave_list_target
    call ilog2 (n_leaves, depth, n_deep)
    n_deep = n_deep * 2
    n_merge = 0
    this%depth = depth
    node => this
    OUTER: do
       do while (depth > 1)
          depth = depth - 1
          allocate (node%left)
          node%left%up => node
          node => node%left
          node%depth = depth
       end do
       node%left => leave_list%leave
       node%down => leave_list%leave%down
       leave_list => leave_list%next
       node%right => leave_list%leave
       content => leave_list%leave
       leave_list => leave_list%next
       n_merge = n_merge + 2
       INNER: do
          if (associated (node%up)) then
             if (node%is_left_child ()) then
                if (n_merge == n_deep .and. depth == 1) then
                   node => node%up
                   node%right => leave_list%leave
                   node%right%up => node
                   node%down => content%down
                   content => leave_list%leave
                   leave_list => leave_list%next
                   n_merge = n_merge + 1
                   cycle
                end if
                exit INNER
             else
                node => node%up
                depth = depth + 1
             end if
          else
             exit OUTER
          end if
       end do INNER
       node => node%up
       node%down => content%down
       allocate (node%right)
       node%right%up => node
       node => node%right
       if (n_deep == n_merge) then
          depth = depth - 1
       end if
       node%depth = depth
    end do OUTER
    call this%set_leftmost
    call this%set_rightmost
  end subroutine fibonacci_root_list_to_tree
  
\eatline
\nwendcode{}\nwbegindocs{1701} This subroutine has neither been used nor revised for a long time,
so it might be broken. 
\nwenddocs{}\nwbegincode{1702}\moddef{Muli fibonacci: fibonacci root: TBP}\plusendmoddef
  procedure :: merge => fibonacci_root_merge
\nwendcode{}\nwbegincode{1703}\moddef{Muli fibonacci: procedures}\plusendmoddef
  subroutine fibonacci_root_merge(this_tree,that_tree,merge_tree)
    class(fibonacci_root_t), intent(in) :: this_tree
    class(fibonacci_root_t), intent(in) :: that_tree
    class(fibonacci_root_t), pointer, intent(out) :: merge_tree
    class(fibonacci_leave_t), pointer :: this_leave, that_leave, old_leave
    type(fibonacci_leave_list_t), target :: leave_list
    class(fibonacci_leave_list_t), pointer :: last_leave
    integer :: n_leaves
    if (associated (this_tree%leftmost) .and. associated (that_tree%leftmost)) then
       n_leaves = 1
       this_leave => this_tree%leftmost
       that_leave => that_tree%leftmost
       if (this_leave < that_leave) then
          allocate (leave_list%leave, source=this_leave)
          call this_leave%find_right_leave (this_leave)
       else
          allocate (leave_list%leave, source=that_leave)
          call that_leave%find_right_leave (that_leave)
       end if
       last_leave => leave_list
       do while (associated (this_leave) .and. associated (that_leave))
          if (this_leave < that_leave) then
             old_leave => this_leave
             call this_leave%find_right_leave (this_leave)
          else
             old_leave=>that_leave
             call that_leave%find_right_leave (that_leave)
          end if
          allocate (last_leave%next)
          last_leave => last_leave%next
          allocate (last_leave%leave, source=old_leave)
          n_leaves = n_leaves + 1
       end do
       if (associated (this_leave)) then
          old_leave => this_leave
       else
          old_leave => that_leave
       end if
       do while (associated (old_leave))
          allocate (last_leave%next)
          last_leave => last_leave%next
          allocate (last_leave%leave, source=old_leave)
          n_leaves = n_leaves + 1
          call old_leave%find_right_leave (old_leave)
       end do
       allocate (merge_tree)
       call merge_tree%list_to_tree (n_leaves, leave_list)
    else
       n_leaves = 0
    end if
    if (associated (leave_list%next)) then
       last_leave => leave_list%next
       do while (associated (last_leave%next))
          leave_list%next => last_leave%next
          deallocate (last_leave)
          last_leave => leave_list%next
       end do
       deallocate (last_leave)
    end if
  end subroutine fibonacci_root_merge
  
\eatline
\nwendcode{}\nwbegindocs{1704}\nwdocspar
\nwenddocs{}\nwbegincode{1705}\moddef{Muli fibonacci: fibonacci root: TBP}\plusendmoddef
  procedure :: set_leftmost => fibonacci_root_set_leftmost
\nwendcode{}\nwbegincode{1706}\moddef{Muli fibonacci: procedures}\plusendmoddef
  subroutine fibonacci_root_set_leftmost (this)
    class(fibonacci_root_t) :: this
    call this%find_leftmost (this%leftmost)
  end subroutine fibonacci_root_set_leftmost

\eatline
\nwendcode{}\nwbegindocs{1707}\nwdocspar
\nwenddocs{}\nwbegincode{1708}\moddef{Muli fibonacci: fibonacci root: TBP}\plusendmoddef
  procedure :: set_rightmost => fibonacci_root_set_rightmost
\nwendcode{}\nwbegincode{1709}\moddef{Muli fibonacci: procedures}\plusendmoddef
  subroutine fibonacci_root_set_rightmost (this)
    class(fibonacci_root_t) :: this
    call this%find_rightmost (this%rightmost)
  end subroutine fibonacci_root_set_rightmost

\eatline
\nwendcode{}\nwbegindocs{1710}\nwdocspar
\nwenddocs{}\nwbegincode{1711}\moddef{Muli fibonacci: fibonacci root: TBP}\plusendmoddef
  procedure :: init_by_leave => fibonacci_root_init_by_leave
\nwendcode{}\nwbegincode{1712}\moddef{Muli fibonacci: procedures}\plusendmoddef
  subroutine fibonacci_root_init_by_leave (this, left_leave, right_leave)
    class(fibonacci_root_t), target, intent(out) :: this
    class(fibonacci_leave_t), target, intent(in) :: left_leave, right_leave
    if (left_leave <= right_leave) then
       this%left => left_leave
       this%right => right_leave
       this%leftmost => left_leave
       this%rightmost => right_leave
    else
       this%left => right_leave
       this%right => left_leave
       this%leftmost => right_leave
       this%rightmost => left_leave
    end if
    this%left%up => this
    this%right%up => this
    this%down => this%leftmost%down
    this%depth = 1
    this%leftmost%right => this%rightmost
    this%rightmost%left => this%leftmost
    this%is_valid_c = .true.
  end subroutine fibonacci_root_init_by_leave
  
\eatline
\nwendcode{}\nwbegindocs{1713}\nwdocspar
\nwenddocs{}\nwbegincode{1714}\moddef{Muli fibonacci: fibonacci root: TBP}\plusendmoddef
  procedure :: init_by_content => fibonacci_root_init_by_content
\nwendcode{}\nwbegincode{1715}\moddef{Muli fibonacci: procedures}\plusendmoddef
  subroutine fibonacci_root_init_by_content (this, left_content, right_content)
    class(fibonacci_root_t), target, intent(out) :: this
    class(measure_class_t), intent(in), target :: left_content, right_content
    call this%reset ()
    print *, "fibonacci_root_init_by_content: ", left_content%measure (), &
         right_content%measure ()
    if (left_content < right_content) then
       call this%leftmost%set_content (left_content)
       call this%rightmost%set_content (right_content)
    else
       call this%leftmost%set_content (right_content)
       call this%rightmost%set_content (left_content)
    end if
    this%down => this%leftmost%down
    this%is_valid_c = .true.
  end subroutine fibonacci_root_init_by_content
  
\eatline
\nwendcode{}\nwbegindocs{1716}\nwdocspar
\nwenddocs{}\nwbegincode{1717}\moddef{Muli fibonacci: fibonacci root: TBP}\plusendmoddef
  procedure :: reset => fibonacci_root_reset
\nwendcode{}\nwbegindocs{1718}\nwdocspar
\nwenddocs{}\nwbegincode{1719}\moddef{Muli fibonacci: procedures}\plusendmoddef
  subroutine fibonacci_root_reset (this)
    class(fibonacci_root_t), target, intent(inout) :: this
    call this%deallocate_tree ()
    allocate (this%leftmost)
    allocate (this%rightmost) 
    this%depth = 1
    this%leftmost%depth = 0
    this%rightmost%depth = 0
    this%left => this%leftmost
    this%right => this%rightmost
    this%left%up => this
    this%right%up => this
    this%leftmost%right => this%rightmost
    this%rightmost%left => this%leftmost
  end subroutine fibonacci_root_reset
  
\eatline
\nwendcode{}\nwbegindocs{1720}\nwdocspar
\nwenddocs{}\nwbegincode{1721}\moddef{Muli fibonacci: fibonacci root: TBP}\plusendmoddef
  procedure :: deallocate_tree => fibonacci_root_deallocate_tree
\nwendcode{}\nwbegincode{1722}\moddef{Muli fibonacci: procedures}\plusendmoddef
  recursive subroutine fibonacci_root_deallocate_tree (this)
    class(fibonacci_root_t), intent(inout) :: this
    call this%deallocate_tree ()
    nullify (this%leftmost)
    nullify (this%rightmost)
  end subroutine fibonacci_root_deallocate_tree

\eatline
\nwendcode{}\nwbegindocs{1723}\nwdocspar
\nwenddocs{}\nwbegincode{1724}\moddef{Muli fibonacci: fibonacci root: TBP}\plusendmoddef
  procedure :: deallocate_all => fibonacci_root_deallocate_all  
\nwendcode{}\nwbegindocs{1725}\nwdocspar
\nwenddocs{}\nwbegincode{1726}\moddef{Muli fibonacci: procedures}\plusendmoddef
  recursive subroutine fibonacci_root_deallocate_all (this)
    class(fibonacci_root_t), intent(inout) :: this
    call this%deallocate_all ()
    nullify (this%leftmost)
    nullify (this%rightmost)
  end subroutine fibonacci_root_deallocate_all

\eatline
\nwendcode{}\nwbegindocs{1727}\nwdocspar
\nwenddocs{}\nwbegincode{1728}\moddef{Muli fibonacci: types}\plusendmoddef
  ! class(serializable_ref_type), pointer :: ref_list  
\nwendcode{}\nwbegindocs{1729}\nwdocspar
\nwenddocs{}\nwbegincode{1730}\moddef{Muli fibonacci: types}\plusendmoddef
  type, extends (fibonacci_root_t) :: fibonacci_stub_t
   contains
   \LA{}Muli fibonacci: fibonacci stub: TBP\RA{}
  end type fibonacci_stub_t

\eatline
\nwendcode{}\nwbegindocs{1731}\nwdocspar
\nwenddocs{}\nwbegincode{1732}\moddef{Muli fibonacci: fibonacci stub: TBP}\endmoddef
  procedure :: write_to_marker => fibonacci_stub_write_to_marker
\nwendcode{}\nwbegincode{1733}\moddef{Muli fibonacci: procedures}\plusendmoddef
  subroutine fibonacci_stub_write_to_marker (this, marker, status)
    class(fibonacci_stub_t), intent(in) :: this
    class(marker_t), intent(inout) :: marker
    integer(dik), intent(out) :: status
  end subroutine fibonacci_stub_write_to_marker
  
\eatline
\nwendcode{}\nwbegindocs{1734}\nwdocspar
\nwenddocs{}\nwbegincode{1735}\moddef{Muli fibonacci: fibonacci stub: TBP}\plusendmoddef
  procedure :: read_target_from_marker => fibonacci_stub_read_target_from_marker
\nwendcode{}\nwbegincode{1736}\moddef{Muli fibonacci: procedures}\plusendmoddef
  subroutine fibonacci_stub_read_target_from_marker (this, marker, status)
    class(fibonacci_stub_t), target, intent(out) :: this
    class(marker_t), intent(inout) :: marker
    integer(dik), intent(out) :: status
  end subroutine fibonacci_stub_read_target_from_marker

\eatline
\nwendcode{}\nwbegindocs{1737}\nwdocspar
\nwenddocs{}\nwbegincode{1738}\moddef{Muli fibonacci: fibonacci stub: TBP}\plusendmoddef
  ! procedure :: print_to_unit => fibonacci_stub_print_to_unit  
\nwendcode{}\nwbegindocs{1739}\nwdocspar
\nwenddocs{}\nwbegincode{1740}\moddef{Muli fibonacci: fibonacci stub: TBP}\plusendmoddef
  procedure, nopass :: get_type => fibonacci_stub_get_type
\nwendcode{}\nwbegindocs{1741}\nwdocspar
\nwenddocs{}\nwbegincode{1742}\moddef{Muli fibonacci: procedures}\plusendmoddef
  pure subroutine fibonacci_stub_get_type (type)
    character(:), allocatable, intent(out) :: type
    allocate (type, source="fibonacci_stub_t")
  end subroutine fibonacci_stub_get_type

\eatline
\nwendcode{}\nwbegindocs{1743}\nwdocspar
\nwenddocs{}\nwbegincode{1744}\moddef{Muli fibonacci: fibonacci stub: TBP}\plusendmoddef
  procedure :: push_by_content => fibonacci_stub_push_by_content
\nwendcode{}\nwbegincode{1745}\moddef{Muli fibonacci: procedures}\plusendmoddef
  subroutine fibonacci_stub_push_by_content (this, content)
    class(fibonacci_stub_t), target, intent(inout) :: this
    class(measure_class_t), target, intent(in) :: content
    class(fibonacci_leave_t), pointer :: leave
    allocate (leave)
    call leave%set_content (content)
    call this%push_by_leave (leave)
  end subroutine fibonacci_stub_push_by_content

\eatline
\nwendcode{}\nwbegindocs{1746}\nwdocspar
\nwenddocs{}\nwbegincode{1747}\moddef{Muli fibonacci: fibonacci stub: TBP}\plusendmoddef
  procedure :: push_by_leave => fibonacci_stub_push_by_leave
\nwendcode{}\nwbegincode{1748}\moddef{Muli fibonacci: procedures}\plusendmoddef
  subroutine fibonacci_stub_push_by_leave (this, new_leave)
    class(fibonacci_stub_t), target, intent(inout) :: this
    class(fibonacci_leave_t), pointer, intent(inout) :: new_leave
    class(fibonacci_leave_t), pointer :: old_leave
    if (this%depth < 1) then
       if (associated (this%leftmost)) then     
          old_leave => this%leftmost
          call this%init_by_leave (old_leave, new_leave)
       else
          this%leftmost => new_leave
       end if
    else
       call fibonacci_root_push_by_leave (this, new_leave)
    end if
  end subroutine fibonacci_stub_push_by_leave
  
\eatline
\nwendcode{}\nwbegindocs{1749}\nwdocspar
\nwenddocs{}\nwbegincode{1750}\moddef{Muli fibonacci: fibonacci stub: TBP}\plusendmoddef
  procedure :: pop_left => fibonacci_stub_pop_left
\nwendcode{}\nwbegincode{1751}\moddef{Muli fibonacci: procedures}\plusendmoddef
  subroutine fibonacci_stub_pop_left (this, leave)
    class(fibonacci_stub_t), intent(inout), target :: this
    class(fibonacci_leave_t), pointer, intent(out) :: leave
    if (this%depth < 2) then
       if (this%depth == 1) then
          leave => this%leftmost
          this%leftmost => this%rightmost
          nullify (this%rightmost)
          nullify (this%right)
          nullify (this%left)
          this%depth = 0
          this%is_valid_c = .false.
       else
          if (associated (this%leftmost)) then
             leave => this%leftmost
             nullify (this%leftmost)
          end if
       end if
    else
       call fibonacci_root_pop_left (this, leave)
    end if
  end subroutine fibonacci_stub_pop_left

\eatline
\nwendcode{}\nwbegindocs{1752}\nwdocspar
\nwenddocs{}\nwbegincode{1753}\moddef{Muli fibonacci: fibonacci stub: TBP}\plusendmoddef
  procedure :: pop_right => fibonacci_stub_pop_right
\nwendcode{}\nwbegincode{1754}\moddef{Muli fibonacci: procedures}\plusendmoddef
  subroutine fibonacci_stub_pop_right (this, leave)
    class(fibonacci_stub_t), intent(inout), target :: this
    class(fibonacci_leave_t), pointer, intent(out) :: leave
    if (this%depth < 2) then
       if (this%depth == 1) then
          this%is_valid_c = .false.
          if (associated (this%rightmost)) then
             leave => this%rightmost
             nullify (this%rightmost)
             nullify (this%right)
          else
             if (associated (this%leftmost)) then
                leave => this%leftmost
                nullify (this%leftmost)
                nullify (this%left)
             else
                nullify (leave)
             end if
          end if
       end if
    else
       call fibonacci_root_pop_right (this, leave)
    end if
  end subroutine fibonacci_stub_pop_right
     
\eatline
\nwendcode{}\nwbegindocs{1755}\nwdocspar
\nwenddocs{}\nwbegincode{1756}\moddef{Muli fibonacci: public}\plusendmoddef
  public :: fibonacci_leave_list_t
\nwendcode{}\nwbegincode{1757}\moddef{Muli fibonacci: types}\plusendmoddef
  type fibonacci_leave_list_t
     class(fibonacci_leave_t), pointer :: leave => null()
     class(fibonacci_leave_list_t), pointer :: next => null()
  end type fibonacci_leave_list_t
  
\eatline
\nwendcode{}\nwbegindocs{1758}\nwdocspar
\nwenddocs{}\nwbegincode{1759}\moddef{Muli fibonacci: procedures}\plusendmoddef
!  subroutine fibonacci_node_update_value (this, right_value)
!    class(fibonacci_node_t), target :: this
!    class(fibonacci_node_t), pointer:: node
!    real(default), intent(in) :: right_value
!    if (associated (this%left) .and. associated (this%right)) then
!       node => this
!       ! node%value = node%left%value
!       ! right_value = node%right%value
!       INNER: do while (associated (node%up))
!          if (node%is_right_child ()) then
!             node => node%up
!          else
!             node%up%value = right_value
!             exit
!          end if
!       end do INNER
!    end if
!  end subroutine fibonacci_node_update_value

\eatline
\nwendcode{}\nwbegindocs{1760}\nwdocspar
\nwenddocs{}\nwbegincode{1761}\moddef{Muli fibonacci: procedures}\plusendmoddef
!  subroutine fibonacci_root_copy_node (this, primitive)
!    class(fibonacci_root_t), intent(out) :: this
!    type(fibonacci_node_t), intent(in) :: primitive
!    call fibonacci_node_copy_node (this, primitive)
!    call primitive%find_leftmost (this%leftmost)
!    call primitive%find_rightmost (this%rightmost)
!  end subroutine fibonacci_root_copy_node

\eatline
\nwendcode{}\nwbegindocs{1762}\nwdocspar
\nwenddocs{}\nwbegincode{1763}\moddef{Muli fibonacci: procedures}\plusendmoddef
!  subroutine fibonacci_root_push_by_node (this, new_leave)
!    class(fibonacci_root_t), target, intent(inout) :: this
!    class(fibonacci_leave_t), pointer, intent(inout) :: new_leave
!    class(fibonacci_leave_t), pointer :: old_leave
!    if (new_leave <= this%leftmost) then
!       old_leave => this%leftmost
!       this%leftmost => new_leave
!    else
!       if (new_leave > this%rightmost) then
!          old_leave => this%rightmost
!          this%rightmost => new_leave
!       else
!          call this%find (new_leave%measure(), old_leave)
!       end if
!    end if
!    ! call old_leave%insert_leave_by_node (new_leave)
!    call fibonacci_leave_insert_leave_by_node (old_leave, new_leave)
!    call new_leave%up%repair ()
!    ! call new_leave%up%update_value ()
!  end subroutine fibonacci_root_push_by_node

\eatline
\nwendcode{}\nwbegindocs{1764}\nwdocspar
\nwenddocs{}\nwbegincode{1765}\moddef{Muli fibonacci: procedures}\plusendmoddef
  subroutine fibonacci_leave_write_content (this, unit)
    class(fibonacci_leave_t), intent(in), target :: this
    integer,optional, intent(in) :: unit
    call this%down%print_all (unit)
  end subroutine fibonacci_leave_write_content

\eatline
\nwendcode{}\nwbegindocs{1766}\nwdocspar
\nwenddocs{}\nwbegincode{1767}\moddef{Muli fibonacci: procedures}\plusendmoddef
  subroutine fibonacci_leave_write (this, unit)
    class(fibonacci_leave_t), intent(in), target :: this
    integer,optional, intent(in) :: unit
    call this%print_all (unit)
  end subroutine fibonacci_leave_write

\eatline
\nwendcode{}\nwbegindocs{1768}\nwdocspar
\nwenddocs{}\nwbegincode{1769}\moddef{Muli fibonacci: procedures}\plusendmoddef
  subroutine fibonacci_leave_write_value (this, unit)
    class(fibonacci_leave_t), intent(in), target :: this
    integer, intent(in), optional :: unit
    if (present (unit)) then
       write(unit, fmt=*)  this%measure ()
    else
       print *, this%measure ()
    end if
    ! call this%print_little (unit)
  end subroutine fibonacci_leave_write_value

\eatline
\nwendcode{}\nwbegindocs{1770}\nwdocspar
\nwenddocs{}\nwbegincode{1771}\moddef{Muli fibonacci: procedures}\plusendmoddef
  subroutine fibonacci_node_spawn (new_node, left_leave, right_leave, &
       left_left_leave, right_right_leave)
    class(fibonacci_node_t), pointer, intent(out) :: new_node
    class(fibonacci_leave_t), target, intent(inout) :: left_leave, right_leave
    class(fibonacci_node_t), pointer, intent(inout) :: left_left_leave, &
         right_right_leave
    allocate (new_node)
    new_node%depth = 1
    if (associated (left_left_leave)) then
       left_left_leave%right => left_leave
       left_leave%left => left_left_leave
    else
       nullify (left_leave%left)
    end if
    if (associated (right_right_leave)) then
       right_right_leave%left => right_leave
       right_leave%right => right_right_leave
    else
       nullify (right_leave%right)
    end if
    new_node%left => left_leave
    new_node%right => right_leave
    new_node%down => left_leave%down
    new_node%depth = 1
    left_leave%up => new_node
    right_leave%up => new_node
    left_leave%right => right_leave
    right_leave%left => left_leave
  end subroutine fibonacci_node_spawn
  
\eatline
\nwendcode{}\nwbegindocs{1772}\nwdocspar
%%%%%%%%%%%%%%%%%%%%%%%%%%%%%%%%%%%%%%%%%%%%%%%%%%%%%%%%%%%%%%%%%%%%%%%%
\section{Adaptive Quadrature}

This file contains the module {\tt{}muli{\char95}aq} which is an acronym for
adaptive quadrature. The abstract type {\tt{}aq{\char95}class} is extended
and the deferred procedure {\tt{}evaluate} is defined. {\tt{}muli{\char95}aq} calls
{\tt{}evaluate} to evaluate the integrand at any point in the given
range. We have tried to use a procedure pointer instead of the
deferred type-bound procedure, but no compiler was able to handle
procedure pointers plus cuba was not able to handle parameters, to wit
dimensions of the integrand that should not get integrated. So we
switched to this odd way of using inheritance. Meanwhile these
problems got solved and we could go for a more straightforward
solution, but it works fine as it is. 

{\tt{}aq{\char95}class} uses {\tt{}muli{\char95}trapezium} to approximate the integral. It
still has to do the subdivision of segments and has to check whether
the precision goal is reached.

Finally, the result is written do disc using the serialization framework 
defined in {\tt{}muli{\char95}base}. Since QCD is not expected to change frequently,
the only reason to regenerate this function is a change of the used PDF set.
Then you can read the integral from disc each time you run a simulation with
the same PDF set.

\nwenddocs{}\nwbegincode{1773}\moddef{\code{}muli{\char95}aq.f90\edoc{}}\endmoddef
\LA{}File header\RA{}

module muli_aq
  
\LA{}Use kinds\RA{}
  use constants
  use diagnostics
  use muli_base
  use muli_cuba
  use muli_trapezium
  use muli_fibonacci_tree

\LA{}Standard module header\RA{}  

\LA{}Muli aq: public\RA{}

\LA{}Muli aq: types\RA{}

\LA{}Muli aq: interfaces\RA{}

contains

\LA{}Muli aq: procedures\RA{}  

end module muli_aq

\eatline
\nwendcode{}\nwbegindocs{1774}The variables {\tt{}error{\char95}goal}, {\tt{}err{\char95}tree} and {\tt{}int{\char95}list} must be
initialised before the main loop can be called. Additionaly, the nodes
and segments should be preprocessed by {\tt{}first{\char95}run} before the main
loop is called.
\nwenddocs{}\nwbegincode{1775}\moddef{Muli aq: types}\endmoddef
  type, extends (identified_t), abstract :: aq_class
     logical :: is_deferred_initialised = .false.
     logical :: is_error_tree_initialised = .false.
     logical :: is_goal_set = .false.
     logical :: is_initialised = .false.
     logical :: is_run = .false.
     logical :: is_goal_reached = .false.
     logical :: is_integrated = .false.
     integer(dik) :: n_nodes = 0
     integer(dik) :: max_nodes = 10000
     integer :: dim_integral = 1
     real(default) :: abs_error_goal = 0._default
     real(default) :: rel_error_goal = 0.1_default
     real(default) :: scaled_error_goal = 0._default
     real(default) :: integral = 1._default
     real(default) :: integral_error = 0._default
     real(default), dimension(2) :: region = [0._default, 1._default]
     real(default), dimension(:,:), allocatable :: convergence
     real(default) :: total_time = 0
     real(default) :: loop_time = 0
     real(default) :: int_time = 0
     real(default) :: cuba_time = 0
     real(default) :: init_time = 0
     real(default) :: cpu_time = 0
     real(default) :: error_goal = 0._default
     class(fibonacci_root_t), pointer :: err_tree => null()
     class(muli_trapezium_list_t), pointer :: int_list => null()
   contains
   \LA{}Muli aq: aq class: TBP\RA{}  
  end type aq_class
  
\eatline
\nwendcode{}\nwbegindocs{1776}\nwdocspar
\nwenddocs{}\nwbegincode{1777}\moddef{Muli aq: aq class: TBP}\endmoddef
  procedure :: write_to_marker => aq_write_to_marker
\nwendcode{}\nwbegincode{1778}\moddef{Muli aq: procedures}\endmoddef
  subroutine aq_write_to_marker (this, marker, status)
    class(aq_class), intent(in) :: this
    class(marker_t), intent(inout) :: marker
    integer(dik), intent(out) :: status  
    class(ser_class_t), pointer :: ser
    call marker%mark_begin ("aq_class")
    call this%base_write_to_marker (marker, status)
    call marker%mark ("is_deferred_initialised", &
         this%is_deferred_initialised)
    call marker%mark ("is_error_tree_initialised", &
         this%is_error_tree_initialised)
    call marker%mark ("is_goal_set", this%is_goal_set)
    call marker%mark ("is_initialised", this%is_initialised)
    call marker%mark ("is_run", this%is_run)
    call marker%mark ("is_goal_reached", this%is_goal_reached)
    call marker%mark ("is_integrated", this%is_integrated)
    call marker%mark ("n_nodes", this%n_nodes)
    call marker%mark ("max_nodes", this%max_nodes)
    call marker%mark ("dim_integral", this%dim_integral)    
    call marker%mark ("abs_error_goal", this%abs_error_goal)
    call marker%mark ("rel_error_goal", this%rel_error_goal)
    call marker%mark ("scaled_error_goal", this%scaled_error_goal)
    call marker%mark ("error_goal", this%error_goal)
    call marker%mark ("integral", this%integral)
    call marker%mark ("integral_error", this%integral_error)
    call marker%mark ("region", this%region(1:2))
    ser => this%err_tree
    call marker%mark_pointer ("err_tree", ser)
    ser => this%int_list
    call marker%mark_pointer ("int_list", ser)
    call marker%mark_end ("aq_class")
  end subroutine aq_write_to_marker
  
\eatline
\nwendcode{}\nwbegindocs{1779}\nwdocspar
\nwenddocs{}\nwbegincode{1780}\moddef{Muli aq: aq class: TBP}\plusendmoddef
  procedure :: read_from_marker => aq_read_from_marker
\nwendcode{}\nwbegincode{1781}\moddef{Muli aq: procedures}\plusendmoddef
  subroutine aq_read_from_marker (this, marker, status)
    class(aq_class), intent(out) :: this
    class(marker_t), intent(inout) :: marker
    integer(dik), intent(out) :: status  
    class(ser_class_t), pointer :: ser
    call marker%pick_begin ("aq_class", status=status)
    call this%base_read_from_marker (marker, status)
    call marker%pick ("is_deferred_initialised", &
         this%is_deferred_initialised, status)
    call marker%pick ("is_error_tree_initialised", &
         this%is_error_tree_initialised, status)
    call marker%pick ("is_goal_set", this%is_goal_set, status)
    call marker%pick ("is_initialised", this%is_initialised, status)
    call marker%pick ("is_run", this%is_run, status)
    call marker%pick ("is_goal_reached", this%is_goal_reached, status)
    call marker%pick ("is_integrated", this%is_integrated, status)
    call marker%pick ("n_nodes", this%n_nodes, status)
    call marker%pick ("max_nodes", this%max_nodes, status)
    call marker%pick ("dim_integral", this%dim_integral, status)    
    call marker%pick ("abs_error_goal", this%abs_error_goal, status)
    call marker%pick ("rel_error_goal", this%rel_error_goal, status)
    call marker%pick ("scaled_error_goal", this%scaled_error_goal, status)
    call marker%pick ("error_goal", this%error_goal, status)
    call marker%pick ("integral", this%integral, status)
    call marker%pick ("integral_error", this%integral_error, status)
    call marker%pick ("region", this%region(1:2), status)
    call marker%pick_pointer ("err_tree", ser)
    if (associated (ser)) then
       select type (ser)
       class is (fibonacci_root_t)
          this%err_tree => ser
       class default
          nullify (this%err_tree)
       end select
    end if
    call marker%pick_pointer ("int_list", ser)
    if (associated (ser)) then
       select type (ser)
       class is (muli_trapezium_list_t)
          this%int_list => ser
       class default
          nullify (this%int_list)
       end select
    end if
    call marker%pick_end ("aq_class", status)
  end subroutine aq_read_from_marker
  
\eatline
\nwendcode{}\nwbegindocs{1782}\nwdocspar
\nwenddocs{}\nwbegincode{1783}\moddef{Muli aq: aq class: TBP}\plusendmoddef
  procedure :: print_to_unit => aq_print_to_unit  
\nwendcode{}\nwbegincode{1784}\moddef{Muli aq: procedures}\plusendmoddef
  subroutine aq_print_to_unit (this, unit, parents, components, peers)
    class(aq_class), intent(in) :: this
    integer, intent(in) :: unit
    integer(dik), intent(in) :: parents, components, peers
    integer :: ite
    class(ser_class_t), pointer :: ser
    if (parents > 0)  call this%base_print_to_unit &
         (unit, parents-1, components, peers)
    write (unit, "(A)")  "Components of aq_class"
    write (unit, "(A,L1)")  "Deferred class initialised: ", &
         this%is_deferred_initialised
    write (unit, "(A,L1)")  "Error tree initialised:     ", &
         this%is_error_tree_initialised
    write (unit, "(A,L1)")  "Accuracy goal set:          ", this%is_goal_set
    write (unit, "(A,L1)")  "Ready for run:              ", this%is_initialised
    write (unit, "(A,L1)")  "Is run:                     ", this%is_run
    write (unit, "(A,L1)")  "Accuracy goal reached:      ", this%is_goal_reached
    write (unit, "(A,L1)")  "Integral calculated:        ", this%is_integrated
    write (unit, "(A,I10)") "Number of nodes:            ", this%n_nodes
    write (unit, "(A,I10)") "Maximal number of nodes:    ", this%max_nodes
    write (unit, "(A,I10)") "Dimension of integral:      ", this%dim_integral
    write (unit, "(A,E20.10)") "Given abs. error goal: ", this%abs_error_goal
    write (unit, "(A,E20.10)") "Given rel. error goal: ", this%rel_error_goal
    write (unit, "(A,E20.10)") "Guessed abs error goal:", this%scaled_error_goal
    write (unit, "(A,E20.10)") "Actual abs error goal: ", this%error_goal
    write (unit, "(A,E20.10)") "Integral               ", this%integral
    write (unit, "(A,E20.10)") "Estimated abs. error:  ", this%integral_error
    ! if (this%integral == 0) then
    !    write (unit, "(A,E20.10)")  "Estimated rel. error:  ", &
    !       this%integral_error / this%integral
    ! else
    !    write (unit, "(A,E20.10)")  "Estimated rel. error:  INF"
    ! end if
    write (unit, "(A,E10.5,A,E10.5,A)")  "Integration region =  (", &
         this%region(1), " : ", this%region(2), ")"
    ser => this%err_tree
    call serialize_print_comp_pointer &
         (ser, unit, parents, components, peers, "error tree")
    ser => this%int_list
    call serialize_print_comp_pointer &
         (ser, unit, parents, components, peers, "integral list")
  end subroutine aq_print_to_unit
  
\eatline
\nwendcode{}\nwbegindocs{1785}\nwdocspar
\nwenddocs{}\nwbegincode{1786}\moddef{Muli aq: aq class: TBP}\plusendmoddef
  procedure, nopass :: get_type => aq_get_type  
\nwendcode{}\nwbegincode{1787}\moddef{Muli aq: procedures}\plusendmoddef
  pure subroutine aq_get_type (type)
    character(:), allocatable, intent(out) :: type
    allocate (type, source="aq_type")
  end subroutine aq_get_type
    
\eatline
\nwendcode{}\nwbegindocs{1788}\nwdocspar
\nwenddocs{}\nwbegincode{1789}\moddef{Muli aq: aq class: TBP}\plusendmoddef
  procedure :: deserialize_from_marker => aq_deserialize_from_marker
\nwendcode{}\nwbegincode{1790}\moddef{Muli aq: procedures}\plusendmoddef
  subroutine aq_deserialize_from_marker (this, name, marker)
    class(aq_class), intent(out) :: this
    character(*), intent(in) :: name
    class(marker_t), intent(inout) :: marker
    class(ser_class_t), pointer :: ser
    allocate (muli_trapezium_t :: ser)
    call marker%push_reference (ser)
    allocate (fibonacci_root_t :: ser)
    call marker%push_reference (ser)
    allocate (fibonacci_leave_t :: ser)
    call marker%push_reference (ser)
    allocate (fibonacci_node_t :: ser)
    call marker%push_reference (ser)
    call serializable_deserialize_from_marker (this, name, marker)
    call marker%pop_reference (ser)
    deallocate (ser)
    call marker%pop_reference (ser)
    deallocate (ser)
    call marker%pop_reference (ser)
    deallocate (ser)
    call marker%pop_reference (ser)
    deallocate (ser)
  end subroutine aq_deserialize_from_marker
  
\eatline
\nwendcode{}\nwbegindocs{1791}The relative error goal is supposed to be $10^{-4}$.
\nwenddocs{}\nwbegincode{1792}\moddef{Muli aq: aq class: TBP}\plusendmoddef
  generic :: initialize => aq_initialize
  procedure :: aq_initialize
\nwendcode{}\nwbegincode{1793}\moddef{Muli aq: procedures}\plusendmoddef
  subroutine aq_initialize (this, id, name, goal, max_nodes, dim, init)
    class(aq_class), intent(out) :: this
    integer(dik), intent(in) :: id, max_nodes
    integer, intent(in) :: dim
    character, intent(in) :: name
    real(default) :: goal
    real(default), dimension(:), intent(in) :: init
    call this%initialize (id, name)
    this%rel_error_goal = goal
    this%max_nodes = max_nodes
    call this%init_error_tree (dim, init)
  end subroutine aq_initialize
  
\eatline
\nwendcode{}\nwbegindocs{1794}\nwdocspar
\nwenddocs{}\nwbegincode{1795}\moddef{Muli aq: aq class: TBP}\plusendmoddef
  procedure :: print_times => aq_print_times
\nwendcode{}\nwbegincode{1796}\moddef{Muli aq: procedures}\plusendmoddef
  subroutine aq_print_times (this)
    class(aq_class), intent(in) :: this
    write (*, "(A,E20.10)") "Initialization time:  ", this%init_time
    write (*, "(A,E20.10)") "Main loop time:       ", this%loop_time
    write (*, "(A,E20.10)") "Integration time:     ", this%int_time
    write (*, "(A,E20.10)") "Overall run time:     ", this%total_time
    write (*, "(A,E20.10)") "Cuba integration time:", this%cuba_time
  end subroutine aq_print_times
  
\nwendcode{}\nwbegindocs{1797}\nwdocspar
\nwenddocs{}\nwbegincode{1798}\moddef{Muli aq: aq class: TBP}\plusendmoddef
  procedure :: write_convergence => aq_write_convergence   
\nwendcode{}\nwbegincode{1799}\moddef{Muli aq: procedures}\plusendmoddef
  subroutine aq_write_convergence (this, unit)
    class(aq_class), intent(in) :: this
    integer, intent(in) :: unit
    integer, dimension(2) :: s
    integer :: node
    if (allocated (this%convergence)) then
       s = shape (this%convergence)
       do node = 1, s(2)
          write (unit, *) node, this%convergence (1:2, node)
       end do
    end if
  end subroutine aq_write_convergence
  
\eatline
\nwendcode{}\nwbegindocs{1800}\nwdocspar
\nwenddocs{}\nwbegincode{1801}\moddef{Muli aq: aq class: TBP}\plusendmoddef
  procedure :: reset => aq_reset
\nwendcode{}\nwbegincode{1802}\moddef{Muli aq: procedures}\plusendmoddef
  subroutine aq_reset (this)
    class(aq_class) :: this
    this%is_deferred_initialised = .false.
    this%is_error_tree_initialised = .false.
    this%is_goal_set = .false.
    this%is_initialised = .false.
    this%is_run = .false.
    this%is_goal_reached = .false.
    this%is_integrated = .false.
    this%n_nodes = 0
    this%max_nodes = 10000
    this%dim_integral = 1
    this%abs_error_goal = 1._default
    this%rel_error_goal = 0.1_default
    this%scaled_error_goal = 0.0_default
    this%error_goal = 0.0_default
    this%integral = 0.0_default
    this%integral_error = 0.0_default
    this%region = [ 0.0_default, 1._default ]
    this%total_time = 0
    this%loop_time = 0
    this%int_time = 0
    this%init_time = 0
    call this%dealloc_trees ()
  end subroutine aq_reset
  
\eatline
\nwendcode{}\nwbegindocs{1803}\nwdocspar
\nwenddocs{}\nwbegincode{1804}\moddef{Muli aq: aq class: TBP}\plusendmoddef
  procedure :: dealloc_trees => aq_dealloc_trees
  procedure :: finalize => aq_dealloc_trees  
\nwendcode{}\nwbegincode{1805}\moddef{Muli aq: procedures}\plusendmoddef
  subroutine aq_dealloc_trees (this)
    class(aq_class) :: this
    if (associated (this%err_tree)) then
       call this%err_tree%deallocate_all ()
       deallocate (this%err_tree)
    end if
    if (associated (this%int_list)) then
       call this%int_list%finalize ()
       deallocate (this%int_list)
    end if
  end subroutine aq_dealloc_trees
  
\eatline
\nwendcode{}\nwbegindocs{1806}\nwdocspar
\nwenddocs{}\nwbegincode{1807}\moddef{Muli aq: aq class: TBP}\plusendmoddef
  procedure :: init_error_tree => aq_init_error_tree
\nwendcode{}\nwbegincode{1808}\moddef{Muli aq: procedures}\plusendmoddef
  subroutine aq_init_error_tree (this, dim_integral, x_array)
    class(aq_class) :: this
    integer, intent(in) :: dim_integral
    real(default), dimension(:), intent(in) :: x_array
    real(default) :: center
    real(default), dimension(:), allocatable :: l_val, c_val, r_val
    class(muli_trapezium_t), pointer :: left_node => null()
    class(muli_trapezium_t), pointer :: right_node => null()
    integer :: x_size, pos
    ! print '("Entermarker aq_init_error_tree...")'
    call cpu_time (this%init_time)
    this%is_initialised = .false.
    this%integral = 0._default
    this%dim_integral = dim_integral
    x_size = size(x_array)
    if (x_size < 2) then
         call msg_error ("aq_init_error_tree: I need at least two real values")
    else
       allocate (l_val (0:dim_integral-1))
       allocate (c_val (0:dim_integral-1))
       allocate (r_val (0:dim_integral-1))
       this%region = [x_array(1), x_array(x_size)]
       if (x_size < 3) then
          center = (x_array(2) - x_array(1)) / 2._default
          call this%evaluate (x_array(1), l_val)
          call this%evaluate (center, c_val)
          call this%evaluate (x_array(2), r_val)
          allocate (left_node)
          call left_node%initialize (dim=dim_integral, &
               r_position=center, d_position=center-x_array(1))
          call left_node%set_r_value (c_val)
          call left_node%set_d_value (c_val - l_val)
          allocate (right_node)
          call right_node%initialize (dim=dim_integral, &
               r_position=x_array(2), d_position=x_array(2)-center)
          call right_node%set_r_value (r_val)
          call right_node%set_d_value (r_val - c_val)
       else
          call this%evaluate (x_array(1), l_val)
          call this%evaluate (x_array(2), c_val)
          call this%evaluate (x_array(3), r_val)
          allocate (left_node)
          call left_node%initialize (dim=dim_integral, &
               r_position=x_array(2), d_position=x_array(2)-x_array(1))
          call left_node%set_r_value (c_val)
          call left_node%set_d_value (c_val - l_val)
          allocate (right_node)
          call right_node%initialize (dim=dim_integral, &
               r_position=x_array(3), d_position=x_array(3)-x_array(2))
          call right_node%set_r_value (r_val)
          call right_node%set_d_value (r_val - c_val)
       end if
       call left_node%update ()
       call right_node%update ()
       this%integral = sum (left_node%get_d_integral () + &
            right_node%get_d_integral ())
       if (.not. associated (this%err_tree)) then
          allocate(this%err_tree)
       end if
       print *, left_node%measure ()
       print *, right_node%measure ()
       call this%err_tree%init_by_content (left_node, right_node)
       ! call this%err_tree%write_pstricks (11)
       if (x_size > 3) then
          do pos = 4, x_size
             print *, "aq_init_error_tree", pos, "/", x_size
             l_val = right_node%get_r_value_array ()
             call this%evaluate (x_array(pos), r_val)
             c_val = r_val - l_val
             allocate (right_node)
             call right_node%initialize (dim=dim_integral, &
                  r_position=x_array(pos), &
                  d_position=x_array(pos)-x_array(pos-1))
             call right_node%set_r_value (r_val)
             call right_node%set_d_value (c_val)
             call right_node%update ()
             call this%err_tree%push_by_content (right_node)
             ! call this%err_tree%write_pstricks (11)
             this%integral = this%integral + sum (right_node%get_d_integral())
          end do
          this%n_nodes = x_size
       end if
       this%is_error_tree_initialised = .true.
    end if
    call this%set_goal ()
    this%is_initialised = .true.
    call cpu_time (this%cpu_time)
    this%init_time = this%cpu_time - this%init_time
    this%cuba_time = this%init_time
    allocate (this%convergence (2, this%n_nodes:this%max_nodes))
  end subroutine aq_init_error_tree
  
\nwendcode{}\nwbegindocs{1809}\nwdocspar
\nwenddocs{}\nwbegincode{1810}\moddef{Muli aq: aq class: TBP}\plusendmoddef
  procedure :: set_rel_goal => aq_set_rel_goal
\nwendcode{}\nwbegincode{1811}\moddef{Muli aq: procedures}\plusendmoddef
  subroutine aq_set_rel_goal (this, goal)
    class(aq_class) :: this
    real(default) :: goal
    this%rel_error_goal = goal
    call this%set_goal
  end subroutine aq_set_rel_goal

\eatline
\nwendcode{}\nwbegindocs{1812}\nwdocspar
\nwenddocs{}\nwbegincode{1813}\moddef{Muli aq: aq class: TBP}\plusendmoddef
  procedure :: set_abs_goal => aq_set_abs_goal
\nwendcode{}\nwbegincode{1814}\moddef{Muli aq: procedures}\plusendmoddef
  subroutine aq_set_abs_goal (this, goal)
    class(aq_class) :: this
    real(default) :: goal
    this%abs_error_goal = goal
    call this%set_goal
  end subroutine aq_set_abs_goal
  
\eatline
\nwendcode{}\nwbegindocs{1815}\nwdocspar
\nwenddocs{}\nwbegincode{1816}\moddef{Muli aq: aq class: TBP}\plusendmoddef
  procedure :: set_goal => aq_set_goal
\nwendcode{}\nwbegincode{1817}\moddef{Muli aq: procedures}\plusendmoddef
  subroutine aq_set_goal (this)
    class(aq_class) :: this
    this%scaled_error_goal = this%rel_error_goal * abs(this%integral)
    if ((this%scaled_error_goal == zero) .and. &
        (this%abs_error_goal == zero)) then
       this%is_goal_set = .false.
       this%error_goal = zero
    else
       if (this%scaled_error_goal == zero) then
          this%error_goal = this%abs_error_goal
       else
          if (this%abs_error_goal == zero) then
             this%error_goal = this%scaled_error_goal
          else
             this%error_goal = max(this%scaled_error_goal, &
                  this%abs_error_goal)
          end if
       end if
       if (this%error_goal > zero) then
          this%is_goal_set = .true.
       else
          this%is_goal_set = .false.
       end if
    end if
  end subroutine aq_set_goal
  
\eatline
\nwendcode{}\nwbegindocs{1818}\nwdocspar
\nwenddocs{}\nwbegincode{1819}\moddef{Muli aq: aq class: TBP}\plusendmoddef
  procedure :: check_init => aq_check_init  
\nwendcode{}\nwbegincode{1820}\moddef{Muli aq: procedures}\plusendmoddef
  subroutine aq_check_init (this)
    class(aq_class) :: this
    this%is_initialised = this%is_error_tree_initialised .and. &
         this%is_deferred_initialised
  end subroutine aq_check_init
  
\eatline
\nwendcode{}\nwbegindocs{1821} This routine is unsafe, when {\tt{}n{\char95}nodes} $< 4$.
\nwenddocs{}\nwbegincode{1822}\moddef{Muli aq: aq class: TBP}\plusendmoddef
  procedure :: main_loop => aq_main_loop
\nwendcode{}\nwbegincode{1823}\moddef{Muli aq: procedures}\plusendmoddef
  subroutine aq_main_loop (this)
    class(aq_class) :: this
    class(fibonacci_leave_t), pointer :: rightmost
    class(measure_class_t), pointer :: content
    class(muli_trapezium_t), pointer :: new_node !,debug
    logical :: limit = .false.
    real(default) :: center
    real(default), dimension(:), allocatable :: c_val
    allocate (c_val (0:this%dim_integral-1))
    LOOP: do
       call this%err_tree%pop_right (rightmost)
       if (rightmost < this%error_goal / this%n_nodes) then
          this%is_goal_reached = .true.
          exit LOOP
       else
          call rightmost%get_content (content)
          select type (content)
          class is (muli_trapezium_t)
             write (*, "(A,I5,A,E14.7,A,E14.7,A,E14.7,A,E14.7)")  &
                  "nodes: ", this%n_nodes, " error: ", &
                  rightmost%measure() * this%n_nodes, &
                  " goal: ", this%error_goal, " node at: ", &
                  content%get_l_position(), "-", content%get_r_position()
             this%convergence (1, this%n_nodes) = this%error_goal / this%n_nodes
             this%convergence (2, this%n_nodes) = rightmost%measure ()
             center = content%get_r_position () - &
                  content%get_d_position () / two
             call cpu_time (this%cpu_time)
             this%cuba_time = this%cuba_time - this%cpu_time
             call this%evaluate (center, c_val)
             call cpu_time (this%cpu_time)
             this%cuba_time = this%cuba_time + this%cpu_time
             call content%split (c_val, center, new_node)
             call this%err_tree%push_by_leave (rightmost)
             call this%err_tree%push_by_content (new_node)
          end select
          this%n_nodes = this%n_nodes + 1
          if (this%n_nodes > this%max_nodes) then
             limit = .true.             
             exit LOOP
          end if
       end if
    end do LOOP
    call this%err_tree%push_by_leave (rightmost)
  end subroutine aq_main_loop
  
\eatline
\nwendcode{}\nwbegindocs{1824}\nwdocspar
\nwenddocs{}\nwbegincode{1825}\moddef{Muli aq: aq class: TBP}\plusendmoddef
  procedure :: run => aq_run
\nwendcode{}\nwbegincode{1826}\moddef{Muli aq: procedures}\plusendmoddef
  subroutine aq_run (this)
    class(aq_class) :: this
    call cpu_time (this%total_time)
    if (.not. this%is_error_tree_initialised) then
       call this%init_error_tree (this%dim_integral, this%region)
    end if
    this%is_run = .false.
    this%is_goal_reached = .false.
    call this%main_loop ()
    this%is_run = .true.
    call cpu_time (this%cpu_time)
    this%total_time = this%cpu_time - this%total_time
  end subroutine aq_run
  
\eatline
\nwendcode{}\nwbegindocs{1827}\nwdocspar
\nwenddocs{}\nwbegincode{1828}\moddef{Muli aq: aq class: TBP}\plusendmoddef
  procedure :: integrate => aq_integrate  
\nwendcode{}\nwbegincode{1829}\moddef{Muli aq: procedures}\plusendmoddef
  subroutine aq_integrate (this, int_tree)
    class(aq_class) :: this
    class(muli_trapezium_node_class_t), pointer :: node
    type(muli_trapezium_tree_t), intent(out) :: int_tree
    real(default) :: sum
    this%is_integrated = .false.
    this%integral_error = zero
    if (this%is_run) then
    call cpu_time (this%int_time)
       call fibonacci_tree_resort_and_convert_to_trapezium_list &
            (this%err_tree, this%int_list)
       ! call this%int_list%print_all ()
       call this%int_list%integrate (this%integral, this%integral_error)
       call this%int_list%to_tree (int_tree)
       this%is_integrated = .true.
       call cpu_time (this%cpu_time)
       this%int_time = this%cpu_time - this%int_time
    end if
  end subroutine aq_integrate
  
\eatline
\nwendcode{}\nwbegindocs{1830}\nwdocspar
\nwenddocs{}\nwbegincode{1831}\moddef{Muli aq: aq class: TBP}\plusendmoddef
  procedure(evaluate_if), deferred :: evaluate
  ! procedure(evaluate_ratios_if), deferred :: evaluate_ratios  
\nwendcode{}\nwbegincode{1832}\moddef{Muli aq: procedures}\plusendmoddef
\nwendcode{}\nwbegindocs{1833} Usually, the tree is sorted by the sum of errors. Now it shall be
sorted by the right position. When at least one branch of the tree is
itself a tree, i.e. each branch has got at least two leaves, then
process each branch and merge the results. Now we got two sorted
lists. Which one's leftmost node has got the lowest value of
{\tt{}r{\char95}position}? That one shall be the beginning of the merged list
{\tt{}lin{\char95}list}. Everything is prepared for the algorithm: {\tt{}lin{\char95}list}
is the beginning of the sorted list, {\tt{}last{\char95}node} is its
end. {\tt{}left{\char95}node} and {\tt{}right{\char95}node} are the leftmost nodes of the
remainders of {\tt{}left{\char95}list} and {\tt{}right{\char95}list}. The latter will get
stripped from left to right, until one of them ends. Then, either 
{\tt{}left{\char95}list} or {\tt{}right{\char95}list} is completely merged into
{\tt{}lin{\char95}list}. The other one gets appended to {\tt{}lin{\char95}list}. In the
second part of the big if clause, the tree has got two leaves at
most. Is it more than one? There, {\tt{}fib{\char95}tree} is a single leave with
an allocated "{\tt{}content}" component of type {\tt{}muli{\char95}trapezium{\char95}t}. If
"{\tt{}content}" is not type-compatible with {\tt{}muli{\char95}trapezium{\char95}t}, then
this whole conversion cannot succeed. We allocate a new node of type
{\tt{}muli{\char95}trapezium{\char95}list{\char95}t}. This list does not contain the content of
{\tt{}fib{\char95}tree}, it {\em is} a copy of the content, for
{\tt{}muli{\char95}trapezium{\char95}list{\char95}t} is an extension of {\tt{}muli{\char95}trapezium{\char95}t}. In
the next step, each branch of {\tt{}fib{\char95}tree} is a single leave. We could
call this soubroutine for each branch, but we do copy and paste for
each branch instead. Finally we append one list to the other, the
lowest value of {\tt{}r{\char95}position} comes first. 
\nwenddocs{}\nwbegincode{1834}\moddef{Muli aq: procedures}\plusendmoddef
  recursive subroutine fibonacci_tree_resort_and_convert_to_trapezium_list &
       (fib_tree, lin_list)
    class(fibonacci_node_t), intent(in) :: fib_tree
    class(fibonacci_node_t),pointer :: leave
    class(muli_trapezium_list_t), pointer, intent(out) :: lin_list
    class(muli_trapezium_list_t), pointer :: left_list, right_list
    class(muli_trapezium_node_class_t), pointer :: &
         left_node, right_node, last_node
    class(measure_class_t), pointer :: content
    if (fib_tree%depth > 1) then
       ! print *,"3A"
       call fibonacci_tree_resort_and_convert_to_trapezium_list &
            (fib_tree%left, left_list)
       call fibonacci_tree_resort_and_convert_to_trapezium_list &
            (fib_tree%right, right_list)
       if (left_list%is_left_of (right_list)) then
          lin_list => left_list
          call left_list%get_right (left_node)
          right_node => right_list
       else
          lin_list => right_list
          left_node => left_list
          call right_list%get_right (right_node)
       end if
       last_node => lin_list
       do while (associated (left_node) .and. associated (right_node))          
          if (left_node%is_left_of (right_node)) then
             call last_node%append (left_node)
             call last_node%get_right (last_node)
             call left_node%get_right (left_node)
          else
             call last_node%append (right_node)
             call last_node%get_right (last_node)
             call right_node%get_right (right_node)
          end if
       end do
       if (associated (left_node)) then
          call last_node%append (left_node)
       else
          call last_node%append (right_node)
       end if
       !!! It's done.
       ! print *,"3E"
    else
       if (fib_tree%depth == 0) then
          ! print *,"1A"
          select type (fib_tree)
          class is (fibonacci_leave_t)
             call fib_tree%get_content (content)
             select type (content)
             class is (muli_trapezium_t)
                call content%to_node (content%get_r_position(), list=lin_list)
             class default
                call msg_fatal &
                     ("fibonacci_tree_resort_and_convert_to_trapezium_list: " // &
                     "Content of fibonacci_tree is not type compatible " // &
                     "to muli_trapezium_t")
             end select
          end select
          ! print *,"1E"
       else
          ! print *,"2A"
          leave => fib_tree%left
          select type (leave)
          class is (fibonacci_leave_t)
             call leave%get_content (content)
             select type (content)
             class is (muli_trapezium_t)
                call content%to_node (content%get_r_position(), list=left_list)
             class default
                call msg_fatal &
                     ("fibonacci_tree_resort_and_convert_to_trapezium_list: " // &
                     "Content of fibonacci_tree is not type compatible " // &
                     "to muli_trapezium_t")
             end select
          end select
          leave => fib_tree%right
          select type (leave)
          class is (fibonacci_leave_t)
             call leave%get_content (content)
             select type (content)
             class is (muli_trapezium_t)
                call content%to_node (content%get_r_position(), list=right_list)
             class default
                call msg_fatal &
                     ("fibonacci_tree_resort_and_convert_to_trapezium_list: " // &
                     "Content of fibonacci_tree is not type compatible " // &
                     "to muli_trapezium_t")
             end select
          end select
          if (left_list%is_left_of (right_list)) then
             call left_list%append (right_list)
             lin_list => left_list
          else
             call right_list%append (left_list)
             lin_list => right_list
          end if
          ! print *,"2E"
       end if
    end if
    ! call lin_list%print_all ()
    ! call lin_list%check ()
  end subroutine fibonacci_tree_resort_and_convert_to_trapezium_list
  
\eatline
\nwendcode{}\nwbegindocs{1835}\nwdocspar
\nwenddocs{}\nwbegincode{1836}\moddef{Muli aq: interfaces}\endmoddef
  interface
     subroutine evaluate_if (this, x, y)
       use kinds !NODEP!
       import aq_class
       class(aq_class), intent(inout) :: this
       real(default), intent(in) :: x
       real(default), intent(out) , dimension(:) :: y
     end subroutine evaluate_if

     ! subroutine evaluate_ratios_if (this, cont)
     !   use kinds
     !   use lin_approx_tree_module, only: muli_trapezium_t
     !   import aq_class
     !   class(aq_class) :: this
     !   class(muli_trapezium_t), intent(inout), pointer :: cont
     ! end subroutine evaluate_ratios_if
  end interface
  
\eatline
\nwendcode{}\nwbegindocs{1837}\nwdocspar
%%%%%%%%%%%%%%%%%%%%%%%%%%%%%%%%%%%%%%%%%%%%%%%%%%%%%%%%%%%%%%%%%%%%%%%%
\nwenddocs{}\nwbegincode{1838}\moddef{\code{}muli{\char95}parameters.f90\edoc{}}\endmoddef
! This is a dummy for muli_parameters_module
module muli_parameters_module
end module muli_parameters_module

\eatline
\nwendcode{}\nwbegindocs{1839}\nwdocspar
%%%%%%%%%%%%%%%%%%%%%%%%%%%%%%%%%%%%%%%%%%%%%%%%%%%%%%%%%%%%%%%%%%%%%%%%
\section{Integrands for Multiple Interactions}

This file contains the module {\tt{}muli{\char95}dsigma}. Its only type
{\tt{}muli{\char95}dsigma{\char95}t} provides an integrand to {\tt{}aq{\char95}class}. The actual
integrand is the normalized differential cross section of a QCD $2\to
2$ process $1/\sigma_0 \times d^3 \sigma / ( d p_T^2  d x_1  d x_2)$.  
We need a root function of this integrand in terms of $p_T$, so we have
to integrate out $x_1$ and $x_2$ and have to approximate the root
function of the leftover variable $p_T$. Integration of $x_1$ and
$x_2$ is done by CUBA, the root function is approximated by
{\tt{}muli{\char95}aq}. 

\nwenddocs{}\nwbegincode{1840}\moddef{\code{}muli{\char95}dsigma.f90\edoc{}}\endmoddef
\LA{}File header\RA{}

module muli_dsigma
\LA{}Use kinds\RA{}
  use constants
  use muli_momentum
  use muli_interactions
  use muli_base
  use muli_cuba
  use muli_aq

\LA{}Standard module head\RA{}
  
\LA{}Muli dsigma: variables\RA{}

\LA{}Muli dsigma: public\RA{}

\LA{}Muli dsigma: types\RA{}
  
contains
  
\LA{}Muli dsigma: procedures\RA{}  

end module muli_dsigma

\eatline
\nwendcode{}\nwbegindocs{1841}\nwdocspar
\nwenddocs{}\nwbegincode{1842}\moddef{Muli dsigma: variables}\endmoddef
  integer, parameter :: dim_f = 17
  
\eatline
\nwendcode{}\nwbegindocs{1843}\nwdocspar
\nwenddocs{}\nwbegincode{1844}\moddef{Muli dsigma: public}\endmoddef
  public :: muli_dsigma_t
\nwendcode{}\nwbegincode{1845}\moddef{Muli dsigma: types}\endmoddef
  type, extends (aq_class) :: muli_dsigma_t
     private
     type(transverse_mom_t) :: pt
     type(cuba_divonne_t) :: cuba_int
   contains
   \LA{}Muli dsigma: dsigma: TBP\RA{}  
  end type muli_dsigma_t

\eatline
\nwendcode{}\nwbegindocs{1846}\nwdocspar
\nwenddocs{}\nwbegincode{1847}\moddef{Muli dsigma: dsigma: TBP}\endmoddef
  procedure :: write_to_marker => muli_dsigma_write_to_marker
\nwendcode{}\nwbegincode{1848}\moddef{Muli dsigma: procedures}\endmoddef
  subroutine muli_dsigma_write_to_marker (this, marker, status)
    class(muli_dsigma_t), intent(in) :: this
    class(marker_t), intent(inout) :: marker
    integer(dik), intent(out) :: status
    class(ser_class_t), pointer :: ser
    call marker%mark_begin ("muli_dsigma_t")
    call aq_write_to_marker (this, marker, status)
    call this%cuba_int%serialize (marker, "cuba_int")
    call marker%mark_end ("muli_dsigma_t")
  end subroutine muli_dsigma_write_to_marker
  
\nwendcode{}\nwbegindocs{1849} %def muli_dsigma_write_to_marker
\nwenddocs{}\nwbegindocs{1850}\nwdocspar
\nwenddocs{}\nwbegincode{1851}\moddef{Muli dsigma: dsigma: TBP}\plusendmoddef
  procedure :: read_from_marker => muli_dsigma_read_from_marker
\nwendcode{}\nwbegincode{1852}\moddef{Muli dsigma: procedures}\plusendmoddef
  subroutine muli_dsigma_read_from_marker (this, marker, status)
    class(muli_dsigma_t), intent(out) :: this
    class(marker_t), intent(inout) :: marker
    integer(dik), intent(out) :: status
    call marker%pick_begin ("muli_dsigma_t", status=status)
    call aq_read_from_marker (this, marker, status)
    call this%cuba_int%deserialize ("cuba_int", marker)
    call marker%pick_end ("muli_dsigma_t", status)
  end subroutine muli_dsigma_read_from_marker

\eatline
\nwendcode{}\nwbegindocs{1853}\nwdocspar
\nwenddocs{}\nwbegincode{1854}\moddef{Muli dsigma: dsigma: TBP}\plusendmoddef
  procedure :: print_to_unit => muli_dsigma_print_to_unit  
\nwendcode{}\nwbegincode{1855}\moddef{Muli dsigma: procedures}\plusendmoddef
  subroutine muli_dsigma_print_to_unit &
       (this, unit, parents, components, peers)
    class(muli_dsigma_t), intent(in) :: this
    integer, intent(in) :: unit
    integer(dik), intent(in) :: parents, components, peers
    integer :: ite
    if (parents > 0)  &
         call aq_print_to_unit (this, unit, parents-1, components, peers)
    write (unit, "(A)")  "Components of muli_dsigma_t"
    if (components > 0) then
       write (unit, "(A)")  "Printing components of cuba_int:"
       call this%cuba_int%print_to_unit (unit, parents, components-1, peers)
    else
       write (unit, "(A)")  "Skipping components of cuba_int:"
    end if
  end subroutine muli_dsigma_print_to_unit
  
\eatline
\nwendcode{}\nwbegindocs{1856}\nwdocspar
\nwenddocs{}\nwbegincode{1857}\moddef{Muli dsigma: dsigma: TBP}\plusendmoddef
  procedure, nopass :: get_type => muli_dsigma_get_type    
\nwendcode{}\nwbegindocs{1858}\nwdocspar
\nwenddocs{}\nwbegincode{1859}\moddef{Muli dsigma: procedures}\plusendmoddef
  pure subroutine muli_dsigma_get_type (type)
    character(:), allocatable, intent(out) :: type
    allocate (type, source="muli_dsigma_t")
  end subroutine muli_dsigma_get_type

\eatline
\nwendcode{}\nwbegindocs{1860}\nwdocspar
\nwenddocs{}\nwbegincode{1861}\moddef{Muli dsigma: dsigma: TBP}\plusendmoddef
  procedure :: generate => muli_dsigma_generate
\nwendcode{}\nwbegincode{1862}\moddef{Muli dsigma: procedures}\plusendmoddef
  subroutine muli_dsigma_generate (this, gev2_scale_cutoff, gev2_s, int_tree)
    class(muli_dsigma_t), intent(inout) :: this
    real(default), intent(in) :: gev2_scale_cutoff, gev2_s
    type(muli_trapezium_tree_t), intent(out) :: int_tree
    real(default), dimension(ceiling (log (gev2_s/gev2_scale_cutoff)/two)) :: &
         initial_values
    integer :: n
    print *, gev2_s/gev2_scale_cutoff, &
         ceiling (log (gev2_s/gev2_scale_cutoff)/two)
    ! allocate (initial_values (ceiling (-log (gev2_scale_cutoff))/2))
    ! allocate (real(default), &
    !    dimension (ceiling (log(gev2_scale_cutoff))/2) :: initial_values)
    initial_values(1) = sqrt(gev2_scale_cutoff/gev2_s) * two
    do n = 2, size(initial_values) - 1
       initial_values(n) = initial_values(n-1) * euler
    end do
    initial_values(n) = one
    print *, initial_values
    ! stop
    call this%initialize (i_one, "dsigma")
    call this%pt%initialize (gev2_s)
    this%abs_error_goal = zero
    this%rel_error_goal = scale(one, -12) !-12
    this%max_nodes = 1000
    call this%cuba_int%set_common (dim_f=dim_f, dim_x=2, &
         eps_rel=scale(this%rel_error_goal,-8), flags = 0)
    call this%cuba_int%set_deferred (xgiven_flat = [1.E-2_default, &
         5.E-1_default + epsilon(1._default), 1.E-2_default, &
         5.E-1_default - epsilon(1._default)])
    print *, "muli_dsigma_generate:"
    ! print *, "Cuba Error Goal:    ", this%cuba_int%eps_rel
    print *, "Overall Error Goal: ", this%rel_error_goal
    call this%init_error_tree (dim_f, initial_values)
    call this%run ()
    call this%integrate (int_tree)
    call this%err_tree%deallocate_all ()
    deallocate (this%err_tree)
    nullify (this%int_list)
  end subroutine muli_dsigma_generate

\nwendcode{}\nwbegindocs{1863} %def muli_dsigma_generate  
\nwenddocs{}\nwbegindocs{1864}\nwdocspar
\nwenddocs{}\nwbegincode{1865}\moddef{Muli dsigma: dsigma: TBP}\plusendmoddef
  procedure :: evaluate => muli_dsigma_evaluate
\nwendcode{}\nwbegincode{1866}\moddef{Muli dsigma: procedures}\plusendmoddef
  subroutine muli_dsigma_evaluate (this, x, y)
    class(muli_dsigma_t), intent(inout) :: this
    real(default), intent(in) :: x
    real(default), intent(out), dimension(:):: y
    call this%pt%set_unit_scale (x)
    ! print *, "muli_dsigma_evaluate x=", x
    ! call this%cuba_int%integrate_userdata &
    !       (interactions_proton_proton_integrand_param_17_reg, this%pt)
    ! if (this%cuba_int%fail == 0) then
    ! call this%cuba_int%print_all ()
    call this%cuba_int%get_integral_array (y)
    ! else
    !    print *, "muli_dsigma_evaluate: failed."
    !    stop
    ! end if
  end subroutine muli_dsigma_evaluate
  
\eatline
\nwendcode{}\nwbegindocs{1867}\nwdocspar
\nwenddocs{}\nwbegincode{1868}\moddef{Muli dsigma: dsigma: TBP}\plusendmoddef
  generic :: initialize => muli_dsigma_initialize
  procedure :: muli_dsigma_initialize
\nwendcode{}\nwbegincode{1869}\moddef{Muli dsigma: procedures}\plusendmoddef
  subroutine muli_dsigma_initialize &
       (this, id, name, goal, max_nodes, dim, cuba_goal)
    class(muli_dsigma_t), intent(inout) :: this
    integer(dik), intent(in) :: id, max_nodes
    integer, intent(in) :: dim
    character(*), intent(in) :: name
    real(default), intent(in) :: goal, cuba_goal
    call this%initialize (id,name)
    ! 1E-4      
    this%rel_error_goal = goal
    this%max_nodes = max_nodes
    call this%cuba_int%set_common (dim_f=dim, dim_x=2, &
         ! 1E-6           
         eps_rel=cuba_goal, flags = 0)
    call this%cuba_int%set_deferred (xgiven_flat = [1.E-2_default, &
         5.E-1_default + epsilon(1._default), &
         1.E-2_default, 5.E-1_default - epsilon(1._default)])
    ! call aq_initialize (this, id, name, d_goal, max_nodes, dim_f, &
    !        [8E-1_default/7E3_default, 2E-3_default, 1E-2_default, &
    !         1E-1_default, one])
    call this%init_error_tree (dim, [8.E-1_default/7.E3_default, &
         2.E-3_default, 1.E-2_default, 1.E-1_default, &
         1._default])
    this%is_deferred_initialised = .true.
  end subroutine muli_dsigma_initialize

\eatline
\nwendcode{}\nwbegindocs{1870}\nwdocspar
\nwenddocs{}\nwbegincode{1871}\moddef{Muli dsigma: dsigma: TBP}\plusendmoddef
  ! procedure :: reset => muli_dsigma_reset  
\nwendcode{}\nwbegincode{1872}\moddef{Muli dsigma: procedures}\plusendmoddef
  ! subroutine muli_dsigma_reset (this)
  !   class(muli_dsigma_t), intent(inout) :: this
  !   call aq_reset (this)
  !   call this%initialize &
  !      (id, name, d_goal, max_nodes, dim_f, init, cuba_goal)
  ! end subroutine muli_dsigma_reset
  
\eatline
\nwendcode{}\nwbegindocs{1873}\nwdocspar
%%%%%%%%%%%%%%%%%%%%%%%%%%%%%%%%%%%%%%%%%%%%%%%%%%%%%%%%%%%%%%%%%%%%%%%%
\section{MC Integrations for QCD $2\to 2$ processes}


This file contains the module {\tt{}muli{\char95}mcint} which is the Monte Carlo
generator for QCD $2\to 2$ interactions at given evolution parameter
and given stratus. While {\tt{}muli{\char95}t} takes care of generating the
evolution parameter and the stratus, this module is about bookkeeping
the strati and implementing a downstream importance sampling. The
evolution parameter is a measure of transferred momentum and a ``stratus''
is a PDF category or, to be more precise, is whether the incoming partons
are gluons or sea quarks or valence quarks.

The importance sampling then subdivides the phase space of variables
$\left\{ x_1, x_2, p_T \right\}$ into $n^3$ regions such that each
region holds approximately $n$ interactions. Thus, we can generate a
phase space point very quickly just by randomly picking a region,
randomly picking a point within this region and comparing its exact 
cross section with the mean cross section for this actual evolution
parameter and the actual stratus times the area of the picked phase space
region.

The mean values must be generated in the module {\tt{}muli{\char95}dsigma} before
and are given to the procedure
{\tt{}sample{\char95}inclusive{\char95}generate{\char95}hit}. Finally the generated subregions 
should be written to a file via {\tt{}write{\char95}to{\char95}marker} and then reused
for each later \whizard\ run. 

The type {\tt{}sample{\char95}inclusive{\char95}t} holds the 16 strati, while the type
{\tt{}sample{\char95}int{\char95}kind{\char95}t} represents a single stratus, {\tt{}sample{\char95}3d{\char95}t} is
the whole $\left\{ x_1, x_2, p_T \right\}$ phase space for each  
stratus, {\tt{}sample{\char95}2d{\char95}t} is the $\left\{ x_1, x_2 \right\}$ plane with
a slice of $p_T$ and {\tt{}sample{\char95}region{\char95}t} finally is a phase space
region. 

\nwenddocs{}\nwbegincode{1874}\moddef{\code{}muli{\char95}mcint.f90\edoc{}}\endmoddef
\LA{}File header\RA{}

module muli_mcint
\LA{}Use kinds\RA{}  
  use constants
  use muli_base
  use tao_random_numbers !NODEP!
  use muli_interactions

\LA{}Standard module head\RA{}  

\LA{}Muli MC int: variables\RA{}

\LA{}Muli MC int: public\RA{}

\LA{}Muli MC int: types\RA{}
  
contains

\LA{}Muli MC int: procedures\RA{}

end module muli_mcint

\eatline
\nwendcode{}\nwbegindocs{1875}\nwdocspar
\nwenddocs{}\nwbegincode{1876}\moddef{Muli MC int: variables}\endmoddef
  integer, parameter :: max_n = 2**30
\nwendcode{}\nwbegincode{1877}\moddef{Muli MC int: variables}\plusendmoddef
  real(default), parameter :: max_d = one * max_n
\nwendcode{}\nwbegincode{1878}\moddef{Muli MC int: variables}\plusendmoddef
  real(default), parameter, dimension(2,2) :: &
       unit_square = reshape([zero,zero,one,one], [2,2])

\eatline
\nwendcode{}\nwbegindocs{1879}\nwdocspar
\nwenddocs{}\nwbegincode{1880}\moddef{Muli MC int: public}\endmoddef
  public :: sample_region_t
\nwendcode{}\nwbegincode{1881}\moddef{Muli MC int: types}\endmoddef
  type, extends (ser_class_t) :: sample_region_t
     integer :: n_hits = 0
     integer :: n_alloc = 0
     real(default), dimension(2,2) :: corners = unit_square
     real(default), dimension(:,:), allocatable :: hyp_hits
   contains
   \LA{}Muli MC int: sample region: TBP\RA{}  
  end type sample_region_t
  
\eatline
\nwendcode{}\nwbegindocs{1882}\nwdocspar
\nwenddocs{}\nwbegincode{1883}\moddef{Muli MC int: sample region: TBP}\endmoddef
  procedure :: write_to_marker => sample_region_write_to_marker
\nwendcode{}\nwbegincode{1884}\moddef{Muli MC int: procedures}\endmoddef
  subroutine sample_region_write_to_marker (this, marker, status)
    class(sample_region_t), intent(in) :: this
    class(marker_t), intent(inout) :: marker
    integer(dik), intent(out) :: status
    integer :: n
    call marker%mark_begin ("sample_region_t")
    call marker%mark ("n_hits", this%n_hits)
    call marker%mark ("n_alloc", this%n_alloc)
    call marker%mark ("lower_corner", this%corners(1:2,1))
    call marker%mark ("upper_corner", this%corners(1:2,2))
    if (allocated (this%hyp_hits)) then
       call marker%mark ("hyp_hits", this%hyp_hits (1:3,:this%n_hits))
    else
       call marker%mark_nothing ("hyp_hits")
    end if
    call marker%mark_end ("sample_region_t")
  end subroutine sample_region_write_to_marker
  
\eatline
\nwendcode{}\nwbegindocs{1885}\nwdocspar
\nwenddocs{}\nwbegincode{1886}\moddef{Muli MC int: sample region: TBP}\plusendmoddef
  procedure :: read_from_marker => sample_region_read_from_marker
\nwendcode{}\nwbegincode{1887}\moddef{Muli MC int: procedures}\plusendmoddef
  subroutine sample_region_read_from_marker (this, marker, status)
    class(sample_region_t), intent(out) :: this
    class(marker_t), intent(inout) :: marker
    integer(dik), intent(out) :: status  
    integer :: n
    call marker%pick_begin ("sample_region_t", status=status)
    call marker%pick ("n_hits", this%n_hits, status)
    call marker%pick ("n_alloc", this%n_alloc, status)
    call marker%pick ("lower_corner", this%corners(1:2,1), status)
    call marker%pick ("upper_corner", this%corners(1:2,2), status)
    if (allocated (this%hyp_hits))  deallocate (this%hyp_hits)
    call marker%verify_nothing ("hyp_hits", status)
    if (.not. status == serialize_nothing) then
       allocate (this%hyp_hits (3,this%n_alloc))
       call marker%pick ("hyp_hits", this%hyp_hits (1:3,:this%n_hits), status)
    end if
    call marker%pick_end ("sample_region_t", status)
  end subroutine sample_region_read_from_marker
  
\eatline
\nwendcode{}\nwbegindocs{1888}\nwdocspar
\nwenddocs{}\nwbegincode{1889}\moddef{Muli MC int: sample region: TBP}\plusendmoddef
  procedure :: print_to_unit => sample_region_print_to_unit  
\nwendcode{}\nwbegincode{1890}\moddef{Muli MC int: procedures}\plusendmoddef
  subroutine sample_region_print_to_unit &
       (this, unit, parents, components, peers)
    class(sample_region_t), intent(in) :: this
    integer, intent(in) :: unit
    integer(dik), intent(in) :: parents, components, peers
    write (unit, "(1x,A)")      "components of sample_region_t"
    write (unit, "(3x,A,I10)")  "n_hits:           ", this%n_hits
    write (unit, "(3x,A,I10)")  "n_alloc:          ", this%n_alloc
    write (unit, "(3x,4(E20.10))") "corners:          ", this%corners
    if (allocated (this%hyp_hits) .and. this%n_hits > 0) then
       if (components > 0) then
          write (unit,"(3x,A)")  "hits:"
          print *, shape (this%hyp_hits)
          write (unit, "(3(e20.10))")  this%hyp_hits (1:3, this%n_hits)
       else
          write (unit, "(3x,A)")  "skipping hits."
       end if
    else
       write (unit, "(3x,A)")  "hits are not allocated."
    end if
  end subroutine sample_region_print_to_unit
  
\eatline
\nwendcode{}\nwbegindocs{1891}\nwdocspar
\nwenddocs{}\nwbegincode{1892}\moddef{Muli MC int: sample region: TBP}\plusendmoddef
  procedure, nopass :: get_type => sample_region_get_type  
\nwendcode{}\nwbegincode{1893}\moddef{Muli MC int: procedures}\plusendmoddef
  pure subroutine sample_region_get_type (type)
    character(:), allocatable, intent(out) :: type
    allocate (type, source="sample_region_t")
  end subroutine sample_region_get_type

\eatline
\nwendcode{}\nwbegindocs{1894}\nwdocspar
\nwenddocs{}\nwbegincode{1895}\moddef{Muli MC int: sample region: TBP}\plusendmoddef
  procedure :: initialize => sample_region_initialize
\nwendcode{}\nwbegincode{1896}\moddef{Muli MC int: procedures}\plusendmoddef
  subroutine sample_region_initialize (this, n_alloc)
    class(sample_region_t), intent(out) :: this
    integer, intent(in) :: n_alloc
    if (allocated (this%hyp_hits))  deallocate (this%hyp_hits)
    allocate (this%hyp_hits (3,n_alloc))
    this%n_alloc = n_alloc
  end subroutine sample_region_initialize
  
\eatline
\nwendcode{}\nwbegindocs{1897}\nwdocspar
\nwenddocs{}\nwbegincode{1898}\moddef{Muli MC int: sample region: TBP}\plusendmoddef
  procedure :: generate_hit => sample_region_generate_hit
\nwendcode{}\nwbegincode{1899}\moddef{Muli MC int: procedures}\plusendmoddef
  pure subroutine sample_region_generate_hit (this, rnd, area, hit)
    class(sample_region_t), intent(in) :: this
    integer, intent(in), dimension(2) :: rnd
    real(default), dimension(2), intent(out) :: hit
    real(default), intent(out) :: area
    call muli_mcint_generate_hit (rnd, this%corners, hit)
    area = this%area ()
  end subroutine sample_region_generate_hit
  
\eatline
\nwendcode{}\nwbegindocs{1900}\nwdocspar
\nwenddocs{}\nwbegincode{1901}\moddef{Muli MC int: sample region: TBP}\plusendmoddef
  procedure :: confirm_hit => sample_region_confirm_hit
\nwendcode{}\nwbegincode{1902}\moddef{Muli MC int: procedures}\plusendmoddef
  subroutine sample_region_confirm_hit (this, hit)
    class(sample_region_t), intent(inout) :: this
    real(default), dimension(3), intent(in) :: hit
    ! print *,"sample_region_confirm_hit: ", this%n_hits, this%n_alloc, hit
    this%n_hits = this%n_hits + 1
    if (this%n_hits <= this%n_alloc) then
       this%hyp_hits (1:3, this%n_hits) = hit
    else
       print *, "sample_region_confirm_hit: Region is already full."
    end if
  end subroutine sample_region_confirm_hit
  
\eatline
\nwendcode{}\nwbegindocs{1903}\nwdocspar
\nwenddocs{}\nwbegincode{1904}\moddef{Muli MC int: sample region: TBP}\plusendmoddef
  procedure :: split => sample_region_split  
\nwendcode{}\nwbegincode{1905}\moddef{Muli MC int: procedures}\plusendmoddef
  subroutine sample_region_split (this, pos, dimX, n_alloc, lower, upper)
    class(sample_region_t), intent(in) :: this
    type(sample_region_t), intent(out) :: lower, upper
    real(default), dimension(3) :: hit
    real(default), intent(in) :: pos
    integer, intent(in) :: dimX, n_alloc
    integer :: n_hit    
    call lower%initialize (n_alloc)
    call upper%initialize (n_alloc)
    do n_hit = 1, this%n_hits
       hit = this%hyp_hits (1:3, n_hit)
       if (hit(dimX) < pos) then
          call lower%confirm_hit (hit)
       else
          call upper%confirm_hit (hit)
       end if
    end do
    lower%corners = this%corners
    upper%corners = this%corners
    if (dimX < 3) then
       lower%corners(dimX,2) = pos
       upper%corners(dimX,1) = pos
    end if
  end subroutine sample_region_split
  
\eatline
\nwendcode{}\nwbegindocs{1906}\nwdocspar
\nwenddocs{}\nwbegincode{1907}\moddef{Muli MC int: sample region: TBP}\plusendmoddef
  procedure :: write_hits => sample_region_write_hits
\nwendcode{}\nwbegincode{1908}\moddef{Muli MC int: procedures}\plusendmoddef
  subroutine sample_region_write_hits (this, unit)
    class(sample_region_t), intent(in) :: this
    integer, intent(in) :: unit
    integer :: n
    do n = 1, this%n_hits
       write (unit, *)  this%hyp_hits (1:3,n)
    end do
  end subroutine sample_region_write_hits
  
\eatline
\nwendcode{}\nwbegindocs{1909}\nwdocspar
\nwenddocs{}\nwbegincode{1910}\moddef{Muli MC int: sample region: TBP}\plusendmoddef
  procedure :: is_full => sample_region_is_full  
\nwendcode{}\nwbegincode{1911}\moddef{Muli MC int: procedures}\plusendmoddef
  elemental logical function sample_region_is_full (this)
    class(sample_region_t), intent(in) :: this
    sample_region_is_full = this%n_alloc == this%n_hits
  end function sample_region_is_full
  
\eatline
\nwendcode{}\nwbegindocs{1912}\nwdocspar
\nwenddocs{}\nwbegincode{1913}\moddef{Muli MC int: sample region: TBP}\plusendmoddef
  procedure :: move_components => sample_region_move_components
\nwendcode{}\nwbegincode{1914}\moddef{Muli MC int: procedures}\plusendmoddef
  subroutine sample_region_move_components (this, that)
    class(sample_region_t), intent(inout) :: this
    class(sample_region_t), intent(out) :: that
    that%n_alloc = this%n_alloc
    that%n_hits = this%n_hits
    that%corners = this%corners
    call move_alloc (this%hyp_hits, that%hyp_hits)
    this%n_alloc = 0
    this%n_hits = 0
  end subroutine sample_region_move_components
  
\eatline
\nwendcode{}\nwbegindocs{1915}\nwdocspar
\nwenddocs{}\nwbegincode{1916}\moddef{Muli MC int: sample region: TBP}\plusendmoddef
  procedure :: mean => sample_region_mean
\nwendcode{}\nwbegincode{1917}\moddef{Muli MC int: procedures}\plusendmoddef
  elemental function sample_region_mean (this, dim)
    real(default) :: sample_region_mean
    class(sample_region_t), intent(in) :: this
    integer, intent(in) :: dim
    sample_region_mean = sum (this%hyp_hits (dim,1:this%n_hits)) / this%n_hits
  end function sample_region_mean
  
\eatline
\nwendcode{}\nwbegindocs{1918}\nwdocspar
\nwenddocs{}\nwbegincode{1919}\moddef{Muli MC int: sample region: TBP}\plusendmoddef
  procedure :: area => sample_region_area  
\nwendcode{}\nwbegincode{1920}\moddef{Muli MC int: procedures}\plusendmoddef
  elemental function sample_region_area (this)
    real(default) :: sample_region_area
    class(sample_region_t), intent(in) :: this
    sample_region_area = product (this%corners(1:2,2) - this%corners(1:2,1))
  end function sample_region_area
  
\eatline
\nwendcode{}\nwbegindocs{1921}\nwdocspar
\nwenddocs{}\nwbegincode{1922}\moddef{Muli MC int: sample region: TBP}\plusendmoddef
  procedure :: density => sample_region_density
\nwendcode{}\nwbegincode{1923}\moddef{Muli MC int: procedures}\plusendmoddef
  elemental function sample_region_density (this)
    real(default) :: sample_region_density
    class(sample_region_t), intent(in) :: this
    sample_region_density = this%n_hits / this%area ()
  end function sample_region_density

\eatline
\nwendcode{}\nwbegindocs{1924}\nwdocspar
\nwenddocs{}\nwbegincode{1925}\moddef{Muli MC int: sample region: TBP}\plusendmoddef
  procedure :: contains => sample_region_contains
\nwendcode{}\nwbegincode{1926}\moddef{Muli MC int: procedures}\plusendmoddef
  pure logical function sample_region_contains (this, hit)
    class(sample_region_t), intent(in) :: this
    real(default), intent(in), dimension(3) :: hit
    sample_region_contains = (this%corners(1,1) <= hit(1) .and. &
         hit(1) <= this%corners(1,2) .and. &
         this%corners(2,1) <= hit(2) .and. &
         hit(2) <= this%corners(2,2))
  end function sample_region_contains
  
\eatline
\nwendcode{}\nwbegindocs{1927}\nwdocspar
\nwenddocs{}\nwbegincode{1928}\moddef{Muli MC int: sample region: TBP}\plusendmoddef
  procedure :: to_generator => sample_region_to_generator  
\nwendcode{}\nwbegincode{1929}\moddef{Muli MC int: procedures}\plusendmoddef
  subroutine sample_region_to_generator (this)
    class(sample_region_t), intent(inout) :: this
    if (allocated (this%hyp_hits))  deallocate (this%hyp_hits)
    this%n_alloc = 0
  end subroutine sample_region_to_generator
  
\eatline
\nwendcode{}\nwbegindocs{1930}\nwdocspar
\nwenddocs{}\nwbegincode{1931}\moddef{Muli MC int: public}\plusendmoddef
  public :: sample_2d_t
\nwendcode{}\nwbegincode{1932}\moddef{Muli MC int: types}\plusendmoddef
  type, extends (ser_class_t) :: sample_2d_t
     integer :: n_regions = 0
     integer :: n_alloc = 0
     integer :: n_hits = 0
     real(default), dimension(2) :: range = [0,1]
     type(sample_region_t), dimension(:), allocatable :: regions
   contains
   \LA{}Muli MC int: sample 2D: TBP\RA{}
  end type sample_2d_t
  
\eatline
\nwendcode{}\nwbegindocs{1933}\nwdocspar
\nwenddocs{}\nwbegincode{1934}\moddef{Muli MC int: sample 2D: TBP}\endmoddef
  procedure :: write_to_marker => sample_2d_write_to_marker
\nwendcode{}\nwbegincode{1935}\moddef{Muli MC int: procedures}\plusendmoddef
  subroutine sample_2d_write_to_marker (this, marker, status)
    class(sample_2d_t), intent(in) :: this
    class(marker_t), intent(inout) :: marker
    integer(dik), intent(out) :: status
    integer :: n
    call marker%mark_begin ("sample_2d_t")
    call marker%mark ("n_regions", this%n_regions)
    call marker%mark ("n_alloc", this%n_alloc)
    call marker%mark ("n_hits", this%n_hits)
    call marker%mark ("range", this%range)
    if (this%n_regions > 0) then
       call marker%mark_instance_begin &
            (this%regions(1), name="sample_2d_t", shape=shape (this%regions))
       do n = 1, this%n_regions
          call sample_region_write_to_marker (this%regions(n), marker, status)
       end do
       call marker%mark_instance_end ()
    end if
    call marker%mark_end ("sample_2d_t")
  end subroutine sample_2d_write_to_marker
  
\eatline
\nwendcode{}\nwbegindocs{1936}\nwdocspar
\nwenddocs{}\nwbegincode{1937}\moddef{Muli MC int: sample 2D: TBP}\plusendmoddef
  procedure :: read_from_marker => sample_2d_read_from_marker
\nwendcode{}\nwbegincode{1938}\moddef{Muli MC int: procedures}\plusendmoddef
  subroutine sample_2d_read_from_marker (this, marker, status)
    class(sample_2d_t), intent(out) :: this
    class(marker_t), intent(inout) :: marker
    integer(dik), intent(out) :: status  
    integer :: n
    call marker%pick_begin ("sample_2d_t", status=status)
    call marker%pick ("n_regions", this%n_regions, status)
    call marker%pick ("n_alloc", this%n_alloc, status)
    call marker%pick ("n_hits", this%n_hits, status)
    call marker%pick ("range", this%range, status)
    if (this%n_regions > 0) then
       call marker%pick_begin ("regions", status=status)
       allocate (this%regions (this%n_regions))
       do n = 1, this%n_regions
          call sample_region_read_from_marker (this%regions(n), marker, status)
       end do
       call marker%pick_end ("regions", status)
    end if
    call marker%pick_end ("sample_2d_t", status)
  end subroutine sample_2d_read_from_marker
  
\eatline
\nwendcode{}\nwbegindocs{1939}\nwdocspar
\nwenddocs{}\nwbegincode{1940}\moddef{Muli MC int: sample 2D: TBP}\plusendmoddef
  procedure :: print_to_unit => sample_2d_print_to_unit
\nwendcode{}\nwbegincode{1941}\moddef{Muli MC int: procedures}\plusendmoddef
  subroutine sample_2d_print_to_unit (this, unit, parents, components, peers)
    class(sample_2d_t), intent(in) :: this
    integer, intent(in) :: unit
    integer(dik), intent(in) :: parents, components, peers
    integer :: n
    write (unit, "(1x,A)")           "components of sample_2d_t"
    write (unit, "(3x,A,I10)")       "n_regions:        ", this%n_regions
    write (unit, "(3x,A,I10)")       "n_alloc:          ", this%n_alloc
    write (unit, "(3x,A,2(E20.10))") "range:            ", this%range
    if (allocated (this%regions)) then
       if (components > 0) then
          write (unit, "(3x,A)")  "regions:"
          do n = 1, this%n_regions
             call this%regions(n)%print_to_unit &
                  (unit, parents, components-1, peers)
          end do
       else
          write (unit, "(3x,A)")  "skipping regions."
       end if
    else
       write (unit, "(3x,A)")  "regions are not allocated."
    end if
  end subroutine sample_2d_print_to_unit
  
\eatline
\nwendcode{}\nwbegindocs{1942}\nwdocspar
\nwenddocs{}\nwbegincode{1943}\moddef{Muli MC int: sample 2D: TBP}\plusendmoddef
  procedure, nopass :: get_type => sample_2d_get_type    
\nwendcode{}\nwbegincode{1944}\moddef{Muli MC int: procedures}\plusendmoddef
  pure subroutine sample_2d_get_type (type)
    character(:), allocatable, intent(out) :: type
    allocate (type, source="sample_2d_t")
  end subroutine sample_2d_get_type
  
\eatline
\nwendcode{}\nwbegindocs{1945}\nwdocspar
\nwenddocs{}\nwbegincode{1946}\moddef{Muli MC int: sample 2D: TBP}\plusendmoddef
  procedure :: initialize => sample_2d_initialize
\nwendcode{}\nwbegincode{1947}\moddef{Muli MC int: procedures}\plusendmoddef
  subroutine sample_2d_initialize (this, n_alloc)
    class(sample_2d_t), intent(out) :: this
    integer, intent(in) :: n_alloc
    integer :: n
    if (allocated (this%regions))  deallocate (this%regions)
    allocate (this%regions (n_alloc))
    this%n_alloc = n_alloc
    this%n_regions = 1
    call this%regions(1)%initialize (n_alloc)
    ! do n = 1, n_alloc
    !    call this%regions(n)%initialize (n_alloc)
    ! end do
  end subroutine sample_2d_initialize
  
\eatline
\nwendcode{}\nwbegindocs{1948}\nwdocspar
\nwenddocs{}\nwbegincode{1949}\moddef{Muli MC int: sample 2D: TBP}\plusendmoddef
  procedure :: contains => sample_2d_contains
\nwendcode{}\nwbegincode{1950}\moddef{Muli MC int: procedures}\plusendmoddef
  pure logical function sample_2d_contains (this, pts2)
    class(sample_2d_t), intent(in) :: this
    real(default), intent(in) :: pts2
    sample_2d_contains = this%range(1) <= pts2 .and. pts2 <= this%range(2)
  end function sample_2d_contains
  
\eatline
\nwendcode{}\nwbegindocs{1951}\nwdocspar
\nwenddocs{}\nwbegincode{1952}\moddef{Muli MC int: sample 2D: TBP}\plusendmoddef
  procedure :: generate_hit => sample_2d_generate_hit
\nwendcode{}\nwbegincode{1953}\moddef{Muli MC int: procedures}\plusendmoddef
  pure subroutine sample_2d_generate_hit (this, rnd, boost, hit, region)
    class(sample_2d_t), intent(in) :: this
    integer, dimension(3), intent(in) :: rnd
    integer, intent(out) :: region
    integer :: n, sum
    real(default), dimension(2), intent(out) :: hit
    real(default), intent(out) :: boost
    if (0 < this%n_hits .and. this%n_hits < 10) then
       !!! this should be improved
       sum = modulo(rnd(1),this%n_hits) + 1  
       region = 0
       do while (sum > 0)
          region = region + 1
          sum = sum - this%regions(region)%n_hits
       end do
       call this%regions(region)%generate_hit (rnd(2:3), boost, hit)
       boost = boost * this%n_hits / this%regions(region)%n_hits
    else
       if (this%n_regions > 1) then
          !!! this should be improved
          region = modulo(rnd(1), this%n_regions) + 1 
          call this%regions(region)%generate_hit (rnd(2:3), boost, hit)
          boost = boost * this%n_regions
       else
          region = 1
          call this%regions(1)%generate_hit (rnd(2:3), boost, hit)
       end if
    end if
  end subroutine sample_2d_generate_hit

!  pure subroutine sample_2d_generate_hit (this, rnd, boost, hit, region)
!    class(sample_2d_t), intent(in) :: this
!    integer, dimension(3), intent(in) :: rnd
!    integer, intent(out) :: region
!    real(double), dimension(2), intent(out) :: hit
!    real(double), intent(out) :: boost
!    region = modulo(rnd(1), this%n_regions) + 1   !!! this should be improved
!    call this%regions(region)%generate_hit (rnd(2:3), boost, hit)
!    boost = boost * this%n_regions
!  end subroutine sample_2d_generate_hit
    
\eatline
\nwendcode{}\nwbegindocs{1954}\nwdocspar
\nwenddocs{}\nwbegincode{1955}\moddef{Muli MC int: sample 2D: TBP}\plusendmoddef
  procedure :: confirm_hit => sample_2d_confirm_hit
\nwendcode{}\nwbegincode{1956}\moddef{Muli MC int: procedures}\plusendmoddef
  subroutine sample_2d_confirm_hit (this, hit, region, full)
    class(sample_2d_t), intent(inout) :: this
    integer, intent(in) :: region
    real(default), dimension(3), intent(in) :: hit
    type(sample_region_t), allocatable :: old_region
    real(default), dimension(2) :: mean, var, diff, cm, cv, c
    integer :: n, n_alloc, dim 
    logical, intent(out) :: full
    this%n_hits = this%n_hits + 1
    if (region <= this%n_alloc) then
       full = .false.
       call this%regions(region)%confirm_hit (hit)
       n_alloc = this%regions(region)%n_alloc
       if (this%regions(region)%is_full()) then
          if (this%is_full()) then
             full = .true.
          else
             this%n_regions = this%n_regions + 1
             allocate (old_region)
             call this%regions(region)%move_components (old_region)
             mean = sum (old_region%hyp_hits(1:2,:), dim=2) / n_alloc
             var = 0
             do n = 1, n_alloc
                var = var + abs(mean-old_region%hyp_hits(1:2,n))
             end do
             var = var / n_alloc
             diff = old_region%corners(1:2,2) - old_region%corners(1:2,1)
             cm = abs ([0.5_default,0.5_default] - &
                  (old_region%corners(1:2,2) - mean) / diff)
             cv = abs(2*([0.25_default,0.25_default] - var / diff))
             c = max(cm,cv)
             if (c(1) < c(2)) then
                dim = 2
             else
                dim = 1
             end if
             call old_region%split (mean(dim), dim, this%n_alloc, &
                  this%regions(region), this%regions(this%n_regions))
          end if
       end if
    else
       write (*,*) "sample_2d_confirm_hit: Region ", region, &
            " not allocated."
    end if
  end subroutine sample_2d_confirm_hit
  
\eatline
\nwendcode{}\nwbegindocs{1957}\nwdocspar
\nwenddocs{}\nwbegincode{1958}\moddef{Muli MC int: sample 2D: TBP}\plusendmoddef
  procedure :: split => sample_2d_split
\nwendcode{}\nwbegincode{1959}\moddef{Muli MC int: procedures}\plusendmoddef
  recursive subroutine sample_2d_split (this, n_alloc, pos, lower, upper)
    class(sample_2d_t), intent(in) :: this
    integer, intent(in) :: n_alloc
    real(default), intent(in) :: pos
    type(sample_2d_t), intent(out) :: lower, upper
    integer :: n_r, n_h
    real(default), dimension(3) :: hit
    ! print *,"sample_2d_split: ", pos, this%range
    call lower%initialize (4*n_alloc)
    call upper%initialize (4*n_alloc)
    do n_r = this%n_regions, 1, -1
       do n_h = 1, this%regions(n_r)%n_hits
          hit = this%regions(n_r)%hyp_hits (1:3,n_h)
          if (hit(3) > pos) then
             call upper%push (hit)
          else
             call lower%push (hit)             
          end if
       end do
    end do
    lower%range = [this%range(1), pos]
    upper%range = [pos, this%range(2)]
  end subroutine sample_2d_split

!  subroutine sample_2d_split (this, n_alloc, pos, lower, upper)
!    class(sample_2d_t), intent(in) :: this
!    integer, intent(in) :: n_alloc
!    real(default), intent(in) :: pos
!    type(sample_2d_t), intent(out) :: lower, upper
!    integer :: n, n_hit
!    real(default), dimension(3) :: hit
!    allocate (lower%regions (n_alloc))
!    allocate (upper%regions (n_alloc))
!    !$OMP PARALLEL DO FIRSTPRIVATE (this, pos, n_alloc) SHARED (lower, upper)
!    do n = 1, this%n_regions
!      call sample_region_split (this%regions(n), pos, 3, n_alloc, &
!         lower%regions(n),upper%regions(n))
!    end do
!    !$OMP END PARALLEL DO
!    lower%n_regions = this%n_regions
!    upper%n_regions = this%n_regions
!    lower%n_alloc = n_alloc
!    upper%n_alloc = n_alloc
!    lower%range = [this%range(1), pos]
!    upper%range = [pos, this%range(2)]
!  end subroutine sample_2d_split
    
\eatline
\nwendcode{}\nwbegindocs{1960}\nwdocspar
\nwenddocs{}\nwbegincode{1961}\moddef{Muli MC int: sample 2D: TBP}\plusendmoddef
  procedure :: push => sample_2d_push
\nwendcode{}\nwbegincode{1962}\moddef{Muli MC int: procedures}\plusendmoddef
  subroutine sample_2d_push (this, hit)
    class(sample_2d_t), intent(inout) :: this
    real(default), dimension(3), intent(in) :: hit
    integer :: region
    logical :: full
    do region = 1, this%n_regions
       if (this%regions(region)%contains (hit)) then
          call this%confirm_hit (hit, region, full)
          ! call this%regions(region)%confirm_hit (hit)
          if (full) print *,"sample_2d_push: region is full now"
          exit
       end if
    end do
    if (region > this%n_regions) & 
         print *, "sample_2d_push: no region contains ", hit
  end subroutine sample_2d_push
  
\eatline
\nwendcode{}\nwbegindocs{1963}\nwdocspar
\nwenddocs{}\nwbegincode{1964}\moddef{Muli MC int: sample 2D: TBP}\plusendmoddef
  procedure :: write_hits => sample_2d_write_hits
\nwendcode{}\nwbegincode{1965}\moddef{Muli MC int: procedures}\plusendmoddef
  subroutine sample_2d_write_hits (this, unit)
    class(sample_2d_t), intent(in) :: this
    integer, intent(in) :: unit
    integer :: n
    do n = 1, this%n_regions
       call this%regions(n)%write_hits (unit)
    end do
  end subroutine sample_2d_write_hits
  
\eatline
\nwendcode{}\nwbegindocs{1966}\nwdocspar
\nwenddocs{}\nwbegincode{1967}\moddef{Muli MC int: sample 2D: TBP}\plusendmoddef
  procedure :: is_full => sample_2d_is_full
\nwendcode{}\nwbegincode{1968}\moddef{Muli MC int: procedures}\plusendmoddef
  elemental logical function sample_2d_is_full (this)
    class(sample_2d_t), intent(in) :: this
    sample_2d_is_full = this%n_alloc == this%n_regions
  end function sample_2d_is_full
  
\eatline
\nwendcode{}\nwbegindocs{1969}\nwdocspar
\nwenddocs{}\nwbegincode{1970}\moddef{Muli MC int: sample 2D: TBP}\plusendmoddef
  procedure :: move_components => sample_2d_move_components
\nwendcode{}\nwbegincode{1971}\moddef{Muli MC int: procedures}\plusendmoddef
  subroutine sample_2d_move_components (this, that)
    class(sample_2d_t), intent(inout) :: this
    class(sample_2d_t), intent(out) :: that
    that%n_alloc = this%n_alloc
    that%n_regions = this%n_regions
    that%n_hits = this%n_hits
    that%range = this%range
    call move_alloc (this%regions, that%regions)
    this%n_alloc = 0
    this%n_regions = 0
    this%n_hits = 0
    this%range = [zero,zero]
  end subroutine sample_2d_move_components
  
\eatline
\nwendcode{}\nwbegindocs{1972}\nwdocspar
\nwenddocs{}\nwbegincode{1973}\moddef{Muli MC int: sample 2D: TBP}\plusendmoddef
  procedure :: thickness => sample_2d_thickness
\nwendcode{}\nwbegincode{1974}\moddef{Muli MC int: procedures}\plusendmoddef
  elemental function sample_2d_thickness (this)
    class(sample_2d_t), intent(in) :: this
    real(default) :: sample_2d_thickness
    sample_2d_thickness = this%range(2) - this%range(1)
  end function sample_2d_thickness
  
\nwendcode{}\nwbegindocs{1975}\nwdocspar
\nwenddocs{}\nwbegincode{1976}\moddef{Muli MC int: sample 2D: TBP}\plusendmoddef
  procedure :: analyse => sample_2d_analyse
\nwendcode{}\nwbegincode{1977}\moddef{Muli MC int: procedures}\plusendmoddef
  subroutine sample_2d_analyse (this, dir, file)
    class(sample_2d_t), intent(in) :: this
    character(*), intent(in) :: dir, file
    integer :: u
    real(default), dimension(1:2,0:100,0:100) :: grid
    integer, dimension(0:100,0:100) :: i_grid
    integer :: r, x, y
    integer, dimension(2,2) :: i
    call generate_unit (u)
    print *, "sample_2d_analyse: ", dir // "/" // file
    open (u, file=dir//"/"//file)
    do x = 0, 100
       do y = 0, 100
          grid(1:2,x,y) = [-one,-one]
       end do
    end do
    do r = 1, this%n_regions
       i = int(this%regions(r)%corners*1E2_default)
       do x = i(1,1), i(1,2)
          do y = i(2,1), i(2,2)
             i_grid(x,y) = this%regions(r)%n_hits
             grid(1,x,y) = one / this%regions(r)%area ()
             grid(2,x,y) = this%regions(r)%density ()
          end do
       end do
    end do
    do x = 0, 100
       do y = 0, 100
          write (u, *) x, y, i_grid(x,y), grid(1:2,x,y)
       end do
       write (u, *)
    end do
    close (u)
  end subroutine sample_2d_analyse
  
\eatline
\nwendcode{}\nwbegindocs{1978}\nwdocspar
\nwenddocs{}\nwbegincode{1979}\moddef{Muli MC int: sample 2D: TBP}\plusendmoddef
  procedure :: to_generator => sample_2d_to_generator
\nwendcode{}\nwbegincode{1980}\moddef{Muli MC int: procedures}\plusendmoddef
  subroutine sample_2d_to_generator (this)
    class(sample_2d_t), intent(inout) :: this
    integer :: region
    do region = 1, this%n_regions
       call this%regions(region)%to_generator ()
    end do
  end subroutine sample_2d_to_generator
  
\eatline
\nwendcode{}\nwbegindocs{1981}\nwdocspar
\nwenddocs{}\nwbegincode{1982}\moddef{Muli MC int: sample 2D: TBP}\plusendmoddef
  procedure :: mean => sample_2d_mean  
\nwendcode{}\nwbegincode{1983}\moddef{Muli MC int: procedures}\plusendmoddef
  elemental function sample_2d_mean (this, dim) result (mean)
    class(sample_2d_t), intent(in) :: this
    integer, intent(in) :: dim
    real(default) :: mean
    integer :: region, hit
    mean = zero
    do region = 1, this%n_regions
       do hit = 1, this%regions(region)%n_hits
          mean = mean + this%regions(region)%hyp_hits (dim, hit)
       end do
    end do
    mean = mean / this%n_hits
  end function sample_2d_mean
  
\eatline
\nwendcode{}\nwbegindocs{1984}\nwdocspar
\nwenddocs{}\nwbegincode{1985}\moddef{Muli MC int: public}\plusendmoddef
  public :: sample_3d_t
\nwendcode{}\nwbegincode{1986}\moddef{Muli MC int: types}\plusendmoddef
  type, extends (ser_class_t) :: sample_3d_t
     integer::n_slices=0
     integer::n_alloc=0
     type(sample_2d_t), dimension(:),allocatable::slices     
   contains
   \LA{}Muli MC int: sample 3D: TBP\RA{}
  end type sample_3d_t
  
\eatline
\nwendcode{}\nwbegindocs{1987}\nwdocspar
\nwenddocs{}\nwbegincode{1988}\moddef{Muli MC int: sample 3D: TBP}\endmoddef
  procedure :: write_to_marker => sample_3d_write_to_marker
\nwendcode{}\nwbegincode{1989}\moddef{Muli MC int: procedures}\plusendmoddef
  subroutine sample_3d_write_to_marker (this, marker, status)
    class(sample_3d_t), intent(in) :: this
    class(marker_t), intent(inout) :: marker
    integer(dik), intent(out) :: status 
    integer :: n
    call marker%mark_begin ("sample_3d_t")
    call marker%mark ("n_slices", this%n_slices)
    call marker%mark ("n_alloc", this%n_alloc)
    if (this%n_slices > 0) then
       call marker%mark_instance_begin &
            (this%slices(1), "slices", shape=shape(this%slices))
       do n = 1, this%n_slices
          call sample_2d_write_to_marker (this%slices(n), marker, status)
       end do
       call marker%mark_instance_end ()
    end if
    call marker%mark_end ("sample_3d_t")
  end subroutine sample_3d_write_to_marker
  
\eatline
\nwendcode{}\nwbegindocs{1990}\nwdocspar
\nwenddocs{}\nwbegincode{1991}\moddef{Muli MC int: sample 3D: TBP}\plusendmoddef
  procedure :: read_from_marker => sample_3d_read_from_marker
\nwendcode{}\nwbegincode{1992}\moddef{Muli MC int: procedures}\plusendmoddef
  subroutine sample_3d_read_from_marker (this, marker, status)
    class(sample_3d_t), intent(out) :: this
    class(marker_t), intent(inout) :: marker
    integer(dik), intent(out) :: status  
    integer :: n
    call marker%pick_begin ("sample_3d_t", status=status)
    call marker%pick ("n_slices", this%n_slices, status)
    call marker%pick ("n_alloc", this%n_alloc, status)
    if (this%n_slices > 0) then
       call marker%pick_instance_begin ("slices", status=status)
       allocate(this%slices (this%n_slices))
       do n = 1, this%n_slices
          call sample_2d_read_from_marker (this%slices(n), marker, status)
       end do
       call marker%pick_instance_end (status)
    end if
    call marker%pick_end ("sample_3d_t", status)
  end subroutine sample_3d_read_from_marker
  
\eatline
\nwendcode{}\nwbegindocs{1993}\nwdocspar
\nwenddocs{}\nwbegincode{1994}\moddef{Muli MC int: sample 3D: TBP}\plusendmoddef
  procedure :: print_to_unit => sample_3d_print_to_unit
\nwendcode{}\nwbegincode{1995}\moddef{Muli MC int: procedures}\plusendmoddef
  subroutine sample_3d_print_to_unit (this, unit, parents, components, peers)
    class(sample_3d_t), intent(in) :: this
    integer, intent(in) :: unit
    integer(dik), intent(in) :: parents, components, peers
    integer :: n
    write (unit, "(1x,A)")     "components of sample_3d_t"
    write (unit, "(3x,A,I10)") "n_slices:         ", this%n_slices
    write (unit, "(3x,A,I10)") "n_alloc:          ", this%n_alloc
    if (allocated (this%slices)) then
       if (components > 0) then
          do n = 1, this%n_slices
             call this%slices(n)%print_to_unit (unit, parents, components-1, peers)
          end do
       else
          write (unit, "(3x,A)")  "skipping slices."
       end if
    else
       write (unit, "(3x,A)")  "slices are not allocated."
    end if
  end subroutine sample_3d_print_to_unit
  
\eatline
\nwendcode{}\nwbegindocs{1996}\nwdocspar
\nwenddocs{}\nwbegincode{1997}\moddef{Muli MC int: sample 3D: TBP}\plusendmoddef
  procedure, nopass :: get_type => sample_3d_get_type  
\nwendcode{}\nwbegincode{1998}\moddef{Muli MC int: procedures}\plusendmoddef
  pure subroutine sample_3d_get_type (type)
    character(:), allocatable, intent(out) :: type
    allocate (type, source="sample_3d_t")
  end subroutine sample_3d_get_type
  
\eatline
\nwendcode{}\nwbegindocs{1999}\nwdocspar
\nwenddocs{}\nwbegincode{2000}\moddef{Muli MC int: sample 3D: TBP}\plusendmoddef
  procedure :: measure => sample_3d_measure  
\nwendcode{}\nwbegincode{2001}\moddef{Muli MC int: procedures}\plusendmoddef
  elemental function sample_3d_measure (this)
    real(default) :: sample_3d_measure
    class(sample_3d_t), intent(in) :: this
    sample_3d_measure = one
  end function sample_3d_measure
  
\eatline
\nwendcode{}\nwbegindocs{2002}\nwdocspar
\nwenddocs{}\nwbegincode{2003}\moddef{Muli MC int: sample 3D: TBP}\plusendmoddef
  procedure :: to_generator => sample_3d_to_generator
\nwendcode{}\nwbegincode{2004}\moddef{Muli MC int: procedures}\plusendmoddef
  subroutine sample_3d_to_generator(this)
    class(sample_3d_t), intent(inout)::this
    integer::slice
    do slice=1,this%n_slices
       call this%slices(slice)%to_generator()
    end do
  end subroutine sample_3d_to_generator
  
\eatline
\nwendcode{}\nwbegindocs{2005}\nwdocspar
\nwenddocs{}\nwbegincode{2006}\moddef{Muli MC int: sample 3D: TBP}\plusendmoddef
  generic :: initialize => sample_3d_initialize  
  procedure :: sample_3d_initialize
\nwendcode{}\nwbegincode{2007}\moddef{Muli MC int: procedures}\plusendmoddef
  subroutine sample_3d_initialize (this, n_alloc)
    class(sample_3d_t), intent(out) :: this
    integer, intent(in) :: n_alloc
    if (allocated (this%slices))  deallocate (this%slices)
    if (n_alloc > 0) then
       allocate (this%slices (n_alloc))
       this%n_alloc = n_alloc
       this%n_slices = 1
       call this%slices(1)%initialize (n_alloc)
    else
       this%n_alloc = 0
    end if
  end subroutine sample_3d_initialize
  
\eatline
\nwendcode{}\nwbegindocs{2008}\nwdocspar
\nwenddocs{}\nwbegincode{2009}\moddef{Muli MC int: sample 3D: TBP}\plusendmoddef
  procedure :: sample_3d_generate_hit
  generic :: generate_hit => sample_3d_generate_hit    
\nwendcode{}\nwbegincode{2010}\moddef{Muli MC int: procedures}\plusendmoddef
  pure subroutine sample_3d_generate_hit &
       (this, rnd, pts2, boost, hit, region, slice)
    class(sample_3d_t), intent(in) :: this
    integer, intent(in), dimension(3) :: rnd
    real(default), intent(in) :: pts2
    integer, intent(out) :: slice, region
    real(default), dimension(3), intent(out) :: hit
    real(default), intent(out) :: boost
    if (this%n_slices == 0) then
       call muli_mcint_generate_hit (rnd, unit_square, hit(1:2))
       boost = 1._default
       slice = 1
       region = 1
    else
       do slice = 1, this%n_slices
          if (this%slices(slice)%contains (pts2))  exit          
       end do
       call this%slices(slice)%generate_hit (rnd, boost, hit(1:2), region)
    end if
    hit(3) = pts2
  end subroutine sample_3d_generate_hit
  
\eatline
\nwendcode{}\nwbegindocs{2011}\nwdocspar
\nwenddocs{}\nwbegincode{2012}\moddef{Muli MC int: sample 3D: TBP}\plusendmoddef
  procedure :: sample_3d_confirm_hit
  generic :: confirm_hit => sample_3d_confirm_hit    
\nwendcode{}\nwbegincode{2013}\moddef{Muli MC int: procedures}\plusendmoddef
  subroutine sample_3d_confirm_hit (this, hit, region, slice)
    class(sample_3d_t), intent(inout) :: this
    integer, intent(in) :: slice, region
    real(default), intent(in), dimension(3) :: hit
    type(sample_2d_t), allocatable :: old_slice
    integer :: n
    logical :: full
    if (this%n_alloc < slice) then
       print *, "sample_3d_confirm_hit: Slice ", slice, " not allocated."
    else
    ! if (.not. allocated (this%slices))  call this%initialize (2)
       call this%slices(slice)%confirm_hit (hit, region, full)
       if (full) then
          if (this%n_alloc == this%n_slices)  call this%enlarge ()
          this%n_slices = this%n_slices + 1
          allocate (old_slice)
          call this%slices(slice)%move_components (old_slice)
          call sample_2d_split (old_slice, this%n_alloc, &
               old_slice%mean(3), this%slices(slice), &
               this%slices(this%n_slices))
       end if
    end if
  end subroutine sample_3d_confirm_hit
  
\eatline
\nwendcode{}\nwbegindocs{2014}\nwdocspar
\nwenddocs{}\nwbegincode{2015}\moddef{Muli MC int: sample 3D: TBP}\plusendmoddef
  procedure :: enlarge => sample_3d_enlarge
\nwendcode{}\nwbegincode{2016}\moddef{Muli MC int: procedures}\plusendmoddef
  subroutine sample_3d_enlarge (this)
    class(sample_3d_t), intent(inout) :: this
    type(sample_2d_t), allocatable, dimension(:) :: old_slices
    integer :: n
    print *, "sample_3d_enlarge"
    call move_alloc (this%slices, old_slices)
    this%n_alloc = this%n_alloc * 2
    allocate (this%slices (this%n_alloc))
    do n = 1, size(old_slices)
       call old_slices(n)%move_components(this%slices(n))
    end do
  end subroutine sample_3d_enlarge
  
\eatline
\nwendcode{}\nwbegindocs{2017}\nwdocspar
\nwenddocs{}\nwbegincode{2018}\moddef{Muli MC int: public}\plusendmoddef
  public :: sample_int_kind_t
\nwendcode{}\nwbegincode{2019}\moddef{Muli MC int: types}\plusendmoddef
  type, extends (sample_3d_t) :: sample_int_kind_t
     integer :: n_proc = 0
     integer(kind=i64) :: n_tries = 0
     integer :: n_hits = 0
     integer :: n_over = 0
     integer, dimension(:), allocatable :: hits, weights, processes
     real(default) :: overall_boost = 1E-1_default
   contains
   \LA{}Muli MC int: sample int: TBP\RA{}  
  end type sample_int_kind_t
  
\eatline
\nwendcode{}\nwbegindocs{2020}\nwdocspar
\nwenddocs{}\nwbegincode{2021}\moddef{Muli MC int: sample int: TBP}\endmoddef
  procedure :: write_to_marker => sample_int_kind_write_to_marker
\nwendcode{}\nwbegincode{2022}\moddef{Muli MC int: procedures}\plusendmoddef
  subroutine sample_int_kind_write_to_marker (this, marker, status)
    class(sample_int_kind_t), intent(in) :: this
    class(marker_t), intent(inout) :: marker
    integer(dik), intent(out) :: status
    call marker%mark_begin ("sample_int_kind_t")
    call sample_3d_write_to_marker (this, marker, status)
    call marker%mark ("n_hits", this%n_hits)
    call marker%mark ("n_proc", this%n_proc)
    call marker%mark ("boost", this%overall_boost)
    if (this%n_hits > 0) then
       call marker%mark ("hits", this%hits)
    end if
    if (this%n_proc > 0) then
       call marker%mark ("processes", this%processes)
       call marker%mark ("weights", this%weights)
    end if
    call marker%mark_end ("sample_int_kind_t")
  end subroutine sample_int_kind_write_to_marker
  
\eatline
\nwendcode{}\nwbegindocs{2023}\nwdocspar
\nwenddocs{}\nwbegincode{2024}\moddef{Muli MC int: sample int: TBP}\plusendmoddef
  procedure :: read_from_marker => sample_int_kind_read_from_marker
\nwendcode{}\nwbegincode{2025}\moddef{Muli MC int: procedures}\plusendmoddef
  subroutine sample_int_kind_read_from_marker (this, marker, status)
    class(sample_int_kind_t), intent(out) :: this
    class(marker_t), intent(inout) :: marker
    integer(dik), intent(out) :: status  
    call marker%pick_begin ("sample_int_kind_t", status=status)
    call sample_3d_read_from_marker (this, marker, status)
    call marker%pick ("n_hits", this%n_hits, status)
    call marker%pick ("n_proc", this%n_proc, status)
    call marker%pick ("boost", this%overall_boost, status)
    if (this%n_hits > 0) then
       allocate (this%hits (this%n_hits))
       call marker%pick ("hits", this%hits, status)
    end if
    if (this%n_proc > 0) then
       allocate (this%processes (this%n_proc))
       call marker%pick ("processes", this%processes, status)
       allocate (this%weights (this%n_proc))
       call marker%pick ("weights", this%weights, status)
    end if
    call marker%pick_end ("sample_int_kind_t", status)
  end subroutine sample_int_kind_read_from_marker
  
\eatline
\nwendcode{}\nwbegindocs{2026}\nwdocspar
\nwenddocs{}\nwbegincode{2027}\moddef{Muli MC int: sample int: TBP}\plusendmoddef
  procedure :: print_to_unit => sample_int_kind_print_to_unit
\nwendcode{}\nwbegincode{2028}\moddef{Muli MC int: procedures}\plusendmoddef

  subroutine sample_int_kind_print_to_unit (this, unit, parents, components, peers)
    class(sample_int_kind_t), intent(in) :: this
    integer, intent(in) :: unit
    integer(dik), intent(in) :: parents, components, peers
    integer :: n
    if (parents > 0)  call sample_3d_print_to_unit &
         (this, unit, parents, components, peers)
    write (unit, "(1x,A)")       "components of sample_int_kind_t"
    write (unit, "(3x,A,I10)")   "n_hits:          ", this%n_hits
    write (unit, "(3x,A,I10)")   "n_proc:          ", this%n_proc
    write (unit, "(3x,A,E14.7)") "overall_boost:   ", this%overall_boost
    write (unit, "(3x,A)")       "hits:"
    write (unit, "(3x,10(I0,1x))")  this%hits(1:this%n_hits)
    write (unit, "(3x,A)")       "weights:"
    write (unit, "(3x,10(I0,1x))")  this%weights
    write (unit, "(3x,A)")  "processes:"
    write (unit, "(3x,2(I0,1x))") this%processes
  end subroutine sample_int_kind_print_to_unit
  
\eatline
\nwendcode{}\nwbegindocs{2029}\nwdocspar
\nwenddocs{}\nwbegincode{2030}\moddef{Muli MC int: sample int: TBP}\plusendmoddef
  procedure, nopass :: get_type => sample_int_kind_get_type
\nwendcode{}\nwbegincode{2031}\moddef{Muli MC int: procedures}\plusendmoddef
  pure subroutine sample_int_kind_get_type (type)
    character(:),allocatable, intent(out) :: type
    allocate (type, source="sample_int_kind_t")
  end subroutine sample_int_kind_get_type
  
\eatline
\nwendcode{}\nwbegindocs{2032}\nwdocspar
\nwenddocs{}\nwbegincode{2033}\moddef{Muli MC int: sample int: TBP}\plusendmoddef
  procedure :: to_generator => sample_int_kind_to_generator  
\nwendcode{}\nwbegincode{2034}\moddef{Muli MC int: procedures}\plusendmoddef
  subroutine sample_int_kind_to_generator(this)
    class(sample_int_kind_t), intent(inout)::this
    integer::int_kind
    if (allocated(this%hits))deallocate(this%hits)
    call sample_3d_to_generator(this)
  end subroutine sample_int_kind_to_generator

\eatline
\nwendcode{}\nwbegindocs{2035}\nwdocspar
\nwenddocs{}\nwbegincode{2036}\moddef{Muli MC int: sample int: TBP}\plusendmoddef
  procedure :: process_id => sample_int_kind_process_id
\nwendcode{}\nwbegincode{2037}\moddef{Muli MC int: procedures}\plusendmoddef
  elemental integer function sample_int_kind_process_id (this, subprocess)
    class(sample_int_kind_t), intent(in) :: this
    integer, intent(in) :: subprocess
    sample_int_kind_process_id = this%processes(subprocess)
  end function sample_int_kind_process_id
  
\eatline
\nwendcode{}\nwbegindocs{2038}\nwdocspar
\nwenddocs{}\nwbegincode{2039}\moddef{Muli MC int: sample int: TBP}\plusendmoddef
  procedure :: sample_int_kind_initialize
  generic :: initialize => sample_int_kind_initialize  
\nwendcode{}\nwbegincode{2040}\moddef{Muli MC int: procedures}\plusendmoddef
  subroutine sample_int_kind_initialize (this, n_alloc, processes, overall_boost)
    class(sample_int_kind_t), intent(out) :: this
    integer, intent(in) :: n_alloc
    integer, intent(in), dimension(:) :: processes
    real(default), optional, intent(in) :: overall_boost
    integer :: s, n
    s = size(processes)
    call sample_3d_initialize (this, n_alloc)
    if (allocated (this%hits))  deallocate (this%hits)
    allocate (this%hits (n_alloc))
    if (allocated (this%weights))  deallocate (this%weights)
    allocate (this%weights(s))
    if (allocated (this%processes))  deallocate (this%processes)
    allocate (this%processes(s), source=processes)
    do n = 1, s
       this%weights(n) = 0
    end do
    this%n_alloc = n_alloc
    this%n_hits = 0
    this%n_proc = s
    if (present (overall_boost))  this%overall_boost = overall_boost
    this%overall_boost = this%overall_boost * this%n_proc
    ! print *, this%weights
  end subroutine sample_int_kind_initialize
  
\nwendcode{}\nwbegindocs{2041}%sample_int_kind_initialize
\nwenddocs{}\nwbegindocs{2042}\nwdocspar
\nwenddocs{}\nwbegincode{2043}\moddef{Muli MC int: sample int: TBP}\plusendmoddef
  procedure :: sample_int_kind_generate_hit
\nwendcode{}\nwbegincode{2044}\moddef{Muli MC int: procedures}\plusendmoddef
  pure subroutine sample_int_kind_generate_hit &
       (this, rnd, pts2, boost, hit, region, slice, subprocess)
    class(sample_int_kind_t), intent(in) :: this
    integer, dimension(4), intent(in) :: rnd
    real(default), intent(in) :: pts2
    real(default), dimension(3), intent(out) :: hit
    integer, intent(out) :: region, slice, subprocess
    real(default), intent(out) :: boost
    integer :: n_n
    ! print *, rnd, pts2, boost, hit, region, slice, subprocess
    call sample_3d_generate_hit &
         (this, rnd(2:4), pts2, boost, hit, region, slice)
    n_n = modulo(rnd(1), this%n_hits + size(this%weights)) + 1
    if (n_n > this%n_hits) then
       subprocess = n_n - this%n_hits
    else
       subprocess = this%hits(n_n)
    end if
    boost = boost * this%overall_boost * (this%n_proc + this%n_hits) / &
         (this%n_proc * (this%weights(subprocess) + 1))
  end subroutine sample_int_kind_generate_hit
  
\eatline
\nwendcode{}\nwbegindocs{2045}\nwdocspar
\nwenddocs{}\nwbegincode{2046}\moddef{Muli MC int: sample int: TBP}\plusendmoddef
  procedure :: mcgenerate_hit => sample_int_kind_mcgenerate_hit
  generic :: generate_hit => sample_int_kind_generate_hit
\nwendcode{}\nwbegincode{2047}\moddef{Muli MC int: procedures}\plusendmoddef
  subroutine sample_int_kind_mcgenerate_hit (this, pts2, mean, &
       integrand_kind, tao_rnd, process_id, cart_hit)
    class(sample_int_kind_t), intent(inout) :: this
    integer, intent(in) :: integrand_kind
    real(default), intent(in) :: pts2, mean
    type(tao_random_state), intent(inout) :: tao_rnd
    real(default), dimension(3), intent(out) :: cart_hit
    integer, intent(out) :: process_id
    real(default) :: boost
    integer :: region, slice, subprocess
    integer, dimension(4) :: i_rnd
    real(default) :: dddsigma, d_rnd
    real(default), dimension(3) :: hyp_hit
    MC: do 
       this%n_tries = this%n_tries + 1
       call tao_random_number (tao_rnd, i_rnd)
       call tao_random_number (tao_rnd, d_rnd)
       ! print *, pts2, mean, integrand_kind, process_id, cart_hit
       call this%generate_hit (i_rnd, pts2, boost, hyp_hit, region, &
            slice, subprocess)
       process_id = this%process_id(subprocess)
       call interactions_dddsigma_reg (process_id, integrand_kind, &
            hyp_hit, cart_hit, dddsigma)
       dddsigma = dddsigma * boost
       if (d_rnd*mean < dddsigma) then
          exit MC
       end if
    end do MC
    if (mean < dddsigma) then
       call this%confirm_hit (hyp_hit, region, slice, subprocess, .true.)
    else
       call this%confirm_hit (hyp_hit, region, slice, subprocess, .false.)
    end if
  end subroutine sample_int_kind_mcgenerate_hit
  
\eatline
\nwendcode{}\nwbegindocs{2048}\nwdocspar
\nwenddocs{}\nwbegincode{2049}\moddef{Muli MC int: sample int: TBP}\plusendmoddef
  procedure :: sample_int_kind_confirm_hit
  generic :: confirm_hit => sample_int_kind_confirm_hit    
\nwendcode{}\nwbegincode{2050}\moddef{Muli MC int: procedures}\plusendmoddef
  subroutine sample_int_kind_confirm_hit &
       (this, hit, region, slice, subprocess, over)
    class(sample_int_kind_t), intent(inout) :: this
    real(default), dimension(3), intent(in) :: hit
    integer, intent(in) :: region, slice, subprocess
    integer, dimension(:), allocatable :: tmp_hits
    logical, optional, intent(in) :: over
    this%n_hits = this%n_hits + 1
    if (present(over)) then
       if (over) then
          this%n_over = this%n_over + 1
          this%overall_boost = this%overall_boost / 1.1_default
       else
          this%overall_boost = this%overall_boost * 1.0001_default
       end if
    end if
    if (0 < size(this%hits)) then
       if (this%n_hits > size(this%hits)) then
          call move_alloc (this%hits, tmp_hits)
          allocate (this%hits (2*size(tmp_hits)))
          this%hits (1:size(tmp_hits)) = tmp_hits
       end if
       this%hits(this%n_hits) = subprocess
    end if
    this%weights(subprocess) = this%weights(subprocess) + 1
    call sample_3d_confirm_hit (this, hit, region, slice)
  end subroutine sample_int_kind_confirm_hit
  
\eatline
\nwendcode{}\nwbegindocs{2051}\nwdocspar
\nwenddocs{}\nwbegincode{2052}\moddef{Muli MC int: sample int: TBP}\plusendmoddef
  procedure :: analyse => sample_int_kind_analyse
\nwendcode{}\nwbegincode{2053}\moddef{Muli MC int: procedures}\plusendmoddef
  subroutine sample_int_kind_analyse (this, dir, prefix)
    class(sample_int_kind_t), intent(in) :: this
    character(*), intent(in) :: dir, prefix
    integer :: slices_unit, subprocs_unit
    integer :: n, slice
    character(3) :: slice_name
    integer, dimension(:), allocatable :: int_a
    real(default), dimension(:), allocatable :: real_a
    call generate_unit (slices_unit)
    print *, "sample_int_kind_analyse: ", dir // "/" // prefix // &
         "slice_distribution.plot"
    open (slices_unit, file=dir // "/" // prefix // "slice_distribution.plot")
    call generate_unit (subprocs_unit)
    print *, "sample_int_kind_analyse: ", dir // "/" // prefix // &
         "subproc_distribution.plot"
    open (subprocs_unit, file=dir // "/" // prefix // &
         "subproc_distribution.plot")
    allocate (real_a (this%n_slices))
    allocate (int_a (this%n_slices))
    do n = 1, this%n_slices
       real_a(n) = this%slices(n)%range(1)
    end do
    call misc_sort (real_a, int_a)
    do n = 1, size (this%weights)
       if (this%n_hits > 0) then
          write (subprocs_unit, fmt=*) real(this%weights(n)), &
               real(this%weights(n)+1) / this%n_hits
       else
          write (subprocs_unit, fmt=*) 0, 0
       end if
    end do
    do n = 1, this%n_slices       
       slice = int_a (n)
       call integer_with_leading_zeros (n, 3, slice_name)
       call sample_2d_analyse (this%slices(slice), dir, prefix // &
            slice_name // ".plot")  
       print *, this%n_hits, this%slices(slice)%range(2) - &
            this%slices(slice)%range(1)
       if (this%n_hits > 0) then
          write (slices_unit, *)  this%slices(slice)%range(1), &
               this%slices(slice)%range(2), this%slices(slice)%n_hits, &
               real (this%slices(slice)%n_hits) / (this%n_hits * &
               (this%slices(slice)%range(2) - this%slices(slice)%range(1)))
       else
          write (slices_unit, *)  this%slices(slice)%range(1), &
               this%slices(slice)%range(2), this%slices(slice)%n_hits, zero
       end if
    end do
    write (slices_unit, *) one, zero, zero, zero
    close (slices_unit)
    close (subprocs_unit)
  end subroutine sample_int_kind_analyse
  
\eatline
\nwendcode{}\nwbegindocs{2054}\nwdocspar
\nwenddocs{}\nwbegincode{2055}\moddef{Muli MC int: public}\plusendmoddef
  public :: sample_inclusive_t
\nwendcode{}\nwbegincode{2056}\moddef{Muli MC int: types}\plusendmoddef
  type, extends (ser_class_t) :: sample_inclusive_t
     integer :: n_alloc = 0
     integer(kind=i64) :: n_tries_sum = i_zero
     integer(kind=i64) :: n_over_sum = i_zero
     integer(kind=i64) :: n_hits_sum = i_zero
     type(sample_int_kind_t), dimension(:), allocatable :: int_kinds
   contains
   \LA{}Muli MC int: sample inclusive: TBP\RA{}  
  end type sample_inclusive_t
  
\eatline
\nwendcode{}\nwbegindocs{2057}\nwdocspar
\nwenddocs{}\nwbegincode{2058}\moddef{Muli MC int: sample inclusive: TBP}\endmoddef
  procedure :: write_to_marker => sample_inclusive_write_to_marker
\nwendcode{}\nwbegincode{2059}\moddef{Muli MC int: procedures}\plusendmoddef
  subroutine sample_inclusive_write_to_marker (this, marker, status)
    class(sample_inclusive_t), intent(in) :: this
    class(marker_t), intent(inout) :: marker
    integer(dik), intent(out) :: status
    integer :: n
    call marker%mark_begin ("sample_inclusive_t")
    call marker%mark ("n_alloc", this%n_alloc)
    if (allocated (this%int_kinds)) then
       call marker%mark_begin (tag="int_kinds", shape=shape(this%int_kinds))
       do n = 1, size(this%int_kinds)
          call this%int_kinds(n)%write_to_marker (marker, status)
       end do
       call marker%mark_instance_end ()
    else
       call marker%mark_empty (tag="int_kinds", shape=[0])
    end if
    call marker%mark_end ("sample_inclusive_t")
  end subroutine sample_inclusive_write_to_marker
  
\eatline
\nwendcode{}\nwbegindocs{2060}\nwdocspar
\nwenddocs{}\nwbegincode{2061}\moddef{Muli MC int: sample inclusive: TBP}\plusendmoddef
  procedure :: read_from_marker => sample_inclusive_read_from_marker
\nwendcode{}\nwbegincode{2062}\moddef{Muli MC int: procedures}\plusendmoddef
  subroutine sample_inclusive_read_from_marker (this, marker, status)
    class(sample_inclusive_t), intent(out) :: this
    class(marker_t), intent(inout) :: marker
    integer(dik), intent(out) :: status
    integer :: n
    integer, dimension(:), allocatable :: s
    call marker%pick_begin ("sample_inclusive_t", status=status)
    call marker%pick ("n_alloc", this%n_alloc, status)
    call marker%pick_begin ("int_kinds", shape=s, status=status)
    if (s(1) > 0) then
       do n = 1, size(this%int_kinds)
          call this%int_kinds(n)%read_from_marker (marker, status)
       end do
       call marker%pick_end ("int_kinds",status)
    end if
    call marker%pick_end ("sample_inclusive_t", status)
  end subroutine sample_inclusive_read_from_marker
  
\eatline
\nwendcode{}\nwbegindocs{2063}\nwdocspar
\nwenddocs{}\nwbegincode{2064}\moddef{Muli MC int: sample inclusive: TBP}\plusendmoddef
  procedure :: print_to_unit => sample_inclusive_print_to_unit
\nwendcode{}\nwbegincode{2065}\moddef{Muli MC int: procedures}\plusendmoddef
  subroutine sample_inclusive_print_to_unit (this, unit, parents, components, peers)
    class(sample_inclusive_t), intent(in) :: this
    integer, intent(in) :: unit
    integer(dik), intent(in) :: parents, components, peers
    integer :: n
    write (unit, "(1x,A)")  "components of sample_inclusive_t"
    write (unit, "(3x,A,I10)")  "n_alloc:          ", this%n_alloc
    if (allocated (this%int_kinds)) then
       if (components > 0) then
          write (unit, "(3x,A)")  "int_kinds:"
          do n = 1, this%n_alloc
             call this%int_kinds(n)%print_to_unit &
                  (unit, parents, components-1, peers)
          end do
       else
          write (unit, "(3x,A)")  "skipping int_kinds."
       end if
    else
       write (unit, "(3x,A)")  "int_kinds are not allocated."
    end if
  end subroutine sample_inclusive_print_to_unit
  
\eatline
\nwendcode{}\nwbegindocs{2066}\nwdocspar
\nwenddocs{}\nwbegincode{2067}\moddef{Muli MC int: sample inclusive: TBP}\plusendmoddef
  procedure, nopass :: get_type => sample_inclusive_get_type  
\nwendcode{}\nwbegincode{2068}\moddef{Muli MC int: procedures}\plusendmoddef
  pure subroutine sample_inclusive_get_type (type)
    character(:), allocatable, intent(out) :: type
    allocate (type, source = "sample_inclusive_t")
  end subroutine sample_inclusive_get_type
  
\eatline
\nwendcode{}\nwbegindocs{2069}\nwdocspar
\nwenddocs{}\nwbegincode{2070}\moddef{Muli MC int: sample inclusive: TBP}\plusendmoddef
  procedure :: process_id => sample_inclusive_process_id 
\nwendcode{}\nwbegincode{2071}\moddef{Muli MC int: procedures}\plusendmoddef
  elemental integer function sample_inclusive_process_id &
       (this, subprocess, int_kind)
    class(sample_inclusive_t), intent(in) :: this
    integer, intent(in) :: subprocess, int_kind
    sample_inclusive_process_id = &
         this%int_kinds(int_kind)%processes (subprocess)
  end function sample_inclusive_process_id
  
\eatline
\nwendcode{}\nwbegindocs{2072}\nwdocspar
\nwenddocs{}\nwbegincode{2073}\moddef{Muli MC int: sample inclusive: TBP}\plusendmoddef
  procedure :: initialize => sample_inclusive_initialize
\nwendcode{}\nwbegincode{2074}\moddef{Muli MC int: procedures}\plusendmoddef
  subroutine sample_inclusive_initialize &
       (this, n_alloc, sizes, processes, overall_boost)
    class(sample_inclusive_t), intent(out) :: this
    integer, intent(in) :: n_alloc
    integer, dimension(:), intent(in) :: sizes, processes
    real(default), optional, intent(in) :: overall_boost
    integer :: n, sum
    this%n_tries_sum = i_zero
    this%n_over_sum = 0
    this%n_alloc = size(sizes)
    if (allocated(this%int_kinds))  deallocate (this%int_kinds)
    allocate (this%int_kinds (this%n_alloc))
    sum = 0
    do n = 1, this%n_alloc
       call this%int_kinds(n)%initialize (n_alloc, &
            processes(sum+1:sum+sizes(n)), overall_boost)
       sum = sum + sizes(n)
    end do
  end subroutine sample_inclusive_initialize
  
\eatline
\nwendcode{}\nwbegindocs{2075}\nwdocspar
\nwenddocs{}\nwbegincode{2076}\moddef{Muli MC int: sample inclusive: TBP}\plusendmoddef
  procedure :: finalize => sample_inclusive_finalize
\nwendcode{}\nwbegincode{2077}\moddef{Muli MC int: procedures}\plusendmoddef
  subroutine sample_inclusive_finalize (this)
    class(sample_inclusive_t), intent(inout) :: this
    deallocate (this%int_kinds)
    this%n_alloc = 0
  end subroutine sample_inclusive_finalize
  
\eatline
\nwendcode{}\nwbegindocs{2078}\nwdocspar
\nwenddocs{}\nwbegincode{2079}\moddef{Muli MC int: sample inclusive: TBP}\plusendmoddef
  procedure :: generate_hit => sample_inclusive_generate_hit
\nwendcode{}\nwbegincode{2080}\moddef{Muli MC int: procedures}\plusendmoddef
  pure subroutine sample_inclusive_generate_hit &
       (this, rnd, pts2, int_kind, hit, region, boost, slice, process)
    class(sample_inclusive_t), intent(in) :: this
    integer, dimension(4), intent(in) :: rnd
    real(default), intent(in) :: pts2
    integer, intent(in) :: int_kind
    real(default), dimension(3), intent(out) :: hit
    integer, intent(out) :: region, slice, process
    real(default), intent(out) :: boost
    call this%int_kinds(int_kind)%generate_hit &
         (rnd, pts2, boost, hit, region, slice, process)
  end subroutine sample_inclusive_generate_hit
  
\eatline
\nwendcode{}\nwbegindocs{2081}\nwdocspar
\nwenddocs{}\nwbegincode{2082}\moddef{Muli MC int: sample inclusive: TBP}\plusendmoddef
  procedure :: mcgenerate_hit => sample_inclusive_mcgenerate_hit
\nwendcode{}\nwbegincode{2083}\moddef{Muli MC int: procedures}\plusendmoddef
  subroutine sample_inclusive_mcgenerate_hit &
       (this, pts2, mean, integrand_kind, tao_rnd, process_id, cart_hit)
    class(sample_inclusive_t), intent(inout) :: this
    real(default), intent(in) :: pts2, mean
    integer, intent(in) :: integrand_kind
    type(tao_random_state), intent(inout) :: tao_rnd
    real(default), dimension(3), intent(out) :: cart_hit
    integer, intent(out) :: process_id
    ! print *, "sample_inclusive_mcgenerate_hit &
    !   (this,",pts2,mean,integrand_kind,process_id,cart_hit,")"
    ! print *, allocated (this%int_kinds)
    call sample_int_kind_mcgenerate_hit (this%int_kinds(integrand_kind), &
         pts2, mean, integrand_kind, tao_rnd, process_id, cart_hit)
  end subroutine sample_inclusive_mcgenerate_hit
  
\eatline
\nwendcode{}\nwbegindocs{2084}\nwdocspar
\nwenddocs{}\nwbegincode{2085}\moddef{Muli MC int: sample inclusive: TBP}\plusendmoddef
  procedure :: confirm_hit => sample_inclusive_confirm_hit
\nwendcode{}\nwbegincode{2086}\moddef{Muli MC int: procedures}\plusendmoddef
  subroutine sample_inclusive_confirm_hit &
       (this, hit, int_kind, region, slice, process, over)
    class(sample_inclusive_t), intent(inout) :: this
    real(default), dimension(3), intent(in) :: hit
    integer, intent(in) :: int_kind, region, slice, process
    logical, optional, intent(in) :: over
    call this%int_kinds(int_kind)%confirm_hit &
         (hit, region, slice, process, over)
  end subroutine sample_inclusive_confirm_hit
  
\eatline
\nwendcode{}\nwbegindocs{2087}\nwdocspar
\nwenddocs{}\nwbegincode{2088}\moddef{Muli MC int: sample inclusive: TBP}\plusendmoddef
  procedure :: sum_up => sample_inclusive_sum_up
\nwendcode{}\nwbegincode{2089}\moddef{Muli MC int: procedures}\plusendmoddef
  subroutine sample_inclusive_sum_up (this)
    class(sample_inclusive_t), intent(inout) :: this
    integer :: n
    this%n_tries_sum = i_zero
    this%n_hits_sum = i_zero
    this%n_over_sum = i_zero
    do n = 1, this%n_alloc
       this%n_tries_sum = this%n_tries_sum+this%int_kinds(n)%n_tries
       this%n_hits_sum = this%n_hits_sum+this%int_kinds(n)%n_hits
       this%n_over_sum = this%n_over_sum+this%int_kinds(n)%n_over
    end do
  end subroutine sample_inclusive_sum_up
  
\eatline
\nwendcode{}\nwbegindocs{2090}\nwdocspar
\nwenddocs{}\nwbegincode{2091}\moddef{Muli MC int: sample inclusive: TBP}\plusendmoddef
  procedure :: analyse => sample_inclusive_analyse
\nwendcode{}\nwbegincode{2092}\moddef{Muli MC int: procedures}\plusendmoddef
  subroutine sample_inclusive_analyse (this, dir, subdirs)
    class(sample_inclusive_t), intent(in) :: this
    character(*), intent(in) :: dir
    logical, intent(in) :: subdirs
    integer :: inclusive_unit    
    integer :: n, n_hits
    character(2) :: sample_name
    call generate_unit (inclusive_unit)
    open (inclusive_unit, file = dir // "/int_kinds.plot")
    n_hits = 0
    do n = 1, size(this%int_kinds)
       n_hits = n_hits + this%int_kinds(n)%n_hits
    end do
    do n = 1, size(this%int_kinds)
       write (inclusive_unit, *)  n, real(this%int_kinds(n)%n_hits) / n_hits
       call integer_with_leading_zeros (n, 2, sample_name)
       if (subdirs) then
          call sample_int_kind_analyse (this%int_kinds(n), &
               dir // "/" // sample_name, "")
       else
          call sample_int_kind_analyse (this%int_kinds(n), &
               dir, sample_name // "_")
       end if
    end do
    close (inclusive_unit)
  end subroutine sample_inclusive_analyse
  
\eatline
\nwendcode{}\nwbegindocs{2093}\nwdocspar
\nwenddocs{}\nwbegincode{2094}\moddef{Muli MC int: sample inclusive: TBP}\plusendmoddef
  procedure :: to_generator => sample_inclusive_to_generator
\nwendcode{}\nwbegincode{2095}\moddef{Muli MC int: procedures}\plusendmoddef
  subroutine sample_inclusive_to_generator (this)
    class(sample_inclusive_t), intent(inout) :: this
    integer :: int_kind
    do int_kind = 1, size(this%int_kinds)
       call this%int_kinds(int_kind)%to_generator ()
    end do
  end subroutine sample_inclusive_to_generator
  
\eatline
\nwendcode{}\nwbegindocs{2096}\nwdocspar
\nwenddocs{}\nwbegincode{2097}\moddef{Muli MC int: sample inclusive: TBP}\plusendmoddef
  procedure :: allocate => sample_inclusive_allocate  
\nwendcode{}\nwbegincode{2098}\moddef{Muli MC int: procedures}\plusendmoddef
  subroutine sample_inclusive_allocate (this, n_alloc)
    class(sample_inclusive_t), intent(out) :: this
    integer, intent(in) :: n_alloc
    allocate (this%int_kinds (n_alloc))
    this%n_alloc = n_alloc
  end subroutine sample_inclusive_allocate
  
\eatline
\nwendcode{}\nwbegindocs{2099}\nwdocspar
\nwenddocs{}\nwbegincode{2100}\moddef{Muli MC int: procedures}\plusendmoddef
  pure subroutine muli_mcint_generate_hit (rnd, corners, hit)
    real(default), dimension(2), intent(out) :: hit
    integer, intent(in), dimension(2) :: rnd
    real(default), dimension(2,2), intent(in) :: corners
    ! print *, hit
    ! print *, corners
    ! print *, (corners(1:2,2) - corners(1:2,1))
    hit = (rnd / max_d) * (corners(1:2,2) - corners(1:2,1)) + corners(1:2,1)
  end subroutine muli_mcint_generate_hit
  
\eatline
\nwendcode{}\nwbegindocs{2101}\nwdocspar
\nwenddocs{}\nwbegincode{2102}\moddef{Muli MC int: procedures}\plusendmoddef
  subroutine plot_pstvue3d (unit, corners, density)
    integer, intent(in) :: unit
    real(default), dimension(2,2), intent(in) :: corners
    real(default), intent(in) :: density
    real(default), dimension(2) :: width, mean
    real(default), dimension(3,3) :: plot
    width = (corners(:,2) - corners(:,1)) / two
    mean = (corners(:,1) + corners(:,2)) / two
    plot(1,1) = width(1)
    plot(2,1) = width(2)
    plot(3,1) = density / two
    plot(1,2) = mean(1)
    plot(2,2) = mean(2)
    plot(3,2) = density / two
    call log_color_code (density, plot(1:3,3))
    if (density > one) then                      
       write (unit, fmt='("\\mybigcube\{",F14.7,"\}\{",F14.7,"\}\{",F14.7,"\} &
            & \{",F14.7,"\}\{",F14.7,"\}\{",F14.7,"\}\{",F14.7,"\}\{",F14.7,"\} &
            & \{",F14.7,"\}")') plot
       return
    end if
    write (unit, fmt='("\\mycube\{",F14.7,"\}\{",F14.7,"\}\{",F14.7,"\} &
         & \{",F14.7,"\}\{",F14.7,"\}\{",F14.7,"\}\{",F14.7,"\}\{",F14.7,"\} &
         & \{",F14.7,"\}")') plot
  end subroutine plot_pstvue3d

\eatline
\nwendcode{}\nwbegindocs{2103}\nwdocspar
\nwenddocs{}\nwbegincode{2104}\moddef{Muli MC int: procedures}\plusendmoddef
  subroutine log_color_code (number, rgb)
    real(default), intent(in) :: number
    real(default), dimension(3), intent(out) :: rgb
    if (number < exp(-five)) then
       rgb = [zero, zero, exp(five)*number]
    else
       if (number < exp(-four)) then
          rgb = [zero, (number-exp(-five))/(exp(-four)-exp(-five)), one]
       else
          if (number < exp(-three)) then
             rgb = [zero, one, one-((number-exp(-four))/(exp(-three)-exp(-four)))]
          else
             if (number < exp(-two)) then
                rgb = [(number-exp(-three))/(exp(-two)-exp(-three)), one, zero]
             else
                if (number < exp(-one)) then                   
                   rgb = [one, one-(number-exp(-two))/(exp(-one)-exp(-two)), zero]
                else
                   if (number < one) then
                      rgb = [one, zero, (number-exp(-three))/(one-exp(-three))]
                   else                      
                      rgb = [exp(one), one, one] * exp(-number)
                      return
                   end if
                end if
             end if
          end if
       end if
    end if
  end subroutine log_color_code

\eatline
\nwendcode{}\nwbegindocs{2105}\nwdocspar
\nwenddocs{}\nwbegincode{2106}\moddef{Muli MC int: procedures}\plusendmoddef
  recursive subroutine misc_sort (in, out)
    real(default), dimension(:), intent(in) :: in
    integer, dimension(:), intent(out) :: out
    integer, dimension(:), allocatable :: tmp
    integer :: n, k, l, cut
    if (size(in) == 1) then
       out = [1]
    else
       if (size(in) == 2) then
          if (in(1) <= in(2)) then
             out = [1,2]
          else
             out = [2,1]
          end if
       else
          cut = size(in) / 2
          k = 1
          l = cut + 1
          allocate (tmp (size(in)))
          call misc_sort (in(1:cut), tmp(1:cut))
          call misc_sort (in(cut+1:), tmp(cut+1:))
          do n = cut + 1, size(in)
             tmp(n) = tmp(n) + cut
          end do
          do n = 1, size(in)
             if (k > cut) then
                out(n) = tmp(l)
                l = l + 1
             else
                if (l > size(tmp)) then
                   out(n) = tmp(k)
                   k = k + 1
                else
                   if (in(tmp(k)) < in(tmp(l))) then
                      out(n) = tmp(k)
                      k = k + 1
                   else
                      out(n) = tmp(l)
                      l = l + 1
                   end if
                end if
             end if
          end do
       end if
    end if
  end subroutine misc_sort
  
\eatline
\nwendcode{}\nwbegindocs{2107}\nwdocspar
%%%%%%%%%%%%%%%%%%%%%%%%%%%%%%%%%%%%%%%%%%%%%%%%%%%%%%%%%%%%%%%%%%%%%%%%
\section{Proton remnants}

This file contains the module {\tt{}muli{\char95}remnant}. All bookkeeping of the
proton remnants and twin quarks is done here. Furthermore, reweighting of the
PDFs to derive remnant PDFs is done here.

\nwenddocs{}\nwbegincode{2108}\moddef{\code{}muli{\char95}remnant.f90\edoc{}}\endmoddef
\LA{}File header\RA{}

module muli_remnant
  use, intrinsic :: iso_fortran_env
\LA{}Use kinds with double\RA{}
\LA{}Use strings\RA{}
  use string_utils
  use constants
  use diagnostics
  use pdf_builtin !NODEP!
  use tao_random_numbers !NODEP!
  use muli_base
  use muli_interactions
  use muli_momentum
!  use sf_lhapdf !NODEP!

\LA{}Standard module head\RA{}

\LA{}Muli remnant: variables\RA{}

\LA{}Muli remnant: public\RA{}

\LA{}Muli remnant: types\RA{}

\LA{}Muli remnant: interfaces\RA{}

contains

\LA{}Muli remnant: procedures\RA{}  
  
  subroutine pp_remnant_momentum_pdf (this,x_proton,gev2_scale,n,pdg_f,pdf)
    class(pp_remnant_t), intent(in) :: this
    real(default), intent(in) :: x_proton, gev2_scale
    integer, intent(in) :: n, pdg_f
    real(default), intent(out) :: pdf
    if (n==1 .or. n==2) then
       if (x_proton<=this%proton(n)%momentum_fraction) then
          if (pdg_f==PDG_FLAVOR_g) then
             call this%proton(n)%momentum_flavor_pdf (sqrt(GeV2_scale), &
                  x_proton / this%proton(n)%momentum_fraction, &
                  LHA_FLAVOR_g, pdf)
          else
             call this%proton(n)%momentum_flavor_pdf (sqrt(GeV2_scale), &
                  x_proton / this%proton(n)%momentum_fraction, &
                  pdg_f, pdf)
          end if
          pdf = pdf * this%proton(n)%momentum_fraction
       else
          pdf = zero
       end if
    else
       call msg_error ("pp_remnant_momentum_pdf: n must be either 1 or 2, " &
            // "but it is " // char (str (n)) // ".")
       stop
    end if
  end subroutine pp_remnant_momentum_pdf

  subroutine pp_remnant_apply_interaction(this,qcd_2_2)
    class(pp_remnant_t), intent(inout)::this
    class(qcd_2_2_class), intent(in)::qcd_2_2
    integer, dimension(4)::lha_f
    integer, dimension(2)::int_k
    real(default) :: gev_pt
    real(default), dimension(2) :: mom_f
    integer::n
    mom_f=qcd_2_2%get_remnant_momentum_fractions()
    lha_f=qcd_2_2%get_lha_flavors()
    int_k=qcd_2_2%get_pdf_int_kinds()
    gev_pt=qcd_2_2%get_gev_scale()
    !print *,"pp_remnant_apply_interaction",mom_f,qcd_2_2%get_parton_id(1),qcd_2_2%get_parton_id(2),lha_f
    do n=1,2
       select case (int_k(n))
       case (PDF_VALENCE_DOWN)
          call this%proton(n)%remove_valence_down_quark &
               (qcd_2_2%get_parton_id(n), gev_pt, mom_f(n))
       case (PDF_VALENCE_UP)
          call this%proton(n)%remove_valence_up_quark(qcd_2_2%get_parton_id(n),gev_pt,mom_f(n))
       case (PDF_SEA)
          call this%proton(n)%remove_sea_quark(qcd_2_2%get_parton_id(n),gev_pt,mom_f(n),lha_f(n))
       case (PDF_GLUON)
          call this%proton(n)%remove_gluon(qcd_2_2%get_parton_id(n),gev_pt,mom_f(n))
       end select
    end do
    this%X=this%proton(1)%momentum_fraction*this%proton(2)%momentum_fraction
  end subroutine pp_remnant_apply_interaction

  pure function pp_remnant_get_pdf_int_weights &
       (this, pdf_int_kinds) result (weight)
    class(pp_remnant_t), intent(in) :: this
    real(default) :: weight
    integer, dimension(2), intent(in) :: pdf_int_kinds ! pdf_int_kind
    weight = this%proton(1)%pdf_int_weight (pdf_int_kinds(1)) * &
         this%proton(2)%pdf_int_weight (pdf_int_kinds(2)) !*((this%x)**2)
  end function pp_remnant_get_pdf_int_weights

  elemental function pp_remnant_get_pdf_int_weight(this,kind1,kind2) result(weight)
    class(pp_remnant_t), intent(in)::this
    real(double)::weight
    integer, intent(in)::kind1,kind2 ! pdf_int_kind
    weight=this%proton(1)%pdf_int_weight(kind1)*this%proton(2)%pdf_int_weight(kind2)!*((this%x)**2)
  end function pp_remnant_get_pdf_int_weight

  subroutine pp_remnant_set_pdf_weight (this, weights)
    class(pp_remnant_t), intent(inout) :: this
    real(default), dimension(10), intent(in) :: weights
    this%proton(1)%pdf_int_weight = weights(1:5)
    this%proton(2)%pdf_int_weight = weights(6:10)
  end subroutine pp_remnant_set_pdf_weight

  elemental function pp_remnant_get_gev_initial_cme(this) result(cme)
    class(pp_remnant_t), intent(in)::this
    real(double)::cme
    cme=this%gev_initial_cme
  end function pp_remnant_get_gev_initial_cme

  elemental function pp_remnant_get_gev_actual_cme(this) result(cme)
    class(pp_remnant_t), intent(in)::this
    real(double)::cme
    cme=this%gev_initial_cme*this%X
  end function pp_remnant_get_gev_actual_cme

  elemental function pp_remnant_get_cme_fraction(this) result(cme)
    class(pp_remnant_t), intent(in)::this
    real(double)::cme
    cme=this%X
  end function pp_remnant_get_cme_fraction

  pure function pp_remnant_get_proton_remnant_momentum_fractions(this) result(fractions)
    class(pp_remnant_t), intent(in)::this
    real(double), dimension(2)::fractions
    fractions=[this%proton(1)%get_momentum_fraction(),this%proton(2)%get_momentum_fraction()]
  end function pp_remnant_get_proton_remnant_momentum_fractions

  subroutine pp_remnant_get_proton_remnants(this,proton1,proton2)
    class(pp_remnant_t),target, intent(in)::this
    class(proton_remnant_t), intent(out),pointer::proton1,proton2
    proton1=>this%proton(1)
    proton2=>this%proton(2)
  end subroutine pp_remnant_get_proton_remnants

  subroutine pp_remnant_get_remnant_parton_flavor_pdf_arrays &
       (this, GeV_scale, momentum1, momentum2, pdf1, pdf2)
    class(pp_remnant_t), intent(in) :: this
    real(default), intent(in) :: GeV_scale, momentum1, momentum2
    real(double), dimension(-6:6), intent(out)::pdf1,pdf2
    call this%proton(1)%parton_flavor_pdf_array (GeV_scale, momentum1, pdf1)
    call this%proton(2)%parton_flavor_pdf_array (GeV_scale, momentum2, pdf2)
  end subroutine pp_remnant_get_remnant_parton_flavor_pdf_arrays

  pure function remnant_dglap_splitting_gqq (z) result(p)
    real(default) :: p
    real(default), intent(in) :: z
    p = (z**2 + (1-z)**2) / two
  end function remnant_dglap_splitting_gqq
  
  pure function remnant_gluon_pdf_approx (x,p) result (g)
    real(default) :: g
    integer, intent(in) :: p
    real(default), intent(in) :: x
    g = ((1-x)**p) / x
  end function remnant_gluon_pdf_approx

  pure function remnant_norm_0 (xs) result (c0)
    real(default) :: c0
    real(default), intent(in) :: xs
    c0 = 6*xs / (2 - xs * (3 - 3*xs + 2*xs**2))
  end function remnant_norm_0

  pure function remnant_norm_1 (xs) result (c1)
    real(default) :: c1
    real(default), intent(in)::xs
    c1 = 3*xs / (2 - xs**2 * (3-xs) + 3*xs*log(xs))
  end function remnant_norm_1

  pure function remnant_norm_4 (xs) result (c4)
    real(default) :: c4
    real(default), intent(in) :: xs
    real(default) :: y
    if ((one-xs) > 1E-3_default) then
       c4 = 3*xs / (1 + 11*xs + 6*xs*log(xs) + 12*xs**3*log(xs) + &
            18*xs**2*log(xs) + 9*xs**2 - 19*xs**3 - 2*xs**4)
    else
       y=one / (one-xs)
       c4= 1130._default / 11907._default &
            - 10._default * y**5 &
            - 40._default * y**4 / three &
            -160._default * y**3 / 63._default &
            + 50._default * y**2 / 189._default &
            -565._default * y / 3969._default &
            -186170._default * (one-xs) / 2750517._default
    end if
  end function remnant_norm_4

  pure function remnant_norm (xs, p) result (c)
    real(default) :: c
    real(default), intent(in) :: xs
    integer, intent(in) :: p
    select case (p)
    case (0)
       c = remnant_norm_0 (xs)
    case (1)
       c = remnant_norm_1 (xs)
    case default
       c = remnant_norm_4 (xs)
    end select
  end function remnant_norm

  pure function remnant_twin_pdf_p (x, xs, p) result (qc)
    real(default) :: qc
    real(default), intent(in) :: x, xs
    integer, intent(in) :: p
    qc = remnant_norm (xs, p) * remnant_gluon_pdf_approx (xs + x, p) &
         * remnant_dglap_splitting_gqq (xs / (xs+x)) / (xs+x)
  end function remnant_twin_pdf_p
  
  elemental function remnant_twin_momentum_4 (xs) result (p)
    real(default) :: p
    real(default), intent(in) :: xs
    if (xs < 0.99_default) then
       p = (-9 * (-1+xs) * xs * (1+xs) * (5+xs*(24+xs)) + &
            12*xs*(1+2*xs)*(1+2*xs*(5+2*xs))*Log(xs))/&
            (8*(1+2*xs)*((-1+xs)*(1+xs*(10+xs))-6*xs*(1+xs)*Log(xs)))
    else
       p = (1-xs) / 6 - (5*(-1+xs)**2)/63+(5*(-1+xs)**3) / 216
    end if
  end function remnant_twin_momentum_4

  subroutine gnuplot_integrated_pdf (this, momentum_unit, parton_unit)
    class(proton_remnant_t), intent(in)::this
    integer, intent(in)::momentum_unit,parton_unit
!    real(double), intent(in)::gev_scale
    integer, parameter :: x_grid = 1000000
    integer, parameter :: q_grid = 100
    integer::n,m,mem
    real(default) :: x, q, dx, dq, overall_sum, xmin, xmax, &
         q2min, q2max, qmin, qmax
    real(double) :: q2min_dbl, q2max_dbl, xmin_dbl, xmax_dbl
    real(double), dimension(-6:6) :: sea_pdf, sea_momentum_pdf_sum, &
         sea_parton_pdf_sum
    real(default), dimension(2) :: valence_pdf, valence_momentum_pdf_sum,&
         valence_parton_pdf_sum
    real(default), allocatable, dimension(:) :: twin_momentum_pdf_sum
    class(muli_parton_t), pointer :: tmp_twin
    mem=1
    call GetXmin (mem, xmin_dbl)
    call GetXmax (mem, xmax_dbl)
    call GetQ2max (mem, q2max_dbl)
    call GetQ2min (mem, q2min_dbl)
    xmin = xmin_dbl
    xmax = xmax_dbl
    q2min = q2min_dbl
    q2max = q2max_dbl
    qmin=sqrt(q2min)
    qmax=sqrt(q2max)
    print *,"qmin=",qmin,"GeV"
    print *,"qmax=",qmax,"GeV"
    dx=(xmax-xmin)/x_grid
    dq=(qmax-qmin)/q_grid
    q=qmin+dq/2D0
    tmp_twin=>this%twin_partons%next
    n=0
    if (this%n_twins>0) then
       allocate(twin_momentum_pdf_sum(this%n_twins))
       do while(associated(tmp_twin))
          n=n+1
          twin_momentum_pdf_sum(n)=tmp_twin%momentum
          tmp_twin=>tmp_twin%next
       end do
    end if
    do m=1,q_grid
       valence_momentum_pdf_sum=[0D0,0D0]
       valence_parton_pdf_sum=[0D0,0D0]
       sea_momentum_pdf_sum=[0D0,0D0,0D0,0D0,0D0,0D0,0D0,0D0,0D0,0D0,0D0,0D0,0D0]    
       sea_parton_pdf_sum=[0D0,0D0,0D0,0D0,0D0,0D0,0D0,0D0,0D0,0D0,0D0,0D0,0D0]    
       x=xmin+dx/2D0
       do n=1,x_grid
          call this%parton_kind_pdf_array (Q, x, valence_pdf, sea_pdf)
          valence_parton_pdf_sum=valence_parton_pdf_sum+valence_pdf
          sea_parton_pdf_sum=sea_parton_pdf_sum+sea_pdf
          call this%momentum_kind_pdf_array (Q, x, valence_pdf, sea_pdf)
          valence_momentum_pdf_sum=valence_momentum_pdf_sum+valence_pdf
          sea_momentum_pdf_sum=sea_momentum_pdf_sum+sea_pdf
          x=x+dx
       end do
       valence_parton_pdf_sum=valence_parton_pdf_sum*dx
       sea_parton_pdf_sum=sea_parton_pdf_sum*dx
       valence_momentum_pdf_sum=valence_momentum_pdf_sum*dx
       sea_momentum_pdf_sum=sea_momentum_pdf_sum*dx
       if (this%n_twins>0) then
          write(momentum_unit,fmt=*)q,&
               sum(valence_momentum_pdf_sum)+sum(sea_momentum_pdf_sum)+sum(twin_momentum_pdf_sum),&
               valence_momentum_pdf_sum,&
               sea_momentum_pdf_sum,&
               twin_momentum_pdf_sum
       else
          write(momentum_unit,fmt=*)q,&
               sum(valence_momentum_pdf_sum)+sum(sea_momentum_pdf_sum),&
               valence_momentum_pdf_sum,&
               sea_momentum_pdf_sum
       end if
       write(parton_unit,fmt=*)q,&
            sum(valence_parton_pdf_sum)+sum(sea_parton_pdf_sum),&
            valence_parton_pdf_sum,&
            sea_parton_pdf_sum
       q=q+dq
    end do
  end subroutine gnuplot_integrated_pdf

end module muli_remnant

\eatline
\nwendcode{}\nwbegindocs{2109}\nwdocspar
\nwenddocs{}\nwbegincode{2110}\moddef{Muli remnant: public}\endmoddef
  public :: pdfnorm_t
\nwendcode{}\nwbegincode{2111}\moddef{Muli remnant: types}\endmoddef
  type, extends (ser_class_t) :: pdfnorm_t
     real(default) :: qmin, qmax, dq
     real(default), dimension(-6:6, 0:nq) :: pdf_int
     real(default), dimension(0:4, 0:nq) :: pdf_norm
   contains
   \LA{}Muli remnant: pdfnorm: TBP\RA{}
  end type pdfnorm_t
  
\eatline
\nwendcode{}\nwbegindocs{2112}\nwdocspar
\nwenddocs{}\nwbegincode{2113}\moddef{Muli remnant: pdfnorm: TBP}\endmoddef
  procedure :: write_to_marker => pdfnorm_write_to_marker
\nwendcode{}\nwbegincode{2114}\moddef{Muli remnant: procedures}\endmoddef
  subroutine pdfnorm_write_to_marker (this, marker, status)
    class(pdfnorm_t), intent(in) :: this
    class(marker_t), intent(inout) :: marker
    integer(dik), intent(out) :: status
    call marker%mark_begin ("pdfnorm_t")
    call marker%mark ("qmin", this%qmin)
    call marker%mark ("qmax", this%qmax)
    call marker%mark ("dq", this%dq)
    call marker%mark ("pdf_int", this%pdf_int)
    call marker%mark ("pdf_norm", this%pdf_norm)
    call marker%mark_end ("pdfnorm_t")
  end subroutine pdfnorm_write_to_marker
  
\eatline
\nwendcode{}\nwbegindocs{2115}\nwdocspar
\nwenddocs{}\nwbegincode{2116}\moddef{Muli remnant: pdfnorm: TBP}\plusendmoddef
  procedure :: read_from_marker => pdfnorm_read_from_marker
\nwendcode{}\nwbegincode{2117}\moddef{Muli remnant: procedures}\plusendmoddef
  subroutine pdfnorm_read_from_marker (this, marker, status)
    class(pdfnorm_t), intent(out) :: this
    class(marker_t), intent(inout) :: marker
    integer(dik), intent(out) :: status
    character(:), allocatable :: name
    call marker%pick_begin ("pdfnorm_t", status=status)
    call marker%pick ("qmin", this%qmin, status)
    call marker%pick ("qmax", this%qmax, status)
    call marker%pick ("dq", this%dq, status)
    call marker%pick ("pdf_int", this%pdf_int, status)
    call marker%pick ("pdf_norm", this%pdf_norm, status)
    call marker%pick_end ("pdfnorm_t", status=status)
  end subroutine pdfnorm_read_from_marker
  
\eatline
\nwendcode{}\nwbegindocs{2118}\nwdocspar
\nwenddocs{}\nwbegincode{2119}\moddef{Muli remnant: pdfnorm: TBP}\plusendmoddef
  procedure :: print_to_unit => pdfnorm_print_to_unit
\nwendcode{}\nwbegincode{2120}\moddef{Muli remnant: procedures}\plusendmoddef
  recursive subroutine pdfnorm_print_to_unit &
       (this, unit, parents, components, peers)
    class(pdfnorm_t), intent(in) :: this
    integer, intent(in) :: unit
    integer(dik), intent(in) :: parents, components, peers
    write (unit, "(1x,A)")       "Components of pdfnorm_t:"
    write (unit, "(3x,A,F7.6)")  "qmin:    ", this%qmin
    write (unit, "(3x,A,F7.6)")  "qmax:    ", this%qmax
    write (unit, "(3x,A,F7.6)")  "dq:      ", this%dq
    if (components > 0) then
       write(unit, "(3x,A,13(F8.6,1x))") "pdf_int:  ", this%pdf_int
       write(unit, "(3x,A,5(F8.6,1x))")  "pdf_norm: ", this%pdf_norm
    else
       write(unit, "(3x,A)")  "Skipping pdf_int"
       write(unit, "(3x,A)")  "Skipping pdf_norm"
    end if
  end subroutine pdfnorm_print_to_unit
  
\eatline
\nwendcode{}\nwbegindocs{2121}\nwdocspar
\nwenddocs{}\nwbegincode{2122}\moddef{Muli remnant: procedures}\plusendmoddef
\nwendcode{}\nwbegindocs{2123}\nwdocspar
\nwenddocs{}\nwbegincode{2124}\moddef{Muli remnant: pdfnorm: TBP}\plusendmoddef
  procedure, nopass :: get_type => pdfnorm_get_type
\nwendcode{}\nwbegincode{2125}\moddef{Muli remnant: procedures}\plusendmoddef
  pure subroutine pdfnorm_get_type (type)
    character(:), allocatable, intent(out) :: type
    allocate (type, source="pdfnorm_t")
  end subroutine pdfnorm_get_type
  
\eatline
\nwendcode{}\nwbegindocs{2126}\nwdocspar
\nwenddocs{}\nwbegincode{2127}\moddef{Muli remnant: pdfnorm: TBP}\plusendmoddef
  procedure, nopass :: verify_type => pdfnorm_verify_type  
\nwendcode{}\nwbegincode{2128}\moddef{Muli remnant: procedures}\plusendmoddef
  elemental logical function pdfnorm_verify_type (type) result (match)
    character(*), intent(in) :: type
    match = type == "pdfnorm_t"
  end function pdfnorm_verify_type
  
\eatline
\nwendcode{}\nwbegindocs{2129}\nwdocspar
\nwenddocs{}\nwbegincode{2130}\moddef{Muli remnant: pdfnorm: TBP}\plusendmoddef
  procedure :: scan => pdfnorm_scan
\nwendcode{}\nwbegincode{2131}\moddef{Muli remnant: procedures}\plusendmoddef
  subroutine pdfnorm_scan (this)
    class(pdfnorm_t), intent(out) :: this
    integer :: ix, iq
    real(double) :: xmin, xmax, dx
    real(double) :: q, q2min, q2max
    real(double), dimension(-6:6) :: f
    real(double), dimension(0:2) :: x
    call getxmin (0, xmin)
    call getxmax (0, xmax)
    call getq2min (0, q2min)
    call getq2max (0, q2max)
    this%qmin = sqrt(sqrt(q2min))
    this%qmax = sqrt(sqrt(q2max))
    this%dq = (this%qmax - this%qmin) / nq
    xmin = sqrt(xmin)
    xmax = sqrt(xmax)   
    dx= (xmax - xmin) / nx
    do iq = 0, nq
       print *, "iq=", iq, "/", nq
       q = (this%qmin + iq * this%dq)**2
       x(0) = xmin**2
       x(1) = (xmin+dx)**2
       call evolvePDF (x(0), q, f)
       f(1) = f(1) - f(-1)
       f(2) = f(2) - f(-2)
       this%pdf_int(:,iq) = (x(1) - x(0)) * f
       do ix = 2, nx
          x(2) = (xmin + ix*dx)**2
          call evolvePDF (x(1), q, f)
          f(1) = f(1) - f(-1)
          f(2) = f(2) - f(-2)
          this%pdf_int(:,iq) = this%pdf_int(:,iq) + f*(x(2) - x(0))
          x(0) = x(1)
          x(1) = x(2)
       end do
       call evolvePDF (x(1), q, f)
       f(1) = f(1) - f(-1)
       f(2) = f(2) - f(-2)
       this%pdf_int(:,iq) = (this%pdf_int(:,iq) + f*(x(1)-x(0))) / two
       this%pdf_norm(4,iq) = this%pdf_int(2,iq)
       this%pdf_norm(3,iq) = this%pdf_int(1,iq)
       this%pdf_int(2,iq) = this%pdf_int(2,iq) + this%pdf_int(-2,iq)
       this%pdf_int(1,iq) = this%pdf_int(1,iq) + this%pdf_int(-1,iq)
       this%pdf_norm(1,iq) = this%pdf_int(0,iq)
       this%pdf_norm(2,iq) = sum (this%pdf_int(-6:-1,iq)) + &
            sum(this%pdf_int(-2:-1,iq)) + sum(this%pdf_int(3:6,iq))
       this%pdf_norm(0,iq) = sum(this%pdf_int(:,iq))
       this%pdf_norm(1,iq) = this%pdf_norm(1,iq) / this%pdf_norm(0,iq)
       this%pdf_norm(2,iq) = this%pdf_norm(2,iq) / this%pdf_norm(0,iq)
       this%pdf_norm(3,iq) = this%pdf_norm(3,iq) / this%pdf_norm(0,iq)
       this%pdf_norm(4,iq) = this%pdf_norm(4,iq) / this%pdf_norm(0,iq)
       ! print *, this%pdf_norm(0,iq) - one
    end do
  end subroutine pdfnorm_scan
  
\eatline
\nwendcode{}\nwbegindocs{2132}\nwdocspar
\nwenddocs{}\nwbegincode{2133}\moddef{Muli remnant: pdfnorm: TBP}\plusendmoddef
  procedure :: get_norm => pdfnorm_get_norm    
\nwendcode{}\nwbegincode{2134}\moddef{Muli remnant: procedures}\plusendmoddef
  subroutine pdfnorm_get_norm (this, gev_q, dim, kind, norm)
    class(pdfnorm_t), intent(in) :: this
    real(default), intent(in) :: gev_q
    integer, intent(in) :: dim, kind
    real(default), intent(out)::norm
    integer :: iq
    real(default) :: x,q , z0, z1, z2, z3, z4
    norm = -one
    q = sqrt(gev_q) - this%qmin
    iq = floor(q / this%dq)
    x = q / this%dq - iq
    if (iq < 0) then
       print *, "pdfnorm_getnorm: q < q_min ", gev_q, this%qmin**2
       norm = this%pdf_norm (kind, 0)
    else
       if (iq >= nq) then
          print *, "pdfnorm_getnorm: q >= q_max ", gev_q, this%qmax**2
          norm = this%pdf_norm (kind, nq)
       else
          select case (dim)
          case (0)
             norm = this%pdf_norm (kind, iq)
          case (1)
             norm = this%pdf_norm(kind,iq) * (one - x) + &
                  this%pdf_norm(kind,iq+1) * x
          case (2)
             x = x + mod(iq,2)
             iq = iq - mod(iq,2)
             z0 = this%pdf_norm(kind, iq)
             z1 = this%pdf_norm(kind, iq+1)
             z2 = this%pdf_norm(kind, iq+2)
             norm = ((z0 - 2D0*z1 + z2) * x - (three*z0 - four*z1 + z2)) * &
                  x / two + z0
          case (3)
             x = x + mod(iq,3)
             iq = iq - mod(iq,3)
             z0 = this%pdf_norm(kind, iq)
             z1 = this%pdf_norm(kind, iq+1)
             z2 = this%pdf_norm(kind, iq+2)
             z3 = this%pdf_norm(kind, iq+3)
             norm = (( - (z0 - 3*z1 + 3*z2 -z3) * x + 3 * (2*z0 - &
                  5*z1 + 4*z2 - z3))*x - (11*z0 - 18*z1 + 9*z2 - 2*z3)) * &
                  x / 6._default + z0
          case (4)
             x = x + mod(iq,4)
             iq = iq - mod(iq,4)
             z0 = this%pdf_norm(kind, iq)
             z1 = this%pdf_norm(kind, iq+1)
             z2 = this%pdf_norm(kind, iq+2)
             z3 = this%pdf_norm(kind, iq+3)
             z4 = this%pdf_norm(kind, iq+4)
             norm = (((((z0 - 4*z1 + 6*z2 - 4*z3 + z4) * x &
                  -2 * (5*z0 - 18*z1 + 24*z2 - 14*z3 + 3*z4)) * x &
                  + (35*z0 - 104*z1 + 114*z2 - 56*z3 + 11*z4)) * x &
                  -2 * (25*z0 - 48*z1 + 36*z2 - 16*z3 + 3*z4)) * x) / &
                  24._default + z0
          case default
             norm = this%pdf_norm(kind, iq) * (one - x) + &
                  this%pdf_norm(kind, iq+1) * x
          end select
          ! print *, iq, x, norm
       end if
    end if
  end subroutine pdfnorm_get_norm
  
\eatline
\nwendcode{}\nwbegindocs{2135}\nwdocspar
\nwenddocs{}\nwbegincode{2136}\moddef{Muli remnant: variables}\endmoddef
  integer, parameter :: nx = 10000000
\nwendcode{}\nwbegincode{2137}\moddef{Muli remnant: variables}\plusendmoddef
  integer, parameter :: nq = 60
\nwendcode{}\nwbegincode{2138}\moddef{Muli remnant: public}\plusendmoddef
  public :: remnant_weight_model  
\nwendcode{}\nwbegincode{2139}\moddef{Muli remnant: variables}\plusendmoddef
  integer :: remnant_weight_model = 2
\nwendcode{}\nwbegincode{2140}\moddef{Muli remnant: variables}\plusendmoddef
  integer :: gluon_exp = 4
  
\eatline
\nwendcode{}\nwbegindocs{2141}\nwdocspar
\nwenddocs{}\nwbegincode{2142}\moddef{Muli remnant: public}\plusendmoddef
  public :: muli_parton_t
\nwendcode{}\nwbegincode{2143}\moddef{Muli remnant: types}\plusendmoddef
  type, extends (ser_class_t) :: muli_parton_t
     private
     integer :: id = -1
     integer :: lha_flavor
     real(default) :: momentum = -1
     class(muli_parton_t), pointer :: twin => null()
     class(muli_parton_t), pointer :: next => null()
   contains
   \LA{}Muli remnant: muli parton: TBP\RA{}
  end type muli_parton_t

\eatline
\nwendcode{}\nwbegindocs{2144}\nwdocspar
\nwenddocs{}\nwbegincode{2145}\moddef{Muli remnant: muli parton: TBP}\endmoddef
  procedure :: write_to_marker => parton_write_to_marker
\nwendcode{}\nwbegincode{2146}\moddef{Muli remnant: procedures}\plusendmoddef
  subroutine parton_write_to_marker (this, marker, status)
    class(muli_parton_t), intent(in) :: this
    class(marker_t), intent(inout) :: marker
    integer(dik), intent(out) :: status
    call marker%mark_begin ("muli_parton_t")
    call marker%mark ("id", this%id)
    call marker%mark ("lha", this%lha_flavor)
    call marker%mark ("momentum", this%momentum)
    call marker%mark_end ("muli_parton_t")
  end subroutine parton_write_to_marker

\eatline
\nwendcode{}\nwbegindocs{2147}\nwdocspar
\nwenddocs{}\nwbegincode{2148}\moddef{Muli remnant: muli parton: TBP}\plusendmoddef
  procedure :: read_from_marker => parton_read_from_marker
\nwendcode{}\nwbegincode{2149}\moddef{Muli remnant: procedures}\plusendmoddef
  subroutine parton_read_from_marker (this, marker, status)
    class(muli_parton_t), intent(out) :: this
    class(marker_t), intent(inout) :: marker
    integer(dik), intent(out) :: status
    character(:), allocatable :: name
    call marker%pick_begin ("muli_parton_t", status=status)
    call marker%pick ("id", this%id, status)
    call marker%pick ("lha", this%lha_flavor, status)
    call marker%pick ("momentum", this%momentum, status)
    call marker%pick_end ("muli_parton_t", status=status)
  end subroutine parton_read_from_marker
  
\eatline
\nwendcode{}\nwbegindocs{2150}\nwdocspar
\nwenddocs{}\nwbegincode{2151}\moddef{Muli remnant: muli parton: TBP}\plusendmoddef
  procedure :: print_to_unit => parton_print_to_unit
\nwendcode{}\nwbegincode{2152}\moddef{Muli remnant: procedures}\plusendmoddef
  recursive subroutine parton_print_to_unit &
       (this, unit, parents, components, peers)
    class(muli_parton_t), intent(in) :: this
    integer, intent(in) :: unit
    integer(dik), intent(in) :: parents, components, peers
    class(ser_class_t), pointer :: ser
    write (unit, "(1x,A)")      "Components of muli_parton_t:"
    write (unit, "(3x,A,I7)")   "id:         ", this%id
    write (unit, "(3x,A,I7)")   "lha flavor: ", this%lha_flavor
    write (unit, "(3x,A,F7.6)") "momentum:   ", this%momentum
    ser => this%next
    call serialize_print_peer_pointer &
         (ser, unit, parents, components, peers-i_one, "next")
    ser => this%twin
    call serialize_print_comp_pointer &
         (ser, unit, parents, components, peers-i_one, "twin")
  end subroutine parton_print_to_unit

\eatline
\nwendcode{}\nwbegindocs{2153}\nwdocspar
\nwenddocs{}\nwbegincode{2154}\moddef{Muli remnant: muli parton: TBP}\plusendmoddef
  procedure, nopass :: get_type => parton_get_type  
\nwendcode{}\nwbegincode{2155}\moddef{Muli remnant: procedures}\plusendmoddef
  pure subroutine parton_get_type (type)
    character(:), allocatable, intent(out) :: type
    allocate (type, source="muli_parton_t")
  end subroutine parton_get_type
  
\eatline
\nwendcode{}\nwbegindocs{2156}\nwdocspar
\nwenddocs{}\nwbegincode{2157}\moddef{Muli remnant: muli parton: TBP}\plusendmoddef
  procedure :: unweighted_pdf => twin_unweighted_pdf
\nwendcode{}\nwbegincode{2158}\moddef{Muli remnant: procedures}\plusendmoddef
  pure function twin_unweighted_pdf (this, momentum_fraction) result (pdf)
    !parton pdf
    class(muli_parton_t), intent(in) :: this
    real(default), intent(in) :: momentum_fraction
    real(default) :: pdf
    if (momentum_fraction + this%twin%momentum < one) then
       pdf = remnant_twin_pdf_p (momentum_fraction, &
            this%twin%momentum, gluon_exp)
    else
       pdf = zero
    end if
  end function twin_unweighted_pdf
  
\eatline
\nwendcode{}\nwbegindocs{2159}\nwdocspar
\nwenddocs{}\nwbegincode{2160}\moddef{Muli remnant: muli parton: TBP}\plusendmoddef
  procedure :: deallocate => twin_deallocate
\nwendcode{}\nwbegincode{2161}\moddef{Muli remnant: procedures}\plusendmoddef
  recursive subroutine twin_deallocate (this)
    class(muli_parton_t) :: this
    if (associated (this%next)) then
       call this%next%deallocate
       deallocate (this%next)
    end if
  end subroutine twin_deallocate
  
\eatline
\nwendcode{}\nwbegindocs{2162}\nwdocspar
\nwenddocs{}\nwbegincode{2163}\moddef{Muli remnant: muli parton: TBP}\plusendmoddef
  procedure :: push => parton_push
\nwendcode{}\nwbegincode{2164}\moddef{Muli remnant: procedures}\plusendmoddef
  subroutine parton_push (this, parton)
    class(muli_parton_t), intent(inout) :: this
    class(muli_parton_t), intent(inout), pointer :: parton
    ! print *, "parton_push ", parton%id
    parton%next => this%next
    this%next => parton
  end subroutine parton_push
  
\eatline
\nwendcode{}\nwbegindocs{2165}\nwdocspar
\nwenddocs{}\nwbegincode{2166}\moddef{Muli remnant: muli parton: TBP}\plusendmoddef
  generic :: pop => pop_by_id, pop_by_association  
  procedure :: pop_by_id => parton_pop_by_id
  procedure :: pop_by_association => parton_pop_by_association  
\nwendcode{}\nwbegincode{2167}\moddef{Muli remnant: procedures}\plusendmoddef
  subroutine parton_pop_by_id (this, id, parton)
    class(muli_parton_t), target, intent(inout) :: this
    integer, intent(in) :: id
    class(muli_parton_t), intent(out), pointer :: parton
    class(muli_parton_t), pointer :: tmp_parton    
    tmp_parton => this
    do while (associated (tmp_parton%next))
       if (tmp_parton%next%id == id)  exit
       tmp_parton => tmp_parton%next
    end do
    if (associated (tmp_parton%next)) then
       parton => tmp_parton%next
       tmp_parton%next => parton%next
       nullify (parton%next)
       ! print *,"parton_pop ",id,parton%id
    else
       nullify (parton)
       print *,"parton_pop ", id, "NULL"
    end if
  end subroutine parton_pop_by_id
  
\eatline
\nwendcode{}\nwbegindocs{2168}\nwdocspar
\nwenddocs{}\nwbegincode{2169}\moddef{Muli remnant: procedures}\plusendmoddef
  subroutine parton_pop_by_association (this, parton)
    class(muli_parton_t), target, intent(inout) :: this
    class(muli_parton_t), intent(inout), target :: parton
    class(muli_parton_t), pointer :: tmp_parton    
    tmp_parton => this
    do while (associated (tmp_parton%next))
       if (associated (tmp_parton%next, parton))  exit
       tmp_parton=>tmp_parton%next
    end do
    if (associated(tmp_parton%next)) then
       tmp_parton%next => parton%next
       nullify (parton%next)
       ! print *,"parton_pop ", parton%id
    else
       print *, "parton_pop NULL"
    end if
  end subroutine parton_pop_by_association
  
\eatline
\nwendcode{}\nwbegindocs{2170}\nwdocspar
\nwenddocs{}\nwbegincode{2171}\moddef{Muli remnant: public}\plusendmoddef
  public :: proton_remnant_t
\nwendcode{}\nwbegincode{2172}\moddef{Muli remnant: types}\plusendmoddef
  type, extends (ser_class_t) :: proton_remnant_t
     private
     integer, dimension(2) :: valence_content = [1,2]
     integer :: n_twins = 0
     !!! [gluon, sea quark, valence down, valence up, twin]
     real(default), dimension(5) :: pdf_int_weight = [one, one, one, one, one]
     real(default) :: momentum_fraction = one
     real(default) :: twin_norm = one
     type(muli_parton_t) :: twin_partons
     type(muli_parton_t) :: is_partons
     type(muli_parton_t) :: fs_partons
     !!! These pointers shall not be allocated, deallocated, 
     !!!    serialized or deserialized explicitly.
     class(pdfnorm_t), pointer :: pdf_norm => null()
   contains
   \LA{}Muli remnant: proton remnant: TBP\RA{}
  end type proton_remnant_t

\nwendcode{}\nwbegindocs{2173} %def proton_remnant_t
\nwenddocs{}\nwbegindocs{2174}\nwdocspar
\nwenddocs{}\nwbegincode{2175}\moddef{Muli remnant: proton remnant: TBP}\endmoddef
  procedure :: remove_valence_quark => proton_remnant_remove_valence_quark  
\nwendcode{}\nwbegincode{2176}\moddef{Muli remnant: procedures}\plusendmoddef
  subroutine proton_remnant_remove_valence_quark &
       (this, id, GeV_scale, momentum_fraction, lha_flavor)
    class(proton_remnant_t), intent(inout) :: this
    integer, intent(in) :: id
    real(default), intent(in) :: GeV_scale, momentum_fraction
    integer, intent(in) :: lha_flavor  !!! d=1, u=2
    if (lha_flavor == 1 .or. lha_flavor == 2) then
       associate (q => this%valence_content (lha_flavor))
         if (q > 0) then
            q = q - 1
            call this%push_is_parton (id, lha_flavor, momentum_fraction)
            this%momentum_fraction = this%momentum_fraction * (one - momentum_fraction)
            call this%calculate_weight (GeV_scale)
         else
            write (*, "(1x,A,I2,A)")  "proton_remnant_remove_valence_quark: " // &
                 "Cannot remove parton ", lha_flavor, ": There are no such " // &
                 "partons left."
            call this%print_all ()
         end if
       end associate
    else
       write (*, "(1x,A,I2,A)")  "proton_remnant_remove_valence_quark: Cannot " // &
            "remove parton ", lha_flavor, ": There are no such valence partons."
    end if
  end subroutine proton_remnant_remove_valence_quark
    
\eatline
\nwendcode{}\nwbegindocs{2177}\nwdocspar
\nwenddocs{}\nwbegincode{2178}\moddef{Muli remnant: proton remnant: TBP}\plusendmoddef
  procedure :: remove_sea_quark => proton_remnant_remove_sea_quark
\nwendcode{}\nwbegincode{2179}\moddef{Muli remnant: procedures}\plusendmoddef
  subroutine proton_remnant_remove_sea_quark &
       (this, id, GeV_scale, momentum_fraction, lha_flavor)
    class(proton_remnant_t), intent(inout) :: this
    integer, intent(in) :: id
    real(default), intent(in) :: GeV_scale, momentum_fraction
    integer, intent(in) :: lha_flavor
    ! print *, "proton_remnant_remove_sea_quark", momentum_fraction
    if (lha_flavor > -6 .and. lha_flavor < 6 .and. lha_flavor .ne. 0) then
       this%momentum_fraction = this%momentum_fraction * &
            (one - momentum_fraction)
       call this%push_twin (id, lha_flavor, momentum_fraction, GeV_scale)
    end if
  end subroutine proton_remnant_remove_sea_quark
  
\eatline
\nwendcode{}\nwbegindocs{2180}\nwdocspar
\nwenddocs{}\nwbegincode{2181}\moddef{Muli remnant: proton remnant: TBP}\plusendmoddef
  procedure :: remove_gluon => proton_remnant_remove_gluon
\nwendcode{}\nwbegincode{2182}\moddef{Muli remnant: procedures}\plusendmoddef
  subroutine proton_remnant_remove_gluon &
       (this, id, GeV_scale, momentum_fraction)
    class(proton_remnant_t), intent(inout) :: this
    integer, intent(in) :: id
    real(default), intent(in) :: GeV_scale, momentum_fraction
    this%momentum_fraction = this%momentum_fraction * (one - momentum_fraction)
    call this%push_is_parton (id, LHA_FLAVOR_g, momentum_fraction)
  end subroutine proton_remnant_remove_gluon
  
\eatline
\nwendcode{}\nwbegindocs{2183}\nwdocspar
\nwenddocs{}\nwbegincode{2184}\moddef{Muli remnant: proton remnant: TBP}\plusendmoddef
  procedure :: remove_valence_up_quark => proton_remnant_remove_valence_up_quark
\nwendcode{}\nwbegincode{2185}\moddef{Muli remnant: procedures}\plusendmoddef
  subroutine proton_remnant_remove_valence_up_quark &
       (this, id, GeV_scale, momentum_fraction)
    class(proton_remnant_t), intent(inout) :: this
    integer, intent(in) :: id
    real(default), intent(in) :: GeV_scale, momentum_fraction
    associate (q => this%valence_content (LHA_FLAVOR_u))
      if (q > 0) then
         q = q - 1
         call this%push_is_parton (id, LHA_FLAVOR_u, momentum_fraction)
         this%momentum_fraction = this%momentum_fraction * (one - momentum_fraction)
         call this%calculate_weight (GeV_scale)
      else
         write (*, "(1x,A,I2,A)")  "proton_remnant_remove_valence_up_quark: " // &
              "Cannot remove parton ", LHA_FLAVOR_u, ": There are no such " // &
              "partons left."
         call this%print_all
      end if
    end associate
  end subroutine proton_remnant_remove_valence_up_quark
  
\eatline
\nwendcode{}\nwbegindocs{2186}\nwdocspar
\nwenddocs{}\nwbegincode{2187}\moddef{Muli remnant: proton remnant: TBP}\plusendmoddef
  procedure :: remove_valence_down_quark => &
       proton_remnant_remove_valence_down_quark
\nwendcode{}\nwbegincode{2188}\moddef{Muli remnant: procedures}\plusendmoddef
  subroutine proton_remnant_remove_valence_down_quark &
       (this, id, GeV_scale, momentum_fraction)
    class(proton_remnant_t), intent(inout) :: this
    integer, intent(in) :: id
    real(default), intent(in) :: GeV_scale, momentum_fraction
    associate (q => this%valence_content(LHA_FLAVOR_d))
      if (q > 0) then
         q = q - 1
         call this%push_is_parton (id, LHA_FLAVOR_d, momentum_fraction)
         this%momentum_fraction = this%momentum_fraction * &
              (one - momentum_fraction)
         call this%calculate_weight (GeV_scale)
      else
         write (*, "(1x,A,I2,A)")  "proton_remnant_remove_valence_down_quark:" // &
              "Cannot remove parton ", LHA_FLAVOR_d, ": There are no " // &
              "such partons left."
         call this%print_all
      end if
    end associate
  end subroutine proton_remnant_remove_valence_down_quark  

\eatline
\nwendcode{}\nwbegindocs{2189}\nwdocspar
\nwenddocs{}\nwbegincode{2190}\moddef{Muli remnant: proton remnant: TBP}\plusendmoddef
  procedure :: remove_twin => proton_remnant_remove_twin  
\nwendcode{}\nwbegincode{2191}\moddef{Muli remnant: procedures}\plusendmoddef
  subroutine proton_remnant_remove_twin (this, id, GeV_scale)
    class(proton_remnant_t), intent(inout) :: this
    integer, intent(in) :: id
    real(default), intent(in) :: GeV_scale
    class(muli_parton_t), pointer :: twin
    call this%twin_partons%pop (id, twin)
    call this%fs_partons%push (twin)
    this%twin_norm = this%twin_norm - twin%momentum
    this%n_twins = this%n_twins - 1
    call this%calculate_weight (GeV_scale)
  end subroutine proton_remnant_remove_twin
  
\eatline
\nwendcode{}\nwbegindocs{2192}\nwdocspar
\nwenddocs{}\nwbegincode{2193}\moddef{Muli remnant: proton remnant: TBP}\plusendmoddef
  procedure :: momentum_twin_pdf => proton_remnant_momentum_twin_pdf
\nwendcode{}\nwbegincode{2194}\moddef{Muli remnant: procedures}\plusendmoddef
  subroutine proton_remnant_momentum_twin_pdf &
       (this, lha_flavor, momentum_fraction,pdf)
    class(proton_remnant_t), intent(in) :: this
    integer, intent(in) :: lha_flavor
    real(default), intent(in) :: momentum_fraction
    real(default), intent(out) :: pdf
    call this%parton_twin_pdf (lha_flavor, momentum_fraction, pdf)
    pdf = pdf * momentum_fraction
  end subroutine proton_remnant_momentum_twin_pdf
  
\eatline
\nwendcode{}\nwbegindocs{2195}\nwdocspar
\nwenddocs{}\nwbegincode{2196}\moddef{Muli remnant: proton remnant: TBP}\plusendmoddef
  procedure :: momentum_twin_pdf_array => proton_remnant_momentum_twin_pdf_array
\nwendcode{}\nwbegincode{2197}\moddef{Muli remnant: procedures}\plusendmoddef
  subroutine proton_remnant_momentum_twin_pdf_array &
       (this, momentum_fraction, pdf)
    class(proton_remnant_t), intent(in) :: this
    real(default), intent(in) :: momentum_fraction
    real(default), dimension(this%n_twins), intent(out) :: pdf
    call this%parton_twin_pdf_array (momentum_fraction, pdf)
    pdf = pdf * momentum_fraction
  end subroutine proton_remnant_momentum_twin_pdf_array
  
\eatline
\nwendcode{}\nwbegindocs{2198}\nwdocspar
\nwenddocs{}\nwbegincode{2199}\moddef{Muli remnant: proton remnant: TBP}\plusendmoddef
  procedure :: momentum_kind_pdf => proton_remnant_momentum_kind_pdf
\nwendcode{}\nwbegincode{2200}\moddef{Muli remnant: procedures}\plusendmoddef
  
  subroutine proton_remnant_momentum_kind_pdf (this, GeV_scale, &
       momentum_fraction, lha_flavor, valence_pdf, sea_pdf, twin_pdf)
    class(proton_remnant_t), intent(in) :: this
    real(default), intent(in) :: GeV_scale, momentum_fraction
    integer, intent(in) :: lha_flavor              !!! g, u, d, etc.
    real(default), intent(out) :: valence_pdf, sea_pdf, twin_pdf
    real(double), dimension(-6:6) :: pdf_array
    call evolvePDF (momentum_fraction, GeV_scale, pdf_array)
    select case (lha_flavor)
    case (0) !!! gluon
       valence_pdf = zero
       sea_pdf=pdf_array(0)
    case (1) !!! down
       valence_pdf = this%get_valence_down_weight() * (pdf_array(1) - pdf_array(-1))
       sea_pdf = pdf_array(-1)
    case (2) !!! up
       valence_pdf = this%get_valence_up_weight() * (pdf_array(2) - pdf_array(-2))
       sea_pdf = pdf_array(-2)
    case default
       valence_pdf = zero
       sea_pdf = pdf_array (lha_flavor)
    end select
    sea_pdf = sea_pdf * this%get_sea_weight()
    call this%momentum_twin_pdf (lha_flavor, momentum_fraction, twin_pdf)
  end subroutine proton_remnant_momentum_kind_pdf
  
\eatline
\nwendcode{}\nwbegindocs{2201}\nwdocspar
\nwenddocs{}\nwbegincode{2202}\moddef{Muli remnant: proton remnant: TBP}\plusendmoddef
  procedure :: momentum_flavor_pdf => proton_remnant_momentum_flavor_pdf
\nwendcode{}\nwbegincode{2203}\moddef{Muli remnant: procedures}\plusendmoddef
  subroutine proton_remnant_momentum_flavor_pdf (this, GeV_scale, &
       momentum_fraction, lha_flavor, pdf)
    class(proton_remnant_t), intent(in) :: this
    real(default), intent(in) :: GeV_scale, momentum_fraction
    integer, intent(in) :: lha_flavor     !!! g, u, d, etc.
    real(default), intent(out) :: pdf
    real(default) :: valence_pdf, sea_pdf, twin_pdf
    call this%momentum_kind_pdf (GeV_scale, momentum_fraction, &
         lha_flavor, valence_pdf, sea_pdf, twin_pdf)
    pdf = valence_pdf + sea_pdf + twin_pdf
  end subroutine proton_remnant_momentum_flavor_pdf
  
\eatline
\nwendcode{}\nwbegindocs{2204}\nwdocspar
\nwenddocs{}\nwbegincode{2205}\moddef{Muli remnant: proton remnant: TBP}\plusendmoddef
  procedure :: momentum_kind_pdf_array => &
       proton_remnant_momentum_kind_pdf_array  
\nwendcode{}\nwbegincode{2206}\moddef{Muli remnant: procedures}\plusendmoddef
  subroutine proton_remnant_momentum_kind_pdf_array (this, GeV_scale, &
       momentum_fraction, valence_pdf, sea_pdf)
    class(proton_remnant_t), intent(in) :: this
    real(default), intent(in) :: GeV_scale, momentum_fraction
    real(default), dimension(2), intent(out) :: valence_pdf
    real(double), dimension(-6:6), intent(out) :: sea_pdf
    call evolvePDF (momentum_fraction, GeV_scale, sea_pdf)
    valence_pdf(1) = (sea_pdf(1) - sea_pdf(-1)) * &
         this%pdf_int_weight (PDF_VALENCE_DOWN)
    valence_pdf(2) = (sea_pdf(2) - sea_pdf(-2)) * &
         this%pdf_int_weight (PDF_VALENCE_UP)
    sea_pdf(1) = sea_pdf(-1)
    sea_pdf(2) = sea_pdf(-2)
    sea_pdf = sea_pdf * this%get_sea_weight()
    !!! no twin yet
  end subroutine proton_remnant_momentum_kind_pdf_array
  
\eatline
\nwendcode{}\nwbegindocs{2207}\nwdocspar
\nwenddocs{}\nwbegincode{2208}\moddef{Muli remnant: proton remnant: TBP}\plusendmoddef
  procedure :: momentum_flavor_pdf_array => proton_remnant_momentum_flavor_pdf_array
\nwendcode{}\nwbegincode{2209}\moddef{Muli remnant: procedures}\plusendmoddef
  subroutine proton_remnant_momentum_flavor_pdf_array &
       (this, GeV_scale, momentum_fraction, pdf)
    class(proton_remnant_t), intent(in) :: this
    real(default), intent(in) :: GeV_scale, momentum_fraction
    real(double), dimension(-6:6), intent(out) :: pdf
    real(default), dimension(2) :: valence_pdf
    call this%momentum_kind_pdf_array &
         (GeV_scale, momentum_fraction, valence_pdf, pdf)
    pdf(1:2) = pdf(1:2) + valence_pdf
    !!! no twin yet
  end subroutine proton_remnant_momentum_flavor_pdf_array
  
\eatline
\nwendcode{}\nwbegindocs{2210}\nwdocspar
\nwenddocs{}\nwbegincode{2211}\moddef{Muli remnant: proton remnant: TBP}\plusendmoddef
  procedure :: parton_twin_pdf => proton_remnant_parton_twin_pdf
\nwendcode{}\nwbegincode{2212}\moddef{Muli remnant: procedures}\plusendmoddef
  subroutine proton_remnant_parton_twin_pdf &
       (this, lha_flavor, momentum_fraction, pdf)
    class(proton_remnant_t), intent(in) :: this
    integer, intent(in) :: lha_flavor
    real(default), intent(in) :: momentum_fraction
    real(default) :: pdf
    class(muli_parton_t), pointer :: tmp_twin
    pdf = 0D0
    tmp_twin => this%twin_partons%next
    do while (associated (tmp_twin))
       if (tmp_twin%lha_flavor == lha_flavor) pdf = pdf &
            + tmp_twin%unweighted_pdf (momentum_fraction)
       tmp_twin => tmp_twin%next
    end do
    pdf = pdf * this%get_twin_weight ()
  end subroutine proton_remnant_parton_twin_pdf
  
\eatline
\nwendcode{}\nwbegindocs{2213}\nwdocspar
\nwenddocs{}\nwbegincode{2214}\moddef{Muli remnant: proton remnant: TBP}\plusendmoddef
  procedure :: parton_twin_pdf_array => proton_remnant_parton_twin_pdf_array
\nwendcode{}\nwbegincode{2215}\moddef{Muli remnant: procedures}\plusendmoddef
  subroutine proton_remnant_parton_twin_pdf_array &
       (this, momentum_fraction, pdf)
    class(proton_remnant_t), intent(in) :: this
    real(default), intent(in) :: momentum_fraction
    real(default), dimension(this%n_twins), intent(out) :: pdf
    class(muli_parton_t), pointer :: tmp_twin
    integer :: l
    tmp_twin => this%twin_partons%next
    l = 0
    do while (associated (tmp_twin))
       l = l + 1
       pdf(l) = tmp_twin%unweighted_pdf (momentum_fraction) * this%twin_norm
       tmp_twin => tmp_twin%next
    end do
  end subroutine proton_remnant_parton_twin_pdf_array
  
\eatline
\nwendcode{}\nwbegindocs{2216}\nwdocspar
\nwenddocs{}\nwbegincode{2217}\moddef{Muli remnant: proton remnant: TBP}\plusendmoddef
  procedure :: parton_kind_pdf => proton_remnant_parton_kind_pdf
\nwendcode{}\nwbegincode{2218}\moddef{Muli remnant: procedures}\plusendmoddef
  subroutine proton_remnant_parton_kind_pdf (this, GeV_scale, momentum_fraction, &
       lha_flavor, valence_pdf, sea_pdf, twin_pdf)
    class(proton_remnant_t), intent(in) :: this
    real(default), intent(in) :: GeV_scale, momentum_fraction
    integer, intent(in) :: lha_flavor              !!! g, u, d, etc.
    real(default), intent(out) :: valence_pdf, sea_pdf, twin_pdf
    call this%momentum_kind_pdf (GeV_scale, momentum_fraction, &
         lha_flavor, valence_pdf, sea_pdf, twin_pdf)
    valence_pdf = valence_pdf / momentum_fraction
    sea_pdf = sea_pdf / momentum_fraction
    twin_pdf = twin_pdf / momentum_fraction
  end subroutine proton_remnant_parton_kind_pdf
  
\eatline
\nwendcode{}\nwbegindocs{2219}\nwdocspar
\nwenddocs{}\nwbegincode{2220}\moddef{Muli remnant: proton remnant: TBP}\plusendmoddef
  procedure :: parton_flavor_pdf => proton_remnant_parton_flavor_pdf
\nwendcode{}\nwbegincode{2221}\moddef{Muli remnant: procedures}\plusendmoddef
  subroutine proton_remnant_parton_flavor_pdf (this, GeV_scale, &
       momentum_fraction, lha_flavor, pdf)
    class(proton_remnant_t), intent(in) :: this
    real(default), intent(in) :: GeV_scale, momentum_fraction
    integer, intent(in) :: lha_flavor     !g,u,d,etc.
    real(default), intent(out) :: pdf
    call this%momentum_flavor_pdf (GeV_scale, momentum_fraction, &
         lha_flavor, pdf)
    pdf = pdf / momentum_fraction
  end subroutine proton_remnant_parton_flavor_pdf
    
\eatline
\nwendcode{}\nwbegindocs{2222}\nwdocspar
\nwenddocs{}\nwbegincode{2223}\moddef{Muli remnant: proton remnant: TBP}\plusendmoddef
  procedure :: parton_kind_pdf_array => proton_remnant_parton_kind_pdf_array
\nwendcode{}\nwbegincode{2224}\moddef{Muli remnant: procedures}\plusendmoddef
  subroutine proton_remnant_parton_kind_pdf_array &
       (this, GeV_scale, momentum_fraction, valence_pdf, sea_pdf)
    class(proton_remnant_t), intent(in) :: this
    real(default), intent(in) :: GeV_scale, momentum_fraction
    real(default), dimension(2), intent(out) :: valence_pdf
    real(double), dimension(-6:6), intent(out) :: sea_pdf
    call evolvePDF (momentum_fraction, GeV_scale, sea_pdf)
    sea_pdf = sea_pdf / momentum_fraction
    valence_pdf(1) = (sea_pdf(1) - sea_pdf(-1)) * this%valence_content(1)
    valence_pdf(2) = (sea_pdf(2) - sea_pdf(-2)) * (this%valence_content(2) / two)
    sea_pdf(1) = sea_pdf(-1)
    sea_pdf(2) = sea_pdf(-2)
    valence_pdf = valence_pdf * this%get_valence_weight()
    sea_pdf = sea_pdf * this%get_sea_weight()
    !!! no twin yet
  end subroutine proton_remnant_parton_kind_pdf_array
  
\eatline
\nwendcode{}\nwbegindocs{2225}\nwdocspar
\nwenddocs{}\nwbegincode{2226}\moddef{Muli remnant: proton remnant: TBP}\plusendmoddef
  procedure :: parton_flavor_pdf_array => proton_remnant_parton_flavor_pdf_array  
\nwendcode{}\nwbegincode{2227}\moddef{Muli remnant: procedures}\plusendmoddef
  subroutine proton_remnant_parton_flavor_pdf_array &
       (this, GeV_scale, momentum_fraction, pdf)
    class(proton_remnant_t), intent(in) :: this
    real(default), intent(in) :: GeV_scale, momentum_fraction
    real(double), dimension(-6:6), intent(out) :: pdf
    real(double), dimension(2) :: valence_pdf
    real(double), dimension(-6:6) :: twin_pdf
    call msg_error ("proton_remnant_flavor_pdf_array: Not yet implemented.")
  end subroutine proton_remnant_parton_flavor_pdf_array
  
\eatline
\nwendcode{}\nwbegindocs{2228}\nwdocspar
\nwenddocs{}\nwbegincode{2229}\moddef{Muli remnant: proton remnant: TBP}\plusendmoddef
  procedure :: get_pdf_int_weight => proton_remnant_get_pdf_int_weight
\nwendcode{}\nwbegincode{2230}\moddef{Muli remnant: procedures}\plusendmoddef
  pure function proton_remnant_get_pdf_int_weight (this) result(weight)
    class(proton_remnant_t), intent(in) :: this
    real(default), dimension(5) :: weight
    weight = this%pdf_int_weight
  end function proton_remnant_get_pdf_int_weight
  
\eatline
\nwendcode{}\nwbegindocs{2231}\nwdocspar
\nwenddocs{}\nwbegincode{2232}\moddef{Muli remnant: proton remnant: TBP}\plusendmoddef
  procedure :: get_valence_down_weight => proton_remnant_get_valence_down_weight
\nwendcode{}\nwbegincode{2233}\moddef{Muli remnant: procedures}\plusendmoddef
  elemental function proton_remnant_get_valence_down_weight (this) result (weight)
    class(proton_remnant_t), intent(in) :: this
    real(default) :: weight
    weight = this%pdf_int_weight (PDF_VALENCE_DOWN)
  end function proton_remnant_get_valence_down_weight
  
\eatline
\nwendcode{}\nwbegindocs{2234}\nwdocspar
\nwenddocs{}\nwbegincode{2235}\moddef{Muli remnant: proton remnant: TBP}\plusendmoddef
  procedure :: get_valence_up_weight => proton_remnant_get_valence_up_weight
\nwendcode{}\nwbegincode{2236}\moddef{Muli remnant: procedures}\plusendmoddef
  elemental function proton_remnant_get_valence_up_weight (this) result (weight)
    class(proton_remnant_t), intent(in) :: this
    real(default) :: weight
    weight = this%pdf_int_weight (PDF_VALENCE_UP)
  end function proton_remnant_get_valence_up_weight
  
\eatline
\nwendcode{}\nwbegindocs{2237}\nwdocspar
\nwenddocs{}\nwbegincode{2238}\moddef{Muli remnant: proton remnant: TBP}\plusendmoddef
  procedure :: get_valence_weight => proton_remnant_get_valence_weight
\nwendcode{}\nwbegincode{2239}\moddef{Muli remnant: procedures}\plusendmoddef
  pure function proton_remnant_get_valence_weight (this) result (weight)
    class(proton_remnant_t), intent(in) :: this
    real(default), dimension(2) :: weight
    weight = this%pdf_int_weight(3:4)
  end function proton_remnant_get_valence_weight
  
\eatline
\nwendcode{}\nwbegindocs{2240}\nwdocspar
\nwenddocs{}\nwbegincode{2241}\moddef{Muli remnant: proton remnant: TBP}\plusendmoddef
  procedure :: get_gluon_weight => proton_remnant_get_gluon_weight
\nwendcode{}\nwbegincode{2242}\moddef{Muli remnant: procedures}\plusendmoddef
  elemental function proton_remnant_get_gluon_weight (this) result (weight)
    class(proton_remnant_t), intent(in) :: this
    real(default) :: weight
    weight = this%pdf_int_weight (PDF_GLUON)
  end function proton_remnant_get_gluon_weight

\eatline
\nwendcode{}\nwbegindocs{2243}\nwdocspar
\nwenddocs{}\nwbegincode{2244}\moddef{Muli remnant: proton remnant: TBP}\plusendmoddef
  procedure :: get_sea_weight => proton_remnant_get_sea_weight
\nwendcode{}\nwbegincode{2245}\moddef{Muli remnant: procedures}\plusendmoddef
  elemental function proton_remnant_get_sea_weight (this) result (weight)
    class(proton_remnant_t), intent(in) :: this
    real(default) :: weight
    weight = this%pdf_int_weight (PDF_SEA)
  end function proton_remnant_get_sea_weight
  
\eatline
\nwendcode{}\nwbegindocs{2246}\nwdocspar
\nwenddocs{}\nwbegincode{2247}\moddef{Muli remnant: proton remnant: TBP}\plusendmoddef
  procedure :: get_twin_weight => proton_remnant_get_twin_weight
\nwendcode{}\nwbegindocs{2248}\nwdocspar
\nwenddocs{}\nwbegincode{2249}\moddef{Muli remnant: procedures}\plusendmoddef
  elemental function proton_remnant_get_twin_weight (this) result (weight)
    class(proton_remnant_t), intent(in) :: this
    real(default) :: weight
    weight = this%pdf_int_weight (PDF_TWIN)
  end function proton_remnant_get_twin_weight
  
\eatline
\nwendcode{}\nwbegindocs{2250}\nwdocspar
\nwenddocs{}\nwbegincode{2251}\moddef{Muli remnant: proton remnant: TBP}\plusendmoddef
  procedure :: get_valence_content => proton_remnant_get_valence_content 
\nwendcode{}\nwbegincode{2252}\moddef{Muli remnant: procedures}\plusendmoddef
  pure function proton_remnant_get_valence_content (this) result (valence)
    class(proton_remnant_t), intent(in) :: this
    integer, dimension(2) :: valence
    valence = this%valence_content
  end function proton_remnant_get_valence_content

\eatline
\nwendcode{}\nwbegindocs{2253}\nwdocspar
\nwenddocs{}\nwbegincode{2254}\moddef{Muli remnant: proton remnant: TBP}\plusendmoddef
  procedure :: get_momentum_fraction => proton_remnant_get_momentum_fraction   
\nwendcode{}\nwbegincode{2255}\moddef{Muli remnant: procedures}\plusendmoddef
  elemental function proton_remnant_get_momentum_fraction (this) result (momentum)
    class(proton_remnant_t), intent(in) :: this
    real(default) :: momentum
    momentum = this%momentum_fraction
  end function proton_remnant_get_momentum_fraction
  
\eatline
\nwendcode{}\nwbegindocs{2256}\nwdocspar
\nwenddocs{}\nwbegincode{2257}\moddef{Muli remnant: proton remnant: TBP}\plusendmoddef
  procedure :: deallocate => proton_remnant_deallocate
\nwendcode{}\nwbegincode{2258}\moddef{Muli remnant: procedures}\plusendmoddef
  subroutine proton_remnant_deallocate (this)
    class(proton_remnant_t), intent(inout) :: this
    call this%is_partons%deallocate
    call this%fs_partons%deallocate
    call this%twin_partons%deallocate
    this%twin_norm = zero
    this%n_twins = 0
  end subroutine proton_remnant_deallocate
  
\eatline
\nwendcode{}\nwbegindocs{2259}\nwdocspar
\nwenddocs{}\nwbegincode{2260}\moddef{Muli remnant: proton remnant: TBP}\plusendmoddef
  procedure :: initialize => proton_remnant_initialize
\nwendcode{}\nwbegincode{2261}\moddef{Muli remnant: procedures}\plusendmoddef
  subroutine proton_remnant_initialize (this, pdf_norm)
    class(proton_remnant_t), intent(out) :: this
    class(pdfnorm_t),target, intent(in) :: pdf_norm
    this%pdf_norm => pdf_norm
  end subroutine proton_remnant_initialize
  
\eatline
\nwendcode{}\nwbegindocs{2262}\nwdocspar
\nwenddocs{}\nwbegincode{2263}\moddef{Muli remnant: proton remnant: TBP}\plusendmoddef
  procedure :: finalize => proton_remnant_finalize
\nwendcode{}\nwbegincode{2264}\moddef{Muli remnant: procedures}\plusendmoddef
  subroutine proton_remnant_finalize (this)
    class(proton_remnant_t), intent(inout) :: this
    call this%deallocate ()
    nullify (this%pdf_norm)
  end subroutine proton_remnant_finalize
  
\eatline
\nwendcode{}\nwbegindocs{2265}\nwdocspar
\nwenddocs{}\nwbegincode{2266}\moddef{Muli remnant: proton remnant: TBP}\plusendmoddef
  procedure :: apply_initial_splitting => proton_remnant_apply_initial_splitting
\nwendcode{}\nwbegincode{2267}\moddef{Muli remnant: procedures}\plusendmoddef
  subroutine proton_remnant_apply_initial_splitting &
       (this, id, pdg_flavor, x, gev_scale, rnd)
    class(proton_remnant_t), intent(inout) :: this
    integer, intent(in) :: id, pdg_flavor
    real(default), intent(in) :: x, gev_scale, rnd
    real(default) :: valence_pdf, sea_pdf, twin_pdf
    select case (pdg_flavor)
    case (PDG_FLAVOR_g)
       call this%remove_gluon (id, gev_scale, x)
    case (PDG_FLAVOR_u)
       call this%parton_kind_pdf (gev_scale, x, pdg_flavor, &
            valence_pdf, sea_pdf, twin_pdf)
       if (valence_pdf / (valence_pdf + sea_pdf) < rnd) then
          call this%remove_sea_quark (id, gev_scale, x, pdg_flavor)
       else
          call this%remove_valence_up_quark (id, gev_scale, x)
       end if
    case (PDG_FLAVOR_d)
       call this%parton_kind_pdf (gev_scale, x, &
            pdg_flavor, valence_pdf, sea_pdf, twin_pdf)
       if (valence_pdf / (valence_pdf + sea_pdf) < rnd) then
          call this%remove_sea_quark (id, gev_scale, x, pdg_flavor)
       else
          call this%remove_valence_down_quark (id, gev_scale, x)
       end if
    case default
       call this%remove_sea_quark (id, gev_scale, x, pdg_flavor)
    end select
    this%momentum_fraction = (one - x)
  end subroutine proton_remnant_apply_initial_splitting
  
\eatline
\nwendcode{}\nwbegindocs{2268}\nwdocspar
\nwenddocs{}\nwbegincode{2269}\moddef{Muli remnant: proton remnant: TBP}\plusendmoddef
  procedure :: reset => proton_remnant_reset
\nwendcode{}\nwbegincode{2270}\moddef{Muli remnant: procedures}\plusendmoddef
  subroutine proton_remnant_reset (this)
    class(proton_remnant_t), intent(inout) :: this
    call this%deallocate ()
    this%valence_content = [1, 2]
    this%pdf_int_weight = [one, one, one, one, one]
    this%momentum_fraction = one
  end subroutine proton_remnant_reset
  
\eatline
\nwendcode{}\nwbegindocs{2271}\nwdocspar
\nwenddocs{}\nwbegincode{2272}\moddef{Muli remnant: proton remnant: TBP}\plusendmoddef
  procedure :: calculate_weight => proton_remnant_calculate_weight
\nwendcode{}\nwbegincode{2273}\moddef{Muli remnant: procedures}\plusendmoddef
  subroutine proton_remnant_calculate_weight (this, GeV_scale)
    class(proton_remnant_t), intent(inout) :: this
    real(default), intent(in) :: GeV_scale
    real(default) :: all, gluon, sea, vu, vd, valence, twin, weight
    call this%pdf_norm%get_norm (GeV_scale, 1, 0, all)
    call this%pdf_norm%get_norm (GeV_scale, 1, PDF_GLUON, gluon)
    call this%pdf_norm%get_norm (GeV_scale, 1, PDF_SEA, sea)
    call this%pdf_norm%get_norm (GeV_scale, 1, PDF_VALENCE_DOWN, vd)
    call this%pdf_norm%get_norm (GeV_scale, 1, PDF_VALENCE_UP, vu)
    valence = vd * this%valence_content (LHA_FLAVOR_d) + &
              vu * this%valence_content (LHA_FLAVOR_u) / two
    twin = this%twin_norm / all
    ! print *, all, gluon + sea + valence + twin, gluon, sea, valence, twin
    ! pdf_rescale = (one - n_d_valence * mean_d1 - n_u_valence * mean_u2) / &
    !     (1.E-1_default * mean_d1 - two * mean_u2) !!! pythia
    select case (remnant_weight_model)
    case (0) !!! no reweighting
       this%pdf_int_weight = [one, one, one, one, one]
    case (2) !!! pythia-like, only sea
       weight = (one - valence - twin) / (sea + gluon)
       this%pdf_int_weight = [weight, weight, one, one, one]
    case (3) !!! only valence and twin
       weight = (one - sea - gluon) / (valence + twin)
       this%pdf_int_weight = [one, one, weight, weight, weight]
    case (4) !!! only sea and twin
       weight = (one - valence) / (sea + gluon + twin)
       this%pdf_int_weight = [one, weight, one, one, weight]
    case default !!! equal weight
       weight = one / (valence + sea + gluon + twin)
       this%pdf_int_weight = [weight, weight, weight, weight, weight]
    end select
    this%pdf_int_weight(PDF_VALENCE_DOWN) = &
         this%pdf_int_weight(PDF_VALENCE_DOWN) * this%valence_content(1)
    this%pdf_int_weight(PDF_VALENCE_UP) = &
         this%pdf_int_weight(PDF_VALENCE_UP) * this%valence_content(2) * &
         5.E-1_default
    ! print('("New rescale factors are: ",2(I10),7(E14.7))'),&
    !         this%valence_content,&
    !         this%pdf_int_weight,&
    !         sea_norm,&
    !         valence_norm,&
    !         this%twin_norm
  end subroutine proton_remnant_calculate_weight
  
\eatline
\nwendcode{}\nwbegindocs{2274}\nwdocspar
\nwenddocs{}\nwbegincode{2275}\moddef{Muli remnant: proton remnant: TBP}\plusendmoddef
  procedure :: push_is_parton => proton_remnant_push_is_parton
\nwendcode{}\nwbegincode{2276}\moddef{Muli remnant: procedures}\plusendmoddef
  subroutine proton_remnant_push_is_parton &
       (this, id, lha_flavor, momentum_fraction)
    class(proton_remnant_t), intent(inout) :: this
    integer, intent(in) :: id, lha_flavor
    real(default), intent(in) :: momentum_fraction
    class(muli_parton_t), pointer :: tmp_parton
    allocate (tmp_parton)
    tmp_parton%id = id
    tmp_parton%lha_flavor = lha_flavor
    tmp_parton%momentum = momentum_fraction
    call this%is_partons%push (tmp_parton)
  end subroutine proton_remnant_push_is_parton
  
\eatline
\nwendcode{}\nwbegindocs{2277}\nwdocspar
\nwenddocs{}\nwbegincode{2278}\moddef{Muli remnant: proton remnant: TBP}\plusendmoddef
  procedure :: push_twin => proton_remnant_push_twin
\nwendcode{}\nwbegincode{2279}\moddef{Muli remnant: procedures}\plusendmoddef
  subroutine proton_remnant_push_twin &
       (this, id, lha_flavor, momentum_fraction, gev_scale)
    class(proton_remnant_t), intent(inout) :: this
    integer, intent(in) :: id, lha_flavor !!! of IS parton
    real(default), intent(in) :: momentum_fraction !!! of IS parton
    real(default), intent(in) :: GeV_scale
    class(muli_parton_t), pointer :: new_is, new_twin
    real(default) :: norm
    ! print *, "proton_remnant_push_twin", momentum_fraction
    allocate (new_is)    
    allocate (new_twin)
    !!! IS initialization
    new_is%id = id
    new_is%lha_flavor = lha_flavor
    new_is%momentum = momentum_fraction
    new_is%twin => new_twin
    !!! twin initialization
    new_twin%id = -id
    new_twin%lha_flavor = -lha_flavor
    new_twin%momentum = remnant_twin_momentum_4 (momentum_fraction)
    new_twin%twin => new_is
    !!! remnant update
    this%n_twins = this%n_twins + 1
    this%twin_norm = this%twin_norm + new_twin%momentum
    call this%is_partons%push (new_is)
    call this%twin_partons%push (new_twin)
    call this%calculate_weight (GeV_scale)
  end subroutine proton_remnant_push_twin
  
\eatline
\nwendcode{}\nwbegindocs{2280}\nwdocspar
\nwenddocs{}\nwbegincode{2281}\moddef{Muli remnant: proton remnant: TBP}\plusendmoddef
  procedure :: calculate_twin_norm => proton_remnant_calculate_twin_norm
\nwendcode{}\nwbegincode{2282}\moddef{Muli remnant: procedures}\plusendmoddef
  subroutine proton_remnant_calculate_twin_norm (this)
    class(proton_remnant_t), intent(inout) :: this
    class(muli_parton_t), pointer :: twin
    integer :: n
    if (associated (this%twin_partons%next)) then
       this%twin_norm = zero
       twin => this%twin_partons%next
       do while (associated (twin))
          this%twin_norm = this%twin_norm + twin%momentum
          twin => twin%next
       end do
    else
       this%twin_norm = zero
    end if
  end subroutine proton_remnant_calculate_twin_norm
  
\eatline
\nwendcode{}\nwbegindocs{2283}\nwdocspar
\nwenddocs{}\nwbegincode{2284}\moddef{Muli remnant: proton remnant: TBP}\plusendmoddef
  procedure :: replace_is_parton => proton_remnant_replace_is_parton  
\nwendcode{}\nwbegincode{2285}\moddef{Muli remnant: procedures}\plusendmoddef
  subroutine proton_remnant_replace_is_parton &
       (this, old_id, new_id, pdg_f, x_proton, gev_scale)
    class(proton_remnant_t), intent(inout) :: this
    integer, intent(in) :: old_id, new_id, pdg_f
    real(default), intent(in) :: x_proton, gev_scale
    class(muli_parton_t), pointer :: old_is_parton
    integer :: lha_flavor
    real(default) :: momentum_fraction
    momentum_fraction = x_proton / this%momentum_fraction
    !!! convert PDG flavor numbers to LHA flavor numbers
    if (pdg_f == PDG_FLAVOR_g) then
       lha_flavor = LHA_FLAVOR_g
    else
       lha_flavor = pdg_f
    end if
    !!! we remove the old initial state parton from initial state stack.
    call this%is_partons%pop (old_id, old_is_parton)
    !!! this check has no physical meaning, it's just a check for consistency.
    if (associated (old_is_parton)) then
       !!! do we emit a gluon?
       if (lha_flavor == old_is_parton%lha_flavor) then
          !!! has the old initial state parton been a sea quark?
          if (associated (old_is_parton%twin)) then
             !!! the connection of the old IS parton with its twin was 
             !!!    provisional. We remove it now
             call this%twin_partons%pop (old_is_parton%twin)
             call this%fs_partons%push (old_is_parton%twin)
             this%n_twins = this%n_twins - 1
             !!! and generate a new initial state parton - twin pair.
             call this%push_twin &
                  (new_id, lha_flavor, momentum_fraction, gev_scale)
          else
             !!! there is no twin, so we just insert the new initial state parton.
             call this%push_is_parton (new_id, lha_flavor, momentum_fraction)
          end if
       else
          ! we emit a quark. is this a g->qqbar splitting?
          if (lha_flavor==LHA_FLAVOR_g) then
             !!! we insert the new initial state gloun.
             call this%push_is_parton (new_id, lha_flavor, momentum_fraction)
             !!! has the old initial state quark got a twin?
             if (associated (old_is_parton%twin)) then
                !!! we assume that this twin is the second splitting particle. 
                !!! so the twin becomes a final state particle now and 
                !!! must be removed from the is stack.
                call this%remove_twin (-old_id, GeV_scale)
             else
                !!! the old initial state quark has been a valence quark. 
                !!! what should we do now? is this splitting sensible at all? 
                !!! we don't know but allow these splittings. The most trivial
                !!! treatment is to restore the former valence quark.
                this%valence_content (old_is_parton%lha_flavor) = &
                     this%valence_content (old_is_parton%lha_flavor) + 1
             end if
          else
             !!! this is a q->qg splitting. the new initial state quark emits 
             !!! the preceding initial state gluon. yeah, backward evolution is 
             !!! confusing! the new initial state quark is not part of the 
             !!! proton remnant any longer. how do we remove a quark from 
             !!! the remnant? we add a conjugated twin parton and assume, 
             !!! that this twin is created in a not yet resolved g->qqbar splitting.
             call this%push_twin (new_id, lha_flavor, momentum_fraction, gev_scale)
          end if
       end if
       !!! everything is done. what shall we do with the old initial state 
       !!! parton? we don't need it any more but we store it anyway.
       call this%fs_partons%push (old_is_parton)
       !!! the new initial state parton has taken away momentum, so we update 
       !!! the remnant momentum fraction.
       this%momentum_fraction = this%momentum_fraction * (1 - &
            momentum_fraction) / (1 - old_is_parton%momentum)
    else
       !!! this is a bug.
       print *, "proton_remnant_replace_is_parton: parton #", old_id, &
            " not found on ISR stack."
       if (associated (this%is_partons%next)) then
          print *, "actual content of isr stack:"
          call this%is_partons%next%print_peers ()
       else
          print *, "ISR stack is not associated."
       end if
       stop
    end if
  end subroutine proton_remnant_replace_is_parton
  
\eatline
\nwendcode{}\nwbegindocs{2286}\nwdocspar
\nwenddocs{}\nwbegincode{2287}\moddef{Muli remnant: proton remnant: TBP}\plusendmoddef
  procedure :: write_to_marker => proton_remnant_write_to_marker
\nwendcode{}\nwbegincode{2288}\moddef{Muli remnant: procedures}\plusendmoddef
  subroutine proton_remnant_write_to_marker (this, marker, status)
    class(proton_remnant_t), intent(in) :: this
    class(marker_t), intent(inout) :: marker
    integer(dik), intent(out) :: status
    call marker%mark_begin ("proton_remnant_t")
    call marker%mark ("valence_content", this%valence_content)
    call marker%mark ("momentum_fraction", this%momentum_fraction)
    call marker%mark ("pdf_int_weight", this%pdf_int_weight)
    call marker%mark_end ("proton_remnant_t")
  end subroutine proton_remnant_write_to_marker
  
\eatline
\nwendcode{}\nwbegindocs{2289}\nwdocspar
\nwenddocs{}\nwbegincode{2290}\moddef{Muli remnant: proton remnant: TBP}\plusendmoddef
  procedure :: read_from_marker => proton_remnant_read_from_marker
\nwendcode{}\nwbegincode{2291}\moddef{Muli remnant: procedures}\plusendmoddef

  subroutine proton_remnant_read_from_marker (this, marker, status)
    class(proton_remnant_t), intent(out) :: this
    class(marker_t), intent(inout) :: marker
    integer(dik), intent(out) :: status
    character(:), allocatable :: name
    call marker%pick_begin ("proton_remnant_t", status=status)
    call marker%pick ("valence_content", this%valence_content, status)
    call marker%pick ("momentum_fraction", this%momentum_fraction, status)
    call marker%pick ("pdf_int_weight", this%pdf_int_weight, status)
    call marker%pick_end ("proton_remnant_t", status=status)
  end subroutine proton_remnant_read_from_marker
  
\eatline
\nwendcode{}\nwbegindocs{2292}\nwdocspar
\nwenddocs{}\nwbegincode{2293}\moddef{Muli remnant: proton remnant: TBP}\plusendmoddef
  procedure :: print_to_unit => proton_remnant_print_to_unit
\nwendcode{}\nwbegincode{2294}\moddef{Muli remnant: procedures}\plusendmoddef
  subroutine proton_remnant_print_to_unit (this, unit, parents, components, peers)
    class(proton_remnant_t), intent(in) :: this
    integer, intent(in) :: unit
    integer(dik), intent(in) :: parents, components, peers
    write (unit, '("Components of proton_remnant_t:")')
    write (unit, '("Valence Content:           ",I1,":",I1)')this&
         &%valence_content
    write (unit, "(1x,A,I1)") "N Twins:                   ", this%n_twins
    write (unit, "(1x,A,5(F7.3))") "INT weights [g,s,d,u,t]    ", &
         this%pdf_int_weight
    write (unit, "(1x,A,F7.3)") "Total Momentum Fraction:   ", &
         this%momentum_fraction
    write (unit, "(1x,A,F7.3)") "Twin Norm:            ", this%twin_norm
  end subroutine proton_remnant_print_to_unit
  
\eatline
\nwendcode{}\nwbegindocs{2295}\nwdocspar
\nwenddocs{}\nwbegincode{2296}\moddef{Muli remnant: proton remnant: TBP}\plusendmoddef
  procedure, nopass :: get_type => proton_remnant_get_type  
\nwendcode{}\nwbegincode{2297}\moddef{Muli remnant: procedures}\plusendmoddef
  pure subroutine proton_remnant_get_type (type)
    character(:), allocatable, intent(out) :: type
    allocate (type, source="proton_remnant_t")
  end subroutine proton_remnant_get_type
  
\eatline
\nwendcode{}\nwbegindocs{2298}\nwdocspar
\nwenddocs{}\nwbegincode{2299}\moddef{Muli remnant: proton remnant: TBP}\plusendmoddef
  procedure :: gnuplot_momentum_kind_pdf_array => &
       proton_remnant_gnuplot_momentum_kind_pdf_array    
\nwendcode{}\nwbegincode{2300}\moddef{Muli remnant: procedures}\plusendmoddef
  subroutine proton_remnant_gnuplot_momentum_kind_pdf_array &
       (this, momentum_unit, parton_unit, GeV_scale)
    class(proton_remnant_t), intent(in) :: this
    integer, intent(in) :: momentum_unit, parton_unit
    real(default), intent(in) :: GeV_scale
    real(default), dimension(2) :: valence_pdf
    real(double), dimension(-6:6) :: sea_pdf
    real(default), dimension(this%n_twins) :: twin_pdf
    integer :: x
    real(default) :: momentum_fraction
    do x = 1, 100
       momentum_fraction = x * 1E-2_default
       call this%momentum_kind_pdf_array (GeV_scale, momentum_fraction, &
            valence_pdf, sea_pdf)
       call this%momentum_twin_pdf_array (momentum_fraction, twin_pdf)
       write (momentum_unit, fmt=*) momentum_fraction, &
            sum(valence_pdf) + sum(sea_pdf) + sum(twin_pdf), &
            valence_pdf, sea_pdf, twin_pdf
       call this%parton_kind_pdf_array (GeV_scale, momentum_fraction, &
            valence_pdf, sea_pdf)
       call this%parton_twin_pdf_array (momentum_fraction, twin_pdf)
       write(parton_unit,fmt=* )momentum_fraction, &
            sum(valence_pdf) + sum(sea_pdf) + sum(twin_pdf), &
            valence_pdf, sea_pdf, twin_pdf       
    end do
  end subroutine proton_remnant_gnuplot_momentum_kind_pdf_array
  
\eatline
\nwendcode{}\nwbegindocs{2301}\nwdocspar
\nwenddocs{}\nwbegincode{2302}\moddef{Muli remnant: public}\plusendmoddef
  public :: pp_remnant_t
\nwendcode{}\nwbegincode{2303}\moddef{Muli remnant: types}\plusendmoddef
  type, extends (ser_class_t) :: pp_remnant_t
     logical :: initialized = .false.
     real(default), private :: gev_initial_cme = gev_cme_tot
     real(default), private :: X = one
     type(proton_remnant_t), dimension(2) :: proton
     class(pdfnorm_t), pointer :: pdf_norm
   contains
   \LA{}Muli remnant: PP remnant: TBP\RA{}
  end type pp_remnant_t
  
\eatline
\nwendcode{}\nwbegindocs{2304}\nwdocspar
\nwenddocs{}\nwbegincode{2305}\moddef{Muli remnant: PP remnant: TBP}\endmoddef
  procedure :: write_to_marker => pp_remnant_write_to_marker
\nwendcode{}\nwbegincode{2306}\moddef{Muli remnant: procedures}\plusendmoddef
  subroutine pp_remnant_write_to_marker (this, marker, status)
    class(pp_remnant_t), intent(in) :: this
    class(marker_t), intent(inout) :: marker
    integer(dik), intent(out) :: status
    call marker%mark_begin ("pp_remnant_t")
    call marker%mark ("gev_initial_cme", this%gev_initial_cme)
    call marker%mark ("X", this%X)
    call this%proton(1)%write_to_marker (marker, status)
    call this%proton(2)%write_to_marker (marker, status)
    call marker%mark_end ("pp_remnant_t")
  end subroutine pp_remnant_write_to_marker
  
\eatline
\nwendcode{}\nwbegindocs{2307}\nwdocspar
\nwenddocs{}\nwbegincode{2308}\moddef{Muli remnant: PP remnant: TBP}\plusendmoddef
  procedure :: read_from_marker => pp_remnant_read_from_marker
\nwendcode{}\nwbegincode{2309}\moddef{Muli remnant: procedures}\plusendmoddef
  subroutine pp_remnant_read_from_marker (this, marker, status)
    class(pp_remnant_t), intent(out) :: this
    class(marker_t), intent(inout) :: marker
    integer(dik), intent(out) :: status
    character(:), allocatable :: name
    call marker%pick_begin ("pp_remnant_t", status=status)
    call marker%pick ("gev_initial_cme", this%gev_initial_cme, status)
    call marker%pick ("X", this%X, status)
    call this%proton(1)%read_from_marker (marker, status)
    call this%proton(2)%read_from_marker (marker, status)
    call marker%pick_end ("pp_remnant_t", status=status)
  end subroutine pp_remnant_read_from_marker
  
\eatline
\nwendcode{}\nwbegindocs{2310}\nwdocspar
\nwenddocs{}\nwbegincode{2311}\moddef{Muli remnant: PP remnant: TBP}\plusendmoddef
  procedure :: print_to_unit => pp_remnant_print_to_unit
\nwendcode{}\nwbegincode{2312}\moddef{Muli remnant: procedures}\plusendmoddef
  subroutine pp_remnant_print_to_unit (this, unit, parents, components, peers)
    class(pp_remnant_t), intent(in) :: this
    integer, intent(in) :: unit
    integer(dik), intent(in) :: parents, components, peers
    write (unit, "(1x,A)")  "Components of pp_remnant_t:"
    write (unit, "(1x,A,F10.3)") "Initial center of mass energy: ", &
         this%gev_initial_cme
    write (unit, "(1x,A,F10.3)") "Actual center of mass energy:  ", &
         this%get_gev_actual_cme ()
    write (unit, "(1x,A,F10.3)") "Total Momentum Fraction is:    ", this%X
    if (components>0) then
       write (unit, "(3x,A)") "Proton 1:"
       call this%proton(1)%print_to_unit (unit, parents, components-1, peers)
       write (unit, "(3x,A)") "Proton 2:"
       call this%proton(2)%print_to_unit (unit, parents, components-1, peers)
    end if
    ! write (unit, "(1x,A,F7.2)") "Total Momentum Fraction:   ", &
    !    this%momentum_fraction
  end subroutine pp_remnant_print_to_unit
  
\eatline
\nwendcode{}\nwbegindocs{2313}\nwdocspar
\nwenddocs{}\nwbegincode{2314}\moddef{Muli remnant: PP remnant: TBP}\plusendmoddef
  procedure, nopass :: get_type => pp_remnant_get_type  
\nwendcode{}\nwbegincode{2315}\moddef{Muli remnant: procedures}\plusendmoddef
  pure subroutine pp_remnant_get_type (type)
    character(:), allocatable, intent(out) :: type
    allocate (type, source="pp_remnant_t")
  end subroutine pp_remnant_get_type
  
\eatline
\nwendcode{}\nwbegindocs{2316}\nwdocspar
\nwenddocs{}\nwbegincode{2317}\moddef{Muli remnant: PP remnant: TBP}\plusendmoddef
  procedure :: initialize => pp_remnant_initialize
\nwendcode{}\nwbegincode{2318}\moddef{Muli remnant: procedures}\plusendmoddef
  subroutine pp_remnant_initialize (this, muli_dir, &
       lhapdf_dir, lhapdf_file, lhapdf_member)
    class(pp_remnant_t), intent(out) :: this
    character(*), intent(in) :: muli_dir, lhapdf_dir, lhapdf_file
    integer, intent(in) :: lhapdf_member
    logical :: exist
    allocate (this%pdf_norm)
    ! call InitPDFset (lhapdf_dir // lhapdf_file)
    ! call InitPDF (lhapdf_member)
    print *, "looking for previously generated pdf integrals..."
    inquire (file=muli_dir // "/pdf_norm_" // lhapdf_file // &
         ".xml", exist=exist)
    if (exist) then
       print *, "found. Starting deserialization..."
       call this%pdf_norm%deserialize (name="pdf_norm_" // lhapdf_file, &
            file=muli_dir // "/pdf_norm_" // lhapdf_file // ".xml")
       print *, "done."
    else
       print *, "No integrals found. Starting generation..."
       call this%pdf_norm%scan ()
       print *, "done."
       call this%pdf_norm%serialize (name="pdf_norm_" // lhapdf_file, &
            file=muli_dir // "/pdf_norm_" // lhapdf_file // ".xml")
    end if
    call this%proton(1)%initialize (this%pdf_norm)
    call this%proton(2)%initialize (this%pdf_norm)
    this%initialized = .true.
    ! call this%print_all ()
  end subroutine pp_remnant_initialize
  
\eatline
\nwendcode{}\nwbegindocs{2319}\nwdocspar
\nwenddocs{}\nwbegincode{2320}\moddef{Muli remnant: PP remnant: TBP}\plusendmoddef
  procedure :: finalize => pp_remnant_finalize
\nwendcode{}\nwbegincode{2321}\moddef{Muli remnant: procedures}\plusendmoddef
  subroutine pp_remnant_finalize (this)
    class(pp_remnant_t), intent(inout) :: this
    call this%proton(1)%finalize ()
    call this%proton(2)%finalize ()
    deallocate (this%pdf_norm)
  end subroutine pp_remnant_finalize
  
\eatline
\nwendcode{}\nwbegindocs{2322}\nwdocspar
\nwenddocs{}\nwbegincode{2323}\moddef{Muli remnant: PP remnant: TBP}\plusendmoddef
  procedure :: apply_initial_interaction => &
       pp_remnant_apply_initial_interaction
\nwendcode{}\nwbegincode{2324}\moddef{Muli remnant: procedures}\plusendmoddef
  subroutine pp_remnant_apply_initial_interaction (this, gev_cme, x1, &
       x2, pdg_f1, pdg_f2, n1, n2, gev_scale, rnd1, rnd2)
    class(pp_remnant_t), intent(inout) :: this
    real(default), intent(in)::gev_cme,x1,x2,gev_scale,rnd1,rnd2
    integer, intent(in) :: pdg_f1, pdg_f2, n1, n2
    if (this%initialized) then
       call this%proton(1)%apply_initial_splitting &
            (n1, pdg_f1, x1, gev_scale, rnd1)
       call this%proton(2)%apply_initial_splitting &
            (n2, pdg_f2, x2, gev_scale, rnd2)
       this%X = (one - x1) * (one - x2)
       this%gev_initial_cme = gev_cme
       ! call this%print_all ()
    else
       call msg_error ("pp_remnant_apply_initial_interaction: Not yet " // &
            "initialized, call pp_remnant_initialize first!")
    end if
  end subroutine pp_remnant_apply_initial_interaction
 
\eatline
\nwendcode{}\nwbegindocs{2325}\nwdocspar
\nwenddocs{}\nwbegincode{2326}\moddef{Muli remnant: PP remnant: TBP}\plusendmoddef
  procedure :: reset => pp_remnant_reset  
\nwendcode{}\nwbegincode{2327}\moddef{Muli remnant: procedures}\plusendmoddef
  subroutine pp_remnant_reset (this)
    class(pp_remnant_t), intent(inout) :: this
    call this%proton(1)%reset ()
    call this%proton(2)%reset ()
    this%X = one
  end subroutine pp_remnant_reset
  
\eatline
\nwendcode{}\nwbegindocs{2328}\nwdocspar
\nwenddocs{}\nwbegincode{2329}\moddef{Muli remnant: PP remnant: TBP}\plusendmoddef
  procedure :: replace_parton => pp_remnant_replace_parton
\nwendcode{}\nwbegincode{2330}\moddef{Muli remnant: procedures}\plusendmoddef
  subroutine pp_remnant_replace_parton (this, proton_id, old_id, &
       new_id, pdg_f, x_proton, gev_scale)
    class(pp_remnant_t), intent(inout) :: this
    integer, intent(in) :: proton_id, old_id, new_id, pdg_f
    real(default), intent(in) :: x_proton, gev_scale
    call this%proton(proton_id)%replace_is_parton &
         (old_id, new_id, pdg_f, x_proton, gev_scale)
  end subroutine pp_remnant_replace_parton
  
\eatline
\nwendcode{}\nwbegindocs{2331}\nwdocspar
\nwenddocs{}\nwbegincode{2332}\moddef{Muli remnant: PP remnant: TBP}\plusendmoddef
  procedure :: momentum_pdf => pp_remnant_momentum_pdf
\nwendcode{}\nwbegincode{2333}\moddef{Muli remnant: PP remnant: TBP}\plusendmoddef
  procedure :: parton_pdf => pp_remnant_parton_pdf
\nwendcode{}\nwbegincode{2334}\moddef{Muli remnant: procedures}\plusendmoddef
  subroutine pp_remnant_parton_pdf (this, x_proton, gev2_scale, n, pdg_f, pdf)
    class(pp_remnant_t), intent(in) :: this
    real(default), intent(in) :: x_proton, gev2_scale
    integer, intent(in) :: n, pdg_f
    real(default), intent(out) :: pdf
    if (n==1 .or. n==2) then
       if (x_proton <= this%proton(n)%momentum_fraction) then
          if (pdg_f == PDG_FLAVOR_g) then
             call this%proton(n)%parton_flavor_pdf (sqrt(GeV2_scale), &
                  x_proton/this%proton(n)%momentum_fraction, LHA_FLAVOR_g, &
                  pdf)
          else
             call this%proton(n)%parton_flavor_pdf (sqrt(GeV2_scale), &
                  x_proton/this%proton(n)%momentum_fraction,pdg_f, pdf)
          end if
          pdf = pdf * this%proton(n)%momentum_fraction
       else
          pdf = zero
       end if
    else
       call msg_error ("pp_remnant_parton_pdf: n must be either 1 or 2, " &
            // "but it is " // char (str (n)) // ".")
       stop
    end if
  end subroutine pp_remnant_parton_pdf
  
\eatline
\nwendcode{}\nwbegindocs{2335}\nwdocspar
\nwenddocs{}\nwbegincode{2336}\moddef{Muli remnant: PP remnant: TBP}\plusendmoddef
  procedure :: apply_interaction => pp_remnant_apply_interaction
\nwendcode{}\nwbegincode{2337}\moddef{Muli remnant: PP remnant: TBP}\plusendmoddef
  procedure :: get_pdf_int_weights => pp_remnant_get_pdf_int_weights
\nwendcode{}\nwbegincode{2338}\moddef{Muli remnant: PP remnant: TBP}\plusendmoddef
  procedure :: get_pdf_int_weight => pp_remnant_get_pdf_int_weight
\nwendcode{}\nwbegincode{2339}\moddef{Muli remnant: PP remnant: TBP}\plusendmoddef
  procedure :: set_pdf_weight => pp_remnant_set_pdf_weight
\nwendcode{}\nwbegincode{2340}\moddef{Muli remnant: PP remnant: TBP}\plusendmoddef
  procedure :: get_gev_initial_cme => pp_remnant_get_gev_initial_cme
\nwendcode{}\nwbegincode{2341}\moddef{Muli remnant: PP remnant: TBP}\plusendmoddef
  procedure :: get_gev_actual_cme => pp_remnant_get_gev_actual_cme
\nwendcode{}\nwbegincode{2342}\moddef{Muli remnant: PP remnant: TBP}\plusendmoddef
  procedure :: get_cme_fraction => pp_remnant_get_cme_fraction
\nwendcode{}\nwbegincode{2343}\moddef{Muli remnant: PP remnant: TBP}\plusendmoddef
  procedure :: get_proton_remnant_momentum_fractions => &
       pp_remnant_get_proton_remnant_momentum_fractions
\nwendcode{}\nwbegincode{2344}\moddef{Muli remnant: PP remnant: TBP}\plusendmoddef
  procedure :: get_proton_remnants => pp_remnant_get_proton_remnants
\nwendcode{}\nwbegincode{2345}\moddef{Muli remnant: PP remnant: TBP}\plusendmoddef
  procedure :: get_remnant_parton_flavor_pdf_arrays => &
       pp_remnant_get_remnant_parton_flavor_pdf_arrays     
\nwendcode{}\nwbegincode{2346}\moddef{Muli remnant: interfaces}\endmoddef
  interface 
     subroutine getxmin (mem, xmin)
       integer, intent(in) :: mem
       double precision, intent(out) :: xmin
     end subroutine getxmin
  end interface
\eatline
\nwendcode{}\nwbegindocs{2347}\nwdocspar
\nwenddocs{}\nwbegincode{2348}\moddef{Muli remnant: interfaces}\plusendmoddef
  interface 
     subroutine getxmax (mem, xmax)
       integer, intent(in) :: mem
       double precision, intent(out) :: xmax
     end subroutine getxmax
  end interface
\eatline
\nwendcode{}\nwbegindocs{2349}\nwdocspar
\nwenddocs{}\nwbegincode{2350}\moddef{Muli remnant: interfaces}\plusendmoddef
  interface
     subroutine getq2min (mem, q2min)
       integer, intent(in) :: mem
       double precision, intent(out) :: q2min
     end subroutine getq2min
  end interface
\eatline
\nwendcode{}\nwbegindocs{2351}\nwdocspar
\nwenddocs{}\nwbegincode{2352}\moddef{Muli remnant: interfaces}\plusendmoddef
  interface
     subroutine getq2max (mem, q2max)
       integer, intent(in) :: mem
       double precision, intent(out) :: q2max
     end subroutine getq2max
  end interface
  
\eatline
\nwendcode{}\nwbegindocs{2353}\nwdocspar
%%%%%%%%%%%%%%%%%%%%%%%%%%%%%%%%%%%%%%%%%%%%%%%%%%%%%%%%%%%%%%%%%%%%%%%%
\nwenddocs{}\nwbegincode{2354}\moddef{\code{}muli.f90\edoc{}}\endmoddef
!!! module: muli
!!! This code is part of my Ph.D studies.
!!! 
!!! Copyright (C) 2011 Hans-Werner Boschmann <boschmann@tp1.physik.uni-siegen.de>
!!! 
!!! This program is free software; you can redistribute it and/or modify it
!!! under the terms of the GNU General Public License as published by the Free 
!!! Software Foundation; either version 3 of the License, or (at your option) 
!!! any later version.
!!! 
!!! This program is distributed in the hope that it will be useful, but WITHOUT
!!! ANY WARRANTY; without even the implied warranty of MERCHANTABILITY or 
!!! FITNESS FOR A PARTICULAR PURPOSE. See the GNU General Public License for
!!! more details.
!!! 
!!! You should have received a copy of the GNU General Public License along
!!! with this program; if not, see <http://www.gnu.org/licenses/>.
!!! 
!!! Latest Change: 2011-08-03 11:12:37 CEST(+0200)
!!! 
!!!!!!!!!!!!!!!!!!!!!!!!!!!!!!!!!!!!!!!!!!!!!!!!!!!!!!!!!!!!!!!!!!!!!!!!!!!!!!!!

!!! This file contains the module "muli" which is the multiple parton 
!!! interactions interface module to WHIZARD. muli is supposed to run together
!!! with initial state radiation. Both share a momentum evolution variable and 
!!! compete for beam momentum, so the generation of this scale variable must
!!! be fully transparent to WHIZARD. That's why the corresponding procedures
!!! are here, while all other work like phase space integration, flavour
!!! generation and treatment of the beam remnant is put into muli_dsigma, 
!!! muli_mcint and muli_remnant respectively.

!!! qcd_2_2_type is a container class for properties of qcd 2->2 interactions.
!!! It holds a very condensed internal representation and offers a convenient
!!! set of TBP to query all aspects without the burden of the generator 
!!! internals.

!!! muli_t then is an extension of qcd_2_2_type that adds generator
!!! internals like random number generator, integrated cross-sections, the
!!! actual Monte Carlo generator for flavour genertation and beam remnants in
!!! tao_rnd, dsigma, samples and beam respectively.

module muli
  use, intrinsic :: iso_fortran_env
\LA{}Use kinds\RA{}  
  use constants
  use tao_random_numbers !NODEP!
  use muli_base
  use muli_momentum
  use muli_trapezium
  use muli_interactions
  use muli_dsigma
  use muli_mcint
  use muli_remnant

\LA{}Standard module head\RA{}

  logical,parameter::muli_default_modify_pdfs=.true.
  integer,parameter::muli_default_lhapdf_member=0
  character(*),parameter::muli_default_lhapdf_file="cteq6ll.LHpdf"
  
  type, extends(qcd_2_2_class) :: qcd_2_2_type
     private
     integer::process_id=-1
     integer::integrand_id=-1
     integer, dimension(2)::parton_ids=[0,0]
     integer, dimension(4)::flow=[0,0,0,0]
     real(default), dimension(3) :: momentum_fractions = [-one, -one, -one]
     real(default), dimension(3) :: hyperbolic_fractions = [-one ,- one,- one]
   contains
     !overridden ser_class_t procedures
     procedure,public::write_to_marker=>qcd_2_2_write_to_marker
     procedure,public::read_from_marker=>qcd_2_2_read_from_marker
     procedure,public::print_to_unit=>qcd_2_2_print_to_unit
     procedure,public,nopass::get_type=>qcd_2_2_get_type
     ! new type-bound-procedures
     procedure,public::get_process_id=>qcd_2_2_get_process_id
     procedure,public::get_integrand_id=>qcd_2_2_get_integrand_id
     procedure,public::get_diagram_kind=>qcd_2_2_get_diagram_kind
     procedure,public::get_diagram_color_kind=>qcd_2_2_get_diagram_color_kind
     procedure,public::get_io_kind=>qcd_2_2_get_io_kind
     procedure,public::get_lha_flavors=>qcd_2_2_get_lha_flavors
     procedure,public::get_pdg_flavors=>qcd_2_2_get_pdg_flavors
     procedure,public::get_parton_id=>qcd_2_2_get_parton_id
     procedure,public::get_parton_kinds=>qcd_2_2_get_parton_kinds
     procedure,public::get_pdf_int_kinds=>qcd_2_2_get_pdf_int_kinds
     procedure,public::get_momentum_boost=>qcd_2_2_get_momentum_boost
     procedure,public::get_hyperbolic_fractions=>qcd_2_2_get_hyperbolic_fractions
     procedure,public::get_remnant_momentum_fractions=>qcd_2_2_get_remnant_momentum_fractions
     procedure,public::get_total_momentum_fractions=>qcd_2_2_get_total_momentum_fractions
     procedure,public::get_color_flow=>qcd_2_2_get_color_flow
     procedure,public::get_color_correlations=>qcd_2_2_get_color_correlations
     procedure,public::qcd_2_2_initialize
     generic,public::initialize=>qcd_2_2_initialize
  end type qcd_2_2_type

  public :: muli_t
  type,extends(qcd_2_2_type)::muli_t
!     private
     ! model parameters
     real(default)::GeV2_scale_cutoff     
     logical::modify_pdfs=muli_default_modify_pdfs
     ! pt chain status
     logical::finished=.false.
     logical::exceeded=.false.
     ! timers
     real(default) :: init_time = zero
     real(default) :: pt_time = zero
     real(default) :: partons_time = zero
     real(default) :: confirm_time = zero
     ! generator internals
     logical::initialized=.false.
     logical::initial_interaction_given=.false.
     real(default) :: mean = one
     real(default), dimension(0:16) :: start_integrals = &
          [0,0,0,0,0,0,0,0,0,0,0,0,0,0,0,0,0]
     type(tao_random_state)::tao_rnd
     type(muli_trapezium_tree_t)::dsigma
     type(sample_inclusive_t)::samples
     type(pp_remnant_t)::beam
     ! these pointers shall not be allocated, deallocated, serialized or deserialized explicitly.
     class(muli_trapezium_node_class_t),pointer::node=>null()
   contains
     ! overridden ser_class_t procedures
     procedure,public::write_to_marker=>muli_write_to_marker
     procedure,public::read_from_marker=>muli_read_from_marker
     procedure,public::print_to_unit=>muli_print_to_unit
     procedure,public,nopass::get_type=>muli_get_type
     ! init / final
     procedure,public::muli_initialize
     procedure,public::apply_initial_interaction=>muli_apply_initial_interaction
     procedure,public::finalize=>muli_finalize
     procedure,public::stop_trainer=>muli_stop_trainer
     procedure,public::reset_timer=>muli_reset_timer
     procedure,public::restart=>muli_restart
     generic,public::  initialize=>muli_initialize
     ! status query
     procedure,public::is_initialized=>muli_is_initialized
     procedure,public::is_initial_interaction_given=>muli_is_initial_interaction_given
     procedure,public::is_finished=>muli_is_finished
     ! user interface
     procedure,public::enable_remnant_pdf=>muli_enable_remnant_pdf
     procedure,public::disable_remnant_pdf=>muli_disable_remnant_pdf
     procedure,public::generate_gev2_pt2=>muli_generate_gev2_pt2
     procedure,public::generate_partons=>muli_generate_partons
     procedure,public::generate_flow=>muli_generate_flow
     procedure,public::replace_parton=>muli_replace_parton
     procedure,public::get_parton_pdf=>muli_get_parton_pdf
     procedure,public::get_momentum_pdf=>muli_get_momentum_pdf
     procedure,public::print_timer=>muli_print_timer
     procedure,public::generate_samples=>muli_generate_samples
     ! beam test
     procedure,public::fake_interaction=>muli_fake_interaction
     ! private procedures
     procedure,private::generate_next_scale=>muli_generate_next_scale
     procedure,private::confirm=>muli_confirm
  end type muli_t

contains

!!!!!!!!!!!!!!!!!!!!!!!!!!!!!!!!!!!!!!!!!!!!!!
!!! type bound procedures for qcd_2_2_type !!!
!!!!!!!!!!!!!!!!!!!!!!!!!!!!!!!!!!!!!!!!!!!!!!
  
  !overridden ser_class_t procedures

  subroutine qcd_2_2_write_to_marker(this,marker,status)
    class(qcd_2_2_type), intent(in)::this
    class(marker_t), intent(inout)::marker
    integer(dik), intent(out)::status
    call marker%mark_begin("qcd_2_2_type")
    call this%mom_write_to_marker (marker, status)
    call marker%mark("process_id",this%process_id)
    call marker%mark("integrand_id",this%integrand_id)
    call marker%mark("momentum_fractions",this%momentum_fractions)
    call marker%mark("hyperbolic_fractions",this%hyperbolic_fractions)
    call marker%mark_end("qcd_2_2_type")
  end subroutine qcd_2_2_write_to_marker

  subroutine qcd_2_2_read_from_marker(this,marker,status)
    class(qcd_2_2_type), intent(out)::this
    class(marker_t), intent(inout)::marker
    integer(dik), intent(out)::status
    call marker%pick_begin("qcd_2_2_type",status=status)
    call this%mom_read_from_marker (marker, status)
    call marker%pick("process_id",this%process_id,status)
    call marker%pick("integrand_id",this%integrand_id,status)
    call marker%pick("momentum_fractions",this%momentum_fractions,status)
    call marker%pick("hyperbolic_fractions",this%hyperbolic_fractions,status)
    call marker%pick_end("qcd_2_2_type",status=status)
  end subroutine qcd_2_2_read_from_marker

  subroutine qcd_2_2_print_to_unit(this,unit,parents,components,peers)
    class(qcd_2_2_type), intent(in)::this
    integer, intent(in)::unit
    integer(dik), intent(in)::parents,components,peers
    integer, dimension(2,4)::flow
    integer::index
    if (parents > i_zero) & 
         call this%mom_print_to_unit (unit, parents-1, components, peers)
    write(unit,'("Components of qcd_2_2_type:")')
    write(unit,'("Process id is:       ",I3)')this%get_process_id()
    write(unit,'("Integrand id is:     ",I3)')this%get_integrand_id()
    if (this%get_integrand_id()>0) then
       write(unit,'("LHA Flavors are:     ",4(I3))')this%get_lha_flavors()
       write(unit,'("PDG Flavors are:     ",4(I3))')this%get_pdg_flavors()
       write(unit,'("Parton kinds are:    ",2(I3))')this%get_parton_kinds()
       write(unit,'("PDF int kinds are:   ",2(I3))')this%get_pdf_int_kinds()
       write(unit,'("Diagram kind is:     ",2(I3))')this%get_diagram_kind()
    end if
    call this%get_color_correlations(1,index,flow)
    write(unit,'("Color Permutations:  ",4(I0))')this%flow
    write(unit,'("Color Connections:")')
    write(unit,'("(",I0,",",I0,")+(",I0,",",I0,")->(",I0,",",I0,")+(",I0,",",I0,")")')flow
    write(unit,'("Evolution scale is:  ",E14.7)')this%get_unit2_scale()
    write(unit,'("Momentum boost is:   ",E14.7)')this%get_momentum_boost()
    write(unit,'("Remant momentum fractions are: ",2(E14.7))')this%get_remnant_momentum_fractions()
    write(unit,'("Total momentum fractions are:  ",2(E14.7))')this%get_total_momentum_fractions()
  end subroutine qcd_2_2_print_to_unit
  
  pure subroutine qcd_2_2_get_type(type)
    character(:),allocatable, intent(out)::type
    allocate(type,source="qcd_2_2_type")
  end subroutine qcd_2_2_get_type

  elemental function qcd_2_2_get_process_id(this) result(id)
    class(qcd_2_2_type), intent(in)::this
    integer::id
    id=this%process_id
  end function qcd_2_2_get_process_id

  elemental function qcd_2_2_get_integrand_id(this) result(id)
    class(qcd_2_2_type), intent(in)::this
    integer::id
    id=this%integrand_id
  end function qcd_2_2_get_integrand_id

  pure function qcd_2_2_get_lha_flavors(this) result(lha)
    class(qcd_2_2_type), intent(in)::this
    integer, dimension(4)::lha
    lha = valid_processes(1:4,this%process_id)
  end function qcd_2_2_get_lha_flavors

  pure function qcd_2_2_get_pdg_flavors(this) result(pdg)
    class(qcd_2_2_type), intent(in)::this
    integer, dimension(4)::pdg
    pdg=this%get_lha_flavors()
    where(pdg==0) pdg=21
  end function qcd_2_2_get_pdg_flavors

  pure function qcd_2_2_get_pdf_int_kinds(this) result(kinds)
    class(qcd_2_2_type), intent(in)::this
    integer, dimension(2)::kinds
    kinds=double_pdf_kinds(1:2,this%integrand_id)
  end function qcd_2_2_get_pdf_int_kinds

  pure function qcd_2_2_get_parton_id(this,n) result(id)
    class(qcd_2_2_type), intent(in)::this
    integer, intent(in)::n
    integer::id
    id=this%parton_ids(n)
  end function qcd_2_2_get_parton_id

  pure function qcd_2_2_get_parton_kinds(this) result(kinds)
    class(qcd_2_2_type), intent(in)::this
    integer, dimension(2)::kinds
    kinds=this%get_pdf_int_kinds()
    kinds(1)=parton_kind_of_int_kind(kinds(1))
    kinds(2)=parton_kind_of_int_kind(kinds(2))
  end function qcd_2_2_get_parton_kinds

  elemental function qcd_2_2_get_io_kind(this) result(kind)
    class(qcd_2_2_type), intent(in)::this
    integer::kind
    kind=valid_processes(5,this%process_id)
  end function qcd_2_2_get_io_kind

  elemental function qcd_2_2_get_diagram_kind(this) result(kind)
    class(qcd_2_2_type), intent(in)::this
    integer::kind
    kind=valid_processes(6,this%process_id)
  end function qcd_2_2_get_diagram_kind

  pure function qcd_2_2_get_diagram_color_kind(this) result(kind)
    ! This is one more hack. Before merging into the interleaved algorithm, muli has only cared for 
    ! summed cross sections, but not in specific color flows. So two different diagrams with equal 
    ! cross sections were summed up to diagram kind 1. 
    ! Now muli also generates color flows, so we must devide diagram kind 1 into diagram color kind 0
    ! and diagram color kind 1.
    class(qcd_2_2_type), intent(in)::this
    integer::kind
    kind=valid_processes(6,this%process_id)
    if (kind==1) then
       if (product(valid_processes(1:2,this%process_id))>0) then
          kind=0
       end if
    end if
  end function qcd_2_2_get_diagram_color_kind

  elemental function qcd_2_2_get_momentum_boost (this) result (boost)
    class(qcd_2_2_type), intent(in) :: this
    real(default) :: boost
    boost = -one
    ! print('("qcd_2_2_get_momentum_boost: not yet implemented.")')
    ! boost = this%momentum_boost
  end function qcd_2_2_get_momentum_boost

  pure function qcd_2_2_get_hyperbolic_fractions(this) result(fractions)
    class(qcd_2_2_type), intent(in)::this
    real(double), dimension(3)::fractions
    fractions=this%hyperbolic_fractions
  end function qcd_2_2_get_hyperbolic_fractions

  pure function qcd_2_2_get_remnant_momentum_fractions (this) result (fractions)
    class(qcd_2_2_type), intent(in) :: this
    real(default), dimension(2) :: fractions
    fractions = this%momentum_fractions(1:2)
  end function qcd_2_2_get_remnant_momentum_fractions

  pure function qcd_2_2_get_total_momentum_fractions (this) result (fractions)
    class(qcd_2_2_type), intent(in) :: this
    real(default), dimension(2) :: fractions
    fractions = [-one, -one]
    ! fractions = this%hyperbolic_fractions(1:2) * &
    !    this%beam%get_proton_remnant_momentum_fractions()
  end function qcd_2_2_get_total_momentum_fractions

  pure function qcd_2_2_get_color_flow(this) result(flow)
    class(qcd_2_2_type), intent(in)::this
    integer, dimension(4)::flow
    flow=this%flow
  end function qcd_2_2_get_color_flow

  subroutine qcd_2_2_get_color_correlations(this,start_index,final_index,flow)
    class(qcd_2_2_type), intent(in)::this
    integer, intent(in)::start_index
    integer, intent(out)::final_index
    integer, dimension(2,4), intent(out)::flow
    integer::pos,f_end,f_beginning
    final_index=start_index
    ! we set all flows to i_zero. i_zero means no connection.
    flow=reshape([0,0,0,0,0,0,0,0],[2,4])
    ! look at all four possible ends of color lines
    do f_end=1,4
       ! the beginning of of this potential line is stored in flow. i_zero means no line.
       f_beginning=this%flow(f_end)
       ! is there a line beginning at f_beginning and ending at f_end?
       if (f_beginning>0) then
          ! yes it is. we get a new number for this new line
          final_index=final_index+1
          ! is this line beginning in the initial state?
          if (f_beginning<3) then
             ! yes it is. lets connect the color entry of f_begin.            
             flow(1,f_beginning)=final_index
          else
             ! no, it's the final state. lets connect the anticolor entry of f_begin.
             flow(2,f_beginning)=final_index
          end if
          ! is this line ending in the final state?
          if (f_end>2) then
             ! yes it is. lets connect the color entry of f_end.
             flow(1,f_end)=final_index
          else
             ! no, it's the initial state. lets connect the anticolor entry of f_end.
             flow(2,f_end)=final_index
          end if
       end if
    end do
  end subroutine qcd_2_2_get_color_correlations

  subroutine qcd_2_2_initialize (this, gev2_s, process_id, &
       integrand_id, parton_ids, flow,hyp, cart)
    class(qcd_2_2_type), intent(out) :: this
    real(default), intent(in) :: gev2_s
    integer, intent(in) :: process_id, integrand_id
    integer, dimension(2), intent(in) :: parton_ids
    integer, dimension(4), intent(in) :: flow
    real(default), dimension(3), intent(in)::hyp
    real(default), dimension(3), intent(in), optional :: cart
    call this%initialize (gev2_s)
    this%process_id = process_id
    this%integrand_id = integrand_id
    this%parton_ids = parton_ids
    this%flow = flow
    this%hyperbolic_fractions = hyp
    if (present(cart)) then 
       this%momentum_fractions = cart
    else
       this%momentum_fractions = h_to_c_param (hyp)
    end if
  end subroutine qcd_2_2_initialize

  !!!!!!!!!!!!!!!!!!!!!!!!!!!!!!!!!!!!!!!!!!!
  !!! type bound procedures for muli_t !!!
  !!!!!!!!!!!!!!!!!!!!!!!!!!!!!!!!!!!!!!!!!!!

  ! overridden ser_class_t procedures

  subroutine muli_write_to_marker(this,marker,status)
    class(muli_t), intent(in) :: this
    class(marker_t), intent(inout)::marker
    integer(dik), intent(out)::status
    call marker%mark_begin("muli_t")
    call qcd_2_2_write_to_marker (this,marker,status)
    call marker%mark("modify_pdfs",this%modify_pdfs)
    call marker%mark("initialized",this%initialized)
    call marker%mark("initial_interaction_given",this%initial_interaction_given)
    call marker%mark("finished",this%finished)
    call marker%mark("init_time",this%init_time)
    call marker%mark("pt_time",this%pt_time)
    call marker%mark("partons_time",this%partons_time)
    call marker%mark("confirm_time",this%confirm_time)
!    call marker%mark_instance(this%start_values,"start_values")
    call marker%mark_instance(this%dsigma,"dsigma")
    call marker%mark_instance(this%samples,"samples")
    call marker%mark_instance(this%beam,"beam")
    call marker%mark_end("muli_t")
  end subroutine muli_write_to_marker
  
  subroutine muli_read_from_marker(this,marker,status)
    class(muli_t), intent(out) :: this
    class(marker_t), intent(inout)::marker
    integer(dik), intent(out)::status
    call marker%pick_begin("muli_t",status=status)
    call qcd_2_2_read_from_marker(this,marker,status)
    call marker%pick("modify_pdfs",this%modify_pdfs,status)
    call marker%pick("initialized",this%initialized,status)
    call marker%pick("initial_interaction_given",this%initial_interaction_given,status)
    call marker%pick("finished",this%finished,status)
    call marker%pick("init_time",this%init_time,status)
    call marker%pick("pt_time",this%pt_time,status)
    call marker%pick("partons_time",this%partons_time,status)
    call marker%pick("confirm_time",this%confirm_time,status)
!    call marker%pick_instance("start_values",this%start_values,status=status)
    call marker%pick_instance("dsigma",this%dsigma,status=status)
    call marker%pick_instance("samples",this%samples,status=status)
    call marker%pick_instance("beam",this%beam,status=status)
    call marker%pick_end("muli_t",status)
  end subroutine muli_read_from_marker

  subroutine muli_print_to_unit(this,unit,parents,components,peers)
    class(muli_t), intent(in)::this
    integer, intent(in)::unit
    integer(dik), intent(in)::parents,components,peers
    if (parents>0)call qcd_2_2_print_to_unit(this,unit,parents-1,components,peers)
    write(unit,fmt="(a)")"Components of muli_t :"
    write(unit,'("Model Parameters:")')
    write(unit,'("GeV2_scale_cutoff : ",E20.10)')this%GeV2_scale_cutoff
    write(unit,'("Modify PDF        : ",L1)')this%modify_pdfs
    write(unit,'("PT Chain Status:")')
    write(unit,'("Initialized       : ",L1)')this%initialized
    write(unit,'("initial_interaction_given: ",L1)')this%initial_interaction_given
    write(unit,'("Finished          : ",L1)')this%finished
    write(unit,'("Exceeded          : ",L1)')this%exceeded
    write(unit,'("Generator Internals:")')
    write(unit,'("Mean Value        : ",E20.10)')this%mean
    if (components>i_zero) then
       write(unit,'("Start Integrals   : ",16(E20.10))')this%start_integrals(1:16)
!       write(unit,'("start_values Component:")')
!       call this%start_values%print_to_unit(unit,parents,components-1,peers)
       write(unit,'("dsigma Component:")')
       call this%dsigma%print_to_unit(unit,parents,components-1,peers)
       write(unit,'("samples Component:")')
       call this%samples%print_to_unit(unit,parents,components-1,peers)
       write(unit,'("beam Component:")')
       call this%beam%print_to_unit(unit,parents,components-1,peers)
    else
       write(unit,'("Skipping Derived-Type Components.")')
    end if
!    call print_comp_pointer(this%start_node,unit,i_zero,min(components-1,i_one),i_zero,"start_node")
!    call serialize_print_comp_pointer(this%node,unit,i_zero,min(components-1,i_one),i_zero,"node")
  end subroutine muli_print_to_unit

  pure subroutine muli_get_type(type)
    character(:),allocatable, intent(out)::type
    allocate(type,source="muli_t")
  end subroutine muli_get_type

  ! user interface

  subroutine muli_apply_initial_interaction (this, GeV2_s, &
       x1, x2, pdg_f1, pdg_f2, n1, n2)
    class(muli_t), intent(inout) :: this
    real(default), intent(in) :: Gev2_s, x1, x2
    integer, intent(in):: pdg_f1, pdg_f2, n1, n2
    real(default) :: rnd1, rnd2, time
    if (this%initialized) then
       call cpu_time(time)
       this%init_time=this%init_time-time
       print *,"muli_apply_initial_interaction:"
       print *,"gev2_s=",gev2_s
       print *,"x1=",x1
       print *,"x2=",x2
       print *,"pdg_f1=",pdg_f1
       print *,"pdg_f2=",pdg_f2
       print *,"n1=",n1
       print *,"n2=",n2
       call tao_random_number (this%tao_rnd, rnd1)
       call tao_random_number (this%tao_rnd, rnd2)
       call cpu_time(time)
       this%init_time=this%init_time+time
       call this%beam%apply_initial_interaction &
                 (sqrt (gev2_s), x1, x2, pdg_f1, pdg_f2, n1, n2,&
            ! This is a hack: We should give the pt scale of the initial interaction.
            ! Unfortunately, we only know the invarient mass shat. shat/2 is the upper
            ! bound of pt, so we use it for now.
            sqrt(gev2_s)*x1*x2/2D0,&
            rnd1,rnd2)
       this%initial_interaction_given=.true.
    else
       print *,"muli_apply_initial_interaction: call muli_initialize first. STOP"
       STOP
    end if
  end subroutine muli_apply_initial_interaction

  subroutine muli_initialize(this,&
       GeV2_scale_cutoff,&
       gev2_s,&              
       muli_dir,&
       random_seed)
    class(muli_t), intent(out)::this
    real(kind=default), intent(in)::gev2_s,GeV2_scale_cutoff
    character(*), intent(in)::muli_dir
    integer, intent(in),optional::random_seed
    real(double)::time
    logical::exist
    type(muli_dsigma_t)::dsigma_aq
    character(3)::lhapdf_member_c
    call cpu_time(time)
    this%init_time=this%init_time-time
    print *,"muli_initialize: The MULI modules are still not fully populated, so MULI might &
         &generate some dummy values instead of real Monte Carlo generated interactions."
    print *,"Given Parameters:"
    print *,"GeV2_scale_cutoff=",GeV2_scale_cutoff
    print *,"muli_dir=",muli_dir
    print *,"lhapdf_dir=",""
    print *,"lhapdf_file=",muli_default_lhapdf_file
    print *,"lhapdf_member=",muli_default_lhapdf_member
    print *,""
    call this%transverse_mom_t%initialize (gev2_s)
    call this%beam%initialize(&
         muli_dir,&
         lhapdf_dir="",&
         lhapdf_file=muli_default_lhapdf_file,&
         lhapdf_member=muli_default_lhapdf_member)
    this%GeV2_scale_cutoff=GeV2_scale_cutoff
    if (present(random_seed)) then
       call tao_random_create(this%tao_rnd,random_seed)
    else
       call tao_random_create(this%tao_rnd,1)
    end if
    print *,"looking for previously generated root function..."
    call integer_with_leading_zeros(muli_default_lhapdf_member,3,lhapdf_member_c)
    inquire(file=muli_dir//"/dsigma_"//muli_default_lhapdf_file//".xml",exist=exist)
    if (exist) then
       print *,"found. Starting deserialization..."
       call this%dsigma%deserialize(&
            name="dsigma_"//muli_default_lhapdf_file//"_"//lhapdf_member_c,&
            file=muli_dir//"/dsigma_"//muli_default_lhapdf_file//".xml")
!       call this%dsigma%print_all()
       print *,"done. Starting generation of plots..."
       call this%dsigma%gnuplot(muli_dir)
       print *,"done."
    else
       print *,"No root function found. Starting generation of root function..."
       call dsigma_aq%generate(GeV2_scale_cutoff,gev2_s,this%dsigma)
       print *,"done. Starting serialization of root function..."
       call this%dsigma%serialize(& 
            name="dsigma_"//muli_default_lhapdf_file//"_"//lhapdf_member_c,&
            file=muli_dir//"/dsigma_"//muli_default_lhapdf_file//".xml")
       print *,"done. Starting serialization of generator..."
       call dsigma_aq%serialize(&
            name="dsigma_aq_"//muli_default_lhapdf_file//"_"//lhapdf_member_c,&
            file=muli_dir//"/dsigma_aq_"//muli_default_lhapdf_file//".xml")
       print *,"done. Starting generation of plots..."
       call this%dsigma%gnuplot(muli_dir)
       print *,"done."
    end if
    print *,""
    print *,"looking for previously generated samples..."
    inquire(file=muli_dir//"/samples.xml",exist=exist)
    if (exist) then
       print *,"found. Starting deserialization..."
       call this%samples%deserialize("samples",muli_dir//"/samples.xml")
    else
       print *,"No samples found. Starting with default initialization."
       call this%samples%initialize(4, int_sizes_all, int_all, 1E-2_default)
    end if
    call this%restart()
    this%initialized=.true.
    call cpu_time(time)
    this%init_time=this%init_time+time
  end subroutine muli_initialize

  subroutine muli_finalize(this)
    class(muli_t), intent(inout)::this
    print *,"muli_finalize"
    nullify(this%node)
    call this%dsigma%finalize()
    call this%samples%finalize()
    call this%beam%finalize()
  end subroutine muli_finalize

  subroutine muli_stop_trainer(this)
    class(muli_t), intent(inout)::this
    print *,"muli_stop_trainer: DUMMY!"
  end subroutine muli_stop_trainer

  subroutine muli_reset_timer(this)
    class(muli_t), intent(inout)::this
    this%init_time=0D0
    this%pt_time=0D0
    this%partons_time=0D0
    this%confirm_time=0D0
  end subroutine muli_reset_timer

  subroutine muli_restart(this)
    class(muli_t), intent(inout)::this
    call this%dsigma%get_rightmost(this%node)
    call this%beam%reset()
!    print *,associated(this%node)
!    nullify(this%node)
    this%finished=.false.
    this%process_id=-1
    this%integrand_id=-1
    this%momentum_fractions=[-1D0,-1D0,1D0]
    this%hyperbolic_fractions=[-1D0,-1D0,1D0]
 !   this%start_values%process_id=-1
 !   this%start_values%integrand_id=-1
!    this%start_values%momentum_fractions=[-1D0,-1D0,1D0]
!    this%start_values%hyperbolic_fractions=[-1D0,-1D0,1D0]
    this%start_integrals=[0D0,0D0,0D0,0D0,0D0,0D0,0D0,0D0,0D0,0D0,0D0,0D0,0D0,0D0,0D0,0D0,0D0]
  end subroutine muli_restart

  elemental function muli_is_initialized(this) result(res)
    logical::res
    class(muli_t), intent(in) :: this
    res=this%initialized
  end function muli_is_initialized
  
  elemental function muli_is_initial_interaction_given(this) result(res)
    logical::res
    class(muli_t), intent(in) :: this
    res=this%initial_interaction_given
  end function muli_is_initial_interaction_given

  elemental function muli_is_finished(this) result(res)
    logical::res
    class(muli_t), intent(in) :: this
    res=this%finished
  end function muli_is_finished

  subroutine muli_enable_remnant_pdf(this)
    class(muli_t), intent(inout)::this
    this%modify_pdfs=.true.
  end subroutine muli_enable_remnant_pdf

  subroutine muli_disable_remnant_pdf(this)
    class(muli_t), intent(inout)::this
    this%modify_pdfs=.false.
  end subroutine muli_disable_remnant_pdf

  subroutine muli_generate_gev2_pt2 (this, gev2_start_scale, gev2_new_scale)
    class(muli_t), intent(inout) :: this
    real(kind=default), intent(in) :: gev2_start_scale
    real(kind=default), intent(out) :: gev2_new_scale
    real(double) :: time
    call cpu_time (time)
    this%pt_time = this%pt_time - time
    call this%set_gev2_scale (gev2_start_scale)
    this%start_integrals=this%node%approx_integral(this%get_unit_scale())
    call this%generate_next_scale()
    gev2_new_scale=this%get_gev2_scale()
    call cpu_time(time)
    this%pt_time=this%pt_time+time
  end subroutine muli_generate_gev2_pt2
  
  subroutine muli_generate_flow(this)
    class(muli_t), intent(inout)::this
    integer::rnd
    integer::m,n
    logical, dimension(3)::t
    integer, dimension(4)::tmp_flow, tmp_array
    ! we initialize with zeros. a i_zero means no line ends here.
    this%flow=[0,0,0,0]
    ! we randomly pick a color flow
    call tao_random_number(this%tao_rnd,rnd)
    ! the third position of muli_flow_stats is the sum of all flow wheights of stratum diagram_kind.
    ! so we generate a random number 0 <= m < sum(weights)
    m=modulo(rnd,muli_flow_stats(3,this%get_diagram_color_kind()))
    ! lets visit all color flows of stratum diagram_kind. the first and second position of muli_flow_stats
    ! tells us the index of the first and the last valid color flow.
    do n=muli_flow_stats(1,this%get_diagram_color_kind()),muli_flow_stats(2,this%get_diagram_color_kind())
       ! now we remove the weight of flow n from our random number.
       m=m-muli_flows(0,n)
       ! this is how we pick a flow.
       if (m<0) then
          ! the actual flow
          this%flow=muli_flows(1:4,n)
          exit
       end if
    end do
    ! the diagram kind contains a primitive diagram and all diagramms which can be deriven by
    ! (1) global charge conjugation
    ! (2) permutation of the initial state particles
    ! (3) permutation of the final state particles
    ! lets see, what transformations we have got in our actual interaction.
    tmp_array = this%get_lha_flavors ()
    t = muli_get_state_transformations (this%get_diagram_color_kind (), &
         tmp_array)
    !     this%get_lha_flavors ())    
    ! now we have to apply these transformations to our flow.
    ! (1) means: swap beginning and end of a line. flow is a permutation that maps
    ! ends to their beginnings, so we apply flow to itself:
!!$    print *,"(0)",this%flow
    if (t(1)) then
       tmp_flow=this%flow
       this%flow=[0,0,0,0]
       do n=1,4
          if (tmp_flow(n)>0)this%flow(tmp_flow(n))=n
       end do
!!$       print *,"(1)",this%flow
    end if
    if (t(2)) then
       ! we swap the particles 1 and 2
       tmp_flow(1)=this%flow(2)
       tmp_flow(2)=this%flow(1)
       tmp_flow(3:4)=this%flow(3:4)
!!$       print *,"(2)",tmp_flow
       ! we swap the beginnings assigned to particle 1 and 2
       where(tmp_flow==1)
          this%flow=2
       elsewhere(tmp_flow==2)
          this%flow=1
       elsewhere
          this%flow=tmp_flow
       end where
!!$       print *,"(2)",this%flow
    end if
    if (t(3)) then
       ! we swap the particles 3 and 4
       tmp_flow(1:2)=this%flow(1:2)
       tmp_flow(3)=this%flow(4)
       tmp_flow(4)=this%flow(3)
!!$       print *,"(3)",tmp_flow
       ! we swap the beginnings assigned to particle 3 and 4
       where(tmp_flow==3)
          this%flow=4
       elsewhere(tmp_flow==4)
          this%flow=3
       elsewhere
          this%flow=tmp_flow
       end where
!!$       print *,"(3)",this%flow
    end if
  end subroutine muli_generate_flow

  subroutine muli_generate_partons (this, n1, n2, x_proton_1, x_proton_2, &
       pdg_f1, pdg_f2, pdg_f3, pdg_f4)
    class(muli_t), intent(inout) :: this
    integer, intent(in) :: n1, n2
    real(kind=default), intent(out) :: x_proton_1, x_proton_2
    integer, intent(out) :: pdg_f1, pdg_f2, pdg_f3, pdg_f4
    integer, dimension(4) :: pdg_f
    real(double) :: time
    !print *,"muli_generate_partons: n1=",n1," n2=",n2
    this%parton_ids(1) = n1
    this%parton_ids(2) = n2
    call cpu_time(time)
    this%partons_time=this%partons_time-time
    this%mean = this%node%approx_value_n (this%get_unit_scale(), &
         this%integrand_id)
    call this%samples%mcgenerate_hit(&
         this%get_unit2_scale(),&
         this%mean,&
         this%integrand_id,&
         this%tao_rnd,&
         this%process_id,&
         this%momentum_fractions)
    !print *,"muli_generate_partons",this%momentum_fractions
    call this%generate_flow()
    if (this%modify_pdfs) then
       call cpu_time(time)
       this%partons_time=this%partons_time+time
       this%confirm_time=this%confirm_time-time
       call this%beam%apply_interaction(this)
       call cpu_time(time)
       this%confirm_time=this%confirm_time+time
       this%partons_time=this%partons_time-time
    end if
    x_proton_1=this%momentum_fractions(1)
    x_proton_2=this%momentum_fractions(2)
    pdg_f=this%get_pdg_flavors()
    pdg_f1=pdg_f(1)
    pdg_f2=pdg_f(2)
    pdg_f3=pdg_f(3)
    pdg_f4=pdg_f(4)
    call cpu_time(time)
    this%partons_time=this%partons_time-time
    call qcd_2_2_print_to_unit(this,output_unit,100_dik,100_dik,100_dik)
  end subroutine muli_generate_partons

  subroutine muli_replace_parton(this,proton_id,old_id,new_id,pdg_f,x_proton,gev_scale)
    class(muli_t), intent(inout)::this
    integer, intent(in)::proton_id,old_id,new_id,pdg_f
    real(kind=default), intent(in)::x_proton,gev_scale
    !print *,"muli_replace_parton(",proton_id,old_id,new_id,pdg_f,x_proton,gev_scale,")"
    if (proton_id==1.or.proton_id==2) then
       call this%beam%replace_parton(proton_id,old_id,new_id,pdg_f,x_proton,gev_scale)
    else
       print *,"muli_replace_parton: proton_id must be 1 or 2, but ",proton_id," was given."
       STOP
    end if
  end subroutine muli_replace_parton

  function muli_get_momentum_pdf &
       (this, x_proton, gev2_scale, n, pdg_f) result (pdf)
    real(default) :: pdf
    class(muli_t), intent(in) :: this
    real(default), intent(in) :: x_proton, gev2_scale
    integer, intent(in) :: n, pdg_f
    call this%beam%momentum_pdf (x_proton, gev2_scale, n, pdg_f, pdf)
  end function muli_get_momentum_pdf

  function muli_get_parton_pdf &
       (this, x_proton, gev2_scale, n, pdg_f) result(pdf)
    real(default) :: pdf
    class(muli_t), intent(in) :: this
    real(default), intent(in) :: x_proton, gev2_scale
    integer, intent(in) :: n, pdg_f
    call this%beam%parton_pdf (x_proton, gev2_scale, n, pdg_f, pdf)
  end function muli_get_parton_pdf

  subroutine muli_print_timer(this)
    class(muli_t), intent(in) :: this
    print('("Init time:    ",E20.10)'),this%init_time
    print('("PT gen time:  ",E20.10)'),this%pt_time
    print('("Partons time: ",E20.10)'),this%partons_time
    print('("Confirm time: ",E20.10)'),this%confirm_time
    print('("Overall time: ",E20.10)'),this%init_time+this%pt_time+this%partons_time+this%confirm_time
  end subroutine muli_print_timer

  subroutine muli_generate_next_scale (this, integrand_kind)
    class(muli_t), intent(inout) :: this
    integer, intent(in), optional :: integrand_kind
    real(default) :: pts,tmp_pts,rnd
    integer :: tmp_int_kind
    class(muli_trapezium_node_class_t), pointer :: tmp_node
      pts = -one
      if (present(integrand_kind)) then
         call tao_random_number (this%tao_rnd, rnd)
         call generate_single_pts(&
              integrand_kind,&
              this%start_integrals(integrand_kind),&
              this%beam%get_pdf_int_weights(double_pdf_kinds(1:2,integrand_kind)),&
              rnd,&
              this%dsigma,&
              pts,&
              this%node)
      else
         do tmp_int_kind = 1, 16
            call tao_random_number (this%tao_rnd, rnd)
            call generate_single_pts(&
                 tmp_int_kind,&
                 this%start_integrals(tmp_int_kind),&
                 this%beam%get_pdf_int_weights(double_pdf_kinds(1:2,tmp_int_kind)),&
                 rnd,&
                 this%dsigma,&
                 tmp_pts,&
                 tmp_node)
            if (tmp_pts>pts) then
               pts=tmp_pts
               this%integrand_id=tmp_int_kind
               this%node=>tmp_node
            end if
         end do
      end if
      if (pts>0) then
         call this%set_unit_scale(pts)
      else
         this%finished=.true.
      end if
!      print *,this%finished,this%integrand_id
  contains
    subroutine generate_single_pts &
         (int_kind, start_int, weight, rnd, int_tree, pts, node)
      integer, intent(in) :: int_kind
      real(default), intent(in) :: start_int, weight, rnd
      type(muli_trapezium_tree_t), intent(in) :: int_tree
      real(default), intent(out) :: pts
      class(muli_trapezium_node_class_t),pointer, intent(out) :: node
      real(default) :: arg
      ! print *,int_kind,start_int,weight,rnd
      if (weight>0D0) then
         arg=start_int-log(rnd)/weight
         call int_tree%find_decreasing (arg, int_kind, node)
         if (node%get_l_integral(int_kind)>arg) then
            pts=node%approx_position_by_integral(int_kind,arg)
         else
            pts=-1D0
         end if
      else
         pts=-1D0
      end if
    end subroutine generate_single_pts
  end subroutine muli_generate_next_scale

  subroutine muli_confirm (this)
    class(muli_t), intent(inout) :: this
    this%mean=this%node%approx_value_n(this%get_unit_scale(),this%integrand_id)
    this%start_integrals=this%node%approx_integral(this%get_unit_scale())
  end subroutine muli_confirm

  subroutine muli_generate_samples(this,n_total,n_print,integrand_kind,muli_dir,analyse)
    class(muli_t), intent(inout)::this
    integer(dik), intent(in)::n_total,n_print
    integer, intent(in)::integrand_kind
    character(*), intent(in)::muli_dir
    logical, intent(in)::analyse
    integer(dik)::n_inner

    class(muli_trapezium_node_class_t),pointer::start_node=>null()
    class(muli_trapezium_node_class_t),pointer,save::s_node=>null()
    class(muli_trapezium_node_class_t),pointer,save::node=>null()   

    character(2)::prefix
    integer,save::t_slice,t_region,t_proc,t_subproc,t_max_n=0
    integer(dik)::n_t,n_p,n_m
    integer::n,m,u,unit=0
    integer(dik)::n_tries=0
    integer(dik)::n_hits=0
    integer(dik)::n_over=0
    integer(dik)::n_miss=0
    real(default),save, dimension(3)::cart_hit
    integer,save, dimension(4)::t_i_rnd
    !  integer,save, dimension(5)::r_n_proc
    real(default), dimension(16)::d_rnd
    real(default),save::t_area,t_dddsigma,t_rnd,t_weight,t_arg
    real(default)::mean=0D0
    real(default)::time=0D0
    real(default)::timepa=0D0
    real(default)::timept=0D0
    real(default)::timet=0D0
    real(default)::pts,s_pts=1D0
    real(default)::pts2=1D0
    real(default)::rnd
    logical::running
    character(3)::num
    integer::success=-1
    integer::chain_length=0
    integer::int_kind
    integer::process_id
    real(double), dimension(0:16)::integral
    call this%print_parents()
    n_tries=one
    n_inner=n_total/n_print
    n_t=i_zero
    print:do while(n_t<n_total)
       call cpu_time(time)
       timet=-time
       n_p=i_zero
       inner:do while(n_p<n_print)
          chain_length=0
!          print *,"new chain"
          call this%restart()
          this%integrand_id=integrand_kind
          call cpu_time(time)
          timept=timept-time
          call this%generate_next_scale(integrand_kind)
          call cpu_time(time)
          timept=timept+time
          chain:do while(.not.this%is_finished())
             chain_length=chain_length+1
             n_p=n_p+1
             call this%confirm()
             call cpu_time(time)
             timepa=timepa-time
             !             print *,this%get_unit2_scale()
             call this%samples%mcgenerate_hit(&
                  this%get_unit2_scale(),&
                  this%mean,&
                  this%integrand_id,&
                  this%tao_rnd,&
                  this%process_id,&
                  this%momentum_fractions)
             call cpu_time(time)
             timepa=timepa+time
             timept=timept-time
             call this%generate_next_scale(integrand_kind)
             call cpu_time(time)
             timept=timept+time
          end do chain
!          print *,"chain length=",chain_length
       end do inner
       n_t=n_t+n_p
       call this%samples%sum_up()
       call cpu_time(time)
       timet=timet+time
       print *,n_t,"/",n_total
       print *,"time: ",timet
       print *,"pt time: ",timept
       print *,"pa time: ",timepa       
       print *,this%samples%n_tries_sum,this%samples%n_hits_sum,this%samples%n_over_sum
       if (this%samples%n_hits_sum>0) then
          print *,(this%samples%n_hits_sum*10000)/this%samples%n_tries_sum,&
               (this%samples%n_over_sum*10000)/this%samples%n_hits_sum
       else
          print *,"no hits"
       end if
!       print('(7(I11," "),5(E14.7," "))')&
!            ,n_p,n_print,n_tries,n_hits,n_over,&
!            int((n_hits*1D3)/n_tries),int((n_over*1D6)/n_tries),&
!            n_hits/real(n_over),time1,time2,time3,&
!            this%samples%int_kinds(integrand_kind)%overall_boost
    end do print
    call integer_with_leading_zeros(integrand_kind,2,prefix)
    if (analyse) then
       call this%samples%int_kinds(integrand_kind)%analyse(muli_dir,prefix//"_")
       call this%samples%int_kinds(integrand_kind)%serialize(&
            "sample_int_kind_"//prefix,&
            muli_dir//"/sample_int_kind/"//prefix//".xml")
    end if
    call this%samples%int_kinds(integrand_kind)%serialize(&
         "sample_int_kind_"//prefix,&
         muli_dir//"/sample_int_kind/"//prefix//".xml")
  end subroutine muli_generate_samples

  subroutine muli_fake_interaction (this, GeV2_scale, x1, x2, &
       process_id, integrand_id, n1, n2, flow)
    class(muli_t), intent(inout) :: this
    real(default), intent(in) :: Gev2_scale, x1, x2
    integer, intent(in) :: process_id, integrand_id, n1, n2
    integer, dimension(4), intent(in), optional :: flow
    call this%set_gev2_scale (Gev2_scale)
    this%process_id = process_id
    this%integrand_id = integrand_id
    this%parton_ids = [n1, n2]
    if (present (flow)) then
       this%flow = flow
    else
       this%flow = [0,0,0,0]
    end if
    this%momentum_fractions = [x1, x2, this%get_unit2_scale()]
    call this%beam%apply_interaction (this)
    call this%beam%print_all ()
  end subroutine muli_fake_interaction

!!!!!!!!!!!!!!!!!!!!!!!!!!
!!  private procedures  !!
!!!!!!!!!!!!!!!!!!!!!!!!!!

end module muli

\nwendcode{}\nwbegindocs{2355}\nwdocspar
\nwenddocs{}\nwbegincode{2356}\moddef{\code{}muli{\char95}cross{\char95}sections.f90\edoc{}}\endmoddef
! This is a dummy for muli_cross_sections_module
module muli_cross_sections_module
end module muli_cross_sections_module

\nwendcode{}\nwbegindocs{2357}\nwdocspar
\nwenddocs{}\nwbegincode{2358}\moddef{\code{}muli{\char95}sampling.f90\edoc{}}\endmoddef
! This is a dummy for muli_sampling_module
module muli_sampling_module
end module muli_sampling_module

\nwendcode{}\nwbegindocs{2359}\nwdocspar
\nwenddocs{}\nwbegindocs{2360}\nwdocspar
\nwenddocs{}\nwfilename{../../../src/user/user.nw}\nwbegindocs{0}% -*- ess-noweb-default-code-mode: f90-mode; noweb-default-code-mode: f90-mode; -*- 
% WHIZARD code as NOWEB source: matrix elements and process libraries

\nwenddocs{}\nwbegindocs{1}\nwdocspar
%%%%%%%%%%%%%%%%%%%%%%%%%%%%%%%%%%%%%%%%%%%%%%%%%%%%%%%%%%%%%%%%%%%%%%%%
\chapter{User Plugin Support}

Here we collect interface code that enables the user to inject his own
code into the WHIZARD workflow.  The code uses data types defined elsewhere, and
is referenced in the {\tt{}eval{\char95}trees} module.

These are the modules:
\begin{description}
\item[user\_code\_interface]
  Generic support and specific additions.
\item[sf\_user]
  Handle user-defined structure functions.
\end{description}

\clearpage
%%%%%%%%%%%%%%%%%%%%%%%%%%%%%%%%%%%%%%%%%%%%%%%%%%%%%%%%%%%%%%%%%%%%%%%%
\section{User Code Interface}

\nwenddocs{}\nwbegincode{2}\moddef{\code{}user{\char95}code{\char95}interface.f90\edoc{}}\endmoddef
\LA{}File header\RA{}

module user_code_interface

  use iso_c_binding !NODEP!
\LA{}Use kinds\RA{}
\LA{}Use strings\RA{}
  use diagnostics
  use c_particles
  use os_interface

\LA{}Standard module head\RA{}

\LA{}User Code: public\RA{}

\LA{}User Code: variables\RA{}

\LA{}User Code: interfaces\RA{}

contains

\LA{}User Code: procedures\RA{}

end module user_code_interface
\eatline
\nwendcode{}\nwbegindocs{3}\nwdocspar
\subsection{User Code Management}
\nwenddocs{}\nwbegindocs{4}\nwdocspar
This data structure globally holds the filehandle of the user-code
library:
\nwenddocs{}\nwbegincode{5}\moddef{User Code: public}\endmoddef
  public :: has_user_lib
\nwendcode{}\nwbegincode{6}\moddef{User Code: variables}\endmoddef
  type(dlaccess_t), save :: user_lib_handle
  logical, save :: has_user_lib = .false.
  type(string_t), save :: user

\eatline
\nwendcode{}\nwbegindocs{7}\nwdocspar
Compile, link and load user code files.  Dlopen all user-provided
libraries, included the one just compiled (if any).

By default, we are looking for a library {\tt{}user.so/dylib}.  If this is not
present, try {\tt{}user.f90} and compile it.  This can be overridden.

In detail: First, compile all sources explicitly specified on the
command line.  Then collect all libraries specified on the command
line, including {\tt{}user.so} if it was generated.  If there is still no
code, check for an existing {\tt{}user.f90} and compile this.  Link
everything into a {\tt{}user.la} libtool library.  When 
done, dlopen all libraries that we have so far.
\nwenddocs{}\nwbegincode{8}\moddef{User Code: public}\plusendmoddef
  public :: user_code_init
\nwendcode{}\nwbegincode{9}\moddef{User Code: procedures}\endmoddef
  subroutine user_code_init (user_src, user_lib, user_target, rebuild, os_data)
    type(string_t), dimension(:), intent(in) :: user_src, user_lib
    type(string_t), intent(in) :: user_target
    logical, intent(in) :: rebuild
    type(os_data_t), intent(in) :: os_data
    type(string_t) :: user_src_file, user_obj_files, user_lib_file
    logical :: exist
    type(c_funptr) :: fptr
    integer :: i
    call msg_message ("Initializing user code")
    user = user_target;  if (user == "")  user = "user"
    user_obj_files = ""
    inquire (file = char (user) // ".la", exist = exist)
    if (rebuild .or. .not. exist) then
       do i = 1, size (user_src)
          user_src_file = user_src(i) // os_data%fc_src_ext
          inquire (file = char (user_src_file), exist = exist)
          if (exist) then
             call msg_message ("Found user-code source '" &
                  // char (user_src_file) // "'.")
             call compile_user_src (user_src_file, user_obj_files)
          else
             call msg_fatal ("User-code source '" // char (user_src_file) &
                  // "' not found")
          end if
       end do
       do i = 1, size (user_lib)
          user_lib_file = user_lib(i) // ".la"
          inquire (file = char (user_lib_file), exist = exist)
          if (exist) then
             call msg_message ("Found user-code library '" &
                  // char (user_lib_file) // "'.")
          else
             user_lib_file = user_lib(i) // os_data%shlib_ext
             inquire (file = char (user_lib_file), exist = exist)
             if (exist) then
                call msg_message ("Found user-code library '" &
                     // char (user_lib_file) // "'.")
             else
                call msg_fatal ("User-code library '" // char (user_lib(i)) &
                     // "' not found")
             end if
          end if
          user_obj_files = user_obj_files // " " // user_lib_file
       end do
       if (user_obj_files == "") then
          user_src_file = user // os_data%fc_src_ext
          inquire (file = char (user_src_file), exist = exist)
          if (exist) then
             call msg_message ("Found user-code source '" &
                  // char (user_src_file) // "'.")
             call compile_user_src (user_src_file, user_obj_files)
          else
             call msg_fatal ("User-code source '" // char (user_src_file) &
                  // "' not found")
          end if
       end if
       if (user_obj_files /= "") then
          call link_user (char (user), user_obj_files)
       end if
    end if
    call dlaccess_init &
         (user_lib_handle, var_str ("."), &
          user // os_data%shlib_ext, os_data)
    if (dlaccess_has_error (user_lib_handle)) then
       call msg_error (char (dlaccess_get_error (user_lib_handle)))
       call msg_fatal ("Loading user code library '" // char (user) &
            // ".la' failed")
    else
       call msg_message ("User code library '" // char (user) &
            // ".la' successfully loaded")
       has_user_lib = .true.
    end if
  contains
    subroutine compile_user_src (user_src_file, user_obj_files)
      type(string_t), intent(in) :: user_src_file
      type(string_t), intent(inout) :: user_obj_files
      type(string_t) :: basename, ext
      logical :: exist
      basename = user_src_file
      call split (basename, ext, ".", back=.true.)
      if ("." // ext == os_data%fc_src_ext) then
         inquire (file = char (user_src_file), exist = exist)
         if (exist) then
            call msg_message ("Compiling user code file '" &
                 // char (user_src_file) // "'")
            call os_compile_shared (basename, os_data)
            user_obj_files = user_obj_files // " " // basename // ".lo"
         else
            call msg_error ("User code file '" // char (user_src_file) &
                 // "' not found.")
         end if
      else
         call msg_error ("User code file '" // char (user_src_file) &
              // "' should have file extension '" &
              // char (os_data%fc_src_ext) // "'")
      end if
    end subroutine compile_user_src
    subroutine link_user (user_lib, user_obj_files)
      character(*), intent(in) :: user_lib
      type(string_t), intent(in) :: user_obj_files
      call msg_message ("Linking user code library '" &
           // user_lib // char (os_data%shlib_ext) // "'")
      call os_link_shared (user_obj_files, var_str (user_lib), os_data)
    end subroutine link_user
  end subroutine user_code_init

\eatline
\nwendcode{}\nwbegindocs{10}Unload all user-code libraries.
\nwenddocs{}\nwbegincode{11}\moddef{User Code: public}\plusendmoddef
  public :: user_code_final
\nwendcode{}\nwbegincode{12}\moddef{User Code: procedures}\plusendmoddef
  subroutine user_code_final ()
    if (has_user_lib) then
       call dlaccess_final (user_lib_handle)
       has_user_lib = .false.
    end if
  end subroutine user_code_final

\eatline
\nwendcode{}\nwbegindocs{13}Try to load the possible user-defined procedures from the dlopened
libraries.  If a procedure is not found, do nothing.
\nwenddocs{}\nwbegincode{14}\moddef{User Code: public}\plusendmoddef
  public :: user_code_find_proc
\nwendcode{}\nwbegincode{15}\moddef{User Code: procedures}\plusendmoddef
  function user_code_find_proc (name) result (fptr)
    type(string_t), intent(in) :: name
    type(c_funptr) :: fptr
    integer :: i
    fptr = c_null_funptr
    !!! Ticket #529
    ! fptr = libmanager_get_c_funptr (char (user), char (name))
    if (.not. c_associated (fptr)) then
       if (has_user_lib) then
          fptr = dlaccess_get_c_funptr (user_lib_handle, name)
          if (.not. c_associated (fptr)) then      
             call msg_fatal ("User procedure '" // char (name) // "' not found")
          end if
       else
          call msg_fatal ("User procedure '" // char (name) &
               // "' called without user library (missing -u flag?)")
       end if
    end if
  end function user_code_find_proc

\nwendcode{}\nwbegindocs{16} %def user_code_find_proc
\nwenddocs{}\nwbegindocs{17}\nwdocspar
\subsection{Interfaces for user-defined functions}
The following functions represent user-defined real observables.  There may
be one or two particles as argument, the result is a real value.
\nwenddocs{}\nwbegincode{18}\moddef{User Code: public}\plusendmoddef
  public :: user_obs_int_unary
  public :: user_obs_int_binary
  public :: user_obs_real_unary
  public :: user_obs_real_binary
\nwendcode{}\nwbegincode{19}\moddef{User Code: interfaces}\endmoddef
  abstract interface
     function user_obs_int_unary (prt1) result (ival) bind(C)
       use iso_c_binding !NODEP!
       use c_particles !NODEP!
       type(c_prt_t), intent(in) :: prt1
       integer(c_int) :: ival
     end function user_obs_int_unary
  end interface

  abstract interface
     function user_obs_int_binary (prt1, prt2) result (ival) bind(C)
       use iso_c_binding !NODEP!
       use c_particles !NODEP!
       type(c_prt_t), intent(in) :: prt1, prt2
       integer(c_int) :: ival
     end function user_obs_int_binary
  end interface

  abstract interface
     function user_obs_real_unary (prt1) result (rval) bind(C)
       use iso_c_binding !NODEP!
       use c_particles !NODEP!
       type(c_prt_t), intent(in) :: prt1
       real(c_double) :: rval
     end function user_obs_real_unary
  end interface

  abstract interface
     function user_obs_real_binary (prt1, prt2) result (rval) bind(C)
       use iso_c_binding !NODEP!
       use c_particles !NODEP!
       type(c_prt_t), intent(in) :: prt1, prt2
       real(c_double) :: rval
     end function user_obs_real_binary
  end interface

\eatline
\eatline
\nwendcode{}\nwbegindocs{20}\nwdocspar
The following function takes an array of C-compatible particles and
return a single value.  The particle array represents a subevent.  For
C interoperability, we have to use an assumed-size array, hence the
array size has to be transferred explicitly.

The cut function returns an {\tt{}int}, which we should interpret as a
logical value (nonzero=true).
\nwenddocs{}\nwbegincode{21}\moddef{User Code: public}\plusendmoddef
  public :: user_cut_fun
\nwendcode{}\nwbegincode{22}\moddef{User Code: interfaces}\plusendmoddef
  abstract interface
     function user_cut_fun (prt, n_prt) result (iflag) bind(C)
       use iso_c_binding !NODEP!
       use c_particles !NODEP!
       type(c_prt_t), dimension(*), intent(in) :: prt
       integer(c_int), intent(in) :: n_prt
       integer(c_int) :: iflag
     end function user_cut_fun
  end interface

\eatline
\nwendcode{}\nwbegindocs{23}The event-shape function returns a real value.
\nwenddocs{}\nwbegincode{24}\moddef{User Code: public}\plusendmoddef
  public :: user_event_shape_fun
\nwendcode{}\nwbegincode{25}\moddef{User Code: interfaces}\plusendmoddef
  abstract interface
     function user_event_shape_fun (prt, n_prt) result (rval) bind(C)
       use iso_c_binding !NODEP!
       use c_particles !NODEP!
       type(c_prt_t), dimension(*), intent(in) :: prt
       integer(c_int), intent(in) :: n_prt
       real(c_double) :: rval
     end function user_event_shape_fun
  end interface

\eatline
\nwendcode{}\nwbegindocs{26}\nwdocspar
\subsection{Interfaces for user-defined interactions}
The following procedure interfaces pertain to user-defined
interactions, e.g., spectra or structure functions.

This subroutine retrieves the basic information for setting up the
interaction and event generation.  All parameters are
{\tt{}intent(inout)}, so we can provide default values.  {\tt{}n{\char95}in} and
{\tt{}n{\char95}out} are the number of incoming and outgoing particles,
respectively.  {\tt{}n{\char95}states} is the total number of distinct states
that should be generated (counting all states of the incoming
particles).  {\tt{}n{\char95}col} is the maximal number of color entries a
particle can have.  {\tt{}n{\char95}dim} is the number of input parameters, i.e.,
integration dimensions, that the structure function call requires for
computing kinematics and matrix elements.
{\tt{}n{\char95}var} is the number of variables (e.g., momentum fractions) that
the structure function call has to transfer from the kinematics to the
dynamics evaluation.
\nwenddocs{}\nwbegincode{27}\moddef{User Code: public}\plusendmoddef
  public :: user_int_info
\nwendcode{}\nwbegincode{28}\moddef{User Code: interfaces}\plusendmoddef
  abstract interface
     subroutine user_int_info (n_in, n_out, n_states, n_col, n_dim, n_var) &
          bind(C)
       use iso_c_binding !NODEP!
       integer(c_int), intent(inout) :: n_in, n_out, n_states, n_col
       integer(c_int), intent(inout) :: n_dim, n_var
     end subroutine user_int_info
  end interface

\nwendcode{}\nwbegindocs{29} %def user_int_info
\nwenddocs{}\nwbegindocs{30}This subroutine retrieves the settings for the quantum number mask
of a given particle index in the interaction.  A nonzero value
indicates that the corresponding quantum number is to be ignored.  The
lock index is the index of a particle that the current particle is
related to.  The relation applies if quantum numbers of one of the
locked particles are summed over.  (This is intended for helicity.)
\nwenddocs{}\nwbegincode{31}\moddef{User Code: public}\plusendmoddef
  public :: user_int_mask
\nwendcode{}\nwbegincode{32}\moddef{User Code: interfaces}\plusendmoddef
  abstract interface
     subroutine user_int_mask (i_prt, m_flv, m_hel, m_col, i_lock) bind(C)
       use iso_c_binding !NODEP!
       integer(c_int), intent(in) :: i_prt
       integer(c_int), intent(inout) :: m_flv, m_hel, m_col, i_lock
     end subroutine user_int_mask
  end interface

\eatline
\nwendcode{}\nwbegindocs{33}This subroutine retrieves the quantum numbers for the particle
index {\tt{}i{\char95}prt} in state {\tt{}i{\char95}state}.  The {\tt{}flv} value is a PDG
code.  The {\tt{}hel} value is an integer helicity (twice the helicity
for fermions).  The {\tt{}col} array is an array which has at most
{\tt{}n{\char95}col} entries (see above).  All parameters are {\tt{}intent(inout)}
since default values exist.  In particular, if a mask entry is set by
the previous procedure, the corresponding quantum number is ignored
anyway.
\nwenddocs{}\nwbegincode{34}\moddef{User Code: public}\plusendmoddef
  public :: user_int_state
\nwendcode{}\nwbegincode{35}\moddef{User Code: interfaces}\plusendmoddef
  abstract interface
     subroutine user_int_state (i_state, i_prt, flv, hel, col) bind(C)
       use iso_c_binding !NODEP!
       integer(c_int), intent(in) :: i_state, i_prt
       integer(c_int), intent(inout) :: flv, hel
       integer(c_int), dimension(*), intent(inout) :: col
     end subroutine user_int_state
  end interface

\eatline
\nwendcode{}\nwbegindocs{36}This subroutine takes an array of particle objects
with array length {\tt{}n{\char95}in} and an array of input parameters between 0 and 1
with array length {\tt{}n{\char95}dim}.  It returns an array of particle objects 
with array length {\tt{}n{\char95}out}.  In addition, it returns an array of
internal variables (e.g., momentum fractions, Jacobian) with array
length {\tt{}n{\char95}var} that is used by the
following subroutine for evaluating the dynamics, i.e., the matrix
elements.
\nwenddocs{}\nwbegincode{37}\moddef{User Code: public}\plusendmoddef
  public :: user_int_kinematics
\nwendcode{}\nwbegincode{38}\moddef{User Code: interfaces}\plusendmoddef
  abstract interface
     subroutine user_int_kinematics (prt_in, rval, prt_out, xval) bind(C)
       use iso_c_binding !NODEP!
       use c_particles !NODEP!
       type(c_prt_t), dimension(*), intent(in) :: prt_in
       real(c_double), dimension(*), intent(in) :: rval
       type(c_prt_t), dimension(*), intent(inout) :: prt_out
       real(c_double), dimension(*), intent(out) :: xval
     end subroutine user_int_kinematics
  end interface

\eatline
\nwendcode{}\nwbegindocs{39}This subroutine takes the array of variables (e.g., momentum
fractions) with length {\tt{}n{\char95}var} which has been generated by the
previous subroutine and a real variable, the energy scale of the
event.  It returns an array of matrix-element values, one entry for
each quantum state {\tt{}n{\char95}states}.  The ordering of matrix elements must
correspond to the ordering of states.
\nwenddocs{}\nwbegincode{40}\moddef{User Code: public}\plusendmoddef
  public :: user_int_evaluate
\nwendcode{}\nwbegincode{41}\moddef{User Code: interfaces}\plusendmoddef
  abstract interface
     subroutine user_int_evaluate (xval, scale, fval) bind(C)
       use iso_c_binding !NODEP!
       real(c_double), dimension(*), intent(in) :: xval
       real(c_double), intent(in) :: scale
       real(c_double), dimension(*), intent(out) :: fval
     end subroutine user_int_evaluate
  end interface

\eatline
\nwendcode{}\nwbegindocs{42}\nwdocspar
%%%%%%%%%%%%%%%%%%%%%%%%%%%%%%%%%%%%%%%%%%%%%%%%%%%%%%%%%%%%%%%%%%%%%%%%
\section{User Plugin for Structure Functions}
This variant gives access to user-defined structure functions or spectra.

\subsection{The module}
\nwenddocs{}\nwbegincode{43}\moddef{\code{}sf{\char95}user.f90\edoc{}}\endmoddef
\LA{}File header\RA{}

module sf_user

  use, intrinsic :: iso_c_binding !NODEP!
\LA{}Use kinds\RA{}
\LA{}Use strings\RA{}
  use io_units
  use format_defs, only: FMT_17
  use diagnostics
  use c_particles
  use lorentz
  use subevents
  use user_code_interface
  use pdg_arrays
  use model_data
  use flavors
  use helicities
  use colors
  use quantum_numbers
  use state_matrices
  use polarizations
  use interactions
  use sf_aux
  use sf_base
  
\LA{}Standard module head\RA{}

\LA{}SF user: public\RA{}

\LA{}SF user: types\RA{}

contains

\LA{}SF user: procedures\RA{}

end module sf_user
\eatline
\nwendcode{}\nwbegindocs{44} 
\subsection{The user structure function data block} 
The data block holds the procedure pointers that are used for retrieving
static information, as well as the actual evaluation.
\nwenddocs{}\nwbegincode{45}\moddef{SF user: public}\endmoddef
  public :: user_data_t
\nwendcode{}\nwbegincode{46}\moddef{SF user: types}\endmoddef
  type, extends(sf_data_t) :: user_data_t
     private
     type(string_t) :: name
     integer :: n_in
     integer :: n_out
     integer :: n_tot
     integer :: n_states
     integer :: n_col
     integer :: n_dim
     integer :: n_var
     integer, dimension(2) :: pdg_in
     class(model_data_t), pointer :: model => null ()
     procedure(user_int_info), nopass, pointer :: info => null ()
     procedure(user_int_mask), nopass, pointer :: mask => null ()
     procedure(user_int_state), nopass, pointer :: state => null ()
     procedure(user_int_kinematics), nopass, pointer :: kinematics => null ()
     procedure(user_int_evaluate), nopass, pointer :: evaluate => null ()
   contains
   \LA{}SF user: user data: TBP\RA{}  
  end type user_data_t

\eatline
\nwendcode{}\nwbegindocs{47}Assign procedure pointers from a dynamically loaded library, given the
specified {\tt{}name}.  

We have to distinguish three cases: (1) Both beams are affected, and
the user spectrum implements both beams.  There is a single data
object. (2) Both beams are
affected, and the user spectrum applies to single beams.  Fill two
different objects.  (3) A single beam is affected.
\nwenddocs{}\nwbegincode{48}\moddef{SF User: public}\endmoddef
  public :: sf_user_data_init
\nwendcode{}\nwbegincode{49}\moddef{SF User: procedures}\endmoddef
  subroutine sf_user_data_init (data, name, flv, model)
    type(sf_user_data_t), intent(out) :: data
    type(string_t), intent(in) :: name
    type(flavor_t), dimension(2), intent(in) :: flv
    class(model_data_t), intent(in), target :: model
    integer(c_int) :: n_in
    integer(c_int) :: n_out
    integer(c_int) :: n_states
    integer(c_int) :: n_col
    integer(c_int) :: n_dim
    integer(c_int) :: n_var
    data%name = name
    data%pdg_in = flavor_get_pdg (flv)
    data%model => model
    call c_f_procpointer (user_code_find_proc (name // "_info"), data%info)
    call c_f_procpointer (user_code_find_proc (name // "_mask"), data%mask)
    call c_f_procpointer (user_code_find_proc (name // "_state"), data%state)
    call c_f_procpointer &
         (user_code_find_proc (name // "_kinematics"), data%kinematics)
    call c_f_procpointer &
         (user_code_find_proc (name // "_evaluate"), data%evaluate)
    n_in = 1
    n_out = 2
    n_states = 1
    n_col = 2
    n_dim = 1
    n_var = 1
    call data%info (n_in, n_out, n_states, n_col, n_dim, n_var)
    data%n_in = n_in
    data%n_out = n_out
    data%n_tot = n_in + n_out
    data%n_states = n_states
    data%n_col = n_col
    data%n_dim = n_dim
    data%n_var = n_var
  end subroutine sf_user_data_init

\eatline
\nwendcode{}\nwbegindocs{50}Output
\nwenddocs{}\nwbegincode{51}\moddef{SF user: user data: TBP}\endmoddef
  procedure :: write => user_data_write
\nwendcode{}\nwbegincode{52}\moddef{SF user: procedures}\endmoddef
  subroutine user_data_write (data, unit, verbose) 
    class(user_data_t), intent(in) :: data
    integer, intent(in), optional :: unit
    logical, intent(in), optional :: verbose        
    integer :: u
    u = given_output_unit (unit);  if (u < 0)  return
    write (u, "(1x,A,A)") "User structure function: ", char (data%name)
  end subroutine user_data_write

\eatline
\nwendcode{}\nwbegindocs{53}Retrieving contents
\nwenddocs{}\nwbegincode{54}\moddef{SF User: public}\plusendmoddef
  public :: sf_user_data_get_name
\nwendcode{}\nwbegincode{55}\moddef{SF User: procedures}\plusendmoddef
  function sf_user_data_get_name (data) result (name)
    type(string_t) :: name
    type(sf_user_data_t), intent(in) :: data
    name = data%name
  end function sf_user_data_get_name

\eatline
\nwendcode{}\nwbegincode{56}\moddef{SF User: public}\plusendmoddef
  public :: sf_user_data_get_n_in
  public :: sf_user_data_get_n_out
  public :: sf_user_data_get_n_tot
  public :: sf_user_data_get_n_dim
  public :: sf_user_data_get_n_var
\nwendcode{}\nwbegincode{57}\moddef{SF User: procedures}\plusendmoddef
  function sf_user_data_get_n_in (data) result (n_in)
    integer :: n_in
    type(sf_user_data_t), intent(in) :: data
    n_in = data%n_in
  end function sf_user_data_get_n_in

   function sf_user_data_get_n_out (data) result (n_out)
    integer :: n_out
    type(sf_user_data_t), intent(in) :: data
    n_out = data%n_out
  end function sf_user_data_get_n_out

  function sf_user_data_get_n_tot (data) result (n_tot)
    integer :: n_tot
    type(sf_user_data_t), intent(in) :: data
    n_tot = data%n_tot
  end function sf_user_data_get_n_tot

  function sf_user_data_get_n_dim (data) result (n_dim)
    integer :: n_dim
    type(sf_user_data_t), intent(in) :: data
    n_dim = data%n_dim
  end function sf_user_data_get_n_dim

 function sf_user_data_get_n_var (data) result (n_var)
    integer :: n_var
    type(sf_user_data_t), intent(in) :: data
    n_var = data%n_var
  end function sf_user_data_get_n_var

\eatline
\eatline
\eatline
\eatline
\eatline
\nwendcode{}\nwbegindocs{58}\nwdocspar
\subsection{The interaction}
We fill the interaction by looking up the table of states using the interface
functions.

For particles which have a known flavor (as indicated by the mask), we
compute the mass squared, so we can use it for the invariant mass of
the particle objects.
\nwenddocs{}\nwbegincode{59}\moddef{SF user: user: TBP}\endmoddef
  procedure :: init => user_init
\nwendcode{}\nwbegincode{60}\moddef{SF user: procedures}\plusendmoddef
  subroutine user_init (sf_int, data)
    !!! JRR: WK please check (#529)
    class(user_t), intent(out) :: sf_int
    class(sf_data_t), intent(in), target :: data
    type(quantum_numbers_mask_t), dimension(:), allocatable :: mask
    integer, dimension(:), allocatable :: hel_lock
    integer(c_int) :: m_flv, m_hel, m_col, i_lock
    type(quantum_numbers_t), dimension(:), allocatable :: qn
    integer(c_int) :: f, h
    integer(c_int), dimension(:), allocatable :: c
    type(flavor_t) :: flv
    type(helicity_t) :: hel
    type(color_t) :: col
    integer :: i, s
    integer(c_int) :: i_prt, i_state
    select type (data)
    type is (user_data_t)   
       allocate (mask (data%n_tot))
       allocate (hel_lock (data%n_tot))
       allocate (qn (data%n_tot))
       allocate (c (data%n_col))
       do i = 1, size (mask)
          i_prt = i
          m_flv = 0;  m_col = 0;  m_hel = 0;  i_lock = 0
          call data%mask (i_prt, m_flv, m_col, m_hel, i_lock)
          mask(i) = &
               quantum_numbers_mask (m_flv /= 0, m_col /= 0, m_hel /= 0)
          hel_lock(i) = i_lock
       end do
       !!! JRR: WK please check (#529)
       !!! Will have to be filled in later.
       ! call sf_int%base_init (mask, &
       !      hel_lock = hel_lock)
       call interaction_init &
            (sf_int%interaction_t, data%n_in, 0, data%n_out, mask=mask, &
            hel_lock=hel_lock, set_relations=.true.)       
       do s = 1, data%n_states
          i_state = s
          do i = 1, data%n_tot
             i_prt = i
             f = 0;  h = 0;  c = 0
             call data%state (i_state, i_prt, f, h, c)
             if (m_flv == 0) then
                call flv%init (int (f), data%model)
             else
                call flv%init ()
             end if
             if (m_hel == 0) then
                call hel%init (int (h))
             else
                call hel%init ()
             end if
             if (m_col == 0) then
                call color_init_from_array (col, int (c))
             else
                call col%init ()
             end if
             call qn(i)%init (flv, col, hel)
          end do
          call interaction_add_state (sf_int%interaction_t, qn)
       end do
       call interaction_freeze (sf_int%interaction_t)
       !!! JRR: WK please check (#529)
       !!! What has to be inserted here?
       ! call sf_int%set_incoming (??)
       ! call sf_int%set_radiated (??)
       ! call sf_int%set_outgoing (??)
       sf_int%status = SF_INITIAL
    end select
  end subroutine user_init

\eatline
\nwendcode{}\nwbegindocs{61}\nwdocspar
\nwenddocs{}\nwbegindocs{62}Allocate the interaction record.
\nwenddocs{}\nwbegincode{63}\moddef{SF user: user data: TBP}\plusendmoddef
  procedure :: allocate_sf_int => user_data_allocate_sf_int
\nwendcode{}\nwbegincode{64}\moddef{SF user: procedures}\plusendmoddef
  subroutine user_data_allocate_sf_int (data, sf_int)
    class(user_data_t), intent(in) :: data
    class(sf_int_t), intent(inout), allocatable :: sf_int
    allocate (user_t :: sf_int)
  end subroutine user_data_allocate_sf_int
  
\eatline
\nwendcode{}\nwbegindocs{65}The number of parameters is one.  We do not generate transverse momentum.
\nwenddocs{}\nwbegincode{66}\moddef{SF user: user data: TBP}\plusendmoddef
  procedure :: get_n_par => user_data_get_n_par
\nwendcode{}\nwbegincode{67}\moddef{SF user: procedures}\plusendmoddef
  function user_data_get_n_par (data) result (n)
    class(user_data_t), intent(in) :: data
    integer :: n
    n = data%n_var
  end function user_data_get_n_par
  
\eatline
\nwendcode{}\nwbegindocs{68}\nwdocspar
\nwenddocs{}\nwbegindocs{69}Return the outgoing particle PDG codes. This has to be inferred from
the states (right?). JRR: WK please check.
\nwenddocs{}\nwbegincode{70}\moddef{SF user: user data: TBP}\plusendmoddef
  procedure :: get_pdg_out => user_data_get_pdg_out
\nwendcode{}\nwbegincode{71}\moddef{SF user: procedures}\plusendmoddef
  subroutine user_data_get_pdg_out (data, pdg_out)
    class(user_data_t), intent(in) :: data
    type(pdg_array_t), dimension(:), intent(inout) :: pdg_out
    !!! JRR: WK please check (#529)
    !!! integer :: n, np, i
    !!! n = count (data%mask)
    !!! np = 0;  if (data%has_photon .and. data%mask_photon)  np = 1
    !!! allocate (pdg_out (n + np))
    !!! pdg_out(1:n) = pack ([(i, i = -6, 6)], data%mask)
    !!! if (np == 1)  pdg_out(n+np) = PHOTON
  end subroutine user_data_get_pdg_out
  
\eatline
\nwendcode{}\nwbegindocs{72}\subsection{The user structure function}
For maximal flexibility, user structure functions separate kinematics from
dynamics just as the PDF interface does.  (JRR: Ok, I guess this now
done for all structure functions, right?) We create {\tt{}c{\char95}prt{\char95}t}
particle objects from the incoming momenta (all other quantum numbers
are irrelevant) and call the user-supplied kinematics function to
compute the outgoing momenta, along with other variables that will be
needed for matrix element evaluation.  If known, we use the mass
squared computed above.
!!! JRR: WK please check (\#529)
I don't know actually whether this really fits into the setup done by
WK. 
\nwenddocs{}\nwbegincode{73}\moddef{SF user: types}\plusendmoddef
  !!! JRR: WK please check (#529)
  type, extends (sf_int_t) :: user_t
     type(user_data_t), pointer :: data => null ()
     real(default) :: x = 0
     real(default) :: q = 0
   contains
   \LA{}SF user: user: TBP\RA{}
  end type user_t 
  
\eatline
\nwendcode{}\nwbegindocs{74}Type string: display the name of the user structure function.
\nwenddocs{}\nwbegincode{75}\moddef{SF user: user: TBP}\plusendmoddef
  procedure :: type_string => user_type_string
\nwendcode{}\nwbegincode{76}\moddef{SF user: procedures}\plusendmoddef
  function user_type_string (object) result (string)
    class(user_t), intent(in) :: object
    type(string_t) :: string
    if (associated (object%data)) then
       string = "User structure function: " // object%data%name
    else
       string = "User structure function: [undefined]"
    end if
  end function user_type_string
  
\eatline
\nwendcode{}\nwbegindocs{77}Output.  Call the interaction routine after displaying the configuration.
\nwenddocs{}\nwbegincode{78}\moddef{SF user: user: TBP}\plusendmoddef
  procedure :: write => user_write
\nwendcode{}\nwbegincode{79}\moddef{SF user: procedures}\plusendmoddef
  subroutine user_write (object, unit, testflag)
    !!! JRR: WK please check (#529)
    !!! Guess these variables do not exist for user strfun (?)
    class(user_t), intent(in) :: object
    integer, intent(in), optional :: unit
    logical, intent(in), optional :: testflag
    integer :: u
    u = given_output_unit (unit)
    if (associated (object%data)) then
       call object%data%write (u)
       if (object%status >= SF_DONE_KINEMATICS) then
          write (u, "(1x,A)")  "SF parameters:"
          write (u, "(3x,A," // FMT_17 // ")")  "x =", object%x
          if (object%status >= SF_FAILED_EVALUATION) then
             write (u, "(3x,A," // FMT_17 // ")")  "Q =", object%q
          end if
       end if
       call object%base_write (u, testflag)
    else
       write (u, "(1x,A)")  "User structure function data: [undefined]"
    end if
  end subroutine user_write
    
\eatline
\nwendcode{}\nwbegindocs{80}\nwdocspar
\subsection{Kinematics}
Set kinematics.  If {\tt{}map} is unset, the $r$ and $x$ values
coincide, and the Jacobian $f(r)$ is trivial.

If {\tt{}map} is set, we are asked to provide an efficient mapping.
For the test case, we set $x=r^2$ and consequently $f(r)=2r$.
\nwenddocs{}\nwbegincode{81}\moddef{SF user: user: TBP}\plusendmoddef
  procedure :: complete_kinematics => user_complete_kinematics
\nwendcode{}\nwbegincode{82}\moddef{SF user: procedures}\plusendmoddef
  subroutine user_complete_kinematics (sf_int, x, f, r, rb, map)
    !!! JRR: WK please check (#529)
    !!! This cannot be correct, as the CIRCE1 structure function has 
    !!! twice the variables (2->4 instead of 1->2 splitting)
    class(user_t), intent(inout) :: sf_int
    real(default), dimension(:), intent(out) :: x
    real(default), intent(out) :: f
    real(default), dimension(:), intent(in) :: r
    real(default), dimension(:), intent(in) :: rb
    logical, intent(in) :: map
    real(default) :: xb1
    if (map) then
       call msg_fatal ("User structure function: map flag not supported")
    else
       x(1) = r(1)
       f = 1
    end if
    xb1 = 1 - x(1)
    call sf_int%split_momentum (x, xb1)
    select case (sf_int%status)
    case (SF_DONE_KINEMATICS)
       sf_int%x = x(1)
    case (SF_FAILED_KINEMATICS)
       sf_int%x = 0
       f = 0
    end select
  end subroutine user_complete_kinematics

\eatline
\nwendcode{}\nwbegindocs{83}Compute inverse kinematics.  Here, we start with the $x$ array and
compute the ``input'' $r$ values and the Jacobian $f$.  After this, we
can set momenta by the same formula as for normal kinematics.
\nwenddocs{}\nwbegincode{84}\moddef{SF user: user: TBP}\plusendmoddef
  procedure :: inverse_kinematics => user_inverse_kinematics
\nwendcode{}\nwbegincode{85}\moddef{SF user: procedures}\plusendmoddef
  subroutine user_inverse_kinematics (sf_int, x, f, r, rb, map, set_momenta)
    !!! JRR: WK please check (#529)
    !!! This cannot be correct, as the CIRCE1 structure function has 
    !!! twice the variables (2->4 instead of 1->2 splitting)
    class(user_t), intent(inout) :: sf_int
    real(default), dimension(:), intent(in) :: x
    real(default), intent(out) :: f
    real(default), dimension(:), intent(out) :: r
    real(default), dimension(:), intent(out) :: rb
    logical, intent(in) :: map
    logical, intent(in), optional :: set_momenta
    real(default) :: xb1
    logical :: set_mom
    set_mom = .false.;  if (present (set_momenta))  set_mom = set_momenta
    if (map) then
       call msg_fatal ("User structure function: map flag not supported")
    else
       r(1) = x(1)
       f = 1
    end if
    xb1 = 1 - x(1)
    rb = 1 - r
    if (set_mom) then
       call sf_int%split_momentum (x, xb1)
       select case (sf_int%status)
       case (SF_DONE_KINEMATICS)
          sf_int%x = x(1)
       case (SF_FAILED_KINEMATICS)
          sf_int%x = 0
          f = 0
       end select
    end if
  end subroutine user_inverse_kinematics

\eatline
\nwendcode{}\nwbegindocs{86}\nwdocspar
\nwenddocs{}\nwbegincode{87}\moddef{SF User: public}\plusendmoddef
  public :: interaction_set_kinematics_sf_user
\nwendcode{}\nwbegincode{88}\moddef{SF User: procedures}\plusendmoddef
  subroutine interaction_set_kinematics_sf_user (int, x, r, data)
    type(interaction_t), intent(inout) :: int
    real(default), dimension(:), intent(out) :: x
    real(default), dimension(:), intent(in) :: r
    type(sf_user_data_t), intent(in) :: data
    type(vector4_t), dimension(data%n_in) :: p_in
    type(vector4_t), dimension(data%n_out) :: p_out
    type(c_prt_t), dimension(data%n_in) :: prt_in
    type(c_prt_t), dimension(data%n_out) :: prt_out
    real(c_double), dimension(data%n_var) :: xval
    call interaction_get_momenta_sub (int, p_in, outgoing=.false.)
    prt_in = vector4_to_c_prt (p_in)
    prt_in%type = PRT_INCOMING
    call data%kinematics (prt_in, real (r, c_double), prt_out, xval)
    x = xval
    p_out = vector4_from_c_prt (prt_out)
    call interaction_set_momenta (int, p_out, outgoing=.true.)
  end subroutine interaction_set_kinematics_sf_user

\eatline
\nwendcode{}\nwbegindocs{89}The matrix-element evaluation may require a scale parameter, therefore this
routine is separate.  We take the variables computed above together
with the event energy scale and call the user function that computes
the matrix elements.
\nwenddocs{}\nwbegincode{90}\moddef{SF user: user: TBP}\plusendmoddef
  procedure :: apply => user_apply
\nwendcode{}\nwbegincode{91}\moddef{SF user: procedures}\plusendmoddef
  subroutine user_apply (sf_int, scale) !, x, data)    
    !!! JRR: WK please check (#529)
    class(user_t), intent(inout) :: sf_int
    real(default), intent(in) :: scale
    real(default), dimension(:), allocatable :: x
    real(c_double), dimension(sf_int%data%n_states) :: fval
    complex(default), dimension(sf_int%data%n_states) :: fc
    associate (data => sf_int%data)
      !!! This is wrong, has to be replaced
      ! allocate (x, size (sf_int%x)))
      x = sf_int%x
      call data%evaluate (real (x, c_double), real (scale, c_double), fval)
      fc = fval
      call interaction_set_matrix_element (sf_int%interaction_t, fc)
    end associate
    sf_int%status = SF_EVALUATED
  end subroutine user_apply

\eatline
\nwendcode{}\nwfilename{../../../src/blha/blha.nw}\nwbegindocs{0}% -*- ess-noweb-default-code-mode: f90-mode; noweb-default-code-mode: f90-mode; -*- 
% WHIZARD code as NOWEB source: matrix elements and process libraries

\nwenddocs{}\nwbegindocs{1}\nwdocspar
%%%%%%%%%%%%%%%%%%%%%%%%%%%%%%%%%%%%%%%%%%%%%%%%%%%%%%%%%%%%%%%%%%%%%%%%
\chapter{BLHA Interface}

The code in this chapter implements support for the BLHA record that
communicates data for NLO processes.

These are the modules:
\begin{description}
\item[blha\_config]
\item[blha\_interface]
\item[blha\_driver]
\end{description}

\nwenddocs{}\nwbegindocs{2}\nwdocspar
%%%%%%%%%%%%%%%%%%%%%%%%%%%%%%%%%%%%%%%%%%%%%%%%%%%%%%%%%%%%%%%%%%%%%%%%
\section{Module definition}

These modules implement the communication with one loop matrix element providers
according to the Binoth LesHouches Accord Interface. The actual matrix
element(s) are loaded as a dynamic library.

The module is split into a configuration interface which manages configuration
and handles the request and contract files, a module which interfaces the OLP
matrix elements and a driver.

\nwenddocs{}\nwbegincode{3}\moddef{\code{}blha{\char95}config.f90\edoc{}}\endmoddef
\LA{}File header\RA{}

module blha_config

  use kinds
\LA{}Use strings\RA{}
  use io_units
  use constants
  use string_utils
  use system_defs, only: EOF
  use diagnostics
  use md5
  use model_data
  use flavors
  use quantum_numbers
  use pdg_arrays
  use sorting
  use lexers
  use parser
  use syntax_rules
  use ifiles

\LA{}Standard module head\RA{}

\LA{}BLHA config: public\RA{}

\LA{}BLHA config: parameters\RA{}

\LA{}BLHA config: types\RA{}

\LA{}BLHA config: variables\RA{}

\LA{}BLHA config: interfaces\RA{}

contains

\LA{}BLHA config: procedures\RA{}

end module blha_config

\eatline
\nwendcode{}\nwbegindocs{4}\nwdocspar
\nwenddocs{}\nwbegincode{5}\moddef{\code{}blha{\char95}interface.f90\edoc{}}\endmoddef
\LA{}File header\RA{}

module blha_interface

\LA{}Use strings\RA{}
  use io_units
  use string_utils
  use constants
  use diagnostics
  use sm_physics
  use md5
  use lorentz
  use model_data
  use flavors
  use quantum_numbers
  use interactions
  use evaluators
  use particles
  use quantum_numbers
  use blha_config
  use, intrinsic :: iso_c_binding !NODEP!
  use os_interface

\LA{}Standard module head\RA{}

\LA{}BLHA interface: public\RA{}

\LA{}BLHA interface: types\RA{}

\LA{}BLHA interface: interfaces\RA{}

contains

\LA{}BLHA interface: procedures\RA{}

end module blha_interface
\eatline
\nwendcode{}\nwbegindocs{6}%
\nwenddocs{}\nwbegincode{7}\moddef{\code{}blha{\char95}driver.f90\edoc{}}\endmoddef
\LA{}File header\RA{}

module blha_driver

\LA{}Use strings\RA{}
  use io_units
  use constants
  use diagnostics
  use sm_physics
  use md5
  use lorentz
  use model_data
  use flavors
  use quantum_numbers
  use interactions
  use evaluators
  use particles
  use quantum_numbers
  use blha_config
  use blha_interface

\LA{}Standard module head\RA{}

\LA{}BLHA driver: public\RA{}

contains

\LA{}BLHA driver: procedures\RA{}

end module blha_driver
\eatline
\nwendcode{}\nwbegindocs{8}%

\section{Configuration}

Parameters to enumerate the different options in the order.
\nwenddocs{}\nwbegincode{9}\moddef{BLHA config: parameters}\endmoddef
  integer, public, parameter :: &
       BLHA_MEST_SUM=1, BLHA_MEST_AVG=2, BLHA_MEST_OTHER=3
  integer, public, parameter :: &
       BLHA_CT_QCD=1, BLHA_CT_EW=2, BLHA_CT_QED=3, BLHA_CT_OTHER=4
  integer, public, parameter :: &
       BLHA_IRREG_CDR=1, BLHA_IRREG_DRED=2, BLHA_IRREG_THV=3, &
       BLHA_IRREG_MREG=4, BLHA_IRREG_OTHER=5
  integer, public, parameter :: &
       BLHA_SUBMODE_NONE = 1, BLHA_SUBMODE_OTHER = 2
  integer, public, parameter :: &
       BLHA_MPS_ONSHELL=1, BLHA_MPS_OTHER=2
  integer, public, parameter :: &
       BLHA_MODE_GOSAM=1, BLHA_MODE_FEYNARTS = 2, BLHA_MODE_GENERIC=3
  integer, public, parameter :: &
       BLHA_OM_NONE=1, BLHA_OM_NOCPL=2, BLHA_OM_OTHER=3
   integer, public, parameter :: &
       BLHA_VERSION_1 = 1, BLHA_VERSION_2 = 2
  integer, public, parameter :: &
       BLHA_AMP_LOOP = 1, BLHA_AMP_CC = 2, BLHA_AMP_SC = 3, &
       BLHA_AMP_TREE = 4, BLHA_AMP_LOOPINDUCED = 5
  integer, public, parameter :: &
       BLHA_EW_GF = 1, BLHA_EW_MZ = 2, BLHA_EW_MSBAR = 3, &
       BLHA_EW_0 = 4, BLHA_EW_RUN = 5, BLHA_EW_DEFAULT = 6
  integer, public, parameter :: &
       BLHA_WIDTH_COMPLEX = 1, BLHA_WIDTH_FIXED = 2, &
       BLHA_WIDTH_RUNNING = 3, BLHA_WIDTH_POLE = 4, &
       BLHA_WIDTH_DEFAULT = 5 

\eatline
\eatline
\eatline
\eatline
\eatline
\eatline
\eatline
\eatline
\nwendcode{}\nwbegindocs{10}\nwdocspar
This type encapsulates a BLHA request.
\nwenddocs{}\nwbegincode{11}\moddef{BLHA config: public}\endmoddef
  public :: blha_configuration_t
  public :: blha_cfg_process_node_t
\nwendcode{}\nwbegincode{12}\moddef{BLHA config: types}\endmoddef
  type :: blha_cfg_process_node_t
     integer, dimension(:), allocatable :: pdg_in, pdg_out
     integer, dimension(:), allocatable :: fingerprint
     integer :: nsub
     integer, dimension(:), allocatable :: ids
     integer :: amplitude_type
     type(blha_cfg_process_node_t), pointer :: next => null ()
  end type blha_cfg_process_node_t

  type :: blha_configuration_t
     type(string_t) :: name
     class(model_data_t), pointer :: model => null ()
     type(string_t) :: md5
     integer :: version = 2
     logical :: dirty = .false.
     integer :: n_proc = 0
     real(default) :: accuracy_target
     logical :: debug_unstable
     integer :: mode = BLHA_MODE_GENERIC
     type(blha_cfg_process_node_t), pointer :: processes => null ()
     integer, dimension(2) :: matrix_element_square_type = BLHA_MEST_SUM
     !!! !!! !!! Triggers gfortran 4.9.0 ICE     
     ! type(string_t), dimension (2) :: matrix_element_square_type_other
     integer :: correction_type = BLHA_CT_QCD
     type(string_t) :: correction_type_other
     integer :: irreg = BLHA_IRREG_THV
     type(string_t) :: irreg_other
     integer :: massive_particle_scheme = BLHA_MPS_ONSHELL
     type(string_t) :: massive_particle_scheme_other
     integer :: subtraction_mode = BLHA_SUBMODE_NONE
     type(string_t) :: subtraction_mode_other
     type(string_t) :: model_file
     logical :: subdivide_subprocesses = .false.
     integer :: alphas_power = -1, alpha_power = -1
     integer :: ew_scheme = BLHA_EW_DEFAULT
     integer :: width_scheme = BLHA_WIDTH_DEFAULT
     integer :: operation_mode = BLHA_OM_NONE
     type(string_t) :: operation_mode_other
  end type blha_configuration_t

\eatline
\nwendcode{}\nwbegindocs{13}\nwdocspar
This types control the creation of BLHA-interface files
\nwenddocs{}\nwbegincode{14}\moddef{BLHA config: public}\plusendmoddef
  public :: blha_flv_state_t
  public :: blha_master_t
\nwendcode{}\nwbegincode{15}\moddef{BLHA config: types}\plusendmoddef
  type:: blha_flv_state_t
    integer, dimension(:), allocatable :: flavors
    integer :: flv_mult
    logical :: flv_real = .false.
  end type blha_flv_state_t

  type :: blha_master_t
    integer, dimension(:,:), allocatable :: flv_born, flv_real
    integer :: n_in, n_out
    logical :: compute_loops = .true.
    logical :: compute_correlations = .false.
    logical :: compute_real_trees = .false.
    integer :: alpha_power, alphas_power
    type(string_t) :: basename
    type(string_t), dimension(:), allocatable :: suffix
    type(blha_configuration_t), dimension(:), allocatable :: blha_cfg
    integer :: n_files = 0
  contains
  \LA{}BLHA config: blha master: TBP\RA{}
  end type blha_master_t

\eatline
\nwendcode{}\nwbegindocs{16}Master-Routines
\nwenddocs{}\nwbegincode{17}\moddef{BLHA config: blha master: TBP}\endmoddef
  procedure :: init => blha_master_init 
\nwendcode{}\nwbegincode{18}\moddef{BLHA config: procedures}\endmoddef
  subroutine blha_master_init (master, basename, model, &
                               n_in, n_out, &
                               cmp_loops, cmp_corr, cmp_real, &
                               alpha_power, alphas_power, &
                               flv_born, flv_real)
    class(blha_master_t), intent(inout) :: master
    type(string_t), intent(in) :: basename
    class(model_data_t), intent(in), target :: model
    integer, intent(in) :: n_in, n_out
    logical, intent(in) :: cmp_loops, cmp_corr, cmp_real
    integer, intent(in) :: alpha_power, alphas_power
    integer, dimension(:,:), allocatable, intent(in) :: &
                             flv_born, flv_real
    integer :: n_proc, n_proc_real, n_flv
    type(blha_flv_state_t), dimension(:), allocatable :: blha_flavor
    integer :: i, i_flv
    integer :: i_file
    integer, dimension(:), allocatable :: amp_type

    n_proc = 0; n_flv = 1; n_proc_real = 0
    master%compute_real_trees = cmp_real
    if (cmp_loops) master%n_files = master%n_files+1
    if (cmp_corr) master%n_files = master%n_files+1
    if (cmp_real) then
       if (allocated (flv_real)) then
         master%n_files = master%n_files+1
         n_proc_real = size (flv_real, 2)
         n_flv = n_flv + n_proc_real
       end if
    end if
    allocate (master%blha_cfg (master%n_files))
    allocate (master%suffix (master%n_files))
    i_file = 1
    if (cmp_loops) then
       if (allocated (flv_born)) then
          allocate (blha_flavor (size (flv_born, 2)))
          do i_flv = 1, size (flv_born, 2)
             blha_flavor(i_flv)%flavors = flv_born(:,i_flv)
             blha_flavor(i_flv)%flv_mult = 2
          end do
          master%suffix(i_file) = "_LOOP"
          call blha_init_virtual (master%blha_cfg(i_file), blha_flavor, &
                                  alpha_power, alphas_power, &
                                  basename, model)
          i_file = i_file+1
        else
          call msg_fatal ("GoSam Loops requested but " &
                           // "Born flavor not existing")
        end if
    end if
    if (allocated (blha_flavor)) deallocate (blha_flavor)
    if (cmp_corr) then
       if (allocated (flv_born)) then
          allocate (blha_flavor (size (flv_born, 2)))
          do i_flv = 1, size (flv_born, 2)
             blha_flavor(i_flv)%flavors = flv_born(:,i_flv)
             blha_flavor(i_flv)%flv_mult = 3
          end do
          master%suffix(i_file) = "_SUB"
          call blha_init_subtraction (master%blha_cfg(i_file), blha_flavor, &
                                  alpha_power, alphas_power, &
                                  basename, model)
          i_file = i_file+1
       else
          call msg_fatal ("GoSam Correlations requested but "&
                           // "Born flavor not existing")
       end if
    end if
    if (allocated (blha_flavor)) deallocate (blha_flavor)
    if (cmp_real) then
       if (allocated (flv_real)) then
          allocate (blha_flavor (size (flv_real, 2)))
          do i_flv = 1, size (flv_real, 2)
             blha_flavor(i_flv)%flavors = flv_real(:,i_flv)
             blha_flavor(i_flv)%flv_mult = 1
          end do
          master%suffix(i_file) = "_REAL"
          call blha_init_real (master%blha_cfg(i_file), blha_flavor, &
                                  alpha_power, alphas_power, &
                                  basename, model)
       else
          call msg_fatal ("GoSam Trees requested but "&
                           // "Real flavor not existing")
       end if
    end if
  end subroutine blha_master_init

\eatline
\nwendcode{}\nwbegindocs{19}\nwdocspar
\nwenddocs{}\nwbegincode{20}\moddef{BLHA config: procedures}\plusendmoddef
  subroutine blha_init_virtual (blha_cfg, blha_flavor, ap, asp, basename, model)
    type(blha_configuration_t), intent(inout) :: blha_cfg
    type(blha_flv_state_t), intent(in), dimension(:) :: blha_flavor
    integer, intent(in) :: ap, asp
    type(string_t), intent(in) :: basename
    type(model_data_t), intent(in), target :: model
    integer, dimension(:), allocatable :: amp_type
    integer :: i

    allocate (amp_type (size (blha_flavor)*2))
    do i = 1, size (blha_flavor)
       amp_type(2*i-1) = BLHA_AMP_LOOP
       amp_type(2*i) = BLHA_AMP_CC
    end do
    call blha_configuration_init (blha_cfg, basename // "_VIRT" , &
                                  model, BLHA_MODE_GOSAM)
    call blha_configuration_append_processes (blha_cfg, &
                                  blha_flavor, amp_type)
    call blha_configuration_set (blha_cfg, 2, &
                                 correction_type = BLHA_CT_QCD, &
                                 irreg = BLHA_IRREG_CDR, &
                                 alphas_power = asp, &
                                 alpha_power = ap, &
                                 ew_scheme = BLHA_EW_GF, &
                                 debug = .true.)
  end subroutine blha_init_virtual

  subroutine blha_init_subtraction (blha_cfg, blha_flavor, ap, asp, basename, model)
    type(blha_configuration_t), intent(inout) :: blha_cfg
    type(blha_flv_state_t), intent(in), dimension(:) :: blha_flavor
    integer, intent(in) :: ap, asp
    type(string_t), intent(in) :: basename
    type(model_data_t), intent(in), target :: model
    integer, dimension(:), allocatable :: amp_type
    integer :: i

    allocate (amp_type (size (blha_flavor)*3))
    do i = 1, size (blha_flavor)
       amp_type(3*i-2) = BLHA_AMP_TREE
       amp_type(3*i-1) = BLHA_AMP_CC
       amp_type(3*i) = BLHA_AMP_SC
    end do
    call blha_configuration_init (blha_cfg, basename // "_SUB" , &
                                  model, BLHA_MODE_GOSAM)
    call blha_configuration_append_processes (blha_cfg, &
                                  blha_flavor, amp_type)
    call blha_configuration_set (blha_cfg, 2, &
                                 correction_type = BLHA_CT_QCD, &
                                 irreg = BLHA_IRREG_CDR, &
                                 alphas_power = asp, &
                                 alpha_power = ap, &
                                 ew_scheme = BLHA_EW_GF, &
                                 debug = .true.)
  end subroutine blha_init_subtraction

  subroutine blha_init_real (blha_cfg, blha_flavor, ap, asp, basename, model)
    type(blha_configuration_t), intent(inout) :: blha_cfg
    type(blha_flv_state_t), intent(in), dimension(:) :: blha_flavor
    integer, intent(in) :: ap, asp
    type(string_t), intent(in) :: basename
    type(model_data_t), intent(in), target :: model
    integer, dimension(:), allocatable :: amp_type
    integer :: i

    allocate (amp_type (size (blha_flavor)))
    do i = 1, size (blha_flavor)
       amp_type(i) = BLHA_AMP_TREE
    end do
    call blha_configuration_init (blha_cfg, basename // "_REAL" , &
                                  model, BLHA_MODE_GOSAM)
    call blha_configuration_append_processes (blha_cfg, &
                                  blha_flavor, amp_type)
    call blha_configuration_set (blha_cfg, 2, &
                                 correction_type = BLHA_CT_QCD, &
                                 irreg = BLHA_IRREG_CDR, &
                                 alphas_power = asp+1, &
                                 alpha_power = ap, &
                                 ew_scheme = BLHA_EW_GF, &
                                 debug = .true.)
  end subroutine blha_init_real

\eatline
\nwendcode{}\nwbegindocs{21}\nwdocspar
\nwenddocs{}\nwbegincode{22}\moddef{BLHA config: blha master: TBP}\plusendmoddef
  procedure :: generate => blha_master_generate
\nwendcode{}\nwbegincode{23}\moddef{BLHA config: procedures}\plusendmoddef
  subroutine blha_master_generate (master, basename)
    class(blha_master_t), intent(in) :: master
    type(string_t), intent(in) :: basename
    integer :: unit
    type(string_t) :: filename
    integer :: i_file
    do i_file = 1, master%n_files
       filename = basename // master%suffix(i_file) // ".olp"
       unit = free_unit ()
       open (unit, file = char (filename), status = 'replace', action = 'write')
       call blha_configuration_write (master%blha_cfg(i_file), unit)
       close (unit)
    end do
  end subroutine blha_master_generate

\eatline
\nwendcode{}\nwbegindocs{24}   
Initializer.
\nwenddocs{}\nwbegincode{25}\moddef{BLHA config: public}\plusendmoddef
  public :: blha_configuration_init
\nwendcode{}\nwbegincode{26}\moddef{BLHA config: procedures}\plusendmoddef
  subroutine blha_configuration_init (cfg, name, model, mode)
    type(blha_configuration_t), intent(inout) :: cfg
    type(string_t), intent(in) :: name
    class(model_data_t), target, intent(in) :: model
    integer, intent(in), optional :: mode
    if (.not. associated (cfg%model)) then
       cfg%name = name
       cfg%model => model
    end if
    if (present (mode)) cfg%mode = mode
  end subroutine blha_configuration_init

\eatline
\nwendcode{}\nwbegindocs{27}\nwdocspar
Finalizer.
\nwenddocs{}\nwbegincode{28}\moddef{BLHA config: public}\plusendmoddef
  public :: blha_configuration_final
\nwendcode{}\nwbegincode{29}\moddef{BLHA config: procedures}\plusendmoddef
  subroutine blha_configuration_final (cfg)
    type(blha_configuration_t), intent(inout) :: cfg
    type(blha_cfg_process_node_t), pointer :: cur, next
    cur => cfg%processes
    do while (associated (cur))
       next => cur%next
       deallocate (cur)
       nullify (cur)
       cur => next
    end do
  end subroutine blha_configuration_final

\eatline
\nwendcode{}\nwbegindocs{30}\nwdocspar
Merge sort a process list w.r.t. to the process fingerprints. This is
necessary for canonicalizing the process list prior to calculating the
MD5 sum. 
\nwenddocs{}\nwbegincode{31}\moddef{BLHA config: procedures}\plusendmoddef
  subroutine sort_processes (list, n)
    type(blha_cfg_process_node_t), pointer :: list
    integer, intent(in), optional :: n
    type :: pnode
       type(blha_cfg_process_node_t), pointer :: p
    end type pnode
    type(pnode), dimension(:), allocatable :: array
    integer :: count, i, s, i1, i2, i3
    type(blha_cfg_process_node_t), pointer :: node
    if (present (n)) then
       count = n
    else
       node => list
       count = 0
       do while (associated (node))
          node => node%next
          count = count + 1
       end do
    end if
    ! Store list nodes into an array
    if (count == 0) return
    allocate (array(count))
    i = 1
    node => list
    do i = 1, count
       array(i)%p => node
       node => node%next
    end do
    s = 1
    ! Merge sort the array
    do while (s < count)
       i = 0
       i1 = 1
       i2 = s
       do while (i2 < count)
          i3 = min (s*(i+2), count)
          array(i1:i3) = merge (array(i1:i2), array(i2+1:i3))
          i = i + 2
          i1 = s*i+1
          i2 = s*(i+1)
       end do
       s = s * 2
    end do
    ! Relink according to their new order
    list => array(1)%p
    nullify (array(count)%p%next)
    node => list
    do i = 2, count
       node%next => array(i)%p
       node => node%next
    end do

  contains

    ! .le. comparision
    function lt (n1, n2) result (predicate)
      type(blha_cfg_process_node_t), intent(in) :: n1, n2
      logical :: predicate
      integer :: i
      predicate = .true.
      do i = 1, size (n1%fingerprint)
         if (n1%fingerprint(i) < n2%fingerprint(i)) return
         if (n1%fingerprint(i) > n2%fingerprint(i)) then
            predicate = .false.
            return
         end if
      end do
    end function lt

    ! Sorting core --- merge two sorted chunks
    function merge (l1, l2) result (lo)
      type(pnode), dimension(:), intent(in) :: l1, l2
      type(pnode), dimension(size (l1) + size (l2)) :: lo
      integer :: i, i1, i2
      i1 = 1
      i2 = 1
      do i = 1, size (lo)
         if (i1 > size (l1)) then
            lo(i)%p => l2(i2)%p
            i2 = i2 + 1
         elseif (i2 > size (l2)) then
            lo(i)%p => l1(i1)%p
            i1 = i1 + 1
         elseif (lt (l1(i1)%p, l2(i2)%p)) then
            lo(i)%p => l1(i1)%p
            i1 = i1 + 1
         else
            lo(i)%p => l2(i2)%p
            i2 = i2 + 1
         end if
      end do
    end function merge

  end subroutine sort_processes
  
\eatline
\nwendcode{}\nwbegindocs{32}\nwdocspar
Create an array of massive particle indices, to be used by the 
"MassiveParticle"-statement of the order file.
\nwenddocs{}\nwbegincode{33}\moddef{BLHA config: procedures}\plusendmoddef
  subroutine blha_configuration_get_massive_particles &
             (cfg, massive, i_massive)
    type(blha_configuration_t), intent(in) :: cfg
    logical, intent(out) :: massive
    integer, intent(out), dimension(:), allocatable :: i_massive
    integer, parameter :: max_particles = 10
    integer, dimension(max_particles) :: i_massive_tmp
    integer, dimension(max_particles) :: checked
    type(blha_cfg_process_node_t), pointer :: current_process 
    integer :: k
    integer :: n_massive
    n_massive = 0; k = 1
    checked = 0
    if (associated (cfg%processes)) then
       current_process => cfg%processes
    else
       call msg_fatal ("BLHA, massive particles: " // &
                       "No processes allocated!")
    end if
    do
       call check_pdg_list (current_process%pdg_in)
       call check_pdg_list (current_process%pdg_out)
       if (k > max_particles) &
          call msg_fatal ("BLHA, massive particles: " // &
                          "Max. number of particles exceeded!")
       if (associated (current_process%next)) then
          current_process => current_process%next
       else
          exit
       end if
    end do       
    if (n_massive > 0) then
       allocate (i_massive (n_massive))
       i_massive = i_massive_tmp (1:n_massive)
       massive = .true.
    else
       massive = .false.
    end if
  contains
    subroutine check_pdg_list (pdg_list)
       integer, dimension(:), intent(in) :: pdg_list
       integer :: i, i_pdg
       type(flavor_t) :: flv
       do i = 1, size (pdg_list)
          i_pdg = abs (pdg_list(i))
          call flv%init (i_pdg, cfg%model)
          if (flv%get_mass () > 0._default) then 
             !!! Avoid duplicates in output
             if (.not. any (checked == i_pdg)) then
                i_massive_tmp(k) = i_pdg
                checked(k) = i_pdg  
                k=k+1
                n_massive=n_massive+1
             end if
          end if
       end do
    end subroutine check_pdg_list
  end subroutine blha_configuration_get_massive_particles

\eatline
\nwendcode{}\nwbegindocs{34}\nwdocspar
Append a process. This expands the flavor sum, sorts it and then
eliminates any duplicates.
\nwenddocs{}\nwbegincode{35}\moddef{BLHA config: public}\plusendmoddef
  public :: blha_configuration_append_process
\nwendcode{}\nwbegincode{36}\moddef{BLHA config: procedures}\plusendmoddef
  subroutine blha_configuration_append_process (cfg, pdg_in, pdg_out, nsub, ids)
    type(blha_configuration_t), intent(inout) :: cfg
    type(pdg_array_t), dimension(:), intent(in) :: pdg_in, pdg_out
    integer, optional, intent(in) :: nsub
    integer, optional, dimension(:), intent(in) :: ids
    type(blha_cfg_process_node_t), pointer :: root, node, tmp
    ! Multiindex for counting through the PDG numbers
    integer, dimension(size (pdg_in)) :: i_in
    integer, dimension(size (pdg_out)) :: i_out
    ! Handle the list of lists
    type :: ilist
       integer, dimension(:), allocatable :: i
    end type ilist
    type(ilist), dimension(size (pdg_in)) :: ilist_i
    type(ilist), dimension(size (pdg_out)) :: ilist_o
    integer :: i, j, nproc
    logical :: inc
    ! Extract PDGs into integer lists
    do i = 1, size (pdg_in)
       ilist_i(i)%i = pdg_in(i)
    end do
    do i = 1, size (pdg_out)
       ilist_o(i)%i = pdg_out(i)
    end do
    i_in = 1
    i_out = 1
    allocate (root)
    node => root
    ! Perform the expansion
    nproc = 0
    EXPAND: do
       ! Transfer the PDG selection...
       allocate (node%pdg_in(size (pdg_in)))
       allocate (node%pdg_out(size (pdg_out)))
       allocate (node%fingerprint (size (pdg_in) + size (pdg_out)))
       if (present (nsub)) node%nsub = nsub
       if (present (ids)) then
          allocate (node%ids(size (ids)))
          node%ids = ids
       end if
       forall (j=1:size(ilist_i)) &
          node%pdg_in(j) = ilist_i(j)%i(i_in(j))
       forall (j=1:size(ilist_o)) &
          node%pdg_out(j) = ilist_o(j)%i(i_out(j))
       node%fingerprint = [ node%pdg_in, sort (node%pdg_out) ]
       nproc = nproc + 1
       inc = .false.
       ! ... and increment the multiindex
       do j = 1, size (i_out)
          if (i_out(j) < size (ilist_o(j)%i)) then
             i_out(j) = i_out(j) + 1
             inc = .true.
             exit
          else
             i_out(j) = 1
          end if
       end do
       if (.not. inc) then
          do j = 1, size (i_in)
             if (i_in(j) < size (ilist_i(j)%i)) then
                i_in(j) = i_in(j) + 1
                inc = .true.
                exit
             else
                i_in(j) = 1
             end if
          end do
       end if
       if (.not. inc) exit EXPAND
       allocate (node%next)
       node => node%next
    end do EXPAND
    ! Do the sorting
    call sort_processes (root, nproc)
    ! Kill duplicates
    node => root
    do while (associated (node))
       if (.not. associated (node%next)) exit
       if (all (node%fingerprint == node%next%fingerprint)) then
          tmp => node%next%next
          deallocate (node%next)
          node%next => tmp
          nproc = nproc - 1
       else
          node => node%next
       end if
    end do
    ! Append the remaining list
    if (associated (cfg%processes)) then
       node => cfg%processes
       do while (associated (node%next))
          node => node%next
       end do
       node%next => root
    else
       cfg%processes => root
    end if
    cfg%n_proc = cfg%n_proc + nproc
    cfg%dirty = .true.
    
  end subroutine blha_configuration_append_process

\eatline
\nwendcode{}\nwbegindocs{37}\nwdocspar
\nwenddocs{}\nwbegincode{38}\moddef{BLHA config: public}\plusendmoddef
  public :: blha_configuration_append_processes
\nwendcode{}\nwbegincode{39}\moddef{BLHA config: procedures}\plusendmoddef
  subroutine blha_configuration_append_processes (cfg, flavor, amp_type)
    type(blha_configuration_t), intent(inout) :: cfg
    type(blha_flv_state_t), dimension(:), intent(in) :: flavor
    integer, dimension(:), intent(in), optional :: amp_type
    integer :: n_tot
    type(blha_cfg_process_node_t), pointer :: current_node
    integer :: i_process, i_flv
    integer, dimension(:), allocatable :: pdg_in, pdg_out
    integer, dimension(:), allocatable :: flavor_state
    integer :: proc_offset, n_proc_tot
    proc_offset = 0; n_proc_tot = 0
    do i_flv = 1, size (flavor)
       n_proc_tot = n_proc_tot + flavor(i_flv)%flv_mult
    end do
    if (.not. associated (cfg%processes)) &
      allocate (cfg%processes)
    current_node => cfg%processes
    do i_flv = 1, size (flavor)
       n_tot = size (flavor(i_flv)%flavors)
       allocate (pdg_in (2), pdg_out (n_tot - 2))
       allocate (flavor_state (n_tot))
       flavor_state = flavor(i_flv)%flavors
       do i_process = 1, flavor(i_flv)%flv_mult
          pdg_in = flavor_state (1:2)
          pdg_out = flavor_state (3:)
          current_node%pdg_in = pdg_in
          current_node%pdg_out = pdg_out
          current_node%amplitude_type = amp_type (proc_offset+i_process)
          if (proc_offset+i_process /= n_proc_tot) then
            allocate (current_node%next)
            current_node => current_node%next
          end if
          if (i_process == flavor(i_flv)%flv_mult) &
             proc_offset = proc_offset + flavor(i_flv)%flv_mult
       end do
       deallocate (pdg_in, pdg_out)
       deallocate (flavor_state)
    end do
  end subroutine blha_configuration_append_processes

\eatline
\nwendcode{}\nwbegindocs{40}Change parameter(s).
\nwenddocs{}\nwbegincode{41}\moddef{BLHA config: public}\plusendmoddef
  public :: blha_configuration_set
\nwendcode{}\nwbegincode{42}\moddef{BLHA config: procedures}\plusendmoddef
!  subroutine blha_configuration_set ( cfg, &
!       matrix_element_square_type_hel, matrix_element_square_type_hel_other, &
!       matrix_element_square_type_col, matrix_element_square_type_col_other, &
!       correction_type, correction_type_other, &
!       irreg, irreg_other, &
!       massive_particle_scheme, massive_particle_scheme_other, &
!       subtraction_mode, subtraction_mode_other, &
!       model_file, subdivide_subprocesses, alphas_power, alpha_power, &
!       operation_mode, operation_mode_other)
  subroutine blha_configuration_set (cfg, &
!       version, amplitude_type, correction_type, irreg, massive_particle_scheme, &
       version, correction_type, irreg, massive_particle_scheme, &
       model_file, alphas_power, alpha_power, ew_scheme, width_scheme, &
       accuracy, debug)
    type(blha_configuration_t), intent(inout) :: cfg
    integer, optional, intent(in) :: version
!    integer, optional, intent(in) :: amplitude_type
!    integer, optional, intent(in) :: matrix_element_square_type_hel
!    type(string_t), optional, intent(in) :: matrix_element_square_type_hel_other
!    integer, optional, intent(in) :: matrix_element_square_type_col
!    type(string_t), optional, intent(in) :: matrix_element_square_type_col_other
    integer, optional, intent(in) :: correction_type
!    type(string_t), optional, intent(in) :: correction_type_other
    integer, optional, intent(in) :: irreg
!    type(string_t), optional, intent(in) :: irreg_other
    integer, optional, intent(in) :: massive_particle_scheme
!    type(string_t), optional, intent(in) :: massive_particle_scheme_other
!    integer, optional, intent(in) :: subtraction_mode
!    type(string_t), optional, intent(in) :: subtraction_mode_other
    type(string_t), optional, intent(in) :: model_file
!    logical, optional, intent(in) :: subdivide_subprocesses
    integer, optional, intent(in) :: alphas_power, alpha_power
    integer, optional, intent(in) :: ew_scheme
    integer, optional, intent(in) :: width_scheme
    real(default), optional, intent(in) :: accuracy
    logical, optional, intent(in) :: debug
!    integer, intent(in), optional :: operation_mode
!    type(string_t), intent(in), optional :: operation_mode_other
!    if (present (matrix_element_square_type_hel)) &
!       cfg%matrix_element_square_type(1) = matrix_element_square_type_hel
!    if (present (matrix_element_square_type_hel_other)) &
!       cfg%matrix_element_square_type_other(1) = matrix_element_square_type_hel_other
!    if (present (matrix_element_square_type_col)) &
!       cfg%matrix_element_square_type(2) = matrix_element_square_type_col
!    if (present (matrix_element_square_type_col_other)) &
!       cfg%matrix_element_square_type_other(2) = matrix_element_square_type_col_other
    if (present (version)) &
       cfg%version = version
!    if (present (amplitude_type)) &
!       cfg%amplitude_type = amplitude_type
    if (present (correction_type)) &
       cfg%correction_type = correction_type
!    if (present (correction_type_other)) &
!       cfg%correction_type_other = correction_type_other
    if (present (irreg)) &
       cfg%irreg = irreg
!    if (present (irreg_other)) &
!       cfg%irreg_other = irreg_other
    if (present (massive_particle_scheme)) &
       cfg%massive_particle_scheme = massive_particle_scheme
!    if (present (massive_particle_scheme_other)) &
!       cfg%massive_particle_scheme_other = massive_particle_scheme_other
!    if (present (subtraction_mode)) &
!       cfg%subtraction_mode = subtraction_mode
!    if (present (subtraction_mode_other)) &
!       cfg%subtraction_mode_other = subtraction_mode_other
    if (present (model_file)) &
       cfg%model_file = model_file
!    if (present (subdivide_subprocesses)) &
!       cfg%subdivide_subprocesses = subdivide_subprocesses
    if (present (alphas_power)) &
       cfg%alphas_power = alphas_power
    if (present (alpha_power)) &
       cfg%alpha_power = alpha_power
!    if (present (operation_mode)) &
!       cfg%operation_mode = operation_mode
!    if (present (operation_mode_other)) &
!       cfg%operation_mode_other = operation_mode_other
    if (present (ew_scheme)) &
       cfg%ew_scheme = ew_scheme
    if (present (width_scheme)) &
       cfg%width_scheme = width_scheme
    if (present (accuracy)) &
       cfg%accuracy_target = accuracy
    if (present (debug)) &
       cfg%debug_unstable = debug
    cfg%dirty = .false.
  end subroutine blha_configuration_set

\eatline
\nwendcode{}\nwbegindocs{43}\nwdocspar
\nwenddocs{}\nwbegincode{44}\moddef{BLHA config: public}\plusendmoddef
  public :: blha_configuration_get_n_proc
\nwendcode{}\nwbegincode{45}\moddef{BLHA config: procedures}\plusendmoddef
  function blha_configuration_get_n_proc (cfg) result (n_proc)
    type(blha_configuration_t), intent(in) :: cfg
    integer :: n_proc
    n_proc = cfg%n_proc
  end function blha_configuration_get_n_proc

\eatline
\nwendcode{}\nwbegindocs{46}\nwdocspar
Write the BLHA file. Internal mode is intented for md5summing only. 
\nwenddocs{}\nwbegincode{47}\moddef{BLHA config: public}\plusendmoddef
  public :: blha_configuration_write
\nwendcode{}\nwbegincode{48}\moddef{BLHA config: procedures}\plusendmoddef
  subroutine blha_configuration_write (cfg, unit, internal)
    type(blha_configuration_t), intent(in) :: cfg
    integer, intent(in), optional :: unit
    logical, intent(in), optional :: internal
    integer :: u
    logical :: full
    type(string_t) :: buf
    type(blha_cfg_process_node_t), pointer :: node
    integer :: i
    character(3) :: pdg_char
    character(6) :: accuracy
    character(len=25), parameter :: pad = ""
    integer, dimension(:), allocatable :: i_massive
    logical :: massive_particles
    u = given_output_unit (unit)
    full = .true.; if (present (internal)) full = .not. internal
    if (full .and. cfg%dirty) call msg_bug ( &
       "BUG: attempted to write out a dirty BLHA configuration")
    if (full) then
       write (u,'(A)') "# BLHA order written by WHIZARD \LA{}Version\RA{}"
       write (u,'(A)')
    end if
    select case (cfg%mode)
       case (BLHA_MODE_GOSAM); buf = "GoSam"
       case default; buf = "vanilla"
    end select
    write (u,'(A)') "# BLHA interface mode: " // char (buf)
    write (u,'(A)') "# process: " // char (cfg%name)
    write (u,'(A)') "# model: " // char (cfg%model%get_name ())
!    if (full) then
!       write (u,'(A)')
!       write (u,'(A)') '#@WO MD5 "' // char (cfg%md5) // '"'
!       write (u,'(A)')
!    end if
!    if (all (cfg%matrix_element_square_type == BLHA_MEST_SUM)) then
!       buf = "CHsummed"
!    elseif (all (cfg%matrix_element_square_type == BLHA_MEST_AVG)) then
!       buf = "CHaveraged"
!    else
!       buf = (render_mest ("H", cfg%matrix_element_square_type(1), &
!             cfg%matrix_element_square_type_other(1)) // " ") // &
!          render_mest ("C", cfg%matrix_element_square_type(2), &
!             cfg%matrix_element_square_type_other(2))
!    end if
!    write (u,'(A25,A)') "MatrixElementSquareType" // pad, char (buf)
    select case (cfg%version)
       case (1); buf = "BLHA1"
       case (2); buf = "BLHA2"
    end select
    write (u, '(A25,A)') "InterfaceVersion " // pad, char (buf)
    select case (cfg%correction_type)
       case (BLHA_CT_QCD); buf = "QCD"
       case (BLHA_CT_EW); buf = "EW"
       case (BLHA_CT_QED); buf = "QED"
       case default; buf = cfg%correction_type_other
    end select
    write (u,'(A25,A)') "CorrectionType" // pad, char (buf)
    select case (cfg%irreg)
       case (BLHA_IRREG_CDR); buf = "CDR"
       case (BLHA_IRREG_DRED); buf = "DRED"
       case (BLHA_IRREG_THV); buf = "tHV"
       case (BLHA_IRREG_MREG); buf = "MassReg"
       case default; buf = cfg%irreg_other
    end select
    write (u,'(A25,A)') "IRregularisation" // pad, char (buf)
    select case (cfg%massive_particle_scheme)
       case (BLHA_MPS_ONSHELL); buf = "OnShell"
       case default; buf = cfg%massive_particle_scheme_other
    end select
    write (u,'(A25,A)') "MassiveParticleScheme" // pad, char (buf)
!    select case (cfg%subtraction_mode)
!       case (BLHA_SUBMODE_NONE); buf = "None"
!       case default; buf = cfg%subtraction_mode_other
!    end select
!    write (u,'(A25,A)') "IRsubtractionMethod" // pad, char (buf)
!    write (u,'(A25,A)') "ModelFile" // pad, char (cfg%model_file)
!    if (cfg%subdivide_subprocesses) then
!       write (u,'(A25,A)') "SubdivideSubprocesses" // pad, "yes"
!    else
!       write (u,'(A25,A)') "SubdivideSubprocess" // pad, "no"
!    end if
    if (cfg%alphas_power >= 0) write (u,'(A25,A)') &
       "AlphasPower" // pad, int2char (cfg%alphas_power)
    if (cfg%alpha_power >= 0) write (u,'(A25,A)') &
       "AlphaPower " // pad, int2char (cfg%alpha_power)
    select case (cfg%ew_scheme)
       case (BLHA_EW_GF); buf = "alphaGF"
       case (BLHA_EW_MZ); buf = "alphaMZ"
       case (BLHA_EW_MSBAR); buf = "alphaMSbar"
       case (BLHA_EW_0); buf = "alpha0"
       case (BLHA_EW_RUN); buf = "alphaRUN"
       case (BLHA_EW_DEFAULT); buf = "OLPDefined"
    end select
    write (u, '(A25, A)') "EWScheme " // pad, char (buf)
    call blha_configuration_get_massive_particles &
         (cfg, massive_particles, i_massive)
    if (massive_particles) then
       write (u, '(A16,10I2)') 'MassiveParticles', i_massive
    end if
    if (full) then
       write (u,'(A)')
       write (u,'(A)') "# Process definitions"
       write (u,'(A)')
    end if
!    if (cfg%accuracy_target /= 0) then
!      write (accuracy, '(f6.5)') cfg%accuracy_target
!      write (u, '(A25,A)') "AccuracyTarget " // pad , accuracy 
!    end if
    if (cfg%debug_unstable) then
      buf = "True"
    else
      buf = "False"
    end if
    write (u, '(A25,A)') "DebugUnstable " // pad, char (buf)
    write (u, *)
    node => cfg%processes
    do while (associated (node))
       select case (node%amplitude_type)
         case (BLHA_AMP_LOOP); buf = "Loop"
         case (BLHA_AMP_CC); buf = "ccTree"
         case (BLHA_AMP_SC); buf = "scTree"
         case (BLHA_AMP_TREE); buf = "Tree"
         case (BLHA_AMP_LOOPINDUCED); buf = "LoopInduced"
       end select
       write (u, '(A25, A)') "AmplitudeType " // pad, char (buf)

       buf = ""
       do i = 1, size (node%pdg_in)
          write (pdg_char,'(I3)') node%pdg_in(i)
          buf = (buf // pdg_char) // " "
       end do
       buf = buf // "-> "
       do i = 1, size (node%pdg_out)
          write (pdg_char,'(I3)') node%pdg_out(i)
          buf = (buf // pdg_char) // " "
       end do
       write (u,'(A)') char (trim (buf))
       write (u, *)
       node => node%next
    end do

  contains

    function render_mest (prefix, mest, other) result (tag)
      character, intent(in) :: prefix
      integer, intent(in) :: mest
      type(string_t), intent(in) :: other
      type(string_t) :: tag
      select case (mest)
      case (BLHA_MEST_AVG); tag = prefix // "averaged"
      case (BLHA_MEST_SUM); tag = prefix // "summed"
      case default; tag = other
      end select
    end function render_mest

  end subroutine blha_configuration_write

\eatline
\nwendcode{}\nwbegindocs{49}\nwdocspar
``Freeze'' the configuration by calculating the MD5 sum.
\nwenddocs{}\nwbegincode{50}\moddef{BLHA config: public}\plusendmoddef
  public :: blha_configuration_freeze
\nwendcode{}\nwbegincode{51}\moddef{BLHA config: procedures}\plusendmoddef
  subroutine blha_configuration_freeze (cfg)
    type(blha_configuration_t), intent(inout) :: cfg
    integer :: u
    if (.not. cfg%dirty) return
    call sort_processes (cfg%processes)
    u = free_unit ()
    open (unit=u, status="scratch", action="readwrite")
    call blha_configuration_write (cfg, u, internal=.true.)
    rewind (u)
    cfg%md5 = md5sum (u)
    cfg%dirty = .false.
    close (u)
  end subroutine blha_configuration_freeze

\eatline
\nwendcode{}\nwbegindocs{52}\nwdocspar
Read a contract file, again creating a {\tt{}blha{\char95}configuration{\char95}t}
object. 
\nwenddocs{}\nwbegincode{53}\moddef{BLHA config: public}\plusendmoddef
  public :: blha_read_contract
\nwendcode{}\nwbegincode{54}\moddef{BLHA config: interfaces}\endmoddef
  interface blha_read_contract
     module procedure blha_read_contract_unit, &
          blha_read_contract_file
  end interface blha_read_contract 
  
\nwendcode{}\nwbegincode{55}\moddef{BLHA config: procedures}\plusendmoddef
  subroutine blha_read_contract_file (cfg, ok, fname, success)
    type(blha_configuration_t), intent(inout) :: cfg
    logical, intent(out) :: ok
    type(string_t), intent(in) :: fname
    logical, intent(out), optional :: success
    integer :: u, stat
    u = free_unit ()
    open (u, file=char (fname), status="old", action="read", iostat=stat)
    if (stat /= 0) then
       if (present (success)) then
          success = .false.
          return
       else
          call msg_bug ('Unable to open contract file "' // char (fname) // '"')
       end if
    end if
    call blha_read_contract_unit (cfg, ok, u, success)
    close (u)
  end subroutine blha_read_contract_file

  subroutine blha_read_contract_unit (cfg, ok, u, success)
    type(blha_configuration_t), intent(inout) :: cfg
    logical, intent(out) :: ok
    integer, intent(in) :: u
    logical, intent(out), optional :: success
    type(stream_t) :: stream
    type(ifile_t) :: preprocessed
    type(lexer_t) :: lexer
    type(parse_tree_t) :: parse_tree
    type(string_t) :: md5
    call stream_init (stream, u)
    call contract_preprocess (stream, preprocessed)
    call stream_final (stream)
    call stream_init (stream, preprocessed)
    call blha_init_lexer (lexer)
    call lexer_assign_stream (lexer, stream)
    call parse_tree_init (parse_tree, syntax_blha_contract, lexer)
    call blha_transfer_contract (cfg, ok, parse_tree, success)
    call blha_configuration_write (cfg, internal=.true.)
    call lexer_final (lexer)
    call stream_final (stream)
    call ifile_final (preprocessed)
    if (ok) then
       md5 = cfg%md5
       call blha_configuration_freeze (cfg)
       if (char (trim (md5 )) /= "") then
          if (md5 /= cfg%md5) then
             call msg_warning ("BLHA contract does not match the recorded " &
                // "checksum --- this counts as an error!")
             ok = .false.
          end if
       else
          call msg_warning ("It seems the OLP scrubbed our checksum, unable " &
             // "to check contract consistency.")
       end if
    end if
  end subroutine blha_read_contract_unit

\eatline
\eatline
\nwendcode{}\nwbegindocs{56}\nwdocspar
Walk the parse tree and transfer the results to the
{\tt{}blha{\char95}configuration} object. The {\tt{}goto} is a poor man's
replacement for exceptions which would be an appropiate error handling  
mechanism here. 
\nwenddocs{}\nwbegincode{57}\moddef{BLHA config: procedures}\plusendmoddef
  subroutine blha_transfer_contract (cfg, ok, parse_tree, success)
    type(blha_configuration_t), intent(inout) :: cfg
    logical, intent(out) :: ok
    type(parse_tree_t), intent(in), target :: parse_tree
    logical, intent(out), optional :: success
    type(parse_node_t), pointer :: pn_root, pn_line, pn_request, &
         pn_result, pn_key, pn_opt, pn_state_in, pn_state_out, pn_pdg
    type(string_t) :: emsg
    integer :: nopt, i, nsub
    integer, dimension(:), allocatable :: ids
    logical, dimension(2) :: flags
    type(pdg_array_t), dimension(:), allocatable :: pdg_in, pdg_out
    ok = .true.
    pn_root => parse_tree_get_root_ptr (parse_tree)
    pn_line => parse_node_get_sub_ptr (pn_root)
    do while (associated (pn_line))
       pn_request => parse_node_get_sub_ptr (pn_line)
       if (.not. associated (pn_request)) cycle
       if (char (parse_node_get_rule_key (pn_request)) == "process") then
          pn_result => parse_node_get_sub_ptr (pn_line, 2)
          pn_state_in => parse_node_get_sub_ptr (pn_request, 1)
          pn_state_out => parse_node_get_sub_ptr (pn_request, 3)
          allocate (pdg_in (parse_node_get_n_sub (pn_state_in)))
          allocate (pdg_out (parse_node_get_n_sub (pn_state_out)))
          i = 1
          pn_pdg => parse_node_get_sub_ptr (pn_state_in)
          do while (associated (pn_pdg))
             pdg_in(i) = [get_int (pn_pdg)]
             pn_pdg => parse_node_get_next_ptr (pn_pdg)
             i = i + 1
          end do
          i = 1
          pn_pdg => parse_node_get_sub_ptr (pn_state_out)
          do while (associated (pn_pdg))
             pdg_out(i) = [get_int (pn_pdg)]
             pn_pdg => parse_node_get_next_ptr (pn_pdg)
             i = i + 1
          end do
          i = parse_node_get_n_sub (pn_result)
          emsg = "broken process line"
          if (i < 2) goto 10
          pn_opt => parse_node_get_sub_ptr (pn_result, 2)
          do while (associated (pn_opt))
             if (char (parse_node_get_rule_key (pn_opt)) == "string") then
                call msg_warning ("While reading the BLHA contract: " // &
                   'the OLP returned an error for a process: "' // &
                   char (parse_node_get_string (pn_opt)) // '"')
                ok = .false.
                return
             end if
             pn_opt => parse_node_get_next_ptr (pn_opt)
          end do
          pn_opt => parse_node_get_sub_ptr (pn_result, 2)
          nsub = get_int (pn_opt)
          if (nsub /= i - 2) goto 10
          allocate (ids(nsub))
          i = 1
          pn_opt => parse_node_get_next_ptr (pn_opt)
          do while (associated (pn_opt))
             ids(i) = get_int (pn_opt)
             pn_opt => parse_node_get_next_ptr (pn_opt)
          end do
          call blha_configuration_append_process (cfg, pdg_in, pdg_out, &
             nsub=nsub, ids=ids)
          deallocate (pdg_in, pdg_out, ids)
       else
          pn_result => parse_node_get_sub_ptr (parse_node_get_next_ptr (pn_request), 2)
          pn_key => parse_node_get_sub_ptr (pn_request)
          pn_opt => parse_node_get_next_ptr (pn_key)
          nopt = parse_node_get_n_sub (pn_request) - 1
          select case (char (parse_node_get_rule_key (pn_key)))
             case ("md5")
                cfg%md5 = parse_node_get_string (pn_opt)
             case ("modelfile")
                cfg%model_file = get_fname (pn_opt)
                call check_result (pn_result, "ModelFile")
             case ("irregularisation")
                select case (lower_case (char (parse_node_get_string (pn_opt))))
                   case ("cdr"); cfg%irreg = BLHA_IRREG_CDR
                   case ("dred"); cfg%irreg = BLHA_IRREG_DRED
                   case ("thv"); cfg%irreg = BLHA_IRREG_THV
                   case ("mreg"); cfg%irreg = BLHA_IRREG_MREG
                   case default
                      cfg%irreg = BLHA_IRREG_OTHER 
                      cfg%irreg_other = parse_node_get_string (pn_opt)
                end select
                call check_result (pn_result, "IRRegularisation")
             case ("irsubtractionmethod")
                select case (lower_case (char (parse_node_get_string (pn_opt))))
                   case ("none"); cfg%subtraction_mode = BLHA_SUBMODE_NONE
                   case default
                      cfg%subtraction_mode = BLHA_SUBMODE_OTHER
                      cfg%subtraction_mode_other = parse_node_get_string(pn_opt)
                end select
                call check_result (pn_result, "IRSubtractionMethod")
             case ("massiveparticlescheme")
                select case (lower_case (char (parse_node_get_string (pn_opt))))
                   case ("onshell")
                      cfg%massive_particle_scheme = BLHA_MPS_ONSHELL
                   case default
                      cfg%massive_particle_scheme = BLHA_MPS_OTHER
                      cfg%massive_particle_scheme_other = &
                         parse_node_get_string (pn_opt)
                end select
                call check_result (pn_result, "MassiveParticleScheme")
             case ("matrixelementsquaretype")
                select case (nopt)
                   case (1)
                      select case (lower_case (char (parse_node_get_string (pn_opt))))
                         case ("chsummed")
                            cfg%matrix_element_square_type = BLHA_MEST_SUM
                         case ("chaveraged")
                            cfg%matrix_element_square_type = BLHA_MEST_AVG
                         case default
                            emsg = "invalid MatrixElementSquareType: " // &
                               parse_node_get_string (pn_opt)
                            goto 10
                      end select
                   case (2)
                      do i = 1, 2
                         pn_opt => parse_node_get_next_ptr (pn_key, i)
                         select case (lower_case (char (parse_node_get_string ( &
                               pn_opt))))
                            case ("csummed")
                               cfg%matrix_element_square_type(2) = BLHA_MEST_SUM
                               flags(2) = .true.
                            case ("caveraged")
                               cfg%matrix_element_square_type(2) = BLHA_MEST_AVG
                               flags(2) = .true.
                            case ("hsummed")
                               cfg%matrix_element_square_type(1) = BLHA_MEST_SUM
                               flags(1) = .true.
                            case ("haveraged")
                               cfg%matrix_element_square_type(1) = BLHA_MEST_AVG
                               flags(1) = .true.
                            case default
                               emsg = "invalid MatrixElementSquareType: " // &
                                  parse_node_get_string (pn_opt)
                               goto 10
                         end select
                      end do
                      if (.not. all (flags)) then
                         emsg = "MatrixElementSquareType: setup not exhaustive"
                         goto 10
                      end if
                   case default
                      emsg = "MatrixElementSquareType: too many options"
                      goto 10
                end select
                call check_result (pn_result, "MatrixElementSquareType")
             case ("correctiontype")
                select case (lower_case (char (parse_node_get_string (pn_opt))))
                   case ("qcd"); cfg%correction_type = BLHA_CT_QCD
                   case ("qed"); cfg%correction_type = BLHA_CT_QED
                   case ("ew"); cfg%correction_type = BLHA_CT_EW
                   case default
                      cfg%correction_type = BLHA_CT_OTHER
                      cfg%correction_type_other = parse_node_get_string (pn_opt)
                end select
                call check_result (pn_result, "CorrectionType")
             case ("alphaspower")
                cfg%alphas_power = get_int (pn_opt)
                call check_result (pn_result, "AlphasPower")
             case ("alphapower")
                cfg%alpha_power = get_int (pn_opt)
                call check_result (pn_result, "AlphaPower")
             case ("subdividesubprocess")
                select case (lower_case (char (parse_node_get_string (pn_opt))))
                   case ("yes"); cfg%subdivide_subprocesses = .true.
                   case ("no"); cfg%subdivide_subprocesses = .false.
                   case default
                      emsg = 'SubdivideSubprocess: invalid argument "' // &
                         parse_node_get_string (pn_opt) // '"'
                      goto 10
                end select
                call check_result (pn_result, "SubdivideSubprocess")
             case default
                emsg = "unknown statement: " // parse_node_get_rule_key (pn_key)
                goto 10
          end select
       end if
       pn_line => parse_node_get_next_ptr (pn_line)
    end do
    if (present (success)) success = .true.
    return
10  continue
    if (present (success)) then
       call msg_error ("Error reading BLHA contract: " // char (emsg))
       success = .false.
       return
    else
       call msg_fatal ("Error reading BLHA contract: " // char (emsg))
    end if

  contains

    function get_int (pn) result (i)
      type(parse_node_t), pointer :: pn
      integer :: i
      if (char (parse_node_get_rule_key (pn)) == "integer") then
         i = parse_node_get_integer (pn)
      else
         i = parse_node_get_integer (parse_node_get_sub_ptr (pn, 2))
         if (char (parse_node_get_rule_key (parse_node_get_sub_ptr (pn))) &
              == "-") i = -i
      end if
    end function get_int

    subroutine check_result (pn, step)
      type(parse_node_t), pointer :: pn
      character(*), intent(in) :: step
      type(string_t) :: res
      res = parse_node_get_string (pn)
      if (char (trim (res)) == "") then
         call msg_warning ("BLHA contract file: " // step // &
              ": OLP didn't return a status --- assuming an error")
         ok = .false.
      elseif (char (upper_case (res)) /= "OK") then
         call msg_warning ("BLHA contract file: " // step // &
              ': OLP error "' // char (res) // '"')
         ok = .false.
      end if
    end subroutine check_result

    function get_fname (pn) result (fname)
      type(parse_node_t), pointer :: pn
      type(string_t) :: fname
      type(parse_node_t), pointer :: pn_component
      if (char (parse_node_get_rule_key (pn)) == "string") then
         fname = parse_node_get_string (pn)
      else
         fname = ""
         pn_component => parse_node_get_sub_ptr (pn)
         do while (associated (pn_component))
            if (char (parse_node_get_rule_key (pn_component)) == "id") then
               fname = fname // parse_node_get_string (pn_component)
            else
               fname = fname // parse_node_get_key (pn_component)
            end if
            pn_component => parse_node_get_next_ptr (pn_component)
         end do
      end if
    end function get_fname

  end subroutine blha_transfer_contract

\eatline
\nwendcode{}\nwbegindocs{58}\nwdocspar
Initialize the lexer.
\nwenddocs{}\nwbegincode{59}\moddef{BLHA config: procedures}\plusendmoddef
  subroutine blha_init_lexer (lexer)
    type(lexer_t), intent(inout) :: lexer
    call lexer_init (lexer, &
       comment_chars = "#", &
       quote_chars = '"', &
       quote_match = '"', &
       single_chars = '\{\}|./\\:', &
       special_class = ["->"], &
       keyword_list = syntax_get_keyword_list_ptr (syntax_blha_contract), &
       upper_case_keywords = .false. &
       ) 
  end subroutine blha_init_lexer

\eatline
\nwendcode{}\nwbegindocs{60}\nwdocspar
Define the parser syntax table.
\nwenddocs{}\nwbegincode{61}\moddef{BLHA config: variables}\endmoddef
  type(syntax_t), target, save :: syntax_blha_contract
\nwendcode{}\nwbegincode{62}\moddef{BLHA config: public}\plusendmoddef
  public :: syntax_blha_contract_init
\nwendcode{}\nwbegincode{63}\moddef{BLHA config: procedures}\plusendmoddef
  subroutine syntax_blha_contract_init ()
    type(ifile_t) :: ifile
    call ifile_append (ifile, "SEQ contract = line*")
    call ifile_append (ifile, "KEY '->'")
    call ifile_append (ifile, "KEY '.'")
    call ifile_append (ifile, "KEY '/'")
    call ifile_append (ifile, "KEY '\\'")
    call ifile_append (ifile, "KEY '+'")
    call ifile_append (ifile, "KEY '-'")
    call ifile_append (ifile, "KEY '|'")
    call ifile_append (ifile, "KEY ':'")
    call ifile_append (ifile, "IDE id")
    call ifile_append (ifile, "INT integer")
    call ifile_append (ifile, "ALT sign = '+' | '-'")
    call ifile_append (ifile, "SEQ signed_integer = sign integer")
    call ifile_append (ifile, "QUO string = '""'...'""'")
    call ifile_append (ifile, "GRO line = '\{' line_contents '\}'")
    call ifile_append (ifile, "SEQ line_contents = request result?")
    call ifile_append (ifile, "ALT request = definition | process")
    call ifile_append (ifile, "ALT definition = option_unary | option_nary | " &
       // "option_path | option_numeric")
    call ifile_append (ifile, "KEY matrixelementsquaretype")
    call ifile_append (ifile, "KEY irregularisation")
    call ifile_append (ifile, "KEY massiveparticlescheme")
    call ifile_append (ifile, "KEY irsubtractionmethod")
    call ifile_append (ifile, "KEY modelfile")
    call ifile_append (ifile, "KEY operationmode")
    call ifile_append (ifile, "KEY subdividesubprocess")
    call ifile_append (ifile, "KEY alphaspower")
    call ifile_append (ifile, "KEY alphapower")
    call ifile_append (ifile, "KEY correctiontype")
    call ifile_append (ifile, "KEY md5")
    call ifile_append (ifile, "SEQ option_unary = key_unary arg")
    call ifile_append (ifile, "SEQ option_nary = key_nary arg+")
    call ifile_append (ifile, "SEQ option_path = key_path arg_path")
    call ifile_append (ifile, "SEQ option_numeric = key_numeric arg_numeric")
    call ifile_append (ifile, "ALT key_unary = irregularisation | " &
       // "massiveparticlescheme | irsubtractionmethod | subdividesubprocess | " &
       // "correctiontype | md5")
    call ifile_append (ifile, "ALT key_nary = matrixelementsquaretype | " &
       // "operationmode")
    call ifile_append (ifile, "ALT key_numeric = alphaspower | alphapower")
    call ifile_append (ifile, "ALT key_path = modelfile")
    call ifile_append (ifile, "ALT arg = id | string")
    call ifile_append (ifile, "ALT arg_numeric = integer | signed_integer")
    call ifile_append (ifile, "ALT arg_path = filename | string")
    call ifile_append (ifile, "SEQ filename = filename_atom+")
    call ifile_append (ifile, "ALT filename_atom = id | '.' | '/' | '\\' | ':'")
    call ifile_append (ifile, "SEQ process = state '->' state")
    call ifile_append (ifile, "SEQ state = pdg+")
    call ifile_append (ifile, "ALT pdg = integer | signed_integer")
    call ifile_append (ifile, "SEQ result = '|' result_atom+")
    call ifile_append (ifile, "ALT result_atom = integer | string")
    call syntax_init (syntax_blha_contract, ifile)
    call ifile_final (ifile)
  end subroutine syntax_blha_contract_init

\eatline
\nwendcode{}\nwbegindocs{64}\nwdocspar
\nwenddocs{}\nwbegincode{65}\moddef{BLHA config: public}\plusendmoddef
  public :: syntax_blha_contract_final
\nwendcode{}\nwbegincode{66}\moddef{BLHA config: procedures}\plusendmoddef
  subroutine syntax_blha_contract_final
    call syntax_final (syntax_blha_contract)
  end subroutine syntax_blha_contract_final

\eatline
\nwendcode{}\nwbegindocs{67}\nwdocspar
As the contract file is line-oriented, we apply a preprocessing step which
reformats the file in a way suitable for our free-form parser.
\nwenddocs{}\nwbegincode{68}\moddef{BLHA config: procedures}\plusendmoddef
  subroutine contract_preprocess (stream, ifile)
    type(stream_t), intent(inout) :: stream
    type(ifile_t), intent(out) :: ifile
    type(string_t) :: buf, reg, transformed
    integer :: stat, n
    buf = ""
    LINES: do
       call stream_get_record (stream, reg, stat)
       select case (stat)
          case (0)
          case (EOF); exit LINES
          case default
             call msg_bug ("I/O error while reading BLHA contract file")
       end select
       buf = buf // trim (reg)
       ! Take care of continuation lines
       if (char (extract (buf, len (buf), len(buf))) == '&') then
          buf = extract (buf, 1, len (buf) - 1) // " "
          cycle LINES
       end if
       buf = adjustl (buf)
       ! Transform #@WO comments into ordinary statements
       if (char (extract (buf, 1, 4)) == "#@WO") &
          buf = extract (buf, 5)
       ! Kill comments and blank lines
       if ((char (trim (buf)) == "") .or. &
          (char (extract (buf, 1, 1)) == "#")) then
             buf = ""
             cycle LINES
          end if
       ! Chop off any end-of-line comments
       call split (buf, reg, "#")
       ! Split line into order and result
       call split (reg, buf, "|")
       reg = trim (adjustl (reg))
       buf = trim (adjustl (buf))
       ! Check whether the order is a process definition
       n = scan (buf, ">")
       if (n == 0) then
          ! No -> quote result
          reg = ('"' // reg) // '"'
       else
          ! Yes -> leave any numbers as they are, quote any leftovers
          n = scan (reg, "0123456789", back=.true.)
          if (n < len (reg)) &
             reg = char (extract (reg, 1, n)) // ' "' // &
                char (trim (adjustl (extract (reg, n+1)))) // '"'
       end if
       ! Enclose the line into curly brackets
       transformed = "\{" // char (buf) // " | " // char (reg) // "\}"
       call ifile_append (ifile, transformed)
       buf = ""
    end do LINES
  end subroutine contract_preprocess

\eatline
\nwendcode{}\nwbegindocs{69}\nwdocspar
Test.
\nwenddocs{}\nwbegincode{70}\moddef{BLHA config: public}\plusendmoddef
  public :: blha_config_test
\nwendcode{}\nwbegincode{71}\moddef{BLHA config: procedures}\plusendmoddef
  subroutine blha_config_test (model, cfg, ok)
    type(pdg_array_t), dimension(2) :: pdg_in
    type(pdg_array_t), dimension(4) :: pdg_out
    class(model_data_t), pointer :: model
    type(blha_configuration_t), intent(out) :: cfg
    logical, intent(out) :: ok
    integer :: u
    logical :: flag
    ok = .false.
    pdg_in(1) = [1, 2, -1, -2]
    pdg_in(2) = pdg_in(1)
    pdg_out(1) = pdg_in(1)
    pdg_out(2) = [11]
    pdg_out(3) = [-11]
    pdg_out(4) = pdg_out(1)
    call blha_configuration_init (cfg, var_str ("test"), model)
    call blha_configuration_set (cfg, alphas_power = 2, alpha_power = 3)
    call blha_configuration_append_process (cfg, pdg_in, pdg_out)
    call blha_configuration_freeze (cfg)
    print *
    call blha_configuration_write (cfg)
    print *
    call blha_configuration_final (cfg)
    call blha_configuration_init (cfg, var_str ("test"), model, &
       mode=BLHA_MODE_GOSAM)
    call blha_configuration_set (cfg, alphas_power = 0, &
       model_file = var_str ("test.slha"))
    pdg_in(1) = [1]
    pdg_in(2) = [-1]
    pdg_out(1) = [22]
    pdg_out(2) = [22]   
    call blha_configuration_append_process (cfg, pdg_in, pdg_out(1:2))
    call blha_configuration_freeze (cfg)
    u = free_unit ()
    open (u, file="test.blha.order", action="write", status="replace")
    call blha_configuration_write (cfg, u)
    call blha_configuration_final (cfg)
    inquire (file="test.blha.contract", exist=flag)
    if (.not. flag) return
    call blha_configuration_init (cfg, var_str ("test"), model, mode=BLHA_MODE_GOSAM)
    call blha_read_contract (cfg, ok, var_str ("test.blha.contract"), success=flag)
    print *, "Reading back processed configuration: success? ", ok
  end subroutine blha_config_test

\eatline
\nwendcode{}\nwbegindocs{72}\nwdocspar

\section{OLP matrix element interface}

The prototypes of the OLP functions.
\nwenddocs{}\nwbegincode{73}\moddef{BLHA interface: interfaces}\endmoddef
  abstract interface
     subroutine ext_olp_start (file, status) bind(c)
       import
       character(c_char), dimension(*), intent(in) :: file
       integer(c_int), intent(out) :: status
     end subroutine ext_olp_Start

     subroutine ext_olp_evalsubprocess &
          (label, momenta, scale, parameters, amp) bind(c)
       import
       integer(c_int), intent(in), value :: label
       real(c_double), dimension(*), intent(in) :: momenta
       real(c_double), intent(in), value :: scale
       real(c_double), dimension(*), intent(in) :: parameters
       real(c_double), dimension(*), intent(out) :: amp
     end subroutine ext_olp_evalsubprocess

     subroutine ext_olp_finalize () bind(c)
     end subroutine ext_olp_finalize

     subroutine ext_olp_option (assignment, status) bind(c)
       import
       character(c_char), dimension(*), intent(in) :: assignment
       integer(c_int), intent(out) :: status
     end subroutine ext_olp_option
  end interface

\eatline
\eatline
\nwendcode{}\nwbegindocs{74}\nwdocspar
The OLP library is encapsulated together with the configuration in
derived type: 
\nwenddocs{}\nwbegincode{75}\moddef{BLHA interface: public}\endmoddef
  public :: blha_olp_t
\nwendcode{}\nwbegincode{76}\moddef{BLHA interface: types}\endmoddef
  type :: blha_olp_t
     private
     type(blha_configuration_t) :: cfg
     type(string_t) :: library
     integer :: n_in, n_out, n_flv, n_hel, n_col
     integer, dimension(:,:), allocatable :: flv_state
     logical :: color_summed = .true., flavor_summed = .true.
     logical :: loaded = .false.
     type(dlaccess_t) :: lib_handle
     procedure(ext_olp_start), pointer, nopass :: olp_start => null ()
     procedure(ext_olp_evalsubprocess), pointer, nopass :: &
        olp_evalsubprocess => null ()
     procedure(ext_olp_finalize), pointer, nopass :: olp_finalize => null ()
     procedure(ext_olp_option), pointer, nopass :: olp_option => null ()
  end type blha_olp_t

\eatline
\nwendcode{}\nwbegindocs{77}\nwdocspar
Init the {\tt{}blha{\char95}olp{\char95}t} object and try to dlopen the library. 
\nwenddocs{}\nwbegincode{78}\moddef{BLHA interface: public}\plusendmoddef
  public :: blha_olp_init
\nwendcode{}\nwbegincode{79}\moddef{BLHA interface: procedures}\endmoddef
  subroutine blha_olp_init (olp, cfg, library, success)
    type(blha_olp_t), intent(out) :: olp
    type(string_t), intent(in), optional :: library
    type(blha_configuration_t), intent(in) :: cfg
    logical, intent(out), optional :: success
    type(blha_cfg_process_node_t), pointer :: node
    type(string_t) :: prefix, libname
    type(c_funptr) :: fptr
    integer :: olp_status
    success = .true.
    node => cfg%processes
    if (.not. associated (node)) then
       call error ("blha_interface_init: empty process list")
       return
    end if
    olp%n_in = size (node%pdg_in)
    olp%n_out = size (node%pdg_out)
    do while (associated (node))
       if ((olp%n_in /= size (node%pdg_in)) .or. &
             (olp%n_out /= size (node%pdg_out))) then
          call error ("blha_interface_init: inconsistent process list")
          return
       end if
       node => node%next
    end do
    if (present (library)) then
       olp%library = library
    else
       olp%library = cfg%name // ".so"
    end if
    if (char (extract (olp%library, 1, 1)) == "/") then
       prefix = ""
       libname = extract (olp%library, 2)
    else
       prefix = "."
       libname = olp%library
    end if
    call dlaccess_init (olp%lib_handle, prefix, libname)
    if (dlaccess_has_error (olp%lib_handle)) then
       call error ("blha_interface_init: error opening library: " // &
          char (dlaccess_get_error (olp%lib_handle)))
       call dlaccess_final (olp%lib_handle)
       return
    end if
    fptr = dlaccess_get_c_funptr (olp%lib_handle, var_str ("OLP_Start"))
    if (.not. check_dlstate ()) return
    call c_f_procpointer (fptr, olp%olp_start)
    fptr = dlaccess_get_c_funptr (olp%lib_handle, var_str ("OLP_EvalSubProcess"))
    if (.not. check_dlstate ()) return
    call c_f_procpointer (fptr, olp%olp_evalsubprocess)
    if (olp%cfg%mode == BLHA_MODE_GOSAM) then
       fptr = dlaccess_get_c_funptr (olp%lib_handle, var_str ("OLP_Finalize"))
       if (.not. check_dlstate ()) return
       call c_f_procpointer (fptr, olp%olp_finalize)
       fptr = dlaccess_get_c_funptr (olp%lib_handle, var_str ("OLP_Option"))
       if (.not. check_dlstate ()) return
       call c_f_procpointer (fptr, olp%olp_option)
    end if
    call olp%olp_start (string_f2c (cfg%model_file), olp_status)
    if (olp_status /= 1) then
       call error ("blha_interface_init: OLP initialization failed")
       call dlaccess_final (olp%lib_handle)
    end if
    success = .true.
    olp%loaded = .true.

  contains

    function check_dlstate () result (ok)
      logical :: ok
      ok = .not. dlaccess_has_error (olp%lib_handle)
      if (.not. ok) then
         call error ("blha_interface_init: error loading library: " // &
              char (dlaccess_get_error (olp%lib_handle)))
         call dlaccess_final (olp%lib_handle)
      end if
    end function check_dlstate
   
    subroutine error (msg)
      character(*), intent(in) :: msg
      if (present (success)) then
         call msg_error (msg)
         success = .false.
      else
         call msg_fatal (msg)
      end if
    end subroutine error

  end subroutine blha_olp_init

\eatline
\nwendcode{}\nwbegindocs{80}Finalizer.
\nwenddocs{}\nwbegincode{81}\moddef{BLHA interface: public}\plusendmoddef
  public :: blha_olp_final
\nwendcode{}\nwbegincode{82}\moddef{BLHA interface: procedures}\plusendmoddef
  subroutine blha_olp_final (olp)
    type(blha_olp_t), intent(inout) :: olp
    if (.not. olp%loaded) return
    if (associated (olp%olp_finalize)) call olp%olp_finalize
    call dlaccess_final (olp%lib_handle)
    olp%loaded = .false.
  end subroutine blha_olp_final

\eatline
\nwendcode{}\nwbegindocs{83}Test.
\nwenddocs{}\nwbegincode{84}\moddef{BLHA interface: public}\plusendmoddef
  public :: blha_interface_test
\nwendcode{}\nwbegincode{85}\moddef{BLHA interface: procedures}\plusendmoddef
  subroutine blha_interface_test (cfg, ok)
    type(blha_configuration_t), intent(inout) :: cfg
    type(blha_olp_t) :: olp
    logical, intent(out) :: ok
    call blha_olp_init (olp, cfg, library=var_str ("blha_test.so"), success=ok)
    print *, "loading OLP library: success?", ok
    call blha_olp_final (olp)
  end subroutine blha_interface_test

\eatline
\nwendcode{}\nwbegindocs{86}\nwdocspar
\section{OLP driver}

\nwenddocs{}\nwbegincode{87}\moddef{BLHA driver: public}\endmoddef
  public :: blha_test
\nwendcode{}\nwbegincode{88}\moddef{BLHA driver: procedures}\endmoddef
  subroutine blha_test (model)
    class(model_data_t), pointer :: model
    type (blha_configuration_t) :: cfg
    logical :: ok
    call blha_config_test (model, cfg, ok)
    if (ok) call blha_interface_test (cfg, ok)
  end subroutine blha_test

\eatline
\nwendcode{}\nwfilename{../../../src/gosam/gosam.nw}\nwbegindocs{0}% -*- ess-noweb-default-code-mode: f90-mode; noweb-default-code-mode: f90-mode; -*- 
% WHIZARD code as NOWEB source: matrix elements and process libraries

\nwenddocs{}\nwbegindocs{1}\nwdocspar
%%%%%%%%%%%%%%%%%%%%%%%%%%%%%%%%%%%%%%%%%%%%%%%%%%%%%%%%%%%%%%%%%%%%%%%%
\chapter{GoSam Interface}

The code in this chapter makes amplitudes accessible to \whizard\ that
are generated and computed by the GoSam package.

These are the modules:
\begin{description}
\item[loop\_archive]
  Provide some useful extra functionality.
\item[prc\_gosam]
  The actual interface, following the \whizard\ conventions for
  matrix-element generator methods.
\end{description}

\clearpage
%%%%%%%%%%%%%%%%%%%%%%%%%%%%%%%%%%%%%%%%%%%%%%%%%%%%%%%%%%%%%%%%%%%%%%%%
\section{Save previously generated loop libraries}
\nwenddocs{}\nwbegincode{2}\moddef{\code{}loop{\char95}archive.f90\edoc{}}\endmoddef
\LA{}File header\RA{}

module loop_archive

  use io_units
\LA{}Use strings\RA{}
  use diagnostics
  use os_interface

\LA{}Loop archive: public\RA{}

\LA{}Loop archive: types\RA{}

contains

\LA{}Loop archive: procedures\RA{}

end module loop_archive
\eatline
\nwendcode{}\nwbegindocs{3}\nwdocspar
\nwenddocs{}\nwbegincode{4}\moddef{Loop archive: public}\endmoddef
  public :: loop_archive_t
\nwendcode{}\nwbegincode{5}\moddef{Loop archive: types}\endmoddef
  type :: loop_archive_t
    logical :: active = .false.
    type(string_t) :: name
    type(string_t) :: current_prefix
  contains
  \LA{}Loop archive: loop archive: TBP\RA{}
  end type loop_archive_t

\eatline
\nwendcode{}\nwbegindocs{6}\nwdocspar
\nwenddocs{}\nwbegincode{7}\moddef{Loop archive: loop archive: TBP}\endmoddef
  procedure :: activate => loop_archive_activate
\nwendcode{}\nwbegincode{8}\moddef{Loop archive: procedures}\endmoddef
  subroutine loop_archive_activate (archive, name)
    class(loop_archive_t), intent(inout) :: archive
    type(string_t), intent(in) :: name
    integer :: status, success
    type(string_t) :: prefix
    archive%name = name
    call os_system_call ('test -d "' // name // &
                         '"', status = status, verbose = .true.)
    if (status /= 0) then
      call os_system_call ('mkdir ' // name, &
                            status = success, verbose = .true.)
      if (success /= 0) call msg_fatal ("Creation of loop archive failed!")
    end if
    archive%active = .true.
  end subroutine loop_archive_activate

\eatline
\nwendcode{}\nwbegindocs{9}\nwdocspar
\nwenddocs{}\nwbegincode{10}\moddef{Loop archive: loop archive: TBP}\plusendmoddef
  procedure :: record => loop_archive_record
\nwendcode{}\nwbegincode{11}\moddef{Loop archive: procedures}\plusendmoddef
  subroutine loop_archive_record (archive, olp_file, config_file, lib)
    class(loop_archive_t), intent(inout) :: archive
    type(string_t), intent(in) :: olp_file, config_file, lib
    type(string_t) :: current_prefix
    type(string_t) :: filename
    
! Copy, rename and move olp-file
    filename = archive%current_prefix // '.olp'
    call os_system_call ('cp ' // olp_file // ' ' // &
                         filename)
    call os_system_call ('mv ' // filename // ' ' // &
                         archive%name)
! Do the same with the loop-library and the config file
    filename = archive%current_prefix // '_libgolem_olp.so'
    call os_system_call ('cp ' // lib // ' ' // filename)
    call os_system_call ('mv ' // filename // ' ' // &
                         archive%name)
    filename = archive%current_prefix // '_golem.in'
    call os_system_call ('cp ' // config_file // ' ' // &
                         filename)
    call os_system_call ('mv ' // filename // ' ' // &
                         archive%name)

  end subroutine loop_archive_record

\eatline
\nwendcode{}\nwbegindocs{12}\nwdocspar
\nwenddocs{}\nwbegincode{13}\moddef{Loop archive: loop archive: TBP}\plusendmoddef
  procedure :: search => loop_archive_search
\nwendcode{}\nwbegincode{14}\moddef{Loop archive: procedures}\plusendmoddef
  subroutine loop_archive_search (archive, files, found)
    class(loop_archive_t), intent(inout) :: archive
    type(string_t), dimension(3), intent(in) :: files
    logical, intent(out) :: found
    type(string_t) :: current_olp, current_config, current_lib
    character(len=3) :: prefix
    integer :: counter
    logical, dimension(2) :: exist
    integer, dimension(2) :: identical
    integer :: i

    counter = 1
    do
       write(prefix,"(A,I2.2)") 'V', counter
       current_olp = archive%name // '/' // var_str (prefix) // '.olp'
       current_config = archive%name // '/' // var_str (prefix) // '_golem.in'
       current_lib = archive%name // '/' // var_str (prefix) // '_libgolem_olp.so'
       inquire (file = char (current_olp), exist = exist(1))
       inquire (file = char (current_config), exist = exist(2))
       if (all (exist)) then
          call os_system_call ('diff ' // current_olp // ' ' // files(1) // &
                               ' > /dev/null', status=identical(1))
          call os_system_call ('diff ' // current_config // ' ' // files(2) // &
                               ' > /dev/null', status=identical(2))
            if (all (identical == 0)) then
               found = .true.
               exit
            else
               counter = counter+1
            end if
       else
          found = .false.
          exit
       end if
       if (counter >= 100) call msg_fatal ("Maximum number of loop-libraries exceeded!")
    end do
    write(prefix,"(A,I2.2)") 'V', counter
    archive%current_prefix = var_str (prefix)    
  end subroutine loop_archive_search

\eatline
\nwendcode{}\nwbegindocs{15}\nwdocspar
\nwenddocs{}\nwbegincode{16}\moddef{Loop archive: loop archive: TBP}\plusendmoddef
  procedure :: restore => loop_archive_restore
\nwendcode{}\nwbegincode{17}\moddef{Loop archive: procedures}\plusendmoddef
  subroutine loop_archive_restore (archive, olp_orig, path)
    class(loop_archive_t), intent(inout) :: archive
    type(string_t), intent(in) :: olp_orig, path
    type(string_t) :: olp_file, config_file, lib
    
    olp_file = archive%current_prefix // '.olp'
    config_file = archive%current_prefix // '_golem.in'
    lib = archive%current_prefix // '_libgolem_olp.so'

    call os_system_call ('cp ' // archive%name // '/' // olp_file // ' .')
    call os_system_call ('mv ' // olp_file // ' ' // olp_orig)
    
    call os_system_call ('cp ' // archive%name // '/' // config_file // ' .')
    call os_system_call ('mv ' // config_file // ' golem.in')

    call os_system_call ('cp ' // archive%name // '/' // lib // ' .')
    call os_system_call ('mv ' // lib // ' libgolem_olp.so')
    call os_system_call ('cp libgolem_olp.so ' // path)
    
  end subroutine loop_archive_restore

\eatline
\nwendcode{}\nwbegindocs{18}\nwdocspar
\clearpage
%%%%%%%%%%%%%%%%%%%%%%%%%%%%%%%%%%%%%%%%%%%%%%%%%%%%%%%%%%%%%%%%%%%%%%%%
\section{Gosam Interface}
\nwenddocs{}\nwbegincode{19}\moddef{\code{}prc{\char95}gosam.f90\edoc{}}\endmoddef
\LA{}File header\RA{}

module prc_gosam
  
  use, intrinsic :: iso_c_binding !NODEP!
  use, intrinsic :: iso_fortran_env

  use kinds
\LA{}Use strings\RA{}  
  use io_units
  use constants
  use system_defs, only: TAB
  use system_dependencies
  use file_utils
  use string_utils
  use physics_defs
  use diagnostics
  use os_interface
  use lorentz
  use interactions
  use pdg_arrays
  use sm_qcd
  use flavors
  use model_data
  
  use process_constants
  use prclib_interfaces
  use prc_core_def
  use prc_core

  use blha_config
  use loop_archive

\LA{}Standard module head\RA{}

\LA{}prc gosam: constants\RA{}

\LA{}prc gosam: public\RA{}

\LA{}prc gosam: types\RA{}

\LA{}prc gosam: interfaces\RA{}

contains

\LA{}prc gosam: procedures\RA{}

end module prc_gosam

\nwendcode{}\nwbegindocs{20}\nwdocspar
%

\nwenddocs{}\nwbegincode{21}\moddef{prc gosam: constants}\endmoddef
  character(10), dimension(5), parameter ::  &
             lib_suffix = [character(10) :: &
             '.a', '.la', '.so', '.so.0', '.so.0.0.0']

\nwendcode{}\nwbegindocs{22}The GoSam interface requires the momentum and result array to be of a fixed length.
\nwenddocs{}\nwbegincode{23}\moddef{prc gosam: constants}\plusendmoddef
  integer, parameter :: GOSAM_PARAMETER_LIMIT = 10
  integer, parameter :: GOSAM_MOMENTUM_LIMIT = 50
  integer, parameter :: GOSAM_RESULTS_LIMIT = 60

\nwendcode{}\nwbegindocs{24}\nwdocspar
\nwenddocs{}\nwbegincode{25}\moddef{prc gosam: constants}\plusendmoddef
  integer, parameter :: Q_TO_QG = 1
  integer, parameter :: G_TO_GG = 2
  integer, parameter :: G_TO_QQ = 3

\nwendcode{}\nwbegindocs{26}\nwdocspar
\nwenddocs{}\nwbegincode{27}\moddef{prc gosam: public}\endmoddef
  public :: gosam_writer_t
\nwendcode{}\nwbegincode{28}\moddef{prc gosam: types}\endmoddef
  type, extends (prc_writer_f_module_t) :: gosam_writer_t
    type(string_t) :: gosam_dir
    type(string_t) :: golem_dir
    type(string_t) :: samurai_dir
    type(string_t) :: ninja_dir
    type(string_t) :: form_dir
    type(string_t) :: qgraf_dir
    type(blha_configuration_t) :: blha_cfg
    type(string_t) :: model_name
    type(string_t) :: process_mode
    type(string_t) :: process_string
  contains
  \LA{}prc gosam: gosam writer: TBP\RA{}
  end type gosam_writer_t

\nwendcode{}\nwbegindocs{29}\nwdocspar
%

\nwenddocs{}\nwbegincode{30}\moddef{prc gosam: public}\plusendmoddef
  public :: gosam_def_t
\nwendcode{}\nwbegincode{31}\moddef{prc gosam: types}\plusendmoddef
  type, extends (prc_core_def_t) :: gosam_def_t
    type(string_t) :: basename
    type(string_t) :: suffix
    logical :: execute_olp = .true.
  contains
  \LA{}prc gosam: gosam def: TBP\RA{}
  end type gosam_def_t

\nwendcode{}\nwbegindocs{32}\nwdocspar
%
\nwenddocs{}\nwbegincode{33}\moddef{prc gosam: types}\plusendmoddef
  type, extends (prc_core_driver_t) :: gosam_driver_t 
    type(string_t) :: gosam_dir
    type(string_t) :: olp_file
    type(string_t) :: olc_file
    type(string_t) :: olp_dir
    type(loop_archive_t) :: loop_archive
    procedure(olp_start),nopass,  pointer :: &
              gosam_olp_start => null ()
    procedure(olp_eval), nopass, pointer :: &
              gosam_olp_eval => null()
    procedure(olp_info), nopass, pointer :: &
              gosam_olp_info => null ()
    procedure(olp_set_parameter), nopass, pointer :: &
              gosam_olp_set_parameter => null ()
    procedure(olp_eval2), nopass, pointer :: &
              gosam_olp_eval2 => null ()
    procedure(olp_option), nopass, pointer :: &
              gosam_olp_option => null ()
    procedure(olp_polvec), nopass, pointer :: &
              gosam_olp_polvec => null ()
    procedure(olp_finalize), nopass, pointer :: &
              gosam_olp_finalize => null ()
    procedure(olp_print_parameter), nopass, pointer :: &
              gosam_olp_print_parameter => null ()
\nwendcode{}\nwbegindocs{34}We have to add two O'Mega-routines to the GoSam driver to ensure
  proper process setup.
\nwenddocs{}\nwbegincode{35}\moddef{prc gosam: types}\plusendmoddef
    procedure(omega_update_alpha_s), nopass, pointer :: &
              update_alpha_s => null ()
    procedure(omega_is_allowed), nopass, pointer :: &
              is_allowed => null ()

  contains
  \LA{}prc gosam: gosam driver: TBP\RA{}
  end type gosam_driver_t

\nwendcode{}\nwbegindocs{36}\nwdocspar
%
\nwenddocs{}\nwbegincode{37}\moddef{prc gosam: public}\plusendmoddef
  public :: prc_gosam_t
\nwendcode{}\nwbegincode{38}\moddef{prc gosam: types}\plusendmoddef
  type, extends (prc_core_t) :: prc_gosam_t
    type(qcd_t) :: qcd
    integer :: n_flv
    real(default) :: maximum_accuracy = 10000.0
    logical :: initialized = .false.
    integer, dimension(:), allocatable :: i_born, i_sc, i_cc
    integer, dimension(:), allocatable :: i_real
    integer, dimension(:), allocatable :: i_virt
  contains
  \LA{}prc gosam: prc gosam: TBP\RA{}
  end type prc_gosam_t

\nwendcode{}\nwbegindocs{39}\nwdocspar
%
\nwenddocs{}\nwbegincode{40}\moddef{prc gosam: types}\plusendmoddef
  type, extends (workspace_t) :: gosam_state_t
    logical :: new_kinematics = .true.
    real(default) :: alpha_qcd = -1
  contains
  \LA{}prc gosam: gosam state: TBP\RA{}
  end type gosam_state_t

\eatline
\nwendcode{}\nwbegindocs{41}\nwdocspar
\nwenddocs{}\nwbegincode{42}\moddef{prc gosam: interfaces}\endmoddef
  interface 
    subroutine olp_start (contract_file_name, ierr) bind (C,name="OLP_Start")
      import
      character(kind=c_char, len=1), intent(in) :: contract_file_name
      integer(kind=c_int), intent(out) :: ierr
    end subroutine olp_start
  end interface

\eatline
\nwendcode{}\nwbegindocs{43}\nwdocspar
\nwenddocs{}\nwbegincode{44}\moddef{prc gosam: interfaces}\plusendmoddef
  interface
    subroutine olp_eval (label, momenta, mu, parameters, res) &
         bind (C,name="OLP_EvalSubProcess")
      import
      integer(kind=c_int), value, intent(in) :: label
      real(kind=c_double), value, intent(in) :: mu
      real(kind=c_double), dimension(GOSAM_MOMENTUM_LIMIT), intent(in) :: momenta
      real(kind=c_double), dimension(GOSAM_PARAMETER_LIMIT), intent(in) :: parameters
      real(kind=c_double), dimension(GOSAM_RESULTS_LIMIT), intent(out) :: res
    end subroutine olp_eval
  end interface

\eatline
\nwendcode{}\nwbegindocs{45}\nwdocspar
\nwenddocs{}\nwbegincode{46}\moddef{prc gosam: interfaces}\plusendmoddef
  interface
    subroutine olp_info (olp_file, olp_version, message) &
         bind(C,name="OLP_Info")
      import
      character(kind=c_char), intent(inout), dimension(15) :: olp_file
      character(kind=c_char), intent(inout), dimension(15) :: olp_version
      character(kind=c_char), intent(inout), dimension(255) :: message
    end subroutine olp_info
  end interface

\eatline
\nwendcode{}\nwbegindocs{47}\nwdocspar
\nwenddocs{}\nwbegincode{48}\moddef{prc gosam: interfaces}\plusendmoddef
  interface
    subroutine olp_set_parameter &
         (variable_name, real_part, complex_part, success) &
            bind(C,name="OLP_SetParameter")
      import
      character(kind=c_char,len=1), intent(in) :: variable_name
      real(kind=c_double), intent(in) :: real_part, complex_part
      integer(kind=c_int), intent(out) :: success
    end subroutine olp_set_parameter
  end interface

\eatline
\nwendcode{}\nwbegindocs{49}\nwdocspar
\nwenddocs{}\nwbegincode{50}\moddef{prc gosam: interfaces}\plusendmoddef
  interface
    subroutine olp_eval2 (label, momenta, mu, res, acc) &
         bind(C,name="OLP_EvalSubProcess2")
      import
      integer(kind=c_int), intent(in) :: label
      real(kind=c_double), intent(in) :: mu
      real(kind=c_double), dimension(GOSAM_MOMENTUM_LIMIT), intent(in) :: momenta
      real(kind=c_double), dimension(GOSAM_RESULTS_LIMIT), intent(out) :: res
      real(kind=c_double), intent(out) :: acc
    end subroutine olp_eval2
  end interface

\eatline
\nwendcode{}\nwbegindocs{51}\nwdocspar
\nwenddocs{}\nwbegincode{52}\moddef{prc gosam: interfaces}\plusendmoddef
  abstract interface
     subroutine omega_update_alpha_s (alpha_s) bind(C)
       import
       real(c_default_float), intent(in) :: alpha_s
     end subroutine omega_update_alpha_s
  end interface
  
  abstract interface
     subroutine omega_is_allowed (flv, hel, col, flag) bind(C)
       import
       integer(c_int), intent(in) :: flv, hel, col
       logical(c_bool), intent(out) :: flag
     end subroutine omega_is_allowed
  end interface

\eatline
\nwendcode{}\nwbegindocs{53}\nwdocspar
\nwenddocs{}\nwbegincode{54}\moddef{prc gosam: interfaces}\plusendmoddef
  interface
    subroutine olp_option (line, stat) bind(C,name="OLP_Option")
      import
      character(kind=c_char, len=1), intent(in) :: line
      integer(kind=c_int), intent(out) :: stat
    end subroutine
  end interface

\eatline
\nwendcode{}\nwbegindocs{55}\nwdocspar
\nwenddocs{}\nwbegincode{56}\moddef{prc gosam: interfaces}\plusendmoddef
  interface
    subroutine olp_polvec (p, q, eps) bind(C,name="OLP_Polvec")
      import
      real(kind=c_double), dimension(0:3), intent(in) :: p, q
      real(kind=c_double), dimension(0:7), intent(out) :: eps
    end subroutine
  end interface

\eatline
\nwendcode{}\nwbegindocs{57}\nwdocspar
\nwenddocs{}\nwbegincode{58}\moddef{prc gosam: interfaces}\plusendmoddef
  interface
    subroutine olp_finalize () bind(C,name="OLP_Finalize")
      import
    end subroutine olp_finalize
  end interface

\eatline
\nwendcode{}\nwbegindocs{59}\nwdocspar
\nwenddocs{}\nwbegincode{60}\moddef{prc gosam: interfaces}\plusendmoddef
  interface
    subroutine olp_print_parameter (filename) bind(C,name="OLP_PrintParameter")
      import
      character(kind=c_char, len=1), intent(in) :: filename
    end subroutine olp_print_parameter
  end interface

\eatline
\nwendcode{}\nwbegindocs{61}\nwdocspar
\nwenddocs{}\nwbegincode{62}\moddef{prc gosam: gosam def: TBP}\endmoddef
  procedure :: init => gosam_def_init
\nwendcode{}\nwbegincode{63}\moddef{prc gosam: procedures}\endmoddef
  subroutine gosam_def_init (object, basename, model_name, &
                             prt_in, prt_out, nlo_type)
    class(gosam_def_t), intent(inout) :: object
    type(string_t), intent(in) :: model_name
    type(string_t), intent(in) :: basename
    type(string_t), dimension(:), intent(in) :: prt_in, prt_out
    type(string_t), intent(in) :: nlo_type
    object%basename = basename
    allocate (gosam_writer_t :: object%writer)
    select case (char (nlo_type))
    case ('Real')
       object%suffix = '_REAL'
    case ('Virtual')
       object%suffix = '_LOOP'
    case ('Subtraction')
       object%suffix = '_SUB'
    end select
    select type (writer => object%writer)
    type is (gosam_writer_t)
      call writer%init (model_name, prt_in, prt_out)
    end select
  end subroutine gosam_def_init

\eatline
\nwendcode{}\nwbegindocs{64}\nwdocspar
\nwenddocs{}\nwbegincode{65}\moddef{prc gosam: gosam def: TBP}\plusendmoddef
  procedure, nopass :: needs_code => gosam_def_needs_code
\nwendcode{}\nwbegincode{66}\moddef{prc gosam: procedures}\plusendmoddef
  function gosam_def_needs_code () result (flag)
    logical :: flag
    flag = .true.
  end function gosam_def_needs_code

\nwendcode{}\nwbegindocs{67} %def gosam_def_needs_code
\nwenddocs{}\nwbegindocs{68}\nwdocspar
\nwenddocs{}\nwbegincode{69}\moddef{prc gosam: gosam def: TBP}\plusendmoddef
  procedure, nopass :: type_string => gosam_def_type_string
\nwendcode{}\nwbegincode{70}\moddef{prc gosam: procedures}\plusendmoddef
  function gosam_def_type_string () result (string)
    type(string_t) :: string
    string = "gosam"
  end function gosam_def_type_string

\nwendcode{}\nwbegindocs{71}\nwdocspar
%
\nwenddocs{}\nwbegincode{72}\moddef{prc gosam: gosam def: TBP}\plusendmoddef
  procedure :: write => gosam_def_write
\nwendcode{}\nwbegincode{73}\moddef{prc gosam: procedures}\plusendmoddef
  subroutine gosam_def_write (object, unit)
    class(gosam_def_t), intent(in) :: object
    integer, intent(in) :: unit
    select type (writer => object%writer)
    type is (gosam_writer_t)
      call writer%write (unit)
    end select
  end subroutine gosam_def_write

\nwendcode{}\nwbegindocs{74}\nwdocspar
%
\nwenddocs{}\nwbegincode{75}\moddef{prc gosam: gosam def: TBP}\plusendmoddef
  procedure :: read => gosam_def_read
\nwendcode{}\nwbegincode{76}\moddef{prc gosam: procedures}\plusendmoddef
  subroutine gosam_def_read (object, unit)
    class(gosam_def_t), intent(out) :: object
    integer, intent(in) :: unit
    call msg_bug ("GoSam process definition: input not supported yet")
  end subroutine gosam_def_read


\eatline
\nwendcode{}\nwbegindocs{77}\nwdocspar
\nwenddocs{}\nwbegincode{78}\moddef{prc gosam: gosam def: TBP}\plusendmoddef
  procedure :: allocate_driver => gosam_def_allocate_driver
\nwendcode{}\nwbegincode{79}\moddef{prc gosam: procedures}\plusendmoddef
  subroutine gosam_def_allocate_driver (object, driver, basename)
    class(gosam_def_t), intent(in) :: object
    class(prc_core_driver_t), intent(out), allocatable :: driver
    type(string_t), intent(in) :: basename
    if (.not. allocated (driver)) allocate (gosam_driver_t :: driver)
  end subroutine gosam_def_allocate_driver

\nwendcode{}\nwbegindocs{80}\nwdocspar
%
\nwenddocs{}\nwbegincode{81}\moddef{prc gosam: gosam def: TBP}\plusendmoddef
  procedure, nopass :: get_features => gosam_def_get_features
\nwendcode{}\nwbegincode{82}\moddef{prc gosam: procedures}\plusendmoddef
  subroutine gosam_def_get_features (features)
    type(string_t), dimension(:), allocatable, intent(out) :: features
    allocate (features (6))
    features = [ &
         var_str ("init"), &
         var_str ("update_alpha_s"), &
         var_str ("reset_helicity_selection"), &
         var_str ("is_allowed"), &
         var_str ("new_event"), &
         var_str ("get_amplitude")]
  end subroutine gosam_def_get_features 

\nwendcode{}\nwbegindocs{83}\nwdocspar
%
\nwenddocs{}\nwbegincode{84}\moddef{prc gosam: gosam def: TBP}\plusendmoddef
  procedure :: connect => gosam_def_connect
\nwendcode{}\nwbegincode{85}\moddef{prc gosam: procedures}\plusendmoddef
  subroutine gosam_def_connect (def, lib_driver, i, proc_driver)   
    class(gosam_def_t), intent(in) :: def
    class(prclib_driver_t), intent(in) :: lib_driver
    integer, intent(in) :: i
    integer :: pid, fid
    class(prc_core_driver_t), intent(inout) :: proc_driver
    type(c_funptr) :: fptr
    select type (proc_driver)
    type is (gosam_driver_t)       
       pid = i
       fid = 2
       call lib_driver%get_fptr (pid, fid, fptr)
       call c_f_procpointer (fptr, proc_driver%update_alpha_s)
       fid = 4
       call lib_driver%get_fptr (pid, fid, fptr)
       call c_f_procpointer (fptr, proc_driver%is_allowed)
    end select
  end subroutine gosam_def_connect

\nwendcode{}\nwbegindocs{86}\nwdocspar
%
\nwenddocs{}\nwbegincode{87}\moddef{prc gosam: gosam writer: TBP}\endmoddef
  procedure :: write_wrapper => gosam_write_wrapper
\nwendcode{}\nwbegincode{88}\moddef{prc gosam: procedures}\plusendmoddef
  subroutine gosam_write_wrapper (writer, unit, id, feature)
    class(gosam_writer_t), intent(in) :: writer
    integer, intent(in) :: unit
    type(string_t), intent(in) :: id, feature    
    type(string_t) :: name
    name = writer%get_c_procname (id, feature)
    write (unit, *)
    select case (char (feature))
    case ("init")
       write (unit, "(9A)")  "subroutine ", char (name), " (par) bind(C)"
       write (unit, "(2x,9A)")  "use iso_c_binding"
       write (unit, "(2x,9A)")  "use kinds"
       write (unit, "(2x,9A)")  "use opr_", char (id)
       write (unit, "(2x,9A)")  "real(c_default_float), dimension(*), &
            &intent(in) :: par"
       if (c_default_float == default) then
          write (unit, "(2x,9A)")  "call ", char (feature), " (par)"
       end if
       write (unit, "(9A)")  "end subroutine ", char (name)
    case ("update_alpha_s")
       write (unit, "(9A)")  "subroutine ", char (name), " (alpha_s) bind(C)"
       write (unit, "(2x,9A)")  "use iso_c_binding"
       write (unit, "(2x,9A)")  "use kinds"
       write (unit, "(2x,9A)")  "use opr_", char (id)
       if (c_default_float == default) then
          write (unit, "(2x,9A)")  "real(c_default_float), intent(in) &
               &:: alpha_s"
          write (unit, "(2x,9A)")  "call ", char (feature), " (alpha_s)"
       end if
       write (unit, "(9A)")  "end subroutine ", char (name)
    case ("reset_helicity_selection")
       write (unit, "(9A)")  "subroutine ", char (name), &
            " (threshold, cutoff) bind(C)"
       write (unit, "(2x,9A)")  "use iso_c_binding"
       write (unit, "(2x,9A)")  "use kinds"
       write (unit, "(2x,9A)")  "use opr_", char (id)
       if (c_default_float == default) then
          write (unit, "(2x,9A)")  "real(c_default_float), intent(in) &
               &:: threshold"
          write (unit, "(2x,9A)")  "integer(c_int), intent(in) :: cutoff"
          write (unit, "(2x,9A)")  "call ", char (feature), &
               " (threshold, int (cutoff))"
       end if
       write (unit, "(9A)")  "end subroutine ", char (name)
    case ("is_allowed")
       write (unit, "(9A)")  "subroutine ", char (name), &
            " (flv, hel, col, flag) bind(C)"
       write (unit, "(2x,9A)")  "use iso_c_binding"
       write (unit, "(2x,9A)")  "use kinds"
       write (unit, "(2x,9A)")  "use opr_", char (id)
       write (unit, "(2x,9A)")  "integer(c_int), intent(in) :: flv, hel, col"
       write (unit, "(2x,9A)")  "logical(c_bool), intent(out) :: flag"    
       write (unit, "(2x,9A)")  "flag = ", char (feature), &
            " (int (flv), int (hel), int (col))"
       write (unit, "(9A)")  "end subroutine ", char (name)
    case ("new_event")
       write (unit, "(9A)")  "subroutine ", char (name), " (p) bind(C)"
       write (unit, "(2x,9A)")  "use iso_c_binding"
       write (unit, "(2x,9A)")  "use kinds"
       write (unit, "(2x,9A)")  "use opr_", char (id)
       if (c_default_float == default) then
          write (unit, "(2x,9A)")  "real(c_default_float), dimension(0:3,*), &
               &intent(in) :: p"
          write (unit, "(2x,9A)")  "call ", char (feature), " (p)"
       end if
       write (unit, "(9A)")  "end subroutine ", char (name)
    case ("get_amplitude")
       write (unit, "(9A)")  "subroutine ", char (name), &
            " (flv, hel, col, amp) bind(C)"
       write (unit, "(2x,9A)")  "use iso_c_binding"
       write (unit, "(2x,9A)")  "use kinds"
       write (unit, "(2x,9A)")  "use opr_", char (id)
       write (unit, "(2x,9A)")  "integer(c_int), intent(in) :: flv, hel, col"
       write (unit, "(2x,9A)")  "complex(c_default_complex), intent(out) &
            &:: amp"    
       write (unit, "(2x,9A)")  "amp = ", char (feature), &
            " (int (flv), int (hel), int (col))"
       write (unit, "(9A)")  "end subroutine ", char (name)
    end select

  end subroutine gosam_write_wrapper

\nwendcode{}\nwbegindocs{89}\nwdocspar
%
\nwenddocs{}\nwbegincode{90}\moddef{prc gosam: gosam writer: TBP}\plusendmoddef
  procedure :: write_interface => gosam_write_interface
\nwendcode{}\nwbegincode{91}\moddef{prc gosam: procedures}\plusendmoddef
  subroutine gosam_write_interface (writer, unit, id, feature)
    class(gosam_writer_t), intent(in) :: writer
    integer, intent(in) :: unit
    type(string_t), intent(in) :: id
    type(string_t), intent(in) :: feature
    type(string_t) :: name
    name = writer%get_c_procname (id, feature)
    write (unit, "(2x,9A)")  "interface"
    select case (char (feature))
    case ("init")
       write (unit, "(5x,9A)")  "subroutine ", char (name), " (par) bind(C)"
       write (unit, "(7x,9A)")  "import"
       write (unit, "(7x,9A)")  "real(c_default_float), dimension(*), &
            &intent(in) :: par"
       write (unit, "(5x,9A)")  "end subroutine ", char (name)
    case ("update_alpha_s")
       write (unit, "(5x,9A)")  "subroutine ", char (name), " (alpha_s) bind(C)"
       write (unit, "(7x,9A)")  "import"
       write (unit, "(7x,9A)")  "real(c_default_float), intent(in) :: alpha_s"
       write (unit, "(5x,9A)")  "end subroutine ", char (name)
    case ("reset_helicity_selection")
       write (unit, "(5x,9A)")  "subroutine ", char (name), " &
            &(threshold, cutoff) bind(C)"
       write (unit, "(7x,9A)")  "import"
       write (unit, "(7x,9A)")  "real(c_default_float), intent(in) :: threshold"
       write (unit, "(7x,9A)")  "integer(c_int), intent(in) :: cutoff"
       write (unit, "(5x,9A)")  "end subroutine ", char (name)
    case ("is_allowed")
       write (unit, "(5x,9A)")  "subroutine ", char (name), " &
            &(flv, hel, col, flag) bind(C)"
       write (unit, "(7x,9A)")  "import"
       write (unit, "(7x,9A)")  "integer(c_int), intent(in) :: flv, hel, col"
       write (unit, "(7x,9A)")  "logical(c_bool), intent(out) :: flag"    
       write (unit, "(5x,9A)")  "end subroutine ", char (name)
    case ("new_event")
       write (unit, "(5x,9A)")  "subroutine ", char (name), " (p) bind(C)"
       write (unit, "(7x,9A)")  "import"
       write (unit, "(7x,9A)")  "real(c_default_float), dimension(0:3,*), &
            &intent(in) :: p"
       write (unit, "(5x,9A)")  "end subroutine ", char (name)
    case ("get_amplitude")
       write (unit, "(5x,9A)")  "subroutine ", char (name), " &
            &(flv, hel, col, amp) bind(C)"
       write (unit, "(7x,9A)")  "import"
       write (unit, "(7x,9A)")  "integer(c_int), intent(in) :: flv, hel, col"
       write (unit, "(7x,9A)")  "complex(c_default_complex), intent(out) &
            &:: amp"    
       write (unit, "(5x,9A)")  "end subroutine ", char (name)
    end select
    write (unit, "(2x,9A)")  "end interface"
  end subroutine gosam_write_interface

\nwendcode{}\nwbegindocs{92}\nwdocspar
%
\nwenddocs{}\nwbegincode{93}\moddef{prc gosam: gosam writer: TBP}\plusendmoddef
  procedure :: write_source_code => gosam_write_source_code
\nwendcode{}\nwbegincode{94}\moddef{prc gosam: procedures}\plusendmoddef
  subroutine gosam_write_source_code (writer, id)
    class(gosam_writer_t), intent(in) :: writer
    type(string_t), intent(in) :: id
  end subroutine gosam_write_source_code

\nwendcode{}\nwbegindocs{95}\nwdocspar
%
\nwenddocs{}\nwbegincode{96}\moddef{prc gosam: gosam writer: TBP}\plusendmoddef
  procedure :: write_makefile_code => gosam_write_makefile_code
\nwendcode{}\nwbegincode{97}\moddef{prc gosam: procedures}\plusendmoddef
  subroutine gosam_write_makefile_code (writer, unit, id, os_data, testflag)
    class(gosam_writer_t), intent(in) :: writer
    integer, intent(in) :: unit
    type(string_t), intent(in) :: id
    type(os_data_t), intent(in) :: os_data
    logical, intent(in), optional :: testflag
    type(string_t) :: omega_binary, omega_path
    omega_binary = "omega_" // writer%model_name // ".opt"
    omega_path = os_data%whizard_omega_binpath // "/" // omega_binary
    write (unit, "(5A)")  "OBJECTS += ", char (id), ".lo"
    write (unit, "(5A)")  char (id), ".f90:"
    write (unit, "(99A)")  TAB, char (omega_path), &
         " -o ", char (id), ".f90", &
         " -target:whizard", &
         " -target:parameter_module parameters_", char (writer%model_name), &
         " -target:module opr_", char (id), &
         " -target:md5sum '", writer%md5sum, "'", &
         char (writer%process_mode), char (writer%process_string)
    write (unit, "(5A)")  "clean-", char (id), ":"
    write (unit, "(5A)")  TAB, "rm -f ", char (id), ".f90"
    write (unit, "(5A)")  TAB, "rm -f opr_", char (id), ".mod"
    write (unit, "(5A)")  TAB, "rm -f ", char (id), ".lo"
    write (unit, "(5A)")  "CLEAN_SOURCES += ", char (id), ".f90"    
    write (unit, "(5A)")  "CLEAN_OBJECTS += opr_", char (id), ".mod"       
    write (unit, "(5A)")  "CLEAN_OBJECTS += ", char (id), ".lo"
    write (unit, "(5A)")  char (id), ".lo: ", char (id), ".f90"
    write (unit, "(5A)")  TAB, "$(LTFCOMPILE) $<"

  end subroutine gosam_write_makefile_code

\nwendcode{}\nwbegindocs{98}\nwdocspar
%
\nwenddocs{}\nwbegincode{99}\moddef{prc gosam: gosam writer: TBP}\plusendmoddef
  procedure, nopass:: get_procname => gosam_writer_get_procname
\nwendcode{}\nwbegincode{100}\moddef{prc gosam: procedures}\plusendmoddef
  function gosam_writer_get_procname (feature) result (name)
    type(string_t) :: name
    type(string_t), intent(in) :: feature
    select case (char (feature))
    case ("n_in");   name = "number_particles_in"
    case ("n_out");  name = "number_particles_out"
    case ("n_flv");  name = "number_flavor_states"
    case ("n_hel");  name = "number_spin_states"
    case ("n_col");  name = "number_color_flows"
    case ("n_cin");  name = "number_color_indices"
    case ("n_cf");   name = "number_color_factors"
    case ("flv_state");  name = "flavor_states"
    case ("hel_state");  name = "spin_states"
    case ("col_state");  name = "color_flows"
    case default
       name = feature
    end select
  end function gosam_writer_get_procname

\eatline
\nwendcode{}\nwbegindocs{101}\nwdocspar
\nwenddocs{}\nwbegincode{102}\moddef{prc gosam: gosam writer: TBP}\plusendmoddef
  procedure, nopass :: get_module_name => gosam_writer_get_module_name
\nwendcode{}\nwbegincode{103}\moddef{prc gosam: procedures}\plusendmoddef
  function gosam_writer_get_module_name (id) result (name)
    type(string_t) :: name
    type(string_t), intent(in) :: id
    name = "opr_" // id
  end function gosam_writer_get_module_name

\eatline
\nwendcode{}\nwbegindocs{104}\nwdocspar
\nwenddocs{}\nwbegincode{105}\moddef{prc gosam: gosam writer: TBP}\plusendmoddef
  procedure, nopass :: type_name => gosam_writer_type_name
\nwendcode{}\nwbegincode{106}\moddef{prc gosam: procedures}\plusendmoddef
  function gosam_writer_type_name () result (string)
    type(string_t) :: string
    string = "gosam"
  end function gosam_writer_type_name

\nwendcode{}\nwbegindocs{107}\nwdocspar
%
\nwenddocs{}\nwbegincode{108}\moddef{prc gosam: gosam writer: TBP}\plusendmoddef
  procedure :: write => gosam_writer_write
\nwendcode{}\nwbegincode{109}\moddef{prc gosam: procedures}\plusendmoddef
  subroutine gosam_writer_write (writer, unit)
    class(gosam_writer_t), intent(in) :: writer
    integer, intent(in) :: unit    
    write (unit, "(1x,A)")  char (writer%get_process_string ())
  end subroutine gosam_writer_write

\nwendcode{}\nwbegindocs{110}\nwdocspar
%
\nwenddocs{}\nwbegincode{111}\moddef{prc gosam: gosam writer: TBP}\plusendmoddef
  procedure :: init => gosam_writer_init
\nwendcode{}\nwbegincode{112}\moddef{prc gosam: procedures}\plusendmoddef
  subroutine gosam_writer_init (writer,model_name, prt_in, prt_out) 
    class(gosam_writer_t), intent(inout) :: writer
    type(string_t), intent(in) :: model_name
    type(string_t), dimension(:), intent(in) :: prt_in, prt_out
    integer :: i, unit

    writer%gosam_dir = GOSAM_DIR
    writer%golem_dir = GOLEM_DIR 
    writer%samurai_dir = SAMURAI_DIR
    writer%ninja_dir = NINJA_DIR
    writer%form_dir = FORM_DIR
    writer%qgraf_dir = QGRAF_DIR

    writer%model_name = model_name

    select case (size (prt_in))
      case (1); writer%process_mode = " -decay"
      case (2); writer%process_mode = " -scatter"
    end select
    associate (s => writer%process_string)
      s = " '" 
      do i = 1, size (prt_in)
         if (i > 1) s = s // " "
         s = s // prt_in(i)
      end do
      s = s // " ->"
      do i = 1, size (prt_out)
         s = s // " " // prt_out(i)
      end do
      s = s // "'"
    end associate

    unit = free_unit ()
    open (unit, file = "golem.in", status = "replace", action = "write")
    call writer%generate_configuration_file (unit)
    close(unit)
  end subroutine gosam_writer_init  

\eatline
\nwendcode{}\nwbegindocs{113}\nwdocspar
\nwenddocs{}\nwbegincode{114}\moddef{prc gosam: gosam driver: TBP}\endmoddef
  procedure, nopass :: type_name => gosam_driver_type_name
\nwendcode{}\nwbegincode{115}\moddef{prc gosam: procedures}\plusendmoddef
  function gosam_driver_type_name () result (string)
    type(string_t) :: string
    string = "gosam"
  end function gosam_driver_type_name

\nwendcode{}\nwbegindocs{116}\nwdocspar
%
\nwenddocs{}\nwbegincode{117}\moddef{prc gosam: gosam driver: TBP}\plusendmoddef
  procedure :: init_gosam => gosam_driver_init_gosam
\nwendcode{}\nwbegincode{118}\moddef{prc gosam: procedures}\plusendmoddef
  subroutine gosam_driver_init_gosam (object, os_data, olp_file, &
                                olc_file, olp_dir)
    class(gosam_driver_t), intent(inout) :: object
    type(os_data_t), intent(in) :: os_data
    type(string_t), intent(in) :: olp_file, olc_file, olp_dir
    object%gosam_dir = GOSAM_DIR
    object%olp_file = olp_file
    object%olc_file = olc_file
    object%olp_dir = olp_dir
    call object%loop_archive%activate (var_str ('Generated_Loops'))
  end subroutine gosam_driver_init_gosam

\eatline
\nwendcode{}\nwbegindocs{119}\nwdocspar
\nwenddocs{}\nwbegincode{120}\moddef{prc gosam: gosam writer: TBP}\plusendmoddef
  procedure :: generate_configuration_file => &
            gosam_writer_generate_configuration_file
\nwendcode{}\nwbegincode{121}\moddef{prc gosam: procedures}\plusendmoddef
  subroutine gosam_writer_generate_configuration_file &
          (object, unit)
      class(gosam_writer_t), intent(in) :: object
      integer, intent(in) :: unit
      type(string_t) :: fc_bin
      type(string_t) :: form_bin, qgraf_bin, haggies_bin
      type(string_t) :: fcflags_golem, ldflags_golem
      type(string_t) :: fcflags_samurai, ldflags_samurai
      type(string_t) :: fcflags_ninja, ldflags_ninja
      type(string_t) :: ldflags_avh_olo, ldflags_qcdloop
      fc_bin = DEFAULT_FC
      form_bin = object%form_dir // '/bin/tform'
      qgraf_bin = object%qgraf_dir // '/bin/qgraf'
      if (object%gosam_dir /= "") then
        haggies_bin = '/usr/bin/java -jar ' // object%gosam_dir // &
                       '/share/golem/haggies/haggies.jar'
      else
        call msg_fatal ("generate_configuration_file: At least " // &
             "the GoSam Directory has to be specified!")
      end if
      if (object%golem_dir /= "") then
        fcflags_golem = "-I" // object%golem_dir // "/include/golem95"
        ldflags_golem = "-L" // object%golem_dir // "/lib -lgolem"
      end if
      if (object%samurai_dir /= "") then
        fcflags_samurai = "-I" // object%samurai_dir // "/include/samurai"
        ldflags_samurai = "-L" // object%samurai_dir // "/lib -lsamurai"
        ldflags_avh_olo = "-L" // object%samurai_dir // "/lib -lavh_olo"
        ldflags_qcdloop = "-L" // object%samurai_dir // "/lib -lqcdloop"
      end if
      if (object%ninja_dir /= "") then
        fcflags_ninja = "-I" // object%ninja_dir // "/include/ninja " &
                        // "-I" // object%ninja_dir // "/include"
        ldflags_ninja = "-L" // object%ninja_dir // "/lib -lninja"
      end if
      write (unit, "(A)") "+avh_olo.ldflags=" &
            // char (ldflags_avh_olo) 
      write (unit, "(A)") "reduction_programs=golem95, samurai, ninja"
      write (unit, "(A)") "extensions=autotools"
      write (unit, "(A)") "+qcdloop.ldflags=" &
            // char (ldflags_qcdloop)
      write (unit, "(A)") "+zzz.extensions=qcdloop, avh_olo"
      write (unit, "(A)") "fc.bin=" // char (fc_bin)
      write (unit, "(A)") "form.bin=" // char (form_bin)
      write (unit, "(A)") "qgraf.bin=" // char (qgraf_bin)
      write (unit, "(A)") "golem95.fcflags=" // char (fcflags_golem)
      write (unit, "(A)") "golem95.ldflags=" // char (ldflags_golem)
      write (unit, "(A)") "haggies.bin=" // char (haggies_bin)
      write (unit, "(A)") "samurai.fcflags=" // char (fcflags_samurai)
      write (unit, "(A)") "samurai.ldflags=" // char (ldflags_samurai)
      write (unit, "(A)") "ninja.fcflags=" // char (fcflags_ninja)
      write (unit, "(A)") "ninja.ldflags=" // char (ldflags_ninja)
      !!! This might collide with the mass-setup in the order-file
      !!! write (unit, "(A)") "zero=mU,mD,mC,mS,mB"
      write (unit, "(A)") "PSP_check=False"
      write (unit, "(A)") "filter.lo=lambda d: d.iprop(H) == 0 and d.iprop(chi) == 0"
      write (unit, "(A)") "filter.nlo=lambda d: d.iprop(H) == 0 and d.iprop(chi) == 0"
  end subroutine gosam_writer_generate_configuration_file

\eatline
\nwendcode{}\nwbegindocs{122}\nwdocspar
\nwenddocs{}\nwbegincode{123}\moddef{prc gosam: gosam writer: TBP}\plusendmoddef
  procedure :: get_process_string => gosam_writer_get_process_string
\nwendcode{}\nwbegincode{124}\moddef{prc gosam: procedures}\plusendmoddef
  function gosam_writer_get_process_string (writer) result (s_proc)
    class(gosam_writer_t), intent(in) :: writer
    type(string_t) :: s_proc
  end function gosam_writer_get_process_string

\eatline
\nwendcode{}\nwbegindocs{125}\nwdocspar
\nwenddocs{}\nwbegincode{126}\moddef{prc gosam: gosam writer: TBP}\plusendmoddef
  procedure :: get_n_proc => gosam_writer_get_n_proc
\nwendcode{}\nwbegincode{127}\moddef{prc gosam: procedures}\plusendmoddef
  function gosam_writer_get_n_proc (writer) result (n_proc)
    class(gosam_writer_t), intent(in) :: writer
    integer :: n_proc
    n_proc = blha_configuration_get_n_proc (writer%blha_cfg)
  end function gosam_writer_get_n_proc

\eatline
\nwendcode{}\nwbegindocs{128}\nwdocspar
\nwenddocs{}\nwbegincode{129}\moddef{prc gosam: gosam driver: TBP}\plusendmoddef
  procedure :: execute_olp_file => gosam_driver_execute_olp_file
\nwendcode{}\nwbegincode{130}\moddef{prc gosam: procedures}\plusendmoddef
  subroutine gosam_driver_execute_olp_file (object)
    class(gosam_driver_t), intent(in) :: object
    type(string_t) :: command
    command = object%gosam_dir // "/bin" // &
              '/gosam.py --olp ' // object%olp_file // &
              ' --destination=' // object%olp_dir // &
              ' -f -z'
    call os_system_call (command, verbose = .true.)
    call msg_message ("Configure olp_modules")
    command = "cd " // object%olp_dir // "; ./autogen.sh --prefix=$(pwd)"
    call os_system_call (command, verbose = .true.)
    call msg_message ("Installing olp_modules")
    command = "cd " // object%olp_dir // ";  make install"
    call os_system_call (command, verbose = .true.)
  end subroutine gosam_driver_execute_olp_file

\eatline
\nwendcode{}\nwbegindocs{131}\nwdocspar
\nwenddocs{}\nwbegincode{132}\moddef{prc gosam: gosam driver: TBP}\plusendmoddef
  procedure :: set_alpha_s => gosam_driver_set_alpha_s
\nwendcode{}\nwbegincode{133}\moddef{prc gosam: procedures}\plusendmoddef
  subroutine gosam_driver_set_alpha_s (driver, alpha_s)
     class(gosam_driver_t), intent(inout) :: driver
     real(default), intent(in) :: alpha_s
     integer :: ierr
     call driver%gosam_olp_set_parameter &
              (c_char_'alphaS'//c_null_char, &
               dble (alpha_s), 0._double, ierr)
  end subroutine gosam_driver_set_alpha_s

\eatline
\nwendcode{}\nwbegindocs{134}\nwdocspar
\nwenddocs{}\nwbegincode{135}\moddef{prc gosam: gosam driver: TBP}\plusendmoddef
  procedure :: set_mass_and_width => gosam_driver_set_mass_and_width
\nwendcode{}\nwbegincode{136}\moddef{prc gosam: procedures}\plusendmoddef
  subroutine gosam_driver_set_mass_and_width (driver, &
                                       i_pdg, mass, width)
    class(gosam_driver_t), intent(inout) :: driver
    integer, intent(in) :: i_pdg
    real(default), intent(in), optional :: mass
    real(default), intent(in), optional :: width
    type(string_t) :: buf
    character(kind=c_char,len=20) :: c_string
    integer :: ierr
    if (present (mass)) then
       buf = 'mass(' // integer_to_string (abs(i_pdg)) // ')'
       c_string = char(buf)//c_null_char
       call driver%gosam_olp_set_parameter &
                (c_string, dble(mass), 0._double, ierr)
       if (ierr == 0) then
          buf = "GoSam: Attempt to set mass of particle " // &
                integer_to_string (abs(i_pdg)) // "failed"
          call msg_fatal (char(buf))
       end if
    end if
    if (present (width)) then
       buf = 'width(' // integer_to_string (abs(i_pdg)) // ')'
       c_string = char(buf)//c_null_char
       call driver%gosam_olp_set_parameter &
                (c_string, dble(width), 0._double, ierr)
       if (ierr == 0) then
          buf = "GoSam: Attempt to set width of particle " // &
                integer_to_string (abs(i_pdg)) // "failed"
          call msg_fatal (char(buf))
       end if
    end if
  end subroutine gosam_driver_set_mass_and_width

\eatline
\nwendcode{}\nwbegindocs{137}\nwdocspar
\nwenddocs{}\nwbegincode{138}\moddef{prc gosam: gosam driver: TBP}\plusendmoddef
  procedure :: load => gosam_driver_load
\nwendcode{}\nwbegincode{139}\moddef{prc gosam: procedures}\plusendmoddef
  subroutine gosam_driver_load (object, os_data, store, success)
    class(gosam_driver_t), intent(inout) :: object
    type(os_data_t), intent(in) :: os_data
    logical, intent(in) :: store
    logical, intent(out) :: success
    type(dlaccess_t) :: dlaccess
    type(string_t) ::  path
    type(c_funptr) :: c_fptr

    path = object%olp_dir // '/.libs'
    call dlaccess_init &
          (dlaccess, var_str ("."), path // '/libgolem_olp.so', os_data)
    if (os_file_exist (path // '/libgolem_olp.so')) then
      c_fptr = dlaccess_get_c_funptr (dlaccess, var_str ("OLP_Start"))
      call c_f_procpointer (c_fptr, object%gosam_olp_start)

      c_fptr = dlaccess_get_c_funptr (dlaccess, var_str ("OLP_EvalSubProcess"))
      call c_f_procpointer (c_fptr, object%gosam_olp_eval)

      c_fptr = dlaccess_get_c_funptr (dlaccess, var_str ("OLP_Info"))
      call c_f_procpointer (c_fptr, object%gosam_olp_info)

      c_fptr = dlaccess_get_c_funptr (dlaccess, var_str ("OLP_SetParameter"))
      call c_f_procpointer (c_fptr, object%gosam_olp_set_parameter)

      c_fptr = dlaccess_get_c_funptr (dlaccess, var_str ("OLP_EvalSubProcess2"))
      call c_f_procpointer (c_fptr, object%gosam_olp_eval2)

      c_fptr = dlaccess_get_c_funptr (dlaccess, var_str ("OLP_Option"))
      call c_f_procpointer (c_fptr, object%gosam_olp_option)

      c_fptr = dlaccess_get_c_funptr (dlaccess, var_str ("OLP_Polvec"))
      call c_f_procpointer (c_fptr, object%gosam_olp_polvec)

      c_fptr = dlaccess_get_c_funptr (dlaccess, var_str ("OLP_Finalize"))
      call c_f_procpointer (c_fptr, object%gosam_olp_finalize)

      c_fptr = dlaccess_get_c_funptr (dlaccess, var_str ("OLP_PrintParameter"))
      call c_f_procpointer (c_fptr, object%gosam_olp_print_parameter)

      success = .not. dlaccess_has_error (dlaccess)
      if (store .and. success) call object%loop_archive%record (object%olp_file, &
                                  var_str ('golem.in'), path // '/libgolem_olp.so')
    else
      success = .false.
    end if
  end subroutine gosam_driver_load

\eatline
\nwendcode{}\nwbegindocs{140}\nwdocspar
\nwenddocs{}\nwbegincode{141}\moddef{prc gosam: gosam driver: TBP}\plusendmoddef
  procedure :: read_olc_file => gosam_driver_read_olc_file
\nwendcode{}\nwbegincode{142}\moddef{prc gosam: procedures}\plusendmoddef
  subroutine gosam_driver_read_olc_file (driver, flavors, amp_type, flv_index, label)
    class(gosam_driver_t), intent(inout) :: driver
    integer, intent(in), dimension(:,:) :: flavors
    integer, intent(out), dimension(20) :: amp_type, flv_index, label
    integer :: unit, filestat
    character(len=100):: rd_line 
    character(len=3) :: rd_flavor
    logical :: read_flavor, born_found
    integer, dimension(:,:,:), allocatable :: flv_position
    integer, dimension(:,:), allocatable :: flv_read
    integer :: k, i_flv, i_part, current_flavor
    integer :: pos_label

    flv_position = compute_flavor_positions (flavors)
    amp_type = -1; flv_index = -1; label = -1
    allocate (flv_read (size (flavors, 1), &
                        size (flavors, 2)))
    unit = free_unit ()
    open (unit, file=char(driver%olc_file), status="old") 
    read_flavor=.false.
    k = 1
    do
      read (unit, '(A)', iostat = filestat) rd_line
      if (filestat == iostat_end) then
         exit
      else
         if (rd_line(1:13) == 'AmplitudeType') then
            if (rd_line(15:19) == 'Loop') then
               amp_type(k) = BLHA_AMP_LOOP
            else if (rd_line(15:19) == 'Tree') then
               amp_type(k) = BLHA_AMP_TREE
            else if (rd_line(15:21) == 'ccTree') then
               amp_type(k) = BLHA_AMP_CC
            else if (rd_line(15:21) == 'scTree') then
               amp_type(k) = BLHA_AMP_SC
            else
               call msg_fatal ("AmplitudeType present but &
                               &AmpType not known!")
            end if
            read_flavor = .true.
         else if (read_flavor) then
            do i_flv = 1, size (flavors, 2)
               do i_part = 1, size (flavors, 1)
                  read (rd_line (flv_position (i_part, i_flv, 1): &
                                 flv_position (i_part, i_flv, 2)), '(I3)') &
                                 flv_read (i_part, i_flv)
               end do
            end do
            born_found = .false.
            do i_flv = 1, size (flavors, 2)
               if (all (flv_read (:,i_flv) == flavors (:,i_flv))) then
                  flv_index (k) = i_flv
                  pos_label = maxval (flv_position (:,i_flv,:)) + 6
                  read (rd_line (pos_label:pos_label), '(I2)') label(k)
                  born_found = .true.
                  k = k+1
                  exit
               end if
            end do
            if (.not. born_found) call msg_fatal & 
                     ("No underlying Born found")
         end if   
      end if
    end do
    close(unit)
  contains
    function compute_flavor_positions (flavors) result (pos)
      integer, intent(in), dimension(:,:) :: flavors
      integer, dimension(:,:,:), allocatable :: pos
      integer :: i_flv, i_part
      integer :: i_first, i_last
      allocate (pos (size (flavors, 1), &
                     size (flavors, 2), 2))
      do i_flv = 1, size (flavors, 2)
         i_first = 1; i_last = 0
         do i_part = 1, size (flavors, 1)
            if (i_last > 0 .and. i_part /= 3) then
               i_first = i_last+2
            else if (i_last > 0 .and. i_part == 3) then
               i_first = i_first + 6
            end if
            i_last = i_first + flavor_digits (flavors (i_part, i_flv)) - 1
            pos (i_part, i_flv, 1) = i_first
            pos (i_part, i_flv, 2) = i_last
         end do
      end do
    end function compute_flavor_positions             
            
    function flavor_digits (flavor) result (ndigits)
      integer, intent(in) :: flavor
      integer :: ndigits
      if (abs(flavor) >= 10) then
         ndigits = 2
      else
         ndigits = 1
      end if
      if (flavor < 0) ndigits = ndigits+1
    end function flavor_digits
  end subroutine gosam_driver_read_olc_file

\eatline
\nwendcode{}\nwbegindocs{143}\nwdocspar
\nwenddocs{}\nwbegincode{144}\moddef{prc gosam: prc gosam: TBP}\endmoddef
  procedure :: execute_olp_file => prc_gosam_execute_olp_file
\nwendcode{}\nwbegincode{145}\moddef{prc gosam: procedures}\plusendmoddef
  subroutine prc_gosam_execute_olp_file (object, os_data)
    class(prc_gosam_t), intent(inout) :: object
    type(os_data_t), intent(in) :: os_data
    logical :: found, dl_success
    integer :: ierr
    select type (driver => object%driver)
    type is (gosam_driver_t)
       call driver%loop_archive%search ([driver%olp_file, &
                        var_str ('golem.in'), driver%olp_dir // &
                        '/.libs/libgolem_olp.so'], found)
       if (found) then
          call driver%loop_archive%restore (driver%olp_file, driver%olp_dir)
       else
         call driver%execute_olp_file ()
       end if
       call driver%load (os_data, .not.found, dl_success)
       if (.not. dl_success) &
          call msg_fatal ('Error linking GoSam libraries')
       call driver%gosam_olp_start (char (driver%olc_file), ierr)
    end select
  end subroutine prc_gosam_execute_olp_file
       
\eatline
\nwendcode{}\nwbegindocs{146}\nwdocspar
\nwenddocs{}\nwbegincode{147}\moddef{prc gosam: prc gosam: TBP}\plusendmoddef
  procedure :: write => prc_gosam_write
\nwendcode{}\nwbegincode{148}\moddef{prc gosam: procedures}\plusendmoddef
  subroutine prc_gosam_write (object, unit)
    class(prc_gosam_t), intent(in) :: object
    integer, intent(in), optional :: unit
    call msg_message ("GOSAM")
  end subroutine prc_gosam_write

\nwendcode{}\nwbegindocs{149}\nwdocspar
%
\nwenddocs{}\nwbegincode{150}\moddef{prc gosam: prc gosam: TBP}\plusendmoddef
  procedure :: needs_mcset => prc_gosam_needs_mcset
\nwendcode{}\nwbegincode{151}\moddef{prc gosam: procedures}\plusendmoddef
  function prc_gosam_needs_mcset (object) result (flag)
    class(prc_gosam_t), intent(in) :: object
    logical :: flag
    flag = .true.
  end function prc_gosam_needs_mcset

\nwendcode{}\nwbegindocs{152}\nwdocspar
%
\nwenddocs{}\nwbegincode{153}\moddef{prc gosam: prc gosam: TBP}\plusendmoddef
  procedure :: get_n_terms => prc_gosam_get_n_terms
\nwendcode{}\nwbegincode{154}\moddef{prc gosam: procedures}\plusendmoddef
  function prc_gosam_get_n_terms (object) result (n)
    class(prc_gosam_t), intent(in) :: object
    integer :: n
    n = 1
  end function prc_gosam_get_n_terms

\nwendcode{}\nwbegindocs{155} 
%
\nwenddocs{}\nwbegincode{156}\moddef{prc gosam: prc gosam: TBP}\plusendmoddef
  procedure :: is_allowed => prc_gosam_is_allowed
\nwendcode{}\nwbegincode{157}\moddef{prc gosam: procedures}\plusendmoddef
  function prc_gosam_is_allowed (object, i_term, f, h, c) result (flag)
    class(prc_gosam_t), intent(in) :: object
    integer, intent(in) :: i_term, f, h, c
    logical :: flag
    logical(c_bool) :: cflag
    select type (driver => object%driver)
    type is (gosam_driver_t)
       call driver%is_allowed (f, h, c, cflag)
       flag = cflag
    class default
       call msg_fatal ("Gosam instance created, but driver is not a GoSam driver!")
    end select
  end function prc_gosam_is_allowed

\nwendcode{}\nwbegindocs{158}\nwdocspar
%
\nwenddocs{}\nwbegincode{159}\moddef{prc gosam: prc gosam: TBP}\plusendmoddef
  procedure :: update_alpha_s => prc_gosam_update_alpha_s
\nwendcode{}\nwbegincode{160}\moddef{prc gosam: procedures}\plusendmoddef
  subroutine prc_gosam_update_alpha_s (object, tmp, fac_scale) 
    class(prc_gosam_t), intent(in) :: object
    class(workspace_t), intent(inout), allocatable :: tmp
    real(default), intent(in) :: fac_scale
    real(default) :: alpha_qcd
    if (allocated (object%qcd%alpha)) then
       alpha_qcd = object%qcd%alpha%get (fac_scale)
       select type (driver => object%driver)
       type is (gosam_driver_t)
          call driver%update_alpha_s (alpha_qcd)
       end select 
       if (allocated (tmp)) then
          select type (tmp)
          type is (gosam_state_t)
             tmp%alpha_qcd = alpha_qcd
          end select
       end if
    end if
  end subroutine prc_gosam_update_alpha_s

\eatline
\nwendcode{}\nwbegindocs{161}\nwdocspar
\nwenddocs{}\nwbegincode{162}\moddef{prc gosam: prc gosam: TBP}\plusendmoddef
  procedure :: compute_hard_kinematics => prc_gosam_compute_hard_kinematics
\nwendcode{}\nwbegincode{163}\moddef{prc gosam: procedures}\plusendmoddef
  subroutine prc_gosam_compute_hard_kinematics &
       (object, p_seed, i_term, int_hard, tmp)
    class(prc_gosam_t), intent(in) :: object
    type(vector4_t), dimension(:), intent(in) :: p_seed
    integer, intent(in) :: i_term
    type(interaction_t), intent(inout) :: int_hard
    class(workspace_t), intent(inout), allocatable :: tmp 
    call interaction_set_momenta (int_hard, p_seed)
    if (allocated (tmp)) then
      select type (tmp)
      type is (gosam_state_t); tmp%new_kinematics = .true.
      end select
    end if
  end subroutine prc_gosam_compute_hard_kinematics

\nwendcode{}\nwbegindocs{164}\nwdocspar
%
\nwenddocs{}\nwbegincode{165}\moddef{prc gosam: prc gosam: TBP}\plusendmoddef
  procedure :: compute_eff_kinematics => prc_gosam_compute_eff_kinematics
\nwendcode{}\nwbegincode{166}\moddef{prc gosam: procedures}\plusendmoddef
  subroutine prc_gosam_compute_eff_kinematics &
       (object, i_term, int_hard, int_eff, tmp)
    class(prc_gosam_t), intent(in) :: object
    integer, intent(in) :: i_term
    type(interaction_t), intent(in) :: int_hard
    type(interaction_t), intent(inout) :: int_eff
    class(workspace_t), intent(inout), allocatable :: tmp
  end subroutine prc_gosam_compute_eff_kinematics

\nwendcode{}\nwbegindocs{167}\nwdocspar
%
\nwenddocs{}\nwbegincode{168}\moddef{prc gosam: prc gosam: TBP}\plusendmoddef
  procedure :: compute_amplitude => prc_gosam_compute_amplitude
\nwendcode{}\nwbegincode{169}\moddef{prc gosam: procedures}\plusendmoddef
  function prc_gosam_compute_amplitude &
       (object, j, p, f, h, c, fac_scale, ren_scale, alpha_qcd_forced, tmp) &
       result (amp)
    class(prc_gosam_t), intent(in) :: object
    integer, intent(in) :: j
    type(vector4_t), dimension(:), intent(in) :: p
    integer, intent(in) :: f, h, c
    real(default), intent(in) :: fac_scale, ren_scale
    real(default), intent(in), allocatable :: alpha_qcd_forced
    class(workspace_t), intent(inout), allocatable, optional :: tmp
    complex(default) :: amp
    select type (tmp)
    type is (gosam_state_t)
      tmp%alpha_qcd = object%qcd%alpha%get (fac_scale)
    end select
    amp = 0.0
  end function prc_gosam_compute_amplitude

\nwendcode{}\nwbegindocs{170}\nwdocspar
%
\nwenddocs{}\nwbegincode{171}\moddef{prc gosam: prc gosam: TBP}\plusendmoddef
  procedure :: recover_kinematics => prc_gosam_recover_kinematics
\nwendcode{}\nwbegincode{172}\moddef{prc gosam: procedures}\plusendmoddef
  subroutine prc_gosam_recover_kinematics &
       (object, p_seed, int_hard, int_eff, tmp)
    class(prc_gosam_t), intent(in) :: object
    type(vector4_t), dimension(:), intent(inout) :: p_seed
    type(interaction_t), intent(inout) :: int_hard, int_eff
    class(workspace_t), intent(inout), allocatable :: tmp
    integer :: n_in
    n_in = interaction_get_n_in (int_eff)
    call interaction_set_momenta (int_eff, p_seed(1:n_in), outgoing = .false.)
    p_seed(n_in+1:) = interaction_get_momenta (int_eff, outgoing = .true.)
  end subroutine prc_gosam_recover_kinematics

\nwendcode{}\nwbegindocs{173}\nwdocspar
%
\nwenddocs{}\nwbegincode{174}\moddef{prc gosam: prc gosam: TBP}\plusendmoddef
   procedure :: init_gosam => prc_gosam_init_gosam
\nwendcode{}\nwbegincode{175}\moddef{prc gosam: procedures}\plusendmoddef
  subroutine prc_gosam_init_gosam (object, gosam_template)
    !!! Not prc_gosam_init to avoid name-clash
    class(prc_gosam_t), intent(inout) :: object
    type(gosam_template_t), intent(inout) :: gosam_template
    integer :: i_flv

    object%n_flv = size (object%data%flv_state, 2)
   
    if (gosam_template%compute_loop ()) then
       allocate (object%i_virt (object%n_flv), &
                 object%i_cc (object%n_flv))
    else if (gosam_template%compute_subtraction ()) then
       allocate (object%i_born (object%n_flv), &
                 object%i_cc (object%n_flv) , &
                 object%i_sc (object%n_flv))
    else if (gosam_template%compute_real_trees ()) then
       allocate (object%i_real (object%n_flv))
    end if
  end subroutine prc_gosam_init_gosam

\eatline
\nwendcode{}\nwbegindocs{176}\nwdocspar
\nwenddocs{}\nwbegincode{177}\moddef{prc gosam: prc gosam: TBP}\plusendmoddef
  procedure :: get_nflv => prc_gosam_get_nflv
\nwendcode{}\nwbegincode{178}\moddef{prc gosam: procedures}\plusendmoddef
  function prc_gosam_get_nflv (object) result (n_flv)
    class(prc_gosam_t), intent(in) :: object
    integer :: n_flv
    n_flv = object%n_flv
  end function prc_gosam_get_nflv

\eatline
\nwendcode{}\nwbegindocs{179}\nwdocspar
\nwenddocs{}\nwbegincode{180}\moddef{prc gosam: prc gosam: TBP}\plusendmoddef
  procedure :: init_driver => prc_gosam_init_driver
\nwendcode{}\nwbegincode{181}\moddef{prc gosam: procedures}\plusendmoddef
  subroutine prc_gosam_init_driver (object, os_data)
    class(prc_gosam_t), intent(inout) :: object
    type(os_data_t), intent(in) :: os_data
    logical :: dl_success
    type(string_t) :: olp_file, olc_file, olp_dir
    integer(c_int) :: success
    logical :: found = .false.

    select type (def => object%def)
    type is (gosam_def_t)
      olp_file = def%basename // def%suffix // '.olp'
      olc_file = def%basename // def%suffix // '.olc'
      olp_dir = def%basename // def%suffix // '_olp_modules'
    class default
      call msg_bug ("prc_gosam_init_driver: core_def should be of gosam-type")
    end select

    select type(driver => object%driver)
    type is (gosam_driver_t)
      call driver%init_gosam (os_data, olp_file, olc_file, olp_dir)
    end select
  end subroutine prc_gosam_init_driver

\eatline
\nwendcode{}\nwbegindocs{182}\nwdocspar
\nwenddocs{}\nwbegincode{183}\moddef{prc gosam: prc gosam: TBP}\plusendmoddef
  procedure :: set_initialized => prc_gosam_set_initialized
\nwendcode{}\nwbegincode{184}\moddef{prc gosam: procedures}\plusendmoddef
  subroutine prc_gosam_set_initialized (prc_gosam)
    class(prc_gosam_t), intent(inout) :: prc_gosam
    prc_gosam%initialized = .true.
  end subroutine prc_gosam_set_initialized 

\eatline
\nwendcode{}\nwbegindocs{185}\nwdocspar
\nwenddocs{}\nwbegincode{186}\moddef{prc gosam: prc gosam: TBP}\plusendmoddef
  procedure :: set_parameters => prc_gosam_set_parameters
\nwendcode{}\nwbegincode{187}\moddef{prc gosam: procedures}\plusendmoddef
  subroutine prc_gosam_set_parameters (object, qcd, use_color_factors)
    class(prc_gosam_t), intent(inout) :: object
    type(qcd_t), intent(in) :: qcd
    logical, intent(in) :: use_color_factors
    object%qcd = qcd
    object%use_color_factors = use_color_factors

  end subroutine prc_gosam_set_parameters

\eatline
\nwendcode{}\nwbegindocs{188}\nwdocspar
\nwenddocs{}\nwbegincode{189}\moddef{prc gosam: procedures}\plusendmoddef
  function create_blha_momentum_array (p) result (mom)
    type(vector4_t), intent(in), dimension(:) :: p
    real(double), dimension(GOSAM_MOMENTUM_LIMIT) :: mom
    integer :: n, i, k

    n = size (p)
    if (n > 10) call msg_fatal ("Number of external particles exceeeds" &
                                 // "size of GoSam-internal momentum array")
    k = 1
    do i = 1, n
       mom(k:k+3) = vector4_get_components (p(i))
       mom(k+4) = invariant_mass (p(i))
       k = k+5
    end do
    mom (k:50) = 0.0
  end function create_blha_momentum_array

\eatline
\nwendcode{}\nwbegindocs{190}\nwdocspar
\nwenddocs{}\nwbegincode{191}\moddef{prc gosam: prc gosam: TBP}\plusendmoddef
  procedure :: compute_sqme_virt => prc_gosam_compute_sqme_virt
\nwendcode{}\nwbegincode{192}\moddef{prc gosam: procedures}\plusendmoddef
  subroutine prc_gosam_compute_sqme_virt (object, &
                i_flv, p, ren_scale, alpha_s, sqme, bad_point)
    class(prc_gosam_t), intent(inout) :: object
    integer, intent(in) :: i_flv
    type(vector4_t), dimension(:), intent(in) :: p
    real(default), intent(in) :: ren_scale
    real(default), intent(in) :: alpha_s
    logical, intent(out) :: bad_point
    real(default), dimension(4), intent(out) :: sqme
    real(double) :: mu
    real(double), dimension(GOSAM_MOMENTUM_LIMIT) :: mom
    real(double), dimension(GOSAM_RESULTS_LIMIT) :: r
    real(double) :: acc_dble
    real(default) :: acc

    mom = create_blha_momentum_array (p)
    if (ren_scale == 0.0) then
      mu = sqrt (2* (p(1)*p(2)))
    else
      mu = ren_scale
    end if
    select type (driver => object%driver)
    type is (gosam_driver_t)
      call driver%set_alpha_s (alpha_s)
      call driver%gosam_olp_eval2 (object%i_virt(i_flv), &
                                   mom, mu, r, acc_dble) 
    end select
    acc = acc_dble
    sqme = r(1:4)
    if (acc > object%maximum_accuracy) then
       bad_point = .true.
    else
       bad_point = .false.
    end if
  end subroutine prc_gosam_compute_sqme_virt

\eatline
\nwendcode{}\nwbegindocs{193}\nwdocspar
\nwenddocs{}\nwbegincode{194}\moddef{prc gosam: prc gosam: TBP}\plusendmoddef
  procedure :: compute_sqme_real => prc_gosam_compute_sqme_real
\nwendcode{}\nwbegincode{195}\moddef{prc gosam: procedures}\plusendmoddef
  subroutine prc_gosam_compute_sqme_real &
         (object, i_flv, p, ren_scale, alpha_s, sqme, bad_point)
    class(prc_gosam_t), intent(inout) :: object
    integer, intent(in) :: i_flv
    type(vector4_t), intent(in), dimension(:) :: p
    real(default), intent(in) :: ren_scale
    real(default), intent(in) :: alpha_s
    real(default), intent(out) :: sqme
    logical, intent(out) :: bad_point
    real(double), dimension(GOSAM_MOMENTUM_LIMIT) :: mom
    real(double), dimension(GOSAM_RESULTS_LIMIT) :: r
    real(double) :: mu
    real(double) :: acc_dble
    real(default) :: acc
 
    mom = create_blha_momentum_array (p)
    if (ren_scale == 0.0) then
       mu = sqrt (2*p(1)*p(2))
    else
      mu = ren_scale
    end if
    select type (driver => object%driver)
    type is (gosam_driver_t)
       call driver%set_alpha_s (alpha_s)
       call driver%gosam_olp_eval2 (object%i_real(i_flv), mom, &
                                    mu, r, acc_dble)
       sqme = r(4)
    end select
    acc = acc_dble
    if (acc > object%maximum_accuracy) bad_point = .true.
  end subroutine prc_gosam_compute_sqme_real

\eatline
\nwendcode{}\nwbegindocs{196}\nwdocspar
\nwenddocs{}\nwbegincode{197}\moddef{prc gosam: prc gosam: TBP}\plusendmoddef
  procedure :: compute_sqme_cc => prc_gosam_compute_sqme_cc
\nwendcode{}\nwbegincode{198}\moddef{prc gosam: procedures}\plusendmoddef
  subroutine prc_gosam_compute_sqme_cc &
         (object, i_flv, p, ren_scale, alpha_s, &
          born_out, born_cc, bad_point)
    class(prc_gosam_t), intent(inout) :: object
    integer, intent(in) :: i_flv
    type(vector4_t), intent(in), dimension(:) :: p
    real(default), intent(in) :: ren_scale
    real(default), intent(in) :: alpha_s
    real(default), intent(out), optional :: born_out
    real(default), intent(inout), dimension(:,:) :: born_cc
    logical, intent(out) :: bad_point
    real(double), dimension(GOSAM_MOMENTUM_LIMIT) :: mom
    real(double), dimension(GOSAM_RESULTS_LIMIT) :: r
    real(double) :: mu
    integer :: i, j, pos
    integer :: im1, jm1
    real(double) :: acc_dble1, acc_dble2
    real(default) :: acc1, acc2
    real(default) :: born

    mom = create_blha_momentum_array (p)
    if (ren_scale == 0.0) then
       mu = sqrt (2*p(1)*p(2))
    else
       mu = ren_scale
    end if
    select type (driver => object%driver)
    type is (gosam_driver_t)
       call driver%set_alpha_s (alpha_s)
       if (allocated (object%i_born)) then
          call driver%gosam_olp_eval2 (object%i_born(i_flv), &
                                       mom, mu, r, acc_dble1)
          born = r(4)
       end if
       if (present (born_out)) born_out = born
       call driver%gosam_olp_eval2 (object%i_cc(i_flv), &
                                    mom, mu, r, acc_dble2)
    end select
    do j = 1, size (p)
      do i = 1, j
        if (i <= 2 .or. j <= 2) then
          born_cc (i,j) = 0._default
        else if (i == j) then
          born_cc (i,j) = -cf*born
        else
          im1 = i-1; jm1 = j-1
          pos = im1 + jm1*(jm1-1)/2 + 1
          born_cc (i,j) = -r(pos)
        end if
        born_cc (j,i) = born_cc (i,j)
      end do
    end do
    acc1 = acc_dble1; acc2 = acc_dble2
    if (acc1 > object%maximum_accuracy .or. &
        acc2 > object%maximum_accuracy) then
      bad_point = .true.
    end if
  end subroutine prc_gosam_compute_sqme_cc

\eatline
\nwendcode{}\nwbegindocs{199}\nwdocspar
\nwenddocs{}\nwbegincode{200}\moddef{prc gosam: prc gosam: TBP}\plusendmoddef
  procedure :: compute_sqme_sc => prc_gosam_compute_sqme_sc
\nwendcode{}\nwbegincode{201}\moddef{prc gosam: procedures}\plusendmoddef
  subroutine prc_gosam_compute_sqme_sc (object, &
                i_flv, em, p, ren_scale_in, alpha_s, &
            me_sc, bad_point)
    class(prc_gosam_t), intent(inout) :: object
    integer, intent(in) :: i_flv
    integer, intent(in) :: em
    type(vector4_t), intent(in), dimension(:) :: p
    real(default), intent(in) :: ren_scale_in
    real(default), intent(in) :: alpha_s
    complex(default), intent(out) :: me_sc
    logical, intent(out) :: bad_point
    real(double), dimension(GOSAM_MOMENTUM_LIMIT) :: mom
    real(double), dimension(GOSAM_RESULTS_LIMIT) :: r
    real(double) :: ren_scale
    integer :: i, igm1, n
    integer :: pos_real, pos_imag
    real(double) :: acc_dble
    real(default) :: acc

    me_sc = cmplx(0,0,default)
    mom = create_blha_momentum_array (p)
    if (ren_scale_in == 0.0) then
      ren_scale = sqrt (2*p(1)*p(2))
    else
      ren_scale = ren_scale_in
    end if
    select type (driver => object%driver)
    type is (gosam_driver_t)
       call driver%set_alpha_s (alpha_s)
       call driver%gosam_olp_eval2 (object%i_sc(i_flv), &
                                    mom, ren_scale, r, acc_dble)
    end select

\nwendcode{}\nwbegindocs{202}The BLHA-interface conventions require the quantity $S_{ij} = \langle M_{i,+}|T_iT_j|M_{i,-}\rangle$ to be
  produced, where $i$ is the position of the splitting gluon. However, $\tilde{M} = \langle M_{i,-}|M_{i,+}\rangle$ 
  is needed. This can be obtained using color conservation, $\sum_{j} T_j|M\rangle = 0$, so that
  \begin{equation*}
    \sum_{j \neq i} S_{ij} = -\langle M_{i,+}|T_i^2|M_{i,-}\rangle = -C_A \langle M_{i,+}|M_{i,-}\rangle = -C_A \tilde{M}^*
  \end{equation*}
  According to BLHA conventions, the real part of $S_{ij}$ is located at positions $2i + 2nj$ in the output array, where 
  $n$ denotes the number of external particles and the enumeration of particles starts at zero.
  The subsequent position, i.e. $2i + 2nj + 1$ is designated to the imaginary part of $S_{ij}$. Note that, since the first 
  array position is 1, the implemented position association deviates from the above one in the addition of 1.
\nwenddocs{}\nwbegincode{203}\moddef{prc gosam: procedures}\plusendmoddef
    igm1 = em-1
    n = size(p)
    do i = 0, n-1
      pos_real = 2*igm1 + 2*n*i + 1
      pos_imag = pos_real + 1
      me_sc = me_sc + cmplx (r(pos_real), r(pos_imag), default)
    end do

    me_sc = -conjg(me_sc)/CA
  
    acc = acc_dble
    if (acc > object%maximum_accuracy) bad_point = .true.
  end subroutine prc_gosam_compute_sqme_sc

\eatline
\nwendcode{}\nwbegindocs{204}\nwdocspar
\nwenddocs{}\nwbegincode{205}\moddef{prc gosam: prc gosam: TBP}\plusendmoddef
  procedure :: allocate_workspace => prc_gosam_allocate_workspace
\nwendcode{}\nwbegincode{206}\moddef{prc gosam: procedures}\plusendmoddef
  subroutine prc_gosam_allocate_workspace (object, tmp)
    class(prc_gosam_t), intent(in) :: object
    class(workspace_t), intent(inout), allocatable :: tmp
    allocate (gosam_state_t :: tmp)
  end subroutine prc_gosam_allocate_workspace

\eatline
\nwendcode{}\nwbegindocs{207}\nwdocspar
\nwenddocs{}\nwbegincode{208}\moddef{prc gosam: prc gosam: TBP}\plusendmoddef
  procedure :: get_alpha_s => prc_gosam_get_alpha_s
\nwendcode{}\nwbegincode{209}\moddef{prc gosam: procedures}\plusendmoddef
  function prc_gosam_get_alpha_s (object, tmp) result (alpha)
    class(prc_gosam_t), intent(in) :: object
    class(workspace_t), intent(in), allocatable :: tmp
    real(default) :: alpha
    if (allocated (tmp)) then
      select type (tmp)
      type is (gosam_state_t)
        alpha = tmp%alpha_qcd
      end select
    else
      alpha = 0._default 
    end if
  end function prc_gosam_get_alpha_s

\eatline
\nwendcode{}\nwbegindocs{210}\nwdocspar
\nwenddocs{}\nwbegincode{211}\moddef{prc gosam: prc gosam: TBP}\plusendmoddef
  procedure :: read_olc_file => prc_gosam_read_olc_file
\nwendcode{}\nwbegincode{212}\moddef{prc gosam: procedures}\plusendmoddef
  subroutine prc_gosam_read_olc_file (object, flavors)
    class(prc_gosam_t), intent(inout) :: object
    integer, intent(in), dimension(:,:) :: flavors
    integer, dimension(20) :: amp_type, flv_index, label
    integer :: i_proc
    select type (driver => object%driver)
    type is (gosam_driver_t)
       call driver%read_olc_file (flavors, amp_type, flv_index, label)
    end select
    do i_proc = 1, size (amp_type)
       if (amp_type (i_proc) < 0) exit
       select case (amp_type (i_proc))
       case (BLHA_AMP_TREE)
          if (allocated (object%i_born)) then
             object%i_born(flv_index(i_proc)) = label(i_proc)
          
          else if (allocated (object%i_real)) then
             object%i_real(flv_index(i_proc)) = label(i_proc)
          else 
             call msg_fatal ("Tree matrix element present, &
                             &but neither Born nor real indices are allocated!")
          end if
       case (BLHA_AMP_CC)
          if (allocated (object%i_cc)) then
             object%i_cc(flv_index(i_proc)) = label(i_proc)
          else
             call msg_fatal ("Color-correlated matrix element present, &
                              &but cc-indices are not allocated!")
          end if
       case (BLHA_AMP_SC)
          if (allocated (object%i_sc)) then
             object%i_sc(flv_index(i_proc)) = label(i_proc)
          else
             call msg_fatal ("Spin-correlated matrix element present, &
                             &but sc-indices are not allocated!")
          end if
       case (BLHA_AMP_LOOP)
          if (allocated (object%i_virt)) then
             object%i_virt(flv_index(i_proc)) = label(i_proc)
          else
             call msg_fatal ("Loop matrix element present, &
                             &but virt-indices are not allocated!")
          end if
       case default
          call msg_fatal ("Undefined amplitude type")
       end select
    end do
  end subroutine prc_gosam_read_olc_file

\eatline
\nwendcode{}\nwbegindocs{213}\nwdocspar
\nwenddocs{}\nwbegincode{214}\moddef{prc gosam: prc gosam: TBP}\plusendmoddef
  procedure :: set_particle_properties => prc_gosam_set_particle_properties
\nwendcode{}\nwbegincode{215}\moddef{prc gosam: procedures}\plusendmoddef
  subroutine prc_gosam_set_particle_properties (object, model) 
    class(prc_gosam_t), intent(inout) :: object
    class(model_data_t), intent(in), target :: model
    type(flavor_t), dimension(:), allocatable :: flv
    real(default), dimension(:), allocatable :: masses
    real(default), dimension(:), allocatable :: widths
    integer :: i_part, i_flv
    integer :: n_part, n_flv
    n_part = size (object%data%flv_state,1)
    n_flv = size (object%data%flv_state,2)
    allocate (flv (n_part), masses (n_part), widths (n_part))
    do i_flv = 1, n_flv
       associate (i_pdg => object%data%flv_state (:,i_flv))
          call flv%init (i_pdg, model)
          masses = flv%get_mass ()
          widths = flv%get_width ()
          select type (driver => object%driver)
          type is (gosam_driver_t)
             do i_part = 1, n_part
                if (masses(i_part) /= 0._default) &
                   call driver%set_mass_and_width (abs(i_pdg(i_part)), &
                                                   mass=masses(i_part))
                if (abs(i_pdg(i_part)) == 6) &
                   call driver%set_mass_and_width (abs(i_pdg(i_part)), &
                                                   width=widths(i_part))
             end do
          end select
       end associate
    end do
  end subroutine prc_gosam_set_particle_properties

\eatline
\nwendcode{}\nwbegindocs{216}          
\nwenddocs{}\nwbegincode{217}\moddef{prc gosam: gosam state: TBP}\endmoddef
  procedure :: write => gosam_state_write
\nwendcode{}\nwbegincode{218}\moddef{prc gosam: procedures}\plusendmoddef
  subroutine gosam_state_write (object, unit)
    class(gosam_state_t), intent(in) :: object
    integer, intent(in), optional :: unit
    call msg_bug ("gosam_state_write: What to write?")
  end subroutine gosam_state_write

\eatline
\nwendcode{}\nwbegindocs{219}\nwdocspar
\nwenddocs{}\nwbegincode{220}\moddef{prc gosam: gosam state: TBP}\plusendmoddef
  procedure :: reset_new_kinematics => gosam_state_reset_new_kinematics
\nwendcode{}\nwbegincode{221}\moddef{prc gosam: procedures}\plusendmoddef
  subroutine gosam_state_reset_new_kinematics (object)
    class(gosam_state_t), intent(inout) :: object
    object%new_kinematics = .true.
  end subroutine gosam_state_reset_new_kinematics

\nwendcode{}\nwbegindocs{222}\nwdocspar
%
\nwenddocs{}\nwbegincode{223}\moddef{prc gosam: procedures}\plusendmoddef
  function check_golem_installation (os_data) result (res)
     type(os_data_t), intent(in) :: os_data
     logical :: res
     type(string_t) :: libavh_olo, libgolem
     integer :: i
     res = .true.
     do i = 1, size (lib_suffix)
       libavh_olo = os_data%prefix // '/lib/libavh_olo' // &
            trim (lib_suffix (i))
       libgolem = os_data%prefix // '/lib/libgolem' // trim (lib_suffix (i))
       res = res .and. (os_file_exist (libavh_olo) .and. &
            os_file_exist (libgolem))
       if (.not. res) exit
     end do
   end function check_golem_installation

\eatline
\nwendcode{}\nwbegindocs{224}\nwdocspar
\nwenddocs{}\nwbegincode{225}\moddef{prc gosam: procedures}\plusendmoddef
  function check_samurai_installation (os_data) result (res)
    type(os_data_t), intent(in) :: os_data
    logical :: res
    type(string_t) :: libsamurai
    integer :: i
    res = .true.
    do i = 1, size (lib_suffix)
      libsamurai = os_data%prefix // '/lib/libsamurai' // trim (lib_suffix(i))
      res = os_file_exist (libsamurai)
      if (.not. res) exit
    end do
  end function check_samurai_installation 

\eatline
\nwendcode{}\nwbegindocs{226}\nwdocspar
\nwenddocs{}\nwbegincode{227}\moddef{prc gosam: procedures}\plusendmoddef
  function check_ninja_installation (os_data) result (res)
    type(os_data_t), intent(in) :: os_data
    logical :: res
    type(string_t) :: libninja
    integer :: i
    res = .true.
    do i = 1, size (lib_suffix)
      libninja = os_data%prefix // '/lib/libninja' // trim (lib_suffix(i))
      res = os_file_exist (libninja)
      if (.not. res) exit
    end do
  end function check_ninja_installation

\eatline
\nwendcode{}\nwbegindocs{228}\nwdocspar
\nwenddocs{}\nwbegincode{229}\moddef{prc gosam: public}\plusendmoddef
  public :: gosam_template_t
\nwendcode{}\nwbegincode{230}\moddef{prc gosam: types}\plusendmoddef
  type :: gosam_template_t
    integer :: I_REAL = 1
    integer :: I_LOOP = 2
    integer :: I_SUB = 3
    logical, dimension(3) :: compute_component
  contains
  \LA{}prc gosam: gosam template: TBP\RA{}
  end type gosam_template_t

\eatline
\nwendcode{}\nwbegindocs{231}\nwdocspar
\nwenddocs{}\nwbegincode{232}\moddef{prc gosam: gosam template: TBP}\endmoddef
  procedure :: init => gosam_template_init
\nwendcode{}\nwbegincode{233}\moddef{prc gosam: procedures}\plusendmoddef
  subroutine gosam_template_init (template)
    class(gosam_template_t), intent(inout) :: template
    template%compute_component = .false.
  end subroutine gosam_template_init

\eatline
\nwendcode{}\nwbegindocs{234}\nwdocspar
\nwenddocs{}\nwbegincode{235}\moddef{prc gosam: gosam template: TBP}\plusendmoddef
  procedure :: set_loop => gosam_template_set_loop
  procedure :: set_subtraction => gosam_template_set_subtraction
  procedure :: set_real_trees => gosam_template_set_real_trees
\nwendcode{}\nwbegincode{236}\moddef{prc gosam: procedures}\plusendmoddef
  subroutine gosam_template_set_loop (template, val)
    class(gosam_template_t), intent(inout) :: template
    logical, intent(in) :: val
    template%compute_component(template%I_LOOP) = val
  end subroutine gosam_template_set_loop

  subroutine gosam_template_set_subtraction (template, val)
    class(gosam_template_t), intent(inout) :: template
    logical, intent(in) :: val
    template%compute_component (template%I_SUB) = val
  end subroutine gosam_template_set_subtraction

  subroutine gosam_template_set_real_trees (template, val)
    class(gosam_template_t), intent(inout) :: template
    logical, intent(in) :: val
    template%compute_component (template%I_REAL) = val
  end subroutine gosam_template_set_real_trees

\eatline
\nwendcode{}\nwbegindocs{237}\nwdocspar
\nwenddocs{}\nwbegincode{238}\moddef{prc gosam: gosam template: TBP}\plusendmoddef
  procedure :: compute_loop => gosam_template_compute_loop
  procedure :: compute_subtraction => gosam_template_compute_subtraction
  procedure :: compute_real_trees => gosam_template_compute_real_trees
\nwendcode{}\nwbegincode{239}\moddef{prc gosam: procedures}\plusendmoddef
  function gosam_template_compute_loop (template) result (val)
    class(gosam_template_t), intent(in) :: template
    logical :: val
    val = template%compute_component (template%I_LOOP)
  end function gosam_template_compute_loop  

  function gosam_template_compute_subtraction (template) result (val)
    class(gosam_template_t), intent(in) :: template
    logical :: val
    val = template%compute_component (template%I_SUB)
  end function gosam_template_compute_subtraction

  function gosam_template_compute_real_trees (template) result (val)
    class(gosam_template_t), intent(in) :: template
    logical :: val
    val = template%compute_component (template%I_REAL)
  end function gosam_template_compute_real_trees

\eatline
\nwendcode{}\nwbegincode{240}\moddef{prc gosam: gosam template: TBP}\plusendmoddef
  procedure :: check => gosam_template_check
\nwendcode{}\nwbegincode{241}\moddef{prc gosam: procedures}\plusendmoddef
  function gosam_template_check (template) result (val)
    class(gosam_template_t), intent(in) :: template
    logical :: val
    val = count (template%compute_component) == 1
  end function gosam_template_check

\eatline
\nwendcode{}\nwbegindocs{242}\nwdocspar
\nwenddocs{}\nwbegincode{243}\moddef{prc gosam: gosam template: TBP}\plusendmoddef
  procedure :: reset => gosam_template_reset
\nwendcode{}\nwbegincode{244}\moddef{prc gosam: procedures}\plusendmoddef
  subroutine gosam_template_reset (template)
    class(gosam_template_t), intent(inout) :: template
    template%compute_component = .false.
  end subroutine gosam_template_reset

\eatline
\nwendcode{}\nwbegindocs{245}\nwdocspar
\nwenddocs{}\nwfilename{../../../src/fks/fks.nw}\nwbegindocs{0}% -*- ess-noweb-default-code-mode: f90-mode; noweb-default-code-mode: f90-mode; -*-
% WHIZARD code as NOWEB source: matrix elements and process libraries

\nwenddocs{}\nwbegindocs{1}\nwdocspar
%%%%%%%%%%%%%%%%%%%%%%%%%%%%%%%%%%%%%%%%%%%%%%%%%%%%%%%%%%%%%%%%%%%%%%%%
\chapter{FKS Subtraction Scheme}

The code in this chapter implements the FKS subtraction scheme for use
with \whizard.

These are the modules:
\begin{description}
\item[fks\_regions]
  Given a process definition, identify singular regions in the
  associated phase space.
\item[virtual]
  Handle the virtual correction matrix element.
\item[real\_subtraction]
  Handle the real-subtraction matrix element.
\item[nlo\_data]
  Manage the subtraction objects.
\item[phs\_fks]
  Phase-space parameterization with modifications for the FKS scheme.
\end{description}

This chapter deals with next-to-leading order contributions to cross sections.
Basically, there are three major issues to be adressed: The creation
of the $N+1$-particle flavor structure, the construction of the
$N+1$-particle phase space and the actual calculation of the real- and
virtual-subtracted matrix elements. The first is dealt with using the
{\tt{}auto{\char95}components} class, and it will be shown that the second
and third issue are connected in FKS subtraction. 

\section{Brief outline of FKS subtraction}
{\em In the current state, this discussion is only concerned with
  lepton collisions. For hadron collisions, renormalization of parton
  distributions has to be taken into account. Further, for QCD
  corrections, initial-state radiation is necessarily
  present. However, most quantities have so far been only constructed
  for final-state emissions}

The aim is to calculate the next-to-leading order cross section
according to 
\begin{equation*}
  d\sigma_{\rm{NLO}} = \mathcal{B} + \mathcal{V} +
  \mathcal{R}d\Phi_{\rm{rad}}. 
\end{equation*}
Analytically, the divergences, in terms of poles in the complex
quantity $\varepsilon = 2-d/2$, cancel. However, this is in general
only valid in an arbitrary, comlex number of dimensions. This is,
roughly, the content of the KLN-theorem. \whizard, as any
other numerical program, is confined to four dimensions. We will
assume that the KLN-theorem is valid and that there exist subtraction
terms $\mathcal{C}$ such that 
\begin{equation*}
  d\sigma_{\rm{NLO}} = \mathcal{B} + \underbrace{\mathcal{V} +
    \mathcal{C}}_{\text{finite}} + \underbrace{\mathcal{R} -
    \mathcal{C}}_{\text{finite}}, 
\end{equation*}
i.e. the subtraction terms correspond to the divergent limits of the
real and virtual matrix element.

Because $\mathcal{C}$ subtracts the divergences of $\mathcal{R}$ as
well as those of $\mathcal{V}$, it suffices to consider one of them,
so we focus on $\mathcal{R}$. For this purpose, $\mathcal{R}$ is
rewritten, 
\begin{equation*}
  \mathcal{R} = \frac{1}{\xi^2}\frac{1}{1-y} \left(\xi^2
    (1-y)\mathcal{R}\right) =
  \frac{1}{\xi^2}\frac{1}{1-y}\tilde{\mathcal{R}}, 
\end{equation*}
with $\xi = \left(2k_{\rm{rad}}^0\right)/\sqrt{s}$ and $y =
\cos\theta$, where $k_{\rm{rad}}^0$ denotes the energy of the radiated
parton and $\theta$ is the angle between emitter and radiated
parton. $\tilde{\mathcal{R}}$ is finite, therefore the whole
singularity structure is contained in the prefactor
$\xi^{-2}(1-y)^{-1}$. Combined with the d-dimensional phase space
element, 
\begin{equation*}
  \frac{d^{d-1}k}{2k^0(2\pi)^{d-1}} =
  \frac{s^{1-\varepsilon}}{(4\pi)^{d-1}}\xi^{1-2\varepsilon}\left(1-y^2\right)^{-\varepsilon} 
  d\xi dy d\Omega^{d-2}, 
\end{equation*}
this yields 
\begin{equation*}
  d\Phi_{\rm{rad}} \mathcal{R} = dy (1-y)^{-1-\varepsilon} d\xi
  \xi^{-1-2\varepsilon} \tilde{R}. 
\end{equation*}
This can further be rewritten in terms of plus-distributions, 
\begin{align*}
\xi^{-1-2\varepsilon} &= -\frac{1}{2\varepsilon}\delta(\xi) +
\left(\frac{1}{\xi}\right)_+ -
2\varepsilon\left(\frac{\log\xi}{\xi}\right)_+ +
\mathcal{O}(\varepsilon^2),\\ 
(1-y)^{-1-\varepsilon} &= -\frac{2^{-\varepsilon}}{\varepsilon}
\delta(1-y) + \left(\frac{1}{1-y}\right)_+ - \varepsilon
\left(\frac{1}{1-y}\right)_+\log(1-y) + \mathcal{O}(\varepsilon^2), 
\end{align*}
(imagine that all this is written inside of integrals, which are
spared for ease of notation) such that 
\begin{align*}
d\Phi_{\rm{rad}} \mathcal{R} &= -\frac{1}{2\varepsilon} dy
(1-y)^{-1-\varepsilon}\tilde{R} (0,y) -
d\xi\left[\frac{2^{-\varepsilon}}{\varepsilon}\left(\frac{1}{\xi}\right)_+
  - 2\left(\frac{\log\xi}{\xi}\right)_+\right] \tilde{R}(\xi,1) \\ 
                              &+ dy d\xi \left(\frac{1}{\xi}\right)_+
                              \left(\frac{1}{1-y}\right)_+
                              \tilde{R}(\xi, y) +
                              \mathcal{O}(\varepsilon).\\ 
\end{align*}
The summand in the second line is of order $\mathcal{O}(1)$ and is the
only one to reproduce $\mathcal{R}{\xi,y}$. It thus constitutes the
sum of the real matrix element and the corresponding counterterms. 
The first summand consequently consists of the subtraction terms to
the virtual matrix elements. Above formula thus allows to calculate
all quantities to render the matrix elements finite.

\clearpage
%%%%%%%%%%%%%%%%%%%%%%%%%%%%%%%%%%%%%%%%%%%%%%%%%%%%%%%%%%%%%%%%%%%%%%%%
\section{Identifying singular regions}
In the FKS subtraction scheme, the phase space is decomposed into
disjoint singular regions, such that 
\begin{equation}
\label{eq:S_complete}
  \sum_i \mathcal{S}_i + \sum_{ij}\mathcal{S}_{ij} = 1.
\end{equation}
The quantities $\mathcal{S}_i$ and $\mathcal{S}_{ij}$ are functions of
phase space corresponding to a pair of particles indices which can
make up a divergent phase space region. We call such an index pair a
fundamental tuple. For example, the process $e^+ \, e^- \rightarrow u
\, \bar{u} \, g$ has two singular regions, $(3,5)$ and $(4,5)$,
indicating that the gluon can be soft or collinear with respect to
either the quark or the anti-quark. Therefore, the functions $S_{ij}$
have to be chosen in such a way that their contribution makes up most
of \eqref{eq:S_complete} in phase-space configurations where
(final-state) particle $j$ is collinear to particle $i$ or/and
particle $j$ is soft. The functions $S_i$ is the corresponding
quantity for initial-state divergences.

As a singular region we understand the collection of real flavor
structures associated with an emitter and a list of all possible
fundamental tuples. As an example, consider the process $e^+ \, e^-
\rightarrow u \, \bar{u} \, g$. At next-to-leading order, processes
with an additionally radiated particle have to be considered. In this
case, these are $e^+ \, e^- \rightarrow u \, \bar{u}, \, g \, g$, 
and $e^+ \, e^- \rightarrow u \, \bar{u} \, u \, \bar{u}$ (or the same
process with any other quark). Table \ref{table:singular regions} sums
up all possible singular regions for this problem. 
\begin{table}
\label{table:singular regions}
\begin{tabular}{|c|c|c|c|}
  \hline
  \texttt{alr} & \texttt{flst\_alr} & \texttt{emi} &
  \texttt{ftuple\_list}\\ \hline 
  1 & [-11,11,2,-2,21,21] & 3 & {(3,5), (3,6), (4,5), (4,6), (5,6)} \\ \hline
  2 & [-11,11,2,-2,21,21] & 4 & {(3,5), (3,6), (4,5), (4,6), (5,6)} \\ \hline
  3 & [-11,11,2,-2,21,21] & 5 & {(3,5), (3,6), (4,5), (4,6), (5,6)} \\ \hline
  4 & [-11,11,2,-2,2,-2]  & 5 & {(5,6)} \\ 
  \hline
\end{tabular}
\caption{List of singular regions. The particles are represented by
  their PDG codes. The third column contains the emitter for the
  specific singular region. For the process involving an additional
  gluon, the gluon can either be emitted from one of the quarks or
  from the first gluon. Each emitter yields the same list of
  fundamental tuples, five in total. The last singular region
  corresponds to the process where the gluon splits up into two
  quarks. Here, there is only one fundamental tuple, corresponding to
  a singular configuration of the momenta of the additional quarks.} 
\end{table}
\\
\begin{table}
\label{table:ftuples and flavors}
\begin{tabular}{|c|c|c|c|}
  \hline
  \texttt{alr} & \texttt{ftuple} & \texttt{emitter} &
  \texttt{flst\_alr} \\ \hline 
  1 & $(3,5)$ & 5 & [-11,11,-2,21,2,21] \\ \hline
  2 & $(4,5)$ & 5 & [-11,11,2,21,-2,21] \\ \hline
  3 & $(3,6)$ & 5 & [-11,11,-2,21,2,21] \\ \hline
  4 & $(4,6)$ & 5 & [-11,11,2,21,-2,21] \\ \hline
  5 & $(5,6)$ & 5 & [-11,11,2,-2,21,21] \\ \hline
  6 & $(5,6)$ & 5 & [-11,11,2,-2,2,-2] \\ \hline
\end{tabular}
\caption{Initial list of singular regions}
\end{table}
Thus, during the preparation of a NLO-calculation, the possible
singular regions have to be identified. {\tt{}fks{\char95}regions.f90} deals 
with this issue.

\clearpage
%%%%%%%%%%%%%%%%%%%%%%%%%%%%%%%%%%%%%%%%%%%%%%%%%%%%%%%%%%%%%%%%%%%%%%%%
\section{FKS Regions}
\nwenddocs{}\nwbegincode{2}\moddef{\code{}fks{\char95}regions.f90\edoc{}}\endmoddef
\LA{}File header\RA{}

module fks_regions

  use kinds
  use io_units
\LA{}Use strings\RA{}
  use constants
  use diagnostics
  use flavors
  use process_constants
  use lorentz
  use pdg_arrays
  use model_data
  use physics_defs

\LA{}Standard module head\RA{}

\LA{}fks regions: public\RA{}

\LA{}fks regions: parameters\RA{}

\LA{}fks regions: types\RA{}

\LA{}fks regions: interfaces\RA{}

contains

\LA{}fks regions: procedures\RA{}

end module fks_regions
\eatline
\nwendcode{}\nwbegindocs{3}There are three fundamental splitting types: $q \rightarrow qg$, $g \rightarrow gg$ and 
$g \rightarrow qq$.
\nwenddocs{}\nwbegincode{4}\moddef{fks regions: parameters}\endmoddef
  integer, parameter :: UNDEFINED_SPLITTING = 0
  integer, parameter :: Q_TO_QG = 1
  integer, parameter :: G_TO_GG = 2
  integer, parameter :: G_TO_QQ = 3

\nwendcode{}\nwbegindocs{5}\nwdocspar
\nwenddocs{}\nwbegindocs{6}We group the indices of the emitting and the radiated particle in
the {\tt{}ftuple}-object. 
\nwenddocs{}\nwbegincode{7}\moddef{fks regions: public}\endmoddef
  public :: ftuple_t
\nwendcode{}\nwbegincode{8}\moddef{fks regions: types}\endmoddef
  type :: ftuple_t 
    integer, dimension(2) :: ireg
    integer :: splitting_type
  contains
  \LA{}fks regions: ftuple: TBP\RA{} 
  end type ftuple_t
  
\eatline
\nwendcode{}\nwbegindocs{9}\nwdocspar
\nwenddocs{}\nwbegincode{10}\moddef{fks regions: ftuple: TBP}\endmoddef
  procedure :: write => ftuple_write
\nwendcode{}\nwbegincode{11}\moddef{fks regions: procedures}\endmoddef
  subroutine ftuple_write (ftuple, unit)
    class(ftuple_t), intent(in) :: ftuple
    integer, intent(in), optional :: unit
    integer :: u
    u = given_output_unit (unit); if (u < 0) return
    write (u, "(A1,I1,A1,I1,A1)") &
         '(', ftuple%ireg(1), ',', ftuple%ireg(2), ')'
  end subroutine ftuple_write

\eatline
\nwendcode{}\nwbegindocs{12}\nwdocspar
\nwenddocs{}\nwbegincode{13}\moddef{fks regions: ftuple: TBP}\plusendmoddef
  procedure :: get => ftuple_get
\nwendcode{}\nwbegincode{14}\moddef{fks regions: procedures}\plusendmoddef
  subroutine ftuple_get (ftuple, pos1, pos2)
    class(ftuple_t), intent(in) :: ftuple
    integer, intent(out) :: pos1, pos2
    pos1 = ftuple%ireg(1)
    pos2 = ftuple%ireg(2)
  end subroutine ftuple_get

\eatline
\nwendcode{}\nwbegindocs{15}\nwdocspar
\nwenddocs{}\nwbegincode{16}\moddef{fks regions: ftuple: TBP}\plusendmoddef
  procedure :: set => ftuple_set
\nwendcode{}\nwbegincode{17}\moddef{fks regions: procedures}\plusendmoddef
  subroutine ftuple_set (ftuple, pos1, pos2)
    class(ftuple_t), intent(inout) :: ftuple
    integer, intent(in) ::  pos1, pos2
    ftuple%ireg(1) = pos1
    ftuple%ireg(2) = pos2
  end subroutine ftuple_set

\eatline
\nwendcode{}\nwbegindocs{18}\nwdocspar
\nwenddocs{}\nwbegincode{19}\moddef{fks regions: ftuple: TBP}\plusendmoddef
  procedure :: determine_splitting_type => ftuple_determine_splitting_type
\nwendcode{}\nwbegincode{20}\moddef{fks regions: procedures}\plusendmoddef
  subroutine ftuple_determine_splitting_type (ftuple, flv, i, j)
    class(ftuple_t), intent(inout) :: ftuple
    type(flv_structure_t), intent(in) :: flv
    integer, intent(in) :: i, j
    associate (flst => flv%flst)
       if (flst(i) == GLUON .and. flst(j) == GLUON) then
          ftuple%splitting_type = G_TO_GG
       else if (flst(i)+flst(j) == 0 &
             .and. is_quark (abs(flst(i)))) then
          ftuple%splitting_type = G_TO_QQ
       else if (is_quark(abs(flst(i))) .and. flst(j) == GLUON &
             .or. is_quark(abs(flst(j))) .and. flst(i) == GLUON) then
          ftuple%splitting_type = Q_TO_QG
       else
          ftuple%splitting_type = UNDEFINED_SPLITTING
       end if
    end associate
  end subroutine ftuple_determine_splitting_type

\eatline
\nwendcode{}\nwbegindocs{21}Each singular region can have a different number of
emitter-radiation pairs. This is coped with using the linked list
{\tt{}ftuple{\char95}list}. 
\nwenddocs{}\nwbegincode{22}\moddef{fks regions: types}\plusendmoddef
  type :: ftuple_list_t
    integer :: index
    type(ftuple_t) :: ftuple
    type(ftuple_list_t), pointer :: next
    type(ftuple_list_t), pointer :: prev
    type(ftuple_list_t), pointer :: equiv
  contains
   \LA{}fks regions: ftuple list: TBP\RA{}
  end type ftuple_list_t

\eatline
\nwendcode{}\nwbegindocs{23}\nwdocspar
\nwenddocs{}\nwbegincode{24}\moddef{fks regions: ftuple list: TBP}\endmoddef
  procedure :: init => ftuple_list_init
\nwendcode{}\nwbegincode{25}\moddef{fks regions: procedures}\plusendmoddef
  subroutine ftuple_list_init (list)
    class(ftuple_list_t), intent(inout) :: list
    list%index = 0
    nullify (list%next)
    nullify (list%prev)
    nullify (list%equiv)
  end subroutine ftuple_list_init

\eatline
\nwendcode{}\nwbegindocs{26}\nwdocspar
\nwenddocs{}\nwbegincode{27}\moddef{fks regions: ftuple list: TBP}\plusendmoddef
  procedure :: write => ftuple_list_write
\nwendcode{}\nwbegincode{28}\moddef{fks regions: procedures}\plusendmoddef
  subroutine ftuple_list_write (list)
    class(ftuple_list_t), intent(in), target :: list
    type(ftuple_list_t), pointer :: current
    select type (list)
    type is (ftuple_list_t)
    current => list
    do
      call current%ftuple%write
      if (associated (current%next)) then
        current => current%next
      else
        exit
      end if
    end do
    end select
  end subroutine ftuple_list_write

\eatline
\nwendcode{}\nwbegindocs{29}\nwdocspar
\nwenddocs{}\nwbegincode{30}\moddef{fks regions: ftuple list: TBP}\plusendmoddef
  procedure :: append => ftuple_list_append
\nwendcode{}\nwbegincode{31}\moddef{fks regions: procedures}\plusendmoddef
  subroutine ftuple_list_append (list, ftuple) 
   class(ftuple_list_t), intent(inout), target :: list
   type(ftuple_t), intent(in) :: ftuple
   type(ftuple_list_t), pointer :: current

   select type (list)
   type is (ftuple_list_t)
   if (list%index == 0) then
      nullify(list%next)
      list%index = 1
      list%ftuple = ftuple
   else
      current => list
      do
       if (associated (current%next)) then
         current => current%next
       else
         allocate (current%next)
         nullify (current%next%next)
         nullify (current%next%equiv)
         current%next%prev => current
         current%next%index = current%index + 1
         current%next%ftuple = ftuple
         exit
       end if
     end do
   end if
   end select
  end subroutine ftuple_list_append

\eatline
\nwendcode{}\nwbegindocs{32}\nwdocspar
\nwenddocs{}\nwbegincode{33}\moddef{fks regions: ftuple list: TBP}\plusendmoddef
  procedure :: get_n_tuples => ftuple_list_get_n_tuples
\nwendcode{}\nwbegincode{34}\moddef{fks regions: procedures}\plusendmoddef
  function ftuple_list_get_n_tuples (list) result(n_tuples)
    class(ftuple_list_t), intent(inout), target :: list
    integer :: n_tuples
    type(ftuple_list_t), pointer :: current
    select type (list)
    type is (ftuple_list_t)
      current => list
      n_tuples = 1
      do 
        if (associated (current%next)) then
          current => current%next
          n_tuples = n_tuples + 1
        else
          exit
        end if
       end do
    end select
  end function ftuple_list_get_n_tuples

\eatline
\nwendcode{}\nwbegindocs{35}\nwdocspar
\nwenddocs{}\nwbegincode{36}\moddef{fks regions: ftuple list: TBP}\plusendmoddef
  procedure :: get_entry => ftuple_list_get_entry
\nwendcode{}\nwbegincode{37}\moddef{fks regions: procedures}\plusendmoddef
  function ftuple_list_get_entry(list, index) result(entry)
   class(ftuple_list_t), intent(inout), target :: list
   integer, intent(in) :: index
   type(ftuple_list_t), pointer :: entry
   type(ftuple_list_t), pointer :: current
   integer :: i
   select type (list)
   type is (ftuple_list_t)
   current => list
   if (index <= list%get_n_tuples ()) then
   if (index == 1) then
     entry => current
   else
     do i=1,index-1
       current => current%next
     end do
     entry => current
   end if
   else
     call msg_fatal &
          ("Index must be smaller or equal than the total number of regions!")
   end if
   end select  
  end function ftuple_list_get_entry

\eatline
\nwendcode{}\nwbegindocs{38}\nwdocspar
\nwenddocs{}\nwbegincode{39}\moddef{fks regions: ftuple list: TBP}\plusendmoddef
  procedure :: get_ftuple => ftuple_list_get_ftuple
\nwendcode{}\nwbegincode{40}\moddef{fks regions: procedures}\plusendmoddef
  function ftuple_list_get_ftuple (list, index)  result (ftuple)
    class(ftuple_list_t), intent(inout) :: list
    integer, intent(in) :: index
    type(ftuple_t) :: ftuple
    type(ftuple_list_t) :: entry
    entry = list%get_entry (index)
    ftuple = entry%ftuple
  end function ftuple_list_get_ftuple

\eatline
\nwendcode{}\nwbegindocs{41}\nwdocspar
\nwenddocs{}\nwbegincode{42}\moddef{fks regions: ftuple list: TBP}\plusendmoddef
  procedure :: set_equiv => ftuple_list_set_equiv
\nwendcode{}\nwbegincode{43}\moddef{fks regions: procedures}\plusendmoddef
  subroutine ftuple_list_set_equiv (list, i1, i2)
    class(ftuple_list_t), intent(inout) :: list
    integer, intent(in) :: i1, i2
    type(ftuple_list_t), pointer :: list1, list2
    select type (list)
    type is (ftuple_list_t)
       list1 => list%get_entry (i1)
       list2 => list%get_entry (i2)
       do
          if (associated (list1%equiv)) then
             list1 => list1%equiv
          else
             exit
          end if
       end do
       list1%equiv => list2
    end select
  end subroutine ftuple_list_set_equiv

\eatline
\nwendcode{}\nwbegindocs{44}\nwdocspar
\nwenddocs{}\nwbegincode{45}\moddef{fks regions: ftuple list: TBP}\plusendmoddef
  procedure :: check_equiv => ftuple_list_check_equiv
\nwendcode{}\nwbegincode{46}\moddef{fks regions: procedures}\plusendmoddef
  function ftuple_list_check_equiv(list, i1, i2) result(eq)
    class(ftuple_list_t), intent(inout) :: list
    integer, intent(in) :: i1, i2
    logical :: eq
    type(ftuple_list_t), pointer :: current
    select type (list)
    type is (ftuple_list_t)
      current => list%get_entry (i1)
      do
        if (associated (current%equiv)) then
          current => current%equiv
          if (current%index == i2) then
            eq = .true.
            exit
          end if
        else
          eq = .false.
          exit
        end if
      end do
    end select
  end function ftuple_list_check_equiv

\eatline
\nwendcode{}\nwbegindocs{47}Class for working with the flavor specification arrays.
\nwenddocs{}\nwbegincode{48}\moddef{fks regions: public}\plusendmoddef
  public :: flv_structure_t
\nwendcode{}\nwbegincode{49}\moddef{fks regions: types}\plusendmoddef
  type :: flv_structure_t
    integer, dimension(:), allocatable :: flst
    integer :: nlegs
    logical, dimension(:), allocatable :: massive
    logical, dimension(:), allocatable :: colored
  contains
  \LA{}fks regions: flv structure: TBP\RA{}
  end type flv_structure_t

\eatline
\nwendcode{}\nwbegindocs{50}\nwdocspar
Returns \texttt{true} if the two particles at position \texttt{i}
and \texttt{j} in the flavor array can originate from the same
splitting. For this purpose, the function first checks whether the splitting is
allowed at all. If this is the case, the emitter is removed from the
flavor array. If the resulting array is equivalent to the Born flavor
structure \texttt{flv\_born}, the pair is accepted as a valid
splitting. 
\nwenddocs{}\nwbegincode{51}\moddef{fks regions: flv structure: TBP}\endmoddef
  procedure :: valid_pair => flv_structure_valid_pair
\nwendcode{}\nwbegincode{52}\moddef{fks regions: procedures}\plusendmoddef
  function flv_structure_valid_pair &
       (flv_real,i,j, flv_born, model) result (valid)
    class(flv_structure_t), intent(in) :: flv_real
    integer, intent(in) :: i,j
    type(flv_structure_t), intent(in) :: flv_born
    type(model_data_t), intent(in) :: model
    logical :: valid
    integer :: k, n_orig
    type(flv_structure_t) :: flv_test
    integer, dimension(:), allocatable :: flv_orig, flv_orig2
    valid = .false.
\nwendcode{}\nwbegindocs{53}\nwdocspar
First check whether the splitting is possible. The array
{\tt{}flv{\char95}orig} contains all particles which share a vertex with the
particles at position {\tt{}i} and {\tt{}j}. If its size is equal to zero,
no splitting is possible and the subroutine is exited. 
\nwenddocs{}\nwbegincode{54}\moddef{fks regions: procedures}\plusendmoddef
    call model%match_vertex &
         (flv_real%flst(i), flv_real%flst(j), flv_orig)
    n_orig = size (flv_orig)
    if (n_orig == 0) then
      return
    else 
\nwendcode{}\nwbegindocs{55}\nwdocspar
For a quark emitting a gluon, {\tt{}flv{\char95}orig} contains the PDG code of
the anti-quark. To be on the safe side, a second array is created,
which contains both the positively and negatively signed PDG
codes. Then, the origial tuple $(i,j)$ is removed from the real flavor
structure and the particles in {\tt{}flv{\char95}orig2} are inserted. 
If the resulting Born configuration is equal to the underlying Born
configuration, up to a permutation of final-state particles, the tuple
$(i,j)$ is accepted as valid. 
\nwenddocs{}\nwbegincode{56}\moddef{fks regions: procedures}\plusendmoddef
      allocate (flv_orig2 (2*n_orig))
      flv_orig2 (1:n_orig) = flv_orig
      flv_orig2 (n_orig+1:2*n_orig) = -flv_orig
      do k = 1, 2*n_orig
        flv_test = flv_real%insert_particle (i,j,flv_orig2(k))
        valid = flv_born == flv_test
        if (valid) return
      end do
    end if
  end function flv_structure_valid_pair

\eatline
\nwendcode{}\nwbegindocs{57}This function checks whether two flavor arrays are the same up to a
permutation of the final-state particles 
\nwenddocs{}\nwbegincode{58}\moddef{fks regions: procedures}\plusendmoddef
  function flv_structure_equivalent (flv1, flv2) result(equiv)
    type(flv_structure_t), intent(in) :: flv1, flv2
    logical :: equiv
    integer :: i, j, n
    integer :: f1, f2
    logical, dimension(:), allocatable :: present, checked
    n = size (flv1%flst)
    equiv = .true.
    if (n /= size (flv2%flst)) then
      call msg_fatal &
           ('flv_structure_equivalent: flavor arrays do not have equal lengths')
    else
      allocate (present(n))
      allocate (checked(n))
        do i=1,n
           present(i) = .false.
           checked(i) = .false.
        end do
        do i=1,n
          do j=1,n
          if (flv1%flst(i) == flv2%flst(j) .and. .not. checked(j)) then 
              present(i) = .true.
              checked(j) = .true.
              exit
            end if
          end do
        end do
        do i=1,n
          if(.not.present(i)) equiv = .false.
        end do
    end if      
  end function flv_structure_equivalent

\eatline
\nwendcode{}\nwbegindocs{59}Returs a new flavor array with the particle at position
\texttt{index} removed. 
\nwenddocs{}\nwbegincode{60}\moddef{fks regions: flv structure: TBP}\plusendmoddef
  procedure :: remove_particle => flv_structure_remove_particle
\nwendcode{}\nwbegincode{61}\moddef{fks regions: procedures}\plusendmoddef
  function flv_structure_remove_particle (flv1, index) result(flv2)
    class(flv_structure_t), intent(in) :: flv1
    integer, intent(in) :: index   
    type(flv_structure_t) :: flv2
    integer :: n1, n2
    n1 = size (flv1%flst)
    n2 = n1-1
    if (allocated (flv2%flst)) then
      deallocate (flv2%flst)
    end if
    allocate (flv2%flst (n2))
    if (index == 1) then
      flv2%flst(1:n2) = flv1%flst(2:n1)
    else if (index == n1) then
      flv2%flst(1:n2) = flv1%flst(1:n2)
    else
      flv2%flst(1:index-1) = flv1%flst(1:index-1)
      flv2%flst(index:n2) = flv1%flst(index+1:n1)
    end if
  end function flv_structure_remove_particle

\eatline
\nwendcode{}\nwbegindocs{62}Removes the paritcles at position i1 and i2 and inserts a new
particle at position i1. 
\nwenddocs{}\nwbegincode{63}\moddef{fks regions: flv structure: TBP}\plusendmoddef
  procedure :: insert_particle => flv_structure_insert_particle
\nwendcode{}\nwbegincode{64}\moddef{fks regions: procedures}\plusendmoddef
  function flv_structure_insert_particle (flv1, i1, i2, particle) result (flv2)
    class(flv_structure_t), intent(in) :: flv1
    integer, intent(in) :: i1, i2, particle
    type(flv_structure_t) :: flv2
    type(flv_structure_t) :: flv_tmp
    integer :: n1, n2
    n1 = size (flv1%flst)
    n2 = n1-1
    allocate (flv2%flst(n2))
    if (i1 < i2) then
      flv_tmp = flv1%remove_particle (i1)
      flv_tmp = flv_tmp%remove_particle (i2-1)
    else if(i2 < i1) then
      flv_tmp = flv1%remove_particle(i2)
      flv_tmp = flv_tmp%remove_particle(i1-1)
    else
      call msg_fatal ("Trying to set ftuple with i1 = i2!")
    end if
    if (i1 == 1) then
      flv2%flst(1) = particle
      flv2%flst(2:n2) = flv_tmp%flst(1:n2-1)
    else if (i1 == n1 .or. i1 == n2) then
      flv2%flst(1:n2-1) = flv_tmp%flst(1:n2-1)
      flv2%flst(n2) = particle
    else
      flv2%flst(1:i1-1) = flv_tmp%flst(1:i1-1)
      flv2%flst(i1) = particle
      flv2%flst(i1+1:n2) = flv_tmp%flst(i1:n2-1)
    end if
  end function flv_structure_insert_particle 

\eatline
\nwendcode{}\nwbegindocs{65}Returns the number of particles in a flavor array
\nwenddocs{}\nwbegincode{66}\moddef{fks regions: flv structure: TBP}\plusendmoddef
  procedure :: get_nlegs => flv_structure_get_nlegs
\nwendcode{}\nwbegincode{67}\moddef{fks regions: procedures}\plusendmoddef
  function flv_structure_get_nlegs (flv) result(n)
    class(flv_structure_t), intent(in) :: flv
    integer :: n
    n = size (flv%flst) 
  end function flv_structure_get_nlegs

\eatline
\nwendcode{}\nwbegindocs{68}Counts the number of occurances of a particle in a 
flavor array
\nwenddocs{}\nwbegincode{69}\moddef{fks regions: flv structure: TBP}\plusendmoddef
  procedure :: count_particle => flv_structure_count_particle
\nwendcode{}\nwbegincode{70}\moddef{fks regions: procedures}\plusendmoddef
  function flv_structure_count_particle (flv, part) result (n)
    class(flv_structure_t), intent(in) :: flv
    integer, intent(in) :: part
    integer :: n
    n = count (flv%flst == part)
  end function flv_structure_count_particle

\eatline
\nwendcode{}\nwbegindocs{71}Initializer for flavor structures
\nwenddocs{}\nwbegincode{72}\moddef{fks regions: flv structure: TBP}\plusendmoddef
  procedure :: init => flv_structure_init
\nwendcode{}\nwbegincode{73}\moddef{fks regions: procedures}\plusendmoddef
  subroutine flv_structure_init (flv, aval)
    class(flv_structure_t), intent(inout) :: flv
    integer, intent(in), dimension(:) :: aval
    integer :: n
    n = size (aval)
    allocate (flv%flst (n))
    flv%flst(1:n) = aval(1:n)
    flv%nlegs = n
  end subroutine flv_structure_init

\eatline
\nwendcode{}\nwbegindocs{74}\nwdocspar
\nwenddocs{}\nwbegincode{75}\moddef{fks regions: flv structure: TBP}\plusendmoddef
  procedure :: write => flv_structure_write
\nwendcode{}\nwbegincode{76}\moddef{fks regions: procedures}\plusendmoddef
  subroutine flv_structure_write (flv, unit)
    class(flv_structure_t), intent(in) :: flv
    integer, intent(in), optional :: unit
    integer :: i, u
    u = given_output_unit (unit); if (u < 0) return
    write (u, '(A1)',advance = 'no') '['
    do i = 1, size(flv%flst)-1
      write (u, '(I3,A1)', advance = 'no') flv%flst(i), ','
    end do
    write (u, '(I3,A1)') flv%flst(i), ']'
  end subroutine flv_structure_write

\eatline
\nwendcode{}\nwbegindocs{77}Creates the underlying Born flavor structure for a given real flavor
structure if the particle at position \texttt{emitter} is removed 
\nwenddocs{}\nwbegincode{78}\moddef{fks regions: flv structure: TBP}\plusendmoddef
  procedure :: create_uborn => flv_structure_create_uborn
\nwendcode{}\nwbegincode{79}\moddef{fks regions: procedures}\plusendmoddef
  function flv_structure_create_uborn (flst_alr, emitter) result(flst_uborn)
    class(flv_structure_t), intent(in) :: flst_alr
    integer, intent(in) :: emitter
    type(flv_structure_t) :: flst_uborn
    integer n_alr, n_uborn
    n_alr = size(flst_alr%flst)
    n_uborn = n_alr-1
    allocate (flst_uborn%flst (n_uborn))
    if (emitter > 2) then
      if (flst_alr%flst(n_alr) == 21) then
         !!! Emitted particle is a gluon => just remove it
         flst_uborn = flst_alr%remove_particle(n_alr)
         !!! Emission type is a gluon splitting into two quars
      else if (is_quark (abs(flst_alr%flst(n_alr))) .and. &
               is_quark (abs(flst_alr%flst(n_alr-1))) .and. &
               flst_alr%flst(n_alr) + flst_alr%flst(n_alr-1) == 0) then
         flst_uborn = flst_alr%insert_particle(n_alr-1,n_alr,21)
      end if
     else
        if (flst_alr%flst(n_alr) == 21) then
           flst_uborn = flst_alr%remove_particle(n_alr)
        else if (is_quark (abs(flst_alr%flst(n_alr))) .and. &
                  is_gluon (abs(flst_alr%flst(emitter)))) then
           flst_uborn = &
                flst_alr%insert_particle (emitter,n_alr,-flst_alr%flst(n_alr))
        else if (is_quark (abs(flst_alr%flst(n_alr))) .and. &
                  is_quark (abs(flst_alr%flst(emitter))) .and. &
                  flst_alr%flst(n_alr) == flst_alr%flst(emitter)) then
           flst_uborn = flst_alr%insert_particle(emitter,n_alr,21)
        end if
     end if
  end function flv_structure_create_uborn

\eatline
\nwendcode{}\nwbegindocs{80}\nwdocspar
\nwenddocs{}\nwbegincode{81}\moddef{fks regions: flv structure: TBP}\plusendmoddef
  procedure :: evaluate => flv_structure_evaluate
\nwendcode{}\nwbegincode{82}\moddef{fks regions: procedures}\plusendmoddef
  subroutine flv_structure_evaluate (flv, n, model)
    class(flv_structure_t), intent(inout) :: flv
    integer, intent(in) :: n
    type(model_data_t), intent(in) :: model
    integer :: i
    type(flavor_t) :: flavor
    allocate (flv%massive (n), flv%colored(n))
    do i = 1, n
       call flavor%init (flv%flst(i), model)
       flv%massive(i) = flavor%get_mass () > 0
       flv%colored(i) = is_quark (abs(flv%flst(i))) .or. &
                        is_gluon (flv%flst(i))
    end do
  end subroutine flv_structure_evaluate 

\eatline
\nwendcode{}\nwbegindocs{83}\nwdocspar
\nwenddocs{}\nwbegincode{84}\moddef{fks regions: public}\plusendmoddef
  public :: singular_region_t
\nwendcode{}\nwbegincode{85}\moddef{fks regions: types}\plusendmoddef
  type :: singular_region_t
    integer :: alr
    type(flv_structure_t) :: flst_real
    type(flv_structure_t) :: flst_uborn
    integer :: mult
    integer :: emitter
    integer :: nregions
    integer :: real_index
    type(ftuple_t), dimension(:), allocatable :: flst_allreg
    integer :: uborn_index
    logical :: double_fsr
    logical :: soft_divergence
    logical :: coll_divergence
  contains
  \LA{}fks regions: singular region: TBP\RA{}
  end type singular_region_t

\eatline
\nwendcode{}\nwbegindocs{86}In case of a $g \rightarrow gg$ or $g \rightarrow qq$ splitting, the factor
\begin{equation*}
 \frac{2E_{\rm{em}}}{E_{\rm{em}} + E_{\rm{rad}}}
\end{equation*}
is multiplied to the real matrix element. This way, the symmetry of the splitting is used
and only one singular region has to be taken into account. However, the factor ensures that
there is only a soft singularity if the radiated parton becomes soft.
\nwenddocs{}\nwbegincode{87}\moddef{fks regions: singular region: TBP}\endmoddef
  procedure :: set_splitting_info => singular_region_set_splitting_info
\nwendcode{}\nwbegincode{88}\moddef{fks regions: procedures}\plusendmoddef
  subroutine singular_region_set_splitting_info (region) 
    class(singular_region_t), intent(inout) :: region
    integer :: i1, i2
    integer :: reg
    region%double_fsr = .false.
    associate (ftuple => region%flst_allreg)
       do reg = 1, region%nregions
          call ftuple(reg)%get (i1, i2)
          if (i1 /= region%emitter) then
             cycle
          else
             region%soft_divergence = &
                ftuple(reg)%splitting_type /= G_TO_QQ

             region%coll_divergence = &
                    .not. region%flst_real%massive(i1)

             if (ftuple(reg)%splitting_type > 1) then
                region%double_fsr = .true.
                exit
             else if (ftuple(reg)%splitting_type == UNDEFINED_SPLITTING) then
                call msg_fatal ("All splittings should be defined!") 
             end if
          end if
       end do
     end associate
  end subroutine singular_region_set_splitting_info

\eatline
\nwendcode{}\nwbegindocs{89}\nwdocspar
\nwenddocs{}\nwbegincode{90}\moddef{fks regions: singular region: TBP}\plusendmoddef
  procedure :: double_fsr_factor => singular_region_double_fsr_factor
\nwendcode{}\nwbegincode{91}\moddef{fks regions: procedures}\plusendmoddef
  function singular_region_double_fsr_factor (region, p) result (val)
    class(singular_region_t), intent(in) :: region
    type(vector4_t), intent(in), dimension(:) :: p
    real(default) :: val
    real(default) :: E_rad, E_em
    if (region%double_fsr) then
      E_em = energy (p(region%emitter))
      E_rad = energy (p(region%flst_real%get_nlegs()))
      val = 2*E_em/(E_em + E_rad)
    else
      val = 1._default
    end if
  end function singular_region_double_fsr_factor

\eatline
\nwendcode{}\nwbegindocs{92}\nwdocspar
\nwenddocs{}\nwbegincode{93}\moddef{fks regions: singular region: TBP}\plusendmoddef
  procedure :: has_soft_divergence => singular_region_has_soft_divergence
\nwendcode{}\nwbegincode{94}\moddef{fks regions: procedures}\plusendmoddef
  function singular_region_has_soft_divergence (region) result (div)
     class(singular_region_t), intent(in) :: region
     logical :: div
     div = region%soft_divergence
  end function singular_region_has_soft_divergence

\eatline
\nwendcode{}\nwbegindocs{95}\nwdocspar
\nwenddocs{}\nwbegincode{96}\moddef{fks regions: singular region: TBP}\plusendmoddef
  procedure :: has_collinear_divergence => &
            singular_region_has_collinear_divergence
\nwendcode{}\nwbegincode{97}\moddef{fks regions: procedures}\plusendmoddef
  function singular_region_has_collinear_divergence (region) result (div)
    class(singular_region_t), intent(in) :: region
    logical :: div
    div = region%coll_divergence
  end function singular_region_has_collinear_divergence

\eatline
\nwendcode{}\nwbegindocs{98}\nwdocspar
\nwenddocs{}\nwbegincode{99}\moddef{fks regions: types}\plusendmoddef
  type, abstract :: fks_mapping_t
     real(default) :: sumdij
     real(default) :: sumdij_soft
  contains
  \LA{}fks regions: fks mapping: TBP\RA{}
  end type fks_mapping_t

\eatline
\nwendcode{}\nwbegindocs{100}\nwdocspar
\nwenddocs{}\nwbegincode{101}\moddef{fks regions: public}\plusendmoddef
  public :: fks_mapping_default_t
\nwendcode{}\nwbegincode{102}\moddef{fks regions: types}\plusendmoddef
  type, extends (fks_mapping_t) :: fks_mapping_default_t
    real(default) :: exp_1, exp_2
  contains
  \LA{}fks regions: fks mapping 1: TBP\RA{}
  end type fks_mapping_default_t

\eatline
\nwendcode{}\nwbegindocs{103}\nwdocspar
\nwenddocs{}\nwbegincode{104}\moddef{fks regions: interfaces}\endmoddef
  interface operator(==)
    module procedure flv_structure_equivalent
  end interface

\eatline
\nwendcode{}\nwbegindocs{105}\nwdocspar
\nwenddocs{}\nwbegincode{106}\moddef{fks regions: public}\plusendmoddef
  public :: region_data_t
\nwendcode{}\nwbegincode{107}\moddef{fks regions: types}\plusendmoddef
  type :: region_data_t
    type(singular_region_t), dimension(:), allocatable :: regions
    type(flv_structure_t), dimension(:), allocatable :: flv_born
    type(flv_structure_t), dimension(:), allocatable :: flv_real
    integer, dimension(:), allocatable :: emitters
    integer :: n_regions
    integer :: n_emitters
    integer :: n_flv_born
    integer :: n_flv_real
    integer :: nlegs_born
    integer :: nlegs_real
    type(flavor_t) :: flv_extra
    class(fks_mapping_t), allocatable :: fks_mapping
  contains
  \LA{}fks regions: reg data: TBP\RA{}
  end type region_data_t

\eatline
\nwendcode{}\nwbegindocs{108}\nwdocspar
\nwenddocs{}\nwbegincode{109}\moddef{fks regions: reg data: TBP}\endmoddef
  procedure :: init => region_data_init
\nwendcode{}\nwbegincode{110}\moddef{fks regions: procedures}\plusendmoddef
  subroutine region_data_init (reg_data, model, flavor_born, &
                               flavor_real, mapping_type)
    class(region_data_t), intent(inout) :: reg_data
    type(model_data_t), intent(in) :: model
    integer, intent(inout), dimension(:,:), allocatable :: &
         flavor_born, flavor_real
    integer, intent(in) :: mapping_type
    integer, dimension(:), allocatable :: current_flavor
    type(ftuple_list_t), dimension(:), allocatable :: ftuples
    integer, dimension(:), allocatable :: emitter
    type(flv_structure_t), dimension(:), allocatable :: flst_alr
    integer :: i
    reg_data%n_flv_born = size(flavor_born(1,:))
    reg_data%n_flv_real = size(flavor_real(1,:))
    reg_data%nlegs_born = size(flavor_born(:,1))
    reg_data%nlegs_real = reg_data%nlegs_born + 1
    allocate (reg_data%flv_born (reg_data%n_flv_born))
    allocate (reg_data%flv_real (reg_data%n_flv_real))
    allocate (current_flavor (reg_data%n_flv_born))
    do i = 1, reg_data%n_flv_born
      current_flavor = flavor_born(:,i)
      call reg_data%flv_born(i)%init (current_flavor)
    end do
    deallocate (current_flavor)
    allocate (current_flavor (reg_data%n_flv_real))
    do i = 1, reg_data%n_flv_real
      current_flavor = flavor_real(:,i)
      call reg_data%flv_real(i)%init (current_flavor)
    end do   

    select case (mapping_type)
    case (1)
       allocate (fks_mapping_default_t :: reg_data%fks_mapping)
    case default
       call msg_fatal ("Init region_data: FKS mapping not implemented!")
    end select

    call reg_data%flv_extra%init &
         (reg_data%flv_real(1)%flst(reg_data%nlegs_real), &
         model)
    call reg_data%find_regions (model, ftuples, emitter, flst_alr)
    call reg_data%init_regions (ftuples, emitter, flst_alr)
    call reg_data%evaluate_flavors (model)
    call reg_data%set_splitting_info ()
    call reg_data%find_emitters ()
    call reg_data%write ()
  end subroutine region_data_init

\eatline
\nwendcode{}\nwbegindocs{111}\nwdocspar
\nwenddocs{}\nwbegincode{112}\moddef{fks regions: reg data: TBP}\plusendmoddef
  procedure :: evaluate_flavors => region_data_evaluate_flavors
\nwendcode{}\nwbegincode{113}\moddef{fks regions: procedures}\plusendmoddef
  subroutine region_data_evaluate_flavors (reg_data, model)
    class(region_data_t), intent(inout) :: reg_data
    type(model_data_t), intent(in) :: model
    integer :: i
    do i = 1, reg_data%n_regions
       associate (region => reg_data%regions(i))
          call region%flst_uborn%evaluate (reg_data%nlegs_born, model)
          call region%flst_real%evaluate (reg_data%nlegs_real, model)
       end associate
    end do
    do i = 1, reg_data%n_flv_born
       call reg_data%flv_born(i)%evaluate (reg_data%nlegs_born, model)
    end do
    do i = 1, reg_data%n_flv_real
       call reg_data%flv_real(i)%evaluate (reg_data%nlegs_real, model)
    end do
  end subroutine region_data_evaluate_flavors

\eatline
\nwendcode{}\nwbegindocs{114}Creates a list containing the emitter of each singular region.
\nwenddocs{}\nwbegincode{115}\moddef{fks regions: reg data: TBP}\plusendmoddef
  procedure :: get_emitters => region_data_get_emitters
\nwendcode{}\nwbegincode{116}\moddef{fks regions: procedures}\plusendmoddef
  function region_data_get_emitters (reg_data) result(emitters)
    class(region_data_t), intent(inout) :: reg_data
    integer, dimension(:), allocatable :: emitters
    integer :: i
    allocate (emitters (reg_data%n_regions))
    do i = 1, reg_data%n_regions
       emitters(i) = reg_data%regions(i)%emitter
    end do
  end function region_data_get_emitters

\eatline
\nwendcode{}\nwbegindocs{117}Returns $S_i = \frac{1}{\mathcal{D}d_i}$ or $S_{ij} =
\frac{1}{\mathcal{D}d_{ij}}$ for one particular singular region.  At
this point, the flavor array should be rearranged in such a way that
the emitted particle is at the last position of 
the flavor structure list.
\nwenddocs{}\nwbegincode{118}\moddef{fks regions: reg data: TBP}\plusendmoddef
  procedure :: get_svalue => region_data_get_svalue
\nwendcode{}\nwbegincode{119}\moddef{fks regions: procedures}\plusendmoddef
  function region_data_get_svalue (reg_data, p, alr, emitter) result (sval)
    class(region_data_t), intent(inout) :: reg_data
    type(vector4_t), intent(in), dimension(:) :: p
    integer, intent(in) :: alr, emitter
    real(default) :: sval
    associate (map => reg_data%fks_mapping) 
      map%sumdij = map%compute_sumdij (reg_data%regions(alr), p)
      sval = map%svalue (p, emitter, reg_data%nlegs_real)
    end associate
  end function region_data_get_svalue

\eatline
\nwendcode{}\nwbegindocs{120}The same as above, but for the soft limit.
\nwenddocs{}\nwbegincode{121}\moddef{fks regions: reg data: TBP}\plusendmoddef
  procedure :: get_svalue_soft => region_data_get_svalue_soft
\nwendcode{}\nwbegincode{122}\moddef{fks regions: procedures}\plusendmoddef
  function region_data_get_svalue_soft &
       (reg_data, p, p_soft, alr, emitter) result (sval)
    class(region_data_t), intent(inout) :: reg_data
    type(vector4_t), intent(in), dimension(:) :: p
    type(vector4_t), intent(in) :: p_soft
    integer, intent(in) :: alr, emitter
    real(default) :: sval
    associate (map => reg_data%fks_mapping)
      map%sumdij_soft = &
      map%compute_sumdij_soft (reg_data%regions(alr), p, p_soft)
      sval = map%svalue_soft (p, p_soft, emitter)
    end associate
  end function region_data_get_svalue_soft

\eatline
\nwendcode{}\nwbegindocs{123}This subroutine starts with a specification of $N$- and
$N+1$-particle configurations, {\tt{}flst{\char95}born} and {\tt{}flst{\char95}real}, saved
in {\tt{}reg{\char95}data}. From these, it creates a list of fundamental tuples,
a list of emitters and a list containing the $N+1$-particle
configuration, rearranged  in such a way that the emitter-radiation
pair is last ({\tt{}flst{\char95}alr}). For the $e^+ \, e^- \, \rightarrow u \,
\bar{u} \, g$- example, the generated objects are shown in table
\ref{table:ftuples and flavors}. Note that at this point, {\tt{}flst{\char95}alr}
is arranged in such a way that the emitter can only be equal to
$n_{legs}-1$ for final-state radiation or 0, 1, or 2 for initial-state
radiation. Further, it occurs that regions can be equivalent. For
example in table \ref{table:ftuples and flavors} the regions
correpsonding to \texttt{alr} = 1 and \texttt{alr} = 3 as well as
\texttt{alr} = 2 and \texttt{alr} = 4 describe the same physics and
are therefore equivalent. 
\nwenddocs{}\nwbegindocs{124}\nwdocspar
\nwenddocs{}\nwbegincode{125}\moddef{fks regions: reg data: TBP}\plusendmoddef
  procedure :: find_regions => region_data_find_regions
\nwendcode{}\nwbegincode{126}\moddef{fks regions: procedures}\plusendmoddef
  subroutine region_data_find_regions &
       (reg_data, model, ftuples, emitter, flst_alr)
    class(region_data_t), intent(in) :: reg_data
    type(model_data_t), intent(in) :: model
    type(ftuple_list_t), intent(out), dimension(:), allocatable :: ftuples
    integer, intent(out), dimension(:), allocatable :: emitter
    type(flv_structure_t), intent(out), dimension(:), allocatable :: flst_alr
    type(ftuple_t) :: current_ftuple
    integer, dimension(:), allocatable :: emitter_tmp
    type(flv_structure_t), dimension(:), allocatable :: flst_alr_tmp
    integer :: nreg, nborn, nreal
    integer :: nlegreal
    integer :: n_gluon_real, n_gluon_born
    integer, parameter :: maxnregions = 100
    integer :: i, j, k, l, n

    associate (flv_born => reg_data%flv_born)
      associate (flv_real => reg_data%flv_real)
        nborn = size (flv_born)
        nreal = size (flv_real)
        nlegreal = size (flv_real(1)%flst)
        allocate (ftuples (nreal))
        allocate (emitter_tmp (maxnregions))
        allocate (flst_alr_tmp (maxnregions))
        n = 0

        ITERATE_REAL_FLAVOR: do l = 1, nreal
           call ftuples(l)%init     
\LA{}fks: check final state emissions\RA{}  
\LA{}fks: check initial state emissions\RA{}
        end do ITERATE_REAL_FLAVOR

        nreg = n

      end associate
    end associate

    allocate (flst_alr (nreg))
    allocate (emitter (nreg))
    flst_alr(1:nreg) = flst_alr_tmp(1:nreg)
    emitter(1:nreg) = emitter_tmp(1:nreg)
  end subroutine region_data_find_regions 

\eatline
\nwendcode{}\nwbegindocs{127}First check final state emissions
\nwenddocs{}\nwbegincode{128}\moddef{fks: check final state emissions}\endmoddef
           do i = 3, nlegreal
             do j = i+1, nlegreal
               do k = 1, nborn
                 if (flv_real(l)%valid_pair(i,j, flv_born(k), model) &
                     .or. flv_real(l)%valid_pair(j,i,flv_born(k), model)) then
                   n = n+1
                   n_gluon_real = flv_real(l)%count_particle (GLUON)
                   n_gluon_born = flv_born(k)%count_particle (GLUON)
                   if (n_gluon_born - n_gluon_real < 0) then
                     if(flv_real(l)%valid_pair(i,j, flv_born(k), model)) then
                        flst_alr_tmp(n) = create_alr (flv_real(l),i,j)
                     else
                        flst_alr_tmp(n) = create_alr (flv_real(l),j,i)
                     end if
                   else
                     flst_alr_tmp(n) = flv_real(l)
                   end if
                   call current_ftuple%set (i,j)
                   call current_ftuple%determine_splitting_type (flv_real(l), i, j)
                   call ftuples(l)%append (current_ftuple)
                   emitter_tmp(n) = nlegreal - 1
                   exit
                 end if
               end do
             end do
\nwendcode{}\nwbegindocs{129}Check initial-state emissions. It suffices to just check the
final-state of the first and the initial-state of the second array. 
\nwenddocs{}\nwbegincode{130}\moddef{fks: check initial state emissions}\endmoddef
             do k = 1, nborn
               if (flv_real(l)%valid_pair(1,i, flv_born(k), model) &
                   .and. flv_real(l)%valid_pair(2,i, flv_born(k), model)) then
                 n = n + 1
                 call current_ftuple%set (0,i)
                 call ftuples(l)%append (current_ftuple)
                 emitter_tmp(n) = 0
                 flst_alr_tmp(n) = create_alr (flv_real(l),0,i)
                 exit
               else if (flv_real(l)%valid_pair(1,i, flv_born(k), model) &
                        .and. .not. &
                        flv_real(l)%valid_pair(2,i, flv_born(k), model)) then
                 n = n+1
                 call current_ftuple%set (1,i)
                 call ftuples(l)%append (current_ftuple)
                 emitter_tmp(n) = 1
                 flst_alr_tmp(n) = create_alr (flv_real(l),1,i)
                 exit
               else if (flv_real(l)%valid_pair(2,i, flv_born(k), model) &
                       .and. .not. &
                       flv_real(l)%valid_pair(1,i, flv_born(k), model)) then
                 n = n+1
                 call current_ftuple%set(2,i)
                 call ftuples(l)%append (current_ftuple)
                 emitter_tmp(n) = 2
                 flst_alr_tmp(n) = create_alr (flv_real(l),2,i)
                 exit
               end if
             end do
           end do
\nwendcode{}\nwbegindocs{131}Creates singular regions according to table \ref{table:singular
regions}. It scans all regions in table \ref{table:ftuples and
flavors} and records the real flavor structures. If they are
equivalent, the flavor structure is not recorded, but the multiplicity
of the present one is increased.
\nwenddocs{}\nwbegincode{132}\moddef{fks regions: reg data: TBP}\plusendmoddef
  procedure :: init_regions => region_data_init_singular_regions
\nwendcode{}\nwbegincode{133}\moddef{fks regions: procedures}\plusendmoddef
  subroutine region_data_init_singular_regions &
       (reg_data, ftuples, emitter, flst_alr)
    class(region_data_t), intent(inout) :: reg_data
    type(ftuple_list_t), intent(inout), dimension(:), allocatable :: ftuples
    type(ftuple_list_t) :: current_region
    integer, intent(in), dimension(:), allocatable :: emitter
    type(flv_structure_t), intent(in), dimension(:), allocatable :: flst_alr
    type(flv_structure_t), dimension(:), allocatable :: flst_uborn, flst_alr2
    integer, dimension(:), allocatable :: mult
    integer, dimension(:), allocatable :: flst_emitter
    integer :: nregions, maxregions
    integer, dimension(:,:), allocatable :: perm_list
    integer, dimension(:), allocatable :: index
    integer :: i, j, k, l
    integer :: nlegs 
    logical :: equiv
    integer :: nreg, i1, i2
    integer :: i_first, j_first
    integer, dimension(:), allocatable :: &
         region_to_ftuple, ftuple_limits, k_index
    type(flv_structure_t) :: flst_save

    maxregions = size (emitter)
    nlegs = size(flst_alr(1)%flst)

    allocate (flst_uborn (maxregions))
    allocate (flst_alr2 (maxregions))
    allocate (mult (maxregions))
    allocate (flst_emitter (maxregions))
    allocate (index (maxregions))
    allocate (region_to_ftuple (maxregions))
    allocate (ftuple_limits (size (ftuples)))
    allocate (k_index (maxregions))

    mult = 0

    do i = 1, size(ftuples)
      ftuple_limits(i) = ftuples(i)%get_n_tuples ()
    end do
    if (.not. (sum (ftuple_limits) == maxregions)) &
         call msg_fatal ("Too many regions!")
    k = 1
    do j =1, size(ftuples)
      do i = 1, ftuple_limits(j)
        region_to_ftuple(k) = i
        k = k + 1
      end do
    end do
    i_first = 1
    j_first = 1
    j = 1
    SCAN_REGIONS: do l = 1, size(ftuples)
    SCAN_FTUPLES: do i = i_first, i_first + ftuple_limits (l) -1 
      equiv = .false.
      if (i==i_first) then
        flst_alr2(j)%flst = flst_alr(i)%flst
        mult(j) = mult(j) + 1
        flst_uborn(j) = flst_alr(i)%create_uborn (emitter(i))
        flst_emitter(j) = emitter(i)
        index (j) = region_to_index(ftuples, i)
        k_index (j) = region_to_ftuple(i)
        j = j+1
      else
        !!! Check for equivalent flavor structures
        do k =j_first ,j-1
           if (emitter(i) == emitter(k) .and. emitter(i) > 2) then
             if (flst_alr(i) == flst_alr2(k) .and. &
                 flst_alr(i)%flst(nlegs-1) == flst_alr2(k)%flst(nlegs-1) &
                 .and. flst_alr(i)%flst(nlegs) == flst_alr2(k)%flst(nlegs)) then
                   mult(k) = mult(k) + 1
                   equiv = .true.
                   call ftuples (region_to_index(ftuples, i))%set_equiv &
                        (k_index(k), region_to_ftuple(i))
                   exit
              end if
           else if (emitter(i) == emitter(k) .and. emitter(i) <= 2) then
             if (flst_alr(i) == flst_alr2(k)) then
               mult(k) = mult(k) + 1
               equiv = .true.
               call ftuples (region_to_index(ftuples,i))%set_equiv &
                    (k_index(k), region_to_ftuple(i))
               exit
             end if
          end if
        end do
        if (.not.equiv) then
          flst_alr2(j)%flst = flst_alr(i)%flst
          mult(j) = mult(j) + 1
          flst_uborn(j) = flst_alr(i)%create_uborn (emitter(i))
          flst_emitter(j) = emitter(i)
          index (j) = region_to_index (ftuples, i)
          k_index (j) = region_to_ftuple(i)
          j = j+1
        end if
      end if
    end do SCAN_FTUPLES
    i_first = i_first + ftuple_limits(l)
    j_first = j_first + j - 1
    end do SCAN_REGIONS
    nregions = j-1
    allocate (reg_data%regions (nregions))
    do j = 1, nregions
      do i = 1, reg_data%n_flv_born
        if (reg_data%flv_born (i) == flst_uborn (j)) then
          if (allocated (perm_list)) deallocate (perm_list)
          call fks_permute_born &
               (reg_data%flv_born (i), flst_uborn (j), perm_list)
          call fks_apply_perm (flst_alr2(j), flst_emitter(j), perm_list)
        end if
      end do
    end do
    !!! Check if new emitters require a rearrangement of ftuples
    do i = 1, nregions
      reg_data%regions(i)%alr = i
      reg_data%regions(i)%flst_real = flst_alr2(i)
      reg_data%regions(i)%mult = mult(i)
      reg_data%regions(i)%flst_uborn = flst_uborn(i)
      reg_data%regions(i)%emitter = flst_emitter(i)
      nreg = ftuples (index(i))%get_n_tuples ()
      reg_data%regions(i)%nregions = nreg
      allocate (reg_data%regions(i)%flst_allreg (nreg))
      do j = 1, nreg
        current_region = ftuples (index(i))%get_entry (j)
        if (.not. associated (current_region%equiv)) then
          call current_region%ftuple%get (i1, i2)
          if (i2 /= nlegs) &
             call current_region%ftuple%set (i1, nlegs)
        end if
        reg_data%regions(i)%flst_allreg (j) = current_region%ftuple
      end do
    end do
    !!! Find underlying Born index
    do j = 1, nregions
      do i = 1, reg_data%n_flv_born
        if (reg_data%flv_born (i) == reg_data%regions(j)%flst_uborn) then
          reg_data%regions(j)%uborn_index = i
          exit
        end if
      end do
    end do

    k = 1
    associate (regions => reg_data%regions)
       do i = 1, nregions
          if (i==1) then
             regions(i)%real_index = 1
             flst_save = flst_alr2(1)
             cycle
          end if
          if (flst_alr2(i) == flst_save) then
             regions(i)%real_index = k
          else
             k = k+1
             regions(i)%real_index = k
             flst_save = flst_alr2(i)
          end if
       end do
    end associate
    reg_data%n_regions = size (reg_data%regions)

  end subroutine region_data_init_singular_regions

\eatline
\nwendcode{}\nwbegindocs{134}Create an array containing all emitters of a singular region.
\nwenddocs{}\nwbegincode{135}\moddef{fks regions: reg data: TBP}\plusendmoddef
  procedure :: find_emitters => region_data_find_emitters
\nwendcode{}\nwbegincode{136}\moddef{fks regions: procedures}\plusendmoddef
  subroutine region_data_find_emitters (reg_data)
    class(region_data_t), intent(inout) :: reg_data
    integer :: i, j, n
    integer :: em
    integer, dimension(10) :: em_count
    em_count = 0
    n = 0

\eatline
\nwendcode{}\nwbegindocs{137}Count the number of different emitters.
\nwenddocs{}\nwbegincode{138}\moddef{fks regions: procedures}\plusendmoddef
    do i = 1, reg_data%n_regions
      em = reg_data%regions(i)%emitter
      if (.not. any (em_count == em)) then
        n = n+1
        em_count(i) = em
      end if
    end do

    if (n < 1) call msg_fatal ("region_data_find_emitters: No emitters found")
    reg_data%n_emitters = n
    allocate (reg_data%emitters (reg_data%n_emitters))
    reg_data%emitters = 0

    j = 1
    do i = 1, size(reg_data%regions)
      em = reg_data%regions(i)%emitter
      if (.not. any (reg_data%emitters == em)) then
        reg_data%emitters(j) = em
        j = j+1
      end if
    end do
  end subroutine region_data_find_emitters

\eatline
\nwendcode{}\nwbegindocs{139}\nwdocspar
\nwenddocs{}\nwbegincode{140}\moddef{fks regions: reg data: TBP}\plusendmoddef
  procedure :: set_splitting_info => region_data_set_splitting_info
\nwendcode{}\nwbegincode{141}\moddef{fks regions: procedures}\plusendmoddef
  subroutine region_data_set_splitting_info (reg_data)
    class(region_data_t), intent(inout) :: reg_data
    integer :: reg
    do reg = 1, reg_data%n_regions
       call reg_data%regions(reg)%set_splitting_info ()
    end do
  end subroutine region_data_set_splitting_info

\eatline
\nwendcode{}\nwbegindocs{142}\nwdocspar
\nwenddocs{}\nwbegincode{143}\moddef{fks regions: reg data: TBP}\plusendmoddef
  procedure :: write_regions => region_data_write_regions
\nwendcode{}\nwbegincode{144}\moddef{fks regions: procedures}\plusendmoddef
  subroutine region_data_write_regions (reg_data)
    class(region_data_t), intent(inout) :: reg_data
    integer :: n, i, j
    n = size(reg_data%regions)
    associate (regions => reg_data%regions)
      do i = 1, n
        print *, i, '//', regions(i)%flst_real%flst, '//', &
             regions(i)%mult ,'//', regions(i)%flst_uborn%flst , &
                    '//', regions(i)%emitter
        do j = 1, size (regions(i)%flst_allreg)
          call regions(i)%flst_allreg(j)%write
        end do
      end do
    end associate
  end subroutine region_data_write_regions

\eatline
\nwendcode{}\nwbegindocs{145}Creates a table with information about all singular regions and
writes it to a file. 
\nwenddocs{}\nwbegindocs{146}Returns the index of the real flavor structure an ftuple belongs to.
\nwenddocs{}\nwbegincode{147}\moddef{fks regions: reg data: TBP}\plusendmoddef
  procedure :: write => region_data_write
\nwendcode{}\nwbegincode{148}\moddef{fks regions: procedures}\plusendmoddef
  subroutine region_data_write (reg_data, proc)
    class(region_data_t), intent(inout) :: reg_data
    type(string_t), intent(inout), optional :: proc
    integer :: u, i, j
    integer :: nreal, nborn
    integer :: i1, i2, nreg
    integer :: maxnregions, nreg_diff
    integer :: nleft, nright
    type(singular_region_t) :: region
    type(string_t) :: flst_title, ftuple_title
    character(len=7) :: flst_format = "(I3,A1)"
    character(len=16) :: ireg_format = "(A1,I3,A1,I3,A2)"
    character(len=7) :: ireg_space_format = "(7X,A1)"
    u = free_unit ()
    open (u, file="region_data.log", action = "write", status="replace")
    maxnregions = 1
    do j = 1, reg_data%n_regions
      if (size (reg_data%regions(j)%flst_allreg) > maxnregions) &
           maxnregions = reg_data%regions(j)%nregions
    end do
    flst_title = '(A' // flst_title_format(reg_data%nlegs_real) // ')'
    ftuple_title = '(A' // ftuple_title_format() // ')'
    write (u,*) 'Total number of regions: ', size(reg_data%regions)
    write (u, '(A3)', advance = 'no') 'alr'
    call write_separator ()
    write (u, char (flst_title), advance = 'no') 'flst_real'
    call write_separator ()
    write (u, '(A3)', advance = 'no') 'em'
    call write_separator ()
    write (u, '(A3)', advance = 'no') 'mult'
    call write_separator ()
    write (u, '(A4)', advance = 'no') 'nreg'
    call write_separator ()
    write (u, char (ftuple_title), advance = 'no') 'ftuples'
    call write_separator ()
    flst_title = '(A' // flst_title_format(reg_data%nlegs_born) // ')'
    write (u, char (flst_title), advance = 'no') 'flst_born'
    call write_separator ()
    write (u, '(A7)') 'i_uborn'
    do j = 1, reg_data%n_regions
      region = reg_data%regions(j)
      nreal = size (region%flst_real%flst)
      nborn = size (region%flst_uborn%flst)
      write (u, '(I3)', advance = 'no') j
      call write_separator ()
      write (u, '(A1)', advance = 'no') '['
      do i = 1, nreal-1
        write (u, flst_format, advance = 'no') region%flst_real%flst(i), ','
      end do
      write (u, flst_format, advance = 'no') region%flst_real%flst(nreal), ']'
      call write_separator ()
      write (u, '(I3)', advance = 'no') region%emitter
      call write_separator ()
      write (u, '(I3)', advance = 'no') region%mult
      call write_separator ()
      write (u, '(I4)', advance = 'no') region%nregions
      call write_separator ()
      !!! write ftuples
      nreg = region%nregions
      if (nreg == maxnregions) then
        nleft = 0
        nright = 0
      else
        nreg_diff = maxnregions - nreg
        nleft = nreg_diff/2
        if (mod(nreg_diff,2) == 0) then
          nright = nleft
        else
          nright = nleft + 1
        end if
      end if
      if (nleft > 0) then
        do i=1,nleft
          write(u,ireg_space_format, advance='no') ' '
        end do
      end if
      write (u,'(A1)', advance = 'no') '\{'
      if (nreg > 1) then
        do i=1,nreg-1
          call region%flst_allreg(i)%get (i1, i2)
          write(u,ireg_format,advance = 'no') '(', i1, ',', i2, '),'
        end do
      end if
      call region%flst_allreg(nreg)%get (i1, i2) 
      write (u,ireg_format,advance = 'no') '(', i1, ',', i2, ')\}' 
      if (nright > 0) then
        do i=1,nright
          write(u,ireg_space_format, advance='no') ' '
        end do
      end if
      !!! end write ftuples
      call write_separator ()
      write (u,'(A1)',advance = 'no') '['
      do i=1,nborn-1
        write(u,flst_format,advance = 'no') region%flst_uborn%flst(i), ','
      end do
      write (u,flst_format, advance = 'no') region%flst_uborn%flst(nborn), ']'
      call write_separator ()
      write (u, '(I7)', advance = 'no') region%uborn_index
      write(u,*) ''
    end do
    close (u)
  contains
    function flst_title_format (n) result (frmt)
      integer, intent(in) :: n
      type(string_t) :: frmt
      character(len=2) :: frmt_char
      write (frmt_char, '(I2)') 4*n+1
      frmt = var_str (frmt_char)
    end function flst_title_format

    function ftuple_title_format () result (frmt)
      type(string_t) :: frmt
      character(len=2) :: frmt_char
      write (frmt_char, '(I2)') 10*maxnregions+1
      frmt = var_str (frmt_char)
    end function ftuple_title_format

    subroutine write_separator ()
      character(len=10) :: sep_format = "(1X,A2,1X)"
      write (u, sep_format, advance = 'no') '||'
    end subroutine write_separator 

  end subroutine region_data_write

\eatline
\nwendcode{}\nwbegindocs{149}Returns the index of the real flavor structure an ftuple belogs to.
\nwenddocs{}\nwbegincode{150}\moddef{fks regions: procedures}\plusendmoddef
  function region_to_index (list, i) result(index)
    type(ftuple_list_t), intent(inout), dimension(:), allocatable :: list
    integer, intent(in) :: i
    integer :: index
    integer :: nlist
    integer :: j
    integer, dimension(:), allocatable :: nreg
    nlist = size(list)
    allocate (nreg (nlist))
    do j = 1, nlist
      if (j == 1) then
        nreg(j) = list(j)%get_n_tuples ()
      else
        nreg(j) = nreg(j-1) + list(j)%get_n_tuples ()
      end if
    end do
    do j = 1, nlist
      if (j == 1) then
        if (i <= nreg(j)) then
          index = j
          exit
        end if
      else
        if (i > nreg(j-1) .and. i <= nreg(j)) then
          index = j
          exit
        end if
      end if
    end do
  end function region_to_index

\eatline
\nwendcode{}\nwbegindocs{151}Rearrange the flavor array in such a way that the emitted particle
is last and the emitter is second last. {\tt{}i1} is the index of the
emitter, {\tt{}i2} is the index of the emitted particle. Only works for
final-state emitters. 
\nwenddocs{}\nwbegincode{152}\moddef{fks regions: procedures}\plusendmoddef
  function create_alr (flv1,i1,i2) result(flv2)
    type(flv_structure_t), intent(in) :: flv1
    integer, intent(in) :: i1, i2
    type(flv_structure_t) :: flv2
    integer :: n, i, j
    n = size (flv1%flst)
    allocate (flv2%flst (n))
    if (i1 > 2) then
      flv2%flst(1:2) = flv1%flst(1:2)
      flv2%flst(n-1) = flv1%flst(i1)
      flv2%flst(n) = flv1%flst(i2)
\nwendcode{}\nwbegindocs{153}Order remaining particles according to their original position
\nwenddocs{}\nwbegincode{154}\moddef{fks regions: procedures}\plusendmoddef
      j = 3
      do i = 3,n
        if (i /= i1 .and. i /= i2) then
          flv2%flst(j) = flv1%flst(i)
          j = j+1
        end if
      end do
    else
      call msg_fatal ("Create alr: Only works for final-state emissions!")
    end if
  end function create_alr

\eatline
\nwendcode{}\nwbegindocs{155}Explain
\nwenddocs{}\nwbegincode{156}\moddef{fks regions: procedures}\plusendmoddef
  subroutine fks_permute_born (flv_in, flv_out, perm_list)
    type(flv_structure_t), intent(in) :: flv_in
    type(flv_structure_t), intent(inout) :: flv_out
    integer, intent(out), dimension(:,:), allocatable :: perm_list
    integer, dimension(:,:), allocatable :: perm_list_tmp
    integer :: n_perms, n_perms_max
    integer :: nlegs
    integer :: flv1, flv2, tmp
    integer :: i, j, j_min
    n_perms_max = 100 
    !!! actually (n-1)!, but there seems to be no intrinsic function 
    !!! of this type in fortran
    if (allocated (perm_list_tmp)) deallocate (perm_list_tmp)
    allocate (perm_list_tmp (n_perms_max,2))
    n_perms = 0
    j_min = 3
    nlegs = size (flv_in%flst)
      do i = 3, nlegs
        flv1 = flv_in%flst(i)
        do j = j_min, nlegs
          flv2 = flv_out%flst(j)
          if (flv1 == flv2 .and. i /= j) then
            n_perms = n_perms + 1
            tmp = flv_out%flst(i)
            flv_out%flst(i) = flv2
            flv_out%flst(j) = tmp
            perm_list_tmp (n_perms, 1) = j
            perm_list_tmp (n_perms, 2) = i
            j_min = j_min + 1
            exit
          end if
        end do
      end do
      allocate (perm_list (n_perms, 2))
      perm_list (1:n_perms, :) = perm_list_tmp (1:n_perms, :)
  end subroutine fks_permute_born

\eatline
\nwendcode{}\nwbegindocs{157}Explain
\nwenddocs{}\nwbegincode{158}\moddef{fks regions: procedures}\plusendmoddef
  subroutine fks_apply_perm (flv, emitter, perm_list)
    type(flv_structure_t), intent(inout) :: flv
    integer, intent(inout) :: emitter
    integer, intent(in), dimension(:,:), allocatable :: perm_list
    integer :: i
    integer :: i1, i2
    integer :: tmp
    do i = 1, size (perm_list (:,1))
      i1 = perm_list (i,1)
      i2 = perm_list (i,2)
      tmp = flv%flst (i1)
      flv%flst (i1) = flv%flst (i2)
      flv%flst (i2) = tmp
      if (i1 == emitter) emitter = i2
    end do
  end subroutine fks_apply_perm

\eatline
\nwendcode{}\nwbegindocs{159}Translates the tree code of the emitter branch into the position of
the emitter in the flavor structure array. 
\nwenddocs{}\nwbegincode{160}\moddef{fks regions: public}\plusendmoddef
  public :: fks_tree_to_position
\nwendcode{}\nwbegincode{161}\moddef{fks regions: procedures}\plusendmoddef
  function fks_tree_to_position (k, n_tot) result(pos)
    integer, intent(in) :: k, n_tot
    integer :: pos
    integer :: k_tot
    k_tot = 2**(n_tot - 1)
    !!! Inital-state particles
    if (k == k_tot) then
       pos = 1
    else if (k == k_tot/2) then
       pos = 2
       !!! Final-state particles
    else
       ! pos = 3 + nint(log(k)/log(2))
       pos = 3 + dual_log (k)
    end if
  contains
    recursive function dual_log (x) result (ld) 
      integer, intent(in) :: x
      integer :: ld
      if (x == 1) then
        ld = 0
      else
        ld = 1 + dual_log (x/2)
      end if
    end function dual_log 
  end function fks_tree_to_position
\eatline
\nwendcode{}\nwbegindocs{162}\nwdocspar
\nwenddocs{}\nwbegincode{163}\moddef{fks regions: fks mapping: TBP}\endmoddef
  procedure (fks_mapping_dij), deferred :: dij
\nwendcode{}\nwbegincode{164}\moddef{fks regions: interfaces}\plusendmoddef
  abstract interface
    function fks_mapping_dij (map, p, i, j) result (d)
      import
      class(fks_mapping_t), intent(in) :: map
      type(vector4_t), intent(in), dimension(:) :: p
      integer, intent(in) :: i, j
      real(default) :: d
    end function fks_mapping_dij
  end interface

\eatline
\nwendcode{}\nwbegindocs{165}\nwdocspar
\nwenddocs{}\nwbegincode{166}\moddef{fks regions: fks mapping: TBP}\plusendmoddef
  procedure (fks_mapping_compute_sumdij), deferred :: compute_sumdij
\nwendcode{}\nwbegincode{167}\moddef{fks regions: interfaces}\plusendmoddef
  abstract interface
    function fks_mapping_compute_sumdij (map, sregion, p) result (d)
      import
      class(fks_mapping_t), intent(in) :: map
      type(singular_region_t), intent(inout) :: sregion
      type(vector4_t), intent(in), dimension(:) :: p
      real(default) :: d
    end function fks_mapping_compute_sumdij
  end interface

\eatline
\nwendcode{}\nwbegindocs{168}\nwdocspar
\nwenddocs{}\nwbegincode{169}\moddef{fks regions: fks mapping: TBP}\plusendmoddef
  procedure (fks_mapping_svalue), deferred :: svalue
\nwendcode{}\nwbegincode{170}\moddef{fks regions: interfaces}\plusendmoddef
  abstract interface
    function fks_mapping_svalue (map, p, i, j) result (value)
      import
      class(fks_mapping_t), intent(in) :: map
      type(vector4_t), intent(in), dimension(:) :: p
      integer, intent(in) :: i, j
      real(default) :: value
    end function fks_mapping_svalue
  end interface

\eatline
\nwendcode{}\nwbegincode{171}\moddef{fks regions: fks mapping: TBP}\plusendmoddef
  procedure (fks_mapping_dij_soft), deferred :: dij_soft
\nwendcode{}\nwbegincode{172}\moddef{fks regions: interfaces}\plusendmoddef
  abstract interface
    function fks_mapping_dij_soft (map, p_born, p_soft, em) result (d)
      import
      class(fks_mapping_t), intent(in) :: map
      type(vector4_t), intent(in), dimension(:) :: p_born
      type(vector4_t), intent(in) :: p_soft
      integer, intent(in) :: em
      real(default) :: d
    end function fks_mapping_dij_soft
  end interface

\eatline
\nwendcode{}\nwbegindocs{173}\nwdocspar
\nwenddocs{}\nwbegincode{174}\moddef{fks regions: fks mapping: TBP}\plusendmoddef
  procedure (fks_mapping_compute_sumdij_soft), deferred :: compute_sumdij_soft
\nwendcode{}\nwbegincode{175}\moddef{fks regions: interfaces}\plusendmoddef
  abstract interface
    function fks_mapping_compute_sumdij_soft (map, sregion, p_born, p_soft) result (d)
      import
      class(fks_mapping_t), intent(in) :: map
      type(singular_region_t), intent(inout) :: sregion
      type(vector4_t), intent(in), dimension(:) :: p_born
      type(vector4_t), intent(in) :: p_soft
      real(default) :: d
    end function
  end interface
\eatline
\nwendcode{}\nwbegindocs{176}\nwdocspar
\nwenddocs{}\nwbegincode{177}\moddef{fks regions: fks mapping: TBP}\plusendmoddef
  procedure (fks_mapping_svalue_soft), deferred :: svalue_soft
\nwendcode{}\nwbegincode{178}\moddef{fks regions: interfaces}\plusendmoddef
  abstract interface
    function fks_mapping_svalue_soft (map, p_born, p_soft, em) result (value)
      import
      class(fks_mapping_t), intent(in) :: map
      type(vector4_t), intent(in), dimension(:) :: p_born
      type(vector4_t), intent(in) :: p_soft
      integer, intent(in) :: em
      real(default) :: value
    end function fks_mapping_svalue_soft
  end interface

\eatline
\nwendcode{}\nwbegindocs{179}\nwdocspar
\nwenddocs{}\nwbegincode{180}\moddef{fks regions: fks mapping 1: TBP}\endmoddef
  procedure :: set_parameter => fks_mapping_default_set_parameter
\nwendcode{}\nwbegincode{181}\moddef{fks regions: procedures}\plusendmoddef
  subroutine fks_mapping_default_set_parameter (map, dij_exp1, dij_exp2)
    class(fks_mapping_default_t), intent(inout) :: map
    real(default), intent(in) :: dij_exp1, dij_exp2
    map%exp_1 = dij_exp1
    map%exp_2 = dij_exp2
  end subroutine fks_mapping_default_set_parameter

\eatline
\nwendcode{}\nwbegindocs{182}\nwdocspar
\nwenddocs{}\nwbegincode{183}\moddef{fks regions: fks mapping 1: TBP}\plusendmoddef
  procedure :: dij => fks_mapping_default_dij
\nwendcode{}\nwbegincode{184}\moddef{fks regions: procedures}\plusendmoddef
  function fks_mapping_default_dij (map, p, i, j) result (d)
    class(fks_mapping_default_t), intent(in) :: map
    type(vector4_t), intent(in), dimension(:) :: p
    integer, intent(in) :: i, j
    real(default) :: d
    real(default) :: sqrts
    real :: y
    real :: E1, E2

    if (i /= j .and. (i > 2 .or. j > 2)) then
      if (i == 0 .or. j == 0) then
        if (j == 0) then
          E1 = energy (p(i))
          y = polar_angle_ct (p(i))
        else
          E1 = energy (p(j))
          y = polar_angle_ct(p(j))
        end if
        d = (E1**2 * (1-y**2))**map%exp_2
      else
        E1 = energy(p(i))
        E2 = energy(p(j))
        y = enclosed_angle_ct (p(i), p(j))
        sqrts = (p(1)+p(2))**1
        d = (2*p(i)*p(j) * E1*E2 / (E1 + E2)**2)**map%exp_1
      end if  
    else if (i == j) then
      call msg_fatal ("Invalid FKS region: Emitter equals FKS parton!")
    else
      !!! case i,j <= 2 not yet implemented
      d = 0
    end if
  end function fks_mapping_default_dij

\nwendcode{}\nwbegindocs{185} %def fks_mapping_default_dij
\nwenddocs{}\nwbegindocs{186}\nwdocspar
\nwenddocs{}\nwbegincode{187}\moddef{fks regions: fks mapping 1: TBP}\plusendmoddef
  procedure :: compute_sumdij => fks_mapping_default_compute_sumdij
\nwendcode{}\nwbegincode{188}\moddef{fks regions: procedures}\plusendmoddef
  function fks_mapping_default_compute_sumdij (map, sregion, p) result (d)
    class(fks_mapping_default_t), intent(in) :: map
    type(singular_region_t), intent(inout) :: sregion
    type(vector4_t), intent(in), dimension(:) :: p
    real(default) :: d
    integer :: i, k, l

    associate (ftuples => sregion%flst_allreg)
      d = 0
      do i = 1, sregion%nregions
        call ftuples(i)%get (k, l)
        d = d + 1.0/map%dij (p, k, l)
      end do
    end associate

  end function fks_mapping_default_compute_sumdij

\eatline
\nwendcode{}\nwbegindocs{189}\nwdocspar
\nwenddocs{}\nwbegincode{190}\moddef{fks regions: fks mapping 1: TBP}\plusendmoddef
  procedure :: svalue => fks_mapping_default_svalue
\nwendcode{}\nwbegincode{191}\moddef{fks regions: procedures}\plusendmoddef
  function fks_mapping_default_svalue (map, p, i, j) result (value)
    class(fks_mapping_default_t), intent(in) :: map
    type(vector4_t), intent(in), dimension(:) :: p
    integer, intent(in) :: i, j
    real(default) :: value
    value = 1._default / (map%dij (p, i, j) * map%sumdij)
  end function fks_mapping_default_svalue

\eatline
\nwendcode{}\nwbegindocs{192}\nwdocspar
\nwenddocs{}\nwbegincode{193}\moddef{fks regions: fks mapping 1: TBP}\plusendmoddef
  procedure :: dij_soft => fks_mapping_default_dij_soft
\nwendcode{}\nwbegincode{194}\moddef{fks regions: procedures}\plusendmoddef
  function fks_mapping_default_dij_soft (map, p_born, p_soft, em) result (d)
    class(fks_mapping_default_t), intent(in) :: map
    type(vector4_t), intent(in), dimension(:) :: p_born
    type(vector4_t), intent(in) :: p_soft
    integer, intent(in) :: em
    real(default) :: d
    d = (2*p_born(em)*p_soft / energy(p_born(em)))**map%exp_1
  end function fks_mapping_default_dij_soft

\eatline
\nwendcode{}\nwbegindocs{195}\nwdocspar
\nwenddocs{}\nwbegincode{196}\moddef{fks regions: fks mapping 1: TBP}\plusendmoddef
  procedure :: compute_sumdij_soft => fks_mapping_default_compute_sumdij_soft
\nwendcode{}\nwbegincode{197}\moddef{fks regions: procedures}\plusendmoddef
  function fks_mapping_default_compute_sumdij_soft (map, sregion, p_born, p_soft) result (d)
    class(fks_mapping_default_t), intent(in) :: map
    type(singular_region_t), intent(inout) :: sregion
    type(vector4_t), intent(in), dimension(:) :: p_born
    type(vector4_t), intent(in) :: p_soft
    real(default) :: d
    integer :: i, k, l
    integer :: nlegs
    d = 0
    nlegs = size (sregion%flst_real%flst)
    associate (ftuples => sregion%flst_allreg)
      do i = 1, sregion%nregions
        call ftuples(i)%get (k,l)
        if (l == nlegs) then
          d = d + 1._default/map%dij_soft (p_born, p_soft, k)
        end if
      end do
    end associate
  end function fks_mapping_default_compute_sumdij_soft

\eatline
\nwendcode{}\nwbegindocs{198}\nwdocspar
\nwenddocs{}\nwbegincode{199}\moddef{fks regions: fks mapping 1: TBP}\plusendmoddef
  procedure :: svalue_soft => fks_mapping_default_svalue_soft
\nwendcode{}\nwbegincode{200}\moddef{fks regions: procedures}\plusendmoddef
  function fks_mapping_default_svalue_soft (map, p_born, p_soft, em) result (value)
    class(fks_mapping_default_t), intent(in) :: map
    type(vector4_t), intent(in), dimension(:) :: p_born
    type(vector4_t), intent(in) :: p_soft
    integer, intent(in) :: em
    real(default) :: value
    value = 1._default/(map%sumdij_soft*map%dij_soft (p_born, p_soft, em))
  end function fks_mapping_default_svalue_soft

\eatline
\nwendcode{}\nwbegindocs{201}\nwdocspar
\clearpage
%%%%%%%%%%%%%%%%%%%%%%%%%%%%%%%%%%%%%%%%%%%%%%%%%%%%%%%%%%%%%%%%%%%%%%%%
\section{Virtual contribution to the cross section}
\nwenddocs{}\nwbegincode{202}\moddef{\code{}virtual.f90\edoc{}}\endmoddef
\LA{}File header\RA{}

module virtual

\LA{}Use kinds\RA{}
\LA{}Use strings\RA{}
  use constants, only: pi, twopi
  use diagnostics
  use pdg_arrays
  use model_data
  use physics_defs
  use sm_physics
  use lorentz
  use flavors
  use fks_regions

\LA{}Standard module head\RA{}

\LA{}virtual: public\RA{}

\LA{}virtual: types\RA{}

contains
 
\LA{}virtual: procedures\RA{}

end module virtual
\eatline
\nwendcode{}\nwbegincode{203}\moddef{virtual: public}\endmoddef
  public :: virtual_t
\nwendcode{}\nwbegincode{204}\moddef{virtual: types}\endmoddef
  type :: virtual_t
    real(default) :: Q
    real(default), dimension(:,:), allocatable :: I
    real(default) :: vfin
    real(default) :: sqme_cc
    real(default) :: sqme_virt
    real(default), dimension(:,:), allocatable :: gamma_0, gamma_p, c_flv
    real(default) :: ren_scale2
    integer, dimension(:), allocatable :: n_is_neutrinos
    integer :: nlegs, nflv
    logical :: bad_point
    logical :: use_internal_color_correlations
  contains
  \LA{}virtual: virtual: TBP\RA{}
  end type virtual_t
  
\eatline
\nwendcode{}\nwbegincode{205}\moddef{virtual: virtual: TBP}\endmoddef
  procedure :: init => virtual_init
\nwendcode{}\nwbegincode{206}\moddef{virtual: procedures}\endmoddef
 subroutine virtual_init (object, flv_born)
    class(virtual_t), intent(inout) :: object
    integer, intent(in), dimension(:,:) :: flv_born
    integer :: nlegs, nflv
    integer :: i_flv
    type(flavor_t) :: flv
    object%nlegs = size (flv_born, 1); object%nflv = size (flv_born, 2)
    allocate (object%I (object%nlegs, object%nlegs))
    allocate (object%gamma_0 (object%nlegs, object%nflv), &
              object%gamma_p (object%nlegs, object%nflv), &
              object%c_flv (object%nlegs, object%nflv))
    call object%init_constants (flv_born)
    allocate (object%n_is_neutrinos (object%nflv))
    object%n_is_neutrinos = 0
    do i_flv = 1, object%nflv
       if (is_neutrino (flv_born(1, i_flv))) &
          object%n_is_neutrinos(i_flv) = object%n_is_neutrinos(i_flv) + 1
       if (is_neutrino (flv_born(2, i_flv))) &
          object%n_is_neutrinos(i_flv) = object%n_is_neutrinos(i_flv) + 1
    end do 
  contains
    function is_neutrino (flv) result (neutrino)
      integer, intent(in) :: flv
      logical :: neutrino
      neutrino = (abs(flv)==12 .or. abs(flv)==14 .or. abs(flv)==16)
    end function is_neutrino
  end subroutine virtual_init

\eatline
\nwendcode{}\nwbegindocs{207}Write down constant definition somewhere
\nwenddocs{}\nwbegincode{208}\moddef{virtual: virtual: TBP}\plusendmoddef
  procedure :: init_constants => virtual_init_constants
\nwendcode{}\nwbegincode{209}\moddef{virtual: procedures}\plusendmoddef
  subroutine virtual_init_constants (object, flv_born)
    class(virtual_t), intent(inout) :: object
    integer, intent(in), dimension(:,:) :: flv_born
    integer :: i_part, i_flv
    integer, parameter :: nf = 1 
    do i_flv = 1, size (flv_born, 2)
       do i_part = 1, size (flv_born, 1)
          if (is_gluon (flv_born(i_part, i_flv))) then
             object%gamma_0(i_part, i_flv) = (11*ca - 2*nf)/6
             object%gamma_p(i_part, i_flv) = (67.0/9 - 2*pi**2/3)*ca - 23.0/18*nf
             object%c_flv(i_part, i_flv) = ca
          else if (is_quark (abs(flv_born(i_part, i_flv)))) then
             object%gamma_0(i_part, i_flv) = 1.5*cf
             object%gamma_p(i_part, i_flv) = (6.5 - 2*pi**2/3)*cf
             object%c_flv(i_part, i_flv) = cf
          else
             object%gamma_0(i_part, i_flv) = 0
             object%gamma_p(i_part, i_flv) = 0
             object%c_flv(i_part, i_flv) = 0
          end if
       end do
    end do
  end subroutine virtual_init_constants
  
\eatline
\nwendcode{}\nwbegindocs{210}Set the renormalization scale. If the input is zero, use the
center-of-mass energy. 
\nwenddocs{}\nwbegincode{211}\moddef{virtual: virtual: TBP}\plusendmoddef
  procedure :: set_ren_scale => virtual_set_ren_scale
\nwendcode{}\nwbegincode{212}\moddef{virtual: procedures}\plusendmoddef
  subroutine virtual_set_ren_scale (object, p, ren_scale)
    class(virtual_t), intent(inout) :: object
    type(vector4_t), dimension(:), intent(in) :: p
    real(default), intent(in) :: ren_scale
    if (ren_scale > 0) then
      object%ren_scale2 = ren_scale**2
    else
      object%ren_scale2 = (p(1)+p(2))**2
    end if
  end subroutine virtual_set_ren_scale
  
\eatline
\nwendcode{}\nwbegindocs{213}The virtual-subtracted matrix element is given by the equation
\begin{equation}
  \label{virt_sub}
  \mathcal{V} = \frac{\alpha_s}{2\pi}\left(\mathcal{Q}\mathcal{B} +
  \sum \mathcal{I}_{ij}\mathcal{B}_{ij} + \mathcal{V}_{fin}\right),
\end{equation} 
where the quantity $\mathcal{Q}$ is given by
\begin{equation}
  \begin{split}
    \mathcal{Q} = \sum_{i=1}^n&\left[\gamma'_{f_i} - \log\frac{s}{Q^2}
                  \left(\gamma_{f_i} - 2C_{f_i}\log\frac{2E_i}{\sqrt{s}}\right)\right.\\
                  &\left. +2C_{f_i}\log^2\frac{2E_i}{\sqrt{s}} - 2\gamma_{f_i}\frac{2E_i}{\sqrt{s}}\right]\\
                  &- \log\frac{\mu_F^2}{Q^2}\left(\gamma_{f+} + \gamma_{f-}\right).
  \end{split}
  \label{virt_Q}
\end{equation} 
The expressions for $\mathcal{I}_{ij}$ can be found in equations (\ref{I_00}), (\ref{I_mm}), (\ref{I_0m}), 
depending on whether the particles involved in the radiation process are massive or massless. 
\nwenddocs{}\nwbegincode{214}\moddef{virtual: virtual: TBP}\plusendmoddef
  procedure :: evaluate => virtual_evaluate
\nwendcode{}\nwbegincode{215}\moddef{virtual: procedures}\plusendmoddef
  subroutine virtual_evaluate &
       (object, reg_data, i_flv, alpha_s, p_born, born, b_ij)
    class(virtual_t), intent(inout) :: object
    type(region_data_t), intent(in) :: reg_data
    integer, intent(in) :: i_flv
    real(default), intent(in) :: alpha_s
    type(vector4_t), intent(inout), dimension(:), allocatable :: p_born
    real(default), intent(in) :: born
    real(default), intent(in), dimension(:,:,:), allocatable :: b_ij
    integer :: i, j, alr
    integer :: nlegs
    real(default) :: BI
    if (object%bad_point) then
       object%sqme_virt = 0
    else
       BI = 0
       alr = find_first_matching_uborn (reg_data, i_flv)
       associate (flst_born => reg_data%regions(alr)%flst_uborn)
         call object%compute_Q (p_born, i_flv, flst_born%massive)
         do i = 1, object%nlegs
           do j = 1, object%nlegs
             if (i /= j) then
                if (flst_born%colored(i) .and. flst_born%colored(j)) then
                  call object%compute_I (p_born, flst_born%massive, i, j)
                  BI = BI + b_ij (i,j,reg_data%regions(alr)%uborn_index) * &
                                  object%I(i,j)
               end if
             end if
           end do
         end do
       end associate
       if (object%use_internal_color_correlations) BI = BI*born
       !!! A factor of alpha_s/twopi is assumed to be included in vfin
       object%sqme_virt = alpha_s/twopi * (object%Q*born + BI) + object%vfin
       if (object%n_is_neutrinos(i_flv) > 0) &
          object%sqme_virt = object%sqme_virt * object%n_is_neutrinos(i_flv)*2
    end if
  contains
    function find_first_matching_uborn (reg_data, i_proc) result (alr_out)
       type(region_data_t), intent(in) :: reg_data
       integer, intent(in) :: i_proc
       integer :: alr_out
       integer :: k
       alr_out = 0
       do k = 1, reg_data%n_regions
          alr_out = alr_out+1
          if (reg_data%regions(k)%uborn_index == i_proc) exit
       end do
    end function find_first_matching_uborn
  end subroutine virtual_evaluate

\eatline
\nwendcode{}\nwbegindocs{216}\nwdocspar
\nwenddocs{}\nwbegincode{217}\moddef{virtual: virtual: TBP}\plusendmoddef
  procedure :: compute_vfin_test => virtual_compute_vfin_test
\nwendcode{}\nwbegincode{218}\moddef{virtual: procedures}\plusendmoddef
  subroutine virtual_compute_vfin_test (object, p_born, sqme_born)
    class(virtual_t), intent(inout) :: object
    type(vector4_t), intent(in), dimension(:) :: p_born
    real(default), intent(in) :: sqme_born
    real(default) :: s
    s = (p_born(1)+p_born(2))**2
    !!! ----NOTE: Test implementation for e+ e- -> u ubar
    object%vfin = sqme_born * cf * &
         (pi**2 - 8 + 3*log(s/object%ren_scale2) - log(s/object%ren_scale2)**2)
    object%bad_point = .false.
  end subroutine virtual_compute_vfin_test

\eatline
\nwendcode{}\nwbegindocs{219}\nwdocspar
\nwenddocs{}\nwbegincode{220}\moddef{virtual: virtual: TBP}\plusendmoddef
  procedure :: set_vfin => virtual_set_vfin
\nwendcode{}\nwbegincode{221}\moddef{virtual: procedures}\plusendmoddef
  subroutine virtual_set_vfin (object, vfin)
    class(virtual_t), intent(inout) :: object
    real(default) :: vfin
    object%vfin = vfin
  end subroutine virtual_set_vfin

\eatline
\nwendcode{}\nwbegindocs{222} 
\nwenddocs{}\nwbegincode{223}\moddef{virtual: virtual: TBP}\plusendmoddef
  procedure :: set_bad_point => virtual_set_bad_point
\nwendcode{}\nwbegincode{224}\moddef{virtual: procedures}\plusendmoddef
  subroutine virtual_set_bad_point (object, value)
     class(virtual_t), intent(inout) :: object
     logical, intent(in) :: value
     object%bad_point = value
  end subroutine virtual_set_bad_point

\nwendcode{}\nwbegindocs{225} %def virtual_set_bad_point
\nwenddocs{}\nwbegindocs{226}\nwdocspar
\nwenddocs{}\nwbegincode{227}\moddef{virtual: virtual: TBP}\plusendmoddef
  procedure :: compute_Q => virtual_compute_Q
\nwendcode{}\nwbegincode{228}\moddef{virtual: procedures}\plusendmoddef
  subroutine virtual_compute_Q (object, p_born, i_flv, massive)
    class(virtual_t), intent(inout) :: object
    type(vector4_t), intent(in), dimension(:), allocatable :: p_born
    integer, intent(in) :: i_flv
    logical, dimension(:), intent(in) :: massive
    real(default) :: sqrts, E
    real(default) :: s1, s2, s3, s4
    integer :: i
    sqrts = sqrt ((p_born(1)+p_born(2))**2)
    !!! ---- NOTE: Implementation only works for lepton collisions. 
    !!!            This implies that both the summand containing log(s/q**2) 
    !!!            and (γ_fp + γ…vfm) vanish. 
    !!!            Also, s = (p_born(1)+p_born(2))**2
    object%Q = 0
    do i = 1, object%nlegs
       if (.not. massive (i)) then
          s1 = object%gamma_p(i, i_flv)
          E = vector4_get_component (p_born(i), 0)
          s2 = log(sqrts**2/object%ren_scale2)*&
               (object%gamma_0(i, i_flv)-2 * object%c_flv(i, i_flv) * log(2*E/sqrts))
          s3 = 2*log(2*E/sqrts)**2*object%c_flv(i, i_flv)
          s4 = 2*log(2*E/sqrts)*object%gamma_0(i, i_flv)
          object%Q = object%Q + s1 - s2 + s3 - s4
       else
          s1 = log(sqrts**2/object%ren_scale2)
          s2 = 0.5*I_m_eps (p_born(i))
          object%Q = object%Q - object%c_flv(i, i_flv) * (s1-s2)
       end if
    end do
  end subroutine virtual_compute_Q
  
\eatline
\nwendcode{}\nwbegindocs{229}The following code implements the $\mathcal{I}_{ij}$-function appearing in eq. blub. 
They are defined as follows:
\begin{itemize}
\item[Massles-Massles Case]
  \begin{equation}
    \begin{split}
    \mathcal{I}_{ij} &= \frac{1}{2}\log^2\frac{s}{Q^2} + \log\frac{s}{Q^2}\log\frac{k_ik_j}{2E_iE_j}
                                                       - \rm{Li}_2\left(\frac{k_ik_j}{2E_iE_j}\right) \\
                     &+ \frac{1}{2}\log^2\frac{k_ik_j}{2E_iE_j} - \log\left(1-\frac{k_ik_j}{2E_iE_j}\right) 
                        \log\frac{k_ik_j}{2E_iE_j}.
    \end{split}
    \label{I_00}
  \end{equation}
\item[Massive-Massive Case]
  \begin{equation}
     \mathcal{I}_{ij} = -\frac{1}{2}I_0(k_i, k_j)\log\frac{Q^2}{s} - \frac{1}{2}I_\epsilon(k_i,k_j)
     \label{I_mm}
  \end{equation}
  with
  \begin{equation*}
    I_0(k_i, k_j) = \frac{1}{\beta}\log\frac{1+\beta}{1-\beta}, \qquad 
                    \beta = \sqrt{1-\frac{k_i^2k_j^2}{(k_i \cdot k_j)^2}}
  \end{equation*}
  and a rather involved expression for $I_\epsilon$:
  \begin{align*}
    I_\epsilon(k_i, k_j) &= \left(K(z_j)-K(z_i)\right) \frac{1-\vec{\beta_i}\cdot\vec{\beta_j}}{\sqrt{a(1-b)}}, \\
    \vec{\beta_i} &= \frac{\vec{k}_i}{k_i^0}, \\
    a &= \beta_i^2 + \beta_j^2 - 2\vec{\beta}_i \cdot \vec{\beta}_j, \\
    x_i &= \frac{\beta_i^2 -\vec{\beta}_i \cdot \vec{\beta}_j}{a}, \\
    x_j &= \frac{\beta_j^2 -\vec{\beta}_i \cdot \vec{\beta}_j}{a} = 1-x_j, \\
    b &= \frac{\beta_i^2\beta_j^2 - (\vec{\beta}_i\cdot\vec{\beta}_j)^2}{a}, \\
    c &= \sqrt{\frac{b}{4a}}, \\
    z_+ &= \frac{1+\sqrt{1-b}}{\sqrt{b}}, \\
    z_- &= \frac{1-\sqrt{1-b}}{\sqrt{b}}, \\
    z_i &= \frac{\sqrt{x_i^2 + 4c^2} - x_i}{2c}, \\
    z_j &= \frac{\sqrt{x_j^2 + 4c^2} + x_j}{2c}, \\
    K(z) = &-\frac{1}{2}\log^2\frac{(z-z_-)(z_+-z)}{(z_++z)(z_-+z)} - 2Li_2\left(\frac{2z_-(z_+-z)}{(z_+-z_-)(z_-+z)}\right) \\
           &-2Li_2\left(-\frac{2z_+(z_-+z)}{(z_+-z_-)(z_+-z)}\right)
  \end{align*}
  
\item[Massive-Massless Case]
  \begin{equation}
     \mathcal{I}_{ij} = \frac{1}{2}\left[\log^2\frac{Q}{s} - \frac{\pi^2}{6}\right] -\frac{1}{2}I_0(k_i,k_j)\log\frac{Q^2}{s}
                       -\frac{1}{2}I_\epsilon(k_i,k_j)
     \label{I_0m}
  \end{equation}
  with
  \begin{align*}
     I_0(p,k) &= \log\frac{(\hat{p}\cdot\hat{k})^2}{\hat{k}^2}, \\
     I_\varepsilon(p,k) &= -2\left[\frac{1}{4}\log^2\frac{1-\beta}{1+\beta} + \log\frac{\hat{p}\cdot\hat{k}}{1+\beta}\log\frac{\hat{p}\cdot\hat{k}}{1-\beta} 
                           +Li_2\left(1-\frac{\hat{p}\cdot\hat{k}}{1+\beta}\right) + Li_2\left(1-\frac{\hat{p}\cdot\hat{k}}{1-\beta}\right)\right]
  \end{align*}
  using
  \begin{equation*}
     \hat{p} = \frac{p}{p^0}, \quad \hat{k} = \frac{k}{k^0}, \quad \beta = \frac{|\vec{k}|}{k_0}.
  \end{equation*} 
                           
\end{itemize}

\nwenddocs{}\nwbegincode{230}\moddef{virtual: virtual: TBP}\plusendmoddef
  procedure :: compute_I => virtual_compute_I
\nwendcode{}\nwbegincode{231}\moddef{virtual: procedures}\plusendmoddef
  subroutine virtual_compute_I (object, p_born, massive, i, j)
    class(virtual_t), intent(inout) :: object
    type(vector4_t), intent(in), dimension(:) :: p_born
    logical, dimension(:), intent(in) :: massive
    integer, intent(in) :: i, j
    real(default) :: somu2
    somu2 = (p_born(1)+p_born(2))**2/object%ren_scale2
    if (massive(i) .and. massive(j)) then
       object%I(i,j) = compute_Imm (p_born(i), p_born(j), somu2)
    else if (.not.massive(i) .and. massive(j)) then
       object%I(i,j) = compute_I0m (p_born(i), p_born(j), somu2)
    else if (massive(i) .and. .not.massive(j)) then
       object%I(i,j) = compute_I0m (p_born(j), p_born(i), somu2)
    else
       object%I(i,j) = compute_I00 (p_born(i), p_born(j), somu2)
    end if
  end subroutine virtual_compute_I

  function compute_I00 (pi, pj, somu2) result (I)
    type(vector4_t), intent(in) :: pi, pj
    real(default), intent(in) :: somu2
    real(default) :: I
    real(default) :: Ei, Ej
    real(default) :: pij, Eij
    real(default) :: s
    real(default) :: s1, s2, s3, s4, s5
    real(default) :: arglog
    real(default), parameter :: tiny_value = epsilon(1.0)
    !!! ----NOTE: As above, only lepton collisions. Therefore, the 
    !!!           first and second summand are not taken into account.

    s1 = 0; s2 = 0; s3 = 0; s4 = 0; s5 = 0
    Ei = vector4_get_component (pi, 0)
    Ej = vector4_get_component (pj, 0)
    pij = pi*pj; Eij = Ei*Ej
    s1 = 0.5*log(somu2)**2
    s2 = log(somu2)*log(pij/(2*Eij))
    s3 = Li2 (pij / (2*Eij))
    s4 = 0.5*log (pij / (2*Eij))**2
    arglog = 1 - pij/(2*Eij)
    if (arglog > tiny_value) then
      s5 = log(arglog) * log(pij / (2*Eij))
    else
      s5 = 0
    end if
    I = s1 + s2 - s3 + s4 - s5
  end function compute_I00

  function compute_I0m (ki, kj, somu2) result (I)
    type(vector4_t), intent(in) :: ki, kj
    real(default), intent(in) :: somu2
    real(default) :: I
    real(default) :: logsomu
    real(default) :: s1, s2, s3
    s1 = 0; s2 = 0; s3 = 0
    logsomu = log(somu2)
    s1 = 0.5*(0.5*logsomu - pi**2/6)
    s2 = 0.5*I_0m_0 (ki, kj)*logsomu
    s3 = 0.5*I_0m_eps (ki, kj)
    I = s1 + s2 - s3
  end function compute_I0m

  function compute_Imm (pi, pj, somu2) result (I)
    type(vector4_t), intent(in) :: pi, pj
    real(default), intent(in) :: somu2
    real(default) :: I
    real(default) :: s1, s2
    s1 = 0.5*log(somu2)*I_mm_0(pi, pj)
    s2 = 0.5*I_mm_eps(pi, pj)
    I = s1 - s2
  end function compute_Imm
  
\eatline
\nwendcode{}\nwbegindocs{232}\nwdocspar
\nwenddocs{}\nwbegincode{233}\moddef{virtual: procedures}\plusendmoddef
  function I_m_eps (p) result (I)
    type(vector4_t), intent(in) :: p
    real(default) :: I
    real(default) :: beta
    beta = space_part_norm (p)/p%p(0)
    I = 2*log((1+beta)/(1-beta))/beta
  end function I_m_eps

\eatline
\nwendcode{}\nwbegindocs{234}@For $p^2=0$ and $k^2 \neq 0$, this computes the expression
\begin{equation*}
  I_\epsilon(p,k) = -2\left[\frac{1}{4}\log^2\frac{1-\beta}{1+\beta}
                    +\log \frac{\hat{p} \cdot k}{1+\beta}\log{\hat{p} \cdot k}{1-\beta} 
                    +Li_2\left(1-\frac{\hat{p} \cdot k}{1+\beta}\right)
                    +Li_2\left(1-\frac{\hat{p} \cdot k}{1-\beta}\right)\right],
\end{equation*}
with $\hat{p} = \frac{p}{p^0}$ and $\beta = \frac{|\vec{k}|}{k^0}$.
\nwenddocs{}\nwbegincode{235}\moddef{virtual: procedures}\plusendmoddef
  function I_0m_eps (p, k) result (I)
    type(vector4_t), intent(in) :: p, k
    real(default) :: I
    type(vector4_t) :: pp, kp
    real(default) :: beta

    pp = p/energy(p); kp = k/energy(k) 
  
    beta = sqrt (1-pp*kp)
    I = -2*(log((1-beta)/(1+beta))**2/4 + log((pp*kp)/(1+beta))*log((pp*kp)/(1-beta)) &
        + Li2(1-(pp*kp)/(1+beta)) + Li2(1-(pp*kp)/(1-beta)))
  end function I_0m_eps
    
\eatline
\nwendcode{}\nwbegindocs{236}@For $p^2=0$ and $k^2 \neq 0$, computes the expression
\begin{equation*}
  I_0(p,k) = \log\frac{(\hat{p}\cdot\hat{k})^2}{\hat{k}^2}
\end{equation*}
\nwenddocs{}\nwbegincode{237}\moddef{virtual: procedures}\plusendmoddef
  function I_0m_0 (p, k) result (I)
    type(vector4_t), intent(in) :: p, k
    real(default) :: I
    type(vector4_t) :: pp, kp
    
    pp = p/energy(p); kp = k/energy(k)
    I = log((pp*kp)**2/kp**2)
  end function I_0m_0

\eatline
\nwendcode{}\nwbegindocs{238}For $k_1^2 \neq 0$ and $k_2^2 \neq 0$, computes the expression
\begin{equation*}
  I_\epsilon(k_1, k_2) = \left[K(z_2) - K(z_1)\right]
                         \frac{1-\vec{\beta}_1\cdot\vec{\beta}_2}{\sqrt{a(1-b)}},
\end{equation*}
where $\vec{\beta}_i = \frac{\vec{k}_i}{k_i^0}$. Further
\begin{align*}
  a &= \beta_1^2 + \beta_2^2 - 2\vec{\beta}_1\cdot\vec{\beta}_2, \\
  b &= \frac{\beta_1^2\beta_2^2 - \left(\vec{\beta}_1\cdot\vec{\beta}_2\right)^2}{a}, \\
  c &= \sqrt{\frac{b}{4a}}, \\
  K(z) = &-\frac{1}{2}\log^2\frac{(z-z_-)(z_+-z)}{(z_++z)(z_-+z)} 
          - 2Li_2\left(\frac{2z_-(z_+-z)}{(z_+-z_-)(z_-+z)}\right) \\
         &-2Li_2\left(-\frac{2z_+(z_-+z)}{(z_+-z_-)(z_+-z)}\right),\\
  z_+ &= \frac{1+\sqrt{1-b}}{\sqrt{b}}, \\
  z_- &= \frac{1-\sqrt{1-b}}{\sqrt{b}}, \\
  x_1 &= \frac{\beta_1^2-\vec{\beta}_1\cdot \vec{\beta}_2}{a}, \\
  x_2 &= 1-x_1
  z_1 &= \frac{\sqrt{x_1^2+4c^2}-x_1}{2c}, \\
  z_2 &= \frac{\sqrt{x_2^2+4c^2}+x_2}{2c}.
\end{align*}
\nwenddocs{}\nwbegincode{239}\moddef{virtual: procedures}\plusendmoddef
  function I_mm_eps (p1, p2) result (I)
    type(vector4_t), intent(in) :: p1, p2
    real(default) :: I
    type(vector3_t) :: beta1, beta2
    real(default) :: a, b
    real(default) :: zp, zm, z1, z2, x1, x2
    real(default) :: zmb, z1b
    real(default) :: K1, K2
   
    beta1 = space_part (p1)/energy(p1)
    beta2 = space_part (p2)/energy(p2)
    a = beta1**2 + beta2**2 - 2*beta1*beta2
    b = beta1**2*beta2**2 - (beta1*beta2)**2
    x1 = beta1**2 - beta1*beta2
    x2 = beta2**2 - beta1*beta2
    zp = sqrt(a) + sqrt(a-b)
    zm = sqrt(a) - sqrt(a-b)
    zmb = 1/zp
    z1 = sqrt(x1**2+b) - x1
    z2 = sqrt(x2**2+b) + x2
    z1b = 1/(sqrt(x1**2+b)+x1)
    K1 = -0.5*log(((z1b-zmb)*(zp-z1))/((zp+z1)*(z1b+zmb)))**2 &
          -2*Li2((2*zmb*(zp-z1))/((zp-zm)*(zmb+z1b))) &
          -2*Li2((-2*zp*(zm+z1))/((zp-zm)*(zp-z1)))
    K2 = -0.5*log(((z2-zm)*(zp-z2))/((zp+z2)*(z2+zm)))**2 &
          -2*Li2((2*zm*(zp-z2))/((zp-zm)*(zm+z2))) &
          -2*Li2((-2*zp*(zm+z2))/((zp-zm)*(zp-z2)))
    I = (K2 - K1) * (1-beta1*beta2)/sqrt(a-b)
  end function I_mm_eps

\eatline
\nwendcode{}\nwbegindocs{240}@For $k_1^2 \neq 0$ and $k_2^2 \neq 0$ ,computes the expression
\begin{equation*}
  I_0(k_1, k_2) = \frac{1}{\beta} \log\frac{1+\beta}{1-\beta},
  \quad
  \beta = \sqrt{1-\frac{k_1^2k_2^2}{(k_1 \cdot k_2)^2}}
\end{equation*}
\nwenddocs{}\nwbegincode{241}\moddef{virtual: procedures}\plusendmoddef
  function I_mm_0 (k1, k2) result (I)
    type(vector4_t), intent(in) :: k1, k2
    real(default) :: I
    real(default) :: beta
    beta = sqrt (1-k1**2*k2**2/(k1*k2)**2)
    I = log((1+beta)/(1-beta))/beta
  end function I_mm_0

\eatline
\nwendcode{}\nwbegindocs{242}\nwdocspar
\clearpage
%%%%%%%%%%%%%%%%%%%%%%%%%%%%%%%%%%%%%%%%%%%%%%%%%%%%%%%%%%%%%%%%%%%%%%%%
\section{Real Subtraction}
\nwenddocs{}\nwbegincode{243}\moddef{\code{}real{\char95}subtraction.f90\edoc{}}\endmoddef
\LA{}File header\RA{}

module real_subtraction

\LA{}Use kinds\RA{}
\LA{}Use strings\RA{}
  use constants
  use diagnostics
  use pdg_arrays
  use model_data
  use physics_defs
  use sm_physics
  use lorentz
  use flavors
  use fks_regions

\LA{}Standard module head\RA{}

\LA{}real subtraction: public\RA{}

\LA{}real subtraction: types\RA{}

contains
 
\LA{}real subtraction: procedures\RA{}

end module real_subtraction
\eatline
\nwendcode{}\nwbegindocs{244}\nwdocspar
\subsubsection{Soft subtraction terms}
In the soft limit, the real matrix element behaves as
\begin{equation*}
  \mathcal{R}_{\rm{soft}} = 4\pi\alpha_s \left[\sum_{i \neq j} \mathcal{B}_{ij} \frac{k_i \cdot k_j}{(k_i \cdot k)(k_j \cdot k)} - \mathcal{B} \sum_{i} \frac{k_i^2}{(k_i \cdot k)^2}C_i\right],
\end{equation*}
where $k$ denotes the momentum of the emitted parton. The quantity $\mathcal{B}_{ij}$ is called the color-correlated Born matrix element defined as
\begin{equation*}
  \mathcal{B}_{ij} = \frac{1}{2s} \sum_{\stackrel{colors}{spins}} \mathcal{M}_{\{c_k\}}\left(\mathcal{M}^\dagger_{\{c_k\}}\right)_{\stackrel{c_i \rightarrow c_i'}{c_j \rightarrow c_j'}} T^a_{c_i,c_i'} T^a_{c_j,c_j'}.
\end{equation*}
\nwenddocs{}\nwbegincode{245}\moddef{real subtraction: types}\endmoddef
  type :: soft_subtraction_t
    real(default), dimension(:), allocatable :: value
    type(region_data_t) :: reg_data
    integer :: nlegs_born, nlegs_real
    real(default), dimension(:,:), allocatable :: momentum_matrix
    logical :: use_internal_color_correlations = .true.
    logical :: use_internal_spin_correlations = .false.
  contains
  \LA{}real subtraction: soft sub: TBP\RA{}
  end type soft_subtraction_t
  
\eatline
\nwendcode{}\nwbegincode{246}\moddef{real subtraction: soft sub: TBP}\endmoddef
  procedure :: init => soft_subtraction_init
\nwendcode{}\nwbegincode{247}\moddef{real subtraction: procedures}\endmoddef
  subroutine soft_subtraction_init (sub_soft, reg_data, nlegs_born, &
                                    nlegs_real)
    class(soft_subtraction_t), intent(inout) :: sub_soft
    type(region_data_t), intent(in) :: reg_data
    integer, intent(in) :: nlegs_born, nlegs_real
    sub_soft%reg_data = reg_data
    sub_soft%nlegs_born = nlegs_born
    sub_soft%nlegs_real = nlegs_real
    allocate (sub_soft%value (reg_data%n_regions))
    allocate (sub_soft%momentum_matrix &
              (nlegs_born, nlegs_born))
  end subroutine soft_subtraction_init
 
\eatline
\nwendcode{}\nwbegindocs{248}The treatment of the momentum $k$ follows the discussion about the soft limit of the partition functions (ref????). 
The parton momentum is pulled out, $k = E \hat{k}$. In fact, we will substitute $\hat{k}$ for $k$ throughout the code, 
because the energy will factor out of the equation when the soft $\mathcal{S}$-function is multiplied.
The soft momentum is a unit vector, because $k^2 = \left(k^0\right)^2 - \left(k^0\right)^2\hat{\vec{k}}^2 = 0$.
\nwenddocs{}\nwbegincode{249}\moddef{real subtraction: soft sub: TBP}\plusendmoddef
  procedure :: create_softvec_fsr => soft_subtraction_create_softvec_fsr
\nwendcode{}\nwbegincode{250}\moddef{real subtraction: procedures}\plusendmoddef
  function soft_subtraction_create_softvec_fsr &
       (sub_soft, p_born, y, phi, emitter) result (p_soft)
    class(soft_subtraction_t), intent(inout) :: sub_soft
    type(vector4_t), intent(in), dimension(:) :: p_born
    real(default), intent(in) :: y, phi
    integer, intent(in) :: emitter
    real(default) :: cpsi
    type(vector4_t) :: p_soft
    type(vector3_t) :: dir
    type(lorentz_transformation_t) :: rot
\nwendcode{}\nwbegindocs{251}The soft momentum is constructed by first creating a unit vector parallel to the emitter's Born momentum. This unit vector
is then rotated about the corresponding angles $y$ and $\phi$.
\nwenddocs{}\nwbegincode{252}\moddef{real subtraction: procedures}\plusendmoddef
    p_soft%p(0) = 1._default
    p_soft%p(1:3) = p_born(emitter)%p(1:3) / space_part_norm (p_born(emitter))
    dir = create_orthogonal (space_part (p_born(emitter)))
    rot = rotation (y, sqrt(1-y**2), dir)
    p_soft = rot*p_soft
    if (phi /= 0) then
      dir = space_part (p_born(emitter)) / &
            space_part_norm (p_born(emitter))
      rot = rotation (cos(phi), sin(phi), dir)
      p_soft = rot*p_soft
    end if
  end function soft_subtraction_create_softvec_fsr
  
\eatline
\nwendcode{}\nwbegindocs{253}Computation of $\mathcal{R}_{\rm{soft}}$:
\nwenddocs{}\nwbegincode{254}\moddef{real subtraction: soft sub: TBP}\plusendmoddef
  procedure :: compute => soft_subtraction_compute
\nwendcode{}\nwbegincode{255}\moddef{real subtraction: procedures}\plusendmoddef
  subroutine soft_subtraction_compute (sub_soft, p_born, p_real, &
       sqme_born, born_ij, y, y_soft, phi, alpha_s_born, alr, emitter)
    class(soft_subtraction_t), intent(inout) :: sub_soft
    type(vector4_t), intent(in), dimension(:) :: p_born
    type(vector4_t), intent(in), dimension(:) :: p_real
    real(default), intent(in) :: sqme_born
    real(default), intent(in), dimension(:,:) :: born_ij
    real(default), intent(in) :: y, y_soft, phi
    real(default), intent(in) :: alpha_s_born
    integer, intent(in) :: alr, emitter
    type(vector4_t) :: p_soft
    real(default) :: s_alpha_soft
    real(default) :: q02
    real(default) :: kb
    integer :: i, j

    p_soft = sub_soft%create_softvec_fsr (p_born, y_soft, phi, emitter)
    s_alpha_soft = sub_soft%reg_data%get_svalue_soft &
         (p_born, p_soft, alr, emitter)
    call sub_soft%compute_momentum_matrix (p_born, p_soft)
    sub_soft%value(alr) = 4*pi*alpha_s_born * s_alpha_soft
    kb = 0._default
    do i = 1, size (p_born)
       do j = 1, size (p_born)
          kb = kb + sub_soft%momentum_matrix (i,j) * &
               born_ij (i,j)
       end do 
    end do
    sub_soft%value(alr) = sub_soft%value(alr)*kb
    q02 = 4* vector4_get_component (p_born(1), 0) * &
         vector4_get_component (p_born(2), 0)
    sub_soft%value(alr) = 4/q02 * (1-y) * sub_soft%value(alr) 
  end subroutine soft_subtraction_compute

\eatline
\nwendcode{}\nwbegindocs{256}We have to multiply this with $\xi^2(1-y)$. Further, when applying
the soft $\mathcal{S}$-function, the energy of the radiated particle
is factored out. Thus we have $\xi^2/E_{em}^2(1-y) = 4/q_0^2(1-y)$. 
\nwenddocs{}\nwbegincode{257}\moddef{nlo data: soft subtraction computation}\endmoddef
\nwendcode{}\nwbegindocs{258}Computes the quantity $\mathcal{K}_{ij} = \frac{k_i \cdot
k_j}{(k_i\cdot k)(k_j\cdot k)}$. 
\nwenddocs{}\nwbegincode{259}\moddef{real subtraction: soft sub: TBP}\plusendmoddef
  procedure :: compute_momentum_matrix => &
       soft_subtraction_compute_momentum_matrix
\nwendcode{}\nwbegincode{260}\moddef{real subtraction: procedures}\plusendmoddef
  subroutine soft_subtraction_compute_momentum_matrix &
       (sub_soft, p_born, p_soft)
    class(soft_subtraction_t), intent(inout) :: sub_soft
    type(vector4_t), intent(in), dimension(:) :: p_born
    type(vector4_t), intent(in) :: p_soft
    real(default) :: num, deno1, deno2
    integer :: i, j
    do i = 1, sub_soft%nlegs_born
      do j = 1, sub_soft%nlegs_born
        if (i <= j) then
          num = p_born(i) * p_born(j)
          deno1 = p_born(i)*p_soft
          deno2 = p_born(j)*p_soft
          sub_soft%momentum_matrix(i,j) = num/(deno1*deno2)
        else
           !!! momentum matrix is symmetric.
          sub_soft%momentum_matrix(i,j) = sub_soft%momentum_matrix(j,i)
        end if
      end do
    end do
  end subroutine soft_subtraction_compute_momentum_matrix
  
\eatline
\nwendcode{}\nwbegindocs{261}\nwdocspar
\subsection{Collinear and soft-collinear subtraction terms}
This data type deals with the calculation of the collinear and
soft-collinear contribution to the cross section. 
\nwenddocs{}\nwbegincode{262}\moddef{real subtraction: types}\plusendmoddef
  type :: coll_subtraction_t
    real(default), dimension(:), allocatable :: value
    real(default), dimension(:), allocatable :: value_soft
    integer :: n_alr
    real(default), dimension(0:3,0:3) :: b_munu
    real(default) , dimension(0:3,0:3) :: k_perp_matrix
  contains
  \LA{}real subtraction: coll sub: TBP\RA{}
  end type coll_subtraction_t
  
\eatline
\nwendcode{}\nwbegindocs{263}\nwdocspar
\nwenddocs{}\nwbegincode{264}\moddef{real subtraction: coll sub: TBP}\endmoddef
  procedure :: init => coll_subtraction_init
\nwendcode{}\nwbegincode{265}\moddef{real subtraction: procedures}\plusendmoddef
  subroutine coll_subtraction_init (coll_sub, n_alr)
    class(coll_subtraction_t), intent(inout) :: coll_sub
    integer, intent(in) :: n_alr
    coll_sub%n_alr = n_alr
    allocate (coll_sub%value (n_alr))
    allocate (coll_sub%value_soft (n_alr))
  end subroutine coll_subtraction_init

\eatline
\nwendcode{}\nwbegindocs{266}To compute the collinear limit of $\mathcal{R}$, we follow the original FKS-paper. Here, the real amplitude is supposed to factorize
in the collinear limit. considering a splitting $g \rightarrow g(i)g(j)$,
\begin{equation*}
  \mathcal{A}^{(n)}(h_i,h_j,\{h_l\} \stackrel{i \parallel j}{\rightarrow} g_s \sum_{d_e} \sum_{h_e} C(d_e, b, c) 
                                                                          S_{gg}^{h_eh_ih_j}(z)\mathcal{A}_{d_e}^{(n-1)} (h_e, \{h_l\})
\end{equation*}
\textbf{Explain quantities.} Evaluating this expression leads to
\begin{equation*}
  |\mathcal{A}^{(n)}(h_i,h_j,\{h_l\})|^2 \stackrel{i \parallel j}{\rightarrow}  |\mathcal{N}^{(n)}(h_i,h_j,\{h_l\}|^2 + \mathcal{R}(h_i,h_j,\{h_l\})|^2,
\end{equation*}
where $\mathcal{R}$ contains all contributions with spin-correlated amplitudes, i.e. terms like $\mathcal{A}_{d_e}^{(n-1)}(+,\{h_l\})\left(\mathcal{A}_{d_e}^{(n-1)}(-,\{h_l\})\right)^*$. Explicitly,
\begin{equation*}
  \sum_{h_i,h_j,\{h_l\}}\sum_{\{d_l\}} |\mathcal{N}^{(n)}(h_i,h_j,\{h_l\})|^2 = \frac{4\pi\alpha_s}{k_i \cdot k_j} P_{gg}(z)|\mathcal{A}^{(n-1)}|^2,
\end{equation*}
and
\begin{align*}
  \sum_{h_i,h_j,\{h_l\}}\sum_{\{d_l\}} &\mathcal{R}(h_i,h_j,\{h_l\}) = 32\pi\alpha_s C_A z(1-z) \\
                        &\times \underbrace{Re\left\{\frac{\langle ij \rangle}{[ij]}\sum_{\{h_l\}}\sum_{d_e, \{d_l\}} \mathcal{A}^{(n-1)}_{d_e}(+,\{h_l\})\left(\mathcal{A}^{(n-1)}_{d_e}(-,\{h_l\})\right)^*\right\}}
                        _{-\tilde{\mathcal{M}}^{(n-1)}/(2k_i\cdot k_j)}.
\end{align*}
This yields
\begin{equation}
  \mathcal{M}^{(n)} \stackrel{i \parallel j}{\rightarrow} \frac{4\pi\alpha_s}{k_i \cdot k_j} P_{gg}(z)\mathcal{M}^{(n-1)} - \frac{16\pi\alpha_s}{k_i \cdot k_j} C_A z(1-z)\tilde{\mathcal{M}}^{(n-1)}.
  \label{coll1}
\end{equation}
The equivalent expression for a $g \rightarrow qq$-splitting is given by
\begin{equation}  
  \mathcal{M}^{(n)} \stackrel{i \parallel j}{\rightarrow} \frac{4\pi\alpha_s}{k_i \cdot k_j} P_{qg}(z)\mathcal{M}^{(n-1)} + \frac{16\pi\alpha_s}{k_i \cdot k_j} T_F z(1-z)\tilde{\mathcal{M}}^{(n-1)}.
  \label{coll2}
\end{equation}
\nwenddocs{}\nwbegincode{267}\moddef{real subtraction: coll sub: TBP}\plusendmoddef
  procedure :: compute => coll_subtraction_compute
\nwendcode{}\nwbegincode{268}\moddef{real subtraction: procedures}\plusendmoddef
  subroutine coll_subtraction_compute &
       (coll_sub, sregion, p_born, sqme_born, sqme_born_sc, &
        xi, alpha_s, alr, soft_in)
    class(coll_subtraction_t), intent(inout) :: coll_sub
    type(singular_region_t), intent(in) :: sregion
    type(vector4_t), intent(in), dimension(:) :: p_born
    real(default), intent(in) :: sqme_born
    real(default), intent(in) :: sqme_born_sc
    real(default), intent(in) :: xi, alpha_s
    integer, intent(in) :: alr
\nwendcode{}\nwbegindocs{269}The function obtains a flag to indicate whether the limit to be computed is also soft. Moreover, this consideration
explains the structure of the code. In the soft limit, we find $z \rightarrow 0$ as well as $\xi \rightarrow 0$. However,
the quantity $z/\xi$ is finite, because 
\begin{equation*}
  \frac{z}{\xi} = \frac{p_{rad}^0}{\bar{p}_{em}^0} \frac{q^0}{2p_{em}^0} = \frac{q^0}{2\bar{p}^0_{em}}.
\end{equation*}
Thus, all expressions are written in terms of this well-behaved quantity. Recalling that there is a prefactor of $(\xi^2z)^{-1}$,
it is necessary to expand with $z$ once.
\nwenddocs{}\nwbegincode{270}\moddef{real subtraction: procedures}\plusendmoddef
    logical, intent(in), optional :: soft_in
    real(default) :: res
    real(default) :: q0, z, p0
    real(default) :: zoxi, onemz
    real(default) :: pggz, pqgz
    integer :: nlegs, emitter
    integer :: flv_em, flv_rad
    logical :: soft
    if (present (soft_in)) then
      soft = soft_in
    else
      soft = .false.
    end if
    nlegs = size (sregion%flst_real%flst)
    emitter = sregion%emitter
    flv_rad = sregion%flst_real%flst(nlegs)
    flv_em = sregion%flst_real%flst(emitter)
    p0 = vector4_get_component (p_born(emitter),0)
    if (sregion%emitter <= 2) then
      coll_sub%value(alr) = 0   
      return
    else
      q0 = vector4_get_component (p_born(1), 0) + &
           vector4_get_component (p_born(2), 0)
      !!! Here, z corresponds to 1-z in the formulas of arXiv:1002.2581; 
      !!! the integrand is symmetric under this variable change
      zoxi = q0/(2*p0)
      z = xi*zoxi; onemz = 1-z
      if (is_gluon(flv_em) .and. is_gluon(flv_rad)) then
\nwendcode{}\nwbegindocs{271}Implementation of equation (\ref{coll1}). Note that an additional factor $z$, so that in the last step, the whole expression is divided by $z/\xi$.
\nwenddocs{}\nwbegincode{272}\moddef{real subtraction: procedures}\plusendmoddef
         pggz = 2*CA*(z**2*onemz + z**2/onemz + onemz)
         res = pggz*sqme_born - 4*CA*z**2*onemz*sqme_born_sc
         res = res/zoxi
      else if (is_quark(abs(flv_em)) .and. is_quark (abs(flv_rad))) then
\nwendcode{}\nwbegindocs{273}Equation \ref{coll2}
\nwenddocs{}\nwbegincode{274}\moddef{real subtraction: procedures}\plusendmoddef
         pqgz = TR*z*(1-2*z*onemz)
         res = pqgz*sqme_born + 4*TR*z**2*onemz*sqme_born_sc
         res = res/zoxi
      else if (is_quark (abs(flv_em)) .and. is_gluon (flv_rad)) then
         res = sqme_born*CF*(1+onemz**2)/zoxi
      else
         call msg_fatal ('Impossible flavor structure in collinear counterterm!') 
      end if
    end if
    res = res /(p0**2*onemz*zoxi)
    res = res * 4*pi*alpha_s 

    if (soft) then
      coll_sub%value_soft (alr) = res
    else
      coll_sub%value (alr) = res
    end if
  end subroutine coll_subtraction_compute
  
\eatline
\nwendcode{}\nwbegindocs{275}Here, $\xi = 0$ is already required.
\nwenddocs{}\nwbegincode{276}\moddef{real subtraction: coll sub: TBP}\plusendmoddef
  procedure :: compute_soft_limit => coll_subtraction_compute_soft_limit
\nwendcode{}\nwbegincode{277}\moddef{real subtraction: procedures}\plusendmoddef
  subroutine coll_subtraction_compute_soft_limit &
       (coll_sub, sregion, p_born, sqme_born, &
        sqme_born_sc, xi, alpha_s, alr)
    class(coll_subtraction_t), intent(inout) :: coll_sub
    type(singular_region_t), intent(in) :: sregion
    type(vector4_t), intent(in), dimension(:) :: p_born
    real(default), intent(in) :: sqme_born
    real(default), intent(in) :: sqme_born_sc
    real(default) :: xi, alpha_s
    integer, intent(in) :: alr
    call coll_sub%compute (sregion, p_born, sqme_born, &
                           sqme_born_sc, xi, alpha_s, alr, .true.)
  end subroutine coll_subtraction_compute_soft_limit
  
\eatline
\nwendcode{}\nwbegindocs{278}\nwdocspar
\subsection{Real Subtraction}
Just a container for the real kinematics variables.
\nwenddocs{}\nwbegincode{279}\moddef{real subtraction: public}\endmoddef
  public :: real_kinematics_t
\nwendcode{}\nwbegincode{280}\moddef{real subtraction: types}\plusendmoddef
  type :: real_kinematics_t
    real(default) :: xi_tilde
    real(default) :: phi
    real(default), dimension(:), allocatable :: xi_max, y
    real(default), dimension(3) :: jac
    type(vector4_t), dimension(:), allocatable :: p_real
    real(default), dimension(:), allocatable :: jac_rand
    real(default), dimension(:), allocatable :: y_soft
  end type real_kinematics_t

\eatline
\nwendcode{}\nwbegincode{281}\moddef{real subtraction: public}\plusendmoddef
  public :: real_subtraction_t
\nwendcode{}\nwbegincode{282}\moddef{real subtraction: types}\plusendmoddef
  type :: real_subtraction_t
    type(region_data_t) :: reg_data
    type(vector4_t), dimension(:), allocatable :: p_born, p_real
    type(real_kinematics_t) :: fks_kinematics
    integer :: current_alr = 0
    real(default), dimension(:), allocatable :: sqme_real_bare
    real(default), dimension(:), allocatable :: sqme_born
    real(default), dimension(:,:,:), allocatable :: sqme_born_cc
    complex(default) :: me_sc
    type(soft_subtraction_t) :: sub_soft
    type(coll_subtraction_t) :: sub_coll
    logical, dimension(:), allocatable :: sc_required
  contains
  \LA{}real subtraction: real subtraction: TBP\RA{}
  end type real_subtraction_t

\eatline
\nwendcode{}\nwbegindocs{283}Initializer
\nwenddocs{}\nwbegincode{284}\moddef{real subtraction: real subtraction: TBP}\endmoddef
  procedure :: init => real_subtraction_init
\nwendcode{}\nwbegincode{285}\moddef{real subtraction: procedures}\plusendmoddef
  subroutine real_subtraction_init (rsub, reg_data, nlegs_born, &
                                    nlegs_real)
    class(real_subtraction_t), intent(inout) :: rsub
    type(region_data_t), intent(in) :: reg_data
    integer, intent(in) :: nlegs_born, nlegs_real
    integer :: alr, i_uborn
    rsub%reg_data = reg_data
    allocate (rsub%p_born (nlegs_born))
    allocate (rsub%p_real (nlegs_real))
    allocate (rsub%sqme_real_bare (reg_data%n_flv_real))
    allocate (rsub%sqme_born (reg_data%n_flv_born))
    allocate (rsub%sqme_born_cc (nlegs_born, nlegs_born, &
                                 reg_data%n_flv_born))
    allocate (rsub%sc_required (reg_data%n_regions))
    do alr = 1, reg_data%n_regions
       i_uborn = reg_data%regions(alr)%uborn_index
       rsub%sc_required(alr) = &
          reg_data%flv_born(i_uborn)%count_particle (GLUON) > 0
    end do
     
    call rsub%sub_soft%init (reg_data, nlegs_born, nlegs_real)
    call rsub%sub_coll%init (reg_data%n_regions)
  end subroutine real_subtraction_init

\eatline
\nwendcode{}\nwbegindocs{286}\nwdocspar
\nwenddocs{}\nwbegincode{287}\moddef{real subtraction: real subtraction: TBP}\plusendmoddef
  procedure :: set_momenta => real_subtraction_set_momenta
\nwendcode{}\nwbegincode{288}\moddef{real subtraction: procedures}\plusendmoddef
  subroutine real_subtraction_set_momenta (rsub, p_born, p_real)
    class(real_subtraction_t), intent(inout) :: rsub
    type(vector4_t), intent(in), dimension(:) :: p_born, p_real
    rsub%p_born = p_born; rsub%p_real = p_real
  end subroutine real_subtraction_set_momenta

\eatline
\nwendcode{}\nwbegindocs{289}\nwdocspar
\nwenddocs{}\nwbegincode{290}\moddef{real subtraction: real subtraction: TBP}\plusendmoddef
  procedure :: set_fks_kinematics => real_subtraction_set_fks_kinematics
\nwendcode{}\nwbegincode{291}\moddef{real subtraction: procedures}\plusendmoddef
  subroutine real_subtraction_set_fks_kinematics (rsub, fks_kinematics)
    class(real_subtraction_t), intent(inout) :: rsub
    type(real_kinematics_t), intent(in) :: fks_kinematics
    rsub%fks_kinematics = fks_kinematics
  end subroutine real_subtraction_set_fks_kinematics

\eatline
\nwendcode{}\nwbegindocs{292}\nwdocspar
\nwenddocs{}\nwbegincode{293}\moddef{real subtraction: real subtraction: TBP}\plusendmoddef
  procedure :: set_alr => real_subtraction_set_alr
\nwendcode{}\nwbegincode{294}\moddef{real subtraction: procedures}\plusendmoddef
  subroutine real_subtraction_set_alr (rsub, alr)
    class(real_subtraction_t), intent(inout) :: rsub
    integer, intent(in) :: alr
    rsub%current_alr = alr
  end subroutine real_subtraction_set_alr

\eatline
\nwendcode{}\nwbegindocs{295}@\subsection{The real contribution to the cross section}
In each singular region $\alpha$, the real contribution to $\sigma$ is
given by the second summand of eqn. \ref{fks: sub: complete}, 
\begin{equation}
  \sigma^\alpha_{\text{real}} = \int d\Phi_n \int_0^{2\pi} d\phi
  \int_{-1}^1 dy \int_0^{\xi_{\text{max}}} d\xi
  \left(\frac{1}{\xi}\right)_+ \left(\frac{1}{1-y}\right)_+
  \underbrace{\frac{J(\Phi_n, \xi, y, \phi)}{\xi}
    \left[(1-y)\xi^2\mathcal{R}^\alpha(\Phi_{n+1})\right]}_{g^\alpha(\xi,y)}. 
\end{equation}
Writing out the plus-distribution and introducing $\tilde{\xi} =
\xi/\xi_{\text{max}}$ to set the upper integration limit to 1,  this
turns out to be equal to 
\begin{equation}
  \begin{split}
    \sigma^\alpha_{\rm{real}} &= \int d\Phi_n \int_0^{2/pi}d\phi
    \int_{-1}^1 \frac{dy}{1-y} \Bigg\{\int_0^1
    d\tilde{\xi}\Bigg[\frac{g^\alpha(\tilde{\xi}\xi_{\rm{max}},y)}{\tilde{\xi}}
    - \underbrace{\frac{g^\alpha(0,y)}{\tilde{\xi}}}_{\text{soft}} -
    \underbrace{\frac{g^\alpha(\tilde{\xi}\xi_{\rm{max}},1)}{\tilde{\xi}}}_{\text{coll.}}
    +
    \underbrace{\frac{g^\alpha(0,1)}{\tilde{\xi}}}_{\text{coll.+soft}}\Bigg]
    \\ 
&+ \left[\log\xi_{\rm{max}}(y)g^\alpha(0,y) - \log\xi_{\rm{max}}(1)g^\alpha(0,1)\right]\Bigg\}.
  \end{split}
\end{equation}
This formula is implemented in \texttt{compute\_sqme\_real\_fin}
\nwenddocs{}\nwbegincode{296}\moddef{real subtraction: real subtraction: TBP}\plusendmoddef
  procedure :: compute => real_subtraction_compute
\nwendcode{}\nwbegincode{297}\moddef{real subtraction: procedures}\plusendmoddef
  function real_subtraction_compute (rsub, emitter, alpha_s) result (sqme)
    class(real_subtraction_t), intent(inout) :: rsub 
    integer, intent(in) :: emitter
    real(default) :: alpha_s
    real(default) :: sqme
    integer :: alr

    sqme = 0._default
    !!! Loop over all singular regions
    do alr = 1, size (rsub%reg_data%regions)
         if (emitter == rsub%reg_data%regions(alr)%emitter) then
            call rsub%set_alr (alr)
            sqme = sqme + rsub%evaluate_region (emitter, alpha_s) 
         end if
    end do
    sqme = sqme*rsub%get_phs_factor ()
  end function real_subtraction_compute

\eatline
\nwendcode{}\nwbegindocs{298}\nwdocspar
\nwenddocs{}\nwbegincode{299}\moddef{real subtraction: real subtraction: TBP}\plusendmoddef
  procedure :: evaluate_region => real_subtraction_evaluate_region
\nwendcode{}\nwbegincode{300}\moddef{real subtraction: procedures}\plusendmoddef
  function real_subtraction_evaluate_region (rsub, emitter, &
                                             alpha_s) result (sqme)
    class(real_subtraction_t), intent(inout) :: rsub
    integer, intent(in) :: emitter
    real(default), intent(in) :: alpha_s
    real(default) :: sqme
    integer :: i_real
    real(default) :: sqme0, sqme_soft, sqme_coll, sqme_cs, sqme_remn
    real(default) :: s_alpha
    real(default) :: xi, xi_max, xi_tilde, y, phi
    xi_tilde = rsub%fks_kinematics%xi_tilde
    xi_max = rsub%fks_kinematics%xi_max(emitter)
    xi = xi_tilde * xi_max
    y = rsub%fks_kinematics%y(emitter)
    phi = rsub%fks_kinematics%phi
    associate (region => rsub%reg_data%regions(rsub%current_alr))
      i_real = region%real_index
      sqme0 = rsub%sqme_real_bare (i_real)
      s_alpha = rsub%reg_data%get_svalue (rsub%p_real, rsub%current_alr, emitter)
      sqme0 = sqme0 * s_alpha
      sqme0 = sqme0 * region%mult
      sqme0 = sqme0 * region%double_fsr_factor (rsub%p_real)
      call rsub%compute_sub_soft (emitter, alpha_s)
      call rsub%compute_sub_coll (emitter, alpha_s)
      call rsub%compute_sub_coll_soft (emitter, alpha_s)
      call rsub%get_limit_sqme (sqme_soft, sqme_coll, sqme_cs)
\nwendcode{}\nwbegindocs{301}Supply the $\tilde{\xi}^{-1}(1-y)^{-1}$-factor; note that this is not done for the pure real matrix element because 
it directly cancels out with the numerator. Further, the jacobians are supplied in their appropriate limits.
\nwenddocs{}\nwbegincode{302}\moddef{real subtraction: procedures}\plusendmoddef
      associate (jac => rsub%fks_kinematics%jac)
         sqme0 = sqme0 * xi**2/xi_tilde * jac(1)
         sqme_soft = sqme_soft/(1-y)/xi_tilde * jac(2)
         sqme_coll = sqme_coll/(1-y)/xi_tilde * jac(3)
         sqme_cs = sqme_cs/(1-y)/xi_tilde * jac(2)
      end associate
      sqme_remn = (sqme_soft - sqme_cs)*log(xi_max)*xi_tilde
      sqme = sqme0 - sqme_soft - sqme_coll + sqme_cs + sqme_remn
      sqme = sqme * rsub%fks_kinematics%jac_rand (emitter)
    end associate
    !!! This occurs if the result is NaN.
    if (sqme /= sqme) & 
       call write_computation_status ()
  contains 
    subroutine write_computation_status ()
       integer :: i_uborn
       i_uborn = rsub%reg_data%regions(rsub%current_alr)%uborn_index
       print *, 'alr: ', rsub%current_alr
       print *, 'emitter: ', emitter
       print *, 'xi_max: ', xi_max
       print *, 'xi: ', xi, 'y: ', y
       print *, 'sqme_born: ', rsub%sqme_born(i_uborn), &
       'sqme_real: ', sqme0, &
       'sqme_soft: ', sqme_soft, &
       'sqme_coll: ', sqme_coll, &
       'sqme_coll-soft: ', sqme_cs, &
       'sqme_remn: ', sqme_remn
    end subroutine write_computation_status

  end function real_subtraction_evaluate_region

\eatline
\nwendcode{}\nwbegindocs{303}\nwdocspar
\nwenddocs{}\nwbegincode{304}\moddef{real subtraction: real subtraction: TBP}\plusendmoddef
  procedure :: get_limit_sqme => real_subtraction_get_limit_sqme
\nwendcode{}\nwbegincode{305}\moddef{real subtraction: procedures}\plusendmoddef
  subroutine real_subtraction_get_limit_sqme (rsub, &
                  sqme_soft, sqme_coll, sqme_cs)
    class(real_subtraction_t), intent(in) :: rsub
    real(default), intent(out) :: sqme_soft, sqme_coll, sqme_cs
    integer :: alr
    alr = rsub%current_alr
    sqme_soft = rsub%sub_soft%value(alr)
    sqme_coll = rsub%sub_coll%value(alr)
    sqme_cs = rsub%sub_coll%value_soft(alr)
  end subroutine real_subtraction_get_limit_sqme

\eatline
\nwendcode{}\nwbegindocs{306}\nwdocspar
\nwenddocs{}\nwbegincode{307}\moddef{real subtraction: real subtraction: TBP}\plusendmoddef
  procedure :: get_phs_factor => real_subtraction_get_phs_factor
\nwendcode{}\nwbegincode{308}\moddef{real subtraction: procedures}\plusendmoddef
  function real_subtraction_get_phs_factor (rsub) result (factor)
    class(real_subtraction_t), intent(in) :: rsub
    real(default) :: factor
    real(default) :: s
    s = (rsub%p_born(1)+rsub%p_born(2))**2
    factor = s / (8*twopi2) 
  end function real_subtraction_get_phs_factor

\eatline
\nwendcode{}\nwbegindocs{309}\nwdocspar
\subsection{xxx}
\nwenddocs{}\nwbegincode{310}\moddef{real subtraction: real subtraction: TBP}\plusendmoddef
  procedure :: compute_sub_soft => real_subtraction_compute_sub_soft
\nwendcode{}\nwbegincode{311}\moddef{real subtraction: procedures}\plusendmoddef
  subroutine real_subtraction_compute_sub_soft &
                             (rsub, emitter, alpha_s)
    class(real_subtraction_t), intent(inout) :: rsub
    integer, intent(in) :: emitter
    real(default), intent(in) :: alpha_s
    integer :: alr
    alr = rsub%current_alr
    associate (sregion => rsub%reg_data%regions(alr))
       if (sregion%has_soft_divergence ()) then
             call rsub%sub_soft%compute (rsub%p_born, rsub%p_real, &
                                    rsub%sqme_born(sregion%uborn_index), &
                                    rsub%sqme_born_cc(:,:,sregion%uborn_index), &
                                    rsub%fks_kinematics%y(emitter), &
                                    rsub%fks_kinematics%y_soft(emitter), &
                                    rsub%fks_kinematics%phi, &
                                    alpha_s, alr, emitter)
       else
          rsub%sub_soft%value(alr) = 0._default
       end if
    end associate
  end subroutine real_subtraction_compute_sub_soft

\eatline
\nwendcode{}\nwbegindocs{312}\nwdocspar
\nwenddocs{}\nwbegincode{313}\moddef{real subtraction: real subtraction: TBP}\plusendmoddef
  procedure :: compute_sub_coll => real_subtraction_compute_sub_coll 
\nwendcode{}\nwbegincode{314}\moddef{real subtraction: procedures}\plusendmoddef
  subroutine real_subtraction_compute_sub_coll (rsub, em, alpha_s)
    class(real_subtraction_t), intent(inout) :: rsub
    integer, intent(in) :: em
    real(default), intent(in) :: alpha_s
    real(default) :: xi
    real(default) :: sqme_sc
    complex(default) :: prod1, prod2
    integer :: alr
    alr = rsub%current_alr
    xi = rsub%fks_kinematics%xi_tilde * rsub%fks_kinematics%xi_max (em)
    associate (sregion => rsub%reg_data%regions(alr)) 
       if (sregion%has_collinear_divergence ()) then
          if (rsub%sc_required(alr)) then
             call spinor_product (rsub%p_real(em), rsub%p_real(rsub%reg_data%nlegs_real), &
                                  prod1, prod2)
             sqme_sc = real (prod1/prod2*rsub%me_sc)
          else
             sqme_sc = 0._default
          end if
          call rsub%sub_coll%compute (sregion, rsub%p_born, &
                                      rsub%sqme_born(sregion%uborn_index), &
                                      sqme_sc, xi, alpha_s, alr)
       else
          rsub%sub_coll%value(alr) = 0._default
       end if
    end associate
  end subroutine real_subtraction_compute_sub_coll

\eatline
\nwendcode{}\nwbegindocs{315}\nwdocspar
\nwenddocs{}\nwbegincode{316}\moddef{real subtraction: real subtraction: TBP}\plusendmoddef
  procedure :: compute_sub_coll_soft => real_subtraction_compute_sub_coll_soft
\nwendcode{}\nwbegincode{317}\moddef{real subtraction: procedures}\plusendmoddef
  subroutine real_subtraction_compute_sub_coll_soft (rsub, em, alpha_s)
    class(real_subtraction_t), intent(inout) :: rsub
    integer, intent(in) :: em
    real(default), intent(in) :: alpha_s
    complex(default) :: prod1, prod2
    real(default) :: sqme_sc
    real(default), parameter :: xi = 0
    integer :: alr
    alr = rsub%current_alr
    associate (sregion => rsub%reg_data%regions(alr))
       if (sregion%has_collinear_divergence ()) then
          if (rsub%sc_required(alr)) then
             call spinor_product (rsub%p_real(em), rsub%p_real(rsub%reg_data%nlegs_real), &
                                  prod1, prod2)
             sqme_sc = real (prod1/prod2*rsub%me_sc)
          else
             sqme_sc = 0._default
          end if
          call rsub%sub_coll%compute_soft_limit (sregion, rsub%p_born, &
                             rsub%sqme_born(sregion%uborn_index), &
                             sqme_sc, xi, alpha_s, alr)
       else
          rsub%sub_coll%value_soft(alr) = 0._default
       end if
    end associate
  end subroutine real_subtraction_compute_sub_coll_soft

\eatline
\nwendcode{}\nwbegindocs{318}\nwdocspar
\clearpage
%%%%%%%%%%%%%%%%%%%%%%%%%%%%%%%%%%%%%%%%%%%%%%%%%%%%%%%%%%%%%%%%%%%%%%%%
\section{Combining the FKS Pieces}
\nwenddocs{}\nwbegincode{319}\moddef{\code{}nlo{\char95}data.f90\edoc{}}\endmoddef
\LA{}File header\RA{}

module nlo_data

\LA{}Use kinds\RA{}
\LA{}Use strings\RA{}
  use io_units
  use constants, only: pi, twopi
  use unit_tests
  use diagnostics
  use physics_defs
  use process_constants !NODEP!
  use sm_physics
  use os_interface
  use model_data
  use parser
  use pdg_arrays
  use particle_specifiers
  use phs_single
  use state_matrices
  use interactions
  use lorentz
  use prc_core
  use pdg_arrays
  use sf_base
  use colors
  use flavors
  use fks_regions
  use phs_fks
  use virtual
  use real_subtraction

\LA{}Standard module head\RA{}

\LA{}nlo data: public\RA{}

\LA{}nlo data: types\RA{}

contains
 
\LA{}nlo data: procedures\RA{}

end module nlo_data
\eatline
\nwendcode{}\nwbegindocs{320}We need to associate singular regions with compatible color flows.
\nwenddocs{}\nwbegincode{321}\moddef{nlo data: public}\endmoddef
  public :: fks_template_t
\nwendcode{}\nwbegincode{322}\moddef{nlo data: types}\endmoddef
  type :: fks_template_t
    integer :: mapping_type
    logical :: count_kinematics = .false.
    real(default) :: fks_dij_exp1
    real(default) :: fks_dij_exp2
  contains
  \LA{}nlo data: fks template: TBP\RA{}
  end type fks_template_t

\eatline
\nwendcode{}\nwbegindocs{323}\nwdocspar
\nwenddocs{}\nwbegincode{324}\moddef{nlo data: fks template: TBP}\endmoddef
  procedure :: write => fks_template_write
\nwendcode{}\nwbegincode{325}\moddef{nlo data: procedures}\endmoddef
  subroutine fks_template_write (object, unit)
    class(fks_template_t), intent(in) :: object
    integer, intent(in), optional :: unit
    integer :: u
    u = given_output_unit (unit)
    write (u,'(1x,A)') 'FKS Template: '
    write (u,'(1x,A,I0)') 'Mapping Type: ', object%mapping_type
    write (u,'(1x,A,ES4.3,ES4.3)') 'd_ij exponentials: ', object%fks_dij_exp1, object%fks_dij_exp2
  end subroutine fks_template_write

\eatline
\nwendcode{}\nwbegindocs{326}\nwdocspar
\nwenddocs{}\nwbegincode{327}\moddef{nlo data: fks template: TBP}\plusendmoddef
  procedure :: set_dij_exp => fks_template_set_dij_exp
\nwendcode{}\nwbegincode{328}\moddef{nlo data: procedures}\plusendmoddef
  subroutine fks_template_set_dij_exp (object, exp1, exp2)
    class(fks_template_t), intent(inout) :: object
    real(default), intent(in) :: exp1, exp2
    object%fks_dij_exp1 = exp1
    object%fks_dij_exp2 = exp2
  end subroutine fks_template_set_dij_exp

\eatline
\nwendcode{}\nwbegindocs{329}\nwdocspar
\nwenddocs{}\nwbegincode{330}\moddef{nlo data: fks template: TBP}\plusendmoddef
  procedure :: set_mapping_type => fks_template_set_mapping_type
\nwendcode{}\nwbegincode{331}\moddef{nlo data: procedures}\plusendmoddef
  subroutine fks_template_set_mapping_type (object, val)
    class(fks_template_t), intent(inout) :: object
    integer, intent(in) :: val
    object%mapping_type = val
  end subroutine fks_template_set_mapping_type

\eatline
\nwendcode{}\nwbegindocs{332}\nwdocspar
\nwenddocs{}\nwbegincode{333}\moddef{nlo data: fks template: TBP}\plusendmoddef
  procedure :: set_counter => fks_template_set_counter
\nwendcode{}\nwbegincode{334}\moddef{nlo data: procedures}\plusendmoddef
  subroutine fks_template_set_counter (object)
    class(fks_template_t), intent(inout) :: object
    object%count_kinematics = .true.
  end subroutine fks_template_set_counter

\eatline
\nwendcode{}\nwbegindocs{335}\nwdocspar
\nwenddocs{}\nwbegincode{336}\moddef{nlo data: types}\plusendmoddef
  type :: ftuple_color_map_t
    integer :: index
    integer :: color_index
    type(ftuple_color_map_t), pointer :: next
    type(ftuple_color_map_t), pointer :: prev
  contains
  \LA{}nlo data: color map: TBP\RA{}  
  end type ftuple_color_map_t

\eatline
\nwendcode{}\nwbegindocs{337}\nwdocspar
\nwenddocs{}\nwbegincode{338}\moddef{nlo data: color map: TBP}\endmoddef
  procedure :: init => ftuple_color_map_init
\nwendcode{}\nwbegincode{339}\moddef{nlo data: procedures}\plusendmoddef
  subroutine ftuple_color_map_init (icm)
    class(ftuple_color_map_t), intent(inout), target :: icm
    icm%index = 0
    icm%color_index = 0
    nullify (icm%next)
    nullify (icm%prev)
  end subroutine ftuple_color_map_init

\eatline
\nwendcode{}\nwbegindocs{340}Explain
\nwenddocs{}\nwbegincode{341}\moddef{nlo data: color map: TBP}\plusendmoddef
  procedure :: present => ftuple_color_map_present  
\nwendcode{}\nwbegincode{342}\moddef{nlo data: procedures}\plusendmoddef
  function ftuple_color_map_present (icm, color_index) result(pres)
    class(ftuple_color_map_t), intent(in), target :: icm
    integer, intent(in) :: color_index
    logical :: pres
    type(ftuple_color_map_t), pointer :: current
    select type (icm)
    type is (ftuple_color_map_t)
    current => icm
    pres = .false.
    do
      if (current%color_index == color_index) then
        pres = .true.
        exit
      else
        if (associated (current%next)) then
          current => current%next
        else
          exit
        end if
      end if
    end do
    end select
  end function ftuple_color_map_present        

\eatline
\nwendcode{}\nwbegindocs{343}Appends a color flow to the list
\nwenddocs{}\nwbegincode{344}\moddef{nlo data: color map: TBP}\plusendmoddef
  procedure :: append => ftuple_color_map_append  
\nwendcode{}\nwbegincode{345}\moddef{nlo data: procedures}\plusendmoddef
  subroutine ftuple_color_map_append (icm, val)
    class(ftuple_color_map_t), intent(inout), target :: icm
    integer, intent(in) :: val
    type(ftuple_color_map_t), pointer :: current
    select type (icm)
    type is (ftuple_color_map_t)
    if (.not. icm%present (val)) then 
      if (icm%index == 0) then
        nullify(icm%next)
        icm%index = 1
        icm%color_index = val
      else
        current => icm
        do 
          if (associated (current%next)) then
            current => current%next
          else
            allocate (current%next)
            nullify (current%next%next)
            current%next%prev => current 
            current%next%index = current%index + 1
            current%next%color_index = val
            exit
          end if
        end do
      end if
    end if
    end select
  end subroutine ftuple_color_map_append

\eatline
\nwendcode{}\nwbegindocs{346}\nwdocspar
\nwenddocs{}\nwbegincode{347}\moddef{nlo data: color map: TBP}\plusendmoddef
  procedure :: get_n_entries => ftuple_color_map_get_n_entries
\nwendcode{}\nwbegincode{348}\moddef{nlo data: procedures}\plusendmoddef
  function ftuple_color_map_get_n_entries (icm) result(n_entries)
    class(ftuple_color_map_t), intent(in), target :: icm
    integer :: n_entries
    type(ftuple_color_map_t), pointer :: current
    select type (icm)
    type is (ftuple_color_map_t)
    current => icm
    n_entries = 0
    do 
      if (associated (current%next)) then
        current => current%next
      else
        n_entries = current%index
        exit
      end if
    end do
    end select
  end function ftuple_color_map_get_n_entries

\eatline
\nwendcode{}\nwbegindocs{349}Explain
\nwenddocs{}\nwbegincode{350}\moddef{nlo data: color map: TBP}\plusendmoddef
  procedure :: get_index_array => ftuple_color_map_get_index_array  
\nwendcode{}\nwbegincode{351}\moddef{nlo data: procedures}\plusendmoddef
  function ftuple_color_map_get_index_array (icm) result(iarr)
    class(ftuple_color_map_t), intent(in), target :: icm
    integer, dimension(:), allocatable :: iarr
    type(ftuple_color_map_t), pointer :: current
    integer :: n_entries
    integer :: i
    select type (icm)
    type is (ftuple_color_map_t)
    n_entries = icm%get_n_entries ()
    allocate (iarr(n_entries))
    do i = 1, n_entries
      if (i == 1) then
        current => icm
      else
        current => current%next
      end if
      iarr(i) = current%color_index
    end do
    end select
  end function ftuple_color_map_get_index_array

\eatline
\nwendcode{}\nwbegindocs{352}\nwdocspar
\nwenddocs{}\nwbegincode{353}\moddef{nlo data: color map: TBP}\plusendmoddef
  procedure :: get_entry => ftuple_color_map_get_entry  
\nwendcode{}\nwbegincode{354}\moddef{nlo data: procedures}\plusendmoddef
  function ftuple_color_map_get_entry (icm, index) result(entry)
    class(ftuple_color_map_t), intent(in), target :: icm
    integer, intent(in) :: index
    integer :: entry
    type(ftuple_color_map_t), pointer :: current
    integer :: i
    select type (icm)
    type is (ftuple_color_map_t)
    if (index <= icm%get_n_entries ()) then
      do i = 1, icm%get_n_entries ()
        if (i == 1) then
          current => icm
        else
          current => current%next
        end if
        if (i == index) entry = current%color_index
      end do
    else
      entry = 0
    end if
    end select 
  end function ftuple_color_map_get_entry

\eatline
\nwendcode{}\nwbegindocs{355}Explain
\nwenddocs{}\nwbegincode{356}\moddef{nlo data: color map: TBP}\plusendmoddef
  procedure :: create_map => ftuple_color_map_create_map
\nwendcode{}\nwbegincode{357}\moddef{nlo data: procedures}\plusendmoddef
  recursive subroutine ftuple_color_map_create_map (icm, flst, &
       emitter, allreg, color_states_born, color_states_real, p_rad_in)
    class(ftuple_color_map_t), intent(inout) :: icm
    type(flv_structure_t), intent(in) :: flst
    integer, intent(in) :: emitter
    type(ftuple_t), intent(in), dimension(:) :: allreg
    integer, intent(in), dimension(:,:,:) :: color_states_born
    integer, intent(in), dimension(:,:,:) :: color_states_real
    integer, intent(in), optional :: p_rad_in
    integer :: nreg, region
    integer :: p1, p2, p_rad
    integer :: flv_em, flv_rad
    integer :: n_col_real, n_col_born
    integer, dimension(2) :: col_em, col_rad
    integer :: i
    !!! splitting type: 1 - q -> qg
    !!!                 2 - g -> qq
    !!!                 3 - g -> gg
    integer :: splitting_type_flv, splitting_type_col
    nreg = size (allreg)
    n_col_real = size (color_states_real (1,1,:))
    n_col_born = size (color_states_born (1,1,:))
    do region = 1, nreg
      call allreg(region)%get (p1, p2)
      if (p1 == emitter .or. p2 == emitter .or. present (p_rad_in)) then
        if (.not. present (p_rad_in)) then
          if (p1 == emitter) then
            p_rad = p2
          else
            p_rad = p1
          end if
        else
          p_rad = p_rad_in
        end if
        if (emitter /= 0) then
          flv_em = flst%flst (emitter)
        else
          call icm%create_map &
               (flst, 1, allreg, color_states_born, color_states_real, p_rad)
          call icm%create_map &
               (flst, 2, allreg, color_states_born, color_states_real, p_rad)
        end if
        flv_rad = flst%flst (p_rad)
        if (is_quark (abs(flv_em)) .and. is_gluon (flv_rad)) then
           splitting_type_flv = 1
        else if (is_quark (abs(flv_em)) .and. flv_em + flv_rad == 0) then
           splitting_type_flv = 2
        else if (is_gluon (flv_em) .and. is_gluon (flv_rad)) then
           splitting_type_flv = 3
        else 
          splitting_type_flv = 0
        end if 
        do i = 1, n_col_real
           col_em = color_states_real(:,emitter,i)
           col_rad = color_states_real(:,p_rad,i)
          if (is_color_singlet (col_em(1), col_em(2)) &
              .and. (is_color_doublet (col_rad(1), col_rad(2)) &
              .or. is_color_ghost (col_rad(1), col_rad(2)))) then
            splitting_type_col = 1
          else if (is_color_singlet (col_em(1), col_em(2)) .and. &
                   is_color_singlet (col_rad(1), col_rad(2))) then
            splitting_type_col = 2
          else if (is_color_doublet (col_em(1), col_em(2)) .and. &
                   is_color_doublet (col_rad(1), col_rad(2))) then
            splitting_type_col = 3
          else
            splitting_type_col = 0
          end if
          if (splitting_type_flv == splitting_type_col .and. &
              splitting_type_flv /= 0) then
            call icm%append (i)
          end if
        end do    
      end if  
    end do
  contains
    function is_color_singlet (c1, c2) result (singlet)
      integer, intent(in) :: c1, c2
      logical :: singlet
      singlet = (c1 == 0 .and. c2 /= 0) .or. (c1 /= 0 .and. c2 == 0) 
    end function is_color_singlet
    function is_color_doublet (c1, c2) result (doublet)
      integer, intent(in) :: c1, c2
      logical :: doublet
      doublet = c1 /= 0 .and. c2 /= 0
    end function is_color_doublet
    function is_color_ghost (c1, c2) result (ghost)
      integer, intent(in) :: c1, c2
      logical :: ghost
      ghost = c1 == 0 .and. c2 == 0
    end function is_color_ghost
  end subroutine ftuple_color_map_create_map

\eatline
\nwendcode{}\nwbegindocs{358}This data type contains color information, necessary for both soft
and virtual counterterms. 
\nwenddocs{}\nwbegincode{359}\moddef{nlo data: types}\plusendmoddef
  type color_data_t
    type(ftuple_color_map_t), dimension(:), allocatable :: icm
    integer, dimension(:,:,:), allocatable :: col_state_born, col_state_real
    logical, dimension(:,:), allocatable :: ghost_flag_born, ghost_flag_real
    integer :: n_col_born, n_col_real
    type(color_t), dimension(:,:), allocatable :: color_real, color_born
    integer, dimension(:), allocatable :: col_born
    complex(default), dimension(:), allocatable :: color_factors_born
    integer, dimension(:,:), allocatable :: cf_index_real
    real(default), dimension(:,:,:), allocatable :: beta_ij
    logical :: color_is_conserved
  contains
  \LA{}nlo data: color data: TBP\RA{}
  end type color_data_t

\eatline
\nwendcode{}\nwbegindocs{360}\nwdocspar
\nwenddocs{}\nwbegincode{361}\moddef{nlo data: color data: TBP}\endmoddef
  procedure :: init => color_data_init
\nwendcode{}\nwbegincode{362}\moddef{nlo data: procedures}\plusendmoddef
  subroutine color_data_init (color_data, reg_data, prc_constants)
    class(color_data_t), intent(inout) :: color_data
    type(region_data_t), intent(inout) :: reg_data
    type(process_constants_t), intent(inout), dimension(2) :: prc_constants
    integer :: nlegs_born, nlegs_real
    integer :: i
    nlegs_born = reg_data%nlegs_born
    nlegs_real = reg_data%nlegs_real
    call prc_constants(1)%get_col_state (color_data%col_state_born)
    call prc_constants(2)%get_col_state (color_data%col_state_real)
    call prc_constants(2)%get_cf_index (color_data%cf_index_real)
    call prc_constants(1)%get_color_factors (color_data%color_factors_born)
    color_data%n_col_born = size (color_data%col_state_born(1,1,:))
    color_data%n_col_real = size (color_data%col_state_real(1,1,:))
    color_data%ghost_flag_born = prc_constants(1)%get_ghost_flag ()
    color_data%ghost_flag_real = prc_constants(2)%get_ghost_flag ()
    allocate (color_data%color_real (nlegs_real, color_data%n_col_real))
    allocate (color_data%icm (reg_data%n_regions))
    do i = 1, color_data%n_col_real
      call color_init_from_array (color_data%color_real (:,i), &
           color_data%col_state_real (:,:,i), &
           color_data%ghost_flag_real (:,i))
    end do
    do i = 1, size(reg_data%regions)
      call color_data%icm(i)%init
      associate (region => reg_data%regions(i))
        call color_data%icm(i)%create_map (region%flst_real, region%emitter, &
             region%flst_allreg, color_data%col_state_born, &
             color_data%col_state_real)
      end associate
    end do
    call color_data%init_betaij (reg_data)
  end subroutine color_data_init

\eatline
\nwendcode{}\nwbegindocs{363}Allocate and compute $\beta_{ij}$:
\nwenddocs{}\nwbegincode{364}\moddef{nlo data: color data: TBP}\plusendmoddef
  procedure :: init_betaij => color_data_init_betaij
\nwendcode{}\nwbegincode{365}\moddef{nlo data: procedures}\plusendmoddef
  subroutine color_data_init_betaij (color_data, reg_data)
    class(color_data_t), intent(inout) :: color_data
    type(region_data_t), intent(inout) :: reg_data
    integer :: i
    allocate (color_data%beta_ij (reg_data%nlegs_born, &
         reg_data%nlegs_born, reg_data%n_flv_born))
    do i = 1, reg_data%n_flv_born
      call color_data%fill_betaij_matrix (reg_data%nlegs_born, i, &
           reg_data%regions(1)%flst_real, reg_data)
    end do
  end subroutine color_data_init_betaij

\eatline
\nwendcode{}\nwbegindocs{366}Actual computation of $\beta_{ij}$.
\nwenddocs{}\nwbegincode{367}\moddef{nlo data: color data: TBP}\plusendmoddef
  procedure :: fill_betaij_matrix => color_data_fill_betaij_matrix
\nwendcode{}\nwbegincode{368}\moddef{nlo data: procedures}\plusendmoddef
  subroutine color_data_fill_betaij_matrix &
       (color_data, n_legs, uborn_index, flst_real, reg_data)
    class(color_data_t), intent(inout) :: color_data
    integer, intent(in) :: n_legs, uborn_index
    type(flv_structure_t), intent(in) :: flst_real
    type(region_data_t), intent(inout) :: reg_data
    integer :: em1, em2
    associate (flv_born => reg_data%flv_born (uborn_index))
    do em1 = 1, n_legs
      do em2 = 1, n_legs
        if (flv_born%colored(em1) .and. flv_born%colored(em2)) then
          if (em1 < em2) then
             color_data%beta_ij (em1, em2, uborn_index) &
                = color_data%compute_bij &
                     (reg_data, uborn_index, flst_real, em1, em2)
          else if (em1 > em2) then
             !!! B_ij is symmetric
             color_data%beta_ij (em1, em2, uborn_index) = &
                  color_data%beta_ij (em2, em1, uborn_index)
          else
            if (is_quark (abs (flv_born%flst (em1)))) then
              color_data%beta_ij (em1, em2, uborn_index) = -cf
            else
              color_data%beta_ij (em1, em2, uborn_index) = -ca
            end if
          end if
        else
          color_data%beta_ij (em1, em2, uborn_index) = 0.0
        end if
      end do
    end do
    end associate
    call check_color_conservation (color_data%beta_ij (:,:,uborn_index), &
         n_legs, color_data%color_is_conserved)
  contains
    subroutine check_color_conservation (bij_matrix, n_legs, success)
      real(default), intent(in), dimension(:,:) :: bij_matrix
      integer, intent(in) :: n_legs
      logical, intent(out) :: success
      logical, dimension(:), allocatable :: check
      integer :: i, j
      real(default) :: bcheck
      real(default), parameter :: tol = 0.0001_default
      allocate (check (n_legs))
      do i = 1, n_legs
        bcheck = 0.0
        do j = 1, n_legs
          if (i /= j) bcheck = bcheck + bij_matrix (i, j)
        end do
        if (is_quark (abs(flst_real%flst (i))) .or. &
            is_gluon (flst_real%flst (i))) then
          if (is_quark (abs(flst_real%flst (i))) .and. &
               (bcheck - cf) < tol) then
             check (i) = .true.
          else if (is_gluon (flst_real%flst (i)) .and. &
               (bcheck - ca) < tol) then
            check (i) = .true.
          else
            check (i) = .false.
          end if
        else 
          if (bcheck < tol) then
            check (i) = .true.
          else
            check (i) = .false.
          end if
        end if
      end do
      if (.not. all (check)) then
        success = .false.
        ! call msg_fatal ("Color conservation violated!") 
      else
        success = .true.
      end if
    end subroutine check_color_conservation
  end subroutine color_data_fill_betaij_matrix

\eatline
\nwendcode{}\nwbegindocs{369}Explain
\nwenddocs{}\nwbegincode{370}\moddef{nlo data: color data: TBP}\plusendmoddef
  procedure :: compute_bij => color_data_compute_bij
\nwendcode{}\nwbegincode{371}\moddef{nlo data: procedures}\plusendmoddef
  function color_data_compute_bij &
       (color_data, reg_data, uborn_index, flst_real, em1, em2) result (bij)
    class(color_data_t), intent(inout) :: color_data
    type(region_data_t), intent(inout) :: reg_data
    integer, intent(in) :: uborn_index
    type(flv_structure_t), intent(in) :: flst_real
    integer, intent(in) :: em1, em2
    real(default) :: bij
    logical, dimension(:,:), allocatable :: cf_present
    type(singular_region_t), dimension(2,100) :: reg
    integer ::  i, j, k, l
    type(ftuple_color_map_t) :: icm1, icm2
    integer :: i1, i2
    real(default) :: color_factor, color_factor_born
    integer, dimension(2) :: i_reg
    logical , dimension(2) :: found
    integer, dimension(2,100) :: map_em_col_tmp
    integer, dimension(:), allocatable :: map_em_col1, map_em_col2
    integer, dimension(2) :: col1, col2
    integer, dimension(:), allocatable :: iarray1, iarray2
    integer, dimension(:), allocatable :: iisec1, iisec2
    integer :: sign
    color_factor = 0.0; color_factor_born = 0.0
    found = .false.
    !!! Include distinction between Born flavors
    do i = 1, size (color_data%color_factors_born)
      color_factor_born = color_factor_born + color_data%color_factors_born (i)
    end do
    i1 = 1
    i2 = 1
    !!! Catch case em = 0
    if (em1 == 0 .or. em2 == 0) then
       !!! What to do?
       bij = 0.0
    else
       do i = 1, color_data%n_col_real
          col1 = color_data%col_state_real (:, em1, i)
          col2 = color_data%col_state_real (:, reg_data%nlegs_real, i)
          if (share_line (col1, col2)) then
             map_em_col_tmp(1,i1) = i
             i1 = i1+1
          end if
          col1 = color_data%col_state_real (:, em2, i)
          if (share_line (col1, col2)) then 
             map_em_col_tmp(2,i2) = i
             i2 = i2 + 1
          end if
       end do
       allocate (map_em_col1 (i1), map_em_col2 (i2))
       map_em_col1 = map_em_col_tmp (1,1:i1-1)
       map_em_col2 = map_em_col_tmp (2,1:i2-1)

       i_reg = 1
    
       do i = 1, reg_data%n_regions
           if (uborn_index == reg_data%regions(i)%uborn_index) then
             if (em1 == reg_data%regions(i)%emitter .or. &
                 (em1 <= 2 .and. reg_data%regions(i)%emitter == 0)) then
               reg(1,i_reg(1)) = reg_data%regions(i)
               i_reg(1) = i_reg(1)+1
               found(1) = .true.
             end if
             if (em2 == reg_data%regions(i)%emitter .or. &
                 (em2 <= 2 .and. reg_data%regions(i)%emitter == 0)) then
               reg(2,i_reg(2)) = reg_data%regions(i)
               i_reg(2) = i_reg(2)+1
               found(2) = .true.
             end if
           end if
       end do
       if (.not. (found(1).and.found(2))) then
         bij = 0
         return
       end if

       do i = 1, i_reg(1)-1
         do j = 1, i_reg(2)-1
           icm1 = color_data%icm (reg(1,i)%alr)
           icm2 = color_data%icm (reg(2,j)%alr)
       
           iarray1 = icm1%get_index_array ()
           iarray2 = icm2%get_index_array ()
       
           iisec1 = pack (iarray1, [ (any(iarray1(i) == map_em_col1), &
                i = 1, size(iarray1)) ])
           iisec2 = pack (iarray2, [ (any(iarray2(i) == map_em_col2), &
                i = 1, size(iarray2)) ])
       
           cf_present = color_index_present (color_data%cf_index_real)
       
           do k = 1, size (iisec1)
             do l = 1, size (iisec2)
               i1 = iisec1(k)
               i2 = iisec2(l)
               if (cf_present (i1, i2)) then
                 if (is_gluon (flst_real%flst (em1)) .or. &
                     is_gluon (flst_real%flst (em2))) then
                   sign = get_sign (color_data%col_state_real (:,:,i1)) * &
                        get_sign (color_data%col_state_real (:,:,i2))
                 else
                   sign = 1
                 end if
                   color_factor = color_factor + sign*compute_color_factor &
                         (color_data%color_real(:,i1), &
                          color_data%color_real(:,i2))
               end if
             end do
           end do       
         end do
       end do
       !!! The real color factor always differs from the Born one 
       !!! by one vertex factor. Thus, apply the factor 1/2  
       bij = color_factor / (2 * color_factor_born) 
    end if

  contains
    function share_line (col1, col2) result (share)
      integer, intent(in), dimension(2) :: col1, col2
      logical :: share
      logical :: id1, id2, id3
      id1 = (abs(col1(1)) == abs(col2(1)) .and. col1(1) /= 0) .or. &
            (abs(col1(2)) == abs(col2(2)) .and. col1(2) /= 0)
      id2 = (abs(col1(1)) == abs(col2(2)) .and. col1(1) /= 0) .or. &
            (abs(col1(2)) == abs(col2(1)) .and. col1(2) /= 0)
      id3 = col2(1) == 0 .and. col2(2) == 0
      if (id1 .or. id2 .or. id3) then
        share = .true.
      else
        share = .false.
      end if
    end function share_line

    function get_sign (col) result (sign)
      integer, intent(in), dimension(:,:) :: col
      integer :: sign
      integer, dimension(:), allocatable :: iref, iperm
      integer :: iref1, iperm1
      integer :: n, i, i_first, j
      integer :: i1, i2
      integer :: p1, p2
      p1 = 2; p2 = 2
      iref1 = 0; iperm1 = 0; i_first = 0
      do i = 1, size(col(1,:))
        if (.not. all (col(:,i) == 0)) then 
          if (col(1,i) == 0) then
            i1 = col(2,i)
            iref1 = i; iperm1 = i
            i_first = i
          else
            i1 = col(1,i)
            iref1 = i; iperm1 = i
            i_first = i
          end if
          exit
        end if
      end do
      if (iref1 == 0 .or. iperm1 == 0 .or. i_first == 0) &
         call msg_fatal ("Invalid color structure")
      n = size(col(1,:)) - i_first + 1
      allocate (iref(n), iperm(n))
      iref(1) = iref1; iperm(1) = iperm1
      do i = i_first+1, size(col(1,:))
        if (all (col(:,i) == 0)) cycle
        if (i == size(col(1,:))) then
          iref(p1) = i_first + 1
        else
          iref(p1) = i + 1
          p1 = p1 + 1
        end if
        do j = i_first+1, size(col(1,:))
          if (col(1,j) == -i1) then
            i1 = col(2,j)
            iperm(p2) = j
            p2 = p2 + 1
            exit
          else if (col(2,j) == -i1) then
            i1 = col(1,j)
            iperm(p2) = j
            p2 = p2 + 1
            exit
          end if
        end do
      end do
      sign = 1
      do i = 1, n
        if (iperm(i) == iref(i)) then
          cycle
        else
          do j = i+1, n
            if (iperm(j) == iref(i)) then
              i1 = j
              exit
            end if
          end do
          i2 = iperm(i)
          iperm(i) = iperm(i1) 
          iperm(i1) = i2
          sign = -sign
        end if
      end do 
    end function get_sign
 
    function color_index_present (cf_index) result (cf_present) 
      integer, intent(in), dimension(:,:), allocatable :: cf_index
      logical, dimension(:,:), allocatable :: cf_present
      integer :: n_col
      integer :: c, i1, i2
      n_col = size (cf_index(1,:))
      allocate (cf_present (n_col, n_col))
      cf_present = .false.
      do c = 1, n_col
        i1 = cf_index (1, c)
        i2 = cf_index (2, c)
        cf_present (i1, i2) = .true.
        if (i1 /= i2) cf_present(i2, i1) = .true.
      end do
    end function color_index_present
  end function color_data_compute_bij

\eatline
\nwendcode{}\nwbegindocs{372}\nwdocspar
\nwenddocs{}\nwbegincode{373}\moddef{nlo data: color data: TBP}\plusendmoddef
  procedure :: write => color_data_write
\nwendcode{}\nwbegincode{374}\moddef{nlo data: procedures}\plusendmoddef
  subroutine color_data_write (color_data, unit)
    class(color_data_t), intent(in) :: color_data
    integer, intent(in), optional :: unit
    integer :: u, i, i1, i2
    integer :: n_legs
    u = given_output_unit (unit); if (u < 0) return
    n_legs = size (color_data%beta_ij, dim=2)
    write (u, "(1x,A)") "Color information: "
    write (u, "(1x,A,1x,I1)") "Number of Born color states: ", &
         color_data%n_col_born
    write (u, "(1x,A,1x,I1)") "Number of real color states: ", &
         color_data%n_col_real
    write (u, "(1x,A)") "Color correlation: "
    do i = 1, size (color_data%beta_ij, dim=3)
      write (u, "(1x,A,1x,I1)") "State nr. ", i
      write (u, "(1x,A)") "-------------"
      write (u, "(1x,A,1x,A,1x,A)") "i1", "i2", "color factor"
      do i1 = 1, n_legs
        do i2 = 1, i1
          write (u, "(1x,I1,1x,I1,1x,F5.2)") &
               i1, i2, color_data%beta_ij (i1,i2,i)
        end do
      end do
      write (u, "(1x,A)") "========================================"
    end do
    if (color_data%color_is_conserved) then
      write (u, "(1x,A)") "Color is conserved."
    else
      write (u, "(1x,A)") "Fatal error: Color conversation is violated."
    end if
  end subroutine color_data_write

\eatline
\nwendcode{}\nwbegindocs{375}\nwdocspar
\nwenddocs{}\nwbegincode{376}\moddef{nlo data: nlo data: TBP}\endmoddef
  procedure :: compute_sqme_real_fin => nlo_data_compute_sqme_real_fin
\nwendcode{}\nwbegincode{377}\moddef{nlo data: procedures}\plusendmoddef
  function nlo_data_compute_sqme_real_fin (nlo_data, weight, p_real) result (sqme_fin)
    class(nlo_data_t), intent(inout) :: nlo_data
    real(default), intent(in) :: weight
    type(vector4_t), intent(inout), dimension(:), allocatable :: p_real
    type(vector4_t), dimension(:), allocatable :: p_born
    real(default) :: sqme_fin
    integer :: emitter
    if (.not. nlo_data%alpha_s_born_set) &
      call msg_fatal ("Strong coupling not set for real calculation - abort")
    emitter = nlo_data%get_active_emitter ()
    p_born = interaction_get_momenta (nlo_data%int_born)
    call nlo_data%real_terms%set_momenta (p_born, p_real)
    call nlo_data%real_terms%set_fks_kinematics (nlo_data%real_kinematics)
    sqme_fin = nlo_data%real_terms%compute (emitter, &
                                            nlo_data%alpha_s_born)
    sqme_fin = sqme_fin * weight
  end function nlo_data_compute_sqme_real_fin

\eatline
\nwendcode{}\nwbegindocs{378}Check if there are massive emitters. Since the mass-structure of all
underlying Born configurations have to be the same, we just use the first
one to determine this.
\nwenddocs{}\nwbegincode{379}\moddef{nlo data: nlo data: TBP}\plusendmoddef
  procedure :: has_massive_emitter => nlo_data_has_massive_emitter
\nwendcode{}\nwbegincode{380}\moddef{nlo data: procedures}\plusendmoddef
  function nlo_data_has_massive_emitter (nlo_data) result (val)
    class(nlo_data_t), intent(inout) :: nlo_data
    logical :: val
    integer :: n_tot, i
    val = .false.
    n_tot = nlo_data%n_in + nlo_data%n_out_born
    do i = nlo_data%n_in+1, n_tot
       if (any (i == nlo_data%reg_data%emitters)) &
          val = val .or. nlo_data%reg_data%flv_born(1)%massive(i)
    end do
  end function nlo_data_has_massive_emitter

\eatline
\nwendcode{}\nwbegindocs{381}\nwdocspar
\subsection{Putting it together}
This data type governs the whole calculation. It contains information
about color, spin and flavor as well as the information about the Born
process. 
\nwenddocs{}\nwbegincode{382}\moddef{nlo data: public}\plusendmoddef
  public :: nlo_data_t
\nwendcode{}\nwbegincode{383}\moddef{nlo data: types}\plusendmoddef
  type :: nlo_data_t
    type(string_t) :: nlo_type
    type(region_data_t) :: reg_data
    integer :: n_alr
    integer :: n_in
    integer :: n_out_born
    integer :: n_out_real
    integer :: n_allowed_born
    integer :: n_flv_born 
    integer :: n_flv_real
    integer :: active_emitter
    complex(default), dimension(:), allocatable :: amp_born
    type(color_data_t) :: color_data
    type(real_kinematics_t) :: real_kinematics
    type(virtual_t) :: virtual_terms
    type(real_subtraction_t) :: real_terms
    integer, dimension(:,:), allocatable :: flv_state_born
    integer, dimension(:,:), allocatable :: flv_state_real
    integer, dimension(:), allocatable :: flv_born
    integer, dimension(:), allocatable :: hel_born
    integer, dimension(:), allocatable :: col_born
    real(default) :: alpha_s_born
    logical :: alpha_s_born_set
    real(default), dimension(:), allocatable :: sqme_born
    real(default), dimension(:), allocatable :: sqme_real_non_sub
    real(default), dimension(:,:,:), allocatable :: sqme_born_cc
    complex(default) :: me_sc
    real(default), public :: sqme_real
    real(default), public :: sqme_virt
    type(interaction_t), public :: int_born
    type(kinematics_counter_t), public :: counter
    logical, public :: counter_exists = .false.
    logical :: use_internal_color_correlations = .true.
    logical :: use_internal_spin_correlations = .false.
  contains
  \LA{}nlo data: nlo data: TBP\RA{}
  end type nlo_data_t

\eatline
\nwendcode{}\nwbegindocs{384}\nwdocspar
\nwenddocs{}\nwbegincode{385}\moddef{nlo data: public}\plusendmoddef
  public :: nlo_pointer_t
\nwendcode{}\nwbegincode{386}\moddef{nlo data: types}\plusendmoddef
  type :: nlo_pointer_t
    type(nlo_data_t), public, pointer :: nlo_data => null ()
  end type nlo_pointer_t

\eatline
\nwendcode{}\nwbegindocs{387}\nwdocspar
\nwenddocs{}\nwbegincode{388}\moddef{nlo data: nlo data: TBP}\plusendmoddef
  procedure :: init => nlo_data_init
\nwendcode{}\nwbegincode{389}\moddef{nlo data: procedures}\plusendmoddef
  subroutine nlo_data_init (nlo_data, nlo_type, prc_constants, template, model)
     class(nlo_data_t), intent(inout) :: nlo_data
     type(string_t), intent(in) :: nlo_type
     type(process_constants_t), intent(inout), dimension(2) :: prc_constants
     type(fks_template_t), intent(in) :: template
     type(model_data_t), intent(in) :: model
     integer :: i
     nlo_data%nlo_type = nlo_type
     select case (char (nlo_type))
     case ('Real')
       nlo_data%flv_state_born = prc_constants(1)%get_flv_state ()
       nlo_data%flv_state_real = prc_constants(2)%get_flv_state ()
       call nlo_data%reg_data%init (model, &
            nlo_data%flv_state_born, &
            nlo_data%flv_state_real, &
            template%mapping_type)
       select type (mapping => nlo_data%reg_data%fks_mapping)
       type is (fks_mapping_default_t)
         call mapping%set_parameter (template%fks_dij_exp1, template%fks_dij_exp2)
       end select
       nlo_data%n_flv_born = nlo_data%reg_data%n_flv_born
       nlo_data%n_flv_real = nlo_data%reg_data%n_flv_real
       allocate (nlo_data%sqme_born (nlo_data%n_flv_born))
       allocate (nlo_data%sqme_real_non_sub (nlo_data%n_flv_real))
       nlo_data%sqme_born = 0.0
       nlo_data%n_in = prc_constants(2)%n_in
       nlo_data%n_out_born = prc_constants(1)%n_out
       nlo_data%n_out_real = prc_constants(2)%n_out
       allocate (nlo_data%sqme_born_cc (nlo_data%n_in + nlo_data%n_out_born, &
                                        nlo_data%n_in + nlo_data%n_out_born, &
                                        nlo_data%n_flv_born))
       if (nlo_data%use_internal_color_correlations) &
          call nlo_data%color_data%init (nlo_data%reg_data, prc_constants)
       nlo_data%alpha_s_born_set = .false.
       call nlo_data%init_real_kinematics
       call nlo_data%real_terms%init (nlo_data%reg_data, &
                                      nlo_data%n_in + nlo_data%n_out_born, &
                                      nlo_data%n_in + nlo_data%n_out_born)
       nlo_data%counter_exists = template%count_kinematics
       if (nlo_data%counter_exists) call nlo_data%counter%init(20)
     end select
   end subroutine nlo_data_init

\eatline
\nwendcode{}\nwbegindocs{390}\nwdocspar
\nwenddocs{}\nwbegincode{391}\moddef{nlo data: nlo data: TBP}\plusendmoddef
  procedure :: copy_matrix_elements => nlo_data_copy_matrix_elements
\nwendcode{}\nwbegincode{392}\moddef{nlo data: procedures}\plusendmoddef
  subroutine nlo_data_copy_matrix_elements (nlo_data)
    class(nlo_data_t), intent(inout) :: nlo_data
     nlo_data%real_terms%sqme_born = nlo_data%sqme_born
     nlo_data%real_terms%sqme_born_cc = nlo_data%sqme_born_cc
     nlo_data%real_terms%sqme_real_bare = nlo_data%sqme_real_non_sub
     nlo_data%real_terms%me_sc = nlo_data%me_sc
  end subroutine nlo_data_copy_matrix_elements

\eatline
\nwendcode{}\nwbegindocs{393}\nwdocspar
\nwenddocs{}\nwbegincode{394}\moddef{nlo data: nlo data: TBP}\plusendmoddef
  procedure :: write => nlo_data_write
\nwendcode{}\nwbegincode{395}\moddef{nlo data: procedures}\plusendmoddef
  !!! Incomplete
  subroutine nlo_data_write (nlo_data, unit)
    class(nlo_data_t), intent(in) :: nlo_data
    integer, intent(in), optional :: unit
    integer :: i, u
    u = given_output_unit (unit); if (u < 0) return
    write (u, "(1x,A,I0)") "n_alr          = ", nlo_data%n_alr
    write (u, "(1x,A,I0)") "n_in           = ", nlo_data%n_in
    write (u, "(1x,A,I0)") "n_out_born     = ", nlo_data%n_out_born
    write (u, "(1x,A,I0)") "n_out_real     = ", nlo_data%n_out_real
    write (u, "(1x,A,I0)") "n_allowed_born = ", nlo_data%n_allowed_born
  end subroutine nlo_data_write

\eatline
\nwendcode{}\nwbegindocs{396}\nwdocspar
\nwenddocs{}\nwbegincode{397}\moddef{nlo data: nlo data: TBP}\plusendmoddef
  procedure :: init_born_amps => nlo_data_init_born_amps
\nwendcode{}\nwbegincode{398}\moddef{nlo data: procedures}\plusendmoddef
  subroutine nlo_data_init_born_amps (nlo_data, n, internal_correlations)
    class(nlo_data_t), intent(inout) :: nlo_data
    integer, intent(in) :: n
    logical, intent(in) :: internal_correlations
    nlo_data%n_allowed_born = n
    allocate (nlo_data%amp_born (n))
  end subroutine nlo_data_init_born_amps

\eatline
\nwendcode{}\nwbegindocs{399}\nwdocspar
\nwenddocs{}\nwbegincode{400}\moddef{nlo data: nlo data: TBP}\plusendmoddef
  procedure :: set_internal_procedures => nlo_data_set_internal_procedures
\nwendcode{}\nwbegincode{401}\moddef{nlo data: procedures}\plusendmoddef
  subroutine nlo_data_set_internal_procedures (nlo_data, flag_color, flag_spin)
    class(nlo_data_t), intent(inout) :: nlo_data
    logical, intent(in) :: flag_color, flag_spin
    nlo_data%use_internal_color_correlations = flag_color
    nlo_data%real_terms%sub_soft%use_internal_color_correlations = flag_color
    nlo_data%virtual_terms%use_internal_color_correlations = flag_color
    nlo_data%use_internal_spin_correlations = flag_spin
  end subroutine nlo_data_set_internal_procedures

\eatline
\nwendcode{}\nwbegindocs{402}\nwdocspar
\nwenddocs{}\nwbegincode{403}\moddef{nlo data: nlo data: TBP}\plusendmoddef
  procedure :: init_virtual => nlo_data_init_virtual
\nwendcode{}\nwbegincode{404}\moddef{nlo data: procedures}\plusendmoddef
  subroutine nlo_data_init_virtual (nlo_data)
    class(nlo_data_t), intent(inout) :: nlo_data
    call nlo_data%virtual_terms%init (nlo_data%flv_state_born)
  end subroutine nlo_data_init_virtual

\eatline
\nwendcode{}\nwbegindocs{405}\nwdocspar
\nwenddocs{}\nwbegincode{406}\moddef{nlo data: nlo data: TBP}\plusendmoddef
  procedure :: set_nlo_type => nlo_data_set_nlo_type
\nwendcode{}\nwbegincode{407}\moddef{nlo data: procedures}\plusendmoddef
  subroutine nlo_data_set_nlo_type (nlo_data, nlo_type)
   class(nlo_data_t), intent(inout) :: nlo_data
   type(string_t), intent(in) :: nlo_type
   nlo_data%nlo_type = nlo_type
  end subroutine nlo_data_set_nlo_type

\eatline
\nwendcode{}\nwbegindocs{408}\nwdocspar
\nwenddocs{}\nwbegincode{409}\moddef{nlo data: nlo data: TBP}\plusendmoddef
  procedure :: get_nlo_type => nlo_data_get_nlo_type 
\nwendcode{}\nwbegincode{410}\moddef{nlo data: procedures}\plusendmoddef
  pure function nlo_data_get_nlo_type (nlo_data) result(nlo_type)
    class(nlo_data_t), intent(in) :: nlo_data
    type(string_t) :: nlo_type
    nlo_type = nlo_data%nlo_type
  end function nlo_data_get_nlo_type

\eatline
\nwendcode{}\nwbegindocs{411}\nwdocspar
\nwenddocs{}\nwbegincode{412}\moddef{nlo data: nlo data: TBP}\plusendmoddef
  procedure :: get_emitters => nlo_data_get_emitters
\nwendcode{}\nwbegincode{413}\moddef{nlo data: procedures}\plusendmoddef
  function nlo_data_get_emitters (nlo_data) result(emitters)
    class(nlo_data_t), intent(inout) :: nlo_data
    integer, dimension(:), allocatable :: emitters
    emitters = nlo_data%reg_data%get_emitters ()
  end function nlo_data_get_emitters

\eatline
\nwendcode{}\nwbegindocs{414}\nwdocspar
\nwenddocs{}\nwbegincode{415}\moddef{nlo data: nlo data: TBP}\plusendmoddef
  procedure :: set_active_emitter => nlo_data_set_active_emitter
\nwendcode{}\nwbegincode{416}\moddef{nlo data: procedures}\plusendmoddef
  subroutine nlo_data_set_active_emitter (nlo_data, emitter)
    class(nlo_data_t), intent(inout) :: nlo_data
    integer, intent(in) :: emitter
    nlo_data%active_emitter = emitter
  end subroutine nlo_data_set_active_emitter

\eatline
\nwendcode{}\nwbegindocs{417}\nwdocspar
\nwenddocs{}\nwbegincode{418}\moddef{nlo data: nlo data: TBP}\plusendmoddef
  procedure :: get_active_emitter => nlo_data_get_active_emitter
\nwendcode{}\nwbegincode{419}\moddef{nlo data: procedures}\plusendmoddef
  function nlo_data_get_active_emitter (nlo_data) result(emitter)
    class(nlo_data_t), intent(inout) :: nlo_data
    integer :: emitter
    emitter = nlo_data%active_emitter
  end function nlo_data_get_active_emitter

\eatline
\nwendcode{}\nwbegindocs{420}\nwdocspar
\nwenddocs{}\nwbegincode{421}\moddef{nlo data: nlo data: TBP}\plusendmoddef
  procedure :: set_flv_born => nlo_data_set_flv_born
\nwendcode{}\nwbegincode{422}\moddef{nlo data: procedures}\plusendmoddef
  subroutine nlo_data_set_flv_born (nlo_data, flv_in)
    class(nlo_data_t), intent(inout) :: nlo_data
    integer, intent(in), dimension(:), allocatable :: flv_in
    allocate (nlo_data%flv_born (size (flv_in)))
    nlo_data%flv_born = flv_in
  end subroutine nlo_data_set_flv_born 

\eatline
\nwendcode{}\nwbegindocs{423}\nwdocspar
\nwenddocs{}\nwbegincode{424}\moddef{nlo data: nlo data: TBP}\plusendmoddef
  procedure :: set_hel_born => nlo_data_set_hel_born
\nwendcode{}\nwbegincode{425}\moddef{nlo data: procedures}\plusendmoddef
  subroutine nlo_data_set_hel_born (nlo_data, hel_in)
    class(nlo_data_t), intent(inout) :: nlo_data
    integer, intent(in), dimension(:), allocatable :: hel_in
    allocate (nlo_data%hel_born (size (hel_in)))
    nlo_data%hel_born = hel_in
  end subroutine nlo_data_set_hel_born 

\eatline
\nwendcode{}\nwbegindocs{426}\nwdocspar
\nwenddocs{}\nwbegincode{427}\moddef{nlo data: nlo data: TBP}\plusendmoddef
  procedure :: set_col_born => nlo_data_set_col_born
\nwendcode{}\nwbegincode{428}\moddef{nlo data: procedures}\plusendmoddef
  subroutine nlo_data_set_col_born (nlo_data, col_in)
    class(nlo_data_t), intent(inout) :: nlo_data
    integer, intent(in), dimension(:), allocatable :: col_in
    allocate (nlo_data%col_born (size (col_in)))
    nlo_data%col_born = col_in
  end subroutine nlo_data_set_col_born

\eatline
\nwendcode{}\nwbegindocs{429}\nwdocspar
\nwenddocs{}\nwbegincode{430}\moddef{nlo data: nlo data: TBP}\plusendmoddef
  procedure :: set_alpha_s_born => nlo_data_set_alpha_s_born
\nwendcode{}\nwbegincode{431}\moddef{nlo data: procedures}\plusendmoddef
  subroutine nlo_data_set_alpha_s_born (nlo_data, as_born)
    class (nlo_data_t), intent(inout) :: nlo_data
    real(default), intent(in) :: as_born
    nlo_data%alpha_s_born = as_born
    nlo_data%alpha_s_born_set = .true.
  end subroutine nlo_data_set_alpha_s_born

\eatline
\nwendcode{}\nwbegindocs{432}\nwdocspar
\nwenddocs{}\nwbegincode{433}\moddef{nlo data: nlo data: TBP}\plusendmoddef
  procedure :: init_real_kinematics => nlo_data_init_real_kinematics
\nwendcode{}\nwbegincode{434}\moddef{nlo data: procedures}\plusendmoddef
  subroutine nlo_data_init_real_kinematics (nlo_data)
    class(nlo_data_t), intent(inout) :: nlo_data
    integer :: n_tot
    n_tot = nlo_data%n_in + nlo_data%n_out_born
    allocate (nlo_data%real_kinematics%xi_max (n_tot))
    allocate (nlo_data%real_kinematics%y (n_tot))
    allocate (nlo_data%real_kinematics%y_soft (n_tot))
    allocate (nlo_data%real_kinematics%jac_rand (n_tot))
    nlo_data%real_kinematics%xi_tilde = 0
    nlo_data%real_kinematics%y = 0
    nlo_data%real_kinematics%xi_max = 0
    nlo_data%real_kinematics%phi = 0
  end subroutine nlo_data_init_real_kinematics

\eatline
\nwendcode{}\nwbegindocs{435}\nwdocspar
\nwenddocs{}\nwbegincode{436}\moddef{nlo data: nlo data: TBP}\plusendmoddef
  procedure :: set_real_kinematics => nlo_data_set_real_kinematics 
\nwendcode{}\nwbegincode{437}\moddef{nlo data: procedures}\plusendmoddef
  subroutine nlo_data_set_real_kinematics (nlo_data, xi_tilde, y, phi, xi_max, &
                                           jac, jac_rand)
    class(nlo_data_t), intent(inout) :: nlo_data
    real(default), dimension(:), allocatable :: xi_max, y
    real(default), intent(in) :: xi_tilde
    real(default), intent(in) :: phi
    real(default), intent(in), dimension(3) :: jac
    real(default), intent(in), dimension(:), allocatable :: jac_rand
    nlo_data%real_kinematics%xi_tilde = xi_tilde
    nlo_data%real_kinematics%y = y
    nlo_data%real_kinematics%phi = phi
    nlo_data%real_kinematics%xi_max = xi_max
    nlo_data%real_kinematics%jac = jac
    nlo_data%real_kinematics%jac_rand = jac_rand
  end subroutine nlo_data_set_real_kinematics

\eatline
\nwendcode{}\nwbegindocs{438}\nwdocspar
\nwenddocs{}\nwbegincode{439}\moddef{nlo data: nlo data: TBP}\plusendmoddef
  procedure :: get_real_kinematics => nlo_data_get_real_kinematics
\nwendcode{}\nwbegincode{440}\moddef{nlo data: procedures}\plusendmoddef
  subroutine nlo_data_get_real_kinematics &
       (nlo_data, em, xi_tilde, y, xi_max, jac, phi, jac_rand)
    class(nlo_data_t), intent(in) :: nlo_data
    integer, intent(in) :: em
    real(default), intent(out) :: xi_tilde, y, xi_max
    real(default), intent(out), dimension(3), optional :: jac
    !!! For most applications, phi is not relevant. Thus, it is not 
    !!! always transferred as a dummy-variable
    real(default), intent(out), optional :: phi
    real(default), intent(out), dimension(:), optional :: jac_rand
    xi_tilde = nlo_data%real_kinematics%xi_tilde
    y = nlo_data%real_kinematics%y(em)
    xi_max = nlo_data%real_kinematics%xi_max (em)
    if (present (jac)) jac = nlo_data%real_kinematics%jac
    if (present (phi)) phi = nlo_data%real_kinematics%phi
    if (present (jac_rand)) jac_rand = nlo_data%real_kinematics%jac_rand
  end subroutine nlo_data_get_real_kinematics

\eatline
\nwendcode{}\nwbegindocs{441}\nwdocspar
\nwenddocs{}\nwbegincode{442}\moddef{nlo data: nlo data: TBP}\plusendmoddef
  procedure :: set_real_momenta => nlo_data_set_real_momenta 
\nwendcode{}\nwbegincode{443}\moddef{nlo data: procedures}\plusendmoddef
  subroutine nlo_data_set_real_momenta (nlo_data, p)
    class(nlo_data_t), intent(inout) :: nlo_data
    type(vector4_t), intent(in), dimension(:), allocatable :: p
    nlo_data%real_kinematics%p_real = p
  end subroutine nlo_data_set_real_momenta

\eatline
\nwendcode{}\nwbegincode{444}\moddef{nlo data: nlo data: TBP}\plusendmoddef
  procedure :: get_real_momenta => nlo_data_get_real_momenta
\nwendcode{}\nwbegincode{445}\moddef{nlo data: procedures}\plusendmoddef
  function nlo_data_get_real_momenta (nlo_data) result (p)
    class(nlo_data_t), intent(inout) :: nlo_data
    type(vector4_t), dimension(:), allocatable :: p
    p = nlo_data%real_kinematics%p_real
  end function nlo_data_get_real_momenta

\eatline
\nwendcode{}\nwbegindocs{446}\nwdocspar
\nwenddocs{}\nwbegincode{447}\moddef{nlo data: nlo data: TBP}\plusendmoddef
  procedure :: set_y_soft => nlo_data_set_y_soft
\nwendcode{}\nwbegincode{448}\moddef{nlo data: procedures}\plusendmoddef
  subroutine nlo_data_set_y_soft (nlo_data, y_soft, emitter)
    class(nlo_data_t), intent(inout) :: nlo_data
    real(default), intent(in) :: y_soft
    integer, intent(in) :: emitter
    nlo_data%real_kinematics%y_soft(emitter) = y_soft
  end subroutine nlo_data_set_y_soft

\eatline
\nwendcode{}\nwbegindocs{449}\nwdocspar
\nwenddocs{}\nwbegincode{450}\moddef{nlo data: nlo data: TBP}\plusendmoddef
  procedure :: set_jacobian => nlo_data_set_jacobian
\nwendcode{}\nwbegincode{451}\moddef{nlo data: procedures}\plusendmoddef
  subroutine nlo_data_set_jacobian (nlo_data, jac)
    class(nlo_data_t), intent(inout) :: nlo_data
    real(default), intent(in), dimension(3) :: jac
    nlo_data%real_kinematics%jac = jac
  end subroutine nlo_data_set_jacobian

\eatline
\nwendcode{}\nwbegindocs{452}\nwdocspar
\subsection{xxx}
\nwenddocs{}\nwbegincode{453}\moddef{nlo data: nlo data: TBP}\plusendmoddef
  procedure :: compute_virt => nlo_data_compute_virt
\nwendcode{}\nwbegincode{454}\moddef{nlo data: procedures}\plusendmoddef
  function nlo_data_compute_virt (nlo_data, i_flv, int_born) result(sqme_virt)
    class(nlo_data_t), intent(inout) :: nlo_data
    integer, intent(in) :: i_flv
    type(interaction_t), intent(in) :: int_born
    real(default) :: sqme_virt
    type(vector4_t), dimension(:), allocatable :: p_born
    p_born = interaction_get_momenta (int_born)
    if (nlo_data%use_internal_color_correlations) then
      call nlo_data%virtual_terms%evaluate (nlo_data%reg_data, &
                                            i_flv, &
                                            nlo_data%alpha_s_born, &
                                            p_born, &
                                            nlo_data%sqme_born(i_flv), &
                                            nlo_data%color_data%beta_ij)
    else
      call nlo_data%virtual_terms%evaluate (nlo_data%reg_data, &
                                            i_flv, &
                                            nlo_data%alpha_s_born, &
                                            p_born, &
                                            nlo_data%sqme_born(i_flv), &
                                            nlo_data%sqme_born_cc)
    end if
    sqme_virt = nlo_data%virtual_terms%sqme_virt
  end function nlo_data_compute_virt

\eatline
\nwendcode{}\nwbegindocs{455}\nwdocspar
\nwenddocs{}\nwbegincode{456}\moddef{nlo data: nlo data: TBP}\plusendmoddef
  procedure :: requires_spin_correlation => &
                  nlo_data_requires_spin_correlation
\nwendcode{}\nwbegincode{457}\moddef{nlo data: procedures}\plusendmoddef
  function nlo_data_requires_spin_correlation (nlo_data, i_flv) result (val)
    class(nlo_data_t), intent(in) :: nlo_data
    integer, intent(in) :: i_flv
    logical :: val
    val = nlo_data%real_terms%sc_required (i_flv)
  end function nlo_data_requires_spin_correlation

\eatline
\nwendcode{}\nwbegindocs{458}\nwdocspar
\clearpage
%%%%%%%%%%%%%%%%%%%%%%%%%%%%%%%%%%%%%%%%%%%%%%%%%%%%%%%%%%%%%%%%%%%%%%%
\section{The FKS phase space}
To handle the FKS phase space and adapt it to multi-channel
integration, some extensions have to be made for pre-existing types. 
The \texttt{phs\_tree}-type has an attribute containing the tree-code
of the emitting branch.  
\nwenddocs{}\nwbegincode{459}\moddef{PHS trees: phs tree: TBP}\endmoddef
  procedure :: get_emitter => phs_tree_get_emitter
\nwendcode{}\nwbegincode{460}\moddef{PHS trees: procedures}\plusendmoddef
  function phs_tree_get_emitter (tree) result (emitter)
    class(phs_tree_t), intent(in) :: tree
    integer :: emitter
    emitter = tree%emitter
  end function phs_tree_get_emitter

\eatline
\nwendcode{}\nwbegindocs{461}\nwdocspar
\nwenddocs{}\nwbegincode{462}\moddef{\code{}phs{\char95}fks.f90\edoc{}}\endmoddef
\LA{}File header\RA{}

module phs_fks
  
\LA{}Use kinds\RA{}
\LA{}Use strings\RA{}
  use constants
  use diagnostics
  use lorentz
  use flavors
  use sf_mappings
  use sf_base
  use phs_base
  use phs_wood
  use process_constants

\LA{}Standard module head\RA{}

\LA{}phs fks: public\RA{}

\LA{}phs fks: types\RA{}

\LA{}phs fks: interfaces\RA{}

contains

\LA{}phs fks: procedures\RA{}

end module phs_fks

\eatline
\nwendcode{}\nwbegindocs{463}\nwdocspar
\nwenddocs{}\nwbegincode{464}\moddef{phs fks: public}\endmoddef
  public :: phs_fks_config_t
\nwendcode{}\nwbegincode{465}\moddef{phs fks: types}\endmoddef
  type, extends (phs_wood_config_t) :: phs_fks_config_t
  contains
  \LA{}phs fks: fks config: TBP\RA{}
  end type phs_fks_config_t

\eatline
\nwendcode{}\nwbegincode{466}\moddef{phs fks: fks config: TBP}\endmoddef
  procedure :: final => phs_fks_config_final
\nwendcode{}\nwbegincode{467}\moddef{phs fks: procedures}\endmoddef
  subroutine phs_fks_config_final (object)
    class(phs_fks_config_t), intent(inout) :: object
!    call object%phs_wood_config_t%final ()
  end subroutine phs_fks_config_final
 
\eatline
\nwendcode{}\nwbegindocs{468}\nwdocspar
\nwenddocs{}\nwbegincode{469}\moddef{phs fks: fks config: TBP}\plusendmoddef
  procedure :: write => phs_fks_config_write
\nwendcode{}\nwbegincode{470}\moddef{phs fks: procedures}\plusendmoddef
  subroutine phs_fks_config_write (object, unit)
    class(phs_fks_config_t), intent(in) :: object
    integer, intent(in), optional :: unit
    call object%phs_wood_config_t%write
  end subroutine phs_fks_config_write

\eatline
\nwendcode{}\nwbegindocs{471}\nwdocspar
\nwenddocs{}\nwbegincode{472}\moddef{phs fks: fks config: TBP}\plusendmoddef
  procedure :: configure => phs_fks_config_configure
\nwendcode{}\nwbegincode{473}\moddef{phs fks: procedures}\plusendmoddef
  subroutine phs_fks_config_configure (phs_config, sqrts, &
        sqrts_fixed, cm_frame, azimuthal_dependence, rebuild, &
        ignore_mismatch, nlo_type)
    class(phs_fks_config_t), intent(inout) :: phs_config
    real(default), intent(in) :: sqrts
    logical, intent(in), optional :: sqrts_fixed
    logical, intent(in), optional :: cm_frame
    logical, intent(in), optional :: azimuthal_dependence
    logical, intent(in), optional :: rebuild
    logical, intent(in), optional :: ignore_mismatch
    type(string_t), intent(inout), optional :: nlo_type
    if (present (nlo_type)) then
      if (nlo_type /= 'Real') & 
        call msg_fatal ("FKS config has to be called with nlo_type = 'Real'")
    end if 
    if (.not. phs_config%extended_phs) &
       phs_config%n_par = phs_config%n_par + 3
!!! Channel equivalences not accessible yet
    phs_config%provides_equivalences = .false.
  end subroutine phs_fks_config_configure

\eatline
\nwendcode{}\nwbegindocs{474}\nwdocspar
\nwenddocs{}\nwbegincode{475}\moddef{phs fks: fks config: TBP}\plusendmoddef
  procedure :: startup_message => phs_fks_config_startup_message
\nwendcode{}\nwbegincode{476}\moddef{phs fks: procedures}\plusendmoddef
  subroutine phs_fks_config_startup_message (phs_config, unit)
    class(phs_fks_config_t), intent(in) :: phs_config
    integer, intent(in), optional :: unit
    call phs_config%phs_wood_config_t%startup_message
  end subroutine phs_fks_config_startup_message

\eatline
\nwendcode{}\nwbegindocs{477}\nwdocspar
\nwenddocs{}\nwbegincode{478}\moddef{phs fks: fks config: TBP}\plusendmoddef
  procedure :: increase_n_par => phs_fks_config_increase_n_par
\nwendcode{}\nwbegincode{479}\moddef{phs fks: procedures}\plusendmoddef
  subroutine phs_fks_config_increase_n_par (phs_config)
    class(phs_fks_config_t), intent(inout) :: phs_config
    phs_config%n_par = phs_config%n_par + 3
  end subroutine phs_fks_config_increase_n_par

\eatline
\nwendcode{}\nwbegindocs{480}\nwdocspar
\nwenddocs{}\nwbegincode{481}\moddef{phs fks: fks config: TBP}\plusendmoddef
  procedure, nopass :: allocate_instance => phs_fks_config_allocate_instance
\nwendcode{}\nwbegincode{482}\moddef{phs fks: procedures}\plusendmoddef
  subroutine phs_fks_config_allocate_instance (phs)
    class(phs_t), intent(inout), pointer :: phs
    allocate (phs_fks_t :: phs)
  end subroutine phs_fks_config_allocate_instance

\eatline
\nwendcode{}\nwbegindocs{483}\nwdocspar
\nwenddocs{}\nwbegincode{484}\moddef{phs fks: fks config: TBP}\plusendmoddef
  procedure :: set_born_config => phs_fks_config_set_born_config
\nwendcode{}\nwbegincode{485}\moddef{phs fks: procedures}\plusendmoddef
  subroutine phs_fks_config_set_born_config (phs_config, phs_cfg_born)
    class(phs_fks_config_t), intent(inout) :: phs_config
    type(phs_wood_config_t), intent(in), target :: phs_cfg_born
    phs_config%forest = phs_cfg_born%forest
    phs_config%n_channel = phs_cfg_born%n_channel
    phs_config%n_par = phs_cfg_born%n_par
    phs_config%sqrts = phs_cfg_born%sqrts
    phs_config%par = phs_cfg_born%par
    phs_config%sqrts_fixed = phs_cfg_born%sqrts_fixed
    phs_config%azimuthal_dependence = phs_cfg_born%azimuthal_dependence
    phs_config%provides_chains = phs_cfg_born%provides_chains
    phs_config%chain = phs_cfg_born%chain
  end subroutine phs_fks_config_set_born_config

\eatline
\nwendcode{}\nwbegindocs{486}Keep score about the real kinematics.
\nwenddocs{}\nwbegincode{487}\moddef{phs fks: public}\plusendmoddef
  public :: kinematics_counter_t
\nwendcode{}\nwbegincode{488}\moddef{phs fks: types}\plusendmoddef
  type :: kinematics_counter_t
    integer :: n_bins = 0
    integer, dimension(:), allocatable :: histo_xi
    integer, dimension(:), allocatable :: histo_xi_tilde
    integer, dimension(:), allocatable :: histo_xi_max 
    integer, dimension(:), allocatable :: histo_y
    integer, dimension(:), allocatable :: histo_phi
  contains
  \LA{}phs fks: kinematics counter: TBP\RA{}
  end type kinematics_counter_t

\eatline
\nwendcode{}\nwbegindocs{489}\nwdocspar
\nwenddocs{}\nwbegincode{490}\moddef{phs fks: kinematics counter: TBP}\endmoddef
  procedure :: init => kinematics_counter_init
\nwendcode{}\nwbegincode{491}\moddef{phs fks: procedures}\plusendmoddef
  subroutine kinematics_counter_init (counter, n_bins)
    class(kinematics_counter_t), intent(inout) :: counter
    integer, intent(in) :: n_bins
    counter%n_bins = n_bins
    allocate (counter%histo_xi (n_bins), counter%histo_xi_tilde (n_bins))
    allocate (counter%histo_y (n_bins), counter%histo_phi (n_bins))
    allocate (counter%histo_xi_max (n_bins))
    counter%histo_xi = 0
    counter%histo_xi_tilde = 0
    counter%histo_xi_max = 0
    counter%histo_y = 0
    counter%histo_phi = 0
  end subroutine kinematics_counter_init

\eatline
\nwendcode{}\nwbegindocs{492}\nwdocspar
\nwenddocs{}\nwbegincode{493}\moddef{phs fks: kinematics counter: TBP}\plusendmoddef
  procedure :: record => kinematics_counter_record
\nwendcode{}\nwbegincode{494}\moddef{phs fks: procedures}\plusendmoddef
  subroutine kinematics_counter_record (counter, xi, xi_tilde, &
                                        xi_max, y, phi)
     class(kinematics_counter_t), intent(inout) :: counter
     real(default), intent(in), optional :: xi, xi_tilde, xi_max
     real(default), intent(in), optional :: y, phi
 
     if (counter%n_bins > 0) then
       if (present (xi)) then
          call fill_histogram (counter%histo_xi, xi, &
                               0.0_default, 1.0_default)
       end if
       if (present (xi_tilde)) then
          call fill_histogram (counter%histo_xi_tilde, xi_tilde, &
                               0.0_default, 1.0_default)
       end if
       if (present (xi_max)) then
          call fill_histogram (counter%histo_xi_max, xi_max, &
                               0.0_default, 1.0_default)
       end if
       if (present (y)) then
          call fill_histogram (counter%histo_y, y, -1.0_default, 1.0_default)
       end if
       if (present (phi)) then
          call fill_histogram (counter%histo_phi, phi, 0.0_default, twopi)
       end if
     end if
  contains
     subroutine fill_histogram (histo, value, val_min, val_max)
        integer, dimension(:), allocatable :: histo
        real(default), intent(in) :: value, val_min, val_max
        real(default) :: step, lo, hi
        integer :: bin
        step = (val_max-val_min) / counter%n_bins
        do bin = 1, counter%n_bins
           lo = (bin-1) * step
           hi = bin * step
           if (value >= lo .and. value < hi) then
               histo (bin) = histo (bin) + 1
               exit
           end if
        end do
     end subroutine fill_histogram
  end subroutine kinematics_counter_record

\eatline
\nwendcode{}\nwbegincode{495}\moddef{phs fks: kinematics counter: TBP}\plusendmoddef
  procedure :: display => kinematics_counter_display
\nwendcode{}\nwbegincode{496}\moddef{phs fks: procedures}\plusendmoddef
  subroutine kinematics_counter_display (counter)
     class(kinematics_counter_t), intent(in) :: counter
     print *, 'xi: ', counter%histo_xi
     print *, 'xi_tilde: ', counter%histo_xi_tilde
     print *, 'xi_max: ', counter%histo_xi_max
     print *, 'y: ', counter%histo_y
     print *, 'phi: ', counter%histo_phi
  end subroutine kinematics_counter_display

\eatline
\nwendcode{}\nwbegindocs{497}The fks phase space type contains the wood phase space and
separately the in- and outcoming momenta for the real process and the
corresponding Born momenta. Additionally, there are the variables
$\xi$,$\xi_{max}$, $y$ and $\phi$ which are used to create the real
phase space, as well as the jacobian and its corresponding soft and
collinear limit. Lastly, the array \texttt{ch\_to\_em} connects each
channel with an emitter.
\nwenddocs{}\nwbegincode{498}\moddef{phs fks: public}\plusendmoddef
  public :: phs_fks_t
\nwendcode{}\nwbegincode{499}\moddef{phs fks: types}\plusendmoddef
  type, extends (phs_wood_t) :: phs_fks_t
    type(vector4_t), dimension(:), allocatable :: p_born
    type(vector4_t), dimension(:), allocatable :: q_born
    type(vector4_t), dimension(:), allocatable :: p_real
    type(vector4_t), dimension(:), allocatable :: q_real
    type(vector4_t), dimension(:), allocatable :: p_born_tot
    real(default) :: E_beam_tot
    real(default) :: E_gluon
    real(default) :: mrec2
    real(default), dimension(:), allocatable :: m2
    real(default) :: xi_tilde, phi
    real(default), dimension(:), allocatable :: xi_max
    real(default), dimension(:), allocatable :: y
    real(default) :: xi_min, y_max
    real(default) :: y_soft
    real(default), dimension(:), allocatable :: jac_rand
    real(default), dimension(3) :: jac
    integer, dimension(:), allocatable :: emitters
    type(kinematics_counter_t) :: counter
    logical :: massive_phsp = .false.
    logical :: perform_generation = .true.

  contains
  \LA{}phs fks: phs fks: TBP\RA{}
  end type phs_fks_t

\eatline
\nwendcode{}\nwbegindocs{500}\nwdocspar
\nwenddocs{}\nwbegincode{501}\moddef{phs fks: interfaces}\endmoddef

  interface compute_beta
    module procedure compute_beta_massless
    module procedure compute_beta_massive
  end interface

  interface get_xi_max_fsr
    module procedure get_xi_max_fsr_massless
    module procedure get_xi_max_fsr_massive
  end interface

\eatline
\nwendcode{}\nwbegindocs{502}Initializer for the phase space. Calls the initialization of the
corresponding Born phase space, sets up the
channel-emitter-association and allocates space for the momenta.  
\nwenddocs{}\nwbegincode{503}\moddef{phs fks: phs fks: TBP}\endmoddef
  procedure :: init => phs_fks_init
\nwendcode{}\nwbegincode{504}\moddef{phs fks: procedures}\plusendmoddef
  subroutine phs_fks_init (phs, phs_config)
    class(phs_fks_t), intent(out) :: phs
    class(phs_config_t), intent(in), target :: phs_config

    call phs%base_init (phs_config)
    select type (phs_config)
    type is (phs_fks_config_t)
       phs%config => phs_config
       phs%forest = phs_config%forest
    end select

    select type(phs)
    type is (phs_fks_t)
      phs%n_r_born = phs%config%n_par - 3
      call phs%init_momenta (phs_config)
      allocate (phs%m2 (phs_config%n_tot))
      allocate (phs%xi_max (phs_config%n_tot))
      allocate (phs%y (phs_config%n_tot))
      allocate (phs%jac_rand (phs_config%n_tot))
      phs%xi_max = 0._default
      phs%y = 0._default
      phs%jac_rand = 1._default
      phs%jac = 1._default
      phs%xi_min = 1d-5
      phs%y_max = 1._default
    end select
  end subroutine phs_fks_init

\eatline
\nwendcode{}\nwbegindocs{505}\nwdocspar
\nwenddocs{}\nwbegincode{506}\moddef{phs fks: phs fks: TBP}\plusendmoddef
  procedure :: final => phs_fks_final
\nwendcode{}\nwbegincode{507}\moddef{phs fks: procedures}\plusendmoddef
  subroutine phs_fks_final (object)
    class(phs_fks_t), intent(inout) :: object
  end subroutine phs_fks_final

\eatline
\nwendcode{}\nwbegindocs{508}\nwdocspar
\nwenddocs{}\nwbegincode{509}\moddef{phs fks: phs fks: TBP}\plusendmoddef
  procedure :: init_momenta => phs_fks_init_momenta
\nwendcode{}\nwbegincode{510}\moddef{phs fks: procedures}\plusendmoddef
  subroutine phs_fks_init_momenta (phs, phs_config)
    class(phs_fks_t), intent(inout) :: phs
    class(phs_config_t), intent(in) :: phs_config
    allocate (phs%p_born (phs_config%n_in))
    allocate (phs%q_born (phs_config%n_out-1))
    allocate (phs%p_real (phs_config%n_in))
    allocate (phs%q_real (phs_config%n_out-1))
    allocate (phs%p_born_tot (phs%config%n_in + phs%config%n_out-1))
  end subroutine phs_fks_init_momenta

\eatline
\nwendcode{}\nwbegindocs{511}Evaluate selected channel. First, the subroutine calls the
evaluation procedure of the underlying Born phase space, using $n_r -
3$ random numbers. Then, the remaining three random numbers are used
to create $\xi$, $y$ and $\phi$, from which the real momenta are
calculated from the Born momenta. 
\nwenddocs{}\nwbegincode{512}\moddef{phs fks: phs fks: TBP}\plusendmoddef
  procedure :: evaluate_selected_channel => phs_fks_evaluate_selected_channel
\nwendcode{}\nwbegincode{513}\moddef{phs fks: procedures}\plusendmoddef
  subroutine phs_fks_evaluate_selected_channel (phs, c_in, r_in)
    class(phs_fks_t), intent(inout) :: phs
    integer, intent(in) :: c_in
    real(default), intent(in), dimension(:) :: r_in
    integer :: em 
    real(default), dimension(:), allocatable :: r_born
    integer :: n_r_born
    integer :: n_in
 
    n_in = phs%config%n_in
    call phs%phs_wood_t%evaluate_selected_channel (c_in, r_in)
    phs%p_born = phs%phs_wood_t%p
    phs%q_born = phs%phs_wood_t%q
    phs%p_born_tot (1:n_in) = phs%p_born
    phs%p_born_tot (n_in+1:) = phs%q_born
    phs%r(:,c_in) = r_in

    if (phs%perform_generation) call phs%generate_real_kinematics &
                                             (r_in(phs%n_r_born+1:phs%n_r_born+3))
  end subroutine phs_fks_evaluate_selected_channel

\eatline
\nwendcode{}\nwbegindocs{514}\nwdocspar
\nwenddocs{}\nwbegincode{515}\moddef{phs fks: phs fks: TBP}\plusendmoddef
  procedure :: evaluate_other_channels => phs_fks_evaluate_other_channels
\nwendcode{}\nwbegincode{516}\moddef{phs fks: procedures}\plusendmoddef
  subroutine phs_fks_evaluate_other_channels (phs, c_in)
    class(phs_fks_t), intent(inout) :: phs
    integer, intent(in) :: c_in
    integer :: c
    call phs%phs_wood_t%evaluate_other_channels (c_in)
    phs%r_defined = .true.
  end subroutine phs_fks_evaluate_other_channels

\eatline
\nwendcode{}\nwbegindocs{517}\nwdocspar
\nwenddocs{}\nwbegincode{518}\moddef{phs fks: phs fks: TBP}\plusendmoddef
  procedure :: get_mcpar => phs_fks_get_mcpar
\nwendcode{}\nwbegincode{519}\moddef{phs fks: procedures}\plusendmoddef
  subroutine phs_fks_get_mcpar (phs, c, r)
    class(phs_fks_t), intent(in) :: phs
    integer, intent(in) :: c
    real(default), dimension(:), intent(out) :: r
    integer :: n_r_born 
    r(1:phs%n_r_born) = phs%r(1:phs%n_r_born,c)
    r(phs%n_r_born+1:) = phs%r_real
  end subroutine phs_fks_get_mcpar

\eatline
\nwendcode{}\nwbegindocs{520}\nwdocspar
\nwenddocs{}\nwbegincode{521}\moddef{phs fks: phs fks: TBP}\plusendmoddef
  procedure :: get_real_kinematics => phs_fks_get_real_kinematics
\nwendcode{}\nwbegincode{522}\moddef{phs fks: procedures}\plusendmoddef
  subroutine phs_fks_get_real_kinematics (phs, xit, y, phi, xi_max, jac, jac_rand)
    class(phs_fks_t), intent(inout) :: phs
    real(default), intent(out), dimension(:), allocatable :: xi_max
    real(default), intent(out) :: xit
    real(default), intent(out), dimension(:), allocatable :: y
    real(default), intent(out) :: phi
    real(default), intent(out), dimension(3) :: jac
    real(default), intent(out), dimension(:), allocatable :: jac_rand
    xit = phs%xi_tilde
    y = phs%y
    phi = phs%phi
    xi_max = phs%xi_max
    jac = phs%jac
    jac_rand = phs%jac_rand
  end subroutine phs_fks_get_real_kinematics

\eatline
\nwendcode{}\nwbegindocs{523}\nwdocspar
\nwenddocs{}\nwbegincode{524}\moddef{phs fks: phs fks: TBP}\plusendmoddef
  procedure :: enable_massive_emitters => &
                     phs_fks_enable_massive_emitters
\nwendcode{}\nwbegincode{525}\moddef{phs fks: procedures}\plusendmoddef
  subroutine phs_fks_enable_massive_emitters (phs_fks)
    class(phs_fks_t), intent(inout) :: phs_fks
    phs_fks%massive_phsp = .true.
  end subroutine phs_fks_enable_massive_emitters

\eatline
\nwendcode{}\nwbegindocs{526}\nwdocspar
\nwenddocs{}\nwbegincode{527}\moddef{phs fks: phs fks: TBP}\plusendmoddef
  procedure :: set_emitters => phs_fks_set_emitters
\nwendcode{}\nwbegincode{528}\moddef{phs fks: procedures}\plusendmoddef
  subroutine phs_fks_set_emitters (phs, emitters)
    class(phs_fks_t), intent(inout) :: phs
    integer, intent(in), dimension(:), allocatable :: emitters
    phs%emitters = emitters
  end subroutine phs_fks_set_emitters

\eatline
\nwendcode{}\nwbegindocs{529}\nwdocspar
\nwenddocs{}\nwbegincode{530}\moddef{phs fks: phs fks: TBP}\plusendmoddef
  procedure :: get_born_momenta => phs_fks_get_born_momenta
\nwendcode{}\nwbegincode{531}\moddef{phs fks: procedures}\plusendmoddef
  subroutine phs_fks_get_born_momenta (phs, p)
    class(phs_fks_t), intent(inout) :: phs
    type(vector4_t), intent(out), dimension(:) :: p
    p(1:phs%config%n_in) = phs%p_born
    p(phs%config%n_in+1:) = phs%q_born
  end subroutine phs_fks_get_born_momenta

\eatline
\nwendcode{}\nwbegindocs{532}\nwdocspar
\nwenddocs{}\nwbegincode{533}\moddef{phs fks: phs fks: TBP}\plusendmoddef
  procedure :: get_outgoing_momenta => phs_fks_get_outgoing_momenta
\nwendcode{}\nwbegincode{534}\moddef{phs fks: procedures}\plusendmoddef
  subroutine phs_fks_get_outgoing_momenta (phs, q)
    class(phs_fks_t), intent(in) :: phs
    type(vector4_t), intent(out), dimension(:) :: q
    q = phs%q_real
  end subroutine phs_fks_get_outgoing_momenta

\eatline
\nwendcode{}\nwbegindocs{535}\nwdocspar
\nwenddocs{}\nwbegincode{536}\moddef{phs fks: phs fks: TBP}\plusendmoddef
  procedure :: get_incoming_momenta => phs_fks_get_incoming_momenta
\nwendcode{}\nwbegincode{537}\moddef{phs fks: procedures}\plusendmoddef
  subroutine phs_fks_get_incoming_momenta (phs, p)
    class(phs_fks_t), intent(in) :: phs
    type(vector4_t), intent(inout), dimension(:), allocatable :: p
    p = phs%p_real
  end subroutine phs_fks_get_incoming_momenta

\eatline
\nwendcode{}\nwbegindocs{538}\nwdocspar
\nwenddocs{}\nwbegincode{539}\moddef{phs fks: phs fks: TBP}\plusendmoddef
  procedure :: display_kinematics => phs_fks_display_kinematics
\nwendcode{}\nwbegincode{540}\moddef{phs fks: procedures}\plusendmoddef
  subroutine phs_fks_display_kinematics (phs)
     class(phs_fks_t), intent(in) :: phs
!     call phs%counter%display ()
  end subroutine phs_fks_display_kinematics

\eatline
\nwendcode{}\nwbegindocs{541}\nwdocspar
\nwenddocs{}\nwbegincode{542}\moddef{phs fks: phs fks: TBP}\plusendmoddef
  procedure :: generate_real_kinematics => phs_fks_generate_real_kinematics
\nwendcode{}\nwbegincode{543}\moddef{phs fks: procedures}\plusendmoddef
  subroutine phs_fks_generate_real_kinematics (phs, r_in)
    class(phs_fks_t), intent(inout) :: phs
    real(default), intent(in), dimension(:) :: r_in
    integer :: em
    real(default) :: beta

    if (size (r_in) /= 3) call msg_fatal &
                        ("Real kinematics need to be generated using three random numbers!")
    phs%E_beam_tot = phs%p_born(1)%p(0) + phs%p_born(2)%p(0)

    !!! Jacobian corresponding to the transformation rand -> (xi, y, phi)
    phs%jac_rand = 1.0
    !!! Produce real momentum
    phs%phi = r_in(3)*twopi
    phs%xi_tilde = phs%xi_min + r_in(1)*(1-phs%xi_min)
    phs%jac_rand = phs%jac_rand * (1-phs%xi_min)       
    do em = 1, phs%config%n_tot
      if (any (phs%emitters == em)) then
         if (phs%massive_phsp) then
            phs%m2(em) = phs%p_born_tot(em)**2
            beta = beta_emitter (phs%E_beam_tot, phs%p_born_tot(em))
            phs%y(em) = 1.0/beta * (1-(1+beta) * &
                            exp(-r_in(2)*log((1+beta)/(1-beta))))
            phs%jac_rand(em) = phs%jac_rand(em) * &
                               (1-beta*phs%y(em))*log((1+beta)/(1-beta))/beta
            phs%xi_max (em) = get_xi_max_fsr &
                                   (phs%p_born_tot, em, phs%m2(em), phs%y(em))
         else
            phs%m2(em) = 0._default
            phs%xi_max (em) = get_xi_max_fsr (phs%p_born_tot, em)
         end if
      end if
    end do
    if (.not. phs%massive_phsp) then
       phs%y = (1-2*r_in(2))*phs%y_max
       phs%jac_rand = phs%jac_rand*3*(1-phs%y(1)**2)
       phs%y = 1.5_default*(phs%y - phs%y**3/3)
    end if

    phs%r_real = r_in
  contains
    function beta_emitter (q0, p) result (beta)
      real(default), intent(in) :: q0
      type(vector4_t), intent(in) :: p
      real(default) :: beta
      real(default) :: m2, mrec2, k0_max
      m2 = p**2
      mrec2 = (q0-p%p(0))**2 - p%p(1)**2 - p%p(2)**2 - p%p(3)**2
      k0_max = (q0**2-mrec2+m2)/(2*q0)
      beta = sqrt(1-m2/k0_max**2)
    end function beta_emitter 
  end subroutine phs_fks_generate_real_kinematics

\eatline
\nwendcode{}\nwbegindocs{544}\nwdocspar
\subsection{Creation of the real phase space - FSR}
At this point, the Born phase space has been generated, as well as the
three random variables $\xi$, $y$ and $\phi$. The question is how the
real phase space is generated for a final-state emission
configuration. We work with two different sets of momenta, the Born
configuration $\Bigl\{ \bar{k}_{\oplus}, \bar{k}_{\ominus}, \bar{k}_{1}, ...,
\bar{k}_{n} \Bigr\}$ and the real configuration $\Bigl\{ k_{\oplus},
k_{\ominus}, k_1,..., k_n, k_{n+1} \Bigr\}$. We define the momentum of
the emitter to be on the $n$-th position and the momentum of the
radiated particle to be at position $n+1$. The magnitude of the
spatial component of k is denoted by $\underline{k}$.

For final-state emissions, it is $\bar{k}_\oplus = k_\oplus$ and
$\bar{k}_\ominus = k_\ominus$. Thus, the center-of-mass systems
coincide and it is 
\begin{equation}
  q = \sum_{i=1}^n \bar{k}_i = \sum_{i=1}^{n+1} k_i,
\end{equation}
with $\vec{q} = 0$ and $q^2 = \left(q^0\right)^2$. 

We want to construct the real phase space from the Born phase space
using three random numbers. They are defined as follows: 
\begin{itemize}
\item $\xi = \frac{2k_{n+1}^0}{\sqrt{s}} \in [0, \xi_{max}]$, where
  $k_{n+1}$ denotes the four-momentum of the radiated particle. 
\item $y = \cos\theta = \frac{\vec{k}_n \cdot
    \vec{k}_{n+1}}{\underline{k}_n \underline{k}_{n+1}}$ is the
  splitting angle. 
\item The angle between tho two splitting particles in the transversal
  plane, $phi \in [0,2\pi]$.  
\end{itemize}
Further, $k_{rec} = \sum_{i=1}^{n-1} k_i$ denotes the sum of all
recoiling momenta. 
%\begin{dubious}
%  Note that this calculation only works for massless particles. 
%\end{dubious}
\nwenddocs{}\nwbegincode{545}\moddef{phs fks: phs fks: TBP}\plusendmoddef
  procedure :: generate_fsr => phs_fks_generate_fsr
\nwendcode{}\nwbegincode{546}\moddef{phs fks: procedures}\plusendmoddef
  subroutine phs_fks_generate_fsr (phs, emitter, p_born, p_real)
    !!! Important: Momenta must be input in the center-of-mass frame
    class(phs_fks_t), intent(inout) :: phs
    integer, intent(in) :: emitter
    type(vector4_t), intent(in), dimension(:) :: p_born 
    type(vector4_t), intent(out), dimension(:), allocatable :: p_real
    integer nlegborn, nlegreal
    type(vector4_t) :: q
    real(default) :: q0, q2, uk_np1, uk_n
    real(default) :: uk_rec, k_rec0
    type(vector3_t) :: k_n_born, k
    real(default) :: uk_n_born
    real(default) :: uk, k2, k0_n
    real(default) :: cpsi, beta
    type(vector3_t) :: vec, vec_orth
    type(lorentz_transformation_t) :: rot, lambda
    integer :: i
    real(default) :: xi, y, phi

    xi = phs%xi_tilde * phs%xi_max(emitter)
    y = phs%y(emitter)
    phi = phs%phi
    nlegreal = phs%config%n_in + phs%config%n_out
    nlegborn = nlegreal-1
    if (emitter <= 2 .or. emitter > nlegborn) then
      call msg_fatal ("Generate FSR phase space: Invalid emitter!")
    end if
    allocate (p_real (nlegreal))

    p_real(1) = p_born(1)
    p_real(2) = p_born(2)
    q = p_born(1) + p_born(2)
    q0 = phs%E_beam_tot
    q2 = q**2

    phs%E_gluon = q0*xi/2
    uk_np1 = phs%E_gluon
    k_n_born = p_born(emitter)%p(1:3)
    uk_n_born = k_n_born**1

    phs%mrec2 = (q-p_born(emitter))**2
    if (phs%massive_phsp) then
       call phs%compute_emitter_kinematics (emitter, k0_n, uk_n, uk)
    else
       call phs%compute_emitter_kinematics (emitter, uk_n, uk)
       phs%y_soft = y
       k0_n = uk_n
    end if
    
    vec = uk_n / uk_n_born * k_n_born
    vec_orth = create_orthogonal (vec)
    p_real(emitter)%p(0) = k0_n
    p_real(emitter)%p(1:3) = vec%p(1:3)
    cpsi = (uk_n**2 + uk**2 - uk_np1**2) / (2*(uk_n * uk))
    !!! This is to catch the case where cpsi = 1, but numerically 
    !!! turns out to be slightly larger than 1. 
    if (cpsi > 1._default) cpsi = 1._default
    rot = rotation (cpsi, -sqrt (1-cpsi**2), vec_orth)
    p_real(emitter) = rot*p_real(emitter)
    vec = uk_np1 / uk_n_born * k_n_born
    vec_orth = create_orthogonal (vec)
    p_real(nlegreal)%p(0) = uk_np1
    p_real(nlegreal)%p(1:3) = vec%p(1:3)
    cpsi = (uk_np1**2 + uk**2 - uk_n**2) / (2*(uk_np1 * uk))
    rot = rotation (cpsi, sqrt (1-cpsi**2), vec_orth)
    p_real(nlegreal) = rot*p_real(nlegreal)
\nwendcode{}\nwbegindocs{547}Construction of the recoiling momenta. The reshuffling of momenta
must not change the invariant mass of the recoiling system, which
means $k_{\rm{rec}}^2 = \bar{k_{\rm{rec}}}^2$. Therefore, the momenta
are related by a boost, $\bar{k}_i = \Lambda k_i$. The boost parameter
is 
\begin{equation*}
  \beta = \frac{q^2 - (k_{\rm{rec}}^0 +
    \underline{k}_{\rm{rec}})^2}{q^2 + (k_{\rm{rec}}^0 +
    \underline{k}_{\rm{rec}})^2} 
\end{equation*}
\nwenddocs{}\nwbegincode{548}\moddef{phs fks: procedures}\plusendmoddef
    k_rec0 = q0 - p_real(emitter)%p(0) - p_real(nlegreal)%p(0)
    uk_rec = sqrt (k_rec0**2 - phs%mrec2)
    if (phs%massive_phsp) then
       beta = compute_beta (q2, k_rec0, uk_rec, &
                            p_born(emitter)%p(0), uk_n_born)
    else
       beta = compute_beta (q2, k_rec0, uk_rec)
    end if   
    k = p_real(emitter)%p(1:3) + p_real(nlegreal)%p(1:3)
    vec%p(1:3) = 1/uk*k%p(1:3)
    lambda = boost (beta/sqrt(1-beta**2), vec)
    do i = 3, nlegborn
      if (i /= emitter) then
        p_real(i) = lambda * p_born(i)
      end if
    end do
    vec%p(1:3) = p_born(emitter)%p(1:3)/uk_n_born
    rot = rotation (cos(phi), sin(phi), vec)
    p_real(nlegreal) = rot * p_real(nlegreal)
    p_real(emitter) = rot * p_real(emitter)
\nwendcode{}\nwbegindocs{549}The factor $\frac{q^2}{(4\pi)^3}$ is not included here since it is
supplied during phase space generation. Also, we already divide by
$\xi$. 
\nwenddocs{}\nwbegincode{550}\moddef{phs fks: procedures}\plusendmoddef
    if (phs%massive_phsp) then 
       phs%jac(1) = phs%jac(1)*4/q0/uk_n_born/xi
    else
       k2 = 2*uk_n*uk_np1*(1-y)
       phs%jac(1) = uk_n**2/uk_n_born / (uk_n - k2/(2*q0))
    end if
    !!! Soft jacobian
    phs%jac(2) = 1.0
    !!! Collinear jacobian
    phs%jac(3) = 1-xi/2*q0/uk_n_born

  end subroutine phs_fks_generate_fsr

\eatline
\nwendcode{}\nwbegincode{551}\moddef{phs fks: phs fks: TBP}\plusendmoddef
  generic :: compute_emitter_kinematics => &
                    compute_emitter_kinematics_massless, &
                    compute_emitter_kinematics_massive 
  procedure :: compute_emitter_kinematics_massless => &
                    phs_fks_compute_emitter_kinematics_massless
  procedure :: compute_emitter_kinematics_massive => &
                    phs_fks_compute_emitter_kinematics_massive
\nwendcode{}\nwbegincode{552}\moddef{phs fks: procedures}\plusendmoddef
  subroutine phs_fks_compute_emitter_kinematics_massless &
                                    (phs, em, uk_em, uk)
    class(phs_fks_t), intent(inout) :: phs
    integer, intent(in) :: em
    real(default), intent(out) :: uk_em, uk
    real(default) :: y, k0_np1, q0, q2

    y = phs%y(em); k0_np1 = phs%E_gluon
    q0 = phs%E_beam_tot; q2 = q0**2

    uk_em = (q2 - phs%mrec2 - 2*q0*k0_np1) / (2*(q0 - k0_np1*(1-y))) 
    uk = sqrt (uk_em**2 + k0_np1**2 + 2*uk_em*k0_np1*y)
  end subroutine phs_fks_compute_emitter_kinematics_massless

  subroutine phs_fks_compute_emitter_kinematics_massive &
                                    (phs, em, k0_em, uk_em, uk)
    class(phs_fks_t), intent(inout) :: phs
    integer, intent(in) :: em
    real(default), intent(inout) :: k0_em, uk_em, uk
    real(default) :: y, k0_np1, q0, q2, mrec2, m2
    real(default) :: k0_rec_max, k0_em_max, k0_rec, uk_rec
    real(default) :: z, z1, z2

    y = phs%y(em); k0_np1 = phs%E_gluon
    q0 = phs%E_beam_tot; q2 = q0**2
    mrec2 = phs%mrec2; m2 = phs%m2(em)

    k0_rec_max = (q2-m2+mrec2)/(2*q0)
    k0_em_max = (q2+m2-mrec2)/(2*q0)
    z1 = (k0_rec_max+sqrt (k0_rec_max**2-mrec2))/q0
    z2 = (k0_rec_max-sqrt (k0_rec_max**2-mrec2))/q0
    z = z2 - (z2-z1)*(1+y)/2
    k0_em = k0_em_max - k0_np1*z
    k0_rec = q0 - k0_np1 - k0_em
    uk_em = sqrt(k0_em**2-m2)
    uk_rec = sqrt(k0_rec**2 - mrec2)
    uk = uk_rec
    phs%jac = q0*(z1-z2)/4*k0_np1
    phs%y_soft = (2*q2*z-q2-mrec2+m2)/(sqrt(k0_em_max**2-m2)*q0)/2
  end subroutine phs_fks_compute_emitter_kinematics_massive

\eatline
\nwendcode{}\nwbegindocs{553}\nwdocspar
\nwenddocs{}\nwbegincode{554}\moddef{phs fks: procedures}\plusendmoddef
  function compute_beta_massless (q2, k0_rec, uk_rec) result (beta)
    real(default), intent(in) :: q2, k0_rec, uk_rec
    real(default) :: beta
    beta = (q2 - (k0_rec + uk_rec)**2) / (q2 + (k0_rec + uk_rec)**2)
  end function compute_beta_massless

  function compute_beta_massive (q2, k0_rec, uk_rec, &
                                 k0_em_born, uk_em_born) result (beta)
    real(default), intent(in) :: q2, k0_rec, uk_rec
    real(default), intent(in) :: k0_em_born, uk_em_born
    real(default) :: beta
    real(default) :: k0_rec_born, uk_rec_born, alpha
    k0_rec_born = sqrt(q2) - k0_em_born
    uk_rec_born = uk_em_born
    alpha = (k0_rec+uk_rec)/(k0_rec_born+uk_rec_born)
    beta = (1-alpha**2)/(1+alpha**2)
  end function compute_beta_massive

\eatline
\nwendcode{}\nwbegindocs{555}The momentum of the radiated particle is computed according to 
\begin{equation}
  \label{eq:phs fks:compute k_n}
  \underline{k}_n = \frac{q^2 - M_{\rm{rec}}^2 -
    2q^0\underline{k}_{n+1}}{2(q^0 - \underline{k}_{n+1}(1-y))}, 
\end{equation}
with $k = k_n + k_{n+1}$ and $M_{\rm{rec}}^2 = k_{\rm{rec}}^2 =
\left(q-k\right)^2$. Because of $\boldsymbol{\bar{k}}_n \parallel
\boldsymbol{k}_n + \boldsymbol{k}_{n+1}$ we find $M_{\rm{rec}}^2 =
\left(q-\bar{k}_n\right)^2$. 
Equation \ref{eq:phs fks: compute k_n} follows from the fact that
$\left(\boldsymbol{k} - \boldsymbol{k}_n\right)^2 =
\boldsymbol{k}_{n+1}^2$, which is equivalent to $\boldsymbol{k}_n
\cdot \boldsymbol{k} = \frac{1}{2} \left(\underline{k}_n^2 +
  \underline{k}^2 - \underline{k}_{n+1}^2\right)$.\\ 
$\boldsymbol{k}_n$ and $\boldsymbol{k}_{n+1}$ are obtained by first
setting up vectors parallel to $\boldsymbol{\bar{k}}_n$,  
\begin{equation*}
  \boldsymbol{k}_n' = \underline{k}_n
  \frac{\bar{\pmb{k}}_n}{\underline{\bar{k}}_n}, \quad \pmb{k}_{n+1}'
  = \underline{k}_{n+1}\frac{\bar{\pmb{k}}_n}{\underline{\bar{k}}_n}, 
\end{equation*}
and then rotating these vectors by an amount of $\cos\psi_n =
\frac{\boldsymbol{k}_n\cdot\pmb{k}}{\underline{k}_n \underline{k}}$. 
\nwenddocs{}\nwbegindocs{556}The emitted particle cannot have more momentum than the emitter has
in the Born phase space. Thus, there is an upper bound for $\xi$,
determined by the condition $k_{n+1}^0 = \underline{\bar{k}}_n$, which
is equal to 
\begin{equation*}
\xi_{\rm{max}} = \frac{2}{\underline{\bar{k}}_n}{q^0}.
\end{equation*}
\nwenddocs{}\nwbegincode{557}\moddef{phs fks: procedures}\plusendmoddef
  function get_xi_max_fsr_massless (p_born, emitter) result (xi_max)
    type(vector4_t), intent(in), dimension(:) :: p_born
    integer, intent(in) :: emitter
    real(default) :: xi_max
    real(default) :: uk_n_born, q0
    q0 = vector4_get_component (p_born(1) + p_born(2), 0)
    uk_n_born = space_part_norm (p_born(emitter))
    xi_max = 2*uk_n_born / q0
  end function get_xi_max_fsr_massless
  
\eatline
\nwendcode{}\nwbegindocs{558}\nwdocspar
\nwenddocs{}\nwbegincode{559}\moddef{phs fks: procedures}\plusendmoddef
  function get_xi_max_fsr_massive (p_born, emitter, m2, y) result (xi_max)
    type(vector4_t), intent(in), dimension(:) :: p_born
    integer, intent(in) :: emitter
    real(default), intent(in) :: m2, y
    real(default) :: xi_max
    real(default) :: q0, mrec2, k0_rec_max
    real(default) :: z, z1, z2
    real(default) :: k_np1_max
    q0 = 2*p_born(1)%p(0)
    associate (p => p_born(emitter)%p)
       mrec2 = (q0-p(0))**2 - p(1)**2 - p(2)**2 - p(3)**2
    end associate
    k0_rec_max = (q0**2-m2+mrec2)/(2*q0)
    z1 = (k0_rec_max+sqrt(k0_rec_max**2-mrec2))/q0
    z2 = (k0_rec_max-sqrt(k0_rec_max**2-mrec2))/q0
    z = z2 - (z2-z1)*(1+y)/2
    k_np1_max = -(q0**2*z**2 - 2*q0*k0_rec_max*z + mrec2)/(2*q0*z*(1-z))
    xi_max = 2*k_np1_max/q0
  end function get_xi_max_fsr_massive

\eatline
\nwendcode{}\nwbegindocs{560}\nwdocspar
\subsection{Creation of the real phase space - ISR}
\nwenddocs{}\nwbegincode{561}\moddef{phs fks: phs fks: TBP}\plusendmoddef
  procedure :: generate_isr => phs_fks_generate_isr
\nwendcode{}\nwbegincode{562}\moddef{phs fks: procedures}\plusendmoddef
  subroutine phs_fks_generate_isr &
       (phs, emitter, p_born, p_real)
    !!! Important: Import momenta in the lab frame
    class(phs_fks_t), intent(inout) :: phs
    integer, intent(in) :: emitter
    type(vector4_t), intent(in) , dimension(:), allocatable :: p_born
    type(vector4_t), intent(out), dimension(:), allocatable :: p_real
    real(default) :: xi, y, phi
    integer :: nlegborn, nlegreal
    real(default) :: sqrts_beam
    real(default) :: sqrts_born
    real(default) :: k0_np1
    type(vector3_t) :: beta_l, vec_t
    real(default) :: beta_t, beta_gamma_l, beta_gamma_t
    real(default) :: k_tmp, k_p, k_m, k_p0, k_m0
    real(default) :: pt_rad
    type(lorentz_transformation_t) :: lambda_t, lambda_l, lambda_l_inv
    real(default) :: x_plus, x_minus, barx_plus, barx_minus
    integer :: i

    xi = phs%xi_tilde * phs%xi_max(emitter)
    y = phs%y(emitter)
    phi = phs%phi
    nlegreal = phs%config%n_in + phs%config%n_out
    nlegborn = nlegreal - 1
    sqrts_beam = phs%config%sqrts
    sqrts_born = sqrt ((p_born(1) + p_born(2))**2)
    allocate (p_real (nlegreal))
    !!! Create radiation momentum
    k0_np1 = sqrts_born*xi/2 
    !!! There must be the real cm-energy, not the Born one! 
    !!!    s_real = s_born / (1-xi)
    !!! Build radiation momentum in the rest frame of the real momenta
    p_real(nlegreal)%p(0) = k0_np1
    p_real(nlegreal)%p(1) = k0_np1*sqrt(1-y**2)*sin(phi)
    p_real(nlegreal)%p(2) = k0_np1*sqrt(1-y**2)*cos(phi)
    p_real(nlegreal)%p(3) = k0_np1*y

    !!! Boost to lab frame missing
    pt_rad = transverse_part (p_real(nlegreal))
    beta_t = sqrt (1 + sqrts_born**2 * (1-xi) / pt_rad**2)
    beta_gamma_t = 1/sqrt(beta_t)
    k_p0 = vector4_get_component (p_born(1), 0)
    k_m0 = vector4_get_component (p_born(2), 0)
    
    beta_l%p(1:3) = (p_born(1)%p(1:3) + p_born(2)%p(1:3)) / &
                    (p_born(1)%p(0) + p_born(2)%p(0))
    beta_gamma_l = beta_l**1
    beta_l = beta_l / beta_gamma_l
    beta_gamma_l = beta_gamma_l / sqrt (1 - beta_gamma_l**2)
    vec_t%p(1:2) = p_real(nlegreal)%p(1:2)
    vec_t%p(3) = 0._default
    call normalize (vec_t)

    lambda_l = boost(beta_gamma_l, beta_l)
    lambda_t = boost(-beta_gamma_t, vec_t)
    lambda_l_inv = boost(-beta_gamma_l, beta_l)
    forall (i=3:nlegborn) &
         p_real(i) = lambda_l_inv * lambda_t * lambda_l * p_born(i)
    !!! Now need access to the x-variables of the IS-partons
    barx_plus = 2*vector4_get_component(p_born(1), 0)/sqrts_beam
    barx_minus = 2*vector4_get_component(p_born(2), 0)/sqrts_beam
    x_plus = barx_plus/sqrt(1-xi) * sqrt ((2-xi*(1-y)) / (2-xi*(1+y)))
    x_minus = barx_minus/sqrt(1-xi) * sqrt ((2-xi*(1+y)) / (2-xi*(1-y)))
    p_real(1) = x_plus/barx_plus * p_born(1)
    p_real(2) = x_minus/barx_minus * p_born(2)
    !!! Total nonsense
    phs%jac(1) = 1
  end subroutine phs_fks_generate_isr
\eatline
\nwendcode{}\nwfilename{../../../src/model_features/model_features.nw}\nwbegindocs{0}% -*- ess-noweb-default-code-mode: f90-mode; noweb-default-code-mode: f90-mode; -*- 
\nwenddocs{}\nwbegindocs{1}\nwdocspar
%%%%%%%%%%%%%%%%%%%%%%%%%%%%%%%%%%%%%%%%%%%%%%%%%%%%%%%%%%%%%%%%%%%%%%%%
\chapter{Model Handling and Features}

These modules deal with process definitions and physics models.

These modules use the {\tt{}model{\char95}data} methods to automatically generate
process definitions.
\begin{description}
\item[auto\_components]
  Generic process-definition generator.  We can specify a basic
  process or initial particle(s) and some rules to extend this
  process, given a model definition with particle names and vertex
  structures.
\item[radiation\_generator]
  Applies the generic generator to the specific problem of generating
  NLO corrections in a restricted setup.
\end{description}

Model construction:
\begin{description}
\item[eval\_trees]
  Implementation of the generic {\tt{}expr{\char95}t} type for the concrete
  evaluation of expressions that access user variables.
  
  This module is actually part of the Sindarin language implementation, and
  should be moved elsewhere.  Currently, the {\tt{}models} module relies
  on it.
\item[models]
  Extends the {\tt{}model{\char95}data{\char95}t} structure by user-variable objects for
  easy access, and provides the means to read a model definition from file.
\item[slha\_interface]
  Read/write a SUSY model in the standardized SLHA format.  The format
  defines fields and parameters, but no vertices.
\end{description}
  
\clearpage
%%%%%%%%%%%%%%%%%%%%%%%%%%%%%%%%%%%%%%%%%%%%%%%%%%%%%%%%%%%%%%%%%%%%%%%%
\section{Automatic generation of process components}

This module provides the functionality for automatically generating radiation
corrections or decays, provided as lists of PDG codes.
\nwenddocs{}\nwbegincode{2}\moddef{\code{}auto{\char95}components.f90\edoc{}}\endmoddef
\LA{}File header\RA{}

module auto_components

\LA{}Use kinds\RA{}
\LA{}Use strings\RA{}
  use io_units
  use unit_tests
  use diagnostics
  use model_data
  use pdg_arrays
  use model_testbed

\LA{}Standard module head\RA{}

\LA{}Auto components: public\RA{}

\LA{}Auto components: types\RA{}

\LA{}Auto components: interfaces\RA{}

contains

\LA{}Auto components: procedures\RA{}

\LA{}Auto components: tests\RA{}

end module auto_components
\eatline
\nwendcode{}\nwbegindocs{3}\nwdocspar
\subsection{Constraints: Abstract types}
An abstract type that denotes a constraint on the automatically generated
states.  The concrete objects are applied as visitor objects at certain hooks
during the splitting algorithm.
\nwenddocs{}\nwbegincode{4}\moddef{Auto components: types}\endmoddef
  type, abstract :: split_constraint_t
   contains
   \LA{}Auto components: split constraint: TBP\RA{}
  end type split_constraint_t

\eatline
\nwendcode{}\nwbegindocs{5}By default, all checks return true.
\nwenddocs{}\nwbegincode{6}\moddef{Auto components: split constraint: TBP}\endmoddef
  procedure :: check_before_split  => split_constraint_check_before_split
  procedure :: check_before_insert => split_constraint_check_before_insert
  procedure :: check_before_record => split_constraint_check_before_record
\nwendcode{}\nwbegincode{7}\moddef{Auto components: procedures}\endmoddef
  subroutine split_constraint_check_before_split (c, table, pl, k, passed)
    class(split_constraint_t), intent(in) :: c
    class(ps_table_t), intent(in) :: table
    type(pdg_list_t), intent(in) :: pl
    integer, intent(in) :: k
    logical, intent(out) :: passed
    passed = .true.
  end subroutine split_constraint_check_before_split
  
  subroutine split_constraint_check_before_insert (c, table, pa, pl, passed)
    class(split_constraint_t), intent(in) :: c
    class(ps_table_t), intent(in) :: table
    type(pdg_array_t), intent(in) :: pa
    type(pdg_list_t), intent(inout) :: pl
    logical, intent(out) :: passed
    passed = .true.
  end subroutine split_constraint_check_before_insert
  
  subroutine split_constraint_check_before_record (c, table, pl, n_loop, passed)
    class(split_constraint_t), intent(in) :: c
    class(ps_table_t), intent(in) :: table
    type(pdg_list_t), intent(in) :: pl
    integer, intent(in) :: n_loop
    logical, intent(out) :: passed
    passed = .true.
  end subroutine split_constraint_check_before_record
  
\eatline
\eatline
\eatline
\nwendcode{}\nwbegindocs{8}A transparent wrapper, so we can collect constraints of different type.
\nwenddocs{}\nwbegincode{9}\moddef{Auto components: types}\plusendmoddef
  type :: split_constraint_wrap_t
     class(split_constraint_t), allocatable :: c
  end type split_constraint_wrap_t
  
\eatline
\nwendcode{}\nwbegindocs{10}A collection of constraints.
\nwenddocs{}\nwbegincode{11}\moddef{Auto components: public}\endmoddef
  public :: split_constraints_t
\nwendcode{}\nwbegincode{12}\moddef{Auto components: types}\plusendmoddef
  type :: split_constraints_t
     class(split_constraint_wrap_t), dimension(:), allocatable :: cc
   contains
   \LA{}Auto components: split constraints: TBP\RA{}
  end type split_constraints_t
  
\eatline
\nwendcode{}\nwbegindocs{13}Initialize the constraints set with a specific number of elements.
\nwenddocs{}\nwbegincode{14}\moddef{Auto components: split constraints: TBP}\endmoddef
  procedure :: init => split_constraints_init
\nwendcode{}\nwbegincode{15}\moddef{Auto components: procedures}\plusendmoddef
  subroutine split_constraints_init (constraints, n)
    class(split_constraints_t), intent(out) :: constraints
    integer, intent(in) :: n
    allocate (constraints%cc (n))
  end subroutine split_constraints_init
  
\eatline
\nwendcode{}\nwbegindocs{16}Set a constraint.
\nwenddocs{}\nwbegincode{17}\moddef{Auto components: split constraints: TBP}\plusendmoddef
  procedure :: set => split_constraints_set
\nwendcode{}\nwbegincode{18}\moddef{Auto components: procedures}\plusendmoddef
  subroutine split_constraints_set (constraints, i, c)
    class(split_constraints_t), intent(inout) :: constraints
    integer, intent(in) :: i
    class(split_constraint_t), intent(in) :: c
    allocate (constraints%cc(i)%c, source = c)
  end subroutine split_constraints_set
  
\eatline
\nwendcode{}\nwbegindocs{19}Apply checks.

{\tt{}check{\char95}before{\char95}split} is applied to the particle list that we want
to split.

{\tt{}check{\char95}before{\char95}insert} is applied to the particle list {\tt{}pl} that is to
replace the particle {\tt{}pa} that is split.  This check may transform the
particle list.

{\tt{}check{\char95}before{\char95}record} is applied to the complete new particle list that
results from splitting before it is recorded.
\nwenddocs{}\nwbegincode{20}\moddef{Auto components: split constraints: TBP}\plusendmoddef
  procedure :: check_before_split  => split_constraints_check_before_split
  procedure :: check_before_insert => split_constraints_check_before_insert
  procedure :: check_before_record => split_constraints_check_before_record
\nwendcode{}\nwbegincode{21}\moddef{Auto components: procedures}\plusendmoddef
  subroutine split_constraints_check_before_split &
       (constraints, table, pl, k, passed)
    class(split_constraints_t), intent(in) :: constraints
    class(ps_table_t), intent(in) :: table
    type(pdg_list_t), intent(in) :: pl
    integer, intent(in) :: k
    logical, intent(out) :: passed
    integer :: i
    passed = .true.
    do i = 1, size (constraints%cc)
       call constraints%cc(i)%c%check_before_split (table, pl, k, passed)
       if (.not. passed)  return
    end do
  end subroutine split_constraints_check_before_split
    
  subroutine split_constraints_check_before_insert &
       (constraints, table, pa, pl, passed)
    class(split_constraints_t), intent(in) :: constraints
    class(ps_table_t), intent(in) :: table
    type(pdg_array_t), intent(in) :: pa
    type(pdg_list_t), intent(inout) :: pl
    logical, intent(out) :: passed
    integer :: i
    passed = .true.
    do i = 1, size (constraints%cc)
       call constraints%cc(i)%c%check_before_insert (table, pa, pl, passed)
       if (.not. passed)  return
    end do
  end subroutine split_constraints_check_before_insert
    
  subroutine split_constraints_check_before_record &
       (constraints, table, pl, n_loop, passed)
    class(split_constraints_t), intent(in) :: constraints
    class(ps_table_t), intent(in) :: table
    type(pdg_list_t), intent(in) :: pl
    integer, intent(in) :: n_loop
    logical, intent(out) :: passed
    integer :: i
    passed = .true.
    do i = 1, size (constraints%cc)
       call constraints%cc(i)%c%check_before_record (table, pl, n_loop, passed)
       if (.not. passed)  return
    end do
  end subroutine split_constraints_check_before_record
    
\eatline
\eatline
\eatline
\nwendcode{}\nwbegindocs{22}\nwdocspar
\subsection{Specific constraints}
\subsubsection{Number of particles}
Specific constraint: The number of particles plus the number of loops, if
any, must remain less than the given limit.  Note that the number of loops is
defined only when we are recording the entry.
\nwenddocs{}\nwbegincode{23}\moddef{Auto components: types}\plusendmoddef
  type, extends (split_constraint_t) :: constraint_n_tot
     private
     integer :: n_max = 0
   contains
     procedure :: check_before_split => constraint_n_tot_check_before_split
     procedure :: check_before_record => constraint_n_tot_check_before_record
  end type constraint_n_tot
  
\eatline
\nwendcode{}\nwbegincode{24}\moddef{Auto components: public}\plusendmoddef
  public :: constrain_n_tot
\nwendcode{}\nwbegincode{25}\moddef{Auto components: procedures}\plusendmoddef
  function constrain_n_tot (n_max) result (c)
    integer, intent(in) :: n_max
    type(constraint_n_tot) :: c
    c%n_max = n_max
  end function constrain_n_tot
  
  subroutine constraint_n_tot_check_before_split (c, table, pl, k, passed)
    class(constraint_n_tot), intent(in) :: c
    class(ps_table_t), intent(in) :: table
    type(pdg_list_t), intent(in) :: pl
    integer, intent(in) :: k
    logical, intent(out) :: passed
    passed = pl%get_size () < c%n_max
  end subroutine constraint_n_tot_check_before_split

  subroutine constraint_n_tot_check_before_record (c, table, pl, n_loop, passed)
    class(constraint_n_tot), intent(in) :: c
    class(ps_table_t), intent(in) :: table
    type(pdg_list_t), intent(in) :: pl
    integer, intent(in) :: n_loop
    logical, intent(out) :: passed
    passed = pl%get_size () + n_loop <= c%n_max
  end subroutine constraint_n_tot_check_before_record

\eatline
\eatline
\nwendcode{}\nwbegindocs{26}\nwdocspar
\subsubsection{Number of loops}
Specific constraint: The number of loops is limited, independent of the
total number of particles.
\nwenddocs{}\nwbegincode{27}\moddef{Auto components: types}\plusendmoddef
  type, extends (split_constraint_t) :: constraint_n_loop
     private
     integer :: n_loop_max = 0
   contains
     procedure :: check_before_record => constraint_n_loop_check_before_record
  end type constraint_n_loop
  
\eatline
\nwendcode{}\nwbegincode{28}\moddef{Auto components: public}\plusendmoddef
  public :: constrain_n_loop
\nwendcode{}\nwbegincode{29}\moddef{Auto components: procedures}\plusendmoddef
  function constrain_n_loop (n_loop_max) result (c)
    integer, intent(in) :: n_loop_max
    type(constraint_n_loop) :: c
    c%n_loop_max = n_loop_max
  end function constrain_n_loop

  subroutine constraint_n_loop_check_before_record &
       (c, table, pl, n_loop, passed)
    class(constraint_n_loop), intent(in) :: c
    class(ps_table_t), intent(in) :: table
    type(pdg_list_t), intent(in) :: pl
    integer, intent(in) :: n_loop
    logical, intent(out) :: passed
    passed = n_loop <= c%n_loop_max
  end subroutine constraint_n_loop_check_before_record

\eatline
\eatline
\nwendcode{}\nwbegindocs{30}\nwdocspar
\subsubsection{Particles allowed in splitting}
Specific constraint: The entries in the particle list ready for insertion
are matched to a given list of particle patterns.  If a match occurs, the
entry is replaced by the corresponding pattern.  If there is no match, the
check fails.
\nwenddocs{}\nwbegincode{31}\moddef{Auto components: types}\plusendmoddef
  type, extends (split_constraint_t) :: constraint_insert
     private
     type(pdg_list_t) :: pl_match
   contains
     procedure :: check_before_insert => constraint_insert_check_before_insert
  end type constraint_insert
  
\eatline
\nwendcode{}\nwbegincode{32}\moddef{Auto components: public}\plusendmoddef
  public :: constrain_insert
\nwendcode{}\nwbegincode{33}\moddef{Auto components: procedures}\plusendmoddef
  function constrain_insert (pl_match) result (c)
    type(pdg_list_t), intent(in) :: pl_match
    type(constraint_insert) :: c
    c%pl_match = pl_match
  end function constrain_insert
  
  subroutine constraint_insert_check_before_insert (c, table, pa, pl, passed)
    class(constraint_insert), intent(in) :: c
    class(ps_table_t), intent(in) :: table
    type(pdg_array_t), intent(in) :: pa
    type(pdg_list_t), intent(inout) :: pl
    logical, intent(out) :: passed
    call pl%match_replace (c%pl_match, passed)
  end subroutine constraint_insert_check_before_insert

\eatline
\eatline
\nwendcode{}\nwbegindocs{34}\nwdocspar
\subsubsection{Particles required in final state}
Specific constraint: The entries in the recorded state must be a superset of
the entries in the given list (for instance, the lowest-order
state).
\nwenddocs{}\nwbegincode{35}\moddef{Auto components: types}\plusendmoddef
  type, extends (split_constraint_t) :: constraint_require
     private
     type(pdg_list_t) :: pl
   contains
     procedure :: check_before_record => constraint_require_check_before_record
  end type constraint_require
  
\eatline
\nwendcode{}\nwbegindocs{36}We check the current state by matching all particle entries against the
stored particle list, and crossing out the particles in the latter list when a
match is found.  The constraint passed if all entries have been crossed out.

For an {\tt{}if{\char95}table} in particular, we check the final state only.
\nwenddocs{}\nwbegincode{37}\moddef{Auto components: public}\plusendmoddef
  public :: constrain_require
\nwendcode{}\nwbegincode{38}\moddef{Auto components: procedures}\plusendmoddef
  function constrain_require (pl) result (c)
    type(pdg_list_t), intent(in) :: pl
    type(constraint_require) :: c
    c%pl = pl
  end function constrain_require
  
  subroutine constraint_require_check_before_record &
       (c, table, pl, n_loop, passed)
    class(constraint_require), intent(in) :: c
    class(ps_table_t), intent(in) :: table
    type(pdg_list_t), intent(in) :: pl
    integer, intent(in) :: n_loop
    logical, intent(out) :: passed
    logical, dimension(:), allocatable :: mask
    integer :: i, k, n_in
    select type (table)
    type is (if_table_t)
       n_in = 2
    class default
       n_in = 0
    end select
    allocate (mask (c%pl%get_size ()), source = .true.)
    do i = n_in + 1, pl%get_size ()
       k = c%pl%find_match (pl%get (i), mask)
       if (k /= 0)  mask(k) = .false.
    end do
    passed = .not. any (mask)
  end subroutine constraint_require_check_before_record

\eatline
\eatline
\nwendcode{}\nwbegindocs{39}\nwdocspar
\subsubsection{Radiation}
Specific constraint: We have radiation pattern if the original particle
matches an entry in the list of particles that should replace it.  The
constraint prohibits this situation.
\nwenddocs{}\nwbegincode{40}\moddef{Auto components: public}\plusendmoddef
  public :: constrain_radiation
\nwendcode{}\nwbegincode{41}\moddef{Auto components: types}\plusendmoddef
  type, extends (split_constraint_t) :: constraint_radiation
     private
   contains
     procedure :: check_before_insert => &
          constraint_radiation_check_before_insert
  end type constraint_radiation
  
\eatline
\nwendcode{}\nwbegincode{42}\moddef{Auto components: procedures}\plusendmoddef
  function constrain_radiation () result (c)
    type(constraint_radiation) :: c
  end function constrain_radiation
  
  subroutine constraint_radiation_check_before_insert (c, table, pa, pl, passed)
    class(constraint_radiation), intent(in) :: c
    class(ps_table_t), intent(in) :: table
    type(pdg_array_t), intent(in) :: pa
    type(pdg_list_t), intent(inout) :: pl
    logical, intent(out) :: passed
    passed = .not. (pl .match. pa)
  end subroutine constraint_radiation_check_before_insert

\eatline
\eatline
\nwendcode{}\nwbegindocs{43}\nwdocspar
\subsubsection{Mass sum}
Specific constraint: The sum of masses within the particle list must
be smaller than a given limit.  For in/out state combinations, we
check initial and final state separately.

If we specify {\tt{}margin} in the initialization, the sum must be
strictly less than the limit minus the given margin (which may be
zero).  If not, equality is allowed.
\nwenddocs{}\nwbegincode{44}\moddef{Auto components: public}\plusendmoddef
  public :: constrain_mass_sum
\nwendcode{}\nwbegincode{45}\moddef{Auto components: types}\plusendmoddef
  type, extends (split_constraint_t) :: constraint_mass_sum
     private
     real(default) :: mass_limit = 0
     logical :: strictly_less = .false.
     real(default) :: margin = 0
   contains
     procedure :: check_before_record => constraint_mass_sum_check_before_record
  end type constraint_mass_sum
  
\eatline
\nwendcode{}\nwbegincode{46}\moddef{Auto components: procedures}\plusendmoddef
  function constrain_mass_sum (mass_limit, margin) result (c)
    real(default), intent(in) :: mass_limit
    real(default), intent(in), optional :: margin
    type(constraint_mass_sum) :: c
    c%mass_limit = mass_limit
    if (present (margin)) then
       c%strictly_less = .true.
       c%margin = margin
    end if
  end function constrain_mass_sum
  
  subroutine constraint_mass_sum_check_before_record &
       (c, table, pl, n_loop, passed)
    class(constraint_mass_sum), intent(in) :: c
    class(ps_table_t), intent(in) :: table
    type(pdg_list_t), intent(in) :: pl
    integer, intent(in) :: n_loop
    logical, intent(out) :: passed
    real(default) :: limit
    if (c%strictly_less) then
       limit = c%mass_limit - c%margin
       select type (table)
       type is (if_table_t)
          passed = mass_sum (pl, 1, 2, table%model) < limit  &
               .and. mass_sum (pl, 3, pl%get_size (), table%model) < limit
       class default
          passed = mass_sum (pl, 1, pl%get_size (), table%model) < limit
       end select
    else
       limit = c%mass_limit
       select type (table)
       type is (if_table_t)
          passed = mass_sum (pl, 1, 2, table%model) <= limit &
               .and. mass_sum (pl, 3, pl%get_size (), table%model) <= limit
       class default
          passed = mass_sum (pl, 1, pl%get_size (), table%model) <= limit
       end select
    end if
  end subroutine constraint_mass_sum_check_before_record

\eatline
\eatline
\nwendcode{}\nwbegindocs{47}\nwdocspar
\subsubsection{Initial state particles}
Specific constraint: The two incoming particles must both match the given
particle list.  This is checked for the generated particle list, just before
it is recorded.
\nwenddocs{}\nwbegincode{48}\moddef{Auto components: public}\plusendmoddef
  public :: constrain_in_state
\nwendcode{}\nwbegincode{49}\moddef{Auto components: types}\plusendmoddef
  type, extends (split_constraint_t) :: constraint_in_state
     private
     type(pdg_list_t) :: pl
   contains
     procedure :: check_before_record => constraint_in_state_check_before_record
  end type constraint_in_state

\eatline
\nwendcode{}\nwbegincode{50}\moddef{Auto components: procedures}\plusendmoddef
  function constrain_in_state (pl) result (c)
    type(pdg_list_t), intent(in) :: pl
    type(constraint_in_state) :: c
    c%pl = pl
  end function constrain_in_state

  subroutine constraint_in_state_check_before_record &
       (c, table, pl, n_loop, passed)
    class(constraint_in_state), intent(in) :: c
    class(ps_table_t), intent(in) :: table
    type(pdg_list_t), intent(in) :: pl
    integer, intent(in) :: n_loop
    logical, intent(out) :: passed
    integer :: i
    select type (table)
    type is (if_table_t)
       passed = .false.
       do i = 1, 2
          if (.not. (c%pl .match. pl%get (i)))  return
       end do
    end select
    passed = .true.
  end subroutine constraint_in_state_check_before_record
  
\eatline
\eatline
\nwendcode{}\nwbegindocs{51}\nwdocspar
\subsection{Tables of states}
Automatically generate a list of possible process components for a given
initial set (a single massive particle or a preset list of states).

The set of process components are generated by recursive splitting, applying
constraints on the fly that control and limit the process.  The generated
states are accumulated in a table that we can read out after completion.
\nwenddocs{}\nwbegincode{52}\moddef{Auto components: types}\plusendmoddef
  type, extends (pdg_list_t) :: ps_entry_t
     integer :: n_loop = 0
     integer :: n_rad = 0
     type(ps_entry_t), pointer :: previous => null ()
     type(ps_entry_t), pointer :: next => null ()
  end type ps_entry_t

\eatline
\nwendcode{}\nwbegindocs{53}\nwdocspar
This is the wrapper type for the decay tree for the list of final
states and the final array.  First, an abstract base type:
\nwenddocs{}\nwbegincode{54}\moddef{Auto components: public}\plusendmoddef
  public :: ps_table_t
\nwendcode{}\nwbegincode{55}\moddef{Auto components: types}\plusendmoddef
  type, abstract :: ps_table_t
     private
     class(model_data_t), pointer :: model => null ()
     logical :: loops = .false.
     type(ps_entry_t), pointer :: first => null ()
     type(ps_entry_t), pointer :: last => null ()
   contains
   \LA{}Auto components: ps table: TBP\RA{}
  end type ps_table_t
     
\eatline
\nwendcode{}\nwbegindocs{56}The extensions: one for decay, one for generic final states.  The decay-state
table stores the initial particle.  The final-state table is
indifferent, and the initial/final state table treats the first two
particles in its list as incoming antiparticles.
\nwenddocs{}\nwbegincode{57}\moddef{Auto components: public}\plusendmoddef
  public :: ds_table_t
  public :: if_table_t
\nwendcode{}\nwbegincode{58}\moddef{Auto components: types}\plusendmoddef
  type, extends (ps_table_t) :: ds_table_t
     private
     integer :: pdg_in = 0
   contains
   \LA{}Auto components: ds table: TBP\RA{}
  end type ds_table_t

  type, extends (ps_table_t) :: fs_table_t
   contains
   \LA{}Auto components: fs table: TBP\RA{}
  end type fs_table_t

  type, extends (fs_table_t) :: if_table_t
   contains
   \LA{}Auto components: if table: TBP\RA{}
  end type if_table_t

\eatline
\nwendcode{}\nwbegindocs{59}Finalizer: we must deallocate the embedded list.
\nwenddocs{}\nwbegincode{60}\moddef{Auto components: ps table: TBP}\endmoddef
  procedure :: final => ps_table_final
\nwendcode{}\nwbegincode{61}\moddef{Auto components: procedures}\plusendmoddef
  subroutine ps_table_final (object)
    class(ps_table_t), intent(inout) :: object
    type(ps_entry_t), pointer :: current
    do while (associated (object%first))
       current => object%first
       object%first => current%next
       deallocate (current)
    end do
    nullify (object%last)
  end subroutine ps_table_final
  
\eatline
\nwendcode{}\nwbegindocs{62}Write the table.  A base writer for the body and specific writers
for the headers.
\nwenddocs{}\nwbegincode{63}\moddef{Auto components: ps table: TBP}\plusendmoddef
  procedure :: base_write => ps_table_base_write
  procedure (ps_table_write), deferred :: write
\nwendcode{}\nwbegincode{64}\moddef{Auto components: interfaces}\endmoddef
  interface
     subroutine ps_table_write (object, unit)
       import
       class(ps_table_t), intent(in) :: object
       integer, intent(in), optional :: unit
     end subroutine ps_table_write
  end interface
\nwendcode{}\nwbegincode{65}\moddef{Auto components: ds table: TBP}\endmoddef
  procedure :: write => ds_table_write
\nwendcode{}\nwbegincode{66}\moddef{Auto components: fs table: TBP}\endmoddef
  procedure :: write => fs_table_write
\nwendcode{}\nwbegincode{67}\moddef{Auto components: if table: TBP}\endmoddef
  procedure :: write => if_table_write
\nwendcode{}\nwbegindocs{68}The first {\tt{}n{\char95}in} particles will be replaced by antiparticles in
the output, and we write an arrow if {\tt{}n{\char95}in} is present.
\nwenddocs{}\nwbegincode{69}\moddef{Auto components: procedures}\plusendmoddef
  subroutine ps_table_base_write (object, unit, n_in)
    class(ps_table_t), intent(in) :: object
    integer, intent(in), optional :: unit
    integer, intent(in), optional :: n_in
    integer, dimension(:), allocatable :: pdg
    type(ps_entry_t), pointer :: entry
    type(field_data_t), pointer :: prt
    integer :: u, i, j, n0
    u = given_output_unit (unit)
    entry => object%first
    do while (associated (entry))
       write (u, "(2x)", advance = "no")
       if (present (n_in)) then
          do i = 1, n_in
             write (u, "(1x)", advance = "no")
             pdg = entry%get (i)
             do j = 1, size (pdg)
                prt => object%model%get_field_ptr (pdg(j))
                if (j > 1)  write (u, "(':')", advance = "no")
                write (u, "(A)", advance = "no") &
                     char (prt%get_name (pdg(j) >= 0))
             end do
          end do
          write (u, "(1x,A)", advance = "no")  "=>"
          n0 = n_in + 1
       else
          n0 = 1
       end if
       do i = n0, entry%get_size ()
          write (u, "(1x)", advance = "no")
          pdg = entry%get (i)
          do j = 1, size (pdg)
             prt => object%model%get_field_ptr (pdg(j))
             if (j > 1)  write (u, "(':')", advance = "no")
             write (u, "(A)", advance = "no") &
                  char (prt%get_name (pdg(j) < 0))
          end do
       end do
       if (object%loops) then
          write (u, "(2x,'[',I0,',',I0,']')")  entry%n_loop, entry%n_rad
       else
          write (u, *)
       end if
       entry => entry%next
    end do
  end subroutine ps_table_base_write
          
  subroutine ds_table_write (object, unit)
    class(ds_table_t), intent(in) :: object
    integer, intent(in), optional :: unit
    type(field_data_t), pointer :: prt
    integer :: u
    u = given_output_unit (unit)
    prt => object%model%get_field_ptr (object%pdg_in)
    write (u, "(1x,A,1x,A)")  "Decays for particle:", &
         char (prt%get_name (object%pdg_in < 0))
    call object%base_write (u)
  end subroutine ds_table_write
          
  subroutine fs_table_write (object, unit)
    class(fs_table_t), intent(in) :: object
    integer, intent(in), optional :: unit
    integer :: u
    u = given_output_unit (unit)
    write (u, "(1x,A)")  "Table of final states:"
    call object%base_write (u)
  end subroutine fs_table_write
          
  subroutine if_table_write (object, unit)
    class(if_table_t), intent(in) :: object
    integer, intent(in), optional :: unit
    integer :: u
    u = given_output_unit (unit)
    write (u, "(1x,A)")  "Table of in/out states:"
    call object%base_write (u, n_in = 2)
  end subroutine if_table_write
          
\eatline
\nwendcode{}\nwbegindocs{70}Obtain a particle string for a given index in the pdg list
\nwenddocs{}\nwbegincode{71}\moddef{Auto components: ps table: TBP}\plusendmoddef
  procedure :: get_particle_string => ps_table_get_particle_string
\nwendcode{}\nwbegincode{72}\moddef{Auto components: procedures}\plusendmoddef
  subroutine ps_table_get_particle_string (object, index, n_in, prt_in, prt_out)
    class(ps_table_t), intent(in) :: object
    integer, intent(in) :: index
    integer, intent(in) :: n_in
    type(string_t), intent(out), dimension(:), allocatable :: prt_in, prt_out
    type(string_t) :: prt_tmp
    type(field_data_t), pointer :: prt
    type(ps_entry_t), pointer :: entry
    integer, dimension(:), allocatable :: pdg
    integer :: n0
    integer :: i, j 
    entry => object%first
    i = 1
    do while (i < index)
      if (associated (entry%next)) then
        entry => entry%next
        i=i+1
      else
        call msg_fatal ("ps_table: entry with requested index does not exist!")
      end if
    end do
    n0 = n_in + 1
    allocate (prt_in (n_in), prt_out (entry%get_size () - n_in))
    do i = 1, n_in
      prt_in(i) = ""
      pdg = entry%get(i)
      do j = 1, size(pdg)
        prt => object%model%get_field_ptr (pdg(j))
        prt_in(i) = prt_in(i) // prt%get_name (pdg(j) >= 0)
        if (j /= size(pdg)) &
           prt_in(i) = prt_in(i) // ":"
      end do
    end do
    do i = n0, entry%get_size ()
      prt_out(i-n_in) = ""
      pdg = entry%get(i)
      do j = 1, size(pdg)
         prt => object%model%get_field_ptr (pdg(j))
         prt_out(i-n_in) = prt_out(i-n_in) // prt%get_name (pdg(j) < 0)
         if (j /= size(pdg)) &
            prt_out(i-n_in) = prt_out(i-n_in) // ":"
      end do
    end do
  end subroutine ps_table_get_particle_string
                        
\eatline
\nwendcode{}\nwbegindocs{73}Initialize with a predefined set of final states, or in/out state lists.
\nwenddocs{}\nwbegincode{74}\moddef{Auto components: ps table: TBP}\plusendmoddef
  generic :: init => ps_table_init
  procedure, private :: ps_table_init
\nwendcode{}\nwbegincode{75}\moddef{Auto components: if table: TBP}\plusendmoddef
  generic :: init => if_table_init
  procedure, private :: if_table_init
\nwendcode{}\nwbegincode{76}\moddef{Auto components: procedures}\plusendmoddef
  subroutine ps_table_init (table, model, pl, constraints)
    class(ps_table_t), intent(out) :: table
    class(model_data_t), intent(in), target :: model
    type(pdg_list_t), dimension(:), intent(in) :: pl
    type(split_constraints_t), intent(in) :: constraints
    logical :: passed
    integer :: i
    table%model => model
    do i = 1, size (pl)
       call table%record (pl(i), 0, 0, constraints, passed)
       if (.not. passed) then
          call msg_fatal ("Registering process components: constraint failed")
       end if
    end do
  end subroutine ps_table_init
    
  subroutine if_table_init (table, model, pl_in, pl_out, constraints)
    class(if_table_t), intent(out) :: table
    class(model_data_t), intent(in), target :: model
    type(pdg_list_t), dimension(:), intent(in) :: pl_in, pl_out
    type(split_constraints_t), intent(in) :: constraints
    integer :: i, j, k, p, n_in, n_out
    type(pdg_array_t), dimension(:), allocatable :: pa_in
    type(pdg_list_t), dimension(:), allocatable :: pl
    allocate (pl (size (pl_in) * size (pl_out)))
    k = 0
    do i = 1, size (pl_in)
       n_in = pl_in(i)%get_size ()
       allocate (pa_in (n_in))
       do p = 1, n_in
          pa_in(p) = pl_in(i)%get (p)
       end do
       do j = 1, size (pl_out)
          n_out = pl_out(j)%get_size ()
          k = k + 1
          call pl(k)%init (n_in + n_out)
          do p = 1, n_in
             call pl(k)%set (p, invert_pdg_array (pa_in(p), model))
          end do
          do p = 1, n_out
             call pl(k)%set (n_in + p, pl_out(j)%get (p))
          end do
       end do
       deallocate (pa_in)
    end do
    call table%init (model, pl, constraints)
  end subroutine if_table_init
    
\eatline
\nwendcode{}\nwbegindocs{77}Enable loops for the table.  This affects both splitting and output.
\nwenddocs{}\nwbegincode{78}\moddef{Auto components: ps table: TBP}\plusendmoddef
  procedure :: enable_loops => ps_table_enable_loops
\nwendcode{}\nwbegincode{79}\moddef{Auto components: procedures}\plusendmoddef
  subroutine ps_table_enable_loops (table)
    class(ps_table_t), intent(inout) :: table
    table%loops = .true.
  end subroutine ps_table_enable_loops
    
\eatline
\nwendcode{}\nwbegindocs{80}\nwdocspar
\subsection{Top-level methods}
Create a table for a single-particle decay.  Construct all possible final
states from a single particle with PDG code {\tt{}pdg{\char95}in}.  The construction is
limited by the given {\tt{}constraints}.
\nwenddocs{}\nwbegincode{81}\moddef{Auto components: ds table: TBP}\plusendmoddef
  procedure :: make => ds_table_make
\nwendcode{}\nwbegincode{82}\moddef{Auto components: procedures}\plusendmoddef
  subroutine ds_table_make (table, model, pdg_in, constraints)
    class(ds_table_t), intent(out) :: table
    class(model_data_t), intent(in), target :: model
    integer, intent(in) :: pdg_in
    type(split_constraints_t), intent(in) :: constraints
    type(pdg_list_t) :: pl_in
    type(pdg_list_t), dimension(0) :: pl
    logical :: passed
    call table%init (model, pl, constraints)
    table%pdg_in = pdg_in
    call pl_in%init (1)
    call pl_in%set (1, [pdg_in])
    call table%split (pl_in, 0, constraints)
  end subroutine ds_table_make
    
\eatline
\nwendcode{}\nwbegindocs{83}Split all entries in a growing table, starting from a table that may already
contain states.  Add and record split states on the fly.
\nwenddocs{}\nwbegincode{84}\moddef{Auto components: fs table: TBP}\plusendmoddef
  procedure :: radiate => fs_table_radiate
\nwendcode{}\nwbegincode{85}\moddef{Auto components: procedures}\plusendmoddef
  subroutine fs_table_radiate (table, constraints)
    class(fs_table_t), intent(inout) :: table
    type(split_constraints_t) :: constraints
    type(ps_entry_t), pointer :: current
    integer :: i, s
    logical :: passed
    current => table%first
    do while (associated (current))
       call table%split (current, 0, constraints, record = .true.)
       current => current%next
    end do
  end subroutine fs_table_radiate

\eatline
\nwendcode{}\nwbegindocs{86}\nwdocspar
\subsection{Splitting algorithm}
Recursive splitting.  First of all, we record the current {\tt{}pdg{\char95}list} in
the table, subject to {\tt{}constraints}, if requested.  We also record copies of
the list marked as loop corrections.

When we record a particle list, we sort it first.

If there is room for splitting, We take a PDG array list and the index of an
element, and split this element in all possible ways.  The split entry is
inserted into the list, which we split further.

The recursion terminates whenever the split array would have a length
greater than $n_\text{max}$.
\nwenddocs{}\nwbegincode{87}\moddef{Auto components: ps table: TBP}\plusendmoddef
  procedure :: split => ps_table_split
\nwendcode{}\nwbegincode{88}\moddef{Auto components: procedures}\plusendmoddef
  recursive subroutine ps_table_split (table, pl, n_rad, constraints, &
        record)
    class(ps_table_t), intent(inout) :: table
    class(pdg_list_t), intent(in) :: pl
    integer, intent(in) :: n_rad
    type(split_constraints_t), intent(in) :: constraints
    logical, intent(in), optional :: record
    integer :: n_loop, i
    logical :: passed
    type(vertex_iterator_t) :: vit
    integer, dimension(:), allocatable :: pdg1
    integer, dimension(:), allocatable :: pdg2
    if (present (record)) then
       if (record) then
          n_loop = 0
          INCR_LOOPS: do
             call table%record_sorted (pl, n_loop, n_rad, constraints, passed)
             if (.not. passed)  exit INCR_LOOPS
             if (.not. table%loops)  exit INCR_LOOPS
             n_loop = n_loop + 1
          end do INCR_LOOPS
       end if
    end if
    do i = 1, pl%get_size ()
       call constraints%check_before_split (table, pl, i, passed)
       if (passed) then
          pdg1 = pl%get (i)
          call vit%init (table%model, pdg1)
          SCAN_VERTICES: do
             call vit%get_next_match (pdg2)
             if (allocated (pdg2)) then
                call table%insert (pl, n_rad, i, pdg2, constraints)
             else
                exit SCAN_VERTICES
             end if
          end do SCAN_VERTICES
       end if
    end do
  end subroutine ps_table_split
    
\eatline
\nwendcode{}\nwbegindocs{89}The worker part: insert the list of particles found by vertex matching in
place of entry {\tt{}i} in the PDG list.  Then split/record further.

The {\tt{}n{\char95}in} parameter tells the replacement routine to insert the new
particles after entry {\tt{}n{\char95}in}.  Otherwise, they follow index {\tt{}i}.
\nwenddocs{}\nwbegincode{90}\moddef{Auto components: ps table: TBP}\plusendmoddef
  procedure :: insert => ps_table_insert
\nwendcode{}\nwbegincode{91}\moddef{Auto components: procedures}\plusendmoddef
  recursive subroutine ps_table_insert &
       (table, pl, n_rad, i, pdg, constraints, n_in)
    class(ps_table_t), intent(inout) :: table
    class(pdg_list_t), intent(in) :: pl
    integer, intent(in) :: n_rad, i
    integer, dimension(:), intent(in) :: pdg
    type(split_constraints_t), intent(in) :: constraints
    integer, intent(in), optional :: n_in
    type(pdg_list_t) :: pl_insert
    logical :: passed
    integer :: k, s
    s = size (pdg)
    call pl_insert%init (s)
    do k = 1, s
       call pl_insert%set (k, pdg(k))
    end do
    call constraints%check_before_insert (table, pl%get (i), pl_insert, passed)
    if (passed) then
       call table%split (pl%replace (i, pl_insert, n_in), n_rad + s - 1, &
            constraints, record = .true.)
    end if
  end subroutine ps_table_insert
    
\eatline
\nwendcode{}\nwbegindocs{92}Special case:
If we are splitting an initial particle, there is slightly more to
do.  We loop over the particles from the vertex match and replace the
initial particle by each of them in turn.  The remaining particles
must be appended after the second initial particle, so they will end
up in the out state.  This is done by providing the {\tt{}n{\char95}in} argument
to the base method as an optional argument.

Note that we must call the base-method procedure explicitly, so the
{\tt{}table} argument keeps its dynamic type as {\tt{}if{\char95}table} inside this
procedure.
\nwenddocs{}\nwbegincode{93}\moddef{Auto components: if table: TBP}\plusendmoddef
  procedure :: insert => if_table_insert
\nwendcode{}\nwbegincode{94}\moddef{Auto components: procedures}\plusendmoddef
  recursive subroutine if_table_insert  &
       (table, pl, n_rad, i, pdg, constraints, n_in)
    class(if_table_t), intent(inout) :: table
    class(pdg_list_t), intent(in) :: pl
    integer, intent(in) :: n_rad, i
    integer, dimension(:), intent(in) :: pdg
    type(split_constraints_t), intent(in) :: constraints
    integer, intent(in), optional :: n_in
    integer, dimension(:), allocatable :: pdg_work
    integer :: p
    if (i > 2) then
       call ps_table_insert (table, pl, n_rad, i, pdg, constraints)
    else
       allocate (pdg_work (size (pdg)))
       do p = 1, size (pdg)
          pdg_work(1) = pdg(p)
          pdg_work(2:p) = pdg(1:p-1)
          pdg_work(p+1:) = pdg(p+1:)
          call ps_table_insert (table, &
               pl, n_rad, i, pdg_work, constraints, n_in = 2)
       end do
    end if
  end subroutine if_table_insert

\eatline
\nwendcode{}\nwbegindocs{95}Sort before recording.  In the case of the {\tt{}if{\char95}table}, we do not
sort the first {\tt{}n{\char95}in} particle entries.  Instead, we check whether they are
allowed in the {\tt{}pl{\char95}beam} PDG list, if that is provided.
\nwenddocs{}\nwbegincode{96}\moddef{Auto components: ps table: TBP}\plusendmoddef
  procedure :: record_sorted => ps_table_record_sorted
\nwendcode{}\nwbegincode{97}\moddef{Auto components: if table: TBP}\plusendmoddef
  procedure :: record_sorted => if_table_record_sorted
\nwendcode{}\nwbegincode{98}\moddef{Auto components: procedures}\plusendmoddef
  subroutine ps_table_record_sorted &
       (table, pl, n_loop, n_rad, constraints, passed)
    class(ps_table_t), intent(inout) :: table
    type(pdg_list_t), intent(in) :: pl
    integer, intent(in) :: n_loop, n_rad
    type(split_constraints_t), intent(in) :: constraints
    logical, intent(out) :: passed
    call table%record (pl%sort_abs (), n_loop, n_rad, constraints, passed)
  end subroutine ps_table_record_sorted
  
  subroutine if_table_record_sorted &
       (table, pl, n_loop, n_rad, constraints, passed)
    class(if_table_t), intent(inout) :: table
    type(pdg_list_t), intent(in) :: pl
    integer, intent(in) :: n_loop, n_rad
    type(split_constraints_t), intent(in) :: constraints
    logical, intent(out) :: passed
    integer :: i
    call table%record (pl%sort_abs (2), n_loop, n_rad, constraints, passed)
  end subroutine if_table_record_sorted

\eatline
\nwendcode{}\nwbegindocs{99}Record an entry: insert into the list.  Check the ordering and
insert it at the correct place, unless it is already there.

We record an array only if its mass sum is less than the total
available energy.  This restriction is removed by setting
{\tt{}constrained} to false.
\nwenddocs{}\nwbegincode{100}\moddef{Auto components: ps table: TBP}\plusendmoddef
  procedure :: record => ps_table_record
\nwendcode{}\nwbegincode{101}\moddef{Auto components: procedures}\plusendmoddef
  subroutine ps_table_record (table, pl, n_loop, n_rad, constraints, passed)
    class(ps_table_t), intent(inout) :: table
    type(pdg_list_t), intent(in) :: pl
    integer, intent(in) :: n_loop, n_rad
    type(split_constraints_t), intent(in) :: constraints
    logical, intent(out) :: passed
    type(ps_entry_t), pointer :: current
    logical :: mass_limit
    if (.not. pl%is_regular ())  return
    call constraints%check_before_record (table, pl, n_loop, passed)
    if (.not. passed)  return
    current => table%first
    do while (associated (current))
       if (pl == current) then
          if (n_loop == current%n_loop)  return
       else if (pl < current) then
          call insert
          return
       end if
       current => current%next
    end do
    call insert
  contains
    subroutine insert ()
      type(ps_entry_t), pointer :: entry
      allocate (entry)
      entry%pdg_list_t = pl
      entry%n_loop = n_loop
      entry%n_rad = n_rad
      if (associated (current)) then
         if (associated (current%previous)) then
            current%previous%next => entry
            entry%previous => current%previous
         else
            table%first => entry
         end if
         entry%next => current
         current%previous => entry
      else
         if (associated (table%last)) then
            table%last%next => entry
            entry%previous => table%last
         else
            table%first => entry
         end if
         table%last => entry
      end if
    end subroutine insert
  end subroutine ps_table_record
    
\eatline
\nwendcode{}\nwbegindocs{102}\nwdocspar
\subsection{Tools}
Compute the mass sum for a PDG list object, counting the entries with indices
between (including) {\tt{}n1} and {\tt{}n2}.  Rely on the requirement that
if an entry is a PDG array, this array must be degenerate in mass.
\nwenddocs{}\nwbegincode{103}\moddef{Auto components: procedures}\plusendmoddef
  function mass_sum (pl, n1, n2, model) result (m)
    type(pdg_list_t), intent(in) :: pl
    integer, intent(in) :: n1, n2
    class(model_data_t), intent(in), target :: model
    integer, dimension(:), allocatable :: pdg
    real(default) :: m
    type(field_data_t), pointer :: prt
    integer :: i
    m = 0
    do i = n1, n2
       pdg = pl%get (i)
       prt => model%get_field_ptr (pdg(1))
       m = m + prt%get_mass ()
    end do
  end function mass_sum
  
\eatline
\nwendcode{}\nwbegindocs{104}Invert a PDG array, replacing particles by antiparticles.  This
depends on the model.
\nwenddocs{}\nwbegincode{105}\moddef{Auto components: procedures}\plusendmoddef
  function invert_pdg_array (pa, model) result (pa_inv)
    type(pdg_array_t), intent(in) :: pa
    class(model_data_t), intent(in), target :: model
    type(pdg_array_t) :: pa_inv
    type(field_data_t), pointer :: prt
    integer :: i, pdg
    pa_inv = pa
    do i = 1, pa_inv%get_length ()
       pdg = pa_inv%get (i)
       prt => model%get_field_ptr (pdg)
       if (prt%has_antiparticle ())  call pa_inv%set (i, -pdg)
    end do
  end function invert_pdg_array
          
\eatline
\nwendcode{}\nwbegindocs{106}\nwdocspar
\subsection{Access results}
Return the number of generated decays.
\nwenddocs{}\nwbegincode{107}\moddef{Auto components: ps table: TBP}\plusendmoddef
  procedure :: get_length => ps_table_get_length
\nwendcode{}\nwbegincode{108}\moddef{Auto components: procedures}\plusendmoddef
  function ps_table_get_length (ps_table) result (n)
    class(ps_table_t), intent(in) :: ps_table
    integer :: n
    type(ps_entry_t), pointer :: entry
    n = 0
    entry => ps_table%first
    do while (associated (entry))
       n = n + 1
       entry => entry%next
    end do
  end function ps_table_get_length

\eatline
\nwendcode{}\nwbegindocs{109}Return an allocated array of decay products (PDG codes).  If
requested, return also the loop and radiation order count.
\nwenddocs{}\nwbegincode{110}\moddef{Auto components: ps table: TBP}\plusendmoddef
  procedure :: get_pdg_out => ps_table_get_pdg_out
\nwendcode{}\nwbegincode{111}\moddef{Auto components: procedures}\plusendmoddef
  subroutine ps_table_get_pdg_out (ps_table, i, pa_out, n_loop, n_rad)
    class(ps_table_t), intent(in) :: ps_table
    integer, intent(in) :: i
    type(pdg_array_t), dimension(:), allocatable, intent(out) :: pa_out
    integer, intent(out), optional :: n_loop, n_rad
    type(ps_entry_t), pointer :: entry
    integer :: n, j
    n = 0
    entry => ps_table%first
    FIND_ENTRY: do while (associated (entry))
       n = n + 1
       if (n == i) then
          allocate (pa_out (entry%get_size ()))
          do j = 1, entry%get_size ()
             pa_out(j) = entry%get (j)
             if (present (n_loop))  n_loop = entry%n_loop
             if (present (n_rad))  n_rad = entry%n_rad
          end do
          exit FIND_ENTRY
       end if
       entry => entry%next
    end do FIND_ENTRY
  end subroutine ps_table_get_pdg_out
  
\eatline
\nwendcode{}\nwbegindocs{112}\nwdocspar
\subsection{Unit tests}
\nwenddocs{}\nwbegincode{113}\moddef{Auto components: public}\plusendmoddef
  public :: auto_components_test
\nwendcode{}\nwbegincode{114}\moddef{Auto components: tests}\endmoddef
  subroutine auto_components_test (u, results)
    integer, intent(in) :: u
    type(test_results_t), intent(inout) :: results
  \LA{}Auto components: execute tests\RA{}
  end subroutine auto_components_test

\eatline
\nwendcode{}\nwbegindocs{115}\nwdocspar
\subsubsection{Generate Decay Table}
Determine all kinematically allowed decay channels for a Higgs boson,
using default parameter values.
\nwenddocs{}\nwbegincode{116}\moddef{Auto components: execute tests}\endmoddef
  call test (auto_components_1, "auto_components_1", &
       "generate decay table", &
       u, results)
\nwendcode{}\nwbegincode{117}\moddef{Auto components: tests}\plusendmoddef
  subroutine auto_components_1 (u)
    integer, intent(in) :: u
    class(model_data_t), pointer :: model
    type(field_data_t), pointer :: prt
    type(ds_table_t) :: ds_table
    type(split_constraints_t) :: constraints

    write (u, "(A)")  "* Test output: auto_components_1"
    write (u, "(A)")  "*   Purpose: determine Higgs decay table"
    write (u, *)

    write (u, "(A)")  "* Read Standard Model"

    call prepare_model (model, var_str ("SM"))

    prt => model%get_field_ptr (25)

    write (u, *)
    write (u, "(A)")  "* Higgs decays n = 2"
    write (u, *)

    call constraints%init (2)
    call constraints%set (1, constrain_n_tot (2))
    call constraints%set (2, constrain_mass_sum (prt%get_mass ()))

    call ds_table%make (model, 25, constraints)
    call ds_table%write (u)
    call ds_table%final ()

    write (u, *)
    write (u, "(A)")  "* Higgs decays n = 3 (w/o radiative)"
    write (u, *)

    call constraints%init (3)
    call constraints%set (1, constrain_n_tot (3))
    call constraints%set (2, constrain_mass_sum (prt%get_mass ()))
    call constraints%set (3, constrain_radiation ())

    call ds_table%make (model, 25, constraints)
    call ds_table%write (u)
    call ds_table%final ()

    write (u, *)
    write (u, "(A)")  "* Higgs decays n = 3 (w/ radiative)"
    write (u, *)

    call constraints%init (2)
    call constraints%set (1, constrain_n_tot (3))
    call constraints%set (2, constrain_mass_sum (prt%get_mass ()))

    call ds_table%make (model, 25, constraints)
    call ds_table%write (u)
    call ds_table%final ()

    write (u, *)
    write (u, "(A)")  "* Cleanup"
    
    call cleanup_model (model)

    write (u, *)
    write (u, "(A)")  "* Test output end: auto_components_1"

  end subroutine auto_components_1

\eatline
\nwendcode{}\nwbegindocs{118}\nwdocspar
\subsubsection{Generate radiation}
Given a final state, add radiation (NLO and NNLO).  We provide a list
of particles that is allowed to occur in the generated final states.
\nwenddocs{}\nwbegincode{119}\moddef{Auto components: execute tests}\plusendmoddef
  call test (auto_components_2, "auto_components_2", &
       "generate NLO corrections, final state", &
       u, results)
\nwendcode{}\nwbegincode{120}\moddef{Auto components: tests}\plusendmoddef
  subroutine auto_components_2 (u)
    integer, intent(in) :: u
    class(model_data_t), pointer :: model
    type(pdg_list_t), dimension(:), allocatable :: pl, pl_zzh
    type(pdg_list_t) :: pl_match
    type(fs_table_t) :: fs_table
    type(split_constraints_t) :: constraints
    real(default) :: sqrts
    integer :: i

    write (u, "(A)")  "* Test output: auto_components_2"
    write (u, "(A)")  "*   Purpose: generate radiation (NLO)"
    write (u, *)


    write (u, "(A)")  "* Read Standard Model"

    call prepare_model (model, var_str ("SM"))

    write (u, *)
    write (u, "(A)")  "* LO final state"
    write (u, *)
    
    allocate (pl (2))
    call pl(1)%init (2)
    call pl(1)%set (1, 1)
    call pl(1)%set (2, -1)
    call pl(2)%init (2)
    call pl(2)%set (1, 21)
    call pl(2)%set (2, 21)
    do i = 1, 2
       call pl(i)%write (u);  write (u, *)
    end do

    write (u, *)
    write (u, "(A)")  "* Initialize FS table"
    write (u, *)
    
    call constraints%init (1)
    call constraints%set (1, constrain_n_tot (3))

    call fs_table%init (model, pl, constraints)
    call fs_table%write (u)
    
    write (u, *)
    write (u, "(A)")  "* Generate NLO corrections, unconstrained"
    write (u, *)
    
    call fs_table%radiate (constraints)
    call fs_table%write (u)
    call fs_table%final ()
    
    write (u, *)
    write (u, "(A)")  "* Generate NLO corrections, &
         &complete but mass-constrained"
    write (u, *)
    
    sqrts = 50

    call constraints%init (2)
    call constraints%set (1, constrain_n_tot (3))
    call constraints%set (2, constrain_mass_sum (sqrts))

    call fs_table%init (model, pl, constraints)
    call fs_table%radiate (constraints)
    call fs_table%write (u)
    call fs_table%final ()
    
    write (u, *)
    write (u, "(A)")  "* Generate NLO corrections, restricted"
    write (u, *)
    
    call pl_match%init ([1, -1, 21])
    
    call constraints%init (2)
    call constraints%set (1, constrain_n_tot (3))
    call constraints%set (2, constrain_insert (pl_match))

    call fs_table%init (model, pl, constraints)
    call fs_table%radiate (constraints)
    call fs_table%write (u)
    call fs_table%final ()
    
    write (u, *)
    write (u, "(A)")  "* Generate NNLO corrections, restricted, with one loop"
    write (u, *)
    
    call pl_match%init ([1, -1, 21])
    
    call constraints%init (3)
    call constraints%set (1, constrain_n_tot (4))
    call constraints%set (2, constrain_n_loop (1))
    call constraints%set (3, constrain_insert (pl_match))

    call fs_table%init (model, pl, constraints)
    call fs_table%enable_loops ()
    call fs_table%radiate (constraints)
    call fs_table%write (u)
    call fs_table%final ()
    
    write (u, *)
    write (u, "(A)")  "* Generate NNLO corrections, restricted, with loops"
    write (u, *)
    
    call constraints%init (2)
    call constraints%set (1, constrain_n_tot (4))
    call constraints%set (2, constrain_insert (pl_match))

    call fs_table%init (model, pl, constraints)
    call fs_table%enable_loops ()
    call fs_table%radiate (constraints)
    call fs_table%write (u)
    call fs_table%final ()
    
    write (u, *)
    write (u, "(A)")  "* Generate NNLO corrections, restricted, to Z Z H, &
         &no loops"
    write (u, *)
    
    allocate (pl_zzh (1))
    call pl_zzh(1)%init (3)
    call pl_zzh(1)%set (1, 23)
    call pl_zzh(1)%set (2, 23)
    call pl_zzh(1)%set (3, 25)

    call constraints%init (3)
    call constraints%set (1, constrain_n_tot (5))
    call constraints%set (2, constrain_mass_sum (500._default))
    call constraints%set (3, constrain_require (pl_zzh(1)))

    call fs_table%init (model, pl_zzh, constraints)
    call fs_table%radiate (constraints)
    call fs_table%write (u)
    call fs_table%final ()
    
    call cleanup_model (model)

    write (u, *)
    write (u, "(A)")  "* Test output end: auto_components_2"

  end subroutine auto_components_2

\eatline
\nwendcode{}\nwbegindocs{121}\nwdocspar
\subsubsection{Generate radiation from initial and final state}
Given a process, add radiation (NLO and NNLO).  We provide a list
of particles that is allowed to occur in the generated final states.
\nwenddocs{}\nwbegincode{122}\moddef{Auto components: execute tests}\plusendmoddef
  call test (auto_components_3, "auto_components_3", &
       "generate NLO corrections, in and out", &
       u, results)
\nwendcode{}\nwbegincode{123}\moddef{Auto components: tests}\plusendmoddef
  subroutine auto_components_3 (u)
    integer, intent(in) :: u
    class(model_data_t), pointer :: model
    type(pdg_list_t), dimension(:), allocatable :: pl_in, pl_out
    type(pdg_list_t) :: pl_match, pl_beam
    type(if_table_t) :: if_table
    type(split_constraints_t) :: constraints
    real(default) :: sqrts
    integer :: i

    write (u, "(A)")  "* Test output: auto_components_3"
    write (u, "(A)")  "*   Purpose: generate radiation (NLO)"
    write (u, *)

    write (u, "(A)")  "* Read Standard Model"

    call prepare_model (model, var_str ("SM"))

    write (u, *)
    write (u, "(A)")  "* LO initial state"
    write (u, *)
    
    allocate (pl_in (2))
    call pl_in(1)%init (2)
    call pl_in(1)%set (1, 1)
    call pl_in(1)%set (2, -1)
    call pl_in(2)%init (2)
    call pl_in(2)%set (1, -1)
    call pl_in(2)%set (2, 1)
    do i = 1, 2
       call pl_in(i)%write (u);  write (u, *)
    end do

    write (u, *)
    write (u, "(A)")  "* LO final state"
    write (u, *)
    
    allocate (pl_out (1))
    call pl_out(1)%init (1)
    call pl_out(1)%set (1, 23)
    call pl_out(1)%write (u);  write (u, *)

    write (u, *)
    write (u, "(A)")  "* Initialize FS table"
    write (u, *)
    
    call constraints%init (1)
    call constraints%set (1, constrain_n_tot (4))
    
    call if_table%init (model, pl_in, pl_out, constraints)
    call if_table%write (u)
    
    write (u, *)
    write (u, "(A)")  "* Generate NLO corrections, unconstrained"
    write (u, *)
    
    call if_table%radiate (constraints)
    call if_table%write (u)
    call if_table%final ()
    
    write (u, *)
    write (u, "(A)")  "* Generate NLO corrections, &
         &complete but mass-constrained"
    write (u, *)
    
    sqrts = 100
    call constraints%init (2)
    call constraints%set (1, constrain_n_tot (4))
    call constraints%set (2, constrain_mass_sum (sqrts))
    
    call if_table%init (model, pl_in, pl_out, constraints)
    call if_table%radiate (constraints)
    call if_table%write (u)
    call if_table%final ()
    
    write (u, *)
    write (u, "(A)")  "* Generate NLO corrections, &
         &mass-constrained, restricted beams"
    write (u, *)
    
    call pl_beam%init (3)
    call pl_beam%set (1, 1)
    call pl_beam%set (2, -1)
    call pl_beam%set (3, 21)

    call constraints%init (3)
    call constraints%set (1, constrain_n_tot (4))
    call constraints%set (2, constrain_in_state (pl_beam))
    call constraints%set (3, constrain_mass_sum (sqrts))

    call if_table%init (model, pl_in, pl_out, constraints)
    call if_table%radiate (constraints)
    call if_table%write (u)
    call if_table%final ()
    
    write (u, *)
    write (u, "(A)")  "* Generate NLO corrections, restricted"
    write (u, *)
    
    call pl_match%init ([1, -1, 21])
    
    call constraints%init (4)
    call constraints%set (1, constrain_n_tot (4))
    call constraints%set (2, constrain_in_state (pl_beam))
    call constraints%set (3, constrain_mass_sum (sqrts))
    call constraints%set (4, constrain_insert (pl_match)) 

    call if_table%init (model, pl_in, pl_out, constraints)
    call if_table%radiate (constraints)
    call if_table%write (u)
    call if_table%final ()
    
    write (u, *)
    write (u, "(A)")  "* Generate NNLO corrections, restricted, Z preserved, &
         &with loops"
    write (u, *)
    
    call constraints%init (5)
    call constraints%set (1, constrain_n_tot (5))
    call constraints%set (2, constrain_in_state (pl_beam))
    call constraints%set (3, constrain_mass_sum (sqrts))
    call constraints%set (4, constrain_insert (pl_match)) 
    call constraints%set (5, constrain_require (pl_out(1)))

    call if_table%init (model, pl_in, pl_out, constraints)
    call if_table%enable_loops ()
    call if_table%radiate (constraints)
    call if_table%write (u)
    call if_table%final ()
    
    call cleanup_model (model)

    write (u, *)
    write (u, "(A)")  "* Test output end: auto_components_3"

  end subroutine auto_components_3

\eatline
\nwendcode{}\nwbegindocs{124}\nwdocspar
%%%%%%%%%%%%%%%%%%%%%%%%%%%%%%%%%%%%%%%%%%%%%%%%%%%%%%%%%%%%%%%%%%%%%%%%
\section{Creating the real flavor structure}
\nwenddocs{}\nwbegincode{125}\moddef{\code{}radiation{\char95}generator.f90\edoc{}}\endmoddef
\LA{}File header\RA{}

module radiation_generator

\LA{}Use kinds\RA{}
\LA{}Use strings\RA{}
  use diagnostics
  use os_interface
  use pdg_arrays
  use particle_specifiers
  use model_data
  use auto_components
  use physics_defs

  implicit none
  private

\LA{}radiation generator: public\RA{}
  
\LA{}radiation generator: types\RA{}

\LA{}radiation generator: interfaces\RA{}

contains

\LA{}radiation generator: procedures\RA{}

end module radiation_generator
\eatline
\nwendcode{}\nwbegindocs{126}\nwdocspar
\nwenddocs{}\nwbegincode{127}\moddef{radiation generator: types}\endmoddef
  type :: pdg_sorter_t
     integer :: pdg
  end type pdg_sorter_t

  type, extends (pdg_sorter_t) :: pdg_sorter_born_t
     logical :: checked
  end type pdg_sorter_born_t

  type, extends (pdg_sorter_t) :: pdg_sorter_real_t
     integer :: associated_born 
  end type pdg_sorter_real_t

\eatline
\nwendcode{}\nwbegindocs{128}\nwdocspar
\nwenddocs{}\nwbegincode{129}\moddef{radiation generator: types}\plusendmoddef
  type :: pdg_states_t
    type(pdg_array_t), dimension(:), allocatable :: pdg
    type(pdg_states_t), pointer :: next
    integer :: n_particles
  contains
  \LA{}radiation generator: pdg states: TBP\RA{}
  end type pdg_states_t

\eatline
\nwendcode{}\nwbegindocs{130}\nwdocspar
\nwenddocs{}\nwbegincode{131}\moddef{radiation generator: public}\endmoddef
  public :: radiation_generator_init
\nwendcode{}\nwbegincode{132}\moddef{radiation generator: interfaces}\endmoddef
  interface radiation_generator_init
     module procedure radiation_generator_init_pdg_list
     module procedure radiation_generator_init_pdg_array
  end interface

\eatline
\nwendcode{}\nwbegindocs{133}\nwdocspar
\nwenddocs{}\nwbegincode{134}\moddef{radiation generator: public}\plusendmoddef
  public :: radiation_generator_t
\nwendcode{}\nwbegincode{135}\moddef{radiation generator: types}\plusendmoddef
  type :: radiation_generator_t
    logical :: qcd_enabled = .false.
    logical :: qed_enabled = .false.
    logical :: is_gluon = .false.
    logical :: fs_gluon = .false.
    type(pdg_list_t) :: pl_in, pl_out
    type(split_constraints_t) :: constraints
    integer :: n_tot
    integer :: n_out
    integer :: n_loops
    integer :: n_light_quarks
    real(default) :: mass_sum
    class(model_data_t), pointer :: radiation_model
    type(pdg_states_t) :: pdg_raw
    type(pdg_array_t), dimension(:), allocatable :: pdg_in_born, pdg_out_born
  contains
  \LA{}radiation generator: radiation generator: TBP\RA{}
  end type radiation_generator_t

\nwendcode{}\nwbegindocs{136}\nwdocspar
%
\nwenddocs{}\nwbegincode{137}\moddef{radiation generator: pdg states: TBP}\endmoddef
  procedure :: init => pdg_states_init
\nwendcode{}\nwbegincode{138}\moddef{radiation generator: procedures}\endmoddef
  subroutine pdg_states_init (states)
    class(pdg_states_t), intent(inout) :: states
    nullify (states%next)
  end subroutine pdg_states_init

\eatline
\nwendcode{}\nwbegindocs{139}\nwdocspar
\nwenddocs{}\nwbegincode{140}\moddef{radiation generator: pdg states: TBP}\plusendmoddef
  procedure :: add => pdg_states_add
\nwendcode{}\nwbegincode{141}\moddef{radiation generator: procedures}\plusendmoddef
  subroutine pdg_states_add (states, pdg)
    class(pdg_states_t), intent(inout), target :: states
    type(pdg_array_t), dimension(:), intent(in) :: pdg
    type(pdg_states_t), pointer :: current_state
    select type (states)
    type is (pdg_states_t)
      current_state => states
      do
        if (associated (current_state%next)) then
          current_state => current_state%next
        else
          allocate (current_state%next)
          nullify(current_state%next%next)
          current_state%pdg = pdg
          exit
        end if
      end do
    end select
  end subroutine pdg_states_add

\eatline
\nwendcode{}\nwbegindocs{142}\nwdocspar
\nwenddocs{}\nwbegincode{143}\moddef{radiation generator: pdg states: TBP}\plusendmoddef
  procedure :: get_n_states => pdg_states_get_n_states
\nwendcode{}\nwbegincode{144}\moddef{radiation generator: procedures}\plusendmoddef
  function pdg_states_get_n_states (states) result (n)
    class(pdg_states_t), intent(in), target :: states
    integer :: n
    type(pdg_states_t), pointer :: current_state
    n = 0
    select type(states)
    type is (pdg_states_t)
      current_state => states
      do
        if (associated (current_state%next)) then
          n = n+1
          current_state => current_state%next
        else
          exit
        end if
      end do
    end select
  end function pdg_states_get_n_states

\eatline
\nwendcode{}\nwbegindocs{145}\nwdocspar
\nwenddocs{}\nwbegincode{146}\moddef{radiation generator: procedures}\plusendmoddef
  subroutine radiation_generator_init_pdg_list &
       (generator, qcd, qed, pl_in, pl_out)
    type(radiation_generator_t), intent(inout) :: generator
    logical, intent(in), optional :: qcd, qed
    type(pdg_list_t), intent(in) :: pl_in, pl_out
    if (present (qcd))  generator%qcd_enabled = qcd
    if (present (qed))  generator%qed_enabled = qed
    generator%pl_in = pl_in
    generator%pl_out = pl_out
    generator%is_gluon = pl_in%search_for_particle (GLUON)
    generator%fs_gluon = pl_out%search_for_particle (GLUON)
    generator%mass_sum = 0._default
    call generator%pdg_raw%init ()
  end subroutine radiation_generator_init_pdg_list

  subroutine radiation_generator_init_pdg_array &
       (generator, qcd, qed, pdg_in, pdg_out)
    type(radiation_generator_t), intent(inout) :: generator
    logical, intent(in), optional :: qcd, qed
    type(pdg_array_t), intent(in), dimension(:), allocatable :: pdg_in, pdg_out
    type(pdg_list_t) :: pl_in, pl_out
    integer :: i
    call pl_in%init(size (pdg_in))    
    call pl_out%init(size (pdg_out))
    do i = 1, size (pdg_in)
       call pl_in%set (i, pdg_in(i))
    end do
    do i = 1, size (pdg_out)
       call pl_out%set (i, pdg_out(i))
    end do
    call radiation_generator_init (generator, qcd, qed, pl_in, pl_out)
  end subroutine radiation_generator_init_pdg_array

\eatline
\nwendcode{}\nwbegindocs{147}\nwdocspar
\nwenddocs{}\nwbegincode{148}\moddef{radiation generator: radiation generator: TBP}\endmoddef
  procedure :: init_radiation_model => &
                    radiation_generator_init_radiation_model
\nwendcode{}\nwbegincode{149}\moddef{radiation generator: procedures}\plusendmoddef
  subroutine radiation_generator_init_radiation_model (generator, model)
    class(radiation_generator_t), intent(inout) :: generator
    class(model_data_t), intent(in), target :: model
    generator%radiation_model => model
  end subroutine radiation_generator_init_radiation_model

\eatline
\nwendcode{}\nwbegindocs{150}\nwdocspar
\nwenddocs{}\nwbegincode{151}\moddef{radiation generator: radiation generator: TBP}\plusendmoddef
  procedure :: set_n => radiation_generator_set_n
\nwendcode{}\nwbegincode{152}\moddef{radiation generator: procedures}\plusendmoddef
  subroutine radiation_generator_set_n (generator, n_in, n_out, n_loops)
    class(radiation_generator_t), intent(inout) :: generator
    integer, intent(in) :: n_in, n_out, n_loops
    generator%n_tot = n_in + n_out + 1
    generator%n_out = n_out
    generator%n_loops = n_loops
  end subroutine radiation_generator_set_n

\eatline
\nwendcode{}\nwbegindocs{153}\nwdocspar
\nwenddocs{}\nwbegincode{154}\moddef{radiation generator: radiation generator: TBP}\plusendmoddef
  procedure :: set_constraints => radiation_generator_set_constraints
\nwendcode{}\nwbegincode{155}\moddef{radiation generator: procedures}\plusendmoddef
  subroutine radiation_generator_set_constraints &
       (generator, set_n_loop, set_mass_sum, &
        set_selected_particles, set_required_particles)
    class(radiation_generator_t), intent(inout), target :: generator
    logical, intent(in) :: set_n_loop   
    logical, intent(in) :: set_mass_sum
    logical, intent(in) :: set_selected_particles
    logical, intent(in) :: set_required_particles
    integer :: i, n, n_constraints
    type(pdg_list_t) :: pl_req, pl_insert
    type(pdg_array_t) :: pdg_gluon, pdg_photon
    type(pdg_array_t) :: pdg_add, pdg_tmp
    integer :: i_skip
    i_skip = -1
    
    n_constraints = 1 + count([set_n_loop, set_mass_sum, &
         set_selected_particles, set_required_particles])
    associate (constraints => generator%constraints)
      n = 1
      call constraints%init (n_constraints)
      call constraints%set (n, constrain_n_tot (generator%n_tot))
      n = n+1
      if (set_n_loop) then
         call constraints%set (n, constrain_n_loop(generator%n_loops))
         n = n+1
      end if 
      if (set_mass_sum) then
        call constraints%set (n, constrain_mass_sum(generator%mass_sum))
        n = n+1
      end if
      if (set_required_particles) then
        if (generator%fs_gluon) then
           do i = 1, generator%n_out
              pdg_tmp = generator%pl_out%get(i)
              if (pdg_tmp%search_for_particle (GLUON)) then
                 i_skip = i
                 exit
              end if
           end do
           call pl_req%init (generator%n_out-1)
        else
           call pl_req%init (generator%n_out)
        end if
        do i = 1, generator%n_out
           if (i == i_skip) cycle
           call pl_req%set (i, generator%pl_out%get(i))
        end do          
        call constraints%set (n, constrain_require (pl_req))
        n = n+1
      end if
      if (set_selected_particles) then
        call pl_insert%init (generator%n_out+1)
        do i = 1, generator%n_out
           call pl_insert%set(i, generator%pl_out%get(i))
        end do
        pdg_gluon = GLUON; pdg_photon = PHOTON
        if (generator%fs_gluon .and. generator%qcd_enabled) then
           do i = 1, generator%n_light_quarks
              pdg_tmp = i
              pdg_add = pdg_add // pdg_tmp
           end do
        end if
        if (generator%qcd_enabled) pdg_add = pdg_add // pdg_gluon
        if (generator%qed_enabled) pdg_add = pdg_add // pdg_photon
        call pl_insert%set (generator%n_out+1, pdg_add)
        call constraints%set (n, constrain_insert (pl_insert))
      end if
    end associate

  end subroutine radiation_generator_set_constraints

\eatline
\nwendcode{}\nwbegindocs{156}\nwdocspar
\nwenddocs{}\nwbegincode{157}\moddef{radiation generator: radiation generator: TBP}\plusendmoddef
  procedure :: generate => radiation_generator_generate
\nwendcode{}\nwbegincode{158}\moddef{radiation generator: procedures}\plusendmoddef
  subroutine radiation_generator_generate (generator, prt_tot_in, prt_tot_out)
    type :: prt_array_t
       type(string_t), dimension(:), allocatable :: prt
    end type
    integer, parameter :: n_flv_max = 10
    class(radiation_generator_t), intent(inout) :: generator
    type(string_t), intent(out), dimension(:), allocatable :: prt_tot_in, prt_tot_out
    type(prt_array_t), dimension(n_flv_max) :: prt_in, prt_out
    type(prt_array_t), dimension(n_flv_max) :: prt_out0
    type(pdg_array_t), dimension(:), allocatable :: pdg_tmp, pdg_out, pdg_in, pdg_raw
    type(if_table_t) :: if_table
    type(pdg_list_t), dimension(:), allocatable :: pl_in, pl_out
    integer :: i, j
    integer, dimension(:), allocatable :: reshuffle_list
    logical :: found
    integer :: flv
    integer :: n_out
    type(string_t), dimension(:), allocatable :: buf

    allocate (pl_in (1), pl_out (1))
    found = .false.
    flv = 0
 
    pl_in(1) = generator%pl_in
    pl_out(1) = generator%pl_out
    
    call pl_in(1)%create_pdg_array (pdg_in)
    call pl_out(1)%create_pdg_array (pdg_out)

    call if_table%init &
         (generator%radiation_model, pl_in, pl_out, generator%constraints)
    call if_table%radiate (generator%constraints)

    allocate (pdg_raw (generator%n_tot))

    do i = 1, if_table%get_length ()
      call if_table%get_pdg_out (i, pdg_tmp)
      if (size (pdg_tmp) == generator%n_tot) then
         call if_table%get_particle_string (i, 2, prt_tot_in, prt_out0(flv+1)%prt)
         call pdg_reshuffle (pdg_out, pdg_tmp, reshuffle_list)
         pdg_raw(1:2) = pdg_in(1:2)
         do j = 1, size (reshuffle_list)
            pdg_raw(reshuffle_list(j)+2) = pdg_tmp(j+2)
         end do
         found = .true.
         flv = flv+1
      end if
    end do

    if (found) then
      do i = 1, flv
         allocate (prt_out(i)%prt (generator%n_tot-2))
      end do
      allocate (prt_tot_out (generator%n_tot-2))
      allocate (buf (generator%n_tot-2))
      buf = ""

      do j = 1, flv
         do i = 1, size (reshuffle_list)
            prt_out(j)%prt(reshuffle_list(i)) = prt_out0(j)%prt(i)
            buf(i) = buf(i) // prt_out(j)%prt(i)
            if (j /= flv) buf(i) = buf(i) // ":"
         end do
      end do
      prt_tot_out = buf
    else
      call msg_fatal ("No QCD corrections for this process!")
    end if
    deallocate (pdg_raw)
  contains
    subroutine pdg_reshuffle (pdg_born, pdg_real, list)
      type(pdg_array_t), intent(in), dimension(:) :: pdg_born, pdg_real
      integer, intent(out), dimension(:), allocatable :: list
      type(pdg_sorter_born_t), dimension(:), allocatable :: sort_born
      type(pdg_sorter_real_t), dimension(:), allocatable :: sort_real
      integer :: i, i_min
      integer :: n_born, n_real
      integer :: ib, ir
      logical :: check
      integer, parameter :: n_in = 2
 
      n_born = size (pdg_born); n_real = size (pdg_real)
      allocate (list (n_real-n_in))
      allocate (sort_born (n_born))
      allocate (sort_real (n_real-n_in))

      sort_born%pdg = pdg_born%get ()
      sort_real%pdg = pdg_real(3:n_real)%get ()
      sort_born%checked = .false.
      sort_real%associated_born = 0

      do ib = 1, n_born
         sort_born(ib)%checked = any (sort_born(ib)%pdg == sort_real%pdg)
         if (sort_born(ib)%checked) then
            do ir = 1, n_real-2
               if (sort_born(ib)%pdg == sort_real(ir)%pdg) then
                  sort_real(ir)%associated_born = ib
                  exit
               end if
            end do
         end if
      end do

      i_min = maxval (sort_real%associated_born) + 1

      do ir = 1, n_real-2
         if (sort_real(ir)%associated_born == 0) then
            sort_real(ir)%associated_born = i_min
            i_min = i_min+1
         end if
      end do

      list = sort_real%associated_born

    end subroutine pdg_reshuffle 
         
  end subroutine radiation_generator_generate

\eatline
\nwendcode{}\nwbegindocs{159}\nwdocspar
\nwenddocs{}\nwbegincode{160}\moddef{radiation generator: radiation generator: TBP}\plusendmoddef
  procedure :: get_raw_states => radiation_generator_get_raw_states
\nwendcode{}\nwbegincode{161}\moddef{radiation generator: procedures}\plusendmoddef
  function radiation_generator_get_raw_states (generator) result (raw_states)
    class(radiation_generator_t), intent(in), target :: generator
    integer, dimension(:,:), allocatable :: raw_states
    type(pdg_states_t), pointer :: state
    integer :: n_states, n_particles
    integer :: i_state
    integer :: j
    state => generator%pdg_raw
    n_states = generator%pdg_raw%get_n_states ()
    n_particles = size (generator%pdg_raw%pdg)
    allocate (raw_states (n_particles, n_states))
    do i_state = 1, n_states
      do j = 1, n_particles
        raw_states (j, i_state) = state%pdg(j)%get ()
      end do
        state => state%next
    end do
  end function radiation_generator_get_raw_states

\eatline
\nwendcode{}\nwbegindocs{162}\nwdocspar
\nwenddocs{}\nwbegincode{163}\moddef{radiation generator: radiation generator: TBP}\plusendmoddef
  procedure :: save_born_raw => radiation_generator_save_born_raw
\nwendcode{}\nwbegincode{164}\moddef{radiation generator: procedures}\plusendmoddef
  subroutine radiation_generator_save_born_raw (generator, pdg_in, pdg_out)
    class(radiation_generator_t), intent(inout) :: generator
    type(pdg_array_t), dimension(:), allocatable, intent(in) :: pdg_in, pdg_out
    integer :: i 
    !!! !!! !!! Explicit allocation due to gfortran 4.7.4  
    allocate (generator%pdg_in_born (size (pdg_in))) 
    do i = 1, size (pdg_in) 
       generator%pdg_in_born(i) = pdg_in(i) 
    end do 
    allocate (generator%pdg_out_born (size (pdg_out))) 
    do i = 1, size (pdg_out) 
       generator%pdg_out_born(i) = pdg_out(i) 
    end do 
  end subroutine radiation_generator_save_born_raw
\eatline
\nwendcode{}\nwbegindocs{165}\nwdocspar
\nwenddocs{}\nwbegincode{166}\moddef{radiation generator: radiation generator: TBP}\plusendmoddef
  procedure :: get_born_raw => radiation_generator_get_born_raw
\nwendcode{}\nwbegincode{167}\moddef{radiation generator: procedures}\plusendmoddef
  function radiation_generator_get_born_raw (generator) result (flv_born)
    class(radiation_generator_t), intent(in) :: generator
    integer, dimension(:,:), allocatable :: flv_born
    integer :: i_part, n_particles
    n_particles = size (generator%pdg_in_born) + size (generator%pdg_out_born)
    allocate (flv_born (n_particles, 1))
    flv_born(1,1) = generator%pdg_in_born(1)%get ()
    flv_born(2,1) = generator%pdg_in_born(2)%get ()
    do i_part = 3, n_particles
      flv_born(i_part, 1) = generator%pdg_out_born(i_part-2)%get ()
    end do
  end function radiation_generator_get_born_raw

\eatline
\nwendcode{}\nwbegindocs{168}\nwdocspar
\clearpage
%%%%%%%%%%%%%%%%%%%%%%%%%%%%%%%%%%%%%%%%%%%%%%%%%%%%%%%%%%%%%%%%%%%%%%%%
\section{Sindarin Expression Implementation}

This module defines expressions of all kinds, represented in
a tree structure, for repeated evaluation.  This provides an
implementation of the {\tt{}expr{\char95}base} abstract type.

We have two flavors of expressions: one with particles and one without
particles.  The latter version is used for defining cut/selection
criteria and for online analysis.
\nwenddocs{}\nwbegincode{169}\moddef{\code{}eval{\char95}trees.f90\edoc{}}\endmoddef
\LA{}File header\RA{}

module eval_trees

  use, intrinsic :: iso_c_binding !NODEP!
\LA{}Use kinds\RA{}
\LA{}Use strings\RA{}
  use io_units
  use constants
  use format_defs, only: FMT_19
  use unit_tests
  use diagnostics
  use lorentz
  use md5
  use formats
  use sorting
  use ifiles
  use lexers
  use syntax_rules
  use parser
  use analysis
  use jets
  use pdg_arrays
  use subevents
  use user_code_interface
  use var_base
  use expr_base
  use variables
  use observables

\LA{}Standard module head\RA{}

\LA{}Eval trees: public\RA{}

\LA{}Eval trees: types\RA{}

\LA{}Eval trees: interfaces\RA{}

\LA{}Eval trees: variables\RA{}

contains

\LA{}Eval trees: procedures\RA{}

\LA{}Eval trees: tests\RA{}

end module eval_trees
\eatline
\nwendcode{}\nwbegindocs{170}\nwdocspar
\subsection{Tree nodes}
The evaluation tree consists of branch nodes (unary and binary) and of
leaf nodes, originating from a common root.  The node object should be
polymorphic.  For the time being, polymorphism is emulated here.  This
means that we have to maintain all possibilities that the node may
hold, including associated procedures as pointers.

The following parameter values characterize the node.  Unary and
binary operators have sub-nodes.  The other are leaf nodes.  Possible
leafs are literal constants or named-parameter references.
\nwenddocs{}\nwbegincode{171}\moddef{Eval trees: types}\endmoddef
  integer, parameter :: EN_UNKNOWN = 0, EN_UNARY = 1, EN_BINARY = 2
  integer, parameter :: EN_CONSTANT = 3, EN_VARIABLE = 4
  integer, parameter :: EN_CONDITIONAL = 5, EN_BLOCK = 6
  integer, parameter :: EN_RECORD_CMD = 7
  integer, parameter :: EN_OBS1_INT = 11, EN_OBS2_INT = 12
  integer, parameter :: EN_OBS1_REAL = 21, EN_OBS2_REAL = 22
  integer, parameter :: EN_UOBS1_INT = 31, EN_UOBS2_INT = 32
  integer, parameter :: EN_UOBS1_REAL = 41, EN_UOBS2_REAL = 42
  integer, parameter :: EN_PRT_FUN_UNARY = 101, EN_PRT_FUN_BINARY = 102
  integer, parameter :: EN_EVAL_FUN_UNARY = 111, EN_EVAL_FUN_BINARY = 112
  integer, parameter :: EN_LOG_FUN_UNARY = 121, EN_LOG_FUN_BINARY = 122
  integer, parameter :: EN_INT_FUN_UNARY = 131, EN_INT_FUN_BINARY = 132
  integer, parameter :: EN_REAL_FUN_UNARY = 141, EN_REAL_FUN_BINARY = 142
  integer, parameter :: EN_FORMAT_STR = 161

\eatline
\eatline
\eatline
\eatline
\eatline
\eatline
\eatline
\eatline
\eatline
\eatline
\nwendcode{}\nwbegincode{172}\moddef{Eval trees: types}\plusendmoddef
  type :: eval_node_t
     private
     type(string_t) :: tag
     integer :: type = EN_UNKNOWN
     integer :: result_type = V_NONE
     type(var_list_t), pointer :: var_list => null ()
     type(string_t) :: var_name
     logical, pointer :: value_is_known => null ()
     logical,           pointer :: lval => null ()
     integer,           pointer :: ival => null ()
     real(default),     pointer :: rval => null ()
     complex(default),  pointer :: cval => null ()
     type(subevt_t),  pointer :: pval => null ()
     type(pdg_array_t), pointer :: aval => null ()
     type(string_t),    pointer :: sval => null ()
     type(eval_node_t), pointer :: arg0 => null ()
     type(eval_node_t), pointer :: arg1 => null ()
     type(eval_node_t), pointer :: arg2 => null ()
     type(eval_node_t), pointer :: arg3 => null ()
     type(eval_node_t), pointer :: arg4 => null ()
     procedure(obs_unary_int),   nopass, pointer :: obs1_int  => null ()
     procedure(obs_unary_real),  nopass, pointer :: obs1_real => null ()
     procedure(obs_binary_int),  nopass, pointer :: obs2_int  => null ()
     procedure(obs_binary_real), nopass, pointer :: obs2_real => null ()
     integer, pointer :: prt_type => null ()
     integer, pointer :: index => null ()
     real(default), pointer :: tolerance => null ()
     integer, pointer :: jet_algorithm => null ()
     real(default), pointer :: jet_r => null ()
     type(prt_t), pointer :: prt1 => null ()
     type(prt_t), pointer :: prt2 => null ()
     procedure(unary_log),  nopass, pointer :: op1_log  => null ()
     procedure(unary_int),  nopass, pointer :: op1_int  => null ()
     procedure(unary_real), nopass, pointer :: op1_real => null ()
     procedure(unary_cmplx), nopass, pointer :: op1_cmplx => null ()
     procedure(unary_pdg),  nopass, pointer :: op1_pdg  => null ()
     procedure(unary_sev),  nopass, pointer :: op1_sev  => null ()
     procedure(unary_str),  nopass, pointer :: op1_str  => null ()
     procedure(unary_cut),  nopass, pointer :: op1_cut  => null ()
     procedure(unary_evi),  nopass, pointer :: op1_evi  => null ()
     procedure(unary_evr),  nopass, pointer :: op1_evr  => null ()
     procedure(binary_log),  nopass, pointer :: op2_log  => null ()
     procedure(binary_int),  nopass, pointer :: op2_int  => null ()
     procedure(binary_real), nopass, pointer :: op2_real => null ()
     procedure(binary_cmplx), nopass, pointer :: op2_cmplx => null ()
     procedure(binary_pdg),  nopass, pointer :: op2_pdg  => null ()
     procedure(binary_sev),  nopass, pointer :: op2_sev  => null ()
     procedure(binary_str),  nopass, pointer :: op2_str  => null ()
     procedure(binary_cut),  nopass, pointer :: op2_cut  => null ()
     procedure(binary_evi),  nopass, pointer :: op2_evi  => null ()
     procedure(binary_evr),  nopass, pointer :: op2_evr  => null ()
  end type eval_node_t

\eatline
\nwendcode{}\nwbegindocs{173}Finalize a node recursively.  Allocated constants are deleted,
pointers are ignored.
\nwenddocs{}\nwbegincode{174}\moddef{Eval trees: procedures}\endmoddef
  recursive subroutine eval_node_final_rec (node)
    type(eval_node_t), intent(inout) :: node
    select case (node%type)
    case (EN_UNARY)
       call eval_node_final_rec (node%arg1)
    case (EN_BINARY)
       call eval_node_final_rec (node%arg1)
       call eval_node_final_rec (node%arg2)
    case (EN_CONDITIONAL)
       call eval_node_final_rec (node%arg0)
       call eval_node_final_rec (node%arg1)
       call eval_node_final_rec (node%arg2)
    case (EN_BLOCK)
       call eval_node_final_rec (node%arg0)
       call eval_node_final_rec (node%arg1)
    case (EN_PRT_FUN_UNARY, EN_EVAL_FUN_UNARY, &
          EN_LOG_FUN_UNARY, EN_INT_FUN_UNARY, EN_REAL_FUN_UNARY)
       if (associated (node%arg0))  call eval_node_final_rec (node%arg0)
       call eval_node_final_rec (node%arg1)
       deallocate (node%index)
       deallocate (node%prt1)
    case (EN_PRT_FUN_BINARY, EN_EVAL_FUN_BINARY, &
          EN_LOG_FUN_BINARY, EN_INT_FUN_BINARY, EN_REAL_FUN_BINARY)
       if (associated (node%arg0))  call eval_node_final_rec (node%arg0)
       call eval_node_final_rec (node%arg1)
       call eval_node_final_rec (node%arg2)
       deallocate (node%index)
       deallocate (node%prt1)
       deallocate (node%prt2)
    case (EN_FORMAT_STR)
       if (associated (node%arg0))  call eval_node_final_rec (node%arg0)
       if (associated (node%arg1))  call eval_node_final_rec (node%arg1)
       deallocate (node%ival)
    case (EN_RECORD_CMD)
       if (associated (node%arg0))  call eval_node_final_rec (node%arg0)
       if (associated (node%arg1))  call eval_node_final_rec (node%arg1)
       if (associated (node%arg2))  call eval_node_final_rec (node%arg2)
       if (associated (node%arg3))  call eval_node_final_rec (node%arg3)
       if (associated (node%arg4))  call eval_node_final_rec (node%arg4)
    end select
    select case (node%type)
    case (EN_UNARY, EN_BINARY, EN_CONDITIONAL, EN_CONSTANT, EN_BLOCK, &
          EN_PRT_FUN_UNARY, EN_PRT_FUN_BINARY, &
          EN_EVAL_FUN_UNARY, EN_EVAL_FUN_BINARY, &
          EN_LOG_FUN_UNARY, EN_LOG_FUN_BINARY, &
          EN_INT_FUN_UNARY, EN_INT_FUN_BINARY, &
          EN_REAL_FUN_UNARY, EN_REAL_FUN_BINARY, &
          EN_FORMAT_STR, EN_RECORD_CMD)
       select case (node%result_type)
       case (V_LOG);  deallocate (node%lval)
       case (V_INT);  deallocate (node%ival)
       case (V_REAL); deallocate (node%rval)
       case (V_CMPLX); deallocate (node%cval)
       case (V_SEV);  deallocate (node%pval)
       case (V_PDG);  deallocate (node%aval)
       case (V_STR);  deallocate (node%sval)
       end select
       deallocate (node%value_is_known)
    end select
  end subroutine eval_node_final_rec

\eatline
\nwendcode{}\nwbegindocs{175}\nwdocspar
\subsubsection{Leaf nodes}
Initialize a leaf node with a literal constant.
\nwenddocs{}\nwbegincode{176}\moddef{Eval trees: procedures}\plusendmoddef
  subroutine eval_node_init_log (node, lval)
    type(eval_node_t), intent(out) :: node
    logical, intent(in) :: lval
    node%type = EN_CONSTANT
    node%result_type = V_LOG
    allocate (node%lval, node%value_is_known)
    node%lval = lval
    node%value_is_known = .true.
  end subroutine eval_node_init_log

  subroutine eval_node_init_int (node, ival)
    type(eval_node_t), intent(out) :: node
    integer, intent(in) :: ival
    node%type = EN_CONSTANT
    node%result_type = V_INT
    allocate (node%ival, node%value_is_known)
    node%ival = ival
    node%value_is_known = .true.
  end subroutine eval_node_init_int

  subroutine eval_node_init_real (node, rval)
    type(eval_node_t), intent(out) :: node
    real(default), intent(in) :: rval
    node%type = EN_CONSTANT
    node%result_type = V_REAL
    allocate (node%rval, node%value_is_known)
    node%rval = rval
    node%value_is_known = .true.
  end subroutine eval_node_init_real
    
  subroutine eval_node_init_cmplx (node, cval)
    type(eval_node_t), intent(out) :: node
    complex(default), intent(in) :: cval
    node%type = EN_CONSTANT
    node%result_type = V_CMPLX
    allocate (node%cval, node%value_is_known)
    node%cval = cval
    node%value_is_known = .true.
  end subroutine eval_node_init_cmplx

  subroutine eval_node_init_subevt (node, pval)
    type(eval_node_t), intent(out) :: node
    type(subevt_t), intent(in) :: pval
    node%type = EN_CONSTANT
    node%result_type = V_SEV
    allocate (node%pval, node%value_is_known)
    node%pval = pval
    node%value_is_known = .true.
  end subroutine eval_node_init_subevt

  subroutine eval_node_init_pdg_array (node, aval)
    type(eval_node_t), intent(out) :: node
    type(pdg_array_t), intent(in) :: aval
    node%type = EN_CONSTANT
    node%result_type = V_PDG
    allocate (node%aval, node%value_is_known)
    node%aval = aval
    node%value_is_known = .true.
  end subroutine eval_node_init_pdg_array

  subroutine eval_node_init_string (node, sval)
    type(eval_node_t), intent(out) :: node
    type(string_t), intent(in) :: sval
    node%type = EN_CONSTANT
    node%result_type = V_STR
    allocate (node%sval, node%value_is_known)
    node%sval = sval
    node%value_is_known = .true.
  end subroutine eval_node_init_string

\eatline
\eatline
\eatline
\nwendcode{}\nwbegindocs{177}Initialize a leaf node with a pointer to a named parameter
\nwenddocs{}\nwbegincode{178}\moddef{Eval trees: procedures}\plusendmoddef
  subroutine eval_node_init_log_ptr (node, name, lval, is_known)
    type(eval_node_t), intent(out) :: node
    type(string_t), intent(in) :: name
    logical, intent(in), target :: lval
    logical, intent(in), target :: is_known
    node%type = EN_VARIABLE
    node%tag = name
    node%result_type = V_LOG
    node%lval => lval
    node%value_is_known => is_known
  end subroutine eval_node_init_log_ptr

  subroutine eval_node_init_int_ptr (node, name, ival, is_known)
    type(eval_node_t), intent(out) :: node
    type(string_t), intent(in) :: name
    integer, intent(in), target :: ival
    logical, intent(in), target :: is_known
    node%type = EN_VARIABLE
    node%tag = name
    node%result_type = V_INT
    node%ival => ival
    node%value_is_known => is_known
  end subroutine eval_node_init_int_ptr

  subroutine eval_node_init_real_ptr (node, name, rval, is_known)
    type(eval_node_t), intent(out) :: node
    type(string_t), intent(in) :: name
    real(default), intent(in), target :: rval
    logical, intent(in), target :: is_known
    node%type = EN_VARIABLE
    node%tag = name
    node%result_type = V_REAL
    node%rval => rval
    node%value_is_known => is_known
  end subroutine eval_node_init_real_ptr
  
  subroutine eval_node_init_cmplx_ptr (node, name, cval, is_known)
    type(eval_node_t), intent(out) :: node
    type(string_t), intent(in) :: name
    complex(default), intent(in), target :: cval
    logical, intent(in), target :: is_known
    node%type = EN_VARIABLE
    node%tag = name
    node%result_type = V_CMPLX
    node%cval => cval
    node%value_is_known => is_known
  end subroutine eval_node_init_cmplx_ptr

  subroutine eval_node_init_subevt_ptr (node, name, pval, is_known)
    type(eval_node_t), intent(out) :: node
    type(string_t), intent(in) :: name
    type(subevt_t), intent(in), target :: pval
    logical, intent(in), target :: is_known
    node%type = EN_VARIABLE
    node%tag = name
    node%result_type = V_SEV
    node%pval => pval
    node%value_is_known => is_known
  end subroutine eval_node_init_subevt_ptr

  subroutine eval_node_init_pdg_array_ptr (node, name, aval, is_known)
    type(eval_node_t), intent(out) :: node
    type(string_t), intent(in) :: name
    type(pdg_array_t), intent(in), target :: aval
    logical, intent(in), target :: is_known
    node%type = EN_VARIABLE
    node%tag = name
    node%result_type = V_PDG
    node%aval => aval
    node%value_is_known => is_known
  end subroutine eval_node_init_pdg_array_ptr

  subroutine eval_node_init_string_ptr (node, name, sval, is_known)
    type(eval_node_t), intent(out) :: node
    type(string_t), intent(in) :: name
    type(string_t), intent(in), target :: sval
    logical, intent(in), target :: is_known
    node%type = EN_VARIABLE
    node%tag = name
    node%result_type = V_STR
    node%sval => sval
    node%value_is_known => is_known
  end subroutine eval_node_init_string_ptr

\eatline
\eatline
\eatline
\nwendcode{}\nwbegindocs{179}The procedure-pointer cases:
\nwenddocs{}\nwbegincode{180}\moddef{Eval trees: procedures}\plusendmoddef
  subroutine eval_node_init_obs1_int_ptr (node, name, obs1_iptr, p1)
    type(eval_node_t), intent(out) :: node
    type(string_t), intent(in) :: name
    procedure(obs_unary_int), intent(in), pointer :: obs1_iptr
    type(prt_t), intent(in), target :: p1
    node%type = EN_OBS1_INT
    node%tag = name
    node%result_type = V_INT
    node%obs1_int => obs1_iptr
    node%prt1 => p1
    allocate (node%ival, node%value_is_known)
    node%value_is_known = .false.
  end subroutine eval_node_init_obs1_int_ptr

  subroutine eval_node_init_obs2_int_ptr (node, name, obs2_iptr, p1, p2)
    type(eval_node_t), intent(out) :: node
    type(string_t), intent(in) :: name
    procedure(obs_binary_int), intent(in), pointer :: obs2_iptr
    type(prt_t), intent(in), target :: p1, p2
    node%type = EN_OBS2_INT
    node%tag = name
    node%result_type = V_INT
    node%obs2_int => obs2_iptr
    node%prt1 => p1
    node%prt2 => p2
    allocate (node%ival, node%value_is_known)
    node%value_is_known = .false.
  end subroutine eval_node_init_obs2_int_ptr

  subroutine eval_node_init_obs1_real_ptr (node, name, obs1_rptr, p1)
    type(eval_node_t), intent(out) :: node
    type(string_t), intent(in) :: name
    procedure(obs_unary_real), intent(in), pointer :: obs1_rptr
    type(prt_t), intent(in), target :: p1
    node%type = EN_OBS1_REAL
    node%tag = name
    node%result_type = V_REAL
    node%obs1_real => obs1_rptr
    node%prt1 => p1
    allocate (node%rval, node%value_is_known)
    node%value_is_known = .false.
  end subroutine eval_node_init_obs1_real_ptr

  subroutine eval_node_init_obs2_real_ptr (node, name, obs2_rptr, p1, p2)
    type(eval_node_t), intent(out) :: node
    type(string_t), intent(in) :: name
    procedure(obs_binary_real), intent(in), pointer :: obs2_rptr
    type(prt_t), intent(in), target :: p1, p2
    node%type = EN_OBS2_REAL
    node%tag = name
    node%result_type = V_REAL
    node%obs2_real => obs2_rptr
    node%prt1 => p1
    node%prt2 => p2
    allocate (node%rval, node%value_is_known)
    node%value_is_known = .false.
  end subroutine eval_node_init_obs2_real_ptr

\eatline
\eatline
\eatline
\eatline
\nwendcode{}\nwbegindocs{181}These nodes refer to user-defined procedures.
\nwenddocs{}\nwbegincode{182}\moddef{Eval trees: procedures}\plusendmoddef
  subroutine eval_node_init_uobs1_int (node, name, arg)
    type(eval_node_t), intent(out) :: node
    type(string_t), intent(in) :: name
    type(eval_node_t), intent(in), target :: arg
    node%type = EN_UOBS1_INT
    node%tag = name
    node%result_type = V_INT
    allocate (node%ival, node%value_is_known)
    node%value_is_known = .false.
    node%arg0 => arg
  end subroutine eval_node_init_uobs1_int

  subroutine eval_node_init_uobs2_int (node, name, arg)
    type(eval_node_t), intent(out) :: node
    type(string_t), intent(in) :: name
    type(eval_node_t), intent(in), target :: arg
    node%type = EN_UOBS2_INT
    node%tag = name
    node%result_type = V_INT
    allocate (node%ival, node%value_is_known)
    node%value_is_known = .false.
    node%arg0 => arg
  end subroutine eval_node_init_uobs2_int

  subroutine eval_node_init_uobs1_real (node, name, arg)
    type(eval_node_t), intent(out) :: node
    type(string_t), intent(in) :: name
    type(eval_node_t), intent(in), target :: arg
    node%type = EN_UOBS1_REAL
    node%tag = name
    node%result_type = V_REAL
    allocate (node%rval, node%value_is_known)
    node%value_is_known = .false.
    node%arg0 => arg
  end subroutine eval_node_init_uobs1_real

  subroutine eval_node_init_uobs2_real (node, name, arg)
    type(eval_node_t), intent(out) :: node
    type(string_t), intent(in) :: name
    type(eval_node_t), intent(in), target :: arg
    node%type = EN_UOBS2_REAL
    node%tag = name
    node%result_type = V_REAL
    allocate (node%rval, node%value_is_known)
    node%value_is_known = .false.
    node%arg0 => arg
  end subroutine eval_node_init_uobs2_real

\eatline
\eatline
\eatline
\eatline
\nwendcode{}\nwbegindocs{183}\nwdocspar
\subsubsection{Branch nodes}
Initialize a branch node, sub-nodes are given.
\nwenddocs{}\nwbegincode{184}\moddef{Eval trees: procedures}\plusendmoddef
  subroutine eval_node_init_branch (node, tag, result_type, arg1, arg2)
    type(eval_node_t), intent(out) :: node
    type(string_t), intent(in) :: tag
    integer, intent(in) :: result_type
    type(eval_node_t), intent(in), target :: arg1
    type(eval_node_t), intent(in), target, optional :: arg2
    if (present (arg2)) then
       node%type = EN_BINARY
    else
       node%type = EN_UNARY
    end if
    node%tag = tag
    node%result_type = result_type
    call eval_node_allocate_value (node)
    node%arg1 => arg1
    if (present (arg2))  node%arg2 => arg2
  end subroutine eval_node_init_branch

\eatline
\nwendcode{}\nwbegindocs{185}Allocate the node value according to the result type.
\nwenddocs{}\nwbegincode{186}\moddef{Eval trees: procedures}\plusendmoddef
  subroutine eval_node_allocate_value (node)
    type(eval_node_t), intent(inout) :: node
    select case (node%result_type)
    case (V_LOG);  allocate (node%lval)
    case (V_INT);  allocate (node%ival)
    case (V_REAL); allocate (node%rval)
    case (V_CMPLX); allocate (node%cval)
    case (V_PDG);  allocate (node%aval)
    case (V_SEV);  allocate (node%pval)
       call subevt_init (node%pval)
    case (V_STR);  allocate (node%sval)
    end select
    allocate (node%value_is_known)
    node%value_is_known = .false.
  end subroutine eval_node_allocate_value

\eatline
\nwendcode{}\nwbegindocs{187}Initialize a block node which contains, in addition to the
expression to be evaluated, a variable definition.  The result type is
not yet assigned, because we can compile the enclosed expression only
after the var list is set up.

Note that the node always allocates a new variable list and appends it to the
current one.  Thus, if the variable redefines an existing one, it only shadows
it but does not reset it.  Any side-effects are therefore absent and need not
be undone outside the block.

If the flag {\tt{}new} is set, a variable is (re)declared.  This must not be done
for intrinsic variables.  Vice versa, if the variable is not existent, the
{\tt{}new} flag is required.
\nwenddocs{}\nwbegincode{188}\moddef{Eval trees: procedures}\plusendmoddef
  subroutine eval_node_init_block (node, name, type, var_def, var_list)
    type(eval_node_t), intent(out), target :: node
    type(string_t), intent(in) :: name
    integer, intent(in) :: type
    type(eval_node_t), intent(in), target :: var_def
    type(var_list_t), intent(in), target :: var_list
    node%type = EN_BLOCK
    node%tag = "var_def"
    node%var_name = name
    node%arg1 => var_def
    allocate (node%var_list)
    call node%var_list%link (var_list)
    if (var_def%type == EN_CONSTANT) then
       select case (type)
       case (V_LOG)
          call var_list_append_log  (node%var_list, name, var_def%lval)
       case (V_INT)
          call var_list_append_int  (node%var_list, name, var_def%ival)
       case (V_REAL)
          call var_list_append_real (node%var_list, name, var_def%rval)
       case (V_CMPLX)
          call var_list_append_cmplx (node%var_list, name, var_def%cval)
       case (V_PDG)
          call var_list_append_pdg_array &
               (node%var_list, name, var_def%aval)
       case (V_SEV)
          call var_list_append_subevt &
               (node%var_list, name, var_def%pval)
       case (V_STR)
          call var_list_append_string (node%var_list, name, var_def%sval)
       end select
    else
       select case (type)
       case (V_LOG);  call var_list_append_log_ptr &
            (node%var_list, name, var_def%lval, var_def%value_is_known)
       case (V_INT);  call var_list_append_int_ptr &
            (node%var_list, name, var_def%ival, var_def%value_is_known)
       case (V_REAL); call var_list_append_real_ptr &
            (node%var_list, name, var_def%rval, var_def%value_is_known)
       case (V_CMPLX); call var_list_append_cmplx_ptr &
            (node%var_list, name, var_def%cval, var_def%value_is_known)
       case (V_PDG);  call var_list_append_pdg_array_ptr &
            (node%var_list, name, var_def%aval, var_def%value_is_known)
       case (V_SEV); call var_list_append_subevt_ptr &
            (node%var_list, name, var_def%pval, var_def%value_is_known)
       case (V_STR); call var_list_append_string_ptr &
            (node%var_list, name, var_def%sval, var_def%value_is_known)
       end select
    end if
  end subroutine eval_node_init_block

\eatline
\nwendcode{}\nwbegindocs{189}Complete block initialization by assigning the expression to
evaluate to {\tt{}arg0}.
\nwenddocs{}\nwbegincode{190}\moddef{Eval trees: procedures}\plusendmoddef
  subroutine eval_node_set_expr (node, arg, result_type)
    type(eval_node_t), intent(inout) :: node
    type(eval_node_t), intent(in), target :: arg
    integer, intent(in), optional :: result_type
    if (present (result_type)) then
       node%result_type = result_type
    else
       node%result_type = arg%result_type
    end if
    call eval_node_allocate_value (node)
    node%arg0 => arg
  end subroutine eval_node_set_expr

\eatline
\nwendcode{}\nwbegindocs{191}Initialize a conditional.  There are three branches: the condition
(evaluates to logical) and the two alternatives (evaluate both to the
same arbitrary type).
\nwenddocs{}\nwbegincode{192}\moddef{Eval trees: procedures}\plusendmoddef
  subroutine eval_node_init_conditional (node, result_type, cond, arg1, arg2)
    type(eval_node_t), intent(out) :: node
    integer, intent(in) :: result_type
    type(eval_node_t), intent(in), target :: cond, arg1, arg2
    node%type = EN_CONDITIONAL
    node%tag = "cond"
    node%result_type = result_type
    call eval_node_allocate_value (node)
    node%arg0 => cond
    node%arg1 => arg1
    node%arg2 => arg2
  end subroutine eval_node_init_conditional

\eatline
\nwendcode{}\nwbegindocs{193}Initialize a recording command (which evaluates to a logical
constant).  The first branch is the ID of the analysis object to be
filled, the optional branches 1 to 4 are the values to be recorded.

If the event-weight pointer is null, we record values with unit weight.
Otherwise, we use the value pointed to as event weight.

There can be up to four arguments which represent $x$, $y$, $\Delta y$,
$\Delta x$.  Therefore, this is the only node type that may fill four
sub-nodes.
\nwenddocs{}\nwbegincode{194}\moddef{Eval trees: procedures}\plusendmoddef
  subroutine eval_node_init_record_cmd &
      (node, event_weight, id, arg1, arg2, arg3, arg4)
    type(eval_node_t), intent(out) :: node
    real(default), pointer :: event_weight
    type(eval_node_t), intent(in), target :: id
    type(eval_node_t), intent(in), optional, target :: arg1, arg2, arg3, arg4
    call eval_node_init_log (node, .true.)
    node%type = EN_RECORD_CMD
    node%rval => event_weight
    node%tag = "record_cmd"
    node%arg0 => id
    if (present (arg1)) then
       node%arg1 => arg1
       if (present (arg2)) then
          node%arg2 => arg2
          if (present (arg3)) then
             node%arg3 => arg3
             if (present (arg4)) then
                node%arg4 => arg4
             end if
          end if
       end if
    end if
  end subroutine eval_node_init_record_cmd
    
\eatline
\nwendcode{}\nwbegindocs{195}Initialize a node for operations on subevents.  The particle
lists (one or two) are inserted as {\tt{}arg1} and {\tt{}arg2}.  We
allocated particle pointers as temporaries for iterating over particle
lists.  The procedure pointer which holds the function to evaluate for
the subevents (e.g., combine, select) is also initialized.  
\nwenddocs{}\nwbegincode{196}\moddef{Eval trees: procedures}\plusendmoddef
  subroutine eval_node_init_prt_fun_unary (node, arg1, name, proc)
    type(eval_node_t), intent(out) :: node
    type(eval_node_t), intent(in), target :: arg1
    type(string_t), intent(in) :: name
    procedure(unary_sev) :: proc
    node%type = EN_PRT_FUN_UNARY
    node%tag = name
    node%result_type = V_SEV
    call eval_node_allocate_value (node)
    node%arg1 => arg1
    allocate (node%index, source = 0)
    allocate (node%prt1)
    node%op1_sev => proc
  end subroutine eval_node_init_prt_fun_unary

  subroutine eval_node_init_prt_fun_binary (node, arg1, arg2, name, proc)
    type(eval_node_t), intent(out) :: node
    type(eval_node_t), intent(in), target :: arg1, arg2
    type(string_t), intent(in) :: name
    procedure(binary_sev) :: proc
    node%type = EN_PRT_FUN_BINARY
    node%tag = name
    node%result_type = V_SEV
    call eval_node_allocate_value (node)
    node%arg1 => arg1
    node%arg2 => arg2
    allocate (node%index, source = 0)
    allocate (node%prt1)
    allocate (node%prt2)
    node%op2_sev => proc
  end subroutine eval_node_init_prt_fun_binary

\eatline
\nwendcode{}\nwbegindocs{197}Similar, but for particle-list functions that evaluate to a real
value.
\nwenddocs{}\nwbegincode{198}\moddef{Eval trees: procedures}\plusendmoddef
  subroutine eval_node_init_eval_fun_unary (node, arg1, name)
    type(eval_node_t), intent(out) :: node
    type(eval_node_t), intent(in), target :: arg1
    type(string_t), intent(in) :: name
    node%type = EN_EVAL_FUN_UNARY
    node%tag = name
    node%result_type = V_REAL
    call eval_node_allocate_value (node)
    node%arg1 => arg1
    allocate (node%index, source = 0)
    allocate (node%prt1)
  end subroutine eval_node_init_eval_fun_unary

  subroutine eval_node_init_eval_fun_binary (node, arg1, arg2, name)
    type(eval_node_t), intent(out) :: node
    type(eval_node_t), intent(in), target :: arg1, arg2
    type(string_t), intent(in) :: name
    node%type = EN_EVAL_FUN_BINARY
    node%tag = name
    node%result_type = V_REAL
    call eval_node_allocate_value (node)
    node%arg1 => arg1
    node%arg2 => arg2
    allocate (node%index, source = 0)
    allocate (node%prt1)
    allocate (node%prt2)
  end subroutine eval_node_init_eval_fun_binary

\eatline
\nwendcode{}\nwbegindocs{199}These are for particle-list functions that evaluate to a logical
value.
\nwenddocs{}\nwbegincode{200}\moddef{Eval trees: procedures}\plusendmoddef
  subroutine eval_node_init_log_fun_unary (node, arg1, name, proc)
    type(eval_node_t), intent(out) :: node
    type(eval_node_t), intent(in), target :: arg1
    type(string_t), intent(in) :: name
    procedure(unary_cut) :: proc
    node%type = EN_LOG_FUN_UNARY
    node%tag = name
    node%result_type = V_LOG
    call eval_node_allocate_value (node)
    node%arg1 => arg1
    allocate (node%index, source = 0)
    allocate (node%prt1)
    node%op1_cut => proc
  end subroutine eval_node_init_log_fun_unary

  subroutine eval_node_init_log_fun_binary (node, arg1, arg2, name, proc)
    type(eval_node_t), intent(out) :: node
    type(eval_node_t), intent(in), target :: arg1, arg2
    type(string_t), intent(in) :: name
    procedure(binary_cut) :: proc
    node%type = EN_LOG_FUN_BINARY
    node%tag = name
    node%result_type = V_LOG
    call eval_node_allocate_value (node)
    node%arg1 => arg1
    node%arg2 => arg2
    allocate (node%index, source = 0)
    allocate (node%prt1)
    allocate (node%prt2)
    node%op2_cut => proc
  end subroutine eval_node_init_log_fun_binary

\eatline
\nwendcode{}\nwbegindocs{201}These are for particle-list functions that evaluate to an integer
value.
\nwenddocs{}\nwbegincode{202}\moddef{Eval trees: procedures}\plusendmoddef
  subroutine eval_node_init_int_fun_unary (node, arg1, name, proc)
    type(eval_node_t), intent(out) :: node
    type(eval_node_t), intent(in), target :: arg1
    type(string_t), intent(in) :: name
    procedure(unary_evi) :: proc
    node%type = EN_INT_FUN_UNARY
    node%tag = name
    node%result_type = V_INT
    call eval_node_allocate_value (node)
    node%arg1 => arg1
    allocate (node%index, source = 0)
    allocate (node%prt1)
    node%op1_evi => proc
  end subroutine eval_node_init_int_fun_unary

  subroutine eval_node_init_int_fun_binary (node, arg1, arg2, name, proc)
    type(eval_node_t), intent(out) :: node
    type(eval_node_t), intent(in), target :: arg1, arg2
    type(string_t), intent(in) :: name
    procedure(binary_evi) :: proc
    node%type = EN_INT_FUN_BINARY
    node%tag = name
    node%result_type = V_INT
    call eval_node_allocate_value (node)
    node%arg1 => arg1
    node%arg2 => arg2
    allocate (node%index, source = 0)
    allocate (node%prt1)
    allocate (node%prt2)
    node%op2_evi => proc
  end subroutine eval_node_init_int_fun_binary

\eatline
\nwendcode{}\nwbegindocs{203}These are for particle-list functions that evaluate to a real
value.
\nwenddocs{}\nwbegincode{204}\moddef{Eval trees: procedures}\plusendmoddef
  subroutine eval_node_init_real_fun_unary (node, arg1, name, proc)
    type(eval_node_t), intent(out) :: node
    type(eval_node_t), intent(in), target :: arg1
    type(string_t), intent(in) :: name
    procedure(unary_evr) :: proc
    node%type = EN_REAL_FUN_UNARY
    node%tag = name
    node%result_type = V_INT
    call eval_node_allocate_value (node)
    node%arg1 => arg1
    allocate (node%index, source = 0)
    allocate (node%prt1)
    node%op1_evr => proc
  end subroutine eval_node_init_real_fun_unary

  subroutine eval_node_init_real_fun_binary (node, arg1, arg2, name, proc)
    type(eval_node_t), intent(out) :: node
    type(eval_node_t), intent(in), target :: arg1, arg2
    type(string_t), intent(in) :: name
    procedure(binary_evr) :: proc
    node%type = EN_REAL_FUN_BINARY
    node%tag = name
    node%result_type = V_INT
    call eval_node_allocate_value (node)
    node%arg1 => arg1
    node%arg2 => arg2
    allocate (node%index, source = 0)
    allocate (node%prt1)
    allocate (node%prt2)
    node%op2_evr => proc
  end subroutine eval_node_init_real_fun_binary

\eatline
\nwendcode{}\nwbegindocs{205}Initialize a node for a string formatting function (sprintf).
\nwenddocs{}\nwbegincode{206}\moddef{Eval trees: procedures}\plusendmoddef
  subroutine eval_node_init_format_string (node, fmt, arg, name, n_args)
    type(eval_node_t), intent(out) :: node
    type(eval_node_t), pointer :: fmt, arg
    type(string_t), intent(in) :: name
    integer, intent(in) :: n_args
    node%type = EN_FORMAT_STR
    node%tag = name
    node%result_type = V_STR
    call eval_node_allocate_value (node)
    node%arg0 => fmt
    node%arg1 => arg
    allocate (node%ival)
    node%ival = n_args
  end subroutine eval_node_init_format_string

\eatline
\nwendcode{}\nwbegindocs{207}If particle functions depend upon a condition (or an expression is
evaluated), the observables that can be evaluated for the given
particles have to be thrown on the local variable stack.  This is done
here.  Each observable is initialized with the particle pointers which
have been allocated for the node.

The integer variable that is referred to by the {\tt{}Index}
pseudo-observable is always known when it is referred to.
\nwenddocs{}\nwbegincode{208}\moddef{Eval trees: procedures}\plusendmoddef
  subroutine eval_node_set_observables (node, var_list)
    type(eval_node_t), intent(inout) :: node
    type(var_list_t), intent(in), target :: var_list
    logical, save, target :: known = .true.
    allocate (node%var_list)
    call node%var_list%link (var_list)
    allocate (node%index, source = 0)
    call var_list_append_int_ptr &
         (node%var_list, var_str ("Index"), node%index, known, intrinsic=.true.)
    if (.not. associated (node%prt2)) then
       call var_list_set_observables_unary &
            (node%var_list, node%prt1)
    else
       call var_list_set_observables_binary &
            (node%var_list, node%prt1, node%prt2)
    end if
  end subroutine eval_node_set_observables

\eatline
\nwendcode{}\nwbegindocs{209}\nwdocspar
\subsubsection{Output}
\nwenddocs{}\nwbegincode{210}\moddef{Eval trees: procedures}\plusendmoddef
  subroutine eval_node_write (node, unit, indent)
    type(eval_node_t), intent(in) :: node
    integer, intent(in), optional :: unit
    integer, intent(in), optional :: indent
    integer :: u, ind
    u = given_output_unit (unit);  if (u < 0)  return
    ind = 0;  if (present (indent)) ind = indent
    write (u, "(A)", advance="no")  repeat ("|  ", ind) // "o "
    select case (node%type)
    case (EN_UNARY, EN_BINARY, EN_CONDITIONAL, &
          EN_PRT_FUN_UNARY, EN_PRT_FUN_BINARY, &
          EN_EVAL_FUN_UNARY, EN_EVAL_FUN_BINARY, &
          EN_LOG_FUN_UNARY, EN_LOG_FUN_BINARY, &
          EN_INT_FUN_UNARY, EN_INT_FUN_BINARY, &
          EN_REAL_FUN_UNARY, EN_REAL_FUN_BINARY)
       write (u, "(A)", advance="no")  "[" // char (node%tag) // "] ="
    case (EN_CONSTANT)
       write (u, "(A)", advance="no")  "[const] ="
    case (EN_VARIABLE)
       write (u, "(A)", advance="no")  char (node%tag) // " =>"
    case (EN_OBS1_INT, EN_OBS2_INT, EN_OBS1_REAL, EN_OBS2_REAL, &
         EN_UOBS1_INT, EN_UOBS2_INT, EN_UOBS1_REAL, EN_UOBS2_REAL)
       write (u, "(A)", advance="no")  char (node%tag) // " ="
    case (EN_BLOCK)
       write (u, "(A)", advance="no")  "[" // char (node%tag) // "]" // &
            char (node%var_name) // " [expr] = "
    case default
       write (u, "(A)", advance="no")  "[???] ="
    end select
    select case (node%result_type)
    case (V_LOG)
       if (node%value_is_known) then
          if (node%lval) then
             write (u, "(1x,A)") "true"
          else
             write (u, "(1x,A)") "false"
          end if
       else
          write (u, "(1x,A)") "[unknown logical]"
       end if
    case (V_INT)
       if (node%value_is_known) then
          write (u, "(1x,I0)")  node%ival
       else
          write (u, "(1x,A)") "[unknown integer]"
       end if
    case (V_REAL)
       if (node%value_is_known) then
          write (u, "(1x," // FMT_19 // ")") node%rval
       else
          write (u, "(1x,A)") "[unknown real]"
       end if
   case (V_CMPLX)
       if (node%value_is_known) then
          write (u, "(1x,'('," // FMT_19 // ",','," // &
               FMT_19 // ",')')") node%cval
       else
          write (u, "(1x,A)") "[unknown complex]"
       end if
    case (V_SEV)
       if (char (node%tag) == "@evt") then
          write (u, "(1x,A)") "[event subevent]"
       else if (node%value_is_known) then
          call subevt_write &
               (node%pval, unit, prefix = repeat ("|  ", ind + 1))
       else
          write (u, "(1x,A)") "[unknown subevent]"
       end if
    case (V_PDG)
       write (u, "(1x)", advance="no")
       call pdg_array_write (node%aval, u);  write (u, *)
    case (V_STR)
       if (node%value_is_known) then
          write (u, "(A)")  '"' // char (node%sval) // '"'
       else
          write (u, "(1x,A)") "[unknown string]"
       end if
    case default
       write (u, "(1x,A)") "[empty]"
    end select
    select case (node%type)
    case (EN_OBS1_INT, EN_OBS1_REAL, EN_UOBS1_INT, EN_UOBS1_REAL)
       write (u, "(A,6x,A)", advance="no")  repeat ("|  ", ind), "prt1 ="
       call prt_write (node%prt1, unit)
    case (EN_OBS2_INT, EN_OBS2_REAL, EN_UOBS2_INT, EN_UOBS2_REAL)
       write (u, "(A,6x,A)", advance="no")  repeat ("|  ", ind), "prt1 ="
       call prt_write (node%prt1, unit)
       write (u, "(A,6x,A)", advance="no")  repeat ("|  ", ind), "prt2 ="
       call prt_write (node%prt2, unit)
    end select
  end subroutine eval_node_write

  recursive subroutine eval_node_write_rec (node, unit, indent)
    type(eval_node_t), intent(in) :: node
    integer, intent(in), optional :: unit
    integer, intent(in), optional :: indent
    integer :: u, ind
    u = given_output_unit (unit);  if (u < 0)  return
    ind = 0;  if (present (indent))  ind = indent
    call eval_node_write (node, unit, indent)
    select case (node%type)
    case (EN_UNARY)
       if (associated (node%arg0)) &
            call eval_node_write_rec (node%arg0, unit, ind+1)
       call eval_node_write_rec (node%arg1, unit, ind+1)
    case (EN_BINARY)
       if (associated (node%arg0)) &
            call eval_node_write_rec (node%arg0, unit, ind+1)
       call eval_node_write_rec (node%arg1, unit, ind+1)
       call eval_node_write_rec (node%arg2, unit, ind+1)
    case (EN_BLOCK)
       call eval_node_write_rec (node%arg1, unit, ind+1)
       call eval_node_write_rec (node%arg0, unit, ind+1)
    case (EN_CONDITIONAL)
       call eval_node_write_rec (node%arg0, unit, ind+1)
       call eval_node_write_rec (node%arg1, unit, ind+1)
       call eval_node_write_rec (node%arg2, unit, ind+1)
    case (EN_PRT_FUN_UNARY, EN_EVAL_FUN_UNARY, &
          EN_LOG_FUN_UNARY, EN_INT_FUN_UNARY, EN_REAL_FUN_UNARY)
       if (associated (node%arg0)) &
            call eval_node_write_rec (node%arg0, unit, ind+1)
       call eval_node_write_rec (node%arg1, unit, ind+1)
    case (EN_PRT_FUN_BINARY, EN_EVAL_FUN_BINARY, &
          EN_LOG_FUN_BINARY, EN_INT_FUN_BINARY, EN_REAL_FUN_BINARY)
       if (associated (node%arg0)) &
            call eval_node_write_rec (node%arg0, unit, ind+1)
       call eval_node_write_rec (node%arg1, unit, ind+1)
       call eval_node_write_rec (node%arg2, unit, ind+1)
    case (EN_RECORD_CMD)
       if (associated (node%arg1)) then
          call eval_node_write_rec (node%arg1, unit, ind+1)
          if (associated (node%arg2)) then
             call eval_node_write_rec (node%arg2, unit, ind+1)
             if (associated (node%arg3)) then
                call eval_node_write_rec (node%arg3, unit, ind+1)
                if (associated (node%arg4)) then
                   call eval_node_write_rec (node%arg4, unit, ind+1)
                end if
             end if
          end if
       end if
    end select
  end subroutine eval_node_write_rec

\eatline
\nwendcode{}\nwbegindocs{211}\nwdocspar
\subsection{Operation types}
For the operations associated to evaluation tree nodes, we define
abstract interfaces for all cases.

Particles/subevents are transferred by-reference, to avoid
unnecessary copying.  Therefore, subroutines instead of functions.
(Furthermore, the function version of {\tt{}unary{\char95}prt} triggers an obscure
bug in nagfor 5.2(649) [invalid C code].)
\nwenddocs{}\nwbegincode{212}\moddef{Eval trees: interfaces}\endmoddef
  abstract interface
     logical function unary_log (arg)
       import eval_node_t
       type(eval_node_t), intent(in) :: arg
     end function unary_log
  end interface
  abstract interface
     integer function unary_int (arg)
       import eval_node_t
       type(eval_node_t), intent(in) :: arg
     end function unary_int
  end interface
  abstract interface
     real(default) function unary_real (arg)
       import default
       import eval_node_t
       type(eval_node_t), intent(in) :: arg
     end function unary_real
  end interface
  abstract interface
     complex(default) function unary_cmplx (arg)
       import default
       import eval_node_t
       type(eval_node_t), intent(in) :: arg
     end function unary_cmplx
  end interface
  abstract interface
     subroutine unary_pdg (pdg_array, arg)
       import pdg_array_t
       import eval_node_t
       type(pdg_array_t), intent(out) :: pdg_array
       type(eval_node_t), intent(in) :: arg
     end subroutine unary_pdg
  end interface
  abstract interface
     subroutine unary_sev (subevt, arg, arg0)
       import subevt_t
       import eval_node_t
       type(subevt_t), intent(inout) :: subevt
       type(eval_node_t), intent(in) :: arg
       type(eval_node_t), intent(inout), optional :: arg0
     end subroutine unary_sev
  end interface
  abstract interface
     subroutine unary_str (string, arg)
       import string_t
       import eval_node_t
       type(string_t), intent(out) :: string
       type(eval_node_t), intent(in) :: arg
     end subroutine unary_str
  end interface
  abstract interface
     logical function unary_cut (arg1, arg0)
       import eval_node_t
       type(eval_node_t), intent(in) :: arg1
       type(eval_node_t), intent(inout) :: arg0
     end function unary_cut
  end interface
  abstract interface
     subroutine unary_evi (ival, arg1, arg0)
       import eval_node_t
       integer, intent(out) :: ival
       type(eval_node_t), intent(in) :: arg1
       type(eval_node_t), intent(inout), optional :: arg0
     end subroutine unary_evi
  end interface
  abstract interface
     subroutine unary_evr (rval, arg1, arg0)
       import eval_node_t, default
       real(default), intent(out) :: rval
       type(eval_node_t), intent(in) :: arg1
       type(eval_node_t), intent(inout), optional :: arg0
     end subroutine unary_evr
  end interface
  abstract interface
     logical function binary_log (arg1, arg2)
       import eval_node_t
       type(eval_node_t), intent(in) :: arg1, arg2
     end function binary_log
  end interface
  abstract interface
     integer function binary_int (arg1, arg2)
       import eval_node_t
       type(eval_node_t), intent(in) :: arg1, arg2
     end function binary_int
  end interface
  abstract interface
     real(default) function binary_real (arg1, arg2)
       import default
       import eval_node_t
       type(eval_node_t), intent(in) :: arg1, arg2
     end function binary_real
  end interface
  abstract interface
     complex(default) function binary_cmplx (arg1, arg2)
       import default
       import eval_node_t
       type(eval_node_t), intent(in) :: arg1, arg2
     end function binary_cmplx
  end interface
  abstract interface
     subroutine binary_pdg (pdg_array, arg1, arg2)
       import pdg_array_t
       import eval_node_t
       type(pdg_array_t), intent(out) :: pdg_array
       type(eval_node_t), intent(in) :: arg1, arg2
     end subroutine binary_pdg
  end interface
  abstract interface
     subroutine binary_sev (subevt, arg1, arg2, arg0)
       import subevt_t
       import eval_node_t
       type(subevt_t), intent(inout) :: subevt
       type(eval_node_t), intent(in) :: arg1, arg2
       type(eval_node_t), intent(inout), optional :: arg0
     end subroutine binary_sev
  end interface
  abstract interface
     subroutine binary_str (string, arg1, arg2)
       import string_t
       import eval_node_t
       type(string_t), intent(out) :: string
       type(eval_node_t), intent(in) :: arg1, arg2
     end subroutine binary_str
  end interface
  abstract interface
     logical function binary_cut (arg1, arg2, arg0)
       import eval_node_t
       type(eval_node_t), intent(in) :: arg1, arg2
       type(eval_node_t), intent(inout) :: arg0
     end function binary_cut
  end interface
  abstract interface
     subroutine binary_evi (ival, arg1, arg2, arg0)
       import eval_node_t
       integer, intent(out) :: ival
       type(eval_node_t), intent(in) :: arg1, arg2
       type(eval_node_t), intent(inout), optional :: arg0
     end subroutine binary_evi
  end interface
  abstract interface
     subroutine binary_evr (rval, arg1, arg2, arg0)
       import eval_node_t, default
       real(default), intent(out) :: rval
       type(eval_node_t), intent(in) :: arg1, arg2
       type(eval_node_t), intent(inout), optional :: arg0
     end subroutine binary_evr
  end interface

\nwendcode{}\nwbegindocs{213}The following subroutines set the procedure pointer:
\nwenddocs{}\nwbegincode{214}\moddef{Eval trees: procedures}\plusendmoddef
  subroutine eval_node_set_op1_log (en, op)
    type(eval_node_t), intent(inout) :: en
    procedure(unary_log) :: op
    en%op1_log => op
    end subroutine eval_node_set_op1_log
  
  subroutine eval_node_set_op1_int (en, op)
    type(eval_node_t), intent(inout) :: en
    procedure(unary_int) :: op
    en%op1_int => op
  end subroutine eval_node_set_op1_int

  subroutine eval_node_set_op1_real (en, op)
    type(eval_node_t), intent(inout) :: en
    procedure(unary_real) :: op
    en%op1_real => op
  end subroutine eval_node_set_op1_real
  
  subroutine eval_node_set_op1_cmplx (en, op)
    type(eval_node_t), intent(inout) :: en
    procedure(unary_cmplx) :: op
    en%op1_cmplx => op
  end subroutine eval_node_set_op1_cmplx

  subroutine eval_node_set_op1_pdg (en, op)
    type(eval_node_t), intent(inout) :: en
    procedure(unary_pdg) :: op
    en%op1_pdg => op
  end subroutine eval_node_set_op1_pdg

  subroutine eval_node_set_op1_sev (en, op)
    type(eval_node_t), intent(inout) :: en
    procedure(unary_sev) :: op
    en%op1_sev => op
  end subroutine eval_node_set_op1_sev

  subroutine eval_node_set_op1_str (en, op)
    type(eval_node_t), intent(inout) :: en
    procedure(unary_str) :: op
    en%op1_str => op
  end subroutine eval_node_set_op1_str

  subroutine eval_node_set_op2_log (en, op)
    type(eval_node_t), intent(inout) :: en
    procedure(binary_log) :: op
    en%op2_log => op
  end subroutine eval_node_set_op2_log

  subroutine eval_node_set_op2_int (en, op)
    type(eval_node_t), intent(inout) :: en
    procedure(binary_int) :: op
    en%op2_int => op
  end subroutine eval_node_set_op2_int

  subroutine eval_node_set_op2_real (en, op)
    type(eval_node_t), intent(inout) :: en
    procedure(binary_real) :: op
    en%op2_real => op
  end subroutine eval_node_set_op2_real
  
  subroutine eval_node_set_op2_cmplx (en, op)
    type(eval_node_t), intent(inout) :: en
    procedure(binary_cmplx) :: op
    en%op2_cmplx => op
  end subroutine eval_node_set_op2_cmplx

  subroutine eval_node_set_op2_pdg (en, op)
    type(eval_node_t), intent(inout) :: en
    procedure(binary_pdg) :: op
    en%op2_pdg => op
  end subroutine eval_node_set_op2_pdg

  subroutine eval_node_set_op2_sev (en, op)
    type(eval_node_t), intent(inout) :: en
    procedure(binary_sev) :: op
    en%op2_sev => op
  end subroutine eval_node_set_op2_sev

  subroutine eval_node_set_op2_str (en, op)
    type(eval_node_t), intent(inout) :: en
    procedure(binary_str) :: op
    en%op2_str => op
  end subroutine eval_node_set_op2_str

\eatline
\nwendcode{}\nwbegindocs{215}\nwdocspar
\subsection{Specific operators}

Our expression syntax contains all Fortran functions that make sense.
These functions have to be provided in a form that they can be used in
procedures pointers, and have the abstract interfaces above.
For some intrinsic functions, we could use specific versions provided
by Fortran directly.  However, this has two drawbacks: (i) We should
work with the values instead of the eval-nodes as argument, which
complicates the interface; (ii) more importantly, the {\tt{}default} real
type need not be equivalent to double precision.  This would, at
least, introduce system dependencies.  Finally, for operators there
are no specific versions.

Therefore, we write wrappers for all possible functions, at the
expense of some overhead.
\subsubsection{Binary numerical functions}
\nwenddocs{}\nwbegincode{216}\moddef{Eval trees: procedures}\plusendmoddef
  integer function add_ii (en1, en2) result (y)
    type(eval_node_t), intent(in) :: en1, en2
    y = en1%ival + en2%ival
  end function add_ii
  real(default) function add_ir (en1, en2) result (y)
    type(eval_node_t), intent(in) :: en1, en2
    y = en1%ival + en2%rval
  end function add_ir
  complex(default) function add_ic (en1, en2) result (y)
    type(eval_node_t), intent(in) :: en1, en2
    y = en1%ival + en2%cval
  end function add_ic
  real(default) function add_ri (en1, en2) result (y)
    type(eval_node_t), intent(in) :: en1, en2
    y = en1%rval + en2%ival
  end function add_ri
  complex(default) function add_ci (en1, en2) result (y)
    type(eval_node_t), intent(in) :: en1, en2
    y = en1%cval + en2%ival
  end function add_ci
  complex(default) function add_cr (en1, en2) result (y)
    type(eval_node_t), intent(in) :: en1, en2
    y = en1%cval + en2%rval
  end function add_cr
  complex(default) function add_rc (en1, en2) result (y)
    type(eval_node_t), intent(in) :: en1, en2
    y = en1%rval + en2%cval
  end function add_rc
  real(default) function add_rr (en1, en2) result (y)
    type(eval_node_t), intent(in) :: en1, en2
    y = en1%rval + en2%rval
  end function add_rr
  complex(default) function add_cc (en1, en2) result (y)
    type(eval_node_t), intent(in) :: en1, en2
    y = en1%cval + en2%cval
  end function add_cc

  integer function sub_ii (en1, en2) result (y)
    type(eval_node_t), intent(in) :: en1, en2
    y = en1%ival - en2%ival
  end function sub_ii
  real(default) function sub_ir (en1, en2) result (y)
    type(eval_node_t), intent(in) :: en1, en2
    y = en1%ival - en2%rval
  end function sub_ir
  real(default) function sub_ri (en1, en2) result (y)
    type(eval_node_t), intent(in) :: en1, en2
    y = en1%rval - en2%ival
  end function sub_ri  
  complex(default) function sub_ic (en1, en2) result (y)
    type(eval_node_t), intent(in) :: en1, en2
    y = en1%ival - en2%cval
  end function sub_ic
  complex(default) function sub_ci (en1, en2) result (y)
    type(eval_node_t), intent(in) :: en1, en2
    y = en1%cval - en2%ival
  end function sub_ci  
  complex(default) function sub_cr (en1, en2) result (y)
    type(eval_node_t), intent(in) :: en1, en2
    y = en1%cval - en2%rval
  end function sub_cr
  complex(default) function sub_rc (en1, en2) result (y)
    type(eval_node_t), intent(in) :: en1, en2
    y = en1%rval - en2%cval
  end function sub_rc
  real(default) function sub_rr (en1, en2) result (y)
    type(eval_node_t), intent(in) :: en1, en2
    y = en1%rval - en2%rval
  end function sub_rr
  complex(default) function sub_cc (en1, en2) result (y)
    type(eval_node_t), intent(in) :: en1, en2
    y = en1%cval - en2%cval
  end function sub_cc

  integer function mul_ii (en1, en2) result (y)
    type(eval_node_t), intent(in) :: en1, en2
    y = en1%ival * en2%ival
  end function mul_ii
  real(default) function mul_ir (en1, en2) result (y)
    type(eval_node_t), intent(in) :: en1, en2
    y = en1%ival * en2%rval
  end function mul_ir
  real(default) function mul_ri (en1, en2) result (y)
    type(eval_node_t), intent(in) :: en1, en2
    y = en1%rval * en2%ival
  end function mul_ri
  complex(default) function mul_ic (en1, en2) result (y)
    type(eval_node_t), intent(in) :: en1, en2
    y = en1%ival * en2%cval
  end function mul_ic
  complex(default) function mul_ci (en1, en2) result (y)
    type(eval_node_t), intent(in) :: en1, en2
    y = en1%cval * en2%ival
  end function mul_ci
  complex(default) function mul_rc (en1, en2) result (y)
    type(eval_node_t), intent(in) :: en1, en2
    y = en1%rval * en2%cval
  end function mul_rc
  complex(default) function mul_cr (en1, en2) result (y)
    type(eval_node_t), intent(in) :: en1, en2
    y = en1%cval * en2%rval
  end function mul_cr
  real(default) function mul_rr (en1, en2) result (y)
    type(eval_node_t), intent(in) :: en1, en2
    y = en1%rval * en2%rval
  end function mul_rr
  complex(default) function mul_cc (en1, en2) result (y)
    type(eval_node_t), intent(in) :: en1, en2
    y = en1%cval * en2%cval
  end function mul_cc

  integer function div_ii (en1, en2) result (y)
    type(eval_node_t), intent(in) :: en1, en2
    if (en2%ival == 0) then
       if (en1%ival >= 0) then
          call msg_warning ("division by zero: " // int2char (en1%ival) // &
             " / 0 ; result set to 0")
       else
          call msg_warning ("division by zero: (" // int2char (en1%ival) // &
             ") / 0 ; result set to 0")
       end if
       y = 0
       return
    end if
    y = en1%ival / en2%ival
  end function div_ii
  real(default) function div_ir (en1, en2) result (y)
    type(eval_node_t), intent(in) :: en1, en2
    y = en1%ival / en2%rval
  end function div_ir
  real(default) function div_ri (en1, en2) result (y)
    type(eval_node_t), intent(in) :: en1, en2
    y = en1%rval / en2%ival
  end function div_ri
  complex(default) function div_ic (en1, en2) result (y)
    type(eval_node_t), intent(in) :: en1, en2
    y = en1%ival / en2%cval
  end function div_ic
  complex(default) function div_ci (en1, en2) result (y)
    type(eval_node_t), intent(in) :: en1, en2
    y = en1%cval / en2%ival
  end function div_ci
  complex(default) function div_rc (en1, en2) result (y)
    type(eval_node_t), intent(in) :: en1, en2
    y = en1%rval / en2%cval
  end function div_rc
  complex(default) function div_cr (en1, en2) result (y)
    type(eval_node_t), intent(in) :: en1, en2
    y = en1%cval / en2%rval
  end function div_cr
  real(default) function div_rr (en1, en2) result (y)
    type(eval_node_t), intent(in) :: en1, en2
    y = en1%rval / en2%rval
  end function div_rr
  complex(default) function div_cc (en1, en2) result (y)
    type(eval_node_t), intent(in) :: en1, en2
    y = en1%cval / en2%cval
  end function div_cc

  integer function pow_ii (en1, en2) result (y)
    type(eval_node_t), intent(in) :: en1, en2
    integer :: a, b
    real(default) :: rres
    a = en1%ival
    b = en2%ival
    if ((a == 0) .and. (b < 0)) then
       call msg_warning ("division by zero: " // int2char (a) // &
          " ^ (" // int2char (b) // ") ; result set to 0")
       y = 0
       return
    end if
    rres = real(a, default) ** b
    y = rres
    if (real(y, default) /= rres) then
       if (b < 0) then
          call msg_warning ("result of all-integer operation " // &
                int2char (a) // " ^ (" // int2char (b) // &
                ") has been trucated to "// int2char (y), &
             [  var_str ("Chances are that you want to use " // &
                "reals instead of integers at this point.") ])
       else
          call msg_warning ("integer overflow in " // int2char (a) // &
                " ^ " // int2char (b) // " ; result is " // int2char (y), &
             [  var_str ("Using reals instead of integers might help.")])
       end if
    end if
  end function pow_ii
  real(default) function pow_ri (en1, en2) result (y)
    type(eval_node_t), intent(in) :: en1, en2
    y = en1%rval ** en2%ival
  end function pow_ri
  complex(default) function pow_ci (en1, en2) result (y)
    type(eval_node_t), intent(in) :: en1, en2
    y = en1%cval ** en2%ival
  end function pow_ci
  real(default) function pow_ir (en1, en2) result (y)
    type(eval_node_t), intent(in) :: en1, en2
    y = en1%ival ** en2%rval
  end function pow_ir
  real(default) function pow_rr (en1, en2) result (y)
    type(eval_node_t), intent(in) :: en1, en2
    y = en1%rval ** en2%rval
  end function pow_rr
  complex(default) function pow_cr (en1, en2) result (y)
    type(eval_node_t), intent(in) :: en1, en2
    y = en1%cval ** en2%rval
  end function pow_cr
  complex(default) function pow_ic (en1, en2) result (y)
    type(eval_node_t), intent(in) :: en1, en2
    y = en1%ival ** en2%cval
  end function pow_ic
  complex(default) function pow_rc (en1, en2) result (y)
    type(eval_node_t), intent(in) :: en1, en2
    y = en1%rval ** en2%cval
  end function pow_rc
  complex(default) function pow_cc (en1, en2) result (y)
    type(eval_node_t), intent(in) :: en1, en2
    y = en1%cval ** en2%cval
  end function pow_cc

  integer function max_ii (en1, en2) result (y)
    type(eval_node_t), intent(in) :: en1, en2
    y = max (en1%ival, en2%ival)
  end function max_ii
  real(default) function max_ir (en1, en2) result (y)
    type(eval_node_t), intent(in) :: en1, en2
    y = max (real (en1%ival, default), en2%rval)
  end function max_ir
  real(default) function max_ri (en1, en2) result (y)
    type(eval_node_t), intent(in) :: en1, en2
    y = max (en1%rval, real (en2%ival, default))
  end function max_ri
  real(default) function max_rr (en1, en2) result (y)
    type(eval_node_t), intent(in) :: en1, en2
    y = max (en1%rval, en2%rval)
  end function max_rr
  integer function min_ii (en1, en2) result (y)
    type(eval_node_t), intent(in) :: en1, en2
    y = min (en1%ival, en2%ival)
  end function min_ii
  real(default) function min_ir (en1, en2) result (y)
    type(eval_node_t), intent(in) :: en1, en2
    y = min (real (en1%ival, default), en2%rval)
  end function min_ir
  real(default) function min_ri (en1, en2) result (y)
    type(eval_node_t), intent(in) :: en1, en2
    y = min (en1%rval, real (en2%ival, default))
  end function min_ri
  real(default) function min_rr (en1, en2) result (y)
    type(eval_node_t), intent(in) :: en1, en2
    y = min (en1%rval, en2%rval)
  end function min_rr

  integer function mod_ii (en1, en2) result (y)
    type(eval_node_t), intent(in) :: en1, en2
    y = mod (en1%ival, en2%ival)
  end function mod_ii
  real(default) function mod_ir (en1, en2) result (y)
    type(eval_node_t), intent(in) :: en1, en2
    y = mod (real (en1%ival, default), en2%rval)
  end function mod_ir
  real(default) function mod_ri (en1, en2) result (y)
    type(eval_node_t), intent(in) :: en1, en2
    y = mod (en1%rval, real (en2%ival, default))
  end function mod_ri
  real(default) function mod_rr (en1, en2) result (y)
    type(eval_node_t), intent(in) :: en1, en2
    y = mod (en1%rval, en2%rval)
  end function mod_rr
  integer function modulo_ii (en1, en2) result (y)
    type(eval_node_t), intent(in) :: en1, en2
    y = modulo (en1%ival, en2%ival)
  end function modulo_ii
  real(default) function modulo_ir (en1, en2) result (y)
    type(eval_node_t), intent(in) :: en1, en2
    y = modulo (real (en1%ival, default), en2%rval)
  end function modulo_ir
  real(default) function modulo_ri (en1, en2) result (y)
    type(eval_node_t), intent(in) :: en1, en2
    y = modulo (en1%rval, real (en2%ival, default))
  end function modulo_ri
  real(default) function modulo_rr (en1, en2) result (y)
    type(eval_node_t), intent(in) :: en1, en2
    y = modulo (en1%rval, en2%rval)
  end function modulo_rr

\nwendcode{}\nwbegindocs{217}\nwdocspar
\subsubsection{Unary numeric functions}
\nwenddocs{}\nwbegincode{218}\moddef{Eval trees: procedures}\plusendmoddef
  real(default) function real_i (en) result (y)
    type(eval_node_t), intent(in) :: en
    y = en%ival
  end function real_i
  real(default) function real_c (en) result (y)
    type(eval_node_t), intent(in) :: en
    y = en%cval
  end function real_c
  integer function int_r (en) result (y)
    type(eval_node_t), intent(in) :: en
    y = en%rval
  end function int_r
  complex(default) function cmplx_i (en) result (y)
    type(eval_node_t), intent(in) :: en
    y = en%ival
  end function cmplx_i
  integer function int_c (en) result (y)
    type(eval_node_t), intent(in) :: en
    y = en%cval
  end function int_c 
  complex(default) function cmplx_r (en) result (y)
    type(eval_node_t), intent(in) :: en
    y = en%rval
  end function cmplx_r  
  integer function nint_r (en) result (y)
    type(eval_node_t), intent(in) :: en
    y = nint (en%rval)
  end function nint_r
  integer function floor_r (en) result (y)
    type(eval_node_t), intent(in) :: en
    y = floor (en%rval)
  end function floor_r
  integer function ceiling_r (en) result (y)
    type(eval_node_t), intent(in) :: en
    y = ceiling (en%rval)
  end function ceiling_r

  integer function neg_i (en) result (y)
    type(eval_node_t), intent(in) :: en
    y = - en%ival
  end function neg_i
  real(default) function neg_r (en) result (y)
    type(eval_node_t), intent(in) :: en
    y = - en%rval
  end function neg_r
  complex(default) function neg_c (en) result (y)
    type(eval_node_t), intent(in) :: en
    y = - en%cval
  end function neg_c
  integer function abs_i (en) result (y)
    type(eval_node_t), intent(in) :: en
    y = abs (en%ival)
  end function abs_i
  real(default) function abs_r (en) result (y)
    type(eval_node_t), intent(in) :: en
    y = abs (en%rval)
  end function abs_r
  real(default) function abs_c (en) result (y)
    type(eval_node_t), intent(in) :: en
    y = abs (en%cval)
  end function abs_c
  integer function sgn_i (en) result (y)
    type(eval_node_t), intent(in) :: en
    y = sign (1, en%ival)
  end function sgn_i
  real(default) function sgn_r (en) result (y)
    type(eval_node_t), intent(in) :: en
    y = sign (1._default, en%rval)
  end function sgn_r

  real(default) function sqrt_r (en) result (y)
    type(eval_node_t), intent(in) :: en
    y = sqrt (en%rval)
  end function sqrt_r
  real(default) function exp_r (en) result (y)
    type(eval_node_t), intent(in) :: en
    y = exp (en%rval)
  end function exp_r
  real(default) function log_r (en) result (y)
    type(eval_node_t), intent(in) :: en
    y = log (en%rval)
  end function log_r
  real(default) function log10_r (en) result (y)
    type(eval_node_t), intent(in) :: en
    y = log10 (en%rval)
  end function log10_r
  
  complex(default) function sqrt_c (en) result (y)
    type(eval_node_t), intent(in) :: en
    y = sqrt (en%cval)
  end function sqrt_c
  complex(default) function exp_c (en) result (y)
    type(eval_node_t), intent(in) :: en
    y = exp (en%cval)
  end function exp_c
  complex(default) function log_c (en) result (y)
    type(eval_node_t), intent(in) :: en
    y = log (en%cval)
  end function log_c

  real(default) function sin_r (en) result (y)
    type(eval_node_t), intent(in) :: en
    y = sin (en%rval)
  end function sin_r
  real(default) function cos_r (en) result (y)
    type(eval_node_t), intent(in) :: en
    y = cos (en%rval)
  end function cos_r
  real(default) function tan_r (en) result (y)
    type(eval_node_t), intent(in) :: en
    y = tan (en%rval)
  end function tan_r
  real(default) function asin_r (en) result (y)
    type(eval_node_t), intent(in) :: en
    y = asin (en%rval)
  end function asin_r
  real(default) function acos_r (en) result (y)
    type(eval_node_t), intent(in) :: en
    y = acos (en%rval)
  end function acos_r
  real(default) function atan_r (en) result (y)
    type(eval_node_t), intent(in) :: en
    y = atan (en%rval)
  end function atan_r
  
  complex(default) function sin_c (en) result (y)
    type(eval_node_t), intent(in) :: en
    y = sin (en%cval)
  end function sin_c
  complex(default) function cos_c (en) result (y)
    type(eval_node_t), intent(in) :: en
    y = cos (en%cval)
  end function cos_c
  
  real(default) function sinh_r (en) result (y)
    type(eval_node_t), intent(in) :: en
    y = sinh (en%rval)
  end function sinh_r
  real(default) function cosh_r (en) result (y)
    type(eval_node_t), intent(in) :: en
    y = cosh (en%rval)
  end function cosh_r
  real(default) function tanh_r (en) result (y)
    type(eval_node_t), intent(in) :: en
    y = tanh (en%rval)
  end function tanh_r
!!! These are F2008 additions but accepted by nagfor 5.3 and gfortran 4.6+
!!! Currently not used.
!   real(default) function asinh_r (en) result (y)
!     type(eval_node_t), intent(in) :: en
!     y = asinh (en%rval)
!   end function asinh_r
!   real(default) function acosh_r (en) result (y)
!     type(eval_node_t), intent(in) :: en
!     y = acosh (en%rval)
!   end function acosh_r
!   real(default) function atanh_r (en) result (y)
!     type(eval_node_t), intent(in) :: en
!     y = atanh (en%rval)
!   end function atanh_r
  
\nwendcode{}\nwbegindocs{219}\nwdocspar
\subsubsection{Binary logical functions}
Logical expressions:
\nwenddocs{}\nwbegincode{220}\moddef{Eval trees: procedures}\plusendmoddef
  logical function ignore_first_ll (en1, en2) result (y)
    type(eval_node_t), intent(in) :: en1, en2
    y = en2%lval
  end function ignore_first_ll
  logical function or_ll (en1, en2) result (y)
    type(eval_node_t), intent(in) :: en1, en2
    y = en1%lval .or. en2%lval
  end function or_ll
  logical function and_ll (en1, en2) result (y)
    type(eval_node_t), intent(in) :: en1, en2
    y = en1%lval .and. en2%lval
  end function and_ll

\nwendcode{}\nwbegindocs{221}Comparisons:
\nwenddocs{}\nwbegincode{222}\moddef{Eval trees: procedures}\plusendmoddef
  logical function comp_lt_ii (en1, en2) result (y)
    type(eval_node_t), intent(in) :: en1, en2
    y = en1%ival < en2%ival
  end function comp_lt_ii
  logical function comp_lt_ir (en1, en2) result (y)
    type(eval_node_t), intent(in) :: en1, en2
    y = en1%ival < en2%rval
  end function comp_lt_ir
  logical function comp_lt_ri (en1, en2) result (y)
    type(eval_node_t), intent(in) :: en1, en2
    y = en1%rval < en2%ival
  end function comp_lt_ri
  logical function comp_lt_rr (en1, en2) result (y)
    type(eval_node_t), intent(in) :: en1, en2
    y = en1%rval < en2%rval
  end function comp_lt_rr
  
  logical function comp_gt_ii (en1, en2) result (y)
    type(eval_node_t), intent(in) :: en1, en2
    y = en1%ival > en2%ival
  end function comp_gt_ii
  logical function comp_gt_ir (en1, en2) result (y)
    type(eval_node_t), intent(in) :: en1, en2
    y = en1%ival > en2%rval
  end function comp_gt_ir
  logical function comp_gt_ri (en1, en2) result (y)
    type(eval_node_t), intent(in) :: en1, en2
    y = en1%rval > en2%ival
  end function comp_gt_ri
  logical function comp_gt_rr (en1, en2) result (y)
    type(eval_node_t), intent(in) :: en1, en2
    y = en1%rval > en2%rval
  end function comp_gt_rr
  
  logical function comp_le_ii (en1, en2) result (y)
    type(eval_node_t), intent(in) :: en1, en2
    y = en1%ival <= en2%ival
  end function comp_le_ii
  logical function comp_le_ir (en1, en2) result (y)
    type(eval_node_t), intent(in) :: en1, en2
    y = en1%ival <= en2%rval
  end function comp_le_ir
  logical function comp_le_ri (en1, en2) result (y)
    type(eval_node_t), intent(in) :: en1, en2
    y = en1%rval <= en2%ival
  end function comp_le_ri
  logical function comp_le_rr (en1, en2) result (y)
    type(eval_node_t), intent(in) :: en1, en2
    y = en1%rval <= en2%rval
  end function comp_le_rr
  
  logical function comp_ge_ii (en1, en2) result (y)
    type(eval_node_t), intent(in) :: en1, en2
    y = en1%ival >= en2%ival
  end function comp_ge_ii
  logical function comp_ge_ir (en1, en2) result (y)
    type(eval_node_t), intent(in) :: en1, en2
    y = en1%ival >= en2%rval
  end function comp_ge_ir
  logical function comp_ge_ri (en1, en2) result (y)
    type(eval_node_t), intent(in) :: en1, en2
    y = en1%rval >= en2%ival
  end function comp_ge_ri
  logical function comp_ge_rr (en1, en2) result (y)
    type(eval_node_t), intent(in) :: en1, en2
    y = en1%rval >= en2%rval
  end function comp_ge_rr
  
  logical function comp_eq_ii (en1, en2) result (y)
    type(eval_node_t), intent(in) :: en1, en2
    y = en1%ival == en2%ival
  end function comp_eq_ii
  logical function comp_eq_ir (en1, en2) result (y)
    type(eval_node_t), intent(in) :: en1, en2
    y = en1%ival == en2%rval
  end function comp_eq_ir
  logical function comp_eq_ri (en1, en2) result (y)
    type(eval_node_t), intent(in) :: en1, en2
    y = en1%rval == en2%ival
  end function comp_eq_ri
  logical function comp_eq_rr (en1, en2) result (y)
    type(eval_node_t), intent(in) :: en1, en2
    y = en1%rval == en2%rval
  end function comp_eq_rr
  logical function comp_eq_ss (en1, en2) result (y)
    type(eval_node_t), intent(in) :: en1, en2
    y = en1%sval == en2%sval
  end function comp_eq_ss
  
  logical function comp_ne_ii (en1, en2) result (y)
    type(eval_node_t), intent(in) :: en1, en2
    y = en1%ival /= en2%ival
  end function comp_ne_ii
  logical function comp_ne_ir (en1, en2) result (y)
    type(eval_node_t), intent(in) :: en1, en2
    y = en1%ival /= en2%rval
  end function comp_ne_ir
  logical function comp_ne_ri (en1, en2) result (y)
    type(eval_node_t), intent(in) :: en1, en2
    y = en1%rval /= en2%ival
  end function comp_ne_ri
  logical function comp_ne_rr (en1, en2) result (y)
    type(eval_node_t), intent(in) :: en1, en2
    y = en1%rval /= en2%rval
  end function comp_ne_rr
  logical function comp_ne_ss (en1, en2) result (y)
    type(eval_node_t), intent(in) :: en1, en2
    y = en1%sval /= en2%sval
  end function comp_ne_ss
  
\nwendcode{}\nwbegindocs{223}Comparisons with tolerance:
\nwenddocs{}\nwbegincode{224}\moddef{Eval trees: procedures}\plusendmoddef
  logical function comp_se_ii (en1, en2) result (y)
    type(eval_node_t), intent(in) :: en1, en2
    if (associated (en1%tolerance)) then
       y = abs (en1%ival - en2%ival) <= en1%tolerance
    else
       y = en1%ival == en2%ival
    end if
  end function comp_se_ii
  logical function comp_se_ri (en1, en2) result (y)
    type(eval_node_t), intent(in) :: en1, en2
    if (associated (en1%tolerance)) then
       y = abs (en1%rval - en2%ival) <= en1%tolerance
    else
       y = en1%rval == en2%ival
    end if
  end function comp_se_ri
  logical function comp_se_ir (en1, en2) result (y)
    type(eval_node_t), intent(in) :: en1, en2
    if (associated (en1%tolerance)) then
       y = abs (en1%ival - en2%rval) <= en1%tolerance
    else
       y = en1%ival == en2%rval
    end if
  end function comp_se_ir
  logical function comp_se_rr (en1, en2) result (y)
    type(eval_node_t), intent(in) :: en1, en2
    if (associated (en1%tolerance)) then
       y = abs (en1%rval - en2%rval) <= en1%tolerance
    else
       y = en1%rval == en2%rval
    end if
  end function comp_se_rr
  logical function comp_ns_ii (en1, en2) result (y)
    type(eval_node_t), intent(in) :: en1, en2
    if (associated (en1%tolerance)) then
       y = abs (en1%ival - en2%ival) > en1%tolerance
    else
       y = en1%ival /= en2%ival
    end if
  end function comp_ns_ii
  logical function comp_ns_ri (en1, en2) result (y)
    type(eval_node_t), intent(in) :: en1, en2
    if (associated (en1%tolerance)) then
       y = abs (en1%rval - en2%ival) > en1%tolerance
    else
       y = en1%rval /= en2%ival
    end if
  end function comp_ns_ri
  logical function comp_ns_ir (en1, en2) result (y)
    type(eval_node_t), intent(in) :: en1, en2
    if (associated (en1%tolerance)) then
       y = abs (en1%ival - en2%rval) > en1%tolerance
    else
       y = en1%ival /= en2%rval
    end if
  end function comp_ns_ir
  logical function comp_ns_rr (en1, en2) result (y)
    type(eval_node_t), intent(in) :: en1, en2
    if (associated (en1%tolerance)) then
       y = abs (en1%rval - en2%rval) > en1%tolerance
    else
       y = en1%rval /= en2%rval
    end if
  end function comp_ns_rr
  
  logical function comp_ls_ii (en1, en2) result (y)
    type(eval_node_t), intent(in) :: en1, en2
    if (associated (en1%tolerance)) then
       y = en1%ival <= en2%ival + en1%tolerance
    else
       y = en1%ival <= en2%ival
    end if
  end function comp_ls_ii
  logical function comp_ls_ri (en1, en2) result (y)
    type(eval_node_t), intent(in) :: en1, en2
    if (associated (en1%tolerance)) then
       y = en1%rval <= en2%ival + en1%tolerance
    else
       y = en1%rval <= en2%ival
    end if
  end function comp_ls_ri
  logical function comp_ls_ir (en1, en2) result (y)
    type(eval_node_t), intent(in) :: en1, en2
    if (associated (en1%tolerance)) then
       y = en1%ival <= en2%rval + en1%tolerance
    else
       y = en1%ival <= en2%rval
    end if
  end function comp_ls_ir
  logical function comp_ls_rr (en1, en2) result (y)
    type(eval_node_t), intent(in) :: en1, en2
    if (associated (en1%tolerance)) then
       y = en1%rval <= en2%rval + en1%tolerance
    else
       y = en1%rval <= en2%rval
    end if
  end function comp_ls_rr

  logical function comp_ll_ii (en1, en2) result (y)
    type(eval_node_t), intent(in) :: en1, en2
    if (associated (en1%tolerance)) then
       y = en1%ival < en2%ival - en1%tolerance
    else
       y = en1%ival < en2%ival
    end if
  end function comp_ll_ii
  logical function comp_ll_ri (en1, en2) result (y)
    type(eval_node_t), intent(in) :: en1, en2
    if (associated (en1%tolerance)) then
       y = en1%rval < en2%ival - en1%tolerance
    else
       y = en1%rval < en2%ival
    end if
  end function comp_ll_ri
  logical function comp_ll_ir (en1, en2) result (y)
    type(eval_node_t), intent(in) :: en1, en2
    if (associated (en1%tolerance)) then
       y = en1%ival < en2%rval - en1%tolerance
    else
       y = en1%ival < en2%rval
    end if
  end function comp_ll_ir
  logical function comp_ll_rr (en1, en2) result (y)
    type(eval_node_t), intent(in) :: en1, en2
    if (associated (en1%tolerance)) then
       y = en1%rval < en2%rval - en1%tolerance
    else
       y = en1%rval < en2%rval
    end if
  end function comp_ll_rr

  logical function comp_gs_ii (en1, en2) result (y)
    type(eval_node_t), intent(in) :: en1, en2
    if (associated (en1%tolerance)) then
       y = en1%ival >= en2%ival - en1%tolerance
    else
       y = en1%ival >= en2%ival
    end if
  end function comp_gs_ii
  logical function comp_gs_ri (en1, en2) result (y)
    type(eval_node_t), intent(in) :: en1, en2
    if (associated (en1%tolerance)) then
       y = en1%rval >= en2%ival - en1%tolerance
    else
       y = en1%rval >= en2%ival
    end if
  end function comp_gs_ri
  logical function comp_gs_ir (en1, en2) result (y)
    type(eval_node_t), intent(in) :: en1, en2
    if (associated (en1%tolerance)) then
       y = en1%ival >= en2%rval - en1%tolerance
    else
       y = en1%ival >= en2%rval
    end if
  end function comp_gs_ir
  logical function comp_gs_rr (en1, en2) result (y)
    type(eval_node_t), intent(in) :: en1, en2
    if (associated (en1%tolerance)) then
       y = en1%rval >= en2%rval - en1%tolerance
    else
       y = en1%rval >= en2%rval
    end if
  end function comp_gs_rr

  logical function comp_gg_ii (en1, en2) result (y)
    type(eval_node_t), intent(in) :: en1, en2
    if (associated (en1%tolerance)) then
       y = en1%ival > en2%ival + en1%tolerance
    else
       y = en1%ival > en2%ival
    end if
  end function comp_gg_ii
  logical function comp_gg_ri (en1, en2) result (y)
    type(eval_node_t), intent(in) :: en1, en2
    if (associated (en1%tolerance)) then
       y = en1%rval > en2%ival + en1%tolerance
    else
       y = en1%rval > en2%ival
    end if
  end function comp_gg_ri
  logical function comp_gg_ir (en1, en2) result (y)
    type(eval_node_t), intent(in) :: en1, en2
    if (associated (en1%tolerance)) then
       y = en1%ival > en2%rval + en1%tolerance
    else
       y = en1%ival > en2%rval
    end if
  end function comp_gg_ir
  logical function comp_gg_rr (en1, en2) result (y)
    type(eval_node_t), intent(in) :: en1, en2
    if (associated (en1%tolerance)) then
       y = en1%rval > en2%rval + en1%tolerance
    else
       y = en1%rval > en2%rval
    end if
  end function comp_gg_rr

\nwendcode{}\nwbegindocs{225}\nwdocspar
\subsubsection{Unary logical functions}
\nwenddocs{}\nwbegincode{226}\moddef{Eval trees: procedures}\plusendmoddef
  logical function not_l (en) result (y)
    type(eval_node_t), intent(in) :: en
    y = .not. en%lval
  end function not_l

\nwendcode{}\nwbegindocs{227}\nwdocspar
\subsubsection{Unary PDG-array functions}
Make a PDG-array object from an integer.
\nwenddocs{}\nwbegincode{228}\moddef{Eval trees: procedures}\plusendmoddef
  subroutine pdg_i (pdg_array, en)
    type(pdg_array_t), intent(out) :: pdg_array
    type(eval_node_t), intent(in) :: en
    pdg_array = en%ival
  end subroutine pdg_i

\nwendcode{}\nwbegindocs{229}\nwdocspar
\subsubsection{Binary PDG-array functions}
Concatenate two PDG-array objects.
\nwenddocs{}\nwbegincode{230}\moddef{Eval trees: procedures}\plusendmoddef
  subroutine concat_cc (pdg_array, en1, en2)
    type(pdg_array_t), intent(out) :: pdg_array
    type(eval_node_t), intent(in) :: en1, en2
    pdg_array = en1%aval // en2%aval
  end subroutine concat_cc

\nwendcode{}\nwbegindocs{231}\nwdocspar
\subsubsection{Unary particle-list functions}
Combine all particles of the first argument.  If {\tt{}en0} is present,
create a mask which is true only for those particles that pass the
test.
\nwenddocs{}\nwbegincode{232}\moddef{Eval trees: procedures}\plusendmoddef
  subroutine collect_p (subevt, en1, en0)
    type(subevt_t), intent(inout) :: subevt
    type(eval_node_t), intent(in) :: en1
    type(eval_node_t), intent(inout), optional :: en0
    logical, dimension(:), allocatable :: mask1
    integer :: n, i
    n = subevt_get_length (en1%pval)
    allocate (mask1 (n))
    if (present (en0)) then
       do i = 1, n
          en0%index = i
          en0%prt1 = subevt_get_prt (en1%pval, i)
          call eval_node_evaluate (en0)
          mask1(i) = en0%lval
       end do
    else
       mask1 = .true.
    end if
    call subevt_collect (subevt, en1%pval, mask1)
  end subroutine collect_p

\eatline
\nwendcode{}\nwbegindocs{233}Cluster the particles of the first argument.  If {\tt{}en0} is present,
create a mask which is true only for those particles that pass the
test.
\nwenddocs{}\nwbegincode{234}\moddef{Eval trees: procedures}\plusendmoddef
  subroutine cluster_p (subevt, en1, en0)
    type(subevt_t), intent(inout) :: subevt
    type(eval_node_t), intent(in) :: en1
    type(eval_node_t), intent(inout), optional :: en0
    logical, dimension(:), allocatable :: mask1
    integer :: n, i
    ! Should not be hardcoded!
    type(jet_definition_t) :: jet_def
    call jet_def%init (en1%jet_algorithm, en1%jet_r)
    n = subevt_get_length (en1%pval)
    allocate (mask1 (n))
    if (present (en0)) then
       do i = 1, n
          en0%index = i
          en0%prt1 = subevt_get_prt (en1%pval, i)
          call eval_node_evaluate (en0)
          mask1(i) = en0%lval
       end do
    else
       mask1 = .true.
    end if
    call subevt_cluster (subevt, en1%pval, mask1, jet_def)
    call jet_def%final ()
  end subroutine cluster_p

\eatline
\nwendcode{}\nwbegindocs{235}Select all particles of the first argument.  If {\tt{}en0} is present,
create a mask which is true only for those particles that pass the
test.
\nwenddocs{}\nwbegincode{236}\moddef{Eval trees: procedures}\plusendmoddef
  subroutine select_p (subevt, en1, en0)
    type(subevt_t), intent(inout) :: subevt
    type(eval_node_t), intent(in) :: en1
    type(eval_node_t), intent(inout), optional :: en0
    logical, dimension(:), allocatable :: mask1
    integer :: n, i
    n = subevt_get_length (en1%pval)
    allocate (mask1 (n))
    if (present (en0)) then
       do i = 1, subevt_get_length (en1%pval)
          en0%index = i
          en0%prt1 = subevt_get_prt (en1%pval, i)
          call eval_node_evaluate (en0)
          mask1(i) = en0%lval
       end do
    else
       mask1 = .true.
    end if
    call subevt_select (subevt, en1%pval, mask1)
  end subroutine select_p

\eatline
\nwendcode{}\nwbegindocs{237}Extract the particle with index given by {\tt{}en0} from the argument
list.  Negative indices count from the end.  If {\tt{}en0} is absent,
extract the first particle.  The result is a list with a single entry,
or no entries if the original list was empty or if the index is out of
range.

This function has no counterpart with two arguments.
\nwenddocs{}\nwbegincode{238}\moddef{Eval trees: procedures}\plusendmoddef
  subroutine extract_p (subevt, en1, en0)
    type(subevt_t), intent(inout) :: subevt
    type(eval_node_t), intent(in) :: en1
    type(eval_node_t), intent(inout), optional :: en0
    integer :: index
    if (present (en0)) then
       call eval_node_evaluate (en0)
       select case (en0%result_type)
       case (V_INT);  index = en0%ival
       case default
          call eval_node_write (en0)
          call msg_fatal (" Index parameter of 'extract' must be integer.")
       end select
    else
       index = 1
    end if
    call subevt_extract (subevt, en1%pval, index)
  end subroutine extract_p

\eatline
\nwendcode{}\nwbegindocs{239}Sort the subevent according to the result of evaluating
{\tt{}en0}.  If {\tt{}en0} is absent, sort by default method (PDG code,
particles before antiparticles).
\nwenddocs{}\nwbegincode{240}\moddef{Eval trees: procedures}\plusendmoddef
  subroutine sort_p (subevt, en1, en0)
    type(subevt_t), intent(inout) :: subevt
    type(eval_node_t), intent(in) :: en1
    type(eval_node_t), intent(inout), optional :: en0
    integer, dimension(:), allocatable :: ival
    real(default), dimension(:), allocatable :: rval
    integer :: i, n
    n = subevt_get_length (en1%pval)
    if (present (en0)) then
       select case (en0%result_type)
       case (V_INT);  allocate (ival (n))
       case (V_REAL); allocate (rval (n))
       end select
       do i = 1, n
          en0%index = i
          en0%prt1 = subevt_get_prt (en1%pval, i)
          call eval_node_evaluate (en0)
          select case (en0%result_type)
          case (V_INT);  ival(i) = en0%ival
          case (V_REAL); rval(i) = en0%rval
          end select
       end do
       select case (en0%result_type)
       case (V_INT);  call subevt_sort (subevt, en1%pval, ival)
       case (V_REAL); call subevt_sort (subevt, en1%pval, rval)
       end select
    else
       call subevt_sort (subevt, en1%pval)
    end if
  end subroutine sort_p

\eatline
\nwendcode{}\nwbegindocs{241}The following functions return a logical value.  {\tt{}all} evaluates
to true if the condition {\tt{}en0} is true for all elements of the
subevent.  {\tt{}any} and {\tt{}no} are analogous.
\nwenddocs{}\nwbegincode{242}\moddef{Eval trees: procedures}\plusendmoddef
  function all_p (en1, en0) result (lval)
    logical :: lval
    type(eval_node_t), intent(in) :: en1
    type(eval_node_t), intent(inout) :: en0
    integer :: i, n
    n = subevt_get_length (en1%pval)
    lval = .true.
    do i = 1, n
       en0%index = i
       en0%prt1 = subevt_get_prt (en1%pval, i)
       call eval_node_evaluate (en0)
       lval = en0%lval
       if (.not. lval)  exit
    end do
  end function all_p
       
  function any_p (en1, en0) result (lval)
    logical :: lval
    type(eval_node_t), intent(in) :: en1
    type(eval_node_t), intent(inout) :: en0
    integer :: i, n
    n = subevt_get_length (en1%pval)
    lval = .false.
    do i = 1, n
       en0%index = i
       en0%prt1 = subevt_get_prt (en1%pval, i)
       call eval_node_evaluate (en0)
       lval = en0%lval
       if (lval)  exit
    end do
  end function any_p
       
  function no_p (en1, en0) result (lval)
    logical :: lval
    type(eval_node_t), intent(in) :: en1
    type(eval_node_t), intent(inout) :: en0
    integer :: i, n
    n = subevt_get_length (en1%pval)
    lval = .true.
    do i = 1, n
       en0%index = i
       en0%prt1 = subevt_get_prt (en1%pval, i)
       call eval_node_evaluate (en0)
       lval = .not. en0%lval
       if (lval)  exit
    end do
  end function no_p
       
\eatline
\nwendcode{}\nwbegindocs{243}This is the interface to user-supplied observables.  The node {\tt{}en0}
evaluates to a string that indicates the procedure name.  We search for the
procedure in the dynamic library and load it into the procedure pointer which
is then called.  {\tt{}en1} is the subevent on which the external code
operates.  The external function returns a {\tt{}c{\char95}int}, which we translate
into a real value.
\nwenddocs{}\nwbegincode{244}\moddef{Eval trees: procedures}\plusendmoddef
  function user_obs_int_p (en0, prt1) result (ival)
    integer :: ival
    type(eval_node_t), intent(inout) :: en0
    type(prt_t), intent(in) :: prt1
    type(string_t) :: name
    procedure(user_obs_int_unary), pointer :: user_obs
    call eval_node_evaluate (en0)
    if (en0%value_is_known) then
       select case (en0%result_type)
       case (V_STR);  name = en0%sval
       case default
          call msg_bug ("user_obs: procedure name must be a string")
          name = ""
       end select
       call c_f_procpointer (user_code_find_proc (name), user_obs)
       ival = user_obs (c_prt (prt1))
    else
       call eval_node_write_rec (en0)
       call msg_fatal ("User observable name is undefined")
    end if
  end function user_obs_int_p

  function user_obs_real_p (en0, prt1) result (rval)
    real(default) :: rval
    type(eval_node_t), intent(inout) :: en0
    type(prt_t), intent(in) :: prt1
    type(string_t) :: name
    procedure(user_obs_real_unary), pointer :: user_obs
    call eval_node_evaluate (en0)
    if (en0%value_is_known) then
       select case (en0%result_type)
       case (V_STR);  name = en0%sval
       case default
          call msg_bug ("user_obs: procedure name must be a string")
          name = ""
       end select
       call c_f_procpointer (user_code_find_proc (name), user_obs)
       rval = user_obs (c_prt (prt1))
    else
       call eval_node_write_rec (en0)
       call msg_fatal ("User observable name is undefined")
    end if
  end function user_obs_real_p

\eatline
\eatline
\nwendcode{}\nwbegindocs{245}This is the interface to user-supplied cut code.  The node {\tt{}en0}
evaluates to a string that indicates the procedure name.
\nwenddocs{}\nwbegincode{246}\moddef{Eval trees: procedures}\plusendmoddef
  function user_cut_p (en1, en0) result (lval)
    logical :: lval
    type(eval_node_t), intent(in) :: en1
    type(eval_node_t), intent(inout) :: en0
    type(string_t) :: name
    procedure(user_cut_fun), pointer :: user_cut
    call eval_node_evaluate (en0)
    select case (en0%result_type)
    case (V_STR);  name = en0%sval
    case default
       call msg_bug ("user_cut: procedure name must be a string")
       name = ""
    end select
    call c_f_procpointer (user_code_find_proc (name), user_cut)
    lval = user_cut (c_prt (en1%pval), &
                     int (subevt_get_length (en1%pval), kind=c_int)) &
           /= 0
  end function user_cut_p

\eatline
\nwendcode{}\nwbegindocs{247}The following function returns an integer value, namely the number
of particles for which the condition is true.  If there is no
condition, it returns simply the length of the subevent.

A function would be more natural.  Making it a subroutine avoids
another compiler bug (internal error in nagfor 5.2 (649)).  (See the
interface {\tt{}unary{\char95}evi}.)
\nwenddocs{}\nwbegincode{248}\moddef{Eval trees: procedures}\plusendmoddef
  subroutine count_a (ival, en1, en0)
    integer, intent(out) :: ival
    type(eval_node_t), intent(in) :: en1
    type(eval_node_t), intent(inout), optional :: en0
    integer :: i, n, count
    n = subevt_get_length (en1%pval)
    if (present (en0)) then
       count = 0
       do i = 1, n
          en0%index = i
          en0%prt1 = subevt_get_prt (en1%pval, i)
          call eval_node_evaluate (en0)
          if (en0%lval)  count = count + 1
       end do
       ival = count
    else
       ival = n
    end if
  end subroutine count_a
       
\eatline
\nwendcode{}\nwbegindocs{249}This evaluates a user-defined event-shape observable for the current
subevent.
\nwenddocs{}\nwbegincode{250}\moddef{Eval trees: procedures}\plusendmoddef
  subroutine user_event_shape_a (rval, en1, en0)
    real(default), intent(out) :: rval
    type(eval_node_t), intent(in) :: en1
    type(eval_node_t), intent(inout), optional :: en0
    type(string_t) :: name
    procedure(user_event_shape_fun), pointer :: user_event_shape
    if (.not. present (en0))  call msg_bug &
         ("user_event_shape called without procedure name")
    call eval_node_evaluate (en0)
    select case (en0%result_type)
    case (V_STR);  name = en0%sval
    case default
       call msg_bug ("user_event_shape: procedure name must be a string")
       name = ""
    end select
    call c_f_procpointer (user_code_find_proc (name), user_event_shape)
    rval = user_event_shape (c_prt (en1%pval), &
                             int (subevt_get_length (en1%pval), kind=c_int))
  end subroutine user_event_shape_a
       
\eatline
\nwendcode{}\nwbegindocs{251}\nwdocspar
\subsubsection{Binary particle-list functions}
This joins two subevents, stored in the evaluation nodes {\tt{}en1}
and {\tt{}en2}.  If {\tt{}en0} is also present, it amounts to a logical test
returning true or false for every pair of particles.  A particle of
the second list gets a mask entry only if it passes the test for all
particles of the first list.
\nwenddocs{}\nwbegincode{252}\moddef{Eval trees: procedures}\plusendmoddef
  subroutine join_pp (subevt, en1, en2, en0)
    type(subevt_t), intent(inout) :: subevt
    type(eval_node_t), intent(in) :: en1, en2
    type(eval_node_t), intent(inout), optional :: en0
    logical, dimension(:), allocatable :: mask2
    integer :: i, j, n1, n2
    n1 = subevt_get_length (en1%pval)
    n2 = subevt_get_length (en2%pval)
    allocate (mask2 (n2))
    mask2 = .true.
    if (present (en0)) then
       do i = 1, n1
          en0%index = i
          en0%prt1 = subevt_get_prt (en1%pval, i)
          do j = 1, n2
             en0%prt2 = subevt_get_prt (en2%pval, j)
             call eval_node_evaluate (en0)
             mask2(j) = mask2(j) .and. en0%lval
          end do
       end do
    end if
    call subevt_join (subevt, en1%pval, en2%pval, mask2)
  end subroutine join_pp

\eatline
\nwendcode{}\nwbegindocs{253}Combine two subevents, i.e., make a list of composite particles
built from all possible particle pairs from the two lists.  If {\tt{}en0}
is present, create a mask which is true only for those pairs that pass
the test.
\nwenddocs{}\nwbegincode{254}\moddef{Eval trees: procedures}\plusendmoddef
  subroutine combine_pp (subevt, en1, en2, en0)
    type(subevt_t), intent(inout) :: subevt
    type(eval_node_t), intent(in) :: en1, en2
    type(eval_node_t), intent(inout), optional :: en0
    logical, dimension(:,:), allocatable :: mask12
    integer :: i, j, n1, n2
    n1 = subevt_get_length (en1%pval)
    n2 = subevt_get_length (en2%pval)
    if (present (en0)) then
       allocate (mask12 (n1, n2))
       do i = 1, n1
          en0%index = i
          en0%prt1 = subevt_get_prt (en1%pval, i)
          do j = 1, n2
             en0%prt2 = subevt_get_prt (en2%pval, j)
             call eval_node_evaluate (en0)
             mask12(i,j) = en0%lval
          end do
       end do
       call subevt_combine (subevt, en1%pval, en2%pval, mask12)
    else
       call subevt_combine (subevt, en1%pval, en2%pval)
    end if
  end subroutine combine_pp

\eatline
\nwendcode{}\nwbegindocs{255}Combine all particles of the first argument.  If {\tt{}en0} is present,
create a mask which is true only for those particles that pass the
test w.r.t. all particles in the second argument.  If {\tt{}en0} is
absent, the second argument is ignored.
\nwenddocs{}\nwbegincode{256}\moddef{Eval trees: procedures}\plusendmoddef
  subroutine collect_pp (subevt, en1, en2, en0)
    type(subevt_t), intent(inout) :: subevt
    type(eval_node_t), intent(in) :: en1, en2
    type(eval_node_t), intent(inout), optional :: en0
    logical, dimension(:), allocatable :: mask1
    integer :: i, j, n1, n2
    n1 = subevt_get_length (en1%pval)
    n2 = subevt_get_length (en2%pval)
    allocate (mask1 (n1))
    mask1 = .true.
    if (present (en0)) then
       do i = 1, n1
          en0%index = i
          en0%prt1 = subevt_get_prt (en1%pval, i)
          do j = 1, n2
             en0%prt2 = subevt_get_prt (en2%pval, j)
             call eval_node_evaluate (en0)
             mask1(i) = mask1(i) .and. en0%lval
          end do
       end do
    end if
    call subevt_collect (subevt, en1%pval, mask1)
  end subroutine collect_pp

\eatline
\nwendcode{}\nwbegindocs{257}Select all particles of the first argument.  If {\tt{}en0} is present,
create a mask which is true only for those particles that pass the
test w.r.t. all particles in the second argument.  If {\tt{}en0} is
absent, the second argument is ignored, and the first argument is
transferred unchanged.  (This case is not very useful, of course.)
\nwenddocs{}\nwbegincode{258}\moddef{Eval trees: procedures}\plusendmoddef
  subroutine select_pp (subevt, en1, en2, en0)
    type(subevt_t), intent(inout) :: subevt
    type(eval_node_t), intent(in) :: en1, en2
    type(eval_node_t), intent(inout), optional :: en0
    logical, dimension(:), allocatable :: mask1
    integer :: i, j, n1, n2
    n1 = subevt_get_length (en1%pval)
    n2 = subevt_get_length (en2%pval)
    allocate (mask1 (n1))
    mask1 = .true.
    if (present (en0)) then
       do i = 1, n1
          en0%index = i
          en0%prt1 = subevt_get_prt (en1%pval, i)
          do j = 1, n2
             en0%prt2 = subevt_get_prt (en2%pval, j)
             call eval_node_evaluate (en0)
             mask1(i) = mask1(i) .and. en0%lval
          end do
       end do
    end if
    call subevt_select (subevt, en1%pval, mask1)
  end subroutine select_pp

\eatline
\nwendcode{}\nwbegindocs{259}Sort the first subevent according to the result of evaluating
{\tt{}en0}.  From the second subevent, only the first element is
taken as reference.  If {\tt{}en0} is absent, we sort by default method
(PDG code, particles before antiparticles).
\nwenddocs{}\nwbegincode{260}\moddef{Eval trees: procedures}\plusendmoddef
  subroutine sort_pp (subevt, en1, en2, en0)
    type(subevt_t), intent(inout) :: subevt
    type(eval_node_t), intent(in) :: en1, en2
    type(eval_node_t), intent(inout), optional :: en0
    integer, dimension(:), allocatable :: ival
    real(default), dimension(:), allocatable :: rval
    integer :: i, n1
    n1 = subevt_get_length (en1%pval)
    if (present (en0)) then
       select case (en0%result_type)
       case (V_INT);  allocate (ival (n1))
       case (V_REAL); allocate (rval (n1))
       end select
       do i = 1, n1
          en0%index = i
          en0%prt1 = subevt_get_prt (en1%pval, i)
          en0%prt2 = subevt_get_prt (en2%pval, 1)
          call eval_node_evaluate (en0)
          select case (en0%result_type)
          case (V_INT);  ival(i) = en0%ival
          case (V_REAL); rval(i) = en0%rval
          end select
       end do
       select case (en0%result_type)
       case (V_INT);  call subevt_sort (subevt, en1%pval, ival)
       case (V_REAL); call subevt_sort (subevt, en1%pval, rval)
       end select
    else
       call subevt_sort (subevt, en1%pval)
    end if
  end subroutine sort_pp

\eatline
\nwendcode{}\nwbegindocs{261}The following functions return a logical value.  {\tt{}all} evaluates
to true if the condition {\tt{}en0} is true for all valid element pairs
of both subevents.  Invalid pairs (with common {\tt{}src} entry) are
ignored.

{\tt{}any} and {\tt{}no} are analogous.
\nwenddocs{}\nwbegincode{262}\moddef{Eval trees: procedures}\plusendmoddef
  function all_pp (en1, en2, en0) result (lval)
    logical :: lval
    type(eval_node_t), intent(in) :: en1, en2
    type(eval_node_t), intent(inout) :: en0
    integer :: i, j, n1, n2
    n1 = subevt_get_length (en1%pval)
    n2 = subevt_get_length (en2%pval)
    lval = .true.
    LOOP1: do i = 1, n1
       en0%index = i
       en0%prt1 = subevt_get_prt (en1%pval, i)
       do j = 1, n2
          en0%prt2 = subevt_get_prt (en2%pval, j)
          if (are_disjoint (en0%prt1, en0%prt2)) then
             call eval_node_evaluate (en0)
             lval = en0%lval
             if (.not. lval)  exit LOOP1
          end if
       end do
    end do LOOP1
  end function all_pp
       
  function any_pp (en1, en2, en0) result (lval)
    logical :: lval
    type(eval_node_t), intent(in) :: en1, en2
    type(eval_node_t), intent(inout) :: en0
    integer :: i, j, n1, n2
    n1 = subevt_get_length (en1%pval)
    n2 = subevt_get_length (en2%pval)
    lval = .false.
    LOOP1: do i = 1, n1
       en0%index = i
       en0%prt1 = subevt_get_prt (en1%pval, i)
       do j = 1, n2
          en0%prt2 = subevt_get_prt (en2%pval, j)
          if (are_disjoint (en0%prt1, en0%prt2)) then
             call eval_node_evaluate (en0)
             lval = en0%lval
             if (lval)  exit LOOP1
          end if
       end do
    end do LOOP1
  end function any_pp
       
  function no_pp (en1, en2, en0) result (lval)
    logical :: lval
    type(eval_node_t), intent(in) :: en1, en2
    type(eval_node_t), intent(inout) :: en0
    integer :: i, j, n1, n2
    n1 = subevt_get_length (en1%pval)
    n2 = subevt_get_length (en2%pval)
    lval = .true.
    LOOP1: do i = 1, n1
       en0%index = i
       en0%prt1 = subevt_get_prt (en1%pval, i)
       do j = 1, n2
          en0%prt2 = subevt_get_prt (en2%pval, j)
          if (are_disjoint (en0%prt1, en0%prt2)) then
             call eval_node_evaluate (en0)
             lval = .not. en0%lval
             if (lval)  exit LOOP1
          end if
       end do
    end do LOOP1
  end function no_pp
       
\eatline
\nwendcode{}\nwbegindocs{263}This function evaluates an observable for a pair of particles.  From the two
particle lists, we take the first pair without {\tt{}src} overlap.  If there is
no valid pair, we revert the status of the value to unknown.
\nwenddocs{}\nwbegincode{264}\moddef{Eval trees: procedures}\plusendmoddef
  subroutine eval_pp (en1, en2, en0, rval, is_known)
    type(eval_node_t), intent(in) :: en1, en2
    type(eval_node_t), intent(inout) :: en0
    real(default), intent(out) :: rval
    logical, intent(out) :: is_known
    integer :: i, j, n1, n2
    n1 = subevt_get_length (en1%pval)
    n2 = subevt_get_length (en2%pval)
    rval = 0
    is_known = .false.
    LOOP1: do i = 1, n1
       en0%index = i
       en0%prt1 = subevt_get_prt (en1%pval, i)
       do j = 1, n2
          en0%prt2 = subevt_get_prt (en2%pval, j)
          if (are_disjoint (en0%prt1, en0%prt2)) then
             call eval_node_evaluate (en0)
             rval = en0%rval
             is_known = .true.
             exit LOOP1
          end if
       end do
    end do LOOP1
  end subroutine eval_pp

\eatline
\nwendcode{}\nwbegindocs{265}This is the interface to user-supplied observables.  The node {\tt{}en0}
evaluates to a string that indicates the procedure name.  We search for the
procedure in the dynamic library and load it into the procedure pointer which
is then called.  {\tt{}en1} is the subevent on which the external code
operates.  The external function returns a {\tt{}c{\char95}int}, which we translate
into a real value.
\nwenddocs{}\nwbegincode{266}\moddef{Eval trees: procedures}\plusendmoddef
  function user_obs_int_pp (en0, prt1, prt2) result (ival)
    integer :: ival
    type(eval_node_t), intent(inout) :: en0
    type(prt_t), intent(in) :: prt1, prt2
    type(string_t) :: name
    procedure(user_obs_int_binary), pointer :: user_obs
    call eval_node_evaluate (en0)
    if (en0%value_is_known) then
       select case (en0%result_type)
       case (V_STR);  name = en0%sval
       case default
          call msg_bug ("user_obs: procedure name must be a string")
          name = ""
       end select
       call c_f_procpointer (user_code_find_proc (name), user_obs)
       ival = user_obs (c_prt (prt1), c_prt (prt2))
    else
       call eval_node_write_rec (en0)
       call msg_fatal ("User observable name is undefined")
    end if
  end function user_obs_int_pp

  function user_obs_real_pp (en0, prt1, prt2) result (rval)
    real(default) :: rval
    type(eval_node_t), intent(inout) :: en0
    type(prt_t), intent(in) :: prt1, prt2
    type(string_t) :: name
    procedure(user_obs_real_binary), pointer :: user_obs
    call eval_node_evaluate (en0)
    if (en0%value_is_known) then
       select case (en0%result_type)
       case (V_STR);  name = en0%sval
       case default
          call msg_bug ("user_obs: procedure name must be a string")
          name = ""
       end select
       call c_f_procpointer (user_code_find_proc (name), user_obs)
       rval = user_obs (c_prt (prt1), c_prt (prt2))
    else
       call eval_node_write_rec (en0)
       call msg_fatal ("User observable name is undefined")
    end if
  end function user_obs_real_pp

\eatline
\eatline
\nwendcode{}\nwbegindocs{267}The following function returns an integer value, namely the number
of valid particle-pairs from both lists for which the condition is
true.  Invalid pairs (with common {\tt{}src} entry) are ignored.  If
there is no condition, it returns the number of valid particle pairs.

A function would be more natural.  Making it a subroutine avoids
another compiler bug (internal error in nagfor 5.2 (649)).  (See the
interface {\tt{}binary{\char95}num}.)
\nwenddocs{}\nwbegincode{268}\moddef{Eval trees: procedures}\plusendmoddef
  subroutine count_pp (ival, en1, en2, en0)
    integer, intent(out) :: ival
    type(eval_node_t), intent(in) :: en1, en2
    type(eval_node_t), intent(inout), optional :: en0
    integer :: i, j, n1, n2, count
    n1 = subevt_get_length (en1%pval)
    n2 = subevt_get_length (en2%pval)
    if (present (en0)) then
       count = 0
       do i = 1, n1
          en0%index = i
          en0%prt1 = subevt_get_prt (en1%pval, i)
          do j = 1, n2
             en0%prt2 = subevt_get_prt (en2%pval, j)
             if (are_disjoint (en0%prt1, en0%prt2)) then
                call eval_node_evaluate (en0)
                if (en0%lval)  count = count + 1
             end if
          end do
       end do
    else
       count = 0
       do i = 1, n1
          do j = 1, n2
             if (are_disjoint (subevt_get_prt (en1%pval, i), &
                               subevt_get_prt (en2%pval, j))) then
                count = count + 1
             end if
          end do
       end do
    end if
    ival = count
  end subroutine count_pp
       
\eatline
\nwendcode{}\nwbegindocs{269}This function makes up a subevent from the second argument
which consists only of particles which match the PDG code array (first
argument).
\nwenddocs{}\nwbegincode{270}\moddef{Eval trees: procedures}\plusendmoddef
  subroutine select_pdg_ca (subevt, en1, en2, en0)
    type(subevt_t), intent(inout) :: subevt
    type(eval_node_t), intent(in) :: en1, en2
    type(eval_node_t), intent(inout), optional :: en0
    if (present (en0)) then
       call subevt_select_pdg_code (subevt, en1%aval, en2%pval, en0%ival)
    else
       call subevt_select_pdg_code (subevt, en1%aval, en2%pval)
    end if
  end subroutine select_pdg_ca

\eatline
\nwendcode{}\nwbegindocs{271}\nwdocspar
\subsubsection{Binary string functions}
Currently, the only string operation is concatenation.
\nwenddocs{}\nwbegincode{272}\moddef{Eval trees: procedures}\plusendmoddef
  subroutine concat_ss (string, en1, en2)
    type(string_t), intent(out) :: string
    type(eval_node_t), intent(in) :: en1, en2
    string = en1%sval // en2%sval
  end subroutine concat_ss

\eatline
\nwendcode{}\nwbegindocs{273}\nwdocspar
\subsection{Compiling the parse tree}
The evaluation tree is built recursively by following a parse tree.
Debug option:
\nwenddocs{}\nwbegincode{274}\moddef{Eval trees: variables}\endmoddef
  logical, parameter :: debug = .false.

\eatline
\nwendcode{}\nwbegindocs{275}Evaluate an expression.  The requested type is given as an optional
argument; default is numeric (integer or real).
\nwenddocs{}\nwbegincode{276}\moddef{Eval trees: procedures}\plusendmoddef
  recursive subroutine eval_node_compile_genexpr &
       (en, pn, var_list, result_type)
    type(eval_node_t), pointer :: en
    type(parse_node_t), intent(in) :: pn
    type(var_list_t), intent(in), target :: var_list
    integer, intent(in), optional :: result_type
    if (debug) then
       print *, "read genexpr";  call parse_node_write (pn)
    end if
    if (present (result_type)) then
       select case (result_type)
       case (V_INT, V_REAL, V_CMPLX)
          call eval_node_compile_expr  (en, pn, var_list)
       case (V_LOG)
          call eval_node_compile_lexpr (en, pn, var_list)
       case (V_SEV)
          call eval_node_compile_pexpr (en, pn, var_list)
       case (V_PDG)
          call eval_node_compile_cexpr (en, pn, var_list)
       case (V_STR)
          call eval_node_compile_sexpr (en, pn, var_list)
       end select
    else
       call eval_node_compile_expr  (en, pn, var_list)
    end if
    if (debug) then
       call eval_node_write (en)
       print *, "done genexpr"
    end if
  end subroutine eval_node_compile_genexpr
  
\eatline
\nwendcode{}\nwbegindocs{277}\nwdocspar
\subsubsection{Numeric expressions}
This procedure compiles a numerical expression.  This is a single term
or a sum or difference of terms.  We have to account for all
combinations of integer and real arguments.  If both are constant, we
immediately do the calculation and allocate a constant node.
\nwenddocs{}\nwbegincode{278}\moddef{Eval trees: procedures}\plusendmoddef
  recursive subroutine eval_node_compile_expr (en, pn, var_list)
    type(eval_node_t), pointer :: en
    type(parse_node_t), intent(in) :: pn
    type(var_list_t), intent(in), target :: var_list
    type(parse_node_t), pointer :: pn_term, pn_addition, pn_op, pn_arg
    type(eval_node_t), pointer :: en1, en2
    type(string_t) :: key
    integer :: t1, t2, t
    if (debug) then
       print *, "read expr";  call parse_node_write (pn)
    end if
    pn_term => parse_node_get_sub_ptr (pn)
    select case (char (parse_node_get_rule_key (pn_term)))
    case ("term")
       call eval_node_compile_term (en, pn_term, var_list)
       pn_addition => parse_node_get_next_ptr (pn_term, tag="addition")
    case ("addition")
       en => null ()
       pn_addition => pn_term
    case default
       call parse_node_mismatch ("term|addition", pn)
    end select
    do while (associated (pn_addition))
       pn_op => parse_node_get_sub_ptr (pn_addition)
       pn_arg => parse_node_get_next_ptr (pn_op, tag="term")
       call eval_node_compile_term (en2, pn_arg, var_list)
       t2 = en2%result_type
       if (associated (en)) then
          en1 => en
          t1 = en1%result_type
       else
          allocate (en1)
          select case (t2)
          case (V_INT);  call eval_node_init_int  (en1, 0)
          case (V_REAL); call eval_node_init_real (en1, 0._default)
          case (V_CMPLX); call eval_node_init_cmplx (en1, cmplx &
                         (0._default, 0._default, kind=default))
          end select
          t1 = t2
       end if
       t = numeric_result_type (t1, t2)
       allocate (en)
       key = parse_node_get_key (pn_op)
       if (en1%type == EN_CONSTANT .and. en2%type == EN_CONSTANT) then
          select case (char (key))
          case ("+")
             select case (t1)
             case (V_INT)
                select case (t2)
                case (V_INT);  call eval_node_init_int  (en, add_ii (en1, en2))
                case (V_REAL); call eval_node_init_real (en, add_ir (en1, en2))
                case (V_CMPLX); call eval_node_init_cmplx (en, add_ic (en1, en2))
                end select
             case (V_REAL)
                select case (t2)
                case (V_INT);  call eval_node_init_real (en, add_ri (en1, en2))
                case (V_REAL); call eval_node_init_real (en, add_rr (en1, en2))
                case (V_CMPLX); call eval_node_init_cmplx (en, add_rc (en1, en2))
                end select
              case (V_CMPLX)
                select case (t2)
                case (V_INT);  call eval_node_init_cmplx (en, add_ci (en1, en2))
                case (V_REAL); call eval_node_init_cmplx (en, add_cr (en1, en2))
                case (V_CMPLX); call eval_node_init_cmplx (en, add_cc (en1, en2))
                end select
             end select
          case ("-")
             select case (t1)
             case (V_INT)
                select case (t2)
                case (V_INT);  call eval_node_init_int  (en, sub_ii (en1, en2))
                case (V_REAL); call eval_node_init_real (en, sub_ir (en1, en2))
                case (V_CMPLX); call eval_node_init_cmplx (en, sub_ic (en1, en2))                
                end select
             case (V_REAL)
                select case (t2)
                case (V_INT);  call eval_node_init_real (en, sub_ri (en1, en2))
                case (V_REAL); call eval_node_init_real (en, sub_rr (en1, en2))
                case (V_CMPLX); call eval_node_init_cmplx (en, sub_rc (en1, en2))                
                end select
             case (V_CMPLX)
                select case (t2)
                case (V_INT);  call eval_node_init_cmplx (en, sub_ci (en1, en2))
                case (V_REAL); call eval_node_init_cmplx (en, sub_cr (en1, en2))
                case (V_CMPLX); call eval_node_init_cmplx (en, sub_cc (en1, en2))                
                end select
             end select
          end select
          call eval_node_final_rec (en1)
          call eval_node_final_rec (en2)
          deallocate (en1, en2)
       else   
          call eval_node_init_branch (en, key, t, en1, en2)
          select case (char (key))
          case ("+")
             select case (t1)
             case (V_INT)
                select case (t2)
                case (V_INT);   call eval_node_set_op2_int  (en, add_ii)
                case (V_REAL);  call eval_node_set_op2_real (en, add_ir)
                case (V_CMPLX);  call eval_node_set_op2_cmplx (en, add_ic)
                end select
             case (V_REAL)
                select case (t2)
                case (V_INT);   call eval_node_set_op2_real (en, add_ri)
                case (V_REAL);  call eval_node_set_op2_real (en, add_rr)
                case (V_CMPLX);  call eval_node_set_op2_cmplx (en, add_rc)
                end select
             case (V_CMPLX)
                select case (t2)
                case (V_INT);   call eval_node_set_op2_cmplx (en, add_ci)
                case (V_REAL);  call eval_node_set_op2_cmplx (en, add_cr)
                case (V_CMPLX);  call eval_node_set_op2_cmplx (en, add_cc)
                end select
             end select
          case ("-")
             select case (t1)
             case (V_INT)
                select case (t2)
                case (V_INT);   call eval_node_set_op2_int  (en, sub_ii)
                case (V_REAL);  call eval_node_set_op2_real (en, sub_ir)
                case (V_CMPLX);  call eval_node_set_op2_cmplx (en, sub_ic)
                end select
             case (V_REAL)
                select case (t2)
                case (V_INT);   call eval_node_set_op2_real (en, sub_ri)
                case (V_REAL);  call eval_node_set_op2_real (en, sub_rr)
                case (V_CMPLX);  call eval_node_set_op2_cmplx (en, sub_rc)
                end select
             case (V_CMPLX)
                select case (t2)
                case (V_INT);   call eval_node_set_op2_cmplx (en, sub_ci)
                case (V_REAL);  call eval_node_set_op2_cmplx (en, sub_cr)
                case (V_CMPLX);  call eval_node_set_op2_cmplx (en, sub_cc)
                end select
             end select
          end select
       end if
       pn_addition => parse_node_get_next_ptr (pn_addition)
    end do
    if (debug) then
       call eval_node_write (en)
       print *, "done expr"
    end if
  end subroutine eval_node_compile_expr

\eatline
\nwendcode{}\nwbegincode{279}\moddef{Eval trees: procedures}\plusendmoddef
  recursive subroutine eval_node_compile_term (en, pn, var_list)
    type(eval_node_t), pointer :: en
    type(parse_node_t), intent(in) :: pn
    type(var_list_t), intent(in), target :: var_list
    type(parse_node_t), pointer :: pn_factor, pn_multiplication, pn_op, pn_arg
    type(eval_node_t), pointer :: en1, en2
    type(string_t) :: key
    integer :: t1, t2, t
    if (debug) then
       print *, "read term";  call parse_node_write (pn)
    end if
    pn_factor => parse_node_get_sub_ptr (pn, tag="factor")
    call eval_node_compile_factor (en, pn_factor, var_list)
    pn_multiplication => &
         parse_node_get_next_ptr (pn_factor, tag="multiplication")
    do while (associated (pn_multiplication))
       pn_op => parse_node_get_sub_ptr (pn_multiplication)
       pn_arg => parse_node_get_next_ptr (pn_op, tag="factor")
       en1 => en
       call eval_node_compile_factor (en2, pn_arg, var_list)
       t1 = en1%result_type
       t2 = en2%result_type
       t = numeric_result_type (t1, t2)
       allocate (en)
       key = parse_node_get_key (pn_op)
       if (en1%type == EN_CONSTANT .and. en2%type == EN_CONSTANT) then
          select case (char (key))
          case ("*")
             select case (t1)
             case (V_INT)
                select case (t2)
                case (V_INT);  call eval_node_init_int  (en, mul_ii (en1, en2))
                case (V_REAL); call eval_node_init_real (en, mul_ir (en1, en2))
                case (V_CMPLX); call eval_node_init_cmplx (en, mul_ic (en1, en2))
                end select
             case (V_REAL)
                select case (t2)
                case (V_INT);  call eval_node_init_real (en, mul_ri (en1, en2))
                case (V_REAL); call eval_node_init_real (en, mul_rr (en1, en2))
                case (V_CMPLX); call eval_node_init_cmplx (en, mul_rc (en1, en2))
                end select
              case (V_CMPLX)
                select case (t2)
                case (V_INT);  call eval_node_init_cmplx (en, mul_ci (en1, en2))
                case (V_REAL); call eval_node_init_cmplx (en, mul_cr (en1, en2))
                case (V_CMPLX); call eval_node_init_cmplx (en, mul_cc (en1, en2))
                end select
             end select
          case ("/")
             select case (t1)
             case (V_INT)
                select case (t2)
                case (V_INT);  call eval_node_init_int  (en, div_ii (en1, en2))
                case (V_REAL); call eval_node_init_real (en, div_ir (en1, en2))
                case (V_CMPLX); call eval_node_init_real (en, div_ir (en1, en2))
                end select
             case (V_REAL)
                select case (t2)
                case (V_INT);  call eval_node_init_real (en, div_ri (en1, en2))
                case (V_REAL); call eval_node_init_real (en, div_rr (en1, en2))
                case (V_CMPLX); call eval_node_init_cmplx (en, div_rc (en1, en2))
                end select
             case (V_CMPLX)
                select case (t2)
                case (V_INT);  call eval_node_init_cmplx (en, div_ci (en1, en2))
                case (V_REAL); call eval_node_init_cmplx (en, div_cr (en1, en2))
                case (V_CMPLX); call eval_node_init_cmplx (en, div_cc (en1, en2))
                end select
             end select
          end select
          call eval_node_final_rec (en1)
          call eval_node_final_rec (en2)
          deallocate (en1, en2)
       else
          call eval_node_init_branch (en, key, t, en1, en2)
          select case (char (key))
          case ("*")
             select case (t1)
             case (V_INT)
                select case (t2)
                case (V_INT);   call eval_node_set_op2_int  (en, mul_ii)
                case (V_REAL);  call eval_node_set_op2_real (en, mul_ir)
                case (V_CMPLX);  call eval_node_set_op2_cmplx (en, mul_ic)
                end select
             case (V_REAL)
                select case (t2)
                case (V_INT);   call eval_node_set_op2_real (en, mul_ri)
                case (V_REAL);  call eval_node_set_op2_real (en, mul_rr)
                case (V_CMPLX);  call eval_node_set_op2_cmplx (en, mul_rc)                
                end select
             case (V_CMPLX)
                select case (t2)
                case (V_INT);   call eval_node_set_op2_cmplx (en, mul_ci)
                case (V_REAL);  call eval_node_set_op2_cmplx (en, mul_cr)
                case (V_CMPLX);  call eval_node_set_op2_cmplx (en, mul_cc)                
                end select
             end select
          case ("/")
             select case (t1)
             case (V_INT)
                select case (t2)
                case (V_INT);   call eval_node_set_op2_int  (en, div_ii)
                case (V_REAL);  call eval_node_set_op2_real (en, div_ir)
                case (V_CMPLX);  call eval_node_set_op2_cmplx (en, div_ic)                
                end select
             case (V_REAL)
                select case (t2)
                case (V_INT);   call eval_node_set_op2_real (en, div_ri)
                case (V_REAL);  call eval_node_set_op2_real (en, div_rr)
                case (V_CMPLX);  call eval_node_set_op2_cmplx (en, div_rc)                
                end select
             case (V_CMPLX)
                select case (t2)
                case (V_INT);   call eval_node_set_op2_cmplx (en, div_ci)
                case (V_REAL);  call eval_node_set_op2_cmplx (en, div_cr)
                case (V_CMPLX);  call eval_node_set_op2_cmplx (en, div_cc)                
                end select
             end select
          end select
       end if
       pn_multiplication => parse_node_get_next_ptr (pn_multiplication)
    end do
    if (debug) then
       call eval_node_write (en)
       print *, "done term"
    end if
  end subroutine eval_node_compile_term

\eatline
\nwendcode{}\nwbegincode{280}\moddef{Eval trees: procedures}\plusendmoddef
  recursive subroutine eval_node_compile_factor (en, pn, var_list)
    type(eval_node_t), pointer :: en
    type(parse_node_t), intent(in) :: pn
    type(var_list_t), intent(in), target :: var_list
    type(parse_node_t), pointer :: pn_value, pn_exponentiation, pn_op, pn_arg
    type(eval_node_t), pointer :: en1, en2
    type(string_t) :: key
    integer :: t1, t2, t
    if (debug) then
       print *, "read factor";  call parse_node_write (pn)
    end if
    pn_value => parse_node_get_sub_ptr (pn)
    call eval_node_compile_signed_value (en, pn_value, var_list)
    pn_exponentiation => &
         parse_node_get_next_ptr (pn_value, tag="exponentiation")
    if (associated (pn_exponentiation)) then
       pn_op => parse_node_get_sub_ptr (pn_exponentiation)
       pn_arg => parse_node_get_next_ptr (pn_op)
       en1 => en
       call eval_node_compile_signed_value (en2, pn_arg, var_list)
       t1 = en1%result_type
       t2 = en2%result_type
       t = numeric_result_type (t1, t2)
       allocate (en)
       key = parse_node_get_key (pn_op)
       if (en1%type == EN_CONSTANT .and. en2%type == EN_CONSTANT) then
          select case (t1)
          case (V_INT)
             select case (t2)
             case (V_INT);   call eval_node_init_int   (en, pow_ii (en1, en2))
             case (V_REAL);  call eval_node_init_real  (en, pow_ir (en1, en2))
             case (V_CMPLX); call eval_node_init_cmplx (en, pow_ic (en1, en2))
             end select
          case (V_REAL)
             select case (t2)
             case (V_INT);   call eval_node_init_real  (en, pow_ri (en1, en2))
             case (V_REAL);  call eval_node_init_real  (en, pow_rr (en1, en2))
             case (V_CMPLX); call eval_node_init_cmplx (en, pow_rc (en1, en2))
             end select
          case (V_CMPLX)
             select case (t2)
             case (V_INT);   call eval_node_init_cmplx (en, pow_ci (en1, en2))
             case (V_REAL);  call eval_node_init_cmplx (en, pow_cr (en1, en2))
             case (V_CMPLX); call eval_node_init_cmplx (en, pow_cc (en1, en2))
             end select
          end select
          call eval_node_final_rec (en1)
          call eval_node_final_rec (en2)
          deallocate (en1, en2)
       else
          call eval_node_init_branch (en, key, t, en1, en2)
          select case (t1)
          case (V_INT)
             select case (t2)
             case (V_INT);   call eval_node_set_op2_int  (en, pow_ii)
             case (V_REAL,V_CMPLX);  call eval_type_error (pn, "exponentiation", t1)
             end select
          case (V_REAL)
             select case (t2)
             case (V_INT);   call eval_node_set_op2_real (en, pow_ri)
             case (V_REAL);  call eval_node_set_op2_real (en, pow_rr)
             case (V_CMPLX);  call eval_type_error (pn, "exponentiation", t1)
             end select
          case (V_CMPLX)
             select case (t2)
             case (V_INT);   call eval_node_set_op2_cmplx (en, pow_ci)
             case (V_REAL);  call eval_node_set_op2_cmplx (en, pow_cr)
             case (V_CMPLX);  call eval_node_set_op2_cmplx (en, pow_cc)
             end select
          end select
       end if
    end if
    if (debug) then
       call eval_node_write (en)
       print *, "done factor"
    end if
  end subroutine eval_node_compile_factor

\eatline
\nwendcode{}\nwbegincode{281}\moddef{Eval trees: procedures}\plusendmoddef
  recursive subroutine eval_node_compile_signed_value (en, pn, var_list)
    type(eval_node_t), pointer :: en
    type(parse_node_t), intent(in) :: pn
    type(var_list_t), intent(in), target :: var_list
    type(parse_node_t), pointer :: pn_arg
    type(eval_node_t), pointer :: en1
    integer :: t
    if (debug) then
       print *, "read signed value";  call parse_node_write (pn)
    end if
    select case (char (parse_node_get_rule_key (pn)))
    case ("signed_value")
       pn_arg => parse_node_get_sub_ptr (pn, 2)
       call eval_node_compile_value (en1, pn_arg, var_list)
       t = en1%result_type
       allocate (en)
       if (en1%type == EN_CONSTANT) then
          select case (t)
          case (V_INT);  call eval_node_init_int  (en, neg_i (en1))
          case (V_REAL); call eval_node_init_real (en, neg_r (en1))
          case (V_CMPLX); call eval_node_init_cmplx (en, neg_c (en1))          
          end select
          call eval_node_final_rec (en1)
          deallocate (en1)
       else
          call eval_node_init_branch (en, var_str ("-"), t, en1)
          select case (t)
          case (V_INT);  call eval_node_set_op1_int  (en, neg_i)
          case (V_REAL); call eval_node_set_op1_real (en, neg_r)
          case (V_CMPLX); call eval_node_set_op1_cmplx (en, neg_c)          
          end select
       end if
    case default
       call eval_node_compile_value (en, pn, var_list)
    end select
    if (debug) then
       call eval_node_write (en)
       print *, "done signed value"
    end if
  end subroutine eval_node_compile_signed_value

\eatline
\nwendcode{}\nwbegindocs{282}Integer, real and complex values have an optional unit.  The unit is
extracted and applied immediately.  An integer with unit evaluates to
a real constant.
\nwenddocs{}\nwbegincode{283}\moddef{Eval trees: procedures}\plusendmoddef
  recursive subroutine eval_node_compile_value (en, pn, var_list)
    type(eval_node_t), pointer :: en
    type(parse_node_t), intent(in) :: pn
    type(var_list_t), intent(in), target :: var_list
    if (debug) then
       print *, "read value";  call parse_node_write (pn)
    end if
    select case (char (parse_node_get_rule_key (pn)))
    case ("integer_value", "real_value", "complex_value")
       call eval_node_compile_numeric_value (en, pn)
    case ("pi")
       call eval_node_compile_constant (en, pn)
    case ("I")
       call eval_node_compile_constant (en, pn)       
    case ("variable")
       call eval_node_compile_variable (en, pn, var_list)
    case ("result")
       call eval_node_compile_result (en, pn, var_list)
    case ("user_observable")
       call eval_node_compile_user_observable (en, pn, var_list)
    case ("expr")
       call eval_node_compile_expr (en, pn, var_list)
    case ("block_expr")
       call eval_node_compile_block_expr (en, pn, var_list)
    case ("conditional_expr")
       call eval_node_compile_conditional (en, pn, var_list)
    case ("unary_function")
       call eval_node_compile_unary_function (en, pn, var_list)
    case ("binary_function")
       call eval_node_compile_binary_function (en, pn, var_list)
    case ("eval_fun")
       call eval_node_compile_eval_function (en, pn, var_list)
    case ("count_fun", "user_event_fun")
       call eval_node_compile_numeric_function (en, pn, var_list)
    case default
       call parse_node_mismatch &
            ("integer|real|complex|constant|variable|" // &
             "expr|block_expr|conditional_expr|" // &
             "unary_function|binary_function|numeric_pexpr", pn)
    end select
    if (debug) then
       call eval_node_write (en)
       print *, "done value"
    end if
  end subroutine eval_node_compile_value

\eatline
\nwendcode{}\nwbegindocs{284}Real, complex and integer values are numeric literals with an
optional unit attached.  In case of an integer, the unit actually
makes it a real value in disguise.  The signed version of real values
is not possible in generic expressions; it is a special case for
numeric constants in model files (see below). We do not introduce
signed versions of complex values.
\nwenddocs{}\nwbegincode{285}\moddef{Eval trees: procedures}\plusendmoddef
  subroutine eval_node_compile_numeric_value (en, pn)
    type(eval_node_t), pointer :: en
    type(parse_node_t), intent(in), target :: pn
    type(parse_node_t), pointer :: pn_val, pn_unit
    allocate (en)
    pn_val => parse_node_get_sub_ptr (pn)
    pn_unit => parse_node_get_next_ptr (pn_val)
    select case (char (parse_node_get_rule_key (pn)))
    case ("integer_value")
       if (associated (pn_unit)) then
          call eval_node_init_real (en, &
               parse_node_get_integer (pn_val) * parse_node_get_unit (pn_unit))
       else
          call eval_node_init_int (en, parse_node_get_integer (pn_val))
       end if
    case ("real_value")
       if (associated (pn_unit)) then
          call eval_node_init_real (en, &
               parse_node_get_real (pn_val) * parse_node_get_unit (pn_unit))
       else
          call eval_node_init_real (en, parse_node_get_real (pn_val))
       end if
    case ("complex_value")
       if (associated (pn_unit)) then
          call eval_node_init_cmplx (en, &
               parse_node_get_cmplx (pn_val) * parse_node_get_unit (pn_unit))
       else
          call eval_node_init_cmplx (en, parse_node_get_cmplx (pn_val))
       end if
    case ("neg_real_value")
       pn_val => parse_node_get_sub_ptr (parse_node_get_sub_ptr (pn, 2))
       pn_unit => parse_node_get_next_ptr (pn_val)
       if (associated (pn_unit)) then
          call eval_node_init_real (en, &
               - parse_node_get_real (pn_val) * parse_node_get_unit (pn_unit))
       else
          call eval_node_init_real (en, - parse_node_get_real (pn_val))
       end if
    case ("pos_real_value")
       pn_val => parse_node_get_sub_ptr (parse_node_get_sub_ptr (pn, 2))
       pn_unit => parse_node_get_next_ptr (pn_val)
       if (associated (pn_unit)) then
          call eval_node_init_real (en, &
               parse_node_get_real (pn_val) * parse_node_get_unit (pn_unit))
       else
          call eval_node_init_real (en, parse_node_get_real (pn_val))
       end if
    case default
       call parse_node_mismatch &
       ("integer_value|real_value|complex_value|neg_real_value|pos_real_value", pn)
    end select
  end subroutine eval_node_compile_numeric_value
    
\eatline
\nwendcode{}\nwbegindocs{286}These are the units, predefined and hardcoded.  The default energy
unit is GeV, the default angular unit is radians.  We include units
for observables of dimension energy squared.  Luminosities are
normalized in inverse femtobarns.
\nwenddocs{}\nwbegincode{287}\moddef{Eval trees: procedures}\plusendmoddef
  function parse_node_get_unit (pn) result (factor)
    real(default) :: factor
    real(default) :: unit
    type(parse_node_t), intent(in) :: pn
    type(parse_node_t), pointer :: pn_unit, pn_unit_power 
    type(parse_node_t), pointer :: pn_frac, pn_num, pn_int, pn_div, pn_den
    integer :: num, den
    pn_unit => parse_node_get_sub_ptr (pn)
    select case (char (parse_node_get_key (pn_unit)))
    case ("TeV");  unit = 1.e3_default
    case ("GeV");  unit = 1
    case ("MeV");  unit = 1.e-3_default
    case ("keV");  unit = 1.e-6_default
    case ("eV");   unit = 1.e-9_default
    case ("meV");   unit = 1.e-12_default
    case ("nbarn");  unit = 1.e6_default
    case ("pbarn");  unit = 1.e3_default
    case ("fbarn");  unit = 1
    case ("abarn");  unit = 1.e-3_default
    case ("rad");     unit = 1
    case ("mrad");    unit = 1.e-3_default
    case ("degree");  unit = degree
    case ("%");  unit = 1.e-2_default
    case default
       call msg_bug (" Unit '" // &
            char (parse_node_get_key (pn)) // "' is undefined.")
    end select
    pn_unit_power => parse_node_get_next_ptr (pn_unit)
    if (associated (pn_unit_power)) then
       pn_frac => parse_node_get_sub_ptr (pn_unit_power, 2)
       pn_num => parse_node_get_sub_ptr (pn_frac)
       select case (char (parse_node_get_rule_key (pn_num)))
       case ("neg_int")
          pn_int => parse_node_get_sub_ptr (pn_num, 2)
          num = - parse_node_get_integer (pn_int)
       case ("pos_int")
          pn_int => parse_node_get_sub_ptr (pn_num, 2)
          num = parse_node_get_integer (pn_int)
       case ("integer_literal")
          num = parse_node_get_integer (pn_num)
       case default
          call parse_node_mismatch ("neg_int|pos_int|integer_literal", pn_num)
       end select
       pn_div => parse_node_get_next_ptr (pn_num)
       if (associated (pn_div)) then
          pn_den => parse_node_get_sub_ptr (pn_div, 2)
          den = parse_node_get_integer (pn_den)
       else
          den = 1
       end if
    else
       num = 1
       den = 1
    end if
    factor = unit ** (real (num, default) / den)
  end function parse_node_get_unit

\eatline
\nwendcode{}\nwbegindocs{288}There are only two predefined constants, but more can be added easily.
\nwenddocs{}\nwbegincode{289}\moddef{Eval trees: procedures}\plusendmoddef
  subroutine eval_node_compile_constant (en, pn)
    type(eval_node_t), pointer :: en
    type(parse_node_t), intent(in) :: pn
    if (debug) then
       print *, "read constant";  call parse_node_write (pn)
    end if
    allocate (en)
    select case (char (parse_node_get_key (pn)))
    case ("pi");     call eval_node_init_real (en, pi)
    case ("I");      call eval_node_init_cmplx (en, imago)    
    case default
       call parse_node_mismatch ("pi or I", pn)
    end select
    if (debug) then
       call eval_node_write (en)
       print *, "done constant"
    end if
  end subroutine eval_node_compile_constant

\eatline
\nwendcode{}\nwbegindocs{290}Compile a variable, with or without a specified type.
Take the list of variables, look for the name and make a node with a
pointer to the value.  If no type is provided, the variable is
numeric, and the stored value determines whether it is real or integer.

We explicitly demand that the variable is defined, so we do not accidentally
point to variables that are declared only later in the script but have come
into existence in a previous compilation pass.

Variables may actually be anonymous, these are expressions in disguise.  In
that case, the expression replaces the variable name in the parse tree, and we
allocate an ordinary expression node in the eval tree.

Variables of type {\tt{}V{\char95}PDG} (pdg-code array) are not treated here.
They are handled by {\tt{}eval{\char95}node{\char95}compile{\char95}cvariable}.
\nwenddocs{}\nwbegincode{291}\moddef{Eval trees: procedures}\plusendmoddef
  recursive subroutine eval_node_compile_variable (en, pn, var_list, var_type)
    type(eval_node_t), pointer :: en
    type(parse_node_t), intent(in), target :: pn
    type(var_list_t), intent(in), target :: var_list
    integer, intent(in), optional :: var_type
    type(parse_node_t), pointer :: pn_name
    type(string_t) :: var_name
    logical, target, save :: no_lval
    real(default), target, save :: no_rval
    type(subevt_t), target, save :: no_pval
    type(string_t), target, save :: no_sval
    logical, target, save :: unknown = .false.
    integer :: type
    logical :: defined
    logical, pointer :: known
    logical, pointer :: lptr
    integer, pointer :: iptr
    real(default), pointer :: rptr
    complex(default), pointer :: cptr
    type(pdg_array_t), pointer :: aptr
    type(subevt_t), pointer :: pptr
    type(string_t), pointer :: sptr
    procedure(obs_unary_int), pointer :: obs1_iptr
    procedure(obs_unary_real), pointer :: obs1_rptr
    procedure(obs_binary_int), pointer :: obs2_iptr
    procedure(obs_binary_real), pointer :: obs2_rptr
    type(prt_t), pointer :: p1, p2
    if (debug) then
       print *, "read variable";  call parse_node_write (pn)
    end if
    if (present (var_type)) then
       select case (var_type)
       case (V_REAL, V_OBS1_REAL, V_OBS2_REAL, V_INT, V_OBS1_INT, &
                V_OBS2_INT, V_CMPLX)
          pn_name => pn
       case default
          pn_name => parse_node_get_sub_ptr (pn, 2)
       end select
    else
       pn_name => pn
    end if
    select case (char (parse_node_get_rule_key (pn_name)))
    case ("expr")
       call eval_node_compile_expr (en, pn_name, var_list)
    case ("lexpr")
       call eval_node_compile_lexpr (en, pn_name, var_list)
    case ("sexpr")
       call eval_node_compile_sexpr (en, pn_name, var_list)
    case ("pexpr")
       call eval_node_compile_pexpr (en, pn_name, var_list)
    case ("variable")
       var_name = parse_node_get_string (pn_name)
       if (present (var_type)) then
          select case (var_type)
          case (V_LOG);  var_name = "?" // var_name
          case (V_SEV);  var_name = "@" // var_name
          case (V_STR);  var_name = "$" // var_name   ! $ sign
          end select
       end if
       call var_list%get_var_properties &
            (var_name, req_type=var_type, type=type, is_defined=defined)
       allocate (en)
       if (defined) then
          select case (type)
          case (V_LOG)
             call var_list%get_lptr (var_name, lptr, known)
             call eval_node_init_log_ptr (en, var_name, lptr, known)
          case (V_INT)
             call var_list%get_iptr (var_name, iptr, known)
             call eval_node_init_int_ptr (en, var_name, iptr, known)
          case (V_REAL)
             call var_list%get_rptr (var_name, rptr, known)
             call eval_node_init_real_ptr (en, var_name, rptr, known)
          case (V_CMPLX)
             call var_list%get_cptr (var_name, cptr, known)
             call eval_node_init_cmplx_ptr (en, var_name, cptr, known)
          case (V_SEV)
             call var_list%get_pptr (var_name, pptr, known)
             call eval_node_init_subevt_ptr (en, var_name, pptr, known)
          case (V_STR)
             call var_list%get_sptr (var_name, sptr, known)
             call eval_node_init_string_ptr (en, var_name, sptr, known)
          case (V_OBS1_INT)
             call var_list%get_obs1_iptr (var_name, obs1_iptr, p1)
             call eval_node_init_obs1_int_ptr (en, var_name, obs1_iptr, p1)
          case (V_OBS2_INT)
             call var_list%get_obs2_iptr (var_name, obs2_iptr, p1, p2)
             call eval_node_init_obs2_int_ptr (en, var_name, obs2_iptr, p1, p2)
          case (V_OBS1_REAL)
             call var_list%get_obs1_rptr (var_name, obs1_rptr, p1)
             call eval_node_init_obs1_real_ptr (en, var_name, obs1_rptr, p1)
          case (V_OBS2_REAL)
             call var_list%get_obs2_rptr (var_name, obs2_rptr, p1, p2)
             call eval_node_init_obs2_real_ptr (en, var_name, obs2_rptr, p1, p2)
          case default
             call parse_node_write (pn)
             call msg_fatal ("Variable of this type " // &
                  "is not allowed in the present context")
             if (present (var_type)) then
                select case (var_type)
                case (V_LOG)
                   call eval_node_init_log_ptr (en, var_name, no_lval, unknown)
                case (V_SEV)
                   call eval_node_init_subevt_ptr &
                        (en, var_name, no_pval, unknown)
                case (V_STR)
                   call eval_node_init_string_ptr &
                        (en, var_name, no_sval, unknown)
                end select
             else
                call eval_node_init_real_ptr (en, var_name, no_rval, unknown)
             end if
          end select
       else
          call parse_node_write (pn)
          call msg_error ("This variable is undefined at this point")
          if (present (var_type)) then
             select case (var_type)
             case (V_LOG)
                call eval_node_init_log_ptr (en, var_name, no_lval, unknown)
             case (V_SEV)
                call eval_node_init_subevt_ptr &
                     (en, var_name, no_pval, unknown)
             case (V_STR)
                call eval_node_init_string_ptr (en, var_name, no_sval, unknown)
             end select
          else
             call eval_node_init_real_ptr (en, var_name, no_rval, unknown)
          end if
       end if
    end select
    if (debug) then
       call eval_node_write (en)
       print *, "done variable"
    end if
  end subroutine eval_node_compile_variable

\eatline
\nwendcode{}\nwbegindocs{292}In a given context, a variable has to have a certain type.
\nwenddocs{}\nwbegincode{293}\moddef{Eval trees: procedures}\plusendmoddef
  subroutine check_var_type (pn, ok, type_actual, type_requested)
    type(parse_node_t), intent(in) :: pn
    logical, intent(out) :: ok
    integer, intent(in) :: type_actual
    integer, intent(in), optional :: type_requested
    if (present (type_requested)) then
       select case (type_requested)
       case (V_LOG)
          select case (type_actual)
          case (V_LOG)
          case default
             call parse_node_write (pn)
             call msg_fatal ("Variable type is invalid (should be logical)")
             ok = .false.
          end select
       case (V_SEV)
          select case (type_actual)
          case (V_SEV)
          case default
             call parse_node_write (pn)
             call msg_fatal &
                  ("Variable type is invalid (should be particle set)")
             ok = .false.
          end select
       case (V_PDG)
          select case (type_actual)
          case (V_PDG)
          case default
             call parse_node_write (pn)
             call msg_fatal &
                  ("Variable type is invalid (should be PDG array)")
             ok = .false.
          end select
       case (V_STR)
          select case (type_actual)
          case (V_STR)
          case default
             call parse_node_write (pn)
             call msg_fatal &
                  ("Variable type is invalid (should be string)")
             ok = .false.
          end select
       case default
          call parse_node_write (pn)
          call msg_bug ("Variable type is unknown")
       end select
    else
       select case (type_actual)
       case (V_REAL, V_OBS1_REAL, V_OBS2_REAL, V_INT, V_OBS1_INT, &
                V_OBS2_INT, V_CMPLX)
       case default
          call parse_node_write (pn)
          call msg_fatal ("Variable type is invalid (should be numeric)")
          ok = .false.
       end select
    end if
    ok = .true.
  end subroutine check_var_type

\eatline
\nwendcode{}\nwbegindocs{294}Retrieve the result of an integration.  If the requested process has
been integrated, the results are available as special variables.  (The
variables cannot be accessed in the usual way since they contain
brackets in their names.)

Since this compilation step may occur before the processes have been
loaded, we have to initialize the required variables before they are
used.
\nwenddocs{}\nwbegincode{295}\moddef{Eval trees: procedures}\plusendmoddef
  subroutine eval_node_compile_result (en, pn, var_list)
    type(eval_node_t), pointer :: en
    type(parse_node_t), intent(in), target :: pn
    type(var_list_t), intent(in), target :: var_list
    type(parse_node_t), pointer :: pn_key, pn_prc_id
    type(string_t) :: key, prc_id, var_name
    integer, pointer :: iptr
    real(default), pointer :: rptr
    logical, pointer :: known
    if (debug) then
       print *, "read result";  call parse_node_write (pn)
    end if
    pn_key => parse_node_get_sub_ptr (pn)
    pn_prc_id => parse_node_get_next_ptr (pn_key)
    key = parse_node_get_key (pn_key)
    prc_id = parse_node_get_string (pn_prc_id)
    var_name = key // "(" // prc_id // ")"
    if (var_list%contains (var_name)) then
       allocate (en)
       select case (char(key))
       case ("num_id", "n_calls")
          call var_list%get_iptr (var_name, iptr, known)
          call eval_node_init_int_ptr (en, var_name, iptr, known)
       case ("integral", "error")
          call var_list%get_rptr (var_name, rptr, known)
          call eval_node_init_real_ptr (en, var_name, rptr, known)
       end select
    else
       call msg_fatal ("Result variable '" // char (var_name) &
            // "' is undefined (call 'integrate' before use)")
    end if
    if (debug) then
       call eval_node_write (en)
       print *, "done result"
    end if
  end subroutine eval_node_compile_result

\eatline
\nwendcode{}\nwbegindocs{296}This user observable behaves like a variable.  We link the node to
the generic user-observable entry in the variable list.  The syntax
element has an argument which provides the name of the user variable,
this is stored as an eval-node alongside with the variable.  When the
variable value is used, the user-supplied external function is called
and provides the (real) result value.
\nwenddocs{}\nwbegincode{297}\moddef{Eval trees: procedures}\plusendmoddef
  subroutine eval_node_compile_user_observable (en, pn, var_list)
    type(eval_node_t), pointer :: en
    type(parse_node_t), intent(in), target :: pn
    type(var_list_t), intent(in), target :: var_list
    type(parse_node_t), pointer :: pn_key, pn_arg, pn_obs
    type(eval_node_t), pointer :: en0
    integer :: res_type
    type(string_t) :: var_name
    integer :: type
    logical :: defined
    if (debug) then
       print *, "read user observable";  call parse_node_write (pn)
    end if
    pn_key => parse_node_get_sub_ptr (pn)
    select case (char (parse_node_get_key (pn_key)))
    case ("user_obs")
       res_type = V_REAL
    case default
       call parse_node_write (pn_key)
       call msg_bug ("user_observable: wrong keyword")
    end select
    pn_arg => parse_node_get_next_ptr (pn_key)
    pn_obs => parse_node_get_sub_ptr (pn_arg)
    call eval_node_compile_sexpr (en0, pn_obs, var_list)
    select case (res_type)
    case (V_INT);  var_name = "_User_obs_int"
    case (V_REAL); var_name = "_User_obs_real"
    end select
    call var_list%get_var_properties (var_name, type=type, is_defined=defined)
    allocate (en)
    if (defined) then
       select case (type)
       case (V_UOBS1_INT)
          call eval_node_init_uobs1_int (en, var_name, en0)
       case (V_UOBS2_INT)
          call eval_node_init_uobs2_int (en, var_name, en0)
       case (V_UOBS1_REAL)
          call eval_node_init_uobs1_real (en, var_name, en0)
       case (V_UOBS2_REAL)
          call eval_node_init_uobs2_real (en, var_name, en0)
       end select
    else
       call parse_node_write (pn)
       call msg_error ("This variable is undefined at this point")
    end if
    if (debug) then
       call eval_node_write (en)
       print *, "done user observable"
    end if
  end subroutine eval_node_compile_user_observable

\eatline
\nwendcode{}\nwbegindocs{298}Functions with a single argument.  For non-constant arguments, watch
for functions which convert their argument to a different type.
\nwenddocs{}\nwbegincode{299}\moddef{Eval trees: procedures}\plusendmoddef
  recursive subroutine eval_node_compile_unary_function (en, pn, var_list)
    type(eval_node_t), pointer :: en
    type(parse_node_t), intent(in) :: pn
    type(var_list_t), intent(in), target :: var_list
    type(parse_node_t), pointer :: pn_fname, pn_arg
    type(eval_node_t), pointer :: en1
    type(string_t) :: key
    integer :: t
    if (debug) then
       print *, "read unary function";  call parse_node_write (pn)
    end if
    pn_fname => parse_node_get_sub_ptr (pn)
    pn_arg => parse_node_get_next_ptr (pn_fname, tag="function_arg1")
    call eval_node_compile_expr &
         (en1, parse_node_get_sub_ptr (pn_arg, tag="expr"), var_list)
    t = en1%result_type
    allocate (en)
    key = parse_node_get_key (pn_fname)
    if (en1%type == EN_CONSTANT) then
       select case (char (key))
       case ("complex")
          select case (t)
          case (V_INT);  call eval_node_init_cmplx (en, cmplx_i (en1))
          case (V_REAL); call eval_node_init_cmplx (en, cmplx_r (en1))
          case (V_CMPLX); deallocate (en);  en => en1;  en1 => null ()
          case default;  call eval_type_error (pn, char (key), t)          
          end select
       case ("real")
          select case (t)
          case (V_INT);  call eval_node_init_real (en, real_i (en1))
          case (V_REAL); deallocate (en);  en => en1;  en1 => null ()
          case (V_CMPLX); call eval_node_init_real (en, real_c (en1))
          case default;  call eval_type_error (pn, char (key), t)          
          end select
       case ("int")
          select case (t)
          case (V_INT);  deallocate (en);  en => en1;  en1 => null ()
          case (V_REAL); call eval_node_init_int  (en, int_r (en1))
          case (V_CMPLX); call eval_node_init_int  (en, int_c (en1))
          end select
       case ("nint")
          select case (t)
          case (V_INT);  deallocate (en);  en => en1;  en1 => null ()
          case (V_REAL); call eval_node_init_int  (en, nint_r (en1))          
          case default;  call eval_type_error (pn, char (key), t)
          end select
       case ("floor")
          select case (t)
          case (V_INT);  deallocate (en);  en => en1;  en1 => null ()
          case (V_REAL); call eval_node_init_int  (en, floor_r (en1))
          case default;  call eval_type_error (pn, char (key), t)          
          end select
       case ("ceiling")
          select case (t)
          case (V_INT);  deallocate (en);  en => en1;  en1 => null ()
          case (V_REAL); call eval_node_init_int  (en, ceiling_r (en1))
          case default;  call eval_type_error (pn, char (key), t)
          end select
       case ("abs")
          select case (t)
          case (V_INT);  call eval_node_init_int  (en, abs_i (en1))
          case (V_REAL); call eval_node_init_real (en, abs_r (en1))
          case (V_CMPLX); call eval_node_init_real (en, abs_c (en1))          
          end select
       case ("sgn")
          select case (t)
          case (V_INT);  call eval_node_init_int  (en, sgn_i (en1))
          case (V_REAL); call eval_node_init_real (en, sgn_r (en1))
          case default;  call eval_type_error (pn, char (key), t)
          end select
       case ("sqrt")
          select case (t)
          case (V_REAL); call eval_node_init_real (en, sqrt_r (en1))
          case (V_CMPLX); call eval_node_init_cmplx (en, sqrt_c (en1))          
          case default;  call eval_type_error (pn, char (key), t)
          end select
       case ("exp")
          select case (t)
          case (V_REAL); call eval_node_init_real (en, exp_r (en1))
          case (V_CMPLX); call eval_node_init_cmplx (en, exp_c (en1))
          case default;  call eval_type_error (pn, char (key), t)
          end select
       case ("log")
          select case (t)
          case (V_REAL); call eval_node_init_real (en, log_r (en1))
          case (V_CMPLX); call eval_node_init_cmplx (en, log_c (en1))          
          case default;  call eval_type_error (pn, char (key), t)
          end select
       case ("log10")
          select case (t)
          case (V_REAL); call eval_node_init_real (en, log10_r (en1))
          case default;  call eval_type_error (pn, char (key), t)
          end select
       case ("sin")
          select case (t)
          case (V_REAL); call eval_node_init_real (en, sin_r (en1))
          case (V_CMPLX); call eval_node_init_cmplx (en, sin_c (en1))          
          case default;  call eval_type_error (pn, char (key), t)
          end select
       case ("cos")
          select case (t)
          case (V_REAL); call eval_node_init_real (en, cos_r (en1))
          case (V_CMPLX); call eval_node_init_cmplx (en, cos_c (en1))          
          case default;  call eval_type_error (pn, char (key), t)
          end select
       case ("tan")
          select case (t)
          case (V_REAL); call eval_node_init_real (en, tan_r (en1))
          case default;  call eval_type_error (pn, char (key), t)
          end select
       case ("asin")
          select case (t)
          case (V_REAL); call eval_node_init_real (en, asin_r (en1))
          case default;  call eval_type_error (pn, char (key), t)
          end select
       case ("acos")
          select case (t)
          case (V_REAL); call eval_node_init_real (en, acos_r (en1))
          case default;  call eval_type_error (pn, char (key), t)
          end select
       case ("atan")
          select case (t)
          case (V_REAL); call eval_node_init_real (en, atan_r (en1))
          case default;  call eval_type_error (pn, char (key), t)
          end select
       case ("sinh")
          select case (t)
          case (V_REAL); call eval_node_init_real (en, sinh_r (en1))
          case default;  call eval_type_error (pn, char (key), t)
          end select
       case ("cosh")
          select case (t)
          case (V_REAL); call eval_node_init_real (en, cosh_r (en1))
          case default;  call eval_type_error (pn, char (key), t)
          end select
       case ("tanh")
          select case (t)
          case (V_REAL); call eval_node_init_real (en, tanh_r (en1))
          case default;  call eval_type_error (pn, char (key), t)
          end select
       case default
          call parse_node_mismatch ("function name", pn_fname)
       end select
       if (associated (en1)) then
          call eval_node_final_rec (en1)
          deallocate (en1)
       end if
    else
       select case (char (key))
       case ("complex")
          call eval_node_init_branch (en, key, V_CMPLX, en1)
       case ("real")
          call eval_node_init_branch (en, key, V_REAL, en1)
       case ("int", "nint", "floor", "ceiling")
          call eval_node_init_branch (en, key, V_INT, en1)
       case default
          call eval_node_init_branch (en, key, t, en1)
       end select
       select case (char (key))
       case ("complex")
          select case (t)
          case (V_INT);  call eval_node_set_op1_cmplx (en, cmplx_i)
          case (V_REAL); call eval_node_set_op1_cmplx (en, cmplx_r)
          case (V_CMPLX); deallocate (en);  en => en1
          case default;  call eval_type_error (pn, char (key), t)
          end select
       case ("real")
          select case (t)
          case (V_INT);  call eval_node_set_op1_real (en, real_i)
          case (V_REAL); deallocate (en);  en => en1
          case (V_CMPLX); call eval_node_set_op1_real (en, real_c)
          case default;  call eval_type_error (pn, char (key), t)
          end select
       case ("int")
          select case (t)
          case (V_INT);  deallocate (en);  en => en1
          case (V_REAL); call eval_node_set_op1_int (en, int_r)
          case (V_CMPLX); call eval_node_set_op1_int (en, int_c)          
          end select
       case ("nint")
          select case (t)
          case (V_INT);  deallocate (en);  en => en1
          case (V_REAL); call eval_node_set_op1_int (en, nint_r)
          case default;  call eval_type_error (pn, char (key), t)
          end select
       case ("floor")
          select case (t)
          case (V_INT);  deallocate (en);  en => en1
          case (V_REAL); call eval_node_set_op1_int (en, floor_r)
          case default;  call eval_type_error (pn, char (key), t)
          end select
       case ("ceiling")
          select case (t)
          case (V_INT);  deallocate (en);  en => en1
          case (V_REAL); call eval_node_set_op1_int (en, ceiling_r)
          case default;  call eval_type_error (pn, char (key), t)
          end select
       case ("abs")
          select case (t)
          case (V_INT);  call eval_node_set_op1_int  (en, abs_i)
          case (V_REAL); call eval_node_set_op1_real (en, abs_r)
          case (V_CMPLX); 
             call eval_node_init_branch (en, key, V_REAL, en1)
             call eval_node_set_op1_real (en, abs_c)
          end select
       case ("sgn")
          select case (t)
          case (V_INT);  call eval_node_set_op1_int  (en, sgn_i)
          case (V_REAL); call eval_node_set_op1_real (en, sgn_r)
          case default;  call eval_type_error (pn, char (key), t)
          end select
       case ("sqrt")
          select case (t)
          case (V_REAL); call eval_node_set_op1_real (en, sqrt_r)
          case (V_CMPLX); call eval_node_set_op1_cmplx (en, sqrt_c)          
          case default;  call eval_type_error (pn, char (key), t)
          end select
       case ("exp")
          select case (t)
          case (V_REAL); call eval_node_set_op1_real (en, exp_r)
          case (V_CMPLX); call eval_node_set_op1_cmplx (en, exp_c)          
          case default;  call eval_type_error (pn, char (key), t)
          end select
       case ("log")
          select case (t)
          case (V_REAL); call eval_node_set_op1_real (en, log_r)
          case (V_CMPLX); call eval_node_set_op1_cmplx (en, log_c)          
          case default;  call eval_type_error (pn, char (key), t)
          end select
       case ("log10")
          select case (t)
          case (V_REAL); call eval_node_set_op1_real (en, log10_r)
          case default;  call eval_type_error (pn, char (key), t)
          end select
       case ("sin")
          select case (t)
          case (V_REAL); call eval_node_set_op1_real (en, sin_r)
          case (V_CMPLX); call eval_node_set_op1_cmplx (en, sin_c)          
          case default;  call eval_type_error (pn, char (key), t)
          end select
       case ("cos")
          select case (t)
          case (V_REAL); call eval_node_set_op1_real (en, cos_r)
          case (V_CMPLX); call eval_node_set_op1_cmplx (en, cos_c)          
          case default;  call eval_type_error (pn, char (key), t)
          end select
       case ("tan")
          select case (t)
          case (V_REAL); call eval_node_set_op1_real (en, tan_r)
          case default;  call eval_type_error (pn, char (key), t)
          end select
       case ("asin")
          select case (t)
          case (V_REAL); call eval_node_set_op1_real (en, asin_r)
          case default;  call eval_type_error (pn, char (key), t)
          end select
       case ("acos")
          select case (t)
          case (V_REAL); call eval_node_set_op1_real (en, acos_r)
          case default;  call eval_type_error (pn, char (key), t)
          end select
       case ("atan")
          select case (t)
          case (V_REAL); call eval_node_set_op1_real (en, atan_r)
          case default;  call eval_type_error (pn, char (key), t)
          end select
       case ("sinh")
          select case (t)
          case (V_REAL); call eval_node_set_op1_real (en, sinh_r)
          case default;  call eval_type_error (pn, char (key), t)
          end select
       case ("cosh")
          select case (t)
          case (V_REAL); call eval_node_set_op1_real (en, cosh_r)
          case default;  call eval_type_error (pn, char (key), t)
          end select
       case ("tanh")
          select case (t)
          case (V_REAL); call eval_node_set_op1_real (en, tanh_r)
          case default;  call eval_type_error (pn, char (key), t)
          end select
       case default
          call parse_node_mismatch ("function name", pn_fname)
       end select
    end if
    if (debug) then
       call eval_node_write (en)
       print *, "done function"
    end if
  end subroutine eval_node_compile_unary_function

\eatline
\nwendcode{}\nwbegindocs{300}Functions with two arguments.
\nwenddocs{}\nwbegincode{301}\moddef{Eval trees: procedures}\plusendmoddef
  recursive subroutine eval_node_compile_binary_function (en, pn, var_list)
    type(eval_node_t), pointer :: en
    type(parse_node_t), intent(in) :: pn
    type(var_list_t), intent(in), target :: var_list
    type(parse_node_t), pointer :: pn_fname, pn_arg, pn_arg1, pn_arg2
    type(eval_node_t), pointer :: en1, en2
    type(string_t) :: key
    integer :: t1, t2
    if (debug) then
       print *, "read binary function";  call parse_node_write (pn)
    end if
    pn_fname => parse_node_get_sub_ptr (pn)
    pn_arg => parse_node_get_next_ptr (pn_fname, tag="function_arg2")
    pn_arg1 => parse_node_get_sub_ptr (pn_arg, tag="expr")
    pn_arg2 => parse_node_get_next_ptr (pn_arg1, tag="expr")
    call eval_node_compile_expr (en1, pn_arg1, var_list)
    call eval_node_compile_expr (en2, pn_arg2, var_list)
    t1 = en1%result_type
    t2 = en2%result_type
    allocate (en)
    key = parse_node_get_key (pn_fname)
    if (en1%type == EN_CONSTANT .and. en2%type == EN_CONSTANT) then
       select case (char (key))
       case ("max")
          select case (t1)
          case (V_INT)
             select case (t2)
             case (V_INT);  call eval_node_init_int  (en, max_ii (en1, en2))
             case (V_REAL); call eval_node_init_real (en, max_ir (en1, en2))
             case default;  call eval_type_error (pn, char (key), t2)      
             end select
          case (V_REAL)
             select case (t2)
             case (V_INT);  call eval_node_init_real (en, max_ri (en1, en2))
             case (V_REAL); call eval_node_init_real (en, max_rr (en1, en2))
             case default;  call eval_type_error (pn, char (key), t2) 
             end select
           case default;  call eval_type_error (pn, char (key), t1)             
         end select
       case ("min")
          select case (t1)
          case (V_INT)
             select case (t2)
             case (V_INT);  call eval_node_init_int  (en, min_ii (en1, en2))
             case (V_REAL); call eval_node_init_real (en, min_ir (en1, en2))
             case default;  call eval_type_error (pn, char (key), t2)             
             end select
          case (V_REAL)
             select case (t2)
             case (V_INT);  call eval_node_init_real (en, min_ri (en1, en2))
             case (V_REAL); call eval_node_init_real (en, min_rr (en1, en2))
             case default;  call eval_type_error (pn, char (key), t2)                          
             end select
           case default;  call eval_type_error (pn, char (key), t1)             
         end select
       case ("mod")
          select case (t1)
          case (V_INT)
             select case (t2)
             case (V_INT);  call eval_node_init_int  (en, mod_ii (en1, en2))
             case (V_REAL); call eval_node_init_real (en, mod_ir (en1, en2))
             case default;  call eval_type_error (pn, char (key), t2)             
             end select
          case (V_REAL)
             select case (t2)
             case (V_INT);  call eval_node_init_real (en, mod_ri (en1, en2))
             case (V_REAL); call eval_node_init_real (en, mod_rr (en1, en2))
             case default;  call eval_type_error (pn, char (key), t2)                          
             end select
           case default;  call eval_type_error (pn, char (key), t1)             
          end select
       case ("modulo")
          select case (t1)
          case (V_INT)
             select case (t2)
             case (V_INT);  call eval_node_init_int  (en, modulo_ii (en1, en2))
             case (V_REAL); call eval_node_init_real (en, modulo_ir (en1, en2))
             case default;  call eval_type_error (pn, char (key), t2)             
             end select
          case (V_REAL)
             select case (t2)
             case (V_INT);  call eval_node_init_real (en, modulo_ri (en1, en2))
             case (V_REAL); call eval_node_init_real (en, modulo_rr (en1, en2))
             case default;  call eval_type_error (pn, char (key), t2)                          
             end select
           case default;  call eval_type_error (pn, char (key), t2)             
         end select
       case default
          call parse_node_mismatch ("function name", pn_fname)
       end select
       call eval_node_final_rec (en1)
       deallocate (en1)
    else
       call eval_node_init_branch (en, key, t1, en1, en2)
       select case (char (key))
       case ("max")
          select case (t1)
          case (V_INT)
             select case (t2)
             case (V_INT);  call eval_node_set_op2_int  (en, max_ii)
             case (V_REAL); call eval_node_set_op2_real (en, max_ir)
             case default;  call eval_type_error (pn, char (key), t2)             
             end select
          case (V_REAL)
             select case (t2)
             case (V_INT);  call eval_node_set_op2_real (en, max_ri)
             case (V_REAL); call eval_node_set_op2_real (en, max_rr)
             case default;  call eval_type_error (pn, char (key), t2)                          
             end select
           case default;  call eval_type_error (pn, char (key), t2)               
         end select
       case ("min")
          select case (t1)
          case (V_INT)
             select case (t2)
             case (V_INT);  call eval_node_set_op2_int  (en, min_ii)
             case (V_REAL); call eval_node_set_op2_real (en, min_ir)
             case default;  call eval_type_error (pn, char (key), t2)             
             end select
          case (V_REAL)
             select case (t2)
             case (V_INT);  call eval_node_set_op2_real (en, min_ri)
             case (V_REAL); call eval_node_set_op2_real (en, min_rr)
             case default;  call eval_type_error (pn, char (key), t2)                          
             end select
           case default;  call eval_type_error (pn, char (key), t2)            
         end select
       case ("mod")
          select case (t1)
          case (V_INT)
             select case (t2)
             case (V_INT);  call eval_node_set_op2_int  (en, mod_ii)
             case (V_REAL); call eval_node_set_op2_real (en, mod_ir)
             case default;  call eval_type_error (pn, char (key), t2)                          
             end select
          case (V_REAL)
             select case (t2)
             case (V_INT);  call eval_node_set_op2_real (en, mod_ri)
             case (V_REAL); call eval_node_set_op2_real (en, mod_rr)
             case default;  call eval_type_error (pn, char (key), t2)                          
             end select
           case default;  call eval_type_error (pn, char (key), t2)                          
        end select
       case ("modulo")
          select case (t1)
          case (V_INT)
             select case (t2)
             case (V_INT);  call eval_node_set_op2_int  (en, modulo_ii)
             case (V_REAL); call eval_node_set_op2_real (en, modulo_ir)
             case default;  call eval_type_error (pn, char (key), t2)                          
             end select
          case (V_REAL)
             select case (t2)
             case (V_INT);  call eval_node_set_op2_real (en, modulo_ri)
             case (V_REAL); call eval_node_set_op2_real (en, modulo_rr)
             case default;  call eval_type_error (pn, char (key), t2)                          
             end select
           case default;  call eval_type_error (pn, char (key), t2)                          
         end select
       case default
          call parse_node_mismatch ("function name", pn_fname)
       end select
    end if
    if (debug) then
       call eval_node_write (en)
       print *, "done function"
    end if
  end subroutine eval_node_compile_binary_function

\eatline
\nwendcode{}\nwbegindocs{302}\nwdocspar
\subsubsection{Variable definition}
A block expression contains a variable definition (first argument) and
an expression where the definition can be used (second argument).  The
{\tt{}result{\char95}type} decides which type of expression is expected for the
second argument.  For numeric variables, if there is a mismatch
between real and integer type, insert an extra node for type
conversion.
\nwenddocs{}\nwbegincode{303}\moddef{Eval trees: procedures}\plusendmoddef
  recursive subroutine eval_node_compile_block_expr &
       (en, pn, var_list, result_type)
    type(eval_node_t), pointer :: en
    type(parse_node_t), intent(in) :: pn
    type(var_list_t), intent(in), target :: var_list
    integer, intent(in), optional :: result_type
    type(parse_node_t), pointer :: pn_var_spec, pn_var_subspec
    type(parse_node_t), pointer :: pn_var_type, pn_var_name, pn_var_expr
    type(parse_node_t), pointer :: pn_expr
    type(string_t) :: var_name
    type(eval_node_t), pointer :: en1, en2
    integer :: var_type
    logical :: new
    if (debug) then
       print *, "read block expr";  call parse_node_write (pn)
    end if
    new = .false.
    pn_var_spec => parse_node_get_sub_ptr (pn, 2)
    select case (char (parse_node_get_rule_key (pn_var_spec)))
    case ("var_num");      var_type = V_NONE
       pn_var_name => parse_node_get_sub_ptr (pn_var_spec)
    case ("var_int");      var_type = V_INT
       new = .true.
       pn_var_name => parse_node_get_sub_ptr (pn_var_spec, 2)
    case ("var_real");     var_type = V_REAL
       new = .true.
       pn_var_name => parse_node_get_sub_ptr (pn_var_spec, 2)
    case ("var_cmplx");     var_type = V_CMPLX
       new = .true.
       pn_var_name => parse_node_get_sub_ptr (pn_var_spec, 2)
    case ("var_logical_new");  var_type = V_LOG
       new = .true.
       pn_var_subspec => parse_node_get_sub_ptr (pn_var_spec, 2)
       pn_var_name => parse_node_get_sub_ptr (pn_var_subspec, 2)
    case ("var_logical_spec");  var_type = V_LOG
       pn_var_name => parse_node_get_sub_ptr (pn_var_spec, 2)
    case ("var_plist_new");    var_type = V_SEV
       new = .true.
       pn_var_subspec => parse_node_get_sub_ptr (pn_var_spec, 2)
       pn_var_name => parse_node_get_sub_ptr (pn_var_subspec, 2)
    case ("var_plist_spec");    var_type = V_SEV
       new = .true.
       pn_var_name => parse_node_get_sub_ptr (pn_var_spec, 2)
    case ("var_alias");    var_type = V_PDG
       new = .true.
       pn_var_name => parse_node_get_sub_ptr (pn_var_spec, 2)
    case ("var_string_new");   var_type = V_STR
       new = .true.
       pn_var_subspec => parse_node_get_sub_ptr (pn_var_spec, 2)
       pn_var_name => parse_node_get_sub_ptr (pn_var_subspec, 2)
    case ("var_string_spec");   var_type = V_STR
       pn_var_name => parse_node_get_sub_ptr (pn_var_spec, 2)
    case default
       call parse_node_mismatch &
            ("logical|int|real|plist|alias", pn_var_type)
    end select
    pn_var_expr => parse_node_get_next_ptr (pn_var_name, 2)
    pn_expr => parse_node_get_next_ptr (pn_var_spec, 2)
    var_name = parse_node_get_string (pn_var_name)
    select case (var_type)
    case (V_LOG);  var_name = "?" // var_name
    case (V_SEV);  var_name = "@" // var_name
    case (V_STR);  var_name = "$" // var_name    ! $ sign
    end select
    call var_list_check_user_var (var_list, var_name, var_type, new)
    call eval_node_compile_genexpr (en1, pn_var_expr, var_list, var_type)
    call insert_conversion_node (en1, var_type)
    allocate (en)
    call eval_node_init_block (en, var_name, var_type, en1, var_list)
    call eval_node_compile_genexpr (en2, pn_expr, en%var_list, result_type)
    call eval_node_set_expr (en, en2)
    if (debug) then
       call eval_node_write (en)
       print *, "done block expr"
    end if
  end subroutine eval_node_compile_block_expr
    
\eatline
\nwendcode{}\nwbegindocs{304}Insert a conversion node for integer/real/complex transformation if necessary.
What shall we do for the complex to integer/real conversion?
\nwenddocs{}\nwbegincode{305}\moddef{Eval trees: procedures}\plusendmoddef
  subroutine insert_conversion_node (en, result_type)
    type(eval_node_t), pointer :: en
    integer, intent(in) :: result_type
    type(eval_node_t), pointer :: en_conv
    select case (en%result_type)
    case (V_INT)
       select case (result_type)
       case (V_REAL)
          allocate (en_conv)
          call eval_node_init_branch (en_conv, var_str ("real"), V_REAL, en)
          call eval_node_set_op1_real (en_conv, real_i)
          en => en_conv
       case (V_CMPLX)
          allocate (en_conv)
          call eval_node_init_branch (en_conv, var_str ("complex"), V_CMPLX, en)
          call eval_node_set_op1_cmplx (en_conv, cmplx_i)
          en => en_conv          
       end select
    case (V_REAL)
       select case (result_type)
       case (V_INT)
          allocate (en_conv)
          call eval_node_init_branch (en_conv, var_str ("int"), V_INT, en)
          call eval_node_set_op1_int (en_conv, int_r)
          en => en_conv
       case (V_CMPLX)
          allocate (en_conv)
          call eval_node_init_branch (en_conv, var_str ("complex"), V_CMPLX, en)
          call eval_node_set_op1_cmplx (en_conv, cmplx_r)
          en => en_conv
       end select          
    case (V_CMPLX)
       select case (result_type)
       case (V_INT)
          allocate (en_conv)
          call eval_node_init_branch (en_conv, var_str ("int"), V_INT, en)
          call eval_node_set_op1_int (en_conv, int_c)
          en => en_conv
       case (V_REAL)
          allocate (en_conv)
          call eval_node_init_branch (en_conv, var_str ("real"), V_REAL, en)
          call eval_node_set_op1_real (en_conv, real_c)
          en => en_conv
       end select          
     case default     
     end select    
  end subroutine insert_conversion_node

\eatline
\nwendcode{}\nwbegindocs{306}\nwdocspar
\subsubsection{Conditionals}
A conditional has the structure if lexpr then expr else expr.  So we
first evaluate the logical expression, then depending on the result
the first or second expression.  Note that the second expression is
mandatory.

The {\tt{}result{\char95}type}, if present, defines the requested type of the
{\tt{}then} and {\tt{}else} clauses.  Default is numeric (int/real).  If
there is a mismatch between real and integer result types, insert
conversion nodes.
\nwenddocs{}\nwbegincode{307}\moddef{Eval trees: procedures}\plusendmoddef
  recursive subroutine eval_node_compile_conditional &
       (en, pn, var_list, result_type)
    type(eval_node_t), pointer :: en
    type(parse_node_t), intent(in) :: pn
    type(var_list_t), intent(in), target :: var_list
    integer, intent(in), optional :: result_type
    type(parse_node_t), pointer :: pn_condition, pn_expr
    type(parse_node_t), pointer :: pn_maybe_elsif, pn_elsif_branch
    type(parse_node_t), pointer :: pn_maybe_else, pn_else_branch, pn_else_expr
    type(eval_node_t), pointer :: en0, en1, en2
    integer :: restype
    if (debug) then
       print *, "read conditional";  call parse_node_write (pn)
    end if
    pn_condition => parse_node_get_sub_ptr (pn, 2, tag="lexpr")
    pn_expr => parse_node_get_next_ptr (pn_condition, 2)
    call eval_node_compile_lexpr (en0, pn_condition, var_list)
    call eval_node_compile_genexpr (en1, pn_expr, var_list, result_type)
    if (present (result_type)) then
       restype = major_result_type (result_type, en1%result_type)
    else
       restype = en1%result_type
    end if
    pn_maybe_elsif => parse_node_get_next_ptr (pn_expr)
    select case (char (parse_node_get_rule_key (pn_maybe_elsif)))
    case ("maybe_elsif_expr", &
          "maybe_elsif_lexpr", &
          "maybe_elsif_pexpr", &
          "maybe_elsif_cexpr", &
          "maybe_elsif_sexpr")
       pn_elsif_branch => parse_node_get_sub_ptr (pn_maybe_elsif)
       pn_maybe_else => parse_node_get_next_ptr (pn_maybe_elsif)
       select case (char (parse_node_get_rule_key (pn_maybe_else)))
       case ("maybe_else_expr", &
          "maybe_else_lexpr", &
          "maybe_else_pexpr", &
          "maybe_else_cexpr", &
          "maybe_else_sexpr")
          pn_else_branch => parse_node_get_sub_ptr (pn_maybe_else)
          pn_else_expr => parse_node_get_sub_ptr (pn_else_branch, 2)
       case default
          pn_else_expr => null ()
       end select
       call eval_node_compile_elsif &
            (en2, pn_elsif_branch, pn_else_expr, var_list, restype)
    case ("maybe_else_expr", &
          "maybe_else_lexpr", &
          "maybe_else_pexpr", &
          "maybe_else_cexpr", &
          "maybe_else_sexpr")
       pn_maybe_else => pn_maybe_elsif
       pn_maybe_elsif => null ()
       pn_else_branch => parse_node_get_sub_ptr (pn_maybe_else)
       pn_else_expr => parse_node_get_sub_ptr (pn_else_branch, 2)
       call eval_node_compile_genexpr &
            (en2, pn_else_expr, var_list, restype)
    case ("endif")
       call eval_node_compile_default_else (en2, restype)
    case default
       call msg_bug ("Broken conditional: unexpected " &
            // char (parse_node_get_rule_key (pn_maybe_elsif)))
    end select
    call eval_node_create_conditional (en, en0, en1, en2, restype)
    call conditional_insert_conversion_nodes (en, restype)
    if (debug) then
       call eval_node_write (en)
       print *, "done conditional"
    end if
  end subroutine eval_node_compile_conditional

\eatline
\nwendcode{}\nwbegindocs{308}This recursively generates 'elsif' conditionals as a chain of sub-nodes of
the main conditional.
\nwenddocs{}\nwbegincode{309}\moddef{Eval trees: procedures}\plusendmoddef
  recursive subroutine eval_node_compile_elsif &
       (en, pn, pn_else_expr, var_list, result_type)
    type(eval_node_t), pointer :: en
    type(parse_node_t), intent(in), target :: pn
    type(parse_node_t), pointer :: pn_else_expr
    type(var_list_t), intent(in), target :: var_list
    integer, intent(inout) :: result_type
    type(parse_node_t), pointer :: pn_next, pn_condition, pn_expr
    type(eval_node_t), pointer :: en0, en1, en2
    pn_condition => parse_node_get_sub_ptr (pn, 2, tag="lexpr")
    pn_expr => parse_node_get_next_ptr (pn_condition, 2)
    call eval_node_compile_lexpr (en0, pn_condition, var_list)
    call eval_node_compile_genexpr (en1, pn_expr, var_list, result_type)
    result_type = major_result_type (result_type, en1%result_type)
    pn_next => parse_node_get_next_ptr (pn)
    if (associated (pn_next)) then
       call eval_node_compile_elsif &
            (en2, pn_next, pn_else_expr, var_list, result_type)
       result_type = major_result_type (result_type, en2%result_type)
    else if (associated (pn_else_expr)) then
       call eval_node_compile_genexpr &
            (en2, pn_else_expr, var_list, result_type)
       result_type = major_result_type (result_type, en2%result_type)
    else
       call eval_node_compile_default_else (en2, result_type)
    end if
    call eval_node_create_conditional (en, en0, en1, en2, result_type)
  end subroutine eval_node_compile_elsif       

\eatline
\nwendcode{}\nwbegindocs{310}This makes a default 'else' branch in case it was omitted.  The default value
just depends on the expected type.
\nwenddocs{}\nwbegincode{311}\moddef{Eval trees: procedures}\plusendmoddef
  subroutine eval_node_compile_default_else (en, result_type)
    type(eval_node_t), pointer :: en
    integer, intent(in) :: result_type
    type(subevt_t) :: pval_empty
    type(pdg_array_t) :: aval_undefined
    allocate (en)
    select case (result_type)
    case (V_LOG);  call eval_node_init_log (en, .false.)
    case (V_INT);  call eval_node_init_int (en, 0)
    case (V_REAL);  call eval_node_init_real (en, 0._default)
    case (V_CMPLX)
         call eval_node_init_cmplx (en, (0._default, 0._default))
    case (V_SEV)
       call subevt_init (pval_empty)
       call eval_node_init_subevt (en, pval_empty)
    case (V_PDG)
       call eval_node_init_pdg_array  (en, aval_undefined)
    case (V_STR)
       call eval_node_init_string (en, var_str (""))
    case default
       call msg_bug ("Undefined type for 'else' branch in conditional")
    end select
  end subroutine eval_node_compile_default_else

\eatline
\nwendcode{}\nwbegindocs{312}If the logical expression is constant, we can simplify the conditional node
by replacing it with the selected branch.  Otherwise, we initialize a true
branching.
\nwenddocs{}\nwbegincode{313}\moddef{Eval trees: procedures}\plusendmoddef
  subroutine eval_node_create_conditional (en, en0, en1, en2, result_type)
    type(eval_node_t), pointer :: en, en0, en1, en2
    integer, intent(in) :: result_type
    if (en0%type == EN_CONSTANT) then
       if (en0%lval) then
          en => en1
          call eval_node_final_rec (en2)
          deallocate (en2)
       else
          en => en2
          call eval_node_final_rec (en1)
          deallocate (en1)
       end if
    else
       allocate (en)
       call eval_node_init_conditional (en, result_type, en0, en1, en2)
    end if
  end subroutine eval_node_create_conditional

\eatline
\nwendcode{}\nwbegindocs{314}Return the numerical result type which should be used for the combination of
the two result types.
\nwenddocs{}\nwbegincode{315}\moddef{Eval trees: procedures}\plusendmoddef
  function major_result_type (t1, t2) result (t)
    integer :: t
    integer, intent(in) :: t1, t2
    select case (t1)
    case (V_INT)
       select case (t2)
       case (V_INT, V_REAL, V_CMPLX)
          t = t2
       case default
          call type_mismatch ()
       end select
    case (V_REAL)
       select case (t2)
       case (V_INT)
          t = t1
       case (V_REAL, V_CMPLX)
          t = t2
       case default
          call type_mismatch ()
       end select
    case (V_CMPLX)
       select case (t2)
       case (V_INT, V_REAL, V_CMPLX)
          t = t1
       case default
          call type_mismatch ()
       end select
    case default
       if (t1 == t2) then
          t = t1
       else
          call type_mismatch ()
       end if
    end select
  contains
    subroutine type_mismatch ()
      call msg_bug ("Type mismatch in branches of a conditional expression")
    end subroutine type_mismatch
  end function major_result_type

\eatline
\nwendcode{}\nwbegindocs{316}Recursively insert conversion nodes where necessary.
\nwenddocs{}\nwbegincode{317}\moddef{Eval trees: procedures}\plusendmoddef
  recursive subroutine conditional_insert_conversion_nodes (en, result_type)
    type(eval_node_t), intent(inout), target :: en
    integer, intent(in) :: result_type
    select case (result_type)
    case (V_INT, V_REAL, V_CMPLX)
       call insert_conversion_node (en%arg1, result_type)
       if (en%arg2%type == EN_CONDITIONAL) then
          call conditional_insert_conversion_nodes (en%arg2, result_type)
       else
          call insert_conversion_node (en%arg2, result_type)
       end if
    end select
  end subroutine conditional_insert_conversion_nodes

\eatline
\nwendcode{}\nwbegindocs{318}\nwdocspar
\subsubsection{Logical expressions}
A logical expression consists of one or more singlet logical expressions
concatenated by {\tt{};}.  This is for allowing side-effects, only the last value
is used.
\nwenddocs{}\nwbegincode{319}\moddef{Eval trees: procedures}\plusendmoddef
  recursive subroutine eval_node_compile_lexpr (en, pn, var_list)
    type(eval_node_t), pointer :: en
    type(parse_node_t), intent(in) :: pn
    type(var_list_t), intent(in), target :: var_list
    type(parse_node_t), pointer :: pn_term, pn_sequel, pn_arg
    type(eval_node_t), pointer :: en1, en2
    if (debug) then
       print *, "read lexpr";  call parse_node_write (pn)
    end if
    pn_term => parse_node_get_sub_ptr (pn, tag="lsinglet")
    call eval_node_compile_lsinglet (en, pn_term, var_list)
    pn_sequel => parse_node_get_next_ptr (pn_term, tag="lsequel")
    do while (associated (pn_sequel))
       pn_arg => parse_node_get_sub_ptr (pn_sequel, 2, tag="lsinglet")
       en1 => en
       call eval_node_compile_lsinglet (en2, pn_arg, var_list)
       allocate (en)
       if (en1%type == EN_CONSTANT .and. en2%type == EN_CONSTANT) then
          call eval_node_init_log (en, ignore_first_ll (en1, en2))
          call eval_node_final_rec (en1)
          call eval_node_final_rec (en2)
          deallocate (en1, en2)
       else   
          call eval_node_init_branch &
               (en, var_str ("lsequel"), V_LOG, en1, en2)
          call eval_node_set_op2_log (en, ignore_first_ll)
       end if
       pn_sequel => parse_node_get_next_ptr (pn_sequel)
    end do
    if (debug) then
       call eval_node_write (en)
       print *, "done lexpr"
    end if
  end subroutine eval_node_compile_lexpr

\eatline
\nwendcode{}\nwbegindocs{320}A logical singlet expression consists of one or more logical terms
concatenated by {\tt{}or}.
\nwenddocs{}\nwbegincode{321}\moddef{Eval trees: procedures}\plusendmoddef
  recursive subroutine eval_node_compile_lsinglet (en, pn, var_list)
    type(eval_node_t), pointer :: en
    type(parse_node_t), intent(in) :: pn
    type(var_list_t), intent(in), target :: var_list
    type(parse_node_t), pointer :: pn_term, pn_alternative, pn_arg
    type(eval_node_t), pointer :: en1, en2
    if (debug) then
       print *, "read lsinglet";  call parse_node_write (pn)
    end if
    pn_term => parse_node_get_sub_ptr (pn, tag="lterm")
    call eval_node_compile_lterm (en, pn_term, var_list)
    pn_alternative => parse_node_get_next_ptr (pn_term, tag="alternative")
    do while (associated (pn_alternative))
       pn_arg => parse_node_get_sub_ptr (pn_alternative, 2, tag="lterm")
       en1 => en
       call eval_node_compile_lterm (en2, pn_arg, var_list)
       allocate (en)
       if (en1%type == EN_CONSTANT .and. en2%type == EN_CONSTANT) then
          call eval_node_init_log (en, or_ll (en1, en2))
          call eval_node_final_rec (en1)
          call eval_node_final_rec (en2)
          deallocate (en1, en2)
       else   
          call eval_node_init_branch &
               (en, var_str ("alternative"), V_LOG, en1, en2)
          call eval_node_set_op2_log (en, or_ll)
       end if
       pn_alternative => parse_node_get_next_ptr (pn_alternative)
    end do
    if (debug) then
       call eval_node_write (en)
       print *, "done lsinglet"
    end if
  end subroutine eval_node_compile_lsinglet

\eatline
\nwendcode{}\nwbegindocs{322}A logical term consists of one or more logical values
concatenated by {\tt{}and}.
\nwenddocs{}\nwbegincode{323}\moddef{Eval trees: procedures}\plusendmoddef
  recursive subroutine eval_node_compile_lterm (en, pn, var_list)
    type(eval_node_t), pointer :: en
    type(parse_node_t), intent(in) :: pn
    type(var_list_t), intent(in), target :: var_list
    type(parse_node_t), pointer :: pn_term, pn_coincidence, pn_arg
    type(eval_node_t), pointer :: en1, en2
    if (debug) then
       print *, "read lterm";  call parse_node_write (pn)
    end if
    pn_term => parse_node_get_sub_ptr (pn)
    call eval_node_compile_lvalue (en, pn_term, var_list)
    pn_coincidence => parse_node_get_next_ptr (pn_term, tag="coincidence")
    do while (associated (pn_coincidence))
       pn_arg => parse_node_get_sub_ptr (pn_coincidence, 2)
       en1 => en
       call eval_node_compile_lvalue (en2, pn_arg, var_list)
       allocate (en)
       if (en1%type == EN_CONSTANT .and. en2%type == EN_CONSTANT) then
          call eval_node_init_log (en, and_ll (en1, en2))
          call eval_node_final_rec (en1)
          call eval_node_final_rec (en2)
          deallocate (en1, en2)
       else   
          call eval_node_init_branch &
               (en, var_str ("coincidence"), V_LOG, en1, en2)
          call eval_node_set_op2_log (en, and_ll)
       end if
       pn_coincidence => parse_node_get_next_ptr (pn_coincidence)
    end do
    if (debug) then
       call eval_node_write (en)
       print *, "done lterm"
    end if
  end subroutine eval_node_compile_lterm

\eatline
\nwendcode{}\nwbegindocs{324}Logical variables are disabled, because they are confused with the
l.h.s.\ of compared expressions.
\nwenddocs{}\nwbegincode{325}\moddef{Eval trees: procedures}\plusendmoddef
  recursive subroutine eval_node_compile_lvalue (en, pn, var_list)
    type(eval_node_t), pointer :: en
    type(parse_node_t), intent(in) :: pn
    type(var_list_t), intent(in), target :: var_list
    if (debug) then
       print *, "read lvalue";  call parse_node_write (pn)
    end if
    select case (char (parse_node_get_rule_key (pn)))
    case ("true")
       allocate (en)
       call eval_node_init_log (en, .true.)
    case ("false")
       allocate (en)
       call eval_node_init_log (en, .false.)
    case ("negation")
       call eval_node_compile_negation (en, pn, var_list)
    case ("lvariable")
       call eval_node_compile_variable (en, pn, var_list, V_LOG)
    case ("lexpr")
       call eval_node_compile_lexpr (en, pn, var_list)
    case ("block_lexpr")
       call eval_node_compile_block_expr (en, pn, var_list, V_LOG)
    case ("conditional_lexpr")
       call eval_node_compile_conditional (en, pn, var_list, V_LOG)
    case ("compared_expr")
       call eval_node_compile_compared_expr (en, pn, var_list, V_REAL)
    case ("compared_sexpr")
       call eval_node_compile_compared_expr (en, pn, var_list, V_STR)
    case ("all_fun", "any_fun", "no_fun", "user_cut_fun")
       call eval_node_compile_log_function (en, pn, var_list)
    case ("record_cmd")
       call eval_node_compile_record_cmd (en, pn, var_list)
    case default
       call parse_node_mismatch &
            ("true|false|negation|lvariable|" // &
             "lexpr|block_lexpr|conditional_lexpr|" // &
             "compared_expr|compared_sexpr|logical_pexpr", pn)
    end select
    if (debug) then
       call eval_node_write (en)
       print *, "done lvalue"
    end if
  end subroutine eval_node_compile_lvalue

\eatline
\nwendcode{}\nwbegindocs{326}A negation consists of the keyword {\tt{}not} and a logical value.
\nwenddocs{}\nwbegincode{327}\moddef{Eval trees: procedures}\plusendmoddef
  recursive subroutine eval_node_compile_negation (en, pn, var_list)
    type(eval_node_t), pointer :: en
    type(parse_node_t), intent(in) :: pn
    type(var_list_t), intent(in), target :: var_list
    type(parse_node_t), pointer :: pn_arg
    type(eval_node_t), pointer :: en1
    if (debug) then
       print *, "read negation";  call parse_node_write (pn)
    end if
    pn_arg => parse_node_get_sub_ptr (pn, 2)
    call eval_node_compile_lvalue (en1, pn_arg, var_list)
    allocate (en)
    if (en1%type == EN_CONSTANT) then
       call eval_node_init_log (en, not_l (en1))
       call eval_node_final_rec (en1)
       deallocate (en1)
    else
       call eval_node_init_branch (en, var_str ("not"), V_LOG, en1)
       call eval_node_set_op1_log (en, not_l)
    end if
    if (debug) then
       call eval_node_write (en)
       print *, "done negation"
    end if
  end subroutine eval_node_compile_negation

\eatline
\nwendcode{}\nwbegindocs{328}\nwdocspar
\subsubsection{Comparisons}
Up to the loop, this is easy.  There is always at least one
comparison.  This is evaluated, and the result is the logical node
{\tt{}en}.  If it is constant, we keep its second sub-node as {\tt{}en2}.
(Thus, at the very end {\tt{}en2} has to be deleted if {\tt{}en} is (still)
constant.)

If there is another comparison, we first check if the first comparison
was constant.  In that case, there are two possibilities: (i) it was
true.  Then, its right-hand side is compared with the new right-hand
side, and the result replaces the previous one which is deleted.  (ii)
it was false.  In this case, the result of the whole comparison is
false, and we can exit the loop without evaluating anything else.

Now assume that the first comparison results in a valid branch, its
second sub-node kept as {\tt{}en2}.  We first need a copy of this, which
becomes the new left-hand side.  If {\tt{}en2} is constant, we make an
identical constant node {\tt{}en1}.  Otherwise, we make {\tt{}en1} an
appropriate pointer node.  Next, the first branch is saved as {\tt{}en0}
and we evaluate the comparison between {\tt{}en1} and the a right-hand
side.  If this turns out to be constant, there are again two
possibilities: (i) true, then we revert to the previous result.  (ii)
false, then the wh
\nwenddocs{}\nwbegincode{329}\moddef{Eval trees: procedures}\plusendmoddef
  recursive subroutine eval_node_compile_compared_expr (en, pn, var_list, type)
    type(eval_node_t), pointer :: en
    type(parse_node_t), intent(in) :: pn
    type(var_list_t), intent(in), target :: var_list
    integer, intent(in) :: type
    type(parse_node_t), pointer :: pn_comparison, pn_expr1
    type(eval_node_t), pointer :: en0, en1, en2
    if (debug) then
       print *, "read comparison";  call parse_node_write (pn)
    end if
    select case (type)
    case (V_INT, V_REAL)
       pn_expr1 => parse_node_get_sub_ptr (pn, tag="expr")
       call eval_node_compile_expr (en1, pn_expr1, var_list)
       pn_comparison => parse_node_get_next_ptr (pn_expr1, tag="comparison")
    case (V_STR)
       pn_expr1 => parse_node_get_sub_ptr (pn, tag="sexpr")
       call eval_node_compile_sexpr (en1, pn_expr1, var_list)
       pn_comparison => parse_node_get_next_ptr (pn_expr1, tag="str_comparison")
    end select
    call eval_node_compile_comparison &
         (en, en1, en2, pn_comparison, var_list, type)
    pn_comparison => parse_node_get_next_ptr (pn_comparison)
    SCAN_FURTHER: do while (associated (pn_comparison))
       if (en%type == EN_CONSTANT) then
          if (en%lval) then
             en1 => en2
             call eval_node_final_rec (en);  deallocate (en)
             call eval_node_compile_comparison &
                  (en, en1, en2, pn_comparison, var_list, type)
          else
             exit SCAN_FURTHER
          end if
       else
          allocate (en1)
          if (en2%type == EN_CONSTANT) then
             select case (en2%result_type)
             case (V_INT);  call eval_node_init_int    (en1, en2%ival)
             case (V_REAL); call eval_node_init_real   (en1, en2%rval)
             case (V_STR);  call eval_node_init_string (en1, en2%sval)
             end select
          else
             select case (en2%result_type)
             case (V_INT);  call eval_node_init_int_ptr &
                  (en1, var_str ("(previous)"), en2%ival, en2%value_is_known)
             case (V_REAL); call eval_node_init_real_ptr &
                  (en1, var_str ("(previous)"), en2%rval, en2%value_is_known)
             case (V_STR);  call eval_node_init_string_ptr &
                  (en1, var_str ("(previous)"), en2%sval, en2%value_is_known)
             end select
          end if
          en0 => en
          call eval_node_compile_comparison &
               (en, en1, en2, pn_comparison, var_list, type)
          if (en%type == EN_CONSTANT) then
             if (en%lval) then
                call eval_node_final_rec (en);  deallocate (en)
                en => en0
             else
                call eval_node_final_rec (en0);  deallocate (en0)
                exit SCAN_FURTHER
             end if
          else
             en1 => en
             allocate (en)
             call eval_node_init_branch (en, var_str ("and"), V_LOG, en0, en1)
             call eval_node_set_op2_log (en, and_ll)
          end if
       end if
       pn_comparison => parse_node_get_next_ptr (pn_comparison)
    end do SCAN_FURTHER
    if (en%type == EN_CONSTANT .and. associated (en2)) then
       call eval_node_final_rec (en2);  deallocate (en2)
    end if
    if (debug) then
       call eval_node_write (en)
       print *, "done compared_expr"
    end if
  end subroutine eval_node_compile_compared_expr

\nwendcode{}\nwbegindocs{330}%dev eval_node_compile_compared_expr
\nwenddocs{}\nwbegindocs{331}This takes two extra arguments: {\tt{}en1}, the left-hand-side of the
comparison, is already allocated and evaluated.  {\tt{}en2} (the
right-hand side) and {\tt{}en} (the result) are allocated by the
routine.  {\tt{}pn} is the parse node which contains the operator and the
right-hand side as subnodes.

If the result of the comparison is constant, {\tt{}en1} is deleted but
{\tt{}en2} is kept, because it may be used in a subsequent comparison.
{\tt{}en} then becomes a constant.  If the result is variable, {\tt{}en}
becomes a branch node which refers to {\tt{}en1} and {\tt{}en2}.
\nwenddocs{}\nwbegincode{332}\moddef{Eval trees: procedures}\plusendmoddef
  recursive subroutine eval_node_compile_comparison &
       (en, en1, en2, pn, var_list, type)
    type(eval_node_t), pointer :: en, en1, en2
    type(parse_node_t), intent(in) :: pn
    type(var_list_t), intent(in), target :: var_list
    integer, intent(in) :: type
    type(parse_node_t), pointer :: pn_op, pn_arg
    type(string_t) :: key
    integer :: t1, t2
    real(default), pointer :: tolerance_ptr
    pn_op => parse_node_get_sub_ptr (pn)
    key = parse_node_get_key (pn_op)
    select case (type)
    case (V_INT, V_REAL)
       pn_arg => parse_node_get_next_ptr (pn_op, tag="expr")
       call eval_node_compile_expr (en2, pn_arg, var_list)
    case (V_STR)
       pn_arg => parse_node_get_next_ptr (pn_op, tag="sexpr")
       call eval_node_compile_sexpr (en2, pn_arg, var_list)
    end select
    t1 = en1%result_type
    t2 = en2%result_type
    allocate (en)
    if (en1%type == EN_CONSTANT .and. en2%type == EN_CONSTANT) then
       call var_list%get_rptr (var_str ("tolerance"), tolerance_ptr)
       en1%tolerance => tolerance_ptr
       select case (char (key))
       case ("<")
          select case (t1)
          case (V_INT)
             select case (t2)
             case (V_INT);  call eval_node_init_log (en, comp_lt_ii (en1, en2))
             case (V_REAL); call eval_node_init_log (en, comp_ll_ir (en1, en2))
             end select
          case (V_REAL)
             select case (t2)
             case (V_INT);  call eval_node_init_log (en, comp_ll_ri (en1, en2))
             case (V_REAL); call eval_node_init_log (en, comp_ll_rr (en1, en2))
             end select
          end select
       case (">")
          select case (t1)
          case (V_INT)
             select case (t2)
             case (V_INT);  call eval_node_init_log (en, comp_gt_ii (en1, en2))
             case (V_REAL); call eval_node_init_log (en, comp_gg_ir (en1, en2))
             end select
          case (V_REAL)
             select case (t2)
             case (V_INT);  call eval_node_init_log (en, comp_gg_ri (en1, en2))
             case (V_REAL); call eval_node_init_log (en, comp_gg_rr (en1, en2))
             end select
          end select
       case ("<=")
          select case (t1)
          case (V_INT)
             select case (t2)
             case (V_INT);  call eval_node_init_log (en, comp_le_ii (en1, en2))
             case (V_REAL); call eval_node_init_log (en, comp_ls_ir (en1, en2))
             end select
          case (V_REAL)
             select case (t2)
             case (V_INT);  call eval_node_init_log (en, comp_ls_ri (en1, en2))
             case (V_REAL); call eval_node_init_log (en, comp_ls_rr (en1, en2))
             end select
          end select
       case (">=")
          select case (t1)
          case (V_INT)
             select case (t2)
             case (V_INT);  call eval_node_init_log (en, comp_ge_ii (en1, en2))
             case (V_REAL); call eval_node_init_log (en, comp_gs_ir (en1, en2))
             end select
          case (V_REAL)
             select case (t2)
             case (V_INT);  call eval_node_init_log (en, comp_gs_ri (en1, en2))
             case (V_REAL); call eval_node_init_log (en, comp_gs_rr (en1, en2))
             end select
          end select
       case ("==")
          select case (t1)
          case (V_INT)
             select case (t2)
             case (V_INT);  call eval_node_init_log (en, comp_eq_ii (en1, en2))
             case (V_REAL); call eval_node_init_log (en, comp_se_ir (en1, en2))
             end select
          case (V_REAL)
             select case (t2)
             case (V_INT);  call eval_node_init_log (en, comp_se_ri (en1, en2))
             case (V_REAL); call eval_node_init_log (en, comp_se_rr (en1, en2))
             end select
          case (V_STR)
             select case (t2)
             case (V_STR);  call eval_node_init_log (en, comp_eq_ss (en1, en2))
             end select
          end select
       case ("<>")
          select case (t1)
          case (V_INT)
             select case (t2)
             case (V_INT);  call eval_node_init_log (en, comp_ne_ii (en1, en2))
             case (V_REAL); call eval_node_init_log (en, comp_ns_ir (en1, en2))
             end select
          case (V_REAL)
             select case (t2)
             case (V_INT);  call eval_node_init_log (en, comp_ns_ri (en1, en2))
             case (V_REAL); call eval_node_init_log (en, comp_ns_rr (en1, en2))
             end select
          case (V_STR)
             select case (t2)
             case (V_STR);  call eval_node_init_log (en, comp_ne_ss (en1, en2))
             end select
          end select
       end select
       call eval_node_final_rec (en1)
       deallocate (en1)
    else
       call eval_node_init_branch (en, key, V_LOG, en1, en2)
       select case (char (key))
       case ("<")
          select case (t1)
          case (V_INT)
             select case (t2)
             case (V_INT);  call eval_node_set_op2_log (en, comp_lt_ii)
             case (V_REAL); call eval_node_set_op2_log (en, comp_ll_ir)
             end select
          case (V_REAL)
             select case (t2)
             case (V_INT);  call eval_node_set_op2_log (en, comp_ll_ri)
             case (V_REAL); call eval_node_set_op2_log (en, comp_ll_rr)
             end select
          end select
       case (">")
          select case (t1)
          case (V_INT)
             select case (t2)
             case (V_INT);  call eval_node_set_op2_log (en, comp_gt_ii)
             case (V_REAL); call eval_node_set_op2_log (en, comp_gg_ir)
             end select
          case (V_REAL)
             select case (t2)
             case (V_INT);  call eval_node_set_op2_log (en, comp_gg_ri)
             case (V_REAL); call eval_node_set_op2_log (en, comp_gg_rr)
             end select
          end select
       case ("<=")
          select case (t1)
          case (V_INT)
             select case (t2)
             case (V_INT);  call eval_node_set_op2_log (en, comp_le_ii)
             case (V_REAL); call eval_node_set_op2_log (en, comp_ls_ir)
             end select
          case (V_REAL)
             select case (t2)
             case (V_INT);  call eval_node_set_op2_log (en, comp_ls_ri)
             case (V_REAL); call eval_node_set_op2_log (en, comp_ls_rr)
             end select
          end select
       case (">=")
          select case (t1)
          case (V_INT)
             select case (t2)
             case (V_INT);  call eval_node_set_op2_log (en, comp_ge_ii)
             case (V_REAL); call eval_node_set_op2_log (en, comp_gs_ir)
             end select
          case (V_REAL)
             select case (t2)
             case (V_INT);  call eval_node_set_op2_log (en, comp_gs_ri)
             case (V_REAL); call eval_node_set_op2_log (en, comp_gs_rr)
             end select
          end select
       case ("==")
          select case (t1)
          case (V_INT)
             select case (t2)
             case (V_INT);  call eval_node_set_op2_log (en, comp_eq_ii)
             case (V_REAL); call eval_node_set_op2_log (en, comp_se_ir)
             end select
          case (V_REAL)
             select case (t2)
             case (V_INT);  call eval_node_set_op2_log (en, comp_se_ri)
             case (V_REAL); call eval_node_set_op2_log (en, comp_se_rr)
             end select
          case (V_STR)
             select case (t2)
             case (V_STR);  call eval_node_set_op2_log (en, comp_eq_ss)
             end select
          end select
       case ("<>")
          select case (t1)
          case (V_INT)
             select case (t2)
             case (V_INT);  call eval_node_set_op2_log (en, comp_ne_ii)
             case (V_REAL); call eval_node_set_op2_log (en, comp_ns_ir)
             end select
          case (V_REAL)
             select case (t2)
             case (V_INT);  call eval_node_set_op2_log (en, comp_ns_ri)
             case (V_REAL); call eval_node_set_op2_log (en, comp_ns_rr)
             end select
          case (V_STR)
             select case (t2)
             case (V_STR);  call eval_node_set_op2_log (en, comp_ne_ss)
             end select
          end select
       end select
       call var_list%get_rptr (var_str ("tolerance"), tolerance_ptr)
       en1%tolerance => tolerance_ptr
    end if
  end subroutine eval_node_compile_comparison

\eatline
\nwendcode{}\nwbegindocs{333}\nwdocspar
\subsubsection{Recording analysis data}
The {\tt{}record} command is actually a logical expression which always
evaluates {\tt{}true}.
\nwenddocs{}\nwbegincode{334}\moddef{Eval trees: procedures}\plusendmoddef
  recursive subroutine eval_node_compile_record_cmd (en, pn, var_list)
    type(eval_node_t), pointer :: en
    type(parse_node_t), intent(in) :: pn
    type(var_list_t), intent(in), target :: var_list
    type(parse_node_t), pointer :: pn_key, pn_tag, pn_arg
    type(parse_node_t), pointer :: pn_arg1, pn_arg2, pn_arg3, pn_arg4
    type(eval_node_t), pointer :: en0, en1, en2, en3, en4
    real(default), pointer :: event_weight
    if (debug) then
       print *, "read record_cmd";  call parse_node_write (pn)
    end if
    pn_key => parse_node_get_sub_ptr (pn)
    pn_tag => parse_node_get_next_ptr (pn_key)
    pn_arg => parse_node_get_next_ptr (pn_tag)
    select case (char (parse_node_get_key (pn_key)))
    case ("record")
       call var_list%get_rptr (var_str ("event_weight"), event_weight)
    case ("record_unweighted")
       event_weight => null ()
    case ("record_excess")
       call var_list%get_rptr (var_str ("event_excess"), event_weight)
    end select
    select case (char (parse_node_get_rule_key (pn_tag)))
    case ("analysis_id")
       allocate (en0)
       call eval_node_init_string (en0, parse_node_get_string (pn_tag))
    case default
       call eval_node_compile_sexpr (en0, pn_tag, var_list)
    end select
    allocate (en)
    if (associated (pn_arg)) then
       pn_arg1 => parse_node_get_sub_ptr (pn_arg)
       call eval_node_compile_expr (en1, pn_arg1, var_list)
       if (en1%result_type == V_INT) &
            call insert_conversion_node (en1, V_REAL)
       pn_arg2 => parse_node_get_next_ptr (pn_arg1)
       if (associated (pn_arg2)) then
          call eval_node_compile_expr (en2, pn_arg2, var_list)
          if (en2%result_type == V_INT) &
               call insert_conversion_node (en2, V_REAL)
          pn_arg3 => parse_node_get_next_ptr (pn_arg2)
          if (associated (pn_arg3)) then
             call eval_node_compile_expr (en3, pn_arg3, var_list)
             if (en3%result_type == V_INT) &
                  call insert_conversion_node (en3, V_REAL)
             pn_arg4 => parse_node_get_next_ptr (pn_arg3)
             if (associated (pn_arg4)) then
                call eval_node_compile_expr (en4, pn_arg4, var_list)
                if (en4%result_type == V_INT) &
                     call insert_conversion_node (en4, V_REAL)
                call eval_node_init_record_cmd &
                     (en, event_weight, en0, en1, en2, en3, en4)
             else
                call eval_node_init_record_cmd &
                     (en, event_weight, en0, en1, en2, en3)
             end if
          else
             call eval_node_init_record_cmd (en, event_weight, en0, en1, en2)
          end if
       else
          call eval_node_init_record_cmd (en, event_weight, en0, en1)
       end if
    else
       call eval_node_init_record_cmd (en, event_weight, en0)
    end if
    if (debug) then
       call eval_node_write (en)
       print *, "done record_cmd"
    end if
  end subroutine eval_node_compile_record_cmd

\eatline
\nwendcode{}\nwbegindocs{335}\nwdocspar
\subsubsection{Particle-list expressions}
A particle expression is a subevent or a concatenation of
particle-list terms (using \verb|join|).
\nwenddocs{}\nwbegincode{336}\moddef{Eval trees: procedures}\plusendmoddef
  recursive subroutine eval_node_compile_pexpr (en, pn, var_list)
    type(eval_node_t), pointer :: en
    type(parse_node_t), intent(in) :: pn
    type(var_list_t), intent(in), target :: var_list
    type(parse_node_t), pointer :: pn_pterm, pn_concatenation, pn_op, pn_arg
    type(eval_node_t), pointer :: en1, en2
    type(subevt_t) :: subevt
    if (debug) then
       print *, "read pexpr";  call parse_node_write (pn)
    end if
    pn_pterm => parse_node_get_sub_ptr (pn)
    call eval_node_compile_pterm (en, pn_pterm, var_list)
    pn_concatenation => &
         parse_node_get_next_ptr (pn_pterm, tag="pconcatenation")
    do while (associated (pn_concatenation))
       pn_op => parse_node_get_sub_ptr (pn_concatenation)
       pn_arg => parse_node_get_next_ptr (pn_op)
       en1 => en
       call eval_node_compile_pterm (en2, pn_arg, var_list)
       allocate (en)
       if (en1%type == EN_CONSTANT .and. en2%type == EN_CONSTANT) then
          call subevt_join (subevt, en1%pval, en2%pval)
          call eval_node_init_subevt (en, subevt)
          call eval_node_final_rec (en1)
          call eval_node_final_rec (en2)
          deallocate (en1, en2)
       else   
          call eval_node_init_branch &
               (en, var_str ("join"), V_SEV, en1, en2)
          call eval_node_set_op2_sev (en, join_pp)
       end if
       pn_concatenation => parse_node_get_next_ptr (pn_concatenation)
    end do
    if (debug) then
       call eval_node_write (en)
       print *, "done pexpr"
    end if
  end subroutine eval_node_compile_pexpr

\eatline
\nwendcode{}\nwbegindocs{337}A particle term is a subevent or a combination of
particle-list values (using \verb|combine|).
\nwenddocs{}\nwbegincode{338}\moddef{Eval trees: procedures}\plusendmoddef
  recursive subroutine eval_node_compile_pterm (en, pn, var_list)
    type(eval_node_t), pointer :: en
    type(parse_node_t), intent(in) :: pn
    type(var_list_t), intent(in), target :: var_list
    type(parse_node_t), pointer :: pn_pvalue, pn_combination, pn_op, pn_arg
    type(eval_node_t), pointer :: en1, en2
    type(subevt_t) :: subevt
    if (debug) then
       print *, "read pterm";  call parse_node_write (pn)
    end if
    pn_pvalue => parse_node_get_sub_ptr (pn)
    call eval_node_compile_pvalue (en, pn_pvalue, var_list)
    pn_combination => &
         parse_node_get_next_ptr (pn_pvalue, tag="pcombination")
    do while (associated (pn_combination))
       pn_op => parse_node_get_sub_ptr (pn_combination)
       pn_arg => parse_node_get_next_ptr (pn_op)
       en1 => en
       call eval_node_compile_pvalue (en2, pn_arg, var_list)
       allocate (en)
       if (en1%type == EN_CONSTANT .and. en2%type == EN_CONSTANT) then
          call subevt_combine (subevt, en1%pval, en2%pval)
          call eval_node_init_subevt (en, subevt)
          call eval_node_final_rec (en1)
          call eval_node_final_rec (en2)
          deallocate (en1, en2)
       else   
          call eval_node_init_branch &
               (en, var_str ("combine"), V_SEV, en1, en2)
          call eval_node_set_op2_sev (en, combine_pp)
       end if
       pn_combination => parse_node_get_next_ptr (pn_combination)
    end do
    if (debug) then
       call eval_node_write (en)
       print *, "done pterm"
    end if
  end subroutine eval_node_compile_pterm

\eatline
\nwendcode{}\nwbegindocs{339}A particle-list value is a PDG-code array, a particle identifier, a
variable, a (grouped) pexpr, a block pexpr, a conditional, or a
particle-list function.

The {\tt{}cexpr} node is responsible for transforming a constant PDG-code
array into a subevent.  It takes the code array as its first
argument, the event subevent as its second argument, and the
requested particle type (incoming/outgoing) as its zero-th argument.
The result is the list of particles in the event that match the code
array.
\nwenddocs{}\nwbegincode{340}\moddef{Eval trees: procedures}\plusendmoddef
  recursive subroutine eval_node_compile_pvalue (en, pn, var_list)
    type(eval_node_t), pointer :: en
    type(parse_node_t), intent(in) :: pn
    type(var_list_t), intent(in), target :: var_list
    type(parse_node_t), pointer :: pn_prefix_cexpr
    type(eval_node_t), pointer :: en1, en2, en0
    type(string_t) :: key
    type(subevt_t), pointer :: evt_ptr
    logical, pointer :: known
    if (debug) then
       print *, "read pvalue";  call parse_node_write (pn)
    end if
    select case (char (parse_node_get_rule_key (pn)))
    case ("pexpr_src")
       call eval_node_compile_prefix_cexpr (en1, pn, var_list)
       allocate (en2)
       if (var_list%contains (var_str ("@evt"))) then
          call var_list%get_pptr (var_str ("@evt"), evt_ptr, known)
          call eval_node_init_subevt_ptr (en2, var_str ("@evt"), evt_ptr, known)
          allocate (en)
          call eval_node_init_branch &
               (en, var_str ("prt_selection"), V_SEV, en1, en2)
          call eval_node_set_op2_sev (en, select_pdg_ca)
          allocate (en0)
          pn_prefix_cexpr => parse_node_get_sub_ptr (pn)
          key = parse_node_get_rule_key (pn_prefix_cexpr)
          select case (char (key))
          case ("incoming_prt")
             call eval_node_init_int (en0, PRT_INCOMING)
             en%arg0 => en0
          case ("outgoing_prt")
             call eval_node_init_int (en0, PRT_OUTGOING)
             en%arg0 => en0
          end select
       else
          call parse_node_write (pn)
          call msg_bug (" Missing event data while compiling pvalue")
       end if
    case ("pvariable")
       call eval_node_compile_variable (en, pn, var_list, V_SEV)
    case ("pexpr")
       call eval_node_compile_pexpr (en, pn, var_list)
    case ("block_pexpr")
       call eval_node_compile_block_expr (en, pn, var_list, V_SEV)
    case ("conditional_pexpr")
       call eval_node_compile_conditional (en, pn, var_list, V_SEV)
    case ("join_fun", "combine_fun", "collect_fun", "cluster_fun", &
          "select_fun", "extract_fun", "sort_fun")
       call eval_node_compile_prt_function (en, pn, var_list)
    case default
       call parse_node_mismatch &
            ("prefix_cexpr|pvariable|" // &
             "grouped_pexpr|block_pexpr|conditional_pexpr|" // &
             "prt_function", pn)
    end select
    if (debug) then
       call eval_node_write (en)
       print *, "done pvalue"
    end if
  end subroutine eval_node_compile_pvalue

\eatline
\nwendcode{}\nwbegindocs{341}\nwdocspar
\subsubsection{Particle functions}
This combines the treatment of 'join', 'combine', 'collect', 'cluster',
'select', 
and 'extract' which all have the same syntax.  The one or two argument
nodes are allocated.  If there is a condition, the condition node is
also allocated as a logical expression, for which the variable list is
augmented by the appropriate (unary/binary) observables.
\nwenddocs{}\nwbegincode{342}\moddef{Eval trees: procedures}\plusendmoddef
  recursive subroutine eval_node_compile_prt_function (en, pn, var_list)
    type(eval_node_t), pointer :: en
    type(parse_node_t), intent(in) :: pn
    type(var_list_t), intent(in), target :: var_list
    type(parse_node_t), pointer :: pn_clause, pn_key, pn_cond, pn_args
    type(parse_node_t), pointer :: pn_arg0, pn_arg1, pn_arg2
    type(eval_node_t), pointer :: en0, en1, en2
    type(string_t) :: key
    if (debug) then
       print *, "read prt_function";  call parse_node_write (pn)
    end if
    pn_clause => parse_node_get_sub_ptr (pn)
    pn_key  => parse_node_get_sub_ptr (pn_clause)
    pn_cond => parse_node_get_next_ptr (pn_key)
    if (associated (pn_cond)) &
         pn_arg0 => parse_node_get_sub_ptr (pn_cond, 2)
    pn_args => parse_node_get_next_ptr (pn_clause)
    pn_arg1 => parse_node_get_sub_ptr (pn_args)
    pn_arg2 => parse_node_get_next_ptr (pn_arg1)
    key = parse_node_get_key (pn_key)
    call eval_node_compile_pexpr (en1, pn_arg1, var_list)
    allocate (en)
    if (.not. associated (pn_arg2)) then
       select case (char (key))
       case ("collect")
          call eval_node_init_prt_fun_unary (en, en1, key, collect_p)
       case ("cluster")
          if (fastjet_available ()) then
             call fastjet_init ()
          else
             call msg_fatal &
               ("'cluster' function requires FastJet, which is not enabled")
          end if
          call eval_node_init_prt_fun_unary (en, en1, key, cluster_p)
          call var_list%get_iptr (var_str ("jet_algorithm"), en1%jet_algorithm)
          call var_list%get_rptr (var_str ("jet_r"), en1%jet_r)
       case ("select")
          call eval_node_init_prt_fun_unary (en, en1, key, select_p)
       case ("extract")
          call eval_node_init_prt_fun_unary (en, en1, key, extract_p)
       case ("sort")
          call eval_node_init_prt_fun_unary (en, en1, key, sort_p)
       case default
          call msg_bug (" Unary particle function '" // char (key) // &
               "' undefined")
       end select
    else
       call eval_node_compile_pexpr (en2, pn_arg2, var_list)
       select case (char (key))
       case ("join")
          call eval_node_init_prt_fun_binary (en, en1, en2, key, join_pp)
       case ("combine")
          call eval_node_init_prt_fun_binary (en, en1, en2, key, combine_pp)
       case ("collect")
          call eval_node_init_prt_fun_binary (en, en1, en2, key, collect_pp)
       case ("select")
          call eval_node_init_prt_fun_binary (en, en1, en2, key, select_pp)
       case ("sort")
          call eval_node_init_prt_fun_binary (en, en1, en2, key, sort_pp)
       case default
          call msg_bug (" Binary particle function '" // char (key) // &
               "' undefined")
       end select
    end if
    if (associated (pn_cond)) then
       call eval_node_set_observables (en, var_list)
       select case (char (key))
       case ("extract", "sort")
          call eval_node_compile_expr (en0, pn_arg0, en%var_list)
       case default
          call eval_node_compile_lexpr (en0, pn_arg0, en%var_list)
       end select
       en%arg0 => en0
    end if
    if (debug) then
       call eval_node_write (en)
       print *, "done prt_function"
    end if
  end subroutine eval_node_compile_prt_function

\eatline
\nwendcode{}\nwbegindocs{343}The {\tt{}eval} expression is similar, but here the expression {\tt{}arg0}
is mandatory, and the whole thing evaluates to a numeric value.
\nwenddocs{}\nwbegincode{344}\moddef{Eval trees: procedures}\plusendmoddef
  recursive subroutine eval_node_compile_eval_function (en, pn, var_list)
    type(eval_node_t), pointer :: en
    type(parse_node_t), intent(in) :: pn
    type(var_list_t), intent(in), target :: var_list
    type(parse_node_t), pointer :: pn_key, pn_arg0, pn_args, pn_arg1, pn_arg2
    type(eval_node_t), pointer :: en0, en1, en2
    type(string_t) :: key
    if (debug) then
       print *, "read eval_function";  call parse_node_write (pn)
    end if
    pn_key => parse_node_get_sub_ptr (pn)
    pn_arg0 => parse_node_get_next_ptr (pn_key)
    pn_args => parse_node_get_next_ptr (pn_arg0)
    pn_arg1 => parse_node_get_sub_ptr (pn_args)
    pn_arg2 => parse_node_get_next_ptr (pn_arg1)
    key = parse_node_get_key (pn_key)
    call eval_node_compile_pexpr (en1, pn_arg1, var_list)
    allocate (en)
    if (.not. associated (pn_arg2)) then
       call eval_node_init_eval_fun_unary (en, en1, key)
    else
       call eval_node_compile_pexpr (en2, pn_arg2, var_list)
       call eval_node_init_eval_fun_binary (en, en1, en2, key)
    end if
    call eval_node_set_observables (en, var_list)
    call eval_node_compile_expr (en0, pn_arg0, en%var_list)
    if (en0%result_type /= V_REAL) &
         call msg_fatal (" 'eval' function does not result in real value")
    call eval_node_set_expr (en, en0)
    if (debug) then
       call eval_node_write (en)
       print *, "done eval_function"
    end if
  end subroutine eval_node_compile_eval_function

\eatline
\nwendcode{}\nwbegindocs{345}Logical functions of subevents.
\nwenddocs{}\nwbegincode{346}\moddef{Eval trees: procedures}\plusendmoddef
  recursive subroutine eval_node_compile_log_function (en, pn, var_list)
    type(eval_node_t), pointer :: en
    type(parse_node_t), intent(in) :: pn
    type(var_list_t), intent(in), target :: var_list
    type(parse_node_t), pointer :: pn_key, pn_str
    type(parse_node_t), pointer :: pn_arg0, pn_args, pn_arg1, pn_arg2
    type(eval_node_t), pointer :: en0, en1, en2
    type(string_t) :: key
    if (debug) then
       print *, "read log_function";  call parse_node_write (pn)
    end if
    select case (char (parse_node_get_rule_key (pn)))
    case ("all_fun", "any_fun", "no_fun")
       pn_key => parse_node_get_sub_ptr (pn)
       pn_arg0 => parse_node_get_next_ptr (pn_key)
       pn_args => parse_node_get_next_ptr (pn_arg0)
    case ("user_cut_fun")
       pn_key => parse_node_get_sub_ptr (pn)
       pn_str => parse_node_get_next_ptr (pn_key)
       pn_arg0 => parse_node_get_sub_ptr (pn_str)
       pn_args => parse_node_get_next_ptr (pn_str)
    case default
       call parse_node_mismatch &
            ("all_fun|any_fun|no_fun|user_cut_fun", &
            pn)
    end select
    pn_arg1 => parse_node_get_sub_ptr (pn_args)
    pn_arg2 => parse_node_get_next_ptr (pn_arg1)
    key = parse_node_get_key (pn_key)
    call eval_node_compile_pexpr (en1, pn_arg1, var_list)
    allocate (en)
    if (.not. associated (pn_arg2)) then
       select case (char (key))
       case ("all")
          call eval_node_init_log_fun_unary (en, en1, key, all_p)
       case ("any")
          call eval_node_init_log_fun_unary (en, en1, key, any_p)
       case ("no")
          call eval_node_init_log_fun_unary (en, en1, key, no_p)
       case ("user_cut")
          call eval_node_init_log_fun_unary (en, en1, key, user_cut_p)
       case default
          call msg_bug ("Unary logical particle function '" // char (key) // &
               "' undefined")
       end select
    else
       call eval_node_compile_pexpr (en2, pn_arg2, var_list)
       select case (char (key))
       case ("all")
          call eval_node_init_log_fun_binary (en, en1, en2, key, all_pp)
       case ("any")
          call eval_node_init_log_fun_binary (en, en1, en2, key, any_pp)
       case ("no")
          call eval_node_init_log_fun_binary (en, en1, en2, key, no_pp)
       case default
          call msg_bug ("Binary logical particle function '" // char (key) // &
               "' undefined")
       end select
    end if
    if (associated (pn_arg0)) then
       call eval_node_set_observables (en, var_list)
       select case (char (key))
       case ("all", "any", "no")
          call eval_node_compile_lexpr (en0, pn_arg0, en%var_list)
       case ("user_cut")
          call eval_node_compile_sexpr (en0, pn_arg0, en%var_list)
       case default                    
          call msg_bug ("Compiling logical particle function: missing mode")
       end select
       call eval_node_set_expr (en, en0, V_LOG)
    end if
    if (debug) then
       call eval_node_write (en)
       print *, "done log_function"
    end if
  end subroutine eval_node_compile_log_function

\eatline
\nwendcode{}\nwbegindocs{347}Numeric functions of subevents.
\nwenddocs{}\nwbegincode{348}\moddef{Eval trees: procedures}\plusendmoddef
  recursive subroutine eval_node_compile_numeric_function (en, pn, var_list)
    type(eval_node_t), pointer :: en
    type(parse_node_t), intent(in) :: pn
    type(var_list_t), intent(in), target :: var_list
    type(parse_node_t), pointer :: pn_clause, pn_key, pn_cond, pn_args
    type(parse_node_t), pointer :: pn_arg0, pn_arg1, pn_arg2
    type(eval_node_t), pointer :: en0, en1, en2
    type(string_t) :: key
    if (debug) then
       print *, "read numeric_function";  call parse_node_write (pn)
    end if
    select case (char (parse_node_get_rule_key (pn)))
    case ("count_fun")
       pn_clause => parse_node_get_sub_ptr (pn)
       pn_key  => parse_node_get_sub_ptr (pn_clause)
       pn_cond => parse_node_get_next_ptr (pn_key)
       if (associated (pn_cond)) then
          pn_arg0 => parse_node_get_sub_ptr (pn_cond, 2)
       else
          pn_arg0 => null ()
       end if
       pn_args => parse_node_get_next_ptr (pn_clause)
    case ("user_event_fun")
       pn_key => parse_node_get_sub_ptr (pn)
       pn_cond => parse_node_get_next_ptr (pn_key)
       pn_arg0 => parse_node_get_sub_ptr (pn_cond)
       pn_args => parse_node_get_next_ptr (pn_cond)
    end select
    pn_arg1 => parse_node_get_sub_ptr (pn_args)
    pn_arg2 => parse_node_get_next_ptr (pn_arg1)
    key = parse_node_get_key (pn_key)
    call eval_node_compile_pexpr (en1, pn_arg1, var_list)
    allocate (en)
    if (.not. associated (pn_arg2)) then
       select case (char (key))
       case ("count")
          call eval_node_init_int_fun_unary (en, en1, key, count_a)
       case ("user_event_shape")
          call eval_node_init_real_fun_unary (en, en1, key, user_event_shape_a)
       case default
          call msg_bug ("Unary subevent function '" // char (key) // &
               "' undefined")
       end select
    else
       call eval_node_compile_pexpr (en2, pn_arg2, var_list)
       select case (char (key))
       case ("count")
          call eval_node_init_int_fun_binary (en, en1, en2, key, count_pp)
       case default
          call msg_bug ("Binary subevent function '" // char (key) // &
               "' undefined")
       end select
    end if
    if (associated (pn_arg0)) then
       call eval_node_set_observables (en, var_list)
       select case (char (key))
       case ("count")
          call eval_node_compile_lexpr (en0, pn_arg0, en%var_list)
          call eval_node_set_expr (en, en0, V_INT)
       case ("user_event_shape")
          call eval_node_compile_sexpr (en0, pn_arg0, en%var_list)
          call eval_node_set_expr (en, en0, V_REAL)
       end select
    end if
    if (debug) then
       call eval_node_write (en)
       print *, "done numeric_function"
    end if
  end subroutine eval_node_compile_numeric_function

\eatline
\nwendcode{}\nwbegindocs{349}\nwdocspar
\subsubsection{PDG-code arrays}
A PDG-code expression is either prefixed by {\tt{}incoming} or
{\tt{}outgoing}, a block, or a conditional.  In any case, it evaluates to
a constant.
\nwenddocs{}\nwbegincode{350}\moddef{Eval trees: procedures}\plusendmoddef
  recursive subroutine eval_node_compile_prefix_cexpr (en, pn, var_list)
    type(eval_node_t), pointer :: en
    type(parse_node_t), intent(in) :: pn
    type(var_list_t), intent(in), target :: var_list
    type(parse_node_t), pointer :: pn_avalue, pn_prt
    type(string_t) :: key
    if (debug) then
       print *, "read prefix_cexpr";  call parse_node_write (pn)
    end if
    pn_avalue => parse_node_get_sub_ptr (pn)
    key = parse_node_get_rule_key (pn_avalue)
    select case (char (key))
    case ("incoming_prt")
       pn_prt => parse_node_get_sub_ptr (pn_avalue, 2)
       call eval_node_compile_cexpr (en, pn_prt, var_list)
    case ("outgoing_prt")
       pn_prt => parse_node_get_sub_ptr (pn_avalue, 1)
       call eval_node_compile_cexpr (en, pn_prt, var_list)
    case default
       call parse_node_mismatch &
            ("incoming_prt|outgoing_prt", &
             pn_avalue)
    end select
    if (debug) then
       call eval_node_write (en)
       print *, "done prefix_cexpr"
    end if
  end subroutine eval_node_compile_prefix_cexpr

\eatline
\nwendcode{}\nwbegindocs{351}A PDG array is a string of PDG code definitions (or aliases),
concatenated by ':'.  The code definitions may be variables which are
not defined at compile time, so we have to allocate sub-nodes.  This
analogous to {\tt{}eval{\char95}node{\char95}compile{\char95}term}.
\nwenddocs{}\nwbegincode{352}\moddef{Eval trees: procedures}\plusendmoddef
   recursive subroutine eval_node_compile_cexpr (en, pn, var_list)
    type(eval_node_t), pointer :: en
    type(parse_node_t), intent(in) :: pn
    type(var_list_t), intent(in), target :: var_list
    type(parse_node_t), pointer :: pn_prt, pn_concatenation
    type(eval_node_t), pointer :: en1, en2
    type(pdg_array_t) :: aval
    if (debug) then
       print *, "read cexpr";  call parse_node_write (pn)
    end if
    pn_prt => parse_node_get_sub_ptr (pn)
    call eval_node_compile_avalue (en, pn_prt, var_list)
    pn_concatenation => parse_node_get_next_ptr (pn_prt)
    do while (associated (pn_concatenation))
       pn_prt => parse_node_get_sub_ptr (pn_concatenation, 2)
       en1 => en
       call eval_node_compile_avalue (en2, pn_prt, var_list)
       allocate (en)
       if (en1%type == EN_CONSTANT .and. en2%type == EN_CONSTANT) then
          call concat_cc (aval, en1, en2)
          call eval_node_init_pdg_array (en, aval)
          call eval_node_final_rec (en1)
          call eval_node_final_rec (en2)
          deallocate (en1, en2)
       else
          call eval_node_init_branch (en, var_str (":"), V_PDG, en1, en2)
          call eval_node_set_op2_pdg (en, concat_cc)
       end if
       pn_concatenation => parse_node_get_next_ptr (pn_concatenation)
    end do
    if (debug) then
       call eval_node_write (en)
       print *, "done cexpr"
    end if
  end subroutine eval_node_compile_cexpr

\eatline
\nwendcode{}\nwbegindocs{353}Compile a PDG-code type value.  It may be either an integer expression
or a variable of type PDG array, optionally quoted.
\nwenddocs{}\nwbegincode{354}\moddef{Eval trees: procedures}\plusendmoddef
  recursive subroutine eval_node_compile_avalue (en, pn, var_list)
    type(eval_node_t), pointer :: en
    type(parse_node_t), intent(in) :: pn
    type(var_list_t), intent(in), target :: var_list
    if (debug) then
       print *, "read avalue";  call parse_node_write (pn)
    end if
    select case (char (parse_node_get_rule_key (pn)))
    case ("pdg_code")
       call eval_node_compile_pdg_code (en, pn, var_list)
    case ("cvariable", "variable", "prt_name")
       call eval_node_compile_cvariable (en, pn, var_list)
    case ("cexpr")
       call eval_node_compile_cexpr (en, pn, var_list)
    case ("block_cexpr")
       call eval_node_compile_block_expr (en, pn, var_list, V_PDG)
    case ("conditional_cexpr")
       call eval_node_compile_conditional (en, pn, var_list, V_PDG)
    case default
       call parse_node_mismatch &
            ("grouped_cexpr|block_cexpr|conditional_cexpr|" // &
             "pdg_code|cvariable|prt_name", pn)
    end select
    if (debug) then
       call eval_node_write (en)
       print *, "done avalue"
    end if
  end subroutine eval_node_compile_avalue

\eatline
\nwendcode{}\nwbegindocs{355}Compile a PDG-code expression, which is the key {\tt{}PDG} with an
integer expression as argument.  The procedure is analogous to
{\tt{}eval{\char95}node{\char95}compile{\char95}unary{\char95}function}.
\nwenddocs{}\nwbegincode{356}\moddef{Eval trees: procedures}\plusendmoddef
  subroutine eval_node_compile_pdg_code (en, pn, var_list)
    type(eval_node_t), pointer :: en
    type(parse_node_t), intent(in), target :: pn
    type(var_list_t), intent(in), target :: var_list
    type(parse_node_t), pointer :: pn_arg
    type(eval_node_t), pointer :: en1
    type(string_t) :: key
    type(pdg_array_t) :: aval
    integer :: t
    if (debug) then
       print *, "read PDG code";  call parse_node_write (pn)
    end if
    pn_arg => parse_node_get_sub_ptr (pn, 2)
    call eval_node_compile_expr &
         (en1, parse_node_get_sub_ptr (pn_arg, tag="expr"), var_list)
    t = en1%result_type
    allocate (en)
    key = "PDG"
    if (en1%type == EN_CONSTANT) then
       select case (t)
       case (V_INT)
          call pdg_i (aval, en1)
          call eval_node_init_pdg_array (en, aval)
       case default;  call eval_type_error (pn, char (key), t)
       end select
       call eval_node_final_rec (en1)
       deallocate (en1)
    else
       select case (t)
       case (V_INT);  call eval_node_set_op1_pdg (en, pdg_i)
       case default;  call eval_type_error (pn, char (key), t)
       end select
    end if
    if (debug) then
       call eval_node_write (en)
       print *, "done function"
    end if
  end subroutine eval_node_compile_pdg_code

\eatline
\nwendcode{}\nwbegindocs{357}This is entirely analogous to {\tt{}eval{\char95}node{\char95}compile{\char95}variable}.
However, PDG-array variables occur in different contexts.

To avoid name clashes between PDG-array variables and ordinary
variables, we prepend a character ({\tt{}*}).  This is not visible to the
user.
\nwenddocs{}\nwbegincode{358}\moddef{Eval trees: procedures}\plusendmoddef
  subroutine eval_node_compile_cvariable (en, pn, var_list)
    type(eval_node_t), pointer :: en
    type(parse_node_t), intent(in), target :: pn
    type(var_list_t), intent(in), target :: var_list
    type(parse_node_t), pointer :: pn_name
    type(string_t) :: var_name
    type(pdg_array_t), pointer :: aptr
    type(pdg_array_t), target, save :: no_aval
    logical, pointer :: known
    logical, target, save :: unknown = .false.
    if (debug) then
       print *, "read cvariable";  call parse_node_write (pn)
    end if
    pn_name => pn
    var_name = parse_node_get_string (pn_name)
    allocate (en)
    if (var_list%contains (var_name)) then
       call var_list%get_aptr (var_name, aptr, known)
       call eval_node_init_pdg_array_ptr (en, var_name, aptr, known)
    else
       call parse_node_write (pn)
       call msg_error ("This PDG-array variable is undefined at this point")
       call eval_node_init_pdg_array_ptr (en, var_name, no_aval, unknown)
    end if
    if (debug) then
       call eval_node_write (en)
       print *, "done cvariable"
    end if
  end subroutine eval_node_compile_cvariable

\eatline
\nwendcode{}\nwbegindocs{359}\nwdocspar
\subsubsection{String expressions}
A string expression is either a string value or a concatenation of
string values. 
\nwenddocs{}\nwbegincode{360}\moddef{Eval trees: procedures}\plusendmoddef
  recursive subroutine eval_node_compile_sexpr (en, pn, var_list)
    type(eval_node_t), pointer :: en
    type(parse_node_t), intent(in) :: pn
    type(var_list_t), intent(in), target :: var_list
    type(parse_node_t), pointer :: pn_svalue, pn_concatenation, pn_op, pn_arg
    type(eval_node_t), pointer :: en1, en2
    type(string_t) :: string
    if (debug) then
       print *, "read sexpr";  call parse_node_write (pn)
    end if
    pn_svalue => parse_node_get_sub_ptr (pn)
    call eval_node_compile_svalue (en, pn_svalue, var_list)
    pn_concatenation => &
         parse_node_get_next_ptr (pn_svalue, tag="str_concatenation")
    do while (associated (pn_concatenation))
       pn_op => parse_node_get_sub_ptr (pn_concatenation)
       pn_arg => parse_node_get_next_ptr (pn_op)
       en1 => en
       call eval_node_compile_svalue (en2, pn_arg, var_list)
       allocate (en)
       if (en1%type == EN_CONSTANT .and. en2%type == EN_CONSTANT) then
          call concat_ss (string, en1, en2)
          call eval_node_init_string (en, string)
          call eval_node_final_rec (en1)
          call eval_node_final_rec (en2)
          deallocate (en1, en2)
       else   
          call eval_node_init_branch &
               (en, var_str ("concat"), V_STR, en1, en2)
          call eval_node_set_op2_str (en, concat_ss)
       end if
       pn_concatenation => parse_node_get_next_ptr (pn_concatenation)
    end do
    if (debug) then
       call eval_node_write (en)
       print *, "done sexpr"
    end if
  end subroutine eval_node_compile_sexpr

\eatline
\nwendcode{}\nwbegindocs{361}A string value is a string literal, a 
variable, a (grouped) sexpr, a block sexpr, or a conditional.
\nwenddocs{}\nwbegincode{362}\moddef{Eval trees: procedures}\plusendmoddef
  recursive subroutine eval_node_compile_svalue (en, pn, var_list)
    type(eval_node_t), pointer :: en
    type(parse_node_t), intent(in) :: pn
    type(var_list_t), intent(in), target :: var_list
    if (debug) then
       print *, "read svalue";  call parse_node_write (pn)
    end if
    select case (char (parse_node_get_rule_key (pn)))
    case ("svariable")
       call eval_node_compile_variable (en, pn, var_list, V_STR)
    case ("sexpr")
       call eval_node_compile_sexpr (en, pn, var_list)
    case ("block_sexpr")
       call eval_node_compile_block_expr (en, pn, var_list, V_STR)
    case ("conditional_sexpr")
       call eval_node_compile_conditional (en, pn, var_list, V_STR)
    case ("sprintf_fun")
       call eval_node_compile_sprintf (en, pn, var_list)
    case ("string_literal")
       allocate (en)
       call eval_node_init_string (en, parse_node_get_string (pn))
    case default
       call parse_node_mismatch &
            ("svariable|" // &
             "grouped_sexpr|block_sexpr|conditional_sexpr|" // &
             "string_function|string_literal", pn)
    end select
    if (debug) then
       call eval_node_write (en)
       print *, "done svalue"
    end if
  end subroutine eval_node_compile_svalue

\eatline
\nwendcode{}\nwbegindocs{363}There is currently one string function, {\tt{}sprintf}.  For
{\tt{}sprintf}, the first argument (no brackets) is the format string, the
optional arguments in brackets are the expressions or variables to be
formatted.
\nwenddocs{}\nwbegincode{364}\moddef{Eval trees: procedures}\plusendmoddef
  recursive subroutine eval_node_compile_sprintf (en, pn, var_list)
    type(eval_node_t), pointer :: en
    type(parse_node_t), intent(in) :: pn
    type(var_list_t), intent(in), target :: var_list
    type(parse_node_t), pointer :: pn_clause, pn_key, pn_args
    type(parse_node_t), pointer :: pn_arg0
    type(eval_node_t), pointer :: en0, en1
    integer :: n_args
    type(string_t) :: key
    if (debug) then
       print *, "read sprintf_fun";  call parse_node_write (pn)
    end if
    pn_clause => parse_node_get_sub_ptr (pn)
    pn_key  => parse_node_get_sub_ptr (pn_clause)
    pn_arg0 => parse_node_get_next_ptr (pn_key)
    pn_args => parse_node_get_next_ptr (pn_clause)
    call eval_node_compile_sexpr (en0, pn_arg0, var_list)
    if (associated (pn_args)) then
       call eval_node_compile_sprintf_args (en1, pn_args, var_list, n_args)
    else
       n_args = 0
       en1 => null ()
    end if
    allocate (en)
    key = parse_node_get_key (pn_key)
    call eval_node_init_format_string (en, en0, en1, key, n_args)
    if (debug) then
       call eval_node_write (en)
       print *, "done sprintf_fun"
    end if
  end subroutine eval_node_compile_sprintf

\eatline
\nwendcode{}\nwbegincode{365}\moddef{Eval trees: procedures}\plusendmoddef
  subroutine eval_node_compile_sprintf_args (en, pn, var_list, n_args)
    type(eval_node_t), pointer :: en
    type(parse_node_t), intent(in) :: pn
    type(var_list_t), intent(in), target :: var_list
    integer, intent(out) :: n_args
    type(parse_node_t), pointer :: pn_arg
    integer :: i
    type(eval_node_t), pointer :: en1, en2
    n_args = parse_node_get_n_sub (pn)
    en => null ()
    do i = n_args, 1, -1
       pn_arg => parse_node_get_sub_ptr (pn, i)
       select case (char (parse_node_get_rule_key (pn_arg)))
       case ("lvariable")
          call eval_node_compile_variable (en1, pn_arg, var_list, V_LOG)
       case ("svariable")
          call eval_node_compile_variable (en1, pn_arg, var_list, V_STR)
       case ("expr")
          call eval_node_compile_expr (en1, pn_arg, var_list)
       case default
          call parse_node_mismatch ("variable|svariable|lvariable|expr", pn_arg)
       end select
       if (associated (en)) then
          en2 => en
          allocate (en)
          call eval_node_init_branch &
               (en, var_str ("sprintf_arg"), V_NONE, en1, en2)
       else
          allocate (en)
          call eval_node_init_branch &
               (en, var_str ("sprintf_arg"), V_NONE, en1)
       end if
    end do
  end subroutine eval_node_compile_sprintf_args

\eatline
\nwendcode{}\nwbegindocs{366}Evaluation.  We allocate the argument list and apply the Fortran wrapper for
the {\tt{}sprintf} function. 
\nwenddocs{}\nwbegincode{367}\moddef{Eval trees: procedures}\plusendmoddef
  subroutine evaluate_sprintf (string, n_args, en_fmt, en_arg)
    type(string_t), intent(out) :: string
    integer, intent(in) :: n_args
    type(eval_node_t), pointer :: en_fmt
    type(eval_node_t), intent(in), optional, target :: en_arg
    type(eval_node_t), pointer :: en_branch, en_var
    type(sprintf_arg_t), dimension(:), allocatable :: arg
    type(string_t) :: fmt
    logical :: autoformat
    integer :: i, j, sprintf_argc
    autoformat = .not. associated (en_fmt)
    if (autoformat) fmt = ""
    if (present (en_arg)) then
       sprintf_argc = 0
       en_branch => en_arg
       do i = 1, n_args
          select case (en_branch%arg1%result_type)
             case (V_CMPLX); sprintf_argc = sprintf_argc + 2
             case default  ; sprintf_argc = sprintf_argc + 1
          end select
          en_branch => en_branch%arg2
       end do
       allocate (arg (sprintf_argc))
       j = 1
       en_branch => en_arg
       do i = 1, n_args
          en_var => en_branch%arg1
          select case (en_var%result_type)
          case (V_LOG)
             call sprintf_arg_init (arg(j), en_var%lval)
             if (autoformat) fmt = fmt // "%s "
          case (V_INT);
             call sprintf_arg_init (arg(j), en_var%ival)
             if (autoformat) fmt = fmt // "%i "
          case (V_REAL);
             call sprintf_arg_init (arg(j), en_var%rval)
             if (autoformat) fmt = fmt // "%g "
          case (V_STR)
             call sprintf_arg_init (arg(j), en_var%sval)
             if (autoformat) fmt = fmt // "%s "
          case (V_CMPLX)
             call sprintf_arg_init (arg(j), real (en_var%cval, default))
             j = j + 1
             call sprintf_arg_init (arg(j), aimag (en_var%cval))
             if (autoformat) fmt = fmt // "(%g + %g * I) "
          case default
             call eval_node_write (en_var)
             call msg_error ("sprintf is implemented " &
                  // "for logical, integer, real, and string values only")
          end select
          j = j + 1
          en_branch => en_branch%arg2
       end do
    else
       allocate (arg(0))
    end if
    if (autoformat) then
       string = sprintf (trim (fmt), arg)
    else
       string = sprintf (en_fmt%sval, arg)
    end if
  end subroutine evaluate_sprintf

\eatline
\nwendcode{}\nwbegindocs{368}\nwdocspar
\subsection{Auxiliary functions for the compiler}
Issue an error that the current node could not be compiled because of
type mismatch:
\nwenddocs{}\nwbegincode{369}\moddef{Eval trees: procedures}\plusendmoddef
  subroutine eval_type_error (pn, string, t)
    type(parse_node_t), intent(in) :: pn
    character(*), intent(in) :: string
    integer, intent(in) :: t
    type(string_t) :: type
    select case (t)
    case (V_NONE); type = "(none)"
    case (V_LOG);  type = "'logical'"
    case (V_INT);  type = "'integer'"
    case (V_REAL); type = "'real'"
    case (V_CMPLX); type = "'complex'"    
    case default;  type = "(unknown)"
    end select
    call parse_node_write (pn)
    call msg_fatal (" The " // string // &
         " operation is not defined for the given argument type " // &
         char (type))
  end subroutine eval_type_error

\eatline
\nwendcode{}\nwbegindocs{370}\nwdocspar
If two numerics are combined, the result is integer if both
arguments are integer, if one is integer and the other real or both
are real, than its argument is real, otherwise complex.
\nwenddocs{}\nwbegincode{371}\moddef{Eval trees: procedures}\plusendmoddef
  function numeric_result_type (t1, t2) result (t)
    integer, intent(in) :: t1, t2
    integer :: t
    if (t1 == V_INT .and. t2 == V_INT) then
       t = V_INT
    else if (t1 == V_INT .and. t2 == V_REAL) then
       t = V_REAL
    else if (t1 == V_REAL .and. t2 == V_INT) then
       t = V_REAL
    else if (t1 == V_REAL .and. t2 == V_REAL) then
       t = V_REAL   
    else
       t = V_CMPLX
    end if
  end function numeric_result_type

\eatline
\nwendcode{}\nwbegindocs{372}\nwdocspar
\subsection{Evaluation}
Evaluation is done recursively.  For leaf nodes nothing is to be done.

Evaluating particle-list functions: First, we evaluate the particle
lists.  If a condition is present, we assign the particle pointers of
the condition node to the allocated particle entries in the parent
node, keeping in mind that the observables in the variable stack used
for the evaluation of the condition also contain pointers to these
entries.  Then, the assigned procedure is evaluated, which sets the
subevent in the parent node.  If required, the procedure
evaluates the condition node once for each (pair of) particles to
determine the result.
\nwenddocs{}\nwbegincode{373}\moddef{Eval trees: procedures}\plusendmoddef
  recursive subroutine eval_node_evaluate (en)
    type(eval_node_t), intent(inout) :: en
    logical :: exist
    select case (en%type)
    case (EN_UNARY)
       if (associated (en%arg1)) then
          call eval_node_evaluate (en%arg1)
          en%value_is_known = en%arg1%value_is_known
       else
          en%value_is_known = .false.
       end if
       if (en%value_is_known) then
          select case (en%result_type)
          case (V_LOG);  en%lval = en% op1_log  (en%arg1)
          case (V_INT);  en%ival = en% op1_int  (en%arg1)
          case (V_REAL); en%rval = en% op1_real (en%arg1)
          case (V_CMPLX); en%cval = en% op1_cmplx (en%arg1)          
          case (V_PDG);  
             call en% op1_pdg  (en%aval, en%arg1)
          case (V_SEV)
             if (associated (en%arg0)) then
                call en% op1_sev (en%pval, en%arg1, en%arg0)
             else
                call en% op1_sev (en%pval, en%arg1)
             end if
          case (V_STR)
             call en% op1_str (en%sval, en%arg1)
          end select
       end if
    case (EN_BINARY)
       if (associated (en%arg1) .and. associated (en%arg2)) then
          call eval_node_evaluate (en%arg1)
          call eval_node_evaluate (en%arg2)
          en%value_is_known = &
               en%arg1%value_is_known .and. en%arg2%value_is_known
       else
          en%value_is_known = .false.
       end if
       if (en%value_is_known) then
          select case (en%result_type)
          case (V_LOG);  en%lval = en% op2_log  (en%arg1, en%arg2)
          case (V_INT);  en%ival = en% op2_int  (en%arg1, en%arg2)
          case (V_REAL); en%rval = en% op2_real (en%arg1, en%arg2)
          case (V_CMPLX); en%cval = en% op2_cmplx (en%arg1, en%arg2)          
          case (V_PDG)
             call en% op2_pdg  (en%aval, en%arg1, en%arg2)
          case (V_SEV)
             if (associated (en%arg0)) then
                call en% op2_sev (en%pval, en%arg1, en%arg2, en%arg0)
             else
                call en% op2_sev (en%pval, en%arg1, en%arg2)
             end if
          case (V_STR)
             call en% op2_str (en%sval, en%arg1, en%arg2)
          end select
       end if
    case (EN_BLOCK)
       if (associated (en%arg1) .and. associated (en%arg0)) then
          call eval_node_evaluate (en%arg1)
          call eval_node_evaluate (en%arg0)
          en%value_is_known = en%arg0%value_is_known
       else
          en%value_is_known = .false.
       end if
       if (en%value_is_known) then
          select case (en%result_type)
          case (V_LOG);  en%lval = en%arg0%lval
          case (V_INT);  en%ival = en%arg0%ival
          case (V_REAL); en%rval = en%arg0%rval
          case (V_CMPLX); en%cval = en%arg0%cval          
          case (V_PDG);  en%aval = en%arg0%aval
          case (V_SEV);  en%pval = en%arg0%pval
          case (V_STR);  en%sval = en%arg0%sval
          end select
       end if
    case (EN_CONDITIONAL)
       if (associated (en%arg0)) then
          call eval_node_evaluate (en%arg0)
          en%value_is_known = en%arg0%value_is_known
       else
          en%value_is_known = .false.
       end if
       if (en%arg0%value_is_known) then
          if (en%arg0%lval) then
             call eval_node_evaluate (en%arg1)
             en%value_is_known = en%arg1%value_is_known
             if (en%value_is_known) then
                select case (en%result_type)
                case (V_LOG);  en%lval = en%arg1%lval
                case (V_INT);  en%ival = en%arg1%ival
                case (V_REAL); en%rval = en%arg1%rval
                case (V_CMPLX); en%cval = en%arg1%cval                
                case (V_PDG);  en%aval = en%arg1%aval
                case (V_SEV);  en%pval = en%arg1%pval
                case (V_STR);  en%sval = en%arg1%sval
                end select
             end if
          else
             call eval_node_evaluate (en%arg2)
             en%value_is_known = en%arg2%value_is_known
             if (en%value_is_known) then
                select case (en%result_type)
                case (V_LOG);  en%lval = en%arg2%lval
                case (V_INT);  en%ival = en%arg2%ival
                case (V_REAL); en%rval = en%arg2%rval
                case (V_CMPLX); en%cval = en%arg2%cval                
                case (V_PDG);  en%aval = en%arg2%aval
                case (V_SEV);  en%pval = en%arg2%pval
                case (V_STR);  en%sval = en%arg2%sval
                end select
             end if
          end if
       end if
    case (EN_RECORD_CMD)
       exist = .true.
       en%lval = .false.
       call eval_node_evaluate (en%arg0)
       if (en%arg0%value_is_known) then
          if (associated (en%arg1)) then
             call eval_node_evaluate (en%arg1)
             if (en%arg1%value_is_known) then
                if (associated (en%arg2)) then
                   call eval_node_evaluate (en%arg2)
                   if (en%arg2%value_is_known) then
                      if (associated (en%arg3)) then
                         call eval_node_evaluate (en%arg3)
                         if (en%arg3%value_is_known) then
                            if (associated (en%arg4)) then
                               call eval_node_evaluate (en%arg4)
                               if (en%arg4%value_is_known) then
                                  if (associated (en%rval)) then
                                     call analysis_record_data (en%arg0%sval, &
                                          en%arg1%rval, en%arg2%rval, &
                                          en%arg3%rval, en%arg4%rval, &
                                          weight=en%rval, exist=exist, &
                                          success=en%lval)
                                  else
                                     call analysis_record_data (en%arg0%sval, &
                                          en%arg1%rval, en%arg2%rval, &
                                          en%arg3%rval, en%arg4%rval, &
                                          exist=exist, success=en%lval)
                                  end if
                               end if
                            else                                   
                               if (associated (en%rval)) then
                                  call analysis_record_data (en%arg0%sval, &
                                       en%arg1%rval, en%arg2%rval, &
                                       en%arg3%rval, &
                                       weight=en%rval, exist=exist, &
                                       success=en%lval)
                               else
                                  call analysis_record_data (en%arg0%sval, &
                                       en%arg1%rval, en%arg2%rval, &
                                       en%arg3%rval, &
                                       exist=exist, success=en%lval)
                               end if
                            end if
                         end if
                      else
                         if (associated (en%rval)) then
                            call analysis_record_data (en%arg0%sval, &
                                 en%arg1%rval, en%arg2%rval, &
                                 weight=en%rval, exist=exist, &
                                 success=en%lval)
                         else
                            call analysis_record_data (en%arg0%sval, &
                                 en%arg1%rval, en%arg2%rval, &
                                 exist=exist, success=en%lval)
                         end if
                      end if
                   end if
                else
                   if (associated (en%rval)) then
                      call analysis_record_data (en%arg0%sval, &
                           en%arg1%rval, &
                           weight=en%rval, exist=exist, success=en%lval)
                   else
                      call analysis_record_data (en%arg0%sval, &
                           en%arg1%rval, &
                           exist=exist, success=en%lval)
                   end if
                end if
             end if
          else
             if (associated (en%rval)) then
                call analysis_record_data (en%arg0%sval, 1._default, &
                     weight=en%rval, exist=exist, success=en%lval)
             else
                call analysis_record_data (en%arg0%sval, 1._default, &
                     exist=exist, success=en%lval)
             end if
          end if
          if (.not. exist) then
             call msg_error ("Analysis object '" // char (en%arg0%sval) &
                  // "' is undefined")
             en%arg0%value_is_known = .false.
          end if
       end if
    case (EN_OBS1_INT)
       en%ival = en% obs1_int (en%prt1)
       en%value_is_known = .true.
    case (EN_OBS2_INT)
       en%ival = en% obs2_int (en%prt1, en%prt2)
       en%value_is_known = .true.
    case (EN_OBS1_REAL)
       en%rval = en% obs1_real (en%prt1)
       en%value_is_known = .true.
    case (EN_OBS2_REAL)
       en%rval = en% obs2_real (en%prt1, en%prt2)
       en%value_is_known = .true.
    case (EN_UOBS1_INT)
       en%ival = user_obs_int_p (en%arg0, en%prt1)
       en%value_is_known = .true.
    case (EN_UOBS2_INT)
       en%ival = user_obs_int_pp (en%arg0, en%prt1, en%prt2)
       en%value_is_known = .true.
    case (EN_UOBS1_REAL)
       en%rval = user_obs_real_p (en%arg0, en%prt1)
       en%value_is_known = .true.
    case (EN_UOBS2_REAL)
       en%rval = user_obs_real_pp (en%arg0, en%prt1, en%prt2)
       en%value_is_known = .true.
    case (EN_PRT_FUN_UNARY)
       call eval_node_evaluate (en%arg1)
       en%value_is_known = en%arg1%value_is_known
       if (en%value_is_known) then
          if (associated (en%arg0)) then
             en%arg0%index => en%index
             en%arg0%prt1 => en%prt1
             call en% op1_sev (en%pval, en%arg1, en%arg0)
          else
             call en% op1_sev (en%pval, en%arg1)
          end if
       end if
    case (EN_PRT_FUN_BINARY)
       call eval_node_evaluate (en%arg1)
       call eval_node_evaluate (en%arg2)
       en%value_is_known = &
            en%arg1%value_is_known .and. en%arg2%value_is_known
       if (en%value_is_known) then
          if (associated (en%arg0)) then
             en%arg0%index => en%index
             en%arg0%prt1 => en%prt1
             en%arg0%prt2 => en%prt2
             call en% op2_sev (en%pval, en%arg1, en%arg2, en%arg0)
          else
             call en% op2_sev (en%pval, en%arg1, en%arg2)
          end if
       end if
    case (EN_EVAL_FUN_UNARY)
       call eval_node_evaluate (en%arg1)
       en%value_is_known = subevt_is_nonempty (en%arg1%pval)
       if (en%value_is_known) then
          en%arg0%index => en%index
          en%index = 1
          en%arg0%prt1 => en%prt1
          en%prt1 = subevt_get_prt (en%arg1%pval, 1)
          call eval_node_evaluate (en%arg0)
          en%rval = en%arg0%rval
       end if
    case (EN_EVAL_FUN_BINARY)
       call eval_node_evaluate (en%arg1)
       call eval_node_evaluate (en%arg2)
       en%value_is_known = &
            subevt_is_nonempty (en%arg1%pval) .and. &
            subevt_is_nonempty (en%arg2%pval)
       if (en%value_is_known) then
          en%arg0%index => en%index
          en%arg0%prt1 => en%prt1
          en%arg0%prt2 => en%prt2
          en%index = 1
          call eval_pp (en%arg1, en%arg2, en%arg0, en%rval, en%value_is_known)
       end if
    case (EN_LOG_FUN_UNARY)
       call eval_node_evaluate (en%arg1)
       en%value_is_known = .true.
       if (en%value_is_known) then
          en%arg0%index => en%index
          en%arg0%prt1 => en%prt1
          en%lval = en% op1_cut (en%arg1, en%arg0)
       end if
    case (EN_LOG_FUN_BINARY)
       call eval_node_evaluate (en%arg1)
       call eval_node_evaluate (en%arg2)
       en%value_is_known = .true.
       if (en%value_is_known) then
          en%arg0%index => en%index
          en%arg0%prt1 => en%prt1
          en%arg0%prt2 => en%prt2
          en%lval = en% op2_cut (en%arg1, en%arg2, en%arg0)
       end if
    case (EN_INT_FUN_UNARY)
       call eval_node_evaluate (en%arg1)
       en%value_is_known = en%arg1%value_is_known
       if (en%value_is_known) then
          if (associated (en%arg0)) then
             en%arg0%index => en%index
             en%arg0%prt1 => en%prt1
             call en% op1_evi (en%ival, en%arg1, en%arg0)
          else
             call en% op1_evi (en%ival, en%arg1)
          end if
       end if
    case (EN_INT_FUN_BINARY)
       call eval_node_evaluate (en%arg1)
       call eval_node_evaluate (en%arg2)
       en%value_is_known = &
            en%arg1%value_is_known .and. &
            en%arg2%value_is_known
       if (en%value_is_known) then
          if (associated (en%arg0)) then
             en%arg0%index => en%index
             en%arg0%prt1 => en%prt1
             en%arg0%prt2 => en%prt2
             call en% op2_evi (en%ival, en%arg1, en%arg2, en%arg0)
          else
             call en% op2_evi (en%ival, en%arg1, en%arg2) 
          end if
       end if
    case (EN_REAL_FUN_UNARY)
       call eval_node_evaluate (en%arg1)
       en%value_is_known = en%arg1%value_is_known
       if (en%value_is_known) then
          if (associated (en%arg0)) then
             en%arg0%index => en%index
             en%arg0%prt1 => en%prt1
             call en% op1_evr (en%rval, en%arg1, en%arg0)
          else
             call en% op1_evr (en%rval, en%arg1)
          end if
       end if
    case (EN_REAL_FUN_BINARY)
       call eval_node_evaluate (en%arg1)
       call eval_node_evaluate (en%arg2)
       en%value_is_known = &
            en%arg1%value_is_known .and. &
            en%arg2%value_is_known
       if (en%value_is_known) then
          if (associated (en%arg0)) then
             en%arg0%index => en%index
             en%arg0%prt1 => en%prt1
             en%arg0%prt2 => en%prt2
             call en% op2_evr (en%rval, en%arg1, en%arg2, en%arg0)
          else
             call en% op2_evr (en%rval, en%arg1, en%arg2) 
          end if
       end if
    case (EN_FORMAT_STR)
       if (associated (en%arg0)) then
          call eval_node_evaluate (en%arg0)
          en%value_is_known = en%arg0%value_is_known
       else
          en%value_is_known = .true.
       end if
       if (associated (en%arg1)) then
          call eval_node_evaluate (en%arg1)
          en%value_is_known = &
               en%value_is_known .and. en%arg1%value_is_known
          if (en%value_is_known) then
             call evaluate_sprintf (en%sval, en%ival, en%arg0, en%arg1)
          end if
       else
          if (en%value_is_known) then
             call evaluate_sprintf (en%sval, en%ival, en%arg0)
          end if
       end if
    end select
    if (debug) then
       print *, "evaluated"
       call eval_node_write (en)
    end if
  end subroutine eval_node_evaluate

\eatline
\nwendcode{}\nwbegindocs{374}\nwdocspar
\subsection{Evaluation syntax}
We have two different flavors of the syntax: with and without particles.
\nwenddocs{}\nwbegincode{375}\moddef{Eval trees: variables}\plusendmoddef
  type(syntax_t), target, save :: syntax_expr
  type(syntax_t), target, save :: syntax_pexpr

\eatline
\nwendcode{}\nwbegindocs{376}These are for testing only and may be removed:
\nwenddocs{}\nwbegincode{377}\moddef{Eval trees: public}\endmoddef
  public :: syntax_expr_init
  public :: syntax_pexpr_init
\nwendcode{}\nwbegincode{378}\moddef{Eval trees: procedures}\plusendmoddef
  subroutine syntax_expr_init ()
    type(ifile_t) :: ifile
    call define_expr_syntax (ifile, particles=.false., analysis=.false.)
    call syntax_init (syntax_expr, ifile)
    call ifile_final (ifile)
  end subroutine syntax_expr_init

  subroutine syntax_pexpr_init ()
    type(ifile_t) :: ifile
    call define_expr_syntax (ifile, particles=.true., analysis=.false.)
    call syntax_init (syntax_pexpr, ifile)
    call ifile_final (ifile)
  end subroutine syntax_pexpr_init

\eatline
\nwendcode{}\nwbegincode{379}\moddef{Eval trees: public}\plusendmoddef
  public :: syntax_expr_final
  public :: syntax_pexpr_final
\nwendcode{}\nwbegincode{380}\moddef{Eval trees: procedures}\plusendmoddef
  subroutine syntax_expr_final ()
    call syntax_final (syntax_expr)
  end subroutine syntax_expr_final

  subroutine syntax_pexpr_final ()
    call syntax_final (syntax_pexpr)
  end subroutine syntax_pexpr_final

\eatline
\nwendcode{}\nwbegincode{381}\moddef{Eval trees: public}\plusendmoddef
  public :: syntax_pexpr_write
\nwendcode{}\nwbegincode{382}\moddef{Eval trees: procedures}\plusendmoddef
  subroutine syntax_pexpr_write (unit)
    integer, intent(in), optional :: unit
    call syntax_write (syntax_pexpr, unit)
  end subroutine syntax_pexpr_write

\eatline
\nwendcode{}\nwbegincode{383}\moddef{Eval trees: public}\plusendmoddef
  public :: define_expr_syntax
\nwendcode{}\nwbegindocs{384}Numeric expressions.
\nwenddocs{}\nwbegincode{385}\moddef{Eval trees: procedures}\plusendmoddef
  subroutine define_expr_syntax (ifile, particles, analysis)
    type(ifile_t), intent(inout) :: ifile
    logical, intent(in) :: particles, analysis
    type(string_t) :: numeric_pexpr
    type(string_t) :: var_plist, var_alias
    if (particles) then
       numeric_pexpr = " | numeric_pexpr"
       var_plist = " | var_plist"
       var_alias = " | var_alias"
    else
       numeric_pexpr = ""
       var_plist = ""
       var_alias = ""
    end if
    call ifile_append (ifile, "SEQ expr = subexpr addition*")
    call ifile_append (ifile, "ALT subexpr = addition | term")
    call ifile_append (ifile, "SEQ addition = plus_or_minus term")
    call ifile_append (ifile, "SEQ term = factor multiplication*")
    call ifile_append (ifile, "SEQ multiplication = times_or_over factor")
    call ifile_append (ifile, "SEQ factor = value exponentiation?")
    call ifile_append (ifile, "SEQ exponentiation = to_the value")
    call ifile_append (ifile, "ALT plus_or_minus = '+' | '-'")
    call ifile_append (ifile, "ALT times_or_over = '*' | '/'")
    call ifile_append (ifile, "ALT to_the = '^' | '**'")
    call ifile_append (ifile, "KEY '+'")
    call ifile_append (ifile, "KEY '-'")
    call ifile_append (ifile, "KEY '*'")
    call ifile_append (ifile, "KEY '/'")
    call ifile_append (ifile, "KEY '^'")
    call ifile_append (ifile, "KEY '**'")
    call ifile_append (ifile, "ALT value = signed_value | unsigned_value")
    call ifile_append (ifile, "SEQ signed_value = '-' unsigned_value")
    call ifile_append (ifile, "ALT unsigned_value = " // &
         "numeric_value | constant | variable | " // &
         "result | user_observable | " // &
         "grouped_expr | block_expr | conditional_expr | " // &
         "unary_function | binary_function" // &
         numeric_pexpr)
    call ifile_append (ifile, "ALT numeric_value = integer_value | " &
         // "real_value | complex_value")
    call ifile_append (ifile, "SEQ integer_value = integer_literal unit_expr?")
    call ifile_append (ifile, "SEQ real_value = real_literal unit_expr?")
    call ifile_append (ifile, "SEQ complex_value = complex_literal unit_expr?")    
    call ifile_append (ifile, "INT integer_literal")
    call ifile_append (ifile, "REA real_literal")
    call ifile_append (ifile, "COM complex_literal")
    call ifile_append (ifile, "SEQ unit_expr = unit unit_power?")
    call ifile_append (ifile, "ALT unit = " // &
         "TeV | GeV | MeV | keV | eV | meV | " // &
         "nbarn | pbarn | fbarn | abarn | " // &
         "rad | mrad | degree | '%'")
    call ifile_append (ifile, "KEY TeV")
    call ifile_append (ifile, "KEY GeV")
    call ifile_append (ifile, "KEY MeV")
    call ifile_append (ifile, "KEY keV")
    call ifile_append (ifile, "KEY eV")
    call ifile_append (ifile, "KEY meV")
    call ifile_append (ifile, "KEY nbarn")
    call ifile_append (ifile, "KEY pbarn")
    call ifile_append (ifile, "KEY fbarn")
    call ifile_append (ifile, "KEY abarn")
    call ifile_append (ifile, "KEY rad")
    call ifile_append (ifile, "KEY mrad")
    call ifile_append (ifile, "KEY degree")
    call ifile_append (ifile, "KEY '%'")
    call ifile_append (ifile, "SEQ unit_power = '^' frac_expr")
    call ifile_append (ifile, "ALT frac_expr = frac | grouped_frac")
    call ifile_append (ifile, "GRO grouped_frac = ( frac_expr )")
    call ifile_append (ifile, "SEQ frac = signed_int div?")
    call ifile_append (ifile, "ALT signed_int = " &
         // "neg_int | pos_int | integer_literal")
    call ifile_append (ifile, "SEQ neg_int = '-' integer_literal")
    call ifile_append (ifile, "SEQ pos_int = '+' integer_literal")
    call ifile_append (ifile, "SEQ div = '/' integer_literal")
    call ifile_append (ifile, "ALT constant = pi | I")
    call ifile_append (ifile, "KEY pi")
    call ifile_append (ifile, "KEY I")
    call ifile_append (ifile, "IDE variable")
    call ifile_append (ifile, "SEQ result = result_key result_arg")
    call ifile_append (ifile, "ALT result_key = " // &
         "num_id | integral | error")
    call ifile_append (ifile, "SEQ user_observable = user_obs user_arg")
    call ifile_append (ifile, "KEY user_obs")
    call ifile_append (ifile, "ARG user_arg = ( sexpr )")
    call ifile_append (ifile, "KEY num_id")
    call ifile_append (ifile, "KEY integral")
    call ifile_append (ifile, "KEY error")
    call ifile_append (ifile, "GRO result_arg = ( process_id )")
    call ifile_append (ifile, "IDE process_id")
    call ifile_append (ifile, "SEQ unary_function = fun_unary function_arg1")
    call ifile_append (ifile, "SEQ binary_function = fun_binary function_arg2")
    call ifile_append (ifile, "ALT fun_unary = " // &
         "complex | real | int | nint | floor | ceiling | abs | sgn | " // &
         "sqrt | exp | log | log10 | " // &
         "sin | cos | tan | asin | acos | atan | " // &
         "sinh | cosh | tanh")
    call ifile_append (ifile, "KEY complex")
    call ifile_append (ifile, "KEY real")
    call ifile_append (ifile, "KEY int")
    call ifile_append (ifile, "KEY nint")
    call ifile_append (ifile, "KEY floor")
    call ifile_append (ifile, "KEY ceiling")
    call ifile_append (ifile, "KEY abs")
    call ifile_append (ifile, "KEY sgn")
    call ifile_append (ifile, "KEY sqrt")
    call ifile_append (ifile, "KEY exp")
    call ifile_append (ifile, "KEY log")
    call ifile_append (ifile, "KEY log10")
    call ifile_append (ifile, "KEY sin")
    call ifile_append (ifile, "KEY cos")
    call ifile_append (ifile, "KEY tan")
    call ifile_append (ifile, "KEY asin")
    call ifile_append (ifile, "KEY acos")
    call ifile_append (ifile, "KEY atan")
    call ifile_append (ifile, "KEY sinh")
    call ifile_append (ifile, "KEY cosh")
    call ifile_append (ifile, "KEY tanh")
    call ifile_append (ifile, "ALT fun_binary = max | min | mod | modulo")
    call ifile_append (ifile, "KEY max")
    call ifile_append (ifile, "KEY min")
    call ifile_append (ifile, "KEY mod")
    call ifile_append (ifile, "KEY modulo")
    call ifile_append (ifile, "ARG function_arg1 = ( expr )")
    call ifile_append (ifile, "ARG function_arg2 = ( expr, expr )")
    call ifile_append (ifile, "GRO grouped_expr = ( expr )")
    call ifile_append (ifile, "SEQ block_expr = let var_spec in expr")
    call ifile_append (ifile, "KEY let")
    call ifile_append (ifile, "ALT var_spec = " // &
         "var_num | var_int | var_real | var_complex | " // &
         "var_logical" // var_plist // var_alias // " | var_string")
    call ifile_append (ifile, "SEQ var_num = var_name '=' expr")
    call ifile_append (ifile, "SEQ var_int = int var_name '=' expr")
    call ifile_append (ifile, "SEQ var_real = real var_name '=' expr")
    call ifile_append (ifile, "SEQ var_complex = complex var_name '=' complex_expr")
    call ifile_append (ifile, "ALT complex_expr = " // &
         "cexpr_real | cexpr_complex")
    call ifile_append (ifile, "ARG cexpr_complex = ( expr, expr )")
    call ifile_append (ifile, "SEQ cexpr_real = expr")
    call ifile_append (ifile, "IDE var_name")
    call ifile_append (ifile, "KEY '='")
    call ifile_append (ifile, "KEY in")
    call ifile_append (ifile, "SEQ conditional_expr = " // &
         "if lexpr then expr maybe_elsif_expr maybe_else_expr endif")
    call ifile_append (ifile, "SEQ maybe_elsif_expr = elsif_expr*")
    call ifile_append (ifile, "SEQ maybe_else_expr = else_expr?")
    call ifile_append (ifile, "SEQ elsif_expr = elsif lexpr then expr")
    call ifile_append (ifile, "SEQ else_expr = else expr")
    call ifile_append (ifile, "KEY if")
    call ifile_append (ifile, "KEY then")
    call ifile_append (ifile, "KEY elsif")
    call ifile_append (ifile, "KEY else")
    call ifile_append (ifile, "KEY endif")
    call define_lexpr_syntax (ifile, particles, analysis)
    call define_sexpr_syntax (ifile)
    if (particles) then
       call define_pexpr_syntax (ifile)
       call define_cexpr_syntax (ifile)
       call define_var_plist_syntax (ifile)
       call define_var_alias_syntax (ifile)
       call define_numeric_pexpr_syntax (ifile)
       call define_logical_pexpr_syntax (ifile)
    end if

  end subroutine define_expr_syntax

\eatline
\nwendcode{}\nwbegindocs{386}Logical expressions.
\nwenddocs{}\nwbegincode{387}\moddef{Eval trees: procedures}\plusendmoddef
  subroutine define_lexpr_syntax (ifile, particles, analysis)
    type(ifile_t), intent(inout) :: ifile
    logical, intent(in) :: particles, analysis
    type(string_t) :: logical_pexpr, record_cmd
    if (particles) then
       logical_pexpr = " | logical_pexpr"
    else
       logical_pexpr = ""
    end if
    if (analysis) then
       record_cmd = " | record_cmd"
    else
       record_cmd = ""
    end if
    call ifile_append (ifile, "SEQ lexpr = lsinglet lsequel*")
    call ifile_append (ifile, "SEQ lsequel = ';' lsinglet")
    call ifile_append (ifile, "SEQ lsinglet = lterm alternative*")
    call ifile_append (ifile, "SEQ alternative = or lterm")
    call ifile_append (ifile, "SEQ lterm = lvalue coincidence*")
    call ifile_append (ifile, "SEQ coincidence = and lvalue")
    call ifile_append (ifile, "KEY ';'")
    call ifile_append (ifile, "KEY or")
    call ifile_append (ifile, "KEY and")
    call ifile_append (ifile, "ALT lvalue = " // &
         "true | false | lvariable | negation | " // &
         "grouped_lexpr | block_lexpr | conditional_lexpr | " // &
         "compared_expr | compared_sexpr" // &
         logical_pexpr //  record_cmd)
    call ifile_append (ifile, "KEY true")
    call ifile_append (ifile, "KEY false")
    call ifile_append (ifile, "SEQ lvariable = '?' alt_lvariable")
    call ifile_append (ifile, "KEY '?'")
    call ifile_append (ifile, "ALT alt_lvariable = variable | grouped_lexpr")
    call ifile_append (ifile, "SEQ negation = not lvalue")
    call ifile_append (ifile, "KEY not")
    call ifile_append (ifile, "GRO grouped_lexpr = ( lexpr )")
    call ifile_append (ifile, "SEQ block_lexpr = let var_spec in lexpr")
    call ifile_append (ifile, "ALT var_logical = " // &
         "var_logical_new | var_logical_spec")
    call ifile_append (ifile, "SEQ var_logical_new = logical var_logical_spec")
    call ifile_append (ifile, "KEY logical")
    call ifile_append (ifile, "SEQ var_logical_spec = '?' var_name = lexpr")
    call ifile_append (ifile, "SEQ conditional_lexpr = " // &
         "if lexpr then lexpr maybe_elsif_lexpr maybe_else_lexpr endif")
    call ifile_append (ifile, "SEQ maybe_elsif_lexpr = elsif_lexpr*")
    call ifile_append (ifile, "SEQ maybe_else_lexpr = else_lexpr?")
    call ifile_append (ifile, "SEQ elsif_lexpr = elsif lexpr then lexpr")
    call ifile_append (ifile, "SEQ else_lexpr = else lexpr")
    call ifile_append (ifile, "SEQ compared_expr = expr comparison+")
    call ifile_append (ifile, "SEQ comparison = compare expr")
    call ifile_append (ifile, "ALT compare = " // &
         "'<' | '>' | '<=' | '>=' | '==' | '<>'")
    call ifile_append (ifile, "KEY '<'")
    call ifile_append (ifile, "KEY '>'")
    call ifile_append (ifile, "KEY '<='")
    call ifile_append (ifile, "KEY '>='")
    call ifile_append (ifile, "KEY '=='")
    call ifile_append (ifile, "KEY '<>'")
    call ifile_append (ifile, "SEQ compared_sexpr = sexpr str_comparison+")
    call ifile_append (ifile, "SEQ str_comparison = str_compare sexpr")
    call ifile_append (ifile, "ALT str_compare = '==' | '<>'")
    if (analysis) then
       call ifile_append (ifile, "SEQ record_cmd = " // &
            "record_key analysis_tag record_arg?")
       call ifile_append (ifile, "ALT record_key = " // &
            "record | record_unweighted | record_excess")
       call ifile_append (ifile, "KEY record")
       call ifile_append (ifile, "KEY record_unweighted")
       call ifile_append (ifile, "KEY record_excess")
       call ifile_append (ifile, "ALT analysis_tag = analysis_id | sexpr")
       call ifile_append (ifile, "IDE analysis_id")
       call ifile_append (ifile, "ARG record_arg = ( expr+ )")
    end if
  end subroutine define_lexpr_syntax

\eatline
\nwendcode{}\nwbegindocs{388}String expressions.
\nwenddocs{}\nwbegincode{389}\moddef{Eval trees: procedures}\plusendmoddef
  subroutine define_sexpr_syntax (ifile)
    type(ifile_t), intent(inout) :: ifile
    call ifile_append (ifile, "SEQ sexpr = svalue str_concatenation*")
    call ifile_append (ifile, "SEQ str_concatenation = '&' svalue")
    call ifile_append (ifile, "KEY '&'")
    call ifile_append (ifile, "ALT svalue = " // &
         "grouped_sexpr | block_sexpr | conditional_sexpr | " // &
         "svariable | string_function | string_literal")
    call ifile_append (ifile, "GRO grouped_sexpr = ( sexpr )")
    call ifile_append (ifile, "SEQ block_sexpr = let var_spec in sexpr")
    call ifile_append (ifile, "SEQ conditional_sexpr = " // &
         "if lexpr then sexpr maybe_elsif_sexpr maybe_else_sexpr endif")
    call ifile_append (ifile, "SEQ maybe_elsif_sexpr = elsif_sexpr*")
    call ifile_append (ifile, "SEQ maybe_else_sexpr = else_sexpr?")
    call ifile_append (ifile, "SEQ elsif_sexpr = elsif lexpr then sexpr")
    call ifile_append (ifile, "SEQ else_sexpr = else sexpr")
    call ifile_append (ifile, "SEQ svariable = '$' alt_svariable")
    call ifile_append (ifile, "KEY '$'")
    call ifile_append (ifile, "ALT alt_svariable = variable | grouped_sexpr")
    call ifile_append (ifile, "ALT var_string = " // &
         "var_string_new | var_string_spec")
    call ifile_append (ifile, "SEQ var_string_new = string var_string_spec")
    call ifile_append (ifile, "KEY string")
    call ifile_append (ifile, "SEQ var_string_spec = '$' var_name = sexpr") ! $
    call ifile_append (ifile, "ALT string_function = sprintf_fun")
    call ifile_append (ifile, "SEQ sprintf_fun = sprintf_clause sprintf_args?")
    call ifile_append (ifile, "SEQ sprintf_clause = sprintf sexpr")
    call ifile_append (ifile, "KEY sprintf")
    call ifile_append (ifile, "ARG sprintf_args = ( sprintf_arg* )")
    call ifile_append (ifile, "ALT sprintf_arg = " &
         // "lvariable | svariable | expr")
    call ifile_append (ifile, "QUO string_literal = '""'...'""'")
  end subroutine define_sexpr_syntax

\eatline
\nwendcode{}\nwbegindocs{390}Eval trees that evaluate to subevents.
\nwenddocs{}\nwbegincode{391}\moddef{Eval trees: procedures}\plusendmoddef
  subroutine define_pexpr_syntax (ifile)
    type(ifile_t), intent(inout) :: ifile
    call ifile_append (ifile, "SEQ pexpr = pterm pconcatenation*")
    call ifile_append (ifile, "SEQ pconcatenation = '&' pterm")
    ! call ifile_append (ifile, "KEY '&'")   !!! (Key exists already)
    call ifile_append (ifile, "SEQ pterm = pvalue pcombination*")
    call ifile_append (ifile, "SEQ pcombination = '+' pvalue")
    ! call ifile_append (ifile, "KEY '+'")   !!! (Key exists already)
    call ifile_append (ifile, "ALT pvalue = " // &
         "pexpr_src | pvariable | " // &
         "grouped_pexpr | block_pexpr | conditional_pexpr | " // &
         "prt_function")
    call ifile_append (ifile, "SEQ pexpr_src = prefix_cexpr")
    call ifile_append (ifile, "ALT prefix_cexpr = " // &
         "incoming_prt | outgoing_prt")
    call ifile_append (ifile, "SEQ incoming_prt = incoming cexpr")
    call ifile_append (ifile, "KEY incoming")
    call ifile_append (ifile, "SEQ outgoing_prt = cexpr")
    call ifile_append (ifile, "SEQ pvariable = '@' alt_pvariable")
    call ifile_append (ifile, "KEY '@'")
    call ifile_append (ifile, "ALT alt_pvariable = variable | grouped_pexpr")
    call ifile_append (ifile, "GRO grouped_pexpr = '[' pexpr ']'")
    call ifile_append (ifile, "SEQ block_pexpr = let var_spec in pexpr")
    call ifile_append (ifile, "SEQ conditional_pexpr = " // &
         "if lexpr then pexpr maybe_elsif_pexpr maybe_else_pexpr endif")
    call ifile_append (ifile, "SEQ maybe_elsif_pexpr = elsif_pexpr*")
    call ifile_append (ifile, "SEQ maybe_else_pexpr = else_pexpr?")
    call ifile_append (ifile, "SEQ elsif_pexpr = elsif lexpr then pexpr")
    call ifile_append (ifile, "SEQ else_pexpr = else pexpr")
    call ifile_append (ifile, "ALT prt_function = " // &
         "join_fun | combine_fun | collect_fun | cluster_fun | " // &
         "select_fun | extract_fun | sort_fun")
    call ifile_append (ifile, "SEQ join_fun = join_clause pargs2")
    call ifile_append (ifile, "SEQ combine_fun = combine_clause pargs2")
    call ifile_append (ifile, "SEQ collect_fun = collect_clause pargs1")
    call ifile_append (ifile, "SEQ cluster_fun = cluster_clause pargs1")
    call ifile_append (ifile, "SEQ select_fun = select_clause pargs1")
    call ifile_append (ifile, "SEQ extract_fun = extract_clause pargs1")
    call ifile_append (ifile, "SEQ sort_fun = sort_clause pargs1")
    call ifile_append (ifile, "SEQ join_clause = join condition?")
    call ifile_append (ifile, "SEQ combine_clause = combine condition?")
    call ifile_append (ifile, "SEQ collect_clause = collect condition?")
    call ifile_append (ifile, "SEQ cluster_clause = cluster condition?")
    call ifile_append (ifile, "SEQ select_clause = select condition?")
    call ifile_append (ifile, "SEQ extract_clause = extract position?")
    call ifile_append (ifile, "SEQ sort_clause = sort criterion?")
    call ifile_append (ifile, "KEY join")
    call ifile_append (ifile, "KEY combine")
    call ifile_append (ifile, "KEY collect")
    call ifile_append (ifile, "KEY cluster")
    call ifile_append (ifile, "KEY select")
    call ifile_append (ifile, "SEQ condition = if lexpr")
    call ifile_append (ifile, "KEY extract")
    call ifile_append (ifile, "SEQ position = index expr")
    call ifile_append (ifile, "KEY sort")
    call ifile_append (ifile, "SEQ criterion = by expr")
    call ifile_append (ifile, "KEY index")
    call ifile_append (ifile, "KEY by")
    call ifile_append (ifile, "ARG pargs2 = '[' pexpr, pexpr ']'")
    call ifile_append (ifile, "ARG pargs1 = '[' pexpr, pexpr? ']'")
  end subroutine define_pexpr_syntax

\eatline
\nwendcode{}\nwbegindocs{392}Eval trees that evaluate to PDG-code arrays.
\nwenddocs{}\nwbegincode{393}\moddef{Eval trees: procedures}\plusendmoddef
  subroutine define_cexpr_syntax (ifile)
    type(ifile_t), intent(inout) :: ifile
    call ifile_append (ifile, "SEQ cexpr = avalue concatenation*")
    call ifile_append (ifile, "SEQ concatenation = ':' avalue")
    call ifile_append (ifile, "KEY ':'")
    call ifile_append (ifile, "ALT avalue = " // &
         "grouped_cexpr | block_cexpr | conditional_cexpr | " // &
         "variable | pdg_code | prt_name")
    call ifile_append (ifile, "GRO grouped_cexpr = ( cexpr )")
    call ifile_append (ifile, "SEQ block_cexpr = let var_spec in cexpr")
    call ifile_append (ifile, "SEQ conditional_cexpr = " // &
         "if lexpr then cexpr maybe_elsif_cexpr maybe_else_cexpr endif")
    call ifile_append (ifile, "SEQ maybe_elsif_cexpr = elsif_cexpr*")
    call ifile_append (ifile, "SEQ maybe_else_cexpr = else_cexpr?")
    call ifile_append (ifile, "SEQ elsif_cexpr = elsif lexpr then cexpr")
    call ifile_append (ifile, "SEQ else_cexpr = else cexpr")
    call ifile_append (ifile, "SEQ pdg_code = pdg pdg_arg")
    call ifile_append (ifile, "KEY pdg")
    call ifile_append (ifile, "ARG pdg_arg = ( expr )")
    call ifile_append (ifile, "QUO prt_name = '""'...'""'")
  end subroutine define_cexpr_syntax

\eatline
\nwendcode{}\nwbegindocs{394}Extra variable types.
\nwenddocs{}\nwbegincode{395}\moddef{Eval trees: procedures}\plusendmoddef
  subroutine define_var_plist_syntax (ifile)
    type(ifile_t), intent(inout) :: ifile
    call ifile_append (ifile, "ALT var_plist = var_plist_new | var_plist_spec")
    call ifile_append (ifile, "SEQ var_plist_new = subevt var_plist_spec")
    call ifile_append (ifile, "KEY subevt")
    call ifile_append (ifile, "SEQ var_plist_spec = '@' var_name '=' pexpr")
  end subroutine define_var_plist_syntax

  subroutine define_var_alias_syntax (ifile)
    type(ifile_t), intent(inout) :: ifile
    call ifile_append (ifile, "SEQ var_alias = alias var_name '=' cexpr")
    call ifile_append (ifile, "KEY alias")
  end subroutine define_var_alias_syntax

\eatline
\nwendcode{}\nwbegindocs{396}Particle-list expressions that evaluate to numeric values
\nwenddocs{}\nwbegincode{397}\moddef{Eval trees: procedures}\plusendmoddef
  subroutine define_numeric_pexpr_syntax (ifile)
    type(ifile_t), intent(inout) :: ifile
    call ifile_append (ifile, "ALT numeric_pexpr = " &
         // "eval_fun | count_fun | event_shape_fun")
    call ifile_append (ifile, "SEQ eval_fun = eval expr pargs1")
    call ifile_append (ifile, "SEQ count_fun = count_clause pargs1")
    call ifile_append (ifile, "SEQ count_clause = count condition?")
    call ifile_append (ifile, "KEY eval")
    call ifile_append (ifile, "KEY count")
    call ifile_append (ifile, "ALT event_shape_fun = user_event_fun")
    call ifile_append (ifile, "SEQ user_event_fun = " &
         // "user_event_shape user_arg pargs1")
    call ifile_append (ifile, "KEY user_event_shape")
  end subroutine define_numeric_pexpr_syntax

\eatline
\nwendcode{}\nwbegindocs{398}Particle-list functions that evaluate to logical values.
\nwenddocs{}\nwbegincode{399}\moddef{Eval trees: procedures}\plusendmoddef
  subroutine define_logical_pexpr_syntax (ifile)
    type(ifile_t), intent(inout) :: ifile
    call ifile_append (ifile, "ALT logical_pexpr = " // &
         "all_fun | any_fun | no_fun | user_cut_fun")
    call ifile_append (ifile, "SEQ all_fun = all lexpr pargs1")
    call ifile_append (ifile, "SEQ any_fun = any lexpr pargs1")
    call ifile_append (ifile, "SEQ no_fun = no lexpr pargs1")
    call ifile_append (ifile, "KEY all")
    call ifile_append (ifile, "KEY any")
    call ifile_append (ifile, "KEY no")
    call ifile_append (ifile, "SEQ user_cut_fun = user_cut user_arg pargs1")
    call ifile_append (ifile, "KEY user_cut")
  end subroutine define_logical_pexpr_syntax

\eatline
\nwendcode{}\nwbegindocs{400}All characters that can occur in expressions (apart from alphanumeric).
\nwenddocs{}\nwbegincode{401}\moddef{Eval trees: procedures}\plusendmoddef
  subroutine lexer_init_eval_tree (lexer, particles)
    type(lexer_t), intent(out) :: lexer
    logical, intent(in) :: particles
    type(keyword_list_t), pointer :: keyword_list
    if (particles) then
       keyword_list => syntax_get_keyword_list_ptr (syntax_pexpr)
    else
       keyword_list => syntax_get_keyword_list_ptr (syntax_expr)
    end if
    call lexer_init (lexer, &
         comment_chars = "#!", &
         quote_chars = '"', &
         quote_match = '"', &
         single_chars = "()[],;:&%?$@", &
         special_class = [ "+-*/^", "<>=~ " ] , &
         keyword_list = keyword_list)
  end subroutine lexer_init_eval_tree

\eatline
\nwendcode{}\nwbegindocs{402}\nwdocspar
\subsection{Set up appropriate parse trees}
Parse an input stream as a specific flavor of expression.  The
appropriate expression syntax has to be available.
\nwenddocs{}\nwbegincode{403}\moddef{Eval trees: public}\plusendmoddef
  public :: parse_tree_init_expr
  public :: parse_tree_init_lexpr
  public :: parse_tree_init_pexpr
  public :: parse_tree_init_cexpr
  public :: parse_tree_init_sexpr
\nwendcode{}\nwbegincode{404}\moddef{Eval trees: procedures}\plusendmoddef
  subroutine parse_tree_init_expr (parse_tree, stream, particles)
    type(parse_tree_t), intent(out) :: parse_tree
    type(stream_t), intent(inout), target :: stream
    logical, intent(in) :: particles
    type(lexer_t) :: lexer
    call lexer_init_eval_tree (lexer, particles)
    call lexer_assign_stream (lexer, stream)
    if (particles) then
       call parse_tree_init &
            (parse_tree, syntax_pexpr, lexer, var_str ("expr"))
    else
       call parse_tree_init &
            (parse_tree, syntax_expr, lexer, var_str ("expr"))
    end if
    call lexer_final (lexer)
  end subroutine parse_tree_init_expr

  subroutine parse_tree_init_lexpr (parse_tree, stream, particles)
    type(parse_tree_t), intent(out) :: parse_tree
    type(stream_t), intent(inout), target :: stream
    logical, intent(in) :: particles
    type(lexer_t) :: lexer
    call lexer_init_eval_tree (lexer, particles)
    call lexer_assign_stream (lexer, stream)
    if (particles) then
       call parse_tree_init &
            (parse_tree, syntax_pexpr, lexer, var_str ("lexpr"))
    else
       call parse_tree_init &
            (parse_tree, syntax_expr, lexer, var_str ("lexpr"))
    end if
    call lexer_final (lexer)
  end subroutine parse_tree_init_lexpr

  subroutine parse_tree_init_pexpr (parse_tree, stream)
    type(parse_tree_t), intent(out) :: parse_tree
    type(stream_t), intent(inout), target :: stream
    type(lexer_t) :: lexer
    call lexer_init_eval_tree (lexer, .true.)
    call lexer_assign_stream (lexer, stream)
    call parse_tree_init &
         (parse_tree, syntax_pexpr, lexer, var_str ("pexpr"))
    call lexer_final (lexer)
  end subroutine parse_tree_init_pexpr

  subroutine parse_tree_init_cexpr (parse_tree, stream)
    type(parse_tree_t), intent(out) :: parse_tree
    type(stream_t), intent(inout), target :: stream
    type(lexer_t) :: lexer
    call lexer_init_eval_tree (lexer, .true.)
    call lexer_assign_stream (lexer, stream)
    call parse_tree_init &
         (parse_tree, syntax_pexpr, lexer, var_str ("cexpr"))
    call lexer_final (lexer)
  end subroutine parse_tree_init_cexpr

  subroutine parse_tree_init_sexpr (parse_tree, stream, particles)
    type(parse_tree_t), intent(out) :: parse_tree
    type(stream_t), intent(inout), target :: stream
    logical, intent(in) :: particles
    type(lexer_t) :: lexer
    call lexer_init_eval_tree (lexer, particles)
    call lexer_assign_stream (lexer, stream)
    if (particles) then
       call parse_tree_init &
            (parse_tree, syntax_pexpr, lexer, var_str ("sexpr"))
    else
       call parse_tree_init &
            (parse_tree, syntax_expr, lexer, var_str ("sexpr"))
    end if
    call lexer_final (lexer)
  end subroutine parse_tree_init_sexpr

\eatline
\eatline
\eatline
\eatline
\eatline
\nwendcode{}\nwbegindocs{405}\nwdocspar
\subsection{The evaluation tree}
The evaluation tree contains the initial variable list and the root node.
\nwenddocs{}\nwbegincode{406}\moddef{Eval trees: public}\plusendmoddef
  public :: eval_tree_t
\nwendcode{}\nwbegincode{407}\moddef{Eval trees: types}\plusendmoddef
  type, extends (expr_t) :: eval_tree_t
     private
     type(parse_node_t), pointer :: pn => null ()
     type(var_list_t) :: var_list
     type(eval_node_t), pointer :: root => null ()
   contains
   \LA{}Eval trees: eval tree: TBP\RA{}
  end type eval_tree_t

\eatline
\nwendcode{}\nwbegindocs{408}Init from stream, using a temporary parse tree.
\nwenddocs{}\nwbegincode{409}\moddef{Eval trees: procedures}\plusendmoddef
  subroutine eval_tree_init_stream &
       (eval_tree, stream, var_list, subevt, result_type)
    type(eval_tree_t), intent(out), target :: eval_tree
    type(stream_t), intent(inout), target :: stream
    type(var_list_t), intent(in), target :: var_list
    type(subevt_t), intent(in), target, optional :: subevt
    integer, intent(in), optional :: result_type
    type(parse_tree_t) :: parse_tree
    type(parse_node_t), pointer :: nd_root
    integer :: type
    type = V_REAL;  if (present (result_type))  type = result_type
    select case (type)
    case (V_INT, V_REAL, V_CMPLX)
       call parse_tree_init_expr (parse_tree, stream, present (subevt))
    case (V_LOG)
       call parse_tree_init_lexpr (parse_tree, stream, present (subevt))
    case (V_SEV)
       call parse_tree_init_pexpr (parse_tree, stream)
    case (V_PDG)
       call parse_tree_init_cexpr (parse_tree, stream)
    case (V_STR)
       call parse_tree_init_sexpr (parse_tree, stream, present (subevt))
    end select
    nd_root => parse_tree_get_root_ptr (parse_tree)
    if (associated (nd_root)) then
       select case (type)
       case (V_INT, V_REAL, V_CMPLX)
          call eval_tree_init_expr (eval_tree, nd_root, var_list, subevt)
       case (V_LOG)
          call eval_tree_init_lexpr (eval_tree, nd_root, var_list, subevt)
       case (V_SEV)
          call eval_tree_init_pexpr (eval_tree, nd_root, var_list, subevt)
       case (V_PDG)
          call eval_tree_init_cexpr (eval_tree, nd_root, var_list, subevt)
       case (V_STR)
          call eval_tree_init_sexpr (eval_tree, nd_root, var_list, subevt)
       end select
    end if
    call parse_tree_final (parse_tree)
  end subroutine eval_tree_init_stream

\eatline
\nwendcode{}\nwbegindocs{410}API (to be superseded by the methods below): Init from a given parse-tree
node.  If we evaluate an expression that contains particle-list references,
the original subevent has to be supplied.  The initial variable list is
optional.
\nwenddocs{}\nwbegincode{411}\moddef{Eval trees: eval tree: TBP}\endmoddef
  procedure :: init_expr  => eval_tree_init_expr
  procedure :: init_lexpr => eval_tree_init_lexpr
  procedure :: init_pexpr => eval_tree_init_pexpr
  procedure :: init_cexpr => eval_tree_init_cexpr
  procedure :: init_sexpr => eval_tree_init_sexpr
\nwendcode{}\nwbegincode{412}\moddef{Eval trees: procedures}\plusendmoddef
  subroutine eval_tree_init_expr &
      (expr, parse_node, var_list, subevt)
    class(eval_tree_t), intent(out), target :: expr
    type(parse_node_t), intent(in), target :: parse_node
    type(var_list_t), intent(in), target :: var_list
    type(subevt_t), intent(in), optional, target :: subevt
    call eval_tree_link_var_list (expr, var_list)
    if (present (subevt))  call eval_tree_set_subevt (expr, subevt)
    call eval_node_compile_expr &
         (expr%root, parse_node, expr%var_list)
  end subroutine eval_tree_init_expr
    
  subroutine eval_tree_init_lexpr &
      (expr, parse_node, var_list, subevt)
    class(eval_tree_t), intent(out), target :: expr
    type(parse_node_t), intent(in), target :: parse_node
    type(var_list_t), intent(in), target :: var_list
    type(subevt_t), intent(in), optional, target :: subevt
    call eval_tree_link_var_list (expr, var_list)
    if (present (subevt))  call eval_tree_set_subevt (expr, subevt)
    call eval_node_compile_lexpr &
         (expr%root, parse_node, expr%var_list)
  end subroutine eval_tree_init_lexpr

  subroutine eval_tree_init_pexpr &
      (expr, parse_node, var_list, subevt)
    class(eval_tree_t), intent(out), target :: expr
    type(parse_node_t), intent(in), target :: parse_node
    type(var_list_t), intent(in), target :: var_list
    type(subevt_t), intent(in), optional, target :: subevt
    call eval_tree_link_var_list (expr, var_list)
    if (present (subevt))  call eval_tree_set_subevt (expr, subevt)
    call eval_node_compile_pexpr &
         (expr%root, parse_node, expr%var_list)
  end subroutine eval_tree_init_pexpr

  subroutine eval_tree_init_cexpr &
      (expr, parse_node, var_list, subevt)
    class(eval_tree_t), intent(out), target :: expr
    type(parse_node_t), intent(in), target :: parse_node
    type(var_list_t), intent(in), target :: var_list
    type(subevt_t), intent(in), optional, target :: subevt
    call eval_tree_link_var_list (expr, var_list)
    if (present (subevt))  call eval_tree_set_subevt (expr, subevt)
    call eval_node_compile_cexpr &
         (expr%root, parse_node, expr%var_list)
  end subroutine eval_tree_init_cexpr

  subroutine eval_tree_init_sexpr &
      (expr, parse_node, var_list, subevt)
    class(eval_tree_t), intent(out), target :: expr
    type(parse_node_t), intent(in), target :: parse_node
    type(var_list_t), intent(in), target :: var_list
    type(subevt_t), intent(in), optional, target :: subevt
    call eval_tree_link_var_list (expr, var_list)
    if (present (subevt))  call eval_tree_set_subevt (expr, subevt)
    call eval_node_compile_sexpr &
         (expr%root, parse_node, expr%var_list)
  end subroutine eval_tree_init_sexpr

\eatline
\eatline
\eatline
\eatline
\eatline
\nwendcode{}\nwbegindocs{413}Alternative: set up the expression using the parse node that has already
been stored.  We assume that the {\tt{}subevt} or any other variable that
may be referred to has already been added to the local variable list.
\nwenddocs{}\nwbegincode{414}\moddef{Eval trees: eval tree: TBP}\plusendmoddef
  procedure :: setup_expr  => eval_tree_setup_expr
  procedure :: setup_lexpr => eval_tree_setup_lexpr
  procedure :: setup_pexpr => eval_tree_setup_pexpr
  procedure :: setup_cexpr => eval_tree_setup_cexpr
  procedure :: setup_sexpr => eval_tree_setup_sexpr
\nwendcode{}\nwbegincode{415}\moddef{Eval trees: procedures}\plusendmoddef
  subroutine eval_tree_setup_expr (expr, vars)
    class(eval_tree_t), intent(inout), target :: expr
    class(vars_t), intent(in), target :: vars
    call eval_tree_link_var_list (expr, vars)
    call eval_node_compile_expr (expr%root, expr%pn, expr%var_list)
  end subroutine eval_tree_setup_expr
    
  subroutine eval_tree_setup_lexpr (expr, vars)
    class(eval_tree_t), intent(inout), target :: expr
    class(vars_t), intent(in), target :: vars
    call eval_tree_link_var_list (expr, vars)
    call eval_node_compile_lexpr (expr%root, expr%pn, expr%var_list)
  end subroutine eval_tree_setup_lexpr

  subroutine eval_tree_setup_pexpr (expr, vars)
    class(eval_tree_t), intent(inout), target :: expr
    class(vars_t), intent(in), target :: vars
    call eval_tree_link_var_list (expr, vars)
    call eval_node_compile_pexpr (expr%root, expr%pn, expr%var_list)
  end subroutine eval_tree_setup_pexpr

  subroutine eval_tree_setup_cexpr (expr, vars)
    class(eval_tree_t), intent(inout), target :: expr
    class(vars_t), intent(in), target :: vars
    call eval_tree_link_var_list (expr, vars)
    call eval_node_compile_cexpr (expr%root, expr%pn, expr%var_list)
  end subroutine eval_tree_setup_cexpr

  subroutine eval_tree_setup_sexpr (expr, vars)
    class(eval_tree_t), intent(inout), target :: expr
    class(vars_t), intent(in), target :: vars
    call eval_tree_link_var_list (expr, vars)
    call eval_node_compile_sexpr (expr%root, expr%pn, expr%var_list)
  end subroutine eval_tree_setup_sexpr

\eatline
\eatline
\eatline
\eatline
\eatline
\nwendcode{}\nwbegindocs{416}This extra API function handles numerical constant expressions only.
The only nontrivial part is the optional unit.
\nwenddocs{}\nwbegincode{417}\moddef{Eval trees: eval tree: TBP}\plusendmoddef
  procedure :: init_numeric_value => eval_tree_init_numeric_value
\nwendcode{}\nwbegincode{418}\moddef{Eval trees: procedures}\plusendmoddef
  subroutine eval_tree_init_numeric_value (eval_tree, parse_node)
    class(eval_tree_t), intent(out), target :: eval_tree
    type(parse_node_t), intent(in), target :: parse_node
    call eval_node_compile_numeric_value (eval_tree%root, parse_node)
  end subroutine eval_tree_init_numeric_value

\eatline
\nwendcode{}\nwbegindocs{419}Initialize the variable list, linking it to a context variable list.
\nwenddocs{}\nwbegincode{420}\moddef{Eval trees: procedures}\plusendmoddef
  subroutine eval_tree_link_var_list (eval_tree, vars)
    type(eval_tree_t), intent(inout), target :: eval_tree
    class(vars_t), intent(in), target :: vars
    call eval_tree%var_list%link (vars)
  end subroutine eval_tree_link_var_list

\eatline
\nwendcode{}\nwbegindocs{421}Include a subevent object in the initialization.  We add a pointer
to this as variable {\tt{}@evt} in the local variable list.
\nwenddocs{}\nwbegincode{422}\moddef{Eval trees: procedures}\plusendmoddef
  subroutine eval_tree_set_subevt (eval_tree, subevt)
    type(eval_tree_t), intent(inout), target :: eval_tree
    type(subevt_t), intent(in), target :: subevt
    logical, save, target :: known = .true.
    call var_list_append_subevt_ptr &
         (eval_tree%var_list, var_str ("@evt"), subevt, known, &
         intrinsic=.true.)
  end subroutine eval_tree_set_subevt

\eatline
\nwendcode{}\nwbegindocs{423}Finalizer.
\nwenddocs{}\nwbegincode{424}\moddef{Eval trees: eval tree: TBP}\plusendmoddef
  procedure :: final => eval_tree_final
\nwendcode{}\nwbegincode{425}\moddef{Eval trees: procedures}\plusendmoddef
  subroutine eval_tree_final (expr)
    class(eval_tree_t), intent(inout) :: expr
    call expr%var_list%final ()
    if (associated (expr%root)) then
       call eval_node_final_rec (expr%root)
       deallocate (expr%root)
    end if
  end subroutine eval_tree_final

\eatline
\nwendcode{}\nwbegindocs{426}\nwdocspar
\nwenddocs{}\nwbegincode{427}\moddef{Eval trees: eval tree: TBP}\plusendmoddef
  procedure :: evaluate => eval_tree_evaluate
\nwendcode{}\nwbegincode{428}\moddef{Eval trees: procedures}\plusendmoddef
  subroutine eval_tree_evaluate (expr)
    class(eval_tree_t), intent(inout) :: expr
    if (associated (expr%root)) then
       call eval_node_evaluate (expr%root)
    end if
  end subroutine eval_tree_evaluate

\eatline
\nwendcode{}\nwbegindocs{429}Check if the eval tree is allocated.
\nwenddocs{}\nwbegincode{430}\moddef{Eval trees: procedures}\plusendmoddef
  function eval_tree_is_defined (eval_tree) result (flag)
    logical :: flag
    type(eval_tree_t), intent(in) :: eval_tree
    flag = associated (eval_tree%root)
  end function eval_tree_is_defined

\eatline
\nwendcode{}\nwbegindocs{431}Check if the eval tree result is constant.
\nwenddocs{}\nwbegincode{432}\moddef{Eval trees: procedures}\plusendmoddef
  function eval_tree_is_constant (eval_tree) result (flag)
    logical :: flag
    type(eval_tree_t), intent(in) :: eval_tree
    if (associated (eval_tree%root)) then
       flag = eval_tree%root%type == EN_CONSTANT
    else
       flag = .false.
    end if
  end function eval_tree_is_constant

\eatline
\nwendcode{}\nwbegindocs{433}Insert a conversion node at the root, if necessary (only for
real/int conversion)
\nwenddocs{}\nwbegincode{434}\moddef{Eval trees: procedures}\plusendmoddef
  subroutine eval_tree_convert_result (eval_tree, result_type)
    type(eval_tree_t), intent(inout) :: eval_tree
    integer, intent(in) :: result_type
    if (associated (eval_tree%root)) then
       call insert_conversion_node (eval_tree%root, result_type)
    end if
  end subroutine eval_tree_convert_result

\eatline
\nwendcode{}\nwbegindocs{435}Return the value of the top node, after evaluation.  If the tree is
empty, return the type of {\tt{}V{\char95}NONE}.  When extracting the value, no
check for existence is done.  For numeric values, the functions are
safe against real/integer mismatch.
\nwenddocs{}\nwbegincode{436}\moddef{Eval trees: eval tree: TBP}\plusendmoddef
  procedure :: is_known => eval_tree_result_is_known
  procedure :: get_log => eval_tree_get_log
  procedure :: get_int => eval_tree_get_int
  procedure :: get_real => eval_tree_get_real
  procedure :: get_cmplx => eval_tree_get_cmplx
  procedure :: get_pdg_array => eval_tree_get_pdg_array
  procedure :: get_subevt => eval_tree_get_subevt
  procedure :: get_string => eval_tree_get_string
\nwendcode{}\nwbegincode{437}\moddef{Eval trees: procedures}\plusendmoddef
  function eval_tree_get_result_type (expr) result (type)
    integer :: type
    class(eval_tree_t), intent(in) :: expr
    if (associated (expr%root)) then
       type = expr%root%result_type
    else
       type = V_NONE
    end if
  end function eval_tree_get_result_type
    
  function eval_tree_result_is_known (expr) result (flag)
    logical :: flag
    class(eval_tree_t), intent(in) :: expr
    if (associated (expr%root)) then
       select case (expr%root%result_type)
       case (V_LOG, V_INT, V_REAL)
          flag = expr%root%value_is_known
       case default
          flag = .true.
       end select
    else
       flag = .false.
    end if
  end function eval_tree_result_is_known

  function eval_tree_result_is_known_ptr (expr) result (ptr)
    logical, pointer :: ptr
    class(eval_tree_t), intent(in) :: expr
    logical, target, save :: known = .true.
    if (associated (expr%root)) then
       select case (expr%root%result_type)
       case (V_LOG, V_INT, V_REAL)
          ptr => expr%root%value_is_known
       case default
          ptr => known
       end select
    else
       ptr => null ()
    end if
  end function eval_tree_result_is_known_ptr

  function eval_tree_get_log (expr) result (lval)
    logical :: lval
    class(eval_tree_t), intent(in) :: expr
    if (associated (expr%root))  lval = expr%root%lval
  end function eval_tree_get_log
    
  function eval_tree_get_int (expr) result (ival)
    integer :: ival
    class(eval_tree_t), intent(in) :: expr
    if (associated (expr%root)) then
       select case (expr%root%result_type)
       case (V_INT);  ival = expr%root%ival
       case (V_REAL); ival = expr%root%rval
       case (V_CMPLX); ival = expr%root%cval       
       end select
    end if
  end function eval_tree_get_int
    
  function eval_tree_get_real (expr) result (rval)
    real(default) :: rval
    class(eval_tree_t), intent(in) :: expr
    if (associated (expr%root)) then
       select case (expr%root%result_type)
       case (V_REAL); rval = expr%root%rval
       case (V_INT);  rval = expr%root%ival
       case (V_CMPLX);  rval = expr%root%cval       
       end select
    end if
  end function eval_tree_get_real
      
  function eval_tree_get_cmplx (expr) result (cval)
    complex(default) :: cval
    class(eval_tree_t), intent(in) :: expr
    if (associated (expr%root)) then
       select case (expr%root%result_type)
       case (V_CMPLX); cval = expr%root%cval
       case (V_REAL); cval = expr%root%rval
       case (V_INT);  cval = expr%root%ival
       end select
    end if
  end function eval_tree_get_cmplx

  function eval_tree_get_pdg_array (expr) result (aval)
    type(pdg_array_t) :: aval
    class(eval_tree_t), intent(in) :: expr
    if (associated (expr%root)) then
       aval = expr%root%aval
    end if
  end function eval_tree_get_pdg_array

  function eval_tree_get_subevt (expr) result (pval)
    type(subevt_t) :: pval
    class(eval_tree_t), intent(in) :: expr
    if (associated (expr%root)) then
       pval = expr%root%pval
    end if
  end function eval_tree_get_subevt

  function eval_tree_get_string (expr) result (sval)
    type(string_t) :: sval
    class(eval_tree_t), intent(in) :: expr
    if (associated (expr%root)) then
       sval = expr%root%sval
    end if
  end function eval_tree_get_string

\eatline
\eatline
\eatline
\eatline
\eatline
\eatline
\eatline
\eatline
\nwendcode{}\nwbegindocs{438}Return a pointer to the value of the top node.
\nwenddocs{}\nwbegincode{439}\moddef{Eval trees: procedures}\plusendmoddef
  function eval_tree_get_log_ptr (eval_tree) result (lval)
    logical, pointer :: lval
    type(eval_tree_t), intent(in) :: eval_tree
    if (associated (eval_tree%root)) then
       lval => eval_tree%root%lval
    else
       lval => null ()
    end if
  end function eval_tree_get_log_ptr
    
  function eval_tree_get_int_ptr (eval_tree) result (ival)
    integer, pointer :: ival
    type(eval_tree_t), intent(in) :: eval_tree
    if (associated (eval_tree%root)) then
       ival => eval_tree%root%ival
    else
       ival => null ()
    end if
  end function eval_tree_get_int_ptr
    
  function eval_tree_get_real_ptr (eval_tree) result (rval)
    real(default), pointer :: rval
    type(eval_tree_t), intent(in) :: eval_tree
    if (associated (eval_tree%root)) then
       rval => eval_tree%root%rval
    else
       rval => null ()
    end if
  end function eval_tree_get_real_ptr
    
  function eval_tree_get_cmplx_ptr (eval_tree) result (cval)
    complex(default), pointer :: cval
    type(eval_tree_t), intent(in) :: eval_tree
    if (associated (eval_tree%root)) then
       cval => eval_tree%root%cval
    else
       cval => null ()
    end if
  end function eval_tree_get_cmplx_ptr

  function eval_tree_get_subevt_ptr (eval_tree) result (pval)
    type(subevt_t), pointer :: pval
    type(eval_tree_t), intent(in) :: eval_tree
    if (associated (eval_tree%root)) then
       pval => eval_tree%root%pval
    else
       pval => null ()
    end if
  end function eval_tree_get_subevt_ptr
    
  function eval_tree_get_pdg_array_ptr (eval_tree) result (aval)
    type(pdg_array_t), pointer :: aval
    type(eval_tree_t), intent(in) :: eval_tree
    if (associated (eval_tree%root)) then
       aval => eval_tree%root%aval
    else
       aval => null ()
    end if
  end function eval_tree_get_pdg_array_ptr
    
  function eval_tree_get_string_ptr (eval_tree) result (sval)
    type(string_t), pointer :: sval
    type(eval_tree_t), intent(in) :: eval_tree
    if (associated (eval_tree%root)) then
       sval => eval_tree%root%sval
    else
       sval => null ()
    end if
  end function eval_tree_get_string_ptr
    
\eatline
\eatline
\eatline
\eatline
\nwendcode{}\nwbegincode{440}\moddef{Eval trees: eval tree: TBP}\plusendmoddef
  procedure :: write => eval_tree_write
\nwendcode{}\nwbegincode{441}\moddef{Eval trees: procedures}\plusendmoddef
  subroutine eval_tree_write (expr, unit, write_vars)
    class(eval_tree_t), intent(in) :: expr
    integer, intent(in), optional :: unit
    logical, intent(in), optional :: write_vars
    integer :: u
    logical :: vl
    u = given_output_unit (unit);  if (u < 0)  return
    vl = .false.;  if (present (write_vars))  vl = write_vars
    write (u, "(1x,A)") "Evaluation tree:"
    if (associated (expr%root)) then
       call eval_node_write_rec (expr%root, unit)
    else
       write (u, "(3x,A)") "[empty]"
    end if
    if (vl)  call var_list_write (expr%var_list, unit)
  end subroutine eval_tree_write

\eatline
\nwendcode{}\nwbegindocs{442}Use the written representation for generating an MD5 sum:
\nwenddocs{}\nwbegincode{443}\moddef{Eval trees: procedures}\plusendmoddef
  function eval_tree_get_md5sum (eval_tree) result (md5sum_et)
    character(32) :: md5sum_et
    type(eval_tree_t), intent(in) :: eval_tree
    integer :: u
    u = free_unit ()
    open (unit = u, status = "scratch", action = "readwrite")
    call eval_tree_write (eval_tree, unit=u)
    rewind (u)
    md5sum_et = md5sum (u)
    close (u)
  end function eval_tree_get_md5sum

\eatline
\nwendcode{}\nwbegindocs{444}\nwdocspar
\subsection{Direct evaluation}
These procedures create an eval tree and evaluate it on-the-fly, returning
only the final value.  The evaluation must yield a well-defined value, unless
the {\tt{}is{\char95}known} flag is present, which will be set accordingly.
\nwenddocs{}\nwbegincode{445}\moddef{Eval trees: public}\plusendmoddef
  public :: eval_log
  public :: eval_int
  public :: eval_real
  public :: eval_cmplx
  public :: eval_subevt
  public :: eval_pdg_array
  public :: eval_string
\nwendcode{}\nwbegincode{446}\moddef{Eval trees: procedures}\plusendmoddef
  function eval_log &
       (parse_node, var_list, subevt, is_known) result (lval)
    logical :: lval
    type(parse_node_t), intent(in), target :: parse_node
    type(var_list_t), intent(in), target :: var_list
    type(subevt_t), intent(in), optional, target :: subevt
    logical, intent(out), optional :: is_known
    type(eval_tree_t), target :: eval_tree
    call eval_tree_init_lexpr &
         (eval_tree, parse_node, var_list, subevt)
    call eval_tree_evaluate (eval_tree)
    if (eval_tree_result_is_known (eval_tree)) then
       if (present (is_known))  is_known = .true.
       lval = eval_tree_get_log (eval_tree)
    else if (present (is_known)) then
       is_known = .false.
    else
       call eval_tree_unknown (eval_tree, parse_node)
       lval = .false.
    end if
    call eval_tree_final (eval_tree)
  end function eval_log

  function eval_int &
       (parse_node, var_list, subevt, is_known) result (ival)
    integer :: ival
    type(parse_node_t), intent(in), target :: parse_node
    type(var_list_t), intent(in), target :: var_list
    type(subevt_t), intent(in), optional, target :: subevt
    logical, intent(out), optional :: is_known
    type(eval_tree_t), target :: eval_tree
    call eval_tree_init_expr &
         (eval_tree, parse_node, var_list, subevt)
    call eval_tree_evaluate (eval_tree)
    if (eval_tree_result_is_known (eval_tree)) then
       if (present (is_known))  is_known = .true.
       ival = eval_tree_get_int (eval_tree)
    else if (present (is_known)) then
       is_known = .false.
    else
       call eval_tree_unknown (eval_tree, parse_node)
       ival = 0
    end if
    call eval_tree_final (eval_tree)
  end function eval_int

  function eval_real &
       (parse_node, var_list, subevt, is_known) result (rval)
    real(default) :: rval
    type(parse_node_t), intent(in), target :: parse_node
    type(var_list_t), intent(in), target :: var_list
    type(subevt_t), intent(in), optional, target :: subevt
    logical, intent(out), optional :: is_known
    type(eval_tree_t), target :: eval_tree
    call eval_tree_init_expr &
         (eval_tree, parse_node, var_list, subevt)
    call eval_tree_evaluate (eval_tree)
    if (eval_tree_result_is_known (eval_tree)) then
       if (present (is_known))  is_known = .true.
       rval = eval_tree_get_real (eval_tree)
    else if (present (is_known)) then
       is_known = .false.
    else 
       call eval_tree_unknown (eval_tree, parse_node)
       rval = 0
    end if
    call eval_tree_final (eval_tree)
  end function eval_real

  function eval_cmplx &
       (parse_node, var_list, subevt, is_known) result (cval)
    complex(default) :: cval
    type(parse_node_t), intent(in), target :: parse_node
    type(var_list_t), intent(in), target :: var_list
    type(subevt_t), intent(in), optional, target :: subevt
    logical, intent(out), optional :: is_known
    type(eval_tree_t), target :: eval_tree
    call eval_tree_init_expr &
         (eval_tree, parse_node, var_list, subevt)
    call eval_tree_evaluate (eval_tree)
    if (eval_tree_result_is_known (eval_tree)) then
       if (present (is_known))  is_known = .true.
       cval = eval_tree_get_cmplx (eval_tree)
    else if (present (is_known)) then
       is_known = .false.
    else
       call eval_tree_unknown (eval_tree, parse_node)
       cval = 0
    end if
    call eval_tree_final (eval_tree)
  end function eval_cmplx

  function eval_subevt &
       (parse_node, var_list, subevt, is_known) result (pval)
    type(subevt_t) :: pval
    type(parse_node_t), intent(in), target :: parse_node
    type(var_list_t), intent(in), target :: var_list
    type(subevt_t), intent(in), optional, target :: subevt
    logical, intent(out), optional :: is_known
    type(eval_tree_t), target :: eval_tree
    call eval_tree_init_pexpr &
         (eval_tree, parse_node, var_list, subevt)
    call eval_tree_evaluate (eval_tree)
    if (eval_tree_result_is_known (eval_tree)) then
       if (present (is_known))  is_known = .true.
       pval = eval_tree_get_subevt (eval_tree)
    else if (present (is_known)) then
       is_known = .false.
    else
       call eval_tree_unknown (eval_tree, parse_node)
    end if
    call eval_tree_final (eval_tree)
  end function eval_subevt

  function eval_pdg_array &
       (parse_node, var_list, subevt, is_known) result (aval)
    type(pdg_array_t) :: aval
    type(parse_node_t), intent(in), target :: parse_node
    type(var_list_t), intent(in), target :: var_list
    type(subevt_t), intent(in), optional, target :: subevt
    logical, intent(out), optional :: is_known
    type(eval_tree_t), target :: eval_tree
    call eval_tree_init_cexpr &
         (eval_tree, parse_node, var_list, subevt)
    call eval_tree_evaluate (eval_tree)
    if (eval_tree_result_is_known (eval_tree)) then
       if (present (is_known))  is_known = .true.
       aval = eval_tree_get_pdg_array (eval_tree)
    else if (present (is_known)) then
       is_known = .false.
    else
       call eval_tree_unknown (eval_tree, parse_node)
    end if
    call eval_tree_final (eval_tree)
  end function eval_pdg_array

  function eval_string &
       (parse_node, var_list, subevt, is_known) result (sval)
    type(string_t) :: sval
    type(parse_node_t), intent(in), target :: parse_node
    type(var_list_t), intent(in), target :: var_list
    type(subevt_t), intent(in), optional, target :: subevt
    logical, intent(out), optional :: is_known
    type(eval_tree_t), target :: eval_tree
    call eval_tree_init_sexpr &
         (eval_tree, parse_node, var_list, subevt)
    call eval_tree_evaluate (eval_tree)
    if (eval_tree_result_is_known (eval_tree)) then
       if (present (is_known))  is_known = .true.
       sval = eval_tree_get_string (eval_tree)
    else if (present (is_known)) then
       is_known = .false.
    else
       call eval_tree_unknown (eval_tree, parse_node)
       sval = ""
    end if
    call eval_tree_final (eval_tree)
  end function eval_string

\eatline
\eatline
\eatline
\nwendcode{}\nwbegindocs{447}Here is a variant that returns numeric values of all possible kinds, the
appropriate kind to be selected later:
\nwenddocs{}\nwbegincode{448}\moddef{Eval trees: public}\plusendmoddef
  public :: eval_numeric
\nwendcode{}\nwbegincode{449}\moddef{Eval trees: procedures}\plusendmoddef
  subroutine eval_numeric &
       (parse_node, var_list, subevt, ival, rval, cval, &
        is_known, result_type)
    type(parse_node_t), intent(in), target :: parse_node
    type(var_list_t), intent(in), target :: var_list
    type(subevt_t), intent(in), optional, target :: subevt
    integer, intent(out), optional :: ival
    real(default), intent(out), optional :: rval
    complex(default), intent(out), optional :: cval
    logical, intent(out), optional :: is_known
    integer, intent(out), optional :: result_type
    type(eval_tree_t), target :: eval_tree
    call eval_tree_init_expr &
         (eval_tree, parse_node, var_list, subevt)
    call eval_tree_evaluate (eval_tree)
    if (eval_tree_result_is_known (eval_tree)) then
       if (present (ival))  ival = eval_tree_get_int (eval_tree)
       if (present (rval))  rval = eval_tree_get_real (eval_tree)
       if (present (cval))  cval = eval_tree_get_cmplx (eval_tree)
       if (present (is_known))  is_known = .true.
    else
       call eval_tree_unknown (eval_tree, parse_node)
       if (present (ival))  ival = 0
       if (present (rval))  rval = 0
       if (present (cval))  cval = 0
       if (present (is_known))  is_known = .false.
    end if
    if (present (result_type))  &
         result_type = eval_tree_get_result_type (eval_tree)
    call eval_tree_final (eval_tree)
  end subroutine eval_numeric

\eatline
\nwendcode{}\nwbegindocs{450}Error message with debugging info:
\nwenddocs{}\nwbegincode{451}\moddef{Eval trees: procedures}\plusendmoddef
  subroutine eval_tree_unknown (eval_tree, parse_node)
    type(eval_tree_t), intent(in) :: eval_tree
    type(parse_node_t), intent(in) :: parse_node
    call parse_node_write_rec (parse_node)
    call eval_tree_write (eval_tree)
    call msg_error ("Evaluation yields an undefined result, inserting default")
  end subroutine eval_tree_unknown

\eatline
\nwendcode{}\nwbegindocs{452}\nwdocspar
\subsection{Factory Type}
Since {\tt{}eval{\char95}tree{\char95}t} is an implementation of {\tt{}expr{\char95}t}, we also need a
matching factory type and build method.
\nwenddocs{}\nwbegincode{453}\moddef{Eval trees: public}\plusendmoddef
  public :: eval_tree_factory_t
\nwendcode{}\nwbegincode{454}\moddef{Eval trees: types}\plusendmoddef
  type, extends (expr_factory_t) :: eval_tree_factory_t
     private
     type(parse_node_t), pointer :: pn => null ()
   contains
  \LA{}Eval trees: eval tree factory: TBP\RA{}
  end type eval_tree_factory_t
  
\eatline
\nwendcode{}\nwbegindocs{455}Output: delegate to the output of the embedded parse node.
\nwenddocs{}\nwbegincode{456}\moddef{Eval trees: eval tree factory: TBP}\endmoddef
  procedure :: write => eval_tree_factory_write
\nwendcode{}\nwbegincode{457}\moddef{Eval trees: procedures}\plusendmoddef
  subroutine eval_tree_factory_write (expr_factory, unit)
    class(eval_tree_factory_t), intent(in) :: expr_factory
    integer, intent(in), optional :: unit
    if (associated (expr_factory%pn)) then
       call parse_node_write_rec (expr_factory%pn, unit)
    end if
  end subroutine eval_tree_factory_write
       
\eatline
\nwendcode{}\nwbegindocs{458}Initializer: take a parse node and hide it thus from the environment.
\nwenddocs{}\nwbegincode{459}\moddef{Eval trees: eval tree factory: TBP}\plusendmoddef
  procedure :: init => eval_tree_factory_init
\nwendcode{}\nwbegincode{460}\moddef{Eval trees: procedures}\plusendmoddef
  subroutine eval_tree_factory_init (expr_factory, pn)
    class(eval_tree_factory_t), intent(out) :: expr_factory
    type(parse_node_t), intent(in), pointer :: pn
    expr_factory%pn => pn
  end subroutine eval_tree_factory_init
  
\eatline
\nwendcode{}\nwbegindocs{461}Factory method: allocate expression with correct eval tree type.  If the
stored parse node is not associate, don't allocate.
\nwenddocs{}\nwbegincode{462}\moddef{Eval trees: eval tree factory: TBP}\plusendmoddef
  procedure :: build => eval_tree_factory_build
\nwendcode{}\nwbegincode{463}\moddef{Eval trees: procedures}\plusendmoddef
  subroutine eval_tree_factory_build (expr_factory, expr)
    class(eval_tree_factory_t), intent(in) :: expr_factory
    class(expr_t), intent(out), allocatable :: expr
    if (associated (expr_factory%pn)) then
       allocate (eval_tree_t :: expr)
       select type (expr)
       type is (eval_tree_t)
          expr%pn => expr_factory%pn
       end select
    end if
  end subroutine eval_tree_factory_build
  
\eatline
\nwendcode{}\nwbegindocs{464}\nwdocspar
\subsection{Unit tests}
\nwenddocs{}\nwbegincode{465}\moddef{Eval trees: public}\plusendmoddef
  public :: expressions_test
\nwendcode{}\nwbegincode{466}\moddef{Eval trees: procedures}\plusendmoddef
  subroutine expressions_test (u, results)
    integer, intent(in) :: u
    type (test_results_t), intent(inout) :: results
  \LA{}Eval trees: execute tests\RA{}  
  end subroutine expressions_test

\eatline
\nwendcode{}\nwbegindocs{467}Testing the routines of the expressions module. First a simple unary
observable and the node evaluation.
\nwenddocs{}\nwbegincode{468}\moddef{Eval trees: execute tests}\endmoddef
  call test (expressions_1, "expressions_1", &
       "check simple observable", &
       u, results) 
\nwendcode{}\nwbegincode{469}\moddef{Eval trees: tests}\endmoddef
  subroutine expressions_1 (u)
    integer, intent(in) :: u
    type(var_list_t), pointer :: var_list => null ()
    type(eval_node_t), pointer :: node => null ()
    type(prt_t), pointer :: prt => null ()
    type(string_t) :: var_name
    procedure(obs_unary_int), pointer :: obs1_iptr
    type(prt_t), pointer :: p1

    write (u, "(A)")  "* Test output: Expressions"
    write (u, "(A)")  "*   Purpose: test simple observable and node evaluation"
    write (u, "(A)")       

    write (u, "(A)")  "* Setting a unary observable:"
    write (u, "(A)")    
    
    allocate (var_list)
    allocate (prt)
    call var_list_set_observables_unary (var_list, prt)
    call var_list_write (var_list, u)
    
    write (u, "(A)")  "* Evaluating the observable node:"
    write (u, "(A)")        
    
    var_name = "PDG"

    allocate (node)
    call var_list%get_obs1_iptr (var_name, obs1_iptr, p1)
    call eval_node_init_obs1_int_ptr (node, var_name, obs1_iptr, p1)
    call eval_node_write (node, u)

    write (u, "(A)")  "* Cleanup"
    write (u, "(A)")        
    
    call eval_node_final_rec (node)
    deallocate (node)
    call var_list%final ()
    deallocate (var_list)
    deallocate (prt)
       
    write (u, "(A)")
    write (u, "(A)")  "* Test output end: expressions_1"        
    
  end subroutine expressions_1
  
\nwendcode{}\nwbegindocs{470} %def expressions_1
\nwenddocs{}\nwbegindocs{471}Parse a complicated expression, transfer it to a parse tree and evaluate.
\nwenddocs{}\nwbegincode{472}\moddef{Eval trees: execute tests}\plusendmoddef
  call test (expressions_2, "expressions_2", &
       "check expression transfer to parse tree", &
       u, results) 
\nwendcode{}\nwbegincode{473}\moddef{Eval trees: tests}\plusendmoddef
  subroutine expressions_2 (u)
    integer, intent(in) :: u
    type(ifile_t) :: ifile
    type(stream_t) :: stream
    type(eval_tree_t) :: eval_tree
    type(string_t) :: expr_text
    type(var_list_t), pointer :: var_list => null ()
    
    write (u, "(A)")  "* Test output: Expressions"
    write (u, "(A)")  "*   Purpose: test parse routines"
    write (u, "(A)")       
    
    call syntax_expr_init ()
    call syntax_write (syntax_expr, u)     
    allocate (var_list)
    call var_list_append_real (var_list, var_str ("tolerance"), 0._default)
    call var_list_append_real (var_list, var_str ("x"), -5._default)
    call var_list_append_int  (var_list, var_str ("foo"), -27)
    call var_list_append_real (var_list, var_str ("mb"), 4._default)
    expr_text = &
         "let real twopi = 2 * pi in" // &
         "  twopi * sqrt (25.d0 - mb^2)" // &
         "  / (let int mb_or_0 = max (mb, 0) in" // &
         "       1 + (if -1 TeV <= x < mb_or_0 then abs(x) else x endif))"
    call ifile_append (ifile, expr_text)
    call stream_init (stream, ifile)
    call var_list_write (var_list, unit = u)
    call eval_tree_init_stream (eval_tree, stream, var_list=var_list)
    call eval_tree_evaluate (eval_tree)
    call eval_tree_write (eval_tree, u)
    
    write (u, "(A)")  "* Input string:"
    write (u, "(A,A)")  "     ", char (expr_text)
    write (u, "(A)") 
    write (u, "(A)")  "* Cleanup"
    
    call stream_final (stream) 
    call ifile_final (ifile)
    call eval_tree_final (eval_tree)
    call var_list%final ()
    deallocate (var_list)
    call syntax_expr_final ()     
    
    write (u, "(A)")
    write (u, "(A)")  "* Test output end: expressions_2"        
    
  end subroutine expressions_2 
  
\nwendcode{}\nwbegindocs{474} %def expressions_2
\nwenddocs{}\nwbegindocs{475}Test a subevent expression.
\nwenddocs{}\nwbegincode{476}\moddef{Eval trees: execute tests}\plusendmoddef
  call test (expressions_3, "expressions_3", &
       "check subevent expressions", &
       u, results) 
\nwendcode{}\nwbegincode{477}\moddef{Eval trees: tests}\plusendmoddef
  subroutine expressions_3 (u)
    integer, intent(in) :: u
    type(subevt_t) :: subevt

    write (u, "(A)")  "* Test output: Expressions"
    write (u, "(A)")  "*   Purpose: test subevent expressions"
    write (u, "(A)")       

    write (u, "(A)")  "* Initialize subevent:"
    write (u, "(A)")       
            
    call subevt_init (subevt)
    call subevt_reset (subevt, 1)
    call subevt_set_incoming (subevt, 1, &
         22, vector4_moving (1.e3_default, 1.e3_default, 1), &
         0._default, [2])
    call subevt_write (subevt, u)
    call subevt_reset (subevt, 4)
    call subevt_reset (subevt, 3)
    call subevt_set_incoming (subevt, 1, &
         21, vector4_moving (1.e3_default, 1.e3_default, 3), &
         0._default, [1])
    call subevt_polarize (subevt, 1, -1)
    call subevt_set_outgoing (subevt, 2, &
         1, vector4_moving (0._default, 1.e3_default, 3), &
         -1.e6_default, [7])
    call subevt_set_composite (subevt, 3, &
         vector4_moving (-1.e3_default, 0._default, 3), &
         [2, 7])
    call subevt_write (subevt, u)
    
    write (u, "(A)")
    write (u, "(A)")  "* Test output end: expressions_3"        
    
  end subroutine expressions_3
  
\nwendcode{}\nwbegindocs{478} %def expressions_3
\nwenddocs{}\nwbegindocs{479}Test expressions from a PDG array.
\nwenddocs{}\nwbegincode{480}\moddef{Eval trees: execute tests}\plusendmoddef
  call test (expressions_4, "expressions_4", &
       "check pdg array expressions", &
       u, results) 
\nwendcode{}\nwbegincode{481}\moddef{Eval trees: tests}\plusendmoddef
  subroutine expressions_4 (u)
    integer, intent(in) :: u
    type(subevt_t), target :: subevt
    type(string_t) :: expr_text
    type(ifile_t) :: ifile
    type(stream_t) :: stream
    type(eval_tree_t) :: eval_tree
    type(var_list_t), pointer :: var_list => null ()
    type(pdg_array_t) :: aval

    write (u, "(A)")  "* Test output: Expressions"
    write (u, "(A)")  "*   Purpose: test pdg array expressions"
    write (u, "(A)")       

    write (u, "(A)")  "* Initialization:"
    write (u, "(A)")       
    
    call syntax_pexpr_init ()
    call syntax_write (syntax_pexpr, u)
    allocate (var_list)
    call var_list_append_real (var_list, var_str ("tolerance"), 0._default)
    aval = 0
    call var_list_append_pdg_array (var_list, var_str ("particle"), aval)
    aval = [11,-11]
    call var_list_append_pdg_array (var_list, var_str ("lepton"), aval)
    aval = 22
    call var_list_append_pdg_array (var_list, var_str ("photon"), aval)
    aval = 1
    call var_list_append_pdg_array (var_list, var_str ("u"), aval)
    call subevt_init (subevt)
    call subevt_reset (subevt, 6)
    call subevt_set_incoming (subevt, 1, &
         1, vector4_moving (1._default, 1._default, 1), 0._default)
    call subevt_set_incoming (subevt, 2, &
         -1, vector4_moving (2._default, 2._default, 1), 0._default)
    call subevt_set_outgoing (subevt, 3, &
         22, vector4_moving (3._default, 3._default, 1), 0._default)
    call subevt_set_outgoing (subevt, 4, &
         22, vector4_moving (4._default, 4._default, 1), 0._default)
    call subevt_set_outgoing (subevt, 5, &
         11, vector4_moving (5._default, 5._default, 1), 0._default)
    call subevt_set_outgoing (subevt, 6, &
         -11, vector4_moving (6._default, 6._default, 1), 0._default)
    write (u, "(A)")
    write (u, "(A)")  "* Expression:"
    expr_text = &
         "let alias quark = pdg(1):pdg(2):pdg(3) in" // &
         "  any E > 3 GeV " // &
         "    [sort by - Pt " // &
         "       [select if Index < 6 " // &
         "          [photon:pdg(-11):pdg(3):quark " // &
         "           & incoming particle]]]" // &
         "  and" // &
         "  eval Theta [extract index -1 [photon]] > 45 degree" // &
         "  and" // &
         "  count [incoming photon] * 3 > 0"
    write (u, "(A,A)")  "     ", char (expr_text)
    write (u, "(A)")
    
    write (u, "(A)")
    write (u, "(A)")  "* Extract the evaluation tree:"
    write (u, "(A)")
    
    call ifile_append (ifile, expr_text)
    call stream_init (stream, ifile)
    call eval_tree_init_stream (eval_tree, stream, var_list, subevt, V_LOG)
    call eval_tree_write (eval_tree, u)
    call eval_tree_evaluate (eval_tree)

    write (u, "(A)")
    write (u, "(A)")  "* Evaluate the tree:"
    write (u, "(A)")

    call eval_tree_write (eval_tree, u)

    write (u, "(A)")
    write (u, "(A)")  "* Cleanup"
    write (u, "(A)")
    
    call stream_final (stream)
    call ifile_final (ifile)
    call eval_tree_final (eval_tree)
    call var_list%final ()
    deallocate (var_list)
    call syntax_pexpr_final ()  
    
    write (u, "(A)")
    write (u, "(A)")  "* Test output end: expressions_4"        
        
  end subroutine expressions_4

\eatline
\nwendcode{}\nwbegindocs{482}\nwdocspar
\clearpage
%%%%%%%%%%%%%%%%%%%%%%%%%%%%%%%%%%%%%%%%%%%%%%%%%%%%%%%%%%%%%%%%%%%%%%%%
\section{Physics Models}

A model object represents a physics model.  It contains a table of particle
data, a list of parameters, and a vertex table.  The list of parameters is a
variable list which includes the real parameters (which are pointers to the
particle data table) and PDG array variables for the particles themselves.
The vertex list is used for phase-space generation, not for calculating the
matrix element.

The actual numeric model data are in the base type {\tt{}model{\char95}data{\char95}t},
as part of the {\tt{}qft} section.  We implement the {\tt{}model{\char95}t} as an
extension of this, for convenient direct access to the base-type
methods via inheritance.  (Alternatively, we could delegate these calls
explicitly.)  The extension contains administrative additions, such as
the methods for recalculating derived data and keeping the parameter
set consistent.  It thus acts as a proxy of the actual model-data
object towards the \whizard\ package.  There are further proxy
objects, such as the {\tt{}parameter{\char95}t} array which provides the
interface to the actual numeric parameters.

Model definitions are read from model files.  Therefore, this module contains
a parser for model files.  The parameter definitions (derived parameters) are
Sindarin expressions.

The models, as read from file, are stored in a model library which is a simple
list of model definitions.  For setting up a process object we should make a
copy (an instance) of a model, which gets the current parameter values from
the global variable list.

\subsection{Module}
\nwenddocs{}\nwbegincode{483}\moddef{\code{}models.f90\edoc{}}\endmoddef
\LA{}File header\RA{}

module models

  use, intrinsic :: iso_c_binding !NODEP!
  
\LA{}Use kinds\RA{}
  use kinds, only: i8, i32
  use kinds, only: c_default_float
\LA{}Use strings\RA{}
  use format_defs, only: FMT_19
  use io_units
  use unit_tests
  use diagnostics
  use md5
  use os_interface
  use physics_defs, only: UNDEFINED, SCALAR, SPINOR
  use model_data
  
  use ifiles
  use syntax_rules
  use lexers
  use parser
  use pdg_arrays
  use variables
  use expr_base
  use eval_trees

\LA{}Standard module head\RA{}

\LA{}Models: public\RA{}

\LA{}Models: parameters\RA{}

\LA{}Models: types\RA{}

\LA{}Models: interfaces\RA{}

\LA{}Models: variables\RA{}

contains

\LA{}Models: procedures\RA{}

\LA{}Models: tests\RA{}

end module models
\eatline
\nwendcode{}\nwbegindocs{484}\nwdocspar
\subsection{Physics Parameters}
A parameter has a name, a value.  Derived parameters also have a
definition in terms of other parameters, which is stored as an
{\tt{}eval{\char95}tree}.  External parameters are set by an external program.

This parameter object should be considered as a proxy object.  The
parameter name and value are stored in a corresponding
{\tt{}modelpar{\char95}data{\char95}t} object which is located in a {\tt{}model{\char95}data{\char95}t}
object.  The latter is a component of the {\tt{}model{\char95}t} handler.
Methods of {\tt{}parameter{\char95}t} can be delegated to the {\tt{}par{\char95}data{\char95}t}
component.

The {\tt{}pn} component is a pointer to the parameter definition inside the
model parse tree.  It allows us to recreate the {\tt{}eval{\char95}tree} when making
copies (instances) of the parameter object.
\nwenddocs{}\nwbegincode{485}\moddef{Models: parameters}\endmoddef
  integer, parameter :: PAR_NONE = 0
  integer, parameter :: PAR_INDEPENDENT = 1, PAR_DERIVED = 2
  integer, parameter :: PAR_EXTERNAL = 3

\eatline
\nwendcode{}\nwbegincode{486}\moddef{Models: types}\endmoddef
  type :: parameter_t
     private
     integer :: type  = PAR_NONE
     class(modelpar_data_t), pointer :: data => null ()
     type(parse_node_t), pointer :: pn => null ()
     class(expr_t), allocatable :: expr
   contains
   \LA{}Models: parameter: TBP\RA{}
  end type parameter_t

\eatline
\nwendcode{}\nwbegindocs{487}Initialization depends on parameter type.  Independent parameters
are initialized by a constant value or a constant numerical expression
(which may contain a unit).  Derived parameters are initialized by an
arbitrary numerical expression, which makes use of the current
variable list.  The expression is evaluated by the function
{\tt{}parameter{\char95}reset}.

This implementation supports only real parameters and real values.
\nwenddocs{}\nwbegincode{488}\moddef{Models: parameter: TBP}\endmoddef
  procedure :: init_independent_value => parameter_init_independent_value
  procedure :: init_independent => parameter_init_independent
  procedure :: init_derived => parameter_init_derived
  procedure :: init_external => parameter_init_external
\nwendcode{}\nwbegincode{489}\moddef{Models: procedures}\endmoddef
  subroutine parameter_init_independent_value (par, par_data, name, value)
    class(parameter_t), intent(out) :: par
    class(modelpar_data_t), intent(in), target :: par_data
    type(string_t), intent(in) :: name
    real(default), intent(in) :: value
    par%type = PAR_INDEPENDENT
    par%data => par_data
    call par%data%init (name, value)
  end subroutine parameter_init_independent_value

  subroutine parameter_init_independent (par, par_data, name, pn)
    class(parameter_t), intent(out) :: par
    class(modelpar_data_t), intent(in), target :: par_data
    type(string_t), intent(in) :: name
    type(parse_node_t), intent(in), target :: pn
    par%type = PAR_INDEPENDENT
    par%pn => pn
    allocate (eval_tree_t :: par%expr)
    select type (expr => par%expr)
    type is (eval_tree_t)
       call expr%init_numeric_value (pn)
    end select
    par%data => par_data
    call par%data%init (name, par%expr%get_real ())
  end subroutine parameter_init_independent

  subroutine parameter_init_derived (par, par_data, name, pn, var_list)
    class(parameter_t), intent(out) :: par
    class(modelpar_data_t), intent(in), target :: par_data
    type(string_t), intent(in) :: name
    type(parse_node_t), intent(in), target :: pn
    type(var_list_t), intent(in), target :: var_list
    par%type = PAR_DERIVED
    par%pn => pn
    allocate (eval_tree_t :: par%expr)
    select type (expr => par%expr)
    type is (eval_tree_t)
       call expr%init_expr (pn, var_list=var_list)
    end select
    par%data => par_data
    call par%expr%evaluate ()
    call par%data%init (name, par%expr%get_real ())
  end subroutine parameter_init_derived

  subroutine parameter_init_external (par, par_data, name)
    class(parameter_t), intent(out) :: par
    class(modelpar_data_t), intent(in), target :: par_data
    type(string_t), intent(in) :: name
    par%type = PAR_EXTERNAL
    par%data => par_data
    call par%data%init (name, 0._default)
  end subroutine parameter_init_external

\eatline
\eatline
\eatline
\eatline
\nwendcode{}\nwbegindocs{490}The finalizer is needed for the evaluation tree in the definition.
\nwenddocs{}\nwbegincode{491}\moddef{Models: parameter: TBP}\plusendmoddef
  procedure :: final => parameter_final
\nwendcode{}\nwbegincode{492}\moddef{Models: procedures}\plusendmoddef
  subroutine parameter_final (par)
    class(parameter_t), intent(inout) :: par
    if (allocated (par%expr)) then
       call par%expr%final ()
    end if
  end subroutine parameter_final

\eatline
\nwendcode{}\nwbegindocs{493}All derived parameters should be recalculated if some independent
parameters have changed:
\nwenddocs{}\nwbegincode{494}\moddef{Models: parameter: TBP}\plusendmoddef
  procedure :: reset_derived => parameter_reset_derived
\nwendcode{}\nwbegincode{495}\moddef{Models: procedures}\plusendmoddef
  subroutine parameter_reset_derived (par)
    class(parameter_t), intent(inout) :: par
    select case (par%type)
    case (PAR_DERIVED)
       call par%expr%evaluate ()
       par%data = par%expr%get_real ()
    end select
  end subroutine parameter_reset_derived

\eatline
\nwendcode{}\nwbegindocs{496}Output.  [We should have a formula format for the eval tree,
suitable for input and output!]
\nwenddocs{}\nwbegincode{497}\moddef{Models: parameter: TBP}\plusendmoddef
  procedure :: write => parameter_write
\nwendcode{}\nwbegincode{498}\moddef{Models: procedures}\plusendmoddef
  subroutine parameter_write (par, unit, write_defs)
    class(parameter_t), intent(in) :: par
    integer, intent(in), optional :: unit
    logical, intent(in), optional :: write_defs
    logical :: defs
    integer :: u
    u = given_output_unit (unit);  if (u < 0)  return
    defs = .false.;  if (present (write_defs))  defs = write_defs
    select case (par%type)
    case (PAR_INDEPENDENT)
       write (u, "(3x,A)", advance="no")  "parameter"
    case (PAR_DERIVED)
       write (u, "(3x,A)", advance="no")  "derived"
    case (PAR_EXTERNAL)
       write (u, "(3x,A)", advance="no")  "external"
    end select
    call par%data%write (u)
    select case (par%type)
    case (PAR_DERIVED)
       if (defs) then
          call par%expr%write (unit)
       else
          write (u, *)
       end if
    case default
       write (u, *)
    end select
  end subroutine parameter_write

\eatline
\nwendcode{}\nwbegindocs{499}Screen output variant.  Restrict output to the given parameter type.
\nwenddocs{}\nwbegincode{500}\moddef{Models: parameter: TBP}\plusendmoddef
  procedure :: show => parameter_show
\nwendcode{}\nwbegincode{501}\moddef{Models: procedures}\plusendmoddef
  subroutine parameter_show (par, l, u, partype)
    class(parameter_t), intent(in) :: par
    integer, intent(in) :: l, u
    integer, intent(in) :: partype
    if (par%type == partype) then
       call par%data%show (l, u)
    end if
  end subroutine parameter_show
    
\eatline
\nwendcode{}\nwbegindocs{502}\nwdocspar
\subsection{Model Object}
A model object holds all information about parameters, particles,
and vertices.  For models that require an external program for
parameter calculation, there is the pointer to a function that does
this calculation, given the set of independent and derived parameters.

As explained above, the type inherits from {\tt{}model{\char95}data{\char95}t}, which is
the actual storage for the model data.

When reading a model, we create a parse tree.  Parameter definitions are
available via parse nodes.  Since we may need those later when making model
instances, we keep the whole parse tree in the model definition (but not in
the instances).
\nwenddocs{}\nwbegincode{503}\moddef{Models: public}\endmoddef
  public :: model_t
\nwendcode{}\nwbegincode{504}\moddef{Models: types}\plusendmoddef
  type, extends (model_data_t) :: model_t
     private
     character(32) :: md5sum = ""
     type(parameter_t), dimension(:), allocatable :: par
     integer :: max_par_name_length = 0
     integer :: max_field_name_length = 0
     type(var_list_t) :: var_list
     type(string_t) :: dlname
     procedure(model_init_external_parameters), nopass, pointer :: &
          init_external_parameters => null ()
     type(dlaccess_t) :: dlaccess
     type(parse_tree_t) :: parse_tree
   contains
   \LA{}Models: model: TBP\RA{}
  end type model_t

\eatline
\nwendcode{}\nwbegindocs{505}This is the interface for a procedure that initializes the
calculation of external parameters, given the array of all
parameters.
\nwenddocs{}\nwbegincode{506}\moddef{Models: interfaces}\endmoddef
  abstract interface
     subroutine model_init_external_parameters (par) bind (C)
       import
       real(c_default_float), dimension(*), intent(inout) :: par
     end subroutine model_init_external_parameters
  end interface

\eatline
\nwendcode{}\nwbegindocs{507}Initialization: Specify the number of parameters, particles,
vertices and allocate memory.  If an associated DL library is
specified, load this library.
\nwenddocs{}\nwbegincode{508}\moddef{Models: model: TBP}\endmoddef
  generic :: init => model_init
  procedure, private :: model_init
\nwendcode{}\nwbegincode{509}\moddef{Models: procedures}\plusendmoddef
  subroutine model_init &
       (model, name, libname, os_data, n_par, n_prt, n_vtx)
    class(model_t), intent(inout) :: model
    type(string_t), intent(in) :: name, libname
    type(os_data_t), intent(in) :: os_data
    integer, intent(in) :: n_par, n_prt, n_vtx
    type(c_funptr) :: c_fptr
    type(string_t) :: libpath
    call model%basic_init (name, n_par, n_prt, n_vtx)
    if (libname /= "") then
       if (.not. os_data%use_testfiles) then
          libpath = os_data%whizard_models_libpath_local
          model%dlname = os_get_dlname ( &
            libpath // "/" // libname, os_data, ignore=.true.)
       end if
       if (model%dlname == "") then
          libpath = os_data%whizard_models_libpath
          model%dlname = os_get_dlname (libpath // "/" // libname, os_data)
       end if
    else
       model%dlname = ""
    end if
    if (model%dlname /= "") then
       if (.not. dlaccess_is_open (model%dlaccess)) then
          if (logging) &
               call msg_message ("Loading model auxiliary library '" &
               // char (libpath) // "/" // char (model%dlname) // "'")
          call dlaccess_init (model%dlaccess, os_data%whizard_models_libpath, &
               model%dlname, os_data)
          if (dlaccess_has_error (model%dlaccess)) then
             call msg_message (char (dlaccess_get_error (model%dlaccess)))
             call msg_fatal ("Loading model auxiliary library '" &
                  // char (model%dlname) // "' failed")
             return
          end if
          c_fptr = dlaccess_get_c_funptr (model%dlaccess, &
               var_str ("init_external_parameters"))
          if (dlaccess_has_error (model%dlaccess)) then
             call msg_message (char (dlaccess_get_error (model%dlaccess)))
             call msg_fatal ("Loading function from auxiliary library '" &
                  // char (model%dlname) // "' failed")
             return
          end if
          call c_f_procpointer (c_fptr, model% init_external_parameters)
       end if
    end if
  end subroutine model_init

\eatline
\nwendcode{}\nwbegindocs{510}For a model instance, we do not attempt to load a DL library.  This is the
core of the full initializer above.
\nwenddocs{}\nwbegincode{511}\moddef{Models: model: TBP}\plusendmoddef
  procedure, private :: basic_init => model_basic_init
\nwendcode{}\nwbegincode{512}\moddef{Models: procedures}\plusendmoddef
  subroutine model_basic_init (model, name, n_par, n_prt, n_vtx)
    class(model_t), intent(inout) :: model
    type(string_t), intent(in) :: name
    integer, intent(in) :: n_par, n_prt, n_vtx
    allocate (model%par (n_par))
    call model%model_data_t%init (name, n_par, 0, n_prt, n_vtx)
  end subroutine model_basic_init
    
\eatline
\nwendcode{}\nwbegindocs{513}Finalization: The variable list contains allocated pointers, also the parse
tree.  We also close the DL access object, if any, that enables external
parameter calculation.
\nwenddocs{}\nwbegincode{514}\moddef{Models: model: TBP}\plusendmoddef
  procedure :: final => model_final
\nwendcode{}\nwbegincode{515}\moddef{Models: procedures}\plusendmoddef
  subroutine model_final (model)
    class(model_t), intent(inout) :: model
    integer :: i
    if (allocated (model%par)) then
       do i = 1, size (model%par)
          call model%par(i)%final ()
       end do
    end if
    call model%var_list%final (follow_link=.false.)
    if (model%dlname /= "")  call dlaccess_final (model%dlaccess)
    call parse_tree_final (model%parse_tree)
    call model%model_data_t%final ()
  end subroutine model_final

\eatline
\nwendcode{}\nwbegindocs{516}Output.  By default, the output is in the form of an input file.  If
{\tt{}verbose} is true, for each derived parameter the definition (eval
tree) is displayed, and the vertex hash table is shown.
\nwenddocs{}\nwbegincode{517}\moddef{Models: model: TBP}\plusendmoddef
  procedure :: write => model_write
\nwendcode{}\nwbegincode{518}\moddef{Models: procedures}\plusendmoddef
  subroutine model_write (model, unit, verbose, &
       show_md5sum, show_variables, show_parameters, &
       show_particles, show_vertices)
    class(model_t), intent(in) :: model
    integer, intent(in), optional :: unit
    logical, intent(in), optional :: verbose
    logical, intent(in), optional :: show_md5sum
    logical, intent(in), optional :: show_variables
    logical, intent(in), optional :: show_parameters
    logical, intent(in), optional :: show_particles
    logical, intent(in), optional :: show_vertices
    logical :: verb, show_md5, show_par, show_var
    integer :: u, i
    u = given_output_unit (unit);  if (u < 0)  return
    verb = .false.;  if (present (verbose))  verb = verbose
    show_md5 = .true.;  if (present (show_md5sum)) &
         show_md5 = show_md5sum
    show_par = .true.;  if (present (show_parameters)) &
         show_par = show_parameters
    show_var = verb;  if (present (show_variables)) &
         show_var = show_variables
    write (u, "(A,A,A)") 'model "', char (model%get_name ()), '"'
    if (show_md5 .and. model%md5sum /= "") &
         write (u, "(1x,A,A,A)") "! md5sum = '", model%md5sum, "'"
    if (show_par) then
       write (u, *)
       do i = 1, size (model%par)
          call model%par(i)%write (u, write_defs=verbose)
       end do
    end if
    call model%model_data_t%write (unit, verbose, &
         show_md5sum, show_variables, &
         show_parameters=.false., &
         show_particles=show_particles, &
         show_vertices=show_vertices)
    if (show_var) then
       write (u, *)
       call var_list_write (model%var_list, unit, follow_link=.false.)
    end if
  end subroutine model_write

\eatline
\nwendcode{}\nwbegindocs{519}Screen output, condensed form.
\nwenddocs{}\nwbegincode{520}\moddef{Models: model: TBP}\plusendmoddef
  procedure :: show => model_show
\nwendcode{}\nwbegincode{521}\moddef{Models: procedures}\plusendmoddef
  subroutine model_show (model, unit)
    class(model_t), intent(in) :: model
    integer, intent(in), optional :: unit
    integer :: i, u, l
    u = given_output_unit (unit)
    write (u, "(A,1x,A)")  "Model:", char (model%get_name ())
    l = model%max_field_name_length
    call model%show_fields (l, u)
    l = model%max_par_name_length
    if (any (model%par%type == PAR_INDEPENDENT)) then
       write (u, "(2x,A)")  "Independent parameters:"
       do i = 1, size (model%par)
          call model%par(i)%show (l, u, PAR_INDEPENDENT)
       end do
    end if
    if (any (model%par%type == PAR_DERIVED)) then
       write (u, "(2x,A)")  "Derived parameters:"
       do i = 1, size (model%par)
          call model%par(i)%show (l, u, PAR_DERIVED)
       end do
    end if
    if (any (model%par%type == PAR_EXTERNAL)) then
       write (u, "(2x,A)")  "External parameters:"
       do i = 1, size (model%par)
          call model%par(i)%show (l, u, PAR_EXTERNAL)
       end do
    end if
  end subroutine model_show

\eatline
\nwendcode{}\nwbegindocs{522}Show all fields/particles.
\nwenddocs{}\nwbegincode{523}\moddef{Models: model: TBP}\plusendmoddef
  procedure :: show_fields => model_show_fields
\nwendcode{}\nwbegincode{524}\moddef{Models: procedures}\plusendmoddef
  subroutine model_show_fields (model, l, unit)
    class(model_t), intent(in), target :: model
    integer, intent(in) :: l
    integer, intent(in), optional :: unit
    type(field_data_t), pointer :: field
    integer :: u, i
    u = given_output_unit (unit)
    write (u, "(2x,A)")  "Particles:"
    do i = 1, model%get_n_field ()
       field => model%get_field_ptr_by_index (i)
       call field%show (l, u)
    end do
  end subroutine model_show_fields
  
\eatline
\nwendcode{}\nwbegindocs{525}Show the list of stable, unstable, polarized, or unpolarized
particles, respectively.
\nwenddocs{}\nwbegincode{526}\moddef{Models: model: TBP}\plusendmoddef
  procedure :: show_stable => model_show_stable
  procedure :: show_unstable => model_show_unstable
  procedure :: show_polarized => model_show_polarized
  procedure :: show_unpolarized => model_show_unpolarized
\nwendcode{}\nwbegincode{527}\moddef{Models: procedures}\plusendmoddef
  subroutine model_show_stable (model, unit)
    class(model_t), intent(in), target :: model
    integer, intent(in), optional :: unit
    type(field_data_t), pointer :: field
    integer :: u, i
    u = given_output_unit (unit)
    write (u, "(A,1x)", advance="no")  "Stable particles:"
    do i = 1, model%get_n_field ()
       field => model%get_field_ptr_by_index (i)
       if (field%is_stable (.false.)) then
          write (u, "(1x,A)", advance="no")  char (field%get_name (.false.))
       end if
       if (field%has_antiparticle ()) then
          if (field%is_stable (.true.)) then
             write (u, "(1x,A)", advance="no")  char (field%get_name (.true.))
          end if
       end if
    end do
    write (u, *)
  end subroutine model_show_stable
  
  subroutine model_show_unstable (model, unit)
    class(model_t), intent(in), target :: model
    integer, intent(in), optional :: unit
    type(field_data_t), pointer :: field
    integer :: u, i
    u = given_output_unit (unit)
    write (u, "(A,1x)", advance="no")  "Unstable particles:"
    do i = 1, model%get_n_field ()
       field => model%get_field_ptr_by_index (i)
       if (.not. field%is_stable (.false.)) then
          write (u, "(1x,A)", advance="no")  char (field%get_name (.false.))
       end if
       if (field%has_antiparticle ()) then
          if (.not. field%is_stable (.true.)) then
             write (u, "(1x,A)", advance="no")  char (field%get_name (.true.))
          end if
       end if
    end do
    write (u, *)
  end subroutine model_show_unstable
  
  subroutine model_show_polarized (model, unit)
    class(model_t), intent(in), target :: model
    integer, intent(in), optional :: unit
    type(field_data_t), pointer :: field
    integer :: u, i
    u = given_output_unit (unit)
    write (u, "(A,1x)", advance="no")  "Polarized particles:"
    do i = 1, model%get_n_field ()
       field => model%get_field_ptr_by_index (i)
       if (field%is_polarized (.false.)) then
          write (u, "(1x,A)", advance="no") char (field%get_name (.false.))
       end if
       if (field%has_antiparticle ()) then
          if (field%is_polarized (.true.)) then
             write (u, "(1x,A)", advance="no") char (field%get_name (.true.))
          end if
       end if
    end do
    write (u, *)
  end subroutine model_show_polarized
  
  subroutine model_show_unpolarized (model, unit)
    class(model_t), intent(in), target :: model
    integer, intent(in), optional :: unit
    type(field_data_t), pointer :: field
    integer :: u, i
    u = given_output_unit (unit)
    write (u, "(A,1x)", advance="no")  "Unpolarized particles:"
    do i = 1, model%get_n_field ()
       field => model%get_field_ptr_by_index (i)
       if (.not. field%is_polarized (.false.)) then
          write (u, "(1x,A)", advance="no") &
               char (field%get_name (.false.))
       end if
       if (field%has_antiparticle ()) then
          if (.not. field%is_polarized (.true.)) then
             write (u, "(1x,A)", advance="no") char (field%get_name (.true.))
          end if
       end if
    end do
    write (u, *)
  end subroutine model_show_unpolarized
  
\eatline
\eatline
\eatline
\eatline
\nwendcode{}\nwbegindocs{528}Retrieve the MD5 sum of a model (actually, of the model file).
\nwenddocs{}\nwbegincode{529}\moddef{Models: model: TBP}\plusendmoddef
  procedure :: get_md5sum => model_get_md5sum
\nwendcode{}\nwbegincode{530}\moddef{Models: procedures}\plusendmoddef
  function model_get_md5sum (model) result (md5sum)
    character(32) :: md5sum
    class(model_t), intent(in) :: model
    md5sum = model%md5sum
  end function model_get_md5sum

\eatline
\nwendcode{}\nwbegindocs{531}Parameters are defined by an expression which may be constant or
arbitrary.

Note: the auxiliary pointer {\tt{}value{\char95}ptr} is a workaround for a bug in gfortran
4.8.1: the target of the function pointer is lost, if the pointer is provided
directly as argument.
\nwenddocs{}\nwbegincode{532}\moddef{Models: model: TBP}\plusendmoddef
  procedure, private :: &
       set_parameter_constant => model_set_parameter_constant
  procedure, private :: &
       set_parameter_parse_node => model_set_parameter_parse_node
  procedure, private :: &
       set_parameter_external => model_set_parameter_external
\nwendcode{}\nwbegincode{533}\moddef{Models: procedures}\plusendmoddef
  subroutine model_set_parameter_constant (model, i, name, value)
    class(model_t), intent(inout), target :: model
    integer, intent(in) :: i
    type(string_t), intent(in) :: name
    real(default), intent(in) :: value
    logical, save, target :: known = .true.
    class(modelpar_data_t), pointer :: par_data
    real(default), pointer :: value_ptr
    par_data => model%get_par_real_ptr (i)
    call model%par(i)%init_independent_value (par_data, name, value)
    value_ptr => par_data%get_real_ptr ()
    call var_list_append_real_ptr (model%var_list, &
         name, value_ptr, &
         is_known=known, intrinsic=.true.)
    model%max_par_name_length = max (model%max_par_name_length, len (name))
  end subroutine model_set_parameter_constant

  subroutine model_set_parameter_parse_node (model, i, name, pn, constant)
    class(model_t), intent(inout), target :: model
    integer, intent(in) :: i
    type(string_t), intent(in) :: name
    type(parse_node_t), intent(in), target :: pn
    logical, intent(in) :: constant
    logical, save, target :: known = .true.
    class(modelpar_data_t), pointer :: par_data
    real(default), pointer :: value_ptr
    par_data => model%get_par_real_ptr (i)
    if (constant) then
       call model%par(i)%init_independent (par_data, name, pn)
    else
       call model%par(i)%init_derived (par_data, name, pn, model%var_list)
    end if
    value_ptr => par_data%get_real_ptr ()
    call var_list_append_real_ptr (model%var_list, &
         name, value_ptr, &
         is_known=known, locked=.not.constant, intrinsic=.true.)
    model%max_par_name_length = max (model%max_par_name_length, len (name))
  end subroutine model_set_parameter_parse_node

  subroutine model_set_parameter_external (model, i, name)
    class(model_t), intent(inout), target :: model
    integer, intent(in) :: i
    type(string_t), intent(in) :: name
    logical, save, target :: known = .true.
    class(modelpar_data_t), pointer :: par_data
    real(default), pointer :: value_ptr
    par_data => model%get_par_real_ptr (i)
    call model%par(i)%init_external (par_data, name)
    value_ptr => par_data%get_real_ptr ()
    call var_list_append_real_ptr (model%var_list, &
         name, value_ptr, &
         is_known=known, locked=.true., intrinsic=.true.)
    model%max_par_name_length = max (model%max_par_name_length, len (name))
  end subroutine model_set_parameter_external

\eatline
\nwendcode{}\nwbegindocs{534}Make a copy of a parameter.  We assume that the {\tt{}model{\char95}data{\char95}t}
parameter arrays have already been copied, so names and values are
available in the current model, and can be used as targets.  The eval
tree should not be copied, since it should refer to the new variable
list.  The safe solution is to make use of the above initializers,
which also take care of the building a new variable list.
\nwenddocs{}\nwbegincode{535}\moddef{Models: model: TBP}\plusendmoddef
  procedure, private :: copy_parameter => model_copy_parameter
\nwendcode{}\nwbegincode{536}\moddef{Models: procedures}\plusendmoddef
  subroutine model_copy_parameter (model, i, par)
    class(model_t), intent(inout), target :: model
    integer, intent(in) :: i
    type(parameter_t), intent(in) :: par
    type(string_t) :: name
    real(default) :: value
    name = par%data%get_name ()
    select case (par%type)
    case (PAR_INDEPENDENT)
       if (associated (par%pn)) then
          call model%set_parameter_parse_node (i, name, par%pn, &
               constant = .true.)
       else
          value = par%data%get_real ()
          call model%set_parameter_constant (i, name, value)
       end if
    case (PAR_DERIVED)
       call model%set_parameter_parse_node (i, name, par%pn, &
            constant = .false.)
    case (PAR_EXTERNAL)
       call model%set_parameter_external (i, name)
    end select
  end subroutine model_copy_parameter
  
\eatline
\nwendcode{}\nwbegindocs{537}Rcalculate all derived parameters.
\nwenddocs{}\nwbegincode{538}\moddef{Models: model: TBP}\plusendmoddef
  procedure :: update_parameters => model_parameters_update
\nwendcode{}\nwbegincode{539}\moddef{Models: procedures}\plusendmoddef
  subroutine model_parameters_update (model)
    class(model_t), intent(inout) :: model
    integer :: i
    real(default), dimension(:), allocatable :: par
    do i = 1, size (model%par)
       call model%par(i)%reset_derived ()
    end do
    if (associated (model%init_external_parameters)) then
       allocate (par (model%get_n_real ()))
       call model%real_parameters_to_c_array (par)
       call model%init_external_parameters (par)
       call model%real_parameters_from_c_array (par)
    end if
  end subroutine model_parameters_update

\eatline
\nwendcode{}\nwbegindocs{540}Initialize field data with PDG long name and PDG code.
\nwenddocs{}\nwbegincode{541}\moddef{Models: model: TBP}\plusendmoddef
  procedure, private :: init_field => model_init_field
\nwendcode{}\nwbegincode{542}\moddef{Models: procedures}\plusendmoddef
  subroutine model_init_field (model, i, longname, pdg)
    class(model_t), intent(inout), target :: model
    integer, intent(in) :: i
    type(string_t), intent(in) :: longname
    integer, intent(in) :: pdg
    type(field_data_t), pointer :: field
    field => model%get_field_ptr_by_index (i)
    call field%init (longname, pdg)
  end subroutine model_init_field
    
\eatline
\nwendcode{}\nwbegindocs{543}Copy field data for index {\tt{}i} from another particle which serves
as a template.  The name should be the unique long name.
\nwenddocs{}\nwbegincode{544}\moddef{Models: model: TBP}\plusendmoddef
  procedure, private :: copy_field => model_copy_field
\nwendcode{}\nwbegincode{545}\moddef{Models: procedures}\plusendmoddef
  subroutine model_copy_field (model, i, name_src)
    class(model_t), intent(inout), target :: model
    integer, intent(in) :: i
    type(string_t), intent(in) :: name_src
    type(field_data_t), pointer :: field_src, field
    field_src => model%get_field_ptr (name_src)
    field => model%get_field_ptr_by_index (i)
    call field%copy_from (field_src)
  end subroutine model_copy_field

\eatline
\nwendcode{}\nwbegindocs{546}\nwdocspar
\subsection{Model Access via Variables}
Write the model variable list.
\nwenddocs{}\nwbegincode{547}\moddef{Models: model: TBP}\plusendmoddef
  procedure :: write_var_list => model_write_var_list
\nwendcode{}\nwbegincode{548}\moddef{Models: procedures}\plusendmoddef
  subroutine model_write_var_list (model, unit, follow_link)
    class(model_t), intent(in) :: model
    integer, intent(in), optional :: unit
    logical, intent(in), optional :: follow_link
    call var_list_write (model%var_list, unit, follow_link)
  end subroutine model_write_var_list
  
\eatline
\nwendcode{}\nwbegindocs{549}Link a variable list to the model variables.
\nwenddocs{}\nwbegincode{550}\moddef{Models: model: TBP}\plusendmoddef
  procedure :: link_var_list => model_link_var_list
\nwendcode{}\nwbegincode{551}\moddef{Models: procedures}\plusendmoddef
  subroutine model_link_var_list (model, var_list)
    class(model_t), intent(inout) :: model
    type(var_list_t), intent(in), target :: var_list
    call model%var_list%link (var_list)
  end subroutine model_link_var_list

\eatline
\nwendcode{}\nwbegindocs{552}\nwdocspar
Check if the model contains a named variable.
\nwenddocs{}\nwbegincode{553}\moddef{Models: model: TBP}\plusendmoddef
  procedure :: var_exists => model_var_exists
\nwendcode{}\nwbegincode{554}\moddef{Models: procedures}\plusendmoddef
  function model_var_exists (model, name) result (flag)
    class(model_t), intent(in) :: model
    type(string_t), intent(in) :: name
    logical :: flag
    flag = model%var_list%contains (name, follow_link=.false.)
  end function model_var_exists
  
\eatline
\nwendcode{}\nwbegindocs{555}Check if the model variable is a derived parameter, i.e., locked.
\nwenddocs{}\nwbegincode{556}\moddef{Models: model: TBP}\plusendmoddef
  procedure :: var_is_locked => model_var_is_locked
\nwendcode{}\nwbegincode{557}\moddef{Models: procedures}\plusendmoddef
  function model_var_is_locked (model, name) result (flag)
    class(model_t), intent(in) :: model
    type(string_t), intent(in) :: name
    logical :: flag
    flag = model%var_list%is_locked (name, follow_link=.false.)
  end function model_var_is_locked
  
\eatline
\nwendcode{}\nwbegindocs{558}Set a model parameter via the named variable.  We assume that the
variable exists and is writable, i.e., non-locked.  We update the
model and variable list, so independent and derived parameters are
always synchronized.
\nwenddocs{}\nwbegincode{559}\moddef{Models: model: TBP}\plusendmoddef
  procedure :: set_real => model_var_set_real
\nwendcode{}\nwbegincode{560}\moddef{Models: procedures}\plusendmoddef
  subroutine model_var_set_real (model, name, rval, verbose, pacified)
    class(model_t), intent(inout) :: model
    type(string_t), intent(in) :: name
    real(default), intent(in) :: rval
    logical, intent(in), optional :: verbose, pacified
    call var_list_set_real (model%var_list, name, rval, &
         is_known=.true., ignore=.false., &
         verbose=verbose, model_name=model%get_name (), pacified=pacified)
    call model%update_parameters ()
  end subroutine model_var_set_real
    
\eatline
\nwendcode{}\nwbegindocs{561}Retrieve a model parameter value.
\nwenddocs{}\nwbegincode{562}\moddef{Models: model: TBP}\plusendmoddef
  procedure :: get_rval => model_var_get_rval
\nwendcode{}\nwbegincode{563}\moddef{Models: procedures}\plusendmoddef
  function model_var_get_rval (model, name) result (rval)
    class(model_t), intent(in) :: model
    type(string_t), intent(in) :: name
    real(default) :: rval
    rval = model%var_list%get_rval (name, follow_link=.false.)
  end function model_var_get_rval
  
\eatline
\nwendcode{}\nwbegindocs{564}\nwdocspar
[To be deleted] Return a pointer to the variable list.
\nwenddocs{}\nwbegincode{565}\moddef{Models: model: TBP}\plusendmoddef
  procedure :: get_var_list_ptr => model_get_var_list_ptr
\nwendcode{}\nwbegincode{566}\moddef{Models: procedures}\plusendmoddef
  function model_get_var_list_ptr (model) result (var_list)
    type(var_list_t), pointer :: var_list
    class(model_t), intent(in), target :: model
    var_list => model%var_list
  end function model_get_var_list_ptr

\eatline
\nwendcode{}\nwbegindocs{567}\nwdocspar
\subsection{Reading models from file}
This procedure defines the model-file syntax for the parser, returning
an internal file (ifile).

Note that arithmetic operators are defined as keywords in the
expression syntax, so we exclude them here.
\nwenddocs{}\nwbegincode{568}\moddef{Models: procedures}\plusendmoddef
  subroutine define_model_file_syntax (ifile)
    type(ifile_t), intent(inout) :: ifile
    call ifile_append (ifile, "SEQ model_def = model_name_def " // &
         "parameters derived_pars external_pars particles vertices")
    call ifile_append (ifile, "SEQ model_name_def = model model_name")
    call ifile_append (ifile, "KEY model")
    call ifile_append (ifile, "QUO model_name = '""'...'""'")
    call ifile_append (ifile, "SEQ parameters = parameter_def*")
    call ifile_append (ifile, "SEQ parameter_def = parameter par_name " // &
         "'=' any_real_value")
    call ifile_append (ifile, "ALT any_real_value = " &
         // "neg_real_value | pos_real_value | real_value")
    call ifile_append (ifile, "SEQ neg_real_value = '-' real_value")
    call ifile_append (ifile, "SEQ pos_real_value = '+' real_value")
    call ifile_append (ifile, "KEY parameter")
    call ifile_append (ifile, "IDE par_name")
    ! call ifile_append (ifile, "KEY '='")          !!! Key already exists
    call ifile_append (ifile, "SEQ derived_pars = derived_def*")
    call ifile_append (ifile, "SEQ derived_def = derived par_name " // &
         "'=' expr")
    call ifile_append (ifile, "KEY derived")
    call ifile_append (ifile, "SEQ external_pars = external_def*")
    call ifile_append (ifile, "SEQ external_def = external par_name")
    call ifile_append (ifile, "KEY external")
    call ifile_append (ifile, "SEQ particles = particle_def*")
    call ifile_append (ifile, "SEQ particle_def = particle prt_longname " // &
         "prt_pdg prt_details")
    call ifile_append (ifile, "KEY particle")
    call ifile_append (ifile, "IDE prt_longname")
    call ifile_append (ifile, "INT prt_pdg")
    call ifile_append (ifile, "ALT prt_details = prt_src | prt_properties")
    call ifile_append (ifile, "SEQ prt_src = like prt_longname prt_properties")
    call ifile_append (ifile, "KEY like")
    call ifile_append (ifile, "SEQ prt_properties = prt_property*")
    call ifile_append (ifile, "ALT prt_property = " // & 
         "parton | invisible | gauge | left | right | " // &
         "prt_name | prt_anti | prt_tex_name | prt_tex_anti | " // &
         "prt_spin | prt_isospin | prt_charge | " // &
         "prt_color | prt_mass | prt_width")
    call ifile_append (ifile, "KEY parton")
    call ifile_append (ifile, "KEY invisible")
    call ifile_append (ifile, "KEY gauge")
    call ifile_append (ifile, "KEY left")
    call ifile_append (ifile, "KEY right")
    call ifile_append (ifile, "SEQ prt_name = name name_def+")
    call ifile_append (ifile, "SEQ prt_anti = anti name_def+")
    call ifile_append (ifile, "SEQ prt_tex_name = tex_name name_def")
    call ifile_append (ifile, "SEQ prt_tex_anti = tex_anti name_def")
    call ifile_append (ifile, "KEY name")
    call ifile_append (ifile, "KEY anti")
    call ifile_append (ifile, "KEY tex_name")
    call ifile_append (ifile, "KEY tex_anti")
    call ifile_append (ifile, "ALT name_def = name_string | name_id")
    call ifile_append (ifile, "QUO name_string = '""'...'""'")
    call ifile_append (ifile, "IDE name_id")
    call ifile_append (ifile, "SEQ prt_spin = spin frac")
    call ifile_append (ifile, "KEY spin")
    call ifile_append (ifile, "SEQ prt_isospin = isospin frac")
    call ifile_append (ifile, "KEY isospin")
    call ifile_append (ifile, "SEQ prt_charge = charge frac")
    call ifile_append (ifile, "KEY charge")
    call ifile_append (ifile, "SEQ prt_color = color integer_literal")
    call ifile_append (ifile, "KEY color")
    call ifile_append (ifile, "SEQ prt_mass = mass par_name")
    call ifile_append (ifile, "KEY mass")
    call ifile_append (ifile, "SEQ prt_width = width par_name")
    call ifile_append (ifile, "KEY width")
    call ifile_append (ifile, "SEQ vertices = vertex_def*")
    call ifile_append (ifile, "SEQ vertex_def = vertex name_def+")
    call ifile_append (ifile, "KEY vertex")
    call define_expr_syntax (ifile, particles=.false., analysis=.false.)
  end subroutine define_model_file_syntax

\eatline
\nwendcode{}\nwbegindocs{569}The model-file syntax and lexer are fixed, therefore stored as
module variables:
\nwenddocs{}\nwbegincode{570}\moddef{Models: variables}\endmoddef
  type(syntax_t), target, save :: syntax_model_file

\eatline
\nwendcode{}\nwbegincode{571}\moddef{Models: public}\plusendmoddef
  public :: syntax_model_file_init
\nwendcode{}\nwbegincode{572}\moddef{Models: procedures}\plusendmoddef
  subroutine syntax_model_file_init ()
    type(ifile_t) :: ifile
    call define_model_file_syntax (ifile)
    call syntax_init (syntax_model_file, ifile)
    call ifile_final (ifile)
  end subroutine syntax_model_file_init

\eatline
\nwendcode{}\nwbegincode{573}\moddef{Models: procedures}\plusendmoddef
  subroutine lexer_init_model_file (lexer)
    type(lexer_t), intent(out) :: lexer
    call lexer_init (lexer, &
         comment_chars = "#!", &
         quote_chars = '"\{', &
         quote_match = '"\}', &
         single_chars = ":()", &
         special_class = [ "+-*/^", "<>=  " ] , &
         keyword_list = syntax_get_keyword_list_ptr (syntax_model_file))
  end subroutine lexer_init_model_file

\eatline
\nwendcode{}\nwbegincode{574}\moddef{Models: public}\plusendmoddef
  public :: syntax_model_file_final
\nwendcode{}\nwbegincode{575}\moddef{Models: procedures}\plusendmoddef
  subroutine syntax_model_file_final ()
    call syntax_final (syntax_model_file)
  end subroutine syntax_model_file_final

\eatline
\nwendcode{}\nwbegincode{576}\moddef{Models: public}\plusendmoddef
  public :: syntax_model_file_write
\nwendcode{}\nwbegincode{577}\moddef{Models: procedures}\plusendmoddef
  subroutine syntax_model_file_write (unit)
    integer, intent(in), optional :: unit
    call syntax_write (syntax_model_file, unit)
  end subroutine syntax_model_file_write

\eatline
\nwendcode{}\nwbegincode{578}\moddef{Models: model: TBP}\plusendmoddef
  procedure :: read => model_read
\nwendcode{}\nwbegincode{579}\moddef{Models: procedures}\plusendmoddef
  subroutine model_read (model, filename, os_data, exist)
    class(model_t), intent(out), target :: model
    type(string_t), intent(in) :: filename
    type(os_data_t), intent(in) :: os_data
    logical, intent(out), optional :: exist
    type(string_t) :: file
    type(stream_t), target :: stream
    type(lexer_t) :: lexer
    integer :: unit
    character(32) :: model_md5sum
    type(parse_node_t), pointer :: nd_model_def, nd_model_name_def
    type(parse_node_t), pointer :: nd_parameters, nd_derived_pars
    type(parse_node_t), pointer :: nd_external_pars
    type(parse_node_t), pointer :: nd_particles, nd_vertices
    type(string_t) :: model_name, lib_name
    integer :: n_par, n_der, n_ext, n_prt, n_vtx
    real(c_default_float), dimension(:), allocatable :: par
    integer :: i
    type(parse_node_t), pointer :: nd_par_def
    type(parse_node_t), pointer :: nd_der_def
    type(parse_node_t), pointer :: nd_ext_def
    type(parse_node_t), pointer :: nd_prt
    type(parse_node_t), pointer :: nd_vtx
    logical :: model_exist
    file = filename
    inquire (file=char(file), exist=model_exist)
    if ((.not. model_exist) .and. (.not. os_data%use_testfiles)) then
       file = os_data%whizard_modelpath_local // "/" // filename
       inquire (file = char (file), exist = model_exist)
    end if
    if (.not. model_exist) then
       file = os_data%whizard_modelpath // "/" // filename
       inquire (file = char (file), exist = model_exist)
    end if
    if (.not. model_exist) then
       call msg_fatal ("Model file '" // char (filename) // "' not found")
       if (present (exist))  exist = .false.
       return
    end if
    if (present (exist))  exist = .true.
    if (logging) call msg_message ("Reading model file '" // char (file) // "'")
    call lexer_init_model_file (lexer)
    unit = free_unit ()
    open (file=char(file), unit=unit, action="read", status="old")
    model_md5sum = md5sum (unit)
    close (unit)
    call stream_init (stream, char (file))
    call lexer_assign_stream (lexer, stream)
    call parse_tree_init (model%parse_tree, syntax_model_file, lexer)
    call stream_final (stream)
    call lexer_final (lexer)
    nd_model_def => parse_tree_get_root_ptr (model%parse_tree)
    nd_model_name_def => parse_node_get_sub_ptr (nd_model_def)
    model_name = parse_node_get_string &
         (parse_node_get_sub_ptr (nd_model_name_def, 2))
    nd_parameters => parse_node_get_next_ptr (nd_model_name_def)
    if (associated (nd_parameters)) then
       if (parse_node_get_rule_key (nd_parameters) == "parameters") then
          n_par = parse_node_get_n_sub (nd_parameters)
          nd_par_def => parse_node_get_sub_ptr (nd_parameters)
          nd_derived_pars => parse_node_get_next_ptr (nd_parameters)
       else
          n_par = 0
          nd_derived_pars => nd_parameters
          nd_parameters => null ()
       end if
    else
       n_par = 0
       nd_derived_pars => null ()
    end if
    if (associated (nd_derived_pars)) then
       if (parse_node_get_rule_key (nd_derived_pars) == "derived_pars") then
          n_der = parse_node_get_n_sub (nd_derived_pars)
          nd_der_def => parse_node_get_sub_ptr (nd_derived_pars)
          nd_external_pars => parse_node_get_next_ptr (nd_derived_pars)
       else
          n_der = 0
          nd_external_pars => nd_derived_pars
          nd_derived_pars => null ()
       end if
    else
       n_der = 0
       nd_external_pars => null ()
    end if
    if (associated (nd_external_pars)) then
       if (parse_node_get_rule_key (nd_external_pars) == "external_pars") then
          n_ext = parse_node_get_n_sub (nd_external_pars)
          lib_name = "external." // model_name
          nd_ext_def => parse_node_get_sub_ptr (nd_external_pars)
          nd_particles => parse_node_get_next_ptr (nd_external_pars)
       else
          n_ext = 0
          lib_name = ""
          nd_particles => nd_external_pars
          nd_external_pars => null ()
       end if
    else
       n_ext = 0
       lib_name = ""
       nd_particles => null ()
    end if
    if (associated (nd_particles)) then
       if (parse_node_get_rule_key (nd_particles) == "particles") then
          n_prt = parse_node_get_n_sub (nd_particles)
          nd_prt => parse_node_get_sub_ptr (nd_particles)
          nd_vertices => parse_node_get_next_ptr (nd_particles)
       else
          n_prt = 0
          nd_vertices => nd_particles
          nd_particles => null ()
       end if
    else
       n_prt = 0
       nd_vertices => null ()
    end if
    if (associated (nd_vertices)) then
       n_vtx = parse_node_get_n_sub (nd_vertices)
       nd_vtx => parse_node_get_sub_ptr (nd_vertices)
    else
       n_vtx = 0
    end if
    call model%init (model_name, lib_name, os_data, &
         n_par + n_der + n_ext, n_prt, n_vtx)
    model%md5sum = model_md5sum
    do i = 1, n_par
       call model%read_parameter (i, nd_par_def)
       nd_par_def => parse_node_get_next_ptr (nd_par_def)
    end do
    do i = n_par + 1, n_par + n_der
       call model%read_derived (i, nd_der_def)
       nd_der_def => parse_node_get_next_ptr (nd_der_def)
    end do
    do i = n_par + n_der + 1, n_par + n_der + n_ext
       call model%read_external (i, nd_ext_def)
       nd_ext_def => parse_node_get_next_ptr (nd_ext_def)
    end do
    if (associated (model% init_external_parameters)) then
       allocate (par (model%get_n_real ()))
       call model%real_parameters_to_c_array (par)
       call model%init_external_parameters (par)
       call model%real_parameters_from_c_array (par)
    end if
    do i = 1, n_prt
       call model%read_field (i, nd_prt)
       nd_prt => parse_node_get_next_ptr (nd_prt)
    end do
    do i = 1, n_vtx
       call model%read_vertex (i, nd_vtx)
       nd_vtx => parse_node_get_next_ptr (nd_vtx)
    end do
    call model%freeze_vertices ()
    call model%append_field_vars ()
  end subroutine model_read

\eatline
\nwendcode{}\nwbegindocs{580}Parameters are real values (literal) with an optional unit.
\nwenddocs{}\nwbegincode{581}\moddef{Models: model: TBP}\plusendmoddef
  procedure, private :: read_parameter => model_read_parameter
\nwendcode{}\nwbegincode{582}\moddef{Models: procedures}\plusendmoddef
  subroutine model_read_parameter (model, i, node)
    class(model_t), intent(inout), target :: model
    integer, intent(in) :: i
    type(parse_node_t), intent(in), target :: node
    type(parse_node_t), pointer :: node_name, node_val
    type(string_t) :: name
    node_name => parse_node_get_sub_ptr (node, 2)
    name = parse_node_get_string (node_name)
    node_val => parse_node_get_next_ptr (node_name, 2)
    call model%set_parameter_parse_node (i, name, node_val, constant=.true.)
  end subroutine model_read_parameter

\eatline
\nwendcode{}\nwbegindocs{583}Derived parameters have any numeric expression as their definition.
\nwenddocs{}\nwbegincode{584}\moddef{Models: model: TBP}\plusendmoddef
  procedure, private :: read_derived => model_read_derived
\nwendcode{}\nwbegincode{585}\moddef{Models: procedures}\plusendmoddef
  subroutine model_read_derived (model, i, node)
    class(model_t), intent(inout), target :: model
    integer, intent(in) :: i
    type(parse_node_t), intent(in), target :: node
    type(string_t) :: name
    type(parse_node_t), pointer :: pn_expr
    name = parse_node_get_string (parse_node_get_sub_ptr (node, 2))
    pn_expr => parse_node_get_sub_ptr (node, 4)
    call model%set_parameter_parse_node (i, name, pn_expr, constant=.false.)
  end subroutine model_read_derived

\eatline
\nwendcode{}\nwbegindocs{586}External parameters have no definition; they are handled by an
external library. 
\nwenddocs{}\nwbegincode{587}\moddef{Models: model: TBP}\plusendmoddef
  procedure, private :: read_external => model_read_external
\nwendcode{}\nwbegincode{588}\moddef{Models: procedures}\plusendmoddef
  subroutine model_read_external (model, i, node)
    class(model_t), intent(inout), target :: model
    integer, intent(in) :: i
    type(parse_node_t), intent(in), target :: node
    type(string_t) :: name
    name = parse_node_get_string (parse_node_get_sub_ptr (node, 2))
    call model%set_parameter_external (i, name)
  end subroutine model_read_external

\eatline
\nwendcode{}\nwbegincode{589}\moddef{Models: model: TBP}\plusendmoddef
  procedure, private :: read_field => model_read_field
\nwendcode{}\nwbegincode{590}\moddef{Models: procedures}\plusendmoddef
  subroutine model_read_field (model, i, node)
    class(model_t), intent(inout), target :: model
    integer, intent(in) :: i
    type(parse_node_t), intent(in) :: node
    type(parse_node_t), pointer :: nd_src, nd_props, nd_prop
    type(string_t) :: longname
    integer :: pdg
    type(string_t) :: name_src
    type(string_t), dimension(:), allocatable :: name
    type(field_data_t), pointer :: field, field_src
    longname = parse_node_get_string (parse_node_get_sub_ptr (node, 2))
    pdg = parse_node_get_integer (parse_node_get_sub_ptr (node, 3)) 
    field => model%get_field_ptr_by_index (i)
    call field%init (longname, pdg)
    nd_src => parse_node_get_sub_ptr (node, 4)
    if (associated (nd_src)) then
       if (parse_node_get_rule_key (nd_src) == "prt_src") then
          name_src = parse_node_get_string (parse_node_get_sub_ptr (nd_src, 2))
          field_src => model%get_field_ptr (name_src, check=.true.)
          call field%copy_from (field_src)
          nd_props => parse_node_get_sub_ptr (nd_src, 3)
       else
          nd_props => nd_src
       end if
       nd_prop => parse_node_get_sub_ptr (nd_props)
       do while (associated (nd_prop))
          select case (char (parse_node_get_rule_key (nd_prop)))
          case ("invisible")
             call field%set (is_visible=.false.)
          case ("parton")
             call field%set (is_parton=.true.)
          case ("gauge")
             call field%set (is_gauge=.true.)
          case ("left")
             call field%set (is_left_handed=.true.)
          case ("right")
             call field%set (is_right_handed=.true.)
          case ("prt_name")
             call read_names (nd_prop, name)
             call field%set (name=name)
          case ("prt_anti")
             call read_names (nd_prop, name)
             call field%set (anti=name)
          case ("prt_tex_name")
             call field%set ( &
                  tex_name = parse_node_get_string &
                  (parse_node_get_sub_ptr (nd_prop, 2)))
          case ("prt_tex_anti")
             call field%set ( &
                  tex_anti = parse_node_get_string &
                  (parse_node_get_sub_ptr (nd_prop, 2)))
          case ("prt_spin")
             call field%set ( &
                  spin_type = read_frac &
                  (parse_node_get_sub_ptr (nd_prop, 2), 2))
          case ("prt_isospin")
             call field%set ( &
                  isospin_type = read_frac &
                  (parse_node_get_sub_ptr (nd_prop, 2), 2))
          case ("prt_charge")
             call field%set ( &
                  charge_type = read_frac &
                  (parse_node_get_sub_ptr (nd_prop, 2), 3))
          case ("prt_color")
             call field%set ( &
                  color_type = parse_node_get_integer &
                  (parse_node_get_sub_ptr (nd_prop, 2)))
          case ("prt_mass")
             call field%set ( &
                  mass_data = model%get_par_data_ptr &
                  (parse_node_get_string &
                  (parse_node_get_sub_ptr (nd_prop, 2))))
          case ("prt_width")
             call field%set ( &
                  width_data = model%get_par_data_ptr &
                  (parse_node_get_string &
                  (parse_node_get_sub_ptr (nd_prop, 2))))
          case default
             call msg_bug (" Unknown particle property '" &
                  // char (parse_node_get_rule_key (nd_prop)) // "'")
          end select
          if (allocated (name))  deallocate (name)
          nd_prop => parse_node_get_next_ptr (nd_prop)
       end do
    end if
    call field%freeze ()
  end subroutine model_read_field

\eatline
\nwendcode{}\nwbegincode{591}\moddef{Models: model: TBP}\plusendmoddef
  procedure, private :: read_vertex => model_read_vertex
\nwendcode{}\nwbegincode{592}\moddef{Models: procedures}\plusendmoddef
  subroutine model_read_vertex (model, i, node)
    class(model_t), intent(inout) :: model
    integer, intent(in) :: i
    type(parse_node_t), intent(in) :: node
    type(string_t), dimension(:), allocatable :: name
    call read_names (node, name)
    call model%set_vertex (i, name)
  end subroutine model_read_vertex

\eatline
\nwendcode{}\nwbegincode{593}\moddef{Models: procedures}\plusendmoddef
  subroutine read_names (node, name)
    type(parse_node_t), intent(in) :: node
    type(string_t), dimension(:), allocatable, intent(inout) :: name
    type(parse_node_t), pointer :: nd_name
    integer :: n_names, i
    n_names = parse_node_get_n_sub (node) - 1
    allocate (name (n_names))
    nd_name => parse_node_get_sub_ptr (node, 2)
    do i = 1, n_names
       name(i) = parse_node_get_string (nd_name)
       nd_name => parse_node_get_next_ptr (nd_name)
    end do
  end subroutine read_names

\eatline
\nwendcode{}\nwbegincode{594}\moddef{Models: procedures}\plusendmoddef
  function read_frac (nd_frac, base) result (qn_type)
    integer :: qn_type
    type(parse_node_t), intent(in) :: nd_frac
    integer, intent(in) :: base
    type(parse_node_t), pointer :: nd_num, nd_den
    integer :: num, den
    nd_num => parse_node_get_sub_ptr (nd_frac)
    nd_den => parse_node_get_next_ptr (nd_num)
    select case (char (parse_node_get_rule_key (nd_num)))
    case ("integer_literal")
       num = parse_node_get_integer (nd_num)
    case ("neg_int")
       num = - parse_node_get_integer (parse_node_get_sub_ptr (nd_num, 2))
    case ("pos_int")
       num = parse_node_get_integer (parse_node_get_sub_ptr (nd_num, 2))
    case default
       call parse_tree_bug (nd_num, "int|neg_int|pos_int")
    end select
    if (associated (nd_den)) then
       den = parse_node_get_integer (parse_node_get_sub_ptr (nd_den, 2))
    else
       den = 1
    end if
    if (den == 1) then
       qn_type = sign (1 + abs (num) * base, num)
    else if (den == base) then
       qn_type = sign (abs (num) + 1, num)
    else
       call parse_node_write_rec (nd_frac)
       call msg_fatal (" Fractional quantum number: wrong denominator")
    end if
  end function read_frac

\eatline
\nwendcode{}\nwbegindocs{595}Append field (PDG-array) variables to the variable list, based on
the field content.
\nwenddocs{}\nwbegincode{596}\moddef{Models: model: TBP}\plusendmoddef
  procedure, private :: append_field_vars => model_append_field_vars
\nwendcode{}\nwbegincode{597}\moddef{Models: procedures}\plusendmoddef
  subroutine model_append_field_vars (model)
    class(model_t), intent(inout) :: model
    type(pdg_array_t) :: aval
    type(field_data_t), dimension(:), pointer :: field_array
    type(field_data_t), pointer :: field
    type(string_t) :: name
    type(string_t), dimension(:), allocatable :: name_array
    integer, dimension(:), allocatable :: pdg
    logical, dimension(:), allocatable :: mask
    integer :: i, j
    field_array => model%get_field_array_ptr ()
    aval = UNDEFINED
    call var_list_append_pdg_array &
         (model%var_list, var_str ("particle"), &
          aval, locked = .true., intrinsic=.true.)
    do i = 1, size (field_array)
       aval = field_array(i)%get_pdg ()
       name = field_array(i)%get_longname ()
       call var_list_append_pdg_array &
            (model%var_list, name, aval, locked=.true., intrinsic=.true.)
       call field_array(i)%get_name_array (.false., name_array)
       do j = 1, size (name_array)
          call var_list_append_pdg_array &
               (model%var_list, name_array(j), &
               aval, locked=.true., intrinsic=.true.)
       end do
       model%max_field_name_length = &
            max (model%max_field_name_length, len (name_array(1)))
       aval = - field_array(i)%get_pdg ()
       call field_array(i)%get_name_array (.true., name_array)
       do j = 1, size (name_array)
          call var_list_append_pdg_array &
               (model%var_list, name_array(j), &
               aval, locked=.true., intrinsic=.true.)
       end do
       if (size (name_array) > 0) then
          model%max_field_name_length = &
               max (model%max_field_name_length, len (name_array(1)))
       end if
    end do
    call model%get_all_pdg (pdg)
    allocate (mask (size (pdg)))
    do i = 1, size (pdg)
       field => model%get_field_ptr (pdg(i))
       mask(i) = field%get_charge_type () /= 1
    end do
    aval = pack (pdg, mask)
    call var_list_append_pdg_array &
         (model%var_list, var_str ("charged"), &
          aval, locked = .true., intrinsic=.true.)
    do i = 1, size (pdg)
       field => model%get_field_ptr (pdg(i))
       mask(i) = field%get_color_type () /= 1
    end do
    aval = pack (pdg, mask)
    call var_list_append_pdg_array &
         (model%var_list, var_str ("colored"), &
          aval, locked = .true., intrinsic=.true.)
  end subroutine model_append_field_vars
    
\eatline
\nwendcode{}\nwbegindocs{598}\nwdocspar
\subsection{Model list}
List of currently active models
\nwenddocs{}\nwbegincode{599}\moddef{Models: types}\plusendmoddef
  type, extends (model_t) :: model_entry_t
     type(model_entry_t), pointer :: next => null ()
  end type model_entry_t

\eatline
\nwendcode{}\nwbegincode{600}\moddef{Models: public}\plusendmoddef
  public :: model_list_t
\nwendcode{}\nwbegincode{601}\moddef{Models: types}\plusendmoddef
  type :: model_list_t
     type(model_entry_t), pointer :: first => null ()
     type(model_entry_t), pointer :: last => null ()
     type(model_list_t), pointer :: context => null ()
   contains
   \LA{}Models: model list: TBP\RA{}
  end type model_list_t

\eatline
\nwendcode{}\nwbegindocs{602}Write an account of the model list.  We write linked lists first, starting
from the global context.
\nwenddocs{}\nwbegincode{603}\moddef{Models: model list: TBP}\endmoddef
  procedure :: write => model_list_write
\nwendcode{}\nwbegincode{604}\moddef{Models: procedures}\plusendmoddef
  recursive subroutine model_list_write (object, unit, verbose, follow_link)
    class(model_list_t), intent(in) :: object
    integer, intent(in), optional :: unit
    logical, intent(in), optional :: verbose
    logical, intent(in), optional :: follow_link
    type(model_entry_t), pointer :: current
    logical :: rec
    integer :: u
    u = given_output_unit (unit);  if (u < 0)  return
    rec = .true.;  if (present (follow_link))  rec = follow_link
    if (rec .and. associated (object%context)) then
       call object%context%write (unit, verbose, follow_link)
    end if
    current => object%first
    if (associated (current)) then
       do while (associated (current))
          call current%write (unit, verbose)
          current => current%next
          if (associated (current))  write (u, *)
       end do
    end if
  end subroutine model_list_write

\eatline
\nwendcode{}\nwbegindocs{605}Link this list to another one.
\nwenddocs{}\nwbegincode{606}\moddef{Models: model list: TBP}\plusendmoddef
  procedure :: link => model_list_link
\nwendcode{}\nwbegincode{607}\moddef{Models: procedures}\plusendmoddef
  subroutine model_list_link (model_list, context)
    class(model_list_t), intent(inout) :: model_list
    type(model_list_t), intent(in), target :: context
    model_list%context => context
  end subroutine model_list_link
  
\eatline
\nwendcode{}\nwbegindocs{608}(Private, used below:)
Append an existing model, for which we have allocated a pointer entry, to
the model list.  The original pointer becomes disassociated, and the model
should now be considered as part of the list.  We assume that this model is
not yet part of the list.  

If we provide a {\tt{}model} argument, this returns a pointer to the new entry.
\nwenddocs{}\nwbegincode{609}\moddef{Models: model list: TBP}\plusendmoddef
  procedure, private :: import => model_list_import
\nwendcode{}\nwbegincode{610}\moddef{Models: procedures}\plusendmoddef
  subroutine model_list_import (model_list, current, model)
    class(model_list_t), intent(inout) :: model_list
    type(model_entry_t), pointer, intent(inout) :: current
    type(model_t), optional, pointer, intent(out) :: model
    if (associated (current)) then
       if (associated (model_list%first)) then
          model_list%last%next => current
       else
          model_list%first => current
       end if
       model_list%last => current
       if (present (model))  model => current%model_t
       current => null ()
    end if
  end subroutine model_list_import
       
\eatline
\nwendcode{}\nwbegindocs{611}(Private, currently test only:) 
Add a new model with given {\tt{}name} to the list, if it does not yet
exist.  If successful, return a pointer to the new model.
\nwenddocs{}\nwbegincode{612}\moddef{Models: model list: TBP}\plusendmoddef
  procedure, private :: add => model_list_add
\nwendcode{}\nwbegincode{613}\moddef{Models: procedures}\plusendmoddef
  subroutine model_list_add (model_list, &
       name, os_data, n_par, n_prt, n_vtx, model)
    class(model_list_t), intent(inout) :: model_list
    type(string_t), intent(in) :: name
    type(os_data_t), intent(in) :: os_data
    integer, intent(in) :: n_par, n_prt, n_vtx
    type(model_t), pointer :: model
    type(model_entry_t), pointer :: current
    if (model_list%model_exists (name, follow_link=.false.)) then
       model => null ()
    else
       allocate (current)
       call current%init (name, var_str (""), os_data, &
            n_par, n_prt, n_vtx)
       call model_list%import (current, model)
    end if
  end subroutine model_list_add

\eatline
\nwendcode{}\nwbegindocs{614}Read a new model from file and add to the list, if it does not yet
exist.  Finalize the model by allocating the vertex table.  Return a
pointer to the new model.  If unsuccessful, return the original
pointer.

The model is always inserted in the last link of a chain of model lists.  This
way, we avoid loading models twice from different contexts.   When a model is
modified, we should first allocate a local copy.
\nwenddocs{}\nwbegincode{615}\moddef{Models: model list: TBP}\plusendmoddef
  procedure :: read_model => model_list_read_model
\nwendcode{}\nwbegincode{616}\moddef{Models: procedures}\plusendmoddef
  subroutine model_list_read_model (model_list, name, filename, os_data, model)
    class(model_list_t), intent(inout), target :: model_list
    type(string_t), intent(in) :: name, filename
    type(os_data_t), intent(in) :: os_data
    type(model_t), pointer, intent(inout) :: model
    class(model_list_t), pointer :: global_model_list
    type(model_entry_t), pointer :: current
    logical :: exist
    if (.not. model_list%model_exists (name, follow_link=.true.)) then
       allocate (current)
       call current%read (filename, os_data, exist)
       if (.not. exist)  return
       if (current%get_name () /= name) then
          call msg_fatal ("Model file '" // char (filename) // &
               "' contains model '" // char (current%get_name ()) // &
               "' instead of '" // char (name) // "'")
          call current%final ();  deallocate (current)
          return
       end if
       global_model_list => model_list
       do while (associated (global_model_list%context))
          global_model_list => global_model_list%context
       end do
       call global_model_list%import (current, model)
    else
       model => model_list%get_model_ptr (name)
    end if
  end subroutine model_list_read_model

\eatline
\nwendcode{}\nwbegindocs{617}Append a copy of an existing model to a model list.  Optionally, return
pointer to the new entry.
\nwenddocs{}\nwbegincode{618}\moddef{Models: model list: TBP}\plusendmoddef
  procedure :: append_copy => model_list_append_copy
\nwendcode{}\nwbegincode{619}\moddef{Models: procedures}\plusendmoddef
  subroutine model_list_append_copy (model_list, orig, model)
    class(model_list_t), intent(inout) :: model_list
    type(model_t), intent(in), target :: orig
    type(model_t), intent(out), pointer, optional :: model
    type(model_entry_t), pointer :: copy
    allocate (copy)
    call copy%init_instance (orig)
    call model_list%import (copy, model)
  end subroutine model_list_append_copy
    
\eatline
\nwendcode{}\nwbegindocs{620}Check if a model exists by examining the list.  Check recursively unless
told otherwise.
\nwenddocs{}\nwbegincode{621}\moddef{Models: model list: TBP}\plusendmoddef
  procedure :: model_exists => model_list_model_exists
\nwendcode{}\nwbegincode{622}\moddef{Models: procedures}\plusendmoddef
  recursive function model_list_model_exists (model_list, name, follow_link) &
       result (exists)
    class(model_list_t), intent(in) :: model_list
    logical :: exists
    type(string_t), intent(in) :: name
    logical, intent(in), optional :: follow_link
    type(model_entry_t), pointer :: current
    logical :: rec
    rec = .true.;  if (present (follow_link))  rec = follow_link
    current => model_list%first
    do while (associated (current))
       if (current%get_name () == name) then
          exists = .true.
          return
       end if
       current => current%next
    end do
    if (rec .and. associated (model_list%context)) then
       exists = model_list%context%model_exists (name, follow_link)
    else
       exists = .false.
    end if
  end function model_list_model_exists

\eatline
\nwendcode{}\nwbegindocs{623}Return a pointer to a named model.  Search recursively unless told otherwise.
\nwenddocs{}\nwbegincode{624}\moddef{Models: model list: TBP}\plusendmoddef
  procedure :: get_model_ptr => model_list_get_model_ptr
\nwendcode{}\nwbegincode{625}\moddef{Models: procedures}\plusendmoddef
  recursive function model_list_get_model_ptr (model_list, name, follow_link) &
       result (model)
    class(model_list_t), intent(in) :: model_list
    type(model_t), pointer :: model
    type(string_t), intent(in) :: name
    logical, intent(in), optional :: follow_link
    type(model_entry_t), pointer :: current
    logical :: rec
    rec = .true.;  if (present (follow_link))  rec = follow_link
    current => model_list%first
    do while (associated (current))
       if (current%get_name () == name) then
          model => current%model_t
          return
       end if
       current => current%next
    end do
    if (rec .and. associated (model_list%context)) then
       model => model_list%context%get_model_ptr (name, follow_link)
    else
       model => null ()
    end if
  end function model_list_get_model_ptr

\eatline
\nwendcode{}\nwbegindocs{626}Delete the list of models.  No recursion.
\nwenddocs{}\nwbegincode{627}\moddef{Models: model list: TBP}\plusendmoddef
  procedure :: final => model_list_final
\nwendcode{}\nwbegincode{628}\moddef{Models: procedures}\plusendmoddef
  subroutine model_list_final (model_list)
    class(model_list_t), intent(inout) :: model_list
    type(model_entry_t), pointer :: current
    model_list%last => null ()
    do while (associated (model_list%first))
       current => model_list%first
       model_list%first => model_list%first%next
       call current%final ()
       deallocate (current)
    end do
  end subroutine model_list_final

\eatline
\nwendcode{}\nwbegindocs{629}\nwdocspar
\subsection{Model instances}
A model instance is a copy of a model object.  The parameters are true
copies.  The particle data and the variable list pointers should point to the
copy, so modifying the parameters has only a local effect.  Hence, we build
them up explicitly.  The vertex array is also rebuilt, it contains particle
pointers.  Finally, the vertex hash table can be copied directly since it
contains no pointers.

The {\tt{}multiplicity} entry depends on the association of the {\tt{}mass{\char95}data}
entry and therefore has to be set at the end.

The instance must carry the {\tt{}target} attribute.

Parameters: the {\tt{}copy{\char95}parameter} method essentially copies the parameter
decorations (parse node, expression etc.).  The current parameter values are
part of the {\tt{}model{\char95}data{\char95}t} base type and are copied afterwards via its
{\tt{}copy{\char95}from} method.
\nwenddocs{}\nwbegincode{630}\moddef{Models: model: TBP}\plusendmoddef
  procedure :: init_instance => model_copy
\nwendcode{}\nwbegincode{631}\moddef{Models: procedures}\plusendmoddef
  subroutine model_copy (model, orig)
    class(model_t), intent(out), target :: model
    type(model_t), intent(in) :: orig
    integer :: n_par, n_prt, n_vtx
    integer :: i
    n_par = size (orig%par)
    n_prt = orig%get_n_field ()
    n_vtx = orig%get_n_vtx ()
    call model%basic_init (orig%get_name (), n_par, n_prt, n_vtx)
    model%md5sum = orig%md5sum
    do i = 1, n_par
       call model%copy_parameter (i, orig%par(i))
    end do
    model%init_external_parameters => orig%init_external_parameters
    call model%copy_from (orig)
    model%max_par_name_length = orig%max_par_name_length
    call model%append_field_vars ()
  end subroutine model_copy
  
\eatline
\nwendcode{}\nwbegindocs{632}(Obsolete:) Turn a pointer to a model into a true copy, i.e., a new instance.
Reset: delete this instance.
\nwenddocs{}\nwbegincode{633}\moddef{Models: public}\plusendmoddef
!   public :: model_pointer_to_instance
!   public :: model_pointer_delete_instance
\nwendcode{}\nwbegincode{634}\moddef{Models: procedures}\plusendmoddef
!   subroutine model_pointer_to_instance (model)
!     type(model_t), pointer, intent(inout) :: model
!     type(model_t), pointer :: model_tmp
!     model_tmp => model
!     allocate (model)
!     call model%init_instance (model_tmp)
!   end subroutine model_pointer_to_instance
    
!   subroutine model_pointer_delete_instance (model)
!     type(model_t), pointer, intent(inout) :: model
!     if (associated (model)) then
!        call model%final ()
!        deallocate (model)
!     end if
!   end subroutine model_pointer_delete_instance
    
\eatline
\eatline
\nwendcode{}\nwbegindocs{635}\nwdocspar
\subsection{Unit tests}
\nwenddocs{}\nwbegincode{636}\moddef{Models: public}\plusendmoddef
  public :: models_test
\nwendcode{}\nwbegincode{637}\moddef{Models: tests}\endmoddef
  subroutine models_test (u, results)
    integer, intent(in) :: u
    type(test_results_t), intent(inout) :: results
  \LA{}Models: execute tests\RA{}
  end subroutine models_test

\eatline
\nwendcode{}\nwbegindocs{638}\nwdocspar
\subsubsection{Construct a Model}
Here, we construct a toy model explicitly without referring to a file.
\nwenddocs{}\nwbegincode{639}\moddef{Models: execute tests}\endmoddef
  call test (models_1, "models_1", &
       "construct model", &
       u, results)
\nwendcode{}\nwbegincode{640}\moddef{Models: tests}\plusendmoddef
  subroutine models_1 (u)
    integer, intent(in) :: u
    type(os_data_t) :: os_data
    type(model_list_t) :: model_list
    type(model_t), pointer :: model
    type(string_t) :: model_name
    type(string_t) :: x_longname
    type(string_t), dimension(2) :: parname
    type(string_t), dimension(2) :: x_name
    type(string_t), dimension(1) :: x_anti
    type(string_t) :: x_tex_name, x_tex_anti
    type(string_t) :: y_longname
    type(string_t), dimension(2) :: y_name
    type(string_t) :: y_tex_name
    type(field_data_t), pointer :: field

    write (u, "(A)")  "* Test output: models_1"
    write (u, "(A)")  "*   Purpose: create a model"
    write (u, *)

    model_name = "Test model"
    call model_list%add (model_name, os_data, 2, 2, 3, model)
    parname(1) = "mx"
    parname(2) = "coup"
    call model%set_parameter_constant (1, parname(1), 10._default)
    call model%set_parameter_constant (2, parname(2), 1.3_default)
    x_longname = "X_LEPTON"
    x_name(1) = "X"
    x_name(2) = "x"
    x_anti(1) = "Xbar"
    x_tex_name = "X^+"
    x_tex_anti = "X^-"
    field => model%get_field_ptr_by_index (1)
    call field%init (x_longname, 99)
    call field%set ( &
         .true., .false., .false., .false., .false., &
         name=x_name, anti=x_anti, tex_name=x_tex_name, tex_anti=x_tex_anti, &
         spin_type=SPINOR, isospin_type=-3, charge_type=2, &
         mass_data=model%get_par_data_ptr (parname(1)))
    y_longname = "Y_COLORON"
    y_name(1) = "Y"
    y_name(2) = "yc"
    y_tex_name = "Y^0"
    field => model%get_field_ptr_by_index (2)
    call field%init (y_longname, 97)
    call field%set ( &
          .false., .false., .true., .false., .false., &
          name=y_name, tex_name=y_tex_name, &
          spin_type=SCALAR, isospin_type=2, charge_type=1, color_type=8)
    call model%set_vertex (1, [99, 99, 99])
    call model%set_vertex (2, [99, 99, 99, 99])
    call model%set_vertex (3, [99, 97, 99])
    call model_list%write (u)

    call model_list%final ()

    write (u, *)
    write (u, "(A)")  "* Test output end: models_1"

  end subroutine models_1

\eatline
\nwendcode{}\nwbegindocs{641}\nwdocspar
\subsubsection{Read a Model}
Read a predefined model from file.
\nwenddocs{}\nwbegincode{642}\moddef{Models: execute tests}\plusendmoddef
  call test (models_2, "models_2", &
       "read model", &
       u, results)
\nwendcode{}\nwbegincode{643}\moddef{Models: tests}\plusendmoddef
  subroutine models_2 (u)
    integer, intent(in) :: u
    type(os_data_t) :: os_data
    type(model_list_t) :: model_list
    type(model_t), pointer :: model

    write (u, "(A)")  "* Test output: models_2"
    write (u, "(A)")  "*   Purpose: read a model from file"
    write (u, *)

    call syntax_model_file_init ()
    call os_data_init (os_data)

    call model_list%read_model (var_str ("Test"), var_str ("Test.mdl"), &
         os_data, model)
    call model_list%write (u)
    
    write (u, *)
    write (u, "(A)")  "* Variable list"
    write (u, *)
    
    call var_list_write (model%var_list, u)

    write (u, *)
    write (u, "(A)")  "* Cleanup"
    
    call model_list%final ()
    call syntax_model_file_final ()

    write (u, *)
    write (u, "(A)")  "* Test output end: models_2"

  end subroutine models_2

\eatline
\nwendcode{}\nwbegindocs{644}\nwdocspar
\subsubsection{Model Instance}
Read a predefined model from file and create an instance.
\nwenddocs{}\nwbegincode{645}\moddef{Models: execute tests}\plusendmoddef
  call test (models_3, "models_3", &
       "model instance", &
       u, results)
\nwendcode{}\nwbegincode{646}\moddef{Models: tests}\plusendmoddef
  subroutine models_3 (u)
    integer, intent(in) :: u
    type(os_data_t) :: os_data
    type(model_list_t) :: model_list
    type(model_t), pointer :: model
    type(model_t), pointer :: instance

    write (u, "(A)")  "* Test output: models_3"
    write (u, "(A)")  "*   Purpose: create a model instance"
    write (u, *)

    call syntax_model_file_init ()
    call os_data_init (os_data)

    call model_list%read_model (var_str ("Test"), var_str ("Test.mdl"), &
         os_data, model)
    allocate (instance)
    call instance%init_instance (model)
    
    call model%write (u)
    
    write (u, *)
    write (u, "(A)")  "* Variable list"
    write (u, *)
    
    call var_list_write (instance%var_list, u)

    write (u, *)
    write (u, "(A)")  "* Cleanup"
    
    call instance%final ()
    deallocate (instance)
    
    call model_list%final ()
    call syntax_model_file_final ()

    write (u, *)
    write (u, "(A)")  "* Test output end: models_3"

  end subroutine models_3

\eatline
\nwendcode{}\nwbegindocs{647}\nwdocspar
\subsubsection{Unstable and Polarized Particles}
Read a predefined model from file and define decays and polarization.
\nwenddocs{}\nwbegincode{648}\moddef{Models: execute tests}\plusendmoddef
  call test (models_4, "models_4", &
       "handle decays and polarization", &
       u, results)
\nwendcode{}\nwbegincode{649}\moddef{Models: tests}\plusendmoddef
  subroutine models_4 (u)
    integer, intent(in) :: u
    type(os_data_t) :: os_data
    type(model_list_t) :: model_list
    type(model_t), pointer :: model, model_instance
    character(32) :: md5sum

    write (u, "(A)")  "* Test output: models_4"
    write (u, "(A)")  "*   Purpose: set and unset decays and polarization"
    write (u, *)

    call syntax_model_file_init ()
    call os_data_init (os_data)

    write (u, "(A)")  "* Read model from file"

    call model_list%read_model (var_str ("Test"), var_str ("Test.mdl"), &
         os_data, model)

    md5sum = model%get_parameters_md5sum ()
    write (u, *)
    write (u, "(1x,3A)")  "MD5 sum (parameters) = '", md5sum, "'"

    write (u, *)
    write (u, "(A)")  "* Set particle decays and polarization"
    write (u, *)

    call model%set_unstable (25, [var_str ("dec1"), var_str ("dec2")])
    call model%set_polarized (6)
    call model%set_unstable (-6, [var_str ("fdec")])

    call model%write (u)

    md5sum = model%get_parameters_md5sum ()
    write (u, *)
    write (u, "(1x,3A)")  "MD5 sum (parameters) = '", md5sum, "'"

    write (u, *)
    write (u, "(A)")  "* Create a model instance"

    allocate (model_instance)
    call model_instance%init_instance (model)

    write (u, *)
    write (u, "(A)")  "* Revert particle decays and polarization"
    write (u, *)

    call model%set_stable (25)
    call model%set_unpolarized (6)
    call model%set_stable (-6)

    call model%write (u)
    
    md5sum = model%get_parameters_md5sum ()
    write (u, *)
    write (u, "(1x,3A)")  "MD5 sum (parameters) = '", md5sum, "'"

    write (u, *)
    write (u, "(A)")  "* Show the model instance"
    write (u, *)

    call model_instance%write (u)

    md5sum = model_instance%get_parameters_md5sum ()
    write (u, *)
    write (u, "(1x,3A)")  "MD5 sum (parameters) = '", md5sum, "'"

    write (u, *)
    write (u, "(A)")  "* Cleanup"
    
    call model_instance%final ()
    deallocate (model_instance)
    call model_list%final ()
    call syntax_model_file_final ()

    write (u, *)
    write (u, "(A)")  "* Test output end: models_4"

  end subroutine models_4

\eatline
\nwendcode{}\nwbegindocs{650}\nwdocspar
\subsubsection{Model Variables}
Read a predefined model from file and modify some parameters.

Note that the MD5 sum is not modified by this.
\nwenddocs{}\nwbegincode{651}\moddef{Models: execute tests}\plusendmoddef
  call test (models_5, "models_5", &
       "handle parameters", &
       u, results)
\nwendcode{}\nwbegincode{652}\moddef{Models: tests}\plusendmoddef
  subroutine models_5 (u)
    integer, intent(in) :: u
    type(os_data_t) :: os_data
    type(model_list_t) :: model_list
    type(model_t), pointer :: model, model_instance
    character(32) :: md5sum

    write (u, "(A)")  "* Test output: models_5"
    write (u, "(A)")  "*   Purpose: access and modify model variables"
    write (u, *)

    call syntax_model_file_init ()
    call os_data_init (os_data)

    write (u, "(A)")  "* Read model from file"

    call model_list%read_model (var_str ("Test"), var_str ("Test.mdl"), &
         os_data, model)

    write (u, *)

    call model%write (u, &
         show_md5sum = .true., &
         show_variables = .true., &
         show_parameters = .true., &
         show_particles = .false., &
         show_vertices = .false.)

    write (u, *)
    write (u, "(A)")  "* Check parameter status"
    write (u, *)

    write (u, "(1x,A,L1)") "xy exists = ", model%var_exists (var_str ("xx"))
    write (u, "(1x,A,L1)") "ff exists = ", model%var_exists (var_str ("ff"))
    write (u, "(1x,A,L1)") "mf exists = ", model%var_exists (var_str ("mf"))
    write (u, "(1x,A,L1)") "ff locked = ", model%var_is_locked (var_str ("ff"))
    write (u, "(1x,A,L1)") "mf locked = ", model%var_is_locked (var_str ("mf"))

    write (u, *)
    write (u, "(1x,A,F6.2)") "ff = ", model%get_rval (var_str ("ff"))
    write (u, "(1x,A,F6.2)") "mf = ", model%get_rval (var_str ("mf"))

    write (u, *)
    write (u, "(A)")  "* Modify parameter"
    write (u, *)
    
    call model%set_real (var_str ("ff"), 1._default)

    call model%write (u, &
         show_md5sum = .true., &
         show_variables = .true., &
         show_parameters = .true., &
         show_particles = .false., &
         show_vertices = .false.)

    write (u, *)
    write (u, "(A)")  "* Cleanup"
    
    call model_list%final ()
    call syntax_model_file_final ()

    write (u, *)
    write (u, "(A)")  "* Test output end: models_5"

  end subroutine models_5

\eatline
\nwendcode{}\nwbegindocs{653}\nwdocspar
\clearpage
%%%%%%%%%%%%%%%%%%%%%%%%%%%%%%%%%%%%%%%%%%%%%%%%%%%%%%%%%%%%%%%%%%%%%%%%
\section{The SUSY Les Houches Accord}

The SUSY Les Houches Accord defines a standard interfaces for storing
the physics data of SUSY models.  Here, we provide the means
for reading, storing, and writing such data.
\nwenddocs{}\nwbegincode{654}\moddef{\code{}slha{\char95}interface.f90\edoc{}}\endmoddef
\LA{}File header\RA{}

module slha_interface

\LA{}Use kinds\RA{}
\LA{}Use strings\RA{}
  use io_units
  use constants
  use string_utils, only: upper_case
  use unit_tests
  use system_defs, only: VERSION_STRING
  use system_defs, only: EOF
  use diagnostics
  use os_interface
  use ifiles
  use lexers
  use syntax_rules
  use parser
  use variables
  use models

\LA{}Standard module head\RA{}

\LA{}SLHA: public\RA{}

\LA{}SLHA: parameters\RA{}

\LA{}SLHA: variables\RA{}

  save

contains

\LA{}SLHA: procedures\RA{}

\LA{}SLHA: tests\RA{}

end module slha_interface
\eatline
\nwendcode{}\nwbegindocs{655}\nwdocspar
\subsection{Preprocessor}
SLHA is a mixed-format standard.  It should be read in assuming free
format (but line-oriented), but it has some fixed-format elements.

To overcome this difficulty, we implement a preprocessing step which
transforms the SLHA into a format that can be swallowed by our generic
free-format lexer and parser.  Each line with a blank first character
is assumed to be a data line.  We prepend a 'DATA' keyword to these lines.
Furthermore, to enforce line-orientation, each line is appended a '\$'
key which is recognized by the parser.  To do this properly, we first
remove trailing comments, and skip lines consisting only of comments.

The preprocessor reads from a stream and puts out an {\tt{}ifile}.
Blocks that are not recognized are skipped.  For some blocks, data
items are quoted, so they can be read as strings if necessary.
\nwenddocs{}\nwbegincode{656}\moddef{SLHA: parameters}\endmoddef
  integer, parameter :: MODE_SKIP = 0, MODE_DATA = 1, MODE_INFO = 2

\eatline
\nwendcode{}\nwbegincode{657}\moddef{SLHA: procedures}\endmoddef
  subroutine slha_preprocess (stream, ifile)
    type(stream_t), intent(inout), target :: stream
    type(ifile_t), intent(out) :: ifile
    type(string_t) :: buffer, line, item
    integer :: iostat
    integer :: mode
    mode = MODE
    SCAN_FILE: do
       call stream_get_record (stream, buffer, iostat)
       select case (iostat)
       case (0)
          call split (buffer, line, "#")
          if (len_trim (line) == 0)  cycle SCAN_FILE
          select case (char (extract (line, 1, 1)))
          case ("B", "b")
             mode = check_block_handling (line)
             call ifile_append (ifile, line // "$")
          case ("D", "d")
             mode = MODE_DATA
             call ifile_append (ifile, line // "$")
          case (" ")
             select case (mode)
             case (MODE_DATA)
                call ifile_append (ifile, "DATA" // line // "$")
             case (MODE_INFO)
                line = adjustl (line)
                call split (line, item, " ")
                call ifile_append (ifile, "INFO" // " " // item // " " &
                     // '"' // trim (adjustl (line)) // '" $')
             end select
          case default
             call msg_message (char (line))
             call msg_fatal ("SLHA: Incomprehensible line")
          end select
       case (EOF)
          exit SCAN_FILE
       case default
          call msg_fatal ("SLHA: I/O error occured while reading SLHA input")
       end select
    end do SCAN_FILE
  end subroutine slha_preprocess

\eatline
\nwendcode{}\nwbegindocs{658}Return the mode that we should treat this block with.  We need
to recognize only those blocks that we actually use.
\nwenddocs{}\nwbegincode{659}\moddef{SLHA: procedures}\plusendmoddef
  function check_block_handling (line) result (mode)
    integer :: mode
    type(string_t), intent(in) :: line
    type(string_t) :: buffer, key, block_name
    buffer = trim (line)
    call split (buffer, key, " ")
    buffer = adjustl (buffer)
    call split (buffer, block_name, " ")
    block_name = trim (adjustl (upper_case (block_name)))
    select case (char (block_name))
    case ("MODSEL", "MINPAR", "SMINPUTS")
       mode = MODE_DATA
    case ("MASS")
       mode = MODE_DATA
    case ("NMIX", "UMIX", "VMIX", "STOPMIX", "SBOTMIX", "STAUMIX")
       mode = MODE_DATA
    case ("NMHMIX", "NMAMIX", "NMNMIX", "NMSSMRUN")
       mode = MODE_DATA
    case ("ALPHA", "HMIX")
       mode = MODE_DATA
    case ("AU", "AD", "AE")
       mode = MODE_DATA
    case ("SPINFO", "DCINFO")
       mode = MODE_INFO
    case default
       mode = MODE_SKIP
    end select
  end function check_block_handling

\eatline
\nwendcode{}\nwbegindocs{660}\nwdocspar
\subsection{Lexer and syntax}
\nwenddocs{}\nwbegincode{661}\moddef{SLHA: variables}\endmoddef
  type(syntax_t), target :: syntax_slha

\eatline
\nwendcode{}\nwbegincode{662}\moddef{SLHA: public}\endmoddef
  public :: syntax_slha_init
\nwendcode{}\nwbegincode{663}\moddef{SLHA: procedures}\plusendmoddef
  subroutine syntax_slha_init ()
    type(ifile_t) :: ifile
    call define_slha_syntax (ifile)
    call syntax_init (syntax_slha, ifile)
    call ifile_final (ifile)
  end subroutine syntax_slha_init

\eatline
\nwendcode{}\nwbegincode{664}\moddef{SLHA: public}\plusendmoddef
  public :: syntax_slha_final
\nwendcode{}\nwbegincode{665}\moddef{SLHA: procedures}\plusendmoddef
  subroutine syntax_slha_final ()
    call syntax_final (syntax_slha)
  end subroutine syntax_slha_final

\eatline
\nwendcode{}\nwbegincode{666}\moddef{SLHA: public}\plusendmoddef
  public :: syntax_slha_write
\nwendcode{}\nwbegincode{667}\moddef{SLHA: procedures}\plusendmoddef
  subroutine syntax_slha_write (unit)
    integer, intent(in), optional :: unit
    call syntax_write (syntax_slha, unit)
  end subroutine syntax_slha_write

\eatline
\nwendcode{}\nwbegincode{668}\moddef{SLHA: procedures}\plusendmoddef
  subroutine define_slha_syntax (ifile)
    type(ifile_t), intent(inout) :: ifile
    call ifile_append (ifile, "SEQ slha = chunk*")
    call ifile_append (ifile, "ALT chunk = block_def | decay_def")
    call ifile_append (ifile, "SEQ block_def = " &
         // "BLOCK block_spec '$' block_line*")
    call ifile_append (ifile, "KEY BLOCK")
    call ifile_append (ifile, "SEQ block_spec = block_name qvalue?")
    call ifile_append (ifile, "IDE block_name")
    call ifile_append (ifile, "SEQ qvalue = qname '=' real")
    call ifile_append (ifile, "IDE qname")
    call ifile_append (ifile, "KEY '='")
    call ifile_append (ifile, "REA real")
    call ifile_append (ifile, "KEY '$'")
    call ifile_append (ifile, "ALT block_line = block_data | block_info")
    call ifile_append (ifile, "SEQ block_data = DATA data_line '$'")
    call ifile_append (ifile, "KEY DATA")
    call ifile_append (ifile, "SEQ data_line = data_item+")
    call ifile_append (ifile, "ALT data_item = signed_number | number")
    call ifile_append (ifile, "SEQ signed_number = sign number")
    call ifile_append (ifile, "ALT sign = '+' | '-'")
    call ifile_append (ifile, "ALT number = integer | real")
    call ifile_append (ifile, "INT integer")
    call ifile_append (ifile, "KEY '-'")
    call ifile_append (ifile, "KEY '+'")
    call ifile_append (ifile, "SEQ block_info = INFO info_line '$'")
    call ifile_append (ifile, "KEY INFO")
    call ifile_append (ifile, "SEQ info_line = integer string_literal")
    call ifile_append (ifile, "QUO string_literal = '""'...'""'")
    call ifile_append (ifile, "SEQ decay_def = " &
         // "DECAY decay_spec '$' decay_data*")
    call ifile_append (ifile, "KEY DECAY")
    call ifile_append (ifile, "SEQ decay_spec = pdg_code data_item")
    call ifile_append (ifile, "ALT pdg_code = signed_integer | integer")
    call ifile_append (ifile, "SEQ signed_integer = sign integer")
    call ifile_append (ifile, "SEQ decay_data = DATA decay_line '$'")
    call ifile_append (ifile, "SEQ decay_line = data_item integer pdg_code+")
  end subroutine define_slha_syntax

\eatline
\nwendcode{}\nwbegindocs{669}The SLHA specification allows for string data items in certain
places.  Currently, we do not interpret them, but the strings, which
are not quoted, must be parsed somehow.  The hack for this problem is
to allow essentially all characters as special characters, so the
string can be read before it is discarded.
\nwenddocs{}\nwbegincode{670}\moddef{SLHA: public}\plusendmoddef
  public :: lexer_init_slha 
\nwendcode{}\nwbegincode{671}\moddef{SLHA: procedures}\plusendmoddef
  subroutine lexer_init_slha (lexer)
    type(lexer_t), intent(out) :: lexer
    call lexer_init (lexer, &
         comment_chars = "#", &
         quote_chars = '"', &
         quote_match = '"', &
         single_chars = "+-=$", &
         special_class = [ "" ], &
         keyword_list = syntax_get_keyword_list_ptr (syntax_slha), &
         upper_case_keywords = .true.)  ! $
  end subroutine lexer_init_slha

\eatline
\nwendcode{}\nwbegindocs{672}\nwdocspar
\subsection{Interpreter}
\subsubsection{Find blocks}
From the parse tree, find the node that represents a particular
block.  If {\tt{}required} is true, issue an error if not found.
Since {\tt{}block{\char95}name} is always invoked with capital letters, we 
have to capitalize {\tt{}pn{\char95}block{\char95}name}. 
\nwenddocs{}\nwbegincode{673}\moddef{SLHA: procedures}\plusendmoddef
  function slha_get_block_ptr &
       (parse_tree, block_name, required) result (pn_block)
    type(parse_node_t), pointer :: pn_block
    type(parse_tree_t), intent(in) :: parse_tree
    type(string_t), intent(in) :: block_name
    logical, intent(in) :: required
    type(parse_node_t), pointer :: pn_root, pn_block_spec, pn_block_name
    pn_root => parse_tree_get_root_ptr (parse_tree)
    pn_block => parse_node_get_sub_ptr (pn_root)
    do while (associated (pn_block))
       select case (char (parse_node_get_rule_key (pn_block)))
       case ("block_def")
          pn_block_spec => parse_node_get_sub_ptr (pn_block, 2)
          pn_block_name => parse_node_get_sub_ptr (pn_block_spec)
          if (trim (adjustl (upper_case (parse_node_get_string &
                   (pn_block_name)))) == block_name) then
             return
          end if
       end select
       pn_block => parse_node_get_next_ptr (pn_block)
    end do
    if (required) then
       call msg_fatal ("SLHA: block '" // char (block_name) // "' not found")
    end if
  end function slha_get_block_ptr

\eatline
\nwendcode{}\nwbegindocs{674}Scan the file for the first/next DECAY block.
\nwenddocs{}\nwbegincode{675}\moddef{SLHA: procedures}\plusendmoddef
  function slha_get_first_decay_ptr (parse_tree) result (pn_decay)
    type(parse_node_t), pointer :: pn_decay
    type(parse_tree_t), intent(in) :: parse_tree
    type(parse_node_t), pointer :: pn_root
    pn_root => parse_tree_get_root_ptr (parse_tree)
    pn_decay => parse_node_get_sub_ptr (pn_root)
    do while (associated (pn_decay))
       select case (char (parse_node_get_rule_key (pn_decay)))
       case ("decay_def")
          return
       end select
       pn_decay => parse_node_get_next_ptr (pn_decay)
    end do
  end function slha_get_first_decay_ptr

  function slha_get_next_decay_ptr (pn_block) result (pn_decay)
    type(parse_node_t), pointer :: pn_decay
    type(parse_node_t), intent(in), target :: pn_block
    pn_decay => parse_node_get_next_ptr (pn_block)
    do while (associated (pn_decay))
       select case (char (parse_node_get_rule_key (pn_decay)))
       case ("decay_def")
          return
       end select
       pn_decay => parse_node_get_next_ptr (pn_decay)
    end do
  end function slha_get_next_decay_ptr

\eatline
\nwendcode{}\nwbegindocs{676}\nwdocspar
\subsubsection{Extract and transfer data from blocks}
Given the parse node of a block, find the parse node of a particular
switch or data line.  Return this node and the node of the data item
following the integer code.
\nwenddocs{}\nwbegincode{677}\moddef{SLHA: procedures}\plusendmoddef
  subroutine slha_find_index_ptr (pn_block, pn_data, pn_item, code)
    type(parse_node_t), intent(in), target :: pn_block
    type(parse_node_t), intent(out), pointer :: pn_data
    type(parse_node_t), intent(out), pointer :: pn_item
    integer, intent(in) :: code
    pn_data => parse_node_get_sub_ptr (pn_block, 4)
    call slha_next_index_ptr (pn_data, pn_item, code)
  end subroutine slha_find_index_ptr

  subroutine slha_find_index_pair_ptr (pn_block, pn_data, pn_item, code1, code2)
    type(parse_node_t), intent(in), target :: pn_block
    type(parse_node_t), intent(out), pointer :: pn_data
    type(parse_node_t), intent(out), pointer :: pn_item
    integer, intent(in) :: code1, code2
    pn_data => parse_node_get_sub_ptr (pn_block, 4)
    call slha_next_index_pair_ptr (pn_data, pn_item, code1, code2)
  end subroutine slha_find_index_pair_ptr

\eatline
\nwendcode{}\nwbegindocs{678}Starting from the pointer to a data line, find a data line with the
given integer code.
\nwenddocs{}\nwbegincode{679}\moddef{SLHA: procedures}\plusendmoddef
  subroutine slha_next_index_ptr (pn_data, pn_item, code)
    type(parse_node_t), intent(inout), pointer :: pn_data
    integer, intent(in) :: code
    type(parse_node_t), intent(out), pointer :: pn_item
    type(parse_node_t), pointer :: pn_line, pn_code
    do while (associated (pn_data))
       pn_line => parse_node_get_sub_ptr (pn_data, 2)
       pn_code => parse_node_get_sub_ptr (pn_line)
       select case (char (parse_node_get_rule_key (pn_code)))
       case ("integer")
          if (parse_node_get_integer (pn_code) == code) then
             pn_item => parse_node_get_next_ptr (pn_code)
             return
          end if
       end select
       pn_data => parse_node_get_next_ptr (pn_data)
    end do
    pn_item => null ()
  end subroutine slha_next_index_ptr

\eatline
\nwendcode{}\nwbegindocs{680}Starting from the pointer to a data line, find a data line with the
given integer code pair.
\nwenddocs{}\nwbegincode{681}\moddef{SLHA: procedures}\plusendmoddef
  subroutine slha_next_index_pair_ptr (pn_data, pn_item, code1, code2)
    type(parse_node_t), intent(inout), pointer :: pn_data
    integer, intent(in) :: code1, code2
    type(parse_node_t), intent(out), pointer :: pn_item
    type(parse_node_t), pointer :: pn_line, pn_code1, pn_code2
    do while (associated (pn_data))
       pn_line => parse_node_get_sub_ptr (pn_data, 2)
       pn_code1 => parse_node_get_sub_ptr (pn_line)
       select case (char (parse_node_get_rule_key (pn_code1)))
       case ("integer")
          if (parse_node_get_integer (pn_code1) == code1) then
             pn_code2 => parse_node_get_next_ptr (pn_code1)
             if (associated (pn_code2)) then
                select case (char (parse_node_get_rule_key (pn_code2)))
                case ("integer")
                   if (parse_node_get_integer (pn_code2) == code2) then
                      pn_item => parse_node_get_next_ptr (pn_code2)
                      return
                   end if
                end select
             end if
          end if
       end select
       pn_data => parse_node_get_next_ptr (pn_data)
    end do
    pn_item => null ()
  end subroutine slha_next_index_pair_ptr

\eatline
\nwendcode{}\nwbegindocs{682}\nwdocspar
\subsubsection{Handle info data}
Return all strings with index {\tt{}i}.  The result is an allocated
string array.  Since we do not know the number of matching entries in
advance, we build an intermediate list which is transferred to the
final array and deleted before exiting.
\nwenddocs{}\nwbegincode{683}\moddef{SLHA: procedures}\plusendmoddef
  subroutine retrieve_strings_in_block (pn_block, code, str_array)
    type(parse_node_t), intent(in), target :: pn_block
    integer, intent(in) :: code
    type(string_t), dimension(:), allocatable, intent(out) :: str_array
    type(parse_node_t), pointer :: pn_data, pn_item
    type :: str_entry_t
       type(string_t) :: str
       type(str_entry_t), pointer :: next => null ()
    end type str_entry_t
    type(str_entry_t), pointer :: first => null ()
    type(str_entry_t), pointer :: current => null ()
    integer :: n
    n = 0
    call slha_find_index_ptr (pn_block, pn_data, pn_item, code)
    if (associated (pn_item)) then
       n = n + 1
       allocate (first)
       first%str = parse_node_get_string (pn_item)
       current => first
       do while (associated (pn_data))
          pn_data => parse_node_get_next_ptr (pn_data)
          call slha_next_index_ptr (pn_data, pn_item, code)
          if (associated (pn_item)) then
             n = n + 1
             allocate (current%next)
             current => current%next
             current%str = parse_node_get_string (pn_item)
          end if
       end do
       allocate (str_array (n))
       n = 0
       do while (associated (first))
          n = n + 1
          current => first
          str_array(n) = current%str
          first => first%next
          deallocate (current)
       end do
    else
       allocate (str_array (0))
    end if
  end subroutine retrieve_strings_in_block

\eatline
\nwendcode{}\nwbegindocs{684}\nwdocspar
\subsubsection{Transfer data from SLHA to variables}
Extract real parameter with index {\tt{}i}.  If it does not
exist, retrieve it from the variable list, using the given name.
\nwenddocs{}\nwbegincode{685}\moddef{SLHA: procedures}\plusendmoddef
  function get_parameter_in_block (pn_block, code, name, var_list) result (var)
    real(default) :: var
    type(parse_node_t), intent(in), target :: pn_block
    integer, intent(in) :: code
    type(string_t), intent(in) :: name
    type(var_list_t), intent(in), target :: var_list
    type(parse_node_t), pointer :: pn_data, pn_item
    call slha_find_index_ptr (pn_block, pn_data, pn_item, code)
    if (associated (pn_item)) then
       var = get_real_parameter (pn_item)
    else
       var = var_list%get_rval (name)
    end if
  end function get_parameter_in_block

\eatline
\nwendcode{}\nwbegindocs{686}Extract a real data item with index {\tt{}i}.  If it
does exist, set it in the variable list, using the given name.  If
the variable is not present in the variable list, ignore it.
\nwenddocs{}\nwbegincode{687}\moddef{SLHA: procedures}\plusendmoddef
  subroutine set_data_item (pn_block, code, name, var_list)
    type(parse_node_t), intent(in), target :: pn_block
    integer, intent(in) :: code
    type(string_t), intent(in) :: name
    type(var_list_t), intent(inout), target :: var_list
    type(parse_node_t), pointer :: pn_data, pn_item
    call slha_find_index_ptr (pn_block, pn_data, pn_item, code)
    if (associated (pn_item)) then
       call var_list_set_real &
            (var_list, name,  get_real_parameter (pn_item), &
             is_known=.true., ignore=.true.)
    end if
  end subroutine set_data_item

\eatline
\nwendcode{}\nwbegindocs{688}Extract a real matrix element with index {\tt{}i,j}.  If it
does exists, set it in the variable list, using the given name.  If
the variable is not present in the variable list, ignore it.
\nwenddocs{}\nwbegincode{689}\moddef{SLHA: procedures}\plusendmoddef
  subroutine set_matrix_element (pn_block, code1, code2, name, var_list)
    type(parse_node_t), intent(in), target :: pn_block
    integer, intent(in) :: code1, code2
    type(string_t), intent(in) :: name
    type(var_list_t), intent(inout), target :: var_list
    type(parse_node_t), pointer :: pn_data, pn_item
    call slha_find_index_pair_ptr (pn_block, pn_data, pn_item, code1, code2)
    if (associated (pn_item)) then
       call var_list_set_real &
            (var_list, name,  get_real_parameter (pn_item), &
             is_known=.true., ignore=.true.)
    end if
  end subroutine set_matrix_element

\eatline
\nwendcode{}\nwbegindocs{690}\nwdocspar
\subsubsection{Transfer data from variables to SLHA}
Get a real/integer parameter with index {\tt{}i} from the variable list and write
it to the current output file.  In the integer case, we account for
the fact that the variable is type real.  If it does not exist, do nothing.
\nwenddocs{}\nwbegincode{691}\moddef{SLHA: procedures}\plusendmoddef
  subroutine write_integer_data_item (u, code, name, var_list, comment)
    integer, intent(in) :: u
    integer, intent(in) :: code
    type(string_t), intent(in) :: name
    type(var_list_t), intent(in) :: var_list
    character(*), intent(in) :: comment
    integer :: item
    if (var_list%contains (name)) then
       item = nint (var_list%get_rval (name))
       call write_integer_parameter (u, code, item, comment)
    end if
  end subroutine write_integer_data_item

  subroutine write_real_data_item (u, code, name, var_list, comment)
    integer, intent(in) :: u
    integer, intent(in) :: code
    type(string_t), intent(in) :: name
    type(var_list_t), intent(in) :: var_list
    character(*), intent(in) :: comment
    real(default) :: item
    if (var_list%contains (name)) then
       item = var_list%get_rval (name)
       call write_real_parameter (u, code, item, comment)
    end if
  end subroutine write_real_data_item

\eatline
\nwendcode{}\nwbegindocs{692}Get a real data item with two integer indices from the variable list
and write it to the current output file.  If it does not exist, do
nothing.
\nwenddocs{}\nwbegincode{693}\moddef{SLHA: procedures}\plusendmoddef
  subroutine write_matrix_element (u, code1, code2, name, var_list, comment)
    integer, intent(in) :: u
    integer, intent(in) :: code1, code2
    type(string_t), intent(in) :: name
    type(var_list_t), intent(in) :: var_list
    character(*), intent(in) :: comment
    real(default) :: item
    if (var_list%contains (name)) then
       item = var_list%get_rval (name)
       call write_real_matrix_element (u, code1, code2, item, comment)
    end if
  end subroutine write_matrix_element

\eatline
\nwendcode{}\nwbegindocs{694}\nwdocspar
\subsection{Auxiliary function}
Write a block header.
\nwenddocs{}\nwbegincode{695}\moddef{SLHA: procedures}\plusendmoddef
  subroutine write_block_header (u, name, comment)
    integer, intent(in) :: u
    character(*), intent(in) :: name, comment
    write (u, "(A,1x,A,3x,'#',1x,A)")  "BLOCK", name, comment
  end subroutine write_block_header

\eatline
\nwendcode{}\nwbegindocs{696}Extract a real parameter that may be defined real or
integer, signed or unsigned.
\nwenddocs{}\nwbegincode{697}\moddef{SLHA: procedures}\plusendmoddef
  function get_real_parameter (pn_item) result (var)
    real(default) :: var
    type(parse_node_t), intent(in), target :: pn_item
    type(parse_node_t), pointer :: pn_sign, pn_var
    integer :: sign
    select case (char (parse_node_get_rule_key (pn_item)))
    case ("signed_number")
       pn_sign => parse_node_get_sub_ptr (pn_item)
       pn_var  => parse_node_get_next_ptr (pn_sign)
       select case (char (parse_node_get_key (pn_sign)))
       case ("+");  sign = +1
       case ("-");  sign = -1
       end select
    case default
       sign = +1
       pn_var => pn_item
    end select
    select case (char (parse_node_get_rule_key (pn_var)))
    case ("integer");  var = sign * parse_node_get_integer (pn_var)
    case ("real");     var = sign * parse_node_get_real (pn_var)
    end select
  end function get_real_parameter

\eatline
\nwendcode{}\nwbegindocs{698}Auxiliary: Extract an integer parameter that may be defined signed
or unsigned.  A real value is an error.
\nwenddocs{}\nwbegincode{699}\moddef{SLHA: procedures}\plusendmoddef
  function get_integer_parameter (pn_item) result (var)
    integer :: var
    type(parse_node_t), intent(in), target :: pn_item
    type(parse_node_t), pointer :: pn_sign, pn_var
    integer :: sign
    select case (char (parse_node_get_rule_key (pn_item)))
    case ("signed_integer")
       pn_sign => parse_node_get_sub_ptr (pn_item)
       pn_var  => parse_node_get_next_ptr (pn_sign)
       select case (char (parse_node_get_key (pn_sign)))
       case ("+");  sign = +1
       case ("-");  sign = -1
       end select
    case ("integer")
       sign = +1
       pn_var => pn_item
    case default
       call parse_node_write (pn_var)
       call msg_error ("SLHA: Integer parameter expected")
       var = 0
       return
    end select
    var = sign * parse_node_get_integer (pn_var)
  end function get_integer_parameter

\eatline
\nwendcode{}\nwbegindocs{700}Write an integer parameter with a single index directly to file,
using the required output format.
\nwenddocs{}\nwbegincode{701}\moddef{SLHA: procedures}\plusendmoddef
  subroutine write_integer_parameter (u, code, item, comment)
    integer, intent(in) :: u
    integer, intent(in) :: code
    integer, intent(in) :: item
    character(*), intent(in) :: comment
1   format (1x, I9, 3x, 3x, I9, 4x, 3x, '#', 1x, A)
    write (u, 1)  code, item, comment
  end subroutine write_integer_parameter

\eatline
\nwendcode{}\nwbegindocs{702}Write a real parameter with two indices directly to file,
using the required output format.
\nwenddocs{}\nwbegincode{703}\moddef{SLHA: procedures}\plusendmoddef
  subroutine write_real_parameter (u, code, item, comment)
    integer, intent(in) :: u
    integer, intent(in) :: code
    real(default), intent(in) :: item
    character(*), intent(in) :: comment
1   format (1x, I9, 3x, 1P, E16.8, 0P, 3x, '#', 1x, A)
    write (u, 1)  code, item, comment
  end subroutine write_real_parameter

\eatline
\nwendcode{}\nwbegindocs{704}Write a real parameter with a single index directly to file,
using the required output format.
\nwenddocs{}\nwbegincode{705}\moddef{SLHA: procedures}\plusendmoddef
  subroutine write_real_matrix_element (u, code1, code2, item, comment)
    integer, intent(in) :: u
    integer, intent(in) :: code1, code2
    real(default), intent(in) :: item
    character(*), intent(in) :: comment
1   format (1x, I2, 1x, I2, 3x, 1P, E16.8, 0P, 3x, '#', 1x, A)
    write (u, 1)  code1, code2, item, comment
  end subroutine write_real_matrix_element

\eatline
\nwendcode{}\nwbegindocs{706}\nwdocspar
\subsubsection{The concrete SLHA interpreter}
SLHA codes for particular physics models
\nwenddocs{}\nwbegincode{707}\moddef{SLHA: parameters}\plusendmoddef
  integer, parameter :: MDL_MSSM = 0
  integer, parameter :: MDL_NMSSM = 1
\eatline
\nwendcode{}\nwbegindocs{708}Take the parse tree and extract relevant data.  Select the correct
model and store all data that is present in the appropriate variable
list.  Finally, update the variable record.
\nwenddocs{}\nwbegincode{709}\moddef{SLHA: procedures}\plusendmoddef
  subroutine slha_interpret_parse_tree &
       (parse_tree, model, input, spectrum, decays)
    type(parse_tree_t), intent(in) :: parse_tree
    type(model_t), intent(inout), target :: model
    logical, intent(in) :: input, spectrum, decays
    logical :: errors
    integer :: mssm_type
    call slha_handle_MODSEL (parse_tree, model, mssm_type)
    if (input) then
       call slha_handle_SMINPUTS (parse_tree, model)
       call slha_handle_MINPAR (parse_tree, model, mssm_type)
    end if
    if (spectrum) then
       call slha_handle_info_block (parse_tree, "SPINFO", errors)
       if (errors)  return
       call slha_handle_MASS (parse_tree, model)
       call slha_handle_matrix_block (parse_tree, "NMIX", "mn_", 4, 4, model)
       call slha_handle_matrix_block (parse_tree, "NMNMIX", "mixn_", 5, 5, model)
       call slha_handle_matrix_block (parse_tree, "UMIX", "mu_", 2, 2, model)
       call slha_handle_matrix_block (parse_tree, "VMIX", "mv_", 2, 2, model)
       call slha_handle_matrix_block (parse_tree, "STOPMIX", "mt_", 2, 2, model)
       call slha_handle_matrix_block (parse_tree, "SBOTMIX", "mb_", 2, 2, model)
       call slha_handle_matrix_block (parse_tree, "STAUMIX", "ml_", 2, 2, model)
       call slha_handle_matrix_block (parse_tree, "NMHMIX", "mixh0_", 3, 3, model)
       call slha_handle_matrix_block (parse_tree, "NMAMIX", "mixa0_", 2, 3, model)
       call slha_handle_ALPHA (parse_tree, model)
       call slha_handle_HMIX (parse_tree, model)
       call slha_handle_NMSSMRUN (parse_tree, model)
       call slha_handle_matrix_block (parse_tree, "AU", "Au_", 3, 3, model)
       call slha_handle_matrix_block (parse_tree, "AD", "Ad_", 3, 3, model)
       call slha_handle_matrix_block (parse_tree, "AE", "Ae_", 3, 3, model)
    end if
    if (decays) then
       call slha_handle_info_block (parse_tree, "DCINFO", errors)
       if (errors)  return
       call slha_handle_decays (parse_tree, model)
    end if
  end subroutine slha_interpret_parse_tree

\eatline
\nwendcode{}\nwbegindocs{710}\nwdocspar
\subsubsection{Info blocks}
Handle the informational blocks SPINFO and DCINFO.  The first two
items are program name and version.  Items with index 3 are warnings.
Items with index 4 are errors.  We reproduce these as WHIZARD warnings
and errors.
\nwenddocs{}\nwbegincode{711}\moddef{SLHA: procedures}\plusendmoddef
  subroutine slha_handle_info_block (parse_tree, block_name, errors)
    type(parse_tree_t), intent(in) :: parse_tree
    character(*), intent(in) :: block_name
    logical, intent(out) :: errors
    type(parse_node_t), pointer :: pn_block
    type(string_t), dimension(:), allocatable :: msg
    integer :: i
    pn_block => slha_get_block_ptr &
         (parse_tree, var_str (block_name), required=.true.)
    if (.not. associated (pn_block)) then
       call msg_error ("SLHA: Missing info block '" &
            // trim (block_name) // "'; ignored.")
       errors = .true.
       return
    end if
    select case (block_name)
    case ("SPINFO")
       call msg_message ("SLHA: SUSY spectrum program info:")
    case ("DCINFO")
       call msg_message ("SLHA: SUSY decay program info:")
    end select
    call retrieve_strings_in_block (pn_block, 1, msg)
    do i = 1, size (msg)
       call msg_message ("SLHA: " // char (msg(i)))
    end do
    call retrieve_strings_in_block (pn_block, 2, msg)
    do i = 1, size (msg)
       call msg_message ("SLHA: " // char (msg(i)))
    end do
    call retrieve_strings_in_block (pn_block, 3, msg)
    do i = 1, size (msg)
       call msg_warning ("SLHA: " // char (msg(i)))
    end do
    call retrieve_strings_in_block (pn_block, 4, msg)
    do i = 1, size (msg)
       call msg_error ("SLHA: " // char (msg(i)))
    end do
    errors = size (msg) > 0
  end subroutine slha_handle_info_block

\eatline
\nwendcode{}\nwbegindocs{712}\nwdocspar
\subsubsection{MODSEL}
Handle the overall model definition.  Only certain models are
recognized.  The soft-breaking model templates that determine the set
of input parameters:
\nwenddocs{}\nwbegincode{713}\moddef{SLHA: parameters}\plusendmoddef
  integer, parameter :: MSSM_GENERIC = 0
  integer, parameter :: MSSM_SUGRA = 1
  integer, parameter :: MSSM_GMSB = 2
  integer, parameter :: MSSM_AMSB = 3

\eatline
\nwendcode{}\nwbegincode{714}\moddef{SLHA: procedures}\plusendmoddef
  subroutine slha_handle_MODSEL (parse_tree, model, mssm_type)
    type(parse_tree_t), intent(in) :: parse_tree
    type(model_t), intent(in), target :: model
    integer, intent(out) :: mssm_type
    type(parse_node_t), pointer :: pn_block, pn_data, pn_item
    type(string_t) :: model_name
    pn_block => slha_get_block_ptr &
         (parse_tree, var_str ("MODSEL"), required=.true.)
    call slha_find_index_ptr (pn_block, pn_data, pn_item, 1)
    if (associated (pn_item)) then
       mssm_type = get_integer_parameter (pn_item)
    else
       mssm_type = MSSM_GENERIC
    end if
    call slha_find_index_ptr (pn_block, pn_data, pn_item, 3)
    if (associated (pn_item)) then
       select case (parse_node_get_integer (pn_item))
       case (MDL_MSSM);  model_name = "MSSM"
       case (MDL_NMSSM); model_name = "NMSSM"
       case default
          call msg_fatal ("SLHA: unknown model code in MODSEL")
          return
       end select
    else
       model_name = "MSSM"
    end if
    call slha_find_index_ptr (pn_block, pn_data, pn_item, 4)
    if (associated (pn_item)) then
      call msg_fatal (" R-parity violation is currently not supported by WHIZARD.")     
    end if
    call slha_find_index_ptr (pn_block, pn_data, pn_item, 5)
    if (associated (pn_item)) then
      call msg_fatal (" CP violation is currently not supported by WHIZARD.")   
    end if
    select case (char (model_name))
    case ("MSSM")
       select case (char (model%get_name ()))
       case ("MSSM","MSSM_CKM","MSSM_Grav","MSSM_Hgg")
          model_name = model%get_name ()
       case default
          call msg_fatal ("Selected model '" &
               // char (model%get_name ()) // "' does not match model '" &
               // char (model_name) // "' in SLHA input file.") 
          return
       end select
    case ("NMSSM")       
       select case (char (model%get_name ()))
       case ("NMSSM","NMSSM_CKM","NMSSM_Hgg")
          model_name = model%get_name ()
       case default
          call msg_fatal ("Selected model '" &
               // char (model%get_name ()) // "' does not match model '" &
               // char (model_name) // "' in SLHA input file.") 
          return
       end select
    case default
       call msg_bug ("SLHA model name '" &
            // char (model_name) // "' not recognized.")
       return
    end select
    call msg_message ("SLHA: Initializing model '" // char (model_name) // "'")
  end subroutine slha_handle_MODSEL

\eatline
\nwendcode{}\nwbegindocs{715}Write a MODSEL block, based on the contents of the current model.
\nwenddocs{}\nwbegincode{716}\moddef{SLHA: procedures}\plusendmoddef
  subroutine slha_write_MODSEL (u, model, mssm_type)
    integer, intent(in) :: u
    type(model_t), intent(in), target :: model
    integer, intent(out) :: mssm_type
    type(var_list_t), pointer :: var_list
    integer :: model_id
    type(string_t) :: mtype_string
    var_list => model%get_var_list_ptr ()
    if (var_list%contains (var_str ("mtype"))) then
       mssm_type = nint (var_list%get_rval (var_str ("mtype")))
    else
       call msg_error ("SLHA: parameter 'mtype' (SUSY breaking scheme) " &
            // "is unknown in current model, no SLHA output possible")
       mssm_type = -1
       return
    end if
    call write_block_header (u, "MODSEL", "SUSY model selection")
    select case (mssm_type)
    case (0);  mtype_string = "Generic MSSM"
    case (1);  mtype_string = "SUGRA"
    case (2);  mtype_string = "GMSB"
    case (3);  mtype_string = "AMSB"
    case default
       mtype_string = "unknown"
    end select
    call write_integer_parameter (u, 1, mssm_type, &
         "SUSY-breaking scheme: " // char (mtype_string))
    select case (char (model%get_name ()))
    case ("MSSM");  model_id = MDL_MSSM
    case ("NMSSM"); model_id = MDL_NMSSM
    case default
       model_id = 0
    end select
    call write_integer_parameter (u, 3, model_id, &
         "SUSY model type: " // char (model%get_name ()))
  end subroutine slha_write_MODSEL
    
\eatline
\nwendcode{}\nwbegindocs{717}\nwdocspar
\subsubsection{SMINPUTS}
Read SM parameters and update the variable list accordingly.  If a
parameter is not defined in the block, we use the previous value from
the model variable list.  For the basic parameters we have to do a
small recalculation, since SLHA uses the $G_F$-$\alpha$-$m_Z$ scheme,
while \whizard\ derives them from $G_F$, $m_W$, and $m_Z$.
\nwenddocs{}\nwbegincode{718}\moddef{SLHA: procedures}\plusendmoddef
  subroutine slha_handle_SMINPUTS (parse_tree, model)
    type(parse_tree_t), intent(in) :: parse_tree
    type(model_t), intent(inout), target :: model
    type(parse_node_t), pointer :: pn_block
    real(default) :: alpha_em_i, GF, alphas, mZ
    real(default) :: ee, vv, cw_sw, cw2, mW
    real(default) :: mb, mtop, mtau
    type(var_list_t), pointer :: var_list
    var_list => model%get_var_list_ptr ()
    pn_block => slha_get_block_ptr &
         (parse_tree, var_str ("SMINPUTS"), required=.true.)
    if (.not. (associated (pn_block)))  return
    alpha_em_i = &
         get_parameter_in_block (pn_block, 1, var_str ("alpha_em_i"), var_list)
    GF = get_parameter_in_block (pn_block, 2, var_str ("GF"), var_list)
    alphas = &
         get_parameter_in_block (pn_block, 3, var_str ("alphas"), var_list)
    mZ   = get_parameter_in_block (pn_block, 4, var_str ("mZ"), var_list)
    mb   = get_parameter_in_block (pn_block, 5, var_str ("mb"), var_list)
    mtop = get_parameter_in_block (pn_block, 6, var_str ("mtop"), var_list)
    mtau = get_parameter_in_block (pn_block, 7, var_str ("mtau"), var_list)
    ee = sqrt (4 * pi / alpha_em_i)
    vv = 1 / sqrt (sqrt (2._default) * GF)
    cw_sw = ee * vv / (2 * mZ)
    if (2*cw_sw <= 1) then
       cw2 = (1 + sqrt (1 - 4 * cw_sw**2)) / 2
       mW = mZ * sqrt (cw2)
       call var_list_set_real (var_list, var_str ("GF"), GF, .true.)
       call var_list_set_real (var_list, var_str ("mZ"), mZ, .true.)
       call var_list_set_real (var_list, var_str ("mW"), mW, .true.)
       call var_list_set_real (var_list, var_str ("mtau"), mtau, .true.)
       call var_list_set_real (var_list, var_str ("mb"), mb, .true.)
       call var_list_set_real (var_list, var_str ("mtop"), mtop, .true.)
       call var_list_set_real (var_list, var_str ("alphas"), alphas, .true.)
    else
       call msg_fatal ("SLHA: Unphysical SM parameter values")
       return
    end if
  end subroutine slha_handle_SMINPUTS

\eatline
\nwendcode{}\nwbegindocs{719}Write a SMINPUTS block.
\nwenddocs{}\nwbegincode{720}\moddef{SLHA: procedures}\plusendmoddef
  subroutine slha_write_SMINPUTS (u, model)
    integer, intent(in) :: u
    type(model_t), intent(in), target :: model
    type(var_list_t), pointer :: var_list
    var_list => model%get_var_list_ptr ()
    call write_block_header (u, "SMINPUTS", "SM input parameters")
    call write_real_data_item (u, 1, var_str ("alpha_em_i"), var_list, &
         "Inverse electromagnetic coupling alpha (Z pole)")
    call write_real_data_item (u, 2, var_str ("GF"), var_list, &
         "Fermi constant")
    call write_real_data_item (u, 3, var_str ("alphas"), var_list, &
         "Strong coupling alpha_s (Z pole)")
    call write_real_data_item (u, 4, var_str ("mZ"), var_list, &
         "Z mass")
    call write_real_data_item (u, 5, var_str ("mb"), var_list, &
         "b running mass (at mb)")
    call write_real_data_item (u, 6, var_str ("mtop"), var_list, &
         "top mass")
    call write_real_data_item (u, 7, var_str ("mtau"), var_list, &
         "tau mass")
  end subroutine slha_write_SMINPUTS

\eatline
\nwendcode{}\nwbegindocs{721}\nwdocspar
\subsubsection{MINPAR}
The block of SUSY input parameters.  They are accessible to WHIZARD,
but they only get used when an external spectrum generator is
invoked.  The precise set of parameters depends on the type of SUSY
breaking, which by itself is one of the parameters.
\nwenddocs{}\nwbegincode{722}\moddef{SLHA: procedures}\plusendmoddef
  subroutine slha_handle_MINPAR (parse_tree, model, mssm_type)
    type(parse_tree_t), intent(in) :: parse_tree
    type(model_t), intent(inout), target :: model
    integer, intent(in) :: mssm_type
    type(var_list_t), pointer :: var_list
    type(parse_node_t), pointer :: pn_block
    var_list => model%get_var_list_ptr ()
    call var_list_set_real (var_list, &
         var_str ("mtype"),  real(mssm_type, default), is_known=.true.)
    pn_block => slha_get_block_ptr &
         (parse_tree, var_str ("MINPAR"), required=.true.)
    select case (mssm_type)
    case (MSSM_SUGRA)
       call set_data_item (pn_block, 1, var_str ("m_zero"), var_list)
       call set_data_item (pn_block, 2, var_str ("m_half"), var_list)
       call set_data_item (pn_block, 3, var_str ("tanb"), var_list)
       call set_data_item (pn_block, 4, var_str ("sgn_mu"), var_list)
       call set_data_item (pn_block, 5, var_str ("A0"), var_list)
    case (MSSM_GMSB)
       call set_data_item (pn_block, 1, var_str ("Lambda"), var_list)
       call set_data_item (pn_block, 2, var_str ("M_mes"), var_list)
       call set_data_item (pn_block, 3, var_str ("tanb"), var_list)
       call set_data_item (pn_block, 4, var_str ("sgn_mu"), var_list)
       call set_data_item (pn_block, 5, var_str ("N_5"), var_list)
       call set_data_item (pn_block, 6, var_str ("c_grav"), var_list)
    case (MSSM_AMSB)
       call set_data_item (pn_block, 1, var_str ("m_zero"), var_list)
       call set_data_item (pn_block, 2, var_str ("m_grav"), var_list)
       call set_data_item (pn_block, 3, var_str ("tanb"), var_list)
       call set_data_item (pn_block, 4, var_str ("sgn_mu"), var_list)
    case default
       call set_data_item (pn_block, 3, var_str ("tanb"), var_list)
    end select
  end subroutine slha_handle_MINPAR

\eatline
\nwendcode{}\nwbegindocs{723}Write a MINPAR block as appropriate for the current model type.
\nwenddocs{}\nwbegincode{724}\moddef{SLHA: procedures}\plusendmoddef
  subroutine slha_write_MINPAR (u, model, mssm_type)
    integer, intent(in) :: u
    type(model_t), intent(in), target :: model
    integer, intent(in) :: mssm_type
    type(var_list_t), pointer :: var_list
    var_list => model%get_var_list_ptr ()
    call write_block_header (u, "MINPAR", "Basic SUSY input parameters")
    select case (mssm_type)
    case (MSSM_SUGRA)
       call write_real_data_item (u, 1, var_str ("m_zero"), var_list, &
            "Common scalar mass")
       call write_real_data_item (u, 2, var_str ("m_half"), var_list, &
            "Common gaugino mass")
       call write_real_data_item (u, 3, var_str ("tanb"), var_list, &
            "tan(beta)")
       call write_integer_data_item (u, 4, &
            var_str ("sgn_mu"), var_list, &
            "Sign of mu")
       call write_real_data_item (u, 5, var_str ("A0"), var_list, &
            "Common trilinear coupling")
    case (MSSM_GMSB)
       call write_real_data_item (u, 1, var_str ("Lambda"), var_list, &
            "Soft-breaking scale")
       call write_real_data_item (u, 2, var_str ("M_mes"), var_list, &
            "Messenger scale")
       call write_real_data_item (u, 3, var_str ("tanb"), var_list, &
            "tan(beta)")
       call write_integer_data_item (u, 4, &
            var_str ("sgn_mu"), var_list, &
            "Sign of mu")
       call write_integer_data_item (u, 5, var_str ("N_5"), var_list, &
            "Messenger index")
       call write_real_data_item (u, 6, var_str ("c_grav"), var_list, &
            "Gravitino mass factor")
    case (MSSM_AMSB)
       call write_real_data_item (u, 1, var_str ("m_zero"), var_list, &
            "Common scalar mass")
       call write_real_data_item (u, 2, var_str ("m_grav"), var_list, &
            "Gravitino mass")
       call write_real_data_item (u, 3, var_str ("tanb"), var_list, &
            "tan(beta)")
       call write_integer_data_item (u, 4, &
            var_str ("sgn_mu"), var_list, &
            "Sign of mu")
    case default
       call write_real_data_item (u, 3, var_str ("tanb"), var_list, &
            "tan(beta)")
    end select
  end subroutine slha_write_MINPAR
    
\eatline
\nwendcode{}\nwbegindocs{725}\nwdocspar
\subsubsection{Mass spectrum}
Set masses.  Since the particles are identified by PDG code, read
the line and try to set the appropriate particle mass in the current
model.  At the end, update parameters, just in case the $W$ or $Z$
mass was included.
\nwenddocs{}\nwbegincode{726}\moddef{SLHA: procedures}\plusendmoddef
  subroutine slha_handle_MASS (parse_tree, model)
    type(parse_tree_t), intent(in) :: parse_tree
    type(model_t), intent(inout), target :: model
    type(parse_node_t), pointer :: pn_block, pn_data, pn_line, pn_code
    type(parse_node_t), pointer :: pn_mass
    integer :: pdg
    real(default) :: mass
    pn_block => slha_get_block_ptr &
         (parse_tree, var_str ("MASS"), required=.true.)
    if (.not. (associated (pn_block)))  return
    pn_data => parse_node_get_sub_ptr (pn_block, 4)
    do while (associated (pn_data))
       pn_line => parse_node_get_sub_ptr (pn_data, 2)
       pn_code => parse_node_get_sub_ptr (pn_line)
       if (associated (pn_code)) then
          pdg = get_integer_parameter (pn_code)
          pn_mass => parse_node_get_next_ptr (pn_code)
          if (associated (pn_mass)) then
             mass = get_real_parameter (pn_mass)
             call model%set_field_mass (pdg, mass)
          else
             call msg_error ("SLHA: Block MASS: Missing mass value")
          end if
       else
          call msg_error ("SLHA: Block MASS: Missing PDG code")
       end if
       pn_data => parse_node_get_next_ptr (pn_data)
    end do
  end subroutine slha_handle_MASS

\eatline
\nwendcode{}\nwbegindocs{727}\nwdocspar
\subsubsection{Widths}
Set widths.  For each DECAY block, extract the header, read the PDG
code and width, and try to set the appropriate particle width in the
current model.
\nwenddocs{}\nwbegincode{728}\moddef{SLHA: procedures}\plusendmoddef
  subroutine slha_handle_decays (parse_tree, model)
    type(parse_tree_t), intent(in) :: parse_tree
    type(model_t), intent(inout), target :: model
    type(parse_node_t), pointer :: pn_decay, pn_decay_spec, pn_code, pn_width
    integer :: pdg
    real(default) :: width
    pn_decay => slha_get_first_decay_ptr (parse_tree)
    do while (associated (pn_decay))
       pn_decay_spec => parse_node_get_sub_ptr (pn_decay, 2)
       pn_code => parse_node_get_sub_ptr (pn_decay_spec)
       pdg = get_integer_parameter (pn_code)
       pn_width => parse_node_get_next_ptr (pn_code)
       width = get_real_parameter (pn_width)
       call model%set_field_width (pdg, width)
       pn_decay => slha_get_next_decay_ptr (pn_decay)
    end do
  end subroutine slha_handle_decays

\eatline
\nwendcode{}\nwbegindocs{729}\nwdocspar
\subsubsection{Mixing matrices}
Read mixing matrices.  We can treat all matrices by a single
procedure if we just know the block name, variable prefix, and matrix
dimension.  The matrix dimension must be less than 10. 
For the pseudoscalar Higgses in NMSSM-type models we need off-diagonal 
matrices, so we generalize the definition.
\nwenddocs{}\nwbegincode{730}\moddef{SLHA: procedures}\plusendmoddef
  subroutine slha_handle_matrix_block &
       (parse_tree, block_name, var_prefix, dim1, dim2, model)
    type(parse_tree_t), intent(in) :: parse_tree
    character(*), intent(in) :: block_name, var_prefix
    integer, intent(in) :: dim1, dim2
    type(model_t), intent(inout), target :: model
    type(parse_node_t), pointer :: pn_block
    type(var_list_t), pointer :: var_list
    integer :: i, j
    character(len=len(var_prefix)+2) :: var_name
    var_list => model%get_var_list_ptr ()
    pn_block => slha_get_block_ptr &
         (parse_tree, var_str (block_name), required=.false.)
    if (.not. (associated (pn_block)))  return
    do i = 1, dim1
       do j = 1, dim2
          write (var_name, "(A,I1,I1)")  var_prefix, i, j
          call set_matrix_element (pn_block, i, j, var_str (var_name), var_list)
       end do
    end do
  end subroutine slha_handle_matrix_block

\eatline
\nwendcode{}\nwbegindocs{731}\nwdocspar
\subsubsection{Higgs data}
Read the block ALPHA which holds just the Higgs mixing angle.
\nwenddocs{}\nwbegincode{732}\moddef{SLHA: procedures}\plusendmoddef
  subroutine slha_handle_ALPHA (parse_tree, model)
    type(parse_tree_t), intent(in) :: parse_tree
    type(model_t), intent(inout), target :: model
    type(parse_node_t), pointer :: pn_block, pn_line, pn_data, pn_item
    type(var_list_t), pointer :: var_list
    real(default) :: al_h
    var_list => model%get_var_list_ptr ()
    pn_block => slha_get_block_ptr &
         (parse_tree, var_str ("ALPHA"), required=.false.)
    if (.not. (associated (pn_block)))  return
    pn_data => parse_node_get_sub_ptr (pn_block, 4)
    pn_line => parse_node_get_sub_ptr (pn_data, 2)
    pn_item => parse_node_get_sub_ptr (pn_line)
    if (associated (pn_item)) then
       al_h = get_real_parameter (pn_item)
       call var_list_set_real (var_list, var_str ("al_h"), al_h, &
            is_known=.true., ignore=.true.)
    end if
  end subroutine slha_handle_ALPHA

\eatline
\nwendcode{}\nwbegindocs{733}Read the block HMIX for the Higgs mixing parameters
\nwenddocs{}\nwbegincode{734}\moddef{SLHA: procedures}\plusendmoddef
  subroutine slha_handle_HMIX (parse_tree, model)
    type(parse_tree_t), intent(in) :: parse_tree
    type(model_t), intent(inout), target :: model
    type(parse_node_t), pointer :: pn_block
    type(var_list_t), pointer :: var_list
    var_list => model%get_var_list_ptr ()
    pn_block => slha_get_block_ptr &
         (parse_tree, var_str ("HMIX"), required=.false.)
    if (.not. (associated (pn_block)))  return
    call set_data_item (pn_block, 1, var_str ("mu_h"), var_list)
    call set_data_item (pn_block, 2, var_str ("tanb_h"), var_list)
  end subroutine slha_handle_HMIX

\eatline
\nwendcode{}\nwbegindocs{735}Read the block NMSSMRUN for the specific NMSSM parameters
\nwenddocs{}\nwbegincode{736}\moddef{SLHA: procedures}\plusendmoddef
  subroutine slha_handle_NMSSMRUN (parse_tree, model)
    type(parse_tree_t), intent(in) :: parse_tree
    type(model_t), intent(inout), target :: model
    type(parse_node_t), pointer :: pn_block
    type(var_list_t), pointer :: var_list
    var_list => model%get_var_list_ptr ()
    pn_block => slha_get_block_ptr &
         (parse_tree, var_str ("NMSSMRUN"), required=.false.)
    if (.not. (associated (pn_block)))  return
    call set_data_item (pn_block, 1, var_str ("ls"), var_list)
    call set_data_item (pn_block, 2, var_str ("ks"), var_list)
    call set_data_item (pn_block, 3, var_str ("a_ls"), var_list)    
    call set_data_item (pn_block, 4, var_str ("a_ks"), var_list)        
    call set_data_item (pn_block, 5, var_str ("nmu"), var_list)    
    end subroutine slha_handle_NMSSMRUN

\eatline
\nwendcode{}\nwbegindocs{737}\nwdocspar
\subsection{Parser}
Read a SLHA file from stream, including preprocessing, and make up a
parse tree.
\nwenddocs{}\nwbegincode{738}\moddef{SLHA: procedures}\plusendmoddef
  subroutine slha_parse_stream (stream, parse_tree)
    type(stream_t), intent(inout), target :: stream
    type(parse_tree_t), intent(out) :: parse_tree
    type(ifile_t) :: ifile
    type(lexer_t) :: lexer
    type(stream_t), target :: stream_tmp
    call slha_preprocess (stream, ifile)
    call stream_init (stream_tmp, ifile)
    call lexer_init_slha (lexer)
    call lexer_assign_stream (lexer, stream_tmp)
    call parse_tree_init (parse_tree, syntax_slha, lexer)
    call lexer_final (lexer)
    call stream_final (stream_tmp)
    call ifile_final (ifile)
  end subroutine slha_parse_stream

\eatline
\nwendcode{}\nwbegindocs{739}Read a SLHA file chosen by name.  Check first the current directory,
then the directory where SUSY input files should be located.
\nwenddocs{}\nwbegincode{740}\moddef{SLHA: procedures}\plusendmoddef
  subroutine slha_parse_file (file, os_data, parse_tree)
    type(string_t), intent(in) :: file
    type(os_data_t), intent(in) :: os_data
    type(parse_tree_t), intent(out) :: parse_tree
    logical :: exist
    type(string_t) :: filename
    type(stream_t), target :: stream
    call msg_message ("Reading SLHA input file '" // char (file) // "'")
    filename = file
    inquire (file=char(filename), exist=exist)
    if (.not. exist) then
       filename = os_data%whizard_susypath // "/" // file
       inquire (file=char(filename), exist=exist)
       if (.not. exist) then
          call msg_fatal ("SLHA input file '" // char (file) // "' not found")
          return
       end if
    end if
    call stream_init (stream, char (filename))
    call slha_parse_stream (stream, parse_tree)
    call stream_final (stream)
  end subroutine slha_parse_file

\eatline
\nwendcode{}\nwbegindocs{741}\nwdocspar
\subsection{API}
Read the SLHA file, parse it, and interpret the parse tree.  The model
parameters retrieved from the file will be inserted into the
appropriate model, which is loaded and modified in the background.
The pointer to this model is returned as the last argument.
\nwenddocs{}\nwbegincode{742}\moddef{SLHA: public}\plusendmoddef
  public :: slha_read_file
\nwendcode{}\nwbegincode{743}\moddef{SLHA: procedures}\plusendmoddef
  subroutine slha_read_file &
       (file, os_data, model, input, spectrum, decays)
    type(string_t), intent(in) :: file
    type(os_data_t), intent(in) :: os_data
    type(model_t), intent(inout), target :: model
    logical, intent(in) :: input, spectrum, decays
    type(parse_tree_t) :: parse_tree
    call slha_parse_file (file, os_data, parse_tree)
    if (associated (parse_tree_get_root_ptr (parse_tree))) then
       call slha_interpret_parse_tree &
            (parse_tree, model, input, spectrum, decays)
       call parse_tree_final (parse_tree)
       call model%update_parameters ()
    end if
  end subroutine slha_read_file

\eatline
\nwendcode{}\nwbegindocs{744}Write the SLHA contents, as far as possible, to external file.
\nwenddocs{}\nwbegincode{745}\moddef{SLHA: public}\plusendmoddef
  public :: slha_write_file
\nwendcode{}\nwbegincode{746}\moddef{SLHA: procedures}\plusendmoddef
  subroutine slha_write_file (file, model, input, spectrum, decays)
    type(string_t), intent(in) :: file
    type(model_t), target, intent(in) :: model
    logical, intent(in) :: input, spectrum, decays
    integer :: mssm_type
    integer :: u
    u = free_unit ()
    call msg_message ("Writing SLHA output file '" // char (file) // "'")
    open (unit=u, file=char(file), action="write", status="replace")
    write (u, "(A)")  "# SUSY Les Houches Accord"
    write (u, "(A)")  "# Output generated by " // trim (VERSION_STRING)
    call slha_write_MODSEL (u, model, mssm_type)
    if (input) then
       call slha_write_SMINPUTS (u, model)
       call slha_write_MINPAR (u, model, mssm_type)
    end if
    if (spectrum) then
       call msg_bug ("SLHA: spectrum output not supported yet")
    end if
    if (decays) then
       call msg_bug ("SLHA: decays output not supported yet")
    end if
    close (u)
  end subroutine slha_write_file

\eatline
\nwendcode{}\nwbegindocs{747}\nwdocspar
\subsection{Unit tests}
\nwenddocs{}\nwbegincode{748}\moddef{SLHA: public}\plusendmoddef
  public :: slha_test
\nwendcode{}\nwbegincode{749}\moddef{SLHA: procedures}\plusendmoddef
  subroutine slha_test (u, results)
    integer, intent(in) :: u
    type(test_results_t), intent(inout) :: results
  \LA{}SLHA: execute tests\RA{}  
  end subroutine slha_test

\nwendcode{}\nwbegindocs{750} %def slha_test
\nwenddocs{}\nwbegindocs{751}Checking the basics of the SLHA interface.
\nwenddocs{}\nwbegincode{752}\moddef{SLHA: execute tests}\endmoddef
  call test (slha_1, "slha_1", &
       "check SLHA interface", &
       u, results)
\nwendcode{}\nwbegincode{753}\moddef{SLHA: tests}\endmoddef
  subroutine slha_1 (u) 
    integer, intent(in) :: u
    type(os_data_t), pointer :: os_data => null ()
    type(parse_tree_t), pointer :: parse_tree => null ()
    integer :: u_file, iostat
    character(80) :: buffer
    character(*), parameter :: file_slha = "slha_test.dat"
    type(model_list_t) :: model_list
    type(model_t), pointer :: model => null ()
    
    write (u, "(A)")  "* Test output: SLHA Interface"
    write (u, "(A)")  "*   Purpose: test SLHA file reading and writing"
    write (u, "(A)")            

    write (u, "(A)")  "* Initializing"
    write (u, "(A)")    
    
    allocate (os_data)
    allocate (parse_tree)
    call os_data_init (os_data)
    call syntax_model_file_init ()
    call model_list%read_model &
         (var_str("MSSM"), var_str("MSSM.mdl"), os_data, model)
    call syntax_slha_init ()

    write (u, "(A)")  "* Reading SLHA file sps1ap_decays.slha"
    write (u, "(A)")    
    
    call slha_parse_file (var_str ("sps1ap_decays.slha"), os_data, parse_tree)

    write (u, "(A)")  "* Writing the parse tree:"
    write (u, "(A)")    

    call parse_tree_write (parse_tree, u)
    
    write (u, "(A)")  "* Interpreting the parse tree"
    write (u, "(A)")    
    
    call slha_interpret_parse_tree (parse_tree, model, &
         input=.true., spectrum=.true., decays=.true.)    
    call parse_tree_final (parse_tree)

    write (u, "(A)")  "* Writing out the list of variables (reals only):"
    write (u, "(A)")    
    
    call var_list_write (model%get_var_list_ptr (), &
         only_type = V_REAL, unit = u)

    write (u, "(A)")
    write (u, "(A)")  "* Writing SLHA output to '" // file_slha // "'"
    write (u, "(A)")    
        
    call slha_write_file (var_str (file_slha), model, input=.true., &
         spectrum=.false., decays=.false.)
    u_file = free_unit ()
    open (u_file, file = file_slha, action = "read", status = "old")
    do
       read (u_file, "(A)", iostat = iostat)  buffer
       if (buffer(1:37) == "# Output generated by WHIZARD version") then
          buffer = "[...]"
       end if
       if (iostat /= 0)  exit
       write (u, "(A)") trim (buffer)
    end do
    close (u_file)

    write (u, "(A)")
    write (u, "(A)")  "* Cleanup"
    write (u, "(A)")
    
    call parse_tree_final (parse_tree)
    deallocate (parse_tree)
    deallocate (os_data)

    write (u, "(A)")  "* Test output end: slha_1"
    write (u, "(A)")
    
  end subroutine slha_1

\eatline
\nwendcode{}\nwbegindocs{754}\nwdocspar
\nwenddocs{}\nwfilename{../../../src/whizard-core/whizard.nw}\nwbegindocs{0}% -*- ess-noweb-default-code-mode: f90-mode; noweb-default-code-mode: f90-mode; -*- 
% WHIZARD main code as NOWEB source

\nwenddocs{}\nwbegindocs{1}\nwdocspar
%%%%%%%%%%%%%%%%%%%%%%%%%%%%%%%%%%%%%%%%%%%%%%%%%%%%%%%%%%%%%%%%%%%%%%%%
\chapter{Integration and Process Objects}

This is the central part of the \whizard\ package.  It provides the
functionality for evaluating structure functions, kinematics and matrix
elements, integration and event generation.  It combines the various
parts that deal with those tasks individually and organizes the data transfer
between them.
\begin{description}
\item[subevt\_expr]
  This enables process observables as (abstract) expressions, to be
  evaluated for each process call.
\item[parton\_states]
  A {\tt{}parton{\char95}state{\char95}t} object represents an elementary partonic
  interaction.  There are two versions: one for the isolated
  elementary process, one for the elementary process convoluted with
  the structure-function chain.  The parton state is an effective
  state.  It needs not coincide with the seed-kinematics state which is
  used in evaluating phase space.
\item[processes]
  Here, all pieces are combined for the purpose of evaluating the
  elementary processes.  The whole algorithm is coded in terms of
  abstract data types as defined in the appropriate modules: {\tt{}prc{\char95}core}
  for matrix-element evaluation, {\tt{}prc{\char95}core{\char95}def} for the associated
  configuration and driver, {\tt{}sf{\char95}base} for beams and structure-functions,
  {\tt{}phs{\char95}base} for phase space, and {\tt{}mci{\char95}base} for integration and event
  generation.
\item[process\_stacks]
  Process stacks collect process objects.
\end{description}

%%%%%%%%%%%%%%%%%%%%%%%%%%%%%%%%%%%%%%%%%%%%%%%%%%%%%%%%%%%%%%%%%%%%%%%%
\section{Process observables}
We define an abstract {\tt{}subevt{\char95}expr{\char95}t} object as an extension of the
{\tt{}subevt{\char95}t} type.  The object contains a local variable list, variable
instances (as targets for pointers in the variable list), and evaluation
trees.  The evaluation trees reference both the variables and the {\tt{}subevt}.

There are two instances of the abstract type: one for process instances, one
for physical events.  Both have a common logical expression {\tt{}selection}
which determines whether the object passes user-defined cuts.

The intention is that we fill the {\tt{}subevt{\char95}t} base object and compute the
variables once we have evaluated a kinematical phase space point (or a
complete event).  We then evaluate the expressions and can use the results in
further calculations.

The {\tt{}process{\char95}expr{\char95}t} extension contains furthermore scale and weight
expressions.  The {\tt{}event{\char95}expr{\char95}t} extension contains a reweighting-factor
expression and a logical expression for event analysis.  In practice, we will
link the variable list of the {\tt{}event{\char95}obs} object to the variable list of the
currently active {\tt{}process{\char95}obs} object, such that the process variables are
available to both objects.  Event variables are meaningful only for physical
events.

Note that there are unit tests, but they are deferred to the
{\tt{}expr{\char95}tests} module.
\nwenddocs{}\nwbegincode{2}\moddef{\code{}subevt{\char95}expr.f90\edoc{}}\endmoddef
\LA{}File header\RA{}
module subevt_expr

\LA{}Use kinds\RA{}
\LA{}Use strings\RA{}
  use io_units
  use format_utils, only: write_separator
  use diagnostics
  use lorentz
  use subevents
  use variables
  use flavors
  use quantum_numbers
  use interactions
  use particles
  use expr_base

\LA{}Standard module head\RA{}

\LA{}Subevt expr: public\RA{}

\LA{}Subevt expr: types\RA{}

\LA{}Subevt expr: interfaces\RA{}

contains
  
\LA{}Subevt expr: procedures\RA{}

end module subevt_expr
\eatline
\nwendcode{}\nwbegindocs{3}\nwdocspar
\subsection{Abstract base type}
\nwenddocs{}\nwbegincode{4}\moddef{Subevt expr: types}\endmoddef
  type, extends (subevt_t), abstract :: subevt_expr_t
     logical :: subevt_filled = .false.
     type(var_list_t) :: var_list
     real(default) :: sqrts_hat = 0
     integer :: n_in = 0
     integer :: n_out = 0
     integer :: n_tot = 0
     logical :: has_selection = .false.
     class(expr_t), allocatable :: selection
   contains
   \LA{}Subevt expr: subevt expr: TBP\RA{}
  end type subevt_expr_t
  
\eatline
\nwendcode{}\nwbegindocs{5}Output: Base and extended version.  We already have a {\tt{}write} routine for
the {\tt{}subevt{\char95}t} parent type.
\nwenddocs{}\nwbegincode{6}\moddef{Subevt expr: subevt expr: TBP}\endmoddef
  procedure :: base_write => subevt_expr_write
\nwendcode{}\nwbegincode{7}\moddef{Subevt expr: procedures}\endmoddef
  subroutine subevt_expr_write (object, unit, pacified)
    class(subevt_expr_t), intent(in) :: object
    integer, intent(in), optional :: unit
    logical, intent(in), optional :: pacified
    integer :: u
    u = given_output_unit (unit)
    write (u, "(1x,A)")  "Local variables:"
    call write_separator (u)
    call var_list_write (object%var_list, u, follow_link=.false.)
    call write_separator (u)
    if (object%subevt_filled) then
       call object%subevt_t%write (u, pacified = pacified)
       if (object%has_selection) then
          call write_separator (u)
          write (u, "(1x,A)")  "Selection expression:"
          call write_separator (u)
          call object%selection%write (u)
       end if
    else
       write (u, "(1x,A)")  "subevt: [undefined]"
    end if
  end subroutine subevt_expr_write
    
\eatline
\nwendcode{}\nwbegindocs{8}Finalizer.
\nwenddocs{}\nwbegincode{9}\moddef{Subevt expr: subevt expr: TBP}\plusendmoddef
  procedure (subevt_expr_final), deferred :: final
  procedure :: base_final => subevt_expr_final
\nwendcode{}\nwbegincode{10}\moddef{Subevt expr: procedures}\plusendmoddef
  subroutine subevt_expr_final (object)
    class(subevt_expr_t), intent(inout) :: object
    call object%var_list%final ()
    if (object%has_selection) then
       call object%selection%final ()
    end if
  end subroutine subevt_expr_final
  
\eatline
\nwendcode{}\nwbegindocs{11}\nwdocspar
\subsection{Initialization}
Initialization: define local variables and establish pointers.

The common variables are {\tt{}sqrts} (the nominal beam energy, fixed),
{\tt{}sqrts{\char95}hat} (the actual energy), {\tt{}n{\char95}in}, {\tt{}n{\char95}out}, and {\tt{}n{\char95}tot} for
the {\tt{}subevt}.  With the exception of {\tt{}sqrts}, all are implemented as
pointers to subobjects.
\nwenddocs{}\nwbegincode{12}\moddef{Subevt expr: subevt expr: TBP}\plusendmoddef
  procedure (subevt_expr_setup_vars), deferred :: setup_vars
  procedure :: base_setup_vars => subevt_expr_setup_vars
\nwendcode{}\nwbegincode{13}\moddef{Subevt expr: procedures}\plusendmoddef
  subroutine subevt_expr_setup_vars (expr, sqrts)
    class(subevt_expr_t), intent(inout), target :: expr
    real(default), intent(in) :: sqrts
    call expr%var_list%final ()
    call var_list_append_real (expr%var_list, &
         var_str ("sqrts"), sqrts, &
         locked = .true., verbose = .false., intrinsic = .true.)
    call var_list_append_real_ptr (expr%var_list, &
         var_str ("sqrts_hat"), expr%sqrts_hat, &
         is_known = expr%subevt_filled, &
         locked = .true., verbose = .false., intrinsic = .true.)
    call var_list_append_int_ptr (expr%var_list, &
         var_str ("n_in"), expr%n_in, &
         is_known = expr%subevt_filled, &
         locked = .true., verbose = .false., intrinsic = .true.)
    call var_list_append_int_ptr (expr%var_list, &
         var_str ("n_out"), expr%n_out, &
         is_known = expr%subevt_filled, &
         locked = .true., verbose = .false., intrinsic = .true.)
    call var_list_append_int_ptr (expr%var_list, &
         var_str ("n_tot"), expr%n_tot, &
         is_known = expr%subevt_filled, &
         locked = .true., verbose = .false., intrinsic = .true.)
  end subroutine subevt_expr_setup_vars
    
\eatline
\nwendcode{}\nwbegindocs{14}Append the subevent expr (its base-type core) itself to the variable
list, if it is not yet present.
\nwenddocs{}\nwbegincode{15}\moddef{Subevt expr: subevt expr: TBP}\plusendmoddef
  procedure :: setup_var_self => subevt_expr_setup_var_self
\nwendcode{}\nwbegincode{16}\moddef{Subevt expr: procedures}\plusendmoddef
  subroutine subevt_expr_setup_var_self (expr)
    class(subevt_expr_t), intent(inout), target :: expr
    if (.not. expr%var_list%contains (var_str ("@evt"))) then
       call var_list_append_subevt_ptr &
            (expr%var_list, &
            var_str ("@evt"), expr%subevt_t, &
            is_known = expr%subevt_filled, &
            locked = .true., verbose = .false., intrinsic=.true.)
    end if
  end subroutine subevt_expr_setup_var_self
  
\eatline
\nwendcode{}\nwbegindocs{17}Link a variable list to the local one.  This could be done event by event,
but before evaluating expressions.
\nwenddocs{}\nwbegincode{18}\moddef{Subevt expr: subevt expr: TBP}\plusendmoddef
  procedure :: link_var_list => subevt_expr_link_var_list
\nwendcode{}\nwbegincode{19}\moddef{Subevt expr: procedures}\plusendmoddef
  subroutine subevt_expr_link_var_list (expr, var_list)
    class(subevt_expr_t), intent(inout) :: expr
    type(var_list_t), intent(in), target :: var_list
    call expr%var_list%link (var_list)
  end subroutine subevt_expr_link_var_list

\eatline
\nwendcode{}\nwbegindocs{20}Compile the selection expression.  If there is no expression, the build
method won't allocate the expression object.
\nwenddocs{}\nwbegincode{21}\moddef{Subevt expr: subevt expr: TBP}\plusendmoddef
  procedure :: setup_selection => subevt_expr_setup_selection
\nwendcode{}\nwbegincode{22}\moddef{Subevt expr: procedures}\plusendmoddef
  subroutine subevt_expr_setup_selection (expr, ef_cuts)
    class(subevt_expr_t), intent(inout), target :: expr
    class(expr_factory_t), intent(in) :: ef_cuts
    call ef_cuts%build (expr%selection)
    if (allocated (expr%selection)) then
       call expr%setup_var_self ()
       call expr%selection%setup_lexpr (expr%var_list)
       expr%has_selection = .true.
    end if
  end subroutine subevt_expr_setup_selection

\eatline
\nwendcode{}\nwbegindocs{23}\nwdocspar
\subsection{Evaluation}
Reset to initial state, i.e., mark the {\tt{}subevt} as invalid.
\nwenddocs{}\nwbegincode{24}\moddef{Subevt expr: subevt expr: TBP}\plusendmoddef
  procedure :: reset => subevt_expr_reset
  procedure :: base_reset => subevt_expr_reset
\nwendcode{}\nwbegincode{25}\moddef{Subevt expr: procedures}\plusendmoddef
  subroutine subevt_expr_reset (expr)
    class(subevt_expr_t), intent(inout) :: expr
    expr%subevt_filled = .false.
  end subroutine subevt_expr_reset
  
\eatline
\nwendcode{}\nwbegindocs{26}Evaluate the selection expression and return the result.  There is also a
deferred version: this should evaluate the remaining expressions if the event
has passed.
\nwenddocs{}\nwbegincode{27}\moddef{Subevt expr: subevt expr: TBP}\plusendmoddef
  procedure :: base_evaluate => subevt_expr_evaluate
\nwendcode{}\nwbegincode{28}\moddef{Subevt expr: procedures}\plusendmoddef
  subroutine subevt_expr_evaluate (expr, passed)
    class(subevt_expr_t), intent(inout) :: expr
    logical, intent(out) :: passed 
    if (expr%has_selection) then
       call expr%selection%evaluate ()
       if (expr%selection%is_known ()) then
          passed = expr%selection%get_log ()
       else
          call msg_error ("Evaluate selection expression: result undefined")
          passed = .false.
       end if
    else
       passed = .true.
    end if
  end subroutine subevt_expr_evaluate
  
\eatline
\nwendcode{}\nwbegindocs{29}\nwdocspar
\subsection{Implementation for partonic events}
This implementation contains the expressions that we can evaluate for the
partonic process during integration.
\nwenddocs{}\nwbegincode{30}\moddef{Subevt expr: public}\endmoddef
  public :: parton_expr_t
\nwendcode{}\nwbegincode{31}\moddef{Subevt expr: types}\plusendmoddef
  type, extends (subevt_expr_t) :: parton_expr_t
     integer, dimension(:), allocatable :: i_beam
     integer, dimension(:), allocatable :: i_in
     integer, dimension(:), allocatable :: i_out
     logical :: has_scale = .false.
     logical :: has_fac_scale = .false.
     logical :: has_ren_scale = .false.
     logical :: has_weight = .false.
     class(expr_t), allocatable :: scale
     class(expr_t), allocatable :: fac_scale
     class(expr_t), allocatable :: ren_scale
     class(expr_t), allocatable :: weight
   contains
   \LA{}Subevt expr: parton expr: TBP\RA{}
  end type parton_expr_t
     
\eatline
\nwendcode{}\nwbegindocs{32}Finalizer.
\nwenddocs{}\nwbegincode{33}\moddef{Subevt expr: parton expr: TBP}\endmoddef
  procedure :: final => parton_expr_final
\nwendcode{}\nwbegincode{34}\moddef{Subevt expr: procedures}\plusendmoddef
  subroutine parton_expr_final (object)
    class(parton_expr_t), intent(inout) :: object
    call object%base_final ()
    if (object%has_scale) then
       call object%scale%final ()
    end if
    if (object%has_fac_scale) then
       call object%fac_scale%final ()
    end if
    if (object%has_ren_scale) then
       call object%ren_scale%final ()
    end if
    if (object%has_weight) then
       call object%weight%final ()
    end if
  end subroutine parton_expr_final

\eatline
\nwendcode{}\nwbegindocs{35}Output: continue writing the active expressions, after the common selection
expression.

Note: the {\tt{}prefix} argument is declared in the {\tt{}write} method of the
{\tt{}subevt{\char95}t} base type.  Here, it is unused.
\nwenddocs{}\nwbegincode{36}\moddef{Subevt expr: parton expr: TBP}\plusendmoddef
  procedure :: write => parton_expr_write
\nwendcode{}\nwbegincode{37}\moddef{Subevt expr: procedures}\plusendmoddef
  subroutine parton_expr_write (object, unit, prefix, pacified)
    class(parton_expr_t), intent(in) :: object
    integer, intent(in), optional :: unit
    character(*), intent(in), optional :: prefix
    logical, intent(in), optional :: pacified
    integer :: u
    u = given_output_unit (unit)
    call object%base_write (u, pacified = pacified)
    if (object%subevt_filled) then
       if (object%has_scale) then
          call write_separator (u)
          write (u, "(1x,A)")  "Scale expression:"
          call write_separator (u)
          call object%scale%write (u)
       end if
       if (object%has_fac_scale) then
          call write_separator (u)
          write (u, "(1x,A)")  "Factorization scale expression:"
          call write_separator (u)
          call object%fac_scale%write (u)
       end if
       if (object%has_ren_scale) then
          call write_separator (u)
          write (u, "(1x,A)")  "Renormalization scale expression:"
          call write_separator (u)
          call object%ren_scale%write (u)
       end if
       if (object%has_weight) then
          call write_separator (u)
          write (u, "(1x,A)")  "Weight expression:"
          call write_separator (u)
          call object%weight%write (u)
       end if
    end if
  end subroutine parton_expr_write
    
\eatline
\nwendcode{}\nwbegindocs{38}Define variables.
\nwenddocs{}\nwbegincode{39}\moddef{Subevt expr: parton expr: TBP}\plusendmoddef
  procedure :: setup_vars => parton_expr_setup_vars
\nwendcode{}\nwbegincode{40}\moddef{Subevt expr: procedures}\plusendmoddef
  subroutine parton_expr_setup_vars (expr, sqrts)
    class(parton_expr_t), intent(inout), target :: expr
    real(default), intent(in) :: sqrts
    call expr%base_setup_vars (sqrts)
  end subroutine parton_expr_setup_vars

\eatline
\nwendcode{}\nwbegindocs{41}Compile the scale expressions.  If a pointer is disassociated, there is
no expression.
\nwenddocs{}\nwbegincode{42}\moddef{Subevt expr: parton expr: TBP}\plusendmoddef
  procedure :: setup_scale => parton_expr_setup_scale
  procedure :: setup_fac_scale => parton_expr_setup_fac_scale
  procedure :: setup_ren_scale => parton_expr_setup_ren_scale
\nwendcode{}\nwbegincode{43}\moddef{Subevt expr: procedures}\plusendmoddef
  subroutine parton_expr_setup_scale (expr, ef_scale)
    class(parton_expr_t), intent(inout), target :: expr
    class(expr_factory_t), intent(in) :: ef_scale
    call ef_scale%build (expr%scale)
    if (allocated (expr%scale)) then
       call expr%setup_var_self ()
       call expr%scale%setup_expr (expr%var_list)
       expr%has_scale = .true.
    end if
  end subroutine parton_expr_setup_scale

  subroutine parton_expr_setup_fac_scale (expr, ef_fac_scale)
    class(parton_expr_t), intent(inout), target :: expr
    class(expr_factory_t), intent(in) :: ef_fac_scale
    call ef_fac_scale%build (expr%fac_scale)
    if (allocated (expr%fac_scale)) then
       call expr%setup_var_self ()
       call expr%fac_scale%setup_expr (expr%var_list)
       expr%has_fac_scale = .true.
    end if
  end subroutine parton_expr_setup_fac_scale

  subroutine parton_expr_setup_ren_scale (expr, ef_ren_scale)
    class(parton_expr_t), intent(inout), target :: expr
    class(expr_factory_t), intent(in) :: ef_ren_scale
    call ef_ren_scale%build (expr%ren_scale)
    if (allocated (expr%ren_scale)) then
       call expr%setup_var_self ()
       call expr%ren_scale%setup_expr (expr%var_list)
       expr%has_ren_scale = .true.
    end if
  end subroutine parton_expr_setup_ren_scale

\eatline
\eatline
\eatline
\nwendcode{}\nwbegindocs{44}Compile the weight expression.
\nwenddocs{}\nwbegincode{45}\moddef{Subevt expr: parton expr: TBP}\plusendmoddef
  procedure :: setup_weight => parton_expr_setup_weight
\nwendcode{}\nwbegincode{46}\moddef{Subevt expr: procedures}\plusendmoddef
  subroutine parton_expr_setup_weight (expr, ef_weight)
    class(parton_expr_t), intent(inout), target :: expr
    class(expr_factory_t), intent(in) :: ef_weight
    call ef_weight%build (expr%weight)
    if (allocated (expr%weight)) then
       call expr%setup_var_self ()
       call expr%weight%setup_expr (expr%var_list)
       expr%has_weight = .true.
    end if
  end subroutine parton_expr_setup_weight

\eatline
\nwendcode{}\nwbegindocs{47}Filling the partonic state consists of two parts.  The first routine
prepares the subevt without assigning momenta.  It takes the particles from an
{\tt{}interaction{\char95}t}.  It needs the indices and flavors for the beam,
incoming, and outgoing particles.

We can assume that the particle content of the subevt does not change.
Therefore, we set the event variables {\tt{}n{\char95}in}, {\tt{}n{\char95}out}, {\tt{}n{\char95}tot} already
in this initialization step.
\nwenddocs{}\nwbegincode{48}\moddef{Subevt expr: parton expr: TBP}\plusendmoddef
  procedure :: setup_subevt => parton_expr_setup_subevt
\nwendcode{}\nwbegincode{49}\moddef{Subevt expr: procedures}\plusendmoddef
  subroutine parton_expr_setup_subevt (expr, int, &
       i_beam, i_in, i_out, f_beam, f_in, f_out)
    class(parton_expr_t), intent(inout) :: expr
    type(interaction_t), intent(in), target :: int
    integer, dimension(:), intent(in) :: i_beam, i_in, i_out
    type(flavor_t), dimension(:), intent(in) :: f_beam, f_in, f_out
    allocate (expr%i_beam (size (i_beam)))
    allocate (expr%i_in (size (i_in)))
    allocate (expr%i_out (size (i_out)))
    expr%i_beam = i_beam
    expr%i_in = i_in
    expr%i_out = i_out
    call interaction_to_subevt (int, &
         expr%i_beam, expr%i_in, expr%i_out, expr%subevt_t)
    call subevt_set_pdg_beam     (expr%subevt_t, f_beam%get_pdg ())
    call subevt_set_pdg_incoming (expr%subevt_t, f_in%get_pdg ())
    call subevt_set_pdg_outgoing (expr%subevt_t, f_out%get_pdg ())
    call subevt_set_p2_beam     (expr%subevt_t, f_beam%get_mass () ** 2)
    call subevt_set_p2_incoming (expr%subevt_t, f_in%get_mass ()   ** 2)
    call subevt_set_p2_outgoing (expr%subevt_t, f_out%get_mass ()  ** 2)
    expr%n_in  = size (i_in)
    expr%n_out = size (i_out)
    expr%n_tot = expr%n_in + expr%n_out
  end subroutine parton_expr_setup_subevt

\eatline
\nwendcode{}\nwbegindocs{50}Transfer PDG codes, masses (initalization) and momenta to a
predefined subevent.  We use the flavor assignment of the first
branch in the interaction state matrix.  Only incoming and outgoing
particles are transferred.  Switch momentum sign for incoming
particles.
\nwenddocs{}\nwbegincode{51}\moddef{Subevt expr: public}\plusendmoddef
!  public :: interaction_to_subevt
!  public :: interaction_momenta_to_subevt
\nwendcode{}\nwbegincode{52}\moddef{Subevt expr: interfaces}\endmoddef
  interface interaction_momenta_to_subevt
     module procedure interaction_momenta_to_subevt_id
     module procedure interaction_momenta_to_subevt_tr
  end interface

\nwendcode{}\nwbegincode{53}\moddef{Subevt expr: procedures}\plusendmoddef
  subroutine interaction_to_subevt (int, j_beam, j_in, j_out, subevt)
    type(interaction_t), intent(in), target :: int
    integer, dimension(:), intent(in) :: j_beam, j_in, j_out
    type(subevt_t), intent(out) :: subevt
    type(flavor_t), dimension(:), allocatable :: flv
    integer :: n_beam, n_in, n_out, i, j
    allocate (flv (interaction_get_n_tot (int)))
    flv = quantum_numbers_get_flavor (interaction_get_quantum_numbers (int, 1))
    n_beam = size (j_beam)
    n_in = size (j_in)
    n_out = size (j_out)
    call subevt_init (subevt, n_beam + n_in + n_out)
    do i = 1, n_beam
       j = j_beam(i)
       call subevt_set_beam (subevt, i, &
            flv(j)%get_pdg (), &
            vector4_null, &
            flv(j)%get_mass () ** 2)
    end do
    do i = 1, n_in
       j = j_in(i)
       call subevt_set_incoming (subevt, n_beam + i, &
            flv(j)%get_pdg (), &
            vector4_null, &
            flv(j)%get_mass () ** 2)
    end do
    do i = 1, n_out
       j = j_out(i)
       call subevt_set_outgoing (subevt, n_beam + n_in + i, &
            flv(j)%get_pdg (), &
            vector4_null, &
            flv(j)%get_mass () ** 2)
    end do
  end subroutine interaction_to_subevt

  subroutine interaction_momenta_to_subevt_id (int, j_beam, j_in, j_out, subevt)
    type(interaction_t), intent(in) :: int
    integer, dimension(:), intent(in) :: j_beam, j_in, j_out
    type(subevt_t), intent(inout) :: subevt
    call subevt_set_p_beam &
         (subevt, - interaction_get_momenta (int, j_beam))
    call subevt_set_p_incoming &
         (subevt, - interaction_get_momenta (int, j_in))
    call subevt_set_p_outgoing &
         (subevt, interaction_get_momenta (int, j_out))
  end subroutine interaction_momenta_to_subevt_id

  subroutine interaction_momenta_to_subevt_tr &
       (int, j_beam, j_in, j_out, lt, subevt)
    type(interaction_t), intent(in) :: int
    integer, dimension(:), intent(in) :: j_beam, j_in, j_out
    type(subevt_t), intent(inout) :: subevt
    type(lorentz_transformation_t), intent(in) :: lt
    call subevt_set_p_beam &
         (subevt, - lt * interaction_get_momenta (int, j_beam))
    call subevt_set_p_incoming &
         (subevt, - lt * interaction_get_momenta (int, j_in))
    call subevt_set_p_outgoing &
         (subevt, lt * interaction_get_momenta (int, j_out))
  end subroutine interaction_momenta_to_subevt_tr

\eatline
\nwendcode{}\nwbegindocs{54}The second part takes the momenta from the interaction object and thus
completes the subevt.  The partonic energy can then be computed.
\nwenddocs{}\nwbegincode{55}\moddef{Subevt expr: parton expr: TBP}\plusendmoddef
  procedure :: fill_subevt => parton_expr_fill_subevt
\nwendcode{}\nwbegincode{56}\moddef{Subevt expr: procedures}\plusendmoddef
  subroutine parton_expr_fill_subevt (expr, int)
    class(parton_expr_t), intent(inout) :: expr
    type(interaction_t), intent(in), target :: int
    call interaction_momenta_to_subevt (int, &
         expr%i_beam, expr%i_in, expr%i_out, expr%subevt_t)
    expr%sqrts_hat = subevt_get_sqrts_hat (expr%subevt_t)
    expr%subevt_filled = .true.
  end subroutine parton_expr_fill_subevt
    
\eatline
\nwendcode{}\nwbegindocs{57}Evaluate, if the event passes the selection.  For absent expressions we take
default values.
\nwenddocs{}\nwbegincode{58}\moddef{Subevt expr: parton expr: TBP}\plusendmoddef
  procedure :: evaluate => parton_expr_evaluate
\nwendcode{}\nwbegincode{59}\moddef{Subevt expr: procedures}\plusendmoddef
  subroutine parton_expr_evaluate &
       (expr, passed, scale, fac_scale, ren_scale, weight, scale_forced)
    class(parton_expr_t), intent(inout) :: expr
    logical, intent(out) :: passed
    real(default), intent(out) :: scale
    real(default), intent(out) :: fac_scale
    real(default), intent(out) :: ren_scale
    real(default), intent(out) :: weight
    real(default), intent(in), allocatable, optional :: scale_forced
    logical :: force_scale
    force_scale = .false.
    if (present (scale_forced))  force_scale = allocated (scale_forced)
    call expr%base_evaluate (passed)
    if (passed) then
       if (force_scale) then
          scale = scale_forced
       else if (expr%has_scale) then
          call expr%scale%evaluate ()
          if (expr%scale%is_known ()) then
             scale = expr%scale%get_real ()
          else
             call msg_error ("Evaluate scale expression: result undefined")
             scale = 0
          end if
       else
          scale = expr%sqrts_hat
       end if
       if (force_scale) then
          fac_scale = scale_forced
       else if (expr%has_fac_scale) then
          call expr%fac_scale%evaluate ()
          if (expr%fac_scale%is_known ()) then
             fac_scale = expr%fac_scale%get_real ()
          else
             call msg_error ("Evaluate factorization scale expression: &
                  &result undefined")
             fac_scale = 0
          end if
       else
          fac_scale = scale
       end if
       if (force_scale) then
          ren_scale = scale_forced
       else if (expr%has_ren_scale) then
          call expr%ren_scale%evaluate ()
          if (expr%ren_scale%is_known ()) then
             ren_scale = expr%ren_scale%get_real ()
          else
             call msg_error ("Evaluate renormalization scale expression: &
                  &result undefined")
             ren_scale = 0
          end if
       else
          ren_scale = scale
       end if
       if (expr%has_weight) then
          call expr%weight%evaluate ()
          if (expr%weight%is_known ()) then
             weight = expr%weight%get_real ()
          else
             call msg_error ("Evaluate weight expression: result undefined")
             weight = 0
          end if
       else
          weight = 1
       end if
    end if
  end subroutine parton_expr_evaluate
  
\eatline
\nwendcode{}\nwbegindocs{60}Return the beam/incoming parton indices.
\nwenddocs{}\nwbegincode{61}\moddef{Subevt expr: parton expr: TBP}\plusendmoddef
  procedure :: get_beam_index => parton_expr_get_beam_index
  procedure :: get_in_index => parton_expr_get_in_index
\nwendcode{}\nwbegincode{62}\moddef{Subevt expr: procedures}\plusendmoddef
  subroutine parton_expr_get_beam_index (expr, i_beam)
    class(parton_expr_t), intent(in) :: expr
    integer, dimension(:), intent(out) :: i_beam
    i_beam = expr%i_beam
  end subroutine parton_expr_get_beam_index
  
  subroutine parton_expr_get_in_index (expr, i_in)
    class(parton_expr_t), intent(in) :: expr
    integer, dimension(:), intent(out) :: i_in
    i_in = expr%i_in
  end subroutine parton_expr_get_in_index
  
\eatline
\eatline
\nwendcode{}\nwbegindocs{63}\nwdocspar
\subsection{Implementation for full events}
This implementation contains the expressions that we can evaluate for the
full event.  It also contains data that pertain to the event, suitable
for communication with external event formats.  These data
simultaneously serve as pointer targets for the variable lists hidden
in the expressions (eval trees).

Squared matrix element and weight values: when reading events from
file, the {\tt{}ref} value is the number in the file, while the {\tt{}prc}
value is the number that we calculate from the momenta in the file,
possibly with different parameters.  When generating events the first
time, or if we do not recalculate, the numbers should coincide.
Furthermore, the array of {\tt{}alt} values is copied from an array of
alternative event records.  These values should represent calculated
values.
\nwenddocs{}\nwbegincode{64}\moddef{Subevt expr: public}\plusendmoddef
  public :: event_expr_t
\nwendcode{}\nwbegincode{65}\moddef{Subevt expr: types}\plusendmoddef
  type, extends (subevt_expr_t) :: event_expr_t
     logical :: has_reweight = .false.
     logical :: has_analysis = .false.
     class(expr_t), allocatable :: reweight
     class(expr_t), allocatable :: analysis
     logical :: has_id = .false.
     type(string_t) :: id
     logical :: has_num_id = .false.
     integer :: num_id = 0
     logical :: has_index = .false.
     integer :: index = 0
     logical :: has_sqme_ref = .false.
     real(default) :: sqme_ref = 0
     logical :: has_sqme_prc = .false.
     real(default) :: sqme_prc = 0
     logical :: has_weight_ref = .false.
     real(default) :: weight_ref = 0
     logical :: has_weight_prc = .false.
     real(default) :: weight_prc = 0
     logical :: has_excess_prc = .false.
     real(default) :: excess_prc = 0
     integer :: n_alt = 0
     logical :: has_sqme_alt = .false.
     real(default), dimension(:), allocatable :: sqme_alt
     logical :: has_weight_alt = .false.
     real(default), dimension(:), allocatable :: weight_alt
   contains
   \LA{}Subevt expr: event expr: TBP\RA{}
  end type event_expr_t
     
\eatline
\nwendcode{}\nwbegindocs{66}Finalizer for the expressions.
\nwenddocs{}\nwbegincode{67}\moddef{Subevt expr: event expr: TBP}\endmoddef
  procedure :: final => event_expr_final
\nwendcode{}\nwbegincode{68}\moddef{Subevt expr: procedures}\plusendmoddef
  subroutine event_expr_final (object)
    class(event_expr_t), intent(inout) :: object
    call object%base_final ()
    if (object%has_reweight) then
       call object%reweight%final ()
    end if
    if (object%has_analysis) then
       call object%analysis%final ()
    end if
  end subroutine event_expr_final

\eatline
\nwendcode{}\nwbegindocs{69}Output: continue writing the active expressions, after the common selection
expression.

Note: the {\tt{}prefix} argument is declared in the {\tt{}write} method of the
{\tt{}subevt{\char95}t} base type.  Here, it is unused.
\nwenddocs{}\nwbegincode{70}\moddef{Subevt expr: event expr: TBP}\plusendmoddef
  procedure :: write => event_expr_write
\nwendcode{}\nwbegincode{71}\moddef{Subevt expr: procedures}\plusendmoddef
  subroutine event_expr_write (object, unit, prefix, pacified)
    class(event_expr_t), intent(in) :: object
    integer, intent(in), optional :: unit
    character(*), intent(in), optional :: prefix
    logical, intent(in), optional :: pacified
    integer :: u
    u = given_output_unit (unit)
    call object%base_write (u, pacified = pacified)
    if (object%subevt_filled) then
       if (object%has_reweight) then
          call write_separator (u)
          write (u, "(1x,A)")  "Reweighting expression:"
          call write_separator (u)
          call object%reweight%write (u)
       end if
       if (object%has_analysis) then
          call write_separator (u)
          write (u, "(1x,A)")  "Analysis expression:"
          call write_separator (u)
          call object%analysis%write (u)
       end if
    end if
  end subroutine event_expr_write
    
\eatline
\nwendcode{}\nwbegindocs{72}Initializer.  This is required only for the {\tt{}sqme{\char95}alt} and
{\tt{}weight{\char95}alt} arrays.
\nwenddocs{}\nwbegincode{73}\moddef{Subevt expr: event expr: TBP}\plusendmoddef
  procedure :: init => event_expr_init
\nwendcode{}\nwbegincode{74}\moddef{Subevt expr: procedures}\plusendmoddef
  subroutine event_expr_init (expr, n_alt)
    class(event_expr_t), intent(out) :: expr
    integer, intent(in), optional :: n_alt
    if (present (n_alt)) then
       expr%n_alt = n_alt
       allocate (expr%sqme_alt (n_alt), source = 0._default)
       allocate (expr%weight_alt (n_alt), source = 0._default)
    end if
  end subroutine event_expr_init
  
\eatline
\nwendcode{}\nwbegindocs{75}Define variables.  We have the variables of the base type plus
specific variables for full events.  There is the event index.
\nwenddocs{}\nwbegincode{76}\moddef{Subevt expr: event expr: TBP}\plusendmoddef
  procedure :: setup_vars => event_expr_setup_vars
\nwendcode{}\nwbegincode{77}\moddef{Subevt expr: procedures}\plusendmoddef
  subroutine event_expr_setup_vars (expr, sqrts)
    class(event_expr_t), intent(inout), target :: expr
    real(default), intent(in) :: sqrts
    call expr%base_setup_vars (sqrts)
    call var_list_append_string_ptr (expr%var_list, &
         var_str ("$process_id"), expr%id, &
         is_known = expr%has_id, &
         locked = .true., verbose = .false., intrinsic = .true.)
    call var_list_append_int_ptr (expr%var_list, &
         var_str ("process_num_id"), expr%num_id, &
         is_known = expr%has_num_id, &
         locked = .true., verbose = .false., intrinsic = .true.)
    call var_list_append_real_ptr (expr%var_list, &
         var_str ("sqme"), expr%sqme_prc, &
         is_known = expr%has_sqme_prc, &
         locked = .true., verbose = .false., intrinsic = .true.)
    call var_list_append_real_ptr (expr%var_list, &
         var_str ("sqme_ref"), expr%sqme_ref, &
         is_known = expr%has_sqme_ref, &
         locked = .true., verbose = .false., intrinsic = .true.)
    call var_list_append_int_ptr (expr%var_list, &
         var_str ("event_index"), expr%index, &
         is_known = expr%has_index, &
         locked = .true., verbose = .false., intrinsic = .true.)
    call var_list_append_real_ptr (expr%var_list, &
         var_str ("event_weight"), expr%weight_prc, &
         is_known = expr%has_weight_prc, &
         locked = .true., verbose = .false., intrinsic = .true.)
    call var_list_append_real_ptr (expr%var_list, &
         var_str ("event_weight_ref"), expr%weight_ref, &
         is_known = expr%has_weight_ref, &
         locked = .true., verbose = .false., intrinsic = .true.)
    call var_list_append_real_ptr (expr%var_list, &
         var_str ("event_excess"), expr%excess_prc, &
         is_known = expr%has_excess_prc, &
         locked = .true., verbose = .false., intrinsic = .true.)
  end subroutine event_expr_setup_vars

\eatline
\nwendcode{}\nwbegindocs{78}Compile the analysis expression.  If the pointer is disassociated, there is
no expression.
\nwenddocs{}\nwbegincode{79}\moddef{Subevt expr: event expr: TBP}\plusendmoddef
  procedure :: setup_analysis => event_expr_setup_analysis
\nwendcode{}\nwbegincode{80}\moddef{Subevt expr: procedures}\plusendmoddef
  subroutine event_expr_setup_analysis (expr, ef_analysis)
    class(event_expr_t), intent(inout), target :: expr
    class(expr_factory_t), intent(in) :: ef_analysis
    call ef_analysis%build (expr%analysis)
    if (allocated (expr%analysis)) then
       call expr%setup_var_self ()
       call expr%analysis%setup_lexpr (expr%var_list)
       expr%has_analysis = .true.
    end if
  end subroutine event_expr_setup_analysis

\eatline
\nwendcode{}\nwbegindocs{81}Compile the reweight expression.
\nwenddocs{}\nwbegincode{82}\moddef{Subevt expr: event expr: TBP}\plusendmoddef
  procedure :: setup_reweight => event_expr_setup_reweight
\nwendcode{}\nwbegincode{83}\moddef{Subevt expr: procedures}\plusendmoddef
  subroutine event_expr_setup_reweight (expr, ef_reweight)
    class(event_expr_t), intent(inout), target :: expr
    class(expr_factory_t), intent(in) :: ef_reweight
    call ef_reweight%build (expr%reweight)
    if (allocated (expr%reweight)) then
       call expr%setup_var_self ()
       call expr%reweight%setup_expr (expr%var_list)
       expr%has_reweight = .true.
    end if
  end subroutine event_expr_setup_reweight

\eatline
\nwendcode{}\nwbegindocs{84}Store the string or numeric process ID.  This should be done during
initialization.
\nwenddocs{}\nwbegincode{85}\moddef{Subevt expr: event expr: TBP}\plusendmoddef
  procedure :: set_process_id => event_expr_set_process_id
  procedure :: set_process_num_id => event_expr_set_process_num_id
\nwendcode{}\nwbegincode{86}\moddef{Subevt expr: procedures}\plusendmoddef
  subroutine event_expr_set_process_id (expr, id)
    class(event_expr_t), intent(inout) :: expr
    type(string_t), intent(in) :: id
    expr%id = id
    expr%has_id = .true.
  end subroutine event_expr_set_process_id
    
  subroutine event_expr_set_process_num_id (expr, num_id)
    class(event_expr_t), intent(inout) :: expr
    integer, intent(in) :: num_id
    expr%num_id = num_id
    expr%has_num_id = .true.
  end subroutine event_expr_set_process_num_id
    
\eatline
\eatline
\nwendcode{}\nwbegindocs{87}Reset / set the data that pertain to a particular event.
\nwenddocs{}\nwbegincode{88}\moddef{Subevt expr: event expr: TBP}\plusendmoddef
  procedure :: reset => event_expr_reset
  procedure :: set => event_expr_set
\nwendcode{}\nwbegincode{89}\moddef{Subevt expr: procedures}\plusendmoddef
  subroutine event_expr_reset (expr)
    class(event_expr_t), intent(inout) :: expr
    call expr%base_reset ()
    expr%has_sqme_ref = .false.
    expr%has_sqme_prc = .false.
    expr%has_sqme_alt = .false.
    expr%has_weight_ref = .false.
    expr%has_weight_prc = .false.
    expr%has_weight_alt = .false.
    expr%has_excess_prc = .false.
  end subroutine event_expr_reset
  
  subroutine event_expr_set (expr, &
       weight_ref, weight_prc, weight_alt, &
       excess_prc, &
       sqme_ref, sqme_prc, sqme_alt)
    class(event_expr_t), intent(inout) :: expr
    real(default), intent(in), optional :: weight_ref, weight_prc
    real(default), intent(in), optional :: excess_prc
    real(default), intent(in), optional :: sqme_ref, sqme_prc
    real(default), dimension(:), intent(in), optional :: sqme_alt, weight_alt
    if (present (sqme_ref)) then
       expr%has_sqme_ref = .true.
       expr%sqme_ref = sqme_ref
    end if
    if (present (sqme_prc)) then
       expr%has_sqme_prc = .true.
       expr%sqme_prc = sqme_prc
    end if 
    if (present (sqme_alt)) then
       expr%has_sqme_alt = .true.
       expr%sqme_alt = sqme_alt
    end if
    if (present (weight_ref)) then
       expr%has_weight_ref = .true.
       expr%weight_ref = weight_ref
    end if
    if (present (weight_prc)) then
       expr%has_weight_prc = .true.
       expr%weight_prc = weight_prc
    end if
    if (present (weight_alt)) then
       expr%has_weight_alt = .true.
       expr%weight_alt = weight_alt
    end if
    if (present (excess_prc)) then
       expr%has_excess_prc = .true.
       expr%excess_prc = excess_prc
    end if
  end subroutine event_expr_set
  
\eatline
\nwendcode{}\nwbegindocs{90}Fill the event expression: take the particle data and kinematics
from a {\tt{}particle{\char95}set} object.

We allow the particle content to change for each event.  Therefore, we set the
event variables each time.

Also increment the event index; initialize it if necessary.
\nwenddocs{}\nwbegincode{91}\moddef{Subevt expr: event expr: TBP}\plusendmoddef
  procedure :: fill_subevt => event_expr_fill_subevt
\nwendcode{}\nwbegincode{92}\moddef{Subevt expr: procedures}\plusendmoddef
  subroutine event_expr_fill_subevt (expr, particle_set)
    class(event_expr_t), intent(inout) :: expr
    type(particle_set_t), intent(in) :: particle_set
    call particle_set%to_subevt (expr%subevt_t)
    expr%sqrts_hat = subevt_get_sqrts_hat (expr%subevt_t)
    expr%n_in  = subevt_get_n_in  (expr%subevt_t)
    expr%n_out = subevt_get_n_out (expr%subevt_t)
    expr%n_tot = expr%n_in + expr%n_out
    expr%subevt_filled = .true.
    if (expr%has_index) then
       expr%index = expr%index + 1
    else
       expr%index = 1
       expr%has_index = .true.
    end if
  end subroutine event_expr_fill_subevt
  
\eatline
\nwendcode{}\nwbegindocs{93}Evaluate, if the event passes the selection.  For absent expressions we take
default values.
\nwenddocs{}\nwbegincode{94}\moddef{Subevt expr: event expr: TBP}\plusendmoddef
  procedure :: evaluate => event_expr_evaluate
\nwendcode{}\nwbegincode{95}\moddef{Subevt expr: procedures}\plusendmoddef
  subroutine event_expr_evaluate (expr, passed, reweight, analysis_flag)
    class(event_expr_t), intent(inout) :: expr
    logical, intent(out) :: passed
    real(default), intent(out) :: reweight
    logical, intent(out) :: analysis_flag
    call expr%base_evaluate (passed)
    if (passed) then
       if (expr%has_reweight) then
          call expr%reweight%evaluate ()
          if (expr%reweight%is_known ()) then
             reweight = expr%reweight%get_real ()
          else
             call msg_error ("Evaluate reweight expression: &
                  &result undefined")
             reweight = 0
          end if
       else
          reweight = 1
       end if
       if (expr%has_analysis) then
          call expr%analysis%evaluate ()
          if (expr%analysis%is_known ()) then
             analysis_flag = expr%analysis%get_log ()
          else
             call msg_error ("Evaluate analysis expression: &
                  &result undefined")
             analysis_flag = .false.
          end if
       else
          analysis_flag = .true.
       end if
    end if
  end subroutine event_expr_evaluate
  
\eatline
\nwendcode{}\nwbegindocs{96}\nwdocspar
%%%%%%%%%%%%%%%%%%%%%%%%%%%%%%%%%%%%%%%%%%%%%%%%%%%%%%%%%%%%%%%%%%%%%%%%
\section{Parton states}
A {\tt{}parton{\char95}state{\char95}t} object contains the effective kinematics and
dynamics of an elementary partonic interaction, with or without the
beam/structure function state included.  The type is abstract and has
two distinct extensions.  The {\tt{}isolated{\char95}state{\char95}t} extension describes
the isolated elementary interaction where the {\tt{}int{\char95}eff} subobject
contains the complex transition amplitude, exclusive in all quantum
numbers.  The particle content and kinematics describe the effective
partonic state.  The {\tt{}connected{\char95}state{\char95}t} extension contains the
partonic {\tt{}subevt} and the expressions for cuts and scales which use
it.

In the isolated state, the effective partonic interaction may either
be identical to the hard interaction, in which case it is just a
pointer to the latter.  Or it may involve a rearrangement of partons,
in which case we allocate it explicitly and flag this by
{\tt{}int{\char95}is{\char95}allocated}.

The {\tt{}trace} evaluator contains the absolute square of the effective
transition amplitude matrix, summed over final states.  It is also summed over
initial states, depending on the the beam setup allows.  The result is used for
integration.

The {\tt{}matrix} evaluator is the counterpart of {\tt{}trace} which is kept
exclusive in all observable quantum numbers.  The {\tt{}flows} evaluator is
furthermore exclusive in colors, but neglecting all color interference.  The
{\tt{}matrix} and {\tt{}flows} evaluators are filled only for sampling points that
become part of physical events.

Note: It would be natural to make the evaluators allocatable.
However, this causes memory corruption in gfortran 4.6.3.  The extra
{\tt{}has{\char95}XXX} flags indicate whether evaluators are active, instead.

This module contains no unit tests.  The tests are covered by the
{\tt{}processes} module below.
\nwenddocs{}\nwbegincode{97}\moddef{\code{}parton{\char95}states.f90\edoc{}}\endmoddef
\LA{}File header\RA{}
module parton_states

\LA{}Use kinds\RA{}
\LA{}Use strings\RA{}
  use io_units
  use format_utils, only: write_separator
  use diagnostics
  use lorentz
  use subevents
  use variables
  use expr_base
  use model_data
  use flavors
  use helicities
  use colors
  use quantum_numbers
  use state_matrices
  use polarizations
  use interactions
  use evaluators

  use beams
  use sf_base
  use process_constants
  use prc_core
  use subevt_expr

\LA{}Standard module head\RA{}

\LA{}Parton states: public\RA{}

\LA{}Parton states: types\RA{}

contains
  
\LA{}Parton states: procedures\RA{}

end module parton_states
\eatline
\nwendcode{}\nwbegindocs{98}\nwdocspar
\subsection{Abstract base type}
The common part are the evaluators, one for the trace (summed over all
quantum numbers), one for the transition matrix (summed only over
unobservable quantum numbers), and one for the flow distribution
(transition matrix without interferences, exclusive in color flow).
\nwenddocs{}\nwbegincode{99}\moddef{Parton states: types}\endmoddef
  type, abstract :: parton_state_t
     logical :: has_trace = .false.
     logical :: has_matrix = .false.
     logical :: has_flows = .false.
     type(evaluator_t) :: trace
     type(evaluator_t) :: matrix
     type(evaluator_t) :: flows
   contains
   \LA{}Parton states: parton state: TBP\RA{}
  end type parton_state_t

\eatline
\nwendcode{}\nwbegindocs{100}The {\tt{}isolated{\char95}state{\char95}t} extension contains the {\tt{}sf{\char95}chain{\char95}eff} object
and the (hard) effective interaction {\tt{}int{\char95}eff}, separately, both
implemented as a pointer.  The evaluators (trace, matrix, flows) apply
to the hard interaction only.

If the effective interaction differs from the hard interaction, the
pointer is allocated explicitly.  Analogously for {\tt{}sf{\char95}chain{\char95}eff}.
\nwenddocs{}\nwbegincode{101}\moddef{Parton states: public}\endmoddef
  public :: isolated_state_t
\nwendcode{}\nwbegincode{102}\moddef{Parton states: types}\plusendmoddef
  type, extends (parton_state_t) :: isolated_state_t
     logical :: sf_chain_is_allocated = .false.
     type(sf_chain_instance_t), pointer :: sf_chain_eff => null ()
     logical :: int_is_allocated = .false.
     type(interaction_t), pointer :: int_eff => null ()
   contains
   \LA{}Parton states: isolated state: TBP\RA{}
  end type isolated_state_t

\eatline
\nwendcode{}\nwbegindocs{103}The {\tt{}connected{\char95}state{\char95}t} extension contains all data that enable
the evaluation of observables for the effective connected state.  The
evaluators connect the (effective) structure-function chain and hard
interaction that were kept separate in the {\tt{}isolated{\char95}state{\char95}t}.

The {\tt{}flows{\char95}sf} evaluator is an extended copy of the
structure-function 

The {\tt{}expr} subobject consists of the {\tt{}subevt}, a simple event record,
expressions for cuts etc.\ which refer to this record, and a {\tt{}var{\char95}list}
which contains event-specific variables, linked to the process variable
list.  Variables used within the expressions are looked up in {\tt{}var{\char95}list}.
\nwenddocs{}\nwbegincode{104}\moddef{Parton states: types}\plusendmoddef
  public :: connected_state_t
\nwendcode{}\nwbegincode{105}\moddef{Parton states: types}\plusendmoddef
  type, extends (parton_state_t) :: connected_state_t
     type(state_flv_content_t) :: state_flv
     logical :: has_flows_sf = .false.
     type(evaluator_t) :: flows_sf
     logical :: has_expr = .false.
     type(parton_expr_t) :: expr
   contains
   \LA{}Parton states: connected state: TBP\RA{}
  end type connected_state_t
     
\eatline
\nwendcode{}\nwbegindocs{106}Output: each evaluator is written only when it is active.  The
{\tt{}sf{\char95}chain} is only written if it is explicitly allocated.
\nwenddocs{}\nwbegincode{107}\moddef{Parton states: parton state: TBP}\endmoddef
  procedure :: write => parton_state_write
\nwendcode{}\nwbegincode{108}\moddef{Parton states: procedures}\endmoddef
  subroutine parton_state_write (state, unit, testflag)
    class(parton_state_t), intent(in) :: state
    integer, intent(in), optional :: unit
    logical, intent(in), optional :: testflag
    integer :: u
    u = given_output_unit (unit)
    select type (state)
    class is (isolated_state_t)
       if (state%sf_chain_is_allocated) then
          call write_separator (u)
          call state%sf_chain_eff%write (u)
       end if
       if (state%int_is_allocated) then
          call write_separator (u)
          write (u, "(1x,A)") &
               "Effective interaction:"
          call write_separator (u)
          call interaction_write (state%int_eff, u, testflag = testflag)
       end if
    class is (connected_state_t)
       if (state%has_flows_sf) then
          call write_separator (u)
          write (u, "(1x,A)") &
               "Evaluator (extension of the beam evaluator &
               &with color contractions):"
          call write_separator (u)
          call state%flows_sf%write (u, testflag = testflag)
       end if
    end select
    if (state%has_trace) then
       call write_separator (u)
       write (u, "(1x,A)") &
            "Evaluator (trace of the squared transition matrix):"
       call write_separator (u)
       call state%trace%write (u, testflag = testflag)
    end if
    if (state%has_matrix) then
       call write_separator (u)
       write (u, "(1x,A)") &
            "Evaluator (squared transition matrix):"
       call write_separator (u)
       call state%matrix%write (u, testflag = testflag)
    end if
    if (state%has_flows) then
       call write_separator (u)
       write (u, "(1x,A)") &
            "Evaluator (squared color-flow matrix):"
       call write_separator (u)
       call state%flows%write (u, testflag = testflag)
    end if
    select type (state)
    class is (connected_state_t)
       if (state%has_expr) then
          call write_separator (u)
          call state%expr%write (u)
       end if
    end select
  end subroutine parton_state_write
    
\eatline
\nwendcode{}\nwbegindocs{109}Finalize interaction and evaluators, but only if allocated.
\nwenddocs{}\nwbegincode{110}\moddef{Parton states: parton state: TBP}\plusendmoddef
  procedure :: final => parton_state_final
\nwendcode{}\nwbegincode{111}\moddef{Parton states: procedures}\plusendmoddef
  subroutine parton_state_final (state)
    class(parton_state_t), intent(inout) :: state
    if (state%has_flows) then
       call evaluator_final (state%flows)
       state%has_flows = .false.
    end if
    if (state%has_matrix) then
       call evaluator_final (state%matrix)
       state%has_matrix = .false.
    end if
    if (state%has_trace) then
       call evaluator_final (state%trace)
       state%has_trace = .false.
    end if
    select type (state)
    class is (connected_state_t)
       if (state%has_flows_sf) then
          call evaluator_final (state%flows_sf)
          state%has_flows_sf = .false.
       end if
       call state%expr%final ()
    class is (isolated_state_t)
       if (state%int_is_allocated) then
          call interaction_final (state%int_eff)
          deallocate (state%int_eff)
          state%int_is_allocated = .false.
       end if
       if (state%sf_chain_is_allocated) then
          call state%sf_chain_eff%final ()
       end if
    end select
  end subroutine parton_state_final
    
\eatline
\nwendcode{}\nwbegindocs{112}\nwdocspar
\subsection{Common Initialization}
Initialize the isolated parton state.  In this version, the
effective structure-function chain {\tt{}sf{\char95}chain{\char95}eff} and the effective
interaction {\tt{}int{\char95}eff} both are trivial pointers to the seed
structure-function chain and to the hard interaction, respectively.
\nwenddocs{}\nwbegincode{113}\moddef{Parton states: isolated state: TBP}\endmoddef
  procedure :: init => isolated_state_init_pointers
\nwendcode{}\nwbegincode{114}\moddef{Parton states: procedures}\plusendmoddef
  subroutine isolated_state_init_pointers (state, sf_chain, int)
    class(isolated_state_t), intent(out) :: state
    type(sf_chain_instance_t), intent(in), target :: sf_chain
    type(interaction_t), intent(in), target :: int
    state%sf_chain_eff => sf_chain
    state%int_eff => int
  end subroutine isolated_state_init_pointers
    
\eatline
\nwendcode{}\nwbegindocs{115}\nwdocspar
\subsection{Evaluator initialization: isolated state}
Create an evaluator for the trace of the squared transition matrix.
The trace goes over all outgoing quantum numbers.  Whether we trace
over incoming quantum numbers other than color, depends on the given
{\tt{}qn{\char95}mask{\char95}in}.

There are two options: explicitly computing the color factor table
({\tt{}use{\char95}cf} false; {\tt{}nc} defined), or taking the color factor
table from the hard matrix element data.
\nwenddocs{}\nwbegincode{116}\moddef{Parton states: isolated state: TBP}\plusendmoddef
  procedure :: setup_square_trace => isolated_state_setup_square_trace
\nwendcode{}\nwbegincode{117}\moddef{Parton states: procedures}\plusendmoddef
  subroutine isolated_state_setup_square_trace (state, core, qn_mask_in, &
       col)
    class(isolated_state_t), intent(inout), target :: state
    class(prc_core_t), intent(in) :: core
    type(quantum_numbers_mask_t), dimension(:), intent(in) :: qn_mask_in
    integer, dimension(:), intent(in) :: col
    type(quantum_numbers_mask_t), dimension(:), allocatable :: qn_mask    
    associate (data => core%data)
      allocate (qn_mask (data%n_in + data%n_out))
      qn_mask(:data%n_in) = &
              quantum_numbers_mask (.false., .true., .false.) &
              .or. qn_mask_in
      qn_mask(data%n_in+1:) = &
           quantum_numbers_mask (.true., .true., .true.)
    if (core%use_color_factors) then
       call evaluator_init_square (state%trace, &
            state%int_eff, qn_mask, &
            col_flow_index = data%cf_index, &
            col_factor = data%color_factors, &
            col_index_hi = col, &
            nc = core%nc)
    else
       call evaluator_init_square (state%trace, &
            state%int_eff, qn_mask, &
            nc = core%nc)
    end if
    end associate
    state%has_trace = .true.
  end subroutine isolated_state_setup_square_trace
    
\eatline
\nwendcode{}\nwbegindocs{118}Setup the evaluator for the transition matrix, exclusive in
helicities where this is requested.  

For all unstable final-state particles we keep polarization according to the
applicable decay options.  If the process is a decay itself, this applies also
to the initial state.

For all polarized final-state particles, we keep polarization including
off-diagonal entries.  We drop helicity completely for unpolarized final-state
particles.

For the initial state, if the particle has not been handled yet, we
apply the provided {\tt{}qn{\char95}mask{\char95}in} which communicates the beam properties.
\nwenddocs{}\nwbegincode{119}\moddef{Parton states: isolated state: TBP}\plusendmoddef
  procedure :: setup_square_matrix => isolated_state_setup_square_matrix
\nwendcode{}\nwbegincode{120}\moddef{Parton states: procedures}\plusendmoddef
  subroutine isolated_state_setup_square_matrix &
       (state, core, model, qn_mask_in, col)
    class(isolated_state_t), intent(inout), target :: state
    class(prc_core_t), intent(in) :: core
    class(model_data_t), intent(in), target :: model
    type(quantum_numbers_mask_t), dimension(:), intent(in) :: qn_mask_in
    integer, dimension(:), intent(in) :: col
    type(quantum_numbers_mask_t), dimension(:), allocatable :: qn_mask
    type(flavor_t), dimension(:), allocatable :: flv
    integer :: i
    logical :: helmask, helmask_hd
    associate (data => core%data)
      allocate (qn_mask (data%n_in + data%n_out))
      allocate (flv (data%n_flv))
      do i = 1, data%n_in + data%n_out      
         call flv%init (data%flv_state(i,:), model)
         if ((data%n_in == 1 .or. i > data%n_in) &
              .and. any (.not. flv%is_stable ())) then
            helmask = all (flv%decays_isotropically ())
            helmask_hd = all (flv%decays_diagonal ())
            qn_mask(i) = quantum_numbers_mask (.false., .true., helmask, &
                 mask_hd = helmask_hd)
         else if (i > data%n_in) then
            helmask = all (.not. flv%is_polarized ())
            qn_mask(i) = quantum_numbers_mask (.false., .true., helmask)
         else
            qn_mask(i) = quantum_numbers_mask (.false., .true., .false.) &
              .or. qn_mask_in(i)
         end if
      end do
    if (core%use_color_factors) then
       call evaluator_init_square (state%matrix, &
            state%int_eff, qn_mask, &
            col_flow_index = data%cf_index, &
            col_factor = data%color_factors, &
            col_index_hi = col, &
            nc = core%nc)
    else
       call evaluator_init_square (state%matrix, state%int_eff, &
            qn_mask, &
            nc = core%nc)
    end if
    end associate
    state%has_matrix = .true.
  end subroutine isolated_state_setup_square_matrix

\eatline
\nwendcode{}\nwbegindocs{121}This procedure initializes the evaluator that computes the
contributions to color flows, neglecting color interference.
The incoming-particle mask can be used to sum over incoming flavor.

Helicity handling: see above.
\nwenddocs{}\nwbegincode{122}\moddef{Parton states: isolated state: TBP}\plusendmoddef
  procedure :: setup_square_flows => isolated_state_setup_square_flows
\nwendcode{}\nwbegincode{123}\moddef{Parton states: procedures}\plusendmoddef
  subroutine isolated_state_setup_square_flows (state, core, model, qn_mask_in)
    class(isolated_state_t), intent(inout), target :: state
    class(prc_core_t), intent(in) :: core
    class(model_data_t), intent(in), target :: model
    type(quantum_numbers_mask_t), dimension(:), intent(in) :: qn_mask_in
    type(quantum_numbers_mask_t), dimension(:), allocatable :: qn_mask
    type(flavor_t), dimension(:), allocatable :: flv
    integer :: i
    logical :: helmask, helmask_hd
    associate (data => core%data)
      allocate (qn_mask (data%n_in + data%n_out))
      allocate (flv (data%n_flv))
      do i = 1, data%n_in + data%n_out
         call flv%init (data%flv_state(i,:), model)
         if ((data%n_in == 1 .or. i > data%n_in) &
              .and. any (.not. flv%is_stable ())) then
            helmask = all (flv%decays_isotropically ())
            helmask_hd = all (flv%decays_diagonal ())
            qn_mask(i) = quantum_numbers_mask (.false., .false., helmask, &
                 mask_hd = helmask_hd)
         else if (i > data%n_in) then
            helmask = all (.not. flv%is_polarized ())
            qn_mask(i) = quantum_numbers_mask (.false., .false., helmask)
         else
            qn_mask(i) = quantum_numbers_mask (.false., .false., .false.) &
              .or. qn_mask_in(i)
         end if
      end do
      call evaluator_init_square (state%flows, state%int_eff, qn_mask, &
           expand_color_flows = .true.)
    end associate
    state%has_flows = .true.
  end subroutine isolated_state_setup_square_flows

\eatline
\nwendcode{}\nwbegindocs{124}\nwdocspar
\subsection{Evaluator initialization: connected state}
Setup a trace evaluator as a product of two evaluators (incoming state,
effective interaction).  In the result, all quantum numbers are summed over.

If the optional {\tt{}int} interaction is provided, use this for the
first factor in the convolution.  Otherwise, use the final interaction
of the stored {\tt{}sf{\char95}chain}.

The {\tt{}resonant} flag applies if we want to construct
a decay chain.  The resonance property can propagate to the final
event output.
\nwenddocs{}\nwbegincode{125}\moddef{Parton states: connected state: TBP}\endmoddef
  procedure :: setup_connected_trace => connected_state_setup_connected_trace
\nwendcode{}\nwbegincode{126}\moddef{Parton states: procedures}\plusendmoddef
  subroutine connected_state_setup_connected_trace &
       (state, isolated, int, resonant)
    class(connected_state_t), intent(inout), target :: state
    type(isolated_state_t), intent(in), target :: isolated
    type(interaction_t), intent(in), optional, target :: int
    logical, intent(in), optional :: resonant
    type(quantum_numbers_mask_t) :: mask
    type(interaction_t), pointer :: src_int
    mask = quantum_numbers_mask (.true., .true., .true.)
    if (present (int)) then
       src_int => int
    else
       src_int => isolated%sf_chain_eff%get_out_int_ptr ()
    end if
    call evaluator_init_product &
         (state%trace, src_int, isolated%trace, &
          qn_mask_conn = mask, &
          qn_mask_rest = mask, &
          connections_are_resonant = resonant)
    state%has_trace = .true.
  end subroutine connected_state_setup_connected_trace
    
\eatline
\nwendcode{}\nwbegindocs{127}Setup a matrix evaluator as a product of two evaluators (incoming
state, effective interation).  In the intermediate state, color and
helicity is summed over.  In the final state, we keep the quantum
numbers which are present in the original evaluators.
\nwenddocs{}\nwbegincode{128}\moddef{Parton states: connected state: TBP}\plusendmoddef
  procedure :: setup_connected_matrix => connected_state_setup_connected_matrix
\nwendcode{}\nwbegincode{129}\moddef{Parton states: procedures}\plusendmoddef
  subroutine connected_state_setup_connected_matrix &
       (state, isolated, int, resonant, qn_filter_conn)
    class(connected_state_t), intent(inout), target :: state
    type(isolated_state_t), intent(in), target :: isolated
    type(interaction_t), intent(in), optional, target :: int
    logical, intent(in), optional :: resonant
    type(quantum_numbers_t), intent(in), optional :: qn_filter_conn
    type(quantum_numbers_mask_t) :: mask
    type(interaction_t), pointer :: src_int
    mask = quantum_numbers_mask (.false., .true., .true.)
    if (present (int)) then
       src_int => int
    else
       src_int => isolated%sf_chain_eff%get_out_int_ptr ()
    end if
    call evaluator_init_product &
         (state%matrix, src_int, isolated%matrix, mask, &
          qn_filter_conn = qn_filter_conn, &
          connections_are_resonant = resonant)
    state%has_matrix = .true.
  end subroutine connected_state_setup_connected_matrix
  
\eatline
\nwendcode{}\nwbegindocs{130}Setup a matrix evaluator as a product of two evaluators (incoming
state, effective interation).  In the intermediate state, only
helicity is summed over.  In the final state, we keep the quantum
numbers which are present in the original evaluators.


If the optional {\tt{}int} interaction is provided, use this for the
first factor in the convolution.  Otherwise, use the final interaction
of the stored {\tt{}sf{\char95}chain}, after creating an intermediate interaction
that includes a correlated color state.  We assume that for a
caller-provided {\tt{}int}, this is not necessary.
\nwenddocs{}\nwbegincode{131}\moddef{Parton states: connected state: TBP}\plusendmoddef
  procedure :: setup_connected_flows => connected_state_setup_connected_flows
\nwendcode{}\nwbegincode{132}\moddef{Parton states: procedures}\plusendmoddef
  subroutine connected_state_setup_connected_flows &
       (state, isolated, int, resonant, qn_filter_conn)
    class(connected_state_t), intent(inout), target :: state
    type(isolated_state_t), intent(in), target :: isolated
    type(interaction_t), intent(in), optional, target :: int
    logical, intent(in), optional :: resonant
    type(quantum_numbers_t), intent(in), optional :: qn_filter_conn
    type(quantum_numbers_mask_t) :: mask
    type(interaction_t), pointer :: src_int
    mask = quantum_numbers_mask (.false., .false., .true.)
    if (present (int)) then
       src_int => int
    else
       src_int => isolated%sf_chain_eff%get_out_int_ptr ()
       call evaluator_init_color_contractions (state%flows_sf, src_int)
       state%has_flows_sf = .true.
       src_int => evaluator_get_int_ptr (state%flows_sf)
    end if
    call evaluator_init_product &
         (state%flows, src_int, isolated%flows, mask, &
          qn_filter_conn = qn_filter_conn, &
          connections_are_resonant = resonant)
    state%has_flows = .true.
  end subroutine connected_state_setup_connected_flows
  
\eatline
\nwendcode{}\nwbegindocs{133}Determine and store the flavor content for the connected state.
This queries the {\tt{}matrix} evaluator component, which should hold the
requested flavor information.
\nwenddocs{}\nwbegincode{134}\moddef{Parton states: connected state: TBP}\plusendmoddef
  procedure :: setup_state_flv => connected_state_setup_state_flv
\nwendcode{}\nwbegincode{135}\moddef{Parton states: procedures}\plusendmoddef
  subroutine connected_state_setup_state_flv (state, n_out_hard)
    class(connected_state_t), intent(inout), target :: state
    integer, intent(in) :: n_out_hard
    type(interaction_t), pointer :: int
    int => evaluator_get_int_ptr (state%matrix)
    call interaction_get_flv_content (int, state%state_flv, n_out_hard)
  end subroutine connected_state_setup_state_flv
  
\eatline
\nwendcode{}\nwbegindocs{136}Return the current flavor state object.
\nwenddocs{}\nwbegincode{137}\moddef{Parton states: connected state: TBP}\plusendmoddef
  procedure :: get_state_flv => connected_state_get_state_flv
\nwendcode{}\nwbegincode{138}\moddef{Parton states: procedures}\plusendmoddef
  function connected_state_get_state_flv (state) result (state_flv)
    class(connected_state_t), intent(in) :: state
    type(state_flv_content_t) :: state_flv
    state_flv = state%state_flv
  end function connected_state_get_state_flv
  
\eatline
\nwendcode{}\nwbegindocs{139}\nwdocspar
\subsection{Cuts and expressions}
Set up the {\tt{}subevt} that corresponds to the connected interaction.
The index arrays refer to the interaction.

We assign the particles as follows: the beam particles are the first
two (decay process: one) entries in the trace evaluator.  The incoming
partons are identified by their link to the outgoing partons of the
structure-function chain.  The outgoing partons are those of the trace
evaluator, which include radiated partons during the
structure-function chain.
\nwenddocs{}\nwbegincode{140}\moddef{Parton states: connected state: TBP}\plusendmoddef
  procedure :: setup_subevt => connected_state_setup_subevt
\nwendcode{}\nwbegincode{141}\moddef{Parton states: procedures}\plusendmoddef
  subroutine connected_state_setup_subevt (state, sf_chain, f_beam, f_in, f_out)
    class(connected_state_t), intent(inout), target :: state
    type(sf_chain_instance_t), intent(in), target :: sf_chain
    type(flavor_t), dimension(:), intent(in) :: f_beam, f_in, f_out
    integer :: n_beam, n_in, n_out, n_vir, n_tot, i, j
    integer, dimension(:), allocatable :: i_beam, i_in, i_out
    integer :: sf_out_i
    type(interaction_t), pointer :: int, sf_int
    int => evaluator_get_int_ptr (state%trace)
    sf_int => sf_chain%get_out_int_ptr ()
    n_beam = size (f_beam)
    n_in = size (f_in)
    n_out = size (f_out)
    n_vir = interaction_get_n_vir (int)
    n_tot = interaction_get_n_tot (int)
    allocate (i_beam (n_beam), i_in (n_in), i_out (n_out))
    i_beam = [(i, i = 1, n_beam)]
    do j = 1, n_in
       sf_out_i = sf_chain%get_out_i (j)
       i_in(j) = interaction_find_link (int, sf_int, sf_out_i)
    end do
    i_out = [(i, i = n_vir + 1, n_tot)]
    call state%expr%setup_subevt (int, &
         i_beam, i_in, i_out, f_beam, f_in, f_out)
    state%has_expr = .true.
  end subroutine connected_state_setup_subevt

\eatline
\nwendcode{}\nwbegindocs{142}Initialize the variable list specific for this state/term.  We insert event
variables ({\tt{}sqrts{\char95}hat}) and link the process variable list.  The variable
list acquires pointers to subobjects of {\tt{}state}, which must therefore have a
{\tt{}target} attribute.
\nwenddocs{}\nwbegincode{143}\moddef{Parton states: connected state: TBP}\plusendmoddef
  procedure :: setup_var_list => connected_state_setup_var_list
\nwendcode{}\nwbegincode{144}\moddef{Parton states: procedures}\plusendmoddef
  subroutine connected_state_setup_var_list (state, process_var_list, beam_data)
    class(connected_state_t), intent(inout), target :: state
    type(var_list_t), intent(in), target :: process_var_list
    type(beam_data_t), intent(in) :: beam_data
    call state%expr%setup_vars (beam_data_get_sqrts (beam_data))
    call state%expr%link_var_list (process_var_list)
  end subroutine connected_state_setup_var_list
  
\eatline
\nwendcode{}\nwbegindocs{145}Allocate the cut expression etc.
\nwenddocs{}\nwbegincode{146}\moddef{Parton states: connected state: TBP}\plusendmoddef
  procedure :: setup_cuts => connected_state_setup_cuts
  procedure :: setup_scale => connected_state_setup_scale
  procedure :: setup_fac_scale => connected_state_setup_fac_scale
  procedure :: setup_ren_scale => connected_state_setup_ren_scale
  procedure :: setup_weight => connected_state_setup_weight
\nwendcode{}\nwbegincode{147}\moddef{Parton states: procedures}\plusendmoddef
  subroutine connected_state_setup_cuts (state, ef_cuts)
    class(connected_state_t), intent(inout), target :: state
    class(expr_factory_t), intent(in) :: ef_cuts
    call state%expr%setup_selection (ef_cuts)
  end subroutine connected_state_setup_cuts
    
  subroutine connected_state_setup_scale (state, ef_scale)
    class(connected_state_t), intent(inout), target :: state
    class(expr_factory_t), intent(in) :: ef_scale
    call state%expr%setup_scale (ef_scale)
  end subroutine connected_state_setup_scale
    
  subroutine connected_state_setup_fac_scale (state, ef_fac_scale)
    class(connected_state_t), intent(inout), target :: state
    class(expr_factory_t), intent(in) :: ef_fac_scale
    call state%expr%setup_fac_scale (ef_fac_scale)
  end subroutine connected_state_setup_fac_scale
    
  subroutine connected_state_setup_ren_scale (state, ef_ren_scale)
    class(connected_state_t), intent(inout), target :: state
    class(expr_factory_t), intent(in) :: ef_ren_scale
    call state%expr%setup_ren_scale (ef_ren_scale)
  end subroutine connected_state_setup_ren_scale
    
  subroutine connected_state_setup_weight (state, ef_weight)
    class(connected_state_t), intent(inout), target :: state
    class(expr_factory_t), intent(in) :: ef_weight
    call state%expr%setup_weight (ef_weight)
  end subroutine connected_state_setup_weight
    
\eatline
\nwendcode{}\nwbegindocs{148}Reset the expression object: invalidate the subevt.
\nwenddocs{}\nwbegincode{149}\moddef{Parton states: connected state: TBP}\plusendmoddef
  procedure :: reset_expressions => connected_state_reset_expressions
\nwendcode{}\nwbegincode{150}\moddef{Parton states: procedures}\plusendmoddef
  subroutine connected_state_reset_expressions (state)
    class(connected_state_t), intent(inout) :: state
    if (state%has_expr)  call state%expr%reset ()
  end subroutine connected_state_reset_expressions
  
\eatline
\nwendcode{}\nwbegindocs{151}\nwdocspar
\subsection{Evaluation}
Transfer momenta to the trace evaluator and fill the {\tt{}subevt} with
this effective kinematics, if applicable.

Note: we may want to apply a boost for the {\tt{}subevt}.
\nwenddocs{}\nwbegincode{152}\moddef{Parton states: parton state: TBP}\plusendmoddef
  procedure :: receive_kinematics => parton_state_receive_kinematics
\nwendcode{}\nwbegincode{153}\moddef{Parton states: procedures}\plusendmoddef
  subroutine parton_state_receive_kinematics (state)
    class(parton_state_t), intent(inout), target :: state
    type(interaction_t), pointer :: int
    if (state%has_trace) then
       call evaluator_receive_momenta (state%trace)
       select type (state)
       class is (connected_state_t)
          if (state%has_expr) then
             int => evaluator_get_int_ptr (state%trace)
             call state%expr%fill_subevt (int)
          end if
       end select
    end if
  end subroutine parton_state_receive_kinematics

\eatline
\nwendcode{}\nwbegindocs{154}Recover kinematics: We assume that the trace evaluator is filled
with momenta.  Send those momenta back to the sources, then fill the
variables and subevent as above.

The incoming momenta of the connected state are not connected to the
isolated state but to the beam interaction.  Therefore, the incoming
momenta within the isolated state do not become defined, yet.
Instead, we reconstruct the beam (and ISR) momentum configuration.
\nwenddocs{}\nwbegincode{155}\moddef{Parton states: parton state: TBP}\plusendmoddef
  procedure :: send_kinematics => parton_state_send_kinematics
\nwendcode{}\nwbegincode{156}\moddef{Parton states: procedures}\plusendmoddef
  subroutine parton_state_send_kinematics (state)
    class(parton_state_t), intent(inout), target :: state
    type(interaction_t), pointer :: int
    if (state%has_trace) then
       call evaluator_send_momenta (state%trace)
       select type (state)
       class is (connected_state_t)
          int => evaluator_get_int_ptr (state%trace)
          call state%expr%fill_subevt (int)
       end select
    end if
  end subroutine parton_state_send_kinematics

\eatline
\nwendcode{}\nwbegindocs{157}Evaluate the expressions.  The routine evaluates first the cut expression.
If the event passes, it evaluates the other expressions.  Where no expressions
are defined, default values are inserted.
\nwenddocs{}\nwbegincode{158}\moddef{Parton states: connected state: TBP}\plusendmoddef
  procedure :: evaluate_expressions => connected_state_evaluate_expressions
\nwendcode{}\nwbegincode{159}\moddef{Parton states: procedures}\plusendmoddef
  subroutine connected_state_evaluate_expressions (state, passed, &
       scale, fac_scale, ren_scale, weight, scale_forced)
    class(connected_state_t), intent(inout) :: state
    logical, intent(out) :: passed
    real(default), intent(out) :: scale, fac_scale, ren_scale, weight
    real(default), intent(in), allocatable, optional :: scale_forced
    if (state%has_expr) then
       call state%expr%evaluate (passed, scale, fac_scale, ren_scale, weight, &
            scale_forced)
    end if
  end subroutine connected_state_evaluate_expressions
    
\eatline
\nwendcode{}\nwbegindocs{160}Evaluate the structure-function chain, if it is allocated
explicitly.  The argument is the factorization scale.

If the chain is merely a pointer, the chain should already be
evaluated at this point.
\nwenddocs{}\nwbegincode{161}\moddef{Parton states: isolated state: TBP}\plusendmoddef
  procedure :: evaluate_sf_chain => isolated_state_evaluate_sf_chain
\nwendcode{}\nwbegincode{162}\moddef{Parton states: procedures}\plusendmoddef
  subroutine isolated_state_evaluate_sf_chain (state, fac_scale)
    class(isolated_state_t), intent(inout) :: state
    real(default), intent(in) :: fac_scale
    if (state%sf_chain_is_allocated) then
       call state%sf_chain_eff%evaluate (fac_scale)
    end if
  end subroutine isolated_state_evaluate_sf_chain
  
\eatline
\nwendcode{}\nwbegindocs{163}Evaluate the trace. 
\nwenddocs{}\nwbegincode{164}\moddef{Parton states: parton state: TBP}\plusendmoddef
  procedure :: evaluate_trace => parton_state_evaluate_trace
\nwendcode{}\nwbegincode{165}\moddef{Parton states: procedures}\plusendmoddef
  subroutine parton_state_evaluate_trace (state)
    class(parton_state_t), intent(inout) :: state
    if (state%has_trace) then
       call state%trace%evaluate ()
    end if
  end subroutine parton_state_evaluate_trace

\eatline
\nwendcode{}\nwbegindocs{166}Evaluate the extra evaluators that we need for physical events. 
\nwenddocs{}\nwbegincode{167}\moddef{Parton states: parton state: TBP}\plusendmoddef
  procedure :: evaluate_event_data => parton_state_evaluate_event_data
\nwendcode{}\nwbegincode{168}\moddef{Parton states: procedures}\plusendmoddef
  subroutine parton_state_evaluate_event_data (state)
    class(parton_state_t), intent(inout) :: state
    select type (state)
    type is (connected_state_t)
       if (state%has_flows_sf) then
          call evaluator_receive_momenta (state%flows_sf)
          call state%flows_sf%evaluate ()
       end if
    end select
    if (state%has_matrix) then
       call evaluator_receive_momenta (state%matrix)
       call state%matrix%evaluate ()
    end if
    if (state%has_flows) then
       call evaluator_receive_momenta (state%flows)
       call state%flows%evaluate ()
    end if
  end subroutine parton_state_evaluate_event_data

\eatline
\nwendcode{}\nwbegindocs{169}Normalize the helicity density matrix by its trace, i.e., factor out
the trace and put it into an overall normalization factor.  The trace
and flow evaluators are unchanged.
\nwenddocs{}\nwbegincode{170}\moddef{Parton states: parton state: TBP}\plusendmoddef
  procedure :: normalize_matrix_by_trace => &
       parton_state_normalize_matrix_by_trace
\nwendcode{}\nwbegincode{171}\moddef{Parton states: procedures}\plusendmoddef
  subroutine parton_state_normalize_matrix_by_trace (state)
    class(parton_state_t), intent(inout) :: state
    if (state%has_matrix) then
       call evaluator_normalize_by_trace (state%matrix)
    end if
  end subroutine parton_state_normalize_matrix_by_trace
  
\eatline
\nwendcode{}\nwbegindocs{172}\nwdocspar
\subsection{Accessing the state}
Three functions return a pointer to the event-relevant interactions.
\nwenddocs{}\nwbegincode{173}\moddef{Parton states: parton state: TBP}\plusendmoddef
  procedure :: get_trace_int_ptr => parton_state_get_trace_int_ptr
  procedure :: get_matrix_int_ptr => parton_state_get_matrix_int_ptr
  procedure :: get_flows_int_ptr => parton_state_get_flows_int_ptr
\nwendcode{}\nwbegincode{174}\moddef{Parton states: procedures}\plusendmoddef
  function parton_state_get_trace_int_ptr (state) result (ptr)
    class(parton_state_t), intent(in), target :: state
    type(interaction_t), pointer :: ptr
    if (state%has_trace) then
       ptr => evaluator_get_int_ptr (state%trace)
    else
       ptr => null ()
    end if
  end function parton_state_get_trace_int_ptr
  
  function parton_state_get_matrix_int_ptr (state) result (ptr)
    class(parton_state_t), intent(in), target :: state
    type(interaction_t), pointer :: ptr
    if (state%has_matrix) then
       ptr => evaluator_get_int_ptr (state%matrix)
    else
       ptr => null ()
    end if
  end function parton_state_get_matrix_int_ptr
  
  function parton_state_get_flows_int_ptr (state) result (ptr)
    class(parton_state_t), intent(in), target :: state
    type(interaction_t), pointer :: ptr
    if (state%has_flows) then
       ptr => evaluator_get_int_ptr (state%flows)
    else
       ptr => null ()
    end if
  end function parton_state_get_flows_int_ptr
  
\eatline
\eatline
\eatline
\nwendcode{}\nwbegindocs{175}Return the indices of the beam particles and the outgoing particles within
the trace (and thus, matrix and flows) evaluator, respectively.
\nwenddocs{}\nwbegincode{176}\moddef{Parton states: connected state: TBP}\plusendmoddef
  procedure :: get_beam_index => connected_state_get_beam_index
  procedure :: get_in_index => connected_state_get_in_index
\nwendcode{}\nwbegincode{177}\moddef{Parton states: procedures}\plusendmoddef
  subroutine connected_state_get_beam_index (state, i_beam)
    class(connected_state_t), intent(in) :: state
    integer, dimension(:), intent(out) :: i_beam
    call state%expr%get_beam_index (i_beam)
  end subroutine connected_state_get_beam_index
  
  subroutine connected_state_get_in_index (state, i_in)
    class(connected_state_t), intent(in) :: state
    integer, dimension(:), intent(out) :: i_in
    call state%expr%get_in_index (i_in)
  end subroutine connected_state_get_in_index
  
\eatline
\eatline
\nwendcode{}\nwbegindocs{178}Return the number of outgoing (hard) particles for the state.
\nwenddocs{}\nwbegincode{179}\moddef{Parton states: parton state: TBP}\plusendmoddef
  procedure :: get_n_out => parton_state_get_n_out
\nwendcode{}\nwbegincode{180}\moddef{Parton states: procedures}\plusendmoddef
  function parton_state_get_n_out (state) result (n)
    class(parton_state_t), intent(in), target :: state
    integer :: n
    type(interaction_t), pointer :: int
    int => evaluator_get_int_ptr (state%trace)
    n = interaction_get_n_out (int)
  end function parton_state_get_n_out
  
\eatline
\nwendcode{}\nwbegindocs{181}\nwdocspar
%%%%%%%%%%%%%%%%%%%%%%%%%%%%%%%%%%%%%%%%%%%%%%%%%%%%%%%%%%%%%%%%%%%%%%%%
\section{Complete Elementary Processes}
This module combines hard interactions, phase space, and (for
scatterings) structure functions and interfaces them to the
integration module.

The process object implements the combination of a fixed beam and
structure-function setup with a number of elementary processes.  The
latter are called process components.  The process object
represents an entity which is supposedly observable.  It should
be meaningful to talk about the cross section of a process.  

The individual components of a process are, technically, processes
themselves, but they may have unphysical cross sections which have to
be added for a physical result.  Process components may be exclusive
tree-level elementary processes, dipole subtraction term, loop
corrections, etc.

The beam and structure function setup is common to all process
components.  Thus, there is only one instance of this part.

The process may be a scattering process or a decay process.  In the
latter case, there are no structure functions, and the beam setup
consists of a single particle.  Otherwise, the two classes are treated
on the same footing.

Once a sampling point has been chosen, a process determines a set of
partons with a correlated density matrix of quantum numbers.  In general, each
sampling point will generate, for each process component, one or more
distinct parton configurations.  This is the {\tt{}computed} state.  The
computed state is the subject of the multi-channel integration
algorithm.

For NLO computations, it is necessary to project the computed states
onto another set of parton configurations (e.g., by recombining
certain pairs).  This is the {\tt{}observed} state.  When computing
partonic observables, the information is taken from the observed
state.

For the purpose of event generation, we will later select one parton
configuration from the observed state and collapse the correlated
quantum state.  This configuration is then dressed by applying parton
shower, decays and hadronization.  The decay chain, in particular,
combines a scattering process with possible subsequent decay processes
on the parton level, which are full-fledged process objects themselves.
\nwenddocs{}\nwbegincode{182}\moddef{\code{}processes.f90\edoc{}}\endmoddef
\LA{}File header\RA{}

module processes

\LA{}Use kinds\RA{}
\LA{}Use strings\RA{}
  use io_units
  use constants, only: twopi4
  use format_utils, only: write_separator
  use unit_tests
  use system_dependencies
  use diagnostics
  use md5
  use cputime
  use os_interface
  use physics_defs, only: conv
  use lorentz
  use sm_qcd
  use pdg_arrays
  use subevents
  use variables
  use expr_base
  use model_data
  use flavors
  use helicities
  use colors
  use quantum_numbers
  use state_matrices
  use polarizations
  use interactions
  use evaluators
  use particles
  use beam_structures
  use beams
  use sf_mappings
  use sf_base
  use process_constants
  use phs_base
  use phs_single
  use rng_base
  use mci_base
  use mci_midpoint
  use mci_vamp
  
  use vamp !NODEP!  

  use process_constants
  use prclib_interfaces
  use prc_core_def
  use process_libraries
  use prc_test

  use integration_results
  use prc_core
  use parton_states

  use nlo_data
  use phs_wood
  use phs_fks
  use prc_gosam
  use prc_omega
  use blha_config
  
\LA{}Standard module head\RA{}

\LA{}Processes: public\RA{}

\LA{}Processes: parameters\RA{}

\LA{}Processes: types\RA{}

\LA{}Processes: process part types\RA{}

\LA{}Processes: process type\RA{}

\LA{}Processes: interfaces\RA{}

\LA{}Processes: test types\RA{}

contains

\LA{}Processes: procedures\RA{}

\LA{}Processes: tests\RA{}

end module processes
\eatline
\nwendcode{}\nwbegindocs{183}\nwdocspar
\subsection{The Process Object}
A process object is the workspace for creating process instances for
event generation.  After initialization, its contents are filled by
integration passes which shape the integration grids and compute cross
sections.  Processes are set up initially from user-level
configuration data.  After calculating integrals and thus developing
integration grid data, the program may use a process
object or a copy of it for the purpose of generating events

The process object consists of several subobjects with their specific
purposes.  The corresponding types are defined below.  (Technically,
the subobject type definitions have to come before the process type
definition, but with NOWEB magic we reverse this order here.)

The {\tt{}type} determines whether we are considering a decay or a
scattering process.

The {\tt{}meta} object describes the process and its environment.  All
contents become fixed when the object is initialized.

The {\tt{}config} object holds physical and technical configuration data
that have been obtained during process initialization, and which are
common to all process components.

The individual process components are configured in the {\tt{}component}
objects.  These objects contain more configuration parameters and
workspace, as needed for the specific process variant.

The {\tt{}term} objects describe parton configurations which are
technically used as phase-space points.  Each process component may
split into several terms with distinct kinematics and particle
content.  Furthermore, each term may project on a different physical
state, e.g., by particle recombination.  The {\tt{}term} object provides
the framework for this projection, for applying cuts, weight, and thus
completing the process calculation.

The {\tt{}beam{\char95}config} object describes the incoming particles, either the
decay mother or the scattering beams.  It also contains the structure-function
information.

The {\tt{}mci{\char95}entry} objects configure a MC input parameter set and integrator,
each.  The number of parameters depends on the process component and on the
beam and structure-function setup.
\nwenddocs{}\nwbegincode{184}\moddef{Processes: public}\endmoddef
  public :: process_t
\nwendcode{}\nwbegincode{185}\moddef{Processes: process type}\endmoddef
  type :: process_t
     private
     type(process_metadata_t) :: &
          meta
     type(process_config_data_t) :: &
          config
     type(process_counter_t) :: &
          counter
     type(process_component_t), dimension(:), allocatable :: &
          component
     type(process_term_t), dimension(:), allocatable :: &
          term
     type(process_beam_config_t) :: &
          beam_config
     type(process_mci_entry_t), dimension(:), allocatable :: &
          mci_entry
     logical :: nlo_process
   contains
   \LA{}Processes: process: TBP\RA{}
  end type process_t

\eatline
\nwendcode{}\nwbegindocs{186}\nwdocspar
This procedure is an important debugging and inspection tool; it is
not used during normal operation.  The process object is written
to a file (identified by unit, which may also be standard output).
Optional flags determine whether we show everything or just the
interesting parts.
\nwenddocs{}\nwbegincode{187}\moddef{Processes: process: TBP}\endmoddef
  procedure :: write => process_write
\nwendcode{}\nwbegindocs{188}\nwdocspar
\nwenddocs{}\nwbegincode{189}\moddef{Processes: procedures}\endmoddef
  subroutine process_write (process, screen, unit, &
       show_all, show_var_list, &
       show_os_data, &
       show_rng_factory, show_model, show_expressions, &
       show_sfchain, &
       show_equivalences, show_history, show_histories, &
       show_forest, show_x, &
       show_subevt, show_evaluators, pacify)
    class(process_t), intent(in) :: process
    logical, intent(in) :: screen
    integer, intent(in), optional :: unit    
    logical, intent(in), optional :: show_all
    logical, intent(in), optional :: show_var_list
    logical, intent(in), optional :: show_os_data
    logical, intent(in), optional :: show_rng_factory
    logical, intent(in), optional :: show_model, show_expressions
    logical, intent(in), optional :: show_sfchain
    logical, intent(in), optional :: show_equivalences
    logical, intent(in), optional :: show_history, show_histories
    logical, intent(in), optional :: show_forest, show_x
    logical, intent(in), optional :: show_subevt, show_evaluators
    logical, intent(in), optional :: pacify
    logical :: all
    logical :: var_list
    logical :: counters
    logical :: os_data
    logical :: rng_factory, model, expressions
    integer :: u, i
    u = given_output_unit (unit)
    if (present (show_all)) then
       all = show_all
    else
       all = .false.
    end if
    var_list = .false.
    counters = .true.
    os_data = .false.
    model = .false.
    rng_factory = .true.
    expressions = .false.
    if (present (show_var_list)) then
       all = .false.; var_list = show_var_list
    end if
    if (present (show_os_data)) then
       all = .false.; os_data = show_os_data
    end if
    if (present (show_rng_factory)) then
       all = .false.; rng_factory = show_rng_factory
    end if
    if (present (show_model)) then
       all = .false.; model = show_model
    end if
    if (present (show_expressions)) then
       all = .false.; expressions = show_expressions
    end if
    if (all) then
       var_list = .true.
       rng_factory = .true.
       model = .true.
       expressions = .true.
    end if
    if (screen) then
       write (msg_buffer, "(A)")  repeat ("-", 72)
       call msg_message ()
    else
       call write_separator (u, 2)
    end if
    call process%meta%write (u, var_list, screen)
    if (process%meta%type == PRC_UNKNOWN) then
       call write_separator (u, 2)
       return
    else
       if (.not. screen)  call write_separator (u)
    end if
    if (screen)  return
    call process%config%write &
         (u, counters, os_data, rng_factory, model, expressions)
    call write_separator (u, 2)
    if (allocated (process%component)) then
       write (u, "(1x,A)") "Process component configuration:"
       do i = 1, size (process%component)
          call write_separator (u)
          call process%component(i)%write (u)
       end do
    else
       write (u, "(1x,A)") "Process component configuration: [undefined]"
    end if
    call write_separator (u, 2)
    if (allocated (process%term)) then
       write (u, "(1x,A)") "Process term configuration:"
       do i = 1, size (process%term)
          call write_separator (u)
          call process%term(i)%write (u)
       end do
    else
       write (u, "(1x,A)") "Process term configuration: [undefined]"
    end if
    call write_separator (u, 2)
    call process%beam_config%write (u)
    call write_separator (u, 2)
    if (allocated (process%mci_entry)) then
       write (u, "(1x,A)") "Multi-channel integrator configurations:"
       do i = 1, size (process%mci_entry)
          call write_separator (u)
          write (u, "(1x,A,I0,A)")  "MCI #", i, ":"
          call process%mci_entry(i)%write (u, pacify)
       end do
    end if
    call write_separator (u, 2)
  end subroutine process_write
      
\eatline
\nwendcode{}\nwbegindocs{190}Screen output.  Write a short account of the process configuration
and the current results.  The verbose version lists the components,
the short version just the results.
\nwenddocs{}\nwbegincode{191}\moddef{Processes: process: TBP}\plusendmoddef
  procedure :: show => process_show
\nwendcode{}\nwbegincode{192}\moddef{Processes: procedures}\plusendmoddef
  subroutine process_show (object, unit, verbose)
    class(process_t), intent(in) :: object
    integer, intent(in), optional :: unit
    logical, intent(in), optional :: verbose
    integer :: u
    logical :: verb
    u = given_output_unit (unit)
    verb = .true.;  if (present (verbose)) verb = verbose
    if (verb) then
       call object%meta%show (u, object%config%model%get_name ())
       select case (object%meta%type)
       case (PRC_DECAY)
          write (u, "(2x,A)", advance="no")  "Computed width ="
       case (PRC_SCATTERING)
          write (u, "(2x,A)", advance="no")  "Computed cross section ="
       case default;  return
       end select
    else 
       write (u, "(A)", advance="no") char (object%meta%id)
       select case (object%meta%num_id)
       case (0)
          write (u, "(':')")
       case default
          write (u, "(1x,'(',I0,')',':')") object%meta%num_id
       end select
       write (u, "(2x)", advance="no")
    end if
    write (u, "(ES14.7,1x,'+-',ES9.2)", advance="no") &
         object%get_integral_tot (), object%get_error_tot ()
    select case (object%meta%type)
    case (PRC_DECAY)
       write (u, "(1x,A)")  "GeV"
    case (PRC_SCATTERING)
       write (u, "(1x,A)")  "fb"
    case default
       write (u, *)
    end select
  end subroutine process_show
    
\eatline
\nwendcode{}\nwbegindocs{193}Finalizer.  Explicitly iterate over all subobjects that may contain
allocated pointers.
\nwenddocs{}\nwbegincode{194}\moddef{Processes: process: TBP}\plusendmoddef
  procedure :: final => process_final
\nwendcode{}\nwbegincode{195}\moddef{Processes: procedures}\plusendmoddef
  subroutine process_final (process)
    class(process_t), intent(inout) :: process
    integer :: i
    call process%meta%final ()
    call process%config%final ()
    if (allocated (process%component)) then
       do i = 1, size (process%component)
          call process%component(i)%final ()
       end do
    end if
    if (allocated (process%term)) then
       do i = 1, size (process%term)
          call process%term(i)%final ()
       end do
    end if
    call process%beam_config%final ()
    if (allocated (process%mci_entry)) then
       do i = 1, size (process%mci_entry)
          call process%mci_entry(i)%final ()
       end do
    end if
  end subroutine process_final
    
\eatline
\nwendcode{}\nwbegindocs{196}\nwdocspar
\subsubsection{Process setup}
Initialize a process.  We need a process library, the name of the process,
and a run ID.  The model is imported via its pointer, and the original pointer
freed.
\nwenddocs{}\nwbegincode{197}\moddef{Processes: process: TBP}\plusendmoddef
  procedure :: init => process_init
\nwendcode{}\nwbegincode{198}\moddef{Processes: procedures}\plusendmoddef
  subroutine process_init &
       (process, proc_id, run_id, lib, os_data, qcd, rng_factory, model)
    class(process_t), intent(out) :: process
    type(string_t), intent(in) :: proc_id
    type(string_t), intent(in) :: run_id
    type(process_library_t), intent(in), target :: lib
    type(os_data_t), intent(in) :: os_data
    type(qcd_t), intent(in) :: qcd
    class(rng_factory_t), intent(inout), allocatable :: rng_factory
    class(model_data_t), intent(inout), pointer :: model
    if (.not. lib%is_active ()) then
       call msg_bug ("Process init: inactive library not handled yet")
    end if
    if (.not. lib%contains (proc_id)) then
       call msg_fatal ("Process library doesn't contain process '" &
            // char (proc_id) // "'")
       return
    end if
    associate (meta => process%meta)
      call meta%init (proc_id, run_id, lib)
      call process%config%init (meta, os_data, qcd, rng_factory, model)
      allocate (process%component (meta%n_components))
    end associate
    process%nlo_process = lib%get_nlo_process (proc_id)
  end subroutine process_init
  
\eatline
\nwendcode{}\nwbegindocs{199}Store a snapshot of the common variable list.
\nwenddocs{}\nwbegincode{200}\moddef{Processes: process: TBP}\plusendmoddef
  procedure :: set_var_list => process_set_var_list
\nwendcode{}\nwbegincode{201}\moddef{Processes: procedures}\plusendmoddef
  subroutine process_set_var_list (process, var_list)
    class(process_t), intent(inout) :: process
    type(var_list_t), intent(in) :: var_list
    call var_list_init_snapshot &
         (process%meta%var_list, var_list, follow_link=.true.)
  end subroutine process_set_var_list
  
\eatline
\nwendcode{}\nwbegindocs{202}Initialize the process components, one by one, using a template for
the process core object.  The template is taken only for allocating
the correct type; the contents are set by extracting the process entry
from the library.
\nwenddocs{}\nwbegincode{203}\moddef{Processes: process: TBP}\plusendmoddef
  procedure :: init_component => process_init_component
\nwendcode{}\nwbegincode{204}\moddef{Processes: procedures}\plusendmoddef
  subroutine process_init_component &
       (process, index, core_template, mci_template, phs_config_template, &
        fks_template, gosam_template)
    class(process_t), intent(inout) :: process
    integer, intent(in) :: index
    class(prc_core_t), intent(in), allocatable :: core_template
    class(mci_t), intent(in), allocatable :: mci_template
    class(phs_config_t), intent(in), allocatable :: phs_config_template
    type(fks_template_t), intent(in), optional :: fks_template
    type(gosam_template_t), intent(inout), optional :: gosam_template
    call process%component(index)%init (index, &
         process%meta, process%config, &
         core_template, mci_template, phs_config_template, &
         gosam_template, fks_template)
    if (.not. process%component(index)%active) then
       call process%meta%deactivate_component(index)
    end if
  end subroutine process_init_component

\eatline
\nwendcode{}\nwbegindocs{205}Determine the process terms for each process component. 
\nwenddocs{}\nwbegincode{206}\moddef{Processes: process: TBP}\plusendmoddef
  procedure :: setup_terms => process_setup_terms
\nwendcode{}\nwbegincode{207}\moddef{Processes: procedures}\plusendmoddef
  subroutine process_setup_terms (process)
    class(process_t), intent(inout) :: process
    class(model_data_t), pointer :: model
    integer :: i, j, k
    integer, dimension(:), allocatable :: n_entry
    integer :: n_components, n_tot
    model => process%config%model
    n_components = process%meta%n_components
    allocate (n_entry (n_components), source = 0)
    do i = 1, n_components
       associate (component => process%component(i))
         if (component%active)  n_entry(i) = component%core%get_n_terms () 
       end associate
    end do
    n_tot = sum (n_entry)
    allocate (process%term (n_tot))
    k = 0
    do i = 1, n_components
       associate (component => process%component(i))
         if (.not. component%active)  cycle
         associate (core => component%core)
           allocate (component%i_term (n_entry(i)))
           do j = 1, n_entry(i)
              component%i_term(j) = k + j
              call process%term(k+j)%init (k+j, i, j, core, model)
           end do
         end associate
       end associate
       k = k + n_entry(i)
    end do
    process%config%n_terms = n_tot
  end subroutine process_setup_terms

\eatline
\nwendcode{}\nwbegindocs{208}Initialize the beam setup.  This is the trivial version where the
incoming state of the matrix element coincides with the initial state
of the process.  For a scattering process, we need the c.m. energy,
all other variables are set to their default values (no polarization,
lab frame and c.m.\ frame coincide, etc.)

We assume that all components consistently describe a scattering
process, i.e., two incoming particles.

Note: The current layout of the {\tt{}beam{\char95}data{\char95}t} record requires that the
flavor for each beam is unique.  For processes with multiple
flavors in the initial state, one has to set up beams explicitly.
This restriction could be removed by extending the code in the
{\tt{}beams} module.
\nwenddocs{}\nwbegincode{209}\moddef{Processes: process: TBP}\plusendmoddef
  procedure :: setup_beams_sqrts => process_setup_beams_sqrts
\nwendcode{}\nwbegincode{210}\moddef{Processes: procedures}\plusendmoddef
  subroutine process_setup_beams_sqrts (process, sqrts, beam_structure)
    class(process_t), intent(inout) :: process
    real(default), intent(in) :: sqrts
    type(beam_structure_t), intent(in), optional :: beam_structure
    type(pdg_array_t), dimension(:,:), allocatable :: pdg_in
    integer, dimension(2) :: pdg_scattering
    type(flavor_t), dimension(2) :: flv_in
    integer :: i, i0
    allocate (pdg_in (2, process%meta%n_components))
    i0 = 0
    do i = 1, process%meta%n_components
       if (process%component(i)%active) then
          pdg_in(:,i) = process%component(i)%get_pdg_in ()
          if (i0 == 0)  i0 = i
       end if
    end do
    do i = 1, process%meta%n_components
       if (.not. process%component(i)%active) then
          pdg_in(:,i) = pdg_in(:,i0)
       end if
    end do
    if (all (pdg_array_get_length (pdg_in) == 1) .and. &
         all (pdg_in(1,:) == pdg_in(1,i0)) .and. &
         all (pdg_in(2,:) == pdg_in(2,i0))) then
       pdg_scattering = pdg_array_get (pdg_in(:,i0), 1)
       call flv_in%init (pdg_scattering, process%config%model)
       call process%beam_config%init_scattering (flv_in, sqrts, beam_structure)
    else
       call msg_fatal ("Setting up process '" // char (process%meta%id) // "':", &
           [var_str ("   --------------------------------------------"), &
            var_str ("Inconsistent initial state. This happens if either "), &
            var_str ("several processes with non-matching initial states "), &
            var_str ("have been added, or for a single process with an "), &
            var_str ("initial state flavor sum. In that case, please set beams "), &
            var_str ("explicitly [singling out a flavor / structure function.]")]) 
    end if
  end subroutine process_setup_beams_sqrts

\eatline
\nwendcode{}\nwbegindocs{211}This is the version that applies to decay processes.  The energy is the
particle mass, hence no extra argument.
\nwenddocs{}\nwbegincode{212}\moddef{Processes: process: TBP}\plusendmoddef
  procedure :: setup_beams_decay => process_setup_beams_decay
\nwendcode{}\nwbegincode{213}\moddef{Processes: procedures}\plusendmoddef
  subroutine process_setup_beams_decay (process, rest_frame, beam_structure)
    class(process_t), intent(inout) :: process
    logical, intent(in), optional :: rest_frame
    type(beam_structure_t), intent(in), optional :: beam_structure
    type(pdg_array_t), dimension(:,:), allocatable :: pdg_in
    integer, dimension(1) :: pdg_decay
    type(flavor_t), dimension(1) :: flv_in
    integer :: i, i0
    allocate (pdg_in (1, process%meta%n_components))
    i0 = 0
    do i = 1, process%meta%n_components
       if (process%component(i)%active) then
          pdg_in(:,i) = process%component(i)%get_pdg_in ()
          if (i0 == 0)  i0 = i
       end if
    end do
    do i = 1, process%meta%n_components
       if (.not. process%component(i)%active) then
          pdg_in(:,i) = pdg_in(:,i0)
       end if
    end do
    if (all (pdg_array_get_length (pdg_in) == 1) &
         .and. all (pdg_in(1,:) == pdg_in(1,i0))) then
       pdg_decay = pdg_array_get (pdg_in(:,i0), 1)
       call flv_in%init (pdg_decay, process%config%model)
       call process%beam_config%init_decay (flv_in, rest_frame, beam_structure)
    else
       call msg_fatal ("Setting up decay '" &
            // char (process%meta%id) // "': decaying particle not unique")
    end if
  end subroutine process_setup_beams_decay
    
\eatline
\nwendcode{}\nwbegindocs{214}We have to make sure that the masses of the various flavors
in a given position in the particle string coincide.
\nwenddocs{}\nwbegincode{215}\moddef{Processes: process: TBP}\plusendmoddef
  procedure :: check_masses => process_check_masses
\nwendcode{}\nwbegincode{216}\moddef{Processes: procedures}\plusendmoddef
  subroutine process_check_masses (process)
       class(process_t), intent(in) :: process
       type(flavor_t), dimension(:), allocatable :: flv
       real(default), dimension(:), allocatable :: mass
       integer :: i, j
       do i = 1, process%meta%n_components
          if (.not. process%component(i)%active)  cycle
          associate (data => process%component(i)%core%data)
            allocate (flv (data%n_flv), mass (data%n_flv))
            do j = 1, data%n_in + data%n_out
               call flv%init (data%flv_state(j,:), process%config%model)
               mass = flv%get_mass ()
               if (any (mass /= mass(1))) then
                  call msg_fatal ("Process '" // char (process%meta%id) // "': " &
                       // "mass values in flavor combination do not coincide. ")
               end if
            end do
            deallocate (flv, mass)
          end associate
       end do
   end subroutine process_check_masses
  
\eatline
\nwendcode{}\nwbegindocs{217}For some structure functions we need to get the list of initial
state flavors.  This is a two-dimensional array.  The first index is
the beam index, the second index is the component index.  Each array
element is itself a PDG array object, which consists of the list of
incoming PDG values for this beam and component.
\nwenddocs{}\nwbegincode{218}\moddef{Processes: process: TBP}\plusendmoddef
  procedure :: get_pdg_in => process_get_pdg_in
\nwendcode{}\nwbegincode{219}\moddef{Processes: procedures}\plusendmoddef
  subroutine process_get_pdg_in (process, pdg_in)
    class(process_t), intent(in) :: process
    type(pdg_array_t), dimension(:,:), allocatable, intent(out) :: pdg_in
    integer :: i
    allocate (pdg_in (process%config%n_in, process%meta%n_components))
    do i = 1, process%meta%n_components
       if (process%component(i)%active) then
          pdg_in(:,i) = process%component(i)%get_pdg_in ()
       end if
    end do    
  end subroutine process_get_pdg_in
  
\eatline
\nwendcode{}\nwbegindocs{220}Initialize from a complete beam setup.
\nwenddocs{}\nwbegincode{221}\moddef{Processes: process: TBP}\plusendmoddef
  procedure :: setup_beams_beam_structure => process_setup_beams_beam_structure
\nwendcode{}\nwbegincode{222}\moddef{Processes: procedures}\plusendmoddef
  subroutine process_setup_beams_beam_structure &
       (process, beam_structure, sqrts, decay_rest_frame)
    class(process_t), intent(inout) :: process
    type(beam_structure_t), intent(in) :: beam_structure
    real(default), intent(in) :: sqrts
    logical, intent(in), optional :: decay_rest_frame
    if (process%get_n_in () == beam_structure%get_n_beam ()) then
       call process%beam_config%init_beam_structure &
            (beam_structure, sqrts, process%get_model_ptr (), decay_rest_frame)
    else if (beam_structure%get_n_beam () == 0) then
       call msg_fatal ("Asymmetric beams: missing beam particle specification")
    else
       call msg_fatal ("Mismatch of process and beam setup (scattering/decay)")
    end if
  end subroutine process_setup_beams_beam_structure
  
\eatline
\nwendcode{}\nwbegindocs{223}Notify the user about beam setup.
\nwenddocs{}\nwbegincode{224}\moddef{Processes: process: TBP}\plusendmoddef
  procedure :: beams_startup_message => process_beams_startup_message
\nwendcode{}\nwbegincode{225}\moddef{Processes: procedures}\plusendmoddef
  subroutine process_beams_startup_message (process, unit, beam_structure)
    class(process_t), intent(in) :: process
    integer, intent(in), optional :: unit
    type(beam_structure_t), intent(in), optional :: beam_structure
    call process%beam_config%startup_message (unit, beam_structure)
  end subroutine process_beams_startup_message
  
\eatline
\nwendcode{}\nwbegindocs{226}We complete the kinematics configuration after the beam setup, but before we
configure the chain of structure functions.  The reason is that we need the
total energy {\tt{}sqrts} for the kinematics, but the structure-function setup
requires the number of channels, which depends on the kinematics
configuration.  For instance, the kinematics module may return the need for
parameterizing an s-channel resonance.
\nwenddocs{}\nwbegincode{227}\moddef{Processes: process: TBP}\plusendmoddef
  procedure :: configure_phs => process_configure_phs
\nwendcode{}\nwbegincode{228}\moddef{Processes: procedures}\plusendmoddef
  subroutine process_configure_phs (process, rebuild, ignore_mismatch, verbose, &
                                    combined_integration)
    class(process_t), intent(inout) :: process
    logical, intent(in), optional :: rebuild
    logical, intent(in), optional :: ignore_mismatch, verbose
    logical, intent(in), optional :: combined_integration
    real(default) :: sqrts
    integer :: i, i_born
    class(phs_config_t), allocatable :: phs_config_born
    sqrts = process%get_sqrts ()
    do i = 1, process%meta%n_components
       associate (component => process%component(i))
         if (component%active) then
            if (process%nlo_process) then
               select case (char (component%config%get_nlo_type ()))
               case ('Born', 'Virtual', 'Subtraction')
                  call component%configure_phs (sqrts, process%beam_config, &
                       rebuild, ignore_mismatch, verbose=verbose)
                  call check_and_extend_phs (component, combined_integration)
               case ('Real')
                  i_born = component%config%get_associated_born ()
                  call check_and_extend_phs (component, combined_integration)
                  call process%component(i_born)%extract_phs_config (phs_config_born)
                  select type (phs_config_born)
                  type is (phs_wood_config_t)
                     call  component%configure_phs (sqrts, &
                          process%beam_config, rebuild, ignore_mismatch, &
                          phs_config_born, verbose=verbose)
                  end select
                  call process%component(i_born)%restore_phs_config (phs_config_born)
               end select
            else
               call component%configure_phs (sqrts, process%beam_config, &
                    rebuild, ignore_mismatch, verbose=verbose)
            end if
         end if
       end associate
    end do
  contains
    subroutine check_and_extend_phs (component, combined_integration)
      type(process_component_t), intent(inout) :: component
      logical, intent(in), optional :: combined_integration
      if (present (combined_integration)) then
         if (combined_integration) then
            select type (phs_config => component%phs_config)
            type is (phs_wood_config_t)
               call phs_config%set_extended_phs ()
               call phs_config%increase_n_par ()
            type is (phs_fks_config_t)
               call phs_config%set_extended_phs ()
            end select
         end if
      end if
    end subroutine check_and_extend_phs
  end subroutine process_configure_phs
         
\eatline
\nwendcode{}\nwbegindocs{229}Insert the structure-function configuration data.  First allocate the
storage, then insert data one by one.  The third procedure declares a
mapping (of the MC input parameters) for a specific channel and
structure-function combination.

We take the number of channels from the corresponding entry in the
{\tt{}config{\char95}data} section.

Otherwise, these a simple wrapper routines.  The extra level in the
call tree may allow for simple addressing of multiple concurrent beam
configurations, not implemented currently.

If we do not want structure functions, we simply do not call those procedures.
\nwenddocs{}\nwbegincode{230}\moddef{Processes: process: TBP}\plusendmoddef
  procedure :: init_sf_chain => process_init_sf_chain
  generic :: set_sf_channel => set_sf_channel_single
  procedure :: set_sf_channel_single => process_set_sf_channel
  generic :: set_sf_channel => set_sf_channel_array
  procedure :: set_sf_channel_array => process_set_sf_channel_array
\nwendcode{}\nwbegincode{231}\moddef{Processes: procedures}\plusendmoddef
  subroutine process_init_sf_chain (process, sf_config, sf_trace_file)
    class(process_t), intent(inout) :: process
    type(sf_config_t), dimension(:), intent(in) :: sf_config
    type(string_t), intent(in), optional :: sf_trace_file
    type(string_t) :: file
    if (present (sf_trace_file)) then
       if (sf_trace_file /= "") then
          file = sf_trace_file
       else
          file = process%get_id () // "_sftrace.dat"
       end if
       call process%beam_config%init_sf_chain (sf_config, file)
    else
       call process%beam_config%init_sf_chain (sf_config)
    end if
  end subroutine process_init_sf_chain

  subroutine process_set_sf_channel (process, c, sf_channel)
    class(process_t), intent(inout) :: process
    integer, intent(in) :: c
    type(sf_channel_t), intent(in) :: sf_channel
    call process%beam_config%set_sf_channel (c, sf_channel)
  end subroutine process_set_sf_channel
  
  subroutine process_set_sf_channel_array (process, sf_channel)
    class(process_t), intent(inout) :: process
    type(sf_channel_t), dimension(:), intent(in) :: sf_channel
    integer :: c
    call process%beam_config%allocate_sf_channels (size (sf_channel))
    do c = 1, size (sf_channel)
       call process%beam_config%set_sf_channel (c, sf_channel(c))
    end do
  end subroutine process_set_sf_channel_array
  
\eatline
\eatline
\nwendcode{}\nwbegindocs{232}Notify about the structure-function setup.
\nwenddocs{}\nwbegincode{233}\moddef{Processes: process: TBP}\plusendmoddef
  procedure :: sf_startup_message => process_sf_startup_message
\nwendcode{}\nwbegincode{234}\moddef{Processes: procedures}\plusendmoddef
  subroutine process_sf_startup_message (process, sf_string, unit)
    class(process_t), intent(in) :: process
    type(string_t), intent(in) :: sf_string
    integer, intent(in), optional :: unit
    call process%beam_config%sf_startup_message (sf_string, unit)
  end subroutine process_sf_startup_message
  
\eatline
\nwendcode{}\nwbegindocs{235}As soon as both the kinematics configuration and the
structure-function setup are complete, we match parameterizations
(channels) for both.  The matching entries are (re)set in the
{\tt{}component} phase-space configuration, while the structure-function
configuration is left intact.
\nwenddocs{}\nwbegincode{236}\moddef{Processes: process: TBP}\plusendmoddef
  procedure :: collect_channels => process_collect_channels
\nwendcode{}\nwbegincode{237}\moddef{Processes: procedures}\plusendmoddef
  subroutine process_collect_channels (process, coll)
    class(process_t), intent(inout) :: process
    type(phs_channel_collection_t), intent(inout) :: coll
    integer :: i
    do i = 1, process%meta%n_components
       associate (component => process%component(i))
         if (component%active) then
            call component%collect_channels (coll)
         end if
       end associate
    end do
  end subroutine process_collect_channels
         
\eatline
\nwendcode{}\nwbegindocs{238}Independently, we should be able to check if any component doesn't
contain phase-space parameters.  Such a process can only be integrated
if there are structure functions.
\nwenddocs{}\nwbegincode{239}\moddef{Processes: process: TBP}\plusendmoddef
  procedure :: contains_trivial_component => process_contains_trivial_component
\nwendcode{}\nwbegincode{240}\moddef{Processes: procedures}\plusendmoddef
  function process_contains_trivial_component (process) result (flag)
    class(process_t), intent(in) :: process
    logical :: flag
    integer :: i
    flag = .true.
    do i = 1, process%meta%n_components
       associate (component => process%component(i))
         if (component%active) then
            if (component%get_n_phs_par () == 0)  return
         end if
       end associate
    end do
    flag = .false.
  end function process_contains_trivial_component
  
\eatline
\nwendcode{}\nwbegindocs{241}Determine the MC parameter set structure and the MCI configuration for each
process component.  We need data from the structure-function and phase-space
setup, so those should be complete before this is called.  We also
make a random-number generator instance for each MCI group.
\nwenddocs{}\nwbegincode{242}\moddef{Processes: process: TBP}\plusendmoddef
  procedure :: setup_mci => process_setup_mci
\nwendcode{}\nwbegincode{243}\moddef{Processes: procedures}\plusendmoddef
  subroutine process_setup_mci (process, combined_integration)
    class(process_t), intent(inout) :: process
    logical, intent(in), optional :: combined_integration
    integer :: n_mci, i_mci
    integer :: i
    n_mci = 0
    do i = 1, process%meta%n_components
       associate (component => process%component(i))
          if (component%needs_mci_entry ()) then
            n_mci = n_mci + 1
            component%i_mci = n_mci
         end if
       end associate
    end do
    process%config%n_mci = n_mci
    if (.not. allocated (process%config%rng_factory)) &
         call msg_bug ("Process setup: rng factory not allocated")
    allocate (process%mci_entry (n_mci))
    i_mci = 0
    do i = 1, process%meta%n_components
       associate (component => process%component(i))
          if (component%needs_mci_entry () .and. &
              component%config%get_nlo_type () /= 'Subtraction') then
            i_mci = i_mci + 1
            associate (mci_entry => process%mci_entry(i_mci))
              call mci_entry%set_combined_integration (combined_integration)
              call mci_entry%init (process%meta%type, &
                   i_mci, i, component, process%beam_config, &
                   process%config%rng_factory)
            end associate
         end if
       end associate
    end do
    do i_mci = 1, size (process%mci_entry)
       call process%mci_entry(i_mci)%set_parameters (process%meta%var_list)
    end do
  end subroutine process_setup_mci
  
\eatline
\nwendcode{}\nwbegindocs{244}Set cuts.  This is a parse node, namely the right-hand side of the {\tt{}cut}
assignment.  When creating an instance, we compile this into an evaluation
tree.  The parse node may be null.
\nwenddocs{}\nwbegincode{245}\moddef{Processes: process: TBP}\plusendmoddef
  procedure :: set_cuts => process_set_cuts
\nwendcode{}\nwbegincode{246}\moddef{Processes: procedures}\plusendmoddef
  subroutine process_set_cuts (process, ef_cuts)
    class(process_t), intent(inout) :: process
    class(expr_factory_t), intent(in) :: ef_cuts
    allocate (process%config%ef_cuts, source = ef_cuts)
  end subroutine process_set_cuts
  
\eatline
\nwendcode{}\nwbegindocs{247}Analogously for the other expressions.
\nwenddocs{}\nwbegincode{248}\moddef{Processes: process: TBP}\plusendmoddef
  procedure :: set_scale => process_set_scale
  procedure :: set_fac_scale => process_set_fac_scale
  procedure :: set_ren_scale => process_set_ren_scale
  procedure :: set_weight => process_set_weight
\nwendcode{}\nwbegincode{249}\moddef{Processes: procedures}\plusendmoddef
  subroutine process_set_scale (process, ef_scale)
    class(process_t), intent(inout) :: process
    class(expr_factory_t), intent(in) :: ef_scale
    allocate (process%config%ef_scale, source = ef_scale)
  end subroutine process_set_scale
  
  subroutine process_set_fac_scale (process, ef_fac_scale)
    class(process_t), intent(inout) :: process
    class(expr_factory_t), intent(in) :: ef_fac_scale
    allocate (process%config%ef_fac_scale, source = ef_fac_scale)
  end subroutine process_set_fac_scale
  
  subroutine process_set_ren_scale (process, ef_ren_scale)
    class(process_t), intent(inout) :: process
    class(expr_factory_t), intent(in) :: ef_ren_scale
    allocate (process%config%ef_ren_scale, source = ef_ren_scale)
  end subroutine process_set_ren_scale
  
  subroutine process_set_weight (process, ef_weight)
    class(process_t), intent(inout) :: process
    class(expr_factory_t), intent(in) :: ef_weight
    allocate (process%config%ef_weight, source = ef_weight)
  end subroutine process_set_weight
  
\eatline
\eatline
\eatline
\eatline
\nwendcode{}\nwbegindocs{250}\nwdocspar
\subsubsection{MD5 sum}
The MD5 sum of the process object should reflect the state completely,
including integration results.  It is used for checking the integrity
of event files.  This global checksum includes checksums for the
various parts.  In particular, the MCI object receives a checksum that
includes the configuration of all configuration parts relevant for an
individual integration.  This checksum is used for checking the
integrity of integration grids.

We do not need MD5 sums for the process terms, since these are
generated from the component definitions.
\nwenddocs{}\nwbegincode{251}\moddef{Processes: process: TBP}\plusendmoddef
  procedure :: compute_md5sum => process_compute_md5sum
\nwendcode{}\nwbegincode{252}\moddef{Processes: procedures}\plusendmoddef
  subroutine process_compute_md5sum (process)
    class(process_t), intent(inout) :: process
    integer :: i
    call process%config%compute_md5sum ()
    do i = 1, process%config%n_components
       associate (component => process%component(i))
         if (component%active) then
            call component%compute_md5sum ()
         end if
       end associate
    end do
    call process%beam_config%compute_md5sum ()
    do i = 1, process%config%n_mci
       call process%mci_entry(i)%compute_md5sum &
            (process%config, process%component, process%beam_config)
    end do
  end subroutine process_compute_md5sum
  
\eatline
\nwendcode{}\nwbegindocs{253}\nwdocspar
\subsubsection{Integration and event generation}
The sampler test should just evaluate the squared matrix element {\tt{}n{\char95}calls}
times, discarding the results, and return.  This can be done before
integration, e.g., for timing estimates.
\nwenddocs{}\nwbegincode{254}\moddef{Processes: process: TBP}\plusendmoddef
  procedure :: sampler_test => process_sampler_test
\nwendcode{}\nwbegincode{255}\moddef{Processes: procedures}\plusendmoddef
  subroutine process_sampler_test (process, instance, i_mci, n_calls)
    class(process_t), intent(inout) :: process
    type(process_instance_t), intent(inout), target :: instance
    integer, intent(in) :: i_mci
    integer, intent(in) :: n_calls
    call process%mci_entry(i_mci)%sampler_test (instance, n_calls)
  end subroutine process_sampler_test

\eatline
\nwendcode{}\nwbegindocs{256}Integrate the process, using a previously initialized process
instance.  We select one of the available MCI integrators by its index
{\tt{}i{\char95}mci} and thus integrate over (structure functions and) phase
space for the associated (group of) process component(s).

The finalizer should be called after all integration passes have been
completed.  It will, for instance, write a summary of the integration
results.

{\tt{}integrate{\char95}dummy} does a ``dummy'' integration in the sense that
nothing is done but just empty integration results appended.
\nwenddocs{}\nwbegincode{257}\moddef{Processes: process: TBP}\plusendmoddef
  procedure :: integrate => process_integrate
  procedure :: final_integration => process_final_integration
  procedure :: integrate_dummy => process_integrate_dummy
\nwendcode{}\nwbegincode{258}\moddef{Processes: procedures}\plusendmoddef
  subroutine process_integrate (process, instance, i_mci, n_it, n_calls, &
       adapt_grids, adapt_weights, final, pacify)
    class(process_t), intent(inout) :: process
    type(process_instance_t), intent(inout) :: instance
    integer, intent(in) :: i_mci
    integer, intent(in) :: n_it
    integer, intent(in) :: n_calls
    logical, intent(in), optional :: adapt_grids
    logical, intent(in), optional :: adapt_weights
    logical, intent(in), optional :: final, pacify
    call process%mci_entry(i_mci)%integrate (instance, n_it, n_calls, &
         adapt_grids, adapt_weights, final, pacify, i_mci)        
    call instance%display_real_kinematics (i_mci)
  end subroutine process_integrate

  subroutine process_final_integration (process, i_mci)
    class(process_t), intent(inout) :: process
    integer, intent(in) :: i_mci
    call process%mci_entry(i_mci)%final_integration ()
  end subroutine process_final_integration
  
  subroutine process_integrate_dummy (process)
    class(process_t), intent(inout) :: process
    type(integration_results_t) :: results
    integer :: u_log
    u_log = logfile_unit ()
    call results%init (process%meta%type)
    call results%display_init (process%meta%type, screen = .true., unit = u_log)
    call results%new_pass ()
    call results%record (1, 0, 0._default, 0._default, 0._default)
    call results%display_final ()
  end subroutine process_integrate_dummy
  
\eatline
\eatline
\eatline
\nwendcode{}\nwbegindocs{259}Display the final results for the sum of all components.  (This is useful,
obviously, only if there is more than one component.)
\nwenddocs{}\nwbegincode{260}\moddef{Processes: process: TBP}\plusendmoddef
  procedure :: display_summed_results => process_display_summed_results
\nwendcode{}\nwbegincode{261}\moddef{Processes: procedures}\plusendmoddef
  subroutine process_display_summed_results (process)
    class(process_t), intent(inout) :: process
    type(integration_results_t) :: results
    integer :: u_log
    u_log = logfile_unit ()
    call results%init (process%meta%type)
    call results%display_init (process%meta%type, screen = .true., unit = u_log)
    call results%new_pass ()
    call results%record (1, 0, &
         process%get_integral (), &
         process%get_error (), &
         process%get_efficiency ())
    if (process%nlo_process) &
       call results%record_correction (process%get_correction (), &
                                        process%get_correction_error ())
    call results%display_final ()
  end subroutine process_display_summed_results

\eatline
\nwendcode{}\nwbegindocs{262}Run LaTeX/Metapost to generate a ps/pdf file for the integration
history.  We (re)write the driver file -- just in case it has been
missed before -- then we compile it.
\nwenddocs{}\nwbegincode{263}\moddef{Processes: process: TBP}\plusendmoddef
  procedure :: display_integration_history => &
       process_display_integration_history
\nwendcode{}\nwbegincode{264}\moddef{Processes: procedures}\plusendmoddef
  subroutine process_display_integration_history &
       (process, i_mci, filename, os_data, eff_reset)
    class(process_t), intent(inout) :: process
    integer, intent(in) :: i_mci
    type(string_t), intent(in) :: filename
    type(os_data_t), intent(in) :: os_data
    logical, intent(in), optional :: eff_reset
    call integration_results_write_driver &
         (process%mci_entry(i_mci)%results, filename, eff_reset)
    call integration_results_compile_driver &
         (process%mci_entry(i_mci)%results, filename, os_data)
  end subroutine process_display_integration_history

\eatline
\nwendcode{}\nwbegindocs{265}Write a complete logfile (with hardcoded name based on the process ID).
We do not write internal data.
\nwenddocs{}\nwbegincode{266}\moddef{Processes: process: TBP}\plusendmoddef
  procedure :: write_logfile => process_write_logfile
\nwendcode{}\nwbegincode{267}\moddef{Processes: procedures}\plusendmoddef
  subroutine process_write_logfile (process, i_mci, filename)
    class(process_t), intent(inout) :: process
    integer, intent(in) :: i_mci
    type(string_t), intent(in) :: filename    
    type(time_t) :: time
    integer :: unit, u
    unit = free_unit ()
    open (unit = unit, file = char (filename), action = "write", &
          status = "replace")
    u = given_output_unit (unit)
    write (u, "(A)")  repeat ("#", 79)
    call process%meta%write (u, .false., .false.)
    write (u, "(A)")  repeat ("#", 79)
    write (u, "(3x,A,ES17.10)")  "Integral   = ", &
         process%mci_entry(i_mci)%get_integral ()
    write (u, "(3x,A,ES17.10)")  "Error      = ", &
         process%mci_entry(i_mci)%get_error ()
    write (u, "(3x,A,ES17.10)")  "Accuracy   = ", &
         process%mci_entry(i_mci)%get_accuracy ()
    write (u, "(3x,A,ES17.10)")  "Chi2       = ", &
         process%mci_entry(i_mci)%get_chi2 ()
    write (u, "(3x,A,ES17.10)")  "Efficiency = ", &
         process%mci_entry(i_mci)%get_efficiency ()
    call process%mci_entry(i_mci)%get_time (time, 10000)
    if (time%is_known ()) then
       write (u, "(3x,A,1x,A)")  "T(10k evt) = ", char (time%to_string_dhms ())
    else
       write (u, "(3x,A)")  "T(10k evt) =  [undefined]"
    end if
    call process%mci_entry(i_mci)%results%write (u)
    write (u, "(A)")  repeat ("#", 79)
    call process%mci_entry(i_mci)%results%write_chain_weights (u)
    write (u, "(A)")  repeat ("#", 79)
    call process%mci_entry(i_mci)%counter%write (u)
    write (u, "(A)")  repeat ("#", 79)
    select type (mci => process%mci_entry(i_mci)%mci)
    type is (mci_midpoint_t)
       write (u, "(1x,A)")  "MC Integrator is Midpoint rule"
    type is (mci_vamp_t)
       write (u, "(1x,A)")  "MC Integrator is VAMP"       
       call write_separator (u)
       call mci%write_history (u)
       call write_separator (u)       
       if (mci%grid_par%use_vamp_equivalences) then
          call vamp_equivalences_write (mci%equivalences, u)          
       else
          write (u, "(3x,A)") "No VAMP equivalences have been used"
       end if
       call write_separator (u)
       call process%mci_entry(i_mci)%write_chain_weights (u) 
    class default
       write (u, "(1x,A)")  "MC Integrator: [unknown]"
    end select
    write (u, "(A)")  repeat ("#", 79)
    call process%beam_config%data%write (u)
    write (u, "(A)")  repeat ("#", 79)
    if (allocated (process%config%ef_cuts)) then
       write (u, "(3x,A)") "Cut expression:"
       call process%config%ef_cuts%write (u)
    else
       write (u, "(3x,A)") "No cuts used."         
    end if
    call write_separator (u)           
    if (allocated (process%config%ef_scale)) then
       write (u, "(3x,A)") "Scale expression:"
       call process%config%ef_scale%write (u)
    else
       write (u, "(3x,A)") "No scale expression was given."
    end if
    call write_separator (u)           
    if (allocated (process%config%ef_fac_scale)) then
       write (u, "(3x,A)") "Factorization scale expression:"
       call process%config%ef_fac_scale%write (u)
    else
       write (u, "(3x,A)") "No factorization scale expression was given."       
    end if
    call write_separator (u)           
    if (allocated (process%config%ef_ren_scale)) then
       write (u, "(3x,A)") "Renormalization scale expression:"
       call process%config%ef_ren_scale%write (u)
    else
       write (u, "(3x,A)") "No renormalization scale expression was given."
    end if
    call write_separator (u)           
    if (allocated (process%config%ef_weight)) then
       call write_separator (u)
       write (u, "(3x,A)") "Weight expression:"
       call process%config%ef_weight%write (u)
    else
       write (u, "(3x,A)") "No weight expression was given."
    end if
    write (u, "(A)")  repeat ("#", 79)
    write (u, "(1x,A)") "Summary of quantum-number states:" 
    write (u, "(1x,A)")  " + sign: allowed and contributing"
    write (u, "(1x,A)")  " no +  : switched off at runtime"
    call process%write_state_summary (u)
    write (u, "(A)")  repeat ("#", 79)
    write (u, "(A)")  "Variable list:"
    call var_list_write (process%meta%var_list, u)
    write (u, "(A)")  repeat ("#", 79)
    close (u)
  end subroutine process_write_logfile

\eatline
\nwendcode{}\nwbegindocs{268}Display the quantum-number combinations of the process components, and their
current status (allowed or switched off).
\nwenddocs{}\nwbegincode{269}\moddef{Processes: process: TBP}\plusendmoddef
  procedure :: write_state_summary => process_write_state_summary
\nwendcode{}\nwbegincode{270}\moddef{Processes: procedures}\plusendmoddef
  subroutine process_write_state_summary (process, unit)
    class(process_t), intent(in) :: process
    integer, intent(in), optional :: unit
    integer :: i, i_component, u
    u = given_output_unit (unit)
    do i = 1, size (process%term)
       call write_separator (u)
       i_component = process%term(i)%i_component
       if (i_component /= 0) then
          call process%term(i)%write_state_summary &
               (process%component(i_component)%core, unit)
       end if
    end do
  end subroutine process_write_state_summary
       
\eatline
\nwendcode{}\nwbegindocs{271}Prepare event generation for the specified MCI entry.  This implies, in
particular, checking the phase-space file.
\nwenddocs{}\nwbegincode{272}\moddef{Processes: process: TBP}\plusendmoddef
  procedure :: prepare_simulation => process_prepare_simulation
\nwendcode{}\nwbegincode{273}\moddef{Processes: procedures}\plusendmoddef
  subroutine process_prepare_simulation (process, i_mci)
    class(process_t), intent(inout) :: process
    integer, intent(in) :: i_mci
    call process%mci_entry(i_mci)%prepare_simulation ()
  end subroutine process_prepare_simulation

\eatline
\nwendcode{}\nwbegindocs{274}\nwdocspar
Generate a weighted event.  We select one of the available MCI
integrators by its index {\tt{}i{\char95}mci} and thus generate an event for the
associated (group of) process component(s).  The arguments exactly
correspond to the initializer and finalizer above.

The resulting event is stored in the {\tt{}prcoess{\char95}instance} object,
which also holds the workspace of the integrator.

Note: The {\tt{}process} object is declared {\tt{}intent(inout)} because it
contains the random-number state, which changes for each event.
Otherwise, all volatile data are inside the {\tt{}instance} object.
\nwenddocs{}\nwbegincode{275}\moddef{Processes: process: TBP}\plusendmoddef
  procedure :: generate_weighted_event => process_generate_weighted_event
  procedure :: generate_unweighted_event => process_generate_unweighted_event
\nwendcode{}\nwbegincode{276}\moddef{Processes: procedures}\plusendmoddef
  subroutine process_generate_weighted_event (process, instance, i_mci)
    class(process_t), intent(inout) :: process
    type(process_instance_t), intent(inout) :: instance
    integer, intent(in) :: i_mci
    call process%mci_entry(i_mci)%generate_weighted_event (instance)
  end subroutine process_generate_weighted_event

  subroutine process_generate_unweighted_event (process, instance, i_mci)
    class(process_t), intent(inout) :: process
    type(process_instance_t), intent(inout) :: instance
    integer, intent(in) :: i_mci
    call process%mci_entry(i_mci)%generate_unweighted_event (instance)
  end subroutine process_generate_unweighted_event

\eatline
\eatline
\nwendcode{}\nwbegindocs{277}\nwdocspar
This replaces the event generation methods for the situation that the
process instance object has been filled by other means (i.e., reading
and/or recalculating its contents).  We just have to fill in missing
MCI data, especially the event weight.
\nwenddocs{}\nwbegincode{278}\moddef{Processes: process: TBP}\plusendmoddef
  procedure :: recover_event => process_recover_event
\nwendcode{}\nwbegincode{279}\moddef{Processes: procedures}\plusendmoddef
  subroutine process_recover_event (process, instance, i_term)
    class(process_t), intent(inout) :: process
    type(process_instance_t), intent(inout) :: instance
    integer, intent(in) :: i_term
    call process%mci_entry(instance%i_mci)%recover_event (instance, i_term)
  end subroutine process_recover_event

\eatline
\nwendcode{}\nwbegindocs{280}\nwdocspar
\subsubsection{Retrieve process data}
Tell whether integral (and error) are known.
\nwenddocs{}\nwbegincode{281}\moddef{Processes: process: TBP}\plusendmoddef
  generic :: has_integral => has_integral_tot, has_integral_mci
  procedure :: has_integral_tot => process_has_integral_tot
  procedure :: has_integral_mci => process_has_integral_mci
\nwendcode{}\nwbegincode{282}\moddef{Processes: procedures}\plusendmoddef
  function process_has_integral_mci (process, i_mci) result (flag)
    class(process_t), intent(in) :: process
    integer, intent(in) :: i_mci
    logical :: flag
    flag = process%mci_entry(i_mci)%has_integral ()
  end function process_has_integral_mci

  function process_has_integral_tot (process) result (flag)
    class(process_t), intent(in) :: process
    logical :: flag
    integer :: i
    flag = .true.
    do i = 1, size (process%mci_entry)
       if (process%is_active_nlo_component (i)) &
          flag = flag .and. process%mci_entry(i)%has_integral ()
    end do
  end function process_has_integral_tot
  
\eatline
\nwendcode{}\nwbegindocs{283}\nwdocspar
Return the current integral and error obtained by the integrator {\tt{}i{\char95}mci}.
\nwenddocs{}\nwbegincode{284}\moddef{Processes: process: TBP}\plusendmoddef
  generic :: get_integral => get_integral_tot, get_integral_mci
  generic :: get_error => get_error_tot, get_error_mci
  generic :: get_efficiency => get_efficiency_tot, get_efficiency_mci
  procedure :: get_integral_tot => process_get_integral_tot
  procedure :: get_integral_mci => process_get_integral_mci
  procedure :: get_error_tot => process_get_error_tot
  procedure :: get_error_mci => process_get_error_mci
  procedure :: get_efficiency_tot => process_get_efficiency_tot
  procedure :: get_efficiency_mci => process_get_efficiency_mci
\nwendcode{}\nwbegincode{285}\moddef{Processes: procedures}\plusendmoddef
  function process_get_integral_mci (process, i_mci) result (integral)
    class(process_t), intent(in) :: process
    integer, intent(in) :: i_mci
    real(default) :: integral
    integral = process%mci_entry(i_mci)%get_integral ()
  end function process_get_integral_mci
  
  function process_get_error_mci (process, i_mci) result (error)
    class(process_t), intent(in) :: process
    integer, intent(in) :: i_mci
    real(default) :: error
    error = process%mci_entry(i_mci)%get_error ()
  end function process_get_error_mci
  
  function process_get_efficiency_mci (process, i_mci) result (efficiency)
    class(process_t), intent(in) :: process
    integer, intent(in) :: i_mci
    real(default) :: efficiency
    efficiency = process%mci_entry(i_mci)%get_efficiency ()
  end function process_get_efficiency_mci
  
  function process_get_integral_tot (process) result (integral)
    class(process_t), intent(in) :: process
    real(default) :: integral
    integer :: i
    integral = 0
    if (allocated (process%mci_entry)) then
       do i = 1, size (process%mci_entry)
          if (process%is_active_nlo_component(i)) &
               integral = integral + process%mci_entry(i)%get_integral ()
       end do
    end if
  end function process_get_integral_tot
  
  function process_get_error_tot (process) result (error)
    class(process_t), intent(in) :: process
    real(default) :: error
    real(default) :: variance
    integer :: i
    variance = 0
    if (allocated (process%mci_entry)) then
       do i = 1, size (process%mci_entry)
          if (process%is_active_nlo_component(i)) &
               variance = variance + process%mci_entry(i)%get_error () ** 2
       end do
    end if
    error = sqrt (variance)
  end function process_get_error_tot
  
  function process_get_efficiency_tot (process) result (efficiency)
    class(process_t), intent(in) :: process
    real(default) :: efficiency
    real(default) :: den, eff, int
    integer :: i
    den = 0
    if (allocated (process%mci_entry)) then
       do i = 1, size (process%mci_entry)
          if (process%is_active_nlo_component(i)) then
             int = process%get_integral (i)
             if (int > 0) then
                eff = process%mci_entry(i)%get_efficiency ()
                if (eff > 0) then
                   den = den + int / eff
                else
                   efficiency = 0
                   return
                end if
             end if
          end if
       end do
    end if
    if (den > 0) then
       efficiency = process%get_integral () / den
    else
       efficiency = 0
    end if
  end function process_get_efficiency_tot
  
\eatline
\nwendcode{}\nwbegincode{286}\moddef{Processes: process: TBP}\plusendmoddef
  procedure :: get_correction => process_get_correction
  procedure :: get_correction_error => process_get_correction_error
\nwendcode{}\nwbegincode{287}\moddef{Processes: procedures}\plusendmoddef
  function process_get_correction (process) result (ratio)
    class(process_t), intent(in) :: process
    real(default) :: ratio
    real(default) :: int_real, int_virt, int_born
    int_born = process%mci_entry(1)%get_integral ()
    int_real = process%mci_entry(2)%get_integral ()
    int_virt = process%mci_entry(3)%get_integral ()
    ratio = (int_real + int_virt) / int_born
  end function process_get_correction

  function process_get_correction_error (process) result (error)
    class(process_t), intent(in) :: process
    real(default) :: error
    real(default) :: int_born, int_real, int_virt
    real(default) :: err_born, err_real, err_virt
    real(default) :: err2
    int_born = process%mci_entry(1)%get_integral ()
    int_real = process%mci_entry(2)%get_integral ()
    int_virt = process%mci_entry(3)%get_integral ()
    err_born = process%mci_entry(1)%get_error ()
    err_real = process%mci_entry(2)%get_error ()
    err_virt = process%mci_entry(3)%get_error ()
    err2 = (int_real + int_virt)**2 / int_born**4 * err_born**2 &
             + (err_real**2 + err_virt**2) / int_born**2
    error = sqrt (err2)
  end function process_get_correction_error

\eatline
\nwendcode{}\nwbegindocs{288}\nwdocspar
\nwenddocs{}\nwbegincode{289}\moddef{Processes: process: TBP}\plusendmoddef
  procedure :: component_reset_helicity => process_component_reset_helicity
\nwendcode{}\nwbegincode{290}\moddef{Processes: procedures}\plusendmoddef
  subroutine process_component_reset_helicity(proc, i)
    class(process_t), intent(inout) :: proc
    integer, intent(in) :: i
    select type (core => proc%component(i)%core)
      type is (prc_omega_t)
        call core%reset_helicity_selection ()
    end select
    if (i == 2) then
      select type (core_sub_born => proc%component(i)%core_sub_born)
      type is (prc_omega_t)
         call core_sub_born%reset_helicity_selection ()
      end select
    end if
  end subroutine process_component_reset_helicity

\eatline
\nwendcode{}\nwbegindocs{291}Initialize the simultaneous Born calculation.
\nwenddocs{}\nwbegincode{292}\moddef{Processes: process: TBP}\plusendmoddef
  procedure :: init_sub_born => process_init_sub_born
\nwendcode{}\nwbegincode{293}\moddef{Processes: procedures}\plusendmoddef
  subroutine process_init_sub_born (proc, i_component, core_born_in)
    class(process_t), intent(inout) :: proc
    integer, intent(in) :: i_component
    class(prc_core_t), intent(inout) :: core_born_in
    associate (component => proc%component(i_component))
      allocate (component%core_sub_born, source = core_born_in)
    end associate
  end subroutine process_init_sub_born

\eatline
\nwendcode{}\nwbegincode{294}\moddef{Processes: process: TBP}\plusendmoddef
  procedure :: is_nlo_calculation => process_is_nlo_calculation
\nwendcode{}\nwbegincode{295}\moddef{Processes: procedures}\plusendmoddef
  function process_is_nlo_calculation (process) result (nlo)
    class(process_t), intent(in) :: process
    logical :: nlo
    nlo = process%nlo_process
  end function process_is_nlo_calculation

\eatline
\nwendcode{}\nwbegindocs{296}Return nlo data of a process component
\nwenddocs{}\nwbegincode{297}\moddef{Processes: process: TBP}\plusendmoddef
  procedure :: get_component_nlo_type => process_get_component_nlo_type
\nwendcode{}\nwbegincode{298}\moddef{Processes: procedures}\plusendmoddef
  function process_get_component_nlo_type (process, i_component) &
           result (nlo_type)
    class(process_t), intent(in) :: process
    integer, intent(in) :: i_component
    type(string_t) :: nlo_type
    nlo_type = process%component(i_component)%config%get_nlo_type ()
  end function process_get_component_nlo_type

\nwendcode{}\nwbegincode{299}\moddef{Processes: process: TBP}\plusendmoddef
  procedure :: get_component_associated_born &
            => process_get_component_associated_born
\nwendcode{}\nwbegincode{300}\moddef{Processes: procedures}\plusendmoddef
  function process_get_component_associated_born (process, i_component) &
           result (i_born)
    class(process_t), intent(in) :: process
    integer, intent(in) :: i_component
    integer :: i_born
    i_born = process%component(i_component)%config%get_associated_born ()
  end function process_get_component_associated_born

\nwendcode{}\nwbegincode{301}\moddef{Processes: process: TBP}\plusendmoddef
  procedure :: is_active_nlo_component => process_is_active_nlo_component 
\nwendcode{}\nwbegincode{302}\moddef{Processes: procedures}\plusendmoddef
  function process_is_active_nlo_component (process, i_component) &
           result (active)
    class(process_t), intent(in) :: process
    integer, intent(in) :: i_component
    logical :: active
    associate (component => process%component(i_component))
       active = component%config%is_active_nlo_component () &
                .and. component%component_type <= COMP_MASTER
    end associate
  end function process_is_active_nlo_component

\eatline
\eatline
\nwendcode{}\nwbegindocs{303}Return the MD5 sums that summarize the process component
definitions.  These values should be independent of parameters, beam
details, expressions, etc.  They can be used for checking the
integrity of a process when reusing an old event file.
\nwenddocs{}\nwbegincode{304}\moddef{Processes: process: TBP}\plusendmoddef
  procedure :: get_md5sum_prc => process_get_md5sum_prc
\nwendcode{}\nwbegincode{305}\moddef{Processes: procedures}\plusendmoddef
  function process_get_md5sum_prc (process, i_component) result (md5sum)
    class(process_t), intent(in) :: process
    integer, intent(in) :: i_component
    character(32) :: md5sum
    if (process%component(i_component)%active) then
       md5sum = process%component(i_component)%config%get_md5sum ()
    else
       md5sum = ""
    end if
  end function process_get_md5sum_prc
    
\eatline
\nwendcode{}\nwbegindocs{306}Return the MD5 sums that summarize the state of the MCI integrators.
These values should encode all process data, integration and phase
space configuration, etc., and the integration results.  They can thus
be used for checking the integrity of an event-generation setup when
reusing an old event file.
\nwenddocs{}\nwbegincode{307}\moddef{Processes: process: TBP}\plusendmoddef
  procedure :: get_md5sum_mci => process_get_md5sum_mci
\nwendcode{}\nwbegincode{308}\moddef{Processes: procedures}\plusendmoddef
  function process_get_md5sum_mci (process, i_mci) result (md5sum)
    class(process_t), intent(in) :: process
    integer, intent(in) :: i_mci
    character(32) :: md5sum
    md5sum = process%mci_entry(i_mci)%get_md5sum ()
  end function process_get_md5sum_mci
    
\eatline
\nwendcode{}\nwbegindocs{309}Return the MD5 sum of the process configuration.  This should encode
the process setup, data, and expressions, but no integration results.
\nwenddocs{}\nwbegincode{310}\moddef{Processes: process: TBP}\plusendmoddef
  procedure :: get_md5sum_cfg => process_get_md5sum_cfg
\nwendcode{}\nwbegincode{311}\moddef{Processes: procedures}\plusendmoddef
  function process_get_md5sum_cfg (process) result (md5sum)
    class(process_t), intent(in) :: process
    character(32) :: md5sum
    md5sum = process%config%md5sum
  end function process_get_md5sum_cfg

\eatline
\nwendcode{}\nwbegindocs{312}\nwdocspar
\subsection{Metadata}
This information describes the process and its environment.  It is
fixed upon initialization.

The {\tt{}id} string is the name of the process object, as given by the
user.  The matrix element generator will use this string for naming
Fortran procedures and types, so it should qualify as a Fortran name.

The {\tt{}num{\char95}id} is meaningful if nonzero.  It is used for communication
with external programs or file standards which do not support string IDs.

The {\tt{}run{\char95}id} string distinguishes among several runs for the same
process.  It identifies process instances with respect to adapted
integration grids and similar run-specific data.  The run ID is kept
when copying processes for creating instances, however, so it does not
distinguish event samples. 

The {\tt{}var{\char95}list} is a snapshot of the variable list, taken at the
point where the process was initialized.

The {\tt{}lib} pointer accesses the process library where the process
definition and the process driver are located.

The {\tt{}lib{\char95}index} is the index of entry in the process library that
corresponds to the current process.

The {\tt{}component{\char95}id} array identifies the individual process components.

The {\tt{}component{\char95}description} is an array of human-readable strings
that characterize the process components, for instance {\tt{}a,\ b\ =>\ c,\ d}.

The {\tt{}active} mask array marks those components which are active.  The others
are skipped.
\nwenddocs{}\nwbegincode{313}\moddef{Processes: process part types}\endmoddef
  type :: process_metadata_t
     private
     integer :: type = PRC_UNKNOWN
     type(string_t) :: id
     integer :: num_id = 0
     type(string_t) :: run_id
     type(var_list_t) :: var_list
     type(process_library_t), pointer :: lib => null ()
     integer :: lib_index = 0
     integer :: n_components = 0
     type(string_t), dimension(:), allocatable :: component_id
     type(string_t), dimension(:), allocatable :: component_description
     logical, dimension(:), allocatable :: active
   contains
   \LA{}Processes: process metadata: TBP\RA{}
  end type process_metadata_t

\eatline
\nwendcode{}\nwbegindocs{314}The local var list is a snapshot and needs a finalizer.
\nwenddocs{}\nwbegincode{315}\moddef{Processes: process metadata: TBP}\endmoddef
  procedure :: final => process_metadata_final
\nwendcode{}\nwbegincode{316}\moddef{Processes: procedures}\plusendmoddef
  subroutine process_metadata_final (meta)
    class(process_metadata_t), intent(inout) :: meta
    call meta%var_list%final (follow_link=.true.)
  end subroutine process_metadata_final
  
\eatline
\nwendcode{}\nwbegindocs{317}Output: ID and run ID.
We write the variable list only upon request.
\nwenddocs{}\nwbegincode{318}\moddef{Processes: process metadata: TBP}\plusendmoddef
  procedure :: write => process_metadata_write
\nwendcode{}\nwbegincode{319}\moddef{Processes: procedures}\plusendmoddef
  subroutine process_metadata_write (meta, u, var_list, screen)
    class(process_metadata_t), intent(in) :: meta
    integer, intent(in) :: u
    logical, intent(in) :: var_list, screen    
    integer :: i
    select case (meta%type)
    case (PRC_UNKNOWN)
       if (screen) then
          write (msg_buffer, "(A)") "Process [undefined]"
       else
          write (u, "(1x,A)") "Process [undefined]"
       end if
       return
    case (PRC_DECAY)
       if (screen) then
          write (msg_buffer, "(A,1x,A,A,A)") "Process [decay]:", & 
               "'", char (meta%id), "'"
       else
          write (u, "(1x,A)", advance="no") "Process [decay]:"
       end if
    case (PRC_SCATTERING)
       if (screen) then
          write (msg_buffer, "(A,1x,A,A,A)") "Process [scattering]:", &
               "'", char (meta%id), "'"
       else
          write (u, "(1x,A)", advance="no") "Process [scattering]:"
       end if
    case default
       call msg_bug ("process_write: undefined process type")
    end select    
    if (screen)  then
       call msg_message ()
    else
       write (u, "(1x,A,A,A)") "'", char (meta%id), "'"
    end if
    if (meta%num_id /= 0) then
       if (screen) then
          write (msg_buffer, "(2x,A,I0)") "ID (num)      = ", meta%num_id
          call msg_message ()
       else
          write (u, "(3x,A,I0)") "ID (num)      = ", meta%num_id            
       end if
    end if
    if (screen) then
       if (meta%run_id /= "") then
          write (msg_buffer, "(2x,A,A,A)") "Run ID        = '", &
               char (meta%run_id), "'"
          call msg_message ()
       end if
    else
       write (u, "(3x,A,A,A)") "Run ID        = '", char (meta%run_id), "'"       
    end if
    if (associated (meta%lib)) then
       if (screen) then
          write (msg_buffer, "(2x,A,A,A)")  "Library name  = '", &
               char (meta%lib%get_name ()), "'"
          call msg_message ()
       else          
          write (u, "(3x,A,A,A)")  "Library name  = '", &
               char (meta%lib%get_name ()), "'"          
       end if
    else
       if (screen) then
          write (msg_buffer, "(2x,A)")  "Library name  = [not associated]"
          call msg_message ()
       else
          write (u, "(3x,A)")  "Library name  = [not associated]"
       end if
    end if
    if (screen) then
       write (msg_buffer, "(2x,A,I0)")  "Process index = ", meta%lib_index
       call msg_message ()
    else
       write (u, "(3x,A,I0)")  "Process index = ", meta%lib_index
    end if
    if (allocated (meta%component_id)) then
       if (screen) then
          if (any (meta%active)) then
             write (msg_buffer, "(2x,A)")  "Process components:"
          else
             write (msg_buffer, "(2x,A)")  "Process components: [none]"
          end if
          call msg_message ()
       else
          write (u, "(3x,A)")  "Process components:"
       end if
       do i = 1, size (meta%component_id)
          if (.not. meta%active(i))  cycle
          if (screen) then
             write (msg_buffer, "(4x,I0,9A)")  i, ": '", &
                  char (meta%component_id (i)), "':   ", &
                  char (meta%component_description (i))
             call msg_message ()
          else
             write (u, "(5x,I0,9A)")  i, ": '", &
                  char (meta%component_id (i)), "':   ", &
                  char (meta%component_description (i))             
          end if
       end do
    end if
    if (screen) then
       write (msg_buffer, "(A)")  repeat ("-", 72)
       call msg_message ()
    else
       call write_separator (u)       
    end if
    if (screen)  return
    if (var_list) then
       write (u, "(1x,A)")  "Variable list:"
       call write_separator (u)
       call var_list_write (meta%var_list, u)
    else
       write (u, "(1x,A)")  "Variable list: [not shown]"
    end if
  end subroutine process_metadata_write

\eatline
\nwendcode{}\nwbegindocs{320}Short output: list components.
\nwenddocs{}\nwbegincode{321}\moddef{Processes: process metadata: TBP}\plusendmoddef
  procedure :: show => process_metadata_show
\nwendcode{}\nwbegincode{322}\moddef{Processes: procedures}\plusendmoddef
  subroutine process_metadata_show (meta, u, model_name)
    class(process_metadata_t), intent(in) :: meta
    integer, intent(in) :: u
    type(string_t), intent(in) :: model_name
    integer :: i
    select case (meta%type)
    case (PRC_UNKNOWN)
       write (u, "(A)") "Process: [undefined]"
       return
    case default
       write (u, "(A)", advance="no") "Process:"
    end select
    write (u, "(1x,A)", advance="no") char (meta%id)
    select case (meta%num_id)
    case (0)
    case default
       write (u, "(1x,'(',I0,')')", advance="no") meta%num_id
    end select
    select case (char (model_name))
    case ("")
    case default
       write (u, "(1x,'[',A,']')", advance="no")  char (model_name)
    end select
    write (u, *)
    if (allocated (meta%component_id)) then
       do i = 1, size (meta%component_id)
          if (meta%active(i)) then
             write (u, "(2x,I0,':',1x,A)")  i, &
                  char (meta%component_description (i))
          end if
       end do
    end if
  end subroutine process_metadata_show

\eatline
\nwendcode{}\nwbegindocs{323}Initialize.
\nwenddocs{}\nwbegincode{324}\moddef{Processes: process metadata: TBP}\plusendmoddef
  procedure :: init => process_metadata_init
\nwendcode{}\nwbegincode{325}\moddef{Processes: procedures}\plusendmoddef
  subroutine process_metadata_init (meta, id, run_id, lib)
    class(process_metadata_t), intent(out) :: meta
    type(string_t), intent(in) :: id
    type(string_t), intent(in) :: run_id
    type(process_library_t), intent(in), target :: lib
    select case (lib%get_n_in (id))
    case (1);  meta%type = PRC_DECAY
    case (2);  meta%type = PRC_SCATTERING
    case default
       call msg_bug ("Process '" // char (id) // "': impossible n_in")
    end select
    meta%id = id
    meta%run_id = run_id
    meta%lib => lib
    meta%lib_index = lib%get_entry_index (id)
    meta%num_id = lib%get_num_id (id)
    call lib%get_component_list (id, meta%component_id)
    meta%n_components = size (meta%component_id)
    call lib%get_component_description_list (id, meta%component_description)
    allocate (meta%active (meta%n_components), source = .true.)
  end subroutine process_metadata_init
  
\eatline
\nwendcode{}\nwbegindocs{326}Mark a component as inactive.
\nwenddocs{}\nwbegincode{327}\moddef{Processes: process metadata: TBP}\plusendmoddef
  procedure :: deactivate_component => process_metadata_deactivate_component
\nwendcode{}\nwbegincode{328}\moddef{Processes: procedures}\plusendmoddef
  subroutine process_metadata_deactivate_component (meta, i)
    class(process_metadata_t), intent(inout) :: meta
    integer, intent(in) :: i
    call msg_message ("Process component '" &
         // char (meta%component_id(i)) // "': matrix element vanishes")
    meta%active(i) = .false.
  end subroutine process_metadata_deactivate_component
  
\eatline
\nwendcode{}\nwbegindocs{329}\nwdocspar
\subsection{Generic Configuration Data}
This information concerns physical and technical properties of the
process.  It is fixed upon initialization, using data from the
process specification and the variable list.

The number {\tt{}n{\char95}in} is the number of incoming beam particles,
simultaneously the number of incoming partons, 1 for a decay and 2 for
a scattering process. (The number of outgoing partons may depend on
the process component.)

The number {\tt{}n{\char95}components} is the number of components that constitute
the current process.

The number {\tt{}n{\char95}terms} is the number of distinct contributions to the
scattering matrix that constitute the current process.  Each component
may generate several terms.

The number {\tt{}n{\char95}mci} is the number of independent MC
integration configurations that this process uses.  Distinct process
components that share a MCI configuration may be combined pointwise.
(Nevertheless, a given MC variable set may correspond to several
``nearby'' kinematical configurations.)  This is also the number of
distinct sampling-function results that this process can generate.
Process components that use distinct variable sets are added only once
after an integration pass has completed.

The {\tt{}model} pointer identifies the physics model and its
parameters.  This is a pointer to an external object.

The {\tt{}rng{\char95}factory} component spawns independent random-number generators for
use in integration, event generation, and event postprocessing.

Various {\tt{}parse{\char95}node{\char95}t} objects are taken from the SINDARIN input.
They encode expressions for evaluating cuts and scales.  The
workspaces for evaluating those expressions are set up in the
{\tt{}effective{\char95}state} subobjects.  Note that these are really pointers,
so the actual nodes are not stored inside the process object.

The {\tt{}md5sum} is taken and used to verify the process configuration
when re-reading data from file.
\nwenddocs{}\nwbegincode{330}\moddef{Processes: process part types}\plusendmoddef
  type :: process_config_data_t
     private
     integer :: n_in = 0
     integer :: n_components = 0
     integer :: n_terms = 0
     integer :: n_mci = 0
     type(os_data_t) :: os_data
     class(rng_factory_t), allocatable :: rng_factory
     type(string_t) :: model_name
     class(model_data_t), pointer :: model => null ()
     type(qcd_t) :: qcd
     class(expr_factory_t), allocatable :: ef_cuts
     class(expr_factory_t), allocatable :: ef_scale
     class(expr_factory_t), allocatable :: ef_fac_scale
     class(expr_factory_t), allocatable :: ef_ren_scale
     class(expr_factory_t), allocatable :: ef_weight
     character(32) :: md5sum = ""
   contains
   \LA{}Processes: process config data: TBP\RA{}
  end type process_config_data_t

\eatline
\nwendcode{}\nwbegindocs{331}Here, we may compress the expressions for cuts etc.
\nwenddocs{}\nwbegincode{332}\moddef{Processes: process config data: TBP}\endmoddef
  procedure :: write => process_config_data_write
\nwendcode{}\nwbegincode{333}\moddef{Processes: procedures}\plusendmoddef
  subroutine process_config_data_write (config, u, &
       counters, os_data, rng_factory, model, expressions)
    class(process_config_data_t), intent(in) :: config
    integer, intent(in) :: u
    logical, intent(in) :: counters
    logical, intent(in) :: os_data
    logical, intent(in) :: rng_factory
    logical, intent(in) :: model
    logical, intent(in) :: expressions
    write (u, "(1x,A)") "Configuration data:"
    if (counters) then
       write (u, "(3x,A,I0)") "Number of incoming particles = ", &
            config%n_in
       write (u, "(3x,A,I0)") "Number of process components = ", &
            config%n_components
       write (u, "(3x,A,I0)") "Number of process terms      = ", &
            config%n_terms
       write (u, "(3x,A,I0)") "Number of MCI configurations = ", &
            config%n_mci
    end if
    if (os_data) then
       call os_data_write (config%os_data, u)
    end if
    if (associated (config%model)) then
       write (u, "(3x,A,A)")  "Model = ", char (config%model_name)
       if (model) then
          call write_separator (u)
          call config%model%write (u)
          call write_separator (u)
       end if
    else
       write (u, "(3x,A,A,A)")  "Model = ", char (config%model_name), &
            " [not associated]"
    end if
    call config%qcd%write (u, show_md5sum = .false.)
    if (rng_factory) then
       if (allocated (config%rng_factory)) then
          write (u, "(2x)", advance = "no")
          call config%rng_factory%write (u)
       end if
    end if
    call write_separator (u)
    if (expressions) then
       if (allocated (config%ef_cuts)) then
          call write_separator (u)
          write (u, "(3x,A)") "Cut expression:"
          call config%ef_cuts%write (u)
       end if
       if (allocated (config%ef_scale)) then
          call write_separator (u)
          write (u, "(3x,A)") "Scale expression:"
          call config%ef_scale%write (u)
       end if
       if (allocated (config%ef_fac_scale)) then
          call write_separator (u)
          write (u, "(3x,A)") "Factorization scale expression:"
          call config%ef_fac_scale%write (u)
       end if
       if (allocated (config%ef_ren_scale)) then
          call write_separator (u)
          write (u, "(3x,A)") "Renormalization scale expression:"
          call config%ef_ren_scale%write (u)
       end if
       if (allocated (config%ef_weight)) then
          call write_separator (u)
          write (u, "(3x,A)") "Weight expression:"
          call config%ef_weight%write (u)
       end if
    else
       call write_separator (u)
       write (u, "(3x,A)") "Expressions (cut, scales, weight): [not shown]"
    end if
    if (config%md5sum /= "") then
       call write_separator (u)
       write (u, "(3x,A,A,A)")  "MD5 sum (config)  = '", config%md5sum, "'"
    end if
  end subroutine process_config_data_write
       
\eatline
\nwendcode{}\nwbegindocs{334}Initialize.  We use information from the process metadata and from
the process library, given the process ID.  We also store the
currently active OS data set.

The model that we import here should be an instance of the global model.  If we
create the instance outside, we can make it an extension of the basic
{\tt{}model{\char95}data{\char95}t}.  This will allow us to synchronize the variable list with
the model.  Using a pointer instead of an allocatable, we can guarantee the
{\tt{}target} attribute.

The RNG factory object is imported by moving the allocation.
\nwenddocs{}\nwbegincode{335}\moddef{Processes: process config data: TBP}\plusendmoddef
  procedure :: init => process_config_data_init
\nwendcode{}\nwbegincode{336}\moddef{Processes: procedures}\plusendmoddef
  subroutine process_config_data_init &
       (config, meta, os_data, qcd, rng_factory, model)
    class(process_config_data_t), intent(out) :: config
    type(process_metadata_t), intent(in) :: meta
    type(os_data_t), intent(in) :: os_data
    type(qcd_t), intent(in) :: qcd
    class(rng_factory_t), intent(inout), allocatable :: rng_factory
    class(model_data_t), intent(inout), pointer :: model
    config%n_in = meta%lib%get_n_in (meta%id)
    config%n_components = size (meta%component_id)
    config%os_data = os_data
    config%qcd = qcd
    call move_alloc (from = rng_factory, to = config%rng_factory)
    config%model_name = model%get_name ()
    config%model => model
    model => null ()
  end subroutine process_config_data_init

\eatline
\nwendcode{}\nwbegindocs{337}Since the captured model is a separate object allocated via a
pointer, we need a finalizer.
\nwenddocs{}\nwbegincode{338}\moddef{Processes: process config data: TBP}\plusendmoddef
  procedure :: final => process_config_data_final
\nwendcode{}\nwbegincode{339}\moddef{Processes: procedures}\plusendmoddef
  subroutine process_config_data_final (config)
    class(process_config_data_t), intent(inout) :: config
    if (associated (config%model)) then
       call config%model%final ()
       deallocate (config%model)
    end if
  end subroutine process_config_data_final
  
\eatline
\nwendcode{}\nwbegindocs{340}Compute the MD5 sum of the configuration data.  This encodes, in
particular, the model and the expressions for cut, scales, weight,
etc.  It should not contain the IDs and number of components, etc.,
since the MD5 sum should be useful for integrating individual
components.

This is done only once.  If the MD5 sum is nonempty, the calculation
is skipped.
\nwenddocs{}\nwbegincode{341}\moddef{Processes: process config data: TBP}\plusendmoddef
  procedure :: compute_md5sum => process_config_data_compute_md5sum
\nwendcode{}\nwbegincode{342}\moddef{Processes: procedures}\plusendmoddef
  subroutine process_config_data_compute_md5sum (config)
    class(process_config_data_t), intent(inout) :: config
    integer :: u
    if (config%md5sum == "") then
       u = free_unit ()
       open (u, status = "scratch", action = "readwrite")
       call config%write (u, counters = .false., os_data = .false., &
            rng_factory = .false., model = .true., expressions = .true.)
       rewind (u)
       config%md5sum = md5sum (u)
       close (u)
    end if
  end subroutine process_config_data_compute_md5sum
  
\eatline
\nwendcode{}\nwbegindocs{343}\nwdocspar
\subsection{Beam configuration}
The object {\tt{}data} holds all details about the initial beam
configuration.  The allocatable array {\tt{}sf} holds the structure-function
configuration blocks.  There are {\tt{}n{\char95}strfun} entries in the
structure-function chain (not counting the initial beam object).  We
maintain {\tt{}n{\char95}channel} independent parameterizations of this chain.
If this is greater than zero, we need a multi-channel sampling
algorithm, where for each point one channel is selected to generate
kinematics.

The number of parameters that are required for generating a
structure-function chain is {\tt{}n{\char95}sfpar}.

The flag {\tt{}azimuthal{\char95}dependence} tells whether the process setup is
symmetric about the beam axis in the c.m.\ system.  This implies that
there is no transversal beam polarization.  The flag {\tt{}lab{\char95}is{\char95}cm{\char95}frame} is
obvious.
\nwenddocs{}\nwbegincode{344}\moddef{Processes: process part types}\plusendmoddef
  type :: process_beam_config_t
     private
     type(beam_data_t) :: data
     integer :: n_strfun = 0
     integer :: n_channel = 1
     integer :: n_sfpar = 0
     type(sf_config_t), dimension(:), allocatable :: sf
     type(sf_channel_t), dimension(:), allocatable :: sf_channel
     logical :: azimuthal_dependence = .false.
     logical :: lab_is_cm_frame = .true.
     character(32) :: md5sum = ""
     logical :: sf_trace = .false.
     type(string_t) :: sf_trace_file
   contains
   \LA{}Processes: process beam config: TBP\RA{}
  end type process_beam_config_t

\eatline
\nwendcode{}\nwbegindocs{345}Here we write beam data only if they are actually used.
\nwenddocs{}\nwbegincode{346}\moddef{Processes: process beam config: TBP}\endmoddef
  procedure :: write => process_beam_config_write
\nwendcode{}\nwbegincode{347}\moddef{Processes: procedures}\plusendmoddef
  subroutine process_beam_config_write (object, u)
    class(process_beam_config_t), intent(in) :: object
    integer, intent(in) :: u
    integer :: i, c
    call object%data%write (u)
    if (object%data%initialized) then
       write (u, "(3x,A,L1)")  "Azimuthal dependence    = ", &
            object%azimuthal_dependence
       write (u, "(3x,A,L1)")  "Lab frame is c.m. frame = ", &
            object%lab_is_cm_frame
       if (object%md5sum /= "") then
          write (u, "(3x,A,A,A)")  "MD5 sum (beams/strf) = '", &
               object%md5sum, "'"
       end if
       if (allocated (object%sf)) then
          do i = 1, size (object%sf)
             call object%sf(i)%write (u)
          end do
          if (any_sf_channel_has_mapping (object%sf_channel)) then
             write (u, "(1x,A,L1)")  "Structure-function mappings per channel:"
             do c = 1, object%n_channel
                write (u, "(3x,I0,':')", advance="no")  c
                call object%sf_channel(c)%write (u)
             end do
          end if
       end if
    end if
  end subroutine process_beam_config_write
  
\eatline
\nwendcode{}\nwbegindocs{348}The beam data have a finalizer.  We assume that there is none for the
structure-function data.
\nwenddocs{}\nwbegincode{349}\moddef{Processes: process beam config: TBP}\plusendmoddef
  procedure :: final => process_beam_config_final
\nwendcode{}\nwbegincode{350}\moddef{Processes: procedures}\plusendmoddef
  subroutine process_beam_config_final (object)
    class(process_beam_config_t), intent(inout) :: object
    call beam_data_final (object%data)
  end subroutine process_beam_config_final

\eatline
\nwendcode{}\nwbegindocs{351}Initialize the beam setup with a given beam structure object.
\nwenddocs{}\nwbegincode{352}\moddef{Processes: process beam config: TBP}\plusendmoddef
  procedure :: init_beam_structure => process_beam_config_init_beam_structure
\nwendcode{}\nwbegincode{353}\moddef{Processes: procedures}\plusendmoddef
  subroutine process_beam_config_init_beam_structure &
       (beam_config, beam_structure, sqrts, model, decay_rest_frame)
    class(process_beam_config_t), intent(out) :: beam_config
    type(beam_structure_t), intent(in) :: beam_structure
    logical, intent(in), optional :: decay_rest_frame
    real(default), intent(in) :: sqrts
    class(model_data_t), intent(in), target :: model
    call beam_data_init_structure (beam_config%data, &
         beam_structure, sqrts, model, decay_rest_frame)
    beam_config%lab_is_cm_frame = beam_data_cm_frame (beam_config%data)
  end subroutine process_beam_config_init_beam_structure
  
\eatline
\nwendcode{}\nwbegindocs{354}Initialize the beam setup for a scattering process with specified
flavor combination, other properties taken from the beam structure
object (if any).
\nwenddocs{}\nwbegincode{355}\moddef{Processes: process beam config: TBP}\plusendmoddef
  procedure :: init_scattering => process_beam_config_init_scattering
\nwendcode{}\nwbegincode{356}\moddef{Processes: procedures}\plusendmoddef
  subroutine process_beam_config_init_scattering &
       (beam_config, flv_in, sqrts, beam_structure)
    class(process_beam_config_t), intent(out) :: beam_config
    type(flavor_t), dimension(2), intent(in) :: flv_in
    real(default), intent(in) :: sqrts
    type(beam_structure_t), intent(in), optional :: beam_structure
    if (present (beam_structure)) then
       if (beam_structure%polarized ()) then
          call beam_data_init_sqrts (beam_config%data, sqrts, flv_in, &
               beam_structure%get_smatrix (), beam_structure%get_pol_f ())
       else
          call beam_data_init_sqrts (beam_config%data, sqrts, flv_in)
       end if
    else
       call beam_data_init_sqrts (beam_config%data, sqrts, flv_in)
    end if
  end subroutine process_beam_config_init_scattering
    
\eatline
\nwendcode{}\nwbegindocs{357}Initialize the beam setup for a decay process with specified flavor,
other properties taken from the beam structure object (if present).

For a cascade decay, we set
{\tt{}rest{\char95}frame} to false, indicating a event-wise varying momentum.
The beam data itself are initialized for the particle at rest.
\nwenddocs{}\nwbegincode{358}\moddef{Processes: process beam config: TBP}\plusendmoddef
  procedure :: init_decay => process_beam_config_init_decay
\nwendcode{}\nwbegincode{359}\moddef{Processes: procedures}\plusendmoddef
  subroutine process_beam_config_init_decay &
       (beam_config, flv_in, rest_frame, beam_structure)
    class(process_beam_config_t), intent(out) :: beam_config
    type(flavor_t), dimension(1), intent(in) :: flv_in
    logical, intent(in), optional :: rest_frame
    type(beam_structure_t), intent(in), optional :: beam_structure
    if (present (beam_structure)) then
       if (beam_structure%polarized ()) then
          call beam_data_init_decay (beam_config%data, flv_in, &
               beam_structure%get_smatrix (), beam_structure%get_pol_f (), &
               rest_frame = rest_frame)
       else
          call beam_data_init_decay (beam_config%data, flv_in, &
               rest_frame = rest_frame)
       end if
    else
       call beam_data_init_decay (beam_config%data, flv_in, &
            rest_frame = rest_frame)
    end if 
    beam_config%lab_is_cm_frame = beam_data_cm_frame (beam_config%data)
  end subroutine process_beam_config_init_decay
    
\eatline
\nwendcode{}\nwbegindocs{360}Print an informative message.
\nwenddocs{}\nwbegincode{361}\moddef{Processes: process beam config: TBP}\plusendmoddef
  procedure :: startup_message => process_beam_config_startup_message
\nwendcode{}\nwbegincode{362}\moddef{Processes: procedures}\plusendmoddef
  subroutine process_beam_config_startup_message &
       (beam_config, unit, beam_structure)
    class(process_beam_config_t), intent(in) :: beam_config
    integer, intent(in), optional :: unit
    type(beam_structure_t), intent(in), optional :: beam_structure
    integer :: u
    u = free_unit ()
    open (u, status="scratch", action="readwrite")
    if (present (beam_structure)) then
       call beam_structure%write (u)
    end if
    call beam_data_write (beam_config%data, u)
    rewind (u)
    do
       read (u, "(1x,A)", end=1)  msg_buffer
       call msg_message ()
    end do
1   continue
    close (u)
  end subroutine process_beam_config_startup_message

\eatline
\nwendcode{}\nwbegindocs{363}Allocate the structure-function array.
\nwenddocs{}\nwbegincode{364}\moddef{Processes: process beam config: TBP}\plusendmoddef
  procedure :: init_sf_chain => process_beam_config_init_sf_chain
\nwendcode{}\nwbegincode{365}\moddef{Processes: procedures}\plusendmoddef
  subroutine process_beam_config_init_sf_chain &
       (beam_config, sf_config, sf_trace_file)
    class(process_beam_config_t), intent(inout) :: beam_config
    type(sf_config_t), dimension(:), intent(in) :: sf_config
    type(string_t), intent(in), optional :: sf_trace_file
    integer :: i
    beam_config%n_strfun = size (sf_config)
    allocate (beam_config%sf (beam_config%n_strfun))
    do i = 1, beam_config%n_strfun
       associate (sf => sf_config(i))
         call beam_config%sf(i)%init (sf%i, sf%data)
         if (.not. sf%data%is_generator ()) then
            beam_config%n_sfpar = beam_config%n_sfpar + sf%data%get_n_par ()
         end if
       end associate
    end do
    if (present (sf_trace_file)) then
       beam_config%sf_trace = .true.
       beam_config%sf_trace_file = sf_trace_file
    end if
  end subroutine process_beam_config_init_sf_chain

\eatline
\nwendcode{}\nwbegindocs{366}Allocate the structure-function mapping channel array, given the
requested number of channels.
\nwenddocs{}\nwbegincode{367}\moddef{Processes: process beam config: TBP}\plusendmoddef
  procedure :: allocate_sf_channels => process_beam_config_allocate_sf_channels
\nwendcode{}\nwbegincode{368}\moddef{Processes: procedures}\plusendmoddef
  subroutine process_beam_config_allocate_sf_channels (beam_config, n_channel)
    class(process_beam_config_t), intent(inout) :: beam_config
    integer, intent(in) :: n_channel
    beam_config%n_channel = n_channel
    call allocate_sf_channels (beam_config%sf_channel, &
         n_channel = n_channel, &
         n_strfun = beam_config%n_strfun)
  end subroutine process_beam_config_allocate_sf_channels
    
\eatline
\nwendcode{}\nwbegindocs{369}Set a structure-function mapping channel for an array of
structure-function entries, for a single channel.  (The default is no mapping.)
\nwenddocs{}\nwbegincode{370}\moddef{Processes: process beam config: TBP}\plusendmoddef
  procedure :: set_sf_channel => process_beam_config_set_sf_channel
\nwendcode{}\nwbegincode{371}\moddef{Processes: procedures}\plusendmoddef
  subroutine process_beam_config_set_sf_channel (beam_config, c, sf_channel)
    class(process_beam_config_t), intent(inout) :: beam_config
    integer, intent(in) :: c
    type(sf_channel_t), intent(in) :: sf_channel
    beam_config%sf_channel(c) = sf_channel
  end subroutine process_beam_config_set_sf_channel
  
\eatline
\nwendcode{}\nwbegindocs{372}Print an informative startup message.
\nwenddocs{}\nwbegincode{373}\moddef{Processes: process beam config: TBP}\plusendmoddef
  procedure :: sf_startup_message => process_beam_config_sf_startup_message
\nwendcode{}\nwbegincode{374}\moddef{Processes: procedures}\plusendmoddef
  subroutine process_beam_config_sf_startup_message &
       (beam_config, sf_string, unit)
    class(process_beam_config_t), intent(in) :: beam_config
    type(string_t), intent(in) :: sf_string
    integer, intent(in), optional :: unit
    if (beam_config%n_strfun > 0) then
       call msg_message ("Beam structure: " // char (sf_string), unit = unit)
       write (msg_buffer, "(A,3(1x,I0,1x,A))") &
            "Beam structure:", &
            beam_config%n_channel, "channels,", &
            beam_config%n_sfpar, "dimensions"
       call msg_message (unit = unit)
       if (beam_config%sf_trace) then
          call msg_message ("Beam structure: tracing &
               &values in '" // char (beam_config%sf_trace_file) // "'")
       end if
    end if
  end subroutine process_beam_config_sf_startup_message
    
\eatline
\nwendcode{}\nwbegindocs{375}Return the PDF set currently in use, if any.  This should be unique,
so we scan the structure functions until we get a nonzero number.

(This implies that if the PDF set is not unique (e.g., proton and
photon structure used together), this doesn't work correctly.)
\nwenddocs{}\nwbegincode{376}\moddef{Processes: process beam config: TBP}\plusendmoddef
  procedure :: get_pdf_set => process_beam_config_get_pdf_set
\nwendcode{}\nwbegincode{377}\moddef{Processes: procedures}\plusendmoddef
  function process_beam_config_get_pdf_set (beam_config) result (pdf_set)
    class(process_beam_config_t), intent(in) :: beam_config
    integer :: pdf_set
    integer :: i
    if (allocated (beam_config%sf)) then
       do i = 1, size (beam_config%sf)
          pdf_set = beam_config%sf(i)%get_pdf_set ()
          if (pdf_set /= 0)  return
       end do
    else
       pdf_set = 0
    end if
  end function process_beam_config_get_pdf_set
  
\eatline
\nwendcode{}\nwbegindocs{378}Compute the MD5 sum for the complete beam setup.  We rely on the
default output of {\tt{}write} to contain all relevant data.

This is done only once, when the MD5 sum is still empty.
\nwenddocs{}\nwbegincode{379}\moddef{Processes: process beam config: TBP}\plusendmoddef
  procedure :: compute_md5sum => process_beam_config_compute_md5sum
\nwendcode{}\nwbegincode{380}\moddef{Processes: procedures}\plusendmoddef
  subroutine process_beam_config_compute_md5sum (beam_config)
    class(process_beam_config_t), intent(inout) :: beam_config
    integer :: u
    if (beam_config%md5sum == "") then
       u = free_unit ()
       open (u, status = "scratch", action = "readwrite")
       call beam_config%write (u)
       rewind (u)
       beam_config%md5sum = md5sum (u)
       close (u)
    end if
  end subroutine process_beam_config_compute_md5sum

\eatline
\nwendcode{}\nwbegindocs{381}\nwdocspar
\subsubsection{Process call statistics}
This object can record process calls, categorized by evaluation
status.  It is a part of the {\tt{}mci{\char95}entry} component below.
\nwenddocs{}\nwbegincode{382}\moddef{Processes: types}\endmoddef
  type :: process_counter_t
     integer :: total = 0
     integer :: failed_kinematics = 0
     integer :: failed_cuts = 0
     integer :: passed = 0
     integer :: evaluated = 0
     integer :: complete = 0
   contains
   \LA{}Processes: process counter: TBP\RA{}
  end type process_counter_t
  
\eatline
\nwendcode{}\nwbegindocs{383}\nwdocspar
\nwenddocs{}\nwbegincode{384}\moddef{Processes: types}\plusendmoddef
  type :: sqme_collector_t
    real(default) :: sqme_born = 0
    real(default) :: sqme_real = 0
    real(default) :: sqme_virt = 0
  contains
  \LA{}Processes: sqme collector: TBP\RA{}
  end type sqme_collector_t

\eatline
\nwendcode{}\nwbegindocs{385}Here are the corresponding numeric codes:
\nwenddocs{}\nwbegincode{386}\moddef{Processes: parameters}\endmoddef
  integer, parameter :: STAT_UNDEFINED = 0
  integer, parameter :: STAT_INITIAL = 1
  integer, parameter :: STAT_ACTIVATED = 2
  integer, parameter :: STAT_BEAM_MOMENTA = 3
  integer, parameter :: STAT_FAILED_KINEMATICS = 4
  integer, parameter :: STAT_SEED_KINEMATICS = 5
  integer, parameter :: STAT_HARD_KINEMATICS = 6
  integer, parameter :: STAT_EFF_KINEMATICS = 7
  integer, parameter :: STAT_FAILED_CUTS = 8
  integer, parameter :: STAT_PASSED_CUTS = 9
  integer, parameter :: STAT_EVALUATED_TRACE = 10
  integer, parameter :: STAT_EVENT_COMPLETE = 11

  integer, parameter, public :: COMP_DEFAULT = 0
  integer, parameter, public :: COMP_MASTER  = 1
  integer, parameter, public :: COMP_VIRT = 2
  integer, parameter, public :: COMP_REAL = 3
  integer, parameter, public :: COMP_SUB = 4
  
\eatline
\eatline
\eatline
\eatline
\nwendcode{}\nwbegindocs{387}Output.
\nwenddocs{}\nwbegincode{388}\moddef{Processes: process counter: TBP}\endmoddef
  procedure :: write => process_counter_write
\nwendcode{}\nwbegincode{389}\moddef{Processes: procedures}\plusendmoddef
  subroutine process_counter_write (object, unit)
    class(process_counter_t), intent(in) :: object
    integer, intent(in), optional :: unit
    integer :: u
    u = given_output_unit (unit)
    if (object%total > 0) then
       write (u, "(1x,A)")  "Call statistics (current run):"
       write (u, "(3x,A,I0)")  "total       = ", object%total
       write (u, "(3x,A,I0)")  "failed kin. = ", object%failed_kinematics
       write (u, "(3x,A,I0)")  "failed cuts = ", object%failed_cuts
       write (u, "(3x,A,I0)")  "passed cuts = ", object%passed
       write (u, "(3x,A,I0)")  "evaluated   = ", object%evaluated
    else
       write (u, "(1x,A)")  "Call statistics (current run): [no calls]"
    end if
  end subroutine process_counter_write
    
\eatline
\nwendcode{}\nwbegindocs{390}Reset.  Just enforce default initialization.
\nwenddocs{}\nwbegincode{391}\moddef{Processes: process counter: TBP}\plusendmoddef
  procedure :: reset => process_counter_reset
\nwendcode{}\nwbegincode{392}\moddef{Processes: procedures}\plusendmoddef
  subroutine process_counter_reset (counter)
    class(process_counter_t), intent(out) :: counter
  end subroutine process_counter_reset

\eatline
\nwendcode{}\nwbegindocs{393}We record an event according to the lowest status code greater or
equal to the actual status.  This is actually done by the process
instance; the process object just copies the instance counter.
\nwenddocs{}\nwbegincode{394}\moddef{Processes: process counter: TBP}\plusendmoddef
  procedure :: record => process_counter_record
\nwendcode{}\nwbegincode{395}\moddef{Processes: procedures}\plusendmoddef
  subroutine process_counter_record (counter, status)
    class(process_counter_t), intent(inout) :: counter
    integer, intent(in) :: status
    if (status <= STAT_FAILED_KINEMATICS) then
       counter%failed_kinematics = counter%failed_kinematics + 1
    else if (status <= STAT_FAILED_CUTS) then
       counter%failed_cuts = counter%failed_cuts + 1
    else if (status <= STAT_PASSED_CUTS) then
       counter%passed = counter%passed + 1
    else
       counter%evaluated = counter%evaluated + 1
    end if
    counter%total = counter%total + 1
  end subroutine process_counter_record
       
\eatline
\nwendcode{}\nwbegindocs{396}\nwdocspar
\nwenddocs{}\nwbegincode{397}\moddef{Processes: sqme collector: TBP}\endmoddef
  procedure :: get_sqme_sum => sqme_collector_get_sqme_sum
\nwendcode{}\nwbegincode{398}\moddef{Processes: procedures}\plusendmoddef
  function sqme_collector_get_sqme_sum (collector) result (sqme)
    class(sqme_collector_t), intent(in) :: collector
    real(default) :: sqme
    sqme = collector%sqme_born + &
           collector%sqme_real + &
           collector%sqme_virt
  end function sqme_collector_get_sqme_sum

\eatline
\nwendcode{}\nwbegindocs{399}\nwdocspar
\subsection{Multi-channel integration}
The {\tt{}process{\char95}mci{\char95}entry{\char95}t} block contains, for each process component that is
integrated independently, the configuration data for its MC input parameters.
Each input parameter set is handled by a {\tt{}mci{\char95}t} integrator.

The MC input parameter set is broken down into the parameters required by the
structure-function chain and the parameters required by the phase space of the
elementary process.

The MD5 sum collects all information about the associated processes
that may affect the integration.  It does not contain the MCI object
itself or integration results.

MC integration is organized in passes.  Each pass may consist of
several iterations, and for each iteration there is a number of
calls.  We store explicitly the values that apply to the current
pass.  Previous values are archived in the {\tt{}results} object.

The {\tt{}counter} receives the counter statistics from the associated
process instance, for diagnostics.

The {\tt{}results} object records results, broken down in passes and iterations.
\nwenddocs{}\nwbegincode{400}\moddef{Processes: process part types}\plusendmoddef
  type :: process_mci_entry_t
     integer :: i_mci = 0
     integer, dimension(:), allocatable :: i_component
     integer :: process_type = PRC_UNKNOWN
     integer :: n_par = 0
     integer :: n_par_sf = 0
     integer :: n_par_phs = 0
     character(32) :: md5sum = ""
     integer :: pass = 0
     integer :: n_it = 0
     integer :: n_calls = 0
     logical :: activate_timer = .false.
     real(default) :: error_threshold = 0
     class(mci_t), allocatable :: mci
     type(process_counter_t) :: counter
     type(integration_results_t) :: results
     logical :: negative_weights
     logical :: combined_integration = .false.
   contains
   \LA{}Processes: process mci entry: TBP\RA{}
  end type process_mci_entry_t

\eatline
\nwendcode{}\nwbegindocs{401}Finalizer for the {\tt{}mci} component.
\nwenddocs{}\nwbegincode{402}\moddef{Processes: process mci entry: TBP}\endmoddef
  procedure :: final => process_mci_entry_final
\nwendcode{}\nwbegincode{403}\moddef{Processes: procedures}\plusendmoddef
  subroutine process_mci_entry_final (object)
    class(process_mci_entry_t), intent(inout) :: object
    if (allocated (object%mci))  call object%mci%final ()
  end subroutine process_mci_entry_final
  
\eatline
\nwendcode{}\nwbegindocs{404}Output.  Write pass/iteration information only if set (the pass
index is nonzero).  Write the MCI block only if it exists (for some
self-tests it does not).  Write results only if there are any.
\nwenddocs{}\nwbegincode{405}\moddef{Processes: process mci entry: TBP}\plusendmoddef
  procedure :: write => process_mci_entry_write
\nwendcode{}\nwbegincode{406}\moddef{Processes: procedures}\plusendmoddef
  subroutine process_mci_entry_write (object, unit, pacify)
    class(process_mci_entry_t), intent(in) :: object
    integer, intent(in), optional :: unit
    logical, intent(in), optional :: pacify
    integer :: u
    u = given_output_unit (unit)
    write (u, "(3x,A,I0)")  "Associated components = ", object%i_component
    write (u, "(3x,A,I0)")  "MC input parameters   = ", object%n_par
    write (u, "(3x,A,I0)")  "MC parameters (SF)    = ", object%n_par_sf
    write (u, "(3x,A,I0)")  "MC parameters (PHS)   = ", object%n_par_phs
    if (object%pass > 0) then
       write (u, "(3x,A,I0)")  "Current pass          = ", object%pass
       write (u, "(3x,A,I0)")  "Number of iterations  = ", object%n_it
       write (u, "(3x,A,I0)")  "Number of calls       = ", object%n_calls
    end if
    if (object%md5sum /= "") then
       write (u, "(3x,A,A,A)") "MD5 sum (components)  = '", object%md5sum, "'"
    end if
    if (allocated (object%mci)) then
       call object%mci%write (u)
    end if
    call object%counter%write (u)
    if (object%results%exist ()) then
       call object%results%write (u, suppress = pacify)
       call object%results%write_chain_weights (u)
    end if
  end subroutine process_mci_entry_write
       
\eatline
\nwendcode{}\nwbegindocs{407}Write the weights that are assigned to the chains of channels (call groves in
the {\tt{}wood} phase-space implementation.  This depends on the implementation
and is delegated to the {\tt{}mci} component.
\nwenddocs{}\nwbegincode{408}\moddef{Processes: process mci entry: TBP}\plusendmoddef
  procedure :: write_chain_weights => process_mci_entry_write_chain_weights
\nwendcode{}\nwbegincode{409}\moddef{Processes: procedures}\plusendmoddef
  subroutine process_mci_entry_write_chain_weights (mci_entry, unit)
    class(process_mci_entry_t), intent(in) :: mci_entry
    integer, intent(in), optional :: unit
    if (allocated (mci_entry%mci)) then
       call mci_entry%mci%write_chain_weights (unit)
    end if
  end subroutine process_mci_entry_write_chain_weights
       
\eatline
\nwendcode{}\nwbegindocs{410}Initialize.  From the existing configuration, we obtain the number of
channels and the number of parameters, separately for the structure-function
chain and for the associated process component.  We assume that the
phase-space object has already been configured.

The allocation of the MCI integrator with the appropriate concrete type is the
duty of the process core.

We assume that there is only one component associated with a MCI entry.  This
restriction should be relaxed.
\nwenddocs{}\nwbegincode{411}\moddef{Processes: process mci entry: TBP}\plusendmoddef
  procedure :: init => process_mci_entry_init
\nwendcode{}\nwbegincode{412}\moddef{Processes: procedures}\plusendmoddef
  subroutine process_mci_entry_init (mci_entry, &
       process_type, i_mci, i_component, component,&
       beam_config, rng_factory)
    class(process_mci_entry_t), intent(inout) :: mci_entry
    integer, intent(in) :: process_type
    integer, intent(in) :: i_mci
    integer, intent(in) :: i_component
    type(process_component_t), intent(in), target :: component
    type(process_beam_config_t), intent(in) :: beam_config
    class(rng_factory_t), intent(inout) :: rng_factory
    class(rng_t), allocatable :: rng
    integer, dimension(4) :: i_list
    associate (phs_config => component%phs_config)
      mci_entry%i_mci = i_mci
      call mci_entry%create_component_list (i_component, component)
      mci_entry%n_par_sf = beam_config%n_sfpar
      mci_entry%n_par_phs = phs_config%get_n_par () 
      mci_entry%n_par = mci_entry%n_par_sf + mci_entry%n_par_phs
      mci_entry%process_type = process_type
      if (allocated (component%mci_template)) then
         allocate (mci_entry%mci, source=component%mci_template)
         call mci_entry%mci%record_index (mci_entry%i_mci)
         call mci_entry%mci%set_dimensions &
              (mci_entry%n_par, phs_config%get_n_channel ())
         call mci_entry%mci%declare_flat_dimensions &
              (phs_config%get_flat_dimensions ())
         if (phs_config%provides_equivalences) then
            call mci_entry%mci%declare_equivalences &
                 (phs_config%channel, mci_entry%n_par_sf)
         end if
         if (phs_config%provides_chains) then
            call mci_entry%mci%declare_chains (phs_config%chain)
         end if
         call rng_factory%make (rng)
         call mci_entry%mci%import_rng (rng)
      end if
    end associate
    call mci_entry%results%init (process_type)
  end subroutine process_mci_entry_init
  
\eatline
\nwendcode{}\nwbegindocs{413}\nwdocspar
\nwenddocs{}\nwbegincode{414}\moddef{Processes: process mci entry: TBP}\plusendmoddef
  procedure :: create_component_list => &
                     process_mci_entry_create_component_list
\nwendcode{}\nwbegincode{415}\moddef{Processes: procedures}\plusendmoddef
  subroutine process_mci_entry_create_component_list (mci_entry, &
                                     i_component, component)
    class (process_mci_entry_t), intent(inout) :: mci_entry
    integer, intent(in) :: i_component
    type(process_component_t), intent(in), target :: component
    integer, dimension(4) :: i_list
    if (mci_entry%combined_integration) then
      allocate (mci_entry%i_component (3))
      i_list = component%config%get_association_list ()
      mci_entry%i_component = i_list(1:3)
    else
      allocate (mci_entry%i_component (1))
      mci_entry%i_component(1) = i_component
    end if
  end subroutine process_mci_entry_create_component_list

\eatline
\nwendcode{}\nwbegindocs{416}\nwdocspar
\nwenddocs{}\nwbegincode{417}\moddef{Processes: process mci entry: TBP}\plusendmoddef
  procedure :: set_combined_integration => &
                     process_mci_entry_set_combined_integration
\nwendcode{}\nwbegincode{418}\moddef{Processes: procedures}\plusendmoddef
  subroutine process_mci_entry_set_combined_integration (mci_entry, value)
    class(process_mci_entry_t), intent(inout) :: mci_entry
    logical, intent(in), optional :: value
    if (present (value)) &
       mci_entry%combined_integration = value
  end subroutine process_mci_entry_set_combined_integration

\eatline
\nwendcode{}\nwbegindocs{419}Set some additional parameters.
\nwenddocs{}\nwbegincode{420}\moddef{Processes: process mci entry: TBP}\plusendmoddef
  procedure :: set_parameters => process_mci_entry_set_parameters
\nwendcode{}\nwbegincode{421}\moddef{Processes: procedures}\plusendmoddef
  subroutine process_mci_entry_set_parameters (mci_entry, var_list)
    class(process_mci_entry_t), intent(inout) :: mci_entry
    type(var_list_t), intent(in) :: var_list
    real(default) :: error_threshold
    error_threshold = &
         var_list%get_rval (var_str ("error_threshold"))
    mci_entry%activate_timer = &
         var_list%get_lval (var_str ("?integration_timer"))
    call mci_entry%results%set_error_threshold (error_threshold)
  end subroutine process_mci_entry_set_parameters
  
\eatline
\nwendcode{}\nwbegindocs{422}Compute a MD5 sum that summarizes all information that could
influence integration results, for the associated process components.
We take the process-configuration MD5 sum which represents parameters,
cuts, etc., the MD5 sums for the process component definitions and
their phase space objects (which should be configured), and the beam
configuration MD5 sum.  (The QCD setup is included in the process
configuration data MD5 sum.)

Done only once, when the MD5 sum is still empty.
\nwenddocs{}\nwbegincode{423}\moddef{Processes: process mci entry: TBP}\plusendmoddef
  procedure :: compute_md5sum => process_mci_entry_compute_md5sum
\nwendcode{}\nwbegincode{424}\moddef{Processes: procedures}\plusendmoddef
  subroutine process_mci_entry_compute_md5sum (mci_entry, &
       config, component, beam_config)
    class(process_mci_entry_t), intent(inout) :: mci_entry
    type(process_config_data_t), intent(in) :: config
    type(process_component_t), dimension(:), intent(in) :: component
    type(process_beam_config_t), intent(in) :: beam_config
    type(string_t) :: buffer
    integer :: i
    if (mci_entry%md5sum == "") then
       buffer = config%md5sum // beam_config%md5sum
       do i = 1, size (component)
          if (component(i)%active) then
             buffer = buffer // component(i)%config%get_md5sum () &
                  // component(i)%md5sum_phs
          end if
       end do
       mci_entry%md5sum = md5sum (char (buffer))
    end if
    if (allocated (mci_entry%mci)) then
       call mci_entry%mci%set_md5sum (mci_entry%md5sum)
    end if
  end subroutine process_mci_entry_compute_md5sum
  
\eatline
\nwendcode{}\nwbegindocs{425}Test the MCI sampler by calling it a given number of time, discarding the
results.  The instance should be initialized.

The {\tt{}mci{\char95}entry} is {\tt{}intent(inout)} because the integrator contains
the random-number state.
\nwenddocs{}\nwbegincode{426}\moddef{Processes: process mci entry: TBP}\plusendmoddef
  procedure :: sampler_test => process_mci_entry_sampler_test
\nwendcode{}\nwbegincode{427}\moddef{Processes: procedures}\plusendmoddef
  subroutine process_mci_entry_sampler_test (mci_entry, instance, n_calls)
    class(process_mci_entry_t), intent(inout) :: mci_entry
    type(process_instance_t), intent(inout), target :: instance
    integer, intent(in) :: n_calls
    call instance%choose_mci (mci_entry%i_mci)
    call instance%reset_counter ()
    call mci_entry%mci%sampler_test (instance, n_calls)
    mci_entry%counter = instance%get_counter ()
  end subroutine process_mci_entry_sampler_test

\eatline
\nwendcode{}\nwbegindocs{428}Integrate.  The instance should be initialized.

The {\tt{}integrate} method counts as an integration pass; the pass count is
increased by one.  We transfer the pass parameters (number of iterations and
number of calls) to the actual integration routine.

The {\tt{}mci{\char95}entry} is {\tt{}intent(inout)} because the integrator contains
the random-number state.

Note: The results are written to screen and to logfile.  This behavior
is hardcoded.
\nwenddocs{}\nwbegincode{429}\moddef{Processes: process mci entry: TBP}\plusendmoddef
  procedure :: integrate => process_mci_entry_integrate
  procedure :: final_integration => process_mci_entry_final_integration
\nwendcode{}\nwbegincode{430}\moddef{Processes: procedures}\plusendmoddef
  subroutine process_mci_entry_integrate (mci_entry, instance, n_it, n_calls, &
       adapt_grids, adapt_weights, final, pacify, &
       i_component)
    class(process_mci_entry_t), intent(inout) :: mci_entry
    type(process_instance_t), intent(inout) :: instance
    integer, intent(in) :: n_it
    integer, intent(in) :: n_calls
    logical, intent(in), optional :: adapt_grids
    logical, intent(in), optional :: adapt_weights
    logical, intent(in), optional :: final, pacify
    integer, intent(in), optional :: i_component
    integer :: u_log
    type(string_t) :: nlo_type
    if (present (i_component)) then
      if (associated (instance%component(i_component)%config)) then
        nlo_type = instance%component(i_component)%config%config%get_nlo_type ()
      end if
    else
      nlo_type = 'Born'
    end if
    u_log = logfile_unit ()
    call instance%choose_mci (mci_entry%i_mci)
    call instance%reset_counter ()
    mci_entry%pass = mci_entry%pass + 1
    mci_entry%n_it = n_it
    mci_entry%n_calls = n_calls
    if (mci_entry%pass == 1)  &
         call mci_entry%mci%startup_message (n_calls = n_calls)
    call mci_entry%mci%set_timer (active = mci_entry%activate_timer)
    call mci_entry%results%display_init &
         (mci_entry%process_type, screen = .true., unit = u_log)
    call mci_entry%results%new_pass ()
    associate (mci_instance => instance%mci_work(mci_entry%i_mci)%mci)
      select case (char (nlo_type))
      case ('Virtual', 'Real')
        select type (mci_instance)
        type is (mci_vamp_instance_t)
          mci_instance%negative_weights = .true.
        end select
      end select
      call mci_entry%mci%add_pass (adapt_grids, adapt_weights, final)
      call mci_entry%mci%start_timer ()
      call mci_entry%mci%integrate (mci_instance, instance, n_it, &
           n_calls, mci_entry%results, pacify = pacify)
      call mci_entry%mci%stop_timer ()
      if (signal_is_pending ())  return
    end associate
    mci_entry%counter = instance%get_counter ()
    call mci_entry%results%display_pass (pacify)
  end subroutine process_mci_entry_integrate

  subroutine process_mci_entry_final_integration (mci_entry)
    class(process_mci_entry_t), intent(inout) :: mci_entry
    call mci_entry%results%display_final ()
    call mci_entry%time_message ()
  end subroutine process_mci_entry_final_integration

\eatline
\eatline
\nwendcode{}\nwbegindocs{431}If appropriate, issue an informative message about the expected time
for an event sample.
\nwenddocs{}\nwbegincode{432}\moddef{Processes: process mci entry: TBP}\plusendmoddef
  procedure :: get_time => process_mci_entry_get_time
  procedure :: time_message => process_mci_entry_time_message
\nwendcode{}\nwbegincode{433}\moddef{Processes: procedures}\plusendmoddef
  subroutine process_mci_entry_get_time (mci_entry, time, sample)
    class(process_mci_entry_t), intent(in) :: mci_entry
    type(time_t), intent(out) :: time
    integer, intent(in) :: sample
    real(default) :: time_last_pass, efficiency, calls
    time_last_pass = mci_entry%mci%get_time ()
    calls = mci_entry%results%get_n_calls ()
    efficiency = mci_entry%mci%get_efficiency ()
    if (time_last_pass > 0 .and. calls > 0 .and. efficiency > 0) then
       time = nint (time_last_pass / calls / efficiency * sample)
    end if
  end subroutine process_mci_entry_get_time   

  subroutine process_mci_entry_time_message (mci_entry)
    class(process_mci_entry_t), intent(in) :: mci_entry
    type(time_t) :: time
    integer :: sample
    sample = 10000
    call mci_entry%get_time (time, sample)
    if (time%is_known ()) then
       call msg_message ("Time estimate for generating 10000 events: " &
            // char (time%to_string_dhms ()))
    end if
  end subroutine process_mci_entry_time_message
  
\eatline
\nwendcode{}\nwbegindocs{434}Prepare event generation.  (For the test integrator, this does nothing.  It
is relevant for the VAMP integrator.)
\nwenddocs{}\nwbegincode{435}\moddef{Processes: process mci entry: TBP}\plusendmoddef
  procedure :: prepare_simulation => process_mci_entry_prepare_simulation
\nwendcode{}\nwbegincode{436}\moddef{Processes: procedures}\plusendmoddef
  subroutine process_mci_entry_prepare_simulation (mci_entry)
    class(process_mci_entry_t), intent(inout) :: mci_entry
    call mci_entry%mci%prepare_simulation ()
  end subroutine process_mci_entry_prepare_simulation

\eatline
\nwendcode{}\nwbegindocs{437}Generate an event.  The instance should be initialized,
otherwise event generation is directed by the {\tt{}mci} integrator
subobject.  The integrator instance is contained in a {\tt{}mci{\char95}work}
subobject of the process instance, which simultaneously serves as the
sampler object.  (We avoid the anti-aliasing rules if we assume that
the sampling itself does not involve the integrator instance contained in the
process instance.)

Regarding weighted events, we only take events which are valid, which
means that they have valid kinematics and have passed cuts.
Therefore, we have a rejection loop.  For unweighted events, the
unweighting routine should already take care of this.
\nwenddocs{}\nwbegincode{438}\moddef{Processes: process mci entry: TBP}\plusendmoddef
  procedure :: generate_weighted_event => &
       process_mci_entry_generate_weighted_event
  procedure :: generate_unweighted_event => &
       process_mci_entry_generate_unweighted_event
  procedure :: recover_event => process_mci_entry_recover_event
\nwendcode{}\nwbegincode{439}\moddef{Processes: procedures}\plusendmoddef
  subroutine process_mci_entry_generate_weighted_event (mci_entry, instance)
    class(process_mci_entry_t), intent(inout) :: mci_entry
    type(process_instance_t), intent(inout) :: instance
    call instance%choose_mci (mci_entry%i_mci)
    associate (mci_instance => instance%mci_work(mci_entry%i_mci)%mci)
      REJECTION: do
         call mci_entry%mci%generate_weighted_event (mci_instance, instance)
         if (signal_is_pending ())  return
         if (instance%is_valid ())  exit REJECTION
      end do REJECTION
    end associate
  end subroutine process_mci_entry_generate_weighted_event
  
  subroutine process_mci_entry_generate_unweighted_event (mci_entry, instance)
    class(process_mci_entry_t), intent(inout) :: mci_entry
    type(process_instance_t), intent(inout) :: instance
    call instance%choose_mci (mci_entry%i_mci)
    associate (mci_instance => instance%mci_work(mci_entry%i_mci)%mci)
      call mci_entry%mci%generate_unweighted_event (mci_instance, instance)
    end associate
  end subroutine process_mci_entry_generate_unweighted_event
  
  subroutine process_mci_entry_recover_event (mci_entry, instance, i_term)
    class(process_mci_entry_t), intent(inout) :: mci_entry
    type(process_instance_t), intent(inout) :: instance
    integer, intent(in) :: i_term
    integer :: channel
    mci_entry%i_mci = instance%i_mci
    channel = instance%get_channel ()
    associate (mci_instance => instance%mci_work(mci_entry%i_mci)%mci)
      call mci_instance%fetch (instance, channel)
    end associate
  end subroutine process_mci_entry_recover_event
  
\eatline
\eatline
\eatline
\nwendcode{}\nwbegindocs{440}Extract results.
\nwenddocs{}\nwbegincode{441}\moddef{Processes: process mci entry: TBP}\plusendmoddef
  procedure :: has_integral => process_mci_entry_has_integral
  procedure :: get_integral => process_mci_entry_get_integral
  procedure :: get_error => process_mci_entry_get_error  
  procedure :: get_accuracy => process_mci_entry_get_accuracy
  procedure :: get_chi2 => process_mci_entry_get_chi2
  procedure :: get_efficiency => process_mci_entry_get_efficiency
\nwendcode{}\nwbegincode{442}\moddef{Processes: procedures}\plusendmoddef
  function process_mci_entry_has_integral (mci_entry) result (flag)
    class(process_mci_entry_t), intent(in) :: mci_entry
    logical :: flag
    flag = mci_entry%results%exist ()
  end function process_mci_entry_has_integral
    
  function process_mci_entry_get_integral (mci_entry) result (integral)
    class(process_mci_entry_t), intent(in) :: mci_entry
    real(default) :: integral
    integral = mci_entry%results%get_integral ()
  end function process_mci_entry_get_integral

  function process_mci_entry_get_error (mci_entry) result (error)
    class(process_mci_entry_t), intent(in) :: mci_entry
    real(default) :: error
    error = mci_entry%results%get_error ()
  end function process_mci_entry_get_error
  
  function process_mci_entry_get_accuracy (mci_entry) result (accuracy)
    class(process_mci_entry_t), intent(in) :: mci_entry
    real(default) :: accuracy
    accuracy = mci_entry%results%get_accuracy ()
  end function process_mci_entry_get_accuracy
  
  function process_mci_entry_get_chi2 (mci_entry) result (chi2)
    class(process_mci_entry_t), intent(in) :: mci_entry
    real(default) :: chi2
    chi2 = mci_entry%results%get_chi2 ()
  end function process_mci_entry_get_chi2

  function process_mci_entry_get_efficiency (mci_entry) result (efficiency)
    class(process_mci_entry_t), intent(in) :: mci_entry
    real(default) :: efficiency
    efficiency = mci_entry%results%get_efficiency ()
  end function process_mci_entry_get_efficiency  
  
\eatline
\eatline
\eatline
\nwendcode{}\nwbegindocs{443}Return the MCI checksum.  This may be the one used for
configuration, but may also incorporate results, if they change the
state of the integrator (adaptation).
\nwenddocs{}\nwbegincode{444}\moddef{Processes: process mci entry: TBP}\plusendmoddef
  procedure :: get_md5sum => process_mci_entry_get_md5sum
\nwendcode{}\nwbegincode{445}\moddef{Processes: procedures}\plusendmoddef
  function process_mci_entry_get_md5sum (entry) result (md5sum)
    class(process_mci_entry_t), intent(in) :: entry
    character(32) :: md5sum
    md5sum = entry%mci%get_md5sum ()
  end function process_mci_entry_get_md5sum
  
\eatline
\nwendcode{}\nwbegindocs{446}\nwdocspar
\subsection{Process Components}
A process component is an individual contribution to a process
(scattering or decay) which needs not be physical.  The sum over all
components should be physical.

The {\tt{}index} indentifies this component within its parent process.

The actual process component is stored in the {\tt{}core} subobject.  We
use a polymorphic subobject instead of an extension of
{\tt{}process{\char95}component{\char95}t}, because the individual entries in the array
of process components can have different types.  In short,
{\tt{}process{\char95}component{\char95}t} is a wrapper for the actual process variants.

If the {\tt{}active} flag is false, we should skip this component.  This happens
if the associated process has vanishing matrix element.

The index array {\tt{}i{\char95}term} points to the individual terms generated by
this component.  The indices refer to the parent process.

The index {\tt{}i{\char95}mci} is the index of the MC integrator and parameter set which
are associated to this process component.
\nwenddocs{}\nwbegincode{447}\moddef{Processes: process part types}\plusendmoddef
  type :: process_component_t
     private
     type(process_component_def_t), pointer :: config => null ()
     integer :: index = 0
     class(prc_core_t), allocatable :: core
     class(prc_core_t), allocatable :: core_sub_born
     logical :: active = .false.
     class(mci_t), allocatable :: mci_template
     integer, dimension(:), allocatable :: i_term
     integer :: i_mci = 0
     class(phs_config_t), allocatable :: phs_config
     character(32) :: md5sum_phs = ""
     type(fks_template_t) :: fks_template
     integer :: component_type = COMP_DEFAULT
   contains
   \LA{}Processes: process component: TBP\RA{}
  end type process_component_t

\eatline
\nwendcode{}\nwbegindocs{448}Finalizer.  The MCI template may (potentially) need a finalizer.  The process
configuration finalizer may include closing an open scratch file.
\nwenddocs{}\nwbegincode{449}\moddef{Processes: process component: TBP}\endmoddef
  procedure :: final => process_component_final
\nwendcode{}\nwbegincode{450}\moddef{Processes: procedures}\plusendmoddef
  subroutine process_component_final (object)
    class(process_component_t), intent(inout) :: object
    if (allocated (object%mci_template)) then
       call object%mci_template%final ()
    end if
    if (allocated (object%phs_config)) then
       call object%phs_config%final ()
    end if
  end subroutine process_component_final
  
\eatline
\nwendcode{}\nwbegindocs{451}The meaning of {\tt{}verbose} depends on the process variant.
\nwenddocs{}\nwbegincode{452}\moddef{Processes: process component: TBP}\plusendmoddef
  procedure :: write => process_component_write
\nwendcode{}\nwbegincode{453}\moddef{Processes: procedures}\plusendmoddef
  subroutine process_component_write (object, unit)
    class(process_component_t), intent(in) :: object
    integer, intent(in), optional :: unit
    integer :: u
    u = given_output_unit (unit)
    if (allocated (object%core)) then
       write (u, "(1x,A,I0)")  "Component #", object%index
       if (associated (object%config)) then
          call object%config%write (u)
          if (object%md5sum_phs /= "") then
             write (u, "(3x,A,A,A)")  "MD5 sum (phs)       = '", &
                  object%md5sum_phs, "'"
          end if
       end if
       write (u, "(1x,A)") "Process core:"
       call object%core%write (u)
    else
       write (u, "(1x,A)") "Process component: [not allocated]"
    end if
    if (.not. object%active) then
       write (u, "(1x,A)") "[Inactive]"
       return
    end if
    write (u, "(1x,A)") "Referenced data:"
    if (allocated (object%i_term)) then
       write (u, "(3x,A,999(1x,I0))") "Terms                    =", &
            object%i_term
    else
       write (u, "(3x,A)") "Terms                    = [undefined]"
    end if
    if (object%i_mci /= 0) then
       write (u, "(3x,A,I0)") "MC dataset               = ", object%i_mci
    else
       write (u, "(3x,A)") "MC dataset               = [undefined]"
    end if
    if (allocated (object%phs_config)) then
       call object%phs_config%write (u)
    end if
  end subroutine process_component_write

\eatline
\nwendcode{}\nwbegindocs{454}Initialize the component.
\nwenddocs{}\nwbegincode{455}\moddef{Processes: process component: TBP}\plusendmoddef
  procedure :: init => process_component_init
\nwendcode{}\nwbegincode{456}\moddef{Processes: procedures}\plusendmoddef
  subroutine process_component_init (component, &
       i_component, meta, config, &
       core_template, mci_template, phs_config_template, &
       gosam_template, fks_template)
    class(process_component_t), intent(out) :: component
    integer, intent(in) :: i_component
    type(process_metadata_t), intent(in) :: meta
    type(process_config_data_t), intent(in) :: config
    class(prc_core_t), intent(in), allocatable :: core_template
    class(mci_t), intent(in), allocatable :: mci_template
    class(phs_config_t), intent(in), allocatable :: phs_config_template
    type(gosam_template_t), intent(inout), optional :: gosam_template
    type(fks_template_t), intent(in), optional :: fks_template
    integer, dimension(3) :: amp_type

    component%index = i_component
    component%config => meta%lib%get_component_def_ptr (meta%id, i_component)
    allocate (component%core, source=core_template)
    call component%core%init (component%config%get_core_def_ptr (), &
         meta%lib, meta%id, i_component)
    if (present (fks_template)) then
      component%fks_template = fks_template
    end if
    select type (core => component%core)
    type is (prc_gosam_t)
      if (present (gosam_template)) then
        call core%init_gosam (gosam_template)
      else
        call msg_bug ("process_component_init: gosam core cannot be initialized - &
                       & missing template")
      end if
      call core%init_driver (config%os_data)
      call core%set_initialized ()
    end select
    component%active = component%core%has_matrix_element ()
    if (component%active) then
       if (allocated (mci_template)) &
            allocate (component%mci_template, source=mci_template)
       allocate (component%phs_config, source=phs_config_template)
       call component%phs_config%init (component%core%data, config%model)
    end if
  end subroutine process_component_init

\eatline
\nwendcode{}\nwbegindocs{457}Finalize the phase-space configuration. 
\nwenddocs{}\nwbegincode{458}\moddef{Processes: process component: TBP}\plusendmoddef
  procedure :: configure_phs => process_component_configure_phs
\nwendcode{}\nwbegincode{459}\moddef{Processes: procedures}\plusendmoddef
  subroutine process_component_configure_phs &
       (component, sqrts, beam_config, rebuild, &
        ignore_mismatch, phs_config, verbose)
    class(process_component_t), intent(inout) :: component
    real(default), intent(in) :: sqrts
    type(process_beam_config_t), intent(in) :: beam_config
    logical, intent(in), optional :: rebuild
    logical, intent(in), optional :: ignore_mismatch
    type(phs_wood_config_t), intent(in), optional, target :: phs_config
    logical, intent(in), optional :: verbose
    logical :: no_strfun, verb
    type(string_t) :: nlo_type
    no_strfun = beam_config%n_strfun == 0
    verb = .true.;  if (present (verbose))  verb = verbose
    nlo_type = component%config%get_nlo_type ()
    if (present (phs_config)) then
       select type (config => component%phs_config)
       type is (phs_fks_config_t)
          call config%set_born_config (phs_config)
       end select
    end if
    call component%phs_config%configure (sqrts, &
         azimuthal_dependence = beam_config%azimuthal_dependence, &
         sqrts_fixed = no_strfun, &
         cm_frame = beam_config%lab_is_cm_frame .and. no_strfun, &
         rebuild = rebuild, ignore_mismatch = ignore_mismatch, &
         nlo_type = nlo_type)
    if (verb)  call component%phs_config%startup_message ()
  end subroutine process_component_configure_phs
    
\eatline
\nwendcode{}\nwbegindocs{460}\nwdocspar
\nwenddocs{}\nwbegincode{461}\moddef{Processes: process component: TBP}\plusendmoddef
  procedure :: increase_phs_par => process_component_increase_phs_par
\nwendcode{}\nwbegincode{462}\moddef{Processes: procedures}\plusendmoddef
  subroutine process_component_increase_phs_par (component)
    class(process_component_t), intent(inout) :: component
    select type (phs_config => component%phs_config)
    type is (phs_wood_config_t)
       call phs_config%increase_n_par ()
    end select
  end subroutine process_component_increase_phs_par

\eatline
\nwendcode{}\nwbegindocs{463}The process component possesses two MD5 sums: the checksum of the
component definition, which should be available when the component is
initialized, and the phase-space MD5 sum, which is available after
configuration.
\nwenddocs{}\nwbegincode{464}\moddef{Processes: process component: TBP}\plusendmoddef
  procedure :: compute_md5sum => process_component_compute_md5sum
\nwendcode{}\nwbegincode{465}\moddef{Processes: procedures}\plusendmoddef
  subroutine process_component_compute_md5sum (component)
    class(process_component_t), intent(inout) :: component
    component%md5sum_phs = component%phs_config%get_md5sum ()
  end subroutine process_component_compute_md5sum
  
\eatline
\nwendcode{}\nwbegindocs{466}Match phase-space channels with structure-function channels, where
applicable.

This calls a method of the {\tt{}phs{\char95}config} phase-space implementation.
\nwenddocs{}\nwbegincode{467}\moddef{Processes: process component: TBP}\plusendmoddef
  procedure :: collect_channels => process_component_collect_channels
\nwendcode{}\nwbegincode{468}\moddef{Processes: procedures}\plusendmoddef
  subroutine process_component_collect_channels (component, coll)
    class(process_component_t), intent(inout) :: component
    type(phs_channel_collection_t), intent(inout) :: coll
    call component%phs_config%collect_channels (coll)
  end subroutine process_component_collect_channels
    
\eatline
\nwendcode{}\nwbegindocs{469}Return the number of phase-space parameters.
\nwenddocs{}\nwbegincode{470}\moddef{Processes: process component: TBP}\plusendmoddef
  procedure :: get_n_phs_par => process_component_get_n_phs_par
\nwendcode{}\nwbegincode{471}\moddef{Processes: procedures}\plusendmoddef
  function process_component_get_n_phs_par (component) result (n_par)
    class(process_component_t), intent(in) :: component
    integer :: n_par
    n_par = component%phs_config%get_n_par ()
  end function process_component_get_n_phs_par
    
\eatline
\nwendcode{}\nwbegindocs{472}Return the incoming flavor combination as an array of PDG arrays,
one for a decay or two for a scattering process.
\nwenddocs{}\nwbegincode{473}\moddef{Processes: process component: TBP}\plusendmoddef
  procedure :: get_pdg_in => process_component_get_pdg_in
\nwendcode{}\nwbegincode{474}\moddef{Processes: procedures}\plusendmoddef
  function process_component_get_pdg_in (component) result (pdg_in)
    class(process_component_t), intent(in) :: component
    type(pdg_array_t), dimension(:), allocatable :: pdg_in
    type(pdg_array_t) :: pdg_tmp
    integer :: i
    associate (data => component%core%data)
      allocate (pdg_in (data%n_in))
      do i = 1, data%n_in
         pdg_tmp = data%flv_state(i,:)
         pdg_in(i) = sort_abs (pdg_tmp, unique = .true.)
      end do
    end associate
  end function process_component_get_pdg_in
  
\eatline
\nwendcode{}\nwbegindocs{475}\nwdocspar
\nwenddocs{}\nwbegincode{476}\moddef{Processes: process component: TBP}\plusendmoddef
  procedure :: extract_phs_config => process_component_extract_phs_config
  procedure :: restore_phs_config => process_component_restore_phs_config
\nwendcode{}\nwbegincode{477}\moddef{Processes: procedures}\plusendmoddef
  subroutine process_component_extract_phs_config (component, phs_config)
    class(process_component_t), intent(inout) :: component
    class(phs_config_t), intent(inout), allocatable :: phs_config
    call move_alloc (from = component%phs_config, to = phs_config)
  end subroutine process_component_extract_phs_config

  subroutine process_component_restore_phs_config (component, phs_config)
    class(process_component_t), intent(inout) :: component
    class(phs_config_t), intent(inout), allocatable :: phs_config
    call move_alloc (from = phs_config, to = component%phs_config)
  end subroutine process_component_restore_phs_config

\eatline
\nwendcode{}\nwbegindocs{478}\nwdocspar
\subsection{Process terms}
For straightforward tree-level calculations, each process component
corresponds to a unique elementary interaction.  However, in the case
of NLO calculations with subtraction terms, a process component may
split into several separate contributions to the scattering, which are
qualified by interactions with distinct kinematics and particle
content.  We represent their configuration as {\tt{}process{\char95}term{\char95}t}
objects, the actual instances will be introduced below as
{\tt{}term{\char95}instance{\char95}t}.  In any case, the process term contains an
elementary interaction with a definite quantum-number and momentum
content.

The index {\tt{}i{\char95}term{\char95}global} identifies the term relative to the
process.

The index {\tt{}i{\char95}component} identifies the process component which
generates this term, relative to the parent process.

The index {\tt{}i{\char95}term} identifies the term relative to the process
component (not the process).

The {\tt{}data} subobject holds all process constants.

The number of allowed flavor/helicity/color combinations is stored as
{\tt{}n{\char95}allowed}.  This is the total number of independent entries in the
density matrix.  For each combination, the index of the flavor,
helicity, and color state is stored in the arrays {\tt{}flv}, {\tt{}hel},
and {\tt{}col}, respectively.

The flag {\tt{}rearrange} is true if we need to rearrange the particles of the
hard interaction, to obtain the effective parton state.

The interaction {\tt{}int} holds the quantum state for the (resolved) hard
interaction, the parent-child relations of the particles, and their momenta.
The momenta are not filled yet; this is postponed to copies of {\tt{}int} which
go into the process instances.

If recombination is in effect, we should allocate {\tt{}int{\char95}eff} to describe the
rearranged partonic state.
\nwenddocs{}\nwbegincode{479}\moddef{Processes: process part types}\plusendmoddef
  type :: process_term_t
     integer :: i_term_global = 0
     integer :: i_component = 0
     integer :: i_term = 0
     integer :: n_allowed = 0
     integer :: n_allowed_sub_born = 0
     type(process_constants_t) :: data
     real(default) :: alpha_s = 0
     integer, dimension(:), allocatable :: flv, hel, col
     logical :: rearrange = .false.
     type(interaction_t) :: int
     type(interaction_t), pointer :: int_eff => null ()
   contains
   \LA{}Processes: process term: TBP\RA{}
  end type process_term_t
  
\eatline
\nwendcode{}\nwbegindocs{480}For the output, we skip the process constants and the tables of
allowed quantum numbers.  Those can also be read off from the
interaction object.
\nwenddocs{}\nwbegincode{481}\moddef{Processes: process term: TBP}\endmoddef
  procedure :: write => process_term_write
\nwendcode{}\nwbegincode{482}\moddef{Processes: procedures}\plusendmoddef
  subroutine process_term_write (term, unit)
    class(process_term_t), intent(in) :: term
    integer, intent(in), optional :: unit
    integer :: u
    u = given_output_unit (unit)
    write (u, "(1x,A,I0)")  "Term #", term%i_term_global
    write (u, "(3x,A,I0)")  "Process component index      = ", &
         term%i_component
    write (u, "(3x,A,I0)")  "Term index w.r.t. component  = ", &
         term%i_term
    write (u, "(3x,A,L1)")  "Rearrange partons            = ", &
         term%rearrange
    call write_separator (u)
    write (u, "(1x,A)")  "Hard interaction:"
    call write_separator (u)
    call interaction_write (term%int, u)
    if (term%rearrange) then
       call write_separator (u)
       write (u, "(1x,A)")  "Rearranged hard interaction:"
       call write_separator (u)
       call interaction_write (term%int_eff, u)
    end if
  end subroutine process_term_write
     
\eatline
\nwendcode{}\nwbegindocs{483}Write an account of all quantum number states and their current status.
\nwenddocs{}\nwbegincode{484}\moddef{Processes: process term: TBP}\plusendmoddef
  procedure :: write_state_summary => process_term_write_state_summary
\nwendcode{}\nwbegincode{485}\moddef{Processes: procedures}\plusendmoddef
  subroutine process_term_write_state_summary (term, core, unit)
    class(process_term_t), intent(in) :: term
    class(prc_core_t), intent(in) :: core
    integer, intent(in), optional :: unit
    integer :: u, i, f, h, c
    type(state_iterator_t) :: it
    character :: sgn
    u = given_output_unit (unit)
    write (u, "(1x,A,I0)")  "Term #", term%i_term_global
    call it%init (interaction_get_state_matrix_ptr (term%int))
    do while (it%is_valid ())
       i = it%get_me_index ()
       f = term%flv(i)
       h = term%hel(i)
       c = term%col(i)
       if (core%is_allowed (term%i_term, f, h, c)) then
          sgn = "+"
       else
          sgn = " "
       end if
       write (u, "(1x,A1,1x,I0,2x)", advance="no")  sgn, i
       call quantum_numbers_write (it%get_quantum_numbers (), u)
       write (u, *)
       call it%advance ()
    end do
  end subroutine process_term_write_state_summary
  
\eatline
\nwendcode{}\nwbegindocs{486}Finalizer: the {\tt{}int} and potentially {\tt{}int{\char95}eff} components have a
finalizer that we must call.
\nwenddocs{}\nwbegincode{487}\moddef{Processes: process term: TBP}\plusendmoddef
  procedure :: final => process_term_final
\nwendcode{}\nwbegincode{488}\moddef{Processes: procedures}\plusendmoddef
  subroutine process_term_final (term)
    class(process_term_t), intent(inout) :: term
    call interaction_final (term%int)
    if (term%rearrange) then
       call interaction_final (term%int_eff)
       deallocate (term%int_eff)
    end if
  end subroutine process_term_final

\eatline
\nwendcode{}\nwbegindocs{489}Initialize the term.  We copy the process constants from the {\tt{}core}
object and set up the {\tt{}int} hard interaction accordingly.

The {\tt{}alpha{\char95}s} value is useful for writing external event records.  This is
the constant value which may be overridden by a event-specific running value.
If the model does not contain the strong coupling, the value is zero.

The {\tt{}rearrange} part is commented out; this or something equivalent
could become relevant for NLO algorithms.
\nwenddocs{}\nwbegincode{490}\moddef{Processes: process term: TBP}\plusendmoddef
  procedure :: init => process_term_init
\nwendcode{}\nwbegincode{491}\moddef{Processes: procedures}\plusendmoddef
  subroutine process_term_init &
       (term, i_term_global, i_component, i_term, core, model)
    class(process_term_t), intent(inout), target :: term
    integer, intent(in) :: i_term_global
    integer, intent(in) :: i_component
    integer, intent(in) :: i_term
    class(prc_core_t), intent(in) :: core
    class(model_data_t), intent(in), target :: model
    class(modelpar_data_t), pointer :: alpha_s_ptr
!    type(var_list_t), pointer :: var_list
    term%i_term_global = i_term_global
    term%i_component = i_component
    term%i_term = i_term
    call core%get_constants (term%data, i_term)
!    var_list => model%get_var_list_ptr ()
!    if (var_list%contains (var_str ("alphas"))) then
!       term%alpha_s = var_list_get_rval (var_list, var_str ("alphas"))
!    else
!       term%alpha_s = -1
!    end if
    alpha_s_ptr => model%get_par_data_ptr (var_str ("alphas"))
    if (associated (alpha_s_ptr)) then
       term%alpha_s = alpha_s_ptr%get_real ()
    else
       term%alpha_s = -1
    end if

    call term%setup_interaction (core, model)
!     if (term%rearrange) then
!       call term%setup_effective_interaction (core, term%int, term%int_eff)
!     end if
  end subroutine process_term_init
    
\eatline
\nwendcode{}\nwbegindocs{492}We fetch the process constants which determine the quantum numbers and
use those to create the interaction.  The interaction contains
incoming and outgoing particles, no virtuals.  The incoming particles
are parents of the outgoing ones.

Keeping previous \whizard\ conventions, we invert the color assignment
(but not flavor or helicity) for the incoming particles.  When the
color-flow square matrix is evaluated, this inversion is done again,
so in the color-flow sequence we get the color assignments of the
matrix element.
\nwenddocs{}\nwbegincode{493}\moddef{Processes: process term: TBP}\plusendmoddef
  procedure :: setup_interaction => process_term_setup_interaction
\nwendcode{}\nwbegincode{494}\moddef{Processes: procedures}\plusendmoddef
  subroutine process_term_setup_interaction (term, core, model)
    class(process_term_t), intent(inout) :: term
    class(prc_core_t), intent(in) :: core
    class(model_data_t), intent(in), target :: model
    integer :: n_tot
    type(flavor_t), dimension(:), allocatable :: flv
    type(color_t), dimension(:), allocatable :: col
    type(helicity_t), dimension(:), allocatable :: hel
    type(quantum_numbers_t), dimension(:), allocatable :: qn
    integer :: i, n, f, h, c
    associate (data => term%data)
      n_tot = data%n_in + data%n_out
      n = 0
      do f = 1, data%n_flv
         do h = 1, data%n_hel
            do c = 1, data%n_col
               if (core%is_allowed (term%i_term, f, h, c))  n = n + 1
            end do
         end do
      end do
      allocate (term%flv (n), term%col (n), term%hel (n))
      term%n_allowed = n
      allocate (flv (n_tot), col (n_tot), hel (n_tot))
      allocate (qn (n_tot))
      call interaction_init &
           (term%int, data%n_in, 0, data%n_out, set_relations=.true.)
      i = 0
      do f = 1, data%n_flv
         do h = 1, data%n_hel
            do c = 1, data%n_col
               if (core%is_allowed (term%i_term, f, h, c)) then
                  i = i + 1
                  term%flv(i) = f
                  term%hel(i) = h
                  term%col(i) = c
                  call flv%init (data%flv_state(:,f), model)
                  call color_init_from_array (col, &
                       data%col_state(:,:,c), &
                       data%ghost_flag(:,c))
                  call col(:data%n_in)%invert ()
                  call hel%init (data%hel_state(:,h))
                  call qn%init (flv, col, hel)
                  call interaction_add_state (term%int, qn)
               end if
            end do
         end do
      end do
      call interaction_freeze (term%int)
    end associate
  end subroutine process_term_setup_interaction
  
\eatline
\nwendcode{}\nwbegindocs{495}\nwdocspar
\nwenddocs{}\nwbegincode{496}\moddef{Processes: process term: TBP}\plusendmoddef
  procedure :: fetch_constants => process_term_fetch_process_constants
\nwendcode{}\nwbegincode{497}\moddef{Processes: procedures}\plusendmoddef
   subroutine process_term_fetch_process_constants &
       (term, prc_constants)
    class(process_term_t), intent(inout) :: term
    type(process_constants_t), intent(out) :: prc_constants
    prc_constants = term%data
  end subroutine process_term_fetch_process_constants

\eatline
\nwendcode{}\nwbegindocs{498}\nwdocspar
\subsection{Default Iterations}
If the user does not specify the passes and iterations for
integration, we should be able to give reasonable defaults.  These
depend on the process, therefore we implement the following procedures
as methods of the process object.  The algorithm is not very
sophisticated yet, it may be improved by looking at the process in
more detail.

We investigate only the first process component, assuming that it
characterizes the complexity of the process reasonable well.

The number of passes is limited to two: one for adaption, one for
integration.
\nwenddocs{}\nwbegincode{499}\moddef{Processes: process: TBP}\plusendmoddef
  procedure :: get_n_pass_default => process_get_n_pass_default
  procedure :: adapt_grids_default => process_adapt_grids_default
  procedure :: adapt_weights_default => process_adapt_weights_default
\nwendcode{}\nwbegincode{500}\moddef{Processes: procedures}\plusendmoddef
  function process_get_n_pass_default (process) result (n_pass)
    class(process_t), intent(in) :: process
    integer :: n_pass
    integer :: n_eff
    type(process_component_def_t), pointer :: config
    config => process%component(1)%config
    n_eff = config%get_n_tot () - 2
    select case (n_eff)
    case (1)
       n_pass = 1
    case default
       n_pass = 2
    end select
  end function process_get_n_pass_default
  
  function process_adapt_grids_default (process, pass) result (flag)
    class(process_t), intent(in) :: process
    integer, intent(in) :: pass
    logical :: flag
    integer :: n_eff
    type(process_component_def_t), pointer :: config
    config => process%component(1)%config
    n_eff = config%get_n_tot () - 2
    select case (n_eff)
    case (1)
       flag = .false.
    case default
       select case (pass)
       case (1);  flag = .true.
       case (2);  flag = .false.
       case default
          call msg_bug ("adapt grids default: impossible pass index")
       end select
    end select
  end function process_adapt_grids_default
  
  function process_adapt_weights_default (process, pass) result (flag)
    class(process_t), intent(in) :: process
    integer, intent(in) :: pass
    logical :: flag
    integer :: n_eff
    type(process_component_def_t), pointer :: config
    config => process%component(1)%config
    n_eff = config%get_n_tot () - 2
    select case (n_eff)
    case (1)
       flag = .false.
    case default
       select case (pass)
       case (1);  flag = .true.
       case (2);  flag = .false.
       case default
          call msg_bug ("adapt weights default: impossible pass index")
       end select
    end select
  end function process_adapt_weights_default
  
\eatline
\eatline
\eatline
\nwendcode{}\nwbegindocs{501}The number of iterations and calls per iteration depends on the
number of outgoing particles.
\nwenddocs{}\nwbegincode{502}\moddef{Processes: process: TBP}\plusendmoddef
  procedure :: get_n_it_default => process_get_n_it_default
  procedure :: get_n_calls_default => process_get_n_calls_default
\nwendcode{}\nwbegincode{503}\moddef{Processes: procedures}\plusendmoddef
  function process_get_n_it_default (process, pass) result (n_it)
    class(process_t), intent(in) :: process
    integer, intent(in) :: pass
    integer :: n_it
    integer :: n_eff
    type(process_component_def_t), pointer :: config
    config => process%component(1)%config
    n_eff = config%get_n_tot () - 2
    select case (pass)
    case (1)
       select case (n_eff)
       case (1);   n_it = 1
       case (2);   n_it = 3
       case (3);   n_it = 5
       case (4:5); n_it = 10
       case (6);   n_it = 15
       case (7:);  n_it = 20
       end select
    case (2)
       select case (n_eff)
       case (:3);   n_it = 3
       case (4:);   n_it = 5
       end select
    end select
  end function process_get_n_it_default
  
  function process_get_n_calls_default (process, pass) result (n_calls)
    class(process_t), intent(in) :: process
    integer, intent(in) :: pass
    integer :: n_calls
    integer :: n_eff
    type(process_component_def_t), pointer :: config
    config => process%component(1)%config
    n_eff = config%get_n_tot () - 2
    select case (pass)
    case (1)
       select case (n_eff)
       case (1);   n_calls =   100
       case (2);   n_calls =  1000
       case (3);   n_calls =  5000
       case (4);   n_calls = 10000
       case (5);   n_calls = 20000
       case (6:);  n_calls = 50000
       end select
    case (2)
       select case (n_eff)
       case (:3);  n_calls =  10000
       case (4);   n_calls =  20000
       case (5);   n_calls =  50000
       case (6);   n_calls = 100000
       case (7:);  n_calls = 200000
       end select
    end select
  end function process_get_n_calls_default
  
\eatline
\eatline
\nwendcode{}\nwbegindocs{504}\nwdocspar
\subsection{Constant process data}
The following methods return basic process data that stay constant
after initialization.

The process and IDs.
\nwenddocs{}\nwbegincode{505}\moddef{Processes: process: TBP}\plusendmoddef
  procedure :: get_id => process_get_id
  procedure :: get_num_id => process_get_num_id
  procedure :: get_run_id => process_get_run_id
  procedure :: get_library_name => process_get_library_name
\nwendcode{}\nwbegincode{506}\moddef{Processes: procedures}\plusendmoddef
  function process_get_id (process) result (id)
    class(process_t), intent(in) :: process
    type(string_t) :: id
    id = process%meta%id
  end function process_get_id

  function process_get_num_id (process) result (id)
    class(process_t), intent(in) :: process
    integer :: id
    id = process%meta%num_id
  end function process_get_num_id

  function process_get_run_id (process) result (id)
    class(process_t), intent(in) :: process
    type(string_t) :: id
    id = process%meta%run_id
  end function process_get_run_id

  function process_get_library_name (process) result (id)
    class(process_t), intent(in) :: process
    type(string_t) :: id
    id = process%meta%lib%get_name ()
  end function process_get_library_name

\eatline
\eatline
\nwendcode{}\nwbegindocs{507}The number of incoming particles.
\nwenddocs{}\nwbegincode{508}\moddef{Processes: process: TBP}\plusendmoddef
  procedure :: get_n_in => process_get_n_in
\nwendcode{}\nwbegincode{509}\moddef{Processes: procedures}\plusendmoddef
  function process_get_n_in (process) result (n)
    class(process_t), intent(in) :: process
    integer :: n
    n = process%config%n_in
  end function process_get_n_in
    
\eatline
\nwendcode{}\nwbegindocs{510}The number of MCI data sets.
\nwenddocs{}\nwbegincode{511}\moddef{Processes: process: TBP}\plusendmoddef
  procedure :: get_n_mci => process_get_n_mci
\nwendcode{}\nwbegincode{512}\moddef{Processes: procedures}\plusendmoddef
  function process_get_n_mci (process) result (n)
    class(process_t), intent(in) :: process
    integer :: n
    n = process%config%n_mci
  end function process_get_n_mci
    
\eatline
\nwendcode{}\nwbegindocs{513}The number of process components, total.
\nwenddocs{}\nwbegincode{514}\moddef{Processes: process: TBP}\plusendmoddef
  procedure :: get_n_components => process_get_n_components
\nwendcode{}\nwbegincode{515}\moddef{Processes: procedures}\plusendmoddef
  function process_get_n_components (process) result (n)
    class(process_t), intent(in) :: process
    integer :: n
    n = process%meta%n_components
  end function process_get_n_components
    
\eatline
\nwendcode{}\nwbegindocs{516}The number of process terms, total.
\nwenddocs{}\nwbegincode{517}\moddef{Processes: process: TBP}\plusendmoddef
  procedure :: get_n_terms => process_get_n_terms
\nwendcode{}\nwbegincode{518}\moddef{Processes: procedures}\plusendmoddef
  function process_get_n_terms (process) result (n)
    class(process_t), intent(in) :: process
    integer :: n
    n = process%config%n_terms
  end function process_get_n_terms
    
\eatline
\nwendcode{}\nwbegindocs{519}Return the indices of the components that belong to a
specific MCI entry.
\nwenddocs{}\nwbegincode{520}\moddef{Processes: process: TBP}\plusendmoddef
  procedure :: get_i_component => process_get_i_component
\nwendcode{}\nwbegincode{521}\moddef{Processes: procedures}\plusendmoddef
  subroutine process_get_i_component (process, i_mci, i_component)
    class(process_t), intent(in) :: process
    integer, intent(in) :: i_mci
    integer, dimension(:), intent(out), allocatable :: i_component
    associate (mci_entry => process%mci_entry(i_mci))
      allocate (i_component (size (mci_entry%i_component)))
      i_component = mci_entry%i_component
    end associate
  end subroutine process_get_i_component

\eatline
\nwendcode{}\nwbegindocs{522}Return the ID of a specific component.
\nwenddocs{}\nwbegincode{523}\moddef{Processes: process: TBP}\plusendmoddef
  procedure :: get_component_id => process_get_component_id
\nwendcode{}\nwbegincode{524}\moddef{Processes: procedures}\plusendmoddef
  function process_get_component_id (process, i_component) result (id)
    class(process_t), intent(in) :: process
    integer, intent(in) :: i_component
    type(string_t) :: id
    id = process%meta%component_id(i_component)
  end function process_get_component_id
    
\eatline
\nwendcode{}\nwbegindocs{525}Return a pointer to the definition of a specific component.
\nwenddocs{}\nwbegincode{526}\moddef{Processes: process: TBP}\plusendmoddef
  procedure :: get_component_def_ptr => process_get_component_def_ptr
\nwendcode{}\nwbegincode{527}\moddef{Processes: procedures}\plusendmoddef
  function process_get_component_def_ptr (process, i_component) result (ptr)
    class(process_t), intent(in) :: process
    integer, intent(in) :: i_component
    type(process_component_def_t), pointer :: ptr
    ptr => process%meta%lib%get_component_def_ptr (process%meta%id, i_component)
  end function process_get_component_def_ptr
  
\eatline
\nwendcode{}\nwbegindocs{528}These procedures extract and restore (by transferring the
allocation) the process core.  This is useful for changing process
parameters from outside this module.
\nwenddocs{}\nwbegincode{529}\moddef{Processes: process: TBP}\plusendmoddef
  procedure :: extract_component_core => process_extract_component_core
  procedure :: restore_component_core => process_restore_component_core
\nwendcode{}\nwbegincode{530}\moddef{Processes: procedures}\plusendmoddef
  subroutine process_extract_component_core (process, i_component, core)
    class(process_t), intent(inout) :: process
    integer, intent(in) :: i_component
    class(prc_core_t), intent(inout), allocatable :: core
    call move_alloc (from = process%component(i_component)%core, to = core)
  end subroutine process_extract_component_core
    
  subroutine process_restore_component_core (process, i_component, core)
    class(process_t), intent(inout) :: process
    integer, intent(in) :: i_component
    class(prc_core_t), intent(inout), allocatable :: core
    call move_alloc (from = core, to = process%component(i_component)%core)
  end subroutine process_restore_component_core
    
\eatline
\eatline
\nwendcode{}\nwbegindocs{531}The block of process constants.
\nwenddocs{}\nwbegincode{532}\moddef{Processes: process: TBP}\plusendmoddef
  procedure :: get_constants => process_get_constants
\nwendcode{}\nwbegincode{533}\moddef{Processes: procedures}\plusendmoddef
  function process_get_constants (process, i) result (data)
    class(process_t), intent(in) :: process
    integer, intent(in) :: i
    type(process_constants_t) :: data
    data = process%component(i)%core%data
  end function process_get_constants
  
\eatline
\nwendcode{}\nwbegindocs{534}Return the set of outgoing flavors that are associated with a particular
term. We deduce this from the effective interaction.
\nwenddocs{}\nwbegincode{535}\moddef{Processes: process: TBP}\plusendmoddef
  procedure :: get_term_flv_out => process_get_term_flv_out
\nwendcode{}\nwbegincode{536}\moddef{Processes: procedures}\plusendmoddef
  subroutine process_get_term_flv_out (process, i_term, flv)
    class(process_t), intent(in), target :: process
    integer, intent(in) :: i_term
    type(flavor_t), dimension(:,:), allocatable, intent(out) :: flv
    type(interaction_t), pointer :: int
    int => process%term(i_term)%int_eff
    if (.not. associated (int))  int => process%term(i_term)%int
    call interaction_get_flv_out (int, flv)
  end subroutine process_get_term_flv_out

\eatline
\nwendcode{}\nwbegindocs{537}Return true if there is any unstable particle in any of the process
terms.  We decide this based on the provided model instance, not the
one that is stored in the process object.
\nwenddocs{}\nwbegincode{538}\moddef{Processes: process: TBP}\plusendmoddef
  procedure :: contains_unstable => process_contains_unstable
\nwendcode{}\nwbegincode{539}\moddef{Processes: procedures}\plusendmoddef
  function process_contains_unstable (process, model) result (flag)
    class(process_t), intent(in) :: process
    class(model_data_t), intent(in), target :: model
    logical :: flag
    integer :: i_term
    type(flavor_t), dimension(:,:), allocatable :: flv
    flag = .false.
    do i_term = 1, process%get_n_terms ()
       call process%get_term_flv_out (i_term, flv)
       call flv%set_model (model)
       flag = .not. all (flv%is_stable ())
       deallocate (flv)
       if (flag)  return
    end do
  end function process_contains_unstable
    
\eatline
\nwendcode{}\nwbegindocs{540}The nominal process energy.
\nwenddocs{}\nwbegincode{541}\moddef{Processes: process: TBP}\plusendmoddef
  procedure :: get_sqrts => process_get_sqrts
\nwendcode{}\nwbegincode{542}\moddef{Processes: procedures}\plusendmoddef
  function process_get_sqrts (process) result (sqrts)
    class(process_t), intent(in) :: process
    real(default) :: sqrts
    sqrts = beam_data_get_sqrts (process%beam_config%data)
  end function process_get_sqrts
  
\eatline
\nwendcode{}\nwbegincode{543}\moddef{Processes: process: TBP}\plusendmoddef
  procedure :: has_matrix_element => process_has_matrix_element
\nwendcode{}\nwbegincode{544}\moddef{Processes: procedures}\plusendmoddef
  function process_has_matrix_element (process, i) result (flag)
    class(process_t), intent(in) :: process
    integer, intent(in), optional :: i
    logical :: flag
    if (present (i)) then
       flag = process%component(i)%active
    else
       flag = any (process%component%active)
    end if
  end function process_has_matrix_element
  
\eatline
\nwendcode{}\nwbegindocs{545}Pointer to the beam data object.
\nwenddocs{}\nwbegincode{546}\moddef{Processes: process: TBP}\plusendmoddef
  procedure :: get_beam_data_ptr => process_get_beam_data_ptr
\nwendcode{}\nwbegincode{547}\moddef{Processes: procedures}\plusendmoddef
  function process_get_beam_data_ptr (process) result (beam_data)
    class(process_t), intent(in), target :: process
    type(beam_data_t), pointer :: beam_data
    beam_data => process%beam_config%data
  end function process_get_beam_data_ptr

\eatline
\nwendcode{}\nwbegindocs{548}Return true if lab and c.m.\ frame coincide for this process.
\nwenddocs{}\nwbegincode{549}\moddef{Processes: process: TBP}\plusendmoddef
  procedure :: cm_frame => process_cm_frame
\nwendcode{}\nwbegincode{550}\moddef{Processes: procedures}\plusendmoddef
  function process_cm_frame (process) result (flag)
    class(process_t), intent(in), target :: process
    logical :: flag
    type(beam_data_t), pointer :: beam_data
    beam_data => process%beam_config%data
    flag = beam_data_cm_frame (beam_data)
  end function process_cm_frame
  
\eatline
\nwendcode{}\nwbegindocs{551}Get the PDF set currently in use, if any.
\nwenddocs{}\nwbegincode{552}\moddef{Processes: process: TBP}\plusendmoddef
  procedure :: get_pdf_set => process_get_pdf_set
\nwendcode{}\nwbegincode{553}\moddef{Processes: procedures}\plusendmoddef
  function process_get_pdf_set (process) result (pdf_set)
    class(process_t), intent(in) :: process
    integer :: pdf_set
    pdf_set = process%beam_config%get_pdf_set ()
  end function process_get_pdf_set
  
\eatline
\nwendcode{}\nwbegindocs{554}Pointer to the process variable list.
\nwenddocs{}\nwbegincode{555}\moddef{Processes: process: TBP}\plusendmoddef
  procedure :: get_var_list_ptr => process_get_var_list_ptr
\nwendcode{}\nwbegincode{556}\moddef{Processes: procedures}\plusendmoddef
  function process_get_var_list_ptr (process) result (ptr)
    class(process_t), intent(in), target :: process
    type(var_list_t), pointer :: ptr
    ptr => process%meta%var_list
  end function process_get_var_list_ptr
  
\eatline
\nwendcode{}\nwbegindocs{557}Pointer to the common model.
\nwenddocs{}\nwbegincode{558}\moddef{Processes: process: TBP}\plusendmoddef
  procedure :: get_model_ptr => process_get_model_ptr
\nwendcode{}\nwbegincode{559}\moddef{Processes: procedures}\plusendmoddef
  function process_get_model_ptr (process) result (ptr)
    class(process_t), intent(in) :: process
    class(model_data_t), pointer :: ptr
    ptr => process%config%model
  end function process_get_model_ptr
  
\eatline
\nwendcode{}\nwbegindocs{560}Use the embedded RNG factory to spawn a new random-number generator
instance.  (This modifies the state of the factory.)
\nwenddocs{}\nwbegincode{561}\moddef{Processes: process: TBP}\plusendmoddef
  procedure :: make_rng => process_make_rng
\nwendcode{}\nwbegincode{562}\moddef{Processes: procedures}\plusendmoddef
  subroutine process_make_rng (process, rng)
    class(process_t), intent(inout) :: process
    class(rng_t), intent(out), allocatable :: rng
    if (allocated (process%config%rng_factory)) then
       call process%config%rng_factory%make (rng)
    else
       call msg_bug ("Process: make rng: factory not allocated")
    end if
  end subroutine process_make_rng
  
\eatline
\nwendcode{}\nwbegindocs{563}\nwdocspar
\subsection{Compute an amplitude}
Each process variant should allow for computing an amplitude value
directly, without generating a process instance. 

The process component is selected by the index {\tt{}i}.  The term within the
process component is selected by {\tt{}j}.  The momentum
combination is transferred as the array {\tt{}p}.  The function sets the specific
quantum state via the indices of a flavor {\tt{}f}, helicity {\tt{}h}, and color
{\tt{}c} combination.  Each index refers to the list of flavor, helicity, and
color states, respectively, as stored in the process data.

Optionally, we may set factorization and renormalization scale.  If unset, the
partonic c.m.\ energy is inserted.

The function checks arguments for validity.
For invalid arguments (quantum states), we return zero.
\nwenddocs{}\nwbegincode{564}\moddef{Processes: process: TBP}\plusendmoddef
  procedure :: compute_amplitude => process_compute_amplitude
\nwendcode{}\nwbegincode{565}\moddef{Processes: procedures}\plusendmoddef
  function process_compute_amplitude &
       (process, i, j, p, f, h, c, fac_scale, ren_scale, alpha_qcd_forced) &
       result (amp)
    class(process_t), intent(in) :: process
    integer, intent(in) :: i, j
    type(vector4_t), dimension(:), intent(in) :: p
    integer, intent(in) :: f, h, c
    real(default), intent(in), optional :: fac_scale, ren_scale
    real(default), intent(in), allocatable, optional :: alpha_qcd_forced
    real(default) :: fscale, rscale
    real(default), allocatable :: aqcd_forced
    complex(default) :: amp
    amp = 0
    if (0 < i .and. i <= process%meta%n_components) then
       if (process%component(i)%active) then
          associate (data => process%component(i)%core%data)
            if (size (p) == data%n_in + data%n_out &
                 .and. 0 < f .and. f <= data%n_flv &
                 .and. 0 < h .and. h <= data%n_hel &
                 .and. 0 < c .and. c <= data%n_col) then
               if (present (fac_scale)) then
                  fscale = fac_scale
               else
                  fscale = sum (p(data%n_in+1:)) ** 1
               end if
               if (present (ren_scale)) then
                  rscale = ren_scale
               else
                  rscale = fscale
               end if
               if (present (alpha_qcd_forced)) then
                  if (allocated (alpha_qcd_forced)) &
                       allocate (aqcd_forced, source = alpha_qcd_forced)
               end if
               amp = process%component(i)%core%compute_amplitude &
                    (j, p, f, h, c, fscale, rscale, aqcd_forced)
            end if
          end associate
       else
          amp = 0
       end if
    end if
  end function process_compute_amplitude

\eatline
\nwendcode{}\nwbegindocs{566}\nwdocspar
\nwenddocs{}\nwbegincode{567}\moddef{Processes: process: TBP}\plusendmoddef
  procedure :: set_component_type => process_set_component_type
\nwendcode{}\nwbegincode{568}\moddef{Processes: procedures}\plusendmoddef
  subroutine process_set_component_type (process, i_component, i_type)
    class(process_t), intent(inout) :: process
    integer, intent(in) :: i_component, i_type
    process%component(i_component)%component_type = i_type
  end subroutine process_set_component_type

\eatline
\nwendcode{}\nwbegindocs{569}This is for suppression of numerical noise in the integration results
stored in the {\tt{}process{\char95}mci{\char95}entry} type. As the error and efficiency
enter the MD5 sum, we recompute it. 
\nwenddocs{}\nwbegincode{570}\moddef{Processes: process: TBP}\plusendmoddef
  procedure :: pacify => process_pacify
\nwendcode{}\nwbegincode{571}\moddef{Processes: procedures}\plusendmoddef
  subroutine process_pacify (process, efficiency_reset, error_reset)
    class(process_t), intent(inout) :: process
    logical, intent(in), optional :: efficiency_reset, error_reset
    logical :: eff_reset, err_reset
    integer :: i
    eff_reset = .false.
    err_reset = .false.
    if (present (efficiency_reset))  eff_reset = efficiency_reset
    if (present (error_reset))  err_reset = error_reset
    if (allocated (process%mci_entry)) then
       do i = 1, size (process%mci_entry)
          call process%mci_entry(i)%results%pacify (efficiency_reset)
          if (allocated (process%mci_entry(i)%mci)) then
             if (process%mci_entry(i)%mci%error_known .and. err_reset) &
               process%mci_entry(i)%mci%error = 0
             if (process%mci_entry(i)%mci%efficiency_known .and. &
               eff_reset)  process%mci_entry(i)%mci%efficiency = 1             
             select type (mci => process%mci_entry(i)%mci) 
             type is (mci_vamp_t)
                call mci%pacify (efficiency_reset, error_reset)
                call mci%compute_md5sum ()
             end select
          end if
       end do
    end if
  end subroutine process_pacify

\eatline
\nwendcode{}\nwbegindocs{572}\nwdocspar
\subsection{Process instances}

\subsubsection{Kinematics instance}
In this data type we combine all objects (instances) necessary for
generating (or recovering) a kinematical configuration.  The
components work together as an implementation of multi-channel phase
space.

{\tt{}sf{\char95}chain} is an instance of the structure-function chain.  It is
used both for generating kinematics and, after the proper scale has
been determined, evaluating the structure function entries.

{\tt{}phs} is an instance of the phase space for the elementary process.

The array {\tt{}f} contains the products of the Jacobians that originate
from parameter mappings in the structure-function chain or in the
phase space.  We allocate this explicitly if either {\tt{}sf{\char95}chain} or
{\tt{}phs} are explicitly allocated, otherwise we can take over a pointer.

All components are implemented as pointers to (anonymous) targets.
For each component, there is a flag that tells whether this component
is to be regarded as a proper component (`owned' by the object) or as
a pointer.
\nwenddocs{}\nwbegincode{573}\moddef{Processes: types}\plusendmoddef
  type :: kinematics_t
     integer :: n_in = 0
     integer :: n_channel = 0
     integer :: selected_channel = 0
     type(sf_chain_instance_t), pointer :: sf_chain => null ()
     class(phs_t), pointer :: phs => null ()
     real(default), dimension(:), pointer :: f => null ()
     real(default) :: phs_factor
     logical :: sf_chain_allocated = .false.
     logical :: phs_allocated = .false.
     logical :: f_allocated = .false.
     type(string_t) :: nlo_type
     integer :: emitter
   contains
   \LA{}Processes: kinematics: TBP\RA{}
  end type kinematics_t

\eatline
\nwendcode{}\nwbegindocs{574}Output.  Show only those components which are marked as owned.
\nwenddocs{}\nwbegincode{575}\moddef{Processes: kinematics: TBP}\endmoddef
  procedure :: write => kinematics_write
\nwendcode{}\nwbegincode{576}\moddef{Processes: procedures}\plusendmoddef
  subroutine kinematics_write (object, unit)
    class(kinematics_t), intent(in) :: object
    integer, intent(in), optional :: unit
    integer :: u, c
    u = given_output_unit (unit)
    if (object%f_allocated) then
       write (u, "(1x,A)")  "Flux * PHS volume:"
       write (u, "(2x,ES19.12)")  object%phs_factor
       write (u, "(1x,A)")  "Jacobian factors per channel:"
       do c = 1, size (object%f)
          write (u, "(3x,I0,':',1x,ES13.7)", advance="no")  c, object%f(c)
          if (c == object%selected_channel) then
             write (u, "(1x,A)")  "[selected]"
          else
             write (u, *)
          end if
       end do
    end if
    if (object%sf_chain_allocated) then
       call write_separator (u)
       call object%sf_chain%write (u)
    end if
    if (object%phs_allocated) then
       call write_separator (u)
       call object%phs%write (u)
    end if
  end subroutine kinematics_write
    
\eatline
\nwendcode{}\nwbegindocs{577}Finalizer.  Delete only those components which are marked as owned.
\nwenddocs{}\nwbegincode{578}\moddef{Processes: kinematics: TBP}\plusendmoddef
  procedure :: final => kinematics_final
\nwendcode{}\nwbegincode{579}\moddef{Processes: procedures}\plusendmoddef
  subroutine kinematics_final (object)
    class(kinematics_t), intent(inout) :: object
    if (object%sf_chain_allocated) then
       call object%sf_chain%final ()
       deallocate (object%sf_chain)
       object%sf_chain_allocated = .false.
    end if
    if (object%phs_allocated) then
       call object%phs%final ()
       deallocate (object%phs)
       object%phs_allocated = .false.
    end if
    if (object%f_allocated) then
       deallocate (object%f)
       object%f_allocated = .false.
    end if
  end subroutine kinematics_final
  
\eatline
\nwendcode{}\nwbegindocs{580}Set the flags indicating whether the phase space shall be set up for the calculation of the real contribution. For this case, also set the emitter.
\nwenddocs{}\nwbegincode{581}\moddef{Processes: kinematics: TBP}\plusendmoddef
  procedure :: set_nlo_data => kinematics_set_nlo_data
\nwendcode{}\nwbegincode{582}\moddef{Processes: procedures}\plusendmoddef
  subroutine kinematics_set_nlo_data (k, nlo_type, emitter)
    class(kinematics_t), intent(inout) :: k
!    integer, intent(in) :: nlo_type
    type(string_t) :: nlo_type
    integer, intent(in), optional :: emitter
    k%nlo_type = nlo_type
    if (present (emitter)) then
      k%emitter = emitter
    end if
  end subroutine kinematics_set_nlo_data

\eatline
\nwendcode{}\nwbegindocs{583}Allocate the structure-function chain instance, initialize it as a
copy of the {\tt{}sf{\char95}chain} template, and prepare it for evaluation.

The {\tt{}sf{\char95}chain} remains a target because the (usually constant) beam momenta
are taken from there.
\nwenddocs{}\nwbegincode{584}\moddef{Processes: kinematics: TBP}\plusendmoddef
  procedure :: init_sf_chain => kinematics_init_sf_chain
\nwendcode{}\nwbegincode{585}\moddef{Processes: procedures}\plusendmoddef
  subroutine kinematics_init_sf_chain (k, core, sf_chain, config, tmp)
    class(kinematics_t), intent(inout) :: k
    class(prc_core_t), intent(in) :: core
    type(sf_chain_t), intent(in), target :: sf_chain
    type(process_beam_config_t), intent(in) :: config
    class(workspace_t), intent(inout), allocatable :: tmp
    integer :: n_strfun, n_channel
    integer :: c
    k%n_in = beam_data_get_n_in (config%data)
    n_strfun = config%n_strfun
    n_channel = config%n_channel
    allocate (k%sf_chain)
    k%sf_chain_allocated = .true.
    call core%init_sf_chain (k%sf_chain, sf_chain, n_channel, tmp)
    if (n_strfun /= 0) then
       do c = 1, n_channel
          call k%sf_chain%set_channel (c, config%sf_channel(c))
       end do
    end if
    call k%sf_chain%link_interactions ()
    call k%sf_chain%exchange_mask ()
    call k%sf_chain%init_evaluators ()
  end subroutine kinematics_init_sf_chain

\eatline
\nwendcode{}\nwbegindocs{586}Allocate and initialize the phase-space part and the array of
Jacobian factors.
\nwenddocs{}\nwbegincode{587}\moddef{Processes: kinematics: TBP}\plusendmoddef
  procedure :: init_phs => kinematics_init_phs
\nwendcode{}\nwbegincode{588}\moddef{Processes: procedures}\plusendmoddef
  subroutine kinematics_init_phs (k, config)
    class(kinematics_t), intent(inout) :: k
    class(phs_config_t), intent(in), target :: config
    k%n_channel = config%get_n_channel ()
    call config%allocate_instance (k%phs)
    call k%phs%init (config)
    k%phs_allocated = .true.
    allocate (k%f (k%n_channel))
    k%f = 0
    k%f_allocated = .true.
  end subroutine kinematics_init_phs
    
\eatline
\nwendcode{}\nwbegindocs{589}Initialize the kinematics in form of simple pointers.  In essence,
this is a shallow copy, but we have to set the flags correctly to
indicate this fact.
\nwenddocs{}\nwbegincode{590}\moddef{Processes: kinematics: TBP}\plusendmoddef
  procedure :: init_ptr => kinematics_init_ptr
\nwendcode{}\nwbegincode{591}\moddef{Processes: procedures}\plusendmoddef
  subroutine kinematics_init_ptr (k, k_in)
    class(kinematics_t), intent(out) :: k
    type(kinematics_t), intent(in) :: k_in
    k%n_in = k_in%n_in
    k%n_channel = k_in%n_channel
    k%sf_chain => k_in%sf_chain
    k%phs => k_in%phs
    k%f => k_in%f
  end subroutine kinematics_init_ptr
  
\eatline
\nwendcode{}\nwbegindocs{592}Generate kinematics, given a phase-space channel and a MC
parameter set. The main result is the momentum array {\tt{}p}, but we
also fill the momentum entries in the structure-function chain and the
Jacobian-factor array {\tt{}f}.  Regarding phase space, We fill only the
parameter arrays for the selected channel.
\nwenddocs{}\nwbegincode{593}\moddef{Processes: kinematics: TBP}\plusendmoddef
  procedure :: compute_selected_channel => kinematics_compute_selected_channel
\nwendcode{}\nwbegincode{594}\moddef{Processes: procedures}\plusendmoddef
  subroutine kinematics_compute_selected_channel &
       (k, mci_work, phs_channel, p, success, nlo_data)
    class(kinematics_t), intent(inout) :: k
    type(mci_work_t), intent(in) :: mci_work
    integer, intent(in) :: phs_channel
    type(vector4_t), dimension(:), intent(out) :: p
    logical, intent(out) :: success
    type(nlo_data_t), intent(inout), optional :: nlo_data
    integer :: sf_channel
    real(default) :: xit, phi
    real(default), dimension(:), allocatable :: xi_max, y
    real(default), dimension(3) :: jac
    real(default), dimension(:), allocatable :: jac_rand
    k%selected_channel = phs_channel
    sf_channel = k%phs%config%get_sf_channel (phs_channel)
    call k%sf_chain%compute_kinematics (sf_channel, mci_work%get_x_strfun ())
    call k%sf_chain%get_out_momenta (p(1:k%n_in))
    call k%phs%set_incoming_momenta (p(1:k%n_in))
    call k%phs%compute_flux ()
    call k%phs%select_channel (phs_channel)
    select type (phs => k%phs)
    type is (phs_fks_t)
      call phs%set_emitters (nlo_data%reg_data%emitters)
      if (nlo_data%has_massive_emitter ()) &
           call phs%enable_massive_emitters ()
      call phs%evaluate_selected_channel &
           (phs_channel, mci_work%get_x_process ())
      if (phs%q_defined) then
         call phs%get_real_kinematics (xit, y, phi, xi_max, jac, jac_rand)
         if (present (nlo_data)) &
              call nlo_data%set_real_kinematics (xit, y, phi, xi_max, jac, jac_rand)
         call phs%get_born_momenta (p)
         k%phs_factor = phs%get_overall_factor ()
         success = .true.
      else
         k%phs_factor = 0
         success = .false.
      end if
    class default
       call phs%evaluate_selected_channel &
            (phs_channel, mci_work%get_x_process ())
      if (phs%q_defined) then
         call k%phs%get_outgoing_momenta (p(k%n_in+1:))
         k%phs_factor = k%phs%get_overall_factor ()
         success = .true.
      else
       k%phs_factor = 0
       success = .false.
      end if
    end select
  end subroutine kinematics_compute_selected_channel
  
\eatline
\nwendcode{}\nwbegindocs{595}Complete kinematics by filling the non-selected phase-space parameter
arrays.
\nwenddocs{}\nwbegincode{596}\moddef{Processes: kinematics: TBP}\plusendmoddef
  procedure :: compute_other_channels => kinematics_compute_other_channels
\nwendcode{}\nwbegincode{597}\moddef{Processes: procedures}\plusendmoddef
  subroutine kinematics_compute_other_channels (k, mci_work, phs_channel)
    class(kinematics_t), intent(inout) :: k
    type(mci_work_t), intent(in) :: mci_work
    integer, intent(in) :: phs_channel
    integer :: c, c_sf
    call k%phs%evaluate_other_channels (phs_channel)
    do c = 1, k%n_channel
       c_sf = k%phs%config%get_sf_channel (c)
       k%f(c) = k%sf_chain%get_f (c_sf) * k%phs%get_f (c)
    end do
  end subroutine kinematics_compute_other_channels
  
\eatline
\nwendcode{}\nwbegindocs{598}Just fetch the outgoing momenta of the {\tt{}sf{\char95}chain} subobject, which
become the incoming (seed) momenta of the hard interaction.

This is a stripped down-version of the above which we use when
recovering kinematics.  Momenta are known, but no MC parameters yet.

(We do not use the {\tt{}get{\char95}out{\char95}momenta} method of the chain, since this
relies on the structure-function interactions, which are not necessary
filled here.  We do rely on the momenta of the last evaluator in the
chain, however.)
\nwenddocs{}\nwbegincode{599}\moddef{Processes: kinematics: TBP}\plusendmoddef
  procedure :: get_incoming_momenta => kinematics_get_incoming_momenta
\nwendcode{}\nwbegincode{600}\moddef{Processes: procedures}\plusendmoddef
  subroutine kinematics_get_incoming_momenta (k, p)
    class(kinematics_t), intent(in) :: k
    type(vector4_t), dimension(:), intent(out) :: p
    type(interaction_t), pointer :: int
    integer :: i
    int => k%sf_chain%get_out_int_ptr ()
    do i = 1, k%n_in
       p(i) = interaction_get_momentum (int, k%sf_chain%get_out_i (i))
    end do
  end subroutine kinematics_get_incoming_momenta
  
\eatline
\nwendcode{}\nwbegindocs{601}This inverts the remainder of the above {\tt{}compute} method.  We know
the momenta and recover the rest, as far as needed.  If we select a
channel, we can complete the inversion and reconstruct the 
MC parameter set.
\nwenddocs{}\nwbegincode{602}\moddef{Processes: kinematics: TBP}\plusendmoddef
  procedure :: recover_mcpar => kinematics_recover_mcpar
\nwendcode{}\nwbegincode{603}\moddef{Processes: procedures}\plusendmoddef
  subroutine kinematics_recover_mcpar (k, mci_work, phs_channel, p)
    class(kinematics_t), intent(inout) :: k
    type(mci_work_t), intent(inout) :: mci_work
    integer, intent(in) :: phs_channel
    type(vector4_t), dimension(:), intent(in) :: p
    integer :: c, c_sf
    real(default), dimension(:), allocatable :: x_sf, x_phs
    c = phs_channel
    c_sf = k%phs%config%get_sf_channel (c)
    k%selected_channel = c
    call k%sf_chain%recover_kinematics (c_sf)
    call k%phs%set_incoming_momenta (p(1:k%n_in))
    call k%phs%compute_flux ()
    call k%phs%set_outgoing_momenta (p(k%n_in+1:))
    call k%phs%inverse ()
    do c = 1, k%n_channel
       c_sf = k%phs%config%get_sf_channel (c)
       k%f(c) = k%sf_chain%get_f (c_sf) * k%phs%get_f (c)
    end do
    k%phs_factor = k%phs%get_overall_factor ()
    c = phs_channel
    c_sf = k%phs%config%get_sf_channel (c)
    allocate (x_sf (k%sf_chain%config%get_n_bound ()))
    allocate (x_phs (k%phs%config%get_n_par ()))
    call k%phs%select_channel (c)
    call k%sf_chain%get_mcpar (c_sf, x_sf)
    call k%phs%get_mcpar (c, x_phs)
    call mci_work%set_x_strfun (x_sf)
    call mci_work%set_x_process (x_phs)
  end subroutine kinematics_recover_mcpar

\eatline
\nwendcode{}\nwbegindocs{604}Retrieve the MC input parameter array for a specific channel.  We assume
that the kinematics is complete, so this is known for all channels.
\nwenddocs{}\nwbegincode{605}\moddef{Processes: kinematics: TBP}\plusendmoddef
  procedure :: get_mcpar => kinematics_get_mcpar
\nwendcode{}\nwbegincode{606}\moddef{Processes: procedures}\plusendmoddef
  subroutine kinematics_get_mcpar (k, phs_channel, r)
    class(kinematics_t), intent(in) :: k
    integer, intent(in) :: phs_channel
    real(default), dimension(:), intent(out) :: r
    integer :: sf_channel, n_par_sf, n_par_phs
    sf_channel = k%phs%config%get_sf_channel (phs_channel)
    n_par_phs = k%phs%config%get_n_par ()
    n_par_sf = k%sf_chain%config%get_n_bound ()
    if (n_par_sf > 0) then
       call k%sf_chain%get_mcpar (sf_channel, r(1:n_par_sf))
    end if
    if (n_par_phs > 0) then
       call k%phs%get_mcpar (phs_channel, r(n_par_sf+1:))
    end if
  end subroutine kinematics_get_mcpar
  
\eatline
\nwendcode{}\nwbegindocs{607}Evaluate the structure function chain, assuming that kinematics is known.

The status must be precisely {\tt{}SF{\char95}DONE{\char95}KINEMATICS}.  We thus avoid
evaluating the chain twice via different pointers to the same target.
\nwenddocs{}\nwbegincode{608}\moddef{Processes: kinematics: TBP}\plusendmoddef
  procedure :: evaluate_sf_chain => kinematics_evaluate_sf_chain
\nwendcode{}\nwbegincode{609}\moddef{Processes: procedures}\plusendmoddef
  subroutine kinematics_evaluate_sf_chain (k, fac_scale)
    class(kinematics_t), intent(inout) :: k
    real(default), intent(in) :: fac_scale
    select case (k%sf_chain%get_status ())
    case (SF_DONE_KINEMATICS)
       call k%sf_chain%evaluate (fac_scale)
    end select
  end subroutine kinematics_evaluate_sf_chain
  
\eatline
\nwendcode{}\nwbegindocs{610}Recover beam momenta, i.e., return the beam momenta stored in the
current {\tt{}sf{\char95}chain} to their source.  This is a side effect.
\nwenddocs{}\nwbegincode{611}\moddef{Processes: kinematics: TBP}\plusendmoddef
  procedure :: return_beam_momenta => kinematics_return_beam_momenta
\nwendcode{}\nwbegincode{612}\moddef{Processes: procedures}\plusendmoddef
  subroutine kinematics_return_beam_momenta (k)
    class(kinematics_t), intent(in) :: k
    call k%sf_chain%return_beam_momenta ()
  end subroutine kinematics_return_beam_momenta
  
\eatline
\nwendcode{}\nwbegindocs{613}\nwdocspar
\subsubsection{Process component instance}
The actual calculation of a sampling point is done from here.

The {\tt{}config} pointer accesses the corresponding
configuration in the {\tt{}process} object.

The {\tt{}active} flag indicates that we are currently computing this
component, together with all other components that share the same MC
parameter set.  Inactive components are using a different MC parameter
set and are not in use for this sampling point.

The {\tt{}k{\char95}seed} subobject contains the kinematics (structure-function
chain, phase space, etc.) that implements the `seed' configuration of
momenta.  This version of the process kinematics is accessed by the MCI
setup.

{\tt{}p{\char95}seed} is the array of momenta that we compute from the MC input
parameters, via the {\tt{}k{\char95}seed} subobject.  Depending on the process
variant, these may or may not coincide with the momenta that enter the
process terms associated to this component.

The {\tt{}tmp} object can be used for storing intermediate results.  Its precise
type and contents depend on the process variant.
\nwenddocs{}\nwbegincode{614}\moddef{Processes: types}\plusendmoddef
  type :: component_instance_t
     type(process_component_t), pointer :: config => null ()
     logical :: active = .false.
     type(kinematics_t) :: k_seed
     type(vector4_t), dimension(:), allocatable :: p_seed
     logical :: sqme_known = .false.
     real(default) :: sqme = 0
     class(workspace_t), allocatable :: tmp
     type(nlo_data_t), pointer :: nlo_data => null ()
   contains
   \LA{}Processes: component instance: TBP\RA{}
  end type component_instance_t
  
\eatline
\nwendcode{}\nwbegindocs{615}In the header, fetch the component index from the configuration record.
{\tt{}process{\char95}component{\char95}t} configuration block.

We write the {\tt{}sf{\char95}chain} subobject only upon request, since its instances
appear elsewhere.
\nwenddocs{}\nwbegincode{616}\moddef{Processes: component instance: TBP}\endmoddef
  procedure :: write => component_instance_write
\nwendcode{}\nwbegincode{617}\moddef{Processes: procedures}\plusendmoddef
  subroutine component_instance_write (object, unit, testflag)
    class(component_instance_t), intent(in) :: object
    integer, intent(in), optional :: unit
    logical, intent(in), optional :: testflag
    integer :: u, i
    u = given_output_unit (unit)
    if (object%active) then
       if (associated (object%config)) then
          write (u, "(1x,A,I0)")  "Component #", object%config%index
       else
          write (u, "(1x,A)")  "Component [undefined]"
       end if
    else
       write (u, "(1x,A,I0,A)")  "Component #", object%config%index, &
            " [inactive]"
    end if
    if (allocated (object%p_seed)) then
       write (u, "(1x,A)")  "Seed momenta:"
       do i = 1, size (object%p_seed)
          call vector4_write (object%p_seed(i), u, testflag = testflag)
       end do
    end if
    write (u, "(1x,A)")  "Squared matrix element:"
    if (object%sqme_known) then
       write (u, "(2x,ES19.12)")  object%sqme
    else
       write (u, "(2x,A)")  "[undefined]"
    end if
    call object%k_seed%write (u)
    if (allocated (object%tmp)) then
       call write_separator (u)
       call object%tmp%write (u)
    end if
  end subroutine component_instance_write
    
\eatline
\nwendcode{}\nwbegindocs{618}Finalizer
\nwenddocs{}\nwbegincode{619}\moddef{Processes: component instance: TBP}\plusendmoddef
  procedure :: final => component_instance_final
\nwendcode{}\nwbegincode{620}\moddef{Processes: procedures}\plusendmoddef
  subroutine component_instance_final (object)
    class(component_instance_t), intent(inout) :: object
    call object%k_seed%final ()
  end subroutine component_instance_final
  
\eatline
\nwendcode{}\nwbegindocs{621}Initialize: associate the configuration pointer.  Also initialize the
process workspace, if there is anything to do.  This initialization is a
matter of the process core.
\nwenddocs{}\nwbegincode{622}\moddef{Processes: component instance: TBP}\plusendmoddef
  procedure :: init => component_instance_init
\nwendcode{}\nwbegincode{623}\moddef{Processes: procedures}\plusendmoddef
  subroutine component_instance_init (component, config)
    class(component_instance_t), intent(out) :: component
    type(process_component_t), intent(in), target :: config
    type(string_t) :: nlo_type
    integer :: n_in, n_tot
    component%config => config
    nlo_type = component%config%config%get_nlo_type ()
    associate (core => component%config%core)
      n_in = core%data%n_in
      n_tot = n_in + core%data%n_out
      select case (char (nlo_type))
      case ('Real')
        allocate (component%p_seed (n_tot -1))
      case ('Born', 'Virtual')
        allocate (component%p_seed (n_tot))
      end select
      call core%allocate_workspace (component%tmp)
    end associate
  end subroutine component_instance_init

\eatline
\nwendcode{}\nwbegindocs{624}Initialize the seed-kinematics configuration.  All subobjects are
allocated explicitly.
\nwenddocs{}\nwbegincode{625}\moddef{Processes: component instance: TBP}\plusendmoddef
  procedure :: setup_kinematics => component_instance_setup_kinematics
\nwendcode{}\nwbegincode{626}\moddef{Processes: procedures}\plusendmoddef
  subroutine component_instance_setup_kinematics (component, sf_chain, config, &
                                                  combined_integration)
    class(component_instance_t), intent(inout) :: component
    type(sf_chain_t), intent(in), target :: sf_chain
    type(process_beam_config_t), intent(in) :: config
    logical, intent(in), optional :: combined_integration
    class(phs_config_t), pointer :: fks_config
    type(string_t) :: nlo_type
    nlo_type = component%config%config%get_nlo_type ()
    if (component%config%phs_config%nlo_type == "") &
      component%config%phs_config%nlo_type = nlo_type
    if (present (combined_integration)) then
       if (combined_integration) then
          select type (phs_config => component%config%phs_config)
          type is (phs_wood_config_t)
             call phs_config%set_extended_phs ()
          end select
       end if
    end if
    call component%k_seed%init_sf_chain &
         (component%config%core, sf_chain, config, component%tmp)
    call component%k_seed%init_phs (component%config%phs_config)
  end subroutine component_instance_setup_kinematics

\eatline
\nwendcode{}\nwbegindocs{627}Setup seed kinematics, starting from the MC parameter set given as
argument.  As a result, the {\tt{}k{\char95}seed} kinematics object is evaluated
(except for the structure-function matrix-element evaluation, which we
postpone until we know the factorization scale), and we have a valid
{\tt{}p{\char95}seed} momentum array.
\nwenddocs{}\nwbegincode{628}\moddef{Processes: component instance: TBP}\plusendmoddef
  procedure :: compute_seed_kinematics => &
       component_instance_compute_seed_kinematics
\nwendcode{}\nwbegincode{629}\moddef{Processes: procedures}\plusendmoddef
  subroutine component_instance_compute_seed_kinematics &
       (component, mci_work, phs_channel, success)
    class(component_instance_t), intent(inout), target :: component
    type(mci_work_t), intent(in) :: mci_work
    integer, intent(in) :: phs_channel
    logical, intent(out) :: success
    select type (phs => component%k_seed%phs)
    type is (phs_fks_t)
      call component%k_seed%compute_selected_channel &
           (mci_work, phs_channel, component%p_seed, success, &
            component%nlo_data)
      call interaction_set_momenta (component%nlo_data%int_born, &
                                    component%p_seed)
    class default
      call component%k_seed%compute_selected_channel &
         (mci_work, phs_channel, component%p_seed, success)
    end select
  end subroutine component_instance_compute_seed_kinematics
    
\eatline
\nwendcode{}\nwbegindocs{630}Inverse: recover missing parts of the kinematics, given a complete
set of seed momenta.  Select a channel and reconstruct the MC parameter set.
\nwenddocs{}\nwbegincode{631}\moddef{Processes: component instance: TBP}\plusendmoddef
  procedure :: recover_mcpar => component_instance_recover_mcpar
\nwendcode{}\nwbegincode{632}\moddef{Processes: procedures}\plusendmoddef
  subroutine component_instance_recover_mcpar (component, mci_work, phs_channel)
    class(component_instance_t), intent(inout), target :: component
    type(mci_work_t), intent(inout) :: mci_work
    integer, intent(in) :: phs_channel
    call component%k_seed%recover_mcpar &
         (mci_work, phs_channel, component%p_seed)
  end subroutine component_instance_recover_mcpar
  
\eatline
\nwendcode{}\nwbegindocs{633}Compute the momenta in the hard interactions, one for each term that
constitutes this process component.  In simple cases this amounts to
just copying momenta.  In more advanced cases, we may generate
distinct sets of momenta from the seed kinematics.

The interactions in the term instances are accessed individually.  We may
choose to calculate all terms at once together with the seed kinematics, use
{\tt{}component{\char37}tmp} for storage, and just fill the interactions here.
\nwenddocs{}\nwbegincode{634}\moddef{Processes: component instance: TBP}\plusendmoddef
  procedure :: compute_hard_kinematics => &
       component_instance_compute_hard_kinematics
\nwendcode{}\nwbegincode{635}\moddef{Processes: procedures}\plusendmoddef
  subroutine component_instance_compute_hard_kinematics &
       (component, term, skip_term, real_phsp)
    class(component_instance_t), intent(inout) :: component
    type(term_instance_t), dimension(:), intent(inout) :: term
    integer, intent(in), optional :: skip_term
    logical, intent(in), optional :: real_phsp
    integer :: j, i
    associate (core => component%config%core)
      associate (i_term => component%config%i_term)
        if (allocated (component%tmp)) then
           call component%tmp%reset_new_kinematics ()
        end if
        do j = 1, size (i_term)
           i = i_term(j)
           if (present (skip_term)) then
              if (i == skip_term)  cycle
           end if
           if (present (real_phsp)) then
             if (real_phsp) then
               call core%compute_hard_kinematics &
                     (component%p_seed, i, term(i)%nlo_data%int_born, component%tmp)
             else 
               call core%compute_hard_kinematics &
                     (component%p_seed, i, term(i)%int_hard, component%tmp)
             end if
          else
             call core%compute_hard_kinematics & 
                 (component%p_seed, i, term(i)%int_hard, component%tmp)
          end if
        end do
      end associate
    end associate
  end subroutine component_instance_compute_hard_kinematics
    
\eatline
\nwendcode{}\nwbegindocs{636}Here, we invert this.  We fetch the incoming momenta which reside
in the appropriate {\tt{}sf{\char95}chain} object, stored within the {\tt{}k{\char95}seed}
subobject.  On the other hand, we have the outgoing momenta of the
effective interaction.  We rely on the process core to compute the
remaining seed momenta and to fill the momenta within the hard
interaction.  (The latter is trivial if hard and effective interaction
coincide.)

After this is done, the incoming momenta in the trace evaluator that
corresponds to the hard (effective) interaction, are still
left undefined.  We remedy this by calling {\tt{}receive{\char95}kinematics} once.
\nwenddocs{}\nwbegincode{637}\moddef{Processes: component instance: TBP}\plusendmoddef
  procedure :: recover_seed_kinematics => &
       component_instance_recover_seed_kinematics
\nwendcode{}\nwbegincode{638}\moddef{Processes: procedures}\plusendmoddef
  subroutine component_instance_recover_seed_kinematics (component, term)
    class(component_instance_t), intent(inout) :: component
    type(term_instance_t), intent(inout) :: term
    integer :: n_in
    n_in = component%k_seed%n_in
    call component%k_seed%get_incoming_momenta (component%p_seed(1:n_in))
    associate (core => component%config%core)
      call core%recover_kinematics &
           (component%p_seed, term%int_hard, term%isolated%int_eff, &
           component%tmp)
      call term%isolated%receive_kinematics ()
    end associate
  end subroutine component_instance_recover_seed_kinematics
  
\eatline
\nwendcode{}\nwbegindocs{639}Compute the integration parameters for all channels except the selected
one.
\nwenddocs{}\nwbegincode{640}\moddef{Processes: component instance: TBP}\plusendmoddef
  procedure :: compute_other_channels => &
       component_instance_compute_other_channels
\nwendcode{}\nwbegincode{641}\moddef{Processes: procedures}\plusendmoddef
  subroutine component_instance_compute_other_channels &
       (component, mci_work, phs_channel)
    class(component_instance_t), intent(inout), target :: component
    type(mci_work_t), intent(in) :: mci_work
    integer, intent(in) :: phs_channel
    call component%k_seed%compute_other_channels (mci_work, phs_channel)
  end subroutine component_instance_compute_other_channels
    
\eatline
\nwendcode{}\nwbegindocs{642}Recover beam momenta, i.e., return the beam momenta as currently
stored in the kinematics subobject to their source.  This is a side effect.
\nwenddocs{}\nwbegincode{643}\moddef{Processes: component instance: TBP}\plusendmoddef
  procedure :: return_beam_momenta => component_instance_return_beam_momenta
\nwendcode{}\nwbegincode{644}\moddef{Processes: procedures}\plusendmoddef
  subroutine component_instance_return_beam_momenta (component)
    class(component_instance_t), intent(in) :: component
    call component%k_seed%return_beam_momenta ()
  end subroutine component_instance_return_beam_momenta
    
\eatline
\nwendcode{}\nwbegindocs{645}\nwdocspar
\nwenddocs{}\nwbegincode{646}\moddef{Processes: component instance: TBP}\plusendmoddef
  procedure :: evaluate_sqme_born => component_instance_evaluate_sqme_born
\nwendcode{}\nwbegincode{647}\moddef{Processes: procedures}\plusendmoddef
  subroutine component_instance_evaluate_sqme_born (component, term)
    class(component_instance_t), intent(inout) :: component
    type(term_instance_t), intent(inout), target :: term
    type(interaction_t), pointer :: int
    real(default) :: sqme
    int => evaluator_get_int_ptr (term%connected%trace)
    sqme = interaction_get_matrix_element (int, 1)
    component%sqme = sqme * term%weight
  end subroutine component_instance_evaluate_sqme_born    

\eatline
\nwendcode{}\nwbegindocs{648}Evaluate the trace of the transition matrix, convoluted with the initial
state, and summed over all terms.  The trace evaluators of the
individual terms have only a single matrix element. We implicitly drop the
imaginary part of the terms, which should be zero anyway.
\nwenddocs{}\nwbegincode{649}\moddef{Processes: component instance: TBP}\plusendmoddef
  procedure :: evaluate_sqme => component_instance_evaluate_sqme
\nwendcode{}\nwbegincode{650}\moddef{Processes: procedures}\plusendmoddef
  subroutine component_instance_evaluate_sqme (component, term)
    class(component_instance_t), intent(inout) :: component
    type(term_instance_t), dimension(:), intent(inout), target :: term
    type(interaction_t), pointer :: int
    type(vector4_t), dimension(:), allocatable :: p_real
    real(default) :: sqme
    real(default), dimension(4) :: sqme_virt
    integer :: j, i

    component%sqme = 0
    call reset_nlo_components (component)
    associate (i_term => component%config%i_term)
      do j = 1, size (i_term)
         i = i_term(j)
         if (term(i)%passed) then
            call component%evaluate_sqme_born (term(i))
            call add_born_to_nlo_component (component, &
                                            component%sqme*term(i)%weight)
            call component%evaluate_sqme_real (term(i))
            call component%evaluate_sqme_virt (term(i))
         end if
      end do
    end associate
    component%sqme_known = .true.
  end subroutine component_instance_evaluate_sqme
  
\eatline
\nwendcode{}\nwbegindocs{651}\nwdocspar
\nwenddocs{}\nwbegincode{652}\moddef{Processes: procedures}\plusendmoddef
  subroutine reset_nlo_components (component)
     type(component_instance_t), intent(inout) :: component
     select case (char (component%config%config%get_nlo_type ()))
     case ('Real', 'Virtual')
        component%nlo_data%sqme_born = 0
        component%nlo_data%sqme_real_non_sub = 0
     end select
  end subroutine reset_nlo_components

\eatline
\nwendcode{}\nwbegindocs{653}\nwdocspar
\nwenddocs{}\nwbegincode{654}\moddef{Processes: procedures}\plusendmoddef
  subroutine add_born_to_nlo_component (component, summand)
     type(component_instance_t), intent(inout) :: component
     real(default), intent(in) :: summand
     select case (char (component%config%config%get_nlo_type ()))
     case ('Real')
        component%nlo_data%sqme_born = summand
     end select
  end subroutine add_born_to_nlo_component

\nwendcode{}\nwbegindocs{655}Perform the matrix-element-evaluation if the component is flagged as virtual
\nwenddocs{}\nwbegincode{656}\moddef{Processes: component instance: TBP}\plusendmoddef
  procedure :: evaluate_sqme_virt => component_instance_evaluate_sqme_virt
\nwendcode{}\nwbegincode{657}\moddef{Processes: procedures}\plusendmoddef
  subroutine component_instance_evaluate_sqme_virt (component, term)
     class(component_instance_t), intent(inout) :: component
     type(term_instance_t), intent(inout) :: term
     real(default), dimension(4) :: sqme_virt 
     integer :: i_flv
     logical :: bad_point

     if (term%nlo_data%nlo_type /= 'Virtual') return
     call term%nlo_data%set_alpha_s_born &
               (component%config%core%get_alpha_s (component%tmp))
     call term%nlo_data%virtual_terms%set_ren_scale &
                (interaction_get_momenta (term%int_hard), term%ren_scale) 
     select type (core => component%config%core)
     type is (prc_gosam_t)
        do i_flv = 1, core%get_nflv()
           call core%compute_sqme_virt (i_flv, &
                      interaction_get_momenta (term%int_hard), &
                      term%ren_scale, &
                      core%get_alpha_s (component%tmp), &
                      sqme_virt, bad_point)
           if (.not. term%nlo_data%use_internal_color_correlations) &
              call core%compute_sqme_cc (i_flv, &
                        interaction_get_momenta (term%int_hard), &
                        term%ren_scale, &
                        core%get_alpha_s (component%tmp), &
                        born_cc = component%nlo_data%sqme_born_cc (:,:,i_flv), &
                        bad_point = bad_point)
           call term%nlo_data%virtual_terms%set_vfin (sqme_virt(3))
           call term%nlo_data%virtual_terms%set_bad_point (bad_point)
           component%nlo_data%sqme_born = sqme_virt (4)             
           component%sqme = component%sqme + &
                            term%nlo_data%compute_virt (i_flv, term%int_hard) * &
                            term%weight
        end do
     end select
  end subroutine component_instance_evaluate_sqme_virt

\eatline
\nwendcode{}\nwbegindocs{658}Perform the matrix-element-evaluation if the component is flagged as real
\nwenddocs{}\nwbegincode{659}\moddef{Processes: component instance: TBP}\plusendmoddef
  procedure :: evaluate_sqme_real => component_instance_evaluate_sqme_real
\nwendcode{}\nwbegincode{660}\moddef{Processes: procedures}\plusendmoddef
  subroutine component_instance_evaluate_sqme_real (component, term)
     class(component_instance_t), intent(inout) :: component
     type(term_instance_t), intent(inout), target :: term
     type(vector4_t), dimension(:), allocatable :: p_real
     type(interaction_t), pointer :: int
     integer :: i_flv
     logical :: bad_point

     if (term%nlo_data%nlo_type /= 'Real') return
     p_real = term%nlo_data%get_real_momenta ()
     select type (core => component%config%core)
     !!! Get the unsubtracted real matrix element, 
     !!! either from the parton states or from GoSam
     type is (prc_omega_t)
        int => evaluator_get_int_ptr (term%connected_real%trace)
        component%nlo_data%sqme_real_non_sub = interaction_get_matrix_element (int, 1)
     type is (prc_gosam_t)
        do i_flv = 1, component%nlo_data%n_flv_real
           call core%compute_sqme_real (i_flv, p_real, 0._default, &
                                        core%get_alpha_s (component%tmp), &
                                        component%nlo_data%sqme_real_non_sub(i_flv), &
                                        bad_point)
        end do
     end select 
     !!! Get the necessary subtraction matrix elements.
     select type (core_born => component%config%core_sub_born)
     type is (prc_gosam_t)
        do i_flv = 1, component%nlo_data%n_flv_born
           if (.not. term%nlo_data%use_internal_color_correlations) then
              call core_born%compute_sqme_cc (i_flv, &
                          interaction_get_momenta (component%nlo_data%int_born), &
                          0._default, core_born%get_alpha_s (component%tmp), &
                          component%nlo_data%sqme_born(i_flv), &
                          component%nlo_data%sqme_born_cc (:,:,i_flv), &
                          bad_point)
           else
              !!! Implementation for color-correlations using color_data
           end if
           if (.not. term%nlo_data%use_internal_spin_correlations  &
               .and. term%nlo_data%requires_spin_correlation (i_flv)) &
              !!! This far only correct for e+ e- -> q qbar gl
              call core_born%compute_sqme_sc (i_flv, 5, &
                             interaction_get_momenta (component%nlo_data%int_born), &
                             0._default, core_born%get_alpha_s (component%tmp), &
                             component%nlo_data%me_sc, &
                             bad_point)
        end do
     type is (prc_omega_t)
        component%nlo_data%sqme_born_cc (:,:,1) = component%nlo_data%sqme_born(1) * &
                                           component%nlo_data%color_data%beta_ij (:,:,1)
     end select
     call term%nlo_data%copy_matrix_elements ()
     call term%nlo_data%set_alpha_s_born &
                        (component%config%core%get_alpha_s (component%tmp))
     component%nlo_data%sqme_real = term%nlo_data%compute_sqme_real_fin &
                                      (term%weight, p_real) 
  end subroutine component_instance_evaluate_sqme_real

\eatline
\nwendcode{}\nwbegindocs{661}\nwdocspar
\nwenddocs{}\nwbegindocs{662}Get nlo-type and emitter of the component.
\nwenddocs{}\nwbegincode{663}\moddef{Processes: component instance: TBP}\plusendmoddef
  procedure :: get_nlo_type => component_instance_get_nlo_type
\nwendcode{}\nwbegincode{664}\moddef{Processes: procedures}\plusendmoddef
  elemental function component_instance_get_nlo_type (component) result (nlo_type)
    class(component_instance_t), intent(in) :: component
    type(string_t) :: nlo_type
    if (associated (component%nlo_data)) then
       if (component%nlo_data%get_nlo_type () == "") then
          nlo_type = 'Born'
       else
           nlo_type = component%nlo_data%get_nlo_type ()
       end if
    else
       nlo_type = 'Born'
    end if
  end function component_instance_get_nlo_type

\eatline
\nwendcode{}\nwbegindocs{665}\nwdocspar
\nwenddocs{}\nwbegincode{666}\moddef{Processes: component instance: TBP}\plusendmoddef
  procedure :: get_component_type => &
                   component_instance_get_component_type
\nwendcode{}\nwbegincode{667}\moddef{Processes: procedures}\plusendmoddef
  function component_instance_get_component_type (component) result (val)
    class(component_instance_t), intent(in) :: component
    integer :: val
    val = component%config%component_type
  end function component_instance_get_component_type

\eatline
\nwendcode{}\nwbegindocs{668}\nwdocspar
\nwenddocs{}\nwbegincode{669}\moddef{Processes: component instance: TBP}\plusendmoddef
  procedure :: set_component_type => &
                   component_instance_set_component_type
\nwendcode{}\nwbegincode{670}\moddef{Processes: procedures}\plusendmoddef
  subroutine component_instance_set_component_type (component, val)
    class(component_instance_t), intent(in) :: component
    integer, intent(in) :: val
    component%config%component_type = val
  end subroutine component_instance_set_component_type

\eatline
\nwendcode{}\nwbegindocs{671}\nwdocspar
\nwenddocs{}\nwbegincode{672}\moddef{Processes: process component: TBP}\plusendmoddef
  procedure :: needs_mci_entry => process_component_needs_mci_entry
\nwendcode{}\nwbegincode{673}\moddef{Processes: procedures}\plusendmoddef
  function process_component_needs_mci_entry (component) result (value)
    class(process_component_t), intent(in) :: component
    logical :: value
    value = component%active .and. component%core%needs_mcset () &
            .and. component%component_type <= COMP_MASTER 
  end function process_component_needs_mci_entry

\eatline
\nwendcode{}\nwbegindocs{674}\nwdocspar
\subsubsection{Term instance}
A {\tt{}term{\char95}instance{\char95}t} object contains all data that describe a term.  Each
process component consists of one or more distinct terms which may differ in
kinematics, but whose squared transition matrices have to be added pointwise.

The {\tt{}active} flag is set when this term is connected to an active
process component.  Inactive terms are skipped for kinematics and evaluation.

The {\tt{}k{\char95}term} object is the instance of the kinematics setup
(structure-function chain, phase space, etc.) that applies
specifically to this term.  In ordinary cases, it consists of straight
pointers to the seed kinematics.

The {\tt{}amp} array stores the amplitude values when we get them from evaluating
the associated matrix-element code.

The {\tt{}int{\char95}hard} interaction describes the elementary hard process.
It receives the momenta and the amplitude entries for each sampling point.

The {\tt{}isolated} object holds the effective parton state for the
elementary interaction.  The amplitude entries are
computed from {\tt{}int{\char95}hard}.  

The {\tt{}connected} evaluator set
convolutes this scattering matrix with the beam (and possibly
structure-function) density matrix.

The {\tt{}checked} flag is set once we have applied cuts on this term.
The result of this is stored in the {\tt{}passed} flag.  Once the term
has passed cuts, we calculate the various scale and weight expressions.
\nwenddocs{}\nwbegincode{675}\moddef{Processes: types}\plusendmoddef
  type :: term_instance_t
     type(process_term_t), pointer :: config => null ()
     logical :: active = .false.
     type(kinematics_t) :: k_term
     complex(default), dimension(:), allocatable :: amp
     type(interaction_t) :: int_hard
     type(interaction_t) :: int_hard_real
     type(isolated_state_t) :: isolated
     type(connected_state_t) :: connected
     type(isolated_state_t) :: isolated_real
     type(connected_state_t) :: connected_real
     logical :: checked = .false.
     logical :: passed = .false.
     real(default) :: scale = 0
     real(default) :: fac_scale = 0
     real(default) :: ren_scale = 0
     real(default), allocatable :: alpha_qcd_forced
     real(default) :: weight = 1
     type(vector4_t), dimension(:), allocatable :: p_hard
     type(nlo_data_t), pointer :: nlo_data => null ()
   contains
   \LA{}Processes: term instance: TBP\RA{}
  end type term_instance_t
  
\eatline
\nwendcode{}\nwbegindocs{676}\nwdocspar
\nwenddocs{}\nwbegincode{677}\moddef{Processes: term instance: TBP}\endmoddef
  procedure :: write => term_instance_write
\nwendcode{}\nwbegincode{678}\moddef{Processes: procedures}\plusendmoddef
  subroutine term_instance_write (term, unit, show_eff_state, testflag)
    class(term_instance_t), intent(in) :: term
    integer, intent(in), optional :: unit
    logical, intent(in), optional :: show_eff_state
    logical, intent(in), optional :: testflag 
    integer :: u
    logical :: state
    u = given_output_unit (unit)
    state = .true.;  if (present (show_eff_state))  state = show_eff_state
    if (term%active) then
       if (associated (term%config)) then
          write (u, "(1x,A,I0,A,I0,A)")  "Term #", term%config%i_term, &
               " (component #", term%config%i_component, ")"
       else
          write (u, "(1x,A)")  "Term [undefined]"
       end if
    else
       write (u, "(1x,A,I0,A)")  "Term #", term%config%i_term, &
            " [inactive]"
    end if
    if (term%checked) then
       write (u, "(3x,A,L1)")      "passed cuts           = ", term%passed
    end if
    if (term%passed) then
       write (u, "(3x,A,ES19.12)")  "overall scale         = ", term%scale
       write (u, "(3x,A,ES19.12)")  "factorization scale   = ", term%fac_scale
       write (u, "(3x,A,ES19.12)")  "renormalization scale = ", term%ren_scale
       if (allocated (term%alpha_qcd_forced)) then
          write (u, "(3x,A,ES19.12)")  "alpha(QCD) forced     = ", &
               term%alpha_qcd_forced
       end if
       write (u, "(3x,A,ES19.12)")  "reweighting factor    = ", term%weight
    end if
    call term%k_term%write (u)
    call write_separator (u)
    write (u, "(1x,A)")  "Amplitude (transition matrix of the &
         &hard interaction):"
    call write_separator (u)
    call interaction_write (term%int_hard, u, testflag = testflag)
    if (state .and. term%isolated%has_trace) then
       call write_separator (u)
       write (u, "(1x,A)")  "Evaluators for the hard interaction:"
       call term%isolated%write (u, testflag = testflag)
    end if
    if (state .and. term%connected%has_trace) then
       call write_separator (u)
       write (u, "(1x,A)")  "Evaluators for the connected process:"
       call term%connected%write (u, testflag = testflag)
    end if
  end subroutine term_instance_write
    
\eatline
\nwendcode{}\nwbegindocs{679}The interactions and evaluators must be finalized.
\nwenddocs{}\nwbegincode{680}\moddef{Processes: term instance: TBP}\plusendmoddef
  procedure :: final => term_instance_final
\nwendcode{}\nwbegincode{681}\moddef{Processes: procedures}\plusendmoddef
  subroutine term_instance_final (term)
    class(term_instance_t), intent(inout) :: term
    call term%k_term%final ()
    call term%connected%final ()
    call term%isolated%final ()
    call interaction_final (term%int_hard)
    select case (char (term%nlo_data%nlo_type))
    case ('Real')
      call term%connected_real%final ()
      call term%isolated_real%final ()
      call interaction_final (term%int_hard_real)
    end select
  end subroutine term_instance_final
  
\eatline
\nwendcode{}\nwbegindocs{682}For initialization, we make use of defined assignment for the
{\tt{}interaction{\char95}t} type.  This creates a deep copy.

The hard interaction (incoming momenta) is linked to the structure
function instance.  In the isolated state, we either set pointers to
both, or we create modified copies ({\tt{}rearrange}) as effective
structure-function chain and interaction, respectively.

Finally, we set up the {\tt{}subevt} component that will be used for
evaluating observables, collecting particles from the trace evaluator
in the effective connected state.  Their quantum numbers must be
determined by following back source links and set explicitly, since
they are already eliminated in that trace.

The {\tt{}rearrange} parts are still commented out; they could become
relevant for a NLO algorithm.
\nwenddocs{}\nwbegincode{683}\moddef{Processes: term instance: TBP}\plusendmoddef
  procedure :: init => term_instance_init
\nwendcode{}\nwbegincode{684}\moddef{Processes: procedures}\plusendmoddef
  subroutine term_instance_init (term, &
       config, k_seed, beam_config, core, process_var_list, nlo_data)
    class(term_instance_t), intent(out), target :: term
    type(process_term_t), intent(in), target :: config
    type(kinematics_t), intent(in) :: k_seed
    type(process_beam_config_t), intent(in) :: beam_config
    type(interaction_t), pointer :: sf_chain_int, trace_int, src_int
    class(prc_core_t), intent(in) :: core
    type(var_list_t), intent(in), target :: process_var_list
    type(nlo_data_t), intent(in), pointer :: nlo_data
    type(quantum_numbers_mask_t), dimension(:), allocatable :: mask_in
    type(state_matrix_t), pointer :: state_matrix
    type(flavor_t), dimension(:), allocatable :: flv_int, flv_src, f_in, f_out
    type(flavor_t), dimension(:), allocatable :: f_out_real
    integer :: n_in, n_vir, n_out, n_tot
    integer :: i, j
    type(interaction_t), pointer, save :: int_sav 
    type(sf_chain_instance_t), pointer, save :: sf_sav
    integer, dimension(:), allocatable, save :: col_sav
    type(prc_omega_t), save :: core_sav

    term%nlo_data => nlo_data

    term%config => config
    if (config%rearrange) then
       ! rearrangement of seed to hard kinematics not implemented yet
       ! allocate k_term distinct from k_seed as needed.
    else
       ! here, k_term trivially accesses k_seed via pointers
       call term%k_term%init_ptr (k_seed)
    end if
    allocate (term%amp (config%n_allowed))
    term%int_hard = config%int
    allocate (term%p_hard (interaction_get_n_tot (term%int_hard)))
    sf_chain_int => term%k_term%sf_chain%get_out_int_ptr ()
    n_in = interaction_get_n_in (term%int_hard)
    do j = 1, n_in
       i = term%k_term%sf_chain%get_out_i (j)
       call interaction_set_source_link (term%int_hard, j, sf_chain_int, i)
    end do
    if (config%rearrange) then
       ! rearrangement hard to effective kinematics not implemented yet
       ! should use term%config%int_eff as template
       ! allocate distinct sf_chain in term%connected as needed
    else
      select type (phs => term%k_term%phs)
        type is (phs_wood_t)
          ! here, int_hard and sf_chain are trivially accessed via pointers
          call term%isolated%init (term%k_term%sf_chain, term%int_hard)
          int_sav => term%int_hard
          sf_sav => term%k_term%sf_chain
        type is (phs_fks_t)
          term%nlo_data%int_born = int_sav
          call term%isolated%init (sf_sav, term%nlo_data%int_born)
          call term%setup_real_interaction (term%int_hard)
        class default
          call term%isolated%init (term%k_term%sf_chain, term%int_hard)
      end select
    end if
    allocate (mask_in (n_in))
    mask_in = term%k_term%sf_chain%get_out_mask ()
    select type (phs => term%k_term%phs)
      type is (phs_wood_t)
      call term%isolated%setup_square_trace (core, mask_in, term%config%col)
      if (.not. allocated (col_sav)) then
        allocate (col_sav (size (term%config%col)))
        col_sav = term%config%col
      end if
      select type (core)
      type is (prc_omega_t)
        core_sav = core
      end select
      type is (phs_fks_t)
        mask_in = sf_sav%get_out_mask ()
        call term%isolated%setup_square_trace (core_sav, mask_in, col_sav)
        call term%init_states (core)
      class default
        call term%isolated%setup_square_trace (core, mask_in, term%config%col)
    end select
    call term%connected%setup_connected_trace (term%isolated)
    associate (int_eff => term%isolated%int_eff)
      state_matrix => interaction_get_state_matrix_ptr (int_eff)
      n_tot = interaction_get_n_tot  (int_eff)
      allocate (flv_int (n_tot))
      flv_int = quantum_numbers_get_flavor &
           (state_matrix%get_quantum_numbers (1))
      allocate (f_in (n_in))
      f_in = flv_int(1:n_in)
      deallocate (flv_int)
    end associate
    trace_int => evaluator_get_int_ptr (term%connected%trace)
    n_in = interaction_get_n_in (trace_int)
    n_vir = interaction_get_n_vir (trace_int)
    n_out = interaction_get_n_out (trace_int)
    allocate (f_out (n_out))
    do j = 1, n_out
       call interaction_find_source (trace_int, n_in + n_vir + j, src_int, i)
       if (associated (src_int)) then
          state_matrix => interaction_get_state_matrix_ptr (src_int)
          allocate (flv_src (interaction_get_n_tot (src_int)))
          flv_src = quantum_numbers_get_flavor &
               (state_matrix%get_quantum_numbers (1))
          f_out(j) = flv_src(i)
          deallocate (flv_src)
       end if
    end do
    call term%connected%setup_subevt (term%isolated%sf_chain_eff, &
         beam_config%data%flv, f_in, f_out)
    call term%connected%setup_var_list (process_var_list, beam_config%data)
    select case (char (term%nlo_data%get_nlo_type ()))
    case ('Real')
      allocate (f_out_real (n_out+1))
      f_out_real (1:n_out) = f_out
      f_out_real (n_out+1) = term%nlo_data%reg_data%flv_extra
      call term%connected_real%setup_subevt (term%isolated_real%sf_chain_eff, &
           beam_config%data%flv, f_in, f_out_real)
      call term%connected_real%setup_var_list (process_var_list, beam_config%data)
    end select

  end subroutine term_instance_init

\eatline
\nwendcode{}\nwbegindocs{685}Set up the state objects.
\nwenddocs{}\nwbegincode{686}\moddef{Processes: term instance: TBP}\plusendmoddef
  procedure :: init_states => term_instance_init_states
\nwendcode{}\nwbegincode{687}\moddef{Processes: procedures}\plusendmoddef
  subroutine term_instance_init_states (term, core)
    class(term_instance_t), intent(inout), target :: term
    class(prc_core_t), intent(in) :: core
    type(quantum_numbers_mask_t), dimension(:), allocatable :: mask_in
    integer :: i
    mask_in = term%k_term%sf_chain%get_out_mask ()
    call term%isolated_real%init (term%k_term%sf_chain, term%int_hard_real)
    call term%isolated_real%setup_square_trace (core, mask_in, term%config%col)
    call term%connected_real%setup_connected_trace (term%isolated_real) 
  end subroutine term_instance_init_states

\eatline
\nwendcode{}\nwbegindocs{688}For initializing the expressions, we need the local variable list and the
parse trees.
\nwenddocs{}\nwbegincode{689}\moddef{Processes: term instance: TBP}\plusendmoddef
  procedure :: setup_expressions => term_instance_setup_expressions
\nwendcode{}\nwbegincode{690}\moddef{Processes: procedures}\plusendmoddef
  subroutine term_instance_setup_expressions (term, meta, config)
    class(term_instance_t), intent(inout), target :: term
    type(process_metadata_t), intent(in), target :: meta
    type(process_config_data_t), intent(in) :: config
    if (allocated (config%ef_cuts)) &
         call term%connected%setup_cuts (config%ef_cuts)
    if (allocated (config%ef_scale)) &
         call term%connected%setup_scale (config%ef_scale)
    if (allocated (config%ef_fac_scale)) &
         call term%connected%setup_fac_scale (config%ef_fac_scale)
    if (allocated (config%ef_ren_scale)) &
         call term%connected%setup_ren_scale (config%ef_ren_scale)
    if (allocated (config%ef_weight)) &
         call term%connected%setup_weight (config%ef_weight)
  end subroutine term_instance_setup_expressions
    
\eatline
\nwendcode{}\nwbegindocs{691}\nwdocspar
\nwenddocs{}\nwbegincode{692}\moddef{Processes: term instance: TBP}\plusendmoddef
  procedure :: setup_expressions_real => term_instance_setup_expressions_real
\nwendcode{}\nwbegincode{693}\moddef{Processes: procedures}\plusendmoddef
  subroutine term_instance_setup_expressions_real (term, meta, config)
    class(term_instance_t), intent(inout), target :: term
    type(process_metadata_t), intent(in), target :: meta
    type(process_config_data_t), intent(in) :: config
    if (allocated (config%ef_cuts)) &
         call term%connected_real%setup_cuts (config%ef_cuts)
    if (allocated (config%ef_scale)) &
         call term%connected_real%setup_scale (config%ef_scale)
    if (allocated (config%ef_fac_scale)) &
         call term%connected_real%setup_fac_scale (config%ef_fac_scale)
    if (allocated (config%ef_ren_scale)) &
         call term%connected_real%setup_ren_scale (config%ef_ren_scale)
    if (allocated (config%ef_weight)) &
         call term%connected_real%setup_weight (config%ef_weight)
  end subroutine term_instance_setup_expressions_real

\eatline
\nwendcode{}\nwbegindocs{694}Prepare the extra evaluators that we need for processing events.

The quantum numbers mask of the incoming particle
\nwenddocs{}\nwbegincode{695}\moddef{Processes: term instance: TBP}\plusendmoddef
  procedure :: setup_event_data => term_instance_setup_event_data
\nwendcode{}\nwbegincode{696}\moddef{Processes: procedures}\plusendmoddef
  subroutine term_instance_setup_event_data (term, core, model)
    class(term_instance_t), intent(inout), target :: term
    class(prc_core_t), intent(in) :: core
    class(model_data_t), intent(in), target :: model
    integer :: n_in
    type(quantum_numbers_mask_t), dimension(:), allocatable :: mask_in
    n_in = interaction_get_n_in (term%int_hard)
    allocate (mask_in (n_in))
    mask_in = term%k_term%sf_chain%get_out_mask ()
    select case (char (term%nlo_data%nlo_type))
    case ('Born', 'Virtual')
      call term%isolated%setup_square_matrix (core, model, mask_in, &
         term%config%col)
      call term%isolated%setup_square_flows (core, model, mask_in)
      call term%connected%setup_connected_matrix (term%isolated)
      call term%connected%setup_connected_flows (term%isolated)
      call term%connected%setup_state_flv (term%isolated%get_n_out ())
     case ('Real')
      call term%isolated_real%setup_square_matrix (core, model, mask_in, &
               term%config%col)
      call term%isolated_real%setup_square_flows (core, model, mask_in)
      call term%connected_real%setup_connected_matrix &
                (term%isolated_real)
      call term%connected_real%setup_connected_flows &
                (term%isolated_real)
      call term%connected_real%setup_state_flv (term%isolated_real%get_n_out ())
    end select
 end subroutine term_instance_setup_event_data
    
\eatline
\nwendcode{}\nwbegindocs{697}Associate the hard real interaction.
\nwenddocs{}\nwbegincode{698}\moddef{Processes: term instance: TBP}\plusendmoddef
  procedure :: setup_real_interaction => term_instance_setup_real_interaction
\nwendcode{}\nwbegincode{699}\moddef{Processes: procedures}\plusendmoddef
  subroutine term_instance_setup_real_interaction (term, int)
    class(term_instance_t), intent(inout) :: term
    type(interaction_t), intent(in) :: int
    integer :: i
    term%int_hard_real = int
  end subroutine term_instance_setup_real_interaction

\eatline
\nwendcode{}\nwbegindocs{700}Reset the term instance: clear the parton-state expressions and deactivate.
\nwenddocs{}\nwbegincode{701}\moddef{Processes: term instance: TBP}\plusendmoddef
  procedure :: reset => term_instance_reset
\nwendcode{}\nwbegincode{702}\moddef{Processes: procedures}\plusendmoddef
  subroutine term_instance_reset (term)
    class(term_instance_t), intent(inout) :: term
    call term%connected%reset_expressions ()
    if (allocated (term%alpha_qcd_forced))  deallocate (term%alpha_qcd_forced)
    term%active = .false.
  end subroutine term_instance_reset
  
\eatline
\nwendcode{}\nwbegindocs{703}Force an $\alpha_s$ value that should be used in the matrix-element
calculation. 
\nwenddocs{}\nwbegincode{704}\moddef{Processes: term instance: TBP}\plusendmoddef
  procedure :: set_alpha_qcd_forced => term_instance_set_alpha_qcd_forced
\nwendcode{}\nwbegincode{705}\moddef{Processes: procedures}\plusendmoddef
  subroutine term_instance_set_alpha_qcd_forced (term, alpha_qcd)
    class(term_instance_t), intent(inout) :: term
    real(default), intent(in) :: alpha_qcd
    if (allocated (term%alpha_qcd_forced)) then
       term%alpha_qcd_forced = alpha_qcd
    else
       allocate (term%alpha_qcd_forced, source = alpha_qcd)
    end if
  end subroutine term_instance_set_alpha_qcd_forced
  
\eatline
\nwendcode{}\nwbegindocs{706}Complete the kinematics computation for the effective parton states.

We assume that the {\tt{}compute{\char95}hard{\char95}kinematics} method of the process
component instance has already been called, so the {\tt{}int{\char95}hard}
contains the correct hard kinematics.  The duty of this procedure is
first to compute the effective kinematics and store this in the
{\tt{}int{\char95}eff} effective interaction inside the {\tt{}isolated} parton
state.  The effective kinematics may differ from the kinematics in the hard
interaction.  It may involve parton recombination or parton splitting.
The {\tt{}rearrange{\char95}partons} method is responsible for this part.  

We may also call a method to compute the effective structure-function
chain at this point.  This is not implemented yet.

In the simple case that no rearrangement is necessary, as indicated by
the {\tt{}rearrange} flag, the effective interaction is a pointer to the
hard interaction, and we can skip the rearrangement method.  Similarly
for the effective structure-function chain.  (If we have an algorithm
that uses rarrangement, it should evaluate {\tt{}k{\char95}term} explicitly.)

The final step of kinematics setup is to transfer the effective
kinematics to the evaluators and to the {\tt{}subevt}.
\nwenddocs{}\nwbegincode{707}\moddef{Processes: term instance: TBP}\plusendmoddef
  procedure :: compute_eff_kinematics => &
       term_instance_compute_eff_kinematics
\nwendcode{}\nwbegincode{708}\moddef{Processes: procedures}\plusendmoddef
  subroutine term_instance_compute_eff_kinematics (term, component)
    class(term_instance_t), intent(inout) :: term
    type(component_instance_t), dimension(:), intent(inout) :: component
    integer :: i_component, i_term
    term%checked = .false.
    term%passed = .false.
    if (term%config%rearrange) then
       ! should evaluate k_term first if allocated separately, not impl. yet
       i_component = term%config%i_component
       i_term = term%config%i_term
      associate (core => component(i_component)%config%core)
         select case (char (term%nlo_data%nlo_type))
         case ('Born', 'Virtual')
           call core%compute_eff_kinematics &
              (i_term, term%int_hard, term%isolated%int_eff, &
              component(i_component)%tmp)
         case ('Real')
           call component(i_component)%config%core_sub_born%compute_eff_kinematics & 
                (i_term, term%nlo_data%int_born, term%isolated%int_eff, &
                component(i_component)%tmp)
         end select
        select case (char (term%nlo_data%nlo_type))
        case ('Real')
          call core%compute_eff_kinematics &
                (i_term, term%int_hard_real, term%isolated_real%int_eff, &
                 component(i_component)%tmp)
        end select
      end associate
    end if
    call term%isolated%receive_kinematics ()
    call term%connected%receive_kinematics ()
  end subroutine term_instance_compute_eff_kinematics
    
\eatline
\nwendcode{}\nwbegindocs{709}Inverse.  Reconstruct the connected state from the momenta in the
trace evaluator (which we assume to be set), then reconstruct the
isolated state as far as possible.  The second part finalizes the
momentum configuration, using the incoming seed momenta
\nwenddocs{}\nwbegincode{710}\moddef{Processes: term instance: TBP}\plusendmoddef
  procedure :: recover_hard_kinematics => &
       term_instance_recover_hard_kinematics
\nwendcode{}\nwbegincode{711}\moddef{Processes: procedures}\plusendmoddef
  subroutine term_instance_recover_hard_kinematics (term, component)
    class(term_instance_t), intent(inout) :: term
    type(component_instance_t), dimension(:), intent(inout) :: component
    term%checked = .false.
    term%passed = .false.
    call term%connected%send_kinematics ()
    call term%isolated%send_kinematics ()
  end subroutine term_instance_recover_hard_kinematics

\eatline
\nwendcode{}\nwbegindocs{712}Check the term whether it passes cuts and, if successful, evaluate
scales and weights.  The factorization scale is also given to the term
kinematics, enabling structure-function evaluation.
\nwenddocs{}\nwbegincode{713}\moddef{Processes: term instance: TBP}\plusendmoddef
  procedure :: evaluate_expressions => &
       term_instance_evaluate_expressions
\nwendcode{}\nwbegincode{714}\moddef{Processes: procedures}\plusendmoddef
  subroutine term_instance_evaluate_expressions (term, scale_forced)
    class(term_instance_t), intent(inout) :: term
    real(default), intent(in), allocatable, optional :: scale_forced
    call term%connected%evaluate_expressions (term%passed, &
         term%scale, term%fac_scale, term%ren_scale, term%weight, &
         scale_forced)
    term%checked = .true.
  end subroutine term_instance_evaluate_expressions
       
\eatline
\nwendcode{}\nwbegindocs{715}Evaluate the trace: first evaluate the hard interaction, then the trace
evaluator.  We use the {\tt{}evaluate{\char95}interaction} method of the process
component which generated this term.  The {\tt{}subevt} and cut expressions are
not yet filled.

The {\tt{}component} argument is intent(inout) because the {\tt{}compute{\char95}amplitude}
method may modify the {\tt{}tmp} workspace object.
\nwenddocs{}\nwbegincode{716}\moddef{Processes: term instance: TBP}\plusendmoddef
  procedure :: evaluate_interaction => term_instance_evaluate_interaction
\nwendcode{}\nwbegincode{717}\moddef{Processes: procedures}\plusendmoddef
  subroutine term_instance_evaluate_interaction (term, component)
    class(term_instance_t), intent(inout) :: term
    type(component_instance_t), dimension(:), intent(inout) :: component
    integer :: i_component, i_term, i
    i_component = term%config%i_component
    i_term = term%config%i_term
    term%p_hard = interaction_get_momenta (term%int_hard)
    select case (char (term%nlo_data%nlo_type))
    case ('Born', 'Virtual')
       associate (core => component(i_component)%config%core)
          do i = 1, term%config%n_allowed
             term%amp(i) = core%compute_amplitude (i_term, term%p_hard, &
                term%config%flv(i), term%config%hel(i), term%config%col(i), &
                term%fac_scale, term%ren_scale, term%alpha_qcd_forced, &
                component(i_component)%tmp)
          end do
          call interaction_set_matrix_element (term%int_hard, term%amp)
       end associate 
    case ('Real')
       call term%evaluate_interaction_real (component(i_component), i_term)
    end select  
  end subroutine term_instance_evaluate_interaction
  
\eatline
\nwendcode{}\nwbegindocs{718}Evaluate the trace.  First evaluate the
structure-function chain (i.e., the density matrix of the incoming
partons).  Do this twice, in case the sf-chain instances within
{\tt{}k{\char95}term} and {\tt{}isolated} differ.  Next, evaluate the hard
interaction, then compute the convolution with the initial state.
\nwenddocs{}\nwbegincode{719}\moddef{Processes: term instance: TBP}\plusendmoddef
  procedure :: evaluate_trace => term_instance_evaluate_trace
\nwendcode{}\nwbegincode{720}\moddef{Processes: procedures}\plusendmoddef
  subroutine term_instance_evaluate_trace (term)
    class(term_instance_t), intent(inout) :: term
    call term%k_term%evaluate_sf_chain (term%fac_scale)
    call term%isolated%evaluate_sf_chain (term%fac_scale)
    call term%isolated%evaluate_trace ()
    call term%connected%evaluate_trace ()
    select case (char (term%nlo_data%nlo_type))
    case ('Real')
      call term%isolated_real%evaluate_trace ()
      call term%connected_real%evaluate_trace ()
    end select
  end subroutine term_instance_evaluate_trace
  
\eatline
\nwendcode{}\nwbegindocs{721}Evaluate the extra data that we need for processing the object as a
physical event.
\nwenddocs{}\nwbegincode{722}\moddef{Processes: term instance: TBP}\plusendmoddef
  procedure :: evaluate_event_data => term_instance_evaluate_event_data
\nwendcode{}\nwbegincode{723}\moddef{Processes: procedures}\plusendmoddef
  subroutine term_instance_evaluate_event_data (term)
    class(term_instance_t), intent(inout) :: term
    select case (char (term%nlo_data%nlo_type))
    case ('Born', 'Virtual')
      call term%isolated%evaluate_event_data ()
      call term%connected%evaluate_event_data ()
    case ('Real')
      call term%isolated_real%evaluate_event_data ()
      call term%connected_real%evaluate_event_data ()
    end select
  end subroutine term_instance_evaluate_event_data
  
\eatline
\nwendcode{}\nwbegindocs{724}Evaluate interaction for an $N+1$-particle phase space
\nwenddocs{}\nwbegincode{725}\moddef{Processes: term instance: TBP}\plusendmoddef
  procedure :: evaluate_interaction_real &
               => term_instance_evaluate_interaction_real
\nwendcode{}\nwbegincode{726}\moddef{Processes: procedures}\plusendmoddef
  subroutine term_instance_evaluate_interaction_real (term, component, i_term)
     class(term_instance_t), intent(inout) :: term
     type(component_instance_t), intent(inout) :: component
     integer, intent(in) :: i_term
     integer :: emitter
     type(vector4_t), dimension(:), allocatable :: p_born
     type(vector4_t), dimension(:), allocatable :: p_real
     real(default) :: xi, xi_tilde, xi_max
     real(default) :: y, phi
     integer :: i

     emitter = term%nlo_data%get_active_emitter ()
     p_born = interaction_get_momenta (term%nlo_data%int_born)
     if (emitter > 2) then
        select type (phs => term%k_term%phs)
        type is (phs_fks_t)
           call phs%generate_fsr (emitter, p_born, p_real)
           call term%nlo_data%set_jacobian (phs%jac)
           call term%nlo_data%set_y_soft (phs%y_soft, emitter)
           call term%nlo_data%set_real_momenta (p_real)
        end select
     else
        call msg_fatal ("FKS ISR not implemented yet!")
     end if
     call interaction_set_momenta (term%int_hard_real, p_real)

     call term%isolated_real%receive_kinematics ()
     call term%connected_real%receive_kinematics ()

     call term%connected_real%evaluate_expressions (term%passed, &
          term%scale, term%fac_scale, term%ren_scale, term%weight)

     if (term%passed) then
        call component%tmp%reset_new_kinematics ()
        select type (core => component%config%core)
        type is (prc_omega_t)
           do i = 1, term%config%n_allowed
              term%amp(i) = core%compute_amplitude (i_term, &
                 interaction_get_momenta (term%int_hard_real), &
                 term%config%flv(i), term%config%hel(i), term%config%col(i), &
                 term%fac_scale, term%ren_scale, term%alpha_qcd_forced, &
                 component%tmp)
           end do
        type is (prc_gosam_t)
           call core%update_alpha_s (component%tmp, term%fac_scale)
        end select
        call interaction_set_matrix_element (term%int_hard_real, term%amp)
        call component%tmp%reset_new_kinematics ()
        select type (core_sub_born => component%config%core_sub_born)
        type is (prc_omega_t)
           do i = 1, term%nlo_data%n_allowed_born
               term%nlo_data%amp_born(i) = core_sub_born%compute_amplitude (i_term, &
                  interaction_get_momenta (term%nlo_data%int_born), &
                  term%nlo_data%flv_born(i), term%nlo_data%hel_born(i), &
                  term%nlo_data%col_born(i), term%fac_scale, term%ren_scale, &
                  term%alpha_qcd_forced, &
                  component%tmp)
           end do
        type is (prc_gosam_t)
           call core_sub_born%update_alpha_s (component%tmp, term%fac_scale)
        end select 
        call interaction_set_matrix_element (term%nlo_data%int_born, term%nlo_data%amp_born) 
     end if
  end subroutine term_instance_evaluate_interaction_real

\eatline
\nwendcode{}\nwbegindocs{727}Return data that might be useful for external processing.  The
factorization scale:
\nwenddocs{}\nwbegincode{728}\moddef{Processes: term instance: TBP}\plusendmoddef
  procedure :: get_fac_scale => term_instance_get_fac_scale
\nwendcode{}\nwbegincode{729}\moddef{Processes: procedures}\plusendmoddef
  function term_instance_get_fac_scale (term) result (fac_scale)
    class(term_instance_t), intent(in) :: term
    real(default) :: fac_scale
    fac_scale = term%fac_scale
  end function term_instance_get_fac_scale
  
\eatline
\nwendcode{}\nwbegindocs{730}We take the strong coupling from the process core.  The value is calculated
when a new event is requested, so we should call it only after the event has
been evaluated.  If it is not available there (a negative number is returned),
we take the value stored in the term configuration, which should be determined
by the model.  If the model does not provide a value, the result is zero.
\nwenddocs{}\nwbegincode{731}\moddef{Processes: term instance: TBP}\plusendmoddef
  procedure :: get_alpha_s => term_instance_get_alpha_s
\nwendcode{}\nwbegincode{732}\moddef{Processes: procedures}\plusendmoddef
  function term_instance_get_alpha_s (term, component) result (alpha_s)
    class(term_instance_t), intent(in) :: term
    type(component_instance_t), dimension(:), intent(in) :: component
    real(default) :: alpha_s
    integer :: i_component
    i_component = term%config%i_component
    associate (core => component(i_component)%config%core)
      alpha_s = core%get_alpha_s (component(i_component)%tmp)
    end associate
    if (alpha_s < 0)  alpha_s = term%config%alpha_s
  end function term_instance_get_alpha_s
  
\eatline
\nwendcode{}\nwbegindocs{733}\nwdocspar
\subsubsection{MC parameter set and MCI instance}
For each process component that is associated with a multi-channel integration
(MCI) object, the {\tt{}mci{\char95}work{\char95}t} object contains the currently active
parameter set.  It also holds the implementation of the {\tt{}mci{\char95}instance{\char95}t}
that the integrator needs for doing its work.
\nwenddocs{}\nwbegincode{734}\moddef{Processes: types}\plusendmoddef
  type :: mci_work_t
     type(process_mci_entry_t), pointer :: config => null ()
     real(default), dimension(:), allocatable :: x
     class(mci_instance_t), pointer :: mci => null ()
     type(process_counter_t) :: counter
   contains
   \LA{}Processes: mci work: TBP\RA{}
  end type mci_work_t

\eatline
\nwendcode{}\nwbegindocs{735}First write configuration data, then the current values.
\nwenddocs{}\nwbegincode{736}\moddef{Processes: mci work: TBP}\endmoddef
  procedure :: write => mci_work_write
\nwendcode{}\nwbegincode{737}\moddef{Processes: procedures}\plusendmoddef
  subroutine mci_work_write (mci_work, unit, testflag)
    class(mci_work_t), intent(in) :: mci_work
    integer, intent(in), optional :: unit
    logical, intent(in), optional :: testflag
    integer :: u, i
    u = given_output_unit (unit)
    write (u, "(1x,A,I0,A)")  "Active MCI instance #", &
         mci_work%config%i_mci, " ="
    write (u, "(2x)", advance="no")
    do i = 1, mci_work%config%n_par
       write (u, "(1x,F7.5)", advance="no")  mci_work%x(i)
       if (i == mci_work%config%n_par_sf) &
            write (u, "(1x,'|')", advance="no")
    end do
    write (u, *)
    if (associated (mci_work%mci)) then
       call mci_work%mci%write (u, pacify = testflag)
       call mci_work%counter%write (u)
    end if
  end subroutine mci_work_write
         
\eatline
\nwendcode{}\nwbegindocs{738}The {\tt{}mci} component may require finalization.
\nwenddocs{}\nwbegincode{739}\moddef{Processes: mci work: TBP}\plusendmoddef
  procedure :: final => mci_work_final
\nwendcode{}\nwbegincode{740}\moddef{Processes: procedures}\plusendmoddef
  subroutine mci_work_final (mci_work)
    class(mci_work_t), intent(inout) :: mci_work
    if (associated (mci_work%mci)) then
       call mci_work%mci%final ()
       deallocate (mci_work%mci)
    end if
  end subroutine mci_work_final
  
\eatline
\nwendcode{}\nwbegindocs{741}Initialize with the maximum length that we will need.  Contents are
not initialized.

The integrator inside the {\tt{}mci{\char95}entry} object is responsible for
allocating and initializing its own instance, which is referred to by
a pointer in the {\tt{}mci{\char95}work} object.
\nwenddocs{}\nwbegincode{742}\moddef{Processes: mci work: TBP}\plusendmoddef
  procedure :: init => mci_work_init
\nwendcode{}\nwbegincode{743}\moddef{Processes: procedures}\plusendmoddef
  subroutine mci_work_init (mci_work, mci_entry)
    class(mci_work_t), intent(out) :: mci_work
    type(process_mci_entry_t), intent(in), target :: mci_entry
    mci_work%config => mci_entry
    allocate (mci_work%x (mci_entry%n_par))
    if (allocated (mci_entry%mci)) then
       call mci_entry%mci%allocate_instance (mci_work%mci)
       call mci_work%mci%init (mci_entry%mci)
    end if
  end subroutine mci_work_init
  
\eatline
\nwendcode{}\nwbegindocs{744}Set parameters explicitly, either all at once, or separately for the
structure-function and process parts.
\nwenddocs{}\nwbegincode{745}\moddef{Processes: mci work: TBP}\plusendmoddef
  procedure :: set => mci_work_set
  procedure :: set_x_strfun => mci_work_set_x_strfun
  procedure :: set_x_process => mci_work_set_x_process
\nwendcode{}\nwbegincode{746}\moddef{Processes: procedures}\plusendmoddef
  subroutine mci_work_set (mci_work, x)
    class(mci_work_t), intent(inout) :: mci_work
    real(default), dimension(:), intent(in) :: x
    mci_work%x = x
  end subroutine mci_work_set
    
  subroutine mci_work_set_x_strfun (mci_work, x)
    class(mci_work_t), intent(inout) :: mci_work
    real(default), dimension(:), intent(in) :: x
    mci_work%x(1 : mci_work%config%n_par_sf) = x
  end subroutine mci_work_set_x_strfun
    
  subroutine mci_work_set_x_process (mci_work, x)
    class(mci_work_t), intent(inout) :: mci_work
    real(default), dimension(:), intent(in) :: x
    mci_work%x(mci_work%config%n_par_sf + 1 : mci_work%config%n_par) = x
  end subroutine mci_work_set_x_process
    
\eatline
\eatline
\eatline
\nwendcode{}\nwbegindocs{747}Return the array of active components, i.e., those that correspond
to the currently selected MC parameter set.
\nwenddocs{}\nwbegincode{748}\moddef{Processes: mci work: TBP}\plusendmoddef
  procedure :: get_active_components => mci_work_get_active_components
\nwendcode{}\nwbegincode{749}\moddef{Processes: procedures}\plusendmoddef
  function mci_work_get_active_components (mci_work) result (i_component)
    class(mci_work_t), intent(in) :: mci_work
    integer, dimension(:), allocatable :: i_component
    allocate (i_component (size (mci_work%config%i_component)))
    i_component = mci_work%config%i_component
  end function mci_work_get_active_components

\eatline
\nwendcode{}\nwbegindocs{750}Return the active parameters as a simple array with correct length.
Do this separately for the structure-function parameters and the
process parameters.
\nwenddocs{}\nwbegincode{751}\moddef{Processes: mci work: TBP}\plusendmoddef
  procedure :: get_x_strfun => mci_work_get_x_strfun
  procedure :: get_x_process => mci_work_get_x_process
\nwendcode{}\nwbegincode{752}\moddef{Processes: procedures}\plusendmoddef
  function mci_work_get_x_strfun (mci_work) result (x)
    class(mci_work_t), intent(in) :: mci_work
    real(default), dimension(mci_work%config%n_par_sf) :: x
    x = mci_work%x(1 : mci_work%config%n_par_sf)
  end function mci_work_get_x_strfun

  function mci_work_get_x_process (mci_work) result (x)
    class(mci_work_t), intent(in) :: mci_work
    real(default), dimension(mci_work%config%n_par_phs) :: x
    x = mci_work%x(mci_work%config%n_par_sf + 1 : mci_work%config%n_par)
  end function mci_work_get_x_process

\eatline
\eatline
\nwendcode{}\nwbegindocs{753}Initialize and finalize event generation for the specified MCI
entry.  This also resets the counter.
\nwenddocs{}\nwbegincode{754}\moddef{Processes: mci work: TBP}\plusendmoddef
  procedure :: init_simulation => mci_work_init_simulation
  procedure :: final_simulation => mci_work_final_simulation
\nwendcode{}\nwbegincode{755}\moddef{Processes: procedures}\plusendmoddef
  subroutine mci_work_init_simulation (mci_work, safety_factor)
    class(mci_work_t), intent(inout) :: mci_work
    real(default), intent(in), optional :: safety_factor
    call mci_work%mci%init_simulation (safety_factor)
    call mci_work%counter%reset ()
  end subroutine mci_work_init_simulation

  subroutine mci_work_final_simulation (mci_work)
    class(mci_work_t), intent(inout) :: mci_work
    call mci_work%mci%final_simulation ()
  end subroutine mci_work_final_simulation

\eatline
\eatline
\nwendcode{}\nwbegindocs{756}Counter.
\nwenddocs{}\nwbegincode{757}\moddef{Processes: mci work: TBP}\plusendmoddef
  procedure :: reset_counter => mci_work_reset_counter
  procedure :: record_call => mci_work_record_call
  procedure :: get_counter => mci_work_get_counter
\nwendcode{}\nwbegincode{758}\moddef{Processes: procedures}\plusendmoddef
  subroutine mci_work_reset_counter (mci_work)
    class(mci_work_t), intent(inout) :: mci_work
    call mci_work%counter%reset ()
  end subroutine mci_work_reset_counter
  
  subroutine mci_work_record_call (mci_work, status)
    class(mci_work_t), intent(inout) :: mci_work
    integer, intent(in) :: status
    call mci_work%counter%record (status)
  end subroutine mci_work_record_call
    
  function mci_work_get_counter (mci_work) result (counter)
    class(mci_work_t), intent(in) :: mci_work
    type(process_counter_t) :: counter
    counter = mci_work%counter
  end function mci_work_get_counter
  
\eatline
\eatline
\eatline
\nwendcode{}\nwbegindocs{759}\nwdocspar
\subsubsection{The process instance}
A process instance contains all process data that depend on the
sampling point and thus change often.  In essence, it is an event
record at the elementary (parton) level.  We do not call it such, to
avoid confusion with the actual event records.  If decays are
involved, the latter are compositions of several elementary processes
(i.e., their instances).

We implement the process instance as an extension of the {\tt{}mci{\char95}sampler{\char95}t}
that we need for computing integrals and generate events.

The base type contains: the {\tt{}integrand}, the {\tt{}selected{\char95}channel}, the
two-dimensional array {\tt{}x} of parameters, and the one-dimensional array {\tt{}f}
of Jacobians.  These subobjects are public and used for communicating with the
multi-channel integrator.

The {\tt{}process} pointer accesses the process of which this record is
an instance.  It is required whenever the calculation needs invariant
configuration data, therefore the process should stay in memory for
the whole lifetime of its instances.

The {\tt{}evaluation{\char95}status} code is used to check the current status.
In particular, failure at various stages is recorded there.

The {\tt{}count} object records process evaluations, broken down
according to status.

The {\tt{}sqme} value is the single real number that results from
evaluating and tracing the kinematics and matrix elements.  This
is the number that is handed over to an integration routine.

The {\tt{}weight} value is the event weight.  It is defined when an event
has been generated from the process instance, either weighted or
unweighted.  The value is the {\tt{}sqme} value times Jacobian weights
from the integration, or unity, respectively.

The {\tt{}i{\char95}mci} index chooses a subset of components that are associated with
a common parameter set and integrator, i.e., that are added coherently.

The {\tt{}sf{\char95}chain} subobject is a realization of the beam and
structure-function configuration in the {\tt{}process} object.  It is not
used for calculation directly but serves as the template for the
sf-chain instances that are contained in the {\tt{}component} objects.

The {\tt{}component} subobjects determine the state of each component.

The {\tt{}term} subobjects are workspace for evaluating kinematics,
matrix elements, cuts etc.

The {\tt{}mci{\char95}work} subobject contains the array of real input parameters (random
numbers) that generates the kinematical point.  It also contains the workspace
for the MC integrators.  The active entry of the {\tt{}mci{\char95}work} array is
selected by the {\tt{}i{\char95}mci} index above.
\nwenddocs{}\nwbegincode{760}\moddef{Processes: public}\plusendmoddef
  public :: process_instance_t
\nwendcode{}\nwbegincode{761}\moddef{Processes: types}\plusendmoddef
  type, extends (mci_sampler_t) :: process_instance_t
     type(process_t), pointer :: process => null ()
     integer :: evaluation_status = STAT_UNDEFINED
     real(default) :: sqme = 0
     real(default), dimension(:), allocatable :: sqme_real
     real(default) :: weight = 0
     real(default) :: excess = 0
     integer :: i_mci = 0
     integer :: selected_channel = 0
     type(sf_chain_t) :: sf_chain
     type(component_instance_t), dimension(:), allocatable :: component
     type(term_instance_t), dimension(:), allocatable :: term
     type(mci_work_t), dimension(:), allocatable :: mci_work
     type(sqme_collector_t), pointer :: sqme_collector => null()
   contains
   \LA{}Processes: process instance: TBP\RA{}
  end type process_instance_t
     
\eatline
\nwendcode{}\nwbegindocs{762}The output routine contains a header with the most relevant
information about the process, copied from
{\tt{}process{\char95}metadata{\char95}write}.  We mark the active components by an asterisk.

The next section is the MC parameter input.  The following sections
are written only if the evaluation status is beyond setting the
parameters, or if the {\tt{}verbose} option is set.
\nwenddocs{}\nwbegincode{763}\moddef{Processes: process instance: TBP}\endmoddef
  procedure :: write_header => process_instance_write_header
  procedure :: write => process_instance_write
\nwendcode{}\nwbegincode{764}\moddef{Processes: procedures}\plusendmoddef
  subroutine process_instance_write_header (object, unit, testflag)
    class(process_instance_t), intent(in) :: object
    integer, intent(in), optional :: unit
    logical, intent(in), optional :: testflag
    integer :: u, i
    u = given_output_unit (unit)
    call write_separator (u, 2)
    if (associated (object%process)) then
       associate (meta => object%process%meta)
         select case (meta%type)
         case (PRC_UNKNOWN)
            write (u, "(1x,A)") "Process instance [undefined]"
            return
         case (PRC_DECAY)
            write (u, "(1x,A)", advance="no") "Process instance [decay]:"
         case (PRC_SCATTERING)
            write (u, "(1x,A)", advance="no") "Process instance [scattering]:"
         case default
            call msg_bug ("process_instance_write: undefined process type")
         end select
         write (u, "(1x,A,A,A)") "'", char (meta%id), "'"
         write (u, "(3x,A,A,A)") "Run ID = '", char (meta%run_id), "'"
         if (allocated (meta%component_id)) then
            write (u, "(3x,A)")  "Process components:"
            do i = 1, size (meta%component_id)
               if (object%component(i)%active) then
                  write (u, "(3x,'*')", advance="no")
               else
                  write (u, "(4x)", advance="no")
               end if
               write (u, "(1x,I0,9A)")  i, ": '", &
                    char (meta%component_id (i)), "':   ", &
                    char (meta%component_description (i))
            end do
         end if
       end associate
    else
       write (u, "(1x,A)") "Process instance [undefined process]"
       return
    end if
    write (u, "(3x,A)", advance = "no")  "status = "
    select case (object%evaluation_status)
    case (STAT_INITIAL);            write (u, "(A)")  "initialized"
    case (STAT_ACTIVATED);          write (u, "(A)")  "activated"
    case (STAT_BEAM_MOMENTA);       write (u, "(A)")  "beam momenta set"
    case (STAT_FAILED_KINEMATICS);  write (u, "(A)")  "failed kinematics"
    case (STAT_SEED_KINEMATICS);    write (u, "(A)")  "seed kinematics"
    case (STAT_HARD_KINEMATICS);    write (u, "(A)")  "hard kinematics"
    case (STAT_EFF_KINEMATICS);     write (u, "(A)")  "effective kinematics"
    case (STAT_FAILED_CUTS);        write (u, "(A)")  "failed cuts"
    case (STAT_PASSED_CUTS);        write (u, "(A)")  "passed cuts"
    case (STAT_EVALUATED_TRACE);    write (u, "(A)")  "evaluated trace"
       call write_separator (u)
       write (u, "(3x,A,ES19.12)")  "sqme   = ", object%sqme
    case (STAT_EVENT_COMPLETE);   write (u, "(A)")  "event complete"
       call write_separator (u)
       write (u, "(3x,A,ES19.12)")  "sqme   = ", object%sqme
       write (u, "(3x,A,ES19.12)")  "weight = ", object%weight
       if (object%excess /= 0) &
            write (u, "(3x,A,ES19.12)")  "excess = ", object%excess
    case default;                 write (u, "(A)")  "undefined"
    end select
    if (object%i_mci /= 0) then
       call write_separator (u)
       call object%mci_work(object%i_mci)%write (u, testflag)
    end if
    call write_separator (u, 2)
  end subroutine process_instance_write_header

  subroutine process_instance_write (object, unit, testflag)
    class(process_instance_t), intent(in) :: object
    integer, intent(in), optional :: unit
    logical, intent(in), optional :: testflag
    integer :: u, i
    u = given_output_unit (unit)
    call object%write_header (u)
    if (object%evaluation_status >= STAT_BEAM_MOMENTA) then
       call object%sf_chain%write (u)
       call write_separator (u, 2)
       if (object%evaluation_status >= STAT_SEED_KINEMATICS) then
          write (u, "(1x,A)") "Active components:"
          do i = 1, size (object%component)
             if (object%component(i)%active) then
                call write_separator (u)
                call object%component(i)%write (u, testflag)
             end if
          end do
          if (object%evaluation_status >= STAT_HARD_KINEMATICS) then
             call write_separator (u, 2)
             write (u, "(1x,A)") "Active terms:"
             if (any (object%term%active)) then
                do i = 1, size (object%term)
                   if (object%term(i)%active) then
                      call write_separator (u)
                      call object%term(i)%write (u, &
                           show_eff_state = &
                           object%evaluation_status >= STAT_EFF_KINEMATICS, &
                           testflag = testflag)
                   end if
                end do
             end if
          end if
          call write_separator (u, 2)
       end if
    end if
  end subroutine process_instance_write

\eatline
\eatline
\nwendcode{}\nwbegindocs{765}Finalize all subobjects that may contain allocated pointers.
\nwenddocs{}\nwbegincode{766}\moddef{Processes: process instance: TBP}\plusendmoddef
  procedure :: final => process_instance_final
\nwendcode{}\nwbegincode{767}\moddef{Processes: procedures}\plusendmoddef
  subroutine process_instance_final (instance)
    class(process_instance_t), intent(inout) :: instance
    integer :: i
    if (allocated (instance%mci_work)) then
       do i = 1, size (instance%mci_work)
          call instance%mci_work(i)%final ()
       end do
    end if
    call instance%sf_chain%final ()
    if (allocated (instance%component)) then
       do i = 1, size (instance%component)
          call instance%component(i)%final ()
       end do
    end if
    if (allocated (instance%term)) then
       do i = 1, size (instance%term)
          call instance%term(i)%final ()
       end do
    end if
    instance%evaluation_status = STAT_UNDEFINED
  end subroutine process_instance_final

\eatline
\nwendcode{}\nwbegindocs{768}Revert the process instance to initial state.  We do not deallocate
anything, just reset the state index and deactivate all components and
terms.

We do not reset the choice of the MCI set {\tt{}i{\char95}mci} unless this is
required explicitly.
\nwenddocs{}\nwbegincode{769}\moddef{Processes: process instance: TBP}\plusendmoddef
  procedure :: reset => process_instance_reset
\nwendcode{}\nwbegincode{770}\moddef{Processes: procedures}\plusendmoddef
  subroutine process_instance_reset (instance, reset_mci)
    class(process_instance_t), intent(inout) :: instance
    logical, intent(in), optional :: reset_mci
    integer :: i
    instance%component%active = .false.
    do i = 1, size (instance%term)
       call instance%term(i)%reset ()
    end do
    instance%term%checked = .false.
    instance%term%passed = .false.
    if (present (reset_mci)) then
       if (reset_mci)  instance%i_mci = 0
    end if
    instance%selected_channel = 0
    instance%evaluation_status = STAT_INITIAL
  end subroutine process_instance_reset
  
\eatline
\nwendcode{}\nwbegindocs{771}Activate the components and terms that correspond to a currently
selected MCI parameter set.
\nwenddocs{}\nwbegincode{772}\moddef{Processes: process instance: TBP}\plusendmoddef
  procedure :: activate => process_instance_activate
\nwendcode{}\nwbegincode{773}\moddef{Processes: procedures}\plusendmoddef
  subroutine process_instance_activate (instance)
    class(process_instance_t), intent(inout) :: instance
    integer :: i, j
    associate (mci_work => instance%mci_work(instance%i_mci))
      instance%component(mci_work%get_active_components ())%active &
           = .true.
      do i = 1, size (instance%component)
         associate (component => instance%component(i))
           if (component%active) then
              do j = 1, size (component%config%i_term)
                 instance%term(component%config%i_term(j))%active &
                      = .true.
              end do
           end if
         end associate
      end do
    end associate
    instance%evaluation_status = STAT_ACTIVATED
  end subroutine process_instance_activate
  
\eatline
\nwendcode{}\nwbegindocs{774}Initialization connects the instance with a process.  All initial
information is transferred from the process object.  The process
object contains templates for the interaction subobjects (beam and
term), but no evaluators.  The initialization routine
creates evaluators for the matrix element trace, other evaluators
are left untouched.

The {\tt{}instance} object must have the {\tt{}target} attribute (also in
any caller) since the initialization routine assigns various pointers to
subobject of {\tt{}instance}.
\nwenddocs{}\nwbegincode{775}\moddef{Processes: process instance: TBP}\plusendmoddef
  procedure :: init => process_instance_init
\nwendcode{}\nwbegincode{776}\moddef{Processes: procedures}\plusendmoddef
  subroutine process_instance_init (instance, process, cc_flag, &
                                    combined_integration)
    class(process_instance_t), intent(out), target :: instance
    type(process_t), intent(inout), target :: process
    logical, intent(in), optional :: cc_flag
    logical, intent(in), optional :: combined_integration
    integer :: i, i_component
    type(process_constants_t), dimension(2), save :: prc_constants
    type(string_t) :: nlo_type
    integer :: i_born, i_real
    type(nlo_pointer_t), dimension(:), allocatable :: nlo_p
    logical :: use_internal_cc, use_internal_sc

    if (present (combined_integration)) then
!!       if (combined_integration) allocate (instance%sqme_collector)
    end if
    instance%process => process
    call instance%setup_sf_chain (process%beam_config)
    allocate (instance%mci_work (process%config%n_mci))
    do i = 1, size (instance%mci_work)
       call instance%mci_work(i)%init (process%mci_entry(i))
    end do
    allocate (instance%component (process%config%n_components))
    allocate (nlo_p (process%config%n_components))
    do i = 1, size (nlo_p)
      allocate (nlo_p(i)%nlo_data)
    end do
    do i_component = 1, size (instance%component)
       if (process%component(i_component)%active) then
          associate (component => instance%component(i_component))
            call component%init (process%component(i_component))
            nlo_type = component%config%config%get_nlo_type ()
            select case (char (nlo_type))
            case ('Born')
              call nlo_p(i_component)%nlo_data%set_nlo_type (var_str ('Born'))
              if (associated (instance%sqme_collector)) &
                 call component%set_component_type (COMP_MASTER)
            case ('Real')
              i_born = component%config%config%get_associated_born ()
              call process%term(i_born)%fetch_constants (prc_constants(1))
              call process%term(i_component)%fetch_constants(prc_constants(2))
              if (present (cc_flag)) then
                 use_internal_cc = cc_flag
                 use_internal_sc = .false.
              else
                 use_internal_cc = .true.
                 use_internal_sc = .false.
              end if
              call nlo_p(i_component)%nlo_data%set_internal_procedures &
                                          (use_internal_cc, use_internal_sc)
              call nlo_p(i_component)%nlo_data%init (var_str ('Real'), prc_constants, &
                                               process%component(i_component)%fks_template, &
                                               process%config%model)
              call nlo_p(i_component)%nlo_data%set_flv_born (process%term(i_born)%flv)
              call nlo_p(i_component)%nlo_data%set_col_born (process%term(i_born)%col)
              call nlo_p(i_component)%nlo_data%set_hel_born (process%term(i_born)%hel)
              allocate (instance%sqme_real (size (nlo_p(i_component)%nlo_data%reg_data%regions)))
              if (associated (instance%sqme_collector)) &
                 call component%set_component_type (COMP_REAL)
            case ('Virtual')
              if (associated (instance%sqme_collector)) &
                 call component%set_component_type (COMP_VIRT) 
            end select
            call component%setup_kinematics &
                 (instance%sf_chain, process%beam_config, &
                  combined_integration)
          end associate
       end if
    end do
    allocate (instance%term (process%config%n_terms))
    do i = 1, size (instance%term)
       associate (term => instance%term(i))
         i_component = process%term(i)%i_component
         if (i_component /= 0) then
            associate (component => instance%component(i_component))
              nlo_type = component%config%config%get_nlo_type ()
              call component%k_seed%set_nlo_data (nlo_type)
              select case (char (nlo_type))
              case ('Virtual')
                i_born = component%config%config%get_associated_born ()
                i_real = i_born + process%config%n_components / 3
                nlo_p(i_component)%nlo_data = instance%term(i_real)%nlo_data
                call nlo_p(i_component)%nlo_data%set_nlo_type (var_str ('Virtual'))
                call nlo_p(i_component)%nlo_data%init_virtual ()
              end select

              call term%init (process%term(i), &
                   component%k_seed, &
                   process%beam_config, &
                   process%component(i_component)%core, &
                   process%meta%var_list, nlo_p(i_component)%nlo_data)
              call term%setup_expressions (process%meta, process%config)
              select case (char (nlo_type))
              case ('Real')
                call term%setup_expressions_real (process%meta, process%config)
              end select
              component%nlo_data => term%nlo_data
           end associate
         end if
       end associate
    end do
    instance%evaluation_status = STAT_INITIAL
  end subroutine process_instance_init
  
\eatline
\nwendcode{}\nwbegindocs{777}Subroutine of the initialization above: initialize the beam and
structure-function chain template.  We establish pointers to the
configuration data, so {\tt{}beam{\char95}config} must have a {\tt{}target}
attribute.

The resulting chain is not used directly for calculation.  It will
acquire instances which are stored in the process-component instance
objects.
\nwenddocs{}\nwbegincode{778}\moddef{Processes: process instance: TBP}\plusendmoddef
  procedure :: setup_sf_chain => process_instance_setup_sf_chain
\nwendcode{}\nwbegincode{779}\moddef{Processes: procedures}\plusendmoddef
  subroutine process_instance_setup_sf_chain (instance, config)
    class(process_instance_t), intent(inout) :: instance
    type(process_beam_config_t), intent(in), target :: config
    integer :: n_strfun
    n_strfun = config%n_strfun
    if (n_strfun /= 0) then
       call instance%sf_chain%init (config%data, config%sf)
    else
       call instance%sf_chain%init (config%data)
    end if
    if (config%sf_trace) then
       call instance%sf_chain%setup_tracing (config%sf_trace_file)
    end if
  end subroutine process_instance_setup_sf_chain
    
\eatline
\nwendcode{}\nwbegindocs{780}This initialization routine should be called only for process
instances which we intend as a source for physical events.  It
initializes the evaluators in the parton states of the terms.  They
describe the (semi-)exclusive transition matrix and the distribution
of color flow for the partonic process, convoluted with the beam and
structure-function chain.

If the model is not provided explicitly, we may use the model instance that
belongs to the process.  However, an explicit model allows us to override
particle settings.
\nwenddocs{}\nwbegincode{781}\moddef{Processes: process instance: TBP}\plusendmoddef
  procedure :: setup_event_data => process_instance_setup_event_data
\nwendcode{}\nwbegincode{782}\moddef{Processes: procedures}\plusendmoddef
  subroutine process_instance_setup_event_data (instance, model)
    class(process_instance_t), intent(inout), target :: instance
    class(model_data_t), intent(in), optional, target :: model
    class(model_data_t), pointer :: current_model
    integer :: i, i_component
    if (present (model)) then
       current_model => model
    else
       current_model => instance%process%config%model
    end if
    do i = 1, size (instance%term)
       associate (term => instance%term(i))
         if (associated (term%config)) then
            i_component = term%config%i_component
            associate (component => instance%process%component(i_component))
              call term%setup_event_data (component%core, current_model)
            end associate
         end if
       end associate
    end do
  end subroutine process_instance_setup_event_data

\eatline
\nwendcode{}\nwbegindocs{783}\nwdocspar
\nwenddocs{}\nwbegincode{784}\moddef{Processes: process instance: TBP}\plusendmoddef
  procedure :: init_born_amps => process_instance_init_born_amps
\nwendcode{}\nwbegincode{785}\moddef{Processes: procedures}\plusendmoddef
  subroutine process_instance_init_born_amps (instance, i_component, i_born)
    class(process_instance_t), intent(inout) :: instance
    integer, intent(in) :: i_component, i_born
    select type (core_born => instance%component(i_component)%config%core_sub_born)
    type is (prc_omega_t)
      call instance%component(i_component)%nlo_data%init_born_amps &
                          (instance%term(i_born)%config%n_allowed, .true.)
    type is (prc_gosam_t)
      call instance%component(i_component)%nlo_data%init_born_amps &
                          (instance%term(i_born)%config%n_allowed, .false.)
    end select
  end subroutine process_instance_init_born_amps

\eatline
\nwendcode{}\nwbegindocs{786}  
\nwenddocs{}\nwbegindocs{787}Choose a MC parameter set and the corresponding integrator. 

The choice persists beyond calls of the {\tt{}reset} method above.  This method
is automatically called here.
\nwenddocs{}\nwbegincode{788}\moddef{Processes: process instance: TBP}\plusendmoddef
  procedure :: choose_mci => process_instance_choose_mci
\nwendcode{}\nwbegincode{789}\moddef{Processes: procedures}\plusendmoddef
  subroutine process_instance_choose_mci (instance, i_mci)
    class(process_instance_t), intent(inout) :: instance
    integer, intent(in) :: i_mci
    instance%i_mci = i_mci
    call instance%reset ()
  end subroutine process_instance_choose_mci
    
\eatline
\nwendcode{}\nwbegindocs{790}Explicitly set a MC parameter set.  Works only if we are in initial
state.  We assume that the length of the parameter set is correct.  

After setting the parameters, activate the components and terms that
correspond to the chosen MC parameter set.
\nwenddocs{}\nwbegincode{791}\moddef{Processes: process instance: TBP}\plusendmoddef
  procedure :: set_mcpar => process_instance_set_mcpar
\nwendcode{}\nwbegincode{792}\moddef{Processes: procedures}\plusendmoddef
  subroutine process_instance_set_mcpar (instance, x)
    class(process_instance_t), intent(inout) :: instance
    real(default), dimension(:), intent(in) :: x
    if (instance%evaluation_status == STAT_INITIAL) then
       associate (mci_work => instance%mci_work(instance%i_mci))
         call mci_work%set (x)
       end associate
       call instance%activate ()
    end if
  end subroutine process_instance_set_mcpar

\eatline
\nwendcode{}\nwbegindocs{793}Receive the beam momentum/momenta from a source interaction.  This
applies to a cascade decay process instance, where the `beam' momentum
varies event by event.

The master beam momentum array is contained in the main structure
function chain subobject {\tt{}sf{\char95}chain}.  The sf-chain instance that
reside in the components will take their beam momenta from there.

The procedure transforms the instance status into
{\tt{}STAT{\char95}BEAM{\char95}MOMENTA}.  For process instance with fixed beam, this
intermediate status is skipped.
\nwenddocs{}\nwbegincode{794}\moddef{Processes: process instance: TBP}\plusendmoddef
  procedure :: receive_beam_momenta => process_instance_receive_beam_momenta
\nwendcode{}\nwbegincode{795}\moddef{Processes: procedures}\plusendmoddef
  subroutine process_instance_receive_beam_momenta (instance)
    class(process_instance_t), intent(inout) :: instance
    if (instance%evaluation_status >= STAT_INITIAL) then
       call instance%sf_chain%receive_beam_momenta ()
       instance%evaluation_status = STAT_BEAM_MOMENTA
    end if
  end subroutine process_instance_receive_beam_momenta
    
\eatline
\nwendcode{}\nwbegindocs{796}Set the beam momentum/momenta explicitly.  Otherwise, analogous to
the previous procedure.
\nwenddocs{}\nwbegincode{797}\moddef{Processes: process instance: TBP}\plusendmoddef
  procedure :: set_beam_momenta => process_instance_set_beam_momenta
\nwendcode{}\nwbegincode{798}\moddef{Processes: procedures}\plusendmoddef
  subroutine process_instance_set_beam_momenta (instance, p)
    class(process_instance_t), intent(inout) :: instance
    type(vector4_t), dimension(:), intent(in) :: p
    if (instance%evaluation_status >= STAT_INITIAL) then
       call instance%sf_chain%set_beam_momenta (p)
       instance%evaluation_status = STAT_BEAM_MOMENTA
    end if
  end subroutine process_instance_set_beam_momenta
    
\eatline
\nwendcode{}\nwbegindocs{799}Recover the initial beam momenta (those in the {\tt{}sf{\char95}chain}
component), given a valid (recovered) {\tt{}sf{\char95}chain{\char95}instance} in one of
the active components.  We need to do this only if the lab frame is
not the c.m.\ frame, otherwise those beams would be fixed anyway.
\nwenddocs{}\nwbegincode{800}\moddef{Processes: process instance: TBP}\plusendmoddef
  procedure :: recover_beam_momenta => process_instance_recover_beam_momenta
\nwendcode{}\nwbegincode{801}\moddef{Processes: procedures}\plusendmoddef
  subroutine process_instance_recover_beam_momenta (instance, i_term)
    class(process_instance_t), intent(inout) :: instance
    integer, intent(in) :: i_term
    integer :: i
    if (.not. instance%process%beam_config%lab_is_cm_frame) then
       if (instance%evaluation_status >= STAT_EFF_KINEMATICS) then
          i = instance%term(i_term)%config%i_component
          call instance%component(i)%return_beam_momenta ()
       end if
    end if
  end subroutine process_instance_recover_beam_momenta

\eatline
\nwendcode{}\nwbegindocs{802}Explicitly choose MC integration channel.  We assume here that the channel
count is identical for all active components.
\nwenddocs{}\nwbegincode{803}\moddef{Processes: process instance: TBP}\plusendmoddef
  procedure :: select_channel => process_instance_select_channel
\nwendcode{}\nwbegincode{804}\moddef{Processes: procedures}\plusendmoddef
  subroutine process_instance_select_channel (instance, channel)
    class(process_instance_t), intent(inout) :: instance
    integer, intent(in) :: channel
    instance%selected_channel = channel
  end subroutine process_instance_select_channel
  
\eatline
\nwendcode{}\nwbegindocs{805}First step of process evaluation: set up seed kinematics.  That is, for each
active process component, compute a momentum array from the MC input
parameters.

If {\tt{}skip{\char95}term} is set, we skip the component that accesses this
term.  We can assume that the associated data have already been
recovered, and we are just computing the rest.
\nwenddocs{}\nwbegincode{806}\moddef{Processes: process instance: TBP}\plusendmoddef
  procedure :: compute_seed_kinematics => &
       process_instance_compute_seed_kinematics
\nwendcode{}\nwbegincode{807}\moddef{Processes: procedures}\plusendmoddef
  subroutine process_instance_compute_seed_kinematics (instance, skip_term)
    class(process_instance_t), intent(inout) :: instance
    integer, intent(in), optional :: skip_term
    integer :: channel, skip_component, i
    logical :: success
    channel = instance%selected_channel
    if (channel == 0) then
       call msg_bug ("Compute seed kinematics: undefined integration channel")
    end if
    if (present (skip_term)) then
       skip_component = instance%term(skip_term)%config%i_component
    else
       skip_component = 0
    end if
    if (instance%evaluation_status >= STAT_ACTIVATED) then
       success = .true.
       do i = 1, size (instance%component)
          if (i == skip_component)  cycle
          if (instance%component(i)%active) then
             call instance%component(i)%compute_seed_kinematics &
                  (instance%mci_work(instance%i_mci), channel, success)
             if (.not. success)  exit
          end if
       end do
       if (success) then
          instance%evaluation_status = STAT_SEED_KINEMATICS
       else
          instance%evaluation_status = STAT_FAILED_KINEMATICS
       end if
    end if
  end subroutine process_instance_compute_seed_kinematics

\eatline
\nwendcode{}\nwbegindocs{808}Inverse: recover missing parts of the kinematics from the momentum
configuration, which we know for a single term and component.   Given
a channel, reconstruct the MC parameter set.
\nwenddocs{}\nwbegincode{809}\moddef{Processes: process instance: TBP}\plusendmoddef
  procedure :: recover_mcpar => process_instance_recover_mcpar
\nwendcode{}\nwbegincode{810}\moddef{Processes: procedures}\plusendmoddef
  subroutine process_instance_recover_mcpar (instance, i_term)
    class(process_instance_t), intent(inout) :: instance
    integer, intent(in) :: i_term
    integer :: channel
    integer :: i
    if (instance%evaluation_status >= STAT_EFF_KINEMATICS) then
       channel = instance%selected_channel
       if (channel == 0) then
          call msg_bug ("Recover MC parameters: undefined integration channel")
       end if
       i = instance%term(i_term)%config%i_component
       call instance%component(i)%recover_mcpar &
                  (instance%mci_work(instance%i_mci), channel)
    end if
  end subroutine process_instance_recover_mcpar

\eatline
\nwendcode{}\nwbegindocs{811}Second step of process evaluation: compute all momenta, for all active
components, from the seed kinematics.
\nwenddocs{}\nwbegincode{812}\moddef{Processes: process instance: TBP}\plusendmoddef
  procedure :: compute_hard_kinematics => &
       process_instance_compute_hard_kinematics
\nwendcode{}\nwbegincode{813}\moddef{Processes: procedures}\plusendmoddef
  subroutine process_instance_compute_hard_kinematics (instance, skip_term)
    class(process_instance_t), intent(inout) :: instance
    integer, intent(in), optional :: skip_term
    integer :: i
    logical :: real_phsp
    if (instance%evaluation_status >= STAT_SEED_KINEMATICS) then
       do i = 1, size (instance%component)
          if (instance%component(i)%active) then
            select case (char (instance%component(i)%config%config%get_nlo_type()))
            case ('Born', 'Virtual')
              real_phsp = .false.
            case ('Real')
              real_phsp = .true.
            end select
              call instance%component(i)% &
                  compute_hard_kinematics (instance%term, skip_term, real_phsp)
          end if
       end do
       instance%evaluation_status = STAT_HARD_KINEMATICS
    end if
  end subroutine process_instance_compute_hard_kinematics

\eatline
\nwendcode{}\nwbegindocs{814}Inverse: recover seed kinematics.  We know the beam momentum
configuration and the outgoing momenta of the effective interaction,
for one specific term.
\nwenddocs{}\nwbegincode{815}\moddef{Processes: process instance: TBP}\plusendmoddef
  procedure :: recover_seed_kinematics => &
       process_instance_recover_seed_kinematics
\nwendcode{}\nwbegincode{816}\moddef{Processes: procedures}\plusendmoddef
  subroutine process_instance_recover_seed_kinematics (instance, i_term)
    class(process_instance_t), intent(inout) :: instance
    integer, intent(in) :: i_term
    if (instance%evaluation_status >= STAT_EFF_KINEMATICS) then
       associate (i_component => instance%term(i_term)%config%i_component)
         call instance%component(i_component)% &
              recover_seed_kinematics (instance%term(i_term))
       end associate
    end if
  end subroutine process_instance_recover_seed_kinematics
  
\eatline
\nwendcode{}\nwbegindocs{817}Third step of process evaluation: compute the effective momentum
configurations, for all active terms, from the hard kinematics.
\nwenddocs{}\nwbegincode{818}\moddef{Processes: process instance: TBP}\plusendmoddef
  procedure :: compute_eff_kinematics => &
       process_instance_compute_eff_kinematics
\nwendcode{}\nwbegincode{819}\moddef{Processes: procedures}\plusendmoddef
  subroutine process_instance_compute_eff_kinematics (instance, skip_term)
    class(process_instance_t), intent(inout) :: instance
    integer, intent(in), optional :: skip_term
    integer :: i
    if (instance%evaluation_status >= STAT_HARD_KINEMATICS) then
       do i = 1, size (instance%term)
          if (present (skip_term)) then
             if (i == skip_term)  cycle
          end if
          if (instance%term(i)%active) then
             call instance%term(i)% &
                  compute_eff_kinematics (instance%component)
          end if
       end do
       instance%evaluation_status = STAT_EFF_KINEMATICS
    end if
  end subroutine process_instance_compute_eff_kinematics

\eatline
\nwendcode{}\nwbegindocs{820}Inverse: recover the hard kinematics from effective kinematics for
one term, then compute effective kinematics for the other terms.
\nwenddocs{}\nwbegincode{821}\moddef{Processes: process instance: TBP}\plusendmoddef
  procedure :: recover_hard_kinematics => &
       process_instance_recover_hard_kinematics
\nwendcode{}\nwbegincode{822}\moddef{Processes: procedures}\plusendmoddef
  subroutine process_instance_recover_hard_kinematics (instance, i_term)
    class(process_instance_t), intent(inout) :: instance
    integer, intent(in) :: i_term
    integer :: i
    if (instance%evaluation_status >= STAT_EFF_KINEMATICS) then
       call instance%term(i_term)%recover_hard_kinematics (instance%component)
       do i = 1, size (instance%term)
          if (i /= i_term) then
             if (instance%term(i)%active) then
                call instance%term(i)% &
                     compute_eff_kinematics (instance%component)
             end if
          end if
       end do
       instance%evaluation_status = STAT_EFF_KINEMATICS
    end if
  end subroutine process_instance_recover_hard_kinematics
       
\eatline
\nwendcode{}\nwbegindocs{823}Fourth step of process evaluation: check cuts for all terms.  Where
sucessful, compute any scales and weights.  Otherwise, deactive the term.
If any of the terms has passed, set the state to {\tt{}STAT{\char95}PASSED{\char95}CUTS}.

The argument {\tt{}scale{\char95}forced}, if present, will override the scale calculation
in the term expressions.
\nwenddocs{}\nwbegincode{824}\moddef{Processes: process instance: TBP}\plusendmoddef
  procedure :: evaluate_expressions => &
       process_instance_evaluate_expressions
\nwendcode{}\nwbegincode{825}\moddef{Processes: procedures}\plusendmoddef
  subroutine process_instance_evaluate_expressions (instance, scale_forced)
    class(process_instance_t), intent(inout) :: instance
    real(default), intent(in), allocatable, optional :: scale_forced
    integer :: i
    if (instance%evaluation_status >= STAT_EFF_KINEMATICS) then
       do i = 1, size (instance%term)
          if (instance%term(i)%active) then
             call instance%term(i)%evaluate_expressions (scale_forced)
          end if
       end do
       if (any (instance%term%passed)) then
          instance%evaluation_status = STAT_PASSED_CUTS
       else
          instance%evaluation_status = STAT_FAILED_CUTS
       end if
    end if
  end subroutine process_instance_evaluate_expressions

\eatline
\nwendcode{}\nwbegindocs{826}Fifth step of process evaluation: fill the parameters for the non-selected
,channels, that have not been used for seeding.  We should do this after
evaluating cuts, since we may save some expensive calculations if the phase
space point fails the cuts.

If {\tt{}skip{\char95}term} is set, we skip the component that accesses this
term.  We can assume that the associated data have already been
recovered, and we are just computing the rest.
\nwenddocs{}\nwbegincode{827}\moddef{Processes: process instance: TBP}\plusendmoddef
  procedure :: compute_other_channels => &
       process_instance_compute_other_channels
\nwendcode{}\nwbegincode{828}\moddef{Processes: procedures}\plusendmoddef
  subroutine process_instance_compute_other_channels (instance, skip_term)
    class(process_instance_t), intent(inout) :: instance
    integer, intent(in), optional :: skip_term
    integer :: channel, skip_component, i
    channel = instance%selected_channel
    if (channel == 0) then
       call msg_bug ("Compute other channels: undefined integration channel")
    end if
    if (present (skip_term)) then
       skip_component = instance%term(skip_term)%config%i_component
    else
       skip_component = 0
    end if
    if (instance%evaluation_status >= STAT_PASSED_CUTS) then
       do i = 1, size (instance%component)
          if (i == skip_component)  cycle
          if (instance%component(i)%active) then
             call instance%component(i)%compute_other_channels &
                  (instance%mci_work(instance%i_mci), channel)
          end if
       end do
    end if
  end subroutine process_instance_compute_other_channels

\eatline
\nwendcode{}\nwbegindocs{829}Sixth step of process evaluation: evaluate the matrix elements, and compute
the trace (summed over quantum numbers) for all terms.  Finally, sum up the
terms, iterating over all active process components.
\nwenddocs{}\nwbegincode{830}\moddef{Processes: process instance: TBP}\plusendmoddef
  procedure :: evaluate_trace => process_instance_evaluate_trace
\nwendcode{}\nwbegincode{831}\moddef{Processes: procedures}\plusendmoddef
  subroutine process_instance_evaluate_trace (instance)
    class(process_instance_t), intent(inout) :: instance
    integer :: i, j
    integer :: n_legs
    logical :: passed
    integer :: ireg
    integer :: i_real
    real(default) :: sqme_born
    instance%sqme = 0
    i_real = 0
    if (instance%evaluation_status >= STAT_PASSED_CUTS) then
       do i = 1, size (instance%term)
          associate (term => instance%term(i))
            if (term%active .and. term%passed) then
              select case (char (term%nlo_data%nlo_type))
              case ('Real')
                 if (.not. associated (instance%sqme_collector)) then
                    call instance%evaluate_trace_real (term, i)
                 else
                    i_real = i
                 end if
              case ('Born', 'Virtual')
                 call term%evaluate_interaction (instance%component)
                 call term%evaluate_trace ()
              end select
            end if
          end associate
       end do
       do i = 1, size (instance%component)
          associate (component => instance%component(i))
            if (component%active) then
              select case (char (component%config%config%get_nlo_type()))
              case ('Born', 'Virtual')
                call component%evaluate_sqme (instance%term)
                instance%sqme = instance%sqme + component%sqme
                instance%evaluation_status = STAT_EVALUATED_TRACE
                if (associated (instance%sqme_collector)) then
                   select case (component%get_component_type())
                   case (COMP_MASTER)
                      instance%sqme_collector%sqme_born = component%sqme
                   case (COMP_VIRT)
                      instance%sqme_collector%sqme_virt = component%sqme
                   end select
                end if 
              end select
            end if
          end associate
       end do
       if (associated (instance%sqme_collector)) &
          call instance%evaluate_trace_real (instance%term(i_real), i_real)
    else
       ! failed kinematics, failed cuts: set sqme to zero
       instance%sqme = 0
    end if
  end subroutine process_instance_evaluate_trace

\eatline
\nwendcode{}\nwbegindocs{832}Final step of process evaluation: evaluate the matrix elements, and compute
the trace (summed over quantum numbers) for all terms.  Finally, sum up the
terms, iterating over all active process components.

If {\tt{}weight} is provided, we already know the kinematical event
weight (the MCI weight which depends on the kinematics sampling
algorithm, but not on the matrix element), so we do not need to take
it from the MCI record.
\nwenddocs{}\nwbegincode{833}\moddef{Processes: process instance: TBP}\plusendmoddef
  procedure :: evaluate_event_data => process_instance_evaluate_event_data
\nwendcode{}\nwbegincode{834}\moddef{Processes: procedures}\plusendmoddef
  subroutine process_instance_evaluate_event_data (instance, weight)
    class(process_instance_t), intent(inout) :: instance
    real(default), intent(in), optional :: weight
    integer :: i
    if (instance%evaluation_status >= STAT_EVALUATED_TRACE) then
       do i = 1, size (instance%term)
          associate (term => instance%term(i))
            if (term%active .and. term%passed) then
               call term%evaluate_event_data ()
            end if
          end associate
       end do
       if (present (weight)) then
          instance%weight = weight
       else
          instance%weight = &
               instance%mci_work(instance%i_mci)%mci%get_event_weight ()
          instance%excess = &
               instance%mci_work(instance%i_mci)%mci%get_event_excess ()
       end if
       instance%evaluation_status = STAT_EVENT_COMPLETE
    else
       ! failed kinematics etc.: set weight to zero
       instance%weight = 0
    end if
  end subroutine process_instance_evaluate_event_data

\eatline
\nwendcode{}\nwbegindocs{835}For the real part of the process evaluation at NLO, trace evaluation is a 
  bit more intricate because the phase-spaces for different emitters might
  be treated differently by cuts. A phase-space point is discarded if it fails
  the cuts for at least one emitter. 
\nwenddocs{}\nwbegincode{836}\moddef{Processes: process instance: TBP}\plusendmoddef
  procedure :: evaluate_trace_real => process_instance_evaluate_trace_real
\nwendcode{}\nwbegincode{837}\moddef{Processes: procedures}\plusendmoddef
  subroutine process_instance_evaluate_trace_real (instance, term, i)
     class(process_instance_t), intent(inout) :: instance
     type(term_instance_t), intent(inout) :: term
     integer, intent(in) :: i
     integer :: j, ireg
     integer :: nlegs
     logical :: passed
     real(default) :: sqme_born

     nlegs = term%nlo_data%reg_data%nlegs_real
     ireg = 1
     instance%sqme_real = 0
     do j = 1, nlegs
        if (any (term%nlo_data%reg_data%emitters == j)) then
           call term%nlo_data%set_active_emitter (j)
           call term%evaluate_interaction_real (instance%component(i), i)
           if (.not. term%passed) then
           !!! Cuts failed, leave subroutine
              instance%evaluation_status = STAT_FAILED_CUTS
              instance%sqme_real = 0
              return
           end if
           call term%evaluate_trace ()
           if (instance%component(i)%active) then
              associate (component => instance%component(i))
                 if (.not. associated (instance%sqme_collector)) then
                    call component%evaluate_sqme_born (term)
                    sqme_born = component%sqme
                 else
                    sqme_born = instance%sqme_collector%sqme_born
                 end if
                 call add_born_to_nlo_component (component, sqme_born*term%weight)
                 call component%evaluate_sqme_real (term)
                 instance%sqme_real(ireg) = instance%component(i)%nlo_data%sqme_real
              end associate
           end if
           ireg = ireg + 1
        end if
     end do
     instance%sqme = instance%sqme + sum (instance%sqme_real)
     if (associated (instance%sqme_collector)) &
        instance%sqme_collector%sqme_real = sum (instance%sqme_real)
     instance%evaluation_status = STAT_EVALUATED_TRACE  
  end subroutine process_instance_evaluate_trace_real

\eatline
\nwendcode{}\nwbegindocs{838}For unweighted event generation, we should reset the reported event
weight to unity (signed) or zero.  The latter case is appropriate for
an event which failed for whatever reason.
\nwenddocs{}\nwbegincode{839}\moddef{Processes: process instance: TBP}\plusendmoddef
  procedure :: normalize_weight => process_instance_normalize_weight
\nwendcode{}\nwbegincode{840}\moddef{Processes: procedures}\plusendmoddef
  subroutine process_instance_normalize_weight (instance)
    class(process_instance_t), intent(inout) :: instance
    if (instance%weight /= 0) then
       instance%weight = sign (1._default, instance%weight)
    end if
  end subroutine process_instance_normalize_weight
  
\eatline
\nwendcode{}\nwbegindocs{841}This is a convenience routine that performs the computations of the
steps 1 to 5 in a single step.  The arguments are the input for
{\tt{}set{\char95}mcpar}.  After this, the evaluation status should be either
{\tt{}STAT{\char95}FAILED{\char95}KINEMATICS}, {\tt{}STAT{\char95}FAILED{\char95}CUTS} or {\tt{}STAT{\char95}EVALUATED{\char95}TRACE}.

Before calling this, we should call {\tt{}choose{\char95}mci}.
\nwenddocs{}\nwbegincode{842}\moddef{Processes: process instance: TBP}\plusendmoddef
  procedure :: evaluate_sqme => process_instance_evaluate_sqme
\nwendcode{}\nwbegincode{843}\moddef{Processes: procedures}\plusendmoddef
  subroutine process_instance_evaluate_sqme (instance, channel, x)
    class(process_instance_t), intent(inout) :: instance
    integer, intent(in) :: channel
    real(default), dimension(:), intent(in) :: x
    call instance%reset ()
    call instance%set_mcpar (x)
    call instance%select_channel (channel)
    call instance%compute_seed_kinematics ()
    call instance%compute_hard_kinematics ()
    call instance%compute_eff_kinematics ()
    call instance%evaluate_expressions ()
    call instance%compute_other_channels ()
    call instance%evaluate_trace ()
    if (associated (instance%sqme_collector)) &
       instance%sqme = instance%sqme_collector%get_sqme_sum ()
  end subroutine process_instance_evaluate_sqme
  
\eatline
\nwendcode{}\nwbegindocs{844}This is the inverse.  Assuming that the final trace evaluator
contains a valid momentum configuration, recover kinematics
and recalculate the matrix elements and their trace.

To be precise, we first recover kinematics for the given term and
associated component, then recalculate from that all other terms and
active components.  The {\tt{}channel} is not really required to obtain
the matrix element, but it allows us to reconstruct the exact MC
parameter set that corresponds to the given phase space point.

Before calling this, we should call {\tt{}choose{\char95}mci}.
\nwenddocs{}\nwbegincode{845}\moddef{Processes: process instance: TBP}\plusendmoddef
  procedure :: recover => process_instance_recover
\nwendcode{}\nwbegincode{846}\moddef{Processes: procedures}\plusendmoddef
  subroutine process_instance_recover &
       (instance, channel, i_term, update_sqme, scale_forced)
    class(process_instance_t), intent(inout) :: instance
    integer, intent(in) :: channel
    integer, intent(in) :: i_term
    logical, intent(in) :: update_sqme
    real(default), intent(in), allocatable, optional :: scale_forced
    call instance%activate ()
    instance%evaluation_status = STAT_EFF_KINEMATICS
    call instance%recover_hard_kinematics (i_term)
    call instance%recover_seed_kinematics (i_term)
    call instance%select_channel (channel)
    call instance%recover_mcpar (i_term)
    call instance%recover_beam_momenta (i_term)
    call instance%compute_seed_kinematics (i_term)
    call instance%compute_hard_kinematics (i_term)
    call instance%compute_eff_kinematics (i_term)
    call instance%compute_other_channels (i_term)
    call instance%evaluate_expressions (scale_forced)
    if (update_sqme)  call instance%evaluate_trace ()
  end subroutine process_instance_recover
  
\eatline
\nwendcode{}\nwbegindocs{847}The {\tt{}evaluate} method is required by the {\tt{}sampler{\char95}t} base type of which
the process instance is an extension.

The requirement is that after the process instance is evaluated, the
integrand, the selected channel, the $x$ array, and the $f$ Jacobian array are
exposed by the {\tt{}sampler{\char95}t} object.
\nwenddocs{}\nwbegincode{848}\moddef{Processes: process instance: TBP}\plusendmoddef
  procedure :: evaluate => process_instance_evaluate
\nwendcode{}\nwbegincode{849}\moddef{Processes: procedures}\plusendmoddef
  subroutine process_instance_evaluate (sampler, c, x_in, val, x, f)
    class(process_instance_t), intent(inout) :: sampler
    integer, intent(in) :: c
    real(default), dimension(:), intent(in) :: x_in
    real(default), intent(out) :: val
    real(default), dimension(:,:), intent(out) :: x
    real(default), dimension(:), intent(out) :: f
    call sampler%evaluate_sqme (c, x_in)
    if (sampler%is_valid ())  call sampler%fetch (val, x, f)
    call sampler%record_call ()
  end subroutine process_instance_evaluate

\eatline
\nwendcode{}\nwbegindocs{850}The phase-space point is valid if the event has valid kinematics and
has passed the cuts.
\nwenddocs{}\nwbegincode{851}\moddef{Processes: process instance: TBP}\plusendmoddef
  procedure :: is_valid => process_instance_is_valid
\nwendcode{}\nwbegincode{852}\moddef{Processes: procedures}\plusendmoddef
  function process_instance_is_valid (sampler) result (valid)
    class(process_instance_t), intent(in) :: sampler
    logical :: valid
    valid = sampler%evaluation_status >= STAT_PASSED_CUTS
  end function process_instance_is_valid
  
\eatline
\nwendcode{}\nwbegindocs{853}The {\tt{}rebuild} method should rebuild the kinematics section out of
the {\tt{}x{\char95}in} parameter set.  The integrand value {\tt{}val} should not be
computed, but is provided as input.
\nwenddocs{}\nwbegincode{854}\moddef{Processes: process instance: TBP}\plusendmoddef
  procedure :: rebuild => process_instance_rebuild
\nwendcode{}\nwbegincode{855}\moddef{Processes: procedures}\plusendmoddef
  subroutine process_instance_rebuild (sampler, c, x_in, val, x, f)
    class(process_instance_t), intent(inout) :: sampler
    integer, intent(in) :: c
    real(default), dimension(:), intent(in) :: x_in
    real(default), intent(in) :: val
    real(default), dimension(:,:), intent(out) :: x
    real(default), dimension(:), intent(out) :: f
    call msg_bug ("process_instance_rebuild not implemented yet")
    x = 0
    f = 0
  end subroutine process_instance_rebuild

\eatline
\nwendcode{}\nwbegindocs{856}This is another method required by the {\tt{}sampler{\char95}t} base type:
fetch the data that are relevant for the MCI record.
\nwenddocs{}\nwbegincode{857}\moddef{Processes: process instance: TBP}\plusendmoddef
  procedure :: fetch => process_instance_fetch
\nwendcode{}\nwbegincode{858}\moddef{Processes: procedures}\plusendmoddef
  subroutine process_instance_fetch (sampler, val, x, f)
    class(process_instance_t), intent(in) :: sampler
    real(default), intent(out) :: val
    real(default), dimension(:,:), intent(out) :: x
    real(default), dimension(:), intent(out) :: f
    integer :: i, cc
    integer :: n_channel
    val = 0
    FIND_COMPONENT: do i = 1, size (sampler%component)
       associate (component => sampler%component(i))
         if (component%active) then
            associate (k => component%k_seed)
              n_channel = k%n_channel
              do cc = 1, n_channel
                 call k%get_mcpar (cc, x(:,cc))
              end do
              f = k%f
              val = sampler%sqme * k%phs_factor
            end associate
            exit FIND_COMPONENT
         end if
       end associate
    end do FIND_COMPONENT
  end subroutine process_instance_fetch
  
\eatline
\nwendcode{}\nwbegindocs{859}Initialize and finalize event generation for the specified MCI
entry.
\nwenddocs{}\nwbegincode{860}\moddef{Processes: process instance: TBP}\plusendmoddef
  procedure :: init_simulation => process_instance_init_simulation
  procedure :: final_simulation => process_instance_final_simulation
\nwendcode{}\nwbegincode{861}\moddef{Processes: procedures}\plusendmoddef
  subroutine process_instance_init_simulation (instance, i_mci, safety_factor)
    class(process_instance_t), intent(inout) :: instance
    integer, intent(in) :: i_mci
    real(default), intent(in), optional :: safety_factor
    call instance%mci_work(i_mci)%init_simulation (safety_factor)
  end subroutine process_instance_init_simulation

  subroutine process_instance_final_simulation (instance, i_mci)
    class(process_instance_t), intent(inout) :: instance
    integer, intent(in) :: i_mci
    call instance%mci_work(i_mci)%final_simulation ()
  end subroutine process_instance_final_simulation

\eatline
\eatline
\nwendcode{}\nwbegindocs{862}\nwdocspar
\subsection{Accessing the process instance}
Once the seed kinematics is complete, we can retrieve the MC input parameters
for all channels, not just the seed channel.

Note: We choose the first active component.  This makes sense only if the seed
kinematics is identical for all active components.
\nwenddocs{}\nwbegincode{863}\moddef{Processes: process instance: TBP}\plusendmoddef
  procedure :: get_mcpar => process_instance_get_mcpar
\nwendcode{}\nwbegincode{864}\moddef{Processes: procedures}\plusendmoddef
  subroutine process_instance_get_mcpar (instance, channel, x)
    class(process_instance_t), intent(inout) :: instance
    integer, intent(in) :: channel
    real(default), dimension(:), intent(out) :: x
    integer :: i
    if (instance%evaluation_status >= STAT_SEED_KINEMATICS) then
       do i = 1, size (instance%component)
          if (instance%component(i)%active) then
             call instance%component(i)%k_seed%get_mcpar (channel, x)
             return
          end if
       end do
       call msg_bug ("Process instance: get_mcpar: no active channels")
    else
       call msg_bug ("Process instance: get_mcpar: no seed kinematics")
    end if
  end subroutine process_instance_get_mcpar

\eatline
\nwendcode{}\nwbegindocs{865}Return true if the {\tt{}sqme} value is known.  This also implies that the
event is kinematically valid and has passed all cuts.
\nwenddocs{}\nwbegincode{866}\moddef{Processes: process instance: TBP}\plusendmoddef
  procedure :: has_evaluated_trace => process_instance_has_evaluated_trace
\nwendcode{}\nwbegincode{867}\moddef{Processes: procedures}\plusendmoddef
  function process_instance_has_evaluated_trace (instance) result (flag)
    class(process_instance_t), intent(in) :: instance
    logical :: flag
    flag = instance%evaluation_status >= STAT_EVALUATED_TRACE
  end function process_instance_has_evaluated_trace
  
\eatline
\nwendcode{}\nwbegindocs{868}Return true if the event is complete.  In particular, the event must
be kinematically valid, passed all cuts, and the event data have been
computed.
\nwenddocs{}\nwbegincode{869}\moddef{Processes: process instance: TBP}\plusendmoddef
  procedure :: is_complete_event => process_instance_is_complete_event
\nwendcode{}\nwbegincode{870}\moddef{Processes: procedures}\plusendmoddef
  function process_instance_is_complete_event (instance) result (flag)
    class(process_instance_t), intent(in) :: instance
    logical :: flag
    flag = instance%evaluation_status >= STAT_EVENT_COMPLETE
  end function process_instance_is_complete_event
  
\eatline
\nwendcode{}\nwbegindocs{871}Select a term for the process instance which is to provide the event
record.

Note: this should be done using random numbers and applying probabilities for
the various terms and components that belong to the selected MCI group.  The
current implementation simply selects the first term for the first active
component.
\nwenddocs{}\nwbegincode{872}\moddef{Processes: process instance: TBP}\plusendmoddef
  procedure :: select_i_term => process_instance_select_i_term
\nwendcode{}\nwbegincode{873}\moddef{Processes: procedures}\plusendmoddef
  subroutine process_instance_select_i_term (instance, i_term)
    class(process_instance_t), intent(in) :: instance
    integer, intent(out) :: i_term
    integer :: i_mci, i_component
    i_mci = instance%i_mci
    i_component = instance%process%mci_entry(i_mci)%i_component(1)
    i_term = instance%process%component(i_component)%i_term(1)
  end subroutine process_instance_select_i_term
    
\eatline
\nwendcode{}\nwbegindocs{874}Return pointer to the master beam interaction.
\nwenddocs{}\nwbegincode{875}\moddef{Processes: process instance: TBP}\plusendmoddef
  procedure :: get_beam_int_ptr => process_instance_get_beam_int_ptr
\nwendcode{}\nwbegincode{876}\moddef{Processes: procedures}\plusendmoddef
  function process_instance_get_beam_int_ptr (instance) result (ptr)
    class(process_instance_t), intent(in), target :: instance
    type(interaction_t), pointer :: ptr
    ptr => instance%sf_chain%get_beam_int_ptr ()
  end function process_instance_get_beam_int_ptr
  
\eatline
\nwendcode{}\nwbegindocs{877}Return pointers to the matrix and flows interactions, given a term index.
\nwenddocs{}\nwbegincode{878}\moddef{Processes: process instance: TBP}\plusendmoddef
  procedure :: get_trace_int_ptr => process_instance_get_trace_int_ptr
  procedure :: get_matrix_int_ptr => process_instance_get_matrix_int_ptr
  procedure :: get_flows_int_ptr => process_instance_get_flows_int_ptr
\nwendcode{}\nwbegincode{879}\moddef{Processes: procedures}\plusendmoddef
  function process_instance_get_trace_int_ptr (instance, i_term) result (ptr)
    class(process_instance_t), intent(in), target :: instance
    integer, intent(in) :: i_term
    type(interaction_t), pointer :: ptr
    ptr => instance%term(i_term)%connected%get_trace_int_ptr ()
  end function process_instance_get_trace_int_ptr
  
  function process_instance_get_matrix_int_ptr (instance, i_term) result (ptr)
    class(process_instance_t), intent(in), target :: instance
    integer, intent(in) :: i_term
    type(interaction_t), pointer :: ptr
    select case (char (instance%term(i_term)%nlo_data%nlo_type))
    case ('Born', 'Virtual')
      ptr => instance%term(i_term)%connected%get_matrix_int_ptr ()
    case ('Real')
      ptr => instance%term(i_term)%connected_real%get_matrix_int_ptr ()
    end select
  end function process_instance_get_matrix_int_ptr
  
  function process_instance_get_flows_int_ptr (instance, i_term) result (ptr)
    class(process_instance_t), intent(in), target :: instance
    integer, intent(in) :: i_term
    type(interaction_t), pointer :: ptr
    select case (char (instance%term(i_term)%nlo_data%nlo_type))
    case ('Born', 'Virtual')
      ptr => instance%term(i_term)%connected%get_flows_int_ptr ()
    case ('Real')
      ptr => instance%term(i_term)%connected_real%get_flows_int_ptr ()
    end select
  end function process_instance_get_flows_int_ptr
  
\eatline
\eatline
\eatline
\nwendcode{}\nwbegindocs{880}Return the complete account of flavor combinations in the underlying
interaction object, including beams, radiation, and hard interaction.
\nwenddocs{}\nwbegincode{881}\moddef{Processes: process instance: TBP}\plusendmoddef
  procedure :: get_state_flv => process_instance_get_state_flv
\nwendcode{}\nwbegincode{882}\moddef{Processes: procedures}\plusendmoddef
  function process_instance_get_state_flv (instance, i_term) result (state_flv)
    class(process_instance_t), intent(in) :: instance
    integer, intent(in) :: i_term
    type(state_flv_content_t) :: state_flv
    state_flv = instance%term(i_term)%connected%get_state_flv ()
  end function process_instance_get_state_flv
  
\eatline
\nwendcode{}\nwbegindocs{883}Return pointers to the parton states of a selected term.
\nwenddocs{}\nwbegincode{884}\moddef{Processes: process instance: TBP}\plusendmoddef
  procedure :: get_isolated_state_ptr => &
       process_instance_get_isolated_state_ptr
  procedure :: get_connected_state_ptr => &
       process_instance_get_connected_state_ptr
\nwendcode{}\nwbegincode{885}\moddef{Processes: procedures}\plusendmoddef
  function process_instance_get_isolated_state_ptr (instance, i_term) &
       result (ptr)
    class(process_instance_t), intent(in), target :: instance
    integer, intent(in) :: i_term
    type(isolated_state_t), pointer :: ptr
    ptr => instance%term(i_term)%isolated
  end function process_instance_get_isolated_state_ptr
  
  function process_instance_get_connected_state_ptr (instance, i_term) &
       result (ptr)
    class(process_instance_t), intent(in), target :: instance
    integer, intent(in) :: i_term
    type(connected_state_t), pointer :: ptr
    ptr => instance%term(i_term)%connected
  end function process_instance_get_connected_state_ptr
  
\eatline
\eatline
\nwendcode{}\nwbegindocs{886}Return the indices of the beam particles and incoming partons within the
currently active state matrix, respectively.
\nwenddocs{}\nwbegincode{887}\moddef{Processes: process instance: TBP}\plusendmoddef
  procedure :: get_beam_index => process_instance_get_beam_index
  procedure :: get_in_index => process_instance_get_in_index
\nwendcode{}\nwbegincode{888}\moddef{Processes: procedures}\plusendmoddef
  subroutine process_instance_get_beam_index (instance, i_term, i_beam)
    class(process_instance_t), intent(in) :: instance
    integer, intent(in) :: i_term
    integer, dimension(:), intent(out) :: i_beam
    call instance%term(i_term)%connected%get_beam_index (i_beam)
  end subroutine process_instance_get_beam_index
  
  subroutine process_instance_get_in_index (instance, i_term, i_in)
    class(process_instance_t), intent(in) :: instance
    integer, intent(in) :: i_term
    integer, dimension(:), intent(out) :: i_in
    call instance%term(i_term)%connected%get_in_index (i_in)
  end subroutine process_instance_get_in_index
  
\eatline
\eatline
\nwendcode{}\nwbegindocs{889}Return squared matrix element and event weight, and event weight
excess where applicable.
\nwenddocs{}\nwbegincode{890}\moddef{Processes: process instance: TBP}\plusendmoddef
  procedure :: get_sqme => process_instance_get_sqme
  procedure :: get_weight => process_instance_get_weight
  procedure :: get_excess => process_instance_get_excess
\nwendcode{}\nwbegincode{891}\moddef{Processes: procedures}\plusendmoddef
  function process_instance_get_sqme (instance) result (sqme)
    class(process_instance_t), intent(in) :: instance
    real(default) :: sqme
    if (instance%evaluation_status >= STAT_EVALUATED_TRACE) then
       sqme = instance%sqme
    else
       sqme = 0
    end if
  end function process_instance_get_sqme
  
  function process_instance_get_weight (instance) result (weight)
    class(process_instance_t), intent(in) :: instance
    real(default) :: weight
    if (instance%evaluation_status >= STAT_EVENT_COMPLETE) then
       weight = instance%weight
    else
       weight = 0
    end if
  end function process_instance_get_weight
  
  function process_instance_get_excess (instance) result (excess)
    class(process_instance_t), intent(in) :: instance
    real(default) :: excess
    if (instance%evaluation_status >= STAT_EVENT_COMPLETE) then
       excess = instance%excess
    else
       excess = 0
    end if
  end function process_instance_get_excess
  
\eatline
\eatline
\eatline
\nwendcode{}\nwbegindocs{892}Return the currently selected MCI channel.
\nwenddocs{}\nwbegincode{893}\moddef{Processes: process instance: TBP}\plusendmoddef
  procedure :: get_channel => process_instance_get_channel
\nwendcode{}\nwbegincode{894}\moddef{Processes: procedures}\plusendmoddef
  function process_instance_get_channel (instance) result (channel)
    class(process_instance_t), intent(in) :: instance
    integer :: channel
    channel = instance%selected_channel
  end function process_instance_get_channel

\eatline
\nwendcode{}\nwbegindocs{895}Return factorization scale and strong coupling.  We have to select a
term instance.
\nwenddocs{}\nwbegincode{896}\moddef{Processes: process instance: TBP}\plusendmoddef
  procedure :: get_fac_scale => process_instance_get_fac_scale
  procedure :: get_alpha_s => process_instance_get_alpha_s
\nwendcode{}\nwbegincode{897}\moddef{Processes: procedures}\plusendmoddef
  function process_instance_get_fac_scale (instance, i_term) result (fac_scale)
    class(process_instance_t), intent(in) :: instance
    integer, intent(in) :: i_term
    real(default) :: fac_scale
    fac_scale = instance%term(i_term)%get_fac_scale ()
  end function process_instance_get_fac_scale
  
  function process_instance_get_alpha_s (instance, i_term) result (alpha_s)
    class(process_instance_t), intent(in) :: instance
    integer, intent(in) :: i_term
    real(default) :: alpha_s
    alpha_s = instance%term(i_term)%get_alpha_s (instance%component)
  end function process_instance_get_alpha_s
  
\eatline
\eatline
\nwendcode{}\nwbegindocs{898}Counter.
\nwenddocs{}\nwbegincode{899}\moddef{Processes: process instance: TBP}\plusendmoddef
  procedure :: reset_counter => process_instance_reset_counter
  procedure :: record_call => process_instance_record_call
  procedure :: get_counter => process_instance_get_counter
\nwendcode{}\nwbegincode{900}\moddef{Processes: procedures}\plusendmoddef
  subroutine process_instance_reset_counter (process_instance)
    class(process_instance_t), intent(inout) :: process_instance
    call process_instance%mci_work(process_instance%i_mci)%reset_counter ()
  end subroutine process_instance_reset_counter
  
  subroutine process_instance_record_call (process_instance)
    class(process_instance_t), intent(inout) :: process_instance
    call process_instance%mci_work(process_instance%i_mci)%record_call &
         (process_instance%evaluation_status)
  end subroutine process_instance_record_call
    
  function process_instance_get_counter (process_instance) result (counter)
    class(process_instance_t), intent(in) :: process_instance
    type(process_counter_t) :: counter
    counter = process_instance%mci_work(process_instance%i_mci)%get_counter ()
  end function process_instance_get_counter
  
\eatline
\eatline
\eatline
\nwendcode{}\nwbegindocs{901}\nwdocspar
\subsection{Particle sets}
Here we provide two procedures that convert the process instance
from/to a particle set.  The conversion applies to the trace evaluator
which has no quantum-number information, thus it involves only the
momenta and the parent-child relations.  We keep virtual particles.

Nevertheless, it is possible to reconstruct the complete structure
from a particle set.  The reconstruction implies a re-evaluation of
the structure function and matrix-element codes.

The {\tt{}i{\char95}term} index is needed for both input and output, to select
among different active trace evaluators.

In both cases, the {\tt{}instance} object must be properly initialized.

NB: The {\tt{}recover{\char95}beams} option should be used only when the particle
set originates from an external event file, and the user has asked for
it.  It should be switched off when reading from raw event file.
\nwenddocs{}\nwbegincode{902}\moddef{Processes: process instance: TBP}\plusendmoddef
  procedure :: get_trace => process_instance_get_trace
  procedure :: set_trace => process_instance_set_trace
\nwendcode{}\nwbegincode{903}\moddef{Processes: procedures}\plusendmoddef
  subroutine process_instance_get_trace (instance, pset, i_term)
    class(process_instance_t), intent(in), target :: instance
    type(particle_set_t), intent(out) :: pset
    integer, intent(in) :: i_term
    type(interaction_t), pointer :: int
    logical :: ok
    int => instance%get_trace_int_ptr (i_term)
    call pset%init (ok, int, int, FM_IGNORE_HELICITY, &
         [0._default, 0._default], .false., .true.)
  end subroutine process_instance_get_trace
    
  subroutine process_instance_set_trace &
       (instance, pset, i_term, recover_beams, check_match)
    class(process_instance_t), intent(inout), target :: instance
    type(particle_set_t), intent(in) :: pset
    integer, intent(in) :: i_term
    logical, intent(in), optional :: recover_beams, check_match
    type(interaction_t), pointer :: int
    integer :: n_in
    int => instance%get_trace_int_ptr (i_term)
    n_in = instance%process%get_n_in ()
    call pset%fill_interaction (int, n_in, &
         recover_beams = recover_beams, &
         check_match = check_match, &
         state_flv = instance%get_state_flv (i_term))
  end subroutine process_instance_set_trace

\eatline
\eatline
\nwendcode{}\nwbegindocs{904}This procedure allows us to override any QCD setting of the WHIZARD process
and directly set the coupling value that comes together with a particle set.
\nwenddocs{}\nwbegincode{905}\moddef{Processes: process instance: TBP}\plusendmoddef
  procedure :: set_alpha_qcd_forced => process_instance_set_alpha_qcd_forced
\nwendcode{}\nwbegincode{906}\moddef{Processes: procedures}\plusendmoddef
  subroutine process_instance_set_alpha_qcd_forced (instance, i_term, alpha_qcd)
    class(process_instance_t), intent(inout) :: instance
    integer, intent(in) :: i_term
    real(default), intent(in) :: alpha_qcd
    call instance%term(i_term)%set_alpha_qcd_forced (alpha_qcd)
  end subroutine process_instance_set_alpha_qcd_forced
  
\eatline
\nwendcode{}\nwbegindocs{907}\nwdocspar
\nwenddocs{}\nwbegincode{908}\moddef{Processes: process instance: TBP}\plusendmoddef
  procedure :: display_real_kinematics => &
       process_instance_display_real_kinematics
\nwendcode{}\nwbegincode{909}\moddef{Processes: procedures}\plusendmoddef
  subroutine process_instance_display_real_kinematics (instance, i)
     class(process_instance_t), intent(in) :: instance
     integer, intent(in) :: i
     if (instance%term(i)%nlo_data%counter_exists) &
         call instance%term(i)%nlo_data%counter%display ()
  end subroutine process_instance_display_real_kinematics

\eatline
\nwendcode{}\nwbegindocs{910}\nwdocspar
\nwenddocs{}\nwbegincode{911}\moddef{Processes: process instance: TBP}\plusendmoddef
  procedure :: has_nlo_component => process_instance_has_nlo_component
\nwendcode{}\nwbegincode{912}\moddef{Processes: procedures}\plusendmoddef
  function process_instance_has_nlo_component (instance) result (nlo)
    class(process_instance_t), intent(in) :: instance
    logical :: nlo
    nlo = any (instance%component%get_nlo_type() /= 'Born')
  end function process_instance_has_nlo_component

\eatline
\nwendcode{}\nwbegindocs{913}\nwdocspar
\nwenddocs{}\nwbegincode{914}\moddef{Processes: process instance: TBP}\plusendmoddef
  procedure :: create_blha_interface => process_instance_create_blha_interface
\nwendcode{}\nwbegincode{915}\moddef{Processes: procedures}\plusendmoddef
  subroutine process_instance_create_blha_interface (instance)
    class(process_instance_t), intent(inout) :: instance
    logical :: use_gosam_loops
    logical :: use_gosam_correlations
    logical :: use_gosam_real_trees
    integer :: alpha_power, alphas_power
    integer, dimension(:,:), allocatable :: flv_born, flv_real
    integer :: i
    type(blha_master_t) :: blha_master

    associate (process => instance%process)
       associate (var_list => process%meta%var_list)
          use_gosam_loops  = var_list%get_lval (&
                                           var_str ('?use_gosam_loops'))
          use_gosam_correlations = var_list%get_lval (&
                                           var_str ('?use_gosam_correlations'))
          use_gosam_real_trees= var_list%get_lval (&
                                           var_str ('?use_gosam_real_trees'))
          alpha_power = var_list%get_ival (&
                                           var_str ('alpha_power'))
          alphas_power = var_list%get_ival (&
                                           var_str ('alphas_power'))
       end associate
       do i = 1, size (process%term)
          if (instance%component(i)%get_nlo_type() == 'Born') then
             flv_born = process%term(i)%data%flv_state
          else if (instance%component(i)%get_nlo_type() == 'Real') then
             flv_real = process%term(i)%data%flv_state
          end if
       end do
       call blha_master%init (process%meta%id, process%config%model, &
                              2, size (flv_born,1)-2, use_gosam_loops, &
                              use_gosam_correlations, use_gosam_real_trees, &
                              alpha_power, alphas_power, &
                              flv_born, flv_real)
       call blha_master%generate (process%meta%id)
    end associate 
  end subroutine process_instance_create_blha_interface 

\eatline
\nwendcode{}\nwbegindocs{916}\nwdocspar
\nwenddocs{}\nwbegincode{917}\moddef{Processes: process instance: TBP}\plusendmoddef
  procedure :: load_blha_libraries => process_instance_load_blha_libraries
\nwendcode{}\nwbegincode{918}\moddef{Processes: procedures}\plusendmoddef
  subroutine process_instance_load_blha_libraries (instance, os_data)
    class(process_instance_t), intent(inout) :: instance
    type(os_data_t), intent(in) :: os_data
    integer :: i
    do i = 1, size (instance%component)
       select type (core => instance%component(i)%config%core)
       type is (prc_gosam_t)
          call core%execute_olp_file (os_data)
          call core%read_olc_file (instance%process%term(i)%data%flv_state)
          call core%set_particle_properties (instance%process%config%model)
       end select
    end do
  end subroutine process_instance_load_blha_libraries

\eatline
\nwendcode{}\nwbegindocs{919}\nwdocspar
The {\tt{}pacify} subroutine has the purpose of setting numbers to zero
which are (by comparing with a {\tt{}tolerance} parameter) considered
equivalent with zero.  We do this in some unit tests.  Here, we a
apply this to the phase space subobject of the process instance.
\nwenddocs{}\nwbegincode{920}\moddef{Processes: public}\plusendmoddef
  public :: pacify
\nwendcode{}\nwbegincode{921}\moddef{Processes: interfaces}\endmoddef
  interface pacify
     module procedure pacify_process_instance
  end interface pacify
  
\nwendcode{}\nwbegincode{922}\moddef{Processes: procedures}\plusendmoddef
  subroutine pacify_process_instance (instance)
    type(process_instance_t), intent(inout) :: instance
    integer :: i
    do i = 1, size (instance%component)
       call pacify (instance%component(i)%k_seed%phs)
    end do
  end subroutine pacify_process_instance
    
\eatline
\nwendcode{}\nwbegindocs{923}\nwdocspar
\subsection{Unit tests}
\nwenddocs{}\nwbegincode{924}\moddef{Processes: public}\plusendmoddef
  public :: processes_test
\nwendcode{}\nwbegincode{925}\moddef{Processes: tests}\endmoddef
  subroutine processes_test (u, results)
    integer, intent(in) :: u
    type(test_results_t), intent(inout) :: results
  \LA{}Processes: execute tests\RA{}
  end subroutine processes_test
  
\eatline
\nwendcode{}\nwbegindocs{926}\nwdocspar
\subsubsection{Test process type}
For the following tests, we define a simple implementation of the abstract
{\tt{}prc{\char95}core{\char95}t}, designed such as to complement the {\tt{}prc{\char95}test{\char95}t}
process definition type.

Note that it is not given that the actual process is defined as
{\tt{}prc{\char95}test{\char95}t} type.  We enforce this by calling
{\tt{}prc{\char95}test{\char95}create{\char95}library}.  The driver component in the process core
will then become of type {\tt{}prc{\char95}test{\char95}t}.
\nwenddocs{}\nwbegincode{927}\moddef{Processes: public}\plusendmoddef
  public :: test_t
\nwendcode{}\nwbegincode{928}\moddef{Processes: test types}\endmoddef
  type, extends (prc_core_t) :: test_t
   contains
     procedure :: write => test_write
     procedure :: needs_mcset => test_needs_mcset
     procedure :: get_n_terms => test_get_n_terms
     procedure :: is_allowed => test_is_allowed
     procedure :: compute_hard_kinematics => test_compute_hard_kinematics
     procedure :: compute_eff_kinematics => test_compute_eff_kinematics
     procedure :: recover_kinematics => test_recover_kinematics
     procedure :: compute_amplitude => test_compute_amplitude
  end type test_t

\eatline
\nwendcode{}\nwbegincode{929}\moddef{Processes: tests}\plusendmoddef
  subroutine test_write (object, unit)
    class(test_t), intent(in) :: object
    integer, intent(in), optional :: unit
    integer :: u
    u = given_output_unit (unit)
    write (u, "(3x,A)")  "test type implementing prc_test"
  end subroutine test_write
  
\eatline
\nwendcode{}\nwbegindocs{930}This process type always needs a MC parameter set and a
single term.  This only state is always allowed.
\nwenddocs{}\nwbegincode{931}\moddef{Processes: tests}\plusendmoddef
  function test_needs_mcset (object) result (flag)
    class(test_t), intent(in) :: object
    logical :: flag
    flag = .true.
  end function test_needs_mcset
  
  function test_get_n_terms (object) result (n)
    class(test_t), intent(in) :: object
    integer :: n
    n = 1
  end function test_get_n_terms
  
  function test_is_allowed (object, i_term, f, h, c) result (flag)
    class(test_t), intent(in) :: object
    integer, intent(in) :: i_term, f, h, c
    logical :: flag
    flag = .true.
  end function test_is_allowed
  
\eatline
\eatline
\eatline
\nwendcode{}\nwbegindocs{932}Transfer the generated momenta directly to the hard interaction in
the (only) term.  We assume that everything has been set up correctly,
so the array fits.
\nwenddocs{}\nwbegincode{933}\moddef{Processes: tests}\plusendmoddef
  subroutine test_compute_hard_kinematics &
       (object, p_seed, i_term, int_hard, tmp)
    class(test_t), intent(in) :: object
    type(vector4_t), dimension(:), intent(in) :: p_seed
    integer, intent(in) :: i_term
    type(interaction_t), intent(inout) :: int_hard
    class(workspace_t), intent(inout), allocatable :: tmp
    call interaction_set_momenta (int_hard, p_seed)
  end subroutine test_compute_hard_kinematics
  
\eatline
\nwendcode{}\nwbegindocs{934}This procedure is not called for {\tt{}test{\char95}t}, just a placeholder.
\nwenddocs{}\nwbegincode{935}\moddef{Processes: tests}\plusendmoddef
  subroutine test_compute_eff_kinematics &
       (object, i_term, int_hard, int_eff, tmp)
    class(test_t), intent(in) :: object
    integer, intent(in) :: i_term
    type(interaction_t), intent(in) :: int_hard
    type(interaction_t), intent(inout) :: int_eff
    class(workspace_t), intent(inout), allocatable :: tmp
  end subroutine test_compute_eff_kinematics
  
\eatline
\nwendcode{}\nwbegindocs{936}Transfer the incoming momenta of {\tt{}p{\char95}seed} directly to the
effective interaction, and vice versa for the outgoing momenta.

{\tt{}int{\char95}hard} is left untouched since {\tt{}int{\char95}eff} is an alias (via
pointer) to it.
\nwenddocs{}\nwbegincode{937}\moddef{Processes: tests}\plusendmoddef
  subroutine test_recover_kinematics &
       (object, p_seed, int_hard, int_eff, tmp)
    class(test_t), intent(in) :: object
    type(vector4_t), dimension(:), intent(inout) :: p_seed
    type(interaction_t), intent(inout) :: int_hard
    type(interaction_t), intent(inout) :: int_eff
    class(workspace_t), intent(inout), allocatable :: tmp
    integer :: n_in
    n_in = interaction_get_n_in (int_eff)
    call interaction_set_momenta (int_eff, p_seed(1:n_in), outgoing = .false.)
    p_seed(n_in+1:) = interaction_get_momenta (int_eff, outgoing = .true.)
  end subroutine test_recover_kinematics
    
\eatline
\nwendcode{}\nwbegindocs{938}Compute the amplitude.  The driver ignores all quantum numbers and,
in fact, returns a constant.  Nevertheless, we properly transfer the
momentum vectors.
\nwenddocs{}\nwbegincode{939}\moddef{Processes: tests}\plusendmoddef
  function test_compute_amplitude &
       (object, j, p, f, h, c, fac_scale, ren_scale, alpha_qcd_forced, tmp) &
       result (amp)
    class(test_t), intent(in) :: object
    integer, intent(in) :: j
    type(vector4_t), dimension(:), intent(in) :: p
    integer, intent(in) :: f, h, c
    real(default), intent(in) :: fac_scale, ren_scale
    real(default), intent(in), allocatable :: alpha_qcd_forced
    class(workspace_t), intent(inout), allocatable, optional :: tmp
    complex(default) :: amp
    real(default), dimension(:,:), allocatable :: parray
    integer :: i, n_tot
    select type (driver => object%driver)
    type is (prc_test_t)
       if (driver%scattering) then
          n_tot = 4
       else
          n_tot = 3
       end if
       allocate (parray (0:3,n_tot))
       forall (i = 1:n_tot)  parray(:,i) = vector4_get_components (p(i))
       amp = driver%get_amplitude (parray)
    end select
  end function test_compute_amplitude
    
\eatline
\nwendcode{}\nwbegindocs{940}\nwdocspar
\subsubsection{Write an empty process object}
The most trivial test is to write an uninitialized process object.
\nwenddocs{}\nwbegincode{941}\moddef{Processes: execute tests}\endmoddef
  call test (processes_1, "processes_1", &
       "write an empty process object", &
       u, results)
\nwendcode{}\nwbegincode{942}\moddef{Processes: tests}\plusendmoddef
  subroutine processes_1 (u)
    integer, intent(in) :: u
    type(process_t) :: process

    write (u, "(A)")  "* Test output: processes_1"
    write (u, "(A)")  "*   Purpose: display an empty process object"
    write (u, "(A)")

    call process%write (.false., u)
    
    write (u, "(A)")
    write (u, "(A)")  "* Test output end: processes_1"
    
  end subroutine processes_1
  
\eatline
\nwendcode{}\nwbegindocs{943}\nwdocspar
\subsubsection{Initialize a process object}
Initialize a process and display it.
\nwenddocs{}\nwbegincode{944}\moddef{Processes: execute tests}\plusendmoddef
  call test (processes_2, "processes_2", &
       "initialize a simple process object", &
       u, results)
\nwendcode{}\nwbegincode{945}\moddef{Processes: tests}\plusendmoddef
  subroutine processes_2 (u)
    integer, intent(in) :: u
    type(process_library_t), target :: lib
    type(string_t) :: libname
    type(string_t) :: procname
    type(string_t) :: run_id
    type(os_data_t) :: os_data
    type(qcd_t) :: qcd
    class(rng_factory_t), allocatable :: rng_factory
    class(model_data_t), pointer :: model
    type(process_t), allocatable :: process
    class(prc_core_t), allocatable :: core_template
    class(mci_t), allocatable :: mci_template
    class(phs_config_t), allocatable :: phs_config_template

    write (u, "(A)")  "* Test output: processes_2"
    write (u, "(A)")  "*   Purpose: initialize a simple process object"
    write (u, "(A)")

    write (u, "(A)")  "* Build and load a test library with one process"
    write (u, "(A)")

    libname = "processes2"
    procname = libname
    run_id = "run2"
    call os_data_init (os_data)
    allocate (rng_test_factory_t :: rng_factory)
    call prc_test_create_library (libname, lib)

    write (u, "(A)")  "* Initialize a process object"
    write (u, "(A)")

    allocate (model)
    call model%init_test ()

    allocate (process)
    call process%init (procname, run_id, lib, os_data, qcd, rng_factory, model)
    
    allocate (test_t :: core_template)
    allocate (phs_test_config_t :: phs_config_template)
    call process%init_component &
         (1, core_template, mci_template, phs_config_template)

    call process%setup_mci ()

    call process%write (.false., u)
    
    write (u, "(A)")
    write (u, "(A)")  "* Cleanup"

    call process%final ()
    deallocate (process)
    
    write (u, "(A)")
    write (u, "(A)")  "* Test output end: processes_2"
    
  end subroutine processes_2
  
\eatline
\nwendcode{}\nwbegindocs{946}\nwdocspar
\subsubsection{Compute a trivial matrix element}
Initialize a process, retrieve some information and compute a matrix
element.

We use the same trivial process as for the previous test.  All
momentum and state dependence is trivial, so we just test basic
functionality. 
\nwenddocs{}\nwbegincode{947}\moddef{Processes: execute tests}\plusendmoddef
  call test (processes_3, "processes_3", &
       "retrieve a trivial matrix element", &
       u, results)
\nwendcode{}\nwbegincode{948}\moddef{Processes: tests}\plusendmoddef
  subroutine processes_3 (u)
    integer, intent(in) :: u
    type(process_library_t), target :: lib
    type(string_t) :: libname
    type(string_t) :: procname
    type(string_t) :: run_id
    type(os_data_t) :: os_data
    type(qcd_t) :: qcd
    class(rng_factory_t), allocatable :: rng_factory
    class(model_data_t), pointer :: model
    type(process_t), allocatable :: process
    class(prc_core_t), allocatable :: core_template
    class(mci_t), allocatable :: mci_template
    class(phs_config_t), allocatable :: phs_config_template
    type(process_constants_t) :: data
    type(vector4_t), dimension(:), allocatable :: p

    write (u, "(A)")  "* Test output: processes_3"
    write (u, "(A)")  "*   Purpose: create a process &
         &and compute a matrix element"
    write (u, "(A)")

    write (u, "(A)")  "* Build and load a test library with one process"
    write (u, "(A)")

    libname = "processes3"
    procname = libname
    run_id = "run3"
    call os_data_init (os_data)
    allocate (rng_test_factory_t :: rng_factory)
    call prc_test_create_library (libname, lib)

    allocate (model)
    call model%init_test ()

    allocate (process)
    call process%init (procname, run_id, lib, os_data, qcd, rng_factory, model)
    
    allocate (test_t :: core_template)
    allocate (mci_test_t :: mci_template)
    select type (mci_template)
    type is (mci_test_t)
       call mci_template%set_dimensions (2, 2)
       call mci_template%set_divisions (100)
    end select
    allocate (phs_test_config_t :: phs_config_template)
    call process%init_component &
         (1, core_template, mci_template, phs_config_template)

    write (u, "(A)")  "* Return the number of process components"
    write (u, "(A)")

    write (u, "(A,I0)")  "n_components = ", process%get_n_components ()
    
    write (u, "(A)")
    write (u, "(A)")  "* Return the number of flavor states"
    write (u, "(A)")

    data = process%get_constants (1)
    
    write (u, "(A,I0)")  "n_flv(1) = ", data%n_flv
    
    write (u, "(A)")
    write (u, "(A)")  "* Return the first flavor state"
    write (u, "(A)")

    write (u, "(A,4(1x,I0))")  "flv_state(1) =", data%flv_state (:,1)

    write (u, "(A)")
    write (u, "(A)")  "* Set up kinematics &
         &[arbitrary, the matrix element is constant]"
    
    allocate (p (4))

    write (u, "(A)")
    write (u, "(A)")  "* Retrieve the matrix element"
    write (u, "(A)")

    write (u, "(A,F5.3,' + ',F5.3,' I')")  "me (1, p, 1, 1, 1) = ", &
         process%compute_amplitude (1, 1, p, 1, 1, 1)
    

    write (u, "(A)")
    write (u, "(A)")  "* Cleanup"
    
    call process%final ()
    deallocate (process)

    write (u, "(A)")
    write (u, "(A)")  "* Test output end: processes_3"
    
  end subroutine processes_3
  
\eatline
\nwendcode{}\nwbegindocs{949}\nwdocspar
\subsubsection{Generate a process instance}
Initialize a process and process instance, choose a sampling point and
fill the process instance.

We use the same trivial process as for the previous test.  All
momentum and state dependence is trivial, so we just test basic
functionality. 
\nwenddocs{}\nwbegincode{950}\moddef{Processes: execute tests}\plusendmoddef
  call test (processes_4, "processes_4", &
       "create and fill a process instance (partonic event)", &
       u, results)
\nwendcode{}\nwbegincode{951}\moddef{Processes: tests}\plusendmoddef
  subroutine processes_4 (u)
    integer, intent(in) :: u
    type(process_library_t), target :: lib
    type(string_t) :: libname
    type(string_t) :: procname
    type(string_t) :: run_id
    type(os_data_t) :: os_data
    type(qcd_t) :: qcd
    class(rng_factory_t), allocatable :: rng_factory
    class(model_data_t), pointer :: model
    type(process_t), allocatable, target :: process
    class(prc_core_t), allocatable :: core_template
    class(mci_t), allocatable :: mci_template
    class(phs_config_t), allocatable :: phs_config_template
    real(default) :: sqrts
    type(process_instance_t), allocatable, target :: process_instance
    type(particle_set_t) :: pset

    write (u, "(A)")  "* Test output: processes_4"
    write (u, "(A)")  "*   Purpose: create a process &
         &and fill a process instance"
    write (u, "(A)")

    write (u, "(A)")  "* Build and initialize a test process"
    write (u, "(A)")

    libname = "processes4"
    procname = libname
    run_id = "run4"
    call os_data_init (os_data)
    allocate (rng_test_factory_t :: rng_factory)
    call prc_test_create_library (libname, lib)

    call reset_interaction_counter ()

    allocate (model)
    call model%init_test ()

    allocate (process)
    call process%init (procname, run_id, lib, os_data, qcd, rng_factory, model) 
   
    allocate (test_t :: core_template)
    allocate (phs_test_config_t :: phs_config_template)
    call process%init_component &
         (1, core_template, mci_template, phs_config_template)

    write (u, "(A)")  "* Prepare a trivial beam setup"
    write (u, "(A)")
    
    sqrts = 1000
    call process%setup_beams_sqrts (sqrts)
    call process%configure_phs ()
    call process%setup_mci ()

    write (u, "(A)")  "* Complete process initialization"
    write (u, "(A)")

    call process%setup_terms ()
    call process%write (.false., u)

    write (u, "(A)")
    write (u, "(A)")  "* Create a process instance"
    write (u, "(A)")

    allocate (process_instance)
    call process_instance%init (process)
    call process_instance%write (u)

    write (u, "(A)")
    write (u, "(A)")  "* Inject a set of random numbers"
    write (u, "(A)")
     
    call process_instance%choose_mci (1)
    call process_instance%set_mcpar ([0._default, 0._default])
    call process_instance%write (u)

    write (u, "(A)")
    write (u, "(A)")  "* Set up hard kinematics"
    write (u, "(A)")
  
    call process_instance%select_channel (1)
    call process_instance%compute_seed_kinematics ()
    call process_instance%compute_hard_kinematics ()
    call process_instance%compute_eff_kinematics ()
    call process_instance%evaluate_expressions ()
    call process_instance%compute_other_channels ()

    write (u, "(A)")  "* Evaluate matrix element and square"
    write (u, "(A)")
  
    call process_instance%evaluate_trace ()
    call process_instance%write (u)

    call process_instance%get_trace (pset, 1)
    call process_instance%final ()
    deallocate (process_instance)

    write (u, "(A)")
    write (u, "(A)")  "* Particle content:"
    write (u, "(A)")
    
    call write_separator (u)
    call pset%write (u)
    call write_separator (u)

    write (u, "(A)")
    write (u, "(A)")  "* Recover process instance"
    write (u, "(A)")

    allocate (process_instance)
    call process_instance%init (process)
    call process_instance%choose_mci (1)
    call process_instance%set_trace (pset, 1, check_match = .false.)

    call process_instance%activate ()
    process_instance%evaluation_status = STAT_EFF_KINEMATICS
    call process_instance%recover_hard_kinematics (i_term = 1)
    call process_instance%recover_seed_kinematics (i_term = 1)
    call process_instance%select_channel (1)
    call process_instance%recover_mcpar (i_term = 1)

    call process_instance%compute_seed_kinematics (skip_term = 1)
    call process_instance%compute_hard_kinematics (skip_term = 1)
    call process_instance%compute_eff_kinematics (skip_term = 1)

    call process_instance%evaluate_expressions ()
    call process_instance%compute_other_channels (skip_term = 1)
    call process_instance%evaluate_trace ()
    call process_instance%write (u)

    write (u, "(A)")
    write (u, "(A)")  "* Cleanup"
    
    call pset%final ()
    call process_instance%final ()
    deallocate (process_instance)
    
    call process%final ()
    deallocate (process)

    write (u, "(A)")
    write (u, "(A)")  "* Test output end: processes_4"
    
  end subroutine processes_4
  
\eatline
\nwendcode{}\nwbegindocs{952}\nwdocspar
\subsubsection{Structure function configuration}
Configure structure functions (multi-channel) in a process object.
\nwenddocs{}\nwbegincode{953}\moddef{Processes: execute tests}\plusendmoddef
  call test (processes_7, "processes_7", &
       "process configuration with structure functions", &
       u, results)
\nwendcode{}\nwbegincode{954}\moddef{Processes: tests}\plusendmoddef
  subroutine processes_7 (u)
    integer, intent(in) :: u
    type(process_library_t), target :: lib
    type(string_t) :: libname
    type(string_t) :: procname
    type(string_t) :: run_id
    type(os_data_t) :: os_data
    type(qcd_t) :: qcd
    class(rng_factory_t), allocatable :: rng_factory
    class(model_data_t), pointer :: model
    type(process_t), allocatable, target :: process
    class(prc_core_t), allocatable :: core_template
    class(mci_t), allocatable :: mci_template
    class(phs_config_t), allocatable :: phs_config_template
    real(default) :: sqrts
    type(pdg_array_t) :: pdg_in
    class(sf_data_t), allocatable, target :: data
    type(sf_config_t), dimension(:), allocatable :: sf_config
    type(sf_channel_t), dimension(2) :: sf_channel

    write (u, "(A)")  "* Test output: processes_7"
    write (u, "(A)")  "*   Purpose: initialize a process with &
         &structure functions"
    write (u, "(A)")

    write (u, "(A)")  "* Build and initialize a process object"
    write (u, "(A)")

    libname = "processes7"
    procname = libname
    run_id = "run7"
    call os_data_init (os_data)
    allocate (rng_test_factory_t :: rng_factory)
    call prc_test_create_library (libname, lib)

    allocate (model)
    call model%init_test ()

    allocate (process)
    call process%init (procname, run_id, lib, os_data, qcd, rng_factory, model)
    
    allocate (test_t :: core_template)
    allocate (phs_test_config_t :: phs_config_template)
    call process%init_component &
         (1, core_template, mci_template, phs_config_template)

    write (u, "(A)")  "* Set beam, structure functions, and mappings"
    write (u, "(A)")

    sqrts = 1000
    call process%setup_beams_sqrts (sqrts)
    call process%configure_phs ()
    
    pdg_in = 25
    allocate (sf_test_data_t :: data)
    select type (data)
    type is (sf_test_data_t)
       model => process%get_model_ptr ()
       call data%init (model, pdg_in)
    end select

    allocate (sf_config (2))
    call sf_config(1)%init ([1], data)
    call sf_config(2)%init ([2], data)
    call process%init_sf_chain (sf_config)
    deallocate (sf_config)

    call process%beam_config%allocate_sf_channels (3)

    call sf_channel(1)%init (2)
    call sf_channel(1)%activate_mapping ([1,2])
    call process%set_sf_channel (2, sf_channel(1))
    
    call sf_channel(2)%init (2)
    call sf_channel(2)%set_s_mapping ([1,2])
    call process%set_sf_channel (3, sf_channel(2))
    
    call process%setup_mci ()

    call process%write (.false., u)
    
    write (u, "(A)")
    write (u, "(A)")  "* Cleanup"

    call process%final ()
    deallocate (process)
    
    write (u, "(A)")
    write (u, "(A)")  "* Test output end: processes_7"
    
  end subroutine processes_7
  
\eatline
\nwendcode{}\nwbegindocs{955}\nwdocspar
\subsubsection{Evaluating a process with structure function}
Configure structure functions (single-channel) in a process object,
create an instance, compute kinematics and evaluate.

Note the order of operations when setting up structure functions and
phase space.  The beams are first, they determine the {\tt{}sqrts} value.
We can also set up the chain of structure functions.  We then
configure the phase space.  From this, we can obtain information about
special configurations (resonances, etc.), which we need for
allocating the possible structure-function channels (parameterizations
and mappings).  Finally, we match phase-space channels onto
structure-function channels.

In the current example, this matching is trivial; we only have one
structure-function channel.
\nwenddocs{}\nwbegincode{956}\moddef{Processes: execute tests}\plusendmoddef
  call test (processes_8, "processes_8", &
       "process evaluation with structure functions", &
       u, results)
\nwendcode{}\nwbegincode{957}\moddef{Processes: tests}\plusendmoddef
  subroutine processes_8 (u)
    integer, intent(in) :: u
    type(process_library_t), target :: lib
    type(string_t) :: libname
    type(string_t) :: procname
    type(string_t) :: run_id
    type(os_data_t) :: os_data
    type(qcd_t) :: qcd
    class(rng_factory_t), allocatable :: rng_factory
    class(model_data_t), pointer :: model
    type(process_t), allocatable, target :: process
    class(prc_core_t), allocatable :: core_template
    class(mci_t), allocatable :: mci_template
    class(phs_config_t), allocatable :: phs_config_template
    real(default) :: sqrts
    type(process_instance_t), allocatable, target :: process_instance
    type(pdg_array_t) :: pdg_in
    class(sf_data_t), allocatable, target :: data
    type(sf_config_t), dimension(:), allocatable :: sf_config
    type(sf_channel_t) :: sf_channel
    type(particle_set_t) :: pset

    write (u, "(A)")  "* Test output: processes_8"
    write (u, "(A)")  "*   Purpose: evaluate a process with &
         &structure functions"
    write (u, "(A)")

    write (u, "(A)")  "* Build and initialize a process object"
    write (u, "(A)")

    libname = "processes8"
    procname = libname
    run_id = "run8"
    call os_data_init (os_data)
    allocate (rng_test_factory_t :: rng_factory)
    call prc_test_create_library (libname, lib)

    call reset_interaction_counter ()

    allocate (model)
    call model%init_test ()

    allocate (process)
    call process%init (procname, run_id, lib, os_data, qcd, rng_factory, model)
    
    allocate (test_t :: core_template)
    allocate (phs_test_config_t :: phs_config_template)
    call process%init_component &
         (1, core_template, mci_template, phs_config_template)

    write (u, "(A)")  "* Set beam, structure functions, and mappings"
    write (u, "(A)")

    sqrts = 1000
    call process%setup_beams_sqrts (sqrts)

    pdg_in = 25
    allocate (sf_test_data_t :: data)
    select type (data)
    type is (sf_test_data_t)
       model => process%get_model_ptr ()
       call data%init (model, pdg_in)
    end select

    allocate (sf_config (2))
    call sf_config(1)%init ([1], data)
    call sf_config(2)%init ([2], data)
    call process%init_sf_chain (sf_config)
    deallocate (sf_config)
    
    call process%configure_phs ()
    
    call process%beam_config%allocate_sf_channels (1)

    call sf_channel%init (2)
    call sf_channel%activate_mapping ([1,2])
    call process%set_sf_channel (1, sf_channel)
    
    write (u, "(A)")  "* Complete process initialization"
    write (u, "(A)")

    call process%setup_mci ()
    call process%setup_terms ()

    call process%write (.false., u)
    
    write (u, "(A)")
    write (u, "(A)")  "* Create a process instance"
    write (u, "(A)")

    allocate (process_instance)
    call process_instance%init (process)

    write (u, "(A)")  "* Set up kinematics and evaluate"
    write (u, "(A)")

    call process_instance%choose_mci (1)
    call process_instance%evaluate_sqme (1, &
         [0.8_default, 0.8_default, 0.1_default, 0.2_default])
    call process_instance%write (u)

    call process_instance%get_trace (pset, 1)
    call process_instance%final ()
    deallocate (process_instance)

    write (u, "(A)")
    write (u, "(A)")  "* Particle content:"
    write (u, "(A)")
    
    call write_separator (u)
    call pset%write (u)
    call write_separator (u)

    write (u, "(A)")
    write (u, "(A)")  "* Recover process instance"
    write (u, "(A)")

    call reset_interaction_counter (2)

    allocate (process_instance)
    call process_instance%init (process)

    call process_instance%choose_mci (1)
    call process_instance%set_trace (pset, 1, check_match = .false.)
    call process_instance%recover &
         (channel = 1, i_term = 1, update_sqme = .true.)
    call process_instance%write (u)

    write (u, "(A)")
    write (u, "(A)")  "* Cleanup"

    call pset%final ()

    call process_instance%final ()
    deallocate (process_instance)
    
    call process%final ()
    deallocate (process)

    write (u, "(A)")
    write (u, "(A)")  "* Test output end: processes_8"
    
  end subroutine processes_8
  
\eatline
\nwendcode{}\nwbegindocs{958}\nwdocspar
\subsubsection{Multi-channel phase space and structure function}
This is an extension of the previous example.  This time, we have two
distinct structure-function channels which are matched to the two
distinct phase-space channels.
\nwenddocs{}\nwbegincode{959}\moddef{Processes: execute tests}\plusendmoddef
  call test (processes_9, "processes_9", &
       "multichannel kinematics and structure functions", &
       u, results)
\nwendcode{}\nwbegincode{960}\moddef{Processes: tests}\plusendmoddef
  subroutine processes_9 (u)
    integer, intent(in) :: u
    type(process_library_t), target :: lib
    type(string_t) :: libname
    type(string_t) :: procname
    type(string_t) :: run_id
    type(os_data_t) :: os_data
    type(qcd_t) :: qcd
    class(rng_factory_t), allocatable :: rng_factory
    class(model_data_t), pointer :: model
    type(process_t), allocatable, target :: process
    class(prc_core_t), allocatable :: core_template
    class(mci_t), allocatable :: mci_template
    class(phs_config_t), allocatable :: phs_config_template
    real(default) :: sqrts
    type(process_instance_t), allocatable, target :: process_instance
    type(pdg_array_t) :: pdg_in
    class(sf_data_t), allocatable, target :: data
    type(sf_config_t), dimension(:), allocatable :: sf_config
    type(sf_channel_t) :: sf_channel
    real(default), dimension(4) :: x_saved
    type(particle_set_t) :: pset

    write (u, "(A)")  "* Test output: processes_9"
    write (u, "(A)")  "*   Purpose: evaluate a process with &
         &structure functions"
    write (u, "(A)")  "*            in a multi-channel configuration"
    write (u, "(A)")

    write (u, "(A)")  "* Build and initialize a process object"
    write (u, "(A)")

    libname = "processes9"
    procname = libname
    run_id = "run9"
    call os_data_init (os_data)
    allocate (rng_test_factory_t :: rng_factory)
    call prc_test_create_library (libname, lib)

    call reset_interaction_counter ()

    allocate (model)
    call model%init_test ()

    allocate (process)
    call process%init (procname, run_id, lib, os_data, qcd, rng_factory, model)
    
    allocate (test_t :: core_template)
    allocate (phs_test_config_t :: phs_config_template)
    call process%init_component &
         (1, core_template, mci_template, phs_config_template)

    write (u, "(A)")  "* Set beam, structure functions, and mappings"
    write (u, "(A)")

    sqrts = 1000
    call process%setup_beams_sqrts (sqrts)

    pdg_in = 25
    allocate (sf_test_data_t :: data)
    select type (data)
    type is (sf_test_data_t)
       model => process%get_model_ptr ()
       call data%init (model, pdg_in)
    end select

    allocate (sf_config (2))
    call sf_config(1)%init ([1], data)
    call sf_config(2)%init ([2], data)
    call process%init_sf_chain (sf_config)
    deallocate (sf_config)
    
    call process%configure_phs ()
    
    call process%beam_config%allocate_sf_channels (2)

    call sf_channel%init (2)
    call process%set_sf_channel (1, sf_channel)
    
    call sf_channel%init (2)
    call sf_channel%activate_mapping ([1,2])
    call process%set_sf_channel (2, sf_channel)
    
    call process%component(1)%phs_config%set_sf_channel ([1, 2])

    write (u, "(A)")  "* Complete process initialization"
    write (u, "(A)")

    call process%setup_mci ()
    call process%setup_terms ()

    call process%write (.false., u)
    
    write (u, "(A)")
    write (u, "(A)")  "* Create a process instance"
    write (u, "(A)")

    allocate (process_instance)
    call process_instance%init (process)

    write (u, "(A)")  "* Set up kinematics in channel 1 and evaluate"
    write (u, "(A)")

    call process_instance%choose_mci (1)
    call process_instance%evaluate_sqme (1, &
         [0.8_default, 0.8_default, 0.1_default, 0.2_default])
    call process_instance%write (u)

    write (u, "(A)")
    write (u, "(A)")  "* Extract MC input parameters"
    write (u, "(A)")
    
    write (u, "(A)")  "Channel 1:"
    call process_instance%get_mcpar (1, x_saved)
    write (u, "(2x,9(1x,F7.5))")  x_saved

    write (u, "(A)")  "Channel 2:"
    call process_instance%get_mcpar (2, x_saved)
    write (u, "(2x,9(1x,F7.5))")  x_saved

    write (u, "(A)")
    write (u, "(A)")  "* Set up kinematics in channel 2 and evaluate"
    write (u, "(A)")

    call process_instance%evaluate_sqme (2, x_saved)
    call process_instance%write (u)

    call process_instance%get_trace (pset, 1)
    call process_instance%final ()
    deallocate (process_instance)

    write (u, "(A)")
    write (u, "(A)")  "* Recover process instance for channel 2"
    write (u, "(A)")

    call reset_interaction_counter (2)

    allocate (process_instance)
    call process_instance%init (process)

    call process_instance%choose_mci (1)
    call process_instance%set_trace (pset, 1, check_match = .false.)
    call process_instance%recover &
         (channel = 2, i_term = 1, update_sqme = .true.)
    call process_instance%write (u)

    write (u, "(A)")
    write (u, "(A)")  "* Cleanup"

    call pset%final ()

    call process_instance%final ()
    deallocate (process_instance)
    
    call process%final ()
    deallocate (process)

    write (u, "(A)")
    write (u, "(A)")  "* Test output end: processes_9"
    
  end subroutine processes_9
  
\eatline
\nwendcode{}\nwbegindocs{961}\nwdocspar
\subsubsection{Event generation}
Activate the MC integrator for the process object and use it to
generate a single event.  Note that the test integrator does not
require integration in preparation for generating events.
\nwenddocs{}\nwbegincode{962}\moddef{Processes: execute tests}\plusendmoddef
  call test (processes_10, "processes_10", &
       "event generation", &
       u, results)
\nwendcode{}\nwbegincode{963}\moddef{Processes: tests}\plusendmoddef
  subroutine processes_10 (u)
    integer, intent(in) :: u
    type(process_library_t), target :: lib
    type(string_t) :: libname
    type(string_t) :: procname
    type(string_t) :: run_id
    type(os_data_t) :: os_data
    type(qcd_t) :: qcd
    class(rng_factory_t), allocatable :: rng_factory
    class(model_data_t), pointer :: model
    type(process_t), allocatable, target :: process
    class(prc_core_t), allocatable :: core_template
    class(mci_t), allocatable :: mci_template
    class(phs_config_t), allocatable :: phs_config_template
    real(default) :: sqrts
    type(process_instance_t), allocatable, target :: process_instance

    write (u, "(A)")  "* Test output: processes_10"
    write (u, "(A)")  "*   Purpose: generate events for a process without &
         &structure functions"
    write (u, "(A)")  "*            in a multi-channel configuration"
    write (u, "(A)")

    write (u, "(A)")  "* Build and initialize a process object"
    write (u, "(A)")

    libname = "processes10"
    procname = libname
    run_id = "run10"
    call os_data_init (os_data)
    allocate (rng_test_factory_t :: rng_factory)
    call prc_test_create_library (libname, lib)

    call reset_interaction_counter ()

    allocate (model)
    call model%init_test ()

    allocate (process)
    call process%init (procname, run_id, lib, os_data, qcd, rng_factory, model)
    
    allocate (test_t :: core_template)
    allocate (mci_test_t :: mci_template)
    select type (mci_template)
    type is (mci_test_t);  call mci_template%set_divisions (100)
    end select
    allocate (phs_test_config_t :: phs_config_template)
    call process%init_component &
         (1, core_template, mci_template, phs_config_template)

    write (u, "(A)")  "* Prepare a trivial beam setup"
    write (u, "(A)")

    sqrts = 1000
    call process%setup_beams_sqrts (sqrts)
    call process%configure_phs ()

    call process%setup_mci ()
    
    write (u, "(A)")  "* Complete process initialization"
    write (u, "(A)")

    call process%setup_terms ()
    call process%write (.false., u)
    
    write (u, "(A)")
    write (u, "(A)")  "* Create a process instance"
    write (u, "(A)")

    allocate (process_instance)
    call process_instance%init (process)

    write (u, "(A)")  "* Generate weighted event"
    write (u, "(A)")

    select type (mci => process%mci_entry(1)%mci)
    type is (mci_test_t)
       ! This ensures that the next 'random' numbers are 0.3, 0.5, 0.7
       call mci%rng%init (3)
       ! Include the constant PHS factor in the stored maximum of the integrand
       call mci%set_max_factor (conv * twopi4 &
            / (2 * sqrt (lambda (sqrts **2, 125._default**2, 125._default**2))))
    end select

    call process%generate_weighted_event (process_instance, 1)
    call process_instance%write (u)

    write (u, "(A)")
    write (u, "(A)")  "* Generate unweighted event"
    write (u, "(A)")

    call process%generate_unweighted_event (process_instance, 1)
    select type (mci => process%mci_entry(1)%mci)
    type is (mci_test_t)
       write (u, "(A,I0)")  " Success in try ", mci%tries
       write (u, "(A)")
    end select
    
    call process_instance%write (u)

    write (u, "(A)")
    write (u, "(A)")  "* Cleanup"
    
    call process_instance%final ()
    deallocate (process_instance)
    
    call process%final ()
    deallocate (process)

    write (u, "(A)")
    write (u, "(A)")  "* Test output end: processes_10"
    
  end subroutine processes_10
  
\eatline
\nwendcode{}\nwbegindocs{964}\nwdocspar
\subsubsection{Integration}
Activate the MC integrator for the process object and use it to
integrate over phase space.
\nwenddocs{}\nwbegincode{965}\moddef{Processes: execute tests}\plusendmoddef
  call test (processes_11, "processes_11", &
       "integration", &
       u, results)
\nwendcode{}\nwbegincode{966}\moddef{Processes: tests}\plusendmoddef
  subroutine processes_11 (u)
    integer, intent(in) :: u
    type(process_library_t), target :: lib
    type(string_t) :: libname
    type(string_t) :: procname
    type(string_t) :: run_id
    type(os_data_t) :: os_data
    type(qcd_t) :: qcd
    class(rng_factory_t), allocatable :: rng_factory
    class(model_data_t), pointer :: model
    type(process_t), allocatable, target :: process
    class(prc_core_t), allocatable :: core_template
    class(mci_t), allocatable :: mci_template
    class(phs_config_t), allocatable :: phs_config_template
    real(default) :: sqrts
    type(process_instance_t), allocatable, target :: process_instance

    write (u, "(A)")  "* Test output: processes_11"
    write (u, "(A)")  "*   Purpose: integrate a process without &
         &structure functions"
    write (u, "(A)")  "*            in a multi-channel configuration"
    write (u, "(A)")

    write (u, "(A)")  "* Build and initialize a process object"
    write (u, "(A)")

    libname = "processes11"
    procname = libname
    run_id = "run11"
    call os_data_init (os_data)
    allocate (rng_test_factory_t :: rng_factory)
    call prc_test_create_library (libname, lib)

    call reset_interaction_counter ()

    allocate (model)
    call model%init_test ()

    allocate (process)
    call process%init (procname, run_id, lib, os_data, qcd, rng_factory, model)
    
    allocate (test_t :: core_template)
    allocate (mci_test_t :: mci_template)
    select type (mci_template)
    type is (mci_test_t)
       call mci_template%set_divisions (100)
    end select
    allocate (phs_test_config_t :: phs_config_template)
    call process%init_component &
         (1, core_template, mci_template, phs_config_template)

    write (u, "(A)")  "* Prepare a trivial beam setup"
    write (u, "(A)")

    sqrts = 1000
    call process%setup_beams_sqrts (sqrts)
    call process%configure_phs ()

    call process%setup_mci ()
    
    write (u, "(A)")  "* Complete process initialization"
    write (u, "(A)")

    call process%setup_terms ()
    call process%write (.false., u)
    
    write (u, "(A)")
    write (u, "(A)")  "* Create a process instance"
    write (u, "(A)")

    allocate (process_instance)
    call process_instance%init (process)

    write (u, "(A)")  "* Integrate with default test parameters"
    write (u, "(A)")

    call process%integrate (process_instance, 1, n_it=1, n_calls=10000)
    call process%final_integration (1)
    
    call process%write (.false., u)

    write (u, "(A)")
    write (u, "(A,ES13.7)")  " Integral divided by phs factor = ", &
         process%get_integral (1) &
         / process_instance%component(1)%k_seed%phs_factor

    write (u, "(A)")
    write (u, "(A)")  "* Cleanup"
    
    call process_instance%final ()
    deallocate (process_instance)
    
    call process%final ()
    deallocate (process)

    write (u, "(A)")
    write (u, "(A)")  "* Test output end: processes_11"
    
  end subroutine processes_11
  
\eatline
\nwendcode{}\nwbegindocs{967}\nwdocspar
\subsubsection{Complete events}
For the purpose of simplifying further tests, we implement a
convenience routine that initializes a process and prepares a single
event.  This is a wrapup of the test {\tt{}processes{\char95}10}.
\nwenddocs{}\nwbegincode{968}\moddef{Processes: public}\plusendmoddef
  public :: prepare_test_process
\nwendcode{}\nwbegincode{969}\moddef{Processes: procedures}\plusendmoddef
  subroutine prepare_test_process (process, process_instance, model)
    type(process_t), intent(out), target :: process
    type(process_instance_t), intent(out), target :: process_instance
    class(model_data_t), intent(in), target :: model
    type(process_library_t), target :: lib
    type(string_t) :: libname
    type(string_t) :: procname
    type(string_t) :: run_id
    type(os_data_t) :: os_data
    class(model_data_t), pointer :: process_model
    type(qcd_t) :: qcd
    class(rng_factory_t), allocatable :: rng_factory
    class(prc_core_t), allocatable :: core_template
    class(mci_t), allocatable :: mci_template
    class(phs_config_t), allocatable :: phs_config_template
    real(default) :: sqrts
    libname = "processes_test"
    procname = libname
    run_id = "run_test"
    call os_data_init (os_data)
    allocate (rng_test_factory_t :: rng_factory)
    call prc_test_create_library (libname, lib)
    call reset_interaction_counter ()
    allocate (process_model)
    call process_model%init (model%get_name (), &
         model%get_n_real (), &
         model%get_n_complex (), &
         model%get_n_field (), &
         model%get_n_vtx ())
    call process_model%copy_from (model)
    call process%init (procname, run_id, &
         lib, os_data, qcd, rng_factory, process_model)
    allocate (test_t :: core_template)
    allocate (mci_test_t :: mci_template)
    select type (mci_template)
    type is (mci_test_t);  call mci_template%set_divisions (100)
    end select
    allocate (phs_test_config_t :: phs_config_template)
    call process%init_component &
         (1, core_template, mci_template, phs_config_template)
    sqrts = 1000
    call process%setup_beams_sqrts (sqrts)
    call process%configure_phs ()
    call process%setup_mci ()
    call process%setup_terms ()
    call process_instance%init (process)
    select type (mci => process%mci_entry(1)%mci)
    type is (mci_test_t)
       ! This ensures that the next 'random' numbers are 0.3, 0.5, 0.7
       call mci%rng%init (3)
       ! Include the constant PHS factor in the stored maximum of the integrand
       call mci%set_max_factor (conv * twopi4 &
            / (2 * sqrt (lambda (sqrts **2, 125._default**2, 125._default**2))))
    end select
  end subroutine prepare_test_process

\eatline
\nwendcode{}\nwbegindocs{970}Here we do the cleanup of the process and process instance emitted
by the previous routine.
\nwenddocs{}\nwbegincode{971}\moddef{Processes: public}\plusendmoddef
  public :: cleanup_test_process
\nwendcode{}\nwbegincode{972}\moddef{Processes: procedures}\plusendmoddef
  subroutine cleanup_test_process (process, process_instance)
    type(process_t), intent(inout) :: process
    type(process_instance_t), intent(inout) :: process_instance
    call process_instance%final ()
    call process%final ()
  end subroutine cleanup_test_process
    
\eatline
\nwendcode{}\nwbegindocs{973}\nwdocspar
This is the actual test.  Prepare the test process and event, fill
all evaluators, and display the results.  Use a particle set as
temporary storage, read kinematics and recalculate the event.
\nwenddocs{}\nwbegincode{974}\moddef{Processes: execute tests}\plusendmoddef
  call test (processes_12, "processes_12", &
       "event post-processing", &
       u, results)
\nwendcode{}\nwbegincode{975}\moddef{Processes: tests}\plusendmoddef
  subroutine processes_12 (u)
    integer, intent(in) :: u
    type(process_t), allocatable, target :: process
    type(process_instance_t), allocatable, target :: process_instance
    type(particle_set_t) :: pset
    type(model_data_t), target :: model

    write (u, "(A)")  "* Test output: processes_12"
    write (u, "(A)")  "*   Purpose: generate a complete partonic event"
    write (u, "(A)")

    call model%init_test ()

    write (u, "(A)")  "* Build and initialize process and process instance &
         &and generate event"
    write (u, "(A)")

    allocate (process)
    allocate (process_instance)
    call prepare_test_process (process, process_instance, model)
    call process_instance%setup_event_data ()

    call process%prepare_simulation (1)
    call process_instance%init_simulation (1)
    call process%generate_weighted_event (process_instance, 1)
    call process_instance%evaluate_event_data ()

    call process_instance%write (u)

    call process_instance%get_trace (pset, 1)

    call process_instance%final_simulation (1)
    call process_instance%final ()
    deallocate (process_instance)
    
    write (u, "(A)")
    write (u, "(A)")  "* Recover kinematics and recalculate"
    write (u, "(A)")

    call reset_interaction_counter (2)

    allocate (process_instance)
    call process_instance%init (process)
    call process_instance%setup_event_data ()

    call process_instance%choose_mci (1)
    call process_instance%set_trace (pset, 1, check_match = .false.)
    call process_instance%recover &
         (channel = 1, i_term = 1, update_sqme = .true.)

    call process%recover_event (process_instance, 1)
    call process_instance%evaluate_event_data ()

    call process_instance%write (u)
    
    
    write (u, "(A)")
    write (u, "(A)")  "* Cleanup"
    
    call cleanup_test_process (process, process_instance)
    deallocate (process_instance)
    deallocate (process)
    
    call model%final ()

    write (u, "(A)")
    write (u, "(A)")  "* Test output end: processes_12"
    
  end subroutine processes_12
  
\eatline
\nwendcode{}\nwbegindocs{976}\nwdocspar
\subsubsection{Colored interaction}
This test specifically checks the transformation of process data
(flavor, helicity, and color) into an interaction in a process term.

We use the {\tt{}test{\char95}t} process core (which has no nontrivial
particles), but call only the {\tt{}is{\char95}allowed} method, which always
returns true.
\nwenddocs{}\nwbegincode{977}\moddef{Processes: execute tests}\plusendmoddef
  call test (processes_13, "processes_13", &
       "colored interaction", &
       u, results)
\nwendcode{}\nwbegincode{978}\moddef{Processes: tests}\plusendmoddef
  subroutine processes_13 (u)
    integer, intent(in) :: u
    type(os_data_t) :: os_data
    type(model_data_t), target :: model
    type(process_term_t) :: term
    class(prc_core_t), allocatable :: core
    
    write (u, "(A)")  "* Test output: processes_13"
    write (u, "(A)")  "*   Purpose: initialized a colored interaction"
    write (u, "(A)")

    write (u, "(A)")  "* Set up a process constants block"
    write (u, "(A)")

    call os_data_init (os_data)
    call model%init_sm_test ()
!    call model_list%read_model (var_str ("QCD"), var_str ("QCD.mdl"), &
!         os_data, model)
    allocate (test_t :: core)

    associate (data => term%data)
      data%n_in = 2
      data%n_out = 3
      data%n_flv = 2
      data%n_hel = 2
      data%n_col = 2
      data%n_cin = 2

      allocate (data%flv_state (5, 2))
      data%flv_state (:,1) = [ 1, 21, 1, 21, 21]
      data%flv_state (:,2) = [ 2, 21, 2, 21, 21]

      allocate (data%hel_state (5, 2))
      data%hel_state (:,1) = [1, 1, 1, 1, 0]
      data%hel_state (:,2) = [1,-1, 1,-1, 0]

      allocate (data%col_state (2, 5, 2))
      data%col_state (:,:,1) = &
           reshape ([[1, 0], [2,-1], [3, 0], [2,-3], [0,0]], [2,5])
      data%col_state (:,:,2) = &
           reshape ([[1, 0], [2,-3], [3, 0], [2,-1], [0,0]], [2,5])

      allocate (data%ghost_flag (5, 2))
      data%ghost_flag(1:4,:) = .false.
      data%ghost_flag(5,:) = .true.
     
    end associate
    
    write (u, "(A)")  "* Set up the interaction"
    write (u, "(A)")
    
    call reset_interaction_counter ()
    call term%setup_interaction (core, model)
    call interaction_write (term%int, u)
    
    call model%final ()

    write (u, "(A)")
    write (u, "(A)")  "* Test output end: processes_13"
  end subroutine processes_13
    
\eatline
\nwendcode{}\nwbegindocs{979}\nwdocspar
\subsubsection{MD5 sums}
Configure a process with structure functions (multi-channel) and
compute MD5 sums
\nwenddocs{}\nwbegincode{980}\moddef{Processes: execute tests}\plusendmoddef
  call test (processes_14, "processes_14", &
       "process configuration and MD5 sum", &
       u, results)
\nwendcode{}\nwbegincode{981}\moddef{Processes: tests}\plusendmoddef
  subroutine processes_14 (u)
    integer, intent(in) :: u
    type(process_library_t), target :: lib
    type(string_t) :: libname
    type(string_t) :: procname
    type(string_t) :: run_id
    type(os_data_t) :: os_data
    type(qcd_t) :: qcd
    class(rng_factory_t), allocatable :: rng_factory
    class(model_data_t), pointer :: model
    type(process_t), allocatable, target :: process
    class(prc_core_t), allocatable :: core_template
    class(mci_t), allocatable :: mci_template
    class(phs_config_t), allocatable :: phs_config_template
    real(default) :: sqrts
    type(pdg_array_t) :: pdg_in
    class(sf_data_t), allocatable, target :: data
    type(sf_config_t), dimension(:), allocatable :: sf_config
    type(sf_channel_t), dimension(3) :: sf_channel

    write (u, "(A)")  "* Test output: processes_14"
    write (u, "(A)")  "*   Purpose: initialize a process with &
         &structure functions"
    write (u, "(A)")  "*            and compute MD5 sum"
    write (u, "(A)")

    write (u, "(A)")  "* Build and initialize a process object"
    write (u, "(A)")

    libname = "processes7"
    procname = libname
    run_id = "run7"
    call os_data_init (os_data)
    allocate (rng_test_factory_t :: rng_factory)
    call prc_test_create_library (libname, lib)
    call lib%compute_md5sum ()
    
    allocate (model)
    call model%init_test ()

    allocate (process)
    call process%init (procname, run_id, lib, os_data, qcd, rng_factory, model)
    
    allocate (test_t :: core_template)
    allocate (phs_test_config_t :: phs_config_template)
    call process%init_component &
         (1, core_template, mci_template, phs_config_template)

    write (u, "(A)")  "* Set beam, structure functions, and mappings"
    write (u, "(A)")

    sqrts = 1000
    call process%setup_beams_sqrts (sqrts)
    call process%configure_phs ()
    
    pdg_in = 25
    allocate (sf_test_data_t :: data)
    select type (data)
    type is (sf_test_data_t)
       model => process%get_model_ptr ()
       call data%init (model, pdg_in)
    end select

    call process%beam_config%allocate_sf_channels (3)

    allocate (sf_config (2))
    call sf_config(1)%init ([1], data)
    call sf_config(2)%init ([2], data)
    call process%init_sf_chain (sf_config)
    deallocate (sf_config)

    call sf_channel(1)%init (2)
    call process%set_sf_channel (1, sf_channel(1))

    call sf_channel(2)%init (2)
    call sf_channel(2)%activate_mapping ([1,2])
    call process%set_sf_channel (2, sf_channel(2))
    
    call sf_channel(3)%init (2)
    call sf_channel(3)%set_s_mapping ([1,2])
    call process%set_sf_channel (3, sf_channel(3))
    
    call process%setup_mci ()

    call process%compute_md5sum ()

    call process%write (.false., u)
    
    write (u, "(A)")
    write (u, "(A)")  "* Cleanup"

    call process%final ()
    deallocate (process)
    
    write (u, "(A)")
    write (u, "(A)")  "* Test output end: processes_14"
    
  end subroutine processes_14
  
\eatline
\nwendcode{}\nwbegindocs{982}\nwdocspar
\subsubsection{Decay Process Evaluation}
Initialize an evaluate a decay process.
\nwenddocs{}\nwbegincode{983}\moddef{Processes: execute tests}\plusendmoddef
  call test (processes_15, "processes_15", &
       "decay process", &
       u, results)
\nwendcode{}\nwbegincode{984}\moddef{Processes: tests}\plusendmoddef
  subroutine processes_15 (u)
    integer, intent(in) :: u
    type(process_library_t), target :: lib
    type(string_t) :: libname
    type(string_t) :: procname
    type(string_t) :: run_id
    type(os_data_t) :: os_data
    type(qcd_t) :: qcd
    class(rng_factory_t), allocatable :: rng_factory
    class(model_data_t), pointer :: model
    type(process_t), allocatable, target :: process
    class(prc_core_t), allocatable :: core_template
    class(mci_t), allocatable :: mci_template
    class(phs_config_t), allocatable :: phs_config_template
    type(process_instance_t), allocatable, target :: process_instance
    type(particle_set_t) :: pset

    write (u, "(A)")  "* Test output: processes_15"
    write (u, "(A)")  "*   Purpose: initialize a decay process object"
    write (u, "(A)")

    write (u, "(A)")  "* Build and load a test library with one process"
    write (u, "(A)")

    libname = "processes15"
    procname = libname
    run_id = "run15"
    call os_data_init (os_data)
    allocate (rng_test_factory_t :: rng_factory)
    call prc_test_create_library (libname, lib, scattering = .false., &
         decay = .true.)

    allocate (model)
    call model%init_test ()
    call model%set_par (var_str ("ff"), 0.4_default)
    call model%set_par (var_str ("mf"), &
         model%get_real (var_str ("ff")) * model%get_real (var_str ("ms")))

    write (u, "(A)")  "* Initialize a process object"
    write (u, "(A)")

    allocate (process)
    call process%init (procname, run_id, lib, os_data, qcd, rng_factory, model)
    
    allocate (test_t :: core_template)
    allocate (phs_single_config_t :: phs_config_template)
    call process%init_component &
         (1, core_template, mci_template, phs_config_template)

    write (u, "(A)")  "* Prepare a trivial beam setup"
    write (u, "(A)")
    
    call process%setup_beams_decay ()
    call process%configure_phs ()
    call process%setup_mci ()

    write (u, "(A)")  "* Complete process initialization"
    write (u, "(A)")

    call process%setup_terms ()
    call process%write (.false., u)
    
    write (u, "(A)")
    write (u, "(A)")  "* Create a process instance"
    write (u, "(A)")

    call reset_interaction_counter (3)

    allocate (process_instance)
    call process_instance%init (process)
    call process_instance%write (u)

    write (u, "(A)")
    write (u, "(A)")  "* Inject a set of random numbers"
    write (u, "(A)")
     
    call process_instance%choose_mci (1)
    call process_instance%set_mcpar ([0._default, 0._default])
    call process_instance%write (u)

    write (u, "(A)")
    write (u, "(A)")  "* Set up hard kinematics"
    write (u, "(A)")
  
    call process_instance%select_channel (1)
    call process_instance%compute_seed_kinematics ()
    call process_instance%compute_hard_kinematics ()

    write (u, "(A)")  "* Evaluate matrix element and square"
    write (u, "(A)")
  
    call process_instance%compute_eff_kinematics ()
    call process_instance%evaluate_expressions ()
    call process_instance%compute_other_channels ()
    call process_instance%evaluate_trace ()
    call process_instance%write (u)

    call process_instance%get_trace (pset, 1)
    call process_instance%final ()
    deallocate (process_instance)

    write (u, "(A)")
    write (u, "(A)")  "* Particle content:"
    write (u, "(A)")
    
    call write_separator (u)
    call pset%write (u)
    call write_separator (u)

    write (u, "(A)")
    write (u, "(A)")  "* Recover process instance"
    write (u, "(A)")

    call reset_interaction_counter (3)

    allocate (process_instance)
    call process_instance%init (process)
    call process_instance%choose_mci (1)
    call process_instance%set_trace (pset, 1, check_match = .false.)
    call process_instance%recover (1, 1, .true.)
    call process_instance%write (u)

    write (u, "(A)")
    write (u, "(A)")  "* Cleanup"

    call pset%final ()
    call process_instance%final ()
    deallocate (process_instance)
    
    call process%final ()
    deallocate (process)
    
    write (u, "(A)")
    write (u, "(A)")  "* Test output end: processes_15"
    
  end subroutine processes_15
  
\eatline
\nwendcode{}\nwbegindocs{985}\nwdocspar
\subsubsection{Integration: decay}
Activate the MC integrator for the decay object and use it to
integrate over phase space.
\nwenddocs{}\nwbegincode{986}\moddef{Processes: execute tests}\plusendmoddef
  call test (processes_16, "processes_16", &
       "decay integration", &
       u, results)
\nwendcode{}\nwbegincode{987}\moddef{Processes: tests}\plusendmoddef
  subroutine processes_16 (u)
    integer, intent(in) :: u
    type(process_library_t), target :: lib
    type(string_t) :: libname
    type(string_t) :: procname
    type(string_t) :: run_id
    type(os_data_t) :: os_data
    type(qcd_t) :: qcd
    class(rng_factory_t), allocatable :: rng_factory
    class(model_data_t), pointer :: model
    type(process_t), allocatable, target :: process
    class(prc_core_t), allocatable :: core_template
    class(mci_t), allocatable :: mci_template
    class(phs_config_t), allocatable :: phs_config_template
    type(process_instance_t), allocatable, target :: process_instance

    write (u, "(A)")  "* Test output: processes_16"
    write (u, "(A)")  "*   Purpose: integrate a process without &
         &structure functions"
    write (u, "(A)")  "*            in a multi-channel configuration"
    write (u, "(A)")

    write (u, "(A)")  "* Build and initialize a process object"
    write (u, "(A)")

    libname = "processes16"
    procname = libname
    run_id = "run16"
    call os_data_init (os_data)
    allocate (rng_test_factory_t :: rng_factory)
    call prc_test_create_library (libname, lib, scattering = .false., &
         decay = .true.)

    call reset_interaction_counter ()

    allocate (model)
    call model%init_test ()
    call model%set_par (var_str ("ff"), 0.4_default)
    call model%set_par (var_str ("mf"), &
         model%get_real (var_str ("ff")) * model%get_real (var_str ("ms")))

    allocate (process)
    call process%init (procname, run_id, lib, os_data, qcd, rng_factory, model)
    
    allocate (test_t :: core_template)
    allocate (mci_midpoint_t :: mci_template)
    allocate (phs_single_config_t :: phs_config_template)
    call process%init_component &
         (1, core_template, mci_template, phs_config_template)

    write (u, "(A)")  "* Prepare a trivial beam setup"
    write (u, "(A)")

    call process%setup_beams_decay ()
    call process%configure_phs ()

    call process%setup_mci ()
    
    write (u, "(A)")  "* Complete process initialization"
    write (u, "(A)")

    call process%setup_terms ()
    call process%write (.false., u)
    
    write (u, "(A)")
    write (u, "(A)")  "* Create a process instance"
    write (u, "(A)")

    allocate (process_instance)
    call process_instance%init (process)

    write (u, "(A)")  "* Integrate with default test parameters"
    write (u, "(A)")

    call process%integrate (process_instance, 1, n_it=1, n_calls=10000)
    call process%final_integration (1)
    
    call process%write (.false., u)

    write (u, "(A)")
    write (u, "(A,ES13.7)")  " Integral divided by phs factor = ", &
         process%get_integral (1) &
         / process_instance%component(1)%k_seed%phs_factor

    write (u, "(A)")
    write (u, "(A)")  "* Cleanup"
    
    call process_instance%final ()
    deallocate (process_instance)
    
    call process%final ()
    deallocate (process)

    write (u, "(A)")
    write (u, "(A)")  "* Test output end: processes_16"
    
  end subroutine processes_16
  
\eatline
\nwendcode{}\nwbegindocs{988}\nwdocspar
\subsubsection{Decay Process Evaluation}
Initialize an evaluate a decay process for a moving particle.
\nwenddocs{}\nwbegincode{989}\moddef{Processes: execute tests}\plusendmoddef
  call test (processes_17, "processes_17", &
       "decay of moving particle", &
       u, results)
\nwendcode{}\nwbegincode{990}\moddef{Processes: tests}\plusendmoddef
  subroutine processes_17 (u)
    integer, intent(in) :: u
    type(process_library_t), target :: lib
    type(string_t) :: libname
    type(string_t) :: procname
    type(string_t) :: run_id
    type(os_data_t) :: os_data
    type(qcd_t) :: qcd
    class(rng_factory_t), allocatable :: rng_factory
    class(model_data_t), pointer :: model
    type(var_list_t), pointer :: var_list
    type(process_t), allocatable, target :: process
    class(prc_core_t), allocatable :: core_template
    class(mci_t), allocatable :: mci_template
    class(phs_config_t), allocatable :: phs_config_template
    type(process_instance_t), allocatable, target :: process_instance
    type(particle_set_t) :: pset
    type(flavor_t) :: flv_beam
    real(default) :: m, p, E

    write (u, "(A)")  "* Test output: processes_17"
    write (u, "(A)")  "*   Purpose: initialize a decay process object"
    write (u, "(A)")

    write (u, "(A)")  "* Build and load a test library with one process"
    write (u, "(A)")

    libname = "processes17"
    procname = libname
    run_id = "run17"
    call os_data_init (os_data)
    allocate (rng_test_factory_t :: rng_factory)

    call prc_test_create_library (libname, lib, scattering = .false., &
         decay = .true.)

    write (u, "(A)")  "* Initialize a process object"
    write (u, "(A)")

    allocate (model)
    call model%init_test ()
    call model%set_par (var_str ("ff"), 0.4_default)
    call model%set_par (var_str ("mf"), &
         model%get_real (var_str ("ff")) * model%get_real (var_str ("ms")))

    allocate (process)
    call process%init (procname, run_id, lib, os_data, qcd, rng_factory, model)
    
    allocate (test_t :: core_template)
    allocate (phs_single_config_t :: phs_config_template)
    call process%init_component &
         (1, core_template, mci_template, phs_config_template)

    write (u, "(A)")  "* Prepare a trivial beam setup"
    write (u, "(A)")
    
    call process%setup_beams_decay (rest_frame = .false.)
    call process%configure_phs ()
    call process%setup_mci ()

    write (u, "(A)")  "* Complete process initialization"
    write (u, "(A)")

    call process%setup_terms ()
    call process%write (.false., u)
    
    write (u, "(A)")
    write (u, "(A)")  "* Create a process instance"
    write (u, "(A)")

    call reset_interaction_counter (3)

    allocate (process_instance)
    call process_instance%init (process)
    call process_instance%write (u)

    write (u, "(A)")
    write (u, "(A)")  "* Set parent momentum and random numbers"
    write (u, "(A)")
     
    call process_instance%choose_mci (1)
    call process_instance%set_mcpar ([0._default, 0._default])

    model => process%get_model_ptr ()
    call flv_beam%init (25, model)
    m = flv_beam%get_mass ()
    p = 3 * m / 4
    E = sqrt (m**2 + p**2)
    call process_instance%set_beam_momenta ([vector4_moving (E, p, 3)])
     
    call process_instance%write (u)

    write (u, "(A)")
    write (u, "(A)")  "* Set up hard kinematics"
    write (u, "(A)")
  
    call process_instance%select_channel (1)
    call process_instance%compute_seed_kinematics ()
    call process_instance%compute_hard_kinematics ()

    write (u, "(A)")  "* Evaluate matrix element and square"
    write (u, "(A)")
  
    call process_instance%compute_eff_kinematics ()
    call process_instance%evaluate_expressions ()
    call process_instance%compute_other_channels ()
    call process_instance%evaluate_trace ()
    call process_instance%write (u)

    call process_instance%get_trace (pset, 1)
    call process_instance%final ()
    deallocate (process_instance)

    write (u, "(A)")
    write (u, "(A)")  "* Particle content:"
    write (u, "(A)")
    
    call write_separator (u)
    call pset%write (u)
    call write_separator (u)

    write (u, "(A)")
    write (u, "(A)")  "* Recover process instance"
    write (u, "(A)")

    call reset_interaction_counter (3)

    allocate (process_instance)
    call process_instance%init (process)

    call process_instance%choose_mci (1)
    call process_instance%set_trace (pset, 1, check_match = .false.)
    call process_instance%recover (1, 1, .true.)
    call process_instance%write (u)

    write (u, "(A)")
    write (u, "(A)")  "* Cleanup"

    call pset%final ()
    call process_instance%final ()
    deallocate (process_instance)
    
    call process%final ()
    deallocate (process)
    
    write (u, "(A)")
    write (u, "(A)")  "* Test output end: processes_17"
    
  end subroutine processes_17
  
\eatline
\nwendcode{}\nwbegindocs{991}\nwdocspar
%%%%%%%%%%%%%%%%%%%%%%%%%%%%%%%%%%%%%%%%%%%%%%%%%%%%%%%%%%%%%%%%%%%%%%%%
\section{Process Stacks}

For storing and handling multiple processes, we define process stacks.
These are ordinary stacks where new process entries are pushed onto
the top.  We allow for multiple entries with identical process ID, but
distinct run ID.

The implementation is essentially identical to the {\tt{}prclib{\char95}stacks} module
above.  Unfortunately, Fortran supports no generic programming, so we do not
make use of this fact.

When searching for a specific process ID, we will get (a pointer to)
the topmost process entry with that ID on the stack, which was entered
last.  Usually, this is the best version of the process (in terms of
integral, etc.)  Thus the stack terminology makes sense.
\nwenddocs{}\nwbegincode{992}\moddef{\code{}process{\char95}stacks.f90\edoc{}}\endmoddef
\LA{}File header\RA{}

module process_stacks
  
\LA{}Use strings\RA{}
  use io_units
  use format_utils, only: write_separator
  use unit_tests
  use diagnostics
  use os_interface
  use sm_qcd
  use model_data
  use rng_base
  use variables
  use observables
  use process_libraries
  use prc_test
  use processes

\LA{}Standard module head\RA{}

\LA{}Process stacks: public\RA{}

\LA{}Process stacks: types\RA{}

contains
  
\LA{}Process stacks: procedures\RA{}

\LA{}Process stacks: tests\RA{}

end module process_stacks
\eatline
\nwendcode{}\nwbegindocs{993}\nwdocspar
\subsection{The process entry type}
A process entry is a process object, augmented by a pointer to the
next entry.  We do not need specific methods, all relevant methods are
inherited.

On higher level, processes should be prepared as process entry objects.
\nwenddocs{}\nwbegincode{994}\moddef{Process stacks: public}\endmoddef
  public :: process_entry_t
\nwendcode{}\nwbegincode{995}\moddef{Process stacks: types}\endmoddef
  type, extends (process_t) :: process_entry_t
     type(process_entry_t), pointer :: next => null ()
  end type process_entry_t

\eatline
\nwendcode{}\nwbegindocs{996}\nwdocspar
\subsection{The process stack type}
For easy conversion and lookup it is useful to store the filling
number in the object.  The content is stored as a linked list.

The {\tt{}var{\char95}list} component stores process-specific results, so they
can be retrieved as (pseudo) variables.

The process stack can be linked to another one.  This allows us to
work with stacks of local scope.
\nwenddocs{}\nwbegincode{997}\moddef{Process stacks: public}\plusendmoddef
  public :: process_stack_t
\nwendcode{}\nwbegincode{998}\moddef{Process stacks: types}\plusendmoddef
  type :: process_stack_t
     integer :: n = 0
     type(process_entry_t), pointer :: first => null ()
     type(var_list_t), pointer :: var_list => null ()
     type(process_stack_t), pointer :: next => null ()
   contains
   \LA{}Process stacks: process stack: TBP\RA{}
  end type process_stack_t
  
\eatline
\nwendcode{}\nwbegindocs{999}Finalize partly: deallocate the process stack and variable list
entries, but keep the variable list as an empty object.  This way, the
variable list links are kept.
\nwenddocs{}\nwbegincode{1000}\moddef{Process stacks: process stack: TBP}\endmoddef
  procedure :: clear => process_stack_clear
\nwendcode{}\nwbegincode{1001}\moddef{Process stacks: procedures}\endmoddef
  subroutine process_stack_clear (stack)
    class(process_stack_t), intent(inout) :: stack
    type(process_entry_t), pointer :: process
    if (associated (stack%var_list)) then
       call stack%var_list%final ()
    end if
    do while (associated (stack%first))
       process => stack%first
       stack%first => process%next
       call process%final ()
       deallocate (process)
    end do
    stack%n = 0
  end subroutine process_stack_clear
  
\eatline
\nwendcode{}\nwbegindocs{1002}Finalizer.  Clear and deallocate the variable list.
\nwenddocs{}\nwbegincode{1003}\moddef{Process stacks: process stack: TBP}\plusendmoddef
  procedure :: final => process_stack_final
\nwendcode{}\nwbegincode{1004}\moddef{Process stacks: procedures}\plusendmoddef
  subroutine process_stack_final (object)
    class(process_stack_t), intent(inout) :: object
    type(process_entry_t), pointer :: process
    call object%clear ()
    if (associated (object%var_list)) then
       deallocate (object%var_list)
    end if
  end subroutine process_stack_final
  
\eatline
\nwendcode{}\nwbegindocs{1005}Output.  The processes on the stack will be ordered LIFO, i.e.,
backwards.
\nwenddocs{}\nwbegincode{1006}\moddef{Process stacks: process stack: TBP}\plusendmoddef
  procedure :: write => process_stack_write
\nwendcode{}\nwbegincode{1007}\moddef{Process stacks: procedures}\plusendmoddef
  recursive subroutine process_stack_write (object, unit, pacify)
    class(process_stack_t), intent(in) :: object
    integer, intent(in), optional :: unit
    logical, intent(in), optional :: pacify
    type(process_entry_t), pointer :: process
    integer :: u
    u = given_output_unit (unit)
    call write_separator (u, 2)
    select case (object%n)
    case (0)
       write (u, "(1x,A)")  "Process stack: [empty]"
       call write_separator (u, 2)
    case default
       write (u, "(1x,A)")  "Process stack:"
       process => object%first
       do while (associated (process))
          call process%write (.false., u, pacify = pacify)
          process => process%next
       end do
    end select
    if (associated (object%next)) then
       write (u, "(1x,A)")  "[Processes from context environment:]"
       call object%next%write (u, pacify)
    end if
  end subroutine process_stack_write

\eatline
\nwendcode{}\nwbegindocs{1008}The variable list is printed by a separate routine, since
it should be linked to the global variable list, anyway.
\nwenddocs{}\nwbegincode{1009}\moddef{Process stacks: process stack: TBP}\plusendmoddef
  procedure :: write_var_list => process_stack_write_var_list
\nwendcode{}\nwbegincode{1010}\moddef{Process stacks: procedures}\plusendmoddef
  subroutine process_stack_write_var_list (object, unit)
    class(process_stack_t), intent(in) :: object
    integer, intent(in), optional :: unit
    if (associated (object%var_list)) then
       call var_list_write (object%var_list, unit)
    end if
  end subroutine process_stack_write_var_list

\eatline
\nwendcode{}\nwbegindocs{1011}Short output.
\nwenddocs{}\nwbegincode{1012}\moddef{Process stacks: process stack: TBP}\plusendmoddef
  procedure :: show => process_stack_show
\nwendcode{}\nwbegincode{1013}\moddef{Process stacks: procedures}\plusendmoddef
  recursive subroutine process_stack_show (object, unit)
    class(process_stack_t), intent(in) :: object
    integer, intent(in), optional :: unit
    type(process_entry_t), pointer :: process
    integer :: u
    u = given_output_unit (unit)
    select case (object%n)
    case (0)
    case default
       process => object%first
       do while (associated (process))
          call process%show (u, verbose=.false.)
          process => process%next
       end do
    end select
    if (associated (object%next))  call object%next%show ()
  end subroutine process_stack_show

\eatline
\nwendcode{}\nwbegindocs{1014}\nwdocspar
\subsection{Link}
Link the current process stack to a global one.
\nwenddocs{}\nwbegincode{1015}\moddef{Process stacks: process stack: TBP}\plusendmoddef
  procedure :: link => process_stack_link
\nwendcode{}\nwbegincode{1016}\moddef{Process stacks: procedures}\plusendmoddef
  subroutine process_stack_link (local_stack, global_stack)
    class(process_stack_t), intent(inout) :: local_stack
    type(process_stack_t), intent(in), target :: global_stack
    local_stack%next => global_stack
  end subroutine process_stack_link

\eatline
\nwendcode{}\nwbegindocs{1017}Initialize the process variable list and link the main variable list
to it.
\nwenddocs{}\nwbegincode{1018}\moddef{Process stacks: process stack: TBP}\plusendmoddef
  procedure :: init_var_list => process_stack_init_var_list
\nwendcode{}\nwbegincode{1019}\moddef{Process stacks: procedures}\plusendmoddef
  subroutine process_stack_init_var_list (stack, var_list)
    class(process_stack_t), intent(inout) :: stack
    type(var_list_t), intent(inout), optional :: var_list
    allocate (stack%var_list)
    if (present (var_list))  call var_list%link (stack%var_list)
  end subroutine process_stack_init_var_list
  
\eatline
\nwendcode{}\nwbegindocs{1020}Link the process variable list to a global
variable list.
\nwenddocs{}\nwbegincode{1021}\moddef{Process stacks: process stack: TBP}\plusendmoddef
  procedure :: link_var_list => process_stack_link_var_list
\nwendcode{}\nwbegincode{1022}\moddef{Process stacks: procedures}\plusendmoddef
  subroutine process_stack_link_var_list (stack, var_list)
    class(process_stack_t), intent(inout) :: stack
    type(var_list_t), intent(in), target :: var_list
    call stack%var_list%link (var_list)
  end subroutine process_stack_link_var_list
  
\eatline
\nwendcode{}\nwbegindocs{1023}\nwdocspar
\subsection{Push}
We take a process pointer and push it onto the stack.  The previous
pointer is nullified.  Subsequently, the process is `owned' by the
stack and will be finalized when the stack is deleted.
\nwenddocs{}\nwbegincode{1024}\moddef{Process stacks: process stack: TBP}\plusendmoddef
  procedure :: push => process_stack_push
\nwendcode{}\nwbegincode{1025}\moddef{Process stacks: procedures}\plusendmoddef
  subroutine process_stack_push (stack, process)
    class(process_stack_t), intent(inout) :: stack
    type(process_entry_t), intent(inout), pointer :: process
    process%next => stack%first
    stack%first => process
    process => null ()
    stack%n = stack%n + 1
  end subroutine process_stack_push
  
\eatline
\nwendcode{}\nwbegindocs{1026}Initialize process variables for a given process ID, without setting
values.
\nwenddocs{}\nwbegincode{1027}\moddef{Process stacks: process stack: TBP}\plusendmoddef
  procedure :: init_result_vars => process_stack_init_result_vars
\nwendcode{}\nwbegincode{1028}\moddef{Process stacks: procedures}\plusendmoddef
  subroutine process_stack_init_result_vars (stack, id)
    class(process_stack_t), intent(inout) :: stack
    type(string_t), intent(in) :: id
    call var_list_init_num_id (stack%var_list, id)
    call var_list_init_process_results (stack%var_list, id)
  end subroutine process_stack_init_result_vars
  
\eatline
\nwendcode{}\nwbegindocs{1029}Fill process variables with values.  This is executed after the
integration pass.

Note: We set only integral and error.  With multiple MCI records
possible, the results for {\tt{}n{\char95}calls}, {\tt{}chi2} etc. are not
necessarily unique.  (We might set the efficiency, though.)
\nwenddocs{}\nwbegincode{1030}\moddef{Process stacks: process stack: TBP}\plusendmoddef
  procedure :: fill_result_vars => process_stack_fill_result_vars
\nwendcode{}\nwbegincode{1031}\moddef{Process stacks: procedures}\plusendmoddef
  subroutine process_stack_fill_result_vars (stack, id)
    class(process_stack_t), intent(inout) :: stack
    type(string_t), intent(in) :: id
    type(process_t), pointer :: process
    process => stack%get_process_ptr (id)
    if (associated (process)) then
       call var_list_init_num_id (stack%var_list, id, process%get_num_id ())
       if (process%has_integral ()) then
          call var_list_init_process_results (stack%var_list, id, &
               integral = process%get_integral (), &
               error = process%get_error ())
       end if
    else
       call msg_bug ("process_stack_fill_result_vars: unknown process ID")
    end if
  end subroutine process_stack_fill_result_vars

\eatline
\nwendcode{}\nwbegindocs{1032}\nwdocspar
\subsection{Data Access}
Tell if a process exists.
\nwenddocs{}\nwbegincode{1033}\moddef{Process stacks: process stack: TBP}\plusendmoddef
  procedure :: exists => process_stack_exists
\nwendcode{}\nwbegincode{1034}\moddef{Process stacks: procedures}\plusendmoddef
  function process_stack_exists (stack, id) result (flag)
    class(process_stack_t), intent(in) :: stack
    type(string_t), intent(in) :: id
    logical :: flag
    type(process_t), pointer :: process
    process => stack%get_process_ptr (id)
    flag = associated (process)
  end function process_stack_exists

\eatline
\nwendcode{}\nwbegindocs{1035}Return a pointer to a process with specific ID.  Look also at a
linked stack, if necessary.
\nwenddocs{}\nwbegincode{1036}\moddef{Process stacks: process stack: TBP}\plusendmoddef
  procedure :: get_process_ptr => process_stack_get_process_ptr
\nwendcode{}\nwbegincode{1037}\moddef{Process stacks: procedures}\plusendmoddef
  recursive function process_stack_get_process_ptr (stack, id) result (ptr)
    class(process_stack_t), intent(in) :: stack
    type(string_t), intent(in) :: id
    type(process_t), pointer :: ptr
    type(process_entry_t), pointer :: entry
    ptr => null ()
    entry => stack%first
    do while (associated (entry))
       if (entry%get_id () == id) then
          ptr => entry%process_t
          return
       end if
       entry => entry%next
    end do
    if (associated (stack%next))  ptr => stack%next%get_process_ptr (id)
  end function process_stack_get_process_ptr

\eatline
\nwendcode{}\nwbegindocs{1038}\nwdocspar
\subsection{Unit tests}
\nwenddocs{}\nwbegincode{1039}\moddef{Process stacks: public}\plusendmoddef
  public :: process_stacks_test
\nwendcode{}\nwbegincode{1040}\moddef{Process stacks: tests}\endmoddef
  subroutine process_stacks_test (u, results)
    integer, intent(in) :: u
    type(test_results_t), intent(inout) :: results
  \LA{}Process stacks: execute tests\RA{}
  end subroutine process_stacks_test
  
\eatline
\nwendcode{}\nwbegindocs{1041}\nwdocspar
\subsubsection{Write an empty process stack}
The most trivial test is to write an uninitialized process stack.
\nwenddocs{}\nwbegincode{1042}\moddef{Process stacks: execute tests}\endmoddef
  call test (process_stacks_1, "process_stacks_1", &
       "write an empty process stack", &
       u, results)
\nwendcode{}\nwbegincode{1043}\moddef{Process stacks: tests}\plusendmoddef
  subroutine process_stacks_1 (u)
    integer, intent(in) :: u
    type(process_stack_t) :: stack

    write (u, "(A)")  "* Test output: process_stacks_1"
    write (u, "(A)")  "*   Purpose: display an empty process stack"
    write (u, "(A)")

    call stack%write (u)
    
    write (u, "(A)")
    write (u, "(A)")  "* Test output end: process_stacks_1"
    
  end subroutine process_stacks_1
  
\eatline
\nwendcode{}\nwbegindocs{1044}\nwdocspar
\subsubsection{Fill a process stack}
Fill a process stack with two (identical) processes.
\nwenddocs{}\nwbegincode{1045}\moddef{Process stacks: execute tests}\plusendmoddef
  call test (process_stacks_2, "process_stacks_2", &
       "fill a process stack", &
       u, results)
\nwendcode{}\nwbegincode{1046}\moddef{Process stacks: tests}\plusendmoddef
  subroutine process_stacks_2 (u)
    integer, intent(in) :: u
    type(process_stack_t) :: stack
    type(process_library_t), target :: lib
    type(string_t) :: libname
    type(string_t) :: procname
    type(string_t) :: run_id
    type(os_data_t) :: os_data
    type(qcd_t) :: qcd
    class(rng_factory_t), allocatable :: rng_factory
    class(model_data_t), pointer :: model
    type(process_entry_t), pointer :: process => null ()

    write (u, "(A)")  "* Test output: process_stacks_2"
    write (u, "(A)")  "*   Purpose: fill a process stack"
    write (u, "(A)")

    write (u, "(A)")  "* Build, initialize and store two test processes"
    write (u, "(A)")

    libname = "process_stacks2"
    procname = libname
    call os_data_init (os_data)
    allocate (rng_test_factory_t :: rng_factory)
    call prc_test_create_library (libname, lib)

    allocate (model)
    call model%init_test ()

    allocate (process)
    run_id = "run1"
    call process%init &
         (procname, run_id, lib, os_data, qcd, rng_factory, model) 
    call stack%push (process)
    
    allocate (model)
    call model%init_test ()

    allocate (process)
    run_id = "run2"
    call process%init &
         (procname, run_id, lib, os_data, qcd, rng_factory, model) 
    call stack%push (process)
    
    call stack%write (u)
    
    write (u, "(A)")
    write (u, "(A)")  "* Cleanup"

    call stack%final ()

    write (u, "(A)")
    write (u, "(A)")  "* Test output end: process_stacks_2"
    
  end subroutine process_stacks_2
  
\eatline
\nwendcode{}\nwbegindocs{1047}\nwdocspar
\subsubsection{Fill a process stack}
Fill a process stack with two (identical) processes.
\nwenddocs{}\nwbegincode{1048}\moddef{Process stacks: execute tests}\plusendmoddef
  call test (process_stacks_3, "process_stacks_3", &
       "process variables", &
       u, results)
\nwendcode{}\nwbegincode{1049}\moddef{Process stacks: tests}\plusendmoddef
  subroutine process_stacks_3 (u)
    integer, intent(in) :: u
    type(process_stack_t) :: stack
    type(model_data_t), target :: model
    type(string_t) :: procname
    type(process_entry_t), pointer :: process => null ()
    type(process_instance_t), target :: process_instance

    write (u, "(A)")  "* Test output: process_stacks_3"
    write (u, "(A)")  "*   Purpose: setup process variables"
    write (u, "(A)")

    write (u, "(A)")  "* Initialize process variables"
    write (u, "(A)")

    procname = "processes_test"
    call model%init_test ()

    write (u, "(A)")  "* Initialize process variables"
    write (u, "(A)")

    call stack%init_var_list ()
    call stack%init_result_vars (procname)
    call stack%write_var_list (u)

    write (u, "(A)")
    write (u, "(A)")  "* Build and integrate a test process"
    write (u, "(A)")

    allocate (process)
    call prepare_test_process (process%process_t, process_instance, model)
    call process%integrate (process_instance, 1, 1, 1000)
    call process_instance%final ()
    call process%final_integration (1)
    call stack%push (process)
    
    write (u, "(A)")  "* Fill process variables"
    write (u, "(A)")

    call stack%fill_result_vars (procname)
    call stack%write_var_list (u)
    
    write (u, "(A)")
    write (u, "(A)")  "* Cleanup"

    call stack%final ()

    call model%final ()
    
    write (u, "(A)")
    write (u, "(A)")  "* Test output end: process_stacks_3"
    
  end subroutine process_stacks_3
  
\eatline
\nwendcode{}\nwbegindocs{1050}\nwdocspar
\subsubsection{Linked a process stack}
Fill two process stack, linked to each other.
\nwenddocs{}\nwbegincode{1051}\moddef{Process stacks: execute tests}\plusendmoddef
  call test (process_stacks_4, "process_stacks_4", &
       "linked stacks", &
       u, results)
\nwendcode{}\nwbegincode{1052}\moddef{Process stacks: tests}\plusendmoddef
  subroutine process_stacks_4 (u)
    integer, intent(in) :: u
    type(process_library_t), target :: lib
    type(process_stack_t), target :: stack1, stack2
    class(model_data_t), pointer :: model
    type(string_t) :: libname
    type(string_t) :: procname
    type(string_t) :: run_id
    type(os_data_t) :: os_data
    type(qcd_t) :: qcd
    class(rng_factory_t), allocatable :: rng_factory
    type(process_entry_t), pointer :: process => null ()

    write (u, "(A)")  "* Test output: process_stacks_4"
    write (u, "(A)")  "*   Purpose: link process stacks"
    write (u, "(A)")

    write (u, "(A)")  "* Initialize process variables"
    write (u, "(A)")

    libname = "process_stacks_4_lib"
    procname = "process_stacks_4a"

    call os_data_init (os_data)
    allocate (rng_test_factory_t :: rng_factory)

    write (u, "(A)")  "* Initialize first process"
    write (u, "(A)")

    call prc_test_create_library (procname, lib)

    allocate (model)
    call model%init_test ()

    allocate (process)
    run_id = "run1"
    call process%init &
         (procname, run_id, lib, os_data, qcd, rng_factory, model) 
    call stack1%push (process)
    
    write (u, "(A)")  "* Initialize second process"
    write (u, "(A)")

    call stack2%link (stack1)

    procname = "process_stacks_4b"
    call prc_test_create_library (procname, lib)

    allocate (model)
    call model%init_test ()

    allocate (process)
    run_id = "run2"
    call process%init &
         (procname, run_id, lib, os_data, qcd, rng_factory, model) 
    call stack2%push (process)
    
    write (u, "(A)")  "* Show linked stacks"
    write (u, "(A)")

    call stack2%write (u)
    
    write (u, "(A)")
    write (u, "(A)")  "* Cleanup"

    call stack2%final ()
    call stack1%final ()

    write (u, "(A)")
    write (u, "(A)")  "* Test output end: process_stacks_4"
    
  end subroutine process_stacks_4
  
\eatline
\nwendcode{}\nwbegindocs{1053}\nwdocspar
%%%%%%%%%%%%%%%%%%%%%%%%%%%%%%%%%%%%%%%%%%%%%%%%%%%%%%%%%%%%%%%%%%%%%%%%
\chapter{Event Implementation}

With a process object and the associated methods at hand, we can
generate events for elementary processes and, by subsequent
transformation, for complete physical processes.

We have the following modules:
\begin{description}
\item[event\_transforms]
  Abstract base type for transforming a physical process with process
  instance and included evaluators, etc., into a new
  object.  The following modules extend this base type.
\item[decays]
  Combine the elementary process with elementary decay processes and
  thus transform the elementary event into a decayed event, still at
  the parton level.
\item[showers]
  Create QED/QCD showers out of the partons that are emitted by
  elementary processes.  This should be interleaved with showering of
  radiated particles (structure functions) and multiple interactions.
\item[hadrons]
  (not implemented yet) Apply hadronization to the partonic events,
  interleaved with hadron decays.  (The current setup relies on
  hadronizing partonic events externally.)
\item[events]
  Combine all pieces to generate full events.
\item[eio\_raw]
  Raw I/O for complete events.
\end{description}

\nwenddocs{}\nwbegindocs{1054}\nwdocspar
%%%%%%%%%%%%%%%%%%%%%%%%%%%%%%%%%%%%%%%%%%%%%%%%%%%%%%%%%%%%%%%%%%%%%%%%
\section{Abstract Event Transforms}
\nwenddocs{}\nwbegincode{1055}\moddef{\code{}event{\char95}transforms.f90\edoc{}}\endmoddef
\LA{}File header\RA{}

module event_transforms

\LA{}Use kinds\RA{}
\LA{}Use strings\RA{}
  use io_units
  use format_utils, only: write_separator
  use unit_tests
  use diagnostics
  use os_interface
  use sm_qcd
  use model_data
  use state_matrices
  use interactions
  use particles
  use subevents
  use process_libraries
  use prc_core
  use prc_test
  use rng_base
  use mci_base
  use mci_midpoint
  use phs_base
  use phs_single
  use processes
  use process_stacks
  
\LA{}Standard module head\RA{}

\LA{}Event transforms: public\RA{}

\LA{}Event transforms: types\RA{}

\LA{}Event transforms: interfaces\RA{}

contains

\LA{}Event transforms: procedures\RA{}

\LA{}Event transforms: tests\RA{}

end module event_transforms

\eatline
\nwendcode{}\nwbegindocs{1056}\nwdocspar
\subsection{Abstract base type}
Essentially, all methods are abstract, but some get minimal base
versions.  We know that there will be a random-number generator at top
level, and that we will relate to an elementary process.

The model is stored separately.  It may contain modified setting that differ
from the model instance stored in the process object.

Each event transform contains a particle set that it can fill for
further use.  There is a flag that indicates this.

We will collect event transforms in a list, therefore we include
{\tt{}previous} and {\tt{}next} pointers.
\nwenddocs{}\nwbegincode{1057}\moddef{Event transforms: public}\endmoddef
  public :: evt_t
\nwendcode{}\nwbegincode{1058}\moddef{Event transforms: types}\endmoddef
  type, abstract :: evt_t
     type(process_t), pointer :: process => null ()
     type(process_instance_t), pointer :: process_instance => null ()
     class(model_data_t), pointer :: model => null ()
     class(rng_t), allocatable :: rng
     integer :: rejection_count = 0
     logical :: particle_set_exists = .false.
     type(particle_set_t) :: particle_set
     class(evt_t), pointer :: previous => null ()
     class(evt_t), pointer :: next => null ()
   contains
   \LA{}Event transforms: evt: TBP\RA{}
  end type evt_t
  
\eatline
\nwendcode{}\nwbegindocs{1059}Finalizer.  In any case, we finalize the r.n.g.  The process
instance is a pointer and should not be finalized here.
\nwenddocs{}\nwbegincode{1060}\moddef{Event transforms: evt: TBP}\endmoddef
  procedure :: final => evt_final
  procedure :: base_final => evt_final
\nwendcode{}\nwbegincode{1061}\moddef{Event transforms: procedures}\endmoddef
  subroutine evt_final (object)
    class(evt_t), intent(inout) :: object
    if (allocated (object%rng))  call object%rng%final ()
    if (object%particle_set_exists) &
         call object%particle_set%final ()
  end subroutine evt_final
  
\eatline
\nwendcode{}\nwbegindocs{1062}Output.  We can print r.n.g. info.
\nwenddocs{}\nwbegincode{1063}\moddef{Event transforms: evt: TBP}\plusendmoddef
  procedure :: base_write => evt_write
\nwendcode{}\nwbegincode{1064}\moddef{Event transforms: procedures}\plusendmoddef
  subroutine evt_write (object, unit, verbose, testflag)
    class(evt_t), intent(in) :: object
    integer, intent(in), optional :: unit
    logical, intent(in), optional :: verbose, testflag
    integer :: u
    u = given_output_unit (unit)
    if (associated (object%process)) then
       write (u, "(3x,A,A,A)")   "Associated process: '", &
            char (object%process%get_id ()), "'"
    end if
    if (allocated (object%rng)) then
       call object%rng%write (u, 1)
       write (u, "(3x,A,I0)")  "Number of tries = ", object%rejection_count
    end if
    if (object%particle_set_exists) then
       call write_separator (u)
       call object%particle_set%write (u, testflag)
    end if
  end subroutine evt_write
  
\eatline
\nwendcode{}\nwbegindocs{1065}Connect the transform with a process instance (and thus with the
associated process).  Use this to allocate the master random-number
generator.

This is not an initializer; we may initialize the transform by
implementation-specific methods.
\nwenddocs{}\nwbegincode{1066}\moddef{Event transforms: evt: TBP}\plusendmoddef
  procedure :: connect => evt_connect
  procedure :: base_connect => evt_connect
\nwendcode{}\nwbegincode{1067}\moddef{Event transforms: procedures}\plusendmoddef
  subroutine evt_connect (evt, process_instance, model, process_stack)
    class(evt_t), intent(inout), target :: evt
    type(process_instance_t), intent(in), target :: process_instance
    class(model_data_t), intent(in), target :: model
    type(process_stack_t), intent(in), optional :: process_stack
    evt%process => process_instance%process
    evt%process_instance => process_instance
    evt%model => model
    call evt%process%make_rng (evt%rng)
  end subroutine evt_connect
  
\eatline
\nwendcode{}\nwbegindocs{1068}Reset internal state.
\nwenddocs{}\nwbegincode{1069}\moddef{Event transforms: evt: TBP}\plusendmoddef
  procedure :: reset => evt_reset
  procedure :: base_reset => evt_reset
\nwendcode{}\nwbegincode{1070}\moddef{Event transforms: procedures}\plusendmoddef
  subroutine evt_reset (evt)
    class(evt_t), intent(inout) :: evt
    evt%rejection_count = 0
    evt%particle_set_exists = .false.
  end subroutine evt_reset
  
\eatline
\nwendcode{}\nwbegindocs{1071}Prepare for a new event: reset internal state, if necessary.  We
provide MCI and term index of the parent process.
\nwenddocs{}\nwbegincode{1072}\moddef{Event transforms: evt: TBP}\plusendmoddef
  procedure (evt_prepare_new_event), deferred :: prepare_new_event
\nwendcode{}\nwbegincode{1073}\moddef{Event transforms: interfaces}\endmoddef
  interface
     subroutine evt_prepare_new_event (evt, i_mci, i_term)
       import
       class(evt_t), intent(inout) :: evt
       integer, intent(in) :: i_mci, i_term
     end subroutine evt_prepare_new_event
  end interface
       
\eatline
\nwendcode{}\nwbegindocs{1074}Generate a weighted event, using a valid initiator event in the
process instance, and the random-number generator.  The returned event
probability should be a number between zero and one that we can use for
rejection.
\nwenddocs{}\nwbegincode{1075}\moddef{Event transforms: evt: TBP}\plusendmoddef
  procedure (evt_generate_weighted), deferred :: generate_weighted
\nwendcode{}\nwbegincode{1076}\moddef{Event transforms: interfaces}\plusendmoddef
  abstract interface
     subroutine evt_generate_weighted (evt, probability)
       import
       class(evt_t), intent(inout) :: evt
       real(default), intent(out) :: probability
     end subroutine evt_generate_weighted
  end interface
  
\eatline
\nwendcode{}\nwbegindocs{1077}The unweighted event generation routine is actually implemented.  It
uses the random-number generator for simple rejection.  Of course, the
implementation may override this and implement a different way of
generating an unweighted event.
\nwenddocs{}\nwbegincode{1078}\moddef{Event transforms: evt: TBP}\plusendmoddef
  procedure :: generate_unweighted => evt_generate_unweighted
  procedure :: base_generate_unweighted => evt_generate_unweighted
\nwendcode{}\nwbegincode{1079}\moddef{Event transforms: procedures}\plusendmoddef
  subroutine evt_generate_unweighted (evt)
    class(evt_t), intent(inout) :: evt
    real(default) :: p, x
    evt%rejection_count = 0
    REJECTION: do
       evt%rejection_count = evt%rejection_count + 1
       call evt%generate_weighted (p)
       if (signal_is_pending ())  return
       call evt%rng%generate (x)
       if (x < p)  exit REJECTION
    end do REJECTION
  end subroutine evt_generate_unweighted
    
\eatline
\nwendcode{}\nwbegindocs{1080}Make a particle set.  This should take the most recent evaluator (or
whatever stores the event), factorize the density matrix if necessary,
and store as a particle set.

If applicable, the factorization should make use of the
{\tt{}factorization{\char95}mode} and {\tt{}keep{\char95}correlations} settings.

The values {\tt{}r}, if set, should control the factorization in more
detail, e.g., bypassing the random-number generator.
\nwenddocs{}\nwbegincode{1081}\moddef{Event transforms: evt: TBP}\plusendmoddef
  procedure (evt_make_particle_set), deferred :: make_particle_set
\nwendcode{}\nwbegincode{1082}\moddef{Event transforms: interfaces}\plusendmoddef
  interface
     subroutine evt_make_particle_set &
          (evt, factorization_mode, keep_correlations, r)
       import
       class(evt_t), intent(inout) :: evt
       integer, intent(in) :: factorization_mode
       logical, intent(in) :: keep_correlations
       real(default), dimension(:), intent(in), optional :: r
     end subroutine evt_make_particle_set
  end interface
       
\eatline
\nwendcode{}\nwbegindocs{1083}Copy an existing particle set into the event record.  This bypasses
all methods to evaluate the internal state, but may be sufficient for
further processing.
\nwenddocs{}\nwbegincode{1084}\moddef{Event transforms: evt: TBP}\plusendmoddef
  procedure :: set_particle_set => evt_set_particle_set
\nwendcode{}\nwbegincode{1085}\moddef{Event transforms: procedures}\plusendmoddef
  subroutine evt_set_particle_set (evt, particle_set, i_mci, i_term)
    class(evt_t), intent(inout) :: evt
    type(particle_set_t), intent(in) :: particle_set
    integer, intent(in) :: i_term, i_mci
    call evt%prepare_new_event (i_mci, i_term)
    evt%particle_set = particle_set
    evt%particle_set_exists = .true.
  end subroutine evt_set_particle_set
    
\eatline
\nwendcode{}\nwbegindocs{1086}This procedure can help in the previous task, if the particles are
available in the form of an interaction object.  (We need two
interactions, one with color summed over, and one with the probability
distributed among flows.)

We use the two values from the random number generator for factorizing
the state.  For testing purposes, we can provide those numbers explicitly.
\nwenddocs{}\nwbegincode{1087}\moddef{Event transforms: evt: TBP}\plusendmoddef
  procedure :: factorize_interactions => evt_factorize_interactions
\nwendcode{}\nwbegincode{1088}\moddef{Event transforms: procedures}\plusendmoddef
  subroutine evt_factorize_interactions &
       (evt, int_matrix, int_flows, factorization_mode, keep_correlations, r)
    class(evt_t), intent(inout) :: evt
    type(interaction_t), intent(in), target :: int_matrix, int_flows
    integer, intent(in) :: factorization_mode
    logical, intent(in) :: keep_correlations
    real(default), dimension(:), intent(in), optional :: r
    real(default), dimension(2) :: x
    if (present (r)) then
       if (size (r) == 2) then
          x = r
       else
          call msg_bug ("event factorization: size of r array must be 2")
       end if
    else
       call evt%rng%generate (x)
    end if
    call evt%particle_set%init (evt%particle_set_exists, &
         int_matrix, int_flows, factorization_mode, x, &
         keep_correlations, keep_virtual=.true.)
    evt%particle_set_exists = .true.
  end subroutine evt_factorize_interactions
  
\eatline
\nwendcode{}\nwbegindocs{1089}Mark the incoming particles as incoming in the particle set.  This
is necessary because in the interaction objects they are usually
marked as virtual.

In the inquiry functions we set the term index to one; the indices of
beams and incoming particles should be identical for all process
terms.

We use the initial elementary process for obtaining the indices.
Thus, we implicitly assume that the beam and incoming indices stay the
same across event transforms.  If this is not true for a transform
(say, MPI), it should override this method.
\nwenddocs{}\nwbegincode{1090}\moddef{Event transforms: evt: TBP}\plusendmoddef
  procedure :: tag_incoming => evt_tag_incoming
\nwendcode{}\nwbegincode{1091}\moddef{Event transforms: procedures}\plusendmoddef
  subroutine evt_tag_incoming (evt)
    class(evt_t), intent(inout) :: evt
    integer :: i_term, n_in
    integer, dimension(:), allocatable :: beam_index, in_index
    n_in = evt%process%get_n_in ()
    i_term = 1
    allocate (beam_index (n_in))
    call evt%process_instance%get_beam_index (i_term, beam_index)
    call evt%particle_set%reset_status (beam_index, PRT_BEAM)
    allocate (in_index (n_in))
    call evt%process_instance%get_in_index (i_term, in_index)
    call evt%particle_set%reset_status (in_index, PRT_INCOMING)
  end subroutine evt_tag_incoming

\eatline
\nwendcode{}\nwbegindocs{1092}\nwdocspar
\subsection{Implementation: Trivial transform}
This transform contains just a pointer to process and process
instance.  The {\tt{}generate} methods do nothing.
\nwenddocs{}\nwbegincode{1093}\moddef{Event transforms: public}\plusendmoddef
  public :: evt_trivial_t
\nwendcode{}\nwbegincode{1094}\moddef{Event transforms: types}\plusendmoddef
  type, extends (evt_t) :: evt_trivial_t
   contains
   \LA{}Event transforms: evt trivial: TBP\RA{}
  end type evt_trivial_t

\eatline
\nwendcode{}\nwbegindocs{1095}The finalizer is trivial.  Some output:
\nwenddocs{}\nwbegincode{1096}\moddef{Event transforms: evt trivial: TBP}\endmoddef
  procedure :: write => evt_trivial_write
\nwendcode{}\nwbegincode{1097}\moddef{Event transforms: procedures}\plusendmoddef
  subroutine evt_trivial_write (object, unit, verbose, testflag)
    class(evt_trivial_t), intent(in) :: object
    integer, intent(in), optional :: unit
    logical, intent(in), optional :: verbose, testflag
    integer :: u
    u = given_output_unit (unit)
    call write_separator (u, 2)
    write (u, "(1x,A)")  "Event transform: trivial (hard process)"
    call write_separator (u)
    call object%base_write (u, testflag = testflag)
  end subroutine evt_trivial_write
  
\eatline
\nwendcode{}\nwbegindocs{1098}Nothing to do here:
\nwenddocs{}\nwbegincode{1099}\moddef{Event transforms: evt trivial: TBP}\plusendmoddef
  procedure :: prepare_new_event => evt_trivial_prepare_new_event
\nwendcode{}\nwbegincode{1100}\moddef{Event transforms: procedures}\plusendmoddef
  subroutine evt_trivial_prepare_new_event (evt, i_mci, i_term)
    class(evt_trivial_t), intent(inout) :: evt
    integer, intent(in) :: i_mci, i_term
    call evt%reset ()
  end subroutine evt_trivial_prepare_new_event
  
\eatline
\nwendcode{}\nwbegindocs{1101}The weighted generator is, surprisingly, trivial.
\nwenddocs{}\nwbegincode{1102}\moddef{Event transforms: evt trivial: TBP}\plusendmoddef
  procedure :: generate_weighted => evt_trivial_generate_weighted
\nwendcode{}\nwbegincode{1103}\moddef{Event transforms: procedures}\plusendmoddef
  subroutine evt_trivial_generate_weighted (evt, probability)
    class(evt_trivial_t), intent(inout) :: evt
    real(default), intent(out) :: probability
    probability = 1
  end subroutine evt_trivial_generate_weighted
    
\eatline
\nwendcode{}\nwbegindocs{1104}This routine makes a particle set, using the associated process
instance as-is.
\nwenddocs{}\nwbegincode{1105}\moddef{Event transforms: evt trivial: TBP}\plusendmoddef
  procedure :: make_particle_set => evt_trivial_make_particle_set
\nwendcode{}\nwbegincode{1106}\moddef{Event transforms: procedures}\plusendmoddef
  subroutine evt_trivial_make_particle_set &
       (evt, factorization_mode, keep_correlations, r)
    class(evt_trivial_t), intent(inout) :: evt
    integer, intent(in) :: factorization_mode
    logical, intent(in) :: keep_correlations
    real(default), dimension(:), intent(in), optional :: r
    integer :: i_term
    type(interaction_t), pointer :: int_matrix, int_flows
    if (evt%process_instance%is_complete_event ()) then
       call evt%process_instance%select_i_term (i_term)
       int_matrix => evt%process_instance%get_matrix_int_ptr (i_term)
       int_flows  => evt%process_instance%get_flows_int_ptr (i_term)
       call evt%factorize_interactions (int_matrix, int_flows, &
            factorization_mode, keep_correlations, r)
       call evt%tag_incoming ()
    else
       call msg_bug ("Event factorization: event is incomplete")
    end if
  end subroutine evt_trivial_make_particle_set
    
\eatline
\nwendcode{}\nwbegindocs{1107}\nwdocspar
\subsection{Unit tests}
\nwenddocs{}\nwbegincode{1108}\moddef{Event transforms: public}\plusendmoddef
  public :: event_transforms_test
\nwendcode{}\nwbegincode{1109}\moddef{Event transforms: tests}\endmoddef
  subroutine event_transforms_test (u, results)
    integer, intent(in) :: u
    type(test_results_t), intent(inout) :: results
  \LA{}Event transforms: execute tests\RA{}
  end subroutine event_transforms_test
  
\eatline
\nwendcode{}\nwbegindocs{1110}\nwdocspar
\subsubsection{Test trivial event transform}
The trivial transform, as an instance of the abstract transform, does
nothing but to trigger event generation for an elementary process.
\nwenddocs{}\nwbegincode{1111}\moddef{Event transforms: execute tests}\endmoddef
  call test (event_transforms_1, "event_transforms_1", &
       "trivial event transform", &
       u, results)
\nwendcode{}\nwbegincode{1112}\moddef{Event transforms: tests}\plusendmoddef
  subroutine event_transforms_1 (u)
    integer, intent(in) :: u
    type(os_data_t) :: os_data
    type(qcd_t) :: qcd
    class(rng_factory_t), allocatable :: rng_factory
    class(model_data_t), pointer :: model
    type(process_library_t), target :: lib
    type(string_t) :: libname, procname1, run_id
    class(prc_core_t), allocatable :: core_template
    class(mci_t), allocatable :: mci_template
    class(phs_config_t), allocatable :: phs_config_template
    real(default) :: sqrts
    type(process_t), allocatable, target :: process
    type(process_instance_t), allocatable, target :: process_instance
    class(evt_t), allocatable :: evt
    integer :: factorization_mode
    logical :: keep_correlations

    write (u, "(A)")  "* Test output: event_transforms_1"
    write (u, "(A)")  "*   Purpose: handle trivial transform"
    write (u, "(A)")

    write (u, "(A)")  "* Initialize environment and parent process"
    write (u, "(A)")

    call os_data_init (os_data)
    allocate (rng_test_factory_t :: rng_factory)

    libname = "event_transforms_1_lib"
    procname1 = "event_transforms_1_p"
    run_id = "event_transforms_1"

    call prc_test_create_library (libname, lib, &
         scattering = .true., procname1 = procname1)
    call reset_interaction_counter ()

    allocate (model)
    call model%init_test ()

    allocate (process)
    call process%init &
         (procname1, run_id, lib, os_data, qcd, rng_factory, model)
    
    allocate (test_t :: core_template)
    allocate (mci_midpoint_t :: mci_template)
    allocate (phs_single_config_t :: phs_config_template)
    call process%init_component &
         (1, core_template, mci_template, phs_config_template)

    sqrts = 1000
    call process%setup_beams_sqrts (sqrts)
    call process%configure_phs ()
    call process%setup_mci ()
    call process%setup_terms ()

    allocate (process_instance)
    call process_instance%init (process)
    call process%integrate (process_instance, 1, n_it=1, n_calls=100)
    call process%final_integration (1)
    call process_instance%final ()
    deallocate (process_instance)

    allocate (process_instance)
    call process_instance%init (process)
    call process_instance%setup_event_data ()
    call process_instance%init_simulation (1)

    write (u, "(A)")  "* Initialize trivial event transform"
    write (u, "(A)")

    allocate (evt_trivial_t :: evt)
    model => process%get_model_ptr ()
    call evt%connect (process_instance, model)
    
    write (u, "(A)")  "* Generate event and subsequent transform"
    write (u, "(A)")
    
    call process%generate_unweighted_event (process_instance, 1)
    call process_instance%evaluate_event_data ()
    
    call evt%prepare_new_event (1, 1)
    call evt%generate_unweighted ()

    select type (evt)
    type is (evt_trivial_t)
       call write_separator (u, 2)
       call evt%write (u)
       call write_separator (u, 2)
    end select

    write (u, "(A)")
    write (u, "(A)")  "* Obtain particle set"
    write (u, "(A)")
    
    factorization_mode = FM_IGNORE_HELICITY
    keep_correlations = .false.
    
    call evt%make_particle_set (factorization_mode, keep_correlations)

    select type (evt)
    type is (evt_trivial_t)
       call write_separator (u, 2)
       call evt%write (u)
       call write_separator (u, 2)
    end select

    write (u, "(A)")
    write (u, "(A)")  "* Cleanup"

    call evt%final ()
    call process_instance%final ()
    deallocate (process_instance)
    
    write (u, "(A)")
    write (u, "(A)")  "* Test output end: event_transforms_1"
    
  end subroutine event_transforms_1
  
\eatline
\nwendcode{}\nwbegindocs{1113}\nwdocspar
%%%%%%%%%%%%%%%%%%%%%%%%%%%%%%%%%%%%%%%%%%%%%%%%%%%%%%%%%%%%%%%%%%%%%%%%
\section{Decays}
\nwenddocs{}\nwbegincode{1114}\moddef{\code{}decays.f90\edoc{}}\endmoddef
\LA{}File header\RA{}

module decays

\LA{}Use kinds\RA{}
\LA{}Use strings\RA{}
  use io_units
  use format_utils, only: write_indent, write_separator
  use unit_tests
  use diagnostics
  use os_interface
  use sm_qcd
  use flavors
  use helicities
  use quantum_numbers
  use state_matrices
  use interactions
  use evaluators
  use variables
  use model_data
  use rng_base
  use selectors
  use prc_core
  use prc_test
  use process_libraries
  use mci_base
  use mci_midpoint
  use phs_base
  use phs_single
  use parton_states
  use processes
  use process_stacks
  use event_transforms
  
\LA{}Standard module head\RA{}

\LA{}Decays: public\RA{}

\LA{}Decays: types\RA{}

\LA{}Decays: interfaces\RA{}

contains

\LA{}Decays: procedures\RA{}

\LA{}Decays: tests\RA{}

end module decays

\eatline
\nwendcode{}\nwbegindocs{1115}\nwdocspar
\subsection{Final-State Particle Configuration}
A final-state particle may be either stable or unstable.  Here is an
empty abstract type as the parent of both, with holds just the flavor
information.
\nwenddocs{}\nwbegincode{1116}\moddef{Decays: types}\endmoddef
  type, abstract :: any_config_t
   contains
   \LA{}Decays: any config: TBP\RA{}
  end type any_config_t

\eatline
\nwendcode{}\nwbegindocs{1117}Finalizer, depends on the implementation.
\nwenddocs{}\nwbegincode{1118}\moddef{Decays: any config: TBP}\endmoddef
  procedure (any_config_final), deferred :: final
\nwendcode{}\nwbegincode{1119}\moddef{Decays: interfaces}\endmoddef
  interface
     subroutine any_config_final (object)
       import
       class(any_config_t), intent(inout) :: object
     end subroutine any_config_final
  end interface
  
\eatline
\nwendcode{}\nwbegindocs{1120}The output is also deferred:
\nwenddocs{}\nwbegincode{1121}\moddef{Decays: any config: TBP}\plusendmoddef
  procedure (any_config_write), deferred :: write
\nwendcode{}\nwbegincode{1122}\moddef{Decays: interfaces}\plusendmoddef
  interface
     subroutine any_config_write (object, unit, indent, verbose)
       import
       class(any_config_t), intent(in) :: object
       integer, intent(in), optional :: unit, indent
       logical, intent(in), optional :: verbose
     end subroutine any_config_write
  end interface

\eatline
\nwendcode{}\nwbegindocs{1123}This is a container for a stable or unstable particle configurator.
We need this wrapper for preparing arrays that mix stable and unstable
particles.
\nwenddocs{}\nwbegincode{1124}\moddef{Decays: types}\plusendmoddef
  type :: particle_config_t
     class(any_config_t), allocatable :: c
  end type particle_config_t
  
\eatline
\nwendcode{}\nwbegindocs{1125}\nwdocspar
\subsection{Final-State Particle}
In theory, for the particle instance we only need to consider the
unstable case.  However, it is more straightforward to treat
configuration and instance on the same footing, and to introduce a
wrapper for particle objects as above.  This also works
around a compiler bug in gfortran.
\nwenddocs{}\nwbegincode{1126}\moddef{Decays: types}\plusendmoddef
  type, abstract :: any_t
   contains
   \LA{}Decays: any: TBP\RA{}
  end type any_t
  
\eatline
\nwendcode{}\nwbegindocs{1127}Finalizer, depends on the implementation.
\nwenddocs{}\nwbegincode{1128}\moddef{Decays: any: TBP}\endmoddef
  procedure (any_final), deferred :: final
\nwendcode{}\nwbegincode{1129}\moddef{Decays: interfaces}\plusendmoddef
  interface
     subroutine any_final (object)
       import
       class(any_t), intent(inout) :: object
     end subroutine any_final
  end interface
  
\eatline
\nwendcode{}\nwbegindocs{1130}The output is also deferred:
\nwenddocs{}\nwbegincode{1131}\moddef{Decays: any: TBP}\plusendmoddef
  procedure (any_write), deferred :: write
\nwendcode{}\nwbegincode{1132}\moddef{Decays: interfaces}\plusendmoddef
  interface
     subroutine any_write (object, unit, indent)
       import
       class(any_t), intent(in) :: object
       integer, intent(in), optional :: unit, indent
     end subroutine any_write
  end interface

\eatline
\nwendcode{}\nwbegindocs{1133}This is a container for a stable or unstable outgoing particle.
We need this wrapper for preparing arrays that mix stable and unstable
particles.
\nwenddocs{}\nwbegincode{1134}\moddef{Decays: types}\plusendmoddef
  type :: particle_out_t
     class(any_t), allocatable :: c
  end type particle_out_t
  
\eatline
\nwendcode{}\nwbegindocs{1135}\nwdocspar
\subsection{Decay Term Configuration}
A decay term is a distinct final state, corresponding to a process
term.  Each decay process may give rise to several terms with,
possibly, differing flavor content.
\nwenddocs{}\nwbegincode{1136}\moddef{Decays: types}\plusendmoddef
  type :: decay_term_config_t
     type(particle_config_t), dimension(:), allocatable :: prt
   contains
   \LA{}Decays: decay term config: TBP\RA{}
  end type decay_term_config_t
  
\eatline
\nwendcode{}\nwbegindocs{1137}Finalizer, recursive.
\nwenddocs{}\nwbegincode{1138}\moddef{Decays: decay term config: TBP}\endmoddef
  procedure :: final => decay_term_config_final
\nwendcode{}\nwbegincode{1139}\moddef{Decays: procedures}\endmoddef
  recursive subroutine decay_term_config_final (object)
    class(decay_term_config_t), intent(inout) :: object
    integer :: i
    if (allocated (object%prt)) then
       do i = 1, size (object%prt)
          if (allocated (object%prt(i)%c))  call object%prt(i)%c%final ()
       end do
    end if
  end subroutine decay_term_config_final
  
\eatline
\nwendcode{}\nwbegindocs{1140}Output, with optional indentation
\nwenddocs{}\nwbegincode{1141}\moddef{Decays: decay term config: TBP}\plusendmoddef
  procedure :: write => decay_term_config_write
\nwendcode{}\nwbegincode{1142}\moddef{Decays: procedures}\plusendmoddef
  recursive subroutine decay_term_config_write (object, unit, indent, verbose)
    class(decay_term_config_t), intent(in) :: object
    integer, intent(in), optional :: unit, indent
    logical, intent(in), optional :: verbose
    integer :: i, j, u, ind
    logical :: verb
    u = given_output_unit (unit)
    ind = 0;  if (present (indent))  ind = indent
    verb = .true.;  if (present (verbose))  verb = verbose
    call write_indent (u, ind)
    write (u, "(1x,A)", advance="no")  "Final state:"
    do i = 1, size (object%prt)
       select type (prt_config => object%prt(i)%c)
       type is (stable_config_t)
          write (u, "(1x,A)", advance="no") &
               char (prt_config%flv(1)%get_name ())
          do j = 2, size (prt_config%flv)
             write (u, "(':',A)", advance="no") &
                  char (prt_config%flv(j)%get_name ())
          end do
       type is (unstable_config_t)
          write (u, "(1x,A)", advance="no") &
               char (prt_config%flv%get_name ())
       end select
    end do
    write (u, *)
    if (verb) then
       do i = 1, size (object%prt)
          call object%prt(i)%c%write (u, ind)
       end do
    end if
  end subroutine decay_term_config_write

\eatline
\nwendcode{}\nwbegindocs{1143}Initialize, given a set of flavors.  For each flavor, we must indicate
whether the particle is stable.  The second index of the flavor array runs
over alternatives for each decay product; alternatives are allowed only if the
decay product is itself stable.
\nwenddocs{}\nwbegincode{1144}\moddef{Decays: decay term config: TBP}\plusendmoddef
  procedure :: init => decay_term_config_init
\nwendcode{}\nwbegincode{1145}\moddef{Decays: procedures}\plusendmoddef
  recursive subroutine decay_term_config_init &
       (term, flv, stable, model, process_stack)
    class(decay_term_config_t), intent(out) :: term
    type(flavor_t), dimension(:,:), intent(in) :: flv
    logical, dimension(:), intent(in) :: stable
    class(model_data_t), intent(in), target :: model
    type(process_stack_t), intent(in), optional :: process_stack
    type(string_t), dimension(:), allocatable :: decay
    integer :: i
    allocate (term%prt (size (flv, 1)))
    do i = 1, size (flv, 1)
       associate (prt => term%prt(i))
         if (stable(i)) then
            allocate (stable_config_t :: prt%c)
         else
            allocate (unstable_config_t :: prt%c)
         end if
         select type (prt_config => prt%c)
         type is (stable_config_t)
            call prt_config%init (flv(i,:))
         type is (unstable_config_t)
            if (all (flv(i,:) == flv(i,1))) then
               call prt_config%init (flv(i,1))
               call flv(i,1)%get_decays (decay)
               call prt_config%init_decays (decay, model, process_stack)
            else
               call prt_config%write ()
               call msg_fatal ("Decay configuration: &
                    &unstable product must be unique")
            end if
         end select
       end associate
    end do
  end subroutine decay_term_config_init
  
\eatline
\nwendcode{}\nwbegindocs{1146}Recursively compute widths and branching ratios for all unstable particles.
\nwenddocs{}\nwbegincode{1147}\moddef{Decays: decay term config: TBP}\plusendmoddef
  procedure :: compute => decay_term_config_compute
\nwendcode{}\nwbegincode{1148}\moddef{Decays: procedures}\plusendmoddef
  recursive subroutine decay_term_config_compute (term)
    class(decay_term_config_t), intent(inout) :: term
    integer :: i
    do i = 1, size (term%prt)
       select type (unstable_config => term%prt(i)%c)
       type is (unstable_config_t)
          call unstable_config%compute ()
       end select
    end do
  end subroutine decay_term_config_compute
  
\eatline
\nwendcode{}\nwbegindocs{1149}\nwdocspar
\subsection{Decay Term}
A decay term instance is selected when we generate an event for the associated
process instance.  When evaluated, it triggers further decays down the chain.

Only unstable products are allocated as child particles.
\nwenddocs{}\nwbegincode{1150}\moddef{Decays: types}\plusendmoddef
  type :: decay_term_t
     type(decay_term_config_t), pointer :: config => null ()
     type(particle_out_t), dimension(:), allocatable :: particle_out
   contains
   \LA{}Decays: decay term: TBP\RA{}
  end type decay_term_t

\eatline
\nwendcode{}\nwbegindocs{1151}Finalizer.
\nwenddocs{}\nwbegincode{1152}\moddef{Decays: decay term: TBP}\endmoddef
  procedure :: final => decay_term_final
\nwendcode{}\nwbegincode{1153}\moddef{Decays: procedures}\plusendmoddef
  recursive subroutine decay_term_final (object)
    class(decay_term_t), intent(inout) :: object
    integer :: i
    if (allocated (object%particle_out)) then
       do i = 1, size (object%particle_out)
          call object%particle_out(i)%c%final ()
       end do
    end if
  end subroutine decay_term_final
  
\eatline
\nwendcode{}\nwbegindocs{1154}Output.
\nwenddocs{}\nwbegincode{1155}\moddef{Decays: decay term: TBP}\plusendmoddef
  procedure :: write => decay_term_write
\nwendcode{}\nwbegincode{1156}\moddef{Decays: procedures}\plusendmoddef
  recursive subroutine decay_term_write (object, unit, indent)
    class(decay_term_t), intent(in) :: object
    integer, intent(in), optional :: unit, indent
    integer :: i, u, ind
    u = given_output_unit (unit)
    ind = 0;  if (present (indent))  ind = indent
    call object%config%write (u, ind, verbose = .false.)
    do i = 1, size (object%particle_out)
       call object%particle_out(i)%c%write (u, ind)
    end do
  end subroutine decay_term_write

\eatline
\nwendcode{}\nwbegindocs{1157}Recursively write the embedded process instances.
\nwenddocs{}\nwbegincode{1158}\moddef{Decays: decay term: TBP}\plusendmoddef
  procedure :: write_process_instances => decay_term_write_process_instances
\nwendcode{}\nwbegincode{1159}\moddef{Decays: procedures}\plusendmoddef
  recursive subroutine decay_term_write_process_instances (term, unit, verbose)
    class(decay_term_t), intent(in) :: term
    integer, intent(in), optional :: unit
    logical, intent(in), optional :: verbose
    integer :: i
    do i = 1, size (term%particle_out)
       select type (unstable => term%particle_out(i)%c)
       type is (unstable_t)
          call unstable%write_process_instances (unit, verbose)
       end select
    end do
  end subroutine decay_term_write_process_instances
  
\eatline
\nwendcode{}\nwbegindocs{1160}Initialization, using the configuration object.  We allocate
particle objects in parallel to the particle configuration objects
which we use to initialize them, one at a time.
\nwenddocs{}\nwbegincode{1161}\moddef{Decays: decay term: TBP}\plusendmoddef
  procedure :: init => decay_term_init
\nwendcode{}\nwbegincode{1162}\moddef{Decays: procedures}\plusendmoddef
  recursive subroutine decay_term_init (term, config)
    class(decay_term_t), intent(out) :: term
    type(decay_term_config_t), intent(in), target :: config
    integer :: i
    term%config => config
    allocate (term%particle_out (size (config%prt)))
    do i = 1, size (config%prt)
       select type (prt_config => config%prt(i)%c)
       type is (stable_config_t)
          allocate (stable_t :: term%particle_out(i)%c)
          select type (stable => term%particle_out(i)%c)
          type is (stable_t)
             call stable%init (prt_config)
          end select
       type is (unstable_config_t)
          allocate (unstable_t :: term%particle_out(i)%c)
          select type (unstable => term%particle_out(i)%c)
          type is (unstable_t)
             call unstable%init (prt_config)
          end select
       end select
    end do
  end subroutine decay_term_init

\eatline
\nwendcode{}\nwbegindocs{1163}Implement a RNG instance, spawned by the process object.
\nwenddocs{}\nwbegincode{1164}\moddef{Decays: decay term: TBP}\plusendmoddef
  procedure :: make_rng => decay_term_make_rng
\nwendcode{}\nwbegincode{1165}\moddef{Decays: procedures}\plusendmoddef
  subroutine decay_term_make_rng (term, process)
    class(decay_term_t), intent(inout) :: term
    type(process_t), intent(inout) :: process
    class(rng_t), allocatable :: rng
    integer :: i
    do i = 1, size (term%particle_out)
       select type (unstable => term%particle_out(i)%c)
       type is (unstable_t)
          call process%make_rng (rng)
          call unstable%import_rng (rng)
       end select
    end do
  end subroutine decay_term_make_rng
    
\eatline
\nwendcode{}\nwbegindocs{1166}Link the interactions for unstable decay products to the
interaction of the parent process.
\nwenddocs{}\nwbegincode{1167}\moddef{Decays: decay term: TBP}\plusendmoddef
  procedure :: link_interactions => decay_term_link_interactions
\nwendcode{}\nwbegincode{1168}\moddef{Decays: procedures}\plusendmoddef
  recursive subroutine decay_term_link_interactions (term, trace)
    class(decay_term_t), intent(inout) :: term
    type(interaction_t), intent(in), target :: trace
    integer :: i
    do i = 1, size (term%particle_out)
       select type (unstable => term%particle_out(i)%c)
       type is (unstable_t)
          call unstable%link_interactions (i, trace)
       end select
    end do
  end subroutine decay_term_link_interactions
  
\eatline
\nwendcode{}\nwbegindocs{1169}Recursively generate a decay chain, for each of the unstable
particles in the final state.
\nwenddocs{}\nwbegincode{1170}\moddef{Decays: decay term: TBP}\plusendmoddef
  procedure :: select_chain => decay_term_select_chain
\nwendcode{}\nwbegincode{1171}\moddef{Decays: procedures}\plusendmoddef
  recursive subroutine decay_term_select_chain (term)
    class(decay_term_t), intent(inout) :: term
    integer :: i
    do i = 1, size (term%particle_out)
       select type (unstable => term%particle_out(i)%c)
       type is (unstable_t)
          call unstable%select_chain ()
       end select
    end do
  end subroutine decay_term_select_chain

\eatline
\nwendcode{}\nwbegindocs{1172}Recursively generate a decay event, for each of the unstable
particles in the final state.
\nwenddocs{}\nwbegincode{1173}\moddef{Decays: decay term: TBP}\plusendmoddef
  procedure :: generate => decay_term_generate
\nwendcode{}\nwbegincode{1174}\moddef{Decays: procedures}\plusendmoddef
  recursive subroutine decay_term_generate (term)
    class(decay_term_t), intent(inout) :: term
    integer :: i
    do i = 1, size (term%particle_out)
       select type (unstable => term%particle_out(i)%c)
       type is (unstable_t)
          call unstable%generate ()
       end select
    end do
  end subroutine decay_term_generate

\eatline
\nwendcode{}\nwbegindocs{1175}\nwdocspar
\subsection{Decay Root Configuration}
At the root of a decay chain, there is a parent process.  The decay root
stores a pointer to the parent process and the set of decay configurations.
\nwenddocs{}\nwbegincode{1176}\moddef{Decays: types}\plusendmoddef
  type :: decay_root_config_t
     type(string_t) :: process_id
     type(process_t), pointer :: process => null ()
     class(model_data_t), pointer :: model => null ()
     type(decay_term_config_t), dimension(:), allocatable :: term_config
   contains
   \LA{}Decays: decay root config: TBP\RA{}
  end type decay_root_config_t
  
\eatline
\nwendcode{}\nwbegindocs{1177}The finalizer is recursive since there may be cascade decays.
\nwenddocs{}\nwbegincode{1178}\moddef{Decays: decay root config: TBP}\endmoddef
  procedure :: final => decay_root_config_final
\nwendcode{}\nwbegincode{1179}\moddef{Decays: procedures}\plusendmoddef
  recursive subroutine decay_root_config_final (object)
    class(decay_root_config_t), intent(inout) :: object
    integer :: i
    if (allocated (object%term_config)) then
       do i = 1, size (object%term_config)
          call object%term_config(i)%final ()
       end do
    end if
  end subroutine decay_root_config_final
  
\eatline
\nwendcode{}\nwbegindocs{1180}The output routine is also recursive, and it contains an adjustable
indentation.
\nwenddocs{}\nwbegincode{1181}\moddef{Decays: decay root config: TBP}\plusendmoddef
  procedure :: write => decay_root_config_write
  procedure :: write_header => decay_root_config_write_header
  procedure :: write_terms => decay_root_config_write_terms
\nwendcode{}\nwbegincode{1182}\moddef{Decays: procedures}\plusendmoddef
  recursive subroutine decay_root_config_write (object, unit, indent, verbose)
    class(decay_root_config_t), intent(in) :: object
    integer, intent(in), optional :: unit, indent
    logical, intent(in), optional :: verbose
    integer :: u, ind
    u = given_output_unit (unit)
    ind = 0;  if (present (indent))  ind = indent
    call write_indent (u, ind)
    write (u, "(1x,A)")  "Final-state decay tree:"
    call object%write_header (unit, indent)
    call object%write_terms (unit, indent, verbose)
  end subroutine decay_root_config_write

  subroutine decay_root_config_write_header (object, unit, indent)
    class(decay_root_config_t), intent(in) :: object
    integer, intent(in), optional :: unit, indent
    integer :: u, ind
    u = given_output_unit (unit)
    ind = 0;  if (present (indent))  ind = indent
    call write_indent (u, ind)
    if (associated (object%process)) then
       write (u, 3)  "process ID      =", char (object%process_id), "*"
    else
       write (u, 3)  "process ID      =", char (object%process_id)
    end if
3   format (3x,A,2(1x,A))
  end subroutine decay_root_config_write_header
    
  recursive subroutine decay_root_config_write_terms &
       (object, unit, indent, verbose)
    class(decay_root_config_t), intent(in) :: object
    integer, intent(in), optional :: unit, indent
    logical, intent(in), optional :: verbose
    integer :: i, u, ind
    logical :: verb
    u = given_output_unit (unit)
    ind = 0;  if (present (indent))  ind = indent
    verb = .true.;  if (present (verbose))  verb = verbose
    if (verb .and. allocated (object%term_config)) then
       do i = 1, size (object%term_config)
          call object%term_config(i)%write (u, ind + 1)
       end do
    end if
  end subroutine decay_root_config_write_terms
    
\eatline
\nwendcode{}\nwbegindocs{1183}Initialize for a named process and (optionally) a pre-determined
number of terms.
\nwenddocs{}\nwbegincode{1184}\moddef{Decays: decay root config: TBP}\plusendmoddef
  procedure :: init => decay_root_config_init
\nwendcode{}\nwbegincode{1185}\moddef{Decays: procedures}\plusendmoddef
  subroutine decay_root_config_init (decay, model, process_id, n_terms)
    class(decay_root_config_t), intent(out) :: decay
    class(model_data_t), intent(in), target :: model
    type(string_t), intent(in) :: process_id
    integer, intent(in), optional :: n_terms
    decay%model => model
    decay%process_id = process_id
    if (present (n_terms)) then
       allocate (decay%term_config (n_terms))
    end if
  end subroutine decay_root_config_init
       
\eatline
\nwendcode{}\nwbegindocs{1186}Declare a decay term, given an array of flavors.
\nwenddocs{}\nwbegincode{1187}\moddef{Decays: decay root config: TBP}\plusendmoddef
  procedure :: init_term => decay_root_config_init_term
\nwendcode{}\nwbegincode{1188}\moddef{Decays: procedures}\plusendmoddef
  recursive subroutine decay_root_config_init_term &
       (decay, i, flv, stable, model, process_stack)
    class(decay_root_config_t), intent(inout) :: decay
    integer, intent(in) :: i
    type(flavor_t), dimension(:,:), intent(in) :: flv
    logical, dimension(:), intent(in) :: stable
    class(model_data_t), intent(in), target :: model
    type(process_stack_t), intent(in), optional :: process_stack
    call decay%term_config(i)%init (flv, stable, model, process_stack)
  end subroutine decay_root_config_init_term
  
\eatline
\nwendcode{}\nwbegindocs{1189}Connect the decay root configuration with a process object (which should
represent the parent process).  This includes initialization, therefore
intent(out).

The flavor state is retrieved from the process term object.  However, we have
to be careful: the flavor object points to the model instance that is stored
in the process object.  This model instance may not contain the current
setting for unstable particles and decay.  Therefore, we assign the model
directly.

If the {\tt{}process{\char95}instance} argument is provided, we use this for the
flavor state.  This applies to the decay root only, where the process
can be entangled with a beam setup, and the latter contains beam
remnants as further outgoing particles.  These must be included in the
set of outgoing flavors, since the decay application is also done on
the connected state.

Infer stability from the particle properties, using the first row in the set
of flavor states.  For unstable particles, we look for decays,
recursively, available from the process stack (if present).
\nwenddocs{}\nwbegincode{1190}\moddef{Decays: decay root config: TBP}\plusendmoddef
  procedure :: connect => decay_root_config_connect
\nwendcode{}\nwbegincode{1191}\moddef{Decays: procedures}\plusendmoddef
  recursive subroutine decay_root_config_connect &
       (decay, process, model, process_stack, process_instance)
    class(decay_root_config_t), intent(out) :: decay
    type(process_t), intent(in), target :: process
    class(model_data_t), intent(in), target :: model
    type(process_stack_t), intent(in), optional :: process_stack
    type(process_instance_t), intent(in), optional, target :: process_instance
    type(connected_state_t), pointer :: connected_state
    type(interaction_t), pointer :: int
    type(flavor_t), dimension(:,:), allocatable :: flv
    logical, dimension(:), allocatable :: stable
    integer :: i
    call decay%init (model, process%get_id (), process%get_n_terms ())
    do i = 1, size (decay%term_config)
       if (present (process_instance)) then
          connected_state => process_instance%get_connected_state_ptr (i)
          int => connected_state%get_matrix_int_ptr ()
          call interaction_get_flv_out (int, flv)
       else
          call process%get_term_flv_out (i, flv)
       end if
       call flv%set_model (model)
       allocate (stable (size (flv, 1)))
       stable = flv(:,1)%is_stable ()
       call decay%init_term (i, flv, stable, model, process_stack)
       deallocate (flv, stable)
    end do
    decay%process => process
  end subroutine decay_root_config_connect

\eatline
\nwendcode{}\nwbegindocs{1192}Recursively compute widths, errors, and branching ratios.
\nwenddocs{}\nwbegincode{1193}\moddef{Decays: decay root config: TBP}\plusendmoddef
  procedure :: compute => decay_root_config_compute
\nwendcode{}\nwbegincode{1194}\moddef{Decays: procedures}\plusendmoddef
  recursive subroutine decay_root_config_compute (decay)
    class(decay_root_config_t), intent(inout) :: decay
    integer :: i
    do i = 1, size (decay%term_config)
       call decay%term_config(i)%compute ()
    end do
  end subroutine decay_root_config_compute
  
\eatline
\nwendcode{}\nwbegindocs{1195}\nwdocspar
\subsection{Decay Root Instance}
This is the common parent type for decay and decay root.  The process instance
points to the parent process.  The model pointer is separate because particle
settings may be updated w.r.t.\ the parent process object.
\nwenddocs{}\nwbegincode{1196}\moddef{Decays: types}\plusendmoddef
  type, abstract :: decay_gen_t
     type(decay_term_t), dimension(:), allocatable :: term
     type(process_instance_t), pointer :: process_instance => null ()
     integer :: selected_mci = 0
     integer :: selected_term = 0
   contains
   \LA{}Decays: decay gen: TBP\RA{}
  end type decay_gen_t
  
\eatline
\nwendcode{}\nwbegindocs{1197}\nwdocspar
The decay root represents the parent process.  When an event is generated, the
generator selects the term to which the decay chain applies (if possible).

The process instance is just a pointer.
\nwenddocs{}\nwbegincode{1198}\moddef{Decays: types}\plusendmoddef
  type, extends (decay_gen_t) :: decay_root_t
     type(decay_root_config_t), pointer :: config => null ()
   contains
   \LA{}Decays: decay root: TBP\RA{}
  end type decay_root_t
  
\eatline
\nwendcode{}\nwbegindocs{1199}The finalizer has to recursively finalize the terms, but we can skip the
process instance which is not explicitly allocated.
\nwenddocs{}\nwbegincode{1200}\moddef{Decays: decay gen: TBP}\endmoddef
  procedure :: base_final => decay_gen_final
\nwendcode{}\nwbegincode{1201}\moddef{Decays: procedures}\plusendmoddef
  recursive subroutine decay_gen_final (object)
    class(decay_gen_t), intent(inout) :: object
    integer :: i
    if (allocated (object%term)) then
       do i = 1, size (object%term)
          call object%term(i)%final ()
       end do
    end if
  end subroutine decay_gen_final    
  
\eatline
\nwendcode{}\nwbegindocs{1202}No extra finalization for the decay root.
\nwenddocs{}\nwbegincode{1203}\moddef{Decays: decay root: TBP}\endmoddef
  procedure :: final => decay_root_final
\nwendcode{}\nwbegincode{1204}\moddef{Decays: procedures}\plusendmoddef
  subroutine decay_root_final (object)
    class(decay_root_t), intent(inout) :: object
    call object%base_final ()
  end subroutine decay_root_final    
  
\eatline
\nwendcode{}\nwbegindocs{1205}Output.
\nwenddocs{}\nwbegincode{1206}\moddef{Decays: decay root: TBP}\plusendmoddef
  procedure :: write => decay_root_write
\nwendcode{}\nwbegincode{1207}\moddef{Decays: procedures}\plusendmoddef
  subroutine decay_root_write (object, unit)
    class(decay_root_t), intent(in) :: object
    integer, intent(in), optional :: unit
    integer :: u
    u = given_output_unit (unit)
    if (associated (object%config)) then
       call object%config%write (unit, verbose = .false.)
    else
       write (u, "(1x,A)")  "Final-state decay tree: [not configured]"
    end if
    if (object%selected_mci > 0) then
       write (u, "(3x,A,I0)")  "Selected MCI    = ", object%selected_mci
    else
       write (u, "(3x,A)")  "Selected MCI    = [undefined]"
    end if
    if (object%selected_term > 0) then
       write (u, "(3x,A,I0)")  "Selected term   = ", object%selected_term
       call object%term(object%selected_term)%write (u, 1)
    else
       write (u, "(3x,A)")  "Selected term   = [undefined]"
    end if
  end subroutine decay_root_write

\eatline
\nwendcode{}\nwbegindocs{1208}Write the process instances, recursively.
\nwenddocs{}\nwbegincode{1209}\moddef{Decays: decay gen: TBP}\plusendmoddef
  procedure :: write_process_instances => decay_gen_write_process_instances
\nwendcode{}\nwbegincode{1210}\moddef{Decays: procedures}\plusendmoddef
  recursive subroutine decay_gen_write_process_instances (decay, unit, verbose)
    class(decay_gen_t), intent(in) :: decay
    integer, intent(in), optional :: unit
    logical, intent(in), optional :: verbose
    logical :: verb
    verb = .true.;  if (present (verbose))  verb = verbose
    if (associated (decay%process_instance)) then
       if (verb) then
          call decay%process_instance%write (unit)
       else
          call decay%process_instance%write_header (unit)
       end if
    end if
    if (decay%selected_term > 0) then
       call decay%term(decay%selected_term)%write_process_instances (unit, verb)
    end if
  end subroutine decay_gen_write_process_instances
    
\eatline
\nwendcode{}\nwbegindocs{1211}Generic initializer.  All can be done recursively.
\nwenddocs{}\nwbegincode{1212}\moddef{Decays: decay gen: TBP}\plusendmoddef
  procedure :: base_init => decay_gen_init
\nwendcode{}\nwbegincode{1213}\moddef{Decays: procedures}\plusendmoddef
  recursive subroutine decay_gen_init (decay, term_config)
    class(decay_gen_t), intent(out) :: decay
    type(decay_term_config_t), dimension(:), intent(in), target :: term_config
    integer :: i
    allocate (decay%term (size (term_config)))
    do i = 1, size (decay%term)
       call decay%term(i)%init (term_config(i))
    end do
  end subroutine decay_gen_init

\eatline
\nwendcode{}\nwbegindocs{1214}Specific initializer.  We assign the configuration object, which should
correspond to a completely initialized decay configuration tree.  We
also connect to an existing process instance.  Then, we recursively
link the child interactions to the parent process.
\nwenddocs{}\nwbegincode{1215}\moddef{Decays: decay root: TBP}\plusendmoddef
  procedure :: init => decay_root_init
\nwendcode{}\nwbegincode{1216}\moddef{Decays: procedures}\plusendmoddef
  subroutine decay_root_init (decay_root, config, process_instance)
    class(decay_root_t), intent(out) :: decay_root
    type(decay_root_config_t), intent(in), target :: config
    type(process_instance_t), intent(in), target :: process_instance
    call decay_root%base_init (config%term_config)
    decay_root%config => config
    decay_root%process_instance => process_instance
    call decay_root%make_term_rng (config%process)
    call decay_root%link_term_interactions ()
  end subroutine decay_root_init

\eatline
\nwendcode{}\nwbegindocs{1217}Implement random-number generators for unstable decay selection in
all terms.  This is not recursive.

We also make use of the fact that {\tt{}process} is a pointer; the (state
of the RNG factory inside the) target process will be modified by the
rng-spawning method, but not the pointer.
\nwenddocs{}\nwbegincode{1218}\moddef{Decays: decay gen: TBP}\plusendmoddef
  procedure :: make_term_rng => decay_gen_make_term_rng
\nwendcode{}\nwbegincode{1219}\moddef{Decays: procedures}\plusendmoddef
  subroutine decay_gen_make_term_rng (decay, process)
    class(decay_gen_t), intent(inout) :: decay
    type(process_t), intent(in), pointer :: process
    integer :: i
    do i = 1, size (decay%term)
       call decay%term(i)%make_rng (process)
    end do
  end subroutine decay_gen_make_term_rng
    
\eatline
\nwendcode{}\nwbegindocs{1220}Recursively link interactions of the enclosed decay terms to the
corresponding terms in the current process instance.

Note: A bug in nagfor requires the extra {\tt{}i{\char95}term} variable.
\nwenddocs{}\nwbegincode{1221}\moddef{Decays: decay gen: TBP}\plusendmoddef
  procedure :: link_term_interactions => decay_gen_link_term_interactions
\nwendcode{}\nwbegincode{1222}\moddef{Decays: procedures}\plusendmoddef
  recursive subroutine decay_gen_link_term_interactions (decay)
    class(decay_gen_t), intent(inout) :: decay
    integer :: i, i_term
    type(interaction_t), pointer :: trace
    associate (instance => decay%process_instance)
      do i = 1, size (decay%term)
         i_term = i
         trace => instance%get_trace_int_ptr (i_term)
         call decay%term(i_term)%link_interactions (trace)
      end do
    end associate
  end subroutine decay_gen_link_term_interactions

\eatline
\nwendcode{}\nwbegindocs{1223}Select a decay chain: decay modes and process components.
\nwenddocs{}\nwbegincode{1224}\moddef{Decays: decay root: TBP}\plusendmoddef
  procedure :: select_chain => decay_root_select_chain
\nwendcode{}\nwbegincode{1225}\moddef{Decays: procedures}\plusendmoddef
  subroutine decay_root_select_chain (decay_root)
    class(decay_root_t), intent(inout) :: decay_root
    if (decay_root%selected_term > 0) then
       call decay_root%term(decay_root%selected_term)%select_chain ()
    else
       call msg_bug ("Decays: no term selected for parent process")
    end if
  end subroutine decay_root_select_chain

\eatline
\nwendcode{}\nwbegindocs{1226}Generate a decay tree, i.e., for the selected term in the parent
process, recursively generate a decay event for all unstable
particles.

Factor out the trace of the connected state of the parent process.
This trace should not be taken into account for unweighting the decay
chain, since it was already used for unweighting the parent event, or
it determines the overall event weight.
\nwenddocs{}\nwbegincode{1227}\moddef{Decays: decay root: TBP}\plusendmoddef
  procedure :: generate => decay_root_generate
\nwendcode{}\nwbegincode{1228}\moddef{Decays: procedures}\plusendmoddef
  subroutine decay_root_generate (decay_root)
    class(decay_root_t), intent(inout) :: decay_root
    type(connected_state_t), pointer :: connected_state
    if (decay_root%selected_term > 0) then
       connected_state => decay_root%process_instance%get_connected_state_ptr &
            (decay_root%selected_term)
       call connected_state%normalize_matrix_by_trace ()
       call decay_root%term(decay_root%selected_term)%generate ()
    else
       call msg_bug ("Decays: no term selected for parent process")
    end if
  end subroutine decay_root_generate

\eatline
\nwendcode{}\nwbegindocs{1229}\nwdocspar
\subsection{Decay Configuration}
A decay configuration describes a distinct decay mode of a particle.  Each
decay mode may include several terms, which correspond to the terms in the
associated process.  In addition to the base type, the decay configuration
object contains the integral of the parent process and the selector for the
MCI group inside this process.

The flavor component should be identical to the flavor component of the parent
particle ({\tt{}unstable} object).
\nwenddocs{}\nwbegincode{1230}\moddef{Decays: types}\plusendmoddef
  type, extends (decay_root_config_t) :: decay_config_t
     type(flavor_t) :: flv
     real(default) :: weight = 0
     real(default) :: integral = 0
     real(default) :: abs_error = 0
     real(default) :: rel_error = 0
     type(selector_t) :: mci_selector
   contains
   \LA{}Decays: decay config: TBP\RA{}
  end type decay_config_t
  
\eatline
\nwendcode{}\nwbegindocs{1231}The output routine extends the decay-root writer by listing numerical
component values.
\nwenddocs{}\nwbegincode{1232}\moddef{Decays: decay config: TBP}\endmoddef
  procedure :: write => decay_config_write
\nwendcode{}\nwbegincode{1233}\moddef{Decays: procedures}\plusendmoddef
  recursive subroutine decay_config_write (object, unit, indent, verbose)
    class(decay_config_t), intent(in) :: object
    integer, intent(in), optional :: unit, indent
    logical, intent(in), optional :: verbose
    integer :: u, ind
    u = given_output_unit (unit)
    ind = 0;  if (present (indent))  ind = indent
    call write_indent (u, ind)
    write (u, "(1x,A)")  "Decay:"
    call object%write_header (unit, indent)
    call write_indent (u, ind)
    write (u, 2)  "branching ratio =", object%weight * 100
    call write_indent (u, ind)
    write (u, 1)  "partial width   =", object%integral
    call write_indent (u, ind)
    write (u, 1)  "error (abs)     =", object%abs_error
    call write_indent (u, ind)
    write (u, 1)  "error (rel)     =", object%rel_error
1   format (3x,A,ES19.12)
2   format (3x,A,F11.6,1x,'%')
    call object%write_terms (unit, indent, verbose)
  end subroutine decay_config_write
    
\eatline
\nwendcode{}\nwbegindocs{1234}Connect a decay configuration with a process object (which should
represent the decay).  This includes initialization, therefore
intent(out).  We first connect the process itself, then do initializations
that are specific for this decay.

Infer stability from the particle properties, using the first row in the set
of flavor states.  Once we can deal with predetermined decay chains, they
should be used instead.
\nwenddocs{}\nwbegincode{1235}\moddef{Decays: decay config: TBP}\plusendmoddef
  procedure :: connect => decay_config_connect
\nwendcode{}\nwbegincode{1236}\moddef{Decays: procedures}\plusendmoddef
  recursive subroutine decay_config_connect &
       (decay, process, model, process_stack, process_instance)
    class(decay_config_t), intent(out) :: decay
    type(process_t), intent(in), target :: process
    class(model_data_t), intent(in), target :: model
    type(process_stack_t), intent(in), optional :: process_stack
    type(process_instance_t), intent(in), optional, target :: process_instance
    real(default), dimension(:), allocatable :: integral_mci
    integer :: i, n_mci
    call decay%decay_root_config_t%connect (process, model, process_stack)
    decay%integral = process%get_integral ()
    decay%abs_error = process%get_error ()
    if (process%cm_frame ()) then
       call msg_fatal ("Decay process " // char (process%get_id ()) &
            // ": unusable because rest frame is fixed.")
    end if
    n_mci = process%get_n_mci ()
    allocate (integral_mci (n_mci))
    do i = 1, n_mci
       integral_mci(i) = process%get_integral_mci (i)
    end do
    call decay%mci_selector%init (integral_mci)
  end subroutine decay_config_connect

\eatline
\nwendcode{}\nwbegindocs{1237}Set the flavor entry, which repeats the flavor of the parent unstable
particle.
\nwenddocs{}\nwbegincode{1238}\moddef{Decays: decay config: TBP}\plusendmoddef
  procedure :: set_flv => decay_config_set_flv
\nwendcode{}\nwbegincode{1239}\moddef{Decays: procedures}\plusendmoddef
  subroutine decay_config_set_flv (decay, flv)
    class(decay_config_t), intent(inout) :: decay
    type(flavor_t), intent(in) :: flv
    decay%flv = flv
  end subroutine decay_config_set_flv
  
\eatline
\nwendcode{}\nwbegindocs{1240}Compute embedded branchings and the relative error.  This method does not
apply to the decay root.
\nwenddocs{}\nwbegincode{1241}\moddef{Decays: decay config: TBP}\plusendmoddef
  procedure :: compute => decay_config_compute
\nwendcode{}\nwbegincode{1242}\moddef{Decays: procedures}\plusendmoddef
  recursive subroutine decay_config_compute (decay)
    class(decay_config_t), intent(inout) :: decay
    call decay%decay_root_config_t%compute ()
    if (decay%integral == 0) then
       call decay%write ()
       call msg_fatal ("Decay configuration: partial width is zero")
    end if
    decay%rel_error = decay%abs_error / decay%integral
  end subroutine decay_config_compute
  
\eatline
\nwendcode{}\nwbegindocs{1243}\nwdocspar
\subsection{Decay Instance}
The decay contains a collection of terms.  One of them is selected when the
decay is evaluated.  This is similar to the decay root, but we implement it
independently.

The process instance object is allocated via a pointer, so it automatically
behaves as a target.
\nwenddocs{}\nwbegincode{1244}\moddef{Decays: types}\plusendmoddef
  type, extends (decay_gen_t) :: decay_t
     type(decay_config_t), pointer :: config => null ()
     class(rng_t), allocatable :: rng
   contains
   \LA{}Decays: decay: TBP\RA{}
  end type decay_t
     
\eatline
\nwendcode{}\nwbegindocs{1245}The finalizer is recursive.
\nwenddocs{}\nwbegincode{1246}\moddef{Decays: decay: TBP}\endmoddef
  procedure :: final => decay_final
\nwendcode{}\nwbegincode{1247}\moddef{Decays: procedures}\plusendmoddef
  recursive subroutine decay_final (object)
    class(decay_t), intent(inout) :: object
    integer :: i
    call object%base_final ()
    do i = 1, object%config%process%get_n_mci ()
       call object%process_instance%final_simulation (i)
    end do
    call object%process_instance%final ()
    deallocate (object%process_instance)
  end subroutine decay_final
  
\eatline
\nwendcode{}\nwbegindocs{1248}Output.
\nwenddocs{}\nwbegincode{1249}\moddef{Decays: decay: TBP}\plusendmoddef
  procedure :: write => decay_write
\nwendcode{}\nwbegincode{1250}\moddef{Decays: procedures}\plusendmoddef
  recursive subroutine decay_write (object, unit, indent, recursive)
    class(decay_t), intent(in) :: object
    integer, intent(in), optional :: unit, indent, recursive
    integer :: u, ind
    u = given_output_unit (unit)
    ind = 0;  if (present (indent))  ind = indent
    call object%config%write (unit, indent, verbose = .false.)
    if (allocated (object%rng)) then
       call object%rng%write (u, ind + 1)
    end if
    call write_indent (u, ind)
    if (object%selected_mci > 0) then
       write (u, "(3x,A,I0)")  "Selected MCI    = ", object%selected_mci
    else
       write (u, "(3x,A)")  "Selected MCI    = [undefined]"
    end if
    call write_indent (u, ind)
    if (object%selected_term > 0) then
       write (u, "(3x,A,I0)")  "Selected term   = ", object%selected_term
       call object%term(object%selected_term)%write (u, ind + 1)
    else
       write (u, "(3x,A)")  "Selected term   = [undefined]"
    end if
  end subroutine decay_write

\eatline
\nwendcode{}\nwbegindocs{1251}Initializer.  Base initialization is done recursively.  Then, we
prepare the current process instance and allocate a random-number
generator for term selection.  For all unstable particles, we also
allocate a r.n.g. as spawned by the current process.
\nwenddocs{}\nwbegincode{1252}\moddef{Decays: decay: TBP}\plusendmoddef
  procedure :: init => decay_init
\nwendcode{}\nwbegincode{1253}\moddef{Decays: procedures}\plusendmoddef
  recursive subroutine decay_init (decay, config)
    class(decay_t), intent(out) :: decay
    type(decay_config_t), intent(in), target :: config
    integer :: i
    call decay%base_init (config%term_config)
    decay%config => config
    allocate (decay%process_instance)
    call decay%process_instance%init (decay%config%process)
    call decay%process_instance%setup_event_data (decay%config%model)
    do i = 1, decay%config%process%get_n_mci ()
       call decay%process_instance%init_simulation (i)
    end do
    call decay%config%process%make_rng (decay%rng)
    call decay%make_term_rng (decay%config%process)
  end subroutine decay_init

\eatline
\nwendcode{}\nwbegindocs{1254}Link interactions to the parent process.  {\tt{}i{\char95}prt} is the index of
the current outgoing particle in the parent interaction, for which we
take the trace evaluator.  We link it to the beam particle in the beam
interaction of the decay process instance.  Then, repeat the procedure
for the outgoing particles.
\nwenddocs{}\nwbegincode{1255}\moddef{Decays: decay: TBP}\plusendmoddef
  procedure :: link_interactions => decay_link_interactions
\nwendcode{}\nwbegincode{1256}\moddef{Decays: procedures}\plusendmoddef
  recursive subroutine decay_link_interactions (decay, i_prt, trace)
    class(decay_t), intent(inout) :: decay
    integer, intent(in) :: i_prt
    type(interaction_t), intent(in), target :: trace
    type(interaction_t), pointer :: beam_int
    integer :: n_in, n_vir
    beam_int => decay%process_instance%get_beam_int_ptr ()
    n_in = interaction_get_n_in (trace)
    n_vir = interaction_get_n_vir (trace)
    call interaction_set_source_link (beam_int, 1, trace, &
         n_in + n_vir + i_prt)
    call decay%link_term_interactions ()
  end subroutine decay_link_interactions
    
\eatline
\nwendcode{}\nwbegindocs{1257}Determine a decay chain.  For each unstable particle we select one
of the possible decay modes, and for each decay process we select one
of the possible decay MCI components, calling the random-number
generators.  We do not generate momenta, yet.
\nwenddocs{}\nwbegincode{1258}\moddef{Decays: decay: TBP}\plusendmoddef
  procedure :: select_chain => decay_select_chain
\nwendcode{}\nwbegincode{1259}\moddef{Decays: procedures}\plusendmoddef
  recursive subroutine decay_select_chain (decay)
    class(decay_t), intent(inout) :: decay
    real(default) :: x
    integer :: i
    call decay%rng%generate (x)
    decay%selected_mci = decay%config%mci_selector%select (x)
    call decay%process_instance%choose_mci (decay%selected_mci)
    call decay%process_instance%select_i_term (decay%selected_term)
    do i = 1, size (decay%term)
       call decay%term(i)%select_chain ()
    end do
  end subroutine decay_select_chain
  
\eatline
\nwendcode{}\nwbegindocs{1260}Generate a decay.  We first receive the beam momenta from the parent
process (assuming that this is properly linked), then call the
associated process object for a new event.

Factor out the trace of the helicity density matrix of the isolated
state (the one that will be used for the decay chain).  The trace is
taken into account for unweighting the individual decay event and
should therefore be ignored for unweighting the correlated decay
chain afterwards.
\nwenddocs{}\nwbegincode{1261}\moddef{Decays: decay: TBP}\plusendmoddef
  procedure :: generate => decay_generate
\nwendcode{}\nwbegincode{1262}\moddef{Decays: procedures}\plusendmoddef
  recursive subroutine decay_generate (decay)
    class(decay_t), intent(inout) :: decay
    type(isolated_state_t), pointer :: isolated_state
    integer :: i
    call decay%process_instance%receive_beam_momenta ()
    call decay%config%process%generate_unweighted_event &
         (decay%process_instance, decay%selected_mci)
    if (signal_is_pending ())  return
    call decay%process_instance%evaluate_event_data ()
    isolated_state => &
         decay%process_instance%get_isolated_state_ptr (decay%selected_term)
    call isolated_state%normalize_matrix_by_trace ()
    do i = 1, size (decay%term)
       call decay%term(i)%generate ()
       if (signal_is_pending ())  return
    end do
  end subroutine decay_generate
    
\eatline
\nwendcode{}\nwbegindocs{1263}\nwdocspar
\subsection{Stable Particles}
This is a stable particle.  The flavor can be ambiguous (e.g., partons).
\nwenddocs{}\nwbegincode{1264}\moddef{Decays: types}\plusendmoddef
  type, extends (any_config_t) :: stable_config_t
     type(flavor_t), dimension(:), allocatable :: flv
   contains
   \LA{}Decays: stable config: TBP\RA{}
  end type stable_config_t

\eatline
\nwendcode{}\nwbegindocs{1265}The finalizer is empty:
\nwenddocs{}\nwbegincode{1266}\moddef{Decays: stable config: TBP}\endmoddef
  procedure :: final => stable_config_final
\nwendcode{}\nwbegincode{1267}\moddef{Decays: procedures}\plusendmoddef
  subroutine stable_config_final (object)
    class(stable_config_t), intent(inout) :: object
  end subroutine stable_config_final
  
\eatline
\nwendcode{}\nwbegindocs{1268}Output.
\nwenddocs{}\nwbegincode{1269}\moddef{Decays: stable config: TBP}\plusendmoddef
  procedure :: write => stable_config_write
\nwendcode{}\nwbegincode{1270}\moddef{Decays: procedures}\plusendmoddef
  recursive subroutine stable_config_write (object, unit, indent, verbose)
    class(stable_config_t), intent(in) :: object
    integer, intent(in), optional :: unit, indent
    logical, intent(in), optional :: verbose
    integer :: u, i, ind
    u = given_output_unit (unit)
    ind = 0;  if (present (indent))  ind = indent
    call write_indent (u, ind)
    write (u, "(1x,'+',1x,A)", advance = "no")  "Stable:"
    write (u, "(1x,A)", advance = "no")  char (object%flv(1)%get_name ())
    do i = 2, size (object%flv)
       write (u, "(':',A)", advance = "no") &
            char (object%flv(i)%get_name ())
    end do
    write (u, *)
  end subroutine stable_config_write
  
\eatline
\nwendcode{}\nwbegindocs{1271}Initializer.  We are presented with an array of flavors, but there may be
double entries which we remove, so we store only the distinct flavors.
\nwenddocs{}\nwbegincode{1272}\moddef{Decays: stable config: TBP}\plusendmoddef
  procedure :: init => stable_config_init
\nwendcode{}\nwbegincode{1273}\moddef{Decays: procedures}\plusendmoddef
  subroutine stable_config_init (config, flv)
    class(stable_config_t), intent(out) :: config
    type(flavor_t), dimension(:), intent(in) :: flv
    integer, dimension (size (flv)) :: pdg
    logical, dimension (size (flv)) :: mask
    integer :: i
    pdg = flv%get_pdg ()
    mask(1) = .true.
    forall (i = 2 : size (pdg))
       mask(i) = all (pdg(i) /= pdg(1:i-1))
    end forall
    allocate (config%flv (count (mask)))
    config%flv = pack (flv, mask)
  end subroutine stable_config_init
  
\eatline
\nwendcode{}\nwbegindocs{1274}Here is the corresponding object instance.  Except for the pointer
to the configuration, there is no content.
\nwenddocs{}\nwbegincode{1275}\moddef{Decays: types}\plusendmoddef
  type, extends (any_t) :: stable_t
     type(stable_config_t), pointer :: config => null ()
   contains
   \LA{}Decays: stable: TBP\RA{}
  end type stable_t

\eatline
\nwendcode{}\nwbegindocs{1276}The finalizer does nothing.
\nwenddocs{}\nwbegincode{1277}\moddef{Decays: stable: TBP}\endmoddef
  procedure :: final => stable_final
\nwendcode{}\nwbegincode{1278}\moddef{Decays: procedures}\plusendmoddef
  subroutine stable_final (object)
    class(stable_t), intent(inout) :: object
  end subroutine stable_final

\eatline
\nwendcode{}\nwbegindocs{1279}We can delegate output to the configuration object.
\nwenddocs{}\nwbegincode{1280}\moddef{Decays: stable: TBP}\plusendmoddef
  procedure :: write => stable_write
\nwendcode{}\nwbegincode{1281}\moddef{Decays: procedures}\plusendmoddef
  subroutine stable_write (object, unit, indent)
    class(stable_t), intent(in) :: object
    integer, intent(in), optional :: unit, indent
    call object%config%write (unit, indent)
  end subroutine stable_write
  
\eatline
\nwendcode{}\nwbegindocs{1282}Initializer: just assign the configuration.
\nwenddocs{}\nwbegincode{1283}\moddef{Decays: stable: TBP}\plusendmoddef
  procedure :: init => stable_init
\nwendcode{}\nwbegincode{1284}\moddef{Decays: procedures}\plusendmoddef
  subroutine stable_init (stable, config)
    class(stable_t), intent(out) :: stable
    type(stable_config_t), intent(in), target :: config
    stable%config => config
  end subroutine stable_init
  
\eatline
\nwendcode{}\nwbegindocs{1285}\nwdocspar
\subsection{Unstable Particles}
A branching configuration enables us to select among distinct decay
modes of a particle.  We store the particle flavor (with its implicit
link to a model), an array of decay configurations, and a selector object.

The total width, absolute and relative error are stored as
{\tt{}integral}, {\tt{}abs{\char95}error}, and {\tt{}rel{\char95}error}, respectively.

The flavor must be unique in this case.
\nwenddocs{}\nwbegincode{1286}\moddef{Decays: types}\plusendmoddef
  type, extends (any_config_t) :: unstable_config_t
     type(flavor_t) :: flv
     real(default) :: integral = 0
     real(default) :: abs_error = 0
     real(default) :: rel_error = 0
     type(selector_t) :: selector
     type(decay_config_t), dimension(:), allocatable :: decay_config
   contains
   \LA{}Decays: unstable config: TBP\RA{}
  end type unstable_config_t
  
\eatline
\nwendcode{}\nwbegindocs{1287}Finalizer.  The branching configuration can be a recursive structure.
\nwenddocs{}\nwbegincode{1288}\moddef{Decays: unstable config: TBP}\endmoddef
  procedure :: final => unstable_config_final
\nwendcode{}\nwbegincode{1289}\moddef{Decays: procedures}\plusendmoddef
  recursive subroutine unstable_config_final (object)
    class(unstable_config_t), intent(inout) :: object
    integer :: i
    if (allocated (object%decay_config)) then
       do i = 1, size (object%decay_config)
          call object%decay_config(i)%final ()
       end do
    end if
  end subroutine unstable_config_final

\eatline
\nwendcode{}\nwbegindocs{1290}Output.  Since this may be recursive, we include indentation.
\nwenddocs{}\nwbegincode{1291}\moddef{Decays: unstable config: TBP}\plusendmoddef
  procedure :: write => unstable_config_write
\nwendcode{}\nwbegincode{1292}\moddef{Decays: procedures}\plusendmoddef
  recursive subroutine unstable_config_write (object, unit, indent, verbose)
    class(unstable_config_t), intent(in) :: object
    integer, intent(in), optional :: unit, indent
    logical, intent(in), optional :: verbose
    integer :: u, i, ind
    logical :: verb
    u = given_output_unit (unit)
    ind = 0;  if (present (indent))  ind = indent
    verb = .true.;  if (present (verbose))  verb = verbose
    call write_indent (u, ind)
    write (u, "(1x,'+',1x,A,1x,A)")  "Unstable:", &
         char (object%flv%get_name ())
    call write_indent (u, ind)
    write (u, 1)  "total width =", object%integral
    call write_indent (u, ind)
    write (u, 1)  "error (abs) =", object%abs_error
    call write_indent (u, ind)
    write (u, 1)  "error (rel) =", object%rel_error
1   format (5x,A,ES19.12)
    if (verb .and. allocated (object%decay_config)) then
       do i = 1, size (object%decay_config)
          call object%decay_config(i)%write (u, ind + 1)
       end do
    end if
  end subroutine unstable_config_write
  
\eatline
\nwendcode{}\nwbegindocs{1293}Initializer.  For the unstable particle, the flavor is unique.
\nwenddocs{}\nwbegincode{1294}\moddef{Decays: unstable config: TBP}\plusendmoddef
  procedure :: init => unstable_config_init
\nwendcode{}\nwbegincode{1295}\moddef{Decays: procedures}\plusendmoddef
  subroutine unstable_config_init (config, flv)
    class(unstable_config_t), intent(out) :: config
    type(flavor_t), intent(in) :: flv
    config%flv = flv
  end subroutine unstable_config_init
  
\eatline
\nwendcode{}\nwbegindocs{1296}Further initialization: determine the number of decay modes.  We can assume
that the flavor of the particle has been set already.

If the process stack is given, we can delve recursively into actually
assigning decay processes.  Otherwise, we just initialize with decay
process names.
\nwenddocs{}\nwbegincode{1297}\moddef{Decays: unstable config: TBP}\plusendmoddef
  procedure :: init_decays => unstable_config_init_decays
\nwendcode{}\nwbegincode{1298}\moddef{Decays: procedures}\plusendmoddef
  recursive subroutine unstable_config_init_decays &
       (unstable, decay_id, model, process_stack)
    class(unstable_config_t), intent(inout) :: unstable
    type(string_t), dimension(:), intent(in) :: decay_id
    class(model_data_t), intent(in), target :: model
    type(process_stack_t), intent(in), optional :: process_stack
    integer :: i
    allocate (unstable%decay_config (size (decay_id)))
    do i = 1, size (decay_id)
       associate (decay => unstable%decay_config(i))
         if (present (process_stack)) then
            call decay%connect (process_stack%get_process_ptr (decay_id(i)), &
                 model, process_stack)
         else
            call decay%init (model, decay_id(i))
         end if
         call decay%set_flv (unstable%flv)
       end associate
    end do
  end subroutine unstable_config_init_decays
  
\eatline
\nwendcode{}\nwbegindocs{1299}Compute the total width and branching ratios, initializing the decay
selector.
\nwenddocs{}\nwbegincode{1300}\moddef{Decays: unstable config: TBP}\plusendmoddef
  procedure :: compute => unstable_config_compute
\nwendcode{}\nwbegincode{1301}\moddef{Decays: procedures}\plusendmoddef
  recursive subroutine unstable_config_compute (unstable)
    class(unstable_config_t), intent(inout) :: unstable
    integer :: i
    do i = 1, size (unstable%decay_config)
       call unstable%decay_config(i)%compute ()
    end do
    unstable%integral = sum (unstable%decay_config%integral)
    if (unstable%integral <= 0) then
       call unstable%write ()
       call msg_fatal ("Decay configuration: computed total width is zero")
    end if
    unstable%abs_error = sqrt (sum (unstable%decay_config%abs_error ** 2))
    unstable%rel_error = unstable%abs_error / unstable%integral
    call unstable%selector%init (unstable%decay_config%integral)
    do i = 1, size (unstable%decay_config)
       unstable%decay_config(i)%weight &
            = unstable%selector%get_weight (i)
    end do
  end subroutine unstable_config_compute
    
\eatline
\nwendcode{}\nwbegindocs{1302}\nwdocspar
Now we define the instance of an unstable particle.
\nwenddocs{}\nwbegincode{1303}\moddef{Decays: types}\plusendmoddef
  type, extends (any_t) :: unstable_t
     type(unstable_config_t), pointer :: config => null ()
     class(rng_t), allocatable :: rng
     integer :: selected_decay = 0
     type(decay_t), dimension(:), allocatable :: decay
   contains
   \LA{}Decays: unstable: TBP\RA{}
  end type unstable_t
  
\eatline
\nwendcode{}\nwbegindocs{1304}Recursive finalizer.
\nwenddocs{}\nwbegincode{1305}\moddef{Decays: unstable: TBP}\endmoddef
  procedure :: final => unstable_final
\nwendcode{}\nwbegincode{1306}\moddef{Decays: procedures}\plusendmoddef
  recursive subroutine unstable_final (object)
    class(unstable_t), intent(inout) :: object
    integer :: i
    if (allocated (object%decay)) then
       do i = 1, size (object%decay)
          call object%decay(i)%final ()
       end do
    end if
  end subroutine unstable_final
  
\eatline
\nwendcode{}\nwbegindocs{1307}Output.
\nwenddocs{}\nwbegincode{1308}\moddef{Decays: unstable: TBP}\plusendmoddef
  procedure :: write => unstable_write
\nwendcode{}\nwbegincode{1309}\moddef{Decays: procedures}\plusendmoddef
  recursive subroutine unstable_write (object, unit, indent)
    class(unstable_t), intent(in) :: object
    integer, intent(in), optional :: unit, indent
    integer :: u, ind
    u = given_output_unit (unit)
    ind = 0;  if (present (indent))  ind = indent
    call object%config%write (u, ind, verbose=.false.)
    if (allocated (object%rng)) then
       call object%rng%write (u, ind + 2)
    end if
    call write_indent (u, ind)
    if (object%selected_decay > 0) then
       write (u, "(5x,A,I0)") "Sel. decay  = ", object%selected_decay
       call object%decay(object%selected_decay)%write (u, ind + 1)
    else
       write (u, "(5x,A)")  "Sel. decay  = [undefined]"
    end if
  end subroutine unstable_write
    
\eatline
\nwendcode{}\nwbegindocs{1310}Write the embedded process instances.
\nwenddocs{}\nwbegincode{1311}\moddef{Decays: unstable: TBP}\plusendmoddef
  procedure :: write_process_instances => unstable_write_process_instances
\nwendcode{}\nwbegincode{1312}\moddef{Decays: procedures}\plusendmoddef
  recursive subroutine unstable_write_process_instances &
       (unstable, unit, verbose)
    class(unstable_t), intent(in) :: unstable
    integer, intent(in), optional :: unit
    logical, intent(in), optional :: verbose
    if (unstable%selected_decay > 0) then
       call unstable%decay(unstable%selected_decay)% &
            write_process_instances (unit, verbose)
    end if
  end subroutine unstable_write_process_instances
  
\eatline
\nwendcode{}\nwbegindocs{1313}Initialization, using the configuration object.
\nwenddocs{}\nwbegincode{1314}\moddef{Decays: unstable: TBP}\plusendmoddef
  procedure :: init => unstable_init
\nwendcode{}\nwbegincode{1315}\moddef{Decays: procedures}\plusendmoddef
  recursive subroutine unstable_init (unstable, config)
    class(unstable_t), intent(out) :: unstable
    type(unstable_config_t), intent(in), target :: config
    integer :: i
    unstable%config => config
    allocate (unstable%decay (size (config%decay_config)))
    do i = 1, size (config%decay_config)
       call unstable%decay(i)%init (config%decay_config(i))
    end do
  end subroutine unstable_init
  
\eatline
\nwendcode{}\nwbegindocs{1316}Recursively link interactions to the parent process.  {\tt{}i{\char95}prt} is
the index of the current outgoing particle in the parent interaction.
\nwenddocs{}\nwbegincode{1317}\moddef{Decays: unstable: TBP}\plusendmoddef
  procedure :: link_interactions => unstable_link_interactions
\nwendcode{}\nwbegincode{1318}\moddef{Decays: procedures}\plusendmoddef
  recursive subroutine unstable_link_interactions (unstable, i_prt, trace)
    class(unstable_t), intent(inout) :: unstable
    integer, intent(in) :: i_prt
    type(interaction_t), intent(in), target :: trace
    integer :: i
    do i = 1, size (unstable%decay)
       call unstable%decay(i)%link_interactions (i_prt, trace)
    end do
  end subroutine unstable_link_interactions
    
\eatline
\nwendcode{}\nwbegindocs{1319}Import the random-number generator state.
\nwenddocs{}\nwbegincode{1320}\moddef{Decays: unstable: TBP}\plusendmoddef
  procedure :: import_rng => unstable_import_rng
\nwendcode{}\nwbegincode{1321}\moddef{Decays: procedures}\plusendmoddef
  subroutine unstable_import_rng (unstable, rng)
    class(unstable_t), intent(inout) :: unstable
    class(rng_t), intent(inout), allocatable :: rng
    call move_alloc (from = rng, to = unstable%rng)
  end subroutine unstable_import_rng
  
\eatline
\nwendcode{}\nwbegindocs{1322}Generate a decay chain.  First select a decay mode, then call the
{\tt{}select{\char95}chain} method of the selected mode.
\nwenddocs{}\nwbegincode{1323}\moddef{Decays: unstable: TBP}\plusendmoddef
  procedure :: select_chain => unstable_select_chain
\nwendcode{}\nwbegincode{1324}\moddef{Decays: procedures}\plusendmoddef
  recursive subroutine unstable_select_chain (unstable)
    class(unstable_t), intent(inout) :: unstable
    real(default) :: x
    call unstable%rng%generate (x)
    unstable%selected_decay = unstable%config%selector%select (x)
    call unstable%decay(unstable%selected_decay)%select_chain ()
  end subroutine unstable_select_chain
    
\eatline
\nwendcode{}\nwbegindocs{1325}Generate a decay event.
\nwenddocs{}\nwbegincode{1326}\moddef{Decays: unstable: TBP}\plusendmoddef
  procedure :: generate => unstable_generate
\nwendcode{}\nwbegincode{1327}\moddef{Decays: procedures}\plusendmoddef
  recursive subroutine unstable_generate (unstable)
    class(unstable_t), intent(inout) :: unstable
    call unstable%decay(unstable%selected_decay)%generate ()
  end subroutine unstable_generate
    
\eatline
\nwendcode{}\nwbegindocs{1328}\nwdocspar
\subsection{Decay Chain}
While the decay configuration tree and the decay tree are static
entities (during a simulation run), the decay chain is dynamically
generated for each event.  The reason is that with the possibility of
several decay modes for each particle, and several terms for each
process, the total number of distinct decay chains is not under control.

Each entry in the decay chain is a connected parton state.  The origin
of the chain is a connected state in the parent process (not part of
the chain itself).  For each decay, mode and term chosen, we convolute
this with the isolated (!) state of the current decay, to generate a
new connected state.  We accumulate this chain by recursively
traversing the allocated decay tree.  Whenever a particle decays, it
becomes virtual and is replaced by its decay product, while all other
particles stay in the parton state as spectators.

Technically, we implement the decay chain as a stack structure and
include information from the associated decay object for easier
debugging.  This is a decay chain entry:
\nwenddocs{}\nwbegincode{1329}\moddef{Decays: types}\plusendmoddef
  type, extends (connected_state_t) :: decay_chain_entry_t
     integer :: index = 0
     type(decay_config_t), pointer :: config => null ()
     integer :: selected_mci = 0
     integer :: selected_term = 0
     type(decay_chain_entry_t), pointer :: previous => null ()
  end type decay_chain_entry_t
     
\eatline
\nwendcode{}\nwbegindocs{1330}This is the complete chain; we need just a pointer to the last
entry.  We also include a pointer to the master process instance,
which serves as the seed for the decay chain.  

The evaluator {\tt{}correlated{\char95}trace} traces over all quantum numbers
for the final spin-correlated (but color-summed) evaluator of the
decay chain.  This allows us to compute the probability for a momentum
configuration, given that all individual density matrices (of the
initial process and the subsequent decays) have been normalized to one.

Note: This trace is summed over color, so color is treated exactly
when computing spin correlations.  However, we do not keep
non-diagonal color correlations.  When an event is accepted, we
compute probabilities for all color states and can choose one of them.
\nwenddocs{}\nwbegincode{1331}\moddef{Decays: types}\plusendmoddef
  type :: decay_chain_t
     type(process_instance_t), pointer :: process_instance => null ()
     integer :: selected_term = 0
     type(evaluator_t) :: correlated_trace
     type(decay_chain_entry_t), pointer :: last => null ()
   contains
   \LA{}Decays: decay chain: TBP\RA{}
  end type decay_chain_t
  
\eatline
\nwendcode{}\nwbegindocs{1332}The finalizer recursively deletes and deallocates the entries.
\nwenddocs{}\nwbegincode{1333}\moddef{Decays: decay chain: TBP}\endmoddef
  procedure :: final => decay_chain_final
\nwendcode{}\nwbegincode{1334}\moddef{Decays: procedures}\plusendmoddef
  subroutine decay_chain_final (object)
    class(decay_chain_t), intent(inout) :: object
    type(decay_chain_entry_t), pointer :: entry
    do while (associated (object%last))
       entry => object%last
       object%last => entry%previous
       call entry%final ()
       deallocate (entry)
    end do
    call evaluator_final (object%correlated_trace)
  end subroutine decay_chain_final
  
\eatline
\nwendcode{}\nwbegindocs{1335}Doing output recursively allows us to display the chain in
chronological order.
\nwenddocs{}\nwbegincode{1336}\moddef{Decays: decay chain: TBP}\plusendmoddef
  procedure :: write => decay_chain_write
\nwendcode{}\nwbegincode{1337}\moddef{Decays: procedures}\plusendmoddef
  subroutine decay_chain_write (object, unit)
    class(decay_chain_t), intent(in) :: object
    integer, intent(in), optional :: unit
    integer :: u
    u = given_output_unit (unit)
    call write_separator (u, 2)
    write (u, "(1x,A)")  "Decay chain:"
    call write_entries (object%last)
    call write_separator (u, 2)
    write (u, "(1x,A)")  "Evaluator (correlated trace of the decay chain):"
    call write_separator (u)
    call object%correlated_trace%write (u)
    call write_separator (u, 2)
  contains
    recursive subroutine write_entries (entry)
      type(decay_chain_entry_t), intent(in), pointer :: entry
      if (associated (entry)) then
         call write_entries (entry%previous)
         call write_separator (u, 2)
         write (u, "(1x,A,I0)")  "Decay #", entry%index
         call entry%config%write_header (u)
         write (u, "(3x,A,I0)")  "Selected MCI    = ", entry%selected_mci
         write (u, "(3x,A,I0)")  "Selected term   = ", entry%selected_term
         call entry%config%term_config(entry%selected_term)%write (u, indent=1)
         call entry%write (u)
      end if
    end subroutine write_entries
  end subroutine decay_chain_write
    
\eatline
\nwendcode{}\nwbegindocs{1338}\nwdocspar
Build a decay chain, recursively following the selected decays and
terms in a decay tree.  Before start, we finalize the chain, deleting
any previous contents.
\nwenddocs{}\nwbegincode{1339}\moddef{Decays: decay chain: TBP}\plusendmoddef
  procedure :: build => decay_chain_build
\nwendcode{}\nwbegincode{1340}\moddef{Decays: procedures}\plusendmoddef
  subroutine decay_chain_build (chain, decay_root)
    class(decay_chain_t), intent(inout), target :: chain
    type(decay_root_t), intent(in) :: decay_root
    type(quantum_numbers_mask_t), dimension(:), allocatable :: qn_mask
    type(interaction_t), pointer :: int_last_decay
    call chain%final ()
    if (decay_root%selected_term > 0) then
       chain%process_instance => decay_root%process_instance
       chain%selected_term = decay_root%selected_term
       call chain%build_term_entries (decay_root%term(decay_root%selected_term))
    end if
    int_last_decay => chain%last%get_matrix_int_ptr ()
    allocate (qn_mask (interaction_get_n_tot (int_last_decay)))
    call qn_mask%init (mask_f = .true., mask_c = .true., mask_h = .true.)
    call evaluator_init_qn_sum (chain%correlated_trace, int_last_decay, qn_mask)
  end subroutine decay_chain_build
    
\eatline
\nwendcode{}\nwbegindocs{1341}Build the entries that correspond to a decay term.  We have to scan
all unstable particles.
\nwenddocs{}\nwbegincode{1342}\moddef{Decays: decay chain: TBP}\plusendmoddef
  procedure :: build_term_entries => decay_chain_build_term_entries
\nwendcode{}\nwbegincode{1343}\moddef{Decays: procedures}\plusendmoddef
  recursive subroutine decay_chain_build_term_entries (chain, term)
    class(decay_chain_t), intent(inout) :: chain
    type(decay_term_t), intent(in) :: term
    integer :: i
    do i = 1, size (term%particle_out)
       select type (unstable => term%particle_out(i)%c)
       type is (unstable_t)
          if (unstable%selected_decay > 0) then
             call chain%build_decay_entries &
                  (unstable%decay(unstable%selected_decay))
          end if
       end select
    end do
  end subroutine decay_chain_build_term_entries
  
\eatline
\nwendcode{}\nwbegindocs{1344}Build the entries that correspond to a specific decay.  The
decay term should have been determined, so we allocate a decay chain
entry and fill it, then proceed to child decays.

For the first entry, we convolute the connected state of the parent process
instance with the isolated state of the current
decay (which does not contain an extra beam entry for the parent).
For subsequent entries, we take the previous entry as first factor.

In principle, each chain entry (as a parton state) is capable of
holding a subevent object and associated expressions.  We currently do
not make use of that feature.

Before generating the decays, factor out the trace of the helicity
density matrix of the parent parton state.  This trace has
been used for unweighting the original event (unweighted case) or it
determines the overall weight, so it should not be taken into account
in the decay chain generation.
\nwenddocs{}\nwbegincode{1345}\moddef{Decays: decay chain: TBP}\plusendmoddef
  procedure :: build_decay_entries => decay_chain_build_decay_entries
\nwendcode{}\nwbegincode{1346}\moddef{Decays: procedures}\plusendmoddef
  recursive subroutine decay_chain_build_decay_entries (chain, decay)
    class(decay_chain_t), intent(inout) :: chain
    type(decay_t), intent(in) :: decay
    type(decay_chain_entry_t), pointer :: entry
    type(connected_state_t), pointer :: previous_state
    type(isolated_state_t), pointer :: current_decay
    type(helicity_t) :: hel
    type(quantum_numbers_t) :: qn_filter_conn
    allocate (entry)
    if (associated (chain%last)) then
       entry%previous => chain%last
       entry%index = entry%previous%index + 1
       previous_state => entry%previous%connected_state_t
    else
       entry%index = 1
       previous_state => &
            chain%process_instance%get_connected_state_ptr (chain%selected_term)
    end if
    entry%config => decay%config
    entry%selected_mci = decay%selected_mci
    entry%selected_term = decay%selected_term
    current_decay => decay%process_instance%get_isolated_state_ptr &
         (decay%selected_term)
    call entry%setup_connected_trace &
         (current_decay, previous_state%get_trace_int_ptr (), resonant=.true.)
    if (entry%config%flv%has_decay_helicity ()) then
       call hel%init (entry%config%flv%get_decay_helicity ())
       call qn_filter_conn%init (hel)
       call entry%setup_connected_matrix &
            (current_decay, previous_state%get_matrix_int_ptr (), &
            resonant=.true., qn_filter_conn = qn_filter_conn)
       call entry%setup_connected_flows &
            (current_decay, previous_state%get_flows_int_ptr (), &
            resonant=.true., qn_filter_conn = qn_filter_conn)
    else
       call entry%setup_connected_matrix &
            (current_decay, previous_state%get_matrix_int_ptr (), &
            resonant=.true.)
       call entry%setup_connected_flows &
            (current_decay, previous_state%get_flows_int_ptr (), &
            resonant=.true.)
    end if
    chain%last => entry
    call chain%build_term_entries (decay%term(decay%selected_term))
  end subroutine decay_chain_build_decay_entries
    
\eatline
\nwendcode{}\nwbegindocs{1347}Recursively fill the decay chain with momenta and evaluate the
matrix elements.  Since all evaluators should have correct source
entries at this point, momenta are automatically retrieved from the
appropriate process instance.


Like we did above for the parent process, factor out the trace for
each subsequent decay (the helicity density matrix in the isolated
state, which is taken for the convolution).
\nwenddocs{}\nwbegincode{1348}\moddef{Decays: decay chain: TBP}\plusendmoddef
  procedure :: evaluate => decay_chain_evaluate
\nwendcode{}\nwbegincode{1349}\moddef{Decays: procedures}\plusendmoddef
  subroutine decay_chain_evaluate (chain)
    class(decay_chain_t), intent(inout) :: chain
    call evaluate (chain%last)
    call evaluator_receive_momenta (chain%correlated_trace)
    call chain%correlated_trace%evaluate ()
  contains
    recursive subroutine evaluate (entry)
      type(decay_chain_entry_t), intent(inout), pointer :: entry
      if (associated (entry)) then
         call evaluate (entry%previous)
         call entry%receive_kinematics ()
         call entry%evaluate_trace ()
         call entry%evaluate_event_data ()
      end if
    end subroutine evaluate
  end subroutine decay_chain_evaluate
  
\eatline
\nwendcode{}\nwbegindocs{1350}Return the probability of a decay chain.  This is given as the trace
of the density matrix with intermediate helicity correlations,
normalized by the product of the uncorrelated density matrix traces.  This
works only if an event has been evaluated and the {\tt{}correlated{\char95}trace}
evaluator is filled.  By definition, this evaluator has only one
matrix element, and this must be real.
\nwenddocs{}\nwbegincode{1351}\moddef{Decays: decay chain: TBP}\plusendmoddef
  procedure :: get_probability => decay_chain_get_probability
\nwendcode{}\nwbegincode{1352}\moddef{Decays: procedures}\plusendmoddef
  function decay_chain_get_probability (chain) result (x)
    class(decay_chain_t), intent(in) :: chain
    real(default) :: x
    x = evaluator_get_matrix_element (chain%correlated_trace, 1)
  end function decay_chain_get_probability
  
\eatline
\nwendcode{}\nwbegindocs{1353}\nwdocspar
\subsection{Decay as Event Transform}
The {\tt{}evt{\char95}decay} object combines decay configuration, decay tree, and
chain in a single object, as an implementation of the {\tt{}evt} (event
transform) abstract type.
\nwenddocs{}\nwbegincode{1354}\moddef{Decays: public}\endmoddef
  public :: evt_decay_t
\nwendcode{}\nwbegincode{1355}\moddef{Decays: types}\plusendmoddef
  type, extends (evt_t) :: evt_decay_t
     type(decay_root_config_t) :: decay_root_config
     type(decay_root_t) :: decay_root
     type(decay_chain_t) :: decay_chain
   contains
   \LA{}Decays: evt decay: TBP\RA{}
  end type evt_decay_t
  
\eatline
\nwendcode{}\nwbegindocs{1356}Output.  We display the currently selected decay tree, which
includes configuration data, and the decay chain, i.e., the evaluators.
\nwenddocs{}\nwbegincode{1357}\moddef{Decays: evt decay: TBP}\endmoddef
  procedure :: write => evt_decay_write
\nwendcode{}\nwbegincode{1358}\moddef{Decays: procedures}\plusendmoddef
  subroutine evt_decay_write (object, unit, &
       show_decay_tree, &
       show_processes, &
       verbose)
    class(evt_decay_t), intent(in) :: object
    integer, intent(in), optional :: unit
    logical, intent(in), optional :: show_decay_tree, show_processes, verbose
    logical :: dec, prc, verb
    integer :: u
    u = given_output_unit (unit)
    dec = .true.;  if (present (show_decay_tree))  dec = show_decay_tree
    prc = .false.;  if (present (show_processes))  prc = show_processes
    verb = .false.;  if (present (verbose))  verb = verbose
    call write_separator (u, 2)
    write (u, "(1x,A)")  "Event transform: partonic decays"
    call write_separator (u, 2)
    call object%base_write (u)
    if (dec) then
       call write_separator (u)
       call object%decay_root%write (u)
       if (verb) then
          call object%decay_chain%write (u)
       end if
       if (prc) then
          call object%decay_root%write_process_instances (u, verb)
       end if
    else
       call write_separator (u, 2)
    end if
  end subroutine evt_decay_write
    
\eatline
\nwendcode{}\nwbegindocs{1359}Connect with a process instance and process.  This initializes the
decay configuration.  The process stack is used to look for process
objects that implement daughter decays.

When all processes are assigned, configure the decay tree instance, using the
decay tree configuration.  First obtain the branching ratios, then allocate
the decay tree.  This is done once for all events.
\nwenddocs{}\nwbegincode{1360}\moddef{Decays: evt decay: TBP}\plusendmoddef
  procedure :: connect => evt_decay_connect
\nwendcode{}\nwbegincode{1361}\moddef{Decays: procedures}\plusendmoddef
  subroutine evt_decay_connect (evt, process_instance, model, process_stack)
    class(evt_decay_t), intent(inout), target :: evt
    type(process_instance_t), intent(in), target :: process_instance
    class(model_data_t), intent(in), target :: model
    type(process_stack_t), intent(in), optional :: process_stack
    call evt%base_connect (process_instance, model)
    call evt%decay_root_config%connect (process_instance%process, &
         model, process_stack, process_instance)
    call evt%decay_root_config%compute ()
    call evt%decay_root%init (evt%decay_root_config, evt%process_instance)
  end subroutine evt_decay_connect
  
\eatline
\nwendcode{}\nwbegindocs{1362}Prepare a new event: Select a decay chain and build the corresponding chain
object.
\nwenddocs{}\nwbegincode{1363}\moddef{Decays: evt decay: TBP}\plusendmoddef
  procedure :: prepare_new_event => evt_decay_prepare_new_event
\nwendcode{}\nwbegincode{1364}\moddef{Decays: procedures}\plusendmoddef
  subroutine evt_decay_prepare_new_event (evt, i_mci, i_term)
    class(evt_decay_t), intent(inout) :: evt
    integer, intent(in) :: i_mci, i_term
    call evt%reset ()
    evt%decay_root%selected_mci = i_mci
    evt%decay_root%selected_term = i_term
    call evt%decay_root%select_chain ()
    call evt%decay_chain%build (evt%decay_root)
  end subroutine evt_decay_prepare_new_event
  
\eatline
\nwendcode{}\nwbegindocs{1365}Generate a weighted event and assign the resulting weight
(probability).  We use a chain initialized by the preceding
subroutine, fill it with momenta and evaluate.
\nwenddocs{}\nwbegincode{1366}\moddef{Decays: evt decay: TBP}\plusendmoddef
  procedure :: generate_weighted => evt_decay_generate_weighted
\nwendcode{}\nwbegincode{1367}\moddef{Decays: procedures}\plusendmoddef
  subroutine evt_decay_generate_weighted (evt, probability)
    class(evt_decay_t), intent(inout) :: evt
    real(default), intent(out) :: probability
    call evt%decay_root%generate ()
    if (signal_is_pending ())  return
    call evt%decay_chain%evaluate ()
    probability = evt%decay_chain%get_probability ()
  end subroutine evt_decay_generate_weighted
  
\eatline
\nwendcode{}\nwbegindocs{1368}To create a usable event, we have to transform the interaction into a
particle set; this requires factorization for the correlated density matrix,
according to the factorization mode.
\nwenddocs{}\nwbegincode{1369}\moddef{Decays: evt decay: TBP}\plusendmoddef
  procedure :: make_particle_set => evt_decay_make_particle_set
\nwendcode{}\nwbegincode{1370}\moddef{Decays: procedures}\plusendmoddef
  subroutine evt_decay_make_particle_set &
       (evt, factorization_mode, keep_correlations, r)
    class(evt_decay_t), intent(inout) :: evt
    integer, intent(in) :: factorization_mode
    logical, intent(in) :: keep_correlations
    real(default), dimension(:), intent(in), optional :: r
    type(interaction_t), pointer :: int_matrix, int_flows
    type(decay_chain_entry_t), pointer :: last_entry
    last_entry => evt%decay_chain%last
    int_matrix => last_entry%get_matrix_int_ptr ()
    int_flows  => last_entry%get_flows_int_ptr ()
    call evt%factorize_interactions (int_matrix, int_flows, &
         factorization_mode, keep_correlations, r)
    call evt%tag_incoming ()
  end subroutine evt_decay_make_particle_set
    
\eatline
\nwendcode{}\nwbegindocs{1371}\nwdocspar
Eliminate numerical noise for the associated process instances.
\nwenddocs{}\nwbegincode{1372}\moddef{Decays: public}\plusendmoddef
  public :: pacify
\nwendcode{}\nwbegincode{1373}\moddef{Decays: interfaces}\plusendmoddef
  interface pacify
     module procedure pacify_decay
  end interface pacify
\nwendcode{}\nwbegincode{1374}\moddef{Decays: procedures}\plusendmoddef
  subroutine pacify_decay (evt)
    class(evt_decay_t), intent(inout) :: evt
    call pacify_decay_gen (evt%decay_root)
  contains
    recursive subroutine pacify_decay_gen (decay)
      class(decay_gen_t), intent(inout) :: decay
      if (associated (decay%process_instance)) then
         call pacify (decay%process_instance)
      end if
      if (decay%selected_term > 0) then
         call pacify_term (decay%term(decay%selected_term))
      end if
    end subroutine pacify_decay_gen
    recursive subroutine pacify_term (term)
      class(decay_term_t), intent(inout) :: term
      integer :: i
      do i = 1, size (term%particle_out)
         select type (unstable => term%particle_out(i)%c)
         type is (unstable_t);  call pacify_unstable (unstable)
         end select
      end do
    end subroutine pacify_term
    recursive subroutine pacify_unstable (unstable)
      class(unstable_t), intent(inout) :: unstable
      if (unstable%selected_decay > 0) then
         call pacify_decay_gen (unstable%decay(unstable%selected_decay))
      end if
    end subroutine pacify_unstable
  end subroutine pacify_decay
  
\eatline
\nwendcode{}\nwbegindocs{1375}\nwdocspar
\subsection{Unit tests}
\nwenddocs{}\nwbegincode{1376}\moddef{Decays: public}\plusendmoddef
  public :: decays_test
\nwendcode{}\nwbegincode{1377}\moddef{Decays: tests}\endmoddef
  subroutine decays_test (u, results)
    integer, intent(in) :: u
    type(test_results_t), intent(inout) :: results
  \LA{}Decays: execute tests\RA{}
  end subroutine decays_test
  
\eatline
\nwendcode{}\nwbegindocs{1378}\nwdocspar
\subsubsection{Testbed}
As a variation of the {\tt{}prepare{\char95}test{\char95}process} routine used elsewhere, we
define here a routine that creates two processes (scattering $ss\to ss$ and
decay $s\to f\bar f$), compiles and integrates them and prepares for event
generation. 
\nwenddocs{}\nwbegincode{1379}\moddef{Decays: public}\plusendmoddef
  public :: prepare_testbed
\nwendcode{}\nwbegincode{1380}\moddef{Decays: tests}\plusendmoddef
  subroutine prepare_testbed &
       (lib, process_stack, prefix, os_data, &
        scattering, decay, decay_rest_frame)
    type(process_library_t), intent(out), target :: lib
    type(process_stack_t), intent(out) :: process_stack
    type(string_t), intent(in) :: prefix
    type(os_data_t), intent(in) :: os_data
    logical, intent(in) :: scattering, decay
    logical, intent(in), optional :: decay_rest_frame

    type(model_data_t), target :: model
    class(model_data_t), pointer :: model_copy
    type(string_t) :: libname, procname1, procname2, run_id
    type(qcd_t) :: qcd
    class(rng_factory_t), allocatable :: rng_factory
    type(process_entry_t), pointer :: process
    type(process_instance_t), allocatable, target :: process_instance
    class(prc_core_t), allocatable :: core_template
    class(mci_t), allocatable :: mci_template
    class(phs_config_t), allocatable :: phs_config_template
    type(field_data_t), pointer :: field_data
    real(default) :: sqrts

    libname = prefix // "_lib"
    procname1 = prefix // "_p"
    procname2 = prefix // "_d"
    run_id = prefix

    call model%init_test ()
    call model%set_par (var_str ("ff"), 0.4_default)
    call model%set_par (var_str ("mf"), &
         model%get_real (var_str ("ff")) * model%get_real (var_str ("ms")))

    if (scattering .and. decay) then
       field_data => model%get_field_ptr (25)
       call field_data%set (p_is_stable = .false.)
    end if
    
    call prc_test_create_library (libname, lib, &
         scattering = .true., decay = .true., &
         procname1 = procname1, procname2 = procname2)

    call reset_interaction_counter ()

    allocate (test_t :: core_template)
    allocate (mci_midpoint_t :: mci_template)
    allocate (phs_single_config_t :: phs_config_template)

    if (scattering) then

       allocate (rng_test_factory_t :: rng_factory)
       allocate (model_copy)
       call model_copy%init (model%get_name (), &
            model%get_n_real (), &
            model%get_n_complex (), &
            model%get_n_field (), &
            model%get_n_vtx ())
       call model_copy%copy_from (model)

       allocate (process)
       call process%init (procname1, &
            run_id, lib, os_data, qcd, rng_factory, model_copy)
       call process%init_component &
            (1, core_template, mci_template, phs_config_template)
       sqrts = 1000
       call process%setup_beams_sqrts (sqrts)
       call process%configure_phs ()
       call process%setup_mci ()
       call process%setup_terms ()

       allocate (process_instance)
       call process_instance%init (process%process_t)
       call process%integrate (process_instance, 1, n_it=1, n_calls=100)
       call process%final_integration (1)
       call process_instance%final ()
       deallocate (process_instance)

       call process%prepare_simulation (1)
       call process_stack%push (process)
    end if
    
    if (decay) then
       allocate (rng_test_factory_t :: rng_factory)
       allocate (model_copy)
       call model_copy%init (model%get_name (), &
            model%get_n_real (), &
            model%get_n_complex (), &
            model%get_n_field (), &
            model%get_n_vtx ())
       call model_copy%copy_from (model)

       allocate (process)
       call process%init (procname2, &
            run_id, lib, os_data, qcd, rng_factory, model_copy)
       call process%init_component &
            (1, core_template, mci_template, phs_config_template)
       if (present (decay_rest_frame)) then
          call process%setup_beams_decay (rest_frame = decay_rest_frame)
       else
          call process%setup_beams_decay (rest_frame = .not. scattering)
       end if
       call process%configure_phs ()
       call process%setup_mci ()
       call process%setup_terms ()

       allocate (process_instance)
       call process_instance%init (process%process_t)
       call process%integrate (process_instance, 1, n_it=1, n_calls=100)
       call process%final_integration (1)
       call process_instance%final ()
       deallocate (process_instance)

       call process%prepare_simulation (1)
       call process_stack%push (process)
    end if
    
    call model%final ()

  end subroutine prepare_testbed

\eatline
\nwendcode{}\nwbegindocs{1381}\nwdocspar
\subsubsection{Simple decay configuration}
We define a branching configuration with two decay modes.  We set the
integral values by hand, so we do not need to evaluate processes, yet.
\nwenddocs{}\nwbegincode{1382}\moddef{Decays: execute tests}\endmoddef
  call test (decays_1, "decays_1", &
       "branching and decay configuration", &
       u, results)
\nwendcode{}\nwbegincode{1383}\moddef{Decays: tests}\plusendmoddef
  subroutine decays_1 (u)
    integer, intent(in) :: u
    type(os_data_t) :: os_data
    type(model_data_t), target :: model
!    type(model_t), pointer :: model
!    type(model_list_t) :: model_list
    type(flavor_t) :: flv_h
    type(flavor_t), dimension(2,1) :: flv_hbb, flv_hgg
    type(unstable_config_t), allocatable :: unstable

    write (u, "(A)")  "* Test output: decays_1"
    write (u, "(A)")  "*   Purpose: Set up branching and decay configuration"
    write (u, "(A)")

    write (u, "(A)")  "* Initialize environment"
    write (u, "(A)")

    call os_data_init (os_data)
    call model%init_sm_test ()
!    call model_list%read_model (var_str ("SM"), &
!         var_str ("SM.mdl"), os_data, model)

    call flv_h%init (25, model)
    call flv_hbb(:,1)%init ([5, -5], model)
    call flv_hgg(:,1)%init ([22, 22], model)

    write (u, "(A)")  "* Set up branching and decay"
    write (u, "(A)")

    allocate (unstable)
    unstable%flv = flv_h
    call unstable%init_decays ([var_str ("h_bb"), var_str ("h_gg")], model)
    
    associate (decay => unstable%decay_config(1))
      allocate (decay%term_config (1))
      call decay%init_term (1, flv_hbb, stable = [.true., .true.], model=model)
      decay%integral = 1.234e-3_default
      decay%abs_error = decay%integral * .02_default
    end associate
    
    associate (decay => unstable%decay_config(2))
      allocate (decay%term_config (1))
      call decay%init_term (1, flv_hgg, stable = [.true., .true.], model=model)
      decay%integral = 3.085e-4_default
      decay%abs_error = decay%integral * .08_default
    end associate
    
    call unstable%compute ()
    call unstable%write (u)
    
    write (u, "(A)")
    write (u, "(A)")  "* Cleanup"

    call unstable%final ()
    call model%final ()
    
    write (u, "(A)")
    write (u, "(A)")  "* Test output end: decays_1"
    
  end subroutine decays_1
  
\eatline
\nwendcode{}\nwbegindocs{1384}\nwdocspar
\subsubsection{Cascade decay configuration}
We define a branching configuration with one decay, which is followed
by another branching.
\nwenddocs{}\nwbegincode{1385}\moddef{Decays: execute tests}\plusendmoddef
  call test (decays_2, "decays_2", &
       "cascade decay configuration", &
       u, results)
\nwendcode{}\nwbegincode{1386}\moddef{Decays: tests}\plusendmoddef
  subroutine decays_2 (u)
    integer, intent(in) :: u
    type(os_data_t) :: os_data
!    type(model_t), pointer :: model
!    type(model_list_t) :: model_list
    type(model_data_t), target :: model
    type(flavor_t) :: flv_h, flv_wp, flv_wm
    type(flavor_t), dimension(2,1) :: flv_hww, flv_wud, flv_wen
    type(unstable_config_t), allocatable :: unstable
    type(string_t), dimension(:), allocatable :: decay

    write (u, "(A)")  "* Test output: decays_2"
    write (u, "(A)")  "*   Purpose: Set up cascade branching"
    write (u, "(A)")

    write (u, "(A)")  "* Initialize environment"
    write (u, "(A)")

!    call syntax_model_file_init ()
    call os_data_init (os_data)
!    call model_list%read_model (var_str ("SM"), &
!         var_str ("SM.mdl"), os_data, model)
    call model%init_sm_test ()

    call model%set_unstable (25, [var_str ("h_ww")])
    call model%set_unstable (24, [var_str ("w_ud"), var_str ("w_en")])

    call flv_h%init (25, model)
    call flv_hww(:,1)%init ([24, -24], model)
    call flv_wp%init (24, model)
    call flv_wm%init (-24, model)
    call flv_wud(:,1)%init ([2, -1], model)
    call flv_wen(:,1)%init ([-11, 12], model)
    

    write (u, "(A)")  "* Set up branching and decay"
    write (u, "(A)")

    allocate (unstable)
    unstable%flv = flv_h
    call unstable%flv%get_decays (decay)
    call unstable%init_decays (decay, model)
    
    associate (decay => unstable%decay_config(1))

      decay%integral = 1.e-3_default
      decay%abs_error = decay%integral * .01_default

      allocate (decay%term_config (1))
      call decay%init_term (1, flv_hww, stable = [.false., .true.], model=model)

      select type (w => decay%term_config(1)%prt(1)%c)
      type is (unstable_config_t)

         associate (w_decay => w%decay_config(1))
           w_decay%integral = 2._default
           allocate (w_decay%term_config (1))
           call w_decay%init_term (1, flv_wud, stable = [.true., .true.], &
                model=model)
         end associate
         associate (w_decay => w%decay_config(2))
           w_decay%integral = 1._default
           allocate (w_decay%term_config (1))
           call w_decay%init_term (1, flv_wen, stable = [.true., .true.], &
                model=model)
         end associate
         call w%compute ()

      end select
      
    end associate
    
    call unstable%compute ()
    call unstable%write (u)
    
    write (u, "(A)")
    write (u, "(A)")  "* Cleanup"

    call unstable%final ()
    call model%final ()
!    call model_list%final ()
!    call syntax_model_file_final ()
    
    write (u, "(A)")
    write (u, "(A)")  "* Test output end: decays_2"
    
  end subroutine decays_2
  
\eatline
\nwendcode{}\nwbegindocs{1387}\nwdocspar
\subsubsection{Decay and Process Object}
We define a branching configuration with one decay and connect this
with an actual process object.
\nwenddocs{}\nwbegincode{1388}\moddef{Decays: execute tests}\plusendmoddef
  call test (decays_3, "decays_3", &
       "associate process", &
       u, results)
\nwendcode{}\nwbegincode{1389}\moddef{Decays: tests}\plusendmoddef
  subroutine decays_3 (u)
    integer, intent(in) :: u
    type(os_data_t) :: os_data
    class(model_data_t), pointer :: model
    type(process_library_t), target :: lib
    type(string_t) :: prefix
    type(string_t) :: procname2
    type(process_stack_t) :: process_stack
    type(process_t), pointer :: process
    type(unstable_config_t), allocatable :: unstable

    write (u, "(A)")  "* Test output: decays_3"
    write (u, "(A)")  "*   Purpose: Connect a decay configuration &
         &with a process"
    write (u, "(A)")

    write (u, "(A)")  "* Initialize environment and integrate process"
    write (u, "(A)")

    call os_data_init (os_data)

    prefix = "decays_3"
    call prepare_testbed &
         (lib, process_stack, prefix, os_data, &
         scattering=.false., decay=.true., decay_rest_frame=.false.)

    procname2 = prefix // "_d"
    process => process_stack%get_process_ptr (procname2)
    model => process%get_model_ptr ()
    call process%write (.false., u)
    
    write (u, "(A)")
    write (u, "(A)")  "* Set up branching and decay"
    write (u, "(A)")

    allocate (unstable)
    call unstable%flv%init (25, model)
    call unstable%init_decays ([procname2], model)
    
    write (u, "(A)")  "* Connect decay with process object"
    write (u, "(A)")

    associate (decay => unstable%decay_config(1))
      call decay%connect (process, model)
    end associate
    
    call unstable%compute ()
    call unstable%write (u)
    
    write (u, "(A)")
    write (u, "(A)")  "* Cleanup"

    call unstable%final ()
    call process_stack%final ()

    write (u, "(A)")
    write (u, "(A)")  "* Test output end: decays_3"
    
  end subroutine decays_3

\eatline
\nwendcode{}\nwbegindocs{1390}\nwdocspar
\subsubsection{Decay and Process Object}
Building upon the previous test, we set up a decay instance and generate a
decay event.
\nwenddocs{}\nwbegincode{1391}\moddef{Decays: execute tests}\plusendmoddef
  call test (decays_4, "decays_4", &
       "decay instance", &
       u, results)
\nwendcode{}\nwbegincode{1392}\moddef{Decays: tests}\plusendmoddef
  subroutine decays_4 (u)
    integer, intent(in) :: u
    type(os_data_t) :: os_data
    class(model_data_t), pointer :: model
    type(process_library_t), target :: lib
    type(string_t) :: prefix, procname2
    class(rng_t), allocatable :: rng
    type(process_stack_t) :: process_stack
    type(process_t), pointer :: process
    type(unstable_config_t), allocatable, target :: unstable
    type(unstable_t), allocatable :: instance

    write (u, "(A)")  "* Test output: decays_4"
    write (u, "(A)")  "*   Purpose: Create a decay process and evaluate &
         &an instance"
    write (u, "(A)")

    write (u, "(A)")  "* Initialize environment, process, &
         &and decay configuration"
    write (u, "(A)")

    call os_data_init (os_data)

    prefix = "decays_4"
    call prepare_testbed &
         (lib, process_stack, prefix, os_data, &
         scattering=.false., decay=.true., decay_rest_frame = .false.)

    procname2 = prefix // "_d"
    process => process_stack%get_process_ptr (procname2)
    model => process%get_model_ptr ()

    allocate (unstable)
    call unstable%flv%init (25, model)
    call unstable%init_decays ([procname2], model)
    
    call model%set_unstable (25, [procname2])

    associate (decay => unstable%decay_config(1))
      call decay%connect (process, model)
    end associate
    
    call unstable%compute ()

    allocate (rng_test_t :: rng)

    allocate (instance)
    call instance%init (unstable)
    call instance%import_rng (rng)

    call instance%select_chain ()
    call instance%generate ()
    call instance%write (u)

    write (u, *)
    call instance%write_process_instances (u)
    
    write (u, "(A)")
    write (u, "(A)")  "* Cleanup"

    call instance%final ()
    call process_stack%final ()
    call unstable%final ()
    
    write (u, "(A)")
    write (u, "(A)")  "* Test output end: decays_4"
    
  end subroutine decays_4

\eatline
\nwendcode{}\nwbegindocs{1393}\nwdocspar
\subsubsection{Decay with Parent Process}
We define a scattering process $ss\to ss$ and subsequent decays $s\to f\bar
f$.
\nwenddocs{}\nwbegincode{1394}\moddef{Decays: execute tests}\plusendmoddef
  call test (decays_5, "decays_5", &
       "parent process and decay", &
       u, results)
\nwendcode{}\nwbegincode{1395}\moddef{Decays: tests}\plusendmoddef
  subroutine decays_5 (u)
    integer, intent(in) :: u
    type(os_data_t) :: os_data
    class(model_data_t), pointer :: model
    type(process_library_t), target :: lib
    type(string_t) :: prefix, procname1, procname2
    type(process_stack_t) :: process_stack
    type(process_t), pointer :: process
    type(process_instance_t), allocatable, target :: process_instance
    type(decay_root_config_t), target :: decay_root_config
    type(decay_root_t) :: decay_root
    type(decay_chain_t) :: decay_chain
    integer :: i

    write (u, "(A)")  "* Test output: decays_5"
    write (u, "(A)")  "*   Purpose: Handle a process with subsequent decays"
    write (u, "(A)")

    write (u, "(A)")  "* Initialize environment and parent process"
    write (u, "(A)")

    call os_data_init (os_data)

    prefix = "decays_5"
    procname1 = prefix // "_p"
    procname2 = prefix // "_d"
    call prepare_testbed &
         (lib, process_stack, prefix, os_data, &
         scattering=.true., decay=.true.)

    write (u, "(A)")  "* Initialize decay process"
    write (u, "(A)")

    process => process_stack%get_process_ptr (procname1)
    model => process%get_model_ptr ()
    call model%set_unstable (25, [procname2])

    write (u, "(A)")  "* Initialize decay tree configuration"
    write (u, "(A)")

    call decay_root_config%connect (process, model, process_stack)
    call decay_root_config%compute ()
    call decay_root_config%write (u)

    write (u, "(A)")
    write (u, "(A)")  "* Initialize decay tree"

    allocate (process_instance)
    call process_instance%init (process)
    call process_instance%setup_event_data ()
    call process_instance%init_simulation (1)

    call decay_root%init (decay_root_config, process_instance)

    write (u, "(A)")
    write (u, "(A)")  "* Select decay chain"
    write (u, "(A)")

    decay_root%selected_mci = 1
    !!! Not yet implemented; there is only one term anyway:
    ! call process_instance%select_i_term (decay_root%selected_term)
    decay_root%selected_term = 1
    call decay_root%select_chain ()

    call decay_chain%build (decay_root)
    
    call decay_root%write (u)

    write (u, "(A)")
    write (u, "(A)")  "* Generate event"
    write (u, "(A)")

    call process%generate_unweighted_event (process_instance, &
         decay_root%selected_mci)
    call process_instance%evaluate_event_data ()
    
    call decay_root%generate ()
    
    associate (term => decay_root%term(1))
      do i = 1, 2
         select type (unstable => term%particle_out(i)%c)
         type is (unstable_t)
            associate (decay => unstable%decay(1))
              call pacify (decay%process_instance)
            end associate
         end select
      end do
    end associate
    
    write (u, "(A)")  "* Process instances"
    write (u, "(A)")
    
    call decay_root%write_process_instances (u)

    write (u, "(A)")
    write (u, "(A)")  "* Generate decay chain"
    write (u, "(A)")
    
    call decay_chain%evaluate ()
    call decay_chain%write (u)

    write (u, *)
    write (u, "(A,ES19.12)")  "chain probability =", &
         decay_chain%get_probability ()

    write (u, "(A)")
    write (u, "(A)")  "* Cleanup"

    call decay_chain%final ()
    call decay_root%final ()
    call decay_root_config%final ()
    call process_instance%final ()
    deallocate (process_instance)
    
    call process%final ()
    call process_stack%final ()
    
    write (u, "(A)")
    write (u, "(A)")  "* Test output end: decays_5"
    
  end subroutine decays_5

\eatline
\nwendcode{}\nwbegindocs{1396}\nwdocspar
\subsubsection{Decay as Event Transform}
Again, we define a scattering process $ss\to ss$ and subsequent decays
$s\to f\bar f$.
\nwenddocs{}\nwbegincode{1397}\moddef{Decays: execute tests}\plusendmoddef
  call test (decays_6, "decays_6", &
       "evt_decay object", &
       u, results)
\nwendcode{}\nwbegincode{1398}\moddef{Decays: tests}\plusendmoddef
  subroutine decays_6 (u)
    integer, intent(in) :: u
    type(os_data_t) :: os_data
    class(model_data_t), pointer :: model
    type(process_library_t), target :: lib
    type(string_t) :: prefix, procname1, procname2
    type(process_stack_t) :: process_stack
    type(process_t), pointer :: process
    type(process_instance_t), allocatable, target :: process_instance
    type(evt_decay_t), target :: evt_decay
    integer :: factorization_mode
    logical :: keep_correlations

    write (u, "(A)")  "* Test output: decays_6"
    write (u, "(A)")  "*   Purpose: Handle a process with subsequent decays"
    write (u, "(A)")

    write (u, "(A)")  "* Initialize environment and parent process"
    write (u, "(A)")

    call os_data_init (os_data)

    prefix = "decays_6"
    procname1 = prefix // "_p"
    procname2 = prefix // "_d"
    call prepare_testbed &
         (lib, process_stack, prefix, os_data, &
         scattering=.true., decay=.true.)

    write (u, "(A)")  "* Initialize decay process"

    process => process_stack%get_process_ptr (procname1)
    model => process%get_model_ptr ()
    call model%set_unstable (25, [procname2])

    allocate (process_instance)
    call process_instance%init (process)
    call process_instance%setup_event_data ()
    call process_instance%init_simulation (1)
    
    write (u, "(A)")
    write (u, "(A)")  "* Initialize decay object"

    call evt_decay%connect (process_instance, model, process_stack)

    write (u, "(A)")
    write (u, "(A)")  "* Generate scattering event"

    call process%generate_unweighted_event (process_instance, 1)
    call process_instance%evaluate_event_data ()

    write (u, "(A)")
    write (u, "(A)")  "* Select decay chain and generate event"
    write (u, "(A)")

    call evt_decay%prepare_new_event (1, 1)
    call evt_decay%generate_unweighted ()

    factorization_mode = FM_IGNORE_HELICITY
    keep_correlations = .false.
    call evt_decay%make_particle_set (factorization_mode, keep_correlations)

    call evt_decay%write (u, verbose = .true.)

    write (u, "(A)")
    write (u, "(A)")  "* Cleanup"

    call evt_decay%final ()
    call process_instance%final ()
    deallocate (process_instance)
    
    call process_stack%final ()
    
    write (u, "(A)")
    write (u, "(A)")  "* Test output end: decays_6"
    
  end subroutine decays_6

\eatline
\nwendcode{}\nwbegindocs{1399}%%%%%%%%%%%%%%%%%%%%%%%%%%%%%%%%%%%%%%%%%%%%%%%%%%%%%%%%%%%%%%%%%%%%%%
\section{Shower}
\nwenddocs{}\nwbegincode{1400}\moddef{\code{}shower.f90\edoc{}}\endmoddef
\LA{}File header\RA{}

module shower

\LA{}Use kinds with double\RA{}
\LA{}Use strings\RA{}
  use io_units
  use constants, only: pi, twopi, zero
  use format_utils, only: write_separator
  use unit_tests
  use system_defs, only: LF
  use os_interface
  use xml
  use diagnostics
  use lorentz
  use system_dependencies, only: LHAPDF5_AVAILABLE
  use system_dependencies, only: LHAPDF6_AVAILABLE
  use lhapdf !NODEP!
  use pdf_builtin !NODEP!
  
  use shower_base
  use shower_partons
  use shower_core
  use shower_topythia
  use muli, only: muli_t
  use mlm_matching
  use ckkw_pseudo_weights
  use ckkw_matching
  
  use sm_qcd
  use flavors
  use colors
  use particles
  use state_matrices
  use subevents
  use model_data
  use variables
  use beam_structures
  use process_libraries
  use rng_base
  use mci_base
  use phs_base

  use event_transforms
  ! For model_hadrons and testing
  use models
  ! For Pythia6 communication
  use hep_common
  ! For PDF information
  use processes
  ! For integration tests
  use rng_tao
  use mci_midpoint
  use phs_single
  use prc_core
  use prc_omega

\LA{}Standard module head\RA{}

\LA{}Shower: public\RA{}

\LA{}Shower: types\RA{}

contains

\LA{}Shower: procedures\RA{}

\LA{}Shower: tests\RA{}

end module shower
\eatline
\nwendcode{}\nwbegindocs{1401}\nwdocspar
\subsection{Configuration Parameters}
\nwenddocs{}\nwbegincode{1402}\moddef{Shower: public}\endmoddef
  public :: shower_settings_t
\nwendcode{}\nwbegincode{1403}\moddef{Shower: types}\endmoddef
  type :: shower_settings_t
     logical :: ps_isr_active = .false.
     logical :: ps_fsr_active = .false.
     logical :: ps_use_PYTHIA_shower = .false.
     logical :: hadronization_active = .false.
     logical :: mlm_matching = .false.
     logical :: ckkw_matching = .false.
     logical :: muli_active = .false.

     logical :: ps_PYTHIA_verbose = .false.
     type(string_t) :: ps_PYTHIA_PYGIVE

     !!! values present in PYTHIA and WHIZARDs PS, 
     !!! comments denote corresponding PYTHIA values
     real(default) :: ps_mass_cutoff = 1._default      ! PARJ(82)
     real(default) :: ps_fsr_lambda = 0.29_default     ! PARP(72)
     real(default) :: ps_isr_lambda = 0.29_default     ! PARP(61)
     integer :: ps_max_n_flavors = 5                   ! MSTJ(45)
     logical :: ps_isr_alpha_s_running = .true.        ! MSTP(64)
     logical :: ps_fsr_alpha_s_running = .true.        ! MSTJ(44)
     real(default) :: ps_fixed_alpha_s = 0._default    ! PARU(111)
     logical :: ps_isr_pt_ordered = .false.
     logical :: ps_isr_angular_ordered = .true.        ! MSTP(62)
     real(default) :: ps_isr_primordial_kt_width = 0._default  ! PARP(91)
     real(default) :: ps_isr_primordial_kt_cutoff = 5._default ! PARP(93)
     real(default) :: ps_isr_z_cutoff = 0.999_default  ! 1-PARP(66)
     real(default) :: ps_isr_minenergy = 2._default    ! PARP(65)
     real(default) :: ps_isr_tscalefactor = 1._default
     logical :: ps_isr_only_onshell_emitted_partons = .true.   ! MSTP(63)

     !!! MLM settings
     type(mlm_matching_settings_t) :: ms

     !!! CKKW Matching
     type(ckkw_matching_settings_t) :: ckkw_settings
     type(ckkw_pseudo_shower_weights_t) :: ckkw_weights
   contains 
   \LA{}Shower: shower settings: TBP\RA{}  
  end type shower_settings_t

\eatline
\nwendcode{}\nwbegindocs{1404}Read in the shower, matching and hadronization settings
\nwenddocs{}\nwbegincode{1405}\moddef{Shower: shower settings: TBP}\endmoddef
  procedure :: init => shower_settings_init
\nwendcode{}\nwbegincode{1406}\moddef{Shower: procedures}\endmoddef
  subroutine shower_settings_init (shower_settings, var_list)
    class(shower_settings_t), intent(out) :: shower_settings
    type(var_list_t), intent(in) :: var_list

    shower_settings%ps_isr_active = &
         var_list%get_lval (var_str ("?ps_isr_active"))
    shower_settings%ps_fsr_active = &
         var_list%get_lval (var_str ("?ps_fsr_active"))
    shower_settings%hadronization_active = &
         var_list%get_lval (var_str ("?hadronization_active"))
    shower_settings%mlm_matching = &
         var_list%get_lval (var_str ("?mlm_matching"))
    shower_settings%ckkw_matching = & 
         var_list%get_lval (var_str ("?ckkw_matching"))
    shower_settings%muli_active = &
         var_list%get_lval (var_str ("?muli_active"))

    if (.not. shower_settings%ps_fsr_active .and. &
        .not. shower_settings%ps_isr_active .and. &
        .not. shower_settings%hadronization_active .and. &
        .not. shower_settings%mlm_matching) then
       return
    end if

    shower_settings%ps_use_PYTHIA_shower = &
         var_list%get_lval (var_str ("?ps_use_PYTHIA_shower"))
    shower_settings%ps_PYTHIA_verbose = &
         var_list%get_lval (var_str ("?ps_PYTHIA_verbose"))
    shower_settings%ps_PYTHIA_PYGIVE = &
         var_list%get_sval (var_str ("$ps_PYTHIA_PYGIVE"))
    shower_settings%ps_mass_cutoff = &
         var_list%get_rval (var_str ("ps_mass_cutoff"))
    shower_settings%ps_fsr_lambda = &
         var_list%get_rval (var_str ("ps_fsr_lambda"))
    shower_settings%ps_isr_lambda = &
         var_list%get_rval (var_str ("ps_isr_lambda"))
    shower_settings%ps_max_n_flavors = &
         var_list%get_ival (var_str ("ps_max_n_flavors"))
    shower_settings%ps_isr_alpha_s_running = &
         var_list%get_lval (var_str ("?ps_isr_alpha_s_running"))
    shower_settings%ps_fsr_alpha_s_running = &
         var_list%get_lval (var_str ("?ps_fsr_alpha_s_running"))
    shower_settings%ps_fixed_alpha_s = &
         var_list%get_rval (var_str ("ps_fixed_alpha_s"))
    shower_settings%ps_isr_pt_ordered = &
         var_list%get_lval (var_str ("?ps_isr_pt_ordered"))
    shower_settings%ps_isr_angular_ordered = &
         var_list%get_lval (var_str ("?ps_isr_angular_ordered"))
    shower_settings%ps_isr_primordial_kt_width = &
         var_list%get_rval (var_str ("ps_isr_primordial_kt_width"))
    shower_settings%ps_isr_primordial_kt_cutoff = &
         var_list%get_rval (var_str ("ps_isr_primordial_kt_cutoff"))
    shower_settings%ps_isr_z_cutoff = &
         var_list%get_rval (var_str ("ps_isr_z_cutoff"))
    shower_settings%ps_isr_minenergy = &
         var_list%get_rval (var_str ("ps_isr_minenergy"))
    shower_settings%ps_isr_tscalefactor = &
         var_list%get_rval (var_str ("ps_isr_tscalefactor"))
    shower_settings%ps_isr_only_onshell_emitted_partons = &
         var_list%get_lval (&
         var_str ("?ps_isr_only_onshell_emitted_partons"))

    !!! MLM matching
    shower_settings%ms%mlm_Qcut_ME = &
         var_list%get_rval (var_str ("mlm_Qcut_ME"))
    shower_settings%ms%mlm_Qcut_PS = &
         var_list%get_rval (var_str ("mlm_Qcut_PS"))
    shower_settings%ms%mlm_ptmin = &
         var_list%get_rval (var_str ("mlm_ptmin"))
    shower_settings%ms%mlm_etamax = &
         var_list%get_rval (var_str ("mlm_etamax"))
    shower_settings%ms%mlm_Rmin = &
         var_list%get_rval (var_str ("mlm_Rmin"))
    shower_settings%ms%mlm_Emin = &
         var_list%get_rval (var_str ("mlm_Emin"))
    shower_settings%ms%mlm_nmaxMEjets = &
         var_list%get_ival (var_str ("mlm_nmaxMEjets"))

    shower_settings%ms%mlm_ETclusfactor = &
         var_list%get_rval (var_str ("mlm_ETclusfactor"))
    shower_settings%ms%mlm_ETclusminE = &
         var_list%get_rval (var_str ("mlm_ETclusminE"))
    shower_settings%ms%mlm_etaclusfactor = &
         var_list%get_rval (var_str ("mlm_etaclusfactor"))
    shower_settings%ms%mlm_Rclusfactor = &
         var_list%get_rval (var_str ("mlm_Rclusfactor"))
    shower_settings%ms%mlm_Eclusfactor = &
         var_list%get_rval (var_str ("mlm_Eclusfactor"))

    !!! CKKW matching
    ! TODO
  end subroutine shower_settings_init

\eatline
\nwendcode{}\nwbegindocs{1407} 
\nwenddocs{}\nwbegincode{1408}\moddef{Shower: shower settings: TBP}\plusendmoddef
  procedure :: write => shower_settings_write
\nwendcode{}\nwbegincode{1409}\moddef{Shower: procedures}\plusendmoddef
  subroutine shower_settings_write (object, unit)
    class(shower_settings_t), intent(in) :: object
    integer, intent(in), optional :: unit
    integer :: u
    u = given_output_unit (unit);  if (u < 0)  return
    write (u, "(1x,A)")  "Shower settings:"
    call write_separator (u)
    write (u, "(1x,A)")  "Master switches:"
    write (u, "(3x,A,1x,L1)") &
         "ps_isr_active                = ", object%ps_isr_active
    write (u, "(3x,A,1x,L1)") &
         "ps_fsr_active                = ", object%ps_fsr_active
    write (u, "(1x,A)")  "General settings:"
    if (object%ps_isr_active .or. object%ps_fsr_active) then
       write (u, "(3x,A,1x,L1)") &
            "ps_use_PYTHIA_shower         = ", object%ps_use_PYTHIA_shower
       write (u, "(3x,A,ES19.12)") &
            "ps_mass_cutoff               = ", object%ps_mass_cutoff
       write (u, "(3x,A,1x,I1)") &
            "ps_max_n_flavors             = ", object%ps_max_n_flavors
    else
       write (u, "(3x,A)") " [ISR and FSR off]"
    end if
    if (object%ps_isr_active) then
       write (u, "(1x,A)")  "ISR settings:"
       write (u, "(3x,A,1x,L1)") &
            "ps_isr_pt_ordered            = ", object%ps_isr_pt_ordered
       write (u, "(3x,A,ES19.12)") &
            "ps_isr_lambda                = ", object%ps_isr_lambda
       write (u, "(3x,A,1x,L1)") &
            "ps_isr_alpha_s_running       = ", object%ps_isr_alpha_s_running
       write (u, "(3x,A,ES19.12)") &
            "ps_isr_primordial_kt_width   = ", object%ps_isr_primordial_kt_width
       write (u, "(3x,A,ES19.12)") &
            "ps_isr_primordial_kt_cutoff  = ", &
            object%ps_isr_primordial_kt_cutoff
       write (u, "(3x,A,ES19.12)") &
            "ps_isr_z_cutoff              = ", object%ps_isr_z_cutoff
       write (u, "(3x,A,ES19.12)") &
            "ps_isr_minenergy             = ", object%ps_isr_minenergy
       write (u, "(3x,A,ES19.12)") &
            "ps_isr_tscalefactor          = ", object%ps_isr_tscalefactor
    else if (object%ps_fsr_active) then
       write (u, "(3x,A)") " [ISR off]"       
    end if
    if (object%ps_fsr_active) then
       write (u, "(1x,A)")  "FSR settings:"
       write (u, "(3x,A,ES19.12)") &
            "ps_fsr_lambda                = ", object%ps_fsr_lambda
       write (u, "(3x,A,1x,L1)") &
            "ps_fsr_alpha_s_running       = ", object%ps_fsr_alpha_s_running
    else if (object%ps_isr_active) then
       write (u, "(3x,A)") " [FSR off]"              
    end if
    write (u, "(1x,A)")  "Hadronization settings:"
    write (u, "(3x,A,1x,L1)") &
         "hadronization_active         = ", object%hadronization_active
    write (u, "(1x,A)")  "Matching Settings:"
    write (u, "(3x,A,1x,L1)") &
         "mlm_matching                 = ", object%mlm_matching
    if (object%mlm_matching) then
       call mlm_matching_settings_write (object%ms, u)
    end if
    write (u, "(3x,A,1x,L1)") &
         "ckkw_matching                = ", object%ckkw_matching
    if (object%ckkw_matching) then
       ! TODO ckkw settings etc.
    end if
    write (u, "(1x,A)")  "PYTHIA specific settings:"
    write (u, "(3x,A,1x,L1)") &
         "ps_PYTHIA_verbose            = ", object%ps_PYTHIA_verbose
    write (u, "(3x,A,A,A)") &
         "ps_PYTHIA_PYGIVE             =  '", &
         char(object%ps_PYTHIA_PYGIVE), "'"
  end subroutine shower_settings_write

\eatline
\nwendcode{}\nwbegindocs{1410}\nwdocspar
\subsection{Shower Application Interface}
The wrapper subroutine for the shower and the hadronization
interface and MLM matching. This should be the only subroutine called
from WHIZARD. 
\nwenddocs{}\nwbegincode{1411}\moddef{Shower: public}\plusendmoddef
  public :: apply_shower_particle_set
\nwendcode{}\nwbegincode{1412}\moddef{Shower: procedures}\plusendmoddef
  subroutine apply_shower_particle_set & 
       (evt, valid, vetoed)
    class(evt_shower_t), intent(inout) :: evt
    type(particle_t), dimension(1:2) :: prt_in
    logical, intent(inout) :: valid
    logical, intent(inout) :: vetoed
    real(kind=double) :: pdftest       
    logical, parameter :: debug = .false., to_file = .false.
    integer :: i

    type(mlm_matching_data_t) :: mlm_matching_data
    logical, save :: matching_disabled = .false.

    if (.not. evt%settings%ps_fsr_active .and. &
        .not. evt%settings%ps_isr_active .and. &
        .not. evt%settings%hadronization_active .and. &
        .not. evt%settings%mlm_matching) then
       ! return if nothing to do
       return
    end if

    ! return if already invalid or vetoed
    if (.not. valid .or. vetoed)  return

    if (signal_is_pending ()) return    

    do i = 1, 2
       prt_in(i) = evt%particle_set%get_particle (i)
    end do
    
    ! ensure that lhapdf is initialized
    if (evt%pdf_type .eq. STRF_LHAPDF5) then
       if (evt%settings%ps_isr_active .and. &
            (all (abs (prt_in%get_pdg ()) >= 1000))) then
          call GetQ2max (0, pdftest)
          if (pdftest < epsilon(pdftest)) then
             call msg_fatal ("ISR QCD shower enabled, but LHAPDF not" // &
                  "initialized," // LF // "     aborting simulation")
             return
          end if
       end if
    else if (evt%pdf_type == STRF_PDF_BUILTIN) then
       if (evt%settings%ps_use_PYTHIA_shower) then
          call msg_fatal ("Builtin PDFs cannot be used for PYTHIA showers," &
               // LF // "     aborting simulation")
          return
       end if
    end if
    if (evt%settings%mlm_matching .and. evt%settings%ckkw_matching) then
       call msg_fatal ("Both MLM and CKKW matching activated," // &
            LF // "     aborting simulation")
       return      
    end if

    if (debug)  call evt%settings%write ()
    
    if (evt%settings%ps_use_PYTHIA_shower .or. &
         evt%settings%hadronization_active) then
       if (.not. evt%settings%ps_PYTHIA_verbose) then
          call PYGIVE ('MSTU(12)=12345')
          call PYGIVE ('MSTU(13)=0')
       else 
          call PYGIVE ('MSTU(13)=1')
       end if
    end if

    if (debug)  print *, "Shower: beam checks for mlm_matching"

    if (.not. matching_disabled) then
       !!! Check if the beams are hadrons
       if (all (abs (prt_in%get_pdg ()) <= 18)) then
          mlm_matching_data%is_hadron_collision = .false.
       else if (all (abs (prt_in%get_pdg ()) >= 1000)) then
          mlm_matching_data%is_hadron_collision = .true.
       else 
          call msg_error (" Matching didn't recognize beams setup," // &
               LF // "     disabling matching")
          matching_disabled = .true.
          return
       end if
    end if
    
    if (debug)  print *, "Shower: apply shower"

    !!! SHOWER
    if (evt%settings%ps_use_PYTHIA_shower .or. &
         (.not. evt%settings%ps_fsr_active .and. &
          .not. evt%settings%ps_isr_active .and. &
           evt%settings%hadronization_active)) then
       call apply_PYTHIAshower_particle_set (evt%particle_set, &
            evt%settings, mlm_matching_data%P_ME, evt%model, evt%model_hadrons, &
            valid)
       if (debug)  call pylist(2)
    else
       call apply_WHIZARDshower_particle_set (evt%particle_set, &
            evt%settings, mlm_matching_data%P_ME, evt%model, evt%model_hadrons, &
            evt%os_data, evt%pdf_type, evt%pdf_set, evt%pdf, &
            evt%xmin, evt%xmax, evt%qmin, evt%qmax, valid, vetoed)
       if (vetoed) return
    end if
    if (debug) then
       print *, " after SHOWER"
       call evt%particle_set%write ()
    end if
       
    if (evt%settings%mlm_matching .and. &
         (matching_disabled.eqv..false.)) then
    !!! MLM stage 2 -> PS jets and momenta
       call matching_transfer_PS &
            (mlm_matching_data, evt%particle_set, evt%settings)
    !!! MLM stage 3 -> reconstruct and possible reject
       call mlm_matching_apply (mlm_matching_data, evt%settings%ms, vetoed)
       if (vetoed) then
          call mlm_matching_data_final (mlm_matching_data)
          return
       end if
    endif

!!! HADRONIZATION
    if (evt%settings%hadronization_active) then
       !! Assume that the event record is still in the PYTHIA COMMON BLOCKS
       !! transferred there by one of the shower routines
       if (valid) then
          call apply_PYTHIAhadronization (evt%particle_set, &
               evt%settings, evt%model, evt%model_hadrons, valid)
       end if
    end if
!!! FINAL

    call mlm_matching_data_final (mlm_matching_data)

    if (debug)  print *, "SHOWER+HADRONIZATION+MATCHING finished"

  contains
    
    subroutine shower_set_PYTHIA_error (mstu23)
      ! PYTHIA common blocks
      IMPLICIT DOUBLE PRECISION(A-H, O-Z)
      IMPLICIT INTEGER(I-N)
      COMMON/PYDAT1/MSTU(200),PARU(200),MSTJ(200),PARJ(200)
      SAVE/PYDAT1/

      integer, intent(in) :: mstu23

      MSTU(23) = mstu23
    end subroutine shower_set_PYTHIA_error

    function shower_get_PYTHIA_error () result (mstu23)
      ! PYTHIA common blocks
      IMPLICIT DOUBLE PRECISION(A-H, O-Z)
      IMPLICIT INTEGER(I-N)
      COMMON/PYDAT1/MSTU(200),PARU(200),MSTJ(200),PARJ(200)
      SAVE/PYDAT1/

      integer :: mstu23

      mstu23 = MSTU(23)
    end function shower_get_PYTHIA_error

    subroutine apply_PYTHIAshower_particle_set &
         (particle_set, shower_settings, JETS_ME, model, model_hadrons, valid)

      type(particle_set_t), intent(inout) :: particle_set
      type(particle_set_t) :: pset_reduced
      type(shower_settings_t), intent(in) :: shower_settings
      type(vector4_t), dimension(:), allocatable, intent(inout) :: JETS_ME
      class(model_data_t), intent(in), target :: model
      class(model_data_t), intent(in), target :: model_hadrons
      logical, intent(inout) :: valid
      real(kind=default) :: rand
      
      ! units for transfer from WHIZARD to PYTHIA and back
      integer :: u_W2P, u_P2W
      integer, save :: pythia_initialized_for_NPRUP = 0
      logical, save :: pythia_warning_given = .false.
      logical, save :: msg_written = .false.
      type(string_t) :: remaining_PYGIVE, partial_PYGIVE
      character(len=10) :: buffer

      if (signal_is_pending ()) return      
      
      if (debug) then
         print *, "debugging the shower"
         print *, IDBMUP(1), IDBMUP(2)
         print *, EBMUP, PDFGUP, PDFSUP, IDWTUP
         print *, "NPRUP = ", NPRUP
      end if
      
      ! check if the beam particles are quarks
      if (any (abs(IDBMUP) <= 8)) then
         ! PYTHIA doesn't support these settings
         if (.not. pythia_warning_given) then
            call msg_error ("PYTHIA doesn't support quarks as beam particles," &
                 // LF // "     neglecting ISR, FSR and hadronization")
            pythia_warning_given = .true.
         end if
         return
      end if
      
      call particle_set%reduce (pset_reduced)
      call hepeup_from_particle_set (pset_reduced)
      
      call hepeup_set_event_parameters (proc_id=1)

      u_W2P = free_unit ()
      if (debug .and. to_file) then
         open (unit=u_W2P, status="replace", &
              file="whizardout1.lhe", action="readwrite")  
      else
         open (unit=u_W2P, status="scratch", action="readwrite")
      end if
      call shower_W2P_write_event (u_W2P)
      rewind (u_W2P)
      if (signal_is_pending ()) return          
      write (buffer, "(I10)")  u_W2P
      call PYGIVE ("MSTP(161)="//buffer)
      call PYGIVE ("MSTP(162)="//buffer)
      if (debug)  write (*, "(A)")  buffer
      if (shower_settings%ps_isr_active) then
         call PYGIVE ("MSTP(61)=1")  
      else
         call PYGIVE ("MSTP(61)=0")  !!! switch off ISR
      end if
      if (shower_settings%ps_fsr_active) then
         call PYGIVE ("MSTP(71)=1")  
      else
         call PYGIVE ("MSTP(71)=0")   !!! switch off FSR
      end if
      call PYGIVE ("MSTP(111)=0")     !!! switch off hadronization

      if (pythia_initialized_for_NPRUP >= NPRUP) then
         if (debug)  print *, "calling upinit"
         call upinit
         if (debug)  print *, "returned from upinit"
      else
         write (buffer, "(F10.5)") shower_settings%ps_mass_cutoff
         call PYGIVE ("PARJ(82)="//buffer)
         write (buffer, "(F10.5)") shower_settings%ps_isr_tscalefactor
         call PYGIVE ("PARP(71)="//buffer)    

         write (buffer, "(F10.5)") shower_settings%ps_fsr_lambda
         call PYGIVE ("PARP(72)="//buffer)
         write(buffer, "(F10.5)") shower_settings%ps_isr_lambda
         call PYGIVE ("PARP(61)="//buffer)
         write (buffer, "(I10)") shower_settings%ps_max_n_flavors
         call PYGIVE ("MSTJ(45)="//buffer)
         if (shower_settings%ps_isr_alpha_s_running) then
            call PYGIVE ("MSTP(64)=2")
         else
            call PYGIVE ("MSTP(64)=0")
         end if
         if (shower_settings%ps_fsr_alpha_s_running) then
            call PYGIVE ("MSTJ(44)=2")
         else
            call PYGIVE ("MSTJ(44)=0")
         end if
         write (buffer, "(F10.5)") shower_settings%ps_fixed_alpha_s
         call PYGIVE ("PARU(111)="//buffer)
         write (buffer, "(F10.5)") shower_settings%ps_isr_primordial_kt_width
         call PYGIVE ("PARP(91)="//buffer)
         write (buffer, "(F10.5)") shower_settings%ps_isr_primordial_kt_cutoff
         call PYGIVE ("PARP(93)="//buffer)
         write (buffer, "(F10.5)") 1._double - shower_settings%ps_isr_z_cutoff
         call PYGIVE ("PARP(66)="//buffer)
         write (buffer, "(F10.5)") shower_settings%ps_isr_minenergy
         call PYGIVE ("PARP(65)="//buffer)
         if (shower_settings%ps_isr_only_onshell_emitted_partons) then
            call PYGIVE ("MSTP(63)=0")
         else
            call PYGIVE ("MSTP(63)=2")
         end if
         if (shower_settings%mlm_matching) then
            CALL PYGIVE ("MSTP(62)=2")
            CALL PYGIVE ("MSTP(67)=0")
         end if
         if (debug)  print *, "calling pyinit"
         call PYINIT ("USER", "", "", 0D0)

         call evt%rng%generate (rand)
         write (buffer, "(I10)") floor (rand*900000000)
         call pygive ("MRPY(1)="//buffer)
         call pygive ("MRPY(2)=0")

         if (len(shower_settings%ps_PYTHIA_PYGIVE) > 0) then
            remaining_PYGIVE = shower_settings%ps_PYTHIA_PYGIVE
            do while (len (remaining_PYGIVE)>0)
               call split (remaining_PYGIVE, partial_PYGIVE, ";")
               call PYGIVE (char (partial_PYGIVE))
            end do
            if (shower_get_PYTHIA_error() /= 0) then
               call msg_fatal &
                    (" PYTHIA did not recognize ps_PYTHIA_PYGIVE setting.")
            end if
         end if

         pythia_initialized_for_NPRUP = NPRUP
      end if

      if (.not. msg_written) then
         call msg_message ("Using PYTHIA interface for parton showers")
         msg_written = .true.
      end if      
      call PYEVNT ()
            
      if (debug)  write (*, "(A)")  "called pyevnt"
      
      u_P2W = free_unit ()
      write (buffer, "(I10)")  u_P2W
      call PYGIVE ("MSTP(163)="//buffer)
      if (debug .and. to_file) then
         open (unit = u_P2W, file="pythiaout.lhe", status="replace", &
              action="readwrite")  
      else
         open (unit = u_P2W, status="scratch", action="readwrite")
      end if
      if (debug)  write (*, "(A)")  "calling PYLHEO"
      !!! convert pythia /PYJETS/ to lhef given in MSTU(163)=u_P2W
      call PYLHEO
      !!! read and add lhef from u_P2W
      if (signal_is_pending ()) return          
      call shower_add_lhef_to_particle_set &
           (particle_set, u_P2W, model, model_hadrons)
      close (unit=u_P2W)
      
      !!! Transfer momenta of the partons in the final state of 
      !!!     the hard initeraction
      if (shower_settings%mlm_matching)  &
           call get_ME_momenta_from_PYTHIA (JETS_ME)

      if (shower_get_PYTHIA_error () > 0) then
         !!! clean up, discard shower and exit
         call shower_set_PYTHIA_error (0)
         valid = .false.
      end if
      close (unit=u_W2P)
    end subroutine apply_PYTHIAshower_particle_set

    subroutine apply_WHIZARDshower_particle_set & 
         (particle_set, shower_settings, JETS_ME, model, model_hadrons, &
         os_data, pdf_type, pdf_set, pdf, &
         xmin, xmax, qmin, qmax, valid, vetoed)
      type(particle_set_t), intent(inout) :: particle_set
      type(shower_settings_t), intent(in) :: shower_settings
      type(vector4_t), dimension(:), allocatable, intent(inout) :: JETS_ME
      class(model_data_t), intent(in), target :: model
      class(model_data_t), intent(in), target :: model_hadrons
      type(os_data_t), intent(in) :: os_data
      type(lhapdf_pdf_t), intent(inout) :: pdf
      integer, intent(in) :: pdf_set, pdf_type
      real(double), intent(in) :: xmin, xmax, qmin, qmax
      logical, intent(inout) :: valid
      logical, intent(out) :: vetoed

      type(muli_t), save :: mi
      type(shower_t) :: shower
      type(parton_t), dimension(:), allocatable, target :: partons, hadrons
      type(parton_pointer_t), dimension(:), allocatable :: &
           parton_pointers, final_ME_partons
      real(kind=default) :: mi_scale, ps_scale, shat, phi
      type(parton_pointer_t) :: temppp
      type(particle_t) :: prt
      integer, dimension(:), allocatable :: connections
      integer :: n_loop, i, j, k
      integer :: n_hadrons, n_in, n_out
      integer :: n_int
      integer :: max_color_nr
      integer, dimension(2) :: col_array
      integer, dimension(1) :: parent
      logical, save :: msg_written = .false.
      integer, dimension(2,4) :: color_corr
      integer :: u_S2W      

      vetoed = .false.

      if (signal_is_pending ()) return          
      
      if (debug) print *, "Transfer settings from shower_settings to shower"
      call shower_set_D_Min_t (shower_settings%ps_mass_cutoff**2)
      call shower_set_D_Lambda_fsr (shower_settings%ps_fsr_lambda)
      call shower_set_D_Lambda_isr (shower_settings%ps_isr_lambda)
      call shower_set_D_Nf (shower_settings%ps_max_n_flavors)
      call shower_set_D_running_alpha_s_fsr &
           (shower_settings%ps_fsr_alpha_s_running)
      call shower_set_D_running_alpha_s_isr &
           (shower_settings%ps_isr_alpha_s_running)
      call shower_set_D_constantalpha_s &
           (shower_settings%ps_fixed_alpha_s)
      call shower_set_isr_pt_ordered &
           (shower_settings%ps_isr_pt_ordered)
      Call shower_set_primordial_kt_width &
           (shower_settings%ps_isr_primordial_kt_width)
      call shower_set_primordial_kt_cutoff &
           (shower_settings%ps_isr_primordial_kt_cutoff)
      shower%isr_angular_ordered = shower_settings%ps_isr_angular_ordered
      shower%pdf_set = pdf_set
      shower%pdf_type = pdf_type
      shower%maxz_isr = shower_settings%ps_isr_z_cutoff
      shower%minenergy_timelike = shower_settings%ps_isr_minenergy
      shower%tscalefactor_isr = shower_settings%ps_isr_tscalefactor
      call shower_set_rng(evt%rng)
      
      if (.not. msg_written) then
         call msg_message ("Using WHIZARD's internal showering")
         msg_written = .true.
      end if

      n_loop = 0
      TRY_SHOWER: do ! just a loop to be able to discard events
         n_loop = n_loop + 1
         ! TODO: (bcn 2014-12-02) Can this loop even occur twice?
         if (n_loop > 1000) call msg_fatal &
              ("Shower: too many loops (try_shower)")
         if (pdf_type == STRF_LHAPDF6) then
            call shower%create (xmin, xmax, qmin, qmax, pdf)
         else
            call shower%create (xmin, xmax, qmin, qmax)
         end if
         if (signal_is_pending ()) return             
         max_color_nr = 0

         n_hadrons = 0
         n_in = 0
         n_out = 0
         do i = 1, particle_set%get_n_tot ()
            prt = particle_set%get_particle (i)
            if (prt%get_status () == PRT_BEAM) &
                 n_hadrons = n_hadrons + 1
            if (prt%get_status () == PRT_INCOMING) &
                 n_in = n_in + 1
            if (prt%get_status () == PRT_OUTGOING) &
                 n_out = n_out + 1
         end do

         allocate (connections (1:particle_set%get_n_tot ()))
         connections = 0

         allocate (hadrons (1:2))
         allocate (partons (1:n_in+n_out))
         allocate (parton_pointers (1:n_in+n_out))

         j=0
         if (n_hadrons > 0) then
            if (debug) print *, "Transfer hadrons from particle_set to hadrons"
            do i = 1, particle_set%get_n_tot ()
               prt = particle_set%get_particle (i)
               if (prt%get_status () == PRT_BEAM) then
                  j = j + 1
                  hadrons(j)%nr = shower%get_next_free_nr ()
                  hadrons(j)%momentum = prt%get_momentum ()
                  hadrons(j)%t = hadrons(j)%momentum**2
                  hadrons(j)%type = prt%get_pdg ()
                  col_array = prt%get_color ()
                  hadrons(j)%c1 = col_array(1)
                  hadrons(j)%c2 = col_array(2)
                  max_color_nr = max (max_color_nr, abs(hadrons(j)%c1), &
                       abs(hadrons(j)%c2))
                  hadrons(j)%interactionnr = 1
                  connections(i) = j
               end if
            end do
         end if

         j = 0
         if (debug) print *, "Transfer incoming partons from particle_set to partons"
         do i = 1, particle_set%get_n_tot ()
            prt = particle_set%get_particle (i)
            if (prt%get_status () == PRT_INCOMING) then
               j = j+1
               partons(j)%nr = shower%get_next_free_nr ()
               partons(j)%momentum = prt%get_momentum ()
               partons(j)%t = partons(j)%momentum**2
               partons(j)%type = prt%get_pdg ()
               col_array = prt%get_color ()
               partons(j)%c1 = col_array (1)
               partons(j)%c2 = col_array (2)
               parton_pointers(j)%p => partons(j)
               max_color_nr = max (max_color_nr, abs (partons(j)%c1), &
                    abs (partons(j)%c2))
               connections(i)=j
               ! insert dependences on hadrons
               if (prt%get_n_parents () == 1) then
                  parent = prt%get_parents ()
                  partons(j)%initial => hadrons (connections (parent(1)))
                  partons(j)%x = space_part_norm (partons(j)%momentum) / &
                                 space_part_norm (partons(j)%initial%momentum)
               end if
            end if
         end do
         if (signal_is_pending ()) return             
         if (debug) print *, "Transfer outgoing partons from particle_set to partons"
         do i = 1, particle_set%get_n_tot ()
            prt = particle_set%get_particle (i)
            if (prt%get_status () == PRT_OUTGOING) then               
               j = j + 1
               partons(j)%nr = shower%get_next_free_nr ()
               partons(j)%momentum = prt%get_momentum ()
               partons(j)%t = partons(j)%momentum**2
               partons(j)%type = prt%get_pdg ()
               col_array = prt%get_color ()
               partons(j)%c1 = col_array(1)
               partons(j)%c2 = col_array(2)
               parton_pointers(j)%p => partons(j)
               max_color_nr = max (max_color_nr, abs &
                    (partons(j)%c1), abs (partons(j)%c2))
               connections(i) = j
            end if
         end do

         deallocate (connections)

         if (debug) print *, "Insert partons in shower"
         call shower%set_next_color_nr (1 + max_color_nr)
         call shower%add_interaction_2ton_CKKW &
              (parton_pointers, shower_settings%ckkw_weights)

         if (signal_is_pending ()) return             
         if (shower_settings%muli_active) then
            if (debug) print *, "Activate multiple interactions"
            !!! Initialize muli pdf sets, unless initialized
            if (mi%is_initialized ()) then
               call mi%restart ()
            else
               call mi%initialize (&
                    GeV2_scale_cutoff=D_Min_t, &
                    GeV2_s=shower_interaction_get_s &
                    (shower%interactions(1)%i), &
                    muli_dir=char(os_data%whizard_mulipath))
            end if

            !!! initial interaction
            call mi%apply_initial_interaction ( &
                 GeV2_s=shower_interaction_get_s(shower%interactions(1)%i), &
                 x1=shower%interactions(1)%i%partons(1)%p%parent%x, &
                 x2=shower%interactions(1)%i%partons(2)%p%parent%x, &
                 pdg_f1=shower%interactions(1)%i%partons(1)%p%parent%type, &
                 pdg_f2=shower%interactions(1)%i%partons(2)%p%parent%type, &
                 n1=shower%interactions(1)%i%partons(1)%p%parent%nr, &
                 n2=shower%interactions(1)%i%partons(2)%p%parent%nr)
         end if

         if (signal_is_pending ()) return             
         
         if (shower_settings%ckkw_matching) then
            if (debug) print *, "Apply CKKW matching"
            call ckkw_matching_apply (shower, &
                 shower_settings%ckkw_settings, &
                 shower_settings%ckkw_weights, vetoed)
            if (vetoed) then
               return
            end if
         end if

         if (shower_settings%ps_isr_active) then
            i = 0
            BRANCHINGS: do
               i = i+1
               if (signal_is_pending ()) return                   
               if (shower_settings%muli_active) then
                  call mi%generate_gev2_pt2 &
                       (shower%get_ISR_scale (), mi_scale)
               else
                  mi_scale = 0.0
               end if

               !!! Shower: debugging
               !!! shower%generate_next_isr_branching returns a pointer to 
               !!! the parton with the next ISR-branching, this parton's 
               !!! scale is the scale of the next branching 
               ! temppp=shower%generate_next_isr_branching_veto ()
               temppp = shower%generate_next_isr_branching ()
                  
               if (.not. associated (temppp%p) .and. &
                    mi_scale < D_Min_t) then
                  exit BRANCHINGS
               end if
               !!! check if branching or interaction occurs next
               if (associated (temppp%p)) then
                  ps_scale = abs(temppp%p%t)
               else
                  ps_scale = 0._default
               end if
               if (mi_scale > ps_scale) then
                  !!! discard branching evolution lower than mi_scale
                  call shower%set_max_ISR_scale (mi_scale)
                  if (associated (temppp%p)) &
                       call parton_set_simulated(temppp%p, .false.)
                 
                  !!! execute new interaction
                  deallocate (partons)
                  deallocate (parton_pointers)
                  allocate (partons(1:4))
                  allocate (parton_pointers(1:4))
                  do j = 1, 4
                     partons(j)%nr = shower%get_next_free_nr ()
                     partons(j)%belongstointeraction = .true.
                     parton_pointers(j)%p => partons(j)
                  end do
                  call mi%generate_partons (partons(1)%nr, partons(2)%nr, &
                       partons(1)%x, partons(2)%x, &
                       partons(1)%type, partons(2)%type, &
                       partons(3)%type, partons(4)%type)
                  !!! calculate momenta
                  shat = partons(1)%x *partons(2)%x * &
                       shower_interaction_get_s(shower%interactions(1)%i)
                  partons(1)%momentum = [0.5_default * sqrt(shat), &
                       zero, zero, 0.5_default*sqrt(shat)]
                  partons(2)%momentum = [0.5_default * sqrt(shat), &
                       zero, zero, -0.5_default*sqrt(shat)]
                  call parton_set_initial (partons(1), &
                       shower%interactions(1)%i%partons(1)%p%initial)
                  call parton_set_initial (partons(2), &
                       shower%interactions(1)%i%partons(2)%p%initial)
                  partons(1)%belongstoFSR = .false.
                  partons(2)%belongstoFSR = .false.
                  !!! calculate color connection
                  call mi%get_color_correlations &
                      (shower%get_next_color_nr (), &
                      max_color_nr,color_corr)
                  call shower%set_next_color_nr (max_color_nr)

                  partons(1)%c1 = color_corr(1,1)
                  partons(1)%c2 = color_corr(2,1)
                  partons(2)%c1 = color_corr(1,2)
                  partons(2)%c2 = color_corr(2,2)
                  partons(3)%c1 = color_corr(1,3)
                  partons(3)%c2 = color_corr(2,3)
                  partons(4)%c1 = color_corr(1,4)
                  partons(4)%c2 = color_corr(2,4)

                  call evt%rng%generate (phi)
                  phi = 2 * pi * phi
                  partons(3)%momentum = [0.5_default*sqrt(shat), &
                       sqrt(mi_scale)*cos(phi), &
                       sqrt(mi_scale)*sin(phi), &
                       sqrt(0.25_default*shat - mi_scale)]
                  partons(4)%momentum = [ 0.5_default*sqrt(shat), &
                       -sqrt(mi_scale)*cos(phi), &
                       -sqrt(mi_scale)*sin(phi), &
                       -sqrt(0.25_default*shat - mi_scale)]
                  partons(3)%belongstoFSR = .true.
                  partons(4)%belongstoFSR = .true.

                  call shower%add_interaction_2ton (parton_pointers)
                  n_int = size (shower%interactions)
                  do k = 1, 2
                     call mi%replace_parton &
                       (shower%interactions(n_int)%i%partons(k)%p%initial%nr, &
                        shower%interactions(n_int)%i%partons(k)%p%nr, &
                        shower%interactions(n_int)%i%partons(k)%p%parent%nr, &
                        shower%interactions(n_int)%i%partons(k)%p%type, &
                        shower%interactions(n_int)%i%partons(k)%p%x, &
                        mi_scale)
                  end do
                  call shower%write ()
               else
                  !!! execute the next branching 'found' in the previous step
                  call shower%execute_next_isr_branching (temppp)
                  if (shower_settings%muli_active) then
                     call mi%replace_parton (temppp%p%initial%nr, &
                          temppp%p%child1%nr, temppp%p%nr, &
                          temppp%p%type, temppp%p%x, ps_scale)
                  end if

               end if
            end do BRANCHINGS
               
            call shower%generate_fsr_for_isr_partons ()
         else
            if (signal_is_pending ()) return                
            call shower%simulate_no_isr_shower ()
         end if

         !!! some bookkeeping, needed after the shower is done
         call shower%boost_to_labframe ()
         call shower%generate_primordial_kt ()
         call shower%update_beamremnants ()
         !!! clean-up muli: we should finalize the muli pdf sets when 
         !!!      all runs are done. 
         ! call mi%finalize ()

         if (shower_settings%ps_fsr_active) then
            do i = 1, size (shower%interactions)
               if (signal_is_pending ()) return                   
               call shower%interaction_generate_fsr_2ton &
                    (shower%interactions(i)%i)
            end do
         else
            call shower%simulate_no_fsr_shower ()
         end if
         if (debug) then
            write (*, "(A)")  "SHOWER_FINISHED: "
            call shower%write ()
         end if
            
         if (shower_settings%mlm_matching) then
            !!! transfer momenta of the partons in the final state of 
            !!!        the hard initeraction
            if (signal_is_pending ()) return                
            if (allocated (JETS_ME))  deallocate (JETS_ME)
            call shower%get_final_colored_ME_partons (final_ME_partons)
            if (allocated (final_ME_partons)) then
               allocate (JETS_ME(1:size (final_ME_partons)))
               do i = 1, size (final_ME_partons)
                  !!! transfer
                  JETS_ME(i) = final_ME_partons(i)%p%momentum
               end do
               deallocate (final_ME_partons)
            end if
         end if

         u_S2W = free_unit ()
         if (debug .and. to_file) then
            open (unit=u_S2W, file="showerout.lhe", &
                 status="replace", action="readwrite")  
         else
            open (unit=u_S2W, status="scratch", action="readwrite")
         end if
         call shower%write_lhef (u_S2W)
         call shower_add_lhef_to_particle_set &
              (particle_set, u_S2W, model, model_hadrons)
         close (u_S2W)
         
         !!! move the particle data to the PYTHIA COMMON BLOCKS in case 
         !!! hadronization is active
         if (shower_settings%hadronization_active) then
            if (signal_is_pending ()) return                
            call shower_converttopythia (shower)
         end if
         deallocate (partons)
         deallocate (parton_pointers)
         exit TRY_SHOWER
      end do TRY_SHOWER
      if (debug) then
         call particle_set%write ()
         print *, &
           "----------------------apply_shower_particle_set------------------"
         print *, &
           "-----------------------------------------------------------------"         
         if (size (shower%interactions) >= 2) then
            call shower%write ()
         end if
      end if         
         
      call shower%final ()
      call shower_get_rng(evt%rng)
      !!! clean-up muli: we should finalize the muli pdf sets 
      !!!      when all runs are done. 
      ! call mi%finalize()
      return
    end subroutine apply_WHIZARDshower_particle_set

    subroutine get_ME_momenta_from_PYTHIA (JETS_ME)
      IMPLICIT DOUBLE PRECISION(A-H, O-Z)
      IMPLICIT INTEGER(I-N)
      COMMON/PYJETS/N,NPAD,K(4000,5),P(4000,5),V(4000,5)
      SAVE /PYJETS/

      type(vector4_t), dimension(:), allocatable :: JETS_ME
      real(kind=default), dimension(:,:), allocatable :: pdum 
      integer :: i, j, n_jets

      if (allocated (JETS_ME))  deallocate (JETS_ME)
      if (allocated (pdum))  deallocate (pdum)

      if (signal_is_pending ()) return          
      !!! final ME partons start in 7th row of event record
      i = 7
      !!! find number of jets
      n_jets = 0
      do
         if (K(I,1) /= 21) exit
         if ((K(I,2) == 21) .or. (abs(K(I,2)) <= 6)) then
            n_jets = n_jets + 1
         end if
         i = i + 1 
      end do

      if (n_jets == 0) return
      allocate (JETS_ME(1:n_jets))
      allocate (pdum(1:n_jets,4))

      !!! transfer jets
      i = 7
      j = 1
      pdum = p
      do
         if (K(I,1) /= 21) exit
         if ((K(I,2) == 21) .or. (abs(K(I,2)).le.6)) then
            JETS_ME(j)= vector4_moving (pdum(I,4), & 
              vector3_moving ( [pdum(I,1),pdum(I,2),pdum(I,3)] ))
            j = j + 1
         end if
         i = i + 1
      end do
    end subroutine get_ME_momenta_from_PYTHIA
    
    subroutine matching_transfer_PS &
         (data, particle_set, settings)
      !!! transfer partons after parton shower to data%P_PS
      type(mlm_matching_data_t), intent(inout) :: data
      type(particle_set_t), intent(in) :: particle_set
      type(shower_settings_t), intent(in) :: settings
      integer :: i, j, n_jets_PS
      integer, dimension(2) :: col
      type(particle_t) :: tempprt
      real(double) :: eta
      type(vector4_t) :: p_tmp      

      !!! loop over particles and extract final colored ones with eta<etamax
      n_jets_PS = 0
      do i = 1, particle_set%get_n_tot ()
         if (signal_is_pending ()) return             
         tempprt = particle_set%get_particle (i)
         if (tempprt%get_status () /= PRT_OUTGOING) cycle
         col = tempprt%get_color ()
         if (all (col == 0)) cycle
         if (data%is_hadron_collision) then
            p_tmp = tempprt%get_momentum ()
            if (energy (p_tmp) - longitudinal_part (p_tmp) < 1.E-10_default .or. &
                energy (p_tmp) + longitudinal_part (p_tmp) < 1.E-10_default) then
               eta = pseudorapidity (p_tmp)
            else
               eta = rapidity (p_tmp)
            end if
            if (eta > settings%ms%mlm_etaClusfactor * &
                 settings%ms%mlm_etamax)  then
               if (debug) then
                  print *, "REJECTING"
                  call tempprt%write ()
               end if
               cycle
            end if
         end if
         n_jets_PS = n_jets_PS + 1
      end do

      allocate (data%P_PS(1:n_jets_PS))
      if (debug)  write (*, "(A,1x,I0)")  "n_jets_ps =", n_jets_ps

      j = 1
      do i = 1, particle_set%get_n_tot ()
         tempprt = particle_set%get_particle (i)
         if (tempprt%get_status () /= PRT_OUTGOING) cycle
         col = tempprt%get_color ()
         if(all(col == 0)) cycle
         if (data%is_hadron_collision) then
            p_tmp = tempprt%get_momentum ()
            if (energy (p_tmp) - longitudinal_part (p_tmp) < 1.E-10_default .or. &
                energy (p_tmp) + longitudinal_part (p_tmp) < 1.E-10_default) then
               eta = pseudorapidity (p_tmp)
            else
               eta = rapidity (p_tmp)
            end if
            if (eta > settings%ms%mlm_etaClusfactor * &
                 settings%ms%mlm_etamax) cycle
         end if
         data%P_PS(j) = tempprt%get_momentum ()
         j = j + 1
      end do
    end subroutine matching_transfer_PS
    
    subroutine apply_PYTHIAhadronization &
         (particle_set, shower_settings, model, model_hadrons, valid)
      type(particle_set_t), intent(inout) :: particle_set
      type(shower_settings_t), intent(in) :: shower_settings
      class(model_data_t), intent(in), target :: model
      class(model_data_t), intent(in), target :: model_hadrons
      logical, intent(inout) :: valid
      integer :: u_W2P, u_P2W
      type(string_t) :: remaining_PYGIVE, partial_PYGIVE
      logical, save :: msg_written = .false.
      character(len=10) :: buffer

      if (.not. shower_settings%hadronization_active)  return
      if (.not. valid) return
      if (signal_is_pending ()) return          

      u_W2P = free_unit ()
      if (debug) then
         open (unit=u_W2P, status="replace", file="whizardout.lhe", &
              action="readwrite") 
      else
         open (unit=u_W2P, status="scratch", action="readwrite")
      end if
      call shower_W2P_write_event (u_W2P)
      rewind (u_W2P)
      write (buffer, "(I10)")  u_W2P
      call PYGIVE ("MSTP(161)=" // buffer)
      call PYGIVE ("MSTP(162)=" // buffer)

      !!! Assume that the event is still present in the PYTHIA common blocks
      ! call pygive ("MSTP(61)=0")  ! switch off ISR
      ! call pygive ("MSTP(71)=0")  ! switch off FSR

      if (.not. shower_settings%ps_use_PYTHIA_shower .and. &
           len(shower_settings%ps_PYTHIA_PYGIVE) > 0) then
         remaining_PYGIVE = shower_settings%ps_PYTHIA_PYGIVE
         do while (len(remaining_PYGIVE) > 0)
            if (signal_is_pending ()) return                            
            call split (remaining_PYGIVE, partial_PYGIVE, ";")
            call PYGIVE (char (partial_PYGIVE))
         end do
         if (shower_get_PYTHIA_error () /= 0) then
            call msg_fatal ("PYTHIA didn't recognize ps_PYTHIA_PYGIVE setting")
         end if
      end if

      if (.not. (shower_settings%ps_use_PYTHIA_shower .and. &
           (shower_settings%ps_isr_active.or. &
           shower_settings%ps_fsr_active))) then
         if (len(shower_settings%ps_PYTHIA_PYGIVE) > 0) then
            remaining_PYGIVE = shower_settings%ps_PYTHIA_PYGIVE
            do while (len(remaining_PYGIVE) > 0)
               if (signal_is_pending ()) return                   
               call split (remaining_PYGIVE, partial_PYGIVE, ";")
               call PYGIVE (char(partial_PYGIVE))
            end do
            if (shower_get_PYTHIA_error () /= 0) then
               call msg_fatal &
                    ("PYTHIA did not recognize ps_PYTHIA_PYGIVE setting")
            end if
         end if
      end if

      if (.not.msg_written) then
         call msg_message &
              ("Using PYTHIA interface for hadronization and decays")
         msg_written = .true.
      end if

      call PYGIVE ("MSTP(111)=1") !!! switch on hadronization
      if (signal_is_pending ()) return          
      call PYEXEC

      if (shower_get_PYTHIA_error () > 0) then
         !!! clean up, discard shower and exit
         call shower_set_PYTHIA_error (0)
         close (u_W2P)
         valid = .false.
      else
         !!! convert back
         u_P2W = free_unit ()
         write (buffer, "(I10)")  u_P2W
         call PYGIVE ("MSTP(163)=" // buffer)
         if (debug .and. to_file) then
            open (unit=u_P2W, file="pythiaout2.lhe", status="replace", &
                 action="readwrite")  
         else            
            open (unit=u_P2W, status="scratch", action="readwrite")
         end if
         !!! convert pythia /PYJETS/ to lhef given in MSTU(163)=u1
         call pylheo
         !!! read and add lhef from u_P2W
         if (signal_is_pending ()) return             
         call shower_add_lhef_to_particle_set &
              (particle_set, u_P2W, model, model_hadrons)
         close (u_W2P)
         close (u_P2W)
         valid = .true.
      end if
    end subroutine apply_PYTHIAhadronization

    subroutine shower_W2P_write_event (unit)
      integer, intent(in) :: unit
      type(xml_tag_t), allocatable :: tag_lhef, tag_head, tag_init, &
           tag_event, tag_gen_n, tag_gen_v
      allocate (tag_lhef, tag_head, tag_init, tag_event, &
           tag_gen_n, tag_gen_v)
      call tag_lhef%init (var_str ("LesHouchesEvents"), &
         [xml_attribute (var_str ("version"), var_str ("1.0"))], .true.)
      call tag_head%init (var_str ("header"), .true.)
      call tag_init%init (var_str ("init"), .true.)
      call tag_event%init (var_str ("event"), .true.)
      call tag_gen_n%init (var_str ("generator_name"), .true.)
      call tag_gen_v%init (var_str ("generator_version"), .true.)      
      call tag_lhef%write (unit); write (unit, *)
      call tag_head%write (unit); write (unit, *)
      write (unit, "(2x)", advance = "no")
      call tag_gen_n%write (var_str ("WHIZARD"), unit)
      write (unit, *)
      write (unit, "(2x)", advance = "no")      
      call tag_gen_v%write (var_str ("\LA{}Version\RA{}"), unit)
      write (unit, *)
      call tag_head%close (unit); write (unit, *)
      call tag_init%write (unit); write (unit, *)
      call heprup_write_lhef (unit)
      call tag_init%close (unit); write (unit, *)
      call tag_event%write (unit); write (unit, *)
      call hepeup_write_lhef (unit)
      call tag_event%close (unit); write (unit, *)
      call tag_lhef%close (unit); write (unit, *)
      deallocate (tag_lhef, tag_head, tag_init, tag_event, &
           tag_gen_n, tag_gen_v)
    end subroutine shower_W2P_write_event
  
  end subroutine apply_shower_particle_set
\eatline
\nwendcode{}\nwbegindocs{1413}Read in the LHE file opened in unit {\tt{}u} and add the final
particles to the {\tt{}particle{\char95}set}, the outgoing particles of the existing
{\tt{}particle{\char95}set} are set to be virtual before the outgoing particles from the 
LHEF are added to the {\tt{}particle{\char95}set} as outgoing particles. At the same 
time, it is searched for pairs of existing and new particles with
equal momenta and PDG code, these pairs are marked as connected. All 
remaining existing former outgoing particles are set to be the mothers of all 
remaining new outgoing particles.
\nwenddocs{}\nwbegincode{1414}\moddef{Shower: procedures}\plusendmoddef
  subroutine shower_add_lhef_to_particle_set &
       (particle_set, u, model_in, model_hadrons)
    type(particle_set_t), intent(inout) :: particle_set
    integer, intent(in) :: u
    class(model_data_t), intent(in), target :: model_in
    class(model_data_t), intent(in), target :: model_hadrons
    type(flavor_t) :: flv
    type(color_t) :: col
    class(model_data_t), pointer :: model
    type(particle_t), dimension(:), allocatable :: prt_tmp, prt
    integer :: i, j
    type(vector4_t) :: mom, d_mom
    integer, PARAMETER :: MAXLEN=200
    character(len=maxlen) :: string
    integer :: ibeg, n_tot, n_entries
    integer, dimension(:), allocatable :: relations, mothers
    INTEGER :: NUP,IDPRUP,IDUP,ISTUP
    real(kind=double) :: XWGTUP,SCALUP,AQEDUP,AQCDUP,VTIMUP,SPINUP
    integer :: MOTHUP(1:2), ICOLUP(1:2)
    real(kind=double) :: PUP(1:5)
    real(kind=default) :: pup_dum(1:5)
    character(len=5) :: buffer
    character(len=6) :: strfmt
    logical :: not_found
    STRFMT='(A000)'
    WRITE (STRFMT(3:5),'(I3)') MAXLEN

    rewind (u)

    do
       read (u,*, END=501, ERR=502) STRING
       IBEG = 0
       do
          if (signal_is_pending ()) return              
          IBEG = IBEG + 1
          ! Allow indentation.
          IF (STRING (IBEG:IBEG) .EQ. ' ' .and. IBEG < MAXLEN-6) cycle
          exit
       end do
       IF (string(IBEG:IBEG+6) /= '<event>' .and. &
            string(IBEG:IBEG+6) /= '<event ') cycle
       exit
    end do
    !!! Read first line of event info -> number of entries
    read (u, *, END=503, ERR=504) NUP, IDPRUP, XWGTUP, SCALUP, AQEDUP, AQCDUP
    n_tot = particle_set%get_n_tot ()
    allocate (prt_tmp (1:n_tot+NUP))
    allocate (relations (1:NUP), mothers (1:NUP))
    do i = 1, n_tot
       if (signal_is_pending ()) return           
       prt_tmp (i) = particle_set%get_particle (i)
       if (prt_tmp(i)%get_status () == PRT_OUTGOING .or. &
            prt_tmp(i)%get_status () == PRT_BEAM_REMNANT) then
          call prt_tmp(i)%reset_status (PRT_VIRTUAL)
       end if
    end do

    !!! transfer particles from lhef to particle_set
    !!!...Read NUP subsequent lines with information on each particle.
    n_entries = 1
    mothers = 0
    relations = 0    
    PARTICLE_LOOP: do I = 1, NUP
       read (u,*, END=200, ERR=505) IDUP, ISTUP, MOTHUP(1), MOTHUP(2), &
            ICOLUP(1), ICOLUP(2), (PUP (J),J=1,5), VTIMUP, SPINUP
       if (model_in%test_field (IDUP)) then
          model => model_in
       else if (model_hadrons%test_field (IDUP)) then
          model => model_hadrons
       else
          write (buffer, "(I5)") IDUP
          call msg_error ("Parton " // buffer // &
               " found neither in given model file nor in SM_hadrons")
          return
       end if
       call flv%init (IDUP, model)
       if (IABS(IDUP) == 2212 .or. IABS(IDUP) == 2112) then
          ! PYTHIA sometimes sets color indices for protons and neutrons (?)
          ICOLUP (1) = 0
          ICOLUP (2) = 0
       end if
       call col%init_col_acl (ICOLUP (1), ICOLUP (2))
       !!! Settings for unpolarized particles
       ! particle_set%prt (oldsize+i)%hel = ??
       ! particle_set%prt (oldsize+i)%pol = ??
       if (MOTHUP(1) /= 0) then
          mothers(i) = MOTHUP(1)
       end if
       pup_dum = PUP
       if (pup_dum(4) < 1E-10_default)  cycle
       mom = vector4_moving (pup_dum (4), &
            vector3_moving ([pup_dum (1), pup_dum (2), pup_dum (3)]))
       not_found = .true.
       SCAN_PARTICLES: do j = 1, n_tot
          d_mom = prt_tmp(j)%get_momentum () - mom
          if (abs(d_mom**1) < 1E-8_default .and. &               
                (prt_tmp(j)%get_pdg () == IDUP)) then
             not_found = .false.             
             if (.not. prt_tmp(j)%get_status () == PRT_BEAM .or. &
                  .not. prt_tmp(j)%get_status () == PRT_BEAM_REMNANT) &
                  relations(i) = j
          end if
       end do SCAN_PARTICLES               
       if (not_found) then
          call prt_tmp(n_tot+n_entries)%set_flavor (flv)    
          call prt_tmp(n_tot+n_entries)%set_color (col)
          call prt_tmp(n_tot+n_entries)%set_momentum (mom)
          if (MOTHUP(1) /= 0) then 
             if (relations(MOTHUP(1)) /= 0) then
                call prt_tmp(n_tot+n_entries)%set_parents &
                     ([relations(MOTHUP(1))])             
                call prt_tmp(relations(MOTHUP(1)))%add_child (n_tot+n_entries)
                if (prt_tmp(relations(MOTHUP(1)))%get_status () &
                     == PRT_OUTGOING) &
                     call prt_tmp(relations(MOTHUP(1)))%reset_status &
                     (PRT_VIRTUAL)
             end if
          end if
          call prt_tmp(n_tot+n_entries)%set_status (PRT_OUTGOING)
          n_entries = n_entries + 1
       end if
    end do PARTICLE_LOOP

    allocate (prt (1:n_tot+n_entries-1))
    prt = prt_tmp (1:n_tot+n_entries-1)
    ! transfer to particle_set
    call particle_set%replace (prt)
    deallocate (prt, prt_tmp)

200 continue
    return

501 write(*,*) "READING LHEF failed 501"
    return
502 write(*,*) "READING LHEF failed 502"
    return
503 write(*,*) "READING LHEF failed 503"
    return
504 write(*,*) "READING LHEF failed 504"
    return
505 write(*,*) "READING LHEF failed 505"
    return
  end subroutine shower_add_lhef_to_particle_set
\eatline
\nwendcode{}\nwbegindocs{1415}\nwdocspar
\nwenddocs{}\nwbegincode{1416}\moddef{Shower: procedures}\plusendmoddef
!!!!!!!!!!PYTHIA STYLE!!!!!!!!!!!!!
!!! originally PYLHEF subroutine from PYTHIA 6.4.22

  !C...Write out the showered event to a Les Houches Event File.
  !C...Take MSTP(161) as the input for <init>...</init>

  subroutine pylheo ()

  !C...Double precision and integer declarations.
    IMPLICIT DOUBLE PRECISION(A-H, O-Z)
    IMPLICIT INTEGER(I-N)

    !C...PYTHIA commonblock: only used to provide read/write units and version.
    common /PYPARS/ MSTP(200), PARP(200), MSTI(200), PARI(200)
    common /PYJETS/ N, NPAD, K(4000,5), P(4000,5), V(4000,5)
    save /PYPARS/
    save /PYJETS/

    !C...User process initialization commonblock.
    !C...User process event common block.    
    integer, parameter :: MAXPUP = 100, MAXNUP = 500
    integer :: IDBMUP, PDFGUP, PDFSUP, IDWTUP, NPRUP, LPRUP
    integer :: NUP, IDPRUP, IDUP, ISTUP, MOTHUP, ICOLUP
    real(double) :: EBMUP, XSECUP, XERRUP, XMAXUP
    real(double) :: XWGTUP, SCALUP, AQEDUP, AQCDUP, PUP, VTIMUP, SPINUP    
    integer, parameter :: KSUSY1 = 1000000, KSUSY2 = 2000000
    common /HEPRUP/ &
         IDBMUP(2), EBMUP(2), PDFGUP(2), PDFSUP(2), IDWTUP, NPRUP, &
         XSECUP(MAXPUP), XERRUP(MAXPUP), XMAXUP(MAXPUP), LPRUP(MAXPUP)
    save /HEPRUP/
    common /HEPEUP/ &
         NUP, IDPRUP, XWGTUP, SCALUP, AQEDUP, AQCDUP, IDUP(MAXNUP), &
         ISTUP(MAXNUP), MOTHUP(2,MAXNUP), ICOLUP(2,MAXNUP), &
         PUP(5,MAXNUP), VTIMUP(MAXNUP), SPINUP(MAXNUP)
    save /HEPEUP/
    
    !C...Lines to read in assumed never longer than 200 characters.
    PARAMETER (MAXLEN=200)
    character(len=maxlen) :: string
    
    integer :: LEN, ndangling_color, ndangling_antic, ncolor
    
    !C...Format for reading lines.
    character(len=6) :: strfmt
    STRFMT='(A000)'
    write (STRFMT(3:5),'(I3)') MAXLEN

    !C...Rewind initialization and event files.
    rewind MSTP(161)
    rewind MSTP(162)

    !C...Write header info.
    write (MSTP(163), "(A)")  '<LesHouchesEvents version="1.0">'
    write (MSTP(163), "(A)")  "<!--"
    write (MSTP(163), "(A,I1,A1,I3)")  "File generated with PYTHIA ", &
         MSTP(181), ".", MSTP(182)
    write (MSTP(163), "(A)")  " and the WHIZARD2 interface"
    write (MSTP(163), "(A)")  "-->"

    !C...Loop until finds line beginning with "<init>" or "<init ".
100 READ(MSTP(161),STRFMT,END=400,ERR=400) STRING
    IBEG=0
110 IBEG=IBEG+1
    !C...Allow indentation.
    IF(STRING(IBEG:IBEG).EQ.' '.AND.IBEG.LT.MAXLEN-5) GOTO 110
    IF(STRING(IBEG:IBEG+5).NE.'<init>'.AND.STRING(IBEG:IBEG+5).NE.'<init ') GOTO 100
    
    !C...Read first line of initialization info and get number of processes.
    READ(MSTP(161),'(A)',END=400,ERR=400) STRING
    READ(STRING,*,ERR=400) IDBMUP(1),IDBMUP(2),EBMUP(1),EBMUP(2),PDFGUP(1),PDFGUP(2),PDFSUP(1),PDFSUP(2),IDWTUP,NPRUP

    !C...Copy initialization lines, omitting trailing blanks.
    !C...Embed in <init> ... </init> block.
    WRITE(MSTP(163),'(A)') '<init>'
    do IPR = 0, NPRUP
       IF(IPR.GT.0) READ(MSTP(161),'(A)',END=400,ERR=400) STRING
       LEN=MAXLEN+1
120    LEN=LEN-1
       IF(LEN.GT.1.AND.STRING(LEN:LEN).EQ.' ') GOTO 120
       WRITE(MSTP(163),'(A)',ERR=400) STRING(1:LEN)
    end DO
    write (MSTP(163), "(A)")  "</init>"
    
    !!! Find the numbers of entries of the <event block>
    NENTRIES = 0
    do I = 1, N
       if (K(I,1) == 1 .or. K(I,1) == 2 .or. K(I,1) == 21) then
          NENTRIES = NENTRIES + 1
       end if
    end do
    
    !C...Begin an <event> block. Copy event lines, omitting trailing blanks.
    write (MSTP(163), "(A)")  "<event>"
    write (MSTP(163), *)  NENTRIES, IDPRUP, XWGTUP, SCALUP, AQEDUP, AQCDUP
    
    ndangling_color = 0
    ncolor = 0
    ndangling_antic = 0
    NANTIC = 0
    NNEXTC = 1   ! TODO find next free color number ??
    do I = 1, N
       if (signal_is_pending ()) return             
       if ((K(I,1) >= 1 .and. K(I,1) <= 15) .or. (K(I,1) == 21)) then
          if ((K(I,2).eq.21) .or. (IABS(K(I,2)) <= 8) .or. &
               (IABS(K(I,2)) >= KSUSY1+1 .and. IABS(K(I,2)) <= KSUSY1+8) &
               .or. &
               (IABS(K(I,2)) >= KSUSY2+1 .and. IABS(K(I,2)) <= KSUSY2+8) .or. &
               (IABS(K(I,2)) >= 1000 .and. IABS(K(I,2)) <= 9999) ) then
             if (ndangling_color.eq.0 .and. ndangling_antic.eq.0) then
                ! new color string
                ! Gluon and gluino only color octets implemented so far
                if (K(I,2).eq.21 .or. K(I,2).eq.1000021) then  
                   ncolor = NNEXTC
                   ndangling_color = ncolor
                   NNEXTC = NNEXTC + 1
                   NANTIC = NNEXTC
                   ndangling_antic = NANTIC
                   NNEXTC = NNEXTC + 1
                else if (K(I,2) .gt. 0) then  ! particles to have color
                   ncolor = NNEXTC
                   ndangling_color = ncolor
                   NANTIC = 0
                   NNEXTC = NNEXTC + 1
                else if (K(I,2) .lt. 0) then  ! antiparticles to have anticolor
                   NANTIC = NNEXTC
                   ndangling_antic = NANTIC
                   ncolor = 0
                   NNEXTC = NNEXTC + 1
                end if
             else if(K(I,1).eq.1) then
                ! end of string
                ncolor = ndangling_antic
                NANTIC = ndangling_color
                ndangling_color = 0
                ndangling_antic = 0
             else
                ! inside the string
                if(ndangling_color .ne. 0) then
                   NANTIC = ndangling_color
                   ncolor = NNEXTC
                   ndangling_color = NNEXTC
                   NNEXTC = NNEXTC +1
                else if(ndangling_antic .ne. 0) then
                   ncolor = ndangling_antic
                   NANTIC = NNEXTC
                   ndangling_antic = NNEXTC
                   NNEXTC = NNEXTC +1
                else
                   print *, "ERROR IN PYLHEO"
                end if
             end if
          else
             ncolor = 0
             NANTIC = 0
          end if
       !!! As no intermediate are given out here, assume the 
       !!!   incoming partons to be the mothers
          write (MSTP(163),*)  K(I,2), K(I,1), K(I,3), K(I,3), &
               ncolor, NANTIC, (P(I,J),J=1,5), 0, -9
       end if
    end do
    
    !C..End the <event> block. Loop back to look for next event.
    write (MSTP(163), "(A)")  "</event>"
      
    !C...Successfully reached end of event loop: write closing tag
    !C...and remove temporary intermediate files (unless asked not to).
    write (MSTP(163), "(A)")  "</LesHouchesEvents>"
    return
      
    !!C...Error exit.
400 write(*,*) ' PYLHEO file joining failed!'
    
    return
  end subroutine pylheo

\nwendcode{}\nwbegindocs{1417}\nwdocspar
\subsection{Event Transform}
\nwenddocs{}\nwbegincode{1418}\moddef{Shower: public}\plusendmoddef
  public :: evt_shower_t
\nwendcode{}\nwbegincode{1419}\moddef{Shower: types}\plusendmoddef
  type, extends (evt_t) :: evt_shower_t
     type(shower_settings_t) :: settings
     type(model_t), pointer :: model_hadrons => null ()
     type(lhapdf_pdf_t) :: pdf
     real(double) :: xmin = 0, xmax = 0
     real(double) :: qmin = 0, qmax = 0
     type(os_data_t) :: os_data     
     integer :: pdf_type = STRF_NONE
     integer :: pdf_set = 0
   contains
   \LA{}Shower: evt shower: TBP\RA{}
  end type evt_shower_t
  
\eatline
\nwendcode{}\nwbegindocs{1420}Output.
\nwenddocs{}\nwbegincode{1421}\moddef{Shower: evt shower: TBP}\endmoddef
  procedure :: write => evt_shower_write
\nwendcode{}\nwbegincode{1422}\moddef{Shower: procedures}\plusendmoddef
  subroutine evt_shower_write (object, unit, testflag)
    class(evt_shower_t), intent(in) :: object
    integer, intent(in), optional :: unit
    logical, intent(in), optional :: testflag
    integer :: u
    u = given_output_unit (unit)
    call write_separator (u, 2)
    write (u, "(1x,A)")  "Event transform: shower"
    call write_separator (u)
    call object%base_write (u, testflag = testflag)
    call write_separator (u)
    call object%settings%write (u)
  end subroutine evt_shower_write
    
\eatline
\nwendcode{}\nwbegindocs{1423}Initialize the parameters.  The {\tt{}model{\char95}hadrons} is supposed to be
the SM variant that contains all hadrons that may be generated in the
shower.
\nwenddocs{}\nwbegincode{1424}\moddef{Shower: evt shower: TBP}\plusendmoddef
  procedure :: init => evt_shower_init
\nwendcode{}\nwbegincode{1425}\moddef{Shower: procedures}\plusendmoddef
  subroutine evt_shower_init (evt, settings, model_hadrons, os_data)
    class(evt_shower_t), intent(out) :: evt
    type(shower_settings_t), intent(in) :: settings
    type(model_t), intent(in), target :: model_hadrons
    type(os_data_t), intent(in) :: os_data
    evt%settings = settings
    evt%os_data = os_data
    evt%model_hadrons => model_hadrons
  end subroutine evt_shower_init
  
\eatline
\nwendcode{}\nwbegindocs{1426}Use the process and beam structure to determine PDF parameters.
\nwenddocs{}\nwbegincode{1427}\moddef{Shower: evt shower: TBP}\plusendmoddef
  procedure :: setup_pdf => evt_shower_setup_pdf
\nwendcode{}\nwbegincode{1428}\moddef{Shower: procedures}\plusendmoddef
  subroutine evt_shower_setup_pdf (evt, process, beam_structure)
    class(evt_shower_t), intent(inout) :: evt
    type(process_t), intent(in) :: process
    type(beam_structure_t), intent(in) :: beam_structure
    if (beam_structure%contains ("lhapdf")) then
       if (LHAPDF6_AVAILABLE) then
          evt%pdf_type = STRF_LHAPDF6
       else if (LHAPDF5_AVAILABLE) then
          evt%pdf_type = STRF_LHAPDF5
       end if
       evt%pdf_set = process%get_pdf_set ()
       write (msg_buffer, "(A,I0)")  "Shower: interfacing LHAPDF set #", &
            evt%pdf_set
       call msg_message ()
    else if (beam_structure%contains ("pdf_builtin")) then
       evt%pdf_type = STRF_PDF_BUILTIN
       evt%pdf_set = process%get_pdf_set ()
       write (msg_buffer, "(A,I0)")  "Shower: interfacing PDF builtin set #", &
            evt%pdf_set
       call msg_message ()
    end if
  end subroutine evt_shower_setup_pdf
    
\eatline
\nwendcode{}\nwbegindocs{1429}Prepare: nothing to do here.
\nwenddocs{}\nwbegincode{1430}\moddef{Shower: evt shower: TBP}\plusendmoddef
  procedure :: prepare_new_event => evt_shower_prepare_new_event
\nwendcode{}\nwbegincode{1431}\moddef{Shower: procedures}\plusendmoddef
  subroutine evt_shower_prepare_new_event (evt, i_mci, i_term)
    class(evt_shower_t), intent(inout) :: evt
    integer, intent(in) :: i_mci, i_term
    call evt%reset ()
  end subroutine evt_shower_prepare_new_event

\eatline
\nwendcode{}\nwbegindocs{1432}Here we take the particle set from the previous event transform (assuming
that there is always one) and apply the shower algorithm. The result is
stored in the event transform of the current object. We always return a
probability of unity as we don't have the analytic weight of the combination of
shower, MLM matching and hadronization. A subdivision into multiple event
transformations might be useful. Invalid or vetoed events have to be discarded
by the caller which is why we mark the particle set as invalid. 
\nwenddocs{}\nwbegincode{1433}\moddef{Shower: evt shower: TBP}\plusendmoddef
  procedure :: generate_weighted => evt_shower_generate_weighted
\nwendcode{}\nwbegincode{1434}\moddef{Shower: procedures}\plusendmoddef
  subroutine evt_shower_generate_weighted (evt, probability)
    class(evt_shower_t), intent(inout) :: evt
    real(default), intent(out) :: probability
    logical :: valid, vetoed
    valid = .true.
    vetoed = .false.
    if (evt%previous%particle_set_exists) then
       evt%particle_set = evt%previous%particle_set
       if (evt%settings%ps_use_PYTHIA_shower .or. &
           evt%settings%hadronization_active) then
          call evt%assure_heprup ()
       end if
       if (evt%settings%ckkw_matching) then
          call ckkw_pseudo_shower_weights_init (evt%settings%ckkw_weights)
          call ckkw_fake_pseudo_shower_weights (evt%settings%ckkw_settings, &
               evt%settings%ckkw_weights, evt%particle_set)
       end if
       call apply_shower_particle_set (evt, valid, vetoed)
       probability = 1
       !!! BCN: WK please check: In 2.1.1 vetoed events reduced sim%n_events by
       ! one while invalid events did not. This bookkeeping should be reenabled
       ! to have the correct cross section / luminosity.
       if (valid .and. .not. vetoed) then
          evt%particle_set_exists = .true.
       else
          evt%particle_set_exists = .false.
       end if
    else
       call msg_bug ("Shower: input particle set does not exist")
    end if
  end subroutine evt_shower_generate_weighted

\eatline
\nwendcode{}\nwbegindocs{1435}The particle set was already made by the routine above (or the event was
vetoed).  So we do nothing here.  The factorization parameters are irrelevant.
\nwenddocs{}\nwbegincode{1436}\moddef{Shower: evt shower: TBP}\plusendmoddef
  procedure :: make_particle_set => evt_shower_make_particle_set
\nwendcode{}\nwbegincode{1437}\moddef{Shower: procedures}\plusendmoddef
  subroutine evt_shower_make_particle_set &
       (evt, factorization_mode, keep_correlations, r)
    class(evt_shower_t), intent(inout) :: evt
    integer, intent(in) :: factorization_mode
    logical, intent(in) :: keep_correlations
    real(default), dimension(:), intent(in), optional :: r
  end subroutine evt_shower_make_particle_set

\eatline
\nwendcode{}\nwbegindocs{1438}\nwdocspar
\subsubsection{Matching and hadronization functions}
The HEPRUP (event) common block is needed for the interface to the shower.
Filling of it is triggered by some output file formats.  If these are not
present, the common block is filled with some dummy information. Be generous
with the number of processes in HEPRUP so that PYTHIA only rarely needs to be
reinitialized in case events with higher process ids are generated.
\nwenddocs{}\nwbegincode{1439}\moddef{Shower: evt shower: TBP}\plusendmoddef
  procedure :: assure_heprup => event_shower_assure_heprup
\nwendcode{}\nwbegincode{1440}\moddef{Shower: procedures}\plusendmoddef
  subroutine event_shower_assure_heprup (evt)
    class(evt_shower_t), intent(in) :: evt
    type(particle_t), dimension(2) :: prt
    integer :: i, num_id
    integer, parameter :: min_processes = 10

    num_id = 1
    if (LPRUP (num_id) /= 0)  return

    do i = 1, 2
       prt(i) = evt%particle_set%get_particle (i)
    end do
    call heprup_init ( &
         [ prt(1)%get_pdg (), prt(2)%get_pdg () ] , &
         [ vector4_get_component (prt(1)%get_momentum (), 0), &
           vector4_get_component (prt(2)%get_momentum (), 0) ], &
           num_id, .false., .false.)
    do i = 1, (num_id / min_processes + 1) * min_processes
       call heprup_set_process_parameters (i = i, process_id = &
            i, cross_section = 1._default, error = 1._default)
    end do
  end subroutine event_shower_assure_heprup

\eatline
\nwendcode{}\nwbegindocs{1441}Generate fake ckkw weights. This can be dropped, once information
from the matrix element generation is available.
\nwenddocs{}\nwbegincode{1442}\moddef{Shower: public}\plusendmoddef
  public :: ckkw_fake_pseudo_shower_weights
\nwendcode{}\nwbegincode{1443}\moddef{Shower: procedures}\plusendmoddef
  subroutine ckkw_fake_pseudo_shower_weights &
       (ckkw_pseudo_shower_settings, &
        ckkw_pseudo_shower_weights, particle_set)
    type(ckkw_matching_settings_t), intent(inout) :: &
         ckkw_pseudo_shower_settings
    type(ckkw_pseudo_shower_weights_t), intent(inout) :: &
         ckkw_pseudo_shower_weights
    type(particle_set_t), intent(in) :: particle_set
    type(particle_t) :: prt
    integer :: i, j
    integer :: n
    type(vector4_t) :: momentum

    ckkw_pseudo_shower_settings%alphaS = 1.0_default
    ckkw_pseudo_shower_settings%Qmin = 1.0_default
    ckkw_pseudo_shower_settings%n_max_jets = 3

    n = 2**particle_set%get_n_tot()
    if (allocated (ckkw_pseudo_shower_weights%weights)) then 
       deallocate (ckkw_pseudo_shower_weights%weights)
    end if
    allocate (ckkw_pseudo_shower_weights%weights (1:n))
    do i = 1, n
       momentum = vector4_null
       do j = 1, particle_set%get_n_tot ()
          if (btest (i,j-1)) then
             prt = particle_set%get_particle (j) 
             momentum = momentum + prt%get_momentum ()
          end if
       end do
       if (momentum**1 > 0.0) then
          ckkw_pseudo_shower_weights%weights(i) = 1.0 / (momentum**2)
       end if
    end do
    ! equally distribute the weights by type
    if (allocated (ckkw_pseudo_shower_weights%weights_by_type)) then
       deallocate (ckkw_pseudo_shower_weights%weights_by_type)
    end if
    allocate (ckkw_pseudo_shower_weights%weights_by_type (1:n, 0:4))
    do i = 1, n
       do j = 0, 4
          ckkw_pseudo_shower_weights%weights_by_type(i,j) = &
               0.2 * ckkw_pseudo_shower_weights%weights(i)
       end do
    end do
  end subroutine ckkw_fake_pseudo_shower_weights
  
\eatline
\nwendcode{}\nwbegindocs{1444}\nwdocspar
\subsection{Unit Tests}
\nwenddocs{}\nwbegincode{1445}\moddef{Shower: public}\plusendmoddef
  public :: shower_test
\nwendcode{}\nwbegincode{1446}\moddef{Shower: tests}\endmoddef
  subroutine shower_test (u, results)
    integer, intent(in) :: u
    type(test_results_t), intent(inout) :: results
  \LA{}Shower: execute tests\RA{}
  end subroutine shower_test
  
\eatline
\nwendcode{}\nwbegindocs{1447}\nwdocspar
\subsubsection{Testbed}
This sequence sets up a two-jet process, ready for generating events.
\nwenddocs{}\nwbegincode{1448}\moddef{Shower: tests}\plusendmoddef
  subroutine setup_testbed &
       (prefix, os_data, lib, model_list, process, process_instance)
    type(string_t), intent(in) :: prefix
    type(os_data_t), intent(out) :: os_data
    type(process_library_t), intent(out), target :: lib
    type(model_list_t), intent(out) :: model_list
    class(model_data_t), pointer :: model
    type(model_t), pointer :: model_tmp
    type(process_t), target, intent(out) :: process
    type(process_instance_t), target, intent(out) :: process_instance
    type(var_list_t), pointer :: model_vars
    type(string_t) :: model_name, libname, procname, run_id
    type(process_def_entry_t), pointer :: entry
    type(string_t), dimension(:), allocatable :: prt_in, prt_out
    type(qcd_t) :: qcd
    class(rng_factory_t), allocatable :: rng_factory
    class(prc_core_t), allocatable :: core_template
    class(mci_t), allocatable :: mci_template
    class(phs_config_t), allocatable :: phs_config_template
    real(default) :: sqrts

    model_name = "SM"
    libname = prefix // "_lib"
    procname = prefix // "p"
    run_id = "1"
    
    call os_data_init (os_data)
    allocate (rng_tao_factory_t :: rng_factory)
    allocate (model_tmp)
    call model_list%read_model (model_name, model_name // ".mdl", &
         os_data, model_tmp)
    model_vars => model_tmp%get_var_list_ptr ()
    call var_list_set_real (model_vars, var_str ("me"), 0._default, &
         is_known = .true.)
    model => model_tmp

    call lib%init (libname)

    allocate (prt_in (2), source = [var_str ("e-"), var_str ("e+")])
    allocate (prt_out (2), source = [var_str ("d"), var_str ("dbar")])

    allocate (entry)
    call entry%init (procname, model, n_in = 2, n_components = 1)
    call omega_make_process_component (entry, 1, &
         model_name, prt_in, prt_out, &
         report_progress=.true.)
    call lib%append (entry)

    call lib%configure (os_data)
    call lib%write_makefile (os_data, force = .true.)
    call lib%clean (os_data, distclean = .false.)
    call lib%write_driver (force = .true.)
    call lib%load (os_data)
    
    call process%init (procname, run_id, lib, os_data, &
         qcd, rng_factory, model)
    
    allocate (prc_omega_t :: core_template)
    allocate (mci_midpoint_t :: mci_template)
    allocate (phs_single_config_t :: phs_config_template)

    model => process%get_model_ptr ()

    select type (core_template)
    type is (prc_omega_t)
       call core_template%set_parameters (model = model)
    end select
    call process%init_component &
         (1, core_template, mci_template, phs_config_template)

    sqrts = 1000
    call process%setup_beams_sqrts (sqrts)
    call process%configure_phs ()
    call process%setup_mci ()
    call process%setup_terms ()
    
    call process_instance%init (process)
    call process%integrate (process_instance, 1, 1, 1000)
    call process%final_integration (1)

    call process_instance%setup_event_data ()
    call process_instance%init_simulation (1)
    call process%generate_weighted_event (process_instance, 1)
    call process_instance%evaluate_event_data ()

  end subroutine setup_testbed

\eatline
\nwendcode{}\nwbegindocs{1449}\nwdocspar
\subsubsection{FSR Shower}
We generate a two-jet event and shower it using default settings for
the WHIZARD FSR shower.
\nwenddocs{}\nwbegincode{1450}\moddef{Shower: execute tests}\endmoddef
  call test (shower_1, "shower_1", &
       "disabled shower", &
       u, results)
\nwendcode{}\nwbegincode{1451}\moddef{Shower: tests}\plusendmoddef
  subroutine shower_1 (u)
    integer, intent(in) :: u
    type(os_data_t) :: os_data
    type(process_library_t), target :: lib
    type(model_list_t) :: model_list
    class(model_data_t), pointer :: model
    type(model_t), pointer :: model_hadrons
    type(process_t), target :: process
    type(process_instance_t), target :: process_instance
    integer :: factorization_mode
    logical :: keep_correlations
    class(evt_t), allocatable, target :: evt_trivial
    class(evt_t), allocatable, target :: evt_shower
    type(shower_settings_t) :: settings

    write (u, "(A)")  "* Test output: shower_1"
    write (u, "(A)")  "*   Purpose: Two-jet event with disabled shower"
    write (u, "(A)")

    write (u, "(A)")  "* Initialize environment"
    write (u, "(A)")

    call syntax_model_file_init ()
    call os_data_init (os_data)
    call model_list%read_model (var_str ("SM_hadrons"), var_str ("SM_hadrons.mdl"), &
         os_data, model_hadrons)
    call setup_testbed (var_str ("shower_1"), &
         os_data, lib, model_list, process, process_instance)

    write (u, "(A)")  "* Set up trivial transform"
    write (u, "(A)")
    
    allocate (evt_trivial_t :: evt_trivial)
    model => process%get_model_ptr ()
    call evt_trivial%connect (process_instance, model)
    call evt_trivial%prepare_new_event (1, 1)
    call evt_trivial%generate_unweighted ()

    factorization_mode = FM_IGNORE_HELICITY
    keep_correlations = .false.
    call evt_trivial%make_particle_set (factorization_mode, keep_correlations)

    select type (evt_trivial)
    type is (evt_trivial_t)
       call evt_trivial%write (u)
       call write_separator (u, 2)
    end select

    write (u, "(A)")
    write (u, "(A)")  "* Set up shower event transform"
    write (u, "(A)")

    allocate (evt_shower_t :: evt_shower)
    select type (evt_shower)
    type is (evt_shower_t)
       call evt_shower%init (settings, model_hadrons, os_data)
    end select

    call evt_shower%connect (process_instance, model)
    evt_trivial%next => evt_shower
    evt_shower%previous => evt_trivial

    call evt_shower%prepare_new_event (1, 1)
    call evt_shower%generate_unweighted ()
    call evt_shower%make_particle_set (factorization_mode, keep_correlations)

    select type (evt_shower)
    type is (evt_shower_t)
       call evt_shower%write (u)
       call write_separator (u, 2)
    end select

    write (u, "(A)")
    write (u, "(A)")  "* Cleanup"

    call evt_shower%final ()
    call evt_trivial%final ()
    call process_instance%final ()
    call process%final ()
    call lib%final ()
    call model_hadrons%final ()
    deallocate (model_hadrons)
    !    call model_list%final ()    ! no, would deallocate model twice
    call syntax_model_file_final ()
    
    write (u, "(A)")
    write (u, "(A)")  "* Test output end: shower_1"
    
  end subroutine shower_1
  
\eatline
\nwendcode{}\nwbegindocs{1452}\nwdocspar
\subsubsection{Trivial Test}
We generate a two-jet event and shower it using default settings.
Since the default settings have all 
\nwenddocs{}\nwbegincode{1453}\moddef{Shower: execute tests}\plusendmoddef
  call test (shower_2, "shower_2", &
       "final-state shower", &
       u, results)
\nwendcode{}\nwbegincode{1454}\moddef{Shower: tests}\plusendmoddef
  subroutine shower_2 (u)
    integer, intent(in) :: u
    type(os_data_t) :: os_data
    type(process_library_t), target :: lib
    type(model_list_t) :: model_list
    type(model_t), pointer :: model_hadrons
    class(model_data_t), pointer :: model
    type(process_t), target :: process
    type(process_instance_t), target :: process_instance
    integer :: factorization_mode
    logical :: keep_correlations
    class(evt_t), allocatable, target :: evt_trivial
    class(evt_t), allocatable, target :: evt_shower
    type(shower_settings_t) :: settings

    write (u, "(A)")  "* Test output: shower_2"
    write (u, "(A)")  "*   Purpose: Two-jet event with FSR shower"
    write (u, "(A)")

    write (u, "(A)")  "* Initialize environment"
    write (u, "(A)")

    call syntax_model_file_init ()
    call os_data_init (os_data)
    call model_list%read_model (var_str ("SM_hadrons"), var_str ("SM_hadrons.mdl"), &
         os_data, model_hadrons)
    call setup_testbed (var_str ("shower_2"), &
         os_data, lib, model_list, process, process_instance)
    model => process%get_model_ptr ()
    
    write (u, "(A)")  "* Set up trivial transform"
    write (u, "(A)")
    
    allocate (evt_trivial_t :: evt_trivial)
    call evt_trivial%connect (process_instance, model)
    call evt_trivial%prepare_new_event (1, 1)
    call evt_trivial%generate_unweighted ()

    factorization_mode = FM_IGNORE_HELICITY
    keep_correlations = .false.
    call evt_trivial%make_particle_set (factorization_mode, keep_correlations)

    select type (evt_trivial)
    type is (evt_trivial_t)
       call evt_trivial%write (u)
       call write_separator (u, 2)
    end select

    write (u, "(A)")
    write (u, "(A)")  "* Set up shower event transform"
    write (u, "(A)")

    settings%ps_fsr_active = .true.

    allocate (evt_shower_t :: evt_shower)
    select type (evt_shower)
    type is (evt_shower_t)
       call evt_shower%init (settings, model_hadrons, os_data)
    end select

    call evt_shower%connect (process_instance, model)
    evt_trivial%next => evt_shower
    evt_shower%previous => evt_trivial

    call evt_shower%prepare_new_event (1, 1)
    call evt_shower%generate_unweighted ()
    call evt_shower%make_particle_set (factorization_mode, keep_correlations)

    select type (evt_shower)
    type is (evt_shower_t)
       call evt_shower%write (u, testflag = .true.)
       call write_separator (u, 2)
    end select

    write (u, "(A)")
    write (u, "(A)")  "* Cleanup"

    call evt_shower%final ()
    call evt_trivial%final ()
    call process_instance%final ()
    call process%final ()
    call lib%final ()
    call model_hadrons%final ()
    deallocate (model_hadrons)
    !    call model_list%final ()    ! no, would deallocate model twice
    call syntax_model_file_final ()
    
    write (u, "(A)")
    write (u, "(A)")  "* Test output end: shower_2"
    
  end subroutine shower_2
  
\eatline
\nwendcode{}\nwbegindocs{1455}\nwdocspar
%%%%%%%%%%%%%%%%%%%%%%%%%%%%%%%%%%%%%%%%%%%%%%%%%%%%%%%%%%%%%%%%%%%%%%%%
\section{Complete Events}

This module combines hard processes with decay chains, shower, and
hadronization (not implemented yet) to complete events.  It also
manages the input and output of event records in various formats.
\nwenddocs{}\nwbegincode{1456}\moddef{\code{}events.f90\edoc{}}\endmoddef
\LA{}File header\RA{}

module events
  
\LA{}Use kinds\RA{}
\LA{}Use strings\RA{}
  use io_units
  use format_utils, only: write_separator
  use format_defs, only: FMT_12
  use unit_tests
  use diagnostics
  use os_interface
  use subevents
  use variables
  use expr_base
  use model_data
  use state_matrices
  use particles
  use interactions
  use subevt_expr
  use rng_base
  use process_libraries
  use processes
  use process_stacks
  use event_base
  use event_transforms
  use decays
  use shower

\LA{}Standard module head\RA{}

\LA{}Events: public\RA{}

\LA{}Events: types\RA{}

\LA{}Events: interfaces\RA{}

contains
  
\LA{}Events: procedures\RA{}

\LA{}Events: tests\RA{}

end module events
\eatline
\nwendcode{}\nwbegindocs{1457}\nwdocspar
\subsection{Event configuration}
The parameters govern the transformation of an event to a particle set.

The {\tt{}safety{\char95}factor} reduces the acceptance probability for
unweighting.  If greater than one, excess events become less likely,
but the reweighting efficiency also drops.

The {\tt{}sigma} and {\tt{}n} values, if nontrivial, allow for reweighting
the events according to the requested {\tt{}norm{\char95}mode}.

Various {\tt{}parse{\char95}node{\char95}t} objects are taken from the SINDARIN input.
They encode expressions that apply to the current event.  The
workspaces for evaluating those expressions are set up in the
{\tt{}event{\char95}expr{\char95}t} objects.  Note that these are really pointers,
so the actual nodes are not stored inside the event object.
\nwenddocs{}\nwbegincode{1458}\moddef{Events: types}\endmoddef
  type :: event_config_t
     logical :: unweighted = .false.
     integer :: norm_mode = NORM_UNDEFINED
     integer :: factorization_mode = FM_IGNORE_HELICITY
     logical :: keep_correlations = .false.
     real(default) :: sigma = 1
     integer :: n = 1
     real(default) :: safety_factor = 1
     class(expr_factory_t), allocatable :: ef_selection
     class(expr_factory_t), allocatable :: ef_reweight
     class(expr_factory_t), allocatable :: ef_analysis
   contains
   \LA{}Events: event config: TBP\RA{}
  end type event_config_t

\eatline
\nwendcode{}\nwbegindocs{1459}Output.
\nwenddocs{}\nwbegincode{1460}\moddef{Events: event config: TBP}\endmoddef
  procedure :: write => event_config_write
\nwendcode{}\nwbegincode{1461}\moddef{Events: procedures}\endmoddef
  subroutine event_config_write (object, unit, show_expressions)
    class(event_config_t), intent(in) :: object
    integer, intent(in), optional :: unit
    logical, intent(in), optional :: show_expressions
    integer :: u
    u = given_output_unit (unit)
    write (u, "(3x,A,L1)")  "Unweighted         = ", object%unweighted
    write (u, "(3x,A,A)")   "Normalization      = ", &
         char (event_normalization_string (object%norm_mode))
    write (u, "(3x,A)", advance="no")  "Helicity handling  = "
    select case (object%factorization_mode)
    case (FM_IGNORE_HELICITY)
       write (u, "(A)")  "drop"
    case (FM_SELECT_HELICITY)
       write (u, "(A)")  "select"
    case (FM_FACTOR_HELICITY)
       write (u, "(A)")  "factorize"
    end select
    write (u, "(3x,A,L1)")  "Keep correlations  = ", object%keep_correlations
    if (object%safety_factor /= 1) then
       write (u, "(3x,A," // FMT_12 // ")")  &
            "Safety factor      = ", object%safety_factor
    end if
    if (present (show_expressions)) then
       if (show_expressions) then
          if (allocated (object%ef_selection)) then
             call write_separator (u)
             write (u, "(3x,A)") "Event selection expression:"
             call object%ef_selection%write (u)
          end if
          if (allocated (object%ef_reweight)) then
             call write_separator (u)
             write (u, "(3x,A)") "Event reweighting expression:"
             call object%ef_reweight%write (u)
          end if
          if (allocated (object%ef_analysis)) then
             call write_separator (u)
             write (u, "(3x,A)") "Analysis expression:"
             call object%ef_analysis%write (u)
          end if
       end if
    end if
  end subroutine event_config_write
  
\eatline
\nwendcode{}\nwbegindocs{1462}\nwdocspar
\subsection{The event type}
This is the concrete implementation of the {\tt{}generic{\char95}event{\char95}t} core
that is defined above in the {\tt{}event{\char95}base} module.  The core manages
the main (dressed) particle set pointer and the current values for
weights and sqme.  The implementation adds configuration data,
expressions, process references, and event transforms.

Each event refers to a single elementary process.  This process may be
dressed by a shower, a decay chain etc.  We maintain pointers to a
process instance.

A list of event transforms (class {\tt{}evt{\char95}t}) transform the connected
interactions of the process instance into the final particle set.  In
this list, the first transform is always the trivial one, which just
factorizes the process instance.  Subsequent transforms may apply
decays, etc.  The {\tt{}particle{\char95}set} pointer identifies the particle set
that we want to be analyzed and returned by the event, usually the
last one.

Squared matrix element and weight values: when reading events from
file, the {\tt{}ref} value is the number in the file, while the {\tt{}prc}
value is the number that we calculate from the momenta in the file,
possibly with different parameters.  When generating events the first
time, or if we do not recalculate, the numbers should coincide.
Furthermore, the array of {\tt{}alt} values is copied from an array of
alternative event records.  These values should represent calculated
values.

The {\tt{}sqme} and {\tt{}weight} values mirror corresponding values in the
{\tt{}expr} subobject.  The idea is that when generating or reading
events, the event record is filled first, then the {\tt{}expr} object
acquires copies.  These copies are used for writing events and as targets
for pointer variables in the analysis expression.

All data that involve user-provided expressions (selection, reweighting,
analysis) are handled by the {\tt{}expr} subobject.  In particular, evaluating
the event-selection expression sets the {\tt{}passed} flag.  Furthermore,
the {\tt{}expr} subobject collects data that can be used in the analysis
and should be written to file, including copies of {\tt{}sqme} and {\tt{}weight}.
\nwenddocs{}\nwbegincode{1463}\moddef{Events: public}\endmoddef
  public :: event_t
\nwendcode{}\nwbegincode{1464}\moddef{Events: types}\plusendmoddef
  type, extends (generic_event_t) :: event_t
     type(event_config_t) :: config
     type(process_t), pointer :: process => null ()
     type(process_instance_t), pointer :: instance => null ()
     class(rng_t), allocatable :: rng
     integer :: selected_i_mci = 0
     integer :: selected_i_term = 0
     integer :: selected_channel = 0
     logical :: is_complete = .false.
     class(evt_t), pointer :: transform_first => null ()
     class(evt_t), pointer :: transform_last => null ()
     type(event_expr_t) :: expr
     logical :: selection_evaluated = .false.
     logical :: passed = .false.
     real(default), allocatable :: alpha_qcd_forced
     real(default), allocatable :: scale_forced
     real(default) :: reweight = 1
     logical :: analysis_flag = .false.
   contains
   \LA{}Events: event: TBP\RA{}
  end type event_t

\eatline
\nwendcode{}\nwbegindocs{1465}Finalizer: the list of event transforms is deleted iteratively.
\nwenddocs{}\nwbegincode{1466}\moddef{Events: event: TBP}\endmoddef
  procedure :: final => event_final
\nwendcode{}\nwbegincode{1467}\moddef{Events: procedures}\plusendmoddef
  subroutine event_final (object)
    class(event_t), intent(inout) :: object
    class(evt_t), pointer :: evt
    if (allocated (object%rng))  call object%rng%final ()
    call object%expr%final ()
    do while (associated (object%transform_first))
       evt => object%transform_first
       object%transform_first => evt%next
       call evt%final ()
       deallocate (evt)
    end do
  end subroutine event_final
    
\eatline
\nwendcode{}\nwbegindocs{1468}Output.

Particle set: this is a pointer to one of the event transforms, so it
should suffice to print the latter.
\nwenddocs{}\nwbegincode{1469}\moddef{Events: event: TBP}\plusendmoddef
  procedure :: write => event_write
\nwendcode{}\nwbegincode{1470}\moddef{Events: procedures}\plusendmoddef
  subroutine event_write (object, unit, show_process, show_transforms, &
       show_decay, verbose, testflag)
    class(event_t), intent(in) :: object
    integer, intent(in), optional :: unit
    logical, intent(in), optional :: show_process, show_transforms, show_decay
    logical, intent(in), optional :: verbose
    logical, intent(in), optional :: testflag
    logical :: prc, trans, dec, verb
    class(evt_t), pointer :: evt
    integer :: u, i
    u = given_output_unit (unit)
    prc = .true.;  if (present (show_process))  prc = show_process
    trans = .true.;  if (present (show_transforms))  trans = show_transforms
    dec = .true.;  if (present (show_decay))  dec = show_decay
    verb = .false.;  if (present (verbose))  verb = verbose
    call write_separator (u, 2)
    if (object%is_complete) then
       write (u, "(1x,A)")  "Event"
    else
       write (u, "(1x,A)")  "Event [incomplete]"
    end if
    call write_separator (u)
    call object%config%write (u)
    if (object%sqme_ref_is_known () .or. object%weight_ref_is_known ()) then
       call write_separator (u)
    end if
    if (object%sqme_ref_is_known ()) then
       write (u, "(3x,A,ES19.12)") &
            "Squared matrix el. = ", object%get_sqme_ref ()
       if (object%sqme_alt_is_known ()) then
          do i = 1, object%get_n_alt ()
             write (u, "(5x,A,ES19.12,1x,I0)")  &
                  "alternate sqme   = ", object%get_sqme_alt(i), i
          end do
       end if
    end if
    if (object%weight_ref_is_known ()) then
       write (u, "(3x,A,ES19.12)") &
            "Event weight       = ", object%get_weight_ref ()
       if (object%weight_alt_is_known ()) then
          do i = 1, object%get_n_alt ()
             write (u, "(5x,A,ES19.12,1x,I0)")  &
                  "alternate weight = ", object%get_weight_alt(i), i
          end do
       end if
    end if
    if (object%selected_i_mci /= 0) then
       call write_separator (u)
       write (u, "(3x,A,I0)")  "Selected MCI group = ", object%selected_i_mci
       write (u, "(3x,A,I0)")  "Selected term      = ", object%selected_i_term
       write (u, "(3x,A,I0)")  "Selected channel   = ", object%selected_channel
    end if
    if (object%selection_evaluated) then
       call write_separator (u)
       write (u, "(3x,A,L1)")  "Passed selection   = ", object%passed
       if (object%passed) then
          write (u, "(3x,A,ES19.12)") &
               "Reweighting factor = ", object%reweight
          write (u, "(3x,A,L1)") &
               "Analysis flag      = ", object%analysis_flag
       end if
    end if
    if (associated (object%instance)) then
       if (prc) then
          if (verb) then
             call object%instance%write (u, testflag)
          else
             call object%instance%write_header (u)
          end if
       end if
       if (trans) then
          evt => object%transform_first
          do while (associated (evt))
             select type (evt)
             type is (evt_trivial_t)
                call evt%write (u, verbose = verbose, testflag = testflag)
             type is (evt_decay_t)
                call evt%write (u, show_decay_tree = dec, &
                     show_processes = dec .and. verb, verbose = verb)
             type is (evt_shower_t)
                call evt%write (u)
             end select
             call write_separator (u, 2)
             evt => evt%next
          end do
       else
          call write_separator (u, 2)
       end if
       if (object%expr%subevt_filled) then
          call object%expr%write (u, pacified = testflag)
          call write_separator (u, 2)
       end if
    else
       call write_separator (u, 2)
       write (u, "(1x,A)")  "Process instance: [undefined]"
       call write_separator (u, 2)
    end if
  end subroutine event_write

\eatline
\nwendcode{}\nwbegindocs{1471}\nwdocspar
\subsection{Initialization}
Initialize: set configuration parameters, using a variable list.  We
do not call this {\tt{}init}, because this method name will be used by a type
extension.

The default normalization is {\tt{}NORM{\char95}SIGMA}, since the default
generation mode is weighted.

For unweighted events, we may want to a apply a safety factor to event
rejection.  (By default, this factor is unity and can be ignored.)

We also allocate the trivial event transform, which is always the
first one.
\nwenddocs{}\nwbegincode{1472}\moddef{Events: event: TBP}\plusendmoddef
  procedure :: basic_init => event_init
\nwendcode{}\nwbegincode{1473}\moddef{Events: procedures}\plusendmoddef
  subroutine event_init (event, var_list, n_alt)
    class(event_t), intent(out) :: event
    type(var_list_t), intent(in), optional :: var_list
    integer, intent(in), optional :: n_alt
    type(string_t) :: norm_string
    logical :: polarized_events
    if (present (n_alt)) then
       call event%base_init (n_alt)
       call event%expr%init (n_alt)
    else
       call event%base_init (0)
    end if
    if (present (var_list)) then
       event%config%unweighted = var_list%get_lval (&
            var_str ("?unweighted"))
       norm_string = var_list%get_sval (&
            var_str ("$sample_normalization"))
       event%config%norm_mode = &
            event_normalization_mode (norm_string, event%config%unweighted)
       polarized_events = &
            var_list%get_lval (var_str ("?polarized_events"))
       if (polarized_events) then
          event%config%factorization_mode = FM_SELECT_HELICITY
       else
          event%config%factorization_mode = FM_IGNORE_HELICITY
       end if
       if (event%config%unweighted) then
          event%config%safety_factor = var_list%get_rval (&
               var_str ("safety_factor"))
       end if
    else
       event%config%norm_mode = NORM_SIGMA
    end if
    allocate (evt_trivial_t :: event%transform_first)
    event%transform_last => event%transform_first
  end subroutine event_init
    
\eatline
\nwendcode{}\nwbegindocs{1474}Set the {\tt{}sigma} and {\tt{}n} values in the configuration record that
determine non-standard event normalizations.  If these numbers are not
set explicitly, the default value for both is unity, and event
renormalization has no effect.
\nwenddocs{}\nwbegincode{1475}\moddef{Events: event: TBP}\plusendmoddef
  procedure :: set_sigma => event_set_sigma
  procedure :: set_n => event_set_n
\nwendcode{}\nwbegincode{1476}\moddef{Events: procedures}\plusendmoddef
  elemental subroutine event_set_sigma (event, sigma)
    class(event_t), intent(inout) :: event
    real(default), intent(in) :: sigma
    event%config%sigma = sigma
  end subroutine event_set_sigma

  elemental subroutine event_set_n (event, n)
    class(event_t), intent(inout) :: event
    integer, intent(in) :: n
    event%config%n = n
  end subroutine event_set_n
  
\eatline
\nwendcode{}\nwbegindocs{1477}Append an event transform (decays, etc.).  The transform is not yet
connected to a process.  The transform is then considered to belong to
the event object, and will be finalized together with it.  The
original pointer is removed.

We can assume that the trivial transform is already present in the
event object, at least.
\nwenddocs{}\nwbegincode{1478}\moddef{Events: event: TBP}\plusendmoddef
  procedure :: import_transform => event_import_transform
\nwendcode{}\nwbegincode{1479}\moddef{Events: procedures}\plusendmoddef
  subroutine event_import_transform (event, evt)
    class(event_t), intent(inout) :: event
    class(evt_t), intent(inout), pointer :: evt
    event%transform_last%next => evt
    evt%previous => event%transform_last
    event%transform_last => evt
    evt => null ()
  end subroutine event_import_transform
    
\eatline
\nwendcode{}\nwbegindocs{1480}\nwdocspar
We link the event to an existing process instance.  This
includes the variable list, which is linked to the process variable
list.  Note that this is not necessarily identical to the variable
list used for event initialization.

The variable list will contain pointers to {\tt{}event} subobjects, therefore the
{\tt{}target} attribute.

Once we have a process connected, we can use it to obtain an event
generator instance.

The model and process stack may be needed by event transforms.  The
current model setting may be different from the model in the process
(regarding unstable particles, etc.).  The process stack can be used
for assigning extra processes that we need for the event transforms.
\nwenddocs{}\nwbegincode{1481}\moddef{Events: event: TBP}\plusendmoddef
  procedure :: connect => event_connect
\nwendcode{}\nwbegincode{1482}\moddef{Events: procedures}\plusendmoddef
  subroutine event_connect (event, process_instance, model, process_stack)
    class(event_t), intent(inout), target :: event
    type(process_instance_t), intent(in), target :: process_instance
    class(model_data_t), intent(in), target :: model
    type(process_stack_t), intent(in), optional :: process_stack
    type(string_t) :: id
    integer :: num_id
    class(evt_t), pointer :: evt
    event%process => process_instance%process
    event%instance => process_instance
    id = event%process%get_id ()
    if (id /= "")  call event%expr%set_process_id (id)
    num_id = event%process%get_num_id ()
    if (num_id /= 0)  call event%expr%set_process_num_id (num_id)
    call event%expr%setup_vars (event%process%get_sqrts ())
    call event%expr%link_var_list (event%process%get_var_list_ptr ())
    call event%process%make_rng (event%rng)
    evt => event%transform_first
    do while (associated (evt))
       call evt%connect (process_instance, model, process_stack)
       evt => evt%next
    end do
  end subroutine event_connect

\eatline
\nwendcode{}\nwbegindocs{1483}Set the parse nodes for the associated expressions, individually.  The
parse-node pointers may be null.
\nwenddocs{}\nwbegincode{1484}\moddef{Events: event: TBP}\plusendmoddef
  procedure :: set_selection => event_set_selection
  procedure :: set_reweight => event_set_reweight
  procedure :: set_analysis => event_set_analysis
\nwendcode{}\nwbegincode{1485}\moddef{Events: procedures}\plusendmoddef
  subroutine event_set_selection (event, ef_selection)
    class(event_t), intent(inout) :: event
    class(expr_factory_t), intent(in) :: ef_selection
    allocate (event%config%ef_selection, source = ef_selection)
  end subroutine event_set_selection

  subroutine event_set_reweight (event, ef_reweight)
    class(event_t), intent(inout) :: event
    class(expr_factory_t), intent(in) :: ef_reweight
    allocate (event%config%ef_reweight, source = ef_reweight)
  end subroutine event_set_reweight

  subroutine event_set_analysis (event, ef_analysis)
    class(event_t), intent(inout) :: event
    class(expr_factory_t), intent(in) :: ef_analysis
    allocate (event%config%ef_analysis, source = ef_analysis)
  end subroutine event_set_analysis
  
\eatline
\eatline
\eatline
\nwendcode{}\nwbegindocs{1486}Create evaluation trees from the parse trees.  The {\tt{}target} attribute is
required because the expressions contain pointers to event subobjects.
\nwenddocs{}\nwbegincode{1487}\moddef{Events: event: TBP}\plusendmoddef
  procedure :: setup_expressions => event_setup_expressions
\nwendcode{}\nwbegincode{1488}\moddef{Events: procedures}\plusendmoddef
  subroutine event_setup_expressions (event)
    class(event_t), intent(inout), target :: event
    call event%expr%setup_selection (event%config%ef_selection)
    call event%expr%setup_analysis (event%config%ef_analysis)
    call event%expr%setup_reweight (event%config%ef_reweight)
  end subroutine event_setup_expressions
  
\eatline
\nwendcode{}\nwbegindocs{1489}\nwdocspar
\subsection{Evaluation}
To fill the {\tt{}particle{\char95}set}, i.e., the event record proper, we have
to apply all event transforms in order.  The last transform should
fill its associated particle set, factorizing the state matrix
according to the current settings.  There are several parameters
in the event configuration that control this.

We always fill the particle set for the first transform (the hard
process) and the last transform, if different from the first (the
fully dressed process).

Each event transform is an event generator of its own.  We choose to
generate an \emph{unweighted} event for each of them, even if the master
event is assumed to be weighted.  Thus, the overall event weight is
the one of the hard process only.  (There may be more options in future
extensions.)

We can generate the two random numbers that the factorization needs.
For testing purpose, we allow for providing them explicitly, as an option.
\nwenddocs{}\nwbegincode{1490}\moddef{Events: event: TBP}\plusendmoddef
  procedure :: evaluate_transforms => event_evaluate_transforms
\nwendcode{}\nwbegincode{1491}\moddef{Events: procedures}\plusendmoddef
  subroutine event_evaluate_transforms (event, r)
    class(event_t), intent(inout) :: event
    real(default), dimension(:), intent(in), optional :: r
    class(evt_t), pointer :: evt
    integer :: i_term
    call event%discard_particle_set ()
    call event%check ()
    if (event%instance%is_complete_event ()) then
       call event%instance%select_i_term (i_term)
       event%selected_i_term = i_term
       evt => event%transform_first
       do while (associated (evt))
          call evt%prepare_new_event &
               (event%selected_i_mci, event%selected_i_term)
          evt => evt%next
       end do
       evt => event%transform_first
       do while (associated (evt))
          call evt%generate_unweighted ()
          if (signal_is_pending ())  return
          call evt%make_particle_set (event%config%factorization_mode, &
               event%config%keep_correlations)
          if (signal_is_pending ())  return
          if (.not. evt%particle_set_exists)  exit
          evt => evt%next
       end do
       evt => event%transform_last
       if (associated (evt) .and. evt%particle_set_exists) then
          call event%link_particle_set (evt%particle_set)
       end if
    end if
  end subroutine event_evaluate_transforms
    
\eatline
\nwendcode{}\nwbegindocs{1492}\nwdocspar
Evaluate the event-related expressions, given a valid
{\tt{}particle{\char95}set}.  If {\tt{}update{\char95}sqme} is set, we use the process
instance for the {\tt{}sqme{\char95}prc} value.  The {\tt{}sqme{\char95}ref} value is
always taken from the event record.

Note: without the explicit {\tt{}particle{\char95}set} pointer, some gfortran 4.8
version corrupts its memory.
\nwenddocs{}\nwbegincode{1493}\moddef{Events: event: TBP}\plusendmoddef
  procedure :: evaluate_expressions => event_evaluate_expressions
\nwendcode{}\nwbegincode{1494}\moddef{Events: procedures}\plusendmoddef
  subroutine event_evaluate_expressions (event)
    class(event_t), intent(inout) :: event
    type(particle_set_t), pointer :: particle_set
    if (event%has_valid_particle_set ()) then
       particle_set => event%get_particle_set_ptr ()
       call event%expr%fill_subevt (particle_set)
    end if
    if (event%weight_ref_is_known ()) then
       call event%expr%set (weight_ref = event%get_weight_ref ())
    end if
    if (event%weight_prc_is_known ()) then
       call event%expr%set (weight_prc = event%get_weight_prc ())
    end if
    if (event%excess_prc_is_known ()) then
       call event%expr%set (excess_prc = event%get_excess_prc ())
    end if
    if (event%sqme_ref_is_known ()) then
       call event%expr%set (sqme_ref = event%get_sqme_ref ())
    end if
    if (event%sqme_prc_is_known ()) then
       call event%expr%set (sqme_prc = event%get_sqme_prc ())
    end if
    if (event%has_valid_particle_set ()) then
       call event%expr%evaluate &
            (event%passed, event%reweight, event%analysis_flag)
       event%selection_evaluated = .true.
    end if
  end subroutine event_evaluate_expressions
  
\eatline
\nwendcode{}\nwbegindocs{1495}Report the result of the {\tt{}selection} evaluation.
\nwenddocs{}\nwbegincode{1496}\moddef{Events: event: TBP}\plusendmoddef
  procedure :: passed_selection => event_passed_selection
\nwendcode{}\nwbegincode{1497}\moddef{Events: procedures}\plusendmoddef
  function event_passed_selection (event) result (flag)
    class(event_t), intent(in) :: event
    logical :: flag
    flag = event%passed
  end function event_passed_selection
  
\eatline
\nwendcode{}\nwbegindocs{1498}Set alternate sqme and weight arrays.  This should be merged with
the previous routine, if the expressions are allowed to refer to these
values.
\nwenddocs{}\nwbegincode{1499}\moddef{Events: event: TBP}\plusendmoddef
  procedure :: store_alt_values => event_store_alt_values
\nwendcode{}\nwbegincode{1500}\moddef{Events: procedures}\plusendmoddef
  subroutine event_store_alt_values (event)
    class(event_t), intent(inout) :: event
    if (event%weight_alt_is_known ()) then
       call event%expr%set (weight_alt = event%get_weight_alt ())
    end if
    if (event%sqme_alt_is_known ()) then
       call event%expr%set (sqme_alt = event%get_sqme_alt ())
    end if
  end subroutine event_store_alt_values
  
\eatline
\nwendcode{}\nwbegindocs{1501}\nwdocspar
\subsection{Reset to empty state}
Applying this, current event contents are marked as incomplete but
are not deleted.  In particular, the initialization is kept.
\nwenddocs{}\nwbegincode{1502}\moddef{Events: event: TBP}\plusendmoddef
  procedure :: reset => event_reset
\nwendcode{}\nwbegincode{1503}\moddef{Events: procedures}\plusendmoddef
  subroutine event_reset (event)
    class(event_t), intent(inout) :: event
    class(evt_t), pointer :: evt
    call event%base_reset ()
    event%selected_i_mci = 0
    event%selected_i_term = 0
    event%selected_channel = 0
    event%is_complete = .false.
    call event%expr%reset ()
    event%selection_evaluated = .false.
    event%passed = .false.
    event%analysis_flag = .false.
    if (associated (event%instance)) then
       call event%instance%reset (reset_mci = .true.)
    end if
    if (allocated (event%alpha_qcd_forced))  deallocate (event%alpha_qcd_forced)
    if (allocated (event%scale_forced))  deallocate (event%scale_forced)
    evt => event%transform_first
    do while (associated (evt))
       call evt%reset ()
       evt => evt%next
    end do
  end subroutine event_reset
  
\eatline
\nwendcode{}\nwbegindocs{1504}\nwdocspar
\subsection{Squared Matrix Element and Weight}
Transfer the result of the process instance calculation to the
event record header.
\nwenddocs{}\nwbegincode{1505}\moddef{Events: event: TBP}\plusendmoddef
  procedure :: import_instance_results => event_import_instance_results
\nwendcode{}\nwbegincode{1506}\moddef{Events: procedures}\plusendmoddef
  subroutine event_import_instance_results (event)
    class(event_t), intent(inout) :: event
    if (associated (event%instance)) then
       if (event%instance%has_evaluated_trace ()) then
          call event%set ( &
               sqme_prc = event%instance%get_sqme (), &
               weight_prc = event%instance%get_weight (), &
               excess_prc = event%instance%get_excess () &
               )
       end if
    end if
  end subroutine event_import_instance_results
  
\eatline
\nwendcode{}\nwbegindocs{1507}Duplicate the instance result / the reference result in the event
record.
\nwenddocs{}\nwbegincode{1508}\moddef{Events: event: TBP}\plusendmoddef
  procedure :: accept_sqme_ref => event_accept_sqme_ref
  procedure :: accept_sqme_prc => event_accept_sqme_prc
  procedure :: accept_weight_ref => event_accept_weight_ref
  procedure :: accept_weight_prc => event_accept_weight_prc
\nwendcode{}\nwbegincode{1509}\moddef{Events: procedures}\plusendmoddef
  subroutine event_accept_sqme_ref (event)
    class(event_t), intent(inout) :: event
    if (event%sqme_ref_is_known ()) then
       call event%set (sqme_prc = event%get_sqme_ref ())
    end if
  end subroutine event_accept_sqme_ref
  
  subroutine event_accept_sqme_prc (event)
    class(event_t), intent(inout) :: event
    if (event%sqme_prc_is_known ()) then
       call event%set (sqme_ref = event%get_sqme_prc ())
    end if
  end subroutine event_accept_sqme_prc
  
  subroutine event_accept_weight_ref (event)
    class(event_t), intent(inout) :: event
    if (event%weight_ref_is_known ()) then
       call event%set (weight_prc = event%get_weight_ref ())
    end if
  end subroutine event_accept_weight_ref
  
  subroutine event_accept_weight_prc (event)
    class(event_t), intent(inout) :: event
    if (event%weight_prc_is_known ()) then
       call event%set (weight_ref = event%get_weight_prc ())
    end if
  end subroutine event_accept_weight_prc
  
\eatline
\eatline
\eatline
\eatline
\nwendcode{}\nwbegindocs{1510}Update the weight normalization, just after generation.  Unweighted
and weighted events are generated with a different default
normalization.  The intended normalization is stored in the
configuration record.
\nwenddocs{}\nwbegincode{1511}\moddef{Events: event: TBP}\plusendmoddef
  procedure :: update_normalization => event_update_normalization
\nwendcode{}\nwbegincode{1512}\moddef{Events: procedures}\plusendmoddef
  subroutine event_update_normalization (event, mode_ref)
    class(event_t), intent(inout) :: event
    integer, intent(in), optional :: mode_ref
    integer :: mode_old
    real(default) :: weight, excess
    if (present (mode_ref)) then
       mode_old = mode_ref
    else if (event%config%unweighted) then
       mode_old = NORM_UNIT
    else
       mode_old = NORM_SIGMA
    end if
    weight = event%get_weight_prc ()
    call event_normalization_update (weight, &
         event%config%sigma, event%config%n, &
         mode_new = event%config%norm_mode, &
         mode_old = mode_old)
    call event%set_weight_prc (weight)
    excess = event%get_excess_prc ()
    call event_normalization_update (excess, &
         event%config%sigma, event%config%n, &
         mode_new = event%config%norm_mode, &
         mode_old = mode_old)
    call event%set_excess_prc (excess)
  end subroutine event_update_normalization
  
\eatline
\nwendcode{}\nwbegindocs{1513}\nwdocspar
The event is complete if it has a particle set plus valid entries for
the sqme and weight values.
\nwenddocs{}\nwbegincode{1514}\moddef{Events: event: TBP}\plusendmoddef
  procedure :: check => event_check
\nwendcode{}\nwbegincode{1515}\moddef{Events: procedures}\plusendmoddef
  subroutine event_check (event)
    class(event_t), intent(inout) :: event
    event%is_complete = event%has_valid_particle_set () &
         .and. event%sqme_ref_is_known () &
         .and. event%sqme_prc_is_known () &
         .and. event%weight_ref_is_known () &
         .and. event%weight_prc_is_known ()
    if (event%get_n_alt () /= 0) then
       event%is_complete = event%is_complete &
            .and. event%sqme_alt_is_known () &
            .and. event%weight_alt_is_known ()
    end if
  end subroutine event_check
  
\eatline
\nwendcode{}\nwbegindocs{1516}\nwdocspar
\nwenddocs{}\nwbegindocs{1517}\nwdocspar
\subsection{Generation}
Assuming that we have a valid process associated to the event, we
generate an event.  We complete the event data, then factorize the
spin density matrix and transfer it to the particle set.

When done, we retrieve squared matrix element and weight.  In case of
explicit generation, the reference values coincide with the process
values, so we {\tt{}accept} the latter.

The explicit random number argument {\tt{}r} should be generated by a
random-number generator.  It is taken for the factorization algorithm,
bypassing the event-specific random-number generator.  This is useful
for deterministic testing.
\nwenddocs{}\nwbegincode{1518}\moddef{Events: event: TBP}\plusendmoddef
  procedure :: generate => event_generate
\nwendcode{}\nwbegincode{1519}\moddef{Events: procedures}\plusendmoddef
  subroutine event_generate (event, i_mci, r)
    class(event_t), intent(inout) :: event
    integer, intent(in) :: i_mci
    real(default), dimension(:), intent(in), optional :: r
    call event%reset ()
    event%selected_i_mci = i_mci
    if (event%config%unweighted) then
       call event%process%generate_unweighted_event (event%instance, i_mci)
       if (signal_is_pending ()) return
       call event%instance%evaluate_event_data ()
       call event%instance%normalize_weight ()
    else
       call event%process%generate_weighted_event (event%instance, i_mci)
       if (signal_is_pending ()) return
       call event%instance%evaluate_event_data ()
    end if
    event%selected_channel = event%instance%get_channel ()
    call event%import_instance_results ()
    call event%accept_sqme_prc ()
    call event%update_normalization ()
    call event%accept_weight_prc ()
    call event%evaluate_transforms (r)
    if (signal_is_pending ())  return
    call event%check ()
  end subroutine event_generate
  
\eatline
\nwendcode{}\nwbegindocs{1520}Get a copy of the particle set belonging to the hard process.
\nwenddocs{}\nwbegincode{1521}\moddef{Events: event: TBP}\plusendmoddef
  procedure :: get_hard_particle_set => event_get_hard_particle_set
\nwendcode{}\nwbegincode{1522}\moddef{Events: procedures}\plusendmoddef
  subroutine event_get_hard_particle_set (event, pset)
    class(event_t), intent(in) :: event
    type(particle_set_t), intent(out) :: pset
    class(evt_t), pointer :: evt
    evt => event%transform_first
    pset = evt%particle_set
  end subroutine event_get_hard_particle_set
    
\eatline
\nwendcode{}\nwbegindocs{1523}\nwdocspar
\subsection{Recovering an event}
Select MC group, term, and integration channel.
\nwenddocs{}\nwbegincode{1524}\moddef{Events: event: TBP}\plusendmoddef
  procedure :: select => event_select
\nwendcode{}\nwbegincode{1525}\moddef{Events: procedures}\plusendmoddef
  subroutine event_select (event, i_mci, i_term, channel)
    class(event_t), intent(inout) :: event
    integer, intent(in) :: i_mci, i_term, channel
    if (associated (event%instance)) then
       event%selected_i_mci = i_mci
       event%selected_i_term = i_term
       event%selected_channel = channel
    else
       call msg_bug ("Event: select term: process instance undefined")
    end if
  end subroutine event_select

\eatline
\nwendcode{}\nwbegindocs{1526}\nwdocspar
Copy a particle set into the event record.

We deliberately use the first (the trivial) transform for this, i.e.,
the hard process.  The event reader may either read in the transformed
event separately, or apply all event
transforms to the hard particle set to (re)generate a fully dressed
event.

Since this makes all subsequent event transforms invalid, we call
{\tt{}reset} on them.
\nwenddocs{}\nwbegincode{1527}\moddef{Events: event: TBP}\plusendmoddef
  procedure :: set_hard_particle_set => event_set_hard_particle_set
\nwendcode{}\nwbegincode{1528}\moddef{Events: procedures}\plusendmoddef
  subroutine event_set_hard_particle_set (event, particle_set)
    class(event_t), intent(inout) :: event
    type(particle_set_t), intent(in) :: particle_set
    class(evt_t), pointer :: evt
    evt => event%transform_first
    call evt%set_particle_set (particle_set, &
         event%selected_i_mci, event%selected_i_term)
    call event%link_particle_set (evt%particle_set)
    evt => evt%next
    do while (associated (evt))
       call evt%reset ()
       evt => evt%next
    end do
  end subroutine event_set_hard_particle_set

\eatline
\nwendcode{}\nwbegindocs{1529}\nwdocspar
Set the $\alpha_s$ value that should be used in a recalculation.  This should
be called only if we explicitly want to override the QCD setting of the
process core.
\nwenddocs{}\nwbegincode{1530}\moddef{Events: event: TBP}\plusendmoddef
  procedure :: set_alpha_qcd_forced => event_set_alpha_qcd_forced
\nwendcode{}\nwbegincode{1531}\moddef{Events: procedures}\plusendmoddef
  subroutine event_set_alpha_qcd_forced (event, alpha_qcd)
    class(event_t), intent(inout) :: event
    real(default), intent(in) :: alpha_qcd
    if (allocated (event%alpha_qcd_forced)) then
       event%alpha_qcd_forced = alpha_qcd
    else
       allocate (event%alpha_qcd_forced, source = alpha_qcd)
    end if
  end subroutine event_set_alpha_qcd_forced
  
\eatline
\nwendcode{}\nwbegindocs{1532}\nwdocspar
Analogously, for the common scale.  This forces also renormalization and
factorization scale.
\nwenddocs{}\nwbegincode{1533}\moddef{Events: event: TBP}\plusendmoddef
  procedure :: set_scale_forced => event_set_scale_forced
\nwendcode{}\nwbegincode{1534}\moddef{Events: procedures}\plusendmoddef
  subroutine event_set_scale_forced (event, scale)
    class(event_t), intent(inout) :: event
    real(default), intent(in) :: scale
    if (allocated (event%scale_forced)) then
       event%scale_forced = scale
    else
       allocate (event%scale_forced, source = scale)
    end if
  end subroutine event_set_scale_forced
  
\eatline
\nwendcode{}\nwbegindocs{1535}\nwdocspar
Here we try to recover an event from the {\tt{}particle{\char95}set} subobject
and recalculate the structure functions and matrix elements.  We
have the appropriate {\tt{}process} object and an initialized
{\tt{}process{\char95}instance} at hand, so beam and configuration data are
known.  From the {\tt{}particle{\char95}set}, we get the momenta.

The quantum-number information may be incomplete, e.g., helicity
information may be partial or absent.  We recover the event just from
the momentum configuration.

We do not transfer the matrix element from the process instance to the
event record, as we do when generating an event.  The event record may
contain the matrix element as read from file, and the current
calculation may use different parameters.  We thus can compare old and
new values.

The event {\tt{}weight} may also be known already.  If yes, we pass it to the
{\tt{}evaluate{\char95}event{\char95}data} procedure.  It should already be normalized.  If we
have an {\tt{}weight{\char95}factor} value, we obtain the event weight by multiplying the
computed {\tt{}sqme} by this factor.  Otherwise, we make use of the MCI setup
(which should be valid then) to compute the event weight, and we should
normalize the result just as when generating events.

Evaluating event expressions must also be done separately.
\nwenddocs{}\nwbegincode{1536}\moddef{Events: event: TBP}\plusendmoddef
  procedure :: recalculate => event_recalculate
\nwendcode{}\nwbegincode{1537}\moddef{Events: procedures}\plusendmoddef
  subroutine event_recalculate &
       (event, update_sqme, weight_factor, recover_beams)
    class(event_t), intent(inout) :: event
    logical, intent(in) :: update_sqme
    real(default), intent(in), optional :: weight_factor
    logical, intent(in), optional :: recover_beams
    type(particle_set_t), pointer :: particle_set
    integer :: i_mci, i_term, channel
    if (event%has_valid_particle_set ()) then
       particle_set => event%get_particle_set_ptr ()
       i_mci = event%selected_i_mci
       i_term = event%selected_i_term
       channel = event%selected_channel
       if (i_mci == 0 .or. i_term == 0 .or. channel == 0) then
          call msg_bug ("Event: recalculate: undefined selection parameters")
       end if
       call event%instance%choose_mci (i_mci)
       call event%instance%set_trace (particle_set, i_term, recover_beams)
       if (allocated (event%alpha_qcd_forced)) then
          call event%instance%set_alpha_qcd_forced &
               (i_term, event%alpha_qcd_forced)
       end if
       call event%instance%recover (channel, i_term, update_sqme, &
            event%scale_forced) 
       if (signal_is_pending ())  return
       if (update_sqme .and. present (weight_factor)) then
          call event%instance%evaluate_event_data &
               (weight = event%instance%get_sqme () * weight_factor)
       else if (event%weight_ref_is_known ()) then
          call event%instance%evaluate_event_data &
               (weight = event%get_weight_ref ())
       else
          call event%process%recover_event (event%instance, i_term)
          if (signal_is_pending ())  return
          call event%instance%evaluate_event_data ()
          if (event%config%unweighted) then
             call event%instance%normalize_weight ()
          end if
       end if
       if (signal_is_pending ())  return
       if (update_sqme) then
          call event%import_instance_results ()
       else
          call event%accept_sqme_ref ()
          call event%accept_weight_ref ()
       end if
    else
       call msg_bug ("Event: can't recalculate, particle set is undefined")
    end if
  end subroutine event_recalculate
  
\eatline
\nwendcode{}\nwbegindocs{1538}\nwdocspar
\subsection{Access content}
Pointer to the associated process object (the associated model).
\nwenddocs{}\nwbegincode{1539}\moddef{Events: event: TBP}\plusendmoddef
  procedure :: get_process_ptr => event_get_process_ptr
  procedure :: get_process_instance_ptr => event_get_process_instance_ptr
  procedure :: get_model_ptr => event_get_model_ptr
\nwendcode{}\nwbegincode{1540}\moddef{Events: procedures}\plusendmoddef
  function event_get_process_ptr (event) result (ptr)
    class(event_t), intent(in) :: event
    type(process_t), pointer :: ptr
    ptr => event%process
  end function event_get_process_ptr

  function event_get_process_instance_ptr (event) result (ptr)
    class(event_t), intent(in) :: event
    type(process_instance_t), pointer :: ptr
    ptr => event%instance
  end function event_get_process_instance_ptr

  function event_get_model_ptr (event) result (model)
    class(event_t), intent(in) :: event
    class(model_data_t), pointer :: model
    model => event%process%get_model_ptr ()
  end function event_get_model_ptr

\eatline
\eatline
\eatline
\nwendcode{}\nwbegindocs{1541}Return the current values of indices: the MCI group of components, the term
index (different terms corresponding, potentially, to different effective
kinematics), and the MC integration channel.  The {\tt{}i{\char95}mci} call is delegated
to the current process instance.
\nwenddocs{}\nwbegincode{1542}\moddef{Events: event: TBP}\plusendmoddef
  procedure :: get_i_mci => event_get_i_mci
  procedure :: get_i_term => event_get_i_term
  procedure :: get_channel => event_get_channel
\nwendcode{}\nwbegincode{1543}\moddef{Events: procedures}\plusendmoddef
  function event_get_i_mci (event) result (i_mci)
    class(event_t), intent(in) :: event
    integer :: i_mci
    i_mci = event%selected_i_mci
  end function event_get_i_mci
  
  function event_get_i_term (event) result (i_term)
    class(event_t), intent(in) :: event
    integer :: i_term
    i_term = event%selected_i_term
  end function event_get_i_term
  
  function event_get_channel (event) result (channel)
    class(event_t), intent(in) :: event
    integer :: channel
    channel = event%selected_channel
  end function event_get_channel
  
\eatline
\eatline
\eatline
\nwendcode{}\nwbegindocs{1544}This flag tells us whether the event consists just of a hard process
(i.e., holds at most the first, trivial transform), or is a dressed
events with additional transforms.
\nwenddocs{}\nwbegincode{1545}\moddef{Events: event: TBP}\plusendmoddef
  procedure :: has_transform => event_has_transform
\nwendcode{}\nwbegincode{1546}\moddef{Events: procedures}\plusendmoddef
  function event_has_transform (event) result (flag)
    class(event_t), intent(in) :: event
    logical :: flag
    if (associated (event%transform_first)) then
       flag = associated (event%transform_first%next)
    else
       flag = .false.
    end if
  end function event_has_transform
  
\eatline
\nwendcode{}\nwbegindocs{1547}Return the currently selected normalization mode, or alternate
normalization mode.
\nwenddocs{}\nwbegincode{1548}\moddef{Events: event: TBP}\plusendmoddef
  procedure :: get_norm_mode => event_get_norm_mode
\nwendcode{}\nwbegincode{1549}\moddef{Events: procedures}\plusendmoddef
  elemental function event_get_norm_mode (event) result (norm_mode)
    class(event_t), intent(in) :: event
    integer :: norm_mode
    norm_mode = event%config%norm_mode
  end function event_get_norm_mode
  
\eatline
\nwendcode{}\nwbegindocs{1550}Return the kinematical weight, defined as the ratio of event weight
and squared matrix element.
\nwenddocs{}\nwbegincode{1551}\moddef{Events: event: TBP}\plusendmoddef
  procedure :: get_kinematical_weight => event_get_kinematical_weight
\nwendcode{}\nwbegincode{1552}\moddef{Events: procedures}\plusendmoddef
  function event_get_kinematical_weight (event) result (f)
    class(event_t), intent(in) :: event
    real(default) :: f
    if (event%sqme_ref_is_known () .and. event%weight_ref_is_known () &
         .and. event%get_sqme_ref () /= 0) then
       f = event%get_weight_ref () / event%get_sqme_ref ()
    else
       f = 0
    end if
  end function event_get_kinematical_weight
    
\eatline
\nwendcode{}\nwbegindocs{1553}Return data used by external event formats.
\nwenddocs{}\nwbegincode{1554}\moddef{Events: event: TBP}\plusendmoddef
  procedure :: get_index => event_get_index
  procedure :: get_fac_scale => event_get_fac_scale
  procedure :: get_alpha_s => event_get_alpha_s
\nwendcode{}\nwbegincode{1555}\moddef{Events: procedures}\plusendmoddef
  function event_get_index (event) result (index)
    class(event_t), intent(in) :: event
    integer :: index
    index = event%expr%index
  end function event_get_index
    
  function event_get_fac_scale (event) result (fac_scale)
    class(event_t), intent(in) :: event
    real(default) :: fac_scale
    fac_scale = event%instance%get_fac_scale (event%selected_i_term)
  end function event_get_fac_scale
    
  function event_get_alpha_s (event) result (alpha_s)
    class(event_t), intent(in) :: event
    real(default) :: alpha_s
    alpha_s = event%instance%get_alpha_s (event%selected_i_term)
  end function event_get_alpha_s
    
\eatline
\eatline
\eatline
\nwendcode{}\nwbegindocs{1556}\nwdocspar
Eliminate numerical noise in the {\tt{}subevt} expression and in the event
transforms (which includes associated process instances).
\nwenddocs{}\nwbegincode{1557}\moddef{Events: public}\plusendmoddef
  public :: pacify
\nwendcode{}\nwbegincode{1558}\moddef{Events: interfaces}\endmoddef
  interface pacify
     module procedure pacify_event
  end interface pacify
\nwendcode{}\nwbegincode{1559}\moddef{Events: procedures}\plusendmoddef
  subroutine pacify_event (event)
    class(event_t), intent(inout) :: event
    class(evt_t), pointer :: evt
    call event%pacify_particle_set ()
    if (event%expr%subevt_filled)  call pacify (event%expr)
    evt => event%transform_first
    do while (associated (evt))
       select type (evt)
       type is (evt_decay_t);  call pacify (evt)
       end select
       evt => evt%next
    end do
  end subroutine pacify_event
  
\eatline
\nwendcode{}\nwbegindocs{1560}\nwdocspar
\subsection{Unit tests}
\nwenddocs{}\nwbegincode{1561}\moddef{Events: public}\plusendmoddef
  public :: events_test
\nwendcode{}\nwbegincode{1562}\moddef{Events: tests}\endmoddef
  subroutine events_test (u, results)
    integer, intent(in) :: u
    type(test_results_t), intent(inout) :: results
  \LA{}Events: execute tests\RA{}
  end subroutine events_test
  
\eatline
\nwendcode{}\nwbegindocs{1563}\nwdocspar
\subsubsection{Empty event record}
\nwenddocs{}\nwbegincode{1564}\moddef{Events: execute tests}\endmoddef
  call test (events_1, "events_1", &
       "empty event record", &
       u, results)
\nwendcode{}\nwbegincode{1565}\moddef{Events: tests}\plusendmoddef
  subroutine events_1 (u)
    integer, intent(in) :: u
    type(event_t), target :: event

    write (u, "(A)")  "* Test output: events_1"
    write (u, "(A)")  "*   Purpose: display an empty event object"
    write (u, "(A)")

    call event%write (u)
    
    write (u, "(A)")
    write (u, "(A)")  "* Test output end: events_1"
    
  end subroutine events_1
  
\eatline
\nwendcode{}\nwbegindocs{1566}\nwdocspar
\subsubsection{Simple event}
\nwenddocs{}\nwbegincode{1567}\moddef{Events: execute tests}\plusendmoddef
  call test (events_2, "events_2", &
       "generate event", &
       u, results)
\nwendcode{}\nwbegincode{1568}\moddef{Events: tests}\plusendmoddef
  subroutine events_2 (u)
    integer, intent(in) :: u
    type(event_t), allocatable, target :: event
    type(process_t), allocatable, target :: process
    type(process_instance_t), allocatable, target :: process_instance
    type(model_data_t), target :: model

    write (u, "(A)")  "* Test output: events_2"
    write (u, "(A)")  "*   Purpose: generate and display an event"
    write (u, "(A)")

    call model%init_test ()

    write (u, "(A)")  "* Generate test process event"

    allocate (process)
    allocate (process_instance)
    call prepare_test_process (process, process_instance, model)
    call process_instance%setup_event_data ()

    write (u, "(A)")
    write (u, "(A)")  "* Initialize event object"

    allocate (event)
    call event%basic_init ()
    call event%connect (process_instance, process%get_model_ptr ())
    
    write (u, "(A)")
    write (u, "(A)")  "* Generate test process event"

    call process%generate_weighted_event (process_instance, 1)

    write (u, "(A)")
    write (u, "(A)")  "* Fill event object"
    write (u, "(A)")

    call event%generate (1, [0.4_default, 0.4_default])
    call event%evaluate_expressions ()
    call event%write (u)

    write (u, "(A)")
    write (u, "(A)")  "* Cleanup"

    call event%final ()
    deallocate (event)

    call cleanup_test_process (process, process_instance)
    deallocate (process_instance)
    deallocate (process)
    
!    call model_list%final ()
!    call syntax_model_file_final ()
    call model%final ()
    
    write (u, "(A)")
    write (u, "(A)")  "* Test output end: events_2"
    
  end subroutine events_2
  
\eatline
\nwendcode{}\nwbegindocs{1569}\nwdocspar
\subsubsection{Recovering an event}
Generate an event and store the particle set.  Then reset the event
record, recall the particle set, and recover the event from that.

Note: The extra {\tt{}particle{\char95}set{\char95}ptr} auxiliary is a workaround for
memory corruption in gfortran 4.7.
\nwenddocs{}\nwbegincode{1570}\moddef{Events: execute tests}\plusendmoddef
  call test (events_4, "events_4", &
       "recover event", &
       u, results)
\nwendcode{}\nwbegincode{1571}\moddef{Events: tests}\plusendmoddef
  subroutine events_4 (u)
    integer, intent(in) :: u
    type(event_t), allocatable, target :: event
    type(process_t), allocatable, target :: process
    type(process_instance_t), allocatable, target :: process_instance
    type(particle_set_t) :: particle_set
    type(particle_set_t), pointer :: particle_set_ptr
    type(model_data_t), target :: model

    write (u, "(A)")  "* Test output: events_4"
    write (u, "(A)")  "*   Purpose: generate and recover an event"
    write (u, "(A)")

    call model%init_test ()

    write (u, "(A)")  "* Generate test process event and save particle set"
    write (u, "(A)")

    allocate (process)
    allocate (process_instance)
    call prepare_test_process (process, process_instance, model)
    call process_instance%setup_event_data ()

    allocate (event)
    call event%basic_init ()
    call event%connect (process_instance, process%get_model_ptr ())
    
    call event%generate (1, [0.4_default, 0.4_default])
    call event%evaluate_expressions ()
    call event%write (u)
    
    particle_set_ptr => event%get_particle_set_ptr ()
    particle_set = particle_set_ptr

    call event%final ()
    deallocate (event)

    call cleanup_test_process (process, process_instance)
    deallocate (process_instance)
    deallocate (process)

    write (u, "(A)")
    write (u, "(A)")  "* Recover event from particle set"
    write (u, "(A)")
    
    allocate (process)
    allocate (process_instance)
    call prepare_test_process (process, process_instance, model)
    call process_instance%setup_event_data ()

    allocate (event)
    call event%basic_init ()
    call event%connect (process_instance, process%get_model_ptr ())
    
    call event%select (1, 1, 1)
    call event%set_hard_particle_set (particle_set)
    call event%recalculate (update_sqme = .true.)
    call event%write (u)

    write (u, "(A)")
    write (u, "(A)")  "* Transfer sqme and evaluate expressions"
    write (u, "(A)")
    
    call event%accept_sqme_prc ()
    call event%accept_weight_prc ()
    call event%check ()
    call event%evaluate_expressions ()
    call event%write (u)
    
    write (u, "(A)")
    write (u, "(A)")  "* Reset contents"
    write (u, "(A)")

    call event%reset ()
    event%transform_first%particle_set_exists = .false.
    call event%write (u)

    write (u, "(A)")
    write (u, "(A)")  "* Cleanup"

    call particle_set%final ()

    call event%final ()
    deallocate (event)

    call cleanup_test_process (process, process_instance)
    deallocate (process_instance)
    deallocate (process)
    
!    call model_list%final ()
!    call syntax_model_file_final ()
    call model%final ()

    write (u, "(A)")
    write (u, "(A)")  "* Test output end: events_4"
    
  end subroutine events_4
  
\eatline
\nwendcode{}\nwbegindocs{1572}\nwdocspar
\subsubsection{Partially Recovering an event}
Generate an event and store the particle set.  Then reset the event
record, recall the particle set, and recover the event as far as possible
without recomputing the squared matrix element.
\nwenddocs{}\nwbegincode{1573}\moddef{Events: execute tests}\plusendmoddef
  call test (events_5, "events_5", &
       "partially recover event", &
       u, results)
\nwendcode{}\nwbegincode{1574}\moddef{Events: tests}\plusendmoddef
  subroutine events_5 (u)
    integer, intent(in) :: u
    type(event_t), allocatable, target :: event
    type(process_t), allocatable, target :: process
    type(process_instance_t), allocatable, target :: process_instance
    type(particle_set_t) :: particle_set
    type(particle_set_t), pointer :: particle_set_ptr
    real(default) :: sqme, weight
    type(model_data_t), target :: model

    write (u, "(A)")  "* Test output: events_5"
    write (u, "(A)")  "*   Purpose: generate and recover an event"
    write (u, "(A)")

    call model%init_test ()

    write (u, "(A)")  "* Generate test process event and save particle set"
    write (u, "(A)")

    allocate (process)
    allocate (process_instance)
    call prepare_test_process (process, process_instance, model)
    call process_instance%setup_event_data ()

    allocate (event)
    call event%basic_init ()
    call event%connect (process_instance, process%get_model_ptr ())
    
    call event%generate (1, [0.4_default, 0.4_default])
    call event%evaluate_expressions ()
    call event%write (u)
    
    particle_set_ptr => event%get_particle_set_ptr ()
    particle_set = particle_set_ptr
    sqme = event%get_sqme_ref ()
    weight = event%get_weight_ref ()

    call event%final ()
    deallocate (event)

    call cleanup_test_process (process, process_instance)
    deallocate (process_instance)
    deallocate (process)

    write (u, "(A)")
    write (u, "(A)")  "* Recover event from particle set"
    write (u, "(A)")
    
    allocate (process)
    allocate (process_instance)
    call prepare_test_process (process, process_instance, model)
    call process_instance%setup_event_data ()

    allocate (event)
    call event%basic_init ()
    call event%connect (process_instance, process%get_model_ptr ())
    
    call event%select (1, 1, 1)
    call event%set_hard_particle_set (particle_set)
    call event%recalculate (update_sqme = .false.)
    call event%write (u)

    write (u, "(A)")
    write (u, "(A)")  "* Manually set sqme and evaluate expressions"
    write (u, "(A)")
    
    call event%set (sqme_ref = sqme, weight_ref = weight)
    call event%accept_sqme_ref ()
    call event%accept_weight_ref ()
    call event%evaluate_expressions ()
    call event%write (u)
    
    write (u, "(A)")
    write (u, "(A)")  "* Cleanup"

    call particle_set%final ()

    call event%final ()
    deallocate (event)

    call cleanup_test_process (process, process_instance)
    deallocate (process_instance)
    deallocate (process)
    
!    call model_list%final ()
!    call syntax_model_file_final ()
    call model%final ()

    write (u, "(A)")
    write (u, "(A)")  "* Test output end: events_5"
    
  end subroutine events_5
  
\eatline
\nwendcode{}\nwbegindocs{1575}\nwdocspar
\subsubsection{Decays}
Generate an event with subsequent decays.
\nwenddocs{}\nwbegincode{1576}\moddef{Events: execute tests}\plusendmoddef
  call test (events_6, "events_6", &
       "decays", &
       u, results)
\nwendcode{}\nwbegincode{1577}\moddef{Events: tests}\plusendmoddef
  subroutine events_6 (u)
    integer, intent(in) :: u
    type(os_data_t) :: os_data
    class(model_data_t), pointer :: model
    type(string_t) :: prefix, procname1, procname2
    type(process_library_t), target :: lib
    type(process_stack_t) :: process_stack
    class(evt_t), pointer :: evt_decay
    type(event_t), allocatable, target :: event
    type(process_t), pointer :: process
    type(process_instance_t), allocatable, target :: process_instance

    write (u, "(A)")  "* Test output: events_6"
    write (u, "(A)")  "*   Purpose: generate an event with subsequent decays"
    write (u, "(A)")

    write (u, "(A)")  "* Generate test process and decay"
    write (u, "(A)")

    call os_data_init (os_data)

    prefix = "events_6"
    procname1 = prefix // "_p"
    procname2 = prefix // "_d"
    call prepare_testbed &
         (lib, process_stack, prefix, os_data, &
         scattering=.true., decay=.true.)

    write (u, "(A)")  "* Initialize decay process"

    process => process_stack%get_process_ptr (procname1)
    model => process%get_model_ptr ()
    call model%set_unstable (25, [procname2])

    allocate (process_instance)
    call process_instance%init (process)
    call process_instance%setup_event_data ()
    call process_instance%init_simulation (1)

    write (u, "(A)")
    write (u, "(A)")  "* Initialize event transform: decay"

    allocate (evt_decay_t :: evt_decay)
    call evt_decay%connect (process_instance, model, process_stack)
    
    write (u, "(A)")
    write (u, "(A)")  "* Initialize event object"
    write (u, "(A)")

    allocate (event)
    call event%basic_init ()
    call event%connect (process_instance, model)
    call event%import_transform (evt_decay)
    
    call event%write (u, show_decay = .true.)
    
    write (u, "(A)")
    write (u, "(A)")  "* Generate event"
    write (u, "(A)")

    call event%generate (1, [0.4_default, 0.4_default])
    call event%evaluate_expressions ()
    call event%write (u)

    write (u, "(A)")
    write (u, "(A)")  "* Cleanup"

    call event%final ()
    deallocate (event)

    call process_instance%final ()
    deallocate (process_instance)

    call process_stack%final ()
    
    write (u, "(A)")
    write (u, "(A)")  "* Test output end: events_6"
    
  end subroutine events_6
  
\eatline
\nwendcode{}\nwbegindocs{1578}\nwdocspar
\subsubsection{Decays}
Generate a decay event with varying options.
\nwenddocs{}\nwbegincode{1579}\moddef{Events: execute tests}\plusendmoddef
  call test (events_7, "events_7", &
       "decay options", &
       u, results)
\nwendcode{}\nwbegincode{1580}\moddef{Events: tests}\plusendmoddef
  subroutine events_7 (u)
    integer, intent(in) :: u
    type(os_data_t) :: os_data
    class(model_data_t), pointer :: model
    type(string_t) :: prefix, procname2
    type(process_library_t), target :: lib
    type(process_stack_t) :: process_stack
    type(process_t), pointer :: process
    type(process_instance_t), allocatable, target :: process_instance

    write (u, "(A)")  "* Test output: events_7"
    write (u, "(A)")  "*   Purpose: check decay options"
    write (u, "(A)")

    write (u, "(A)")  "* Prepare test process"
    write (u, "(A)")

    call os_data_init (os_data)

    prefix = "events_7"
    procname2 = prefix // "_d"
    call prepare_testbed &
         (lib, process_stack, prefix, os_data, &
         scattering=.false., decay=.true.)

    write (u, "(A)")  "* Generate decay event, default options"
    write (u, "(A)")

    process => process_stack%get_process_ptr (procname2)
    model => process%get_model_ptr ()
    call model%set_unstable (25, [procname2])

    allocate (process_instance)
    call process_instance%init (process)
    call process_instance%setup_event_data (model)
    call process_instance%init_simulation (1)

    call process%generate_weighted_event (process_instance, 1)
    call process_instance%write (u)

    call process_instance%final ()
    deallocate (process_instance)

    write (u, "(A)")
    write (u, "(A)")  "* Generate decay event, helicity-diagonal decay"
    write (u, "(A)")

    process => process_stack%get_process_ptr (procname2)
    model => process%get_model_ptr ()
    call model%set_unstable (25, [procname2], diagonal = .true.)

    allocate (process_instance)
    call process_instance%init (process)
    call process_instance%setup_event_data (model)
    call process_instance%init_simulation (1)

    call process%generate_weighted_event (process_instance, 1)
    call process_instance%write (u)

    call process_instance%final ()
    deallocate (process_instance)

    write (u, "(A)")
    write (u, "(A)")  "* Generate decay event, isotropic decay, &
         &polarized final state"
    write (u, "(A)")

    process => process_stack%get_process_ptr (procname2)
    model => process%get_model_ptr ()
    call model%set_unstable (25, [procname2], isotropic = .true.)
    call model%set_polarized (6)
    call model%set_polarized (-6)

    allocate (process_instance)
    call process_instance%init (process)
    call process_instance%setup_event_data (model)
    call process_instance%init_simulation (1)

    call process%generate_weighted_event (process_instance, 1)
    call process_instance%write (u)

    call process_instance%final ()
    deallocate (process_instance)

    write (u, "(A)")
    write (u, "(A)")  "* Cleanup"
    
    call process_stack%final ()
    
    write (u, "(A)")
    write (u, "(A)")  "* Test output end: events_7"
    
  end subroutine events_7
  
\eatline
\nwendcode{}\nwbegindocs{1581}\nwdocspar
%%%%%%%%%%%%%%%%%%%%%%%%%%%%%%%%%%%%%%%%%%%%%%%%%%%%%%%%%%%%%%%%%%%%%%%% 
\section{Raw Event I/O}
The raw format is for internal use only.  All data are stored
unformatted, so they can be efficiently be re-read on the same
machine, but not necessarily on another machine.

This module explicitly depends on the {\tt{}events} module which provides the
concrete implementation of {\tt{}event{\char95}base}.  The other I/O formats
access only the methods that are defined in {\tt{}event{\char95}base}.
\nwenddocs{}\nwbegincode{1582}\moddef{\code{}eio{\char95}raw.f90\edoc{}}\endmoddef
\LA{}File header\RA{}

module eio_raw
  
  use kinds
  use io_units
\LA{}Use strings\RA{}
  use unit_tests
  use diagnostics

  use lorentz
  use variables
  use model_data
  use particles
  use event_base
  use eio_data
  use eio_base
  use events

\LA{}Standard module head\RA{}

\LA{}EIO raw: public\RA{}

\LA{}EIO raw: parameters\RA{}

\LA{}EIO raw: types\RA{}

contains
  
\LA{}EIO raw: procedures\RA{}

\LA{}EIO raw: tests\RA{}

end module eio_raw
\eatline
\nwendcode{}\nwbegindocs{1583}\nwdocspar
\subsection{File Format Version}
This is the current default file version. 
\nwenddocs{}\nwbegincode{1584}\moddef{EIO raw: parameters}\endmoddef
  integer, parameter :: CURRENT_FILE_VERSION = 2

\eatline
\nwendcode{}\nwbegindocs{1585}The user may change this number; this should force some
compatibility mode for reading and writing.  In any case, the file
version stored in a event file that we read has to match the expected
file version.

History of version numbers:
\begin{enumerate}
\item
  Format for WHIZARD 2.2.0 to 2.2.3.  No version number stored in the raw file.
\item
  Format from 2.2.4 on.  File contains version number.  The file
  contains the transformed particle set (if applicable) after the
  hard-process particle set.
\end{enumerate}
\nwenddocs{}\nwbegindocs{1586}\nwdocspar
\subsection{Type}
Note the file version number.  The default may be reset during
initialization, which should enforce some compatibility mode.
\nwenddocs{}\nwbegincode{1587}\moddef{EIO raw: public}\endmoddef
  public :: eio_raw_t
\nwendcode{}\nwbegincode{1588}\moddef{EIO raw: types}\endmoddef
  type, extends (eio_t) :: eio_raw_t
     logical :: reading = .false.
     logical :: writing = .false.
     integer :: unit = 0
     integer :: norm_mode = NORM_UNDEFINED
     real(default) :: sigma = 1
     integer :: n = 1
     integer :: n_alt = 0
     logical :: check = .false.
     integer :: file_version = CURRENT_FILE_VERSION
   contains
   \LA{}EIO raw: eio raw: TBP\RA{}
  end type eio_raw_t
  
\eatline
\nwendcode{}\nwbegindocs{1589}Output.  This is not the actual event format, but a readable account
of the current object status.
\nwenddocs{}\nwbegincode{1590}\moddef{EIO raw: eio raw: TBP}\endmoddef
  procedure :: write => eio_raw_write
\nwendcode{}\nwbegincode{1591}\moddef{EIO raw: procedures}\endmoddef
  subroutine eio_raw_write (object, unit)
    class(eio_raw_t), intent(in) :: object
    integer, intent(in), optional :: unit
    integer :: u
    u = given_output_unit (unit)
    write (u, "(1x,A)")  "Raw event stream:"
    write (u, "(3x,A,L1)")  "Check MD5 sum     = ", object%check
    if (object%n_alt > 0) then
       write (u, "(3x,A,I0)")  "Alternate weights = ", object%n_alt
    end if
    if (object%reading) then
       write (u, "(3x,A,A)")  "Reading from file = ", char (object%filename)
    else if (object%writing) then
       write (u, "(3x,A,A)")  "Writing to file   = ", char (object%filename)
    else
       write (u, "(3x,A)")  "[closed]"
    end if
  end subroutine eio_raw_write
  
\eatline
\nwendcode{}\nwbegindocs{1592}Finalizer: close any open file.
\nwenddocs{}\nwbegincode{1593}\moddef{EIO raw: eio raw: TBP}\plusendmoddef
  procedure :: final => eio_raw_final
\nwendcode{}\nwbegincode{1594}\moddef{EIO raw: procedures}\plusendmoddef
  subroutine eio_raw_final (object)
    class(eio_raw_t), intent(inout) :: object
    if (object%reading .or. object%writing) then
       write (msg_buffer, "(A,A,A)")  "Events: closing raw file '", &
            char (object%filename), "'"
       call msg_message ()
       close (object%unit)
       object%reading = .false.
       object%writing = .false.
    end if
  end subroutine eio_raw_final
  
\eatline
\nwendcode{}\nwbegindocs{1595}Set the {\tt{}check} flag which determines whether we compare checksums on input.
\nwenddocs{}\nwbegincode{1596}\moddef{EIO raw: eio raw: TBP}\plusendmoddef
  procedure :: set_parameters => eio_raw_set_parameters
\nwendcode{}\nwbegincode{1597}\moddef{EIO raw: procedures}\plusendmoddef
  subroutine eio_raw_set_parameters (eio, check, version_string, extension)
    class(eio_raw_t), intent(inout) :: eio
    logical, intent(in), optional :: check
    type(string_t), intent(in), optional :: version_string
    type(string_t), intent(in), optional :: extension 
    if (present (check))  eio%check = check
    if (present (version_string)) then
       select case (char (version_string))
       case ("", "2.2.4")
          eio%file_version = CURRENT_FILE_VERSION
       case ("2.2")
          eio%file_version = 1
       case default
          call msg_fatal ("Raw event I/O: unsupported version '" &
               // char (version_string) // "'")
          eio%file_version = 0
       end select
    end if
    if (present (extension)) then
       eio%extension = extension
    else
       eio%extension = "evx"
    end if
  end subroutine eio_raw_set_parameters
    
\eatline
\nwendcode{}\nwbegindocs{1598}Initialize event writing.
\nwenddocs{}\nwbegincode{1599}\moddef{EIO raw: eio raw: TBP}\plusendmoddef
  procedure :: init_out => eio_raw_init_out
\nwendcode{}\nwbegincode{1600}\moddef{EIO raw: procedures}\plusendmoddef
  subroutine eio_raw_init_out (eio, sample, data, success, extension)
    class(eio_raw_t), intent(inout) :: eio
    type(string_t), intent(in) :: sample
    type(event_sample_data_t), intent(in), optional :: data
    logical, intent(out), optional :: success
    type(string_t), intent(in), optional :: extension
    character(32) :: md5sum_prc, md5sum_cfg
    character(32), dimension(:), allocatable :: md5sum_alt
    integer :: i
    if (present (extension)) then
       eio%extension  = extension
    else
       eio%extension = "evx"
    end if
    eio%filename = sample // "." // eio%extension
    eio%unit = free_unit ()
    write (msg_buffer, "(A,A,A)")  "Events: writing to raw file '", &
         char (eio%filename), "'"
    call msg_message ()
    eio%writing = .true.
    if (present (data)) then
       md5sum_prc = data%md5sum_prc
       md5sum_cfg = data%md5sum_cfg
       eio%norm_mode = data%norm_mode
       eio%sigma = data%total_cross_section
       eio%n = data%n_evt
       eio%n_alt = data%n_alt
       if (eio%n_alt > 0) then
          allocate (md5sum_alt (data%n_alt), source = data%md5sum_alt)
       end if
    else
       md5sum_prc = ""
       md5sum_cfg = ""
    end if
    open (eio%unit, file = char (eio%filename), form = "unformatted", &
         action = "write", status = "replace")
    select case (eio%file_version)
    case (2:);  write (eio%unit)  eio%file_version
    end select
    write (eio%unit)  md5sum_prc
    write (eio%unit)  md5sum_cfg
    write (eio%unit)  eio%norm_mode
    write (eio%unit)  eio%n_alt
    do i = 1, eio%n_alt
       write (eio%unit)  md5sum_alt(i)
    end do
    if (present (success))  success = .true.
  end subroutine eio_raw_init_out
    
\eatline
\nwendcode{}\nwbegindocs{1601}Initialize event reading.
\nwenddocs{}\nwbegincode{1602}\moddef{EIO raw: eio raw: TBP}\plusendmoddef
  procedure :: init_in => eio_raw_init_in
\nwendcode{}\nwbegincode{1603}\moddef{EIO raw: procedures}\plusendmoddef
  subroutine eio_raw_init_in (eio, sample, data, success, extension)
    class(eio_raw_t), intent(inout) :: eio
    type(string_t), intent(in) :: sample
    type(event_sample_data_t), intent(inout), optional :: data
    logical, intent(out), optional :: success
    type(string_t), intent(in), optional :: extension
    character(32) :: md5sum_prc, md5sum_cfg
    character(32), dimension(:), allocatable :: md5sum_alt
    integer :: i, file_version
    if (present (success))  success = .true.
    if (present (extension)) then
       eio%extension = extension
    else
       eio%extension = "evx"
    end if
    eio%filename = sample // "." // eio%extension
    eio%unit = free_unit ()
    if (present (data)) then
       eio%sigma = data%total_cross_section
       eio%n = data%n_evt
    end if
    write (msg_buffer, "(A,A,A)")  "Events: reading from raw file '", &
         char (eio%filename), "'"
    call msg_message ()
    eio%reading = .true.
    open (eio%unit, file = char (eio%filename), form = "unformatted", &
         action = "read", status = "old")
    select case (eio%file_version)
    case (2:);  read (eio%unit)  file_version
    case default;  file_version = 1
    end select
    if (file_version /= eio%file_version) then
       call msg_error ("Reading event file: raw-file version mismatch.")
       if (present (success))  success = .false.
       return
    else if (file_version /= CURRENT_FILE_VERSION) then
       call msg_warning ("Reading event file: compatibility mode.")
    end if
    read (eio%unit)  md5sum_prc
    read (eio%unit)  md5sum_cfg
    read (eio%unit)  eio%norm_mode
    read (eio%unit)  eio%n_alt
    if (present (data)) then
       if (eio%n_alt /= data%n_alt) then
          if (present (success))  success = .false. !
          return
       end if
    end if
    allocate (md5sum_alt (eio%n_alt))
    do i = 1, eio%n_alt
       read (eio%unit)  md5sum_alt(i)
    end do
    if (present (success)) then
       if (present (data)) then
          if (eio%check) then
             if (data%md5sum_prc /= "") then
                success = success .and. md5sum_prc == data%md5sum_prc
             end if
             if (data%md5sum_cfg /= "") then
                success = success .and. md5sum_cfg == data%md5sum_cfg
             end if
             do i = 1, eio%n_alt
                if (data%md5sum_alt(i) /= "") then
                   success = success .and. md5sum_alt(i) == data%md5sum_alt(i)
                end if
             end do
          else
             call msg_warning ("Reading event file: MD5 sum check disabled")
          end if
       end if
    end if
  end subroutine eio_raw_init_in
    
\eatline
\nwendcode{}\nwbegindocs{1604}Switch from input to output: reopen the file for reading.
\nwenddocs{}\nwbegincode{1605}\moddef{EIO raw: eio raw: TBP}\plusendmoddef
  procedure :: switch_inout => eio_raw_switch_inout
\nwendcode{}\nwbegincode{1606}\moddef{EIO raw: procedures}\plusendmoddef
  subroutine eio_raw_switch_inout (eio, success)
    class(eio_raw_t), intent(inout) :: eio
    logical, intent(out), optional :: success
    write (msg_buffer, "(A,A,A)")  "Events: appending to raw file '", &
         char (eio%filename), "'"
    call msg_message ()
    close (eio%unit, status = "keep")
    eio%reading = .false.
    open (eio%unit, file = char (eio%filename), form = "unformatted", &
         action = "write", position = "append", status = "old")
    eio%writing = .true.
    if (present (success))  success = .true.
  end subroutine eio_raw_switch_inout
  
\eatline
\nwendcode{}\nwbegindocs{1607}Output an event.  Write first the event indices, then weight and
squared matrix element, then the particle set.

We always write the particle set of the hard process.  (Note: this
should be reconsidered.)  We do make a physical copy.

On output, we write the {\tt{}prc} values for weight and sqme, since
these are the values just computed.  On input, we store the values as
{\tt{}ref} values.  The caller can then decide whether to recompute
values and thus obtain distinct {\tt{}prc} values, or just accept them.

The {\tt{}passed} flag is not written.  This allow us to apply different
selection criteria upon rereading.
\nwenddocs{}\nwbegincode{1608}\moddef{EIO raw: eio raw: TBP}\plusendmoddef
  procedure :: output => eio_raw_output
\nwendcode{}\nwbegincode{1609}\moddef{EIO raw: procedures}\plusendmoddef
  subroutine eio_raw_output (eio, event, i_prc, reading, passed, pacify)
    class(eio_raw_t), intent(inout) :: eio
    class(generic_event_t), intent(in), target :: event
    logical, intent(in), optional :: reading, passed, pacify
    integer, intent(in) :: i_prc
    type(particle_set_t), pointer :: pset
    integer :: i
    if (eio%writing) then
       if (event%has_valid_particle_set ()) then
          select type (event)
          type is (event_t)
             write (eio%unit)  i_prc
             write (eio%unit)  event%get_i_mci ()
             write (eio%unit)  event%get_i_term ()
             write (eio%unit)  event%get_channel ()
             write (eio%unit)  event%expr%weight_prc
             write (eio%unit)  event%expr%excess_prc
             write (eio%unit)  event%expr%sqme_prc
             do i = 1, eio%n_alt
                write (eio%unit)  event%expr%weight_alt(i)
                write (eio%unit)  event%expr%sqme_alt(i)
             end do
             allocate (pset)
             call event%get_hard_particle_set (pset)
             call pset%write_raw (eio%unit)
             call pset%final ()
             deallocate (pset)
             select case (eio%file_version)
             case (2:)
                if (event%has_transform ()) then
                   write (eio%unit)  .true.
                   pset => event%get_particle_set_ptr ()
                   call pset%write_raw (eio%unit)
                else
                   write (eio%unit)  .false.
                end if
             end select
          class default
             call msg_bug ("Event: write raw: defined only for full event_t")
          end select
       else
          call msg_bug ("Event: write raw: particle set is undefined")
       end if
    else
       call eio%write ()
       call msg_fatal ("Raw event file is not open for writing")
    end if
  end subroutine eio_raw_output

\eatline
\nwendcode{}\nwbegindocs{1610}Input an event.

Note: the particle set is physically copied.  If there is a
performance issue, we might choose to pointer-assign it instead, with
a different version of {\tt{}event{\char37}set{\char95}hard{\char95}particle{\char95}set}.
\nwenddocs{}\nwbegincode{1611}\moddef{EIO raw: eio raw: TBP}\plusendmoddef
  procedure :: input_i_prc => eio_raw_input_i_prc
  procedure :: input_event => eio_raw_input_event
\nwendcode{}\nwbegincode{1612}\moddef{EIO raw: procedures}\plusendmoddef
  subroutine eio_raw_input_i_prc (eio, i_prc, iostat)
    class(eio_raw_t), intent(inout) :: eio
    integer, intent(out) :: i_prc
    integer, intent(out) :: iostat
    if (eio%reading) then
       read (eio%unit, iostat = iostat)  i_prc
    else
       call eio%write ()
       call msg_fatal ("Raw event file is not open for reading")
    end if
  end subroutine eio_raw_input_i_prc

  subroutine eio_raw_input_event (eio, event, iostat)
    class(eio_raw_t), intent(inout) :: eio
    class(generic_event_t), intent(inout), target :: event
    integer, intent(out) :: iostat
    integer :: i_mci, i_term, channel, i
    real(default) :: weight, excess, sqme
    real(default), dimension(:), allocatable :: weight_alt, sqme_alt
    logical :: has_transform
    type(particle_set_t), pointer :: pset
    class(model_data_t), pointer :: model
    if (eio%reading) then
       select type (event)
       type is (event_t)
          read (eio%unit, iostat = iostat)  i_mci
          if (iostat /= 0)  return
          read (eio%unit, iostat = iostat)  i_term
          if (iostat /= 0)  return
          read (eio%unit, iostat = iostat)  channel
          if (iostat /= 0)  return
          read (eio%unit, iostat = iostat)  weight
          if (iostat /= 0)  return
          read (eio%unit, iostat = iostat)  excess
          if (iostat /= 0)  return
          read (eio%unit, iostat = iostat)  sqme
          if (iostat /= 0)  return
          call event%reset ()
          call event%select (i_mci, i_term, channel)
          if (eio%norm_mode /= NORM_UNDEFINED) then
             call event_normalization_update (weight, &
                  eio%sigma, eio%n, event%get_norm_mode (), eio%norm_mode)
             call event_normalization_update (excess, &
                  eio%sigma, eio%n, event%get_norm_mode (), eio%norm_mode)
          end if
          call event%set (sqme_ref = sqme, weight_ref = weight, &
               excess_prc = excess)
          if (eio%n_alt /= 0) then
             allocate (sqme_alt (eio%n_alt), weight_alt (eio%n_alt))
             do i = 1, eio%n_alt
                read (eio%unit, iostat = iostat)  weight_alt(i)
                if (iostat /= 0)  return
                read (eio%unit, iostat = iostat)  sqme_alt(i)
                if (iostat /= 0)  return
             end do
             call event%set (sqme_alt = sqme_alt, weight_alt = weight_alt)
          end if
          model => null ()
          if (associated (event%process)) then
             model => event%process%get_model_ptr ()
          end if
          allocate (pset)
          call pset%read_raw (eio%unit, iostat)
          if (iostat /= 0)  return
          if (associated (model))  call pset%set_model (model)
          call event%set_hard_particle_set (pset)
          call pset%final ()
          deallocate (pset)
          select case (eio%file_version)
          case (2:)
             read (eio%unit, iostat = iostat)  has_transform
             if (iostat /= 0)  return
             if (has_transform) then
                allocate (pset)
                call pset%read_raw (eio%unit, iostat)
                if (iostat /= 0)  return
                if (associated (model)) &
                     call pset%set_model (model)
                call event%link_particle_set (pset)
             end if
          end select
       class default
          call msg_bug ("Event: read raw: defined only for full event_t")
       end select
    else
       call eio%write ()
       call msg_fatal ("Raw event file is not open for reading")
    end if
  end subroutine eio_raw_input_event

\eatline
\eatline
\nwendcode{}\nwbegindocs{1613}\nwdocspar
\subsection{Unit tests}
\nwenddocs{}\nwbegincode{1614}\moddef{EIO raw: public}\plusendmoddef
  public :: eio_raw_test
\nwendcode{}\nwbegincode{1615}\moddef{EIO raw: tests}\endmoddef
  subroutine eio_raw_test (u, results)
    integer, intent(in) :: u
    type(test_results_t), intent(inout) :: results
  \LA{}EIO raw: execute tests\RA{}
  end subroutine eio_raw_test
  
\eatline
\nwendcode{}\nwbegindocs{1616}\nwdocspar
\subsubsection{Test I/O methods}
We test the implementation of all I/O methods.
\nwenddocs{}\nwbegincode{1617}\moddef{EIO raw: execute tests}\endmoddef
  call test (eio_raw_1, "eio_raw_1", &
       "read and write event contents", &
       u, results)
\nwendcode{}\nwbegincode{1618}\moddef{EIO raw: tests}\plusendmoddef
  subroutine eio_raw_1 (u)
    use processes
    integer, intent(in) :: u
    type(model_data_t), target :: model
    type(event_t), allocatable, target :: event
    type(process_t), allocatable, target :: process
    type(process_instance_t), allocatable, target :: process_instance
    class(eio_t), allocatable :: eio
    integer :: i_prc, iostat
    type(string_t) :: sample

    write (u, "(A)")  "* Test output: eio_raw_1"
    write (u, "(A)")  "*   Purpose: generate and read/write an event"
    write (u, "(A)")

    write (u, "(A)")  "* Initialize test process"
 
    call model%init_test ()

    allocate (process)
    allocate (process_instance)
    call prepare_test_process (process, process_instance, model)
    call process_instance%setup_event_data ()
 
    allocate (event)
    call event%basic_init ()
    call event%connect (process_instance, process%get_model_ptr ())
    
    write (u, "(A)")
    write (u, "(A)")  "* Generate and write an event"
    write (u, "(A)")
 
    sample = "eio_raw_1"
 
    allocate (eio_raw_t :: eio)
    
    call eio%init_out (sample)
    call event%generate (1, [0._default, 0._default])
    call event%evaluate_expressions ()
    call event%write (u)
    write (u, "(A)")

    call eio%output (event, i_prc = 42)
    call eio%write (u)
    call eio%final ()

    call event%final ()
    deallocate (event)
    call process_instance%final ()
    deallocate (process_instance)
    
    write (u, "(A)")
    write (u, "(A)")  "* Re-read the event"
    write (u, "(A)")
    
    call eio%init_in (sample)

    allocate (process_instance)
    call process_instance%init (process)
    call process_instance%setup_event_data ()
    allocate (event)
    call event%basic_init ()
    call event%connect (process_instance, process%get_model_ptr ())
    
    call eio%input_i_prc (i_prc, iostat)
    if (iostat /= 0)  write (u, "(A,I0)")  "I/O error (i_prc):", iostat
    call eio%input_event (event, iostat)
    if (iostat /= 0)  write (u, "(A,I0)")  "I/O error (event):", iostat
    call eio%write (u)
    
    write (u, "(A)")
    write (u, "(1x,A,I0)")  "i_prc = ", i_prc
    write (u, "(A)")
    call event%write (u)

    write (u, "(A)")
    write (u, "(A)")  "* Generate and append another event"
    write (u, "(A)")
    
    call eio%switch_inout ()
    call event%generate (1, [0._default, 0._default])
    call event%evaluate_expressions ()
    call event%write (u)
    write (u, "(A)")

    call eio%output (event, i_prc = 5)
    call eio%write (u)
    call eio%final ()
    
    call event%final ()
    deallocate (event)
    call process_instance%final ()
    deallocate (process_instance)
    
    write (u, "(A)")
    write (u, "(A)")  "* Re-read both events"
    write (u, "(A)")
    
    call eio%init_in (sample)

    allocate (process_instance)
    call process_instance%init (process)
    call process_instance%setup_event_data ()
    allocate (event)
    call event%basic_init ()
    call event%connect (process_instance, process%get_model_ptr ())
    
    call eio%input_i_prc (i_prc, iostat)
    if (iostat /= 0)  write (u, "(A,I0)")  "I/O error (i_prc/1):", iostat
    call eio%input_event (event, iostat)
    if (iostat /= 0)  write (u, "(A,I0)")  "I/O error (event/1):", iostat
    call eio%input_i_prc (i_prc, iostat)
    if (iostat /= 0)  write (u, "(A,I0)")  "I/O error (i_prc/2):", iostat
    call eio%input_event (event, iostat)
    if (iostat /= 0)  write (u, "(A,I0)")  "I/O error (event/2):", iostat
    call eio%write (u)
    
    write (u, "(A)")
    write (u, "(1x,A,I0)")  "i_prc = ", i_prc
    write (u, "(A)")
    call event%write (u)

    write (u, "(A)")
    write (u, "(A)")  "* Cleanup"
 
    call eio%final ()
    deallocate (eio)
 
    call event%final ()
    deallocate (event)
 
    call cleanup_test_process (process, process_instance)
    deallocate (process_instance)
    deallocate (process)
    
    call model%final ()

    write (u, "(A)")
    write (u, "(A)")  "* Test output end: eio_raw_1"
    
  end subroutine eio_raw_1
  
\eatline
\nwendcode{}\nwbegindocs{1619}\nwdocspar
\subsubsection{Test I/O methods}
We test the implementation of all I/O methods.
\nwenddocs{}\nwbegincode{1620}\moddef{EIO raw: execute tests}\plusendmoddef
  call test (eio_raw_2, "eio_raw_2", &
       "handle multiple weights", &
       u, results)
\nwendcode{}\nwbegincode{1621}\moddef{EIO raw: tests}\plusendmoddef
  subroutine eio_raw_2 (u)
    use processes
    integer, intent(in) :: u
    type(model_data_t), target :: model
    type(var_list_t) :: var_list
    type(event_t), allocatable, target :: event
    type(process_t), allocatable, target :: process
    type(process_instance_t), allocatable, target :: process_instance
    type(event_sample_data_t) :: data
    class(eio_t), allocatable :: eio
    integer :: i_prc, iostat
    type(string_t) :: sample

    write (u, "(A)")  "* Test output: eio_raw_2"
    write (u, "(A)")  "*   Purpose: generate and read/write an event"
    write (u, "(A)")  "*            with multiple weights"
    write (u, "(A)")

    call model%init_test ()

    write (u, "(A)")  "* Initialize test process"
 
    allocate (process)
    allocate (process_instance)
    call prepare_test_process (process, process_instance, model)
    call process_instance%setup_event_data ()
 
    call data%init (n_proc = 1, n_alt = 2)

    call var_list_append_log (var_list, var_str ("?unweighted"), .false., &
         intrinsic = .true.)
    call var_list_append_string (var_list, var_str ("$sample_normalization"), &
         var_str ("auto"), intrinsic = .true.)
    call var_list_append_real (var_list, var_str ("safety_factor"), &
         1._default, intrinsic = .true.)

    allocate (event)
    call event%basic_init (var_list, n_alt = 2)
    call event%connect (process_instance, process%get_model_ptr ())
    
    write (u, "(A)")
    write (u, "(A)")  "* Generate and write an event"
    write (u, "(A)")
 
    sample = "eio_raw_2"
 
    allocate (eio_raw_t :: eio)
    
    call eio%init_out (sample, data)
    call event%generate (1, [0._default, 0._default])
    call event%evaluate_expressions ()
    call event%set (sqme_alt = [2._default, 3._default])
    call event%set (weight_alt = &
         [2 * event%get_weight_ref (), 3 * event%get_weight_ref ()])
    call event%store_alt_values ()
    call event%check ()

    call event%write (u)
    write (u, "(A)")

    call eio%output (event, i_prc = 42)
    call eio%write (u)
    call eio%final ()

    call event%final ()
    deallocate (event)
    call process_instance%final ()
    deallocate (process_instance)
    
    write (u, "(A)")
    write (u, "(A)")  "* Re-read the event"
    write (u, "(A)")
    
    call eio%init_in (sample, data)

    allocate (process_instance)
    call process_instance%init (process)
    call process_instance%setup_event_data ()
    allocate (event)
    call event%basic_init (var_list, n_alt = 2)
    call event%connect (process_instance, process%get_model_ptr ())
    
    call eio%input_i_prc (i_prc, iostat)
    if (iostat /= 0)  write (u, "(A,I0)")  "I/O error (i_prc):", iostat
    call eio%input_event (event, iostat)
    if (iostat /= 0)  write (u, "(A,I0)")  "I/O error (event):", iostat
    call eio%write (u)
    
    write (u, "(A)")
    write (u, "(1x,A,I0)")  "i_prc = ", i_prc
    write (u, "(A)")
    call event%write (u)

    write (u, "(A)")
    write (u, "(A)")  "* Cleanup"
 
    call eio%final ()
    deallocate (eio)
 
    call event%final ()
    deallocate (event)
 
    call cleanup_test_process (process, process_instance)
    deallocate (process_instance)
    deallocate (process)
    
    call model%final ()

    write (u, "(A)")
    write (u, "(A)")  "* Test output end: eio_raw_2"
    
  end subroutine eio_raw_2
  
\eatline
\nwendcode{}\nwbegindocs{1622}\nwdocspar
%%%%%%%%%%%%%%%%%%%%%%%%%%%%%%%%%%%%%%%%%%%%%%%%%%%%%%%%%%%%%%%%%%%%%%%%
\chapter{Integration and simulation}

This layer of modules is just below the top-level API.  We lay out specific
data types for integration and simulation and implement the corresponding
algorithms as methods acting on them.  This helps to keep the command-level
implementation concise and simple.

Generic modules for connecting Sindarin with \whizard:
\begin{description}
\item[user\_files]
  Manage files that are accessed by user-written Sindarin code.
\item[rt\_data]
  The main runtime data block, giving access to all features of
  \whizard via contained data structures or pointers.
\item[dispatch]
  Dynamic dispatch for various \whizard\ objects that have multiple
  implementations.  The dispatcher uses data from the {\tt{}rt{\char95}data}
  block to decide which path to follow.
\end{description}

More specific modules, the workhorses of \whizard.  These modules use
{\tt{}rt{\char95}data} components and the dispatcher.  They define dedicated
object types for bundling lower-level data and procedures that together
provide the desired functionality.
\begin{description}
\item[process\_configurations]
  Interpret process definitions to generate matrix-element
  objects and external code.
\item[compilations]
  Compile, link and access external matrix-element code.
\item[integrations]
  Manage integration runs.
\item[event\_streams]
  Manage external event I/O.
\item[simulations]
  Manage simulation runs.
\end{description}

\clearpage
%%%%%%%%%%%%%%%%%%%%%%%%%%%%%%%%%%%%%%%%%%%%%%%%%%%%%%%%%%%%%%%%%%%%%%%%
\section{User-controlled File I/O}

The SINDARIN language includes commands that write output to file (input may
be added later).  We identify files by their name, and manage the unit
internally.  We need procedures for opening, closing, and printing files.

\nwenddocs{}\nwbegincode{1623}\moddef{\code{}user{\char95}files.f90\edoc{}}\endmoddef
\LA{}File header\RA{}

module user_files

\LA{}Use strings\RA{}
  use io_units
  use diagnostics
  use ifiles
  use analysis

\LA{}Standard module head\RA{}

\LA{}User files: public\RA{}

\LA{}User files: types\RA{}

\LA{}User files: interfaces\RA{}

contains

\LA{}User files: procedures\RA{}

end module user_files
\eatline
\nwendcode{}\nwbegindocs{1624}\nwdocspar
\subsection{The file type}
This is a type that describes an open user file and its properties.  The entry
is part of a doubly-linked list.
\nwenddocs{}\nwbegincode{1625}\moddef{User files: types}\endmoddef
  type :: file_t
     private
     type(string_t) :: name
     integer :: unit = -1
     logical :: reading = .false.
     logical :: writing = .false.
     type(file_t), pointer :: prev => null ()
     type(file_t), pointer :: next => null ()
  end type file_t

\eatline
\nwendcode{}\nwbegindocs{1626}The initializer opens the file.
\nwenddocs{}\nwbegincode{1627}\moddef{User files: procedures}\endmoddef
  subroutine file_init (file, name, action, status, position)
    type(file_t), intent(out) :: file
    type(string_t), intent(in) :: name
    character(len=*), intent(in) :: action, status, position
    file%unit = free_unit ()
    file%name = name
    open (unit = file%unit, file = char (file%name), &
          action = action, status = status, position = position)
    select case (action)
    case ("read")
       file%reading = .true.
    case ("write")
       file%writing = .true.
    case ("readwrite")
       file%reading = .true.    
       file%writing = .true.    
    end select
  end subroutine file_init

\eatline
\nwendcode{}\nwbegindocs{1628}The finalizer closes it.
\nwenddocs{}\nwbegincode{1629}\moddef{User files: procedures}\plusendmoddef
  subroutine file_final (file)
    type(file_t), intent(inout) :: file
    close (unit = file%unit)
    file%unit = -1
  end subroutine file_final

\eatline
\nwendcode{}\nwbegindocs{1630}Check if a file is open with correct status.
\nwenddocs{}\nwbegincode{1631}\moddef{User files: procedures}\plusendmoddef
  function file_is_open (file, action) result (flag)
    logical :: flag
    type(file_t), intent(in) :: file
    character(*), intent(in) :: action
    select case (action)
    case ("read")
       flag = file%reading
    case ("write")
       flag = file%writing
    case ("readwrite")
       flag = file%reading .and. file%writing
    case default
       call msg_bug ("Checking file '" // char (file%name) &
            // "': illegal action specifier")
    end select
  end function file_is_open

\eatline
\nwendcode{}\nwbegindocs{1632}Write to the file.  Error if in wrong mode.  If there is no string, just
write an empty record.  If there is a string, respect the {\tt{}advancing}
option.
\nwenddocs{}\nwbegincode{1633}\moddef{User files: procedures}\plusendmoddef
  subroutine file_write_string (file, string, advancing)
    type(file_t), intent(in) :: file
    type(string_t), intent(in), optional :: string
    logical, intent(in), optional :: advancing
    if (file%writing) then
       if (present (string)) then
          if (present (advancing)) then
             if (advancing) then
                write (file%unit, "(A)")  char (string)
             else
                write (file%unit, "(A)", advance="no")  char (string)
             end if
          else
             write (file%unit, "(A)")  char (string)
          end if
       else
          write (file%unit, *)
       end if
    else
       call msg_error ("Writing to file: File '" // char (file%name) &
            // "' is not open for writing.")
    end if
  end subroutine file_write_string

\eatline
\nwendcode{}\nwbegindocs{1634}Write a whole ifile, line by line.
\nwenddocs{}\nwbegincode{1635}\moddef{User files: procedures}\plusendmoddef
  subroutine file_write_ifile (file, ifile)
    type(file_t), intent(in) :: file
    type(ifile_t), intent(in) :: ifile
    type(line_p) :: line
    call line_init (line, ifile)
    do while (line_is_associated (line))
       call file_write_string (file, line_get_string_advance (line))
    end do
  end subroutine file_write_ifile

\eatline
\nwendcode{}\nwbegindocs{1636}Write an analysis object (or all objects) to an open file.
\nwenddocs{}\nwbegincode{1637}\moddef{User files: procedures}\plusendmoddef
  subroutine file_write_analysis (file, tag)
    type(file_t), intent(in) :: file
    type(string_t), intent(in), optional :: tag
    if (file%writing) then
       if (present (tag)) then
          call analysis_write (tag, unit = file%unit)
       else
          call analysis_write (unit = file%unit)
       end if
    else
       call msg_error ("Writing analysis to file: File '" // char (file%name) &
            // "' is not open for writing.")
    end if
  end subroutine file_write_analysis

\eatline
\nwendcode{}\nwbegindocs{1638}\nwdocspar
\subsection{The file list}
We maintain a list of all open files and their attributes.  The list must be
doubly-linked because we may delete entries.
\nwenddocs{}\nwbegincode{1639}\moddef{User files: public}\endmoddef
  public :: file_list_t
\nwendcode{}\nwbegincode{1640}\moddef{User files: types}\plusendmoddef
  type :: file_list_t
     type(file_t), pointer :: first => null ()
     type(file_t), pointer :: last => null ()
  end type file_list_t

\eatline
\nwendcode{}\nwbegindocs{1641}There is no initialization routine, but a finalizer which deletes all:
\nwenddocs{}\nwbegincode{1642}\moddef{User files: public}\plusendmoddef
  public :: file_list_final
\nwendcode{}\nwbegincode{1643}\moddef{User files: procedures}\plusendmoddef
  subroutine file_list_final (file_list)
    type(file_list_t), intent(inout) :: file_list
    type(file_t), pointer :: current
    do while (associated (file_list%first))
       current => file_list%first
       file_list%first => current%next
       call file_final (current)
       deallocate (current)
    end do
    file_list%last => null ()
  end subroutine file_list_final

\eatline
\nwendcode{}\nwbegindocs{1644}Find an entry in the list.  Return null pointer on failure.
\nwenddocs{}\nwbegincode{1645}\moddef{User files: procedures}\plusendmoddef
  function file_list_get_file_ptr (file_list, name) result (current)
    type(file_t), pointer :: current
    type(file_list_t), intent(in) :: file_list
    type(string_t), intent(in) :: name
    current => file_list%first
    do while (associated (current))
       if (current%name == name)  return
       current => current%next
    end do
  end function file_list_get_file_ptr

\eatline
\nwendcode{}\nwbegindocs{1646}Check if a file is open, public version:
\nwenddocs{}\nwbegincode{1647}\moddef{User files: public}\plusendmoddef
  public :: file_list_is_open
\nwendcode{}\nwbegincode{1648}\moddef{User files: procedures}\plusendmoddef
  function file_list_is_open (file_list, name, action) result (flag)
    logical :: flag
    type(file_list_t), intent(in) :: file_list
    type(string_t), intent(in) :: name
    character(len=*), intent(in) :: action
    type(file_t), pointer :: current
    current => file_list_get_file_ptr (file_list, name)
    if (associated (current)) then
       flag = file_is_open (current, action)
    else
       flag = .false.
    end if
  end function file_list_is_open

\eatline
\nwendcode{}\nwbegindocs{1649}Append a new file entry, i.e., open this file.  Error if it is
already open.
\nwenddocs{}\nwbegincode{1650}\moddef{User files: public}\plusendmoddef
  public :: file_list_open
\nwendcode{}\nwbegincode{1651}\moddef{User files: procedures}\plusendmoddef
  subroutine file_list_open (file_list, name, action, status, position)
    type(file_list_t), intent(inout) :: file_list
    type(string_t), intent(in) :: name
    character(len=*), intent(in) :: action, status, position
    type(file_t), pointer :: current
    if (.not. associated (file_list_get_file_ptr (file_list, name))) then
       allocate (current)
       call msg_message ("Opening file '" // char (name) // "' for output")
       call file_init (current, name, action, status, position)
       if (associated (file_list%last)) then
          file_list%last%next => current
          current%prev => file_list%last
       else
          file_list%first => current
       end if
       file_list%last => current
    else
       call msg_error ("Opening file: File '" // char (name) &
            // "' is already open.")
    end if
  end subroutine file_list_open

\eatline
\nwendcode{}\nwbegindocs{1652}Delete a file entry, i.e., close this file.  Error if it is not open.
\nwenddocs{}\nwbegincode{1653}\moddef{User files: public}\plusendmoddef
  public :: file_list_close
\nwendcode{}\nwbegincode{1654}\moddef{User files: procedures}\plusendmoddef
  subroutine file_list_close (file_list, name)
    type(file_list_t), intent(inout) :: file_list
    type(string_t), intent(in) :: name
    type(file_t), pointer :: current
    current => file_list_get_file_ptr (file_list, name)
    if (associated (current)) then
       if (associated (current%prev)) then
          current%prev%next => current%next
       else
          file_list%first => current%next
       end if
       if (associated (current%next)) then
          current%next%prev => current%prev
       else
          file_list%last => current%prev
       end if
       call msg_message ("Closing file '" // char (name) // "' for output")
       call file_final (current)
       deallocate (current)
    else
       call msg_error ("Closing file: File '" // char (name) &
            // "' is not open.")
    end if
  end subroutine file_list_close

\eatline
\nwendcode{}\nwbegindocs{1655}Write a string to file.  Error if it is not open.
\nwenddocs{}\nwbegincode{1656}\moddef{User files: public}\plusendmoddef
  public :: file_list_write
\nwendcode{}\nwbegincode{1657}\moddef{User files: interfaces}\endmoddef
  interface file_list_write
     module procedure file_list_write_string
     module procedure file_list_write_ifile
  end interface
\nwendcode{}\nwbegincode{1658}\moddef{User files: procedures}\plusendmoddef
  subroutine file_list_write_string (file_list, name, string, advancing)
    type(file_list_t), intent(in) :: file_list
    type(string_t), intent(in) :: name
    type(string_t), intent(in), optional :: string
    logical, intent(in), optional :: advancing
    type(file_t), pointer :: current
    current => file_list_get_file_ptr (file_list, name)
    if (associated (current)) then
       call file_write_string (current, string, advancing)
    else
       call msg_error ("Writing to file: File '" // char (name) &
            // "'is not open.")
    end if
  end subroutine file_list_write_string

  subroutine file_list_write_ifile (file_list, name, ifile)
    type(file_list_t), intent(in) :: file_list
    type(string_t), intent(in) :: name
    type(ifile_t), intent(in) :: ifile
    type(line_p) :: line
    type(file_t), pointer :: current
    current => file_list_get_file_ptr (file_list, name)
    if (associated (current)) then
       call file_write_ifile (current, ifile)
    else
       call msg_error ("Writing to file: File '" // char (name) &
            // "'is not open.")
    end if
  end subroutine file_list_write_ifile

\eatline
\nwendcode{}\nwbegindocs{1659}Write an analysis object or all objects to data file.  Error if it is not
open.  If the file name is empty, write to standard output.
\nwenddocs{}\nwbegincode{1660}\moddef{User files: public}\plusendmoddef
  public :: file_list_write_analysis
\nwendcode{}\nwbegincode{1661}\moddef{User files: procedures}\plusendmoddef
  subroutine file_list_write_analysis (file_list, name, tag)
    type(file_list_t), intent(in) :: file_list
    type(string_t), intent(in) :: name
    type(string_t), intent(in), optional :: tag
    type(file_t), pointer :: current
    if (name == "") then
       if (present (tag)) then
          call analysis_write (tag)
       else
          call analysis_write
       end if
    else
       current => file_list_get_file_ptr (file_list, name)
       if (associated (current)) then
          call file_write_analysis (current, tag)
       else
          call msg_error ("Writing analysis to file: File '" // char (name) &
               // "' is not open.")
       end if
    end if
  end subroutine file_list_write_analysis

\eatline
\nwendcode{}\nwbegindocs{1662}\nwdocspar
\clearpage
%%%%%%%%%%%%%%%%%%%%%%%%%%%%%%%%%%%%%%%%%%%%%%%%%%%%%%%%%%%%%%%%%%%%%%%%
\section{Runtime data}

\nwenddocs{}\nwbegincode{1663}\moddef{\code{}rt{\char95}data.f90\edoc{}}\endmoddef
\LA{}File header\RA{}

module rt_data

\LA{}Use kinds\RA{}
\LA{}Use strings\RA{}
  use io_units
  use format_utils, only: write_separator
  use unit_tests
  use system_dependencies
  use diagnostics
  use pdf_builtin !NODEP!
  use sf_lhapdf !NODEP!
  use os_interface
  use ifiles
  use lexers
  use parser
  use models
  use flavors
  use jets
  use subevents
  use pdg_arrays
  use variables
  use eval_trees
  use polarizations
  use beams
  use process_libraries
  use prclib_stacks
  use prc_core
  use beam_structures
  use user_files
  use process_stacks
  use iterations

\LA{}Standard module head\RA{}

\LA{}RT data: public\RA{}

\LA{}RT data: types\RA{}

contains

\LA{}RT data: procedures\RA{}

\LA{}RT data: tests\RA{}

end module rt_data
\eatline
\nwendcode{}\nwbegindocs{1664}\nwdocspar
\subsection{Strategy for models and variables}
The program manages its data via a main {\tt{}rt{\char95}data{\char95}t} object.  During program
flow, various commands create and use local {\tt{}rt{\char95}data{\char95}t} objects.  Those
transient blocks contain either pointers to global object or local copies
which are deleted after use.

Each {\tt{}rt{\char95}data{\char95}t} object contains a variable list component.  This lists
holds (local copies of) all kinds of intrinsic or user-defined variables.  The
variable list is linked to the variable list contained in the local process
library.  This, in turn, is linked to the variable list of the {\tt{}rt{\char95}data{\char95}t}
context, and so on.  

A variable lookup will thus be recursively delegated to the linked variable
lists, until a match is found.  When modifying a variable which is not yet
local, the program creates a local copy and uses this afterwards.  Thus, when
the local {\tt{}rt{\char95}data{\char95}t} object is deleted, the context value is recovered.

Models are kept in a model list which is separate from the variable list.
Otherwise, they are treated in a similar manner: the local list is linked to
the context model list.  Model lookup is thus recursively delegated.  When a
model or any part of it is modified, the model is copied to the local
{\tt{}rt{\char95}data{\char95}t} object, so the context model is not modified.  Commands such as
{\tt{}integrate} will create their own copy of the current model (and of the
current variable list) at the point where they are executed.

When a model is encountered for the first time, it is read from file.  The
reading is automatically delegated to the global context.  Thus, this master
copy survives until the main {\tt{}rt{\char95}data{\char95}t} object is deleted, at program
completion.

If there is a currently active model, its variable list is linked to the main
variable list.  Variable lookups will then start from the model variable
list.  When the current model is switched, the new active model will get this
link instead.  Consequently, a change to the current model is kept as long as
this model has a local copy; it survives local model switches.  On the other
hand, a parameter change in the current model doesn't affect any other model,
even if the parameter name is identical.
\nwenddocs{}\nwbegindocs{1665}\nwdocspar
\subsection{Container for parse nodes}
The runtime data set contains a bunch of parse nodes (chunks of code
that have not been compiled into evaluation trees but saved for later
use).  We collect them here.

This implementation has the useful effect that an assignment between two
objects of this type will establish a pointer-target relationship for
all components.
\nwenddocs{}\nwbegincode{1666}\moddef{RT data: types}\endmoddef
  type :: rt_parse_nodes_t
     type(parse_node_t), pointer :: cuts_lexpr => null ()
     type(parse_node_t), pointer :: scale_expr => null ()     
     type(parse_node_t), pointer :: fac_scale_expr => null ()
     type(parse_node_t), pointer :: ren_scale_expr => null ()     
     type(parse_node_t), pointer :: weight_expr => null ()
     type(parse_node_t), pointer :: selection_lexpr => null ()
     type(parse_node_t), pointer :: reweight_expr => null ()
     type(parse_node_t), pointer :: analysis_lexpr => null ()
     type(parse_node_p), dimension(:), allocatable :: alt_setup
   contains
   \LA{}RT data: rt parse nodes: TBP\RA{}
  end type rt_parse_nodes_t
     
\eatline
\nwendcode{}\nwbegindocs{1667}Clear individual components.  The parse nodes are nullified.  No
finalization needed since the pointer targets are part of the global
parse tree.
\nwenddocs{}\nwbegincode{1668}\moddef{RT data: rt parse nodes: TBP}\endmoddef
  procedure :: clear => rt_parse_nodes_clear
\nwendcode{}\nwbegincode{1669}\moddef{RT data: procedures}\endmoddef
  subroutine rt_parse_nodes_clear (rt_pn, name)
    class(rt_parse_nodes_t), intent(inout) :: rt_pn
    type(string_t), intent(in) :: name
    select case (char (name))
    case ("cuts")
       rt_pn%cuts_lexpr => null ()
    case ("scale")
       rt_pn%scale_expr => null ()
    case ("factorization_scale")
       rt_pn%fac_scale_expr => null ()
    case ("renormalization_scale")
       rt_pn%ren_scale_expr => null ()
    case ("weight")
       rt_pn%weight_expr => null ()
    case ("selection")
       rt_pn%selection_lexpr => null ()
    case ("reweight")
       rt_pn%reweight_expr => null ()
    case ("analysis")
       rt_pn%analysis_lexpr => null ()
    end select
  end subroutine rt_parse_nodes_clear
  
\eatline
\nwendcode{}\nwbegindocs{1670}Output for the parse nodes.
\nwenddocs{}\nwbegincode{1671}\moddef{RT data: rt parse nodes: TBP}\plusendmoddef
  procedure :: write => rt_parse_nodes_write
\nwendcode{}\nwbegincode{1672}\moddef{RT data: procedures}\plusendmoddef
  subroutine rt_parse_nodes_write (object, unit)
    class(rt_parse_nodes_t), intent(in) :: object
    integer, intent(in), optional :: unit
    integer :: u, i
    u = given_output_unit (unit)
    call wrt ("Cuts", object%cuts_lexpr)
    call write_separator (u)
    call wrt ("Scale", object%scale_expr)
    call write_separator (u)
    call wrt ("Factorization scale", object%fac_scale_expr)
    call write_separator (u)
    call wrt ("Renormalization scale", object%ren_scale_expr)
    call write_separator (u)
    call wrt ("Weight", object%weight_expr)
    call write_separator (u, 2)
    call wrt ("Event selection", object%selection_lexpr)
    call write_separator (u)
    call wrt ("Event reweighting factor", object%reweight_expr)
    call write_separator (u)
    call wrt ("Event analysis", object%analysis_lexpr)
    if (allocated (object%alt_setup)) then
       call write_separator (u, 2)
       write (u, "(1x,A,':')")  "Alternative setups"
       do i = 1, size (object%alt_setup)
          call write_separator (u)
          call wrt ("Commands", object%alt_setup(i)%ptr)
       end do
    end if
  contains
    subroutine wrt (title, pn)
      character(*), intent(in) :: title
      type(parse_node_t), intent(in), pointer :: pn
      if (associated (pn)) then
         write (u, "(1x,A,':')")  title
         call write_separator (u)
         call parse_node_write_rec (pn, u)
      else
         write (u, "(1x,A,':',1x,A)")  title, "[undefined]"
      end if
    end subroutine wrt
  end subroutine rt_parse_nodes_write
    
\eatline
\nwendcode{}\nwbegindocs{1673}Screen output for individual components.  (This should eventually be more
condensed, currently we print the internal representation tree.)
\nwenddocs{}\nwbegincode{1674}\moddef{RT data: rt parse nodes: TBP}\plusendmoddef
  procedure :: show => rt_parse_nodes_show
\nwendcode{}\nwbegincode{1675}\moddef{RT data: procedures}\plusendmoddef
  subroutine rt_parse_nodes_show (rt_pn, name, unit)
    class(rt_parse_nodes_t), intent(in) :: rt_pn
    type(string_t), intent(in) :: name
    integer, intent(in), optional :: unit
    type(parse_node_t), pointer :: pn
    integer :: u
    u = given_output_unit (unit)
    select case (char (name))
    case ("cuts")
       pn => rt_pn%cuts_lexpr
    case ("scale")
       pn => rt_pn%scale_expr
    case ("factorization_scale")
       pn => rt_pn%fac_scale_expr
    case ("renormalization_scale")
       pn => rt_pn%ren_scale_expr
    case ("weight")
       pn => rt_pn%weight_expr
    case ("selection")
       pn => rt_pn%selection_lexpr
    case ("reweight")
       pn => rt_pn%reweight_expr
    case ("analysis")
       pn => rt_pn%analysis_lexpr
    end select
    if (associated (pn)) then
       write (u, "(A,1x,A,1x,A)")  "Expression:", char (name), "(parse tree):"
       call parse_node_write_rec (pn, u)
    else
       write (u, "(A,1x,A,A)")  "Expression:", char (name), ": [undefined]"
    end if
  end subroutine rt_parse_nodes_show
  
\eatline
\nwendcode{}\nwbegindocs{1676}\nwdocspar
\subsection{The data type}
This is a big data container which contains everything that is used and
modified during the command flow.  A local copy of this can be used to
temporarily override defaults.  The data set is transparent.
\nwenddocs{}\nwbegincode{1677}\moddef{RT data: public}\endmoddef
  public :: rt_data_t
\nwendcode{}\nwbegincode{1678}\moddef{RT data: types}\plusendmoddef
  type :: rt_data_t
     type(lexer_t), pointer :: lexer => null ()
     type(rt_data_t), pointer :: context => null ()
     type(var_list_t) :: var_list
     type(iterations_list_t) :: it_list
     type(os_data_t) :: os_data
     type(model_list_t) :: model_list
     type(model_t), pointer :: model => null ()
     logical :: model_is_copy = .false.
     type(model_t), pointer :: preload_model => null ()
     type(model_t), pointer :: fallback_model => null ()
     type(model_t), pointer :: radiation_model => null ()
     type(prclib_stack_t) :: prclib_stack
     type(process_library_t), pointer :: prclib => null ()
     type(beam_structure_t) :: beam_structure
     type(rt_parse_nodes_t) :: pn
     type(process_stack_t) :: process_stack
     type(string_t), dimension(:), allocatable :: sample_fmt
     type(file_list_t), pointer :: out_files => null ()
     logical :: quit = .false.
     integer :: quit_code = 0
     type(string_t) :: logfile 
     logical :: nlo_calculation = .false.
     logical, dimension(3) :: active_nlo_components
   contains
   \LA{}RT data: rt data: TBP\RA{}
  end type rt_data_t

\eatline
\nwendcode{}\nwbegindocs{1679}\nwdocspar
\subsection{Output}
\nwenddocs{}\nwbegincode{1680}\moddef{RT data: rt data: TBP}\endmoddef
  procedure :: write => rt_data_write
\nwendcode{}\nwbegincode{1681}\moddef{RT data: procedures}\plusendmoddef
  subroutine rt_data_write (object, unit, vars, pacify)
    class(rt_data_t), intent(in) :: object
    integer, intent(in), optional :: unit
    type(string_t), dimension(:), intent(in), optional :: vars
    logical, intent(in), optional :: pacify
    integer :: u, i
    u = given_output_unit (unit)
    call write_separator (u, 2)
    write (u, "(1x,A)")  "Runtime data:"
    if (present (vars)) then
       if (size (vars) /= 0) then
          call write_separator (u, 2)
          write (u, "(1x,A)")  "Selected variables:"
          call write_separator (u)
          call object%write_vars (u, vars)
       end if
    else
       call write_separator (u, 2)
       if (associated (object%model)) then
          call object%model%write_var_list (u, follow_link=.true.)
       else
          call var_list_write (object%var_list, u, follow_link=.true.)
       end if
    end if
    if (object%it_list%get_n_pass () > 0) then
       call write_separator (u, 2)
       write (u, "(1x)", advance="no")
       call object%it_list%write (u)
    end if
    if (associated (object%model)) then
       call write_separator (u, 2)
       call object%model%write (u)
    end if
    call object%prclib_stack%write (u)
    call object%beam_structure%write (u)
    call write_separator (u, 2)
    call object%pn%write (u)
    if (allocated (object%sample_fmt)) then
       call write_separator (u)
       write (u, "(1x,A)", advance="no")  "Event sample formats = "
       do i = 1, size (object%sample_fmt)
          if (i > 1)  write (u, "(A,1x)", advance="no")  ","
          write (u, "(A)", advance="no")  char (object%sample_fmt(i))
       end do
       write (u, "(A)")
    end if
    call object%process_stack%write (u, pacify)
    write (u, "(1x,A,1x,L1)")  "quit     :", object%quit
    write (u, "(1x,A,1x,I0)")  "quit_code:", object%quit_code
    call write_separator (u, 2)
    write (u, "(1x,A,1x,A)")   "Logfile  :", "'" // trim (char (object%logfile)) // "'"
    call write_separator (u, 2)
  end subroutine rt_data_write
  
\eatline
\nwendcode{}\nwbegindocs{1682}Write only selected variables.
\nwenddocs{}\nwbegincode{1683}\moddef{RT data: rt data: TBP}\plusendmoddef
  procedure :: write_vars => rt_data_write_vars
\nwendcode{}\nwbegincode{1684}\moddef{RT data: procedures}\plusendmoddef
  subroutine rt_data_write_vars (object, unit, vars)
    class(rt_data_t), intent(in), target :: object
    integer, intent(in), optional :: unit
    type(string_t), dimension(:), intent(in), optional :: vars
    type(var_list_t), pointer :: var_list
    integer :: u, i
    u = given_output_unit (unit)
    if (present (vars)) then
       var_list => object%get_var_list_ptr ()
       do i = 1, size (vars)
          associate (var => vars(i))
            if (var_list%contains (var, follow_link=.true.)) then
               call var_list_write_var (var_list, var, unit = u, &
                    follow_link = .true.)
            end if
          end associate
       end do
    end if
  end subroutine rt_data_write_vars
  
\eatline
\nwendcode{}\nwbegindocs{1685}Write only the model list.
\nwenddocs{}\nwbegincode{1686}\moddef{RT data: rt data: TBP}\plusendmoddef
  procedure :: write_model_list => rt_data_write_model_list
\nwendcode{}\nwbegincode{1687}\moddef{RT data: procedures}\plusendmoddef
  subroutine rt_data_write_model_list (object, unit)
    class(rt_data_t), intent(in) :: object
    integer, intent(in), optional :: unit
    integer :: u
    u = given_output_unit (unit)
    call object%model_list%write (u)
  end subroutine rt_data_write_model_list

\eatline
\nwendcode{}\nwbegindocs{1688}Write only the library stack.
\nwenddocs{}\nwbegincode{1689}\moddef{RT data: rt data: TBP}\plusendmoddef
  procedure :: write_libraries => rt_data_write_libraries
\nwendcode{}\nwbegincode{1690}\moddef{RT data: procedures}\plusendmoddef
  subroutine rt_data_write_libraries (object, unit, libpath)
    class(rt_data_t), intent(in) :: object
    integer, intent(in), optional :: unit
    logical, intent(in), optional :: libpath
    integer :: u
    u = given_output_unit (unit)
    call object%prclib_stack%write (u, libpath)
  end subroutine rt_data_write_libraries

\eatline
\nwendcode{}\nwbegindocs{1691}Write only the beam data.
\nwenddocs{}\nwbegincode{1692}\moddef{RT data: rt data: TBP}\plusendmoddef
  procedure :: write_beams => rt_data_write_beams
\nwendcode{}\nwbegincode{1693}\moddef{RT data: procedures}\plusendmoddef
  subroutine rt_data_write_beams (object, unit)
    class(rt_data_t), intent(in) :: object
    integer, intent(in), optional :: unit
    integer :: u
    u = given_output_unit (unit)
    call write_separator (u, 2)
    call object%beam_structure%write (u)
    call write_separator (u, 2)
  end subroutine rt_data_write_beams

\eatline
\nwendcode{}\nwbegindocs{1694}Write only the process and event expressions.
\nwenddocs{}\nwbegincode{1695}\moddef{RT data: rt data: TBP}\plusendmoddef
  procedure :: write_expr => rt_data_write_expr
\nwendcode{}\nwbegincode{1696}\moddef{RT data: procedures}\plusendmoddef
  subroutine rt_data_write_expr (object, unit)
    class(rt_data_t), intent(in) :: object
    integer, intent(in), optional :: unit
    integer :: u
    u = given_output_unit (unit)
    call write_separator (u, 2)
    call object%pn%write (u)
    call write_separator (u, 2)
  end subroutine rt_data_write_expr
  
\eatline
\nwendcode{}\nwbegindocs{1697}Write only the process stack.
\nwenddocs{}\nwbegincode{1698}\moddef{RT data: rt data: TBP}\plusendmoddef
  procedure :: write_process_stack => rt_data_write_process_stack
\nwendcode{}\nwbegincode{1699}\moddef{RT data: procedures}\plusendmoddef
  subroutine rt_data_write_process_stack (object, unit)
    class(rt_data_t), intent(in) :: object
    integer, intent(in), optional :: unit
    call object%process_stack%write (unit)
  end subroutine rt_data_write_process_stack
  
\eatline
\nwendcode{}\nwbegindocs{1700}\nwdocspar
\subsection{Clear}
The {\tt{}clear} command can remove the contents of various subobjects.
The objects themselves should stay.
\nwenddocs{}\nwbegincode{1701}\moddef{RT data: rt data: TBP}\plusendmoddef
  procedure :: clear_beams => rt_data_clear_beams
\nwendcode{}\nwbegincode{1702}\moddef{RT data: procedures}\plusendmoddef
  subroutine rt_data_clear_beams (global)
    class(rt_data_t), intent(inout) :: global
    call global%beam_structure%final_sf ()
    call global%beam_structure%final_pol ()
    call global%beam_structure%final_mom ()
  end subroutine rt_data_clear_beams
  
\eatline
\nwendcode{}\nwbegindocs{1703}\nwdocspar
\subsection{Initialization}
Initialize runtime data.  This defines special variables such as
{\tt{}sqrts}, and should be done only for the instance that is actually
global.   Local copies will inherit the special variables.

We link the global variable list to the process stack variable list,
so the latter is always available (and kept global).
\nwenddocs{}\nwbegincode{1704}\moddef{RT data: rt data: TBP}\plusendmoddef
  procedure :: global_init => rt_data_global_init
\nwendcode{}\nwbegincode{1705}\moddef{RT data: procedures}\plusendmoddef
  subroutine rt_data_global_init (global, paths, logfile)
    class(rt_data_t), intent(out), target :: global
    type(paths_t), intent(in), optional :: paths
    type(string_t), intent(in), optional :: logfile
    logical, target, save :: known = .true.
    integer :: seed
    real(default), parameter :: real_specimen = 1.
    call os_data_init (global%os_data, paths)
    if (present (logfile)) then
       global%logfile = logfile
    else
       global%logfile = ""
    end if
    allocate (global%out_files)
    call system_clock (seed)
    call var_list_append_log_ptr &
         (global%var_list, var_str ("?logging"), logging, known, &
         intrinsic=.true.)
    call var_list_append_int &
         (global%var_list, var_str ("seed"), seed, &
          intrinsic=.true.)
    call var_list_append_string &
         (global%var_list, var_str ("$model_name"), &
          intrinsic=.true.)
    call var_list_append_int &
         (global%var_list, var_str ("process_num_id"), &
         intrinsic=.true.)        
    call var_list_append_string &
         (global%var_list, var_str ("$method"), var_str ("omega"), &
         intrinsic=.true.)        
    call var_list_append_log &
         (global%var_list, var_str ("?report_progress"), .true., &
          intrinsic=.true.)
    call var_list_append_string &
         (global%var_list, var_str ("$restrictions"), var_str (""), &
         intrinsic=.true.)
    call var_list_append_string &
         (global%var_list, var_str ("$omega_flags"), var_str (""), &
         intrinsic=.true.)
    call var_list_append_log &
         (global%var_list, var_str ("?read_color_factors"), .true., &
          intrinsic=.true.)
!!! JRR: WK please check (#529)    
!     call var_list_append_string &
!          (global%var_list, var_str ("$user_procs_cut"), var_str (""), &
!           intrinsic=.true.)     
!     call var_list_append_string &
!          (global%var_list, var_str ("$user_procs_event_shape"), var_str (""), &
!           intrinsic=.true.)     
!     call var_list_append_string &
!          (global%var_list, var_str ("$user_procs_obs1"), var_str (""), &
!           intrinsic=.true.)     
!     call var_list_append_string &
!          (global%var_list, var_str ("$user_procs_obs2"), var_str (""), &
!           intrinsic=.true.)     
!     call var_list_append_string &
!          (global%var_list, var_str ("$user_procs_sf"), var_str (""), &
!           intrinsic=.true.)     
    call var_list_append_log &
         (global%var_list, var_str ("?slha_read_input"), .true., &
          intrinsic=.true.)
    call var_list_append_log &
         (global%var_list, var_str ("?slha_read_spectrum"), .true., &
          intrinsic=.true.)
    call var_list_append_log &
         (global%var_list, var_str ("?slha_read_decays"), .false., &
          intrinsic=.true.)
    call var_list_append_string &
         (global%var_list, var_str ("$library_name"), &
          intrinsic=.true.)
    call var_list_append_real &
         (global%var_list, var_str ("sqrts"), &
          intrinsic=.true.)
    call var_list_append_real &
         (global%var_list, var_str ("luminosity"), 0._default, &
          intrinsic=.true.)
    call var_list_append_log &
         (global%var_list, var_str ("?sf_trace"), .false., &
          intrinsic=.true.)
    call var_list_append_string &
         (global%var_list, var_str ("$sf_trace_file"), var_str (""), &
          intrinsic=.true.)
    call var_list_append_log &
         (global%var_list, var_str ("?sf_allow_s_mapping"), .true., &
          intrinsic=.true.)
    if (present (paths)) then
       call var_list_append_string &
            (global%var_list, var_str ("$lhapdf_dir"), paths%lhapdfdir, &
             intrinsic=.true.)
    else
       call var_list_append_string &
            (global%var_list, var_str ("$lhapdf_dir"), var_str(""), &
             intrinsic=.true.)
    end if 
    call var_list_append_string &
         (global%var_list, var_str ("$lhapdf_file"), var_str (""), &
          intrinsic=.true.)
    call var_list_append_string &
         (global%var_list, var_str ("$lhapdf_photon_file"), var_str (""), &
          intrinsic=.true.)    
    call var_list_append_int &
         (global%var_list, var_str ("lhapdf_member"), 0, &
          intrinsic=.true.)
    call var_list_append_int &
         (global%var_list, var_str ("lhapdf_photon_scheme"), 0, &
          intrinsic=.true.)
    call var_list_append_log &
         (global%var_list, var_str ("?hoppet_b_matching"), .false., &
          intrinsic=.true.)
    call var_list_append_real &
         (global%var_list, var_str ("isr_alpha"), 0._default, &
          intrinsic=.true.)
    call var_list_append_real &
         (global%var_list, var_str ("isr_q_max"), 0._default, &
          intrinsic=.true.)
    call var_list_append_real &
         (global%var_list, var_str ("isr_mass"), 0._default, &
          intrinsic=.true.)
    call var_list_append_int &
         (global%var_list, var_str ("isr_order"), 3, &
          intrinsic=.true.)
    call var_list_append_log &
         (global%var_list, var_str ("?isr_recoil"), .false., &
          intrinsic=.true.)
    call var_list_append_real &
         (global%var_list, var_str ("epa_alpha"), 0._default, &
          intrinsic=.true.)
    call var_list_append_real &
         (global%var_list, var_str ("epa_x_min"), 0._default, &
          intrinsic=.true.)
    call var_list_append_real &
         (global%var_list, var_str ("epa_q_min"), 0._default, &
          intrinsic=.true.)
    call var_list_append_real &
         (global%var_list, var_str ("epa_e_max"), 0._default, &
          intrinsic=.true.)
    call var_list_append_real &
         (global%var_list, var_str ("epa_mass"), 0._default, &
          intrinsic=.true.)
    call var_list_append_log &
         (global%var_list, var_str ("?epa_recoil"), .false., &
          intrinsic=.true.)
    call var_list_append_real &
         (global%var_list, var_str ("ewa_x_min"), 0._default, &
          intrinsic=.true.)
    call var_list_append_real &
         (global%var_list, var_str ("ewa_pt_max"), 0._default, &
          intrinsic=.true.)
    call var_list_append_real &
         (global%var_list, var_str ("ewa_mass"), 0._default, &
          intrinsic=.true.)
    call var_list_append_log &
         (global%var_list, var_str ("?ewa_keep_momentum"), .false., &
          intrinsic=.true.)       
    call var_list_append_log &
         (global%var_list, var_str ("?ewa_keep_energy"), .false., &
          intrinsic=.true.)               
    call var_list_append_log &
         (global%var_list, var_str ("?circe1_photon1"), .false., &
          intrinsic=.true.)     
    call var_list_append_log &
         (global%var_list, var_str ("?circe1_photon2"), .false., &
          intrinsic=.true.)     
    call var_list_append_real &
         (global%var_list, var_str ("circe1_sqrts"), &
          intrinsic=.true.)       
    call var_list_append_log &
         (global%var_list, var_str ("?circe1_generate"), .true., &
          intrinsic=.true.)               
    call var_list_append_log &
         (global%var_list, var_str ("?circe1_map"), .true., &
          intrinsic=.true.)     
    call var_list_append_real &
         (global%var_list, var_str ("circe1_mapping_slope"), 2._default, &
          intrinsic=.true.)     
    call var_list_append_real &
         (global%var_list, var_str ("circe1_eps"), 1e-5_default, &
          intrinsic=.true.)               
    call var_list_append_int &
         (global%var_list, var_str ("circe1_ver"), 0, intrinsic=.true.)
    call var_list_append_int &
         (global%var_list, var_str ("circe1_rev"), 0, intrinsic=.true.)
    call var_list_append_string &
         (global%var_list, var_str ("$circe1_acc"), var_str ("SBAND"), &
          intrinsic=.true.)
    call var_list_append_int &
         (global%var_list, var_str ("circe1_chat"), 0, intrinsic=.true.)
    call var_list_append_log &
         (global%var_list, var_str ("?circe2_polarized"), .true., &
          intrinsic=.true.)               
    call var_list_append_string &    
         (global%var_list, var_str ("$circe2_file"), &
          intrinsic=.true.)      
    call var_list_append_string &    
         (global%var_list, var_str ("$circe2_design"), var_str ("*"), &
          intrinsic=.true.)               
    call var_list_append_string &    
         (global%var_list, var_str ("$beam_events_file"), &
          intrinsic=.true.)      
    call var_list_append_log &
         (global%var_list, var_str ("?beam_events_warn_eof"), .true., &
          intrinsic=.true.)
    call var_list_append_log &
         (global%var_list, var_str ("?energy_scan_normalize"), .false., &
          intrinsic=.true.)
    call var_list_append_log &
         (global%var_list, var_str ("?alpha_s_is_fixed"), .true., &
          intrinsic=.true.)
    call var_list_append_log &
         (global%var_list, var_str ("?alpha_s_from_lhapdf"), .false., &
          intrinsic=.true.)
    call var_list_append_log &
         (global%var_list, var_str ("?alpha_s_from_pdf_builtin"), .false., &
          intrinsic=.true.)
    call var_list_append_int &
         (global%var_list, var_str ("alpha_s_order"), 0, &
          intrinsic=.true.)
    call var_list_append_int &
         (global%var_list, var_str ("alpha_s_nf"), 5, &
          intrinsic=.true.)
    call var_list_append_log &
         (global%var_list, var_str ("?alpha_s_from_mz"), .false., &
          intrinsic=.true.)
    call var_list_append_log &
         (global%var_list, var_str ("?alpha_s_from_lambda_qcd"), .false., &
          intrinsic=.true.)
    call var_list_append_real &
         (global%var_list, var_str ("lambda_qcd"), 200.e-3_default, &
          intrinsic=.true.)
    call var_list_append_log &
         (global%var_list, var_str ("?fatal_beam_decay"), .true., &
          intrinsic=.true.)          
    call var_list_append_log &
         (global%var_list, var_str ("?helicity_selection_active"), .true., &
          intrinsic=.true.)
    call var_list_append_real &
         (global%var_list, var_str ("helicity_selection_threshold"), &
          1E10_default, &
          intrinsic=.true.)
    call var_list_append_int &
         (global%var_list, var_str ("helicity_selection_cutoff"), 1000, &
          intrinsic=.true.)
    call var_list_append_string &
         (global%var_list, var_str ("$rng_method"), var_str ("tao"), &
          intrinsic=.true.)
    call var_list_append_string &
         (global%var_list, var_str ("$integration_method"), var_str ("vamp"), &
          intrinsic=.true.)
    call var_list_append_int &
         (global%var_list, var_str ("threshold_calls"), 10, &
          intrinsic=.true.)
    call var_list_append_int &
         (global%var_list, var_str ("min_calls_per_channel"), 10, &
          intrinsic=.true.)
    call var_list_append_int &
         (global%var_list, var_str ("min_calls_per_bin"), 10, &
          intrinsic=.true.)
    call var_list_append_int &
         (global%var_list, var_str ("min_bins"), 3, &
          intrinsic=.true.)
    call var_list_append_int &
         (global%var_list, var_str ("max_bins"), 20, &
          intrinsic=.true.)
    call var_list_append_log &
         (global%var_list, var_str ("?stratified"), .true., &
          intrinsic=.true.)
    call var_list_append_log &
         (global%var_list, var_str ("?use_vamp_equivalences"), .true., &
          intrinsic=.true.)
    call var_list_append_log &
         (global%var_list, var_str ("?vamp_verbose"), .false., &
         intrinsic=.true.)
    call var_list_append_log &
         (global%var_list, var_str ("?vamp_history_global"), &
         .true., intrinsic=.true.)
    call var_list_append_log &
         (global%var_list, var_str ("?vamp_history_global_verbose"), &
         .false., intrinsic=.true.)
    call var_list_append_log &
         (global%var_list, var_str ("?vamp_history_channels"), &
         .false., intrinsic=.true.)
    call var_list_append_log &
         (global%var_list, var_str ("?vamp_history_channels_verbose"), &
         .false., intrinsic=.true.)
    call var_list_append_real &
         (global%var_list, var_str ("channel_weights_power"), 0.25_default, &
          intrinsic=.true.)
    call var_list_append_string &
         (global%var_list, var_str ("$phs_method"), var_str ("default"), &
          intrinsic=.true.)
    call var_list_append_log &
         (global%var_list, var_str ("?vis_channels"), .false., &
          intrinsic=.true.)       
    call var_list_append_log &
         (global%var_list, var_str ("?check_phs_file"), .true., &
          intrinsic=.true.)
    call var_list_append_string &
         (global%var_list, var_str ("$phs_file"), var_str (""), &
          intrinsic=.true.)
    call var_list_append_log &
         (global%var_list, var_str ("?phs_only"), .false., &
          intrinsic=.true.)
    call var_list_append_real &
         (global%var_list, var_str ("phs_threshold_s"), 50._default, &
          intrinsic=.true.)
    call var_list_append_real &
         (global%var_list, var_str ("phs_threshold_t"), 100._default, &
          intrinsic=.true.)
    call var_list_append_int &
         (global%var_list, var_str ("phs_off_shell"), 2, &
          intrinsic=.true.)
    call var_list_append_int &
         (global%var_list, var_str ("phs_t_channel"), 6, &
          intrinsic=.true.)
    call var_list_append_real &
         (global%var_list, var_str ("phs_e_scale"), 10._default, &
          intrinsic=.true.)
    call var_list_append_real &
         (global%var_list, var_str ("phs_m_scale"), 10._default, &
          intrinsic=.true.)
    call var_list_append_real &
         (global%var_list, var_str ("phs_q_scale"), 10._default, &
          intrinsic=.true.)
    call var_list_append_log &
         (global%var_list, var_str ("?phs_keep_nonresonant"), .true., &
          intrinsic=.true.)
    call var_list_append_log &
         (global%var_list, var_str ("?phs_step_mapping"), .true., &
          intrinsic=.true.)
    call var_list_append_log &
         (global%var_list, var_str ("?phs_step_mapping_exp"), .true., &
          intrinsic=.true.)
    call var_list_append_log &
         (global%var_list, var_str ("?phs_s_mapping"), .true., &
          intrinsic=.true.)
    call var_list_append_string &
         (global%var_list, var_str ("$run_id"), var_str (""), &
          intrinsic=.true.)
    call var_list_append_int &
         (global%var_list, var_str ("n_calls_test"), 0, &
          intrinsic=.true.)
    call var_list_append_log &
         (global%var_list, var_str ("?integration_timer"), .true., &
          intrinsic=.true.)
    call var_list_append_log &
         (global%var_list, var_str ("?check_grid_file"), .true., &
          intrinsic=.true.)
    call var_list_append_real &
         (global%var_list, var_str ("accuracy_goal"), 0._default, &
          intrinsic=.true.)
    call var_list_append_real &
         (global%var_list, var_str ("error_goal"), 0._default, &
          intrinsic=.true.)
    call var_list_append_real &
         (global%var_list, var_str ("relative_error_goal"), 0._default, &
          intrinsic=.true.)
    call var_list_append_real &
         (global%var_list, var_str ("error_threshold"), &
         0._default, intrinsic=.true.)
    call var_list_append_log &
         (global%var_list, var_str ("?vis_history"), .true., &
          intrinsic=.true.)       
    call var_list_append_log &
         (global%var_list, var_str ("?diags"), .false., &
          intrinsic=.true.)
    call var_list_append_log &
         (global%var_list, var_str ("?diags_color"), .false., &
          intrinsic=.true.)    
    call var_list_append_log &
         (global%var_list, var_str ("?check_event_file"), .true., &
          intrinsic=.true.)
    call var_list_append_string &
         (global%var_list, var_str ("$event_file_version"), var_str (""), &
          intrinsic=.true.)
    call var_list_append_int &
         (global%var_list, var_str ("n_events"), 0, &
          intrinsic=.true.)
    call var_list_append_log &
         (global%var_list, var_str ("?unweighted"), .true., &
          intrinsic=.true.)       
    call var_list_append_real &
         (global%var_list, var_str ("safety_factor"), 1._default, &
          intrinsic=.true.)       
    call var_list_append_log &
         (global%var_list, var_str ("?negative_weights"), .false., &
          intrinsic=.true.)       
    call var_list_append_log &
         (global%var_list, var_str ("?keep_beams"), .false., &
          intrinsic=.true.)
    call var_list_append_log &
         (global%var_list, var_str ("?keep_remnants"), .true., &
          intrinsic=.true.)
    call var_list_append_log &
         (global%var_list, var_str ("?recover_beams"), .true., &
          intrinsic=.true.)
    call var_list_append_log &
         (global%var_list, var_str ("?update_event"), .false., &
          intrinsic=.true.)       
    call var_list_append_log &
         (global%var_list, var_str ("?update_sqme"), .false., &
          intrinsic=.true.)       
    call var_list_append_log &
         (global%var_list, var_str ("?update_weight"), .false., &
          intrinsic=.true.)       
    call var_list_append_log &
         (global%var_list, var_str ("?use_alpha_s_from_file"), .false., &
          intrinsic=.true.)       
    call var_list_append_log &
         (global%var_list, var_str ("?use_scale_from_file"), .false., &
          intrinsic=.true.)       
    call var_list_append_log &
         (global%var_list, var_str ("?allow_decays"), .true., &
          intrinsic=.true.)       
    call var_list_append_log &
         (global%var_list, var_str ("?auto_decays"), .false., &
          intrinsic=.true.)       
    call var_list_append_int &
         (global%var_list, var_str ("auto_decays_multiplicity"), 2, &
          intrinsic=.true.)       
    call var_list_append_log &
         (global%var_list, var_str ("?auto_decays_radiative"), .false., &
          intrinsic=.true.)       
    call var_list_append_log &
         (global%var_list, var_str ("?decay_rest_frame"), .false., &
          intrinsic=.true.)       
    call var_list_append_log &
         (global%var_list, var_str ("?isotropic_decay"), .false., &
          intrinsic=.true.)
    call var_list_append_log &
         (global%var_list, var_str ("?diagonal_decay"), .false., &
          intrinsic=.true.)
    call var_list_append_int &
         (global%var_list, var_str ("decay_helicity"), &
          intrinsic=.true.)
    call var_list_append_string &
         (global%var_list, var_str ("$sample"), var_str (""), &
          intrinsic=.true.)
    call var_list_append_string &
         (global%var_list, var_str ("$sample_normalization"), var_str ("auto"),&
          intrinsic=.true.)       
    call var_list_append_log &
         (global%var_list, var_str ("?sample_pacify"), .false., &
          intrinsic=.true.)
    call var_list_append_log &
         (global%var_list, var_str ("?sample_select"), .true., &
          intrinsic=.true.)
    call var_list_append_int &
         (global%var_list, var_str ("sample_max_tries"), 10000, &
         intrinsic = .true.)
    call var_list_append_int &
         (global%var_list, var_str ("sample_split_n_evt"), 0, &
         intrinsic = .true.)
    call var_list_append_int &
         (global%var_list, var_str ("sample_split_index"), 0, &
         intrinsic = .true.)
    call var_list_append_string &
         (global%var_list, var_str ("$rescan_input_format"), var_str ("raw"), &
          intrinsic=.true.)
    call var_list_append_log &
         (global%var_list, var_str ("?read_raw"), .true., &
          intrinsic=.true.)
    call var_list_append_log &
         (global%var_list, var_str ("?write_raw"), .true., &
          intrinsic=.true.)
    call var_list_append_string &
         (global%var_list, var_str ("$extension_raw"), var_str ("evx"), &
         intrinsic=.true.)
    call var_list_append_string &
         (global%var_list, var_str ("$extension_default"), var_str ("evt"), &
         intrinsic=.true.)
    call var_list_append_string &
         (global%var_list, var_str ("$debug_extension"), var_str ("debug"), &
          intrinsic=.true.)
    call var_list_append_log &
         (global%var_list, var_str ("?debug_process"), .true., &
          intrinsic=.true.)
    call var_list_append_log &
         (global%var_list, var_str ("?debug_transforms"), .true., &
          intrinsic=.true.)
    call var_list_append_log &
         (global%var_list, var_str ("?debug_decay"), .true., &
          intrinsic=.true.)
    call var_list_append_log &
         (global%var_list, var_str ("?debug_verbose"), .true., &
          intrinsic=.true.)
    call var_list_append_string &
         (global%var_list, var_str ("$extension_hepevt"), var_str ("hepevt"), &
          intrinsic=.true.)
    call var_list_append_string &
         (global%var_list, var_str ("$extension_ascii_short"), &
          var_str ("short.evt"), intrinsic=.true.)
    call var_list_append_string &
         (global%var_list, var_str ("$extension_ascii_long"), &
          var_str ("long.evt"), intrinsic=.true.)        
    call var_list_append_string &
         (global%var_list, var_str ("$extension_athena"), &
          var_str ("athena.evt"), intrinsic=.true.) 
    call var_list_append_string &
          (global%var_list, var_str ("$extension_mokka"), &
           var_str ("mokka.evt"), intrinsic=.true.)       
    call var_list_append_string &
         (global%var_list, var_str ("$lhef_version"), var_str ("2.0"), &
         intrinsic = .true.)
    call var_list_append_string &
         (global%var_list, var_str ("$lhef_extension"), var_str ("lhe"), &
          intrinsic=.true.)
    call var_list_append_log &
         (global%var_list, var_str ("?lhef_write_sqme_prc"), .true., &
         intrinsic = .true.)
    call var_list_append_log &
         (global%var_list, var_str ("?lhef_write_sqme_ref"), .false., &
         intrinsic = .true.)
    call var_list_append_log &
         (global%var_list, var_str ("?lhef_write_sqme_alt"), .true., &
         intrinsic = .true.)
    call var_list_append_string &
         (global%var_list, var_str ("$extension_lha"), var_str ("lha"), &
          intrinsic=.true.)
    call var_list_append_string &
         (global%var_list, var_str ("$extension_hepmc"), var_str ("hepmc"), &
          intrinsic=.true.)
    call var_list_append_string &
         (global%var_list, var_str ("$extension_lcio"), var_str ("slcio"), &
          intrinsic=.true.)    
    call var_list_append_string &
         (global%var_list, var_str ("$extension_stdhep"), var_str ("hep"), &
          intrinsic=.true.)
    call var_list_append_string &
         (global%var_list, var_str ("$extension_stdhep_up"), &
          var_str ("up.hep"), intrinsic=.true.)
    call var_list_append_string &
         (global%var_list, var_str ("$extension_hepevt_verb"), &
          var_str ("hepevt.verb"), intrinsic=.true.)
    call var_list_append_string &
         (global%var_list, var_str ("$extension_lha_verb"), &
          var_str ("lha.verb"), intrinsic=.true.)
    call var_list_append_int (global%var_list, &
         var_str ("n_bins"), 20, &
         intrinsic=.true.)
    call var_list_append_log (global%var_list, &
         var_str ("?normalize_bins"), .false., &
         intrinsic=.true.)
    call var_list_append_string (global%var_list, &
         var_str ("$obs_label"), var_str (""), &
         intrinsic=.true.)
    call var_list_append_string (global%var_list, &
         var_str ("$obs_unit"), var_str (""), &
          intrinsic=.true.)
    call var_list_append_string (global%var_list, &
         var_str ("$title"), var_str (""), &
          intrinsic=.true.)
    call var_list_append_string (global%var_list, &
         var_str ("$description"), var_str (""), &
          intrinsic=.true.)
    call var_list_append_string (global%var_list, &
         var_str ("$x_label"), var_str (""), &
          intrinsic=.true.)
    call var_list_append_string (global%var_list, &
         var_str ("$y_label"), var_str (""), &
          intrinsic=.true.)
    call var_list_append_int &
         (global%var_list, var_str ("graph_width_mm"), 130, &
          intrinsic=.true.)
    call var_list_append_int &
         (global%var_list, var_str ("graph_height_mm"), 90, &
          intrinsic=.true.)
    call var_list_append_log &
         (global%var_list, var_str ("?y_log"), .false., &
          intrinsic=.true.)
    call var_list_append_log &
         (global%var_list, var_str ("?x_log"), .false., &
          intrinsic=.true.)
    call var_list_append_real &
         (global%var_list, var_str ("x_min"),  &
          intrinsic=.true.)
    call var_list_append_real &
         (global%var_list, var_str ("x_max"),  &
          intrinsic=.true.)
    call var_list_append_real &
         (global%var_list, var_str ("y_min"),  &
          intrinsic=.true.)
    call var_list_append_real &
         (global%var_list, var_str ("y_max"),  &
          intrinsic=.true.)
    call var_list_append_string &
         (global%var_list, var_str ("$gmlcode_bg"), var_str (""), &
          intrinsic=.true.)
    call var_list_append_string &
         (global%var_list, var_str ("$gmlcode_fg"), var_str (""), &
          intrinsic=.true.)
    call var_list_append_log &
         (global%var_list, var_str ("?draw_histogram"), &
          intrinsic=.true.)
    call var_list_append_log &
         (global%var_list, var_str ("?draw_base"), &
          intrinsic=.true.)
    call var_list_append_log &
         (global%var_list, var_str ("?draw_piecewise"), &
          intrinsic=.true.)
    call var_list_append_log &
         (global%var_list, var_str ("?fill_curve"), &
          intrinsic=.true.)
    call var_list_append_log &
         (global%var_list, var_str ("?draw_curve"), &
          intrinsic=.true.)
    call var_list_append_log &
         (global%var_list, var_str ("?draw_errors"), &
          intrinsic=.true.)
    call var_list_append_log &
         (global%var_list, var_str ("?draw_symbols"), &
          intrinsic=.true.)
    call var_list_append_string &
         (global%var_list, var_str ("$fill_options"), &
          intrinsic=.true.)
    call var_list_append_string &
         (global%var_list, var_str ("$draw_options"), &
          intrinsic=.true.)
    call var_list_append_string &
         (global%var_list, var_str ("$err_options"), &
          intrinsic=.true.)
    call var_list_append_string &
         (global%var_list, var_str ("$symbol"), &
          intrinsic=.true.)
    call var_list_append_real (global%var_list, &
         var_str ("tolerance"), 0._default, &
          intrinsic=.true.)
    call var_list_append_int (global%var_list, &
         var_str ("checkpoint"), 0, &
         intrinsic = .true.)
    call var_list_append_log &
         (global%var_list, var_str ("?pacify"), .false., &
         intrinsic=.true.)
    call var_list_append_string &
         (global%var_list, var_str ("$out_file"), var_str (""), &
          intrinsic=.true.)
    call var_list_append_log &
         (global%var_list, var_str ("?out_advance"), .true., &
          intrinsic=.true.)
!!! JRR: WK please check (#542)    
!     call var_list_append_log &
!          (global%var_list, var_str ("?out_custom"), .false., &
!           intrinsic=.true.)
!     call var_list_append_string &
!          (global%var_list, var_str ("$out_comment"), var_str ("# "), &
!           intrinsic=.true.)
!     call var_list_append_log &
!          (global%var_list, var_str ("?out_header"), .true., &
!           intrinsic=.true.)
!     call var_list_append_log &
!          (global%var_list, var_str ("?out_yerr"), .true., &
!           intrinsic=.true.)
!     call var_list_append_log &
!          (global%var_list, var_str ("?out_xerr"), .true., &
!           intrinsic=.true.)
    call var_list_append_int (global%var_list, var_str ("real_range"), &
         range (real_specimen), intrinsic = .true., locked = .true.)
    call var_list_append_int (global%var_list, var_str ("real_precision"), &
         precision (real_specimen), intrinsic = .true., locked = .true.)
    call var_list_append_real (global%var_list, var_str ("real_epsilon"), &
         epsilon (real_specimen), intrinsic = .true., locked = .true.)
    call var_list_append_real (global%var_list, var_str ("real_tiny"), &
         tiny (real_specimen), intrinsic = .true., locked = .true.)
    !!! FastJet parameters
    call var_list_append_int (global%var_list, &
         var_str ("kt_algorithm"), &
         kt_algorithm, &
         intrinsic = .true., locked = .true.)
    call var_list_append_int (global%var_list, &
         var_str ("cambridge_algorithm"), &
         cambridge_algorithm, intrinsic = .true., locked = .true.)
    call var_list_append_int (global%var_list, &
         var_str ("antikt_algorithm"), &
         antikt_algorithm, &
         intrinsic = .true., locked = .true.)
    call var_list_append_int (global%var_list, &
         var_str ("genkt_algorithm"), &
         genkt_algorithm, &
         intrinsic = .true., locked = .true.)
    call var_list_append_int (global%var_list, &
         var_str ("cambridge_for_passive_algorithm"), &
         cambridge_for_passive_algorithm, &
         intrinsic = .true., locked = .true.)
    call var_list_append_int (global%var_list, &
         var_str ("genkt_for_passive_algorithm"), &
         genkt_for_passive_algorithm, &
         intrinsic = .true., locked = .true.)
    call var_list_append_int (global%var_list, &
         var_str ("ee_kt_algorithm"), &
         ee_kt_algorithm, &
         intrinsic = .true., locked = .true.)
    call var_list_append_int (global%var_list, &
         var_str ("ee_genkt_algorithm"), &
         ee_genkt_algorithm, &
         intrinsic = .true., locked = .true.)
    call var_list_append_int (global%var_list, &
         var_str ("plugin_algorithm"), &
         plugin_algorithm, &
         intrinsic = .true., locked = .true.)
    call var_list_append_int (global%var_list, &
         var_str ("undefined_jet_algorithm"), &
         undefined_jet_algorithm, &
         intrinsic = .true., locked = .true.)
    call var_list_append_int (global%var_list, &
         var_str ("jet_algorithm"), undefined_jet_algorithm, &
         intrinsic = .true.)
    call var_list_append_real (global%var_list, &
         var_str ("jet_r"), 0._default, &
         intrinsic = .true.)
    call var_list_append_log &
         (global%var_list, var_str ("?polarized_events"), .false., &
            intrinsic=.true.)
    !!! Default settings for shower
    call var_list_append_log &
         (global%var_list, var_str ("?allow_shower"), .true., &
            intrinsic=.true.)
    call var_list_append_log &
         (global%var_list, var_str ("?ps_fsr_active"), .false., &
            intrinsic=.true.)
    call var_list_append_log &
         (global%var_list, var_str ("?ps_use_PYTHIA_shower"), .false., &
            intrinsic=.true.)
    call var_list_append_log &
         (global%var_list, var_str ("?ps_PYTHIA_verbose"), .false., &
            intrinsic=.true.)
    call var_list_append_string &
         (global%var_list, var_str ("$ps_PYTHIA_PYGIVE"), var_str (""), &
          intrinsic=.true.)
    call var_list_append_log &
         (global%var_list, var_str ("?ps_isr_active"), .false., &
            intrinsic=.true.)
    call var_list_append_real (global%var_list, &
         var_str ("ps_mass_cutoff"), 1._default, intrinsic = .true.)
    call var_list_append_real (global%var_list, &
         var_str ("ps_fsr_lambda"), 0.29_default, intrinsic = .true.)
    call var_list_append_real (global%var_list, &
         var_str ("ps_isr_lambda"), 0.29_default, intrinsic = .true.)
    call var_list_append_int (global%var_list, &
         var_str ("ps_max_n_flavors"), 5, intrinsic = .true.)
    call var_list_append_log &
         (global%var_list, var_str ("?ps_isr_alpha_s_running"), .true., &
            intrinsic=.true.)
    call var_list_append_log &
         (global%var_list, var_str ("?ps_fsr_alpha_s_running"), .true., &
            intrinsic=.true.)
    call var_list_append_real (global%var_list, var_str ("ps_fixed_alpha_s"), &
         0._default, intrinsic = .true.)
    call var_list_append_log &
         (global%var_list, var_str ("?ps_isr_pt_ordered"), .false., &
            intrinsic=.true.)
    call var_list_append_log &
         (global%var_list, var_str ("?ps_isr_angular_ordered"), .true., &
            intrinsic=.true.)
    call var_list_append_real (global%var_list, var_str &
         ("ps_isr_primordial_kt_width"), 0._default, intrinsic = .true.)
    call var_list_append_real (global%var_list, var_str &
         ("ps_isr_primordial_kt_cutoff"), 5._default, intrinsic = .true.)
    call var_list_append_real (global%var_list, var_str &
         ("ps_isr_z_cutoff"), 0.999_default, intrinsic = .true.)
    call var_list_append_real (global%var_list, var_str &
         ("ps_isr_minenergy"), 1._default, intrinsic = .true.)
    call var_list_append_real (global%var_list, var_str &
         ("ps_isr_tscalefactor"), 1._default, intrinsic = .true.)
    call var_list_append_log (global%var_list, var_str &
         ("?ps_isr_only_onshell_emitted_partons"), .false., intrinsic=.true.)
    !!! Default settings for hadronization
    call var_list_append_log &
         (global%var_list, var_str ("?hadronization_active"), .false., &
            intrinsic=.true.)
    !!! Setting for mlm matching
    call var_list_append_log &
         (global%var_list, var_str ("?mlm_matching"), .false., &
            intrinsic=.true.)
    call var_list_append_real (global%var_list, var_str &
         ("mlm_Qcut_ME"), 0._default, intrinsic = .true.)
    call var_list_append_real (global%var_list, var_str &
         ("mlm_Qcut_PS"), 0._default, intrinsic = .true.)
    call var_list_append_real (global%var_list, var_str &
         ("mlm_ptmin"), 0._default, intrinsic = .true.)
    call var_list_append_real (global%var_list, var_str &
         ("mlm_etamax"), 0._default, intrinsic = .true.)
    call var_list_append_real (global%var_list, var_str &
         ("mlm_Rmin"), 0._default, intrinsic = .true.)
    call var_list_append_real (global%var_list, var_str &
         ("mlm_Emin"), 0._default, intrinsic = .true.)
    call var_list_append_int (global%var_list, var_str &
         ("mlm_nmaxMEjets"), 0, intrinsic = .true.)
    call var_list_append_real (global%var_list, var_str &
         ("mlm_ETclusfactor"), 0.2_default, intrinsic = .true.)
    call var_list_append_real (global%var_list, var_str &
         ("mlm_ETclusminE"), 5._default, intrinsic = .true.)
    call var_list_append_real (global%var_list, var_str &
         ("mlm_etaclusfactor"), 1._default, intrinsic = .true.)
    call var_list_append_real (global%var_list, var_str &
         ("mlm_Rclusfactor"), 1._default, intrinsic = .true.)
    call var_list_append_real (global%var_list, var_str &
         ("mlm_Eclusfactor"), 1._default, intrinsic = .true.)
    call var_list_append_log &
         (global%var_list, var_str ("?ckkw_matching"), .false., &
            intrinsic=.true.)
    call var_list_append_log &
         (global%var_list, var_str ("?muli_active"), .false., &
            intrinsic=.true.)
    call var_list_append_string &
         (global%var_list, var_str ("$pdf_builtin_set"), var_str ("CTEQ6L"), &
         intrinsic=.true.)
    call var_list_append_log &
         (global%var_list, var_str ("?omega_openmp"), &
         openmp_is_active (), &
         intrinsic=.true.)
    call var_list_append_log &
         (global%var_list, var_str ("?openmp_is_active"), &
         openmp_is_active (), &
         locked=.true., intrinsic=.true.)
    call var_list_append_int &
         (global%var_list, var_str ("openmp_num_threads_default"), &
         openmp_get_default_max_threads (), &
         locked=.true., intrinsic=.true.)
    call var_list_append_int &
         (global%var_list, var_str ("openmp_num_threads"), &        
         openmp_get_max_threads (), &
         intrinsic=.true.)
    call var_list_append_log &
         (global%var_list, var_str ("?openmp_logging"), &
         .true., intrinsic=.true.)    
    call var_list_append_log &
        (global%var_list, var_str ("?use_gosam_loops"), &
         .true., intrinsic = .true.)
    call var_list_append_log &
        (global%var_list, var_str ("?use_gosam_correlations"), &
         .false., intrinsic = .true.)
    call var_list_append_log &
        (global%var_list, var_str ("?use_gosam_real_trees"), &
         .false., intrinsic = .true.)
    call var_list_append_real &
        (global%var_list, var_str ("fks_dij_exp1"), &
         1._default, intrinsic = .true.)
    call var_list_append_real &
        (global%var_list, var_str ("fks_dij_exp2"), &
         1._default, intrinsic = .true.)
    call var_list_append_int &
        (global%var_list, var_str ("fks_mapping_type"), &
         1, intrinsic = .true.)
    call var_list_append_int &
        (global%var_list, var_str ("alpha_power"), &
         2, intrinsic = .true.)
    call var_list_append_int &
        (global%var_list, var_str ("alphas_power"), &
         0, intrinsic = .true.)
    call var_list_append_log &
        (global%var_list, var_str ("?combined_nlo_integration"), &
         .false., intrinsic = .true.)
    call global%init_pointer_variables ()
    call global%process_stack%init_var_list (global%var_list)
  end subroutine rt_data_global_init

\eatline
\nwendcode{}\nwbegindocs{1706}\nwdocspar
\subsection{Local copies}
This is done at compile time when a local copy of runtime data is
needed: Link the variable list and initialize all derived parameters.
This allows for synchronizing them with local variable changes without
affecting global data.

Also re-initialize pointer variables, so they point to local copies of
their targets.
\nwenddocs{}\nwbegincode{1707}\moddef{RT data: rt data: TBP}\plusendmoddef
  procedure :: local_init => rt_data_local_init
\nwendcode{}\nwbegincode{1708}\moddef{RT data: procedures}\plusendmoddef
  subroutine rt_data_local_init (local, global, env)
    class(rt_data_t), intent(inout), target :: local
    type(rt_data_t), intent(in), target :: global
    integer, intent(in), optional :: env
    local%context => global
    call local%process_stack%link (global%process_stack)
    call local%process_stack%init_var_list (local%var_list)
    call local%process_stack%link_var_list (global%var_list)
    call var_list_append_string &
         (local%var_list, var_str ("$model_name"), var_str (""), &
          intrinsic=.true.)
    call local%init_pointer_variables ()
    local%fallback_model => global%fallback_model
    local%radiation_model => global%radiation_model
    local%os_data = global%os_data
    local%logfile = global%logfile
    call local%model_list%link (global%model_list)
    local%model => global%model
    if (associated (local%model)) then
       call local%model%link_var_list (local%var_list)
    end if
  end subroutine rt_data_local_init

\eatline
\nwendcode{}\nwbegindocs{1709}These variables point to objects which get local copies:
\nwenddocs{}\nwbegincode{1710}\moddef{RT data: rt data: TBP}\plusendmoddef
  procedure :: init_pointer_variables => rt_data_init_pointer_variables
\nwendcode{}\nwbegincode{1711}\moddef{RT data: procedures}\plusendmoddef
  subroutine rt_data_init_pointer_variables (local)
    class(rt_data_t), intent(inout), target :: local
    logical, target, save :: known = .true.
    call var_list_append_string_ptr &
         (local%var_list, var_str ("$fc"), local%os_data%fc, known, &
          intrinsic=.true.)
    call var_list_append_string_ptr &
         (local%var_list, var_str ("$fcflags"), local%os_data%fcflags, known, &
         intrinsic=.true.)
  end subroutine rt_data_init_pointer_variables

\eatline
\nwendcode{}\nwbegindocs{1712}This is done at execution time: Copy data, transfer pointers.
{\tt{}local} has intent(inout) because its local variable list has
already been prepared by the previous routine.

To be pedantic, the local pointers to model and library should point
to the entries in the local copies.  (However, as long as these are
just shallow copies with identical content, this is actually
irrelevant.)

The process library and process stacks behave as global objects.  The
copies of the process library and process stacks should be shallow
copies, so the contents stay identical.  Since objects may be pushed
on the stack in the local environment, upon restoring the global
environment, we should reverse the assignment.  Then the added stack
elements will end up on the global stack.  (This should be
reconsidered in a parallel environment.)
\nwenddocs{}\nwbegincode{1713}\moddef{RT data: rt data: TBP}\plusendmoddef
  procedure :: activate => rt_data_activate
\nwendcode{}\nwbegincode{1714}\moddef{RT data: procedures}\plusendmoddef
  subroutine rt_data_activate (local)
    class(rt_data_t), intent(inout), target :: local
    class(rt_data_t), pointer :: global
    global => local%context
    if (associated (global)) then
       local%lexer => global%lexer
       call global%copy_globals (local)
       local%os_data = global%os_data
       local%logfile = global%logfile
       if (associated (global%prclib)) then
          local%prclib => &
               local%prclib_stack%get_library_ptr (global%prclib%get_name ())
       end if
       call local%import_values ()
       call local%process_stack%link (global%process_stack)
       local%it_list = global%it_list
       local%beam_structure = global%beam_structure
       local%pn = global%pn
       if (allocated (local%sample_fmt))  deallocate (local%sample_fmt)
       if (allocated (global%sample_fmt)) then
          allocate (local%sample_fmt (size (global%sample_fmt)), &
               source = global%sample_fmt)
       end if
       local%out_files => global%out_files
       local%model => global%model
       local%model_is_copy = .false.
    else if (.not. associated (local%model)) then
       local%model => local%preload_model
       local%model_is_copy = .false.
    end if
    if (associated (local%model)) then
       call local%model%link_var_list (local%var_list)
       call var_list_set_string (local%var_list, var_str ("$model_name"), &
            local%model%get_name (), is_known = .true.)
    else
       call var_list_set_string (local%var_list, var_str ("$model_name"), &
            var_str (""), is_known = .false.)
    end if
  end subroutine rt_data_activate

\eatline
\nwendcode{}\nwbegindocs{1715}Restore the previous state of data, without actually finalizing the local
environment.  We also clear the local process stack.  Some local modifications
(model list and process library stack) are communicated to the global context,
if there is any. 

If the {\tt{}keep{\char95}local} flag is set, we want to retain current settings in
the local environment.  In particular, we create an instance of the currently
selected model (which thus becomes separated from the model library!).
The local variables are also kept.
\nwenddocs{}\nwbegincode{1716}\moddef{RT data: rt data: TBP}\plusendmoddef
  procedure :: deactivate => rt_data_deactivate
\nwendcode{}\nwbegincode{1717}\moddef{RT data: procedures}\plusendmoddef
  subroutine rt_data_deactivate (local, global, keep_local)
    class(rt_data_t), intent(inout), target :: local
    class(rt_data_t), intent(inout), optional, target :: global
    logical, intent(in), optional :: keep_local
    type(string_t) :: global_model, local_model
    logical :: same_model, delete
    delete = .true.;  if (present (keep_local))  delete = .not. keep_local
    if (present (global)) then
       if (associated (global%model) .and. associated (local%model)) then 
          global_model = global%model%get_name ()
          local_model = local%model%get_name ()
          same_model = global_model == local_model
       else
          same_model = .false.
       end if
       if (delete) then
          call local%process_stack%clear ()
          call local%unselect_model ()
          call local%unset_values ()
       else if (associated (local%model)) then
          call local%ensure_model_copy ()
       end if
       if (.not. same_model .and. global_model /= "") then
          call msg_message ("Restoring model '" // char (global_model) // "'")
       end if
       if (associated (global%model)) then
          call global%model%link_var_list (global%var_list)
       end if
       call global%restore_globals (local)
    else
       call local%unselect_model ()
    end if
  end subroutine rt_data_deactivate

\eatline
\nwendcode{}\nwbegindocs{1718}This imports the global objects for which local modifications
should be kept.  Currently, this is only the process library stack.
\nwenddocs{}\nwbegincode{1719}\moddef{RT data: rt data: TBP}\plusendmoddef
  procedure :: copy_globals => rt_data_copy_globals
\nwendcode{}\nwbegincode{1720}\moddef{RT data: procedures}\plusendmoddef
  subroutine rt_data_copy_globals (global, local)
    class(rt_data_t), intent(in) :: global
    class(rt_data_t), intent(inout) :: local
    local%prclib_stack = global%prclib_stack
  end subroutine rt_data_copy_globals
 
\eatline
\nwendcode{}\nwbegindocs{1721}This restores global objects, for which local modifications
should be kept.
\nwenddocs{}\nwbegincode{1722}\moddef{RT data: rt data: TBP}\plusendmoddef
  procedure :: restore_globals => rt_data_restore_globals
\nwendcode{}\nwbegincode{1723}\moddef{RT data: procedures}\plusendmoddef
  subroutine rt_data_restore_globals (global, local)
    class(rt_data_t), intent(inout) :: global
    class(rt_data_t), intent(in) :: local
    global%prclib_stack = local%prclib_stack
  end subroutine rt_data_restore_globals
 
\eatline
\nwendcode{}\nwbegindocs{1724}\nwdocspar
\subsection{Finalization}
Finalizer for the variable list and the structure-function list.
This is done only for the global RT dataset; local copies contain
pointers to this and do not need a finalizer.
\nwenddocs{}\nwbegincode{1725}\moddef{RT data: rt data: TBP}\plusendmoddef
  procedure :: final => rt_data_global_final
\nwendcode{}\nwbegincode{1726}\moddef{RT data: procedures}\plusendmoddef
  subroutine rt_data_global_final (global)
    class(rt_data_t), intent(inout) :: global
    call global%process_stack%final ()
    call global%prclib_stack%final ()
!    call global%delete_model_copy ()
    call global%model_list%final ()
    call global%var_list%final (follow_link=.false.)
    if (associated (global%out_files)) then
       call file_list_final (global%out_files)
       deallocate (global%out_files)
    end if
  end subroutine rt_data_global_final

\eatline
\nwendcode{}\nwbegindocs{1727}The local copy needs a finalizer for the variable list, which consists
of local copies.  This finalizer is called only when the local
environment is finally discarded.  (Note that the process stack should
already have been cleared after execution, which can occur many times
for the same local environment.)
\nwenddocs{}\nwbegincode{1728}\moddef{RT data: rt data: TBP}\plusendmoddef
  procedure :: local_final => rt_data_local_final
\nwendcode{}\nwbegincode{1729}\moddef{RT data: procedures}\plusendmoddef
  subroutine rt_data_local_final (local)
    class(rt_data_t), intent(inout) :: local
    call local%process_stack%clear ()
!    call local%delete_model_copy ()
    call local%model_list%final ()
    call local%var_list%final (follow_link=.false.)
  end subroutine rt_data_local_final

\eatline
\nwendcode{}\nwbegindocs{1730}\nwdocspar
\subsection{Model Management}
Read a model, so it becomes available for activation.  No variables or model
copies, this is just initialization.

If this is a local environment, the model will be automatically read into the
global context.
\nwenddocs{}\nwbegincode{1731}\moddef{RT data: rt data: TBP}\plusendmoddef
  procedure :: read_model => rt_data_read_model
\nwendcode{}\nwbegincode{1732}\moddef{RT data: procedures}\plusendmoddef
  subroutine rt_data_read_model (global, name, model)
    class(rt_data_t), intent(inout) :: global
    type(string_t), intent(in) :: name
    type(model_t), pointer, intent(out) :: model
    type(string_t) :: filename
    filename = name // ".mdl"
    call global%model_list%read_model &
         (name, filename, global%os_data, model)
  end subroutine rt_data_read_model
    
\eatline
\nwendcode{}\nwbegindocs{1733}Initialize the fallback model.  This model is used
whenever the current model does not describe all physical particles
(hadrons, mainly).  It is not supposed to be modified, and the pointer
should remain linked to this model.
\nwenddocs{}\nwbegincode{1734}\moddef{RT data: rt data: TBP}\plusendmoddef
  procedure :: init_fallback_model => rt_data_init_fallback_model
\nwendcode{}\nwbegincode{1735}\moddef{RT data: procedures}\plusendmoddef
  subroutine rt_data_init_fallback_model (global, name, filename)
    class(rt_data_t), intent(inout) :: global
    type(string_t), intent(in) :: name, filename
    call global%model_list%read_model &
         (name, filename, global%os_data, global%fallback_model)
  end subroutine rt_data_init_fallback_model
  
\eatline
\nwendcode{}\nwbegindocs{1736}Initialize the radiation model.  This model is used by a
radiation-generator algorithm which is part of the NLO machinery.  It is not
supposed to be modified, and the pointer 
should remain linked to this model.
\nwenddocs{}\nwbegincode{1737}\moddef{RT data: rt data: TBP}\plusendmoddef
  procedure :: init_radiation_model => rt_data_init_radiation_model
\nwendcode{}\nwbegincode{1738}\moddef{RT data: procedures}\plusendmoddef
  subroutine rt_data_init_radiation_model (global, name, filename)
    class(rt_data_t), intent(inout) :: global
    type(string_t), intent(in) :: name, filename
    call global%model_list%read_model &
         (name, filename, global%os_data, global%radiation_model)
  end subroutine rt_data_init_radiation_model
  
\eatline
\nwendcode{}\nwbegindocs{1739}\nwdocspar
Activate a model: assign the current-model pointer and set the model name in
the variable list.  If necessary, read the model from file.  Link the global
variable list to the model variable list.
\nwenddocs{}\nwbegincode{1740}\moddef{RT data: rt data: TBP}\plusendmoddef
  procedure :: select_model => rt_data_select_model
\nwendcode{}\nwbegincode{1741}\moddef{RT data: procedures}\plusendmoddef
  subroutine rt_data_select_model (global, name)
    class(rt_data_t), intent(inout), target :: global
    type(string_t), intent(in) :: name
    logical :: same_model
    if (associated (global%model)) then
       same_model = global%model%get_name () == name
    else
       same_model = .false.
    end if
    if (.not. same_model) then
!       call global%delete_model_copy () 
       global%model => global%model_list%get_model_ptr (name)
       if (.not. associated (global%model)) then
          call global%read_model (name, global%model)
          global%model_is_copy = .false.
       else if (associated (global%context)) then
          global%model_is_copy = &
               global%model_list%model_exists (name, follow_link=.false.)
       else
          global%model_is_copy = .false.
       end if
    end if
    if (associated (global%model)) then
       call global%model%link_var_list (global%var_list)
       call var_list_set_string (global%var_list, var_str ("$model_name"), &
            name, is_known = .true.)
       call msg_message ("Switching to model '" // char (name) // "'")
    else
       call var_list_set_string (global%var_list, var_str ("$model_name"), &
            var_str (""), is_known = .false.)
    end if
  end subroutine rt_data_select_model
  
\eatline
\nwendcode{}\nwbegindocs{1742}\nwdocspar
Remove the model link and unset the model name variable.
\nwenddocs{}\nwbegincode{1743}\moddef{RT data: rt data: TBP}\plusendmoddef
  procedure :: unselect_model => rt_data_unselect_model
\nwendcode{}\nwbegincode{1744}\moddef{RT data: procedures}\plusendmoddef
  subroutine rt_data_unselect_model (global)
    class(rt_data_t), intent(inout), target :: global
    if (associated (global%model)) then
!       call global%delete_model_copy ()
       global%model => null ()
       global%model_is_copy = .false.
       call var_list_set_string (global%var_list, var_str ("$model_name"), &
            var_str (""), is_known = .false.)
    end if
  end subroutine rt_data_unselect_model
  
\eatline
\nwendcode{}\nwbegindocs{1745}\nwdocspar
Create a copy of the currently selected model and append it to the local model
list.  The model pointer is redirected to the copy.
(Not applicable for the global model list, those models will be
modified in-place.)
%Take/delete a snapshot of the currently active model, if applicable.  The
%destructor also de-assigns the model if it just a link.
\nwenddocs{}\nwbegincode{1746}\moddef{RT data: rt data: TBP}\plusendmoddef
  procedure :: ensure_model_copy => rt_data_ensure_model_copy
!   procedure :: delete_model_copy => rt_data_delete_model_copy
\nwendcode{}\nwbegincode{1747}\moddef{RT data: procedures}\plusendmoddef
  subroutine rt_data_ensure_model_copy (global)
    class(rt_data_t), intent(inout), target :: global
    if (associated (global%context)) then
       if (.not. global%model_is_copy) then
          call global%model_list%append_copy (global%model, global%model)
          global%model_is_copy = .true.
          call global%model%link_var_list (global%var_list)
       end if
    end if
  end subroutine rt_data_ensure_model_copy

!   subroutine rt_data_delete_model_copy (global)
!     class(rt_data_t), intent(inout), target :: global
!     if (global%model_is_copy) then
!        call model_pointer_delete_instance (global%model)
!        global%model_is_copy = .false.
!     else
!        global%model => null ()
!     end if
!   end subroutine rt_data_delete_model_copy

\eatline
\eatline
\nwendcode{}\nwbegindocs{1748}\nwdocspar
Modify a model variable.  The update mechanism will ensure that the model
parameter set remains consistent.  This has to take place in a local copy
of the current model.  If there is none yet, create one.
\nwenddocs{}\nwbegincode{1749}\moddef{RT data: rt data: TBP}\plusendmoddef
  procedure :: model_set_real => rt_data_model_set_real
\nwendcode{}\nwbegincode{1750}\moddef{RT data: procedures}\plusendmoddef
  subroutine rt_data_model_set_real (global, name, rval, verbose, pacified)
    class(rt_data_t), intent(inout), target :: global
    type(string_t), intent(in) :: name
    real(default), intent(in) :: rval
    logical, intent(in), optional :: verbose, pacified
    call global%ensure_model_copy ()
    call global%model%set_real (name, rval, verbose, pacified)
  end subroutine rt_data_model_set_real

\eatline
\nwendcode{}\nwbegindocs{1751}\nwdocspar
Modify particle properties.  This has to take place in a local copy
of the current model.  If there is none yet, create one.
\nwenddocs{}\nwbegincode{1752}\moddef{RT data: rt data: TBP}\plusendmoddef
  procedure :: modify_particle => rt_data_modify_particle
\nwendcode{}\nwbegincode{1753}\moddef{RT data: procedures}\plusendmoddef
  subroutine rt_data_modify_particle &
       (global, pdg, polarized, stable, decay, &
       isotropic_decay, diagonal_decay, decay_helicity)
    class(rt_data_t), intent(inout), target :: global
    integer, intent(in) :: pdg
    logical, intent(in), optional :: polarized, stable
    logical, intent(in), optional :: isotropic_decay, diagonal_decay
    integer, intent(in), optional :: decay_helicity
    type(string_t), dimension(:), intent(in), optional :: decay
    call global%ensure_model_copy ()
    if (present (polarized)) then
       if (polarized) then
          call global%model%set_polarized (pdg)
       else
          call global%model%set_unpolarized (pdg)
       end if
    end if
    if (present (stable)) then
       if (stable) then
          call global%model%set_stable (pdg)
       else if (present (decay)) then
          call global%model%set_unstable &
               (pdg, decay, isotropic_decay, diagonal_decay, decay_helicity)
       else
          call msg_bug ("Setting particle unstable: missing decay processes")
       end if
    end if
  end subroutine rt_data_modify_particle

\eatline
\nwendcode{}\nwbegindocs{1754}\nwdocspar
\subsection{Managing Variables}
Return a pointer to the currently active variable list.  If there is no model,
this is the global variable list.  If there is one, it is the model variable
list, which should be linked to the former.
\nwenddocs{}\nwbegincode{1755}\moddef{RT data: rt data: TBP}\plusendmoddef
  procedure :: get_var_list_ptr => rt_data_get_var_list_ptr
\nwendcode{}\nwbegincode{1756}\moddef{RT data: procedures}\plusendmoddef
  function rt_data_get_var_list_ptr (global) result (var_list)
    class(rt_data_t), intent(in), target :: global
    type(var_list_t), pointer :: var_list
    if (associated (global%model)) then
       var_list => global%model%get_var_list_ptr ()
    else
       var_list => global%var_list
    end if
  end function rt_data_get_var_list_ptr

\eatline
\nwendcode{}\nwbegindocs{1757}Initialize a local variable: append it to the current variable list.  No
initial value, yet.
\nwenddocs{}\nwbegincode{1758}\moddef{RT data: rt data: TBP}\plusendmoddef
  procedure :: append_log => rt_data_append_log
  procedure :: append_int => rt_data_append_int
  procedure :: append_real => rt_data_append_real
  procedure :: append_cmplx => rt_data_append_cmplx
  procedure :: append_subevt => rt_data_append_subevt
  procedure :: append_pdg_array => rt_data_append_pdg_array
  procedure :: append_string => rt_data_append_string
\nwendcode{}\nwbegincode{1759}\moddef{RT data: procedures}\plusendmoddef
  subroutine rt_data_append_log (local, name, lval, intrinsic, user)
    class(rt_data_t), intent(inout) :: local
    type(string_t), intent(in) :: name
    logical, intent(in), optional :: lval
    logical, intent(in), optional :: intrinsic, user
    call var_list_append_log (local%var_list, name, lval, &
         intrinsic = intrinsic, user = user)
  end subroutine rt_data_append_log
  
  subroutine rt_data_append_int (local, name, ival, intrinsic, user)
    class(rt_data_t), intent(inout) :: local
    type(string_t), intent(in) :: name
    integer, intent(in), optional :: ival
    logical, intent(in), optional :: intrinsic, user
    call var_list_append_int (local%var_list, name, ival, &
         intrinsic = intrinsic, user = user)
  end subroutine rt_data_append_int
  
  subroutine rt_data_append_real (local, name, rval, intrinsic, user)
    class(rt_data_t), intent(inout) :: local
    type(string_t), intent(in) :: name
    real(default), intent(in), optional :: rval
    logical, intent(in), optional :: intrinsic, user
    call var_list_append_real (local%var_list, name, rval, &
         intrinsic = intrinsic, user = user)
  end subroutine rt_data_append_real
  
  subroutine rt_data_append_cmplx (local, name, cval, intrinsic, user)
    class(rt_data_t), intent(inout) :: local
    type(string_t), intent(in) :: name
    complex(default), intent(in), optional :: cval
    logical, intent(in), optional :: intrinsic, user
    call var_list_append_cmplx (local%var_list, name, cval, &
         intrinsic = intrinsic, user = user)
  end subroutine rt_data_append_cmplx
  
  subroutine rt_data_append_subevt (local, name, pval, intrinsic, user)
    class(rt_data_t), intent(inout) :: local
    type(string_t), intent(in) :: name
    type(subevt_t), intent(in), optional :: pval
    logical, intent(in) :: intrinsic, user
    call var_list_append_subevt (local%var_list, name, &
         intrinsic = intrinsic, user = user)
  end subroutine rt_data_append_subevt
  
  subroutine rt_data_append_pdg_array (local, name, aval, intrinsic, user)
    class(rt_data_t), intent(inout) :: local
    type(string_t), intent(in) :: name
    type(pdg_array_t), intent(in), optional :: aval
    logical, intent(in), optional :: intrinsic, user
    call var_list_append_pdg_array (local%var_list, name, aval, &
         intrinsic = intrinsic, user = user)
  end subroutine rt_data_append_pdg_array
  
  subroutine rt_data_append_string (local, name, sval, intrinsic, user)
    class(rt_data_t), intent(inout) :: local
    type(string_t), intent(in) :: name
    type(string_t), intent(in), optional :: sval
    logical, intent(in), optional :: intrinsic, user
    call var_list_append_string (local%var_list, name, sval, &
         intrinsic = intrinsic, user = user)
  end subroutine rt_data_append_string
  
\eatline
\eatline
\eatline
\eatline
\eatline
\eatline
\eatline
\nwendcode{}\nwbegindocs{1760}Import values for all local variables, given a global context environment
where these variables are defined.
\nwenddocs{}\nwbegincode{1761}\moddef{RT data: rt data: TBP}\plusendmoddef
  procedure :: import_values => rt_data_import_values
\nwendcode{}\nwbegincode{1762}\moddef{RT data: procedures}\plusendmoddef
  subroutine rt_data_import_values (local)
    class(rt_data_t), intent(inout) :: local
    type(string_t) :: name
    integer :: type
    type(rt_data_t), pointer :: global
    global => local%context
    if (associated (global)) then
       call var_list_import (local%var_list, global%var_list)
    end if
  end subroutine rt_data_import_values
    
\eatline
\nwendcode{}\nwbegindocs{1763}Unset all variable values.
\nwenddocs{}\nwbegincode{1764}\moddef{RT data: rt data: TBP}\plusendmoddef
  procedure :: unset_values => rt_data_unset_values
\nwendcode{}\nwbegincode{1765}\moddef{RT data: procedures}\plusendmoddef
  subroutine rt_data_unset_values (global)
    class(rt_data_t), intent(inout) :: global
    call var_list_undefine (global%var_list, follow_link=.false.)
  end subroutine rt_data_unset_values
  
\eatline
\nwendcode{}\nwbegindocs{1766}Set a variable.  (Not a model variable, these are handled separately.)  We
can assume that the variable has been initialized.
\nwenddocs{}\nwbegincode{1767}\moddef{RT data: rt data: TBP}\plusendmoddef
  procedure :: set_log => rt_data_set_log
  procedure :: set_int => rt_data_set_int
  procedure :: set_real => rt_data_set_real
  procedure :: set_cmplx => rt_data_set_cmplx
  procedure :: set_subevt => rt_data_set_subevt
  procedure :: set_pdg_array => rt_data_set_pdg_array
  procedure :: set_string => rt_data_set_string
\nwendcode{}\nwbegincode{1768}\moddef{RT data: procedures}\plusendmoddef
  subroutine rt_data_set_log (global, name, lval, is_known, verbose)
    class(rt_data_t), intent(inout) :: global
    type(string_t), intent(in) :: name
    logical, intent(in) :: lval
    logical, intent(in) :: is_known
    logical, intent(in), optional :: verbose
    call var_list_set_log (global%var_list, name, lval, is_known, &
         verbose=verbose)
  end subroutine rt_data_set_log

  subroutine rt_data_set_int (global, name, ival, is_known, verbose)
    class(rt_data_t), intent(inout) :: global
    type(string_t), intent(in) :: name
    integer, intent(in) :: ival
    logical, intent(in) :: is_known
    logical, intent(in), optional :: verbose
    call var_list_set_int (global%var_list, name, ival, is_known, &
         verbose=verbose)
  end subroutine rt_data_set_int

  subroutine rt_data_set_real (global, name, rval, is_known, verbose, pacified)
    class(rt_data_t), intent(inout) :: global
    type(string_t), intent(in) :: name
    real(default), intent(in) :: rval
    logical, intent(in) :: is_known
    logical, intent(in), optional :: verbose, pacified
    call var_list_set_real (global%var_list, name, rval, is_known, &
         verbose=verbose, pacified=pacified)
  end subroutine rt_data_set_real

  subroutine rt_data_set_cmplx (global, name, cval, is_known, verbose, pacified)
    class(rt_data_t), intent(inout) :: global
    type(string_t), intent(in) :: name
    complex(default), intent(in) :: cval
    logical, intent(in) :: is_known
    logical, intent(in), optional :: verbose, pacified
    call var_list_set_cmplx (global%var_list, name, cval, is_known, &
         verbose=verbose, pacified=pacified)
  end subroutine rt_data_set_cmplx

  subroutine rt_data_set_subevt (global, name, pval, is_known, verbose)
    class(rt_data_t), intent(inout) :: global
    type(string_t), intent(in) :: name
    type(subevt_t), intent(in) :: pval
    logical, intent(in) :: is_known
    logical, intent(in), optional :: verbose
    call var_list_set_subevt (global%var_list, name, pval, is_known, &
         verbose=verbose)
  end subroutine rt_data_set_subevt

  subroutine rt_data_set_pdg_array (global, name, aval, is_known, verbose)
    class(rt_data_t), intent(inout) :: global
    type(string_t), intent(in) :: name
    type(pdg_array_t), intent(in) :: aval
    logical, intent(in) :: is_known
    logical, intent(in), optional :: verbose
    call var_list_set_pdg_array (global%var_list, name, aval, is_known, &
         verbose=verbose)
  end subroutine rt_data_set_pdg_array

  subroutine rt_data_set_string (global, name, sval, is_known, verbose)
    class(rt_data_t), intent(inout) :: global
    type(string_t), intent(in) :: name
    type(string_t), intent(in) :: sval
    logical, intent(in) :: is_known
    logical, intent(in), optional :: verbose
    call var_list_set_string (global%var_list, name, sval, is_known, &
         verbose=verbose)
  end subroutine rt_data_set_string

\eatline
\eatline
\eatline
\eatline
\eatline
\eatline
\eatline
\nwendcode{}\nwbegindocs{1769}Return the value of a variable, assuming that the type is correct.  
\nwenddocs{}\nwbegincode{1770}\moddef{RT data: rt data: TBP}\plusendmoddef
  procedure :: get_lval => rt_data_get_lval
  procedure :: get_ival => rt_data_get_ival
  procedure :: get_rval => rt_data_get_rval
  procedure :: get_cval => rt_data_get_cval
  procedure :: get_pval => rt_data_get_pval
  procedure :: get_aval => rt_data_get_aval
  procedure :: get_sval => rt_data_get_sval
\nwendcode{}\nwbegincode{1771}\moddef{RT data: procedures}\plusendmoddef
  function rt_data_get_lval (global, name) result (lval)
    logical :: lval
    class(rt_data_t), intent(in), target :: global
    type(string_t), intent(in) :: name
    type(var_list_t), pointer :: var_list
    var_list => global%get_var_list_ptr ()
    lval = var_list%get_lval (name)
  end function rt_data_get_lval
  
  function rt_data_get_ival (global, name) result (ival)
    integer :: ival
    class(rt_data_t), intent(in), target :: global
    type(string_t), intent(in) :: name
    type(var_list_t), pointer :: var_list
    var_list => global%get_var_list_ptr ()
    ival = var_list%get_ival (name)
  end function rt_data_get_ival
  
  function rt_data_get_rval (global, name) result (rval)
    real(default) :: rval
    class(rt_data_t), intent(in), target :: global
    type(string_t), intent(in) :: name
    type(var_list_t), pointer :: var_list
    var_list => global%get_var_list_ptr ()
    rval = var_list%get_rval (name)
  end function rt_data_get_rval
    
  function rt_data_get_cval (global, name) result (cval)
    complex(default) :: cval
    class(rt_data_t), intent(in), target :: global
    type(string_t), intent(in) :: name
    type(var_list_t), pointer :: var_list
    var_list => global%get_var_list_ptr ()
    cval = var_list%get_cval (name)
  end function rt_data_get_cval

  function rt_data_get_aval (global, name) result (aval)
    type(pdg_array_t) :: aval
    class(rt_data_t), intent(in), target :: global
    type(string_t), intent(in) :: name
    type(var_list_t), pointer :: var_list
    var_list => global%get_var_list_ptr ()
    aval = var_list%get_aval (name)
  end function rt_data_get_aval
  
  function rt_data_get_pval (global, name) result (pval)
    type(subevt_t) :: pval
    class(rt_data_t), intent(in), target :: global
    type(string_t), intent(in) :: name
    type(var_list_t), pointer :: var_list
    var_list => global%get_var_list_ptr ()
    pval = var_list%get_pval (name)
  end function rt_data_get_pval
  
  function rt_data_get_sval (global, name) result (sval)
    type(string_t) :: sval
    class(rt_data_t), intent(in), target :: global
    type(string_t), intent(in) :: name
    type(var_list_t), pointer :: var_list
    var_list => global%get_var_list_ptr ()
    sval = var_list%get_sval (name)
  end function rt_data_get_sval
  
\eatline
\eatline
\eatline
\eatline
\eatline
\eatline
\eatline
\nwendcode{}\nwbegindocs{1772}\nwdocspar
\subsection{Further Content}
Add a library (available via a pointer of type {\tt{}prclib{\char95}entry{\char95}t}) to
the stack and update the pointer and variable list to the current
library.  The pointer association of {\tt{}prclib{\char95}entry} will be discarded.
\nwenddocs{}\nwbegincode{1773}\moddef{RT data: rt data: TBP}\plusendmoddef
  procedure :: add_prclib => rt_data_add_prclib
\nwendcode{}\nwbegincode{1774}\moddef{RT data: procedures}\plusendmoddef
  subroutine rt_data_add_prclib (global, prclib_entry)
    class(rt_data_t), intent(inout) :: global
    type(prclib_entry_t), intent(inout), pointer :: prclib_entry
    call global%prclib_stack%push (prclib_entry)
    call global%update_prclib (global%prclib_stack%get_first_ptr ())
  end subroutine rt_data_add_prclib
  
\eatline
\nwendcode{}\nwbegindocs{1775}Given a pointer to a process library, make this the currently active
library.
\nwenddocs{}\nwbegincode{1776}\moddef{RT data: rt data: TBP}\plusendmoddef
  procedure :: update_prclib => rt_data_update_prclib
\nwendcode{}\nwbegincode{1777}\moddef{RT data: procedures}\plusendmoddef
  subroutine rt_data_update_prclib (global, lib)
    class(rt_data_t), intent(inout) :: global
    type(process_library_t), intent(in), target :: lib
    global%prclib => lib
    if (global%var_list%contains (&
         var_str ("$library_name"), follow_link = .false.)) then
       call var_list_set_string (global%var_list, &
            var_str ("$library_name"), &
            global%prclib%get_name (), is_known=.true.)
    else
       call var_list_append_string (global%var_list, &
            var_str ("$library_name"), global%prclib%get_name (), &
            intrinsic = .true.)
    end if
  end subroutine rt_data_update_prclib

\eatline
\nwendcode{}\nwbegindocs{1778}\nwdocspar
\subsection{Miscellaneous}
The helicity selection data are distributed among several parameters.  Here,
we collect them in a single record.
\nwenddocs{}\nwbegincode{1779}\moddef{RT data: rt data: TBP}\plusendmoddef
  procedure :: get_helicity_selection => rt_data_get_helicity_selection
\nwendcode{}\nwbegincode{1780}\moddef{RT data: procedures}\plusendmoddef
  function rt_data_get_helicity_selection (rt_data) result (helicity_selection)
    class(rt_data_t), intent(in) :: rt_data
    type(helicity_selection_t) :: helicity_selection
    associate (var_list => rt_data%var_list)
      helicity_selection%active = var_list%get_lval (&
           var_str ("?helicity_selection_active"))
      if (helicity_selection%active) then
         helicity_selection%threshold = var_list%get_rval (&
              var_str ("helicity_selection_threshold"))
         helicity_selection%cutoff = var_list%get_ival (&
              var_str ("helicity_selection_cutoff"))
      end if
    end associate
  end function rt_data_get_helicity_selection

\eatline
\nwendcode{}\nwbegindocs{1781}Show the beam setup: beam structure and relevant global variables.
\nwenddocs{}\nwbegincode{1782}\moddef{RT data: rt data: TBP}\plusendmoddef
  procedure :: show_beams => rt_data_show_beams
\nwendcode{}\nwbegincode{1783}\moddef{RT data: procedures}\plusendmoddef
  subroutine rt_data_show_beams (rt_data, unit)
    class(rt_data_t), intent(in) :: rt_data
    integer, intent(in), optional :: unit
    type(string_t) :: s
    integer :: u
    u = given_output_unit (unit)
    associate (beams => rt_data%beam_structure, var_list => rt_data%var_list)
      call beams%write (u)
      if (.not. beams%asymmetric () .and. beams%get_n_beam () == 2) then
         write (u, "(2x,A,ES19.12,1x,'GeV')") "sqrts =", &
              var_list%get_rval (var_str ("sqrts"))         
      end if
      if (beams%contains ("pdf_builtin")) then
         s = var_list%get_sval (var_str ("$pdf_builtin_set"))
         if (s /= "") then
            write (u, "(2x,A,1x,3A)")  "PDF set =", '"', char (s), '"'
         else
            write (u, "(2x,A,1x,A)")  "PDF set =", "[undefined]"
         end if
      end if
      if (beams%contains ("lhapdf")) then
         s = var_list%get_sval (var_str ("$lhapdf_dir"))
         if (s /= "") then
            write (u, "(2x,A,1x,3A)")  "LHAPDF dir    =", '"', char (s), '"'
         end if
         s = var_list%get_sval (var_str ("$lhapdf_file"))
         if (s /= "") then
            write (u, "(2x,A,1x,3A)")  "LHAPDF file   =", '"', char (s), '"'
            write (u, "(2x,A,1x,I0)") "LHAPDF member =", &
                 var_list%get_ival (var_str ("lhapdf_member"))
         else
            write (u, "(2x,A,1x,A)")  "LHAPDF file   =", "[undefined]"
         end if
      end if
      if (beams%contains ("lhapdf_photon")) then
         s = var_list%get_sval (var_str ("$lhapdf_dir"))
         if (s /= "") then
            write (u, "(2x,A,1x,3A)")  "LHAPDF dir    =", '"', char (s), '"'
         end if
         s = var_list%get_sval (var_str ("$lhapdf_photon_file"))
         if (s /= "") then
            write (u, "(2x,A,1x,3A)")  "LHAPDF file   =", '"', char (s), '"'
            write (u, "(2x,A,1x,I0)") "LHAPDF member =", &
                 var_list%get_ival (var_str ("lhapdf_member"))
            write (u, "(2x,A,1x,I0)") "LHAPDF scheme =", &
                 var_list%get_ival (&
                 var_str ("lhapdf_photon_scheme"))
         else
            write (u, "(2x,A,1x,A)")  "LHAPDF file   =", "[undefined]"
         end if
      end if
      if (beams%contains ("isr")) then
         write (u, "(2x,A,ES19.12)") "ISR alpha =", &
              var_list%get_rval (var_str ("isr_alpha"))
         write (u, "(2x,A,ES19.12)") "ISR Q max =", &
              var_list%get_rval (var_str ("isr_q_max"))
         write (u, "(2x,A,ES19.12)") "ISR mass  =", &
              var_list%get_rval (var_str ("isr_mass"))
         write (u, "(2x,A,1x,I0)") "ISR order  =", &
              var_list%get_ival (var_str ("isr_order"))
         write (u, "(2x,A,1x,L1)") "ISR recoil =", &
              var_list%get_lval (var_str ("?isr_recoil"))
      end if
      if (beams%contains ("epa")) then
         write (u, "(2x,A,ES19.12)") "EPA alpha  =", &
              var_list%get_rval (var_str ("epa_alpha"))
         write (u, "(2x,A,ES19.12)") "EPA x min  =", &
              var_list%get_rval (var_str ("epa_x_min"))
         write (u, "(2x,A,ES19.12)") "EPA Q min  =", &
              var_list%get_rval (var_str ("epa_q_min"))
         write (u, "(2x,A,ES19.12)") "EPA E max  =", &
              var_list%get_rval (var_str ("epa_e_max"))
         write (u, "(2x,A,ES19.12)") "EPA mass   =", &
              var_list%get_rval (var_str ("epa_mass"))
         write (u, "(2x,A,1x,L1)") "EPA recoil =", &
              var_list%get_lval (var_str ("?epa_recoil"))
      end if
      if (beams%contains ("ewa")) then
         write (u, "(2x,A,ES19.12)") "EWA x min       =", &
              var_list%get_rval (var_str ("ewa_x_min"))
         write (u, "(2x,A,ES19.12)") "EWA Pt max      =", &
              var_list%get_rval (var_str ("ewa_pt_max"))
         write (u, "(2x,A,ES19.12)") "EWA mass        =", &
              var_list%get_rval (var_str ("ewa_mass"))
         write (u, "(2x,A,1x,L1)") "EWA mom cons.   =", &
              var_list%get_lval (&
              var_str ("?ewa_keep_momentum"))
         write (u, "(2x,A,1x,L1)") "EWA energ. cons. =", &
              var_list%get_lval (&
              var_str ("ewa_keep_energy"))
      end if
      if (beams%contains ("circe1")) then
         write (u, "(2x,A,1x,I0)") "CIRCE1 version    =", &
              var_list%get_ival (var_str ("circe1_ver"))
         write (u, "(2x,A,1x,I0)") "CIRCE1 revision   =", &
              var_list%get_ival (var_str ("circe1_rev")) 
         s = var_list%get_sval (var_str ("$circe1_acc"))
         write (u, "(2x,A,1x,A)") "CIRCE1 acceler.   =", char (s)
         write (u, "(2x,A,1x,I0)") "CIRCE1 chattin.   =", &
              var_list%get_ival (var_str ("circe1_chat"))
         write (u, "(2x,A,ES19.12)") "CIRCE1 sqrts      =", &
              var_list%get_rval (var_str ("circe1_sqrts"))
         write (u, "(2x,A,ES19.12)") "CIRCE1 epsil.     =", &
              var_list%get_rval (var_str ("circe1_eps"))
         write (u, "(2x,A,1x,L1)") "CIRCE1 phot. 1  =", &
              var_list%get_lval (var_str ("?circe1_photon1"))
         write (u, "(2x,A,1x,L1)") "CIRCE1 phot. 2  =", &
              var_list%get_lval (var_str ("?circe1_photon2"))
         write (u, "(2x,A,1x,L1)") "CIRCE1 generat. =", &
              var_list%get_lval (var_str ("?circe1_generate"))
         write (u, "(2x,A,1x,L1)") "CIRCE1 mapping  =", &
              var_list%get_lval (var_str ("?circe1_map"))
         write (u, "(2x,A,ES19.12)") "CIRCE1 map. slope =", &
              var_list%get_rval (var_str ("circe1_mapping_slope"))
      end if
      if (beams%contains ("circe2")) then
         s = var_list%get_sval (var_str ("$circe2_design"))
         write (u, "(2x,A,1x,A)") "CIRCE2 design   =", char (s) 
         s = var_list%get_sval (var_str ("$circe2_file"))
         write (u, "(2x,A,1x,A)") "CIRCE2 file     =", char (s)
         write (u, "(2x,A,1x,L1)") "CIRCE2 polarized =", &
              var_list%get_lval (var_str ("?circe2_polarized"))
      end if
      if (beams%contains ("beam_events")) then
         s = var_list%get_sval (var_str ("$beam_events_file"))
         write (u, "(2x,A,1x,A)") "Beam events file     =", char (s)
         write (u, "(2x,A,1x,L1)") "Beam events EOF warn =", &
              var_list%get_lval (var_str ("?beam_events_warn_eof"))
      end if
    end associate
  end subroutine rt_data_show_beams
  
\eatline
\nwendcode{}\nwbegindocs{1784}Return the collision energy as determined by the current beam
settings.  Without beam setup, this is the {\tt{}sqrts} variable.

If the value is meaningless for a setup, the function returns zero.
\nwenddocs{}\nwbegincode{1785}\moddef{RT data: rt data: TBP}\plusendmoddef
  procedure :: get_sqrts => rt_data_get_sqrts
\nwendcode{}\nwbegincode{1786}\moddef{RT data: procedures}\plusendmoddef
  function rt_data_get_sqrts (rt_data) result (sqrts)
    class(rt_data_t), intent(in) :: rt_data
    real(default) :: sqrts
    sqrts = rt_data%var_list%get_rval (var_str ("sqrts"))
  end function rt_data_get_sqrts
    
\eatline
\nwendcode{}\nwbegindocs{1787}For testing purposes, the {\tt{}rt{\char95}data{\char95}t} contents can be pacified to
suppress numerical fluctuations in (constant) test matrix elements.
\nwenddocs{}\nwbegincode{1788}\moddef{RT data: rt data: TBP}\plusendmoddef
  procedure :: pacify => rt_data_pacify
\nwendcode{}\nwbegincode{1789}\moddef{RT data: procedures}\plusendmoddef
  subroutine rt_data_pacify (rt_data, efficiency_reset, error_reset)
    class(rt_data_t), intent(inout) :: rt_data
    logical, intent(in), optional :: efficiency_reset, error_reset
    type(process_entry_t), pointer :: process
    process => rt_data%process_stack%first
    do while (associated (process))
       call process%pacify (efficiency_reset, error_reset)
       process => process%next
    end do    
  end subroutine rt_data_pacify

\eatline
\nwendcode{}\nwbegindocs{1790}Switch the matrix element entry to Gosam and back
\nwenddocs{}\nwbegincode{1791}\moddef{RT data: rt data: TBP}\plusendmoddef
  procedure :: change_to_gosam => rt_data_change_to_gosam
\nwendcode{}\nwbegincode{1792}\moddef{RT data: procedures}\plusendmoddef
  subroutine rt_data_change_to_gosam (global, success)
    class(rt_data_t), intent(inout) :: global
    logical, intent(out) :: success
    success = global%var_list%contains (var_str ("$method"))
    if (global%var_list%contains (var_str ("$method"))) then
       call global%var_list%set_sval (var_str ("$method"), var_str ("gosam"))
    end if
!     call var_list_replace_string &
!          (global%var_list, var_str ("$method"), var_str ("gosam"), &
!           intrinsic = .true., success = success)

  end subroutine rt_data_change_to_gosam

\nwendcode{}\nwbegincode{1793}\moddef{RT data: rt data: TBP}\plusendmoddef
  procedure :: change_to_omega => rt_data_change_to_omega
\nwendcode{}\nwbegincode{1794}\moddef{RT data: procedures}\plusendmoddef
  subroutine rt_data_change_to_omega (global)
    class(rt_data_t), intent(inout) :: global
    logical :: success
    success = global%var_list%contains (var_str ("$method"))
    if (global%var_list%contains (var_str ("$method"))) then
       call global%var_list%set_sval (var_str ("$method"), var_str ("omega"))
    end if
!     call var_list_replace_string &
!          (global%var_list, var_str ("$method"), var_str ("omega"), &
!           intrinsic = .true., success = success)
  end subroutine rt_data_change_to_omega

\eatline
\nwendcode{}\nwbegindocs{1795}\nwdocspar
\subsection{Test}
This is the master for calling self-test procedures.
\nwenddocs{}\nwbegincode{1796}\moddef{RT data: public}\plusendmoddef
  public :: rt_data_test
\nwendcode{}\nwbegincode{1797}\moddef{RT data: tests}\endmoddef
  subroutine rt_data_test (u, results)
    integer, intent(in) :: u
    type(test_results_t), intent(inout) :: results
  \LA{}RT data: execute tests\RA{}
  end subroutine rt_data_test

\eatline
\nwendcode{}\nwbegindocs{1798}\nwdocspar
\subsubsection{Initial content}
Force system-dependent objects to well-defined values.  Some of the
variables are locked and therefore must be addressed directly.

This is, of course, only required for testing purposes. In principle,
the {\tt{}real{\char95}specimen} variables could be set to their values in 
{\tt{}rt{\char95}data{\char95}t}, but this depends on the precision again, so we set
them to some dummy values.
\nwenddocs{}\nwbegincode{1799}\moddef{RT data: rt data: TBP}\plusendmoddef
  procedure :: fix_system_dependencies => rt_data_fix_system_dependencies
\nwendcode{}\nwbegincode{1800}\moddef{RT data: tests}\plusendmoddef
  subroutine rt_data_fix_system_dependencies (rt_data)
    class(rt_data_t), intent(inout), target :: rt_data
    type(var_list_t), pointer :: var_list
    var_list => rt_data%var_list
    call var_list_set_log (var_list, var_str ("?omega_openmp"), &
         .false., is_known = .true., force=.true.) 
    call var_list_set_log (var_list, var_str ("?openmp_is_active"), &
         .false., is_known = .true., force=.true.)
    call var_list_set_int (var_list, var_str ("openmp_num_threads_default"), &
         1, is_known = .true., force=.true.)
    call var_list_set_int (var_list, var_str ("openmp_num_threads"), &
         1, is_known = .true., force=.true.)        
    call var_list_set_int (var_list, var_str ("real_range"), &
         307, is_known = .true., force=.true.)
    call var_list_set_int (var_list, var_str ("real_precision"), &
         15, is_known = .true., force=.true.)    
    call var_list_set_real (var_list, var_str ("real_epsilon"), &
         1.e-16_default, is_known = .true., force=.true.)
    call var_list_set_real (var_list, var_str ("real_tiny"), &
         1.e-300_default, is_known = .true., force=.true.)     
    
    rt_data%os_data%fc = "Fortran-compiler"
    rt_data%os_data%fcflags = "Fortran-flags"
        
  end subroutine rt_data_fix_system_dependencies
  
\eatline
\nwendcode{}\nwbegindocs{1801}\nwdocspar
Display the RT data in the state just after (global) initialization.
\nwenddocs{}\nwbegincode{1802}\moddef{RT data: execute tests}\endmoddef
  call test (rt_data_1, "rt_data_1", &
       "initialize", &
       u, results)
\nwendcode{}\nwbegincode{1803}\moddef{RT data: tests}\plusendmoddef
  subroutine rt_data_1 (u)
    integer, intent(in) :: u
    type(rt_data_t), target :: rt_data
    
    write (u, "(A)")  "* Test output: rt_data_1"
    write (u, "(A)")  "*   Purpose: initialize global runtime data"
    write (u, "(A)")

    call rt_data%global_init (logfile = var_str ("rt_data.log"))

    call rt_data%fix_system_dependencies ()
    call rt_data%set_int (var_str ("seed"), 0, is_known=.true.)            

    call rt_data%it_list%init ([2, 3], [5000, 20000])

    call rt_data%write (u)

    call rt_data%final ()
    
    write (u, "(A)")
    write (u, "(A)")  "* Test output end: rt_data_1"
    
  end subroutine rt_data_1
  
\eatline
\nwendcode{}\nwbegindocs{1804}\nwdocspar
\subsubsection{Fill values}
Fill in empty slots in the runtime data block.
\nwenddocs{}\nwbegincode{1805}\moddef{RT data: execute tests}\plusendmoddef
  call test (rt_data_2, "rt_data_2", &
       "fill", &
       u, results)
\nwendcode{}\nwbegincode{1806}\moddef{RT data: tests}\plusendmoddef
  subroutine rt_data_2 (u)
    integer, intent(in) :: u
    type(rt_data_t), target :: rt_data
    type(flavor_t), dimension(2) :: flv
    type(string_t) :: cut_expr_text
    type(ifile_t) :: ifile
    type(stream_t) :: stream
    type(parse_tree_t) :: parse_tree
    
    write (u, "(A)")  "* Test output: rt_data_2"
    write (u, "(A)")  "*   Purpose: initialize global runtime data &
         &and fill contents"
    write (u, "(A)")

    call syntax_model_file_init ()

    call rt_data%global_init ()
    call rt_data%fix_system_dependencies ()

    call rt_data%select_model (var_str ("Test"))

    call rt_data%set_real (var_str ("sqrts"), &
         1000._default, is_known = .true.)
    call rt_data%set_int (var_str ("seed"), &
         0, is_known=.true.)        
    call flv%init ([25,25], rt_data%model)
    
    call rt_data%set_string (var_str ("$run_id"), &
         var_str ("run1"), is_known = .true.)
    call rt_data%set_real (var_str ("luminosity"), &
         33._default, is_known = .true.)
    
    call syntax_pexpr_init ()
    cut_expr_text = "all Pt > 100 [s]"
    call ifile_append (ifile, cut_expr_text)
    call stream_init (stream, ifile)
    call parse_tree_init_lexpr (parse_tree, stream, .true.)
    rt_data%pn%cuts_lexpr => parse_tree_get_root_ptr (parse_tree)
    
    allocate (rt_data%sample_fmt (2))
    rt_data%sample_fmt(1) = "foo_fmt"
    rt_data%sample_fmt(2) = "bar_fmt"
    
    call rt_data%write (u)

    call parse_tree_final (parse_tree)
    call stream_final (stream)
    call ifile_final (ifile)
    call syntax_pexpr_final ()

    call rt_data%final ()
    call syntax_model_file_final ()
    
    write (u, "(A)")
    write (u, "(A)")  "* Test output end: rt_data_2"
    
  end subroutine rt_data_2
  
\eatline
\nwendcode{}\nwbegindocs{1807}\nwdocspar
\subsubsection{Save and restore}
Set up a local runtime data block, change some contents, restore the
global block.
\nwenddocs{}\nwbegincode{1808}\moddef{RT data: execute tests}\plusendmoddef
  call test (rt_data_3, "rt_data_3", &
       "save/restore", &
       u, results)
\nwendcode{}\nwbegincode{1809}\moddef{RT data: tests}\plusendmoddef
  subroutine rt_data_3 (u)
    integer, intent(in) :: u
    type(rt_data_t), target :: rt_data, local
    type(var_list_t), pointer :: model_vars
    type(flavor_t), dimension(2) :: flv
    type(string_t) :: cut_expr_text
    type(ifile_t) :: ifile
    type(stream_t) :: stream
    type(parse_tree_t) :: parse_tree
    type(prclib_entry_t), pointer :: lib
    
    write (u, "(A)")  "* Test output: rt_data_3"
    write (u, "(A)")  "*   Purpose: initialize global runtime data &
         &and fill contents;"
    write (u, "(A)")  "*            copy to local block and back"
    write (u, "(A)")

    write (u, "(A)")  "* Init global data"
    write (u, "(A)")

    call syntax_model_file_init ()

    call rt_data%global_init ()
    call rt_data%fix_system_dependencies ()
    call rt_data%set_int (var_str ("seed"), &
         0, is_known=.true.)        

    call rt_data%select_model (var_str ("Test"))

    call rt_data%set_real (var_str ("sqrts"),&
         1000._default, is_known = .true.)
    call flv%init ([25,25], rt_data%model)
    
    call rt_data%beam_structure%init_sf (flv%get_name (), [1])
    call rt_data%beam_structure%set_sf (1, 1, var_str ("pdf_builtin"))

    call rt_data%set_string (var_str ("$run_id"), &
         var_str ("run1"), is_known = .true.)
    call rt_data%set_real (var_str ("luminosity"), &
         33._default, is_known = .true.)
    
    call syntax_pexpr_init ()
    cut_expr_text = "all Pt > 100 [s]"
    call ifile_append (ifile, cut_expr_text)
    call stream_init (stream, ifile)
    call parse_tree_init_lexpr (parse_tree, stream, .true.)
    rt_data%pn%cuts_lexpr => parse_tree_get_root_ptr (parse_tree)
    
    allocate (rt_data%sample_fmt (2))
    rt_data%sample_fmt(1) = "foo_fmt"
    rt_data%sample_fmt(2) = "bar_fmt"

    allocate (lib)
    call lib%init (var_str ("library_1"))
    call rt_data%add_prclib (lib)

    write (u, "(A)")  "* Init and modify local data"
    write (u, "(A)")

    call local%local_init (rt_data)
    call local%append_string (var_str ("$integration_method"), intrinsic=.true.)
    call local%append_string (var_str ("$phs_method"), intrinsic=.true.)

    call local%activate ()

    write (u, "(1x,A,L1)")  "model associated   = ", associated (local%model)
    write (u, "(1x,A,L1)")  "library associated = ", associated (local%prclib)
    write (u, *)

    call local%model_set_real (var_str ("ms"), 150._default)
    call local%set_string (var_str ("$integration_method"), &
         var_str ("midpoint"), is_known = .true.)
    call local%set_string (var_str ("$phs_method"), &
         var_str ("single"), is_known = .true.)

    local%os_data%fc = "Local compiler"
    
    allocate (lib)
    call lib%init (var_str ("library_2"))
    call local%add_prclib (lib)

    call local%write (u)

    write (u, "(A)")
    write (u, "(A)")  "* Restore global data"
    write (u, "(A)")
    
    call local%deactivate (rt_data)

    write (u, "(1x,A,L1)")  "model associated   = ", associated (rt_data%model)
    write (u, "(1x,A,L1)")  "library associated = ", associated (rt_data%prclib)
    write (u, *)

    call rt_data%write (u)

    write (u, "(A)")
    write (u, "(A)")  "* Cleanup"
    
    call parse_tree_final (parse_tree)
    call stream_final (stream)
    call ifile_final (ifile)
    call syntax_pexpr_final ()

    call rt_data%final ()
    call syntax_model_file_final ()
    
    write (u, "(A)")
    write (u, "(A)")  "* Test output end: rt_data_3"
    
  end subroutine rt_data_3
  
\eatline
\nwendcode{}\nwbegindocs{1810}\nwdocspar
\subsubsection{Show variables}
Display selected variables in the global record.
\nwenddocs{}\nwbegincode{1811}\moddef{RT data: execute tests}\plusendmoddef
  call test (rt_data_4, "rt_data_4", &
       "show variables", &
       u, results)
\nwendcode{}\nwbegincode{1812}\moddef{RT data: tests}\plusendmoddef
  subroutine rt_data_4 (u)
    integer, intent(in) :: u
    type(rt_data_t), target :: rt_data
    
    type(string_t), dimension(0) :: empty_string_array

    write (u, "(A)")  "* Test output: rt_data_4"
    write (u, "(A)")  "*   Purpose: display selected variables"
    write (u, "(A)")

    call rt_data%global_init ()

    write (u, "(A)")  "* No variables:"
    write (u, "(A)")

    call rt_data%write_vars (u, empty_string_array)

    write (u, "(A)")  "* Two variables:"
    write (u, "(A)")

    call rt_data%write_vars (u, &
         [var_str ("?unweighted"), var_str ("$phs_method")])
    
    write (u, "(A)")
    write (u, "(A)")  "* Display whole record with selected variables"
    write (u, "(A)")

    call rt_data%write (u, &
         vars = [var_str ("?unweighted"), var_str ("$phs_method")])

    call rt_data%final ()
    
    write (u, "(A)")
    write (u, "(A)")  "* Test output end: rt_data_4"
    
  end subroutine rt_data_4
  
\eatline
\nwendcode{}\nwbegindocs{1813}\nwdocspar
\subsubsection{Show parts}
Display only selected parts in the state just after (global) initialization.
\nwenddocs{}\nwbegincode{1814}\moddef{RT data: execute tests}\plusendmoddef
  call test (rt_data_5, "rt_data_5", &
       "show parts", &
       u, results)
\nwendcode{}\nwbegincode{1815}\moddef{RT data: tests}\plusendmoddef
  subroutine rt_data_5 (u)
    integer, intent(in) :: u
    type(rt_data_t), target :: rt_data
    
    write (u, "(A)")  "* Test output: rt_data_5"
    write (u, "(A)")  "*   Purpose: display parts of rt data"
    write (u, "(A)")

    call rt_data%global_init ()
    call rt_data%write_libraries (u)

    write (u, "(A)")

    call rt_data%write_beams (u)

    write (u, "(A)")

    call rt_data%write_process_stack (u)

    call rt_data%final ()
    
    write (u, "(A)")
    write (u, "(A)")  "* Test output end: rt_data_5"
    
  end subroutine rt_data_5
  
\eatline
\nwendcode{}\nwbegindocs{1816}\nwdocspar
\subsubsection{Local Model}
Locally modify a model and restore the global one.  We need an auxiliary
function to determine the status of a model particle:
\nwenddocs{}\nwbegincode{1817}\moddef{RT data: tests}\plusendmoddef
  function is_stable (pdg, rt_data) result (flag)
    integer, intent(in) :: pdg
    type(rt_data_t), intent(in) :: rt_data
    logical :: flag
    type(flavor_t) :: flv
    call flv%init (pdg, rt_data%model)
    flag = flv%is_stable ()
  end function is_stable
   
  function is_polarized (pdg, rt_data) result (flag)
    integer, intent(in) :: pdg
    type(rt_data_t), intent(in) :: rt_data
    logical :: flag
    type(flavor_t) :: flv
    call flv%init (pdg, rt_data%model)
    flag = flv%is_polarized ()
  end function is_polarized
    
\eatline
\nwendcode{}\nwbegincode{1818}\moddef{RT data: execute tests}\plusendmoddef
  call test (rt_data_6, "rt_data_6", &
       "local model", &
       u, results)
\nwendcode{}\nwbegincode{1819}\moddef{RT data: tests}\plusendmoddef
  subroutine rt_data_6 (u)
    integer, intent(in) :: u
    type(rt_data_t), target :: rt_data, local
    type(var_list_t), pointer :: model_vars
    type(string_t) :: var_name

    write (u, "(A)")  "* Test output: rt_data_6"
    write (u, "(A)")  "*   Purpose: apply and keep local modifications to model"
    write (u, "(A)")

    call syntax_model_file_init ()

    call rt_data%global_init ()
    call rt_data%select_model (var_str ("Test"))
    
    write (u, "(A)")  "* Original model"
    write (u, "(A)")

    call rt_data%write_model_list (u)
    write (u, *)
    write (u, "(A,L1)")  "s is stable    = ", is_stable (25, rt_data)
    write (u, "(A,L1)")  "f is polarized = ", is_polarized (6, rt_data)

    write (u, *)

    var_name = "ff"

    write (u, "(A)", advance="no")  "Global model variable: "
    model_vars => rt_data%model%get_var_list_ptr ()
    call var_list_write_var (model_vars, var_name, u)

    write (u, "(A)")
    write (u, "(A)")  "* Apply local modifications: unstable"
    write (u, "(A)")

    call local%local_init (rt_data)
    call local%activate ()

    call local%model_set_real (var_name, 0.4_default)
    call local%modify_particle (25, stable = .false., decay = [var_str ("d1")])
    call local%modify_particle (6, stable = .false., &
         decay = [var_str ("f1")], isotropic_decay = .true.)
    call local%modify_particle (-6, stable = .false., &
         decay = [var_str ("f2"), var_str ("f3")], diagonal_decay = .true.)

    call local%model%write (u)

    write (u, "(A)")
    write (u, "(A)")  "* Further modifications"
    write (u, "(A)")

    call local%modify_particle (6, stable = .false., &
         decay = [var_str ("f1")], &
         diagonal_decay = .true., isotropic_decay = .false.)
    call local%modify_particle (-6, stable = .false., &
         decay = [var_str ("f2"), var_str ("f3")], &
         diagonal_decay = .false., isotropic_decay = .true.)
    call local%model%write (u)

    write (u, "(A)")
    write (u, "(A)")  "* Further modifications: f stable but polarized"
    write (u, "(A)")

    call local%modify_particle (6, stable = .true., polarized = .true.)
    call local%modify_particle (-6, stable = .true.)
    call local%model%write (u)

    write (u, "(A)")
    write (u, "(A)")  "* Global model"
    write (u, "(A)")

    call rt_data%model%write (u)
    write (u, *)
    write (u, "(A,L1)")  "s is stable    = ", is_stable (25, rt_data)
    write (u, "(A,L1)")  "f is polarized = ", is_polarized (6, rt_data)

    write (u, "(A)")
    write (u, "(A)")  "* Local model"
    write (u, "(A)")

    call local%model%write (u)
    write (u, *)
    write (u, "(A,L1)")  "s is stable    = ", is_stable (25, local)
    write (u, "(A,L1)")  "f is polarized = ", is_polarized (6, local)

    write (u, *)

    write (u, "(A)", advance="no")  "Global model variable: "
    model_vars => rt_data%model%get_var_list_ptr ()
    call var_list_write_var (model_vars, var_name, u)

    write (u, "(A)", advance="no")  "Local model variable: "
    call var_list_write_var (local%model%get_var_list_ptr (), &
         var_name, u)

    write (u, "(A)")
    write (u, "(A)")  "* Restore global"

    call local%deactivate (rt_data, keep_local = .true.)

    write (u, "(A)")
    write (u, "(A)")  "* Global model"
    write (u, "(A)")

    call rt_data%model%write (u)
    write (u, *)
    write (u, "(A,L1)")  "s is stable    = ", is_stable (25, rt_data)
    write (u, "(A,L1)")  "f is polarized = ", is_polarized (6, rt_data)

    write (u, "(A)")
    write (u, "(A)")  "* Local model"
    write (u, "(A)")

    call local%model%write (u)
    write (u, *)
    write (u, "(A,L1)")  "s is stable    = ", is_stable (25, local)
    write (u, "(A,L1)")  "f is polarized = ", is_polarized (6, local)

    write (u, *)

    write (u, "(A)", advance="no")  "Global model variable: "
    model_vars => rt_data%model%get_var_list_ptr ()
    call var_list_write_var (model_vars, var_name, u)

    write (u, "(A)", advance="no")  "Local model variable: "
    call var_list_write_var (local%model%get_var_list_ptr (), &
         var_name, u)

    write (u, "(A)")
    write (u, "(A)")  "* Cleanup"

    call local%model%final ()
    deallocate (local%model)
    
    call rt_data%final ()
    call syntax_model_file_final ()
    
    write (u, "(A)")
    write (u, "(A)")  "* Test output end: rt_data_6"
    
  end subroutine rt_data_6
  
\eatline
\nwendcode{}\nwbegindocs{1820}\nwdocspar
\subsubsection{Result variables}
Initialize result variables and check that they are accessible via the
global variable list.
\nwenddocs{}\nwbegincode{1821}\moddef{RT data: execute tests}\plusendmoddef
  call test (rt_data_7, "rt_data_7", &
       "result variables", &
       u, results)
\nwendcode{}\nwbegincode{1822}\moddef{RT data: tests}\plusendmoddef
  subroutine rt_data_7 (u)
    integer, intent(in) :: u
    type(rt_data_t), target :: rt_data

    write (u, "(A)")  "* Test output: rt_data_7"
    write (u, "(A)")  "*   Purpose: set and access result variables"
    write (u, "(A)")

    write (u, "(A)")  "* Initialize process variables"
    write (u, "(A)")

    call rt_data%global_init ()
    call rt_data%process_stack%init_result_vars (var_str ("testproc"))
    
    call var_list_write_var (rt_data%var_list, &
         var_str ("integral(testproc)"), u)
    call var_list_write_var (rt_data%var_list, &
         var_str ("error(testproc)"), u)

    write (u, "(A)")
    write (u, "(A)")  "* Cleanup"
    
    call rt_data%final ()
    
    write (u, "(A)")
    write (u, "(A)")  "* Test output end: rt_data_7"
    
  end subroutine rt_data_7
  
\eatline
\nwendcode{}\nwbegindocs{1823}\nwdocspar
\subsubsection{Beam energy}
If beam parameters are set, the variable {\tt{}sqrts} is not necessarily
the collision energy.  The method {\tt{}get{\char95}sqrts} fetches the correct value.
\nwenddocs{}\nwbegincode{1824}\moddef{RT data: execute tests}\plusendmoddef
  call test (rt_data_8, "rt_data_8", &
       "beam energy", &
       u, results)
\nwendcode{}\nwbegincode{1825}\moddef{RT data: tests}\plusendmoddef
  subroutine rt_data_8 (u)
    integer, intent(in) :: u
    type(rt_data_t), target :: rt_data

    write (u, "(A)")  "* Test output: rt_data_8"
    write (u, "(A)")  "*   Purpose: get correct collision energy"
    write (u, "(A)")

    write (u, "(A)")  "* Initialize"
    write (u, "(A)")

    call rt_data%global_init ()

    write (u, "(A)")  "* Set sqrts"
    write (u, "(A)")

    call var_list_set_real (rt_data%var_list, var_str ("sqrts"), &
         1000._default, is_known = .true.)
    write (u, "(1x,A,ES19.12)")  "sqrts =", rt_data%get_sqrts ()

    write (u, "(A)")
    write (u, "(A)")  "* Cleanup"
    
    call rt_data%final ()
    
    write (u, "(A)")
    write (u, "(A)")  "* Test output end: rt_data_8"
    
  end subroutine rt_data_8
  
\eatline
\nwendcode{}\nwbegindocs{1826}\nwdocspar
\subsubsection{Local variable modifications}
\nwenddocs{}\nwbegincode{1827}\moddef{RT data: execute tests}\plusendmoddef
  call test (rt_data_9, "rt_data_9", &
       "local variables", &
       u, results)
\nwendcode{}\nwbegincode{1828}\moddef{RT data: tests}\plusendmoddef
  subroutine rt_data_9 (u)
    integer, intent(in) :: u
    type(rt_data_t), target :: global, local
    type(var_list_t), pointer :: var_list
    type(string_t) :: var_name

    write (u, "(A)")  "* Test output: rt_data_9"
    write (u, "(A)")  "*   Purpose: handle local variables"
    write (u, "(A)")

    call syntax_model_file_init ()

    write (u, "(A)")  "* Initialize global record and set some variables"
    write (u, "(A)")

    call global%global_init ()
    call global%select_model (var_str ("Test"))
    
    call global%set_real (var_str ("sqrts"), 17._default, is_known = .true.)
    call global%set_real (var_str ("luminosity"), 2._default, is_known = .true.)
    call global%model_set_real (var_str ("ff"), 0.5_default)
    call global%model_set_real (var_str ("gy"), 1.2_default)

    var_list => global%get_var_list_ptr ()

    call var_list_write_var (var_list, var_str ("sqrts"), u)
    call var_list_write_var (var_list, var_str ("luminosity"), u)
    call var_list_write_var (var_list, var_str ("ff"), u)
    call var_list_write_var (var_list, var_str ("gy"), u)
    call var_list_write_var (var_list, var_str ("mf"), u)
    call var_list_write_var (var_list, var_str ("x"), u)

    write (u, "(A)")
    
    write (u, "(1x,A,1x,F5.2)")  "sqrts      = ", &
         global%get_rval (var_str ("sqrts"))
    write (u, "(1x,A,1x,F5.2)")  "luminosity = ", &
         global%get_rval (var_str ("luminosity"))
    write (u, "(1x,A,1x,F5.2)")  "ff         = ", &
         global%get_rval (var_str ("ff"))
    write (u, "(1x,A,1x,F5.2)")  "gy         = ", &
         global%get_rval (var_str ("gy"))
    write (u, "(1x,A,1x,F5.2)")  "mf         = ", &
         global%get_rval (var_str ("mf"))
    write (u, "(1x,A,1x,F5.2)")  "x          = ", &
         global%get_rval (var_str ("x"))

    write (u, "(A)")
    write (u, "(A)")  "* Create local record with local variables"
    write (u, "(A)")

    call local%local_init (global)

    call local%append_real (var_str ("luminosity"), intrinsic = .true.)
    call local%append_real (var_str ("x"), user = .true.)

    call local%activate ()

    var_list => local%get_var_list_ptr ()

    call var_list_write_var (var_list, var_str ("sqrts"), u)
    call var_list_write_var (var_list, var_str ("luminosity"), u)
    call var_list_write_var (var_list, var_str ("ff"), u)
    call var_list_write_var (var_list, var_str ("gy"), u)
    call var_list_write_var (var_list, var_str ("mf"), u)
    call var_list_write_var (var_list, var_str ("x"), u)

    write (u, "(A)")
    
    write (u, "(1x,A,1x,F5.2)")  "sqrts      = ", &
         local%get_rval (var_str ("sqrts"))
    write (u, "(1x,A,1x,F5.2)")  "luminosity = ", &
         local%get_rval (var_str ("luminosity"))
    write (u, "(1x,A,1x,F5.2)")  "ff         = ", &
         local%get_rval (var_str ("ff"))
    write (u, "(1x,A,1x,F5.2)")  "gy         = ", &
         local%get_rval (var_str ("gy"))
    write (u, "(1x,A,1x,F5.2)")  "mf         = ", &
         local%get_rval (var_str ("mf"))
    write (u, "(1x,A,1x,F5.2)")  "x          = ", &
         local%get_rval (var_str ("x"))

    write (u, "(A)")
    write (u, "(A)")  "* Modify some local variables"
    write (u, "(A)")

    call local%set_real (var_str ("luminosity"), 42._default, is_known=.true.)
    call local%set_real (var_str ("x"), 6.66_default, is_known=.true.)
    call local%model_set_real (var_str ("ff"), 0.7_default)

    var_list => local%get_var_list_ptr ()

    call var_list_write_var (var_list, var_str ("sqrts"), u)
    call var_list_write_var (var_list, var_str ("luminosity"), u)
    call var_list_write_var (var_list, var_str ("ff"), u)
    call var_list_write_var (var_list, var_str ("gy"), u)
    call var_list_write_var (var_list, var_str ("mf"), u)
    call var_list_write_var (var_list, var_str ("x"), u)

    write (u, "(A)")
    
    write (u, "(1x,A,1x,F5.2)")  "sqrts      = ", &
         local%get_rval (var_str ("sqrts"))
    write (u, "(1x,A,1x,F5.2)")  "luminosity = ", &
         local%get_rval (var_str ("luminosity"))
    write (u, "(1x,A,1x,F5.2)")  "ff         = ", &
         local%get_rval (var_str ("ff"))
    write (u, "(1x,A,1x,F5.2)")  "gy         = ", &
         local%get_rval (var_str ("gy"))
    write (u, "(1x,A,1x,F5.2)")  "mf         = ", &
         local%get_rval (var_str ("mf"))
    write (u, "(1x,A,1x,F5.2)")  "x          = ", &
         local%get_rval (var_str ("x"))

    write (u, "(A)")
    write (u, "(A)")  "* Restore globals"
    write (u, "(A)")

    call local%deactivate (global)
    
    var_list => global%get_var_list_ptr ()

    call var_list_write_var (var_list, var_str ("sqrts"), u)
    call var_list_write_var (var_list, var_str ("luminosity"), u)
    call var_list_write_var (var_list, var_str ("ff"), u)
    call var_list_write_var (var_list, var_str ("gy"), u)
    call var_list_write_var (var_list, var_str ("mf"), u)
    call var_list_write_var (var_list, var_str ("x"), u)

    write (u, "(A)")
    
    write (u, "(1x,A,1x,F5.2)")  "sqrts      = ", &
         global%get_rval (var_str ("sqrts"))
    write (u, "(1x,A,1x,F5.2)")  "luminosity = ", &
         global%get_rval (var_str ("luminosity"))
    write (u, "(1x,A,1x,F5.2)")  "ff         = ", &
         global%get_rval (var_str ("ff"))
    write (u, "(1x,A,1x,F5.2)")  "gy         = ", &
         global%get_rval (var_str ("gy"))
    write (u, "(1x,A,1x,F5.2)")  "mf         = ", &
         global%get_rval (var_str ("mf"))
    write (u, "(1x,A,1x,F5.2)")  "x          = ", &
         global%get_rval (var_str ("x"))

    write (u, "(A)")
    write (u, "(A)")  "* Cleanup"

    call local%local_final ()

    call global%final ()
    call syntax_model_file_final ()
    
    write (u, "(A)")
    write (u, "(A)")  "* Test output end: rt_data_9"
    
  end subroutine rt_data_9
  
\eatline
\nwendcode{}\nwbegindocs{1829}\nwdocspar
\clearpage
%%%%%%%%%%%%%%%%%%%%%%%%%%%%%%%%%%%%%%%%%%%%%%%%%%%%%%%%%%%%%%%%%%%%%%%%
\section{Select implementations}
For abstract types (process core, integrator, phase space, etc.), we need a
way to dynamically select a concrete type, using either data given by the user
or a previous selection of a concrete type.  This is done by subroutines in
the current module.

This module provides no new types, just procedures.
\nwenddocs{}\nwbegincode{1830}\moddef{\code{}dispatch.f90\edoc{}}\endmoddef
\LA{}File header\RA{}

module dispatch
  
\LA{}Use kinds\RA{}
  use kinds, only: i16, double
\LA{}Use strings\RA{}
  use constants, only: PI
  use system_dependencies, only: LHAPDF5_AVAILABLE
  use system_dependencies, only: LHAPDF6_AVAILABLE
  use io_units
  use format_utils, only: write_separator
  use unit_tests
  use diagnostics
  use os_interface
  use physics_defs, only: MZ_REF, ALPHA_QCD_MZ_REF
  use physics_defs, only: PROTON, PHOTON, ELECTRON
  use sm_qcd
  use pdg_arrays
  use variables
  use model_data
  use flavors
  use interactions
  use models
  use process_constants
  use pdf_builtin !NODEP!
  use lhapdf !NODEP!  
  use sf_mappings
  use sf_base
  use sf_pdf_builtin
  use sf_lhapdf
  use sf_circe1
  use sf_circe2
  use sf_isr
  use sf_epa
  use sf_ewa
  use sf_escan
  use sf_beam_events
  use sf_user
  use rng_base
  use rng_tao
  use mci_base
  use mci_midpoint
  use mci_vamp
  use phs_base
  use mappings
  use phs_forests
  use phs_single
  use phs_wood
  use prc_core_def
  use prc_test
  use beams
  use prc_omega
  use prc_template_me
  use prc_core
  use processes
  use shower_base !NODEP!
  use shower
  use event_transforms
  use decays
  use beam_structures
  use eio_base
  use eio_raw
  use eio_checkpoints
  use eio_lhef
  use eio_hepmc
  use eio_lcio
  use eio_stdhep
  use eio_ascii
  use eio_weights
  use rt_data
  use prc_gosam
  use phs_fks
  use nlo_data

\LA{}Standard module head\RA{}

\LA{}Dispatch: public\RA{}

\LA{}Dispatch: types\RA{}

\LA{}Dispatch: interfaces\RA{}

contains
  
\LA{}Dispatch: procedures\RA{}

\LA{}Dispatch: tests\RA{}

end module dispatch
\eatline
\nwendcode{}\nwbegindocs{1831}\nwdocspar
\subsection{Interfaces}
\nwenddocs{}\nwbegincode{1832}\moddef{Dispatch: interfaces}\endmoddef
  interface
     subroutine GetXminM (set, mem, xmin)
       integer, intent(in) :: set, mem
       double precision, intent(out) :: xmin
     end subroutine GetXminM
  end interface

\eatline
\nwendcode{}\nwbegincode{1833}\moddef{Dispatch: interfaces}\plusendmoddef
  interface
     subroutine GetXmaxM (set, mem, xmax)
       integer, intent(in) :: set, mem
       double precision, intent(out) :: xmax
     end subroutine GetXmaxM
  end interface

\eatline
\nwendcode{}\nwbegincode{1834}\moddef{Dispatch: interfaces}\plusendmoddef
  interface
     subroutine GetQ2minM (set, mem, q2min)
       integer, intent(in) :: set, mem
       double precision, intent(out) :: q2min
     end subroutine GetQ2minM
  end interface

\eatline
\nwendcode{}\nwbegincode{1835}\moddef{Dispatch: interfaces}\plusendmoddef
  interface
     subroutine GetQ2maxM (set, mem, q2max)
       integer, intent(in) :: set, mem
       double precision, intent(out) :: q2max
     end subroutine GetQ2maxM
  end interface

\eatline
\nwendcode{}\nwbegindocs{1836}\nwdocspar
\subsection{Process Core Definition}
The {\tt{}prc{\char95}core{\char95}def{\char95}t} abstract type can be instantiated by providing a
{\tt{}{\char36}method} string variable.

Note: {\tt{}core{\char95}def} has intent(inout) because gfortran 4.7.1 crashes for
intent(out).
\nwenddocs{}\nwbegincode{1837}\moddef{Dispatch: public}\endmoddef
  public :: dispatch_core_def
\nwendcode{}\nwbegincode{1838}\moddef{Dispatch: procedures}\endmoddef
  subroutine dispatch_core_def (core_def, prt_in, prt_out, &
                                global, id, nlo_type)
    class(prc_core_def_t), allocatable, intent(inout) :: core_def
    type(string_t), dimension(:), intent(in) :: prt_in
    type(string_t), dimension(:), intent(in) :: prt_out
    type(rt_data_t), intent(in) :: global
    type(string_t), intent(in), optional :: id
    type(string_t), intent(in), optional :: nlo_type
    type(string_t) :: method
    type(string_t) :: model_name
    type(string_t) :: restrictions
    logical :: openmp_support
    logical :: report_progress
    logical :: diags, diags_color
    type(string_t) :: extra_options
    type(model_t), pointer :: model
    model => global%model
    associate (var_list => global%get_var_list_ptr ())
      method = var_list%get_sval (var_str ("$method"))
      if (associated (model)) then
         model_name = model%get_name ()
      else
         model_name = ""
      end if
      select case (char (method))
      case ("unit_test")
         allocate (prc_test_def_t :: core_def)
         select type (core_def)
         type is (prc_test_def_t)
            call core_def%init (model_name, prt_in, prt_out)
         end select
      case ("template")
         allocate (template_me_def_t :: core_def)
         select type (core_def)
         type is (template_me_def_t)
            call core_def%init (model, prt_in, prt_out, unity = .false.)
         end select
      case ("template_unity")
         allocate (template_me_def_t :: core_def)
         select type (core_def)
         type is (template_me_def_t)
            call core_def%init (model, prt_in, prt_out, unity = .true.)
         end select                  
      case ("omega")
         diags = var_list%get_lval (&
              var_str ("?diags"))
         diags_color = var_list%get_lval (&
              var_str ("?diags_color"))         
         restrictions = var_list%get_sval (&
              var_str ("$restrictions"))
         openmp_support = var_list%get_lval (&
              var_str ("?omega_openmp"))
         report_progress = var_list%get_lval (&
              var_str ("?report_progress"))
         extra_options = var_list%get_sval (&
              var_str ("$omega_flags"))
         allocate (omega_omega_def_t :: core_def)
         select type (core_def)
         type is (omega_omega_def_t)
            call core_def%init (model_name, prt_in, prt_out, &
                 restrictions, openmp_support, report_progress, &
                 extra_options, diags, diags_color)
         end select
      case ("ovm")
         diags = var_list%get_lval (&
              var_str ("?diags"))
         diags_color = var_list%get_lval (&
              var_str ("?diags_color"))         
         restrictions = var_list%get_sval (&
              var_str ("$restrictions"))
         openmp_support = var_list%get_lval (&
              var_str ("?omega_openmp"))
         report_progress = var_list%get_lval (&
              var_str ("?report_progress"))
         extra_options = var_list%get_sval (&
              var_str ("$omega_flags"))
         allocate (omega_ovm_def_t :: core_def)
         select type (core_def)
         type is (omega_ovm_def_t)
            call core_def%init (model_name, prt_in, prt_out, &
                 restrictions, openmp_support, report_progress, &
                 extra_options, diags, diags_color)
         end select         
      case ("gosam")
        allocate (gosam_def_t :: core_def)
        select type (core_def)
        type is (gosam_def_t)
          if (present (id)) then
             call core_def%init (id, model_name, prt_in, &
                                 prt_out, nlo_type)
          else
             call msg_fatal ("Dispatch GoSam def: No id!")
          end if
        end select
      case default
         call msg_fatal ("Process configuration: method '" &
              // char (method) // "' not implemented")
      end select
    end associate
  end subroutine dispatch_core_def
    
\eatline
\nwendcode{}\nwbegindocs{1839}\nwdocspar
\subsection{Process core allocation}
Here we allocate an object of abstract type {\tt{}prc{\char95}core{\char95}t} with a concrete
type that matches a process definition.  The {\tt{}prc{\char95}omega{\char95}t} extension
will require the current parameter set, so we take the opportunity to
grab it from the model.
\nwenddocs{}\nwbegincode{1840}\moddef{Dispatch: public}\plusendmoddef
  public :: dispatch_core
\nwendcode{}\nwbegincode{1841}\moddef{Dispatch: procedures}\plusendmoddef
  subroutine dispatch_core (core, core_def, model, &
       helicity_selection, qcd, use_color_factors)
    class(prc_core_t), allocatable, intent(inout) :: core
    class(prc_core_def_t), intent(in) :: core_def
    class(model_data_t), intent(in), target, optional :: model
    type(helicity_selection_t), intent(in), optional :: helicity_selection
    type(qcd_t), intent(in), optional :: qcd
    logical, intent(in), optional :: use_color_factors
    select type (core_def)
    type is (prc_test_def_t)
       allocate (test_t :: core)
    type is (template_me_def_t)
       allocate (prc_template_me_t :: core)
       select type (core)
       type is (prc_template_me_t)
          call core%set_parameters (model) 
       end select       
    class is (omega_def_t)
       if (.not. allocated (core)) allocate (prc_omega_t :: core)
       select type (core)
       type is (prc_omega_t)
          call core%set_parameters (model, & 
               helicity_selection, qcd, use_color_factors)
       end select
    type is (gosam_def_t)
      if (.not. allocated (core)) allocate (prc_gosam_t :: core)
      select type (core)
      type is (prc_gosam_t)
        call core%set_parameters (qcd, use_color_factors)
      end select
    class default
       call msg_bug ("Process core: unexpected process definition type")
    end select
  end subroutine dispatch_core

\eatline
\nwendcode{}\nwbegindocs{1842}\nwdocspar
\subsection{Process core update and restoration}
Here we take an existing object of abstract type {\tt{}prc{\char95}core{\char95}t} and
update the parameters as given by the current state of {\tt{}model}.
Optionally, we can save the previous state as {\tt{}saved{\char95}core}.  The
second routine restores the original from the save.

(In the test case, there is no possible update.)
\nwenddocs{}\nwbegincode{1843}\moddef{Dispatch: public}\plusendmoddef
  public :: dispatch_core_update
  public :: dispatch_core_restore
\nwendcode{}\nwbegincode{1844}\moddef{Dispatch: procedures}\plusendmoddef
  subroutine dispatch_core_update (core, model, helicity_selection, qcd, &
       saved_core)
    class(prc_core_t), allocatable, intent(inout) :: core
    class(model_data_t), intent(in), optional, target :: model
    type(helicity_selection_t), intent(in), optional :: helicity_selection
    type(qcd_t), intent(in), optional :: qcd
    class(prc_core_t), allocatable, intent(inout), optional :: saved_core
    if (present (saved_core)) then
       allocate (saved_core, source = core)
    end if
    select type (core)
    type is (test_t)
    type is (prc_omega_t)
       call core%set_parameters (model, helicity_selection, qcd)
       call core%activate_parameters ()
    type is (prc_gosam_t)
      call msg_message ("dispatch core restore: Gosam implementation not present yet!")
    class default
       call msg_bug ("Process core update: unexpected process definition type")
    end select
  end subroutine dispatch_core_update

  subroutine dispatch_core_restore (core, saved_core)
    class(prc_core_t), allocatable, intent(inout) :: core
    class(prc_core_t), allocatable, intent(inout) :: saved_core
    call move_alloc (from = saved_core, to = core)
    select type (core)
    type is (test_t)
    type is (prc_omega_t)
       call core%activate_parameters ()
    class default
       call msg_bug ("Process core restore: unexpected process definition type")
    end select
  end subroutine dispatch_core_restore

\eatline
\nwendcode{}\nwbegindocs{1845}\nwdocspar
\subsection{Integrator allocation}
Allocate an integrator according to the variable {\tt{}{\char36}integration{\char95}method}.
\nwenddocs{}\nwbegincode{1846}\moddef{Dispatch: public}\plusendmoddef
  public :: dispatch_mci
\nwendcode{}\nwbegincode{1847}\moddef{Dispatch: procedures}\plusendmoddef
  subroutine dispatch_mci (mci, global, process_id)
    class(mci_t), allocatable, intent(inout) :: mci
    type(rt_data_t), intent(in) :: global
    type(string_t), intent(in) :: process_id
    type(string_t) :: run_id
    type(string_t) :: integration_method
    type(grid_parameters_t) :: grid_par
    type(history_parameters_t) :: history_par
    logical :: rebuild_grids, check_grid_file, negative_weights, verbose
    integration_method = &
         global%var_list%get_sval (var_str ("$integration_method"))
    select case (char (integration_method))
    case ("midpoint")
       allocate (mci_midpoint_t :: mci)
    case ("vamp", "default")
       associate (var_list => global%get_var_list_ptr ())
         grid_par%threshold_calls = &
              var_list%get_ival (var_str ("threshold_calls"))
         grid_par%min_calls_per_channel = &
              var_list%get_ival (var_str ("min_calls_per_channel"))
         grid_par%min_calls_per_bin = &
              var_list%get_ival (var_str ("min_calls_per_bin"))
         grid_par%min_bins = &
              var_list%get_ival (var_str ("min_bins"))
         grid_par%max_bins = &
              var_list%get_ival (var_str ("max_bins"))
         grid_par%stratified = &
              var_list%get_lval (var_str ("?stratified"))
         grid_par%use_vamp_equivalences = &
              var_list%get_lval (var_str ("?use_vamp_equivalences"))
         grid_par%channel_weights_power = &
              var_list%get_rval (var_str ("channel_weights_power"))
         grid_par%accuracy_goal = &
              var_list%get_rval (var_str ("accuracy_goal"))
         grid_par%error_goal = &
              var_list%get_rval (var_str ("error_goal"))
         grid_par%rel_error_goal = &
              var_list%get_rval (var_str ("relative_error_goal"))
         history_par%global = &
              var_list%get_lval (var_str ("?vamp_history_global"))
         history_par%global_verbose = &
              var_list%get_lval (var_str ("?vamp_history_global_verbose"))
         history_par%channel = &
              var_list%get_lval (var_str ("?vamp_history_channels"))
         history_par%channel_verbose = &
              var_list%get_lval (var_str ("?vamp_history_channels_verbose"))
         verbose = &
              var_list%get_lval (var_str ("?vamp_verbose"))
         check_grid_file = &
              var_list%get_lval (var_str ("?check_grid_file"))
         run_id = &
              var_list%get_sval (var_str ("$run_id"))
         rebuild_grids = &
              var_list%get_lval (var_str ("?rebuild_grids"))
         negative_weights = &
              var_list%get_lval (var_str ("?negative_weights"))
       end associate
       allocate (mci_vamp_t :: mci)
       select type (mci)
       type is (mci_vamp_t)
          call mci%set_grid_parameters (grid_par)
          if (run_id /= "") then
             call mci%set_grid_filename (process_id, run_id)
          else
             call mci%set_grid_filename (process_id)
          end if
          call mci%set_history_parameters (history_par)
          call mci%set_rebuild_flag (rebuild_grids, check_grid_file)
          mci%negative_weights = negative_weights
          mci%verbose = verbose
       end select
    case default
       call msg_fatal ("Integrator '" &
            // char (integration_method) // "' not implemented")
    end select
  end subroutine dispatch_mci
  
\eatline
\nwendcode{}\nwbegindocs{1848}\nwdocspar
\subsection{Phase-space allocation}
Allocate a phase-space object according to the variable {\tt{}{\char36}phs{\char95}method}.
\nwenddocs{}\nwbegincode{1849}\moddef{Dispatch: public}\plusendmoddef
  public :: dispatch_phs
\nwendcode{}\nwbegincode{1850}\moddef{Dispatch: procedures}\plusendmoddef
  subroutine dispatch_phs (phs, global, process_id, mapping_defaults, phs_par, &
                           phs_method_in)
    class(phs_config_t), allocatable, intent(inout) :: phs
    type(rt_data_t), intent(in) :: global
    type(string_t), intent(in) :: process_id
    type(mapping_defaults_t), intent(in), optional :: mapping_defaults
    type(phs_parameters_t), intent(in), optional :: phs_par
    type(string_t), intent(in), optional :: phs_method_in
    type(string_t) :: phs_method, phs_file, run_id
    logical :: use_equivalences, vis_channels, fatal_beam_decay
    integer :: u_phs
    logical :: exist
    if (present (phs_method_in)) then
       phs_method = phs_method_in
    else
       phs_method = &
            global%var_list%get_sval (var_str ("$phs_method"))
    end if
    phs_file = &
         global%var_list%get_sval (var_str ("$phs_file"))
    use_equivalences = &
         global%var_list%get_lval (var_str ("?use_vamp_equivalences"))
    vis_channels = &
         global%var_list%get_lval (var_str ("?vis_channels"))
    fatal_beam_decay = &
         global%var_list%get_lval (var_str ("?fatal_beam_decay"))
    run_id = &
         global%var_list%get_sval (var_str ("$run_id"))    
    select case (char (phs_method))
    case ("single")
       allocate (phs_single_config_t :: phs)
       if (vis_channels) then
          call msg_warning ("Visualizing phase space channels not " // &
               "available for method 'single'.")
       end if
    case ("fks")
      allocate (phs_fks_config_t :: phs)
    case ("wood", "default")
       allocate (phs_wood_config_t :: phs)
       select type (phs)
       type is (phs_wood_config_t)
          if (phs_file /= "") then
             inquire (file = char (phs_file), exist = exist)
             if (exist) then
                call msg_message ("Phase space: reading configuration from '" &
                     // char (phs_file) // "'")
                u_phs = free_unit ()
                open (u_phs, file = char (phs_file), &
                     action = "read", status = "old")
                call phs%set_input (u_phs)
             else
                call msg_fatal ("Phase space: configuration file '" &
                     // char (phs_file) // "' not found")
             end if
          end if
          if (present (phs_par)) &
               call phs%set_parameters (phs_par)
          if (use_equivalences) &
               call phs%enable_equivalences ()
          if (present (mapping_defaults)) &
               call phs%set_mapping_defaults (mapping_defaults)
          phs%vis_channels = vis_channels
          phs%fatal_beam_decay = fatal_beam_decay
          phs%os_data = global%os_data
          phs%run_id = run_id
       end select
    case default
       call msg_fatal ("Phase space: parameterization method '" &
            // char (phs_method) // "' not implemented")
    end select
  end subroutine dispatch_phs
  
\eatline
\nwendcode{}\nwbegindocs{1851}\nwdocspar
\subsection{FKS parameter}
Initialize parameters used to optimize FKS calculations.
\nwenddocs{}\nwbegincode{1852}\moddef{Dispatch: public}\plusendmoddef
  public :: dispatch_fks
\nwendcode{}\nwbegincode{1853}\moddef{Dispatch: procedures}\plusendmoddef
  subroutine dispatch_fks (fks_template, global)
    type(fks_template_t), intent(inout) :: fks_template
    type(rt_data_t), intent(in) :: global
    real(default) :: fks_dij_exp1, fks_dij_exp2
    integer :: fks_mapping_type
    
    fks_dij_exp1 = &
         global%var_list%get_rval (var_str ("fks_dij_exp1"))
    fks_dij_exp2 = &
         global%var_list%get_rval (var_str ("fks_dij_exp2")) 
    fks_mapping_type = &
         global%var_list%get_ival (var_str ("fks_mapping_type"))

    call fks_template%set_dij_exp (fks_dij_exp1, fks_dij_exp2)
    call fks_template%set_mapping_type (fks_mapping_type)
    
  end subroutine dispatch_fks

\eatline
\nwendcode{}\nwbegindocs{1854}\nwdocspar
\subsection{Random-number generator}
Allocate a random-number generator factory according to the variable
{\tt{}{\char36}rng{\char95}method}, using the current seed in the global record.  We take
only the lower 15 bits of the seed, so the actual value fits into a
positive 16-bit signed integer.

Since we want to guarantee that all random-number generators in a run
are independent, we increment the global seed by one after creating
the rng factory.  If the user wants to have identical sequences, he
can always set the seed manually, before it is used.
\nwenddocs{}\nwbegincode{1855}\moddef{Dispatch: public}\plusendmoddef
  public :: dispatch_rng_factory
\nwendcode{}\nwbegincode{1856}\moddef{Dispatch: procedures}\plusendmoddef
  subroutine dispatch_rng_factory (rng_factory, global, local_input)
    class(rng_factory_t), allocatable, intent(inout) :: rng_factory
    type(rt_data_t), intent(inout), target :: global
    type(rt_data_t), intent(in), target, optional :: local_input
    type(rt_data_t), pointer :: local
    type(string_t) :: rng_method
    integer :: seed
    character(30) :: buffer
    integer(i16) :: s
    if (present (local_input)) then
       local => local_input
    else
       local => global
    end if
    rng_method = &
         local%var_list%get_sval (var_str ("$rng_method"))
    seed = &
         local%var_list%get_ival (var_str ("seed"))
    s = mod (seed, 32768)
    select case (char (rng_method))
    case ("unit_test")
       allocate (rng_test_factory_t :: rng_factory)
       call msg_message ("RNG: Initializing Test random-number generator")
    case ("tao")
       allocate (rng_tao_factory_t :: rng_factory)
       call msg_message ("RNG: Initializing TAO random-number generator")       
    case default
       call msg_fatal ("Random-number generator '" &
            // char (rng_method) // "' not implemented")
    end select
    write (buffer, "(I0)")  s
    call msg_message ("RNG: Setting seed for random-number generator to " &
            // trim (buffer))
    call rng_factory%init (s)
    call var_list_set_int (global%var_list, var_str ("seed"), seed + 1, &
         is_known = .true.)
  end subroutine dispatch_rng_factory
  
\eatline
\nwendcode{}\nwbegindocs{1857}\nwdocspar
\subsection{Structure function configuration data}
This data type is a container for transferring structure-function
specific data from the {\tt{}dispatch{\char95}sf{\char95}data} to the
{\tt{}dispatch{\char95}sf{\char95}channels} subroutine.
\nwenddocs{}\nwbegincode{1858}\moddef{Dispatch: public}\plusendmoddef
  public :: sf_prop_t
\nwendcode{}\nwbegincode{1859}\moddef{Dispatch: types}\endmoddef
  type :: sf_prop_t
     real(default), dimension(2) :: isr_eps = 1
  end type sf_prop_t
  
\eatline
\nwendcode{}\nwbegindocs{1860}\nwdocspar
Allocate a structure-function configuration object according to the
{\tt{}sf{\char95}method} string.

The {\tt{}sf{\char95}prop} object can be used to transfer structure-function
specific data up and to the {\tt{}dispatch{\char95}sf{\char95}channels} subroutine below,
so they can be used for particular mappings.

The {\tt{}global} object is intent(inout) because the RNG generator seed may change
during initialization.

The {\tt{}pdg{\char95}in} array is the array of incoming flavors, corresponding
to the upstream structure function or the beam array.  This will be
checked for the structure function in question and replaced by the
outgoing flavors.  The {\tt{}pdg{\char95}prc} array is the array of incoming
flavors (beam index, component index) for the hard process.
\nwenddocs{}\nwbegincode{1861}\moddef{Dispatch: public}\plusendmoddef
  public :: dispatch_sf_data
\nwendcode{}\nwbegincode{1862}\moddef{Dispatch: procedures}\plusendmoddef
  subroutine dispatch_sf_data (data, sf_method, i_beam, sf_prop, global, &
       pdg_in, pdg_prc)
    class(sf_data_t), allocatable, intent(inout) :: data
    type(string_t), intent(in) :: sf_method
    integer, dimension(:), intent(in) :: i_beam
    type(pdg_array_t), dimension(:), intent(inout) :: pdg_in
    type(pdg_array_t), dimension(:,:), intent(in) :: pdg_prc
    type(sf_prop_t), intent(inout) :: sf_prop
    type(rt_data_t), intent(inout) :: global
    type(model_t), pointer :: model
    type(pdg_array_t), dimension(:), allocatable :: pdg_out
    real(default) :: sqrts, isr_alpha, isr_q_max, isr_mass
    integer :: isr_order
    logical :: isr_recoil
    real(default) :: epa_alpha, epa_x_min, epa_q_min, epa_e_max, epa_mass
    logical :: epa_recoil
    real(default) :: ewa_x_min, ewa_pt_max, ewa_mass
    logical :: ewa_keep_momentum, ewa_keep_energy   
    type(pdg_array_t), dimension(:), allocatable :: pdg_prc1
    integer :: ewa_id
    type(string_t) :: pdf_name
    type(string_t) :: lhapdf_dir, lhapdf_file
    type(string_t), dimension(13) :: lhapdf_photon_sets
    integer :: lhapdf_member, lhapdf_photon_scheme
    logical :: hoppet_b_matching
    class(rng_factory_t), allocatable :: rng_factory
    logical :: circe1_photon1, circe1_photon2, circe1_generate
    real(default) :: circe1_sqrts, circe1_eps
    integer :: circe1_version, circe1_chattiness, &
         circe1_revision
    character(6) :: circe1_accelerator
    logical :: circe2_polarized
    type(string_t) :: circe2_design, circe2_file
    logical :: beam_events_warn_eof
    type(string_t) :: beam_events_dir, beam_events_file
    logical :: escan_normalize
    lhapdf_photon_sets = [var_str ("DOG0.LHgrid"), var_str ("DOG1.LHgrid"), &
         var_str ("DGG.LHgrid"), var_str ("LACG.LHgrid"), &
         var_str ("GSG0.LHgrid"), var_str ("GSG1.LHgrid"), &
         var_str ("GSG960.LHgrid"), var_str ("GSG961.LHgrid"), &
         var_str ("GRVG0.LHgrid"), var_str ("GRVG1.LHgrid"), &
         var_str ("ACFGPG.LHgrid"), var_str ("WHITG.LHgrid"), &
         var_str ("SASG.LHgrid")]
    model => global%model
    sqrts = global%get_sqrts ()
    associate (var_list => global%get_var_list_ptr ())
      select case (char (sf_method))
      case ("sf_test_0", "sf_test_1")
         allocate (sf_test_data_t :: data)
         select type (data)
         type is (sf_test_data_t)
            select case (char (sf_method))
            case ("sf_test_0");  call data%init (model, pdg_in(i_beam(1)))
            case ("sf_test_1");  call data%init (model, pdg_in(i_beam(1)), &
                 mode = 1)
            end select
         end select
      case ("pdf_builtin")
         allocate (pdf_builtin_data_t :: data)
         select type (data)
         type is (pdf_builtin_data_t)
            pdf_name = &
                 var_list%get_sval (var_str ("$pdf_builtin_set"))
            hoppet_b_matching = &
                 var_list%get_lval (var_str ("?hoppet_b_matching"))
            call data%init ( &
                 model, pdg_in(i_beam(1)), &
                 name = pdf_name, &
                 path = global%os_data%pdf_builtin_datapath, &
                 hoppet_b_matching = hoppet_b_matching)
         end select
      case ("pdf_builtin_photon")
         call msg_fatal ("Currently, there are no photon PDFs built into WHIZARD,", &
              [var_str ("for the photon content inside a proton or neutron use"), &
               var_str ("the 'lhapdf_photon' structure function.")])
      case ("lhapdf")
         allocate (lhapdf_data_t :: data)
         if (pdg_array_get (pdg_in(i_beam(1)), 1) == PHOTON) then
            call msg_fatal ("The 'lhapdf' structure is intended only for protons and", &
                 [var_str ("pions, please use 'lhapdf_photon' for photon beams.")])
         end if         
         lhapdf_dir = &
              var_list%get_sval (var_str ("$lhapdf_dir"))  
         lhapdf_file = &
              var_list%get_sval (var_str ("$lhapdf_file")) 
         lhapdf_member = &
              var_list%get_ival (var_str ("lhapdf_member"))
         lhapdf_photon_scheme = &
              var_list%get_ival (var_str ("lhapdf_photon_scheme"))
         hoppet_b_matching = &
              var_list%get_lval (var_str ("?hoppet_b_matching"))
         select type (data)
         type is (lhapdf_data_t)
            call data%init &
                 (model, pdg_in(i_beam(1)), &
                  lhapdf_dir, lhapdf_file, lhapdf_member, &
                  lhapdf_photon_scheme, hoppet_b_matching)
         end select
      case ("lhapdf_photon")
         allocate (lhapdf_data_t :: data)
         if (pdg_array_get_length (pdg_in(i_beam(1))) /= 1 .or. &
              pdg_array_get (pdg_in(i_beam(1)), 1) /= PHOTON) then
            call msg_fatal ("The 'lhapdf_photon' structure function is exclusively for", &
                 [var_str ("photon PDFs, i.e. for photons as beam particles")])
         end if
         lhapdf_dir = &
              var_list%get_sval (var_str ("$lhapdf_dir"))  
         lhapdf_file = &
              var_list%get_sval (var_str ("$lhapdf_photon_file")) 
         lhapdf_member = &
              var_list%get_ival (var_str ("lhapdf_member"))
         lhapdf_photon_scheme = &
              var_list%get_ival (var_str ("lhapdf_photon_scheme"))
         if (.not. any (lhapdf_photon_sets == lhapdf_file)) then
            call msg_fatal ("This PDF set is not supported or not " // & 
                 "intended for photon beams.")
         end if
         select type (data)
         type is (lhapdf_data_t)
            call data%init &
                 (model, pdg_in(i_beam(1)), &
                  lhapdf_dir, lhapdf_file, lhapdf_member, &
                  lhapdf_photon_scheme)
         end select         
      case ("isr")
         allocate (isr_data_t :: data)
         isr_alpha = &
              var_list%get_rval (var_str ("isr_alpha"))
         if (isr_alpha == 0) then
            isr_alpha = (var_list%get_rval (var_str ("ee"))) &
                 ** 2 / (4 * PI)
         end if
         isr_q_max = &
              var_list%get_rval (var_str ("isr_q_max"))
         if (isr_q_max == 0) then
            isr_q_max = sqrts
         end if
         isr_mass   = var_list%get_rval (var_str ("isr_mass"))
         isr_order  = var_list%get_ival (var_str ("isr_order"))
         isr_recoil = var_list%get_lval (var_str ("?isr_recoil")) 
         select type (data)
         type is (isr_data_t)
            call data%init &
                 (model, pdg_in (i_beam(1)), isr_alpha, isr_q_max, &
                 isr_mass, isr_order, isr_recoil)
            call data%check ()
            sf_prop%isr_eps(i_beam(1)) = data%get_eps ()
         end select
      case ("epa")
         allocate (epa_data_t :: data)
         epa_alpha = var_list%get_rval (var_str ("epa_alpha"))
         if (epa_alpha == 0) then
            epa_alpha = (var_list%get_rval (var_str ("ee"))) &
                 ** 2 / (4 * PI)
         end if         
         epa_x_min = var_list%get_rval (var_str ("epa_x_min"))
         epa_q_min = var_list%get_rval (var_str ("epa_q_min"))
         epa_e_max = var_list%get_rval (var_str ("epa_e_max"))
         if (epa_e_max == 0) then
            epa_e_max = sqrts
         end if
         epa_mass   = var_list%get_rval (var_str ("epa_mass"))
         epa_recoil = var_list%get_lval (var_str ("?epa_recoil"))
         select type (data)            
         type is (epa_data_t)
            call data%init &
                 (model, pdg_in (i_beam(1)), epa_alpha, epa_x_min, &
                 epa_q_min, epa_e_max, epa_mass, epa_recoil)
            call data%check ()
         end select
      case ("ewa")
         allocate (ewa_data_t :: data)
         allocate (pdg_prc1 (size (pdg_prc, 2)))
         pdg_prc1 = pdg_prc(i_beam(1),:)
         if (any (pdg_array_get_length (pdg_prc1) /= 1) &
              .or. any (pdg_prc1 /= pdg_prc1(1))) then
            call msg_fatal &
                 ("EWA: process incoming particle (W/Z) must be unique")
         end if
         ewa_id = abs (pdg_array_get (pdg_prc1(1), 1))
         ewa_x_min = var_list%get_rval (var_str ("ewa_x_min"))
         ewa_pt_max = var_list%get_rval (var_str ("ewa_pt_max")) 
         if (ewa_pt_max == 0) then
            ewa_pt_max = sqrts
         end if
         ewa_mass = var_list%get_rval (var_str ("ewa_mass"))  
         ewa_keep_momentum = var_list%get_lval (&
              var_str ("?ewa_keep_momentum"))
         ewa_keep_energy = var_list%get_lval (&
              var_str ("?ewa_keep_energy"))                  
         if (ewa_keep_momentum .and. ewa_keep_energy) &
              call msg_fatal (" EWA cannot conserve both energy " &
                 // "and momentum.")          
         select type (data)
         type is (ewa_data_t)
            call data%init &
                 (model, pdg_in (i_beam(1)), ewa_x_min, &
                 ewa_pt_max, sqrts, ewa_keep_momentum, &
                 ewa_keep_energy, ewa_mass)
            call data%set_id (ewa_id)
            call data%check ()
         end select
      case ("circe1")
         allocate (circe1_data_t :: data)
         select type (data)
         type is (circe1_data_t)
            circe1_photon1 = &
                 var_list%get_lval (var_str ("?circe1_photon1"))        
            circe1_photon2 = &
                 var_list%get_lval (var_str ("?circe1_photon2"))        
            circe1_sqrts = &
                 var_list%get_rval (var_str ("circe1_sqrts"))
            circe1_eps = &
                 var_list%get_rval (var_str ("circe1_eps"))
            if (circe1_sqrts <= 0)  circe1_sqrts = sqrts
            circe1_generate = &
                 var_list%get_lval (var_str ("?circe1_generate"))
            circe1_version = &
                 var_list%get_ival (var_str ("circe1_ver"))
            circe1_revision = &
                 var_list%get_ival (var_str ("circe1_rev"))
            circe1_accelerator = &
                 char (var_list%get_sval (var_str ("$circe1_acc")))
            circe1_chattiness = &
                 var_list%get_ival (var_str ("circe1_chat"))
            call data%init (model, pdg_in, circe1_sqrts, circe1_eps, &
                 [circe1_photon1, circe1_photon2], &
                 circe1_version, circe1_revision, circe1_accelerator, &
                 circe1_chattiness)
            if (circe1_generate) then
               call msg_message ("Circe1: activating generator mode")
               call dispatch_rng_factory (rng_factory, global)
               call data%set_generator_mode (rng_factory)
            end if
         end select
      case ("circe2")
         allocate (circe2_data_t :: data)
         select type (data)
         type is (circe2_data_t)
            circe2_polarized = &
                 var_list%get_lval (var_str ("?circe2_polarized"))
            circe2_file = &
                 var_list%get_sval (var_str ("$circe2_file"))
            circe2_design = &
                 var_list%get_sval (var_str ("$circe2_design"))
            call data%init (global%os_data, model, pdg_in, sqrts, &
                 circe2_polarized, circe2_file, circe2_design)
            call msg_message ("Circe2: activating generator mode")
            call dispatch_rng_factory (rng_factory, global)
            call data%set_generator_mode (rng_factory)
         end select
      case ("beam_events")
         allocate (beam_events_data_t :: data)
         select type (data)
         type is (beam_events_data_t)
            beam_events_dir = global%os_data%whizard_beamsimpath
            beam_events_file = var_list%get_sval (&
                 var_str ("$beam_events_file"))
            beam_events_warn_eof = var_list%get_lval (&
                 var_str ("?beam_events_warn_eof"))
            call data%init (model, pdg_in, &
                    beam_events_dir, beam_events_file, beam_events_warn_eof)  
         end select
      case ("energy_scan")
         escan_normalize = &
              var_list%get_lval (var_str ("?energy_scan_normalize"))
         allocate (escan_data_t :: data)
         select type (data) 
         type is (escan_data_t)
            if (escan_normalize) then
               call data%init (model, pdg_in)  
            else
               call data%init (model, pdg_in, sqrts)  
            end if
         end select
      case default
         call msg_bug ("Structure function '" &
              // char (sf_method) // "' not implemented yet")
      end select
    end associate
    if (allocated (data)) then
       allocate (pdg_out (size (pdg_prc, 1)))
       call data%get_pdg_out (pdg_out)
       pdg_in(i_beam) = pdg_out
    end if
  end subroutine dispatch_sf_data
  
\eatline
\nwendcode{}\nwbegindocs{1863}This is an auxiliary procedure, used by the beam-structure
expansion: tell for a given structure function name, whether it
corresponds to a pair spectrum ($n=2$), a single-particle structure
function ($n=1$), or nothing ($n=0$). Though {\tt{}energy{\char95}scan} can
in principle also be a pair spectrum, it always has only one
parameter.
\nwenddocs{}\nwbegincode{1864}\moddef{Dispatch: procedures}\plusendmoddef
  function strfun_mode (name) result (n)
    type(string_t), intent(in) :: name
    integer :: n
    select case (char (name))
    case ("none")
       n = 0
    case ("sf_test_0", "sf_test_1")
       n = 1
    case ("pdf_builtin","pdf_builtin_photon", &
          "lhapdf","lhapdf_photon")
       n = 1
    case ("isr","epa","ewa")
       n = 1
    case ("circe1", "circe2")
       n = 2
    case ("beam_events")
       n = 2
    case ("energy_scan")
       n = 2
    case default
       call msg_bug ("Structure function '" // char (name) &
            // "' not supported yet")
    end select
  end function strfun_mode
    
\eatline
\nwendcode{}\nwbegindocs{1865}Dispatch a whole structure-function chain, given beam data and beam
structure data.

This could be done generically, but we should look at the specific
combination of structure functions in order to select appropriate mappings.

The {\tt{}beam{\char95}structure} argument gets copied because
we want to expand it to canonical form (one valid structure-function
entry per record) before proceeding further.

The {\tt{}pdg{\char95}prc} argument is the array of incoming flavors.  The first
index is the beam index, the second one the process component index.
Each element is itself a PDG array, notrivial if there is a flavor sum
for the incoming state of this component.

The dispatcher is divided in two parts.  The first part configures the
structure function data themselves.  After this, we can configure the
phase space for the elementary process.
\nwenddocs{}\nwbegincode{1866}\moddef{Dispatch: public}\plusendmoddef
  public :: dispatch_sf_config
\nwendcode{}\nwbegincode{1867}\moddef{Dispatch: procedures}\plusendmoddef
  subroutine dispatch_sf_config (sf_config, sf_prop, global, pdg_prc)
    type(sf_config_t), dimension(:), allocatable, intent(out) :: sf_config
    type(sf_prop_t), intent(out) :: sf_prop
    type(rt_data_t), intent(inout) :: global
    type(beam_structure_t) :: beam_structure
    class(sf_data_t), allocatable :: sf_data
    type(pdg_array_t), dimension(:,:), intent(in) :: pdg_prc
    type(string_t), dimension(:), allocatable :: prt_in
    type(pdg_array_t), dimension(:), allocatable :: pdg_in
    type(flavor_t) :: flv_in
    integer :: n_beam, n_record, i
    beam_structure = global%beam_structure
    call beam_structure%expand (strfun_mode)
    n_record = beam_structure%get_n_record ()
    allocate (sf_config (n_record))
    n_beam = beam_structure%get_n_beam ()
    if (n_beam > 0) then
       allocate (prt_in (n_beam), pdg_in (n_beam))
       prt_in = beam_structure%get_prt ()
       do i = 1, n_beam
          call flv_in%init (prt_in(i), global%model)
          pdg_in(i) = flv_in%get_pdg ()
       end do
    else
       n_beam = size (pdg_prc, 1)
       allocate (pdg_in (n_beam))
       pdg_in = pdg_prc(:,1)
    end if
    do i = 1, n_record
       call dispatch_sf_data (sf_data, &
            beam_structure%get_name (i), &
            beam_structure%get_i_entry (i), &
            sf_prop, global, pdg_in, pdg_prc)
       call sf_config(i)%init (beam_structure%get_i_entry (i), sf_data)
       deallocate (sf_data)
    end do
  end subroutine dispatch_sf_config
    
\eatline
\nwendcode{}\nwbegindocs{1868}The second part configures channel mappings, using some conditions
from the phase space configuration.  If there are no structure
functions, we enable a default setup with a single (dummy)
structure-function channel.  Otherwise, we look at the channel
collection that we got from the phase-space configuration step.  Each
entry should be translated into an independent structure-function
channel, where typically there is one default entry, which could be
mapped using a standard s-channel mapping if the structure function
setup recommends this, and other entries with s-channel resonances.
The latter need to be translated into global mappings from the
structure-function chain.
\nwenddocs{}\nwbegincode{1869}\moddef{Dispatch: public}\plusendmoddef
  public :: dispatch_sf_channels
\nwendcode{}\nwbegincode{1870}\moddef{Dispatch: procedures}\plusendmoddef
  subroutine dispatch_sf_channels (sf_channel, sf_string, sf_prop, coll, global)
    type(sf_channel_t), dimension(:), allocatable, intent(out) :: sf_channel
    type(string_t), intent(out) :: sf_string
    type(sf_prop_t), intent(in) :: sf_prop
    type(phs_channel_collection_t), intent(in) :: coll
    type(rt_data_t), intent(in) :: global
    type(beam_structure_t) :: beam_structure
    class(channel_prop_t), allocatable :: prop
    integer :: n_strfun, n_sf_channel, i
    logical :: sf_allow_s_mapping, circe1_map, circe1_generate
    logical :: s_mapping_enable, endpoint_mapping, power_mapping
    integer, dimension(:), allocatable :: s_mapping, single_mapping
    real(default) :: sqrts, s_mapping_power
    real(default) :: circe1_mapping_slope, endpoint_mapping_slope
    real(default) :: power_mapping_eps
    sqrts = global%get_sqrts ()
    beam_structure = global%beam_structure
    call beam_structure%expand (strfun_mode)
    n_strfun = beam_structure%get_n_record ()
    sf_string = beam_structure%to_string (sf_only = .true.)
    sf_allow_s_mapping = &
         global%var_list%get_lval (var_str ("?sf_allow_s_mapping"))
    circe1_generate = &
         global%var_list%get_lval (var_str ("?circe1_generate"))
    circe1_map = &
         global%var_list%get_lval (var_str ("?circe1_map"))
    circe1_mapping_slope = &
         global%var_list%get_rval (var_str ("circe1_mapping_slope"))
    s_mapping_enable = .false.
    s_mapping_power = 1
    endpoint_mapping = .false.
    endpoint_mapping_slope = 1
    power_mapping = .false.
    select case (char (sf_string))
    case ("", "[any particles]")
    case ("pdf_builtin, none", &
         "pdf_builtin_photon, none", &
         "none, pdf_builtin", &
         "none, pdf_builtin_photon", &
         "lhapdf, none", &
         "lhapdf_photon, none", &
         "none, lhapdf", &
         "none, lhapdf_photon")
    case ("pdf_builtin, none => none, pdf_builtin", &
          "pdf_builtin, none => none, pdf_builtin_photon", &
          "pdf_builtin_photon, none => none, pdf_builtin", &
          "pdf_builtin_photon, none => none, pdf_builtin_photon", &
          "lhapdf, none => none, lhapdf", &
          "lhapdf, none => none, lhapdf_photon", &
          "lhapdf_photon, none => none, lhapdf", &
          "lhapdf_photon, none => none, lhapdf_photon")
       allocate (s_mapping (2), source = [1, 2])
       s_mapping_enable = .true.
       s_mapping_power = 2
    case ("pdf_builtin, none => none, pdf_builtin => epa, none => none, epa", &
          "pdf_builtin, none => none, pdf_builtin => ewa, none => none, ewa", &
          "pdf_builtin, none => none, pdf_builtin => ewa, none => none, epa", &
          "pdf_builtin, none => none, pdf_builtin => epa, none => none, ewa")
       allocate (s_mapping (2), source = [1, 2])
       s_mapping_enable = .true.
       s_mapping_power = 2
    case ("isr, none", &
         "none, isr")
       allocate (single_mapping (1), source = [1])
    case ("isr, none => none, isr")
       allocate (s_mapping (2), source = [1, 2])
       power_mapping = .true.
       power_mapping_eps = minval (sf_prop%isr_eps)
    case ("isr, none => none, isr => epa, none => none, epa", &
          "isr, none => none, isr => ewa, none => none, ewa", &
          "isr, none => none, isr => ewa, none => none, epa", &
          "isr, none => none, isr => epa, none => none, ewa")
       allocate (s_mapping (2), source = [1, 2])
       power_mapping = .true.
       power_mapping_eps = minval (sf_prop%isr_eps)
    case ("circe1 => isr, none => none, isr => epa, none => none, epa", &
          "circe1 => isr, none => none, isr => ewa, none => none, ewa", &
          "circe1 => isr, none => none, isr => ewa, none => none, epa", &
          "circe1 => isr, none => none, isr => epa, none => none, ewa")
       if (circe1_generate) then
          allocate (s_mapping (2), source = [2, 3])
       else
          allocate (s_mapping (3), source = [1, 2, 3])
          endpoint_mapping = .true.
          endpoint_mapping_slope = circe1_mapping_slope
       end if
       power_mapping = .true.
       power_mapping_eps = minval (sf_prop%isr_eps)       
    case ("pdf_builtin, none => none, isr", &
         "pdf_builtin_photon, none => none, isr", &
         "lhapdf, none => none, isr", &
         "lhapdf_photon, none => none, isr")
       allocate (single_mapping (1), source = [2])
    case ("isr, none => none, pdf_builtin", &
         "isr, none => none, pdf_builtin_photon", &
         "isr, none => none, lhapdf", &
         "isr, none => none, lhapdf_photon")
       allocate (single_mapping (1), source = [1])
    case ("epa, none", &
          "none, epa")
       allocate (single_mapping (1), source = [1])
    case ("epa, none => none, epa")
       allocate (single_mapping (2), source = [1, 2])
    case ("epa, none => none, isr", &
         "isr, none => none, epa", &
         "ewa, none => none, isr", &
         "isr, none => none, ewa")
       allocate (single_mapping (2), source = [1, 2])
    case ("pdf_builtin, none => none, epa", &
         "pdf_builtin_photon, none => none, epa", &
         "lhapdf, none => none, epa", &
         "lhapdf_photon, none => none, epa")
       allocate (single_mapping (1), source = [2])
    case ("pdf_builtin, none => none, ewa", &
         "pdf_builtin_photon, none => none, ewa", &
         "lhapdf, none => none, ewa", &
         "lhapdf_photon, none => none, ewa")
       allocate (single_mapping (1), source = [2])       
    case ("epa, none => none, pdf_builtin", &
         "epa, none => none, pdf_builtin_photon", &
         "epa, none => none, lhapdf", &
         "epa, none => none, lhapdf_photon")
       allocate (single_mapping (1), source = [1])
    case ("ewa, none => none, pdf_builtin", &
         "ewa, none => none, pdf_builtin_photon", &
         "ewa, none => none, lhapdf", &
         "ewa, none => none, lhapdf_photon")
       allocate (single_mapping (1), source = [1])       
    case ("ewa, none", &
          "none, ewa")
       allocate (single_mapping (1), source = [1])
    case ("ewa, none => none, ewa")
       allocate (single_mapping (2), source = [1, 2])
    case ("energy_scan, none => none, energy_scan")
       allocate (s_mapping (2), source = [1, 2])
    case ("sf_test_1, none => none, sf_test_1")
       allocate (s_mapping (2), source = [1, 2])
    case ("circe1")
       if (circe1_generate) then
          !!! no mapping
       else if (circe1_map) then
          allocate (s_mapping (1), source = [1])
          endpoint_mapping = .true.
          endpoint_mapping_slope = circe1_mapping_slope
       else
          allocate (s_mapping (1), source = [1])
          s_mapping_enable = .true.
       end if
    case ("circe1 => isr, none => none, isr")
       if (circe1_generate) then
          allocate (s_mapping (2), source = [2, 3])
       else
          allocate (s_mapping (3), source = [1, 2, 3])
          endpoint_mapping = .true.
          endpoint_mapping_slope = circe1_mapping_slope
       end if
       power_mapping = .true.
       power_mapping_eps = minval (sf_prop%isr_eps)
    case ("circe1 => isr, none", &
         "circe1 => none, isr")
       allocate (single_mapping (1), source = [2])
    case ("circe1 => epa, none => none, epa")
       if (circe1_generate) then
          allocate (single_mapping (2), source = [2, 3])
       else
          call msg_fatal ("Circe/EPA: supported with ?circe1_generate=true &
               &only")
       end if
    case ("circe1 => ewa, none => none, ewa")
       if (circe1_generate) then
          allocate (single_mapping (2), source = [2, 3])
       else 
          call msg_fatal ("Circe/EWA: supported with ?circe1_generate=true &
               &only")
       end if
    case ("circe1 => epa, none", &
         "circe1 => none, epa")
       if (circe1_generate) then
          allocate (single_mapping (1), source = [2])
       else
          call msg_fatal ("Circe/EPA: supported with ?circe1_generate=true &
               &only")
       end if
    case ("circe1 => epa, none => none, isr", &
         "circe1 => isr, none => none, epa", &
         "circe1 => ewa, none => none, isr", &
         "circe1 => isr, none => none, ewa")
       if (circe1_generate) then
          allocate (single_mapping (2), source = [2, 3])
       else
          call msg_fatal ("Circe/EPA: supported with ?circe1_generate=true &
               &only")
       end if
    case ("circe2", &
         "beam_events")
       !!! no mapping
    case ("circe2 => isr, none => none, isr", &
       "beam_events => isr, none => none, isr")
       allocate (s_mapping (2), source = [2, 3])
       power_mapping = .true.
       power_mapping_eps = minval (sf_prop%isr_eps)
    case ("circe2 => isr, none", &
         "circe2 => none, isr", &
         "beam_events => isr, none", &
         "beam_events => none, isr")
       allocate (single_mapping (1), source = [2])
    case ("circe2 => epa, none => none, epa", &
         "beam_events => epa, none => none, epa")
       allocate (single_mapping (2), source = [2, 3])
    case ("circe2 => epa, none", &
         "circe2 => none, epa", &
         "circe2 => ewa, none", &
         "circe2 => none, ewa", &
         "beam_events => epa, none", &
         "beam_events => none, epa", &
         "beam_events => ewa, none", &
         "beam_events => none, ewa")
       allocate (single_mapping (1), source = [2])
    case ("circe2 => epa, none => none, isr", &
         "circe2 => isr, none => none, epa", &
         "circe2 => ewa, none => none, isr", &
         "circe2 => isr, none => none, ewa", &
         "beam_events => epa, none => none, isr", &
         "beam_events => isr, none => none, epa", &
         "beam_events => ewa, none => none, isr", &
         "beam_events => isr, none => none, ewa")
       allocate (single_mapping (2), source = [2, 3])
    case ("energy_scan")
    case default
       call msg_fatal ("Beam structure: " &
            // char (sf_string) // " not supported")
    end select
    if (sf_allow_s_mapping .and. coll%n > 0) then
       n_sf_channel = coll%n
       allocate (sf_channel (n_sf_channel))
       do i = 1, n_sf_channel
          call sf_channel(i)%init (n_strfun)
          if (allocated (single_mapping)) then
             call sf_channel(i)%activate_mapping (single_mapping)
          end if
          if (allocated (prop))  deallocate (prop)
          call coll%get_entry (i, prop)
          if (allocated (prop)) then
             if (endpoint_mapping .and. power_mapping) then
                select type (prop)
                type is (resonance_t)
                   call sf_channel(i)%set_eir_mapping (s_mapping, &
                        a = endpoint_mapping_slope, eps = power_mapping_eps, &
                        m = prop%mass / sqrts, w = prop%width / sqrts)
                type is (on_shell_t)
                   call sf_channel(i)%set_eio_mapping (s_mapping, &
                        a = endpoint_mapping_slope, eps = power_mapping_eps, &
                        m = prop%mass / sqrts)
                end select
             else if (endpoint_mapping) then
                select type (prop)
                type is (resonance_t)
                   call sf_channel(i)%set_epr_mapping (s_mapping, &
                        a = endpoint_mapping_slope, &
                        m = prop%mass / sqrts, w = prop%width / sqrts)
                type is (on_shell_t)
                   call sf_channel(i)%set_epo_mapping (s_mapping, &
                        a = endpoint_mapping_slope, &
                        m = prop%mass / sqrts)
                end select
             else if (power_mapping) then
                select type (prop)
                type is (resonance_t)
                   call sf_channel(i)%set_ipr_mapping (s_mapping, &
                        eps = power_mapping_eps, &
                        m = prop%mass / sqrts, w = prop%width / sqrts)
                type is (on_shell_t)
                   call sf_channel(i)%set_ipo_mapping (s_mapping, &
                        eps = power_mapping_eps, &
                        m = prop%mass / sqrts)
                end select
             else if (allocated (s_mapping)) then
                select type (prop)
                type is (resonance_t)
                   call sf_channel(i)%set_res_mapping (s_mapping, &
                        m = prop%mass / sqrts, w = prop%width / sqrts)
                type is (on_shell_t)
                   call sf_channel(i)%set_os_mapping (s_mapping, &
                        m = prop%mass / sqrts)
                end select
             else if (allocated (single_mapping)) then
                select type (prop)
                type is (resonance_t)
                   call sf_channel(i)%set_res_mapping (single_mapping, &
                        m = prop%mass / sqrts, w = prop%width / sqrts)
                type is (on_shell_t)
                   call sf_channel(i)%set_os_mapping (single_mapping, &
                        m = prop%mass / sqrts)
                end select
             end if
          else if (endpoint_mapping .and. power_mapping) then
             call sf_channel(i)%set_ei_mapping (s_mapping, &
                  a = endpoint_mapping_slope, eps = power_mapping_eps)
          else if (endpoint_mapping .and. .not. allocated (single_mapping)) then
             call sf_channel(i)%set_ep_mapping (s_mapping, &
                  a = endpoint_mapping_slope)
          else if (power_mapping .and. .not. allocated (single_mapping)) then
             call sf_channel(i)%set_ip_mapping (s_mapping, &
                  eps = power_mapping_eps)
          else if (s_mapping_enable .and. .not. allocated (single_mapping)) then
             call sf_channel(i)%set_s_mapping (s_mapping, &
                  power = s_mapping_power)
          end if
       end do
    else if (sf_allow_s_mapping) then
       allocate (sf_channel (1))
       call sf_channel(1)%init (n_strfun)
       if (allocated (single_mapping)) then
          call sf_channel(1)%activate_mapping (single_mapping)
       else if (endpoint_mapping .and. power_mapping) then
          call sf_channel(i)%set_ei_mapping (s_mapping, &
               a = endpoint_mapping_slope, eps = power_mapping_eps)
       else if (endpoint_mapping) then
          call sf_channel(1)%set_ep_mapping (s_mapping, &
                  a = endpoint_mapping_slope)
       else if (power_mapping) then
          call sf_channel(1)%set_ip_mapping (s_mapping, &
                  eps = power_mapping_eps)
       else if (s_mapping_enable) then
          call sf_channel(1)%set_s_mapping (s_mapping, &
               power = s_mapping_power)
       end if
    else
       allocate (sf_channel (1))
       call sf_channel(1)%init (n_strfun)
       if (allocated (single_mapping)) then
          call sf_channel(1)%activate_mapping (single_mapping)
       end if
    end if
  end subroutine dispatch_sf_channels
    
\eatline
\nwendcode{}\nwbegindocs{1871}\nwdocspar
\subsection{Event I/O stream}
\nwenddocs{}\nwbegincode{1872}\moddef{Dispatch: public}\plusendmoddef
  public :: dispatch_eio
\nwendcode{}\nwbegincode{1873}\moddef{Dispatch: procedures}\plusendmoddef
  subroutine dispatch_eio (eio, method, global)
    class(eio_t), intent(inout), allocatable :: eio
    type(string_t), intent(in) :: method
    type(rt_data_t), intent(in) :: global
    logical :: check, keep_beams, keep_remnants, recover_beams
    logical :: use_alpha_s_from_file, use_scale_from_file
    logical :: write_sqme_prc, write_sqme_ref, write_sqme_alt
    type(string_t) :: lhef_version, lhef_extension, raw_version
    type(string_t) :: extension_default, debug_extension, extension_hepmc, &
         extension_lha, extension_hepevt, extension_ascii_short, &
         extension_ascii_long, extension_athena, extension_mokka, &
         extension_stdhep, extension_stdhep_up, extension_raw, &
         extension_hepevt_verb, extension_lha_verb, extension_lcio
    integer :: checkpoint
    logical :: show_process, show_transforms, show_decay, verbose, pacify
    select case (char (method))
    case ("raw")
       allocate (eio_raw_t :: eio)
       select type (eio)
       type is (eio_raw_t)
          check = &
               global%var_list%get_lval (var_str ("?check_event_file"))
          raw_version = &
               global%var_list%get_sval (var_str ("$event_file_version"))
          extension_raw = &
               global%var_list%get_sval (var_str ("$extension_raw"))
          call eio%set_parameters (check, raw_version, extension_raw)
       end select
    case ("checkpoint")
       allocate (eio_checkpoints_t :: eio)
       select type (eio)
       type is (eio_checkpoints_t)
          checkpoint = &
               global%var_list%get_ival (var_str ("checkpoint"))
          pacify = &
               global%var_list%get_lval (var_str ("?pacify"))
          call eio%set_parameters (checkpoint, blank = pacify)
       end select
    case ("lhef")
       allocate (eio_lhef_t :: eio)
       select type (eio)
       type is (eio_lhef_t)
          keep_beams = &
               global%var_list%get_lval (var_str ("?keep_beams"))
          keep_remnants = &
               global%var_list%get_lval (var_str ("?keep_remnants"))
          recover_beams = &
               global%var_list%get_lval (var_str ("?recover_beams"))
          use_alpha_s_from_file = &
               global%var_list%get_lval (var_str ("?use_alpha_s_from_file"))
          use_scale_from_file = &
               global%var_list%get_lval (var_str ("?use_scale_from_file"))
          lhef_version = &
               global%var_list%get_sval (var_str ("$lhef_version"))
          lhef_extension = &
               global%var_list%get_sval (var_str ("$lhef_extension"))
          write_sqme_prc = &
               global%var_list%get_lval (var_str ("?lhef_write_sqme_prc"))
          write_sqme_ref = &
               global%var_list%get_lval (var_str ("?lhef_write_sqme_ref"))
          write_sqme_alt = &
               global%var_list%get_lval (var_str ("?lhef_write_sqme_alt"))
          call eio%set_parameters ( &
               keep_beams, keep_remnants, recover_beams, &
               use_alpha_s_from_file, use_scale_from_file, &
               char (lhef_version), lhef_extension, &
               write_sqme_ref, write_sqme_prc, write_sqme_alt)
       end select
    case ("hepmc")
       allocate (eio_hepmc_t :: eio)
       select type (eio)
       type is (eio_hepmc_t)
          ! keep_beams = &
          !      global%var_list%get_lval (var_str ("?keep_beams"))
          use_alpha_s_from_file = &
               global%var_list%get_lval (var_str ("?use_alpha_s_from_file"))
          use_scale_from_file = &
               global%var_list%get_lval (var_str ("?use_scale_from_file"))
          recover_beams = &
               global%var_list%get_lval (var_str ("?recover_beams"))
          extension_hepmc = &
               global%var_list%get_sval (var_str ("$extension_hepmc"))          
          ! call eio%set_parameters (keep_beams, recover_beams, extension_hepmc)
          call eio%set_parameters (recover_beams, &
               use_alpha_s_from_file, use_scale_from_file, &
               extension_hepmc)
       end select
    case ("lcio")
       allocate (eio_lcio_t :: eio)
       select type (eio)
       type is (eio_lcio_t)
          ! keep_beams = &
          !      global%var_list%get_lval (var_str ("?keep_beams"))
          use_alpha_s_from_file = &
               global%var_list%get_lval (var_str ("?use_alpha_s_from_file"))
          use_scale_from_file = &
               global%var_list%get_lval (var_str ("?use_scale_from_file"))
          recover_beams = &
               global%var_list%get_lval (var_str ("?recover_beams"))
          extension_lcio = &
               global%var_list%get_sval (var_str ("$extension_lcio"))
          ! call eio%set_parameters &
          !      (keep_beams, keep_remnants, recover_beams, extension_lcio)
          call eio%set_parameters (recover_beams, &
               use_alpha_s_from_file, use_scale_from_file, &
               extension_lcio)
       end select       
    case ("stdhep")
       allocate (eio_stdhep_hepevt_t :: eio)
       select type (eio)
       type is (eio_stdhep_hepevt_t)                   
          keep_beams = &
               global%var_list%get_lval (var_str ("?keep_beams"))
          keep_remnants = &
               global%var_list%get_lval (var_str ("?keep_remnants"))
          use_alpha_s_from_file = &
               global%var_list%get_lval (var_str ("?use_alpha_s_from_file"))
          use_scale_from_file = &
               global%var_list%get_lval (var_str ("?use_scale_from_file"))
          recover_beams = &
               global%var_list%get_lval (var_str ("?recover_beams"))          
          extension_stdhep = &
               global%var_list%get_sval (var_str ("$extension_stdhep"))
          call eio%set_parameters &
               (keep_beams, keep_remnants, recover_beams, &
                use_alpha_s_from_file, use_scale_from_file, extension_stdhep)
       end select
    case ("stdhep_up")
       allocate (eio_stdhep_hepeup_t :: eio)
       select type (eio)
       type is (eio_stdhep_hepeup_t)          
          keep_beams = &
               global%var_list%get_lval (var_str ("?keep_beams"))
          keep_remnants = &
               global%var_list%get_lval (var_str ("?keep_remnants"))
          use_alpha_s_from_file = &
               global%var_list%get_lval (var_str ("?use_alpha_s_from_file"))
          use_scale_from_file = &
               global%var_list%get_lval (var_str ("?use_scale_from_file"))
          recover_beams = &
               global%var_list%get_lval (var_str ("?recover_beams"))          
          extension_stdhep_up = &
               global%var_list%get_sval (var_str ("$extension_stdhep_up")) 
          call eio%set_parameters (keep_beams, keep_remnants, &
               recover_beams, use_alpha_s_from_file, &
               use_scale_from_file, extension_stdhep_up)          
       end select       
    case ("ascii")   
       allocate (eio_ascii_ascii_t :: eio)
       select type (eio)
       type is (eio_ascii_ascii_t)
          keep_beams = &
               global%var_list%get_lval (var_str ("?keep_beams"))
          keep_remnants = &
               global%var_list%get_lval (var_str ("?keep_remnants"))
          extension_default = &
               global%var_list%get_sval (var_str ("$extension_default"))
          call eio%set_parameters &
               (keep_beams, keep_remnants, extension_default)
       end select       
    case ("athena")   
       allocate (eio_ascii_athena_t :: eio)
       select type (eio)
       type is (eio_ascii_athena_t)
          keep_beams = &
               global%var_list%get_lval (var_str ("?keep_beams"))
          keep_remnants = &
               global%var_list%get_lval (var_str ("?keep_remnants"))
          extension_athena = &
               global%var_list%get_sval (var_str ("$extension_athena"))
          call eio%set_parameters &
               (keep_beams, keep_remnants, extension_athena)
       end select              
    case ("debug")   
       allocate (eio_ascii_debug_t :: eio)
       select type (eio)
       type is (eio_ascii_debug_t)
          debug_extension = &
               global%var_list%get_sval (var_str ("$debug_extension"))          
          show_process = &
               global%var_list%get_lval (var_str ("?debug_process"))
          show_transforms = &
               global%var_list%get_lval (var_str ("?debug_transforms"))
          show_decay = &
               global%var_list%get_lval (var_str ("?debug_decay"))
          verbose = &
               global%var_list%get_lval (var_str ("?debug_verbose"))
          call eio%set_parameters ( &
               extension = debug_extension, &
               show_process = show_process, &
               show_transforms = show_transforms, &
               show_decay = show_decay, &
               verbose = verbose)
       end select
    case ("hepevt")   
       allocate (eio_ascii_hepevt_t :: eio)
       select type (eio)
       type is (eio_ascii_hepevt_t)
          keep_beams = &
               global%var_list%get_lval (var_str ("?keep_beams"))
          keep_remnants = &
               global%var_list%get_lval (var_str ("?keep_remnants"))
          extension_hepevt = &
               global%var_list%get_sval (var_str ("$extension_hepevt"))
          call eio%set_parameters &
               (keep_beams, keep_remnants, extension_hepevt)
       end select              
    case ("hepevt_verb")   
       allocate (eio_ascii_hepevt_verb_t :: eio)
       select type (eio)
       type is (eio_ascii_hepevt_verb_t)
          keep_beams = &
               global%var_list%get_lval (var_str ("?keep_beams"))
          keep_remnants = &
               global%var_list%get_lval (var_str ("?keep_remnants"))
          extension_hepevt_verb = &
               global%var_list%get_sval (var_str ("$extension_hepevt_verb"))
          call eio%set_parameters &
               (keep_beams, keep_remnants, extension_hepevt_verb)
       end select                     
    case ("lha")   
       allocate (eio_ascii_lha_t :: eio)
       select type (eio)
       type is (eio_ascii_lha_t)
          keep_beams = &
               global%var_list%get_lval (var_str ("?keep_beams"))
          keep_remnants = &
               global%var_list%get_lval (var_str ("?keep_remnants"))
          extension_lha = &
               global%var_list%get_sval (var_str ("$extension_lha"))
          call eio%set_parameters &
               (keep_beams, keep_remnants, extension_lha)
       end select                     
    case ("lha_verb")   
       allocate (eio_ascii_lha_verb_t :: eio)
       select type (eio)
       type is (eio_ascii_lha_verb_t)
          keep_beams = &
               global%var_list%get_lval (var_str ("?keep_beams"))
          keep_remnants = &
               global%var_list%get_lval (var_str ("?keep_remnants"))
          extension_lha_verb = global%var_list%get_sval ( &
               var_str ("$extension_lha_verb"))          
          call eio%set_parameters &
               (keep_beams, keep_remnants, extension_lha_verb)
       end select                            
    case ("long")   
       allocate (eio_ascii_long_t :: eio)
       select type (eio)
       type is (eio_ascii_long_t)
          keep_beams = &
               global%var_list%get_lval (var_str ("?keep_beams"))
          keep_remnants = &
               global%var_list%get_lval (var_str ("?keep_remnants"))
          extension_ascii_long = &
               global%var_list%get_sval (var_str ("$extension_ascii_long"))
          call eio%set_parameters &
               (keep_beams, keep_remnants, extension_ascii_long)
       end select              
    case ("mokka")   
       allocate (eio_ascii_mokka_t :: eio)
       select type (eio)
       type is (eio_ascii_mokka_t)
          keep_beams = &
               global%var_list%get_lval (var_str ("?keep_beams"))
          keep_remnants = &
               global%var_list%get_lval (var_str ("?keep_remnants"))
          extension_mokka = &
               global%var_list%get_sval (var_str ("$extension_mokka"))          
          call eio%set_parameters &
               (keep_beams, keep_remnants, extension_mokka)
       end select                     
    case ("short")   
       allocate (eio_ascii_short_t :: eio)
       select type (eio)
       type is (eio_ascii_short_t)
          keep_beams = &
               global%var_list%get_lval (var_str ("?keep_beams"))
          keep_remnants = &
               global%var_list%get_lval (var_str ("?keep_remnants"))
          extension_ascii_short = &
               global%var_list%get_sval (var_str ("$extension_ascii_short"))
          call eio%set_parameters &
               (keep_beams, keep_remnants, extension_ascii_short)
       end select                     
    case ("weight_stream")
       allocate (eio_weights_t :: eio)
       select type (eio)
       type is (eio_weights_t)
          pacify = &
               global%var_list%get_lval (var_str ("?pacify"))
          call eio%set_parameters (pacify = pacify)       
       end select
    case default
       call msg_fatal ("Event I/O method '" // char (method) &
            // "' not implemented")
    end select
    call eio%set_fallback_model (global%fallback_model)
  end subroutine dispatch_eio
  
\eatline
\nwendcode{}\nwbegindocs{1874}\nwdocspar
\subsection{QCD coupling}
Allocate the {\tt{}alpha} (running coupling) component of the {\tt{}qcd} block with
a concrete implementation, depending on the variable settings in the
{\tt{}global} record.

If a fixed $\alpha_s$ is requested, we do not allocate the
{\tt{}qcd{\char37}alpha} object.  In this case, the matrix element code will just take
the model parameter as-is, which implies fixed $\alpha_s$.  If the
object is allocated, the $\alpha_s$ value is computed and updated for
each matrix-element call.
\nwenddocs{}\nwbegincode{1875}\moddef{Dispatch: public}\plusendmoddef
  public :: dispatch_qcd
\nwendcode{}\nwbegincode{1876}\moddef{Dispatch: procedures}\plusendmoddef
  subroutine dispatch_qcd (qcd, global)
    type(qcd_t), intent(inout) :: qcd
    type(rt_data_t), intent(in), target :: global
    type(var_list_t), pointer :: var_list
    logical :: fixed, from_mz, from_pdf_builtin, from_lhapdf, from_lambda_qcd
    real(default) :: mz, alpha_val, lambda
    integer :: nf, order, lhapdf_member
    type(string_t) :: pdfset, lhapdf_dir, lhapdf_file
    var_list => global%get_var_list_ptr ()
    fixed = &
         var_list%get_lval (var_str ("?alpha_s_is_fixed"))
    from_mz = &
         var_list%get_lval (var_str ("?alpha_s_from_mz"))
    from_pdf_builtin = &
         var_list%get_lval (var_str ("?alpha_s_from_pdf_builtin"))
    from_lhapdf = &
         var_list%get_lval (var_str ("?alpha_s_from_lhapdf"))
    from_lambda_qcd = &
         var_list%get_lval (var_str ("?alpha_s_from_lambda_qcd"))
    pdfset = &
         var_list%get_sval (var_str ("$pdf_builtin_set"))    
    lambda = &
         var_list%get_rval (var_str ("lambda_qcd"))
    nf = &
         var_list%get_ival (var_str ("alpha_s_nf"))
    order = &
         var_list%get_ival (var_str ("alpha_s_order"))
    lhapdf_dir = &
         var_list%get_sval (var_str ("$lhapdf_dir"))
    lhapdf_file = &
         var_list%get_sval (var_str ("$lhapdf_file"))
    lhapdf_member = &
         var_list%get_ival (var_str ("lhapdf_member"))         
    var_list => global%get_var_list_ptr ()
    if (var_list%contains (var_str ("mZ"))) then
       mz = var_list%get_rval (var_str ("mZ"))
    else
       mz = MZ_REF
    end if
    if (var_list%contains (var_str ("alphas"))) then
       alpha_val = var_list%get_rval (var_str ("alphas"))
    else
       alpha_val = ALPHA_QCD_MZ_REF
    end if
    if (allocated (qcd%alpha))  deallocate (qcd%alpha)
    if (from_lhapdf .and. from_pdf_builtin) then
        call msg_fatal (" Mixing alphas evolution",  &
             [var_str (" from LHAPDF and builtin PDF is not permitted")])
    end if 
    select case (count ([from_mz, from_pdf_builtin, from_lhapdf, from_lambda_qcd]))
    case (0)
       if (fixed) then 
          allocate (alpha_qcd_fixed_t :: qcd%alpha)          
       else
          call msg_fatal ("QCD alpha: no calculation mode set")
       end if
    case (2:)
       call msg_fatal ("QCD alpha: calculation mode is ambiguous")
    case (1)
       if (fixed) then
          call msg_fatal ("QCD alpha: use '?alpha_s_is_fixed = false' for " // &
               "running alphas")          
       else if (from_mz) then
          allocate (alpha_qcd_from_scale_t :: qcd%alpha)
       else if (from_pdf_builtin) then
          allocate (alpha_qcd_pdf_builtin_t :: qcd%alpha)
       else if (from_lhapdf) then
          allocate (alpha_qcd_lhapdf_t :: qcd%alpha)
       else if (from_lambda_qcd) then
          allocate (alpha_qcd_from_lambda_t :: qcd%alpha)
       end if
       call msg_message ("QCD alpha: using a running strong coupling")
    end select
    select type (alpha => qcd%alpha)
    type is (alpha_qcd_fixed_t)
       alpha%val = alpha_val
    type is (alpha_qcd_from_scale_t)
       alpha%mu_ref = mz
       alpha%ref = alpha_val
       alpha%order = order
       alpha%nf = nf
    type is (alpha_qcd_from_lambda_t)
       alpha%lambda = lambda
       alpha%order = order
       alpha%nf = nf
    type is (alpha_qcd_pdf_builtin_t)
       call alpha%init (pdfset, &
            global%os_data%pdf_builtin_datapath)
    type is (alpha_qcd_lhapdf_t)
       call alpha%init (lhapdf_file, lhapdf_member, lhapdf_dir)
    end select
  end subroutine dispatch_qcd
  
\eatline
\nwendcode{}\nwbegindocs{1877}\nwdocspar
\subsection{Parton Shower}
Transfer the settings and parameters for shower and hadronization from
the {\tt{}global} object to the {\tt{}shower{\char95}settings} object.
\nwenddocs{}\nwbegincode{1878}\moddef{Dispatch: public}\plusendmoddef
  public :: dispatch_shower
\nwendcode{}\nwbegincode{1879}\moddef{Dispatch: procedures}\plusendmoddef
  subroutine dispatch_shower (shower_settings, global)
    class(shower_settings_t), intent(inout) :: shower_settings
    type(rt_data_t), intent(in), target :: global
    type(var_list_t), pointer :: var_list
    var_list => global%get_var_list_ptr ()
    call shower_settings%init (var_list)
  end subroutine dispatch_shower

\eatline
\nwendcode{}\nwbegindocs{1880}\nwdocspar
\subsection{Event Transform}
An event transform is responsible for dressing a partonic event.
Since event transforms are not mutually exclusive but are
concatenated, we provide individual dispatchers for each of them.
\nwenddocs{}\nwbegincode{1881}\moddef{Dispatch: public}\plusendmoddef
  public :: dispatch_evt_decay
  public :: dispatch_evt_shower
\nwendcode{}\nwbegincode{1882}\moddef{Dispatch: procedures}\plusendmoddef
  subroutine dispatch_evt_decay (evt, global)
    class(evt_t), intent(out), pointer :: evt
    type(rt_data_t), intent(in) :: global
    logical :: allow_decays
    allow_decays = &
         global%var_list%get_lval (var_str ("?allow_decays"))
    if (allow_decays) then
       allocate (evt_decay_t :: evt)
       call msg_message ("Simulate: activating decays")
    else
       evt => null ()
    end if
  end subroutine dispatch_evt_decay

  subroutine dispatch_evt_shower (evt, global, process)
    class(evt_t), intent(out), pointer :: evt
    type(rt_data_t), intent(in), target :: global
    type(process_t), intent(in), optional, target :: process
    logical :: allow_shower
    type(var_list_t), pointer :: var_list
    type(string_t) :: lhapdf_file, lhapdf_dir
    integer :: lhapdf_member
    double precision :: xmin, xmax, q2min, q2max
    type(shower_settings_t) :: settings
    external :: GetXminM, GetXmaxM, GetQ2minM, GetQ2maxM    
    var_list => global%get_var_list_ptr ()    
    allow_shower = &
         var_list%get_lval (var_str ("?allow_shower"))
    lhapdf_dir = &
         var_list%get_sval (var_str ("$lhapdf_dir"))    
    lhapdf_file = &
         var_list%get_sval (var_str ("$lhapdf_file"))
    lhapdf_member = &
         var_list%get_ival (var_str ("lhapdf_member"))             
    if (allow_shower) then
       allocate (evt_shower_t :: evt)
       call msg_message ("Simulate: activating parton shower")
       call dispatch_shower (settings, global)
       if (settings%mlm_matching) &
            call msg_message ("Simulate: applying MLM matching")
       if (settings%ckkw_matching) &
            call msg_warning ("Simulate: CKKW(-L) matching not yet supported")
       if (settings%hadronization_active) &
            call msg_message ("Simulate: applying hadronization")
       select type (evt)
       type is (evt_shower_t)
          call evt%init (settings, global%fallback_model, global%os_data)
          if (LHAPDF6_AVAILABLE) then
             call lhapdf_initialize &
                  (1, lhapdf_dir, lhapdf_file, lhapdf_member, evt%pdf)
          end if
          if (present (process)) &
               call evt%setup_pdf (process, global%beam_structure)
          select case (evt%pdf_type)
          case (STRF_LHAPDF6)
             evt%xmin = evt%pdf%getxmin ()
             evt%xmax = evt%pdf%getxmax ()
             evt%qmin = sqrt(evt%pdf%getq2min ())
             evt%qmax = sqrt(evt%pdf%getq2max ())
          case (STRF_LHAPDF5)
             if (LHAPDF5_AVAILABLE) then
                call GetXminM (1, lhapdf_member, xmin)
                call GetXmaxM (1, lhapdf_member, xmax)
                call GetQ2minM (1, lhapdf_member, q2min)
                call GetQ2maxM (1, lhapdf_member, q2max)
                evt%xmin = xmin
                evt%xmax = xmax
                evt%qmin = sqrt(q2min)
                evt%qmax = sqrt(q2max)
             end if
          end select
       end select
    else
       evt => null ()
    end if
  end subroutine dispatch_evt_shower

\eatline
\eatline
\nwendcode{}\nwbegindocs{1883}\nwdocspar
\subsection{SLHA Interface}
This is a container for the settings for the SLHA Interface
\nwenddocs{}\nwbegincode{1884}\moddef{Dispatch: public}\plusendmoddef
  public :: dispatch_slha
\nwendcode{}\nwbegincode{1885}\moddef{Dispatch: procedures}\plusendmoddef
  subroutine dispatch_slha (global, input, spectrum, decays)
    type(rt_data_t), intent(inout), target :: global
    logical, intent(out) :: input, spectrum, decays
    input = &
         global%var_list%get_lval (var_str ("?slha_read_input"))
    spectrum = &
         global%var_list%get_lval (var_str ("?slha_read_spectrum"))
    decays = &
         global%var_list%get_lval (var_str ("?slha_read_decays"))    
  end subroutine dispatch_slha

\eatline
\nwendcode{}\nwbegindocs{1886}\nwdocspar
\subsection{Test}
This is the master for calling self-test procedures.
\nwenddocs{}\nwbegincode{1887}\moddef{Dispatch: public}\plusendmoddef
  public :: dispatch_test
\nwendcode{}\nwbegincode{1888}\moddef{Dispatch: tests}\endmoddef
  subroutine dispatch_test (u, results)
    integer, intent(in) :: u
    type(test_results_t), intent(inout) :: results
  \LA{}Dispatch: execute tests\RA{}
  end subroutine dispatch_test

\eatline
\nwendcode{}\nwbegindocs{1889}\nwdocspar
\subsubsection{Select type: process definition}
\nwenddocs{}\nwbegincode{1890}\moddef{Dispatch: execute tests}\endmoddef
  call test (dispatch_1, "dispatch_1", &
       "process configuration method", &
       u, results)
\nwendcode{}\nwbegincode{1891}\moddef{Dispatch: tests}\plusendmoddef
  subroutine dispatch_1 (u)
    integer, intent(in) :: u
    type(string_t), dimension(2) :: prt_in, prt_out
    type(rt_data_t), target :: global
    class(prc_core_def_t), allocatable :: core_def
    
    write (u, "(A)")  "* Test output: dispatch_1"
    write (u, "(A)")  "*   Purpose: select process configuration method"
    write (u, "(A)")

    call global%global_init ()
    
    call var_list_set_log (global%var_list, var_str ("?omega_openmp"), &
         .false., is_known = .true.)

    prt_in = [var_str ("a"), var_str ("b")]
    prt_out = [var_str ("c"), var_str ("d")]

    write (u, "(A)")  "* Allocate core_def as prc_test_def"

    call var_list_set_string (global%var_list, var_str ("$method"), &
         var_str ("unit_test"), is_known = .true.)
    call dispatch_core_def (core_def, prt_in, prt_out, global)
    select type (core_def)
    type is (prc_test_def_t)
       call core_def%write (u)
    end select
    
    deallocate (core_def)

    write (u, "(A)")
    write (u, "(A)")  "* Allocate core_def as omega_def"
    write (u, "(A)")

    call var_list_set_string (global%var_list, var_str ("$method"), &
         var_str ("omega"), is_known = .true.)
    call dispatch_core_def (core_def, prt_in, prt_out, global)
    select type (core_def)
    type is (omega_omega_def_t)
       call core_def%write (u)
    end select
    
    call global%final ()
    
    write (u, "(A)")
    write (u, "(A)")  "* Test output end: dispatch_1"
    
  end subroutine dispatch_1
  
\eatline
\nwendcode{}\nwbegindocs{1892}\nwdocspar
\subsubsection{Select type: process core}
\nwenddocs{}\nwbegincode{1893}\moddef{Dispatch: execute tests}\plusendmoddef
  call test (dispatch_2, "dispatch_2", &
       "process core", &
       u, results)
\nwendcode{}\nwbegincode{1894}\moddef{Dispatch: tests}\plusendmoddef
  subroutine dispatch_2 (u)
    integer, intent(in) :: u
    type(string_t), dimension(2) :: prt_in, prt_out
    type(rt_data_t), target :: global
    class(prc_core_def_t), allocatable :: core_def
    class(prc_core_t), allocatable :: core
    
    write (u, "(A)")  "* Test output: dispatch_2"
    write (u, "(A)")  "*   Purpose: select process configuration method"
    write (u, "(A)")  "             and allocate process core"
    write (u, "(A)")

    call syntax_model_file_init ()
    call global%global_init ()

    prt_in = [var_str ("a"), var_str ("b")]
    prt_out = [var_str ("c"), var_str ("d")]

    write (u, "(A)")  "* Allocate core as test_t"
    write (u, "(A)")

    call var_list_set_string (global%var_list, var_str ("$method"), &
         var_str ("unit_test"), is_known = .true.)
    call dispatch_core_def (core_def, prt_in, prt_out, global)
    call dispatch_core (core, core_def)
    select type (core)
    type is (test_t)
       call core%write (u)
    end select
    
    deallocate (core)
    deallocate (core_def)

    write (u, "(A)")
    write (u, "(A)")  "* Allocate core as prc_omega_t"
    write (u, "(A)")

    call var_list_set_string (global%var_list, var_str ("$method"), &
         var_str ("omega"), is_known = .true.)
    call dispatch_core_def (core_def, prt_in, prt_out, global)

    call global%select_model (var_str ("Test"))

    call var_list_set_log (global%var_list, &
         var_str ("?helicity_selection_active"), &
         .true., is_known = .true.)
    call var_list_set_real (global%var_list, &
         var_str ("helicity_selection_threshold"), &
         1e9_default, is_known = .true.)
    call var_list_set_int (global%var_list, &
         var_str ("helicity_selection_cutoff"), &
         10, is_known = .true.)
    
    call dispatch_core (core, core_def, &
         global%model, global%get_helicity_selection ())
    call core_def%allocate_driver (core%driver, var_str (""))

    select type (core)
    type is (prc_omega_t)
       call core%write (u)
    end select
    
    call global%final ()
    call syntax_model_file_final ()
    
    write (u, "(A)")
    write (u, "(A)")  "* Test output end: dispatch_2"
    
  end subroutine dispatch_2
  
\eatline
\nwendcode{}\nwbegindocs{1895}\nwdocspar
\subsubsection{Select type: integrator core}
\nwenddocs{}\nwbegincode{1896}\moddef{Dispatch: execute tests}\plusendmoddef
  call test (dispatch_3, "dispatch_3", &
       "integration method", &
       u, results)
\nwendcode{}\nwbegincode{1897}\moddef{Dispatch: tests}\plusendmoddef
  subroutine dispatch_3 (u)
    integer, intent(in) :: u
    type(rt_data_t), target :: global
    class(mci_t), allocatable :: mci
    type(string_t) :: process_id
    
    write (u, "(A)")  "* Test output: dispatch_3"
    write (u, "(A)")  "*   Purpose: select integration method"
    write (u, "(A)")

    call global%global_init ()
    process_id = "dispatch_3"

    write (u, "(A)")  "* Allocate MCI as midpoint_t"
    write (u, "(A)")

    call var_list_set_string (global%var_list, &
         var_str ("$integration_method"), &
         var_str ("midpoint"), is_known = .true.)
    call dispatch_mci (mci, global, process_id)
    select type (mci)
    type is (mci_midpoint_t)
       call mci%write (u)
    end select

    call mci%final ()
    deallocate (mci)
    
    write (u, "(A)")
    write (u, "(A)")  "* Allocate MCI as vamp_t"
    write (u, "(A)")

    call var_list_set_string (global%var_list, &
         var_str ("$integration_method"), &
         var_str ("vamp"), is_known = .true.)
    call var_list_set_int (global%var_list, var_str ("threshold_calls"), &
         1, is_known = .true.)
    call var_list_set_int (global%var_list, var_str ("min_calls_per_channel"), &
         2, is_known = .true.)
    call var_list_set_int (global%var_list, var_str ("min_calls_per_bin"), &
         3, is_known = .true.)
    call var_list_set_int (global%var_list, var_str ("min_bins"), &
         4, is_known = .true.)
    call var_list_set_int (global%var_list, var_str ("max_bins"), &
         5, is_known = .true.)
    call var_list_set_log (global%var_list, var_str ("?stratified"), &
         .false., is_known = .true.)
    call var_list_set_log (global%var_list, var_str ("?use_vamp_equivalences"),&
         .false., is_known = .true.)
    call var_list_set_real (global%var_list, var_str ("channel_weights_power"),&
         4._default, is_known = .true.)
    call var_list_set_log (global%var_list, &
         var_str ("?vamp_history_global_verbose"), &
         .true., is_known = .true.)
    call var_list_set_log (global%var_list, &
         var_str ("?vamp_history_channels"), &
         .true., is_known = .true.)
    call var_list_set_log (global%var_list, &
         var_str ("?vamp_history_channels_verbose"), &
         .true., is_known = .true.)
    call var_list_set_log (global%var_list, var_str ("?stratified"), &
         .false., is_known = .true.)

    call dispatch_mci (mci, global, process_id)
    select type (mci)
    type is (mci_vamp_t)
       call mci%write (u)
       call mci%write_history_parameters (u)
    end select

    call mci%final ()
    deallocate (mci)
    
    write (u, "(A)")
    write (u, "(A)")  "* Allocate MCI as vamp_t, allow for negative weights"
    write (u, "(A)")    
    
    call var_list_set_string (global%var_list, &
         var_str ("$integration_method"), &
         var_str ("vamp"), is_known = .true.)
    call var_list_set_log (global%var_list, var_str ("?negative_weights"), &
         .true., is_known = .true.)
    
    call dispatch_mci (mci, global, process_id)
    select type (mci)       
    type is (mci_vamp_t)
       call mci%write (u)
       call mci%write_history_parameters (u)
    end select
    
    call mci%final ()
    deallocate (mci)
    
    call global%final ()

    write (u, "(A)")
    write (u, "(A)")  "* Test output end: dispatch_3"
    
  end subroutine dispatch_3
  
\eatline
\nwendcode{}\nwbegindocs{1898}\nwdocspar
\subsubsection{Select type: phase-space configuration object}
\nwenddocs{}\nwbegincode{1899}\moddef{Dispatch: execute tests}\plusendmoddef
  call test (dispatch_4, "dispatch_4", &
       "phase-space configuration", &
       u, results)
\nwendcode{}\nwbegincode{1900}\moddef{Dispatch: tests}\plusendmoddef
  subroutine dispatch_4 (u)
    integer, intent(in) :: u
    type(rt_data_t), target :: global
    class(phs_config_t), allocatable :: phs
    type(phs_parameters_t) :: phs_par
    type(mapping_defaults_t) :: mapping_defs
    
    write (u, "(A)")  "* Test output: dispatch_4"
    write (u, "(A)")  "*   Purpose: select phase-space configuration method"
    write (u, "(A)")

    call global%global_init ()

    write (u, "(A)")  "* Allocate PHS as phs_single_t"
    write (u, "(A)")

    call var_list_set_string (global%var_list, &
         var_str ("$phs_method"), &
         var_str ("single"), is_known = .true.)
    call dispatch_phs (phs, global, var_str ("dispatch_4"))
    call phs%write (u)

    call phs%final ()
    deallocate (phs)
    
    write (u, "(A)")
    write (u, "(A)")  "* Allocate PHS as phs_wood_t"
    write (u, "(A)")

    call var_list_set_string (global%var_list, &
         var_str ("$phs_method"), &
         var_str ("wood"), is_known = .true.)
    call dispatch_phs (phs, global, var_str ("dispatch_4"))
    call phs%write (u)
          
    call phs%final ()
    deallocate (phs)
    
    write (u, "(A)")
    write (u, "(A)")  "* Setting parameters for phs_wood_t"
    write (u, "(A)")        

    phs_par%m_threshold_s = 123
    phs_par%m_threshold_t = 456
    phs_par%t_channel = 42
    phs_par%off_shell = 17
    phs_par%keep_nonresonant = .false.    
    mapping_defs%energy_scale = 987
    mapping_defs%invariant_mass_scale = 654
    mapping_defs%momentum_transfer_scale = 321
    mapping_defs%step_mapping = .false.   
    mapping_defs%step_mapping_exp = .false.       
    mapping_defs%enable_s_mapping = .true.       
    call dispatch_phs (phs, global, var_str ("dispatch_4"), &
         mapping_defs, phs_par)    
    call phs%write (u)    
        
    call phs%final ()

    call global%final ()
    
    write (u, "(A)")
    write (u, "(A)")  "* Test output end: dispatch_4"
    
  end subroutine dispatch_4
  
\eatline
\nwendcode{}\nwbegindocs{1901}\nwdocspar
\subsubsection{Select type: random number generator}
\nwenddocs{}\nwbegincode{1902}\moddef{Dispatch: execute tests}\plusendmoddef
  call test (dispatch_5, "dispatch_5", &
       "random-number generator", &
       u, results)
\nwendcode{}\nwbegincode{1903}\moddef{Dispatch: tests}\plusendmoddef
  subroutine dispatch_5 (u)
    integer, intent(in) :: u
    type(rt_data_t), target :: global
    class(rng_factory_t), allocatable :: rng_factory
    
    write (u, "(A)")  "* Test output: dispatch_5"
    write (u, "(A)")  "*   Purpose: select random-number generator"
    write (u, "(A)")

    call global%global_init ()

    write (u, "(A)")  "* Allocate RNG factory as rng_test_factory_t"
    write (u, "(A)")

    call var_list_set_string (global%var_list, &
         var_str ("$rng_method"), &
         var_str ("unit_test"), is_known = .true.)
    call var_list_set_int (global%var_list, &
         var_str ("seed"), 1, is_known = .true.)
    call dispatch_rng_factory (rng_factory, global)
    call rng_factory%write (u)
    deallocate (rng_factory)
    
    write (u, "(A)")
    write (u, "(A)")  "* Allocate RNG factory as rng_tao_factory_t"
    write (u, "(A)")

    call var_list_set_string (global%var_list, &
         var_str ("$rng_method"), &
         var_str ("tao"), is_known = .true.)
    call dispatch_rng_factory (rng_factory, global)
    call rng_factory%write (u)
    deallocate (rng_factory)
    
    call global%final ()
    
    write (u, "(A)")
    write (u, "(A)")  "* Test output end: dispatch_5"
    
  end subroutine dispatch_5
  
\eatline
\nwendcode{}\nwbegindocs{1904}\nwdocspar
\subsubsection{Phase-space configuration with file}
\nwenddocs{}\nwbegincode{1905}\moddef{Dispatch: execute tests}\plusendmoddef
  call test (dispatch_6, "dispatch_6", &
       "configure phase space using file", &
       u, results)
\nwendcode{}\nwbegincode{1906}\moddef{Dispatch: tests}\plusendmoddef
  subroutine dispatch_6 (u)
    integer, intent(in) :: u
    type(rt_data_t), target :: global
    type(os_data_t) :: os_data
    type(process_constants_t) :: process_data
    class(phs_config_t), allocatable :: phs
    integer :: u_phs
    
    write (u, "(A)")  "* Test output: dispatch_6"
    write (u, "(A)")  "*   Purpose: select 'wood' phase-space &
         &for a test process"
    write (u, "(A)")  "*            and read phs configuration from file"
    write (u, "(A)")

    write (u, "(A)")  "* Initialize a process"
    write (u, "(A)")

    call global%global_init ()

    call os_data_init (os_data)
    call syntax_model_file_init ()
    call global%select_model (var_str ("Test"))

    call syntax_phs_forest_init ()
    
    call init_test_process_data (var_str ("dispatch_6"), process_data)

    write (u, "(A)")  "* Write phase-space file"

    u_phs = free_unit ()
    open (u_phs, file = "dispatch_6.phs", action = "write", status = "replace")
    call write_test_phs_file (u_phs, var_str ("dispatch_6"))
    close (u_phs)

    write (u, "(A)")
    write (u, "(A)")  "* Allocate PHS as phs_wood_t"
    write (u, "(A)")

    call var_list_set_string (global%var_list, &
         var_str ("$phs_method"), &
         var_str ("wood"), is_known = .true.)
    call var_list_set_string (global%var_list, &
         var_str ("$phs_file"), &
         var_str ("dispatch_6.phs"), is_known = .true.)
    call dispatch_phs (phs, global, var_str ("dispatch_6"))

    call phs%init (process_data, global%model)
    call phs%configure (sqrts = 1000._default)

    call phs%write (u)
    write (u, "(A)")
    select type (phs)
    type is (phs_wood_config_t)
       call phs%write_forest (u)
    end select

    call phs%final ()

    call global%final ()
    call syntax_model_file_final ()
    
    write (u, "(A)")
    write (u, "(A)")  "* Test output end: dispatch_6"
    
  end subroutine dispatch_6
  
\eatline
\nwendcode{}\nwbegindocs{1907}\nwdocspar
\subsubsection{Select type: structure-function data}
\nwenddocs{}\nwbegincode{1908}\moddef{Dispatch: execute tests}\plusendmoddef
  call test (dispatch_7, "dispatch_7", &
       "structure-function data", &
       u, results)
\nwendcode{}\nwbegincode{1909}\moddef{Dispatch: tests}\plusendmoddef
  subroutine dispatch_7 (u)
    integer, intent(in) :: u
    type(rt_data_t), target :: global
    type(os_data_t) :: os_data
    type(string_t) :: prt, sf_method
    type(sf_prop_t) :: sf_prop
    class(sf_data_t), allocatable :: data
    type(pdg_array_t), dimension(1) :: pdg_in
    type(pdg_array_t), dimension(1,1) :: pdg_prc
    type(pdg_array_t), dimension(1) :: pdg_out
    integer, dimension(:), allocatable :: pdg1
    
    write (u, "(A)")  "* Test output: dispatch_7"
    write (u, "(A)")  "*   Purpose: select and configure &
         &structure function data"
    write (u, "(A)")

    call global%global_init ()
    
    call os_data_init (os_data)
    call syntax_model_file_init ()
    call global%select_model (var_str ("QCD"))
    
    call reset_interaction_counter ()
    call var_list_set_real (global%var_list, var_str ("sqrts"), &
         14000._default, is_known = .true.)
    prt = "p"
    call global%beam_structure%init_sf ([prt, prt], [1])
    pdg_in = 2212
    
    write (u, "(A)")  "* Allocate data as sf_pdf_builtin_t"
    write (u, "(A)")

    sf_method = "pdf_builtin"
    call dispatch_sf_data &
         (data, sf_method, [1], sf_prop, global, pdg_in, pdg_prc)
    call data%write (u)

    call data%get_pdg_out (pdg_out)
    pdg1 = pdg_out(1)
    write (u, "(A)")
    write (u, "(1x,A,99(1x,I0))")  "PDG(out) = ", pdg1

    deallocate (data)
    
    write (u, "(A)")
    write (u, "(A)")  "* Allocate data for different PDF set"
    write (u, "(A)")

    pdg_in = 2212
    
    call var_list_set_string (global%var_list, var_str ("$pdf_builtin_set"), &
         var_str ("CTEQ6M"), is_known = .true.)
    sf_method = "pdf_builtin"
    call dispatch_sf_data &
         (data, sf_method, [1], sf_prop, global, pdg_in, pdg_prc)
    call data%write (u)

    call data%get_pdg_out (pdg_out)
    pdg1 = pdg_out(1)
    write (u, "(A)")
    write (u, "(1x,A,99(1x,I0))")  "PDG(out) = ", pdg1

    deallocate (data)
    
    call global%final ()
    call syntax_model_file_final ()
    
    write (u, "(A)")
    write (u, "(A)")  "* Test output end: dispatch_7"
    
  end subroutine dispatch_7
  
\eatline
\nwendcode{}\nwbegindocs{1910}\nwdocspar
\subsubsection{Beam structure}
\nwenddocs{}\nwbegincode{1911}\moddef{Dispatch: execute tests}\plusendmoddef
  call test (dispatch_8, "dispatch_8", &
       "beam structure", &
       u, results)
\nwendcode{}\nwbegincode{1912}\moddef{Dispatch: tests}\plusendmoddef
  subroutine dispatch_8 (u)
    integer, intent(in) :: u
    type(rt_data_t), target :: global
    type(os_data_t) :: os_data
    type(flavor_t), dimension(2) :: flv
    type(sf_config_t), dimension(:), allocatable :: sf_config
    type(sf_prop_t) :: sf_prop
    type(sf_channel_t), dimension(:), allocatable :: sf_channel
    type(phs_channel_collection_t) :: coll
    type(string_t) :: sf_string
    integer :: i
    type(pdg_array_t), dimension (2,1) :: pdg_prc
    
    write (u, "(A)")  "* Test output: dispatch_8"
    write (u, "(A)")  "*   Purpose: configure a structure-function chain"
    write (u, "(A)")

    call global%global_init ()
    
    call os_data_init (os_data)
    call syntax_model_file_init ()
    call global%select_model (var_str ("QCD"))
    
    write (u, "(A)")  "* Allocate LHC beams with PDF builtin"
    write (u, "(A)")

    call flv(1)%init (PROTON, global%model)
    call flv(2)%init (PROTON, global%model)

    call reset_interaction_counter ()
    call var_list_set_real (global%var_list, var_str ("sqrts"), &
         14000._default, is_known = .true.)
         
    call global%beam_structure%init_sf (flv%get_name (), [1])
    call global%beam_structure%set_sf (1, 1, var_str ("pdf_builtin"))
    
    call dispatch_sf_config (sf_config, sf_prop, global, pdg_prc)
    do i = 1, size (sf_config)
       call sf_config(i)%write (u)
    end do

    call dispatch_sf_channels (sf_channel, sf_string, sf_prop, coll, global)
    write (u, "(1x,A)")  "Mapping configuration:"
    do i = 1, size (sf_channel)
       write (u, "(2x)", advance = "no")
       call sf_channel(i)%write (u)
    end do

    write (u, "(A)")
    write (u, "(A)")  "* Allocate ILC beams with CIRCE1"
    write (u, "(A)")

    call global%select_model (var_str ("QED"))
    call flv(1)%init ( ELECTRON, global%model)
    call flv(2)%init (-ELECTRON, global%model)

    call reset_interaction_counter ()
    call var_list_set_real (global%var_list, var_str ("sqrts"), &
         500._default, is_known = .true.)
    call var_list_set_log (global%var_list, var_str ("?circe1_generate"), &
         .false., is_known = .true.)
         
    call global%beam_structure%init_sf (flv%get_name (), [1])
    call global%beam_structure%set_sf (1, 1, var_str ("circe1"))
    
    call dispatch_sf_config (sf_config, sf_prop, global, pdg_prc)
    do i = 1, size (sf_config)
       call sf_config(i)%write (u)
    end do

    call dispatch_sf_channels (sf_channel, sf_string, sf_prop, coll, global)
    write (u, "(1x,A)")  "Mapping configuration:"
    do i = 1, size (sf_channel)
       write (u, "(2x)", advance = "no")
       call sf_channel(i)%write (u)
    end do

    write (u, "(A)")
    write (u, "(A)")  "* Cleanup"

    call global%final ()
    call syntax_model_file_final ()
    
    write (u, "(A)")
    write (u, "(A)")  "* Test output end: dispatch_8"
    
  end subroutine dispatch_8
  
\eatline
\nwendcode{}\nwbegindocs{1913}\nwdocspar
\subsubsection{Event I/O}
\nwenddocs{}\nwbegincode{1914}\moddef{Dispatch: execute tests}\plusendmoddef
  call test (dispatch_9, "dispatch_9", &
       "event I/O", &
       u, results)
\nwendcode{}\nwbegincode{1915}\moddef{Dispatch: tests}\plusendmoddef
  subroutine dispatch_9 (u)
    integer, intent(in) :: u
    type(rt_data_t), target :: global
    class(eio_t), allocatable :: eio
    
    write (u, "(A)")  "* Test output: dispatch_9"
    write (u, "(A)")  "*   Purpose: allocate an event I/O (eio) stream"
    write (u, "(A)")

    call syntax_model_file_init ()
    call global%global_init ()
    call global%init_fallback_model &
         (var_str ("SM_hadrons"), var_str ("SM_hadrons.mdl"))
    
    write (u, "(A)")  "* Allocate as raw"
    write (u, "(A)")
    
    call dispatch_eio (eio, var_str ("raw"), global)

    call eio%write (u)

    call eio%final ()
    deallocate (eio)

    write (u, "(A)")
    write (u, "(A)")  "* Allocate as checkpoints:"
    write (u, "(A)")
    
    call dispatch_eio (eio, var_str ("checkpoint"), global)

    call eio%write (u)

    call eio%final ()
    deallocate (eio)

    write (u, "(A)")
    write (u, "(A)")  "* Allocate as LHEF:"
    write (u, "(A)")
    
    call var_list_set_string (global%var_list, var_str ("$lhef_extension"), &
         var_str ("lhe_custom"), is_known = .true.)
    call dispatch_eio (eio, var_str ("lhef"), global)

    call eio%write (u)

    call eio%final ()
    deallocate (eio)

    write (u, "(A)")
    write (u, "(A)")  "* Allocate as HepMC:"
    write (u, "(A)")
    
    call dispatch_eio (eio, var_str ("hepmc"), global)

    call eio%write (u)

    call eio%final ()
    deallocate (eio)

    write (u, "(A)")
    write (u, "(A)")  "* Allocate as weight_stream"
    write (u, "(A)")
    
    call dispatch_eio (eio, var_str ("weight_stream"), global)

    call eio%write (u)

    call eio%final ()
    deallocate (eio)

    write (u, "(A)")
    write (u, "(A)")  "* Allocate as debug format"
    write (u, "(A)")
    
    call var_list_set_log (global%var_list, var_str ("?debug_verbose"), &
         .false., is_known = .true.)
    call dispatch_eio (eio, var_str ("debug"), global)

    call eio%write (u)

    write (u, "(A)")
    write (u, "(A)")  "* Cleanup"

    call eio%final ()
    call global%final ()
    call syntax_model_file_final ()
    
    write (u, "(A)")
    write (u, "(A)")  "* Test output end: dispatch_9"
    
  end subroutine dispatch_9
  
\eatline
\nwendcode{}\nwbegindocs{1916}\nwdocspar
\subsubsection{Update process core parameters}
This test dispatches a process core, temporarily modifies parameters,
then restores the original.
\nwenddocs{}\nwbegincode{1917}\moddef{Dispatch: execute tests}\plusendmoddef
  call test (dispatch_10, "dispatch_10", &
       "process core update", &
       u, results)
\nwendcode{}\nwbegincode{1918}\moddef{Dispatch: tests}\plusendmoddef
  subroutine dispatch_10 (u)
    integer, intent(in) :: u
    type(string_t), dimension(2) :: prt_in, prt_out
    type(rt_data_t), target :: global
    class(prc_core_def_t), allocatable :: core_def
    class(prc_core_t), allocatable :: core, saved_core
    type(var_list_t), pointer :: var_list
    
    write (u, "(A)")  "* Test output: dispatch_10"
    write (u, "(A)")  "*   Purpose: select process configuration method,"
    write (u, "(A)")  "             allocate process core,"
    write (u, "(A)")  "             temporarily reset parameters"
    write (u, "(A)")

    call syntax_model_file_init ()
    call global%global_init ()

    prt_in = [var_str ("a"), var_str ("b")]
    prt_out = [var_str ("c"), var_str ("d")]

    write (u, "(A)")  "* Allocate core as prc_omega_t"
    write (u, "(A)")

    call var_list_set_string (global%var_list, var_str ("$method"), &
         var_str ("omega"), is_known = .true.)
    call dispatch_core_def (core_def, prt_in, prt_out, global)

    call global%select_model (var_str ("Test"))

    call dispatch_core (core, core_def, global%model)
    call core_def%allocate_driver (core%driver, var_str (""))

    select type (core)
    type is (prc_omega_t)
       call core%write (u)
    end select
    
    write (u, "(A)")
    write (u, "(A)")  "* Update core with modified model and helicity selection"
    write (u, "(A)")

    var_list => global%get_var_list_ptr ()
    call var_list_set_real (var_list, var_str ("gy"), 2._default, &
         is_known = .true.)
    call global%model%update_parameters ()

    call var_list_set_log (global%var_list, &
         var_str ("?helicity_selection_active"), &
         .true., is_known = .true.)
    call var_list_set_real (global%var_list, &
         var_str ("helicity_selection_threshold"), &
         2e10_default, is_known = .true.)
    call var_list_set_int (global%var_list, &
         var_str ("helicity_selection_cutoff"), &
         5, is_known = .true.)
    
    call dispatch_core_update (core, global%model, &
         global%get_helicity_selection (), &
         saved_core = saved_core)
    select type (core)
    type is (prc_omega_t)
       call core%write (u)
    end select
    
    write (u, "(A)")
    write (u, "(A)")  "* Restore core from save"
    write (u, "(A)")

    call dispatch_core_restore (core, saved_core)
    select type (core)
    type is (prc_omega_t)
       call core%write (u)
    end select
    
    call global%final ()
    call syntax_model_file_final ()
    
    write (u, "(A)")
    write (u, "(A)")  "* Test output end: dispatch_10"
    
  end subroutine dispatch_10
  
\eatline
\nwendcode{}\nwbegindocs{1919}\nwdocspar
\subsubsection{QCD Coupling}
This test dispatches an {\tt{}qcd} object, which is used to compute the
(running) coupling by one of several possible methods.
\nwenddocs{}\nwbegincode{1920}\moddef{Dispatch: execute tests}\plusendmoddef
  call test (dispatch_11, "dispatch_11", &
       "QCD coupling", &
       u, results)
\nwendcode{}\nwbegincode{1921}\moddef{Dispatch: tests}\plusendmoddef
  subroutine dispatch_11 (u)
    integer, intent(in) :: u
    type(rt_data_t), target :: global
    type(qcd_t) :: qcd
    type(var_list_t), pointer :: model_vars
    
    write (u, "(A)")  "* Test output: dispatch_11"
    write (u, "(A)")  "*   Purpose: select QCD coupling formula"
    write (u, "(A)")

    call syntax_model_file_init ()
    call global%global_init ()
    call global%select_model (var_str ("SM"))
    model_vars => global%get_var_list_ptr ()

    write (u, "(A)")  "* Allocate alpha_s as fixed"
    write (u, "(A)")

    call var_list_set_log (global%var_list, var_str ("?alpha_s_is_fixed"), &
         .true., is_known = .true.)
    call dispatch_qcd (qcd, global)
    call qcd%write (u)

    write (u, "(A)")
    write (u, "(A)")  "* Allocate alpha_s as running (built-in)"
    write (u, "(A)")
    
    call var_list_set_log (global%var_list, var_str ("?alpha_s_is_fixed"), &
         .false., is_known = .true.)
    call var_list_set_log (global%var_list, var_str ("?alpha_s_from_mz"), &
         .true., is_known = .true.)
    call var_list_set_int &
         (global%var_list, var_str ("alpha_s_order"), 1, is_known = .true.)
    call var_list_set_real &
         (model_vars, var_str ("alphas"), 0.1234_default, &
          is_known=.true.)
    call var_list_set_real &
         (model_vars, var_str ("mZ"), 91.234_default, &
          is_known=.true.)
    call dispatch_qcd (qcd, global)
    call qcd%write (u)
    
    write (u, "(A)")
    write (u, "(A)")  "* Allocate alpha_s as running (built-in, Lambda defined)"
    write (u, "(A)")
    
    call var_list_set_log (global%var_list, var_str ("?alpha_s_from_mz"), &
         .false., is_known = .true.)
    call var_list_set_log (global%var_list, &
         var_str ("?alpha_s_from_lambda_qcd"), &
         .true., is_known = .true.)
    call var_list_set_real &
         (global%var_list, var_str ("lambda_qcd"), 250.e-3_default, &
          is_known=.true.)
    call var_list_set_int &
         (global%var_list, var_str ("alpha_s_order"), 2, is_known = .true.)
    call var_list_set_int &
         (global%var_list, var_str ("alpha_s_nf"), 4, is_known = .true.)
    call dispatch_qcd (qcd, global)
    call qcd%write (u)
    
    write (u, "(A)")
    write (u, "(A)")  "* Allocate alpha_s as running (using builtin PDF set)"
    write (u, "(A)")
    
    call var_list_set_log (global%var_list, &
         var_str ("?alpha_s_from_lambda_qcd"), &
         .false., is_known = .true.)
    call var_list_set_log &
         (global%var_list, var_str ("?alpha_s_from_pdf_builtin"), &
         .true., is_known = .true.)
    call dispatch_qcd (qcd, global)
    call qcd%write (u)
    
    call global%final ()
    call syntax_model_file_final ()
    
    write (u, "(A)")
    write (u, "(A)")  "* Test output end: dispatch_11"
    
  end subroutine dispatch_11
  
\eatline
\nwendcode{}\nwbegindocs{1922}\nwdocspar
\subsubsection{Shower settings}
This test dispatches an {\tt{}shower{\char95}settings} object, which is used 
to steer the initial and final state.
\nwenddocs{}\nwbegincode{1923}\moddef{Dispatch: execute tests}\plusendmoddef
  call test (dispatch_12, "dispatch_12", &
       "Shower settings", &
       u, results)
\nwendcode{}\nwbegincode{1924}\moddef{Dispatch: tests}\plusendmoddef
  subroutine dispatch_12 (u)
    integer, intent(in) :: u
    type(rt_data_t), target :: global
    type(shower_settings_t) :: shower_settings
    
    write (u, "(A)")  "* Test output: dispatch_12"
    write (u, "(A)")  "*   Purpose: setting ISR/FSR shower"
    write (u, "(A)")

    write (u, "(A)")  "* Default settings"    
    write (u, "(A)")    
    
    call global%global_init ()
    call var_list_set_log (global%var_list, var_str ("?alpha_s_is_fixed"), &
         .true., is_known = .true.)
    call dispatch_shower (shower_settings, global)
    call write_separator (u)
    call shower_settings%write (u)
    call write_separator (u)

    write (u, "(A)")
    write (u, "(A)")  "* Switch on ISR/FSR showers, hadronization"
    write (u, "(A)")  "      and MLM matching"
    write (u, "(A)")
    
    call var_list_set_log (global%var_list, var_str ("?ps_fsr_active"), &
         .true., is_known = .true.)
    call var_list_set_log (global%var_list, var_str ("?ps_isr_active"), &
         .true., is_known = .true.)
    call var_list_set_log (global%var_list, var_str ("?hadronization_active"), &
         .true., is_known = .true.)    
    call var_list_set_log (global%var_list, var_str ("?mlm_matching"), &
         .true., is_known = .true.)        
    call var_list_set_int &
         (global%var_list, var_str ("ps_max_n_flavors"), 4, is_known = .true.)
    call var_list_set_real &
         (global%var_list, var_str ("ps_isr_z_cutoff"), 0.1234_default, &
          is_known=.true.)
    call var_list_set_real (global%var_list, &
         var_str ("mlm_etamax"), 3.456_default, is_known=.true.)
    call var_list_set_string (global%var_list, &
         var_str ("$ps_PYTHIA_PYGIVE"), var_str ("abcdefgh"), is_known=.true.)    
    call dispatch_shower (shower_settings, global)
    call write_separator (u)
    call shower_settings%write (u)
    call write_separator (u)
    
    call global%final ()
    
    write (u, "(A)")
    write (u, "(A)")  "* Test output end: dispatch_12"
    
  end subroutine dispatch_12
  
\eatline
\nwendcode{}\nwbegindocs{1925}\nwdocspar
\subsubsection{Event transforms}
This test dispatches an {\tt{}evt} (event transform) object.  Currently,
the only nontrivial transform is the partonic decay chain..
\nwenddocs{}\nwbegincode{1926}\moddef{Dispatch: execute tests}\plusendmoddef
  call test (dispatch_13, "dispatch_13", &
       "event transforms", &
       u, results)
\nwendcode{}\nwbegincode{1927}\moddef{Dispatch: tests}\plusendmoddef
  subroutine dispatch_13 (u)
    integer, intent(in) :: u
    type(rt_data_t), target :: global
    class(evt_t), pointer :: evt
    type(var_list_t), pointer :: model_vars
    
    write (u, "(A)")  "* Test output: dispatch_13"
    write (u, "(A)")  "*   Purpose: configure event transform"
    write (u, "(A)")

    call syntax_model_file_init ()
    call global%global_init ()
    call global%init_fallback_model &
         (var_str ("SM_hadrons"), var_str ("SM_hadrons.mdl"))
    model_vars => global%get_var_list_ptr ()    

    write (u, "(A)")  "* Partonic decays"
    write (u, "(A)")

    call dispatch_evt_decay (evt, global)
    select type (evt)
    type is (evt_decay_t)
       call evt%write (u, show_decay_tree = .true., verbose = .true.)
    end select

    call evt%final ()
    deallocate (evt)

    write (u, "(A)")
    write (u, "(A)")  "* Shower"
    write (u, "(A)")

    call var_list_set_log (global%var_list, var_str ("?allow_shower"), .true., &
         is_known = .true.)
    call dispatch_evt_shower (evt, global)
    select type (evt)
    type is (evt_shower_t)
       call evt%write (u)
       call write_separator (u, 2)
    end select

    call evt%final ()
    deallocate (evt)

    call global%final ()
    call syntax_model_file_final ()
    
    write (u, "(A)")
    write (u, "(A)")  "* Test output end: dispatch_13"
    
  end subroutine dispatch_13
  
\eatline
\nwendcode{}\nwbegindocs{1928}\nwdocspar
\subsubsection{SLHA interface}
This rather trivial sets all input values for the SLHA interface
to {\tt{}false}.
\nwenddocs{}\nwbegincode{1929}\moddef{Dispatch: execute tests}\plusendmoddef
  call test (dispatch_14, "dispatch_14", &
       "SLHA interface", &
       u, results)
\nwendcode{}\nwbegincode{1930}\moddef{Dispatch: tests}\plusendmoddef
  subroutine dispatch_14 (u)
    integer, intent(in) :: u
    type(rt_data_t), target :: global
    logical :: input, spectrum, decays
    
    write (u, "(A)")  "* Test output: dispatch_14"
    write (u, "(A)")  "*   Purpose: SLHA interface settings"
    write (u, "(A)")

    write (u, "(A)")  "* Default settings"    
    write (u, "(A)")    
    
    call global%global_init ()
    call dispatch_slha (global, &
         input = input, spectrum = spectrum, decays = decays)

    write (u, "(A,1x,L1)")  " slha_read_input     =", input
    write (u, "(A,1x,L1)")  " slha_read_spectrum  =", spectrum   
    write (u, "(A,1x,L1)")  " slha_read_decays    =", decays

    call global%final ()
    call global%global_init ()
    
    write (u, "(A)")
    write (u, "(A)")  "* Set all entries to [false]"    
    write (u, "(A)")        
            
    call var_list_set_log (global%var_list, var_str ("?slha_read_input"), &
         .false., is_known = .true.)
    call var_list_set_log (global%var_list, var_str ("?slha_read_spectrum"), &
         .false., is_known = .true.)
    call var_list_set_log (global%var_list, var_str ("?slha_read_decays"), &
         .false., is_known = .true.)    

    call dispatch_slha (global, &
         input = input, spectrum = spectrum, decays = decays)

    write (u, "(A,1x,L1)")  " slha_read_input     =", input
    write (u, "(A,1x,L1)")  " slha_read_spectrum  =", spectrum   
    write (u, "(A,1x,L1)")  " slha_read_decays    =", decays
    
    call global%final ()
    
    write (u, "(A)")
    write (u, "(A)")  "* Test output end: dispatch_14"
    
  end subroutine dispatch_14
  
\eatline
\nwendcode{}\nwbegindocs{1931}\nwdocspar
\clearpage
%%%%%%%%%%%%%%%%%%%%%%%%%%%%%%%%%%%%%%%%%%%%%%%%%%%%%%%%%%%%%%%%%%%%%%%%
\section{Process Configuration}
This module communicates between the toplevel command structure with
its runtime data set and the process-library handling modules which
collect the definition of individual processes.  Its primary purpose
is to select from the available matrix-element generating methods and
configure the entry in the process library accordingly.
\nwenddocs{}\nwbegincode{1932}\moddef{\code{}process{\char95}configurations.f90\edoc{}}\endmoddef
\LA{}File header\RA{}

module process_configurations
  
\LA{}Use strings\RA{}
  use unit_tests
  use diagnostics

  use variables
  use models
  use prc_core_def
  use particle_specifiers
  use process_libraries
  use prclib_stacks
  use prc_test
  use prc_omega
  use rt_data
  use dispatch

  use prc_gosam
  
\LA{}Standard module head\RA{}

\LA{}Process configurations: public\RA{}

\LA{}Process configurations: types\RA{}

contains
  
\LA{}Process configurations: procedures\RA{}

\LA{}Process configurations: tests\RA{}

end module process_configurations
\eatline
\nwendcode{}\nwbegindocs{1933}\nwdocspar
\subsection{Data Type}
\nwenddocs{}\nwbegincode{1934}\moddef{Process configurations: public}\endmoddef
  public :: process_configuration_t
\nwendcode{}\nwbegincode{1935}\moddef{Process configurations: types}\endmoddef
  type :: process_configuration_t
     type(process_def_entry_t), pointer :: entry => null ()
     type(string_t) :: id
     integer :: num_id = 0
   contains
   \LA{}Process configurations: process configuration: TBP\RA{}
  end type process_configuration_t
  
\eatline
\nwendcode{}\nwbegindocs{1936}Initialize a process.  We only need the name, the number of incoming
particles, and the number of components.
\nwenddocs{}\nwbegincode{1937}\moddef{Process configurations: process configuration: TBP}\endmoddef
  procedure :: init => process_configuration_init
\nwendcode{}\nwbegincode{1938}\moddef{Process configurations: procedures}\endmoddef
  subroutine process_configuration_init &
       (config, prc_name, n_in, n_components, global)
    class(process_configuration_t), intent(out) :: config
    type(string_t), intent(in) :: prc_name
    integer, intent(in) :: n_in
    integer, intent(in) :: n_components 
    type(rt_data_t), intent(in) :: global
    type(model_t), pointer :: model
    model => global%model
    config%id = prc_name
    allocate (config%entry)
    if (global%var_list%is_known (var_str ("process_num_id"))) then
       config%num_id = &
            global%var_list%get_ival (var_str ("process_num_id"))
       call config%entry%init (prc_name, &
            model = model, n_in = n_in, n_components = n_components, &
            num_id = config%num_id, nlo_process = global%nlo_calculation)
    else
       call config%entry%init (prc_name, &
            model = model, n_in = n_in, n_components = n_components, &
            nlo_process = global%nlo_calculation)
    end if
  end subroutine process_configuration_init
    
\eatline
\nwendcode{}\nwbegindocs{1939}Initialize a process component.  The details depend on the process method,
which determines the type of the process component core.  We set the incoming
and outgoing particles (as strings, to be interpreted by the process driver).
All other information is taken from the variable list.

The dispatcher gets only the names of the particles.  The process
component definition gets the complete specifiers which contains a
polarization flag and names of decay processes, where applicable.
\nwenddocs{}\nwbegincode{1940}\moddef{Process configurations: process configuration: TBP}\plusendmoddef
  procedure :: setup_component => process_configuration_setup_component
\nwendcode{}\nwbegincode{1941}\moddef{Process configurations: procedures}\plusendmoddef
  subroutine process_configuration_setup_component &
       (config, i_component, prt_in, prt_out, global, &
        nlo_type, active_in)
    class(process_configuration_t), intent(inout) :: config
    integer, intent(in) :: i_component
    type(prt_spec_t), dimension(:), intent(in) :: prt_in
    type(prt_spec_t), dimension(:), intent(in) :: prt_out
    type(rt_data_t), intent(inout) :: global
    type(string_t), intent(in), optional :: nlo_type
    logical, intent(in), optional :: active_in
    type(string_t), dimension(:), allocatable :: prt_str_in
    type(string_t), dimension(:), allocatable :: prt_str_out
    class(prc_core_def_t), allocatable :: core_def
    type(string_t) :: method
    integer :: i
    logical :: active

    allocate (prt_str_in  (size (prt_in)))
    allocate (prt_str_out (size (prt_out)))
    forall (i = 1:size (prt_in))  prt_str_in(i)  = prt_in(i)% get_name ()
    forall (i = 1:size (prt_out)) prt_str_out(i) = prt_out(i)%get_name ()
    if (present (active_in)) then
      active = active_in
    else
      active = .true.
    end if

    call dispatch_core_def (core_def, prt_str_in, prt_str_out, &
                            global, config%id, nlo_type)
    method = &
         global%var_list%get_sval (var_str ("$method"))
    call config%entry%import_component (i_component, &
       n_out = size (prt_out), &
       prt_in = prt_in, &
       prt_out = prt_out, &
       method = method, &
       variant = core_def, &
       nlo_type = nlo_type, &
       active = active)
  end subroutine process_configuration_setup_component
  
\eatline
\nwendcode{}\nwbegindocs{1942}\nwdocspar
\nwenddocs{}\nwbegincode{1943}\moddef{Process configurations: process configuration: TBP}\plusendmoddef
  procedure :: set_component_associations => &
                   process_configuration_set_component_associations
\nwendcode{}\nwbegincode{1944}\moddef{Process configurations: procedures}\plusendmoddef
  subroutine process_configuration_set_component_associations &
       (config, i_list)
    class(process_configuration_t), intent(inout) :: config
    integer, intent(in), dimension(4) :: i_list 
    integer :: i_component
    do i_component = 1, config%entry%get_n_components ()
       if (any (i_list == i_component)) then
          call config%entry%set_associated_components (i_component, &
                 i_list(1), i_list(2), i_list(3), i_list(4))
       end if
    end do
  end subroutine process_configuration_set_component_associations

\eatline
\nwendcode{}\nwbegindocs{1945}Record a process configuration: append it to the currently selected process
definition library.
\nwenddocs{}\nwbegincode{1946}\moddef{Process configurations: process configuration: TBP}\plusendmoddef
  procedure :: record => process_configuration_record
\nwendcode{}\nwbegincode{1947}\moddef{Process configurations: procedures}\plusendmoddef
  subroutine process_configuration_record (config, global)
    class(process_configuration_t), intent(inout) :: config
    type(rt_data_t), intent(inout) :: global
    if (associated (global%prclib)) then
       call global%prclib%open ()
       call global%prclib%append (config%entry)
       if (config%num_id /= 0) then
          write (msg_buffer, "(5A,I0,A)") "Process library '", &
               char (global%prclib%get_name ()), &
               "': recorded process '", char (config%id), "' (", &
               config%num_id, ")"
       else
          write (msg_buffer, "(5A)") "Process library '", &
               char (global%prclib%get_name ()), &
               "': recorded process '", char (config%id), "'"          
       end if
       call msg_message ()
    else
       call msg_fatal ("Recording process '" // char (config%id) &
            // "': active process library undefined")
    end if
  end subroutine process_configuration_record
  
\eatline
\nwendcode{}\nwbegindocs{1948}\nwdocspar
\subsection{Test}
This is the master for calling self-test procedures.
\nwenddocs{}\nwbegincode{1949}\moddef{Process configurations: public}\plusendmoddef
  public :: process_configurations_test
\nwendcode{}\nwbegincode{1950}\moddef{Process configurations: tests}\endmoddef
  subroutine process_configurations_test (u, results)
    integer, intent(in) :: u
    type(test_results_t), intent(inout) :: results
  \LA{}Process configurations: execute tests\RA{}
  end subroutine process_configurations_test

\eatline
\nwendcode{}\nwbegindocs{1951}\nwdocspar
\subsubsection{Minimal setup}
The workflow for setting up a minimal process configuration with the
test matrix element method.

We wrap this in a public procedure, so we can reuse it in later modules.
The procedure prepares a process definition list for two processes
(one {\tt{}prc{\char95}test} and one {\tt{}omega} type) and appends this to the
process library stack in the global data set.

The {\tt{}mode} argument determines which processes to build.

The {\tt{}procname} argument replaces the predefined procname(s).
\nwenddocs{}\nwbegincode{1952}\moddef{Process configurations: public}\plusendmoddef
  public :: prepare_test_library
\nwendcode{}\nwbegincode{1953}\moddef{Process configurations: tests}\plusendmoddef
  subroutine prepare_test_library (global, libname, mode, procname)
    type(rt_data_t), intent(inout), target :: global
    type(string_t), intent(in) :: libname
    integer, intent(in) :: mode
    type(string_t), intent(in), dimension(:), optional :: procname
    type(prclib_entry_t), pointer :: lib
    type(string_t) :: prc_name
    type(string_t), dimension(:), allocatable :: prt_in, prt_out
    integer :: n_components
    type(process_configuration_t) :: prc_config

    if (.not. associated (global%prclib_stack%get_first_ptr ())) then
       allocate (lib)
       call lib%init (libname)
       call global%add_prclib (lib)
    end if

    if (btest (mode, 0)) then

       call global%select_model (var_str ("Test"))

       if (present (procname)) then
          prc_name = procname(1)
       else
          prc_name = "prc_config_a"
       end if
       n_components = 1
       allocate (prt_in (2), prt_out (2))
       prt_in = [var_str ("s"), var_str ("s")]
       prt_out = [var_str ("s"), var_str ("s")]

       call var_list_set_string (global%var_list, var_str ("$method"),&
            var_str ("unit_test"), is_known = .true.)

       call prc_config%init (prc_name, size (prt_in), n_components, global)
       call prc_config%setup_component (1, &
            new_prt_spec (prt_in), new_prt_spec (prt_out), global)
       call prc_config%record (global)

       deallocate (prt_in, prt_out)
       
    end if

    if (btest (mode, 1)) then

       call global%select_model (var_str ("QED"))

       if (present (procname)) then
          prc_name = procname(2)
       else
          prc_name = "prc_config_b"
       end if
       n_components = 1
       allocate (prt_in (2), prt_out (2))
       prt_in = [var_str ("e+"), var_str ("e-")]
       prt_out = [var_str ("m+"), var_str ("m-")]

       call var_list_set_string (global%var_list, var_str ("$method"),&
            var_str ("omega"), is_known = .true.)

       call prc_config%init (prc_name, size (prt_in), n_components, global)
       call prc_config%setup_component (1, &
            new_prt_spec (prt_in), new_prt_spec (prt_out), global)
       call prc_config%record (global)

       deallocate (prt_in, prt_out)
       
    end if
    
    if (btest (mode, 2)) then

       call global%select_model (var_str ("Test"))

       if (present (procname)) then
          prc_name = procname(1)
       else
          prc_name = "prc_config_a"
       end if
       n_components = 1
       allocate (prt_in (1), prt_out (2))
       prt_in = [var_str ("s")]
       prt_out = [var_str ("f"), var_str ("fbar")]

       call var_list_set_string (global%var_list, var_str ("$method"),&
            var_str ("unit_test"), is_known = .true.)

       call prc_config%init (prc_name, size (prt_in), n_components, global)
       call prc_config%setup_component (1, &
            new_prt_spec (prt_in), new_prt_spec (prt_out), global)
       call prc_config%record (global)

       deallocate (prt_in, prt_out)
       
    end if

  end subroutine prepare_test_library
    
\eatline
\nwendcode{}\nwbegindocs{1954}The actual test: the previous procedure with some prelude and postlude.
In the global variable list, just before printing we reset the
variables where the value may depend on the system and run environment.
\nwenddocs{}\nwbegincode{1955}\moddef{Process configurations: execute tests}\endmoddef
  call test (process_configurations_1, "process_configurations_1", &
       "test processes", &
       u, results)
\nwendcode{}\nwbegincode{1956}\moddef{Process configurations: tests}\plusendmoddef
  subroutine process_configurations_1 (u)
    integer, intent(in) :: u
    type(rt_data_t), target :: global
    
    write (u, "(A)")  "* Test output: process_configurations_1"
    write (u, "(A)")  "*   Purpose: configure test processes"
    write (u, "(A)")

    call syntax_model_file_init ()

    call global%global_init ()
    call var_list_set_log (global%var_list, var_str ("?omega_openmp"), &
         .false., is_known = .true.)

    write (u, "(A)")  "* Configure processes as prc_test, model Test"
    write (u, "(A)")  "*                     and omega, model QED"
    write (u, *)

    call var_list_set_int (global%var_list, var_str ("process_num_id"), &
         42, is_known = .true.)
    call prepare_test_library (global, var_str ("prc_config_lib_1"), 3)

    global%os_data%fc = "Fortran-compiler"
    global%os_data%fcflags = "Fortran-flags"

    call global%write_libraries (u)

    call global%final ()
    call syntax_model_file_final ()
    
    write (u, "(A)")
    write (u, "(A)")  "* Test output end: process_configurations_1"
    
  end subroutine process_configurations_1
  
\eatline
\nwendcode{}\nwbegindocs{1957}\nwdocspar
\subsubsection{\oMega\ options}
Slightly extended example where we pass \oMega\ options to the
library.  The {\tt{}prepare{\char95}test{\char95}library} contents are spelled out.
\nwenddocs{}\nwbegincode{1958}\moddef{Process configurations: execute tests}\plusendmoddef
  call test (process_configurations_2, "process_configurations_2", &
       "omega options", &
       u, results)
\nwendcode{}\nwbegincode{1959}\moddef{Process configurations: tests}\plusendmoddef
  subroutine process_configurations_2 (u)
    integer, intent(in) :: u
    type(rt_data_t), target :: global
    
    type(string_t) :: libname
    type(prclib_entry_t), pointer :: lib
    type(string_t) :: prc_name
    type(string_t), dimension(:), allocatable :: prt_in, prt_out
    integer :: n_components
    type(process_configuration_t) :: prc_config

    write (u, "(A)")  "* Test output: process_configurations_2"
    write (u, "(A)")  "*   Purpose: configure test processes with options"
    write (u, "(A)")

    call syntax_model_file_init ()

    call global%global_init ()
    
    write (u, "(A)")  "* Configure processes as omega, model QED"
    write (u, *)

    libname = "prc_config_lib_2"
    
    allocate (lib)
    call lib%init (libname)
    call global%add_prclib (lib)

    call global%select_model (var_str ("QED"))

    prc_name = "prc_config_c"
    n_components = 2
    allocate (prt_in (2), prt_out (2))
    prt_in = [var_str ("e+"), var_str ("e-")]
    prt_out = [var_str ("m+"), var_str ("m-")]

    call var_list_set_string (global%var_list, var_str ("$method"),&
         var_str ("omega"), is_known = .true.)
    call var_list_set_log (global%var_list, var_str ("?omega_openmp"), &
         .false., is_known = .true.)

    call prc_config%init (prc_name, size (prt_in), n_components, global)

    call var_list_set_log (global%var_list, var_str ("?report_progress"), &
         .true., is_known = .true.)
    call prc_config%setup_component (1, &
         new_prt_spec (prt_in), new_prt_spec (prt_out), global)

    call var_list_set_log (global%var_list, var_str ("?report_progress"), &
         .false., is_known = .true.)
    call var_list_set_log (global%var_list, var_str ("?omega_openmp"), &
         .true., is_known = .true.)
    call var_list_set_string (global%var_list, var_str ("$restrictions"),&
         var_str ("3+4~A"), is_known = .true.)
    call var_list_set_string (global%var_list, var_str ("$omega_flags"), &
         var_str ("-fusion:progress_file omega_prc_config.log"), &
         is_known = .true.)
    call prc_config%setup_component (2, &
         new_prt_spec (prt_in), new_prt_spec (prt_out), global)
    
    call prc_config%record (global)

    deallocate (prt_in, prt_out)
    
    global%os_data%fc = "Fortran-compiler"
    global%os_data%fcflags = "Fortran-flags"

    call global%write_vars (u, [ &
         var_str ("$model_name"), &
         var_str ("$method"), &
         var_str ("?report_progress"), &
         var_str ("$restrictions"), &
         var_str ("$omega_flags")])
    write (u, "(A)")
    call global%write_libraries (u)

    call global%final ()
    call syntax_model_file_final ()
    
    write (u, "(A)")
    write (u, "(A)")  "* Test output end: process_configurations_2"
    
  end subroutine process_configurations_2
  
\eatline
\nwendcode{}\nwbegindocs{1960}\nwdocspar
\clearpage
%%%%%%%%%%%%%%%%%%%%%%%%%%%%%%%%%%%%%%%%%%%%%%%%%%%%%%%%%%%%%%%%%%%%%%%%
\section{Compilation}
This module manages compilation and loading of of process libraries.  It is
needed as a separate module because integration depends on it.
\nwenddocs{}\nwbegincode{1961}\moddef{\code{}compilations.f90\edoc{}}\endmoddef
\LA{}File header\RA{}

module compilations

\LA{}Use strings\RA{}
  use io_units
  use unit_tests
  use system_defs, only: TAB
  use diagnostics
  use os_interface
  use md5
  use variables
  use model_data
  use models
  use process_libraries
  use prclib_stacks
  use rt_data
  use process_configurations

\LA{}Standard module head\RA{}

\LA{}Compilations: public\RA{}

\LA{}Compilations: types\RA{}

contains

\LA{}Compilations: procedures\RA{}

\LA{}Compilations: tests\RA{}

end module compilations
\eatline
\nwendcode{}\nwbegindocs{1962}\nwdocspar
\subsection{The data type}
The compilation item handles the compilation and loading of a single
process library.
\nwenddocs{}\nwbegincode{1963}\moddef{Compilations: types}\endmoddef
  type :: compilation_item_t
     type(string_t) :: libname
     type(string_t) :: static_external_tag
     type(process_library_t), pointer :: lib => null ()
     logical :: recompile_library = .false.
   contains
   \LA{}Compilations: compilation item: TBP\RA{}
  end type compilation_item_t

\eatline
\nwendcode{}\nwbegindocs{1964}Initialize:
\nwenddocs{}\nwbegincode{1965}\moddef{Compilations: compilation item: TBP}\endmoddef
  procedure :: init => compilation_item_init
\nwendcode{}\nwbegincode{1966}\moddef{Compilations: procedures}\endmoddef
  subroutine compilation_item_init (comp, libname, stack, var_list)
    class(compilation_item_t), intent(out) :: comp
    type(string_t), intent(in) :: libname
    type(prclib_stack_t), intent(inout) :: stack
    type(var_list_t), intent(in) :: var_list
    comp%libname = libname
    comp%lib => stack%get_library_ptr (comp%libname)
    if (.not. associated (comp%lib)) then
       call msg_fatal ("Process library '" // char (comp%libname) &
            // "' has not been declared.")
    end if
    comp%recompile_library = &
         var_list%get_lval (var_str ("?recompile_library"))
  end subroutine compilation_item_init

\eatline
\nwendcode{}\nwbegindocs{1967}Compile the current library.  The {\tt{}force} flag has the
effect that we first delete any previous files, as far as accessible
by the current makefile.  It also guarantees that previous files not
accessible by a makefile will be overwritten.
\nwenddocs{}\nwbegincode{1968}\moddef{Compilations: compilation item: TBP}\plusendmoddef
  procedure :: compile => compilation_item_compile
\nwendcode{}\nwbegincode{1969}\moddef{Compilations: procedures}\plusendmoddef
  subroutine compilation_item_compile (comp, model, os_data, force, recompile)
    class(compilation_item_t), intent(inout) :: comp
    class(model_data_t), intent(in), target :: model
    type(os_data_t), intent(in) :: os_data
    logical, intent(in) :: force, recompile
    if (associated (comp%lib)) then
       call msg_message ("Process library '" &
            // char (comp%libname) // "': compiling ...")
       call comp%lib%configure (os_data)
       if (signal_is_pending ())  return
       call comp%lib%compute_md5sum (model)
       call comp%lib%write_makefile (os_data, force)
       if (signal_is_pending ())  return
       if (force) then
          call comp%lib%clean (os_data, distclean = .false.)
          if (signal_is_pending ())  return
       end if
       call comp%lib%write_driver (force)
       if (signal_is_pending ())  return
       if (recompile) then
          call comp%lib%load (os_data, keep_old_source = .true.)
          if (signal_is_pending ())  return
       end if
       call comp%lib%update_status (os_data)
    end if
  end subroutine compilation_item_compile

\eatline
\nwendcode{}\nwbegindocs{1970}Load the current library, just after compiling it.
\nwenddocs{}\nwbegincode{1971}\moddef{Compilations: compilation item: TBP}\plusendmoddef
  procedure :: load => compilation_item_load
\nwendcode{}\nwbegincode{1972}\moddef{Compilations: procedures}\plusendmoddef
  subroutine compilation_item_load (comp, os_data)
    class(compilation_item_t), intent(inout) :: comp
    type(os_data_t), intent(in) :: os_data
    if (associated (comp%lib)) then
       call comp%lib%load (os_data)
    end if
  end subroutine compilation_item_load

\eatline
\nwendcode{}\nwbegindocs{1973}Message as a separate call:
\nwenddocs{}\nwbegincode{1974}\moddef{Compilations: compilation item: TBP}\plusendmoddef
  procedure :: success => compilation_item_success
\nwendcode{}\nwbegincode{1975}\moddef{Compilations: procedures}\plusendmoddef
  subroutine compilation_item_success (comp)
    class(compilation_item_t), intent(in) :: comp
    if (associated (comp%lib)) then
       call msg_message ("Process library '" // char (comp%libname) &
            // "': ... success.")
    else
       call msg_fatal ("Process library '" // char (comp%libname) &
            // "': ... failure.")
    end if
  end subroutine compilation_item_success

\eatline
\eatline
\nwendcode{}\nwbegindocs{1976}\nwdocspar
\subsection{API for library compilation and loading}
This is a shorthand for compiling and loading a single library.  The
{\tt{}compilation{\char95}item} object is used only internally.

The {\tt{}global} data set may actually be local to the caller.  The
compilation affects the library specified by its name if it is on the
stack, but it does not reset the currently selected library.
\nwenddocs{}\nwbegincode{1977}\moddef{Compilations: public}\endmoddef
  public :: compile_library
\nwendcode{}\nwbegincode{1978}\moddef{Compilations: procedures}\plusendmoddef
  subroutine compile_library (libname, global)
    type(string_t), intent(in) :: libname
    type(rt_data_t), intent(inout), target :: global
    type(compilation_item_t) :: comp
    logical :: force, recompile
    force = &
         global%var_list%get_lval (var_str ("?rebuild_library"))
    recompile = &
         global%var_list%get_lval (var_str ("?recompile_library"))
    if (associated (global%model)) then
       call comp%init (libname, global%prclib_stack, global%var_list)
       call comp%compile (global%model, global%os_data, force, recompile)
       if (signal_is_pending ())  return
       call comp%load (global%os_data)
       if (signal_is_pending ())  return
    else
       call msg_fatal ("Process library compilation: " &
            // " model is undefined.")
    end if
    call comp%success ()
  end subroutine compile_library

\eatline
\nwendcode{}\nwbegindocs{1979}\nwdocspar
\subsection{Compiling static executable}
This object handles the creation of a static executable which should
contain a set of static process libraries.
\nwenddocs{}\nwbegincode{1980}\moddef{Compilations: types}\plusendmoddef
  type :: compilation_t
     private
     type(string_t) :: exe_name
     type(string_t), dimension(:), allocatable :: lib_name
   contains
   \LA{}Compilations: compilation: TBP\RA{}
  end type compilation_t

\eatline
\nwendcode{}\nwbegindocs{1981}Output.
\nwenddocs{}\nwbegincode{1982}\moddef{Compilations: compilation: TBP}\endmoddef
  procedure :: write => compilation_write
\nwendcode{}\nwbegincode{1983}\moddef{Compilations: procedures}\plusendmoddef
  subroutine compilation_write (object, unit)
    class(compilation_t), intent(in) :: object
    integer, intent(in), optional :: unit
    integer :: u, i
    u = given_output_unit (unit)
    write (u, "(1x,A)")  "Compilation object:"
    write (u, "(3x,3A)")  "executable        = '", &
         char (object%exe_name), "'"
    write (u, "(3x,A)", advance="no")  "process libraries ="
    do i = 1, size (object%lib_name)
       write (u, "(1x,3A)", advance="no")  "'", char (object%lib_name(i)), "'"
    end do
    write (u, *)
  end subroutine compilation_write
  
\eatline
\nwendcode{}\nwbegindocs{1984}Initialize: we know the names of the executable and of the libraries.
\nwenddocs{}\nwbegincode{1985}\moddef{Compilations: compilation: TBP}\plusendmoddef
  procedure :: init => compilation_init
\nwendcode{}\nwbegincode{1986}\moddef{Compilations: procedures}\plusendmoddef
  subroutine compilation_init (compilation, exe_name, lib_name)
    class(compilation_t), intent(out) :: compilation
    type(string_t), intent(in) :: exe_name
    type(string_t), dimension(:), intent(in) :: lib_name
    compilation%exe_name = exe_name
    allocate (compilation%lib_name (size (lib_name)))
    compilation%lib_name = lib_name
  end subroutine compilation_init
    
\eatline
\nwendcode{}\nwbegindocs{1987}Write the dispatcher subroutine for the compiled libraries.  Also
write a subroutine which returns the names of the compiled libraries.
\nwenddocs{}\nwbegincode{1988}\moddef{Compilations: compilation: TBP}\plusendmoddef
  procedure :: write_dispatcher => compilation_write_dispatcher
\nwendcode{}\nwbegincode{1989}\moddef{Compilations: procedures}\plusendmoddef
  subroutine compilation_write_dispatcher (compilation)
    class(compilation_t), intent(in) :: compilation
    type(string_t) :: file
    integer :: u, i
    file = compilation%exe_name // "_prclib_dispatcher.f90"
    call msg_message ("Static executable '" // char (compilation%exe_name) &
         // "': writing library dispatcher")
    u = free_unit ()
    open (u, file = char (file), status="replace", action="write")
    write (u, "(3A)")  "! Whizard: process libraries for executable '", &
         char (compilation%exe_name), "'"
    write (u, "(A)")  "! Automatically generated file, do not edit"
    write (u, "(A)")  "subroutine dispatch_prclib_static " // &
         "(driver, basename, modellibs_ldflags)"
    write (u, "(A)")  "  use iso_varying_string, string_t => varying_string"
    write (u, "(A)")  "  use prclib_interfaces"
    do i = 1, size (compilation%lib_name)
       associate (lib_name => compilation%lib_name(i))
         write (u, "(A)")  "  use " // char (lib_name) // "_driver"
       end associate
    end do
    write (u, "(A)")  "  implicit none"
    write (u, "(A)")  "  class(prclib_driver_t), intent(inout), allocatable &
         &:: driver"
    write (u, "(A)")  "  type(string_t), intent(in) :: basename"
    write (u, "(A)")  "  logical, intent(in), optional :: " // &
         "modellibs_ldflags"
    write (u, "(A)")  "  select case (char (basename))"
    do i = 1, size (compilation%lib_name)
       associate (lib_name => compilation%lib_name(i))
         write (u, "(3A)")  "  case ('", char (lib_name), "')"
         write (u, "(3A)")  "     allocate (", char (lib_name), "_driver_t &
              &:: driver)"
       end associate
    end do
    write (u, "(A)")  "  end select"
    write (u, "(A)")  "end subroutine dispatch_prclib_static"
    write (u, *)
    write (u, "(A)")  "subroutine get_prclib_static (libname)"
    write (u, "(A)")  "  use iso_varying_string, string_t => varying_string"
    write (u, "(A)")  "  implicit none"
    write (u, "(A)")  "  type(string_t), dimension(:), intent(inout), &
         &allocatable :: libname"
    write (u, "(A,I0,A)")  "  allocate (libname (", &
         size (compilation%lib_name), "))"
    do i = 1, size (compilation%lib_name)
       associate (lib_name => compilation%lib_name(i))
         write (u, "(A,I0,A,A,A)")  "  libname(", i, ") = '", &
              char (lib_name), "'"
       end associate
    end do
    write (u, "(A)")  "end subroutine get_prclib_static"
    close (u)
  end subroutine compilation_write_dispatcher
    
\eatline
\nwendcode{}\nwbegindocs{1990}Write the Makefile subroutine for the compiled libraries.
\nwenddocs{}\nwbegincode{1991}\moddef{Compilations: compilation: TBP}\plusendmoddef
  procedure :: write_makefile => compilation_write_makefile
\nwendcode{}\nwbegincode{1992}\moddef{Compilations: procedures}\plusendmoddef
  subroutine compilation_write_makefile (compilation, os_data, ext_libtag)
    class(compilation_t), intent(in) :: compilation
    type(os_data_t), intent(in) :: os_data
    type(string_t), intent(in), optional :: ext_libtag
    type(string_t) :: file, ext_tag
    integer :: u, i
    if (present (ext_libtag)) then
       ext_tag = ext_libtag
    else
       ext_tag = ""
    end if
    file = compilation%exe_name // ".makefile"
    call msg_message ("Static executable '" // char (compilation%exe_name) &
         // "': writing makefile")
    u = free_unit ()
    open (u, file = char (file), status="replace", action="write")
    write (u, "(3A)")  "# WHIZARD: Makefile for executable '", &
         char (compilation%exe_name), "'"
    write (u, "(A)")  "# Automatically generated file, do not edit"
    write (u, "(A)") ""
    write (u, "(A)") "# Executable name"
    write (u, "(A)") "EXE = " // char (compilation%exe_name)
    write (u, "(A)") ""
    write (u, "(A)") "# Compiler"
    write (u, "(A)") "FC = " // char (os_data%fc)
    write (u, "(A)") ""
    write (u, "(A)") "# Included libraries"
    write (u, "(A)") "FCINCL = " // char (os_data%whizard_includes)
    write (u, "(A)") ""
    write (u, "(A)") "# Compiler flags"
    write (u, "(A)") "FCFLAGS = " // char (os_data%fcflags)
    write (u, "(A)") "LDFLAGS = " // char (os_data%ldflags)
    write (u, "(A)") "LDFLAGS_STATIC = " // char (os_data%ldflags_static)   
    write (u, "(A)") "LDFLAGS_HEPMC = " // char (os_data%ldflags_hepmc)
    write (u, "(A)") "LDFLAGS_HOPPET = " // char (os_data%ldflags_hoppet)    
    write (u, "(A)") "LDWHIZARD = " // char (os_data%whizard_ldflags)
    write (u, "(A)") ""
    write (u, "(A)") "# Libtool"
    write (u, "(A)") "LIBTOOL = " // char (os_data%whizard_libtool)
    write (u, "(A)") "FCOMPILE = $(LIBTOOL) --tag=FC --mode=compile"
    write (u, "(A)") "LINK = $(LIBTOOL) --tag=FC --mode=link"
    write (u, "(A)") ""
    write (u, "(A)") "# Compile commands (default)"
    write (u, "(A)") "LTFCOMPILE = $(FCOMPILE) $(FC) -c $(FCINCL) $(FCFLAGS)"
    write (u, "(A)") ""
    write (u, "(A)") "# Default target"
    write (u, "(A)") "all: link"
    write (u, "(A)") ""
    write (u, "(A)") "# Libraries"
    do i = 1, size (compilation%lib_name)
       associate (lib_name => compilation%lib_name(i))
         write (u, "(A)") "LIBRARIES += " // char (lib_name) // ".la"
         write (u, "(A)") char (lib_name) // ".la:"
         write (u, "(A)") TAB // "$(MAKE) -f " // char (lib_name) // ".makefile"
       end associate
    end do
    write (u, "(A)") ""
    write (u, "(A)") "# Library dispatcher"
    write (u, "(A)") "DISP = $(EXE)_prclib_dispatcher"
    write (u, "(A)") "$(DISP).lo: $(DISP).f90 $(LIBRARIES)"
    write (u, "(A)") TAB // "$(LTFCOMPILE) $<"
    write (u, "(A)") ""
    write (u, "(A)") "# Executable"
    write (u, "(A)") "$(EXE): $(DISP).lo $(LIBRARIES)"
    write (u, "(A)") TAB // "$(LINK) $(FC) -static-libtool-libs $(FCFLAGS) \\"
    write (u, "(A)") TAB // "   $(LDWHIZARD) $(LDFLAGS) \\" 
    write (u, "(A)") TAB // "   -o $(EXE) $^ \\"
    write (u, "(A)") TAB // "   $(LDFLAGS_HEPMC) $(LDFLAGS_HOPPET) \\"
    write (u, "(A)") TAB // "   $(LDFLAGS_STATIC)" // char (ext_tag)
    write (u, "(A)") ""
    write (u, "(A)") "# Main targets"
    write (u, "(A)") "link: compile $(EXE)"
    write (u, "(A)") "compile: $(LIBRARIES) $(DISP).lo"
    write (u, "(A)") ".PHONY: link compile"
    write (u, "(A)") ""
    write (u, "(A)") "# Cleanup targets"
    write (u, "(A)") "clean-exe:"
    write (u, "(A)") TAB // "rm -f $(EXE)"
    write (u, "(A)") "clean-objects:"
    write (u, "(A)") TAB // "rm -f $(DISP).lo"
    write (u, "(A)") "clean-source:"
    write (u, "(A)") TAB // "rm -f $(DISP).f90"
    write (u, "(A)") "clean-makefile:"
    write (u, "(A)") TAB // "rm -f $(EXE).makefile"
    write (u, "(A)") ""
    write (u, "(A)") "clean: clean-exe clean-objects clean-source"
    write (u, "(A)") "distclean: clean clean-makefile"
    write (u, "(A)") ".PHONY: clean distclean"
    close (u)
  end subroutine compilation_write_makefile
    
\eatline
\nwendcode{}\nwbegindocs{1993}Compile the dispatcher source code.
\nwenddocs{}\nwbegincode{1994}\moddef{Compilations: compilation: TBP}\plusendmoddef
  procedure :: make_compile => compilation_make_compile
\nwendcode{}\nwbegincode{1995}\moddef{Compilations: procedures}\plusendmoddef
  subroutine compilation_make_compile (compilation, os_data)
    class(compilation_t), intent(in) :: compilation
    type(os_data_t), intent(in) :: os_data
    call os_system_call ("make compile " // os_data%makeflags &
         // " -f " // compilation%exe_name // ".makefile")
  end subroutine compilation_make_compile
  
\eatline
\nwendcode{}\nwbegindocs{1996}Link the dispatcher together with all matrix-element code and the
\whizard\ and \oMega\ main libraries, to generate a static executable.
\nwenddocs{}\nwbegincode{1997}\moddef{Compilations: compilation: TBP}\plusendmoddef
  procedure :: make_link => compilation_make_link
\nwendcode{}\nwbegincode{1998}\moddef{Compilations: procedures}\plusendmoddef
  subroutine compilation_make_link (compilation, os_data)
    class(compilation_t), intent(in) :: compilation
    type(os_data_t), intent(in) :: os_data
    call os_system_call ("make link " // os_data%makeflags &
         // " -f " // compilation%exe_name // ".makefile")
  end subroutine compilation_make_link
  
\eatline
\nwendcode{}\nwbegindocs{1999}Cleanup.
\nwenddocs{}\nwbegincode{2000}\moddef{Compilations: compilation: TBP}\plusendmoddef
  procedure :: make_clean_exe => compilation_make_clean_exe
\nwendcode{}\nwbegincode{2001}\moddef{Compilations: procedures}\plusendmoddef
  subroutine compilation_make_clean_exe (compilation, os_data)
    class(compilation_t), intent(in) :: compilation
    type(os_data_t), intent(in) :: os_data
    call os_system_call ("make clean-exe " // os_data%makeflags &
         // " -f " // compilation%exe_name // ".makefile")
  end subroutine compilation_make_clean_exe
  
\eatline
\nwendcode{}\nwbegindocs{2002}\nwdocspar
\subsection{API for executable compilation}
This is a shorthand for compiling and loading an executable, including
the enclosed libraries.  The {\tt{}compilation} object is used only internally.

The {\tt{}global} data set may actually be local to the caller.  The
compilation affects the library specified by its name if it is on the
stack, but it does not reset the currently selected library.
\nwenddocs{}\nwbegincode{2003}\moddef{Compilations: public}\plusendmoddef
  public :: compile_executable
\nwendcode{}\nwbegincode{2004}\moddef{Compilations: procedures}\plusendmoddef
  subroutine compile_executable (exename, libname, global)
    type(string_t), intent(in) :: exename
    type(string_t), dimension(:), intent(in) :: libname    
    type(rt_data_t), intent(inout), target :: global
    type(compilation_t) :: compilation
    type(compilation_item_t) :: item
    type(string_t) :: ext_libtag
    logical :: force, recompile
    integer :: i
    ext_libtag = ""
    force = &
         global%var_list%get_lval (var_str ("?rebuild_library"))
    recompile = &
         global%var_list%get_lval (var_str ("?recompile_library"))
    call compilation%init (exename, [libname])
    if (signal_is_pending ())  return
    call compilation%write_dispatcher ()
    if (signal_is_pending ())  return
    do i = 1, size (libname)
       call item%init (libname(i), global%prclib_stack, global%var_list)
       call item%compile (global%model, global%os_data, &
            force=force, recompile=recompile)
       ext_libtag = "" // item%lib%get_static_modelname (global%os_data)       
       if (signal_is_pending ())  return
       call item%success ()
    end do
    call compilation%write_makefile (global%os_data, ext_libtag)
    if (signal_is_pending ())  return    
    call compilation%make_compile (global%os_data)
    if (signal_is_pending ())  return
    call compilation%make_link (global%os_data)
  end subroutine compile_executable

\eatline
\nwendcode{}\nwbegindocs{2005}\nwdocspar
\subsection{Test}
This is the master for calling self-test procedures.
\nwenddocs{}\nwbegincode{2006}\moddef{Compilations: public}\plusendmoddef
  public :: compilations_test
\nwendcode{}\nwbegincode{2007}\moddef{Compilations: tests}\endmoddef
  subroutine compilations_test (u, results)
    integer, intent(in) :: u
    type(test_results_t), intent(inout) :: results
  \LA{}Compilations: execute tests\RA{}
end subroutine compilations_test

\eatline
\nwendcode{}\nwbegindocs{2008}\nwdocspar
\subsubsection{Intrinsic Matrix Element}
Compile an intrinsic test matrix element ({\tt{}prc{\char95}test} type).

Note: In this and the following test, we reset the Fortran compiler and flag
variables immediately before they are printed, so the test is portable.
\nwenddocs{}\nwbegincode{2009}\moddef{Compilations: execute tests}\endmoddef
  call test (compilations_1, "compilations_1", &
       "intrinsic test processes", &
       u, results)
\nwendcode{}\nwbegincode{2010}\moddef{Compilations: tests}\plusendmoddef
  subroutine compilations_1 (u)
    integer, intent(in) :: u
    type(string_t) :: libname, procname
    type(rt_data_t), target :: global
    
    write (u, "(A)")  "* Test output: compilations_1"
    write (u, "(A)")  "*   Purpose: configure and compile test process"
    write (u, "(A)")

    call syntax_model_file_init ()

    call global%global_init ()

    libname = "compilation_1"
    procname = "prc_comp_1"
    call prepare_test_library (global, libname, 1, [procname])

    call compile_library (libname, global)
    
    call global%write_libraries (u)

    call global%final ()
    call syntax_model_file_final ()
    
    write (u, "(A)")
    write (u, "(A)")  "* Test output end: compilations_1"
    
  end subroutine compilations_1
  
\eatline
\nwendcode{}\nwbegindocs{2011}\nwdocspar
\subsubsection{External Matrix Element}
Compile an external test matrix element ({\tt{}omega} type)
\nwenddocs{}\nwbegincode{2012}\moddef{Compilations: execute tests}\plusendmoddef
  call test (compilations_2, "compilations_2", &
       "external process (omega)", &
       u, results)
\nwendcode{}\nwbegincode{2013}\moddef{Compilations: tests}\plusendmoddef
  subroutine compilations_2 (u)
    integer, intent(in) :: u
    type(string_t) :: libname, procname
    type(rt_data_t), target :: global
    
    write (u, "(A)")  "* Test output: compilations_2"
    write (u, "(A)")  "*   Purpose: configure and compile test process"
    write (u, "(A)")

    call syntax_model_file_init ()

    call global%global_init ()
    call var_list_set_log (global%var_list, var_str ("?omega_openmp"), &
         .false., is_known = .true.)

    libname = "compilation_2"
    procname = "prc_comp_2"
    call prepare_test_library (global, libname, 2, [procname,procname])

    call compile_library (libname, global)
    
    call global%write_libraries (u, libpath = .false.)

    call global%final ()
    call syntax_model_file_final ()
    
    write (u, "(A)")
    write (u, "(A)")  "* Test output end: compilations_2"
    
  end subroutine compilations_2
  
\eatline
\nwendcode{}\nwbegindocs{2014}\nwdocspar
\subsubsection{External Matrix Element}
Compile an external test matrix element ({\tt{}omega} type) and
create driver files for a static executable.
\nwenddocs{}\nwbegincode{2015}\moddef{Compilations: execute tests}\plusendmoddef
  call test (compilations_3, "compilations_3", &
       "static executable: driver", &
       u, results)
\nwendcode{}\nwbegincode{2016}\moddef{Compilations: tests}\plusendmoddef
  subroutine compilations_3 (u)
    integer, intent(in) :: u
    type(string_t) :: libname, procname, exename
    type(rt_data_t), target :: global
    type(compilation_t) :: compilation
    integer :: u_file
    character(80) :: buffer

    write (u, "(A)")  "* Test output: compilations_3"
    write (u, "(A)")  "*   Purpose: make static executable"
    write (u, "(A)")

    write (u, "(A)")  "* Initialize library"
    write (u, "(A)")

    call syntax_model_file_init ()

    call global%global_init ()
    call var_list_set_log (global%var_list, var_str ("?omega_openmp"), &
         .false., is_known = .true.)

    libname = "compilations_3_lib"
    procname = "prc_comp_3"
    exename = "compilations_3"
    
    call prepare_test_library (global, libname, 2, [procname,procname])

    call compilation%init (exename, [libname])
    call compilation%write (u)

    write (u, "(A)")
    write (u, "(A)")  "* Write dispatcher"
    write (u, "(A)")

    call compilation%write_dispatcher ()

    u_file = free_unit ()
    open (u_file, file = char (exename) // "_prclib_dispatcher.f90", &
         status = "old", action = "read")
    do
       read (u_file, "(A)", end = 1)  buffer
       write (u, "(A)")  trim (buffer)
    end do
1   close (u_file)

    write (u, "(A)")
    write (u, "(A)")  "* Write Makefile"
    write (u, "(A)")

    associate (os_data => global%os_data)
      os_data%fc = "fortran-compiler"
      os_data%whizard_includes = "my-includes"
      os_data%fcflags = "my-fcflags"
      os_data%ldflags = "my-ldflags"
      os_data%ldflags_static = "my-ldflags-static"
      os_data%ldflags_hepmc = "my-ldflags-hepmc"
      os_data%ldflags_hoppet = "my-ldflags-hoppet"      
      os_data%whizard_ldflags = "my-ldwhizard"
      os_data%whizard_libtool = "my-libtool"
    end associate

    call compilation%write_makefile (global%os_data)

    open (u_file, file = char (exename) // ".makefile", &
         status = "old", action = "read")
    do
       read (u_file, "(A)", end = 2)  buffer
       write (u, "(A)")  trim (buffer)
    end do
2   close (u_file)

    write (u, "(A)")
    write (u, "(A)")  "* Cleanup"

    call global%final ()
    call syntax_model_file_final ()
    
    write (u, "(A)")
    write (u, "(A)")  "* Test output end: compilations_3"
    
  end subroutine compilations_3
  
\eatline
\nwendcode{}\nwbegindocs{2017}\nwdocspar
\subsection{Test static build}
The tests for building a static executable are separate, since they
should be skipped if the \whizard\ build itself has static libraries
disabled.
\nwenddocs{}\nwbegincode{2018}\moddef{Compilations: public}\plusendmoddef
  public :: compilations_static_test
\nwendcode{}\nwbegincode{2019}\moddef{Compilations: tests}\plusendmoddef
  subroutine compilations_static_test (u, results)
    integer, intent(in) :: u
    type(test_results_t), intent(inout) :: results
  \LA{}Compilations: static tests\RA{}
end subroutine compilations_static_test

\eatline
\nwendcode{}\nwbegindocs{2020}\nwdocspar
\subsubsection{External Matrix Element}
Compile an external test matrix element ({\tt{}omega} type) and
incorporate this in a new static WHIZARD executable.
\nwenddocs{}\nwbegincode{2021}\moddef{Compilations: static tests}\endmoddef
  call test (compilations_static_1, "compilations_static_1", &
       "static executable: compilation", &
       u, results)
\nwendcode{}\nwbegincode{2022}\moddef{Compilations: tests}\plusendmoddef
  subroutine compilations_static_1 (u)
    integer, intent(in) :: u
    type(string_t) :: libname, procname, exename
    type(rt_data_t), target :: global
    type(compilation_item_t) :: item
    type(compilation_t) :: compilation
    logical :: exist

    write (u, "(A)")  "* Test output: compilations_static_1"
    write (u, "(A)")  "*   Purpose: make static executable"
    write (u, "(A)")

    write (u, "(A)")  "* Initialize library"

    call syntax_model_file_init ()

    call global%global_init ()
    call var_list_set_log (global%var_list, var_str ("?omega_openmp"), &
         .false., is_known = .true.)

    libname = "compilations_static_1_lib"
    procname = "prc_comp_stat_1"
    exename = "compilations_static_1"
    
    call prepare_test_library (global, libname, 2, [procname,procname])

    call compilation%init (exename, [libname])

    write (u, "(A)")
    write (u, "(A)")  "* Write dispatcher"

    call compilation%write_dispatcher ()

    write (u, "(A)")
    write (u, "(A)")  "* Write Makefile"

    call compilation%write_makefile (global%os_data)

    write (u, "(A)")
    write (u, "(A)")  "* Build libraries"

    call item%init (libname, global%prclib_stack, global%var_list)
    call item%compile &
         (global%model, global%os_data, force=.true., recompile=.false.)
    call item%success ()

    write (u, "(A)")
    write (u, "(A)")  "* Check executable (should be absent)"
    write (u, "(A)")
    
    call compilation%make_clean_exe (global%os_data)
    inquire (file = char (exename), exist = exist)
    write (u, "(A,A,L1)")  char (exename), " exists = ", exist

    write (u, "(A)")
    write (u, "(A)")  "* Build executable"
    write (u, "(A)")

    call compilation%make_compile (global%os_data)
    call compilation%make_link (global%os_data)

    write (u, "(A)")  "* Check executable (should be present)"
    write (u, "(A)")
    
    inquire (file = char (exename), exist = exist)
    write (u, "(A,A,L1)")  char (exename), " exists = ", exist

    write (u, "(A)")
    write (u, "(A)")  "* Cleanup"

    call compilation%make_clean_exe (global%os_data)

    call global%final ()
    call syntax_model_file_final ()
    
    write (u, "(A)")
    write (u, "(A)")  "* Test output end: compilations_static_1"
    
  end subroutine compilations_static_1
  
\eatline
\nwendcode{}\nwbegindocs{2023}\nwdocspar
\subsubsection{External Matrix Element}
Compile an external test matrix element ({\tt{}omega} type) and
incorporate this in a new static WHIZARD executable.  In this version,
we use the wrapper {\tt{}compile{\char95}executable} procedure.
\nwenddocs{}\nwbegincode{2024}\moddef{Compilations: static tests}\plusendmoddef
  call test (compilations_static_2, "compilations_static_2", &
       "static executable: shortcut", &
       u, results)
\nwendcode{}\nwbegincode{2025}\moddef{Compilations: tests}\plusendmoddef
  subroutine compilations_static_2 (u)
    integer, intent(in) :: u
    type(string_t) :: libname, procname, exename
    type(rt_data_t), target :: global
    type(compilation_item_t) :: item
    type(compilation_t) :: compilation
    logical :: exist
    integer :: u_file

    write (u, "(A)")  "* Test output: compilations_static_2"
    write (u, "(A)")  "*   Purpose: make static executable"
    write (u, "(A)")

    write (u, "(A)")  "* Initialize library and compile"
    write (u, "(A)")

    call syntax_model_file_init ()

    call global%global_init ()
    call var_list_set_log (global%var_list, var_str ("?omega_openmp"), &
         .false., is_known = .true.)

    libname = "compilations_static_2_lib"
    procname = "prc_comp_stat_2"
    exename = "compilations_static_2"
    
    call prepare_test_library (global, libname, 2, [procname,procname])

    call compile_executable (exename, [libname], global)

    write (u, "(A)")  "* Check executable (should be present)"
    write (u, "(A)")
    
    inquire (file = char (exename), exist = exist)
    write (u, "(A,A,L1)")  char (exename), " exists = ", exist

    write (u, "(A)")
    write (u, "(A)")  "* Cleanup"

    u_file = free_unit ()
    open (u_file, file = char (exename), status = "old", action = "write")
    close (u_file, status = "delete")

    call global%final ()
    call syntax_model_file_final ()
    
    write (u, "(A)")
    write (u, "(A)")  "* Test output end: compilations_static_2"
    
  end subroutine compilations_static_2
  
\eatline
\nwendcode{}\nwbegindocs{2026}\nwdocspar
\clearpage
%%%%%%%%%%%%%%%%%%%%%%%%%%%%%%%%%%%%%%%%%%%%%%%%%%%%%%%%%%%%%%%%%%%%%%%%
\section{Integration}
This module manages phase space setup, matrix-element evaluation and
integration, as far as it is not done by lower-level routines, in particular
in the {\tt{}processes} module.
\nwenddocs{}\nwbegincode{2027}\moddef{\code{}integrations.f90\edoc{}}\endmoddef
\LA{}File header\RA{}

module integrations

\LA{}Use kinds\RA{}
\LA{}Use strings\RA{}
  use io_units
  use unit_tests
  use diagnostics
  use os_interface
  use cputime
  use sm_qcd
  use ifiles
  use lexers
  use parser
  use model_data
  use flavors
  use pdg_arrays
  use variables
  use expr_base
  use eval_trees
  use models
  use interactions
  use sf_mappings
  use sf_base
  use phs_base
  use mappings
  use phs_forests
  use phs_wood
  use rng_base
  use mci_base
  use process_libraries
  use prc_core
  use processes
  use process_stacks
  use models
  use iterations
  use rt_data
  use dispatch
  use process_configurations
  use compilations

  use process_constants
  use prc_omega
  use prc_gosam
  use nlo_data

\LA{}Standard module head\RA{}

\LA{}Integrations: public\RA{}

\LA{}Integrations: types\RA{}

contains

\LA{}Integrations: procedures\RA{}

\LA{}Integrations: tests\RA{}

end module integrations
\eatline
\nwendcode{}\nwbegindocs{2028}\nwdocspar
\subsection{The integration type}
This type holds all relevant data, the integration methods operates on this.
In contrast to the {\tt{}simulation{\char95}t} introduced later, the {\tt{}integration{\char95}t}
applies to a single process.
\nwenddocs{}\nwbegincode{2029}\moddef{Integrations: public}\endmoddef
  public :: integration_t
\nwendcode{}\nwbegincode{2030}\moddef{Integrations: types}\endmoddef
  type :: integration_t
    private
    type(string_t) :: process_id
    type(string_t) :: run_id
    type(process_t), pointer :: process => null ()
    type(var_list_t), pointer :: model_vars => null ()
    type(qcd_t) :: qcd
    logical :: rebuild_phs = .false.
    logical :: ignore_phs_mismatch = .false.
    logical :: phs_only = .false.
    logical :: process_has_me = .true.
    integer :: n_calls_test = 0
    logical :: vis_history = .true.
    type(string_t) :: history_filename
    type(string_t) :: log_filename
    logical :: combined_integration = .false.
   contains
   \LA{}Integrations: integration: TBP\RA{} 
  end type integration_t

\eatline
\nwendcode{}\nwbegindocs{2031}\nwdocspar
\nwenddocs{}\nwbegindocs{2032}\nwdocspar
\subsection{Initialization}
Initialization, first part: Create a process entry.
Push it on the stack if the {\tt{}global} environment is supplied.
\nwenddocs{}\nwbegincode{2033}\moddef{Integrations: integration: TBP}\endmoddef
  procedure :: create_process => integration_create_process
\nwendcode{}\nwbegincode{2034}\moddef{Integrations: procedures}\endmoddef
  subroutine integration_create_process (intg, process_id, global)
    class(integration_t), intent(out) :: intg
    type(rt_data_t), intent(inout), optional, target :: global
    type(string_t), intent(in) :: process_id
    type(process_entry_t), pointer :: process_entry
    intg%process_id = process_id
    if (present (global)) then
       allocate (process_entry)
       intg%process => process_entry%process_t
       call global%process_stack%push (process_entry)
    else
       allocate (process_t :: intg%process)
    end if
    intg%model_vars => null ()
  end subroutine integration_create_process

\eatline
\nwendcode{}\nwbegindocs{2035}Initialization, second part: Initialize the process object, using the local
environment.  We allocate a RNG factory and a QCD object.
We also fetch a pointer to the model that the process uses.  The
process initializer will create a snapshot of that model.

This procedure
does not modify the {\tt{}local} stack directly.  The intent(inout) attribute for
the {\tt{}local} data set is due to the random generator seed which may be
incremented during initialization.

NOTE: Changes to model parameters within the current context are respected
only if the process model coincides with the current model.  This is the usual
case.  If not, we read
the model from the global model library, which has default parameters.  To
become more flexible, we should implement a local model library which records
local changes to currently inactive models.
\nwenddocs{}\nwbegincode{2036}\moddef{Integrations: integration: TBP}\plusendmoddef
  procedure :: init_process => integration_init_process
\nwendcode{}\nwbegincode{2037}\moddef{Integrations: procedures}\plusendmoddef
  subroutine integration_init_process (intg, local)
    class(integration_t), intent(inout) :: intg
    type(rt_data_t), intent(inout), target :: local
    type(string_t) :: model_name
    type(model_t), pointer :: model
    class(model_data_t), pointer :: model_instance
    class(rng_factory_t), allocatable :: rng_factory
    if (.not. local%prclib%contains (intg%process_id)) then
       call msg_fatal ("Process '" // char (intg%process_id) // "' not found" &
            // " in library '" // char (local%prclib%get_name ()) // "'")
       return
    end if
    intg%run_id = &
         local%var_list%get_sval (var_str ("$run_id"))
    call dispatch_qcd (intg%qcd, local)
    call dispatch_rng_factory (rng_factory, local)
    model_name = local%prclib%get_model_name (intg%process_id)
    if (local%get_sval (var_str ("$model_name")) == model_name) then
       model => local%model
    else
       model => local%model_list%get_model_ptr (model_name)
    end if
    allocate (model_t :: model_instance)
    select type (model_instance)
    type is (model_t)
       call model_instance%init_instance (model)
       intg%model_vars => model_instance%get_var_list_ptr ()
    end select
    call intg%process%init (intg%process_id, intg%run_id, &
         local%prclib, &
         local%os_data, intg%qcd, rng_factory, model_instance)
  end subroutine integration_init_process
    
\eatline
\nwendcode{}\nwbegindocs{2038}Initialization, third part: complete process configuration.   
\nwenddocs{}\nwbegincode{2039}\moddef{Integrations: integration: TBP}\plusendmoddef
  procedure :: setup_process => integration_setup_process
\nwendcode{}\nwbegincode{2040}\moddef{Integrations: procedures}\plusendmoddef
  subroutine integration_setup_process (intg, local, verbose)
    class(integration_t), intent(inout) :: intg
    type(rt_data_t), intent(inout), target :: local
    logical, intent(in), optional :: verbose
    type(var_list_t), pointer :: var_list
    class(prc_core_t), allocatable :: core_template
    class(phs_config_t), allocatable :: phs_config_template
    class(phs_config_t), allocatable :: phs_config_template_real
    type(phs_parameters_t) :: phs_par
    type(mapping_defaults_t) :: mapping_defs
    class(mci_t), allocatable :: mci_template
    integer :: n_components, n_in, i_component
    type(pdg_array_t), dimension(:,:), allocatable :: pdg_prc
    type(process_component_def_t), pointer :: config
    type(helicity_selection_t), allocatable :: helicity_selection
    real(default) :: sqrts
    logical :: decay_rest_frame, use_color_factors
    type(sf_config_t), dimension(:), allocatable :: sf_config
    type(sf_prop_t) :: sf_prop
    type(sf_channel_t), dimension(:), allocatable :: sf_channel
    type(phs_channel_collection_t) :: phs_channel_collection
    logical :: sf_trace
    type(string_t) :: sf_string, sf_trace_file
    logical :: verb
    type(process_constants_t) :: prc_const
    integer :: i_born
    type(fks_template_t) :: fks_template
    type(gosam_template_t) :: gosam_template
    type(eval_tree_factory_t) :: expr_factory
  
    verb = .true.; if (present (verbose))  verb = verbose
    
    call intg%process%set_var_list (local%get_var_list_ptr ())
    var_list => intg%process%get_var_list_ptr ()

    intg%rebuild_phs = &
         var_list%get_lval (var_str ("?rebuild_phase_space"))
    intg%ignore_phs_mismatch = &
         .not. var_list%get_lval (var_str ("?check_phs_file"))
    intg%phs_only = &
         var_list%get_lval (var_str ("?phs_only"))
    phs_par%m_threshold_s = &
         var_list%get_rval (var_str ("phs_threshold_s"))
    phs_par%m_threshold_t = &
         var_list%get_rval (var_str ("phs_threshold_t"))
    phs_par%off_shell = &
         var_list%get_ival (var_str ("phs_off_shell"))
    phs_par%keep_nonresonant = &
         var_list%get_lval (var_str ("?phs_keep_nonresonant"))
    phs_par%t_channel = &
         var_list%get_ival (var_str ("phs_t_channel"))
    mapping_defs%energy_scale = &
         var_list%get_rval (var_str ("phs_e_scale"))
    mapping_defs%invariant_mass_scale = &
         var_list%get_rval (var_str ("phs_m_scale"))
    mapping_defs%momentum_transfer_scale = &
         var_list%get_rval (var_str ("phs_q_scale"))
    mapping_defs%step_mapping = &
         var_list%get_lval (var_str ("?phs_step_mapping"))
    mapping_defs%step_mapping_exp = &
         var_list%get_lval (var_str ("?phs_step_mapping_exp"))
    mapping_defs%enable_s_mapping = &
         var_list%get_lval (var_str ("?phs_s_mapping"))

    call dispatch_phs (phs_config_template, local, &
         intg%process_id, mapping_defs, phs_par)
    
    
    intg%n_calls_test = &
         var_list%get_ival (var_str ("n_calls_test"))

    !!! We avoid two dots in the filename due to a bug in certain MetaPost versions.
    if (intg%run_id /= "") then
       intg%history_filename = intg%process_id // "." // intg%run_id &
            // "-history"
       intg%log_filename = intg%process_id // "." // intg%run_id // ".log"
    else
       intg%history_filename = intg%process_id // "-history"
       intg%log_filename = intg%process_id // ".log"
    end if

    call dispatch_mci (mci_template, local, intg%process_id)

    if (verb) then
       call msg_message ("Initializing integration for process " &
            // char (intg%process_id) // ":")
       if (intg%run_id /= "") then
          call msg_message ("Run ID = " // '"' // char (intg%run_id) // '"')
       end if
    end if
    
    helicity_selection = local%get_helicity_selection ()

    intg%vis_history = &
         var_list%get_lval (var_str ("?vis_history"))
    use_color_factors = var_list%get_lval &
         (var_str ("?read_color_factors"))
    
    n_components = intg%process%get_n_components ()
    n_in = intg%process%get_n_in ()

    call gosam_template%init ()
    intg%combined_integration = var_list%get_lval (&
                                var_str ('?combined_nlo_integration'))
    
    do i_component = 1, n_components
       config => intg%process%get_component_def_ptr (i_component)
       call dispatch_core (core_template, config%get_core_def_ptr (), &
            intg%process%get_model_ptr (), helicity_selection, intg%qcd, &
            use_color_factors)
       select case (char (config%get_nlo_type ()))
       case ('Virtual')
         call gosam_template%set_loop (var_list%get_lval (&
                                       var_str ('?use_gosam_loops')))
         call intg%process%init_component &
            (i_component, core_template, mci_template, phs_config_template, &
             gosam_template = gosam_template)
         if (intg%combined_integration) &
            call intg%process%set_component_type (i_component, COMP_VIRT)
         call gosam_template%reset ()
       case ('Real')
         call gosam_template%set_real_trees (var_list%get_lval (&
                                             var_str ('?use_gosam_real_trees')))
         call dispatch_phs (phs_config_template_real, local, &
             intg%process_id, mapping_defs, phs_par, &
             var_str ('fks'))
         call dispatch_fks (fks_template, local)
         call intg%process%init_component &
            (i_component, core_template, mci_template, &
             phs_config_template_real, fks_template = fks_template, &
             gosam_template = gosam_template)
         if (intg%combined_integration) &
            call intg%process%set_component_type (i_component, COMP_REAL)
         call gosam_template%reset ()
       case ('Born')
         call intg%process%init_component &
            (i_component, core_template, mci_template, phs_config_template)
         i_born = config%get_associated_born ()
         if (intg%combined_integration) &
            call intg%process%set_component_type (i_component, COMP_MASTER)
       case ('Subtraction')
         call gosam_template%set_subtraction (var_list%get_lval (&
                                              var_str ('?use_gosam_correlations')))
         call intg%process%init_component &
             (i_component, core_template, mci_template, phs_config_template, &
              gosam_template = gosam_template)
         if (intg%combined_integration) &
            call intg%process%set_component_type (i_component, COMP_SUB)
         call gosam_template%reset ()
       case default
         call msg_fatal ("setup_process: NLO type not implemented!")
       end select
       deallocate (core_template)
    end do

    if (verb)  call intg%process%write (screen = .true.)
    
    intg%process_has_me = intg%process%has_matrix_element ()
    if (.not. intg%process_has_me) then
       call msg_warning ("Process '" &
            // char (intg%process_id) // "': matrix element vanishes")
    end if
    
    sqrts = local%get_sqrts ()
    decay_rest_frame = &
         var_list%get_lval (var_str ("?decay_rest_frame"))    
    if (intg%process_has_me) then
       if (local%beam_structure%is_set ()) then
          call intg%process%setup_beams_beam_structure &
               (local%beam_structure, sqrts, decay_rest_frame)
        else if (n_in == 2) then
          call intg%process%setup_beams_sqrts &
               (sqrts, local%beam_structure)
       else 
          call intg%process%setup_beams_decay &
               (decay_rest_frame, local%beam_structure)
       end if
    end if
    call intg%process%check_masses ()
    if (verb .and. intg%process_has_me) then
       call intg%process%beams_startup_message &
            (beam_structure = local%beam_structure)
    end if
    
    if (intg%process_has_me) then
       call intg%process%get_pdg_in (pdg_prc)
    else
       allocate (pdg_prc (n_in, n_components))
       pdg_prc = 0
    end if
    call dispatch_sf_config (sf_config, sf_prop, local, pdg_prc)
    sf_trace = &
         var_list%get_lval (var_str ("?sf_trace"))
    sf_trace_file = &
         var_list%get_sval (var_str ("$sf_trace_file"))
    if (sf_trace) then
       call intg%process%init_sf_chain (sf_config, sf_trace_file)
    else
       call intg%process%init_sf_chain (sf_config)
    end if

    if (intg%process_has_me) then
       call intg%process%configure_phs &
            (intg%rebuild_phs, intg%ignore_phs_mismatch, verbose=verbose, &
             combined_integration=intg%combined_integration)
       if (size (sf_config) > 0) then
          call intg%process%collect_channels (phs_channel_collection)
       else if (intg%process%contains_trivial_component ()) then
          call msg_fatal ("Integrate: 2 -> 1 process can't be handled &
               &with fixed-energy beams")
       end if
       call dispatch_sf_channels &
            (sf_channel, sf_string, sf_prop, phs_channel_collection, local)
       if (allocated (sf_channel)) then
          if (size (sf_channel) > 0) then
             call intg%process%set_sf_channel (sf_channel)
          end if
       end if
       call phs_channel_collection%final ()
       if (verb)  call intg%process%sf_startup_message (sf_string)    
    end if
    
    call intg%setup_process_mci ()
    call intg%process%setup_terms ()

    if (associated (local%pn%cuts_lexpr)) then
       if (verb)  call msg_message ("Applying user-defined cuts.")
       call expr_factory%init (local%pn%cuts_lexpr)
       call intg%process%set_cuts (expr_factory)
    else
       if (verb)  call msg_warning ("No cuts have been defined.")
    end if    
    if (associated (local%pn%scale_expr)) then
       if (verb) call msg_message ("Using user-defined general scale.")
       call expr_factory%init (local%pn%scale_expr)
       call intg%process%set_scale (expr_factory)
    end if
    if (associated (local%pn%fac_scale_expr)) then
       if (verb) call msg_message ("Using user-defined factorization scale.")
       call expr_factory%init (local%pn%fac_scale_expr)
       call intg%process%set_fac_scale (expr_factory)
    end if
    if (associated (local%pn%ren_scale_expr)) then
       if (verb) call msg_message ("Using user-defined renormalization scale.")
       call expr_factory%init (local%pn%ren_scale_expr)
       call intg%process%set_ren_scale (expr_factory)
    end if
    if (associated (local%pn%weight_expr)) then
       if (verb) call msg_message ("Using user-defined reweighting factor.")
       call expr_factory%init (local%pn%weight_expr)
       call intg%process%set_weight (expr_factory)
    end if

    call intg%process%compute_md5sum ()
  end subroutine integration_setup_process

\eatline
\nwendcode{}\nwbegindocs{2041}\nwdocspar
\subsection{Integration}
Integrate: do the final integration.  Here, we do a multi-iteration
integration.  Again, we skip iterations that are already on file.
Record the results in the global variable list.
\nwenddocs{}\nwbegincode{2042}\moddef{Integrations: integration: TBP}\plusendmoddef
  procedure :: evaluate => integration_evaluate
\nwendcode{}\nwbegincode{2043}\moddef{Integrations: procedures}\plusendmoddef
  subroutine integration_evaluate &
       (intg, process_instance, i_mci, pass, it_list, pacify)
    class(integration_t), intent(inout) :: intg
    type(process_instance_t), intent(inout), target :: process_instance
    integer, intent(in) :: i_mci
    integer, intent(in) :: pass
    type(iterations_list_t), intent(in) :: it_list
    logical, intent(in), optional :: pacify
    integer :: n_calls, n_it
    logical :: adapt_grids, adapt_weights, final
        
    n_it = it_list%get_n_it (pass)
    n_calls = it_list%get_n_calls (pass)
    adapt_grids = it_list%adapt_grids (pass)
    adapt_weights = it_list%adapt_weights (pass)
    final = pass == it_list%get_n_pass ()
    
    call intg%process%integrate (process_instance, &
         i_mci, n_it, n_calls, adapt_grids, adapt_weights, &
         final, pacify)

  end subroutine integration_evaluate

\eatline
\nwendcode{}\nwbegindocs{2044}In case the user has not provided a list of iterations, make a
reasonable default.  This can depend on the process.  The usual
approach is to define two distinct passes, one for adaptation and one
for integration.
\nwenddocs{}\nwbegincode{2045}\moddef{Integrations: integration: TBP}\plusendmoddef
  procedure :: make_iterations_list => integration_make_iterations_list
\nwendcode{}\nwbegincode{2046}\moddef{Integrations: procedures}\plusendmoddef
  subroutine integration_make_iterations_list (intg, it_list)
    class(integration_t), intent(in) :: intg
    type(iterations_list_t), intent(out) :: it_list
    integer :: pass, n_pass
    integer, dimension(:), allocatable :: n_it, n_calls
    logical, dimension(:), allocatable :: adapt_grids, adapt_weights
    n_pass = intg%process%get_n_pass_default ()
    allocate (n_it (n_pass), n_calls (n_pass))
    allocate (adapt_grids (n_pass), adapt_weights (n_pass))
    do pass = 1, n_pass
       n_it(pass)          = intg%process%get_n_it_default (pass)
       n_calls(pass)       = intg%process%get_n_calls_default (pass)
       adapt_grids(pass)   = intg%process%adapt_grids_default (pass)
       adapt_weights(pass) = intg%process%adapt_weights_default (pass)
    end do
    call it_list%init (n_it, n_calls, &
         adapt_grids = adapt_grids, adapt_weights = adapt_weights)
  end subroutine integration_make_iterations_list
  
\eatline
\nwendcode{}\nwbegindocs{2047}\nwdocspar
\subsection{API for integration objects}
This initializer does everything except assigning cuts/scale/weight
expressions.  
\nwenddocs{}\nwbegincode{2048}\moddef{Integrations: integration: TBP}\plusendmoddef
  procedure :: init => integration_init
\nwendcode{}\nwbegincode{2049}\moddef{Integrations: procedures}\plusendmoddef
  subroutine integration_init (intg, process_id, local, global, local_stack)
    class(integration_t), intent(out) :: intg
    type(string_t), intent(in) :: process_id
    type(rt_data_t), intent(inout), target :: local
    type(rt_data_t), intent(inout), optional, target :: global
    logical, intent(in), optional :: local_stack
    logical :: use_local
    use_local = .false.;  if (present (local_stack))  use_local = local_stack
    if (present (global)) then
       call intg%create_process (process_id, global)
    else if (use_local) then
       call intg%create_process (process_id, local)
    else
       call intg%create_process (process_id)
    end if
    call intg%init_process (local)
    call intg%setup_process (local)
  end subroutine integration_init

\eatline
\nwendcode{}\nwbegindocs{2050}Do the integration for a single process, both warmup and final evaluation.
The {\tt{}eff{\char95}reset} flag is to suppress numerical noise in the graphical output
of the integration history.
\nwenddocs{}\nwbegincode{2051}\moddef{Integrations: integration: TBP}\plusendmoddef
  procedure :: integrate => integration_integrate
\nwendcode{}\nwbegincode{2052}\moddef{Integrations: procedures}\plusendmoddef
  subroutine integration_integrate (intg, local, eff_reset)
    class(integration_t), intent(inout) :: intg
    type(rt_data_t), intent(in), target :: local
    logical, intent(in), optional :: eff_reset
    type(string_t) :: log_filename
    type(var_list_t), pointer :: var_list
    type(process_instance_t), allocatable, target :: process_instance
    type(iterations_list_t) :: it_list
    logical :: pacify
    integer :: pass, i_mci, n_mci, n_pass
    type(string_t) :: nlo_type
    logical :: display_summed
    logical :: use_internal_color_correlations

    var_list => intg%process%get_var_list_ptr ()
    use_internal_color_correlations = &
         .not. var_list%get_lval (var_str ('?use_gosam_correlations'))

    allocate (process_instance)
    call process_instance%init (intg%process, use_internal_color_correlations, &
                                combined_integration = intg%combined_integration)
    if (process_instance%has_nlo_component ()) then
       call process_instance%create_blha_interface ()
       call process_instance%load_blha_libraries (local%os_data)
    end if

    call openmp_set_num_threads_verbose &
         (var_list%get_ival (var_str ("openmp_num_threads")), &
          var_list%get_lval (var_str ("?openmp_logging")))
    pacify = var_list%get_lval (var_str ("?pacify"))

    display_summed = .true.
    n_mci = intg%process%get_n_mci ()
    if (n_mci == 1) then
       write (msg_buffer, "(A,A,A)") &
            "Starting integration for process '", &
            char (intg%process%get_id ()), "'"
       call msg_message ()
    end if
    call intg%setup_component_cores (process_instance)

    do i_mci = 1, n_mci
       if (intg%process%is_active_nlo_component (i_mci)) then
         if (n_mci > 1) then
            write (msg_buffer, "(A,A,A,I0)") &
                 "Starting integration for process '", &
                 char (intg%process%get_id ()), "' part ", i_mci
            call msg_message ()
         end if
         n_pass = local%it_list%get_n_pass ()
         if (n_pass == 0) then
            call msg_message ("Integrate: iterations not specified, &
                 &using default")
            call intg%make_iterations_list (it_list)
            n_pass = it_list%get_n_pass ()
         else
            it_list = local%it_list
         end if
         call msg_message ("Integrate: " // char (it_list%to_string ()))
         do pass = 1, n_pass
            call intg%evaluate (process_instance, i_mci, pass, it_list, pacify)
            if (signal_is_pending ())  return
         end do
         call intg%process%final_integration (i_mci)       
         if (intg%vis_history) then
            call intg%process%display_integration_history &
                 (i_mci, intg%history_filename, local%os_data, eff_reset)
         end if       
         if (local%logfile == intg%log_filename) then
            if (intg%run_id /= "") then
               log_filename = intg%process_id // "." // intg%run_id // &
                    ".var.log"
            else
               log_filename = intg%process_id // ".var.log"
            end if
            call msg_message ("Name clash for global logfile and process log: ", &
                 arr =[var_str ("| Renaming log file from ") // local%logfile, &
                       var_str ("|   to ") // log_filename // var_str (" .")])
         else
            log_filename = intg%log_filename
         end if
         call intg%process%write_logfile (i_mci, log_filename)    
       else
         nlo_type = intg%process%get_component_nlo_type (i_mci)
         if (nlo_type /= "Subtraction") display_summed = .false.
       end if          
    end do

    if (n_mci > 1 .and. display_summed) then
       call msg_message ("Integrate: sum of all components")
       call intg%process%display_summed_results ()
    end if

    call process_instance%final ()
    deallocate (process_instance)

  end subroutine integration_integrate
  
\eatline
\nwendcode{}\nwbegindocs{2053}\nwdocspar
\nwenddocs{}\nwbegincode{2054}\moddef{Integrations: integration: TBP}\plusendmoddef
  procedure :: setup_component_cores => integration_setup_component_cores
\nwendcode{}\nwbegincode{2055}\moddef{Integrations: procedures}\plusendmoddef
  subroutine integration_setup_component_cores (intg, instance)
    class(integration_t), intent(inout) :: intg
    type(process_instance_t), intent(inout) :: instance
    integer :: i_born, i_sub
    integer :: n_components, i_component
    type (string_t) :: nlo_type
    class(prc_core_t), allocatable :: core_born
    associate (process => intg%process)
       n_components = process%get_n_components ()
       do i_component = 1, n_components
          nlo_type = process%get_component_nlo_type (i_component)
          i_born = process%get_component_associated_born (i_component)
          i_sub = i_born + 3
          select case (char (nlo_type))
          case ('Real', 'Virtual')
             call process%extract_component_core (i_sub, core_born)
             call process%init_sub_born (i_component, core_born)
             call process%restore_component_core (i_sub, core_born)
             call instance%init_born_amps (i_component, i_born)
         end select
       end do
    end associate
  end subroutine integration_setup_component_cores

\eatline
\nwendcode{}\nwbegindocs{2056}\nwdocspar
\nwenddocs{}\nwbegincode{2057}\moddef{Integrations: integration: TBP}\plusendmoddef
  procedure :: setup_process_mci => integration_setup_process_mci
\nwendcode{}\nwbegincode{2058}\moddef{Integrations: procedures}\plusendmoddef
  subroutine integration_setup_process_mci (intg)
    class(integration_t), intent(inout) :: intg
    call intg%process%setup_mci (intg%combined_integration)
  end subroutine integration_setup_process_mci

\eatline
\nwendcode{}\nwbegindocs{2059}Do a dummy integration for a process which could not be initialized (e.g.,
has no matrix element).  The result is zero.
\nwenddocs{}\nwbegincode{2060}\moddef{Integrations: integration: TBP}\plusendmoddef
  procedure :: integrate_dummy => integration_integrate_dummy 
\nwendcode{}\nwbegincode{2061}\moddef{Integrations: procedures}\plusendmoddef
  subroutine integration_integrate_dummy (intg)
    class(integration_t), intent(inout) :: intg
    call intg%process%integrate_dummy ()
  end subroutine integration_integrate_dummy
     
\eatline
\nwendcode{}\nwbegindocs{2062}Just sample the matrix element under realistic conditions (but no
cuts); throw away the results.
\nwenddocs{}\nwbegincode{2063}\moddef{Integrations: integration: TBP}\plusendmoddef
  procedure :: sampler_test => integration_sampler_test
\nwendcode{}\nwbegincode{2064}\moddef{Integrations: procedures}\plusendmoddef
  subroutine integration_sampler_test (intg)
    class(integration_t), intent(inout) :: intg
    type(process_instance_t), allocatable, target :: process_instance
    integer :: n_mci, i_mci
    type(timer_t) :: timer_mci, timer_tot
    real(default) :: t_mci, t_tot
    allocate (process_instance)
    call process_instance%init (intg%process)
    n_mci = intg%process%get_n_mci ()
    if (n_mci == 1) then
       write (msg_buffer, "(A,A,A)") &
            "Test: probing process '", &
            char (intg%process%get_id ()), "'"
       call msg_message ()
    end if
    call timer_tot%start ()
    do i_mci = 1, n_mci
       if (n_mci > 1) then
          write (msg_buffer, "(A,A,A,I0)") &
               "Test: probing process '", &
               char (intg%process%get_id ()), "' part ", i_mci
          call msg_message ()
       end if
       call timer_mci%start ()
       call intg%process%sampler_test &
            (process_instance, i_mci, intg%n_calls_test)
       call timer_mci%stop ()
       t_mci = timer_mci
       write (msg_buffer, "(A,ES12.5)")  "Test: " &
            // "time in seconds (wallclock): ", t_mci
       call msg_message ()
    end do
    call timer_tot%stop ()
    t_tot = timer_tot
    if (n_mci > 1) then
       write (msg_buffer, "(A,ES12.5)")  "Test: " &
            // "total time      (wallclock): ", t_tot
       call msg_message ()
    end if
    call process_instance%final ()
  end subroutine integration_sampler_test

\eatline
\nwendcode{}\nwbegindocs{2065}Return the process pointer (needed by simulate):
\nwenddocs{}\nwbegincode{2066}\moddef{Integrations: integration: TBP}\plusendmoddef
  procedure :: get_process_ptr => integration_get_process_ptr
\nwendcode{}\nwbegincode{2067}\moddef{Integrations: procedures}\plusendmoddef
  function integration_get_process_ptr (intg) result (ptr)
    class(integration_t), intent(in) :: intg
    type(process_t), pointer :: ptr
    ptr => intg%process
  end function integration_get_process_ptr

\eatline
\nwendcode{}\nwbegindocs{2068}Simply integrate, do a dummy integration if necessary.  The {\tt{}integration}
object exists only internally.

If the {\tt{}global} environment is provided, the process object is appended to
the global stack.  Otherwise, if {\tt{}local{\char95}stack} is set, we append to the
local process stack.  If this is unset, the {\tt{}process} object is not recorded
permanently.

The {\tt{}init{\char95}only} flag can be used to skip the actual integration part.  We
will end up with a process object that is completely initialized, including
phase space configuration. 

The {\tt{}eff{\char95}reset} flag is to suppress numerical noise in the visualization
of the integration history.
\nwenddocs{}\nwbegincode{2069}\moddef{Integrations: public}\plusendmoddef
  public :: integrate_process
\nwendcode{}\nwbegincode{2070}\moddef{Integrations: procedures}\plusendmoddef
  subroutine integrate_process (process_id, local, global, local_stack, init_only, eff_reset)
    type(string_t), intent(in) :: process_id
    type(rt_data_t), intent(inout), target :: local
    type(rt_data_t), intent(inout), optional, target :: global
    logical, intent(in), optional :: local_stack, init_only, eff_reset
    type(string_t) :: prclib_name
    type(integration_t) :: intg
    character(32) :: buffer

    if (.not. associated (local%prclib)) then
       call msg_fatal ("Integrate: current process library is undefined")
       return
    end if

    if (.not. local%prclib%is_active ()) then
       call msg_message ("Integrate: current process library needs compilation")
       prclib_name = local%prclib%get_name ()
       call compile_library (prclib_name, local)
       if (signal_is_pending ())  return
       call msg_message ("Integrate: compilation done")
    end if

    call intg%init (process_id, local, global, local_stack)
    if (signal_is_pending ())  return

    if (present (init_only)) then
       if (init_only) return
    end if

    if (intg%n_calls_test > 0) then
       write (buffer, "(I0)")  intg%n_calls_test
       call msg_message ("Integrate: test (" // trim (buffer) // " calls) ...")
       call intg%sampler_test ()
       call msg_message ("Integrate: ... test complete.")
       if (signal_is_pending ())  return
    end if

    if (intg%phs_only) then
       call msg_message ("Integrate: phase space only, skipping integration")
    else
       if (intg%process_has_me) then
          call intg%integrate (local, eff_reset)
       else
          call intg%integrate_dummy ()
       end if
    end if

  end subroutine integrate_process

\eatline
\nwendcode{}\nwbegindocs{2071}\nwdocspar
\subsection{Test}
This is the master for calling self-test procedures.
\nwenddocs{}\nwbegincode{2072}\moddef{Integrations: public}\plusendmoddef
  public :: integrations_test
\nwendcode{}\nwbegincode{2073}\moddef{Integrations: tests}\endmoddef
  subroutine integrations_test (u, results)
    integer, intent(in) :: u
    type(test_results_t), intent(inout) :: results
  \LA{}Integrations: execute tests\RA{}
  end subroutine integrations_test

\eatline
\nwendcode{}\nwbegindocs{2074}\nwdocspar
\nwenddocs{}\nwbegincode{2075}\moddef{Integrations: public}\plusendmoddef
  public :: integrations_history_test    
\nwendcode{}\nwbegincode{2076}\moddef{Integrations: tests}\plusendmoddef
  subroutine integrations_history_test (u, results)
    integer, intent(in) :: u
    type(test_results_t), intent(inout) :: results
  \LA{}Integrations: execute history tests\RA{}
  end subroutine integrations_history_test  

\eatline
\nwendcode{}\nwbegindocs{2077}\nwdocspar
\subsubsection{Integration of test process}
Compile and integrate an intrinsic test matrix element ({\tt{}prc{\char95}test}
type).  The phase-space implementation is {\tt{}phs{\char95}single}
(single-particle phase space), the integrator is {\tt{}mci{\char95}midpoint}.

The cross section for the $2\to 2$ process $ss\to ss$ with its
constant matrix element is given by
\begin{equation}
  \sigma = c\times f\times \Phi_2 \times |M|^2.
\end{equation}
$c$ is the conversion constant
\begin{equation}
  c = 0.3894\times 10^{12}\;\mathrm{fb}\,\mathrm{GeV}^2.
\end{equation}
$f$ is the flux of the incoming particles with mass
$m=125\,\mathrm{GeV}$ and energy $\sqrt{s}=1000\,\mathrm{GeV}$
\begin{equation}
  f = \frac{(2\pi)^4}{2\lambda^{1/2}(s,m^2,m^2)}
    = \frac{(2\pi)^4}{2\sqrt{s}\,\sqrt{s - 4m^2}}
    = 8.048\times 10^{-4}\;\mathrm{GeV}^{-2}
\end{equation}
$\Phi_2$ is the volume of the two-particle phase space
\begin{equation}
  \Phi_2 = \frac{1}{4(2\pi)^5} = 2.5529\times 10^{-5}.
\end{equation}
The squared matrix element $|M|^2$ is unity.
Combining everything, we obtain
\begin{equation}
  \sigma = 8000\;\mathrm{fb}
\end{equation}
This number should appear as the final result.

Note: In this and the following test, we reset the Fortran compiler and flag
variables immediately before they are printed, so the test is portable.
\nwenddocs{}\nwbegincode{2078}\moddef{Integrations: execute tests}\endmoddef
  call test (integrations_1, "integrations_1", &
       "intrinsic test process", &
       u, results)
\nwendcode{}\nwbegincode{2079}\moddef{Integrations: tests}\plusendmoddef
  subroutine integrations_1 (u)
    integer, intent(in) :: u
    type(string_t) :: libname, procname
    type(rt_data_t), target :: global
    
    write (u, "(A)")  "* Test output: integrations_1"
    write (u, "(A)")  "*   Purpose: integrate test process"
    write (u, "(A)")

    call syntax_model_file_init ()

    call global%global_init ()

    libname = "integration_1"
    procname = "prc_config_a"
    
    call prepare_test_library (global, libname, 1)
    call compile_library (libname, global)

    call global%set_string (var_str ("$run_id"), &
         var_str ("integrations1"), is_known = .true.)
    call global%set_string (var_str ("$method"), &
         var_str ("unit_test"), is_known = .true.)
    call global%set_string (var_str ("$phs_method"), &
         var_str ("single"), is_known = .true.)
    call global%set_string (var_str ("$integration_method"),&
         var_str ("midpoint"), is_known = .true.)
    call global%set_log (var_str ("?vis_history"),&
         .false., is_known = .true.)    
    call global%set_log (var_str ("?integration_timer"),&
         .false., is_known = .true.) 
    call global%set_int (var_str ("seed"), &
         0, is_known=.true.)    
    
    call global%set_real (var_str ("sqrts"),&
         1000._default, is_known = .true.)

    call global%it_list%init ([1], [1000])

    call reset_interaction_counter ()
    call integrate_process (procname, global, local_stack=.true.)

    call global%write (u, vars = [ &
         var_str ("$method"), &
         var_str ("sqrts"), &
         var_str ("$integration_method"), &
         var_str ("$phs_method"), &
         var_str ("$run_id")])
    
    call global%final ()
    call syntax_model_file_final ()
    
    write (u, "(A)")
    write (u, "(A)")  "* Test output end: integrations_1"
    
  end subroutine integrations_1
  
\eatline
\nwendcode{}\nwbegindocs{2080}\nwdocspar
\subsubsection{Integration with cuts}
Compile and integrate an intrinsic test matrix element ({\tt{}prc{\char95}test}
type) with cuts set.
\nwenddocs{}\nwbegincode{2081}\moddef{Integrations: execute tests}\plusendmoddef
  call test (integrations_2, "integrations_2", &
       "intrinsic test process with cut", &
       u, results)
\nwendcode{}\nwbegincode{2082}\moddef{Integrations: tests}\plusendmoddef
  subroutine integrations_2 (u)
    integer, intent(in) :: u
    type(string_t) :: libname, procname
    type(rt_data_t), target :: global

    type(string_t) :: cut_expr_text
    type(ifile_t) :: ifile
    type(stream_t) :: stream
    type(parse_tree_t) :: parse_tree
    
    type(string_t), dimension(0) :: empty_string_array

    write (u, "(A)")  "* Test output: integrations_2"
    write (u, "(A)")  "*   Purpose: integrate test process with cut"
    write (u, "(A)")

    call syntax_model_file_init ()

    call global%global_init ()

    write (u, "(A)")  "* Prepare a cut expression"
    write (u, "(A)")

    call syntax_pexpr_init ()
    cut_expr_text = "all Pt > 100 [s]"
    call ifile_append (ifile, cut_expr_text)
    call stream_init (stream, ifile)
    call parse_tree_init_lexpr (parse_tree, stream, .true.)
    global%pn%cuts_lexpr => parse_tree_get_root_ptr (parse_tree)
    
    write (u, "(A)")  "* Build and initialize a test process"
    write (u, "(A)")

    libname = "integration_3"
    procname = "prc_config_a"
    
    call prepare_test_library (global, libname, 1)
    call compile_library (libname, global)

    call global%set_string (var_str ("$run_id"), &
         var_str ("integrations1"), is_known = .true.)
    call global%set_string (var_str ("$method"), &
         var_str ("unit_test"), is_known = .true.)
    call global%set_string (var_str ("$phs_method"), &
         var_str ("single"), is_known = .true.)
    call global%set_string (var_str ("$integration_method"),&
         var_str ("midpoint"), is_known = .true.)
    call global%set_log (var_str ("?vis_history"),&
         .false., is_known = .true.)    
    call global%set_log (var_str ("?integration_timer"),&
         .false., is_known = .true.)  
    call global%set_int (var_str ("seed"), &
         0, is_known=.true.)    

    call global%set_real (var_str ("sqrts"),&
         1000._default, is_known = .true.)

    call global%it_list%init ([1], [1000])
    
    call reset_interaction_counter ()
    call integrate_process (procname, global, local_stack=.true.)
    
    call global%write (u, vars = empty_string_array)
    
    call global%final ()
    call syntax_model_file_final ()
    
    write (u, "(A)")
    write (u, "(A)")  "* Test output end: integrations_2"
    
  end subroutine integrations_2
  
\eatline
\nwendcode{}\nwbegindocs{2083}\nwdocspar
\subsubsection{Standard phase space}
Compile and integrate an intrinsic test matrix element ({\tt{}prc{\char95}test}
type) using the default ({\tt{}phs{\char95}wood}) phase-space implementation.  We
use an explicit phase-space configuration file with a single channel
and integrate by {\tt{}mci{\char95}midpoint}.
\nwenddocs{}\nwbegincode{2084}\moddef{Integrations: execute tests}\plusendmoddef
  call test (integrations_3, "integrations_3", &
       "standard phase space", &
       u, results)
\nwendcode{}\nwbegincode{2085}\moddef{Integrations: tests}\plusendmoddef
  subroutine integrations_3 (u)
    integer, intent(in) :: u
    type(string_t) :: libname, procname
    type(rt_data_t), target :: global
    integer :: u_phs
    
    write (u, "(A)")  "* Test output: integrations_3"
    write (u, "(A)")  "*   Purpose: integrate test process"
    write (u, "(A)")

    write (u, "(A)")  "* Initialize process and parameters"
    write (u, "(A)")

    call syntax_model_file_init ()
    call syntax_phs_forest_init ()

    call global%global_init ()

    libname = "integration_3"
    procname = "prc_config_a"
    
    call prepare_test_library (global, libname, 1)
    call compile_library (libname, global)

    call global%set_string (var_str ("$run_id"), &
         var_str ("integrations1"), is_known = .true.)
    call global%set_string (var_str ("$method"), &
         var_str ("unit_test"), is_known = .true.)
    call global%set_string (var_str ("$phs_method"), &
         var_str ("default"), is_known = .true.)
    call global%set_string (var_str ("$integration_method"),&
         var_str ("midpoint"), is_known = .true.)
    call global%set_log (var_str ("?vis_history"),&
         .false., is_known = .true.)    
    call global%set_log (var_str ("?integration_timer"),&
         .false., is_known = .true.)    
    call global%set_log (var_str ("?phs_s_mapping"),&
         .false., is_known = .true.)   
    call global%set_int (var_str ("seed"), &
         0, is_known=.true.)    
    
    call global%set_real (var_str ("sqrts"),&
         1000._default, is_known = .true.)

    write (u, "(A)")  "* Create a scratch phase-space file"
    write (u, "(A)")

    u_phs = free_unit ()
    open (u_phs, file = "integrations_3.phs", &
         status = "replace", action = "write")
    call write_test_phs_file (u_phs, var_str ("prc_config_a_i1"))
    close (u_phs)

    call global%set_string (var_str ("$phs_file"),&
         var_str ("integrations_3.phs"), is_known = .true.)

    call global%it_list%init ([1], [1000])

    write (u, "(A)")  "* Integrate"
    write (u, "(A)")

    call reset_interaction_counter ()
    call integrate_process (procname, global, local_stack=.true.)
    
    call global%write (u, vars = [ &
         var_str ("$phs_method"), &
         var_str ("$phs_file")])
    
    write (u, "(A)")
    write (u, "(A)")  "* Cleanup"

    call global%final ()
    call syntax_phs_forest_final ()
    call syntax_model_file_final ()
    
    write (u, "(A)")
    write (u, "(A)")  "* Test output end: integrations_3"
    
  end subroutine integrations_3
  
\eatline
\nwendcode{}\nwbegindocs{2086}\nwdocspar
\subsubsection{VAMP integration}
Compile and integrate an intrinsic test matrix element ({\tt{}prc{\char95}test}
type) using the single-channel ({\tt{}phs{\char95}single}) phase-space
implementation.  The integration method is {\tt{}vamp}.
\nwenddocs{}\nwbegincode{2087}\moddef{Integrations: execute tests}\plusendmoddef
  call test (integrations_4, "integrations_4", &
       "VAMP integration (one iteration)", &
       u, results)
\nwendcode{}\nwbegincode{2088}\moddef{Integrations: tests}\plusendmoddef
  subroutine integrations_4 (u)
    integer, intent(in) :: u
    type(string_t) :: libname, procname
    type(rt_data_t), target :: global
    
    write (u, "(A)")  "* Test output: integrations_4"
    write (u, "(A)")  "*   Purpose: integrate test process using VAMP"
    write (u, "(A)")

    write (u, "(A)")  "* Initialize process and parameters"
    write (u, "(A)")

    call syntax_model_file_init ()

    call global%global_init ()

    libname = "integrations_4_lib"
    procname = "integrations_4"
    
    call prepare_test_library (global, libname, 1, [procname])
    call compile_library (libname, global)

    call global%append_log (&
         var_str ("?rebuild_grids"), .true., intrinsic = .true.)

    call global%set_string (var_str ("$run_id"), &
         var_str ("r1"), is_known = .true.)
    call global%set_string (var_str ("$method"), &
         var_str ("unit_test"), is_known = .true.)
    call global%set_string (var_str ("$phs_method"), &
         var_str ("single"), is_known = .true.)
    call global%set_string (var_str ("$integration_method"),&
         var_str ("vamp"), is_known = .true.)
    call global%set_log (var_str ("?use_vamp_equivalences"),&
         .false., is_known = .true.)
    call global%set_log (var_str ("?vis_history"),&
         .false., is_known = .true.)    
    call global%set_log (var_str ("?integration_timer"),&
         .false., is_known = .true.)    
    call global%set_int (var_str ("seed"), &
         0, is_known=.true.)    
    
    call global%set_real (var_str ("sqrts"),&
         1000._default, is_known = .true.)

    call global%it_list%init ([1], [1000])

    write (u, "(A)")  "* Integrate"
    write (u, "(A)")

    call reset_interaction_counter ()
    call integrate_process (procname, global, local_stack=.true.)
    
    call global%pacify (efficiency_reset = .true., error_reset = .true.)
    call global%write (u, vars = [var_str ("$integration_method")], &
            pacify = .true.)
    
    write (u, "(A)")
    write (u, "(A)")  "* Cleanup"

    call global%final ()
    call syntax_model_file_final ()
    
    write (u, "(A)")
    write (u, "(A)")  "* Test output end: integrations_4"
    
  end subroutine integrations_4
  
\eatline
\nwendcode{}\nwbegindocs{2089}\nwdocspar
\subsubsection{Multiple iterations integration}
Compile and integrate an intrinsic test matrix element ({\tt{}prc{\char95}test}
type) using the single-channel ({\tt{}phs{\char95}single}) phase-space
implementation.  The integration method is {\tt{}vamp}.  We launch three
iterations.
\nwenddocs{}\nwbegincode{2090}\moddef{Integrations: execute tests}\plusendmoddef
  call test (integrations_5, "integrations_5", &
       "VAMP integration (three iterations)", &
       u, results)
\nwendcode{}\nwbegincode{2091}\moddef{Integrations: tests}\plusendmoddef
  subroutine integrations_5 (u)
    integer, intent(in) :: u
    type(string_t) :: libname, procname
    type(rt_data_t), target :: global
    
    write (u, "(A)")  "* Test output: integrations_5"
    write (u, "(A)")  "*   Purpose: integrate test process using VAMP"
    write (u, "(A)")

    write (u, "(A)")  "* Initialize process and parameters"
    write (u, "(A)")

    call syntax_model_file_init ()

    call global%global_init ()

    libname = "integrations_5_lib"
    procname = "integrations_5"
    
    call prepare_test_library (global, libname, 1, [procname])
    call compile_library (libname, global)

    call global%append_log (&
         var_str ("?rebuild_grids"), .true., intrinsic = .true.)

    call global%set_string (var_str ("$run_id"), &
         var_str ("r1"), is_known = .true.)
    call global%set_string (var_str ("$method"), &
         var_str ("unit_test"), is_known = .true.)
    call global%set_string (var_str ("$phs_method"), &
         var_str ("single"), is_known = .true.)
    call global%set_string (var_str ("$integration_method"),&
         var_str ("vamp"), is_known = .true.)
    call global%set_log (var_str ("?use_vamp_equivalences"),&
         .false., is_known = .true.)
    call global%set_log (var_str ("?vis_history"),&
         .false., is_known = .true.)    
    call global%set_log (var_str ("?integration_timer"),&
         .false., is_known = .true.)    
    call global%set_int (var_str ("seed"), &
         0, is_known=.true.)
    
    call global%set_real (var_str ("sqrts"),&
         1000._default, is_known = .true.)

    call global%it_list%init ([3], [1000])

    write (u, "(A)")  "* Integrate"
    write (u, "(A)")

    call reset_interaction_counter ()
    call integrate_process (procname, global, local_stack=.true.)
    
    call global%pacify (efficiency_reset = .true., error_reset = .true.)
    call global%write (u, vars = [var_str ("$integration_method")], &
            pacify = .true.)
    
    write (u, "(A)")
    write (u, "(A)")  "* Cleanup"

    call global%final ()
    call syntax_model_file_final ()
    
    write (u, "(A)")
    write (u, "(A)")  "* Test output end: integrations_5"
    
  end subroutine integrations_5
  
\eatline
\nwendcode{}\nwbegindocs{2092}\nwdocspar
\subsubsection{Multiple passes integration}
Compile and integrate an intrinsic test matrix element ({\tt{}prc{\char95}test}
type) using the single-channel ({\tt{}phs{\char95}single}) phase-space
implementation.  The integration method is {\tt{}vamp}.  We launch three
passes with three iterations each.
\nwenddocs{}\nwbegincode{2093}\moddef{Integrations: execute tests}\plusendmoddef
  call test (integrations_6, "integrations_6", &
       "VAMP integration (three passes)", &
       u, results)
\nwendcode{}\nwbegincode{2094}\moddef{Integrations: tests}\plusendmoddef
  subroutine integrations_6 (u)
    integer, intent(in) :: u
    type(string_t) :: libname, procname
    type(rt_data_t), target :: global
    type(string_t), dimension(0) :: no_vars
    
    write (u, "(A)")  "* Test output: integrations_6"
    write (u, "(A)")  "*   Purpose: integrate test process using VAMP"
    write (u, "(A)")

    write (u, "(A)")  "* Initialize process and parameters"
    write (u, "(A)")

    call syntax_model_file_init ()

    call global%global_init ()

    libname = "integrations_6_lib"
    procname = "integrations_6"
    
    call prepare_test_library (global, libname, 1, [procname])
    call compile_library (libname, global)

    call global%append_log (&
         var_str ("?rebuild_grids"), .true., intrinsic = .true.)

    call global%set_string (var_str ("$run_id"), &
         var_str ("r1"), is_known = .true.)
    call global%set_string (var_str ("$method"), &
         var_str ("unit_test"), is_known = .true.)
    call global%set_string (var_str ("$phs_method"), &
         var_str ("single"), is_known = .true.)
    call global%set_string (var_str ("$integration_method"),&
         var_str ("vamp"), is_known = .true.)
    call global%set_log (var_str ("?use_vamp_equivalences"),&
         .false., is_known = .true.)
    call global%set_log (var_str ("?vis_history"),&
         .false., is_known = .true.)    
    call global%set_log (var_str ("?integration_timer"),&
         .false., is_known = .true.)    
    call global%set_int (var_str ("seed"), &
         0, is_known=.true.)    
    
    call global%set_real (var_str ("sqrts"),&
         1000._default, is_known = .true.)

    call global%it_list%init ([3, 3, 3], [1000, 1000, 1000], &
         adapt = [.true., .true., .false.], &
         adapt_code = [var_str ("wg"), var_str ("g"), var_str ("")])

    write (u, "(A)")  "* Integrate"
    write (u, "(A)")

    call reset_interaction_counter ()
    call integrate_process (procname, global, local_stack=.true.)
    
    call global%pacify (efficiency_reset = .true., error_reset = .true.)
    call global%write (u, vars = no_vars, pacify = .true.)

    write (u, "(A)")
    write (u, "(A)")  "* Cleanup"

    call global%final ()
    call syntax_model_file_final ()
    
    write (u, "(A)")
    write (u, "(A)")  "* Test output end: integrations_6"
    
  end subroutine integrations_6
  
\eatline
\nwendcode{}\nwbegindocs{2095}\nwdocspar
\subsubsection{VAMP and default phase space}
Compile and integrate an intrinsic test matrix element ({\tt{}prc{\char95}test}
type) using the default ({\tt{}phs{\char95}wood}) phase-space
implementation.  The integration method is {\tt{}vamp}.  We launch three
passes with three iterations each.  We enable channel equivalences and
groves.
\nwenddocs{}\nwbegincode{2096}\moddef{Integrations: execute tests}\plusendmoddef
  call test (integrations_7, "integrations_7", &
       "VAMP integration with wood phase space", &
       u, results)
\nwendcode{}\nwbegincode{2097}\moddef{Integrations: tests}\plusendmoddef
  subroutine integrations_7 (u)
    integer, intent(in) :: u
    type(string_t) :: libname, procname
    type(rt_data_t), target :: global
    type(string_t), dimension(0) :: no_vars
    integer :: iostat, u_phs
    character(95) :: buffer
    type(string_t) :: phs_file
    logical :: exist
    
    write (u, "(A)")  "* Test output: integrations_7"
    write (u, "(A)")  "*   Purpose: integrate test process using VAMP"
    write (u, "(A)")

    write (u, "(A)")  "* Initialize process and parameters"
    write (u, "(A)")

    call syntax_model_file_init ()
    call syntax_phs_forest_init ()

    call global%global_init ()

    libname = "integrations_7_lib"
    procname = "integrations_7"
    
    call prepare_test_library (global, libname, 1, [procname])
    call compile_library (libname, global)

    call global%append_log (&
         var_str ("?rebuild_phase_space"), .true., intrinsic = .true.)
    call global%append_log (&
         var_str ("?rebuild_grids"), .true., intrinsic = .true.)

    call global%set_string (var_str ("$run_id"), &
         var_str ("r1"), is_known = .true.)
    call global%set_string (var_str ("$method"), &
         var_str ("unit_test"), is_known = .true.)
    call global%set_string (var_str ("$phs_method"), &
         var_str ("wood"), is_known = .true.)
    call global%set_string (var_str ("$integration_method"),&
         var_str ("vamp"), is_known = .true.)
    call global%set_log (var_str ("?use_vamp_equivalences"),&
         .true., is_known = .true.)
    call global%set_log (var_str ("?vis_history"),&
         .false., is_known = .true.)    
    call global%set_log (var_str ("?integration_timer"),&
         .false., is_known = .true.)    
    call global%set_log (var_str ("?phs_s_mapping"),&
         .false., is_known = .true.)    
    call global%set_int (var_str ("seed"), &
         0, is_known=.true.)
    
    call global%set_real (var_str ("sqrts"),&
         1000._default, is_known = .true.)

    call global%it_list%init ([3, 3, 3], [1000, 1000, 1000], &
         adapt = [.true., .true., .false.], &
         adapt_code = [var_str ("wg"), var_str ("g"), var_str ("")])

    write (u, "(A)")  "* Integrate"
    write (u, "(A)")

    call reset_interaction_counter ()
    call integrate_process (procname, global, local_stack=.true.)
    
    call global%pacify (efficiency_reset = .true., error_reset = .true.)
    call global%write (u, vars = no_vars, pacify = .true.)

    write (u, "(A)")
    write (u, "(A)")  "* Cleanup"

    call global%final ()
    call syntax_phs_forest_final ()
    call syntax_model_file_final ()
    
    write (u, "(A)")
    write (u, "(A)")  "* Generated phase-space file"
    write (u, "(A)")

    phs_file = procname // "_i1.r1.phs"
    inquire (file = char (phs_file), exist = exist)
    if (exist) then
       u_phs = free_unit ()
       open (u_phs, file = char (phs_file), action = "read", status = "old")
       iostat = 0
       do while (iostat == 0)
          read (u_phs, "(A)", iostat = iostat)  buffer
          if (iostat == 0)  write (u, "(A)")  trim (buffer)
       end do
       close (u_phs)
    else
       write (u, "(A)")  "[file is missing]"
    end if

    write (u, "(A)")
    write (u, "(A)")  "* Test output end: integrations_7"
    
  end subroutine integrations_7
  
\eatline
\nwendcode{}\nwbegindocs{2098}\nwdocspar
\subsubsection{Structure functions}
Compile and integrate an intrinsic test matrix element ({\tt{}prc{\char95}test}
type) using the default ({\tt{}phs{\char95}wood}) phase-space
implementation.  The integration method is {\tt{}vamp}.  There is a structure
function of type {\tt{}unit{\char95}test}.

We use a test structure function $f(x)=x$ for both beams.  Together with the
$1/x_1x_2$ factor from the phase-space flux and a unit matrix element, we
should get the same result as previously for the process without structure
functions.  There is a slight correction due to the $m_s$ mass which we set to
zero here.
\nwenddocs{}\nwbegincode{2099}\moddef{Integrations: execute tests}\plusendmoddef
  call test (integrations_8, "integrations_8", &
       "integration with structure function", &
       u, results)
\nwendcode{}\nwbegincode{2100}\moddef{Integrations: tests}\plusendmoddef
  subroutine integrations_8 (u)
    integer, intent(in) :: u
    type(string_t) :: libname, procname
    type(rt_data_t), target :: global
    type(flavor_t) :: flv
    type(string_t) :: name
    
    write (u, "(A)")  "* Test output: integrations_8"
    write (u, "(A)")  "*   Purpose: integrate test process using VAMP &
         &with structure function"
    write (u, "(A)")

    write (u, "(A)")  "* Initialize process and parameters"
    write (u, "(A)")

    call syntax_model_file_init ()
    call syntax_phs_forest_init ()

    call global%global_init ()

    libname = "integrations_8_lib"
    procname = "integrations_8"
    
    call prepare_test_library (global, libname, 1, [procname])
    call compile_library (libname, global)

    call global%append_log (&
         var_str ("?rebuild_phase_space"), .true., intrinsic = .true.)
    call global%append_log (&
         var_str ("?rebuild_grids"), .true., intrinsic = .true.)

    call global%set_string (var_str ("$run_id"), &
         var_str ("r1"), is_known = .true.)
    call global%set_string (var_str ("$method"), &
         var_str ("unit_test"), is_known = .true.)
    call global%set_string (var_str ("$phs_method"), &
         var_str ("wood"), is_known = .true.)
    call global%set_string (var_str ("$integration_method"),&
         var_str ("vamp"), is_known = .true.)
    call global%set_log (var_str ("?use_vamp_equivalences"),&
         .true., is_known = .true.)
    call global%set_log (var_str ("?vis_history"),&
         .false., is_known = .true.)    
    call global%set_log (var_str ("?integration_timer"),&
         .false., is_known = .true.)    
    call global%set_log (var_str ("?phs_s_mapping"),&
         .false., is_known = .true.)  
    call global%set_int (var_str ("seed"), &
         0, is_known=.true.)    
    
    call global%set_real (var_str ("sqrts"),&
         1000._default, is_known = .true.)
    call global%model_set_real (var_str ("ms"), 0._default)

    call reset_interaction_counter ()

    call flv%init (25, global%model)
         
    name = flv%get_name ()
    call global%beam_structure%init_sf ([name, name], [1])
    call global%beam_structure%set_sf (1, 1, var_str ("sf_test_1"))

    write (u, "(A)")  "* Integrate"
    write (u, "(A)")

    call global%it_list%init ([1], [1000])
    call integrate_process (procname, global, local_stack=.true.)
    
    call global%write (u, vars = [var_str ("ms")])

    write (u, "(A)")
    write (u, "(A)")  "* Cleanup"

    call global%final ()
    call syntax_phs_forest_final ()
    call syntax_model_file_final ()

    write (u, "(A)")
    write (u, "(A)")  "* Test output end: integrations_8"
    
  end subroutine integrations_8
  
\eatline
\nwendcode{}\nwbegindocs{2101}\nwdocspar
\subsubsection{Integration history for VAMP integration with default
  phase space}
This test is only run when event analysis can be done.
\nwenddocs{}\nwbegincode{2102}\moddef{Integrations: execute history tests}\endmoddef
  call test (integrations_history_1, "integrations_history_1", &
       "Test integration history files", &
       u, results)
\nwendcode{}\nwbegincode{2103}\moddef{Integrations: tests}\plusendmoddef
  subroutine integrations_history_1 (u)
    integer, intent(in) :: u
    type(string_t) :: libname, procname
    type(rt_data_t), target :: global
    type(string_t), dimension(0) :: no_vars
    integer :: iostat, u_his
    character(91) :: buffer
    type(string_t) :: his_file, ps_file, pdf_file
    logical :: exist, exist_ps, exist_pdf
    
    write (u, "(A)")  "* Test output: integrations_history_1"
    write (u, "(A)")  "*   Purpose: test integration history files"
    write (u, "(A)")

    write (u, "(A)")  "* Initialize process and parameters"
    write (u, "(A)")

    call syntax_model_file_init ()
    call syntax_phs_forest_init ()

    call global%global_init ()

    libname = "integrations_history_1_lib"
    procname = "integrations_history_1"

    call global%set_log (var_str ("?vis_history"), &
         .true., is_known = .true.)        
    call global%set_log (var_str ("?integration_timer"),&
         .false., is_known = .true.)    
    call global%set_log (var_str ("?phs_s_mapping"),&
         .false., is_known = .true.)    
    
    call prepare_test_library (global, libname, 1, [procname])
    call compile_library (libname, global)

    call global%append_log (&
         var_str ("?rebuild_phase_space"), .true., intrinsic = .true.)
    call global%append_log (&
         var_str ("?rebuild_grids"), .true., intrinsic = .true.)

    call global%set_string (var_str ("$run_id"), &
         var_str ("r1"), is_known = .true.)
    call global%set_string (var_str ("$method"), &
         var_str ("unit_test"), is_known = .true.)
    call global%set_string (var_str ("$phs_method"), &
         var_str ("wood"), is_known = .true.)
    call global%set_string (var_str ("$integration_method"),&
         var_str ("vamp"), is_known = .true.)
    call global%set_log (var_str ("?use_vamp_equivalences"),&
         .true., is_known = .true.)
    call global%set_real (var_str ("error_threshold"),&
         5E-6_default, is_known = .true.)
    call global%set_int (var_str ("seed"), &
         0, is_known=.true.)    

    call global%set_real (var_str ("sqrts"),&
         1000._default, is_known = .true.)

    call global%it_list%init ([2, 2, 2], [1000, 1000, 1000], &
         adapt = [.true., .true., .false.], &
         adapt_code = [var_str ("wg"), var_str ("g"), var_str ("")])

    write (u, "(A)")  "* Integrate"
    write (u, "(A)")

    call reset_interaction_counter ()
    call integrate_process (procname, global, local_stack=.true., &
         eff_reset = .true.)
    
    call global%pacify (efficiency_reset = .true., error_reset = .true.)
    call global%write (u, vars = no_vars, pacify = .true.)
    
    write (u, "(A)")
    write (u, "(A)")  "* Generated history files"
    write (u, "(A)")

    his_file = procname // ".r1-history.tex"
    ps_file  = procname // ".r1-history.ps"
    pdf_file = procname // ".r1-history.pdf"
    inquire (file = char (his_file), exist = exist)
    if (exist) then
       u_his = free_unit ()
       open (u_his, file = char (his_file), action = "read", status = "old")
       iostat = 0
       do while (iostat == 0)
          read (u_his, "(A)", iostat = iostat)  buffer
          if (iostat == 0)  write (u, "(A)")  trim (buffer)
       end do
       close (u_his)
    else
       write (u, "(A)")  "[History LaTeX file is missing]"
    end if
    inquire (file = char (ps_file), exist = exist_ps)
    if (exist_ps) then
       write (u, "(A)")  "[History Postscript file exists and is nonempty]"
    else
       write (u, "(A)")  "[History Postscript file is missing/non-regular]"
    end if
    inquire (file = char (pdf_file), exist = exist_pdf)
    if (exist_pdf) then
       write (u, "(A)")  "[History PDF file exists and is nonempty]"
    else
       write (u, "(A)")  "[History PDF file is missing/non-regular]"
    end if    
    
    write (u, "(A)")
    write (u, "(A)")  "* Cleanup"

    call global%final ()
    call syntax_phs_forest_final ()
    call syntax_model_file_final ()    
    
    write (u, "(A)")
    write (u, "(A)")  "* Test output end: integrations_history_1"
    
  end subroutine integrations_history_1
  
\eatline
\nwendcode{}\nwbegindocs{2104}\nwdocspar
\clearpage
%%%%%%%%%%%%%%%%%%%%%%%%%%%%%%%%%%%%%%%%%%%%%%%%%%%%%%%%%%%%%%%%%%%%%%%%
\section{Event Streams}
This module manages I/O from/to multiple concurrent event streams.
Usually, there is at most one input stream, but several output
streams.  For the latter, we set up an array which can hold {\tt{}eio{\char95}t}
(event I/O) objects of different dynamic types simultaneously.  One of
them may be marked as an input channel.
\nwenddocs{}\nwbegincode{2105}\moddef{\code{}event{\char95}streams.f90\edoc{}}\endmoddef
\LA{}File header\RA{}

module event_streams
  
\LA{}Use kinds\RA{}
\LA{}Use strings\RA{}
  use io_units
  use unit_tests
  use diagnostics
  use variables
  use model_data
  use models
  use processes
  use events
  use eio_data
  use eio_base
  use rt_data
  use dispatch
  
\LA{}Standard module head\RA{}

\LA{}Event streams: public\RA{}

\LA{}Event streams: types\RA{}

contains

\LA{}Event streams: tests\RA{}

\LA{}Event streams: procedures\RA{}

end module event_streams
\eatline
\nwendcode{}\nwbegindocs{2106}\nwdocspar
\subsection{Event Stream Array}
Each entry is an {\tt{}eio{\char95}t} object.  Since the type is dynamic, we need
a wrapper:
\nwenddocs{}\nwbegincode{2107}\moddef{Event streams: types}\endmoddef
  type :: event_stream_entry_t
     class(eio_t), allocatable :: eio
  end type event_stream_entry_t
  
\eatline
\nwendcode{}\nwbegindocs{2108}An array of event-stream entry objects.  If one of the entries is an
input channel, {\tt{}i{\char95}in} is the corresponding index.
\nwenddocs{}\nwbegincode{2109}\moddef{Event streams: public}\endmoddef
  public :: event_stream_array_t
\nwendcode{}\nwbegincode{2110}\moddef{Event streams: types}\plusendmoddef
  type :: event_stream_array_t
     type(event_stream_entry_t), dimension(:), allocatable :: entry
     integer :: i_in = 0
   contains
   \LA{}Event streams: event stream array: TBP\RA{}
  end type event_stream_array_t
  
\eatline
\nwendcode{}\nwbegindocs{2111}Output.
\nwenddocs{}\nwbegincode{2112}\moddef{Event streams: event stream array: TBP}\endmoddef
  procedure :: write => event_stream_array_write
\nwendcode{}\nwbegincode{2113}\moddef{Event streams: procedures}\endmoddef
  subroutine event_stream_array_write (object, unit)
    class(event_stream_array_t), intent(in) :: object
    integer, intent(in), optional :: unit
    integer :: u, i
    u = given_output_unit (unit)
    write (u, "(1x,A)")  "Event stream array:"
    if (allocated (object%entry)) then
       select case (size (object%entry))
       case (0)
          write (u, "(3x,A)")  "[empty]"
       case default
          do i = 1, size (object%entry)
             if (i == object%i_in)  write (u, "(1x,A)")  "Input stream:"
             call object%entry(i)%eio%write (u)
          end do
       end select
    else
       write (u, "(3x,A)")  "[undefined]"
    end if
  end subroutine event_stream_array_write

\eatline
\nwendcode{}\nwbegindocs{2114}Finalize all streams.
\nwenddocs{}\nwbegincode{2115}\moddef{Event streams: event stream array: TBP}\plusendmoddef
  procedure :: final => event_stream_array_final
\nwendcode{}\nwbegincode{2116}\moddef{Event streams: procedures}\plusendmoddef
  subroutine event_stream_array_final (es_array)
    class(event_stream_array_t), intent(inout) :: es_array
    integer :: i
    do i = 1, size (es_array%entry)
       call es_array%entry(i)%eio%final ()
    end do
  end subroutine event_stream_array_final

\eatline
\nwendcode{}\nwbegindocs{2117}Initialization.  We use a generic {\tt{}sample} name, open event I/O
objects for all provided stream types (using the {\tt{}dispatch{\char95}eio}
routine), and initialize for the given list of process pointers.  If
there is an {\tt{}input} argument, this channel is initialized as an input
channel and appended to the array.

The {\tt{}input{\char95}data} or, if not present, {\tt{}data} may be modified.  This
happens if we open a stream for reading and get new information there.
\nwenddocs{}\nwbegincode{2118}\moddef{Event streams: event stream array: TBP}\plusendmoddef
  procedure :: init => event_stream_array_init
\nwendcode{}\nwbegincode{2119}\moddef{Event streams: procedures}\plusendmoddef
  subroutine event_stream_array_init &
       (es_array, sample, stream_fmt, global, &
       data, input, input_sample, input_data, allow_switch, checkpoint, &
       error)
    class(event_stream_array_t), intent(out) :: es_array
    type(string_t), intent(in) :: sample
    type(string_t), dimension(:), intent(in) :: stream_fmt
    type(rt_data_t), intent(in) :: global
    type(event_sample_data_t), intent(inout), optional :: data
    type(string_t), intent(in), optional :: input
    type(string_t), intent(in), optional :: input_sample
    type(event_sample_data_t), intent(inout), optional :: input_data
    logical, intent(in), optional :: allow_switch
    integer, intent(in), optional :: checkpoint
    logical, intent(out), optional :: error
    type(string_t) :: sample_in
    integer :: n, i
    logical :: success, switch
    if (present (input)) then
       n = size (stream_fmt) + 1
    else
       n = size (stream_fmt)
    end if
    if (present (input_sample)) then
       sample_in = input_sample
    else
       sample_in = sample
    end if
    if (present (allow_switch)) then
       switch = allow_switch
    else
       switch = .true.
    end if
    if (present (error)) then
       error = .false.
    end if
    if (present (checkpoint)) then
       allocate (es_array%entry (n + 1))
       call dispatch_eio &
            (es_array%entry(n+1)%eio, var_str ("checkpoint"), global)
       call es_array%entry(n+1)%eio%init_out (sample, data)
    else
       allocate (es_array%entry (n))
    end if
    if (present (input)) then
       call dispatch_eio (es_array%entry(n)%eio, input, global)
       if (present (input_data)) then
          call es_array%entry(n)%eio%init_in (sample_in, input_data, success)
       else
          call es_array%entry(n)%eio%init_in (sample_in, data, success)
       end if
       if (success) then
          es_array%i_in = n
       else if (present (input_sample)) then
          if (present (error)) then
             error = .true.
          else
             call msg_fatal ("Events: &
                  &parameter mismatch in input, aborting")
          end if
       else
          call msg_message ("Events: &
               &parameter mismatch, discarding old event set")
          call es_array%entry(n)%eio%final ()
          if (switch) then
             call msg_message ("Events: generating new events")
             call es_array%entry(n)%eio%init_out (sample, data)
          end if
       end if
    end if
    do i = 1, size (stream_fmt)
       call dispatch_eio (es_array%entry(i)%eio, stream_fmt(i), global)
       call es_array%entry(i)%eio%init_out (sample, data)
    end do
  end subroutine event_stream_array_init
  
\eatline
\nwendcode{}\nwbegindocs{2120}Switch the (only) input channel to an output channel, so further
events are appended to the respective stream.
\nwenddocs{}\nwbegincode{2121}\moddef{Event streams: event stream array: TBP}\plusendmoddef
  procedure :: switch_inout => event_stream_array_switch_inout
\nwendcode{}\nwbegincode{2122}\moddef{Event streams: procedures}\plusendmoddef
  subroutine event_stream_array_switch_inout (es_array)
    class(event_stream_array_t), intent(inout) :: es_array
    integer :: n
    if (es_array%has_input ()) then
       n = es_array%i_in
       call es_array%entry(n)%eio%switch_inout ()
       es_array%i_in = 0
    else
       call msg_bug ("Reading events: switch_inout: no input stream selected")
    end if
  end subroutine event_stream_array_switch_inout
  
\eatline
\nwendcode{}\nwbegindocs{2123}Output an event (with given process number) to all output streams.
If there is no output stream, do nothing.
\nwenddocs{}\nwbegincode{2124}\moddef{Event streams: event stream array: TBP}\plusendmoddef
  procedure :: output => event_stream_array_output
\nwendcode{}\nwbegincode{2125}\moddef{Event streams: procedures}\plusendmoddef
  subroutine event_stream_array_output (es_array, event, i_prc, &
                                        event_index, passed, pacify)
    class(event_stream_array_t), intent(inout) :: es_array
    type(event_t), intent(in), target :: event
    integer, intent(in) :: i_prc, event_index
    logical, intent(in), optional :: passed, pacify
    integer :: i
    do i = 1, size (es_array%entry)
       if (i /= es_array%i_in) then
          associate (eio => es_array%entry(i)%eio)
            if (eio%split) then
               if (event_index > 1 .and. &
                    mod (event_index, eio%split_n_evt) == 1) then
                  call eio%split_out ()
               end if
            end if
            call eio%output (event, i_prc, reading = es_array%i_in /= 0, &
                 passed = passed, &
                 pacify = pacify)
          end associate
       end if
    end do
  end subroutine event_stream_array_output
  
\eatline
\nwendcode{}\nwbegindocs{2126}Input the {\tt{}i{\char95}prc} index which selects the process for the current
event.  This is separated from reading the event, because it
determines which event record to read.  {\tt{}iostat} may indicate an
error or an EOF condition, as usual.
\nwenddocs{}\nwbegincode{2127}\moddef{Event streams: event stream array: TBP}\plusendmoddef
  procedure :: input_i_prc => event_stream_array_input_i_prc
\nwendcode{}\nwbegincode{2128}\moddef{Event streams: procedures}\plusendmoddef
  subroutine event_stream_array_input_i_prc (es_array, i_prc, iostat)
    class(event_stream_array_t), intent(inout) :: es_array
    integer, intent(out) :: i_prc
    integer, intent(out) :: iostat
    integer :: n
    if (es_array%has_input ()) then
       n = es_array%i_in
       call es_array%entry(n)%eio%input_i_prc (i_prc, iostat)
    else
       call msg_fatal ("Reading events: no input stream selected")
    end if
  end subroutine event_stream_array_input_i_prc
  
\eatline
\nwendcode{}\nwbegindocs{2129}Input an event from the selected input stream.  {\tt{}iostat} may indicate an
error or an EOF condition, as usual.
\nwenddocs{}\nwbegincode{2130}\moddef{Event streams: event stream array: TBP}\plusendmoddef
  procedure :: input_event => event_stream_array_input_event
\nwendcode{}\nwbegincode{2131}\moddef{Event streams: procedures}\plusendmoddef
  subroutine event_stream_array_input_event (es_array, event, iostat)
    class(event_stream_array_t), intent(inout) :: es_array
    type(event_t), intent(inout), target :: event
    integer, intent(out) :: iostat
    integer :: n
    if (es_array%has_input ()) then
       n = es_array%i_in
       call es_array%entry(n)%eio%input_event (event, iostat)
    else
       call msg_fatal ("Reading events: no input stream selected")
    end if
  end subroutine event_stream_array_input_event
  
\eatline
\nwendcode{}\nwbegindocs{2132}Return true if there is an input channel among the event streams.
\nwenddocs{}\nwbegincode{2133}\moddef{Event streams: event stream array: TBP}\plusendmoddef
  procedure :: has_input => event_stream_array_has_input
\nwendcode{}\nwbegincode{2134}\moddef{Event streams: procedures}\plusendmoddef
  function event_stream_array_has_input (es_array) result (flag)
    class(event_stream_array_t), intent(in) :: es_array
    logical :: flag
    flag = es_array%i_in /= 0
  end function event_stream_array_has_input
  
\eatline
\nwendcode{}\nwbegindocs{2135}\nwdocspar
\subsection{Unit Tests}
\nwenddocs{}\nwbegincode{2136}\moddef{Event streams: public}\plusendmoddef
  public :: event_streams_test
\nwendcode{}\nwbegincode{2137}\moddef{Event streams: tests}\endmoddef
  subroutine event_streams_test (u, results)
    integer, intent(in) :: u
    type(test_results_t), intent(inout) :: results
  \LA{}Event streams: execute tests\RA{}
  end subroutine event_streams_test
  
\eatline
\nwendcode{}\nwbegindocs{2138}\nwdocspar
\subsubsection{Empty event stream}
This should set up an empty event output stream array, including
initialization, output, and finalization (which are all no-ops).
\nwenddocs{}\nwbegincode{2139}\moddef{Event streams: execute tests}\endmoddef
  call test (event_streams_1, "event_streams_1", &
       "empty event stream array", &
       u, results)
\nwendcode{}\nwbegincode{2140}\moddef{Event streams: tests}\plusendmoddef
  subroutine event_streams_1 (u)
    integer, intent(in) :: u
    type(event_stream_array_t) :: es_array
    type(rt_data_t) :: global
    type(event_t) :: event
    type(string_t) :: sample
    type(string_t), dimension(0) :: empty_string_array

    write (u, "(A)")  "* Test output: event_streams_1"
    write (u, "(A)")  "*   Purpose: handle empty event stream array"
    write (u, "(A)")

    sample = "event_streams_1"

    call es_array%init (sample, empty_string_array, global)
    call es_array%output (event, 42, 1)
    call es_array%write (u)
    call es_array%final ()
    
    write (u, "(A)")
    write (u, "(A)")  "* Test output end: event_streams_1"
    
  end subroutine event_streams_1
  
\eatline
\nwendcode{}\nwbegindocs{2141}\nwdocspar
\subsubsection{Nontrivial event stream}
Here we generate a trivial event and choose {\tt{}raw} output as an entry in
the stream array.
\nwenddocs{}\nwbegincode{2142}\moddef{Event streams: execute tests}\plusendmoddef
  call test (event_streams_2, "event_streams_2", &
       "nontrivial event stream array", &
       u, results)
\nwendcode{}\nwbegincode{2143}\moddef{Event streams: tests}\plusendmoddef
  subroutine event_streams_2 (u)
    integer, intent(in) :: u
    type(event_stream_array_t) :: es_array
    type(rt_data_t) :: global
    type(model_data_t), target :: model
    type(event_t), allocatable, target :: event
    type(process_t), allocatable, target :: process
    type(process_instance_t), allocatable, target :: process_instance
    type(string_t) :: sample
    type(string_t), dimension(0) :: empty_string_array
    integer :: i_prc, iostat

    write (u, "(A)")  "* Test output: event_streams_2"
    write (u, "(A)")  "*   Purpose: handle empty event stream array"
    write (u, "(A)")

    call syntax_model_file_init ()
    call global%global_init ()
    call global%init_fallback_model &
         (var_str ("SM_hadrons"), var_str ("SM_hadrons.mdl"))

    call model%init_test ()

    write (u, "(A)")  "* Generate test process event"
    write (u, "(A)")

    allocate (process)
    allocate (process_instance)
    call prepare_test_process (process, process_instance, model)
    call process_instance%setup_event_data ()

    allocate (event)
    call event%basic_init ()
    call event%connect (process_instance, process%get_model_ptr ())
    call event%generate (1, [0.4_default, 0.4_default])
    call event%evaluate_expressions ()
    call event%write (u)

    write (u, "(A)")
    write (u, "(A)") "* Allocate raw eio stream and write event to file"
    write (u, "(A)")

    sample = "event_streams_2"

    call es_array%init (sample, [var_str ("raw")], global)
    call es_array%output (event, 1, 1)
    call es_array%write (u)
    call es_array%final ()
    
    write (u, "(A)")
    write (u, "(A)") "* Reallocate raw eio stream for reading"
    write (u, "(A)")

    sample = "foo"
    call es_array%init (sample, empty_string_array, global, &
         input = var_str ("raw"), input_sample = var_str ("event_streams_2"))
    call es_array%write (u)

    write (u, "(A)")
    write (u, "(A)") "* Reread event"
    write (u, "(A)")
    
    call es_array%input_i_prc (i_prc, iostat)
    
    write (u, "(1x,A,I0)")  "i_prc = ", i_prc
    write (u, "(A)")
    call es_array%input_event (event, iostat)
    call es_array%final ()
    
    call event%write (u)
    
    call global%final ()

    call model%final ()
    call syntax_model_file_final ()

    write (u, "(A)")
    write (u, "(A)")  "* Test output end: event_streams_2"
    
  end subroutine event_streams_2
  
\eatline
\nwendcode{}\nwbegindocs{2144}\nwdocspar
\subsubsection{Switch in/out}
Here we generate an event file and test switching from writing to
reading when the file is exhausted.
\nwenddocs{}\nwbegincode{2145}\moddef{Event streams: execute tests}\plusendmoddef
  call test (event_streams_3, "event_streams_3", &
       "switch input/output", &
       u, results)
\nwendcode{}\nwbegincode{2146}\moddef{Event streams: tests}\plusendmoddef
  subroutine event_streams_3 (u)
    integer, intent(in) :: u
    type(event_stream_array_t) :: es_array
    type(rt_data_t) :: global
    type(model_data_t), target :: model
    type(event_t), allocatable, target :: event
    type(process_t), allocatable, target :: process
    type(process_instance_t), allocatable, target :: process_instance
    type(string_t) :: sample
    type(string_t), dimension(0) :: empty_string_array
    integer :: i_prc, iostat

    write (u, "(A)")  "* Test output: event_streams_3"
    write (u, "(A)")  "*   Purpose: handle in/out switching"
    write (u, "(A)")

    call syntax_model_file_init ()
    call global%global_init ()
    call global%init_fallback_model &
         (var_str ("SM_hadrons"), var_str ("SM_hadrons.mdl"))

    call model%init_test ()

    write (u, "(A)")  "* Generate test process event"
    write (u, "(A)")

    allocate (process)
    allocate (process_instance)
    call prepare_test_process (process, process_instance, model)
    call process_instance%setup_event_data ()

    allocate (event)
    call event%basic_init ()
    call event%connect (process_instance, process%get_model_ptr ())
    call event%generate (1, [0.4_default, 0.4_default])
    call event%evaluate_expressions ()

    write (u, "(A)") "* Allocate raw eio stream and write event to file"
    write (u, "(A)")

    sample = "event_streams_3"

    call es_array%init (sample, [var_str ("raw")], global)
    call es_array%output (event, 1, 1)
    call es_array%write (u)
    call es_array%final ()
    
    write (u, "(A)")
    write (u, "(A)") "* Reallocate raw eio stream for reading"
    write (u, "(A)")

    call es_array%init (sample, empty_string_array, global, &
         input = var_str ("raw"))
    call es_array%write (u)

    write (u, "(A)")
    write (u, "(A)") "* Reread event"
    write (u, "(A)")
    
    call es_array%input_i_prc (i_prc, iostat)
    call es_array%input_event (event, iostat)

    write (u, "(A)") "* Attempt to read another event (fail), then generate"
    write (u, "(A)")
    
    call es_array%input_i_prc (i_prc, iostat)
    if (iostat < 0) then
       call es_array%switch_inout ()
       call event%generate (1, [0.3_default, 0.3_default])
       call event%evaluate_expressions ()
       call es_array%output (event, 1, 2)
    end if
    call es_array%write (u)
    call es_array%final ()
    
    write (u, "(A)")
    call event%write (u)

    write (u, "(A)")
    write (u, "(A)") "* Reallocate raw eio stream for reading"
    write (u, "(A)")

    call es_array%init (sample, empty_string_array, global, &
         input = var_str ("raw"))
    call es_array%write (u)

    write (u, "(A)")
    write (u, "(A)") "* Reread two events and display 2nd event"
    write (u, "(A)")
    
    call es_array%input_i_prc (i_prc, iostat)
    call es_array%input_event (event, iostat)
    call es_array%input_i_prc (i_prc, iostat)
    
    call es_array%input_event (event, iostat)
    call es_array%final ()

    call event%write (u)
    
    call global%final ()
    
    call model%final ()
    call syntax_model_file_final ()

    write (u, "(A)")
    write (u, "(A)")  "* Test output end: event_streams_3"
    
  end subroutine event_streams_3
  
\eatline
\nwendcode{}\nwbegindocs{2147}\nwdocspar
\subsubsection{Checksum}
Here we generate an event file and repeat twice, once with identical
parameters and once with modified parameters.
\nwenddocs{}\nwbegincode{2148}\moddef{Event streams: execute tests}\plusendmoddef
  call test (event_streams_4, "event_streams_4", &
       "check MD5 sum", &
       u, results)
\nwendcode{}\nwbegincode{2149}\moddef{Event streams: tests}\plusendmoddef
  subroutine event_streams_4 (u)
    integer, intent(in) :: u
    type(event_stream_array_t) :: es_array
    type(rt_data_t) :: global
    type(process_t), allocatable, target :: process
    type(string_t) :: sample
    type(string_t), dimension(0) :: empty_string_array
    type(event_sample_data_t) :: data

    write (u, "(A)")  "* Test output: event_streams_4"
    write (u, "(A)")  "*   Purpose: handle in/out switching"
    write (u, "(A)")

    write (u, "(A)")  "* Generate test process event"
    write (u, "(A)")

    call syntax_model_file_init ()
    call global%global_init ()
    call global%init_fallback_model &
         (var_str ("SM_hadrons"), var_str ("SM_hadrons.mdl"))
    
    call global%set_log (var_str ("?check_event_file"), &
         .true., is_known = .true.)

    allocate (process)

    write (u, "(A)") "* Allocate raw eio stream for writing"
    write (u, "(A)")

    sample = "event_streams_4"
    data%md5sum_cfg = "1234567890abcdef1234567890abcdef"

    call es_array%init (sample, [var_str ("raw")], global, data)
    call es_array%write (u)
    call es_array%final ()
    
    write (u, "(A)")
    write (u, "(A)") "* Reallocate raw eio stream for reading"
    write (u, "(A)")

    call es_array%init (sample, empty_string_array, global, &
         data, input = var_str ("raw"))
    call es_array%write (u)
    call es_array%final ()

    write (u, "(A)")
    write (u, "(A)") "* Reallocate modified raw eio stream for reading (fail)"
    write (u, "(A)")

    data%md5sum_cfg = "1234567890______1234567890______"
    call es_array%init (sample, empty_string_array, global, &
         data, input = var_str ("raw"))
    call es_array%write (u)
    call es_array%final ()
    
    write (u, "(A)")
    write (u, "(A)") "* Repeat ignoring checksum"
    write (u, "(A)")

    call global%set_log (var_str ("?check_event_file"), &
         .false., is_known = .true.)
    call es_array%init (sample, empty_string_array, global, &
         data, input = var_str ("raw"))
    call es_array%write (u)
    call es_array%final ()
    
    call global%final ()
    call syntax_model_file_final ()

    write (u, "(A)")
    write (u, "(A)")  "* Test output end: event_streams_4"
    
  end subroutine event_streams_4
  
\eatline
\nwendcode{}\nwbegindocs{2150}\nwdocspar
\clearpage
%%%%%%%%%%%%%%%%%%%%%%%%%%%%%%%%%%%%%%%%%%%%%%%%%%%%%%%%%%%%%%%%%%%%%%%%
\section{Simulation}
This module manages simulation: event generation and reading/writing of event
files.  The {\tt{}simulation} object is intended to be used (via a pointer)
outside of \whizard, if events are generated individually by an external
driver.
\nwenddocs{}\nwbegincode{2151}\moddef{\code{}simulations.f90\edoc{}}\endmoddef
\LA{}File header\RA{}

module simulations

\LA{}Use kinds\RA{}
\LA{}Use strings\RA{}
  use io_units
  use format_utils, only: write_separator
  use format_defs, only: FMT_19
  use unit_tests
  use diagnostics
  use sm_qcd
  use md5
  use ifiles
  use lexers
  use parser
  use variables
  use eval_trees
  use model_data
  use flavors
  use particles
  use state_matrices
  use interactions
  use models
  use beams
  use phs_forests
  use rng_base
  use selectors
  use prc_core
  use prclib_stacks
  use processes
  use event_base
  use events
  use event_transforms
  use decays
  use eio_data
  use eio_base
  use eio_raw
  use eio_ascii
  use rt_data
  use dispatch
  use process_configurations
  use compilations
  use integrations
  use event_streams

\LA{}Standard module head\RA{}

\LA{}Simulations: public\RA{}

\LA{}Simulations: types\RA{}

\LA{}Simulations: interfaces\RA{}

contains

\LA{}Simulations: procedures\RA{}

\LA{}Simulations: tests\RA{}

end module simulations
\eatline
\nwendcode{}\nwbegindocs{2152}\nwdocspar
\subsection{Event counting}
In this object we collect statistical information about an event
sample or sub-sample.
\nwenddocs{}\nwbegincode{2153}\moddef{Simulations: types}\endmoddef
  type :: counter_t
     integer :: total = 0
     integer :: generated = 0
     integer :: read = 0
     integer :: positive = 0
     integer :: negative = 0
     integer :: zero = 0
     integer :: excess = 0
     real(default) :: max_excess = 0
     real(default) :: sum_excess = 0
   contains
   \LA{}Simulations: counter: TBP\RA{}
  end type counter_t
  
\eatline
\nwendcode{}\nwbegindocs{2154}Output.
\nwenddocs{}\nwbegincode{2155}\moddef{Simulations: counter: TBP}\endmoddef
  procedure :: write => counter_write
\nwendcode{}\nwbegincode{2156}\moddef{Simulations: procedures}\endmoddef
  subroutine counter_write (object, unit)
    class(counter_t), intent(in) :: object
    integer, intent(in), optional :: unit
    integer :: u
    u = given_output_unit (unit)
1   format (3x,A,I0)
2   format (5x,A,I0)
3   format (5x,A,ES19.12)
    write (u, 1)  "Events total      = ", object%total
    write (u, 2)  "generated       = ", object%generated
    write (u, 2)  "read            = ", object%read
    write (u, 2)  "positive weight = ", object%positive
    write (u, 2)  "negative weight = ", object%negative
    write (u, 2)  "zero weight     = ", object%zero
    write (u, 2)  "excess weight   = ", object%excess
    if (object%excess /= 0) then
       write (u, 3)  "max excess      = ", object%max_excess
       write (u, 3)  "avg excess      = ", object%sum_excess / object%total
    end if
  end subroutine counter_write

\eatline
\nwendcode{}\nwbegindocs{2157}This is a screen message: if there was an excess, display statistics.
\nwenddocs{}\nwbegincode{2158}\moddef{Simulations: counter: TBP}\plusendmoddef
  procedure :: show_excess => counter_show_excess
\nwendcode{}\nwbegincode{2159}\moddef{Simulations: procedures}\plusendmoddef
  subroutine counter_show_excess (counter)
    class(counter_t), intent(in) :: counter
    if (counter%excess > 0) then
       write (msg_buffer, "(A,1x,I0,1x,A,1x,'(',F7.3,' %)')") &
            "Encountered events with excess weight:", counter%excess, &
            "events", 100 * counter%excess / real (counter%total)
       call msg_warning ()
       write (msg_buffer, "(A,ES10.3)") &
            "Maximum excess weight =", counter%max_excess
       call msg_message ()
       write (msg_buffer, "(A,ES10.3)") &
            "Average excess weight =", counter%sum_excess / counter%total
       call msg_message ()
    end if
  end subroutine counter_show_excess
    
\eatline
\nwendcode{}\nwbegindocs{2160}Count an event.  The weight and event source are optional; by
default we assume that the event has been generated and has positive
weight.
\nwenddocs{}\nwbegincode{2161}\moddef{Simulations: counter: TBP}\plusendmoddef
  procedure :: record => counter_record
\nwendcode{}\nwbegincode{2162}\moddef{Simulations: procedures}\plusendmoddef
  subroutine counter_record (counter, weight, excess, from_file)
    class(counter_t), intent(inout) :: counter
    real(default), intent(in), optional :: weight, excess
    logical, intent(in), optional :: from_file
    counter%total = counter%total + 1
    if (present (from_file)) then
       if (from_file) then
          counter%read = counter%read + 1
       else
          counter%generated = counter%generated + 1
       end if
    else
       counter%generated = counter%generated + 1
    end if
    if (present (weight)) then
       if (weight > 0) then
          counter%positive = counter%positive + 1
       else if (weight < 0) then
          counter%negative = counter%negative + 1
       else
          counter%zero = counter%zero + 1
       end if
    else
       counter%positive = counter%positive + 1
    end if
    if (present (excess)) then
       if (excess > 0) then
          counter%excess = counter%excess + 1
          counter%max_excess = max (counter%max_excess, excess)
          counter%sum_excess = counter%sum_excess + excess
       end if
    end if
  end subroutine counter_record
    
\eatline
\nwendcode{}\nwbegindocs{2163}\nwdocspar
\subsection{Simulation: component sets}
For each set of process components that share a MCI entry in the
process configuration, we keep a separate event record.
\nwenddocs{}\nwbegincode{2164}\moddef{Simulations: types}\plusendmoddef
  type :: mci_set_t
     private
     integer :: n_components = 0
     integer, dimension(:), allocatable :: i_component
     type(string_t), dimension(:), allocatable :: component_id
     logical :: has_integral = .false.
     real(default) :: integral = 0
     real(default) :: error = 0
     real(default) :: weight_mci = 0
     type(counter_t) :: counter
   contains
   \LA{}Simulations: mci set: TBP\RA{}
  end type mci_set_t
     
\eatline
\nwendcode{}\nwbegindocs{2165}Output.
\nwenddocs{}\nwbegincode{2166}\moddef{Simulations: mci set: TBP}\endmoddef
  procedure :: write => mci_set_write
\nwendcode{}\nwbegincode{2167}\moddef{Simulations: procedures}\plusendmoddef
  subroutine mci_set_write (object, unit)
    class(mci_set_t), intent(in) :: object
    integer, intent(in), optional :: unit
    integer :: u, i
    u = given_output_unit (unit)
    write (u, "(3x,A)")  "Components:"
    do i = 1, object%n_components
       write (u, "(5x,I0,A,A,A)")  object%i_component(i), &
            ": '", char (object%component_id(i)), "'"
    end do
    if (object%has_integral) then
       write (u, "(3x,A," // FMT_19 // ")")  "Integral  = ", object%integral
       write (u, "(3x,A," // FMT_19 // ")")  "Error     = ", object%error
       write (u, "(3x,A,F13.10)")  "Weight    =", object%weight_mci
    else
       write (u, "(3x,A)")  "Integral  = [undefined]"
    end if
    call object%counter%write (u)
  end subroutine mci_set_write
  
\eatline
\nwendcode{}\nwbegindocs{2168}Initialize: Get the indices and names for the process components
that will contribute to this set.
\nwenddocs{}\nwbegincode{2169}\moddef{Simulations: mci set: TBP}\plusendmoddef
  procedure :: init => mci_set_init
\nwendcode{}\nwbegincode{2170}\moddef{Simulations: procedures}\plusendmoddef
  subroutine mci_set_init (object, i_mci, process)
    class(mci_set_t), intent(out) :: object
    integer, intent(in) :: i_mci
    type(process_t), intent(in), target :: process
    integer :: i
    call process%get_i_component (i_mci, object%i_component)
    object%n_components = size (object%i_component)
    allocate (object%component_id (object%n_components))
    do i = 1, size (object%component_id)
       object%component_id(i) = &
            process%get_component_id (object%i_component(i))
    end do
    if (process%has_integral (i_mci)) then
       object%integral = process%get_integral (i_mci)
       object%error = process%get_error (i_mci)
       object%has_integral = .true.
    end if
  end subroutine mci_set_init
    
\eatline
\nwendcode{}\nwbegindocs{2171}\nwdocspar
\subsection{Process-core Safe}
This is an object that temporarily holds a process core object.  We
need this while rescanning a process with modified parameters.  After
the rescan, we want to restore the original state.
\nwenddocs{}\nwbegincode{2172}\moddef{Simulations: types}\plusendmoddef
  type :: core_safe_t
     class(prc_core_t), allocatable :: core
  end type core_safe_t
  
\eatline
\nwendcode{}\nwbegindocs{2173}\nwdocspar
\subsection{Process Object}
The simulation works on process objects.  This subroutine makes a
process object available for simulation.  The process is in the
process stack.  If integration is not yet
done, do it.  Any generated process object should be put on the global stack,
if it is separate from the local one.
\nwenddocs{}\nwbegincode{2174}\moddef{Simulations: procedures}\plusendmoddef
  subroutine prepare_process &
       (process, process_id, integrate, local, global)
    type(process_t), pointer, intent(out) :: process
    type(string_t), intent(in) :: process_id
    logical, intent(in) :: integrate
    type(rt_data_t), intent(inout), target :: local
    type(rt_data_t), intent(inout), optional, target :: global
    if (present (global)) then
       process => global%process_stack%get_process_ptr (process_id)
    else
       process => local%process_stack%get_process_ptr (process_id)
    end if
    if (.not. associated (process)) then
       if (integrate) then
          call msg_message ("Simulate: process '" &
               // char (process_id) // "' needs integration")
       else
          call msg_message ("Simulate: process '" &
               // char (process_id) // "' needs initialization")
       end if
       if (present (global)) then
          call integrate_process (process_id, local, global, &
            init_only = .not. integrate)
       else
          call integrate_process (process_id, local, local_stack=.true., &
            init_only = .not. integrate)
       end if
       if (signal_is_pending ())  return
       process => global%process_stack%get_process_ptr (process_id)
       if (associated (process)) then
          if (integrate) then
             call msg_message ("Simulate: integration done")
             call global%process_stack%fill_result_vars (process_id)
          else
             call msg_message ("Simulate: process initialization done")
          end if
       else
          call msg_fatal ("Simulate: process '" &
               // char (process_id) // "' could not be initialized: aborting")
       end if
    end if
  end subroutine prepare_process
    
\eatline
\nwendcode{}\nwbegindocs{2175}\nwdocspar
\subsection{Simulation entry}
For each process that we consider for event generation, we need a
separate entry.  The entry separately records the process ID and run ID.  The
{\tt{}weight{\char95}mci} array is used for selecting a component set (which
shares a MCI record inside the process container) when generating an
event for the current process.

The simulation entry is an extension of the {\tt{}event{\char95}t} event record.
This core object contains configuration data, pointers to the process
and process instance, the expressions, flags and values that are
evaluated at runtime, and the resulting particle set.

The entry explicitly allocate the {\tt{}process{\char95}instance}, which becomes
the process-specific workspace for the event record.

If entries with differing environments are present simultaneously, we
may need to switch QCD parameters and/or the model event by event.  In
this case, the {\tt{}qcd} and/or {\tt{}model} components are present.
\nwenddocs{}\nwbegincode{2176}\moddef{Simulations: types}\plusendmoddef
  type, extends (event_t) :: entry_t
     private
     type(string_t) :: process_id
     type(string_t) :: library
     type(string_t) :: run_id
     logical :: has_integral = .false.
     real(default) :: integral = 0
     real(default) :: error = 0
     real(default) :: process_weight = 0
     logical :: valid = .false.
     type(counter_t) :: counter
     integer :: n_in = 0
     integer :: n_mci = 0
     type(mci_set_t), dimension(:), allocatable :: mci_set
     type(selector_t) :: mci_selector
     type(core_safe_t), dimension(:), allocatable :: core_safe
     class(model_data_t), pointer :: model => null ()
     type(qcd_t) :: qcd
     logical :: nlo_event = .false.
   contains
   \LA{}Simulations: entry: TBP\RA{}
  end type entry_t

\eatline
\nwendcode{}\nwbegindocs{2177}Output.  Write just the configuration, the event is written by a
separate routine.

The {\tt{}verbose} option is unused, it is required by the interface of
the base-object method.
\nwenddocs{}\nwbegincode{2178}\moddef{Simulations: entry: TBP}\endmoddef
  procedure :: write_config => entry_write_config
\nwendcode{}\nwbegincode{2179}\moddef{Simulations: procedures}\plusendmoddef
  subroutine entry_write_config (object, unit)
    class(entry_t), intent(in) :: object
    integer, intent(in), optional :: unit
    integer :: u, i
    u = given_output_unit (unit)
    write (u, "(3x,A,A,A)")  "Process   = '", char (object%process_id), "'"
    write (u, "(3x,A,A,A)")  "Library   = '", char (object%library), "'"
    write (u, "(3x,A,A,A)")  "Run       = '", char (object%run_id), "'"
    write (u, "(3x,A,L1)")   "is valid  = ", object%valid
    if (object%has_integral) then
       write (u, "(3x,A," // FMT_19 // ")")  "Integral  = ", object%integral
       write (u, "(3x,A," // FMT_19 // ")")  "Error     = ", object%error
       write (u, "(3x,A,F13.10)")  "Weight    =", object%process_weight
    else
       write (u, "(3x,A)")  "Integral  = [undefined]"
    end if
    write (u, "(3x,A,I0)")   "MCI sets  = ", object%n_mci
    call object%counter%write (u)
    do i = 1, size (object%mci_set)
       write (u, "(A)")
       write (u, "(1x,A,I0,A)")  "MCI set #", i, ":"
       call object%mci_set(i)%write (u)
    end do
    if (allocated (object%core_safe)) then
       do i = 1, size (object%core_safe)
          write (u, "(1x,A,I0,A)")  "Saved process-component core #", i, ":"
          call object%core_safe(i)%core%write (u)
       end do
    end if
  end subroutine entry_write_config
  
\eatline
\nwendcode{}\nwbegindocs{2180}Finalizer.  The {\tt{}instance} pointer component of the {\tt{}event{\char95}t}
base type points to a target which we did explicitly allocate in the
{\tt{}entry{\char95}init} procedure.  Therefore, we finalize and explicitly
deallocate it here.  Then we call the finalizer of the base type.
\nwenddocs{}\nwbegincode{2181}\moddef{Simulations: entry: TBP}\plusendmoddef
  procedure :: final => entry_final
\nwendcode{}\nwbegincode{2182}\moddef{Simulations: procedures}\plusendmoddef
  subroutine entry_final (object)
    class(entry_t), intent(inout) :: object
    integer :: i
    if (associated (object%instance)) then
       do i = 1, object%n_mci
          call object%instance%final_simulation (i)
       end do
       call object%instance%final ()
       deallocate (object%instance)
    end if
    call object%event_t%final ()
  end subroutine entry_final
  
\eatline
\nwendcode{}\nwbegindocs{2183}Initialization.  Search for a process entry and allocate a process
instance as an anonymous object, temporarily accessible via the
{\tt{}process{\char95}instance} pointer.  Assign data by looking at the process
object and at the environment.

If {\tt{}n{\char95}alt} is set, we prepare for additional alternate sqme and weight
entries.

If the process object is not found initially: if {\tt{}integrate} is set, attempt
an integration pass and try again.  Otherwise, just initialize the object.

If {\tt{}generate} is set, prepare the MCI objects for generating new events.
For pure rescanning, this is not necessary.

When done, we assign the {\tt{}instance} and {\tt{}process} pointers of the
base type by the {\tt{}connect} method, so we can reference them later.
\nwenddocs{}\nwbegincode{2184}\moddef{Simulations: entry: TBP}\plusendmoddef
  procedure :: init => entry_init
\nwendcode{}\nwbegincode{2185}\moddef{Simulations: procedures}\plusendmoddef
  subroutine entry_init &
       (entry, process_id, integrate, generate, update_sqme, &
       local, global, n_alt)
    class(entry_t), intent(inout), target :: entry
    type(string_t), intent(in) :: process_id
    logical, intent(in) :: integrate, generate, update_sqme
    type(rt_data_t), intent(inout), target :: local
    type(rt_data_t), intent(inout), optional, target :: global
    integer, intent(in), optional :: n_alt
    type(process_t), pointer :: process, master_process
    type(process_instance_t), pointer :: process_instance
    integer :: i
    
    call prepare_process (master_process, process_id, integrate, local, global)
    if (signal_is_pending ())  return

    if (.not. master_process%has_matrix_element ()) then
       entry%has_integral = .true.
       entry%process_id = process_id
       entry%valid = .false.          
       return
    end if
    
    call entry%basic_init (local%var_list, n_alt)

    entry%process_id = process_id
    if (generate .or. integrate) then
       entry%run_id = master_process%get_run_id ()
       process => master_process
    else
       call local%set_log (var_str ("?rebuild_phase_space"), &
            .false., is_known = .true.)
       call local%set_log (var_str ("?check_phs_file"), &
            .false., is_known = .true.)
       call local%set_log (var_str ("?rebuild_grids"), &
            .false., is_known = .true.)
       entry%run_id = &
            local%var_list%get_sval (var_str ("$run_id"))
       if (update_sqme) then
          call prepare_local_process (process, process_id, local)
       else
          process => master_process
       end if
    end if

    call entry%import_process_characteristics (process)

    allocate (entry%mci_set (entry%n_mci))
    do i = 1, size (entry%mci_set)
       call entry%mci_set(i)%init (i, master_process)
    end do
    if (process%is_nlo_calculation ()) then
      entry%nlo_event = .true.
      call entry%combine_mci_sets ()
    end if

    call entry%import_process_results (master_process)
    call entry%prepare_expressions (local)

    call prepare_process_instance (process_instance, process, local%model)
    if (generate) then
       do i = 1, entry%n_mci
          call process%prepare_simulation (i)
          call process_instance%init_simulation (i, entry%config%safety_factor)
       end do
    end if
    call entry%setup_event_transforms (process, local)

    if (present (global)) then
       call entry%connect (process_instance, local%model, global%process_stack)
    else
       call entry%connect (process_instance, local%model, local%process_stack)
    end if
    call entry%setup_expressions ()

    entry%model => process%get_model_ptr ()
    call dispatch_qcd (entry%qcd, local)
    entry%valid = .true.
    
  end subroutine entry_init
    
\eatline
\nwendcode{}\nwbegindocs{2186}Part of simulation-entry initialization: set up a process object for
local use.
\nwenddocs{}\nwbegincode{2187}\moddef{Simulations: procedures}\plusendmoddef
  subroutine prepare_local_process (process, process_id, local)
    type(process_t), pointer, intent(inout) :: process
    type(string_t), intent(in) :: process_id
    type(rt_data_t), intent(inout), target :: local
    type(integration_t) :: intg
    call intg%create_process (process_id)
    call intg%init_process (local)
    call intg%setup_process (local, verbose=.false.)
    process => intg%get_process_ptr ()
  end subroutine prepare_local_process
  
\eatline
\nwendcode{}\nwbegindocs{2188}Part of simulation-entry initialization: set up a process instance
matching the selected process object.

The model that we can provide as an extra argument can modify particle
settings (polarization) in the density matrices that will be constructed.  It
does not affect parameters.
\nwenddocs{}\nwbegincode{2189}\moddef{Simulations: procedures}\plusendmoddef
  subroutine prepare_process_instance (process_instance, process, model)
    type(process_instance_t), pointer, intent(inout) :: process_instance
    type(process_t), intent(inout), target :: process
    class(model_data_t), intent(in), optional :: model
    allocate (process_instance)
    call process_instance%init (process)
    call process_instance%setup_event_data (model)
  end subroutine prepare_process_instance 

\eatline
\nwendcode{}\nwbegindocs{2190}Part of simulation-entry initialization: query the
process for basic information.
\nwenddocs{}\nwbegincode{2191}\moddef{Simulations: entry: TBP}\plusendmoddef
  procedure, private :: import_process_characteristics &
       => entry_import_process_characteristics
\nwendcode{}\nwbegincode{2192}\moddef{Simulations: procedures}\plusendmoddef
  subroutine entry_import_process_characteristics (entry, process)
    class(entry_t), intent(inout) :: entry
    type(process_t), intent(in), target :: process
    entry%library = process%get_library_name ()
    entry%n_in = process%get_n_in ()
    entry%n_mci = process%get_n_mci ()
  end subroutine entry_import_process_characteristics

\eatline
\nwendcode{}\nwbegindocs{2193}Part of simulation-entry initialization: query the
process for integration results.
\nwenddocs{}\nwbegincode{2194}\moddef{Simulations: entry: TBP}\plusendmoddef
  procedure, private :: import_process_results &
       => entry_import_process_results
\nwendcode{}\nwbegincode{2195}\moddef{Simulations: procedures}\plusendmoddef
  subroutine entry_import_process_results (entry, process)
    class(entry_t), intent(inout) :: entry
    type(process_t), intent(in), target :: process
    if (process%has_integral ()) then
       entry%integral = process%get_integral ()
       entry%error = process%get_error ()
       call entry%set_sigma (entry%integral)
       entry%has_integral = .true.
    end if
  end subroutine entry_import_process_results

\eatline
\nwendcode{}\nwbegindocs{2196}Part of simulation-entry initialization: create expression factory
objects and store them.
\nwenddocs{}\nwbegincode{2197}\moddef{Simulations: entry: TBP}\plusendmoddef
  procedure, private :: prepare_expressions &
       => entry_prepare_expressions
\nwendcode{}\nwbegincode{2198}\moddef{Simulations: procedures}\plusendmoddef
  subroutine entry_prepare_expressions (entry, local)
    class(entry_t), intent(inout) :: entry
    type(rt_data_t), intent(in), target :: local
    type(eval_tree_factory_t) :: expr_factory
    call expr_factory%init (local%pn%selection_lexpr)
    call entry%set_selection (expr_factory)
    call expr_factory%init (local%pn%reweight_expr)
    call entry%set_reweight (expr_factory)
    call expr_factory%init (local%pn%analysis_lexpr)
    call entry%set_analysis (expr_factory)
  end subroutine entry_prepare_expressions

\eatline
\nwendcode{}\nwbegindocs{2199}Part of simulation-entry initialization: dispatch event transforms
(decay, shower) as requested.
\nwenddocs{}\nwbegincode{2200}\moddef{Simulations: entry: TBP}\plusendmoddef
  procedure, private :: setup_event_transforms &
       => entry_setup_event_transforms
\nwendcode{}\nwbegincode{2201}\moddef{Simulations: procedures}\plusendmoddef
  subroutine entry_setup_event_transforms (entry, process, local)
    class(entry_t), intent(inout) :: entry
    type(process_t), intent(in), target :: process
    type(rt_data_t), intent(in), target :: local
    class(evt_t), pointer :: evt
    logical :: enable_qcd
    if (process%contains_unstable (local%model)) then
       call dispatch_evt_decay (evt, local)
       if (associated (evt))  call entry%import_transform (evt)
    end if
    enable_qcd = local%get_lval (var_str ("?ps_isr_active")) &
            .or. local%get_lval (var_str ("?ps_fsr_active")) &
            .or. local%get_lval (var_str ("?hadronization_active")) &
            .or. local%get_lval (var_str ("?mlm_matching")) &
            .or. local%get_lval (var_str ("?ckkw_matching")) &
            .or. local%get_lval (var_str ("?muli_active"))
    if (enable_qcd) then 
       call dispatch_evt_shower (evt, local, process)
       if (associated (evt))  call entry%import_transform (evt)
    end if
  end subroutine entry_setup_event_transforms

\eatline
\nwendcode{}\nwbegindocs{2202}Compute weights.  The integral in the argument is the sum of integrals for
all processes in the sample.  After computing the process weights, we repeat
the normalization procedure for the process components.
\nwenddocs{}\nwbegincode{2203}\moddef{Simulations: entry: TBP}\plusendmoddef
  procedure :: init_mci_selector => entry_init_mci_selector
\nwendcode{}\nwbegincode{2204}\moddef{Simulations: procedures}\plusendmoddef
  subroutine entry_init_mci_selector (entry)
    class(entry_t), intent(inout) :: entry
    integer :: i
    if (entry%has_integral) then
       call entry%mci_selector%init (entry%mci_set%integral)
       do i = 1, entry%n_mci
          entry%mci_set(i)%weight_mci = entry%mci_selector%get_weight (i)
       end do
    end if
  end subroutine entry_init_mci_selector
  
\eatline
\nwendcode{}\nwbegindocs{2205}Select a MCI entry, using the embedded random-number generator.
\nwenddocs{}\nwbegincode{2206}\moddef{Simulations: entry: TBP}\plusendmoddef
  procedure :: select_mci => entry_select_mci
\nwendcode{}\nwbegincode{2207}\moddef{Simulations: procedures}\plusendmoddef
  function entry_select_mci (entry) result (i_mci)
    class(entry_t), intent(inout) :: entry
    integer :: i_mci
    call entry%mci_selector%generate (entry%rng, i_mci)
  end function entry_select_mci
  
\eatline
\nwendcode{}\nwbegindocs{2208}Record an event for this entry, i.e., increment the appropriate counters.
\nwenddocs{}\nwbegincode{2209}\moddef{Simulations: entry: TBP}\plusendmoddef
  procedure :: record => entry_record
\nwendcode{}\nwbegincode{2210}\moddef{Simulations: procedures}\plusendmoddef
  subroutine entry_record (entry, i_mci, from_file)
    class(entry_t), intent(inout) :: entry
    integer, intent(in) :: i_mci
    logical, intent(in), optional :: from_file
    real(default) :: weight, excess
    weight = entry%get_weight_prc ()
    excess = entry%get_excess_prc ()
    call entry%counter%record (weight, excess, from_file)
    call entry%mci_set(i_mci)%counter%record (weight, excess)
  end subroutine entry_record
    
\eatline
\nwendcode{}\nwbegindocs{2211}Update and restore the process core that this entry accesses, when
parameters change.  If explicit arguments {\tt{}model}, {\tt{}qcd}, or
{\tt{}helicity{\char95}selection} are provided, use those.  Otherwise use the
parameters stored in the process object.
\nwenddocs{}\nwbegincode{2212}\moddef{Simulations: entry: TBP}\plusendmoddef
  procedure :: update_process => entry_update_process
  procedure :: restore_process => entry_restore_process
\nwendcode{}\nwbegincode{2213}\moddef{Simulations: procedures}\plusendmoddef
  subroutine entry_update_process (entry, model, qcd, helicity_selection)
    class(entry_t), intent(inout) :: entry
    class(model_data_t), intent(in), optional, target :: model
    type(qcd_t), intent(in), optional :: qcd
    type(helicity_selection_t), intent(in), optional :: helicity_selection
    type(process_t), pointer :: process
    class(prc_core_t), allocatable :: core
    integer :: i, n_components
    class(model_data_t), pointer :: model_local
    type(qcd_t) :: qcd_local
    if (present (model)) then
       model_local => model
    else
       model_local => entry%model
    end if
    if (present (qcd)) then
       qcd_local = qcd
    else
       qcd_local = entry%qcd
    end if
    process => entry%get_process_ptr ()
    n_components = process%get_n_components ()
    allocate (entry%core_safe (n_components))
    do i = 1, n_components
       if (process%has_matrix_element (i)) then
          call process%extract_component_core (i, core)
          call dispatch_core_update (core, &
               model_local, helicity_selection, qcd_local, &
               entry%core_safe(i)%core)
          call process%restore_component_core (i, core)
       end if
    end do
  end subroutine entry_update_process
  
  subroutine entry_restore_process (entry)
    class(entry_t), intent(inout) :: entry
    type(process_t), pointer :: process
    class(prc_core_t), allocatable :: core
    integer :: i, n_components
    process => entry%get_process_ptr ()
    n_components = process%get_n_components ()
    do i = 1, n_components
       if (process%has_matrix_element (i)) then
          call process%extract_component_core (i, core)
          call dispatch_core_restore (core, entry%core_safe(i)%core)
          call process%restore_component_core (i, core)
       end if
    end do
    deallocate (entry%core_safe)
  end subroutine entry_restore_process
  
\eatline
\eatline
\nwendcode{}\nwbegincode{2214}\moddef{Simulations: entry: TBP}\plusendmoddef
  procedure :: combine_mci_sets => entry_combine_mci_sets
\nwendcode{}\nwbegincode{2215}\moddef{Simulations: procedures}\plusendmoddef
  subroutine entry_combine_mci_sets (entry)
    class(entry_t), intent(inout) :: entry
    integer :: n_components_lo, i_component, i_virt
    n_components_lo = entry%n_mci / 3
    do i_component = 1, n_components_lo
      i_virt = i_component + 2*n_components_lo
      entry%mci_set(i_component)%integral = &
        entry%mci_set(i_component)%integral + entry%mci_set(i_virt)%integral
      entry%mci_set(i_component)%error = sqrt (&
        entry%mci_set(i_component)%error**2 + entry%mci_set(i_virt)%error**2)
      entry%mci_set(i_virt)%integral = 0._default
      entry%mci_set(i_virt)%has_integral = .false.
    end do
  end subroutine entry_combine_mci_sets
 
\eatline
\nwendcode{}\nwbegindocs{2216}\nwdocspar
\subsection{Entries for alternative environment}
Entries for alternate environments.  [No additional components
anymore, so somewhat redundant.]
\nwenddocs{}\nwbegincode{2217}\moddef{Simulations: types}\plusendmoddef
  type, extends (entry_t) :: alt_entry_t
   contains
   \LA{}Simulations: alt entry: TBP\RA{}
  end type alt_entry_t
  
\eatline
\nwendcode{}\nwbegindocs{2218}The alternative entries are there to re-evaluate the event, given
momenta, in a different context.

Therefore, we allocate a local process object and use this as the
reference for the local process instance, when initializing the entry.
We temporarily import the {\tt{}process} object into an {\tt{}integration{\char95}t}
wrapper, to take advantage of the associated methods.  The local
process object is built in the context of the current environment,
here called {\tt{}global}.  Then, we initialize the process instance.

The {\tt{}master{\char95}process} object contains the integration results to which we
refer when recalculating an event.  Therefore, we use this object instead of
the locally built {\tt{}process} when we extract the integration results.

The locally built {\tt{}process} object should be finalized when done.  It
remains accessible via the {\tt{}event{\char95}t} base object of {\tt{}entry}, which
contains pointers to the process and instance.
\nwenddocs{}\nwbegincode{2219}\moddef{Simulations: alt entry: TBP}\endmoddef
  procedure :: init_alt => alt_entry_init
\nwendcode{}\nwbegincode{2220}\moddef{Simulations: procedures}\plusendmoddef
  subroutine alt_entry_init (entry, process_id, master_process, local)
    class(alt_entry_t), intent(inout), target :: entry
    type(string_t), intent(in) :: process_id
    type(process_t), intent(in), target :: master_process
    type(rt_data_t), intent(inout), target :: local
    class(rng_factory_t), allocatable :: rng_factory
    type(process_t), pointer :: process
    type(process_instance_t), pointer :: process_instance
    type(string_t) :: run_id
    type(integration_t) :: intg
    integer :: i

    call msg_message ("Simulate: initializing alternate process setup ...")

    run_id = &
         local%var_list%get_sval (var_str ("$run_id"))
    call local%set_log (var_str ("?rebuild_phase_space"), &
         .false., is_known = .true.)
    call local%set_log (var_str ("?check_phs_file"), &
         .false., is_known = .true.)
    call local%set_log (var_str ("?rebuild_grids"), &
         .false., is_known = .true.)
    
    call entry%basic_init (local%var_list)
    
    call prepare_local_process (process, process_id, local)
    entry%process_id = process_id
    entry%run_id = run_id

    call entry%import_process_characteristics (process)
    
    allocate (entry%mci_set (entry%n_mci))
    do i = 1, size (entry%mci_set)
       call entry%mci_set(i)%init (i, master_process)
    end do

    call entry%import_process_results (master_process)
    call entry%prepare_expressions (local)

    call prepare_process_instance (process_instance, process, local%model)
    call entry%setup_event_transforms (process, local)

    call entry%connect (process_instance, local%model, local%process_stack)
    call entry%setup_expressions ()

    entry%model => process%get_model_ptr ()

    call msg_message ("...  alternate process setup complete.")

  end subroutine alt_entry_init

\eatline
\nwendcode{}\nwbegindocs{2221}Copy the particle set from the master entry to the alternate entry.
This is the particle set of the hard process.
\nwenddocs{}\nwbegincode{2222}\moddef{Simulations: alt entry: TBP}\plusendmoddef
  procedure :: fill_particle_set => entry_fill_particle_set
\nwendcode{}\nwbegincode{2223}\moddef{Simulations: procedures}\plusendmoddef
  subroutine entry_fill_particle_set (alt_entry, entry)
    class(alt_entry_t), intent(inout) :: alt_entry
    class(entry_t), intent(in), target :: entry
    type(particle_set_t) :: pset
    call entry%get_hard_particle_set (pset)
    call alt_entry%set_hard_particle_set (pset)
    call pset%final ()
  end subroutine entry_fill_particle_set
    
\eatline
\nwendcode{}\nwbegindocs{2224}\nwdocspar
\subsection{The simulation type}
Each simulation object corresponds to an event sample, identified by
the {\tt{}sample{\char95}id}.

The simulation may cover several processes simultaneously.  All
process-specific data, including the event records, are stored in the
{\tt{}entry} subobjects.  The {\tt{}current} index indicates which record
was selected last. {\tt{}version} is foreseen to contain a tag on the \whizard\
event file version. It can be 
\nwenddocs{}\nwbegincode{2225}\moddef{Simulations: public}\endmoddef
  public :: simulation_t
\nwendcode{}\nwbegincode{2226}\moddef{Simulations: types}\plusendmoddef
  type :: simulation_t
     private
     type(string_t) :: sample_id
     logical :: unweighted = .true.
     logical :: negative_weights = .false.
     logical :: respect_selection = .true.
     integer :: norm_mode = NORM_UNDEFINED
     logical :: update_sqme = .false.
     logical :: update_weight = .false.
     logical :: update_event = .false.
     logical :: recover_beams = .false.
     logical :: pacify = .false.
     integer :: n_max_tries = 10000
     integer :: n_prc = 0
     integer :: n_alt = 0
     logical :: has_integral = .false.
     logical :: valid
     real(default) :: integral = 0
     real(default) :: error = 0
     integer :: version = 1
     character(32) :: md5sum_prc = ""
     character(32) :: md5sum_cfg = ""
     character(32), dimension(:), allocatable :: md5sum_alt
     type(entry_t), dimension(:), allocatable :: entry
     type(alt_entry_t), dimension(:,:), allocatable :: alt_entry
     type(selector_t) :: process_selector
     integer :: n_evt_requested = 0
     integer :: split_n_evt = 0
     integer :: split_index = 0
     type(counter_t) :: counter
     class(rng_t), allocatable :: rng
     integer :: i_prc = 0
     integer :: i_mci = 0
     real(default) :: weight = 0
     real(default) :: excess = 0
   contains
   \LA{}Simulations: simulation: TBP\RA{}
  end type simulation_t
  
\eatline
\nwendcode{}\nwbegindocs{2227}Output.  {\tt{}write{\char95}config} writes just the configuration.  {\tt{}write}
as a method of the base type {\tt{}event{\char95}t}
writes the current event and process instance, depending on options.
\nwenddocs{}\nwbegincode{2228}\moddef{Simulations: simulation: TBP}\endmoddef
  procedure :: write => simulation_write
\nwendcode{}\nwbegincode{2229}\moddef{Simulations: procedures}\plusendmoddef
  subroutine simulation_write (object, unit)
    class(simulation_t), intent(in) :: object
    integer, intent(in), optional :: unit
    integer :: u, i
    u = given_output_unit (unit)
    call write_separator (u, 2)
    write (u, "(1x,A,A,A)")  "Event sample: '", char (object%sample_id), "'"
    write (u, "(3x,A,I0)")  "Processes    = ", object%n_prc
    if (object%n_alt > 0) then
       write (u, "(3x,A,I0)")  "Alt.wgts     = ", object%n_alt
    end if
    write (u, "(3x,A,L1)")  "Unweighted   = ", object%unweighted
    write (u, "(3x,A,A)")   "Event norm   = ", &
         char (event_normalization_string (object%norm_mode))
    write (u, "(3x,A,L1)")  "Neg. weights = ", object%negative_weights
    write (u, "(3x,A,L1)")  "Respect sel. = ", object%respect_selection
    write (u, "(3x,A,L1)")  "Update sqme  = ", object%update_sqme
    write (u, "(3x,A,L1)")  "Update wgt   = ", object%update_weight
    write (u, "(3x,A,L1)")  "Update event = ", object%update_event
    write (u, "(3x,A,L1)")  "Recov. beams = ", object%recover_beams
    write (u, "(3x,A,L1)")  "Pacify       = ", object%pacify
    write (u, "(3x,A,I0)")  "Max. tries   = ", object%n_max_tries
    if (object%has_integral) then
       write (u, "(3x,A," // FMT_19 // ")")  "Integral     = ", object%integral
       write (u, "(3x,A," // FMT_19 // ")")  "Error        = ", object%error
    else
       write (u, "(3x,A)")  "Integral     = [undefined]"
    end if
    write (u, "(3x,A,L1)")  "Sim. valid   = ", object%valid
    write (u, "(3x,A,I0)")  "Ev.file ver. = ", object%version
    if (object%md5sum_prc /= "") then
       write (u, "(3x,A,A,A)")  "MD5 sum (proc)   = '", object%md5sum_prc, "'"
    end if
    if (object%md5sum_cfg /= "") then
       write (u, "(3x,A,A,A)")  "MD5 sum (config) = '", object%md5sum_cfg, "'"
    end if
    write (u, "(3x,A,I0)")  "Events requested  = ", object%n_evt_requested
    if (object%split_n_evt > 0) then
       write (u, "(3x,A,I0)")  "Events per file   = ", object%split_n_evt
       write (u, "(3x,A,I0)")  "First file index  = ", object%split_index
    end if
    call object%counter%write (u)
    call write_separator (u)
    if (object%i_prc /= 0) then
       write (u, "(1x,A)")  "Current event:"
       write (u, "(3x,A,I0,A,A)")  "Process #", &
            object%i_prc, ": ", &
            char (object%entry(object%i_prc)%process_id)
       write (u, "(3x,A,I0)")  "MCI set #", object%i_mci
       write (u, "(3x,A," // FMT_19 // ")")  "Weight    = ", object%weight
       if (object%excess /= 0) &
            write (u, "(3x,A," // FMT_19 // ")")  "Excess    = ", object%excess
    else
       write (u, "(1x,A,I0,A,A)")  "Current event: [undefined]"
    end if
    call write_separator (u)
    if (allocated (object%rng)) then
       call object%rng%write (u)
    else
       write (u, "(3x,A)")  "Random-number generator: [undefined]"
    end if
    if (allocated (object%entry)) then
       do i = 1, size (object%entry)
          if (i == 1) then
             call write_separator (u, 2)
          else
             call write_separator (u)
          end if
          write (u, "(1x,A,I0,A)") "Process #", i, ":"
          call object%entry(i)%write_config (u)
       end do
    end if
    call write_separator (u, 2)
  end subroutine simulation_write
  
\eatline
\nwendcode{}\nwbegindocs{2230}Write the current event record.  If an explicit index is given,
write that event record.

We implement writing to {\tt{}unit} (event contents / debugging format)
and writing to an {\tt{}eio} event stream (storage). We include a {\tt{}testflag}
in order to suppress numerical noise in the testsuite.
\nwenddocs{}\nwbegincode{2231}\moddef{Simulations: simulation: TBP}\plusendmoddef
  generic :: write_event => write_event_unit
  procedure :: write_event_unit => simulation_write_event_unit
\nwendcode{}\nwbegincode{2232}\moddef{Simulations: procedures}\plusendmoddef
  subroutine simulation_write_event_unit &
       (object, unit, i_prc, verbose, testflag)
    class(simulation_t), intent(in) :: object
    integer, intent(in), optional :: unit
    logical, intent(in), optional :: verbose
    integer, intent(in), optional :: i_prc
    logical, intent(in), optional :: testflag
    logical :: pacified
    integer :: current
    pacified = .false.
    if (present(testflag)) pacified = testflag
    pacified = pacified .or. object%pacify
    if (present (i_prc)) then
       current = i_prc
    else
       current = object%i_prc
    end if
    if (current > 0) then
       call object%entry(current)%write (unit, verbose = verbose, &
            testflag = pacified)
    else
       call msg_fatal ("Simulation: write event: no process selected")
    end if
  end subroutine simulation_write_event_unit

\eatline
\nwendcode{}\nwbegindocs{2233}This writes one of the alternate events, if allocated.
\nwenddocs{}\nwbegincode{2234}\moddef{Simulations: simulation: TBP}\plusendmoddef
  procedure :: write_alt_event => simulation_write_alt_event
\nwendcode{}\nwbegincode{2235}\moddef{Simulations: procedures}\plusendmoddef
  subroutine simulation_write_alt_event (object, unit, j_alt, i_prc, &
       verbose, testflag)
    class(simulation_t), intent(in) :: object
    integer, intent(in), optional :: unit
    integer, intent(in), optional :: j_alt
    integer, intent(in), optional :: i_prc
    logical, intent(in), optional :: verbose
    logical, intent(in), optional :: testflag
    integer :: i, j
    if (present (j_alt)) then
       j = j_alt
    else
       j = 1
    end if
    if (present (i_prc)) then
       i = i_prc
    else
       i = object%i_prc
    end if
    if (i > 0) then
       if (j> 0 .and. j <= object%n_alt) then
          call object%alt_entry(i,j)%write (unit, verbose = verbose, &
               testflag = testflag)
       else
          call msg_fatal ("Simulation: write alternate event: out of range")
       end if
    else
       call msg_fatal ("Simulation: write alternate event: no process selected")
    end if
  end subroutine simulation_write_alt_event

\eatline
\nwendcode{}\nwbegindocs{2236}Finalizer.
\nwenddocs{}\nwbegincode{2237}\moddef{Simulations: simulation: TBP}\plusendmoddef
  procedure :: final => simulation_final
\nwendcode{}\nwbegincode{2238}\moddef{Simulations: procedures}\plusendmoddef
  subroutine simulation_final (object)
    class(simulation_t), intent(inout) :: object
    integer :: i, j
    if (allocated (object%entry)) then
       do i = 1, size (object%entry)
          call object%entry(i)%final ()
       end do
    end if
    if (allocated (object%alt_entry)) then
       do j = 1, size (object%alt_entry, 2)
          do i = 1, size (object%alt_entry, 1)
             call object%alt_entry(i,j)%final ()
          end do
       end do
    end if
    if (allocated (object%rng))  call object%rng%final ()
  end subroutine simulation_final
  
\eatline
\nwendcode{}\nwbegindocs{2239}Initialization.  We can deduce all data from the given list of
process IDs and the global data set.  The process objects are taken
from the stack.  Once the individual integrals are known, we add them (and the
errors), to get the sample integral.

If there are alternative environments, we suspend initialization for
setting up alternative process objects, then restore the master
process and its parameters.  The generator or rescanner can then
switch rapidly between processes.

If {\tt{}integrate} is set, we make sure that all affected processes are
integrated before simulation.  This is necessary if we want to actually
generate events.  If {\tt{}integrate} is unset, we don't need the integral
because we just rescan existing events.  In that case, we just need compiled
matrix elements.

If {\tt{}generate} is set, we prepare for actually generating events.  Otherwise,
we may only read and rescan events.
\nwenddocs{}\nwbegincode{2240}\moddef{Simulations: simulation: TBP}\plusendmoddef
  procedure :: init => simulation_init
\nwendcode{}\nwbegincode{2241}\moddef{Simulations: procedures}\plusendmoddef
  subroutine simulation_init (simulation, &
       process_id, integrate, generate, local, global, alt_env)
    class(simulation_t), intent(out), target :: simulation
    type(string_t), dimension(:), intent(in) :: process_id
    logical, intent(in) :: integrate, generate
    type(rt_data_t), intent(inout), target :: local
    type(rt_data_t), intent(inout), optional, target :: global
    type(rt_data_t), dimension(:), intent(inout), optional, target :: alt_env
    class(rng_factory_t), allocatable :: rng_factory
    type(string_t) :: norm_string, version_string
    integer :: i, j
    simulation%sample_id = &
         local%get_sval (var_str ("$sample"))
    simulation%unweighted = &
         local%get_lval (var_str ("?unweighted"))
    simulation%negative_weights = &
         local%get_lval (var_str ("?negative_weights"))
    simulation%respect_selection = &
         local%get_lval (var_str ("?sample_select"))
    version_string = &
         local%get_sval (var_str ("$event_file_version"))
    norm_string = &
         local%get_sval (var_str ("$sample_normalization"))
    simulation%norm_mode = &
         event_normalization_mode (norm_string, simulation%unweighted)
    simulation%pacify = &
         local%get_lval (var_str ("?sample_pacify"))
    simulation%n_max_tries = &
         local%get_ival (var_str ("sample_max_tries"))
    simulation%split_n_evt = &
         local%get_ival (var_str ("sample_split_n_evt"))
    simulation%split_index = &
         local%get_ival (var_str ("sample_split_index"))
    simulation%update_sqme = &
         local%get_lval (var_str ("?update_sqme"))
    simulation%update_weight = &
         local%get_lval (var_str ("?update_weight"))
    simulation%update_event = &
         local%get_lval (var_str ("?update_event"))
    simulation%recover_beams = &
         local%get_lval (var_str ("?recover_beams"))
    select case (size (process_id))
    case (0)
       call msg_error ("Simulation: no process selected")
    case (1)
       write (msg_buffer, "(A,A,A)") &
            "Starting simulation for process '", &
            char (process_id(1)), "'"
       call msg_message ()
    case default
       write (msg_buffer, "(A,A,A)") &
            "Starting simulation for processes '", &
            char (process_id(1)), "' etc."
       call msg_message ()
    end select
    select case (char (version_string))
    case ("", "2.2.4")
       simulation%version = 2
    case ("2.2")
       simulation%version = 1
    case default
       simulation%version = 0
    end select
    if (simulation%version == 0) then
       call msg_fatal ("Event file format '" &
            // char (version_string) &
            // "' is not compatible with this version.")
    end if        
    simulation%n_prc = size (process_id)
    allocate (simulation%entry (simulation%n_prc))
    if (present (alt_env)) then
       simulation%n_alt = size (alt_env)
       do i = 1, simulation%n_prc
          call simulation%entry(i)%init (process_id(i), &
               integrate, generate, &
               simulation%update_sqme, &
               local, global, simulation%n_alt)
          if (signal_is_pending ())  return
       end do
       if (.not. any (simulation%entry%valid)) then
          call msg_error ("Simulate: no process has a valid matrix element.")
          simulation%valid = .false.
          return
       end if
       call simulation%update_processes ()
       allocate (simulation%alt_entry (simulation%n_prc, simulation%n_alt))
       allocate (simulation%md5sum_alt (simulation%n_alt))
       simulation%md5sum_alt = ""
       do j = 1, simulation%n_alt
          do i = 1, simulation%n_prc
             call simulation%alt_entry(i,j)%init_alt (process_id(i), &
                  simulation%entry(i)%get_process_ptr (), alt_env(j))
             if (signal_is_pending ())  return
          end do
       end do
       call simulation%restore_processes ()
    else       
       do i = 1, simulation%n_prc
          call simulation%entry(i)%init &
               (process_id(i), &
               integrate, generate, simulation%update_sqme, &
               local, global)
          if (signal_is_pending ())  return          
       end do
       if (.not. any (simulation%entry%valid)) then
          call msg_error ("Simulate: " &
               // "no process has a valid matrix element.") 
          simulation%valid = .false.
          return
       end if
    end if
    call dispatch_rng_factory (rng_factory, local)
    call rng_factory%make (simulation%rng)
    if (all (simulation%entry%has_integral)) then
       simulation%integral = sum (simulation%entry%integral)
       simulation%error = sqrt (sum (simulation%entry%error ** 2))
       simulation%has_integral = .true.
       if (integrate .and. generate) then
          do i = 1, simulation%n_prc
             if (simulation%entry(i)%integral < 0 .and. .not. &
                  simulation%negative_weights) then
                call msg_fatal ("Integral of process '" // &
                     char (process_id (i)) // "'is negative.")
             end if
          end do
       end if
    else
       if (integrate .and. generate) &
            call msg_error ("Simulation contains undefined integrals.")
    end if
    if (simulation%integral > 0 .or. &
         (simulation%integral < 0 .and. simulation%negative_weights)) then
       simulation%valid = .true.
    else if (generate) then
       call msg_error ("Simulate: " &
            // "sum of process integrals must be positive; skipping.")
       simulation%valid = .false.
       return
    end if
    if (simulation%valid)  call simulation%compute_md5sum ()
  end subroutine simulation_init

\eatline
\nwendcode{}\nwbegindocs{2242}The number of events that we want to simulate is determined by the
settings of {\tt{}n{\char95}events}, {\tt{}luminosity}, and {\tt{}?unweighted}.  For
weighted events, we take {\tt{}n{\char95}events} at face value as the number of
matrix element calls.  For unweighted events, if the process is a
decay, {\tt{}n{\char95}events} is the number of unweighted events.  In these
cases, the luminosity setting is ignored.  

For unweighted events with a scattering process, we calculate the
event number that corresponds to the luminosity, given the current
value of the integral.  We then compare this with {\tt{}n{\char95}events} and
choose the larger number.
\nwenddocs{}\nwbegincode{2243}\moddef{Simulations: simulation: TBP}\plusendmoddef
  procedure :: compute_n_events => simulation_compute_n_events
\nwendcode{}\nwbegincode{2244}\moddef{Simulations: procedures}\plusendmoddef
  subroutine simulation_compute_n_events (simulation, n_events, var_list)
    class(simulation_t), intent(in) :: simulation
    integer, intent(out) :: n_events
    type(var_list_t) :: var_list
    real(default) :: lumi, x_events_lumi
    integer :: n_events_lumi
    logical :: is_scattering
    n_events = &
         var_list%get_ival (var_str ("n_events"))
    lumi = &
         var_list%get_rval (var_str ("luminosity"))
    if (simulation%unweighted) then
       is_scattering = simulation%entry(1)%n_in == 2
       if (is_scattering) then
          x_events_lumi = abs (simulation%integral * lumi)
          if (x_events_lumi < huge (n_events)) then
             n_events_lumi = nint (x_events_lumi)
          else
             call msg_message ("Simulation: luminosity too large, &
                  &limiting number of events")
             n_events_lumi = huge (n_events)
          end if
          if (n_events_lumi > n_events) then
             call msg_message ("Simulation: using n_events as computed from &
                  &luminosity value")
             n_events = n_events_lumi
          else
             write (msg_buffer, "(A,1x,I0)") &
                  "Simulation: requested number of events =", n_events
             call msg_message ()
             if (simulation%integral /= 0) then
                write (msg_buffer, "(A,1x,ES11.4)") &
                     "            corr. to luminosity [fb-1] = ", &
                     n_events / simulation%integral        
                call msg_message ()
             end if
          end if
       end if
    end if
  end subroutine simulation_compute_n_events

\eatline
\nwendcode{}\nwbegindocs{2245}Compute the checksum of the process set.  We retrieve the MD5 sums
of all processes.  This depends only on the process definitions, while
parameters are not considered.  The configuration checksum is
retrieved from the MCI records in the process objects and furthermore
includes beams, parameters, integration results, etc., so matching the
latter should guarantee identical physics.
\nwenddocs{}\nwbegincode{2246}\moddef{Simulations: simulation: TBP}\plusendmoddef
  procedure :: compute_md5sum => simulation_compute_md5sum
\nwendcode{}\nwbegincode{2247}\moddef{Simulations: procedures}\plusendmoddef
  subroutine simulation_compute_md5sum (simulation)
    class(simulation_t), intent(inout) :: simulation
    type(process_t), pointer :: process
    type(string_t) :: buffer
    integer :: j, i, n_mci, i_mci, n_component, i_component
    if (simulation%md5sum_prc == "") then
       buffer = ""
       do i = 1, simulation%n_prc
          if (.not. simulation%entry(i)%valid) cycle
          process => simulation%entry(i)%get_process_ptr ()
          n_component = process%get_n_components ()
          do i_component = 1, n_component
             if (process%has_matrix_element (i_component)) then
                buffer = buffer // process%get_md5sum_prc (i_component)
             end if
          end do
       end do
       simulation%md5sum_prc = md5sum (char (buffer))
    end if
    if (simulation%md5sum_cfg == "") then
       buffer = ""
       do i = 1, simulation%n_prc
          if (.not. simulation%entry(i)%valid) cycle          
          process => simulation%entry(i)%get_process_ptr ()
          n_mci = process%get_n_mci ()
          do i_mci = 1, n_mci
             buffer = buffer // process%get_md5sum_mci (i_mci)
          end do
       end do
       simulation%md5sum_cfg = md5sum (char (buffer))
    end if
    do j = 1, simulation%n_alt
       if (simulation%md5sum_alt(j) == "") then
          buffer = ""
          do i = 1, simulation%n_prc
             process => simulation%alt_entry(i,j)%get_process_ptr ()
             buffer = buffer // process%get_md5sum_cfg ()
          end do
          simulation%md5sum_alt(j) = md5sum (char (buffer))
       end if
    end do
  end subroutine simulation_compute_md5sum

\eatline
\nwendcode{}\nwbegindocs{2248}Initialize the process selector, using the entry integrals as process
weights.
\nwenddocs{}\nwbegincode{2249}\moddef{Simulations: simulation: TBP}\plusendmoddef
  procedure :: init_process_selector => simulation_init_process_selector
\nwendcode{}\nwbegincode{2250}\moddef{Simulations: procedures}\plusendmoddef
  subroutine simulation_init_process_selector (simulation)
    class(simulation_t), intent(inout) :: simulation
    integer :: i
    if (simulation%has_integral) then
       call simulation%process_selector%init (simulation%entry%integral)
       do i = 1, simulation%n_prc
          associate (entry => simulation%entry(i))
            if (.not. entry%valid) then
               call msg_warning ("Process '" // char (entry%process_id) // &
                    "': matrix element vanishes, no events can be generated.")
               cycle
            end if
            call entry%init_mci_selector ()
            entry%process_weight = simulation%process_selector%get_weight (i)
          end associate
       end do
    end if
  end subroutine simulation_init_process_selector
    
\eatline
\nwendcode{}\nwbegindocs{2251}Select a process, using the random-number generator.
\nwenddocs{}\nwbegincode{2252}\moddef{Simulations: simulation: TBP}\plusendmoddef
  procedure :: select_prc => simulation_select_prc
\nwendcode{}\nwbegincode{2253}\moddef{Simulations: procedures}\plusendmoddef
  function simulation_select_prc (simulation) result (i_prc)
    class(simulation_t), intent(inout) :: simulation
    integer :: i_prc
    call simulation%process_selector%generate (simulation%rng, i_prc)
  end function simulation_select_prc

\eatline
\nwendcode{}\nwbegindocs{2254}Select a MCI set for the selected process.
\nwenddocs{}\nwbegincode{2255}\moddef{Simulations: simulation: TBP}\plusendmoddef
  procedure :: select_mci => simulation_select_mci
\nwendcode{}\nwbegincode{2256}\moddef{Simulations: procedures}\plusendmoddef
  function simulation_select_mci (simulation) result (i_mci)
    class(simulation_t), intent(inout) :: simulation
    integer :: i_mci
    if (simulation%i_prc /= 0) then
       i_mci = simulation%entry(simulation%i_prc)%select_mci ()
    end if
  end function simulation_select_mci

\eatline
\nwendcode{}\nwbegindocs{2257}Generate a predefined number of events.  First select a process and
a component set, then generate an event for that process and factorize
the quantum state.  The pair of random numbers can be used for
factorization.

When generating events, we drop all configurations where the event is
marked as incomplete.  This happens if the event fails cuts.  In fact,
such events are dropped already by the sampler if unweighting is in
effect, so this can happen only for weighted events.  By setting a
limit given by {\tt{}sample{\char95}max{\char95}tries} (user parameter), we can avoid an
endless loop.

NB: When reading from file, event transforms can't be applied because the
process instance will not be complete.  This should be fixed.
\nwenddocs{}\nwbegincode{2258}\moddef{Simulations: simulation: TBP}\plusendmoddef
  procedure :: generate => simulation_generate
\nwendcode{}\nwbegincode{2259}\moddef{Simulations: procedures}\plusendmoddef
  subroutine simulation_generate (simulation, n, es_array)
    class(simulation_t), intent(inout) :: simulation
    integer, intent(in) :: n
    type(event_stream_array_t), intent(inout), optional :: es_array
    type(string_t) :: str1, str2, str3
    logical :: generate_new, passed
    integer :: i, j
    simulation%n_evt_requested = n
    call simulation%entry%set_n (n)
    if (simulation%n_alt > 0)  call simulation%alt_entry%set_n (n)
    str1 = "Events: generating"
    if (present (es_array)) then
       if (es_array%has_input ())  str1 = "Events: reading"
    end if
    if (simulation%entry(1)%config%unweighted) then
       str2 = "unweighted"
    else
       str2 = "weighted"
    end if
    if (simulation%entry(1)%config%factorization_mode == &
         FM_IGNORE_HELICITY) then
       str3 = ", unpolarized"
    else 
       str3 = ", polarized"
    end if    
    write (msg_buffer, "(A,1x,I0,1x,A,1x,A)")  char (str1), n, &
         char (str2) // char(str3), "events ..."
    call msg_message ()
    write (msg_buffer, "(A,1x,A)") "Events: event normalization mode", &
         char (event_normalization_string (simulation%norm_mode))
    call msg_message ()
    do i = 1, n
       if (present (es_array)) then
          call simulation%read_event (es_array, .true., generate_new)
       else
          generate_new = .true.
       end if
       if (generate_new) then
          simulation%i_prc = simulation%select_prc ()
          simulation%i_mci = simulation%select_mci ()
          associate (entry => simulation%entry(simulation%i_prc))
            do j = 1, simulation%n_max_tries
               if (.not. entry%valid)  call msg_warning &
                       ("Process '" // char (entry%process_id) // "': " // &
                       "matrix element vanishes, no events can be generated.")
               call entry%generate (simulation%i_mci)
               if (signal_is_pending ()) return
               if (entry%has_valid_particle_set ())  exit
            end do
            if (.not. entry%has_valid_particle_set ()) then
               write (msg_buffer, "(A,I0,A)")  "Simulation: failed to &
                    &generate valid event after ", &
                    simulation%n_max_tries, " tries (sample_max_tries)"
               call msg_fatal ()
            end if
            call entry%evaluate_expressions ()
            if (signal_is_pending ()) return
            if (entry%passed_selection ()) then
               simulation%weight = entry%get_weight_ref ()
               simulation%excess = entry%get_excess_prc ()
            end if
            call simulation%counter%record &
                 (simulation%weight, simulation%excess)
            call entry%record (simulation%i_mci)
          end associate
       else
          associate (entry => simulation%entry(simulation%i_prc))
            call entry%accept_sqme_ref ()
            call entry%accept_weight_ref ()
            call entry%check ()
            call entry%evaluate_expressions ()
            if (signal_is_pending ()) return
            if (entry%passed_selection ()) then
               simulation%weight = entry%get_weight_ref ()
               simulation%excess = entry%get_excess_prc ()
            end if
            call simulation%counter%record &
                 (simulation%weight, simulation%excess, from_file=.true.)
            call entry%record (simulation%i_mci, from_file=.true.)
          end associate
       end if
       call simulation%calculate_alt_entries ()
       if (signal_is_pending ()) return
       if (simulation%pacify)  call pacify (simulation)
       if (simulation%respect_selection) then
          passed = simulation%entry(simulation%i_prc)%passed_selection ()
       else
          passed = .true.
       end if
       if (present (es_array)) then
          call simulation%write_event (es_array, passed)
       end if
    end do
    call msg_message ("        ... event sample complete.")
    call simulation%counter%show_excess ()
  end subroutine simulation_generate
  
\eatline
\nwendcode{}\nwbegindocs{2260}Compute the event matrix element and weight for all alternative
environments, given the current event and selected process.  We first
copy the particle set, then temporarily update the process core with
local parameters, recalculate everything, and restore the process
core.

The event weight is obtained by rescaling the original event weight with the
ratio of the new and old {\tt{}sqme} values.  (In particular, if the old
value was zero, the weight will stay zero.)

Note: this may turn out to be inefficient because we always replace
all parameters and recalculate everything, once for each event and
environment.  However, a more fine-grained control requires more
code.  In any case, while we may keep multiple process cores (which
stay constant for a simulation run), we still have to update the
external matrix element parameters event by event.  The matrix element
``object'' is present only once.
\nwenddocs{}\nwbegincode{2261}\moddef{Simulations: simulation: TBP}\plusendmoddef
  procedure :: calculate_alt_entries => simulation_calculate_alt_entries
\nwendcode{}\nwbegincode{2262}\moddef{Simulations: procedures}\plusendmoddef
  subroutine simulation_calculate_alt_entries (simulation)
    class(simulation_t), intent(inout) :: simulation
    real(default) :: factor
    real(default), dimension(:), allocatable :: sqme_alt, weight_alt
    integer :: n_alt, i, j
    i = simulation%i_prc
    n_alt = simulation%n_alt
    if (n_alt == 0)  return
    allocate (sqme_alt (n_alt), weight_alt (n_alt))
    associate (entry => simulation%entry(i))
      do j = 1, n_alt
         if (signal_is_pending ())  return
         factor = entry%get_kinematical_weight ()
         associate (alt_entry => simulation%alt_entry(i,j))
           call alt_entry%update_process ()
           call alt_entry%select &
                (entry%get_i_mci (), entry%get_i_term (), entry%get_channel ())
           call alt_entry%fill_particle_set (entry)
           call alt_entry%recalculate &
                (update_sqme = .true., weight_factor = factor)
           if (signal_is_pending ())  return
           call alt_entry%accept_sqme_prc ()
           call alt_entry%update_normalization ()
           call alt_entry%accept_weight_prc ()
           call alt_entry%check ()
           call alt_entry%evaluate_expressions ()
           if (signal_is_pending ())  return
           call alt_entry%restore_process ()
           sqme_alt(j) = alt_entry%get_sqme_ref ()
           if (alt_entry%passed_selection ()) then
              weight_alt(j) = alt_entry%get_weight_ref ()
           end if
         end associate
      end do
      call entry%set (sqme_alt = sqme_alt, weight_alt = weight_alt)
      call entry%check ()
      call entry%store_alt_values ()
    end associate
  end subroutine simulation_calculate_alt_entries
       
\eatline
\nwendcode{}\nwbegindocs{2263}Rescan an undefined number of events.

If {\tt{}update{\char95}event} or {\tt{}update{\char95}sqme} is set, we have to recalculate the
event, starting from the particle set.  If the latter is set, this includes
the squared matrix element (i.e., the amplitude is evaluated).  Otherwise,
only kinematics and observables derived from it are recovered.

If any of the update flags is set, we will come up with separate
{\tt{}sqme{\char95}prc} and {\tt{}weight{\char95}prc} values.  (The latter is only distinct
if {\tt{}update{\char95}weight} is set.)  Otherwise, we accept the reference values.
\nwenddocs{}\nwbegincode{2264}\moddef{Simulations: simulation: TBP}\plusendmoddef
  procedure :: rescan => simulation_rescan
\nwendcode{}\nwbegincode{2265}\moddef{Simulations: procedures}\plusendmoddef
  subroutine simulation_rescan (simulation, n, es_array, global)
    class(simulation_t), intent(inout) :: simulation
    integer, intent(in) :: n
    type(event_stream_array_t), intent(inout) :: es_array
    type(rt_data_t), intent(inout) :: global
    type(qcd_t) :: qcd
    type(string_t) :: str1, str2, str3
    logical :: complete
    str1 = "Rescanning"
    if (simulation%entry(1)%config%unweighted) then
       str2 = "unweighted"
    else
       str2 = "weighted"
    end if
    simulation%n_evt_requested = n
    call simulation%entry%set_n (n)
    if (simulation%update_sqme .or. simulation%update_weight) then
       call dispatch_qcd (qcd, global)
       call simulation%update_processes &
            (global%model, qcd, global%get_helicity_selection ())
       str3 = "(process parameters updated) "
    else
       str3 = ""
    end if
    write (msg_buffer, "(A,1x,A,1x,A,A,A)")  char (str1), char (str2), &
         "events ", char (str3), "..."
    call msg_message ()
    do
       call simulation%read_event (es_array, .false., complete)
       if (complete)  exit
       if (simulation%update_event &
            .or. simulation%update_sqme &
            .or. simulation%update_weight) then
          call simulation%recalculate ()
          if (signal_is_pending ())  return
          associate (entry => simulation%entry(simulation%i_prc))
            call entry%update_normalization ()
            if (simulation%update_event) then
               call entry%evaluate_transforms ()
            end if
            call entry%check ()
            call entry%evaluate_expressions ()
            if (signal_is_pending ())  return
            simulation%weight = entry%get_weight_prc ()
            call simulation%counter%record (simulation%weight, from_file=.true.)
            call entry%record (simulation%i_mci, from_file=.true.)
          end associate
       else
          associate (entry => simulation%entry(simulation%i_prc))
            call entry%accept_sqme_ref ()
            call entry%accept_weight_ref ()
            call entry%check ()
            call entry%evaluate_expressions ()
            if (signal_is_pending ())  return
            simulation%weight = entry%get_weight_ref ()
            call simulation%counter%record (simulation%weight, from_file=.true.)
            call entry%record (simulation%i_mci, from_file=.true.)
          end associate
       end if
       call simulation%calculate_alt_entries ()
       if (signal_is_pending ())  return
       call simulation%write_event (es_array)
    end do
    if (simulation%update_sqme .or. simulation%update_weight) then
       call simulation%restore_processes ()
    end if
  end subroutine simulation_rescan
  
\eatline
\nwendcode{}\nwbegindocs{2266}These routines take care of temporary parameter redefinitions that
we want to take effect while recalculating the matrix elements.  We
extract the core(s) of the processes that we are simulating, apply the
changes, and make sure that the changes are actually used.  This is
the duty of {\tt{}dispatch{\char95}core{\char95}update}.  When done, we restore the
original versions using {\tt{}dispatch{\char95}core{\char95}restore}.
\nwenddocs{}\nwbegincode{2267}\moddef{Simulations: simulation: TBP}\plusendmoddef
  procedure :: update_processes => simulation_update_processes
  procedure :: restore_processes => simulation_restore_processes
\nwendcode{}\nwbegincode{2268}\moddef{Simulations: procedures}\plusendmoddef
  subroutine simulation_update_processes (simulation, &
       model, qcd, helicity_selection)
    class(simulation_t), intent(inout) :: simulation
    class(model_data_t), intent(in), optional, target :: model
    type(qcd_t), intent(in), optional :: qcd
    type(helicity_selection_t), intent(in), optional :: helicity_selection
    integer :: i
    do i = 1, simulation%n_prc
       call simulation%entry(i)%update_process (model, qcd, helicity_selection)
    end do
  end subroutine simulation_update_processes
  
  subroutine simulation_restore_processes (simulation)
    class(simulation_t), intent(inout) :: simulation
    integer :: i
    do i = 1, simulation%n_prc
       call simulation%entry(i)%restore_process ()
    end do
  end subroutine simulation_restore_processes
  
\eatline
\eatline
\nwendcode{}\nwbegindocs{2269}\nwdocspar
\subsection{Event Stream I/O} 
Write an event to a generic {\tt{}eio} event stream.  The process index
must be selected, or the current index must be available.
\nwenddocs{}\nwbegincode{2270}\moddef{Simulations: simulation: TBP}\plusendmoddef
  generic :: write_event => write_event_eio
  procedure :: write_event_eio => simulation_write_event_eio
\nwendcode{}\nwbegincode{2271}\moddef{Simulations: procedures}\plusendmoddef
  subroutine simulation_write_event_eio (object, eio, i_prc)
    class(simulation_t), intent(in) :: object
    class(eio_t), intent(inout) :: eio
    integer, intent(in), optional :: i_prc
    integer :: current
    if (present (i_prc)) then
       current = i_prc
    else
       current = object%i_prc
    end if
    if (current > 0) then
       if (object%split_n_evt > 0) then
          if (object%counter%total > 1 .and. &
               mod (object%counter%total, object%split_n_evt) == 1) then
             call eio%split_out ()
          end if
       end if
       call eio%output (object%entry(current)%event_t, current, pacify = object%pacify)
    else
       call msg_fatal ("Simulation: write event: no process selected")
    end if
  end subroutine simulation_write_event_eio

\eatline
\nwendcode{}\nwbegindocs{2272}\nwdocspar
Read an event from a generic {\tt{}eio} event stream.  The event stream element
must specify the process within the sample ({\tt{}i{\char95}prc}), the MC group for this
process ({\tt{}i{\char95}mci}), the selected term ({\tt{}i{\char95}term}), the selected MC
integration {\tt{}channel}, and the particle set of the event.

We may encounter EOF, which we indicate by storing 0 for the process index
{\tt{}i{\char95}prc}.  An I/O error will be reported, and we also abort reading.
\nwenddocs{}\nwbegincode{2273}\moddef{Simulations: simulation: TBP}\plusendmoddef
  generic :: read_event => read_event_eio
  procedure :: read_event_eio => simulation_read_event_eio
\nwendcode{}\nwbegincode{2274}\moddef{Simulations: procedures}\plusendmoddef
  subroutine simulation_read_event_eio (object, eio)
    class(simulation_t), intent(inout) :: object
    class(eio_t), intent(inout) :: eio
    integer :: iostat, current
    call eio%input_i_prc (current, iostat)
    select case (iostat)
    case (0)
       object%i_prc = current
       call eio%input_event (object%entry(current)%event_t, iostat)
    end select
    select case (iostat)
    case (:-1)
       object%i_prc = 0
       object%i_mci = 0
    case (1:)
       call msg_error ("Reading events: I/O error, aborting read")
       object%i_prc = 0
       object%i_mci = 0
    case default
       object%i_mci = object%entry(current)%get_i_mci ()
    end select
  end subroutine simulation_read_event_eio

\eatline
\nwendcode{}\nwbegindocs{2275}\nwdocspar
\subsection{Event Stream Array} 
Write an event using an array of event I/O streams.
The process index must be selected, or the current index must be
available.
\nwenddocs{}\nwbegincode{2276}\moddef{Simulations: simulation: TBP}\plusendmoddef
  generic :: write_event => write_event_es_array
  procedure :: write_event_es_array => simulation_write_event_es_array
\nwendcode{}\nwbegincode{2277}\moddef{Simulations: procedures}\plusendmoddef
  subroutine simulation_write_event_es_array (object, es_array, passed)
    class(simulation_t), intent(in) :: object
    class(event_stream_array_t), intent(inout) :: es_array
    logical, intent(in), optional :: passed
    integer :: i_prc, event_index
    i_prc = object%i_prc
    if (i_prc > 0) then
       event_index = object%counter%total
       call es_array%output (object%entry(i_prc)%event_t, i_prc, &
            event_index, passed = passed, pacify = object%pacify)
    else
       call msg_fatal ("Simulation: write event: no process selected")
    end if
  end subroutine simulation_write_event_es_array

\eatline
\nwendcode{}\nwbegindocs{2278}Read an event using an array of event I/O streams.  Reading is
successful if there is an input stream within the array, and if a
valid event can be read from that stream.  If there is a stream, but
EOF is passed when reading the first item, we switch the channel to
output and return failure but no error message, such that new events
can be appended to that stream.
\nwenddocs{}\nwbegincode{2279}\moddef{Simulations: simulation: TBP}\plusendmoddef
  generic :: read_event => read_event_es_array
  procedure :: read_event_es_array => simulation_read_event_es_array
\nwendcode{}\nwbegincode{2280}\moddef{Simulations: procedures}\plusendmoddef
  subroutine simulation_read_event_es_array (object, es_array, enable_switch, &
       fail)
    class(simulation_t), intent(inout) :: object
    class(event_stream_array_t), intent(inout) :: es_array
    logical, intent(in) :: enable_switch
    logical, intent(out) :: fail
    integer :: iostat, i_prc
    if (es_array%has_input ()) then
       fail = .false.
       call es_array%input_i_prc (i_prc, iostat)
       select case (iostat)
       case (0)
          object%i_prc = i_prc
          call es_array%input_event (object%entry(i_prc)%event_t, iostat)
       case (:-1)
          write (msg_buffer, "(A,1x,I0,1x,A)")  &
               "... event file terminates after", &
               object%counter%read, "events."
          call msg_message ()
          if (enable_switch) then
             call es_array%switch_inout ()
             write (msg_buffer, "(A,1x,I0,1x,A)")  &
                  "Generating remaining ", &
                  object%n_evt_requested - object%counter%read, "events ..."
             call msg_message ()
          end if
          fail = .true.
          return
       end select
       select case (iostat)
       case (0)
          object%i_mci = object%entry(i_prc)%get_i_mci ()
       case default
          write (msg_buffer, "(A,1x,I0,1x,A)")  &
               "Reading events: I/O error, aborting read after", &
               object%counter%read, "events."
          call msg_error ()
          object%i_prc = 0
          object%i_mci = 0
          fail = .true.
       end select
    else
       fail = .true.
    end if
  end subroutine simulation_read_event_es_array

\eatline
\nwendcode{}\nwbegindocs{2281}\nwdocspar
\subsection{Recover event}
Recalculate the process instance contents, given an event with known particle
set.  The indices for MC, term, and channel must be already set.  The
{\tt{}recalculate} method of the selected entry will import the result
into {\tt{}sqme{\char95}prc} and {\tt{}weight{\char95}prc}.
\nwenddocs{}\nwbegincode{2282}\moddef{Simulations: simulation: TBP}\plusendmoddef
  procedure :: recalculate => simulation_recalculate
\nwendcode{}\nwbegincode{2283}\moddef{Simulations: procedures}\plusendmoddef
  subroutine simulation_recalculate (simulation)
    class(simulation_t), intent(inout) :: simulation
    integer :: i_prc
    i_prc = simulation%i_prc
    associate (entry => simulation%entry(i_prc))
      if (simulation%update_weight) then
         call simulation%entry(i_prc)%recalculate &
              (update_sqme = simulation%update_sqme, &
              recover_beams = simulation%recover_beams, &
              weight_factor = entry%get_kinematical_weight ())
      else
         call simulation%entry(i_prc)%recalculate &
              (update_sqme = simulation%update_sqme, &
              recover_beams = simulation%recover_beams)
      end if
    end associate
  end subroutine simulation_recalculate

\eatline
\nwendcode{}\nwbegindocs{2284}\nwdocspar
\subsection{Extract contents}
Return the MD5 sum that summarizes configuration and integration
(but not the event file).  Used for initializing the event streams.
\nwenddocs{}\nwbegincode{2285}\moddef{Simulations: simulation: TBP}\plusendmoddef
  procedure :: get_md5sum_prc => simulation_get_md5sum_prc
  procedure :: get_md5sum_cfg => simulation_get_md5sum_cfg
  procedure :: get_md5sum_alt => simulation_get_md5sum_alt
\nwendcode{}\nwbegincode{2286}\moddef{Simulations: procedures}\plusendmoddef
  function simulation_get_md5sum_prc (simulation) result (md5sum)
    class(simulation_t), intent(in) :: simulation
    character(32) :: md5sum
    md5sum = simulation%md5sum_prc
  end function simulation_get_md5sum_prc
    
  function simulation_get_md5sum_cfg (simulation) result (md5sum)
    class(simulation_t), intent(in) :: simulation
    character(32) :: md5sum
    md5sum = simulation%md5sum_cfg
  end function simulation_get_md5sum_cfg
    
  function simulation_get_md5sum_alt (simulation, i) result (md5sum)
    class(simulation_t), intent(in) :: simulation
    integer, intent(in) :: i
    character(32) :: md5sum
    md5sum = simulation%md5sum_alt(i)
  end function simulation_get_md5sum_alt
    
\eatline
\eatline
\nwendcode{}\nwbegindocs{2287}Return data that may be useful for writing event files..
\nwenddocs{}\nwbegincode{2288}\moddef{Simulations: simulation: TBP}\plusendmoddef
  procedure :: get_data => simulation_get_data
\nwendcode{}\nwbegincode{2289}\moddef{Simulations: procedures}\plusendmoddef
  function simulation_get_data (simulation, alt) result (sdata)
    class(simulation_t), intent(in) :: simulation
    logical, intent(in), optional :: alt
    type(event_sample_data_t) :: sdata
    type(process_t), pointer :: process
    type(beam_data_t), pointer :: beam_data
    type(flavor_t), dimension(:), allocatable :: flv
    integer :: n, i
    logical :: enable_alt
    enable_alt = .true.;  if (present (alt))  enable_alt = alt    
    process => simulation%entry(1)%get_process_ptr ()
    beam_data => process%get_beam_data_ptr ()
    if (enable_alt) then
       call sdata%init (simulation%n_prc, simulation%n_alt)
       do i = 1, simulation%n_alt
          sdata%md5sum_alt(i) = simulation%get_md5sum_alt (i)
       end do
    else
       call sdata%init (simulation%n_prc)
    end if
    sdata%unweighted = simulation%unweighted
    sdata%negative_weights = simulation%negative_weights
    sdata%norm_mode = simulation%norm_mode
    n = beam_data_get_n_in (beam_data)
    sdata%n_beam = n
    allocate (flv (n))
    flv = beam_data_get_flavor (beam_data)
    sdata%pdg_beam(:n) = flv%get_pdg ()
    sdata%energy_beam(:n) = beam_data_get_energy (beam_data)
    do i = 1, simulation%n_prc
       if (.not. simulation%entry(i)%valid) cycle
       process => simulation%entry(i)%get_process_ptr ()
       sdata%proc_num_id(i) = process%get_num_id ()
       if (sdata%proc_num_id(i) == 0)  sdata%proc_num_id(i) = i
       if (simulation%entry(i)%has_integral) then
          sdata%cross_section(i) = simulation%entry(i)%integral
          sdata%error(i) = simulation%entry(i)%error
       end if
    end do
    sdata%total_cross_section = sum (sdata%cross_section)
    sdata%md5sum_prc = simulation%get_md5sum_prc ()
    sdata%md5sum_cfg = simulation%get_md5sum_cfg ()
    if (simulation%split_n_evt > 0) then
       sdata%split_n_evt = simulation%split_n_evt
       sdata%split_index = simulation%split_index
    end if
  end function simulation_get_data
    
\eatline
\nwendcode{}\nwbegindocs{2290}Return a default name for the current event sample.  This is the
process ID of the first process.
\nwenddocs{}\nwbegincode{2291}\moddef{Simulations: simulation: TBP}\plusendmoddef
  procedure :: get_default_sample_name => simulation_get_default_sample_name
\nwendcode{}\nwbegincode{2292}\moddef{Simulations: procedures}\plusendmoddef
  function simulation_get_default_sample_name (simulation) result (sample)
    class(simulation_t), intent(in) :: simulation
    type(string_t) :: sample
    type(process_t), pointer :: process
    sample = "whizard"
    if (simulation%n_prc > 0) then
       process => simulation%entry(1)%get_process_ptr ()
       if (associated (process)) then
          sample = process%get_id ()
       end if
    end if
  end function simulation_get_default_sample_name

\eatline
\nwendcode{}\nwbegindocs{2293}\nwdocspar
\nwenddocs{}\nwbegincode{2294}\moddef{Simulations: simulation: TBP}\plusendmoddef
  procedure :: is_valid => simulation_is_valid
\nwendcode{}\nwbegincode{2295}\moddef{Simulations: procedures}\plusendmoddef
  function simulation_is_valid (simulation) result (valid)
    class(simulation_t), intent(inout) :: simulation
    logical :: valid
    valid = simulation%valid
  end function simulation_is_valid

\eatline
\nwendcode{}\nwbegindocs{2296}\nwdocspar
\subsection{Auxiliary}
Call pacify: eliminate numerical noise.
\nwenddocs{}\nwbegincode{2297}\moddef{Simulations: interfaces}\endmoddef
  interface pacify
     module procedure pacify_simulation
  end interface
\nwendcode{}\nwbegincode{2298}\moddef{Simulations: procedures}\plusendmoddef
  subroutine pacify_simulation (simulation)
    class(simulation_t), intent(inout) :: simulation
    integer :: i, j
    i = simulation%i_prc
    if (i > 0) then
       call pacify (simulation%entry(i))
       do j = 1, simulation%n_alt
          call pacify (simulation%alt_entry(i,j))
       end do
    end if
  end subroutine pacify_simulation
  
\eatline
\nwendcode{}\nwbegindocs{2299}\nwdocspar
\subsection{Test}
This is the master for calling self-test procedures.
\nwenddocs{}\nwbegincode{2300}\moddef{Simulations: public}\plusendmoddef
  public :: simulations_test
\nwendcode{}\nwbegincode{2301}\moddef{Simulations: tests}\endmoddef
  subroutine simulations_test (u, results)
    integer, intent(in) :: u
    type(test_results_t), intent(inout) :: results
  \LA{}Simulations: execute tests\RA{}
  end subroutine simulations_test

\eatline
\nwendcode{}\nwbegindocs{2302}\nwdocspar
\subsubsection{Initialization}
Initialize a {\tt{}simulation{\char95}t} object, including the embedded event records.
\nwenddocs{}\nwbegincode{2303}\moddef{Simulations: execute tests}\endmoddef
  call test (simulations_1, "simulations_1", &
       "initialization", &
       u, results)
\nwendcode{}\nwbegincode{2304}\moddef{Simulations: tests}\plusendmoddef
  subroutine simulations_1 (u)
    integer, intent(in) :: u
    type(string_t) :: libname, procname1, procname2
    type(rt_data_t), target :: global
    type(simulation_t), target :: simulation
    
    write (u, "(A)")  "* Test output: simulations_1"
    write (u, "(A)")  "*   Purpose: initialize simulation"
    write (u, "(A)")

    write (u, "(A)")  "* Initialize processes"
    write (u, "(A)")

    call syntax_model_file_init ()

    call global%global_init ()
    call global%set_log (var_str ("?omega_openmp"), &
         .false., is_known = .true.)
    call global%set_int (var_str ("seed"), &
         0, is_known = .true.)    
    
    libname = "simulation_1a"
    procname1 = "simulation_1p"
    
    call prepare_test_library (global, libname, 1, [procname1])
    call compile_library (libname, global)

    call global%set_string (var_str ("$method"), &
         var_str ("unit_test"), is_known = .true.)
    call global%set_string (var_str ("$phs_method"), &
         var_str ("single"), is_known = .true.)
    call global%set_string (var_str ("$integration_method"),&
         var_str ("midpoint"), is_known = .true.)
    call global%set_log (var_str ("?vis_history"),&
         .false., is_known = .true.)    
    call global%set_log (var_str ("?integration_timer"),&
         .false., is_known = .true.)    
    call global%set_log (var_str ("?recover_beams"), &
         .false., is_known = .true.)
    
    call global%set_real (var_str ("sqrts"),&
         1000._default, is_known = .true.)

    call global%it_list%init ([1], [1000])

    call global%set_string (var_str ("$run_id"), &
         var_str ("simulations1"), is_known = .true.)
    call integrate_process (procname1, global, local_stack=.true.)

    procname2 = "sim_extra"
    
    call prepare_test_library (global, libname, 1, [procname2])
    call compile_library (libname, global)
    call global%set_string (var_str ("$run_id"), &
         var_str ("simulations2"), is_known = .true.)


    write (u, "(A)")  "* Initialize event generation"
    write (u, "(A)")

    call global%set_string (var_str ("$sample"), &
         var_str ("sim1"), is_known = .true.)
    call integrate_process (procname2, global, local_stack=.true.)

    call simulation%init ([procname1, procname2], .false., .true., global)
    call simulation%init_process_selector ()
    call simulation%write (u)

    write (u, "(A)")
    write (u, "(A)")  "* Write the event record for the first process"
    write (u, "(A)")
    
    call simulation%write_event (u, i_prc = 1)

    write (u, "(A)")
    write (u, "(A)")  "* Cleanup"

    call simulation%final ()
    call global%final ()
    
    write (u, "(A)")
    write (u, "(A)")  "* Test output end: simulations_1"
    
  end subroutine simulations_1
  
\eatline
\nwendcode{}\nwbegindocs{2305}\nwdocspar
\subsubsection{Weighted events}
Generate events for a single process.
\nwenddocs{}\nwbegincode{2306}\moddef{Simulations: execute tests}\plusendmoddef
  call test (simulations_2, "simulations_2", &
       "weighted events", &
       u, results)
\nwendcode{}\nwbegincode{2307}\moddef{Simulations: tests}\plusendmoddef
  subroutine simulations_2 (u)
    integer, intent(in) :: u
    type(string_t) :: libname, procname1
    type(rt_data_t), target :: global
    type(simulation_t), target :: simulation
    type(event_sample_data_t) :: data
    
    write (u, "(A)")  "* Test output: simulations_2"
    write (u, "(A)")  "*   Purpose: generate events for a single process"
    write (u, "(A)")

    write (u, "(A)")  "* Initialize processes"
    write (u, "(A)")

    call syntax_model_file_init ()

    call global%global_init ()
    call global%set_log (var_str ("?omega_openmp"), &
         .false., is_known = .true.)
    call global%set_int (var_str ("seed"), &
         0, is_known = .true.)    
    
    libname = "simulation_2a"
    procname1 = "simulation_2p"
    
    call prepare_test_library (global, libname, 1, [procname1])
    call compile_library (libname, global)

    call global%append_log (&
         var_str ("?rebuild_events"), .true., intrinsic = .true.)

    call global%set_string (var_str ("$method"), &
         var_str ("unit_test"), is_known = .true.)
    call global%set_string (var_str ("$phs_method"), &
         var_str ("single"), is_known = .true.)
    call global%set_string (var_str ("$integration_method"),&
         var_str ("midpoint"), is_known = .true.)
    call global%set_log (var_str ("?vis_history"),&
         .false., is_known = .true.)    
    call global%set_log (var_str ("?integration_timer"),&
         .false., is_known = .true.)    
    call global%set_log (var_str ("?recover_beams"), &
         .false., is_known = .true.)

    call global%set_real (var_str ("sqrts"),&
         1000._default, is_known = .true.)

    call global%it_list%init ([1], [1000])

    call global%set_string (var_str ("$run_id"), &
         var_str ("simulations1"), is_known = .true.)
    call integrate_process (procname1, global, local_stack=.true.)

    write (u, "(A)")  "* Initialize event generation"
    write (u, "(A)")

    call global%set_log (var_str ("?unweighted"), &
         .false., is_known = .true.)
    call simulation%init ([procname1], .true., .true., global)
    call simulation%init_process_selector ()

    data = simulation%get_data ()
    call data%write (u)

    write (u, "(A)")
    write (u, "(A)")  "* Generate three events"
    write (u, "(A)")

    call simulation%generate (3)
    call simulation%write (u)

    write (u, "(A)")
    write (u, "(A)")  "* Write the event record for the last event"
    write (u, "(A)")
    
    call simulation%write_event (u)

    write (u, "(A)")
    write (u, "(A)")  "* Cleanup"

    call simulation%final ()
    call global%final ()
    
    write (u, "(A)")
    write (u, "(A)")  "* Test output end: simulations_2"
    
  end subroutine simulations_2
  
\eatline
\nwendcode{}\nwbegindocs{2308}\nwdocspar
\subsubsection{Unweighted events}
Generate events for a single process.
\nwenddocs{}\nwbegincode{2309}\moddef{Simulations: execute tests}\plusendmoddef
  call test (simulations_3, "simulations_3", &
       "unweighted events", &
       u, results)
\nwendcode{}\nwbegincode{2310}\moddef{Simulations: tests}\plusendmoddef
  subroutine simulations_3 (u)
    integer, intent(in) :: u
    type(string_t) :: libname, procname1
    type(rt_data_t), target :: global
    type(simulation_t), target :: simulation
    type(event_sample_data_t) :: data
    
    write (u, "(A)")  "* Test output: simulations_3"
    write (u, "(A)")  "*   Purpose: generate unweighted events &
         &for a single process"
    write (u, "(A)")

    write (u, "(A)")  "* Initialize processes"
    write (u, "(A)")

    call syntax_model_file_init ()

    call global%global_init ()
    call global%set_log (var_str ("?omega_openmp"), &
         .false., is_known = .true.)
    call global%set_int (var_str ("seed"), &
         0, is_known = .true.)    
    
    libname = "simulation_3a"
    procname1 = "simulation_3p"
    
    call prepare_test_library (global, libname, 1, [procname1])
    call compile_library (libname, global)

    call global%append_log (&
         var_str ("?rebuild_events"), .true., intrinsic = .true.)

    call global%set_string (var_str ("$method"), &
         var_str ("unit_test"), is_known = .true.)
    call global%set_string (var_str ("$phs_method"), &
         var_str ("single"), is_known = .true.)
    call global%set_string (var_str ("$integration_method"),&
         var_str ("midpoint"), is_known = .true.)
    call global%set_log (var_str ("?vis_history"),&
         .false., is_known = .true.)    
    call global%set_log (var_str ("?integration_timer"),&
         .false., is_known = .true.)    
    call global%set_log (var_str ("?recover_beams"), &
         .false., is_known = .true.)

    call global%set_real (var_str ("sqrts"),&
         1000._default, is_known = .true.)

    call global%it_list%init ([1], [1000])

    call global%set_string (var_str ("$run_id"), &
         var_str ("simulations1"), is_known = .true.)
    call integrate_process (procname1, global, local_stack=.true.)

    write (u, "(A)")  "* Initialize event generation"
    write (u, "(A)")

    call simulation%init ([procname1], .true., .true., global)
    call simulation%init_process_selector ()

    data = simulation%get_data ()
    call data%write (u)

    write (u, "(A)")
    write (u, "(A)")  "* Generate three events"
    write (u, "(A)")

    call simulation%generate (3)
    call simulation%write (u)

    write (u, "(A)")
    write (u, "(A)")  "* Write the event record for the last event"
    write (u, "(A)")
    
    call simulation%write_event (u)

    write (u, "(A)")
    write (u, "(A)")  "* Cleanup"

    call simulation%final ()
    call global%final ()
    
    write (u, "(A)")
    write (u, "(A)")  "* Test output end: simulations_3"
    
  end subroutine simulations_3
  
\eatline
\nwendcode{}\nwbegindocs{2311}\nwdocspar
\subsubsection{Simulating process with structure functions}
Generate events for a single process.
\nwenddocs{}\nwbegincode{2312}\moddef{Simulations: execute tests}\plusendmoddef
  call test (simulations_4, "simulations_4", &
       "process with structure functions", &
       u, results)
\nwendcode{}\nwbegincode{2313}\moddef{Simulations: tests}\plusendmoddef
  subroutine simulations_4 (u)
    integer, intent(in) :: u
    type(string_t) :: libname, procname1
    type(rt_data_t), target :: global
    type(flavor_t) :: flv
    type(string_t) :: name
    type(simulation_t), target :: simulation
    type(event_sample_data_t) :: data
    
    write (u, "(A)")  "* Test output: simulations_4"
    write (u, "(A)")  "*   Purpose: generate events for a single process &
         &with structure functions"
    write (u, "(A)")

    write (u, "(A)")  "* Initialize processes"
    write (u, "(A)")

    call syntax_model_file_init ()
    call syntax_phs_forest_init ()

    call global%global_init ()
    call global%set_log (var_str ("?omega_openmp"), &
         .false., is_known = .true.)
    call global%set_int (var_str ("seed"), &
         0, is_known = .true.)    
    
    libname = "simulation_4a"
    procname1 = "simulation_4p"
    
    call prepare_test_library (global, libname, 1, [procname1])
    call compile_library (libname, global)

    call global%append_log (&
         var_str ("?rebuild_phase_space"), .true., intrinsic = .true.)
    call global%append_log (&
         var_str ("?rebuild_grids"), .true., intrinsic = .true.)
    call global%append_log (&
         var_str ("?rebuild_events"), .true., intrinsic = .true.)

    call global%set_string (var_str ("$run_id"), &
         var_str ("r1"), is_known = .true.)
    call global%set_string (var_str ("$method"), &
         var_str ("unit_test"), is_known = .true.)
    call global%set_string (var_str ("$phs_method"), &
         var_str ("wood"), is_known = .true.)
    call global%set_string (var_str ("$integration_method"),&
         var_str ("vamp"), is_known = .true.)
    call global%set_log (var_str ("?use_vamp_equivalences"),&
         .true., is_known = .true.)
    call global%set_real (var_str ("sqrts"),&
         1000._default, is_known = .true.)
    call global%model_set_real (var_str ("ms"), &
         0._default)
    call global%set_log (var_str ("?vis_history"),&
         .false., is_known = .true.)    
    call global%set_log (var_str ("?integration_timer"),&
         .false., is_known = .true.)    
    call global%set_log (var_str ("?recover_beams"), &
         .false., is_known = .true.)
    
    call reset_interaction_counter ()

    call flv%init (25, global%model)
    name = flv%get_name ()
    
    call global%beam_structure%init_sf ([name, name], [1])
    call global%beam_structure%set_sf (1, 1, var_str ("sf_test_1"))

    write (u, "(A)")  "* Integrate"
    write (u, "(A)")

    call global%it_list%init ([1], [1000])

    call global%set_string (var_str ("$run_id"), &
         var_str ("r1"), is_known = .true.)
    call integrate_process (procname1, global, local_stack=.true.)

    write (u, "(A)")  "* Initialize event generation"
    write (u, "(A)")

    call global%set_log (var_str ("?unweighted"), &
         .false., is_known = .true.)
    call global%set_string (var_str ("$sample"), &
         var_str ("simulations4"), is_known = .true.)
    call simulation%init ([procname1], .true., .true., global)
    call simulation%init_process_selector ()

    data = simulation%get_data ()
    call data%write (u)

    write (u, "(A)")
    write (u, "(A)")  "* Generate three events"
    write (u, "(A)")

    call simulation%generate (3)
    call simulation%write (u)

    write (u, "(A)")
    write (u, "(A)")  "* Write the event record for the last event"
    write (u, "(A)")
    
    call simulation%write_event (u)

    write (u, "(A)")
    write (u, "(A)")  "* Cleanup"

    call simulation%final ()
    call global%final ()
    
    write (u, "(A)")
    write (u, "(A)")  "* Test output end: simulations_4"
    
  end subroutine simulations_4
  
\eatline
\nwendcode{}\nwbegindocs{2314}\nwdocspar
\subsubsection{Event I/O}
Generate event for a test process, write to file and reread.
\nwenddocs{}\nwbegincode{2315}\moddef{Simulations: execute tests}\plusendmoddef
  call test (simulations_5, "simulations_5", &
       "raw event I/O", &
       u, results)
\nwendcode{}\nwbegincode{2316}\moddef{Simulations: tests}\plusendmoddef
  subroutine simulations_5 (u)
    integer, intent(in) :: u
    type(string_t) :: libname, procname1, sample
    type(rt_data_t), target :: global
    class(eio_t), allocatable :: eio
    type(simulation_t), allocatable, target :: simulation
    
    write (u, "(A)")  "* Test output: simulations_5"
    write (u, "(A)")  "*   Purpose: generate events for a single process"
    write (u, "(A)")  "*            write to file and reread"
    write (u, "(A)")

    write (u, "(A)")  "* Initialize processes"
    write (u, "(A)")

    call syntax_model_file_init ()

    call global%global_init ()
    call global%set_log (var_str ("?omega_openmp"), &
         .false., is_known = .true.)
    call global%set_int (var_str ("seed"), &
         0, is_known = .true.)    
    
    libname = "simulation_5a"
    procname1 = "simulation_5p"
    
    call prepare_test_library (global, libname, 1, [procname1])
    call compile_library (libname, global)

    call global%append_log (&
         var_str ("?rebuild_events"), .true., intrinsic = .true.)

    call global%set_string (var_str ("$method"), &
         var_str ("unit_test"), is_known = .true.)
    call global%set_string (var_str ("$phs_method"), &
         var_str ("single"), is_known = .true.)
    call global%set_string (var_str ("$integration_method"),&
         var_str ("midpoint"), is_known = .true.)   
    call global%set_log (var_str ("?vis_history"),&
         .false., is_known = .true.)    
    call global%set_log (var_str ("?integration_timer"),&
         .false., is_known = .true.)    
    call global%set_log (var_str ("?recover_beams"), &
         .false., is_known = .true.)

    call global%set_real (var_str ("sqrts"),&
         1000._default, is_known = .true.)

    call global%it_list%init ([1], [1000])

    call global%set_string (var_str ("$run_id"), &
         var_str ("simulations5"), is_known = .true.)
    call integrate_process (procname1, global, local_stack=.true.)

    write (u, "(A)")  "* Initialize event generation"
    write (u, "(A)")

    call global%set_log (var_str ("?unweighted"), &
         .false., is_known = .true.)
    sample = "simulations5"
    call global%set_string (var_str ("$sample"), &
         sample, is_known = .true.)
    allocate (simulation)
    call simulation%init ([procname1], .true., .true., global)
    call simulation%init_process_selector ()

    write (u, "(A)")  "* Initialize raw event file"
    write (u, "(A)")

    allocate (eio_raw_t :: eio)
    call eio%init_out (sample)
    
    write (u, "(A)")  "* Generate an event"
    write (u, "(A)")

    call simulation%generate (1)
    call simulation%write_event (u)
    call simulation%write_event (eio)

    call eio%final ()
    deallocate (eio)
    call simulation%final ()
    deallocate (simulation)
    
    write (u, "(A)")
    write (u, "(A)")  "* Re-read the event from file"
    write (u, "(A)")
    
    call global%set_log (var_str ("?update_sqme"), &
         .true., is_known = .true.)
    call global%set_log (var_str ("?update_weight"), &
         .true., is_known = .true.)
    call global%set_log (var_str ("?recover_beams"), &
         .false., is_known = .true.)

    allocate (simulation)
    call simulation%init ([procname1], .true., .true., global)
    call simulation%init_process_selector ()
    allocate (eio_raw_t :: eio)
    call eio%init_in (sample)
    
    call simulation%read_event (eio)
    call simulation%write_event (u)

    write (u, "(A)")
    write (u, "(A)")  "* Recalculate process instance"
    write (u, "(A)")

    call simulation%recalculate ()
    call simulation%entry(simulation%i_prc)%evaluate_expressions ()
    call simulation%write_event (u)

    write (u, "(A)")
    write (u, "(A)")  "* Cleanup"

    call eio%final ()
    call simulation%final ()
    call global%final ()
    
    write (u, "(A)")
    write (u, "(A)")  "* Test output end: simulations_5"
    
  end subroutine simulations_5
  
\eatline
\nwendcode{}\nwbegindocs{2317}\nwdocspar
\subsubsection{Event I/O}
Generate event for a real process with structure functions, write to file and
reread.
\nwenddocs{}\nwbegincode{2318}\moddef{Simulations: execute tests}\plusendmoddef
  call test (simulations_6, "simulations_6", &
       "raw event I/O with structure functions", &
       u, results)
\nwendcode{}\nwbegincode{2319}\moddef{Simulations: tests}\plusendmoddef
  subroutine simulations_6 (u)
    integer, intent(in) :: u
    type(string_t) :: libname, procname1, sample
    type(rt_data_t), target :: global
    class(eio_t), allocatable :: eio
    type(simulation_t), allocatable, target :: simulation
    type(flavor_t) :: flv
    type(string_t) :: name
    
    write (u, "(A)")  "* Test output: simulations_6"
    write (u, "(A)")  "*   Purpose: generate events for a single process"
    write (u, "(A)")  "*            write to file and reread"
    write (u, "(A)")

    write (u, "(A)")  "* Initialize process and integrate"
    write (u, "(A)")

    call syntax_model_file_init ()

    call global%global_init ()
    call global%set_log (var_str ("?omega_openmp"), &
         .false., is_known = .true.)
    call global%set_int (var_str ("seed"), &
         0, is_known = .true.)    
    
    libname = "simulation_6"
    procname1 = "simulation_6p"
    
    call prepare_test_library (global, libname, 1, [procname1])
    call compile_library (libname, global)

    call global%append_log (&
         var_str ("?rebuild_phase_space"), .true., intrinsic = .true.)
    call global%append_log (&
         var_str ("?rebuild_grids"), .true., intrinsic = .true.)
    call global%append_log (&
         var_str ("?rebuild_events"), .true., intrinsic = .true.)

    call global%set_string (var_str ("$method"), &
         var_str ("unit_test"), is_known = .true.)
    call global%set_string (var_str ("$phs_method"), &
         var_str ("wood"), is_known = .true.)
    call global%set_string (var_str ("$integration_method"),&
         var_str ("vamp"), is_known = .true.)
    call global%set_log (var_str ("?use_vamp_equivalences"),&
         .true., is_known = .true.)
    call global%set_log (var_str ("?vis_history"),&
         .false., is_known = .true.)    
    call global%set_log (var_str ("?integration_timer"),&
         .false., is_known = .true.)    
    call global%set_log (var_str ("?recover_beams"), &
         .false., is_known = .true.)

    call global%set_real (var_str ("sqrts"),&
         1000._default, is_known = .true.)
    call global%model_set_real (var_str ("ms"), &
         0._default)

    call flv%init (25, global%model)
    name = flv%get_name ()
    
    call global%beam_structure%init_sf ([name, name], [1])
    call global%beam_structure%set_sf (1, 1, var_str ("sf_test_1"))

    call global%it_list%init ([1], [1000])

    call global%set_string (var_str ("$run_id"), &
         var_str ("r1"), is_known = .true.)
    call integrate_process (procname1, global, local_stack=.true.)

    write (u, "(A)")  "* Initialize event generation"
    write (u, "(A)")

    call reset_interaction_counter ()
    
    call global%set_log (var_str ("?unweighted"), &
         .false., is_known = .true.)
    sample = "simulations6"
    call global%set_string (var_str ("$sample"), &
         sample, is_known = .true.)
    allocate (simulation)
    call simulation%init ([procname1], .true., .true., global)
    call simulation%init_process_selector ()

    write (u, "(A)")  "* Initialize raw event file"
    write (u, "(A)")

    allocate (eio_raw_t :: eio)
    call eio%init_out (sample)
    
    write (u, "(A)")  "* Generate an event"
    write (u, "(A)")

    call simulation%generate (1)
    call pacify (simulation%entry(simulation%i_prc))
    call simulation%write_event (u, verbose = .true., testflag = .true.)
    call simulation%write_event (eio)

    call eio%final ()
    deallocate (eio)
    call simulation%final ()
    deallocate (simulation)
    
    write (u, "(A)")
    write (u, "(A)")  "* Re-read the event from file"
    write (u, "(A)")
    
    call reset_interaction_counter ()
    
    call global%set_log (var_str ("?update_sqme"), &
         .true., is_known = .true.)
    call global%set_log (var_str ("?update_weight"), &
         .true., is_known = .true.)

    allocate (simulation)
    call simulation%init ([procname1], .true., .true., global)
    call simulation%init_process_selector ()
    allocate (eio_raw_t :: eio)
    call eio%init_in (sample)
    
    call simulation%read_event (eio)
    call simulation%write_event (u, verbose = .true., testflag = .true.)

    write (u, "(A)")
    write (u, "(A)")  "* Recalculate process instance"
    write (u, "(A)")

    call simulation%recalculate ()
    call simulation%entry(simulation%i_prc)%evaluate_expressions ()
    call simulation%write_event (u, verbose = .true., testflag = .true.)

    write (u, "(A)")
    write (u, "(A)")  "* Cleanup"

    call eio%final ()
    call simulation%final ()
    call global%final ()
    
    write (u, "(A)")
    write (u, "(A)")  "* Test output end: simulations_6"
    
  end subroutine simulations_6
  
\eatline
\nwendcode{}\nwbegindocs{2320}\nwdocspar
\subsubsection{Automatic Event I/O}
Generate events with raw-format event file as cache: generate, reread,
append.
\nwenddocs{}\nwbegincode{2321}\moddef{Simulations: execute tests}\plusendmoddef
  call test (simulations_7, "simulations_7", &
       "automatic raw event I/O", &
       u, results)
\nwendcode{}\nwbegincode{2322}\moddef{Simulations: tests}\plusendmoddef
  subroutine simulations_7 (u)
    integer, intent(in) :: u
    type(string_t) :: libname, procname1, sample
    type(rt_data_t), target :: global
    type(string_t), dimension(0) :: empty_string_array
    type(event_sample_data_t) :: data
    type(event_stream_array_t) :: es_array
    type(simulation_t), allocatable, target :: simulation
    type(flavor_t) :: flv
    type(string_t) :: name
    
    write (u, "(A)")  "* Test output: simulations_7"
    write (u, "(A)")  "*   Purpose: generate events for a single process"
    write (u, "(A)")  "*            write to file and reread"
    write (u, "(A)")

    write (u, "(A)")  "* Initialize process and integrate"
    write (u, "(A)")

    call syntax_model_file_init ()

    call global%global_init ()
    call global%init_fallback_model &
         (var_str ("SM_hadrons"), var_str ("SM_hadrons.mdl"))

    call global%set_log (var_str ("?omega_openmp"), &
         .false., is_known = .true.)
    call global%set_int (var_str ("seed"), &
         0, is_known = .true.)    
    
    libname = "simulation_7"
    procname1 = "simulation_7p"
    
    call prepare_test_library (global, libname, 1, [procname1])
    call compile_library (libname, global)

    call global%append_log (&
         var_str ("?rebuild_phase_space"), .true., intrinsic = .true.)
    call global%append_log (&
         var_str ("?rebuild_grids"), .true., intrinsic = .true.)
    call global%append_log (&
         var_str ("?rebuild_events"), .true., intrinsic = .true.)

    call global%set_string (var_str ("$method"), &
         var_str ("unit_test"), is_known = .true.)
    call global%set_string (var_str ("$phs_method"), &
         var_str ("wood"), is_known = .true.)
    call global%set_string (var_str ("$integration_method"),&
         var_str ("vamp"), is_known = .true.)
    call global%set_log (var_str ("?use_vamp_equivalences"),&
         .true., is_known = .true.)
    call global%set_log (var_str ("?vis_history"),&
         .false., is_known = .true.)    
    call global%set_log (var_str ("?integration_timer"),&
         .false., is_known = .true.)    
    call global%set_log (var_str ("?recover_beams"), &
         .false., is_known = .true.)

    call global%set_real (var_str ("sqrts"),&
         1000._default, is_known = .true.)
    call global%model_set_real (var_str ("ms"), &
         0._default)

    call flv%init (25, global%model)
    name = flv%get_name ()
    
    call global%beam_structure%init_sf ([name, name], [1])
    call global%beam_structure%set_sf (1, 1, var_str ("sf_test_1"))

    call global%it_list%init ([1], [1000])

    call global%set_string (var_str ("$run_id"), &
         var_str ("r1"), is_known = .true.)
    call integrate_process (procname1, global, local_stack=.true.)

    write (u, "(A)")  "* Initialize event generation"
    write (u, "(A)")

    call reset_interaction_counter ()
    
    call global%set_log (var_str ("?unweighted"), &
         .false., is_known = .true.)
    sample = "simulations7"
    call global%set_string (var_str ("$sample"), &
         sample, is_known = .true.)
    allocate (simulation)
    call simulation%init ([procname1], .true., .true., global)
    call simulation%init_process_selector ()

    write (u, "(A)")  "* Initialize raw event file"
    write (u, "(A)")

    data%md5sum_prc = simulation%get_md5sum_prc ()
    data%md5sum_cfg = simulation%get_md5sum_cfg ()
    call es_array%init (sample, [var_str ("raw")], global, data)
    
    write (u, "(A)")  "* Generate an event"
    write (u, "(A)")

    call simulation%generate (1, es_array)

    call es_array%final ()
    call simulation%final ()
    deallocate (simulation)
    
    write (u, "(A)")  "* Re-read the event from file and generate another one"
    write (u, "(A)")
    
    call global%set_log (&
         var_str ("?rebuild_events"), .false., is_known = .true.)

    call reset_interaction_counter ()
    
    allocate (simulation)
    call simulation%init ([procname1], .true., .true., global)
    call simulation%init_process_selector ()

    data%md5sum_prc = simulation%get_md5sum_prc ()
    data%md5sum_cfg = simulation%get_md5sum_cfg ()
    call es_array%init (sample, empty_string_array, global, data, &
         input = var_str ("raw"))
    
    call simulation%generate (2, es_array)
    
    call pacify (simulation%entry(simulation%i_prc))
    call simulation%write_event (u, verbose = .true.)

    call es_array%final ()
    call simulation%final ()
    deallocate (simulation)
    
    
    write (u, "(A)")
    write (u, "(A)")  "* Re-read both events from file"
    write (u, "(A)")
    
    call reset_interaction_counter ()
    
    allocate (simulation)
    call simulation%init ([procname1], .true., .true., global)
    call simulation%init_process_selector ()

    data%md5sum_prc = simulation%get_md5sum_prc ()
    data%md5sum_cfg = simulation%get_md5sum_cfg ()
    call es_array%init (sample, empty_string_array, global, data, &
         input = var_str ("raw"))

    call simulation%generate (2, es_array)
    
    call pacify (simulation%entry(simulation%i_prc))
    call simulation%write_event (u, verbose = .true.)

    write (u, "(A)")
    write (u, "(A)")  "* Cleanup"

    call es_array%final ()
    call simulation%final ()
    call global%final ()
    
    write (u, "(A)")
    write (u, "(A)")  "* Test output end: simulations_7"
    
  end subroutine simulations_7
  
\eatline
\nwendcode{}\nwbegindocs{2323}\nwdocspar
\subsubsection{Rescanning Events}
Generate events and rescan the resulting raw event file.
\nwenddocs{}\nwbegincode{2324}\moddef{Simulations: execute tests}\plusendmoddef
  call test (simulations_8, "simulations_8", &
       "rescan raw event file", &
       u, results)
\nwendcode{}\nwbegincode{2325}\moddef{Simulations: tests}\plusendmoddef
  subroutine simulations_8 (u)
    integer, intent(in) :: u
    type(string_t) :: libname, procname1, sample
    type(rt_data_t), target :: global
    type(string_t), dimension(0) :: empty_string_array
    type(event_sample_data_t) :: data
    type(event_stream_array_t) :: es_array
    type(simulation_t), allocatable, target :: simulation
    type(flavor_t) :: flv
    type(string_t) :: name
    
    write (u, "(A)")  "* Test output: simulations_8"
    write (u, "(A)")  "*   Purpose: generate events for a single process"
    write (u, "(A)")  "*            write to file and rescan"
    write (u, "(A)")

    write (u, "(A)")  "* Initialize process and integrate"
    write (u, "(A)")

    call syntax_model_file_init ()

    call global%global_init ()
    call global%init_fallback_model &
         (var_str ("SM_hadrons"), var_str ("SM_hadrons.mdl"))

    call global%set_log (var_str ("?omega_openmp"), &
         .false., is_known = .true.)
    call global%set_int (var_str ("seed"), &
         0, is_known = .true.)        

    libname = "simulation_8"
    procname1 = "simulation_8p"
    
    call prepare_test_library (global, libname, 1, [procname1])
    call compile_library (libname, global)

    call global%append_log (&
         var_str ("?rebuild_phase_space"), .true., intrinsic = .true.)
    call global%append_log (&
         var_str ("?rebuild_grids"), .true., intrinsic = .true.)
    call global%append_log (&
         var_str ("?rebuild_events"), .true., intrinsic = .true.)

    call global%set_string (var_str ("$method"), &
         var_str ("unit_test"), is_known = .true.)
    call global%set_string (var_str ("$phs_method"), &
         var_str ("wood"), is_known = .true.)
    call global%set_string (var_str ("$integration_method"),&
         var_str ("vamp"), is_known = .true.)
    call global%set_log (var_str ("?use_vamp_equivalences"),&
         .true., is_known = .true.)   
    call global%set_log (var_str ("?vis_history"),&
         .false., is_known = .true.)    
    call global%set_log (var_str ("?integration_timer"),&
         .false., is_known = .true.)    
    call global%set_log (var_str ("?recover_beams"), &
         .false., is_known = .true.)

    call global%set_real (var_str ("sqrts"),&
         1000._default, is_known = .true.)
    call global%model_set_real (var_str ("ms"), &
         0._default)

    call flv%init (25, global%model)
    name = flv%get_name ()
    
    call global%beam_structure%init_sf ([name, name], [1])
    call global%beam_structure%set_sf (1, 1, var_str ("sf_test_1"))

    call global%it_list%init ([1], [1000])

    call global%set_string (var_str ("$run_id"), &
         var_str ("r1"), is_known = .true.)
    call integrate_process (procname1, global, local_stack=.true.)

    write (u, "(A)")  "* Initialize event generation"
    write (u, "(A)")

    call reset_interaction_counter ()
    
    call global%set_log (var_str ("?unweighted"), &
         .false., is_known = .true.)
    sample = "simulations8"
    call global%set_string (var_str ("$sample"), &
         sample, is_known = .true.)
    allocate (simulation)
    call simulation%init ([procname1], .true., .true., global)
    call simulation%init_process_selector ()

    write (u, "(A)")  "* Initialize raw event file"
    write (u, "(A)")

    data%md5sum_prc = simulation%get_md5sum_prc ()
    data%md5sum_cfg = simulation%get_md5sum_cfg ()
    write (u, "(1x,A,A,A)")  "MD5 sum (proc)   = '", data%md5sum_prc, "'"
    write (u, "(1x,A,A,A)")  "MD5 sum (config) = '", data%md5sum_cfg, "'"
    call es_array%init (sample, [var_str ("raw")], global, &
         data)
    
    write (u, "(A)")
    write (u, "(A)")  "* Generate an event"
    write (u, "(A)")

    call simulation%generate (1, es_array)

    call pacify (simulation%entry(simulation%i_prc))
    call simulation%write_event (u, verbose = .true., testflag = .true.)

    call es_array%final ()
    call simulation%final ()
    deallocate (simulation)
    
    write (u, "(A)")
    write (u, "(A)")  "* Re-read the event from file"
    write (u, "(A)")
    
    call reset_interaction_counter ()
    
    allocate (simulation)
    call simulation%init ([procname1], .false., .false., global)
    call simulation%init_process_selector ()

    data%md5sum_prc = simulation%get_md5sum_prc ()
    data%md5sum_cfg = ""
    write (u, "(1x,A,A,A)")  "MD5 sum (proc)   = '", data%md5sum_prc, "'"
    write (u, "(1x,A,A,A)")  "MD5 sum (config) = '", data%md5sum_cfg, "'"
    call es_array%init (sample, empty_string_array, global, data, &
         input = var_str ("raw"), input_sample = sample, allow_switch = .false.)
    
    call simulation%rescan (1, es_array, global = global)
    
    write (u, "(A)")

    call pacify (simulation%entry(simulation%i_prc))
    call simulation%write_event (u, verbose = .true., testflag = .true.)

    call es_array%final ()
    call simulation%final ()
    deallocate (simulation)
    
    write (u, "(A)")
    write (u, "(A)")  "* Re-read again and recalculate"
    write (u, "(A)")
    
    call reset_interaction_counter ()
    
    call global%set_log (var_str ("?update_sqme"), &
         .true., is_known = .true.)
    call global%set_log (var_str ("?update_event"), &
         .true., is_known = .true.)

    allocate (simulation)
    call simulation%init ([procname1], .false., .false., global)
    call simulation%init_process_selector ()

    data%md5sum_prc = simulation%get_md5sum_prc ()
    data%md5sum_cfg = ""
    write (u, "(1x,A,A,A)")  "MD5 sum (proc)   = '", data%md5sum_prc, "'"
    write (u, "(1x,A,A,A)")  "MD5 sum (config) = '", data%md5sum_cfg, "'"
    call es_array%init (sample, empty_string_array, global, data, &
         input = var_str ("raw"), input_sample = sample, allow_switch = .false.)
    
    call simulation%rescan (1, es_array, global = global)
    
    write (u, "(A)")

    call pacify (simulation%entry(simulation%i_prc))
    call simulation%write_event (u, verbose = .true., testflag = .true.)

    write (u, "(A)")
    write (u, "(A)")  "* Cleanup"

    call es_array%final ()
    call simulation%final ()
    call global%final ()
    
    write (u, "(A)")
    write (u, "(A)")  "* Test output end: simulations_8"
    
  end subroutine simulations_8
  
\eatline
\nwendcode{}\nwbegindocs{2326}\nwdocspar
\subsubsection{Rescanning Check}
Generate events and rescan with process mismatch.
\nwenddocs{}\nwbegincode{2327}\moddef{Simulations: execute tests}\plusendmoddef
  call test (simulations_9, "simulations_9", &
       "rescan mismatch", &
       u, results)
\nwendcode{}\nwbegincode{2328}\moddef{Simulations: tests}\plusendmoddef
  subroutine simulations_9 (u)
    integer, intent(in) :: u
    type(string_t) :: libname, procname1, sample
    type(rt_data_t), target :: global
    type(string_t), dimension(0) :: empty_string_array
    type(event_sample_data_t) :: data
    type(event_stream_array_t) :: es_array
    type(simulation_t), allocatable, target :: simulation
    type(flavor_t) :: flv
    type(string_t) :: name
    logical :: error
    
    write (u, "(A)")  "* Test output: simulations_9"
    write (u, "(A)")  "*   Purpose: generate events for a single process"
    write (u, "(A)")  "*            write to file and rescan"
    write (u, "(A)")

    write (u, "(A)")  "* Initialize process and integrate"
    write (u, "(A)")

    call syntax_model_file_init ()

    call global%global_init ()
    call global%init_fallback_model &
         (var_str ("SM_hadrons"), var_str ("SM_hadrons.mdl"))

    call global%set_log (var_str ("?omega_openmp"), &
         .false., is_known = .true.)
    call global%set_int (var_str ("seed"), &
         0, is_known = .true.)    
    
    libname = "simulation_9"
    procname1 = "simulation_9p"
    
    call prepare_test_library (global, libname, 1, [procname1])
    call compile_library (libname, global)

    call global%append_log (&
         var_str ("?rebuild_phase_space"), .true., intrinsic = .true.)
    call global%append_log (&
         var_str ("?rebuild_grids"), .true., intrinsic = .true.)
    call global%append_log (&
         var_str ("?rebuild_events"), .true., intrinsic = .true.)

    call global%set_string (var_str ("$method"), &
         var_str ("unit_test"), is_known = .true.)
    call global%set_string (var_str ("$phs_method"), &
         var_str ("wood"), is_known = .true.)
    call global%set_string (var_str ("$integration_method"),&
         var_str ("vamp"), is_known = .true.)
    call global%set_log (var_str ("?use_vamp_equivalences"),&
         .true., is_known = .true.)
    call global%set_log (var_str ("?vis_history"),&
         .false., is_known = .true.)    
    call global%set_log (var_str ("?integration_timer"),&
         .false., is_known = .true.)    
    call global%set_log (var_str ("?recover_beams"), &
         .false., is_known = .true.)

    call global%set_real (var_str ("sqrts"),&
         1000._default, is_known = .true.)
    call global%model_set_real (var_str ("ms"), &
         0._default)

    call flv%init (25, global%model)
    name = flv%get_name ()
    
    call global%beam_structure%init_sf ([name, name], [1])
    call global%beam_structure%set_sf (1, 1, var_str ("sf_test_1"))

    call global%it_list%init ([1], [1000])

    call global%set_string (var_str ("$run_id"), &
         var_str ("r1"), is_known = .true.)
    call integrate_process (procname1, global, local_stack=.true.)

    write (u, "(A)")  "* Initialize event generation"
    write (u, "(A)")

    call reset_interaction_counter ()
    
    call global%set_log (var_str ("?unweighted"), &
         .false., is_known = .true.)
    sample = "simulations9"
    call global%set_string (var_str ("$sample"), &
         sample, is_known = .true.)
    allocate (simulation)
    call simulation%init ([procname1], .true., .true., global)
    call simulation%init_process_selector ()

    call simulation%write (u)

    write (u, "(A)")
    write (u, "(A)")  "* Initialize raw event file"
    write (u, "(A)")

    data%md5sum_prc = simulation%get_md5sum_prc ()
    data%md5sum_cfg = simulation%get_md5sum_cfg ()
    write (u, "(1x,A,A,A)")  "MD5 sum (proc)   = '", data%md5sum_prc, "'"
    write (u, "(1x,A,A,A)")  "MD5 sum (config) = '", data%md5sum_cfg, "'"
    call es_array%init (sample, [var_str ("raw")], global, &
         data)
    
    write (u, "(A)")
    write (u, "(A)")  "* Generate an event"
    write (u, "(A)")

    call simulation%generate (1, es_array)

    call es_array%final ()
    call simulation%final ()
    deallocate (simulation)
    
    write (u, "(A)")  "* Initialize event generation for different parameters"
    write (u, "(A)")
    
    call reset_interaction_counter ()
    
    allocate (simulation)
    call simulation%init ([procname1, procname1], .false., .false., global)
    call simulation%init_process_selector ()
    
    call simulation%write (u)

    write (u, "(A)")
    write (u, "(A)")  "* Attempt to re-read the events (should fail)"
    write (u, "(A)")

    data%md5sum_prc = simulation%get_md5sum_prc ()
    data%md5sum_cfg = ""
    write (u, "(1x,A,A,A)")  "MD5 sum (proc)   = '", data%md5sum_prc, "'"
    write (u, "(1x,A,A,A)")  "MD5 sum (config) = '", data%md5sum_cfg, "'"
    call es_array%init (sample, empty_string_array, global, data, &
         input = var_str ("raw"), input_sample = sample, &
         allow_switch = .false., error = error)
    
    write (u, "(1x,A,L1)")  "error = ", error
    
    call simulation%rescan (1, es_array, global = global)

    call es_array%final ()
    call simulation%final ()
    call global%final ()
    
    write (u, "(A)")
    write (u, "(A)")  "* Test output end: simulations_9"
    
  end subroutine simulations_9
  
\eatline
\nwendcode{}\nwbegindocs{2329}\nwdocspar
\subsubsection{Alternative weights}
Generate an event for a single process and reweight it in a
simultaneous calculation.
\nwenddocs{}\nwbegincode{2330}\moddef{Simulations: execute tests}\plusendmoddef
  call test (simulations_10, "simulations_10", &
       "alternative weight", &
       u, results)
\nwendcode{}\nwbegincode{2331}\moddef{Simulations: tests}\plusendmoddef
  subroutine simulations_10 (u)
    integer, intent(in) :: u
    type(string_t) :: libname, procname1, expr_text
    type(rt_data_t), target :: global
    type(rt_data_t), dimension(1), target :: alt_env
    type(ifile_t) :: ifile
    type(stream_t) :: stream
    type(parse_tree_t) :: pt_weight
    type(simulation_t), target :: simulation
    type(event_sample_data_t) :: data
    
    write (u, "(A)")  "* Test output: simulations_10"
    write (u, "(A)")  "*   Purpose: reweight event"
    write (u, "(A)")

    write (u, "(A)")  "* Initialize processes"
    write (u, "(A)")

    call syntax_model_file_init ()
    call syntax_pexpr_init ()

    call global%global_init ()
    call global%set_log (var_str ("?omega_openmp"), &
         .false., is_known = .true.)
    call global%set_int (var_str ("seed"), &
         0, is_known = .true.)    
    
    libname = "simulation_10a"
    procname1 = "simulation_10p"
    
    call prepare_test_library (global, libname, 1, [procname1])
    call compile_library (libname, global)

    call global%append_log (&
         var_str ("?rebuild_phase_space"), .true., intrinsic = .true.)
    call global%append_log (&
         var_str ("?rebuild_grids"), .true., intrinsic = .true.)
    call global%append_log (&
         var_str ("?rebuild_events"), .true., intrinsic = .true.)

    call global%set_string (var_str ("$method"), &
         var_str ("unit_test"), is_known = .true.)
    call global%set_string (var_str ("$phs_method"), &
         var_str ("single"), is_known = .true.)
    call global%set_string (var_str ("$integration_method"),&
         var_str ("midpoint"), is_known = .true.)
    call global%set_log (var_str ("?vis_history"),&
         .false., is_known = .true.)    
    call global%set_log (var_str ("?integration_timer"),&
         .false., is_known = .true.)    
    call global%set_log (var_str ("?recover_beams"), &
         .false., is_known = .true.)

    call global%set_real (var_str ("sqrts"),&
         1000._default, is_known = .true.)

    call global%it_list%init ([1], [1000])

    call global%set_string (var_str ("$run_id"), &
         var_str ("simulations1"), is_known = .true.)
    call integrate_process (procname1, global, local_stack=.true.)

    write (u, "(A)")  "* Initialize alternative environment with custom weight"
    write (u, "(A)")
    
    call alt_env(1)%local_init (global)
    call alt_env(1)%activate ()

    expr_text = "2"
    write (u, "(A,A)")  "weight = ", char (expr_text)
    write (u, *)
    
    call ifile_clear (ifile)
    call ifile_append (ifile, expr_text)
    call stream_init (stream, ifile)
    call parse_tree_init_expr (pt_weight, stream, .true.)
    call stream_final (stream)
    alt_env(1)%pn%weight_expr => parse_tree_get_root_ptr (pt_weight)
    call alt_env(1)%write_expr (u)

    write (u, "(A)")
    write (u, "(A)")  "* Initialize event generation"
    write (u, "(A)")

    call global%set_log (var_str ("?unweighted"), &
         .false., is_known = .true.)
    call simulation%init ([procname1], .true., .true., global, alt_env=alt_env)
    call simulation%init_process_selector ()

    data = simulation%get_data ()
    call data%write (u)

    write (u, "(A)")
    write (u, "(A)")  "* Generate an event"
    write (u, "(A)")

    call simulation%generate (1)
    call simulation%write (u)

    write (u, "(A)")
    write (u, "(A)")  "* Write the event record for the last event"
    write (u, "(A)")
    
    call simulation%write_event (u)

    write (u, "(A)")
    write (u, "(A)")  "* Write the event record for the alternative setup"
    write (u, "(A)")
    
    call simulation%write_alt_event (u)

    write (u, "(A)")
    write (u, "(A)")  "* Cleanup"

    call simulation%final ()
    call global%final ()
    
    call syntax_model_file_final ()
    call syntax_pexpr_final ()

    write (u, "(A)")
    write (u, "(A)")  "* Test output end: simulations_10"
    
  end subroutine simulations_10
  
\eatline
\nwendcode{}\nwbegindocs{2332}\nwdocspar
\subsubsection{Decays}
Generate an event with subsequent partonic decays.
\nwenddocs{}\nwbegincode{2333}\moddef{Simulations: execute tests}\plusendmoddef
  call test (simulations_11, "simulations_11", &
       "decay", &
       u, results)
\nwendcode{}\nwbegincode{2334}\moddef{Simulations: tests}\plusendmoddef
  subroutine simulations_11 (u)
    integer, intent(in) :: u
    type(rt_data_t), target :: global
    type(prclib_entry_t), pointer :: lib
    type(string_t) :: prefix, procname1, procname2
    type(process_t), pointer :: process
    type(simulation_t), target :: simulation
    
    write (u, "(A)")  "* Test output: simulations_11"
    write (u, "(A)")  "*   Purpose: apply decay"
    write (u, "(A)")

    write (u, "(A)")  "* Initialize processes"
    write (u, "(A)")

    call syntax_model_file_init ()
        
    call global%global_init ()
    allocate (lib)
    call global%add_prclib (lib)

    call global%set_int (var_str ("seed"), &
         0, is_known = .true.)        
    call global%set_log (var_str ("?recover_beams"), &
         .false., is_known = .true.)
    
    prefix = "simulation_11"
    procname1 = prefix // "_p"
    procname2 = prefix // "_d"
    call prepare_testbed &
         (global%prclib, global%process_stack, &
         prefix, global%os_data, &
         scattering=.true., decay=.true.)

    call global%select_model (var_str ("Test"))
    call global%model%set_par (var_str ("ff"), 0.4_default)
    call global%model%set_par (var_str ("mf"), &
         global%model%get_real (var_str ("ff")) &
         * global%model%get_real (var_str ("ms")))
    call global%model%set_unstable (25, [procname2])

    write (u, "(A)")  "* Initialize simulation object"
    write (u, "(A)")

    call simulation%init ([procname1], .true., .true., global)
    call simulation%init_process_selector ()

    write (u, "(A)")  "* Generate event"
    write (u, "(A)")

    call simulation%generate (1)
    call simulation%write (u)

    write (u, *)
    
    call simulation%write_event (u)

    write (u, "(A)")
    write (u, "(A)")  "* Cleanup"
    write (u, "(A)")

    call simulation%final ()
    call global%final ()
    
    call syntax_model_file_final ()

    write (u, "(A)")
    write (u, "(A)")  "* Test output end: simulations_11"
    
  end subroutine simulations_11
  
\eatline
\nwendcode{}\nwbegindocs{2335}\nwdocspar
\subsubsection{Split Event Files}
Generate event for a real process with structure functions and write to file,
accepting a limit for the number of events per file.
\nwenddocs{}\nwbegincode{2336}\moddef{Simulations: execute tests}\plusendmoddef
  call test (simulations_12, "simulations_12", &
       "split event files", &
       u, results)
\nwendcode{}\nwbegincode{2337}\moddef{Simulations: tests}\plusendmoddef
  subroutine simulations_12 (u)
    integer, intent(in) :: u
    type(string_t) :: libname, procname1, sample
    type(rt_data_t), target :: global
    class(eio_t), allocatable :: eio
    type(simulation_t), allocatable, target :: simulation
    type(flavor_t) :: flv
    integer :: i_evt
    
    write (u, "(A)")  "* Test output: simulations_12"
    write (u, "(A)")  "*   Purpose: generate events for a single process"
    write (u, "(A)")  "*            and write to split event files"
    write (u, "(A)")

    write (u, "(A)")  "* Initialize process and integrate"
    write (u, "(A)")

    call syntax_model_file_init ()

    call global%global_init ()
    call global%set_log (var_str ("?omega_openmp"), &
         .false., is_known = .true.)
    call global%set_int (var_str ("seed"), &
         0, is_known = .true.)    
    
    libname = "simulation_12"
    procname1 = "simulation_12p"
    
    call prepare_test_library (global, libname, 1, [procname1])
    call compile_library (libname, global)

    call global%append_log (&
         var_str ("?rebuild_phase_space"), .true., intrinsic = .true.)
    call global%append_log (&
         var_str ("?rebuild_grids"), .true., intrinsic = .true.)
    call global%append_log (&
         var_str ("?rebuild_events"), .true., intrinsic = .true.)

    call global%set_string (var_str ("$method"), &
         var_str ("unit_test"), is_known = .true.)
    call global%set_string (var_str ("$phs_method"), &
         var_str ("single"), is_known = .true.)
    call global%set_string (var_str ("$integration_method"),&
         var_str ("midpoint"), is_known = .true.)
    call global%set_log (var_str ("?vis_history"),&
         .false., is_known = .true.)    
    call global%set_log (var_str ("?integration_timer"),&
         .false., is_known = .true.)    
    call global%set_log (var_str ("?recover_beams"), &
         .false., is_known = .true.)

    call global%set_real (var_str ("sqrts"),&
         1000._default, is_known = .true.)
    call global%model_set_real (var_str ("ms"), &
         0._default)

    call flv%init (25, global%model)

    call global%it_list%init ([1], [1000])

    call global%set_string (var_str ("$run_id"), &
         var_str ("r1"), is_known = .true.)
    call integrate_process (procname1, global, local_stack=.true.)

    write (u, "(A)")  "* Initialize event generation"
    write (u, "(A)")

    call global%set_log (var_str ("?unweighted"), &
         .false., is_known = .true.)
    sample = "simulations_12"
    call global%set_string (var_str ("$sample"), &
         sample, is_known = .true.)
    call global%set_int (var_str ("sample_split_n_evt"), &
         2, is_known = .true.)
    call global%set_int (var_str ("sample_split_index"), &
         42, is_known = .true.)
    allocate (simulation)
    call simulation%init ([procname1], .true., .true., global)
    call simulation%init_process_selector ()

    call simulation%write (u)

    write (u, "(A)")
    write (u, "(A)")  "* Initialize ASCII event file"
    write (u, "(A)")

    allocate (eio_ascii_short_t :: eio)
    select type (eio)
    class is (eio_ascii_t);  call eio%set_parameters ()
    end select
    call eio%init_out (sample, data = simulation%get_data ())
    
    write (u, "(A)")  "* Generate 5 events, distributed among three files"

    do i_evt = 1, 5
       call simulation%generate (1)
       call simulation%write_event (eio)
    end do

    call eio%final ()
    deallocate (eio)
    call simulation%final ()
    deallocate (simulation)
    
    write (u, *)
    call display_file ("simulations_12.42.short.evt", u)
    write (u, *)
    call display_file ("simulations_12.43.short.evt", u)
    write (u, *)
    call display_file ("simulations_12.44.short.evt", u)

    write (u, "(A)")
    write (u, "(A)")  "* Cleanup"

    call global%final ()
    
    write (u, "(A)")
    write (u, "(A)")  "* Test output end: simulations_12"
    
  end subroutine simulations_12
  
\eatline
\nwendcode{}\nwbegindocs{2338}Auxiliary: display file contents.
\nwenddocs{}\nwbegincode{2339}\moddef{Simulations: tests}\plusendmoddef
  subroutine display_file (file, u)
    character(*), intent(in) :: file
    integer, intent(in) :: u
    character(256) :: buffer
    integer :: u_file
    write (u, "(3A)")  "* Contents of file '", file, "':"
    write (u, *)
    u_file = free_unit ()
    open (u_file, file = file, action = "read", status = "old")
    do
       read (u_file, "(A)", end = 1)  buffer
       write (u, "(A)")  trim (buffer)
    end do
1   continue
  end subroutine display_file

\eatline
\nwendcode{}\nwbegindocs{2340}\nwdocspar
%%%%%%%%%%%%%%%%%%%%%%%%%%%%%%%%%%%%%%%%%%%%%%%%%%%%%%%%%%%%%%%%%%%%%%%%
\chapter{More Unit Tests}

This chapter collects some procedures for testing that can't be
provided at the point where the corresponding modules are defined,
because they use other modules of a different level.

(We should move them back, collecting the high-level functionality in
init/final hooks that we can set at runtime.)


\section{Expression Testing}

Expression objects are part of process and event objects, but the
process and event object modules should not depend on the
implementation of expressions.  Here, we collect unit tests that
depend on expression implementation.
\nwenddocs{}\nwbegincode{2341}\moddef{\code{}expr{\char95}tests.f90\edoc{}}\endmoddef
\LA{}File header\RA{}
module expr_tests

\LA{}Use kinds\RA{}
\LA{}Use strings\RA{}
  use format_defs, only: FMT_12
  use format_utils, only: write_separator
  use unit_tests
  use os_interface
  use ifiles
  use lexers
  use parser
  use lorentz
  use sm_qcd
  use model_data
  use interactions, only: reset_interaction_counter
  use variables
  use expr_base
  use eval_trees
  use models
  use subevents
  use subevt_expr
  use phs_base
  use rng_base
  use mci_base
  use process_libraries
  use prc_test
  use prc_core
  use processes
  use events

\LA{}Standard module head\RA{}

\LA{}Expr tests: public\RA{}

contains
  
\LA{}Expr tests: tests\RA{}

end module expr_tests
\eatline
\nwendcode{}\nwbegindocs{2342}\nwdocspar
\subsection{Test}
This is the master for calling self-test procedures.
\nwenddocs{}\nwbegincode{2343}\moddef{Expr tests: public}\endmoddef
  public :: subevt_expr_test
\nwendcode{}\nwbegincode{2344}\moddef{Expr tests: tests}\endmoddef
  subroutine subevt_expr_test (u, results)
    integer, intent(in) :: u
    type(test_results_t), intent(inout) :: results
  \LA{}Expr tests: execute tests\RA{}
end subroutine subevt_expr_test

\eatline
\nwendcode{}\nwbegindocs{2345}\nwdocspar
\subsubsection{Parton-event expressions}
\nwenddocs{}\nwbegincode{2346}\moddef{Expr tests: execute tests}\endmoddef
  call test (subevt_expr_1, "subevt_expr_1", &
       "parton-event expressions", &
       u, results)
\nwendcode{}\nwbegincode{2347}\moddef{Expr tests: tests}\plusendmoddef
  subroutine subevt_expr_1 (u)
    integer, intent(in) :: u
    type(string_t) :: expr_text
    type(ifile_t) :: ifile
    type(stream_t) :: stream
    type(parse_tree_t) :: pt_cuts, pt_scale, pt_fac_scale, pt_ren_scale
    type(parse_tree_t) :: pt_weight
    type(parse_node_t), pointer :: pn_cuts, pn_scale, pn_fac_scale, pn_ren_scale
    type(parse_node_t), pointer :: pn_weight
    type(eval_tree_factory_t) :: expr_factory
    type(os_data_t) :: os_data
    type(model_t), target :: model
    type(parton_expr_t), target :: expr
    real(default) :: E, Ex, m
    type(vector4_t), dimension(6) :: p
    integer :: i, pdg
    logical :: passed
    real(default) :: scale, fac_scale, ren_scale, weight
    
    write (u, "(A)")  "* Test output: subevt_expr_1"
    write (u, "(A)")  "*   Purpose: Set up a subevt and associated &
         &process-specific expressions"
    write (u, "(A)")

    call syntax_pexpr_init ()

    call syntax_model_file_init ()
    call os_data_init (os_data)
    call model%read (var_str ("Test.mdl"), os_data)

    write (u, "(A)")  "* Expression texts"
    write (u, "(A)")


    expr_text = "all Pt > 100 [s]"
    write (u, "(A,A)")  "cuts = ", char (expr_text)
    call ifile_clear (ifile)
    call ifile_append (ifile, expr_text)
    call stream_init (stream, ifile)
    call parse_tree_init_lexpr (pt_cuts, stream, .true.)
    call stream_final (stream)
    pn_cuts => parse_tree_get_root_ptr (pt_cuts)

    expr_text = "sqrts"
    write (u, "(A,A)")  "scale = ", char (expr_text)
    call ifile_clear (ifile)
    call ifile_append (ifile, expr_text)
    call stream_init (stream, ifile)
    call parse_tree_init_expr (pt_scale, stream, .true.)
    call stream_final (stream)
    pn_scale => parse_tree_get_root_ptr (pt_scale)

    expr_text = "sqrts_hat"
    write (u, "(A,A)")  "fac_scale = ", char (expr_text)
    call ifile_clear (ifile)
    call ifile_append (ifile, expr_text)
    call stream_init (stream, ifile)
    call parse_tree_init_expr (pt_fac_scale, stream, .true.)
    call stream_final (stream)
    pn_fac_scale => parse_tree_get_root_ptr (pt_fac_scale)

    expr_text = "100"
    write (u, "(A,A)")  "ren_scale = ", char (expr_text)
    call ifile_clear (ifile)
    call ifile_append (ifile, expr_text)
    call stream_init (stream, ifile)
    call parse_tree_init_expr (pt_ren_scale, stream, .true.)
    call stream_final (stream)
    pn_ren_scale => parse_tree_get_root_ptr (pt_ren_scale)

    expr_text = "n_tot - n_in - n_out"
    write (u, "(A,A)")  "weight = ", char (expr_text)
    call ifile_clear (ifile)
    call ifile_append (ifile, expr_text)
    call stream_init (stream, ifile)
    call parse_tree_init_expr (pt_weight, stream, .true.)
    call stream_final (stream)
    pn_weight => parse_tree_get_root_ptr (pt_weight)

    call ifile_final (ifile)

    write (u, "(A)")
    write (u, "(A)")  "* Initialize process expr"
    write (u, "(A)")

    call expr%setup_vars (1000._default)
    call var_list_append_real (expr%var_list, var_str ("tolerance"), 0._default)
    call expr%link_var_list (model%get_var_list_ptr ())

    call expr_factory%init (pn_cuts)
    call expr%setup_selection (expr_factory)
    call expr_factory%init (pn_scale)
    call expr%setup_scale (expr_factory)
    call expr_factory%init (pn_fac_scale)
    call expr%setup_fac_scale (expr_factory)
    call expr_factory%init (pn_ren_scale)
    call expr%setup_ren_scale (expr_factory)
    call expr_factory%init (pn_weight)
    call expr%setup_weight (expr_factory)

    call write_separator (u)
    call expr%write (u)
    call write_separator (u)
    
    write (u, "(A)")
    write (u, "(A)")  "* Fill subevt and evaluate expressions"
    write (u, "(A)")

    call subevt_init (expr%subevt_t, 6)
    E = 500._default
    Ex = 400._default
    m = 125._default
    pdg = 25
    p(1) = vector4_moving (E, sqrt (E**2 - m**2), 3)
    p(2) = vector4_moving (E, -sqrt (E**2 - m**2), 3)
    p(3) = vector4_moving (Ex, sqrt (Ex**2 - m**2), 3)
    p(4) = vector4_moving (Ex, -sqrt (Ex**2 - m**2), 3)
    p(5) = vector4_moving (Ex, sqrt (Ex**2 - m**2), 1)
    p(6) = vector4_moving (Ex, -sqrt (Ex**2 - m**2), 1)

    call expr%reset ()
    do i = 1, 2
       call subevt_set_beam (expr%subevt_t, i, pdg, p(i), m**2)
    end do
    do i = 3, 4
       call subevt_set_incoming (expr%subevt_t, i, pdg, p(i), m**2)
    end do
    do i = 5, 6
       call subevt_set_outgoing (expr%subevt_t, i, pdg, p(i), m**2)
    end do
    expr%sqrts_hat = subevt_get_sqrts_hat (expr%subevt_t)
    expr%n_in = 2
    expr%n_out = 2
    expr%n_tot = 4
    expr%subevt_filled = .true.

    call expr%evaluate (passed, scale, fac_scale, ren_scale, weight)
    
    write (u, "(A,L1)")      "Event has passed      = ", passed
    write (u, "(A," // FMT_12 // ")")  "Scale                 = ", scale
    write (u, "(A," // FMT_12 // ")")  "Factorization scale   = ", fac_scale
    write (u, "(A," // FMT_12 // ")")  "Renormalization scale = ", ren_scale
    write (u, "(A," // FMT_12 // ")")  "Weight                = ", weight
    write (u, "(A)")
    
    call write_separator (u)
    call expr%write (u)
    call write_separator (u)
    
    write (u, "(A)")
    write (u, "(A)")  "* Cleanup"
    
    call expr%final ()

    call model%final ()
    call syntax_model_file_final ()
    
    write (u, "(A)")
    write (u, "(A)")  "* Test output end: subevt_expr_1"
    
  end subroutine subevt_expr_1
  
\eatline
\nwendcode{}\nwbegindocs{2348}\nwdocspar
\subsubsection{Parton-event expressions}
\nwenddocs{}\nwbegincode{2349}\moddef{Expr tests: execute tests}\plusendmoddef
  call test (subevt_expr_2, "subevt_expr_2", &
       "parton-event expressions", &
       u, results)
\nwendcode{}\nwbegincode{2350}\moddef{Expr tests: tests}\plusendmoddef
  subroutine subevt_expr_2 (u)
    integer, intent(in) :: u
    type(string_t) :: expr_text
    type(ifile_t) :: ifile
    type(stream_t) :: stream
    type(parse_tree_t) :: pt_selection
    type(parse_tree_t) :: pt_reweight, pt_analysis
    type(parse_node_t), pointer :: pn_selection
    type(parse_node_t), pointer :: pn_reweight, pn_analysis
    type(os_data_t) :: os_data
    type(model_t), target :: model
    type(eval_tree_factory_t) :: expr_factory
    type(event_expr_t), target :: expr
    real(default) :: E, Ex, m
    type(vector4_t), dimension(6) :: p
    integer :: i, pdg
    logical :: passed
    real(default) :: reweight
    logical :: analysis_flag
    
    write (u, "(A)")  "* Test output: subevt_expr_2"
    write (u, "(A)")  "*   Purpose: Set up a subevt and associated &
         &process-specific expressions"
    write (u, "(A)")

    call syntax_pexpr_init ()

    call syntax_model_file_init ()
    call os_data_init (os_data)
    call model%read (var_str ("Test.mdl"), os_data)

    write (u, "(A)")  "* Expression texts"
    write (u, "(A)")


    expr_text = "all Pt > 100 [s]"
    write (u, "(A,A)")  "selection = ", char (expr_text)
    call ifile_clear (ifile)
    call ifile_append (ifile, expr_text)
    call stream_init (stream, ifile)
    call parse_tree_init_lexpr (pt_selection, stream, .true.)
    call stream_final (stream)
    pn_selection => parse_tree_get_root_ptr (pt_selection)

    expr_text = "n_tot - n_in - n_out"
    write (u, "(A,A)")  "reweight = ", char (expr_text)
    call ifile_clear (ifile)
    call ifile_append (ifile, expr_text)
    call stream_init (stream, ifile)
    call parse_tree_init_expr (pt_reweight, stream, .true.)
    call stream_final (stream)
    pn_reweight => parse_tree_get_root_ptr (pt_reweight)

    expr_text = "true"
    write (u, "(A,A)")  "analysis = ", char (expr_text)
    call ifile_clear (ifile)
    call ifile_append (ifile, expr_text)
    call stream_init (stream, ifile)
    call parse_tree_init_lexpr (pt_analysis, stream, .true.)
    call stream_final (stream)
    pn_analysis => parse_tree_get_root_ptr (pt_analysis)

    call ifile_final (ifile)

    write (u, "(A)")
    write (u, "(A)")  "* Initialize process expr"
    write (u, "(A)")

    call expr%setup_vars (1000._default)
    call expr%link_var_list (model%get_var_list_ptr ())
    call var_list_append_real (expr%var_list, var_str ("tolerance"), 0._default)

    call expr_factory%init (pn_selection)
    call expr%setup_selection (expr_factory)
    call expr_factory%init (pn_analysis)
    call expr%setup_analysis (expr_factory)
    call expr_factory%init (pn_reweight)
    call expr%setup_reweight (expr_factory)

    call write_separator (u)
    call expr%write (u)
    call write_separator (u)
    
    write (u, "(A)")
    write (u, "(A)")  "* Fill subevt and evaluate expressions"
    write (u, "(A)")

    call subevt_init (expr%subevt_t, 6)
    E = 500._default
    Ex = 400._default
    m = 125._default
    pdg = 25
    p(1) = vector4_moving (E, sqrt (E**2 - m**2), 3)
    p(2) = vector4_moving (E, -sqrt (E**2 - m**2), 3)
    p(3) = vector4_moving (Ex, sqrt (Ex**2 - m**2), 3)
    p(4) = vector4_moving (Ex, -sqrt (Ex**2 - m**2), 3)
    p(5) = vector4_moving (Ex, sqrt (Ex**2 - m**2), 1)
    p(6) = vector4_moving (Ex, -sqrt (Ex**2 - m**2), 1)

    call expr%reset ()
    do i = 1, 2
       call subevt_set_beam (expr%subevt_t, i, pdg, p(i), m**2)
    end do
    do i = 3, 4
       call subevt_set_incoming (expr%subevt_t, i, pdg, p(i), m**2)
    end do
    do i = 5, 6
       call subevt_set_outgoing (expr%subevt_t, i, pdg, p(i), m**2)
    end do
    expr%sqrts_hat = subevt_get_sqrts_hat (expr%subevt_t)
    expr%n_in = 2
    expr%n_out = 2
    expr%n_tot = 4
    expr%subevt_filled = .true.

    call expr%evaluate (passed, reweight, analysis_flag)
    
    write (u, "(A,L1)")      "Event has passed      = ", passed
    write (u, "(A," // FMT_12 // ")")  "Reweighting factor    = ", reweight
    write (u, "(A,L1)")      "Analysis flag         = ", analysis_flag
    write (u, "(A)")
    
    call write_separator (u)
    call expr%write (u)
    call write_separator (u)
    
    write (u, "(A)")
    write (u, "(A)")  "* Cleanup"
    
    call expr%final ()

    call model%final ()
    call syntax_model_file_final ()
    
    write (u, "(A)")
    write (u, "(A)")  "* Test output end: subevt_expr_2"
    
  end subroutine subevt_expr_2
  
\eatline
\nwendcode{}\nwbegindocs{2351}\nwdocspar
\subsubsection{Processes: handle partonic cuts}
Initialize a process and process instance, choose a sampling point and
fill the process instance, evaluating a given cut configuration.

We use the same trivial process as for the previous test.  All
momentum and state dependence is trivial, so we just test basic
functionality. 
\nwenddocs{}\nwbegincode{2352}\moddef{Expr tests: execute tests}\plusendmoddef
  call test (processes_5, "processes_5", &
       "handle cuts (partonic event)", &
       u, results)
\nwendcode{}\nwbegincode{2353}\moddef{Expr tests: tests}\plusendmoddef
  subroutine processes_5 (u)
    integer, intent(in) :: u
    type(string_t) :: cut_expr_text
    type(ifile_t) :: ifile
    type(stream_t) :: stream
    type(parse_tree_t) :: parse_tree
    type(eval_tree_factory_t) :: expr_factory
    type(process_library_t), target :: lib
    type(string_t) :: libname
    type(string_t) :: procname
    type(string_t) :: run_id
    type(os_data_t) :: os_data
    type(qcd_t) :: qcd
    class(rng_factory_t), allocatable :: rng_factory
    type(model_t), pointer :: model_tmp
    class(model_data_t), pointer :: model
    type(var_list_t), target :: var_list
    type(process_t), allocatable, target :: process
    class(prc_core_t), allocatable :: core_template
    class(mci_t), allocatable :: mci_template
    class(phs_config_t), allocatable :: phs_config_template
    real(default) :: sqrts
    type(process_instance_t), allocatable, target :: process_instance

    write (u, "(A)")  "* Test output: processes_5"
    write (u, "(A)")  "*   Purpose: create a process &
         &and fill a process instance"
    write (u, "(A)")

    write (u, "(A)")  "* Prepare a cut expression"
    write (u, "(A)")

    call syntax_pexpr_init ()
    cut_expr_text = "all Pt > 100 [s]"
    call ifile_append (ifile, cut_expr_text)
    call stream_init (stream, ifile)
    call parse_tree_init_lexpr (parse_tree, stream, .true.)
    
    write (u, "(A)")  "* Build and initialize a test process"
    write (u, "(A)")

    libname = "processes5"
    procname = libname
    run_id = "run5"
    call os_data_init (os_data)
    allocate (rng_test_factory_t :: rng_factory)
    call prc_test_create_library (libname, lib)

    call syntax_model_file_init ()
    allocate (model_tmp)
    call model_tmp%read (var_str ("Test.mdl"), os_data)
    call var_list_init_snapshot (var_list, model_tmp%get_var_list_ptr ())
    model => model_tmp

    call reset_interaction_counter ()

    allocate (process)
    call process%init &
         (procname, run_id, lib, os_data, qcd, rng_factory, model) 

    call var_list_append_real &
         (var_list, var_str ("tolerance"), 0._default)
    call process%set_var_list (var_list)
    call var_list%final ()
   
    allocate (test_t :: core_template)
    allocate (phs_test_config_t :: phs_config_template)
    call process%init_component &
         (1, core_template, mci_template, phs_config_template)

    write (u, "(A)")  "* Prepare a trivial beam setup"
    write (u, "(A)")
    
    sqrts = 1000
    call process%setup_beams_sqrts (sqrts)
    call process%configure_phs ()
    call process%setup_mci ()

    write (u, "(A)")  "* Complete process initialization and set cuts"
    write (u, "(A)")

    call process%setup_terms ()
    call expr_factory%init (parse_tree_get_root_ptr (parse_tree))
    call process%set_cuts (expr_factory)
    call process%write (.false., u, show_var_list=.true., show_expressions=.true.)

    write (u, "(A)")
    write (u, "(A)")  "* Create a process instance"
    write (u, "(A)")

    allocate (process_instance)
    call process_instance%init (process)

    write (u, "(A)")
    write (u, "(A)")  "* Inject a set of random numbers"
    write (u, "(A)")
     
    call process_instance%choose_mci (1)
    call process_instance%set_mcpar ([0._default, 0._default])

    write (u, "(A)")
    write (u, "(A)")  "* Set up kinematics and subevt, check cuts (should fail)"
    write (u, "(A)")

    call process_instance%select_channel (1)
    call process_instance%compute_seed_kinematics ()
    call process_instance%compute_hard_kinematics ()
    call process_instance%compute_eff_kinematics ()
    call process_instance%evaluate_expressions ()
    call process_instance%compute_other_channels ()
   
    call process_instance%write (u)

    write (u, "(A)")
    write (u, "(A)")  "* Evaluate for another set (should succeed)"
    write (u, "(A)")
     
    call process_instance%reset ()
    call process_instance%set_mcpar ([0.5_default, 0.125_default])
    call process_instance%select_channel (1)
    call process_instance%compute_seed_kinematics ()
    call process_instance%compute_hard_kinematics ()
    call process_instance%compute_eff_kinematics ()
    call process_instance%evaluate_expressions ()
    call process_instance%compute_other_channels ()
    call process_instance%evaluate_trace ()

    call process_instance%write (u)

    write (u, "(A)")
    write (u, "(A)")  "* Evaluate for another set using convenience procedure &
         &(failure)"
    write (u, "(A)")
     
    call process_instance%evaluate_sqme (1, [0.0_default, 0.2_default])

    call process_instance%write_header (u)

    write (u, "(A)")
    write (u, "(A)")  "* Evaluate for another set using convenience procedure &
         &(success)"
    write (u, "(A)")
     
    call process_instance%evaluate_sqme (1, [0.1_default, 0.2_default])

    call process_instance%write_header (u)

    write (u, "(A)")
    write (u, "(A)")  "* Cleanup"
    
    call process_instance%final ()
    deallocate (process_instance)
    
    call process%final ()
    deallocate (process)

    call parse_tree_final (parse_tree)
    call stream_final (stream)
    call ifile_final (ifile)
    call syntax_pexpr_final ()
    
    call syntax_model_file_final ()

    write (u, "(A)")
    write (u, "(A)")  "* Test output end: processes_5"
    
  end subroutine processes_5
  
\eatline
\nwendcode{}\nwbegindocs{2354}\nwdocspar
\subsubsection{Processes: scales and such}
Initialize a process and process instance, choose a sampling point and
fill the process instance, evaluating a given cut configuration.

We use the same trivial process as for the previous test.  All
momentum and state dependence is trivial, so we just test basic
functionality. 
\nwenddocs{}\nwbegincode{2355}\moddef{Expr tests: execute tests}\plusendmoddef
  call test (processes_6, "processes_6", &
       "handle scales and weight (partonic event)", &
       u, results)
\nwendcode{}\nwbegincode{2356}\moddef{Expr tests: tests}\plusendmoddef
  subroutine processes_6 (u)
    integer, intent(in) :: u
    type(string_t) :: expr_text
    type(ifile_t) :: ifile
    type(stream_t) :: stream
    type(parse_tree_t) :: pt_scale, pt_fac_scale, pt_ren_scale, pt_weight
    type(process_library_t), target :: lib
    type(string_t) :: libname
    type(string_t) :: procname
    type(string_t) :: run_id
    type(os_data_t) :: os_data
    type(qcd_t) :: qcd
    class(rng_factory_t), allocatable :: rng_factory
    type(model_t), pointer :: model_tmp
    class(model_data_t), pointer :: model
    type(var_list_t), target :: var_list
    type(process_t), allocatable, target :: process
    class(prc_core_t), allocatable :: core_template
    class(mci_t), allocatable :: mci_template
    class(phs_config_t), allocatable :: phs_config_template
    real(default) :: sqrts
    type(process_instance_t), allocatable, target :: process_instance
    type(eval_tree_factory_t) :: expr_factory

    write (u, "(A)")  "* Test output: processes_6"
    write (u, "(A)")  "*   Purpose: create a process &
         &and fill a process instance"
    write (u, "(A)")

    write (u, "(A)")  "* Prepare expressions"
    write (u, "(A)")

    call syntax_pexpr_init ()

    expr_text = "sqrts - 100 GeV"
    write (u, "(A,A)")  "scale = ", char (expr_text)
    call ifile_clear (ifile)
    call ifile_append (ifile, expr_text)
    call stream_init (stream, ifile)
    call parse_tree_init_expr (pt_scale, stream, .true.)
    call stream_final (stream)

    expr_text = "sqrts_hat"
    write (u, "(A,A)")  "fac_scale = ", char (expr_text)
    call ifile_clear (ifile)
    call ifile_append (ifile, expr_text)
    call stream_init (stream, ifile)
    call parse_tree_init_expr (pt_fac_scale, stream, .true.)
    call stream_final (stream)
    
    expr_text = "eval sqrt (M2) [collect [s]]"
    write (u, "(A,A)")  "ren_scale = ", char (expr_text)
    call ifile_clear (ifile)
    call ifile_append (ifile, expr_text)
    call stream_init (stream, ifile)
    call parse_tree_init_expr (pt_ren_scale, stream, .true.)
    call stream_final (stream)
    
    expr_text = "n_tot * n_in * n_out * (eval Phi / pi [s])"
    write (u, "(A,A)")  "weight = ", char (expr_text)
    call ifile_clear (ifile)
    call ifile_append (ifile, expr_text)
    call stream_init (stream, ifile)
    call parse_tree_init_expr (pt_weight, stream, .true.)
    call stream_final (stream)

    call ifile_final (ifile)
    
    write (u, "(A)")
    write (u, "(A)")  "* Build and initialize a test process"
    write (u, "(A)")

    libname = "processes4"
    procname = libname
    run_id = "run4"
    call os_data_init (os_data)
    allocate (rng_test_factory_t :: rng_factory)
    call prc_test_create_library (libname, lib)

    call syntax_model_file_init ()
    allocate (model_tmp)
    call model_tmp%read (var_str ("Test.mdl"), os_data)
    call var_list_init_snapshot (var_list, model_tmp%get_var_list_ptr ())
    model => model_tmp

    call reset_interaction_counter ()

    allocate (process)
    call process%init &
         (procname, run_id, lib, os_data, qcd, rng_factory, model) 

    call process%set_var_list (var_list)
    call var_list%final ()
   
    allocate (test_t :: core_template)
    allocate (phs_test_config_t :: phs_config_template)
    call process%init_component &
         (1, core_template, mci_template, phs_config_template)

    write (u, "(A)")  "* Prepare a trivial beam setup"
    write (u, "(A)")
    
    sqrts = 1000
    call process%setup_beams_sqrts (sqrts)
    call process%configure_phs ()
    call process%setup_mci ()

    write (u, "(A)")  "* Complete process initialization and set cuts"
    write (u, "(A)")

    call process%setup_terms ()
    call expr_factory%init (parse_tree_get_root_ptr (pt_scale))
    call process%set_scale (expr_factory)
    call expr_factory%init (parse_tree_get_root_ptr (pt_fac_scale))
    call process%set_fac_scale (expr_factory)
    call expr_factory%init (parse_tree_get_root_ptr (pt_ren_scale))
    call process%set_ren_scale (expr_factory)
    call expr_factory%init (parse_tree_get_root_ptr (pt_weight))
    call process%set_weight (expr_factory)
    call process%write (.false., u, show_expressions=.true.)

    write (u, "(A)")
    write (u, "(A)")  "* Create a process instance and evaluate"
    write (u, "(A)")

    allocate (process_instance)
    call process_instance%init (process)
    call process_instance%choose_mci (1)
    call process_instance%evaluate_sqme (1, [0.5_default, 0.125_default])

    call process_instance%write (u)

    write (u, "(A)")
    write (u, "(A)")  "* Cleanup"
    
    call process_instance%final ()
    deallocate (process_instance)
    
    call process%final ()
    deallocate (process)

    call parse_tree_final (pt_scale)
    call parse_tree_final (pt_fac_scale)
    call parse_tree_final (pt_ren_scale)
    call parse_tree_final (pt_weight)
    call syntax_pexpr_final ()

    call syntax_model_file_final ()
    
    write (u, "(A)")
    write (u, "(A)")  "* Test output end: processes_6"
    
  end subroutine processes_6
  
\eatline
\nwendcode{}\nwbegindocs{2357}\nwdocspar
\subsubsection{Event expressions}
After generating an event, fill the {\tt{}subevt} and evaluate expressions for
selection, reweighting, and analysis.
\nwenddocs{}\nwbegincode{2358}\moddef{Expr tests: execute tests}\plusendmoddef
  call test (events_3, "events_3", &
       "expression evaluation", &
       u, results)
\nwendcode{}\nwbegincode{2359}\moddef{Expr tests: tests}\plusendmoddef
  subroutine events_3 (u)
    integer, intent(in) :: u
    type(string_t) :: expr_text
    type(ifile_t) :: ifile
    type(stream_t) :: stream
    type(parse_tree_t) :: pt_selection, pt_reweight, pt_analysis
    type(eval_tree_factory_t) :: expr_factory
    type(event_t), allocatable, target :: event
    type(process_t), allocatable, target :: process
    type(process_instance_t), allocatable, target :: process_instance
    type(os_data_t) :: os_data
    type(model_t), pointer :: model
    type(var_list_t), target :: var_list

    write (u, "(A)")  "* Test output: events_3"
    write (u, "(A)")  "*   Purpose: generate an event and evaluate expressions"
    write (u, "(A)")

    call syntax_pexpr_init ()

    write (u, "(A)")  "* Expression texts"
    write (u, "(A)")

    expr_text = "all Pt > 100 [s]"
    write (u, "(A,A)")  "selection = ", char (expr_text)
    call ifile_clear (ifile)
    call ifile_append (ifile, expr_text)
    call stream_init (stream, ifile)
    call parse_tree_init_lexpr (pt_selection, stream, .true.)
    call stream_final (stream)

    expr_text = "1 + sqrts_hat / sqrts"
    write (u, "(A,A)")  "reweight = ", char (expr_text)
    call ifile_clear (ifile)
    call ifile_append (ifile, expr_text)
    call stream_init (stream, ifile)
    call parse_tree_init_expr (pt_reweight, stream, .true.)
    call stream_final (stream)
    
    expr_text = "true"
    write (u, "(A,A)")  "analysis = ", char (expr_text)
    call ifile_clear (ifile)
    call ifile_append (ifile, expr_text)
    call stream_init (stream, ifile)
    call parse_tree_init_lexpr (pt_analysis, stream, .true.)
    call stream_final (stream)

    call ifile_final (ifile)
    
    write (u, "(A)")
    write (u, "(A)")  "* Initialize test process event"

    call os_data_init (os_data)

    call syntax_model_file_init ()
    allocate (model)
    call model%read (var_str ("Test.mdl"), os_data)
    call var_list_init_snapshot (var_list, model%get_var_list_ptr ())

    allocate (process)
    allocate (process_instance)
    call prepare_test_process (process, process_instance, model)

    call process%set_var_list (var_list)
    call var_list%final ()

    call process_instance%setup_event_data ()

    write (u, "(A)")
    write (u, "(A)")  "* Initialize event object and set expressions"

    allocate (event)
    call event%basic_init ()

    call expr_factory%init (parse_tree_get_root_ptr (pt_selection))
    call event%set_selection (expr_factory)
    call expr_factory%init (parse_tree_get_root_ptr (pt_reweight))
    call event%set_reweight (expr_factory)
    call expr_factory%init (parse_tree_get_root_ptr (pt_analysis))
    call event%set_analysis (expr_factory)
    
    call event%connect (process_instance, process%get_model_ptr ())
    call var_list_append_real &
         (event%expr%var_list, var_str ("tolerance"), 0._default)
    call event%setup_expressions ()

    write (u, "(A)")
    write (u, "(A)")  "* Generate test process event"

    call process%generate_weighted_event (process_instance, 1)

    write (u, "(A)")
    write (u, "(A)")  "* Fill event object and evaluate expressions"
    write (u, "(A)")

    call event%generate (1, [0.4_default, 0.4_default])
    call event%evaluate_expressions ()
    call event%write (u)

    write (u, "(A)")
    write (u, "(A)")  "* Cleanup"

    call event%final ()
    deallocate (event)

    call cleanup_test_process (process, process_instance)
    deallocate (process_instance)
    deallocate (process)

    call syntax_model_file_final ()
    
    write (u, "(A)")
    write (u, "(A)")  "* Test output end: events_3"
    
  end subroutine events_3
  
\eatline
\nwendcode{}\nwbegindocs{2360}\nwdocspar
%%%%%%%%%%%%%%%%%%%%%%%%%%%%%%%%%%%%%%%%%%%%%%%%%%%%%%%%%%%%%%%%%%%%%%%%
\chapter{Top Level}

The top level consists of
\begin{description}
\item[commands]
  Defines generic command-list and command objects, and all specific
  implementations.  Each command type provides a specific
  functionality.  Together with the modules that provide expressions
  and variables, this module defines the Sindarin language.
\item[whizard]
  This module interprets streams of various kind in terms of the
  command language.  It also contains the unit-test feature.  We also
  define the externally visible procedures here, for the \whizard\ as
  a library.
\item[main]
  The driver for \whizard\ as a stand-alone program.  Contains the
  command-line interpreter.
\item[whizard\_c\_interface]
  Alternative top-level procedures, for use in the context of a
  C-compatible caller program.
\end{description}

\clearpage
%%%%%%%%%%%%%%%%%%%%%%%%%%%%%%%%%%%%%%%%%%%%%%%%%%%%%%%%%%%%%%%%%%%%%%%%
\section{Commands}
This module defines the command language of the main input file.
\nwenddocs{}\nwbegincode{2361}\moddef{\code{}commands.f90\edoc{}}\endmoddef
\LA{}File header\RA{}

module commands

\LA{}Use kinds\RA{}
\LA{}Use strings\RA{}
  use io_units
  use constants
  use string_utils, only: lower_case
  use format_utils, only: write_indent
  use format_defs, only: FMT_14, FMT_19
  use unit_tests
  use diagnostics
  use sm_qcd
  use pdf_builtin !NODEP!
  use sorting
  use sf_lhapdf
  use os_interface
  use ifiles
  use lexers
  use syntax_rules
  use parser
  use analysis
  use pdg_arrays
  use variables
  use observables
  use eval_trees
  use models
  use auto_components
  use interactions
  use flavors
  use polarizations
  use beams
  use particle_specifiers
  use process_libraries
  use processes
  use prclib_stacks
  use slha_interface
  use user_files
  use eio_data
  use rt_data
  use dispatch
  use process_configurations
  use compilations
  use integrations
  use event_streams
  use simulations
  use radiation_generator
  use blha_config

\LA{}Standard module head\RA{}

\LA{}Commands: public\RA{}

\LA{}Commands: types\RA{}

\LA{}Commands: variables\RA{}

\LA{}Commands: parameters\RA{}

\LA{}Commands: interfaces\RA{}

contains

\LA{}Commands: procedures\RA{}

\LA{}Commands: tests\RA{}

end module commands
\eatline
\nwendcode{}\nwbegindocs{2362}\nwdocspar
\subsection{The command type}
The command type is a generic type that holds any command, compiled
for execution.

Each command may come with its own local environment.  The command list that
determines this environment is allocated as {\tt{}options}, if necessary.  (It
has to be allocated as a pointer because the type definition is recursive.) The
local environment is available as a pointer which either points to the global
environment, or is explicitly allocated and initialized.
\nwenddocs{}\nwbegincode{2363}\moddef{Commands: types}\endmoddef
  type, abstract :: command_t
     type(parse_node_t), pointer :: pn => null ()
     class(command_t), pointer :: next => null ()
     type(parse_node_t), pointer :: pn_opt => null ()
     type(command_list_t), pointer :: options => null ()
     type(rt_data_t), pointer :: local => null ()
   contains
   \LA{}Commands: command: TBP\RA{}
  end type command_t

\eatline
\nwendcode{}\nwbegindocs{2364}Finalizer: If there is an option list, finalize the option list and
deallocate.  If not, the local environment is just a pointer.
\nwenddocs{}\nwbegincode{2365}\moddef{Commands: command: TBP}\endmoddef
  procedure :: final => command_final
\nwendcode{}\nwbegincode{2366}\moddef{Commands: procedures}\endmoddef
  recursive subroutine command_final (cmd)
    class(command_t), intent(inout) :: cmd
    if (associated (cmd%options)) then
       call cmd%options%final ()
       deallocate (cmd%options)
       call cmd%local%local_final ()
       deallocate (cmd%local)
    else
       cmd%local => null ()
    end if
  end subroutine command_final

\eatline
\nwendcode{}\nwbegindocs{2367}Allocate a command with the appropriate concrete type.  Store the
parse node pointer in the command object, so we can reference to it
when compiling.
\nwenddocs{}\nwbegincode{2368}\moddef{Commands: procedures}\plusendmoddef
  subroutine dispatch_command (command, pn)
    class(command_t), intent(inout), pointer :: command
    type(parse_node_t), intent(in), target :: pn
    select case (char (parse_node_get_rule_key (pn)))
    case ("cmd_model")
       allocate (cmd_model_t :: command)
    case ("cmd_library")
       allocate (cmd_library_t :: command)
    case ("cmd_process")
       allocate (cmd_process_t :: command)
    case ("cmd_nlo")
       allocate (cmd_nlo_t :: command)
    case ("cmd_compile")
       allocate (cmd_compile_t :: command)
    case ("cmd_exec")
       allocate (cmd_exec_t :: command)
     case ("cmd_num", "cmd_complex", "cmd_real", "cmd_int", &
           "cmd_log_decl", "cmd_log", "cmd_string", "cmd_string_decl", &
           "cmd_alias", "cmd_result")
       allocate (cmd_var_t :: command)
    case ("cmd_slha")
       allocate (cmd_slha_t :: command)
    case ("cmd_show")
       allocate (cmd_show_t :: command)
    case ("cmd_clear")
       allocate (cmd_clear_t :: command)
    case ("cmd_expect")
       allocate (cmd_expect_t :: command)
    case ("cmd_beams")
       allocate (cmd_beams_t :: command)
    case ("cmd_beams_pol_density")
       allocate (cmd_beams_pol_density_t :: command)
    case ("cmd_beams_pol_fraction")
       allocate (cmd_beams_pol_fraction_t :: command)
    case ("cmd_beams_momentum")
       allocate (cmd_beams_momentum_t :: command)
    case ("cmd_beams_theta")
       allocate (cmd_beams_theta_t :: command)
    case ("cmd_beams_phi")
       allocate (cmd_beams_phi_t :: command)
    case ("cmd_cuts")
       allocate (cmd_cuts_t :: command)
    case ("cmd_scale")
       allocate (cmd_scale_t :: command)
    case ("cmd_fac_scale")
       allocate (cmd_fac_scale_t :: command)
    case ("cmd_ren_scale")
       allocate (cmd_ren_scale_t :: command)
    case ("cmd_weight")
       allocate (cmd_weight_t :: command)
    case ("cmd_selection")
       allocate (cmd_selection_t :: command)
    case ("cmd_reweight")
       allocate (cmd_reweight_t :: command)
    case ("cmd_iterations")
       allocate (cmd_iterations_t :: command)
    case ("cmd_integrate")
       allocate (cmd_integrate_t :: command)
    case ("cmd_observable")
       allocate (cmd_observable_t :: command)
    case ("cmd_histogram")
       allocate (cmd_histogram_t :: command)
    case ("cmd_plot")
       allocate (cmd_plot_t :: command)
    case ("cmd_graph")
       allocate (cmd_graph_t :: command)
    case ("cmd_record")
       allocate (cmd_record_t :: command)
    case ("cmd_analysis")
       allocate (cmd_analysis_t :: command)
    case ("cmd_alt_setup")
       allocate (cmd_alt_setup_t :: command)
    case ("cmd_unstable")
       allocate (cmd_unstable_t :: command)
    case ("cmd_stable")
       allocate (cmd_stable_t :: command)
    case ("cmd_polarized")
       allocate (cmd_polarized_t :: command)
    case ("cmd_unpolarized")
       allocate (cmd_unpolarized_t :: command)
    case ("cmd_sample_format")
       allocate (cmd_sample_format_t :: command)
    case ("cmd_simulate")
       allocate (cmd_simulate_t :: command)
    case ("cmd_rescan")
       allocate (cmd_rescan_t :: command)
    case ("cmd_write_analysis")
       allocate (cmd_write_analysis_t :: command)
    case ("cmd_compile_analysis")
       allocate (cmd_compile_analysis_t :: command)
    case ("cmd_open_out")
       allocate (cmd_open_out_t :: command)
    case ("cmd_close_out")
       allocate (cmd_close_out_t :: command)
    case ("cmd_printf")
       allocate (cmd_printf_t :: command)
    case ("cmd_scan")
       allocate (cmd_scan_t :: command)
    case ("cmd_if")
       allocate (cmd_if_t :: command)
    case ("cmd_include")
       allocate (cmd_include_t :: command)
    case ("cmd_quit")
       allocate (cmd_quit_t :: command)
    case default
       print *, char (parse_node_get_rule_key (pn))
       call msg_bug ("Command not implemented")
    end select
    command%pn => pn
  end subroutine dispatch_command

\eatline
\nwendcode{}\nwbegindocs{2369}Output.  We allow for indentation so we can display a command tree.
\nwenddocs{}\nwbegincode{2370}\moddef{Commands: command: TBP}\plusendmoddef
  procedure (command_write), deferred :: write
\nwendcode{}\nwbegincode{2371}\moddef{Commands: interfaces}\endmoddef
  abstract interface
     subroutine command_write (cmd, unit, indent)
       import
       class(command_t), intent(in) :: cmd
       integer, intent(in), optional :: unit, indent
     end subroutine command_write
  end interface
  
\eatline
\nwendcode{}\nwbegindocs{2372}Compile a command.  The command type is already fixed, so this is a
deferred type-bound procedure.
\nwenddocs{}\nwbegincode{2373}\moddef{Commands: command: TBP}\plusendmoddef
  procedure (command_compile), deferred :: compile
\nwendcode{}\nwbegincode{2374}\moddef{Commands: interfaces}\plusendmoddef
  abstract interface
     subroutine command_compile (cmd, global)
       import
       class(command_t), intent(inout) :: cmd
       type(rt_data_t), intent(inout), target :: global
     end subroutine command_compile
  end interface

\eatline
\nwendcode{}\nwbegindocs{2375}Execute a command.  This will use and/or modify the runtime data
set.  If the {\tt{}quit} flag is set, the caller should terminate command
execution.
\nwenddocs{}\nwbegincode{2376}\moddef{Commands: command: TBP}\plusendmoddef
  procedure (command_execute), deferred :: execute
\nwendcode{}\nwbegincode{2377}\moddef{Commands: interfaces}\plusendmoddef
  abstract interface
     subroutine command_execute (cmd, global)
       import
       class(command_t), intent(inout) :: cmd
       type(rt_data_t), intent(inout), target :: global
     end subroutine command_execute
  end interface

\eatline
\nwendcode{}\nwbegindocs{2378}\nwdocspar
\subsection{Options}
The {\tt{}options} command list is allocated, initialized, and executed, if the
command is associated with an option text in curly braces.  If present, a
separate local runtime data set {\tt{}local} will be allocated and initialized;
otherwise, {\tt{}local} becomes a pointer to the global dataset.

For output, we indent the options list.
\nwenddocs{}\nwbegincode{2379}\moddef{Commands: command: TBP}\plusendmoddef
  procedure :: write_options => command_write_options
\nwendcode{}\nwbegincode{2380}\moddef{Commands: procedures}\plusendmoddef
  recursive subroutine command_write_options (cmd, unit, indent)
    class(command_t), intent(in) :: cmd
    integer, intent(in), optional :: unit, indent
    integer :: ind
    ind = 1;  if (present (indent))  ind = indent + 1
    if (associated (cmd%options))  call cmd%options%write (unit, ind)
  end subroutine command_write_options
  
\eatline
\nwendcode{}\nwbegindocs{2381}Compile the options list, if any.  This implies initialization of the local
environment.  Should be done once the {\tt{}pn{\char95}opt} node has been assigned (if
applicable), but before the actual command compilation.
\nwenddocs{}\nwbegincode{2382}\moddef{Commands: command: TBP}\plusendmoddef
  procedure :: compile_options => command_compile_options
\nwendcode{}\nwbegincode{2383}\moddef{Commands: procedures}\plusendmoddef
  recursive subroutine command_compile_options (cmd, global)
    class(command_t), intent(inout) :: cmd
    type(rt_data_t), intent(inout), target :: global
    if (associated (cmd%pn_opt)) then
       allocate (cmd%local)
       call cmd%local%local_init (global)
       call global%copy_globals (cmd%local)
       allocate (cmd%options)
       call cmd%options%compile (cmd%pn_opt, cmd%local)
       call global%restore_globals (cmd%local)
       call cmd%local%deactivate ()
    else
       cmd%local => global
    end if
  end subroutine command_compile_options
  
\eatline
\nwendcode{}\nwbegindocs{2384}Execute options.  First prepare the local environment, then execute the
command list. 
\nwenddocs{}\nwbegincode{2385}\moddef{Commands: command: TBP}\plusendmoddef
  procedure :: execute_options => cmd_execute_options
\nwendcode{}\nwbegincode{2386}\moddef{Commands: procedures}\plusendmoddef
  recursive subroutine cmd_execute_options (cmd, global)
    class(command_t), intent(inout) :: cmd
    type(rt_data_t), intent(inout), target :: global
    if (associated (cmd%options)) then
       call cmd%local%activate ()
       call cmd%options%execute (cmd%local)
    end if
  end subroutine cmd_execute_options

\eatline
\nwendcode{}\nwbegindocs{2387}This must be called after the parent command has been executed, to undo
temporary modifications to the environment.  Note that some modifications to
{\tt{}global} can become permanent.
\nwenddocs{}\nwbegincode{2388}\moddef{Commands: command: TBP}\plusendmoddef
  procedure :: reset_options => cmd_reset_options
\nwendcode{}\nwbegincode{2389}\moddef{Commands: procedures}\plusendmoddef
  subroutine cmd_reset_options (cmd, global)
    class(command_t), intent(inout) :: cmd
    type(rt_data_t), intent(inout), target :: global
    if (associated (cmd%options)) then
       call cmd%local%deactivate (global)
    end if
  end subroutine cmd_reset_options
  
\eatline
\nwendcode{}\nwbegindocs{2390}\nwdocspar
\subsection{Specific command types}
\subsubsection{Model configuration}
The command declares a model, looks for the specified file and loads
it.
\nwenddocs{}\nwbegincode{2391}\moddef{Commands: types}\plusendmoddef
  type, extends (command_t) :: cmd_model_t
     private
     type(string_t) :: name
   contains
   \LA{}Commands: cmd model: TBP\RA{}
  end type cmd_model_t

\eatline
\nwendcode{}\nwbegindocs{2392}Output
\nwenddocs{}\nwbegincode{2393}\moddef{Commands: cmd model: TBP}\endmoddef
  procedure :: write => cmd_model_write
\nwendcode{}\nwbegincode{2394}\moddef{Commands: procedures}\plusendmoddef
  subroutine cmd_model_write (cmd, unit, indent)
    class(cmd_model_t), intent(in) :: cmd
    integer, intent(in), optional :: unit, indent
    integer :: u
    u = given_output_unit (unit);  if (u < 0)  return
    call write_indent (u, indent)
    write (u, "(1x,A,1x,'""',A,'""')")  "model =", char (cmd%name)
  end subroutine cmd_model_write

\eatline
\nwendcode{}\nwbegindocs{2395}Compile.  Get the model name and read the model from file, so it is
readily available when the command list is executed.

Assign the model pointer in the {\tt{}global} record, so it can be used for
(read-only) variable lookup while compiling further commands.
\nwenddocs{}\nwbegincode{2396}\moddef{Commands: cmd model: TBP}\plusendmoddef
  procedure :: compile => cmd_model_compile
\nwendcode{}\nwbegincode{2397}\moddef{Commands: procedures}\plusendmoddef
  subroutine cmd_model_compile (cmd, global)
    class(cmd_model_t), intent(inout) :: cmd
    type(rt_data_t), intent(inout), target :: global
    type(parse_node_t), pointer :: pn_name
    type(model_t), pointer :: model
    pn_name => parse_node_get_sub_ptr (cmd%pn, 3)
    cmd%name = parse_node_get_string (pn_name)
    model => null ()
    if (associated (global%model)) then
       if (global%model%get_name () == cmd%name)  model => global%model
    end if
    if (.not. associated (model)) then
       if (global%model_list%model_exists (cmd%name)) then
          model => global%model_list%get_model_ptr (cmd%name)
       else
          call global%read_model (cmd%name, model)
       end if
    end if
    global%model => model
    if (associated (global%model)) then
       call global%model%link_var_list (global%var_list)
    end if
  end subroutine cmd_model_compile

\eatline
\nwendcode{}\nwbegindocs{2398}Execute: Insert a pointer into the global data record and reassign
the variable list.
\nwenddocs{}\nwbegincode{2399}\moddef{Commands: cmd model: TBP}\plusendmoddef
  procedure :: execute => cmd_model_execute
\nwendcode{}\nwbegincode{2400}\moddef{Commands: procedures}\plusendmoddef
  subroutine cmd_model_execute (cmd, global)
    class(cmd_model_t), intent(inout) :: cmd
    type(rt_data_t), intent(inout), target :: global
    call global%select_model (cmd%name)
    if (.not. associated (global%model)) &
         call msg_fatal ("Switching to model '" &
         // char (cmd%name) // "': model not found")
  end subroutine cmd_model_execute

\eatline
\nwendcode{}\nwbegindocs{2401}\nwdocspar
\subsubsection{Library configuration}
We configure a process library that should hold the subsequently
defined processes.  If the referenced library exists already, just
make it the currently active one.
\nwenddocs{}\nwbegincode{2402}\moddef{Commands: types}\plusendmoddef
  type, extends (command_t) :: cmd_library_t
     private
     type(string_t) :: name
   contains
   \LA{}Commands: cmd library: TBP\RA{}
  end type cmd_library_t

\eatline
\nwendcode{}\nwbegindocs{2403}Output.
\nwenddocs{}\nwbegincode{2404}\moddef{Commands: cmd library: TBP}\endmoddef
  procedure :: write => cmd_library_write
\nwendcode{}\nwbegincode{2405}\moddef{Commands: procedures}\plusendmoddef
  subroutine cmd_library_write (cmd, unit, indent)
    class(cmd_library_t), intent(in) :: cmd
    integer, intent(in), optional :: unit, indent
    integer :: u
    u = given_output_unit (unit)
    call write_indent (u, indent)
    write (u, "(1x,A,1x,'""',A,'""')")  "library =", char (cmd%name)
  end subroutine cmd_library_write

\eatline
\nwendcode{}\nwbegindocs{2406}Compile.  Get the library name.
\nwenddocs{}\nwbegincode{2407}\moddef{Commands: cmd library: TBP}\plusendmoddef
  procedure :: compile => cmd_library_compile
\nwendcode{}\nwbegincode{2408}\moddef{Commands: procedures}\plusendmoddef
  subroutine cmd_library_compile (cmd, global)
    class(cmd_library_t), intent(inout) :: cmd
    type(rt_data_t), intent(inout), target :: global
    type(parse_node_t), pointer :: pn_name
    pn_name => parse_node_get_sub_ptr (cmd%pn, 3)
    cmd%name = parse_node_get_string (pn_name)
  end subroutine cmd_library_compile

\eatline
\nwendcode{}\nwbegindocs{2409}Execute: Initialize a new library and push it on the library stack
(if it does not yet exist).   Insert a pointer to the library into the
global data record.  Then, try to load the library unless the
{\tt{}rebuild} flag is set.
\nwenddocs{}\nwbegincode{2410}\moddef{Commands: cmd library: TBP}\plusendmoddef
  procedure :: execute => cmd_library_execute
\nwendcode{}\nwbegincode{2411}\moddef{Commands: procedures}\plusendmoddef
  subroutine cmd_library_execute (cmd, global)
    class(cmd_library_t), intent(inout) :: cmd
    type(rt_data_t), intent(inout), target :: global
    type(prclib_entry_t), pointer :: lib_entry
    type(process_library_t), pointer :: lib
    logical :: rebuild_library
    lib => global%prclib_stack%get_library_ptr (cmd%name)
    rebuild_library = &
         global%var_list%get_lval (var_str ("?rebuild_library"))    
    if (.not. (associated (lib))) then
       allocate (lib_entry)
       call lib_entry%init (cmd%name)
       lib => lib_entry%process_library_t
       call global%add_prclib (lib_entry)
    else
       call global%update_prclib (lib)
    end if
    if (associated (lib) .and. .not. rebuild_library) then
       call lib%update_status (global%os_data)
    end if
  end subroutine cmd_library_execute

\eatline
\nwendcode{}\nwbegindocs{2412}\nwdocspar
\subsubsection{Process configuration}
We define a process-configuration command as a specific type.  The
incoming and outgoing particles are given evaluation-trees which we
transform to PDG-code arrays.  For transferring to \oMega, they are
reconverted to strings.

For the incoming particles, we store parse nodes individually.  We do
not yet resolve the outgoing state, so we store just a single parse
node. 

This also includes the choice of method for the corresponding process:
{\tt{}omega} for \oMega\ matrix elements as Fortran code, {\tt{}ovm} for
\oMega\ matrix elements as a bytecode virtual machine, {\tt{}test} for
special processes, {\tt{}unit{\char95}test} for internal test matrix elements
generated by \whizard, {\tt{}template} and {\tt{}template{\char95}unity} for test
matrix elements generated by \whizard\ as Fortran code similar to the 
\oMega\ code. If the one-loop program (OLP) \gosam\ is linked, also
matrix elements from there (at leading and next-to-leading order) can
be generated via {\tt{}gosam}.
\nwenddocs{}\nwbegincode{2413}\moddef{Commands: types}\plusendmoddef
  type, extends (command_t) :: cmd_process_t
     private
     type(string_t) :: id
     integer :: n_in  = 0
     type(parse_node_p), dimension(:), allocatable :: pn_pdg_in
     type(parse_node_t), pointer :: pn_out => null ()
   contains
   \LA{}Commands: cmd process: TBP\RA{}  
  end type cmd_process_t

\eatline
\nwendcode{}\nwbegindocs{2414}Output.  The particle expressions are not resolved, so we just list the
number of incoming particles.
\nwenddocs{}\nwbegincode{2415}\moddef{Commands: cmd process: TBP}\endmoddef
  procedure :: write => cmd_process_write
\nwendcode{}\nwbegincode{2416}\moddef{Commands: procedures}\plusendmoddef
  subroutine cmd_process_write (cmd, unit, indent)
    class(cmd_process_t), intent(in) :: cmd
    integer, intent(in), optional :: unit, indent
    integer :: u
    u = given_output_unit (unit);  if (u < 0)  return
    call write_indent (u, indent)
    write (u, "(1x,A,A,A,I0,A)")  "process: ", char (cmd%id), " (", &
         size (cmd%pn_pdg_in), " -> X)"
    call cmd%write_options (u, indent)
  end subroutine cmd_process_write

\eatline
\nwendcode{}\nwbegindocs{2417}Compile.  Find and assign the parse nodes.
\nwenddocs{}\nwbegincode{2418}\moddef{Commands: cmd process: TBP}\plusendmoddef
  procedure :: compile => cmd_process_compile
\nwendcode{}\nwbegincode{2419}\moddef{Commands: procedures}\plusendmoddef
  subroutine cmd_process_compile (cmd, global)
    class(cmd_process_t), intent(inout) :: cmd
    type(rt_data_t), intent(inout), target :: global
    type(parse_node_t), pointer :: pn_id, pn_in, pn_codes
    integer :: i
    pn_id => parse_node_get_sub_ptr (cmd%pn, 2)
    pn_in  => parse_node_get_next_ptr (pn_id, 2)
    cmd%pn_out => parse_node_get_next_ptr (pn_in, 2)
    cmd%pn_opt => parse_node_get_next_ptr (cmd%pn_out)
    call cmd%compile_options (global)
    cmd%id = parse_node_get_string (pn_id)
    cmd%n_in  = parse_node_get_n_sub (pn_in)
    pn_codes => parse_node_get_sub_ptr (pn_in)
    allocate (cmd%pn_pdg_in (cmd%n_in))
    do i = 1, cmd%n_in
       cmd%pn_pdg_in(i)%ptr => pn_codes
       pn_codes => parse_node_get_next_ptr (pn_codes)
    end do
  end subroutine cmd_process_compile

\eatline
\nwendcode{}\nwbegindocs{2420}Command execution.  Evaluate the subevents, transform PDG codes
into strings, and add the current process configuration to the
process library. 

The initial state will be unique (one or two particles).  For the final state,
we allow for expressions.  The expressions will be expanded until we have a
sum of final states.  Each distinct final state will get its own process
component.

To identify equivalent final states, we transform the final state into an array
of PDG codes, which we sort and compare.  If a particle entry is actually a
PDG array, only the first entry in the array is used for the comparison.
The user should make sure that there is no overlap between different particles
or arrays which would make the expansion ambiguous.
\nwenddocs{}\nwbegincode{2421}\moddef{Commands: cmd process: TBP}\plusendmoddef
  procedure :: execute => cmd_process_execute
\nwendcode{}\nwbegincode{2422}\moddef{Commands: procedures}\plusendmoddef
  subroutine cmd_process_execute (cmd, global)
    class(cmd_process_t), intent(inout) :: cmd
    type(rt_data_t), intent(inout), target :: global
    type(pdg_array_t) :: pdg_in, pdg_out
    type(pdg_array_t), dimension(:), allocatable :: pdg_out_tab
    type(string_t), dimension(:), allocatable :: prt_in
    type(string_t) :: prt_out, prt_out1
    type(process_configuration_t) :: prc_config
    type(prt_expr_t) :: prt_expr_out
    type(prt_spec_t), dimension(:), allocatable :: prt_spec_in
    type(prt_spec_t), dimension(:), allocatable :: prt_spec_out
    type(var_list_t), pointer :: var_list
    integer, dimension(:), allocatable :: pdg
    integer, dimension(:), allocatable :: i_term
    integer :: i, j, n_in, n_out, n_terms, n_components
    logical :: nlo_calc
    logical, dimension(3) :: active_nlo_components
    type(pdg_array_t), dimension(:), allocatable :: pdg_array_in, pdg_array_out
    type(string_t), dimension(:), allocatable :: prt_in_nlo, prt_out_nlo
    type(radiation_generator_t) :: radiation_generator
    type(pdg_list_t) :: pl_in, pl_out
    logical :: method_changed = .false.
    logical :: use_gosam_loops
    logical :: use_gosam_correlations
    logical :: use_gosam_real_trees
    integer , dimension(4) :: i_list
    logical :: combined_nlo_integration
    
    nlo_calc = cmd%local%nlo_calculation
    active_nlo_components = cmd%local%active_nlo_components

    var_list => cmd%local%get_var_list_ptr ()

    n_in = size (cmd%pn_pdg_in)
    allocate (prt_in (n_in), prt_spec_in (n_in))
    do i = 1, n_in
       pdg_in = &
            eval_pdg_array (cmd%pn_pdg_in(i)%ptr, var_list)
       prt_in(i) = make_flavor_string (pdg_in, cmd%local%model)
       prt_spec_in(i) = new_prt_spec (prt_in(i))
    end do
    call compile_prt_expr &
         (prt_expr_out, cmd%pn_out, var_list, cmd%local%model)
    call prt_expr_out%expand ()
    n_terms = prt_expr_out%get_n_terms ()
    allocate (pdg_out_tab (n_terms))
    allocate (i_term (n_terms), source = 0)
    n_components = 0
    SCAN_COMPONENTS: do i = 1, n_terms
       if (allocated (pdg))  deallocate (pdg)
       call prt_expr_out%term_to_array (prt_spec_out, i)
       n_out = size (prt_spec_out)
       allocate (pdg (n_out))
       do j = 1, n_out
          prt_out = prt_spec_out(j)%to_string ()
          call split (prt_out, prt_out1, ":")
          pdg(j) = cmd%local%model%get_pdg (prt_out1)
       end do
       pdg_out = sort (pdg)
       do j = 1, n_components
          if (pdg_out == pdg_out_tab(j))  cycle SCAN_COMPONENTS
       end do
       n_components = n_components + 1
       i_term(n_components) = i
       pdg_out_tab(n_components) = pdg_out
    end do SCAN_COMPONENTS
    if (nlo_calc) then
      call prc_config%init (cmd%id, n_in, n_components*4, cmd%local)
    else
      call prc_config%init (cmd%id, n_in, n_components, cmd%local)
    end if
    do i = 1, n_components
       call prt_expr_out%term_to_array (prt_spec_out, i_term(i))
       if (nlo_calc) then
         associate (active_comp => cmd%local%active_nlo_components)
            i_list(1) = i
            i_list(2) = i + n_components
            i_list(3) = i + 2*n_components
            i_list(4) = i + 3*n_components
            use_gosam_loops = &
                 global%var_list%get_lval (var_str ('?use_gosam_loops'))
            use_gosam_correlations = &
                 global%var_list%get_lval (var_str ('?use_gosam_correlations'))
            use_gosam_real_trees = &
                 global%var_list%get_lval (var_str ('?use_gosam_real_trees'))

            call prc_config%setup_component (i, prt_spec_in, prt_spec_out, &
                                             cmd%local, var_str ('Born'), &
                                             active_in = active_comp (1))
            call split_prt (prt_spec_in, n_in, pl_in)
            call split_prt (prt_spec_out, n_out, pl_out)
            call radiation_generator_init (radiation_generator, .true., .false., &
                                           pl_in, pl_out)
            call radiation_generator%set_n (n_in, n_out, 0)
            call radiation_generator%set_constraints (.false., .false., .true., .true.)
            call radiation_generator%init_radiation_model &
                 (cmd%local%radiation_model)
            call radiation_generator%generate (prt_in_nlo, prt_out_nlo)

            if (use_gosam_real_trees) then
               if (.not. method_changed) &
                 call global%change_to_gosam (method_changed)
            end if
  
            call prc_config%setup_component (n_components + i, &
                            new_prt_spec (prt_in_nlo), &
                            new_prt_spec (prt_out_nlo),&
                            cmd%local, var_str ('Real'), &
                            active_in = active_comp (2))

            if (use_gosam_loops .and..not. method_changed) &
                 call global%change_to_gosam (method_changed)

            call prc_config%setup_component (n_components*2 + i, prt_spec_in, &
                            prt_spec_out, global, var_str ('Virtual'), &
                            active_in = active_comp (3))

            if (.not. use_gosam_correlations .and. method_changed) then
               call global%change_to_omega ()
            else if (use_gosam_correlations .and..not. method_changed) then
               call global%change_to_gosam (method_changed)
            end if

            call prc_config%setup_component (n_components*3 + i, prt_spec_in, &
                            prt_spec_out, global, var_str ('Subtraction'), &
                            .false.)                     
            call prc_config%set_component_associations (i_list)
         end associate
       else
         call prc_config%setup_component (i, prt_spec_in, prt_spec_out, cmd%local)
       end if
    end do
    call prc_config%record (cmd%local)
 
  contains
    subroutine split_prt (prt, n_out, pl)
      type(prt_spec_t), intent(in), dimension(:), allocatable :: prt
      integer, intent(in) :: n_out
      type(pdg_list_t), intent(out) :: pl
      type(pdg_array_t) :: pdg
      type(string_t) :: prt_string, prt_tmp
      integer, dimension(10) :: i_particle
      integer :: i, j, n
      call pl%init (n_out)
      do i = 1, n_out
         n = 1
         prt_string = prt(i)%to_string ()
         do
           call split (prt_string, prt_tmp, ":")
           if (prt_tmp /= "") then
             i_particle(n) = cmd%local%model%get_pdg (prt_tmp)
             n=n+1
           else
             exit
           end if
         end do
         call pdg_array_init (pdg, n-1)
         do j = 1, n-1 
           call pdg%set (j, i_particle(j))
         end do
         call pl%set (i, pdg)
         call pdg_array_delete (pdg)
      end do
    end subroutine split_prt
             
  end subroutine cmd_process_execute

\eatline
\nwendcode{}\nwbegindocs{2423}This is a method of the eval tree, but cannot be coded inside the
{\tt{}expressions} module since it uses the {\tt{}model} and {\tt{}flv} types
which are not available there.
\nwenddocs{}\nwbegincode{2424}\moddef{Commands: procedures}\plusendmoddef
  function make_flavor_string (aval, model) result (prt)
    type(string_t) :: prt
    type(pdg_array_t), intent(in) :: aval
    type(model_t), intent(in), target :: model
    integer, dimension(:), allocatable :: pdg
    type(flavor_t), dimension(:), allocatable :: flv
    integer :: i
    pdg = aval
    allocate (flv (size (pdg)))
    call flv%init (pdg, model)
    if (size (pdg) /= 0) then
       prt = flv(1)%get_name ()
       do i = 2, size (flv)
          prt = prt // ":" // flv(i)%get_name ()
       end do
    else
       prt = "?"
    end if
  end function make_flavor_string

\eatline
\nwendcode{}\nwbegindocs{2425}Create a pdg array from a particle-specification array
\nwenddocs{}\nwbegincode{2426}\moddef{Commands: procedures}\plusendmoddef
  function make_pdg_array (prt, model) result (pdg_array)
    type(prt_spec_t), intent(in), dimension(:) :: prt
    type(model_t), intent(in) :: model
    integer, dimension(:), allocatable :: aval
    type(pdg_array_t) :: pdg_array
    type(flavor_t) :: flv
    integer :: k
    allocate (aval (size (prt)))
    do k = 1, size (prt)
      call flv%init (prt(k)%to_string (), model)
      aval (k) = flv%get_pdg ()
    end do
    pdg_array = aval
  end function make_pdg_array

\eatline
\nwendcode{}\nwbegindocs{2427}Compile a (possible nested) expression, to obtain a
particle-specifier expression which we can process further.
\nwenddocs{}\nwbegincode{2428}\moddef{Commands: procedures}\plusendmoddef
  recursive subroutine compile_prt_expr (prt_expr, pn, var_list, model)
    type(prt_expr_t), intent(out) :: prt_expr
    type(parse_node_t), intent(in), target :: pn
    type(var_list_t), intent(in), target :: var_list
    type(model_t), intent(in), target :: model
    type(parse_node_t), pointer :: pn_entry, pn_term, pn_addition
    type(pdg_array_t) :: pdg
    type(string_t) :: prt_string
    integer :: n_entry, n_term, i
    select case (char (parse_node_get_rule_key (pn)))
    case ("prt_state_list")
       n_entry = parse_node_get_n_sub (pn)
       pn_entry => parse_node_get_sub_ptr (pn)
       if (n_entry == 1) then
          call compile_prt_expr (prt_expr, pn_entry, var_list, model)
       else
          call prt_expr%init_list (n_entry)
          select type (x => prt_expr%x)
          type is (prt_spec_list_t)
             do i = 1, n_entry
                call compile_prt_expr (x%expr(i), pn_entry, var_list, model)
                pn_entry => parse_node_get_next_ptr (pn_entry)
             end do
          end select
       end if
    case ("prt_state_sum")
       n_term = parse_node_get_n_sub (pn)
       pn_term => parse_node_get_sub_ptr (pn)
       pn_addition => pn_term
       if (n_term == 1) then
          call compile_prt_expr (prt_expr, pn_term, var_list, model)
       else
          call prt_expr%init_sum (n_term)
          select type (x => prt_expr%x)
          type is (prt_spec_sum_t)
             do i = 1, n_term
                call compile_prt_expr (x%expr(i), pn_term, var_list, model)
                pn_addition => parse_node_get_next_ptr (pn_addition)
                if (associated (pn_addition)) &
                     pn_term => parse_node_get_sub_ptr (pn_addition, 2)
             end do
          end select
       end if
    case ("cexpr")
       pdg = eval_pdg_array (pn, var_list)
       prt_string = make_flavor_string (pdg, model)
       call prt_expr%init_spec (new_prt_spec (prt_string))
    case default
       call parse_node_write_rec (pn)
       call msg_bug ("compile prt expr: impossible syntax rule")
    end select
  end subroutine compile_prt_expr
          
\eatline
\nwendcode{}\nwbegindocs{2429}\nwdocspar
\subsubsection{Initiating a NLO calculation}
\nwenddocs{}\nwbegincode{2430}\moddef{Commands: types}\plusendmoddef
  type, extends (command_t) :: cmd_nlo_t
    private
    type(parse_node_p), dimension(3) :: pn_components
    logical, dimension(3) :: active_component  
  contains
    \LA{}Commands: cmd nlo: TBP\RA{}
  end type cmd_nlo_t

\eatline
\nwendcode{}\nwbegindocs{2431}\nwdocspar
\nwenddocs{}\nwbegincode{2432}\moddef{Commands: cmd nlo: TBP}\endmoddef
  procedure :: write => cmd_nlo_write
\nwendcode{}\nwbegincode{2433}\moddef{Commands: procedures}\plusendmoddef
  subroutine cmd_nlo_write (cmd, unit, indent)
    class(cmd_nlo_t), intent(in) :: cmd
    integer, intent(in), optional :: unit, indent
  end subroutine cmd_nlo_write

\eatline
\nwendcode{}\nwbegindocs{2434}As it is, the NLO calculation is switched on by putting {nlo} behind the process definition. This should be made nicer in the future.
\nwenddocs{}\nwbegincode{2435}\moddef{Commands: cmd nlo: TBP}\plusendmoddef
  procedure :: compile => cmd_nlo_compile
\nwendcode{}\nwbegincode{2436}\moddef{Commands: procedures}\plusendmoddef
  subroutine cmd_nlo_compile (cmd, global)
    class(cmd_nlo_t), intent(inout) :: cmd
    type(rt_data_t), intent(inout), target :: global
    type(parse_node_t), pointer :: pn_arg, pn_comp
    integer :: i
    pn_arg => parse_node_get_sub_ptr (cmd%pn, 3)
    cmd%pn_components(1)%ptr => parse_node_get_sub_ptr (pn_arg)
    pn_comp => parse_node_get_next_ptr (cmd%pn_components(1)%ptr)
    i = 2
    do 
      if (associated (pn_comp)) then
         cmd%pn_components(i)%ptr => pn_comp
         pn_comp => parse_node_get_next_ptr (cmd%pn_components(i)%ptr)
         i = i+1
      else
         exit
      end if
    end do
  end subroutine cmd_nlo_compile

\eatline
\nwendcode{}\nwbegindocs{2437}\nwdocspar
\nwenddocs{}\nwbegincode{2438}\moddef{Commands: cmd nlo: TBP}\plusendmoddef
  procedure :: execute => cmd_nlo_execute
\nwendcode{}\nwbegincode{2439}\moddef{Commands: procedures}\plusendmoddef
  subroutine cmd_nlo_execute (cmd, global)
    class(cmd_nlo_t), intent(inout) :: cmd
    type(rt_data_t), intent(inout), target :: global
    type(parse_node_t), pointer :: current_component
    type(string_t) :: component_type
    integer :: i

    cmd%active_component = .false.

    current_component => cmd%pn_components(1)%ptr
    i = 2
    do 
      if (associated (current_component)) then
         component_type = eval_string (current_component, global%var_list)
         select case (char (component_type))
         case ('Born')
            cmd%active_component(1) = .true.
         case ('Real')
            cmd%active_component(2) = .true.
         case ('Virtual')
            cmd%active_component(3) = .true.
         case ('Full')
            cmd%active_component = .true.
         end select
         if (i >= 4) exit
         current_component => cmd%pn_components(i)%ptr
         i = i+1
      else
         exit
      end if
    end do   
    ! global%nlo_calculation = cmd%nlo_calc
    global%nlo_calculation = cmd%active_component(2) &
                        .or. cmd%active_component(3)
    global%active_nlo_components = cmd%active_component
  end subroutine cmd_nlo_execute

\eatline
\nwendcode{}\nwbegindocs{2440}\nwdocspar
\subsubsection{Process compilation}
\nwenddocs{}\nwbegincode{2441}\moddef{Commands: types}\plusendmoddef
  type, extends (command_t) :: cmd_compile_t
     private
     type(string_t), dimension(:), allocatable :: libname
     logical :: make_executable = .false.
     type(string_t) :: exec_name
   contains
   \LA{}Commands: cmd compile: TBP\RA{}
  end type cmd_compile_t

\eatline
\nwendcode{}\nwbegindocs{2442}Output: list all libraries to be compiled.
\nwenddocs{}\nwbegincode{2443}\moddef{Commands: cmd compile: TBP}\endmoddef
  procedure :: write => cmd_compile_write
\nwendcode{}\nwbegincode{2444}\moddef{Commands: procedures}\plusendmoddef
  subroutine cmd_compile_write (cmd, unit, indent)
    class(cmd_compile_t), intent(in) :: cmd
    integer, intent(in), optional :: unit, indent
    integer :: u, i
    u = given_output_unit (unit);  if (u < 0)  return
    call write_indent (u, indent)
    write (u, "(1x,A)", advance="no")  "compile ("
    if (allocated (cmd%libname)) then
       do i = 1, size (cmd%libname)
          if (i > 1)  write (u, "(A,1x)", advance="no")  ","
          write (u, "('""',A,'""')", advance="no")  char (cmd%libname(i))
       end do
    end if
    write (u, "(A)")  ")"
  end subroutine cmd_compile_write

\eatline
\nwendcode{}\nwbegindocs{2445}Compile the libraries specified in the argument.  If the argument is
empty, compile all libraries which can be found in the process library stack.
\nwenddocs{}\nwbegincode{2446}\moddef{Commands: cmd compile: TBP}\plusendmoddef
  procedure :: compile => cmd_compile_compile
\nwendcode{}\nwbegincode{2447}\moddef{Commands: procedures}\plusendmoddef
  subroutine cmd_compile_compile (cmd, global)
    class(cmd_compile_t), intent(inout) :: cmd
    type(rt_data_t), intent(inout), target :: global
    type(parse_node_t), pointer :: pn_cmd, pn_clause, pn_arg, pn_lib
    type(parse_node_t), pointer :: pn_exec_name_spec, pn_exec_name
    integer :: n_lib, i
    pn_cmd => parse_node_get_sub_ptr (cmd%pn)
    pn_clause => parse_node_get_sub_ptr (pn_cmd)
    pn_exec_name_spec => parse_node_get_sub_ptr (pn_clause, 2)
    if (associated (pn_exec_name_spec)) then
       pn_exec_name => parse_node_get_sub_ptr (pn_exec_name_spec, 2)
    else
       pn_exec_name => null ()
    end if
    pn_arg => parse_node_get_next_ptr (pn_clause)
    cmd%pn_opt => parse_node_get_next_ptr (pn_cmd)
    call cmd%compile_options (global)
    if (associated (pn_arg)) then
       n_lib = parse_node_get_n_sub (pn_arg)
    else
       n_lib = 0
    end if
    if (n_lib > 0) then
       allocate (cmd%libname (n_lib))
       pn_lib => parse_node_get_sub_ptr (pn_arg)
       do i = 1, n_lib
          cmd%libname(i) = parse_node_get_string (pn_lib)
          pn_lib => parse_node_get_next_ptr (pn_lib)
       end do
    end if
    if (associated (pn_exec_name)) then
       cmd%make_executable = .true.
       cmd%exec_name = parse_node_get_string (pn_exec_name)
    end if
  end subroutine cmd_compile_compile

\eatline
\nwendcode{}\nwbegindocs{2448}Command execution.  Generate code, write driver, compile and link.
Do this for all libraries in the list.

If no library names have been given and stored while compiling this
command, we collect all libraries from the current stack and compile
those.

We can compile static libraries (which actually just loads them).  However, we
can't incorporate in a generated executable.
\nwenddocs{}\nwbegincode{2449}\moddef{Commands: cmd compile: TBP}\plusendmoddef
  procedure :: execute => cmd_compile_execute
\nwendcode{}\nwbegincode{2450}\moddef{Commands: procedures}\plusendmoddef
  subroutine cmd_compile_execute (cmd, global)
    class(cmd_compile_t), intent(inout) :: cmd
    type(rt_data_t), intent(inout), target :: global
    type(string_t), dimension(:), allocatable :: libname, libname_static
    integer :: i
    if (allocated (cmd%libname)) then
       allocate (libname (size (cmd%libname)))
       libname = cmd%libname
    else
       call cmd%local%prclib_stack%get_names (libname)
    end if
    if (cmd%make_executable) then
       call get_prclib_static (libname_static)
       do i = 1, size (libname)
          if (any (libname_static == libname(i))) then
             call msg_fatal ("Compile: can't include static library '" &
                  // char (libname(i)) // "'")
          end if
       end do
       call compile_executable (cmd%exec_name, libname, cmd%local)
    else
       do i = 1, size (libname)
          call compile_library (libname(i), cmd%local)
       end do
    end if
  end subroutine cmd_compile_execute
    
\eatline
\nwendcode{}\nwbegindocs{2451}\nwdocspar
This is the interface to the external procedure which returns the
names of all static libraries which are part of the executable.  (The
default is none.)  The routine must allocate the array.
\nwenddocs{}\nwbegincode{2452}\moddef{Commands: public}\endmoddef
  public :: get_prclib_static
\nwendcode{}\nwbegincode{2453}\moddef{Commands: interfaces}\plusendmoddef
  interface
     subroutine get_prclib_static (libname)
       import
       type(string_t), dimension(:), intent(inout), allocatable :: libname
     end subroutine get_prclib_static
  end interface

\eatline
\nwendcode{}\nwbegindocs{2454}\nwdocspar
\subsubsection{Execute a shell command}
The argument is a string expression.
\nwenddocs{}\nwbegincode{2455}\moddef{Commands: types}\plusendmoddef
  type, extends (command_t) :: cmd_exec_t
     private
     type(parse_node_t), pointer :: pn_command => null ()
   contains
   \LA{}Commands: cmd exec: TBP\RA{}  
  end type cmd_exec_t

\eatline
\nwendcode{}\nwbegindocs{2456}Simply tell the status.
\nwenddocs{}\nwbegincode{2457}\moddef{Commands: cmd exec: TBP}\endmoddef
  procedure :: write => cmd_exec_write
\nwendcode{}\nwbegincode{2458}\moddef{Commands: procedures}\plusendmoddef
  subroutine cmd_exec_write (cmd, unit, indent)
    class(cmd_exec_t), intent(in) :: cmd
    integer, intent(in), optional :: unit, indent
    integer :: u
    u = given_output_unit (unit);  if (u < 0)  return
    call write_indent (u, indent)
    if (associated (cmd%pn_command)) then
       write (u, "(1x,A)")  "exec: [command associated]"
    else
       write (u, "(1x,A)")  "exec: [undefined]"       
    end if
  end subroutine cmd_exec_write

\eatline
\nwendcode{}\nwbegindocs{2459}Compile the exec command.
\nwenddocs{}\nwbegincode{2460}\moddef{Commands: cmd exec: TBP}\plusendmoddef
  procedure :: compile => cmd_exec_compile
\nwendcode{}\nwbegincode{2461}\moddef{Commands: procedures}\plusendmoddef
  subroutine cmd_exec_compile (cmd, global)
    class(cmd_exec_t), intent(inout) :: cmd
    type(rt_data_t), intent(inout), target :: global
    type(parse_node_t), pointer :: pn_arg, pn_command
    pn_arg => parse_node_get_sub_ptr (cmd%pn, 2)
    pn_command => parse_node_get_sub_ptr (pn_arg)
    cmd%pn_command => pn_command
  end subroutine cmd_exec_compile

\eatline
\nwendcode{}\nwbegindocs{2462}Execute the specified shell command.
\nwenddocs{}\nwbegincode{2463}\moddef{Commands: cmd exec: TBP}\plusendmoddef
  procedure :: execute => cmd_exec_execute
\nwendcode{}\nwbegincode{2464}\moddef{Commands: procedures}\plusendmoddef
  subroutine cmd_exec_execute (cmd, global)
    class(cmd_exec_t), intent(inout) :: cmd
    type(rt_data_t), intent(inout), target :: global
    type(string_t) :: command
    logical :: is_known
    integer :: status
    command = eval_string (cmd%pn_command, global%var_list, is_known=is_known)
    if (is_known) then
       if (command /= "") then
          call os_system_call (command, status, verbose=.true.)
          if (status /= 0) then
             write (msg_buffer, "(A,I0)")  "Return code = ", status
             call msg_message ()
             call msg_error ("System command returned with nonzero status code")
          end if
       end if 
    end if
  end subroutine cmd_exec_execute

\eatline
\nwendcode{}\nwbegindocs{2465}\nwdocspar
\subsubsection{Variable declaration}
A variable can have various types.  Hold the definition as an eval
tree.

There are intrinsic variables, user variables, and model variables.
The latter are further divided in independent variables and dependent
variables.

Regarding model variables: When dealing with them, we always look at
two variable lists in parallel.  The global (or local) variable list
contains the user-visible values.  It includes variables that
correspond to variables in the current model's list.  These, in turn,
are pointers to the model's parameter list, so the model is always in
sync, internally.  To keep the global variable list in sync with the
model, the global variables carry the {\tt{}is{\char95}copy} property and contain
a separate pointer to the model variable.  (The pointer is reassigned
whenever the model changes.)  Modifying the global variable changes
two values simultaneously: the visible value and the model variable,
via this extra pointer.  After each modification, we update dependent
parameters in the model variable list and re-synchronize the global
variable list (again, using these pointers) with the model variable
this.  In the last step, modifications in the derived parameters
become visible.

When we integrate a process, we capture the current variable list of
the current model in a separate model instance, which is stored in the
process object.  Thus, the model parameters associated to this process
at this time are preserved for the lifetime of the process object.

When we generate or rescan events, we can again capture a local model
variable list in a model instance.  This allows us to reweight event
by event with different parameter sets simultaneously.
\nwenddocs{}\nwbegincode{2466}\moddef{Commands: types}\plusendmoddef
  type, extends (command_t) :: cmd_var_t
     private
     type(string_t) :: name
     integer :: type = V_NONE
     type(parse_node_t), pointer :: pn_value => null ()
     logical :: is_intrinsic = .false.
     logical :: is_model_var = .false.
   contains
   \LA{}Commands: cmd var: TBP\RA{}
  end type cmd_var_t

\eatline
\nwendcode{}\nwbegindocs{2467}Output.  We know name, type, and properties, but not the value.
\nwenddocs{}\nwbegincode{2468}\moddef{Commands: cmd var: TBP}\endmoddef
  procedure :: write => cmd_var_write
\nwendcode{}\nwbegincode{2469}\moddef{Commands: procedures}\plusendmoddef
  subroutine cmd_var_write (cmd, unit, indent)
    class(cmd_var_t), intent(in) :: cmd
    integer, intent(in), optional :: unit, indent
    integer :: u
    u = given_output_unit (unit);  if (u < 0)  return
    call write_indent (u, indent)
    write (u, "(1x,A,A,A)", advance="no")  "var: ", char (cmd%name), " ("
    select case (cmd%type)
    case (V_NONE)
       write (u, "(A)", advance="no")  "[unknown]"
    case (V_LOG)
       write (u, "(A)", advance="no")  "logical"
    case (V_INT)
       write (u, "(A)", advance="no")  "int"
    case (V_REAL)
       write (u, "(A)", advance="no")  "real"
    case (V_CMPLX)
       write (u, "(A)", advance="no")  "complex"
    case (V_STR)
       write (u, "(A)", advance="no")  "string"
    case (V_PDG)
       write (u, "(A)", advance="no")  "alias"
    end select
    if (cmd%is_intrinsic) then
       write (u, "(A)", advance="no")  ", intrinsic"
    end if
    if (cmd%is_model_var) then
       write (u, "(A)", advance="no")  ", model"
    end if
    write (u, "(A)")  ")"
  end subroutine cmd_var_write

\eatline
\nwendcode{}\nwbegindocs{2470}Compile the lhs and determine the variable name and type.  Check whether
this variable can be created or modified as requested, and append the value to
the variable list, if appropriate.  The value is initially undefined.
The rhs is assigned to a pointer, to be compiled and evaluated when the
command is executed.
\nwenddocs{}\nwbegincode{2471}\moddef{Commands: cmd var: TBP}\plusendmoddef
  procedure :: compile => cmd_var_compile
\nwendcode{}\nwbegincode{2472}\moddef{Commands: procedures}\plusendmoddef
  subroutine cmd_var_compile (cmd, global)
    class(cmd_var_t), intent(inout) :: cmd
    type(rt_data_t), intent(inout), target :: global
    type(parse_node_t), pointer :: pn_var, pn_name
    type(parse_node_t), pointer :: pn_result, pn_proc
    type(string_t) :: var_name
    type(var_list_t), pointer :: model_vars
    integer :: type
    logical :: new
    pn_result => null ()
    new = .false.
    select case (char (parse_node_get_rule_key (cmd%pn)))
    case ("cmd_log_decl");    type = V_LOG
       pn_var => parse_node_get_sub_ptr (cmd%pn, 2)
       if (.not. associated (pn_var)) then   ! handle masked syntax error 
          cmd%type = V_NONE; return
       end if
       pn_name => parse_node_get_sub_ptr (pn_var, 2)
       new = .true.
    case ("cmd_log");         type = V_LOG
       pn_name => parse_node_get_sub_ptr (cmd%pn, 2)
    case ("cmd_int");         type = V_INT
       pn_name => parse_node_get_sub_ptr (cmd%pn, 2)
       new = .true.
    case ("cmd_real");        type = V_REAL
       pn_name => parse_node_get_sub_ptr (cmd%pn, 2)
       new = .true.
    case ("cmd_complex");       type = V_CMPLX
       pn_name => parse_node_get_sub_ptr (cmd%pn, 2)
       new = .true.
    case ("cmd_num");         type = V_NONE
       pn_name => parse_node_get_sub_ptr (cmd%pn)
    case ("cmd_string_decl"); type = V_STR
       pn_var => parse_node_get_sub_ptr (cmd%pn, 2)
       if (.not. associated (pn_var)) then   ! handle masked syntax error 
          cmd%type = V_NONE; return
       end if
       pn_name => parse_node_get_sub_ptr (pn_var, 2)
       new = .true.
    case ("cmd_string");      type = V_STR
       pn_name => parse_node_get_sub_ptr (cmd%pn, 2)
    case ("cmd_alias");       type = V_PDG
       pn_name => parse_node_get_sub_ptr (cmd%pn, 2)
       new = .true.
    case ("cmd_result");      type = V_REAL
       pn_name => parse_node_get_sub_ptr (cmd%pn)
       pn_result => parse_node_get_sub_ptr (pn_name)
       pn_proc => parse_node_get_next_ptr (pn_result)
    case default
       call parse_node_mismatch &
            ("logical|int|real|complex|?|$|alias|var_name", cmd%pn)  ! $
    end select
    if (.not. associated (pn_name)) then   ! handle masked syntax error 
       cmd%type = V_NONE; return
    end if
    if (.not. associated (pn_result)) then
       var_name = parse_node_get_string (pn_name)
    else
       var_name = parse_node_get_key (pn_result) &
            // "(" // parse_node_get_string (pn_proc) // ")"
    end if
    select case (type)
    case (V_LOG);  var_name = "?" // var_name
    case (V_STR);  var_name = "$" // var_name    ! $
    end select
    if (associated (global%model)) then
       model_vars => global%model%get_var_list_ptr ()
    else
       model_vars => null ()
    end if
    call var_list_check_observable (global%var_list, var_name, type)
    call var_list_check_result_var (global%var_list, var_name, type)
    call var_list_check_user_var (global%var_list, var_name, type, new)
    cmd%name = var_name
    cmd%pn_value => parse_node_get_next_ptr (pn_name, 2)
    if (global%var_list%contains (cmd%name, follow_link = .false.)) then
       ! local variable
       cmd%is_intrinsic = &
            global%var_list%is_intrinsic (cmd%name, follow_link = .false.)
       cmd%type = &
            global%var_list%get_type (cmd%name, follow_link = .false.)
    else
       if (new)  cmd%type = type
       if (global%var_list%contains (cmd%name, follow_link = .true.)) then
          ! global variable
          cmd%is_intrinsic = &
               global%var_list%is_intrinsic (cmd%name, follow_link = .true.)
          if (cmd%type == V_NONE) then
             cmd%type = &
                  global%var_list%get_type (cmd%name, follow_link = .true.)
          end if
       else if (associated (model_vars)) then  ! check model variable
          cmd%is_model_var = &
               model_vars%contains (cmd%name)
          if (cmd%type == V_NONE) then
             cmd%type = &
                  model_vars%get_type (cmd%name)
          end if
       end if
       if (cmd%type == V_NONE) then
          call msg_fatal ("Variable '" // char (cmd%name) // "' " &
               // "set without declaration")
          cmd%type = V_NONE;  return
       end if
       if (cmd%is_model_var) then
          if (new) then
             call msg_fatal ("Model variable '" // char (cmd%name) // "' " &
                  // "redeclared")
          else if (model_vars%is_locked (cmd%name)) then
             call msg_fatal ("Model variable '" // char (cmd%name) // "' " &
                  // "is locked")
          end if
       else
          select case (cmd%type)
          case (V_LOG)
             call var_list_append_log (global%var_list, cmd%name, &
                  intrinsic=cmd%is_intrinsic, user=.true.)
          case (V_INT)
             call var_list_append_int (global%var_list, cmd%name, &
                  intrinsic=cmd%is_intrinsic, user=.true.)
          case (V_REAL)
             call var_list_append_real (global%var_list, cmd%name, &
                  intrinsic=cmd%is_intrinsic, user=.true.)
          case (V_CMPLX)
             call var_list_append_cmplx (global%var_list, cmd%name, &
                  intrinsic=cmd%is_intrinsic, user=.true.)
          case (V_PDG)
             call var_list_append_pdg_array (global%var_list, cmd%name, &
                  intrinsic=cmd%is_intrinsic, user=.true.)
          case (V_STR)
             call var_list_append_string (global%var_list, cmd%name, &
                  intrinsic=cmd%is_intrinsic, user=.true.)
          end select
       end if
    end if
  end subroutine cmd_var_compile

\eatline
\nwendcode{}\nwbegindocs{2473}Execute.  Evaluate the definition and assign the variable value.
If the variable is a model variable, take a snapshot of the model if necessary
and set the variable in the local model.
\nwenddocs{}\nwbegincode{2474}\moddef{Commands: cmd var: TBP}\plusendmoddef
  procedure :: execute => cmd_var_execute
\nwendcode{}\nwbegincode{2475}\moddef{Commands: procedures}\plusendmoddef
  subroutine cmd_var_execute (cmd, global)
    class(cmd_var_t), intent(inout) :: cmd
    type(rt_data_t), intent(inout), target :: global
    type(var_list_t), pointer :: var_list
    real(default) :: rval
    logical :: is_known, pacified
    var_list => global%get_var_list_ptr ()
    if (cmd%is_model_var) then
       pacified = var_list%get_lval (var_str ("?pacify"))     
       rval = eval_real (cmd%pn_value, var_list, is_known=is_known)
       call global%model_set_real &
            (cmd%name, rval, verbose=.true., pacified=pacified)
    else if (cmd%type /= V_NONE) then
       call cmd%set_value (var_list, verbose=.true.)
    end if
  end subroutine cmd_var_execute

\eatline
\nwendcode{}\nwbegindocs{2476}Copy the value to the variable list, where the variable should already exist.
\nwenddocs{}\nwbegincode{2477}\moddef{Commands: cmd var: TBP}\plusendmoddef
  procedure :: set_value => cmd_var_set_value
\nwendcode{}\nwbegincode{2478}\moddef{Commands: procedures}\plusendmoddef
  subroutine cmd_var_set_value (var, var_list, verbose, model_name)
    class(cmd_var_t), intent(inout) :: var
    type(var_list_t), intent(inout), target :: var_list
    logical, intent(in), optional :: verbose
    type(string_t), intent(in), optional :: model_name
    logical :: lval, pacified
    integer :: ival
    real(default) :: rval
    complex(default) :: cval
    type(pdg_array_t) :: aval
    type(string_t) :: sval
    logical :: is_known
    pacified = var_list%get_lval (var_str ("?pacify"))     
    select case (var%type)
    case (V_LOG)
       lval = eval_log (var%pn_value, var_list, is_known=is_known)
       call var_list_set_log (var_list, var%name, &
            lval, is_known, verbose=verbose, model_name=model_name)
    case (V_INT)
       ival = eval_int (var%pn_value, var_list, is_known=is_known)
       call var_list_set_int (var_list, var%name, &
            ival, is_known, verbose=verbose, model_name=model_name)
    case (V_REAL)
       rval = eval_real (var%pn_value, var_list, is_known=is_known)
       call var_list_set_real (var_list, var%name, &
            rval, is_known, verbose=verbose, &
            model_name=model_name, pacified = pacified)
    case (V_CMPLX)
       cval = eval_cmplx (var%pn_value, var_list, is_known=is_known)
       call var_list_set_cmplx (var_list, var%name, &
            cval, is_known, verbose=verbose, &
            model_name=model_name, pacified = pacified)
    case (V_PDG)
       aval = eval_pdg_array (var%pn_value, var_list, is_known=is_known)
       call var_list_set_pdg_array (var_list, var%name, &
            aval, is_known, verbose=verbose, model_name=model_name)
    case (V_STR)
       sval = eval_string (var%pn_value, var_list, is_known=is_known)
       call var_list_set_string (var_list, var%name, &
            sval, is_known, verbose=verbose, model_name=model_name)
    end select
  end subroutine cmd_var_set_value
  
\eatline
\nwendcode{}\nwbegindocs{2479}\nwdocspar
\subsubsection{SLHA}
Read a SLHA (SUSY Les Houches Accord) file to fill the appropriate
model parameters.  We do not access the current variable record, but
directly work on the appropriate SUSY model, which is loaded if
necessary.

We may be in read or write mode.  In the latter case, we may write
just input parameters, or the complete spectrum, or the spectrum with
all decays.
\nwenddocs{}\nwbegincode{2480}\moddef{Commands: types}\plusendmoddef
  type, extends (command_t) :: cmd_slha_t
     private
     type(string_t) :: file
     logical :: write_mode = .false.
   contains
   \LA{}Commands: cmd slha: TBP\RA{}  
  end type cmd_slha_t

\eatline
\nwendcode{}\nwbegindocs{2481}Output.
\nwenddocs{}\nwbegincode{2482}\moddef{Commands: cmd slha: TBP}\endmoddef
  procedure :: write => cmd_slha_write
\nwendcode{}\nwbegincode{2483}\moddef{Commands: procedures}\plusendmoddef
  subroutine cmd_slha_write (cmd, unit, indent)
    class(cmd_slha_t), intent(in) :: cmd
    integer, intent(in), optional :: unit, indent
    integer :: u
    u = given_output_unit (unit);  if (u < 0)  return
    call write_indent (u, indent)
    write (u, "(1x,A,A)")  "slha: file name  = ", char (cmd%file)
    write (u, "(1x,A,L1)") "slha: write mode = ", cmd%write_mode
  end subroutine cmd_slha_write

\eatline
\nwendcode{}\nwbegindocs{2484}Compile.  Read the filename and store it.
\nwenddocs{}\nwbegincode{2485}\moddef{Commands: cmd slha: TBP}\plusendmoddef
  procedure :: compile => cmd_slha_compile
\nwendcode{}\nwbegincode{2486}\moddef{Commands: procedures}\plusendmoddef
  subroutine cmd_slha_compile (cmd, global)
    class(cmd_slha_t), intent(inout) :: cmd
    type(rt_data_t), intent(inout), target :: global
    type(parse_node_t), pointer :: pn_key, pn_arg, pn_file
    pn_key => parse_node_get_sub_ptr (cmd%pn)
    pn_arg => parse_node_get_next_ptr (pn_key)
    pn_file => parse_node_get_sub_ptr (pn_arg)
    call cmd%compile_options (global)
    cmd%pn_opt => parse_node_get_next_ptr (pn_arg)
    select case (char (parse_node_get_key (pn_key)))
    case ("read_slha")
       cmd%write_mode = .false.
    case ("write_slha")
       cmd%write_mode = .true.
    case default
       call parse_node_mismatch ("read_slha|write_slha",  cmd%pn)
    end select
    cmd%file = parse_node_get_string (pn_file)
  end subroutine cmd_slha_compile

\eatline
\nwendcode{}\nwbegindocs{2487}Execute.  Read or write the specified SLHA file.  Behind the scenes,
this will first read the WHIZARD model file, then read the SLHA file
and assign the SLHA parameters as far as determined by
{\tt{}dispatch{\char95}slha}.  Finally, the global variables are synchronized
with the model.  This is similar to executing {\tt{}cmd{\char95}model}.
\nwenddocs{}\nwbegincode{2488}\moddef{Commands: cmd slha: TBP}\plusendmoddef
  procedure :: execute => cmd_slha_execute
\nwendcode{}\nwbegincode{2489}\moddef{Commands: procedures}\plusendmoddef
  subroutine cmd_slha_execute (cmd, global)
    class(cmd_slha_t), intent(inout) :: cmd
    type(rt_data_t), intent(inout), target :: global
    logical :: input, spectrum, decays
    type(model_t), pointer :: model
    if (cmd%write_mode) then
       input = .true.
       spectrum = .false.
       decays = .false.
       if (.not. associated (cmd%local%model)) then
          call msg_fatal ("SLHA: local model not associated")
          return
       end if
       call slha_write_file &
            (cmd%file, cmd%local%model, &
             input = input, spectrum = spectrum, decays = decays)
    else
       if (.not. associated (global%model)) then
          call msg_fatal ("SLHA: global model not associated")
          return
       end if
       call dispatch_slha (cmd%local, &
            input = input, spectrum = spectrum, decays = decays)
       call global%ensure_model_copy ()
       call slha_read_file &
            (cmd%file, cmd%local%os_data, global%model, &
             input = input, spectrum = spectrum, decays = decays)
    end if
  end subroutine cmd_slha_execute

\eatline
\nwendcode{}\nwbegindocs{2490}\nwdocspar
\subsubsection{Show values}
This command shows the current values of variables or other objects,
in a suitably condensed form.
\nwenddocs{}\nwbegincode{2491}\moddef{Commands: types}\plusendmoddef
  type, extends (command_t) :: cmd_show_t
     private
     type(string_t), dimension(:), allocatable :: name
   contains
   \LA{}Commands: cmd show: TBP\RA{}
  end type cmd_show_t

\eatline
\nwendcode{}\nwbegindocs{2492}Output: list the object names, not values.
\nwenddocs{}\nwbegincode{2493}\moddef{Commands: cmd show: TBP}\endmoddef
  procedure :: write => cmd_show_write
\nwendcode{}\nwbegincode{2494}\moddef{Commands: procedures}\plusendmoddef
  subroutine cmd_show_write (cmd, unit, indent)
    class(cmd_show_t), intent(in) :: cmd
    integer, intent(in), optional :: unit, indent
    integer :: u, i
    u = given_output_unit (unit);  if (u < 0)  return
    call write_indent (u, indent)  
    write (u, "(1x,A)", advance="no")  "show: "  
    if (allocated (cmd%name)) then
       do i = 1, size (cmd%name)
          write (u, "(1x,A)", advance="no")  char (cmd%name(i))
       end do
       write (u, *)
    else
       write (u, "(5x,A)")  "[undefined]"
    end if
  end subroutine cmd_show_write

\eatline
\nwendcode{}\nwbegindocs{2495}Compile.  Allocate an array which is filled with the names of the
variables to show.
\nwenddocs{}\nwbegincode{2496}\moddef{Commands: cmd show: TBP}\plusendmoddef
  procedure :: compile => cmd_show_compile
\nwendcode{}\nwbegincode{2497}\moddef{Commands: procedures}\plusendmoddef
  subroutine cmd_show_compile (cmd, global)
    class(cmd_show_t), intent(inout) :: cmd
    type(rt_data_t), intent(inout), target :: global
    type(parse_node_t), pointer :: pn_arg, pn_var, pn_prefix, pn_name
    type(string_t) :: key
    integer :: i, n_args
    pn_arg => parse_node_get_sub_ptr (cmd%pn, 2)
    if (associated (pn_arg)) then
       select case (char (parse_node_get_rule_key (pn_arg)))
       case ("show_arg")
          cmd%pn_opt => parse_node_get_next_ptr (pn_arg)
       case default
          cmd%pn_opt => pn_arg
          pn_arg => null ()
       end select
    end if
    call cmd%compile_options (global)
    if (associated (pn_arg)) then
       n_args = parse_node_get_n_sub (pn_arg)
       allocate (cmd%name (n_args))
       pn_var => parse_node_get_sub_ptr (pn_arg)
       i = 0
       do while (associated (pn_var))
          i = i + 1
          select case (char (parse_node_get_rule_key (pn_var)))
          case ("model", "library", "beams", "iterations", &
                "cuts", "weight", "int", "real", "complex", &
                "scale", "factorization_scale", "renormalization_scale", &
                "selection", "reweight", "analysis", "pdg", &
                "stable", "unstable", "polarized", "unpolarized", &
                "results", "expect", "intrinsic", "string", "logical")
             cmd%name(i) = parse_node_get_key (pn_var)
          case ("result_var")
             pn_prefix => parse_node_get_sub_ptr (pn_var)
             pn_name => parse_node_get_next_ptr (pn_prefix)
             if (associated (pn_name)) then
                cmd%name(i) = parse_node_get_key (pn_prefix) &
                     // "(" // parse_node_get_string (pn_name) // ")"
             else
                cmd%name(i) = parse_node_get_key (pn_prefix)
             end if
          case ("log_var", "string_var", "alias_var")
             pn_prefix => parse_node_get_sub_ptr (pn_var)
             pn_name => parse_node_get_next_ptr (pn_prefix)
             key = parse_node_get_key (pn_prefix)
             if (associated (pn_name)) then
                select case (char (parse_node_get_rule_key (pn_name)))
                case ("var_name")
                   select case (char (key))
                   case ("?", "$")  ! $ sign
                      cmd%name(i) = key // parse_node_get_string (pn_name)
                   case ("alias")
                      cmd%name(i) = parse_node_get_string (pn_name)
                   end select
                case default
                   call parse_node_mismatch &
                        ("var_name",  pn_name)
                end select
             else
                cmd%name(i) = key
             end if
          case default
             cmd%name(i) = parse_node_get_string (pn_var)
          end select
          pn_var => parse_node_get_next_ptr (pn_var)
       end do
    else
       allocate (cmd%name (0))
    end if
  end subroutine cmd_show_compile

\eatline
\nwendcode{}\nwbegindocs{2498}Execute.  Scan the list of objects to show.
\nwenddocs{}\nwbegincode{2499}\moddef{Commands: parameters}\endmoddef
  integer, parameter, public :: SHOW_BUFFER_SIZE = 4096
\nwendcode{}\nwbegincode{2500}\moddef{Commands: cmd show: TBP}\plusendmoddef
  procedure :: execute => cmd_show_execute
\nwendcode{}\nwbegincode{2501}\moddef{Commands: procedures}\plusendmoddef
  subroutine cmd_show_execute (cmd, global)
    class(cmd_show_t), intent(inout) :: cmd
    type(rt_data_t), intent(inout), target :: global
    type(var_list_t), pointer :: var_list, model_vars
    type(model_t), pointer :: model
    type(string_t) :: name
    integer :: n, pdg
    type(flavor_t) :: flv
    type(process_library_t), pointer :: prc_lib
    type(process_t), pointer :: process
    logical :: pacified
    character(SHOW_BUFFER_SIZE) :: buffer
    integer :: i, j, u, u_log, u_out
    u = free_unit ()
    var_list => cmd%local%var_list
    if (associated (cmd%local%model)) then
       model_vars => cmd%local%model%get_var_list_ptr ()
    else
       model_vars => null ()
    end if
    pacified = var_list%get_lval (var_str ("?pacify"))
    open (u, status = "scratch", action = "readwrite")
    if (associated (cmd%local%model)) then
       name = cmd%local%model%get_name ()
    end if
    if (size (cmd%name) == 0) then
       if (associated (model_vars)) then
          call var_list_write (model_vars, model_name = name, &
               unit = u, pacified = pacified, follow_link = .false.)
       end if
       call var_list_write (var_list, unit = u, pacified = pacified)
    else
       do i = 1, size (cmd%name)
          select case (char (cmd%name(i)))
          case ("model")
             if (associated (cmd%local%model)) then
                call cmd%local%model%show (u)
             else
                write (u, "(A)")  "Model: [undefined]"
             end if
          case ("library")
             if (associated (cmd%local%prclib)) then
                call cmd%local%prclib%show (u)
             else
                write (u, "(A)")  "Process library: [undefined]"
             end if
          case ("beams")
             call cmd%local%show_beams (u)
          case ("iterations")
             call cmd%local%it_list%write (u)
          case ("results")
             call cmd%local%process_stack%show (u)
          case ("stable")
             call cmd%local%model%show_stable (u)
          case ("polarized")
             call cmd%local%model%show_polarized (u)
          case ("unpolarized")
             call cmd%local%model%show_unpolarized (u)
          case ("unstable")
             model => cmd%local%model
             call model%show_unstable (u)
             n = model%get_n_field ()
             do j = 1, n
                pdg = model%get_pdg (j)
                call flv%init (pdg, model)
                if (.not. flv%is_stable ()) &
                     call show_unstable (cmd%local, pdg, u)
                if (flv%has_antiparticle ()) then
                   associate (anti => flv%anti ())
                     if (.not. anti%is_stable ()) &
                          call show_unstable (cmd%local, -pdg, u)
                   end associate
                end if
             end do
          case ("cuts", "weight", "scale", &
               "factorization_scale", "renormalization_scale", &
               "selection", "reweight", "analysis")
             call cmd%local%pn%show (cmd%name(i), u)
          case ("expect")
             call expect_summary (force = .true.)
          case ("intrinsic")
             call var_list_write (var_list, &
                  intrinsic=.true., unit=u, pacified = pacified)
          case ("logical")
             if (associated (model_vars)) then
                call var_list_write (model_vars, only_type=V_LOG, &
                     model_name = name, unit=u, pacified = pacified, &
                     follow_link=.false.)
             end if
             call var_list_write (var_list, &
                  only_type=V_LOG, unit=u, pacified = pacified)
          case ("int")
             if (associated (model_vars)) then
                call var_list_write (model_vars, only_type=V_INT, &
                     model_name = name, unit=u, pacified = pacified, &
                     follow_link=.false.)
             end if
             call var_list_write (var_list, only_type=V_INT, &
                  unit=u, pacified = pacified)
          case ("real")
             if (associated (model_vars)) then
                call var_list_write (model_vars, only_type=V_REAL, &
                     model_name = name, unit=u, pacified = pacified, &
                     follow_link=.false.)
             end if
             call var_list_write (var_list, only_type=V_REAL, &
                  unit=u, pacified = pacified)
          case ("complex")
             if (associated (model_vars)) then
                call var_list_write (model_vars, only_type=V_CMPLX, &
                     model_name = name, unit=u, pacified = pacified, &
                     follow_link=.false.)
             end if
             call var_list_write (var_list, only_type=V_CMPLX, &
                  unit=u, pacified = pacified)
          case ("pdg")
             if (associated (model_vars)) then
                call var_list_write (model_vars, only_type=V_PDG, &
                     model_name = name, unit=u, pacified = pacified, &
                     follow_link=.false.)
             end if
             call var_list_write (var_list, only_type=V_PDG, &
                  unit=u, pacified = pacified)
          case ("string") 
             if (associated (model_vars)) then
                call var_list_write (model_vars, only_type=V_STR, &
                     model_name = name, unit=u, pacified = pacified, &
                     follow_link=.false.)
             end if
             call var_list_write (var_list, only_type=V_STR, &
                  unit=u, pacified = pacified)
          case default
             if (analysis_exists (cmd%name(i))) then
                call analysis_write (cmd%name(i), u)
             else if (cmd%local%process_stack%exists (cmd%name(i))) then
                process => cmd%local%process_stack%get_process_ptr (cmd%name(i))
                call process%show (u)
             else if (associated (cmd%local%prclib_stack%get_library_ptr &
                  (cmd%name(i)))) then
                prc_lib => cmd%local%prclib_stack%get_library_ptr (cmd%name(i))
                call prc_lib%show (u)
             else if (associated (model_vars)) then
                if (model_vars%contains (cmd%name(i), follow_link=.false.)) then
                   call var_list_write_var (model_vars, cmd%name(i), &
                        unit = u, model_name = name, pacified = pacified)
                else if (var_list%contains (cmd%name(i))) then
                   call var_list_write_var (var_list, cmd%name(i), &
                        unit = u, pacified = pacified)
                else
                   call msg_error ("show: object '" // char (cmd%name(i)) &
                        // "' not found")
                end if
             else if (var_list%contains (cmd%name(i))) then
                call var_list_write_var (var_list, cmd%name(i), &
                     unit = u, pacified = pacified)
             else
                call msg_error ("show: object '" // char (cmd%name(i)) &
                     // "' not found")
             end if
          end select
       end do
    end if
    rewind (u)
    u_log = logfile_unit ()
    u_out = given_output_unit ()
    do
       read (u, "(A)", end = 1)  buffer
       if (u_log > 0)  write (u_log, "(A)")  trim (buffer)
       if (u_out > 0)  write (u_out, "(A)")  trim (buffer)
    end do
1   close (u)
    if (u_log > 0)  flush (u_log)
    if (u_out > 0)  flush (u_out)
  end subroutine cmd_show_execute

\eatline
\nwendcode{}\nwbegindocs{2502}\nwdocspar
\subsubsection{Clear values}
This command clears the current values of variables or other objects,
where this makes sense.  It parallels the {\tt{}show} command.  The
objects are cleared, but not deleted.
\nwenddocs{}\nwbegincode{2503}\moddef{Commands: types}\plusendmoddef
  type, extends (command_t) :: cmd_clear_t
     private
     type(string_t), dimension(:), allocatable :: name
   contains
   \LA{}Commands: cmd clear: TBP\RA{}
  end type cmd_clear_t
     
\eatline
\nwendcode{}\nwbegindocs{2504}Output: list the names of the objects to be cleared.
\nwenddocs{}\nwbegincode{2505}\moddef{Commands: cmd clear: TBP}\endmoddef
  procedure :: write => cmd_clear_write
\nwendcode{}\nwbegincode{2506}\moddef{Commands: procedures}\plusendmoddef
  subroutine cmd_clear_write (cmd, unit, indent)
    class(cmd_clear_t), intent(in) :: cmd
    integer, intent(in), optional :: unit, indent
    integer :: u, i
    u = given_output_unit (unit);  if (u < 0)  return
    call write_indent (u, indent)  
    write (u, "(1x,A)", advance="no")  "clear: "  
    if (allocated (cmd%name)) then
       do i = 1, size (cmd%name)
          write (u, "(1x,A)", advance="no")  char (cmd%name(i))
       end do
       write (u, *)
    else
       write (u, "(5x,A)")  "[undefined]"
    end if
  end subroutine cmd_clear_write

\eatline
\nwendcode{}\nwbegindocs{2507}Compile.  Allocate an array which is filled with the names of the
objects to be cleared.

Note: there is currently no need to account for options, but we
prepare for that possibility.
\nwenddocs{}\nwbegincode{2508}\moddef{Commands: cmd clear: TBP}\plusendmoddef
  procedure :: compile => cmd_clear_compile
\nwendcode{}\nwbegincode{2509}\moddef{Commands: procedures}\plusendmoddef
  subroutine cmd_clear_compile (cmd, global)
    class(cmd_clear_t), intent(inout) :: cmd
    type(rt_data_t), intent(inout), target :: global
    type(parse_node_t), pointer :: pn_arg, pn_var, pn_prefix, pn_name
    type(string_t) :: key
    integer :: i, n_args
    pn_arg => parse_node_get_sub_ptr (cmd%pn, 2)
    if (associated (pn_arg)) then
       select case (char (parse_node_get_rule_key (pn_arg)))
       case ("clear_arg")
          cmd%pn_opt => parse_node_get_next_ptr (pn_arg)
       case default
          cmd%pn_opt => pn_arg
          pn_arg => null ()
       end select
    end if
    call cmd%compile_options (global)
    if (associated (pn_arg)) then
       n_args = parse_node_get_n_sub (pn_arg)
       allocate (cmd%name (n_args))
       pn_var => parse_node_get_sub_ptr (pn_arg)
       i = 0
       do while (associated (pn_var))
          i = i + 1
          select case (char (parse_node_get_rule_key (pn_var)))
          case ("beams", "iterations", &
                "cuts", "weight", &
                "scale", "factorization_scale", "renormalization_scale", &
                "selection", "reweight", "analysis", &
                "unstable", "polarized", &
                "expect")
             cmd%name(i) = parse_node_get_key (pn_var)
          case ("log_var", "string_var")
             pn_prefix => parse_node_get_sub_ptr (pn_var)
             pn_name => parse_node_get_next_ptr (pn_prefix)
             key = parse_node_get_key (pn_prefix)
             if (associated (pn_name)) then
                select case (char (parse_node_get_rule_key (pn_name)))
                case ("var_name")
                   select case (char (key))
                   case ("?", "$")  ! $ sign
                      cmd%name(i) = key // parse_node_get_string (pn_name)
                   end select
                case default
                   call parse_node_mismatch &
                        ("var_name",  pn_name)
                end select
             else
                cmd%name(i) = key
             end if
          case default
             cmd%name(i) = parse_node_get_string (pn_var)
          end select
          pn_var => parse_node_get_next_ptr (pn_var)
       end do
    else
       allocate (cmd%name (0))
    end if
  end subroutine cmd_clear_compile

\eatline
\nwendcode{}\nwbegindocs{2510}Execute.  Scan the list of objects to clear.

Objects that can be shown but not cleared: model, library, results
\nwenddocs{}\nwbegincode{2511}\moddef{Commands: cmd clear: TBP}\plusendmoddef
  procedure :: execute => cmd_clear_execute
\nwendcode{}\nwbegincode{2512}\moddef{Commands: procedures}\plusendmoddef
  subroutine cmd_clear_execute (cmd, global)
    class(cmd_clear_t), intent(inout) :: cmd
    type(rt_data_t), intent(inout), target :: global
    integer :: i
    logical :: success
    type(var_list_t), pointer :: model_vars
    if (size (cmd%name) == 0) then
       call msg_warning ("clear: no object specified")
    else
       do i = 1, size (cmd%name)
          success = .true.
          select case (char (cmd%name(i)))
          case ("beams")
             call cmd%local%clear_beams ()
          case ("iterations")
             call cmd%local%it_list%clear ()
          case ("polarized")
             call cmd%local%model%clear_polarized ()
          case ("unstable")
             call cmd%local%model%clear_unstable ()
          case ("cuts", "weight", "scale", &
               "factorization_scale", "renormalization_scale", &
               "selection", "reweight", "analysis")
             call cmd%local%pn%clear (cmd%name(i))
          case ("expect")
             call expect_clear ()
          case default
             if (analysis_exists (cmd%name(i))) then
                call analysis_clear (cmd%name(i))
             else if (cmd%local%var_list%contains (cmd%name(i))) then
                if (.not. cmd%local%var_list%is_locked (cmd%name(i))) then
                   call cmd%local%var_list%unset (cmd%name(i))
                else
                   call msg_error ("clear: variable '" // char (cmd%name(i)) &
                        // "' is locked and can't be cleared")
                   success = .false.
                end if
             else if (associated (cmd%local%model)) then
                model_vars => cmd%local%model%get_var_list_ptr ()
                if (model_vars%contains (cmd%name(i), follow_link=.false.)) then
                   call msg_error ("clear: variable '" // char (cmd%name(i)) &
                        // "' is a model variable and can't be cleared")
                else
                   call msg_error ("clear: object '" // char (cmd%name(i)) &
                        // "' not found")
                end if
                success = .false.
             else
                call msg_error ("clear: object '" // char (cmd%name(i)) &
                     // "' not found")
                success = .false.
             end if
          end select
          if (success)  call msg_message ("cleared: " // char (cmd%name(i)))
       end do
    end if
  end subroutine cmd_clear_execute

\eatline
\nwendcode{}\nwbegindocs{2513}\nwdocspar
\subsubsection{Compare values of variables to expectation}
The implementation is similar to the {\tt{}show} command.  There are just
two arguments: two values that should be compared.  For providing
local values for the numerical tolerance, the command has a local
argument list.

If the expectation fails, an error condition is recorded.
\nwenddocs{}\nwbegincode{2514}\moddef{Commands: types}\plusendmoddef
  type, extends (command_t) :: cmd_expect_t
     private
     type(parse_node_t), pointer :: pn_lexpr => null ()
   contains
   \LA{}Commands: cmd expect: TBP\RA{}  
  end type cmd_expect_t

\eatline
\nwendcode{}\nwbegindocs{2515}Simply tell the status.
\nwenddocs{}\nwbegincode{2516}\moddef{Commands: cmd expect: TBP}\endmoddef
  procedure :: write => cmd_expect_write
\nwendcode{}\nwbegincode{2517}\moddef{Commands: procedures}\plusendmoddef
  subroutine cmd_expect_write (cmd, unit, indent)
    class(cmd_expect_t), intent(in) :: cmd
    integer, intent(in), optional :: unit, indent
    integer :: u
    u = given_output_unit (unit);  if (u < 0)  return
    call write_indent (u, indent)
    if (associated (cmd%pn_lexpr)) then
       write (u, "(1x,A)")  "expect: [expression associated]"
    else
       write (u, "(1x,A)")  "expect: [undefined]"       
    end if
  end subroutine cmd_expect_write

\eatline
\nwendcode{}\nwbegindocs{2518}Compile.  This merely assigns the parse node, the actual compilation is done
at execution.  This is necessary because the origin of variables
(local/global) may change during execution.
\nwenddocs{}\nwbegincode{2519}\moddef{Commands: cmd expect: TBP}\plusendmoddef
  procedure :: compile => cmd_expect_compile
\nwendcode{}\nwbegincode{2520}\moddef{Commands: procedures}\plusendmoddef
  subroutine cmd_expect_compile (cmd, global)
    class(cmd_expect_t), intent(inout) :: cmd
    type(rt_data_t), intent(inout), target :: global
    type(parse_node_t), pointer :: pn_arg
    pn_arg => parse_node_get_sub_ptr (cmd%pn, 2)
    cmd%pn_opt => parse_node_get_next_ptr (pn_arg)
    cmd%pn_lexpr => parse_node_get_sub_ptr (pn_arg)
    call cmd%compile_options (global)
  end subroutine cmd_expect_compile

\eatline
\nwendcode{}\nwbegindocs{2521}Execute.  Evaluate both arguments, print them and their difference
(if numerical), and whether they agree.  Record the result.
\nwenddocs{}\nwbegincode{2522}\moddef{Commands: cmd expect: TBP}\plusendmoddef
  procedure :: execute => cmd_expect_execute
\nwendcode{}\nwbegincode{2523}\moddef{Commands: procedures}\plusendmoddef
  subroutine cmd_expect_execute (cmd, global)
    class(cmd_expect_t), intent(inout) :: cmd
    type(rt_data_t), intent(inout), target :: global
    type(var_list_t), pointer :: var_list
    logical :: success, is_known
    var_list => cmd%local%get_var_list_ptr ()
    success = eval_log (cmd%pn_lexpr, var_list, is_known=is_known)
    if (is_known) then
       if (success) then
          call msg_message ("expect: success")
       else
          call msg_error ("expect: failure")
       end if
    else
       call msg_error ("expect: undefined result")
       success = .false.
    end if
    call expect_record (success)
  end subroutine cmd_expect_execute

\eatline
\nwendcode{}\nwbegindocs{2524}\nwdocspar
\subsubsection{Beams}
The beam command includes both beam and structure-function
definition.
\nwenddocs{}\nwbegincode{2525}\moddef{Commands: types}\plusendmoddef
  type, extends (command_t) :: cmd_beams_t
     private
     integer :: n_in = 0
     type(parse_node_p), dimension(:), allocatable :: pn_pdg
     integer :: n_sf_record = 0
     integer, dimension(:), allocatable :: n_entry
     type(parse_node_p), dimension(:,:), allocatable :: pn_sf_entry
   contains
   \LA{}Commands: cmd beams: TBP\RA{}
  end type cmd_beams_t

\eatline
\nwendcode{}\nwbegindocs{2526}Output.  The particle expressions are not resolved.
\nwenddocs{}\nwbegincode{2527}\moddef{Commands: cmd beams: TBP}\endmoddef
  procedure :: write => cmd_beams_write
\nwendcode{}\nwbegincode{2528}\moddef{Commands: procedures}\plusendmoddef
  subroutine cmd_beams_write (cmd, unit, indent)
    class(cmd_beams_t), intent(in) :: cmd
    integer, intent(in), optional :: unit, indent
    integer :: u
    u = given_output_unit (unit);  if (u < 0)  return
    call write_indent (u, indent)
    select case (cmd%n_in)
    case (1)
       write (u, "(1x,A)")  "beams: 1 [decay]"
    case (2)
       write (u, "(1x,A)")  "beams: 2 [scattering]"
    case default
       write (u, "(1x,A)")  "beams: [undefined]"
    end select
    if (allocated (cmd%n_entry)) then
       if (cmd%n_sf_record > 0) then
          write (u, "(1x,A,99(1x,I0))")  "structure function entries:", &
               cmd%n_entry
       end if
    end if
  end subroutine cmd_beams_write

\eatline
\nwendcode{}\nwbegindocs{2529}Compile.  Find and assign the parse nodes.

Note: local environments are not yet supported.
\nwenddocs{}\nwbegincode{2530}\moddef{Commands: cmd beams: TBP}\plusendmoddef
  procedure :: compile => cmd_beams_compile
\nwendcode{}\nwbegincode{2531}\moddef{Commands: procedures}\plusendmoddef
  subroutine cmd_beams_compile (cmd, global)
    class(cmd_beams_t), intent(inout) :: cmd
    type(rt_data_t), intent(inout), target :: global
    type(parse_node_t), pointer :: pn_beam_def, pn_beam_spec
    type(parse_node_t), pointer :: pn_beam_list
    type(parse_node_t), pointer :: pn_codes
    type(parse_node_t), pointer :: pn_strfun_seq, pn_strfun_pair
    type(parse_node_t), pointer :: pn_strfun_def
    integer :: i
    pn_beam_def => parse_node_get_sub_ptr (cmd%pn, 3)
    pn_beam_spec => parse_node_get_sub_ptr (pn_beam_def)
    pn_strfun_seq => parse_node_get_next_ptr (pn_beam_spec)
    pn_beam_list => parse_node_get_sub_ptr (pn_beam_spec)
    call cmd%compile_options (global)
    cmd%n_in = parse_node_get_n_sub (pn_beam_list)
    allocate (cmd%pn_pdg (cmd%n_in))
    pn_codes => parse_node_get_sub_ptr (pn_beam_list)
    do i = 1, cmd%n_in
       cmd%pn_pdg(i)%ptr => pn_codes
       pn_codes => parse_node_get_next_ptr (pn_codes)
    end do
    if (associated (pn_strfun_seq)) then
       cmd%n_sf_record = parse_node_get_n_sub (pn_beam_def) - 1
       allocate (cmd%n_entry (cmd%n_sf_record), source = 1)
       allocate (cmd%pn_sf_entry (2, cmd%n_sf_record))
       do i = 1, cmd%n_sf_record
          pn_strfun_pair => parse_node_get_sub_ptr (pn_strfun_seq, 2)
          pn_strfun_def => parse_node_get_sub_ptr (pn_strfun_pair)
          cmd%pn_sf_entry(1,i)%ptr => pn_strfun_def
          pn_strfun_def => parse_node_get_next_ptr (pn_strfun_def)
          cmd%pn_sf_entry(2,i)%ptr => pn_strfun_def
          if (associated (pn_strfun_def))  cmd%n_entry(i) = 2
          pn_strfun_seq => parse_node_get_next_ptr (pn_strfun_seq)
       end do
    else
       allocate (cmd%n_entry (0))
       allocate (cmd%pn_sf_entry (0, 0))
    end if
  end subroutine cmd_beams_compile

\eatline
\nwendcode{}\nwbegindocs{2532}Command execution: Determine beam particles and structure-function
names, if any.  The results are stored in the {\tt{}beam{\char95}structure}
component of the {\tt{}global} data block.
\nwenddocs{}\nwbegincode{2533}\moddef{Commands: cmd beams: TBP}\plusendmoddef
  procedure :: execute => cmd_beams_execute
\nwendcode{}\nwbegincode{2534}\moddef{Commands: procedures}\plusendmoddef
  subroutine cmd_beams_execute (cmd, global)
    class(cmd_beams_t), intent(inout) :: cmd
    type(rt_data_t), intent(inout), target :: global
    type(var_list_t), pointer :: var_list
    type(pdg_array_t) :: pdg_array
    integer, dimension(:), allocatable :: pdg
    type(flavor_t), dimension(:), allocatable :: flv
    type(parse_node_t), pointer :: pn_key
    type(string_t) :: sf_name
    integer :: i, j
    call lhapdf_global_reset ()
    var_list => cmd%local%get_var_list_ptr ()
    allocate (flv (cmd%n_in))
    do i = 1, cmd%n_in
       pdg_array = eval_pdg_array (cmd%pn_pdg(i)%ptr, var_list)
       pdg = pdg_array
       select case (size (pdg))
       case (1)
          call flv(i)%init ( pdg(1), cmd%local%model)
       case default
          call msg_fatal ("Beams: beam particles must be unique")
       end select
    end do
    select case (cmd%n_in)
    case (1)
       if (cmd%n_sf_record > 0) then
          call msg_fatal ("Beam setup: no structure functions allowed &
               &for decay")
       end if
       call global%beam_structure%init_sf (flv%get_name ())
    case (2)
       call global%beam_structure%init_sf (flv%get_name (), cmd%n_entry)
       do i = 1, cmd%n_sf_record
          do j = 1, cmd%n_entry(i)
             pn_key => parse_node_get_sub_ptr (cmd%pn_sf_entry(j,i)%ptr)
             sf_name = parse_node_get_key (pn_key)
             call global%beam_structure%set_sf (i, j, sf_name)
          end do
       end do
    end select
  end subroutine cmd_beams_execute

\eatline
\nwendcode{}\nwbegindocs{2535}\nwdocspar
\subsubsection{Density matrices for beam polarization}
For holding beam polarization, we define a notation and a data
structure for sparse matrices.  The entries (and the index
expressions) are numerical expressions, so we use evaluation trees.

Each entry in the sparse matrix is an n-tuple of expressions.  The first
tuple elements represent index values, the last one is an arbitrary
(complex) number.  Absent expressions are replaced by default-value rules.

Note: Here, and in some other commands, we would like to store an evaluation
tree, not just a parse node pointer.  However, the current expression handler
wants all variables defined, so the evaluation tree can only be built by
{\tt{}evaluate}, i.e., compiled just-in-time and evaluated immediately.
\nwenddocs{}\nwbegincode{2536}\moddef{Commands: types}\plusendmoddef
  type :: sentry_expr_t
     type(parse_node_p), dimension(:), allocatable :: expr
   contains
   \LA{}Commands: sentry expr: TBP\RA{}
  end type sentry_expr_t

\eatline
\nwendcode{}\nwbegindocs{2537}Compile parse nodes into evaluation trees.
\nwenddocs{}\nwbegincode{2538}\moddef{Commands: sentry expr: TBP}\endmoddef
  procedure :: compile => sentry_expr_compile
\nwendcode{}\nwbegincode{2539}\moddef{Commands: procedures}\plusendmoddef
  subroutine sentry_expr_compile (sentry, pn)
    class(sentry_expr_t), intent(out) :: sentry
    type(parse_node_t), intent(in), target :: pn
    type(parse_node_t), pointer :: pn_expr, pn_extra
    integer :: n_expr, i
    n_expr = parse_node_get_n_sub (pn)
    allocate (sentry%expr (n_expr))
    if (n_expr > 0) then
       i = 0
       pn_expr => parse_node_get_sub_ptr (pn)
       pn_extra => parse_node_get_next_ptr (pn_expr)
       do i = 1, n_expr
          sentry%expr(i)%ptr => pn_expr
          if (associated (pn_extra)) then
             pn_expr => parse_node_get_sub_ptr (pn_extra, 2)
             pn_extra => parse_node_get_next_ptr (pn_extra)
          end if
       end do
    end if
  end subroutine sentry_expr_compile
    
\eatline
\nwendcode{}\nwbegindocs{2540}Evaluate the expressions and return an index array of predefined
length together with a complex value.  If the value (as the last expression)
is undefined, set it to unity.  If index values are undefined, repeat
the previous index value.
\nwenddocs{}\nwbegincode{2541}\moddef{Commands: sentry expr: TBP}\plusendmoddef
  procedure :: evaluate => sentry_expr_evaluate
\nwendcode{}\nwbegincode{2542}\moddef{Commands: procedures}\plusendmoddef
  subroutine sentry_expr_evaluate (sentry, index, value, global)
    class(sentry_expr_t), intent(inout) :: sentry
    integer, dimension(:), intent(out) :: index
    complex(default), intent(out) :: value
    type(rt_data_t), intent(in), target :: global
    type(var_list_t), pointer :: var_list
    integer :: i, n_expr, n_index
    type(eval_tree_t) :: eval_tree
    var_list => global%get_var_list_ptr ()
    n_expr = size (sentry%expr) 
    n_index = size (index)
    if (n_expr <= n_index + 1) then
       do i = 1, min (n_expr, n_index)
          associate (expr => sentry%expr(i))
            call eval_tree%init_expr (expr%ptr, var_list)
            call eval_tree%evaluate ()
            if (eval_tree%is_known ()) then
               index(i) = eval_tree%get_int ()
            else
               call msg_fatal ("Evaluating density matrix: undefined index")
            end if
          end associate
       end do
       do i = n_expr + 1, n_index
          index(i) = index(n_expr)
       end do
       if (n_expr == n_index + 1) then
          associate (expr => sentry%expr(n_expr))
            call eval_tree%init_expr (expr%ptr, var_list)
            call eval_tree%evaluate ()
            if (eval_tree%is_known ()) then
               value = eval_tree%get_cmplx ()
            else
               call msg_fatal ("Evaluating density matrix: undefined index")
            end if
            call eval_tree%final ()
          end associate
       else
          value = 1
       end if
    else
       call msg_fatal ("Evaluating density matrix: index expression too long")
    end if
  end subroutine sentry_expr_evaluate

\eatline
\nwendcode{}\nwbegindocs{2543}The sparse matrix itself consists of an arbitrary number of entries.
\nwenddocs{}\nwbegincode{2544}\moddef{Commands: types}\plusendmoddef
  type :: smatrix_expr_t
     type(sentry_expr_t), dimension(:), allocatable :: entry
   contains
   \LA{}Commands: smatrix expr: TBP\RA{}
  end type smatrix_expr_t
  
\eatline
\nwendcode{}\nwbegindocs{2545}Compile: assign sub-nodes to sentry-expressions and compile those.
\nwenddocs{}\nwbegincode{2546}\moddef{Commands: smatrix expr: TBP}\endmoddef
  procedure :: compile => smatrix_expr_compile
\nwendcode{}\nwbegincode{2547}\moddef{Commands: procedures}\plusendmoddef
  subroutine smatrix_expr_compile (smatrix_expr, pn)
    class(smatrix_expr_t), intent(out) :: smatrix_expr
    type(parse_node_t), intent(in), target :: pn
    type(parse_node_t), pointer :: pn_arg, pn_entry
    integer :: n_entry, i
    pn_arg => parse_node_get_sub_ptr (pn, 2)
    if (associated (pn_arg)) then
       n_entry = parse_node_get_n_sub (pn_arg)
       allocate (smatrix_expr%entry (n_entry))
       pn_entry => parse_node_get_sub_ptr (pn_arg)
       do i = 1, n_entry
          call smatrix_expr%entry(i)%compile (pn_entry)
          pn_entry => parse_node_get_next_ptr (pn_entry)
       end do
    else
       allocate (smatrix_expr%entry (0))
    end if
  end subroutine smatrix_expr_compile

\eatline
\nwendcode{}\nwbegindocs{2548}Evaluate the entries and build a new {\tt{}smatrix} object, which
contains just the numerical results.
\nwenddocs{}\nwbegincode{2549}\moddef{Commands: smatrix expr: TBP}\plusendmoddef
  procedure :: evaluate => smatrix_expr_evaluate
\nwendcode{}\nwbegincode{2550}\moddef{Commands: procedures}\plusendmoddef
  subroutine smatrix_expr_evaluate (smatrix_expr, smatrix, global)
    class(smatrix_expr_t), intent(inout) :: smatrix_expr
    type(smatrix_t), intent(out) :: smatrix
    type(rt_data_t), intent(in), target :: global
    integer, dimension(2) :: idx
    complex(default) :: value
    integer :: i, n_entry
    n_entry = size (smatrix_expr%entry)
    call smatrix%init (2, n_entry)
    do i = 1, n_entry
       call smatrix_expr%entry(i)%evaluate (idx, value, global)
       call smatrix%set_entry (i, idx, value)
    end do
  end subroutine smatrix_expr_evaluate
    
\eatline
\nwendcode{}\nwbegindocs{2551}\nwdocspar
\subsubsection{Beam polarization density}

The beam polarization command defines spin density matrix for one or
two beams (scattering or decay).
\nwenddocs{}\nwbegincode{2552}\moddef{Commands: types}\plusendmoddef
  type, extends (command_t) :: cmd_beams_pol_density_t
     private
     integer :: n_in = 0
     type(smatrix_expr_t), dimension(:), allocatable :: smatrix
   contains
   \LA{}Commands: cmd beams pol density: TBP\RA{}
  end type cmd_beams_pol_density_t

\eatline
\nwendcode{}\nwbegindocs{2553}Output.
\nwenddocs{}\nwbegincode{2554}\moddef{Commands: cmd beams pol density: TBP}\endmoddef
  procedure :: write => cmd_beams_pol_density_write
\nwendcode{}\nwbegincode{2555}\moddef{Commands: procedures}\plusendmoddef
  subroutine cmd_beams_pol_density_write (cmd, unit, indent)
    class(cmd_beams_pol_density_t), intent(in) :: cmd
    integer, intent(in), optional :: unit, indent
    integer :: u
    u = given_output_unit (unit);  if (u < 0)  return
    call write_indent (u, indent)
    select case (cmd%n_in)
    case (1)
       write (u, "(1x,A)")  "beams polarization setup: 1 [decay]"
    case (2)
       write (u, "(1x,A)")  "beams polarization setup: 2 [scattering]"
    case default
       write (u, "(1x,A)")  "beams polarization setup: [undefined]"
    end select
  end subroutine cmd_beams_pol_density_write

\eatline
\nwendcode{}\nwbegindocs{2556}Compile.  Find and assign the parse nodes.

Note: local environments are not yet supported.
\nwenddocs{}\nwbegincode{2557}\moddef{Commands: cmd beams pol density: TBP}\plusendmoddef
  procedure :: compile => cmd_beams_pol_density_compile
\nwendcode{}\nwbegincode{2558}\moddef{Commands: procedures}\plusendmoddef
  subroutine cmd_beams_pol_density_compile (cmd, global)
    class(cmd_beams_pol_density_t), intent(inout) :: cmd
    type(rt_data_t), intent(inout), target :: global
    type(parse_node_t), pointer :: pn_pol_spec, pn_smatrix
    integer :: i
    pn_pol_spec => parse_node_get_sub_ptr (cmd%pn, 3)
    call cmd%compile_options (global)
    cmd%n_in = parse_node_get_n_sub (pn_pol_spec)
    allocate (cmd%smatrix (cmd%n_in))
    pn_smatrix => parse_node_get_sub_ptr (pn_pol_spec)
    do i = 1, cmd%n_in
       call cmd%smatrix(i)%compile (pn_smatrix)
       pn_smatrix => parse_node_get_next_ptr (pn_smatrix)
    end do
  end subroutine cmd_beams_pol_density_compile

\eatline
\nwendcode{}\nwbegindocs{2559}Command execution: Fill polarization density matrices.  No check
yet, the matrices are checked and normalized when the actual beam
object is created, just before integration.  For intermediate storage,
we use the {\tt{}beam{\char95}structure} object in the {\tt{}global} data set.
\nwenddocs{}\nwbegincode{2560}\moddef{Commands: cmd beams pol density: TBP}\plusendmoddef
  procedure :: execute => cmd_beams_pol_density_execute
\nwendcode{}\nwbegincode{2561}\moddef{Commands: procedures}\plusendmoddef
  subroutine cmd_beams_pol_density_execute (cmd, global)
    class(cmd_beams_pol_density_t), intent(inout) :: cmd
    type(rt_data_t), intent(inout), target :: global
    type(smatrix_t) :: smatrix
    integer :: i
    call global%beam_structure%init_pol (cmd%n_in)
    do i = 1, cmd%n_in
       call cmd%smatrix(i)%evaluate (smatrix, global)
       call global%beam_structure%set_smatrix (i, smatrix)
    end do
  end subroutine cmd_beams_pol_density_execute

\eatline
\nwendcode{}\nwbegindocs{2562}\nwdocspar
\subsubsection{Beam polarization fraction}
In addition to the polarization density matrix, we can independently
specify the polarization fraction for one or both beams.
\nwenddocs{}\nwbegincode{2563}\moddef{Commands: types}\plusendmoddef
  type, extends (command_t) :: cmd_beams_pol_fraction_t
     private
     integer :: n_in = 0
     type(parse_node_p), dimension(:), allocatable :: expr
   contains
   \LA{}Commands: cmd beams pol fraction: TBP\RA{}
  end type cmd_beams_pol_fraction_t

\eatline
\nwendcode{}\nwbegindocs{2564}Output.
\nwenddocs{}\nwbegincode{2565}\moddef{Commands: cmd beams pol fraction: TBP}\endmoddef
  procedure :: write => cmd_beams_pol_fraction_write
\nwendcode{}\nwbegincode{2566}\moddef{Commands: procedures}\plusendmoddef
  subroutine cmd_beams_pol_fraction_write (cmd, unit, indent)
    class(cmd_beams_pol_fraction_t), intent(in) :: cmd
    integer, intent(in), optional :: unit, indent
    integer :: u
    u = given_output_unit (unit);  if (u < 0)  return
    call write_indent (u, indent)
    select case (cmd%n_in)
    case (1)
       write (u, "(1x,A)")  "beams polarization fraction: 1 [decay]"
    case (2)
       write (u, "(1x,A)")  "beams polarization fraction: 2 [scattering]"
    case default
       write (u, "(1x,A)")  "beams polarization fraction: [undefined]"
    end select
  end subroutine cmd_beams_pol_fraction_write

\eatline
\nwendcode{}\nwbegindocs{2567}Compile.  Find and assign the parse nodes.

Note: local environments are not yet supported.
\nwenddocs{}\nwbegincode{2568}\moddef{Commands: cmd beams pol fraction: TBP}\plusendmoddef
  procedure :: compile => cmd_beams_pol_fraction_compile
\nwendcode{}\nwbegincode{2569}\moddef{Commands: procedures}\plusendmoddef
  subroutine cmd_beams_pol_fraction_compile (cmd, global)
    class(cmd_beams_pol_fraction_t), intent(inout) :: cmd
    type(rt_data_t), intent(inout), target :: global
    type(parse_node_t), pointer :: pn_frac_spec, pn_expr
    integer :: i
    pn_frac_spec => parse_node_get_sub_ptr (cmd%pn, 3)
    call cmd%compile_options (global)
    cmd%n_in = parse_node_get_n_sub (pn_frac_spec)
    allocate (cmd%expr (cmd%n_in))
    pn_expr => parse_node_get_sub_ptr (pn_frac_spec)
    do i = 1, cmd%n_in
       cmd%expr(i)%ptr => pn_expr
       pn_expr => parse_node_get_next_ptr (pn_expr)
    end do
  end subroutine cmd_beams_pol_fraction_compile

\eatline
\nwendcode{}\nwbegindocs{2570}Command execution: Retrieve the numerical values of the beam
polarization fractions.  The results are stored in the
{\tt{}beam{\char95}structure} component of the {\tt{}global} data block.
\nwenddocs{}\nwbegincode{2571}\moddef{Commands: cmd beams pol fraction: TBP}\plusendmoddef
  procedure :: execute => cmd_beams_pol_fraction_execute
\nwendcode{}\nwbegincode{2572}\moddef{Commands: procedures}\plusendmoddef
  subroutine cmd_beams_pol_fraction_execute (cmd, global)
    class(cmd_beams_pol_fraction_t), intent(inout) :: cmd
    type(rt_data_t), intent(inout), target :: global
    type(var_list_t), pointer :: var_list
    real(default), dimension(:), allocatable :: pol_f
    type(eval_tree_t) :: expr
    integer :: i
    var_list => global%get_var_list_ptr ()
    allocate (pol_f (cmd%n_in))
    do i = 1, cmd%n_in
       call expr%init_expr (cmd%expr(i)%ptr, var_list)
       call expr%evaluate ()
       if (expr%is_known ()) then
          pol_f(i) = expr%get_real ()
       else
          call msg_fatal ("beams polarization fraction: undefined value")
       end if
       call expr%final ()
    end do
    call global%beam_structure%set_pol_f (pol_f)
  end subroutine cmd_beams_pol_fraction_execute

\eatline
\nwendcode{}\nwbegindocs{2573}\nwdocspar
\subsubsection{Beam momentum}
This is completely analogous to the previous command, hence we can use
inheritance.
\nwenddocs{}\nwbegincode{2574}\moddef{Commands: types}\plusendmoddef
  type, extends (cmd_beams_pol_fraction_t) :: cmd_beams_momentum_t
   contains
   \LA{}Commands: cmd beams momentum: TBP\RA{}
  end type cmd_beams_momentum_t

\eatline
\nwendcode{}\nwbegindocs{2575}Output.
\nwenddocs{}\nwbegincode{2576}\moddef{Commands: cmd beams momentum: TBP}\endmoddef
  procedure :: write => cmd_beams_momentum_write
\nwendcode{}\nwbegincode{2577}\moddef{Commands: procedures}\plusendmoddef
  subroutine cmd_beams_momentum_write (cmd, unit, indent)
    class(cmd_beams_momentum_t), intent(in) :: cmd
    integer, intent(in), optional :: unit, indent
    integer :: u
    u = given_output_unit (unit);  if (u < 0)  return
    call write_indent (u, indent)
    select case (cmd%n_in)
    case (1)
       write (u, "(1x,A)")  "beams momentum: 1 [decay]"
    case (2)
       write (u, "(1x,A)")  "beams momentum: 2 [scattering]"
    case default
       write (u, "(1x,A)")  "beams momentum: [undefined]"
    end select
  end subroutine cmd_beams_momentum_write

\eatline
\nwendcode{}\nwbegindocs{2578}Compile: inherited.

Command execution: Not inherited, but just the error string and the final
command are changed.
\nwenddocs{}\nwbegincode{2579}\moddef{Commands: cmd beams momentum: TBP}\plusendmoddef
  procedure :: execute => cmd_beams_momentum_execute
\nwendcode{}\nwbegincode{2580}\moddef{Commands: procedures}\plusendmoddef
  subroutine cmd_beams_momentum_execute (cmd, global)
    class(cmd_beams_momentum_t), intent(inout) :: cmd
    type(rt_data_t), intent(inout), target :: global
    type(var_list_t), pointer :: var_list
    real(default), dimension(:), allocatable :: p
    type(eval_tree_t) :: expr
    integer :: i
    var_list => global%get_var_list_ptr ()
    allocate (p (cmd%n_in))
    do i = 1, cmd%n_in
       call expr%init_expr (cmd%expr(i)%ptr, var_list)
       call expr%evaluate ()
       if (expr%is_known ()) then
          p(i) = expr%get_real ()
       else
          call msg_fatal ("beams momentum: undefined value")
       end if
       call expr%final ()
    end do
    call global%beam_structure%set_momentum (p)
  end subroutine cmd_beams_momentum_execute

\eatline
\nwendcode{}\nwbegindocs{2581}\nwdocspar
\subsubsection{Beam angles}
Again, this is analogous.  There are two angles, polar angle $\theta$
and azimuthal angle $\phi$, which can be set independently for both beams.
\nwenddocs{}\nwbegincode{2582}\moddef{Commands: types}\plusendmoddef
  type, extends (cmd_beams_pol_fraction_t) :: cmd_beams_theta_t
   contains
   \LA{}Commands: cmd beams theta: TBP\RA{}
  end type cmd_beams_theta_t

  type, extends (cmd_beams_pol_fraction_t) :: cmd_beams_phi_t
   contains
   \LA{}Commands: cmd beams phi: TBP\RA{}
  end type cmd_beams_phi_t

\eatline
\eatline
\nwendcode{}\nwbegindocs{2583}Output.
\nwenddocs{}\nwbegincode{2584}\moddef{Commands: cmd beams theta: TBP}\endmoddef
  procedure :: write => cmd_beams_theta_write
\nwendcode{}\nwbegincode{2585}\moddef{Commands: cmd beams phi: TBP}\endmoddef
  procedure :: write => cmd_beams_phi_write
\nwendcode{}\nwbegincode{2586}\moddef{Commands: procedures}\plusendmoddef
  subroutine cmd_beams_theta_write (cmd, unit, indent)
    class(cmd_beams_theta_t), intent(in) :: cmd
    integer, intent(in), optional :: unit, indent
    integer :: u
    u = given_output_unit (unit);  if (u < 0)  return
    call write_indent (u, indent)
    select case (cmd%n_in)
    case (1)
       write (u, "(1x,A)")  "beams theta: 1 [decay]"
    case (2)
       write (u, "(1x,A)")  "beams theta: 2 [scattering]"
    case default
       write (u, "(1x,A)")  "beams theta: [undefined]"
    end select
  end subroutine cmd_beams_theta_write

  subroutine cmd_beams_phi_write (cmd, unit, indent)
    class(cmd_beams_phi_t), intent(in) :: cmd
    integer, intent(in), optional :: unit, indent
    integer :: u
    u = given_output_unit (unit);  if (u < 0)  return
    call write_indent (u, indent)
    select case (cmd%n_in)
    case (1)
       write (u, "(1x,A)")  "beams phi: 1 [decay]"
    case (2)
       write (u, "(1x,A)")  "beams phi: 2 [scattering]"
    case default
       write (u, "(1x,A)")  "beams phi: [undefined]"
    end select
  end subroutine cmd_beams_phi_write

\eatline
\eatline
\nwendcode{}\nwbegindocs{2587}Compile: inherited.

Command execution: Not inherited, but just the error string and the final
command are changed.
\nwenddocs{}\nwbegincode{2588}\moddef{Commands: cmd beams theta: TBP}\plusendmoddef
  procedure :: execute => cmd_beams_theta_execute
\nwendcode{}\nwbegincode{2589}\moddef{Commands: cmd beams phi: TBP}\plusendmoddef
  procedure :: execute => cmd_beams_phi_execute
\nwendcode{}\nwbegincode{2590}\moddef{Commands: procedures}\plusendmoddef
  subroutine cmd_beams_theta_execute (cmd, global)
    class(cmd_beams_theta_t), intent(inout) :: cmd
    type(rt_data_t), intent(inout), target :: global
    type(var_list_t), pointer :: var_list
    real(default), dimension(:), allocatable :: theta
    type(eval_tree_t) :: expr
    integer :: i
    var_list => global%get_var_list_ptr ()
    allocate (theta (cmd%n_in))
    do i = 1, cmd%n_in
       call expr%init_expr (cmd%expr(i)%ptr, var_list)
       call expr%evaluate ()
       if (expr%is_known ()) then
          theta(i) = expr%get_real ()
       else
          call msg_fatal ("beams theta: undefined value")
       end if
       call expr%final ()
    end do
    call global%beam_structure%set_theta (theta)
  end subroutine cmd_beams_theta_execute

  subroutine cmd_beams_phi_execute (cmd, global)
    class(cmd_beams_phi_t), intent(inout) :: cmd
    type(rt_data_t), intent(inout), target :: global
    type(var_list_t), pointer :: var_list
    real(default), dimension(:), allocatable :: phi
    type(eval_tree_t) :: expr
    integer :: i
    var_list => global%get_var_list_ptr ()
    allocate (phi (cmd%n_in))
    do i = 1, cmd%n_in
       call expr%init_expr (cmd%expr(i)%ptr, var_list)
       call expr%evaluate ()
       if (expr%is_known ()) then
          phi(i) = expr%get_real ()
       else
          call msg_fatal ("beams phi: undefined value")
       end if
       call expr%final ()
    end do
    call global%beam_structure%set_phi (phi)
  end subroutine cmd_beams_phi_execute

\eatline
\eatline
\nwendcode{}\nwbegindocs{2591}\nwdocspar
\subsubsection{Cuts}
Define a cut expression.  We store the parse tree for the right-hand
side instead of compiling it.  Compilation is deferred to the process
environment where the cut expression is used.
\nwenddocs{}\nwbegincode{2592}\moddef{Commands: types}\plusendmoddef
  type, extends (command_t) :: cmd_cuts_t
     private
     type(parse_node_t), pointer :: pn_lexpr => null ()
   contains
   \LA{}Commands: cmd cuts: TBP\RA{}
  end type cmd_cuts_t

\eatline
\nwendcode{}\nwbegindocs{2593}Output.  Do not print the parse tree, since this may get cluttered.
Just a message that cuts have been defined.
\nwenddocs{}\nwbegincode{2594}\moddef{Commands: cmd cuts: TBP}\endmoddef
  procedure :: write => cmd_cuts_write
\nwendcode{}\nwbegincode{2595}\moddef{Commands: procedures}\plusendmoddef
  subroutine cmd_cuts_write (cmd, unit, indent)
    class(cmd_cuts_t), intent(in) :: cmd
    integer, intent(in), optional :: unit, indent
    integer :: u
    u = given_output_unit (unit);  if (u < 0)  return
    call write_indent (u, indent)
    write (u, "(1x,A)")  "cuts: [defined]"
  end subroutine cmd_cuts_write
  
\eatline
\nwendcode{}\nwbegindocs{2596}Compile.  Simply store the parse (sub)tree.
\nwenddocs{}\nwbegincode{2597}\moddef{Commands: cmd cuts: TBP}\plusendmoddef
  procedure :: compile => cmd_cuts_compile
\nwendcode{}\nwbegincode{2598}\moddef{Commands: procedures}\plusendmoddef
  subroutine cmd_cuts_compile (cmd, global)
    class(cmd_cuts_t), intent(inout) :: cmd
    type(rt_data_t), intent(inout), target :: global
    cmd%pn_lexpr => parse_node_get_sub_ptr (cmd%pn, 3)
  end subroutine cmd_cuts_compile

\eatline
\nwendcode{}\nwbegindocs{2599}Instead of evaluating the cut expression, link the parse tree to the
global data set, such that it is compiled and executed in the
appropriate process context.
\nwenddocs{}\nwbegincode{2600}\moddef{Commands: cmd cuts: TBP}\plusendmoddef
  procedure :: execute => cmd_cuts_execute
\nwendcode{}\nwbegincode{2601}\moddef{Commands: procedures}\plusendmoddef
  subroutine cmd_cuts_execute (cmd, global)
    class(cmd_cuts_t), intent(inout) :: cmd
    type(rt_data_t), intent(inout), target :: global
    global%pn%cuts_lexpr => cmd%pn_lexpr
  end subroutine cmd_cuts_execute

\eatline
\nwendcode{}\nwbegindocs{2602}\nwdocspar
\subsubsection{General, Factorization and Renormalization Scales}
Define a scale expression for either the renormalization or the 
factorization scale.  We store the parse tree for the right-hand
side instead of compiling it.  Compilation is deferred to the process
environment where the expression is used.
\nwenddocs{}\nwbegincode{2603}\moddef{Commands: types}\plusendmoddef
  type, extends (command_t) :: cmd_scale_t
     private
     type(parse_node_t), pointer :: pn_expr => null ()     
   contains
   \LA{}Commands: cmd scale: TBP\RA{}   
  end type cmd_scale_t

\eatline
\nwendcode{}\nwbegincode{2604}\moddef{Commands: types}\plusendmoddef
  type, extends (command_t) :: cmd_fac_scale_t
     private
     type(parse_node_t), pointer :: pn_expr => null ()     
   contains
   \LA{}Commands: cmd fac scale: TBP\RA{}
  end type cmd_fac_scale_t

\eatline
\nwendcode{}\nwbegincode{2605}\moddef{Commands: types}\plusendmoddef
  type, extends (command_t) :: cmd_ren_scale_t
     private
     type(parse_node_t), pointer :: pn_expr => null ()     
   contains
   \LA{}Commands: cmd ren scale: TBP\RA{}
  end type cmd_ren_scale_t

\eatline
\nwendcode{}\nwbegindocs{2606}Output. Do not print the parse tree, since this may get cluttered.
Just a message that scale, renormalization and factorization have been
defined, respectively.
\nwenddocs{}\nwbegincode{2607}\moddef{Commands: cmd scale: TBP}\endmoddef
  procedure :: write => cmd_scale_write
\nwendcode{}\nwbegincode{2608}\moddef{Commands: procedures}\plusendmoddef
  subroutine cmd_scale_write (cmd, unit, indent)
    class(cmd_scale_t), intent(in) :: cmd
    integer, intent(in), optional :: unit, indent
    integer :: u
    u = given_output_unit (unit);  if (u < 0)  return
    call write_indent (u, indent)
    write (u, "(1x,A)")  "scale: [defined]"
  end subroutine cmd_scale_write
  
\eatline
\nwendcode{}\nwbegindocs{2609}\nwdocspar
\nwenddocs{}\nwbegincode{2610}\moddef{Commands: cmd fac scale: TBP}\endmoddef
  procedure :: write => cmd_fac_scale_write
\nwendcode{}\nwbegincode{2611}\moddef{Commands: procedures}\plusendmoddef
  subroutine cmd_fac_scale_write (cmd, unit, indent)
    class(cmd_fac_scale_t), intent(in) :: cmd
    integer, intent(in), optional :: unit, indent
    integer :: u
    u = given_output_unit (unit);  if (u < 0)  return
    call write_indent (u, indent)
    write (u, "(1x,A)")  "factorization scale: [defined]"
  end subroutine cmd_fac_scale_write
  
\eatline
\nwendcode{}\nwbegindocs{2612}\nwdocspar
\nwenddocs{}\nwbegincode{2613}\moddef{Commands: cmd ren scale: TBP}\endmoddef
  procedure :: write => cmd_ren_scale_write
\nwendcode{}\nwbegincode{2614}\moddef{Commands: procedures}\plusendmoddef
  subroutine cmd_ren_scale_write (cmd, unit, indent)
    class(cmd_ren_scale_t), intent(in) :: cmd
    integer, intent(in), optional :: unit, indent
    integer :: u
    u = given_output_unit (unit);  if (u < 0)  return
    call write_indent (u, indent)
    write (u, "(1x,A)")  "renormalization scale: [defined]"
  end subroutine cmd_ren_scale_write
  
\eatline
\nwendcode{}\nwbegindocs{2615}Compile.  Simply store the parse (sub)tree.
\nwenddocs{}\nwbegincode{2616}\moddef{Commands: cmd scale: TBP}\plusendmoddef
  procedure :: compile => cmd_scale_compile 
\nwendcode{}\nwbegincode{2617}\moddef{Commands: procedures}\plusendmoddef
  subroutine cmd_scale_compile (cmd, global)
    class(cmd_scale_t), intent(inout) :: cmd
    type(rt_data_t), intent(inout), target :: global
    cmd%pn_expr => parse_node_get_sub_ptr (cmd%pn, 3)
  end subroutine cmd_scale_compile

\eatline
\nwendcode{}\nwbegindocs{2618}\nwdocspar
\nwenddocs{}\nwbegincode{2619}\moddef{Commands: cmd fac scale: TBP}\plusendmoddef
  procedure :: compile => cmd_fac_scale_compile
\nwendcode{}\nwbegincode{2620}\moddef{Commands: procedures}\plusendmoddef
  subroutine cmd_fac_scale_compile (cmd, global)
    class(cmd_fac_scale_t), intent(inout) :: cmd
    type(rt_data_t), intent(inout), target :: global
    cmd%pn_expr => parse_node_get_sub_ptr (cmd%pn, 3)
  end subroutine cmd_fac_scale_compile

\eatline
\nwendcode{}\nwbegindocs{2621}\nwdocspar
\nwenddocs{}\nwbegincode{2622}\moddef{Commands: cmd ren scale: TBP}\plusendmoddef
  procedure :: compile => cmd_ren_scale_compile
\nwendcode{}\nwbegincode{2623}\moddef{Commands: procedures}\plusendmoddef
  subroutine cmd_ren_scale_compile (cmd, global)
    class(cmd_ren_scale_t), intent(inout) :: cmd
    type(rt_data_t), intent(inout), target :: global
    cmd%pn_expr => parse_node_get_sub_ptr (cmd%pn, 3)
  end subroutine cmd_ren_scale_compile

\eatline
\nwendcode{}\nwbegindocs{2624}Instead of evaluating the scale expression, link the parse tree to the
global data set, such that it is compiled and executed in the
appropriate process context.
\nwenddocs{}\nwbegincode{2625}\moddef{Commands: cmd scale: TBP}\plusendmoddef
  procedure :: execute => cmd_scale_execute
\nwendcode{}\nwbegincode{2626}\moddef{Commands: procedures}\plusendmoddef
  subroutine cmd_scale_execute (cmd, global)
    class(cmd_scale_t), intent(inout) :: cmd
    type(rt_data_t), intent(inout), target :: global
    global%pn%scale_expr => cmd%pn_expr
  end subroutine cmd_scale_execute

\eatline
\nwendcode{}\nwbegindocs{2627}\nwdocspar
\nwenddocs{}\nwbegincode{2628}\moddef{Commands: cmd fac scale: TBP}\plusendmoddef
  procedure :: execute => cmd_fac_scale_execute
\nwendcode{}\nwbegincode{2629}\moddef{Commands: procedures}\plusendmoddef
  subroutine cmd_fac_scale_execute (cmd, global)
    class(cmd_fac_scale_t), intent(inout) :: cmd
    type(rt_data_t), intent(inout), target :: global
    global%pn%fac_scale_expr => cmd%pn_expr
  end subroutine cmd_fac_scale_execute

\eatline
\nwendcode{}\nwbegindocs{2630}\nwdocspar
\nwenddocs{}\nwbegincode{2631}\moddef{Commands: cmd ren scale: TBP}\plusendmoddef
  procedure :: execute => cmd_ren_scale_execute
\nwendcode{}\nwbegincode{2632}\moddef{Commands: procedures}\plusendmoddef
  subroutine cmd_ren_scale_execute (cmd, global)
    class(cmd_ren_scale_t), intent(inout) :: cmd
    type(rt_data_t), intent(inout), target :: global
    global%pn%ren_scale_expr => cmd%pn_expr    
  end subroutine cmd_ren_scale_execute

\eatline
\nwendcode{}\nwbegindocs{2633}\nwdocspar
\subsubsection{Weight}
Define a weight expression. The weight is applied to a process to be
integrated, event by event. We store the parse tree for the right-hand
side instead of compiling it. Compilation is deferred to the process
environment where the expression is used.
\nwenddocs{}\nwbegincode{2634}\moddef{Commands: types}\plusendmoddef
  type, extends (command_t) :: cmd_weight_t
     private
     type(parse_node_t), pointer :: pn_expr => null ()
   contains
   \LA{}Commands: cmd weight: TBP\RA{}  
  end type cmd_weight_t

\eatline
\nwendcode{}\nwbegindocs{2635}Output. Do not print the parse tree, since this may get cluttered.
Just a message that scale, renormalization and factorization have been
defined, respectively.
\nwenddocs{}\nwbegincode{2636}\moddef{Commands: cmd weight: TBP}\endmoddef
  procedure :: write => cmd_weight_write
\nwendcode{}\nwbegincode{2637}\moddef{Commands: procedures}\plusendmoddef
  subroutine cmd_weight_write (cmd, unit, indent)
    class(cmd_weight_t), intent(in) :: cmd
    integer, intent(in), optional :: unit, indent
    integer :: u
    u = given_output_unit (unit);  if (u < 0)  return
    call write_indent (u, indent)
    write (u, "(1x,A)")  "weight expression: [defined]"
  end subroutine cmd_weight_write
  
\eatline
\nwendcode{}\nwbegindocs{2638}Compile.  Simply store the parse (sub)tree.
\nwenddocs{}\nwbegincode{2639}\moddef{Commands: cmd weight: TBP}\plusendmoddef
  procedure :: compile => cmd_weight_compile
\nwendcode{}\nwbegincode{2640}\moddef{Commands: procedures}\plusendmoddef
  subroutine cmd_weight_compile (cmd, global)
    class(cmd_weight_t), intent(inout) :: cmd
    type(rt_data_t), intent(inout), target :: global
    cmd%pn_expr => parse_node_get_sub_ptr (cmd%pn, 3)
  end subroutine cmd_weight_compile

\eatline
\nwendcode{}\nwbegindocs{2641}Instead of evaluating the expression, link the parse tree to the
global data set, such that it is compiled and executed in the
appropriate process context.
\nwenddocs{}\nwbegincode{2642}\moddef{Commands: cmd weight: TBP}\plusendmoddef
  procedure :: execute => cmd_weight_execute
\nwendcode{}\nwbegincode{2643}\moddef{Commands: procedures}\plusendmoddef
  subroutine cmd_weight_execute (cmd, global)
    class(cmd_weight_t), intent(inout) :: cmd
    type(rt_data_t), intent(inout), target :: global
    global%pn%weight_expr => cmd%pn_expr
  end subroutine cmd_weight_execute

\eatline
\nwendcode{}\nwbegindocs{2644}\nwdocspar
\subsubsection{Selection}
Define a selection expression.  This is to be applied upon simulation or
event-file rescanning, event by event.  We store the parse tree for the
right-hand side instead of compiling it.  Compilation is deferred to the
environment where the expression is used.
\nwenddocs{}\nwbegincode{2645}\moddef{Commands: types}\plusendmoddef
  type, extends (command_t) :: cmd_selection_t
     private
     type(parse_node_t), pointer :: pn_expr => null ()
   contains
   \LA{}Commands: cmd selection: TBP\RA{}
  end type cmd_selection_t

\eatline
\nwendcode{}\nwbegindocs{2646}Output. Do not print the parse tree, since this may get cluttered.
Just a message that scale, renormalization and factorization have been
defined, respectively.
\nwenddocs{}\nwbegincode{2647}\moddef{Commands: cmd selection: TBP}\endmoddef
  procedure :: write => cmd_selection_write
\nwendcode{}\nwbegincode{2648}\moddef{Commands: procedures}\plusendmoddef
  subroutine cmd_selection_write (cmd, unit, indent)
    class(cmd_selection_t), intent(in) :: cmd
    integer, intent(in), optional :: unit, indent
    integer :: u
    u = given_output_unit (unit);  if (u < 0)  return
    call write_indent (u, indent)
    write (u, "(1x,A)")  "selection expression: [defined]"
  end subroutine cmd_selection_write
  
\eatline
\nwendcode{}\nwbegindocs{2649}Compile.  Simply store the parse (sub)tree.
\nwenddocs{}\nwbegincode{2650}\moddef{Commands: cmd selection: TBP}\plusendmoddef
  procedure :: compile => cmd_selection_compile
\nwendcode{}\nwbegincode{2651}\moddef{Commands: procedures}\plusendmoddef
  subroutine cmd_selection_compile (cmd, global)
    class(cmd_selection_t), intent(inout) :: cmd
    type(rt_data_t), intent(inout), target :: global
    cmd%pn_expr => parse_node_get_sub_ptr (cmd%pn, 3)
  end subroutine cmd_selection_compile

\eatline
\nwendcode{}\nwbegindocs{2652}Instead of evaluating the expression, link the parse tree to the
global data set, such that it is compiled and executed in the
appropriate process context.
\nwenddocs{}\nwbegincode{2653}\moddef{Commands: cmd selection: TBP}\plusendmoddef
  procedure :: execute => cmd_selection_execute
\nwendcode{}\nwbegincode{2654}\moddef{Commands: procedures}\plusendmoddef
  subroutine cmd_selection_execute (cmd, global)
    class(cmd_selection_t), intent(inout) :: cmd
    type(rt_data_t), intent(inout), target :: global
    global%pn%selection_lexpr => cmd%pn_expr
  end subroutine cmd_selection_execute

\eatline
\nwendcode{}\nwbegindocs{2655}\nwdocspar
\subsubsection{Reweight}
Define a reweight expression.  This is to be applied upon simulation or
event-file rescanning, event by event.  We store the parse tree for the
right-hand side instead of compiling it.  Compilation is deferred to the
environment where the expression is used.
\nwenddocs{}\nwbegincode{2656}\moddef{Commands: types}\plusendmoddef
  type, extends (command_t) :: cmd_reweight_t
     private
     type(parse_node_t), pointer :: pn_expr => null ()
   contains
   \LA{}Commands: cmd reweight: TBP\RA{}  
  end type cmd_reweight_t

\eatline
\nwendcode{}\nwbegindocs{2657}Output. Do not print the parse tree, since this may get cluttered.
Just a message that scale, renormalization and factorization have been
defined, respectively.
\nwenddocs{}\nwbegincode{2658}\moddef{Commands: cmd reweight: TBP}\endmoddef
  procedure :: write => cmd_reweight_write
\nwendcode{}\nwbegincode{2659}\moddef{Commands: procedures}\plusendmoddef
  subroutine cmd_reweight_write (cmd, unit, indent)
    class(cmd_reweight_t), intent(in) :: cmd
    integer, intent(in), optional :: unit, indent
    integer :: u
    u = given_output_unit (unit);  if (u < 0)  return
    call write_indent (u, indent)
    write (u, "(1x,A)")  "reweight expression: [defined]"
  end subroutine cmd_reweight_write
  
\eatline
\nwendcode{}\nwbegindocs{2660}Compile.  Simply store the parse (sub)tree.
\nwenddocs{}\nwbegincode{2661}\moddef{Commands: cmd reweight: TBP}\plusendmoddef
  procedure :: compile => cmd_reweight_compile
\nwendcode{}\nwbegincode{2662}\moddef{Commands: procedures}\plusendmoddef
  subroutine cmd_reweight_compile (cmd, global)
    class(cmd_reweight_t), intent(inout) :: cmd
    type(rt_data_t), intent(inout), target :: global
    cmd%pn_expr => parse_node_get_sub_ptr (cmd%pn, 3)
  end subroutine cmd_reweight_compile

\eatline
\nwendcode{}\nwbegindocs{2663}Instead of evaluating the expression, link the parse tree to the
global data set, such that it is compiled and executed in the
appropriate process context.
\nwenddocs{}\nwbegincode{2664}\moddef{Commands: cmd reweight: TBP}\plusendmoddef
  procedure :: execute => cmd_reweight_execute
\nwendcode{}\nwbegincode{2665}\moddef{Commands: procedures}\plusendmoddef
  subroutine cmd_reweight_execute (cmd, global)
    class(cmd_reweight_t), intent(inout) :: cmd
    type(rt_data_t), intent(inout), target :: global
    global%pn%reweight_expr => cmd%pn_expr
  end subroutine cmd_reweight_execute

\eatline
\nwendcode{}\nwbegindocs{2666}\nwdocspar
\subsubsection{Alternative Simulation Setups}
Together with simulation, we can re-evaluate event weights in the context of
alternative setups.  The {\tt{}cmd{\char95}alt{\char95}setup{\char95}t} object is designed to hold these
setups, which are brace-enclosed command lists.  Compilation is deferred to
the simulation environment where the setup expression is used.
\nwenddocs{}\nwbegincode{2667}\moddef{Commands: types}\plusendmoddef
  type, extends (command_t) :: cmd_alt_setup_t
     private
     type(parse_node_p), dimension(:), allocatable :: setup
   contains
   \LA{}Commands: cmd alt setup: TBP\RA{}
  end type cmd_alt_setup_t

\eatline
\nwendcode{}\nwbegindocs{2668}Output.  Print just a message that the alternative setup list has been
defined.
\nwenddocs{}\nwbegincode{2669}\moddef{Commands: cmd alt setup: TBP}\endmoddef
  procedure :: write => cmd_alt_setup_write
\nwendcode{}\nwbegincode{2670}\moddef{Commands: procedures}\plusendmoddef
  subroutine cmd_alt_setup_write (cmd, unit, indent)
    class(cmd_alt_setup_t), intent(in) :: cmd
    integer, intent(in), optional :: unit, indent
    integer :: u
    u = given_output_unit (unit);  if (u < 0)  return
    call write_indent (u, indent)
    write (u, "(1x,A,I0,A)")  "alt_setup: ", size (cmd%setup), " entries"
  end subroutine cmd_alt_setup_write
  
\eatline
\nwendcode{}\nwbegindocs{2671}Compile.  Store the parse sub-trees in an array.
\nwenddocs{}\nwbegincode{2672}\moddef{Commands: cmd alt setup: TBP}\plusendmoddef
  procedure :: compile => cmd_alt_setup_compile
\nwendcode{}\nwbegincode{2673}\moddef{Commands: procedures}\plusendmoddef
  subroutine cmd_alt_setup_compile (cmd, global)
    class(cmd_alt_setup_t), intent(inout) :: cmd
    type(rt_data_t), intent(inout), target :: global
    type(parse_node_t), pointer :: pn_list, pn_setup
    integer :: i
    pn_list => parse_node_get_sub_ptr (cmd%pn, 3)
    if (associated (pn_list)) then
       allocate (cmd%setup (parse_node_get_n_sub (pn_list)))
       i = 1
       pn_setup => parse_node_get_sub_ptr (pn_list)
       do while (associated (pn_setup))
          cmd%setup(i)%ptr => pn_setup
          i = i + 1
          pn_setup => parse_node_get_next_ptr (pn_setup)
       end do
    else
       allocate (cmd%setup (0))
    end if
  end subroutine cmd_alt_setup_compile

\eatline
\nwendcode{}\nwbegindocs{2674}Execute.  Transfer the array of command lists to the global environment.
\nwenddocs{}\nwbegincode{2675}\moddef{Commands: cmd alt setup: TBP}\plusendmoddef
  procedure :: execute => cmd_alt_setup_execute
\nwendcode{}\nwbegincode{2676}\moddef{Commands: procedures}\plusendmoddef
  subroutine cmd_alt_setup_execute (cmd, global)
    class(cmd_alt_setup_t), intent(inout) :: cmd
    type(rt_data_t), intent(inout), target :: global
    if (allocated (global%pn%alt_setup))  deallocate (global%pn%alt_setup)
    allocate (global%pn%alt_setup (size (cmd%setup)), source = cmd%setup)
  end subroutine cmd_alt_setup_execute

\eatline
\nwendcode{}\nwbegindocs{2677}\nwdocspar
\subsubsection{Integration}
Integrate several processes, consecutively with identical parameters.
\nwenddocs{}\nwbegincode{2678}\moddef{Commands: types}\plusendmoddef
  type, extends (command_t) :: cmd_integrate_t
     private
     integer :: n_proc = 0
     type(string_t), dimension(:), allocatable :: process_id
   contains
   \LA{}Commands: cmd integrate: TBP\RA{}
  end type cmd_integrate_t
     
\eatline
\nwendcode{}\nwbegindocs{2679}Output: we know the process IDs.
\nwenddocs{}\nwbegincode{2680}\moddef{Commands: cmd integrate: TBP}\endmoddef
  procedure :: write => cmd_integrate_write
\nwendcode{}\nwbegincode{2681}\moddef{Commands: procedures}\plusendmoddef
  subroutine cmd_integrate_write (cmd, unit, indent)
    class(cmd_integrate_t), intent(in) :: cmd
    integer, intent(in), optional :: unit, indent
    integer :: u, i
    u = given_output_unit (unit);  if (u < 0)  return
    call write_indent (u, indent)
    write (u, "(1x,A)", advance="no")  "integrate ("
    do i = 1, cmd%n_proc
       if (i > 1)  write (u, "(A,1x)", advance="no")  ","
       write (u, "(A)", advance="no")  char (cmd%process_id(i))
    end do
    write (u, "(A)")  ")"
  end subroutine cmd_integrate_write

\eatline
\nwendcode{}\nwbegindocs{2682}Compile.
\nwenddocs{}\nwbegincode{2683}\moddef{Commands: cmd integrate: TBP}\plusendmoddef
  procedure :: compile => cmd_integrate_compile
\nwendcode{}\nwbegincode{2684}\moddef{Commands: procedures}\plusendmoddef
  subroutine cmd_integrate_compile (cmd, global)
    class(cmd_integrate_t), intent(inout) :: cmd
    type(rt_data_t), intent(inout), target :: global
    type(parse_node_t), pointer :: pn_proclist, pn_proc
    integer :: i
    pn_proclist => parse_node_get_sub_ptr (cmd%pn, 2)
    cmd%pn_opt => parse_node_get_next_ptr (pn_proclist)
    call cmd%compile_options (global)
    cmd%n_proc = parse_node_get_n_sub (pn_proclist)
    allocate (cmd%process_id (cmd%n_proc))
    pn_proc => parse_node_get_sub_ptr (pn_proclist)
    do i = 1, cmd%n_proc
       cmd%process_id(i) = parse_node_get_string (pn_proc)
       call global%process_stack%init_result_vars (cmd%process_id(i))
       pn_proc => parse_node_get_next_ptr (pn_proc)
    end do
  end subroutine cmd_integrate_compile

\eatline
\nwendcode{}\nwbegindocs{2685}Command execution.  Integrate the process(es) with the predefined number
of passes, iterations and calls.  For structure functions, cuts,
weight and scale, use local definitions if present; by default, the local
definitions are initialized with the global ones.

The {\tt{}integrate} procedure should take its input from the currently
active local environment, but produce a process record in the stack of
the global environment.

Since the process acquires a snapshot of the variable list, so if the global
list (or the local one) is deleted, this does no harm.  This implies that
later changes of the variable list do not affect the stored process.
\nwenddocs{}\nwbegincode{2686}\moddef{Commands: cmd integrate: TBP}\plusendmoddef
  procedure :: execute => cmd_integrate_execute
\nwendcode{}\nwbegincode{2687}\moddef{Commands: procedures}\plusendmoddef
  subroutine cmd_integrate_execute (cmd, global)
    class(cmd_integrate_t), intent(inout) :: cmd
    type(rt_data_t), intent(inout), target :: global
    integer :: i
    do i = 1, cmd%n_proc
       call integrate_process (cmd%process_id(i), cmd%local, global)
       call global%process_stack%fill_result_vars (cmd%process_id(i))
       if (signal_is_pending ())  return
    end do
  end subroutine cmd_integrate_execute

\eatline
\nwendcode{}\nwbegindocs{2688}\nwdocspar
\subsubsection{Observables}
Declare an observable.  After the declaration, it can be used to
record data, and at the end one can retrieve average and error.
\nwenddocs{}\nwbegincode{2689}\moddef{Commands: types}\plusendmoddef
  type, extends (command_t) :: cmd_observable_t
     private
     type(string_t) :: id
   contains
   \LA{}Commands: cmd observable: TBP\RA{}
  end type cmd_observable_t
     
\eatline
\nwendcode{}\nwbegindocs{2690}Output.  We know the ID.
\nwenddocs{}\nwbegincode{2691}\moddef{Commands: cmd observable: TBP}\endmoddef
  procedure :: write => cmd_observable_write
\nwendcode{}\nwbegincode{2692}\moddef{Commands: procedures}\plusendmoddef
  subroutine cmd_observable_write (cmd, unit, indent)
    class(cmd_observable_t), intent(in) :: cmd
    integer, intent(in), optional :: unit, indent
    integer :: u
    u = given_output_unit (unit);  if (u < 0)  return
    call write_indent (u, indent)
    write (u, "(1x,A,A)")  "observable: ", char (cmd%id)
  end subroutine cmd_observable_write
  
\eatline
\nwendcode{}\nwbegindocs{2693}Compile.  Just record the observable ID.
\nwenddocs{}\nwbegincode{2694}\moddef{Commands: cmd observable: TBP}\plusendmoddef
  procedure :: compile => cmd_observable_compile
\nwendcode{}\nwbegincode{2695}\moddef{Commands: procedures}\plusendmoddef
  subroutine cmd_observable_compile (cmd, global)
    class(cmd_observable_t), intent(inout) :: cmd
    type(rt_data_t), intent(inout), target :: global
    type(parse_node_t), pointer :: pn_tag
    pn_tag => parse_node_get_sub_ptr (cmd%pn, 2)
    if (associated (pn_tag)) then
       cmd%pn_opt => parse_node_get_next_ptr (pn_tag)
    end if       
    call cmd%compile_options (global)
    select case (char (parse_node_get_rule_key (pn_tag)))
    case ("analysis_id")
       cmd%id = parse_node_get_string (pn_tag)
    case default
       call msg_bug ("observable: name expression not implemented (yet)")
    end select
  end subroutine cmd_observable_compile

\eatline
\nwendcode{}\nwbegindocs{2696}Command execution.  This declares the observable and allocates it in
the analysis store.
\nwenddocs{}\nwbegincode{2697}\moddef{Commands: cmd observable: TBP}\plusendmoddef
  procedure :: execute => cmd_observable_execute
\nwendcode{}\nwbegincode{2698}\moddef{Commands: procedures}\plusendmoddef
  subroutine cmd_observable_execute (cmd, global)
    class(cmd_observable_t), intent(inout) :: cmd
    type(rt_data_t), intent(inout), target :: global
    type(var_list_t), pointer :: var_list
    type(graph_options_t) :: graph_options
    type(string_t) :: label, unit
    var_list => cmd%local%get_var_list_ptr ()
    label = var_list%get_sval (var_str ("$obs_label"))
    unit = var_list%get_sval (var_str ("$obs_unit"))
    call graph_options_init (graph_options)
    call set_graph_options (graph_options, var_list)
    call analysis_init_observable (cmd%id, label, unit, graph_options)
  end subroutine cmd_observable_execute

\eatline
\nwendcode{}\nwbegindocs{2699}\nwdocspar
\subsubsection{Histograms}
Declare a histogram.  At minimum, we have to set lower and upper bound
and bin width.
\nwenddocs{}\nwbegincode{2700}\moddef{Commands: types}\plusendmoddef
  type, extends (command_t) :: cmd_histogram_t
     private
     type(string_t) :: id
     type(parse_node_t), pointer :: pn_lower_bound => null ()
     type(parse_node_t), pointer :: pn_upper_bound => null ()
     type(parse_node_t), pointer :: pn_bin_width => null ()
   contains
   \LA{}Commands: cmd histogram: TBP\RA{}
  end type cmd_histogram_t
     
\eatline
\nwendcode{}\nwbegindocs{2701}Output.  Just print the ID.
\nwenddocs{}\nwbegincode{2702}\moddef{Commands: cmd histogram: TBP}\endmoddef
  procedure :: write => cmd_histogram_write
\nwendcode{}\nwbegincode{2703}\moddef{Commands: procedures}\plusendmoddef
  subroutine cmd_histogram_write (cmd, unit, indent)
    class(cmd_histogram_t), intent(in) :: cmd
    integer, intent(in), optional :: unit, indent
    integer :: u
    u = given_output_unit (unit);  if (u < 0)  return
    call write_indent (u, indent)
    write (u, "(1x,A,A)")  "histogram: ", char (cmd%id)
  end subroutine cmd_histogram_write
  
\eatline
\nwendcode{}\nwbegindocs{2704}Compile.  Record the histogram ID and initialize lower, upper bound
and bin width.
\nwenddocs{}\nwbegincode{2705}\moddef{Commands: cmd histogram: TBP}\plusendmoddef
  procedure :: compile => cmd_histogram_compile
\nwendcode{}\nwbegincode{2706}\moddef{Commands: procedures}\plusendmoddef
  subroutine cmd_histogram_compile (cmd, global)
    class(cmd_histogram_t), intent(inout) :: cmd
    type(rt_data_t), intent(inout), target :: global
    type(parse_node_t), pointer :: pn_tag, pn_args, pn_arg1, pn_arg2, pn_arg3
    character(*), parameter :: e_illegal_use = &
       "illegal usage of 'histogram': insufficient number of arguments"
    pn_tag => parse_node_get_sub_ptr (cmd%pn, 2)
    pn_args => parse_node_get_next_ptr (pn_tag)
    if (associated (pn_args)) then
       pn_arg1 => parse_node_get_sub_ptr (pn_args)
       if (.not. associated (pn_arg1)) call msg_fatal (e_illegal_use)
       pn_arg2 => parse_node_get_next_ptr (pn_arg1)
       if (.not. associated (pn_arg2)) call msg_fatal (e_illegal_use)
       pn_arg3 => parse_node_get_next_ptr (pn_arg2)
       cmd%pn_opt => parse_node_get_next_ptr (pn_args)
    end if       
    call cmd%compile_options (global)
    select case (char (parse_node_get_rule_key (pn_tag)))
    case ("analysis_id")
       cmd%id = parse_node_get_string (pn_tag)
    case default
       call msg_bug ("histogram: name expression not implemented (yet)")
    end select
    cmd%pn_lower_bound => pn_arg1
    cmd%pn_upper_bound => pn_arg2
    cmd%pn_bin_width => pn_arg3
  end subroutine cmd_histogram_compile

\eatline
\nwendcode{}\nwbegindocs{2707}Command execution.  This declares the histogram and allocates it in
the analysis store.
\nwenddocs{}\nwbegincode{2708}\moddef{Commands: cmd histogram: TBP}\plusendmoddef
  procedure :: execute => cmd_histogram_execute
\nwendcode{}\nwbegincode{2709}\moddef{Commands: procedures}\plusendmoddef
  subroutine cmd_histogram_execute (cmd, global)
    class(cmd_histogram_t), intent(inout) :: cmd
    type(rt_data_t), intent(inout), target :: global
    type(var_list_t), pointer :: var_list
    real(default) :: lower_bound, upper_bound, bin_width
    integer :: bin_number
    logical :: bin_width_is_used, normalize_bins
    type(string_t) :: obs_label, obs_unit
    type(graph_options_t) :: graph_options
    type(drawing_options_t) :: drawing_options

    var_list => cmd%local%get_var_list_ptr ()
    lower_bound = eval_real (cmd%pn_lower_bound, var_list)
    upper_bound = eval_real (cmd%pn_upper_bound, var_list)
    if (associated (cmd%pn_bin_width)) then
       bin_width = eval_real (cmd%pn_bin_width, var_list)
       bin_width_is_used = .true.
    else if (var_list%is_known (var_str ("n_bins"))) then
       bin_number = &
            var_list%get_ival (var_str ("n_bins"))
       bin_width_is_used = .false.
    else
       call msg_error ("Cmd '" // char (cmd%id) // &
            "': neither bin width nor number is defined")
    end if
    normalize_bins = &
         var_list%get_lval (var_str ("?normalize_bins"))
    obs_label = &
         var_list%get_sval (var_str ("$obs_label"))
    obs_unit = &
         var_list%get_sval (var_str ("$obs_unit"))

    call graph_options_init (graph_options)
    call set_graph_options (graph_options, var_list)
    call drawing_options_init_histogram (drawing_options)
    call set_drawing_options (drawing_options, var_list)

    if (bin_width_is_used) then
       call analysis_init_histogram &
            (cmd%id, lower_bound, upper_bound, bin_width, &
             normalize_bins, &
             obs_label, obs_unit, &
             graph_options, drawing_options)
    else
       call analysis_init_histogram &
            (cmd%id, lower_bound, upper_bound, bin_number, &
             normalize_bins, &
             obs_label, obs_unit, &
             graph_options, drawing_options)
    end if
  end subroutine cmd_histogram_execute

\eatline
\nwendcode{}\nwbegindocs{2710}Set the graph options from a variable list.
\nwenddocs{}\nwbegincode{2711}\moddef{Commands: procedures}\plusendmoddef
  subroutine set_graph_options (gro, var_list)
    type(graph_options_t), intent(inout) :: gro
    type(var_list_t), intent(in) :: var_list
    call graph_options_set (gro, title = &
         var_list%get_sval (var_str ("$title")))
    call graph_options_set (gro, description = &
         var_list%get_sval (var_str ("$description")))
    call graph_options_set (gro, x_label = &
         var_list%get_sval (var_str ("$x_label")))
    call graph_options_set (gro, y_label = &
         var_list%get_sval (var_str ("$y_label")))
    call graph_options_set (gro, width_mm = &
         var_list%get_ival (var_str ("graph_width_mm")))
    call graph_options_set (gro, height_mm = &
         var_list%get_ival (var_str ("graph_height_mm")))
    call graph_options_set (gro, x_log = &
         var_list%get_lval (var_str ("?x_log")))
    call graph_options_set (gro, y_log = &
         var_list%get_lval (var_str ("?y_log")))
    if (var_list%is_known (var_str ("x_min"))) &
         call graph_options_set (gro, x_min = &
         var_list%get_rval (var_str ("x_min")))
    if (var_list%is_known (var_str ("x_max"))) &
         call graph_options_set (gro, x_max = &
         var_list%get_rval (var_str ("x_max")))
    if (var_list%is_known (var_str ("y_min"))) &
         call graph_options_set (gro, y_min = &
         var_list%get_rval (var_str ("y_min")))
    if (var_list%is_known (var_str ("y_max"))) &
         call graph_options_set (gro, y_max = &
         var_list%get_rval (var_str ("y_max")))
    call graph_options_set (gro, gmlcode_bg = &
         var_list%get_sval (var_str ("$gmlcode_bg")))
    call graph_options_set (gro, gmlcode_fg = &
         var_list%get_sval (var_str ("$gmlcode_fg")))
  end subroutine set_graph_options

\eatline
\nwendcode{}\nwbegindocs{2712}Set the drawing options from a variable list.
\nwenddocs{}\nwbegincode{2713}\moddef{Commands: procedures}\plusendmoddef
  subroutine set_drawing_options (dro, var_list)
    type(drawing_options_t), intent(inout) :: dro
    type(var_list_t), intent(in) :: var_list
    if (var_list%is_known (var_str ("?draw_histogram"))) then
       if (var_list%get_lval (var_str ("?draw_histogram"))) then
          call drawing_options_set (dro, with_hbars = .true.)
       else
          call drawing_options_set (dro, with_hbars = .false., &
               with_base = .false., fill = .false., piecewise = .false.)
       end if
    end if
    if (var_list%is_known (var_str ("?draw_base"))) then
       if (var_list%get_lval (var_str ("?draw_base"))) then
          call drawing_options_set (dro, with_base = .true.)
       else
          call drawing_options_set (dro, with_base = .false., fill = .false.)
       end if
    end if
    if (var_list%is_known (var_str ("?draw_piecewise"))) then
       if (var_list%get_lval (var_str ("?draw_piecewise"))) then
          call drawing_options_set (dro, piecewise = .true.)
       else
          call drawing_options_set (dro, piecewise = .false.)
       end if
    end if
    if (var_list%is_known (var_str ("?fill_curve"))) then
       if (var_list%get_lval (var_str ("?fill_curve"))) then
          call drawing_options_set (dro, fill = .true., with_base = .true.)
       else
          call drawing_options_set (dro, fill = .false.)
       end if
    end if
    if (var_list%is_known (var_str ("?draw_curve"))) then
       if (var_list%get_lval (var_str ("?draw_curve"))) then
          call drawing_options_set (dro, draw = .true.)
       else
          call drawing_options_set (dro, draw = .false.)
       end if
    end if
    if (var_list%is_known (var_str ("?draw_errors"))) then
       if (var_list%get_lval (var_str ("?draw_errors"))) then
          call drawing_options_set (dro, err = .true.)
       else
          call drawing_options_set (dro, err = .false.)
       end if
    end if
    if (var_list%is_known (var_str ("?draw_symbols"))) then
       if (var_list%get_lval (var_str ("?draw_symbols"))) then
          call drawing_options_set (dro, symbols = .true.)
       else
          call drawing_options_set (dro, symbols = .false.)
       end if
    end if
    if (var_list%is_known (var_str ("$fill_options"))) then
       call drawing_options_set (dro, fill_options = &
            var_list%get_sval (var_str ("$fill_options")))
    end if
    if (var_list%is_known (var_str ("$draw_options"))) then
       call drawing_options_set (dro, draw_options = &
            var_list%get_sval (var_str ("$draw_options")))
    end if
    if (var_list%is_known (var_str ("$err_options"))) then
       call drawing_options_set (dro, err_options = &
            var_list%get_sval (var_str ("$err_options")))
    end if
    if (var_list%is_known (var_str ("$symbol"))) then
       call drawing_options_set (dro, symbol = &
            var_list%get_sval (var_str ("$symbol")))
    end if
    if (var_list%is_known (var_str ("$gmlcode_bg"))) then
       call drawing_options_set (dro, gmlcode_bg = &
            var_list%get_sval (var_str ("$gmlcode_bg")))
    end if
    if (var_list%is_known (var_str ("$gmlcode_fg"))) then
       call drawing_options_set (dro, gmlcode_fg = &
            var_list%get_sval (var_str ("$gmlcode_fg")))
    end if
  end subroutine set_drawing_options

\eatline
\nwendcode{}\nwbegindocs{2714}\nwdocspar
\subsubsection{Plots}
Declare a plot.  No mandatory arguments, just options.
\nwenddocs{}\nwbegincode{2715}\moddef{Commands: types}\plusendmoddef
  type, extends (command_t) :: cmd_plot_t
     private
     type(string_t) :: id
   contains
   \LA{}Commands: cmd plot: TBP\RA{}
  end type cmd_plot_t
     
\eatline
\nwendcode{}\nwbegindocs{2716}Output.  Just print the ID.
\nwenddocs{}\nwbegincode{2717}\moddef{Commands: cmd plot: TBP}\endmoddef
  procedure :: write => cmd_plot_write
\nwendcode{}\nwbegincode{2718}\moddef{Commands: procedures}\plusendmoddef
  subroutine cmd_plot_write (cmd, unit, indent)
    class(cmd_plot_t), intent(in) :: cmd
    integer, intent(in), optional :: unit, indent
    integer :: u
    u = given_output_unit (unit);  if (u < 0)  return
    call write_indent (u, indent)
    write (u, "(1x,A,A)")  "plot: ", char (cmd%id)
  end subroutine cmd_plot_write
  
\eatline
\nwendcode{}\nwbegindocs{2719}Compile.  Record the plot ID and initialize lower, upper bound
and bin width.
\nwenddocs{}\nwbegincode{2720}\moddef{Commands: cmd plot: TBP}\plusendmoddef
  procedure :: compile => cmd_plot_compile
\nwendcode{}\nwbegincode{2721}\moddef{Commands: procedures}\plusendmoddef
  subroutine cmd_plot_compile (cmd, global)
    class(cmd_plot_t), intent(inout) :: cmd
    type(rt_data_t), intent(inout), target :: global
    type(parse_node_t), pointer :: pn_tag
    pn_tag => parse_node_get_sub_ptr (cmd%pn, 2)
    cmd%pn_opt => parse_node_get_next_ptr (pn_tag)
    call cmd%init (pn_tag, global)
  end subroutine cmd_plot_compile

\eatline
\nwendcode{}\nwbegindocs{2722}This init routine is separated because it is reused below for graph
initialization. 
\nwenddocs{}\nwbegincode{2723}\moddef{Commands: cmd plot: TBP}\plusendmoddef
  procedure :: init => cmd_plot_init
\nwendcode{}\nwbegincode{2724}\moddef{Commands: procedures}\plusendmoddef
  subroutine cmd_plot_init (plot, pn_tag, global)
    class(cmd_plot_t), intent(inout) :: plot
    type(parse_node_t), intent(in), pointer :: pn_tag
    type(rt_data_t), intent(inout), target :: global
    call plot%compile_options (global)
    select case (char (parse_node_get_rule_key (pn_tag)))
    case ("analysis_id")
       plot%id = parse_node_get_string (pn_tag)
    case default
       call msg_bug ("plot: name expression not implemented (yet)")
    end select
  end subroutine cmd_plot_init

\eatline
\nwendcode{}\nwbegindocs{2725}Command execution.  This declares the plot and allocates it in
the analysis store.
\nwenddocs{}\nwbegincode{2726}\moddef{Commands: cmd plot: TBP}\plusendmoddef
  procedure :: execute => cmd_plot_execute
\nwendcode{}\nwbegincode{2727}\moddef{Commands: procedures}\plusendmoddef
  subroutine cmd_plot_execute (cmd, global)
    class(cmd_plot_t), intent(inout) :: cmd
    type(rt_data_t), intent(inout), target :: global
    type(var_list_t), pointer :: var_list
    type(graph_options_t) :: graph_options
    type(drawing_options_t) :: drawing_options

    var_list => cmd%local%get_var_list_ptr ()
    call graph_options_init (graph_options)
    call set_graph_options (graph_options, var_list)
    call drawing_options_init_plot (drawing_options)
    call set_drawing_options (drawing_options, var_list)

    call analysis_init_plot (cmd%id, graph_options, drawing_options)
  end subroutine cmd_plot_execute

\eatline
\nwendcode{}\nwbegindocs{2728}\nwdocspar
\subsubsection{Graphs}
Declare a graph.  The graph is defined in terms of its contents.  Both the
graph and its contents may carry options.

The graph object contains its own ID as well as the IDs of its elements.  For
the elements, we reuse the {\tt{}cmd{\char95}plot{\char95}t} defined above.
\nwenddocs{}\nwbegincode{2729}\moddef{Commands: types}\plusendmoddef
  type, extends (command_t) :: cmd_graph_t
     private
     type(string_t) :: id
     integer :: n_elements = 0
     type(cmd_plot_t), dimension(:), allocatable :: el
     type(string_t), dimension(:), allocatable :: element_id
   contains
   \LA{}Commands: cmd graph: TBP\RA{}
  end type cmd_graph_t
     
\eatline
\nwendcode{}\nwbegindocs{2730}Output.  Just print the ID.
\nwenddocs{}\nwbegincode{2731}\moddef{Commands: cmd graph: TBP}\endmoddef
  procedure :: write => cmd_graph_write
\nwendcode{}\nwbegincode{2732}\moddef{Commands: procedures}\plusendmoddef
  subroutine cmd_graph_write (cmd, unit, indent)
    class(cmd_graph_t), intent(in) :: cmd
    integer, intent(in), optional :: unit, indent
    integer :: u
    u = given_output_unit (unit);  if (u < 0)  return
    call write_indent (u, indent)
    write (u, "(1x,A,A,A,I0,A)")  "graph: ", char (cmd%id), &
         " (", cmd%n_elements, " entries)"
  end subroutine cmd_graph_write
  
\eatline
\nwendcode{}\nwbegindocs{2733}Compile.  Record the graph ID and initialize lower, upper bound
and bin width.  For compiling the graph element syntax, we use part of the
{\tt{}cmd{\char95}plot{\char95}t} compiler.

Note: currently, we do not respect options, therefore just IDs on the RHS.
\nwenddocs{}\nwbegincode{2734}\moddef{Commands: cmd graph: TBP}\plusendmoddef
  procedure :: compile => cmd_graph_compile
\nwendcode{}\nwbegincode{2735}\moddef{Commands: procedures}\plusendmoddef
  subroutine cmd_graph_compile (cmd, global)
    class(cmd_graph_t), intent(inout) :: cmd
    type(rt_data_t), intent(inout), target :: global
    type(parse_node_t), pointer :: pn_term, pn_tag, pn_def, pn_app
    integer :: i

    pn_term => parse_node_get_sub_ptr (cmd%pn, 2)
    pn_tag => parse_node_get_sub_ptr (pn_term)
    cmd%pn_opt => parse_node_get_next_ptr (pn_tag)
    call cmd%compile_options (global)
    select case (char (parse_node_get_rule_key (pn_tag)))
    case ("analysis_id")
       cmd%id = parse_node_get_string (pn_tag)
    case default
       call msg_bug ("graph: name expression not implemented (yet)")
    end select
    pn_def => parse_node_get_next_ptr (pn_term, 2)
    cmd%n_elements = parse_node_get_n_sub (pn_def)
    allocate (cmd%element_id (cmd%n_elements))
    allocate (cmd%el (cmd%n_elements))
    pn_term => parse_node_get_sub_ptr (pn_def)
    pn_tag => parse_node_get_sub_ptr (pn_term)
    cmd%el(1)%pn_opt => parse_node_get_next_ptr (pn_tag)
    call cmd%el(1)%init (pn_tag, global)
    cmd%element_id(1) = parse_node_get_string (pn_tag)
    pn_app => parse_node_get_next_ptr (pn_term)
    do i = 2, cmd%n_elements
       pn_term => parse_node_get_sub_ptr (pn_app, 2)
       pn_tag => parse_node_get_sub_ptr (pn_term)
       cmd%el(i)%pn_opt => parse_node_get_next_ptr (pn_tag)
       call cmd%el(i)%init (pn_tag, global)
       cmd%element_id(i) = parse_node_get_string (pn_tag)
       pn_app => parse_node_get_next_ptr (pn_app)
    end do

  end subroutine cmd_graph_compile

\eatline
\nwendcode{}\nwbegindocs{2736}Command execution.  This declares the graph, allocates it in
the analysis store, and copies the graph elements.

For the graph, we set graph and default drawing options.  For the elements, we
reset individual drawing options.

This accesses internals of the contained elements of type {\tt{}cmd{\char95}plot{\char95}t}, see
above.  We might disentangle such an interdependency when this code is
rewritten using proper type extension.
\nwenddocs{}\nwbegincode{2737}\moddef{Commands: cmd graph: TBP}\plusendmoddef
  procedure :: execute => cmd_graph_execute
\nwendcode{}\nwbegincode{2738}\moddef{Commands: procedures}\plusendmoddef
  subroutine cmd_graph_execute (cmd, global)
    class(cmd_graph_t), intent(inout) :: cmd
    type(rt_data_t), intent(inout), target :: global
    type(var_list_t), pointer :: var_list
    type(graph_options_t) :: graph_options
    type(drawing_options_t) :: drawing_options
    integer :: i, type

    var_list => cmd%local%get_var_list_ptr ()
    call graph_options_init (graph_options)
    call set_graph_options (graph_options, var_list)
    call analysis_init_graph (cmd%id, cmd%n_elements, graph_options)

    do i = 1, cmd%n_elements
       if (associated (cmd%el(i)%options)) then
          call cmd%el(i)%options%execute (cmd%el(i)%local)
       end if
       type = analysis_store_get_object_type (cmd%element_id(i))
       select case (type)
       case (AN_HISTOGRAM)
          call drawing_options_init_histogram (drawing_options)
       case (AN_PLOT)
          call drawing_options_init_plot (drawing_options)
       end select
       call set_drawing_options (drawing_options, var_list)
       if (associated (cmd%el(i)%options)) then
          call set_drawing_options (drawing_options, cmd%el(i)%local%var_list)
       end if
       call analysis_fill_graph (cmd%id, i, cmd%element_id(i), drawing_options)
    end do
  end subroutine cmd_graph_execute

\eatline
\nwendcode{}\nwbegindocs{2739}\nwdocspar
\subsubsection{Analysis}
Hold the analysis ID either as a string or as an expression:
\nwenddocs{}\nwbegincode{2740}\moddef{Commands: types}\plusendmoddef
  type :: analysis_id_t
    type(string_t) :: tag
    type(parse_node_t), pointer :: pn_sexpr => null ()
  end type analysis_id_t

\eatline
\nwendcode{}\nwbegindocs{2741}Define the analysis expression.  We store the parse tree for the
right-hand side instead of compiling it.  Compilation is deferred to
the process environment where the analysis expression is used.
\nwenddocs{}\nwbegincode{2742}\moddef{Commands: types}\plusendmoddef
  type, extends (command_t) :: cmd_analysis_t
     private
     type(parse_node_t), pointer :: pn_lexpr => null ()
   contains
   \LA{}Commands: cmd analysis: TBP\RA{}
  end type cmd_analysis_t

\eatline
\nwendcode{}\nwbegindocs{2743}Output.  Print just a message that analysis has been defined.
\nwenddocs{}\nwbegincode{2744}\moddef{Commands: cmd analysis: TBP}\endmoddef
  procedure :: write => cmd_analysis_write
\nwendcode{}\nwbegincode{2745}\moddef{Commands: procedures}\plusendmoddef
  subroutine cmd_analysis_write (cmd, unit, indent)
    class(cmd_analysis_t), intent(in) :: cmd
    integer, intent(in), optional :: unit, indent
    integer :: u
    u = given_output_unit (unit);  if (u < 0)  return
    call write_indent (u, indent)
    write (u, "(1x,A)")  "analysis: [defined]"
  end subroutine cmd_analysis_write
  
\eatline
\nwendcode{}\nwbegindocs{2746}Compile.  Simply store the parse (sub)tree.
\nwenddocs{}\nwbegincode{2747}\moddef{Commands: cmd analysis: TBP}\plusendmoddef
  procedure :: compile => cmd_analysis_compile
\nwendcode{}\nwbegincode{2748}\moddef{Commands: procedures}\plusendmoddef
  subroutine cmd_analysis_compile (cmd, global)
    class(cmd_analysis_t), intent(inout) :: cmd
    type(rt_data_t), intent(inout), target :: global
    cmd%pn_lexpr => parse_node_get_sub_ptr (cmd%pn, 3)
  end subroutine cmd_analysis_compile

\eatline
\nwendcode{}\nwbegindocs{2749}Instead of evaluating the cut expression, link the parse tree to the
global data set, such that it is compiled and executed in the
appropriate process context.
\nwenddocs{}\nwbegincode{2750}\moddef{Commands: cmd analysis: TBP}\plusendmoddef
  procedure :: execute => cmd_analysis_execute
\nwendcode{}\nwbegincode{2751}\moddef{Commands: procedures}\plusendmoddef
  subroutine cmd_analysis_execute (cmd, global)
    class(cmd_analysis_t), intent(inout) :: cmd
    type(rt_data_t), intent(inout), target :: global
    global%pn%analysis_lexpr => cmd%pn_lexpr
  end subroutine cmd_analysis_execute

\eatline
\nwendcode{}\nwbegindocs{2752}\nwdocspar
\subsubsection{Write histograms and plots}
The data type encapsulating the command:
\nwenddocs{}\nwbegincode{2753}\moddef{Commands: types}\plusendmoddef
  type, extends (command_t) :: cmd_write_analysis_t
     private
     type(analysis_id_t), dimension(:), allocatable :: id
     type(string_t), dimension(:), allocatable :: tag
   contains
   \LA{}Commands: cmd write analysis: TBP\RA{}
  end type cmd_write_analysis_t

\eatline
\eatline
\nwendcode{}\nwbegindocs{2754}Output.  Just the keyword.
\nwenddocs{}\nwbegincode{2755}\moddef{Commands: cmd write analysis: TBP}\endmoddef
  procedure :: write => cmd_write_analysis_write
\nwendcode{}\nwbegincode{2756}\moddef{Commands: procedures}\plusendmoddef
  subroutine cmd_write_analysis_write (cmd, unit, indent)
    class(cmd_write_analysis_t), intent(in) :: cmd
    integer, intent(in), optional :: unit, indent
    integer :: u
    u = given_output_unit (unit);  if (u < 0)  return
    call write_indent (u, indent)
    write (u, "(1x,A)")  "write_analysis"
  end subroutine cmd_write_analysis_write
  
\eatline
\nwendcode{}\nwbegindocs{2757}Compile.
\nwenddocs{}\nwbegincode{2758}\moddef{Commands: cmd write analysis: TBP}\plusendmoddef
  procedure :: compile => cmd_write_analysis_compile
\nwendcode{}\nwbegincode{2759}\moddef{Commands: procedures}\plusendmoddef
  subroutine cmd_write_analysis_compile (cmd, global)
    class(cmd_write_analysis_t), intent(inout) :: cmd
    type(rt_data_t), intent(inout), target :: global
    type(parse_node_t), pointer :: pn_clause, pn_args, pn_id
    integer :: n, i
    pn_clause => parse_node_get_sub_ptr (cmd%pn)
    pn_args => parse_node_get_sub_ptr (pn_clause, 2)
    cmd%pn_opt => parse_node_get_next_ptr (pn_clause)
    call cmd%compile_options (global)
    if (associated (pn_args)) then
       n = parse_node_get_n_sub (pn_args)
       allocate (cmd%id (n))
       do i = 1, n
           pn_id => parse_node_get_sub_ptr (pn_args, i)
           if (char (parse_node_get_rule_key (pn_id)) == "analysis_id") then
              cmd%id(i)%tag = parse_node_get_string (pn_id)              
           else
              cmd%id(i)%pn_sexpr => pn_id
           end if
       end do
    else
       allocate (cmd%id (0))
    end if
  end subroutine cmd_write_analysis_compile

\eatline
\nwendcode{}\nwbegindocs{2760}The output format for real data values:
\nwenddocs{}\nwbegincode{2761}\moddef{Commands: parameters}\plusendmoddef
  character(*), parameter, public :: &
       DEFAULT_ANALYSIS_FILENAME = "whizard_analysis.dat"
  character(len=1), dimension(2), parameter, public :: &
       FORBIDDEN_ENDINGS1 = [ "o", "a" ]
  character(len=2), dimension(5), parameter, public :: &       
       FORBIDDEN_ENDINGS2 = [ "mp", "ps", "vg", "lo", "la" ]
  character(len=3), dimension(14), parameter, public :: &
       FORBIDDEN_ENDINGS3 = [ "aux", "dvi", "evt", "evx", "f03", "f90", &
          "f95", "log", "ltp", "mpx", "pdf", "phs", "sin", "tex" ]
       
\eatline
\eatline
\eatline
\eatline
\nwendcode{}\nwbegindocs{2762}As this contains a lot of similar code to {\tt{}cmd{\char95}compile{\char95}analysis{\char95}execute}
we outsource the main code to a subroutine.
\nwenddocs{}\nwbegincode{2763}\moddef{Commands: cmd write analysis: TBP}\plusendmoddef
  procedure :: execute => cmd_write_analysis_execute
\nwendcode{}\nwbegincode{2764}\moddef{Commands: procedures}\plusendmoddef
  subroutine cmd_write_analysis_execute (cmd, global)
    class(cmd_write_analysis_t), intent(inout) :: cmd
    type(rt_data_t), intent(inout), target :: global
    type(var_list_t), pointer :: var_list
    var_list => cmd%local%get_var_list_ptr ()
    call write_analysis_wrap (var_list, global%out_files, &
         cmd%id, tag = cmd%tag)
  end subroutine cmd_write_analysis_execute

\eatline
\nwendcode{}\nwbegindocs{2765}If the {\tt{}data{\char95}file} optional argument is present, this is 
called from {\tt{}cmd{\char95}compile{\char95}analysis{\char95}execute}, which needs the file name for
further processing, and requires the default format.  For the moment, 
parameters and macros for custom data processing are disabled.
\nwenddocs{}\nwbegincode{2766}\moddef{Commands: procedures}\plusendmoddef
  subroutine write_analysis_wrap (var_list, out_files, id, tag, data_file)
    type(var_list_t), intent(inout), target :: var_list
    type(file_list_t), intent(inout), target :: out_files
    type(analysis_id_t), dimension(:), intent(in), target :: id 
    type(string_t), dimension(:), allocatable, intent(out) :: tag
    type(string_t), intent(out), optional :: data_file
    type(string_t) :: defaultfile, file
    integer :: i   
    logical :: keep_open !, custom, header, columns    
    type(string_t) :: extension !, comment_prefix, separator 
!!! JRR: WK please check (#542)           
!     integer :: type
!     type(ifile_t) :: ifile
    logical :: one_file !, has_writer
!     type(analysis_iterator_t) :: iterator
!     type(rt_data_t), target :: sandbox
!     type(command_list_t) :: writer    
    defaultfile = var_list%get_sval (var_str ("$out_file"))
    if (present (data_file)) then
       if (defaultfile == "" .or. defaultfile == ".") then
          defaultfile = DEFAULT_ANALYSIS_FILENAME
       else
          if (scan (".", defaultfile) > 0) then
             call split (defaultfile, extension, ".", back=.true.)
             if (any (lower_case (char(extension)) == FORBIDDEN_ENDINGS1) .or. &
                 any (lower_case (char(extension)) == FORBIDDEN_ENDINGS2) .or. &
                 any (lower_case (char(extension)) == FORBIDDEN_ENDINGS3)) & 
                 call msg_fatal ("The ending " // char(extension) // &
                 " is internal and not allowed as data file.")
             if (extension /= "") then
                if (defaultfile /= "") then
                   defaultfile = defaultfile // "." // extension
                else
                   defaultfile = "whizard_analysis." // extension
                end if
             else
                defaultfile = defaultfile // ".dat"
             endif
          else
             defaultfile = defaultfile // ".dat"
          end if
       end if
       data_file = defaultfile
    end if
    one_file = defaultfile /= ""
    if (one_file) then
       file = defaultfile
       keep_open = file_list_is_open (out_files, file, &
            action = "write")
       if (keep_open) then
          if (present (data_file)) then
             call msg_fatal ("Compiling analysis: File '" &
                  // char (data_file) &
                   // "' can't be used, it is already open.")
          else
             call msg_message ("Appending analysis data to file '" &
                  // char (file) // "'")
          end if
       else
          call file_list_open (out_files, file, &
               action = "write", status = "replace", position = "asis")
          call msg_message ("Writing analysis data to file '" &
               // char (file) // "'")          
       end if
    end if    

!!! JRR: WK please check. Custom data output. Ticket #542
!     if (present (data_file)) then
!        custom = .false.
!     else
!        custom = var_list%get_lval (&
!            var_str ("?out_custom"))
!     end if
!     comment_prefix = var_list%get_sval (&
!          var_str ("$out_comment"))
!     header = var_list%get_lval (&
!          var_str ("?out_header"))
!     write_yerr = var_list%get_lval (&
!          var_str ("?out_yerr"))
!     write_xerr = var_list%get_lval (&
!          var_str ("?out_xerr"))

    call get_analysis_tags (tag, id, var_list)       
    do i = 1, size (tag)
       call file_list_write_analysis &
            (out_files, file, tag(i))
    end do
    if (one_file .and. .not. keep_open) then
       call file_list_close (out_files, file)
    end if       

  contains
    
    subroutine get_analysis_tags (analysis_tag, id, var_list)
      type(string_t), dimension(:), intent(out), allocatable :: analysis_tag
      type(analysis_id_t), dimension(:), intent(in) :: id
      type(var_list_t), intent(in), target :: var_list
      if (size (id) /= 0) then
         allocate (analysis_tag (size (id)))
         do i = 1, size (id)
            if (associated (id(i)%pn_sexpr)) then
               analysis_tag(i) = eval_string (id(i)%pn_sexpr, var_list)
            else
               analysis_tag(i) = id(i)%tag
            end if
         end do
      else
         call analysis_store_get_ids (tag)                
      end if
    end subroutine get_analysis_tags
    
  end subroutine write_analysis_wrap
  
\eatline
\nwendcode{}\nwbegindocs{2767}\subsubsection{Compile analysis results}
This command writes files in a form suitable for GAMELAN and executes the
appropriate commands to compile them.  The first part is identical to
{\tt{}cmd{\char95}write{\char95}analysis}.
\nwenddocs{}\nwbegincode{2768}\moddef{Commands: types}\plusendmoddef
  type, extends (command_t) :: cmd_compile_analysis_t
     private
     type(analysis_id_t), dimension(:), allocatable :: id
     type(string_t), dimension(:), allocatable :: tag
   contains
   \LA{}Commands: cmd compile analysis: TBP\RA{}
  end type cmd_compile_analysis_t

\eatline
\nwendcode{}\nwbegindocs{2769}Output.  Just the keyword.
\nwenddocs{}\nwbegincode{2770}\moddef{Commands: cmd compile analysis: TBP}\endmoddef
  procedure :: write => cmd_compile_analysis_write
\nwendcode{}\nwbegincode{2771}\moddef{Commands: procedures}\plusendmoddef
  subroutine cmd_compile_analysis_write (cmd, unit, indent)
    class(cmd_compile_analysis_t), intent(in) :: cmd
    integer, intent(in), optional :: unit, indent
    integer :: u
    u = given_output_unit (unit);  if (u < 0)  return
    call write_indent (u, indent)
    write (u, "(1x,A)")  "compile_analysis"
  end subroutine cmd_compile_analysis_write
  
\eatline
\nwendcode{}\nwbegindocs{2772}Compile.
\nwenddocs{}\nwbegincode{2773}\moddef{Commands: cmd compile analysis: TBP}\plusendmoddef
  procedure :: compile => cmd_compile_analysis_compile
\nwendcode{}\nwbegincode{2774}\moddef{Commands: procedures}\plusendmoddef
  subroutine cmd_compile_analysis_compile (cmd, global)
    class(cmd_compile_analysis_t), intent(inout) :: cmd
    type(rt_data_t), intent(inout), target :: global
    type(parse_node_t), pointer :: pn_clause, pn_args, pn_id
    integer :: n, i
    pn_clause => parse_node_get_sub_ptr (cmd%pn)
    pn_args => parse_node_get_sub_ptr (pn_clause, 2)
    cmd%pn_opt => parse_node_get_next_ptr (pn_clause)
    call cmd%compile_options (global)
    if (associated (pn_args)) then
       n = parse_node_get_n_sub (pn_args)
       allocate (cmd%id (n))
       do i = 1, n
           pn_id => parse_node_get_sub_ptr (pn_args, i)
           if (char (parse_node_get_rule_key (pn_id)) == "analysis_id") then
              cmd%id(i)%tag = parse_node_get_string (pn_id)
           else
              cmd%id(i)%pn_sexpr => pn_id
           end if
       end do
    else
       allocate (cmd%id (0))       
    end if
  end subroutine cmd_compile_analysis_compile

\eatline
\nwendcode{}\nwbegindocs{2775}First write the analysis data to file, then write a GAMELAN driver and
produce MetaPost and \TeX\ output.
\nwenddocs{}\nwbegincode{2776}\moddef{Commands: cmd compile analysis: TBP}\plusendmoddef
  procedure :: execute => cmd_compile_analysis_execute
\nwendcode{}\nwbegincode{2777}\moddef{Commands: procedures}\plusendmoddef
  subroutine cmd_compile_analysis_execute (cmd, global)
    class(cmd_compile_analysis_t), intent(inout) :: cmd
    type(rt_data_t), intent(inout), target :: global
    type(var_list_t), pointer :: var_list
    type(string_t) :: file, basename, extension, driver_file
    integer :: u_driver
    logical :: has_gmlcode
    var_list => cmd%local%get_var_list_ptr ()
    call write_analysis_wrap (var_list, &
         global%out_files, cmd%id, tag = cmd%tag, &
            data_file = file)
    basename = file    
    if (scan (".", basename) > 0) then
      call split (basename, extension, ".", back=.true.)
    else
      extension = ""
    end if
    driver_file = basename // ".tex"
    u_driver = free_unit ()
    open (unit=u_driver, file=char(driver_file), &
          action="write", status="replace")
    if (allocated (cmd%tag)) then
       call analysis_write_driver (file, cmd%tag, unit=u_driver)
       has_gmlcode = analysis_has_plots (cmd%tag)
    else
       call analysis_write_driver (file, unit=u_driver)
       has_gmlcode = analysis_has_plots ()
    end if
    close (u_driver)
    call msg_message ("Compiling analysis results display in '" &
         // char (driver_file) // "'")
    call analysis_compile_tex (basename, has_gmlcode, global%os_data)
  end subroutine cmd_compile_analysis_execute

\eatline
\nwendcode{}\nwbegindocs{2778}\nwdocspar
\subsection{User-controlled output to data files}

\subsubsection{Open file (output)}
Open a file for output.
\nwenddocs{}\nwbegincode{2779}\moddef{Commands: types}\plusendmoddef
  type, extends (command_t) :: cmd_open_out_t
     private
     type(parse_node_t), pointer :: file_expr => null ()
   contains
   \LA{}Commands: cmd open out: TBP\RA{}
  end type cmd_open_out_t

\eatline
\nwendcode{}\nwbegindocs{2780}Finalizer for the embedded eval tree.
\nwenddocs{}\nwbegincode{2781}\moddef{Commands: procedures}\plusendmoddef
  subroutine cmd_open_out_final (object)
    class(cmd_open_out_t), intent(inout) :: object
  end subroutine cmd_open_out_final
  
\eatline
\nwendcode{}\nwbegindocs{2782}Output (trivial here).
\nwenddocs{}\nwbegincode{2783}\moddef{Commands: cmd open out: TBP}\endmoddef
  procedure :: write => cmd_open_out_write
\nwendcode{}\nwbegincode{2784}\moddef{Commands: procedures}\plusendmoddef
  subroutine cmd_open_out_write (cmd, unit, indent)
    class(cmd_open_out_t), intent(in) :: cmd
    integer, intent(in), optional :: unit, indent
    integer :: u
    u = given_output_unit (unit);  if (u < 0)  return
    call write_indent (u, indent)  
    write (u, "(1x,A)", advance="no")  "open_out: <filename>"  
  end subroutine cmd_open_out_write
  
\eatline
\nwendcode{}\nwbegindocs{2785}Compile: create an eval tree for the filename expression.
\nwenddocs{}\nwbegincode{2786}\moddef{Commands: cmd open out: TBP}\plusendmoddef
  procedure :: compile => cmd_open_out_compile
\nwendcode{}\nwbegincode{2787}\moddef{Commands: procedures}\plusendmoddef
  subroutine cmd_open_out_compile (cmd, global)
    class(cmd_open_out_t), intent(inout) :: cmd
    type(rt_data_t), intent(inout), target :: global
    cmd%file_expr => parse_node_get_sub_ptr (cmd%pn, 2)
    if (associated (cmd%file_expr)) then
       cmd%pn_opt => parse_node_get_next_ptr (cmd%file_expr)
    end if
    call cmd%compile_options (global)
  end subroutine cmd_open_out_compile

\eatline
\nwendcode{}\nwbegindocs{2788}Execute: append the file to the global list of open files.
\nwenddocs{}\nwbegincode{2789}\moddef{Commands: cmd open out: TBP}\plusendmoddef
  procedure :: execute => cmd_open_out_execute
\nwendcode{}\nwbegincode{2790}\moddef{Commands: procedures}\plusendmoddef
  subroutine cmd_open_out_execute (cmd, global)
    class(cmd_open_out_t), intent(inout) :: cmd
    type(rt_data_t), intent(inout), target :: global
    type(var_list_t), pointer :: var_list
    type(eval_tree_t) :: file_expr
    type(string_t) :: file
    var_list => cmd%local%get_var_list_ptr ()
    call file_expr%init_sexpr (cmd%file_expr, var_list)
    call file_expr%evaluate ()
    if (file_expr%is_known ()) then
       file = file_expr%get_string ()
       call file_list_open (global%out_files, file, &
            action = "write", status = "replace", position = "asis")
    else
       call msg_fatal ("open_out: file name argument evaluates to unknown")
    end if
    call file_expr%final ()
  end subroutine cmd_open_out_execute

\eatline
\nwendcode{}\nwbegindocs{2791}\nwdocspar
\subsubsection{Open file (output)}
Close an output file.  Except for the {\tt{}execute} method, everything is
analogous to the open command, so we can just inherit.
\nwenddocs{}\nwbegincode{2792}\moddef{Commands: types}\plusendmoddef
  type, extends (cmd_open_out_t) :: cmd_close_out_t
     private
   contains
   \LA{}Commands: cmd close out: TBP\RA{}
  end type cmd_close_out_t

\eatline
\nwendcode{}\nwbegindocs{2793}Execute: remove the file from the global list of output files.
\nwenddocs{}\nwbegincode{2794}\moddef{Commands: cmd close out: TBP}\endmoddef
  procedure :: execute => cmd_close_out_execute
\nwendcode{}\nwbegincode{2795}\moddef{Commands: procedures}\plusendmoddef
  subroutine cmd_close_out_execute (cmd, global)
    class(cmd_close_out_t), intent(inout) :: cmd
    type(rt_data_t), intent(inout), target :: global
    type(var_list_t), pointer :: var_list
    type(eval_tree_t) :: file_expr
    type(string_t) :: file
    var_list => cmd%local%var_list
    call file_expr%init_sexpr (cmd%file_expr, var_list)
    call file_expr%evaluate ()
    if (file_expr%is_known ()) then
       file = file_expr%get_string ()
       call file_list_close (global%out_files, file)
    else
       call msg_fatal ("close_out: file name argument evaluates to unknown")
    end if
    call file_expr%final ()
  end subroutine cmd_close_out_execute

\eatline
\nwendcode{}\nwbegindocs{2796}\nwdocspar
\subsection{Print custom-formatted values}
\nwenddocs{}\nwbegincode{2797}\moddef{Commands: types}\plusendmoddef
  type, extends (command_t) :: cmd_printf_t
     private
     type(parse_node_t), pointer :: sexpr => null ()
     type(parse_node_t), pointer :: sprintf_fun => null ()
     type(parse_node_t), pointer :: sprintf_clause => null ()
     type(parse_node_t), pointer :: sprintf => null ()
   contains
   \LA{}Commands: cmd printf: TBP\RA{}
  end type cmd_printf_t

\eatline
\nwendcode{}\nwbegindocs{2798}Finalize.
\nwenddocs{}\nwbegincode{2799}\moddef{Commands: cmd printf: TBP}\endmoddef
  procedure :: final => cmd_printf_final
\nwendcode{}\nwbegincode{2800}\moddef{Commands: procedures}\plusendmoddef
  subroutine cmd_printf_final (cmd)
    class(cmd_printf_t), intent(inout) :: cmd
    call parse_node_final (cmd%sexpr, recursive = .false.)
    deallocate (cmd%sexpr)
    call parse_node_final (cmd%sprintf_fun, recursive = .false.)
    deallocate (cmd%sprintf_fun)
    call parse_node_final (cmd%sprintf_clause, recursive = .false.)
    deallocate (cmd%sprintf_clause)
    call parse_node_final (cmd%sprintf, recursive = .false.)
    deallocate (cmd%sprintf)
  end subroutine cmd_printf_final

\eatline
\nwendcode{}\nwbegindocs{2801}Output.  Do not print the parse tree, since this may get cluttered.
Just a message that cuts have been defined.
\nwenddocs{}\nwbegincode{2802}\moddef{Commands: cmd printf: TBP}\plusendmoddef
  procedure :: write => cmd_printf_write
\nwendcode{}\nwbegincode{2803}\moddef{Commands: procedures}\plusendmoddef
  subroutine cmd_printf_write (cmd, unit, indent)
    class(cmd_printf_t), intent(in) :: cmd
    integer, intent(in), optional :: unit, indent
    integer :: u
    u = given_output_unit (unit);  if (u < 0)  return
    call write_indent (u, indent)
    write (u, "(1x,A)")  "printf:"
  end subroutine cmd_printf_write
  
\eatline
\nwendcode{}\nwbegindocs{2804}Compile.  We create a fake parse node (subtree) with a {\tt{}sprintf} command
with identical arguments which can then be handled by the corresponding
evaluation procedure.
\nwenddocs{}\nwbegincode{2805}\moddef{Commands: cmd printf: TBP}\plusendmoddef
  procedure :: compile => cmd_printf_compile
\nwendcode{}\nwbegincode{2806}\moddef{Commands: procedures}\plusendmoddef
  subroutine cmd_printf_compile (cmd, global)
    class(cmd_printf_t), intent(inout) :: cmd
    type(rt_data_t), intent(inout), target :: global
    type(parse_node_t), pointer :: pn_cmd, pn_clause, pn_args, pn_format
    pn_cmd => parse_node_get_sub_ptr (cmd%pn)
    pn_clause => parse_node_get_sub_ptr (pn_cmd)
    pn_format => parse_node_get_sub_ptr (pn_clause, 2)
    pn_args => parse_node_get_next_ptr (pn_clause)
    cmd%pn_opt => parse_node_get_next_ptr (pn_cmd)
    call cmd%compile_options (global)
    allocate (cmd%sexpr)
    call parse_node_create_branch (cmd%sexpr, &
         syntax_get_rule_ptr (syntax_cmd_list, var_str ("sexpr")))
    allocate (cmd%sprintf_fun)
    call parse_node_create_branch (cmd%sprintf_fun, &
         syntax_get_rule_ptr (syntax_cmd_list, var_str ("sprintf_fun")))
    allocate (cmd%sprintf_clause)
    call parse_node_create_branch (cmd%sprintf_clause, &
         syntax_get_rule_ptr (syntax_cmd_list, var_str ("sprintf_clause")))
    allocate (cmd%sprintf)
    call parse_node_create_key (cmd%sprintf, &
         syntax_get_rule_ptr (syntax_cmd_list, var_str ("sprintf")))
    call parse_node_append_sub (cmd%sprintf_clause, cmd%sprintf)
    call parse_node_append_sub (cmd%sprintf_clause, pn_format)
    call parse_node_freeze_branch (cmd%sprintf_clause)
    call parse_node_append_sub (cmd%sprintf_fun, cmd%sprintf_clause)
    if (associated (pn_args)) then
       call parse_node_append_sub (cmd%sprintf_fun, pn_args)
    end if
    call parse_node_freeze_branch (cmd%sprintf_fun)
    call parse_node_append_sub (cmd%sexpr, cmd%sprintf_fun)
    call parse_node_freeze_branch (cmd%sexpr)
  end subroutine cmd_printf_compile

\eatline
\nwendcode{}\nwbegindocs{2807}Execute.  Evaluate the string (pretending this is a {\tt{}sprintf} expression)
and print it.
\nwenddocs{}\nwbegincode{2808}\moddef{Commands: cmd printf: TBP}\plusendmoddef
  procedure :: execute => cmd_printf_execute
\nwendcode{}\nwbegincode{2809}\moddef{Commands: procedures}\plusendmoddef
  subroutine cmd_printf_execute (cmd, global)
    class(cmd_printf_t), intent(inout) :: cmd
    type(rt_data_t), intent(inout), target :: global
    type(var_list_t), pointer :: var_list
    type(string_t) :: string, file
    type(eval_tree_t) :: sprintf_expr
    logical :: advance
    var_list => cmd%local%get_var_list_ptr ()
    advance = var_list%get_lval (&
         var_str ("?out_advance"))
    file = var_list%get_sval (&
         var_str ("$out_file"))
    call sprintf_expr%init_sexpr (cmd%sexpr, var_list)
    call sprintf_expr%evaluate ()
    if (sprintf_expr%is_known ()) then
       string = sprintf_expr%get_string ()
       if (len (file) == 0) then
          call msg_result (char (string))
       else
          call file_list_write (global%out_files, file, string, advance)
       end if
    end if
  end subroutine cmd_printf_execute

\eatline
\nwendcode{}\nwbegindocs{2810}\nwdocspar
\subsubsection{Record data}
The expression syntax already contains a {\tt{}record} keyword; this evaluates to
a logical which is always true, but it has the side-effect of recording data
into analysis objects.  Here we define a command as an interface to this
construct.
\nwenddocs{}\nwbegincode{2811}\moddef{Commands: types}\plusendmoddef
  type, extends (command_t) :: cmd_record_t
     private
     type(parse_node_t), pointer :: pn_lexpr => null ()
   contains
   \LA{}Commands: cmd record: TBP\RA{}
  end type cmd_record_t
     
\eatline
\nwendcode{}\nwbegindocs{2812}Output.  With the compile hack below, there is nothing of interest
to print here.
\nwenddocs{}\nwbegincode{2813}\moddef{Commands: cmd record: TBP}\endmoddef
  procedure :: write => cmd_record_write
\nwendcode{}\nwbegincode{2814}\moddef{Commands: procedures}\plusendmoddef
  subroutine cmd_record_write (cmd, unit, indent)
    class(cmd_record_t), intent(in) :: cmd
    integer, intent(in), optional :: unit, indent
    integer :: u
    u = given_output_unit (unit);  if (u < 0)  return
    call write_indent (u, indent)
    write (u, "(1x,A)")  "record"
  end subroutine cmd_record_write
  
\eatline
\nwendcode{}\nwbegindocs{2815}Compile.  This is a hack which transforms the {\tt{}record} command
into a {\tt{}record} expression, which we handle in the {\tt{}expressions}
module.
\nwenddocs{}\nwbegincode{2816}\moddef{Commands: cmd record: TBP}\plusendmoddef
  procedure :: compile => cmd_record_compile
\nwendcode{}\nwbegincode{2817}\moddef{Commands: procedures}\plusendmoddef
  subroutine cmd_record_compile (cmd, global)
    class(cmd_record_t), intent(inout) :: cmd
    type(rt_data_t), intent(inout), target :: global
    type(parse_node_t), pointer :: pn_lexpr, pn_lsinglet, pn_lterm, pn_record
    call parse_node_create_branch (pn_lexpr, &
         syntax_get_rule_ptr (syntax_cmd_list, var_str ("lexpr")))
    call parse_node_create_branch (pn_lsinglet, &
         syntax_get_rule_ptr (syntax_cmd_list, var_str ("lsinglet")))
    call parse_node_append_sub (pn_lexpr, pn_lsinglet)
    call parse_node_create_branch (pn_lterm, &
         syntax_get_rule_ptr (syntax_cmd_list, var_str ("lterm")))
    call parse_node_append_sub (pn_lsinglet, pn_lterm)
    pn_record => parse_node_get_sub_ptr (cmd%pn)
    call parse_node_append_sub (pn_lterm, pn_record)
    cmd%pn_lexpr => pn_lexpr
  end subroutine cmd_record_compile

\eatline
\nwendcode{}\nwbegindocs{2818}Command execution.  Again, transfer this to the embedded expression
and just forget the logical result.
\nwenddocs{}\nwbegincode{2819}\moddef{Commands: cmd record: TBP}\plusendmoddef
  procedure :: execute => cmd_record_execute
\nwendcode{}\nwbegincode{2820}\moddef{Commands: procedures}\plusendmoddef
  subroutine cmd_record_execute (cmd, global)
    class(cmd_record_t), intent(inout) :: cmd
    type(rt_data_t), intent(inout), target :: global
    type(var_list_t), pointer :: var_list
    logical :: lval
    var_list => global%get_var_list_ptr ()
    lval = eval_log (cmd%pn_lexpr, var_list)
  end subroutine cmd_record_execute

\eatline
\nwendcode{}\nwbegindocs{2821}\nwdocspar
\subsubsection{Unstable particles}
Mark a particle as unstable.  For each unstable particle, we store a
number of decay channels and compute their respective BRs.
\nwenddocs{}\nwbegincode{2822}\moddef{Commands: types}\plusendmoddef
  type, extends (command_t) :: cmd_unstable_t
     private
     integer :: n_proc = 0
     type(string_t), dimension(:), allocatable :: process_id
     type(parse_node_t), pointer :: pn_prt_in => null ()
   contains
   \LA{}Commands: cmd unstable: TBP\RA{}
  end type cmd_unstable_t
     
\eatline
\nwendcode{}\nwbegindocs{2823}Output: we know the process IDs.
\nwenddocs{}\nwbegincode{2824}\moddef{Commands: cmd unstable: TBP}\endmoddef
  procedure :: write => cmd_unstable_write
\nwendcode{}\nwbegincode{2825}\moddef{Commands: procedures}\plusendmoddef
  subroutine cmd_unstable_write (cmd, unit, indent)
    class(cmd_unstable_t), intent(in) :: cmd
    integer, intent(in), optional :: unit, indent
    integer :: u, i
    u = given_output_unit (unit);  if (u < 0)  return
    call write_indent (u, indent)
    write (u, "(1x,A,1x,I0,1x,A)", advance="no")  &
         "unstable:", 1, "("
    do i = 1, cmd%n_proc
       if (i > 1)  write (u, "(A,1x)", advance="no")  ","
       write (u, "(A)", advance="no")  char (cmd%process_id(i))
    end do
    write (u, "(A)")  ")"
  end subroutine cmd_unstable_write

\eatline
\nwendcode{}\nwbegindocs{2826}Compile.  Initiate an eval tree for the decaying particle and
determine the decay channel process IDs.
\nwenddocs{}\nwbegincode{2827}\moddef{Commands: cmd unstable: TBP}\plusendmoddef
  procedure :: compile => cmd_unstable_compile
\nwendcode{}\nwbegincode{2828}\moddef{Commands: procedures}\plusendmoddef
  subroutine cmd_unstable_compile (cmd, global)
    class(cmd_unstable_t), intent(inout) :: cmd
    type(rt_data_t), intent(inout), target :: global
    type(parse_node_t), pointer :: pn_list, pn_proc
    integer :: i
    cmd%pn_prt_in => parse_node_get_sub_ptr (cmd%pn, 2)
    pn_list => parse_node_get_next_ptr (cmd%pn_prt_in)
    if (associated (pn_list)) then
       select case (char (parse_node_get_rule_key (pn_list)))
       case ("unstable_arg")
          cmd%n_proc = parse_node_get_n_sub (pn_list)
          cmd%pn_opt => parse_node_get_next_ptr (pn_list)
       case default
          cmd%n_proc = 0
          cmd%pn_opt => pn_list
          pn_list => null ()
       end select
    end if       
    call cmd%compile_options (global)
    if (associated (pn_list)) then
       allocate (cmd%process_id (cmd%n_proc))
       pn_proc => parse_node_get_sub_ptr (pn_list)
       do i = 1, cmd%n_proc
          cmd%process_id(i) = parse_node_get_string (pn_proc)
          call cmd%local%process_stack%init_result_vars (cmd%process_id(i))
          pn_proc => parse_node_get_next_ptr (pn_proc)
       end do
    else
       allocate (cmd%process_id (0))
    end if
  end subroutine cmd_unstable_compile

\eatline
\nwendcode{}\nwbegindocs{2829}Command execution.  Evaluate the decaying particle and mark the decays in
the current model object.
\nwenddocs{}\nwbegincode{2830}\moddef{Commands: cmd unstable: TBP}\plusendmoddef
  procedure :: execute => cmd_unstable_execute
\nwendcode{}\nwbegincode{2831}\moddef{Commands: procedures}\plusendmoddef
  subroutine cmd_unstable_execute (cmd, global)
    class(cmd_unstable_t), intent(inout) :: cmd
    type(rt_data_t), intent(inout), target :: global
    type(var_list_t), pointer :: var_list
    logical :: auto_decays, auto_decays_radiative
    integer :: auto_decays_multiplicity
    logical :: isotropic_decay, diagonal_decay, polarized_decay
    integer :: decay_helicity
    type(pdg_array_t) :: pa_in
    integer :: pdg_in
    type(string_t) :: libname_cur, libname_dec
    type(string_t), dimension(:), allocatable :: auto_id, tmp_id
    integer :: n_proc_user
    integer :: i, u_tmp
    character(80) :: buffer
    var_list => cmd%local%get_var_list_ptr ()
    auto_decays = &
         var_list%get_lval (var_str ("?auto_decays"))
    if (auto_decays) then
       auto_decays_multiplicity = &
            var_list%get_ival (var_str ("auto_decays_multiplicity"))
       auto_decays_radiative = &
            var_list%get_lval (var_str ("?auto_decays_radiative"))
    end if
    isotropic_decay = &
         var_list%get_lval (var_str ("?isotropic_decay"))
    if (isotropic_decay) then
       diagonal_decay = .false.
       polarized_decay = .false.
    else
       diagonal_decay = &
            var_list%get_lval (var_str ("?diagonal_decay"))
       if (diagonal_decay) then
          polarized_decay = .false.
       else
          polarized_decay = &
               var_list%is_known (var_str ("decay_helicity"))
          if (polarized_decay) then
             decay_helicity = var_list%get_ival (var_str ("decay_helicity"))
          end if
       end if
    end if
    pa_in = eval_pdg_array (cmd%pn_prt_in, var_list)
    if (pdg_array_get_length (pa_in) /= 1) &
         call msg_fatal ("Unstable: decaying particle must be unique")
    pdg_in = pdg_array_get (pa_in, 1)
    n_proc_user = cmd%n_proc
    if (auto_decays) then
       call create_auto_decays (pdg_in, &
            auto_decays_multiplicity, auto_decays_radiative, &
            libname_dec, auto_id, cmd%local)
       allocate (tmp_id (cmd%n_proc + size (auto_id)))
       tmp_id(:cmd%n_proc) = cmd%process_id
       tmp_id(cmd%n_proc+1:) = auto_id
       call move_alloc (from = tmp_id, to = cmd%process_id)
       cmd%n_proc = size (cmd%process_id)
    end if
    libname_cur = cmd%local%prclib%get_name ()
    do i = 1, cmd%n_proc
       if (i == n_proc_user + 1) then
          call cmd%local%update_prclib &
               (cmd%local%prclib_stack%get_library_ptr (libname_dec))
       end if
       if (.not. global%process_stack%exists (cmd%process_id(i))) then
          call var_list_set_log (var_list, &
               var_str ("?decay_rest_frame"), .false., is_known = .true.)
          call integrate_process (cmd%process_id(i), cmd%local, global)
          call global%process_stack%fill_result_vars (cmd%process_id(i))
       end if
    end do
    call cmd%local%update_prclib &
         (cmd%local%prclib_stack%get_library_ptr (libname_cur))
    if (cmd%n_proc > 0) then
       if (polarized_decay) then
          call global%modify_particle (pdg_in, stable = .false., &
               decay = cmd%process_id, &
               isotropic_decay = .false., &
               diagonal_decay = .false., &
               decay_helicity = decay_helicity, &
               polarized = .false.)
       else
          call global%modify_particle (pdg_in, stable = .false., &
               decay = cmd%process_id, &
               isotropic_decay = isotropic_decay, &
               diagonal_decay = diagonal_decay, &
               polarized = .false.)
       end if
       u_tmp = free_unit ()
       open (u_tmp, status = "scratch", action = "readwrite")
       call show_unstable (global, pdg_in, u_tmp)
       rewind (u_tmp)
       do
          read (u_tmp, "(A)", end = 1)  buffer
          write (msg_buffer, "(A)")  trim (buffer)
          call msg_message ()
       end do
1      continue
       close (u_tmp)
    else
       call err_unstable (global, pdg_in)
    end if
  end subroutine cmd_unstable_execute

\eatline
\nwendcode{}\nwbegindocs{2832}Show data for the current unstable particle.  This is called both by
the {\tt{}unstable} and by the {\tt{}show} command.
\nwenddocs{}\nwbegincode{2833}\moddef{Commands: procedures}\plusendmoddef
  subroutine show_unstable (global, pdg, u)
    type(rt_data_t), intent(in), target :: global
    integer, intent(in) :: pdg, u
    type(flavor_t) :: flv
    type(string_t), dimension(:), allocatable :: decay
    real(default), dimension(:), allocatable :: br
    real(default) :: width
    type(process_t), pointer :: process
    type(process_component_def_t), pointer :: prc_def
    type(string_t), dimension(:), allocatable :: prt_out, prt_out_str
    integer :: i, j
    call flv%init (pdg, global%model)
    call flv%get_decays (decay)
    if (.not. allocated (decay))  return
    allocate (prt_out_str (size (decay)))
    allocate (br (size (decay)))
    do i = 1, size (br)
       process => global%process_stack%get_process_ptr (decay(i))
       prc_def => process%get_component_def_ptr (1)
       call prc_def%get_prt_out (prt_out)
       prt_out_str(i) = prt_out(1)
       do j = 2, size (prt_out)
          prt_out_str(i) = prt_out_str(i) // ", " // prt_out(j)
       end do
       br(i) = process%get_integral ()
    end do
    if (all (br >= 0)) then
       if (any (br > 0)) then
          width = sum (br)
          br = br / sum (br)
          write (u, "(A)") "Unstable particle " &
               // char (flv%get_name ()) &
               // ": computed branching ratios:"
          do i = 1, size (br)
             write (u, "(2x,A,':'," // FMT_14 // ",3x,A)") &
                  char (decay(i)), br(i), char (prt_out_str(i))
          end do
          write (u, "(2x,'Total width ='," // FMT_14 // ",' GeV (computed)')")  width
          write (u, "(2x,'            ='," // FMT_14 // ",' GeV (preset)')") &
               flv%get_width ()
          if (flv%decays_isotropically ()) then
             write (u, "(2x,A)")  "Decay options: isotropic"
          else if (flv%decays_diagonal ()) then
             write (u, "(2x,A)")  "Decay options: &
                  &projection on diagonal helicity states"
          else if (flv%has_decay_helicity ()) then
             write (u, "(2x,A,1x,I0)")  "Decay options: projection onto helicity =", &
                  flv%get_decay_helicity ()
          else
             write (u, "(2x,A)")  "Decay options: helicity treated exactly"
          end if
       else
          call msg_fatal ("Unstable particle " &
               // char (flv%get_name ()) &
               // ": partial width vanishes for all decay channels")
       end if
    else
       call msg_fatal ("Unstable particle " &
               // char (flv%get_name ()) &
               // ": partial width is negative")
    end if
  end subroutine show_unstable
    
\eatline
\nwendcode{}\nwbegindocs{2834}If no decays have been found, issue a non-fatal error.
\nwenddocs{}\nwbegincode{2835}\moddef{Commands: procedures}\plusendmoddef
  subroutine err_unstable (global, pdg)
    type(rt_data_t), intent(in), target :: global
    integer, intent(in) :: pdg
    type(flavor_t) :: flv
    call flv%init (pdg, global%model)
    call msg_error ("Unstable: no allowed decays found for particle " &
         // char (flv%get_name ()) // ", keeping as stable")
  end subroutine err_unstable
    
\eatline
\nwendcode{}\nwbegindocs{2836}Auto decays: create process IDs and make up process
configurations, using the PDG codes generated by the {\tt{}ds{\char95}table} make
method.

We allocate and use a self-contained process library that contains only the
decay processes of the current particle.  When done, we revert the global
library pointer to the original library but return the name of the new one.
The new library becomes part of the global library stack and can thus be
referred to at any time.
\nwenddocs{}\nwbegincode{2837}\moddef{Commands: procedures}\plusendmoddef
  subroutine create_auto_decays &
       (pdg_in, mult, rad, libname_dec, process_id, global)
    integer, intent(in) :: pdg_in
    integer, intent(in) :: mult
    logical, intent(in) :: rad
    type(string_t), intent(out) :: libname_dec
    type(string_t), dimension(:), allocatable, intent(out) :: process_id
    type(rt_data_t), intent(inout) :: global
    type(prclib_entry_t), pointer :: lib_entry
    type(process_library_t), pointer :: lib
    type(ds_table_t) :: ds_table
    type(split_constraints_t) :: constraints
    type(pdg_array_t), dimension(:), allocatable :: pa_out
    character(80) :: buffer
    character :: p_or_a
    type(string_t) :: process_string, libname_cur
    type(flavor_t) :: flv_in, flv_out
    type(string_t) :: prt_in
    type(string_t), dimension(:), allocatable :: prt_out
    type(process_configuration_t) :: prc_config
    integer :: i, j, k, n
    call flv_in%init (pdg_in, global%model)
    if (rad) then
       call constraints%init (2)
    else
       call constraints%init (3)
       call constraints%set (3, constrain_radiation ())
    end if
    call constraints%set (1, constrain_n_tot (mult))
    call constraints%set (2, &
         constrain_mass_sum (flv_in%get_mass (), margin = 0._default))
    call ds_table%make (global%model, pdg_in, constraints)
    prt_in = flv_in%get_name ()
    if (pdg_in > 0) then
       p_or_a = "p"
    else
       p_or_a = "a"
    end if
    if (ds_table%get_length () == 0) then
       call msg_warning ("Auto-decays: Particle " // char (prt_in) // ": " &
            // "no decays found")
       libname_dec = ""
       allocate (process_id (0))
    else
       call msg_message ("Creating decay process library for particle " &
            // char (prt_in))
       libname_cur = global%prclib%get_name () 
       write (buffer, "(A,A,I0)")  "_d", p_or_a, abs (pdg_in)
       libname_dec = libname_cur // trim (buffer)
       lib => global%prclib_stack%get_library_ptr (libname_dec)
       if (.not. (associated (lib))) then
          allocate (lib_entry)
          call lib_entry%init (libname_dec)
          lib => lib_entry%process_library_t
          call global%add_prclib (lib_entry)
       else
          call global%update_prclib (lib)
       end if
       allocate (process_id (ds_table%get_length ()))
       do i = 1, size (process_id)
          write (buffer, "(A,'_',A,I0,'_',I0)") &
               "decay", p_or_a, abs (pdg_in), i
          process_id(i) = trim (buffer)
          process_string = process_id(i) // ": " // prt_in // " =>"
          call ds_table%get_pdg_out (i, pa_out)
          allocate (prt_out (size (pa_out)))
          do j = 1, size (pa_out)
             do k = 1, pa_out(j)%get_length ()
                call flv_out%init (pa_out(j)%get (k), global%model)
                if (k == 1) then
                   prt_out(j) = flv_out%get_name ()
                else
                   prt_out(j) = prt_out(j) // ":" // flv_out%get_name ()
                end if
             end do
             process_string = process_string // " " // prt_out(j)
          end do
          call msg_message (char (process_string))
          call prc_config%init (process_id(i), 1, 1, global)
          !!! Causes runtime error with gfortran 4.9.1 
          ! call prc_config%setup_component (1, &
          !      new_prt_spec ([prt_in]), new_prt_spec (prt_out), global)
          !!! Workaround:
          call prc_config%setup_component (1, &
               [new_prt_spec (prt_in)], new_prt_spec (prt_out), global)       
          call prc_config%record (global)
          deallocate (prt_out)
          deallocate (pa_out)
       end do
       lib => global%prclib_stack%get_library_ptr (libname_cur)
       call global%update_prclib (lib)
    end if
    call ds_table%final ()
  end subroutine create_auto_decays
    
\eatline
\nwendcode{}\nwbegindocs{2838}\nwdocspar
\subsubsection{(Stable particles}
Revert the unstable declaration for a list of particles.
\nwenddocs{}\nwbegincode{2839}\moddef{Commands: types}\plusendmoddef
  type, extends (command_t) :: cmd_stable_t
     private
     type(parse_node_p), dimension(:), allocatable :: pn_pdg
   contains
   \LA{}Commands: cmd stable: TBP\RA{}
  end type cmd_stable_t
  
\eatline
\nwendcode{}\nwbegindocs{2840}Output: we know only the number of particles.
\nwenddocs{}\nwbegincode{2841}\moddef{Commands: cmd stable: TBP}\endmoddef
  procedure :: write => cmd_stable_write
\nwendcode{}\nwbegincode{2842}\moddef{Commands: procedures}\plusendmoddef
  subroutine cmd_stable_write (cmd, unit, indent)
    class(cmd_stable_t), intent(in) :: cmd
    integer, intent(in), optional :: unit, indent
    integer :: u
    u = given_output_unit (unit);  if (u < 0)  return
    call write_indent (u, indent)
    write (u, "(1x,A,1x,I0)")  "stable:", size (cmd%pn_pdg)
  end subroutine cmd_stable_write

\eatline
\nwendcode{}\nwbegindocs{2843}Compile.  Assign parse nodes for the particle IDs.
\nwenddocs{}\nwbegincode{2844}\moddef{Commands: cmd stable: TBP}\plusendmoddef
  procedure :: compile => cmd_stable_compile
\nwendcode{}\nwbegincode{2845}\moddef{Commands: procedures}\plusendmoddef
  subroutine cmd_stable_compile (cmd, global)
    class(cmd_stable_t), intent(inout) :: cmd
    type(rt_data_t), intent(inout), target :: global
    type(parse_node_t), pointer :: pn_list, pn_prt
    integer :: n, i
    pn_list => parse_node_get_sub_ptr (cmd%pn, 2)
    cmd%pn_opt => parse_node_get_next_ptr (pn_list)
    call cmd%compile_options (global)
    n = parse_node_get_n_sub (pn_list)
    allocate (cmd%pn_pdg (n))
    pn_prt => parse_node_get_sub_ptr (pn_list)
    i = 1
    do while (associated (pn_prt))
       cmd%pn_pdg(i)%ptr => pn_prt
       pn_prt  => parse_node_get_next_ptr (pn_prt)
       i = i + 1
    end do
  end subroutine cmd_stable_compile

\eatline
\nwendcode{}\nwbegindocs{2846}Execute: apply the modifications to the current model.
\nwenddocs{}\nwbegincode{2847}\moddef{Commands: cmd stable: TBP}\plusendmoddef
  procedure :: execute => cmd_stable_execute
\nwendcode{}\nwbegincode{2848}\moddef{Commands: procedures}\plusendmoddef
  subroutine cmd_stable_execute (cmd, global)
    class(cmd_stable_t), intent(inout) :: cmd
    type(rt_data_t), target, intent(inout) :: global
    type(var_list_t), pointer :: var_list
    type(pdg_array_t) :: pa
    integer :: pdg
    type(flavor_t) :: flv
    integer :: i
    var_list => cmd%local%get_var_list_ptr ()
    do i = 1, size (cmd%pn_pdg)
       pa = eval_pdg_array (cmd%pn_pdg(i)%ptr, var_list)
       if (pdg_array_get_length (pa) /= 1) &
            call msg_fatal ("Stable: listed particles must be unique")
       pdg = pdg_array_get (pa, 1)
       call global%modify_particle (pdg, stable = .true., &
         isotropic_decay = .false., &
         diagonal_decay = .false., &
         polarized = .false.)
       call flv%init (pdg, cmd%local%model)
       call msg_message ("Particle " &
            // char (flv%get_name ()) &
            // " declared as stable")
    end do
  end subroutine cmd_stable_execute
  
\eatline
\nwendcode{}\nwbegindocs{2849}\nwdocspar
\subsubsection{Polarized particles}
These commands mark particles as (un)polarized, to be applied in
subsequent simulation passes.  Since this is technically the same as
the {\tt{}stable} command, we take a shortcut and make this an extension,
just overriding methods.
\nwenddocs{}\nwbegincode{2850}\moddef{Commands: types}\plusendmoddef
  type, extends (cmd_stable_t) :: cmd_polarized_t
   contains
   \LA{}Commands: cmd polarized: TBP\RA{}
  end type cmd_polarized_t
  
  type, extends (cmd_stable_t) :: cmd_unpolarized_t
   contains
   \LA{}Commands: cmd unpolarized: TBP\RA{}
  end type cmd_unpolarized_t
  
\eatline
\nwendcode{}\nwbegindocs{2851}Output: we know only the number of particles.
\nwenddocs{}\nwbegincode{2852}\moddef{Commands: cmd polarized: TBP}\endmoddef
  procedure :: write => cmd_polarized_write
\nwendcode{}\nwbegincode{2853}\moddef{Commands: cmd unpolarized: TBP}\endmoddef
  procedure :: write => cmd_unpolarized_write
\nwendcode{}\nwbegincode{2854}\moddef{Commands: procedures}\plusendmoddef
  subroutine cmd_polarized_write (cmd, unit, indent)
    class(cmd_polarized_t), intent(in) :: cmd
    integer, intent(in), optional :: unit, indent
    integer :: u
    u = given_output_unit (unit);  if (u < 0)  return
    call write_indent (u, indent)
    write (u, "(1x,A,1x,I0)")  "polarized:", size (cmd%pn_pdg)
  end subroutine cmd_polarized_write

  subroutine cmd_unpolarized_write (cmd, unit, indent)
    class(cmd_unpolarized_t), intent(in) :: cmd
    integer, intent(in), optional :: unit, indent
    integer :: u
    u = given_output_unit (unit);  if (u < 0)  return
    call write_indent (u, indent)
    write (u, "(1x,A,1x,I0)")  "unpolarized:", size (cmd%pn_pdg)
  end subroutine cmd_unpolarized_write

\eatline
\eatline
\nwendcode{}\nwbegindocs{2855}Compile: accounted for by the base command.

Execute: apply the modifications to the current model.
\nwenddocs{}\nwbegincode{2856}\moddef{Commands: cmd polarized: TBP}\plusendmoddef
  procedure :: execute => cmd_polarized_execute
\nwendcode{}\nwbegincode{2857}\moddef{Commands: cmd unpolarized: TBP}\plusendmoddef
  procedure :: execute => cmd_unpolarized_execute
\nwendcode{}\nwbegincode{2858}\moddef{Commands: procedures}\plusendmoddef
  subroutine cmd_polarized_execute (cmd, global)
    class(cmd_polarized_t), intent(inout) :: cmd
    type(rt_data_t), target, intent(inout) :: global
    type(var_list_t), pointer :: var_list
    type(pdg_array_t) :: pa
    integer :: pdg
    type(flavor_t) :: flv
    integer :: i
    var_list => cmd%local%get_var_list_ptr ()
    do i = 1, size (cmd%pn_pdg)
       pa = eval_pdg_array (cmd%pn_pdg(i)%ptr, var_list)
       if (pdg_array_get_length (pa) /= 1) &
            call msg_fatal ("Polarized: listed particles must be unique")
       pdg = pdg_array_get (pa, 1)
       call global%modify_particle (pdg, polarized = .true., &
            stable = .true., &
            isotropic_decay = .false., &
            diagonal_decay = .false.)
       call flv%init (pdg, cmd%local%model)
       call msg_message ("Particle " &
            // char (flv%get_name ()) &
            // " declared as polarized")
    end do
  end subroutine cmd_polarized_execute
  
  subroutine cmd_unpolarized_execute (cmd, global)
    class(cmd_unpolarized_t), intent(inout) :: cmd
    type(rt_data_t), target, intent(inout) :: global
    type(var_list_t), pointer :: var_list
    type(pdg_array_t) :: pa
    integer :: pdg
    type(flavor_t) :: flv
    integer :: i
    var_list => cmd%local%get_var_list_ptr ()
    do i = 1, size (cmd%pn_pdg)
       pa = eval_pdg_array (cmd%pn_pdg(i)%ptr, var_list)
       if (pdg_array_get_length (pa) /= 1) &
            call msg_fatal ("Unpolarized: listed particles must be unique")
       pdg = pdg_array_get (pa, 1)
       call global%modify_particle (pdg, polarized = .false., &
            stable = .true., &
            isotropic_decay = .false., &
            diagonal_decay = .false.)
       call flv%init (pdg, cmd%local%model)
       call msg_message ("Particle " &
            // char (flv%get_name ()) &
            // " declared as unpolarized")
    end do
  end subroutine cmd_unpolarized_execute
  
\eatline
\eatline
\nwendcode{}\nwbegindocs{2859}\nwdocspar
\subsubsection{Parameters: formats for event-sample output}
Specify all event formats that are to be used for output files in the
subsequent simulation run.  (The raw format is on by default and can be turned
off here.)
\nwenddocs{}\nwbegincode{2860}\moddef{Commands: types}\plusendmoddef
  type, extends (command_t) :: cmd_sample_format_t
     private
     type(string_t), dimension(:), allocatable :: format
   contains
   \LA{}Commands: cmd sample format: TBP\RA{}
  end type cmd_sample_format_t

\eatline
\nwendcode{}\nwbegindocs{2861}Output: here, everything is known.
\nwenddocs{}\nwbegincode{2862}\moddef{Commands: cmd sample format: TBP}\endmoddef
  procedure :: write => cmd_sample_format_write
\nwendcode{}\nwbegincode{2863}\moddef{Commands: procedures}\plusendmoddef
  subroutine cmd_sample_format_write (cmd, unit, indent)
    class(cmd_sample_format_t), intent(in) :: cmd
    integer, intent(in), optional :: unit, indent
    integer :: u, i
    u = given_output_unit (unit);  if (u < 0)  return
    call write_indent (u, indent)
    write (u, "(1x,A)", advance="no")  "sample_format = "
    do i = 1, size (cmd%format)
       if (i > 1)  write (u, "(A,1x)", advance="no")  ","
       write (u, "(A)", advance="no")  char (cmd%format(i))
    end do
    write (u, "(A)")
  end subroutine cmd_sample_format_write
  
\eatline
\nwendcode{}\nwbegindocs{2864}Compile.  Initialize evaluation trees.
\nwenddocs{}\nwbegincode{2865}\moddef{Commands: cmd sample format: TBP}\plusendmoddef
  procedure :: compile => cmd_sample_format_compile
\nwendcode{}\nwbegincode{2866}\moddef{Commands: procedures}\plusendmoddef
  subroutine cmd_sample_format_compile (cmd, global)
    class(cmd_sample_format_t), intent(inout) :: cmd
    type(rt_data_t), intent(inout), target :: global
    type(parse_node_t), pointer :: pn_arg
    type(parse_node_t), pointer :: pn_format
    integer :: i, n_format
    pn_arg => parse_node_get_sub_ptr (cmd%pn, 3)
    if (associated (pn_arg)) then
       n_format = parse_node_get_n_sub (pn_arg)
       allocate (cmd%format (n_format))
       pn_format => parse_node_get_sub_ptr (pn_arg)
       i = 0
       do while (associated (pn_format))
          i = i + 1
          cmd%format(i) = parse_node_get_string (pn_format)
          pn_format => parse_node_get_next_ptr (pn_format)
       end do
    else
       allocate (cmd%format (0))
    end if
  end subroutine cmd_sample_format_compile

\eatline
\nwendcode{}\nwbegindocs{2867}Execute.  Transfer the list of format specifications to the
corresponding array in the runtime data set.
\nwenddocs{}\nwbegincode{2868}\moddef{Commands: cmd sample format: TBP}\plusendmoddef
  procedure :: execute => cmd_sample_format_execute
\nwendcode{}\nwbegincode{2869}\moddef{Commands: procedures}\plusendmoddef
  subroutine cmd_sample_format_execute (cmd, global)
    class(cmd_sample_format_t), intent(inout) :: cmd
    type(rt_data_t), intent(inout), target :: global
    if (allocated (global%sample_fmt))  deallocate (global%sample_fmt)
    allocate (global%sample_fmt (size (cmd%format)), source = cmd%format)
  end subroutine cmd_sample_format_execute

\eatline
\nwendcode{}\nwbegindocs{2870}\nwdocspar
\subsubsection{The simulate command}
This is the actual SINDARIN command.
\nwenddocs{}\nwbegincode{2871}\moddef{Commands: types}\plusendmoddef
  type, extends (command_t) :: cmd_simulate_t
     ! not private anymore as required by the whizard-c-interface
     integer :: n_proc = 0
     type(string_t), dimension(:), allocatable :: process_id
   contains
   \LA{}Commands: cmd simulate: TBP\RA{}
  end type cmd_simulate_t

\eatline
\nwendcode{}\nwbegindocs{2872}Output: we know the process IDs.
\nwenddocs{}\nwbegincode{2873}\moddef{Commands: cmd simulate: TBP}\endmoddef
  procedure :: write => cmd_simulate_write
\nwendcode{}\nwbegincode{2874}\moddef{Commands: procedures}\plusendmoddef
  subroutine cmd_simulate_write (cmd, unit, indent)
    class(cmd_simulate_t), intent(in) :: cmd
    integer, intent(in), optional :: unit, indent
    integer :: u, i
    u = given_output_unit (unit);  if (u < 0)  return
    call write_indent (u, indent)
    write (u, "(1x,A)", advance="no")  "simulate ("
    do i = 1, cmd%n_proc
       if (i > 1)  write (u, "(A,1x)", advance="no")  ","
       write (u, "(A)", advance="no")  char (cmd%process_id(i))
    end do
    write (u, "(A)")  ")"
  end subroutine cmd_simulate_write

\eatline
\nwendcode{}\nwbegindocs{2875}Compile. In contrast to WHIZARD 1 the confusing option to give the
number of unweighted events for weighted events as if unweighting were
to take place has been abandoned. (We both use {\tt{}n{\char95}events} for
weighted and unweighted events, the variable {\tt{}n{\char95}calls} from WHIZARD
1 has been discarded.
\nwenddocs{}\nwbegincode{2876}\moddef{Commands: cmd simulate: TBP}\plusendmoddef
  procedure :: compile => cmd_simulate_compile
\nwendcode{}\nwbegincode{2877}\moddef{Commands: procedures}\plusendmoddef
  subroutine cmd_simulate_compile (cmd, global)
    class(cmd_simulate_t), intent(inout) :: cmd
    type(rt_data_t), intent(inout), target :: global
    type(parse_node_t), pointer :: pn_proclist, pn_proc
    integer :: i
    pn_proclist => parse_node_get_sub_ptr (cmd%pn, 2)
    cmd%pn_opt => parse_node_get_next_ptr (pn_proclist)
    call cmd%compile_options (global)
    cmd%n_proc = parse_node_get_n_sub (pn_proclist)
    allocate (cmd%process_id (cmd%n_proc))
    pn_proc => parse_node_get_sub_ptr (pn_proclist)
    do i = 1, cmd%n_proc
       cmd%process_id(i) = parse_node_get_string (pn_proc)
       call global%process_stack%init_result_vars (cmd%process_id(i))
       pn_proc => parse_node_get_next_ptr (pn_proc)
    end do
  end subroutine cmd_simulate_compile

\eatline
\nwendcode{}\nwbegindocs{2878}Execute command:  Simulate events.  This is done via a {\tt{}simulation{\char95}t}
object and its associated methods.

Signal handling: the {\tt{}generate} method may exit abnormally if there is a
pending signal.  The current logic ensures that the {\tt{}es{\char95}array} output
channels are closed before the {\tt{}execute} routine returns.  The program will
terminate then in {\tt{}command{\char95}list{\char95}execute}.
\nwenddocs{}\nwbegincode{2879}\moddef{Commands: cmd simulate: TBP}\plusendmoddef
  procedure :: execute => cmd_simulate_execute
\nwendcode{}\nwbegincode{2880}\moddef{Commands: procedures}\plusendmoddef
  subroutine cmd_simulate_execute (cmd, global)
    class(cmd_simulate_t), intent(inout) :: cmd
    type(rt_data_t), intent(inout), target :: global
    type(var_list_t), pointer :: var_list
    type(rt_data_t), dimension(:), allocatable, target :: alt_env
    integer :: n_events, n_fmt
    type(string_t) :: sample
    logical :: rebuild_events, read_raw, write_raw
    type(simulation_t), target :: sim
    type(string_t), dimension(:), allocatable :: sample_fmt
    type(event_stream_array_t) :: es_array
    type(event_sample_data_t) :: data
    integer :: i, checkpoint
    var_list => cmd%local%var_list
    if (allocated (cmd%local%pn%alt_setup)) then
       allocate (alt_env (size (cmd%local%pn%alt_setup)))
       do i = 1, size (alt_env)
          call build_alt_setup (alt_env(i), cmd%local, &
               cmd%local%pn%alt_setup(i)%ptr)
       end do
       call sim%init (cmd%process_id, .true., .true., cmd%local, global, &
            alt_env)
    else
       call sim%init (cmd%process_id, .true., .true., cmd%local, global)
    end if
    if (signal_is_pending ())  return
    if (sim%is_valid ()) then
       call sim%init_process_selector ()    
       call openmp_set_num_threads_verbose &
            (var_list%get_ival (var_str ("openmp_num_threads")), &
            var_list%get_lval (var_str ("?openmp_logging")))
       call sim%compute_n_events (n_events, var_list)
       sample = var_list%get_sval (var_str ("$sample"))
       if (sample == "")  sample = sim%get_default_sample_name ()
       rebuild_events = &
            var_list%get_lval (var_str ("?rebuild_events"))
       read_raw = &
            var_list%get_lval (var_str ("?read_raw")) &
            .and. .not. rebuild_events
       write_raw = &
            var_list%get_lval (var_str ("?write_raw"))
       checkpoint = &
            var_list%get_ival (var_str ("checkpoint"))
       if (read_raw) then
          inquire (file = char (sample) // ".evx", exist = read_raw)
       end if
       if (allocated (cmd%local%sample_fmt)) then
          n_fmt = size (cmd%local%sample_fmt)
       else
          n_fmt = 0
       end if
       data = sim%get_data ()
       data%n_evt = n_events
       if (read_raw) then
          allocate (sample_fmt (n_fmt))
          if (n_fmt > 0)  sample_fmt = cmd%local%sample_fmt
          call es_array%init (sample, &
               sample_fmt, cmd%local, &
               data = data, &
               input = var_str ("raw"), &
               allow_switch = write_raw, &
               checkpoint = checkpoint)
          call sim%generate (n_events, es_array)
          call es_array%final ()
       else if (write_raw) then
          allocate (sample_fmt (n_fmt + 1))
          if (n_fmt > 0)  sample_fmt(:n_fmt) = cmd%local%sample_fmt
          sample_fmt(n_fmt+1) = var_str ("raw")
          call es_array%init (sample, &
               sample_fmt, cmd%local, &
               data = data, &
               checkpoint = checkpoint)
          call sim%generate (n_events, es_array)
          call es_array%final ()
       else if (allocated (cmd%local%sample_fmt) .or. checkpoint > 0) then
          allocate (sample_fmt (n_fmt))
          if (n_fmt > 0)  sample_fmt = cmd%local%sample_fmt
          call es_array%init (sample, &
               sample_fmt, cmd%local, &
               data = data, &
               checkpoint = checkpoint)
          call sim%generate (n_events, es_array)
          call es_array%final ()
       else
          call sim%generate (n_events)
       end if
       if (allocated (alt_env)) then
          do i = 1, size (alt_env)
             call alt_env(i)%local_final ()
          end do
       end if
    end if
    call sim%final ()
  end subroutine cmd_simulate_execute

\eatline
\nwendcode{}\nwbegindocs{2881}Build an alternative setup: the parse tree is stored in the global
environment.  We create a temporary command list to compile and execute this;
the result is an alternative local environment {\tt{}alt{\char95}env} which we can hand
over to the {\tt{}simulate} command.
\nwenddocs{}\nwbegincode{2882}\moddef{Commands: procedures}\plusendmoddef
  recursive subroutine build_alt_setup (alt_env, global, pn)
    type(rt_data_t), intent(inout), target :: alt_env
    type(rt_data_t), intent(inout), target :: global
    type(parse_node_t), intent(in), target :: pn
    type(command_list_t), allocatable :: alt_options
    allocate (alt_options)
    call alt_env%local_init (global)
    call alt_env%activate ()
    call alt_options%compile (pn, alt_env)
    call alt_options%execute (alt_env)
    call alt_env%deactivate (global, keep_local = .true.)
    call alt_options%final ()
  end subroutine build_alt_setup
            
\eatline
\nwendcode{}\nwbegindocs{2883}\nwdocspar
\subsubsection{The rescan command}
This is the actual SINDARIN command.
\nwenddocs{}\nwbegincode{2884}\moddef{Commands: types}\plusendmoddef
  type, extends (command_t) :: cmd_rescan_t
     ! private
     type(parse_node_t), pointer :: pn_filename => null ()
     integer :: n_proc = 0
     type(string_t), dimension(:), allocatable :: process_id
   contains
   \LA{}Commands: cmd rescan: TBP\RA{}
  end type cmd_rescan_t

\eatline
\nwendcode{}\nwbegindocs{2885}Output: we know the process IDs.
\nwenddocs{}\nwbegincode{2886}\moddef{Commands: cmd rescan: TBP}\endmoddef
  procedure :: write => cmd_rescan_write
\nwendcode{}\nwbegincode{2887}\moddef{Commands: procedures}\plusendmoddef
  subroutine cmd_rescan_write (cmd, unit, indent)
    class(cmd_rescan_t), intent(in) :: cmd
    integer, intent(in), optional :: unit, indent
    integer :: u, i
    u = given_output_unit (unit);  if (u < 0)  return
    call write_indent (u, indent)
    write (u, "(1x,A)", advance="no")  "rescan ("
    do i = 1, cmd%n_proc
       if (i > 1)  write (u, "(A,1x)", advance="no")  ","
       write (u, "(A)", advance="no")  char (cmd%process_id(i))
    end do
    write (u, "(A)")  ")"
  end subroutine cmd_rescan_write

\eatline
\nwendcode{}\nwbegindocs{2888}Compile.  The command takes a suffix argument, namely the file name
of requested event file.
\nwenddocs{}\nwbegincode{2889}\moddef{Commands: cmd rescan: TBP}\plusendmoddef
  procedure :: compile => cmd_rescan_compile
\nwendcode{}\nwbegincode{2890}\moddef{Commands: procedures}\plusendmoddef
  subroutine cmd_rescan_compile (cmd, global)
    class(cmd_rescan_t), intent(inout) :: cmd
    type(rt_data_t), intent(inout), target :: global
    type(parse_node_t), pointer :: pn_filename, pn_proclist, pn_proc
    integer :: i
    pn_filename => parse_node_get_sub_ptr (cmd%pn, 2)
    pn_proclist => parse_node_get_next_ptr (pn_filename)
    cmd%pn_opt => parse_node_get_next_ptr (pn_proclist)
    call cmd%compile_options (global)
    cmd%pn_filename => pn_filename
    cmd%n_proc = parse_node_get_n_sub (pn_proclist)
    allocate (cmd%process_id (cmd%n_proc))
    pn_proc => parse_node_get_sub_ptr (pn_proclist)
    do i = 1, cmd%n_proc
       cmd%process_id(i) = parse_node_get_string (pn_proc)
       pn_proc => parse_node_get_next_ptr (pn_proc)
    end do
  end subroutine cmd_rescan_compile

\eatline
\nwendcode{}\nwbegindocs{2891}Execute command:  Rescan events.  This is done via a {\tt{}simulation{\char95}t}
object and its associated methods.
\nwenddocs{}\nwbegincode{2892}\moddef{Commands: cmd rescan: TBP}\plusendmoddef
  procedure :: execute => cmd_rescan_execute
\nwendcode{}\nwbegincode{2893}\moddef{Commands: procedures}\plusendmoddef
  subroutine cmd_rescan_execute (cmd, global)
    class(cmd_rescan_t), intent(inout) :: cmd
    type(rt_data_t), intent(inout), target :: global
    type(var_list_t), pointer :: var_list
    type(rt_data_t), dimension(:), allocatable, target :: alt_env
    type(string_t) :: sample
    logical :: exist, write_raw, update_event, update_sqme, update_weight
    type(simulation_t), target :: sim
    type(event_sample_data_t) :: input_data, data
    type(string_t) :: input_sample
    integer :: n_fmt
    type(string_t), dimension(:), allocatable :: sample_fmt
    type(string_t) :: input_format, input_ext, input_file
    type(string_t) :: lhef_extension, extension_hepmc, extension_lcio
    type(event_stream_array_t) :: es_array
    integer :: i, n_events
    var_list => cmd%local%var_list
    if (allocated (cmd%local%pn%alt_setup)) then
       allocate (alt_env (size (cmd%local%pn%alt_setup)))
       do i = 1, size (alt_env)
          call build_alt_setup (alt_env(i), cmd%local, &
               cmd%local%pn%alt_setup(i)%ptr)
       end do
       call sim%init (cmd%process_id, .false., .false., cmd%local, global, &
            alt_env)
    else
       call sim%init (cmd%process_id, .false., .false., cmd%local, global)
    end if
    call sim%compute_n_events (n_events, var_list)
    input_sample = eval_string (cmd%pn_filename, var_list)
    input_format = var_list%get_sval (&
         var_str ("$rescan_input_format"))
    sample = var_list%get_sval (var_str ("$sample"))
    if (sample == "")  sample = sim%get_default_sample_name ()
    write_raw = var_list%get_lval (var_str ("?write_raw"))
    if (allocated (cmd%local%sample_fmt)) then
       n_fmt = size (cmd%local%sample_fmt)
    else
       n_fmt = 0
    end if
    if (write_raw) then
       if (sample == input_sample) then
          call msg_error ("Rescan: ?write_raw = true: " &
               // "suppressing raw event output (filename clashes with input)")
          allocate (sample_fmt (n_fmt))
          if (n_fmt > 0)  sample_fmt = cmd%local%sample_fmt
       else
          allocate (sample_fmt (n_fmt + 1))
          if (n_fmt > 0)  sample_fmt(:n_fmt) = cmd%local%sample_fmt
          sample_fmt(n_fmt+1) = var_str ("raw")
       end if
    else
       allocate (sample_fmt (n_fmt))
       if (n_fmt > 0)  sample_fmt = cmd%local%sample_fmt
    end if
    update_event = &
         var_list%get_lval (var_str ("?update_event"))
    update_sqme = &
         var_list%get_lval (var_str ("?update_sqme"))
    update_weight = &
         var_list%get_lval (var_str ("?update_weight"))
    if (update_event .or. update_sqme) then
       call msg_message ("Recalculating observables")
       if (update_sqme) then
          call msg_message ("Recalculating squared matrix elements")
       end if
    end if
    lhef_extension = &
         var_list%get_sval (var_str ("$lhef_extension"))
    extension_hepmc = &
         var_list%get_sval (var_str ("$extension_hepmc"))
    extension_lcio = &
         var_list%get_sval (var_str ("$extension_lcio"))
    select case (char (input_format))
    case ("raw");  input_ext = "evx"
       call cmd%local%set_log &
            (var_str ("?recover_beams"), .false., is_known=.true.)
    case ("lhef"); input_ext = lhef_extension
    case ("hepmc"); input_ext = extension_hepmc
    case default
       call msg_fatal ("rescan: input sample format '" // char (input_format) &
            // "' not supported")
    end select
    input_file = input_sample // "." // input_ext
    inquire (file = char (input_file), exist = exist)
    if (exist) then
       input_data = sim%get_data (alt = .false.)
       input_data%n_evt = n_events
       data = sim%get_data ()
       data%n_evt = n_events
       input_data%md5sum_cfg = ""
       call es_array%init (sample, &
            sample_fmt, cmd%local, data, &
            input = input_format, input_sample = input_sample, &
            input_data = input_data, &
            allow_switch = .false.)
       call sim%rescan (n_events, es_array, global = cmd%local)
       call es_array%final ()
    else
       call msg_fatal ("Rescan: event file '" &
            // char (input_file) // "' not found")
    end if
    if (allocated (alt_env)) then
       do i = 1, size (alt_env)
          call alt_env(i)%local_final ()
       end do
    end if
    call sim%final ()
  end subroutine cmd_rescan_execute

\eatline
\nwendcode{}\nwbegindocs{2894}\nwdocspar
\subsubsection{Parameters: number of iterations}
Specify number of iterations and number of calls for one integration pass.
\nwenddocs{}\nwbegincode{2895}\moddef{Commands: types}\plusendmoddef
  type, extends (command_t) :: cmd_iterations_t
     private
     integer :: n_pass = 0
     type(parse_node_p), dimension(:), allocatable :: pn_expr_n_it
     type(parse_node_p), dimension(:), allocatable :: pn_expr_n_calls
     type(parse_node_p), dimension(:), allocatable :: pn_sexpr_adapt
   contains
   \LA{}Commands: cmd iterations: TBP\RA{}
  end type cmd_iterations_t

\eatline
\nwendcode{}\nwbegindocs{2896}Output.  Display the number of passes, which is known after compilation.
\nwenddocs{}\nwbegincode{2897}\moddef{Commands: cmd iterations: TBP}\endmoddef
  procedure :: write => cmd_iterations_write
\nwendcode{}\nwbegincode{2898}\moddef{Commands: procedures}\plusendmoddef
  subroutine cmd_iterations_write (cmd, unit, indent)
    class(cmd_iterations_t), intent(in) :: cmd
    integer, intent(in), optional :: unit, indent
    integer :: u
    u = given_output_unit (unit);  if (u < 0)  return
    call write_indent (u, indent)
    select case (cmd%n_pass)
    case (0)
       write (u, "(1x,A)")  "iterations: [empty]"
    case (1)
       write (u, "(1x,A,I0,A)")  "iterations: ", cmd%n_pass, " pass"
    case default
       write (u, "(1x,A,I0,A)")  "iterations: ", cmd%n_pass, " passes"
    end select
  end subroutine cmd_iterations_write

\eatline
\nwendcode{}\nwbegindocs{2899}Compile.  Initialize evaluation trees.
\nwenddocs{}\nwbegincode{2900}\moddef{Commands: cmd iterations: TBP}\plusendmoddef
  procedure :: compile => cmd_iterations_compile
\nwendcode{}\nwbegincode{2901}\moddef{Commands: procedures}\plusendmoddef
  subroutine cmd_iterations_compile (cmd, global)
    class(cmd_iterations_t), intent(inout) :: cmd
    type(rt_data_t), intent(inout), target :: global
    type(parse_node_t), pointer :: pn_arg, pn_n_it, pn_n_calls, pn_adapt
    type(parse_node_t), pointer :: pn_it_spec, pn_calls_spec, pn_adapt_spec
    integer :: i
    pn_arg => parse_node_get_sub_ptr (cmd%pn, 3)
    if (associated (pn_arg)) then
       cmd%n_pass = parse_node_get_n_sub (pn_arg)
       allocate (cmd%pn_expr_n_it (cmd%n_pass))
       allocate (cmd%pn_expr_n_calls (cmd%n_pass))
       allocate (cmd%pn_sexpr_adapt (cmd%n_pass))
       pn_it_spec => parse_node_get_sub_ptr (pn_arg)
       i = 1
       do while (associated (pn_it_spec))
          pn_n_it => parse_node_get_sub_ptr (pn_it_spec)
          pn_calls_spec => parse_node_get_next_ptr (pn_n_it)
          pn_n_calls => parse_node_get_sub_ptr (pn_calls_spec, 2)
          pn_adapt_spec => parse_node_get_next_ptr (pn_calls_spec)
          if (associated (pn_adapt_spec)) then
             pn_adapt => parse_node_get_sub_ptr (pn_adapt_spec, 2)
          else
             pn_adapt => null ()
          end if
          cmd%pn_expr_n_it(i)%ptr => pn_n_it
          cmd%pn_expr_n_calls(i)%ptr => pn_n_calls
          cmd%pn_sexpr_adapt(i)%ptr => pn_adapt
          i = i + 1
          pn_it_spec => parse_node_get_next_ptr (pn_it_spec)
       end do
    else
       allocate (cmd%pn_expr_n_it (0))
       allocate (cmd%pn_expr_n_calls (0))
    end if
  end subroutine cmd_iterations_compile

\eatline
\nwendcode{}\nwbegindocs{2902}Execute.  Evaluate the trees and transfer the results to the iteration
list in the runtime data set.
\nwenddocs{}\nwbegincode{2903}\moddef{Commands: cmd iterations: TBP}\plusendmoddef
  procedure :: execute => cmd_iterations_execute
\nwendcode{}\nwbegincode{2904}\moddef{Commands: procedures}\plusendmoddef
  subroutine cmd_iterations_execute (cmd, global)
    class(cmd_iterations_t), intent(inout) :: cmd
    type(rt_data_t), intent(inout), target :: global
    type(var_list_t), pointer :: var_list
    integer, dimension(cmd%n_pass) :: n_it, n_calls
    logical, dimension(cmd%n_pass) :: custom_adapt
    type(string_t), dimension(cmd%n_pass) :: adapt_code
    integer :: i
    var_list => global%get_var_list_ptr ()
    do i = 1, cmd%n_pass
       n_it(i) = eval_int (cmd%pn_expr_n_it(i)%ptr, var_list)
       n_calls(i) = &
            eval_int (cmd%pn_expr_n_calls(i)%ptr, var_list)
       if (associated (cmd%pn_sexpr_adapt(i)%ptr)) then
          adapt_code(i) = &
               eval_string (cmd%pn_sexpr_adapt(i)%ptr, &
                            var_list, is_known = custom_adapt(i))
       else
          custom_adapt(i) = .false.
       end if        
    end do
    call global%it_list%init (n_it, n_calls, custom_adapt, adapt_code)
  end subroutine cmd_iterations_execute

\eatline
\nwendcode{}\nwbegindocs{2905}\nwdocspar
\subsubsection{Range expressions}
We need a special type for storing and evaluating range expressions.
\nwenddocs{}\nwbegincode{2906}\moddef{Commands: parameters}\plusendmoddef
  integer, parameter :: STEP_NONE = 0
  integer, parameter :: STEP_ADD = 1
  integer, parameter :: STEP_SUB = 2
  integer, parameter :: STEP_MUL = 3
  integer, parameter :: STEP_DIV = 4
  integer, parameter :: STEP_COMP_ADD = 11
  integer, parameter :: STEP_COMP_MUL = 13
\nwendcode{}\nwbegindocs{2907}\nwdocspar
There is an abstract base type and two implementations: scan over integers and
scan over reals.
\nwenddocs{}\nwbegincode{2908}\moddef{Commands: types}\plusendmoddef
  type, abstract :: range_t
     type(parse_node_t), pointer :: pn_expr => null ()
     type(parse_node_t), pointer :: pn_term => null ()
     type(parse_node_t), pointer :: pn_factor => null ()
     type(parse_node_t), pointer :: pn_value => null ()
     type(parse_node_t), pointer :: pn_literal => null ()
     type(parse_node_t), pointer :: pn_beg => null ()
     type(parse_node_t), pointer :: pn_end => null ()
     type(parse_node_t), pointer :: pn_step => null ()
     type(eval_tree_t) :: expr_beg
     type(eval_tree_t) :: expr_end
     type(eval_tree_t) :: expr_step
     integer :: step_mode = 0
     integer :: n_step = 0
   contains
   \LA{}Commands: range: TBP\RA{}
  end type range_t
     
\eatline
\nwendcode{}\nwbegindocs{2909}These are the implementations:
\nwenddocs{}\nwbegincode{2910}\moddef{Commands: types}\plusendmoddef
  type, extends (range_t) :: range_int_t
     integer :: i_beg = 0
     integer :: i_end = 0
     integer :: i_step = 0
   contains
   \LA{}Commands: range int: TBP\RA{}
end type range_int_t
     
  type, extends (range_t) :: range_real_t
     real(default) :: r_beg = 0
     real(default) :: r_end = 0
     real(default) :: r_step = 0
     real(default) :: lr_beg  = 0
     real(default) :: lr_end  = 0
     real(default) :: lr_step = 0
   contains
   \LA{}Commands: range real: TBP\RA{}
end type range_real_t
     
\eatline
\nwendcode{}\nwbegindocs{2911}Finalize the allocated dummy node.  The other nodes are just pointers.
\nwenddocs{}\nwbegincode{2912}\moddef{Commands: range: TBP}\endmoddef
  procedure :: final => range_final
\nwendcode{}\nwbegincode{2913}\moddef{Commands: procedures}\plusendmoddef
  subroutine range_final (object)
    class(range_t), intent(inout) :: object
    if (associated (object%pn_expr)) then
       call parse_node_final (object%pn_expr, recursive = .false.)
       call parse_node_final (object%pn_term, recursive = .false.)
       call parse_node_final (object%pn_factor, recursive = .false.)
       call parse_node_final (object%pn_value, recursive = .false.)
       call parse_node_final (object%pn_literal, recursive = .false.)
       deallocate (object%pn_expr)
       deallocate (object%pn_term)
       deallocate (object%pn_factor)
       deallocate (object%pn_value)
       deallocate (object%pn_literal)
    end if
  end subroutine range_final
  
\eatline
\nwendcode{}\nwbegindocs{2914}Output.
\nwenddocs{}\nwbegincode{2915}\moddef{Commands: range: TBP}\plusendmoddef
  procedure (range_write), deferred :: write
  procedure :: base_write => range_write
\nwendcode{}\nwbegincode{2916}\moddef{Commands: range int: TBP}\endmoddef
  procedure :: write => range_int_write
\nwendcode{}\nwbegincode{2917}\moddef{Commands: range real: TBP}\endmoddef
  procedure :: write => range_real_write
\nwendcode{}\nwbegincode{2918}\moddef{Commands: procedures}\plusendmoddef
  subroutine range_write (object, unit)
    class(range_t), intent(in) :: object
    integer, intent(in), optional :: unit
    integer :: u
    u = given_output_unit (unit)
    write (u, "(1x,A)")  "Range specification:"
    if (associated (object%pn_expr)) then
       write (u, "(1x,A)")  "Dummy value:"
       call parse_node_write_rec (object%pn_expr, u)
    end if
    if (associated (object%pn_beg)) then
       write (u, "(1x,A)")  "Initial value:"
       call parse_node_write_rec (object%pn_beg, u)
       call object%expr_beg%write (u)
       if (associated (object%pn_end)) then
          write (u, "(1x,A)")  "Final value:"
          call parse_node_write_rec (object%pn_end, u)
          call object%expr_end%write (u)
          if (associated (object%pn_step)) then
             write (u, "(1x,A)")  "Step value:"
             call parse_node_write_rec (object%pn_step, u)
             select case (object%step_mode)
             case (STEP_ADD);   write (u, "(1x,A)")  "Step mode: +"
             case (STEP_SUB);   write (u, "(1x,A)")  "Step mode: -"
             case (STEP_MUL);   write (u, "(1x,A)")  "Step mode: *"
             case (STEP_DIV);   write (u, "(1x,A)")  "Step mode: /"
             case (STEP_COMP_ADD);  write (u, "(1x,A)")  "Division mode: +"
             case (STEP_COMP_MUL);  write (u, "(1x,A)")  "Division mode: *"
             end select
          end if
       end if
    else
       write (u, "(1x,A)")  "Expressions: [undefined]"
    end if
  end subroutine range_write

  subroutine range_int_write (object, unit)
    class(range_int_t), intent(in) :: object
    integer, intent(in), optional :: unit
    integer :: u
    u = given_output_unit (unit)
    call object%base_write (unit)
    write (u, "(1x,A)")  "Range parameters:"
    write (u, "(3x,A,I0)")  "i_beg  = ", object%i_beg
    write (u, "(3x,A,I0)")  "i_end  = ", object%i_end
    write (u, "(3x,A,I0)")  "i_step = ", object%i_step
    write (u, "(3x,A,I0)")  "n_step = ", object%n_step
  end subroutine range_int_write
  
  subroutine range_real_write (object, unit)
    class(range_real_t), intent(in) :: object
    integer, intent(in), optional :: unit
    integer :: u
    u = given_output_unit (unit)
    call object%base_write (unit)
    write (u, "(1x,A)")  "Range parameters:"
    write (u, "(3x,A," // FMT_19 // ")")  "r_beg  = ", object%r_beg
    write (u, "(3x,A," // FMT_19 // ")")  "r_end  = ", object%r_end
    write (u, "(3x,A," // FMT_19 // ")")  "r_step = ", object%r_end
    write (u, "(3x,A,I0)")  "n_step = ", object%n_step
  end subroutine range_real_write
  
\eatline
\nwendcode{}\nwbegindocs{2919}Initialize, given a range expression parse node.  This is common to the
implementations.
\nwenddocs{}\nwbegincode{2920}\moddef{Commands: range: TBP}\plusendmoddef
  procedure :: init => range_init
\nwendcode{}\nwbegincode{2921}\moddef{Commands: procedures}\plusendmoddef
  subroutine range_init (range, pn)
    class(range_t), intent(out) :: range
    type(parse_node_t), intent(in), target :: pn
    type(parse_node_t), pointer :: pn_spec, pn_end, pn_step_spec, pn_op
    select case (char (parse_node_get_rule_key (pn)))
    case ("expr")
    case ("range_expr")
       range%pn_beg => parse_node_get_sub_ptr (pn)
       pn_spec => parse_node_get_next_ptr (range%pn_beg)
       if (associated (pn_spec)) then
          pn_end => parse_node_get_sub_ptr (pn_spec, 2)
          range%pn_end => pn_end
          pn_step_spec => parse_node_get_next_ptr (pn_end)
          if (associated (pn_step_spec)) then
             pn_op => parse_node_get_sub_ptr (pn_step_spec)
             range%pn_step => parse_node_get_next_ptr (pn_op)
             select case (char (parse_node_get_rule_key (pn_op)))
             case ("/+");  range%step_mode = STEP_ADD
             case ("/-");  range%step_mode = STEP_SUB
             case ("/*");  range%step_mode = STEP_MUL
             case ("//");  range%step_mode = STEP_DIV
             case ("/+/");  range%step_mode = STEP_COMP_ADD
             case ("/*/");  range%step_mode = STEP_COMP_MUL
             case default
                call range%write ()
                call msg_bug ("Range: step mode not implemented")
             end select
          else
             range%step_mode = STEP_ADD
          end if
       else
          range%step_mode = STEP_NONE
       end if
       call range%create_value_node ()
    case default
       call msg_bug ("range expression: node type '" &
            // char (parse_node_get_rule_key (pn)) &
            // "' not implemented")
    end select
  end subroutine range_init
  
\eatline
\nwendcode{}\nwbegindocs{2922}This method manually creates a parse node (actually, a cascade of parse
nodes) that hold a constant value as a literal.  The idea is that this node is
inserted as the right-hand side of a fake variable assignment, which is
prepended to each scan iteration.  Before the variable
assignment is compiled and executed, we can manually reset the value of the
literal and thus pretend that the loop variable is assigned this value.
\nwenddocs{}\nwbegincode{2923}\moddef{Commands: range: TBP}\plusendmoddef
  procedure :: create_value_node => range_create_value_node
\nwendcode{}\nwbegincode{2924}\moddef{Commands: procedures}\plusendmoddef
  subroutine range_create_value_node (range)
    class(range_t), intent(inout) :: range
    allocate (range%pn_literal)
    allocate (range%pn_value)
    select type (range)
    type is (range_int_t)
       call parse_node_create_value (range%pn_literal, &
            syntax_get_rule_ptr (syntax_cmd_list, var_str ("integer_literal")),&
            ival = 0)
       call parse_node_create_branch (range%pn_value, &
            syntax_get_rule_ptr (syntax_cmd_list, var_str ("integer_value")))
    type is (range_real_t)
       call parse_node_create_value (range%pn_literal, &
            syntax_get_rule_ptr (syntax_cmd_list, var_str ("real_literal")),&
            rval = 0._default)
       call parse_node_create_branch (range%pn_value, &
            syntax_get_rule_ptr (syntax_cmd_list, var_str ("real_value")))
    class default
       call msg_bug ("range: create value node: type not implemented")
    end select
    call parse_node_append_sub (range%pn_value, range%pn_literal)
    call parse_node_freeze_branch (range%pn_value)
    allocate (range%pn_factor)
    call parse_node_create_branch (range%pn_factor, &
         syntax_get_rule_ptr (syntax_cmd_list, var_str ("factor")))
    call parse_node_append_sub (range%pn_factor, range%pn_value)
    call parse_node_freeze_branch (range%pn_factor)
    allocate (range%pn_term)
    call parse_node_create_branch (range%pn_term, &
         syntax_get_rule_ptr (syntax_cmd_list, var_str ("term")))
    call parse_node_append_sub (range%pn_term, range%pn_factor)
    call parse_node_freeze_branch (range%pn_term)
    allocate (range%pn_expr)
    call parse_node_create_branch (range%pn_expr, &
         syntax_get_rule_ptr (syntax_cmd_list, var_str ("expr")))
    call parse_node_append_sub (range%pn_expr, range%pn_term)
    call parse_node_freeze_branch (range%pn_expr)
  end subroutine range_create_value_node
  
\eatline
\nwendcode{}\nwbegindocs{2925}Compile, given an environment.
\nwenddocs{}\nwbegincode{2926}\moddef{Commands: range: TBP}\plusendmoddef
  procedure :: compile => range_compile
\nwendcode{}\nwbegincode{2927}\moddef{Commands: procedures}\plusendmoddef
  subroutine range_compile (range, global)
    class(range_t), intent(inout) :: range
    type(rt_data_t), intent(in), target :: global
    type(var_list_t), pointer :: var_list
    var_list => global%get_var_list_ptr ()
    if (associated (range%pn_beg)) then
       call range%expr_beg%init_expr (range%pn_beg, var_list)
       if (associated (range%pn_end)) then
          call range%expr_end%init_expr (range%pn_end, var_list)
          if (associated (range%pn_step)) then
             call range%expr_step%init_expr (range%pn_step, var_list)
          end if
       end if
    end if
  end subroutine range_compile
  
\eatline
\nwendcode{}\nwbegindocs{2928}Evaluate: compute the actual bounds and parameters that determine the values
that we can iterate.

This is implementation-specific.
\nwenddocs{}\nwbegincode{2929}\moddef{Commands: range: TBP}\plusendmoddef
  procedure (range_evaluate), deferred :: evaluate
\nwendcode{}\nwbegincode{2930}\moddef{Commands: interfaces}\plusendmoddef
  abstract interface
     subroutine range_evaluate (range)
       import
       class(range_t), intent(inout) :: range
     end subroutine range_evaluate
  end interface
  
\eatline
\nwendcode{}\nwbegindocs{2931}The version for an integer variable.  If the step is subtractive, we invert
the sign and treat it as an additive step.  For a multiplicative step, the
step must be greater than one, and the initial and final values must be of
same sign and strictly ordered.  Analogously for a division step.
\nwenddocs{}\nwbegincode{2932}\moddef{Commands: range int: TBP}\plusendmoddef
  procedure :: evaluate => range_int_evaluate
\nwendcode{}\nwbegincode{2933}\moddef{Commands: procedures}\plusendmoddef
  subroutine range_int_evaluate (range)
    class(range_int_t), intent(inout) :: range
    integer :: ival
    if (associated (range%pn_beg)) then
       call range%expr_beg%evaluate ()
       if (range%expr_beg%is_known ()) then
          range%i_beg = range%expr_beg%get_int ()
       else
          call range%write ()
          call msg_fatal &
               ("Range expression: initial value evaluates to unknown")
       end if
       if (associated (range%pn_end)) then
          call range%expr_end%evaluate ()
          if (range%expr_end%is_known ()) then
             range%i_end = range%expr_end%get_int ()
             if (associated (range%pn_step)) then
                call range%expr_step%evaluate ()
                if (range%expr_step%is_known ()) then
                   range%i_step = range%expr_step%get_int ()
                   select case (range%step_mode)
                   case (STEP_SUB);  range%i_step = - range%i_step
                   end select
                else
                   call range%write ()
                   call msg_fatal &
                        ("Range expression: step value evaluates to unknown")
                end if
             else
                range%i_step = 1
             end if
          else
             call range%write ()
             call msg_fatal &
                  ("Range expression: final value evaluates to unknown")
          end if
       else
          range%i_end = range%i_beg
          range%i_step = 1
       end if
       select case (range%step_mode)
       case (STEP_NONE)
          range%n_step = 1
       case (STEP_ADD, STEP_SUB)
          if (range%i_step /= 0) then
             if (range%i_beg == range%i_end) then
                range%n_step = 1
             else if (sign (1, range%i_end - range%i_beg) &
                  == sign (1, range%i_step)) then
                range%n_step = (range%i_end - range%i_beg) / range%i_step + 1
             else
                range%n_step = 0
             end if
          else
             call msg_fatal ("range evaluation (add): step value is zero")
          end if
       case (STEP_MUL)
          if (range%i_step > 1) then
             if (range%i_beg == range%i_end) then
                range%n_step = 1
             else if (range%i_beg == 0) then
                call msg_fatal ("range evaluation (mul): initial value is zero")
             else if (sign (1, range%i_beg) == sign (1, range%i_end) &
                  .and. abs (range%i_beg) < abs (range%i_end)) then
                range%n_step = 0
                ival = range%i_beg
                do while (abs (ival) <= abs (range%i_end))
                   range%n_step = range%n_step + 1
                   ival = ival * range%i_step
                end do
             else
                range%n_step = 0
             end if
          else
             call msg_fatal &
                  ("range evaluation (mult): step value is one or less")
          end if
       case (STEP_DIV)
          if (range%i_step > 1) then
             if (range%i_beg == range%i_end) then
                range%n_step = 1
             else if (sign (1, range%i_beg) == sign (1, range%i_end) &
                  .and. abs (range%i_beg) > abs (range%i_end)) then
                range%n_step = 0
                ival = range%i_beg
                do while (abs (ival) >= abs (range%i_end))
                   range%n_step = range%n_step + 1
                   if (ival == 0)  exit
                   ival = ival / range%i_step
                end do
             else
                range%n_step = 0
             end if
          else
             call msg_fatal &
                  ("range evaluation (div): step value is one or less")
          end if
       case (STEP_COMP_ADD)
          call msg_fatal ("range evaluation: &
               &step mode /+/ not allowed for integer variable")
       case (STEP_COMP_MUL)
          call msg_fatal ("range evaluation: &
               &step mode /*/ not allowed for integer variable")
       case default
          call range%write ()
          call msg_bug ("range evaluation: step mode not implemented")
       end select
    end if
  end subroutine range_int_evaluate

\eatline
\nwendcode{}\nwbegindocs{2934}The version for a real variable.
\nwenddocs{}\nwbegincode{2935}\moddef{Commands: range real: TBP}\plusendmoddef
  procedure :: evaluate => range_real_evaluate
\nwendcode{}\nwbegincode{2936}\moddef{Commands: procedures}\plusendmoddef
  subroutine range_real_evaluate (range)
    class(range_real_t), intent(inout) :: range
    if (associated (range%pn_beg)) then
       call range%expr_beg%evaluate ()
       if (range%expr_beg%is_known ()) then
          range%r_beg = range%expr_beg%get_real ()
       else
          call range%write ()
          call msg_fatal &
               ("Range expression: initial value evaluates to unknown")
       end if
       if (associated (range%pn_end)) then
          call range%expr_end%evaluate ()
          if (range%expr_end%is_known ()) then
             range%r_end = range%expr_end%get_real ()
             if (associated (range%pn_step)) then
                if (range%expr_step%is_known ()) then
                   select case (range%step_mode)
                   case (STEP_ADD, STEP_SUB, STEP_MUL, STEP_DIV)
                      call range%expr_step%evaluate ()
                      range%r_step = range%expr_step%get_real ()
                      select case (range%step_mode)
                      case (STEP_SUB);  range%r_step = - range%r_step
                      end select
                   case (STEP_COMP_ADD, STEP_COMP_MUL)
                      range%n_step = &
                           max (range%expr_step%get_int (), 0)
                   end select
                else
                   call range%write ()
                   call msg_fatal &
                        ("Range expression: step value evaluates to unknown")
                end if
             else
                call range%write ()
                call msg_fatal &
                     ("Range expression (real): step value must be provided")
             end if
          else
             call range%write ()
             call msg_fatal &
                  ("Range expression: final value evaluates to unknown")
          end if
       else
          range%r_end = range%r_beg
          range%r_step = 1
       end if
       select case (range%step_mode)
       case (STEP_NONE)
          range%n_step = 1
       case (STEP_ADD, STEP_SUB)
          if (range%r_step /= 0) then
             if (sign (1._default, range%r_end - range%r_beg) &
                  == sign (1._default, range%r_step)) then
                range%n_step = &
                     nint ((range%r_end - range%r_beg) / range%r_step + 1)
             else
                range%n_step = 0
             end if
          else
             call msg_fatal ("range evaluation (add): step value is zero")
          end if
       case (STEP_MUL)
          if (range%r_step > 1) then
             if (range%r_beg == 0 .or. range%r_end == 0) then
                call msg_fatal ("range evaluation (mul): bound is zero")
             else if (sign (1._default, range%r_beg) &
                  == sign (1._default, range%r_end) &
                  .and. abs (range%r_beg) <= abs (range%r_end)) then
                range%lr_beg = log (abs (range%r_beg))
                range%lr_end = log (abs (range%r_end))
                range%lr_step = log (range%r_step)
                range%n_step = nint &
                     (abs ((range%lr_end - range%lr_beg) / range%lr_step) + 1)
             else
                range%n_step = 0
             end if
          else
             call msg_fatal &
                  ("range evaluation (mult): step value is one or less")
          end if
       case (STEP_DIV)
          if (range%r_step > 1) then
             if (range%r_beg == 0 .or. range%r_end == 0) then
                call msg_fatal ("range evaluation (div): bound is zero")
             else if (sign (1._default, range%r_beg) &
                  == sign (1._default, range%r_end) &
                  .and. abs (range%r_beg) >= abs (range%r_end)) then
                range%lr_beg = log (abs (range%r_beg))
                range%lr_end = log (abs (range%r_end))
                range%lr_step = -log (range%r_step)
                range%n_step = nint &
                     (abs ((range%lr_end - range%lr_beg) / range%lr_step) + 1)
             else
                range%n_step = 0
             end if
          else
             call msg_fatal &
                  ("range evaluation (mult): step value is one or less")
          end if
       case (STEP_COMP_ADD)
          ! Number of steps already known
       case (STEP_COMP_MUL)
          ! Number of steps already known
          if (range%r_beg == 0 .or. range%r_end == 0) then
             call msg_fatal ("range evaluation (mul): bound is zero")
          else if (sign (1._default, range%r_beg) &
               == sign (1._default, range%r_end)) then
             range%lr_beg = log (abs (range%r_beg))
             range%lr_end = log (abs (range%r_end))
          else
             range%n_step = 0
          end if
       case default
          call range%write ()
          call msg_bug ("range evaluation: step mode not implemented")
       end select
    end if
  end subroutine range_real_evaluate

\eatline
\nwendcode{}\nwbegindocs{2937}Return the number of iterations:
\nwenddocs{}\nwbegincode{2938}\moddef{Commands: range: TBP}\plusendmoddef
  procedure :: get_n_iterations => range_get_n_iterations
\nwendcode{}\nwbegincode{2939}\moddef{Commands: procedures}\plusendmoddef
  function range_get_n_iterations (range) result (n)
    class(range_t), intent(in) :: range
    integer :: n
    n = range%n_step
  end function range_get_n_iterations
  
\eatline
\nwendcode{}\nwbegindocs{2940}Compute the value for iteration {\tt{}i} and store it in the embedded token.
\nwenddocs{}\nwbegincode{2941}\moddef{Commands: range: TBP}\plusendmoddef
  procedure (range_set_value), deferred :: set_value
\nwendcode{}\nwbegincode{2942}\moddef{Commands: interfaces}\plusendmoddef
  abstract interface
     subroutine range_set_value (range, i)
       import
       class(range_t), intent(inout) :: range
       integer, intent(in) :: i
     end subroutine range_set_value
  end interface
  
\eatline
\nwendcode{}\nwbegindocs{2943}In the integer case, we compute the value directly for additive step.  For
multiplicative step, we perform a loop in the same way as above, where the
number of iteration was determined.
\nwenddocs{}\nwbegincode{2944}\moddef{Commands: range int: TBP}\plusendmoddef
  procedure :: set_value => range_int_set_value
\nwendcode{}\nwbegincode{2945}\moddef{Commands: procedures}\plusendmoddef
  subroutine range_int_set_value (range, i)
    class(range_int_t), intent(inout) :: range
    integer, intent(in) :: i
    integer :: k, ival
    select case (range%step_mode)
    case (STEP_NONE)
       ival = range%i_beg
    case (STEP_ADD, STEP_SUB)
       ival = range%i_beg + (i - 1) * range%i_step
    case (STEP_MUL)
       ival = range%i_beg
       do k = 1, i - 1
          ival = ival * range%i_step
       end do
    case (STEP_DIV)
       ival = range%i_beg
       do k = 1, i - 1
          ival = ival / range%i_step
       end do
    case default
       call range%write ()
       call msg_bug ("range iteration: step mode not implemented")
    end select
    call parse_node_set_value (range%pn_literal, ival = ival)
  end subroutine range_int_set_value
  
\eatline
\nwendcode{}\nwbegindocs{2946}In the integer case, we compute the value directly for additive step.  For
multiplicative step, we perform a loop in the same way as above, where the
number of iteration was determined.
\nwenddocs{}\nwbegincode{2947}\moddef{Commands: range real: TBP}\plusendmoddef
  procedure :: set_value => range_real_set_value
\nwendcode{}\nwbegincode{2948}\moddef{Commands: procedures}\plusendmoddef
  subroutine range_real_set_value (range, i)
    class(range_real_t), intent(inout) :: range
    integer, intent(in) :: i
    real(default) :: rval, x
    select case (range%step_mode)
    case (STEP_NONE)
       rval = range%r_beg
    case (STEP_ADD, STEP_SUB, STEP_COMP_ADD)
       if (range%n_step > 1) then
          x = real (i - 1, default) / (range%n_step - 1)
       else
          x = 1._default / 2
       end if
       rval = x * range%r_end + (1 - x) * range%r_beg
    case (STEP_MUL, STEP_DIV, STEP_COMP_MUL)
       if (range%n_step > 1) then
          x = real (i - 1, default) / (range%n_step - 1)
       else
          x = 1._default / 2
       end if
       rval = sign &
            (exp (x * range%lr_end + (1 - x) * range%lr_beg), range%r_beg)
    case default
       call range%write ()
       call msg_bug ("range iteration: step mode not implemented")
    end select
    call parse_node_set_value (range%pn_literal, rval = rval)
  end subroutine range_real_set_value
  
\eatline
\nwendcode{}\nwbegindocs{2949}\nwdocspar
\subsubsection{Scan over parameters and other objects}
The scan command allocates a new parse node for the variable
assignment (the lhs).  The rhs of this parse node is assigned from the
available rhs expressions in the scan list, one at a time, so the
compiled parse node can be prepended to the scan body.

Note: for the integer/real range array, the obvious implementation as a
polymorphic array is suspended because in gfortran 4.7, polymorphic arrays are
apparently broken.
\nwenddocs{}\nwbegincode{2950}\moddef{Commands: types}\plusendmoddef
  type, extends (command_t) :: cmd_scan_t
     private
     type(string_t) :: name
     integer :: n_values = 0
     type(parse_node_p), dimension(:), allocatable :: scan_cmd
     !!! !!! gfortran 4.7.x memory corruption
     !!!  class(range_t), dimension(:), allocatable :: range
     type(range_int_t), dimension(:), allocatable :: range_int
     type(range_real_t), dimension(:), allocatable :: range_real
   contains
   \LA{}Commands: cmd scan: TBP\RA{}  
  end type cmd_scan_t

\eatline
\nwendcode{}\nwbegindocs{2951}Finalizer.  

The auxiliary parse nodes that we have constructed have to be treated
carefully: the embedded pointers all point to persistent objects
somewhere else and should not be finalized, so we should not call the
finalizer recursively.
\nwenddocs{}\nwbegincode{2952}\moddef{Commands: cmd scan: TBP}\endmoddef
  procedure :: final => cmd_scan_final
\nwendcode{}\nwbegincode{2953}\moddef{Commands: procedures}\plusendmoddef
  recursive subroutine cmd_scan_final (cmd)
    class(cmd_scan_t), intent(inout) :: cmd
    type(parse_node_t), pointer :: pn_var_single, pn_decl_single
    type(string_t) :: key
    integer :: i
    if (allocated (cmd%scan_cmd)) then
       do i = 1, size (cmd%scan_cmd)
          pn_var_single => parse_node_get_sub_ptr (cmd%scan_cmd(i)%ptr)
          key = parse_node_get_rule_key (pn_var_single)
          select case (char (key))
          case ("scan_string_decl", "scan_log_decl")
             pn_decl_single => parse_node_get_sub_ptr (pn_var_single, 2)
             call parse_node_final (pn_decl_single, recursive=.false.)
             deallocate (pn_decl_single)
          end select
          call parse_node_final (pn_var_single, recursive=.false.)
          deallocate (pn_var_single)
       end do
       deallocate (cmd%scan_cmd)
    end if
    !!! !!! gfortran 4.7.x memory corruption
    !!!  if (allocated (cmd%range)) then
    !!!     do i = 1, size (cmd%range)
    !!!        call cmd%range(i)%final ()
    !!!     end do
    !!!  end if
    if (allocated (cmd%range_int)) then
       do i = 1, size (cmd%range_int)
          call cmd%range_int(i)%final ()
       end do
    end if
    if (allocated (cmd%range_real)) then
       do i = 1, size (cmd%range_real)
          call cmd%range_real(i)%final ()
       end do
    end if
  end subroutine cmd_scan_final

\eatline
\nwendcode{}\nwbegindocs{2954}Output.
\nwenddocs{}\nwbegincode{2955}\moddef{Commands: cmd scan: TBP}\plusendmoddef
  procedure :: write => cmd_scan_write
\nwendcode{}\nwbegincode{2956}\moddef{Commands: procedures}\plusendmoddef
  subroutine cmd_scan_write (cmd, unit, indent)
    class(cmd_scan_t), intent(in) :: cmd
    integer, intent(in), optional :: unit, indent
    integer :: u
    u = given_output_unit (unit);  if (u < 0)  return
    call write_indent (u, indent)
    write (u, "(1x,A,1x,A,1x,'(',I0,')')")  "scan:", char (cmd%name), &
         cmd%n_values
  end subroutine cmd_scan_write

\eatline
\nwendcode{}\nwbegindocs{2957}Compile the scan command.  We construct a new parse node that
implements the variable assignment for a single element on the rhs,
instead of the whole list that we get from the original parse tree.
By simply copying the node, we copy all pointers and inherit the
targets from the original.  During execution, we should replace the
rhs by the stored rhs pointers (the list elements), one by one, then
(re)compile the redefined node.
\nwenddocs{}\nwbegincode{2958}\moddef{Commands: cmd scan: TBP}\plusendmoddef
  procedure :: compile => cmd_scan_compile
\nwendcode{}\nwbegincode{2959}\moddef{Commands: procedures}\plusendmoddef
  recursive subroutine cmd_scan_compile (cmd, global)
    class(cmd_scan_t), intent(inout) :: cmd
    type(rt_data_t), intent(inout), target :: global
    type(var_list_t), pointer :: var_list
    type(parse_node_t), pointer :: pn_var, pn_body, pn_body_first
    type(parse_node_t), pointer :: pn_decl, pn_name
    type(parse_node_t), pointer :: pn_arg, pn_scan_cmd, pn_rhs
    type(parse_node_t), pointer :: pn_decl_single, pn_var_single
    type(syntax_rule_t), pointer :: var_rule_decl, var_rule
    type(string_t) :: key
    integer :: var_type
    integer :: i
    logical, parameter :: debug = .false.
    if (debug) then
       print *, "compile scan"
       call parse_node_write_rec (cmd%pn)
    end if
    pn_var => parse_node_get_sub_ptr (cmd%pn, 2)
    pn_body => parse_node_get_next_ptr (pn_var)
    if (associated (pn_body)) then
       pn_body_first => parse_node_get_sub_ptr (pn_body)
    else
       pn_body_first => null ()
    end if
    key = parse_node_get_rule_key (pn_var)
    select case (char (key))
    case ("scan_num")
       pn_name => parse_node_get_sub_ptr (pn_var)
       cmd%name = parse_node_get_string (pn_name)
       var_rule => syntax_get_rule_ptr (syntax_cmd_list, var_str ("cmd_num"))
       pn_arg => parse_node_get_next_ptr (pn_name, 2)
    case ("scan_int")
       pn_name => parse_node_get_sub_ptr (pn_var, 2)
       cmd%name = parse_node_get_string (pn_name)
       var_rule => syntax_get_rule_ptr (syntax_cmd_list, var_str ("cmd_int"))
       pn_arg => parse_node_get_next_ptr (pn_name, 2)
    case ("scan_real")
       pn_name => parse_node_get_sub_ptr (pn_var, 2)
       cmd%name = parse_node_get_string (pn_name)
       var_rule => syntax_get_rule_ptr (syntax_cmd_list, var_str ("cmd_real"))
       pn_arg => parse_node_get_next_ptr (pn_name, 2)
    case ("scan_complex")
       pn_name => parse_node_get_sub_ptr (pn_var, 2)
       cmd%name = parse_node_get_string (pn_name)
       var_rule => syntax_get_rule_ptr (syntax_cmd_list, var_str("cmd_complex"))
       pn_arg => parse_node_get_next_ptr (pn_name, 2)
    case ("scan_alias")
       pn_name => parse_node_get_sub_ptr (pn_var, 2)
       cmd%name = parse_node_get_string (pn_name)
       var_rule => syntax_get_rule_ptr (syntax_cmd_list, var_str ("cmd_alias"))
       pn_arg => parse_node_get_next_ptr (pn_name, 2)
    case ("scan_string_decl")
       pn_decl => parse_node_get_sub_ptr (pn_var, 2)
       pn_name => parse_node_get_sub_ptr (pn_decl, 2)
       cmd%name = parse_node_get_string (pn_name)
       var_rule_decl => syntax_get_rule_ptr (syntax_cmd_list, &
            var_str ("cmd_string"))
       var_rule => syntax_get_rule_ptr (syntax_cmd_list, &
            var_str ("cmd_string_decl"))
       pn_arg => parse_node_get_next_ptr (pn_name, 2)
    case ("scan_log_decl")
       pn_decl => parse_node_get_sub_ptr (pn_var, 2)
       pn_name => parse_node_get_sub_ptr (pn_decl, 2)
       cmd%name = parse_node_get_string (pn_name)
       var_rule_decl => syntax_get_rule_ptr (syntax_cmd_list, &
            var_str ("cmd_log"))
       var_rule => syntax_get_rule_ptr (syntax_cmd_list, &
            var_str ("cmd_log_decl"))
       pn_arg => parse_node_get_next_ptr (pn_name, 2)
    case ("scan_cuts")
       var_rule => syntax_get_rule_ptr (syntax_cmd_list, &
            var_str ("cmd_cuts"))
       cmd%name = "cuts"
       pn_arg => parse_node_get_sub_ptr (pn_var, 3)
    case ("scan_weight")
       var_rule => syntax_get_rule_ptr (syntax_cmd_list, &
            var_str ("cmd_weight"))
       cmd%name = "weight"
       pn_arg => parse_node_get_sub_ptr (pn_var, 3)
    case ("scan_scale")
       var_rule => syntax_get_rule_ptr (syntax_cmd_list, &
            var_str ("cmd_scale"))
       cmd%name = "scale"
       pn_arg => parse_node_get_sub_ptr (pn_var, 3)
    case ("scan_ren_scale")
       var_rule => syntax_get_rule_ptr (syntax_cmd_list, &
            var_str ("cmd_ren_scale"))
       cmd%name = "renormalization_scale"
       pn_arg => parse_node_get_sub_ptr (pn_var, 3)
    case ("scan_fac_scale")
       var_rule => syntax_get_rule_ptr (syntax_cmd_list, &
            var_str ("cmd_fac_scale"))
       cmd%name = "factorization_scale"
       pn_arg => parse_node_get_sub_ptr (pn_var, 3)
    case ("scan_selection")
       var_rule => syntax_get_rule_ptr (syntax_cmd_list, &
            var_str ("cmd_selection"))
       cmd%name = "selection"
       pn_arg => parse_node_get_sub_ptr (pn_var, 3)
    case ("scan_reweight")
       var_rule => syntax_get_rule_ptr (syntax_cmd_list, &
            var_str ("cmd_reweight"))
       cmd%name = "reweight"
       pn_arg => parse_node_get_sub_ptr (pn_var, 3)
    case ("scan_analysis")
       var_rule => syntax_get_rule_ptr (syntax_cmd_list, &
            var_str ("cmd_analysis"))
       cmd%name = "analysis"
       pn_arg => parse_node_get_sub_ptr (pn_var, 3)
    case ("scan_model")
       var_rule => syntax_get_rule_ptr (syntax_cmd_list, &
            var_str ("cmd_model"))
       cmd%name = "model"
       pn_arg => parse_node_get_sub_ptr (pn_var, 3)
    case ("scan_library")
       var_rule => syntax_get_rule_ptr (syntax_cmd_list, &
            var_str ("cmd_library"))
       cmd%name = "library"
       pn_arg => parse_node_get_sub_ptr (pn_var, 3)
    case default
       call msg_bug ("scan: case '" // char (key) // "' not implemented")
    end select
    if (associated (pn_arg)) then
       cmd%n_values = parse_node_get_n_sub (pn_arg)
    end if
    var_list => global%get_var_list_ptr ()
    allocate (cmd%scan_cmd (cmd%n_values))
    select case (char (key))
    case ("scan_num")
       var_type = &
            var_list%get_type (cmd%name)
       select case (var_type)
       case (V_INT)
          !!! !!! gfortran 4.7.x memory corruption
          !!!  allocate (range_int_t :: cmd%range (cmd%n_values))
          allocate (cmd%range_int (cmd%n_values))
       case (V_REAL)
          !!! !!! gfortran 4.7.x memory corruption          
          !!!  allocate (range_real_t :: cmd%range (cmd%n_values))
          allocate (cmd%range_real (cmd%n_values))
       case (V_CMPLX)
          call msg_fatal ("scan over complex variable not implemented")
       case (V_NONE)
          call msg_fatal ("scan: variable '" // char (cmd%name) //"' undefined")
       case default
          call msg_bug ("scan: impossible variable type")
       end select
    case ("scan_int")
       !!! !!! gfortran 4.7.x memory corruption       
       !!!  allocate (range_int_t :: cmd%range (cmd%n_values))
       allocate (cmd%range_int (cmd%n_values))
    case ("scan_real")
       !!! !!! gfortran 4.7.x memory corruption
       !!!  allocate (range_real_t :: cmd%range (cmd%n_values))
       allocate (cmd%range_real (cmd%n_values))
    case ("scan_complex")
       call msg_fatal ("scan over complex variable not implemented")
    end select
    i = 1
    if (associated (pn_arg)) then
       pn_rhs => parse_node_get_sub_ptr (pn_arg)
    else
       pn_rhs => null ()
    end if
    do while (associated (pn_rhs))
       allocate (pn_scan_cmd)
       call parse_node_create_branch (pn_scan_cmd, &
            syntax_get_rule_ptr (syntax_cmd_list, var_str ("command_list")))
       allocate (pn_var_single)
       pn_var_single = pn_var
       call parse_node_replace_rule (pn_var_single, var_rule)
       select case (char (key))
       case ("scan_num", "scan_int", "scan_real", &
            "scan_complex", "scan_alias", &
            "scan_cuts", "scan_weight", &
            "scan_scale", "scan_ren_scale", "scan_fac_scale", &
            "scan_selection", "scan_reweight", "scan_analysis", &
            "scan_model", "scan_library")
          if (allocated (cmd%range_int)) then
             call cmd%range_int(i)%init (pn_rhs)
             !!! !!! gfortran 4.7.x memory corruption             
             !!!  call cmd%range_int(i)%compile (global)
             call parse_node_replace_last_sub &
                  (pn_var_single, cmd%range_int(i)%pn_expr)
          else if (allocated (cmd%range_real)) then
             call cmd%range_real(i)%init (pn_rhs)
             !!! !!! gfortran 4.7.x memory corruption 
             !!!  call cmd%range_real(i)%compile (global)
             call parse_node_replace_last_sub &
                  (pn_var_single, cmd%range_real(i)%pn_expr)
          else
             call parse_node_replace_last_sub (pn_var_single, pn_rhs)
          end if
       case ("scan_string_decl", "scan_log_decl")
          allocate (pn_decl_single)
          pn_decl_single = pn_decl
          call parse_node_replace_rule (pn_decl_single, var_rule_decl)
          call parse_node_replace_last_sub (pn_decl_single, pn_rhs)
          call parse_node_freeze_branch (pn_decl_single)
          call parse_node_replace_last_sub (pn_var_single, pn_decl_single)
       case default
          call msg_bug ("scan: case '" // char (key)  &
               // "' broken")
       end select
       call parse_node_freeze_branch (pn_var_single)
       call parse_node_append_sub (pn_scan_cmd, pn_var_single)
       call parse_node_append_sub (pn_scan_cmd, pn_body_first)
       call parse_node_freeze_branch (pn_scan_cmd)
       cmd%scan_cmd(i)%ptr => pn_scan_cmd
       i = i + 1
       pn_rhs => parse_node_get_next_ptr (pn_rhs)
    end do
    if (debug) then
       do i = 1, cmd%n_values
          print *, "scan command ", i
          call parse_node_write_rec (cmd%scan_cmd(i)%ptr)
          if (allocated (cmd%range_int))  call cmd%range_int(i)%write ()
          if (allocated (cmd%range_real))  call cmd%range_real(i)%write ()
       end do
       print *, "original"
       call parse_node_write_rec (cmd%pn)
    end if
  end subroutine cmd_scan_compile

\eatline
\nwendcode{}\nwbegindocs{2960}Execute the loop for all values in the step list.  We use the
parse trees with single variable assignment that we have stored, to
iteratively create a local environment, execute the stored commands, and
destroy it again.  When we encounter a range object, we execute the commands
for each value that this object provides.  Computing this value has the side
effect of modifying the rhs of the variable assignment that heads the local
command list, directly in the local parse tree.
\nwenddocs{}\nwbegincode{2961}\moddef{Commands: cmd scan: TBP}\plusendmoddef
  procedure :: execute => cmd_scan_execute
\nwendcode{}\nwbegincode{2962}\moddef{Commands: procedures}\plusendmoddef
  recursive subroutine cmd_scan_execute (cmd, global)
    class(cmd_scan_t), intent(inout) :: cmd
    type(rt_data_t), intent(inout), target :: global
    type(rt_data_t), allocatable :: local
    integer :: i, j
    do i = 1, cmd%n_values
       if (allocated (cmd%range_int)) then
          call cmd%range_int(i)%compile (global)
          call cmd%range_int(i)%evaluate ()
          do j = 1, cmd%range_int(i)%get_n_iterations ()
             call cmd%range_int(i)%set_value (j)
             allocate (local)
             call build_alt_setup (local, global, cmd%scan_cmd(i)%ptr)
             call local%local_final ()
             deallocate (local)
          end do
       else if (allocated (cmd%range_real)) then
          call cmd%range_real(i)%compile (global)
          call cmd%range_real(i)%evaluate ()
          do j = 1, cmd%range_real(i)%get_n_iterations ()
             call cmd%range_real(i)%set_value (j)
             allocate (local)
             call build_alt_setup (local, global, cmd%scan_cmd(i)%ptr)
             call local%local_final ()
             deallocate (local)
          end do
       else
          allocate (local)
          call build_alt_setup (local, global, cmd%scan_cmd(i)%ptr)
          call local%local_final ()
          deallocate (local)
       end if
    end do
  end subroutine cmd_scan_execute

\eatline
\nwendcode{}\nwbegindocs{2963}\nwdocspar
\subsubsection{Conditionals}
Conditionals are implemented as a list that is compiled and evaluated
recursively; this allows for a straightforward representation of
{\tt{}else\ if} constructs.  A {\tt{}cmd{\char95}if{\char95}t} object can hold either an
{\tt{}else{\char95}if} clause which is another object of this type, or an
{\tt{}else{\char95}body}, but not both.

If- or else-bodies are no scoping units, so all data remain global and
no copy-in copy-out is needed.
\nwenddocs{}\nwbegincode{2964}\moddef{Commands: types}\plusendmoddef
  type, extends (command_t) :: cmd_if_t
     private
     type(parse_node_t), pointer :: pn_if_lexpr => null ()
     type(command_list_t), pointer :: if_body => null ()
     type(cmd_if_t), dimension(:), pointer :: elsif_cmd => null ()
     type(command_list_t), pointer :: else_body => null ()
   contains
   \LA{}Commands: cmd if: TBP\RA{}
  end type cmd_if_t

\eatline
\nwendcode{}\nwbegindocs{2965}Finalizer.  There are no local options, therefore we can simply override
the default finalizer.
\nwenddocs{}\nwbegincode{2966}\moddef{Commands: cmd if: TBP}\endmoddef
  procedure :: final => cmd_if_final
\nwendcode{}\nwbegincode{2967}\moddef{Commands: procedures}\plusendmoddef
  recursive subroutine cmd_if_final (cmd)
    class(cmd_if_t), intent(inout) :: cmd
    integer :: i
    if (associated (cmd%if_body)) then
       call command_list_final (cmd%if_body)
       deallocate (cmd%if_body)
    end if
    if (associated (cmd%elsif_cmd)) then
       do i = 1, size (cmd%elsif_cmd)
          call cmd_if_final (cmd%elsif_cmd(i))
       end do
       deallocate (cmd%elsif_cmd)
    end if
    if (associated (cmd%else_body)) then
       call command_list_final (cmd%else_body)
       deallocate (cmd%else_body)
    end if
  end subroutine cmd_if_final

\eatline
\nwendcode{}\nwbegindocs{2968}Output.  Recursively write the command lists.
\nwenddocs{}\nwbegincode{2969}\moddef{Commands: cmd if: TBP}\plusendmoddef
  procedure :: write => cmd_if_write
\nwendcode{}\nwbegincode{2970}\moddef{Commands: procedures}\plusendmoddef
  subroutine cmd_if_write (cmd, unit, indent)
    class(cmd_if_t), intent(in) :: cmd
    integer, intent(in), optional :: unit, indent
    integer :: u, ind, i
    u = given_output_unit (unit);  if (u < 0)  return
    ind = 0;  if (present (indent))  ind = indent
    call write_indent (u, indent)
    write (u, "(A)")  "if <expr> then"
    if (associated (cmd%if_body)) then
       call cmd%if_body%write (unit, ind + 1)
    end if
    if (associated (cmd%elsif_cmd)) then
       do i = 1, size (cmd%elsif_cmd)
          call write_indent (u, indent)
          write (u, "(A)")  "elsif <expr> then"
          if (associated (cmd%elsif_cmd(i)%if_body)) then
             call cmd%elsif_cmd(i)%if_body%write (unit, ind + 1)
          end if
       end do
    end if
    if (associated (cmd%else_body)) then
       call write_indent (u, indent)
       write (u, "(A)")  "else"
       call cmd%else_body%write (unit, ind + 1)
    end if
  end subroutine cmd_if_write
  
\eatline
\nwendcode{}\nwbegindocs{2971}Compile the conditional.
\nwenddocs{}\nwbegincode{2972}\moddef{Commands: cmd if: TBP}\plusendmoddef
  procedure :: compile => cmd_if_compile
\nwendcode{}\nwbegincode{2973}\moddef{Commands: procedures}\plusendmoddef
  recursive subroutine cmd_if_compile (cmd, global)
    class(cmd_if_t), intent(inout) :: cmd
    type(rt_data_t), intent(inout), target :: global
    type(parse_node_t), pointer :: pn_lexpr, pn_body
    type(parse_node_t), pointer :: pn_elsif_clauses, pn_cmd_elsif
    type(parse_node_t), pointer :: pn_else_clause, pn_cmd_else
    integer :: i, n_elsif
    pn_lexpr => parse_node_get_sub_ptr (cmd%pn, 2)
    cmd%pn_if_lexpr => pn_lexpr
    pn_body => parse_node_get_next_ptr (pn_lexpr, 2)
    select case (char (parse_node_get_rule_key (pn_body)))
    case ("command_list")
       allocate (cmd%if_body)
       call cmd%if_body%compile (pn_body, global)
       pn_elsif_clauses => parse_node_get_next_ptr (pn_body)
    case default
       pn_elsif_clauses => pn_body
    end select
    select case (char (parse_node_get_rule_key (pn_elsif_clauses)))
    case ("elsif_clauses")
       n_elsif = parse_node_get_n_sub (pn_elsif_clauses)
       allocate (cmd%elsif_cmd (n_elsif))
       pn_cmd_elsif => parse_node_get_sub_ptr (pn_elsif_clauses)
       do i = 1, n_elsif
          pn_lexpr => parse_node_get_sub_ptr (pn_cmd_elsif, 2)
          cmd%elsif_cmd(i)%pn_if_lexpr => pn_lexpr
          pn_body => parse_node_get_next_ptr (pn_lexpr, 2)
          if (associated (pn_body)) then
             allocate (cmd%elsif_cmd(i)%if_body)
             call cmd%elsif_cmd(i)%if_body%compile (pn_body, global)
          end if
          pn_cmd_elsif => parse_node_get_next_ptr (pn_cmd_elsif)
       end do
       pn_else_clause => parse_node_get_next_ptr (pn_elsif_clauses)
    case default
       pn_else_clause => pn_elsif_clauses
    end select
    select case (char (parse_node_get_rule_key (pn_else_clause)))
    case ("else_clause")
       pn_cmd_else => parse_node_get_sub_ptr (pn_else_clause)
       pn_body => parse_node_get_sub_ptr (pn_cmd_else, 2)
       if (associated (pn_body)) then
          allocate (cmd%else_body)
          call cmd%else_body%compile (pn_body, global)
       end if
    end select
  end subroutine cmd_if_compile

\eatline
\nwendcode{}\nwbegindocs{2974}(Recursively) execute the condition.  Context remains global in all cases.
\nwenddocs{}\nwbegincode{2975}\moddef{Commands: cmd if: TBP}\plusendmoddef
  procedure :: execute => cmd_if_execute
\nwendcode{}\nwbegincode{2976}\moddef{Commands: procedures}\plusendmoddef
  recursive subroutine cmd_if_execute (cmd, global)
    class(cmd_if_t), intent(inout) :: cmd
    type(rt_data_t), intent(inout), target :: global
    type(var_list_t), pointer :: var_list
    logical :: lval, is_known
    integer :: i
    var_list => global%get_var_list_ptr ()
    lval = eval_log (cmd%pn_if_lexpr, var_list, is_known=is_known)
    if (is_known) then
       if (lval) then
          if (associated (cmd%if_body)) then
             call cmd%if_body%execute (global)
          end if
          return
       end if
    else
       call error_undecided ()
       return
    end if
    if (associated (cmd%elsif_cmd)) then
       SCAN_ELSIF: do i = 1, size (cmd%elsif_cmd)
          lval = eval_log (cmd%elsif_cmd(i)%pn_if_lexpr, var_list, &
                is_known=is_known)
          if (is_known) then
             if (lval) then
                if (associated (cmd%elsif_cmd(i)%if_body)) then
                   call cmd%elsif_cmd(i)%if_body%execute (global)
                end if
                return
             end if
          else
             call error_undecided ()
             return
          end if
       end do SCAN_ELSIF
    end if
    if (associated (cmd%else_body)) then
       call cmd%else_body%execute (global)
    end if
  contains
    subroutine error_undecided ()
      call msg_error ("Undefined result of cmditional expression: " &
           // "neither branch will be executed")
    end subroutine error_undecided
  end subroutine cmd_if_execute

\eatline
\nwendcode{}\nwbegindocs{2977}\nwdocspar
\subsubsection{Include another command-list file}
The include command allocates a local parse tree.  This must not be
deleted before the command object itself is deleted, since pointers
may point to subobjects of it.
\nwenddocs{}\nwbegincode{2978}\moddef{Commands: types}\plusendmoddef
  type, extends (command_t) :: cmd_include_t
     private
     type(string_t) :: file
     type(command_list_t), pointer :: command_list => null ()
     type(parse_tree_t) :: parse_tree
   contains
   \LA{}Commands: cmd include: TBP\RA{}
  end type cmd_include_t

\eatline
\nwendcode{}\nwbegindocs{2979}Finalizer: delete the command list.  No options, so we can simply override
the default finalizer.
\nwenddocs{}\nwbegincode{2980}\moddef{Commands: cmd include: TBP}\endmoddef
  procedure :: final => cmd_include_final
\nwendcode{}\nwbegincode{2981}\moddef{Commands: procedures}\plusendmoddef
  subroutine cmd_include_final (cmd)
    class(cmd_include_t), intent(inout) :: cmd
    call parse_tree_final (cmd%parse_tree)
    if (associated (cmd%command_list)) then
       call cmd%command_list%final ()
       deallocate (cmd%command_list)
    end if
  end subroutine cmd_include_final

\eatline
\nwendcode{}\nwbegindocs{2982}Write: display the command list as-is, if allocated.
\nwenddocs{}\nwbegincode{2983}\moddef{Commands: cmd include: TBP}\plusendmoddef
  procedure :: write => cmd_include_write
\nwendcode{}\nwbegincode{2984}\moddef{Commands: procedures}\plusendmoddef
  subroutine cmd_include_write (cmd, unit, indent)
    class(cmd_include_t), intent(in) :: cmd
    integer, intent(in), optional :: unit, indent
    integer :: u, ind
    u = given_output_unit (unit)
    ind = 0;  if (present (indent))  ind = indent
    call write_indent (u, indent)
    write (u, "(A,A,A,A)")  "include ", '"', char (cmd%file), '"'
    if (associated (cmd%command_list)) then
       call cmd%command_list%write (u, ind + 1)
    end if
  end subroutine cmd_include_write
  
\eatline
\nwendcode{}\nwbegindocs{2985}Compile file contents: First parse the file, then immediately
compile its contents.  Use the global data set.
\nwenddocs{}\nwbegincode{2986}\moddef{Commands: cmd include: TBP}\plusendmoddef
  procedure :: compile => cmd_include_compile
\nwendcode{}\nwbegincode{2987}\moddef{Commands: procedures}\plusendmoddef
  subroutine cmd_include_compile (cmd, global)
    class(cmd_include_t), intent(inout) :: cmd
    type(rt_data_t), intent(inout), target :: global
    type(parse_node_t), pointer :: pn_arg, pn_file
    type(string_t) :: file
    logical :: exist
    integer :: u
    type(stream_t), target :: stream
    type(lexer_t) :: lexer
    pn_arg => parse_node_get_sub_ptr (cmd%pn, 2)
    pn_file => parse_node_get_sub_ptr (pn_arg)
    file = parse_node_get_string (pn_file)
    inquire (file=char(file), exist=exist)
    if (exist) then
       cmd%file = file
    else
       cmd%file = global%os_data%whizard_cutspath // "/" // file
       inquire (file=char(cmd%file), exist=exist)
       if (.not. exist) then
          call msg_error ("Include file '" // char (file) // "' not found")
          return
       end if
    end if
    u = free_unit ()
    call lexer_init_cmd_list (lexer, global%lexer)
    call stream_init (stream, char (cmd%file))
    call lexer_assign_stream (lexer, stream)
    call parse_tree_init (cmd%parse_tree, syntax_cmd_list, lexer)
    call stream_final (stream)
    call lexer_final (lexer)
    close (u)
    allocate (cmd%command_list)
    call cmd%command_list%compile (parse_tree_get_root_ptr (cmd%parse_tree), &
         global)
  end subroutine cmd_include_compile

\eatline
\nwendcode{}\nwbegindocs{2988}Execute file contents in the global context.
\nwenddocs{}\nwbegincode{2989}\moddef{Commands: cmd include: TBP}\plusendmoddef
  procedure :: execute => cmd_include_execute
\nwendcode{}\nwbegincode{2990}\moddef{Commands: procedures}\plusendmoddef
  subroutine cmd_include_execute (cmd, global)
    class(cmd_include_t), intent(inout) :: cmd
    type(rt_data_t), intent(inout), target :: global
    if (associated (cmd%command_list)) then
       call msg_message &
            ("Including Sindarin from '" // char (cmd%file) // "'")
       call cmd%command_list%execute (global)
       call msg_message &
            ("End of included '" // char (cmd%file) // "'")
    end if
  end subroutine cmd_include_execute

\eatline
\nwendcode{}\nwbegindocs{2991}\nwdocspar
\subsubsection{Quit command execution}
The code is the return code of the whole program if it is terminated
by this command.
\nwenddocs{}\nwbegincode{2992}\moddef{Commands: types}\plusendmoddef
  type, extends (command_t) :: cmd_quit_t
     private
     logical :: has_code = .false.
     type(parse_node_t), pointer :: pn_code_expr => null ()
   contains
   \LA{}Commands: cmd quit: TBP\RA{}
  end type cmd_quit_t

\eatline
\nwendcode{}\nwbegindocs{2993}Output.
\nwenddocs{}\nwbegincode{2994}\moddef{Commands: cmd quit: TBP}\endmoddef
  procedure :: write => cmd_quit_write
\nwendcode{}\nwbegincode{2995}\moddef{Commands: procedures}\plusendmoddef
  subroutine cmd_quit_write (cmd, unit, indent)
    class(cmd_quit_t), intent(in) :: cmd
    integer, intent(in), optional :: unit, indent
    integer :: u
    u = given_output_unit (unit);  if (u < 0)  return
    call write_indent (u, indent)
    write (u, "(1x,A,L1)")  "quit: has_code = ", cmd%has_code
  end subroutine cmd_quit_write

\eatline
\nwendcode{}\nwbegindocs{2996}Compile: allocate a {\tt{}quit} object which serves as a placeholder.
\nwenddocs{}\nwbegincode{2997}\moddef{Commands: cmd quit: TBP}\plusendmoddef
  procedure :: compile => cmd_quit_compile
\nwendcode{}\nwbegincode{2998}\moddef{Commands: procedures}\plusendmoddef
  subroutine cmd_quit_compile (cmd, global)
    class(cmd_quit_t), intent(inout) :: cmd
    type(rt_data_t), intent(inout), target :: global
    type(parse_node_t), pointer :: pn_arg
    pn_arg => parse_node_get_sub_ptr (cmd%pn, 2)
    if (associated (pn_arg)) then
       cmd%pn_code_expr => parse_node_get_sub_ptr (pn_arg)
       cmd%has_code = .true.
    end if
  end subroutine cmd_quit_compile

\eatline
\nwendcode{}\nwbegindocs{2999}Execute: The quit command does not execute anything, it just stops
command execution.  This is achieved by setting quit flag and quit
code in the global variable list.  However, the return code, if
present, is an expression which has to be evaluated.
\nwenddocs{}\nwbegincode{3000}\moddef{Commands: cmd quit: TBP}\plusendmoddef
  procedure :: execute => cmd_quit_execute
\nwendcode{}\nwbegincode{3001}\moddef{Commands: procedures}\plusendmoddef
  subroutine cmd_quit_execute (cmd, global)
    class(cmd_quit_t), intent(inout) :: cmd
    type(rt_data_t), intent(inout), target :: global
    type(var_list_t), pointer :: var_list
    logical :: is_known
    var_list => global%get_var_list_ptr ()
    if (cmd%has_code) then
       global%quit_code = eval_int (cmd%pn_code_expr, var_list, &
            is_known=is_known)
       if (.not. is_known) then
          call msg_error ("Undefined return code of quit/exit command")
       end if
    end if
    global%quit = .true.
  end subroutine cmd_quit_execute

\eatline
\nwendcode{}\nwbegindocs{3002}\nwdocspar
\subsection{The command list}
The command list holds a list of commands and relevant global data.
\nwenddocs{}\nwbegincode{3003}\moddef{Commands: public}\plusendmoddef
  public :: command_list_t
\nwendcode{}\nwbegincode{3004}\moddef{Commands: types}\plusendmoddef
  type :: command_list_t
     ! not private anymore as required by the whizard-c-interface
     class(command_t), pointer :: first => null ()
     class(command_t), pointer :: last => null ()
   contains
   \LA{}Commands: command list: TBP\RA{}
  end type command_list_t

\eatline
\nwendcode{}\nwbegindocs{3005}Output.
\nwenddocs{}\nwbegincode{3006}\moddef{Commands: command list: TBP}\endmoddef
  procedure :: write => command_list_write
\nwendcode{}\nwbegincode{3007}\moddef{Commands: procedures}\plusendmoddef
  recursive subroutine command_list_write (cmd_list, unit, indent)
    class(command_list_t), intent(in) :: cmd_list
    integer, intent(in), optional :: unit, indent
    class(command_t), pointer :: cmd
    cmd => cmd_list%first
    do while (associated (cmd))
       call cmd%write (unit, indent)
       cmd => cmd%next
    end do
  end subroutine command_list_write
  
\eatline
\nwendcode{}\nwbegindocs{3008}Append a new command to the list and free the original pointer.
\nwenddocs{}\nwbegincode{3009}\moddef{Commands: command list: TBP}\plusendmoddef
  procedure :: append => command_list_append
\nwendcode{}\nwbegincode{3010}\moddef{Commands: procedures}\plusendmoddef
  subroutine command_list_append (cmd_list, command)
    class(command_list_t), intent(inout) :: cmd_list
    class(command_t), intent(inout), pointer :: command
    if (associated (cmd_list%last)) then
       cmd_list%last%next => command
    else
       cmd_list%first => command
    end if
    cmd_list%last => command
    command => null ()
  end subroutine command_list_append

\eatline
\nwendcode{}\nwbegindocs{3011}Finalize.
\nwenddocs{}\nwbegincode{3012}\moddef{Commands: command list: TBP}\plusendmoddef
  procedure :: final => command_list_final
\nwendcode{}\nwbegincode{3013}\moddef{Commands: procedures}\plusendmoddef
  recursive subroutine command_list_final (cmd_list)
    class(command_list_t), intent(inout) :: cmd_list
    class(command_t), pointer :: command
    do while (associated (cmd_list%first))
       command => cmd_list%first
       cmd_list%first => cmd_list%first%next
       call command%final ()
       deallocate (command)
    end do
    cmd_list%last => null ()
  end subroutine command_list_final

\eatline
\nwendcode{}\nwbegindocs{3014}\nwdocspar
\subsection{Compiling the parse tree}
Transform a parse tree into a command list.  Initialization is assumed
to be done.  

After each command, we set a breakpoint.
\nwenddocs{}\nwbegincode{3015}\moddef{Commands: command list: TBP}\plusendmoddef
  procedure :: compile => command_list_compile
\nwendcode{}\nwbegincode{3016}\moddef{Commands: procedures}\plusendmoddef
  recursive subroutine command_list_compile (cmd_list, pn, global)
    class(command_list_t), intent(inout), target :: cmd_list
    type(parse_node_t), intent(in), target :: pn
    type(rt_data_t), intent(inout), target :: global
    type(parse_node_t), pointer :: pn_cmd
    class(command_t), pointer :: command
    integer :: i
    pn_cmd => parse_node_get_sub_ptr (pn)
    do i = 1, parse_node_get_n_sub (pn)
       call dispatch_command (command, pn_cmd)
       call command%compile (global)
       call cmd_list%append (command)
       call terminate_now_if_signal ()
       pn_cmd => parse_node_get_next_ptr (pn_cmd)
    end do
  end subroutine command_list_compile

\eatline
\nwendcode{}\nwbegindocs{3017}\nwdocspar
\subsection{Executing the command list}
Before executing a command we should execute its options (if any).  After
that, reset the options, i.e., remove temporary effects from the global
state.

Also here, after each command we set a breakpoint.
\nwenddocs{}\nwbegincode{3018}\moddef{Commands: command list: TBP}\plusendmoddef
  procedure :: execute => command_list_execute
\nwendcode{}\nwbegincode{3019}\moddef{Commands: procedures}\plusendmoddef
  recursive subroutine command_list_execute (cmd_list, global)
    class(command_list_t), intent(in) :: cmd_list
    type(rt_data_t), intent(inout), target :: global
    class(command_t), pointer :: command
    command => cmd_list%first
    COMMAND_COND: do while (associated (command))
       call command%execute_options (global)
       call command%execute (global)
       call command%reset_options (global)
       call terminate_now_if_signal ()
       if (global%quit)  exit COMMAND_COND
       command => command%next
    end do COMMAND_COND
  end subroutine command_list_execute

\eatline
\nwendcode{}\nwbegindocs{3020}\nwdocspar
\subsection{Command list syntax}
\nwenddocs{}\nwbegincode{3021}\moddef{Commands: public}\plusendmoddef
  public :: syntax_cmd_list
\nwendcode{}\nwbegincode{3022}\moddef{Commands: variables}\endmoddef
  type(syntax_t), target, save :: syntax_cmd_list

\eatline
\nwendcode{}\nwbegincode{3023}\moddef{Commands: public}\plusendmoddef
  public :: syntax_cmd_list_init
\nwendcode{}\nwbegincode{3024}\moddef{Commands: procedures}\plusendmoddef
  subroutine syntax_cmd_list_init ()
    type(ifile_t) :: ifile
    call define_cmd_list_syntax (ifile)
    call syntax_init (syntax_cmd_list, ifile)
    call ifile_final (ifile)
  end subroutine syntax_cmd_list_init

\eatline
\nwendcode{}\nwbegincode{3025}\moddef{Commands: public}\plusendmoddef
  public :: syntax_cmd_list_final
\nwendcode{}\nwbegincode{3026}\moddef{Commands: procedures}\plusendmoddef
  subroutine syntax_cmd_list_final ()
    call syntax_final (syntax_cmd_list)
  end subroutine syntax_cmd_list_final

\eatline
\nwendcode{}\nwbegincode{3027}\moddef{Commands: public}\plusendmoddef
  public :: syntax_cmd_list_write
\nwendcode{}\nwbegincode{3028}\moddef{Commands: procedures}\plusendmoddef
  subroutine syntax_cmd_list_write (unit)
    integer, intent(in), optional :: unit
    call syntax_write (syntax_cmd_list, unit)
  end subroutine syntax_cmd_list_write

\eatline
\nwendcode{}\nwbegincode{3029}\moddef{Commands: procedures}\plusendmoddef
  subroutine define_cmd_list_syntax (ifile)
    type(ifile_t), intent(inout) :: ifile
    call ifile_append (ifile, "SEQ command_list = command*")
    call ifile_append (ifile, "ALT command = " &
         // "cmd_model | cmd_library | cmd_iterations | cmd_sample_format | " &
         // "cmd_var | cmd_slha | " &
         // "cmd_show | cmd_clear | " &
         // "cmd_expect | " &
         // "cmd_cuts | cmd_scale | cmd_fac_scale | cmd_ren_scale | " &
         // "cmd_weight | cmd_selection | cmd_reweight | " &
         // "cmd_beams | cmd_beams_pol_density | cmd_beams_pol_fraction | " &
         // "cmd_beams_momentum | cmd_beams_theta | cmd_beams_phi | " &
         // "cmd_integrate | " &
         // "cmd_observable | cmd_histogram | cmd_plot | cmd_graph | " &
         // "cmd_record | " &
         // "cmd_analysis | cmd_alt_setup | " &
         // "cmd_unstable | cmd_stable | cmd_simulate | cmd_rescan | " &
         // "cmd_process | cmd_compile | cmd_exec | " &
         // "cmd_scan | cmd_if | cmd_include | cmd_quit | " &
         // "cmd_polarized | cmd_unpolarized | " &
         // "cmd_open_out | cmd_close_out | cmd_printf | " &
         // "cmd_write_analysis | cmd_compile_analysis | cmd_nlo | cmd_components") 
    call ifile_append (ifile, "GRO options = '\{' local_command_list '\}'")
    call ifile_append (ifile, "SEQ local_command_list = local_command*")
    call ifile_append (ifile, "ALT local_command = " &
         // "cmd_model | cmd_library | cmd_iterations | cmd_sample_format | " &
         // "cmd_var | cmd_slha | " &
         // "cmd_show | " &
         // "cmd_expect | " &
         // "cmd_cuts | cmd_scale | cmd_fac_scale | cmd_ren_scale | " &
         // "cmd_weight | cmd_selection | cmd_reweight | " &
         // "cmd_beams | cmd_beams_pol_density | cmd_beams_pol_fraction | " &
         // "cmd_beams_momentum | cmd_beams_theta | cmd_beams_phi | " &
         // "cmd_observable | cmd_histogram | cmd_plot | cmd_graph | " &
         // "cmd_clear | cmd_record | " &
         // "cmd_analysis | cmd_alt_setup | " &
         // "cmd_open_out | cmd_close_out | cmd_printf | " &
         // "cmd_write_analysis | cmd_compile_analysis | cmd_nlo | cmd_components")
    call ifile_append (ifile, "SEQ cmd_model = model '=' model_name")
    call ifile_append (ifile, "KEY model")
    call ifile_append (ifile, "ALT model_name = model_id | string_literal")
    call ifile_append (ifile, "IDE model_id")
    call ifile_append (ifile, "SEQ cmd_library = library '=' lib_name")
    call ifile_append (ifile, "KEY library")
    call ifile_append (ifile, "ALT lib_name = lib_id | string_literal")
    call ifile_append (ifile, "IDE lib_id")
    call ifile_append (ifile, "ALT cmd_var = " &
         // "cmd_log_decl | cmd_log | " &
         // "cmd_int | cmd_real | cmd_complex | cmd_num | " &
         // "cmd_string_decl | cmd_string | cmd_alias | " &
         // "cmd_result")
    call ifile_append (ifile, "SEQ cmd_log_decl = logical cmd_log")
    call ifile_append (ifile, "SEQ cmd_log = '?' var_name '=' lexpr")
    call ifile_append (ifile, "SEQ cmd_int = int var_name '=' expr")
    call ifile_append (ifile, "SEQ cmd_real = real var_name '=' expr")
    call ifile_append (ifile, "SEQ cmd_complex = complex var_name '=' expr")
    call ifile_append (ifile, "SEQ cmd_num = var_name '=' expr")
    call ifile_append (ifile, "SEQ cmd_string_decl = string cmd_string")
    call ifile_append (ifile, "SEQ cmd_string = " &
         // "'$' var_name '=' sexpr") ! $
    call ifile_append (ifile, "SEQ cmd_alias = alias var_name '=' cexpr")
    call ifile_append (ifile, "SEQ cmd_result = result '=' expr") 
    call ifile_append (ifile, "SEQ cmd_slha = slha_action slha_arg options?")
    call ifile_append (ifile, "ALT slha_action = " &
         // "read_slha | write_slha")
    call ifile_append (ifile, "KEY read_slha")
    call ifile_append (ifile, "KEY write_slha")
    call ifile_append (ifile, "ARG slha_arg = ( string_literal )")
    call ifile_append (ifile, "SEQ cmd_show = show show_arg options?")
    call ifile_append (ifile, "KEY show")
    call ifile_append (ifile, "ARG show_arg = ( showable* )")
    call ifile_append (ifile, "ALT showable = " &
         // "model | library | beams | iterations | " &
         // "cuts | weight | logical | string | pdg | " &
         // "scale | factorization_scale | renormalization_scale | " & 
         // "selection | reweight | analysis | " &
         // "stable | unstable | polarized | unpolarized | " &
         // "expect | intrinsic | int | real | complex | " &
         // "alias_var | string | results | result_var | " &
         // "log_var | string_var | var_name")
    call ifile_append (ifile, "KEY results")
    call ifile_append (ifile, "KEY intrinsic")    
    call ifile_append (ifile, "SEQ alias_var = alias var_name")
    call ifile_append (ifile, "SEQ result_var = result_key result_arg?")
    call ifile_append (ifile, "SEQ log_var = '?' var_name")
    call ifile_append (ifile, "SEQ string_var = '$' var_name")  ! $
    call ifile_append (ifile, "SEQ cmd_clear = clear clear_arg options?")    
    call ifile_append (ifile, "KEY clear")
    call ifile_append (ifile, "ARG clear_arg = ( clearable* )")
    call ifile_append (ifile, "ALT clearable = " &
         // "beams | iterations | " &
         // "cuts | weight | " &
         // "scale | factorization_scale | renormalization_scale | " & 
         // "selection | reweight | analysis | " &
         // "unstable | polarized | " &
         // "expect | " &
         // "log_var | string_var | var_name")
    call ifile_append (ifile, "SEQ cmd_expect = expect expect_arg options?")
    call ifile_append (ifile, "KEY expect")
    call ifile_append (ifile, "ARG expect_arg = ( lexpr )")
    call ifile_append (ifile, "SEQ cmd_cuts = cuts '=' lexpr")
    call ifile_append (ifile, "SEQ cmd_scale = scale '=' expr")    
    call ifile_append (ifile, "SEQ cmd_fac_scale = " &
         // "factorization_scale '=' expr")
    call ifile_append (ifile, "SEQ cmd_ren_scale = " &
         // "renormalization_scale '=' expr")
    call ifile_append (ifile, "SEQ cmd_weight = weight '=' expr")
    call ifile_append (ifile, "SEQ cmd_selection = selection '=' lexpr")
    call ifile_append (ifile, "SEQ cmd_reweight = reweight '=' expr")
    call ifile_append (ifile, "KEY cuts")
    call ifile_append (ifile, "KEY scale")    
    call ifile_append (ifile, "KEY factorization_scale")
    call ifile_append (ifile, "KEY renormalization_scale")    
    call ifile_append (ifile, "KEY weight")
    call ifile_append (ifile, "KEY selection")
    call ifile_append (ifile, "KEY reweight")
    call ifile_append (ifile, "SEQ cmd_process = process process_id '=' " &
         // "process_prt '=>' prt_state_list options?")
    call ifile_append (ifile, "KEY process")
    call ifile_append (ifile, "KEY '=>'")
    call ifile_append (ifile, "LIS process_prt = cexpr+")
    call ifile_append (ifile, "LIS prt_state_list = prt_state_sum+")
    call ifile_append (ifile, "SEQ prt_state_sum = " &
         // "prt_state prt_state_addition*")
    call ifile_append (ifile, "SEQ prt_state_addition = '+' prt_state")
    call ifile_append (ifile, "ALT prt_state = grouped_prt_state_list | cexpr")
    call ifile_append (ifile, "GRO grouped_prt_state_list = " &
         // "( prt_state_list )")
    call ifile_append (ifile, "SEQ cmd_compile = compile_cmd options?")
    call ifile_append (ifile, "SEQ compile_cmd = compile_clause compile_arg?")
    call ifile_append (ifile, "SEQ compile_clause = compile exec_name_spec?")
    call ifile_append (ifile, "KEY compile")
    call ifile_append (ifile, "SEQ exec_name_spec = as exec_name")
    call ifile_append (ifile, "KEY as")
    call ifile_append (ifile, "ALT exec_name = exec_id | string_literal")
    call ifile_append (ifile, "IDE exec_id")
    call ifile_append (ifile, "ARG compile_arg = ( lib_name* )")
    call ifile_append (ifile, "SEQ cmd_exec = exec exec_arg")
    call ifile_append (ifile, "KEY exec")
    call ifile_append (ifile, "ARG exec_arg = ( sexpr )")
    call ifile_append (ifile, "SEQ cmd_beams = beams '=' beam_def")
    call ifile_append (ifile, "KEY beams")
    call ifile_append (ifile, "SEQ beam_def = beam_spec strfun_seq*")
    call ifile_append (ifile, "SEQ beam_spec = beam_list")
    call ifile_append (ifile, "LIS beam_list = cexpr, cexpr?")
    call ifile_append (ifile, "SEQ cmd_beams_pol_density = " &
         // "beams_pol_density '=' beams_pol_spec")
    call ifile_append (ifile, "KEY beams_pol_density")
    call ifile_append (ifile, "LIS beams_pol_spec = smatrix, smatrix?")
    call ifile_append (ifile, "SEQ smatrix = '@' smatrix_arg")
    ! call ifile_append (ifile, "KEY '@'")     !!! Key already exists
    call ifile_append (ifile, "ARG smatrix_arg = ( sentry* )")
    call ifile_append (ifile, "SEQ sentry = expr extra_sentry*")
    call ifile_append (ifile, "SEQ extra_sentry = ':' expr")
    call ifile_append (ifile, "SEQ cmd_beams_pol_fraction = " &
         // "beams_pol_fraction '=' beams_par_spec")
    call ifile_append (ifile, "KEY beams_pol_fraction")
    call ifile_append (ifile, "SEQ cmd_beams_momentum = " &
         // "beams_momentum '=' beams_par_spec")
    call ifile_append (ifile, "KEY beams_momentum")
    call ifile_append (ifile, "SEQ cmd_beams_theta = " &
         // "beams_theta '=' beams_par_spec")
    call ifile_append (ifile, "KEY beams_theta")
    call ifile_append (ifile, "SEQ cmd_beams_phi = " &
         // "beams_phi '=' beams_par_spec")
    call ifile_append (ifile, "KEY beams_phi")
    call ifile_append (ifile, "LIS beams_par_spec = expr, expr?")
    call ifile_append (ifile, "SEQ strfun_seq = '=>' strfun_pair")
    call ifile_append (ifile, "LIS strfun_pair = strfun_def, strfun_def?")
    call ifile_append (ifile, "SEQ strfun_def = strfun_id")
    call ifile_append (ifile, "ALT strfun_id = " &
          // "none | lhapdf | lhapdf_photon | pdf_builtin | pdf_builtin_photon | " &
          // "isr | epa | ewa | circe1 | circe2 | energy_scan | " &
          // "beam_events | user_sf_spec")
    call ifile_append (ifile, "KEY none")
    call ifile_append (ifile, "KEY lhapdf")
    call ifile_append (ifile, "KEY lhapdf_photon")    
    call ifile_append (ifile, "KEY pdf_builtin")    
    call ifile_append (ifile, "KEY pdf_builtin_photon")        
    call ifile_append (ifile, "KEY isr")
    call ifile_append (ifile, "KEY epa")
    call ifile_append (ifile, "KEY ewa")    
    call ifile_append (ifile, "KEY circe1")        
    call ifile_append (ifile, "KEY circe2")
    call ifile_append (ifile, "KEY energy_scan")
    call ifile_append (ifile, "KEY beam_events")
    call ifile_append (ifile, "SEQ user_sf_spec = user_strfun user_arg")
    call ifile_append (ifile, "KEY user_strfun")
    call ifile_append (ifile, "SEQ cmd_integrate = " &
         // "integrate proc_arg options?") 
    call ifile_append (ifile, "KEY integrate")
    call ifile_append (ifile, "ARG proc_arg = ( proc_id* )")
    call ifile_append (ifile, "IDE proc_id")
    call ifile_append (ifile, "SEQ cmd_iterations = " &
         // "iterations '=' iterations_list")
    call ifile_append (ifile, "KEY iterations")
    call ifile_append (ifile, "LIS iterations_list = iterations_spec+")
    call ifile_append (ifile, "ALT iterations_spec = it_spec")
    call ifile_append (ifile, "SEQ it_spec = expr calls_spec adapt_spec?")
    call ifile_append (ifile, "SEQ calls_spec = ':' expr")
    call ifile_append (ifile, "SEQ adapt_spec = ':' sexpr")
    call ifile_append (ifile, "SEQ cmd_components = " &
         // "active '=' component_list")
    call ifile_append (ifile, "KEY active")
    call ifile_append (ifile, "LIS component_list = sexpr+")
    call ifile_append (ifile, "SEQ cmd_sample_format = " &
         // "sample_format '=' event_format_list")
    call ifile_append (ifile, "KEY sample_format")
    call ifile_append (ifile, "LIS event_format_list = event_format+")
    call ifile_append (ifile, "IDE event_format")
    call ifile_append (ifile, "SEQ cmd_observable = " &
         // "observable analysis_tag options?")
    call ifile_append (ifile, "KEY observable")
    call ifile_append (ifile, "SEQ cmd_histogram = " &
         // "histogram analysis_tag histogram_arg " & 
         // "options?")
    call ifile_append (ifile, "KEY histogram")
    call ifile_append (ifile, "ARG histogram_arg = (expr, expr, expr?)")
    call ifile_append (ifile, "SEQ cmd_plot = plot analysis_tag options?")
    call ifile_append (ifile, "KEY plot")
    call ifile_append (ifile, "SEQ cmd_graph = graph graph_term '=' graph_def")
    call ifile_append (ifile, "KEY graph")
    call ifile_append (ifile, "SEQ graph_term = analysis_tag options?")
    call ifile_append (ifile, "SEQ graph_def = graph_term graph_append*")
    call ifile_append (ifile, "SEQ graph_append = '&' graph_term")
    call ifile_append (ifile, "SEQ cmd_analysis = analysis '=' lexpr")
    call ifile_append (ifile, "KEY analysis")
    call ifile_append (ifile, "SEQ cmd_alt_setup = " &
         // "alt_setup '=' option_list_expr")
    call ifile_append (ifile, "KEY alt_setup")
    call ifile_append (ifile, "ALT option_list_expr = " &
         // "grouped_option_list | option_list")
    call ifile_append (ifile, "GRO grouped_option_list = ( option_list_expr )")
    call ifile_append (ifile, "LIS option_list = options+")
    call ifile_append (ifile, "SEQ cmd_open_out = open_out open_arg options?")
    call ifile_append (ifile, "SEQ cmd_close_out = close_out open_arg options?")
    call ifile_append (ifile, "KEY open_out")
    call ifile_append (ifile, "KEY close_out")
    call ifile_append (ifile, "ARG open_arg = (sexpr)")
    call ifile_append (ifile, "SEQ cmd_printf = printf_cmd options?")
    call ifile_append (ifile, "SEQ printf_cmd = printf_clause sprintf_args?")
    call ifile_append (ifile, "SEQ printf_clause = printf sexpr")
    call ifile_append (ifile, "KEY printf")
    call ifile_append (ifile, "SEQ cmd_record = record_cmd")
    call ifile_append (ifile, "SEQ cmd_unstable = " &
         // "unstable cexpr unstable_arg options?")
    call ifile_append (ifile, "KEY unstable")
    call ifile_append (ifile, "ARG unstable_arg = ( proc_id* )")
    call ifile_append (ifile, "SEQ cmd_stable = stable stable_list options?")
    call ifile_append (ifile, "KEY stable")
    call ifile_append (ifile, "LIS stable_list = cexpr+")
    call ifile_append (ifile, "KEY polarized")
    call ifile_append (ifile, "SEQ cmd_polarized = polarized polarized_list options?")
    call ifile_append (ifile, "LIS polarized_list = cexpr+")
    call ifile_append (ifile, "KEY unpolarized")
    call ifile_append (ifile, "SEQ cmd_unpolarized = unpolarized unpolarized_list options?")
    call ifile_append (ifile, "LIS unpolarized_list = cexpr+")
    call ifile_append (ifile, "SEQ cmd_simulate = " &
         // "simulate proc_arg options?")
    call ifile_append (ifile, "KEY simulate")
    call ifile_append (ifile, "SEQ cmd_rescan = " &
         // "rescan sexpr proc_arg options?")
    call ifile_append (ifile, "KEY rescan")
    call ifile_append (ifile, "SEQ cmd_scan = scan scan_var scan_body?")
    call ifile_append (ifile, "KEY scan")
    call ifile_append (ifile, "ALT scan_var = " &
         // "scan_log_decl | scan_log | " &
         // "scan_int | scan_real | scan_complex | scan_num | " &
         // "scan_string_decl | scan_string | scan_alias | " &
         // "scan_cuts | scan_weight | " &
         // "scan_scale | scan_ren_scale | scan_fac_scale | " &
         // "scan_selection | scan_reweight | scan_analysis | " &
         // "scan_model | scan_library")
    call ifile_append (ifile, "SEQ scan_log_decl = logical scan_log")
    call ifile_append (ifile, "SEQ scan_log = '?' var_name '=' scan_log_arg")
    call ifile_append (ifile, "ARG scan_log_arg = ( lexpr* )")
    call ifile_append (ifile, "SEQ scan_int = int var_name '=' scan_num_arg")
    call ifile_append (ifile, "SEQ scan_real = real var_name '=' scan_num_arg")
    call ifile_append (ifile, "SEQ scan_complex = " &
         // "complex var_name '=' scan_num_arg")
    call ifile_append (ifile, "SEQ scan_num = var_name '=' scan_num_arg")
    call ifile_append (ifile, "ARG scan_num_arg = ( range* )")
    call ifile_append (ifile, "ALT range = grouped_range | range_expr")
    call ifile_append (ifile, "GRO grouped_range = ( range_expr )")
    call ifile_append (ifile, "SEQ range_expr = expr range_spec?")
    call ifile_append (ifile, "SEQ range_spec = '=>' expr step_spec?")
    call ifile_append (ifile, "SEQ step_spec = step_op expr")
    call ifile_append (ifile, "ALT step_op = " &
         // "'/+' | '/-' | '/*' | '//' | '/+/' | '/*/'")
    call ifile_append (ifile, "KEY '/+'")
    call ifile_append (ifile, "KEY '/-'")
    call ifile_append (ifile, "KEY '/*'")
    call ifile_append (ifile, "KEY '//'")
    call ifile_append (ifile, "KEY '/+/'")
    call ifile_append (ifile, "KEY '/*/'")
    call ifile_append (ifile, "SEQ scan_string_decl = string scan_string")
    call ifile_append (ifile, "SEQ scan_string = " &
         // "'$' var_name '=' scan_string_arg")
    call ifile_append (ifile, "ARG scan_string_arg = ( sexpr* )")
    call ifile_append (ifile, "SEQ scan_alias = " &
         // "alias var_name '=' scan_alias_arg")
    call ifile_append (ifile, "ARG scan_alias_arg = ( cexpr* )")
    call ifile_append (ifile, "SEQ scan_cuts = cuts '=' scan_lexpr_arg")
    call ifile_append (ifile, "ARG scan_lexpr_arg = ( lexpr* )")
    call ifile_append (ifile, "SEQ scan_scale = scale '=' scan_expr_arg")
    call ifile_append (ifile, "ARG scan_expr_arg = ( expr* )")
    call ifile_append (ifile, "SEQ scan_fac_scale = " &
         // "factorization_scale '=' scan_expr_arg")
    call ifile_append (ifile, "SEQ scan_ren_scale = " &
         // "renormalization_scale '=' scan_expr_arg")
    call ifile_append (ifile, "SEQ scan_weight = weight '=' scan_expr_arg")
    call ifile_append (ifile, "SEQ scan_selection = selection '=' scan_lexpr_arg")
    call ifile_append (ifile, "SEQ scan_reweight = reweight '=' scan_expr_arg")
    call ifile_append (ifile, "SEQ scan_analysis = analysis '=' scan_lexpr_arg")
    call ifile_append (ifile, "SEQ scan_model = model '=' scan_model_arg")
    call ifile_append (ifile, "ARG scan_model_arg = ( model_name* )")
    call ifile_append (ifile, "SEQ scan_library = library '=' scan_library_arg")
    call ifile_append (ifile, "ARG scan_library_arg = ( lib_name* )")
    call ifile_append (ifile, "GRO scan_body = '\{' command_list '\}'")
    call ifile_append (ifile, "SEQ cmd_if = " &
         // "if lexpr then command_list elsif_clauses else_clause endif")
    call ifile_append (ifile, "SEQ elsif_clauses = cmd_elsif*")
    call ifile_append (ifile, "SEQ cmd_elsif = elsif lexpr then command_list")
    call ifile_append (ifile, "SEQ else_clause = cmd_else?")
    call ifile_append (ifile, "SEQ cmd_else = else command_list")
    call ifile_append (ifile, "SEQ cmd_include = include include_arg")
    call ifile_append (ifile, "KEY include")
    call ifile_append (ifile, "ARG include_arg = ( string_literal )")
    call ifile_append (ifile, "SEQ cmd_quit = quit_cmd quit_arg?")
    call ifile_append (ifile, "ALT quit_cmd = quit | exit")
    call ifile_append (ifile, "KEY quit")
    call ifile_append (ifile, "KEY exit")
    call ifile_append (ifile, "ARG quit_arg = ( expr )")
    call ifile_append (ifile, "SEQ cmd_write_analysis = " &
         // "write_analysis_clause options?")
    call ifile_append (ifile, "SEQ cmd_compile_analysis = " &
         // "compile_analysis_clause options?")
    call ifile_append (ifile, "SEQ write_analysis_clause = " &
         // "write_analysis write_analysis_arg?")
    call ifile_append (ifile, "SEQ compile_analysis_clause = " &
         // "compile_analysis write_analysis_arg?")
    call ifile_append (ifile, "KEY write_analysis")
    call ifile_append (ifile, "KEY compile_analysis")
    call ifile_append (ifile, "ARG write_analysis_arg = ( analysis_tag* )")
    call ifile_append (ifile, "SEQ cmd_nlo = " &
                       // "nlo_calculation '=' nlo_calculation_list")
    call ifile_append (ifile, "KEY nlo_calculation")
    call ifile_append (ifile, "LIS nlo_calculation_list = sexpr ',' sexpr ',' sexpr")
    call define_expr_syntax (ifile, particles=.true., analysis=.true.)
  end subroutine define_cmd_list_syntax

\eatline
\nwendcode{}\nwbegincode{3030}\moddef{Commands: public}\plusendmoddef
  public :: lexer_init_cmd_list
\nwendcode{}\nwbegincode{3031}\moddef{Commands: procedures}\plusendmoddef
  subroutine lexer_init_cmd_list (lexer, parent_lexer)
    type(lexer_t), intent(out) :: lexer
    type(lexer_t), intent(in), optional, target :: parent_lexer
    call lexer_init (lexer, &
         comment_chars = "#!", &
         quote_chars = '"', &
         quote_match = '"', &
         single_chars = "()[]\{\},;:&%?$@", &
         special_class = [ "+-*/^", "<>=~ " ] , &
         keyword_list = syntax_get_keyword_list_ptr (syntax_cmd_list), &
         parent = parent_lexer)
  end subroutine lexer_init_cmd_list

\eatline
\nwendcode{}\nwbegindocs{3032}\nwdocspar
\subsection{Unit Tests}
\nwenddocs{}\nwbegincode{3033}\moddef{Commands: public}\plusendmoddef
  public :: commands_test
\nwendcode{}\nwbegincode{3034}\moddef{Commands: tests}\endmoddef
  subroutine commands_test (u, results)
    integer, intent(in) :: u
    type(test_results_t), intent(inout) :: results
  \LA{}Commands: execute tests\RA{}
  end subroutine commands_test
  
\eatline
\nwendcode{}\nwbegindocs{3035}\nwdocspar
\subsubsection{Prepare Sindarin code}
This routine parses an internal file, prints the parse tree, and
returns a parse node to the root.  We use the routine in the tests
below.
\nwenddocs{}\nwbegincode{3036}\moddef{Commands: tests}\plusendmoddef
  subroutine parse_ifile (ifile, pn_root, u)
    type(ifile_t), intent(in) :: ifile
    type(parse_node_t), pointer, intent(out) :: pn_root
    integer, intent(in), optional :: u
    type(stream_t), target :: stream
    type(lexer_t), target :: lexer
    type(parse_tree_t) :: parse_tree

    call lexer_init_cmd_list (lexer)
    call stream_init (stream, ifile)
    call lexer_assign_stream (lexer, stream)

    call parse_tree_init (parse_tree, syntax_cmd_list, lexer)
    if (present (u))  call parse_tree_write (parse_tree, u)
    pn_root => parse_tree_get_root_ptr (parse_tree)

    call stream_final (stream)
    call lexer_final (lexer)
  end subroutine parse_ifile

\eatline
\nwendcode{}\nwbegindocs{3037}\nwdocspar
\subsubsection{Empty command list}
Compile and execute an empty command list.  Should do nothing but
test the integrity of the workflow.
\nwenddocs{}\nwbegincode{3038}\moddef{Commands: execute tests}\endmoddef
  call test (commands_1, "commands_1", &
       "empty command list", &
       u, results)
\nwendcode{}\nwbegincode{3039}\moddef{Commands: tests}\plusendmoddef
  subroutine commands_1 (u)
    integer, intent(in) :: u
    type(ifile_t) :: ifile
    type(command_list_t), target :: command_list
    type(rt_data_t), target :: global
    type(parse_node_t), pointer :: pn_root

    write (u, "(A)")  "* Test output: commands_1"
    write (u, "(A)")  "*   Purpose: compile and execute empty command list"
    write (u, "(A)")

    write (u, "(A)")  "*  Initialization"
    write (u, "(A)")

    call syntax_cmd_list_init ()
    call global%global_init ()

    write (u, "(A)")  "*  Parse empty file"
    write (u, "(A)")
    
    call parse_ifile (ifile, pn_root, u)

    write (u, "(A)")
    write (u, "(A)")  "* Compile command list"

    if (associated (pn_root)) then
       call command_list%compile (pn_root, global)
    end if

    write (u, "(A)")
    write (u, "(A)")  "* Execute command list"

    call global%activate ()
    call command_list%execute (global)
    call global%deactivate ()

    write (u, "(A)")
    write (u, "(A)")  "* Cleanup"

    call ifile_final (ifile)

    call command_list%final ()
    call syntax_cmd_list_final ()
    call global%final ()

    write (u, "(A)")
    write (u, "(A)")  "* Test output end: commands_1"

  end subroutine commands_1

\eatline
\nwendcode{}\nwbegindocs{3040}\nwdocspar
\subsubsection{Read model}
Execute a {\tt{}model} assignment.
\nwenddocs{}\nwbegincode{3041}\moddef{Commands: execute tests}\plusendmoddef
  call test (commands_2, "commands_2", &
       "model", &
       u, results)
\nwendcode{}\nwbegincode{3042}\moddef{Commands: tests}\plusendmoddef
  subroutine commands_2 (u)
    integer, intent(in) :: u
    type(ifile_t) :: ifile
    type(command_list_t), target :: command_list
    type(rt_data_t), target :: global
    type(parse_node_t), pointer :: pn_root

    write (u, "(A)")  "* Test output: commands_2"
    write (u, "(A)")  "*   Purpose: set model"
    write (u, "(A)")

    write (u, "(A)")  "*  Initialization"
    write (u, "(A)")

    call syntax_cmd_list_init ()
    call syntax_model_file_init ()
    call global%global_init ()

    write (u, "(A)")  "*  Input file"
    write (u, "(A)")
    
    call ifile_append (ifile, 'model = "Test"')
    
    call ifile_write (ifile, u)

    write (u, "(A)")  "*  Parse file"
    write (u, "(A)")
    
    call parse_ifile (ifile, pn_root, u)

    write (u, "(A)")
    write (u, "(A)")  "* Compile command list"
    write (u, "(A)")

    call command_list%compile (pn_root, global)
    call command_list%write (u)

    write (u, "(A)")
    write (u, "(A)")  "* Execute command list"
    write (u, "(A)")

    call command_list%execute (global)

    write (u, "(A)")  "* Cleanup"

    call ifile_final (ifile)

    call command_list%final ()
    call global%final ()
    call syntax_cmd_list_final ()
    call syntax_model_file_final ()

    write (u, "(A)")
    write (u, "(A)")  "* Test output end: commands_2"

  end subroutine commands_2

\eatline
\nwendcode{}\nwbegindocs{3043}\nwdocspar
\subsubsection{Declare Process}
Read a model, then declare a process.  The process library is allocated
explicitly.  For the process definition, We take the default ({\tt{}omega})
method.  Since we do not compile, \oMega\ is not actually called.
\nwenddocs{}\nwbegincode{3044}\moddef{Commands: execute tests}\plusendmoddef
  call test (commands_3, "commands_3", &
       "process declaration", &
       u, results)
\nwendcode{}\nwbegincode{3045}\moddef{Commands: tests}\plusendmoddef
  subroutine commands_3 (u)
    integer, intent(in) :: u
    type(ifile_t) :: ifile
    type(command_list_t), target :: command_list
    type(rt_data_t), target :: global
    type(parse_node_t), pointer :: pn_root
    type(prclib_entry_t), pointer :: lib

    write (u, "(A)")  "* Test output: commands_3"
    write (u, "(A)")  "*   Purpose: define process"
    write (u, "(A)")

    write (u, "(A)")  "*  Initialization"
    write (u, "(A)")

    call syntax_cmd_list_init ()
    call syntax_model_file_init ()
    call global%global_init ()
    call var_list_set_log (global%var_list, var_str ("?omega_openmp"), &
         .false., is_known = .true.)

    allocate (lib)
    call lib%init (var_str ("lib_cmd3"))
    call global%add_prclib (lib)
    
    write (u, "(A)")  "*  Input file"
    write (u, "(A)")
    
    call ifile_append (ifile, 'model = "Test"')
    call ifile_append (ifile, 'process t3 = s, s => s, s')
    
    call ifile_write (ifile, u)

    write (u, "(A)")
    write (u, "(A)")  "*  Parse file"
    write (u, "(A)")
    
    call parse_ifile (ifile, pn_root, u)

    write (u, "(A)")
    write (u, "(A)")  "* Compile command list"
    write (u, "(A)")

    call command_list%compile (pn_root, global)
    call command_list%write (u)

    write (u, "(A)")
    write (u, "(A)")  "* Execute command list"
    write (u, "(A)")

    call command_list%execute (global)

    call global%prclib_stack%write (u)

    write (u, "(A)")
    write (u, "(A)")  "* Cleanup"

    call ifile_final (ifile)

    call command_list%final ()
    call global%final ()
    call syntax_cmd_list_final ()
    call syntax_model_file_final ()

    write (u, "(A)")
    write (u, "(A)")  "* Test output end: commands_3"

  end subroutine commands_3

\eatline
\nwendcode{}\nwbegindocs{3046}\nwdocspar
\subsubsection{Compile Process}
Read a model, then declare a process and compile the library.  The process
library is allocated explicitly.  For the process definition, We take the
default ({\tt{}unit{\char95}test}) method.  There is no external code, so compilation of
the library is merely a formal status change.
\nwenddocs{}\nwbegincode{3047}\moddef{Commands: execute tests}\plusendmoddef
  call test (commands_4, "commands_4", &
       "compilation", &
       u, results)
\nwendcode{}\nwbegincode{3048}\moddef{Commands: tests}\plusendmoddef
  subroutine commands_4 (u)
    integer, intent(in) :: u
    type(ifile_t) :: ifile
    type(command_list_t), target :: command_list
    type(rt_data_t), target :: global
    type(parse_node_t), pointer :: pn_root
    type(prclib_entry_t), pointer :: lib

    write (u, "(A)")  "* Test output: commands_4"
    write (u, "(A)")  "*   Purpose: define process and compile library"
    write (u, "(A)")

    write (u, "(A)")  "*  Initialization"
    write (u, "(A)")

    call syntax_cmd_list_init ()
    call syntax_model_file_init ()
    call global%global_init ()
    call var_list_set_string (global%var_list, var_str ("$method"), &
         var_str ("unit_test"), is_known=.true.)

    allocate (lib)
    call lib%init (var_str ("lib_cmd4"))
    call global%add_prclib (lib)
    
    write (u, "(A)")  "*  Input file"
    write (u, "(A)")
    
    call ifile_append (ifile, 'model = "Test"')
    call ifile_append (ifile, 'process t4 = s, s => s, s')
    call ifile_append (ifile, 'compile ("lib_cmd4")')
    
    call ifile_write (ifile, u)

    write (u, "(A)")
    write (u, "(A)")  "*  Parse file"
    write (u, "(A)")
    
    call parse_ifile (ifile, pn_root, u)

    write (u, "(A)")
    write (u, "(A)")  "* Compile command list"
    write (u, "(A)")

    call command_list%compile (pn_root, global)
    call command_list%write (u)

    write (u, "(A)")
    write (u, "(A)")  "* Execute command list"
    write (u, "(A)")

    call command_list%execute (global)

    call global%prclib_stack%write (u)

    write (u, "(A)")
    write (u, "(A)")  "* Cleanup"

    call ifile_final (ifile)

    call command_list%final ()
    call global%final ()
    call syntax_cmd_list_final ()
    call syntax_model_file_final ()

    write (u, "(A)")
    write (u, "(A)")  "* Test output end: commands_4"

  end subroutine commands_4

\eatline
\nwendcode{}\nwbegindocs{3049}\nwdocspar
\subsubsection{Integrate Process}
Read a model, then declare a process, compile the library, and
integrate over phase space.  We take the
default ({\tt{}unit{\char95}test}) method and use the simplest methods of
phase-space parameterization and integration.
\nwenddocs{}\nwbegincode{3050}\moddef{Commands: execute tests}\plusendmoddef
  call test (commands_5, "commands_5", &
       "integration", &
       u, results)
\nwendcode{}\nwbegincode{3051}\moddef{Commands: tests}\plusendmoddef
  subroutine commands_5 (u)
    integer, intent(in) :: u
    type(ifile_t) :: ifile
    type(command_list_t), target :: command_list
    type(rt_data_t), target :: global
    type(parse_node_t), pointer :: pn_root
    type(prclib_entry_t), pointer :: lib

    write (u, "(A)")  "* Test output: commands_5"
    write (u, "(A)")  "*   Purpose: define process, iterations, and integrate"
    write (u, "(A)")

    write (u, "(A)")  "*  Initialization"
    write (u, "(A)")

    call syntax_cmd_list_init ()
    call syntax_model_file_init ()
    call global%global_init ()
    call var_list_set_string (global%var_list, var_str ("$method"), &
         var_str ("unit_test"), is_known=.true.)
    call var_list_set_string (global%var_list, var_str ("$phs_method"), &
         var_str ("single"), is_known=.true.)
    call var_list_set_string (global%var_list, var_str ("$integration_method"),&
         var_str ("midpoint"), is_known=.true.)
    call var_list_set_log (global%var_list, var_str ("?vis_history"),&
         .false., is_known=.true.)        
    call var_list_set_log (global%var_list, var_str ("?integration_timer"),&
         .false., is_known = .true.)    
    call var_list_set_real (global%var_list, var_str ("sqrts"), &
         1000._default, is_known=.true.)
    call var_list_set_int (global%var_list, var_str ("seed"), &
         0, is_known=.true.)
    
    allocate (lib)
    call lib%init (var_str ("lib_cmd5"))
    call global%add_prclib (lib)
    
    write (u, "(A)")  "*  Input file"
    write (u, "(A)")
    
    call ifile_append (ifile, 'model = "Test"')
    call ifile_append (ifile, 'process t5 = s, s => s, s')
    call ifile_append (ifile, 'compile')
    call ifile_append (ifile, 'iterations = 1:1000')
    call ifile_append (ifile, 'integrate (t5)')
    
    call ifile_write (ifile, u)

    write (u, "(A)")
    write (u, "(A)")  "*  Parse file"
    write (u, "(A)")
    
    call parse_ifile (ifile, pn_root, u)

    write (u, "(A)")
    write (u, "(A)")  "* Compile command list"
    write (u, "(A)")

    call command_list%compile (pn_root, global)
    call command_list%write (u)

    write (u, "(A)")
    write (u, "(A)")  "* Execute command list"
    write (u, "(A)")

    call reset_interaction_counter ()
    call command_list%execute (global)

    call global%it_list%write (u)
    write (u, "(A)")
    call global%process_stack%write (u)

    write (u, "(A)")
    write (u, "(A)")  "* Cleanup"

    call ifile_final (ifile)

    call command_list%final ()
    call global%final ()
    call syntax_cmd_list_final ()
    call syntax_model_file_final ()

    write (u, "(A)")
    write (u, "(A)")  "* Test output end: commands_5"

  end subroutine commands_5

\eatline
\nwendcode{}\nwbegindocs{3052}\nwdocspar
\subsubsection{Variables}
Set intrinsic and user-defined variables.
\nwenddocs{}\nwbegincode{3053}\moddef{Commands: execute tests}\plusendmoddef
  call test (commands_6, "commands_6", &
       "variables", &
       u, results)
\nwendcode{}\nwbegincode{3054}\moddef{Commands: tests}\plusendmoddef
  subroutine commands_6 (u)
    integer, intent(in) :: u
    type(ifile_t) :: ifile
    type(command_list_t), target :: command_list
    type(rt_data_t), target :: global
    type(parse_node_t), pointer :: pn_root

    write (u, "(A)")  "* Test output: commands_6"
    write (u, "(A)")  "*   Purpose: define and set variables"
    write (u, "(A)")

    write (u, "(A)")  "*  Initialization"
    write (u, "(A)")

    call syntax_cmd_list_init ()
    call global%global_init ()
    call global%write_vars (u, [ &
         var_str ("$run_id"), &
         var_str ("?unweighted"), &
         var_str ("sqrts")])

    write (u, "(A)")
    write (u, "(A)")  "*  Input file"
    write (u, "(A)")
    
    call ifile_append (ifile, '$run_id = "run1"')
    call ifile_append (ifile, '?unweighted = false')
    call ifile_append (ifile, 'sqrts = 1000')
    call ifile_append (ifile, 'int j = 10')
    call ifile_append (ifile, 'real x = 1000.')
    call ifile_append (ifile, 'complex z = 5')
    call ifile_append (ifile, 'string $text = "abcd"')
    call ifile_append (ifile, 'logical ?flag = true')
    
    call ifile_write (ifile, u)

    write (u, "(A)")
    write (u, "(A)")  "*  Parse file"
    write (u, "(A)")
    
    call parse_ifile (ifile, pn_root, u)

    write (u, "(A)")
    write (u, "(A)")  "* Compile command list"
    write (u, "(A)")

    call command_list%compile (pn_root, global)
    call command_list%write (u)

    write (u, "(A)")
    write (u, "(A)")  "* Execute command list"
    write (u, "(A)")

    call command_list%execute (global)

    call global%write_vars (u, [ &
         var_str ("$run_id"), &
         var_str ("?unweighted"), &
         var_str ("sqrts"), &
         var_str ("j"), &
         var_str ("x"), &
         var_str ("z"), &
         var_str ("$text"), &
         var_str ("?flag")])


    write (u, "(A)")
    write (u, "(A)")  "* Cleanup"

    call ifile_final (ifile)

    call command_list%final ()
    call syntax_cmd_list_final ()
    call global%final ()

    write (u, "(A)")
    write (u, "(A)")  "* Test output end: commands_6"

  end subroutine commands_6

\eatline
\nwendcode{}\nwbegindocs{3055}\nwdocspar
\subsubsection{Process library}
Open process libraries explicitly.
\nwenddocs{}\nwbegincode{3056}\moddef{Commands: execute tests}\plusendmoddef
  call test (commands_7, "commands_7", &
       "process library", &
       u, results)
\nwendcode{}\nwbegincode{3057}\moddef{Commands: tests}\plusendmoddef
  subroutine commands_7 (u)
    integer, intent(in) :: u
    type(ifile_t) :: ifile
    type(command_list_t), target :: command_list
    type(rt_data_t), target :: global
    type(parse_node_t), pointer :: pn_root

    write (u, "(A)")  "* Test output: commands_7"
    write (u, "(A)")  "*   Purpose: declare process libraries"
    write (u, "(A)")

    write (u, "(A)")  "*  Initialization"
    write (u, "(A)")

    call syntax_cmd_list_init ()
    call global%global_init ()
    call var_list_set_log (global%var_list, var_str ("?omega_openmp"), &
         .false., is_known = .true.)
    global%os_data%fc = "Fortran-compiler"
    global%os_data%fcflags = "Fortran-flags"

    write (u, "(A)")
    write (u, "(A)")  "*  Input file"
    write (u, "(A)")
    
    call ifile_append (ifile, 'library = "lib_cmd7_1"')
    call ifile_append (ifile, 'library = "lib_cmd7_2"')
    call ifile_append (ifile, 'library = "lib_cmd7_1"')
    
    call ifile_write (ifile, u)

    write (u, "(A)")
    write (u, "(A)")  "*  Parse file"
    write (u, "(A)")
    
    call parse_ifile (ifile, pn_root, u)

    write (u, "(A)")
    write (u, "(A)")  "* Compile command list"
    write (u, "(A)")

    call command_list%compile (pn_root, global)
    call command_list%write (u)

    write (u, "(A)")
    write (u, "(A)")  "* Execute command list"
    write (u, "(A)")

    call command_list%execute (global)

    call global%write_libraries (u)

    write (u, "(A)")
    write (u, "(A)")  "* Cleanup"

    call ifile_final (ifile)

    call command_list%final ()
    call syntax_cmd_list_final ()
    call global%final ()

    write (u, "(A)")
    write (u, "(A)")  "* Test output end: commands_7"

  end subroutine commands_7

\eatline
\nwendcode{}\nwbegindocs{3058}\nwdocspar
\subsubsection{Generate events}
Read a model, then declare a process, compile the library, and
generate weighted events.  We take the
default ({\tt{}unit{\char95}test}) method and use the simplest methods of
phase-space parameterization and integration.
\nwenddocs{}\nwbegincode{3059}\moddef{Commands: execute tests}\plusendmoddef
  call test (commands_8, "commands_8", &
       "event generation", &
       u, results)
\nwendcode{}\nwbegincode{3060}\moddef{Commands: tests}\plusendmoddef
  subroutine commands_8 (u)
    integer, intent(in) :: u
    type(ifile_t) :: ifile
    type(command_list_t), target :: command_list
    type(rt_data_t), target :: global
    type(parse_node_t), pointer :: pn_root
    type(prclib_entry_t), pointer :: lib

    write (u, "(A)")  "* Test output: commands_8"
    write (u, "(A)")  "*   Purpose: define process, integrate, generate events"
    write (u, "(A)")

    write (u, "(A)")  "*  Initialization"
    write (u, "(A)")

    call syntax_cmd_list_init ()
    call syntax_model_file_init ()
    call global%global_init ()
    call global%init_fallback_model &
         (var_str ("SM_hadrons"), var_str ("SM_hadrons.mdl"))

    call var_list_set_string (global%var_list, var_str ("$method"), &
         var_str ("unit_test"), is_known=.true.)
    call var_list_set_string (global%var_list, var_str ("$phs_method"), &
         var_str ("single"), is_known=.true.)
    call var_list_set_string (global%var_list, var_str ("$integration_method"),&
         var_str ("midpoint"), is_known=.true.)
    call var_list_set_log (global%var_list, var_str ("?vis_history"),&
         .false., is_known=.true.)        
    call var_list_set_log (global%var_list, var_str ("?integration_timer"),&
         .false., is_known = .true.)    
    call var_list_set_real (global%var_list, var_str ("sqrts"), &
         1000._default, is_known=.true.)

    allocate (lib)
    call lib%init (var_str ("lib_cmd8"))
    call global%add_prclib (lib)
    
    write (u, "(A)")  "*  Input file"
    write (u, "(A)")
    
    call ifile_append (ifile, 'model = "Test"')
    call ifile_append (ifile, 'process commands_8_p = s, s => s, s')
    call ifile_append (ifile, 'compile')
    call ifile_append (ifile, 'iterations = 1:1000')
    call ifile_append (ifile, 'integrate (commands_8_p)')
    call ifile_append (ifile, '?unweighted = false')
    call ifile_append (ifile, 'n_events = 3')
    call ifile_append (ifile, '?read_raw = false')
    call ifile_append (ifile, 'simulate (commands_8_p)')
    
    call ifile_write (ifile, u)

    write (u, "(A)")
    write (u, "(A)")  "*  Parse file"
    write (u, "(A)")
    
    call parse_ifile (ifile, pn_root, u)

    write (u, "(A)")
    write (u, "(A)")  "* Compile command list"
    write (u, "(A)")

    call command_list%compile (pn_root, global)
    call command_list%write (u)

    write (u, "(A)")
    write (u, "(A)")  "* Execute command list"

    call command_list%execute (global)

    write (u, "(A)")
    write (u, "(A)")  "* Cleanup"

    call ifile_final (ifile)

    call command_list%final ()
    call global%final ()
    call syntax_cmd_list_final ()
    call syntax_model_file_final ()

    write (u, "(A)")
    write (u, "(A)")  "* Test output end: commands_8"

  end subroutine commands_8

\eatline
\nwendcode{}\nwbegindocs{3061}\nwdocspar
\subsubsection{Define cuts}
Declare a cut expression.
\nwenddocs{}\nwbegincode{3062}\moddef{Commands: execute tests}\plusendmoddef
  call test (commands_9, "commands_9", &
       "cuts", &
       u, results)
\nwendcode{}\nwbegincode{3063}\moddef{Commands: tests}\plusendmoddef
  subroutine commands_9 (u)
    integer, intent(in) :: u
    type(ifile_t) :: ifile
    type(command_list_t), target :: command_list
    type(rt_data_t), target :: global
    type(parse_node_t), pointer :: pn_root
    type(string_t), dimension(0) :: no_vars

    write (u, "(A)")  "* Test output: commands_9"
    write (u, "(A)")  "*   Purpose: define cuts"
    write (u, "(A)")

    write (u, "(A)")  "*  Initialization"
    write (u, "(A)")

    call syntax_cmd_list_init ()
    call global%global_init ()

    write (u, "(A)")  "*  Input file"
    write (u, "(A)")
    
    call ifile_append (ifile, 'cuts = all Pt > 0 [particle]')
    
    call ifile_write (ifile, u)

    write (u, "(A)")
    write (u, "(A)")  "*  Parse file"
    write (u, "(A)")
    
    call parse_ifile (ifile, pn_root, u)

    write (u, "(A)")
    write (u, "(A)")  "* Compile command list"
    write (u, "(A)")

    call command_list%compile (pn_root, global)
    call command_list%write (u)

    write (u, "(A)")
    write (u, "(A)")  "* Execute command list"
    write (u, "(A)")

    call command_list%execute (global)

    call global%write (u, vars = no_vars)

    write (u, "(A)")
    write (u, "(A)")  "* Cleanup"

    call ifile_final (ifile)

    call command_list%final ()
    call global%final ()
    call syntax_cmd_list_final ()
    call syntax_model_file_final ()

    write (u, "(A)")
    write (u, "(A)")  "* Test output end: commands_9"

  end subroutine commands_9

\eatline
\nwendcode{}\nwbegindocs{3064}\nwdocspar
\subsubsection{Beams}
Define beam setup.
\nwenddocs{}\nwbegincode{3065}\moddef{Commands: execute tests}\plusendmoddef
  call test (commands_10, "commands_10", &
       "beams", &
       u, results)
\nwendcode{}\nwbegincode{3066}\moddef{Commands: tests}\plusendmoddef
  subroutine commands_10 (u)
    integer, intent(in) :: u
    type(ifile_t) :: ifile
    type(command_list_t), target :: command_list
    type(rt_data_t), target :: global
    type(parse_node_t), pointer :: pn_root

    write (u, "(A)")  "* Test output: commands_10"
    write (u, "(A)")  "*   Purpose: define beams"
    write (u, "(A)")

    write (u, "(A)")  "*  Initialization"
    write (u, "(A)")

    call syntax_cmd_list_init ()
    call syntax_model_file_init ()
    call global%global_init ()

    write (u, "(A)")  "*  Input file"
    write (u, "(A)")
    
    call ifile_append (ifile, 'model = QCD')
    call ifile_append (ifile, 'sqrts = 1000')
    call ifile_append (ifile, 'beams = p, p')
    
    call ifile_write (ifile, u)

    write (u, "(A)")
    write (u, "(A)")  "*  Parse file"
    write (u, "(A)")
    
    call parse_ifile (ifile, pn_root, u)

    write (u, "(A)")
    write (u, "(A)")  "* Compile command list"
    write (u, "(A)")

    call command_list%compile (pn_root, global)
    call command_list%write (u)

    write (u, "(A)")
    write (u, "(A)")  "* Execute command list"
    write (u, "(A)")

    call command_list%execute (global)

    call global%write_beams (u)

    write (u, "(A)")
    write (u, "(A)")  "* Cleanup"

    call ifile_final (ifile)

    call command_list%final ()
    call global%final ()
    call syntax_cmd_list_final ()
    call syntax_model_file_final ()

    write (u, "(A)")
    write (u, "(A)")  "* Test output end: commands_10"

  end subroutine commands_10

\eatline
\nwendcode{}\nwbegindocs{3067}\nwdocspar
\subsubsection{Structure functions}
Define beam setup with structure functions
\nwenddocs{}\nwbegincode{3068}\moddef{Commands: execute tests}\plusendmoddef
  call test (commands_11, "commands_11", &
       "structure functions", &
       u, results)
\nwendcode{}\nwbegincode{3069}\moddef{Commands: tests}\plusendmoddef
  subroutine commands_11 (u)
    integer, intent(in) :: u
    type(ifile_t) :: ifile
    type(command_list_t), target :: command_list
    type(rt_data_t), target :: global
    type(parse_node_t), pointer :: pn_root

    write (u, "(A)")  "* Test output: commands_11"
    write (u, "(A)")  "*   Purpose: define beams with structure functions"
    write (u, "(A)")

    write (u, "(A)")  "*  Initialization"
    write (u, "(A)")

    call syntax_cmd_list_init ()
    call syntax_model_file_init ()
    call global%global_init ()

    write (u, "(A)")  "*  Input file"
    write (u, "(A)")
    
    call ifile_append (ifile, 'model = QCD')
    call ifile_append (ifile, 'sqrts = 1100')
    call ifile_append (ifile, 'beams = p, p => lhapdf => pdf_builtin, isr')
    
    call ifile_write (ifile, u)

    write (u, "(A)")
    write (u, "(A)")  "*  Parse file"
    write (u, "(A)")
    
    call parse_ifile (ifile, pn_root, u)

    write (u, "(A)")
    write (u, "(A)")  "* Compile command list"
    write (u, "(A)")

    call command_list%compile (pn_root, global)
    call command_list%write (u)

    write (u, "(A)")
    write (u, "(A)")  "* Execute command list"
    write (u, "(A)")

    call command_list%execute (global)

    call global%write_beams (u)

    write (u, "(A)")
    write (u, "(A)")  "* Cleanup"

    call ifile_final (ifile)

    call command_list%final ()
    call global%final ()
    call syntax_cmd_list_final ()
    call syntax_model_file_final ()

    write (u, "(A)")
    write (u, "(A)")  "* Test output end: commands_11"

  end subroutine commands_11

\eatline
\nwendcode{}\nwbegindocs{3070}\nwdocspar
\subsubsection{Rescan events}
Read a model, then declare a process, compile the library, and
generate weighted events.  We take the
default ({\tt{}unit{\char95}test}) method and use the simplest methods of
phase-space parameterization and integration.  Then, rescan the
generated event sample.
\nwenddocs{}\nwbegincode{3071}\moddef{Commands: execute tests}\plusendmoddef
  call test (commands_12, "commands_12", &
       "event rescanning", &
       u, results)
\nwendcode{}\nwbegincode{3072}\moddef{Commands: tests}\plusendmoddef
  subroutine commands_12 (u)
    integer, intent(in) :: u
    type(ifile_t) :: ifile
    type(command_list_t), target :: command_list
    type(rt_data_t), target :: global
    type(parse_node_t), pointer :: pn_root
    type(prclib_entry_t), pointer :: lib

    write (u, "(A)")  "* Test output: commands_12"
    write (u, "(A)")  "*   Purpose: generate events and rescan"
    write (u, "(A)")

    write (u, "(A)")  "*  Initialization"
    write (u, "(A)")

    call syntax_cmd_list_init ()
    call syntax_model_file_init ()

    call global%global_init ()
    call var_list_append_log (global%var_list, &
         var_str ("?rebuild_phase_space"), .false., &
         intrinsic=.true.)
    call var_list_append_log (global%var_list, &
         var_str ("?rebuild_grids"), .false., &
         intrinsic=.true.)
    call global%init_fallback_model &
         (var_str ("SM_hadrons"), var_str ("SM_hadrons.mdl"))

    call var_list_set_string (global%var_list, var_str ("$method"), &
         var_str ("unit_test"), is_known=.true.)
    call var_list_set_string (global%var_list, var_str ("$phs_method"), &
         var_str ("single"), is_known=.true.)
    call var_list_set_string (global%var_list, var_str ("$integration_method"),&
         var_str ("midpoint"), is_known=.true.)
    call var_list_set_log (global%var_list, var_str ("?vis_history"),&
         .false., is_known=.true.)        
    call var_list_set_log (global%var_list, var_str ("?integration_timer"),&
         .false., is_known = .true.)    
    call var_list_set_real (global%var_list, var_str ("sqrts"), &
         1000._default, is_known=.true.)

    allocate (lib)
    call lib%init (var_str ("lib_cmd12"))
    call global%add_prclib (lib)
    
    write (u, "(A)")  "*  Input file"
    write (u, "(A)")
    
    call ifile_append (ifile, 'model = "Test"')
    call ifile_append (ifile, 'process commands_12_p = s, s => s, s')
    call ifile_append (ifile, 'compile')
    call ifile_append (ifile, 'iterations = 1:1000')
    call ifile_append (ifile, 'integrate (commands_12_p)')
    call ifile_append (ifile, '?unweighted = false')
    call ifile_append (ifile, 'n_events = 3')
    call ifile_append (ifile, '?read_raw = false')
    call ifile_append (ifile, 'simulate (commands_12_p)')
    call ifile_append (ifile, '?write_raw = false')
    call ifile_append (ifile, 'rescan "commands_12_p" (commands_12_p)')
    
    call ifile_write (ifile, u)

    write (u, "(A)")
    write (u, "(A)")  "*  Parse file"
    write (u, "(A)")
    
    call parse_ifile (ifile, pn_root, u)

    write (u, "(A)")
    write (u, "(A)")  "* Compile command list"
    write (u, "(A)")

    call command_list%compile (pn_root, global)
    call command_list%write (u)

    write (u, "(A)")
    write (u, "(A)")  "* Execute command list"

    call command_list%execute (global)

    write (u, "(A)")
    write (u, "(A)")  "* Cleanup"

    call ifile_final (ifile)

    call command_list%final ()
    call global%final ()
    call syntax_cmd_list_final ()
    call syntax_model_file_final ()

    write (u, "(A)")
    write (u, "(A)")  "* Test output end: commands_12"

  end subroutine commands_12

\eatline
\nwendcode{}\nwbegindocs{3073}\nwdocspar
\subsubsection{Event Files}
Set output formats for event files.
\nwenddocs{}\nwbegincode{3074}\moddef{Commands: execute tests}\plusendmoddef
  call test (commands_13, "commands_13", &
       "event output formats", &
       u, results)
\nwendcode{}\nwbegincode{3075}\moddef{Commands: tests}\plusendmoddef
  subroutine commands_13 (u)
    integer, intent(in) :: u
    type(ifile_t) :: ifile
    type(command_list_t), target :: command_list
    type(rt_data_t), target :: global
    type(parse_node_t), pointer :: pn_root
    type(prclib_entry_t), pointer :: lib
    logical :: exist

    write (u, "(A)")  "* Test output: commands_13"
    write (u, "(A)")  "*   Purpose: generate events and rescan"
    write (u, "(A)")

    write (u, "(A)")  "*  Initialization"
    write (u, "(A)")

    call syntax_cmd_list_init ()
    call syntax_model_file_init ()
    call global%global_init ()
    call global%init_fallback_model &
         (var_str ("SM_hadrons"), var_str ("SM_hadrons.mdl"))

    call var_list_set_string (global%var_list, var_str ("$method"), &
         var_str ("unit_test"), is_known=.true.)
    call var_list_set_string (global%var_list, var_str ("$phs_method"), &
         var_str ("single"), is_known=.true.)
    call var_list_set_string (global%var_list, var_str ("$integration_method"),&
         var_str ("midpoint"), is_known=.true.)
    call var_list_set_real (global%var_list, var_str ("sqrts"), &
         1000._default, is_known=.true.)
    call var_list_set_log (global%var_list, var_str ("?vis_history"),&
         .false., is_known=.true.)    
    call var_list_set_log (global%var_list, var_str ("?integration_timer"),&
         .false., is_known = .true.)    

    allocate (lib)
    call lib%init (var_str ("lib_cmd13"))
    call global%add_prclib (lib)
    
    write (u, "(A)")  "*  Input file"
    write (u, "(A)")
    
    call ifile_append (ifile, 'model = "Test"')
    call ifile_append (ifile, 'process commands_13_p = s, s => s, s')
    call ifile_append (ifile, 'compile')
    call ifile_append (ifile, 'iterations = 1:1000')
    call ifile_append (ifile, 'integrate (commands_13_p)')
    call ifile_append (ifile, '?unweighted = false')
    call ifile_append (ifile, 'n_events = 1')
    call ifile_append (ifile, '?read_raw = false')
    call ifile_append (ifile, 'sample_format = weight_stream')
    call ifile_append (ifile, 'simulate (commands_13_p)')
    
    call ifile_write (ifile, u)

    write (u, "(A)")
    write (u, "(A)")  "*  Parse file"
    write (u, "(A)")
    
    call parse_ifile (ifile, pn_root, u)

    write (u, "(A)")
    write (u, "(A)")  "* Compile command list"
    write (u, "(A)")

    call command_list%compile (pn_root, global)
    call command_list%write (u)

    write (u, "(A)")
    write (u, "(A)")  "* Execute command list"

    call command_list%execute (global)

    write (u, "(A)")
    write (u, "(A)")  "* Verify output files"
    write (u, "(A)")

    inquire (file = "commands_13_p.evx", exist = exist)
    if (exist)  write (u, "(1x,A)")  "raw"

    inquire (file = "commands_13_p.weights.dat", exist = exist)
    if (exist)  write (u, "(1x,A)")  "weight_stream"

    write (u, "(A)")
    write (u, "(A)")  "* Cleanup"

    call ifile_final (ifile)

    call command_list%final ()
    call global%final ()
    call syntax_cmd_list_final ()
    call syntax_model_file_final ()

    write (u, "(A)")
    write (u, "(A)")  "* Test output end: commands_13"

  end subroutine commands_13

\eatline
\nwendcode{}\nwbegindocs{3076}\nwdocspar
\subsubsection{Compile Empty Libraries}
(This is a regression test:)  Declare two empty libraries and compile them.
\nwenddocs{}\nwbegincode{3077}\moddef{Commands: execute tests}\plusendmoddef
  call test (commands_14, "commands_14", &
       "empty libraries", &
       u, results)
\nwendcode{}\nwbegincode{3078}\moddef{Commands: tests}\plusendmoddef
  subroutine commands_14 (u)
    integer, intent(in) :: u
    type(ifile_t) :: ifile
    type(command_list_t), target :: command_list
    type(rt_data_t), target :: global
    type(parse_node_t), pointer :: pn_root

    write (u, "(A)")  "* Test output: commands_14"
    write (u, "(A)")  "*   Purpose: define and compile empty libraries"
    write (u, "(A)")

    write (u, "(A)")  "* Initialization"
    write (u, "(A)")

    call syntax_model_file_init ()
    call syntax_cmd_list_init ()

    call global%global_init ()

    write (u, "(A)")  "* Input file"
    write (u, "(A)")
    
    call ifile_append (ifile, 'model = "Test"')
    call ifile_append (ifile, 'library = "lib1"')
    call ifile_append (ifile, 'library = "lib2"')
    call ifile_append (ifile, 'compile ()')
    
    call ifile_write (ifile, u)

    write (u, "(A)")
    write (u, "(A)")  "* Parse file"
    write (u, "(A)")
    
    call parse_ifile (ifile, pn_root)

    write (u, "(A)")  "* Compile command list"
    write (u, "(A)")

    call command_list%compile (pn_root, global)

    write (u, "(A)")  "* Execute command list"
    write (u, "(A)")

    call command_list%execute (global)

    call global%prclib_stack%write (u)

    write (u, "(A)")
    write (u, "(A)")  "* Cleanup"

    call ifile_final (ifile)

    call command_list%final ()
    call global%final ()

    call syntax_cmd_list_final ()
    call syntax_model_file_final ()

    write (u, "(A)")
    write (u, "(A)")  "* Test output end: commands_14"

  end subroutine commands_14

\eatline
\nwendcode{}\nwbegindocs{3079}\nwdocspar
\subsubsection{Compile Process}
Read a model, then declare a process and compile the library.  The process
library is allocated explicitly.  For the process definition, We take the
default ({\tt{}unit{\char95}test}) method.  There is no external code, so compilation of
the library is merely a formal status change.
\nwenddocs{}\nwbegincode{3080}\moddef{Commands: execute tests}\plusendmoddef
  call test (commands_15, "commands_15", &
       "compilation", &
       u, results)
\nwendcode{}\nwbegincode{3081}\moddef{Commands: tests}\plusendmoddef
  subroutine commands_15 (u)
    integer, intent(in) :: u
    type(ifile_t) :: ifile
    type(command_list_t), target :: command_list
    type(rt_data_t), target :: global
    type(parse_node_t), pointer :: pn_root
    type(prclib_entry_t), pointer :: lib

    write (u, "(A)")  "* Test output: commands_15"
    write (u, "(A)")  "*   Purpose: define process and compile library"
    write (u, "(A)")

    write (u, "(A)")  "* Initialization"
    write (u, "(A)")

    call syntax_cmd_list_init ()
    call syntax_model_file_init ()
    call global%global_init ()
    call var_list_set_string (global%var_list, var_str ("$method"), &
         var_str ("unit_test"), is_known=.true.)
    call var_list_set_string (global%var_list, var_str ("$phs_method"), &
         var_str ("single"), is_known=.true.)
    call var_list_set_string (global%var_list, var_str ("$integration_method"),&
         var_str ("midpoint"), is_known=.true.)
    call var_list_set_real (global%var_list, var_str ("sqrts"), &
         1000._default, is_known=.true.)
    call var_list_set_log (global%var_list, var_str ("?vis_history"),&
         .false., is_known=.true.)    
    call var_list_set_log (global%var_list, var_str ("?integration_timer"),&
         .false., is_known = .true.)    
    
    allocate (lib)
    call lib%init (var_str ("lib_cmd15"))
    call global%add_prclib (lib)
    
    write (u, "(A)")  "* Input file"
    write (u, "(A)")
    
    call ifile_append (ifile, 'model = "Test"')
    call ifile_append (ifile, 'process t15 = s, s => s, s')
    call ifile_append (ifile, 'iterations = 1:1000')
    call ifile_append (ifile, 'integrate (t15)')
    
    call ifile_write (ifile, u)

    write (u, "(A)")
    write (u, "(A)")  "* Parse file"
    write (u, "(A)")
    
    call parse_ifile (ifile, pn_root)

    write (u, "(A)")  "* Compile command list"
    write (u, "(A)")

    call command_list%compile (pn_root, global)

    write (u, "(A)")  "* Execute command list"
    write (u, "(A)")

    call command_list%execute (global)

    call global%prclib_stack%write (u)

    write (u, "(A)")
    write (u, "(A)")  "* Cleanup"

    call ifile_final (ifile)

    call command_list%final ()
    call global%final ()
    call syntax_cmd_list_final ()
    call syntax_model_file_final ()

    write (u, "(A)")
    write (u, "(A)")  "* Test output end: commands_15"

  end subroutine commands_15

\eatline
\nwendcode{}\nwbegindocs{3082}\nwdocspar
\subsubsection{Observable}
Declare an observable, fill it and display.
\nwenddocs{}\nwbegincode{3083}\moddef{Commands: execute tests}\plusendmoddef
  call test (commands_16, "commands_16", &
       "observables", &
       u, results)
\nwendcode{}\nwbegincode{3084}\moddef{Commands: tests}\plusendmoddef
  subroutine commands_16 (u)
    integer, intent(in) :: u
    type(ifile_t) :: ifile
    type(command_list_t), target :: command_list
    type(rt_data_t), target :: global
    type(parse_node_t), pointer :: pn_root

    write (u, "(A)")  "* Test output: commands_16"
    write (u, "(A)")  "*   Purpose: declare an observable"
    write (u, "(A)")

    write (u, "(A)")  "* Initialization"
    write (u, "(A)")

    call syntax_cmd_list_init ()
    call global%global_init ()
    
    write (u, "(A)")  "* Input file"
    write (u, "(A)")
    
    call ifile_append (ifile, '$obs_label = "foo"')
    call ifile_append (ifile, '$obs_unit = "cm"')
    call ifile_append (ifile, '$title = "Observable foo"')
    call ifile_append (ifile, '$description = "This is observable foo"')
    call ifile_append (ifile, 'observable foo')
    
    call ifile_write (ifile, u)

    write (u, "(A)")
    write (u, "(A)")  "* Parse file"
    write (u, "(A)")
    
    call parse_ifile (ifile, pn_root)

    write (u, "(A)")  "* Compile command list"
    write (u, "(A)")

    call command_list%compile (pn_root, global)

    call command_list%write (u)

    write (u, "(A)")
    write (u, "(A)")  "* Execute command list"
    write (u, "(A)")

    call command_list%execute (global)

    write (u, "(A)")  "* Record two data items"
    write (u, "(A)")

    call analysis_record_data (var_str ("foo"), 1._default)
    call analysis_record_data (var_str ("foo"), 3._default)

    write (u, "(A)")  "* Display analysis store"
    write (u, "(A)")

    call analysis_write (u, verbose=.true.)

    write (u, "(A)")
    write (u, "(A)")  "* Cleanup"

    call ifile_final (ifile)

    call analysis_final ()
    call command_list%final ()
    call global%final ()
    call syntax_cmd_list_final ()
    call syntax_model_file_final ()

    write (u, "(A)")
    write (u, "(A)")  "* Test output end: commands_16"

  end subroutine commands_16

\eatline
\nwendcode{}\nwbegindocs{3085}\nwdocspar
\subsubsection{Histogram}
Declare a histogram, fill it and display.
\nwenddocs{}\nwbegincode{3086}\moddef{Commands: execute tests}\plusendmoddef
  call test (commands_17, "commands_17", &
       "histograms", &
       u, results)
\nwendcode{}\nwbegincode{3087}\moddef{Commands: tests}\plusendmoddef
  subroutine commands_17 (u)
    integer, intent(in) :: u
    type(ifile_t) :: ifile
    type(command_list_t), target :: command_list
    type(rt_data_t), target :: global
    type(parse_node_t), pointer :: pn_root
    type(string_t), dimension(3) :: name
    integer :: i

    write (u, "(A)")  "* Test output: commands_17"
    write (u, "(A)")  "*   Purpose: declare histograms"
    write (u, "(A)")

    write (u, "(A)")  "* Initialization"
    write (u, "(A)")

    call syntax_cmd_list_init ()
    call global%global_init ()
    
    write (u, "(A)")  "* Input file"
    write (u, "(A)")
    
    call ifile_append (ifile, '$obs_label = "foo"')
    call ifile_append (ifile, '$obs_unit = "cm"')
    call ifile_append (ifile, '$title = "Histogram foo"')
    call ifile_append (ifile, '$description = "This is histogram foo"')
    call ifile_append (ifile, 'histogram foo (0,5,1)')
    call ifile_append (ifile, '$title = "Histogram bar"')
    call ifile_append (ifile, '$description = "This is histogram bar"')
    call ifile_append (ifile, 'n_bins = 2')
    call ifile_append (ifile, 'histogram bar (0,5)')
    call ifile_append (ifile, '$title = "Histogram gee"')
    call ifile_append (ifile, '$description = "This is histogram gee"')
    call ifile_append (ifile, '?normalize_bins = true')
    call ifile_append (ifile, 'histogram gee (0,5)')
    
    call ifile_write (ifile, u)

    write (u, "(A)")
    write (u, "(A)")  "* Parse file"
    write (u, "(A)")
    
    call parse_ifile (ifile, pn_root)

    write (u, "(A)")  "* Compile command list"
    write (u, "(A)")

    call command_list%compile (pn_root, global)

    call command_list%write (u)

    write (u, "(A)")
    write (u, "(A)")  "* Execute command list"
    write (u, "(A)")

    call command_list%execute (global)

    write (u, "(A)")  "* Record two data items"
    write (u, "(A)")

    name(1) = "foo"
    name(2) = "bar"
    name(3) = "gee"
    
    do i = 1, 3
       call analysis_record_data (name(i), 0.1_default, &
            weight = 0.25_default)
       call analysis_record_data (name(i), 3.1_default)
       call analysis_record_data (name(i), 4.1_default, &
            excess = 0.5_default)
       call analysis_record_data (name(i), 7.1_default)
    end do

    write (u, "(A)")  "* Display analysis store"
    write (u, "(A)")

    call analysis_write (u, verbose=.true.)

    write (u, "(A)")
    write (u, "(A)")  "* Cleanup"

    call ifile_final (ifile)

    call analysis_final ()
    call command_list%final ()
    call global%final ()
    call syntax_cmd_list_final ()
    call syntax_model_file_final ()

    write (u, "(A)")
    write (u, "(A)")  "* Test output end: commands_17"

  end subroutine commands_17

\eatline
\nwendcode{}\nwbegindocs{3088}\nwdocspar
\subsubsection{Plot}
Declare a plot, fill it and display contents.
\nwenddocs{}\nwbegincode{3089}\moddef{Commands: execute tests}\plusendmoddef
  call test (commands_18, "commands_18", &
       "plots", &
       u, results)
\nwendcode{}\nwbegincode{3090}\moddef{Commands: tests}\plusendmoddef
  subroutine commands_18 (u)
    integer, intent(in) :: u
    type(ifile_t) :: ifile
    type(command_list_t), target :: command_list
    type(rt_data_t), target :: global
    type(parse_node_t), pointer :: pn_root

    write (u, "(A)")  "* Test output: commands_18"
    write (u, "(A)")  "*   Purpose: declare a plot"
    write (u, "(A)")

    write (u, "(A)")  "* Initialization"
    write (u, "(A)")

    call syntax_cmd_list_init ()
    call global%global_init ()
    
    write (u, "(A)")  "* Input file"
    write (u, "(A)")
    
    call ifile_append (ifile, '$obs_label = "foo"')
    call ifile_append (ifile, '$obs_unit = "cm"')
    call ifile_append (ifile, '$title = "Plot foo"')
    call ifile_append (ifile, '$description = "This is plot foo"')
    call ifile_append (ifile, '$x_label = "x axis"')
    call ifile_append (ifile, '$y_label = "y axis"')
    call ifile_append (ifile, '?x_log = false')
    call ifile_append (ifile, '?y_log = true')
    call ifile_append (ifile, 'x_min = -1')
    call ifile_append (ifile, 'x_max = 1')
    call ifile_append (ifile, 'y_min = 0.1')
    call ifile_append (ifile, 'y_max = 1000')
    call ifile_append (ifile, 'plot foo')
    
    call ifile_write (ifile, u)

    write (u, "(A)")
    write (u, "(A)")  "* Parse file"
    write (u, "(A)")
    
    call parse_ifile (ifile, pn_root)

    write (u, "(A)")  "* Compile command list"
    write (u, "(A)")

    call command_list%compile (pn_root, global)

    call command_list%write (u)

    write (u, "(A)")
    write (u, "(A)")  "* Execute command list"
    write (u, "(A)")

    call command_list%execute (global)

    write (u, "(A)")  "* Record two data items"
    write (u, "(A)")

    call analysis_record_data (var_str ("foo"), 0._default, 20._default, &
         xerr = 0.25_default)
    call analysis_record_data (var_str ("foo"), 0.5_default, 0.2_default, &
         yerr = 0.07_default)
    call analysis_record_data (var_str ("foo"), 3._default, 2._default)

    write (u, "(A)")  "* Display analysis store"
    write (u, "(A)")

    call analysis_write (u, verbose=.true.)

    write (u, "(A)")
    write (u, "(A)")  "* Cleanup"

    call ifile_final (ifile)

    call analysis_final ()
    call command_list%final ()
    call global%final ()
    call syntax_cmd_list_final ()
    call syntax_model_file_final ()

    write (u, "(A)")
    write (u, "(A)")  "* Test output end: commands_18"

  end subroutine commands_18

\eatline
\nwendcode{}\nwbegindocs{3091}\nwdocspar
\subsubsection{Graph}
Combine two (empty) plots to a graph.
\nwenddocs{}\nwbegincode{3092}\moddef{Commands: execute tests}\plusendmoddef
  call test (commands_19, "commands_19", &
       "graphs", &
       u, results)
\nwendcode{}\nwbegincode{3093}\moddef{Commands: tests}\plusendmoddef
  subroutine commands_19 (u)
    integer, intent(in) :: u
    type(ifile_t) :: ifile
    type(command_list_t), target :: command_list
    type(rt_data_t), target :: global
    type(parse_node_t), pointer :: pn_root

    write (u, "(A)")  "* Test output: commands_19"
    write (u, "(A)")  "*   Purpose: combine two plots to a graph"
    write (u, "(A)")

    write (u, "(A)")  "* Initialization"
    write (u, "(A)")

    call syntax_cmd_list_init ()
    call global%global_init ()
    
    write (u, "(A)")  "* Input file"
    write (u, "(A)")
    
    call ifile_append (ifile, 'plot a')
    call ifile_append (ifile, 'plot b')
    call ifile_append (ifile, '$title = "Graph foo"')
    call ifile_append (ifile, '$description = "This is graph foo"')
    call ifile_append (ifile, 'graph foo = a & b')
    
    call ifile_write (ifile, u)

    write (u, "(A)")
    write (u, "(A)")  "* Parse file"
    write (u, "(A)")
    
    call parse_ifile (ifile, pn_root)

    write (u, "(A)")  "* Compile command list"
    write (u, "(A)")

    call command_list%compile (pn_root, global)

    call command_list%write (u)

    write (u, "(A)")
    write (u, "(A)")  "* Execute command list"
    write (u, "(A)")

    call command_list%execute (global)

    write (u, "(A)")  "* Display analysis object"
    write (u, "(A)")

    call analysis_write (var_str ("foo"), u)

    write (u, "(A)")
    write (u, "(A)")  "* Cleanup"

    call ifile_final (ifile)

    call analysis_final ()
    call command_list%final ()
    call global%final ()
    call syntax_cmd_list_final ()
    call syntax_model_file_final ()

    write (u, "(A)")
    write (u, "(A)")  "* Test output end: commands_19"

  end subroutine commands_19

\eatline
\nwendcode{}\nwbegindocs{3094}\nwdocspar
\subsubsection{Record Data}
Record data in previously allocated analysis objects.
\nwenddocs{}\nwbegincode{3095}\moddef{Commands: execute tests}\plusendmoddef
  call test (commands_20, "commands_20", &
       "record data", &
       u, results)
\nwendcode{}\nwbegincode{3096}\moddef{Commands: tests}\plusendmoddef
  subroutine commands_20 (u)
    integer, intent(in) :: u
    type(ifile_t) :: ifile
    type(command_list_t), target :: command_list
    type(rt_data_t), target :: global
    type(parse_node_t), pointer :: pn_root

    write (u, "(A)")  "* Test output: commands_20"
    write (u, "(A)")  "*   Purpose: record data"
    write (u, "(A)")

    write (u, "(A)")  "* Initialization: create observable, histogram, plot"
    write (u, "(A)")

    call syntax_cmd_list_init ()
    call global%global_init ()
    
    call analysis_init_observable (var_str ("o"))
    call analysis_init_histogram (var_str ("h"), 0._default, 1._default, 3, &
         normalize_bins = .false.)
    call analysis_init_plot (var_str ("p"))
    
    write (u, "(A)")  "* Input file"
    write (u, "(A)")
    
    call ifile_append (ifile, 'record o (1.234)')
    call ifile_append (ifile, 'record h (0.5)')
    call ifile_append (ifile, 'record p (1, 2)')
    
    call ifile_write (ifile, u)

    write (u, "(A)")
    write (u, "(A)")  "* Parse file"
    write (u, "(A)")
    
    call parse_ifile (ifile, pn_root)

    write (u, "(A)")  "* Compile command list"
    write (u, "(A)")

    call command_list%compile (pn_root, global)

    call command_list%write (u)

    write (u, "(A)")
    write (u, "(A)")  "* Execute command list"
    write (u, "(A)")

    call command_list%execute (global)

    write (u, "(A)")  "* Display analysis object"
    write (u, "(A)")

    call analysis_write (u, verbose = .true.)

    write (u, "(A)")
    write (u, "(A)")  "* Cleanup"

    call ifile_final (ifile)

    call analysis_final ()
    call command_list%final ()
    call global%final ()
    call syntax_cmd_list_final ()
    call syntax_model_file_final ()

    write (u, "(A)")
    write (u, "(A)")  "* Test output end: commands_20"

  end subroutine commands_20

\eatline
\nwendcode{}\nwbegindocs{3097}\nwdocspar
\subsubsection{Analysis}
Declare an analysis expression and use it to fill an observable during
event generation.
\nwenddocs{}\nwbegincode{3098}\moddef{Commands: execute tests}\plusendmoddef
  call test (commands_21, "commands_21", &
       "analysis expression", &
       u, results)
\nwendcode{}\nwbegincode{3099}\moddef{Commands: tests}\plusendmoddef
  subroutine commands_21 (u)
    integer, intent(in) :: u
    type(ifile_t) :: ifile
    type(command_list_t), target :: command_list
    type(rt_data_t), target :: global
    type(parse_node_t), pointer :: pn_root
    type(prclib_entry_t), pointer :: lib

    write (u, "(A)")  "* Test output: commands_21"
    write (u, "(A)")  "*   Purpose: create and use analysis expression"
    write (u, "(A)")

    write (u, "(A)")  "* Initialization: create observable"
    write (u, "(A)")

    call syntax_cmd_list_init ()
    call syntax_model_file_init ()
    call global%global_init ()
    call global%init_fallback_model &
         (var_str ("SM_hadrons"), var_str ("SM_hadrons.mdl"))

    call var_list_set_string (global%var_list, var_str ("$method"), &
         var_str ("unit_test"), is_known=.true.)
    call var_list_set_string (global%var_list, var_str ("$phs_method"), &
         var_str ("single"), is_known=.true.)
    call var_list_set_string (global%var_list, var_str ("$integration_method"),&
         var_str ("midpoint"), is_known=.true.)
    call var_list_set_log (global%var_list, var_str ("?vis_history"),&
         .false., is_known=.true.)        
    call var_list_set_log (global%var_list, var_str ("?integration_timer"),&
         .false., is_known = .true.)    
    call var_list_set_real (global%var_list, var_str ("sqrts"), &
         1000._default, is_known=.true.)

    allocate (lib)
    call lib%init (var_str ("lib_cmd8"))
    call global%add_prclib (lib)
    
    call analysis_init_observable (var_str ("m"))
    
    write (u, "(A)")  "* Input file"
    write (u, "(A)")
    
    call ifile_append (ifile, 'model = "Test"')
    call ifile_append (ifile, 'process commands_21_p = s, s => s, s')
    call ifile_append (ifile, 'compile')
    call ifile_append (ifile, 'iterations = 1:100')
    call ifile_append (ifile, 'integrate (commands_21_p)')
    call ifile_append (ifile, '?unweighted = true')
    call ifile_append (ifile, 'n_events = 3')
    call ifile_append (ifile, '?read_raw = false')
    call ifile_append (ifile, 'observable m')
    call ifile_append (ifile, 'analysis = record m (eval M [s])')
    call ifile_append (ifile, 'simulate (commands_21_p)')
    
    call ifile_write (ifile, u)

    write (u, "(A)")
    write (u, "(A)")  "* Parse file"
    write (u, "(A)")
    
    call parse_ifile (ifile, pn_root)

    write (u, "(A)")  "* Compile command list"
    write (u, "(A)")

    call command_list%compile (pn_root, global)

    call command_list%write (u)

    write (u, "(A)")
    write (u, "(A)")  "* Execute command list"
    write (u, "(A)")

    call command_list%execute (global)

    write (u, "(A)")  "* Display analysis object"
    write (u, "(A)")

    call analysis_write (u)

    write (u, "(A)")
    write (u, "(A)")  "* Cleanup"

    call ifile_final (ifile)

    call analysis_final ()
    call command_list%final ()
    call global%final ()
    call syntax_cmd_list_final ()
    call syntax_model_file_final ()

    write (u, "(A)")
    write (u, "(A)")  "* Test output end: commands_21"

  end subroutine commands_21

\eatline
\nwendcode{}\nwbegindocs{3100}\nwdocspar
\subsubsection{Write Analysis}
Write accumulated analysis data to file.
\nwenddocs{}\nwbegincode{3101}\moddef{Commands: execute tests}\plusendmoddef
  call test (commands_22, "commands_22", &
       "write analysis", &
       u, results)
\nwendcode{}\nwbegincode{3102}\moddef{Commands: tests}\plusendmoddef
  subroutine commands_22 (u)
    integer, intent(in) :: u
    type(ifile_t) :: ifile
    type(command_list_t), target :: command_list
    type(rt_data_t), target :: global
    type(parse_node_t), pointer :: pn_root
    integer :: u_file, iostat
    logical :: exist
    character(80) :: buffer

    write (u, "(A)")  "* Test output: commands_22"
    write (u, "(A)")  "*   Purpose: write analysis data"
    write (u, "(A)")

    write (u, "(A)")  "* Initialization: create observable"
    write (u, "(A)")

    call syntax_cmd_list_init ()
    call global%global_init ()
    
    call analysis_init_observable (var_str ("m"))
    call analysis_record_data (var_str ("m"), 125._default)
    
    write (u, "(A)")  "* Input file"
    write (u, "(A)")
    
    call ifile_append (ifile, '$out_file = "commands_22.dat"')
    call ifile_append (ifile, 'write_analysis')
    
    call ifile_write (ifile, u)

    write (u, "(A)")
    write (u, "(A)")  "* Parse file"
    write (u, "(A)")
    
    call parse_ifile (ifile, pn_root)

    write (u, "(A)")  "* Compile command list"
    write (u, "(A)")

    call command_list%compile (pn_root, global)

    call command_list%write (u)

    write (u, "(A)")
    write (u, "(A)")  "* Execute command list"
    write (u, "(A)")

    call command_list%execute (global)

    write (u, "(A)")  "* Display analysis data"
    write (u, "(A)")

    inquire (file = "commands_22.dat", exist = exist)
    if (.not. exist) then
       write (u, "(A)")  "ERROR: File commands_22.dat not found"
       return
    end if
    
    u_file = free_unit ()
    open (u_file, file = "commands_22.dat", &
         action = "read", status = "old")
    do
       read (u_file, "(A)", iostat = iostat)  buffer
       if (iostat /= 0)  exit
       write (u, "(A)") trim (buffer)
    end do
    close (u_file)

    write (u, "(A)")
    write (u, "(A)")  "* Cleanup"

    call ifile_final (ifile)

    call analysis_final ()
    call command_list%final ()
    call global%final ()
    call syntax_cmd_list_final ()
    call syntax_model_file_final ()

    write (u, "(A)")
    write (u, "(A)")  "* Test output end: commands_22"

  end subroutine commands_22

\eatline
\nwendcode{}\nwbegindocs{3103}\nwdocspar
\subsubsection{Compile Analysis}
Write accumulated analysis data to file and compile.
\nwenddocs{}\nwbegincode{3104}\moddef{Commands: execute tests}\plusendmoddef
  call test (commands_23, "commands_23", &
       "compile analysis", &
       u, results)
\nwendcode{}\nwbegincode{3105}\moddef{Commands: tests}\plusendmoddef
  subroutine commands_23 (u)
    integer, intent(in) :: u
    type(ifile_t) :: ifile
    type(command_list_t), target :: command_list
    type(rt_data_t), target :: global
    type(parse_node_t), pointer :: pn_root
    integer :: u_file, iostat
    character(256) :: buffer
    logical :: exist
    type(graph_options_t) :: graph_options

    write (u, "(A)")  "* Test output: commands_23"
    write (u, "(A)")  "*   Purpose: write and compile analysis data"
    write (u, "(A)")

    write (u, "(A)")  "* Initialization: create and fill histogram"
    write (u, "(A)")

    call syntax_cmd_list_init ()
    call global%global_init ()
    
    call graph_options_init (graph_options)
    call graph_options_set (graph_options, &
         title = var_str ("Histogram for test: commands 23"), &
         description = var_str ("This is a test."), &
         width_mm = 125, height_mm = 85)
    call analysis_init_histogram (var_str ("h"), &
         0._default, 10._default, 2._default, .false., &
         graph_options = graph_options)
    call analysis_record_data (var_str ("h"), 1._default)
    call analysis_record_data (var_str ("h"), 1._default)
    call analysis_record_data (var_str ("h"), 1._default)
    call analysis_record_data (var_str ("h"), 1._default)
    call analysis_record_data (var_str ("h"), 3._default)
    call analysis_record_data (var_str ("h"), 3._default)
    call analysis_record_data (var_str ("h"), 3._default)
    call analysis_record_data (var_str ("h"), 5._default)
    call analysis_record_data (var_str ("h"), 7._default)
    call analysis_record_data (var_str ("h"), 7._default)
    call analysis_record_data (var_str ("h"), 7._default)
    call analysis_record_data (var_str ("h"), 7._default)
    call analysis_record_data (var_str ("h"), 9._default)
    call analysis_record_data (var_str ("h"), 9._default)
    call analysis_record_data (var_str ("h"), 9._default)
    call analysis_record_data (var_str ("h"), 9._default)
    call analysis_record_data (var_str ("h"), 9._default)
    call analysis_record_data (var_str ("h"), 9._default)
    call analysis_record_data (var_str ("h"), 9._default)
    
    write (u, "(A)")  "* Input file"
    write (u, "(A)")
    
    call ifile_append (ifile, '$out_file = "commands_23.dat"')
    call ifile_append (ifile, 'compile_analysis')
    
    call ifile_write (ifile, u)

    write (u, "(A)")
    write (u, "(A)")  "* Parse file"
    write (u, "(A)")
    
    call parse_ifile (ifile, pn_root)

    write (u, "(A)")  "* Compile command list"
    write (u, "(A)")

    call command_list%compile (pn_root, global)

    call command_list%write (u)

    write (u, "(A)")
    write (u, "(A)")  "* Delete Postscript output"
    write (u, "(A)")
    
    inquire (file = "commands_23.ps", exist = exist)
    if (exist) then
       u_file = free_unit ()
       open (u_file, file = "commands_23.ps", action = "write", status = "old")
       close (u_file, status = "delete")
    end if
    inquire (file = "commands_23.ps", exist = exist)
    write (u, "(1x,A,L1)")  "Postcript output exists = ", exist

    write (u, "(A)")
    write (u, "(A)")  "* Execute command list"
    write (u, "(A)")

    call command_list%execute (global)

    write (u, "(A)")  "* TeX file"
    write (u, "(A)")

    inquire (file = "commands_23.tex", exist = exist)
    if (.not. exist) then
       write (u, "(A)")  "ERROR: File commands_23.tex not found"
       return
    end if
    
    u_file = free_unit ()
    open (u_file, file = "commands_23.tex", &
         action = "read", status = "old")
    do
       read (u_file, "(A)", iostat = iostat)  buffer
       if (iostat /= 0)  exit
       write (u, "(A)") trim (buffer)
    end do
    close (u_file)
    write (u, *)
    
    inquire (file = "commands_23.ps", exist = exist)
    write (u, "(1x,A,L1)")  "Postcript output exists = ", exist

    write (u, "(A)")
    write (u, "(A)")  "* Cleanup"

    call ifile_final (ifile)

    call analysis_final ()
    call command_list%final ()
    call global%final ()
    call syntax_cmd_list_final ()
    call syntax_model_file_final ()

    write (u, "(A)")
    write (u, "(A)")  "* Test output end: commands_23"

  end subroutine commands_23

\eatline
\nwendcode{}\nwbegindocs{3106}\nwdocspar
\subsubsection{Histogram}
Declare a histogram, fill it and display.
\nwenddocs{}\nwbegincode{3107}\moddef{Commands: execute tests}\plusendmoddef
  call test (commands_24, "commands_24", &
       "drawing options", &
       u, results)
\nwendcode{}\nwbegincode{3108}\moddef{Commands: tests}\plusendmoddef
  subroutine commands_24 (u)
    integer, intent(in) :: u
    type(ifile_t) :: ifile
    type(command_list_t), target :: command_list
    type(rt_data_t), target :: global
    type(parse_node_t), pointer :: pn_root

    write (u, "(A)")  "* Test output: commands_24"
    write (u, "(A)")  "*   Purpose: check graph and drawing options"
    write (u, "(A)")

    write (u, "(A)")  "* Initialization"
    write (u, "(A)")

    call syntax_cmd_list_init ()
    call global%global_init ()
    
    write (u, "(A)")  "* Input file"
    write (u, "(A)")
    
    call ifile_append (ifile, '$title = "Title"')
    call ifile_append (ifile, '$description = "Description"')
    call ifile_append (ifile, '$x_label = "X Label"')
    call ifile_append (ifile, '$y_label = "Y Label"')
    call ifile_append (ifile, 'graph_width_mm = 111')
    call ifile_append (ifile, 'graph_height_mm = 222')
    call ifile_append (ifile, 'x_min = -11')
    call ifile_append (ifile, 'x_max = 22')
    call ifile_append (ifile, 'y_min = -33')
    call ifile_append (ifile, 'y_max = 44')
    call ifile_append (ifile, '$gmlcode_bg = "GML Code BG"')
    call ifile_append (ifile, '$gmlcode_fg = "GML Code FG"')
    call ifile_append (ifile, '$fill_options = "Fill Options"')
    call ifile_append (ifile, '$draw_options = "Draw Options"')
    call ifile_append (ifile, '$err_options = "Error Options"')
    call ifile_append (ifile, '$symbol = "Symbol"')
    call ifile_append (ifile, 'histogram foo (0,1)')
    call ifile_append (ifile, 'plot bar')
    
    call ifile_write (ifile, u)

    write (u, "(A)")
    write (u, "(A)")  "* Parse file"
    write (u, "(A)")
    
    call parse_ifile (ifile, pn_root)

    write (u, "(A)")  "* Compile command list"
    write (u, "(A)")

    call command_list%compile (pn_root, global)

    call command_list%write (u)

    write (u, "(A)")
    write (u, "(A)")  "* Execute command list"
    write (u, "(A)")

    call command_list%execute (global)

    write (u, "(A)")  "* Display analysis store"
    write (u, "(A)")

    call analysis_write (u, verbose=.true.)

    write (u, "(A)")
    write (u, "(A)")  "* Cleanup"

    call ifile_final (ifile)

    call analysis_final ()
    call command_list%final ()
    call global%final ()
    call syntax_cmd_list_final ()
    call syntax_model_file_final ()

    write (u, "(A)")
    write (u, "(A)")  "* Test output end: commands_24"

  end subroutine commands_24

\eatline
\nwendcode{}\nwbegindocs{3109}\nwdocspar
\subsubsection{Local Environment}
Declare a local environment.
\nwenddocs{}\nwbegincode{3110}\moddef{Commands: execute tests}\plusendmoddef
  call test (commands_25, "commands_25", &
       "local process environment", &
       u, results)
\nwendcode{}\nwbegincode{3111}\moddef{Commands: tests}\plusendmoddef
  subroutine commands_25 (u)
    integer, intent(in) :: u
    type(ifile_t) :: ifile
    type(command_list_t), target :: command_list
    type(rt_data_t), target :: global
    type(parse_node_t), pointer :: pn_root

    write (u, "(A)")  "* Test output: commands_25"
    write (u, "(A)")  "*   Purpose: declare local environment for process"
    write (u, "(A)")

    call syntax_model_file_init ()
    call syntax_cmd_list_init ()
    call global%global_init ()
    call var_list_set_log (global%var_list, var_str ("?omega_openmp"), &
         .false., is_known = .true.)
    
    write (u, "(A)")  "* Input file"
    write (u, "(A)")
    
    call ifile_append (ifile, 'library = "commands_25_lib"')
    call ifile_append (ifile, 'model = "Test"')
    call ifile_append (ifile, 'process commands_25_p1 = g, g => g, g &
         &\{ model = "QCD" \}')
    
    call ifile_write (ifile, u)

    write (u, "(A)")
    write (u, "(A)")  "* Parse file"
    write (u, "(A)")
    
    call parse_ifile (ifile, pn_root)

    write (u, "(A)")  "* Compile command list"
    write (u, "(A)")

    call command_list%compile (pn_root, global)
    call command_list%write (u)

    write (u, "(A)")
    write (u, "(A)")  "* Execute command list"
    write (u, "(A)")

    call command_list%execute (global)
    call global%write_libraries (u)

    write (u, "(A)")
    write (u, "(A)")  "* Cleanup"

    call ifile_final (ifile)

    call command_list%final ()
    call global%final ()
    call syntax_cmd_list_final ()
    call syntax_model_file_final ()

    write (u, "(A)")
    write (u, "(A)")  "* Test output end: commands_25"

  end subroutine commands_25

\eatline
\nwendcode{}\nwbegindocs{3112}\nwdocspar
\subsubsection{Alternative Setups}
Declare a list of alternative setups.
\nwenddocs{}\nwbegincode{3113}\moddef{Commands: execute tests}\plusendmoddef
  call test (commands_26, "commands_26", &
       "alternative setups", &
       u, results)
\nwendcode{}\nwbegincode{3114}\moddef{Commands: tests}\plusendmoddef
  subroutine commands_26 (u)
    integer, intent(in) :: u
    type(ifile_t) :: ifile
    type(command_list_t), target :: command_list
    type(rt_data_t), target :: global
    type(parse_node_t), pointer :: pn_root

    write (u, "(A)")  "* Test output: commands_26"
    write (u, "(A)")  "*   Purpose: declare alternative setups for simulation"
    write (u, "(A)")

    call syntax_cmd_list_init ()
    call global%global_init ()
    
    write (u, "(A)")  "* Input file"
    write (u, "(A)")
    
    call ifile_append (ifile, 'int i = 0')
    call ifile_append (ifile, 'alt_setup = (\{ i = 1 \}, \{ i = 2 \})')
    
    call ifile_write (ifile, u)

    write (u, "(A)")
    write (u, "(A)")  "* Parse file"
    write (u, "(A)")
    
    call parse_ifile (ifile, pn_root)

    write (u, "(A)")  "* Compile command list"
    write (u, "(A)")

    call command_list%compile (pn_root, global)
    call command_list%write (u)

    write (u, "(A)")
    write (u, "(A)")  "* Execute command list"
    write (u, "(A)")

    call command_list%execute (global)

    call global%write_expr (u)
    
    write (u, "(A)")
    write (u, "(A)")  "* Cleanup"

    call ifile_final (ifile)

    call command_list%final ()
    call global%final ()
    call syntax_cmd_list_final ()

    write (u, "(A)")
    write (u, "(A)")  "* Test output end: commands_26"

  end subroutine commands_26

\eatline
\nwendcode{}\nwbegindocs{3115}\nwdocspar
\subsubsection{Unstable Particle}
Define decay processes and declare a particle as unstable.  Also check
the commands stable, polarized, unpolarized.
\nwenddocs{}\nwbegincode{3116}\moddef{Commands: execute tests}\plusendmoddef
  call test (commands_27, "commands_27", &
       "unstable and polarized particles", &
       u, results)
\nwendcode{}\nwbegincode{3117}\moddef{Commands: tests}\plusendmoddef
  subroutine commands_27 (u)
    integer, intent(in) :: u
    type(ifile_t) :: ifile
    type(command_list_t), target :: command_list
    type(rt_data_t), target :: global
    type(parse_node_t), pointer :: pn_root
    type(prclib_entry_t), pointer :: lib

    write (u, "(A)")  "* Test output: commands_27"
    write (u, "(A)")  "*   Purpose: modify particle properties"
    write (u, "(A)")

    call syntax_cmd_list_init ()
    call syntax_model_file_init ()
    call global%global_init ()
    call var_list_set_string (global%var_list, var_str ("$method"), &
         var_str ("unit_test"), is_known=.true.)
    call var_list_set_string (global%var_list, var_str ("$phs_method"), &
         var_str ("single"), is_known=.true.)
    call var_list_set_string (global%var_list, var_str ("$integration_method"),&
         var_str ("midpoint"), is_known=.true.)
    call var_list_set_log (global%var_list, var_str ("?vis_history"),&
         .false., is_known=.true.)    
    call var_list_set_log (global%var_list, var_str ("?integration_timer"),&
         .false., is_known = .true.)    
    
    allocate (lib)
    call lib%init (var_str ("commands_27_lib"))
    call global%add_prclib (lib)

    write (u, "(A)")  "* Input file"
    write (u, "(A)")
    
    call ifile_append (ifile, 'model = "Test"')
    call ifile_append (ifile, 'ff = 0.4')
    call ifile_append (ifile, 'process d1 = s => f, fbar')
    call ifile_append (ifile, 'unstable s (d1)')
    call ifile_append (ifile, 'polarized f, fbar')
    
    call ifile_write (ifile, u)

    write (u, "(A)")
    write (u, "(A)")  "* Parse file"
    write (u, "(A)")
    
    call parse_ifile (ifile, pn_root)

    write (u, "(A)")  "* Compile command list"
    write (u, "(A)")

    call command_list%compile (pn_root, global)
    call command_list%write (u)

    write (u, "(A)")
    write (u, "(A)")  "* Execute command list"
    write (u, "(A)")

    call command_list%execute (global)

    write (u, "(A)")  "* Show model"
    write (u, "(A)")
    
    call global%model%write (u)

    write (u, "(A)")
    write (u, "(A)")  "* Extra Input"
    write (u, "(A)")
    
    call ifile_final (ifile)
    call ifile_append (ifile, '?diagonal_decay = true')
    call ifile_append (ifile, 'unstable s (d1)')
    
    call ifile_write (ifile, u)

    write (u, "(A)")
    write (u, "(A)")  "* Parse file"
    write (u, "(A)")
    
    call parse_ifile (ifile, pn_root)

    write (u, "(A)")  "* Compile command list"
    write (u, "(A)")

    call command_list%final ()
    call command_list%compile (pn_root, global)
    call command_list%write (u)

    write (u, "(A)")
    write (u, "(A)")  "* Execute command list"
    write (u, "(A)")

    call command_list%execute (global)

    write (u, "(A)")  "* Show model"
    write (u, "(A)")
    
    call global%model%write (u)

    write (u, "(A)")
    write (u, "(A)")  "* Extra Input"
    write (u, "(A)")
    
    call ifile_final (ifile)
    call ifile_append (ifile, '?isotropic_decay = true')
    call ifile_append (ifile, 'unstable s (d1)')
    
    call ifile_write (ifile, u)

    write (u, "(A)")
    write (u, "(A)")  "* Parse file"
    write (u, "(A)")
    
    call parse_ifile (ifile, pn_root)

    write (u, "(A)")  "* Compile command list"
    write (u, "(A)")

    call command_list%final ()
    call command_list%compile (pn_root, global)
    call command_list%write (u)

    write (u, "(A)")
    write (u, "(A)")  "* Execute command list"
    write (u, "(A)")

    call command_list%execute (global)

    write (u, "(A)")  "* Show model"
    write (u, "(A)")
    
    call global%model%write (u)

    write (u, "(A)")
    write (u, "(A)")  "* Extra Input"
    write (u, "(A)")
    
    call ifile_final (ifile)
    call ifile_append (ifile, 'stable s')
    call ifile_append (ifile, 'unpolarized f')
    
    call ifile_write (ifile, u)

    write (u, "(A)")
    write (u, "(A)")  "* Parse file"
    write (u, "(A)")
    
    call parse_ifile (ifile, pn_root)

    write (u, "(A)")  "* Compile command list"
    write (u, "(A)")

    call command_list%final ()
    call command_list%compile (pn_root, global)
    call command_list%write (u)

    write (u, "(A)")
    write (u, "(A)")  "* Execute command list"
    write (u, "(A)")

    call command_list%execute (global)

    write (u, "(A)")  "* Show model"
    write (u, "(A)")
    
    call global%model%write (u)

    write (u, "(A)")
    write (u, "(A)")  "* Cleanup"

    call ifile_final (ifile)

    call command_list%final ()
    call global%final ()
    call syntax_model_file_init ()
    call syntax_cmd_list_final ()

    write (u, "(A)")
    write (u, "(A)")  "* Test output end: commands_27"

  end subroutine commands_27

\eatline
\nwendcode{}\nwbegindocs{3118}\nwdocspar
\subsubsection{Quit the program}
Quit the program.
\nwenddocs{}\nwbegincode{3119}\moddef{Commands: execute tests}\plusendmoddef
  call test (commands_28, "commands_28", &
       "quit", &
       u, results)
\nwendcode{}\nwbegincode{3120}\moddef{Commands: tests}\plusendmoddef
  subroutine commands_28 (u)
    integer, intent(in) :: u
    type(ifile_t) :: ifile
    type(command_list_t), target :: command_list
    type(rt_data_t), target :: global
    type(parse_node_t), pointer :: pn_root1, pn_root2
    type(string_t), dimension(0) :: no_vars

    write (u, "(A)")  "* Test output: commands_28"
    write (u, "(A)")  "*   Purpose: quit the program"
    write (u, "(A)")

    write (u, "(A)")  "*  Initialization"
    write (u, "(A)")

    call syntax_cmd_list_init ()
    call global%global_init ()    
    
    write (u, "(A)")  "*  Input file: quit without code"
    write (u, "(A)")
    
    call ifile_append (ifile, 'quit')    
    
    call ifile_write (ifile, u)

    write (u, "(A)")
    write (u, "(A)")  "*  Parse file"
    write (u, "(A)")
    
    call parse_ifile (ifile, pn_root1, u)

    write (u, "(A)")
    write (u, "(A)")  "* Compile command list"
    write (u, "(A)")

    call command_list%compile (pn_root1, global)
    call command_list%write (u)

    write (u, "(A)")
    write (u, "(A)")  "* Execute command list"
    write (u, "(A)")

    call command_list%execute (global)

    call global%write (u, vars = no_vars)

    write (u, "(A)")
    write (u, "(A)")  "*  Input file: quit with code"
    write (u, "(A)")
    
    call ifile_final (ifile)
    call command_list%final ()
    call ifile_append (ifile, 'quit ( 3 + 4 )')        
   
    call ifile_write (ifile, u)

    write (u, "(A)")
    write (u, "(A)")  "*  Parse file"
    write (u, "(A)")
    
    call parse_ifile (ifile, pn_root2, u)

    write (u, "(A)")
    write (u, "(A)")  "* Compile command list"
    write (u, "(A)")

    call command_list%compile (pn_root2, global)
    call command_list%write (u)

    write (u, "(A)")
    write (u, "(A)")  "* Execute command list"
    write (u, "(A)")

    call command_list%execute (global)

    call global%write (u, vars = no_vars)

    write (u, "(A)")
    write (u, "(A)")  "* Cleanup"

    call ifile_final (ifile)

    call command_list%final ()
    call global%final ()
    call syntax_cmd_list_final ()
    
    write (u, "(A)")
    write (u, "(A)")  "* Test output end: commands_28"

  end subroutine commands_28

\eatline
\nwendcode{}\nwbegindocs{3121}\nwdocspar
\nwenddocs{}\nwbegindocs{3122}\nwdocspar
\subsubsection{SLHA interface}
Testing commands steering the SLHA interface.
\nwenddocs{}\nwbegincode{3123}\moddef{Commands: execute tests}\plusendmoddef
  call test (commands_29, "commands_29", &
       "SLHA interface", &
       u, results)
\nwendcode{}\nwbegincode{3124}\moddef{Commands: tests}\plusendmoddef
  subroutine commands_29 (u)
    integer, intent(in) :: u
    type(ifile_t) :: ifile
    type(command_list_t), target :: command_list
    type(rt_data_t), target :: global
    type(var_list_t), pointer :: model_vars
    type(parse_node_t), pointer :: pn_root

    write (u, "(A)")  "* Test output: commands_29"
    write (u, "(A)")  "*   Purpose: test SLHA interface"
    write (u, "(A)")

    write (u, "(A)")  "*  Initialization"
    write (u, "(A)")

    call syntax_cmd_list_init ()
    call syntax_model_file_init ()
    call syntax_slha_init ()
    call global%global_init ()
    
    write (u, "(A)")  "*  Model MSSM, read SLHA file"
    write (u, "(A)")
    
    call ifile_append (ifile, 'model = "MSSM"')
    call ifile_append (ifile, '?slha_read_decays = true')    
    call ifile_append (ifile, 'read_slha ("sps1ap_decays.slha")')    
    
    call ifile_write (ifile, u)

    write (u, "(A)")
    write (u, "(A)")  "*  Parse file"
    write (u, "(A)")
    
    call parse_ifile (ifile, pn_root, u)

    write (u, "(A)")
    write (u, "(A)")  "* Compile command list"
    write (u, "(A)")

    call command_list%compile (pn_root, global)
    call command_list%write (u)
           
    write (u, "(A)")
    write (u, "(A)")  "* Model MSSM, default values:"
    write (u, "(A)")    
        
    call global%model%write (u, verbose = .false., &
         show_vertices = .false., show_particles = .false.)
    
    write (u, "(A)")
    write (u, "(A)")  "* Selected global variables"
    write (u, "(A)")

    model_vars => global%model%get_var_list_ptr ()

    call var_list_write_var (model_vars, var_str ("mch1"), u)
    call var_list_write_var (model_vars, var_str ("wch1"), u)

    write (u, "(A)")
    write (u, "(A)")  "* Execute command list"
    write (u, "(A)")

    call command_list%execute (global)

    write (u, "(A)")  "* Model MSSM, values from SLHA file"
    write (u, "(A)")
        
    call global%model%write (u, verbose = .false., &
         show_vertices = .false., show_particles = .false.)

    write (u, "(A)")
    write (u, "(A)")  "* Selected global variables"
    write (u, "(A)")

    model_vars => global%model%get_var_list_ptr ()

    call var_list_write_var (model_vars, var_str ("mch1"), u)
    call var_list_write_var (model_vars, var_str ("wch1"), u)

    write (u, "(A)")
    write (u, "(A)")  "* Cleanup"

    call ifile_final (ifile)

    call command_list%final ()
    call global%final ()
    call syntax_slha_final ()
    call syntax_model_file_final ()
    call syntax_cmd_list_final ()
    
    write (u, "(A)")
    write (u, "(A)")  "* Test output end: commands_29"

  end subroutine commands_29

\eatline
\nwendcode{}\nwbegindocs{3125}\nwdocspar
\subsubsection{Expressions for scales}
Declare a scale, factorization scale or factorization scale expression.
\nwenddocs{}\nwbegincode{3126}\moddef{Commands: execute tests}\plusendmoddef
  call test (commands_30, "commands_30", &
       "scales", &
       u, results)
\nwendcode{}\nwbegincode{3127}\moddef{Commands: tests}\plusendmoddef
  subroutine commands_30 (u)
    integer, intent(in) :: u
    type(ifile_t) :: ifile
    type(command_list_t), target :: command_list
    type(rt_data_t), target :: global
    type(parse_node_t), pointer :: pn_root

    write (u, "(A)")  "* Test output: commands_30"
    write (u, "(A)")  "*   Purpose: define scales"
    write (u, "(A)")

    write (u, "(A)")  "*  Initialization"
    write (u, "(A)")

    call syntax_cmd_list_init ()
    call global%global_init ()

    write (u, "(A)")  "*  Input file"
    write (u, "(A)")
    
    call ifile_append (ifile, 'scale = 200 GeV')
    call ifile_append (ifile, &
         'factorization_scale = eval Pt [particle]')
    call ifile_append (ifile, &
         'renormalization_scale = eval E [particle]')
    
    call ifile_write (ifile, u)

    write (u, "(A)")
    write (u, "(A)")  "*  Parse file"
    write (u, "(A)")
    
    call parse_ifile (ifile, pn_root, u)

    write (u, "(A)")
    write (u, "(A)")  "* Compile command list"
    write (u, "(A)")

    call command_list%compile (pn_root, global)
    call command_list%write (u)

    write (u, "(A)")
    write (u, "(A)")  "* Execute command list"
    write (u, "(A)")

    call command_list%execute (global)

    call global%write_expr (u)

    write (u, "(A)")
    write (u, "(A)")  "* Cleanup"

    call ifile_final (ifile)

    call command_list%final ()
    call global%final ()
    call syntax_cmd_list_final ()

    write (u, "(A)")
    write (u, "(A)")  "* Test output end: commands_30"

  end subroutine commands_30

\eatline
\nwendcode{}\nwbegindocs{3128}\nwdocspar
\subsubsection{Weight and reweight expressions}
Declare an expression for event weights and reweighting.
\nwenddocs{}\nwbegincode{3129}\moddef{Commands: execute tests}\plusendmoddef
  call test (commands_31, "commands_31", &
       "event weights/reweighting", &
       u, results)
\nwendcode{}\nwbegincode{3130}\moddef{Commands: tests}\plusendmoddef
  subroutine commands_31 (u)
    integer, intent(in) :: u
    type(ifile_t) :: ifile
    type(command_list_t), target :: command_list
    type(rt_data_t), target :: global
    type(parse_node_t), pointer :: pn_root

    write (u, "(A)")  "* Test output: commands_31"
    write (u, "(A)")  "*   Purpose: define weight/reweight"
    write (u, "(A)")

    write (u, "(A)")  "*  Initialization"
    write (u, "(A)")

    call syntax_cmd_list_init ()
    call global%global_init ()

    write (u, "(A)")  "*  Input file"
    write (u, "(A)")
    
    call ifile_append (ifile, 'weight = eval Pz [particle]')
    call ifile_append (ifile, 'reweight = eval M2 [particle]')    
    
    call ifile_write (ifile, u)

    write (u, "(A)")
    write (u, "(A)")  "*  Parse file"
    write (u, "(A)")
    
    call parse_ifile (ifile, pn_root, u)

    write (u, "(A)")
    write (u, "(A)")  "* Compile command list"
    write (u, "(A)")

    call command_list%compile (pn_root, global)
    call command_list%write (u)

    write (u, "(A)")
    write (u, "(A)")  "* Execute command list"
    write (u, "(A)")

    call command_list%execute (global)

    call global%write_expr (u)

    write (u, "(A)")
    write (u, "(A)")  "* Cleanup"

    call ifile_final (ifile)

    call command_list%final ()
    call global%final ()
    call syntax_cmd_list_final ()
    call syntax_model_file_final ()

    write (u, "(A)")
    write (u, "(A)")  "* Test output end: commands_31"

  end subroutine commands_31

\eatline
\nwendcode{}\nwbegindocs{3131}\nwdocspar
\subsubsection{Selecting events}
Declare an expression for selecting events in an analysis.
\nwenddocs{}\nwbegincode{3132}\moddef{Commands: execute tests}\plusendmoddef
  call test (commands_32, "commands_32", &
       "event selection", &
       u, results)
\nwendcode{}\nwbegincode{3133}\moddef{Commands: tests}\plusendmoddef
  subroutine commands_32 (u)
    integer, intent(in) :: u
    type(ifile_t) :: ifile
    type(command_list_t), target :: command_list
    type(rt_data_t), target :: global
    type(parse_node_t), pointer :: pn_root

    write (u, "(A)")  "* Test output: commands_32"
    write (u, "(A)")  "*   Purpose: define selection"
    write (u, "(A)")

    write (u, "(A)")  "*  Initialization"
    write (u, "(A)")

    call syntax_cmd_list_init ()
    call global%global_init ()

    write (u, "(A)")  "*  Input file"
    write (u, "(A)")
    
    call ifile_append (ifile, 'selection = any PDG == 13 [particle]')
    
    call ifile_write (ifile, u)

    write (u, "(A)")
    write (u, "(A)")  "*  Parse file"
    write (u, "(A)")
    
    call parse_ifile (ifile, pn_root, u)

    write (u, "(A)")
    write (u, "(A)")  "* Compile command list"
    write (u, "(A)")

    call command_list%compile (pn_root, global)
    call command_list%write (u)

    write (u, "(A)")
    write (u, "(A)")  "* Execute command list"
    write (u, "(A)")

    call command_list%execute (global)

    call global%write_expr (u)

    write (u, "(A)")
    write (u, "(A)")  "* Cleanup"

    call ifile_final (ifile)

    call command_list%final ()
    call global%final ()
    call syntax_cmd_list_final ()

    write (u, "(A)")
    write (u, "(A)")  "* Test output end: commands_32"

  end subroutine commands_32

\eatline
\nwendcode{}\nwbegindocs{3134}\nwdocspar
\subsubsection{Executing shell commands}
Execute a shell command.
\nwenddocs{}\nwbegincode{3135}\moddef{Commands: execute tests}\plusendmoddef
  call test (commands_33, "commands_33", &
       "execute shell command", &
       u, results)
\nwendcode{}\nwbegincode{3136}\moddef{Commands: tests}\plusendmoddef
  subroutine commands_33 (u)
    integer, intent(in) :: u
    type(ifile_t) :: ifile
    type(command_list_t), target :: command_list
    type(rt_data_t), target :: global
    type(parse_node_t), pointer :: pn_root
    integer :: u_file, iostat
    character(3) :: buffer

    write (u, "(A)")  "* Test output: commands_33"
    write (u, "(A)")  "*   Purpose: execute shell command"
    write (u, "(A)")

    write (u, "(A)")  "*  Initialization"
    write (u, "(A)")

    call syntax_cmd_list_init ()
    call global%global_init ()

    write (u, "(A)")  "*  Input file"
    write (u, "(A)")
    
    call ifile_append (ifile, 'exec ("echo foo >> bar")')
    
    call ifile_write (ifile, u)

    write (u, "(A)")
    write (u, "(A)")  "*  Parse file"
    write (u, "(A)")
    
    call parse_ifile (ifile, pn_root, u)

    write (u, "(A)")
    write (u, "(A)")  "* Compile command list"
    write (u, "(A)")

    call command_list%compile (pn_root, global)
    call command_list%write (u)

    write (u, "(A)")
    write (u, "(A)")  "* Execute command list"
    write (u, "(A)")

    call command_list%execute (global)
    u_file = free_unit ()
    open (u_file, file = "bar", &
         action = "read", status = "old")
    do 
       read (u_file, "(A)", iostat = iostat)  buffer
       if (iostat /= 0) exit        
    end do
    write (u, "(A,A)")  "should be 'foo': ", trim (buffer)           
    close (u_file)
        
    write (u, "(A)")
    write (u, "(A)")  "* Cleanup"

    call ifile_final (ifile)

    call command_list%final ()
    call global%final ()
    call syntax_cmd_list_final ()

    write (u, "(A)")
    write (u, "(A)")  "* Test output end: commands_33"

  end subroutine commands_33

\eatline
\nwendcode{}\nwbegindocs{3137}\nwdocspar
\clearpage
%%%%%%%%%%%%%%%%%%%%%%%%%%%%%%%%%%%%%%%%%%%%%%%%%%%%%%%%%%%%%%%%%%%%%%%%
\section{Toplevel module WHIZARD}
\nwenddocs{}\nwbegincode{3138}\moddef{\code{}whizard.f90\edoc{}}\endmoddef
\LA{}File header\RA{}

module whizard

  use io_units
\LA{}Use strings\RA{}
  use unit_tests
  use system_defs, only: VERSION_STRING
  use system_defs, only: EOF, BACKSLASH
  use diagnostics
  use os_interface
  use formats
  use md5
  use sorting
  use cputime
  use sm_qcd
  use ifiles
  use lexers
  use parser
  use xml
  use colors
  use state_matrices
  use analysis
  use variables
  use model_testbed
  use user_code_interface
  use eval_trees
  use particles
  use models
  use auto_components
  use evaluators
  use phs_forests
  use beams
  use polarizations
  use sf_aux
  use sf_mappings
  use sf_base
  use sf_pdf_builtin
  use sf_lhapdf
  use sf_circe1
  use sf_circe2
  use sf_isr
  use sf_epa
  use sf_ewa
  use sf_escan
  use sf_beam_events
  use sf_user
  use phs_base
  use phs_single
  use phs_wood
  use rng_base
  use rng_tao
  use selectors
  use mci_base
  use mci_midpoint
  use mci_vamp
  use prclib_interfaces
  use particle_specifiers
  use process_libraries
  use prclib_stacks
  use hepmc_interface
  use lcio_interface
  use jets
  use pdg_arrays
  use interactions
  use slha_interface
  use cascades
  use blha_driver
  use blha_config
  use prc_core
  use prc_test
  use prc_template_me
  use prc_omega
  use subevt_expr
  use processes
  use process_stacks
  use event_transforms
  use decays
  use shower
  use events

  use hep_events
  use eio_data
  use eio_base
  use eio_raw
  use eio_checkpoints
  use eio_lhef
  use eio_hepmc
  use eio_lcio
  use eio_stdhep
  use eio_ascii
  use eio_weights

  use iterations
  use beam_structures
  use rt_data
  use dispatch
  use process_configurations
  use compilations
  use integrations
  use event_streams
  use simulations
  use nlo_data

  use expr_tests

  use commands

\LA{}Standard module head\RA{}

\LA{}WHIZARD: public\RA{}

\LA{}WHIZARD: types\RA{}

  save

contains

\LA{}WHIZARD: procedures\RA{}

end module whizard
\eatline
\nwendcode{}\nwbegindocs{3139}\nwdocspar
\subsection{Options}
Here we introduce a wrapper that holds various user options, so they
can transparently be passed from the main program to the {\tt{}whizard}
object.  Most parameters are used for initializing the {\tt{}global}
state.
\nwenddocs{}\nwbegincode{3140}\moddef{WHIZARD: public}\endmoddef
  public :: whizard_options_t
\nwendcode{}\nwbegincode{3141}\moddef{WHIZARD: types}\endmoddef
  type :: whizard_options_t
     type(string_t) :: preload_model
     type(string_t) :: default_lib
     type(string_t) :: preload_libraries
     logical :: rebuild_library = .false.
     logical :: recompile_library = .false.
     logical :: rebuild_user
     logical :: rebuild_phs = .false.
     logical :: rebuild_grids = .false.
     logical :: rebuild_events = .false.
  end type whizard_options_t
  
\eatline
\nwendcode{}\nwbegindocs{3142}\nwdocspar
\subsection{Parse tree stack}
We collect all parse trees that we generate in the {\tt{}whizard} object.  To
this end, we create a stack of parse trees.  They must not be finalized before
the {\tt{}global} object is finalized, because items such as a cut definition may
contain references to the parse tree from which they were generated.
\nwenddocs{}\nwbegincode{3143}\moddef{WHIZARD: types}\plusendmoddef
  type, extends (parse_tree_t) :: pt_entry_t
     type(pt_entry_t), pointer :: previous => null ()
  end type pt_entry_t
  
\eatline
\nwendcode{}\nwbegindocs{3144}This is the stack.  Since we always prepend, we just need the {\tt{}last}
pointer.
\nwenddocs{}\nwbegincode{3145}\moddef{WHIZARD: types}\plusendmoddef
  type :: pt_stack_t
     type(pt_entry_t), pointer :: last => null ()
   contains
   \LA{}WHIZARD: pt stack: TBP\RA{}
  end type pt_stack_t
  
\eatline
\nwendcode{}\nwbegindocs{3146}The finalizer is called at the very end.
\nwenddocs{}\nwbegincode{3147}\moddef{WHIZARD: pt stack: TBP}\endmoddef
  procedure :: final => pt_stack_final
\nwendcode{}\nwbegincode{3148}\moddef{WHIZARD: procedures}\endmoddef
  subroutine pt_stack_final (pt_stack)
    class(pt_stack_t), intent(inout) :: pt_stack
    type(pt_entry_t), pointer :: current
    do while (associated (pt_stack%last))
       current => pt_stack%last
       pt_stack%last => current%previous
       call parse_tree_final (current%parse_tree_t)
       deallocate (current)
    end do
  end subroutine pt_stack_final

\eatline
\nwendcode{}\nwbegindocs{3149}Create and push a new entry, keeping the previous ones.
\nwenddocs{}\nwbegincode{3150}\moddef{WHIZARD: pt stack: TBP}\plusendmoddef
  procedure :: push => pt_stack_push
\nwendcode{}\nwbegincode{3151}\moddef{WHIZARD: procedures}\plusendmoddef
  subroutine pt_stack_push (pt_stack, parse_tree)
    class(pt_stack_t), intent(inout) :: pt_stack
    type(parse_tree_t), intent(out), pointer :: parse_tree
    type(pt_entry_t), pointer :: current
    allocate (current)
    parse_tree => current%parse_tree_t
    current%previous => pt_stack%last
    pt_stack%last => current
  end subroutine pt_stack_push
  
\eatline
\nwendcode{}\nwbegindocs{3152}\nwdocspar
\subsection{The {\tt{}whizard} object}
An object of type {\tt{}whizard{\char95}t} is the top-level wrapper for a
\whizard\ instance.  The object holds various default
settings and the current state of the generator, the {\tt{}global} object
of type {\tt{}rt{\char95}data{\char95}t}.  This object contains, for instance, the list
of variables and the process libraries.

Since components of the {\tt{}global} subobject are frequently used as
targets, the {\tt{}whizard} object should also consistently carry the
{\tt{}target} attribute.

The various self-tests do no not use this object.  They initialize
only specific subsets of the system, according to their needs.

Note: we intend to allow several concurrent instances.  In the current
implementation, there are still a few obstacles to this: the model
library and the syntax tables are global variables, and the error
handling uses global state.  This should be improved.
\nwenddocs{}\nwbegincode{3153}\moddef{WHIZARD: public}\plusendmoddef
  public :: whizard_t
\nwendcode{}\nwbegincode{3154}\moddef{WHIZARD: types}\plusendmoddef
  type :: whizard_t
     type(whizard_options_t) :: options
     type(rt_data_t) :: global
     type(pt_stack_t) :: pt_stack
   contains
   \LA{}WHIZARD: whizard: TBP\RA{}
  end type whizard_t
  
\eatline
\nwendcode{}\nwbegindocs{3155}\nwdocspar
\subsection{Initialization and finalization}
\nwenddocs{}\nwbegincode{3156}\moddef{WHIZARD: whizard: TBP}\endmoddef
  procedure :: init => whizard_init
\nwendcode{}\nwbegincode{3157}\moddef{WHIZARD: procedures}\plusendmoddef
  subroutine whizard_init (whizard, options, paths, logfile)
    class(whizard_t), intent(out), target :: whizard
    type(whizard_options_t), intent(in) :: options    
    type(paths_t), intent(in), optional :: paths
    type(string_t), intent(in), optional :: logfile
    call init_syntax_tables ()
    whizard%options = options
    call whizard%global%global_init (paths, logfile)
    call whizard%init_rebuild_flags ()
    call whizard%preload_model ()
    call whizard%preload_library ()
    call whizard%global%init_fallback_model &
         (var_str ("SM_hadrons"), var_str ("SM_hadrons.mdl"))
    call whizard%global%init_radiation_model &
         (var_str ("SM_rad"), var_str ("SM_rad.mdl"))
  end subroutine whizard_init
  
\eatline
\nwendcode{}\nwbegindocs{3158}Apart from the global data which have been initialized above, the
process and model lists need to be finalized.
\nwenddocs{}\nwbegincode{3159}\moddef{WHIZARD: whizard: TBP}\plusendmoddef
  procedure :: final => whizard_final
\nwendcode{}\nwbegincode{3160}\moddef{WHIZARD: procedures}\plusendmoddef
  subroutine whizard_final (whizard)
    class(whizard_t), intent(inout), target :: whizard
    call whizard%global%final ()
    call whizard%pt_stack%final ()
!!! JRR: WK please check (#529)
    !    call user_code_final ()
    call final_syntax_tables () 
  end subroutine whizard_final
  
\eatline
\nwendcode{}\nwbegindocs{3161}\nwdocspar
Set the rebuild flags.  They can be specified on the command line and
set the initial value for the associated logical variables.
\nwenddocs{}\nwbegincode{3162}\moddef{WHIZARD: whizard: TBP}\plusendmoddef
  procedure :: init_rebuild_flags => whizard_init_rebuild_flags
\nwendcode{}\nwbegincode{3163}\moddef{WHIZARD: procedures}\plusendmoddef
  subroutine whizard_init_rebuild_flags (whizard)
    class(whizard_t), intent(inout), target :: whizard
    associate (var_list => whizard%global%var_list, options => whizard%options) 
      call var_list_append_log &
           (var_list, var_str ("?rebuild_library"), options%rebuild_library, &
           intrinsic=.true.)
      call var_list_append_log &
           (var_list, var_str ("?recompile_library"), &
           options%recompile_library, &
           intrinsic=.true.)
      call var_list_append_log &
           (var_list, var_str ("?rebuild_phase_space"), options%rebuild_phs, &
           intrinsic=.true.)
      call var_list_append_log &
           (var_list, var_str ("?rebuild_grids"), options%rebuild_grids, &
           intrinsic=.true.)
      call var_list_append_log &
           (var_list, var_str ("?rebuild_events"), options%rebuild_events, &
           intrinsic=.true.)
    end associate
  end subroutine whizard_init_rebuild_flags

\eatline
\nwendcode{}\nwbegindocs{3164}\nwdocspar
This procedure preloads a model, if a model name is given.
\nwenddocs{}\nwbegincode{3165}\moddef{WHIZARD: whizard: TBP}\plusendmoddef
  procedure :: preload_model => whizard_preload_model
\nwendcode{}\nwbegincode{3166}\moddef{WHIZARD: procedures}\plusendmoddef
  subroutine whizard_preload_model (whizard)
    class(whizard_t), intent(inout), target :: whizard
    type(string_t) :: model_name
    model_name = whizard%options%preload_model
    if (model_name /= "") then
       call whizard%global%read_model (model_name, whizard%global%preload_model)
       whizard%global%model => whizard%global%preload_model
       if (associated (whizard%global%model)) then
          call whizard%global%model%link_var_list (whizard%global%var_list)
          call msg_message ("Preloaded model: " &
               // char (model_name))
       else
          call msg_fatal ("Preloading model " // char (model_name) &
               // " failed")
       end if
    else
       call msg_message ("No model preloaded")
    end if
  end subroutine whizard_preload_model
    
\eatline
\nwendcode{}\nwbegindocs{3167}\nwdocspar
This procedure preloads a library, if a library name is given.

Note: This version just opens a new library with that name.  It does not load
(yet) an existing library on file, as previous \whizard\ versions would do.
\nwenddocs{}\nwbegincode{3168}\moddef{WHIZARD: whizard: TBP}\plusendmoddef
  procedure :: preload_library => whizard_preload_library
\nwendcode{}\nwbegincode{3169}\moddef{WHIZARD: procedures}\plusendmoddef
  subroutine whizard_preload_library (whizard)
    class(whizard_t), intent(inout), target :: whizard
    type(string_t) :: library_name, libs
    type(string_t), dimension(:), allocatable :: libname_static
    type(prclib_entry_t), pointer :: lib_entry
    integer :: i
    call get_prclib_static (libname_static)
    do i = 1, size (libname_static)
       allocate (lib_entry)
       call lib_entry%init_static (libname_static(i))
       call whizard%global%add_prclib (lib_entry)
    end do
    libs = adjustl (whizard%options%preload_libraries)
    if (libs == "" .and. whizard%options%default_lib /= "") then
          allocate (lib_entry)
          call lib_entry%init (whizard%options%default_lib)
          call whizard%global%add_prclib (lib_entry)
          call msg_message ("Preloaded library: " // &
               char (whizard%options%default_lib))    
       end if    
    SCAN_LIBS: do while (libs /= "")
       call split (libs, library_name, " ")      
       if (library_name /= "") then
          allocate (lib_entry)
          call lib_entry%init (library_name)
          call whizard%global%add_prclib (lib_entry)
          call msg_message ("Preloaded library: " // char (library_name))
       end if
    end do SCAN_LIBS    
  end subroutine whizard_preload_library
    
\eatline
\nwendcode{}\nwbegindocs{3170}\nwdocspar
\subsection{Initialization and finalization (old version)}
These procedures initialize and finalize global variables.  Most of
them are collected in the {\tt{}global} data record located here, the
others are syntax tables located in various modules, which do not
change during program execution.  Furthermore, there is a global model
list and a global process store, which get filled during program
execution but are finalized here.

During initialization, we can preload a default model and initialize a
default library for setting up processes.  The default library is
loaded if requested by the setup.  Further libraries can be loaded as
specified by command-line flags.
\nwenddocs{}\nwbegindocs{3171}Initialize/finalize the syntax tables used by WHIZARD:
\nwenddocs{}\nwbegincode{3172}\moddef{WHIZARD: public}\plusendmoddef
  public :: init_syntax_tables
  public :: final_syntax_tables
\nwendcode{}\nwbegincode{3173}\moddef{WHIZARD: procedures}\plusendmoddef
  subroutine init_syntax_tables ()
    call syntax_model_file_init ()
    call syntax_phs_forest_init ()
    call syntax_pexpr_init ()
    call syntax_slha_init ()
    call syntax_blha_contract_init ()    
    call syntax_cmd_list_init ()
  end subroutine init_syntax_tables

  subroutine final_syntax_tables ()
    call syntax_model_file_final ()
    call syntax_phs_forest_final ()
    call syntax_pexpr_final ()
    call syntax_slha_final ()
    call syntax_blha_contract_final ()    
    call syntax_cmd_list_final ()
  end subroutine final_syntax_tables

\eatline
\eatline
\nwendcode{}\nwbegindocs{3174}Write the syntax tables to external files.
\nwenddocs{}\nwbegincode{3175}\moddef{WHIZARD: public}\plusendmoddef
  public :: write_syntax_tables
\nwendcode{}\nwbegincode{3176}\moddef{WHIZARD: procedures}\plusendmoddef
  subroutine write_syntax_tables ()
    integer :: unit
    character(*), parameter :: file_model = "whizard.model_file.syntax"
    character(*), parameter :: file_phs = "whizard.phase_space_file.syntax"
    character(*), parameter :: file_pexpr = "whizard.prt_expressions.syntax"
    character(*), parameter :: file_slha = "whizard.slha.syntax"
    character(*), parameter :: file_sindarin = "whizard.sindarin.syntax"
    unit = free_unit ()
    print *, "Writing file '" // file_model // "'"
    open (unit=unit, file=file_model, status="replace", action="write")
    write (unit, "(A)")  VERSION_STRING
    write (unit, "(A)")  "Syntax definition file: " // file_model
    call syntax_model_file_write (unit)
    close (unit)
    print *, "Writing file '" // file_phs // "'"
    open (unit=unit, file=file_phs, status="replace", action="write")
    write (unit, "(A)")  VERSION_STRING
    write (unit, "(A)")  "Syntax definition file: " // file_phs
    call syntax_phs_forest_write (unit)
    close (unit)
    print *, "Writing file '" // file_pexpr // "'"
    open (unit=unit, file=file_pexpr, status="replace", action="write")
    write (unit, "(A)")  VERSION_STRING
    write (unit, "(A)")  "Syntax definition file: " // file_pexpr
    call syntax_pexpr_write (unit)
    close (unit)
    print *, "Writing file '" // file_slha // "'"
    open (unit=unit, file=file_slha, status="replace", action="write")
    write (unit, "(A)")  VERSION_STRING
    write (unit, "(A)")  "Syntax definition file: " // file_slha
    call syntax_slha_write (unit)
    close (unit)
    print *, "Writing file '" // file_sindarin // "'"
    open (unit=unit, file=file_sindarin, status="replace", action="write")
    write (unit, "(A)")  VERSION_STRING
    write (unit, "(A)")  "Syntax definition file: " // file_sindarin
    call syntax_cmd_list_write (unit)
    close (unit)
  end subroutine write_syntax_tables

\eatline
\nwendcode{}\nwbegindocs{3177}\nwdocspar
\subsection{Execute command lists}
Process commands given on the command line, stored as an {\tt{}ifile}.  The whole
input is read, compiled and executed as a whole.
\nwenddocs{}\nwbegincode{3178}\moddef{WHIZARD: whizard: TBP}\plusendmoddef
  procedure :: process_ifile => whizard_process_ifile
\nwendcode{}\nwbegincode{3179}\moddef{WHIZARD: procedures}\plusendmoddef
  subroutine whizard_process_ifile (whizard, ifile, quit, quit_code)
    class(whizard_t), intent(inout), target :: whizard
    type(ifile_t), intent(in) :: ifile
    logical, intent(out) :: quit
    integer, intent(out) :: quit_code
    type(lexer_t), target :: lexer
    type(stream_t), target :: stream
    call msg_message ("Reading commands given on the command line")
    call lexer_init_cmd_list (lexer)
    call stream_init (stream, ifile)
    call whizard%process_stream (stream, lexer, quit, quit_code)
    call stream_final (stream)
    call lexer_final (lexer)
  end subroutine whizard_process_ifile

\eatline
\nwendcode{}\nwbegindocs{3180}Process standard input as a command list.  The whole input is read,
compiled and executed as a whole.
\nwenddocs{}\nwbegincode{3181}\moddef{WHIZARD: whizard: TBP}\plusendmoddef
  procedure :: process_stdin => whizard_process_stdin
\nwendcode{}\nwbegincode{3182}\moddef{WHIZARD: procedures}\plusendmoddef
  subroutine whizard_process_stdin (whizard, quit, quit_code)
    class(whizard_t), intent(inout), target :: whizard
    logical, intent(out) :: quit
    integer, intent(out) :: quit_code
    type(lexer_t), target :: lexer
    type(stream_t), target :: stream
    call msg_message ("Reading commands from standard input")
    call lexer_init_cmd_list (lexer)
    call stream_init (stream, 5)
    call whizard%process_stream (stream, lexer, quit, quit_code)
    call stream_final (stream)
    call lexer_final (lexer)
  end subroutine whizard_process_stdin

\eatline
\nwendcode{}\nwbegindocs{3183}Process a file as a command list.
\nwenddocs{}\nwbegincode{3184}\moddef{WHIZARD: whizard: TBP}\plusendmoddef
  procedure :: process_file => whizard_process_file
\nwendcode{}\nwbegincode{3185}\moddef{WHIZARD: procedures}\plusendmoddef
  subroutine whizard_process_file (whizard, file, quit, quit_code)
    class(whizard_t), intent(inout), target :: whizard
    type(string_t), intent(in) :: file
    logical, intent(out) :: quit
    integer, intent(out) :: quit_code
    type(lexer_t), target :: lexer
    type(stream_t), target :: stream
    logical :: exist
    call msg_message ("Reading commands from file '" // char (file) // "'")
    inquire (file=char(file), exist=exist)
    if (exist) then
       call lexer_init_cmd_list (lexer)
       call stream_init (stream, char (file))
       call whizard%process_stream (stream, lexer, quit, quit_code)
       call stream_final (stream)
       call lexer_final (lexer)
    else
       call msg_error ("File '" // char (file) // "' not found")
    end if
  end subroutine whizard_process_file

\eatline
\nwendcode{}\nwbegindocs{3186}\nwdocspar
\nwenddocs{}\nwbegincode{3187}\moddef{WHIZARD: whizard: TBP}\plusendmoddef
  procedure :: process_stream => whizard_process_stream
\nwendcode{}\nwbegincode{3188}\moddef{WHIZARD: procedures}\plusendmoddef
  subroutine whizard_process_stream (whizard, stream, lexer, quit, quit_code)
    class(whizard_t), intent(inout), target :: whizard
    type(stream_t), intent(inout), target :: stream
    type(lexer_t), intent(inout), target :: lexer
    logical, intent(out) :: quit
    integer, intent(out) :: quit_code
    type(parse_tree_t), pointer :: parse_tree
    type(command_list_t), target :: command_list
    call lexer_assign_stream (lexer, stream)
    call whizard%pt_stack%push (parse_tree)
    call parse_tree_init (parse_tree, syntax_cmd_list, lexer)
    if (associated (parse_tree_get_root_ptr (parse_tree))) then
       whizard%global%lexer => lexer
       call command_list%compile (parse_tree_get_root_ptr (parse_tree), &
            whizard%global)
    end if
    call whizard%global%activate ()
    call command_list%execute (whizard%global)
    call command_list%final ()
    quit = whizard%global%quit
    quit_code = whizard%global%quit_code
  end subroutine whizard_process_stream

\eatline
\nwendcode{}\nwbegindocs{3189}\nwdocspar
\subsection{The WHIZARD shell}
This procedure implements interactive mode.  One line is processed at
a time.
\nwenddocs{}\nwbegincode{3190}\moddef{WHIZARD: whizard: TBP}\plusendmoddef
  procedure :: shell => whizard_shell
\nwendcode{}\nwbegincode{3191}\moddef{WHIZARD: procedures}\plusendmoddef
  subroutine whizard_shell (whizard, quit_code)
    class(whizard_t), intent(inout), target :: whizard
    integer, intent(out) :: quit_code
    type(lexer_t), target :: lexer
    type(stream_t), target :: stream
    type(string_t) :: prompt1
    type(string_t) :: prompt2
    type(string_t) :: input
    type(string_t) :: extra
    integer :: last
    integer :: iostat
    logical :: mask_tmp
    logical :: quit
    call msg_message ("Launching interactive shell")
    call lexer_init_cmd_list (lexer)
    prompt1 = "whish? "
    prompt2 = "     > "
    COMMAND_LOOP: do
       call put (6, prompt1)
       call get (5, input, iostat=iostat)
       if (iostat > 0 .or. iostat == EOF) exit COMMAND_LOOP
       CONTINUE_INPUT: do
          last = len_trim (input)
          if (extract (input, last, last) /= BACKSLASH)  exit CONTINUE_INPUT
          call put (6, prompt2)
          call get (5, extra, iostat=iostat)
          if (iostat > 0) exit COMMAND_LOOP
          input = replace (input, last, extra)
       end do CONTINUE_INPUT
       call stream_init (stream, input)
       mask_tmp = mask_fatal_errors
       mask_fatal_errors = .true.
       call whizard%process_stream (stream, lexer, quit, quit_code)
       msg_count = 0
       mask_fatal_errors = mask_tmp
       call stream_final (stream)
       if (quit)  exit COMMAND_LOOP
    end do COMMAND_LOOP
    print *
    call lexer_final (lexer)
  end subroutine whizard_shell

\eatline
\nwendcode{}\nwbegindocs{3192}\nwdocspar
\subsection{Self-tests}
For those self-tests, we need some auxiliary routines that provide an
enviroment.  The environment depends on things that are not available at the
level of the module that we want to test.

\subsubsection{Testbed for event I/O}
This subroutine prepares a test process with a single event.  All objects are
allocated via anonymous pointers, because we want to recover the pointers and
delete the objects in a separate procedure.
\nwenddocs{}\nwbegincode{3193}\moddef{WHIZARD: procedures}\plusendmoddef
  subroutine prepare_eio_test (event, unweighted, n_alt)
    use variables
    use model_data
    use event_base
    class(generic_event_t), intent(inout), pointer :: event
    logical, intent(in), optional :: unweighted
    integer, intent(in), optional :: n_alt
    type(model_data_t), pointer :: model
    type(var_list_t) :: var_list
    type(process_t), pointer :: process
    type(process_instance_t), pointer :: process_instance

    allocate (model)
    call model%init_test ()

    allocate (process)
    allocate (process_instance)

    call prepare_test_process (process, process_instance, model)
    call process_instance%setup_event_data ()
 
    call model%final ()
    deallocate (model)

    allocate (event_t :: event)
    select type (event)
    type is (event_t)
       if (present (unweighted)) then
          call var_list_append_log (var_list, &
               var_str ("?unweighted"), unweighted, &
               intrinsic = .true.)
       else
          call var_list_append_log (var_list, &
               var_str ("?unweighted"), .true., &
               intrinsic = .true.)
       end if
       call var_list_append_string (var_list, &
            var_str ("$sample_normalization"), &
            var_str ("auto"), intrinsic = .true.)
       call event%basic_init (var_list, n_alt)
       call event%connect (process_instance, process%get_model_ptr ())
       call var_list%final ()
    end select
  end subroutine prepare_eio_test
    
\eatline
\nwendcode{}\nwbegindocs{3194}Recover those pointers, finalize the objects and deallocate.
\nwenddocs{}\nwbegincode{3195}\moddef{WHIZARD: procedures}\plusendmoddef
  subroutine cleanup_eio_test (event)
    use model_data
    use event_base
    class(generic_event_t), intent(inout), pointer :: event
    type(process_t), pointer :: process
    type(process_instance_t), pointer :: process_instance
    select type (event)
    type is (event_t)
       process => event%get_process_ptr ()
       process_instance => event%get_process_instance_ptr ()
       call cleanup_test_process (process, process_instance)
       deallocate (process_instance)
       deallocate (process)
       call event%final ()
    end select
    deallocate (event)
  end subroutine cleanup_eio_test
    
\eatline
\nwendcode{}\nwbegindocs{3196}Assign those procedures to appropriate pointers (module variables) in the
{\tt{}eio{\char95}base} module, so they can be called as if they were module procedures.
\nwenddocs{}\nwbegincode{3197}\moddef{WHIZARD: prepare testbed}\endmoddef
  eio_prepare_test => prepare_eio_test
  eio_cleanup_test => cleanup_eio_test
\nwendcode{}\nwbegindocs{3198}\nwdocspar
\subsubsection{Any Model}
This procedure reads any model from file and, optionally, assigns a
var-list pointer.
\nwenddocs{}\nwbegincode{3199}\moddef{WHIZARD: procedures}\plusendmoddef
  subroutine prepare_whizard_model (model, name, vars)
  \LA{}Use strings\RA{}
    use os_interface
    use model_data
    use var_base
    use models
    class(model_data_t), intent(inout), pointer :: model
    type(string_t), intent(in) :: name
    class(vars_t), pointer, intent(out), optional :: vars
    type(os_data_t) :: os_data
    call syntax_model_file_init ()
    call os_data_init (os_data)
    allocate (model_t :: model)
    select type (model)
    type is (model_t)
       call model%read (name // ".mdl", os_data)
       if (present (vars)) then
          vars => model%get_var_list_ptr ()
       end if
    end select
  end subroutine prepare_whizard_model
    
\eatline
\nwendcode{}\nwbegindocs{3200}Cleanup after use.  Includes deletion of the model-file syntax.
\nwenddocs{}\nwbegincode{3201}\moddef{WHIZARD: procedures}\plusendmoddef
  subroutine cleanup_whizard_model (model)
    use model_data
    use models
    class(model_data_t), intent(inout), pointer :: model
    call model%final ()
    deallocate (model)
    call syntax_model_file_final ()
  end subroutine cleanup_whizard_model
    
\eatline
\nwendcode{}\nwbegindocs{3202}Assign those procedures to appropriate pointers (module variables) in the
{\tt{}model{\char95}testbed} module, so they can be called as if they were module
procedures.
\nwenddocs{}\nwbegincode{3203}\moddef{WHIZARD: prepare testbed}\plusendmoddef
  prepare_model => prepare_whizard_model
  cleanup_model => cleanup_whizard_model
\nwendcode{}\nwbegindocs{3204}\nwdocspar
\subsubsection{Fallback model: hadrons}
Some event format tests require the hadronic SM implementation, which
has to be read from file.  We provide the functionality here, so the
tests do not depend on model I/O.
\nwenddocs{}\nwbegincode{3205}\moddef{WHIZARD: procedures}\plusendmoddef
  subroutine prepare_fallback_model (model)
    use model_data
    class(model_data_t), intent(inout), pointer :: model
    call prepare_whizard_model (model, var_str ("SM_hadrons"))
  end subroutine prepare_fallback_model
    
\eatline
\nwendcode{}\nwbegindocs{3206}Assign those procedures to appropriate pointers (module variables) in the
{\tt{}eio{\char95}base} module, so they can be called as if they were module procedures.
\nwenddocs{}\nwbegincode{3207}\moddef{WHIZARD: prepare testbed}\plusendmoddef
  eio_prepare_fallback_model => prepare_fallback_model
  eio_cleanup_fallback_model => cleanup_model
\nwendcode{}\nwbegindocs{3208}\nwdocspar
\subsubsection{Procedure for Checking}
This is for developers only, but needs a well-defined interface.
\nwenddocs{}\nwbegincode{3209}\moddef{WHIZARD: public}\plusendmoddef
  public :: whizard_check
\nwendcode{}\nwbegincode{3210}\moddef{WHIZARD: procedures}\plusendmoddef
  subroutine whizard_check (check, results)
    type(string_t), intent(in) :: check
    type(test_results_t), intent(inout) :: results
    type(os_data_t) :: os_data
    integer :: u
    call os_data_init (os_data)
    u = free_unit ()
    open (u, file="whizard_check." // char (check) // ".log", &
         action="write", status="replace")
    call msg_message (repeat ('=', 76), 0)
    call msg_message ("Running self-test: " // char (check), 0)
    call msg_message (repeat ('-', 76), 0)
  \LA{}WHIZARD: prepare testbed\RA{}
    select case (char (check))
    case ("analysis")
       call analysis_test (u, results)
    case ("beams")
       call beam_test (u, results)
    case ("cascades")
       call cascade_test (u, results)
    case ("colors")
       call color_test (u, results)
    case ("evaluators")
      call evaluator_test (u, results)
    case ("expressions")
       call expressions_test (u, results)
    case ("formats")
       call format_test (u, results)
    case ("hepmc")
       call hepmc_test (u, results)
    case ("lcio")
       call lcio_test (u, results)
    case ("jets")
       call jets_test (u, results)
    case ("pdg_arrays")
       call pdg_arrays_test (u, results)
    case ("interactions")
       call interaction_test (u, results)
    case ("lexers")
       call lexer_test (u, results)
    case ("os_interface") 
       call os_interface_test (u, results)
    case ("cputime") 
       call cputime_test (u, results)
    case ("parser")
       call parse_test (u, results)       
    case ("sorting")
       call sorting_test (u, results)
    case ("md5")       
       call md5_test (u, results)       
    case ("xml")
       call xml_test (u, results)
    case ("sm_qcd")
       call sm_qcd_test (u, results)
    case ("models")
       call models_test (u, results)
    case ("auto_components")
       call auto_components_test (u, results)
    case ("particles")
       call particles_test (u, results)
    case ("polarizations")
       call polarization_test (u, results)
    case ("sf_aux")
       call sf_aux_test (u, results)
    case ("sf_mappings")
       call sf_mappings_test (u, results)
    case ("sf_base")
       call sf_base_test (u, results)
    case ("sf_pdf_builtin")
       call sf_pdf_builtin_test (u, results)
    case ("sf_lhapdf")
       call sf_lhapdf_test (u, results)
    case ("sf_isr")
       call sf_isr_test (u, results)
    case ("sf_epa")
       call sf_epa_test (u, results)
    case ("sf_ewa")
       call sf_ewa_test (u, results)
    case ("sf_circe1")
       call sf_circe1_test (u, results)
    case ("sf_circe2")
       call sf_circe2_test (u, results)
    case ("sf_beam_events")
       call sf_beam_events_test (u, results)
    case ("sf_escan")
       call sf_escan_test (u, results)
    case ("phs_base")
       call phs_base_test (u, results)
    case ("phs_single")
       call phs_single_test (u, results)
    case ("phs_forests")
       call phs_forest_test (u, results)       
    case ("phs_wood")
       call phs_wood_test (u, results)
    case ("phs_wood_vis")
       call phs_wood_vis_test (u, results)
    case ("mci_base")
       call mci_base_test (u, results)
    case ("rng_base")
       call rng_base_test (u, results)
    case ("rng_tao")
       call rng_tao_test (u, results)
    case ("selectors")
       call selectors_test (u, results)
    case ("mci_midpoint")
       call mci_midpoint_test (u, results)
    case ("mci_vamp")
       call mci_vamp_test (u, results)
    case ("prclib_interfaces")
       call prclib_interfaces_test (u, results)
    case ("particle_specifiers")
       call particle_specifiers_test (u, results)
    case ("process_libraries")
       call process_libraries_test (u, results)
    case ("prclib_stacks")
       call prclib_stacks_test (u, results)
    case ("slha_interface")
       call slha_test (u, results)
    case ("state_matrices")
       call state_matrix_test (u, results)
    case ("prc_test")
       call prc_test_test (u, results)
    case ("subevt_expr")
       call subevt_expr_test (u, results)
    case ("processes")
       call processes_test (u, results)       
    case ("process_stacks")
       call process_stacks_test (u, results)   
    case ("event_transforms")
       call event_transforms_test (u, results)       
    case ("decays")
       call decays_test (u, results)       
    case ("shower")
       call shower_test (u, results)       
    case ("events")
       call events_test (u, results)
    case ("prc_template_me")
       call prc_template_me_test (u, results)              
    case ("prc_omega")
       call prc_omega_test (u, results)
    case ("prc_omega_diags")
       call prc_omega_diags_test (u, results)       
    case ("hep_events")
       call hep_events_test (u, results)       
    case ("eio_data")       
       call eio_data_test (u, results)
    case ("eio_base")
       call eio_base_test (u, results)
    case ("eio_raw")
       call eio_raw_test (u, results)
    case ("eio_checkpoints")
       call eio_checkpoints_test (u, results)
    case ("eio_lhef")
       call eio_lhef_test (u, results)
    case ("eio_hepmc")
       call eio_hepmc_test (u, results)
    case ("eio_lcio")
       call eio_lcio_test (u, results)
    case ("eio_stdhep")
       call eio_stdhep_test (u, results)
    case ("eio_ascii")
       call eio_ascii_test (u, results)       
    case ("eio_weights")
       call eio_weights_test (u, results)
    case ("iterations")
       call iterations_test (u, results)
    case ("beam_structures")
       call beam_structures_test (u, results)
    case ("rt_data")
       call rt_data_test (u, results)
    case ("dispatch")
       call dispatch_test (u, results)
    case ("process_configurations")
       call process_configurations_test (u, results)
    case ("compilations")
       call compilations_test (u, results)
    case ("compilations_static")
       call compilations_static_test (u, results)
    case ("integrations")
       call integrations_test (u, results)
    case ("integrations_history")
       call integrations_history_test (u, results)
    case ("event_streams")
       call event_streams_test (u, results)
    case ("simulations")
       call simulations_test (u, results)
    case ("commands")
       call commands_test (u, results)
    case ("all")
       call analysis_test (u, results)
       call beam_test (u, results)
       call md5_test (u, results)
       call lexer_test (u, results)
       call sorting_test (u, results)
       call parse_test (u, results)
       call color_test (u, results)
       call evaluator_test (u, results)
       call expressions_test (u, results)
       call format_test (u, results)
       call hepmc_test (u, results)
       call lcio_test (u, results)
       call jets_test (u, results)
       call os_interface_test (u, results)
       call cputime_test (u, results)
       call interaction_test (u, results)
       call xml_test (u, results)
       call sm_qcd_test (u, results)
       call models_test (u, results)
       call auto_components_test (u, results)
       call particles_test (u, results)
       call polarization_test (u, results)
       call sf_aux_test (u, results)
       call sf_mappings_test (u, results)
       call sf_base_test (u, results)
       call sf_pdf_builtin_test (u, results)
       call sf_lhapdf_test (u, results)
       call sf_isr_test (u, results)
       call sf_epa_test (u, results)
       call sf_ewa_test (u, results)
       call sf_circe1_test (u, results)
       call sf_circe2_test (u, results)
       call sf_beam_events_test (u, results)
       call sf_escan_test (u, results)
       call phs_base_test (u, results)
       call phs_single_test (u, results)
       call phs_forest_test (u, results)
       call phs_wood_test (u, results)
       call phs_wood_vis_test (u, results)
       call rng_base_test (u, results)
       call cascade_test (u, results)
       call rng_tao_test (u, results)
       call selectors_test (u, results)
       call mci_base_test (u, results)
       call mci_midpoint_test (u, results)
       call mci_vamp_test (u, results)
       call prclib_interfaces_test (u, results)
       call particle_specifiers_test (u, results)
       call process_libraries_test (u, results)
       call prclib_stacks_test (u, results)
       call slha_test (u, results)
       call state_matrix_test (u, results)
       call prc_test_test (u, results)
       call subevt_expr_test (u, results)
       call processes_test (u, results)
       call process_stacks_test (u, results)   
       call event_transforms_test (u, results)
       call decays_test (u, results)
       call shower_test (u, results)
       call events_test (u, results)
       call prc_omega_test (u, results)
       call prc_omega_diags_test (u, results)
       call prc_template_me_test (u, results)
       call hep_events_test (u, results)
       call eio_data_test (u, results)
       call eio_base_test (u, results)
       call eio_raw_test (u, results)
       call eio_checkpoints_test (u, results)
       call eio_lhef_test (u, results)
       call eio_hepmc_test (u, results)
       call eio_lcio_test (u, results)
       call eio_stdhep_test (u, results)
       call eio_ascii_test (u, results)
       call eio_weights_test (u, results)
       call iterations_test (u, results)
       call beam_structures_test (u, results)
       call rt_data_test (u, results)
       call dispatch_test (u, results)
       call process_configurations_test (u, results)
       call compilations_test (u, results)
       call compilations_static_test (u, results)
       call integrations_test (u, results)
       call integrations_history_test (u, results)       
       call event_streams_test (u, results)
       call simulations_test (u, results)
       call commands_test (u, results)
    case default
       call msg_fatal ("Self-test '" // char (check) // "' not implemented.")
    end select
    close (u)
  end subroutine whizard_check

\eatline
\nwendcode{}\nwbegindocs{3211}\nwdocspar
\clearpage
%%%%%%%%%%%%%%%%%%%%%%%%%%%%%%%%%%%%%%%%%%%%%%%%%%%%%%%%%%%%%%%%%%%%%%
\section{Driver program}
The main program handles command options, initializes the environment,
and runs WHIZARD in a particular mode (interactive, file, standard
input).

This is also used in the C interface:
\nwenddocs{}\nwbegincode{3212}\moddef{Main: cmdline arg len declaration}\endmoddef
  integer, parameter :: CMDLINE_ARG_LEN = 1000
\eatline
\nwendcode{}\nwbegindocs{3213}\nwdocspar
The actual main program:
\nwenddocs{}\nwbegincode{3214}\moddef{\code{}main.f90\edoc{}}\endmoddef
\LA{}File header\RA{}

program main

\LA{}Use strings\RA{}
  use unit_tests
  use system_dependencies
  use diagnostics
  use ifiles
  use os_interface
  use whizard

  implicit none

\LA{}Main: cmdline arg len declaration\RA{}

  ! Main program variable declarations
  character(CMDLINE_ARG_LEN) :: arg
  character(2) :: option
  type(string_t) :: long_option, value
  integer :: i, j, arg_len, arg_status
  logical :: look_for_options
  logical :: interactive
  logical :: banner
  type(string_t) :: files, this, model, default_lib, library, libraries
  type(string_t) :: check, checks, logfile
  type(test_results_t) :: test_results
  logical :: success
  logical :: user_code_enable = .false.
  integer :: n_user_src = 0, n_user_lib = 0
  type(string_t) :: user_src, user_lib, user_target
  type(paths_t) :: paths
  logical :: rebuild_library, rebuild_user
  logical :: rebuild_phs, rebuild_grids, rebuild_events
  logical :: recompile_library
  type(ifile_t) :: commands
  type(string_t) :: command

  type(whizard_options_t), allocatable :: options
  type(whizard_t), allocatable, target :: whizard_instance

  ! Exit status
  logical :: quit = .false.
  integer :: quit_code = 0

  ! Initial values
  look_for_options = .true.
  interactive = .false.
  files = ""
  model = "SM"
  default_lib = "default_lib"
  library = ""
  libraries = ""  
  banner = .true.
  logging = .true.
  logfile = "whizard.log"
  check = ""
  checks = ""
  user_src = ""
  user_lib = ""
  user_target = ""
  rebuild_library = .false.
  rebuild_user = .false.
  rebuild_phs = .false.
  rebuild_grids = .false.
  rebuild_events = .false.
  recompile_library = .false.
  call paths_init (paths)

  ! Read and process options
  i = 0
  SCAN_CMDLINE: do
     i = i + 1
     call get_command_argument (i, arg, arg_len, arg_status)
     select case (arg_status)
     case (0)
     case (-1)
        call msg_error (" Command argument truncated: '" // arg // "'")
     case default
        exit SCAN_CMDLINE
     end select
     if (look_for_options) then
        select case (arg(1:2))
        case ("--")
           value = trim (arg)
           call split (value, long_option, "=")
           select case (char (long_option))
           case ("--version")
              call no_option_value (long_option, value)
              call print_version (); stop
           case ("--help")
              call no_option_value (long_option, value)
              call print_usage (); stop
           case ("--prefix")
              paths%prefix = get_option_value (i, long_option, value)
              cycle SCAN_CMDLINE
           case ("--exec-prefix")
              paths%exec_prefix = get_option_value (i, long_option, value)
              cycle SCAN_CMDLINE
           case ("--bindir")
              paths%bindir = get_option_value (i, long_option, value)
              cycle SCAN_CMDLINE
           case ("--libdir")
              paths%libdir = get_option_value (i, long_option, value)
              cycle SCAN_CMDLINE
           case ("--includedir")
              paths%includedir = get_option_value (i, long_option, value)
              cycle SCAN_CMDLINE
           case ("--datarootdir")
              paths%datarootdir = get_option_value (i, long_option, value)
              cycle SCAN_CMDLINE
           case ("--libtool")
              paths%libtool = get_option_value (i, long_option, value)
              cycle SCAN_CMDLINE
           case ("--lhapdfdir")
              paths%lhapdfdir = get_option_value (i, long_option, value)
              cycle SCAN_CMDLINE
           case ("--check")
              check = get_option_value (i, long_option, value)
              checks = checks // " " // check
              cycle SCAN_CMDLINE
           case ("--execute")
              command = get_option_value (i, long_option, value)
              call ifile_append (commands, command)
              cycle SCAN_CMDLINE
           case ("--interactive")
              call no_option_value (long_option, value)
              interactive = .true.
              cycle SCAN_CMDLINE
           case ("--library")
              library = get_option_value (i, long_option, value)
              libraries = libraries // " " // library
              cycle SCAN_CMDLINE
           case ("--no-library")
              call no_option_value (long_option, value)
              default_lib = ""
              library = ""              
              libraries = ""
              cycle SCAN_CMDLINE
           case ("--localprefix")
              paths%localprefix = get_option_value (i, long_option, value)
              cycle SCAN_CMDLINE
           case ("--logfile")
              logfile = get_option_value (i, long_option, value)
              cycle SCAN_CMDLINE
           case ("--no-logfile")
              call no_option_value (long_option, value)
              logfile = ""
              cycle SCAN_CMDLINE
           case ("--logging")
              call no_option_value (long_option, value)
              logging = .true.
              cycle SCAN_CMDLINE
           case ("--no-logging")
              call no_option_value (long_option, value)
              logging = .false.
              cycle SCAN_CMDLINE
           case ("--banner")
              call no_option_value (long_option, value)
              banner = .true.
              cycle SCAN_CMDLINE
           case ("--no-banner")
              call no_option_value (long_option, value)
              banner = .false.
              cycle SCAN_CMDLINE
           case ("--model")
              model = get_option_value (i, long_option, value)
              cycle SCAN_CMDLINE
           case ("--no-model")
              call no_option_value (long_option, value)
              model = ""
              cycle SCAN_CMDLINE
           case ("--rebuild")
              call no_option_value (long_option, value)
              rebuild_library = .true.
              rebuild_user = .true.
              rebuild_phs = .true.
              rebuild_grids = .true.
              rebuild_events = .true.
              cycle SCAN_CMDLINE
           case ("--no-rebuild")
              call no_option_value (long_option, value)
              rebuild_library = .false.
              recompile_library = .false.
              rebuild_user = .false.
              rebuild_phs = .false.
              rebuild_grids = .false.
              rebuild_events = .false.
              cycle SCAN_CMDLINE
           case ("--rebuild-library")
              call no_option_value (long_option, value)
              rebuild_library = .true.
              cycle SCAN_CMDLINE
           case ("--rebuild-user")
              call no_option_value (long_option, value)
              rebuild_user = .true.
              cycle SCAN_CMDLINE
           case ("--rebuild-phase-space")
              call no_option_value (long_option, value)
              rebuild_phs = .true.
              cycle SCAN_CMDLINE
           case ("--rebuild-grids")
              call no_option_value (long_option, value)
              rebuild_grids = .true.
              cycle SCAN_CMDLINE
           case ("--rebuild-events")
              call no_option_value (long_option, value)
              rebuild_events = .true.
              cycle SCAN_CMDLINE
           case ("--recompile")
              call no_option_value (long_option, value)
              recompile_library = .true.
              rebuild_grids = .true.
              cycle SCAN_CMDLINE
           case ("--user")
              user_code_enable = .true.
              cycle SCAN_CMDLINE
           case ("--user-src")
              if (user_src == "") then
                 user_src = get_option_value (i, long_option, value)
              else
                 user_src = user_src // " " &
                      // get_option_value (i, long_option, value)
              end if
              n_user_src = n_user_src + 1
              cycle SCAN_CMDLINE
           case ("--user-lib")
              if (user_lib == "") then
                 user_lib = get_option_value (i, long_option, value)
              else
                 user_lib = user_lib // " " &
                      // get_option_value (i, long_option, value)
              end if
              n_user_lib = n_user_lib + 1
              cycle SCAN_CMDLINE
           case ("--user-target")
              user_target = get_option_value (i, long_option, value)
              cycle SCAN_CMDLINE
           case ("--write-syntax-tables")
              call no_option_value (long_option, value)
              call init_syntax_tables ()
              call write_syntax_tables ()
              call final_syntax_tables ()
              stop
              cycle SCAN_CMDLINE
           case default
              call print_usage ()
              call msg_fatal ("Option '" // trim (arg) // "' not recognized")
           end select
        end select
        select case (arg(1:1))
        case ("-")
           j = 1
           if (len_trim (arg) == 1) then
              look_for_options = .false.
           else
              SCAN_SHORT_OPTIONS: do
                 j = j + 1
                 if (j > len_trim (arg)) exit SCAN_SHORT_OPTIONS
                 option = "-" // arg(j:j)
                 select case (option)
                 case ("-V")
                    call print_version (); stop
                 case ("-?", "-h")
                    call print_usage (); stop
                 case ("-e")
                    command = get_option_value (i, var_str (option))
                    call ifile_append (commands, command)
                    cycle SCAN_CMDLINE
                 case ("-i")
                    interactive = .true.
                    cycle SCAN_SHORT_OPTIONS
                 case ("-l")
                    if (j == len_trim (arg)) then
                       library = get_option_value (i, var_str (option))
                    else
                       library = trim (arg(j+1:))
                    end if
                    libraries = libraries // " " // library                    
                    cycle SCAN_CMDLINE
                 case ("-L")
                    if (j == len_trim (arg)) then
                       logfile = get_option_value (i, var_str (option))
                    else
                       logfile = trim (arg(j+1:))
                    end if
                    cycle SCAN_CMDLINE
                 case ("-m")
                    if (j < len_trim (arg))  call msg_fatal &
                         ("Option '" // option // "' needs a value")
                    model = get_option_value (i, var_str (option))
                    cycle SCAN_CMDLINE
                 case ("-r")
                    rebuild_library = .true.
                    rebuild_user = .true.
                    rebuild_phs = .true.
                    rebuild_grids = .true.
                    rebuild_events = .true.
                    cycle SCAN_SHORT_OPTIONS
                 case ("-u")
                    user_code_enable = .true.
                    cycle SCAN_SHORT_OPTIONS
                 case default
                    call print_usage ()
                    call msg_fatal &
                         ("Option '" // option // "' not recognized")
                 end select
              end do SCAN_SHORT_OPTIONS
           end if
        case default
           files = files // " " // trim (arg)
        end select
     else
        files = files // " " // trim (arg)
     end if
  end do SCAN_CMDLINE

  ! Overall initialization
  if (logfile /= "")  call logfile_init (logfile)
  if (banner)  call msg_banner ()
  
   ! Run any self-checks (and no commands)
   if (checks /= "") then
      checks = trim (adjustl (checks))
      RUN_CHECKS: do while (checks /= "")
         call split (checks, check, " ")
         call whizard_check (check, test_results)
      end do RUN_CHECKS
      call test_results%wrapup (6, success)
      if (.not. success)  quit_code = 7
      quit = .true.
   end if
   
   allocate (options)
   allocate (whizard_instance)

   if (.not. quit) then

      ! Set options and initialize the whizard object
      options%preload_model = model
      options%default_lib = default_lib
      options%preload_libraries = libraries
      options%rebuild_library = rebuild_library
      options%recompile_library = recompile_library
      options%rebuild_user = rebuild_user 
      options%rebuild_phs = rebuild_phs
      options%rebuild_grids = rebuild_grids
      options%rebuild_events = rebuild_events
      
      call whizard_instance%init (options, paths, logfile)
      
      call mask_term_signals ()

   end if

   ! Run commands given on the command line
   if (.not. quit .and. ifile_get_length (commands) > 0) then
      call whizard_instance%process_ifile (commands, quit, quit_code)
   end if

   if (.not. quit) then
      ! Process commands from standard input
      if (.not. interactive .and. files == "") then
         call whizard_instance%process_stdin (quit, quit_code)

         ! ... or process commands from file
      else
         files = trim (adjustl (files))
         SCAN_FILES: do while (files /= "")
            call split (files, this, " ")
            call whizard_instance%process_file (this, quit, quit_code)
            if (quit)  exit SCAN_FILES
         end do SCAN_FILES

      end if
  end if
 
  ! Enter an interactive shell if requested
  if (.not. quit .and. interactive) then
     call whizard_instance%shell (quit_code)
  end if

  ! Overall finalization
  call ifile_final (commands)

  deallocate (options)

  call whizard_instance%final ()
  deallocate (whizard_instance)
  
  call terminate_now_if_signal ()
  call release_term_signals ()
  call msg_terminate (quit_code = quit_code)

!!!!!!!!!!!!!!!!!!!!!!!!!!!!!!!!!!!!!!!!!!!!!!!!!!!!!!!!!!!!!!!!!!!!!!!!
contains

  subroutine no_option_value (option, value)
    type(string_t), intent(in) :: option, value
    if (value /= "") then
       call msg_error (" Option '" // char (option) // "' should have no value")
    end if
  end subroutine no_option_value

  function get_option_value (i, option, value) result (string)
    type(string_t) :: string
    integer, intent(inout) :: i
    type(string_t), intent(in) :: option
    type(string_t), intent(in), optional :: value
    character(CMDLINE_ARG_LEN) :: arg
    character(CMDLINE_ARG_LEN) :: arg_value
    integer :: arg_len, arg_status
    logical :: has_value
    if (present (value)) then
       has_value = value /= ""
    else
       has_value = .false.
    end if
    if (has_value) then
       string = value
    else
       i = i + 1 
       call get_command_argument (i, arg_value, arg_len, arg_status)
       select case (arg_status)
       case (0)
       case (-1)
          call msg_error (" Option value truncated: '" // arg // "'")
       case default
          call print_usage ()
          call msg_fatal (" Option '" // char (option) // "' needs a value")
       end select
       select case (arg(1:1))
       case ("-")
          call print_usage ()
          call msg_fatal (" Option '" // char (option) // "' needs a value")
       end select
       string = trim (arg_value)
    end if
  end function get_option_value

  subroutine print_version ()
    print "(A)", "WHIZARD " // WHIZARD_VERSION 
    print "(A)", "Copyright (C) 1999-2015 Wolfgang Kilian, Thorsten Ohl, Juergen Reuter"
    print "(A)", "              ---------------------------------------                "
    print "(A)", "This is free software; see the source for copying conditions.  There is NO"
    print "(A)", "warranty; not even for MERCHANTABILITY or FITNESS FOR A PARTICULAR PURPOSE."
    print *
  end subroutine print_version

  subroutine print_usage ()
    print "(A)", "WHIZARD " // WHIZARD_VERSION
    print "(A)", "Usage: whizard [OPTIONS] [FILE]"
    print "(A)", "Run WHIZARD with the command list taken from FILE(s)"
    print "(A)", "Options for resetting default directories and tools" &
            // "(GNU naming conventions):"
    print "(A)", "    --prefix DIR"
    print "(A)", "    --exec_prefix DIR"
    print "(A)", "    --bindir DIR"
    print "(A)", "    --libdir DIR"
    print "(A)", "    --includedir DIR"
    print "(A)", "    --datarootdir DIR"
    print "(A)", "    --libtool LOCAL_LIBTOOL"
    print "(A)", "    --lhapdfdir DIR   (PDF sets directory)"
    print "(A)", "Other options:"
    print "(A)", "-h, --help            display this help and exit"
    print "(A)", "    --banner          display banner at startup (default)"
    print "(A)", "-e, --execute CMDS    execute SINDARIN CMDS before reading FILE(s)"
    print "(A)", "-i, --interactive     run interactively after reading FILE(s)"
    print "(A)", "-l, --library         preload process library NAME"
    print "(A)", "    --localprefix DIR"
    print "(A)", "                      search in DIR for local models (default: ~/.whizard)"
    print "(A)", "-L, --logfile FILE    write log to FILE (default: 'whizard.log'"
    print "(A)", "    --logging         switch on logging at startup (default)"
    print "(A)", "-m, --model NAME      preload model NAME (default: 'SM')"
    print "(A)", "    --no-banner       do not display banner at startup"
    print "(A)", "    --no-library      do not preload process library"
    print "(A)", "    --no-logfile      do not write a logfile"
    print "(A)", "    --no-logging      switch off logging at startup"
    print "(A)", "    --no-model        do not preload a model"
    print "(A)", "    --no-rebuild      do not force rebuilding"
    print "(A)", "-r, --rebuild         rebuild all (see below)"
    print "(A)", "    --rebuild-library"
    print "(A)", "                      rebuild process code library"
    print "(A)", "    --rebuild-user    rebuild user-provided code"
    print "(A)", "    --rebuild-phase-space"
    print "(A)", "                      rebuild phase-space configuration"
    print "(A)", "    --rebuild-grids   rebuild integration grids"
    print "(A)", "    --rebuild-events  rebuild event samples"
    print "(A)", "    --recompile       recompile process code"
    print "(A)", "-u  --user            enable user-provided code"
    print "(A)", "    --user-src FILE   user-provided source file"
    print "(A)", "    --user-lib FILE   user-provided library file"
    print "(A)", "    --user-target BN  basename of created user library (default: user)"
    print "(A)", "-V, --version         output version information and exit"
    print "(A)", "    --write-syntax-tables"
    print "(A)", "                      write the internal syntax tables to files and exit"
    print "(A)", "-                     further options are taken as filenames"
    print *
    print "(A)", "With no FILE, read standard input."
  end subroutine print_usage

end program main
\eatline
\nwendcode{}\nwbegindocs{3215}\nwdocspar
\clearpage
%%%%%%%%%%%%%%%%%%%%%%%%%%%%%%%%%%%%%%%%%%%%%%%%%%%%%%%%%%%%%%%%%%%%%% 
\section{Whizard-C-Interface}
\nwenddocs{}\nwbegincode{3216}\moddef{\code{}whizard-c-interface.f90\edoc{}}\endmoddef
\LA{}File header\RA{}

\LA{}Whizard-C-Interface: Internals\RA{}
\LA{}Whizard-C-Interface: Init and Finalize\RA{}
\LA{}Whizard-C-Interface: Interfaced Commads\RA{}
\LA{}Whizard-C-Interface: HepMC\RA{}

\nwendcode{}\nwbegindocs{3217}\nwdocspar
\nwenddocs{}\nwbegincode{3218}\moddef{Whizard-C-Interface: Internals}\endmoddef
  subroutine c_whizard_convert_string (c_string, f_string)
    use, intrinsic :: iso_c_binding
    use iso_varying_string, string_t => varying_string !NODEP!
  
    implicit none
  
    character(kind=c_char), intent(in) :: c_string(*)
    type(string_t), intent(inout) :: f_string
    character(len=1) :: dummy_char
    integer :: dummy_i = 1
  
    f_string = ""
    do
       if (c_string(dummy_i) == c_null_char) then
          exit
       else if (c_string(dummy_i) == c_new_line) then
          dummy_char = CHAR(13)
          f_string = f_string // dummy_char
          dummy_char = CHAR(10)
       else
          dummy_char = c_string (dummy_i)
       end if
       f_string = f_string // dummy_char
       dummy_i = dummy_i + 1
    end do
    dummy_i = 1
  end subroutine c_whizard_convert_string
  
  subroutine c_whizard_commands (w_c_instance, cmds)
    use, intrinsic :: iso_c_binding
    use iso_varying_string, string_t => varying_string !NODEP!
    use commands
    use diagnostics
    use lexers
    use models
    use parser
    use whizard

    type(c_ptr), intent(inout) :: w_c_instance
    type(whizard_t), pointer :: whizard_instance
    type(string_t) :: cmds
    type(parse_tree_t) :: parse_tree
    type(parse_node_t), pointer :: pn_root
    type(stream_t), target :: stream
    type(lexer_t) :: lexer
    type(command_list_t), target :: cmd_list

    call c_f_pointer (w_c_instance, whizard_instance)
    call lexer_init_cmd_list (lexer)
    call syntax_cmd_list_init ()
  
    call stream_init (stream, cmds)
    call lexer_assign_stream (lexer, stream)
    call parse_tree_init (parse_tree, syntax_cmd_list, lexer)
    pn_root => parse_tree_get_root_ptr (parse_tree)
  
    if (associated (pn_root)) then
       call cmd_list%compile (pn_root, whizard_instance%global)
    end if
    call whizard_instance%global%activate ()
    call cmd_list%execute (whizard_instance%global)
    call cmd_list%final ()
  
    call parse_tree_final (parse_tree)
    call stream_final (stream)
    call lexer_final (lexer)
    call syntax_cmd_list_final ()
  end subroutine c_whizard_commands

\nwendcode{}\nwbegindocs{3219}\nwdocspar
\nwenddocs{}\nwbegincode{3220}\moddef{Whizard-C-Interface: Init and Finalize}\endmoddef
  subroutine c_whizard_init (w_c_instance) bind(C)
    use, intrinsic :: iso_c_binding
    use iso_varying_string, string_t => varying_string !NODEP!    
    use system_dependencies
    use diagnostics
    use unit_tests
    use ifiles
    use os_interface
    use whizard
  
    implicit none
  
  \LA{}Main: cmdline arg len declaration\RA{}

    type(c_ptr), intent(out) :: w_c_instance    
    character(CMDLINE_ARG_LEN) :: arg
    character(2) :: option
    type(string_t) :: long_option, value
    integer :: i, j, arg_len, arg_status
    logical :: look_for_options
    logical :: interactive
    logical :: banner
    type(string_t) :: files, this, model, default_lib, library, libraries
    type(string_t) :: check, checks, logfile
    type(test_results_t) :: test_results
    logical :: success
    logical :: user_code_enable = .false.
    integer :: n_user_src = 0, n_user_lib = 0
    type(string_t) :: user_src, user_lib
    type(paths_t) :: paths
    logical :: rebuild_library, rebuild_user
    logical :: rebuild_phs, rebuild_grids, rebuild_events
    logical :: recompile_library
    type(ifile_t) :: commands
    type(string_t) :: command

    type(whizard_options_t), allocatable :: options
    type(whizard_t), pointer :: whizard_instance
    
    ! Exit status
    logical :: quit = .false.
    integer :: quit_code = 0
  
    ! Initial values
    look_for_options = .true.
    interactive = .false.
    files = ""
    model = "SM"
    default_lib = "default_lib"
    library = ""
    libraries = ""    
    banner = .true.
    logging = .true.
    logfile = "whizard.log"
    check = ""
    checks = ""
    user_src = ""
    user_lib = ""
    rebuild_library = .false.
    recompile_library = .false.
    rebuild_user = .false.
    rebuild_phs = .false.
    rebuild_grids = .false.
    rebuild_events = .false.
    call paths_init (paths)
  
    ! Overall initialization
    if (logfile /= "")  call logfile_init (logfile)
    call mask_term_signals ()
    if (banner)  call msg_banner ()
    
    ! Set options and initialize the whizard object
    allocate (options)
    options%preload_model = model
    options%default_lib = default_lib
    options%preload_libraries = libraries
    options%rebuild_library = rebuild_library
    options%rebuild_user = rebuild_user
    options%rebuild_phs = rebuild_phs
    options%rebuild_grids = rebuild_grids
    options%rebuild_events = rebuild_events
    
    allocate (whizard_instance)
    call whizard_instance%init (options, paths)
      
    if (checks /= "") then
       checks = trim (adjustl (checks))
       RUN_CHECKS: do while (checks /= "")
          call split (checks, check, " ")
          call whizard_check (check, test_results)
       end do RUN_CHECKS
       call test_results%wrapup (6, success)
       if (.not. success)  quit_code = 7
       quit = .true.
    end if
    
    w_c_instance = c_loc (whizard_instance)
              
  end subroutine c_whizard_init
  
  subroutine c_whizard_finalize (w_c_instance) bind(C)
    use, intrinsic :: iso_c_binding
    use iso_varying_string, string_t => varying_string !NODEP!
    use system_dependencies
    use diagnostics
    use ifiles
    use os_interface
    use whizard
    
    type(c_ptr), intent(in) :: w_c_instance
    type(whizard_t), pointer :: whizard_instance
    logical :: quit = .false.
    integer :: quit_code = 0
      
    call c_f_pointer (w_c_instance, whizard_instance)
    call whizard_instance%final ()
    deallocate (whizard_instance)
    call terminate_now_if_signal ()
    call release_term_signals ()
    call msg_terminate (quit_code = quit_code)
  end subroutine c_whizard_finalize
  
  subroutine c_whizard_process_string (w_c_instance, c_cmds_in) bind(C)
    use, intrinsic :: iso_c_binding
    use iso_varying_string, string_t => varying_string !NODEP!
  
    implicit none
  
    type(c_ptr), intent(inout) :: w_c_instance
    character(kind=c_char) :: c_cmds_in(*)
    type(string_t) :: f_cmds
  
    call c_whizard_convert_string (c_cmds_in, f_cmds)
    call c_whizard_commands (w_c_instance, f_cmds)
  end subroutine c_whizard_process_string

\nwendcode{}\nwbegindocs{3221}\nwdocspar
\nwenddocs{}\nwbegincode{3222}\moddef{Whizard-C-Interface: Interfaced Commads}\endmoddef
  subroutine c_whizard_model (w_c_instance, c_model) bind(C)
    use, intrinsic :: iso_c_binding
    use iso_varying_string, string_t => varying_string !NODEP!
  
    implicit none
  
    type(c_ptr), intent(inout) :: w_c_instance
    character(kind=c_char) :: c_model(*)
    type(string_t) :: model, mdl_str
  
    call c_whizard_convert_string (c_model, model)
    mdl_str = "model = " // model
    call c_whizard_commands (w_c_instance, mdl_str)
  end subroutine c_whizard_model
  
  subroutine c_whizard_library (w_c_instance, c_library) bind(C)
    use, intrinsic :: iso_c_binding
    use iso_varying_string, string_t => varying_string !NODEP!
  
    implicit none
  
    type(c_ptr), intent(inout) :: w_c_instance
    character(kind=c_char) :: c_library(*)
    type(string_t) :: library, lib_str
  
    call c_whizard_convert_string(c_library, library)
    lib_str = "library = " // library
    call c_whizard_commands (w_c_instance, lib_str)
  end subroutine c_whizard_library
  
  subroutine c_whizard_process (w_c_instance, c_id, c_in, c_out) bind(C)
    use, intrinsic :: iso_c_binding
    use iso_varying_string, string_t => varying_string !NODEP!
  
    implicit none
  
    type(c_ptr), intent(inout) :: w_c_instance
    character(kind=c_char) :: c_id(*), c_in(*), c_out(*)
    type(string_t) :: proc_str, id, in, out
  
    call c_whizard_convert_string (c_id, id)
    call c_whizard_convert_string (c_in, in)
    call c_whizard_convert_string (c_out, out)
    proc_str = "process " // id // " = " // in // " => " // out
    call c_whizard_commands (w_c_instance, proc_str)
  end subroutine c_whizard_process
  
  subroutine c_whizard_compile (w_c_instance) bind(C)
    use, intrinsic :: iso_c_binding
    use iso_varying_string, string_t => varying_string !NODEP!
  
    type(c_ptr), intent(inout) :: w_c_instance
    type(string_t) :: cmp_str
    cmp_str = "compile"
    call c_whizard_commands (w_c_instance, cmp_str)
  end subroutine c_whizard_compile
    
  subroutine c_whizard_beams (w_c_instance, c_specs) bind(C)
    use, intrinsic :: iso_c_binding
    use iso_varying_string, string_t => varying_string !NODEP!
  
    implicit none
  
    type(c_ptr), intent(inout) :: w_c_instance
    character(kind=c_char) :: c_specs(*)
    type(string_t) :: specs, beam_str
  
    call c_whizard_convert_string (c_specs, specs)
    beam_str = "beams = " // specs
    call c_whizard_commands (w_c_instance, beam_str)
  end subroutine c_whizard_beams
  
  subroutine c_whizard_integrate (w_c_instance, c_process) bind(C)
    use, intrinsic :: iso_c_binding
    use iso_varying_string, string_t => varying_string !NODEP!
  
    implicit none
  
    type(c_ptr), intent(inout) :: w_c_instance
    character(kind=c_char) :: c_process(*)
    type(string_t) :: process, int_str
  
    call c_whizard_convert_string (c_process, process)
    int_str = "integrate (" // process //")"
    call c_whizard_commands (w_c_instance, int_str)
  end subroutine c_whizard_integrate
  
  subroutine c_whizard_matrix_element_test &
       (w_c_instance, c_process, n_calls) bind(C)
    use, intrinsic :: iso_c_binding
    use iso_varying_string, string_t => varying_string !NODEP!
  
    implicit none

    type(c_ptr), intent(inout) :: w_c_instance
    integer(kind=c_int) :: n_calls
    character(kind=c_char) :: c_process(*)
    type(string_t) :: process, me_str
    character(len=8) :: buffer
  
    call c_whizard_convert_string (c_process, process)
    write (buffer, "(I0)")  n_calls
    me_str = "integrate (" // process // ") \{ ?phs_only = true" // &
         "  n_calls_test = " // trim (buffer) 
    call c_whizard_commands (w_c_instance, me_str)
  end subroutine c_whizard_matrix_element_test
  
  subroutine c_whizard_simulate (w_c_instance, c_id) bind(C)
    use, intrinsic :: iso_c_binding
    use iso_varying_string, string_t => varying_string !NODEP!
  
    implicit none
  
    type(c_ptr), intent(inout) :: w_c_instance
    character(kind=c_char) :: c_id(*)
    type(string_t) :: sim_str, id
  
    call c_whizard_convert_string(c_id, id)
    sim_str = "simulate (" // id // ")"
    call c_whizard_commands (w_c_instance, sim_str)
  end subroutine c_whizard_simulate
  
  subroutine c_whizard_sqrts (w_c_instance, c_value, c_unit) bind(C)
    use, intrinsic :: iso_c_binding
    use iso_varying_string, string_t => varying_string !NODEP!
  
    implicit none

    type(c_ptr), intent(inout) :: w_c_instance
    character(kind=c_char) :: c_unit(*)
    integer(kind=c_int) :: c_value
    integer :: f_value
    character(len=8) :: f_val
    type(string_t) :: val, unit, sqrts_str
  
    f_value = c_value
    write (f_val,'(i8)') f_value
    val = f_val
    call c_whizard_convert_string (c_unit, unit)
    sqrts_str = "sqrts =" // val // unit
    call c_whizard_commands (w_c_instance, sqrts_str)
  end subroutine c_whizard_sqrts
  
\nwendcode{}\nwbegindocs{3223}\nwdocspar
\nwenddocs{}\nwbegincode{3224}\moddef{Whizard-C-Interface: HepMC}\endmoddef
  type(c_ptr) function c_whizard_hepmc_test &
       (w_c_instance, c_id, c_proc_id, c_event_id) bind(C)
    use, intrinsic :: iso_c_binding
    use iso_varying_string, string_t => varying_string !NODEP!
    use commands
    use diagnostics
    use events
    use hepmc_interface
    use lexers
    use models
    use parser
    use processes
    use rt_data
    use simulations
    use whizard
    use os_interface
   
    implicit none

    type(c_ptr), intent(inout) :: w_c_instance    
    type(string_t) :: sim_str
    type(parse_tree_t) :: parse_tree
    type(parse_node_t), pointer :: pn_root
    type(stream_t), target :: stream
    type(lexer_t) :: lexer
    type(command_list_t), pointer :: cmd_list
    type(whizard_t), pointer :: whizard_instance
   
    integer :: i_evt
    type(simulation_t), target :: sim
   
    character(kind=c_char), intent(in) :: c_id(*)
    type(string_t) :: id
    integer(kind=c_int), value :: c_proc_id, c_event_id
    integer :: proc_id, event_id
  
    type(hepmc_event_t), pointer :: hepmc_event
   
    type(process_t), pointer :: process
    integer :: proc
   
    integer :: factorization_mode, try
   
    call c_f_pointer (w_c_instance, whizard_instance)        
    
    call c_whizard_convert_string (c_id, id)
    sim_str = "simulate (" // id // ")" 
  
    proc_id = c_proc_id
    event_id = c_event_id
  
    allocate (hepmc_event)
    call hepmc_event_init (hepmc_event, c_proc_id, c_event_id)
  
    call syntax_cmd_list_init ()
    call lexer_init_cmd_list (lexer)
   
    call stream_init (stream, sim_str)
    call lexer_assign_stream (lexer, stream)
    call parse_tree_init (parse_tree, syntax_cmd_list, lexer)
    pn_root => parse_tree_get_root_ptr (parse_tree)
   
    allocate (cmd_list)
    if (associated (pn_root)) then
       call cmd_list%compile (pn_root, whizard_instance%global)
    end if   
  
    call sim%init ([id], .true., .true., whizard_instance%global)

    !!! This should generate a HepMC event as hepmc_event_t type
    call msg_message ("Not enabled for the moment.")
  
    call sim%final ()

    call cmd_list%final ()
   
    call parse_tree_final (parse_tree)
    call stream_final (stream)
    call lexer_final (lexer)
    call syntax_cmd_list_final ()
  
    c_whizard_hepmc_test = c_loc(hepmc_event)
    return
  end function c_whizard_hepmc_test
\nwendcode{}\nwbegindocs{3225}%%%%%%%%%%%%%%%%%%%%%%%%%%%%%%%%%%%%%%%%%%%%%%%%%%%%%%%%%%%%%%%%%%%%%% 
\nwenddocs{}\nwfilename{../../../src/noweb-frame/whizard-postlude.nw}\nwbegindocs{0}\chapter{Cross References}
\section{Identifiers}
\nowebindex
\section{Chunks}
\nowebchunks
\InputIfFileExists{\jobname.ind}{}{}
\end{document}
\nwenddocs{}
